\section{Bench Huge Section 0}
% synthetic bench section 0
% This is a sample LaTeX input file.  (Version of 12 August 2004.)
%
% A '%' character causes TeX to ignore all remaining text on the line,
% and is used for comments like this one.

\documentclass{article}      % Specifies the document class

                             % The preamble begins here.
\title{An Example Document}  % Declares the document's title.
\author{Leslie Lamport}      % Declares the author's name.
\date{January 21, 1994}      % Deleting this command produces today's date.

\newcommand{\ip}[2]{(#1, #2)}
                             % Defines \ip{arg1}{arg2} to mean
                             % (arg1, arg2).

%\newcommand{\ip}[2]{\langle #1 | #2\rangle}
                             % This is an alternative definition of
                             % \ip that is commented out.

\begin{document}             % End of preamble and beginning of text.

\maketitle                   % Produces the title.

This is an example input file.  Comparing it with
the output it generates can show you how to
produce a simple document of your own.

\section{Ordinary Text}      % Produces section heading.  Lower-level
                             % sections are begun with similar
                             % \subsection and \subsubsection commands.

The ends  of words and sentences are marked
  by   spaces. It  doesn't matter how many
spaces    you type; one is as good as 100.  The
end of   a line counts as a space.

One   or more   blank lines denote the  end
of  a paragraph.

Since any number of consecutive spaces are treated
like a single one, the formatting of the input
file makes no difference to
      \LaTeX,                % The \LaTeX command generates the LaTeX logo.
but it makes a difference to you.  When you use
\LaTeX, making your input file as easy to read
as possible will be a great help as you write
your document and when you change it.  This sample
file shows how you can add comments to your own input
file.

Because printing is different from typewriting,
there are a number of things that you have to do
differently when preparing an input file than if
you were just typing the document directly.
Quotation marks like
       ``this''
have to be handled specially, as do quotes within
quotes:
       ``\,`this'            % \, separates the double and single quote.
        is what I just
        wrote, not  `that'\,''.

Dashes come in three sizes: an
       intra-word
dash, a medium dash for number ranges like
       1--2,
and a punctuation
       dash---like
this.

A sentence-ending space should be larger than the
space between words within a sentence.  You
sometimes have to type special commands in
conjunction with punctuation characters to get
this right, as in the following sentence.
       Gnats, gnus, etc.\ all  % `\ ' makes an inter-word space.
       begin with G\@.         % \@ marks end-of-sentence punctuation.
You should check the spaces after periods when
reading your output to make sure you haven't
forgotten any special cases.  Generating an
ellipsis
       \ldots\               % `\ ' is needed after `\ldots' because TeX
                             % ignores spaces after command names like \ldots
                             % made from \ + letters.
                             %
                             % Note how a `%' character causes TeX to ignore
                             % the end of the input line, so these blank lines
                             % do not start a new paragraph.
                             %
with the right spacing around the periods requires
a special command.

\LaTeX\ interprets some common characters as
commands, so you must type special commands to
generate them.  These characters include the
following:
       \$ \& \% \# \{ and \}.

In printing, text is usually emphasized with an
       \emph{italic}
type style.

\begin{em}
   A long segment of text can also be emphasized
   in this way.  Text within such a segment can be
   given \emph{additional} emphasis.
\end{em}

It is sometimes necessary to prevent \LaTeX\ from
breaking a line where it might otherwise do so.
This may be at a space, as between the ``Mr.''\ and
``Jones'' in
       ``Mr.~Jones'',        % ~ produces an unbreakable interword space.
or within a word---especially when the word is a
symbol like
       \mbox{\emph{itemnum}}
that makes little sense when hyphenated across
lines.

Footnotes\footnote{This is an example of a footnote.}
pose no problem.

\LaTeX\ is good at typesetting mathematical formulas
like
       \( x-3y + z = 7 \)
or
       \( a_{1} > x^{2n} + y^{2n} > x' \)
or
       \( \ip{A}{B} = \sum_{i} a_{i} b_{i} \).
The spaces you type in a formula are
ignored.  Remember that a letter like
       $x$                   % $ ... $  and  \( ... \)  are equivalent
is a formula when it denotes a mathematical
symbol, and it should be typed as one.

\section{Displayed Text}

Text is displayed by indenting it from the left
margin.  Quotations are commonly displayed.  There
are short quotations
\begin{quote}
   This is a short quotation.  It consists of a
   single paragraph of text.  See how it is formatted.
\end{quote}
and longer ones.
\begin{quotation}
   This is a longer quotation.  It consists of two
   paragraphs of text, neither of which are
   particularly interesting.

   This is the second paragraph of the quotation.  It
   is just as dull as the first paragraph.
\end{quotation}
Another frequently-displayed structure is a list.
The following is an example of an \emph{itemized}
list.
\begin{itemize}
   \item This is the first item of an itemized list.
         Each item in the list is marked with a ``tick''.
         You don't have to worry about what kind of tick
         mark is used.

   \item This is the second item of the list.  It
         contains another list nested inside it.  The inner
         list is an \emph{enumerated} list.
         \begin{enumerate}
            \item This is the first item of an enumerated
                  list that is nested within the itemized list.

            \item This is the second item of the inner list.
                  \LaTeX\ allows you to nest lists deeper than
                  you really should.
         \end{enumerate}
         This is the rest of the second item of the outer
         list.  It is no more interesting than any other
         part of the item.
   \item This is the third item of the list.
\end{itemize}
You can even display poetry.
\begin{verse}
   There is an environment
    for verse \\             % The \\ command separates lines
   Whose features some poets % within a stanza.
   will curse.

                             % One or more blank lines separate stanzas.

   For instead of making\\
   Them do \emph{all} line breaking, \\
   It allows them to put too many words on a line when they'd rather be
   forced to be terse.
\end{verse}

Mathematical formulas may also be displayed.  A
displayed formula
is
one-line long; multiline
formulas require special formatting instructions.
   \[  \ip{\Gamma}{\psi'} = x'' + y^{2} + z_{i}^{n}\]
Don't start a paragraph with a displayed equation,
nor make one a paragraph by itself.

\end{document}               % End of document.

\section{Bench Huge Section 1}
% synthetic bench section 1
% This is a sample LaTeX input file.  (Version of 12 August 2004.)
%
% A '%' character causes TeX to ignore all remaining text on the line,
% and is used for comments like this one.

\documentclass{article}      % Specifies the document class

                             % The preamble begins here.
\title{An Example Document}  % Declares the document's title.
\author{Leslie Lamport}      % Declares the author's name.
\date{January 21, 1994}      % Deleting this command produces today's date.

\newcommand{\ip}[2]{(#1, #2)}
                             % Defines \ip{arg1}{arg2} to mean
                             % (arg1, arg2).

%\newcommand{\ip}[2]{\langle #1 | #2\rangle}
                             % This is an alternative definition of
                             % \ip that is commented out.

\begin{document}             % End of preamble and beginning of text.

\maketitle                   % Produces the title.

This is an example input file.  Comparing it with
the output it generates can show you how to
produce a simple document of your own.

\section{Ordinary Text}      % Produces section heading.  Lower-level
                             % sections are begun with similar
                             % \subsection and \subsubsection commands.

The ends  of words and sentences are marked
  by   spaces. It  doesn't matter how many
spaces    you type; one is as good as 100.  The
end of   a line counts as a space.

One   or more   blank lines denote the  end
of  a paragraph.

Since any number of consecutive spaces are treated
like a single one, the formatting of the input
file makes no difference to
      \LaTeX,                % The \LaTeX command generates the LaTeX logo.
but it makes a difference to you.  When you use
\LaTeX, making your input file as easy to read
as possible will be a great help as you write
your document and when you change it.  This sample
file shows how you can add comments to your own input
file.

Because printing is different from typewriting,
there are a number of things that you have to do
differently when preparing an input file than if
you were just typing the document directly.
Quotation marks like
       ``this''
have to be handled specially, as do quotes within
quotes:
       ``\,`this'            % \, separates the double and single quote.
        is what I just
        wrote, not  `that'\,''.

Dashes come in three sizes: an
       intra-word
dash, a medium dash for number ranges like
       1--2,
and a punctuation
       dash---like
this.

A sentence-ending space should be larger than the
space between words within a sentence.  You
sometimes have to type special commands in
conjunction with punctuation characters to get
this right, as in the following sentence.
       Gnats, gnus, etc.\ all  % `\ ' makes an inter-word space.
       begin with G\@.         % \@ marks end-of-sentence punctuation.
You should check the spaces after periods when
reading your output to make sure you haven't
forgotten any special cases.  Generating an
ellipsis
       \ldots\               % `\ ' is needed after `\ldots' because TeX
                             % ignores spaces after command names like \ldots
                             % made from \ + letters.
                             %
                             % Note how a `%' character causes TeX to ignore
                             % the end of the input line, so these blank lines
                             % do not start a new paragraph.
                             %
with the right spacing around the periods requires
a special command.

\LaTeX\ interprets some common characters as
commands, so you must type special commands to
generate them.  These characters include the
following:
       \$ \& \% \# \{ and \}.

In printing, text is usually emphasized with an
       \emph{italic}
type style.

\begin{em}
   A long segment of text can also be emphasized
   in this way.  Text within such a segment can be
   given \emph{additional} emphasis.
\end{em}

It is sometimes necessary to prevent \LaTeX\ from
breaking a line where it might otherwise do so.
This may be at a space, as between the ``Mr.''\ and
``Jones'' in
       ``Mr.~Jones'',        % ~ produces an unbreakable interword space.
or within a word---especially when the word is a
symbol like
       \mbox{\emph{itemnum}}
that makes little sense when hyphenated across
lines.

Footnotes\footnote{This is an example of a footnote.}
pose no problem.

\LaTeX\ is good at typesetting mathematical formulas
like
       \( x-3y + z = 7 \)
or
       \( a_{1} > x^{2n} + y^{2n} > x' \)
or
       \( \ip{A}{B} = \sum_{i} a_{i} b_{i} \).
The spaces you type in a formula are
ignored.  Remember that a letter like
       $x$                   % $ ... $  and  \( ... \)  are equivalent
is a formula when it denotes a mathematical
symbol, and it should be typed as one.

\section{Displayed Text}

Text is displayed by indenting it from the left
margin.  Quotations are commonly displayed.  There
are short quotations
\begin{quote}
   This is a short quotation.  It consists of a
   single paragraph of text.  See how it is formatted.
\end{quote}
and longer ones.
\begin{quotation}
   This is a longer quotation.  It consists of two
   paragraphs of text, neither of which are
   particularly interesting.

   This is the second paragraph of the quotation.  It
   is just as dull as the first paragraph.
\end{quotation}
Another frequently-displayed structure is a list.
The following is an example of an \emph{itemized}
list.
\begin{itemize}
   \item This is the first item of an itemized list.
         Each item in the list is marked with a ``tick''.
         You don't have to worry about what kind of tick
         mark is used.

   \item This is the second item of the list.  It
         contains another list nested inside it.  The inner
         list is an \emph{enumerated} list.
         \begin{enumerate}
            \item This is the first item of an enumerated
                  list that is nested within the itemized list.

            \item This is the second item of the inner list.
                  \LaTeX\ allows you to nest lists deeper than
                  you really should.
         \end{enumerate}
         This is the rest of the second item of the outer
         list.  It is no more interesting than any other
         part of the item.
   \item This is the third item of the list.
\end{itemize}
You can even display poetry.
\begin{verse}
   There is an environment
    for verse \\             % The \\ command separates lines
   Whose features some poets % within a stanza.
   will curse.

                             % One or more blank lines separate stanzas.

   For instead of making\\
   Them do \emph{all} line breaking, \\
   It allows them to put too many words on a line when they'd rather be
   forced to be terse.
\end{verse}

Mathematical formulas may also be displayed.  A
displayed formula
is
one-line long; multiline
formulas require special formatting instructions.
   \[  \ip{\Gamma}{\psi'} = x'' + y^{2} + z_{i}^{n}\]
Don't start a paragraph with a displayed equation,
nor make one a paragraph by itself.

\end{document}               % End of document.

\section{Bench Huge Section 2}
% synthetic bench section 2
% This is a sample LaTeX input file.  (Version of 12 August 2004.)
%
% A '%' character causes TeX to ignore all remaining text on the line,
% and is used for comments like this one.

\documentclass{article}      % Specifies the document class

                             % The preamble begins here.
\title{An Example Document}  % Declares the document's title.
\author{Leslie Lamport}      % Declares the author's name.
\date{January 21, 1994}      % Deleting this command produces today's date.

\newcommand{\ip}[2]{(#1, #2)}
                             % Defines \ip{arg1}{arg2} to mean
                             % (arg1, arg2).

%\newcommand{\ip}[2]{\langle #1 | #2\rangle}
                             % This is an alternative definition of
                             % \ip that is commented out.

\begin{document}             % End of preamble and beginning of text.

\maketitle                   % Produces the title.

This is an example input file.  Comparing it with
the output it generates can show you how to
produce a simple document of your own.

\section{Ordinary Text}      % Produces section heading.  Lower-level
                             % sections are begun with similar
                             % \subsection and \subsubsection commands.

The ends  of words and sentences are marked
  by   spaces. It  doesn't matter how many
spaces    you type; one is as good as 100.  The
end of   a line counts as a space.

One   or more   blank lines denote the  end
of  a paragraph.

Since any number of consecutive spaces are treated
like a single one, the formatting of the input
file makes no difference to
      \LaTeX,                % The \LaTeX command generates the LaTeX logo.
but it makes a difference to you.  When you use
\LaTeX, making your input file as easy to read
as possible will be a great help as you write
your document and when you change it.  This sample
file shows how you can add comments to your own input
file.

Because printing is different from typewriting,
there are a number of things that you have to do
differently when preparing an input file than if
you were just typing the document directly.
Quotation marks like
       ``this''
have to be handled specially, as do quotes within
quotes:
       ``\,`this'            % \, separates the double and single quote.
        is what I just
        wrote, not  `that'\,''.

Dashes come in three sizes: an
       intra-word
dash, a medium dash for number ranges like
       1--2,
and a punctuation
       dash---like
this.

A sentence-ending space should be larger than the
space between words within a sentence.  You
sometimes have to type special commands in
conjunction with punctuation characters to get
this right, as in the following sentence.
       Gnats, gnus, etc.\ all  % `\ ' makes an inter-word space.
       begin with G\@.         % \@ marks end-of-sentence punctuation.
You should check the spaces after periods when
reading your output to make sure you haven't
forgotten any special cases.  Generating an
ellipsis
       \ldots\               % `\ ' is needed after `\ldots' because TeX
                             % ignores spaces after command names like \ldots
                             % made from \ + letters.
                             %
                             % Note how a `%' character causes TeX to ignore
                             % the end of the input line, so these blank lines
                             % do not start a new paragraph.
                             %
with the right spacing around the periods requires
a special command.

\LaTeX\ interprets some common characters as
commands, so you must type special commands to
generate them.  These characters include the
following:
       \$ \& \% \# \{ and \}.

In printing, text is usually emphasized with an
       \emph{italic}
type style.

\begin{em}
   A long segment of text can also be emphasized
   in this way.  Text within such a segment can be
   given \emph{additional} emphasis.
\end{em}

It is sometimes necessary to prevent \LaTeX\ from
breaking a line where it might otherwise do so.
This may be at a space, as between the ``Mr.''\ and
``Jones'' in
       ``Mr.~Jones'',        % ~ produces an unbreakable interword space.
or within a word---especially when the word is a
symbol like
       \mbox{\emph{itemnum}}
that makes little sense when hyphenated across
lines.

Footnotes\footnote{This is an example of a footnote.}
pose no problem.

\LaTeX\ is good at typesetting mathematical formulas
like
       \( x-3y + z = 7 \)
or
       \( a_{1} > x^{2n} + y^{2n} > x' \)
or
       \( \ip{A}{B} = \sum_{i} a_{i} b_{i} \).
The spaces you type in a formula are
ignored.  Remember that a letter like
       $x$                   % $ ... $  and  \( ... \)  are equivalent
is a formula when it denotes a mathematical
symbol, and it should be typed as one.

\section{Displayed Text}

Text is displayed by indenting it from the left
margin.  Quotations are commonly displayed.  There
are short quotations
\begin{quote}
   This is a short quotation.  It consists of a
   single paragraph of text.  See how it is formatted.
\end{quote}
and longer ones.
\begin{quotation}
   This is a longer quotation.  It consists of two
   paragraphs of text, neither of which are
   particularly interesting.

   This is the second paragraph of the quotation.  It
   is just as dull as the first paragraph.
\end{quotation}
Another frequently-displayed structure is a list.
The following is an example of an \emph{itemized}
list.
\begin{itemize}
   \item This is the first item of an itemized list.
         Each item in the list is marked with a ``tick''.
         You don't have to worry about what kind of tick
         mark is used.

   \item This is the second item of the list.  It
         contains another list nested inside it.  The inner
         list is an \emph{enumerated} list.
         \begin{enumerate}
            \item This is the first item of an enumerated
                  list that is nested within the itemized list.

            \item This is the second item of the inner list.
                  \LaTeX\ allows you to nest lists deeper than
                  you really should.
         \end{enumerate}
         This is the rest of the second item of the outer
         list.  It is no more interesting than any other
         part of the item.
   \item This is the third item of the list.
\end{itemize}
You can even display poetry.
\begin{verse}
   There is an environment
    for verse \\             % The \\ command separates lines
   Whose features some poets % within a stanza.
   will curse.

                             % One or more blank lines separate stanzas.

   For instead of making\\
   Them do \emph{all} line breaking, \\
   It allows them to put too many words on a line when they'd rather be
   forced to be terse.
\end{verse}

Mathematical formulas may also be displayed.  A
displayed formula
is
one-line long; multiline
formulas require special formatting instructions.
   \[  \ip{\Gamma}{\psi'} = x'' + y^{2} + z_{i}^{n}\]
Don't start a paragraph with a displayed equation,
nor make one a paragraph by itself.

\end{document}               % End of document.

\section{Bench Huge Section 3}
% synthetic bench section 3
% This is a sample LaTeX input file.  (Version of 12 August 2004.)
%
% A '%' character causes TeX to ignore all remaining text on the line,
% and is used for comments like this one.

\documentclass{article}      % Specifies the document class

                             % The preamble begins here.
\title{An Example Document}  % Declares the document's title.
\author{Leslie Lamport}      % Declares the author's name.
\date{January 21, 1994}      % Deleting this command produces today's date.

\newcommand{\ip}[2]{(#1, #2)}
                             % Defines \ip{arg1}{arg2} to mean
                             % (arg1, arg2).

%\newcommand{\ip}[2]{\langle #1 | #2\rangle}
                             % This is an alternative definition of
                             % \ip that is commented out.

\begin{document}             % End of preamble and beginning of text.

\maketitle                   % Produces the title.

This is an example input file.  Comparing it with
the output it generates can show you how to
produce a simple document of your own.

\section{Ordinary Text}      % Produces section heading.  Lower-level
                             % sections are begun with similar
                             % \subsection and \subsubsection commands.

The ends  of words and sentences are marked
  by   spaces. It  doesn't matter how many
spaces    you type; one is as good as 100.  The
end of   a line counts as a space.

One   or more   blank lines denote the  end
of  a paragraph.

Since any number of consecutive spaces are treated
like a single one, the formatting of the input
file makes no difference to
      \LaTeX,                % The \LaTeX command generates the LaTeX logo.
but it makes a difference to you.  When you use
\LaTeX, making your input file as easy to read
as possible will be a great help as you write
your document and when you change it.  This sample
file shows how you can add comments to your own input
file.

Because printing is different from typewriting,
there are a number of things that you have to do
differently when preparing an input file than if
you were just typing the document directly.
Quotation marks like
       ``this''
have to be handled specially, as do quotes within
quotes:
       ``\,`this'            % \, separates the double and single quote.
        is what I just
        wrote, not  `that'\,''.

Dashes come in three sizes: an
       intra-word
dash, a medium dash for number ranges like
       1--2,
and a punctuation
       dash---like
this.

A sentence-ending space should be larger than the
space between words within a sentence.  You
sometimes have to type special commands in
conjunction with punctuation characters to get
this right, as in the following sentence.
       Gnats, gnus, etc.\ all  % `\ ' makes an inter-word space.
       begin with G\@.         % \@ marks end-of-sentence punctuation.
You should check the spaces after periods when
reading your output to make sure you haven't
forgotten any special cases.  Generating an
ellipsis
       \ldots\               % `\ ' is needed after `\ldots' because TeX
                             % ignores spaces after command names like \ldots
                             % made from \ + letters.
                             %
                             % Note how a `%' character causes TeX to ignore
                             % the end of the input line, so these blank lines
                             % do not start a new paragraph.
                             %
with the right spacing around the periods requires
a special command.

\LaTeX\ interprets some common characters as
commands, so you must type special commands to
generate them.  These characters include the
following:
       \$ \& \% \# \{ and \}.

In printing, text is usually emphasized with an
       \emph{italic}
type style.

\begin{em}
   A long segment of text can also be emphasized
   in this way.  Text within such a segment can be
   given \emph{additional} emphasis.
\end{em}

It is sometimes necessary to prevent \LaTeX\ from
breaking a line where it might otherwise do so.
This may be at a space, as between the ``Mr.''\ and
``Jones'' in
       ``Mr.~Jones'',        % ~ produces an unbreakable interword space.
or within a word---especially when the word is a
symbol like
       \mbox{\emph{itemnum}}
that makes little sense when hyphenated across
lines.

Footnotes\footnote{This is an example of a footnote.}
pose no problem.

\LaTeX\ is good at typesetting mathematical formulas
like
       \( x-3y + z = 7 \)
or
       \( a_{1} > x^{2n} + y^{2n} > x' \)
or
       \( \ip{A}{B} = \sum_{i} a_{i} b_{i} \).
The spaces you type in a formula are
ignored.  Remember that a letter like
       $x$                   % $ ... $  and  \( ... \)  are equivalent
is a formula when it denotes a mathematical
symbol, and it should be typed as one.

\section{Displayed Text}

Text is displayed by indenting it from the left
margin.  Quotations are commonly displayed.  There
are short quotations
\begin{quote}
   This is a short quotation.  It consists of a
   single paragraph of text.  See how it is formatted.
\end{quote}
and longer ones.
\begin{quotation}
   This is a longer quotation.  It consists of two
   paragraphs of text, neither of which are
   particularly interesting.

   This is the second paragraph of the quotation.  It
   is just as dull as the first paragraph.
\end{quotation}
Another frequently-displayed structure is a list.
The following is an example of an \emph{itemized}
list.
\begin{itemize}
   \item This is the first item of an itemized list.
         Each item in the list is marked with a ``tick''.
         You don't have to worry about what kind of tick
         mark is used.

   \item This is the second item of the list.  It
         contains another list nested inside it.  The inner
         list is an \emph{enumerated} list.
         \begin{enumerate}
            \item This is the first item of an enumerated
                  list that is nested within the itemized list.

            \item This is the second item of the inner list.
                  \LaTeX\ allows you to nest lists deeper than
                  you really should.
         \end{enumerate}
         This is the rest of the second item of the outer
         list.  It is no more interesting than any other
         part of the item.
   \item This is the third item of the list.
\end{itemize}
You can even display poetry.
\begin{verse}
   There is an environment
    for verse \\             % The \\ command separates lines
   Whose features some poets % within a stanza.
   will curse.

                             % One or more blank lines separate stanzas.

   For instead of making\\
   Them do \emph{all} line breaking, \\
   It allows them to put too many words on a line when they'd rather be
   forced to be terse.
\end{verse}

Mathematical formulas may also be displayed.  A
displayed formula
is
one-line long; multiline
formulas require special formatting instructions.
   \[  \ip{\Gamma}{\psi'} = x'' + y^{2} + z_{i}^{n}\]
Don't start a paragraph with a displayed equation,
nor make one a paragraph by itself.

\end{document}               % End of document.

\section{Bench Huge Section 4}
% synthetic bench section 4
% This is a sample LaTeX input file.  (Version of 12 August 2004.)
%
% A '%' character causes TeX to ignore all remaining text on the line,
% and is used for comments like this one.

\documentclass{article}      % Specifies the document class

                             % The preamble begins here.
\title{An Example Document}  % Declares the document's title.
\author{Leslie Lamport}      % Declares the author's name.
\date{January 21, 1994}      % Deleting this command produces today's date.

\newcommand{\ip}[2]{(#1, #2)}
                             % Defines \ip{arg1}{arg2} to mean
                             % (arg1, arg2).

%\newcommand{\ip}[2]{\langle #1 | #2\rangle}
                             % This is an alternative definition of
                             % \ip that is commented out.

\begin{document}             % End of preamble and beginning of text.

\maketitle                   % Produces the title.

This is an example input file.  Comparing it with
the output it generates can show you how to
produce a simple document of your own.

\section{Ordinary Text}      % Produces section heading.  Lower-level
                             % sections are begun with similar
                             % \subsection and \subsubsection commands.

The ends  of words and sentences are marked
  by   spaces. It  doesn't matter how many
spaces    you type; one is as good as 100.  The
end of   a line counts as a space.

One   or more   blank lines denote the  end
of  a paragraph.

Since any number of consecutive spaces are treated
like a single one, the formatting of the input
file makes no difference to
      \LaTeX,                % The \LaTeX command generates the LaTeX logo.
but it makes a difference to you.  When you use
\LaTeX, making your input file as easy to read
as possible will be a great help as you write
your document and when you change it.  This sample
file shows how you can add comments to your own input
file.

Because printing is different from typewriting,
there are a number of things that you have to do
differently when preparing an input file than if
you were just typing the document directly.
Quotation marks like
       ``this''
have to be handled specially, as do quotes within
quotes:
       ``\,`this'            % \, separates the double and single quote.
        is what I just
        wrote, not  `that'\,''.

Dashes come in three sizes: an
       intra-word
dash, a medium dash for number ranges like
       1--2,
and a punctuation
       dash---like
this.

A sentence-ending space should be larger than the
space between words within a sentence.  You
sometimes have to type special commands in
conjunction with punctuation characters to get
this right, as in the following sentence.
       Gnats, gnus, etc.\ all  % `\ ' makes an inter-word space.
       begin with G\@.         % \@ marks end-of-sentence punctuation.
You should check the spaces after periods when
reading your output to make sure you haven't
forgotten any special cases.  Generating an
ellipsis
       \ldots\               % `\ ' is needed after `\ldots' because TeX
                             % ignores spaces after command names like \ldots
                             % made from \ + letters.
                             %
                             % Note how a `%' character causes TeX to ignore
                             % the end of the input line, so these blank lines
                             % do not start a new paragraph.
                             %
with the right spacing around the periods requires
a special command.

\LaTeX\ interprets some common characters as
commands, so you must type special commands to
generate them.  These characters include the
following:
       \$ \& \% \# \{ and \}.

In printing, text is usually emphasized with an
       \emph{italic}
type style.

\begin{em}
   A long segment of text can also be emphasized
   in this way.  Text within such a segment can be
   given \emph{additional} emphasis.
\end{em}

It is sometimes necessary to prevent \LaTeX\ from
breaking a line where it might otherwise do so.
This may be at a space, as between the ``Mr.''\ and
``Jones'' in
       ``Mr.~Jones'',        % ~ produces an unbreakable interword space.
or within a word---especially when the word is a
symbol like
       \mbox{\emph{itemnum}}
that makes little sense when hyphenated across
lines.

Footnotes\footnote{This is an example of a footnote.}
pose no problem.

\LaTeX\ is good at typesetting mathematical formulas
like
       \( x-3y + z = 7 \)
or
       \( a_{1} > x^{2n} + y^{2n} > x' \)
or
       \( \ip{A}{B} = \sum_{i} a_{i} b_{i} \).
The spaces you type in a formula are
ignored.  Remember that a letter like
       $x$                   % $ ... $  and  \( ... \)  are equivalent
is a formula when it denotes a mathematical
symbol, and it should be typed as one.

\section{Displayed Text}

Text is displayed by indenting it from the left
margin.  Quotations are commonly displayed.  There
are short quotations
\begin{quote}
   This is a short quotation.  It consists of a
   single paragraph of text.  See how it is formatted.
\end{quote}
and longer ones.
\begin{quotation}
   This is a longer quotation.  It consists of two
   paragraphs of text, neither of which are
   particularly interesting.

   This is the second paragraph of the quotation.  It
   is just as dull as the first paragraph.
\end{quotation}
Another frequently-displayed structure is a list.
The following is an example of an \emph{itemized}
list.
\begin{itemize}
   \item This is the first item of an itemized list.
         Each item in the list is marked with a ``tick''.
         You don't have to worry about what kind of tick
         mark is used.

   \item This is the second item of the list.  It
         contains another list nested inside it.  The inner
         list is an \emph{enumerated} list.
         \begin{enumerate}
            \item This is the first item of an enumerated
                  list that is nested within the itemized list.

            \item This is the second item of the inner list.
                  \LaTeX\ allows you to nest lists deeper than
                  you really should.
         \end{enumerate}
         This is the rest of the second item of the outer
         list.  It is no more interesting than any other
         part of the item.
   \item This is the third item of the list.
\end{itemize}
You can even display poetry.
\begin{verse}
   There is an environment
    for verse \\             % The \\ command separates lines
   Whose features some poets % within a stanza.
   will curse.

                             % One or more blank lines separate stanzas.

   For instead of making\\
   Them do \emph{all} line breaking, \\
   It allows them to put too many words on a line when they'd rather be
   forced to be terse.
\end{verse}

Mathematical formulas may also be displayed.  A
displayed formula
is
one-line long; multiline
formulas require special formatting instructions.
   \[  \ip{\Gamma}{\psi'} = x'' + y^{2} + z_{i}^{n}\]
Don't start a paragraph with a displayed equation,
nor make one a paragraph by itself.

\end{document}               % End of document.

\section{Bench Huge Section 5}
% synthetic bench section 5
% This is a sample LaTeX input file.  (Version of 12 August 2004.)
%
% A '%' character causes TeX to ignore all remaining text on the line,
% and is used for comments like this one.

\documentclass{article}      % Specifies the document class

                             % The preamble begins here.
\title{An Example Document}  % Declares the document's title.
\author{Leslie Lamport}      % Declares the author's name.
\date{January 21, 1994}      % Deleting this command produces today's date.

\newcommand{\ip}[2]{(#1, #2)}
                             % Defines \ip{arg1}{arg2} to mean
                             % (arg1, arg2).

%\newcommand{\ip}[2]{\langle #1 | #2\rangle}
                             % This is an alternative definition of
                             % \ip that is commented out.

\begin{document}             % End of preamble and beginning of text.

\maketitle                   % Produces the title.

This is an example input file.  Comparing it with
the output it generates can show you how to
produce a simple document of your own.

\section{Ordinary Text}      % Produces section heading.  Lower-level
                             % sections are begun with similar
                             % \subsection and \subsubsection commands.

The ends  of words and sentences are marked
  by   spaces. It  doesn't matter how many
spaces    you type; one is as good as 100.  The
end of   a line counts as a space.

One   or more   blank lines denote the  end
of  a paragraph.

Since any number of consecutive spaces are treated
like a single one, the formatting of the input
file makes no difference to
      \LaTeX,                % The \LaTeX command generates the LaTeX logo.
but it makes a difference to you.  When you use
\LaTeX, making your input file as easy to read
as possible will be a great help as you write
your document and when you change it.  This sample
file shows how you can add comments to your own input
file.

Because printing is different from typewriting,
there are a number of things that you have to do
differently when preparing an input file than if
you were just typing the document directly.
Quotation marks like
       ``this''
have to be handled specially, as do quotes within
quotes:
       ``\,`this'            % \, separates the double and single quote.
        is what I just
        wrote, not  `that'\,''.

Dashes come in three sizes: an
       intra-word
dash, a medium dash for number ranges like
       1--2,
and a punctuation
       dash---like
this.

A sentence-ending space should be larger than the
space between words within a sentence.  You
sometimes have to type special commands in
conjunction with punctuation characters to get
this right, as in the following sentence.
       Gnats, gnus, etc.\ all  % `\ ' makes an inter-word space.
       begin with G\@.         % \@ marks end-of-sentence punctuation.
You should check the spaces after periods when
reading your output to make sure you haven't
forgotten any special cases.  Generating an
ellipsis
       \ldots\               % `\ ' is needed after `\ldots' because TeX
                             % ignores spaces after command names like \ldots
                             % made from \ + letters.
                             %
                             % Note how a `%' character causes TeX to ignore
                             % the end of the input line, so these blank lines
                             % do not start a new paragraph.
                             %
with the right spacing around the periods requires
a special command.

\LaTeX\ interprets some common characters as
commands, so you must type special commands to
generate them.  These characters include the
following:
       \$ \& \% \# \{ and \}.

In printing, text is usually emphasized with an
       \emph{italic}
type style.

\begin{em}
   A long segment of text can also be emphasized
   in this way.  Text within such a segment can be
   given \emph{additional} emphasis.
\end{em}

It is sometimes necessary to prevent \LaTeX\ from
breaking a line where it might otherwise do so.
This may be at a space, as between the ``Mr.''\ and
``Jones'' in
       ``Mr.~Jones'',        % ~ produces an unbreakable interword space.
or within a word---especially when the word is a
symbol like
       \mbox{\emph{itemnum}}
that makes little sense when hyphenated across
lines.

Footnotes\footnote{This is an example of a footnote.}
pose no problem.

\LaTeX\ is good at typesetting mathematical formulas
like
       \( x-3y + z = 7 \)
or
       \( a_{1} > x^{2n} + y^{2n} > x' \)
or
       \( \ip{A}{B} = \sum_{i} a_{i} b_{i} \).
The spaces you type in a formula are
ignored.  Remember that a letter like
       $x$                   % $ ... $  and  \( ... \)  are equivalent
is a formula when it denotes a mathematical
symbol, and it should be typed as one.

\section{Displayed Text}

Text is displayed by indenting it from the left
margin.  Quotations are commonly displayed.  There
are short quotations
\begin{quote}
   This is a short quotation.  It consists of a
   single paragraph of text.  See how it is formatted.
\end{quote}
and longer ones.
\begin{quotation}
   This is a longer quotation.  It consists of two
   paragraphs of text, neither of which are
   particularly interesting.

   This is the second paragraph of the quotation.  It
   is just as dull as the first paragraph.
\end{quotation}
Another frequently-displayed structure is a list.
The following is an example of an \emph{itemized}
list.
\begin{itemize}
   \item This is the first item of an itemized list.
         Each item in the list is marked with a ``tick''.
         You don't have to worry about what kind of tick
         mark is used.

   \item This is the second item of the list.  It
         contains another list nested inside it.  The inner
         list is an \emph{enumerated} list.
         \begin{enumerate}
            \item This is the first item of an enumerated
                  list that is nested within the itemized list.

            \item This is the second item of the inner list.
                  \LaTeX\ allows you to nest lists deeper than
                  you really should.
         \end{enumerate}
         This is the rest of the second item of the outer
         list.  It is no more interesting than any other
         part of the item.
   \item This is the third item of the list.
\end{itemize}
You can even display poetry.
\begin{verse}
   There is an environment
    for verse \\             % The \\ command separates lines
   Whose features some poets % within a stanza.
   will curse.

                             % One or more blank lines separate stanzas.

   For instead of making\\
   Them do \emph{all} line breaking, \\
   It allows them to put too many words on a line when they'd rather be
   forced to be terse.
\end{verse}

Mathematical formulas may also be displayed.  A
displayed formula
is
one-line long; multiline
formulas require special formatting instructions.
   \[  \ip{\Gamma}{\psi'} = x'' + y^{2} + z_{i}^{n}\]
Don't start a paragraph with a displayed equation,
nor make one a paragraph by itself.

\end{document}               % End of document.

\section{Bench Huge Section 6}
% synthetic bench section 6
% This is a sample LaTeX input file.  (Version of 12 August 2004.)
%
% A '%' character causes TeX to ignore all remaining text on the line,
% and is used for comments like this one.

\documentclass{article}      % Specifies the document class

                             % The preamble begins here.
\title{An Example Document}  % Declares the document's title.
\author{Leslie Lamport}      % Declares the author's name.
\date{January 21, 1994}      % Deleting this command produces today's date.

\newcommand{\ip}[2]{(#1, #2)}
                             % Defines \ip{arg1}{arg2} to mean
                             % (arg1, arg2).

%\newcommand{\ip}[2]{\langle #1 | #2\rangle}
                             % This is an alternative definition of
                             % \ip that is commented out.

\begin{document}             % End of preamble and beginning of text.

\maketitle                   % Produces the title.

This is an example input file.  Comparing it with
the output it generates can show you how to
produce a simple document of your own.

\section{Ordinary Text}      % Produces section heading.  Lower-level
                             % sections are begun with similar
                             % \subsection and \subsubsection commands.

The ends  of words and sentences are marked
  by   spaces. It  doesn't matter how many
spaces    you type; one is as good as 100.  The
end of   a line counts as a space.

One   or more   blank lines denote the  end
of  a paragraph.

Since any number of consecutive spaces are treated
like a single one, the formatting of the input
file makes no difference to
      \LaTeX,                % The \LaTeX command generates the LaTeX logo.
but it makes a difference to you.  When you use
\LaTeX, making your input file as easy to read
as possible will be a great help as you write
your document and when you change it.  This sample
file shows how you can add comments to your own input
file.

Because printing is different from typewriting,
there are a number of things that you have to do
differently when preparing an input file than if
you were just typing the document directly.
Quotation marks like
       ``this''
have to be handled specially, as do quotes within
quotes:
       ``\,`this'            % \, separates the double and single quote.
        is what I just
        wrote, not  `that'\,''.

Dashes come in three sizes: an
       intra-word
dash, a medium dash for number ranges like
       1--2,
and a punctuation
       dash---like
this.

A sentence-ending space should be larger than the
space between words within a sentence.  You
sometimes have to type special commands in
conjunction with punctuation characters to get
this right, as in the following sentence.
       Gnats, gnus, etc.\ all  % `\ ' makes an inter-word space.
       begin with G\@.         % \@ marks end-of-sentence punctuation.
You should check the spaces after periods when
reading your output to make sure you haven't
forgotten any special cases.  Generating an
ellipsis
       \ldots\               % `\ ' is needed after `\ldots' because TeX
                             % ignores spaces after command names like \ldots
                             % made from \ + letters.
                             %
                             % Note how a `%' character causes TeX to ignore
                             % the end of the input line, so these blank lines
                             % do not start a new paragraph.
                             %
with the right spacing around the periods requires
a special command.

\LaTeX\ interprets some common characters as
commands, so you must type special commands to
generate them.  These characters include the
following:
       \$ \& \% \# \{ and \}.

In printing, text is usually emphasized with an
       \emph{italic}
type style.

\begin{em}
   A long segment of text can also be emphasized
   in this way.  Text within such a segment can be
   given \emph{additional} emphasis.
\end{em}

It is sometimes necessary to prevent \LaTeX\ from
breaking a line where it might otherwise do so.
This may be at a space, as between the ``Mr.''\ and
``Jones'' in
       ``Mr.~Jones'',        % ~ produces an unbreakable interword space.
or within a word---especially when the word is a
symbol like
       \mbox{\emph{itemnum}}
that makes little sense when hyphenated across
lines.

Footnotes\footnote{This is an example of a footnote.}
pose no problem.

\LaTeX\ is good at typesetting mathematical formulas
like
       \( x-3y + z = 7 \)
or
       \( a_{1} > x^{2n} + y^{2n} > x' \)
or
       \( \ip{A}{B} = \sum_{i} a_{i} b_{i} \).
The spaces you type in a formula are
ignored.  Remember that a letter like
       $x$                   % $ ... $  and  \( ... \)  are equivalent
is a formula when it denotes a mathematical
symbol, and it should be typed as one.

\section{Displayed Text}

Text is displayed by indenting it from the left
margin.  Quotations are commonly displayed.  There
are short quotations
\begin{quote}
   This is a short quotation.  It consists of a
   single paragraph of text.  See how it is formatted.
\end{quote}
and longer ones.
\begin{quotation}
   This is a longer quotation.  It consists of two
   paragraphs of text, neither of which are
   particularly interesting.

   This is the second paragraph of the quotation.  It
   is just as dull as the first paragraph.
\end{quotation}
Another frequently-displayed structure is a list.
The following is an example of an \emph{itemized}
list.
\begin{itemize}
   \item This is the first item of an itemized list.
         Each item in the list is marked with a ``tick''.
         You don't have to worry about what kind of tick
         mark is used.

   \item This is the second item of the list.  It
         contains another list nested inside it.  The inner
         list is an \emph{enumerated} list.
         \begin{enumerate}
            \item This is the first item of an enumerated
                  list that is nested within the itemized list.

            \item This is the second item of the inner list.
                  \LaTeX\ allows you to nest lists deeper than
                  you really should.
         \end{enumerate}
         This is the rest of the second item of the outer
         list.  It is no more interesting than any other
         part of the item.
   \item This is the third item of the list.
\end{itemize}
You can even display poetry.
\begin{verse}
   There is an environment
    for verse \\             % The \\ command separates lines
   Whose features some poets % within a stanza.
   will curse.

                             % One or more blank lines separate stanzas.

   For instead of making\\
   Them do \emph{all} line breaking, \\
   It allows them to put too many words on a line when they'd rather be
   forced to be terse.
\end{verse}

Mathematical formulas may also be displayed.  A
displayed formula
is
one-line long; multiline
formulas require special formatting instructions.
   \[  \ip{\Gamma}{\psi'} = x'' + y^{2} + z_{i}^{n}\]
Don't start a paragraph with a displayed equation,
nor make one a paragraph by itself.

\end{document}               % End of document.

\section{Bench Huge Section 7}
% synthetic bench section 7
% This is a sample LaTeX input file.  (Version of 12 August 2004.)
%
% A '%' character causes TeX to ignore all remaining text on the line,
% and is used for comments like this one.

\documentclass{article}      % Specifies the document class

                             % The preamble begins here.
\title{An Example Document}  % Declares the document's title.
\author{Leslie Lamport}      % Declares the author's name.
\date{January 21, 1994}      % Deleting this command produces today's date.

\newcommand{\ip}[2]{(#1, #2)}
                             % Defines \ip{arg1}{arg2} to mean
                             % (arg1, arg2).

%\newcommand{\ip}[2]{\langle #1 | #2\rangle}
                             % This is an alternative definition of
                             % \ip that is commented out.

\begin{document}             % End of preamble and beginning of text.

\maketitle                   % Produces the title.

This is an example input file.  Comparing it with
the output it generates can show you how to
produce a simple document of your own.

\section{Ordinary Text}      % Produces section heading.  Lower-level
                             % sections are begun with similar
                             % \subsection and \subsubsection commands.

The ends  of words and sentences are marked
  by   spaces. It  doesn't matter how many
spaces    you type; one is as good as 100.  The
end of   a line counts as a space.

One   or more   blank lines denote the  end
of  a paragraph.

Since any number of consecutive spaces are treated
like a single one, the formatting of the input
file makes no difference to
      \LaTeX,                % The \LaTeX command generates the LaTeX logo.
but it makes a difference to you.  When you use
\LaTeX, making your input file as easy to read
as possible will be a great help as you write
your document and when you change it.  This sample
file shows how you can add comments to your own input
file.

Because printing is different from typewriting,
there are a number of things that you have to do
differently when preparing an input file than if
you were just typing the document directly.
Quotation marks like
       ``this''
have to be handled specially, as do quotes within
quotes:
       ``\,`this'            % \, separates the double and single quote.
        is what I just
        wrote, not  `that'\,''.

Dashes come in three sizes: an
       intra-word
dash, a medium dash for number ranges like
       1--2,
and a punctuation
       dash---like
this.

A sentence-ending space should be larger than the
space between words within a sentence.  You
sometimes have to type special commands in
conjunction with punctuation characters to get
this right, as in the following sentence.
       Gnats, gnus, etc.\ all  % `\ ' makes an inter-word space.
       begin with G\@.         % \@ marks end-of-sentence punctuation.
You should check the spaces after periods when
reading your output to make sure you haven't
forgotten any special cases.  Generating an
ellipsis
       \ldots\               % `\ ' is needed after `\ldots' because TeX
                             % ignores spaces after command names like \ldots
                             % made from \ + letters.
                             %
                             % Note how a `%' character causes TeX to ignore
                             % the end of the input line, so these blank lines
                             % do not start a new paragraph.
                             %
with the right spacing around the periods requires
a special command.

\LaTeX\ interprets some common characters as
commands, so you must type special commands to
generate them.  These characters include the
following:
       \$ \& \% \# \{ and \}.

In printing, text is usually emphasized with an
       \emph{italic}
type style.

\begin{em}
   A long segment of text can also be emphasized
   in this way.  Text within such a segment can be
   given \emph{additional} emphasis.
\end{em}

It is sometimes necessary to prevent \LaTeX\ from
breaking a line where it might otherwise do so.
This may be at a space, as between the ``Mr.''\ and
``Jones'' in
       ``Mr.~Jones'',        % ~ produces an unbreakable interword space.
or within a word---especially when the word is a
symbol like
       \mbox{\emph{itemnum}}
that makes little sense when hyphenated across
lines.

Footnotes\footnote{This is an example of a footnote.}
pose no problem.

\LaTeX\ is good at typesetting mathematical formulas
like
       \( x-3y + z = 7 \)
or
       \( a_{1} > x^{2n} + y^{2n} > x' \)
or
       \( \ip{A}{B} = \sum_{i} a_{i} b_{i} \).
The spaces you type in a formula are
ignored.  Remember that a letter like
       $x$                   % $ ... $  and  \( ... \)  are equivalent
is a formula when it denotes a mathematical
symbol, and it should be typed as one.

\section{Displayed Text}

Text is displayed by indenting it from the left
margin.  Quotations are commonly displayed.  There
are short quotations
\begin{quote}
   This is a short quotation.  It consists of a
   single paragraph of text.  See how it is formatted.
\end{quote}
and longer ones.
\begin{quotation}
   This is a longer quotation.  It consists of two
   paragraphs of text, neither of which are
   particularly interesting.

   This is the second paragraph of the quotation.  It
   is just as dull as the first paragraph.
\end{quotation}
Another frequently-displayed structure is a list.
The following is an example of an \emph{itemized}
list.
\begin{itemize}
   \item This is the first item of an itemized list.
         Each item in the list is marked with a ``tick''.
         You don't have to worry about what kind of tick
         mark is used.

   \item This is the second item of the list.  It
         contains another list nested inside it.  The inner
         list is an \emph{enumerated} list.
         \begin{enumerate}
            \item This is the first item of an enumerated
                  list that is nested within the itemized list.

            \item This is the second item of the inner list.
                  \LaTeX\ allows you to nest lists deeper than
                  you really should.
         \end{enumerate}
         This is the rest of the second item of the outer
         list.  It is no more interesting than any other
         part of the item.
   \item This is the third item of the list.
\end{itemize}
You can even display poetry.
\begin{verse}
   There is an environment
    for verse \\             % The \\ command separates lines
   Whose features some poets % within a stanza.
   will curse.

                             % One or more blank lines separate stanzas.

   For instead of making\\
   Them do \emph{all} line breaking, \\
   It allows them to put too many words on a line when they'd rather be
   forced to be terse.
\end{verse}

Mathematical formulas may also be displayed.  A
displayed formula
is
one-line long; multiline
formulas require special formatting instructions.
   \[  \ip{\Gamma}{\psi'} = x'' + y^{2} + z_{i}^{n}\]
Don't start a paragraph with a displayed equation,
nor make one a paragraph by itself.

\end{document}               % End of document.

\section{Bench Huge Section 8}
% synthetic bench section 8
% This is a sample LaTeX input file.  (Version of 12 August 2004.)
%
% A '%' character causes TeX to ignore all remaining text on the line,
% and is used for comments like this one.

\documentclass{article}      % Specifies the document class

                             % The preamble begins here.
\title{An Example Document}  % Declares the document's title.
\author{Leslie Lamport}      % Declares the author's name.
\date{January 21, 1994}      % Deleting this command produces today's date.

\newcommand{\ip}[2]{(#1, #2)}
                             % Defines \ip{arg1}{arg2} to mean
                             % (arg1, arg2).

%\newcommand{\ip}[2]{\langle #1 | #2\rangle}
                             % This is an alternative definition of
                             % \ip that is commented out.

\begin{document}             % End of preamble and beginning of text.

\maketitle                   % Produces the title.

This is an example input file.  Comparing it with
the output it generates can show you how to
produce a simple document of your own.

\section{Ordinary Text}      % Produces section heading.  Lower-level
                             % sections are begun with similar
                             % \subsection and \subsubsection commands.

The ends  of words and sentences are marked
  by   spaces. It  doesn't matter how many
spaces    you type; one is as good as 100.  The
end of   a line counts as a space.

One   or more   blank lines denote the  end
of  a paragraph.

Since any number of consecutive spaces are treated
like a single one, the formatting of the input
file makes no difference to
      \LaTeX,                % The \LaTeX command generates the LaTeX logo.
but it makes a difference to you.  When you use
\LaTeX, making your input file as easy to read
as possible will be a great help as you write
your document and when you change it.  This sample
file shows how you can add comments to your own input
file.

Because printing is different from typewriting,
there are a number of things that you have to do
differently when preparing an input file than if
you were just typing the document directly.
Quotation marks like
       ``this''
have to be handled specially, as do quotes within
quotes:
       ``\,`this'            % \, separates the double and single quote.
        is what I just
        wrote, not  `that'\,''.

Dashes come in three sizes: an
       intra-word
dash, a medium dash for number ranges like
       1--2,
and a punctuation
       dash---like
this.

A sentence-ending space should be larger than the
space between words within a sentence.  You
sometimes have to type special commands in
conjunction with punctuation characters to get
this right, as in the following sentence.
       Gnats, gnus, etc.\ all  % `\ ' makes an inter-word space.
       begin with G\@.         % \@ marks end-of-sentence punctuation.
You should check the spaces after periods when
reading your output to make sure you haven't
forgotten any special cases.  Generating an
ellipsis
       \ldots\               % `\ ' is needed after `\ldots' because TeX
                             % ignores spaces after command names like \ldots
                             % made from \ + letters.
                             %
                             % Note how a `%' character causes TeX to ignore
                             % the end of the input line, so these blank lines
                             % do not start a new paragraph.
                             %
with the right spacing around the periods requires
a special command.

\LaTeX\ interprets some common characters as
commands, so you must type special commands to
generate them.  These characters include the
following:
       \$ \& \% \# \{ and \}.

In printing, text is usually emphasized with an
       \emph{italic}
type style.

\begin{em}
   A long segment of text can also be emphasized
   in this way.  Text within such a segment can be
   given \emph{additional} emphasis.
\end{em}

It is sometimes necessary to prevent \LaTeX\ from
breaking a line where it might otherwise do so.
This may be at a space, as between the ``Mr.''\ and
``Jones'' in
       ``Mr.~Jones'',        % ~ produces an unbreakable interword space.
or within a word---especially when the word is a
symbol like
       \mbox{\emph{itemnum}}
that makes little sense when hyphenated across
lines.

Footnotes\footnote{This is an example of a footnote.}
pose no problem.

\LaTeX\ is good at typesetting mathematical formulas
like
       \( x-3y + z = 7 \)
or
       \( a_{1} > x^{2n} + y^{2n} > x' \)
or
       \( \ip{A}{B} = \sum_{i} a_{i} b_{i} \).
The spaces you type in a formula are
ignored.  Remember that a letter like
       $x$                   % $ ... $  and  \( ... \)  are equivalent
is a formula when it denotes a mathematical
symbol, and it should be typed as one.

\section{Displayed Text}

Text is displayed by indenting it from the left
margin.  Quotations are commonly displayed.  There
are short quotations
\begin{quote}
   This is a short quotation.  It consists of a
   single paragraph of text.  See how it is formatted.
\end{quote}
and longer ones.
\begin{quotation}
   This is a longer quotation.  It consists of two
   paragraphs of text, neither of which are
   particularly interesting.

   This is the second paragraph of the quotation.  It
   is just as dull as the first paragraph.
\end{quotation}
Another frequently-displayed structure is a list.
The following is an example of an \emph{itemized}
list.
\begin{itemize}
   \item This is the first item of an itemized list.
         Each item in the list is marked with a ``tick''.
         You don't have to worry about what kind of tick
         mark is used.

   \item This is the second item of the list.  It
         contains another list nested inside it.  The inner
         list is an \emph{enumerated} list.
         \begin{enumerate}
            \item This is the first item of an enumerated
                  list that is nested within the itemized list.

            \item This is the second item of the inner list.
                  \LaTeX\ allows you to nest lists deeper than
                  you really should.
         \end{enumerate}
         This is the rest of the second item of the outer
         list.  It is no more interesting than any other
         part of the item.
   \item This is the third item of the list.
\end{itemize}
You can even display poetry.
\begin{verse}
   There is an environment
    for verse \\             % The \\ command separates lines
   Whose features some poets % within a stanza.
   will curse.

                             % One or more blank lines separate stanzas.

   For instead of making\\
   Them do \emph{all} line breaking, \\
   It allows them to put too many words on a line when they'd rather be
   forced to be terse.
\end{verse}

Mathematical formulas may also be displayed.  A
displayed formula
is
one-line long; multiline
formulas require special formatting instructions.
   \[  \ip{\Gamma}{\psi'} = x'' + y^{2} + z_{i}^{n}\]
Don't start a paragraph with a displayed equation,
nor make one a paragraph by itself.

\end{document}               % End of document.

\section{Bench Huge Section 9}
% synthetic bench section 9
% This is a sample LaTeX input file.  (Version of 12 August 2004.)
%
% A '%' character causes TeX to ignore all remaining text on the line,
% and is used for comments like this one.

\documentclass{article}      % Specifies the document class

                             % The preamble begins here.
\title{An Example Document}  % Declares the document's title.
\author{Leslie Lamport}      % Declares the author's name.
\date{January 21, 1994}      % Deleting this command produces today's date.

\newcommand{\ip}[2]{(#1, #2)}
                             % Defines \ip{arg1}{arg2} to mean
                             % (arg1, arg2).

%\newcommand{\ip}[2]{\langle #1 | #2\rangle}
                             % This is an alternative definition of
                             % \ip that is commented out.

\begin{document}             % End of preamble and beginning of text.

\maketitle                   % Produces the title.

This is an example input file.  Comparing it with
the output it generates can show you how to
produce a simple document of your own.

\section{Ordinary Text}      % Produces section heading.  Lower-level
                             % sections are begun with similar
                             % \subsection and \subsubsection commands.

The ends  of words and sentences are marked
  by   spaces. It  doesn't matter how many
spaces    you type; one is as good as 100.  The
end of   a line counts as a space.

One   or more   blank lines denote the  end
of  a paragraph.

Since any number of consecutive spaces are treated
like a single one, the formatting of the input
file makes no difference to
      \LaTeX,                % The \LaTeX command generates the LaTeX logo.
but it makes a difference to you.  When you use
\LaTeX, making your input file as easy to read
as possible will be a great help as you write
your document and when you change it.  This sample
file shows how you can add comments to your own input
file.

Because printing is different from typewriting,
there are a number of things that you have to do
differently when preparing an input file than if
you were just typing the document directly.
Quotation marks like
       ``this''
have to be handled specially, as do quotes within
quotes:
       ``\,`this'            % \, separates the double and single quote.
        is what I just
        wrote, not  `that'\,''.

Dashes come in three sizes: an
       intra-word
dash, a medium dash for number ranges like
       1--2,
and a punctuation
       dash---like
this.

A sentence-ending space should be larger than the
space between words within a sentence.  You
sometimes have to type special commands in
conjunction with punctuation characters to get
this right, as in the following sentence.
       Gnats, gnus, etc.\ all  % `\ ' makes an inter-word space.
       begin with G\@.         % \@ marks end-of-sentence punctuation.
You should check the spaces after periods when
reading your output to make sure you haven't
forgotten any special cases.  Generating an
ellipsis
       \ldots\               % `\ ' is needed after `\ldots' because TeX
                             % ignores spaces after command names like \ldots
                             % made from \ + letters.
                             %
                             % Note how a `%' character causes TeX to ignore
                             % the end of the input line, so these blank lines
                             % do not start a new paragraph.
                             %
with the right spacing around the periods requires
a special command.

\LaTeX\ interprets some common characters as
commands, so you must type special commands to
generate them.  These characters include the
following:
       \$ \& \% \# \{ and \}.

In printing, text is usually emphasized with an
       \emph{italic}
type style.

\begin{em}
   A long segment of text can also be emphasized
   in this way.  Text within such a segment can be
   given \emph{additional} emphasis.
\end{em}

It is sometimes necessary to prevent \LaTeX\ from
breaking a line where it might otherwise do so.
This may be at a space, as between the ``Mr.''\ and
``Jones'' in
       ``Mr.~Jones'',        % ~ produces an unbreakable interword space.
or within a word---especially when the word is a
symbol like
       \mbox{\emph{itemnum}}
that makes little sense when hyphenated across
lines.

Footnotes\footnote{This is an example of a footnote.}
pose no problem.

\LaTeX\ is good at typesetting mathematical formulas
like
       \( x-3y + z = 7 \)
or
       \( a_{1} > x^{2n} + y^{2n} > x' \)
or
       \( \ip{A}{B} = \sum_{i} a_{i} b_{i} \).
The spaces you type in a formula are
ignored.  Remember that a letter like
       $x$                   % $ ... $  and  \( ... \)  are equivalent
is a formula when it denotes a mathematical
symbol, and it should be typed as one.

\section{Displayed Text}

Text is displayed by indenting it from the left
margin.  Quotations are commonly displayed.  There
are short quotations
\begin{quote}
   This is a short quotation.  It consists of a
   single paragraph of text.  See how it is formatted.
\end{quote}
and longer ones.
\begin{quotation}
   This is a longer quotation.  It consists of two
   paragraphs of text, neither of which are
   particularly interesting.

   This is the second paragraph of the quotation.  It
   is just as dull as the first paragraph.
\end{quotation}
Another frequently-displayed structure is a list.
The following is an example of an \emph{itemized}
list.
\begin{itemize}
   \item This is the first item of an itemized list.
         Each item in the list is marked with a ``tick''.
         You don't have to worry about what kind of tick
         mark is used.

   \item This is the second item of the list.  It
         contains another list nested inside it.  The inner
         list is an \emph{enumerated} list.
         \begin{enumerate}
            \item This is the first item of an enumerated
                  list that is nested within the itemized list.

            \item This is the second item of the inner list.
                  \LaTeX\ allows you to nest lists deeper than
                  you really should.
         \end{enumerate}
         This is the rest of the second item of the outer
         list.  It is no more interesting than any other
         part of the item.
   \item This is the third item of the list.
\end{itemize}
You can even display poetry.
\begin{verse}
   There is an environment
    for verse \\             % The \\ command separates lines
   Whose features some poets % within a stanza.
   will curse.

                             % One or more blank lines separate stanzas.

   For instead of making\\
   Them do \emph{all} line breaking, \\
   It allows them to put too many words on a line when they'd rather be
   forced to be terse.
\end{verse}

Mathematical formulas may also be displayed.  A
displayed formula
is
one-line long; multiline
formulas require special formatting instructions.
   \[  \ip{\Gamma}{\psi'} = x'' + y^{2} + z_{i}^{n}\]
Don't start a paragraph with a displayed equation,
nor make one a paragraph by itself.

\end{document}               % End of document.

\section{Bench Huge Section 10}
% synthetic bench section 10
% This is a sample LaTeX input file.  (Version of 12 August 2004.)
%
% A '%' character causes TeX to ignore all remaining text on the line,
% and is used for comments like this one.

\documentclass{article}      % Specifies the document class

                             % The preamble begins here.
\title{An Example Document}  % Declares the document's title.
\author{Leslie Lamport}      % Declares the author's name.
\date{January 21, 1994}      % Deleting this command produces today's date.

\newcommand{\ip}[2]{(#1, #2)}
                             % Defines \ip{arg1}{arg2} to mean
                             % (arg1, arg2).

%\newcommand{\ip}[2]{\langle #1 | #2\rangle}
                             % This is an alternative definition of
                             % \ip that is commented out.

\begin{document}             % End of preamble and beginning of text.

\maketitle                   % Produces the title.

This is an example input file.  Comparing it with
the output it generates can show you how to
produce a simple document of your own.

\section{Ordinary Text}      % Produces section heading.  Lower-level
                             % sections are begun with similar
                             % \subsection and \subsubsection commands.

The ends  of words and sentences are marked
  by   spaces. It  doesn't matter how many
spaces    you type; one is as good as 100.  The
end of   a line counts as a space.

One   or more   blank lines denote the  end
of  a paragraph.

Since any number of consecutive spaces are treated
like a single one, the formatting of the input
file makes no difference to
      \LaTeX,                % The \LaTeX command generates the LaTeX logo.
but it makes a difference to you.  When you use
\LaTeX, making your input file as easy to read
as possible will be a great help as you write
your document and when you change it.  This sample
file shows how you can add comments to your own input
file.

Because printing is different from typewriting,
there are a number of things that you have to do
differently when preparing an input file than if
you were just typing the document directly.
Quotation marks like
       ``this''
have to be handled specially, as do quotes within
quotes:
       ``\,`this'            % \, separates the double and single quote.
        is what I just
        wrote, not  `that'\,''.

Dashes come in three sizes: an
       intra-word
dash, a medium dash for number ranges like
       1--2,
and a punctuation
       dash---like
this.

A sentence-ending space should be larger than the
space between words within a sentence.  You
sometimes have to type special commands in
conjunction with punctuation characters to get
this right, as in the following sentence.
       Gnats, gnus, etc.\ all  % `\ ' makes an inter-word space.
       begin with G\@.         % \@ marks end-of-sentence punctuation.
You should check the spaces after periods when
reading your output to make sure you haven't
forgotten any special cases.  Generating an
ellipsis
       \ldots\               % `\ ' is needed after `\ldots' because TeX
                             % ignores spaces after command names like \ldots
                             % made from \ + letters.
                             %
                             % Note how a `%' character causes TeX to ignore
                             % the end of the input line, so these blank lines
                             % do not start a new paragraph.
                             %
with the right spacing around the periods requires
a special command.

\LaTeX\ interprets some common characters as
commands, so you must type special commands to
generate them.  These characters include the
following:
       \$ \& \% \# \{ and \}.

In printing, text is usually emphasized with an
       \emph{italic}
type style.

\begin{em}
   A long segment of text can also be emphasized
   in this way.  Text within such a segment can be
   given \emph{additional} emphasis.
\end{em}

It is sometimes necessary to prevent \LaTeX\ from
breaking a line where it might otherwise do so.
This may be at a space, as between the ``Mr.''\ and
``Jones'' in
       ``Mr.~Jones'',        % ~ produces an unbreakable interword space.
or within a word---especially when the word is a
symbol like
       \mbox{\emph{itemnum}}
that makes little sense when hyphenated across
lines.

Footnotes\footnote{This is an example of a footnote.}
pose no problem.

\LaTeX\ is good at typesetting mathematical formulas
like
       \( x-3y + z = 7 \)
or
       \( a_{1} > x^{2n} + y^{2n} > x' \)
or
       \( \ip{A}{B} = \sum_{i} a_{i} b_{i} \).
The spaces you type in a formula are
ignored.  Remember that a letter like
       $x$                   % $ ... $  and  \( ... \)  are equivalent
is a formula when it denotes a mathematical
symbol, and it should be typed as one.

\section{Displayed Text}

Text is displayed by indenting it from the left
margin.  Quotations are commonly displayed.  There
are short quotations
\begin{quote}
   This is a short quotation.  It consists of a
   single paragraph of text.  See how it is formatted.
\end{quote}
and longer ones.
\begin{quotation}
   This is a longer quotation.  It consists of two
   paragraphs of text, neither of which are
   particularly interesting.

   This is the second paragraph of the quotation.  It
   is just as dull as the first paragraph.
\end{quotation}
Another frequently-displayed structure is a list.
The following is an example of an \emph{itemized}
list.
\begin{itemize}
   \item This is the first item of an itemized list.
         Each item in the list is marked with a ``tick''.
         You don't have to worry about what kind of tick
         mark is used.

   \item This is the second item of the list.  It
         contains another list nested inside it.  The inner
         list is an \emph{enumerated} list.
         \begin{enumerate}
            \item This is the first item of an enumerated
                  list that is nested within the itemized list.

            \item This is the second item of the inner list.
                  \LaTeX\ allows you to nest lists deeper than
                  you really should.
         \end{enumerate}
         This is the rest of the second item of the outer
         list.  It is no more interesting than any other
         part of the item.
   \item This is the third item of the list.
\end{itemize}
You can even display poetry.
\begin{verse}
   There is an environment
    for verse \\             % The \\ command separates lines
   Whose features some poets % within a stanza.
   will curse.

                             % One or more blank lines separate stanzas.

   For instead of making\\
   Them do \emph{all} line breaking, \\
   It allows them to put too many words on a line when they'd rather be
   forced to be terse.
\end{verse}

Mathematical formulas may also be displayed.  A
displayed formula
is
one-line long; multiline
formulas require special formatting instructions.
   \[  \ip{\Gamma}{\psi'} = x'' + y^{2} + z_{i}^{n}\]
Don't start a paragraph with a displayed equation,
nor make one a paragraph by itself.

\end{document}               % End of document.

\section{Bench Huge Section 11}
% synthetic bench section 11
% This is a sample LaTeX input file.  (Version of 12 August 2004.)
%
% A '%' character causes TeX to ignore all remaining text on the line,
% and is used for comments like this one.

\documentclass{article}      % Specifies the document class

                             % The preamble begins here.
\title{An Example Document}  % Declares the document's title.
\author{Leslie Lamport}      % Declares the author's name.
\date{January 21, 1994}      % Deleting this command produces today's date.

\newcommand{\ip}[2]{(#1, #2)}
                             % Defines \ip{arg1}{arg2} to mean
                             % (arg1, arg2).

%\newcommand{\ip}[2]{\langle #1 | #2\rangle}
                             % This is an alternative definition of
                             % \ip that is commented out.

\begin{document}             % End of preamble and beginning of text.

\maketitle                   % Produces the title.

This is an example input file.  Comparing it with
the output it generates can show you how to
produce a simple document of your own.

\section{Ordinary Text}      % Produces section heading.  Lower-level
                             % sections are begun with similar
                             % \subsection and \subsubsection commands.

The ends  of words and sentences are marked
  by   spaces. It  doesn't matter how many
spaces    you type; one is as good as 100.  The
end of   a line counts as a space.

One   or more   blank lines denote the  end
of  a paragraph.

Since any number of consecutive spaces are treated
like a single one, the formatting of the input
file makes no difference to
      \LaTeX,                % The \LaTeX command generates the LaTeX logo.
but it makes a difference to you.  When you use
\LaTeX, making your input file as easy to read
as possible will be a great help as you write
your document and when you change it.  This sample
file shows how you can add comments to your own input
file.

Because printing is different from typewriting,
there are a number of things that you have to do
differently when preparing an input file than if
you were just typing the document directly.
Quotation marks like
       ``this''
have to be handled specially, as do quotes within
quotes:
       ``\,`this'            % \, separates the double and single quote.
        is what I just
        wrote, not  `that'\,''.

Dashes come in three sizes: an
       intra-word
dash, a medium dash for number ranges like
       1--2,
and a punctuation
       dash---like
this.

A sentence-ending space should be larger than the
space between words within a sentence.  You
sometimes have to type special commands in
conjunction with punctuation characters to get
this right, as in the following sentence.
       Gnats, gnus, etc.\ all  % `\ ' makes an inter-word space.
       begin with G\@.         % \@ marks end-of-sentence punctuation.
You should check the spaces after periods when
reading your output to make sure you haven't
forgotten any special cases.  Generating an
ellipsis
       \ldots\               % `\ ' is needed after `\ldots' because TeX
                             % ignores spaces after command names like \ldots
                             % made from \ + letters.
                             %
                             % Note how a `%' character causes TeX to ignore
                             % the end of the input line, so these blank lines
                             % do not start a new paragraph.
                             %
with the right spacing around the periods requires
a special command.

\LaTeX\ interprets some common characters as
commands, so you must type special commands to
generate them.  These characters include the
following:
       \$ \& \% \# \{ and \}.

In printing, text is usually emphasized with an
       \emph{italic}
type style.

\begin{em}
   A long segment of text can also be emphasized
   in this way.  Text within such a segment can be
   given \emph{additional} emphasis.
\end{em}

It is sometimes necessary to prevent \LaTeX\ from
breaking a line where it might otherwise do so.
This may be at a space, as between the ``Mr.''\ and
``Jones'' in
       ``Mr.~Jones'',        % ~ produces an unbreakable interword space.
or within a word---especially when the word is a
symbol like
       \mbox{\emph{itemnum}}
that makes little sense when hyphenated across
lines.

Footnotes\footnote{This is an example of a footnote.}
pose no problem.

\LaTeX\ is good at typesetting mathematical formulas
like
       \( x-3y + z = 7 \)
or
       \( a_{1} > x^{2n} + y^{2n} > x' \)
or
       \( \ip{A}{B} = \sum_{i} a_{i} b_{i} \).
The spaces you type in a formula are
ignored.  Remember that a letter like
       $x$                   % $ ... $  and  \( ... \)  are equivalent
is a formula when it denotes a mathematical
symbol, and it should be typed as one.

\section{Displayed Text}

Text is displayed by indenting it from the left
margin.  Quotations are commonly displayed.  There
are short quotations
\begin{quote}
   This is a short quotation.  It consists of a
   single paragraph of text.  See how it is formatted.
\end{quote}
and longer ones.
\begin{quotation}
   This is a longer quotation.  It consists of two
   paragraphs of text, neither of which are
   particularly interesting.

   This is the second paragraph of the quotation.  It
   is just as dull as the first paragraph.
\end{quotation}
Another frequently-displayed structure is a list.
The following is an example of an \emph{itemized}
list.
\begin{itemize}
   \item This is the first item of an itemized list.
         Each item in the list is marked with a ``tick''.
         You don't have to worry about what kind of tick
         mark is used.

   \item This is the second item of the list.  It
         contains another list nested inside it.  The inner
         list is an \emph{enumerated} list.
         \begin{enumerate}
            \item This is the first item of an enumerated
                  list that is nested within the itemized list.

            \item This is the second item of the inner list.
                  \LaTeX\ allows you to nest lists deeper than
                  you really should.
         \end{enumerate}
         This is the rest of the second item of the outer
         list.  It is no more interesting than any other
         part of the item.
   \item This is the third item of the list.
\end{itemize}
You can even display poetry.
\begin{verse}
   There is an environment
    for verse \\             % The \\ command separates lines
   Whose features some poets % within a stanza.
   will curse.

                             % One or more blank lines separate stanzas.

   For instead of making\\
   Them do \emph{all} line breaking, \\
   It allows them to put too many words on a line when they'd rather be
   forced to be terse.
\end{verse}

Mathematical formulas may also be displayed.  A
displayed formula
is
one-line long; multiline
formulas require special formatting instructions.
   \[  \ip{\Gamma}{\psi'} = x'' + y^{2} + z_{i}^{n}\]
Don't start a paragraph with a displayed equation,
nor make one a paragraph by itself.

\end{document}               % End of document.

\section{Bench Huge Section 12}
% synthetic bench section 12
% This is a sample LaTeX input file.  (Version of 12 August 2004.)
%
% A '%' character causes TeX to ignore all remaining text on the line,
% and is used for comments like this one.

\documentclass{article}      % Specifies the document class

                             % The preamble begins here.
\title{An Example Document}  % Declares the document's title.
\author{Leslie Lamport}      % Declares the author's name.
\date{January 21, 1994}      % Deleting this command produces today's date.

\newcommand{\ip}[2]{(#1, #2)}
                             % Defines \ip{arg1}{arg2} to mean
                             % (arg1, arg2).

%\newcommand{\ip}[2]{\langle #1 | #2\rangle}
                             % This is an alternative definition of
                             % \ip that is commented out.

\begin{document}             % End of preamble and beginning of text.

\maketitle                   % Produces the title.

This is an example input file.  Comparing it with
the output it generates can show you how to
produce a simple document of your own.

\section{Ordinary Text}      % Produces section heading.  Lower-level
                             % sections are begun with similar
                             % \subsection and \subsubsection commands.

The ends  of words and sentences are marked
  by   spaces. It  doesn't matter how many
spaces    you type; one is as good as 100.  The
end of   a line counts as a space.

One   or more   blank lines denote the  end
of  a paragraph.

Since any number of consecutive spaces are treated
like a single one, the formatting of the input
file makes no difference to
      \LaTeX,                % The \LaTeX command generates the LaTeX logo.
but it makes a difference to you.  When you use
\LaTeX, making your input file as easy to read
as possible will be a great help as you write
your document and when you change it.  This sample
file shows how you can add comments to your own input
file.

Because printing is different from typewriting,
there are a number of things that you have to do
differently when preparing an input file than if
you were just typing the document directly.
Quotation marks like
       ``this''
have to be handled specially, as do quotes within
quotes:
       ``\,`this'            % \, separates the double and single quote.
        is what I just
        wrote, not  `that'\,''.

Dashes come in three sizes: an
       intra-word
dash, a medium dash for number ranges like
       1--2,
and a punctuation
       dash---like
this.

A sentence-ending space should be larger than the
space between words within a sentence.  You
sometimes have to type special commands in
conjunction with punctuation characters to get
this right, as in the following sentence.
       Gnats, gnus, etc.\ all  % `\ ' makes an inter-word space.
       begin with G\@.         % \@ marks end-of-sentence punctuation.
You should check the spaces after periods when
reading your output to make sure you haven't
forgotten any special cases.  Generating an
ellipsis
       \ldots\               % `\ ' is needed after `\ldots' because TeX
                             % ignores spaces after command names like \ldots
                             % made from \ + letters.
                             %
                             % Note how a `%' character causes TeX to ignore
                             % the end of the input line, so these blank lines
                             % do not start a new paragraph.
                             %
with the right spacing around the periods requires
a special command.

\LaTeX\ interprets some common characters as
commands, so you must type special commands to
generate them.  These characters include the
following:
       \$ \& \% \# \{ and \}.

In printing, text is usually emphasized with an
       \emph{italic}
type style.

\begin{em}
   A long segment of text can also be emphasized
   in this way.  Text within such a segment can be
   given \emph{additional} emphasis.
\end{em}

It is sometimes necessary to prevent \LaTeX\ from
breaking a line where it might otherwise do so.
This may be at a space, as between the ``Mr.''\ and
``Jones'' in
       ``Mr.~Jones'',        % ~ produces an unbreakable interword space.
or within a word---especially when the word is a
symbol like
       \mbox{\emph{itemnum}}
that makes little sense when hyphenated across
lines.

Footnotes\footnote{This is an example of a footnote.}
pose no problem.

\LaTeX\ is good at typesetting mathematical formulas
like
       \( x-3y + z = 7 \)
or
       \( a_{1} > x^{2n} + y^{2n} > x' \)
or
       \( \ip{A}{B} = \sum_{i} a_{i} b_{i} \).
The spaces you type in a formula are
ignored.  Remember that a letter like
       $x$                   % $ ... $  and  \( ... \)  are equivalent
is a formula when it denotes a mathematical
symbol, and it should be typed as one.

\section{Displayed Text}

Text is displayed by indenting it from the left
margin.  Quotations are commonly displayed.  There
are short quotations
\begin{quote}
   This is a short quotation.  It consists of a
   single paragraph of text.  See how it is formatted.
\end{quote}
and longer ones.
\begin{quotation}
   This is a longer quotation.  It consists of two
   paragraphs of text, neither of which are
   particularly interesting.

   This is the second paragraph of the quotation.  It
   is just as dull as the first paragraph.
\end{quotation}
Another frequently-displayed structure is a list.
The following is an example of an \emph{itemized}
list.
\begin{itemize}
   \item This is the first item of an itemized list.
         Each item in the list is marked with a ``tick''.
         You don't have to worry about what kind of tick
         mark is used.

   \item This is the second item of the list.  It
         contains another list nested inside it.  The inner
         list is an \emph{enumerated} list.
         \begin{enumerate}
            \item This is the first item of an enumerated
                  list that is nested within the itemized list.

            \item This is the second item of the inner list.
                  \LaTeX\ allows you to nest lists deeper than
                  you really should.
         \end{enumerate}
         This is the rest of the second item of the outer
         list.  It is no more interesting than any other
         part of the item.
   \item This is the third item of the list.
\end{itemize}
You can even display poetry.
\begin{verse}
   There is an environment
    for verse \\             % The \\ command separates lines
   Whose features some poets % within a stanza.
   will curse.

                             % One or more blank lines separate stanzas.

   For instead of making\\
   Them do \emph{all} line breaking, \\
   It allows them to put too many words on a line when they'd rather be
   forced to be terse.
\end{verse}

Mathematical formulas may also be displayed.  A
displayed formula
is
one-line long; multiline
formulas require special formatting instructions.
   \[  \ip{\Gamma}{\psi'} = x'' + y^{2} + z_{i}^{n}\]
Don't start a paragraph with a displayed equation,
nor make one a paragraph by itself.

\end{document}               % End of document.

\section{Bench Huge Section 13}
% synthetic bench section 13
% This is a sample LaTeX input file.  (Version of 12 August 2004.)
%
% A '%' character causes TeX to ignore all remaining text on the line,
% and is used for comments like this one.

\documentclass{article}      % Specifies the document class

                             % The preamble begins here.
\title{An Example Document}  % Declares the document's title.
\author{Leslie Lamport}      % Declares the author's name.
\date{January 21, 1994}      % Deleting this command produces today's date.

\newcommand{\ip}[2]{(#1, #2)}
                             % Defines \ip{arg1}{arg2} to mean
                             % (arg1, arg2).

%\newcommand{\ip}[2]{\langle #1 | #2\rangle}
                             % This is an alternative definition of
                             % \ip that is commented out.

\begin{document}             % End of preamble and beginning of text.

\maketitle                   % Produces the title.

This is an example input file.  Comparing it with
the output it generates can show you how to
produce a simple document of your own.

\section{Ordinary Text}      % Produces section heading.  Lower-level
                             % sections are begun with similar
                             % \subsection and \subsubsection commands.

The ends  of words and sentences are marked
  by   spaces. It  doesn't matter how many
spaces    you type; one is as good as 100.  The
end of   a line counts as a space.

One   or more   blank lines denote the  end
of  a paragraph.

Since any number of consecutive spaces are treated
like a single one, the formatting of the input
file makes no difference to
      \LaTeX,                % The \LaTeX command generates the LaTeX logo.
but it makes a difference to you.  When you use
\LaTeX, making your input file as easy to read
as possible will be a great help as you write
your document and when you change it.  This sample
file shows how you can add comments to your own input
file.

Because printing is different from typewriting,
there are a number of things that you have to do
differently when preparing an input file than if
you were just typing the document directly.
Quotation marks like
       ``this''
have to be handled specially, as do quotes within
quotes:
       ``\,`this'            % \, separates the double and single quote.
        is what I just
        wrote, not  `that'\,''.

Dashes come in three sizes: an
       intra-word
dash, a medium dash for number ranges like
       1--2,
and a punctuation
       dash---like
this.

A sentence-ending space should be larger than the
space between words within a sentence.  You
sometimes have to type special commands in
conjunction with punctuation characters to get
this right, as in the following sentence.
       Gnats, gnus, etc.\ all  % `\ ' makes an inter-word space.
       begin with G\@.         % \@ marks end-of-sentence punctuation.
You should check the spaces after periods when
reading your output to make sure you haven't
forgotten any special cases.  Generating an
ellipsis
       \ldots\               % `\ ' is needed after `\ldots' because TeX
                             % ignores spaces after command names like \ldots
                             % made from \ + letters.
                             %
                             % Note how a `%' character causes TeX to ignore
                             % the end of the input line, so these blank lines
                             % do not start a new paragraph.
                             %
with the right spacing around the periods requires
a special command.

\LaTeX\ interprets some common characters as
commands, so you must type special commands to
generate them.  These characters include the
following:
       \$ \& \% \# \{ and \}.

In printing, text is usually emphasized with an
       \emph{italic}
type style.

\begin{em}
   A long segment of text can also be emphasized
   in this way.  Text within such a segment can be
   given \emph{additional} emphasis.
\end{em}

It is sometimes necessary to prevent \LaTeX\ from
breaking a line where it might otherwise do so.
This may be at a space, as between the ``Mr.''\ and
``Jones'' in
       ``Mr.~Jones'',        % ~ produces an unbreakable interword space.
or within a word---especially when the word is a
symbol like
       \mbox{\emph{itemnum}}
that makes little sense when hyphenated across
lines.

Footnotes\footnote{This is an example of a footnote.}
pose no problem.

\LaTeX\ is good at typesetting mathematical formulas
like
       \( x-3y + z = 7 \)
or
       \( a_{1} > x^{2n} + y^{2n} > x' \)
or
       \( \ip{A}{B} = \sum_{i} a_{i} b_{i} \).
The spaces you type in a formula are
ignored.  Remember that a letter like
       $x$                   % $ ... $  and  \( ... \)  are equivalent
is a formula when it denotes a mathematical
symbol, and it should be typed as one.

\section{Displayed Text}

Text is displayed by indenting it from the left
margin.  Quotations are commonly displayed.  There
are short quotations
\begin{quote}
   This is a short quotation.  It consists of a
   single paragraph of text.  See how it is formatted.
\end{quote}
and longer ones.
\begin{quotation}
   This is a longer quotation.  It consists of two
   paragraphs of text, neither of which are
   particularly interesting.

   This is the second paragraph of the quotation.  It
   is just as dull as the first paragraph.
\end{quotation}
Another frequently-displayed structure is a list.
The following is an example of an \emph{itemized}
list.
\begin{itemize}
   \item This is the first item of an itemized list.
         Each item in the list is marked with a ``tick''.
         You don't have to worry about what kind of tick
         mark is used.

   \item This is the second item of the list.  It
         contains another list nested inside it.  The inner
         list is an \emph{enumerated} list.
         \begin{enumerate}
            \item This is the first item of an enumerated
                  list that is nested within the itemized list.

            \item This is the second item of the inner list.
                  \LaTeX\ allows you to nest lists deeper than
                  you really should.
         \end{enumerate}
         This is the rest of the second item of the outer
         list.  It is no more interesting than any other
         part of the item.
   \item This is the third item of the list.
\end{itemize}
You can even display poetry.
\begin{verse}
   There is an environment
    for verse \\             % The \\ command separates lines
   Whose features some poets % within a stanza.
   will curse.

                             % One or more blank lines separate stanzas.

   For instead of making\\
   Them do \emph{all} line breaking, \\
   It allows them to put too many words on a line when they'd rather be
   forced to be terse.
\end{verse}

Mathematical formulas may also be displayed.  A
displayed formula
is
one-line long; multiline
formulas require special formatting instructions.
   \[  \ip{\Gamma}{\psi'} = x'' + y^{2} + z_{i}^{n}\]
Don't start a paragraph with a displayed equation,
nor make one a paragraph by itself.

\end{document}               % End of document.

\section{Bench Huge Section 14}
% synthetic bench section 14
% This is a sample LaTeX input file.  (Version of 12 August 2004.)
%
% A '%' character causes TeX to ignore all remaining text on the line,
% and is used for comments like this one.

\documentclass{article}      % Specifies the document class

                             % The preamble begins here.
\title{An Example Document}  % Declares the document's title.
\author{Leslie Lamport}      % Declares the author's name.
\date{January 21, 1994}      % Deleting this command produces today's date.

\newcommand{\ip}[2]{(#1, #2)}
                             % Defines \ip{arg1}{arg2} to mean
                             % (arg1, arg2).

%\newcommand{\ip}[2]{\langle #1 | #2\rangle}
                             % This is an alternative definition of
                             % \ip that is commented out.

\begin{document}             % End of preamble and beginning of text.

\maketitle                   % Produces the title.

This is an example input file.  Comparing it with
the output it generates can show you how to
produce a simple document of your own.

\section{Ordinary Text}      % Produces section heading.  Lower-level
                             % sections are begun with similar
                             % \subsection and \subsubsection commands.

The ends  of words and sentences are marked
  by   spaces. It  doesn't matter how many
spaces    you type; one is as good as 100.  The
end of   a line counts as a space.

One   or more   blank lines denote the  end
of  a paragraph.

Since any number of consecutive spaces are treated
like a single one, the formatting of the input
file makes no difference to
      \LaTeX,                % The \LaTeX command generates the LaTeX logo.
but it makes a difference to you.  When you use
\LaTeX, making your input file as easy to read
as possible will be a great help as you write
your document and when you change it.  This sample
file shows how you can add comments to your own input
file.

Because printing is different from typewriting,
there are a number of things that you have to do
differently when preparing an input file than if
you were just typing the document directly.
Quotation marks like
       ``this''
have to be handled specially, as do quotes within
quotes:
       ``\,`this'            % \, separates the double and single quote.
        is what I just
        wrote, not  `that'\,''.

Dashes come in three sizes: an
       intra-word
dash, a medium dash for number ranges like
       1--2,
and a punctuation
       dash---like
this.

A sentence-ending space should be larger than the
space between words within a sentence.  You
sometimes have to type special commands in
conjunction with punctuation characters to get
this right, as in the following sentence.
       Gnats, gnus, etc.\ all  % `\ ' makes an inter-word space.
       begin with G\@.         % \@ marks end-of-sentence punctuation.
You should check the spaces after periods when
reading your output to make sure you haven't
forgotten any special cases.  Generating an
ellipsis
       \ldots\               % `\ ' is needed after `\ldots' because TeX
                             % ignores spaces after command names like \ldots
                             % made from \ + letters.
                             %
                             % Note how a `%' character causes TeX to ignore
                             % the end of the input line, so these blank lines
                             % do not start a new paragraph.
                             %
with the right spacing around the periods requires
a special command.

\LaTeX\ interprets some common characters as
commands, so you must type special commands to
generate them.  These characters include the
following:
       \$ \& \% \# \{ and \}.

In printing, text is usually emphasized with an
       \emph{italic}
type style.

\begin{em}
   A long segment of text can also be emphasized
   in this way.  Text within such a segment can be
   given \emph{additional} emphasis.
\end{em}

It is sometimes necessary to prevent \LaTeX\ from
breaking a line where it might otherwise do so.
This may be at a space, as between the ``Mr.''\ and
``Jones'' in
       ``Mr.~Jones'',        % ~ produces an unbreakable interword space.
or within a word---especially when the word is a
symbol like
       \mbox{\emph{itemnum}}
that makes little sense when hyphenated across
lines.

Footnotes\footnote{This is an example of a footnote.}
pose no problem.

\LaTeX\ is good at typesetting mathematical formulas
like
       \( x-3y + z = 7 \)
or
       \( a_{1} > x^{2n} + y^{2n} > x' \)
or
       \( \ip{A}{B} = \sum_{i} a_{i} b_{i} \).
The spaces you type in a formula are
ignored.  Remember that a letter like
       $x$                   % $ ... $  and  \( ... \)  are equivalent
is a formula when it denotes a mathematical
symbol, and it should be typed as one.

\section{Displayed Text}

Text is displayed by indenting it from the left
margin.  Quotations are commonly displayed.  There
are short quotations
\begin{quote}
   This is a short quotation.  It consists of a
   single paragraph of text.  See how it is formatted.
\end{quote}
and longer ones.
\begin{quotation}
   This is a longer quotation.  It consists of two
   paragraphs of text, neither of which are
   particularly interesting.

   This is the second paragraph of the quotation.  It
   is just as dull as the first paragraph.
\end{quotation}
Another frequently-displayed structure is a list.
The following is an example of an \emph{itemized}
list.
\begin{itemize}
   \item This is the first item of an itemized list.
         Each item in the list is marked with a ``tick''.
         You don't have to worry about what kind of tick
         mark is used.

   \item This is the second item of the list.  It
         contains another list nested inside it.  The inner
         list is an \emph{enumerated} list.
         \begin{enumerate}
            \item This is the first item of an enumerated
                  list that is nested within the itemized list.

            \item This is the second item of the inner list.
                  \LaTeX\ allows you to nest lists deeper than
                  you really should.
         \end{enumerate}
         This is the rest of the second item of the outer
         list.  It is no more interesting than any other
         part of the item.
   \item This is the third item of the list.
\end{itemize}
You can even display poetry.
\begin{verse}
   There is an environment
    for verse \\             % The \\ command separates lines
   Whose features some poets % within a stanza.
   will curse.

                             % One or more blank lines separate stanzas.

   For instead of making\\
   Them do \emph{all} line breaking, \\
   It allows them to put too many words on a line when they'd rather be
   forced to be terse.
\end{verse}

Mathematical formulas may also be displayed.  A
displayed formula
is
one-line long; multiline
formulas require special formatting instructions.
   \[  \ip{\Gamma}{\psi'} = x'' + y^{2} + z_{i}^{n}\]
Don't start a paragraph with a displayed equation,
nor make one a paragraph by itself.

\end{document}               % End of document.

\section{Bench Huge Section 15}
% synthetic bench section 15
% This is a sample LaTeX input file.  (Version of 12 August 2004.)
%
% A '%' character causes TeX to ignore all remaining text on the line,
% and is used for comments like this one.

\documentclass{article}      % Specifies the document class

                             % The preamble begins here.
\title{An Example Document}  % Declares the document's title.
\author{Leslie Lamport}      % Declares the author's name.
\date{January 21, 1994}      % Deleting this command produces today's date.

\newcommand{\ip}[2]{(#1, #2)}
                             % Defines \ip{arg1}{arg2} to mean
                             % (arg1, arg2).

%\newcommand{\ip}[2]{\langle #1 | #2\rangle}
                             % This is an alternative definition of
                             % \ip that is commented out.

\begin{document}             % End of preamble and beginning of text.

\maketitle                   % Produces the title.

This is an example input file.  Comparing it with
the output it generates can show you how to
produce a simple document of your own.

\section{Ordinary Text}      % Produces section heading.  Lower-level
                             % sections are begun with similar
                             % \subsection and \subsubsection commands.

The ends  of words and sentences are marked
  by   spaces. It  doesn't matter how many
spaces    you type; one is as good as 100.  The
end of   a line counts as a space.

One   or more   blank lines denote the  end
of  a paragraph.

Since any number of consecutive spaces are treated
like a single one, the formatting of the input
file makes no difference to
      \LaTeX,                % The \LaTeX command generates the LaTeX logo.
but it makes a difference to you.  When you use
\LaTeX, making your input file as easy to read
as possible will be a great help as you write
your document and when you change it.  This sample
file shows how you can add comments to your own input
file.

Because printing is different from typewriting,
there are a number of things that you have to do
differently when preparing an input file than if
you were just typing the document directly.
Quotation marks like
       ``this''
have to be handled specially, as do quotes within
quotes:
       ``\,`this'            % \, separates the double and single quote.
        is what I just
        wrote, not  `that'\,''.

Dashes come in three sizes: an
       intra-word
dash, a medium dash for number ranges like
       1--2,
and a punctuation
       dash---like
this.

A sentence-ending space should be larger than the
space between words within a sentence.  You
sometimes have to type special commands in
conjunction with punctuation characters to get
this right, as in the following sentence.
       Gnats, gnus, etc.\ all  % `\ ' makes an inter-word space.
       begin with G\@.         % \@ marks end-of-sentence punctuation.
You should check the spaces after periods when
reading your output to make sure you haven't
forgotten any special cases.  Generating an
ellipsis
       \ldots\               % `\ ' is needed after `\ldots' because TeX
                             % ignores spaces after command names like \ldots
                             % made from \ + letters.
                             %
                             % Note how a `%' character causes TeX to ignore
                             % the end of the input line, so these blank lines
                             % do not start a new paragraph.
                             %
with the right spacing around the periods requires
a special command.

\LaTeX\ interprets some common characters as
commands, so you must type special commands to
generate them.  These characters include the
following:
       \$ \& \% \# \{ and \}.

In printing, text is usually emphasized with an
       \emph{italic}
type style.

\begin{em}
   A long segment of text can also be emphasized
   in this way.  Text within such a segment can be
   given \emph{additional} emphasis.
\end{em}

It is sometimes necessary to prevent \LaTeX\ from
breaking a line where it might otherwise do so.
This may be at a space, as between the ``Mr.''\ and
``Jones'' in
       ``Mr.~Jones'',        % ~ produces an unbreakable interword space.
or within a word---especially when the word is a
symbol like
       \mbox{\emph{itemnum}}
that makes little sense when hyphenated across
lines.

Footnotes\footnote{This is an example of a footnote.}
pose no problem.

\LaTeX\ is good at typesetting mathematical formulas
like
       \( x-3y + z = 7 \)
or
       \( a_{1} > x^{2n} + y^{2n} > x' \)
or
       \( \ip{A}{B} = \sum_{i} a_{i} b_{i} \).
The spaces you type in a formula are
ignored.  Remember that a letter like
       $x$                   % $ ... $  and  \( ... \)  are equivalent
is a formula when it denotes a mathematical
symbol, and it should be typed as one.

\section{Displayed Text}

Text is displayed by indenting it from the left
margin.  Quotations are commonly displayed.  There
are short quotations
\begin{quote}
   This is a short quotation.  It consists of a
   single paragraph of text.  See how it is formatted.
\end{quote}
and longer ones.
\begin{quotation}
   This is a longer quotation.  It consists of two
   paragraphs of text, neither of which are
   particularly interesting.

   This is the second paragraph of the quotation.  It
   is just as dull as the first paragraph.
\end{quotation}
Another frequently-displayed structure is a list.
The following is an example of an \emph{itemized}
list.
\begin{itemize}
   \item This is the first item of an itemized list.
         Each item in the list is marked with a ``tick''.
         You don't have to worry about what kind of tick
         mark is used.

   \item This is the second item of the list.  It
         contains another list nested inside it.  The inner
         list is an \emph{enumerated} list.
         \begin{enumerate}
            \item This is the first item of an enumerated
                  list that is nested within the itemized list.

            \item This is the second item of the inner list.
                  \LaTeX\ allows you to nest lists deeper than
                  you really should.
         \end{enumerate}
         This is the rest of the second item of the outer
         list.  It is no more interesting than any other
         part of the item.
   \item This is the third item of the list.
\end{itemize}
You can even display poetry.
\begin{verse}
   There is an environment
    for verse \\             % The \\ command separates lines
   Whose features some poets % within a stanza.
   will curse.

                             % One or more blank lines separate stanzas.

   For instead of making\\
   Them do \emph{all} line breaking, \\
   It allows them to put too many words on a line when they'd rather be
   forced to be terse.
\end{verse}

Mathematical formulas may also be displayed.  A
displayed formula
is
one-line long; multiline
formulas require special formatting instructions.
   \[  \ip{\Gamma}{\psi'} = x'' + y^{2} + z_{i}^{n}\]
Don't start a paragraph with a displayed equation,
nor make one a paragraph by itself.

\end{document}               % End of document.

\section{Bench Huge Section 16}
% synthetic bench section 16
% This is a sample LaTeX input file.  (Version of 12 August 2004.)
%
% A '%' character causes TeX to ignore all remaining text on the line,
% and is used for comments like this one.

\documentclass{article}      % Specifies the document class

                             % The preamble begins here.
\title{An Example Document}  % Declares the document's title.
\author{Leslie Lamport}      % Declares the author's name.
\date{January 21, 1994}      % Deleting this command produces today's date.

\newcommand{\ip}[2]{(#1, #2)}
                             % Defines \ip{arg1}{arg2} to mean
                             % (arg1, arg2).

%\newcommand{\ip}[2]{\langle #1 | #2\rangle}
                             % This is an alternative definition of
                             % \ip that is commented out.

\begin{document}             % End of preamble and beginning of text.

\maketitle                   % Produces the title.

This is an example input file.  Comparing it with
the output it generates can show you how to
produce a simple document of your own.

\section{Ordinary Text}      % Produces section heading.  Lower-level
                             % sections are begun with similar
                             % \subsection and \subsubsection commands.

The ends  of words and sentences are marked
  by   spaces. It  doesn't matter how many
spaces    you type; one is as good as 100.  The
end of   a line counts as a space.

One   or more   blank lines denote the  end
of  a paragraph.

Since any number of consecutive spaces are treated
like a single one, the formatting of the input
file makes no difference to
      \LaTeX,                % The \LaTeX command generates the LaTeX logo.
but it makes a difference to you.  When you use
\LaTeX, making your input file as easy to read
as possible will be a great help as you write
your document and when you change it.  This sample
file shows how you can add comments to your own input
file.

Because printing is different from typewriting,
there are a number of things that you have to do
differently when preparing an input file than if
you were just typing the document directly.
Quotation marks like
       ``this''
have to be handled specially, as do quotes within
quotes:
       ``\,`this'            % \, separates the double and single quote.
        is what I just
        wrote, not  `that'\,''.

Dashes come in three sizes: an
       intra-word
dash, a medium dash for number ranges like
       1--2,
and a punctuation
       dash---like
this.

A sentence-ending space should be larger than the
space between words within a sentence.  You
sometimes have to type special commands in
conjunction with punctuation characters to get
this right, as in the following sentence.
       Gnats, gnus, etc.\ all  % `\ ' makes an inter-word space.
       begin with G\@.         % \@ marks end-of-sentence punctuation.
You should check the spaces after periods when
reading your output to make sure you haven't
forgotten any special cases.  Generating an
ellipsis
       \ldots\               % `\ ' is needed after `\ldots' because TeX
                             % ignores spaces after command names like \ldots
                             % made from \ + letters.
                             %
                             % Note how a `%' character causes TeX to ignore
                             % the end of the input line, so these blank lines
                             % do not start a new paragraph.
                             %
with the right spacing around the periods requires
a special command.

\LaTeX\ interprets some common characters as
commands, so you must type special commands to
generate them.  These characters include the
following:
       \$ \& \% \# \{ and \}.

In printing, text is usually emphasized with an
       \emph{italic}
type style.

\begin{em}
   A long segment of text can also be emphasized
   in this way.  Text within such a segment can be
   given \emph{additional} emphasis.
\end{em}

It is sometimes necessary to prevent \LaTeX\ from
breaking a line where it might otherwise do so.
This may be at a space, as between the ``Mr.''\ and
``Jones'' in
       ``Mr.~Jones'',        % ~ produces an unbreakable interword space.
or within a word---especially when the word is a
symbol like
       \mbox{\emph{itemnum}}
that makes little sense when hyphenated across
lines.

Footnotes\footnote{This is an example of a footnote.}
pose no problem.

\LaTeX\ is good at typesetting mathematical formulas
like
       \( x-3y + z = 7 \)
or
       \( a_{1} > x^{2n} + y^{2n} > x' \)
or
       \( \ip{A}{B} = \sum_{i} a_{i} b_{i} \).
The spaces you type in a formula are
ignored.  Remember that a letter like
       $x$                   % $ ... $  and  \( ... \)  are equivalent
is a formula when it denotes a mathematical
symbol, and it should be typed as one.

\section{Displayed Text}

Text is displayed by indenting it from the left
margin.  Quotations are commonly displayed.  There
are short quotations
\begin{quote}
   This is a short quotation.  It consists of a
   single paragraph of text.  See how it is formatted.
\end{quote}
and longer ones.
\begin{quotation}
   This is a longer quotation.  It consists of two
   paragraphs of text, neither of which are
   particularly interesting.

   This is the second paragraph of the quotation.  It
   is just as dull as the first paragraph.
\end{quotation}
Another frequently-displayed structure is a list.
The following is an example of an \emph{itemized}
list.
\begin{itemize}
   \item This is the first item of an itemized list.
         Each item in the list is marked with a ``tick''.
         You don't have to worry about what kind of tick
         mark is used.

   \item This is the second item of the list.  It
         contains another list nested inside it.  The inner
         list is an \emph{enumerated} list.
         \begin{enumerate}
            \item This is the first item of an enumerated
                  list that is nested within the itemized list.

            \item This is the second item of the inner list.
                  \LaTeX\ allows you to nest lists deeper than
                  you really should.
         \end{enumerate}
         This is the rest of the second item of the outer
         list.  It is no more interesting than any other
         part of the item.
   \item This is the third item of the list.
\end{itemize}
You can even display poetry.
\begin{verse}
   There is an environment
    for verse \\             % The \\ command separates lines
   Whose features some poets % within a stanza.
   will curse.

                             % One or more blank lines separate stanzas.

   For instead of making\\
   Them do \emph{all} line breaking, \\
   It allows them to put too many words on a line when they'd rather be
   forced to be terse.
\end{verse}

Mathematical formulas may also be displayed.  A
displayed formula
is
one-line long; multiline
formulas require special formatting instructions.
   \[  \ip{\Gamma}{\psi'} = x'' + y^{2} + z_{i}^{n}\]
Don't start a paragraph with a displayed equation,
nor make one a paragraph by itself.

\end{document}               % End of document.

\section{Bench Huge Section 17}
% synthetic bench section 17
% This is a sample LaTeX input file.  (Version of 12 August 2004.)
%
% A '%' character causes TeX to ignore all remaining text on the line,
% and is used for comments like this one.

\documentclass{article}      % Specifies the document class

                             % The preamble begins here.
\title{An Example Document}  % Declares the document's title.
\author{Leslie Lamport}      % Declares the author's name.
\date{January 21, 1994}      % Deleting this command produces today's date.

\newcommand{\ip}[2]{(#1, #2)}
                             % Defines \ip{arg1}{arg2} to mean
                             % (arg1, arg2).

%\newcommand{\ip}[2]{\langle #1 | #2\rangle}
                             % This is an alternative definition of
                             % \ip that is commented out.

\begin{document}             % End of preamble and beginning of text.

\maketitle                   % Produces the title.

This is an example input file.  Comparing it with
the output it generates can show you how to
produce a simple document of your own.

\section{Ordinary Text}      % Produces section heading.  Lower-level
                             % sections are begun with similar
                             % \subsection and \subsubsection commands.

The ends  of words and sentences are marked
  by   spaces. It  doesn't matter how many
spaces    you type; one is as good as 100.  The
end of   a line counts as a space.

One   or more   blank lines denote the  end
of  a paragraph.

Since any number of consecutive spaces are treated
like a single one, the formatting of the input
file makes no difference to
      \LaTeX,                % The \LaTeX command generates the LaTeX logo.
but it makes a difference to you.  When you use
\LaTeX, making your input file as easy to read
as possible will be a great help as you write
your document and when you change it.  This sample
file shows how you can add comments to your own input
file.

Because printing is different from typewriting,
there are a number of things that you have to do
differently when preparing an input file than if
you were just typing the document directly.
Quotation marks like
       ``this''
have to be handled specially, as do quotes within
quotes:
       ``\,`this'            % \, separates the double and single quote.
        is what I just
        wrote, not  `that'\,''.

Dashes come in three sizes: an
       intra-word
dash, a medium dash for number ranges like
       1--2,
and a punctuation
       dash---like
this.

A sentence-ending space should be larger than the
space between words within a sentence.  You
sometimes have to type special commands in
conjunction with punctuation characters to get
this right, as in the following sentence.
       Gnats, gnus, etc.\ all  % `\ ' makes an inter-word space.
       begin with G\@.         % \@ marks end-of-sentence punctuation.
You should check the spaces after periods when
reading your output to make sure you haven't
forgotten any special cases.  Generating an
ellipsis
       \ldots\               % `\ ' is needed after `\ldots' because TeX
                             % ignores spaces after command names like \ldots
                             % made from \ + letters.
                             %
                             % Note how a `%' character causes TeX to ignore
                             % the end of the input line, so these blank lines
                             % do not start a new paragraph.
                             %
with the right spacing around the periods requires
a special command.

\LaTeX\ interprets some common characters as
commands, so you must type special commands to
generate them.  These characters include the
following:
       \$ \& \% \# \{ and \}.

In printing, text is usually emphasized with an
       \emph{italic}
type style.

\begin{em}
   A long segment of text can also be emphasized
   in this way.  Text within such a segment can be
   given \emph{additional} emphasis.
\end{em}

It is sometimes necessary to prevent \LaTeX\ from
breaking a line where it might otherwise do so.
This may be at a space, as between the ``Mr.''\ and
``Jones'' in
       ``Mr.~Jones'',        % ~ produces an unbreakable interword space.
or within a word---especially when the word is a
symbol like
       \mbox{\emph{itemnum}}
that makes little sense when hyphenated across
lines.

Footnotes\footnote{This is an example of a footnote.}
pose no problem.

\LaTeX\ is good at typesetting mathematical formulas
like
       \( x-3y + z = 7 \)
or
       \( a_{1} > x^{2n} + y^{2n} > x' \)
or
       \( \ip{A}{B} = \sum_{i} a_{i} b_{i} \).
The spaces you type in a formula are
ignored.  Remember that a letter like
       $x$                   % $ ... $  and  \( ... \)  are equivalent
is a formula when it denotes a mathematical
symbol, and it should be typed as one.

\section{Displayed Text}

Text is displayed by indenting it from the left
margin.  Quotations are commonly displayed.  There
are short quotations
\begin{quote}
   This is a short quotation.  It consists of a
   single paragraph of text.  See how it is formatted.
\end{quote}
and longer ones.
\begin{quotation}
   This is a longer quotation.  It consists of two
   paragraphs of text, neither of which are
   particularly interesting.

   This is the second paragraph of the quotation.  It
   is just as dull as the first paragraph.
\end{quotation}
Another frequently-displayed structure is a list.
The following is an example of an \emph{itemized}
list.
\begin{itemize}
   \item This is the first item of an itemized list.
         Each item in the list is marked with a ``tick''.
         You don't have to worry about what kind of tick
         mark is used.

   \item This is the second item of the list.  It
         contains another list nested inside it.  The inner
         list is an \emph{enumerated} list.
         \begin{enumerate}
            \item This is the first item of an enumerated
                  list that is nested within the itemized list.

            \item This is the second item of the inner list.
                  \LaTeX\ allows you to nest lists deeper than
                  you really should.
         \end{enumerate}
         This is the rest of the second item of the outer
         list.  It is no more interesting than any other
         part of the item.
   \item This is the third item of the list.
\end{itemize}
You can even display poetry.
\begin{verse}
   There is an environment
    for verse \\             % The \\ command separates lines
   Whose features some poets % within a stanza.
   will curse.

                             % One or more blank lines separate stanzas.

   For instead of making\\
   Them do \emph{all} line breaking, \\
   It allows them to put too many words on a line when they'd rather be
   forced to be terse.
\end{verse}

Mathematical formulas may also be displayed.  A
displayed formula
is
one-line long; multiline
formulas require special formatting instructions.
   \[  \ip{\Gamma}{\psi'} = x'' + y^{2} + z_{i}^{n}\]
Don't start a paragraph with a displayed equation,
nor make one a paragraph by itself.

\end{document}               % End of document.

\section{Bench Huge Section 18}
% synthetic bench section 18
% This is a sample LaTeX input file.  (Version of 12 August 2004.)
%
% A '%' character causes TeX to ignore all remaining text on the line,
% and is used for comments like this one.

\documentclass{article}      % Specifies the document class

                             % The preamble begins here.
\title{An Example Document}  % Declares the document's title.
\author{Leslie Lamport}      % Declares the author's name.
\date{January 21, 1994}      % Deleting this command produces today's date.

\newcommand{\ip}[2]{(#1, #2)}
                             % Defines \ip{arg1}{arg2} to mean
                             % (arg1, arg2).

%\newcommand{\ip}[2]{\langle #1 | #2\rangle}
                             % This is an alternative definition of
                             % \ip that is commented out.

\begin{document}             % End of preamble and beginning of text.

\maketitle                   % Produces the title.

This is an example input file.  Comparing it with
the output it generates can show you how to
produce a simple document of your own.

\section{Ordinary Text}      % Produces section heading.  Lower-level
                             % sections are begun with similar
                             % \subsection and \subsubsection commands.

The ends  of words and sentences are marked
  by   spaces. It  doesn't matter how many
spaces    you type; one is as good as 100.  The
end of   a line counts as a space.

One   or more   blank lines denote the  end
of  a paragraph.

Since any number of consecutive spaces are treated
like a single one, the formatting of the input
file makes no difference to
      \LaTeX,                % The \LaTeX command generates the LaTeX logo.
but it makes a difference to you.  When you use
\LaTeX, making your input file as easy to read
as possible will be a great help as you write
your document and when you change it.  This sample
file shows how you can add comments to your own input
file.

Because printing is different from typewriting,
there are a number of things that you have to do
differently when preparing an input file than if
you were just typing the document directly.
Quotation marks like
       ``this''
have to be handled specially, as do quotes within
quotes:
       ``\,`this'            % \, separates the double and single quote.
        is what I just
        wrote, not  `that'\,''.

Dashes come in three sizes: an
       intra-word
dash, a medium dash for number ranges like
       1--2,
and a punctuation
       dash---like
this.

A sentence-ending space should be larger than the
space between words within a sentence.  You
sometimes have to type special commands in
conjunction with punctuation characters to get
this right, as in the following sentence.
       Gnats, gnus, etc.\ all  % `\ ' makes an inter-word space.
       begin with G\@.         % \@ marks end-of-sentence punctuation.
You should check the spaces after periods when
reading your output to make sure you haven't
forgotten any special cases.  Generating an
ellipsis
       \ldots\               % `\ ' is needed after `\ldots' because TeX
                             % ignores spaces after command names like \ldots
                             % made from \ + letters.
                             %
                             % Note how a `%' character causes TeX to ignore
                             % the end of the input line, so these blank lines
                             % do not start a new paragraph.
                             %
with the right spacing around the periods requires
a special command.

\LaTeX\ interprets some common characters as
commands, so you must type special commands to
generate them.  These characters include the
following:
       \$ \& \% \# \{ and \}.

In printing, text is usually emphasized with an
       \emph{italic}
type style.

\begin{em}
   A long segment of text can also be emphasized
   in this way.  Text within such a segment can be
   given \emph{additional} emphasis.
\end{em}

It is sometimes necessary to prevent \LaTeX\ from
breaking a line where it might otherwise do so.
This may be at a space, as between the ``Mr.''\ and
``Jones'' in
       ``Mr.~Jones'',        % ~ produces an unbreakable interword space.
or within a word---especially when the word is a
symbol like
       \mbox{\emph{itemnum}}
that makes little sense when hyphenated across
lines.

Footnotes\footnote{This is an example of a footnote.}
pose no problem.

\LaTeX\ is good at typesetting mathematical formulas
like
       \( x-3y + z = 7 \)
or
       \( a_{1} > x^{2n} + y^{2n} > x' \)
or
       \( \ip{A}{B} = \sum_{i} a_{i} b_{i} \).
The spaces you type in a formula are
ignored.  Remember that a letter like
       $x$                   % $ ... $  and  \( ... \)  are equivalent
is a formula when it denotes a mathematical
symbol, and it should be typed as one.

\section{Displayed Text}

Text is displayed by indenting it from the left
margin.  Quotations are commonly displayed.  There
are short quotations
\begin{quote}
   This is a short quotation.  It consists of a
   single paragraph of text.  See how it is formatted.
\end{quote}
and longer ones.
\begin{quotation}
   This is a longer quotation.  It consists of two
   paragraphs of text, neither of which are
   particularly interesting.

   This is the second paragraph of the quotation.  It
   is just as dull as the first paragraph.
\end{quotation}
Another frequently-displayed structure is a list.
The following is an example of an \emph{itemized}
list.
\begin{itemize}
   \item This is the first item of an itemized list.
         Each item in the list is marked with a ``tick''.
         You don't have to worry about what kind of tick
         mark is used.

   \item This is the second item of the list.  It
         contains another list nested inside it.  The inner
         list is an \emph{enumerated} list.
         \begin{enumerate}
            \item This is the first item of an enumerated
                  list that is nested within the itemized list.

            \item This is the second item of the inner list.
                  \LaTeX\ allows you to nest lists deeper than
                  you really should.
         \end{enumerate}
         This is the rest of the second item of the outer
         list.  It is no more interesting than any other
         part of the item.
   \item This is the third item of the list.
\end{itemize}
You can even display poetry.
\begin{verse}
   There is an environment
    for verse \\             % The \\ command separates lines
   Whose features some poets % within a stanza.
   will curse.

                             % One or more blank lines separate stanzas.

   For instead of making\\
   Them do \emph{all} line breaking, \\
   It allows them to put too many words on a line when they'd rather be
   forced to be terse.
\end{verse}

Mathematical formulas may also be displayed.  A
displayed formula
is
one-line long; multiline
formulas require special formatting instructions.
   \[  \ip{\Gamma}{\psi'} = x'' + y^{2} + z_{i}^{n}\]
Don't start a paragraph with a displayed equation,
nor make one a paragraph by itself.

\end{document}               % End of document.

\section{Bench Huge Section 19}
% synthetic bench section 19
% This is a sample LaTeX input file.  (Version of 12 August 2004.)
%
% A '%' character causes TeX to ignore all remaining text on the line,
% and is used for comments like this one.

\documentclass{article}      % Specifies the document class

                             % The preamble begins here.
\title{An Example Document}  % Declares the document's title.
\author{Leslie Lamport}      % Declares the author's name.
\date{January 21, 1994}      % Deleting this command produces today's date.

\newcommand{\ip}[2]{(#1, #2)}
                             % Defines \ip{arg1}{arg2} to mean
                             % (arg1, arg2).

%\newcommand{\ip}[2]{\langle #1 | #2\rangle}
                             % This is an alternative definition of
                             % \ip that is commented out.

\begin{document}             % End of preamble and beginning of text.

\maketitle                   % Produces the title.

This is an example input file.  Comparing it with
the output it generates can show you how to
produce a simple document of your own.

\section{Ordinary Text}      % Produces section heading.  Lower-level
                             % sections are begun with similar
                             % \subsection and \subsubsection commands.

The ends  of words and sentences are marked
  by   spaces. It  doesn't matter how many
spaces    you type; one is as good as 100.  The
end of   a line counts as a space.

One   or more   blank lines denote the  end
of  a paragraph.

Since any number of consecutive spaces are treated
like a single one, the formatting of the input
file makes no difference to
      \LaTeX,                % The \LaTeX command generates the LaTeX logo.
but it makes a difference to you.  When you use
\LaTeX, making your input file as easy to read
as possible will be a great help as you write
your document and when you change it.  This sample
file shows how you can add comments to your own input
file.

Because printing is different from typewriting,
there are a number of things that you have to do
differently when preparing an input file than if
you were just typing the document directly.
Quotation marks like
       ``this''
have to be handled specially, as do quotes within
quotes:
       ``\,`this'            % \, separates the double and single quote.
        is what I just
        wrote, not  `that'\,''.

Dashes come in three sizes: an
       intra-word
dash, a medium dash for number ranges like
       1--2,
and a punctuation
       dash---like
this.

A sentence-ending space should be larger than the
space between words within a sentence.  You
sometimes have to type special commands in
conjunction with punctuation characters to get
this right, as in the following sentence.
       Gnats, gnus, etc.\ all  % `\ ' makes an inter-word space.
       begin with G\@.         % \@ marks end-of-sentence punctuation.
You should check the spaces after periods when
reading your output to make sure you haven't
forgotten any special cases.  Generating an
ellipsis
       \ldots\               % `\ ' is needed after `\ldots' because TeX
                             % ignores spaces after command names like \ldots
                             % made from \ + letters.
                             %
                             % Note how a `%' character causes TeX to ignore
                             % the end of the input line, so these blank lines
                             % do not start a new paragraph.
                             %
with the right spacing around the periods requires
a special command.

\LaTeX\ interprets some common characters as
commands, so you must type special commands to
generate them.  These characters include the
following:
       \$ \& \% \# \{ and \}.

In printing, text is usually emphasized with an
       \emph{italic}
type style.

\begin{em}
   A long segment of text can also be emphasized
   in this way.  Text within such a segment can be
   given \emph{additional} emphasis.
\end{em}

It is sometimes necessary to prevent \LaTeX\ from
breaking a line where it might otherwise do so.
This may be at a space, as between the ``Mr.''\ and
``Jones'' in
       ``Mr.~Jones'',        % ~ produces an unbreakable interword space.
or within a word---especially when the word is a
symbol like
       \mbox{\emph{itemnum}}
that makes little sense when hyphenated across
lines.

Footnotes\footnote{This is an example of a footnote.}
pose no problem.

\LaTeX\ is good at typesetting mathematical formulas
like
       \( x-3y + z = 7 \)
or
       \( a_{1} > x^{2n} + y^{2n} > x' \)
or
       \( \ip{A}{B} = \sum_{i} a_{i} b_{i} \).
The spaces you type in a formula are
ignored.  Remember that a letter like
       $x$                   % $ ... $  and  \( ... \)  are equivalent
is a formula when it denotes a mathematical
symbol, and it should be typed as one.

\section{Displayed Text}

Text is displayed by indenting it from the left
margin.  Quotations are commonly displayed.  There
are short quotations
\begin{quote}
   This is a short quotation.  It consists of a
   single paragraph of text.  See how it is formatted.
\end{quote}
and longer ones.
\begin{quotation}
   This is a longer quotation.  It consists of two
   paragraphs of text, neither of which are
   particularly interesting.

   This is the second paragraph of the quotation.  It
   is just as dull as the first paragraph.
\end{quotation}
Another frequently-displayed structure is a list.
The following is an example of an \emph{itemized}
list.
\begin{itemize}
   \item This is the first item of an itemized list.
         Each item in the list is marked with a ``tick''.
         You don't have to worry about what kind of tick
         mark is used.

   \item This is the second item of the list.  It
         contains another list nested inside it.  The inner
         list is an \emph{enumerated} list.
         \begin{enumerate}
            \item This is the first item of an enumerated
                  list that is nested within the itemized list.

            \item This is the second item of the inner list.
                  \LaTeX\ allows you to nest lists deeper than
                  you really should.
         \end{enumerate}
         This is the rest of the second item of the outer
         list.  It is no more interesting than any other
         part of the item.
   \item This is the third item of the list.
\end{itemize}
You can even display poetry.
\begin{verse}
   There is an environment
    for verse \\             % The \\ command separates lines
   Whose features some poets % within a stanza.
   will curse.

                             % One or more blank lines separate stanzas.

   For instead of making\\
   Them do \emph{all} line breaking, \\
   It allows them to put too many words on a line when they'd rather be
   forced to be terse.
\end{verse}

Mathematical formulas may also be displayed.  A
displayed formula
is
one-line long; multiline
formulas require special formatting instructions.
   \[  \ip{\Gamma}{\psi'} = x'' + y^{2} + z_{i}^{n}\]
Don't start a paragraph with a displayed equation,
nor make one a paragraph by itself.

\end{document}               % End of document.

\section{Bench Huge Section 20}
% synthetic bench section 20
% This is a sample LaTeX input file.  (Version of 12 August 2004.)
%
% A '%' character causes TeX to ignore all remaining text on the line,
% and is used for comments like this one.

\documentclass{article}      % Specifies the document class

                             % The preamble begins here.
\title{An Example Document}  % Declares the document's title.
\author{Leslie Lamport}      % Declares the author's name.
\date{January 21, 1994}      % Deleting this command produces today's date.

\newcommand{\ip}[2]{(#1, #2)}
                             % Defines \ip{arg1}{arg2} to mean
                             % (arg1, arg2).

%\newcommand{\ip}[2]{\langle #1 | #2\rangle}
                             % This is an alternative definition of
                             % \ip that is commented out.

\begin{document}             % End of preamble and beginning of text.

\maketitle                   % Produces the title.

This is an example input file.  Comparing it with
the output it generates can show you how to
produce a simple document of your own.

\section{Ordinary Text}      % Produces section heading.  Lower-level
                             % sections are begun with similar
                             % \subsection and \subsubsection commands.

The ends  of words and sentences are marked
  by   spaces. It  doesn't matter how many
spaces    you type; one is as good as 100.  The
end of   a line counts as a space.

One   or more   blank lines denote the  end
of  a paragraph.

Since any number of consecutive spaces are treated
like a single one, the formatting of the input
file makes no difference to
      \LaTeX,                % The \LaTeX command generates the LaTeX logo.
but it makes a difference to you.  When you use
\LaTeX, making your input file as easy to read
as possible will be a great help as you write
your document and when you change it.  This sample
file shows how you can add comments to your own input
file.

Because printing is different from typewriting,
there are a number of things that you have to do
differently when preparing an input file than if
you were just typing the document directly.
Quotation marks like
       ``this''
have to be handled specially, as do quotes within
quotes:
       ``\,`this'            % \, separates the double and single quote.
        is what I just
        wrote, not  `that'\,''.

Dashes come in three sizes: an
       intra-word
dash, a medium dash for number ranges like
       1--2,
and a punctuation
       dash---like
this.

A sentence-ending space should be larger than the
space between words within a sentence.  You
sometimes have to type special commands in
conjunction with punctuation characters to get
this right, as in the following sentence.
       Gnats, gnus, etc.\ all  % `\ ' makes an inter-word space.
       begin with G\@.         % \@ marks end-of-sentence punctuation.
You should check the spaces after periods when
reading your output to make sure you haven't
forgotten any special cases.  Generating an
ellipsis
       \ldots\               % `\ ' is needed after `\ldots' because TeX
                             % ignores spaces after command names like \ldots
                             % made from \ + letters.
                             %
                             % Note how a `%' character causes TeX to ignore
                             % the end of the input line, so these blank lines
                             % do not start a new paragraph.
                             %
with the right spacing around the periods requires
a special command.

\LaTeX\ interprets some common characters as
commands, so you must type special commands to
generate them.  These characters include the
following:
       \$ \& \% \# \{ and \}.

In printing, text is usually emphasized with an
       \emph{italic}
type style.

\begin{em}
   A long segment of text can also be emphasized
   in this way.  Text within such a segment can be
   given \emph{additional} emphasis.
\end{em}

It is sometimes necessary to prevent \LaTeX\ from
breaking a line where it might otherwise do so.
This may be at a space, as between the ``Mr.''\ and
``Jones'' in
       ``Mr.~Jones'',        % ~ produces an unbreakable interword space.
or within a word---especially when the word is a
symbol like
       \mbox{\emph{itemnum}}
that makes little sense when hyphenated across
lines.

Footnotes\footnote{This is an example of a footnote.}
pose no problem.

\LaTeX\ is good at typesetting mathematical formulas
like
       \( x-3y + z = 7 \)
or
       \( a_{1} > x^{2n} + y^{2n} > x' \)
or
       \( \ip{A}{B} = \sum_{i} a_{i} b_{i} \).
The spaces you type in a formula are
ignored.  Remember that a letter like
       $x$                   % $ ... $  and  \( ... \)  are equivalent
is a formula when it denotes a mathematical
symbol, and it should be typed as one.

\section{Displayed Text}

Text is displayed by indenting it from the left
margin.  Quotations are commonly displayed.  There
are short quotations
\begin{quote}
   This is a short quotation.  It consists of a
   single paragraph of text.  See how it is formatted.
\end{quote}
and longer ones.
\begin{quotation}
   This is a longer quotation.  It consists of two
   paragraphs of text, neither of which are
   particularly interesting.

   This is the second paragraph of the quotation.  It
   is just as dull as the first paragraph.
\end{quotation}
Another frequently-displayed structure is a list.
The following is an example of an \emph{itemized}
list.
\begin{itemize}
   \item This is the first item of an itemized list.
         Each item in the list is marked with a ``tick''.
         You don't have to worry about what kind of tick
         mark is used.

   \item This is the second item of the list.  It
         contains another list nested inside it.  The inner
         list is an \emph{enumerated} list.
         \begin{enumerate}
            \item This is the first item of an enumerated
                  list that is nested within the itemized list.

            \item This is the second item of the inner list.
                  \LaTeX\ allows you to nest lists deeper than
                  you really should.
         \end{enumerate}
         This is the rest of the second item of the outer
         list.  It is no more interesting than any other
         part of the item.
   \item This is the third item of the list.
\end{itemize}
You can even display poetry.
\begin{verse}
   There is an environment
    for verse \\             % The \\ command separates lines
   Whose features some poets % within a stanza.
   will curse.

                             % One or more blank lines separate stanzas.

   For instead of making\\
   Them do \emph{all} line breaking, \\
   It allows them to put too many words on a line when they'd rather be
   forced to be terse.
\end{verse}

Mathematical formulas may also be displayed.  A
displayed formula
is
one-line long; multiline
formulas require special formatting instructions.
   \[  \ip{\Gamma}{\psi'} = x'' + y^{2} + z_{i}^{n}\]
Don't start a paragraph with a displayed equation,
nor make one a paragraph by itself.

\end{document}               % End of document.

\section{Bench Huge Section 21}
% synthetic bench section 21
% This is a sample LaTeX input file.  (Version of 12 August 2004.)
%
% A '%' character causes TeX to ignore all remaining text on the line,
% and is used for comments like this one.

\documentclass{article}      % Specifies the document class

                             % The preamble begins here.
\title{An Example Document}  % Declares the document's title.
\author{Leslie Lamport}      % Declares the author's name.
\date{January 21, 1994}      % Deleting this command produces today's date.

\newcommand{\ip}[2]{(#1, #2)}
                             % Defines \ip{arg1}{arg2} to mean
                             % (arg1, arg2).

%\newcommand{\ip}[2]{\langle #1 | #2\rangle}
                             % This is an alternative definition of
                             % \ip that is commented out.

\begin{document}             % End of preamble and beginning of text.

\maketitle                   % Produces the title.

This is an example input file.  Comparing it with
the output it generates can show you how to
produce a simple document of your own.

\section{Ordinary Text}      % Produces section heading.  Lower-level
                             % sections are begun with similar
                             % \subsection and \subsubsection commands.

The ends  of words and sentences are marked
  by   spaces. It  doesn't matter how many
spaces    you type; one is as good as 100.  The
end of   a line counts as a space.

One   or more   blank lines denote the  end
of  a paragraph.

Since any number of consecutive spaces are treated
like a single one, the formatting of the input
file makes no difference to
      \LaTeX,                % The \LaTeX command generates the LaTeX logo.
but it makes a difference to you.  When you use
\LaTeX, making your input file as easy to read
as possible will be a great help as you write
your document and when you change it.  This sample
file shows how you can add comments to your own input
file.

Because printing is different from typewriting,
there are a number of things that you have to do
differently when preparing an input file than if
you were just typing the document directly.
Quotation marks like
       ``this''
have to be handled specially, as do quotes within
quotes:
       ``\,`this'            % \, separates the double and single quote.
        is what I just
        wrote, not  `that'\,''.

Dashes come in three sizes: an
       intra-word
dash, a medium dash for number ranges like
       1--2,
and a punctuation
       dash---like
this.

A sentence-ending space should be larger than the
space between words within a sentence.  You
sometimes have to type special commands in
conjunction with punctuation characters to get
this right, as in the following sentence.
       Gnats, gnus, etc.\ all  % `\ ' makes an inter-word space.
       begin with G\@.         % \@ marks end-of-sentence punctuation.
You should check the spaces after periods when
reading your output to make sure you haven't
forgotten any special cases.  Generating an
ellipsis
       \ldots\               % `\ ' is needed after `\ldots' because TeX
                             % ignores spaces after command names like \ldots
                             % made from \ + letters.
                             %
                             % Note how a `%' character causes TeX to ignore
                             % the end of the input line, so these blank lines
                             % do not start a new paragraph.
                             %
with the right spacing around the periods requires
a special command.

\LaTeX\ interprets some common characters as
commands, so you must type special commands to
generate them.  These characters include the
following:
       \$ \& \% \# \{ and \}.

In printing, text is usually emphasized with an
       \emph{italic}
type style.

\begin{em}
   A long segment of text can also be emphasized
   in this way.  Text within such a segment can be
   given \emph{additional} emphasis.
\end{em}

It is sometimes necessary to prevent \LaTeX\ from
breaking a line where it might otherwise do so.
This may be at a space, as between the ``Mr.''\ and
``Jones'' in
       ``Mr.~Jones'',        % ~ produces an unbreakable interword space.
or within a word---especially when the word is a
symbol like
       \mbox{\emph{itemnum}}
that makes little sense when hyphenated across
lines.

Footnotes\footnote{This is an example of a footnote.}
pose no problem.

\LaTeX\ is good at typesetting mathematical formulas
like
       \( x-3y + z = 7 \)
or
       \( a_{1} > x^{2n} + y^{2n} > x' \)
or
       \( \ip{A}{B} = \sum_{i} a_{i} b_{i} \).
The spaces you type in a formula are
ignored.  Remember that a letter like
       $x$                   % $ ... $  and  \( ... \)  are equivalent
is a formula when it denotes a mathematical
symbol, and it should be typed as one.

\section{Displayed Text}

Text is displayed by indenting it from the left
margin.  Quotations are commonly displayed.  There
are short quotations
\begin{quote}
   This is a short quotation.  It consists of a
   single paragraph of text.  See how it is formatted.
\end{quote}
and longer ones.
\begin{quotation}
   This is a longer quotation.  It consists of two
   paragraphs of text, neither of which are
   particularly interesting.

   This is the second paragraph of the quotation.  It
   is just as dull as the first paragraph.
\end{quotation}
Another frequently-displayed structure is a list.
The following is an example of an \emph{itemized}
list.
\begin{itemize}
   \item This is the first item of an itemized list.
         Each item in the list is marked with a ``tick''.
         You don't have to worry about what kind of tick
         mark is used.

   \item This is the second item of the list.  It
         contains another list nested inside it.  The inner
         list is an \emph{enumerated} list.
         \begin{enumerate}
            \item This is the first item of an enumerated
                  list that is nested within the itemized list.

            \item This is the second item of the inner list.
                  \LaTeX\ allows you to nest lists deeper than
                  you really should.
         \end{enumerate}
         This is the rest of the second item of the outer
         list.  It is no more interesting than any other
         part of the item.
   \item This is the third item of the list.
\end{itemize}
You can even display poetry.
\begin{verse}
   There is an environment
    for verse \\             % The \\ command separates lines
   Whose features some poets % within a stanza.
   will curse.

                             % One or more blank lines separate stanzas.

   For instead of making\\
   Them do \emph{all} line breaking, \\
   It allows them to put too many words on a line when they'd rather be
   forced to be terse.
\end{verse}

Mathematical formulas may also be displayed.  A
displayed formula
is
one-line long; multiline
formulas require special formatting instructions.
   \[  \ip{\Gamma}{\psi'} = x'' + y^{2} + z_{i}^{n}\]
Don't start a paragraph with a displayed equation,
nor make one a paragraph by itself.

\end{document}               % End of document.

\section{Bench Huge Section 22}
% synthetic bench section 22
% This is a sample LaTeX input file.  (Version of 12 August 2004.)
%
% A '%' character causes TeX to ignore all remaining text on the line,
% and is used for comments like this one.

\documentclass{article}      % Specifies the document class

                             % The preamble begins here.
\title{An Example Document}  % Declares the document's title.
\author{Leslie Lamport}      % Declares the author's name.
\date{January 21, 1994}      % Deleting this command produces today's date.

\newcommand{\ip}[2]{(#1, #2)}
                             % Defines \ip{arg1}{arg2} to mean
                             % (arg1, arg2).

%\newcommand{\ip}[2]{\langle #1 | #2\rangle}
                             % This is an alternative definition of
                             % \ip that is commented out.

\begin{document}             % End of preamble and beginning of text.

\maketitle                   % Produces the title.

This is an example input file.  Comparing it with
the output it generates can show you how to
produce a simple document of your own.

\section{Ordinary Text}      % Produces section heading.  Lower-level
                             % sections are begun with similar
                             % \subsection and \subsubsection commands.

The ends  of words and sentences are marked
  by   spaces. It  doesn't matter how many
spaces    you type; one is as good as 100.  The
end of   a line counts as a space.

One   or more   blank lines denote the  end
of  a paragraph.

Since any number of consecutive spaces are treated
like a single one, the formatting of the input
file makes no difference to
      \LaTeX,                % The \LaTeX command generates the LaTeX logo.
but it makes a difference to you.  When you use
\LaTeX, making your input file as easy to read
as possible will be a great help as you write
your document and when you change it.  This sample
file shows how you can add comments to your own input
file.

Because printing is different from typewriting,
there are a number of things that you have to do
differently when preparing an input file than if
you were just typing the document directly.
Quotation marks like
       ``this''
have to be handled specially, as do quotes within
quotes:
       ``\,`this'            % \, separates the double and single quote.
        is what I just
        wrote, not  `that'\,''.

Dashes come in three sizes: an
       intra-word
dash, a medium dash for number ranges like
       1--2,
and a punctuation
       dash---like
this.

A sentence-ending space should be larger than the
space between words within a sentence.  You
sometimes have to type special commands in
conjunction with punctuation characters to get
this right, as in the following sentence.
       Gnats, gnus, etc.\ all  % `\ ' makes an inter-word space.
       begin with G\@.         % \@ marks end-of-sentence punctuation.
You should check the spaces after periods when
reading your output to make sure you haven't
forgotten any special cases.  Generating an
ellipsis
       \ldots\               % `\ ' is needed after `\ldots' because TeX
                             % ignores spaces after command names like \ldots
                             % made from \ + letters.
                             %
                             % Note how a `%' character causes TeX to ignore
                             % the end of the input line, so these blank lines
                             % do not start a new paragraph.
                             %
with the right spacing around the periods requires
a special command.

\LaTeX\ interprets some common characters as
commands, so you must type special commands to
generate them.  These characters include the
following:
       \$ \& \% \# \{ and \}.

In printing, text is usually emphasized with an
       \emph{italic}
type style.

\begin{em}
   A long segment of text can also be emphasized
   in this way.  Text within such a segment can be
   given \emph{additional} emphasis.
\end{em}

It is sometimes necessary to prevent \LaTeX\ from
breaking a line where it might otherwise do so.
This may be at a space, as between the ``Mr.''\ and
``Jones'' in
       ``Mr.~Jones'',        % ~ produces an unbreakable interword space.
or within a word---especially when the word is a
symbol like
       \mbox{\emph{itemnum}}
that makes little sense when hyphenated across
lines.

Footnotes\footnote{This is an example of a footnote.}
pose no problem.

\LaTeX\ is good at typesetting mathematical formulas
like
       \( x-3y + z = 7 \)
or
       \( a_{1} > x^{2n} + y^{2n} > x' \)
or
       \( \ip{A}{B} = \sum_{i} a_{i} b_{i} \).
The spaces you type in a formula are
ignored.  Remember that a letter like
       $x$                   % $ ... $  and  \( ... \)  are equivalent
is a formula when it denotes a mathematical
symbol, and it should be typed as one.

\section{Displayed Text}

Text is displayed by indenting it from the left
margin.  Quotations are commonly displayed.  There
are short quotations
\begin{quote}
   This is a short quotation.  It consists of a
   single paragraph of text.  See how it is formatted.
\end{quote}
and longer ones.
\begin{quotation}
   This is a longer quotation.  It consists of two
   paragraphs of text, neither of which are
   particularly interesting.

   This is the second paragraph of the quotation.  It
   is just as dull as the first paragraph.
\end{quotation}
Another frequently-displayed structure is a list.
The following is an example of an \emph{itemized}
list.
\begin{itemize}
   \item This is the first item of an itemized list.
         Each item in the list is marked with a ``tick''.
         You don't have to worry about what kind of tick
         mark is used.

   \item This is the second item of the list.  It
         contains another list nested inside it.  The inner
         list is an \emph{enumerated} list.
         \begin{enumerate}
            \item This is the first item of an enumerated
                  list that is nested within the itemized list.

            \item This is the second item of the inner list.
                  \LaTeX\ allows you to nest lists deeper than
                  you really should.
         \end{enumerate}
         This is the rest of the second item of the outer
         list.  It is no more interesting than any other
         part of the item.
   \item This is the third item of the list.
\end{itemize}
You can even display poetry.
\begin{verse}
   There is an environment
    for verse \\             % The \\ command separates lines
   Whose features some poets % within a stanza.
   will curse.

                             % One or more blank lines separate stanzas.

   For instead of making\\
   Them do \emph{all} line breaking, \\
   It allows them to put too many words on a line when they'd rather be
   forced to be terse.
\end{verse}

Mathematical formulas may also be displayed.  A
displayed formula
is
one-line long; multiline
formulas require special formatting instructions.
   \[  \ip{\Gamma}{\psi'} = x'' + y^{2} + z_{i}^{n}\]
Don't start a paragraph with a displayed equation,
nor make one a paragraph by itself.

\end{document}               % End of document.

\section{Bench Huge Section 23}
% synthetic bench section 23
% This is a sample LaTeX input file.  (Version of 12 August 2004.)
%
% A '%' character causes TeX to ignore all remaining text on the line,
% and is used for comments like this one.

\documentclass{article}      % Specifies the document class

                             % The preamble begins here.
\title{An Example Document}  % Declares the document's title.
\author{Leslie Lamport}      % Declares the author's name.
\date{January 21, 1994}      % Deleting this command produces today's date.

\newcommand{\ip}[2]{(#1, #2)}
                             % Defines \ip{arg1}{arg2} to mean
                             % (arg1, arg2).

%\newcommand{\ip}[2]{\langle #1 | #2\rangle}
                             % This is an alternative definition of
                             % \ip that is commented out.

\begin{document}             % End of preamble and beginning of text.

\maketitle                   % Produces the title.

This is an example input file.  Comparing it with
the output it generates can show you how to
produce a simple document of your own.

\section{Ordinary Text}      % Produces section heading.  Lower-level
                             % sections are begun with similar
                             % \subsection and \subsubsection commands.

The ends  of words and sentences are marked
  by   spaces. It  doesn't matter how many
spaces    you type; one is as good as 100.  The
end of   a line counts as a space.

One   or more   blank lines denote the  end
of  a paragraph.

Since any number of consecutive spaces are treated
like a single one, the formatting of the input
file makes no difference to
      \LaTeX,                % The \LaTeX command generates the LaTeX logo.
but it makes a difference to you.  When you use
\LaTeX, making your input file as easy to read
as possible will be a great help as you write
your document and when you change it.  This sample
file shows how you can add comments to your own input
file.

Because printing is different from typewriting,
there are a number of things that you have to do
differently when preparing an input file than if
you were just typing the document directly.
Quotation marks like
       ``this''
have to be handled specially, as do quotes within
quotes:
       ``\,`this'            % \, separates the double and single quote.
        is what I just
        wrote, not  `that'\,''.

Dashes come in three sizes: an
       intra-word
dash, a medium dash for number ranges like
       1--2,
and a punctuation
       dash---like
this.

A sentence-ending space should be larger than the
space between words within a sentence.  You
sometimes have to type special commands in
conjunction with punctuation characters to get
this right, as in the following sentence.
       Gnats, gnus, etc.\ all  % `\ ' makes an inter-word space.
       begin with G\@.         % \@ marks end-of-sentence punctuation.
You should check the spaces after periods when
reading your output to make sure you haven't
forgotten any special cases.  Generating an
ellipsis
       \ldots\               % `\ ' is needed after `\ldots' because TeX
                             % ignores spaces after command names like \ldots
                             % made from \ + letters.
                             %
                             % Note how a `%' character causes TeX to ignore
                             % the end of the input line, so these blank lines
                             % do not start a new paragraph.
                             %
with the right spacing around the periods requires
a special command.

\LaTeX\ interprets some common characters as
commands, so you must type special commands to
generate them.  These characters include the
following:
       \$ \& \% \# \{ and \}.

In printing, text is usually emphasized with an
       \emph{italic}
type style.

\begin{em}
   A long segment of text can also be emphasized
   in this way.  Text within such a segment can be
   given \emph{additional} emphasis.
\end{em}

It is sometimes necessary to prevent \LaTeX\ from
breaking a line where it might otherwise do so.
This may be at a space, as between the ``Mr.''\ and
``Jones'' in
       ``Mr.~Jones'',        % ~ produces an unbreakable interword space.
or within a word---especially when the word is a
symbol like
       \mbox{\emph{itemnum}}
that makes little sense when hyphenated across
lines.

Footnotes\footnote{This is an example of a footnote.}
pose no problem.

\LaTeX\ is good at typesetting mathematical formulas
like
       \( x-3y + z = 7 \)
or
       \( a_{1} > x^{2n} + y^{2n} > x' \)
or
       \( \ip{A}{B} = \sum_{i} a_{i} b_{i} \).
The spaces you type in a formula are
ignored.  Remember that a letter like
       $x$                   % $ ... $  and  \( ... \)  are equivalent
is a formula when it denotes a mathematical
symbol, and it should be typed as one.

\section{Displayed Text}

Text is displayed by indenting it from the left
margin.  Quotations are commonly displayed.  There
are short quotations
\begin{quote}
   This is a short quotation.  It consists of a
   single paragraph of text.  See how it is formatted.
\end{quote}
and longer ones.
\begin{quotation}
   This is a longer quotation.  It consists of two
   paragraphs of text, neither of which are
   particularly interesting.

   This is the second paragraph of the quotation.  It
   is just as dull as the first paragraph.
\end{quotation}
Another frequently-displayed structure is a list.
The following is an example of an \emph{itemized}
list.
\begin{itemize}
   \item This is the first item of an itemized list.
         Each item in the list is marked with a ``tick''.
         You don't have to worry about what kind of tick
         mark is used.

   \item This is the second item of the list.  It
         contains another list nested inside it.  The inner
         list is an \emph{enumerated} list.
         \begin{enumerate}
            \item This is the first item of an enumerated
                  list that is nested within the itemized list.

            \item This is the second item of the inner list.
                  \LaTeX\ allows you to nest lists deeper than
                  you really should.
         \end{enumerate}
         This is the rest of the second item of the outer
         list.  It is no more interesting than any other
         part of the item.
   \item This is the third item of the list.
\end{itemize}
You can even display poetry.
\begin{verse}
   There is an environment
    for verse \\             % The \\ command separates lines
   Whose features some poets % within a stanza.
   will curse.

                             % One or more blank lines separate stanzas.

   For instead of making\\
   Them do \emph{all} line breaking, \\
   It allows them to put too many words on a line when they'd rather be
   forced to be terse.
\end{verse}

Mathematical formulas may also be displayed.  A
displayed formula
is
one-line long; multiline
formulas require special formatting instructions.
   \[  \ip{\Gamma}{\psi'} = x'' + y^{2} + z_{i}^{n}\]
Don't start a paragraph with a displayed equation,
nor make one a paragraph by itself.

\end{document}               % End of document.

\section{Bench Huge Section 24}
% synthetic bench section 24
% This is a sample LaTeX input file.  (Version of 12 August 2004.)
%
% A '%' character causes TeX to ignore all remaining text on the line,
% and is used for comments like this one.

\documentclass{article}      % Specifies the document class

                             % The preamble begins here.
\title{An Example Document}  % Declares the document's title.
\author{Leslie Lamport}      % Declares the author's name.
\date{January 21, 1994}      % Deleting this command produces today's date.

\newcommand{\ip}[2]{(#1, #2)}
                             % Defines \ip{arg1}{arg2} to mean
                             % (arg1, arg2).

%\newcommand{\ip}[2]{\langle #1 | #2\rangle}
                             % This is an alternative definition of
                             % \ip that is commented out.

\begin{document}             % End of preamble and beginning of text.

\maketitle                   % Produces the title.

This is an example input file.  Comparing it with
the output it generates can show you how to
produce a simple document of your own.

\section{Ordinary Text}      % Produces section heading.  Lower-level
                             % sections are begun with similar
                             % \subsection and \subsubsection commands.

The ends  of words and sentences are marked
  by   spaces. It  doesn't matter how many
spaces    you type; one is as good as 100.  The
end of   a line counts as a space.

One   or more   blank lines denote the  end
of  a paragraph.

Since any number of consecutive spaces are treated
like a single one, the formatting of the input
file makes no difference to
      \LaTeX,                % The \LaTeX command generates the LaTeX logo.
but it makes a difference to you.  When you use
\LaTeX, making your input file as easy to read
as possible will be a great help as you write
your document and when you change it.  This sample
file shows how you can add comments to your own input
file.

Because printing is different from typewriting,
there are a number of things that you have to do
differently when preparing an input file than if
you were just typing the document directly.
Quotation marks like
       ``this''
have to be handled specially, as do quotes within
quotes:
       ``\,`this'            % \, separates the double and single quote.
        is what I just
        wrote, not  `that'\,''.

Dashes come in three sizes: an
       intra-word
dash, a medium dash for number ranges like
       1--2,
and a punctuation
       dash---like
this.

A sentence-ending space should be larger than the
space between words within a sentence.  You
sometimes have to type special commands in
conjunction with punctuation characters to get
this right, as in the following sentence.
       Gnats, gnus, etc.\ all  % `\ ' makes an inter-word space.
       begin with G\@.         % \@ marks end-of-sentence punctuation.
You should check the spaces after periods when
reading your output to make sure you haven't
forgotten any special cases.  Generating an
ellipsis
       \ldots\               % `\ ' is needed after `\ldots' because TeX
                             % ignores spaces after command names like \ldots
                             % made from \ + letters.
                             %
                             % Note how a `%' character causes TeX to ignore
                             % the end of the input line, so these blank lines
                             % do not start a new paragraph.
                             %
with the right spacing around the periods requires
a special command.

\LaTeX\ interprets some common characters as
commands, so you must type special commands to
generate them.  These characters include the
following:
       \$ \& \% \# \{ and \}.

In printing, text is usually emphasized with an
       \emph{italic}
type style.

\begin{em}
   A long segment of text can also be emphasized
   in this way.  Text within such a segment can be
   given \emph{additional} emphasis.
\end{em}

It is sometimes necessary to prevent \LaTeX\ from
breaking a line where it might otherwise do so.
This may be at a space, as between the ``Mr.''\ and
``Jones'' in
       ``Mr.~Jones'',        % ~ produces an unbreakable interword space.
or within a word---especially when the word is a
symbol like
       \mbox{\emph{itemnum}}
that makes little sense when hyphenated across
lines.

Footnotes\footnote{This is an example of a footnote.}
pose no problem.

\LaTeX\ is good at typesetting mathematical formulas
like
       \( x-3y + z = 7 \)
or
       \( a_{1} > x^{2n} + y^{2n} > x' \)
or
       \( \ip{A}{B} = \sum_{i} a_{i} b_{i} \).
The spaces you type in a formula are
ignored.  Remember that a letter like
       $x$                   % $ ... $  and  \( ... \)  are equivalent
is a formula when it denotes a mathematical
symbol, and it should be typed as one.

\section{Displayed Text}

Text is displayed by indenting it from the left
margin.  Quotations are commonly displayed.  There
are short quotations
\begin{quote}
   This is a short quotation.  It consists of a
   single paragraph of text.  See how it is formatted.
\end{quote}
and longer ones.
\begin{quotation}
   This is a longer quotation.  It consists of two
   paragraphs of text, neither of which are
   particularly interesting.

   This is the second paragraph of the quotation.  It
   is just as dull as the first paragraph.
\end{quotation}
Another frequently-displayed structure is a list.
The following is an example of an \emph{itemized}
list.
\begin{itemize}
   \item This is the first item of an itemized list.
         Each item in the list is marked with a ``tick''.
         You don't have to worry about what kind of tick
         mark is used.

   \item This is the second item of the list.  It
         contains another list nested inside it.  The inner
         list is an \emph{enumerated} list.
         \begin{enumerate}
            \item This is the first item of an enumerated
                  list that is nested within the itemized list.

            \item This is the second item of the inner list.
                  \LaTeX\ allows you to nest lists deeper than
                  you really should.
         \end{enumerate}
         This is the rest of the second item of the outer
         list.  It is no more interesting than any other
         part of the item.
   \item This is the third item of the list.
\end{itemize}
You can even display poetry.
\begin{verse}
   There is an environment
    for verse \\             % The \\ command separates lines
   Whose features some poets % within a stanza.
   will curse.

                             % One or more blank lines separate stanzas.

   For instead of making\\
   Them do \emph{all} line breaking, \\
   It allows them to put too many words on a line when they'd rather be
   forced to be terse.
\end{verse}

Mathematical formulas may also be displayed.  A
displayed formula
is
one-line long; multiline
formulas require special formatting instructions.
   \[  \ip{\Gamma}{\psi'} = x'' + y^{2} + z_{i}^{n}\]
Don't start a paragraph with a displayed equation,
nor make one a paragraph by itself.

\end{document}               % End of document.

\section{Bench Huge Section 25}
% synthetic bench section 25
% This is a sample LaTeX input file.  (Version of 12 August 2004.)
%
% A '%' character causes TeX to ignore all remaining text on the line,
% and is used for comments like this one.

\documentclass{article}      % Specifies the document class

                             % The preamble begins here.
\title{An Example Document}  % Declares the document's title.
\author{Leslie Lamport}      % Declares the author's name.
\date{January 21, 1994}      % Deleting this command produces today's date.

\newcommand{\ip}[2]{(#1, #2)}
                             % Defines \ip{arg1}{arg2} to mean
                             % (arg1, arg2).

%\newcommand{\ip}[2]{\langle #1 | #2\rangle}
                             % This is an alternative definition of
                             % \ip that is commented out.

\begin{document}             % End of preamble and beginning of text.

\maketitle                   % Produces the title.

This is an example input file.  Comparing it with
the output it generates can show you how to
produce a simple document of your own.

\section{Ordinary Text}      % Produces section heading.  Lower-level
                             % sections are begun with similar
                             % \subsection and \subsubsection commands.

The ends  of words and sentences are marked
  by   spaces. It  doesn't matter how many
spaces    you type; one is as good as 100.  The
end of   a line counts as a space.

One   or more   blank lines denote the  end
of  a paragraph.

Since any number of consecutive spaces are treated
like a single one, the formatting of the input
file makes no difference to
      \LaTeX,                % The \LaTeX command generates the LaTeX logo.
but it makes a difference to you.  When you use
\LaTeX, making your input file as easy to read
as possible will be a great help as you write
your document and when you change it.  This sample
file shows how you can add comments to your own input
file.

Because printing is different from typewriting,
there are a number of things that you have to do
differently when preparing an input file than if
you were just typing the document directly.
Quotation marks like
       ``this''
have to be handled specially, as do quotes within
quotes:
       ``\,`this'            % \, separates the double and single quote.
        is what I just
        wrote, not  `that'\,''.

Dashes come in three sizes: an
       intra-word
dash, a medium dash for number ranges like
       1--2,
and a punctuation
       dash---like
this.

A sentence-ending space should be larger than the
space between words within a sentence.  You
sometimes have to type special commands in
conjunction with punctuation characters to get
this right, as in the following sentence.
       Gnats, gnus, etc.\ all  % `\ ' makes an inter-word space.
       begin with G\@.         % \@ marks end-of-sentence punctuation.
You should check the spaces after periods when
reading your output to make sure you haven't
forgotten any special cases.  Generating an
ellipsis
       \ldots\               % `\ ' is needed after `\ldots' because TeX
                             % ignores spaces after command names like \ldots
                             % made from \ + letters.
                             %
                             % Note how a `%' character causes TeX to ignore
                             % the end of the input line, so these blank lines
                             % do not start a new paragraph.
                             %
with the right spacing around the periods requires
a special command.

\LaTeX\ interprets some common characters as
commands, so you must type special commands to
generate them.  These characters include the
following:
       \$ \& \% \# \{ and \}.

In printing, text is usually emphasized with an
       \emph{italic}
type style.

\begin{em}
   A long segment of text can also be emphasized
   in this way.  Text within such a segment can be
   given \emph{additional} emphasis.
\end{em}

It is sometimes necessary to prevent \LaTeX\ from
breaking a line where it might otherwise do so.
This may be at a space, as between the ``Mr.''\ and
``Jones'' in
       ``Mr.~Jones'',        % ~ produces an unbreakable interword space.
or within a word---especially when the word is a
symbol like
       \mbox{\emph{itemnum}}
that makes little sense when hyphenated across
lines.

Footnotes\footnote{This is an example of a footnote.}
pose no problem.

\LaTeX\ is good at typesetting mathematical formulas
like
       \( x-3y + z = 7 \)
or
       \( a_{1} > x^{2n} + y^{2n} > x' \)
or
       \( \ip{A}{B} = \sum_{i} a_{i} b_{i} \).
The spaces you type in a formula are
ignored.  Remember that a letter like
       $x$                   % $ ... $  and  \( ... \)  are equivalent
is a formula when it denotes a mathematical
symbol, and it should be typed as one.

\section{Displayed Text}

Text is displayed by indenting it from the left
margin.  Quotations are commonly displayed.  There
are short quotations
\begin{quote}
   This is a short quotation.  It consists of a
   single paragraph of text.  See how it is formatted.
\end{quote}
and longer ones.
\begin{quotation}
   This is a longer quotation.  It consists of two
   paragraphs of text, neither of which are
   particularly interesting.

   This is the second paragraph of the quotation.  It
   is just as dull as the first paragraph.
\end{quotation}
Another frequently-displayed structure is a list.
The following is an example of an \emph{itemized}
list.
\begin{itemize}
   \item This is the first item of an itemized list.
         Each item in the list is marked with a ``tick''.
         You don't have to worry about what kind of tick
         mark is used.

   \item This is the second item of the list.  It
         contains another list nested inside it.  The inner
         list is an \emph{enumerated} list.
         \begin{enumerate}
            \item This is the first item of an enumerated
                  list that is nested within the itemized list.

            \item This is the second item of the inner list.
                  \LaTeX\ allows you to nest lists deeper than
                  you really should.
         \end{enumerate}
         This is the rest of the second item of the outer
         list.  It is no more interesting than any other
         part of the item.
   \item This is the third item of the list.
\end{itemize}
You can even display poetry.
\begin{verse}
   There is an environment
    for verse \\             % The \\ command separates lines
   Whose features some poets % within a stanza.
   will curse.

                             % One or more blank lines separate stanzas.

   For instead of making\\
   Them do \emph{all} line breaking, \\
   It allows them to put too many words on a line when they'd rather be
   forced to be terse.
\end{verse}

Mathematical formulas may also be displayed.  A
displayed formula
is
one-line long; multiline
formulas require special formatting instructions.
   \[  \ip{\Gamma}{\psi'} = x'' + y^{2} + z_{i}^{n}\]
Don't start a paragraph with a displayed equation,
nor make one a paragraph by itself.

\end{document}               % End of document.

\section{Bench Huge Section 26}
% synthetic bench section 26
% This is a sample LaTeX input file.  (Version of 12 August 2004.)
%
% A '%' character causes TeX to ignore all remaining text on the line,
% and is used for comments like this one.

\documentclass{article}      % Specifies the document class

                             % The preamble begins here.
\title{An Example Document}  % Declares the document's title.
\author{Leslie Lamport}      % Declares the author's name.
\date{January 21, 1994}      % Deleting this command produces today's date.

\newcommand{\ip}[2]{(#1, #2)}
                             % Defines \ip{arg1}{arg2} to mean
                             % (arg1, arg2).

%\newcommand{\ip}[2]{\langle #1 | #2\rangle}
                             % This is an alternative definition of
                             % \ip that is commented out.

\begin{document}             % End of preamble and beginning of text.

\maketitle                   % Produces the title.

This is an example input file.  Comparing it with
the output it generates can show you how to
produce a simple document of your own.

\section{Ordinary Text}      % Produces section heading.  Lower-level
                             % sections are begun with similar
                             % \subsection and \subsubsection commands.

The ends  of words and sentences are marked
  by   spaces. It  doesn't matter how many
spaces    you type; one is as good as 100.  The
end of   a line counts as a space.

One   or more   blank lines denote the  end
of  a paragraph.

Since any number of consecutive spaces are treated
like a single one, the formatting of the input
file makes no difference to
      \LaTeX,                % The \LaTeX command generates the LaTeX logo.
but it makes a difference to you.  When you use
\LaTeX, making your input file as easy to read
as possible will be a great help as you write
your document and when you change it.  This sample
file shows how you can add comments to your own input
file.

Because printing is different from typewriting,
there are a number of things that you have to do
differently when preparing an input file than if
you were just typing the document directly.
Quotation marks like
       ``this''
have to be handled specially, as do quotes within
quotes:
       ``\,`this'            % \, separates the double and single quote.
        is what I just
        wrote, not  `that'\,''.

Dashes come in three sizes: an
       intra-word
dash, a medium dash for number ranges like
       1--2,
and a punctuation
       dash---like
this.

A sentence-ending space should be larger than the
space between words within a sentence.  You
sometimes have to type special commands in
conjunction with punctuation characters to get
this right, as in the following sentence.
       Gnats, gnus, etc.\ all  % `\ ' makes an inter-word space.
       begin with G\@.         % \@ marks end-of-sentence punctuation.
You should check the spaces after periods when
reading your output to make sure you haven't
forgotten any special cases.  Generating an
ellipsis
       \ldots\               % `\ ' is needed after `\ldots' because TeX
                             % ignores spaces after command names like \ldots
                             % made from \ + letters.
                             %
                             % Note how a `%' character causes TeX to ignore
                             % the end of the input line, so these blank lines
                             % do not start a new paragraph.
                             %
with the right spacing around the periods requires
a special command.

\LaTeX\ interprets some common characters as
commands, so you must type special commands to
generate them.  These characters include the
following:
       \$ \& \% \# \{ and \}.

In printing, text is usually emphasized with an
       \emph{italic}
type style.

\begin{em}
   A long segment of text can also be emphasized
   in this way.  Text within such a segment can be
   given \emph{additional} emphasis.
\end{em}

It is sometimes necessary to prevent \LaTeX\ from
breaking a line where it might otherwise do so.
This may be at a space, as between the ``Mr.''\ and
``Jones'' in
       ``Mr.~Jones'',        % ~ produces an unbreakable interword space.
or within a word---especially when the word is a
symbol like
       \mbox{\emph{itemnum}}
that makes little sense when hyphenated across
lines.

Footnotes\footnote{This is an example of a footnote.}
pose no problem.

\LaTeX\ is good at typesetting mathematical formulas
like
       \( x-3y + z = 7 \)
or
       \( a_{1} > x^{2n} + y^{2n} > x' \)
or
       \( \ip{A}{B} = \sum_{i} a_{i} b_{i} \).
The spaces you type in a formula are
ignored.  Remember that a letter like
       $x$                   % $ ... $  and  \( ... \)  are equivalent
is a formula when it denotes a mathematical
symbol, and it should be typed as one.

\section{Displayed Text}

Text is displayed by indenting it from the left
margin.  Quotations are commonly displayed.  There
are short quotations
\begin{quote}
   This is a short quotation.  It consists of a
   single paragraph of text.  See how it is formatted.
\end{quote}
and longer ones.
\begin{quotation}
   This is a longer quotation.  It consists of two
   paragraphs of text, neither of which are
   particularly interesting.

   This is the second paragraph of the quotation.  It
   is just as dull as the first paragraph.
\end{quotation}
Another frequently-displayed structure is a list.
The following is an example of an \emph{itemized}
list.
\begin{itemize}
   \item This is the first item of an itemized list.
         Each item in the list is marked with a ``tick''.
         You don't have to worry about what kind of tick
         mark is used.

   \item This is the second item of the list.  It
         contains another list nested inside it.  The inner
         list is an \emph{enumerated} list.
         \begin{enumerate}
            \item This is the first item of an enumerated
                  list that is nested within the itemized list.

            \item This is the second item of the inner list.
                  \LaTeX\ allows you to nest lists deeper than
                  you really should.
         \end{enumerate}
         This is the rest of the second item of the outer
         list.  It is no more interesting than any other
         part of the item.
   \item This is the third item of the list.
\end{itemize}
You can even display poetry.
\begin{verse}
   There is an environment
    for verse \\             % The \\ command separates lines
   Whose features some poets % within a stanza.
   will curse.

                             % One or more blank lines separate stanzas.

   For instead of making\\
   Them do \emph{all} line breaking, \\
   It allows them to put too many words on a line when they'd rather be
   forced to be terse.
\end{verse}

Mathematical formulas may also be displayed.  A
displayed formula
is
one-line long; multiline
formulas require special formatting instructions.
   \[  \ip{\Gamma}{\psi'} = x'' + y^{2} + z_{i}^{n}\]
Don't start a paragraph with a displayed equation,
nor make one a paragraph by itself.

\end{document}               % End of document.

\section{Bench Huge Section 27}
% synthetic bench section 27
% This is a sample LaTeX input file.  (Version of 12 August 2004.)
%
% A '%' character causes TeX to ignore all remaining text on the line,
% and is used for comments like this one.

\documentclass{article}      % Specifies the document class

                             % The preamble begins here.
\title{An Example Document}  % Declares the document's title.
\author{Leslie Lamport}      % Declares the author's name.
\date{January 21, 1994}      % Deleting this command produces today's date.

\newcommand{\ip}[2]{(#1, #2)}
                             % Defines \ip{arg1}{arg2} to mean
                             % (arg1, arg2).

%\newcommand{\ip}[2]{\langle #1 | #2\rangle}
                             % This is an alternative definition of
                             % \ip that is commented out.

\begin{document}             % End of preamble and beginning of text.

\maketitle                   % Produces the title.

This is an example input file.  Comparing it with
the output it generates can show you how to
produce a simple document of your own.

\section{Ordinary Text}      % Produces section heading.  Lower-level
                             % sections are begun with similar
                             % \subsection and \subsubsection commands.

The ends  of words and sentences are marked
  by   spaces. It  doesn't matter how many
spaces    you type; one is as good as 100.  The
end of   a line counts as a space.

One   or more   blank lines denote the  end
of  a paragraph.

Since any number of consecutive spaces are treated
like a single one, the formatting of the input
file makes no difference to
      \LaTeX,                % The \LaTeX command generates the LaTeX logo.
but it makes a difference to you.  When you use
\LaTeX, making your input file as easy to read
as possible will be a great help as you write
your document and when you change it.  This sample
file shows how you can add comments to your own input
file.

Because printing is different from typewriting,
there are a number of things that you have to do
differently when preparing an input file than if
you were just typing the document directly.
Quotation marks like
       ``this''
have to be handled specially, as do quotes within
quotes:
       ``\,`this'            % \, separates the double and single quote.
        is what I just
        wrote, not  `that'\,''.

Dashes come in three sizes: an
       intra-word
dash, a medium dash for number ranges like
       1--2,
and a punctuation
       dash---like
this.

A sentence-ending space should be larger than the
space between words within a sentence.  You
sometimes have to type special commands in
conjunction with punctuation characters to get
this right, as in the following sentence.
       Gnats, gnus, etc.\ all  % `\ ' makes an inter-word space.
       begin with G\@.         % \@ marks end-of-sentence punctuation.
You should check the spaces after periods when
reading your output to make sure you haven't
forgotten any special cases.  Generating an
ellipsis
       \ldots\               % `\ ' is needed after `\ldots' because TeX
                             % ignores spaces after command names like \ldots
                             % made from \ + letters.
                             %
                             % Note how a `%' character causes TeX to ignore
                             % the end of the input line, so these blank lines
                             % do not start a new paragraph.
                             %
with the right spacing around the periods requires
a special command.

\LaTeX\ interprets some common characters as
commands, so you must type special commands to
generate them.  These characters include the
following:
       \$ \& \% \# \{ and \}.

In printing, text is usually emphasized with an
       \emph{italic}
type style.

\begin{em}
   A long segment of text can also be emphasized
   in this way.  Text within such a segment can be
   given \emph{additional} emphasis.
\end{em}

It is sometimes necessary to prevent \LaTeX\ from
breaking a line where it might otherwise do so.
This may be at a space, as between the ``Mr.''\ and
``Jones'' in
       ``Mr.~Jones'',        % ~ produces an unbreakable interword space.
or within a word---especially when the word is a
symbol like
       \mbox{\emph{itemnum}}
that makes little sense when hyphenated across
lines.

Footnotes\footnote{This is an example of a footnote.}
pose no problem.

\LaTeX\ is good at typesetting mathematical formulas
like
       \( x-3y + z = 7 \)
or
       \( a_{1} > x^{2n} + y^{2n} > x' \)
or
       \( \ip{A}{B} = \sum_{i} a_{i} b_{i} \).
The spaces you type in a formula are
ignored.  Remember that a letter like
       $x$                   % $ ... $  and  \( ... \)  are equivalent
is a formula when it denotes a mathematical
symbol, and it should be typed as one.

\section{Displayed Text}

Text is displayed by indenting it from the left
margin.  Quotations are commonly displayed.  There
are short quotations
\begin{quote}
   This is a short quotation.  It consists of a
   single paragraph of text.  See how it is formatted.
\end{quote}
and longer ones.
\begin{quotation}
   This is a longer quotation.  It consists of two
   paragraphs of text, neither of which are
   particularly interesting.

   This is the second paragraph of the quotation.  It
   is just as dull as the first paragraph.
\end{quotation}
Another frequently-displayed structure is a list.
The following is an example of an \emph{itemized}
list.
\begin{itemize}
   \item This is the first item of an itemized list.
         Each item in the list is marked with a ``tick''.
         You don't have to worry about what kind of tick
         mark is used.

   \item This is the second item of the list.  It
         contains another list nested inside it.  The inner
         list is an \emph{enumerated} list.
         \begin{enumerate}
            \item This is the first item of an enumerated
                  list that is nested within the itemized list.

            \item This is the second item of the inner list.
                  \LaTeX\ allows you to nest lists deeper than
                  you really should.
         \end{enumerate}
         This is the rest of the second item of the outer
         list.  It is no more interesting than any other
         part of the item.
   \item This is the third item of the list.
\end{itemize}
You can even display poetry.
\begin{verse}
   There is an environment
    for verse \\             % The \\ command separates lines
   Whose features some poets % within a stanza.
   will curse.

                             % One or more blank lines separate stanzas.

   For instead of making\\
   Them do \emph{all} line breaking, \\
   It allows them to put too many words on a line when they'd rather be
   forced to be terse.
\end{verse}

Mathematical formulas may also be displayed.  A
displayed formula
is
one-line long; multiline
formulas require special formatting instructions.
   \[  \ip{\Gamma}{\psi'} = x'' + y^{2} + z_{i}^{n}\]
Don't start a paragraph with a displayed equation,
nor make one a paragraph by itself.

\end{document}               % End of document.

\section{Bench Huge Section 28}
% synthetic bench section 28
% This is a sample LaTeX input file.  (Version of 12 August 2004.)
%
% A '%' character causes TeX to ignore all remaining text on the line,
% and is used for comments like this one.

\documentclass{article}      % Specifies the document class

                             % The preamble begins here.
\title{An Example Document}  % Declares the document's title.
\author{Leslie Lamport}      % Declares the author's name.
\date{January 21, 1994}      % Deleting this command produces today's date.

\newcommand{\ip}[2]{(#1, #2)}
                             % Defines \ip{arg1}{arg2} to mean
                             % (arg1, arg2).

%\newcommand{\ip}[2]{\langle #1 | #2\rangle}
                             % This is an alternative definition of
                             % \ip that is commented out.

\begin{document}             % End of preamble and beginning of text.

\maketitle                   % Produces the title.

This is an example input file.  Comparing it with
the output it generates can show you how to
produce a simple document of your own.

\section{Ordinary Text}      % Produces section heading.  Lower-level
                             % sections are begun with similar
                             % \subsection and \subsubsection commands.

The ends  of words and sentences are marked
  by   spaces. It  doesn't matter how many
spaces    you type; one is as good as 100.  The
end of   a line counts as a space.

One   or more   blank lines denote the  end
of  a paragraph.

Since any number of consecutive spaces are treated
like a single one, the formatting of the input
file makes no difference to
      \LaTeX,                % The \LaTeX command generates the LaTeX logo.
but it makes a difference to you.  When you use
\LaTeX, making your input file as easy to read
as possible will be a great help as you write
your document and when you change it.  This sample
file shows how you can add comments to your own input
file.

Because printing is different from typewriting,
there are a number of things that you have to do
differently when preparing an input file than if
you were just typing the document directly.
Quotation marks like
       ``this''
have to be handled specially, as do quotes within
quotes:
       ``\,`this'            % \, separates the double and single quote.
        is what I just
        wrote, not  `that'\,''.

Dashes come in three sizes: an
       intra-word
dash, a medium dash for number ranges like
       1--2,
and a punctuation
       dash---like
this.

A sentence-ending space should be larger than the
space between words within a sentence.  You
sometimes have to type special commands in
conjunction with punctuation characters to get
this right, as in the following sentence.
       Gnats, gnus, etc.\ all  % `\ ' makes an inter-word space.
       begin with G\@.         % \@ marks end-of-sentence punctuation.
You should check the spaces after periods when
reading your output to make sure you haven't
forgotten any special cases.  Generating an
ellipsis
       \ldots\               % `\ ' is needed after `\ldots' because TeX
                             % ignores spaces after command names like \ldots
                             % made from \ + letters.
                             %
                             % Note how a `%' character causes TeX to ignore
                             % the end of the input line, so these blank lines
                             % do not start a new paragraph.
                             %
with the right spacing around the periods requires
a special command.

\LaTeX\ interprets some common characters as
commands, so you must type special commands to
generate them.  These characters include the
following:
       \$ \& \% \# \{ and \}.

In printing, text is usually emphasized with an
       \emph{italic}
type style.

\begin{em}
   A long segment of text can also be emphasized
   in this way.  Text within such a segment can be
   given \emph{additional} emphasis.
\end{em}

It is sometimes necessary to prevent \LaTeX\ from
breaking a line where it might otherwise do so.
This may be at a space, as between the ``Mr.''\ and
``Jones'' in
       ``Mr.~Jones'',        % ~ produces an unbreakable interword space.
or within a word---especially when the word is a
symbol like
       \mbox{\emph{itemnum}}
that makes little sense when hyphenated across
lines.

Footnotes\footnote{This is an example of a footnote.}
pose no problem.

\LaTeX\ is good at typesetting mathematical formulas
like
       \( x-3y + z = 7 \)
or
       \( a_{1} > x^{2n} + y^{2n} > x' \)
or
       \( \ip{A}{B} = \sum_{i} a_{i} b_{i} \).
The spaces you type in a formula are
ignored.  Remember that a letter like
       $x$                   % $ ... $  and  \( ... \)  are equivalent
is a formula when it denotes a mathematical
symbol, and it should be typed as one.

\section{Displayed Text}

Text is displayed by indenting it from the left
margin.  Quotations are commonly displayed.  There
are short quotations
\begin{quote}
   This is a short quotation.  It consists of a
   single paragraph of text.  See how it is formatted.
\end{quote}
and longer ones.
\begin{quotation}
   This is a longer quotation.  It consists of two
   paragraphs of text, neither of which are
   particularly interesting.

   This is the second paragraph of the quotation.  It
   is just as dull as the first paragraph.
\end{quotation}
Another frequently-displayed structure is a list.
The following is an example of an \emph{itemized}
list.
\begin{itemize}
   \item This is the first item of an itemized list.
         Each item in the list is marked with a ``tick''.
         You don't have to worry about what kind of tick
         mark is used.

   \item This is the second item of the list.  It
         contains another list nested inside it.  The inner
         list is an \emph{enumerated} list.
         \begin{enumerate}
            \item This is the first item of an enumerated
                  list that is nested within the itemized list.

            \item This is the second item of the inner list.
                  \LaTeX\ allows you to nest lists deeper than
                  you really should.
         \end{enumerate}
         This is the rest of the second item of the outer
         list.  It is no more interesting than any other
         part of the item.
   \item This is the third item of the list.
\end{itemize}
You can even display poetry.
\begin{verse}
   There is an environment
    for verse \\             % The \\ command separates lines
   Whose features some poets % within a stanza.
   will curse.

                             % One or more blank lines separate stanzas.

   For instead of making\\
   Them do \emph{all} line breaking, \\
   It allows them to put too many words on a line when they'd rather be
   forced to be terse.
\end{verse}

Mathematical formulas may also be displayed.  A
displayed formula
is
one-line long; multiline
formulas require special formatting instructions.
   \[  \ip{\Gamma}{\psi'} = x'' + y^{2} + z_{i}^{n}\]
Don't start a paragraph with a displayed equation,
nor make one a paragraph by itself.

\end{document}               % End of document.

\section{Bench Huge Section 29}
% synthetic bench section 29
% This is a sample LaTeX input file.  (Version of 12 August 2004.)
%
% A '%' character causes TeX to ignore all remaining text on the line,
% and is used for comments like this one.

\documentclass{article}      % Specifies the document class

                             % The preamble begins here.
\title{An Example Document}  % Declares the document's title.
\author{Leslie Lamport}      % Declares the author's name.
\date{January 21, 1994}      % Deleting this command produces today's date.

\newcommand{\ip}[2]{(#1, #2)}
                             % Defines \ip{arg1}{arg2} to mean
                             % (arg1, arg2).

%\newcommand{\ip}[2]{\langle #1 | #2\rangle}
                             % This is an alternative definition of
                             % \ip that is commented out.

\begin{document}             % End of preamble and beginning of text.

\maketitle                   % Produces the title.

This is an example input file.  Comparing it with
the output it generates can show you how to
produce a simple document of your own.

\section{Ordinary Text}      % Produces section heading.  Lower-level
                             % sections are begun with similar
                             % \subsection and \subsubsection commands.

The ends  of words and sentences are marked
  by   spaces. It  doesn't matter how many
spaces    you type; one is as good as 100.  The
end of   a line counts as a space.

One   or more   blank lines denote the  end
of  a paragraph.

Since any number of consecutive spaces are treated
like a single one, the formatting of the input
file makes no difference to
      \LaTeX,                % The \LaTeX command generates the LaTeX logo.
but it makes a difference to you.  When you use
\LaTeX, making your input file as easy to read
as possible will be a great help as you write
your document and when you change it.  This sample
file shows how you can add comments to your own input
file.

Because printing is different from typewriting,
there are a number of things that you have to do
differently when preparing an input file than if
you were just typing the document directly.
Quotation marks like
       ``this''
have to be handled specially, as do quotes within
quotes:
       ``\,`this'            % \, separates the double and single quote.
        is what I just
        wrote, not  `that'\,''.

Dashes come in three sizes: an
       intra-word
dash, a medium dash for number ranges like
       1--2,
and a punctuation
       dash---like
this.

A sentence-ending space should be larger than the
space between words within a sentence.  You
sometimes have to type special commands in
conjunction with punctuation characters to get
this right, as in the following sentence.
       Gnats, gnus, etc.\ all  % `\ ' makes an inter-word space.
       begin with G\@.         % \@ marks end-of-sentence punctuation.
You should check the spaces after periods when
reading your output to make sure you haven't
forgotten any special cases.  Generating an
ellipsis
       \ldots\               % `\ ' is needed after `\ldots' because TeX
                             % ignores spaces after command names like \ldots
                             % made from \ + letters.
                             %
                             % Note how a `%' character causes TeX to ignore
                             % the end of the input line, so these blank lines
                             % do not start a new paragraph.
                             %
with the right spacing around the periods requires
a special command.

\LaTeX\ interprets some common characters as
commands, so you must type special commands to
generate them.  These characters include the
following:
       \$ \& \% \# \{ and \}.

In printing, text is usually emphasized with an
       \emph{italic}
type style.

\begin{em}
   A long segment of text can also be emphasized
   in this way.  Text within such a segment can be
   given \emph{additional} emphasis.
\end{em}

It is sometimes necessary to prevent \LaTeX\ from
breaking a line where it might otherwise do so.
This may be at a space, as between the ``Mr.''\ and
``Jones'' in
       ``Mr.~Jones'',        % ~ produces an unbreakable interword space.
or within a word---especially when the word is a
symbol like
       \mbox{\emph{itemnum}}
that makes little sense when hyphenated across
lines.

Footnotes\footnote{This is an example of a footnote.}
pose no problem.

\LaTeX\ is good at typesetting mathematical formulas
like
       \( x-3y + z = 7 \)
or
       \( a_{1} > x^{2n} + y^{2n} > x' \)
or
       \( \ip{A}{B} = \sum_{i} a_{i} b_{i} \).
The spaces you type in a formula are
ignored.  Remember that a letter like
       $x$                   % $ ... $  and  \( ... \)  are equivalent
is a formula when it denotes a mathematical
symbol, and it should be typed as one.

\section{Displayed Text}

Text is displayed by indenting it from the left
margin.  Quotations are commonly displayed.  There
are short quotations
\begin{quote}
   This is a short quotation.  It consists of a
   single paragraph of text.  See how it is formatted.
\end{quote}
and longer ones.
\begin{quotation}
   This is a longer quotation.  It consists of two
   paragraphs of text, neither of which are
   particularly interesting.

   This is the second paragraph of the quotation.  It
   is just as dull as the first paragraph.
\end{quotation}
Another frequently-displayed structure is a list.
The following is an example of an \emph{itemized}
list.
\begin{itemize}
   \item This is the first item of an itemized list.
         Each item in the list is marked with a ``tick''.
         You don't have to worry about what kind of tick
         mark is used.

   \item This is the second item of the list.  It
         contains another list nested inside it.  The inner
         list is an \emph{enumerated} list.
         \begin{enumerate}
            \item This is the first item of an enumerated
                  list that is nested within the itemized list.

            \item This is the second item of the inner list.
                  \LaTeX\ allows you to nest lists deeper than
                  you really should.
         \end{enumerate}
         This is the rest of the second item of the outer
         list.  It is no more interesting than any other
         part of the item.
   \item This is the third item of the list.
\end{itemize}
You can even display poetry.
\begin{verse}
   There is an environment
    for verse \\             % The \\ command separates lines
   Whose features some poets % within a stanza.
   will curse.

                             % One or more blank lines separate stanzas.

   For instead of making\\
   Them do \emph{all} line breaking, \\
   It allows them to put too many words on a line when they'd rather be
   forced to be terse.
\end{verse}

Mathematical formulas may also be displayed.  A
displayed formula
is
one-line long; multiline
formulas require special formatting instructions.
   \[  \ip{\Gamma}{\psi'} = x'' + y^{2} + z_{i}^{n}\]
Don't start a paragraph with a displayed equation,
nor make one a paragraph by itself.

\end{document}               % End of document.

\section{Bench Huge Section 30}
% synthetic bench section 30
% This is a sample LaTeX input file.  (Version of 12 August 2004.)
%
% A '%' character causes TeX to ignore all remaining text on the line,
% and is used for comments like this one.

\documentclass{article}      % Specifies the document class

                             % The preamble begins here.
\title{An Example Document}  % Declares the document's title.
\author{Leslie Lamport}      % Declares the author's name.
\date{January 21, 1994}      % Deleting this command produces today's date.

\newcommand{\ip}[2]{(#1, #2)}
                             % Defines \ip{arg1}{arg2} to mean
                             % (arg1, arg2).

%\newcommand{\ip}[2]{\langle #1 | #2\rangle}
                             % This is an alternative definition of
                             % \ip that is commented out.

\begin{document}             % End of preamble and beginning of text.

\maketitle                   % Produces the title.

This is an example input file.  Comparing it with
the output it generates can show you how to
produce a simple document of your own.

\section{Ordinary Text}      % Produces section heading.  Lower-level
                             % sections are begun with similar
                             % \subsection and \subsubsection commands.

The ends  of words and sentences are marked
  by   spaces. It  doesn't matter how many
spaces    you type; one is as good as 100.  The
end of   a line counts as a space.

One   or more   blank lines denote the  end
of  a paragraph.

Since any number of consecutive spaces are treated
like a single one, the formatting of the input
file makes no difference to
      \LaTeX,                % The \LaTeX command generates the LaTeX logo.
but it makes a difference to you.  When you use
\LaTeX, making your input file as easy to read
as possible will be a great help as you write
your document and when you change it.  This sample
file shows how you can add comments to your own input
file.

Because printing is different from typewriting,
there are a number of things that you have to do
differently when preparing an input file than if
you were just typing the document directly.
Quotation marks like
       ``this''
have to be handled specially, as do quotes within
quotes:
       ``\,`this'            % \, separates the double and single quote.
        is what I just
        wrote, not  `that'\,''.

Dashes come in three sizes: an
       intra-word
dash, a medium dash for number ranges like
       1--2,
and a punctuation
       dash---like
this.

A sentence-ending space should be larger than the
space between words within a sentence.  You
sometimes have to type special commands in
conjunction with punctuation characters to get
this right, as in the following sentence.
       Gnats, gnus, etc.\ all  % `\ ' makes an inter-word space.
       begin with G\@.         % \@ marks end-of-sentence punctuation.
You should check the spaces after periods when
reading your output to make sure you haven't
forgotten any special cases.  Generating an
ellipsis
       \ldots\               % `\ ' is needed after `\ldots' because TeX
                             % ignores spaces after command names like \ldots
                             % made from \ + letters.
                             %
                             % Note how a `%' character causes TeX to ignore
                             % the end of the input line, so these blank lines
                             % do not start a new paragraph.
                             %
with the right spacing around the periods requires
a special command.

\LaTeX\ interprets some common characters as
commands, so you must type special commands to
generate them.  These characters include the
following:
       \$ \& \% \# \{ and \}.

In printing, text is usually emphasized with an
       \emph{italic}
type style.

\begin{em}
   A long segment of text can also be emphasized
   in this way.  Text within such a segment can be
   given \emph{additional} emphasis.
\end{em}

It is sometimes necessary to prevent \LaTeX\ from
breaking a line where it might otherwise do so.
This may be at a space, as between the ``Mr.''\ and
``Jones'' in
       ``Mr.~Jones'',        % ~ produces an unbreakable interword space.
or within a word---especially when the word is a
symbol like
       \mbox{\emph{itemnum}}
that makes little sense when hyphenated across
lines.

Footnotes\footnote{This is an example of a footnote.}
pose no problem.

\LaTeX\ is good at typesetting mathematical formulas
like
       \( x-3y + z = 7 \)
or
       \( a_{1} > x^{2n} + y^{2n} > x' \)
or
       \( \ip{A}{B} = \sum_{i} a_{i} b_{i} \).
The spaces you type in a formula are
ignored.  Remember that a letter like
       $x$                   % $ ... $  and  \( ... \)  are equivalent
is a formula when it denotes a mathematical
symbol, and it should be typed as one.

\section{Displayed Text}

Text is displayed by indenting it from the left
margin.  Quotations are commonly displayed.  There
are short quotations
\begin{quote}
   This is a short quotation.  It consists of a
   single paragraph of text.  See how it is formatted.
\end{quote}
and longer ones.
\begin{quotation}
   This is a longer quotation.  It consists of two
   paragraphs of text, neither of which are
   particularly interesting.

   This is the second paragraph of the quotation.  It
   is just as dull as the first paragraph.
\end{quotation}
Another frequently-displayed structure is a list.
The following is an example of an \emph{itemized}
list.
\begin{itemize}
   \item This is the first item of an itemized list.
         Each item in the list is marked with a ``tick''.
         You don't have to worry about what kind of tick
         mark is used.

   \item This is the second item of the list.  It
         contains another list nested inside it.  The inner
         list is an \emph{enumerated} list.
         \begin{enumerate}
            \item This is the first item of an enumerated
                  list that is nested within the itemized list.

            \item This is the second item of the inner list.
                  \LaTeX\ allows you to nest lists deeper than
                  you really should.
         \end{enumerate}
         This is the rest of the second item of the outer
         list.  It is no more interesting than any other
         part of the item.
   \item This is the third item of the list.
\end{itemize}
You can even display poetry.
\begin{verse}
   There is an environment
    for verse \\             % The \\ command separates lines
   Whose features some poets % within a stanza.
   will curse.

                             % One or more blank lines separate stanzas.

   For instead of making\\
   Them do \emph{all} line breaking, \\
   It allows them to put too many words on a line when they'd rather be
   forced to be terse.
\end{verse}

Mathematical formulas may also be displayed.  A
displayed formula
is
one-line long; multiline
formulas require special formatting instructions.
   \[  \ip{\Gamma}{\psi'} = x'' + y^{2} + z_{i}^{n}\]
Don't start a paragraph with a displayed equation,
nor make one a paragraph by itself.

\end{document}               % End of document.

\section{Bench Huge Section 31}
% synthetic bench section 31
% This is a sample LaTeX input file.  (Version of 12 August 2004.)
%
% A '%' character causes TeX to ignore all remaining text on the line,
% and is used for comments like this one.

\documentclass{article}      % Specifies the document class

                             % The preamble begins here.
\title{An Example Document}  % Declares the document's title.
\author{Leslie Lamport}      % Declares the author's name.
\date{January 21, 1994}      % Deleting this command produces today's date.

\newcommand{\ip}[2]{(#1, #2)}
                             % Defines \ip{arg1}{arg2} to mean
                             % (arg1, arg2).

%\newcommand{\ip}[2]{\langle #1 | #2\rangle}
                             % This is an alternative definition of
                             % \ip that is commented out.

\begin{document}             % End of preamble and beginning of text.

\maketitle                   % Produces the title.

This is an example input file.  Comparing it with
the output it generates can show you how to
produce a simple document of your own.

\section{Ordinary Text}      % Produces section heading.  Lower-level
                             % sections are begun with similar
                             % \subsection and \subsubsection commands.

The ends  of words and sentences are marked
  by   spaces. It  doesn't matter how many
spaces    you type; one is as good as 100.  The
end of   a line counts as a space.

One   or more   blank lines denote the  end
of  a paragraph.

Since any number of consecutive spaces are treated
like a single one, the formatting of the input
file makes no difference to
      \LaTeX,                % The \LaTeX command generates the LaTeX logo.
but it makes a difference to you.  When you use
\LaTeX, making your input file as easy to read
as possible will be a great help as you write
your document and when you change it.  This sample
file shows how you can add comments to your own input
file.

Because printing is different from typewriting,
there are a number of things that you have to do
differently when preparing an input file than if
you were just typing the document directly.
Quotation marks like
       ``this''
have to be handled specially, as do quotes within
quotes:
       ``\,`this'            % \, separates the double and single quote.
        is what I just
        wrote, not  `that'\,''.

Dashes come in three sizes: an
       intra-word
dash, a medium dash for number ranges like
       1--2,
and a punctuation
       dash---like
this.

A sentence-ending space should be larger than the
space between words within a sentence.  You
sometimes have to type special commands in
conjunction with punctuation characters to get
this right, as in the following sentence.
       Gnats, gnus, etc.\ all  % `\ ' makes an inter-word space.
       begin with G\@.         % \@ marks end-of-sentence punctuation.
You should check the spaces after periods when
reading your output to make sure you haven't
forgotten any special cases.  Generating an
ellipsis
       \ldots\               % `\ ' is needed after `\ldots' because TeX
                             % ignores spaces after command names like \ldots
                             % made from \ + letters.
                             %
                             % Note how a `%' character causes TeX to ignore
                             % the end of the input line, so these blank lines
                             % do not start a new paragraph.
                             %
with the right spacing around the periods requires
a special command.

\LaTeX\ interprets some common characters as
commands, so you must type special commands to
generate them.  These characters include the
following:
       \$ \& \% \# \{ and \}.

In printing, text is usually emphasized with an
       \emph{italic}
type style.

\begin{em}
   A long segment of text can also be emphasized
   in this way.  Text within such a segment can be
   given \emph{additional} emphasis.
\end{em}

It is sometimes necessary to prevent \LaTeX\ from
breaking a line where it might otherwise do so.
This may be at a space, as between the ``Mr.''\ and
``Jones'' in
       ``Mr.~Jones'',        % ~ produces an unbreakable interword space.
or within a word---especially when the word is a
symbol like
       \mbox{\emph{itemnum}}
that makes little sense when hyphenated across
lines.

Footnotes\footnote{This is an example of a footnote.}
pose no problem.

\LaTeX\ is good at typesetting mathematical formulas
like
       \( x-3y + z = 7 \)
or
       \( a_{1} > x^{2n} + y^{2n} > x' \)
or
       \( \ip{A}{B} = \sum_{i} a_{i} b_{i} \).
The spaces you type in a formula are
ignored.  Remember that a letter like
       $x$                   % $ ... $  and  \( ... \)  are equivalent
is a formula when it denotes a mathematical
symbol, and it should be typed as one.

\section{Displayed Text}

Text is displayed by indenting it from the left
margin.  Quotations are commonly displayed.  There
are short quotations
\begin{quote}
   This is a short quotation.  It consists of a
   single paragraph of text.  See how it is formatted.
\end{quote}
and longer ones.
\begin{quotation}
   This is a longer quotation.  It consists of two
   paragraphs of text, neither of which are
   particularly interesting.

   This is the second paragraph of the quotation.  It
   is just as dull as the first paragraph.
\end{quotation}
Another frequently-displayed structure is a list.
The following is an example of an \emph{itemized}
list.
\begin{itemize}
   \item This is the first item of an itemized list.
         Each item in the list is marked with a ``tick''.
         You don't have to worry about what kind of tick
         mark is used.

   \item This is the second item of the list.  It
         contains another list nested inside it.  The inner
         list is an \emph{enumerated} list.
         \begin{enumerate}
            \item This is the first item of an enumerated
                  list that is nested within the itemized list.

            \item This is the second item of the inner list.
                  \LaTeX\ allows you to nest lists deeper than
                  you really should.
         \end{enumerate}
         This is the rest of the second item of the outer
         list.  It is no more interesting than any other
         part of the item.
   \item This is the third item of the list.
\end{itemize}
You can even display poetry.
\begin{verse}
   There is an environment
    for verse \\             % The \\ command separates lines
   Whose features some poets % within a stanza.
   will curse.

                             % One or more blank lines separate stanzas.

   For instead of making\\
   Them do \emph{all} line breaking, \\
   It allows them to put too many words on a line when they'd rather be
   forced to be terse.
\end{verse}

Mathematical formulas may also be displayed.  A
displayed formula
is
one-line long; multiline
formulas require special formatting instructions.
   \[  \ip{\Gamma}{\psi'} = x'' + y^{2} + z_{i}^{n}\]
Don't start a paragraph with a displayed equation,
nor make one a paragraph by itself.

\end{document}               % End of document.

\section{Bench Huge Section 32}
% synthetic bench section 32
% This is a sample LaTeX input file.  (Version of 12 August 2004.)
%
% A '%' character causes TeX to ignore all remaining text on the line,
% and is used for comments like this one.

\documentclass{article}      % Specifies the document class

                             % The preamble begins here.
\title{An Example Document}  % Declares the document's title.
\author{Leslie Lamport}      % Declares the author's name.
\date{January 21, 1994}      % Deleting this command produces today's date.

\newcommand{\ip}[2]{(#1, #2)}
                             % Defines \ip{arg1}{arg2} to mean
                             % (arg1, arg2).

%\newcommand{\ip}[2]{\langle #1 | #2\rangle}
                             % This is an alternative definition of
                             % \ip that is commented out.

\begin{document}             % End of preamble and beginning of text.

\maketitle                   % Produces the title.

This is an example input file.  Comparing it with
the output it generates can show you how to
produce a simple document of your own.

\section{Ordinary Text}      % Produces section heading.  Lower-level
                             % sections are begun with similar
                             % \subsection and \subsubsection commands.

The ends  of words and sentences are marked
  by   spaces. It  doesn't matter how many
spaces    you type; one is as good as 100.  The
end of   a line counts as a space.

One   or more   blank lines denote the  end
of  a paragraph.

Since any number of consecutive spaces are treated
like a single one, the formatting of the input
file makes no difference to
      \LaTeX,                % The \LaTeX command generates the LaTeX logo.
but it makes a difference to you.  When you use
\LaTeX, making your input file as easy to read
as possible will be a great help as you write
your document and when you change it.  This sample
file shows how you can add comments to your own input
file.

Because printing is different from typewriting,
there are a number of things that you have to do
differently when preparing an input file than if
you were just typing the document directly.
Quotation marks like
       ``this''
have to be handled specially, as do quotes within
quotes:
       ``\,`this'            % \, separates the double and single quote.
        is what I just
        wrote, not  `that'\,''.

Dashes come in three sizes: an
       intra-word
dash, a medium dash for number ranges like
       1--2,
and a punctuation
       dash---like
this.

A sentence-ending space should be larger than the
space between words within a sentence.  You
sometimes have to type special commands in
conjunction with punctuation characters to get
this right, as in the following sentence.
       Gnats, gnus, etc.\ all  % `\ ' makes an inter-word space.
       begin with G\@.         % \@ marks end-of-sentence punctuation.
You should check the spaces after periods when
reading your output to make sure you haven't
forgotten any special cases.  Generating an
ellipsis
       \ldots\               % `\ ' is needed after `\ldots' because TeX
                             % ignores spaces after command names like \ldots
                             % made from \ + letters.
                             %
                             % Note how a `%' character causes TeX to ignore
                             % the end of the input line, so these blank lines
                             % do not start a new paragraph.
                             %
with the right spacing around the periods requires
a special command.

\LaTeX\ interprets some common characters as
commands, so you must type special commands to
generate them.  These characters include the
following:
       \$ \& \% \# \{ and \}.

In printing, text is usually emphasized with an
       \emph{italic}
type style.

\begin{em}
   A long segment of text can also be emphasized
   in this way.  Text within such a segment can be
   given \emph{additional} emphasis.
\end{em}

It is sometimes necessary to prevent \LaTeX\ from
breaking a line where it might otherwise do so.
This may be at a space, as between the ``Mr.''\ and
``Jones'' in
       ``Mr.~Jones'',        % ~ produces an unbreakable interword space.
or within a word---especially when the word is a
symbol like
       \mbox{\emph{itemnum}}
that makes little sense when hyphenated across
lines.

Footnotes\footnote{This is an example of a footnote.}
pose no problem.

\LaTeX\ is good at typesetting mathematical formulas
like
       \( x-3y + z = 7 \)
or
       \( a_{1} > x^{2n} + y^{2n} > x' \)
or
       \( \ip{A}{B} = \sum_{i} a_{i} b_{i} \).
The spaces you type in a formula are
ignored.  Remember that a letter like
       $x$                   % $ ... $  and  \( ... \)  are equivalent
is a formula when it denotes a mathematical
symbol, and it should be typed as one.

\section{Displayed Text}

Text is displayed by indenting it from the left
margin.  Quotations are commonly displayed.  There
are short quotations
\begin{quote}
   This is a short quotation.  It consists of a
   single paragraph of text.  See how it is formatted.
\end{quote}
and longer ones.
\begin{quotation}
   This is a longer quotation.  It consists of two
   paragraphs of text, neither of which are
   particularly interesting.

   This is the second paragraph of the quotation.  It
   is just as dull as the first paragraph.
\end{quotation}
Another frequently-displayed structure is a list.
The following is an example of an \emph{itemized}
list.
\begin{itemize}
   \item This is the first item of an itemized list.
         Each item in the list is marked with a ``tick''.
         You don't have to worry about what kind of tick
         mark is used.

   \item This is the second item of the list.  It
         contains another list nested inside it.  The inner
         list is an \emph{enumerated} list.
         \begin{enumerate}
            \item This is the first item of an enumerated
                  list that is nested within the itemized list.

            \item This is the second item of the inner list.
                  \LaTeX\ allows you to nest lists deeper than
                  you really should.
         \end{enumerate}
         This is the rest of the second item of the outer
         list.  It is no more interesting than any other
         part of the item.
   \item This is the third item of the list.
\end{itemize}
You can even display poetry.
\begin{verse}
   There is an environment
    for verse \\             % The \\ command separates lines
   Whose features some poets % within a stanza.
   will curse.

                             % One or more blank lines separate stanzas.

   For instead of making\\
   Them do \emph{all} line breaking, \\
   It allows them to put too many words on a line when they'd rather be
   forced to be terse.
\end{verse}

Mathematical formulas may also be displayed.  A
displayed formula
is
one-line long; multiline
formulas require special formatting instructions.
   \[  \ip{\Gamma}{\psi'} = x'' + y^{2} + z_{i}^{n}\]
Don't start a paragraph with a displayed equation,
nor make one a paragraph by itself.

\end{document}               % End of document.

\section{Bench Huge Section 33}
% synthetic bench section 33
% This is a sample LaTeX input file.  (Version of 12 August 2004.)
%
% A '%' character causes TeX to ignore all remaining text on the line,
% and is used for comments like this one.

\documentclass{article}      % Specifies the document class

                             % The preamble begins here.
\title{An Example Document}  % Declares the document's title.
\author{Leslie Lamport}      % Declares the author's name.
\date{January 21, 1994}      % Deleting this command produces today's date.

\newcommand{\ip}[2]{(#1, #2)}
                             % Defines \ip{arg1}{arg2} to mean
                             % (arg1, arg2).

%\newcommand{\ip}[2]{\langle #1 | #2\rangle}
                             % This is an alternative definition of
                             % \ip that is commented out.

\begin{document}             % End of preamble and beginning of text.

\maketitle                   % Produces the title.

This is an example input file.  Comparing it with
the output it generates can show you how to
produce a simple document of your own.

\section{Ordinary Text}      % Produces section heading.  Lower-level
                             % sections are begun with similar
                             % \subsection and \subsubsection commands.

The ends  of words and sentences are marked
  by   spaces. It  doesn't matter how many
spaces    you type; one is as good as 100.  The
end of   a line counts as a space.

One   or more   blank lines denote the  end
of  a paragraph.

Since any number of consecutive spaces are treated
like a single one, the formatting of the input
file makes no difference to
      \LaTeX,                % The \LaTeX command generates the LaTeX logo.
but it makes a difference to you.  When you use
\LaTeX, making your input file as easy to read
as possible will be a great help as you write
your document and when you change it.  This sample
file shows how you can add comments to your own input
file.

Because printing is different from typewriting,
there are a number of things that you have to do
differently when preparing an input file than if
you were just typing the document directly.
Quotation marks like
       ``this''
have to be handled specially, as do quotes within
quotes:
       ``\,`this'            % \, separates the double and single quote.
        is what I just
        wrote, not  `that'\,''.

Dashes come in three sizes: an
       intra-word
dash, a medium dash for number ranges like
       1--2,
and a punctuation
       dash---like
this.

A sentence-ending space should be larger than the
space between words within a sentence.  You
sometimes have to type special commands in
conjunction with punctuation characters to get
this right, as in the following sentence.
       Gnats, gnus, etc.\ all  % `\ ' makes an inter-word space.
       begin with G\@.         % \@ marks end-of-sentence punctuation.
You should check the spaces after periods when
reading your output to make sure you haven't
forgotten any special cases.  Generating an
ellipsis
       \ldots\               % `\ ' is needed after `\ldots' because TeX
                             % ignores spaces after command names like \ldots
                             % made from \ + letters.
                             %
                             % Note how a `%' character causes TeX to ignore
                             % the end of the input line, so these blank lines
                             % do not start a new paragraph.
                             %
with the right spacing around the periods requires
a special command.

\LaTeX\ interprets some common characters as
commands, so you must type special commands to
generate them.  These characters include the
following:
       \$ \& \% \# \{ and \}.

In printing, text is usually emphasized with an
       \emph{italic}
type style.

\begin{em}
   A long segment of text can also be emphasized
   in this way.  Text within such a segment can be
   given \emph{additional} emphasis.
\end{em}

It is sometimes necessary to prevent \LaTeX\ from
breaking a line where it might otherwise do so.
This may be at a space, as between the ``Mr.''\ and
``Jones'' in
       ``Mr.~Jones'',        % ~ produces an unbreakable interword space.
or within a word---especially when the word is a
symbol like
       \mbox{\emph{itemnum}}
that makes little sense when hyphenated across
lines.

Footnotes\footnote{This is an example of a footnote.}
pose no problem.

\LaTeX\ is good at typesetting mathematical formulas
like
       \( x-3y + z = 7 \)
or
       \( a_{1} > x^{2n} + y^{2n} > x' \)
or
       \( \ip{A}{B} = \sum_{i} a_{i} b_{i} \).
The spaces you type in a formula are
ignored.  Remember that a letter like
       $x$                   % $ ... $  and  \( ... \)  are equivalent
is a formula when it denotes a mathematical
symbol, and it should be typed as one.

\section{Displayed Text}

Text is displayed by indenting it from the left
margin.  Quotations are commonly displayed.  There
are short quotations
\begin{quote}
   This is a short quotation.  It consists of a
   single paragraph of text.  See how it is formatted.
\end{quote}
and longer ones.
\begin{quotation}
   This is a longer quotation.  It consists of two
   paragraphs of text, neither of which are
   particularly interesting.

   This is the second paragraph of the quotation.  It
   is just as dull as the first paragraph.
\end{quotation}
Another frequently-displayed structure is a list.
The following is an example of an \emph{itemized}
list.
\begin{itemize}
   \item This is the first item of an itemized list.
         Each item in the list is marked with a ``tick''.
         You don't have to worry about what kind of tick
         mark is used.

   \item This is the second item of the list.  It
         contains another list nested inside it.  The inner
         list is an \emph{enumerated} list.
         \begin{enumerate}
            \item This is the first item of an enumerated
                  list that is nested within the itemized list.

            \item This is the second item of the inner list.
                  \LaTeX\ allows you to nest lists deeper than
                  you really should.
         \end{enumerate}
         This is the rest of the second item of the outer
         list.  It is no more interesting than any other
         part of the item.
   \item This is the third item of the list.
\end{itemize}
You can even display poetry.
\begin{verse}
   There is an environment
    for verse \\             % The \\ command separates lines
   Whose features some poets % within a stanza.
   will curse.

                             % One or more blank lines separate stanzas.

   For instead of making\\
   Them do \emph{all} line breaking, \\
   It allows them to put too many words on a line when they'd rather be
   forced to be terse.
\end{verse}

Mathematical formulas may also be displayed.  A
displayed formula
is
one-line long; multiline
formulas require special formatting instructions.
   \[  \ip{\Gamma}{\psi'} = x'' + y^{2} + z_{i}^{n}\]
Don't start a paragraph with a displayed equation,
nor make one a paragraph by itself.

\end{document}               % End of document.

\section{Bench Huge Section 34}
% synthetic bench section 34
% This is a sample LaTeX input file.  (Version of 12 August 2004.)
%
% A '%' character causes TeX to ignore all remaining text on the line,
% and is used for comments like this one.

\documentclass{article}      % Specifies the document class

                             % The preamble begins here.
\title{An Example Document}  % Declares the document's title.
\author{Leslie Lamport}      % Declares the author's name.
\date{January 21, 1994}      % Deleting this command produces today's date.

\newcommand{\ip}[2]{(#1, #2)}
                             % Defines \ip{arg1}{arg2} to mean
                             % (arg1, arg2).

%\newcommand{\ip}[2]{\langle #1 | #2\rangle}
                             % This is an alternative definition of
                             % \ip that is commented out.

\begin{document}             % End of preamble and beginning of text.

\maketitle                   % Produces the title.

This is an example input file.  Comparing it with
the output it generates can show you how to
produce a simple document of your own.

\section{Ordinary Text}      % Produces section heading.  Lower-level
                             % sections are begun with similar
                             % \subsection and \subsubsection commands.

The ends  of words and sentences are marked
  by   spaces. It  doesn't matter how many
spaces    you type; one is as good as 100.  The
end of   a line counts as a space.

One   or more   blank lines denote the  end
of  a paragraph.

Since any number of consecutive spaces are treated
like a single one, the formatting of the input
file makes no difference to
      \LaTeX,                % The \LaTeX command generates the LaTeX logo.
but it makes a difference to you.  When you use
\LaTeX, making your input file as easy to read
as possible will be a great help as you write
your document and when you change it.  This sample
file shows how you can add comments to your own input
file.

Because printing is different from typewriting,
there are a number of things that you have to do
differently when preparing an input file than if
you were just typing the document directly.
Quotation marks like
       ``this''
have to be handled specially, as do quotes within
quotes:
       ``\,`this'            % \, separates the double and single quote.
        is what I just
        wrote, not  `that'\,''.

Dashes come in three sizes: an
       intra-word
dash, a medium dash for number ranges like
       1--2,
and a punctuation
       dash---like
this.

A sentence-ending space should be larger than the
space between words within a sentence.  You
sometimes have to type special commands in
conjunction with punctuation characters to get
this right, as in the following sentence.
       Gnats, gnus, etc.\ all  % `\ ' makes an inter-word space.
       begin with G\@.         % \@ marks end-of-sentence punctuation.
You should check the spaces after periods when
reading your output to make sure you haven't
forgotten any special cases.  Generating an
ellipsis
       \ldots\               % `\ ' is needed after `\ldots' because TeX
                             % ignores spaces after command names like \ldots
                             % made from \ + letters.
                             %
                             % Note how a `%' character causes TeX to ignore
                             % the end of the input line, so these blank lines
                             % do not start a new paragraph.
                             %
with the right spacing around the periods requires
a special command.

\LaTeX\ interprets some common characters as
commands, so you must type special commands to
generate them.  These characters include the
following:
       \$ \& \% \# \{ and \}.

In printing, text is usually emphasized with an
       \emph{italic}
type style.

\begin{em}
   A long segment of text can also be emphasized
   in this way.  Text within such a segment can be
   given \emph{additional} emphasis.
\end{em}

It is sometimes necessary to prevent \LaTeX\ from
breaking a line where it might otherwise do so.
This may be at a space, as between the ``Mr.''\ and
``Jones'' in
       ``Mr.~Jones'',        % ~ produces an unbreakable interword space.
or within a word---especially when the word is a
symbol like
       \mbox{\emph{itemnum}}
that makes little sense when hyphenated across
lines.

Footnotes\footnote{This is an example of a footnote.}
pose no problem.

\LaTeX\ is good at typesetting mathematical formulas
like
       \( x-3y + z = 7 \)
or
       \( a_{1} > x^{2n} + y^{2n} > x' \)
or
       \( \ip{A}{B} = \sum_{i} a_{i} b_{i} \).
The spaces you type in a formula are
ignored.  Remember that a letter like
       $x$                   % $ ... $  and  \( ... \)  are equivalent
is a formula when it denotes a mathematical
symbol, and it should be typed as one.

\section{Displayed Text}

Text is displayed by indenting it from the left
margin.  Quotations are commonly displayed.  There
are short quotations
\begin{quote}
   This is a short quotation.  It consists of a
   single paragraph of text.  See how it is formatted.
\end{quote}
and longer ones.
\begin{quotation}
   This is a longer quotation.  It consists of two
   paragraphs of text, neither of which are
   particularly interesting.

   This is the second paragraph of the quotation.  It
   is just as dull as the first paragraph.
\end{quotation}
Another frequently-displayed structure is a list.
The following is an example of an \emph{itemized}
list.
\begin{itemize}
   \item This is the first item of an itemized list.
         Each item in the list is marked with a ``tick''.
         You don't have to worry about what kind of tick
         mark is used.

   \item This is the second item of the list.  It
         contains another list nested inside it.  The inner
         list is an \emph{enumerated} list.
         \begin{enumerate}
            \item This is the first item of an enumerated
                  list that is nested within the itemized list.

            \item This is the second item of the inner list.
                  \LaTeX\ allows you to nest lists deeper than
                  you really should.
         \end{enumerate}
         This is the rest of the second item of the outer
         list.  It is no more interesting than any other
         part of the item.
   \item This is the third item of the list.
\end{itemize}
You can even display poetry.
\begin{verse}
   There is an environment
    for verse \\             % The \\ command separates lines
   Whose features some poets % within a stanza.
   will curse.

                             % One or more blank lines separate stanzas.

   For instead of making\\
   Them do \emph{all} line breaking, \\
   It allows them to put too many words on a line when they'd rather be
   forced to be terse.
\end{verse}

Mathematical formulas may also be displayed.  A
displayed formula
is
one-line long; multiline
formulas require special formatting instructions.
   \[  \ip{\Gamma}{\psi'} = x'' + y^{2} + z_{i}^{n}\]
Don't start a paragraph with a displayed equation,
nor make one a paragraph by itself.

\end{document}               % End of document.

\section{Bench Huge Section 35}
% synthetic bench section 35
% This is a sample LaTeX input file.  (Version of 12 August 2004.)
%
% A '%' character causes TeX to ignore all remaining text on the line,
% and is used for comments like this one.

\documentclass{article}      % Specifies the document class

                             % The preamble begins here.
\title{An Example Document}  % Declares the document's title.
\author{Leslie Lamport}      % Declares the author's name.
\date{January 21, 1994}      % Deleting this command produces today's date.

\newcommand{\ip}[2]{(#1, #2)}
                             % Defines \ip{arg1}{arg2} to mean
                             % (arg1, arg2).

%\newcommand{\ip}[2]{\langle #1 | #2\rangle}
                             % This is an alternative definition of
                             % \ip that is commented out.

\begin{document}             % End of preamble and beginning of text.

\maketitle                   % Produces the title.

This is an example input file.  Comparing it with
the output it generates can show you how to
produce a simple document of your own.

\section{Ordinary Text}      % Produces section heading.  Lower-level
                             % sections are begun with similar
                             % \subsection and \subsubsection commands.

The ends  of words and sentences are marked
  by   spaces. It  doesn't matter how many
spaces    you type; one is as good as 100.  The
end of   a line counts as a space.

One   or more   blank lines denote the  end
of  a paragraph.

Since any number of consecutive spaces are treated
like a single one, the formatting of the input
file makes no difference to
      \LaTeX,                % The \LaTeX command generates the LaTeX logo.
but it makes a difference to you.  When you use
\LaTeX, making your input file as easy to read
as possible will be a great help as you write
your document and when you change it.  This sample
file shows how you can add comments to your own input
file.

Because printing is different from typewriting,
there are a number of things that you have to do
differently when preparing an input file than if
you were just typing the document directly.
Quotation marks like
       ``this''
have to be handled specially, as do quotes within
quotes:
       ``\,`this'            % \, separates the double and single quote.
        is what I just
        wrote, not  `that'\,''.

Dashes come in three sizes: an
       intra-word
dash, a medium dash for number ranges like
       1--2,
and a punctuation
       dash---like
this.

A sentence-ending space should be larger than the
space between words within a sentence.  You
sometimes have to type special commands in
conjunction with punctuation characters to get
this right, as in the following sentence.
       Gnats, gnus, etc.\ all  % `\ ' makes an inter-word space.
       begin with G\@.         % \@ marks end-of-sentence punctuation.
You should check the spaces after periods when
reading your output to make sure you haven't
forgotten any special cases.  Generating an
ellipsis
       \ldots\               % `\ ' is needed after `\ldots' because TeX
                             % ignores spaces after command names like \ldots
                             % made from \ + letters.
                             %
                             % Note how a `%' character causes TeX to ignore
                             % the end of the input line, so these blank lines
                             % do not start a new paragraph.
                             %
with the right spacing around the periods requires
a special command.

\LaTeX\ interprets some common characters as
commands, so you must type special commands to
generate them.  These characters include the
following:
       \$ \& \% \# \{ and \}.

In printing, text is usually emphasized with an
       \emph{italic}
type style.

\begin{em}
   A long segment of text can also be emphasized
   in this way.  Text within such a segment can be
   given \emph{additional} emphasis.
\end{em}

It is sometimes necessary to prevent \LaTeX\ from
breaking a line where it might otherwise do so.
This may be at a space, as between the ``Mr.''\ and
``Jones'' in
       ``Mr.~Jones'',        % ~ produces an unbreakable interword space.
or within a word---especially when the word is a
symbol like
       \mbox{\emph{itemnum}}
that makes little sense when hyphenated across
lines.

Footnotes\footnote{This is an example of a footnote.}
pose no problem.

\LaTeX\ is good at typesetting mathematical formulas
like
       \( x-3y + z = 7 \)
or
       \( a_{1} > x^{2n} + y^{2n} > x' \)
or
       \( \ip{A}{B} = \sum_{i} a_{i} b_{i} \).
The spaces you type in a formula are
ignored.  Remember that a letter like
       $x$                   % $ ... $  and  \( ... \)  are equivalent
is a formula when it denotes a mathematical
symbol, and it should be typed as one.

\section{Displayed Text}

Text is displayed by indenting it from the left
margin.  Quotations are commonly displayed.  There
are short quotations
\begin{quote}
   This is a short quotation.  It consists of a
   single paragraph of text.  See how it is formatted.
\end{quote}
and longer ones.
\begin{quotation}
   This is a longer quotation.  It consists of two
   paragraphs of text, neither of which are
   particularly interesting.

   This is the second paragraph of the quotation.  It
   is just as dull as the first paragraph.
\end{quotation}
Another frequently-displayed structure is a list.
The following is an example of an \emph{itemized}
list.
\begin{itemize}
   \item This is the first item of an itemized list.
         Each item in the list is marked with a ``tick''.
         You don't have to worry about what kind of tick
         mark is used.

   \item This is the second item of the list.  It
         contains another list nested inside it.  The inner
         list is an \emph{enumerated} list.
         \begin{enumerate}
            \item This is the first item of an enumerated
                  list that is nested within the itemized list.

            \item This is the second item of the inner list.
                  \LaTeX\ allows you to nest lists deeper than
                  you really should.
         \end{enumerate}
         This is the rest of the second item of the outer
         list.  It is no more interesting than any other
         part of the item.
   \item This is the third item of the list.
\end{itemize}
You can even display poetry.
\begin{verse}
   There is an environment
    for verse \\             % The \\ command separates lines
   Whose features some poets % within a stanza.
   will curse.

                             % One or more blank lines separate stanzas.

   For instead of making\\
   Them do \emph{all} line breaking, \\
   It allows them to put too many words on a line when they'd rather be
   forced to be terse.
\end{verse}

Mathematical formulas may also be displayed.  A
displayed formula
is
one-line long; multiline
formulas require special formatting instructions.
   \[  \ip{\Gamma}{\psi'} = x'' + y^{2} + z_{i}^{n}\]
Don't start a paragraph with a displayed equation,
nor make one a paragraph by itself.

\end{document}               % End of document.

\section{Bench Huge Section 36}
% synthetic bench section 36
% This is a sample LaTeX input file.  (Version of 12 August 2004.)
%
% A '%' character causes TeX to ignore all remaining text on the line,
% and is used for comments like this one.

\documentclass{article}      % Specifies the document class

                             % The preamble begins here.
\title{An Example Document}  % Declares the document's title.
\author{Leslie Lamport}      % Declares the author's name.
\date{January 21, 1994}      % Deleting this command produces today's date.

\newcommand{\ip}[2]{(#1, #2)}
                             % Defines \ip{arg1}{arg2} to mean
                             % (arg1, arg2).

%\newcommand{\ip}[2]{\langle #1 | #2\rangle}
                             % This is an alternative definition of
                             % \ip that is commented out.

\begin{document}             % End of preamble and beginning of text.

\maketitle                   % Produces the title.

This is an example input file.  Comparing it with
the output it generates can show you how to
produce a simple document of your own.

\section{Ordinary Text}      % Produces section heading.  Lower-level
                             % sections are begun with similar
                             % \subsection and \subsubsection commands.

The ends  of words and sentences are marked
  by   spaces. It  doesn't matter how many
spaces    you type; one is as good as 100.  The
end of   a line counts as a space.

One   or more   blank lines denote the  end
of  a paragraph.

Since any number of consecutive spaces are treated
like a single one, the formatting of the input
file makes no difference to
      \LaTeX,                % The \LaTeX command generates the LaTeX logo.
but it makes a difference to you.  When you use
\LaTeX, making your input file as easy to read
as possible will be a great help as you write
your document and when you change it.  This sample
file shows how you can add comments to your own input
file.

Because printing is different from typewriting,
there are a number of things that you have to do
differently when preparing an input file than if
you were just typing the document directly.
Quotation marks like
       ``this''
have to be handled specially, as do quotes within
quotes:
       ``\,`this'            % \, separates the double and single quote.
        is what I just
        wrote, not  `that'\,''.

Dashes come in three sizes: an
       intra-word
dash, a medium dash for number ranges like
       1--2,
and a punctuation
       dash---like
this.

A sentence-ending space should be larger than the
space between words within a sentence.  You
sometimes have to type special commands in
conjunction with punctuation characters to get
this right, as in the following sentence.
       Gnats, gnus, etc.\ all  % `\ ' makes an inter-word space.
       begin with G\@.         % \@ marks end-of-sentence punctuation.
You should check the spaces after periods when
reading your output to make sure you haven't
forgotten any special cases.  Generating an
ellipsis
       \ldots\               % `\ ' is needed after `\ldots' because TeX
                             % ignores spaces after command names like \ldots
                             % made from \ + letters.
                             %
                             % Note how a `%' character causes TeX to ignore
                             % the end of the input line, so these blank lines
                             % do not start a new paragraph.
                             %
with the right spacing around the periods requires
a special command.

\LaTeX\ interprets some common characters as
commands, so you must type special commands to
generate them.  These characters include the
following:
       \$ \& \% \# \{ and \}.

In printing, text is usually emphasized with an
       \emph{italic}
type style.

\begin{em}
   A long segment of text can also be emphasized
   in this way.  Text within such a segment can be
   given \emph{additional} emphasis.
\end{em}

It is sometimes necessary to prevent \LaTeX\ from
breaking a line where it might otherwise do so.
This may be at a space, as between the ``Mr.''\ and
``Jones'' in
       ``Mr.~Jones'',        % ~ produces an unbreakable interword space.
or within a word---especially when the word is a
symbol like
       \mbox{\emph{itemnum}}
that makes little sense when hyphenated across
lines.

Footnotes\footnote{This is an example of a footnote.}
pose no problem.

\LaTeX\ is good at typesetting mathematical formulas
like
       \( x-3y + z = 7 \)
or
       \( a_{1} > x^{2n} + y^{2n} > x' \)
or
       \( \ip{A}{B} = \sum_{i} a_{i} b_{i} \).
The spaces you type in a formula are
ignored.  Remember that a letter like
       $x$                   % $ ... $  and  \( ... \)  are equivalent
is a formula when it denotes a mathematical
symbol, and it should be typed as one.

\section{Displayed Text}

Text is displayed by indenting it from the left
margin.  Quotations are commonly displayed.  There
are short quotations
\begin{quote}
   This is a short quotation.  It consists of a
   single paragraph of text.  See how it is formatted.
\end{quote}
and longer ones.
\begin{quotation}
   This is a longer quotation.  It consists of two
   paragraphs of text, neither of which are
   particularly interesting.

   This is the second paragraph of the quotation.  It
   is just as dull as the first paragraph.
\end{quotation}
Another frequently-displayed structure is a list.
The following is an example of an \emph{itemized}
list.
\begin{itemize}
   \item This is the first item of an itemized list.
         Each item in the list is marked with a ``tick''.
         You don't have to worry about what kind of tick
         mark is used.

   \item This is the second item of the list.  It
         contains another list nested inside it.  The inner
         list is an \emph{enumerated} list.
         \begin{enumerate}
            \item This is the first item of an enumerated
                  list that is nested within the itemized list.

            \item This is the second item of the inner list.
                  \LaTeX\ allows you to nest lists deeper than
                  you really should.
         \end{enumerate}
         This is the rest of the second item of the outer
         list.  It is no more interesting than any other
         part of the item.
   \item This is the third item of the list.
\end{itemize}
You can even display poetry.
\begin{verse}
   There is an environment
    for verse \\             % The \\ command separates lines
   Whose features some poets % within a stanza.
   will curse.

                             % One or more blank lines separate stanzas.

   For instead of making\\
   Them do \emph{all} line breaking, \\
   It allows them to put too many words on a line when they'd rather be
   forced to be terse.
\end{verse}

Mathematical formulas may also be displayed.  A
displayed formula
is
one-line long; multiline
formulas require special formatting instructions.
   \[  \ip{\Gamma}{\psi'} = x'' + y^{2} + z_{i}^{n}\]
Don't start a paragraph with a displayed equation,
nor make one a paragraph by itself.

\end{document}               % End of document.

\section{Bench Huge Section 37}
% synthetic bench section 37
% This is a sample LaTeX input file.  (Version of 12 August 2004.)
%
% A '%' character causes TeX to ignore all remaining text on the line,
% and is used for comments like this one.

\documentclass{article}      % Specifies the document class

                             % The preamble begins here.
\title{An Example Document}  % Declares the document's title.
\author{Leslie Lamport}      % Declares the author's name.
\date{January 21, 1994}      % Deleting this command produces today's date.

\newcommand{\ip}[2]{(#1, #2)}
                             % Defines \ip{arg1}{arg2} to mean
                             % (arg1, arg2).

%\newcommand{\ip}[2]{\langle #1 | #2\rangle}
                             % This is an alternative definition of
                             % \ip that is commented out.

\begin{document}             % End of preamble and beginning of text.

\maketitle                   % Produces the title.

This is an example input file.  Comparing it with
the output it generates can show you how to
produce a simple document of your own.

\section{Ordinary Text}      % Produces section heading.  Lower-level
                             % sections are begun with similar
                             % \subsection and \subsubsection commands.

The ends  of words and sentences are marked
  by   spaces. It  doesn't matter how many
spaces    you type; one is as good as 100.  The
end of   a line counts as a space.

One   or more   blank lines denote the  end
of  a paragraph.

Since any number of consecutive spaces are treated
like a single one, the formatting of the input
file makes no difference to
      \LaTeX,                % The \LaTeX command generates the LaTeX logo.
but it makes a difference to you.  When you use
\LaTeX, making your input file as easy to read
as possible will be a great help as you write
your document and when you change it.  This sample
file shows how you can add comments to your own input
file.

Because printing is different from typewriting,
there are a number of things that you have to do
differently when preparing an input file than if
you were just typing the document directly.
Quotation marks like
       ``this''
have to be handled specially, as do quotes within
quotes:
       ``\,`this'            % \, separates the double and single quote.
        is what I just
        wrote, not  `that'\,''.

Dashes come in three sizes: an
       intra-word
dash, a medium dash for number ranges like
       1--2,
and a punctuation
       dash---like
this.

A sentence-ending space should be larger than the
space between words within a sentence.  You
sometimes have to type special commands in
conjunction with punctuation characters to get
this right, as in the following sentence.
       Gnats, gnus, etc.\ all  % `\ ' makes an inter-word space.
       begin with G\@.         % \@ marks end-of-sentence punctuation.
You should check the spaces after periods when
reading your output to make sure you haven't
forgotten any special cases.  Generating an
ellipsis
       \ldots\               % `\ ' is needed after `\ldots' because TeX
                             % ignores spaces after command names like \ldots
                             % made from \ + letters.
                             %
                             % Note how a `%' character causes TeX to ignore
                             % the end of the input line, so these blank lines
                             % do not start a new paragraph.
                             %
with the right spacing around the periods requires
a special command.

\LaTeX\ interprets some common characters as
commands, so you must type special commands to
generate them.  These characters include the
following:
       \$ \& \% \# \{ and \}.

In printing, text is usually emphasized with an
       \emph{italic}
type style.

\begin{em}
   A long segment of text can also be emphasized
   in this way.  Text within such a segment can be
   given \emph{additional} emphasis.
\end{em}

It is sometimes necessary to prevent \LaTeX\ from
breaking a line where it might otherwise do so.
This may be at a space, as between the ``Mr.''\ and
``Jones'' in
       ``Mr.~Jones'',        % ~ produces an unbreakable interword space.
or within a word---especially when the word is a
symbol like
       \mbox{\emph{itemnum}}
that makes little sense when hyphenated across
lines.

Footnotes\footnote{This is an example of a footnote.}
pose no problem.

\LaTeX\ is good at typesetting mathematical formulas
like
       \( x-3y + z = 7 \)
or
       \( a_{1} > x^{2n} + y^{2n} > x' \)
or
       \( \ip{A}{B} = \sum_{i} a_{i} b_{i} \).
The spaces you type in a formula are
ignored.  Remember that a letter like
       $x$                   % $ ... $  and  \( ... \)  are equivalent
is a formula when it denotes a mathematical
symbol, and it should be typed as one.

\section{Displayed Text}

Text is displayed by indenting it from the left
margin.  Quotations are commonly displayed.  There
are short quotations
\begin{quote}
   This is a short quotation.  It consists of a
   single paragraph of text.  See how it is formatted.
\end{quote}
and longer ones.
\begin{quotation}
   This is a longer quotation.  It consists of two
   paragraphs of text, neither of which are
   particularly interesting.

   This is the second paragraph of the quotation.  It
   is just as dull as the first paragraph.
\end{quotation}
Another frequently-displayed structure is a list.
The following is an example of an \emph{itemized}
list.
\begin{itemize}
   \item This is the first item of an itemized list.
         Each item in the list is marked with a ``tick''.
         You don't have to worry about what kind of tick
         mark is used.

   \item This is the second item of the list.  It
         contains another list nested inside it.  The inner
         list is an \emph{enumerated} list.
         \begin{enumerate}
            \item This is the first item of an enumerated
                  list that is nested within the itemized list.

            \item This is the second item of the inner list.
                  \LaTeX\ allows you to nest lists deeper than
                  you really should.
         \end{enumerate}
         This is the rest of the second item of the outer
         list.  It is no more interesting than any other
         part of the item.
   \item This is the third item of the list.
\end{itemize}
You can even display poetry.
\begin{verse}
   There is an environment
    for verse \\             % The \\ command separates lines
   Whose features some poets % within a stanza.
   will curse.

                             % One or more blank lines separate stanzas.

   For instead of making\\
   Them do \emph{all} line breaking, \\
   It allows them to put too many words on a line when they'd rather be
   forced to be terse.
\end{verse}

Mathematical formulas may also be displayed.  A
displayed formula
is
one-line long; multiline
formulas require special formatting instructions.
   \[  \ip{\Gamma}{\psi'} = x'' + y^{2} + z_{i}^{n}\]
Don't start a paragraph with a displayed equation,
nor make one a paragraph by itself.

\end{document}               % End of document.

\section{Bench Huge Section 38}
% synthetic bench section 38
% This is a sample LaTeX input file.  (Version of 12 August 2004.)
%
% A '%' character causes TeX to ignore all remaining text on the line,
% and is used for comments like this one.

\documentclass{article}      % Specifies the document class

                             % The preamble begins here.
\title{An Example Document}  % Declares the document's title.
\author{Leslie Lamport}      % Declares the author's name.
\date{January 21, 1994}      % Deleting this command produces today's date.

\newcommand{\ip}[2]{(#1, #2)}
                             % Defines \ip{arg1}{arg2} to mean
                             % (arg1, arg2).

%\newcommand{\ip}[2]{\langle #1 | #2\rangle}
                             % This is an alternative definition of
                             % \ip that is commented out.

\begin{document}             % End of preamble and beginning of text.

\maketitle                   % Produces the title.

This is an example input file.  Comparing it with
the output it generates can show you how to
produce a simple document of your own.

\section{Ordinary Text}      % Produces section heading.  Lower-level
                             % sections are begun with similar
                             % \subsection and \subsubsection commands.

The ends  of words and sentences are marked
  by   spaces. It  doesn't matter how many
spaces    you type; one is as good as 100.  The
end of   a line counts as a space.

One   or more   blank lines denote the  end
of  a paragraph.

Since any number of consecutive spaces are treated
like a single one, the formatting of the input
file makes no difference to
      \LaTeX,                % The \LaTeX command generates the LaTeX logo.
but it makes a difference to you.  When you use
\LaTeX, making your input file as easy to read
as possible will be a great help as you write
your document and when you change it.  This sample
file shows how you can add comments to your own input
file.

Because printing is different from typewriting,
there are a number of things that you have to do
differently when preparing an input file than if
you were just typing the document directly.
Quotation marks like
       ``this''
have to be handled specially, as do quotes within
quotes:
       ``\,`this'            % \, separates the double and single quote.
        is what I just
        wrote, not  `that'\,''.

Dashes come in three sizes: an
       intra-word
dash, a medium dash for number ranges like
       1--2,
and a punctuation
       dash---like
this.

A sentence-ending space should be larger than the
space between words within a sentence.  You
sometimes have to type special commands in
conjunction with punctuation characters to get
this right, as in the following sentence.
       Gnats, gnus, etc.\ all  % `\ ' makes an inter-word space.
       begin with G\@.         % \@ marks end-of-sentence punctuation.
You should check the spaces after periods when
reading your output to make sure you haven't
forgotten any special cases.  Generating an
ellipsis
       \ldots\               % `\ ' is needed after `\ldots' because TeX
                             % ignores spaces after command names like \ldots
                             % made from \ + letters.
                             %
                             % Note how a `%' character causes TeX to ignore
                             % the end of the input line, so these blank lines
                             % do not start a new paragraph.
                             %
with the right spacing around the periods requires
a special command.

\LaTeX\ interprets some common characters as
commands, so you must type special commands to
generate them.  These characters include the
following:
       \$ \& \% \# \{ and \}.

In printing, text is usually emphasized with an
       \emph{italic}
type style.

\begin{em}
   A long segment of text can also be emphasized
   in this way.  Text within such a segment can be
   given \emph{additional} emphasis.
\end{em}

It is sometimes necessary to prevent \LaTeX\ from
breaking a line where it might otherwise do so.
This may be at a space, as between the ``Mr.''\ and
``Jones'' in
       ``Mr.~Jones'',        % ~ produces an unbreakable interword space.
or within a word---especially when the word is a
symbol like
       \mbox{\emph{itemnum}}
that makes little sense when hyphenated across
lines.

Footnotes\footnote{This is an example of a footnote.}
pose no problem.

\LaTeX\ is good at typesetting mathematical formulas
like
       \( x-3y + z = 7 \)
or
       \( a_{1} > x^{2n} + y^{2n} > x' \)
or
       \( \ip{A}{B} = \sum_{i} a_{i} b_{i} \).
The spaces you type in a formula are
ignored.  Remember that a letter like
       $x$                   % $ ... $  and  \( ... \)  are equivalent
is a formula when it denotes a mathematical
symbol, and it should be typed as one.

\section{Displayed Text}

Text is displayed by indenting it from the left
margin.  Quotations are commonly displayed.  There
are short quotations
\begin{quote}
   This is a short quotation.  It consists of a
   single paragraph of text.  See how it is formatted.
\end{quote}
and longer ones.
\begin{quotation}
   This is a longer quotation.  It consists of two
   paragraphs of text, neither of which are
   particularly interesting.

   This is the second paragraph of the quotation.  It
   is just as dull as the first paragraph.
\end{quotation}
Another frequently-displayed structure is a list.
The following is an example of an \emph{itemized}
list.
\begin{itemize}
   \item This is the first item of an itemized list.
         Each item in the list is marked with a ``tick''.
         You don't have to worry about what kind of tick
         mark is used.

   \item This is the second item of the list.  It
         contains another list nested inside it.  The inner
         list is an \emph{enumerated} list.
         \begin{enumerate}
            \item This is the first item of an enumerated
                  list that is nested within the itemized list.

            \item This is the second item of the inner list.
                  \LaTeX\ allows you to nest lists deeper than
                  you really should.
         \end{enumerate}
         This is the rest of the second item of the outer
         list.  It is no more interesting than any other
         part of the item.
   \item This is the third item of the list.
\end{itemize}
You can even display poetry.
\begin{verse}
   There is an environment
    for verse \\             % The \\ command separates lines
   Whose features some poets % within a stanza.
   will curse.

                             % One or more blank lines separate stanzas.

   For instead of making\\
   Them do \emph{all} line breaking, \\
   It allows them to put too many words on a line when they'd rather be
   forced to be terse.
\end{verse}

Mathematical formulas may also be displayed.  A
displayed formula
is
one-line long; multiline
formulas require special formatting instructions.
   \[  \ip{\Gamma}{\psi'} = x'' + y^{2} + z_{i}^{n}\]
Don't start a paragraph with a displayed equation,
nor make one a paragraph by itself.

\end{document}               % End of document.

\section{Bench Huge Section 39}
% synthetic bench section 39
% This is a sample LaTeX input file.  (Version of 12 August 2004.)
%
% A '%' character causes TeX to ignore all remaining text on the line,
% and is used for comments like this one.

\documentclass{article}      % Specifies the document class

                             % The preamble begins here.
\title{An Example Document}  % Declares the document's title.
\author{Leslie Lamport}      % Declares the author's name.
\date{January 21, 1994}      % Deleting this command produces today's date.

\newcommand{\ip}[2]{(#1, #2)}
                             % Defines \ip{arg1}{arg2} to mean
                             % (arg1, arg2).

%\newcommand{\ip}[2]{\langle #1 | #2\rangle}
                             % This is an alternative definition of
                             % \ip that is commented out.

\begin{document}             % End of preamble and beginning of text.

\maketitle                   % Produces the title.

This is an example input file.  Comparing it with
the output it generates can show you how to
produce a simple document of your own.

\section{Ordinary Text}      % Produces section heading.  Lower-level
                             % sections are begun with similar
                             % \subsection and \subsubsection commands.

The ends  of words and sentences are marked
  by   spaces. It  doesn't matter how many
spaces    you type; one is as good as 100.  The
end of   a line counts as a space.

One   or more   blank lines denote the  end
of  a paragraph.

Since any number of consecutive spaces are treated
like a single one, the formatting of the input
file makes no difference to
      \LaTeX,                % The \LaTeX command generates the LaTeX logo.
but it makes a difference to you.  When you use
\LaTeX, making your input file as easy to read
as possible will be a great help as you write
your document and when you change it.  This sample
file shows how you can add comments to your own input
file.

Because printing is different from typewriting,
there are a number of things that you have to do
differently when preparing an input file than if
you were just typing the document directly.
Quotation marks like
       ``this''
have to be handled specially, as do quotes within
quotes:
       ``\,`this'            % \, separates the double and single quote.
        is what I just
        wrote, not  `that'\,''.

Dashes come in three sizes: an
       intra-word
dash, a medium dash for number ranges like
       1--2,
and a punctuation
       dash---like
this.

A sentence-ending space should be larger than the
space between words within a sentence.  You
sometimes have to type special commands in
conjunction with punctuation characters to get
this right, as in the following sentence.
       Gnats, gnus, etc.\ all  % `\ ' makes an inter-word space.
       begin with G\@.         % \@ marks end-of-sentence punctuation.
You should check the spaces after periods when
reading your output to make sure you haven't
forgotten any special cases.  Generating an
ellipsis
       \ldots\               % `\ ' is needed after `\ldots' because TeX
                             % ignores spaces after command names like \ldots
                             % made from \ + letters.
                             %
                             % Note how a `%' character causes TeX to ignore
                             % the end of the input line, so these blank lines
                             % do not start a new paragraph.
                             %
with the right spacing around the periods requires
a special command.

\LaTeX\ interprets some common characters as
commands, so you must type special commands to
generate them.  These characters include the
following:
       \$ \& \% \# \{ and \}.

In printing, text is usually emphasized with an
       \emph{italic}
type style.

\begin{em}
   A long segment of text can also be emphasized
   in this way.  Text within such a segment can be
   given \emph{additional} emphasis.
\end{em}

It is sometimes necessary to prevent \LaTeX\ from
breaking a line where it might otherwise do so.
This may be at a space, as between the ``Mr.''\ and
``Jones'' in
       ``Mr.~Jones'',        % ~ produces an unbreakable interword space.
or within a word---especially when the word is a
symbol like
       \mbox{\emph{itemnum}}
that makes little sense when hyphenated across
lines.

Footnotes\footnote{This is an example of a footnote.}
pose no problem.

\LaTeX\ is good at typesetting mathematical formulas
like
       \( x-3y + z = 7 \)
or
       \( a_{1} > x^{2n} + y^{2n} > x' \)
or
       \( \ip{A}{B} = \sum_{i} a_{i} b_{i} \).
The spaces you type in a formula are
ignored.  Remember that a letter like
       $x$                   % $ ... $  and  \( ... \)  are equivalent
is a formula when it denotes a mathematical
symbol, and it should be typed as one.

\section{Displayed Text}

Text is displayed by indenting it from the left
margin.  Quotations are commonly displayed.  There
are short quotations
\begin{quote}
   This is a short quotation.  It consists of a
   single paragraph of text.  See how it is formatted.
\end{quote}
and longer ones.
\begin{quotation}
   This is a longer quotation.  It consists of two
   paragraphs of text, neither of which are
   particularly interesting.

   This is the second paragraph of the quotation.  It
   is just as dull as the first paragraph.
\end{quotation}
Another frequently-displayed structure is a list.
The following is an example of an \emph{itemized}
list.
\begin{itemize}
   \item This is the first item of an itemized list.
         Each item in the list is marked with a ``tick''.
         You don't have to worry about what kind of tick
         mark is used.

   \item This is the second item of the list.  It
         contains another list nested inside it.  The inner
         list is an \emph{enumerated} list.
         \begin{enumerate}
            \item This is the first item of an enumerated
                  list that is nested within the itemized list.

            \item This is the second item of the inner list.
                  \LaTeX\ allows you to nest lists deeper than
                  you really should.
         \end{enumerate}
         This is the rest of the second item of the outer
         list.  It is no more interesting than any other
         part of the item.
   \item This is the third item of the list.
\end{itemize}
You can even display poetry.
\begin{verse}
   There is an environment
    for verse \\             % The \\ command separates lines
   Whose features some poets % within a stanza.
   will curse.

                             % One or more blank lines separate stanzas.

   For instead of making\\
   Them do \emph{all} line breaking, \\
   It allows them to put too many words on a line when they'd rather be
   forced to be terse.
\end{verse}

Mathematical formulas may also be displayed.  A
displayed formula
is
one-line long; multiline
formulas require special formatting instructions.
   \[  \ip{\Gamma}{\psi'} = x'' + y^{2} + z_{i}^{n}\]
Don't start a paragraph with a displayed equation,
nor make one a paragraph by itself.

\end{document}               % End of document.

\section{Bench Huge Section 40}
% synthetic bench section 40
% This is a sample LaTeX input file.  (Version of 12 August 2004.)
%
% A '%' character causes TeX to ignore all remaining text on the line,
% and is used for comments like this one.

\documentclass{article}      % Specifies the document class

                             % The preamble begins here.
\title{An Example Document}  % Declares the document's title.
\author{Leslie Lamport}      % Declares the author's name.
\date{January 21, 1994}      % Deleting this command produces today's date.

\newcommand{\ip}[2]{(#1, #2)}
                             % Defines \ip{arg1}{arg2} to mean
                             % (arg1, arg2).

%\newcommand{\ip}[2]{\langle #1 | #2\rangle}
                             % This is an alternative definition of
                             % \ip that is commented out.

\begin{document}             % End of preamble and beginning of text.

\maketitle                   % Produces the title.

This is an example input file.  Comparing it with
the output it generates can show you how to
produce a simple document of your own.

\section{Ordinary Text}      % Produces section heading.  Lower-level
                             % sections are begun with similar
                             % \subsection and \subsubsection commands.

The ends  of words and sentences are marked
  by   spaces. It  doesn't matter how many
spaces    you type; one is as good as 100.  The
end of   a line counts as a space.

One   or more   blank lines denote the  end
of  a paragraph.

Since any number of consecutive spaces are treated
like a single one, the formatting of the input
file makes no difference to
      \LaTeX,                % The \LaTeX command generates the LaTeX logo.
but it makes a difference to you.  When you use
\LaTeX, making your input file as easy to read
as possible will be a great help as you write
your document and when you change it.  This sample
file shows how you can add comments to your own input
file.

Because printing is different from typewriting,
there are a number of things that you have to do
differently when preparing an input file than if
you were just typing the document directly.
Quotation marks like
       ``this''
have to be handled specially, as do quotes within
quotes:
       ``\,`this'            % \, separates the double and single quote.
        is what I just
        wrote, not  `that'\,''.

Dashes come in three sizes: an
       intra-word
dash, a medium dash for number ranges like
       1--2,
and a punctuation
       dash---like
this.

A sentence-ending space should be larger than the
space between words within a sentence.  You
sometimes have to type special commands in
conjunction with punctuation characters to get
this right, as in the following sentence.
       Gnats, gnus, etc.\ all  % `\ ' makes an inter-word space.
       begin with G\@.         % \@ marks end-of-sentence punctuation.
You should check the spaces after periods when
reading your output to make sure you haven't
forgotten any special cases.  Generating an
ellipsis
       \ldots\               % `\ ' is needed after `\ldots' because TeX
                             % ignores spaces after command names like \ldots
                             % made from \ + letters.
                             %
                             % Note how a `%' character causes TeX to ignore
                             % the end of the input line, so these blank lines
                             % do not start a new paragraph.
                             %
with the right spacing around the periods requires
a special command.

\LaTeX\ interprets some common characters as
commands, so you must type special commands to
generate them.  These characters include the
following:
       \$ \& \% \# \{ and \}.

In printing, text is usually emphasized with an
       \emph{italic}
type style.

\begin{em}
   A long segment of text can also be emphasized
   in this way.  Text within such a segment can be
   given \emph{additional} emphasis.
\end{em}

It is sometimes necessary to prevent \LaTeX\ from
breaking a line where it might otherwise do so.
This may be at a space, as between the ``Mr.''\ and
``Jones'' in
       ``Mr.~Jones'',        % ~ produces an unbreakable interword space.
or within a word---especially when the word is a
symbol like
       \mbox{\emph{itemnum}}
that makes little sense when hyphenated across
lines.

Footnotes\footnote{This is an example of a footnote.}
pose no problem.

\LaTeX\ is good at typesetting mathematical formulas
like
       \( x-3y + z = 7 \)
or
       \( a_{1} > x^{2n} + y^{2n} > x' \)
or
       \( \ip{A}{B} = \sum_{i} a_{i} b_{i} \).
The spaces you type in a formula are
ignored.  Remember that a letter like
       $x$                   % $ ... $  and  \( ... \)  are equivalent
is a formula when it denotes a mathematical
symbol, and it should be typed as one.

\section{Displayed Text}

Text is displayed by indenting it from the left
margin.  Quotations are commonly displayed.  There
are short quotations
\begin{quote}
   This is a short quotation.  It consists of a
   single paragraph of text.  See how it is formatted.
\end{quote}
and longer ones.
\begin{quotation}
   This is a longer quotation.  It consists of two
   paragraphs of text, neither of which are
   particularly interesting.

   This is the second paragraph of the quotation.  It
   is just as dull as the first paragraph.
\end{quotation}
Another frequently-displayed structure is a list.
The following is an example of an \emph{itemized}
list.
\begin{itemize}
   \item This is the first item of an itemized list.
         Each item in the list is marked with a ``tick''.
         You don't have to worry about what kind of tick
         mark is used.

   \item This is the second item of the list.  It
         contains another list nested inside it.  The inner
         list is an \emph{enumerated} list.
         \begin{enumerate}
            \item This is the first item of an enumerated
                  list that is nested within the itemized list.

            \item This is the second item of the inner list.
                  \LaTeX\ allows you to nest lists deeper than
                  you really should.
         \end{enumerate}
         This is the rest of the second item of the outer
         list.  It is no more interesting than any other
         part of the item.
   \item This is the third item of the list.
\end{itemize}
You can even display poetry.
\begin{verse}
   There is an environment
    for verse \\             % The \\ command separates lines
   Whose features some poets % within a stanza.
   will curse.

                             % One or more blank lines separate stanzas.

   For instead of making\\
   Them do \emph{all} line breaking, \\
   It allows them to put too many words on a line when they'd rather be
   forced to be terse.
\end{verse}

Mathematical formulas may also be displayed.  A
displayed formula
is
one-line long; multiline
formulas require special formatting instructions.
   \[  \ip{\Gamma}{\psi'} = x'' + y^{2} + z_{i}^{n}\]
Don't start a paragraph with a displayed equation,
nor make one a paragraph by itself.

\end{document}               % End of document.

\section{Bench Huge Section 41}
% synthetic bench section 41
% This is a sample LaTeX input file.  (Version of 12 August 2004.)
%
% A '%' character causes TeX to ignore all remaining text on the line,
% and is used for comments like this one.

\documentclass{article}      % Specifies the document class

                             % The preamble begins here.
\title{An Example Document}  % Declares the document's title.
\author{Leslie Lamport}      % Declares the author's name.
\date{January 21, 1994}      % Deleting this command produces today's date.

\newcommand{\ip}[2]{(#1, #2)}
                             % Defines \ip{arg1}{arg2} to mean
                             % (arg1, arg2).

%\newcommand{\ip}[2]{\langle #1 | #2\rangle}
                             % This is an alternative definition of
                             % \ip that is commented out.

\begin{document}             % End of preamble and beginning of text.

\maketitle                   % Produces the title.

This is an example input file.  Comparing it with
the output it generates can show you how to
produce a simple document of your own.

\section{Ordinary Text}      % Produces section heading.  Lower-level
                             % sections are begun with similar
                             % \subsection and \subsubsection commands.

The ends  of words and sentences are marked
  by   spaces. It  doesn't matter how many
spaces    you type; one is as good as 100.  The
end of   a line counts as a space.

One   or more   blank lines denote the  end
of  a paragraph.

Since any number of consecutive spaces are treated
like a single one, the formatting of the input
file makes no difference to
      \LaTeX,                % The \LaTeX command generates the LaTeX logo.
but it makes a difference to you.  When you use
\LaTeX, making your input file as easy to read
as possible will be a great help as you write
your document and when you change it.  This sample
file shows how you can add comments to your own input
file.

Because printing is different from typewriting,
there are a number of things that you have to do
differently when preparing an input file than if
you were just typing the document directly.
Quotation marks like
       ``this''
have to be handled specially, as do quotes within
quotes:
       ``\,`this'            % \, separates the double and single quote.
        is what I just
        wrote, not  `that'\,''.

Dashes come in three sizes: an
       intra-word
dash, a medium dash for number ranges like
       1--2,
and a punctuation
       dash---like
this.

A sentence-ending space should be larger than the
space between words within a sentence.  You
sometimes have to type special commands in
conjunction with punctuation characters to get
this right, as in the following sentence.
       Gnats, gnus, etc.\ all  % `\ ' makes an inter-word space.
       begin with G\@.         % \@ marks end-of-sentence punctuation.
You should check the spaces after periods when
reading your output to make sure you haven't
forgotten any special cases.  Generating an
ellipsis
       \ldots\               % `\ ' is needed after `\ldots' because TeX
                             % ignores spaces after command names like \ldots
                             % made from \ + letters.
                             %
                             % Note how a `%' character causes TeX to ignore
                             % the end of the input line, so these blank lines
                             % do not start a new paragraph.
                             %
with the right spacing around the periods requires
a special command.

\LaTeX\ interprets some common characters as
commands, so you must type special commands to
generate them.  These characters include the
following:
       \$ \& \% \# \{ and \}.

In printing, text is usually emphasized with an
       \emph{italic}
type style.

\begin{em}
   A long segment of text can also be emphasized
   in this way.  Text within such a segment can be
   given \emph{additional} emphasis.
\end{em}

It is sometimes necessary to prevent \LaTeX\ from
breaking a line where it might otherwise do so.
This may be at a space, as between the ``Mr.''\ and
``Jones'' in
       ``Mr.~Jones'',        % ~ produces an unbreakable interword space.
or within a word---especially when the word is a
symbol like
       \mbox{\emph{itemnum}}
that makes little sense when hyphenated across
lines.

Footnotes\footnote{This is an example of a footnote.}
pose no problem.

\LaTeX\ is good at typesetting mathematical formulas
like
       \( x-3y + z = 7 \)
or
       \( a_{1} > x^{2n} + y^{2n} > x' \)
or
       \( \ip{A}{B} = \sum_{i} a_{i} b_{i} \).
The spaces you type in a formula are
ignored.  Remember that a letter like
       $x$                   % $ ... $  and  \( ... \)  are equivalent
is a formula when it denotes a mathematical
symbol, and it should be typed as one.

\section{Displayed Text}

Text is displayed by indenting it from the left
margin.  Quotations are commonly displayed.  There
are short quotations
\begin{quote}
   This is a short quotation.  It consists of a
   single paragraph of text.  See how it is formatted.
\end{quote}
and longer ones.
\begin{quotation}
   This is a longer quotation.  It consists of two
   paragraphs of text, neither of which are
   particularly interesting.

   This is the second paragraph of the quotation.  It
   is just as dull as the first paragraph.
\end{quotation}
Another frequently-displayed structure is a list.
The following is an example of an \emph{itemized}
list.
\begin{itemize}
   \item This is the first item of an itemized list.
         Each item in the list is marked with a ``tick''.
         You don't have to worry about what kind of tick
         mark is used.

   \item This is the second item of the list.  It
         contains another list nested inside it.  The inner
         list is an \emph{enumerated} list.
         \begin{enumerate}
            \item This is the first item of an enumerated
                  list that is nested within the itemized list.

            \item This is the second item of the inner list.
                  \LaTeX\ allows you to nest lists deeper than
                  you really should.
         \end{enumerate}
         This is the rest of the second item of the outer
         list.  It is no more interesting than any other
         part of the item.
   \item This is the third item of the list.
\end{itemize}
You can even display poetry.
\begin{verse}
   There is an environment
    for verse \\             % The \\ command separates lines
   Whose features some poets % within a stanza.
   will curse.

                             % One or more blank lines separate stanzas.

   For instead of making\\
   Them do \emph{all} line breaking, \\
   It allows them to put too many words on a line when they'd rather be
   forced to be terse.
\end{verse}

Mathematical formulas may also be displayed.  A
displayed formula
is
one-line long; multiline
formulas require special formatting instructions.
   \[  \ip{\Gamma}{\psi'} = x'' + y^{2} + z_{i}^{n}\]
Don't start a paragraph with a displayed equation,
nor make one a paragraph by itself.

\end{document}               % End of document.

\section{Bench Huge Section 42}
% synthetic bench section 42
% This is a sample LaTeX input file.  (Version of 12 August 2004.)
%
% A '%' character causes TeX to ignore all remaining text on the line,
% and is used for comments like this one.

\documentclass{article}      % Specifies the document class

                             % The preamble begins here.
\title{An Example Document}  % Declares the document's title.
\author{Leslie Lamport}      % Declares the author's name.
\date{January 21, 1994}      % Deleting this command produces today's date.

\newcommand{\ip}[2]{(#1, #2)}
                             % Defines \ip{arg1}{arg2} to mean
                             % (arg1, arg2).

%\newcommand{\ip}[2]{\langle #1 | #2\rangle}
                             % This is an alternative definition of
                             % \ip that is commented out.

\begin{document}             % End of preamble and beginning of text.

\maketitle                   % Produces the title.

This is an example input file.  Comparing it with
the output it generates can show you how to
produce a simple document of your own.

\section{Ordinary Text}      % Produces section heading.  Lower-level
                             % sections are begun with similar
                             % \subsection and \subsubsection commands.

The ends  of words and sentences are marked
  by   spaces. It  doesn't matter how many
spaces    you type; one is as good as 100.  The
end of   a line counts as a space.

One   or more   blank lines denote the  end
of  a paragraph.

Since any number of consecutive spaces are treated
like a single one, the formatting of the input
file makes no difference to
      \LaTeX,                % The \LaTeX command generates the LaTeX logo.
but it makes a difference to you.  When you use
\LaTeX, making your input file as easy to read
as possible will be a great help as you write
your document and when you change it.  This sample
file shows how you can add comments to your own input
file.

Because printing is different from typewriting,
there are a number of things that you have to do
differently when preparing an input file than if
you were just typing the document directly.
Quotation marks like
       ``this''
have to be handled specially, as do quotes within
quotes:
       ``\,`this'            % \, separates the double and single quote.
        is what I just
        wrote, not  `that'\,''.

Dashes come in three sizes: an
       intra-word
dash, a medium dash for number ranges like
       1--2,
and a punctuation
       dash---like
this.

A sentence-ending space should be larger than the
space between words within a sentence.  You
sometimes have to type special commands in
conjunction with punctuation characters to get
this right, as in the following sentence.
       Gnats, gnus, etc.\ all  % `\ ' makes an inter-word space.
       begin with G\@.         % \@ marks end-of-sentence punctuation.
You should check the spaces after periods when
reading your output to make sure you haven't
forgotten any special cases.  Generating an
ellipsis
       \ldots\               % `\ ' is needed after `\ldots' because TeX
                             % ignores spaces after command names like \ldots
                             % made from \ + letters.
                             %
                             % Note how a `%' character causes TeX to ignore
                             % the end of the input line, so these blank lines
                             % do not start a new paragraph.
                             %
with the right spacing around the periods requires
a special command.

\LaTeX\ interprets some common characters as
commands, so you must type special commands to
generate them.  These characters include the
following:
       \$ \& \% \# \{ and \}.

In printing, text is usually emphasized with an
       \emph{italic}
type style.

\begin{em}
   A long segment of text can also be emphasized
   in this way.  Text within such a segment can be
   given \emph{additional} emphasis.
\end{em}

It is sometimes necessary to prevent \LaTeX\ from
breaking a line where it might otherwise do so.
This may be at a space, as between the ``Mr.''\ and
``Jones'' in
       ``Mr.~Jones'',        % ~ produces an unbreakable interword space.
or within a word---especially when the word is a
symbol like
       \mbox{\emph{itemnum}}
that makes little sense when hyphenated across
lines.

Footnotes\footnote{This is an example of a footnote.}
pose no problem.

\LaTeX\ is good at typesetting mathematical formulas
like
       \( x-3y + z = 7 \)
or
       \( a_{1} > x^{2n} + y^{2n} > x' \)
or
       \( \ip{A}{B} = \sum_{i} a_{i} b_{i} \).
The spaces you type in a formula are
ignored.  Remember that a letter like
       $x$                   % $ ... $  and  \( ... \)  are equivalent
is a formula when it denotes a mathematical
symbol, and it should be typed as one.

\section{Displayed Text}

Text is displayed by indenting it from the left
margin.  Quotations are commonly displayed.  There
are short quotations
\begin{quote}
   This is a short quotation.  It consists of a
   single paragraph of text.  See how it is formatted.
\end{quote}
and longer ones.
\begin{quotation}
   This is a longer quotation.  It consists of two
   paragraphs of text, neither of which are
   particularly interesting.

   This is the second paragraph of the quotation.  It
   is just as dull as the first paragraph.
\end{quotation}
Another frequently-displayed structure is a list.
The following is an example of an \emph{itemized}
list.
\begin{itemize}
   \item This is the first item of an itemized list.
         Each item in the list is marked with a ``tick''.
         You don't have to worry about what kind of tick
         mark is used.

   \item This is the second item of the list.  It
         contains another list nested inside it.  The inner
         list is an \emph{enumerated} list.
         \begin{enumerate}
            \item This is the first item of an enumerated
                  list that is nested within the itemized list.

            \item This is the second item of the inner list.
                  \LaTeX\ allows you to nest lists deeper than
                  you really should.
         \end{enumerate}
         This is the rest of the second item of the outer
         list.  It is no more interesting than any other
         part of the item.
   \item This is the third item of the list.
\end{itemize}
You can even display poetry.
\begin{verse}
   There is an environment
    for verse \\             % The \\ command separates lines
   Whose features some poets % within a stanza.
   will curse.

                             % One or more blank lines separate stanzas.

   For instead of making\\
   Them do \emph{all} line breaking, \\
   It allows them to put too many words on a line when they'd rather be
   forced to be terse.
\end{verse}

Mathematical formulas may also be displayed.  A
displayed formula
is
one-line long; multiline
formulas require special formatting instructions.
   \[  \ip{\Gamma}{\psi'} = x'' + y^{2} + z_{i}^{n}\]
Don't start a paragraph with a displayed equation,
nor make one a paragraph by itself.

\end{document}               % End of document.

\section{Bench Huge Section 43}
% synthetic bench section 43
% This is a sample LaTeX input file.  (Version of 12 August 2004.)
%
% A '%' character causes TeX to ignore all remaining text on the line,
% and is used for comments like this one.

\documentclass{article}      % Specifies the document class

                             % The preamble begins here.
\title{An Example Document}  % Declares the document's title.
\author{Leslie Lamport}      % Declares the author's name.
\date{January 21, 1994}      % Deleting this command produces today's date.

\newcommand{\ip}[2]{(#1, #2)}
                             % Defines \ip{arg1}{arg2} to mean
                             % (arg1, arg2).

%\newcommand{\ip}[2]{\langle #1 | #2\rangle}
                             % This is an alternative definition of
                             % \ip that is commented out.

\begin{document}             % End of preamble and beginning of text.

\maketitle                   % Produces the title.

This is an example input file.  Comparing it with
the output it generates can show you how to
produce a simple document of your own.

\section{Ordinary Text}      % Produces section heading.  Lower-level
                             % sections are begun with similar
                             % \subsection and \subsubsection commands.

The ends  of words and sentences are marked
  by   spaces. It  doesn't matter how many
spaces    you type; one is as good as 100.  The
end of   a line counts as a space.

One   or more   blank lines denote the  end
of  a paragraph.

Since any number of consecutive spaces are treated
like a single one, the formatting of the input
file makes no difference to
      \LaTeX,                % The \LaTeX command generates the LaTeX logo.
but it makes a difference to you.  When you use
\LaTeX, making your input file as easy to read
as possible will be a great help as you write
your document and when you change it.  This sample
file shows how you can add comments to your own input
file.

Because printing is different from typewriting,
there are a number of things that you have to do
differently when preparing an input file than if
you were just typing the document directly.
Quotation marks like
       ``this''
have to be handled specially, as do quotes within
quotes:
       ``\,`this'            % \, separates the double and single quote.
        is what I just
        wrote, not  `that'\,''.

Dashes come in three sizes: an
       intra-word
dash, a medium dash for number ranges like
       1--2,
and a punctuation
       dash---like
this.

A sentence-ending space should be larger than the
space between words within a sentence.  You
sometimes have to type special commands in
conjunction with punctuation characters to get
this right, as in the following sentence.
       Gnats, gnus, etc.\ all  % `\ ' makes an inter-word space.
       begin with G\@.         % \@ marks end-of-sentence punctuation.
You should check the spaces after periods when
reading your output to make sure you haven't
forgotten any special cases.  Generating an
ellipsis
       \ldots\               % `\ ' is needed after `\ldots' because TeX
                             % ignores spaces after command names like \ldots
                             % made from \ + letters.
                             %
                             % Note how a `%' character causes TeX to ignore
                             % the end of the input line, so these blank lines
                             % do not start a new paragraph.
                             %
with the right spacing around the periods requires
a special command.

\LaTeX\ interprets some common characters as
commands, so you must type special commands to
generate them.  These characters include the
following:
       \$ \& \% \# \{ and \}.

In printing, text is usually emphasized with an
       \emph{italic}
type style.

\begin{em}
   A long segment of text can also be emphasized
   in this way.  Text within such a segment can be
   given \emph{additional} emphasis.
\end{em}

It is sometimes necessary to prevent \LaTeX\ from
breaking a line where it might otherwise do so.
This may be at a space, as between the ``Mr.''\ and
``Jones'' in
       ``Mr.~Jones'',        % ~ produces an unbreakable interword space.
or within a word---especially when the word is a
symbol like
       \mbox{\emph{itemnum}}
that makes little sense when hyphenated across
lines.

Footnotes\footnote{This is an example of a footnote.}
pose no problem.

\LaTeX\ is good at typesetting mathematical formulas
like
       \( x-3y + z = 7 \)
or
       \( a_{1} > x^{2n} + y^{2n} > x' \)
or
       \( \ip{A}{B} = \sum_{i} a_{i} b_{i} \).
The spaces you type in a formula are
ignored.  Remember that a letter like
       $x$                   % $ ... $  and  \( ... \)  are equivalent
is a formula when it denotes a mathematical
symbol, and it should be typed as one.

\section{Displayed Text}

Text is displayed by indenting it from the left
margin.  Quotations are commonly displayed.  There
are short quotations
\begin{quote}
   This is a short quotation.  It consists of a
   single paragraph of text.  See how it is formatted.
\end{quote}
and longer ones.
\begin{quotation}
   This is a longer quotation.  It consists of two
   paragraphs of text, neither of which are
   particularly interesting.

   This is the second paragraph of the quotation.  It
   is just as dull as the first paragraph.
\end{quotation}
Another frequently-displayed structure is a list.
The following is an example of an \emph{itemized}
list.
\begin{itemize}
   \item This is the first item of an itemized list.
         Each item in the list is marked with a ``tick''.
         You don't have to worry about what kind of tick
         mark is used.

   \item This is the second item of the list.  It
         contains another list nested inside it.  The inner
         list is an \emph{enumerated} list.
         \begin{enumerate}
            \item This is the first item of an enumerated
                  list that is nested within the itemized list.

            \item This is the second item of the inner list.
                  \LaTeX\ allows you to nest lists deeper than
                  you really should.
         \end{enumerate}
         This is the rest of the second item of the outer
         list.  It is no more interesting than any other
         part of the item.
   \item This is the third item of the list.
\end{itemize}
You can even display poetry.
\begin{verse}
   There is an environment
    for verse \\             % The \\ command separates lines
   Whose features some poets % within a stanza.
   will curse.

                             % One or more blank lines separate stanzas.

   For instead of making\\
   Them do \emph{all} line breaking, \\
   It allows them to put too many words on a line when they'd rather be
   forced to be terse.
\end{verse}

Mathematical formulas may also be displayed.  A
displayed formula
is
one-line long; multiline
formulas require special formatting instructions.
   \[  \ip{\Gamma}{\psi'} = x'' + y^{2} + z_{i}^{n}\]
Don't start a paragraph with a displayed equation,
nor make one a paragraph by itself.

\end{document}               % End of document.

\section{Bench Huge Section 44}
% synthetic bench section 44
% This is a sample LaTeX input file.  (Version of 12 August 2004.)
%
% A '%' character causes TeX to ignore all remaining text on the line,
% and is used for comments like this one.

\documentclass{article}      % Specifies the document class

                             % The preamble begins here.
\title{An Example Document}  % Declares the document's title.
\author{Leslie Lamport}      % Declares the author's name.
\date{January 21, 1994}      % Deleting this command produces today's date.

\newcommand{\ip}[2]{(#1, #2)}
                             % Defines \ip{arg1}{arg2} to mean
                             % (arg1, arg2).

%\newcommand{\ip}[2]{\langle #1 | #2\rangle}
                             % This is an alternative definition of
                             % \ip that is commented out.

\begin{document}             % End of preamble and beginning of text.

\maketitle                   % Produces the title.

This is an example input file.  Comparing it with
the output it generates can show you how to
produce a simple document of your own.

\section{Ordinary Text}      % Produces section heading.  Lower-level
                             % sections are begun with similar
                             % \subsection and \subsubsection commands.

The ends  of words and sentences are marked
  by   spaces. It  doesn't matter how many
spaces    you type; one is as good as 100.  The
end of   a line counts as a space.

One   or more   blank lines denote the  end
of  a paragraph.

Since any number of consecutive spaces are treated
like a single one, the formatting of the input
file makes no difference to
      \LaTeX,                % The \LaTeX command generates the LaTeX logo.
but it makes a difference to you.  When you use
\LaTeX, making your input file as easy to read
as possible will be a great help as you write
your document and when you change it.  This sample
file shows how you can add comments to your own input
file.

Because printing is different from typewriting,
there are a number of things that you have to do
differently when preparing an input file than if
you were just typing the document directly.
Quotation marks like
       ``this''
have to be handled specially, as do quotes within
quotes:
       ``\,`this'            % \, separates the double and single quote.
        is what I just
        wrote, not  `that'\,''.

Dashes come in three sizes: an
       intra-word
dash, a medium dash for number ranges like
       1--2,
and a punctuation
       dash---like
this.

A sentence-ending space should be larger than the
space between words within a sentence.  You
sometimes have to type special commands in
conjunction with punctuation characters to get
this right, as in the following sentence.
       Gnats, gnus, etc.\ all  % `\ ' makes an inter-word space.
       begin with G\@.         % \@ marks end-of-sentence punctuation.
You should check the spaces after periods when
reading your output to make sure you haven't
forgotten any special cases.  Generating an
ellipsis
       \ldots\               % `\ ' is needed after `\ldots' because TeX
                             % ignores spaces after command names like \ldots
                             % made from \ + letters.
                             %
                             % Note how a `%' character causes TeX to ignore
                             % the end of the input line, so these blank lines
                             % do not start a new paragraph.
                             %
with the right spacing around the periods requires
a special command.

\LaTeX\ interprets some common characters as
commands, so you must type special commands to
generate them.  These characters include the
following:
       \$ \& \% \# \{ and \}.

In printing, text is usually emphasized with an
       \emph{italic}
type style.

\begin{em}
   A long segment of text can also be emphasized
   in this way.  Text within such a segment can be
   given \emph{additional} emphasis.
\end{em}

It is sometimes necessary to prevent \LaTeX\ from
breaking a line where it might otherwise do so.
This may be at a space, as between the ``Mr.''\ and
``Jones'' in
       ``Mr.~Jones'',        % ~ produces an unbreakable interword space.
or within a word---especially when the word is a
symbol like
       \mbox{\emph{itemnum}}
that makes little sense when hyphenated across
lines.

Footnotes\footnote{This is an example of a footnote.}
pose no problem.

\LaTeX\ is good at typesetting mathematical formulas
like
       \( x-3y + z = 7 \)
or
       \( a_{1} > x^{2n} + y^{2n} > x' \)
or
       \( \ip{A}{B} = \sum_{i} a_{i} b_{i} \).
The spaces you type in a formula are
ignored.  Remember that a letter like
       $x$                   % $ ... $  and  \( ... \)  are equivalent
is a formula when it denotes a mathematical
symbol, and it should be typed as one.

\section{Displayed Text}

Text is displayed by indenting it from the left
margin.  Quotations are commonly displayed.  There
are short quotations
\begin{quote}
   This is a short quotation.  It consists of a
   single paragraph of text.  See how it is formatted.
\end{quote}
and longer ones.
\begin{quotation}
   This is a longer quotation.  It consists of two
   paragraphs of text, neither of which are
   particularly interesting.

   This is the second paragraph of the quotation.  It
   is just as dull as the first paragraph.
\end{quotation}
Another frequently-displayed structure is a list.
The following is an example of an \emph{itemized}
list.
\begin{itemize}
   \item This is the first item of an itemized list.
         Each item in the list is marked with a ``tick''.
         You don't have to worry about what kind of tick
         mark is used.

   \item This is the second item of the list.  It
         contains another list nested inside it.  The inner
         list is an \emph{enumerated} list.
         \begin{enumerate}
            \item This is the first item of an enumerated
                  list that is nested within the itemized list.

            \item This is the second item of the inner list.
                  \LaTeX\ allows you to nest lists deeper than
                  you really should.
         \end{enumerate}
         This is the rest of the second item of the outer
         list.  It is no more interesting than any other
         part of the item.
   \item This is the third item of the list.
\end{itemize}
You can even display poetry.
\begin{verse}
   There is an environment
    for verse \\             % The \\ command separates lines
   Whose features some poets % within a stanza.
   will curse.

                             % One or more blank lines separate stanzas.

   For instead of making\\
   Them do \emph{all} line breaking, \\
   It allows them to put too many words on a line when they'd rather be
   forced to be terse.
\end{verse}

Mathematical formulas may also be displayed.  A
displayed formula
is
one-line long; multiline
formulas require special formatting instructions.
   \[  \ip{\Gamma}{\psi'} = x'' + y^{2} + z_{i}^{n}\]
Don't start a paragraph with a displayed equation,
nor make one a paragraph by itself.

\end{document}               % End of document.

\section{Bench Huge Section 45}
% synthetic bench section 45
% This is a sample LaTeX input file.  (Version of 12 August 2004.)
%
% A '%' character causes TeX to ignore all remaining text on the line,
% and is used for comments like this one.

\documentclass{article}      % Specifies the document class

                             % The preamble begins here.
\title{An Example Document}  % Declares the document's title.
\author{Leslie Lamport}      % Declares the author's name.
\date{January 21, 1994}      % Deleting this command produces today's date.

\newcommand{\ip}[2]{(#1, #2)}
                             % Defines \ip{arg1}{arg2} to mean
                             % (arg1, arg2).

%\newcommand{\ip}[2]{\langle #1 | #2\rangle}
                             % This is an alternative definition of
                             % \ip that is commented out.

\begin{document}             % End of preamble and beginning of text.

\maketitle                   % Produces the title.

This is an example input file.  Comparing it with
the output it generates can show you how to
produce a simple document of your own.

\section{Ordinary Text}      % Produces section heading.  Lower-level
                             % sections are begun with similar
                             % \subsection and \subsubsection commands.

The ends  of words and sentences are marked
  by   spaces. It  doesn't matter how many
spaces    you type; one is as good as 100.  The
end of   a line counts as a space.

One   or more   blank lines denote the  end
of  a paragraph.

Since any number of consecutive spaces are treated
like a single one, the formatting of the input
file makes no difference to
      \LaTeX,                % The \LaTeX command generates the LaTeX logo.
but it makes a difference to you.  When you use
\LaTeX, making your input file as easy to read
as possible will be a great help as you write
your document and when you change it.  This sample
file shows how you can add comments to your own input
file.

Because printing is different from typewriting,
there are a number of things that you have to do
differently when preparing an input file than if
you were just typing the document directly.
Quotation marks like
       ``this''
have to be handled specially, as do quotes within
quotes:
       ``\,`this'            % \, separates the double and single quote.
        is what I just
        wrote, not  `that'\,''.

Dashes come in three sizes: an
       intra-word
dash, a medium dash for number ranges like
       1--2,
and a punctuation
       dash---like
this.

A sentence-ending space should be larger than the
space between words within a sentence.  You
sometimes have to type special commands in
conjunction with punctuation characters to get
this right, as in the following sentence.
       Gnats, gnus, etc.\ all  % `\ ' makes an inter-word space.
       begin with G\@.         % \@ marks end-of-sentence punctuation.
You should check the spaces after periods when
reading your output to make sure you haven't
forgotten any special cases.  Generating an
ellipsis
       \ldots\               % `\ ' is needed after `\ldots' because TeX
                             % ignores spaces after command names like \ldots
                             % made from \ + letters.
                             %
                             % Note how a `%' character causes TeX to ignore
                             % the end of the input line, so these blank lines
                             % do not start a new paragraph.
                             %
with the right spacing around the periods requires
a special command.

\LaTeX\ interprets some common characters as
commands, so you must type special commands to
generate them.  These characters include the
following:
       \$ \& \% \# \{ and \}.

In printing, text is usually emphasized with an
       \emph{italic}
type style.

\begin{em}
   A long segment of text can also be emphasized
   in this way.  Text within such a segment can be
   given \emph{additional} emphasis.
\end{em}

It is sometimes necessary to prevent \LaTeX\ from
breaking a line where it might otherwise do so.
This may be at a space, as between the ``Mr.''\ and
``Jones'' in
       ``Mr.~Jones'',        % ~ produces an unbreakable interword space.
or within a word---especially when the word is a
symbol like
       \mbox{\emph{itemnum}}
that makes little sense when hyphenated across
lines.

Footnotes\footnote{This is an example of a footnote.}
pose no problem.

\LaTeX\ is good at typesetting mathematical formulas
like
       \( x-3y + z = 7 \)
or
       \( a_{1} > x^{2n} + y^{2n} > x' \)
or
       \( \ip{A}{B} = \sum_{i} a_{i} b_{i} \).
The spaces you type in a formula are
ignored.  Remember that a letter like
       $x$                   % $ ... $  and  \( ... \)  are equivalent
is a formula when it denotes a mathematical
symbol, and it should be typed as one.

\section{Displayed Text}

Text is displayed by indenting it from the left
margin.  Quotations are commonly displayed.  There
are short quotations
\begin{quote}
   This is a short quotation.  It consists of a
   single paragraph of text.  See how it is formatted.
\end{quote}
and longer ones.
\begin{quotation}
   This is a longer quotation.  It consists of two
   paragraphs of text, neither of which are
   particularly interesting.

   This is the second paragraph of the quotation.  It
   is just as dull as the first paragraph.
\end{quotation}
Another frequently-displayed structure is a list.
The following is an example of an \emph{itemized}
list.
\begin{itemize}
   \item This is the first item of an itemized list.
         Each item in the list is marked with a ``tick''.
         You don't have to worry about what kind of tick
         mark is used.

   \item This is the second item of the list.  It
         contains another list nested inside it.  The inner
         list is an \emph{enumerated} list.
         \begin{enumerate}
            \item This is the first item of an enumerated
                  list that is nested within the itemized list.

            \item This is the second item of the inner list.
                  \LaTeX\ allows you to nest lists deeper than
                  you really should.
         \end{enumerate}
         This is the rest of the second item of the outer
         list.  It is no more interesting than any other
         part of the item.
   \item This is the third item of the list.
\end{itemize}
You can even display poetry.
\begin{verse}
   There is an environment
    for verse \\             % The \\ command separates lines
   Whose features some poets % within a stanza.
   will curse.

                             % One or more blank lines separate stanzas.

   For instead of making\\
   Them do \emph{all} line breaking, \\
   It allows them to put too many words on a line when they'd rather be
   forced to be terse.
\end{verse}

Mathematical formulas may also be displayed.  A
displayed formula
is
one-line long; multiline
formulas require special formatting instructions.
   \[  \ip{\Gamma}{\psi'} = x'' + y^{2} + z_{i}^{n}\]
Don't start a paragraph with a displayed equation,
nor make one a paragraph by itself.

\end{document}               % End of document.

\section{Bench Huge Section 46}
% synthetic bench section 46
% This is a sample LaTeX input file.  (Version of 12 August 2004.)
%
% A '%' character causes TeX to ignore all remaining text on the line,
% and is used for comments like this one.

\documentclass{article}      % Specifies the document class

                             % The preamble begins here.
\title{An Example Document}  % Declares the document's title.
\author{Leslie Lamport}      % Declares the author's name.
\date{January 21, 1994}      % Deleting this command produces today's date.

\newcommand{\ip}[2]{(#1, #2)}
                             % Defines \ip{arg1}{arg2} to mean
                             % (arg1, arg2).

%\newcommand{\ip}[2]{\langle #1 | #2\rangle}
                             % This is an alternative definition of
                             % \ip that is commented out.

\begin{document}             % End of preamble and beginning of text.

\maketitle                   % Produces the title.

This is an example input file.  Comparing it with
the output it generates can show you how to
produce a simple document of your own.

\section{Ordinary Text}      % Produces section heading.  Lower-level
                             % sections are begun with similar
                             % \subsection and \subsubsection commands.

The ends  of words and sentences are marked
  by   spaces. It  doesn't matter how many
spaces    you type; one is as good as 100.  The
end of   a line counts as a space.

One   or more   blank lines denote the  end
of  a paragraph.

Since any number of consecutive spaces are treated
like a single one, the formatting of the input
file makes no difference to
      \LaTeX,                % The \LaTeX command generates the LaTeX logo.
but it makes a difference to you.  When you use
\LaTeX, making your input file as easy to read
as possible will be a great help as you write
your document and when you change it.  This sample
file shows how you can add comments to your own input
file.

Because printing is different from typewriting,
there are a number of things that you have to do
differently when preparing an input file than if
you were just typing the document directly.
Quotation marks like
       ``this''
have to be handled specially, as do quotes within
quotes:
       ``\,`this'            % \, separates the double and single quote.
        is what I just
        wrote, not  `that'\,''.

Dashes come in three sizes: an
       intra-word
dash, a medium dash for number ranges like
       1--2,
and a punctuation
       dash---like
this.

A sentence-ending space should be larger than the
space between words within a sentence.  You
sometimes have to type special commands in
conjunction with punctuation characters to get
this right, as in the following sentence.
       Gnats, gnus, etc.\ all  % `\ ' makes an inter-word space.
       begin with G\@.         % \@ marks end-of-sentence punctuation.
You should check the spaces after periods when
reading your output to make sure you haven't
forgotten any special cases.  Generating an
ellipsis
       \ldots\               % `\ ' is needed after `\ldots' because TeX
                             % ignores spaces after command names like \ldots
                             % made from \ + letters.
                             %
                             % Note how a `%' character causes TeX to ignore
                             % the end of the input line, so these blank lines
                             % do not start a new paragraph.
                             %
with the right spacing around the periods requires
a special command.

\LaTeX\ interprets some common characters as
commands, so you must type special commands to
generate them.  These characters include the
following:
       \$ \& \% \# \{ and \}.

In printing, text is usually emphasized with an
       \emph{italic}
type style.

\begin{em}
   A long segment of text can also be emphasized
   in this way.  Text within such a segment can be
   given \emph{additional} emphasis.
\end{em}

It is sometimes necessary to prevent \LaTeX\ from
breaking a line where it might otherwise do so.
This may be at a space, as between the ``Mr.''\ and
``Jones'' in
       ``Mr.~Jones'',        % ~ produces an unbreakable interword space.
or within a word---especially when the word is a
symbol like
       \mbox{\emph{itemnum}}
that makes little sense when hyphenated across
lines.

Footnotes\footnote{This is an example of a footnote.}
pose no problem.

\LaTeX\ is good at typesetting mathematical formulas
like
       \( x-3y + z = 7 \)
or
       \( a_{1} > x^{2n} + y^{2n} > x' \)
or
       \( \ip{A}{B} = \sum_{i} a_{i} b_{i} \).
The spaces you type in a formula are
ignored.  Remember that a letter like
       $x$                   % $ ... $  and  \( ... \)  are equivalent
is a formula when it denotes a mathematical
symbol, and it should be typed as one.

\section{Displayed Text}

Text is displayed by indenting it from the left
margin.  Quotations are commonly displayed.  There
are short quotations
\begin{quote}
   This is a short quotation.  It consists of a
   single paragraph of text.  See how it is formatted.
\end{quote}
and longer ones.
\begin{quotation}
   This is a longer quotation.  It consists of two
   paragraphs of text, neither of which are
   particularly interesting.

   This is the second paragraph of the quotation.  It
   is just as dull as the first paragraph.
\end{quotation}
Another frequently-displayed structure is a list.
The following is an example of an \emph{itemized}
list.
\begin{itemize}
   \item This is the first item of an itemized list.
         Each item in the list is marked with a ``tick''.
         You don't have to worry about what kind of tick
         mark is used.

   \item This is the second item of the list.  It
         contains another list nested inside it.  The inner
         list is an \emph{enumerated} list.
         \begin{enumerate}
            \item This is the first item of an enumerated
                  list that is nested within the itemized list.

            \item This is the second item of the inner list.
                  \LaTeX\ allows you to nest lists deeper than
                  you really should.
         \end{enumerate}
         This is the rest of the second item of the outer
         list.  It is no more interesting than any other
         part of the item.
   \item This is the third item of the list.
\end{itemize}
You can even display poetry.
\begin{verse}
   There is an environment
    for verse \\             % The \\ command separates lines
   Whose features some poets % within a stanza.
   will curse.

                             % One or more blank lines separate stanzas.

   For instead of making\\
   Them do \emph{all} line breaking, \\
   It allows them to put too many words on a line when they'd rather be
   forced to be terse.
\end{verse}

Mathematical formulas may also be displayed.  A
displayed formula
is
one-line long; multiline
formulas require special formatting instructions.
   \[  \ip{\Gamma}{\psi'} = x'' + y^{2} + z_{i}^{n}\]
Don't start a paragraph with a displayed equation,
nor make one a paragraph by itself.

\end{document}               % End of document.

\section{Bench Huge Section 47}
% synthetic bench section 47
% This is a sample LaTeX input file.  (Version of 12 August 2004.)
%
% A '%' character causes TeX to ignore all remaining text on the line,
% and is used for comments like this one.

\documentclass{article}      % Specifies the document class

                             % The preamble begins here.
\title{An Example Document}  % Declares the document's title.
\author{Leslie Lamport}      % Declares the author's name.
\date{January 21, 1994}      % Deleting this command produces today's date.

\newcommand{\ip}[2]{(#1, #2)}
                             % Defines \ip{arg1}{arg2} to mean
                             % (arg1, arg2).

%\newcommand{\ip}[2]{\langle #1 | #2\rangle}
                             % This is an alternative definition of
                             % \ip that is commented out.

\begin{document}             % End of preamble and beginning of text.

\maketitle                   % Produces the title.

This is an example input file.  Comparing it with
the output it generates can show you how to
produce a simple document of your own.

\section{Ordinary Text}      % Produces section heading.  Lower-level
                             % sections are begun with similar
                             % \subsection and \subsubsection commands.

The ends  of words and sentences are marked
  by   spaces. It  doesn't matter how many
spaces    you type; one is as good as 100.  The
end of   a line counts as a space.

One   or more   blank lines denote the  end
of  a paragraph.

Since any number of consecutive spaces are treated
like a single one, the formatting of the input
file makes no difference to
      \LaTeX,                % The \LaTeX command generates the LaTeX logo.
but it makes a difference to you.  When you use
\LaTeX, making your input file as easy to read
as possible will be a great help as you write
your document and when you change it.  This sample
file shows how you can add comments to your own input
file.

Because printing is different from typewriting,
there are a number of things that you have to do
differently when preparing an input file than if
you were just typing the document directly.
Quotation marks like
       ``this''
have to be handled specially, as do quotes within
quotes:
       ``\,`this'            % \, separates the double and single quote.
        is what I just
        wrote, not  `that'\,''.

Dashes come in three sizes: an
       intra-word
dash, a medium dash for number ranges like
       1--2,
and a punctuation
       dash---like
this.

A sentence-ending space should be larger than the
space between words within a sentence.  You
sometimes have to type special commands in
conjunction with punctuation characters to get
this right, as in the following sentence.
       Gnats, gnus, etc.\ all  % `\ ' makes an inter-word space.
       begin with G\@.         % \@ marks end-of-sentence punctuation.
You should check the spaces after periods when
reading your output to make sure you haven't
forgotten any special cases.  Generating an
ellipsis
       \ldots\               % `\ ' is needed after `\ldots' because TeX
                             % ignores spaces after command names like \ldots
                             % made from \ + letters.
                             %
                             % Note how a `%' character causes TeX to ignore
                             % the end of the input line, so these blank lines
                             % do not start a new paragraph.
                             %
with the right spacing around the periods requires
a special command.

\LaTeX\ interprets some common characters as
commands, so you must type special commands to
generate them.  These characters include the
following:
       \$ \& \% \# \{ and \}.

In printing, text is usually emphasized with an
       \emph{italic}
type style.

\begin{em}
   A long segment of text can also be emphasized
   in this way.  Text within such a segment can be
   given \emph{additional} emphasis.
\end{em}

It is sometimes necessary to prevent \LaTeX\ from
breaking a line where it might otherwise do so.
This may be at a space, as between the ``Mr.''\ and
``Jones'' in
       ``Mr.~Jones'',        % ~ produces an unbreakable interword space.
or within a word---especially when the word is a
symbol like
       \mbox{\emph{itemnum}}
that makes little sense when hyphenated across
lines.

Footnotes\footnote{This is an example of a footnote.}
pose no problem.

\LaTeX\ is good at typesetting mathematical formulas
like
       \( x-3y + z = 7 \)
or
       \( a_{1} > x^{2n} + y^{2n} > x' \)
or
       \( \ip{A}{B} = \sum_{i} a_{i} b_{i} \).
The spaces you type in a formula are
ignored.  Remember that a letter like
       $x$                   % $ ... $  and  \( ... \)  are equivalent
is a formula when it denotes a mathematical
symbol, and it should be typed as one.

\section{Displayed Text}

Text is displayed by indenting it from the left
margin.  Quotations are commonly displayed.  There
are short quotations
\begin{quote}
   This is a short quotation.  It consists of a
   single paragraph of text.  See how it is formatted.
\end{quote}
and longer ones.
\begin{quotation}
   This is a longer quotation.  It consists of two
   paragraphs of text, neither of which are
   particularly interesting.

   This is the second paragraph of the quotation.  It
   is just as dull as the first paragraph.
\end{quotation}
Another frequently-displayed structure is a list.
The following is an example of an \emph{itemized}
list.
\begin{itemize}
   \item This is the first item of an itemized list.
         Each item in the list is marked with a ``tick''.
         You don't have to worry about what kind of tick
         mark is used.

   \item This is the second item of the list.  It
         contains another list nested inside it.  The inner
         list is an \emph{enumerated} list.
         \begin{enumerate}
            \item This is the first item of an enumerated
                  list that is nested within the itemized list.

            \item This is the second item of the inner list.
                  \LaTeX\ allows you to nest lists deeper than
                  you really should.
         \end{enumerate}
         This is the rest of the second item of the outer
         list.  It is no more interesting than any other
         part of the item.
   \item This is the third item of the list.
\end{itemize}
You can even display poetry.
\begin{verse}
   There is an environment
    for verse \\             % The \\ command separates lines
   Whose features some poets % within a stanza.
   will curse.

                             % One or more blank lines separate stanzas.

   For instead of making\\
   Them do \emph{all} line breaking, \\
   It allows them to put too many words on a line when they'd rather be
   forced to be terse.
\end{verse}

Mathematical formulas may also be displayed.  A
displayed formula
is
one-line long; multiline
formulas require special formatting instructions.
   \[  \ip{\Gamma}{\psi'} = x'' + y^{2} + z_{i}^{n}\]
Don't start a paragraph with a displayed equation,
nor make one a paragraph by itself.

\end{document}               % End of document.

\section{Bench Huge Section 48}
% synthetic bench section 48
% This is a sample LaTeX input file.  (Version of 12 August 2004.)
%
% A '%' character causes TeX to ignore all remaining text on the line,
% and is used for comments like this one.

\documentclass{article}      % Specifies the document class

                             % The preamble begins here.
\title{An Example Document}  % Declares the document's title.
\author{Leslie Lamport}      % Declares the author's name.
\date{January 21, 1994}      % Deleting this command produces today's date.

\newcommand{\ip}[2]{(#1, #2)}
                             % Defines \ip{arg1}{arg2} to mean
                             % (arg1, arg2).

%\newcommand{\ip}[2]{\langle #1 | #2\rangle}
                             % This is an alternative definition of
                             % \ip that is commented out.

\begin{document}             % End of preamble and beginning of text.

\maketitle                   % Produces the title.

This is an example input file.  Comparing it with
the output it generates can show you how to
produce a simple document of your own.

\section{Ordinary Text}      % Produces section heading.  Lower-level
                             % sections are begun with similar
                             % \subsection and \subsubsection commands.

The ends  of words and sentences are marked
  by   spaces. It  doesn't matter how many
spaces    you type; one is as good as 100.  The
end of   a line counts as a space.

One   or more   blank lines denote the  end
of  a paragraph.

Since any number of consecutive spaces are treated
like a single one, the formatting of the input
file makes no difference to
      \LaTeX,                % The \LaTeX command generates the LaTeX logo.
but it makes a difference to you.  When you use
\LaTeX, making your input file as easy to read
as possible will be a great help as you write
your document and when you change it.  This sample
file shows how you can add comments to your own input
file.

Because printing is different from typewriting,
there are a number of things that you have to do
differently when preparing an input file than if
you were just typing the document directly.
Quotation marks like
       ``this''
have to be handled specially, as do quotes within
quotes:
       ``\,`this'            % \, separates the double and single quote.
        is what I just
        wrote, not  `that'\,''.

Dashes come in three sizes: an
       intra-word
dash, a medium dash for number ranges like
       1--2,
and a punctuation
       dash---like
this.

A sentence-ending space should be larger than the
space between words within a sentence.  You
sometimes have to type special commands in
conjunction with punctuation characters to get
this right, as in the following sentence.
       Gnats, gnus, etc.\ all  % `\ ' makes an inter-word space.
       begin with G\@.         % \@ marks end-of-sentence punctuation.
You should check the spaces after periods when
reading your output to make sure you haven't
forgotten any special cases.  Generating an
ellipsis
       \ldots\               % `\ ' is needed after `\ldots' because TeX
                             % ignores spaces after command names like \ldots
                             % made from \ + letters.
                             %
                             % Note how a `%' character causes TeX to ignore
                             % the end of the input line, so these blank lines
                             % do not start a new paragraph.
                             %
with the right spacing around the periods requires
a special command.

\LaTeX\ interprets some common characters as
commands, so you must type special commands to
generate them.  These characters include the
following:
       \$ \& \% \# \{ and \}.

In printing, text is usually emphasized with an
       \emph{italic}
type style.

\begin{em}
   A long segment of text can also be emphasized
   in this way.  Text within such a segment can be
   given \emph{additional} emphasis.
\end{em}

It is sometimes necessary to prevent \LaTeX\ from
breaking a line where it might otherwise do so.
This may be at a space, as between the ``Mr.''\ and
``Jones'' in
       ``Mr.~Jones'',        % ~ produces an unbreakable interword space.
or within a word---especially when the word is a
symbol like
       \mbox{\emph{itemnum}}
that makes little sense when hyphenated across
lines.

Footnotes\footnote{This is an example of a footnote.}
pose no problem.

\LaTeX\ is good at typesetting mathematical formulas
like
       \( x-3y + z = 7 \)
or
       \( a_{1} > x^{2n} + y^{2n} > x' \)
or
       \( \ip{A}{B} = \sum_{i} a_{i} b_{i} \).
The spaces you type in a formula are
ignored.  Remember that a letter like
       $x$                   % $ ... $  and  \( ... \)  are equivalent
is a formula when it denotes a mathematical
symbol, and it should be typed as one.

\section{Displayed Text}

Text is displayed by indenting it from the left
margin.  Quotations are commonly displayed.  There
are short quotations
\begin{quote}
   This is a short quotation.  It consists of a
   single paragraph of text.  See how it is formatted.
\end{quote}
and longer ones.
\begin{quotation}
   This is a longer quotation.  It consists of two
   paragraphs of text, neither of which are
   particularly interesting.

   This is the second paragraph of the quotation.  It
   is just as dull as the first paragraph.
\end{quotation}
Another frequently-displayed structure is a list.
The following is an example of an \emph{itemized}
list.
\begin{itemize}
   \item This is the first item of an itemized list.
         Each item in the list is marked with a ``tick''.
         You don't have to worry about what kind of tick
         mark is used.

   \item This is the second item of the list.  It
         contains another list nested inside it.  The inner
         list is an \emph{enumerated} list.
         \begin{enumerate}
            \item This is the first item of an enumerated
                  list that is nested within the itemized list.

            \item This is the second item of the inner list.
                  \LaTeX\ allows you to nest lists deeper than
                  you really should.
         \end{enumerate}
         This is the rest of the second item of the outer
         list.  It is no more interesting than any other
         part of the item.
   \item This is the third item of the list.
\end{itemize}
You can even display poetry.
\begin{verse}
   There is an environment
    for verse \\             % The \\ command separates lines
   Whose features some poets % within a stanza.
   will curse.

                             % One or more blank lines separate stanzas.

   For instead of making\\
   Them do \emph{all} line breaking, \\
   It allows them to put too many words on a line when they'd rather be
   forced to be terse.
\end{verse}

Mathematical formulas may also be displayed.  A
displayed formula
is
one-line long; multiline
formulas require special formatting instructions.
   \[  \ip{\Gamma}{\psi'} = x'' + y^{2} + z_{i}^{n}\]
Don't start a paragraph with a displayed equation,
nor make one a paragraph by itself.

\end{document}               % End of document.

\section{Bench Huge Section 49}
% synthetic bench section 49
% This is a sample LaTeX input file.  (Version of 12 August 2004.)
%
% A '%' character causes TeX to ignore all remaining text on the line,
% and is used for comments like this one.

\documentclass{article}      % Specifies the document class

                             % The preamble begins here.
\title{An Example Document}  % Declares the document's title.
\author{Leslie Lamport}      % Declares the author's name.
\date{January 21, 1994}      % Deleting this command produces today's date.

\newcommand{\ip}[2]{(#1, #2)}
                             % Defines \ip{arg1}{arg2} to mean
                             % (arg1, arg2).

%\newcommand{\ip}[2]{\langle #1 | #2\rangle}
                             % This is an alternative definition of
                             % \ip that is commented out.

\begin{document}             % End of preamble and beginning of text.

\maketitle                   % Produces the title.

This is an example input file.  Comparing it with
the output it generates can show you how to
produce a simple document of your own.

\section{Ordinary Text}      % Produces section heading.  Lower-level
                             % sections are begun with similar
                             % \subsection and \subsubsection commands.

The ends  of words and sentences are marked
  by   spaces. It  doesn't matter how many
spaces    you type; one is as good as 100.  The
end of   a line counts as a space.

One   or more   blank lines denote the  end
of  a paragraph.

Since any number of consecutive spaces are treated
like a single one, the formatting of the input
file makes no difference to
      \LaTeX,                % The \LaTeX command generates the LaTeX logo.
but it makes a difference to you.  When you use
\LaTeX, making your input file as easy to read
as possible will be a great help as you write
your document and when you change it.  This sample
file shows how you can add comments to your own input
file.

Because printing is different from typewriting,
there are a number of things that you have to do
differently when preparing an input file than if
you were just typing the document directly.
Quotation marks like
       ``this''
have to be handled specially, as do quotes within
quotes:
       ``\,`this'            % \, separates the double and single quote.
        is what I just
        wrote, not  `that'\,''.

Dashes come in three sizes: an
       intra-word
dash, a medium dash for number ranges like
       1--2,
and a punctuation
       dash---like
this.

A sentence-ending space should be larger than the
space between words within a sentence.  You
sometimes have to type special commands in
conjunction with punctuation characters to get
this right, as in the following sentence.
       Gnats, gnus, etc.\ all  % `\ ' makes an inter-word space.
       begin with G\@.         % \@ marks end-of-sentence punctuation.
You should check the spaces after periods when
reading your output to make sure you haven't
forgotten any special cases.  Generating an
ellipsis
       \ldots\               % `\ ' is needed after `\ldots' because TeX
                             % ignores spaces after command names like \ldots
                             % made from \ + letters.
                             %
                             % Note how a `%' character causes TeX to ignore
                             % the end of the input line, so these blank lines
                             % do not start a new paragraph.
                             %
with the right spacing around the periods requires
a special command.

\LaTeX\ interprets some common characters as
commands, so you must type special commands to
generate them.  These characters include the
following:
       \$ \& \% \# \{ and \}.

In printing, text is usually emphasized with an
       \emph{italic}
type style.

\begin{em}
   A long segment of text can also be emphasized
   in this way.  Text within such a segment can be
   given \emph{additional} emphasis.
\end{em}

It is sometimes necessary to prevent \LaTeX\ from
breaking a line where it might otherwise do so.
This may be at a space, as between the ``Mr.''\ and
``Jones'' in
       ``Mr.~Jones'',        % ~ produces an unbreakable interword space.
or within a word---especially when the word is a
symbol like
       \mbox{\emph{itemnum}}
that makes little sense when hyphenated across
lines.

Footnotes\footnote{This is an example of a footnote.}
pose no problem.

\LaTeX\ is good at typesetting mathematical formulas
like
       \( x-3y + z = 7 \)
or
       \( a_{1} > x^{2n} + y^{2n} > x' \)
or
       \( \ip{A}{B} = \sum_{i} a_{i} b_{i} \).
The spaces you type in a formula are
ignored.  Remember that a letter like
       $x$                   % $ ... $  and  \( ... \)  are equivalent
is a formula when it denotes a mathematical
symbol, and it should be typed as one.

\section{Displayed Text}

Text is displayed by indenting it from the left
margin.  Quotations are commonly displayed.  There
are short quotations
\begin{quote}
   This is a short quotation.  It consists of a
   single paragraph of text.  See how it is formatted.
\end{quote}
and longer ones.
\begin{quotation}
   This is a longer quotation.  It consists of two
   paragraphs of text, neither of which are
   particularly interesting.

   This is the second paragraph of the quotation.  It
   is just as dull as the first paragraph.
\end{quotation}
Another frequently-displayed structure is a list.
The following is an example of an \emph{itemized}
list.
\begin{itemize}
   \item This is the first item of an itemized list.
         Each item in the list is marked with a ``tick''.
         You don't have to worry about what kind of tick
         mark is used.

   \item This is the second item of the list.  It
         contains another list nested inside it.  The inner
         list is an \emph{enumerated} list.
         \begin{enumerate}
            \item This is the first item of an enumerated
                  list that is nested within the itemized list.

            \item This is the second item of the inner list.
                  \LaTeX\ allows you to nest lists deeper than
                  you really should.
         \end{enumerate}
         This is the rest of the second item of the outer
         list.  It is no more interesting than any other
         part of the item.
   \item This is the third item of the list.
\end{itemize}
You can even display poetry.
\begin{verse}
   There is an environment
    for verse \\             % The \\ command separates lines
   Whose features some poets % within a stanza.
   will curse.

                             % One or more blank lines separate stanzas.

   For instead of making\\
   Them do \emph{all} line breaking, \\
   It allows them to put too many words on a line when they'd rather be
   forced to be terse.
\end{verse}

Mathematical formulas may also be displayed.  A
displayed formula
is
one-line long; multiline
formulas require special formatting instructions.
   \[  \ip{\Gamma}{\psi'} = x'' + y^{2} + z_{i}^{n}\]
Don't start a paragraph with a displayed equation,
nor make one a paragraph by itself.

\end{document}               % End of document.

\section{Bench Huge Section 50}
% synthetic bench section 50
% This is a sample LaTeX input file.  (Version of 12 August 2004.)
%
% A '%' character causes TeX to ignore all remaining text on the line,
% and is used for comments like this one.

\documentclass{article}      % Specifies the document class

                             % The preamble begins here.
\title{An Example Document}  % Declares the document's title.
\author{Leslie Lamport}      % Declares the author's name.
\date{January 21, 1994}      % Deleting this command produces today's date.

\newcommand{\ip}[2]{(#1, #2)}
                             % Defines \ip{arg1}{arg2} to mean
                             % (arg1, arg2).

%\newcommand{\ip}[2]{\langle #1 | #2\rangle}
                             % This is an alternative definition of
                             % \ip that is commented out.

\begin{document}             % End of preamble and beginning of text.

\maketitle                   % Produces the title.

This is an example input file.  Comparing it with
the output it generates can show you how to
produce a simple document of your own.

\section{Ordinary Text}      % Produces section heading.  Lower-level
                             % sections are begun with similar
                             % \subsection and \subsubsection commands.

The ends  of words and sentences are marked
  by   spaces. It  doesn't matter how many
spaces    you type; one is as good as 100.  The
end of   a line counts as a space.

One   or more   blank lines denote the  end
of  a paragraph.

Since any number of consecutive spaces are treated
like a single one, the formatting of the input
file makes no difference to
      \LaTeX,                % The \LaTeX command generates the LaTeX logo.
but it makes a difference to you.  When you use
\LaTeX, making your input file as easy to read
as possible will be a great help as you write
your document and when you change it.  This sample
file shows how you can add comments to your own input
file.

Because printing is different from typewriting,
there are a number of things that you have to do
differently when preparing an input file than if
you were just typing the document directly.
Quotation marks like
       ``this''
have to be handled specially, as do quotes within
quotes:
       ``\,`this'            % \, separates the double and single quote.
        is what I just
        wrote, not  `that'\,''.

Dashes come in three sizes: an
       intra-word
dash, a medium dash for number ranges like
       1--2,
and a punctuation
       dash---like
this.

A sentence-ending space should be larger than the
space between words within a sentence.  You
sometimes have to type special commands in
conjunction with punctuation characters to get
this right, as in the following sentence.
       Gnats, gnus, etc.\ all  % `\ ' makes an inter-word space.
       begin with G\@.         % \@ marks end-of-sentence punctuation.
You should check the spaces after periods when
reading your output to make sure you haven't
forgotten any special cases.  Generating an
ellipsis
       \ldots\               % `\ ' is needed after `\ldots' because TeX
                             % ignores spaces after command names like \ldots
                             % made from \ + letters.
                             %
                             % Note how a `%' character causes TeX to ignore
                             % the end of the input line, so these blank lines
                             % do not start a new paragraph.
                             %
with the right spacing around the periods requires
a special command.

\LaTeX\ interprets some common characters as
commands, so you must type special commands to
generate them.  These characters include the
following:
       \$ \& \% \# \{ and \}.

In printing, text is usually emphasized with an
       \emph{italic}
type style.

\begin{em}
   A long segment of text can also be emphasized
   in this way.  Text within such a segment can be
   given \emph{additional} emphasis.
\end{em}

It is sometimes necessary to prevent \LaTeX\ from
breaking a line where it might otherwise do so.
This may be at a space, as between the ``Mr.''\ and
``Jones'' in
       ``Mr.~Jones'',        % ~ produces an unbreakable interword space.
or within a word---especially when the word is a
symbol like
       \mbox{\emph{itemnum}}
that makes little sense when hyphenated across
lines.

Footnotes\footnote{This is an example of a footnote.}
pose no problem.

\LaTeX\ is good at typesetting mathematical formulas
like
       \( x-3y + z = 7 \)
or
       \( a_{1} > x^{2n} + y^{2n} > x' \)
or
       \( \ip{A}{B} = \sum_{i} a_{i} b_{i} \).
The spaces you type in a formula are
ignored.  Remember that a letter like
       $x$                   % $ ... $  and  \( ... \)  are equivalent
is a formula when it denotes a mathematical
symbol, and it should be typed as one.

\section{Displayed Text}

Text is displayed by indenting it from the left
margin.  Quotations are commonly displayed.  There
are short quotations
\begin{quote}
   This is a short quotation.  It consists of a
   single paragraph of text.  See how it is formatted.
\end{quote}
and longer ones.
\begin{quotation}
   This is a longer quotation.  It consists of two
   paragraphs of text, neither of which are
   particularly interesting.

   This is the second paragraph of the quotation.  It
   is just as dull as the first paragraph.
\end{quotation}
Another frequently-displayed structure is a list.
The following is an example of an \emph{itemized}
list.
\begin{itemize}
   \item This is the first item of an itemized list.
         Each item in the list is marked with a ``tick''.
         You don't have to worry about what kind of tick
         mark is used.

   \item This is the second item of the list.  It
         contains another list nested inside it.  The inner
         list is an \emph{enumerated} list.
         \begin{enumerate}
            \item This is the first item of an enumerated
                  list that is nested within the itemized list.

            \item This is the second item of the inner list.
                  \LaTeX\ allows you to nest lists deeper than
                  you really should.
         \end{enumerate}
         This is the rest of the second item of the outer
         list.  It is no more interesting than any other
         part of the item.
   \item This is the third item of the list.
\end{itemize}
You can even display poetry.
\begin{verse}
   There is an environment
    for verse \\             % The \\ command separates lines
   Whose features some poets % within a stanza.
   will curse.

                             % One or more blank lines separate stanzas.

   For instead of making\\
   Them do \emph{all} line breaking, \\
   It allows them to put too many words on a line when they'd rather be
   forced to be terse.
\end{verse}

Mathematical formulas may also be displayed.  A
displayed formula
is
one-line long; multiline
formulas require special formatting instructions.
   \[  \ip{\Gamma}{\psi'} = x'' + y^{2} + z_{i}^{n}\]
Don't start a paragraph with a displayed equation,
nor make one a paragraph by itself.

\end{document}               % End of document.

\section{Bench Huge Section 51}
% synthetic bench section 51
% This is a sample LaTeX input file.  (Version of 12 August 2004.)
%
% A '%' character causes TeX to ignore all remaining text on the line,
% and is used for comments like this one.

\documentclass{article}      % Specifies the document class

                             % The preamble begins here.
\title{An Example Document}  % Declares the document's title.
\author{Leslie Lamport}      % Declares the author's name.
\date{January 21, 1994}      % Deleting this command produces today's date.

\newcommand{\ip}[2]{(#1, #2)}
                             % Defines \ip{arg1}{arg2} to mean
                             % (arg1, arg2).

%\newcommand{\ip}[2]{\langle #1 | #2\rangle}
                             % This is an alternative definition of
                             % \ip that is commented out.

\begin{document}             % End of preamble and beginning of text.

\maketitle                   % Produces the title.

This is an example input file.  Comparing it with
the output it generates can show you how to
produce a simple document of your own.

\section{Ordinary Text}      % Produces section heading.  Lower-level
                             % sections are begun with similar
                             % \subsection and \subsubsection commands.

The ends  of words and sentences are marked
  by   spaces. It  doesn't matter how many
spaces    you type; one is as good as 100.  The
end of   a line counts as a space.

One   or more   blank lines denote the  end
of  a paragraph.

Since any number of consecutive spaces are treated
like a single one, the formatting of the input
file makes no difference to
      \LaTeX,                % The \LaTeX command generates the LaTeX logo.
but it makes a difference to you.  When you use
\LaTeX, making your input file as easy to read
as possible will be a great help as you write
your document and when you change it.  This sample
file shows how you can add comments to your own input
file.

Because printing is different from typewriting,
there are a number of things that you have to do
differently when preparing an input file than if
you were just typing the document directly.
Quotation marks like
       ``this''
have to be handled specially, as do quotes within
quotes:
       ``\,`this'            % \, separates the double and single quote.
        is what I just
        wrote, not  `that'\,''.

Dashes come in three sizes: an
       intra-word
dash, a medium dash for number ranges like
       1--2,
and a punctuation
       dash---like
this.

A sentence-ending space should be larger than the
space between words within a sentence.  You
sometimes have to type special commands in
conjunction with punctuation characters to get
this right, as in the following sentence.
       Gnats, gnus, etc.\ all  % `\ ' makes an inter-word space.
       begin with G\@.         % \@ marks end-of-sentence punctuation.
You should check the spaces after periods when
reading your output to make sure you haven't
forgotten any special cases.  Generating an
ellipsis
       \ldots\               % `\ ' is needed after `\ldots' because TeX
                             % ignores spaces after command names like \ldots
                             % made from \ + letters.
                             %
                             % Note how a `%' character causes TeX to ignore
                             % the end of the input line, so these blank lines
                             % do not start a new paragraph.
                             %
with the right spacing around the periods requires
a special command.

\LaTeX\ interprets some common characters as
commands, so you must type special commands to
generate them.  These characters include the
following:
       \$ \& \% \# \{ and \}.

In printing, text is usually emphasized with an
       \emph{italic}
type style.

\begin{em}
   A long segment of text can also be emphasized
   in this way.  Text within such a segment can be
   given \emph{additional} emphasis.
\end{em}

It is sometimes necessary to prevent \LaTeX\ from
breaking a line where it might otherwise do so.
This may be at a space, as between the ``Mr.''\ and
``Jones'' in
       ``Mr.~Jones'',        % ~ produces an unbreakable interword space.
or within a word---especially when the word is a
symbol like
       \mbox{\emph{itemnum}}
that makes little sense when hyphenated across
lines.

Footnotes\footnote{This is an example of a footnote.}
pose no problem.

\LaTeX\ is good at typesetting mathematical formulas
like
       \( x-3y + z = 7 \)
or
       \( a_{1} > x^{2n} + y^{2n} > x' \)
or
       \( \ip{A}{B} = \sum_{i} a_{i} b_{i} \).
The spaces you type in a formula are
ignored.  Remember that a letter like
       $x$                   % $ ... $  and  \( ... \)  are equivalent
is a formula when it denotes a mathematical
symbol, and it should be typed as one.

\section{Displayed Text}

Text is displayed by indenting it from the left
margin.  Quotations are commonly displayed.  There
are short quotations
\begin{quote}
   This is a short quotation.  It consists of a
   single paragraph of text.  See how it is formatted.
\end{quote}
and longer ones.
\begin{quotation}
   This is a longer quotation.  It consists of two
   paragraphs of text, neither of which are
   particularly interesting.

   This is the second paragraph of the quotation.  It
   is just as dull as the first paragraph.
\end{quotation}
Another frequently-displayed structure is a list.
The following is an example of an \emph{itemized}
list.
\begin{itemize}
   \item This is the first item of an itemized list.
         Each item in the list is marked with a ``tick''.
         You don't have to worry about what kind of tick
         mark is used.

   \item This is the second item of the list.  It
         contains another list nested inside it.  The inner
         list is an \emph{enumerated} list.
         \begin{enumerate}
            \item This is the first item of an enumerated
                  list that is nested within the itemized list.

            \item This is the second item of the inner list.
                  \LaTeX\ allows you to nest lists deeper than
                  you really should.
         \end{enumerate}
         This is the rest of the second item of the outer
         list.  It is no more interesting than any other
         part of the item.
   \item This is the third item of the list.
\end{itemize}
You can even display poetry.
\begin{verse}
   There is an environment
    for verse \\             % The \\ command separates lines
   Whose features some poets % within a stanza.
   will curse.

                             % One or more blank lines separate stanzas.

   For instead of making\\
   Them do \emph{all} line breaking, \\
   It allows them to put too many words on a line when they'd rather be
   forced to be terse.
\end{verse}

Mathematical formulas may also be displayed.  A
displayed formula
is
one-line long; multiline
formulas require special formatting instructions.
   \[  \ip{\Gamma}{\psi'} = x'' + y^{2} + z_{i}^{n}\]
Don't start a paragraph with a displayed equation,
nor make one a paragraph by itself.

\end{document}               % End of document.

\section{Bench Huge Section 52}
% synthetic bench section 52
% This is a sample LaTeX input file.  (Version of 12 August 2004.)
%
% A '%' character causes TeX to ignore all remaining text on the line,
% and is used for comments like this one.

\documentclass{article}      % Specifies the document class

                             % The preamble begins here.
\title{An Example Document}  % Declares the document's title.
\author{Leslie Lamport}      % Declares the author's name.
\date{January 21, 1994}      % Deleting this command produces today's date.

\newcommand{\ip}[2]{(#1, #2)}
                             % Defines \ip{arg1}{arg2} to mean
                             % (arg1, arg2).

%\newcommand{\ip}[2]{\langle #1 | #2\rangle}
                             % This is an alternative definition of
                             % \ip that is commented out.

\begin{document}             % End of preamble and beginning of text.

\maketitle                   % Produces the title.

This is an example input file.  Comparing it with
the output it generates can show you how to
produce a simple document of your own.

\section{Ordinary Text}      % Produces section heading.  Lower-level
                             % sections are begun with similar
                             % \subsection and \subsubsection commands.

The ends  of words and sentences are marked
  by   spaces. It  doesn't matter how many
spaces    you type; one is as good as 100.  The
end of   a line counts as a space.

One   or more   blank lines denote the  end
of  a paragraph.

Since any number of consecutive spaces are treated
like a single one, the formatting of the input
file makes no difference to
      \LaTeX,                % The \LaTeX command generates the LaTeX logo.
but it makes a difference to you.  When you use
\LaTeX, making your input file as easy to read
as possible will be a great help as you write
your document and when you change it.  This sample
file shows how you can add comments to your own input
file.

Because printing is different from typewriting,
there are a number of things that you have to do
differently when preparing an input file than if
you were just typing the document directly.
Quotation marks like
       ``this''
have to be handled specially, as do quotes within
quotes:
       ``\,`this'            % \, separates the double and single quote.
        is what I just
        wrote, not  `that'\,''.

Dashes come in three sizes: an
       intra-word
dash, a medium dash for number ranges like
       1--2,
and a punctuation
       dash---like
this.

A sentence-ending space should be larger than the
space between words within a sentence.  You
sometimes have to type special commands in
conjunction with punctuation characters to get
this right, as in the following sentence.
       Gnats, gnus, etc.\ all  % `\ ' makes an inter-word space.
       begin with G\@.         % \@ marks end-of-sentence punctuation.
You should check the spaces after periods when
reading your output to make sure you haven't
forgotten any special cases.  Generating an
ellipsis
       \ldots\               % `\ ' is needed after `\ldots' because TeX
                             % ignores spaces after command names like \ldots
                             % made from \ + letters.
                             %
                             % Note how a `%' character causes TeX to ignore
                             % the end of the input line, so these blank lines
                             % do not start a new paragraph.
                             %
with the right spacing around the periods requires
a special command.

\LaTeX\ interprets some common characters as
commands, so you must type special commands to
generate them.  These characters include the
following:
       \$ \& \% \# \{ and \}.

In printing, text is usually emphasized with an
       \emph{italic}
type style.

\begin{em}
   A long segment of text can also be emphasized
   in this way.  Text within such a segment can be
   given \emph{additional} emphasis.
\end{em}

It is sometimes necessary to prevent \LaTeX\ from
breaking a line where it might otherwise do so.
This may be at a space, as between the ``Mr.''\ and
``Jones'' in
       ``Mr.~Jones'',        % ~ produces an unbreakable interword space.
or within a word---especially when the word is a
symbol like
       \mbox{\emph{itemnum}}
that makes little sense when hyphenated across
lines.

Footnotes\footnote{This is an example of a footnote.}
pose no problem.

\LaTeX\ is good at typesetting mathematical formulas
like
       \( x-3y + z = 7 \)
or
       \( a_{1} > x^{2n} + y^{2n} > x' \)
or
       \( \ip{A}{B} = \sum_{i} a_{i} b_{i} \).
The spaces you type in a formula are
ignored.  Remember that a letter like
       $x$                   % $ ... $  and  \( ... \)  are equivalent
is a formula when it denotes a mathematical
symbol, and it should be typed as one.

\section{Displayed Text}

Text is displayed by indenting it from the left
margin.  Quotations are commonly displayed.  There
are short quotations
\begin{quote}
   This is a short quotation.  It consists of a
   single paragraph of text.  See how it is formatted.
\end{quote}
and longer ones.
\begin{quotation}
   This is a longer quotation.  It consists of two
   paragraphs of text, neither of which are
   particularly interesting.

   This is the second paragraph of the quotation.  It
   is just as dull as the first paragraph.
\end{quotation}
Another frequently-displayed structure is a list.
The following is an example of an \emph{itemized}
list.
\begin{itemize}
   \item This is the first item of an itemized list.
         Each item in the list is marked with a ``tick''.
         You don't have to worry about what kind of tick
         mark is used.

   \item This is the second item of the list.  It
         contains another list nested inside it.  The inner
         list is an \emph{enumerated} list.
         \begin{enumerate}
            \item This is the first item of an enumerated
                  list that is nested within the itemized list.

            \item This is the second item of the inner list.
                  \LaTeX\ allows you to nest lists deeper than
                  you really should.
         \end{enumerate}
         This is the rest of the second item of the outer
         list.  It is no more interesting than any other
         part of the item.
   \item This is the third item of the list.
\end{itemize}
You can even display poetry.
\begin{verse}
   There is an environment
    for verse \\             % The \\ command separates lines
   Whose features some poets % within a stanza.
   will curse.

                             % One or more blank lines separate stanzas.

   For instead of making\\
   Them do \emph{all} line breaking, \\
   It allows them to put too many words on a line when they'd rather be
   forced to be terse.
\end{verse}

Mathematical formulas may also be displayed.  A
displayed formula
is
one-line long; multiline
formulas require special formatting instructions.
   \[  \ip{\Gamma}{\psi'} = x'' + y^{2} + z_{i}^{n}\]
Don't start a paragraph with a displayed equation,
nor make one a paragraph by itself.

\end{document}               % End of document.

\section{Bench Huge Section 53}
% synthetic bench section 53
% This is a sample LaTeX input file.  (Version of 12 August 2004.)
%
% A '%' character causes TeX to ignore all remaining text on the line,
% and is used for comments like this one.

\documentclass{article}      % Specifies the document class

                             % The preamble begins here.
\title{An Example Document}  % Declares the document's title.
\author{Leslie Lamport}      % Declares the author's name.
\date{January 21, 1994}      % Deleting this command produces today's date.

\newcommand{\ip}[2]{(#1, #2)}
                             % Defines \ip{arg1}{arg2} to mean
                             % (arg1, arg2).

%\newcommand{\ip}[2]{\langle #1 | #2\rangle}
                             % This is an alternative definition of
                             % \ip that is commented out.

\begin{document}             % End of preamble and beginning of text.

\maketitle                   % Produces the title.

This is an example input file.  Comparing it with
the output it generates can show you how to
produce a simple document of your own.

\section{Ordinary Text}      % Produces section heading.  Lower-level
                             % sections are begun with similar
                             % \subsection and \subsubsection commands.

The ends  of words and sentences are marked
  by   spaces. It  doesn't matter how many
spaces    you type; one is as good as 100.  The
end of   a line counts as a space.

One   or more   blank lines denote the  end
of  a paragraph.

Since any number of consecutive spaces are treated
like a single one, the formatting of the input
file makes no difference to
      \LaTeX,                % The \LaTeX command generates the LaTeX logo.
but it makes a difference to you.  When you use
\LaTeX, making your input file as easy to read
as possible will be a great help as you write
your document and when you change it.  This sample
file shows how you can add comments to your own input
file.

Because printing is different from typewriting,
there are a number of things that you have to do
differently when preparing an input file than if
you were just typing the document directly.
Quotation marks like
       ``this''
have to be handled specially, as do quotes within
quotes:
       ``\,`this'            % \, separates the double and single quote.
        is what I just
        wrote, not  `that'\,''.

Dashes come in three sizes: an
       intra-word
dash, a medium dash for number ranges like
       1--2,
and a punctuation
       dash---like
this.

A sentence-ending space should be larger than the
space between words within a sentence.  You
sometimes have to type special commands in
conjunction with punctuation characters to get
this right, as in the following sentence.
       Gnats, gnus, etc.\ all  % `\ ' makes an inter-word space.
       begin with G\@.         % \@ marks end-of-sentence punctuation.
You should check the spaces after periods when
reading your output to make sure you haven't
forgotten any special cases.  Generating an
ellipsis
       \ldots\               % `\ ' is needed after `\ldots' because TeX
                             % ignores spaces after command names like \ldots
                             % made from \ + letters.
                             %
                             % Note how a `%' character causes TeX to ignore
                             % the end of the input line, so these blank lines
                             % do not start a new paragraph.
                             %
with the right spacing around the periods requires
a special command.

\LaTeX\ interprets some common characters as
commands, so you must type special commands to
generate them.  These characters include the
following:
       \$ \& \% \# \{ and \}.

In printing, text is usually emphasized with an
       \emph{italic}
type style.

\begin{em}
   A long segment of text can also be emphasized
   in this way.  Text within such a segment can be
   given \emph{additional} emphasis.
\end{em}

It is sometimes necessary to prevent \LaTeX\ from
breaking a line where it might otherwise do so.
This may be at a space, as between the ``Mr.''\ and
``Jones'' in
       ``Mr.~Jones'',        % ~ produces an unbreakable interword space.
or within a word---especially when the word is a
symbol like
       \mbox{\emph{itemnum}}
that makes little sense when hyphenated across
lines.

Footnotes\footnote{This is an example of a footnote.}
pose no problem.

\LaTeX\ is good at typesetting mathematical formulas
like
       \( x-3y + z = 7 \)
or
       \( a_{1} > x^{2n} + y^{2n} > x' \)
or
       \( \ip{A}{B} = \sum_{i} a_{i} b_{i} \).
The spaces you type in a formula are
ignored.  Remember that a letter like
       $x$                   % $ ... $  and  \( ... \)  are equivalent
is a formula when it denotes a mathematical
symbol, and it should be typed as one.

\section{Displayed Text}

Text is displayed by indenting it from the left
margin.  Quotations are commonly displayed.  There
are short quotations
\begin{quote}
   This is a short quotation.  It consists of a
   single paragraph of text.  See how it is formatted.
\end{quote}
and longer ones.
\begin{quotation}
   This is a longer quotation.  It consists of two
   paragraphs of text, neither of which are
   particularly interesting.

   This is the second paragraph of the quotation.  It
   is just as dull as the first paragraph.
\end{quotation}
Another frequently-displayed structure is a list.
The following is an example of an \emph{itemized}
list.
\begin{itemize}
   \item This is the first item of an itemized list.
         Each item in the list is marked with a ``tick''.
         You don't have to worry about what kind of tick
         mark is used.

   \item This is the second item of the list.  It
         contains another list nested inside it.  The inner
         list is an \emph{enumerated} list.
         \begin{enumerate}
            \item This is the first item of an enumerated
                  list that is nested within the itemized list.

            \item This is the second item of the inner list.
                  \LaTeX\ allows you to nest lists deeper than
                  you really should.
         \end{enumerate}
         This is the rest of the second item of the outer
         list.  It is no more interesting than any other
         part of the item.
   \item This is the third item of the list.
\end{itemize}
You can even display poetry.
\begin{verse}
   There is an environment
    for verse \\             % The \\ command separates lines
   Whose features some poets % within a stanza.
   will curse.

                             % One or more blank lines separate stanzas.

   For instead of making\\
   Them do \emph{all} line breaking, \\
   It allows them to put too many words on a line when they'd rather be
   forced to be terse.
\end{verse}

Mathematical formulas may also be displayed.  A
displayed formula
is
one-line long; multiline
formulas require special formatting instructions.
   \[  \ip{\Gamma}{\psi'} = x'' + y^{2} + z_{i}^{n}\]
Don't start a paragraph with a displayed equation,
nor make one a paragraph by itself.

\end{document}               % End of document.

\section{Bench Huge Section 54}
% synthetic bench section 54
% This is a sample LaTeX input file.  (Version of 12 August 2004.)
%
% A '%' character causes TeX to ignore all remaining text on the line,
% and is used for comments like this one.

\documentclass{article}      % Specifies the document class

                             % The preamble begins here.
\title{An Example Document}  % Declares the document's title.
\author{Leslie Lamport}      % Declares the author's name.
\date{January 21, 1994}      % Deleting this command produces today's date.

\newcommand{\ip}[2]{(#1, #2)}
                             % Defines \ip{arg1}{arg2} to mean
                             % (arg1, arg2).

%\newcommand{\ip}[2]{\langle #1 | #2\rangle}
                             % This is an alternative definition of
                             % \ip that is commented out.

\begin{document}             % End of preamble and beginning of text.

\maketitle                   % Produces the title.

This is an example input file.  Comparing it with
the output it generates can show you how to
produce a simple document of your own.

\section{Ordinary Text}      % Produces section heading.  Lower-level
                             % sections are begun with similar
                             % \subsection and \subsubsection commands.

The ends  of words and sentences are marked
  by   spaces. It  doesn't matter how many
spaces    you type; one is as good as 100.  The
end of   a line counts as a space.

One   or more   blank lines denote the  end
of  a paragraph.

Since any number of consecutive spaces are treated
like a single one, the formatting of the input
file makes no difference to
      \LaTeX,                % The \LaTeX command generates the LaTeX logo.
but it makes a difference to you.  When you use
\LaTeX, making your input file as easy to read
as possible will be a great help as you write
your document and when you change it.  This sample
file shows how you can add comments to your own input
file.

Because printing is different from typewriting,
there are a number of things that you have to do
differently when preparing an input file than if
you were just typing the document directly.
Quotation marks like
       ``this''
have to be handled specially, as do quotes within
quotes:
       ``\,`this'            % \, separates the double and single quote.
        is what I just
        wrote, not  `that'\,''.

Dashes come in three sizes: an
       intra-word
dash, a medium dash for number ranges like
       1--2,
and a punctuation
       dash---like
this.

A sentence-ending space should be larger than the
space between words within a sentence.  You
sometimes have to type special commands in
conjunction with punctuation characters to get
this right, as in the following sentence.
       Gnats, gnus, etc.\ all  % `\ ' makes an inter-word space.
       begin with G\@.         % \@ marks end-of-sentence punctuation.
You should check the spaces after periods when
reading your output to make sure you haven't
forgotten any special cases.  Generating an
ellipsis
       \ldots\               % `\ ' is needed after `\ldots' because TeX
                             % ignores spaces after command names like \ldots
                             % made from \ + letters.
                             %
                             % Note how a `%' character causes TeX to ignore
                             % the end of the input line, so these blank lines
                             % do not start a new paragraph.
                             %
with the right spacing around the periods requires
a special command.

\LaTeX\ interprets some common characters as
commands, so you must type special commands to
generate them.  These characters include the
following:
       \$ \& \% \# \{ and \}.

In printing, text is usually emphasized with an
       \emph{italic}
type style.

\begin{em}
   A long segment of text can also be emphasized
   in this way.  Text within such a segment can be
   given \emph{additional} emphasis.
\end{em}

It is sometimes necessary to prevent \LaTeX\ from
breaking a line where it might otherwise do so.
This may be at a space, as between the ``Mr.''\ and
``Jones'' in
       ``Mr.~Jones'',        % ~ produces an unbreakable interword space.
or within a word---especially when the word is a
symbol like
       \mbox{\emph{itemnum}}
that makes little sense when hyphenated across
lines.

Footnotes\footnote{This is an example of a footnote.}
pose no problem.

\LaTeX\ is good at typesetting mathematical formulas
like
       \( x-3y + z = 7 \)
or
       \( a_{1} > x^{2n} + y^{2n} > x' \)
or
       \( \ip{A}{B} = \sum_{i} a_{i} b_{i} \).
The spaces you type in a formula are
ignored.  Remember that a letter like
       $x$                   % $ ... $  and  \( ... \)  are equivalent
is a formula when it denotes a mathematical
symbol, and it should be typed as one.

\section{Displayed Text}

Text is displayed by indenting it from the left
margin.  Quotations are commonly displayed.  There
are short quotations
\begin{quote}
   This is a short quotation.  It consists of a
   single paragraph of text.  See how it is formatted.
\end{quote}
and longer ones.
\begin{quotation}
   This is a longer quotation.  It consists of two
   paragraphs of text, neither of which are
   particularly interesting.

   This is the second paragraph of the quotation.  It
   is just as dull as the first paragraph.
\end{quotation}
Another frequently-displayed structure is a list.
The following is an example of an \emph{itemized}
list.
\begin{itemize}
   \item This is the first item of an itemized list.
         Each item in the list is marked with a ``tick''.
         You don't have to worry about what kind of tick
         mark is used.

   \item This is the second item of the list.  It
         contains another list nested inside it.  The inner
         list is an \emph{enumerated} list.
         \begin{enumerate}
            \item This is the first item of an enumerated
                  list that is nested within the itemized list.

            \item This is the second item of the inner list.
                  \LaTeX\ allows you to nest lists deeper than
                  you really should.
         \end{enumerate}
         This is the rest of the second item of the outer
         list.  It is no more interesting than any other
         part of the item.
   \item This is the third item of the list.
\end{itemize}
You can even display poetry.
\begin{verse}
   There is an environment
    for verse \\             % The \\ command separates lines
   Whose features some poets % within a stanza.
   will curse.

                             % One or more blank lines separate stanzas.

   For instead of making\\
   Them do \emph{all} line breaking, \\
   It allows them to put too many words on a line when they'd rather be
   forced to be terse.
\end{verse}

Mathematical formulas may also be displayed.  A
displayed formula
is
one-line long; multiline
formulas require special formatting instructions.
   \[  \ip{\Gamma}{\psi'} = x'' + y^{2} + z_{i}^{n}\]
Don't start a paragraph with a displayed equation,
nor make one a paragraph by itself.

\end{document}               % End of document.

\section{Bench Huge Section 55}
% synthetic bench section 55
% This is a sample LaTeX input file.  (Version of 12 August 2004.)
%
% A '%' character causes TeX to ignore all remaining text on the line,
% and is used for comments like this one.

\documentclass{article}      % Specifies the document class

                             % The preamble begins here.
\title{An Example Document}  % Declares the document's title.
\author{Leslie Lamport}      % Declares the author's name.
\date{January 21, 1994}      % Deleting this command produces today's date.

\newcommand{\ip}[2]{(#1, #2)}
                             % Defines \ip{arg1}{arg2} to mean
                             % (arg1, arg2).

%\newcommand{\ip}[2]{\langle #1 | #2\rangle}
                             % This is an alternative definition of
                             % \ip that is commented out.

\begin{document}             % End of preamble and beginning of text.

\maketitle                   % Produces the title.

This is an example input file.  Comparing it with
the output it generates can show you how to
produce a simple document of your own.

\section{Ordinary Text}      % Produces section heading.  Lower-level
                             % sections are begun with similar
                             % \subsection and \subsubsection commands.

The ends  of words and sentences are marked
  by   spaces. It  doesn't matter how many
spaces    you type; one is as good as 100.  The
end of   a line counts as a space.

One   or more   blank lines denote the  end
of  a paragraph.

Since any number of consecutive spaces are treated
like a single one, the formatting of the input
file makes no difference to
      \LaTeX,                % The \LaTeX command generates the LaTeX logo.
but it makes a difference to you.  When you use
\LaTeX, making your input file as easy to read
as possible will be a great help as you write
your document and when you change it.  This sample
file shows how you can add comments to your own input
file.

Because printing is different from typewriting,
there are a number of things that you have to do
differently when preparing an input file than if
you were just typing the document directly.
Quotation marks like
       ``this''
have to be handled specially, as do quotes within
quotes:
       ``\,`this'            % \, separates the double and single quote.
        is what I just
        wrote, not  `that'\,''.

Dashes come in three sizes: an
       intra-word
dash, a medium dash for number ranges like
       1--2,
and a punctuation
       dash---like
this.

A sentence-ending space should be larger than the
space between words within a sentence.  You
sometimes have to type special commands in
conjunction with punctuation characters to get
this right, as in the following sentence.
       Gnats, gnus, etc.\ all  % `\ ' makes an inter-word space.
       begin with G\@.         % \@ marks end-of-sentence punctuation.
You should check the spaces after periods when
reading your output to make sure you haven't
forgotten any special cases.  Generating an
ellipsis
       \ldots\               % `\ ' is needed after `\ldots' because TeX
                             % ignores spaces after command names like \ldots
                             % made from \ + letters.
                             %
                             % Note how a `%' character causes TeX to ignore
                             % the end of the input line, so these blank lines
                             % do not start a new paragraph.
                             %
with the right spacing around the periods requires
a special command.

\LaTeX\ interprets some common characters as
commands, so you must type special commands to
generate them.  These characters include the
following:
       \$ \& \% \# \{ and \}.

In printing, text is usually emphasized with an
       \emph{italic}
type style.

\begin{em}
   A long segment of text can also be emphasized
   in this way.  Text within such a segment can be
   given \emph{additional} emphasis.
\end{em}

It is sometimes necessary to prevent \LaTeX\ from
breaking a line where it might otherwise do so.
This may be at a space, as between the ``Mr.''\ and
``Jones'' in
       ``Mr.~Jones'',        % ~ produces an unbreakable interword space.
or within a word---especially when the word is a
symbol like
       \mbox{\emph{itemnum}}
that makes little sense when hyphenated across
lines.

Footnotes\footnote{This is an example of a footnote.}
pose no problem.

\LaTeX\ is good at typesetting mathematical formulas
like
       \( x-3y + z = 7 \)
or
       \( a_{1} > x^{2n} + y^{2n} > x' \)
or
       \( \ip{A}{B} = \sum_{i} a_{i} b_{i} \).
The spaces you type in a formula are
ignored.  Remember that a letter like
       $x$                   % $ ... $  and  \( ... \)  are equivalent
is a formula when it denotes a mathematical
symbol, and it should be typed as one.

\section{Displayed Text}

Text is displayed by indenting it from the left
margin.  Quotations are commonly displayed.  There
are short quotations
\begin{quote}
   This is a short quotation.  It consists of a
   single paragraph of text.  See how it is formatted.
\end{quote}
and longer ones.
\begin{quotation}
   This is a longer quotation.  It consists of two
   paragraphs of text, neither of which are
   particularly interesting.

   This is the second paragraph of the quotation.  It
   is just as dull as the first paragraph.
\end{quotation}
Another frequently-displayed structure is a list.
The following is an example of an \emph{itemized}
list.
\begin{itemize}
   \item This is the first item of an itemized list.
         Each item in the list is marked with a ``tick''.
         You don't have to worry about what kind of tick
         mark is used.

   \item This is the second item of the list.  It
         contains another list nested inside it.  The inner
         list is an \emph{enumerated} list.
         \begin{enumerate}
            \item This is the first item of an enumerated
                  list that is nested within the itemized list.

            \item This is the second item of the inner list.
                  \LaTeX\ allows you to nest lists deeper than
                  you really should.
         \end{enumerate}
         This is the rest of the second item of the outer
         list.  It is no more interesting than any other
         part of the item.
   \item This is the third item of the list.
\end{itemize}
You can even display poetry.
\begin{verse}
   There is an environment
    for verse \\             % The \\ command separates lines
   Whose features some poets % within a stanza.
   will curse.

                             % One or more blank lines separate stanzas.

   For instead of making\\
   Them do \emph{all} line breaking, \\
   It allows them to put too many words on a line when they'd rather be
   forced to be terse.
\end{verse}

Mathematical formulas may also be displayed.  A
displayed formula
is
one-line long; multiline
formulas require special formatting instructions.
   \[  \ip{\Gamma}{\psi'} = x'' + y^{2} + z_{i}^{n}\]
Don't start a paragraph with a displayed equation,
nor make one a paragraph by itself.

\end{document}               % End of document.

\section{Bench Huge Section 56}
% synthetic bench section 56
% This is a sample LaTeX input file.  (Version of 12 August 2004.)
%
% A '%' character causes TeX to ignore all remaining text on the line,
% and is used for comments like this one.

\documentclass{article}      % Specifies the document class

                             % The preamble begins here.
\title{An Example Document}  % Declares the document's title.
\author{Leslie Lamport}      % Declares the author's name.
\date{January 21, 1994}      % Deleting this command produces today's date.

\newcommand{\ip}[2]{(#1, #2)}
                             % Defines \ip{arg1}{arg2} to mean
                             % (arg1, arg2).

%\newcommand{\ip}[2]{\langle #1 | #2\rangle}
                             % This is an alternative definition of
                             % \ip that is commented out.

\begin{document}             % End of preamble and beginning of text.

\maketitle                   % Produces the title.

This is an example input file.  Comparing it with
the output it generates can show you how to
produce a simple document of your own.

\section{Ordinary Text}      % Produces section heading.  Lower-level
                             % sections are begun with similar
                             % \subsection and \subsubsection commands.

The ends  of words and sentences are marked
  by   spaces. It  doesn't matter how many
spaces    you type; one is as good as 100.  The
end of   a line counts as a space.

One   or more   blank lines denote the  end
of  a paragraph.

Since any number of consecutive spaces are treated
like a single one, the formatting of the input
file makes no difference to
      \LaTeX,                % The \LaTeX command generates the LaTeX logo.
but it makes a difference to you.  When you use
\LaTeX, making your input file as easy to read
as possible will be a great help as you write
your document and when you change it.  This sample
file shows how you can add comments to your own input
file.

Because printing is different from typewriting,
there are a number of things that you have to do
differently when preparing an input file than if
you were just typing the document directly.
Quotation marks like
       ``this''
have to be handled specially, as do quotes within
quotes:
       ``\,`this'            % \, separates the double and single quote.
        is what I just
        wrote, not  `that'\,''.

Dashes come in three sizes: an
       intra-word
dash, a medium dash for number ranges like
       1--2,
and a punctuation
       dash---like
this.

A sentence-ending space should be larger than the
space between words within a sentence.  You
sometimes have to type special commands in
conjunction with punctuation characters to get
this right, as in the following sentence.
       Gnats, gnus, etc.\ all  % `\ ' makes an inter-word space.
       begin with G\@.         % \@ marks end-of-sentence punctuation.
You should check the spaces after periods when
reading your output to make sure you haven't
forgotten any special cases.  Generating an
ellipsis
       \ldots\               % `\ ' is needed after `\ldots' because TeX
                             % ignores spaces after command names like \ldots
                             % made from \ + letters.
                             %
                             % Note how a `%' character causes TeX to ignore
                             % the end of the input line, so these blank lines
                             % do not start a new paragraph.
                             %
with the right spacing around the periods requires
a special command.

\LaTeX\ interprets some common characters as
commands, so you must type special commands to
generate them.  These characters include the
following:
       \$ \& \% \# \{ and \}.

In printing, text is usually emphasized with an
       \emph{italic}
type style.

\begin{em}
   A long segment of text can also be emphasized
   in this way.  Text within such a segment can be
   given \emph{additional} emphasis.
\end{em}

It is sometimes necessary to prevent \LaTeX\ from
breaking a line where it might otherwise do so.
This may be at a space, as between the ``Mr.''\ and
``Jones'' in
       ``Mr.~Jones'',        % ~ produces an unbreakable interword space.
or within a word---especially when the word is a
symbol like
       \mbox{\emph{itemnum}}
that makes little sense when hyphenated across
lines.

Footnotes\footnote{This is an example of a footnote.}
pose no problem.

\LaTeX\ is good at typesetting mathematical formulas
like
       \( x-3y + z = 7 \)
or
       \( a_{1} > x^{2n} + y^{2n} > x' \)
or
       \( \ip{A}{B} = \sum_{i} a_{i} b_{i} \).
The spaces you type in a formula are
ignored.  Remember that a letter like
       $x$                   % $ ... $  and  \( ... \)  are equivalent
is a formula when it denotes a mathematical
symbol, and it should be typed as one.

\section{Displayed Text}

Text is displayed by indenting it from the left
margin.  Quotations are commonly displayed.  There
are short quotations
\begin{quote}
   This is a short quotation.  It consists of a
   single paragraph of text.  See how it is formatted.
\end{quote}
and longer ones.
\begin{quotation}
   This is a longer quotation.  It consists of two
   paragraphs of text, neither of which are
   particularly interesting.

   This is the second paragraph of the quotation.  It
   is just as dull as the first paragraph.
\end{quotation}
Another frequently-displayed structure is a list.
The following is an example of an \emph{itemized}
list.
\begin{itemize}
   \item This is the first item of an itemized list.
         Each item in the list is marked with a ``tick''.
         You don't have to worry about what kind of tick
         mark is used.

   \item This is the second item of the list.  It
         contains another list nested inside it.  The inner
         list is an \emph{enumerated} list.
         \begin{enumerate}
            \item This is the first item of an enumerated
                  list that is nested within the itemized list.

            \item This is the second item of the inner list.
                  \LaTeX\ allows you to nest lists deeper than
                  you really should.
         \end{enumerate}
         This is the rest of the second item of the outer
         list.  It is no more interesting than any other
         part of the item.
   \item This is the third item of the list.
\end{itemize}
You can even display poetry.
\begin{verse}
   There is an environment
    for verse \\             % The \\ command separates lines
   Whose features some poets % within a stanza.
   will curse.

                             % One or more blank lines separate stanzas.

   For instead of making\\
   Them do \emph{all} line breaking, \\
   It allows them to put too many words on a line when they'd rather be
   forced to be terse.
\end{verse}

Mathematical formulas may also be displayed.  A
displayed formula
is
one-line long; multiline
formulas require special formatting instructions.
   \[  \ip{\Gamma}{\psi'} = x'' + y^{2} + z_{i}^{n}\]
Don't start a paragraph with a displayed equation,
nor make one a paragraph by itself.

\end{document}               % End of document.

\section{Bench Huge Section 57}
% synthetic bench section 57
% This is a sample LaTeX input file.  (Version of 12 August 2004.)
%
% A '%' character causes TeX to ignore all remaining text on the line,
% and is used for comments like this one.

\documentclass{article}      % Specifies the document class

                             % The preamble begins here.
\title{An Example Document}  % Declares the document's title.
\author{Leslie Lamport}      % Declares the author's name.
\date{January 21, 1994}      % Deleting this command produces today's date.

\newcommand{\ip}[2]{(#1, #2)}
                             % Defines \ip{arg1}{arg2} to mean
                             % (arg1, arg2).

%\newcommand{\ip}[2]{\langle #1 | #2\rangle}
                             % This is an alternative definition of
                             % \ip that is commented out.

\begin{document}             % End of preamble and beginning of text.

\maketitle                   % Produces the title.

This is an example input file.  Comparing it with
the output it generates can show you how to
produce a simple document of your own.

\section{Ordinary Text}      % Produces section heading.  Lower-level
                             % sections are begun with similar
                             % \subsection and \subsubsection commands.

The ends  of words and sentences are marked
  by   spaces. It  doesn't matter how many
spaces    you type; one is as good as 100.  The
end of   a line counts as a space.

One   or more   blank lines denote the  end
of  a paragraph.

Since any number of consecutive spaces are treated
like a single one, the formatting of the input
file makes no difference to
      \LaTeX,                % The \LaTeX command generates the LaTeX logo.
but it makes a difference to you.  When you use
\LaTeX, making your input file as easy to read
as possible will be a great help as you write
your document and when you change it.  This sample
file shows how you can add comments to your own input
file.

Because printing is different from typewriting,
there are a number of things that you have to do
differently when preparing an input file than if
you were just typing the document directly.
Quotation marks like
       ``this''
have to be handled specially, as do quotes within
quotes:
       ``\,`this'            % \, separates the double and single quote.
        is what I just
        wrote, not  `that'\,''.

Dashes come in three sizes: an
       intra-word
dash, a medium dash for number ranges like
       1--2,
and a punctuation
       dash---like
this.

A sentence-ending space should be larger than the
space between words within a sentence.  You
sometimes have to type special commands in
conjunction with punctuation characters to get
this right, as in the following sentence.
       Gnats, gnus, etc.\ all  % `\ ' makes an inter-word space.
       begin with G\@.         % \@ marks end-of-sentence punctuation.
You should check the spaces after periods when
reading your output to make sure you haven't
forgotten any special cases.  Generating an
ellipsis
       \ldots\               % `\ ' is needed after `\ldots' because TeX
                             % ignores spaces after command names like \ldots
                             % made from \ + letters.
                             %
                             % Note how a `%' character causes TeX to ignore
                             % the end of the input line, so these blank lines
                             % do not start a new paragraph.
                             %
with the right spacing around the periods requires
a special command.

\LaTeX\ interprets some common characters as
commands, so you must type special commands to
generate them.  These characters include the
following:
       \$ \& \% \# \{ and \}.

In printing, text is usually emphasized with an
       \emph{italic}
type style.

\begin{em}
   A long segment of text can also be emphasized
   in this way.  Text within such a segment can be
   given \emph{additional} emphasis.
\end{em}

It is sometimes necessary to prevent \LaTeX\ from
breaking a line where it might otherwise do so.
This may be at a space, as between the ``Mr.''\ and
``Jones'' in
       ``Mr.~Jones'',        % ~ produces an unbreakable interword space.
or within a word---especially when the word is a
symbol like
       \mbox{\emph{itemnum}}
that makes little sense when hyphenated across
lines.

Footnotes\footnote{This is an example of a footnote.}
pose no problem.

\LaTeX\ is good at typesetting mathematical formulas
like
       \( x-3y + z = 7 \)
or
       \( a_{1} > x^{2n} + y^{2n} > x' \)
or
       \( \ip{A}{B} = \sum_{i} a_{i} b_{i} \).
The spaces you type in a formula are
ignored.  Remember that a letter like
       $x$                   % $ ... $  and  \( ... \)  are equivalent
is a formula when it denotes a mathematical
symbol, and it should be typed as one.

\section{Displayed Text}

Text is displayed by indenting it from the left
margin.  Quotations are commonly displayed.  There
are short quotations
\begin{quote}
   This is a short quotation.  It consists of a
   single paragraph of text.  See how it is formatted.
\end{quote}
and longer ones.
\begin{quotation}
   This is a longer quotation.  It consists of two
   paragraphs of text, neither of which are
   particularly interesting.

   This is the second paragraph of the quotation.  It
   is just as dull as the first paragraph.
\end{quotation}
Another frequently-displayed structure is a list.
The following is an example of an \emph{itemized}
list.
\begin{itemize}
   \item This is the first item of an itemized list.
         Each item in the list is marked with a ``tick''.
         You don't have to worry about what kind of tick
         mark is used.

   \item This is the second item of the list.  It
         contains another list nested inside it.  The inner
         list is an \emph{enumerated} list.
         \begin{enumerate}
            \item This is the first item of an enumerated
                  list that is nested within the itemized list.

            \item This is the second item of the inner list.
                  \LaTeX\ allows you to nest lists deeper than
                  you really should.
         \end{enumerate}
         This is the rest of the second item of the outer
         list.  It is no more interesting than any other
         part of the item.
   \item This is the third item of the list.
\end{itemize}
You can even display poetry.
\begin{verse}
   There is an environment
    for verse \\             % The \\ command separates lines
   Whose features some poets % within a stanza.
   will curse.

                             % One or more blank lines separate stanzas.

   For instead of making\\
   Them do \emph{all} line breaking, \\
   It allows them to put too many words on a line when they'd rather be
   forced to be terse.
\end{verse}

Mathematical formulas may also be displayed.  A
displayed formula
is
one-line long; multiline
formulas require special formatting instructions.
   \[  \ip{\Gamma}{\psi'} = x'' + y^{2} + z_{i}^{n}\]
Don't start a paragraph with a displayed equation,
nor make one a paragraph by itself.

\end{document}               % End of document.

\section{Bench Huge Section 58}
% synthetic bench section 58
% This is a sample LaTeX input file.  (Version of 12 August 2004.)
%
% A '%' character causes TeX to ignore all remaining text on the line,
% and is used for comments like this one.

\documentclass{article}      % Specifies the document class

                             % The preamble begins here.
\title{An Example Document}  % Declares the document's title.
\author{Leslie Lamport}      % Declares the author's name.
\date{January 21, 1994}      % Deleting this command produces today's date.

\newcommand{\ip}[2]{(#1, #2)}
                             % Defines \ip{arg1}{arg2} to mean
                             % (arg1, arg2).

%\newcommand{\ip}[2]{\langle #1 | #2\rangle}
                             % This is an alternative definition of
                             % \ip that is commented out.

\begin{document}             % End of preamble and beginning of text.

\maketitle                   % Produces the title.

This is an example input file.  Comparing it with
the output it generates can show you how to
produce a simple document of your own.

\section{Ordinary Text}      % Produces section heading.  Lower-level
                             % sections are begun with similar
                             % \subsection and \subsubsection commands.

The ends  of words and sentences are marked
  by   spaces. It  doesn't matter how many
spaces    you type; one is as good as 100.  The
end of   a line counts as a space.

One   or more   blank lines denote the  end
of  a paragraph.

Since any number of consecutive spaces are treated
like a single one, the formatting of the input
file makes no difference to
      \LaTeX,                % The \LaTeX command generates the LaTeX logo.
but it makes a difference to you.  When you use
\LaTeX, making your input file as easy to read
as possible will be a great help as you write
your document and when you change it.  This sample
file shows how you can add comments to your own input
file.

Because printing is different from typewriting,
there are a number of things that you have to do
differently when preparing an input file than if
you were just typing the document directly.
Quotation marks like
       ``this''
have to be handled specially, as do quotes within
quotes:
       ``\,`this'            % \, separates the double and single quote.
        is what I just
        wrote, not  `that'\,''.

Dashes come in three sizes: an
       intra-word
dash, a medium dash for number ranges like
       1--2,
and a punctuation
       dash---like
this.

A sentence-ending space should be larger than the
space between words within a sentence.  You
sometimes have to type special commands in
conjunction with punctuation characters to get
this right, as in the following sentence.
       Gnats, gnus, etc.\ all  % `\ ' makes an inter-word space.
       begin with G\@.         % \@ marks end-of-sentence punctuation.
You should check the spaces after periods when
reading your output to make sure you haven't
forgotten any special cases.  Generating an
ellipsis
       \ldots\               % `\ ' is needed after `\ldots' because TeX
                             % ignores spaces after command names like \ldots
                             % made from \ + letters.
                             %
                             % Note how a `%' character causes TeX to ignore
                             % the end of the input line, so these blank lines
                             % do not start a new paragraph.
                             %
with the right spacing around the periods requires
a special command.

\LaTeX\ interprets some common characters as
commands, so you must type special commands to
generate them.  These characters include the
following:
       \$ \& \% \# \{ and \}.

In printing, text is usually emphasized with an
       \emph{italic}
type style.

\begin{em}
   A long segment of text can also be emphasized
   in this way.  Text within such a segment can be
   given \emph{additional} emphasis.
\end{em}

It is sometimes necessary to prevent \LaTeX\ from
breaking a line where it might otherwise do so.
This may be at a space, as between the ``Mr.''\ and
``Jones'' in
       ``Mr.~Jones'',        % ~ produces an unbreakable interword space.
or within a word---especially when the word is a
symbol like
       \mbox{\emph{itemnum}}
that makes little sense when hyphenated across
lines.

Footnotes\footnote{This is an example of a footnote.}
pose no problem.

\LaTeX\ is good at typesetting mathematical formulas
like
       \( x-3y + z = 7 \)
or
       \( a_{1} > x^{2n} + y^{2n} > x' \)
or
       \( \ip{A}{B} = \sum_{i} a_{i} b_{i} \).
The spaces you type in a formula are
ignored.  Remember that a letter like
       $x$                   % $ ... $  and  \( ... \)  are equivalent
is a formula when it denotes a mathematical
symbol, and it should be typed as one.

\section{Displayed Text}

Text is displayed by indenting it from the left
margin.  Quotations are commonly displayed.  There
are short quotations
\begin{quote}
   This is a short quotation.  It consists of a
   single paragraph of text.  See how it is formatted.
\end{quote}
and longer ones.
\begin{quotation}
   This is a longer quotation.  It consists of two
   paragraphs of text, neither of which are
   particularly interesting.

   This is the second paragraph of the quotation.  It
   is just as dull as the first paragraph.
\end{quotation}
Another frequently-displayed structure is a list.
The following is an example of an \emph{itemized}
list.
\begin{itemize}
   \item This is the first item of an itemized list.
         Each item in the list is marked with a ``tick''.
         You don't have to worry about what kind of tick
         mark is used.

   \item This is the second item of the list.  It
         contains another list nested inside it.  The inner
         list is an \emph{enumerated} list.
         \begin{enumerate}
            \item This is the first item of an enumerated
                  list that is nested within the itemized list.

            \item This is the second item of the inner list.
                  \LaTeX\ allows you to nest lists deeper than
                  you really should.
         \end{enumerate}
         This is the rest of the second item of the outer
         list.  It is no more interesting than any other
         part of the item.
   \item This is the third item of the list.
\end{itemize}
You can even display poetry.
\begin{verse}
   There is an environment
    for verse \\             % The \\ command separates lines
   Whose features some poets % within a stanza.
   will curse.

                             % One or more blank lines separate stanzas.

   For instead of making\\
   Them do \emph{all} line breaking, \\
   It allows them to put too many words on a line when they'd rather be
   forced to be terse.
\end{verse}

Mathematical formulas may also be displayed.  A
displayed formula
is
one-line long; multiline
formulas require special formatting instructions.
   \[  \ip{\Gamma}{\psi'} = x'' + y^{2} + z_{i}^{n}\]
Don't start a paragraph with a displayed equation,
nor make one a paragraph by itself.

\end{document}               % End of document.

\section{Bench Huge Section 59}
% synthetic bench section 59
% This is a sample LaTeX input file.  (Version of 12 August 2004.)
%
% A '%' character causes TeX to ignore all remaining text on the line,
% and is used for comments like this one.

\documentclass{article}      % Specifies the document class

                             % The preamble begins here.
\title{An Example Document}  % Declares the document's title.
\author{Leslie Lamport}      % Declares the author's name.
\date{January 21, 1994}      % Deleting this command produces today's date.

\newcommand{\ip}[2]{(#1, #2)}
                             % Defines \ip{arg1}{arg2} to mean
                             % (arg1, arg2).

%\newcommand{\ip}[2]{\langle #1 | #2\rangle}
                             % This is an alternative definition of
                             % \ip that is commented out.

\begin{document}             % End of preamble and beginning of text.

\maketitle                   % Produces the title.

This is an example input file.  Comparing it with
the output it generates can show you how to
produce a simple document of your own.

\section{Ordinary Text}      % Produces section heading.  Lower-level
                             % sections are begun with similar
                             % \subsection and \subsubsection commands.

The ends  of words and sentences are marked
  by   spaces. It  doesn't matter how many
spaces    you type; one is as good as 100.  The
end of   a line counts as a space.

One   or more   blank lines denote the  end
of  a paragraph.

Since any number of consecutive spaces are treated
like a single one, the formatting of the input
file makes no difference to
      \LaTeX,                % The \LaTeX command generates the LaTeX logo.
but it makes a difference to you.  When you use
\LaTeX, making your input file as easy to read
as possible will be a great help as you write
your document and when you change it.  This sample
file shows how you can add comments to your own input
file.

Because printing is different from typewriting,
there are a number of things that you have to do
differently when preparing an input file than if
you were just typing the document directly.
Quotation marks like
       ``this''
have to be handled specially, as do quotes within
quotes:
       ``\,`this'            % \, separates the double and single quote.
        is what I just
        wrote, not  `that'\,''.

Dashes come in three sizes: an
       intra-word
dash, a medium dash for number ranges like
       1--2,
and a punctuation
       dash---like
this.

A sentence-ending space should be larger than the
space between words within a sentence.  You
sometimes have to type special commands in
conjunction with punctuation characters to get
this right, as in the following sentence.
       Gnats, gnus, etc.\ all  % `\ ' makes an inter-word space.
       begin with G\@.         % \@ marks end-of-sentence punctuation.
You should check the spaces after periods when
reading your output to make sure you haven't
forgotten any special cases.  Generating an
ellipsis
       \ldots\               % `\ ' is needed after `\ldots' because TeX
                             % ignores spaces after command names like \ldots
                             % made from \ + letters.
                             %
                             % Note how a `%' character causes TeX to ignore
                             % the end of the input line, so these blank lines
                             % do not start a new paragraph.
                             %
with the right spacing around the periods requires
a special command.

\LaTeX\ interprets some common characters as
commands, so you must type special commands to
generate them.  These characters include the
following:
       \$ \& \% \# \{ and \}.

In printing, text is usually emphasized with an
       \emph{italic}
type style.

\begin{em}
   A long segment of text can also be emphasized
   in this way.  Text within such a segment can be
   given \emph{additional} emphasis.
\end{em}

It is sometimes necessary to prevent \LaTeX\ from
breaking a line where it might otherwise do so.
This may be at a space, as between the ``Mr.''\ and
``Jones'' in
       ``Mr.~Jones'',        % ~ produces an unbreakable interword space.
or within a word---especially when the word is a
symbol like
       \mbox{\emph{itemnum}}
that makes little sense when hyphenated across
lines.

Footnotes\footnote{This is an example of a footnote.}
pose no problem.

\LaTeX\ is good at typesetting mathematical formulas
like
       \( x-3y + z = 7 \)
or
       \( a_{1} > x^{2n} + y^{2n} > x' \)
or
       \( \ip{A}{B} = \sum_{i} a_{i} b_{i} \).
The spaces you type in a formula are
ignored.  Remember that a letter like
       $x$                   % $ ... $  and  \( ... \)  are equivalent
is a formula when it denotes a mathematical
symbol, and it should be typed as one.

\section{Displayed Text}

Text is displayed by indenting it from the left
margin.  Quotations are commonly displayed.  There
are short quotations
\begin{quote}
   This is a short quotation.  It consists of a
   single paragraph of text.  See how it is formatted.
\end{quote}
and longer ones.
\begin{quotation}
   This is a longer quotation.  It consists of two
   paragraphs of text, neither of which are
   particularly interesting.

   This is the second paragraph of the quotation.  It
   is just as dull as the first paragraph.
\end{quotation}
Another frequently-displayed structure is a list.
The following is an example of an \emph{itemized}
list.
\begin{itemize}
   \item This is the first item of an itemized list.
         Each item in the list is marked with a ``tick''.
         You don't have to worry about what kind of tick
         mark is used.

   \item This is the second item of the list.  It
         contains another list nested inside it.  The inner
         list is an \emph{enumerated} list.
         \begin{enumerate}
            \item This is the first item of an enumerated
                  list that is nested within the itemized list.

            \item This is the second item of the inner list.
                  \LaTeX\ allows you to nest lists deeper than
                  you really should.
         \end{enumerate}
         This is the rest of the second item of the outer
         list.  It is no more interesting than any other
         part of the item.
   \item This is the third item of the list.
\end{itemize}
You can even display poetry.
\begin{verse}
   There is an environment
    for verse \\             % The \\ command separates lines
   Whose features some poets % within a stanza.
   will curse.

                             % One or more blank lines separate stanzas.

   For instead of making\\
   Them do \emph{all} line breaking, \\
   It allows them to put too many words on a line when they'd rather be
   forced to be terse.
\end{verse}

Mathematical formulas may also be displayed.  A
displayed formula
is
one-line long; multiline
formulas require special formatting instructions.
   \[  \ip{\Gamma}{\psi'} = x'' + y^{2} + z_{i}^{n}\]
Don't start a paragraph with a displayed equation,
nor make one a paragraph by itself.

\end{document}               % End of document.

\section{Bench Huge Section 60}
% synthetic bench section 60
% This is a sample LaTeX input file.  (Version of 12 August 2004.)
%
% A '%' character causes TeX to ignore all remaining text on the line,
% and is used for comments like this one.

\documentclass{article}      % Specifies the document class

                             % The preamble begins here.
\title{An Example Document}  % Declares the document's title.
\author{Leslie Lamport}      % Declares the author's name.
\date{January 21, 1994}      % Deleting this command produces today's date.

\newcommand{\ip}[2]{(#1, #2)}
                             % Defines \ip{arg1}{arg2} to mean
                             % (arg1, arg2).

%\newcommand{\ip}[2]{\langle #1 | #2\rangle}
                             % This is an alternative definition of
                             % \ip that is commented out.

\begin{document}             % End of preamble and beginning of text.

\maketitle                   % Produces the title.

This is an example input file.  Comparing it with
the output it generates can show you how to
produce a simple document of your own.

\section{Ordinary Text}      % Produces section heading.  Lower-level
                             % sections are begun with similar
                             % \subsection and \subsubsection commands.

The ends  of words and sentences are marked
  by   spaces. It  doesn't matter how many
spaces    you type; one is as good as 100.  The
end of   a line counts as a space.

One   or more   blank lines denote the  end
of  a paragraph.

Since any number of consecutive spaces are treated
like a single one, the formatting of the input
file makes no difference to
      \LaTeX,                % The \LaTeX command generates the LaTeX logo.
but it makes a difference to you.  When you use
\LaTeX, making your input file as easy to read
as possible will be a great help as you write
your document and when you change it.  This sample
file shows how you can add comments to your own input
file.

Because printing is different from typewriting,
there are a number of things that you have to do
differently when preparing an input file than if
you were just typing the document directly.
Quotation marks like
       ``this''
have to be handled specially, as do quotes within
quotes:
       ``\,`this'            % \, separates the double and single quote.
        is what I just
        wrote, not  `that'\,''.

Dashes come in three sizes: an
       intra-word
dash, a medium dash for number ranges like
       1--2,
and a punctuation
       dash---like
this.

A sentence-ending space should be larger than the
space between words within a sentence.  You
sometimes have to type special commands in
conjunction with punctuation characters to get
this right, as in the following sentence.
       Gnats, gnus, etc.\ all  % `\ ' makes an inter-word space.
       begin with G\@.         % \@ marks end-of-sentence punctuation.
You should check the spaces after periods when
reading your output to make sure you haven't
forgotten any special cases.  Generating an
ellipsis
       \ldots\               % `\ ' is needed after `\ldots' because TeX
                             % ignores spaces after command names like \ldots
                             % made from \ + letters.
                             %
                             % Note how a `%' character causes TeX to ignore
                             % the end of the input line, so these blank lines
                             % do not start a new paragraph.
                             %
with the right spacing around the periods requires
a special command.

\LaTeX\ interprets some common characters as
commands, so you must type special commands to
generate them.  These characters include the
following:
       \$ \& \% \# \{ and \}.

In printing, text is usually emphasized with an
       \emph{italic}
type style.

\begin{em}
   A long segment of text can also be emphasized
   in this way.  Text within such a segment can be
   given \emph{additional} emphasis.
\end{em}

It is sometimes necessary to prevent \LaTeX\ from
breaking a line where it might otherwise do so.
This may be at a space, as between the ``Mr.''\ and
``Jones'' in
       ``Mr.~Jones'',        % ~ produces an unbreakable interword space.
or within a word---especially when the word is a
symbol like
       \mbox{\emph{itemnum}}
that makes little sense when hyphenated across
lines.

Footnotes\footnote{This is an example of a footnote.}
pose no problem.

\LaTeX\ is good at typesetting mathematical formulas
like
       \( x-3y + z = 7 \)
or
       \( a_{1} > x^{2n} + y^{2n} > x' \)
or
       \( \ip{A}{B} = \sum_{i} a_{i} b_{i} \).
The spaces you type in a formula are
ignored.  Remember that a letter like
       $x$                   % $ ... $  and  \( ... \)  are equivalent
is a formula when it denotes a mathematical
symbol, and it should be typed as one.

\section{Displayed Text}

Text is displayed by indenting it from the left
margin.  Quotations are commonly displayed.  There
are short quotations
\begin{quote}
   This is a short quotation.  It consists of a
   single paragraph of text.  See how it is formatted.
\end{quote}
and longer ones.
\begin{quotation}
   This is a longer quotation.  It consists of two
   paragraphs of text, neither of which are
   particularly interesting.

   This is the second paragraph of the quotation.  It
   is just as dull as the first paragraph.
\end{quotation}
Another frequently-displayed structure is a list.
The following is an example of an \emph{itemized}
list.
\begin{itemize}
   \item This is the first item of an itemized list.
         Each item in the list is marked with a ``tick''.
         You don't have to worry about what kind of tick
         mark is used.

   \item This is the second item of the list.  It
         contains another list nested inside it.  The inner
         list is an \emph{enumerated} list.
         \begin{enumerate}
            \item This is the first item of an enumerated
                  list that is nested within the itemized list.

            \item This is the second item of the inner list.
                  \LaTeX\ allows you to nest lists deeper than
                  you really should.
         \end{enumerate}
         This is the rest of the second item of the outer
         list.  It is no more interesting than any other
         part of the item.
   \item This is the third item of the list.
\end{itemize}
You can even display poetry.
\begin{verse}
   There is an environment
    for verse \\             % The \\ command separates lines
   Whose features some poets % within a stanza.
   will curse.

                             % One or more blank lines separate stanzas.

   For instead of making\\
   Them do \emph{all} line breaking, \\
   It allows them to put too many words on a line when they'd rather be
   forced to be terse.
\end{verse}

Mathematical formulas may also be displayed.  A
displayed formula
is
one-line long; multiline
formulas require special formatting instructions.
   \[  \ip{\Gamma}{\psi'} = x'' + y^{2} + z_{i}^{n}\]
Don't start a paragraph with a displayed equation,
nor make one a paragraph by itself.

\end{document}               % End of document.

\section{Bench Huge Section 61}
% synthetic bench section 61
% This is a sample LaTeX input file.  (Version of 12 August 2004.)
%
% A '%' character causes TeX to ignore all remaining text on the line,
% and is used for comments like this one.

\documentclass{article}      % Specifies the document class

                             % The preamble begins here.
\title{An Example Document}  % Declares the document's title.
\author{Leslie Lamport}      % Declares the author's name.
\date{January 21, 1994}      % Deleting this command produces today's date.

\newcommand{\ip}[2]{(#1, #2)}
                             % Defines \ip{arg1}{arg2} to mean
                             % (arg1, arg2).

%\newcommand{\ip}[2]{\langle #1 | #2\rangle}
                             % This is an alternative definition of
                             % \ip that is commented out.

\begin{document}             % End of preamble and beginning of text.

\maketitle                   % Produces the title.

This is an example input file.  Comparing it with
the output it generates can show you how to
produce a simple document of your own.

\section{Ordinary Text}      % Produces section heading.  Lower-level
                             % sections are begun with similar
                             % \subsection and \subsubsection commands.

The ends  of words and sentences are marked
  by   spaces. It  doesn't matter how many
spaces    you type; one is as good as 100.  The
end of   a line counts as a space.

One   or more   blank lines denote the  end
of  a paragraph.

Since any number of consecutive spaces are treated
like a single one, the formatting of the input
file makes no difference to
      \LaTeX,                % The \LaTeX command generates the LaTeX logo.
but it makes a difference to you.  When you use
\LaTeX, making your input file as easy to read
as possible will be a great help as you write
your document and when you change it.  This sample
file shows how you can add comments to your own input
file.

Because printing is different from typewriting,
there are a number of things that you have to do
differently when preparing an input file than if
you were just typing the document directly.
Quotation marks like
       ``this''
have to be handled specially, as do quotes within
quotes:
       ``\,`this'            % \, separates the double and single quote.
        is what I just
        wrote, not  `that'\,''.

Dashes come in three sizes: an
       intra-word
dash, a medium dash for number ranges like
       1--2,
and a punctuation
       dash---like
this.

A sentence-ending space should be larger than the
space between words within a sentence.  You
sometimes have to type special commands in
conjunction with punctuation characters to get
this right, as in the following sentence.
       Gnats, gnus, etc.\ all  % `\ ' makes an inter-word space.
       begin with G\@.         % \@ marks end-of-sentence punctuation.
You should check the spaces after periods when
reading your output to make sure you haven't
forgotten any special cases.  Generating an
ellipsis
       \ldots\               % `\ ' is needed after `\ldots' because TeX
                             % ignores spaces after command names like \ldots
                             % made from \ + letters.
                             %
                             % Note how a `%' character causes TeX to ignore
                             % the end of the input line, so these blank lines
                             % do not start a new paragraph.
                             %
with the right spacing around the periods requires
a special command.

\LaTeX\ interprets some common characters as
commands, so you must type special commands to
generate them.  These characters include the
following:
       \$ \& \% \# \{ and \}.

In printing, text is usually emphasized with an
       \emph{italic}
type style.

\begin{em}
   A long segment of text can also be emphasized
   in this way.  Text within such a segment can be
   given \emph{additional} emphasis.
\end{em}

It is sometimes necessary to prevent \LaTeX\ from
breaking a line where it might otherwise do so.
This may be at a space, as between the ``Mr.''\ and
``Jones'' in
       ``Mr.~Jones'',        % ~ produces an unbreakable interword space.
or within a word---especially when the word is a
symbol like
       \mbox{\emph{itemnum}}
that makes little sense when hyphenated across
lines.

Footnotes\footnote{This is an example of a footnote.}
pose no problem.

\LaTeX\ is good at typesetting mathematical formulas
like
       \( x-3y + z = 7 \)
or
       \( a_{1} > x^{2n} + y^{2n} > x' \)
or
       \( \ip{A}{B} = \sum_{i} a_{i} b_{i} \).
The spaces you type in a formula are
ignored.  Remember that a letter like
       $x$                   % $ ... $  and  \( ... \)  are equivalent
is a formula when it denotes a mathematical
symbol, and it should be typed as one.

\section{Displayed Text}

Text is displayed by indenting it from the left
margin.  Quotations are commonly displayed.  There
are short quotations
\begin{quote}
   This is a short quotation.  It consists of a
   single paragraph of text.  See how it is formatted.
\end{quote}
and longer ones.
\begin{quotation}
   This is a longer quotation.  It consists of two
   paragraphs of text, neither of which are
   particularly interesting.

   This is the second paragraph of the quotation.  It
   is just as dull as the first paragraph.
\end{quotation}
Another frequently-displayed structure is a list.
The following is an example of an \emph{itemized}
list.
\begin{itemize}
   \item This is the first item of an itemized list.
         Each item in the list is marked with a ``tick''.
         You don't have to worry about what kind of tick
         mark is used.

   \item This is the second item of the list.  It
         contains another list nested inside it.  The inner
         list is an \emph{enumerated} list.
         \begin{enumerate}
            \item This is the first item of an enumerated
                  list that is nested within the itemized list.

            \item This is the second item of the inner list.
                  \LaTeX\ allows you to nest lists deeper than
                  you really should.
         \end{enumerate}
         This is the rest of the second item of the outer
         list.  It is no more interesting than any other
         part of the item.
   \item This is the third item of the list.
\end{itemize}
You can even display poetry.
\begin{verse}
   There is an environment
    for verse \\             % The \\ command separates lines
   Whose features some poets % within a stanza.
   will curse.

                             % One or more blank lines separate stanzas.

   For instead of making\\
   Them do \emph{all} line breaking, \\
   It allows them to put too many words on a line when they'd rather be
   forced to be terse.
\end{verse}

Mathematical formulas may also be displayed.  A
displayed formula
is
one-line long; multiline
formulas require special formatting instructions.
   \[  \ip{\Gamma}{\psi'} = x'' + y^{2} + z_{i}^{n}\]
Don't start a paragraph with a displayed equation,
nor make one a paragraph by itself.

\end{document}               % End of document.

\section{Bench Huge Section 62}
% synthetic bench section 62
% This is a sample LaTeX input file.  (Version of 12 August 2004.)
%
% A '%' character causes TeX to ignore all remaining text on the line,
% and is used for comments like this one.

\documentclass{article}      % Specifies the document class

                             % The preamble begins here.
\title{An Example Document}  % Declares the document's title.
\author{Leslie Lamport}      % Declares the author's name.
\date{January 21, 1994}      % Deleting this command produces today's date.

\newcommand{\ip}[2]{(#1, #2)}
                             % Defines \ip{arg1}{arg2} to mean
                             % (arg1, arg2).

%\newcommand{\ip}[2]{\langle #1 | #2\rangle}
                             % This is an alternative definition of
                             % \ip that is commented out.

\begin{document}             % End of preamble and beginning of text.

\maketitle                   % Produces the title.

This is an example input file.  Comparing it with
the output it generates can show you how to
produce a simple document of your own.

\section{Ordinary Text}      % Produces section heading.  Lower-level
                             % sections are begun with similar
                             % \subsection and \subsubsection commands.

The ends  of words and sentences are marked
  by   spaces. It  doesn't matter how many
spaces    you type; one is as good as 100.  The
end of   a line counts as a space.

One   or more   blank lines denote the  end
of  a paragraph.

Since any number of consecutive spaces are treated
like a single one, the formatting of the input
file makes no difference to
      \LaTeX,                % The \LaTeX command generates the LaTeX logo.
but it makes a difference to you.  When you use
\LaTeX, making your input file as easy to read
as possible will be a great help as you write
your document and when you change it.  This sample
file shows how you can add comments to your own input
file.

Because printing is different from typewriting,
there are a number of things that you have to do
differently when preparing an input file than if
you were just typing the document directly.
Quotation marks like
       ``this''
have to be handled specially, as do quotes within
quotes:
       ``\,`this'            % \, separates the double and single quote.
        is what I just
        wrote, not  `that'\,''.

Dashes come in three sizes: an
       intra-word
dash, a medium dash for number ranges like
       1--2,
and a punctuation
       dash---like
this.

A sentence-ending space should be larger than the
space between words within a sentence.  You
sometimes have to type special commands in
conjunction with punctuation characters to get
this right, as in the following sentence.
       Gnats, gnus, etc.\ all  % `\ ' makes an inter-word space.
       begin with G\@.         % \@ marks end-of-sentence punctuation.
You should check the spaces after periods when
reading your output to make sure you haven't
forgotten any special cases.  Generating an
ellipsis
       \ldots\               % `\ ' is needed after `\ldots' because TeX
                             % ignores spaces after command names like \ldots
                             % made from \ + letters.
                             %
                             % Note how a `%' character causes TeX to ignore
                             % the end of the input line, so these blank lines
                             % do not start a new paragraph.
                             %
with the right spacing around the periods requires
a special command.

\LaTeX\ interprets some common characters as
commands, so you must type special commands to
generate them.  These characters include the
following:
       \$ \& \% \# \{ and \}.

In printing, text is usually emphasized with an
       \emph{italic}
type style.

\begin{em}
   A long segment of text can also be emphasized
   in this way.  Text within such a segment can be
   given \emph{additional} emphasis.
\end{em}

It is sometimes necessary to prevent \LaTeX\ from
breaking a line where it might otherwise do so.
This may be at a space, as between the ``Mr.''\ and
``Jones'' in
       ``Mr.~Jones'',        % ~ produces an unbreakable interword space.
or within a word---especially when the word is a
symbol like
       \mbox{\emph{itemnum}}
that makes little sense when hyphenated across
lines.

Footnotes\footnote{This is an example of a footnote.}
pose no problem.

\LaTeX\ is good at typesetting mathematical formulas
like
       \( x-3y + z = 7 \)
or
       \( a_{1} > x^{2n} + y^{2n} > x' \)
or
       \( \ip{A}{B} = \sum_{i} a_{i} b_{i} \).
The spaces you type in a formula are
ignored.  Remember that a letter like
       $x$                   % $ ... $  and  \( ... \)  are equivalent
is a formula when it denotes a mathematical
symbol, and it should be typed as one.

\section{Displayed Text}

Text is displayed by indenting it from the left
margin.  Quotations are commonly displayed.  There
are short quotations
\begin{quote}
   This is a short quotation.  It consists of a
   single paragraph of text.  See how it is formatted.
\end{quote}
and longer ones.
\begin{quotation}
   This is a longer quotation.  It consists of two
   paragraphs of text, neither of which are
   particularly interesting.

   This is the second paragraph of the quotation.  It
   is just as dull as the first paragraph.
\end{quotation}
Another frequently-displayed structure is a list.
The following is an example of an \emph{itemized}
list.
\begin{itemize}
   \item This is the first item of an itemized list.
         Each item in the list is marked with a ``tick''.
         You don't have to worry about what kind of tick
         mark is used.

   \item This is the second item of the list.  It
         contains another list nested inside it.  The inner
         list is an \emph{enumerated} list.
         \begin{enumerate}
            \item This is the first item of an enumerated
                  list that is nested within the itemized list.

            \item This is the second item of the inner list.
                  \LaTeX\ allows you to nest lists deeper than
                  you really should.
         \end{enumerate}
         This is the rest of the second item of the outer
         list.  It is no more interesting than any other
         part of the item.
   \item This is the third item of the list.
\end{itemize}
You can even display poetry.
\begin{verse}
   There is an environment
    for verse \\             % The \\ command separates lines
   Whose features some poets % within a stanza.
   will curse.

                             % One or more blank lines separate stanzas.

   For instead of making\\
   Them do \emph{all} line breaking, \\
   It allows them to put too many words on a line when they'd rather be
   forced to be terse.
\end{verse}

Mathematical formulas may also be displayed.  A
displayed formula
is
one-line long; multiline
formulas require special formatting instructions.
   \[  \ip{\Gamma}{\psi'} = x'' + y^{2} + z_{i}^{n}\]
Don't start a paragraph with a displayed equation,
nor make one a paragraph by itself.

\end{document}               % End of document.

\section{Bench Huge Section 63}
% synthetic bench section 63
% This is a sample LaTeX input file.  (Version of 12 August 2004.)
%
% A '%' character causes TeX to ignore all remaining text on the line,
% and is used for comments like this one.

\documentclass{article}      % Specifies the document class

                             % The preamble begins here.
\title{An Example Document}  % Declares the document's title.
\author{Leslie Lamport}      % Declares the author's name.
\date{January 21, 1994}      % Deleting this command produces today's date.

\newcommand{\ip}[2]{(#1, #2)}
                             % Defines \ip{arg1}{arg2} to mean
                             % (arg1, arg2).

%\newcommand{\ip}[2]{\langle #1 | #2\rangle}
                             % This is an alternative definition of
                             % \ip that is commented out.

\begin{document}             % End of preamble and beginning of text.

\maketitle                   % Produces the title.

This is an example input file.  Comparing it with
the output it generates can show you how to
produce a simple document of your own.

\section{Ordinary Text}      % Produces section heading.  Lower-level
                             % sections are begun with similar
                             % \subsection and \subsubsection commands.

The ends  of words and sentences are marked
  by   spaces. It  doesn't matter how many
spaces    you type; one is as good as 100.  The
end of   a line counts as a space.

One   or more   blank lines denote the  end
of  a paragraph.

Since any number of consecutive spaces are treated
like a single one, the formatting of the input
file makes no difference to
      \LaTeX,                % The \LaTeX command generates the LaTeX logo.
but it makes a difference to you.  When you use
\LaTeX, making your input file as easy to read
as possible will be a great help as you write
your document and when you change it.  This sample
file shows how you can add comments to your own input
file.

Because printing is different from typewriting,
there are a number of things that you have to do
differently when preparing an input file than if
you were just typing the document directly.
Quotation marks like
       ``this''
have to be handled specially, as do quotes within
quotes:
       ``\,`this'            % \, separates the double and single quote.
        is what I just
        wrote, not  `that'\,''.

Dashes come in three sizes: an
       intra-word
dash, a medium dash for number ranges like
       1--2,
and a punctuation
       dash---like
this.

A sentence-ending space should be larger than the
space between words within a sentence.  You
sometimes have to type special commands in
conjunction with punctuation characters to get
this right, as in the following sentence.
       Gnats, gnus, etc.\ all  % `\ ' makes an inter-word space.
       begin with G\@.         % \@ marks end-of-sentence punctuation.
You should check the spaces after periods when
reading your output to make sure you haven't
forgotten any special cases.  Generating an
ellipsis
       \ldots\               % `\ ' is needed after `\ldots' because TeX
                             % ignores spaces after command names like \ldots
                             % made from \ + letters.
                             %
                             % Note how a `%' character causes TeX to ignore
                             % the end of the input line, so these blank lines
                             % do not start a new paragraph.
                             %
with the right spacing around the periods requires
a special command.

\LaTeX\ interprets some common characters as
commands, so you must type special commands to
generate them.  These characters include the
following:
       \$ \& \% \# \{ and \}.

In printing, text is usually emphasized with an
       \emph{italic}
type style.

\begin{em}
   A long segment of text can also be emphasized
   in this way.  Text within such a segment can be
   given \emph{additional} emphasis.
\end{em}

It is sometimes necessary to prevent \LaTeX\ from
breaking a line where it might otherwise do so.
This may be at a space, as between the ``Mr.''\ and
``Jones'' in
       ``Mr.~Jones'',        % ~ produces an unbreakable interword space.
or within a word---especially when the word is a
symbol like
       \mbox{\emph{itemnum}}
that makes little sense when hyphenated across
lines.

Footnotes\footnote{This is an example of a footnote.}
pose no problem.

\LaTeX\ is good at typesetting mathematical formulas
like
       \( x-3y + z = 7 \)
or
       \( a_{1} > x^{2n} + y^{2n} > x' \)
or
       \( \ip{A}{B} = \sum_{i} a_{i} b_{i} \).
The spaces you type in a formula are
ignored.  Remember that a letter like
       $x$                   % $ ... $  and  \( ... \)  are equivalent
is a formula when it denotes a mathematical
symbol, and it should be typed as one.

\section{Displayed Text}

Text is displayed by indenting it from the left
margin.  Quotations are commonly displayed.  There
are short quotations
\begin{quote}
   This is a short quotation.  It consists of a
   single paragraph of text.  See how it is formatted.
\end{quote}
and longer ones.
\begin{quotation}
   This is a longer quotation.  It consists of two
   paragraphs of text, neither of which are
   particularly interesting.

   This is the second paragraph of the quotation.  It
   is just as dull as the first paragraph.
\end{quotation}
Another frequently-displayed structure is a list.
The following is an example of an \emph{itemized}
list.
\begin{itemize}
   \item This is the first item of an itemized list.
         Each item in the list is marked with a ``tick''.
         You don't have to worry about what kind of tick
         mark is used.

   \item This is the second item of the list.  It
         contains another list nested inside it.  The inner
         list is an \emph{enumerated} list.
         \begin{enumerate}
            \item This is the first item of an enumerated
                  list that is nested within the itemized list.

            \item This is the second item of the inner list.
                  \LaTeX\ allows you to nest lists deeper than
                  you really should.
         \end{enumerate}
         This is the rest of the second item of the outer
         list.  It is no more interesting than any other
         part of the item.
   \item This is the third item of the list.
\end{itemize}
You can even display poetry.
\begin{verse}
   There is an environment
    for verse \\             % The \\ command separates lines
   Whose features some poets % within a stanza.
   will curse.

                             % One or more blank lines separate stanzas.

   For instead of making\\
   Them do \emph{all} line breaking, \\
   It allows them to put too many words on a line when they'd rather be
   forced to be terse.
\end{verse}

Mathematical formulas may also be displayed.  A
displayed formula
is
one-line long; multiline
formulas require special formatting instructions.
   \[  \ip{\Gamma}{\psi'} = x'' + y^{2} + z_{i}^{n}\]
Don't start a paragraph with a displayed equation,
nor make one a paragraph by itself.

\end{document}               % End of document.

\section{Bench Huge Section 64}
% synthetic bench section 64
% This is a sample LaTeX input file.  (Version of 12 August 2004.)
%
% A '%' character causes TeX to ignore all remaining text on the line,
% and is used for comments like this one.

\documentclass{article}      % Specifies the document class

                             % The preamble begins here.
\title{An Example Document}  % Declares the document's title.
\author{Leslie Lamport}      % Declares the author's name.
\date{January 21, 1994}      % Deleting this command produces today's date.

\newcommand{\ip}[2]{(#1, #2)}
                             % Defines \ip{arg1}{arg2} to mean
                             % (arg1, arg2).

%\newcommand{\ip}[2]{\langle #1 | #2\rangle}
                             % This is an alternative definition of
                             % \ip that is commented out.

\begin{document}             % End of preamble and beginning of text.

\maketitle                   % Produces the title.

This is an example input file.  Comparing it with
the output it generates can show you how to
produce a simple document of your own.

\section{Ordinary Text}      % Produces section heading.  Lower-level
                             % sections are begun with similar
                             % \subsection and \subsubsection commands.

The ends  of words and sentences are marked
  by   spaces. It  doesn't matter how many
spaces    you type; one is as good as 100.  The
end of   a line counts as a space.

One   or more   blank lines denote the  end
of  a paragraph.

Since any number of consecutive spaces are treated
like a single one, the formatting of the input
file makes no difference to
      \LaTeX,                % The \LaTeX command generates the LaTeX logo.
but it makes a difference to you.  When you use
\LaTeX, making your input file as easy to read
as possible will be a great help as you write
your document and when you change it.  This sample
file shows how you can add comments to your own input
file.

Because printing is different from typewriting,
there are a number of things that you have to do
differently when preparing an input file than if
you were just typing the document directly.
Quotation marks like
       ``this''
have to be handled specially, as do quotes within
quotes:
       ``\,`this'            % \, separates the double and single quote.
        is what I just
        wrote, not  `that'\,''.

Dashes come in three sizes: an
       intra-word
dash, a medium dash for number ranges like
       1--2,
and a punctuation
       dash---like
this.

A sentence-ending space should be larger than the
space between words within a sentence.  You
sometimes have to type special commands in
conjunction with punctuation characters to get
this right, as in the following sentence.
       Gnats, gnus, etc.\ all  % `\ ' makes an inter-word space.
       begin with G\@.         % \@ marks end-of-sentence punctuation.
You should check the spaces after periods when
reading your output to make sure you haven't
forgotten any special cases.  Generating an
ellipsis
       \ldots\               % `\ ' is needed after `\ldots' because TeX
                             % ignores spaces after command names like \ldots
                             % made from \ + letters.
                             %
                             % Note how a `%' character causes TeX to ignore
                             % the end of the input line, so these blank lines
                             % do not start a new paragraph.
                             %
with the right spacing around the periods requires
a special command.

\LaTeX\ interprets some common characters as
commands, so you must type special commands to
generate them.  These characters include the
following:
       \$ \& \% \# \{ and \}.

In printing, text is usually emphasized with an
       \emph{italic}
type style.

\begin{em}
   A long segment of text can also be emphasized
   in this way.  Text within such a segment can be
   given \emph{additional} emphasis.
\end{em}

It is sometimes necessary to prevent \LaTeX\ from
breaking a line where it might otherwise do so.
This may be at a space, as between the ``Mr.''\ and
``Jones'' in
       ``Mr.~Jones'',        % ~ produces an unbreakable interword space.
or within a word---especially when the word is a
symbol like
       \mbox{\emph{itemnum}}
that makes little sense when hyphenated across
lines.

Footnotes\footnote{This is an example of a footnote.}
pose no problem.

\LaTeX\ is good at typesetting mathematical formulas
like
       \( x-3y + z = 7 \)
or
       \( a_{1} > x^{2n} + y^{2n} > x' \)
or
       \( \ip{A}{B} = \sum_{i} a_{i} b_{i} \).
The spaces you type in a formula are
ignored.  Remember that a letter like
       $x$                   % $ ... $  and  \( ... \)  are equivalent
is a formula when it denotes a mathematical
symbol, and it should be typed as one.

\section{Displayed Text}

Text is displayed by indenting it from the left
margin.  Quotations are commonly displayed.  There
are short quotations
\begin{quote}
   This is a short quotation.  It consists of a
   single paragraph of text.  See how it is formatted.
\end{quote}
and longer ones.
\begin{quotation}
   This is a longer quotation.  It consists of two
   paragraphs of text, neither of which are
   particularly interesting.

   This is the second paragraph of the quotation.  It
   is just as dull as the first paragraph.
\end{quotation}
Another frequently-displayed structure is a list.
The following is an example of an \emph{itemized}
list.
\begin{itemize}
   \item This is the first item of an itemized list.
         Each item in the list is marked with a ``tick''.
         You don't have to worry about what kind of tick
         mark is used.

   \item This is the second item of the list.  It
         contains another list nested inside it.  The inner
         list is an \emph{enumerated} list.
         \begin{enumerate}
            \item This is the first item of an enumerated
                  list that is nested within the itemized list.

            \item This is the second item of the inner list.
                  \LaTeX\ allows you to nest lists deeper than
                  you really should.
         \end{enumerate}
         This is the rest of the second item of the outer
         list.  It is no more interesting than any other
         part of the item.
   \item This is the third item of the list.
\end{itemize}
You can even display poetry.
\begin{verse}
   There is an environment
    for verse \\             % The \\ command separates lines
   Whose features some poets % within a stanza.
   will curse.

                             % One or more blank lines separate stanzas.

   For instead of making\\
   Them do \emph{all} line breaking, \\
   It allows them to put too many words on a line when they'd rather be
   forced to be terse.
\end{verse}

Mathematical formulas may also be displayed.  A
displayed formula
is
one-line long; multiline
formulas require special formatting instructions.
   \[  \ip{\Gamma}{\psi'} = x'' + y^{2} + z_{i}^{n}\]
Don't start a paragraph with a displayed equation,
nor make one a paragraph by itself.

\end{document}               % End of document.

\section{Bench Huge Section 65}
% synthetic bench section 65
% This is a sample LaTeX input file.  (Version of 12 August 2004.)
%
% A '%' character causes TeX to ignore all remaining text on the line,
% and is used for comments like this one.

\documentclass{article}      % Specifies the document class

                             % The preamble begins here.
\title{An Example Document}  % Declares the document's title.
\author{Leslie Lamport}      % Declares the author's name.
\date{January 21, 1994}      % Deleting this command produces today's date.

\newcommand{\ip}[2]{(#1, #2)}
                             % Defines \ip{arg1}{arg2} to mean
                             % (arg1, arg2).

%\newcommand{\ip}[2]{\langle #1 | #2\rangle}
                             % This is an alternative definition of
                             % \ip that is commented out.

\begin{document}             % End of preamble and beginning of text.

\maketitle                   % Produces the title.

This is an example input file.  Comparing it with
the output it generates can show you how to
produce a simple document of your own.

\section{Ordinary Text}      % Produces section heading.  Lower-level
                             % sections are begun with similar
                             % \subsection and \subsubsection commands.

The ends  of words and sentences are marked
  by   spaces. It  doesn't matter how many
spaces    you type; one is as good as 100.  The
end of   a line counts as a space.

One   or more   blank lines denote the  end
of  a paragraph.

Since any number of consecutive spaces are treated
like a single one, the formatting of the input
file makes no difference to
      \LaTeX,                % The \LaTeX command generates the LaTeX logo.
but it makes a difference to you.  When you use
\LaTeX, making your input file as easy to read
as possible will be a great help as you write
your document and when you change it.  This sample
file shows how you can add comments to your own input
file.

Because printing is different from typewriting,
there are a number of things that you have to do
differently when preparing an input file than if
you were just typing the document directly.
Quotation marks like
       ``this''
have to be handled specially, as do quotes within
quotes:
       ``\,`this'            % \, separates the double and single quote.
        is what I just
        wrote, not  `that'\,''.

Dashes come in three sizes: an
       intra-word
dash, a medium dash for number ranges like
       1--2,
and a punctuation
       dash---like
this.

A sentence-ending space should be larger than the
space between words within a sentence.  You
sometimes have to type special commands in
conjunction with punctuation characters to get
this right, as in the following sentence.
       Gnats, gnus, etc.\ all  % `\ ' makes an inter-word space.
       begin with G\@.         % \@ marks end-of-sentence punctuation.
You should check the spaces after periods when
reading your output to make sure you haven't
forgotten any special cases.  Generating an
ellipsis
       \ldots\               % `\ ' is needed after `\ldots' because TeX
                             % ignores spaces after command names like \ldots
                             % made from \ + letters.
                             %
                             % Note how a `%' character causes TeX to ignore
                             % the end of the input line, so these blank lines
                             % do not start a new paragraph.
                             %
with the right spacing around the periods requires
a special command.

\LaTeX\ interprets some common characters as
commands, so you must type special commands to
generate them.  These characters include the
following:
       \$ \& \% \# \{ and \}.

In printing, text is usually emphasized with an
       \emph{italic}
type style.

\begin{em}
   A long segment of text can also be emphasized
   in this way.  Text within such a segment can be
   given \emph{additional} emphasis.
\end{em}

It is sometimes necessary to prevent \LaTeX\ from
breaking a line where it might otherwise do so.
This may be at a space, as between the ``Mr.''\ and
``Jones'' in
       ``Mr.~Jones'',        % ~ produces an unbreakable interword space.
or within a word---especially when the word is a
symbol like
       \mbox{\emph{itemnum}}
that makes little sense when hyphenated across
lines.

Footnotes\footnote{This is an example of a footnote.}
pose no problem.

\LaTeX\ is good at typesetting mathematical formulas
like
       \( x-3y + z = 7 \)
or
       \( a_{1} > x^{2n} + y^{2n} > x' \)
or
       \( \ip{A}{B} = \sum_{i} a_{i} b_{i} \).
The spaces you type in a formula are
ignored.  Remember that a letter like
       $x$                   % $ ... $  and  \( ... \)  are equivalent
is a formula when it denotes a mathematical
symbol, and it should be typed as one.

\section{Displayed Text}

Text is displayed by indenting it from the left
margin.  Quotations are commonly displayed.  There
are short quotations
\begin{quote}
   This is a short quotation.  It consists of a
   single paragraph of text.  See how it is formatted.
\end{quote}
and longer ones.
\begin{quotation}
   This is a longer quotation.  It consists of two
   paragraphs of text, neither of which are
   particularly interesting.

   This is the second paragraph of the quotation.  It
   is just as dull as the first paragraph.
\end{quotation}
Another frequently-displayed structure is a list.
The following is an example of an \emph{itemized}
list.
\begin{itemize}
   \item This is the first item of an itemized list.
         Each item in the list is marked with a ``tick''.
         You don't have to worry about what kind of tick
         mark is used.

   \item This is the second item of the list.  It
         contains another list nested inside it.  The inner
         list is an \emph{enumerated} list.
         \begin{enumerate}
            \item This is the first item of an enumerated
                  list that is nested within the itemized list.

            \item This is the second item of the inner list.
                  \LaTeX\ allows you to nest lists deeper than
                  you really should.
         \end{enumerate}
         This is the rest of the second item of the outer
         list.  It is no more interesting than any other
         part of the item.
   \item This is the third item of the list.
\end{itemize}
You can even display poetry.
\begin{verse}
   There is an environment
    for verse \\             % The \\ command separates lines
   Whose features some poets % within a stanza.
   will curse.

                             % One or more blank lines separate stanzas.

   For instead of making\\
   Them do \emph{all} line breaking, \\
   It allows them to put too many words on a line when they'd rather be
   forced to be terse.
\end{verse}

Mathematical formulas may also be displayed.  A
displayed formula
is
one-line long; multiline
formulas require special formatting instructions.
   \[  \ip{\Gamma}{\psi'} = x'' + y^{2} + z_{i}^{n}\]
Don't start a paragraph with a displayed equation,
nor make one a paragraph by itself.

\end{document}               % End of document.

\section{Bench Huge Section 66}
% synthetic bench section 66
% This is a sample LaTeX input file.  (Version of 12 August 2004.)
%
% A '%' character causes TeX to ignore all remaining text on the line,
% and is used for comments like this one.

\documentclass{article}      % Specifies the document class

                             % The preamble begins here.
\title{An Example Document}  % Declares the document's title.
\author{Leslie Lamport}      % Declares the author's name.
\date{January 21, 1994}      % Deleting this command produces today's date.

\newcommand{\ip}[2]{(#1, #2)}
                             % Defines \ip{arg1}{arg2} to mean
                             % (arg1, arg2).

%\newcommand{\ip}[2]{\langle #1 | #2\rangle}
                             % This is an alternative definition of
                             % \ip that is commented out.

\begin{document}             % End of preamble and beginning of text.

\maketitle                   % Produces the title.

This is an example input file.  Comparing it with
the output it generates can show you how to
produce a simple document of your own.

\section{Ordinary Text}      % Produces section heading.  Lower-level
                             % sections are begun with similar
                             % \subsection and \subsubsection commands.

The ends  of words and sentences are marked
  by   spaces. It  doesn't matter how many
spaces    you type; one is as good as 100.  The
end of   a line counts as a space.

One   or more   blank lines denote the  end
of  a paragraph.

Since any number of consecutive spaces are treated
like a single one, the formatting of the input
file makes no difference to
      \LaTeX,                % The \LaTeX command generates the LaTeX logo.
but it makes a difference to you.  When you use
\LaTeX, making your input file as easy to read
as possible will be a great help as you write
your document and when you change it.  This sample
file shows how you can add comments to your own input
file.

Because printing is different from typewriting,
there are a number of things that you have to do
differently when preparing an input file than if
you were just typing the document directly.
Quotation marks like
       ``this''
have to be handled specially, as do quotes within
quotes:
       ``\,`this'            % \, separates the double and single quote.
        is what I just
        wrote, not  `that'\,''.

Dashes come in three sizes: an
       intra-word
dash, a medium dash for number ranges like
       1--2,
and a punctuation
       dash---like
this.

A sentence-ending space should be larger than the
space between words within a sentence.  You
sometimes have to type special commands in
conjunction with punctuation characters to get
this right, as in the following sentence.
       Gnats, gnus, etc.\ all  % `\ ' makes an inter-word space.
       begin with G\@.         % \@ marks end-of-sentence punctuation.
You should check the spaces after periods when
reading your output to make sure you haven't
forgotten any special cases.  Generating an
ellipsis
       \ldots\               % `\ ' is needed after `\ldots' because TeX
                             % ignores spaces after command names like \ldots
                             % made from \ + letters.
                             %
                             % Note how a `%' character causes TeX to ignore
                             % the end of the input line, so these blank lines
                             % do not start a new paragraph.
                             %
with the right spacing around the periods requires
a special command.

\LaTeX\ interprets some common characters as
commands, so you must type special commands to
generate them.  These characters include the
following:
       \$ \& \% \# \{ and \}.

In printing, text is usually emphasized with an
       \emph{italic}
type style.

\begin{em}
   A long segment of text can also be emphasized
   in this way.  Text within such a segment can be
   given \emph{additional} emphasis.
\end{em}

It is sometimes necessary to prevent \LaTeX\ from
breaking a line where it might otherwise do so.
This may be at a space, as between the ``Mr.''\ and
``Jones'' in
       ``Mr.~Jones'',        % ~ produces an unbreakable interword space.
or within a word---especially when the word is a
symbol like
       \mbox{\emph{itemnum}}
that makes little sense when hyphenated across
lines.

Footnotes\footnote{This is an example of a footnote.}
pose no problem.

\LaTeX\ is good at typesetting mathematical formulas
like
       \( x-3y + z = 7 \)
or
       \( a_{1} > x^{2n} + y^{2n} > x' \)
or
       \( \ip{A}{B} = \sum_{i} a_{i} b_{i} \).
The spaces you type in a formula are
ignored.  Remember that a letter like
       $x$                   % $ ... $  and  \( ... \)  are equivalent
is a formula when it denotes a mathematical
symbol, and it should be typed as one.

\section{Displayed Text}

Text is displayed by indenting it from the left
margin.  Quotations are commonly displayed.  There
are short quotations
\begin{quote}
   This is a short quotation.  It consists of a
   single paragraph of text.  See how it is formatted.
\end{quote}
and longer ones.
\begin{quotation}
   This is a longer quotation.  It consists of two
   paragraphs of text, neither of which are
   particularly interesting.

   This is the second paragraph of the quotation.  It
   is just as dull as the first paragraph.
\end{quotation}
Another frequently-displayed structure is a list.
The following is an example of an \emph{itemized}
list.
\begin{itemize}
   \item This is the first item of an itemized list.
         Each item in the list is marked with a ``tick''.
         You don't have to worry about what kind of tick
         mark is used.

   \item This is the second item of the list.  It
         contains another list nested inside it.  The inner
         list is an \emph{enumerated} list.
         \begin{enumerate}
            \item This is the first item of an enumerated
                  list that is nested within the itemized list.

            \item This is the second item of the inner list.
                  \LaTeX\ allows you to nest lists deeper than
                  you really should.
         \end{enumerate}
         This is the rest of the second item of the outer
         list.  It is no more interesting than any other
         part of the item.
   \item This is the third item of the list.
\end{itemize}
You can even display poetry.
\begin{verse}
   There is an environment
    for verse \\             % The \\ command separates lines
   Whose features some poets % within a stanza.
   will curse.

                             % One or more blank lines separate stanzas.

   For instead of making\\
   Them do \emph{all} line breaking, \\
   It allows them to put too many words on a line when they'd rather be
   forced to be terse.
\end{verse}

Mathematical formulas may also be displayed.  A
displayed formula
is
one-line long; multiline
formulas require special formatting instructions.
   \[  \ip{\Gamma}{\psi'} = x'' + y^{2} + z_{i}^{n}\]
Don't start a paragraph with a displayed equation,
nor make one a paragraph by itself.

\end{document}               % End of document.

\section{Bench Huge Section 67}
% synthetic bench section 67
% This is a sample LaTeX input file.  (Version of 12 August 2004.)
%
% A '%' character causes TeX to ignore all remaining text on the line,
% and is used for comments like this one.

\documentclass{article}      % Specifies the document class

                             % The preamble begins here.
\title{An Example Document}  % Declares the document's title.
\author{Leslie Lamport}      % Declares the author's name.
\date{January 21, 1994}      % Deleting this command produces today's date.

\newcommand{\ip}[2]{(#1, #2)}
                             % Defines \ip{arg1}{arg2} to mean
                             % (arg1, arg2).

%\newcommand{\ip}[2]{\langle #1 | #2\rangle}
                             % This is an alternative definition of
                             % \ip that is commented out.

\begin{document}             % End of preamble and beginning of text.

\maketitle                   % Produces the title.

This is an example input file.  Comparing it with
the output it generates can show you how to
produce a simple document of your own.

\section{Ordinary Text}      % Produces section heading.  Lower-level
                             % sections are begun with similar
                             % \subsection and \subsubsection commands.

The ends  of words and sentences are marked
  by   spaces. It  doesn't matter how many
spaces    you type; one is as good as 100.  The
end of   a line counts as a space.

One   or more   blank lines denote the  end
of  a paragraph.

Since any number of consecutive spaces are treated
like a single one, the formatting of the input
file makes no difference to
      \LaTeX,                % The \LaTeX command generates the LaTeX logo.
but it makes a difference to you.  When you use
\LaTeX, making your input file as easy to read
as possible will be a great help as you write
your document and when you change it.  This sample
file shows how you can add comments to your own input
file.

Because printing is different from typewriting,
there are a number of things that you have to do
differently when preparing an input file than if
you were just typing the document directly.
Quotation marks like
       ``this''
have to be handled specially, as do quotes within
quotes:
       ``\,`this'            % \, separates the double and single quote.
        is what I just
        wrote, not  `that'\,''.

Dashes come in three sizes: an
       intra-word
dash, a medium dash for number ranges like
       1--2,
and a punctuation
       dash---like
this.

A sentence-ending space should be larger than the
space between words within a sentence.  You
sometimes have to type special commands in
conjunction with punctuation characters to get
this right, as in the following sentence.
       Gnats, gnus, etc.\ all  % `\ ' makes an inter-word space.
       begin with G\@.         % \@ marks end-of-sentence punctuation.
You should check the spaces after periods when
reading your output to make sure you haven't
forgotten any special cases.  Generating an
ellipsis
       \ldots\               % `\ ' is needed after `\ldots' because TeX
                             % ignores spaces after command names like \ldots
                             % made from \ + letters.
                             %
                             % Note how a `%' character causes TeX to ignore
                             % the end of the input line, so these blank lines
                             % do not start a new paragraph.
                             %
with the right spacing around the periods requires
a special command.

\LaTeX\ interprets some common characters as
commands, so you must type special commands to
generate them.  These characters include the
following:
       \$ \& \% \# \{ and \}.

In printing, text is usually emphasized with an
       \emph{italic}
type style.

\begin{em}
   A long segment of text can also be emphasized
   in this way.  Text within such a segment can be
   given \emph{additional} emphasis.
\end{em}

It is sometimes necessary to prevent \LaTeX\ from
breaking a line where it might otherwise do so.
This may be at a space, as between the ``Mr.''\ and
``Jones'' in
       ``Mr.~Jones'',        % ~ produces an unbreakable interword space.
or within a word---especially when the word is a
symbol like
       \mbox{\emph{itemnum}}
that makes little sense when hyphenated across
lines.

Footnotes\footnote{This is an example of a footnote.}
pose no problem.

\LaTeX\ is good at typesetting mathematical formulas
like
       \( x-3y + z = 7 \)
or
       \( a_{1} > x^{2n} + y^{2n} > x' \)
or
       \( \ip{A}{B} = \sum_{i} a_{i} b_{i} \).
The spaces you type in a formula are
ignored.  Remember that a letter like
       $x$                   % $ ... $  and  \( ... \)  are equivalent
is a formula when it denotes a mathematical
symbol, and it should be typed as one.

\section{Displayed Text}

Text is displayed by indenting it from the left
margin.  Quotations are commonly displayed.  There
are short quotations
\begin{quote}
   This is a short quotation.  It consists of a
   single paragraph of text.  See how it is formatted.
\end{quote}
and longer ones.
\begin{quotation}
   This is a longer quotation.  It consists of two
   paragraphs of text, neither of which are
   particularly interesting.

   This is the second paragraph of the quotation.  It
   is just as dull as the first paragraph.
\end{quotation}
Another frequently-displayed structure is a list.
The following is an example of an \emph{itemized}
list.
\begin{itemize}
   \item This is the first item of an itemized list.
         Each item in the list is marked with a ``tick''.
         You don't have to worry about what kind of tick
         mark is used.

   \item This is the second item of the list.  It
         contains another list nested inside it.  The inner
         list is an \emph{enumerated} list.
         \begin{enumerate}
            \item This is the first item of an enumerated
                  list that is nested within the itemized list.

            \item This is the second item of the inner list.
                  \LaTeX\ allows you to nest lists deeper than
                  you really should.
         \end{enumerate}
         This is the rest of the second item of the outer
         list.  It is no more interesting than any other
         part of the item.
   \item This is the third item of the list.
\end{itemize}
You can even display poetry.
\begin{verse}
   There is an environment
    for verse \\             % The \\ command separates lines
   Whose features some poets % within a stanza.
   will curse.

                             % One or more blank lines separate stanzas.

   For instead of making\\
   Them do \emph{all} line breaking, \\
   It allows them to put too many words on a line when they'd rather be
   forced to be terse.
\end{verse}

Mathematical formulas may also be displayed.  A
displayed formula
is
one-line long; multiline
formulas require special formatting instructions.
   \[  \ip{\Gamma}{\psi'} = x'' + y^{2} + z_{i}^{n}\]
Don't start a paragraph with a displayed equation,
nor make one a paragraph by itself.

\end{document}               % End of document.

\section{Bench Huge Section 68}
% synthetic bench section 68
% This is a sample LaTeX input file.  (Version of 12 August 2004.)
%
% A '%' character causes TeX to ignore all remaining text on the line,
% and is used for comments like this one.

\documentclass{article}      % Specifies the document class

                             % The preamble begins here.
\title{An Example Document}  % Declares the document's title.
\author{Leslie Lamport}      % Declares the author's name.
\date{January 21, 1994}      % Deleting this command produces today's date.

\newcommand{\ip}[2]{(#1, #2)}
                             % Defines \ip{arg1}{arg2} to mean
                             % (arg1, arg2).

%\newcommand{\ip}[2]{\langle #1 | #2\rangle}
                             % This is an alternative definition of
                             % \ip that is commented out.

\begin{document}             % End of preamble and beginning of text.

\maketitle                   % Produces the title.

This is an example input file.  Comparing it with
the output it generates can show you how to
produce a simple document of your own.

\section{Ordinary Text}      % Produces section heading.  Lower-level
                             % sections are begun with similar
                             % \subsection and \subsubsection commands.

The ends  of words and sentences are marked
  by   spaces. It  doesn't matter how many
spaces    you type; one is as good as 100.  The
end of   a line counts as a space.

One   or more   blank lines denote the  end
of  a paragraph.

Since any number of consecutive spaces are treated
like a single one, the formatting of the input
file makes no difference to
      \LaTeX,                % The \LaTeX command generates the LaTeX logo.
but it makes a difference to you.  When you use
\LaTeX, making your input file as easy to read
as possible will be a great help as you write
your document and when you change it.  This sample
file shows how you can add comments to your own input
file.

Because printing is different from typewriting,
there are a number of things that you have to do
differently when preparing an input file than if
you were just typing the document directly.
Quotation marks like
       ``this''
have to be handled specially, as do quotes within
quotes:
       ``\,`this'            % \, separates the double and single quote.
        is what I just
        wrote, not  `that'\,''.

Dashes come in three sizes: an
       intra-word
dash, a medium dash for number ranges like
       1--2,
and a punctuation
       dash---like
this.

A sentence-ending space should be larger than the
space between words within a sentence.  You
sometimes have to type special commands in
conjunction with punctuation characters to get
this right, as in the following sentence.
       Gnats, gnus, etc.\ all  % `\ ' makes an inter-word space.
       begin with G\@.         % \@ marks end-of-sentence punctuation.
You should check the spaces after periods when
reading your output to make sure you haven't
forgotten any special cases.  Generating an
ellipsis
       \ldots\               % `\ ' is needed after `\ldots' because TeX
                             % ignores spaces after command names like \ldots
                             % made from \ + letters.
                             %
                             % Note how a `%' character causes TeX to ignore
                             % the end of the input line, so these blank lines
                             % do not start a new paragraph.
                             %
with the right spacing around the periods requires
a special command.

\LaTeX\ interprets some common characters as
commands, so you must type special commands to
generate them.  These characters include the
following:
       \$ \& \% \# \{ and \}.

In printing, text is usually emphasized with an
       \emph{italic}
type style.

\begin{em}
   A long segment of text can also be emphasized
   in this way.  Text within such a segment can be
   given \emph{additional} emphasis.
\end{em}

It is sometimes necessary to prevent \LaTeX\ from
breaking a line where it might otherwise do so.
This may be at a space, as between the ``Mr.''\ and
``Jones'' in
       ``Mr.~Jones'',        % ~ produces an unbreakable interword space.
or within a word---especially when the word is a
symbol like
       \mbox{\emph{itemnum}}
that makes little sense when hyphenated across
lines.

Footnotes\footnote{This is an example of a footnote.}
pose no problem.

\LaTeX\ is good at typesetting mathematical formulas
like
       \( x-3y + z = 7 \)
or
       \( a_{1} > x^{2n} + y^{2n} > x' \)
or
       \( \ip{A}{B} = \sum_{i} a_{i} b_{i} \).
The spaces you type in a formula are
ignored.  Remember that a letter like
       $x$                   % $ ... $  and  \( ... \)  are equivalent
is a formula when it denotes a mathematical
symbol, and it should be typed as one.

\section{Displayed Text}

Text is displayed by indenting it from the left
margin.  Quotations are commonly displayed.  There
are short quotations
\begin{quote}
   This is a short quotation.  It consists of a
   single paragraph of text.  See how it is formatted.
\end{quote}
and longer ones.
\begin{quotation}
   This is a longer quotation.  It consists of two
   paragraphs of text, neither of which are
   particularly interesting.

   This is the second paragraph of the quotation.  It
   is just as dull as the first paragraph.
\end{quotation}
Another frequently-displayed structure is a list.
The following is an example of an \emph{itemized}
list.
\begin{itemize}
   \item This is the first item of an itemized list.
         Each item in the list is marked with a ``tick''.
         You don't have to worry about what kind of tick
         mark is used.

   \item This is the second item of the list.  It
         contains another list nested inside it.  The inner
         list is an \emph{enumerated} list.
         \begin{enumerate}
            \item This is the first item of an enumerated
                  list that is nested within the itemized list.

            \item This is the second item of the inner list.
                  \LaTeX\ allows you to nest lists deeper than
                  you really should.
         \end{enumerate}
         This is the rest of the second item of the outer
         list.  It is no more interesting than any other
         part of the item.
   \item This is the third item of the list.
\end{itemize}
You can even display poetry.
\begin{verse}
   There is an environment
    for verse \\             % The \\ command separates lines
   Whose features some poets % within a stanza.
   will curse.

                             % One or more blank lines separate stanzas.

   For instead of making\\
   Them do \emph{all} line breaking, \\
   It allows them to put too many words on a line when they'd rather be
   forced to be terse.
\end{verse}

Mathematical formulas may also be displayed.  A
displayed formula
is
one-line long; multiline
formulas require special formatting instructions.
   \[  \ip{\Gamma}{\psi'} = x'' + y^{2} + z_{i}^{n}\]
Don't start a paragraph with a displayed equation,
nor make one a paragraph by itself.

\end{document}               % End of document.

\section{Bench Huge Section 69}
% synthetic bench section 69
% This is a sample LaTeX input file.  (Version of 12 August 2004.)
%
% A '%' character causes TeX to ignore all remaining text on the line,
% and is used for comments like this one.

\documentclass{article}      % Specifies the document class

                             % The preamble begins here.
\title{An Example Document}  % Declares the document's title.
\author{Leslie Lamport}      % Declares the author's name.
\date{January 21, 1994}      % Deleting this command produces today's date.

\newcommand{\ip}[2]{(#1, #2)}
                             % Defines \ip{arg1}{arg2} to mean
                             % (arg1, arg2).

%\newcommand{\ip}[2]{\langle #1 | #2\rangle}
                             % This is an alternative definition of
                             % \ip that is commented out.

\begin{document}             % End of preamble and beginning of text.

\maketitle                   % Produces the title.

This is an example input file.  Comparing it with
the output it generates can show you how to
produce a simple document of your own.

\section{Ordinary Text}      % Produces section heading.  Lower-level
                             % sections are begun with similar
                             % \subsection and \subsubsection commands.

The ends  of words and sentences are marked
  by   spaces. It  doesn't matter how many
spaces    you type; one is as good as 100.  The
end of   a line counts as a space.

One   or more   blank lines denote the  end
of  a paragraph.

Since any number of consecutive spaces are treated
like a single one, the formatting of the input
file makes no difference to
      \LaTeX,                % The \LaTeX command generates the LaTeX logo.
but it makes a difference to you.  When you use
\LaTeX, making your input file as easy to read
as possible will be a great help as you write
your document and when you change it.  This sample
file shows how you can add comments to your own input
file.

Because printing is different from typewriting,
there are a number of things that you have to do
differently when preparing an input file than if
you were just typing the document directly.
Quotation marks like
       ``this''
have to be handled specially, as do quotes within
quotes:
       ``\,`this'            % \, separates the double and single quote.
        is what I just
        wrote, not  `that'\,''.

Dashes come in three sizes: an
       intra-word
dash, a medium dash for number ranges like
       1--2,
and a punctuation
       dash---like
this.

A sentence-ending space should be larger than the
space between words within a sentence.  You
sometimes have to type special commands in
conjunction with punctuation characters to get
this right, as in the following sentence.
       Gnats, gnus, etc.\ all  % `\ ' makes an inter-word space.
       begin with G\@.         % \@ marks end-of-sentence punctuation.
You should check the spaces after periods when
reading your output to make sure you haven't
forgotten any special cases.  Generating an
ellipsis
       \ldots\               % `\ ' is needed after `\ldots' because TeX
                             % ignores spaces after command names like \ldots
                             % made from \ + letters.
                             %
                             % Note how a `%' character causes TeX to ignore
                             % the end of the input line, so these blank lines
                             % do not start a new paragraph.
                             %
with the right spacing around the periods requires
a special command.

\LaTeX\ interprets some common characters as
commands, so you must type special commands to
generate them.  These characters include the
following:
       \$ \& \% \# \{ and \}.

In printing, text is usually emphasized with an
       \emph{italic}
type style.

\begin{em}
   A long segment of text can also be emphasized
   in this way.  Text within such a segment can be
   given \emph{additional} emphasis.
\end{em}

It is sometimes necessary to prevent \LaTeX\ from
breaking a line where it might otherwise do so.
This may be at a space, as between the ``Mr.''\ and
``Jones'' in
       ``Mr.~Jones'',        % ~ produces an unbreakable interword space.
or within a word---especially when the word is a
symbol like
       \mbox{\emph{itemnum}}
that makes little sense when hyphenated across
lines.

Footnotes\footnote{This is an example of a footnote.}
pose no problem.

\LaTeX\ is good at typesetting mathematical formulas
like
       \( x-3y + z = 7 \)
or
       \( a_{1} > x^{2n} + y^{2n} > x' \)
or
       \( \ip{A}{B} = \sum_{i} a_{i} b_{i} \).
The spaces you type in a formula are
ignored.  Remember that a letter like
       $x$                   % $ ... $  and  \( ... \)  are equivalent
is a formula when it denotes a mathematical
symbol, and it should be typed as one.

\section{Displayed Text}

Text is displayed by indenting it from the left
margin.  Quotations are commonly displayed.  There
are short quotations
\begin{quote}
   This is a short quotation.  It consists of a
   single paragraph of text.  See how it is formatted.
\end{quote}
and longer ones.
\begin{quotation}
   This is a longer quotation.  It consists of two
   paragraphs of text, neither of which are
   particularly interesting.

   This is the second paragraph of the quotation.  It
   is just as dull as the first paragraph.
\end{quotation}
Another frequently-displayed structure is a list.
The following is an example of an \emph{itemized}
list.
\begin{itemize}
   \item This is the first item of an itemized list.
         Each item in the list is marked with a ``tick''.
         You don't have to worry about what kind of tick
         mark is used.

   \item This is the second item of the list.  It
         contains another list nested inside it.  The inner
         list is an \emph{enumerated} list.
         \begin{enumerate}
            \item This is the first item of an enumerated
                  list that is nested within the itemized list.

            \item This is the second item of the inner list.
                  \LaTeX\ allows you to nest lists deeper than
                  you really should.
         \end{enumerate}
         This is the rest of the second item of the outer
         list.  It is no more interesting than any other
         part of the item.
   \item This is the third item of the list.
\end{itemize}
You can even display poetry.
\begin{verse}
   There is an environment
    for verse \\             % The \\ command separates lines
   Whose features some poets % within a stanza.
   will curse.

                             % One or more blank lines separate stanzas.

   For instead of making\\
   Them do \emph{all} line breaking, \\
   It allows them to put too many words on a line when they'd rather be
   forced to be terse.
\end{verse}

Mathematical formulas may also be displayed.  A
displayed formula
is
one-line long; multiline
formulas require special formatting instructions.
   \[  \ip{\Gamma}{\psi'} = x'' + y^{2} + z_{i}^{n}\]
Don't start a paragraph with a displayed equation,
nor make one a paragraph by itself.

\end{document}               % End of document.

\section{Bench Huge Section 70}
% synthetic bench section 70
% This is a sample LaTeX input file.  (Version of 12 August 2004.)
%
% A '%' character causes TeX to ignore all remaining text on the line,
% and is used for comments like this one.

\documentclass{article}      % Specifies the document class

                             % The preamble begins here.
\title{An Example Document}  % Declares the document's title.
\author{Leslie Lamport}      % Declares the author's name.
\date{January 21, 1994}      % Deleting this command produces today's date.

\newcommand{\ip}[2]{(#1, #2)}
                             % Defines \ip{arg1}{arg2} to mean
                             % (arg1, arg2).

%\newcommand{\ip}[2]{\langle #1 | #2\rangle}
                             % This is an alternative definition of
                             % \ip that is commented out.

\begin{document}             % End of preamble and beginning of text.

\maketitle                   % Produces the title.

This is an example input file.  Comparing it with
the output it generates can show you how to
produce a simple document of your own.

\section{Ordinary Text}      % Produces section heading.  Lower-level
                             % sections are begun with similar
                             % \subsection and \subsubsection commands.

The ends  of words and sentences are marked
  by   spaces. It  doesn't matter how many
spaces    you type; one is as good as 100.  The
end of   a line counts as a space.

One   or more   blank lines denote the  end
of  a paragraph.

Since any number of consecutive spaces are treated
like a single one, the formatting of the input
file makes no difference to
      \LaTeX,                % The \LaTeX command generates the LaTeX logo.
but it makes a difference to you.  When you use
\LaTeX, making your input file as easy to read
as possible will be a great help as you write
your document and when you change it.  This sample
file shows how you can add comments to your own input
file.

Because printing is different from typewriting,
there are a number of things that you have to do
differently when preparing an input file than if
you were just typing the document directly.
Quotation marks like
       ``this''
have to be handled specially, as do quotes within
quotes:
       ``\,`this'            % \, separates the double and single quote.
        is what I just
        wrote, not  `that'\,''.

Dashes come in three sizes: an
       intra-word
dash, a medium dash for number ranges like
       1--2,
and a punctuation
       dash---like
this.

A sentence-ending space should be larger than the
space between words within a sentence.  You
sometimes have to type special commands in
conjunction with punctuation characters to get
this right, as in the following sentence.
       Gnats, gnus, etc.\ all  % `\ ' makes an inter-word space.
       begin with G\@.         % \@ marks end-of-sentence punctuation.
You should check the spaces after periods when
reading your output to make sure you haven't
forgotten any special cases.  Generating an
ellipsis
       \ldots\               % `\ ' is needed after `\ldots' because TeX
                             % ignores spaces after command names like \ldots
                             % made from \ + letters.
                             %
                             % Note how a `%' character causes TeX to ignore
                             % the end of the input line, so these blank lines
                             % do not start a new paragraph.
                             %
with the right spacing around the periods requires
a special command.

\LaTeX\ interprets some common characters as
commands, so you must type special commands to
generate them.  These characters include the
following:
       \$ \& \% \# \{ and \}.

In printing, text is usually emphasized with an
       \emph{italic}
type style.

\begin{em}
   A long segment of text can also be emphasized
   in this way.  Text within such a segment can be
   given \emph{additional} emphasis.
\end{em}

It is sometimes necessary to prevent \LaTeX\ from
breaking a line where it might otherwise do so.
This may be at a space, as between the ``Mr.''\ and
``Jones'' in
       ``Mr.~Jones'',        % ~ produces an unbreakable interword space.
or within a word---especially when the word is a
symbol like
       \mbox{\emph{itemnum}}
that makes little sense when hyphenated across
lines.

Footnotes\footnote{This is an example of a footnote.}
pose no problem.

\LaTeX\ is good at typesetting mathematical formulas
like
       \( x-3y + z = 7 \)
or
       \( a_{1} > x^{2n} + y^{2n} > x' \)
or
       \( \ip{A}{B} = \sum_{i} a_{i} b_{i} \).
The spaces you type in a formula are
ignored.  Remember that a letter like
       $x$                   % $ ... $  and  \( ... \)  are equivalent
is a formula when it denotes a mathematical
symbol, and it should be typed as one.

\section{Displayed Text}

Text is displayed by indenting it from the left
margin.  Quotations are commonly displayed.  There
are short quotations
\begin{quote}
   This is a short quotation.  It consists of a
   single paragraph of text.  See how it is formatted.
\end{quote}
and longer ones.
\begin{quotation}
   This is a longer quotation.  It consists of two
   paragraphs of text, neither of which are
   particularly interesting.

   This is the second paragraph of the quotation.  It
   is just as dull as the first paragraph.
\end{quotation}
Another frequently-displayed structure is a list.
The following is an example of an \emph{itemized}
list.
\begin{itemize}
   \item This is the first item of an itemized list.
         Each item in the list is marked with a ``tick''.
         You don't have to worry about what kind of tick
         mark is used.

   \item This is the second item of the list.  It
         contains another list nested inside it.  The inner
         list is an \emph{enumerated} list.
         \begin{enumerate}
            \item This is the first item of an enumerated
                  list that is nested within the itemized list.

            \item This is the second item of the inner list.
                  \LaTeX\ allows you to nest lists deeper than
                  you really should.
         \end{enumerate}
         This is the rest of the second item of the outer
         list.  It is no more interesting than any other
         part of the item.
   \item This is the third item of the list.
\end{itemize}
You can even display poetry.
\begin{verse}
   There is an environment
    for verse \\             % The \\ command separates lines
   Whose features some poets % within a stanza.
   will curse.

                             % One or more blank lines separate stanzas.

   For instead of making\\
   Them do \emph{all} line breaking, \\
   It allows them to put too many words on a line when they'd rather be
   forced to be terse.
\end{verse}

Mathematical formulas may also be displayed.  A
displayed formula
is
one-line long; multiline
formulas require special formatting instructions.
   \[  \ip{\Gamma}{\psi'} = x'' + y^{2} + z_{i}^{n}\]
Don't start a paragraph with a displayed equation,
nor make one a paragraph by itself.

\end{document}               % End of document.

\section{Bench Huge Section 71}
% synthetic bench section 71
% This is a sample LaTeX input file.  (Version of 12 August 2004.)
%
% A '%' character causes TeX to ignore all remaining text on the line,
% and is used for comments like this one.

\documentclass{article}      % Specifies the document class

                             % The preamble begins here.
\title{An Example Document}  % Declares the document's title.
\author{Leslie Lamport}      % Declares the author's name.
\date{January 21, 1994}      % Deleting this command produces today's date.

\newcommand{\ip}[2]{(#1, #2)}
                             % Defines \ip{arg1}{arg2} to mean
                             % (arg1, arg2).

%\newcommand{\ip}[2]{\langle #1 | #2\rangle}
                             % This is an alternative definition of
                             % \ip that is commented out.

\begin{document}             % End of preamble and beginning of text.

\maketitle                   % Produces the title.

This is an example input file.  Comparing it with
the output it generates can show you how to
produce a simple document of your own.

\section{Ordinary Text}      % Produces section heading.  Lower-level
                             % sections are begun with similar
                             % \subsection and \subsubsection commands.

The ends  of words and sentences are marked
  by   spaces. It  doesn't matter how many
spaces    you type; one is as good as 100.  The
end of   a line counts as a space.

One   or more   blank lines denote the  end
of  a paragraph.

Since any number of consecutive spaces are treated
like a single one, the formatting of the input
file makes no difference to
      \LaTeX,                % The \LaTeX command generates the LaTeX logo.
but it makes a difference to you.  When you use
\LaTeX, making your input file as easy to read
as possible will be a great help as you write
your document and when you change it.  This sample
file shows how you can add comments to your own input
file.

Because printing is different from typewriting,
there are a number of things that you have to do
differently when preparing an input file than if
you were just typing the document directly.
Quotation marks like
       ``this''
have to be handled specially, as do quotes within
quotes:
       ``\,`this'            % \, separates the double and single quote.
        is what I just
        wrote, not  `that'\,''.

Dashes come in three sizes: an
       intra-word
dash, a medium dash for number ranges like
       1--2,
and a punctuation
       dash---like
this.

A sentence-ending space should be larger than the
space between words within a sentence.  You
sometimes have to type special commands in
conjunction with punctuation characters to get
this right, as in the following sentence.
       Gnats, gnus, etc.\ all  % `\ ' makes an inter-word space.
       begin with G\@.         % \@ marks end-of-sentence punctuation.
You should check the spaces after periods when
reading your output to make sure you haven't
forgotten any special cases.  Generating an
ellipsis
       \ldots\               % `\ ' is needed after `\ldots' because TeX
                             % ignores spaces after command names like \ldots
                             % made from \ + letters.
                             %
                             % Note how a `%' character causes TeX to ignore
                             % the end of the input line, so these blank lines
                             % do not start a new paragraph.
                             %
with the right spacing around the periods requires
a special command.

\LaTeX\ interprets some common characters as
commands, so you must type special commands to
generate them.  These characters include the
following:
       \$ \& \% \# \{ and \}.

In printing, text is usually emphasized with an
       \emph{italic}
type style.

\begin{em}
   A long segment of text can also be emphasized
   in this way.  Text within such a segment can be
   given \emph{additional} emphasis.
\end{em}

It is sometimes necessary to prevent \LaTeX\ from
breaking a line where it might otherwise do so.
This may be at a space, as between the ``Mr.''\ and
``Jones'' in
       ``Mr.~Jones'',        % ~ produces an unbreakable interword space.
or within a word---especially when the word is a
symbol like
       \mbox{\emph{itemnum}}
that makes little sense when hyphenated across
lines.

Footnotes\footnote{This is an example of a footnote.}
pose no problem.

\LaTeX\ is good at typesetting mathematical formulas
like
       \( x-3y + z = 7 \)
or
       \( a_{1} > x^{2n} + y^{2n} > x' \)
or
       \( \ip{A}{B} = \sum_{i} a_{i} b_{i} \).
The spaces you type in a formula are
ignored.  Remember that a letter like
       $x$                   % $ ... $  and  \( ... \)  are equivalent
is a formula when it denotes a mathematical
symbol, and it should be typed as one.

\section{Displayed Text}

Text is displayed by indenting it from the left
margin.  Quotations are commonly displayed.  There
are short quotations
\begin{quote}
   This is a short quotation.  It consists of a
   single paragraph of text.  See how it is formatted.
\end{quote}
and longer ones.
\begin{quotation}
   This is a longer quotation.  It consists of two
   paragraphs of text, neither of which are
   particularly interesting.

   This is the second paragraph of the quotation.  It
   is just as dull as the first paragraph.
\end{quotation}
Another frequently-displayed structure is a list.
The following is an example of an \emph{itemized}
list.
\begin{itemize}
   \item This is the first item of an itemized list.
         Each item in the list is marked with a ``tick''.
         You don't have to worry about what kind of tick
         mark is used.

   \item This is the second item of the list.  It
         contains another list nested inside it.  The inner
         list is an \emph{enumerated} list.
         \begin{enumerate}
            \item This is the first item of an enumerated
                  list that is nested within the itemized list.

            \item This is the second item of the inner list.
                  \LaTeX\ allows you to nest lists deeper than
                  you really should.
         \end{enumerate}
         This is the rest of the second item of the outer
         list.  It is no more interesting than any other
         part of the item.
   \item This is the third item of the list.
\end{itemize}
You can even display poetry.
\begin{verse}
   There is an environment
    for verse \\             % The \\ command separates lines
   Whose features some poets % within a stanza.
   will curse.

                             % One or more blank lines separate stanzas.

   For instead of making\\
   Them do \emph{all} line breaking, \\
   It allows them to put too many words on a line when they'd rather be
   forced to be terse.
\end{verse}

Mathematical formulas may also be displayed.  A
displayed formula
is
one-line long; multiline
formulas require special formatting instructions.
   \[  \ip{\Gamma}{\psi'} = x'' + y^{2} + z_{i}^{n}\]
Don't start a paragraph with a displayed equation,
nor make one a paragraph by itself.

\end{document}               % End of document.

\section{Bench Huge Section 72}
% synthetic bench section 72
% This is a sample LaTeX input file.  (Version of 12 August 2004.)
%
% A '%' character causes TeX to ignore all remaining text on the line,
% and is used for comments like this one.

\documentclass{article}      % Specifies the document class

                             % The preamble begins here.
\title{An Example Document}  % Declares the document's title.
\author{Leslie Lamport}      % Declares the author's name.
\date{January 21, 1994}      % Deleting this command produces today's date.

\newcommand{\ip}[2]{(#1, #2)}
                             % Defines \ip{arg1}{arg2} to mean
                             % (arg1, arg2).

%\newcommand{\ip}[2]{\langle #1 | #2\rangle}
                             % This is an alternative definition of
                             % \ip that is commented out.

\begin{document}             % End of preamble and beginning of text.

\maketitle                   % Produces the title.

This is an example input file.  Comparing it with
the output it generates can show you how to
produce a simple document of your own.

\section{Ordinary Text}      % Produces section heading.  Lower-level
                             % sections are begun with similar
                             % \subsection and \subsubsection commands.

The ends  of words and sentences are marked
  by   spaces. It  doesn't matter how many
spaces    you type; one is as good as 100.  The
end of   a line counts as a space.

One   or more   blank lines denote the  end
of  a paragraph.

Since any number of consecutive spaces are treated
like a single one, the formatting of the input
file makes no difference to
      \LaTeX,                % The \LaTeX command generates the LaTeX logo.
but it makes a difference to you.  When you use
\LaTeX, making your input file as easy to read
as possible will be a great help as you write
your document and when you change it.  This sample
file shows how you can add comments to your own input
file.

Because printing is different from typewriting,
there are a number of things that you have to do
differently when preparing an input file than if
you were just typing the document directly.
Quotation marks like
       ``this''
have to be handled specially, as do quotes within
quotes:
       ``\,`this'            % \, separates the double and single quote.
        is what I just
        wrote, not  `that'\,''.

Dashes come in three sizes: an
       intra-word
dash, a medium dash for number ranges like
       1--2,
and a punctuation
       dash---like
this.

A sentence-ending space should be larger than the
space between words within a sentence.  You
sometimes have to type special commands in
conjunction with punctuation characters to get
this right, as in the following sentence.
       Gnats, gnus, etc.\ all  % `\ ' makes an inter-word space.
       begin with G\@.         % \@ marks end-of-sentence punctuation.
You should check the spaces after periods when
reading your output to make sure you haven't
forgotten any special cases.  Generating an
ellipsis
       \ldots\               % `\ ' is needed after `\ldots' because TeX
                             % ignores spaces after command names like \ldots
                             % made from \ + letters.
                             %
                             % Note how a `%' character causes TeX to ignore
                             % the end of the input line, so these blank lines
                             % do not start a new paragraph.
                             %
with the right spacing around the periods requires
a special command.

\LaTeX\ interprets some common characters as
commands, so you must type special commands to
generate them.  These characters include the
following:
       \$ \& \% \# \{ and \}.

In printing, text is usually emphasized with an
       \emph{italic}
type style.

\begin{em}
   A long segment of text can also be emphasized
   in this way.  Text within such a segment can be
   given \emph{additional} emphasis.
\end{em}

It is sometimes necessary to prevent \LaTeX\ from
breaking a line where it might otherwise do so.
This may be at a space, as between the ``Mr.''\ and
``Jones'' in
       ``Mr.~Jones'',        % ~ produces an unbreakable interword space.
or within a word---especially when the word is a
symbol like
       \mbox{\emph{itemnum}}
that makes little sense when hyphenated across
lines.

Footnotes\footnote{This is an example of a footnote.}
pose no problem.

\LaTeX\ is good at typesetting mathematical formulas
like
       \( x-3y + z = 7 \)
or
       \( a_{1} > x^{2n} + y^{2n} > x' \)
or
       \( \ip{A}{B} = \sum_{i} a_{i} b_{i} \).
The spaces you type in a formula are
ignored.  Remember that a letter like
       $x$                   % $ ... $  and  \( ... \)  are equivalent
is a formula when it denotes a mathematical
symbol, and it should be typed as one.

\section{Displayed Text}

Text is displayed by indenting it from the left
margin.  Quotations are commonly displayed.  There
are short quotations
\begin{quote}
   This is a short quotation.  It consists of a
   single paragraph of text.  See how it is formatted.
\end{quote}
and longer ones.
\begin{quotation}
   This is a longer quotation.  It consists of two
   paragraphs of text, neither of which are
   particularly interesting.

   This is the second paragraph of the quotation.  It
   is just as dull as the first paragraph.
\end{quotation}
Another frequently-displayed structure is a list.
The following is an example of an \emph{itemized}
list.
\begin{itemize}
   \item This is the first item of an itemized list.
         Each item in the list is marked with a ``tick''.
         You don't have to worry about what kind of tick
         mark is used.

   \item This is the second item of the list.  It
         contains another list nested inside it.  The inner
         list is an \emph{enumerated} list.
         \begin{enumerate}
            \item This is the first item of an enumerated
                  list that is nested within the itemized list.

            \item This is the second item of the inner list.
                  \LaTeX\ allows you to nest lists deeper than
                  you really should.
         \end{enumerate}
         This is the rest of the second item of the outer
         list.  It is no more interesting than any other
         part of the item.
   \item This is the third item of the list.
\end{itemize}
You can even display poetry.
\begin{verse}
   There is an environment
    for verse \\             % The \\ command separates lines
   Whose features some poets % within a stanza.
   will curse.

                             % One or more blank lines separate stanzas.

   For instead of making\\
   Them do \emph{all} line breaking, \\
   It allows them to put too many words on a line when they'd rather be
   forced to be terse.
\end{verse}

Mathematical formulas may also be displayed.  A
displayed formula
is
one-line long; multiline
formulas require special formatting instructions.
   \[  \ip{\Gamma}{\psi'} = x'' + y^{2} + z_{i}^{n}\]
Don't start a paragraph with a displayed equation,
nor make one a paragraph by itself.

\end{document}               % End of document.

\section{Bench Huge Section 73}
% synthetic bench section 73
% This is a sample LaTeX input file.  (Version of 12 August 2004.)
%
% A '%' character causes TeX to ignore all remaining text on the line,
% and is used for comments like this one.

\documentclass{article}      % Specifies the document class

                             % The preamble begins here.
\title{An Example Document}  % Declares the document's title.
\author{Leslie Lamport}      % Declares the author's name.
\date{January 21, 1994}      % Deleting this command produces today's date.

\newcommand{\ip}[2]{(#1, #2)}
                             % Defines \ip{arg1}{arg2} to mean
                             % (arg1, arg2).

%\newcommand{\ip}[2]{\langle #1 | #2\rangle}
                             % This is an alternative definition of
                             % \ip that is commented out.

\begin{document}             % End of preamble and beginning of text.

\maketitle                   % Produces the title.

This is an example input file.  Comparing it with
the output it generates can show you how to
produce a simple document of your own.

\section{Ordinary Text}      % Produces section heading.  Lower-level
                             % sections are begun with similar
                             % \subsection and \subsubsection commands.

The ends  of words and sentences are marked
  by   spaces. It  doesn't matter how many
spaces    you type; one is as good as 100.  The
end of   a line counts as a space.

One   or more   blank lines denote the  end
of  a paragraph.

Since any number of consecutive spaces are treated
like a single one, the formatting of the input
file makes no difference to
      \LaTeX,                % The \LaTeX command generates the LaTeX logo.
but it makes a difference to you.  When you use
\LaTeX, making your input file as easy to read
as possible will be a great help as you write
your document and when you change it.  This sample
file shows how you can add comments to your own input
file.

Because printing is different from typewriting,
there are a number of things that you have to do
differently when preparing an input file than if
you were just typing the document directly.
Quotation marks like
       ``this''
have to be handled specially, as do quotes within
quotes:
       ``\,`this'            % \, separates the double and single quote.
        is what I just
        wrote, not  `that'\,''.

Dashes come in three sizes: an
       intra-word
dash, a medium dash for number ranges like
       1--2,
and a punctuation
       dash---like
this.

A sentence-ending space should be larger than the
space between words within a sentence.  You
sometimes have to type special commands in
conjunction with punctuation characters to get
this right, as in the following sentence.
       Gnats, gnus, etc.\ all  % `\ ' makes an inter-word space.
       begin with G\@.         % \@ marks end-of-sentence punctuation.
You should check the spaces after periods when
reading your output to make sure you haven't
forgotten any special cases.  Generating an
ellipsis
       \ldots\               % `\ ' is needed after `\ldots' because TeX
                             % ignores spaces after command names like \ldots
                             % made from \ + letters.
                             %
                             % Note how a `%' character causes TeX to ignore
                             % the end of the input line, so these blank lines
                             % do not start a new paragraph.
                             %
with the right spacing around the periods requires
a special command.

\LaTeX\ interprets some common characters as
commands, so you must type special commands to
generate them.  These characters include the
following:
       \$ \& \% \# \{ and \}.

In printing, text is usually emphasized with an
       \emph{italic}
type style.

\begin{em}
   A long segment of text can also be emphasized
   in this way.  Text within such a segment can be
   given \emph{additional} emphasis.
\end{em}

It is sometimes necessary to prevent \LaTeX\ from
breaking a line where it might otherwise do so.
This may be at a space, as between the ``Mr.''\ and
``Jones'' in
       ``Mr.~Jones'',        % ~ produces an unbreakable interword space.
or within a word---especially when the word is a
symbol like
       \mbox{\emph{itemnum}}
that makes little sense when hyphenated across
lines.

Footnotes\footnote{This is an example of a footnote.}
pose no problem.

\LaTeX\ is good at typesetting mathematical formulas
like
       \( x-3y + z = 7 \)
or
       \( a_{1} > x^{2n} + y^{2n} > x' \)
or
       \( \ip{A}{B} = \sum_{i} a_{i} b_{i} \).
The spaces you type in a formula are
ignored.  Remember that a letter like
       $x$                   % $ ... $  and  \( ... \)  are equivalent
is a formula when it denotes a mathematical
symbol, and it should be typed as one.

\section{Displayed Text}

Text is displayed by indenting it from the left
margin.  Quotations are commonly displayed.  There
are short quotations
\begin{quote}
   This is a short quotation.  It consists of a
   single paragraph of text.  See how it is formatted.
\end{quote}
and longer ones.
\begin{quotation}
   This is a longer quotation.  It consists of two
   paragraphs of text, neither of which are
   particularly interesting.

   This is the second paragraph of the quotation.  It
   is just as dull as the first paragraph.
\end{quotation}
Another frequently-displayed structure is a list.
The following is an example of an \emph{itemized}
list.
\begin{itemize}
   \item This is the first item of an itemized list.
         Each item in the list is marked with a ``tick''.
         You don't have to worry about what kind of tick
         mark is used.

   \item This is the second item of the list.  It
         contains another list nested inside it.  The inner
         list is an \emph{enumerated} list.
         \begin{enumerate}
            \item This is the first item of an enumerated
                  list that is nested within the itemized list.

            \item This is the second item of the inner list.
                  \LaTeX\ allows you to nest lists deeper than
                  you really should.
         \end{enumerate}
         This is the rest of the second item of the outer
         list.  It is no more interesting than any other
         part of the item.
   \item This is the third item of the list.
\end{itemize}
You can even display poetry.
\begin{verse}
   There is an environment
    for verse \\             % The \\ command separates lines
   Whose features some poets % within a stanza.
   will curse.

                             % One or more blank lines separate stanzas.

   For instead of making\\
   Them do \emph{all} line breaking, \\
   It allows them to put too many words on a line when they'd rather be
   forced to be terse.
\end{verse}

Mathematical formulas may also be displayed.  A
displayed formula
is
one-line long; multiline
formulas require special formatting instructions.
   \[  \ip{\Gamma}{\psi'} = x'' + y^{2} + z_{i}^{n}\]
Don't start a paragraph with a displayed equation,
nor make one a paragraph by itself.

\end{document}               % End of document.

\section{Bench Huge Section 74}
% synthetic bench section 74
% This is a sample LaTeX input file.  (Version of 12 August 2004.)
%
% A '%' character causes TeX to ignore all remaining text on the line,
% and is used for comments like this one.

\documentclass{article}      % Specifies the document class

                             % The preamble begins here.
\title{An Example Document}  % Declares the document's title.
\author{Leslie Lamport}      % Declares the author's name.
\date{January 21, 1994}      % Deleting this command produces today's date.

\newcommand{\ip}[2]{(#1, #2)}
                             % Defines \ip{arg1}{arg2} to mean
                             % (arg1, arg2).

%\newcommand{\ip}[2]{\langle #1 | #2\rangle}
                             % This is an alternative definition of
                             % \ip that is commented out.

\begin{document}             % End of preamble and beginning of text.

\maketitle                   % Produces the title.

This is an example input file.  Comparing it with
the output it generates can show you how to
produce a simple document of your own.

\section{Ordinary Text}      % Produces section heading.  Lower-level
                             % sections are begun with similar
                             % \subsection and \subsubsection commands.

The ends  of words and sentences are marked
  by   spaces. It  doesn't matter how many
spaces    you type; one is as good as 100.  The
end of   a line counts as a space.

One   or more   blank lines denote the  end
of  a paragraph.

Since any number of consecutive spaces are treated
like a single one, the formatting of the input
file makes no difference to
      \LaTeX,                % The \LaTeX command generates the LaTeX logo.
but it makes a difference to you.  When you use
\LaTeX, making your input file as easy to read
as possible will be a great help as you write
your document and when you change it.  This sample
file shows how you can add comments to your own input
file.

Because printing is different from typewriting,
there are a number of things that you have to do
differently when preparing an input file than if
you were just typing the document directly.
Quotation marks like
       ``this''
have to be handled specially, as do quotes within
quotes:
       ``\,`this'            % \, separates the double and single quote.
        is what I just
        wrote, not  `that'\,''.

Dashes come in three sizes: an
       intra-word
dash, a medium dash for number ranges like
       1--2,
and a punctuation
       dash---like
this.

A sentence-ending space should be larger than the
space between words within a sentence.  You
sometimes have to type special commands in
conjunction with punctuation characters to get
this right, as in the following sentence.
       Gnats, gnus, etc.\ all  % `\ ' makes an inter-word space.
       begin with G\@.         % \@ marks end-of-sentence punctuation.
You should check the spaces after periods when
reading your output to make sure you haven't
forgotten any special cases.  Generating an
ellipsis
       \ldots\               % `\ ' is needed after `\ldots' because TeX
                             % ignores spaces after command names like \ldots
                             % made from \ + letters.
                             %
                             % Note how a `%' character causes TeX to ignore
                             % the end of the input line, so these blank lines
                             % do not start a new paragraph.
                             %
with the right spacing around the periods requires
a special command.

\LaTeX\ interprets some common characters as
commands, so you must type special commands to
generate them.  These characters include the
following:
       \$ \& \% \# \{ and \}.

In printing, text is usually emphasized with an
       \emph{italic}
type style.

\begin{em}
   A long segment of text can also be emphasized
   in this way.  Text within such a segment can be
   given \emph{additional} emphasis.
\end{em}

It is sometimes necessary to prevent \LaTeX\ from
breaking a line where it might otherwise do so.
This may be at a space, as between the ``Mr.''\ and
``Jones'' in
       ``Mr.~Jones'',        % ~ produces an unbreakable interword space.
or within a word---especially when the word is a
symbol like
       \mbox{\emph{itemnum}}
that makes little sense when hyphenated across
lines.

Footnotes\footnote{This is an example of a footnote.}
pose no problem.

\LaTeX\ is good at typesetting mathematical formulas
like
       \( x-3y + z = 7 \)
or
       \( a_{1} > x^{2n} + y^{2n} > x' \)
or
       \( \ip{A}{B} = \sum_{i} a_{i} b_{i} \).
The spaces you type in a formula are
ignored.  Remember that a letter like
       $x$                   % $ ... $  and  \( ... \)  are equivalent
is a formula when it denotes a mathematical
symbol, and it should be typed as one.

\section{Displayed Text}

Text is displayed by indenting it from the left
margin.  Quotations are commonly displayed.  There
are short quotations
\begin{quote}
   This is a short quotation.  It consists of a
   single paragraph of text.  See how it is formatted.
\end{quote}
and longer ones.
\begin{quotation}
   This is a longer quotation.  It consists of two
   paragraphs of text, neither of which are
   particularly interesting.

   This is the second paragraph of the quotation.  It
   is just as dull as the first paragraph.
\end{quotation}
Another frequently-displayed structure is a list.
The following is an example of an \emph{itemized}
list.
\begin{itemize}
   \item This is the first item of an itemized list.
         Each item in the list is marked with a ``tick''.
         You don't have to worry about what kind of tick
         mark is used.

   \item This is the second item of the list.  It
         contains another list nested inside it.  The inner
         list is an \emph{enumerated} list.
         \begin{enumerate}
            \item This is the first item of an enumerated
                  list that is nested within the itemized list.

            \item This is the second item of the inner list.
                  \LaTeX\ allows you to nest lists deeper than
                  you really should.
         \end{enumerate}
         This is the rest of the second item of the outer
         list.  It is no more interesting than any other
         part of the item.
   \item This is the third item of the list.
\end{itemize}
You can even display poetry.
\begin{verse}
   There is an environment
    for verse \\             % The \\ command separates lines
   Whose features some poets % within a stanza.
   will curse.

                             % One or more blank lines separate stanzas.

   For instead of making\\
   Them do \emph{all} line breaking, \\
   It allows them to put too many words on a line when they'd rather be
   forced to be terse.
\end{verse}

Mathematical formulas may also be displayed.  A
displayed formula
is
one-line long; multiline
formulas require special formatting instructions.
   \[  \ip{\Gamma}{\psi'} = x'' + y^{2} + z_{i}^{n}\]
Don't start a paragraph with a displayed equation,
nor make one a paragraph by itself.

\end{document}               % End of document.

\section{Bench Huge Section 75}
% synthetic bench section 75
% This is a sample LaTeX input file.  (Version of 12 August 2004.)
%
% A '%' character causes TeX to ignore all remaining text on the line,
% and is used for comments like this one.

\documentclass{article}      % Specifies the document class

                             % The preamble begins here.
\title{An Example Document}  % Declares the document's title.
\author{Leslie Lamport}      % Declares the author's name.
\date{January 21, 1994}      % Deleting this command produces today's date.

\newcommand{\ip}[2]{(#1, #2)}
                             % Defines \ip{arg1}{arg2} to mean
                             % (arg1, arg2).

%\newcommand{\ip}[2]{\langle #1 | #2\rangle}
                             % This is an alternative definition of
                             % \ip that is commented out.

\begin{document}             % End of preamble and beginning of text.

\maketitle                   % Produces the title.

This is an example input file.  Comparing it with
the output it generates can show you how to
produce a simple document of your own.

\section{Ordinary Text}      % Produces section heading.  Lower-level
                             % sections are begun with similar
                             % \subsection and \subsubsection commands.

The ends  of words and sentences are marked
  by   spaces. It  doesn't matter how many
spaces    you type; one is as good as 100.  The
end of   a line counts as a space.

One   or more   blank lines denote the  end
of  a paragraph.

Since any number of consecutive spaces are treated
like a single one, the formatting of the input
file makes no difference to
      \LaTeX,                % The \LaTeX command generates the LaTeX logo.
but it makes a difference to you.  When you use
\LaTeX, making your input file as easy to read
as possible will be a great help as you write
your document and when you change it.  This sample
file shows how you can add comments to your own input
file.

Because printing is different from typewriting,
there are a number of things that you have to do
differently when preparing an input file than if
you were just typing the document directly.
Quotation marks like
       ``this''
have to be handled specially, as do quotes within
quotes:
       ``\,`this'            % \, separates the double and single quote.
        is what I just
        wrote, not  `that'\,''.

Dashes come in three sizes: an
       intra-word
dash, a medium dash for number ranges like
       1--2,
and a punctuation
       dash---like
this.

A sentence-ending space should be larger than the
space between words within a sentence.  You
sometimes have to type special commands in
conjunction with punctuation characters to get
this right, as in the following sentence.
       Gnats, gnus, etc.\ all  % `\ ' makes an inter-word space.
       begin with G\@.         % \@ marks end-of-sentence punctuation.
You should check the spaces after periods when
reading your output to make sure you haven't
forgotten any special cases.  Generating an
ellipsis
       \ldots\               % `\ ' is needed after `\ldots' because TeX
                             % ignores spaces after command names like \ldots
                             % made from \ + letters.
                             %
                             % Note how a `%' character causes TeX to ignore
                             % the end of the input line, so these blank lines
                             % do not start a new paragraph.
                             %
with the right spacing around the periods requires
a special command.

\LaTeX\ interprets some common characters as
commands, so you must type special commands to
generate them.  These characters include the
following:
       \$ \& \% \# \{ and \}.

In printing, text is usually emphasized with an
       \emph{italic}
type style.

\begin{em}
   A long segment of text can also be emphasized
   in this way.  Text within such a segment can be
   given \emph{additional} emphasis.
\end{em}

It is sometimes necessary to prevent \LaTeX\ from
breaking a line where it might otherwise do so.
This may be at a space, as between the ``Mr.''\ and
``Jones'' in
       ``Mr.~Jones'',        % ~ produces an unbreakable interword space.
or within a word---especially when the word is a
symbol like
       \mbox{\emph{itemnum}}
that makes little sense when hyphenated across
lines.

Footnotes\footnote{This is an example of a footnote.}
pose no problem.

\LaTeX\ is good at typesetting mathematical formulas
like
       \( x-3y + z = 7 \)
or
       \( a_{1} > x^{2n} + y^{2n} > x' \)
or
       \( \ip{A}{B} = \sum_{i} a_{i} b_{i} \).
The spaces you type in a formula are
ignored.  Remember that a letter like
       $x$                   % $ ... $  and  \( ... \)  are equivalent
is a formula when it denotes a mathematical
symbol, and it should be typed as one.

\section{Displayed Text}

Text is displayed by indenting it from the left
margin.  Quotations are commonly displayed.  There
are short quotations
\begin{quote}
   This is a short quotation.  It consists of a
   single paragraph of text.  See how it is formatted.
\end{quote}
and longer ones.
\begin{quotation}
   This is a longer quotation.  It consists of two
   paragraphs of text, neither of which are
   particularly interesting.

   This is the second paragraph of the quotation.  It
   is just as dull as the first paragraph.
\end{quotation}
Another frequently-displayed structure is a list.
The following is an example of an \emph{itemized}
list.
\begin{itemize}
   \item This is the first item of an itemized list.
         Each item in the list is marked with a ``tick''.
         You don't have to worry about what kind of tick
         mark is used.

   \item This is the second item of the list.  It
         contains another list nested inside it.  The inner
         list is an \emph{enumerated} list.
         \begin{enumerate}
            \item This is the first item of an enumerated
                  list that is nested within the itemized list.

            \item This is the second item of the inner list.
                  \LaTeX\ allows you to nest lists deeper than
                  you really should.
         \end{enumerate}
         This is the rest of the second item of the outer
         list.  It is no more interesting than any other
         part of the item.
   \item This is the third item of the list.
\end{itemize}
You can even display poetry.
\begin{verse}
   There is an environment
    for verse \\             % The \\ command separates lines
   Whose features some poets % within a stanza.
   will curse.

                             % One or more blank lines separate stanzas.

   For instead of making\\
   Them do \emph{all} line breaking, \\
   It allows them to put too many words on a line when they'd rather be
   forced to be terse.
\end{verse}

Mathematical formulas may also be displayed.  A
displayed formula
is
one-line long; multiline
formulas require special formatting instructions.
   \[  \ip{\Gamma}{\psi'} = x'' + y^{2} + z_{i}^{n}\]
Don't start a paragraph with a displayed equation,
nor make one a paragraph by itself.

\end{document}               % End of document.

\section{Bench Huge Section 76}
% synthetic bench section 76
% This is a sample LaTeX input file.  (Version of 12 August 2004.)
%
% A '%' character causes TeX to ignore all remaining text on the line,
% and is used for comments like this one.

\documentclass{article}      % Specifies the document class

                             % The preamble begins here.
\title{An Example Document}  % Declares the document's title.
\author{Leslie Lamport}      % Declares the author's name.
\date{January 21, 1994}      % Deleting this command produces today's date.

\newcommand{\ip}[2]{(#1, #2)}
                             % Defines \ip{arg1}{arg2} to mean
                             % (arg1, arg2).

%\newcommand{\ip}[2]{\langle #1 | #2\rangle}
                             % This is an alternative definition of
                             % \ip that is commented out.

\begin{document}             % End of preamble and beginning of text.

\maketitle                   % Produces the title.

This is an example input file.  Comparing it with
the output it generates can show you how to
produce a simple document of your own.

\section{Ordinary Text}      % Produces section heading.  Lower-level
                             % sections are begun with similar
                             % \subsection and \subsubsection commands.

The ends  of words and sentences are marked
  by   spaces. It  doesn't matter how many
spaces    you type; one is as good as 100.  The
end of   a line counts as a space.

One   or more   blank lines denote the  end
of  a paragraph.

Since any number of consecutive spaces are treated
like a single one, the formatting of the input
file makes no difference to
      \LaTeX,                % The \LaTeX command generates the LaTeX logo.
but it makes a difference to you.  When you use
\LaTeX, making your input file as easy to read
as possible will be a great help as you write
your document and when you change it.  This sample
file shows how you can add comments to your own input
file.

Because printing is different from typewriting,
there are a number of things that you have to do
differently when preparing an input file than if
you were just typing the document directly.
Quotation marks like
       ``this''
have to be handled specially, as do quotes within
quotes:
       ``\,`this'            % \, separates the double and single quote.
        is what I just
        wrote, not  `that'\,''.

Dashes come in three sizes: an
       intra-word
dash, a medium dash for number ranges like
       1--2,
and a punctuation
       dash---like
this.

A sentence-ending space should be larger than the
space between words within a sentence.  You
sometimes have to type special commands in
conjunction with punctuation characters to get
this right, as in the following sentence.
       Gnats, gnus, etc.\ all  % `\ ' makes an inter-word space.
       begin with G\@.         % \@ marks end-of-sentence punctuation.
You should check the spaces after periods when
reading your output to make sure you haven't
forgotten any special cases.  Generating an
ellipsis
       \ldots\               % `\ ' is needed after `\ldots' because TeX
                             % ignores spaces after command names like \ldots
                             % made from \ + letters.
                             %
                             % Note how a `%' character causes TeX to ignore
                             % the end of the input line, so these blank lines
                             % do not start a new paragraph.
                             %
with the right spacing around the periods requires
a special command.

\LaTeX\ interprets some common characters as
commands, so you must type special commands to
generate them.  These characters include the
following:
       \$ \& \% \# \{ and \}.

In printing, text is usually emphasized with an
       \emph{italic}
type style.

\begin{em}
   A long segment of text can also be emphasized
   in this way.  Text within such a segment can be
   given \emph{additional} emphasis.
\end{em}

It is sometimes necessary to prevent \LaTeX\ from
breaking a line where it might otherwise do so.
This may be at a space, as between the ``Mr.''\ and
``Jones'' in
       ``Mr.~Jones'',        % ~ produces an unbreakable interword space.
or within a word---especially when the word is a
symbol like
       \mbox{\emph{itemnum}}
that makes little sense when hyphenated across
lines.

Footnotes\footnote{This is an example of a footnote.}
pose no problem.

\LaTeX\ is good at typesetting mathematical formulas
like
       \( x-3y + z = 7 \)
or
       \( a_{1} > x^{2n} + y^{2n} > x' \)
or
       \( \ip{A}{B} = \sum_{i} a_{i} b_{i} \).
The spaces you type in a formula are
ignored.  Remember that a letter like
       $x$                   % $ ... $  and  \( ... \)  are equivalent
is a formula when it denotes a mathematical
symbol, and it should be typed as one.

\section{Displayed Text}

Text is displayed by indenting it from the left
margin.  Quotations are commonly displayed.  There
are short quotations
\begin{quote}
   This is a short quotation.  It consists of a
   single paragraph of text.  See how it is formatted.
\end{quote}
and longer ones.
\begin{quotation}
   This is a longer quotation.  It consists of two
   paragraphs of text, neither of which are
   particularly interesting.

   This is the second paragraph of the quotation.  It
   is just as dull as the first paragraph.
\end{quotation}
Another frequently-displayed structure is a list.
The following is an example of an \emph{itemized}
list.
\begin{itemize}
   \item This is the first item of an itemized list.
         Each item in the list is marked with a ``tick''.
         You don't have to worry about what kind of tick
         mark is used.

   \item This is the second item of the list.  It
         contains another list nested inside it.  The inner
         list is an \emph{enumerated} list.
         \begin{enumerate}
            \item This is the first item of an enumerated
                  list that is nested within the itemized list.

            \item This is the second item of the inner list.
                  \LaTeX\ allows you to nest lists deeper than
                  you really should.
         \end{enumerate}
         This is the rest of the second item of the outer
         list.  It is no more interesting than any other
         part of the item.
   \item This is the third item of the list.
\end{itemize}
You can even display poetry.
\begin{verse}
   There is an environment
    for verse \\             % The \\ command separates lines
   Whose features some poets % within a stanza.
   will curse.

                             % One or more blank lines separate stanzas.

   For instead of making\\
   Them do \emph{all} line breaking, \\
   It allows them to put too many words on a line when they'd rather be
   forced to be terse.
\end{verse}

Mathematical formulas may also be displayed.  A
displayed formula
is
one-line long; multiline
formulas require special formatting instructions.
   \[  \ip{\Gamma}{\psi'} = x'' + y^{2} + z_{i}^{n}\]
Don't start a paragraph with a displayed equation,
nor make one a paragraph by itself.

\end{document}               % End of document.

\section{Bench Huge Section 77}
% synthetic bench section 77
% This is a sample LaTeX input file.  (Version of 12 August 2004.)
%
% A '%' character causes TeX to ignore all remaining text on the line,
% and is used for comments like this one.

\documentclass{article}      % Specifies the document class

                             % The preamble begins here.
\title{An Example Document}  % Declares the document's title.
\author{Leslie Lamport}      % Declares the author's name.
\date{January 21, 1994}      % Deleting this command produces today's date.

\newcommand{\ip}[2]{(#1, #2)}
                             % Defines \ip{arg1}{arg2} to mean
                             % (arg1, arg2).

%\newcommand{\ip}[2]{\langle #1 | #2\rangle}
                             % This is an alternative definition of
                             % \ip that is commented out.

\begin{document}             % End of preamble and beginning of text.

\maketitle                   % Produces the title.

This is an example input file.  Comparing it with
the output it generates can show you how to
produce a simple document of your own.

\section{Ordinary Text}      % Produces section heading.  Lower-level
                             % sections are begun with similar
                             % \subsection and \subsubsection commands.

The ends  of words and sentences are marked
  by   spaces. It  doesn't matter how many
spaces    you type; one is as good as 100.  The
end of   a line counts as a space.

One   or more   blank lines denote the  end
of  a paragraph.

Since any number of consecutive spaces are treated
like a single one, the formatting of the input
file makes no difference to
      \LaTeX,                % The \LaTeX command generates the LaTeX logo.
but it makes a difference to you.  When you use
\LaTeX, making your input file as easy to read
as possible will be a great help as you write
your document and when you change it.  This sample
file shows how you can add comments to your own input
file.

Because printing is different from typewriting,
there are a number of things that you have to do
differently when preparing an input file than if
you were just typing the document directly.
Quotation marks like
       ``this''
have to be handled specially, as do quotes within
quotes:
       ``\,`this'            % \, separates the double and single quote.
        is what I just
        wrote, not  `that'\,''.

Dashes come in three sizes: an
       intra-word
dash, a medium dash for number ranges like
       1--2,
and a punctuation
       dash---like
this.

A sentence-ending space should be larger than the
space between words within a sentence.  You
sometimes have to type special commands in
conjunction with punctuation characters to get
this right, as in the following sentence.
       Gnats, gnus, etc.\ all  % `\ ' makes an inter-word space.
       begin with G\@.         % \@ marks end-of-sentence punctuation.
You should check the spaces after periods when
reading your output to make sure you haven't
forgotten any special cases.  Generating an
ellipsis
       \ldots\               % `\ ' is needed after `\ldots' because TeX
                             % ignores spaces after command names like \ldots
                             % made from \ + letters.
                             %
                             % Note how a `%' character causes TeX to ignore
                             % the end of the input line, so these blank lines
                             % do not start a new paragraph.
                             %
with the right spacing around the periods requires
a special command.

\LaTeX\ interprets some common characters as
commands, so you must type special commands to
generate them.  These characters include the
following:
       \$ \& \% \# \{ and \}.

In printing, text is usually emphasized with an
       \emph{italic}
type style.

\begin{em}
   A long segment of text can also be emphasized
   in this way.  Text within such a segment can be
   given \emph{additional} emphasis.
\end{em}

It is sometimes necessary to prevent \LaTeX\ from
breaking a line where it might otherwise do so.
This may be at a space, as between the ``Mr.''\ and
``Jones'' in
       ``Mr.~Jones'',        % ~ produces an unbreakable interword space.
or within a word---especially when the word is a
symbol like
       \mbox{\emph{itemnum}}
that makes little sense when hyphenated across
lines.

Footnotes\footnote{This is an example of a footnote.}
pose no problem.

\LaTeX\ is good at typesetting mathematical formulas
like
       \( x-3y + z = 7 \)
or
       \( a_{1} > x^{2n} + y^{2n} > x' \)
or
       \( \ip{A}{B} = \sum_{i} a_{i} b_{i} \).
The spaces you type in a formula are
ignored.  Remember that a letter like
       $x$                   % $ ... $  and  \( ... \)  are equivalent
is a formula when it denotes a mathematical
symbol, and it should be typed as one.

\section{Displayed Text}

Text is displayed by indenting it from the left
margin.  Quotations are commonly displayed.  There
are short quotations
\begin{quote}
   This is a short quotation.  It consists of a
   single paragraph of text.  See how it is formatted.
\end{quote}
and longer ones.
\begin{quotation}
   This is a longer quotation.  It consists of two
   paragraphs of text, neither of which are
   particularly interesting.

   This is the second paragraph of the quotation.  It
   is just as dull as the first paragraph.
\end{quotation}
Another frequently-displayed structure is a list.
The following is an example of an \emph{itemized}
list.
\begin{itemize}
   \item This is the first item of an itemized list.
         Each item in the list is marked with a ``tick''.
         You don't have to worry about what kind of tick
         mark is used.

   \item This is the second item of the list.  It
         contains another list nested inside it.  The inner
         list is an \emph{enumerated} list.
         \begin{enumerate}
            \item This is the first item of an enumerated
                  list that is nested within the itemized list.

            \item This is the second item of the inner list.
                  \LaTeX\ allows you to nest lists deeper than
                  you really should.
         \end{enumerate}
         This is the rest of the second item of the outer
         list.  It is no more interesting than any other
         part of the item.
   \item This is the third item of the list.
\end{itemize}
You can even display poetry.
\begin{verse}
   There is an environment
    for verse \\             % The \\ command separates lines
   Whose features some poets % within a stanza.
   will curse.

                             % One or more blank lines separate stanzas.

   For instead of making\\
   Them do \emph{all} line breaking, \\
   It allows them to put too many words on a line when they'd rather be
   forced to be terse.
\end{verse}

Mathematical formulas may also be displayed.  A
displayed formula
is
one-line long; multiline
formulas require special formatting instructions.
   \[  \ip{\Gamma}{\psi'} = x'' + y^{2} + z_{i}^{n}\]
Don't start a paragraph with a displayed equation,
nor make one a paragraph by itself.

\end{document}               % End of document.

\section{Bench Huge Section 78}
% synthetic bench section 78
% This is a sample LaTeX input file.  (Version of 12 August 2004.)
%
% A '%' character causes TeX to ignore all remaining text on the line,
% and is used for comments like this one.

\documentclass{article}      % Specifies the document class

                             % The preamble begins here.
\title{An Example Document}  % Declares the document's title.
\author{Leslie Lamport}      % Declares the author's name.
\date{January 21, 1994}      % Deleting this command produces today's date.

\newcommand{\ip}[2]{(#1, #2)}
                             % Defines \ip{arg1}{arg2} to mean
                             % (arg1, arg2).

%\newcommand{\ip}[2]{\langle #1 | #2\rangle}
                             % This is an alternative definition of
                             % \ip that is commented out.

\begin{document}             % End of preamble and beginning of text.

\maketitle                   % Produces the title.

This is an example input file.  Comparing it with
the output it generates can show you how to
produce a simple document of your own.

\section{Ordinary Text}      % Produces section heading.  Lower-level
                             % sections are begun with similar
                             % \subsection and \subsubsection commands.

The ends  of words and sentences are marked
  by   spaces. It  doesn't matter how many
spaces    you type; one is as good as 100.  The
end of   a line counts as a space.

One   or more   blank lines denote the  end
of  a paragraph.

Since any number of consecutive spaces are treated
like a single one, the formatting of the input
file makes no difference to
      \LaTeX,                % The \LaTeX command generates the LaTeX logo.
but it makes a difference to you.  When you use
\LaTeX, making your input file as easy to read
as possible will be a great help as you write
your document and when you change it.  This sample
file shows how you can add comments to your own input
file.

Because printing is different from typewriting,
there are a number of things that you have to do
differently when preparing an input file than if
you were just typing the document directly.
Quotation marks like
       ``this''
have to be handled specially, as do quotes within
quotes:
       ``\,`this'            % \, separates the double and single quote.
        is what I just
        wrote, not  `that'\,''.

Dashes come in three sizes: an
       intra-word
dash, a medium dash for number ranges like
       1--2,
and a punctuation
       dash---like
this.

A sentence-ending space should be larger than the
space between words within a sentence.  You
sometimes have to type special commands in
conjunction with punctuation characters to get
this right, as in the following sentence.
       Gnats, gnus, etc.\ all  % `\ ' makes an inter-word space.
       begin with G\@.         % \@ marks end-of-sentence punctuation.
You should check the spaces after periods when
reading your output to make sure you haven't
forgotten any special cases.  Generating an
ellipsis
       \ldots\               % `\ ' is needed after `\ldots' because TeX
                             % ignores spaces after command names like \ldots
                             % made from \ + letters.
                             %
                             % Note how a `%' character causes TeX to ignore
                             % the end of the input line, so these blank lines
                             % do not start a new paragraph.
                             %
with the right spacing around the periods requires
a special command.

\LaTeX\ interprets some common characters as
commands, so you must type special commands to
generate them.  These characters include the
following:
       \$ \& \% \# \{ and \}.

In printing, text is usually emphasized with an
       \emph{italic}
type style.

\begin{em}
   A long segment of text can also be emphasized
   in this way.  Text within such a segment can be
   given \emph{additional} emphasis.
\end{em}

It is sometimes necessary to prevent \LaTeX\ from
breaking a line where it might otherwise do so.
This may be at a space, as between the ``Mr.''\ and
``Jones'' in
       ``Mr.~Jones'',        % ~ produces an unbreakable interword space.
or within a word---especially when the word is a
symbol like
       \mbox{\emph{itemnum}}
that makes little sense when hyphenated across
lines.

Footnotes\footnote{This is an example of a footnote.}
pose no problem.

\LaTeX\ is good at typesetting mathematical formulas
like
       \( x-3y + z = 7 \)
or
       \( a_{1} > x^{2n} + y^{2n} > x' \)
or
       \( \ip{A}{B} = \sum_{i} a_{i} b_{i} \).
The spaces you type in a formula are
ignored.  Remember that a letter like
       $x$                   % $ ... $  and  \( ... \)  are equivalent
is a formula when it denotes a mathematical
symbol, and it should be typed as one.

\section{Displayed Text}

Text is displayed by indenting it from the left
margin.  Quotations are commonly displayed.  There
are short quotations
\begin{quote}
   This is a short quotation.  It consists of a
   single paragraph of text.  See how it is formatted.
\end{quote}
and longer ones.
\begin{quotation}
   This is a longer quotation.  It consists of two
   paragraphs of text, neither of which are
   particularly interesting.

   This is the second paragraph of the quotation.  It
   is just as dull as the first paragraph.
\end{quotation}
Another frequently-displayed structure is a list.
The following is an example of an \emph{itemized}
list.
\begin{itemize}
   \item This is the first item of an itemized list.
         Each item in the list is marked with a ``tick''.
         You don't have to worry about what kind of tick
         mark is used.

   \item This is the second item of the list.  It
         contains another list nested inside it.  The inner
         list is an \emph{enumerated} list.
         \begin{enumerate}
            \item This is the first item of an enumerated
                  list that is nested within the itemized list.

            \item This is the second item of the inner list.
                  \LaTeX\ allows you to nest lists deeper than
                  you really should.
         \end{enumerate}
         This is the rest of the second item of the outer
         list.  It is no more interesting than any other
         part of the item.
   \item This is the third item of the list.
\end{itemize}
You can even display poetry.
\begin{verse}
   There is an environment
    for verse \\             % The \\ command separates lines
   Whose features some poets % within a stanza.
   will curse.

                             % One or more blank lines separate stanzas.

   For instead of making\\
   Them do \emph{all} line breaking, \\
   It allows them to put too many words on a line when they'd rather be
   forced to be terse.
\end{verse}

Mathematical formulas may also be displayed.  A
displayed formula
is
one-line long; multiline
formulas require special formatting instructions.
   \[  \ip{\Gamma}{\psi'} = x'' + y^{2} + z_{i}^{n}\]
Don't start a paragraph with a displayed equation,
nor make one a paragraph by itself.

\end{document}               % End of document.

\section{Bench Huge Section 79}
% synthetic bench section 79
% This is a sample LaTeX input file.  (Version of 12 August 2004.)
%
% A '%' character causes TeX to ignore all remaining text on the line,
% and is used for comments like this one.

\documentclass{article}      % Specifies the document class

                             % The preamble begins here.
\title{An Example Document}  % Declares the document's title.
\author{Leslie Lamport}      % Declares the author's name.
\date{January 21, 1994}      % Deleting this command produces today's date.

\newcommand{\ip}[2]{(#1, #2)}
                             % Defines \ip{arg1}{arg2} to mean
                             % (arg1, arg2).

%\newcommand{\ip}[2]{\langle #1 | #2\rangle}
                             % This is an alternative definition of
                             % \ip that is commented out.

\begin{document}             % End of preamble and beginning of text.

\maketitle                   % Produces the title.

This is an example input file.  Comparing it with
the output it generates can show you how to
produce a simple document of your own.

\section{Ordinary Text}      % Produces section heading.  Lower-level
                             % sections are begun with similar
                             % \subsection and \subsubsection commands.

The ends  of words and sentences are marked
  by   spaces. It  doesn't matter how many
spaces    you type; one is as good as 100.  The
end of   a line counts as a space.

One   or more   blank lines denote the  end
of  a paragraph.

Since any number of consecutive spaces are treated
like a single one, the formatting of the input
file makes no difference to
      \LaTeX,                % The \LaTeX command generates the LaTeX logo.
but it makes a difference to you.  When you use
\LaTeX, making your input file as easy to read
as possible will be a great help as you write
your document and when you change it.  This sample
file shows how you can add comments to your own input
file.

Because printing is different from typewriting,
there are a number of things that you have to do
differently when preparing an input file than if
you were just typing the document directly.
Quotation marks like
       ``this''
have to be handled specially, as do quotes within
quotes:
       ``\,`this'            % \, separates the double and single quote.
        is what I just
        wrote, not  `that'\,''.

Dashes come in three sizes: an
       intra-word
dash, a medium dash for number ranges like
       1--2,
and a punctuation
       dash---like
this.

A sentence-ending space should be larger than the
space between words within a sentence.  You
sometimes have to type special commands in
conjunction with punctuation characters to get
this right, as in the following sentence.
       Gnats, gnus, etc.\ all  % `\ ' makes an inter-word space.
       begin with G\@.         % \@ marks end-of-sentence punctuation.
You should check the spaces after periods when
reading your output to make sure you haven't
forgotten any special cases.  Generating an
ellipsis
       \ldots\               % `\ ' is needed after `\ldots' because TeX
                             % ignores spaces after command names like \ldots
                             % made from \ + letters.
                             %
                             % Note how a `%' character causes TeX to ignore
                             % the end of the input line, so these blank lines
                             % do not start a new paragraph.
                             %
with the right spacing around the periods requires
a special command.

\LaTeX\ interprets some common characters as
commands, so you must type special commands to
generate them.  These characters include the
following:
       \$ \& \% \# \{ and \}.

In printing, text is usually emphasized with an
       \emph{italic}
type style.

\begin{em}
   A long segment of text can also be emphasized
   in this way.  Text within such a segment can be
   given \emph{additional} emphasis.
\end{em}

It is sometimes necessary to prevent \LaTeX\ from
breaking a line where it might otherwise do so.
This may be at a space, as between the ``Mr.''\ and
``Jones'' in
       ``Mr.~Jones'',        % ~ produces an unbreakable interword space.
or within a word---especially when the word is a
symbol like
       \mbox{\emph{itemnum}}
that makes little sense when hyphenated across
lines.

Footnotes\footnote{This is an example of a footnote.}
pose no problem.

\LaTeX\ is good at typesetting mathematical formulas
like
       \( x-3y + z = 7 \)
or
       \( a_{1} > x^{2n} + y^{2n} > x' \)
or
       \( \ip{A}{B} = \sum_{i} a_{i} b_{i} \).
The spaces you type in a formula are
ignored.  Remember that a letter like
       $x$                   % $ ... $  and  \( ... \)  are equivalent
is a formula when it denotes a mathematical
symbol, and it should be typed as one.

\section{Displayed Text}

Text is displayed by indenting it from the left
margin.  Quotations are commonly displayed.  There
are short quotations
\begin{quote}
   This is a short quotation.  It consists of a
   single paragraph of text.  See how it is formatted.
\end{quote}
and longer ones.
\begin{quotation}
   This is a longer quotation.  It consists of two
   paragraphs of text, neither of which are
   particularly interesting.

   This is the second paragraph of the quotation.  It
   is just as dull as the first paragraph.
\end{quotation}
Another frequently-displayed structure is a list.
The following is an example of an \emph{itemized}
list.
\begin{itemize}
   \item This is the first item of an itemized list.
         Each item in the list is marked with a ``tick''.
         You don't have to worry about what kind of tick
         mark is used.

   \item This is the second item of the list.  It
         contains another list nested inside it.  The inner
         list is an \emph{enumerated} list.
         \begin{enumerate}
            \item This is the first item of an enumerated
                  list that is nested within the itemized list.

            \item This is the second item of the inner list.
                  \LaTeX\ allows you to nest lists deeper than
                  you really should.
         \end{enumerate}
         This is the rest of the second item of the outer
         list.  It is no more interesting than any other
         part of the item.
   \item This is the third item of the list.
\end{itemize}
You can even display poetry.
\begin{verse}
   There is an environment
    for verse \\             % The \\ command separates lines
   Whose features some poets % within a stanza.
   will curse.

                             % One or more blank lines separate stanzas.

   For instead of making\\
   Them do \emph{all} line breaking, \\
   It allows them to put too many words on a line when they'd rather be
   forced to be terse.
\end{verse}

Mathematical formulas may also be displayed.  A
displayed formula
is
one-line long; multiline
formulas require special formatting instructions.
   \[  \ip{\Gamma}{\psi'} = x'' + y^{2} + z_{i}^{n}\]
Don't start a paragraph with a displayed equation,
nor make one a paragraph by itself.

\end{document}               % End of document.

\section{Bench Huge Section 80}
% synthetic bench section 80
% This is a sample LaTeX input file.  (Version of 12 August 2004.)
%
% A '%' character causes TeX to ignore all remaining text on the line,
% and is used for comments like this one.

\documentclass{article}      % Specifies the document class

                             % The preamble begins here.
\title{An Example Document}  % Declares the document's title.
\author{Leslie Lamport}      % Declares the author's name.
\date{January 21, 1994}      % Deleting this command produces today's date.

\newcommand{\ip}[2]{(#1, #2)}
                             % Defines \ip{arg1}{arg2} to mean
                             % (arg1, arg2).

%\newcommand{\ip}[2]{\langle #1 | #2\rangle}
                             % This is an alternative definition of
                             % \ip that is commented out.

\begin{document}             % End of preamble and beginning of text.

\maketitle                   % Produces the title.

This is an example input file.  Comparing it with
the output it generates can show you how to
produce a simple document of your own.

\section{Ordinary Text}      % Produces section heading.  Lower-level
                             % sections are begun with similar
                             % \subsection and \subsubsection commands.

The ends  of words and sentences are marked
  by   spaces. It  doesn't matter how many
spaces    you type; one is as good as 100.  The
end of   a line counts as a space.

One   or more   blank lines denote the  end
of  a paragraph.

Since any number of consecutive spaces are treated
like a single one, the formatting of the input
file makes no difference to
      \LaTeX,                % The \LaTeX command generates the LaTeX logo.
but it makes a difference to you.  When you use
\LaTeX, making your input file as easy to read
as possible will be a great help as you write
your document and when you change it.  This sample
file shows how you can add comments to your own input
file.

Because printing is different from typewriting,
there are a number of things that you have to do
differently when preparing an input file than if
you were just typing the document directly.
Quotation marks like
       ``this''
have to be handled specially, as do quotes within
quotes:
       ``\,`this'            % \, separates the double and single quote.
        is what I just
        wrote, not  `that'\,''.

Dashes come in three sizes: an
       intra-word
dash, a medium dash for number ranges like
       1--2,
and a punctuation
       dash---like
this.

A sentence-ending space should be larger than the
space between words within a sentence.  You
sometimes have to type special commands in
conjunction with punctuation characters to get
this right, as in the following sentence.
       Gnats, gnus, etc.\ all  % `\ ' makes an inter-word space.
       begin with G\@.         % \@ marks end-of-sentence punctuation.
You should check the spaces after periods when
reading your output to make sure you haven't
forgotten any special cases.  Generating an
ellipsis
       \ldots\               % `\ ' is needed after `\ldots' because TeX
                             % ignores spaces after command names like \ldots
                             % made from \ + letters.
                             %
                             % Note how a `%' character causes TeX to ignore
                             % the end of the input line, so these blank lines
                             % do not start a new paragraph.
                             %
with the right spacing around the periods requires
a special command.

\LaTeX\ interprets some common characters as
commands, so you must type special commands to
generate them.  These characters include the
following:
       \$ \& \% \# \{ and \}.

In printing, text is usually emphasized with an
       \emph{italic}
type style.

\begin{em}
   A long segment of text can also be emphasized
   in this way.  Text within such a segment can be
   given \emph{additional} emphasis.
\end{em}

It is sometimes necessary to prevent \LaTeX\ from
breaking a line where it might otherwise do so.
This may be at a space, as between the ``Mr.''\ and
``Jones'' in
       ``Mr.~Jones'',        % ~ produces an unbreakable interword space.
or within a word---especially when the word is a
symbol like
       \mbox{\emph{itemnum}}
that makes little sense when hyphenated across
lines.

Footnotes\footnote{This is an example of a footnote.}
pose no problem.

\LaTeX\ is good at typesetting mathematical formulas
like
       \( x-3y + z = 7 \)
or
       \( a_{1} > x^{2n} + y^{2n} > x' \)
or
       \( \ip{A}{B} = \sum_{i} a_{i} b_{i} \).
The spaces you type in a formula are
ignored.  Remember that a letter like
       $x$                   % $ ... $  and  \( ... \)  are equivalent
is a formula when it denotes a mathematical
symbol, and it should be typed as one.

\section{Displayed Text}

Text is displayed by indenting it from the left
margin.  Quotations are commonly displayed.  There
are short quotations
\begin{quote}
   This is a short quotation.  It consists of a
   single paragraph of text.  See how it is formatted.
\end{quote}
and longer ones.
\begin{quotation}
   This is a longer quotation.  It consists of two
   paragraphs of text, neither of which are
   particularly interesting.

   This is the second paragraph of the quotation.  It
   is just as dull as the first paragraph.
\end{quotation}
Another frequently-displayed structure is a list.
The following is an example of an \emph{itemized}
list.
\begin{itemize}
   \item This is the first item of an itemized list.
         Each item in the list is marked with a ``tick''.
         You don't have to worry about what kind of tick
         mark is used.

   \item This is the second item of the list.  It
         contains another list nested inside it.  The inner
         list is an \emph{enumerated} list.
         \begin{enumerate}
            \item This is the first item of an enumerated
                  list that is nested within the itemized list.

            \item This is the second item of the inner list.
                  \LaTeX\ allows you to nest lists deeper than
                  you really should.
         \end{enumerate}
         This is the rest of the second item of the outer
         list.  It is no more interesting than any other
         part of the item.
   \item This is the third item of the list.
\end{itemize}
You can even display poetry.
\begin{verse}
   There is an environment
    for verse \\             % The \\ command separates lines
   Whose features some poets % within a stanza.
   will curse.

                             % One or more blank lines separate stanzas.

   For instead of making\\
   Them do \emph{all} line breaking, \\
   It allows them to put too many words on a line when they'd rather be
   forced to be terse.
\end{verse}

Mathematical formulas may also be displayed.  A
displayed formula
is
one-line long; multiline
formulas require special formatting instructions.
   \[  \ip{\Gamma}{\psi'} = x'' + y^{2} + z_{i}^{n}\]
Don't start a paragraph with a displayed equation,
nor make one a paragraph by itself.

\end{document}               % End of document.

\section{Bench Huge Section 81}
% synthetic bench section 81
% This is a sample LaTeX input file.  (Version of 12 August 2004.)
%
% A '%' character causes TeX to ignore all remaining text on the line,
% and is used for comments like this one.

\documentclass{article}      % Specifies the document class

                             % The preamble begins here.
\title{An Example Document}  % Declares the document's title.
\author{Leslie Lamport}      % Declares the author's name.
\date{January 21, 1994}      % Deleting this command produces today's date.

\newcommand{\ip}[2]{(#1, #2)}
                             % Defines \ip{arg1}{arg2} to mean
                             % (arg1, arg2).

%\newcommand{\ip}[2]{\langle #1 | #2\rangle}
                             % This is an alternative definition of
                             % \ip that is commented out.

\begin{document}             % End of preamble and beginning of text.

\maketitle                   % Produces the title.

This is an example input file.  Comparing it with
the output it generates can show you how to
produce a simple document of your own.

\section{Ordinary Text}      % Produces section heading.  Lower-level
                             % sections are begun with similar
                             % \subsection and \subsubsection commands.

The ends  of words and sentences are marked
  by   spaces. It  doesn't matter how many
spaces    you type; one is as good as 100.  The
end of   a line counts as a space.

One   or more   blank lines denote the  end
of  a paragraph.

Since any number of consecutive spaces are treated
like a single one, the formatting of the input
file makes no difference to
      \LaTeX,                % The \LaTeX command generates the LaTeX logo.
but it makes a difference to you.  When you use
\LaTeX, making your input file as easy to read
as possible will be a great help as you write
your document and when you change it.  This sample
file shows how you can add comments to your own input
file.

Because printing is different from typewriting,
there are a number of things that you have to do
differently when preparing an input file than if
you were just typing the document directly.
Quotation marks like
       ``this''
have to be handled specially, as do quotes within
quotes:
       ``\,`this'            % \, separates the double and single quote.
        is what I just
        wrote, not  `that'\,''.

Dashes come in three sizes: an
       intra-word
dash, a medium dash for number ranges like
       1--2,
and a punctuation
       dash---like
this.

A sentence-ending space should be larger than the
space between words within a sentence.  You
sometimes have to type special commands in
conjunction with punctuation characters to get
this right, as in the following sentence.
       Gnats, gnus, etc.\ all  % `\ ' makes an inter-word space.
       begin with G\@.         % \@ marks end-of-sentence punctuation.
You should check the spaces after periods when
reading your output to make sure you haven't
forgotten any special cases.  Generating an
ellipsis
       \ldots\               % `\ ' is needed after `\ldots' because TeX
                             % ignores spaces after command names like \ldots
                             % made from \ + letters.
                             %
                             % Note how a `%' character causes TeX to ignore
                             % the end of the input line, so these blank lines
                             % do not start a new paragraph.
                             %
with the right spacing around the periods requires
a special command.

\LaTeX\ interprets some common characters as
commands, so you must type special commands to
generate them.  These characters include the
following:
       \$ \& \% \# \{ and \}.

In printing, text is usually emphasized with an
       \emph{italic}
type style.

\begin{em}
   A long segment of text can also be emphasized
   in this way.  Text within such a segment can be
   given \emph{additional} emphasis.
\end{em}

It is sometimes necessary to prevent \LaTeX\ from
breaking a line where it might otherwise do so.
This may be at a space, as between the ``Mr.''\ and
``Jones'' in
       ``Mr.~Jones'',        % ~ produces an unbreakable interword space.
or within a word---especially when the word is a
symbol like
       \mbox{\emph{itemnum}}
that makes little sense when hyphenated across
lines.

Footnotes\footnote{This is an example of a footnote.}
pose no problem.

\LaTeX\ is good at typesetting mathematical formulas
like
       \( x-3y + z = 7 \)
or
       \( a_{1} > x^{2n} + y^{2n} > x' \)
or
       \( \ip{A}{B} = \sum_{i} a_{i} b_{i} \).
The spaces you type in a formula are
ignored.  Remember that a letter like
       $x$                   % $ ... $  and  \( ... \)  are equivalent
is a formula when it denotes a mathematical
symbol, and it should be typed as one.

\section{Displayed Text}

Text is displayed by indenting it from the left
margin.  Quotations are commonly displayed.  There
are short quotations
\begin{quote}
   This is a short quotation.  It consists of a
   single paragraph of text.  See how it is formatted.
\end{quote}
and longer ones.
\begin{quotation}
   This is a longer quotation.  It consists of two
   paragraphs of text, neither of which are
   particularly interesting.

   This is the second paragraph of the quotation.  It
   is just as dull as the first paragraph.
\end{quotation}
Another frequently-displayed structure is a list.
The following is an example of an \emph{itemized}
list.
\begin{itemize}
   \item This is the first item of an itemized list.
         Each item in the list is marked with a ``tick''.
         You don't have to worry about what kind of tick
         mark is used.

   \item This is the second item of the list.  It
         contains another list nested inside it.  The inner
         list is an \emph{enumerated} list.
         \begin{enumerate}
            \item This is the first item of an enumerated
                  list that is nested within the itemized list.

            \item This is the second item of the inner list.
                  \LaTeX\ allows you to nest lists deeper than
                  you really should.
         \end{enumerate}
         This is the rest of the second item of the outer
         list.  It is no more interesting than any other
         part of the item.
   \item This is the third item of the list.
\end{itemize}
You can even display poetry.
\begin{verse}
   There is an environment
    for verse \\             % The \\ command separates lines
   Whose features some poets % within a stanza.
   will curse.

                             % One or more blank lines separate stanzas.

   For instead of making\\
   Them do \emph{all} line breaking, \\
   It allows them to put too many words on a line when they'd rather be
   forced to be terse.
\end{verse}

Mathematical formulas may also be displayed.  A
displayed formula
is
one-line long; multiline
formulas require special formatting instructions.
   \[  \ip{\Gamma}{\psi'} = x'' + y^{2} + z_{i}^{n}\]
Don't start a paragraph with a displayed equation,
nor make one a paragraph by itself.

\end{document}               % End of document.

\section{Bench Huge Section 82}
% synthetic bench section 82
% This is a sample LaTeX input file.  (Version of 12 August 2004.)
%
% A '%' character causes TeX to ignore all remaining text on the line,
% and is used for comments like this one.

\documentclass{article}      % Specifies the document class

                             % The preamble begins here.
\title{An Example Document}  % Declares the document's title.
\author{Leslie Lamport}      % Declares the author's name.
\date{January 21, 1994}      % Deleting this command produces today's date.

\newcommand{\ip}[2]{(#1, #2)}
                             % Defines \ip{arg1}{arg2} to mean
                             % (arg1, arg2).

%\newcommand{\ip}[2]{\langle #1 | #2\rangle}
                             % This is an alternative definition of
                             % \ip that is commented out.

\begin{document}             % End of preamble and beginning of text.

\maketitle                   % Produces the title.

This is an example input file.  Comparing it with
the output it generates can show you how to
produce a simple document of your own.

\section{Ordinary Text}      % Produces section heading.  Lower-level
                             % sections are begun with similar
                             % \subsection and \subsubsection commands.

The ends  of words and sentences are marked
  by   spaces. It  doesn't matter how many
spaces    you type; one is as good as 100.  The
end of   a line counts as a space.

One   or more   blank lines denote the  end
of  a paragraph.

Since any number of consecutive spaces are treated
like a single one, the formatting of the input
file makes no difference to
      \LaTeX,                % The \LaTeX command generates the LaTeX logo.
but it makes a difference to you.  When you use
\LaTeX, making your input file as easy to read
as possible will be a great help as you write
your document and when you change it.  This sample
file shows how you can add comments to your own input
file.

Because printing is different from typewriting,
there are a number of things that you have to do
differently when preparing an input file than if
you were just typing the document directly.
Quotation marks like
       ``this''
have to be handled specially, as do quotes within
quotes:
       ``\,`this'            % \, separates the double and single quote.
        is what I just
        wrote, not  `that'\,''.

Dashes come in three sizes: an
       intra-word
dash, a medium dash for number ranges like
       1--2,
and a punctuation
       dash---like
this.

A sentence-ending space should be larger than the
space between words within a sentence.  You
sometimes have to type special commands in
conjunction with punctuation characters to get
this right, as in the following sentence.
       Gnats, gnus, etc.\ all  % `\ ' makes an inter-word space.
       begin with G\@.         % \@ marks end-of-sentence punctuation.
You should check the spaces after periods when
reading your output to make sure you haven't
forgotten any special cases.  Generating an
ellipsis
       \ldots\               % `\ ' is needed after `\ldots' because TeX
                             % ignores spaces after command names like \ldots
                             % made from \ + letters.
                             %
                             % Note how a `%' character causes TeX to ignore
                             % the end of the input line, so these blank lines
                             % do not start a new paragraph.
                             %
with the right spacing around the periods requires
a special command.

\LaTeX\ interprets some common characters as
commands, so you must type special commands to
generate them.  These characters include the
following:
       \$ \& \% \# \{ and \}.

In printing, text is usually emphasized with an
       \emph{italic}
type style.

\begin{em}
   A long segment of text can also be emphasized
   in this way.  Text within such a segment can be
   given \emph{additional} emphasis.
\end{em}

It is sometimes necessary to prevent \LaTeX\ from
breaking a line where it might otherwise do so.
This may be at a space, as between the ``Mr.''\ and
``Jones'' in
       ``Mr.~Jones'',        % ~ produces an unbreakable interword space.
or within a word---especially when the word is a
symbol like
       \mbox{\emph{itemnum}}
that makes little sense when hyphenated across
lines.

Footnotes\footnote{This is an example of a footnote.}
pose no problem.

\LaTeX\ is good at typesetting mathematical formulas
like
       \( x-3y + z = 7 \)
or
       \( a_{1} > x^{2n} + y^{2n} > x' \)
or
       \( \ip{A}{B} = \sum_{i} a_{i} b_{i} \).
The spaces you type in a formula are
ignored.  Remember that a letter like
       $x$                   % $ ... $  and  \( ... \)  are equivalent
is a formula when it denotes a mathematical
symbol, and it should be typed as one.

\section{Displayed Text}

Text is displayed by indenting it from the left
margin.  Quotations are commonly displayed.  There
are short quotations
\begin{quote}
   This is a short quotation.  It consists of a
   single paragraph of text.  See how it is formatted.
\end{quote}
and longer ones.
\begin{quotation}
   This is a longer quotation.  It consists of two
   paragraphs of text, neither of which are
   particularly interesting.

   This is the second paragraph of the quotation.  It
   is just as dull as the first paragraph.
\end{quotation}
Another frequently-displayed structure is a list.
The following is an example of an \emph{itemized}
list.
\begin{itemize}
   \item This is the first item of an itemized list.
         Each item in the list is marked with a ``tick''.
         You don't have to worry about what kind of tick
         mark is used.

   \item This is the second item of the list.  It
         contains another list nested inside it.  The inner
         list is an \emph{enumerated} list.
         \begin{enumerate}
            \item This is the first item of an enumerated
                  list that is nested within the itemized list.

            \item This is the second item of the inner list.
                  \LaTeX\ allows you to nest lists deeper than
                  you really should.
         \end{enumerate}
         This is the rest of the second item of the outer
         list.  It is no more interesting than any other
         part of the item.
   \item This is the third item of the list.
\end{itemize}
You can even display poetry.
\begin{verse}
   There is an environment
    for verse \\             % The \\ command separates lines
   Whose features some poets % within a stanza.
   will curse.

                             % One or more blank lines separate stanzas.

   For instead of making\\
   Them do \emph{all} line breaking, \\
   It allows them to put too many words on a line when they'd rather be
   forced to be terse.
\end{verse}

Mathematical formulas may also be displayed.  A
displayed formula
is
one-line long; multiline
formulas require special formatting instructions.
   \[  \ip{\Gamma}{\psi'} = x'' + y^{2} + z_{i}^{n}\]
Don't start a paragraph with a displayed equation,
nor make one a paragraph by itself.

\end{document}               % End of document.

\section{Bench Huge Section 83}
% synthetic bench section 83
% This is a sample LaTeX input file.  (Version of 12 August 2004.)
%
% A '%' character causes TeX to ignore all remaining text on the line,
% and is used for comments like this one.

\documentclass{article}      % Specifies the document class

                             % The preamble begins here.
\title{An Example Document}  % Declares the document's title.
\author{Leslie Lamport}      % Declares the author's name.
\date{January 21, 1994}      % Deleting this command produces today's date.

\newcommand{\ip}[2]{(#1, #2)}
                             % Defines \ip{arg1}{arg2} to mean
                             % (arg1, arg2).

%\newcommand{\ip}[2]{\langle #1 | #2\rangle}
                             % This is an alternative definition of
                             % \ip that is commented out.

\begin{document}             % End of preamble and beginning of text.

\maketitle                   % Produces the title.

This is an example input file.  Comparing it with
the output it generates can show you how to
produce a simple document of your own.

\section{Ordinary Text}      % Produces section heading.  Lower-level
                             % sections are begun with similar
                             % \subsection and \subsubsection commands.

The ends  of words and sentences are marked
  by   spaces. It  doesn't matter how many
spaces    you type; one is as good as 100.  The
end of   a line counts as a space.

One   or more   blank lines denote the  end
of  a paragraph.

Since any number of consecutive spaces are treated
like a single one, the formatting of the input
file makes no difference to
      \LaTeX,                % The \LaTeX command generates the LaTeX logo.
but it makes a difference to you.  When you use
\LaTeX, making your input file as easy to read
as possible will be a great help as you write
your document and when you change it.  This sample
file shows how you can add comments to your own input
file.

Because printing is different from typewriting,
there are a number of things that you have to do
differently when preparing an input file than if
you were just typing the document directly.
Quotation marks like
       ``this''
have to be handled specially, as do quotes within
quotes:
       ``\,`this'            % \, separates the double and single quote.
        is what I just
        wrote, not  `that'\,''.

Dashes come in three sizes: an
       intra-word
dash, a medium dash for number ranges like
       1--2,
and a punctuation
       dash---like
this.

A sentence-ending space should be larger than the
space between words within a sentence.  You
sometimes have to type special commands in
conjunction with punctuation characters to get
this right, as in the following sentence.
       Gnats, gnus, etc.\ all  % `\ ' makes an inter-word space.
       begin with G\@.         % \@ marks end-of-sentence punctuation.
You should check the spaces after periods when
reading your output to make sure you haven't
forgotten any special cases.  Generating an
ellipsis
       \ldots\               % `\ ' is needed after `\ldots' because TeX
                             % ignores spaces after command names like \ldots
                             % made from \ + letters.
                             %
                             % Note how a `%' character causes TeX to ignore
                             % the end of the input line, so these blank lines
                             % do not start a new paragraph.
                             %
with the right spacing around the periods requires
a special command.

\LaTeX\ interprets some common characters as
commands, so you must type special commands to
generate them.  These characters include the
following:
       \$ \& \% \# \{ and \}.

In printing, text is usually emphasized with an
       \emph{italic}
type style.

\begin{em}
   A long segment of text can also be emphasized
   in this way.  Text within such a segment can be
   given \emph{additional} emphasis.
\end{em}

It is sometimes necessary to prevent \LaTeX\ from
breaking a line where it might otherwise do so.
This may be at a space, as between the ``Mr.''\ and
``Jones'' in
       ``Mr.~Jones'',        % ~ produces an unbreakable interword space.
or within a word---especially when the word is a
symbol like
       \mbox{\emph{itemnum}}
that makes little sense when hyphenated across
lines.

Footnotes\footnote{This is an example of a footnote.}
pose no problem.

\LaTeX\ is good at typesetting mathematical formulas
like
       \( x-3y + z = 7 \)
or
       \( a_{1} > x^{2n} + y^{2n} > x' \)
or
       \( \ip{A}{B} = \sum_{i} a_{i} b_{i} \).
The spaces you type in a formula are
ignored.  Remember that a letter like
       $x$                   % $ ... $  and  \( ... \)  are equivalent
is a formula when it denotes a mathematical
symbol, and it should be typed as one.

\section{Displayed Text}

Text is displayed by indenting it from the left
margin.  Quotations are commonly displayed.  There
are short quotations
\begin{quote}
   This is a short quotation.  It consists of a
   single paragraph of text.  See how it is formatted.
\end{quote}
and longer ones.
\begin{quotation}
   This is a longer quotation.  It consists of two
   paragraphs of text, neither of which are
   particularly interesting.

   This is the second paragraph of the quotation.  It
   is just as dull as the first paragraph.
\end{quotation}
Another frequently-displayed structure is a list.
The following is an example of an \emph{itemized}
list.
\begin{itemize}
   \item This is the first item of an itemized list.
         Each item in the list is marked with a ``tick''.
         You don't have to worry about what kind of tick
         mark is used.

   \item This is the second item of the list.  It
         contains another list nested inside it.  The inner
         list is an \emph{enumerated} list.
         \begin{enumerate}
            \item This is the first item of an enumerated
                  list that is nested within the itemized list.

            \item This is the second item of the inner list.
                  \LaTeX\ allows you to nest lists deeper than
                  you really should.
         \end{enumerate}
         This is the rest of the second item of the outer
         list.  It is no more interesting than any other
         part of the item.
   \item This is the third item of the list.
\end{itemize}
You can even display poetry.
\begin{verse}
   There is an environment
    for verse \\             % The \\ command separates lines
   Whose features some poets % within a stanza.
   will curse.

                             % One or more blank lines separate stanzas.

   For instead of making\\
   Them do \emph{all} line breaking, \\
   It allows them to put too many words on a line when they'd rather be
   forced to be terse.
\end{verse}

Mathematical formulas may also be displayed.  A
displayed formula
is
one-line long; multiline
formulas require special formatting instructions.
   \[  \ip{\Gamma}{\psi'} = x'' + y^{2} + z_{i}^{n}\]
Don't start a paragraph with a displayed equation,
nor make one a paragraph by itself.

\end{document}               % End of document.

\section{Bench Huge Section 84}
% synthetic bench section 84
% This is a sample LaTeX input file.  (Version of 12 August 2004.)
%
% A '%' character causes TeX to ignore all remaining text on the line,
% and is used for comments like this one.

\documentclass{article}      % Specifies the document class

                             % The preamble begins here.
\title{An Example Document}  % Declares the document's title.
\author{Leslie Lamport}      % Declares the author's name.
\date{January 21, 1994}      % Deleting this command produces today's date.

\newcommand{\ip}[2]{(#1, #2)}
                             % Defines \ip{arg1}{arg2} to mean
                             % (arg1, arg2).

%\newcommand{\ip}[2]{\langle #1 | #2\rangle}
                             % This is an alternative definition of
                             % \ip that is commented out.

\begin{document}             % End of preamble and beginning of text.

\maketitle                   % Produces the title.

This is an example input file.  Comparing it with
the output it generates can show you how to
produce a simple document of your own.

\section{Ordinary Text}      % Produces section heading.  Lower-level
                             % sections are begun with similar
                             % \subsection and \subsubsection commands.

The ends  of words and sentences are marked
  by   spaces. It  doesn't matter how many
spaces    you type; one is as good as 100.  The
end of   a line counts as a space.

One   or more   blank lines denote the  end
of  a paragraph.

Since any number of consecutive spaces are treated
like a single one, the formatting of the input
file makes no difference to
      \LaTeX,                % The \LaTeX command generates the LaTeX logo.
but it makes a difference to you.  When you use
\LaTeX, making your input file as easy to read
as possible will be a great help as you write
your document and when you change it.  This sample
file shows how you can add comments to your own input
file.

Because printing is different from typewriting,
there are a number of things that you have to do
differently when preparing an input file than if
you were just typing the document directly.
Quotation marks like
       ``this''
have to be handled specially, as do quotes within
quotes:
       ``\,`this'            % \, separates the double and single quote.
        is what I just
        wrote, not  `that'\,''.

Dashes come in three sizes: an
       intra-word
dash, a medium dash for number ranges like
       1--2,
and a punctuation
       dash---like
this.

A sentence-ending space should be larger than the
space between words within a sentence.  You
sometimes have to type special commands in
conjunction with punctuation characters to get
this right, as in the following sentence.
       Gnats, gnus, etc.\ all  % `\ ' makes an inter-word space.
       begin with G\@.         % \@ marks end-of-sentence punctuation.
You should check the spaces after periods when
reading your output to make sure you haven't
forgotten any special cases.  Generating an
ellipsis
       \ldots\               % `\ ' is needed after `\ldots' because TeX
                             % ignores spaces after command names like \ldots
                             % made from \ + letters.
                             %
                             % Note how a `%' character causes TeX to ignore
                             % the end of the input line, so these blank lines
                             % do not start a new paragraph.
                             %
with the right spacing around the periods requires
a special command.

\LaTeX\ interprets some common characters as
commands, so you must type special commands to
generate them.  These characters include the
following:
       \$ \& \% \# \{ and \}.

In printing, text is usually emphasized with an
       \emph{italic}
type style.

\begin{em}
   A long segment of text can also be emphasized
   in this way.  Text within such a segment can be
   given \emph{additional} emphasis.
\end{em}

It is sometimes necessary to prevent \LaTeX\ from
breaking a line where it might otherwise do so.
This may be at a space, as between the ``Mr.''\ and
``Jones'' in
       ``Mr.~Jones'',        % ~ produces an unbreakable interword space.
or within a word---especially when the word is a
symbol like
       \mbox{\emph{itemnum}}
that makes little sense when hyphenated across
lines.

Footnotes\footnote{This is an example of a footnote.}
pose no problem.

\LaTeX\ is good at typesetting mathematical formulas
like
       \( x-3y + z = 7 \)
or
       \( a_{1} > x^{2n} + y^{2n} > x' \)
or
       \( \ip{A}{B} = \sum_{i} a_{i} b_{i} \).
The spaces you type in a formula are
ignored.  Remember that a letter like
       $x$                   % $ ... $  and  \( ... \)  are equivalent
is a formula when it denotes a mathematical
symbol, and it should be typed as one.

\section{Displayed Text}

Text is displayed by indenting it from the left
margin.  Quotations are commonly displayed.  There
are short quotations
\begin{quote}
   This is a short quotation.  It consists of a
   single paragraph of text.  See how it is formatted.
\end{quote}
and longer ones.
\begin{quotation}
   This is a longer quotation.  It consists of two
   paragraphs of text, neither of which are
   particularly interesting.

   This is the second paragraph of the quotation.  It
   is just as dull as the first paragraph.
\end{quotation}
Another frequently-displayed structure is a list.
The following is an example of an \emph{itemized}
list.
\begin{itemize}
   \item This is the first item of an itemized list.
         Each item in the list is marked with a ``tick''.
         You don't have to worry about what kind of tick
         mark is used.

   \item This is the second item of the list.  It
         contains another list nested inside it.  The inner
         list is an \emph{enumerated} list.
         \begin{enumerate}
            \item This is the first item of an enumerated
                  list that is nested within the itemized list.

            \item This is the second item of the inner list.
                  \LaTeX\ allows you to nest lists deeper than
                  you really should.
         \end{enumerate}
         This is the rest of the second item of the outer
         list.  It is no more interesting than any other
         part of the item.
   \item This is the third item of the list.
\end{itemize}
You can even display poetry.
\begin{verse}
   There is an environment
    for verse \\             % The \\ command separates lines
   Whose features some poets % within a stanza.
   will curse.

                             % One or more blank lines separate stanzas.

   For instead of making\\
   Them do \emph{all} line breaking, \\
   It allows them to put too many words on a line when they'd rather be
   forced to be terse.
\end{verse}

Mathematical formulas may also be displayed.  A
displayed formula
is
one-line long; multiline
formulas require special formatting instructions.
   \[  \ip{\Gamma}{\psi'} = x'' + y^{2} + z_{i}^{n}\]
Don't start a paragraph with a displayed equation,
nor make one a paragraph by itself.

\end{document}               % End of document.

\section{Bench Huge Section 85}
% synthetic bench section 85
% This is a sample LaTeX input file.  (Version of 12 August 2004.)
%
% A '%' character causes TeX to ignore all remaining text on the line,
% and is used for comments like this one.

\documentclass{article}      % Specifies the document class

                             % The preamble begins here.
\title{An Example Document}  % Declares the document's title.
\author{Leslie Lamport}      % Declares the author's name.
\date{January 21, 1994}      % Deleting this command produces today's date.

\newcommand{\ip}[2]{(#1, #2)}
                             % Defines \ip{arg1}{arg2} to mean
                             % (arg1, arg2).

%\newcommand{\ip}[2]{\langle #1 | #2\rangle}
                             % This is an alternative definition of
                             % \ip that is commented out.

\begin{document}             % End of preamble and beginning of text.

\maketitle                   % Produces the title.

This is an example input file.  Comparing it with
the output it generates can show you how to
produce a simple document of your own.

\section{Ordinary Text}      % Produces section heading.  Lower-level
                             % sections are begun with similar
                             % \subsection and \subsubsection commands.

The ends  of words and sentences are marked
  by   spaces. It  doesn't matter how many
spaces    you type; one is as good as 100.  The
end of   a line counts as a space.

One   or more   blank lines denote the  end
of  a paragraph.

Since any number of consecutive spaces are treated
like a single one, the formatting of the input
file makes no difference to
      \LaTeX,                % The \LaTeX command generates the LaTeX logo.
but it makes a difference to you.  When you use
\LaTeX, making your input file as easy to read
as possible will be a great help as you write
your document and when you change it.  This sample
file shows how you can add comments to your own input
file.

Because printing is different from typewriting,
there are a number of things that you have to do
differently when preparing an input file than if
you were just typing the document directly.
Quotation marks like
       ``this''
have to be handled specially, as do quotes within
quotes:
       ``\,`this'            % \, separates the double and single quote.
        is what I just
        wrote, not  `that'\,''.

Dashes come in three sizes: an
       intra-word
dash, a medium dash for number ranges like
       1--2,
and a punctuation
       dash---like
this.

A sentence-ending space should be larger than the
space between words within a sentence.  You
sometimes have to type special commands in
conjunction with punctuation characters to get
this right, as in the following sentence.
       Gnats, gnus, etc.\ all  % `\ ' makes an inter-word space.
       begin with G\@.         % \@ marks end-of-sentence punctuation.
You should check the spaces after periods when
reading your output to make sure you haven't
forgotten any special cases.  Generating an
ellipsis
       \ldots\               % `\ ' is needed after `\ldots' because TeX
                             % ignores spaces after command names like \ldots
                             % made from \ + letters.
                             %
                             % Note how a `%' character causes TeX to ignore
                             % the end of the input line, so these blank lines
                             % do not start a new paragraph.
                             %
with the right spacing around the periods requires
a special command.

\LaTeX\ interprets some common characters as
commands, so you must type special commands to
generate them.  These characters include the
following:
       \$ \& \% \# \{ and \}.

In printing, text is usually emphasized with an
       \emph{italic}
type style.

\begin{em}
   A long segment of text can also be emphasized
   in this way.  Text within such a segment can be
   given \emph{additional} emphasis.
\end{em}

It is sometimes necessary to prevent \LaTeX\ from
breaking a line where it might otherwise do so.
This may be at a space, as between the ``Mr.''\ and
``Jones'' in
       ``Mr.~Jones'',        % ~ produces an unbreakable interword space.
or within a word---especially when the word is a
symbol like
       \mbox{\emph{itemnum}}
that makes little sense when hyphenated across
lines.

Footnotes\footnote{This is an example of a footnote.}
pose no problem.

\LaTeX\ is good at typesetting mathematical formulas
like
       \( x-3y + z = 7 \)
or
       \( a_{1} > x^{2n} + y^{2n} > x' \)
or
       \( \ip{A}{B} = \sum_{i} a_{i} b_{i} \).
The spaces you type in a formula are
ignored.  Remember that a letter like
       $x$                   % $ ... $  and  \( ... \)  are equivalent
is a formula when it denotes a mathematical
symbol, and it should be typed as one.

\section{Displayed Text}

Text is displayed by indenting it from the left
margin.  Quotations are commonly displayed.  There
are short quotations
\begin{quote}
   This is a short quotation.  It consists of a
   single paragraph of text.  See how it is formatted.
\end{quote}
and longer ones.
\begin{quotation}
   This is a longer quotation.  It consists of two
   paragraphs of text, neither of which are
   particularly interesting.

   This is the second paragraph of the quotation.  It
   is just as dull as the first paragraph.
\end{quotation}
Another frequently-displayed structure is a list.
The following is an example of an \emph{itemized}
list.
\begin{itemize}
   \item This is the first item of an itemized list.
         Each item in the list is marked with a ``tick''.
         You don't have to worry about what kind of tick
         mark is used.

   \item This is the second item of the list.  It
         contains another list nested inside it.  The inner
         list is an \emph{enumerated} list.
         \begin{enumerate}
            \item This is the first item of an enumerated
                  list that is nested within the itemized list.

            \item This is the second item of the inner list.
                  \LaTeX\ allows you to nest lists deeper than
                  you really should.
         \end{enumerate}
         This is the rest of the second item of the outer
         list.  It is no more interesting than any other
         part of the item.
   \item This is the third item of the list.
\end{itemize}
You can even display poetry.
\begin{verse}
   There is an environment
    for verse \\             % The \\ command separates lines
   Whose features some poets % within a stanza.
   will curse.

                             % One or more blank lines separate stanzas.

   For instead of making\\
   Them do \emph{all} line breaking, \\
   It allows them to put too many words on a line when they'd rather be
   forced to be terse.
\end{verse}

Mathematical formulas may also be displayed.  A
displayed formula
is
one-line long; multiline
formulas require special formatting instructions.
   \[  \ip{\Gamma}{\psi'} = x'' + y^{2} + z_{i}^{n}\]
Don't start a paragraph with a displayed equation,
nor make one a paragraph by itself.

\end{document}               % End of document.

\section{Bench Huge Section 86}
% synthetic bench section 86
% This is a sample LaTeX input file.  (Version of 12 August 2004.)
%
% A '%' character causes TeX to ignore all remaining text on the line,
% and is used for comments like this one.

\documentclass{article}      % Specifies the document class

                             % The preamble begins here.
\title{An Example Document}  % Declares the document's title.
\author{Leslie Lamport}      % Declares the author's name.
\date{January 21, 1994}      % Deleting this command produces today's date.

\newcommand{\ip}[2]{(#1, #2)}
                             % Defines \ip{arg1}{arg2} to mean
                             % (arg1, arg2).

%\newcommand{\ip}[2]{\langle #1 | #2\rangle}
                             % This is an alternative definition of
                             % \ip that is commented out.

\begin{document}             % End of preamble and beginning of text.

\maketitle                   % Produces the title.

This is an example input file.  Comparing it with
the output it generates can show you how to
produce a simple document of your own.

\section{Ordinary Text}      % Produces section heading.  Lower-level
                             % sections are begun with similar
                             % \subsection and \subsubsection commands.

The ends  of words and sentences are marked
  by   spaces. It  doesn't matter how many
spaces    you type; one is as good as 100.  The
end of   a line counts as a space.

One   or more   blank lines denote the  end
of  a paragraph.

Since any number of consecutive spaces are treated
like a single one, the formatting of the input
file makes no difference to
      \LaTeX,                % The \LaTeX command generates the LaTeX logo.
but it makes a difference to you.  When you use
\LaTeX, making your input file as easy to read
as possible will be a great help as you write
your document and when you change it.  This sample
file shows how you can add comments to your own input
file.

Because printing is different from typewriting,
there are a number of things that you have to do
differently when preparing an input file than if
you were just typing the document directly.
Quotation marks like
       ``this''
have to be handled specially, as do quotes within
quotes:
       ``\,`this'            % \, separates the double and single quote.
        is what I just
        wrote, not  `that'\,''.

Dashes come in three sizes: an
       intra-word
dash, a medium dash for number ranges like
       1--2,
and a punctuation
       dash---like
this.

A sentence-ending space should be larger than the
space between words within a sentence.  You
sometimes have to type special commands in
conjunction with punctuation characters to get
this right, as in the following sentence.
       Gnats, gnus, etc.\ all  % `\ ' makes an inter-word space.
       begin with G\@.         % \@ marks end-of-sentence punctuation.
You should check the spaces after periods when
reading your output to make sure you haven't
forgotten any special cases.  Generating an
ellipsis
       \ldots\               % `\ ' is needed after `\ldots' because TeX
                             % ignores spaces after command names like \ldots
                             % made from \ + letters.
                             %
                             % Note how a `%' character causes TeX to ignore
                             % the end of the input line, so these blank lines
                             % do not start a new paragraph.
                             %
with the right spacing around the periods requires
a special command.

\LaTeX\ interprets some common characters as
commands, so you must type special commands to
generate them.  These characters include the
following:
       \$ \& \% \# \{ and \}.

In printing, text is usually emphasized with an
       \emph{italic}
type style.

\begin{em}
   A long segment of text can also be emphasized
   in this way.  Text within such a segment can be
   given \emph{additional} emphasis.
\end{em}

It is sometimes necessary to prevent \LaTeX\ from
breaking a line where it might otherwise do so.
This may be at a space, as between the ``Mr.''\ and
``Jones'' in
       ``Mr.~Jones'',        % ~ produces an unbreakable interword space.
or within a word---especially when the word is a
symbol like
       \mbox{\emph{itemnum}}
that makes little sense when hyphenated across
lines.

Footnotes\footnote{This is an example of a footnote.}
pose no problem.

\LaTeX\ is good at typesetting mathematical formulas
like
       \( x-3y + z = 7 \)
or
       \( a_{1} > x^{2n} + y^{2n} > x' \)
or
       \( \ip{A}{B} = \sum_{i} a_{i} b_{i} \).
The spaces you type in a formula are
ignored.  Remember that a letter like
       $x$                   % $ ... $  and  \( ... \)  are equivalent
is a formula when it denotes a mathematical
symbol, and it should be typed as one.

\section{Displayed Text}

Text is displayed by indenting it from the left
margin.  Quotations are commonly displayed.  There
are short quotations
\begin{quote}
   This is a short quotation.  It consists of a
   single paragraph of text.  See how it is formatted.
\end{quote}
and longer ones.
\begin{quotation}
   This is a longer quotation.  It consists of two
   paragraphs of text, neither of which are
   particularly interesting.

   This is the second paragraph of the quotation.  It
   is just as dull as the first paragraph.
\end{quotation}
Another frequently-displayed structure is a list.
The following is an example of an \emph{itemized}
list.
\begin{itemize}
   \item This is the first item of an itemized list.
         Each item in the list is marked with a ``tick''.
         You don't have to worry about what kind of tick
         mark is used.

   \item This is the second item of the list.  It
         contains another list nested inside it.  The inner
         list is an \emph{enumerated} list.
         \begin{enumerate}
            \item This is the first item of an enumerated
                  list that is nested within the itemized list.

            \item This is the second item of the inner list.
                  \LaTeX\ allows you to nest lists deeper than
                  you really should.
         \end{enumerate}
         This is the rest of the second item of the outer
         list.  It is no more interesting than any other
         part of the item.
   \item This is the third item of the list.
\end{itemize}
You can even display poetry.
\begin{verse}
   There is an environment
    for verse \\             % The \\ command separates lines
   Whose features some poets % within a stanza.
   will curse.

                             % One or more blank lines separate stanzas.

   For instead of making\\
   Them do \emph{all} line breaking, \\
   It allows them to put too many words on a line when they'd rather be
   forced to be terse.
\end{verse}

Mathematical formulas may also be displayed.  A
displayed formula
is
one-line long; multiline
formulas require special formatting instructions.
   \[  \ip{\Gamma}{\psi'} = x'' + y^{2} + z_{i}^{n}\]
Don't start a paragraph with a displayed equation,
nor make one a paragraph by itself.

\end{document}               % End of document.

\section{Bench Huge Section 87}
% synthetic bench section 87
% This is a sample LaTeX input file.  (Version of 12 August 2004.)
%
% A '%' character causes TeX to ignore all remaining text on the line,
% and is used for comments like this one.

\documentclass{article}      % Specifies the document class

                             % The preamble begins here.
\title{An Example Document}  % Declares the document's title.
\author{Leslie Lamport}      % Declares the author's name.
\date{January 21, 1994}      % Deleting this command produces today's date.

\newcommand{\ip}[2]{(#1, #2)}
                             % Defines \ip{arg1}{arg2} to mean
                             % (arg1, arg2).

%\newcommand{\ip}[2]{\langle #1 | #2\rangle}
                             % This is an alternative definition of
                             % \ip that is commented out.

\begin{document}             % End of preamble and beginning of text.

\maketitle                   % Produces the title.

This is an example input file.  Comparing it with
the output it generates can show you how to
produce a simple document of your own.

\section{Ordinary Text}      % Produces section heading.  Lower-level
                             % sections are begun with similar
                             % \subsection and \subsubsection commands.

The ends  of words and sentences are marked
  by   spaces. It  doesn't matter how many
spaces    you type; one is as good as 100.  The
end of   a line counts as a space.

One   or more   blank lines denote the  end
of  a paragraph.

Since any number of consecutive spaces are treated
like a single one, the formatting of the input
file makes no difference to
      \LaTeX,                % The \LaTeX command generates the LaTeX logo.
but it makes a difference to you.  When you use
\LaTeX, making your input file as easy to read
as possible will be a great help as you write
your document and when you change it.  This sample
file shows how you can add comments to your own input
file.

Because printing is different from typewriting,
there are a number of things that you have to do
differently when preparing an input file than if
you were just typing the document directly.
Quotation marks like
       ``this''
have to be handled specially, as do quotes within
quotes:
       ``\,`this'            % \, separates the double and single quote.
        is what I just
        wrote, not  `that'\,''.

Dashes come in three sizes: an
       intra-word
dash, a medium dash for number ranges like
       1--2,
and a punctuation
       dash---like
this.

A sentence-ending space should be larger than the
space between words within a sentence.  You
sometimes have to type special commands in
conjunction with punctuation characters to get
this right, as in the following sentence.
       Gnats, gnus, etc.\ all  % `\ ' makes an inter-word space.
       begin with G\@.         % \@ marks end-of-sentence punctuation.
You should check the spaces after periods when
reading your output to make sure you haven't
forgotten any special cases.  Generating an
ellipsis
       \ldots\               % `\ ' is needed after `\ldots' because TeX
                             % ignores spaces after command names like \ldots
                             % made from \ + letters.
                             %
                             % Note how a `%' character causes TeX to ignore
                             % the end of the input line, so these blank lines
                             % do not start a new paragraph.
                             %
with the right spacing around the periods requires
a special command.

\LaTeX\ interprets some common characters as
commands, so you must type special commands to
generate them.  These characters include the
following:
       \$ \& \% \# \{ and \}.

In printing, text is usually emphasized with an
       \emph{italic}
type style.

\begin{em}
   A long segment of text can also be emphasized
   in this way.  Text within such a segment can be
   given \emph{additional} emphasis.
\end{em}

It is sometimes necessary to prevent \LaTeX\ from
breaking a line where it might otherwise do so.
This may be at a space, as between the ``Mr.''\ and
``Jones'' in
       ``Mr.~Jones'',        % ~ produces an unbreakable interword space.
or within a word---especially when the word is a
symbol like
       \mbox{\emph{itemnum}}
that makes little sense when hyphenated across
lines.

Footnotes\footnote{This is an example of a footnote.}
pose no problem.

\LaTeX\ is good at typesetting mathematical formulas
like
       \( x-3y + z = 7 \)
or
       \( a_{1} > x^{2n} + y^{2n} > x' \)
or
       \( \ip{A}{B} = \sum_{i} a_{i} b_{i} \).
The spaces you type in a formula are
ignored.  Remember that a letter like
       $x$                   % $ ... $  and  \( ... \)  are equivalent
is a formula when it denotes a mathematical
symbol, and it should be typed as one.

\section{Displayed Text}

Text is displayed by indenting it from the left
margin.  Quotations are commonly displayed.  There
are short quotations
\begin{quote}
   This is a short quotation.  It consists of a
   single paragraph of text.  See how it is formatted.
\end{quote}
and longer ones.
\begin{quotation}
   This is a longer quotation.  It consists of two
   paragraphs of text, neither of which are
   particularly interesting.

   This is the second paragraph of the quotation.  It
   is just as dull as the first paragraph.
\end{quotation}
Another frequently-displayed structure is a list.
The following is an example of an \emph{itemized}
list.
\begin{itemize}
   \item This is the first item of an itemized list.
         Each item in the list is marked with a ``tick''.
         You don't have to worry about what kind of tick
         mark is used.

   \item This is the second item of the list.  It
         contains another list nested inside it.  The inner
         list is an \emph{enumerated} list.
         \begin{enumerate}
            \item This is the first item of an enumerated
                  list that is nested within the itemized list.

            \item This is the second item of the inner list.
                  \LaTeX\ allows you to nest lists deeper than
                  you really should.
         \end{enumerate}
         This is the rest of the second item of the outer
         list.  It is no more interesting than any other
         part of the item.
   \item This is the third item of the list.
\end{itemize}
You can even display poetry.
\begin{verse}
   There is an environment
    for verse \\             % The \\ command separates lines
   Whose features some poets % within a stanza.
   will curse.

                             % One or more blank lines separate stanzas.

   For instead of making\\
   Them do \emph{all} line breaking, \\
   It allows them to put too many words on a line when they'd rather be
   forced to be terse.
\end{verse}

Mathematical formulas may also be displayed.  A
displayed formula
is
one-line long; multiline
formulas require special formatting instructions.
   \[  \ip{\Gamma}{\psi'} = x'' + y^{2} + z_{i}^{n}\]
Don't start a paragraph with a displayed equation,
nor make one a paragraph by itself.

\end{document}               % End of document.

\section{Bench Huge Section 88}
% synthetic bench section 88
% This is a sample LaTeX input file.  (Version of 12 August 2004.)
%
% A '%' character causes TeX to ignore all remaining text on the line,
% and is used for comments like this one.

\documentclass{article}      % Specifies the document class

                             % The preamble begins here.
\title{An Example Document}  % Declares the document's title.
\author{Leslie Lamport}      % Declares the author's name.
\date{January 21, 1994}      % Deleting this command produces today's date.

\newcommand{\ip}[2]{(#1, #2)}
                             % Defines \ip{arg1}{arg2} to mean
                             % (arg1, arg2).

%\newcommand{\ip}[2]{\langle #1 | #2\rangle}
                             % This is an alternative definition of
                             % \ip that is commented out.

\begin{document}             % End of preamble and beginning of text.

\maketitle                   % Produces the title.

This is an example input file.  Comparing it with
the output it generates can show you how to
produce a simple document of your own.

\section{Ordinary Text}      % Produces section heading.  Lower-level
                             % sections are begun with similar
                             % \subsection and \subsubsection commands.

The ends  of words and sentences are marked
  by   spaces. It  doesn't matter how many
spaces    you type; one is as good as 100.  The
end of   a line counts as a space.

One   or more   blank lines denote the  end
of  a paragraph.

Since any number of consecutive spaces are treated
like a single one, the formatting of the input
file makes no difference to
      \LaTeX,                % The \LaTeX command generates the LaTeX logo.
but it makes a difference to you.  When you use
\LaTeX, making your input file as easy to read
as possible will be a great help as you write
your document and when you change it.  This sample
file shows how you can add comments to your own input
file.

Because printing is different from typewriting,
there are a number of things that you have to do
differently when preparing an input file than if
you were just typing the document directly.
Quotation marks like
       ``this''
have to be handled specially, as do quotes within
quotes:
       ``\,`this'            % \, separates the double and single quote.
        is what I just
        wrote, not  `that'\,''.

Dashes come in three sizes: an
       intra-word
dash, a medium dash for number ranges like
       1--2,
and a punctuation
       dash---like
this.

A sentence-ending space should be larger than the
space between words within a sentence.  You
sometimes have to type special commands in
conjunction with punctuation characters to get
this right, as in the following sentence.
       Gnats, gnus, etc.\ all  % `\ ' makes an inter-word space.
       begin with G\@.         % \@ marks end-of-sentence punctuation.
You should check the spaces after periods when
reading your output to make sure you haven't
forgotten any special cases.  Generating an
ellipsis
       \ldots\               % `\ ' is needed after `\ldots' because TeX
                             % ignores spaces after command names like \ldots
                             % made from \ + letters.
                             %
                             % Note how a `%' character causes TeX to ignore
                             % the end of the input line, so these blank lines
                             % do not start a new paragraph.
                             %
with the right spacing around the periods requires
a special command.

\LaTeX\ interprets some common characters as
commands, so you must type special commands to
generate them.  These characters include the
following:
       \$ \& \% \# \{ and \}.

In printing, text is usually emphasized with an
       \emph{italic}
type style.

\begin{em}
   A long segment of text can also be emphasized
   in this way.  Text within such a segment can be
   given \emph{additional} emphasis.
\end{em}

It is sometimes necessary to prevent \LaTeX\ from
breaking a line where it might otherwise do so.
This may be at a space, as between the ``Mr.''\ and
``Jones'' in
       ``Mr.~Jones'',        % ~ produces an unbreakable interword space.
or within a word---especially when the word is a
symbol like
       \mbox{\emph{itemnum}}
that makes little sense when hyphenated across
lines.

Footnotes\footnote{This is an example of a footnote.}
pose no problem.

\LaTeX\ is good at typesetting mathematical formulas
like
       \( x-3y + z = 7 \)
or
       \( a_{1} > x^{2n} + y^{2n} > x' \)
or
       \( \ip{A}{B} = \sum_{i} a_{i} b_{i} \).
The spaces you type in a formula are
ignored.  Remember that a letter like
       $x$                   % $ ... $  and  \( ... \)  are equivalent
is a formula when it denotes a mathematical
symbol, and it should be typed as one.

\section{Displayed Text}

Text is displayed by indenting it from the left
margin.  Quotations are commonly displayed.  There
are short quotations
\begin{quote}
   This is a short quotation.  It consists of a
   single paragraph of text.  See how it is formatted.
\end{quote}
and longer ones.
\begin{quotation}
   This is a longer quotation.  It consists of two
   paragraphs of text, neither of which are
   particularly interesting.

   This is the second paragraph of the quotation.  It
   is just as dull as the first paragraph.
\end{quotation}
Another frequently-displayed structure is a list.
The following is an example of an \emph{itemized}
list.
\begin{itemize}
   \item This is the first item of an itemized list.
         Each item in the list is marked with a ``tick''.
         You don't have to worry about what kind of tick
         mark is used.

   \item This is the second item of the list.  It
         contains another list nested inside it.  The inner
         list is an \emph{enumerated} list.
         \begin{enumerate}
            \item This is the first item of an enumerated
                  list that is nested within the itemized list.

            \item This is the second item of the inner list.
                  \LaTeX\ allows you to nest lists deeper than
                  you really should.
         \end{enumerate}
         This is the rest of the second item of the outer
         list.  It is no more interesting than any other
         part of the item.
   \item This is the third item of the list.
\end{itemize}
You can even display poetry.
\begin{verse}
   There is an environment
    for verse \\             % The \\ command separates lines
   Whose features some poets % within a stanza.
   will curse.

                             % One or more blank lines separate stanzas.

   For instead of making\\
   Them do \emph{all} line breaking, \\
   It allows them to put too many words on a line when they'd rather be
   forced to be terse.
\end{verse}

Mathematical formulas may also be displayed.  A
displayed formula
is
one-line long; multiline
formulas require special formatting instructions.
   \[  \ip{\Gamma}{\psi'} = x'' + y^{2} + z_{i}^{n}\]
Don't start a paragraph with a displayed equation,
nor make one a paragraph by itself.

\end{document}               % End of document.

\section{Bench Huge Section 89}
% synthetic bench section 89
% This is a sample LaTeX input file.  (Version of 12 August 2004.)
%
% A '%' character causes TeX to ignore all remaining text on the line,
% and is used for comments like this one.

\documentclass{article}      % Specifies the document class

                             % The preamble begins here.
\title{An Example Document}  % Declares the document's title.
\author{Leslie Lamport}      % Declares the author's name.
\date{January 21, 1994}      % Deleting this command produces today's date.

\newcommand{\ip}[2]{(#1, #2)}
                             % Defines \ip{arg1}{arg2} to mean
                             % (arg1, arg2).

%\newcommand{\ip}[2]{\langle #1 | #2\rangle}
                             % This is an alternative definition of
                             % \ip that is commented out.

\begin{document}             % End of preamble and beginning of text.

\maketitle                   % Produces the title.

This is an example input file.  Comparing it with
the output it generates can show you how to
produce a simple document of your own.

\section{Ordinary Text}      % Produces section heading.  Lower-level
                             % sections are begun with similar
                             % \subsection and \subsubsection commands.

The ends  of words and sentences are marked
  by   spaces. It  doesn't matter how many
spaces    you type; one is as good as 100.  The
end of   a line counts as a space.

One   or more   blank lines denote the  end
of  a paragraph.

Since any number of consecutive spaces are treated
like a single one, the formatting of the input
file makes no difference to
      \LaTeX,                % The \LaTeX command generates the LaTeX logo.
but it makes a difference to you.  When you use
\LaTeX, making your input file as easy to read
as possible will be a great help as you write
your document and when you change it.  This sample
file shows how you can add comments to your own input
file.

Because printing is different from typewriting,
there are a number of things that you have to do
differently when preparing an input file than if
you were just typing the document directly.
Quotation marks like
       ``this''
have to be handled specially, as do quotes within
quotes:
       ``\,`this'            % \, separates the double and single quote.
        is what I just
        wrote, not  `that'\,''.

Dashes come in three sizes: an
       intra-word
dash, a medium dash for number ranges like
       1--2,
and a punctuation
       dash---like
this.

A sentence-ending space should be larger than the
space between words within a sentence.  You
sometimes have to type special commands in
conjunction with punctuation characters to get
this right, as in the following sentence.
       Gnats, gnus, etc.\ all  % `\ ' makes an inter-word space.
       begin with G\@.         % \@ marks end-of-sentence punctuation.
You should check the spaces after periods when
reading your output to make sure you haven't
forgotten any special cases.  Generating an
ellipsis
       \ldots\               % `\ ' is needed after `\ldots' because TeX
                             % ignores spaces after command names like \ldots
                             % made from \ + letters.
                             %
                             % Note how a `%' character causes TeX to ignore
                             % the end of the input line, so these blank lines
                             % do not start a new paragraph.
                             %
with the right spacing around the periods requires
a special command.

\LaTeX\ interprets some common characters as
commands, so you must type special commands to
generate them.  These characters include the
following:
       \$ \& \% \# \{ and \}.

In printing, text is usually emphasized with an
       \emph{italic}
type style.

\begin{em}
   A long segment of text can also be emphasized
   in this way.  Text within such a segment can be
   given \emph{additional} emphasis.
\end{em}

It is sometimes necessary to prevent \LaTeX\ from
breaking a line where it might otherwise do so.
This may be at a space, as between the ``Mr.''\ and
``Jones'' in
       ``Mr.~Jones'',        % ~ produces an unbreakable interword space.
or within a word---especially when the word is a
symbol like
       \mbox{\emph{itemnum}}
that makes little sense when hyphenated across
lines.

Footnotes\footnote{This is an example of a footnote.}
pose no problem.

\LaTeX\ is good at typesetting mathematical formulas
like
       \( x-3y + z = 7 \)
or
       \( a_{1} > x^{2n} + y^{2n} > x' \)
or
       \( \ip{A}{B} = \sum_{i} a_{i} b_{i} \).
The spaces you type in a formula are
ignored.  Remember that a letter like
       $x$                   % $ ... $  and  \( ... \)  are equivalent
is a formula when it denotes a mathematical
symbol, and it should be typed as one.

\section{Displayed Text}

Text is displayed by indenting it from the left
margin.  Quotations are commonly displayed.  There
are short quotations
\begin{quote}
   This is a short quotation.  It consists of a
   single paragraph of text.  See how it is formatted.
\end{quote}
and longer ones.
\begin{quotation}
   This is a longer quotation.  It consists of two
   paragraphs of text, neither of which are
   particularly interesting.

   This is the second paragraph of the quotation.  It
   is just as dull as the first paragraph.
\end{quotation}
Another frequently-displayed structure is a list.
The following is an example of an \emph{itemized}
list.
\begin{itemize}
   \item This is the first item of an itemized list.
         Each item in the list is marked with a ``tick''.
         You don't have to worry about what kind of tick
         mark is used.

   \item This is the second item of the list.  It
         contains another list nested inside it.  The inner
         list is an \emph{enumerated} list.
         \begin{enumerate}
            \item This is the first item of an enumerated
                  list that is nested within the itemized list.

            \item This is the second item of the inner list.
                  \LaTeX\ allows you to nest lists deeper than
                  you really should.
         \end{enumerate}
         This is the rest of the second item of the outer
         list.  It is no more interesting than any other
         part of the item.
   \item This is the third item of the list.
\end{itemize}
You can even display poetry.
\begin{verse}
   There is an environment
    for verse \\             % The \\ command separates lines
   Whose features some poets % within a stanza.
   will curse.

                             % One or more blank lines separate stanzas.

   For instead of making\\
   Them do \emph{all} line breaking, \\
   It allows them to put too many words on a line when they'd rather be
   forced to be terse.
\end{verse}

Mathematical formulas may also be displayed.  A
displayed formula
is
one-line long; multiline
formulas require special formatting instructions.
   \[  \ip{\Gamma}{\psi'} = x'' + y^{2} + z_{i}^{n}\]
Don't start a paragraph with a displayed equation,
nor make one a paragraph by itself.

\end{document}               % End of document.

\section{Bench Huge Section 90}
% synthetic bench section 90
% This is a sample LaTeX input file.  (Version of 12 August 2004.)
%
% A '%' character causes TeX to ignore all remaining text on the line,
% and is used for comments like this one.

\documentclass{article}      % Specifies the document class

                             % The preamble begins here.
\title{An Example Document}  % Declares the document's title.
\author{Leslie Lamport}      % Declares the author's name.
\date{January 21, 1994}      % Deleting this command produces today's date.

\newcommand{\ip}[2]{(#1, #2)}
                             % Defines \ip{arg1}{arg2} to mean
                             % (arg1, arg2).

%\newcommand{\ip}[2]{\langle #1 | #2\rangle}
                             % This is an alternative definition of
                             % \ip that is commented out.

\begin{document}             % End of preamble and beginning of text.

\maketitle                   % Produces the title.

This is an example input file.  Comparing it with
the output it generates can show you how to
produce a simple document of your own.

\section{Ordinary Text}      % Produces section heading.  Lower-level
                             % sections are begun with similar
                             % \subsection and \subsubsection commands.

The ends  of words and sentences are marked
  by   spaces. It  doesn't matter how many
spaces    you type; one is as good as 100.  The
end of   a line counts as a space.

One   or more   blank lines denote the  end
of  a paragraph.

Since any number of consecutive spaces are treated
like a single one, the formatting of the input
file makes no difference to
      \LaTeX,                % The \LaTeX command generates the LaTeX logo.
but it makes a difference to you.  When you use
\LaTeX, making your input file as easy to read
as possible will be a great help as you write
your document and when you change it.  This sample
file shows how you can add comments to your own input
file.

Because printing is different from typewriting,
there are a number of things that you have to do
differently when preparing an input file than if
you were just typing the document directly.
Quotation marks like
       ``this''
have to be handled specially, as do quotes within
quotes:
       ``\,`this'            % \, separates the double and single quote.
        is what I just
        wrote, not  `that'\,''.

Dashes come in three sizes: an
       intra-word
dash, a medium dash for number ranges like
       1--2,
and a punctuation
       dash---like
this.

A sentence-ending space should be larger than the
space between words within a sentence.  You
sometimes have to type special commands in
conjunction with punctuation characters to get
this right, as in the following sentence.
       Gnats, gnus, etc.\ all  % `\ ' makes an inter-word space.
       begin with G\@.         % \@ marks end-of-sentence punctuation.
You should check the spaces after periods when
reading your output to make sure you haven't
forgotten any special cases.  Generating an
ellipsis
       \ldots\               % `\ ' is needed after `\ldots' because TeX
                             % ignores spaces after command names like \ldots
                             % made from \ + letters.
                             %
                             % Note how a `%' character causes TeX to ignore
                             % the end of the input line, so these blank lines
                             % do not start a new paragraph.
                             %
with the right spacing around the periods requires
a special command.

\LaTeX\ interprets some common characters as
commands, so you must type special commands to
generate them.  These characters include the
following:
       \$ \& \% \# \{ and \}.

In printing, text is usually emphasized with an
       \emph{italic}
type style.

\begin{em}
   A long segment of text can also be emphasized
   in this way.  Text within such a segment can be
   given \emph{additional} emphasis.
\end{em}

It is sometimes necessary to prevent \LaTeX\ from
breaking a line where it might otherwise do so.
This may be at a space, as between the ``Mr.''\ and
``Jones'' in
       ``Mr.~Jones'',        % ~ produces an unbreakable interword space.
or within a word---especially when the word is a
symbol like
       \mbox{\emph{itemnum}}
that makes little sense when hyphenated across
lines.

Footnotes\footnote{This is an example of a footnote.}
pose no problem.

\LaTeX\ is good at typesetting mathematical formulas
like
       \( x-3y + z = 7 \)
or
       \( a_{1} > x^{2n} + y^{2n} > x' \)
or
       \( \ip{A}{B} = \sum_{i} a_{i} b_{i} \).
The spaces you type in a formula are
ignored.  Remember that a letter like
       $x$                   % $ ... $  and  \( ... \)  are equivalent
is a formula when it denotes a mathematical
symbol, and it should be typed as one.

\section{Displayed Text}

Text is displayed by indenting it from the left
margin.  Quotations are commonly displayed.  There
are short quotations
\begin{quote}
   This is a short quotation.  It consists of a
   single paragraph of text.  See how it is formatted.
\end{quote}
and longer ones.
\begin{quotation}
   This is a longer quotation.  It consists of two
   paragraphs of text, neither of which are
   particularly interesting.

   This is the second paragraph of the quotation.  It
   is just as dull as the first paragraph.
\end{quotation}
Another frequently-displayed structure is a list.
The following is an example of an \emph{itemized}
list.
\begin{itemize}
   \item This is the first item of an itemized list.
         Each item in the list is marked with a ``tick''.
         You don't have to worry about what kind of tick
         mark is used.

   \item This is the second item of the list.  It
         contains another list nested inside it.  The inner
         list is an \emph{enumerated} list.
         \begin{enumerate}
            \item This is the first item of an enumerated
                  list that is nested within the itemized list.

            \item This is the second item of the inner list.
                  \LaTeX\ allows you to nest lists deeper than
                  you really should.
         \end{enumerate}
         This is the rest of the second item of the outer
         list.  It is no more interesting than any other
         part of the item.
   \item This is the third item of the list.
\end{itemize}
You can even display poetry.
\begin{verse}
   There is an environment
    for verse \\             % The \\ command separates lines
   Whose features some poets % within a stanza.
   will curse.

                             % One or more blank lines separate stanzas.

   For instead of making\\
   Them do \emph{all} line breaking, \\
   It allows them to put too many words on a line when they'd rather be
   forced to be terse.
\end{verse}

Mathematical formulas may also be displayed.  A
displayed formula
is
one-line long; multiline
formulas require special formatting instructions.
   \[  \ip{\Gamma}{\psi'} = x'' + y^{2} + z_{i}^{n}\]
Don't start a paragraph with a displayed equation,
nor make one a paragraph by itself.

\end{document}               % End of document.

\section{Bench Huge Section 91}
% synthetic bench section 91
% This is a sample LaTeX input file.  (Version of 12 August 2004.)
%
% A '%' character causes TeX to ignore all remaining text on the line,
% and is used for comments like this one.

\documentclass{article}      % Specifies the document class

                             % The preamble begins here.
\title{An Example Document}  % Declares the document's title.
\author{Leslie Lamport}      % Declares the author's name.
\date{January 21, 1994}      % Deleting this command produces today's date.

\newcommand{\ip}[2]{(#1, #2)}
                             % Defines \ip{arg1}{arg2} to mean
                             % (arg1, arg2).

%\newcommand{\ip}[2]{\langle #1 | #2\rangle}
                             % This is an alternative definition of
                             % \ip that is commented out.

\begin{document}             % End of preamble and beginning of text.

\maketitle                   % Produces the title.

This is an example input file.  Comparing it with
the output it generates can show you how to
produce a simple document of your own.

\section{Ordinary Text}      % Produces section heading.  Lower-level
                             % sections are begun with similar
                             % \subsection and \subsubsection commands.

The ends  of words and sentences are marked
  by   spaces. It  doesn't matter how many
spaces    you type; one is as good as 100.  The
end of   a line counts as a space.

One   or more   blank lines denote the  end
of  a paragraph.

Since any number of consecutive spaces are treated
like a single one, the formatting of the input
file makes no difference to
      \LaTeX,                % The \LaTeX command generates the LaTeX logo.
but it makes a difference to you.  When you use
\LaTeX, making your input file as easy to read
as possible will be a great help as you write
your document and when you change it.  This sample
file shows how you can add comments to your own input
file.

Because printing is different from typewriting,
there are a number of things that you have to do
differently when preparing an input file than if
you were just typing the document directly.
Quotation marks like
       ``this''
have to be handled specially, as do quotes within
quotes:
       ``\,`this'            % \, separates the double and single quote.
        is what I just
        wrote, not  `that'\,''.

Dashes come in three sizes: an
       intra-word
dash, a medium dash for number ranges like
       1--2,
and a punctuation
       dash---like
this.

A sentence-ending space should be larger than the
space between words within a sentence.  You
sometimes have to type special commands in
conjunction with punctuation characters to get
this right, as in the following sentence.
       Gnats, gnus, etc.\ all  % `\ ' makes an inter-word space.
       begin with G\@.         % \@ marks end-of-sentence punctuation.
You should check the spaces after periods when
reading your output to make sure you haven't
forgotten any special cases.  Generating an
ellipsis
       \ldots\               % `\ ' is needed after `\ldots' because TeX
                             % ignores spaces after command names like \ldots
                             % made from \ + letters.
                             %
                             % Note how a `%' character causes TeX to ignore
                             % the end of the input line, so these blank lines
                             % do not start a new paragraph.
                             %
with the right spacing around the periods requires
a special command.

\LaTeX\ interprets some common characters as
commands, so you must type special commands to
generate them.  These characters include the
following:
       \$ \& \% \# \{ and \}.

In printing, text is usually emphasized with an
       \emph{italic}
type style.

\begin{em}
   A long segment of text can also be emphasized
   in this way.  Text within such a segment can be
   given \emph{additional} emphasis.
\end{em}

It is sometimes necessary to prevent \LaTeX\ from
breaking a line where it might otherwise do so.
This may be at a space, as between the ``Mr.''\ and
``Jones'' in
       ``Mr.~Jones'',        % ~ produces an unbreakable interword space.
or within a word---especially when the word is a
symbol like
       \mbox{\emph{itemnum}}
that makes little sense when hyphenated across
lines.

Footnotes\footnote{This is an example of a footnote.}
pose no problem.

\LaTeX\ is good at typesetting mathematical formulas
like
       \( x-3y + z = 7 \)
or
       \( a_{1} > x^{2n} + y^{2n} > x' \)
or
       \( \ip{A}{B} = \sum_{i} a_{i} b_{i} \).
The spaces you type in a formula are
ignored.  Remember that a letter like
       $x$                   % $ ... $  and  \( ... \)  are equivalent
is a formula when it denotes a mathematical
symbol, and it should be typed as one.

\section{Displayed Text}

Text is displayed by indenting it from the left
margin.  Quotations are commonly displayed.  There
are short quotations
\begin{quote}
   This is a short quotation.  It consists of a
   single paragraph of text.  See how it is formatted.
\end{quote}
and longer ones.
\begin{quotation}
   This is a longer quotation.  It consists of two
   paragraphs of text, neither of which are
   particularly interesting.

   This is the second paragraph of the quotation.  It
   is just as dull as the first paragraph.
\end{quotation}
Another frequently-displayed structure is a list.
The following is an example of an \emph{itemized}
list.
\begin{itemize}
   \item This is the first item of an itemized list.
         Each item in the list is marked with a ``tick''.
         You don't have to worry about what kind of tick
         mark is used.

   \item This is the second item of the list.  It
         contains another list nested inside it.  The inner
         list is an \emph{enumerated} list.
         \begin{enumerate}
            \item This is the first item of an enumerated
                  list that is nested within the itemized list.

            \item This is the second item of the inner list.
                  \LaTeX\ allows you to nest lists deeper than
                  you really should.
         \end{enumerate}
         This is the rest of the second item of the outer
         list.  It is no more interesting than any other
         part of the item.
   \item This is the third item of the list.
\end{itemize}
You can even display poetry.
\begin{verse}
   There is an environment
    for verse \\             % The \\ command separates lines
   Whose features some poets % within a stanza.
   will curse.

                             % One or more blank lines separate stanzas.

   For instead of making\\
   Them do \emph{all} line breaking, \\
   It allows them to put too many words on a line when they'd rather be
   forced to be terse.
\end{verse}

Mathematical formulas may also be displayed.  A
displayed formula
is
one-line long; multiline
formulas require special formatting instructions.
   \[  \ip{\Gamma}{\psi'} = x'' + y^{2} + z_{i}^{n}\]
Don't start a paragraph with a displayed equation,
nor make one a paragraph by itself.

\end{document}               % End of document.

\section{Bench Huge Section 92}
% synthetic bench section 92
% This is a sample LaTeX input file.  (Version of 12 August 2004.)
%
% A '%' character causes TeX to ignore all remaining text on the line,
% and is used for comments like this one.

\documentclass{article}      % Specifies the document class

                             % The preamble begins here.
\title{An Example Document}  % Declares the document's title.
\author{Leslie Lamport}      % Declares the author's name.
\date{January 21, 1994}      % Deleting this command produces today's date.

\newcommand{\ip}[2]{(#1, #2)}
                             % Defines \ip{arg1}{arg2} to mean
                             % (arg1, arg2).

%\newcommand{\ip}[2]{\langle #1 | #2\rangle}
                             % This is an alternative definition of
                             % \ip that is commented out.

\begin{document}             % End of preamble and beginning of text.

\maketitle                   % Produces the title.

This is an example input file.  Comparing it with
the output it generates can show you how to
produce a simple document of your own.

\section{Ordinary Text}      % Produces section heading.  Lower-level
                             % sections are begun with similar
                             % \subsection and \subsubsection commands.

The ends  of words and sentences are marked
  by   spaces. It  doesn't matter how many
spaces    you type; one is as good as 100.  The
end of   a line counts as a space.

One   or more   blank lines denote the  end
of  a paragraph.

Since any number of consecutive spaces are treated
like a single one, the formatting of the input
file makes no difference to
      \LaTeX,                % The \LaTeX command generates the LaTeX logo.
but it makes a difference to you.  When you use
\LaTeX, making your input file as easy to read
as possible will be a great help as you write
your document and when you change it.  This sample
file shows how you can add comments to your own input
file.

Because printing is different from typewriting,
there are a number of things that you have to do
differently when preparing an input file than if
you were just typing the document directly.
Quotation marks like
       ``this''
have to be handled specially, as do quotes within
quotes:
       ``\,`this'            % \, separates the double and single quote.
        is what I just
        wrote, not  `that'\,''.

Dashes come in three sizes: an
       intra-word
dash, a medium dash for number ranges like
       1--2,
and a punctuation
       dash---like
this.

A sentence-ending space should be larger than the
space between words within a sentence.  You
sometimes have to type special commands in
conjunction with punctuation characters to get
this right, as in the following sentence.
       Gnats, gnus, etc.\ all  % `\ ' makes an inter-word space.
       begin with G\@.         % \@ marks end-of-sentence punctuation.
You should check the spaces after periods when
reading your output to make sure you haven't
forgotten any special cases.  Generating an
ellipsis
       \ldots\               % `\ ' is needed after `\ldots' because TeX
                             % ignores spaces after command names like \ldots
                             % made from \ + letters.
                             %
                             % Note how a `%' character causes TeX to ignore
                             % the end of the input line, so these blank lines
                             % do not start a new paragraph.
                             %
with the right spacing around the periods requires
a special command.

\LaTeX\ interprets some common characters as
commands, so you must type special commands to
generate them.  These characters include the
following:
       \$ \& \% \# \{ and \}.

In printing, text is usually emphasized with an
       \emph{italic}
type style.

\begin{em}
   A long segment of text can also be emphasized
   in this way.  Text within such a segment can be
   given \emph{additional} emphasis.
\end{em}

It is sometimes necessary to prevent \LaTeX\ from
breaking a line where it might otherwise do so.
This may be at a space, as between the ``Mr.''\ and
``Jones'' in
       ``Mr.~Jones'',        % ~ produces an unbreakable interword space.
or within a word---especially when the word is a
symbol like
       \mbox{\emph{itemnum}}
that makes little sense when hyphenated across
lines.

Footnotes\footnote{This is an example of a footnote.}
pose no problem.

\LaTeX\ is good at typesetting mathematical formulas
like
       \( x-3y + z = 7 \)
or
       \( a_{1} > x^{2n} + y^{2n} > x' \)
or
       \( \ip{A}{B} = \sum_{i} a_{i} b_{i} \).
The spaces you type in a formula are
ignored.  Remember that a letter like
       $x$                   % $ ... $  and  \( ... \)  are equivalent
is a formula when it denotes a mathematical
symbol, and it should be typed as one.

\section{Displayed Text}

Text is displayed by indenting it from the left
margin.  Quotations are commonly displayed.  There
are short quotations
\begin{quote}
   This is a short quotation.  It consists of a
   single paragraph of text.  See how it is formatted.
\end{quote}
and longer ones.
\begin{quotation}
   This is a longer quotation.  It consists of two
   paragraphs of text, neither of which are
   particularly interesting.

   This is the second paragraph of the quotation.  It
   is just as dull as the first paragraph.
\end{quotation}
Another frequently-displayed structure is a list.
The following is an example of an \emph{itemized}
list.
\begin{itemize}
   \item This is the first item of an itemized list.
         Each item in the list is marked with a ``tick''.
         You don't have to worry about what kind of tick
         mark is used.

   \item This is the second item of the list.  It
         contains another list nested inside it.  The inner
         list is an \emph{enumerated} list.
         \begin{enumerate}
            \item This is the first item of an enumerated
                  list that is nested within the itemized list.

            \item This is the second item of the inner list.
                  \LaTeX\ allows you to nest lists deeper than
                  you really should.
         \end{enumerate}
         This is the rest of the second item of the outer
         list.  It is no more interesting than any other
         part of the item.
   \item This is the third item of the list.
\end{itemize}
You can even display poetry.
\begin{verse}
   There is an environment
    for verse \\             % The \\ command separates lines
   Whose features some poets % within a stanza.
   will curse.

                             % One or more blank lines separate stanzas.

   For instead of making\\
   Them do \emph{all} line breaking, \\
   It allows them to put too many words on a line when they'd rather be
   forced to be terse.
\end{verse}

Mathematical formulas may also be displayed.  A
displayed formula
is
one-line long; multiline
formulas require special formatting instructions.
   \[  \ip{\Gamma}{\psi'} = x'' + y^{2} + z_{i}^{n}\]
Don't start a paragraph with a displayed equation,
nor make one a paragraph by itself.

\end{document}               % End of document.

\section{Bench Huge Section 93}
% synthetic bench section 93
% This is a sample LaTeX input file.  (Version of 12 August 2004.)
%
% A '%' character causes TeX to ignore all remaining text on the line,
% and is used for comments like this one.

\documentclass{article}      % Specifies the document class

                             % The preamble begins here.
\title{An Example Document}  % Declares the document's title.
\author{Leslie Lamport}      % Declares the author's name.
\date{January 21, 1994}      % Deleting this command produces today's date.

\newcommand{\ip}[2]{(#1, #2)}
                             % Defines \ip{arg1}{arg2} to mean
                             % (arg1, arg2).

%\newcommand{\ip}[2]{\langle #1 | #2\rangle}
                             % This is an alternative definition of
                             % \ip that is commented out.

\begin{document}             % End of preamble and beginning of text.

\maketitle                   % Produces the title.

This is an example input file.  Comparing it with
the output it generates can show you how to
produce a simple document of your own.

\section{Ordinary Text}      % Produces section heading.  Lower-level
                             % sections are begun with similar
                             % \subsection and \subsubsection commands.

The ends  of words and sentences are marked
  by   spaces. It  doesn't matter how many
spaces    you type; one is as good as 100.  The
end of   a line counts as a space.

One   or more   blank lines denote the  end
of  a paragraph.

Since any number of consecutive spaces are treated
like a single one, the formatting of the input
file makes no difference to
      \LaTeX,                % The \LaTeX command generates the LaTeX logo.
but it makes a difference to you.  When you use
\LaTeX, making your input file as easy to read
as possible will be a great help as you write
your document and when you change it.  This sample
file shows how you can add comments to your own input
file.

Because printing is different from typewriting,
there are a number of things that you have to do
differently when preparing an input file than if
you were just typing the document directly.
Quotation marks like
       ``this''
have to be handled specially, as do quotes within
quotes:
       ``\,`this'            % \, separates the double and single quote.
        is what I just
        wrote, not  `that'\,''.

Dashes come in three sizes: an
       intra-word
dash, a medium dash for number ranges like
       1--2,
and a punctuation
       dash---like
this.

A sentence-ending space should be larger than the
space between words within a sentence.  You
sometimes have to type special commands in
conjunction with punctuation characters to get
this right, as in the following sentence.
       Gnats, gnus, etc.\ all  % `\ ' makes an inter-word space.
       begin with G\@.         % \@ marks end-of-sentence punctuation.
You should check the spaces after periods when
reading your output to make sure you haven't
forgotten any special cases.  Generating an
ellipsis
       \ldots\               % `\ ' is needed after `\ldots' because TeX
                             % ignores spaces after command names like \ldots
                             % made from \ + letters.
                             %
                             % Note how a `%' character causes TeX to ignore
                             % the end of the input line, so these blank lines
                             % do not start a new paragraph.
                             %
with the right spacing around the periods requires
a special command.

\LaTeX\ interprets some common characters as
commands, so you must type special commands to
generate them.  These characters include the
following:
       \$ \& \% \# \{ and \}.

In printing, text is usually emphasized with an
       \emph{italic}
type style.

\begin{em}
   A long segment of text can also be emphasized
   in this way.  Text within such a segment can be
   given \emph{additional} emphasis.
\end{em}

It is sometimes necessary to prevent \LaTeX\ from
breaking a line where it might otherwise do so.
This may be at a space, as between the ``Mr.''\ and
``Jones'' in
       ``Mr.~Jones'',        % ~ produces an unbreakable interword space.
or within a word---especially when the word is a
symbol like
       \mbox{\emph{itemnum}}
that makes little sense when hyphenated across
lines.

Footnotes\footnote{This is an example of a footnote.}
pose no problem.

\LaTeX\ is good at typesetting mathematical formulas
like
       \( x-3y + z = 7 \)
or
       \( a_{1} > x^{2n} + y^{2n} > x' \)
or
       \( \ip{A}{B} = \sum_{i} a_{i} b_{i} \).
The spaces you type in a formula are
ignored.  Remember that a letter like
       $x$                   % $ ... $  and  \( ... \)  are equivalent
is a formula when it denotes a mathematical
symbol, and it should be typed as one.

\section{Displayed Text}

Text is displayed by indenting it from the left
margin.  Quotations are commonly displayed.  There
are short quotations
\begin{quote}
   This is a short quotation.  It consists of a
   single paragraph of text.  See how it is formatted.
\end{quote}
and longer ones.
\begin{quotation}
   This is a longer quotation.  It consists of two
   paragraphs of text, neither of which are
   particularly interesting.

   This is the second paragraph of the quotation.  It
   is just as dull as the first paragraph.
\end{quotation}
Another frequently-displayed structure is a list.
The following is an example of an \emph{itemized}
list.
\begin{itemize}
   \item This is the first item of an itemized list.
         Each item in the list is marked with a ``tick''.
         You don't have to worry about what kind of tick
         mark is used.

   \item This is the second item of the list.  It
         contains another list nested inside it.  The inner
         list is an \emph{enumerated} list.
         \begin{enumerate}
            \item This is the first item of an enumerated
                  list that is nested within the itemized list.

            \item This is the second item of the inner list.
                  \LaTeX\ allows you to nest lists deeper than
                  you really should.
         \end{enumerate}
         This is the rest of the second item of the outer
         list.  It is no more interesting than any other
         part of the item.
   \item This is the third item of the list.
\end{itemize}
You can even display poetry.
\begin{verse}
   There is an environment
    for verse \\             % The \\ command separates lines
   Whose features some poets % within a stanza.
   will curse.

                             % One or more blank lines separate stanzas.

   For instead of making\\
   Them do \emph{all} line breaking, \\
   It allows them to put too many words on a line when they'd rather be
   forced to be terse.
\end{verse}

Mathematical formulas may also be displayed.  A
displayed formula
is
one-line long; multiline
formulas require special formatting instructions.
   \[  \ip{\Gamma}{\psi'} = x'' + y^{2} + z_{i}^{n}\]
Don't start a paragraph with a displayed equation,
nor make one a paragraph by itself.

\end{document}               % End of document.

\section{Bench Huge Section 94}
% synthetic bench section 94
% This is a sample LaTeX input file.  (Version of 12 August 2004.)
%
% A '%' character causes TeX to ignore all remaining text on the line,
% and is used for comments like this one.

\documentclass{article}      % Specifies the document class

                             % The preamble begins here.
\title{An Example Document}  % Declares the document's title.
\author{Leslie Lamport}      % Declares the author's name.
\date{January 21, 1994}      % Deleting this command produces today's date.

\newcommand{\ip}[2]{(#1, #2)}
                             % Defines \ip{arg1}{arg2} to mean
                             % (arg1, arg2).

%\newcommand{\ip}[2]{\langle #1 | #2\rangle}
                             % This is an alternative definition of
                             % \ip that is commented out.

\begin{document}             % End of preamble and beginning of text.

\maketitle                   % Produces the title.

This is an example input file.  Comparing it with
the output it generates can show you how to
produce a simple document of your own.

\section{Ordinary Text}      % Produces section heading.  Lower-level
                             % sections are begun with similar
                             % \subsection and \subsubsection commands.

The ends  of words and sentences are marked
  by   spaces. It  doesn't matter how many
spaces    you type; one is as good as 100.  The
end of   a line counts as a space.

One   or more   blank lines denote the  end
of  a paragraph.

Since any number of consecutive spaces are treated
like a single one, the formatting of the input
file makes no difference to
      \LaTeX,                % The \LaTeX command generates the LaTeX logo.
but it makes a difference to you.  When you use
\LaTeX, making your input file as easy to read
as possible will be a great help as you write
your document and when you change it.  This sample
file shows how you can add comments to your own input
file.

Because printing is different from typewriting,
there are a number of things that you have to do
differently when preparing an input file than if
you were just typing the document directly.
Quotation marks like
       ``this''
have to be handled specially, as do quotes within
quotes:
       ``\,`this'            % \, separates the double and single quote.
        is what I just
        wrote, not  `that'\,''.

Dashes come in three sizes: an
       intra-word
dash, a medium dash for number ranges like
       1--2,
and a punctuation
       dash---like
this.

A sentence-ending space should be larger than the
space between words within a sentence.  You
sometimes have to type special commands in
conjunction with punctuation characters to get
this right, as in the following sentence.
       Gnats, gnus, etc.\ all  % `\ ' makes an inter-word space.
       begin with G\@.         % \@ marks end-of-sentence punctuation.
You should check the spaces after periods when
reading your output to make sure you haven't
forgotten any special cases.  Generating an
ellipsis
       \ldots\               % `\ ' is needed after `\ldots' because TeX
                             % ignores spaces after command names like \ldots
                             % made from \ + letters.
                             %
                             % Note how a `%' character causes TeX to ignore
                             % the end of the input line, so these blank lines
                             % do not start a new paragraph.
                             %
with the right spacing around the periods requires
a special command.

\LaTeX\ interprets some common characters as
commands, so you must type special commands to
generate them.  These characters include the
following:
       \$ \& \% \# \{ and \}.

In printing, text is usually emphasized with an
       \emph{italic}
type style.

\begin{em}
   A long segment of text can also be emphasized
   in this way.  Text within such a segment can be
   given \emph{additional} emphasis.
\end{em}

It is sometimes necessary to prevent \LaTeX\ from
breaking a line where it might otherwise do so.
This may be at a space, as between the ``Mr.''\ and
``Jones'' in
       ``Mr.~Jones'',        % ~ produces an unbreakable interword space.
or within a word---especially when the word is a
symbol like
       \mbox{\emph{itemnum}}
that makes little sense when hyphenated across
lines.

Footnotes\footnote{This is an example of a footnote.}
pose no problem.

\LaTeX\ is good at typesetting mathematical formulas
like
       \( x-3y + z = 7 \)
or
       \( a_{1} > x^{2n} + y^{2n} > x' \)
or
       \( \ip{A}{B} = \sum_{i} a_{i} b_{i} \).
The spaces you type in a formula are
ignored.  Remember that a letter like
       $x$                   % $ ... $  and  \( ... \)  are equivalent
is a formula when it denotes a mathematical
symbol, and it should be typed as one.

\section{Displayed Text}

Text is displayed by indenting it from the left
margin.  Quotations are commonly displayed.  There
are short quotations
\begin{quote}
   This is a short quotation.  It consists of a
   single paragraph of text.  See how it is formatted.
\end{quote}
and longer ones.
\begin{quotation}
   This is a longer quotation.  It consists of two
   paragraphs of text, neither of which are
   particularly interesting.

   This is the second paragraph of the quotation.  It
   is just as dull as the first paragraph.
\end{quotation}
Another frequently-displayed structure is a list.
The following is an example of an \emph{itemized}
list.
\begin{itemize}
   \item This is the first item of an itemized list.
         Each item in the list is marked with a ``tick''.
         You don't have to worry about what kind of tick
         mark is used.

   \item This is the second item of the list.  It
         contains another list nested inside it.  The inner
         list is an \emph{enumerated} list.
         \begin{enumerate}
            \item This is the first item of an enumerated
                  list that is nested within the itemized list.

            \item This is the second item of the inner list.
                  \LaTeX\ allows you to nest lists deeper than
                  you really should.
         \end{enumerate}
         This is the rest of the second item of the outer
         list.  It is no more interesting than any other
         part of the item.
   \item This is the third item of the list.
\end{itemize}
You can even display poetry.
\begin{verse}
   There is an environment
    for verse \\             % The \\ command separates lines
   Whose features some poets % within a stanza.
   will curse.

                             % One or more blank lines separate stanzas.

   For instead of making\\
   Them do \emph{all} line breaking, \\
   It allows them to put too many words on a line when they'd rather be
   forced to be terse.
\end{verse}

Mathematical formulas may also be displayed.  A
displayed formula
is
one-line long; multiline
formulas require special formatting instructions.
   \[  \ip{\Gamma}{\psi'} = x'' + y^{2} + z_{i}^{n}\]
Don't start a paragraph with a displayed equation,
nor make one a paragraph by itself.

\end{document}               % End of document.

\section{Bench Huge Section 95}
% synthetic bench section 95
% This is a sample LaTeX input file.  (Version of 12 August 2004.)
%
% A '%' character causes TeX to ignore all remaining text on the line,
% and is used for comments like this one.

\documentclass{article}      % Specifies the document class

                             % The preamble begins here.
\title{An Example Document}  % Declares the document's title.
\author{Leslie Lamport}      % Declares the author's name.
\date{January 21, 1994}      % Deleting this command produces today's date.

\newcommand{\ip}[2]{(#1, #2)}
                             % Defines \ip{arg1}{arg2} to mean
                             % (arg1, arg2).

%\newcommand{\ip}[2]{\langle #1 | #2\rangle}
                             % This is an alternative definition of
                             % \ip that is commented out.

\begin{document}             % End of preamble and beginning of text.

\maketitle                   % Produces the title.

This is an example input file.  Comparing it with
the output it generates can show you how to
produce a simple document of your own.

\section{Ordinary Text}      % Produces section heading.  Lower-level
                             % sections are begun with similar
                             % \subsection and \subsubsection commands.

The ends  of words and sentences are marked
  by   spaces. It  doesn't matter how many
spaces    you type; one is as good as 100.  The
end of   a line counts as a space.

One   or more   blank lines denote the  end
of  a paragraph.

Since any number of consecutive spaces are treated
like a single one, the formatting of the input
file makes no difference to
      \LaTeX,                % The \LaTeX command generates the LaTeX logo.
but it makes a difference to you.  When you use
\LaTeX, making your input file as easy to read
as possible will be a great help as you write
your document and when you change it.  This sample
file shows how you can add comments to your own input
file.

Because printing is different from typewriting,
there are a number of things that you have to do
differently when preparing an input file than if
you were just typing the document directly.
Quotation marks like
       ``this''
have to be handled specially, as do quotes within
quotes:
       ``\,`this'            % \, separates the double and single quote.
        is what I just
        wrote, not  `that'\,''.

Dashes come in three sizes: an
       intra-word
dash, a medium dash for number ranges like
       1--2,
and a punctuation
       dash---like
this.

A sentence-ending space should be larger than the
space between words within a sentence.  You
sometimes have to type special commands in
conjunction with punctuation characters to get
this right, as in the following sentence.
       Gnats, gnus, etc.\ all  % `\ ' makes an inter-word space.
       begin with G\@.         % \@ marks end-of-sentence punctuation.
You should check the spaces after periods when
reading your output to make sure you haven't
forgotten any special cases.  Generating an
ellipsis
       \ldots\               % `\ ' is needed after `\ldots' because TeX
                             % ignores spaces after command names like \ldots
                             % made from \ + letters.
                             %
                             % Note how a `%' character causes TeX to ignore
                             % the end of the input line, so these blank lines
                             % do not start a new paragraph.
                             %
with the right spacing around the periods requires
a special command.

\LaTeX\ interprets some common characters as
commands, so you must type special commands to
generate them.  These characters include the
following:
       \$ \& \% \# \{ and \}.

In printing, text is usually emphasized with an
       \emph{italic}
type style.

\begin{em}
   A long segment of text can also be emphasized
   in this way.  Text within such a segment can be
   given \emph{additional} emphasis.
\end{em}

It is sometimes necessary to prevent \LaTeX\ from
breaking a line where it might otherwise do so.
This may be at a space, as between the ``Mr.''\ and
``Jones'' in
       ``Mr.~Jones'',        % ~ produces an unbreakable interword space.
or within a word---especially when the word is a
symbol like
       \mbox{\emph{itemnum}}
that makes little sense when hyphenated across
lines.

Footnotes\footnote{This is an example of a footnote.}
pose no problem.

\LaTeX\ is good at typesetting mathematical formulas
like
       \( x-3y + z = 7 \)
or
       \( a_{1} > x^{2n} + y^{2n} > x' \)
or
       \( \ip{A}{B} = \sum_{i} a_{i} b_{i} \).
The spaces you type in a formula are
ignored.  Remember that a letter like
       $x$                   % $ ... $  and  \( ... \)  are equivalent
is a formula when it denotes a mathematical
symbol, and it should be typed as one.

\section{Displayed Text}

Text is displayed by indenting it from the left
margin.  Quotations are commonly displayed.  There
are short quotations
\begin{quote}
   This is a short quotation.  It consists of a
   single paragraph of text.  See how it is formatted.
\end{quote}
and longer ones.
\begin{quotation}
   This is a longer quotation.  It consists of two
   paragraphs of text, neither of which are
   particularly interesting.

   This is the second paragraph of the quotation.  It
   is just as dull as the first paragraph.
\end{quotation}
Another frequently-displayed structure is a list.
The following is an example of an \emph{itemized}
list.
\begin{itemize}
   \item This is the first item of an itemized list.
         Each item in the list is marked with a ``tick''.
         You don't have to worry about what kind of tick
         mark is used.

   \item This is the second item of the list.  It
         contains another list nested inside it.  The inner
         list is an \emph{enumerated} list.
         \begin{enumerate}
            \item This is the first item of an enumerated
                  list that is nested within the itemized list.

            \item This is the second item of the inner list.
                  \LaTeX\ allows you to nest lists deeper than
                  you really should.
         \end{enumerate}
         This is the rest of the second item of the outer
         list.  It is no more interesting than any other
         part of the item.
   \item This is the third item of the list.
\end{itemize}
You can even display poetry.
\begin{verse}
   There is an environment
    for verse \\             % The \\ command separates lines
   Whose features some poets % within a stanza.
   will curse.

                             % One or more blank lines separate stanzas.

   For instead of making\\
   Them do \emph{all} line breaking, \\
   It allows them to put too many words on a line when they'd rather be
   forced to be terse.
\end{verse}

Mathematical formulas may also be displayed.  A
displayed formula
is
one-line long; multiline
formulas require special formatting instructions.
   \[  \ip{\Gamma}{\psi'} = x'' + y^{2} + z_{i}^{n}\]
Don't start a paragraph with a displayed equation,
nor make one a paragraph by itself.

\end{document}               % End of document.

\section{Bench Huge Section 96}
% synthetic bench section 96
% This is a sample LaTeX input file.  (Version of 12 August 2004.)
%
% A '%' character causes TeX to ignore all remaining text on the line,
% and is used for comments like this one.

\documentclass{article}      % Specifies the document class

                             % The preamble begins here.
\title{An Example Document}  % Declares the document's title.
\author{Leslie Lamport}      % Declares the author's name.
\date{January 21, 1994}      % Deleting this command produces today's date.

\newcommand{\ip}[2]{(#1, #2)}
                             % Defines \ip{arg1}{arg2} to mean
                             % (arg1, arg2).

%\newcommand{\ip}[2]{\langle #1 | #2\rangle}
                             % This is an alternative definition of
                             % \ip that is commented out.

\begin{document}             % End of preamble and beginning of text.

\maketitle                   % Produces the title.

This is an example input file.  Comparing it with
the output it generates can show you how to
produce a simple document of your own.

\section{Ordinary Text}      % Produces section heading.  Lower-level
                             % sections are begun with similar
                             % \subsection and \subsubsection commands.

The ends  of words and sentences are marked
  by   spaces. It  doesn't matter how many
spaces    you type; one is as good as 100.  The
end of   a line counts as a space.

One   or more   blank lines denote the  end
of  a paragraph.

Since any number of consecutive spaces are treated
like a single one, the formatting of the input
file makes no difference to
      \LaTeX,                % The \LaTeX command generates the LaTeX logo.
but it makes a difference to you.  When you use
\LaTeX, making your input file as easy to read
as possible will be a great help as you write
your document and when you change it.  This sample
file shows how you can add comments to your own input
file.

Because printing is different from typewriting,
there are a number of things that you have to do
differently when preparing an input file than if
you were just typing the document directly.
Quotation marks like
       ``this''
have to be handled specially, as do quotes within
quotes:
       ``\,`this'            % \, separates the double and single quote.
        is what I just
        wrote, not  `that'\,''.

Dashes come in three sizes: an
       intra-word
dash, a medium dash for number ranges like
       1--2,
and a punctuation
       dash---like
this.

A sentence-ending space should be larger than the
space between words within a sentence.  You
sometimes have to type special commands in
conjunction with punctuation characters to get
this right, as in the following sentence.
       Gnats, gnus, etc.\ all  % `\ ' makes an inter-word space.
       begin with G\@.         % \@ marks end-of-sentence punctuation.
You should check the spaces after periods when
reading your output to make sure you haven't
forgotten any special cases.  Generating an
ellipsis
       \ldots\               % `\ ' is needed after `\ldots' because TeX
                             % ignores spaces after command names like \ldots
                             % made from \ + letters.
                             %
                             % Note how a `%' character causes TeX to ignore
                             % the end of the input line, so these blank lines
                             % do not start a new paragraph.
                             %
with the right spacing around the periods requires
a special command.

\LaTeX\ interprets some common characters as
commands, so you must type special commands to
generate them.  These characters include the
following:
       \$ \& \% \# \{ and \}.

In printing, text is usually emphasized with an
       \emph{italic}
type style.

\begin{em}
   A long segment of text can also be emphasized
   in this way.  Text within such a segment can be
   given \emph{additional} emphasis.
\end{em}

It is sometimes necessary to prevent \LaTeX\ from
breaking a line where it might otherwise do so.
This may be at a space, as between the ``Mr.''\ and
``Jones'' in
       ``Mr.~Jones'',        % ~ produces an unbreakable interword space.
or within a word---especially when the word is a
symbol like
       \mbox{\emph{itemnum}}
that makes little sense when hyphenated across
lines.

Footnotes\footnote{This is an example of a footnote.}
pose no problem.

\LaTeX\ is good at typesetting mathematical formulas
like
       \( x-3y + z = 7 \)
or
       \( a_{1} > x^{2n} + y^{2n} > x' \)
or
       \( \ip{A}{B} = \sum_{i} a_{i} b_{i} \).
The spaces you type in a formula are
ignored.  Remember that a letter like
       $x$                   % $ ... $  and  \( ... \)  are equivalent
is a formula when it denotes a mathematical
symbol, and it should be typed as one.

\section{Displayed Text}

Text is displayed by indenting it from the left
margin.  Quotations are commonly displayed.  There
are short quotations
\begin{quote}
   This is a short quotation.  It consists of a
   single paragraph of text.  See how it is formatted.
\end{quote}
and longer ones.
\begin{quotation}
   This is a longer quotation.  It consists of two
   paragraphs of text, neither of which are
   particularly interesting.

   This is the second paragraph of the quotation.  It
   is just as dull as the first paragraph.
\end{quotation}
Another frequently-displayed structure is a list.
The following is an example of an \emph{itemized}
list.
\begin{itemize}
   \item This is the first item of an itemized list.
         Each item in the list is marked with a ``tick''.
         You don't have to worry about what kind of tick
         mark is used.

   \item This is the second item of the list.  It
         contains another list nested inside it.  The inner
         list is an \emph{enumerated} list.
         \begin{enumerate}
            \item This is the first item of an enumerated
                  list that is nested within the itemized list.

            \item This is the second item of the inner list.
                  \LaTeX\ allows you to nest lists deeper than
                  you really should.
         \end{enumerate}
         This is the rest of the second item of the outer
         list.  It is no more interesting than any other
         part of the item.
   \item This is the third item of the list.
\end{itemize}
You can even display poetry.
\begin{verse}
   There is an environment
    for verse \\             % The \\ command separates lines
   Whose features some poets % within a stanza.
   will curse.

                             % One or more blank lines separate stanzas.

   For instead of making\\
   Them do \emph{all} line breaking, \\
   It allows them to put too many words on a line when they'd rather be
   forced to be terse.
\end{verse}

Mathematical formulas may also be displayed.  A
displayed formula
is
one-line long; multiline
formulas require special formatting instructions.
   \[  \ip{\Gamma}{\psi'} = x'' + y^{2} + z_{i}^{n}\]
Don't start a paragraph with a displayed equation,
nor make one a paragraph by itself.

\end{document}               % End of document.

\section{Bench Huge Section 97}
% synthetic bench section 97
% This is a sample LaTeX input file.  (Version of 12 August 2004.)
%
% A '%' character causes TeX to ignore all remaining text on the line,
% and is used for comments like this one.

\documentclass{article}      % Specifies the document class

                             % The preamble begins here.
\title{An Example Document}  % Declares the document's title.
\author{Leslie Lamport}      % Declares the author's name.
\date{January 21, 1994}      % Deleting this command produces today's date.

\newcommand{\ip}[2]{(#1, #2)}
                             % Defines \ip{arg1}{arg2} to mean
                             % (arg1, arg2).

%\newcommand{\ip}[2]{\langle #1 | #2\rangle}
                             % This is an alternative definition of
                             % \ip that is commented out.

\begin{document}             % End of preamble and beginning of text.

\maketitle                   % Produces the title.

This is an example input file.  Comparing it with
the output it generates can show you how to
produce a simple document of your own.

\section{Ordinary Text}      % Produces section heading.  Lower-level
                             % sections are begun with similar
                             % \subsection and \subsubsection commands.

The ends  of words and sentences are marked
  by   spaces. It  doesn't matter how many
spaces    you type; one is as good as 100.  The
end of   a line counts as a space.

One   or more   blank lines denote the  end
of  a paragraph.

Since any number of consecutive spaces are treated
like a single one, the formatting of the input
file makes no difference to
      \LaTeX,                % The \LaTeX command generates the LaTeX logo.
but it makes a difference to you.  When you use
\LaTeX, making your input file as easy to read
as possible will be a great help as you write
your document and when you change it.  This sample
file shows how you can add comments to your own input
file.

Because printing is different from typewriting,
there are a number of things that you have to do
differently when preparing an input file than if
you were just typing the document directly.
Quotation marks like
       ``this''
have to be handled specially, as do quotes within
quotes:
       ``\,`this'            % \, separates the double and single quote.
        is what I just
        wrote, not  `that'\,''.

Dashes come in three sizes: an
       intra-word
dash, a medium dash for number ranges like
       1--2,
and a punctuation
       dash---like
this.

A sentence-ending space should be larger than the
space between words within a sentence.  You
sometimes have to type special commands in
conjunction with punctuation characters to get
this right, as in the following sentence.
       Gnats, gnus, etc.\ all  % `\ ' makes an inter-word space.
       begin with G\@.         % \@ marks end-of-sentence punctuation.
You should check the spaces after periods when
reading your output to make sure you haven't
forgotten any special cases.  Generating an
ellipsis
       \ldots\               % `\ ' is needed after `\ldots' because TeX
                             % ignores spaces after command names like \ldots
                             % made from \ + letters.
                             %
                             % Note how a `%' character causes TeX to ignore
                             % the end of the input line, so these blank lines
                             % do not start a new paragraph.
                             %
with the right spacing around the periods requires
a special command.

\LaTeX\ interprets some common characters as
commands, so you must type special commands to
generate them.  These characters include the
following:
       \$ \& \% \# \{ and \}.

In printing, text is usually emphasized with an
       \emph{italic}
type style.

\begin{em}
   A long segment of text can also be emphasized
   in this way.  Text within such a segment can be
   given \emph{additional} emphasis.
\end{em}

It is sometimes necessary to prevent \LaTeX\ from
breaking a line where it might otherwise do so.
This may be at a space, as between the ``Mr.''\ and
``Jones'' in
       ``Mr.~Jones'',        % ~ produces an unbreakable interword space.
or within a word---especially when the word is a
symbol like
       \mbox{\emph{itemnum}}
that makes little sense when hyphenated across
lines.

Footnotes\footnote{This is an example of a footnote.}
pose no problem.

\LaTeX\ is good at typesetting mathematical formulas
like
       \( x-3y + z = 7 \)
or
       \( a_{1} > x^{2n} + y^{2n} > x' \)
or
       \( \ip{A}{B} = \sum_{i} a_{i} b_{i} \).
The spaces you type in a formula are
ignored.  Remember that a letter like
       $x$                   % $ ... $  and  \( ... \)  are equivalent
is a formula when it denotes a mathematical
symbol, and it should be typed as one.

\section{Displayed Text}

Text is displayed by indenting it from the left
margin.  Quotations are commonly displayed.  There
are short quotations
\begin{quote}
   This is a short quotation.  It consists of a
   single paragraph of text.  See how it is formatted.
\end{quote}
and longer ones.
\begin{quotation}
   This is a longer quotation.  It consists of two
   paragraphs of text, neither of which are
   particularly interesting.

   This is the second paragraph of the quotation.  It
   is just as dull as the first paragraph.
\end{quotation}
Another frequently-displayed structure is a list.
The following is an example of an \emph{itemized}
list.
\begin{itemize}
   \item This is the first item of an itemized list.
         Each item in the list is marked with a ``tick''.
         You don't have to worry about what kind of tick
         mark is used.

   \item This is the second item of the list.  It
         contains another list nested inside it.  The inner
         list is an \emph{enumerated} list.
         \begin{enumerate}
            \item This is the first item of an enumerated
                  list that is nested within the itemized list.

            \item This is the second item of the inner list.
                  \LaTeX\ allows you to nest lists deeper than
                  you really should.
         \end{enumerate}
         This is the rest of the second item of the outer
         list.  It is no more interesting than any other
         part of the item.
   \item This is the third item of the list.
\end{itemize}
You can even display poetry.
\begin{verse}
   There is an environment
    for verse \\             % The \\ command separates lines
   Whose features some poets % within a stanza.
   will curse.

                             % One or more blank lines separate stanzas.

   For instead of making\\
   Them do \emph{all} line breaking, \\
   It allows them to put too many words on a line when they'd rather be
   forced to be terse.
\end{verse}

Mathematical formulas may also be displayed.  A
displayed formula
is
one-line long; multiline
formulas require special formatting instructions.
   \[  \ip{\Gamma}{\psi'} = x'' + y^{2} + z_{i}^{n}\]
Don't start a paragraph with a displayed equation,
nor make one a paragraph by itself.

\end{document}               % End of document.

\section{Bench Huge Section 98}
% synthetic bench section 98
% This is a sample LaTeX input file.  (Version of 12 August 2004.)
%
% A '%' character causes TeX to ignore all remaining text on the line,
% and is used for comments like this one.

\documentclass{article}      % Specifies the document class

                             % The preamble begins here.
\title{An Example Document}  % Declares the document's title.
\author{Leslie Lamport}      % Declares the author's name.
\date{January 21, 1994}      % Deleting this command produces today's date.

\newcommand{\ip}[2]{(#1, #2)}
                             % Defines \ip{arg1}{arg2} to mean
                             % (arg1, arg2).

%\newcommand{\ip}[2]{\langle #1 | #2\rangle}
                             % This is an alternative definition of
                             % \ip that is commented out.

\begin{document}             % End of preamble and beginning of text.

\maketitle                   % Produces the title.

This is an example input file.  Comparing it with
the output it generates can show you how to
produce a simple document of your own.

\section{Ordinary Text}      % Produces section heading.  Lower-level
                             % sections are begun with similar
                             % \subsection and \subsubsection commands.

The ends  of words and sentences are marked
  by   spaces. It  doesn't matter how many
spaces    you type; one is as good as 100.  The
end of   a line counts as a space.

One   or more   blank lines denote the  end
of  a paragraph.

Since any number of consecutive spaces are treated
like a single one, the formatting of the input
file makes no difference to
      \LaTeX,                % The \LaTeX command generates the LaTeX logo.
but it makes a difference to you.  When you use
\LaTeX, making your input file as easy to read
as possible will be a great help as you write
your document and when you change it.  This sample
file shows how you can add comments to your own input
file.

Because printing is different from typewriting,
there are a number of things that you have to do
differently when preparing an input file than if
you were just typing the document directly.
Quotation marks like
       ``this''
have to be handled specially, as do quotes within
quotes:
       ``\,`this'            % \, separates the double and single quote.
        is what I just
        wrote, not  `that'\,''.

Dashes come in three sizes: an
       intra-word
dash, a medium dash for number ranges like
       1--2,
and a punctuation
       dash---like
this.

A sentence-ending space should be larger than the
space between words within a sentence.  You
sometimes have to type special commands in
conjunction with punctuation characters to get
this right, as in the following sentence.
       Gnats, gnus, etc.\ all  % `\ ' makes an inter-word space.
       begin with G\@.         % \@ marks end-of-sentence punctuation.
You should check the spaces after periods when
reading your output to make sure you haven't
forgotten any special cases.  Generating an
ellipsis
       \ldots\               % `\ ' is needed after `\ldots' because TeX
                             % ignores spaces after command names like \ldots
                             % made from \ + letters.
                             %
                             % Note how a `%' character causes TeX to ignore
                             % the end of the input line, so these blank lines
                             % do not start a new paragraph.
                             %
with the right spacing around the periods requires
a special command.

\LaTeX\ interprets some common characters as
commands, so you must type special commands to
generate them.  These characters include the
following:
       \$ \& \% \# \{ and \}.

In printing, text is usually emphasized with an
       \emph{italic}
type style.

\begin{em}
   A long segment of text can also be emphasized
   in this way.  Text within such a segment can be
   given \emph{additional} emphasis.
\end{em}

It is sometimes necessary to prevent \LaTeX\ from
breaking a line where it might otherwise do so.
This may be at a space, as between the ``Mr.''\ and
``Jones'' in
       ``Mr.~Jones'',        % ~ produces an unbreakable interword space.
or within a word---especially when the word is a
symbol like
       \mbox{\emph{itemnum}}
that makes little sense when hyphenated across
lines.

Footnotes\footnote{This is an example of a footnote.}
pose no problem.

\LaTeX\ is good at typesetting mathematical formulas
like
       \( x-3y + z = 7 \)
or
       \( a_{1} > x^{2n} + y^{2n} > x' \)
or
       \( \ip{A}{B} = \sum_{i} a_{i} b_{i} \).
The spaces you type in a formula are
ignored.  Remember that a letter like
       $x$                   % $ ... $  and  \( ... \)  are equivalent
is a formula when it denotes a mathematical
symbol, and it should be typed as one.

\section{Displayed Text}

Text is displayed by indenting it from the left
margin.  Quotations are commonly displayed.  There
are short quotations
\begin{quote}
   This is a short quotation.  It consists of a
   single paragraph of text.  See how it is formatted.
\end{quote}
and longer ones.
\begin{quotation}
   This is a longer quotation.  It consists of two
   paragraphs of text, neither of which are
   particularly interesting.

   This is the second paragraph of the quotation.  It
   is just as dull as the first paragraph.
\end{quotation}
Another frequently-displayed structure is a list.
The following is an example of an \emph{itemized}
list.
\begin{itemize}
   \item This is the first item of an itemized list.
         Each item in the list is marked with a ``tick''.
         You don't have to worry about what kind of tick
         mark is used.

   \item This is the second item of the list.  It
         contains another list nested inside it.  The inner
         list is an \emph{enumerated} list.
         \begin{enumerate}
            \item This is the first item of an enumerated
                  list that is nested within the itemized list.

            \item This is the second item of the inner list.
                  \LaTeX\ allows you to nest lists deeper than
                  you really should.
         \end{enumerate}
         This is the rest of the second item of the outer
         list.  It is no more interesting than any other
         part of the item.
   \item This is the third item of the list.
\end{itemize}
You can even display poetry.
\begin{verse}
   There is an environment
    for verse \\             % The \\ command separates lines
   Whose features some poets % within a stanza.
   will curse.

                             % One or more blank lines separate stanzas.

   For instead of making\\
   Them do \emph{all} line breaking, \\
   It allows them to put too many words on a line when they'd rather be
   forced to be terse.
\end{verse}

Mathematical formulas may also be displayed.  A
displayed formula
is
one-line long; multiline
formulas require special formatting instructions.
   \[  \ip{\Gamma}{\psi'} = x'' + y^{2} + z_{i}^{n}\]
Don't start a paragraph with a displayed equation,
nor make one a paragraph by itself.

\end{document}               % End of document.

\section{Bench Huge Section 99}
% synthetic bench section 99
% This is a sample LaTeX input file.  (Version of 12 August 2004.)
%
% A '%' character causes TeX to ignore all remaining text on the line,
% and is used for comments like this one.

\documentclass{article}      % Specifies the document class

                             % The preamble begins here.
\title{An Example Document}  % Declares the document's title.
\author{Leslie Lamport}      % Declares the author's name.
\date{January 21, 1994}      % Deleting this command produces today's date.

\newcommand{\ip}[2]{(#1, #2)}
                             % Defines \ip{arg1}{arg2} to mean
                             % (arg1, arg2).

%\newcommand{\ip}[2]{\langle #1 | #2\rangle}
                             % This is an alternative definition of
                             % \ip that is commented out.

\begin{document}             % End of preamble and beginning of text.

\maketitle                   % Produces the title.

This is an example input file.  Comparing it with
the output it generates can show you how to
produce a simple document of your own.

\section{Ordinary Text}      % Produces section heading.  Lower-level
                             % sections are begun with similar
                             % \subsection and \subsubsection commands.

The ends  of words and sentences are marked
  by   spaces. It  doesn't matter how many
spaces    you type; one is as good as 100.  The
end of   a line counts as a space.

One   or more   blank lines denote the  end
of  a paragraph.

Since any number of consecutive spaces are treated
like a single one, the formatting of the input
file makes no difference to
      \LaTeX,                % The \LaTeX command generates the LaTeX logo.
but it makes a difference to you.  When you use
\LaTeX, making your input file as easy to read
as possible will be a great help as you write
your document and when you change it.  This sample
file shows how you can add comments to your own input
file.

Because printing is different from typewriting,
there are a number of things that you have to do
differently when preparing an input file than if
you were just typing the document directly.
Quotation marks like
       ``this''
have to be handled specially, as do quotes within
quotes:
       ``\,`this'            % \, separates the double and single quote.
        is what I just
        wrote, not  `that'\,''.

Dashes come in three sizes: an
       intra-word
dash, a medium dash for number ranges like
       1--2,
and a punctuation
       dash---like
this.

A sentence-ending space should be larger than the
space between words within a sentence.  You
sometimes have to type special commands in
conjunction with punctuation characters to get
this right, as in the following sentence.
       Gnats, gnus, etc.\ all  % `\ ' makes an inter-word space.
       begin with G\@.         % \@ marks end-of-sentence punctuation.
You should check the spaces after periods when
reading your output to make sure you haven't
forgotten any special cases.  Generating an
ellipsis
       \ldots\               % `\ ' is needed after `\ldots' because TeX
                             % ignores spaces after command names like \ldots
                             % made from \ + letters.
                             %
                             % Note how a `%' character causes TeX to ignore
                             % the end of the input line, so these blank lines
                             % do not start a new paragraph.
                             %
with the right spacing around the periods requires
a special command.

\LaTeX\ interprets some common characters as
commands, so you must type special commands to
generate them.  These characters include the
following:
       \$ \& \% \# \{ and \}.

In printing, text is usually emphasized with an
       \emph{italic}
type style.

\begin{em}
   A long segment of text can also be emphasized
   in this way.  Text within such a segment can be
   given \emph{additional} emphasis.
\end{em}

It is sometimes necessary to prevent \LaTeX\ from
breaking a line where it might otherwise do so.
This may be at a space, as between the ``Mr.''\ and
``Jones'' in
       ``Mr.~Jones'',        % ~ produces an unbreakable interword space.
or within a word---especially when the word is a
symbol like
       \mbox{\emph{itemnum}}
that makes little sense when hyphenated across
lines.

Footnotes\footnote{This is an example of a footnote.}
pose no problem.

\LaTeX\ is good at typesetting mathematical formulas
like
       \( x-3y + z = 7 \)
or
       \( a_{1} > x^{2n} + y^{2n} > x' \)
or
       \( \ip{A}{B} = \sum_{i} a_{i} b_{i} \).
The spaces you type in a formula are
ignored.  Remember that a letter like
       $x$                   % $ ... $  and  \( ... \)  are equivalent
is a formula when it denotes a mathematical
symbol, and it should be typed as one.

\section{Displayed Text}

Text is displayed by indenting it from the left
margin.  Quotations are commonly displayed.  There
are short quotations
\begin{quote}
   This is a short quotation.  It consists of a
   single paragraph of text.  See how it is formatted.
\end{quote}
and longer ones.
\begin{quotation}
   This is a longer quotation.  It consists of two
   paragraphs of text, neither of which are
   particularly interesting.

   This is the second paragraph of the quotation.  It
   is just as dull as the first paragraph.
\end{quotation}
Another frequently-displayed structure is a list.
The following is an example of an \emph{itemized}
list.
\begin{itemize}
   \item This is the first item of an itemized list.
         Each item in the list is marked with a ``tick''.
         You don't have to worry about what kind of tick
         mark is used.

   \item This is the second item of the list.  It
         contains another list nested inside it.  The inner
         list is an \emph{enumerated} list.
         \begin{enumerate}
            \item This is the first item of an enumerated
                  list that is nested within the itemized list.

            \item This is the second item of the inner list.
                  \LaTeX\ allows you to nest lists deeper than
                  you really should.
         \end{enumerate}
         This is the rest of the second item of the outer
         list.  It is no more interesting than any other
         part of the item.
   \item This is the third item of the list.
\end{itemize}
You can even display poetry.
\begin{verse}
   There is an environment
    for verse \\             % The \\ command separates lines
   Whose features some poets % within a stanza.
   will curse.

                             % One or more blank lines separate stanzas.

   For instead of making\\
   Them do \emph{all} line breaking, \\
   It allows them to put too many words on a line when they'd rather be
   forced to be terse.
\end{verse}

Mathematical formulas may also be displayed.  A
displayed formula
is
one-line long; multiline
formulas require special formatting instructions.
   \[  \ip{\Gamma}{\psi'} = x'' + y^{2} + z_{i}^{n}\]
Don't start a paragraph with a displayed equation,
nor make one a paragraph by itself.

\end{document}               % End of document.

\section{Bench Huge Section 100}
% synthetic bench section 100
% This is a sample LaTeX input file.  (Version of 12 August 2004.)
%
% A '%' character causes TeX to ignore all remaining text on the line,
% and is used for comments like this one.

\documentclass{article}      % Specifies the document class

                             % The preamble begins here.
\title{An Example Document}  % Declares the document's title.
\author{Leslie Lamport}      % Declares the author's name.
\date{January 21, 1994}      % Deleting this command produces today's date.

\newcommand{\ip}[2]{(#1, #2)}
                             % Defines \ip{arg1}{arg2} to mean
                             % (arg1, arg2).

%\newcommand{\ip}[2]{\langle #1 | #2\rangle}
                             % This is an alternative definition of
                             % \ip that is commented out.

\begin{document}             % End of preamble and beginning of text.

\maketitle                   % Produces the title.

This is an example input file.  Comparing it with
the output it generates can show you how to
produce a simple document of your own.

\section{Ordinary Text}      % Produces section heading.  Lower-level
                             % sections are begun with similar
                             % \subsection and \subsubsection commands.

The ends  of words and sentences are marked
  by   spaces. It  doesn't matter how many
spaces    you type; one is as good as 100.  The
end of   a line counts as a space.

One   or more   blank lines denote the  end
of  a paragraph.

Since any number of consecutive spaces are treated
like a single one, the formatting of the input
file makes no difference to
      \LaTeX,                % The \LaTeX command generates the LaTeX logo.
but it makes a difference to you.  When you use
\LaTeX, making your input file as easy to read
as possible will be a great help as you write
your document and when you change it.  This sample
file shows how you can add comments to your own input
file.

Because printing is different from typewriting,
there are a number of things that you have to do
differently when preparing an input file than if
you were just typing the document directly.
Quotation marks like
       ``this''
have to be handled specially, as do quotes within
quotes:
       ``\,`this'            % \, separates the double and single quote.
        is what I just
        wrote, not  `that'\,''.

Dashes come in three sizes: an
       intra-word
dash, a medium dash for number ranges like
       1--2,
and a punctuation
       dash---like
this.

A sentence-ending space should be larger than the
space between words within a sentence.  You
sometimes have to type special commands in
conjunction with punctuation characters to get
this right, as in the following sentence.
       Gnats, gnus, etc.\ all  % `\ ' makes an inter-word space.
       begin with G\@.         % \@ marks end-of-sentence punctuation.
You should check the spaces after periods when
reading your output to make sure you haven't
forgotten any special cases.  Generating an
ellipsis
       \ldots\               % `\ ' is needed after `\ldots' because TeX
                             % ignores spaces after command names like \ldots
                             % made from \ + letters.
                             %
                             % Note how a `%' character causes TeX to ignore
                             % the end of the input line, so these blank lines
                             % do not start a new paragraph.
                             %
with the right spacing around the periods requires
a special command.

\LaTeX\ interprets some common characters as
commands, so you must type special commands to
generate them.  These characters include the
following:
       \$ \& \% \# \{ and \}.

In printing, text is usually emphasized with an
       \emph{italic}
type style.

\begin{em}
   A long segment of text can also be emphasized
   in this way.  Text within such a segment can be
   given \emph{additional} emphasis.
\end{em}

It is sometimes necessary to prevent \LaTeX\ from
breaking a line where it might otherwise do so.
This may be at a space, as between the ``Mr.''\ and
``Jones'' in
       ``Mr.~Jones'',        % ~ produces an unbreakable interword space.
or within a word---especially when the word is a
symbol like
       \mbox{\emph{itemnum}}
that makes little sense when hyphenated across
lines.

Footnotes\footnote{This is an example of a footnote.}
pose no problem.

\LaTeX\ is good at typesetting mathematical formulas
like
       \( x-3y + z = 7 \)
or
       \( a_{1} > x^{2n} + y^{2n} > x' \)
or
       \( \ip{A}{B} = \sum_{i} a_{i} b_{i} \).
The spaces you type in a formula are
ignored.  Remember that a letter like
       $x$                   % $ ... $  and  \( ... \)  are equivalent
is a formula when it denotes a mathematical
symbol, and it should be typed as one.

\section{Displayed Text}

Text is displayed by indenting it from the left
margin.  Quotations are commonly displayed.  There
are short quotations
\begin{quote}
   This is a short quotation.  It consists of a
   single paragraph of text.  See how it is formatted.
\end{quote}
and longer ones.
\begin{quotation}
   This is a longer quotation.  It consists of two
   paragraphs of text, neither of which are
   particularly interesting.

   This is the second paragraph of the quotation.  It
   is just as dull as the first paragraph.
\end{quotation}
Another frequently-displayed structure is a list.
The following is an example of an \emph{itemized}
list.
\begin{itemize}
   \item This is the first item of an itemized list.
         Each item in the list is marked with a ``tick''.
         You don't have to worry about what kind of tick
         mark is used.

   \item This is the second item of the list.  It
         contains another list nested inside it.  The inner
         list is an \emph{enumerated} list.
         \begin{enumerate}
            \item This is the first item of an enumerated
                  list that is nested within the itemized list.

            \item This is the second item of the inner list.
                  \LaTeX\ allows you to nest lists deeper than
                  you really should.
         \end{enumerate}
         This is the rest of the second item of the outer
         list.  It is no more interesting than any other
         part of the item.
   \item This is the third item of the list.
\end{itemize}
You can even display poetry.
\begin{verse}
   There is an environment
    for verse \\             % The \\ command separates lines
   Whose features some poets % within a stanza.
   will curse.

                             % One or more blank lines separate stanzas.

   For instead of making\\
   Them do \emph{all} line breaking, \\
   It allows them to put too many words on a line when they'd rather be
   forced to be terse.
\end{verse}

Mathematical formulas may also be displayed.  A
displayed formula
is
one-line long; multiline
formulas require special formatting instructions.
   \[  \ip{\Gamma}{\psi'} = x'' + y^{2} + z_{i}^{n}\]
Don't start a paragraph with a displayed equation,
nor make one a paragraph by itself.

\end{document}               % End of document.

\section{Bench Huge Section 101}
% synthetic bench section 101
% This is a sample LaTeX input file.  (Version of 12 August 2004.)
%
% A '%' character causes TeX to ignore all remaining text on the line,
% and is used for comments like this one.

\documentclass{article}      % Specifies the document class

                             % The preamble begins here.
\title{An Example Document}  % Declares the document's title.
\author{Leslie Lamport}      % Declares the author's name.
\date{January 21, 1994}      % Deleting this command produces today's date.

\newcommand{\ip}[2]{(#1, #2)}
                             % Defines \ip{arg1}{arg2} to mean
                             % (arg1, arg2).

%\newcommand{\ip}[2]{\langle #1 | #2\rangle}
                             % This is an alternative definition of
                             % \ip that is commented out.

\begin{document}             % End of preamble and beginning of text.

\maketitle                   % Produces the title.

This is an example input file.  Comparing it with
the output it generates can show you how to
produce a simple document of your own.

\section{Ordinary Text}      % Produces section heading.  Lower-level
                             % sections are begun with similar
                             % \subsection and \subsubsection commands.

The ends  of words and sentences are marked
  by   spaces. It  doesn't matter how many
spaces    you type; one is as good as 100.  The
end of   a line counts as a space.

One   or more   blank lines denote the  end
of  a paragraph.

Since any number of consecutive spaces are treated
like a single one, the formatting of the input
file makes no difference to
      \LaTeX,                % The \LaTeX command generates the LaTeX logo.
but it makes a difference to you.  When you use
\LaTeX, making your input file as easy to read
as possible will be a great help as you write
your document and when you change it.  This sample
file shows how you can add comments to your own input
file.

Because printing is different from typewriting,
there are a number of things that you have to do
differently when preparing an input file than if
you were just typing the document directly.
Quotation marks like
       ``this''
have to be handled specially, as do quotes within
quotes:
       ``\,`this'            % \, separates the double and single quote.
        is what I just
        wrote, not  `that'\,''.

Dashes come in three sizes: an
       intra-word
dash, a medium dash for number ranges like
       1--2,
and a punctuation
       dash---like
this.

A sentence-ending space should be larger than the
space between words within a sentence.  You
sometimes have to type special commands in
conjunction with punctuation characters to get
this right, as in the following sentence.
       Gnats, gnus, etc.\ all  % `\ ' makes an inter-word space.
       begin with G\@.         % \@ marks end-of-sentence punctuation.
You should check the spaces after periods when
reading your output to make sure you haven't
forgotten any special cases.  Generating an
ellipsis
       \ldots\               % `\ ' is needed after `\ldots' because TeX
                             % ignores spaces after command names like \ldots
                             % made from \ + letters.
                             %
                             % Note how a `%' character causes TeX to ignore
                             % the end of the input line, so these blank lines
                             % do not start a new paragraph.
                             %
with the right spacing around the periods requires
a special command.

\LaTeX\ interprets some common characters as
commands, so you must type special commands to
generate them.  These characters include the
following:
       \$ \& \% \# \{ and \}.

In printing, text is usually emphasized with an
       \emph{italic}
type style.

\begin{em}
   A long segment of text can also be emphasized
   in this way.  Text within such a segment can be
   given \emph{additional} emphasis.
\end{em}

It is sometimes necessary to prevent \LaTeX\ from
breaking a line where it might otherwise do so.
This may be at a space, as between the ``Mr.''\ and
``Jones'' in
       ``Mr.~Jones'',        % ~ produces an unbreakable interword space.
or within a word---especially when the word is a
symbol like
       \mbox{\emph{itemnum}}
that makes little sense when hyphenated across
lines.

Footnotes\footnote{This is an example of a footnote.}
pose no problem.

\LaTeX\ is good at typesetting mathematical formulas
like
       \( x-3y + z = 7 \)
or
       \( a_{1} > x^{2n} + y^{2n} > x' \)
or
       \( \ip{A}{B} = \sum_{i} a_{i} b_{i} \).
The spaces you type in a formula are
ignored.  Remember that a letter like
       $x$                   % $ ... $  and  \( ... \)  are equivalent
is a formula when it denotes a mathematical
symbol, and it should be typed as one.

\section{Displayed Text}

Text is displayed by indenting it from the left
margin.  Quotations are commonly displayed.  There
are short quotations
\begin{quote}
   This is a short quotation.  It consists of a
   single paragraph of text.  See how it is formatted.
\end{quote}
and longer ones.
\begin{quotation}
   This is a longer quotation.  It consists of two
   paragraphs of text, neither of which are
   particularly interesting.

   This is the second paragraph of the quotation.  It
   is just as dull as the first paragraph.
\end{quotation}
Another frequently-displayed structure is a list.
The following is an example of an \emph{itemized}
list.
\begin{itemize}
   \item This is the first item of an itemized list.
         Each item in the list is marked with a ``tick''.
         You don't have to worry about what kind of tick
         mark is used.

   \item This is the second item of the list.  It
         contains another list nested inside it.  The inner
         list is an \emph{enumerated} list.
         \begin{enumerate}
            \item This is the first item of an enumerated
                  list that is nested within the itemized list.

            \item This is the second item of the inner list.
                  \LaTeX\ allows you to nest lists deeper than
                  you really should.
         \end{enumerate}
         This is the rest of the second item of the outer
         list.  It is no more interesting than any other
         part of the item.
   \item This is the third item of the list.
\end{itemize}
You can even display poetry.
\begin{verse}
   There is an environment
    for verse \\             % The \\ command separates lines
   Whose features some poets % within a stanza.
   will curse.

                             % One or more blank lines separate stanzas.

   For instead of making\\
   Them do \emph{all} line breaking, \\
   It allows them to put too many words on a line when they'd rather be
   forced to be terse.
\end{verse}

Mathematical formulas may also be displayed.  A
displayed formula
is
one-line long; multiline
formulas require special formatting instructions.
   \[  \ip{\Gamma}{\psi'} = x'' + y^{2} + z_{i}^{n}\]
Don't start a paragraph with a displayed equation,
nor make one a paragraph by itself.

\end{document}               % End of document.

\section{Bench Huge Section 102}
% synthetic bench section 102
% This is a sample LaTeX input file.  (Version of 12 August 2004.)
%
% A '%' character causes TeX to ignore all remaining text on the line,
% and is used for comments like this one.

\documentclass{article}      % Specifies the document class

                             % The preamble begins here.
\title{An Example Document}  % Declares the document's title.
\author{Leslie Lamport}      % Declares the author's name.
\date{January 21, 1994}      % Deleting this command produces today's date.

\newcommand{\ip}[2]{(#1, #2)}
                             % Defines \ip{arg1}{arg2} to mean
                             % (arg1, arg2).

%\newcommand{\ip}[2]{\langle #1 | #2\rangle}
                             % This is an alternative definition of
                             % \ip that is commented out.

\begin{document}             % End of preamble and beginning of text.

\maketitle                   % Produces the title.

This is an example input file.  Comparing it with
the output it generates can show you how to
produce a simple document of your own.

\section{Ordinary Text}      % Produces section heading.  Lower-level
                             % sections are begun with similar
                             % \subsection and \subsubsection commands.

The ends  of words and sentences are marked
  by   spaces. It  doesn't matter how many
spaces    you type; one is as good as 100.  The
end of   a line counts as a space.

One   or more   blank lines denote the  end
of  a paragraph.

Since any number of consecutive spaces are treated
like a single one, the formatting of the input
file makes no difference to
      \LaTeX,                % The \LaTeX command generates the LaTeX logo.
but it makes a difference to you.  When you use
\LaTeX, making your input file as easy to read
as possible will be a great help as you write
your document and when you change it.  This sample
file shows how you can add comments to your own input
file.

Because printing is different from typewriting,
there are a number of things that you have to do
differently when preparing an input file than if
you were just typing the document directly.
Quotation marks like
       ``this''
have to be handled specially, as do quotes within
quotes:
       ``\,`this'            % \, separates the double and single quote.
        is what I just
        wrote, not  `that'\,''.

Dashes come in three sizes: an
       intra-word
dash, a medium dash for number ranges like
       1--2,
and a punctuation
       dash---like
this.

A sentence-ending space should be larger than the
space between words within a sentence.  You
sometimes have to type special commands in
conjunction with punctuation characters to get
this right, as in the following sentence.
       Gnats, gnus, etc.\ all  % `\ ' makes an inter-word space.
       begin with G\@.         % \@ marks end-of-sentence punctuation.
You should check the spaces after periods when
reading your output to make sure you haven't
forgotten any special cases.  Generating an
ellipsis
       \ldots\               % `\ ' is needed after `\ldots' because TeX
                             % ignores spaces after command names like \ldots
                             % made from \ + letters.
                             %
                             % Note how a `%' character causes TeX to ignore
                             % the end of the input line, so these blank lines
                             % do not start a new paragraph.
                             %
with the right spacing around the periods requires
a special command.

\LaTeX\ interprets some common characters as
commands, so you must type special commands to
generate them.  These characters include the
following:
       \$ \& \% \# \{ and \}.

In printing, text is usually emphasized with an
       \emph{italic}
type style.

\begin{em}
   A long segment of text can also be emphasized
   in this way.  Text within such a segment can be
   given \emph{additional} emphasis.
\end{em}

It is sometimes necessary to prevent \LaTeX\ from
breaking a line where it might otherwise do so.
This may be at a space, as between the ``Mr.''\ and
``Jones'' in
       ``Mr.~Jones'',        % ~ produces an unbreakable interword space.
or within a word---especially when the word is a
symbol like
       \mbox{\emph{itemnum}}
that makes little sense when hyphenated across
lines.

Footnotes\footnote{This is an example of a footnote.}
pose no problem.

\LaTeX\ is good at typesetting mathematical formulas
like
       \( x-3y + z = 7 \)
or
       \( a_{1} > x^{2n} + y^{2n} > x' \)
or
       \( \ip{A}{B} = \sum_{i} a_{i} b_{i} \).
The spaces you type in a formula are
ignored.  Remember that a letter like
       $x$                   % $ ... $  and  \( ... \)  are equivalent
is a formula when it denotes a mathematical
symbol, and it should be typed as one.

\section{Displayed Text}

Text is displayed by indenting it from the left
margin.  Quotations are commonly displayed.  There
are short quotations
\begin{quote}
   This is a short quotation.  It consists of a
   single paragraph of text.  See how it is formatted.
\end{quote}
and longer ones.
\begin{quotation}
   This is a longer quotation.  It consists of two
   paragraphs of text, neither of which are
   particularly interesting.

   This is the second paragraph of the quotation.  It
   is just as dull as the first paragraph.
\end{quotation}
Another frequently-displayed structure is a list.
The following is an example of an \emph{itemized}
list.
\begin{itemize}
   \item This is the first item of an itemized list.
         Each item in the list is marked with a ``tick''.
         You don't have to worry about what kind of tick
         mark is used.

   \item This is the second item of the list.  It
         contains another list nested inside it.  The inner
         list is an \emph{enumerated} list.
         \begin{enumerate}
            \item This is the first item of an enumerated
                  list that is nested within the itemized list.

            \item This is the second item of the inner list.
                  \LaTeX\ allows you to nest lists deeper than
                  you really should.
         \end{enumerate}
         This is the rest of the second item of the outer
         list.  It is no more interesting than any other
         part of the item.
   \item This is the third item of the list.
\end{itemize}
You can even display poetry.
\begin{verse}
   There is an environment
    for verse \\             % The \\ command separates lines
   Whose features some poets % within a stanza.
   will curse.

                             % One or more blank lines separate stanzas.

   For instead of making\\
   Them do \emph{all} line breaking, \\
   It allows them to put too many words on a line when they'd rather be
   forced to be terse.
\end{verse}

Mathematical formulas may also be displayed.  A
displayed formula
is
one-line long; multiline
formulas require special formatting instructions.
   \[  \ip{\Gamma}{\psi'} = x'' + y^{2} + z_{i}^{n}\]
Don't start a paragraph with a displayed equation,
nor make one a paragraph by itself.

\end{document}               % End of document.

\section{Bench Huge Section 103}
% synthetic bench section 103
% This is a sample LaTeX input file.  (Version of 12 August 2004.)
%
% A '%' character causes TeX to ignore all remaining text on the line,
% and is used for comments like this one.

\documentclass{article}      % Specifies the document class

                             % The preamble begins here.
\title{An Example Document}  % Declares the document's title.
\author{Leslie Lamport}      % Declares the author's name.
\date{January 21, 1994}      % Deleting this command produces today's date.

\newcommand{\ip}[2]{(#1, #2)}
                             % Defines \ip{arg1}{arg2} to mean
                             % (arg1, arg2).

%\newcommand{\ip}[2]{\langle #1 | #2\rangle}
                             % This is an alternative definition of
                             % \ip that is commented out.

\begin{document}             % End of preamble and beginning of text.

\maketitle                   % Produces the title.

This is an example input file.  Comparing it with
the output it generates can show you how to
produce a simple document of your own.

\section{Ordinary Text}      % Produces section heading.  Lower-level
                             % sections are begun with similar
                             % \subsection and \subsubsection commands.

The ends  of words and sentences are marked
  by   spaces. It  doesn't matter how many
spaces    you type; one is as good as 100.  The
end of   a line counts as a space.

One   or more   blank lines denote the  end
of  a paragraph.

Since any number of consecutive spaces are treated
like a single one, the formatting of the input
file makes no difference to
      \LaTeX,                % The \LaTeX command generates the LaTeX logo.
but it makes a difference to you.  When you use
\LaTeX, making your input file as easy to read
as possible will be a great help as you write
your document and when you change it.  This sample
file shows how you can add comments to your own input
file.

Because printing is different from typewriting,
there are a number of things that you have to do
differently when preparing an input file than if
you were just typing the document directly.
Quotation marks like
       ``this''
have to be handled specially, as do quotes within
quotes:
       ``\,`this'            % \, separates the double and single quote.
        is what I just
        wrote, not  `that'\,''.

Dashes come in three sizes: an
       intra-word
dash, a medium dash for number ranges like
       1--2,
and a punctuation
       dash---like
this.

A sentence-ending space should be larger than the
space between words within a sentence.  You
sometimes have to type special commands in
conjunction with punctuation characters to get
this right, as in the following sentence.
       Gnats, gnus, etc.\ all  % `\ ' makes an inter-word space.
       begin with G\@.         % \@ marks end-of-sentence punctuation.
You should check the spaces after periods when
reading your output to make sure you haven't
forgotten any special cases.  Generating an
ellipsis
       \ldots\               % `\ ' is needed after `\ldots' because TeX
                             % ignores spaces after command names like \ldots
                             % made from \ + letters.
                             %
                             % Note how a `%' character causes TeX to ignore
                             % the end of the input line, so these blank lines
                             % do not start a new paragraph.
                             %
with the right spacing around the periods requires
a special command.

\LaTeX\ interprets some common characters as
commands, so you must type special commands to
generate them.  These characters include the
following:
       \$ \& \% \# \{ and \}.

In printing, text is usually emphasized with an
       \emph{italic}
type style.

\begin{em}
   A long segment of text can also be emphasized
   in this way.  Text within such a segment can be
   given \emph{additional} emphasis.
\end{em}

It is sometimes necessary to prevent \LaTeX\ from
breaking a line where it might otherwise do so.
This may be at a space, as between the ``Mr.''\ and
``Jones'' in
       ``Mr.~Jones'',        % ~ produces an unbreakable interword space.
or within a word---especially when the word is a
symbol like
       \mbox{\emph{itemnum}}
that makes little sense when hyphenated across
lines.

Footnotes\footnote{This is an example of a footnote.}
pose no problem.

\LaTeX\ is good at typesetting mathematical formulas
like
       \( x-3y + z = 7 \)
or
       \( a_{1} > x^{2n} + y^{2n} > x' \)
or
       \( \ip{A}{B} = \sum_{i} a_{i} b_{i} \).
The spaces you type in a formula are
ignored.  Remember that a letter like
       $x$                   % $ ... $  and  \( ... \)  are equivalent
is a formula when it denotes a mathematical
symbol, and it should be typed as one.

\section{Displayed Text}

Text is displayed by indenting it from the left
margin.  Quotations are commonly displayed.  There
are short quotations
\begin{quote}
   This is a short quotation.  It consists of a
   single paragraph of text.  See how it is formatted.
\end{quote}
and longer ones.
\begin{quotation}
   This is a longer quotation.  It consists of two
   paragraphs of text, neither of which are
   particularly interesting.

   This is the second paragraph of the quotation.  It
   is just as dull as the first paragraph.
\end{quotation}
Another frequently-displayed structure is a list.
The following is an example of an \emph{itemized}
list.
\begin{itemize}
   \item This is the first item of an itemized list.
         Each item in the list is marked with a ``tick''.
         You don't have to worry about what kind of tick
         mark is used.

   \item This is the second item of the list.  It
         contains another list nested inside it.  The inner
         list is an \emph{enumerated} list.
         \begin{enumerate}
            \item This is the first item of an enumerated
                  list that is nested within the itemized list.

            \item This is the second item of the inner list.
                  \LaTeX\ allows you to nest lists deeper than
                  you really should.
         \end{enumerate}
         This is the rest of the second item of the outer
         list.  It is no more interesting than any other
         part of the item.
   \item This is the third item of the list.
\end{itemize}
You can even display poetry.
\begin{verse}
   There is an environment
    for verse \\             % The \\ command separates lines
   Whose features some poets % within a stanza.
   will curse.

                             % One or more blank lines separate stanzas.

   For instead of making\\
   Them do \emph{all} line breaking, \\
   It allows them to put too many words on a line when they'd rather be
   forced to be terse.
\end{verse}

Mathematical formulas may also be displayed.  A
displayed formula
is
one-line long; multiline
formulas require special formatting instructions.
   \[  \ip{\Gamma}{\psi'} = x'' + y^{2} + z_{i}^{n}\]
Don't start a paragraph with a displayed equation,
nor make one a paragraph by itself.

\end{document}               % End of document.

\section{Bench Huge Section 104}
% synthetic bench section 104
% This is a sample LaTeX input file.  (Version of 12 August 2004.)
%
% A '%' character causes TeX to ignore all remaining text on the line,
% and is used for comments like this one.

\documentclass{article}      % Specifies the document class

                             % The preamble begins here.
\title{An Example Document}  % Declares the document's title.
\author{Leslie Lamport}      % Declares the author's name.
\date{January 21, 1994}      % Deleting this command produces today's date.

\newcommand{\ip}[2]{(#1, #2)}
                             % Defines \ip{arg1}{arg2} to mean
                             % (arg1, arg2).

%\newcommand{\ip}[2]{\langle #1 | #2\rangle}
                             % This is an alternative definition of
                             % \ip that is commented out.

\begin{document}             % End of preamble and beginning of text.

\maketitle                   % Produces the title.

This is an example input file.  Comparing it with
the output it generates can show you how to
produce a simple document of your own.

\section{Ordinary Text}      % Produces section heading.  Lower-level
                             % sections are begun with similar
                             % \subsection and \subsubsection commands.

The ends  of words and sentences are marked
  by   spaces. It  doesn't matter how many
spaces    you type; one is as good as 100.  The
end of   a line counts as a space.

One   or more   blank lines denote the  end
of  a paragraph.

Since any number of consecutive spaces are treated
like a single one, the formatting of the input
file makes no difference to
      \LaTeX,                % The \LaTeX command generates the LaTeX logo.
but it makes a difference to you.  When you use
\LaTeX, making your input file as easy to read
as possible will be a great help as you write
your document and when you change it.  This sample
file shows how you can add comments to your own input
file.

Because printing is different from typewriting,
there are a number of things that you have to do
differently when preparing an input file than if
you were just typing the document directly.
Quotation marks like
       ``this''
have to be handled specially, as do quotes within
quotes:
       ``\,`this'            % \, separates the double and single quote.
        is what I just
        wrote, not  `that'\,''.

Dashes come in three sizes: an
       intra-word
dash, a medium dash for number ranges like
       1--2,
and a punctuation
       dash---like
this.

A sentence-ending space should be larger than the
space between words within a sentence.  You
sometimes have to type special commands in
conjunction with punctuation characters to get
this right, as in the following sentence.
       Gnats, gnus, etc.\ all  % `\ ' makes an inter-word space.
       begin with G\@.         % \@ marks end-of-sentence punctuation.
You should check the spaces after periods when
reading your output to make sure you haven't
forgotten any special cases.  Generating an
ellipsis
       \ldots\               % `\ ' is needed after `\ldots' because TeX
                             % ignores spaces after command names like \ldots
                             % made from \ + letters.
                             %
                             % Note how a `%' character causes TeX to ignore
                             % the end of the input line, so these blank lines
                             % do not start a new paragraph.
                             %
with the right spacing around the periods requires
a special command.

\LaTeX\ interprets some common characters as
commands, so you must type special commands to
generate them.  These characters include the
following:
       \$ \& \% \# \{ and \}.

In printing, text is usually emphasized with an
       \emph{italic}
type style.

\begin{em}
   A long segment of text can also be emphasized
   in this way.  Text within such a segment can be
   given \emph{additional} emphasis.
\end{em}

It is sometimes necessary to prevent \LaTeX\ from
breaking a line where it might otherwise do so.
This may be at a space, as between the ``Mr.''\ and
``Jones'' in
       ``Mr.~Jones'',        % ~ produces an unbreakable interword space.
or within a word---especially when the word is a
symbol like
       \mbox{\emph{itemnum}}
that makes little sense when hyphenated across
lines.

Footnotes\footnote{This is an example of a footnote.}
pose no problem.

\LaTeX\ is good at typesetting mathematical formulas
like
       \( x-3y + z = 7 \)
or
       \( a_{1} > x^{2n} + y^{2n} > x' \)
or
       \( \ip{A}{B} = \sum_{i} a_{i} b_{i} \).
The spaces you type in a formula are
ignored.  Remember that a letter like
       $x$                   % $ ... $  and  \( ... \)  are equivalent
is a formula when it denotes a mathematical
symbol, and it should be typed as one.

\section{Displayed Text}

Text is displayed by indenting it from the left
margin.  Quotations are commonly displayed.  There
are short quotations
\begin{quote}
   This is a short quotation.  It consists of a
   single paragraph of text.  See how it is formatted.
\end{quote}
and longer ones.
\begin{quotation}
   This is a longer quotation.  It consists of two
   paragraphs of text, neither of which are
   particularly interesting.

   This is the second paragraph of the quotation.  It
   is just as dull as the first paragraph.
\end{quotation}
Another frequently-displayed structure is a list.
The following is an example of an \emph{itemized}
list.
\begin{itemize}
   \item This is the first item of an itemized list.
         Each item in the list is marked with a ``tick''.
         You don't have to worry about what kind of tick
         mark is used.

   \item This is the second item of the list.  It
         contains another list nested inside it.  The inner
         list is an \emph{enumerated} list.
         \begin{enumerate}
            \item This is the first item of an enumerated
                  list that is nested within the itemized list.

            \item This is the second item of the inner list.
                  \LaTeX\ allows you to nest lists deeper than
                  you really should.
         \end{enumerate}
         This is the rest of the second item of the outer
         list.  It is no more interesting than any other
         part of the item.
   \item This is the third item of the list.
\end{itemize}
You can even display poetry.
\begin{verse}
   There is an environment
    for verse \\             % The \\ command separates lines
   Whose features some poets % within a stanza.
   will curse.

                             % One or more blank lines separate stanzas.

   For instead of making\\
   Them do \emph{all} line breaking, \\
   It allows them to put too many words on a line when they'd rather be
   forced to be terse.
\end{verse}

Mathematical formulas may also be displayed.  A
displayed formula
is
one-line long; multiline
formulas require special formatting instructions.
   \[  \ip{\Gamma}{\psi'} = x'' + y^{2} + z_{i}^{n}\]
Don't start a paragraph with a displayed equation,
nor make one a paragraph by itself.

\end{document}               % End of document.

\section{Bench Huge Section 105}
% synthetic bench section 105
% This is a sample LaTeX input file.  (Version of 12 August 2004.)
%
% A '%' character causes TeX to ignore all remaining text on the line,
% and is used for comments like this one.

\documentclass{article}      % Specifies the document class

                             % The preamble begins here.
\title{An Example Document}  % Declares the document's title.
\author{Leslie Lamport}      % Declares the author's name.
\date{January 21, 1994}      % Deleting this command produces today's date.

\newcommand{\ip}[2]{(#1, #2)}
                             % Defines \ip{arg1}{arg2} to mean
                             % (arg1, arg2).

%\newcommand{\ip}[2]{\langle #1 | #2\rangle}
                             % This is an alternative definition of
                             % \ip that is commented out.

\begin{document}             % End of preamble and beginning of text.

\maketitle                   % Produces the title.

This is an example input file.  Comparing it with
the output it generates can show you how to
produce a simple document of your own.

\section{Ordinary Text}      % Produces section heading.  Lower-level
                             % sections are begun with similar
                             % \subsection and \subsubsection commands.

The ends  of words and sentences are marked
  by   spaces. It  doesn't matter how many
spaces    you type; one is as good as 100.  The
end of   a line counts as a space.

One   or more   blank lines denote the  end
of  a paragraph.

Since any number of consecutive spaces are treated
like a single one, the formatting of the input
file makes no difference to
      \LaTeX,                % The \LaTeX command generates the LaTeX logo.
but it makes a difference to you.  When you use
\LaTeX, making your input file as easy to read
as possible will be a great help as you write
your document and when you change it.  This sample
file shows how you can add comments to your own input
file.

Because printing is different from typewriting,
there are a number of things that you have to do
differently when preparing an input file than if
you were just typing the document directly.
Quotation marks like
       ``this''
have to be handled specially, as do quotes within
quotes:
       ``\,`this'            % \, separates the double and single quote.
        is what I just
        wrote, not  `that'\,''.

Dashes come in three sizes: an
       intra-word
dash, a medium dash for number ranges like
       1--2,
and a punctuation
       dash---like
this.

A sentence-ending space should be larger than the
space between words within a sentence.  You
sometimes have to type special commands in
conjunction with punctuation characters to get
this right, as in the following sentence.
       Gnats, gnus, etc.\ all  % `\ ' makes an inter-word space.
       begin with G\@.         % \@ marks end-of-sentence punctuation.
You should check the spaces after periods when
reading your output to make sure you haven't
forgotten any special cases.  Generating an
ellipsis
       \ldots\               % `\ ' is needed after `\ldots' because TeX
                             % ignores spaces after command names like \ldots
                             % made from \ + letters.
                             %
                             % Note how a `%' character causes TeX to ignore
                             % the end of the input line, so these blank lines
                             % do not start a new paragraph.
                             %
with the right spacing around the periods requires
a special command.

\LaTeX\ interprets some common characters as
commands, so you must type special commands to
generate them.  These characters include the
following:
       \$ \& \% \# \{ and \}.

In printing, text is usually emphasized with an
       \emph{italic}
type style.

\begin{em}
   A long segment of text can also be emphasized
   in this way.  Text within such a segment can be
   given \emph{additional} emphasis.
\end{em}

It is sometimes necessary to prevent \LaTeX\ from
breaking a line where it might otherwise do so.
This may be at a space, as between the ``Mr.''\ and
``Jones'' in
       ``Mr.~Jones'',        % ~ produces an unbreakable interword space.
or within a word---especially when the word is a
symbol like
       \mbox{\emph{itemnum}}
that makes little sense when hyphenated across
lines.

Footnotes\footnote{This is an example of a footnote.}
pose no problem.

\LaTeX\ is good at typesetting mathematical formulas
like
       \( x-3y + z = 7 \)
or
       \( a_{1} > x^{2n} + y^{2n} > x' \)
or
       \( \ip{A}{B} = \sum_{i} a_{i} b_{i} \).
The spaces you type in a formula are
ignored.  Remember that a letter like
       $x$                   % $ ... $  and  \( ... \)  are equivalent
is a formula when it denotes a mathematical
symbol, and it should be typed as one.

\section{Displayed Text}

Text is displayed by indenting it from the left
margin.  Quotations are commonly displayed.  There
are short quotations
\begin{quote}
   This is a short quotation.  It consists of a
   single paragraph of text.  See how it is formatted.
\end{quote}
and longer ones.
\begin{quotation}
   This is a longer quotation.  It consists of two
   paragraphs of text, neither of which are
   particularly interesting.

   This is the second paragraph of the quotation.  It
   is just as dull as the first paragraph.
\end{quotation}
Another frequently-displayed structure is a list.
The following is an example of an \emph{itemized}
list.
\begin{itemize}
   \item This is the first item of an itemized list.
         Each item in the list is marked with a ``tick''.
         You don't have to worry about what kind of tick
         mark is used.

   \item This is the second item of the list.  It
         contains another list nested inside it.  The inner
         list is an \emph{enumerated} list.
         \begin{enumerate}
            \item This is the first item of an enumerated
                  list that is nested within the itemized list.

            \item This is the second item of the inner list.
                  \LaTeX\ allows you to nest lists deeper than
                  you really should.
         \end{enumerate}
         This is the rest of the second item of the outer
         list.  It is no more interesting than any other
         part of the item.
   \item This is the third item of the list.
\end{itemize}
You can even display poetry.
\begin{verse}
   There is an environment
    for verse \\             % The \\ command separates lines
   Whose features some poets % within a stanza.
   will curse.

                             % One or more blank lines separate stanzas.

   For instead of making\\
   Them do \emph{all} line breaking, \\
   It allows them to put too many words on a line when they'd rather be
   forced to be terse.
\end{verse}

Mathematical formulas may also be displayed.  A
displayed formula
is
one-line long; multiline
formulas require special formatting instructions.
   \[  \ip{\Gamma}{\psi'} = x'' + y^{2} + z_{i}^{n}\]
Don't start a paragraph with a displayed equation,
nor make one a paragraph by itself.

\end{document}               % End of document.

\section{Bench Huge Section 106}
% synthetic bench section 106
% This is a sample LaTeX input file.  (Version of 12 August 2004.)
%
% A '%' character causes TeX to ignore all remaining text on the line,
% and is used for comments like this one.

\documentclass{article}      % Specifies the document class

                             % The preamble begins here.
\title{An Example Document}  % Declares the document's title.
\author{Leslie Lamport}      % Declares the author's name.
\date{January 21, 1994}      % Deleting this command produces today's date.

\newcommand{\ip}[2]{(#1, #2)}
                             % Defines \ip{arg1}{arg2} to mean
                             % (arg1, arg2).

%\newcommand{\ip}[2]{\langle #1 | #2\rangle}
                             % This is an alternative definition of
                             % \ip that is commented out.

\begin{document}             % End of preamble and beginning of text.

\maketitle                   % Produces the title.

This is an example input file.  Comparing it with
the output it generates can show you how to
produce a simple document of your own.

\section{Ordinary Text}      % Produces section heading.  Lower-level
                             % sections are begun with similar
                             % \subsection and \subsubsection commands.

The ends  of words and sentences are marked
  by   spaces. It  doesn't matter how many
spaces    you type; one is as good as 100.  The
end of   a line counts as a space.

One   or more   blank lines denote the  end
of  a paragraph.

Since any number of consecutive spaces are treated
like a single one, the formatting of the input
file makes no difference to
      \LaTeX,                % The \LaTeX command generates the LaTeX logo.
but it makes a difference to you.  When you use
\LaTeX, making your input file as easy to read
as possible will be a great help as you write
your document and when you change it.  This sample
file shows how you can add comments to your own input
file.

Because printing is different from typewriting,
there are a number of things that you have to do
differently when preparing an input file than if
you were just typing the document directly.
Quotation marks like
       ``this''
have to be handled specially, as do quotes within
quotes:
       ``\,`this'            % \, separates the double and single quote.
        is what I just
        wrote, not  `that'\,''.

Dashes come in three sizes: an
       intra-word
dash, a medium dash for number ranges like
       1--2,
and a punctuation
       dash---like
this.

A sentence-ending space should be larger than the
space between words within a sentence.  You
sometimes have to type special commands in
conjunction with punctuation characters to get
this right, as in the following sentence.
       Gnats, gnus, etc.\ all  % `\ ' makes an inter-word space.
       begin with G\@.         % \@ marks end-of-sentence punctuation.
You should check the spaces after periods when
reading your output to make sure you haven't
forgotten any special cases.  Generating an
ellipsis
       \ldots\               % `\ ' is needed after `\ldots' because TeX
                             % ignores spaces after command names like \ldots
                             % made from \ + letters.
                             %
                             % Note how a `%' character causes TeX to ignore
                             % the end of the input line, so these blank lines
                             % do not start a new paragraph.
                             %
with the right spacing around the periods requires
a special command.

\LaTeX\ interprets some common characters as
commands, so you must type special commands to
generate them.  These characters include the
following:
       \$ \& \% \# \{ and \}.

In printing, text is usually emphasized with an
       \emph{italic}
type style.

\begin{em}
   A long segment of text can also be emphasized
   in this way.  Text within such a segment can be
   given \emph{additional} emphasis.
\end{em}

It is sometimes necessary to prevent \LaTeX\ from
breaking a line where it might otherwise do so.
This may be at a space, as between the ``Mr.''\ and
``Jones'' in
       ``Mr.~Jones'',        % ~ produces an unbreakable interword space.
or within a word---especially when the word is a
symbol like
       \mbox{\emph{itemnum}}
that makes little sense when hyphenated across
lines.

Footnotes\footnote{This is an example of a footnote.}
pose no problem.

\LaTeX\ is good at typesetting mathematical formulas
like
       \( x-3y + z = 7 \)
or
       \( a_{1} > x^{2n} + y^{2n} > x' \)
or
       \( \ip{A}{B} = \sum_{i} a_{i} b_{i} \).
The spaces you type in a formula are
ignored.  Remember that a letter like
       $x$                   % $ ... $  and  \( ... \)  are equivalent
is a formula when it denotes a mathematical
symbol, and it should be typed as one.

\section{Displayed Text}

Text is displayed by indenting it from the left
margin.  Quotations are commonly displayed.  There
are short quotations
\begin{quote}
   This is a short quotation.  It consists of a
   single paragraph of text.  See how it is formatted.
\end{quote}
and longer ones.
\begin{quotation}
   This is a longer quotation.  It consists of two
   paragraphs of text, neither of which are
   particularly interesting.

   This is the second paragraph of the quotation.  It
   is just as dull as the first paragraph.
\end{quotation}
Another frequently-displayed structure is a list.
The following is an example of an \emph{itemized}
list.
\begin{itemize}
   \item This is the first item of an itemized list.
         Each item in the list is marked with a ``tick''.
         You don't have to worry about what kind of tick
         mark is used.

   \item This is the second item of the list.  It
         contains another list nested inside it.  The inner
         list is an \emph{enumerated} list.
         \begin{enumerate}
            \item This is the first item of an enumerated
                  list that is nested within the itemized list.

            \item This is the second item of the inner list.
                  \LaTeX\ allows you to nest lists deeper than
                  you really should.
         \end{enumerate}
         This is the rest of the second item of the outer
         list.  It is no more interesting than any other
         part of the item.
   \item This is the third item of the list.
\end{itemize}
You can even display poetry.
\begin{verse}
   There is an environment
    for verse \\             % The \\ command separates lines
   Whose features some poets % within a stanza.
   will curse.

                             % One or more blank lines separate stanzas.

   For instead of making\\
   Them do \emph{all} line breaking, \\
   It allows them to put too many words on a line when they'd rather be
   forced to be terse.
\end{verse}

Mathematical formulas may also be displayed.  A
displayed formula
is
one-line long; multiline
formulas require special formatting instructions.
   \[  \ip{\Gamma}{\psi'} = x'' + y^{2} + z_{i}^{n}\]
Don't start a paragraph with a displayed equation,
nor make one a paragraph by itself.

\end{document}               % End of document.

\section{Bench Huge Section 107}
% synthetic bench section 107
% This is a sample LaTeX input file.  (Version of 12 August 2004.)
%
% A '%' character causes TeX to ignore all remaining text on the line,
% and is used for comments like this one.

\documentclass{article}      % Specifies the document class

                             % The preamble begins here.
\title{An Example Document}  % Declares the document's title.
\author{Leslie Lamport}      % Declares the author's name.
\date{January 21, 1994}      % Deleting this command produces today's date.

\newcommand{\ip}[2]{(#1, #2)}
                             % Defines \ip{arg1}{arg2} to mean
                             % (arg1, arg2).

%\newcommand{\ip}[2]{\langle #1 | #2\rangle}
                             % This is an alternative definition of
                             % \ip that is commented out.

\begin{document}             % End of preamble and beginning of text.

\maketitle                   % Produces the title.

This is an example input file.  Comparing it with
the output it generates can show you how to
produce a simple document of your own.

\section{Ordinary Text}      % Produces section heading.  Lower-level
                             % sections are begun with similar
                             % \subsection and \subsubsection commands.

The ends  of words and sentences are marked
  by   spaces. It  doesn't matter how many
spaces    you type; one is as good as 100.  The
end of   a line counts as a space.

One   or more   blank lines denote the  end
of  a paragraph.

Since any number of consecutive spaces are treated
like a single one, the formatting of the input
file makes no difference to
      \LaTeX,                % The \LaTeX command generates the LaTeX logo.
but it makes a difference to you.  When you use
\LaTeX, making your input file as easy to read
as possible will be a great help as you write
your document and when you change it.  This sample
file shows how you can add comments to your own input
file.

Because printing is different from typewriting,
there are a number of things that you have to do
differently when preparing an input file than if
you were just typing the document directly.
Quotation marks like
       ``this''
have to be handled specially, as do quotes within
quotes:
       ``\,`this'            % \, separates the double and single quote.
        is what I just
        wrote, not  `that'\,''.

Dashes come in three sizes: an
       intra-word
dash, a medium dash for number ranges like
       1--2,
and a punctuation
       dash---like
this.

A sentence-ending space should be larger than the
space between words within a sentence.  You
sometimes have to type special commands in
conjunction with punctuation characters to get
this right, as in the following sentence.
       Gnats, gnus, etc.\ all  % `\ ' makes an inter-word space.
       begin with G\@.         % \@ marks end-of-sentence punctuation.
You should check the spaces after periods when
reading your output to make sure you haven't
forgotten any special cases.  Generating an
ellipsis
       \ldots\               % `\ ' is needed after `\ldots' because TeX
                             % ignores spaces after command names like \ldots
                             % made from \ + letters.
                             %
                             % Note how a `%' character causes TeX to ignore
                             % the end of the input line, so these blank lines
                             % do not start a new paragraph.
                             %
with the right spacing around the periods requires
a special command.

\LaTeX\ interprets some common characters as
commands, so you must type special commands to
generate them.  These characters include the
following:
       \$ \& \% \# \{ and \}.

In printing, text is usually emphasized with an
       \emph{italic}
type style.

\begin{em}
   A long segment of text can also be emphasized
   in this way.  Text within such a segment can be
   given \emph{additional} emphasis.
\end{em}

It is sometimes necessary to prevent \LaTeX\ from
breaking a line where it might otherwise do so.
This may be at a space, as between the ``Mr.''\ and
``Jones'' in
       ``Mr.~Jones'',        % ~ produces an unbreakable interword space.
or within a word---especially when the word is a
symbol like
       \mbox{\emph{itemnum}}
that makes little sense when hyphenated across
lines.

Footnotes\footnote{This is an example of a footnote.}
pose no problem.

\LaTeX\ is good at typesetting mathematical formulas
like
       \( x-3y + z = 7 \)
or
       \( a_{1} > x^{2n} + y^{2n} > x' \)
or
       \( \ip{A}{B} = \sum_{i} a_{i} b_{i} \).
The spaces you type in a formula are
ignored.  Remember that a letter like
       $x$                   % $ ... $  and  \( ... \)  are equivalent
is a formula when it denotes a mathematical
symbol, and it should be typed as one.

\section{Displayed Text}

Text is displayed by indenting it from the left
margin.  Quotations are commonly displayed.  There
are short quotations
\begin{quote}
   This is a short quotation.  It consists of a
   single paragraph of text.  See how it is formatted.
\end{quote}
and longer ones.
\begin{quotation}
   This is a longer quotation.  It consists of two
   paragraphs of text, neither of which are
   particularly interesting.

   This is the second paragraph of the quotation.  It
   is just as dull as the first paragraph.
\end{quotation}
Another frequently-displayed structure is a list.
The following is an example of an \emph{itemized}
list.
\begin{itemize}
   \item This is the first item of an itemized list.
         Each item in the list is marked with a ``tick''.
         You don't have to worry about what kind of tick
         mark is used.

   \item This is the second item of the list.  It
         contains another list nested inside it.  The inner
         list is an \emph{enumerated} list.
         \begin{enumerate}
            \item This is the first item of an enumerated
                  list that is nested within the itemized list.

            \item This is the second item of the inner list.
                  \LaTeX\ allows you to nest lists deeper than
                  you really should.
         \end{enumerate}
         This is the rest of the second item of the outer
         list.  It is no more interesting than any other
         part of the item.
   \item This is the third item of the list.
\end{itemize}
You can even display poetry.
\begin{verse}
   There is an environment
    for verse \\             % The \\ command separates lines
   Whose features some poets % within a stanza.
   will curse.

                             % One or more blank lines separate stanzas.

   For instead of making\\
   Them do \emph{all} line breaking, \\
   It allows them to put too many words on a line when they'd rather be
   forced to be terse.
\end{verse}

Mathematical formulas may also be displayed.  A
displayed formula
is
one-line long; multiline
formulas require special formatting instructions.
   \[  \ip{\Gamma}{\psi'} = x'' + y^{2} + z_{i}^{n}\]
Don't start a paragraph with a displayed equation,
nor make one a paragraph by itself.

\end{document}               % End of document.

\section{Bench Huge Section 108}
% synthetic bench section 108
% This is a sample LaTeX input file.  (Version of 12 August 2004.)
%
% A '%' character causes TeX to ignore all remaining text on the line,
% and is used for comments like this one.

\documentclass{article}      % Specifies the document class

                             % The preamble begins here.
\title{An Example Document}  % Declares the document's title.
\author{Leslie Lamport}      % Declares the author's name.
\date{January 21, 1994}      % Deleting this command produces today's date.

\newcommand{\ip}[2]{(#1, #2)}
                             % Defines \ip{arg1}{arg2} to mean
                             % (arg1, arg2).

%\newcommand{\ip}[2]{\langle #1 | #2\rangle}
                             % This is an alternative definition of
                             % \ip that is commented out.

\begin{document}             % End of preamble and beginning of text.

\maketitle                   % Produces the title.

This is an example input file.  Comparing it with
the output it generates can show you how to
produce a simple document of your own.

\section{Ordinary Text}      % Produces section heading.  Lower-level
                             % sections are begun with similar
                             % \subsection and \subsubsection commands.

The ends  of words and sentences are marked
  by   spaces. It  doesn't matter how many
spaces    you type; one is as good as 100.  The
end of   a line counts as a space.

One   or more   blank lines denote the  end
of  a paragraph.

Since any number of consecutive spaces are treated
like a single one, the formatting of the input
file makes no difference to
      \LaTeX,                % The \LaTeX command generates the LaTeX logo.
but it makes a difference to you.  When you use
\LaTeX, making your input file as easy to read
as possible will be a great help as you write
your document and when you change it.  This sample
file shows how you can add comments to your own input
file.

Because printing is different from typewriting,
there are a number of things that you have to do
differently when preparing an input file than if
you were just typing the document directly.
Quotation marks like
       ``this''
have to be handled specially, as do quotes within
quotes:
       ``\,`this'            % \, separates the double and single quote.
        is what I just
        wrote, not  `that'\,''.

Dashes come in three sizes: an
       intra-word
dash, a medium dash for number ranges like
       1--2,
and a punctuation
       dash---like
this.

A sentence-ending space should be larger than the
space between words within a sentence.  You
sometimes have to type special commands in
conjunction with punctuation characters to get
this right, as in the following sentence.
       Gnats, gnus, etc.\ all  % `\ ' makes an inter-word space.
       begin with G\@.         % \@ marks end-of-sentence punctuation.
You should check the spaces after periods when
reading your output to make sure you haven't
forgotten any special cases.  Generating an
ellipsis
       \ldots\               % `\ ' is needed after `\ldots' because TeX
                             % ignores spaces after command names like \ldots
                             % made from \ + letters.
                             %
                             % Note how a `%' character causes TeX to ignore
                             % the end of the input line, so these blank lines
                             % do not start a new paragraph.
                             %
with the right spacing around the periods requires
a special command.

\LaTeX\ interprets some common characters as
commands, so you must type special commands to
generate them.  These characters include the
following:
       \$ \& \% \# \{ and \}.

In printing, text is usually emphasized with an
       \emph{italic}
type style.

\begin{em}
   A long segment of text can also be emphasized
   in this way.  Text within such a segment can be
   given \emph{additional} emphasis.
\end{em}

It is sometimes necessary to prevent \LaTeX\ from
breaking a line where it might otherwise do so.
This may be at a space, as between the ``Mr.''\ and
``Jones'' in
       ``Mr.~Jones'',        % ~ produces an unbreakable interword space.
or within a word---especially when the word is a
symbol like
       \mbox{\emph{itemnum}}
that makes little sense when hyphenated across
lines.

Footnotes\footnote{This is an example of a footnote.}
pose no problem.

\LaTeX\ is good at typesetting mathematical formulas
like
       \( x-3y + z = 7 \)
or
       \( a_{1} > x^{2n} + y^{2n} > x' \)
or
       \( \ip{A}{B} = \sum_{i} a_{i} b_{i} \).
The spaces you type in a formula are
ignored.  Remember that a letter like
       $x$                   % $ ... $  and  \( ... \)  are equivalent
is a formula when it denotes a mathematical
symbol, and it should be typed as one.

\section{Displayed Text}

Text is displayed by indenting it from the left
margin.  Quotations are commonly displayed.  There
are short quotations
\begin{quote}
   This is a short quotation.  It consists of a
   single paragraph of text.  See how it is formatted.
\end{quote}
and longer ones.
\begin{quotation}
   This is a longer quotation.  It consists of two
   paragraphs of text, neither of which are
   particularly interesting.

   This is the second paragraph of the quotation.  It
   is just as dull as the first paragraph.
\end{quotation}
Another frequently-displayed structure is a list.
The following is an example of an \emph{itemized}
list.
\begin{itemize}
   \item This is the first item of an itemized list.
         Each item in the list is marked with a ``tick''.
         You don't have to worry about what kind of tick
         mark is used.

   \item This is the second item of the list.  It
         contains another list nested inside it.  The inner
         list is an \emph{enumerated} list.
         \begin{enumerate}
            \item This is the first item of an enumerated
                  list that is nested within the itemized list.

            \item This is the second item of the inner list.
                  \LaTeX\ allows you to nest lists deeper than
                  you really should.
         \end{enumerate}
         This is the rest of the second item of the outer
         list.  It is no more interesting than any other
         part of the item.
   \item This is the third item of the list.
\end{itemize}
You can even display poetry.
\begin{verse}
   There is an environment
    for verse \\             % The \\ command separates lines
   Whose features some poets % within a stanza.
   will curse.

                             % One or more blank lines separate stanzas.

   For instead of making\\
   Them do \emph{all} line breaking, \\
   It allows them to put too many words on a line when they'd rather be
   forced to be terse.
\end{verse}

Mathematical formulas may also be displayed.  A
displayed formula
is
one-line long; multiline
formulas require special formatting instructions.
   \[  \ip{\Gamma}{\psi'} = x'' + y^{2} + z_{i}^{n}\]
Don't start a paragraph with a displayed equation,
nor make one a paragraph by itself.

\end{document}               % End of document.

\section{Bench Huge Section 109}
% synthetic bench section 109
% This is a sample LaTeX input file.  (Version of 12 August 2004.)
%
% A '%' character causes TeX to ignore all remaining text on the line,
% and is used for comments like this one.

\documentclass{article}      % Specifies the document class

                             % The preamble begins here.
\title{An Example Document}  % Declares the document's title.
\author{Leslie Lamport}      % Declares the author's name.
\date{January 21, 1994}      % Deleting this command produces today's date.

\newcommand{\ip}[2]{(#1, #2)}
                             % Defines \ip{arg1}{arg2} to mean
                             % (arg1, arg2).

%\newcommand{\ip}[2]{\langle #1 | #2\rangle}
                             % This is an alternative definition of
                             % \ip that is commented out.

\begin{document}             % End of preamble and beginning of text.

\maketitle                   % Produces the title.

This is an example input file.  Comparing it with
the output it generates can show you how to
produce a simple document of your own.

\section{Ordinary Text}      % Produces section heading.  Lower-level
                             % sections are begun with similar
                             % \subsection and \subsubsection commands.

The ends  of words and sentences are marked
  by   spaces. It  doesn't matter how many
spaces    you type; one is as good as 100.  The
end of   a line counts as a space.

One   or more   blank lines denote the  end
of  a paragraph.

Since any number of consecutive spaces are treated
like a single one, the formatting of the input
file makes no difference to
      \LaTeX,                % The \LaTeX command generates the LaTeX logo.
but it makes a difference to you.  When you use
\LaTeX, making your input file as easy to read
as possible will be a great help as you write
your document and when you change it.  This sample
file shows how you can add comments to your own input
file.

Because printing is different from typewriting,
there are a number of things that you have to do
differently when preparing an input file than if
you were just typing the document directly.
Quotation marks like
       ``this''
have to be handled specially, as do quotes within
quotes:
       ``\,`this'            % \, separates the double and single quote.
        is what I just
        wrote, not  `that'\,''.

Dashes come in three sizes: an
       intra-word
dash, a medium dash for number ranges like
       1--2,
and a punctuation
       dash---like
this.

A sentence-ending space should be larger than the
space between words within a sentence.  You
sometimes have to type special commands in
conjunction with punctuation characters to get
this right, as in the following sentence.
       Gnats, gnus, etc.\ all  % `\ ' makes an inter-word space.
       begin with G\@.         % \@ marks end-of-sentence punctuation.
You should check the spaces after periods when
reading your output to make sure you haven't
forgotten any special cases.  Generating an
ellipsis
       \ldots\               % `\ ' is needed after `\ldots' because TeX
                             % ignores spaces after command names like \ldots
                             % made from \ + letters.
                             %
                             % Note how a `%' character causes TeX to ignore
                             % the end of the input line, so these blank lines
                             % do not start a new paragraph.
                             %
with the right spacing around the periods requires
a special command.

\LaTeX\ interprets some common characters as
commands, so you must type special commands to
generate them.  These characters include the
following:
       \$ \& \% \# \{ and \}.

In printing, text is usually emphasized with an
       \emph{italic}
type style.

\begin{em}
   A long segment of text can also be emphasized
   in this way.  Text within such a segment can be
   given \emph{additional} emphasis.
\end{em}

It is sometimes necessary to prevent \LaTeX\ from
breaking a line where it might otherwise do so.
This may be at a space, as between the ``Mr.''\ and
``Jones'' in
       ``Mr.~Jones'',        % ~ produces an unbreakable interword space.
or within a word---especially when the word is a
symbol like
       \mbox{\emph{itemnum}}
that makes little sense when hyphenated across
lines.

Footnotes\footnote{This is an example of a footnote.}
pose no problem.

\LaTeX\ is good at typesetting mathematical formulas
like
       \( x-3y + z = 7 \)
or
       \( a_{1} > x^{2n} + y^{2n} > x' \)
or
       \( \ip{A}{B} = \sum_{i} a_{i} b_{i} \).
The spaces you type in a formula are
ignored.  Remember that a letter like
       $x$                   % $ ... $  and  \( ... \)  are equivalent
is a formula when it denotes a mathematical
symbol, and it should be typed as one.

\section{Displayed Text}

Text is displayed by indenting it from the left
margin.  Quotations are commonly displayed.  There
are short quotations
\begin{quote}
   This is a short quotation.  It consists of a
   single paragraph of text.  See how it is formatted.
\end{quote}
and longer ones.
\begin{quotation}
   This is a longer quotation.  It consists of two
   paragraphs of text, neither of which are
   particularly interesting.

   This is the second paragraph of the quotation.  It
   is just as dull as the first paragraph.
\end{quotation}
Another frequently-displayed structure is a list.
The following is an example of an \emph{itemized}
list.
\begin{itemize}
   \item This is the first item of an itemized list.
         Each item in the list is marked with a ``tick''.
         You don't have to worry about what kind of tick
         mark is used.

   \item This is the second item of the list.  It
         contains another list nested inside it.  The inner
         list is an \emph{enumerated} list.
         \begin{enumerate}
            \item This is the first item of an enumerated
                  list that is nested within the itemized list.

            \item This is the second item of the inner list.
                  \LaTeX\ allows you to nest lists deeper than
                  you really should.
         \end{enumerate}
         This is the rest of the second item of the outer
         list.  It is no more interesting than any other
         part of the item.
   \item This is the third item of the list.
\end{itemize}
You can even display poetry.
\begin{verse}
   There is an environment
    for verse \\             % The \\ command separates lines
   Whose features some poets % within a stanza.
   will curse.

                             % One or more blank lines separate stanzas.

   For instead of making\\
   Them do \emph{all} line breaking, \\
   It allows them to put too many words on a line when they'd rather be
   forced to be terse.
\end{verse}

Mathematical formulas may also be displayed.  A
displayed formula
is
one-line long; multiline
formulas require special formatting instructions.
   \[  \ip{\Gamma}{\psi'} = x'' + y^{2} + z_{i}^{n}\]
Don't start a paragraph with a displayed equation,
nor make one a paragraph by itself.

\end{document}               % End of document.

\section{Bench Huge Section 110}
% synthetic bench section 110
% This is a sample LaTeX input file.  (Version of 12 August 2004.)
%
% A '%' character causes TeX to ignore all remaining text on the line,
% and is used for comments like this one.

\documentclass{article}      % Specifies the document class

                             % The preamble begins here.
\title{An Example Document}  % Declares the document's title.
\author{Leslie Lamport}      % Declares the author's name.
\date{January 21, 1994}      % Deleting this command produces today's date.

\newcommand{\ip}[2]{(#1, #2)}
                             % Defines \ip{arg1}{arg2} to mean
                             % (arg1, arg2).

%\newcommand{\ip}[2]{\langle #1 | #2\rangle}
                             % This is an alternative definition of
                             % \ip that is commented out.

\begin{document}             % End of preamble and beginning of text.

\maketitle                   % Produces the title.

This is an example input file.  Comparing it with
the output it generates can show you how to
produce a simple document of your own.

\section{Ordinary Text}      % Produces section heading.  Lower-level
                             % sections are begun with similar
                             % \subsection and \subsubsection commands.

The ends  of words and sentences are marked
  by   spaces. It  doesn't matter how many
spaces    you type; one is as good as 100.  The
end of   a line counts as a space.

One   or more   blank lines denote the  end
of  a paragraph.

Since any number of consecutive spaces are treated
like a single one, the formatting of the input
file makes no difference to
      \LaTeX,                % The \LaTeX command generates the LaTeX logo.
but it makes a difference to you.  When you use
\LaTeX, making your input file as easy to read
as possible will be a great help as you write
your document and when you change it.  This sample
file shows how you can add comments to your own input
file.

Because printing is different from typewriting,
there are a number of things that you have to do
differently when preparing an input file than if
you were just typing the document directly.
Quotation marks like
       ``this''
have to be handled specially, as do quotes within
quotes:
       ``\,`this'            % \, separates the double and single quote.
        is what I just
        wrote, not  `that'\,''.

Dashes come in three sizes: an
       intra-word
dash, a medium dash for number ranges like
       1--2,
and a punctuation
       dash---like
this.

A sentence-ending space should be larger than the
space between words within a sentence.  You
sometimes have to type special commands in
conjunction with punctuation characters to get
this right, as in the following sentence.
       Gnats, gnus, etc.\ all  % `\ ' makes an inter-word space.
       begin with G\@.         % \@ marks end-of-sentence punctuation.
You should check the spaces after periods when
reading your output to make sure you haven't
forgotten any special cases.  Generating an
ellipsis
       \ldots\               % `\ ' is needed after `\ldots' because TeX
                             % ignores spaces after command names like \ldots
                             % made from \ + letters.
                             %
                             % Note how a `%' character causes TeX to ignore
                             % the end of the input line, so these blank lines
                             % do not start a new paragraph.
                             %
with the right spacing around the periods requires
a special command.

\LaTeX\ interprets some common characters as
commands, so you must type special commands to
generate them.  These characters include the
following:
       \$ \& \% \# \{ and \}.

In printing, text is usually emphasized with an
       \emph{italic}
type style.

\begin{em}
   A long segment of text can also be emphasized
   in this way.  Text within such a segment can be
   given \emph{additional} emphasis.
\end{em}

It is sometimes necessary to prevent \LaTeX\ from
breaking a line where it might otherwise do so.
This may be at a space, as between the ``Mr.''\ and
``Jones'' in
       ``Mr.~Jones'',        % ~ produces an unbreakable interword space.
or within a word---especially when the word is a
symbol like
       \mbox{\emph{itemnum}}
that makes little sense when hyphenated across
lines.

Footnotes\footnote{This is an example of a footnote.}
pose no problem.

\LaTeX\ is good at typesetting mathematical formulas
like
       \( x-3y + z = 7 \)
or
       \( a_{1} > x^{2n} + y^{2n} > x' \)
or
       \( \ip{A}{B} = \sum_{i} a_{i} b_{i} \).
The spaces you type in a formula are
ignored.  Remember that a letter like
       $x$                   % $ ... $  and  \( ... \)  are equivalent
is a formula when it denotes a mathematical
symbol, and it should be typed as one.

\section{Displayed Text}

Text is displayed by indenting it from the left
margin.  Quotations are commonly displayed.  There
are short quotations
\begin{quote}
   This is a short quotation.  It consists of a
   single paragraph of text.  See how it is formatted.
\end{quote}
and longer ones.
\begin{quotation}
   This is a longer quotation.  It consists of two
   paragraphs of text, neither of which are
   particularly interesting.

   This is the second paragraph of the quotation.  It
   is just as dull as the first paragraph.
\end{quotation}
Another frequently-displayed structure is a list.
The following is an example of an \emph{itemized}
list.
\begin{itemize}
   \item This is the first item of an itemized list.
         Each item in the list is marked with a ``tick''.
         You don't have to worry about what kind of tick
         mark is used.

   \item This is the second item of the list.  It
         contains another list nested inside it.  The inner
         list is an \emph{enumerated} list.
         \begin{enumerate}
            \item This is the first item of an enumerated
                  list that is nested within the itemized list.

            \item This is the second item of the inner list.
                  \LaTeX\ allows you to nest lists deeper than
                  you really should.
         \end{enumerate}
         This is the rest of the second item of the outer
         list.  It is no more interesting than any other
         part of the item.
   \item This is the third item of the list.
\end{itemize}
You can even display poetry.
\begin{verse}
   There is an environment
    for verse \\             % The \\ command separates lines
   Whose features some poets % within a stanza.
   will curse.

                             % One or more blank lines separate stanzas.

   For instead of making\\
   Them do \emph{all} line breaking, \\
   It allows them to put too many words on a line when they'd rather be
   forced to be terse.
\end{verse}

Mathematical formulas may also be displayed.  A
displayed formula
is
one-line long; multiline
formulas require special formatting instructions.
   \[  \ip{\Gamma}{\psi'} = x'' + y^{2} + z_{i}^{n}\]
Don't start a paragraph with a displayed equation,
nor make one a paragraph by itself.

\end{document}               % End of document.

\section{Bench Huge Section 111}
% synthetic bench section 111
% This is a sample LaTeX input file.  (Version of 12 August 2004.)
%
% A '%' character causes TeX to ignore all remaining text on the line,
% and is used for comments like this one.

\documentclass{article}      % Specifies the document class

                             % The preamble begins here.
\title{An Example Document}  % Declares the document's title.
\author{Leslie Lamport}      % Declares the author's name.
\date{January 21, 1994}      % Deleting this command produces today's date.

\newcommand{\ip}[2]{(#1, #2)}
                             % Defines \ip{arg1}{arg2} to mean
                             % (arg1, arg2).

%\newcommand{\ip}[2]{\langle #1 | #2\rangle}
                             % This is an alternative definition of
                             % \ip that is commented out.

\begin{document}             % End of preamble and beginning of text.

\maketitle                   % Produces the title.

This is an example input file.  Comparing it with
the output it generates can show you how to
produce a simple document of your own.

\section{Ordinary Text}      % Produces section heading.  Lower-level
                             % sections are begun with similar
                             % \subsection and \subsubsection commands.

The ends  of words and sentences are marked
  by   spaces. It  doesn't matter how many
spaces    you type; one is as good as 100.  The
end of   a line counts as a space.

One   or more   blank lines denote the  end
of  a paragraph.

Since any number of consecutive spaces are treated
like a single one, the formatting of the input
file makes no difference to
      \LaTeX,                % The \LaTeX command generates the LaTeX logo.
but it makes a difference to you.  When you use
\LaTeX, making your input file as easy to read
as possible will be a great help as you write
your document and when you change it.  This sample
file shows how you can add comments to your own input
file.

Because printing is different from typewriting,
there are a number of things that you have to do
differently when preparing an input file than if
you were just typing the document directly.
Quotation marks like
       ``this''
have to be handled specially, as do quotes within
quotes:
       ``\,`this'            % \, separates the double and single quote.
        is what I just
        wrote, not  `that'\,''.

Dashes come in three sizes: an
       intra-word
dash, a medium dash for number ranges like
       1--2,
and a punctuation
       dash---like
this.

A sentence-ending space should be larger than the
space between words within a sentence.  You
sometimes have to type special commands in
conjunction with punctuation characters to get
this right, as in the following sentence.
       Gnats, gnus, etc.\ all  % `\ ' makes an inter-word space.
       begin with G\@.         % \@ marks end-of-sentence punctuation.
You should check the spaces after periods when
reading your output to make sure you haven't
forgotten any special cases.  Generating an
ellipsis
       \ldots\               % `\ ' is needed after `\ldots' because TeX
                             % ignores spaces after command names like \ldots
                             % made from \ + letters.
                             %
                             % Note how a `%' character causes TeX to ignore
                             % the end of the input line, so these blank lines
                             % do not start a new paragraph.
                             %
with the right spacing around the periods requires
a special command.

\LaTeX\ interprets some common characters as
commands, so you must type special commands to
generate them.  These characters include the
following:
       \$ \& \% \# \{ and \}.

In printing, text is usually emphasized with an
       \emph{italic}
type style.

\begin{em}
   A long segment of text can also be emphasized
   in this way.  Text within such a segment can be
   given \emph{additional} emphasis.
\end{em}

It is sometimes necessary to prevent \LaTeX\ from
breaking a line where it might otherwise do so.
This may be at a space, as between the ``Mr.''\ and
``Jones'' in
       ``Mr.~Jones'',        % ~ produces an unbreakable interword space.
or within a word---especially when the word is a
symbol like
       \mbox{\emph{itemnum}}
that makes little sense when hyphenated across
lines.

Footnotes\footnote{This is an example of a footnote.}
pose no problem.

\LaTeX\ is good at typesetting mathematical formulas
like
       \( x-3y + z = 7 \)
or
       \( a_{1} > x^{2n} + y^{2n} > x' \)
or
       \( \ip{A}{B} = \sum_{i} a_{i} b_{i} \).
The spaces you type in a formula are
ignored.  Remember that a letter like
       $x$                   % $ ... $  and  \( ... \)  are equivalent
is a formula when it denotes a mathematical
symbol, and it should be typed as one.

\section{Displayed Text}

Text is displayed by indenting it from the left
margin.  Quotations are commonly displayed.  There
are short quotations
\begin{quote}
   This is a short quotation.  It consists of a
   single paragraph of text.  See how it is formatted.
\end{quote}
and longer ones.
\begin{quotation}
   This is a longer quotation.  It consists of two
   paragraphs of text, neither of which are
   particularly interesting.

   This is the second paragraph of the quotation.  It
   is just as dull as the first paragraph.
\end{quotation}
Another frequently-displayed structure is a list.
The following is an example of an \emph{itemized}
list.
\begin{itemize}
   \item This is the first item of an itemized list.
         Each item in the list is marked with a ``tick''.
         You don't have to worry about what kind of tick
         mark is used.

   \item This is the second item of the list.  It
         contains another list nested inside it.  The inner
         list is an \emph{enumerated} list.
         \begin{enumerate}
            \item This is the first item of an enumerated
                  list that is nested within the itemized list.

            \item This is the second item of the inner list.
                  \LaTeX\ allows you to nest lists deeper than
                  you really should.
         \end{enumerate}
         This is the rest of the second item of the outer
         list.  It is no more interesting than any other
         part of the item.
   \item This is the third item of the list.
\end{itemize}
You can even display poetry.
\begin{verse}
   There is an environment
    for verse \\             % The \\ command separates lines
   Whose features some poets % within a stanza.
   will curse.

                             % One or more blank lines separate stanzas.

   For instead of making\\
   Them do \emph{all} line breaking, \\
   It allows them to put too many words on a line when they'd rather be
   forced to be terse.
\end{verse}

Mathematical formulas may also be displayed.  A
displayed formula
is
one-line long; multiline
formulas require special formatting instructions.
   \[  \ip{\Gamma}{\psi'} = x'' + y^{2} + z_{i}^{n}\]
Don't start a paragraph with a displayed equation,
nor make one a paragraph by itself.

\end{document}               % End of document.

\section{Bench Huge Section 112}
% synthetic bench section 112
% This is a sample LaTeX input file.  (Version of 12 August 2004.)
%
% A '%' character causes TeX to ignore all remaining text on the line,
% and is used for comments like this one.

\documentclass{article}      % Specifies the document class

                             % The preamble begins here.
\title{An Example Document}  % Declares the document's title.
\author{Leslie Lamport}      % Declares the author's name.
\date{January 21, 1994}      % Deleting this command produces today's date.

\newcommand{\ip}[2]{(#1, #2)}
                             % Defines \ip{arg1}{arg2} to mean
                             % (arg1, arg2).

%\newcommand{\ip}[2]{\langle #1 | #2\rangle}
                             % This is an alternative definition of
                             % \ip that is commented out.

\begin{document}             % End of preamble and beginning of text.

\maketitle                   % Produces the title.

This is an example input file.  Comparing it with
the output it generates can show you how to
produce a simple document of your own.

\section{Ordinary Text}      % Produces section heading.  Lower-level
                             % sections are begun with similar
                             % \subsection and \subsubsection commands.

The ends  of words and sentences are marked
  by   spaces. It  doesn't matter how many
spaces    you type; one is as good as 100.  The
end of   a line counts as a space.

One   or more   blank lines denote the  end
of  a paragraph.

Since any number of consecutive spaces are treated
like a single one, the formatting of the input
file makes no difference to
      \LaTeX,                % The \LaTeX command generates the LaTeX logo.
but it makes a difference to you.  When you use
\LaTeX, making your input file as easy to read
as possible will be a great help as you write
your document and when you change it.  This sample
file shows how you can add comments to your own input
file.

Because printing is different from typewriting,
there are a number of things that you have to do
differently when preparing an input file than if
you were just typing the document directly.
Quotation marks like
       ``this''
have to be handled specially, as do quotes within
quotes:
       ``\,`this'            % \, separates the double and single quote.
        is what I just
        wrote, not  `that'\,''.

Dashes come in three sizes: an
       intra-word
dash, a medium dash for number ranges like
       1--2,
and a punctuation
       dash---like
this.

A sentence-ending space should be larger than the
space between words within a sentence.  You
sometimes have to type special commands in
conjunction with punctuation characters to get
this right, as in the following sentence.
       Gnats, gnus, etc.\ all  % `\ ' makes an inter-word space.
       begin with G\@.         % \@ marks end-of-sentence punctuation.
You should check the spaces after periods when
reading your output to make sure you haven't
forgotten any special cases.  Generating an
ellipsis
       \ldots\               % `\ ' is needed after `\ldots' because TeX
                             % ignores spaces after command names like \ldots
                             % made from \ + letters.
                             %
                             % Note how a `%' character causes TeX to ignore
                             % the end of the input line, so these blank lines
                             % do not start a new paragraph.
                             %
with the right spacing around the periods requires
a special command.

\LaTeX\ interprets some common characters as
commands, so you must type special commands to
generate them.  These characters include the
following:
       \$ \& \% \# \{ and \}.

In printing, text is usually emphasized with an
       \emph{italic}
type style.

\begin{em}
   A long segment of text can also be emphasized
   in this way.  Text within such a segment can be
   given \emph{additional} emphasis.
\end{em}

It is sometimes necessary to prevent \LaTeX\ from
breaking a line where it might otherwise do so.
This may be at a space, as between the ``Mr.''\ and
``Jones'' in
       ``Mr.~Jones'',        % ~ produces an unbreakable interword space.
or within a word---especially when the word is a
symbol like
       \mbox{\emph{itemnum}}
that makes little sense when hyphenated across
lines.

Footnotes\footnote{This is an example of a footnote.}
pose no problem.

\LaTeX\ is good at typesetting mathematical formulas
like
       \( x-3y + z = 7 \)
or
       \( a_{1} > x^{2n} + y^{2n} > x' \)
or
       \( \ip{A}{B} = \sum_{i} a_{i} b_{i} \).
The spaces you type in a formula are
ignored.  Remember that a letter like
       $x$                   % $ ... $  and  \( ... \)  are equivalent
is a formula when it denotes a mathematical
symbol, and it should be typed as one.

\section{Displayed Text}

Text is displayed by indenting it from the left
margin.  Quotations are commonly displayed.  There
are short quotations
\begin{quote}
   This is a short quotation.  It consists of a
   single paragraph of text.  See how it is formatted.
\end{quote}
and longer ones.
\begin{quotation}
   This is a longer quotation.  It consists of two
   paragraphs of text, neither of which are
   particularly interesting.

   This is the second paragraph of the quotation.  It
   is just as dull as the first paragraph.
\end{quotation}
Another frequently-displayed structure is a list.
The following is an example of an \emph{itemized}
list.
\begin{itemize}
   \item This is the first item of an itemized list.
         Each item in the list is marked with a ``tick''.
         You don't have to worry about what kind of tick
         mark is used.

   \item This is the second item of the list.  It
         contains another list nested inside it.  The inner
         list is an \emph{enumerated} list.
         \begin{enumerate}
            \item This is the first item of an enumerated
                  list that is nested within the itemized list.

            \item This is the second item of the inner list.
                  \LaTeX\ allows you to nest lists deeper than
                  you really should.
         \end{enumerate}
         This is the rest of the second item of the outer
         list.  It is no more interesting than any other
         part of the item.
   \item This is the third item of the list.
\end{itemize}
You can even display poetry.
\begin{verse}
   There is an environment
    for verse \\             % The \\ command separates lines
   Whose features some poets % within a stanza.
   will curse.

                             % One or more blank lines separate stanzas.

   For instead of making\\
   Them do \emph{all} line breaking, \\
   It allows them to put too many words on a line when they'd rather be
   forced to be terse.
\end{verse}

Mathematical formulas may also be displayed.  A
displayed formula
is
one-line long; multiline
formulas require special formatting instructions.
   \[  \ip{\Gamma}{\psi'} = x'' + y^{2} + z_{i}^{n}\]
Don't start a paragraph with a displayed equation,
nor make one a paragraph by itself.

\end{document}               % End of document.

\section{Bench Huge Section 113}
% synthetic bench section 113
% This is a sample LaTeX input file.  (Version of 12 August 2004.)
%
% A '%' character causes TeX to ignore all remaining text on the line,
% and is used for comments like this one.

\documentclass{article}      % Specifies the document class

                             % The preamble begins here.
\title{An Example Document}  % Declares the document's title.
\author{Leslie Lamport}      % Declares the author's name.
\date{January 21, 1994}      % Deleting this command produces today's date.

\newcommand{\ip}[2]{(#1, #2)}
                             % Defines \ip{arg1}{arg2} to mean
                             % (arg1, arg2).

%\newcommand{\ip}[2]{\langle #1 | #2\rangle}
                             % This is an alternative definition of
                             % \ip that is commented out.

\begin{document}             % End of preamble and beginning of text.

\maketitle                   % Produces the title.

This is an example input file.  Comparing it with
the output it generates can show you how to
produce a simple document of your own.

\section{Ordinary Text}      % Produces section heading.  Lower-level
                             % sections are begun with similar
                             % \subsection and \subsubsection commands.

The ends  of words and sentences are marked
  by   spaces. It  doesn't matter how many
spaces    you type; one is as good as 100.  The
end of   a line counts as a space.

One   or more   blank lines denote the  end
of  a paragraph.

Since any number of consecutive spaces are treated
like a single one, the formatting of the input
file makes no difference to
      \LaTeX,                % The \LaTeX command generates the LaTeX logo.
but it makes a difference to you.  When you use
\LaTeX, making your input file as easy to read
as possible will be a great help as you write
your document and when you change it.  This sample
file shows how you can add comments to your own input
file.

Because printing is different from typewriting,
there are a number of things that you have to do
differently when preparing an input file than if
you were just typing the document directly.
Quotation marks like
       ``this''
have to be handled specially, as do quotes within
quotes:
       ``\,`this'            % \, separates the double and single quote.
        is what I just
        wrote, not  `that'\,''.

Dashes come in three sizes: an
       intra-word
dash, a medium dash for number ranges like
       1--2,
and a punctuation
       dash---like
this.

A sentence-ending space should be larger than the
space between words within a sentence.  You
sometimes have to type special commands in
conjunction with punctuation characters to get
this right, as in the following sentence.
       Gnats, gnus, etc.\ all  % `\ ' makes an inter-word space.
       begin with G\@.         % \@ marks end-of-sentence punctuation.
You should check the spaces after periods when
reading your output to make sure you haven't
forgotten any special cases.  Generating an
ellipsis
       \ldots\               % `\ ' is needed after `\ldots' because TeX
                             % ignores spaces after command names like \ldots
                             % made from \ + letters.
                             %
                             % Note how a `%' character causes TeX to ignore
                             % the end of the input line, so these blank lines
                             % do not start a new paragraph.
                             %
with the right spacing around the periods requires
a special command.

\LaTeX\ interprets some common characters as
commands, so you must type special commands to
generate them.  These characters include the
following:
       \$ \& \% \# \{ and \}.

In printing, text is usually emphasized with an
       \emph{italic}
type style.

\begin{em}
   A long segment of text can also be emphasized
   in this way.  Text within such a segment can be
   given \emph{additional} emphasis.
\end{em}

It is sometimes necessary to prevent \LaTeX\ from
breaking a line where it might otherwise do so.
This may be at a space, as between the ``Mr.''\ and
``Jones'' in
       ``Mr.~Jones'',        % ~ produces an unbreakable interword space.
or within a word---especially when the word is a
symbol like
       \mbox{\emph{itemnum}}
that makes little sense when hyphenated across
lines.

Footnotes\footnote{This is an example of a footnote.}
pose no problem.

\LaTeX\ is good at typesetting mathematical formulas
like
       \( x-3y + z = 7 \)
or
       \( a_{1} > x^{2n} + y^{2n} > x' \)
or
       \( \ip{A}{B} = \sum_{i} a_{i} b_{i} \).
The spaces you type in a formula are
ignored.  Remember that a letter like
       $x$                   % $ ... $  and  \( ... \)  are equivalent
is a formula when it denotes a mathematical
symbol, and it should be typed as one.

\section{Displayed Text}

Text is displayed by indenting it from the left
margin.  Quotations are commonly displayed.  There
are short quotations
\begin{quote}
   This is a short quotation.  It consists of a
   single paragraph of text.  See how it is formatted.
\end{quote}
and longer ones.
\begin{quotation}
   This is a longer quotation.  It consists of two
   paragraphs of text, neither of which are
   particularly interesting.

   This is the second paragraph of the quotation.  It
   is just as dull as the first paragraph.
\end{quotation}
Another frequently-displayed structure is a list.
The following is an example of an \emph{itemized}
list.
\begin{itemize}
   \item This is the first item of an itemized list.
         Each item in the list is marked with a ``tick''.
         You don't have to worry about what kind of tick
         mark is used.

   \item This is the second item of the list.  It
         contains another list nested inside it.  The inner
         list is an \emph{enumerated} list.
         \begin{enumerate}
            \item This is the first item of an enumerated
                  list that is nested within the itemized list.

            \item This is the second item of the inner list.
                  \LaTeX\ allows you to nest lists deeper than
                  you really should.
         \end{enumerate}
         This is the rest of the second item of the outer
         list.  It is no more interesting than any other
         part of the item.
   \item This is the third item of the list.
\end{itemize}
You can even display poetry.
\begin{verse}
   There is an environment
    for verse \\             % The \\ command separates lines
   Whose features some poets % within a stanza.
   will curse.

                             % One or more blank lines separate stanzas.

   For instead of making\\
   Them do \emph{all} line breaking, \\
   It allows them to put too many words on a line when they'd rather be
   forced to be terse.
\end{verse}

Mathematical formulas may also be displayed.  A
displayed formula
is
one-line long; multiline
formulas require special formatting instructions.
   \[  \ip{\Gamma}{\psi'} = x'' + y^{2} + z_{i}^{n}\]
Don't start a paragraph with a displayed equation,
nor make one a paragraph by itself.

\end{document}               % End of document.

\section{Bench Huge Section 114}
% synthetic bench section 114
% This is a sample LaTeX input file.  (Version of 12 August 2004.)
%
% A '%' character causes TeX to ignore all remaining text on the line,
% and is used for comments like this one.

\documentclass{article}      % Specifies the document class

                             % The preamble begins here.
\title{An Example Document}  % Declares the document's title.
\author{Leslie Lamport}      % Declares the author's name.
\date{January 21, 1994}      % Deleting this command produces today's date.

\newcommand{\ip}[2]{(#1, #2)}
                             % Defines \ip{arg1}{arg2} to mean
                             % (arg1, arg2).

%\newcommand{\ip}[2]{\langle #1 | #2\rangle}
                             % This is an alternative definition of
                             % \ip that is commented out.

\begin{document}             % End of preamble and beginning of text.

\maketitle                   % Produces the title.

This is an example input file.  Comparing it with
the output it generates can show you how to
produce a simple document of your own.

\section{Ordinary Text}      % Produces section heading.  Lower-level
                             % sections are begun with similar
                             % \subsection and \subsubsection commands.

The ends  of words and sentences are marked
  by   spaces. It  doesn't matter how many
spaces    you type; one is as good as 100.  The
end of   a line counts as a space.

One   or more   blank lines denote the  end
of  a paragraph.

Since any number of consecutive spaces are treated
like a single one, the formatting of the input
file makes no difference to
      \LaTeX,                % The \LaTeX command generates the LaTeX logo.
but it makes a difference to you.  When you use
\LaTeX, making your input file as easy to read
as possible will be a great help as you write
your document and when you change it.  This sample
file shows how you can add comments to your own input
file.

Because printing is different from typewriting,
there are a number of things that you have to do
differently when preparing an input file than if
you were just typing the document directly.
Quotation marks like
       ``this''
have to be handled specially, as do quotes within
quotes:
       ``\,`this'            % \, separates the double and single quote.
        is what I just
        wrote, not  `that'\,''.

Dashes come in three sizes: an
       intra-word
dash, a medium dash for number ranges like
       1--2,
and a punctuation
       dash---like
this.

A sentence-ending space should be larger than the
space between words within a sentence.  You
sometimes have to type special commands in
conjunction with punctuation characters to get
this right, as in the following sentence.
       Gnats, gnus, etc.\ all  % `\ ' makes an inter-word space.
       begin with G\@.         % \@ marks end-of-sentence punctuation.
You should check the spaces after periods when
reading your output to make sure you haven't
forgotten any special cases.  Generating an
ellipsis
       \ldots\               % `\ ' is needed after `\ldots' because TeX
                             % ignores spaces after command names like \ldots
                             % made from \ + letters.
                             %
                             % Note how a `%' character causes TeX to ignore
                             % the end of the input line, so these blank lines
                             % do not start a new paragraph.
                             %
with the right spacing around the periods requires
a special command.

\LaTeX\ interprets some common characters as
commands, so you must type special commands to
generate them.  These characters include the
following:
       \$ \& \% \# \{ and \}.

In printing, text is usually emphasized with an
       \emph{italic}
type style.

\begin{em}
   A long segment of text can also be emphasized
   in this way.  Text within such a segment can be
   given \emph{additional} emphasis.
\end{em}

It is sometimes necessary to prevent \LaTeX\ from
breaking a line where it might otherwise do so.
This may be at a space, as between the ``Mr.''\ and
``Jones'' in
       ``Mr.~Jones'',        % ~ produces an unbreakable interword space.
or within a word---especially when the word is a
symbol like
       \mbox{\emph{itemnum}}
that makes little sense when hyphenated across
lines.

Footnotes\footnote{This is an example of a footnote.}
pose no problem.

\LaTeX\ is good at typesetting mathematical formulas
like
       \( x-3y + z = 7 \)
or
       \( a_{1} > x^{2n} + y^{2n} > x' \)
or
       \( \ip{A}{B} = \sum_{i} a_{i} b_{i} \).
The spaces you type in a formula are
ignored.  Remember that a letter like
       $x$                   % $ ... $  and  \( ... \)  are equivalent
is a formula when it denotes a mathematical
symbol, and it should be typed as one.

\section{Displayed Text}

Text is displayed by indenting it from the left
margin.  Quotations are commonly displayed.  There
are short quotations
\begin{quote}
   This is a short quotation.  It consists of a
   single paragraph of text.  See how it is formatted.
\end{quote}
and longer ones.
\begin{quotation}
   This is a longer quotation.  It consists of two
   paragraphs of text, neither of which are
   particularly interesting.

   This is the second paragraph of the quotation.  It
   is just as dull as the first paragraph.
\end{quotation}
Another frequently-displayed structure is a list.
The following is an example of an \emph{itemized}
list.
\begin{itemize}
   \item This is the first item of an itemized list.
         Each item in the list is marked with a ``tick''.
         You don't have to worry about what kind of tick
         mark is used.

   \item This is the second item of the list.  It
         contains another list nested inside it.  The inner
         list is an \emph{enumerated} list.
         \begin{enumerate}
            \item This is the first item of an enumerated
                  list that is nested within the itemized list.

            \item This is the second item of the inner list.
                  \LaTeX\ allows you to nest lists deeper than
                  you really should.
         \end{enumerate}
         This is the rest of the second item of the outer
         list.  It is no more interesting than any other
         part of the item.
   \item This is the third item of the list.
\end{itemize}
You can even display poetry.
\begin{verse}
   There is an environment
    for verse \\             % The \\ command separates lines
   Whose features some poets % within a stanza.
   will curse.

                             % One or more blank lines separate stanzas.

   For instead of making\\
   Them do \emph{all} line breaking, \\
   It allows them to put too many words on a line when they'd rather be
   forced to be terse.
\end{verse}

Mathematical formulas may also be displayed.  A
displayed formula
is
one-line long; multiline
formulas require special formatting instructions.
   \[  \ip{\Gamma}{\psi'} = x'' + y^{2} + z_{i}^{n}\]
Don't start a paragraph with a displayed equation,
nor make one a paragraph by itself.

\end{document}               % End of document.

\section{Bench Huge Section 115}
% synthetic bench section 115
% This is a sample LaTeX input file.  (Version of 12 August 2004.)
%
% A '%' character causes TeX to ignore all remaining text on the line,
% and is used for comments like this one.

\documentclass{article}      % Specifies the document class

                             % The preamble begins here.
\title{An Example Document}  % Declares the document's title.
\author{Leslie Lamport}      % Declares the author's name.
\date{January 21, 1994}      % Deleting this command produces today's date.

\newcommand{\ip}[2]{(#1, #2)}
                             % Defines \ip{arg1}{arg2} to mean
                             % (arg1, arg2).

%\newcommand{\ip}[2]{\langle #1 | #2\rangle}
                             % This is an alternative definition of
                             % \ip that is commented out.

\begin{document}             % End of preamble and beginning of text.

\maketitle                   % Produces the title.

This is an example input file.  Comparing it with
the output it generates can show you how to
produce a simple document of your own.

\section{Ordinary Text}      % Produces section heading.  Lower-level
                             % sections are begun with similar
                             % \subsection and \subsubsection commands.

The ends  of words and sentences are marked
  by   spaces. It  doesn't matter how many
spaces    you type; one is as good as 100.  The
end of   a line counts as a space.

One   or more   blank lines denote the  end
of  a paragraph.

Since any number of consecutive spaces are treated
like a single one, the formatting of the input
file makes no difference to
      \LaTeX,                % The \LaTeX command generates the LaTeX logo.
but it makes a difference to you.  When you use
\LaTeX, making your input file as easy to read
as possible will be a great help as you write
your document and when you change it.  This sample
file shows how you can add comments to your own input
file.

Because printing is different from typewriting,
there are a number of things that you have to do
differently when preparing an input file than if
you were just typing the document directly.
Quotation marks like
       ``this''
have to be handled specially, as do quotes within
quotes:
       ``\,`this'            % \, separates the double and single quote.
        is what I just
        wrote, not  `that'\,''.

Dashes come in three sizes: an
       intra-word
dash, a medium dash for number ranges like
       1--2,
and a punctuation
       dash---like
this.

A sentence-ending space should be larger than the
space between words within a sentence.  You
sometimes have to type special commands in
conjunction with punctuation characters to get
this right, as in the following sentence.
       Gnats, gnus, etc.\ all  % `\ ' makes an inter-word space.
       begin with G\@.         % \@ marks end-of-sentence punctuation.
You should check the spaces after periods when
reading your output to make sure you haven't
forgotten any special cases.  Generating an
ellipsis
       \ldots\               % `\ ' is needed after `\ldots' because TeX
                             % ignores spaces after command names like \ldots
                             % made from \ + letters.
                             %
                             % Note how a `%' character causes TeX to ignore
                             % the end of the input line, so these blank lines
                             % do not start a new paragraph.
                             %
with the right spacing around the periods requires
a special command.

\LaTeX\ interprets some common characters as
commands, so you must type special commands to
generate them.  These characters include the
following:
       \$ \& \% \# \{ and \}.

In printing, text is usually emphasized with an
       \emph{italic}
type style.

\begin{em}
   A long segment of text can also be emphasized
   in this way.  Text within such a segment can be
   given \emph{additional} emphasis.
\end{em}

It is sometimes necessary to prevent \LaTeX\ from
breaking a line where it might otherwise do so.
This may be at a space, as between the ``Mr.''\ and
``Jones'' in
       ``Mr.~Jones'',        % ~ produces an unbreakable interword space.
or within a word---especially when the word is a
symbol like
       \mbox{\emph{itemnum}}
that makes little sense when hyphenated across
lines.

Footnotes\footnote{This is an example of a footnote.}
pose no problem.

\LaTeX\ is good at typesetting mathematical formulas
like
       \( x-3y + z = 7 \)
or
       \( a_{1} > x^{2n} + y^{2n} > x' \)
or
       \( \ip{A}{B} = \sum_{i} a_{i} b_{i} \).
The spaces you type in a formula are
ignored.  Remember that a letter like
       $x$                   % $ ... $  and  \( ... \)  are equivalent
is a formula when it denotes a mathematical
symbol, and it should be typed as one.

\section{Displayed Text}

Text is displayed by indenting it from the left
margin.  Quotations are commonly displayed.  There
are short quotations
\begin{quote}
   This is a short quotation.  It consists of a
   single paragraph of text.  See how it is formatted.
\end{quote}
and longer ones.
\begin{quotation}
   This is a longer quotation.  It consists of two
   paragraphs of text, neither of which are
   particularly interesting.

   This is the second paragraph of the quotation.  It
   is just as dull as the first paragraph.
\end{quotation}
Another frequently-displayed structure is a list.
The following is an example of an \emph{itemized}
list.
\begin{itemize}
   \item This is the first item of an itemized list.
         Each item in the list is marked with a ``tick''.
         You don't have to worry about what kind of tick
         mark is used.

   \item This is the second item of the list.  It
         contains another list nested inside it.  The inner
         list is an \emph{enumerated} list.
         \begin{enumerate}
            \item This is the first item of an enumerated
                  list that is nested within the itemized list.

            \item This is the second item of the inner list.
                  \LaTeX\ allows you to nest lists deeper than
                  you really should.
         \end{enumerate}
         This is the rest of the second item of the outer
         list.  It is no more interesting than any other
         part of the item.
   \item This is the third item of the list.
\end{itemize}
You can even display poetry.
\begin{verse}
   There is an environment
    for verse \\             % The \\ command separates lines
   Whose features some poets % within a stanza.
   will curse.

                             % One or more blank lines separate stanzas.

   For instead of making\\
   Them do \emph{all} line breaking, \\
   It allows them to put too many words on a line when they'd rather be
   forced to be terse.
\end{verse}

Mathematical formulas may also be displayed.  A
displayed formula
is
one-line long; multiline
formulas require special formatting instructions.
   \[  \ip{\Gamma}{\psi'} = x'' + y^{2} + z_{i}^{n}\]
Don't start a paragraph with a displayed equation,
nor make one a paragraph by itself.

\end{document}               % End of document.

\section{Bench Huge Section 116}
% synthetic bench section 116
% This is a sample LaTeX input file.  (Version of 12 August 2004.)
%
% A '%' character causes TeX to ignore all remaining text on the line,
% and is used for comments like this one.

\documentclass{article}      % Specifies the document class

                             % The preamble begins here.
\title{An Example Document}  % Declares the document's title.
\author{Leslie Lamport}      % Declares the author's name.
\date{January 21, 1994}      % Deleting this command produces today's date.

\newcommand{\ip}[2]{(#1, #2)}
                             % Defines \ip{arg1}{arg2} to mean
                             % (arg1, arg2).

%\newcommand{\ip}[2]{\langle #1 | #2\rangle}
                             % This is an alternative definition of
                             % \ip that is commented out.

\begin{document}             % End of preamble and beginning of text.

\maketitle                   % Produces the title.

This is an example input file.  Comparing it with
the output it generates can show you how to
produce a simple document of your own.

\section{Ordinary Text}      % Produces section heading.  Lower-level
                             % sections are begun with similar
                             % \subsection and \subsubsection commands.

The ends  of words and sentences are marked
  by   spaces. It  doesn't matter how many
spaces    you type; one is as good as 100.  The
end of   a line counts as a space.

One   or more   blank lines denote the  end
of  a paragraph.

Since any number of consecutive spaces are treated
like a single one, the formatting of the input
file makes no difference to
      \LaTeX,                % The \LaTeX command generates the LaTeX logo.
but it makes a difference to you.  When you use
\LaTeX, making your input file as easy to read
as possible will be a great help as you write
your document and when you change it.  This sample
file shows how you can add comments to your own input
file.

Because printing is different from typewriting,
there are a number of things that you have to do
differently when preparing an input file than if
you were just typing the document directly.
Quotation marks like
       ``this''
have to be handled specially, as do quotes within
quotes:
       ``\,`this'            % \, separates the double and single quote.
        is what I just
        wrote, not  `that'\,''.

Dashes come in three sizes: an
       intra-word
dash, a medium dash for number ranges like
       1--2,
and a punctuation
       dash---like
this.

A sentence-ending space should be larger than the
space between words within a sentence.  You
sometimes have to type special commands in
conjunction with punctuation characters to get
this right, as in the following sentence.
       Gnats, gnus, etc.\ all  % `\ ' makes an inter-word space.
       begin with G\@.         % \@ marks end-of-sentence punctuation.
You should check the spaces after periods when
reading your output to make sure you haven't
forgotten any special cases.  Generating an
ellipsis
       \ldots\               % `\ ' is needed after `\ldots' because TeX
                             % ignores spaces after command names like \ldots
                             % made from \ + letters.
                             %
                             % Note how a `%' character causes TeX to ignore
                             % the end of the input line, so these blank lines
                             % do not start a new paragraph.
                             %
with the right spacing around the periods requires
a special command.

\LaTeX\ interprets some common characters as
commands, so you must type special commands to
generate them.  These characters include the
following:
       \$ \& \% \# \{ and \}.

In printing, text is usually emphasized with an
       \emph{italic}
type style.

\begin{em}
   A long segment of text can also be emphasized
   in this way.  Text within such a segment can be
   given \emph{additional} emphasis.
\end{em}

It is sometimes necessary to prevent \LaTeX\ from
breaking a line where it might otherwise do so.
This may be at a space, as between the ``Mr.''\ and
``Jones'' in
       ``Mr.~Jones'',        % ~ produces an unbreakable interword space.
or within a word---especially when the word is a
symbol like
       \mbox{\emph{itemnum}}
that makes little sense when hyphenated across
lines.

Footnotes\footnote{This is an example of a footnote.}
pose no problem.

\LaTeX\ is good at typesetting mathematical formulas
like
       \( x-3y + z = 7 \)
or
       \( a_{1} > x^{2n} + y^{2n} > x' \)
or
       \( \ip{A}{B} = \sum_{i} a_{i} b_{i} \).
The spaces you type in a formula are
ignored.  Remember that a letter like
       $x$                   % $ ... $  and  \( ... \)  are equivalent
is a formula when it denotes a mathematical
symbol, and it should be typed as one.

\section{Displayed Text}

Text is displayed by indenting it from the left
margin.  Quotations are commonly displayed.  There
are short quotations
\begin{quote}
   This is a short quotation.  It consists of a
   single paragraph of text.  See how it is formatted.
\end{quote}
and longer ones.
\begin{quotation}
   This is a longer quotation.  It consists of two
   paragraphs of text, neither of which are
   particularly interesting.

   This is the second paragraph of the quotation.  It
   is just as dull as the first paragraph.
\end{quotation}
Another frequently-displayed structure is a list.
The following is an example of an \emph{itemized}
list.
\begin{itemize}
   \item This is the first item of an itemized list.
         Each item in the list is marked with a ``tick''.
         You don't have to worry about what kind of tick
         mark is used.

   \item This is the second item of the list.  It
         contains another list nested inside it.  The inner
         list is an \emph{enumerated} list.
         \begin{enumerate}
            \item This is the first item of an enumerated
                  list that is nested within the itemized list.

            \item This is the second item of the inner list.
                  \LaTeX\ allows you to nest lists deeper than
                  you really should.
         \end{enumerate}
         This is the rest of the second item of the outer
         list.  It is no more interesting than any other
         part of the item.
   \item This is the third item of the list.
\end{itemize}
You can even display poetry.
\begin{verse}
   There is an environment
    for verse \\             % The \\ command separates lines
   Whose features some poets % within a stanza.
   will curse.

                             % One or more blank lines separate stanzas.

   For instead of making\\
   Them do \emph{all} line breaking, \\
   It allows them to put too many words on a line when they'd rather be
   forced to be terse.
\end{verse}

Mathematical formulas may also be displayed.  A
displayed formula
is
one-line long; multiline
formulas require special formatting instructions.
   \[  \ip{\Gamma}{\psi'} = x'' + y^{2} + z_{i}^{n}\]
Don't start a paragraph with a displayed equation,
nor make one a paragraph by itself.

\end{document}               % End of document.

\section{Bench Huge Section 117}
% synthetic bench section 117
% This is a sample LaTeX input file.  (Version of 12 August 2004.)
%
% A '%' character causes TeX to ignore all remaining text on the line,
% and is used for comments like this one.

\documentclass{article}      % Specifies the document class

                             % The preamble begins here.
\title{An Example Document}  % Declares the document's title.
\author{Leslie Lamport}      % Declares the author's name.
\date{January 21, 1994}      % Deleting this command produces today's date.

\newcommand{\ip}[2]{(#1, #2)}
                             % Defines \ip{arg1}{arg2} to mean
                             % (arg1, arg2).

%\newcommand{\ip}[2]{\langle #1 | #2\rangle}
                             % This is an alternative definition of
                             % \ip that is commented out.

\begin{document}             % End of preamble and beginning of text.

\maketitle                   % Produces the title.

This is an example input file.  Comparing it with
the output it generates can show you how to
produce a simple document of your own.

\section{Ordinary Text}      % Produces section heading.  Lower-level
                             % sections are begun with similar
                             % \subsection and \subsubsection commands.

The ends  of words and sentences are marked
  by   spaces. It  doesn't matter how many
spaces    you type; one is as good as 100.  The
end of   a line counts as a space.

One   or more   blank lines denote the  end
of  a paragraph.

Since any number of consecutive spaces are treated
like a single one, the formatting of the input
file makes no difference to
      \LaTeX,                % The \LaTeX command generates the LaTeX logo.
but it makes a difference to you.  When you use
\LaTeX, making your input file as easy to read
as possible will be a great help as you write
your document and when you change it.  This sample
file shows how you can add comments to your own input
file.

Because printing is different from typewriting,
there are a number of things that you have to do
differently when preparing an input file than if
you were just typing the document directly.
Quotation marks like
       ``this''
have to be handled specially, as do quotes within
quotes:
       ``\,`this'            % \, separates the double and single quote.
        is what I just
        wrote, not  `that'\,''.

Dashes come in three sizes: an
       intra-word
dash, a medium dash for number ranges like
       1--2,
and a punctuation
       dash---like
this.

A sentence-ending space should be larger than the
space between words within a sentence.  You
sometimes have to type special commands in
conjunction with punctuation characters to get
this right, as in the following sentence.
       Gnats, gnus, etc.\ all  % `\ ' makes an inter-word space.
       begin with G\@.         % \@ marks end-of-sentence punctuation.
You should check the spaces after periods when
reading your output to make sure you haven't
forgotten any special cases.  Generating an
ellipsis
       \ldots\               % `\ ' is needed after `\ldots' because TeX
                             % ignores spaces after command names like \ldots
                             % made from \ + letters.
                             %
                             % Note how a `%' character causes TeX to ignore
                             % the end of the input line, so these blank lines
                             % do not start a new paragraph.
                             %
with the right spacing around the periods requires
a special command.

\LaTeX\ interprets some common characters as
commands, so you must type special commands to
generate them.  These characters include the
following:
       \$ \& \% \# \{ and \}.

In printing, text is usually emphasized with an
       \emph{italic}
type style.

\begin{em}
   A long segment of text can also be emphasized
   in this way.  Text within such a segment can be
   given \emph{additional} emphasis.
\end{em}

It is sometimes necessary to prevent \LaTeX\ from
breaking a line where it might otherwise do so.
This may be at a space, as between the ``Mr.''\ and
``Jones'' in
       ``Mr.~Jones'',        % ~ produces an unbreakable interword space.
or within a word---especially when the word is a
symbol like
       \mbox{\emph{itemnum}}
that makes little sense when hyphenated across
lines.

Footnotes\footnote{This is an example of a footnote.}
pose no problem.

\LaTeX\ is good at typesetting mathematical formulas
like
       \( x-3y + z = 7 \)
or
       \( a_{1} > x^{2n} + y^{2n} > x' \)
or
       \( \ip{A}{B} = \sum_{i} a_{i} b_{i} \).
The spaces you type in a formula are
ignored.  Remember that a letter like
       $x$                   % $ ... $  and  \( ... \)  are equivalent
is a formula when it denotes a mathematical
symbol, and it should be typed as one.

\section{Displayed Text}

Text is displayed by indenting it from the left
margin.  Quotations are commonly displayed.  There
are short quotations
\begin{quote}
   This is a short quotation.  It consists of a
   single paragraph of text.  See how it is formatted.
\end{quote}
and longer ones.
\begin{quotation}
   This is a longer quotation.  It consists of two
   paragraphs of text, neither of which are
   particularly interesting.

   This is the second paragraph of the quotation.  It
   is just as dull as the first paragraph.
\end{quotation}
Another frequently-displayed structure is a list.
The following is an example of an \emph{itemized}
list.
\begin{itemize}
   \item This is the first item of an itemized list.
         Each item in the list is marked with a ``tick''.
         You don't have to worry about what kind of tick
         mark is used.

   \item This is the second item of the list.  It
         contains another list nested inside it.  The inner
         list is an \emph{enumerated} list.
         \begin{enumerate}
            \item This is the first item of an enumerated
                  list that is nested within the itemized list.

            \item This is the second item of the inner list.
                  \LaTeX\ allows you to nest lists deeper than
                  you really should.
         \end{enumerate}
         This is the rest of the second item of the outer
         list.  It is no more interesting than any other
         part of the item.
   \item This is the third item of the list.
\end{itemize}
You can even display poetry.
\begin{verse}
   There is an environment
    for verse \\             % The \\ command separates lines
   Whose features some poets % within a stanza.
   will curse.

                             % One or more blank lines separate stanzas.

   For instead of making\\
   Them do \emph{all} line breaking, \\
   It allows them to put too many words on a line when they'd rather be
   forced to be terse.
\end{verse}

Mathematical formulas may also be displayed.  A
displayed formula
is
one-line long; multiline
formulas require special formatting instructions.
   \[  \ip{\Gamma}{\psi'} = x'' + y^{2} + z_{i}^{n}\]
Don't start a paragraph with a displayed equation,
nor make one a paragraph by itself.

\end{document}               % End of document.

\section{Bench Huge Section 118}
% synthetic bench section 118
% This is a sample LaTeX input file.  (Version of 12 August 2004.)
%
% A '%' character causes TeX to ignore all remaining text on the line,
% and is used for comments like this one.

\documentclass{article}      % Specifies the document class

                             % The preamble begins here.
\title{An Example Document}  % Declares the document's title.
\author{Leslie Lamport}      % Declares the author's name.
\date{January 21, 1994}      % Deleting this command produces today's date.

\newcommand{\ip}[2]{(#1, #2)}
                             % Defines \ip{arg1}{arg2} to mean
                             % (arg1, arg2).

%\newcommand{\ip}[2]{\langle #1 | #2\rangle}
                             % This is an alternative definition of
                             % \ip that is commented out.

\begin{document}             % End of preamble and beginning of text.

\maketitle                   % Produces the title.

This is an example input file.  Comparing it with
the output it generates can show you how to
produce a simple document of your own.

\section{Ordinary Text}      % Produces section heading.  Lower-level
                             % sections are begun with similar
                             % \subsection and \subsubsection commands.

The ends  of words and sentences are marked
  by   spaces. It  doesn't matter how many
spaces    you type; one is as good as 100.  The
end of   a line counts as a space.

One   or more   blank lines denote the  end
of  a paragraph.

Since any number of consecutive spaces are treated
like a single one, the formatting of the input
file makes no difference to
      \LaTeX,                % The \LaTeX command generates the LaTeX logo.
but it makes a difference to you.  When you use
\LaTeX, making your input file as easy to read
as possible will be a great help as you write
your document and when you change it.  This sample
file shows how you can add comments to your own input
file.

Because printing is different from typewriting,
there are a number of things that you have to do
differently when preparing an input file than if
you were just typing the document directly.
Quotation marks like
       ``this''
have to be handled specially, as do quotes within
quotes:
       ``\,`this'            % \, separates the double and single quote.
        is what I just
        wrote, not  `that'\,''.

Dashes come in three sizes: an
       intra-word
dash, a medium dash for number ranges like
       1--2,
and a punctuation
       dash---like
this.

A sentence-ending space should be larger than the
space between words within a sentence.  You
sometimes have to type special commands in
conjunction with punctuation characters to get
this right, as in the following sentence.
       Gnats, gnus, etc.\ all  % `\ ' makes an inter-word space.
       begin with G\@.         % \@ marks end-of-sentence punctuation.
You should check the spaces after periods when
reading your output to make sure you haven't
forgotten any special cases.  Generating an
ellipsis
       \ldots\               % `\ ' is needed after `\ldots' because TeX
                             % ignores spaces after command names like \ldots
                             % made from \ + letters.
                             %
                             % Note how a `%' character causes TeX to ignore
                             % the end of the input line, so these blank lines
                             % do not start a new paragraph.
                             %
with the right spacing around the periods requires
a special command.

\LaTeX\ interprets some common characters as
commands, so you must type special commands to
generate them.  These characters include the
following:
       \$ \& \% \# \{ and \}.

In printing, text is usually emphasized with an
       \emph{italic}
type style.

\begin{em}
   A long segment of text can also be emphasized
   in this way.  Text within such a segment can be
   given \emph{additional} emphasis.
\end{em}

It is sometimes necessary to prevent \LaTeX\ from
breaking a line where it might otherwise do so.
This may be at a space, as between the ``Mr.''\ and
``Jones'' in
       ``Mr.~Jones'',        % ~ produces an unbreakable interword space.
or within a word---especially when the word is a
symbol like
       \mbox{\emph{itemnum}}
that makes little sense when hyphenated across
lines.

Footnotes\footnote{This is an example of a footnote.}
pose no problem.

\LaTeX\ is good at typesetting mathematical formulas
like
       \( x-3y + z = 7 \)
or
       \( a_{1} > x^{2n} + y^{2n} > x' \)
or
       \( \ip{A}{B} = \sum_{i} a_{i} b_{i} \).
The spaces you type in a formula are
ignored.  Remember that a letter like
       $x$                   % $ ... $  and  \( ... \)  are equivalent
is a formula when it denotes a mathematical
symbol, and it should be typed as one.

\section{Displayed Text}

Text is displayed by indenting it from the left
margin.  Quotations are commonly displayed.  There
are short quotations
\begin{quote}
   This is a short quotation.  It consists of a
   single paragraph of text.  See how it is formatted.
\end{quote}
and longer ones.
\begin{quotation}
   This is a longer quotation.  It consists of two
   paragraphs of text, neither of which are
   particularly interesting.

   This is the second paragraph of the quotation.  It
   is just as dull as the first paragraph.
\end{quotation}
Another frequently-displayed structure is a list.
The following is an example of an \emph{itemized}
list.
\begin{itemize}
   \item This is the first item of an itemized list.
         Each item in the list is marked with a ``tick''.
         You don't have to worry about what kind of tick
         mark is used.

   \item This is the second item of the list.  It
         contains another list nested inside it.  The inner
         list is an \emph{enumerated} list.
         \begin{enumerate}
            \item This is the first item of an enumerated
                  list that is nested within the itemized list.

            \item This is the second item of the inner list.
                  \LaTeX\ allows you to nest lists deeper than
                  you really should.
         \end{enumerate}
         This is the rest of the second item of the outer
         list.  It is no more interesting than any other
         part of the item.
   \item This is the third item of the list.
\end{itemize}
You can even display poetry.
\begin{verse}
   There is an environment
    for verse \\             % The \\ command separates lines
   Whose features some poets % within a stanza.
   will curse.

                             % One or more blank lines separate stanzas.

   For instead of making\\
   Them do \emph{all} line breaking, \\
   It allows them to put too many words on a line when they'd rather be
   forced to be terse.
\end{verse}

Mathematical formulas may also be displayed.  A
displayed formula
is
one-line long; multiline
formulas require special formatting instructions.
   \[  \ip{\Gamma}{\psi'} = x'' + y^{2} + z_{i}^{n}\]
Don't start a paragraph with a displayed equation,
nor make one a paragraph by itself.

\end{document}               % End of document.

\section{Bench Huge Section 119}
% synthetic bench section 119
% This is a sample LaTeX input file.  (Version of 12 August 2004.)
%
% A '%' character causes TeX to ignore all remaining text on the line,
% and is used for comments like this one.

\documentclass{article}      % Specifies the document class

                             % The preamble begins here.
\title{An Example Document}  % Declares the document's title.
\author{Leslie Lamport}      % Declares the author's name.
\date{January 21, 1994}      % Deleting this command produces today's date.

\newcommand{\ip}[2]{(#1, #2)}
                             % Defines \ip{arg1}{arg2} to mean
                             % (arg1, arg2).

%\newcommand{\ip}[2]{\langle #1 | #2\rangle}
                             % This is an alternative definition of
                             % \ip that is commented out.

\begin{document}             % End of preamble and beginning of text.

\maketitle                   % Produces the title.

This is an example input file.  Comparing it with
the output it generates can show you how to
produce a simple document of your own.

\section{Ordinary Text}      % Produces section heading.  Lower-level
                             % sections are begun with similar
                             % \subsection and \subsubsection commands.

The ends  of words and sentences are marked
  by   spaces. It  doesn't matter how many
spaces    you type; one is as good as 100.  The
end of   a line counts as a space.

One   or more   blank lines denote the  end
of  a paragraph.

Since any number of consecutive spaces are treated
like a single one, the formatting of the input
file makes no difference to
      \LaTeX,                % The \LaTeX command generates the LaTeX logo.
but it makes a difference to you.  When you use
\LaTeX, making your input file as easy to read
as possible will be a great help as you write
your document and when you change it.  This sample
file shows how you can add comments to your own input
file.

Because printing is different from typewriting,
there are a number of things that you have to do
differently when preparing an input file than if
you were just typing the document directly.
Quotation marks like
       ``this''
have to be handled specially, as do quotes within
quotes:
       ``\,`this'            % \, separates the double and single quote.
        is what I just
        wrote, not  `that'\,''.

Dashes come in three sizes: an
       intra-word
dash, a medium dash for number ranges like
       1--2,
and a punctuation
       dash---like
this.

A sentence-ending space should be larger than the
space between words within a sentence.  You
sometimes have to type special commands in
conjunction with punctuation characters to get
this right, as in the following sentence.
       Gnats, gnus, etc.\ all  % `\ ' makes an inter-word space.
       begin with G\@.         % \@ marks end-of-sentence punctuation.
You should check the spaces after periods when
reading your output to make sure you haven't
forgotten any special cases.  Generating an
ellipsis
       \ldots\               % `\ ' is needed after `\ldots' because TeX
                             % ignores spaces after command names like \ldots
                             % made from \ + letters.
                             %
                             % Note how a `%' character causes TeX to ignore
                             % the end of the input line, so these blank lines
                             % do not start a new paragraph.
                             %
with the right spacing around the periods requires
a special command.

\LaTeX\ interprets some common characters as
commands, so you must type special commands to
generate them.  These characters include the
following:
       \$ \& \% \# \{ and \}.

In printing, text is usually emphasized with an
       \emph{italic}
type style.

\begin{em}
   A long segment of text can also be emphasized
   in this way.  Text within such a segment can be
   given \emph{additional} emphasis.
\end{em}

It is sometimes necessary to prevent \LaTeX\ from
breaking a line where it might otherwise do so.
This may be at a space, as between the ``Mr.''\ and
``Jones'' in
       ``Mr.~Jones'',        % ~ produces an unbreakable interword space.
or within a word---especially when the word is a
symbol like
       \mbox{\emph{itemnum}}
that makes little sense when hyphenated across
lines.

Footnotes\footnote{This is an example of a footnote.}
pose no problem.

\LaTeX\ is good at typesetting mathematical formulas
like
       \( x-3y + z = 7 \)
or
       \( a_{1} > x^{2n} + y^{2n} > x' \)
or
       \( \ip{A}{B} = \sum_{i} a_{i} b_{i} \).
The spaces you type in a formula are
ignored.  Remember that a letter like
       $x$                   % $ ... $  and  \( ... \)  are equivalent
is a formula when it denotes a mathematical
symbol, and it should be typed as one.

\section{Displayed Text}

Text is displayed by indenting it from the left
margin.  Quotations are commonly displayed.  There
are short quotations
\begin{quote}
   This is a short quotation.  It consists of a
   single paragraph of text.  See how it is formatted.
\end{quote}
and longer ones.
\begin{quotation}
   This is a longer quotation.  It consists of two
   paragraphs of text, neither of which are
   particularly interesting.

   This is the second paragraph of the quotation.  It
   is just as dull as the first paragraph.
\end{quotation}
Another frequently-displayed structure is a list.
The following is an example of an \emph{itemized}
list.
\begin{itemize}
   \item This is the first item of an itemized list.
         Each item in the list is marked with a ``tick''.
         You don't have to worry about what kind of tick
         mark is used.

   \item This is the second item of the list.  It
         contains another list nested inside it.  The inner
         list is an \emph{enumerated} list.
         \begin{enumerate}
            \item This is the first item of an enumerated
                  list that is nested within the itemized list.

            \item This is the second item of the inner list.
                  \LaTeX\ allows you to nest lists deeper than
                  you really should.
         \end{enumerate}
         This is the rest of the second item of the outer
         list.  It is no more interesting than any other
         part of the item.
   \item This is the third item of the list.
\end{itemize}
You can even display poetry.
\begin{verse}
   There is an environment
    for verse \\             % The \\ command separates lines
   Whose features some poets % within a stanza.
   will curse.

                             % One or more blank lines separate stanzas.

   For instead of making\\
   Them do \emph{all} line breaking, \\
   It allows them to put too many words on a line when they'd rather be
   forced to be terse.
\end{verse}

Mathematical formulas may also be displayed.  A
displayed formula
is
one-line long; multiline
formulas require special formatting instructions.
   \[  \ip{\Gamma}{\psi'} = x'' + y^{2} + z_{i}^{n}\]
Don't start a paragraph with a displayed equation,
nor make one a paragraph by itself.

\end{document}               % End of document.

\section{Bench Huge Section 120}
% synthetic bench section 120
% This is a sample LaTeX input file.  (Version of 12 August 2004.)
%
% A '%' character causes TeX to ignore all remaining text on the line,
% and is used for comments like this one.

\documentclass{article}      % Specifies the document class

                             % The preamble begins here.
\title{An Example Document}  % Declares the document's title.
\author{Leslie Lamport}      % Declares the author's name.
\date{January 21, 1994}      % Deleting this command produces today's date.

\newcommand{\ip}[2]{(#1, #2)}
                             % Defines \ip{arg1}{arg2} to mean
                             % (arg1, arg2).

%\newcommand{\ip}[2]{\langle #1 | #2\rangle}
                             % This is an alternative definition of
                             % \ip that is commented out.

\begin{document}             % End of preamble and beginning of text.

\maketitle                   % Produces the title.

This is an example input file.  Comparing it with
the output it generates can show you how to
produce a simple document of your own.

\section{Ordinary Text}      % Produces section heading.  Lower-level
                             % sections are begun with similar
                             % \subsection and \subsubsection commands.

The ends  of words and sentences are marked
  by   spaces. It  doesn't matter how many
spaces    you type; one is as good as 100.  The
end of   a line counts as a space.

One   or more   blank lines denote the  end
of  a paragraph.

Since any number of consecutive spaces are treated
like a single one, the formatting of the input
file makes no difference to
      \LaTeX,                % The \LaTeX command generates the LaTeX logo.
but it makes a difference to you.  When you use
\LaTeX, making your input file as easy to read
as possible will be a great help as you write
your document and when you change it.  This sample
file shows how you can add comments to your own input
file.

Because printing is different from typewriting,
there are a number of things that you have to do
differently when preparing an input file than if
you were just typing the document directly.
Quotation marks like
       ``this''
have to be handled specially, as do quotes within
quotes:
       ``\,`this'            % \, separates the double and single quote.
        is what I just
        wrote, not  `that'\,''.

Dashes come in three sizes: an
       intra-word
dash, a medium dash for number ranges like
       1--2,
and a punctuation
       dash---like
this.

A sentence-ending space should be larger than the
space between words within a sentence.  You
sometimes have to type special commands in
conjunction with punctuation characters to get
this right, as in the following sentence.
       Gnats, gnus, etc.\ all  % `\ ' makes an inter-word space.
       begin with G\@.         % \@ marks end-of-sentence punctuation.
You should check the spaces after periods when
reading your output to make sure you haven't
forgotten any special cases.  Generating an
ellipsis
       \ldots\               % `\ ' is needed after `\ldots' because TeX
                             % ignores spaces after command names like \ldots
                             % made from \ + letters.
                             %
                             % Note how a `%' character causes TeX to ignore
                             % the end of the input line, so these blank lines
                             % do not start a new paragraph.
                             %
with the right spacing around the periods requires
a special command.

\LaTeX\ interprets some common characters as
commands, so you must type special commands to
generate them.  These characters include the
following:
       \$ \& \% \# \{ and \}.

In printing, text is usually emphasized with an
       \emph{italic}
type style.

\begin{em}
   A long segment of text can also be emphasized
   in this way.  Text within such a segment can be
   given \emph{additional} emphasis.
\end{em}

It is sometimes necessary to prevent \LaTeX\ from
breaking a line where it might otherwise do so.
This may be at a space, as between the ``Mr.''\ and
``Jones'' in
       ``Mr.~Jones'',        % ~ produces an unbreakable interword space.
or within a word---especially when the word is a
symbol like
       \mbox{\emph{itemnum}}
that makes little sense when hyphenated across
lines.

Footnotes\footnote{This is an example of a footnote.}
pose no problem.

\LaTeX\ is good at typesetting mathematical formulas
like
       \( x-3y + z = 7 \)
or
       \( a_{1} > x^{2n} + y^{2n} > x' \)
or
       \( \ip{A}{B} = \sum_{i} a_{i} b_{i} \).
The spaces you type in a formula are
ignored.  Remember that a letter like
       $x$                   % $ ... $  and  \( ... \)  are equivalent
is a formula when it denotes a mathematical
symbol, and it should be typed as one.

\section{Displayed Text}

Text is displayed by indenting it from the left
margin.  Quotations are commonly displayed.  There
are short quotations
\begin{quote}
   This is a short quotation.  It consists of a
   single paragraph of text.  See how it is formatted.
\end{quote}
and longer ones.
\begin{quotation}
   This is a longer quotation.  It consists of two
   paragraphs of text, neither of which are
   particularly interesting.

   This is the second paragraph of the quotation.  It
   is just as dull as the first paragraph.
\end{quotation}
Another frequently-displayed structure is a list.
The following is an example of an \emph{itemized}
list.
\begin{itemize}
   \item This is the first item of an itemized list.
         Each item in the list is marked with a ``tick''.
         You don't have to worry about what kind of tick
         mark is used.

   \item This is the second item of the list.  It
         contains another list nested inside it.  The inner
         list is an \emph{enumerated} list.
         \begin{enumerate}
            \item This is the first item of an enumerated
                  list that is nested within the itemized list.

            \item This is the second item of the inner list.
                  \LaTeX\ allows you to nest lists deeper than
                  you really should.
         \end{enumerate}
         This is the rest of the second item of the outer
         list.  It is no more interesting than any other
         part of the item.
   \item This is the third item of the list.
\end{itemize}
You can even display poetry.
\begin{verse}
   There is an environment
    for verse \\             % The \\ command separates lines
   Whose features some poets % within a stanza.
   will curse.

                             % One or more blank lines separate stanzas.

   For instead of making\\
   Them do \emph{all} line breaking, \\
   It allows them to put too many words on a line when they'd rather be
   forced to be terse.
\end{verse}

Mathematical formulas may also be displayed.  A
displayed formula
is
one-line long; multiline
formulas require special formatting instructions.
   \[  \ip{\Gamma}{\psi'} = x'' + y^{2} + z_{i}^{n}\]
Don't start a paragraph with a displayed equation,
nor make one a paragraph by itself.

\end{document}               % End of document.

\section{Bench Huge Section 121}
% synthetic bench section 121
% This is a sample LaTeX input file.  (Version of 12 August 2004.)
%
% A '%' character causes TeX to ignore all remaining text on the line,
% and is used for comments like this one.

\documentclass{article}      % Specifies the document class

                             % The preamble begins here.
\title{An Example Document}  % Declares the document's title.
\author{Leslie Lamport}      % Declares the author's name.
\date{January 21, 1994}      % Deleting this command produces today's date.

\newcommand{\ip}[2]{(#1, #2)}
                             % Defines \ip{arg1}{arg2} to mean
                             % (arg1, arg2).

%\newcommand{\ip}[2]{\langle #1 | #2\rangle}
                             % This is an alternative definition of
                             % \ip that is commented out.

\begin{document}             % End of preamble and beginning of text.

\maketitle                   % Produces the title.

This is an example input file.  Comparing it with
the output it generates can show you how to
produce a simple document of your own.

\section{Ordinary Text}      % Produces section heading.  Lower-level
                             % sections are begun with similar
                             % \subsection and \subsubsection commands.

The ends  of words and sentences are marked
  by   spaces. It  doesn't matter how many
spaces    you type; one is as good as 100.  The
end of   a line counts as a space.

One   or more   blank lines denote the  end
of  a paragraph.

Since any number of consecutive spaces are treated
like a single one, the formatting of the input
file makes no difference to
      \LaTeX,                % The \LaTeX command generates the LaTeX logo.
but it makes a difference to you.  When you use
\LaTeX, making your input file as easy to read
as possible will be a great help as you write
your document and when you change it.  This sample
file shows how you can add comments to your own input
file.

Because printing is different from typewriting,
there are a number of things that you have to do
differently when preparing an input file than if
you were just typing the document directly.
Quotation marks like
       ``this''
have to be handled specially, as do quotes within
quotes:
       ``\,`this'            % \, separates the double and single quote.
        is what I just
        wrote, not  `that'\,''.

Dashes come in three sizes: an
       intra-word
dash, a medium dash for number ranges like
       1--2,
and a punctuation
       dash---like
this.

A sentence-ending space should be larger than the
space between words within a sentence.  You
sometimes have to type special commands in
conjunction with punctuation characters to get
this right, as in the following sentence.
       Gnats, gnus, etc.\ all  % `\ ' makes an inter-word space.
       begin with G\@.         % \@ marks end-of-sentence punctuation.
You should check the spaces after periods when
reading your output to make sure you haven't
forgotten any special cases.  Generating an
ellipsis
       \ldots\               % `\ ' is needed after `\ldots' because TeX
                             % ignores spaces after command names like \ldots
                             % made from \ + letters.
                             %
                             % Note how a `%' character causes TeX to ignore
                             % the end of the input line, so these blank lines
                             % do not start a new paragraph.
                             %
with the right spacing around the periods requires
a special command.

\LaTeX\ interprets some common characters as
commands, so you must type special commands to
generate them.  These characters include the
following:
       \$ \& \% \# \{ and \}.

In printing, text is usually emphasized with an
       \emph{italic}
type style.

\begin{em}
   A long segment of text can also be emphasized
   in this way.  Text within such a segment can be
   given \emph{additional} emphasis.
\end{em}

It is sometimes necessary to prevent \LaTeX\ from
breaking a line where it might otherwise do so.
This may be at a space, as between the ``Mr.''\ and
``Jones'' in
       ``Mr.~Jones'',        % ~ produces an unbreakable interword space.
or within a word---especially when the word is a
symbol like
       \mbox{\emph{itemnum}}
that makes little sense when hyphenated across
lines.

Footnotes\footnote{This is an example of a footnote.}
pose no problem.

\LaTeX\ is good at typesetting mathematical formulas
like
       \( x-3y + z = 7 \)
or
       \( a_{1} > x^{2n} + y^{2n} > x' \)
or
       \( \ip{A}{B} = \sum_{i} a_{i} b_{i} \).
The spaces you type in a formula are
ignored.  Remember that a letter like
       $x$                   % $ ... $  and  \( ... \)  are equivalent
is a formula when it denotes a mathematical
symbol, and it should be typed as one.

\section{Displayed Text}

Text is displayed by indenting it from the left
margin.  Quotations are commonly displayed.  There
are short quotations
\begin{quote}
   This is a short quotation.  It consists of a
   single paragraph of text.  See how it is formatted.
\end{quote}
and longer ones.
\begin{quotation}
   This is a longer quotation.  It consists of two
   paragraphs of text, neither of which are
   particularly interesting.

   This is the second paragraph of the quotation.  It
   is just as dull as the first paragraph.
\end{quotation}
Another frequently-displayed structure is a list.
The following is an example of an \emph{itemized}
list.
\begin{itemize}
   \item This is the first item of an itemized list.
         Each item in the list is marked with a ``tick''.
         You don't have to worry about what kind of tick
         mark is used.

   \item This is the second item of the list.  It
         contains another list nested inside it.  The inner
         list is an \emph{enumerated} list.
         \begin{enumerate}
            \item This is the first item of an enumerated
                  list that is nested within the itemized list.

            \item This is the second item of the inner list.
                  \LaTeX\ allows you to nest lists deeper than
                  you really should.
         \end{enumerate}
         This is the rest of the second item of the outer
         list.  It is no more interesting than any other
         part of the item.
   \item This is the third item of the list.
\end{itemize}
You can even display poetry.
\begin{verse}
   There is an environment
    for verse \\             % The \\ command separates lines
   Whose features some poets % within a stanza.
   will curse.

                             % One or more blank lines separate stanzas.

   For instead of making\\
   Them do \emph{all} line breaking, \\
   It allows them to put too many words on a line when they'd rather be
   forced to be terse.
\end{verse}

Mathematical formulas may also be displayed.  A
displayed formula
is
one-line long; multiline
formulas require special formatting instructions.
   \[  \ip{\Gamma}{\psi'} = x'' + y^{2} + z_{i}^{n}\]
Don't start a paragraph with a displayed equation,
nor make one a paragraph by itself.

\end{document}               % End of document.

\section{Bench Huge Section 122}
% synthetic bench section 122
% This is a sample LaTeX input file.  (Version of 12 August 2004.)
%
% A '%' character causes TeX to ignore all remaining text on the line,
% and is used for comments like this one.

\documentclass{article}      % Specifies the document class

                             % The preamble begins here.
\title{An Example Document}  % Declares the document's title.
\author{Leslie Lamport}      % Declares the author's name.
\date{January 21, 1994}      % Deleting this command produces today's date.

\newcommand{\ip}[2]{(#1, #2)}
                             % Defines \ip{arg1}{arg2} to mean
                             % (arg1, arg2).

%\newcommand{\ip}[2]{\langle #1 | #2\rangle}
                             % This is an alternative definition of
                             % \ip that is commented out.

\begin{document}             % End of preamble and beginning of text.

\maketitle                   % Produces the title.

This is an example input file.  Comparing it with
the output it generates can show you how to
produce a simple document of your own.

\section{Ordinary Text}      % Produces section heading.  Lower-level
                             % sections are begun with similar
                             % \subsection and \subsubsection commands.

The ends  of words and sentences are marked
  by   spaces. It  doesn't matter how many
spaces    you type; one is as good as 100.  The
end of   a line counts as a space.

One   or more   blank lines denote the  end
of  a paragraph.

Since any number of consecutive spaces are treated
like a single one, the formatting of the input
file makes no difference to
      \LaTeX,                % The \LaTeX command generates the LaTeX logo.
but it makes a difference to you.  When you use
\LaTeX, making your input file as easy to read
as possible will be a great help as you write
your document and when you change it.  This sample
file shows how you can add comments to your own input
file.

Because printing is different from typewriting,
there are a number of things that you have to do
differently when preparing an input file than if
you were just typing the document directly.
Quotation marks like
       ``this''
have to be handled specially, as do quotes within
quotes:
       ``\,`this'            % \, separates the double and single quote.
        is what I just
        wrote, not  `that'\,''.

Dashes come in three sizes: an
       intra-word
dash, a medium dash for number ranges like
       1--2,
and a punctuation
       dash---like
this.

A sentence-ending space should be larger than the
space between words within a sentence.  You
sometimes have to type special commands in
conjunction with punctuation characters to get
this right, as in the following sentence.
       Gnats, gnus, etc.\ all  % `\ ' makes an inter-word space.
       begin with G\@.         % \@ marks end-of-sentence punctuation.
You should check the spaces after periods when
reading your output to make sure you haven't
forgotten any special cases.  Generating an
ellipsis
       \ldots\               % `\ ' is needed after `\ldots' because TeX
                             % ignores spaces after command names like \ldots
                             % made from \ + letters.
                             %
                             % Note how a `%' character causes TeX to ignore
                             % the end of the input line, so these blank lines
                             % do not start a new paragraph.
                             %
with the right spacing around the periods requires
a special command.

\LaTeX\ interprets some common characters as
commands, so you must type special commands to
generate them.  These characters include the
following:
       \$ \& \% \# \{ and \}.

In printing, text is usually emphasized with an
       \emph{italic}
type style.

\begin{em}
   A long segment of text can also be emphasized
   in this way.  Text within such a segment can be
   given \emph{additional} emphasis.
\end{em}

It is sometimes necessary to prevent \LaTeX\ from
breaking a line where it might otherwise do so.
This may be at a space, as between the ``Mr.''\ and
``Jones'' in
       ``Mr.~Jones'',        % ~ produces an unbreakable interword space.
or within a word---especially when the word is a
symbol like
       \mbox{\emph{itemnum}}
that makes little sense when hyphenated across
lines.

Footnotes\footnote{This is an example of a footnote.}
pose no problem.

\LaTeX\ is good at typesetting mathematical formulas
like
       \( x-3y + z = 7 \)
or
       \( a_{1} > x^{2n} + y^{2n} > x' \)
or
       \( \ip{A}{B} = \sum_{i} a_{i} b_{i} \).
The spaces you type in a formula are
ignored.  Remember that a letter like
       $x$                   % $ ... $  and  \( ... \)  are equivalent
is a formula when it denotes a mathematical
symbol, and it should be typed as one.

\section{Displayed Text}

Text is displayed by indenting it from the left
margin.  Quotations are commonly displayed.  There
are short quotations
\begin{quote}
   This is a short quotation.  It consists of a
   single paragraph of text.  See how it is formatted.
\end{quote}
and longer ones.
\begin{quotation}
   This is a longer quotation.  It consists of two
   paragraphs of text, neither of which are
   particularly interesting.

   This is the second paragraph of the quotation.  It
   is just as dull as the first paragraph.
\end{quotation}
Another frequently-displayed structure is a list.
The following is an example of an \emph{itemized}
list.
\begin{itemize}
   \item This is the first item of an itemized list.
         Each item in the list is marked with a ``tick''.
         You don't have to worry about what kind of tick
         mark is used.

   \item This is the second item of the list.  It
         contains another list nested inside it.  The inner
         list is an \emph{enumerated} list.
         \begin{enumerate}
            \item This is the first item of an enumerated
                  list that is nested within the itemized list.

            \item This is the second item of the inner list.
                  \LaTeX\ allows you to nest lists deeper than
                  you really should.
         \end{enumerate}
         This is the rest of the second item of the outer
         list.  It is no more interesting than any other
         part of the item.
   \item This is the third item of the list.
\end{itemize}
You can even display poetry.
\begin{verse}
   There is an environment
    for verse \\             % The \\ command separates lines
   Whose features some poets % within a stanza.
   will curse.

                             % One or more blank lines separate stanzas.

   For instead of making\\
   Them do \emph{all} line breaking, \\
   It allows them to put too many words on a line when they'd rather be
   forced to be terse.
\end{verse}

Mathematical formulas may also be displayed.  A
displayed formula
is
one-line long; multiline
formulas require special formatting instructions.
   \[  \ip{\Gamma}{\psi'} = x'' + y^{2} + z_{i}^{n}\]
Don't start a paragraph with a displayed equation,
nor make one a paragraph by itself.

\end{document}               % End of document.

\section{Bench Huge Section 123}
% synthetic bench section 123
% This is a sample LaTeX input file.  (Version of 12 August 2004.)
%
% A '%' character causes TeX to ignore all remaining text on the line,
% and is used for comments like this one.

\documentclass{article}      % Specifies the document class

                             % The preamble begins here.
\title{An Example Document}  % Declares the document's title.
\author{Leslie Lamport}      % Declares the author's name.
\date{January 21, 1994}      % Deleting this command produces today's date.

\newcommand{\ip}[2]{(#1, #2)}
                             % Defines \ip{arg1}{arg2} to mean
                             % (arg1, arg2).

%\newcommand{\ip}[2]{\langle #1 | #2\rangle}
                             % This is an alternative definition of
                             % \ip that is commented out.

\begin{document}             % End of preamble and beginning of text.

\maketitle                   % Produces the title.

This is an example input file.  Comparing it with
the output it generates can show you how to
produce a simple document of your own.

\section{Ordinary Text}      % Produces section heading.  Lower-level
                             % sections are begun with similar
                             % \subsection and \subsubsection commands.

The ends  of words and sentences are marked
  by   spaces. It  doesn't matter how many
spaces    you type; one is as good as 100.  The
end of   a line counts as a space.

One   or more   blank lines denote the  end
of  a paragraph.

Since any number of consecutive spaces are treated
like a single one, the formatting of the input
file makes no difference to
      \LaTeX,                % The \LaTeX command generates the LaTeX logo.
but it makes a difference to you.  When you use
\LaTeX, making your input file as easy to read
as possible will be a great help as you write
your document and when you change it.  This sample
file shows how you can add comments to your own input
file.

Because printing is different from typewriting,
there are a number of things that you have to do
differently when preparing an input file than if
you were just typing the document directly.
Quotation marks like
       ``this''
have to be handled specially, as do quotes within
quotes:
       ``\,`this'            % \, separates the double and single quote.
        is what I just
        wrote, not  `that'\,''.

Dashes come in three sizes: an
       intra-word
dash, a medium dash for number ranges like
       1--2,
and a punctuation
       dash---like
this.

A sentence-ending space should be larger than the
space between words within a sentence.  You
sometimes have to type special commands in
conjunction with punctuation characters to get
this right, as in the following sentence.
       Gnats, gnus, etc.\ all  % `\ ' makes an inter-word space.
       begin with G\@.         % \@ marks end-of-sentence punctuation.
You should check the spaces after periods when
reading your output to make sure you haven't
forgotten any special cases.  Generating an
ellipsis
       \ldots\               % `\ ' is needed after `\ldots' because TeX
                             % ignores spaces after command names like \ldots
                             % made from \ + letters.
                             %
                             % Note how a `%' character causes TeX to ignore
                             % the end of the input line, so these blank lines
                             % do not start a new paragraph.
                             %
with the right spacing around the periods requires
a special command.

\LaTeX\ interprets some common characters as
commands, so you must type special commands to
generate them.  These characters include the
following:
       \$ \& \% \# \{ and \}.

In printing, text is usually emphasized with an
       \emph{italic}
type style.

\begin{em}
   A long segment of text can also be emphasized
   in this way.  Text within such a segment can be
   given \emph{additional} emphasis.
\end{em}

It is sometimes necessary to prevent \LaTeX\ from
breaking a line where it might otherwise do so.
This may be at a space, as between the ``Mr.''\ and
``Jones'' in
       ``Mr.~Jones'',        % ~ produces an unbreakable interword space.
or within a word---especially when the word is a
symbol like
       \mbox{\emph{itemnum}}
that makes little sense when hyphenated across
lines.

Footnotes\footnote{This is an example of a footnote.}
pose no problem.

\LaTeX\ is good at typesetting mathematical formulas
like
       \( x-3y + z = 7 \)
or
       \( a_{1} > x^{2n} + y^{2n} > x' \)
or
       \( \ip{A}{B} = \sum_{i} a_{i} b_{i} \).
The spaces you type in a formula are
ignored.  Remember that a letter like
       $x$                   % $ ... $  and  \( ... \)  are equivalent
is a formula when it denotes a mathematical
symbol, and it should be typed as one.

\section{Displayed Text}

Text is displayed by indenting it from the left
margin.  Quotations are commonly displayed.  There
are short quotations
\begin{quote}
   This is a short quotation.  It consists of a
   single paragraph of text.  See how it is formatted.
\end{quote}
and longer ones.
\begin{quotation}
   This is a longer quotation.  It consists of two
   paragraphs of text, neither of which are
   particularly interesting.

   This is the second paragraph of the quotation.  It
   is just as dull as the first paragraph.
\end{quotation}
Another frequently-displayed structure is a list.
The following is an example of an \emph{itemized}
list.
\begin{itemize}
   \item This is the first item of an itemized list.
         Each item in the list is marked with a ``tick''.
         You don't have to worry about what kind of tick
         mark is used.

   \item This is the second item of the list.  It
         contains another list nested inside it.  The inner
         list is an \emph{enumerated} list.
         \begin{enumerate}
            \item This is the first item of an enumerated
                  list that is nested within the itemized list.

            \item This is the second item of the inner list.
                  \LaTeX\ allows you to nest lists deeper than
                  you really should.
         \end{enumerate}
         This is the rest of the second item of the outer
         list.  It is no more interesting than any other
         part of the item.
   \item This is the third item of the list.
\end{itemize}
You can even display poetry.
\begin{verse}
   There is an environment
    for verse \\             % The \\ command separates lines
   Whose features some poets % within a stanza.
   will curse.

                             % One or more blank lines separate stanzas.

   For instead of making\\
   Them do \emph{all} line breaking, \\
   It allows them to put too many words on a line when they'd rather be
   forced to be terse.
\end{verse}

Mathematical formulas may also be displayed.  A
displayed formula
is
one-line long; multiline
formulas require special formatting instructions.
   \[  \ip{\Gamma}{\psi'} = x'' + y^{2} + z_{i}^{n}\]
Don't start a paragraph with a displayed equation,
nor make one a paragraph by itself.

\end{document}               % End of document.

\section{Bench Huge Section 124}
% synthetic bench section 124
% This is a sample LaTeX input file.  (Version of 12 August 2004.)
%
% A '%' character causes TeX to ignore all remaining text on the line,
% and is used for comments like this one.

\documentclass{article}      % Specifies the document class

                             % The preamble begins here.
\title{An Example Document}  % Declares the document's title.
\author{Leslie Lamport}      % Declares the author's name.
\date{January 21, 1994}      % Deleting this command produces today's date.

\newcommand{\ip}[2]{(#1, #2)}
                             % Defines \ip{arg1}{arg2} to mean
                             % (arg1, arg2).

%\newcommand{\ip}[2]{\langle #1 | #2\rangle}
                             % This is an alternative definition of
                             % \ip that is commented out.

\begin{document}             % End of preamble and beginning of text.

\maketitle                   % Produces the title.

This is an example input file.  Comparing it with
the output it generates can show you how to
produce a simple document of your own.

\section{Ordinary Text}      % Produces section heading.  Lower-level
                             % sections are begun with similar
                             % \subsection and \subsubsection commands.

The ends  of words and sentences are marked
  by   spaces. It  doesn't matter how many
spaces    you type; one is as good as 100.  The
end of   a line counts as a space.

One   or more   blank lines denote the  end
of  a paragraph.

Since any number of consecutive spaces are treated
like a single one, the formatting of the input
file makes no difference to
      \LaTeX,                % The \LaTeX command generates the LaTeX logo.
but it makes a difference to you.  When you use
\LaTeX, making your input file as easy to read
as possible will be a great help as you write
your document and when you change it.  This sample
file shows how you can add comments to your own input
file.

Because printing is different from typewriting,
there are a number of things that you have to do
differently when preparing an input file than if
you were just typing the document directly.
Quotation marks like
       ``this''
have to be handled specially, as do quotes within
quotes:
       ``\,`this'            % \, separates the double and single quote.
        is what I just
        wrote, not  `that'\,''.

Dashes come in three sizes: an
       intra-word
dash, a medium dash for number ranges like
       1--2,
and a punctuation
       dash---like
this.

A sentence-ending space should be larger than the
space between words within a sentence.  You
sometimes have to type special commands in
conjunction with punctuation characters to get
this right, as in the following sentence.
       Gnats, gnus, etc.\ all  % `\ ' makes an inter-word space.
       begin with G\@.         % \@ marks end-of-sentence punctuation.
You should check the spaces after periods when
reading your output to make sure you haven't
forgotten any special cases.  Generating an
ellipsis
       \ldots\               % `\ ' is needed after `\ldots' because TeX
                             % ignores spaces after command names like \ldots
                             % made from \ + letters.
                             %
                             % Note how a `%' character causes TeX to ignore
                             % the end of the input line, so these blank lines
                             % do not start a new paragraph.
                             %
with the right spacing around the periods requires
a special command.

\LaTeX\ interprets some common characters as
commands, so you must type special commands to
generate them.  These characters include the
following:
       \$ \& \% \# \{ and \}.

In printing, text is usually emphasized with an
       \emph{italic}
type style.

\begin{em}
   A long segment of text can also be emphasized
   in this way.  Text within such a segment can be
   given \emph{additional} emphasis.
\end{em}

It is sometimes necessary to prevent \LaTeX\ from
breaking a line where it might otherwise do so.
This may be at a space, as between the ``Mr.''\ and
``Jones'' in
       ``Mr.~Jones'',        % ~ produces an unbreakable interword space.
or within a word---especially when the word is a
symbol like
       \mbox{\emph{itemnum}}
that makes little sense when hyphenated across
lines.

Footnotes\footnote{This is an example of a footnote.}
pose no problem.

\LaTeX\ is good at typesetting mathematical formulas
like
       \( x-3y + z = 7 \)
or
       \( a_{1} > x^{2n} + y^{2n} > x' \)
or
       \( \ip{A}{B} = \sum_{i} a_{i} b_{i} \).
The spaces you type in a formula are
ignored.  Remember that a letter like
       $x$                   % $ ... $  and  \( ... \)  are equivalent
is a formula when it denotes a mathematical
symbol, and it should be typed as one.

\section{Displayed Text}

Text is displayed by indenting it from the left
margin.  Quotations are commonly displayed.  There
are short quotations
\begin{quote}
   This is a short quotation.  It consists of a
   single paragraph of text.  See how it is formatted.
\end{quote}
and longer ones.
\begin{quotation}
   This is a longer quotation.  It consists of two
   paragraphs of text, neither of which are
   particularly interesting.

   This is the second paragraph of the quotation.  It
   is just as dull as the first paragraph.
\end{quotation}
Another frequently-displayed structure is a list.
The following is an example of an \emph{itemized}
list.
\begin{itemize}
   \item This is the first item of an itemized list.
         Each item in the list is marked with a ``tick''.
         You don't have to worry about what kind of tick
         mark is used.

   \item This is the second item of the list.  It
         contains another list nested inside it.  The inner
         list is an \emph{enumerated} list.
         \begin{enumerate}
            \item This is the first item of an enumerated
                  list that is nested within the itemized list.

            \item This is the second item of the inner list.
                  \LaTeX\ allows you to nest lists deeper than
                  you really should.
         \end{enumerate}
         This is the rest of the second item of the outer
         list.  It is no more interesting than any other
         part of the item.
   \item This is the third item of the list.
\end{itemize}
You can even display poetry.
\begin{verse}
   There is an environment
    for verse \\             % The \\ command separates lines
   Whose features some poets % within a stanza.
   will curse.

                             % One or more blank lines separate stanzas.

   For instead of making\\
   Them do \emph{all} line breaking, \\
   It allows them to put too many words on a line when they'd rather be
   forced to be terse.
\end{verse}

Mathematical formulas may also be displayed.  A
displayed formula
is
one-line long; multiline
formulas require special formatting instructions.
   \[  \ip{\Gamma}{\psi'} = x'' + y^{2} + z_{i}^{n}\]
Don't start a paragraph with a displayed equation,
nor make one a paragraph by itself.

\end{document}               % End of document.

\section{Bench Huge Section 125}
% synthetic bench section 125
% This is a sample LaTeX input file.  (Version of 12 August 2004.)
%
% A '%' character causes TeX to ignore all remaining text on the line,
% and is used for comments like this one.

\documentclass{article}      % Specifies the document class

                             % The preamble begins here.
\title{An Example Document}  % Declares the document's title.
\author{Leslie Lamport}      % Declares the author's name.
\date{January 21, 1994}      % Deleting this command produces today's date.

\newcommand{\ip}[2]{(#1, #2)}
                             % Defines \ip{arg1}{arg2} to mean
                             % (arg1, arg2).

%\newcommand{\ip}[2]{\langle #1 | #2\rangle}
                             % This is an alternative definition of
                             % \ip that is commented out.

\begin{document}             % End of preamble and beginning of text.

\maketitle                   % Produces the title.

This is an example input file.  Comparing it with
the output it generates can show you how to
produce a simple document of your own.

\section{Ordinary Text}      % Produces section heading.  Lower-level
                             % sections are begun with similar
                             % \subsection and \subsubsection commands.

The ends  of words and sentences are marked
  by   spaces. It  doesn't matter how many
spaces    you type; one is as good as 100.  The
end of   a line counts as a space.

One   or more   blank lines denote the  end
of  a paragraph.

Since any number of consecutive spaces are treated
like a single one, the formatting of the input
file makes no difference to
      \LaTeX,                % The \LaTeX command generates the LaTeX logo.
but it makes a difference to you.  When you use
\LaTeX, making your input file as easy to read
as possible will be a great help as you write
your document and when you change it.  This sample
file shows how you can add comments to your own input
file.

Because printing is different from typewriting,
there are a number of things that you have to do
differently when preparing an input file than if
you were just typing the document directly.
Quotation marks like
       ``this''
have to be handled specially, as do quotes within
quotes:
       ``\,`this'            % \, separates the double and single quote.
        is what I just
        wrote, not  `that'\,''.

Dashes come in three sizes: an
       intra-word
dash, a medium dash for number ranges like
       1--2,
and a punctuation
       dash---like
this.

A sentence-ending space should be larger than the
space between words within a sentence.  You
sometimes have to type special commands in
conjunction with punctuation characters to get
this right, as in the following sentence.
       Gnats, gnus, etc.\ all  % `\ ' makes an inter-word space.
       begin with G\@.         % \@ marks end-of-sentence punctuation.
You should check the spaces after periods when
reading your output to make sure you haven't
forgotten any special cases.  Generating an
ellipsis
       \ldots\               % `\ ' is needed after `\ldots' because TeX
                             % ignores spaces after command names like \ldots
                             % made from \ + letters.
                             %
                             % Note how a `%' character causes TeX to ignore
                             % the end of the input line, so these blank lines
                             % do not start a new paragraph.
                             %
with the right spacing around the periods requires
a special command.

\LaTeX\ interprets some common characters as
commands, so you must type special commands to
generate them.  These characters include the
following:
       \$ \& \% \# \{ and \}.

In printing, text is usually emphasized with an
       \emph{italic}
type style.

\begin{em}
   A long segment of text can also be emphasized
   in this way.  Text within such a segment can be
   given \emph{additional} emphasis.
\end{em}

It is sometimes necessary to prevent \LaTeX\ from
breaking a line where it might otherwise do so.
This may be at a space, as between the ``Mr.''\ and
``Jones'' in
       ``Mr.~Jones'',        % ~ produces an unbreakable interword space.
or within a word---especially when the word is a
symbol like
       \mbox{\emph{itemnum}}
that makes little sense when hyphenated across
lines.

Footnotes\footnote{This is an example of a footnote.}
pose no problem.

\LaTeX\ is good at typesetting mathematical formulas
like
       \( x-3y + z = 7 \)
or
       \( a_{1} > x^{2n} + y^{2n} > x' \)
or
       \( \ip{A}{B} = \sum_{i} a_{i} b_{i} \).
The spaces you type in a formula are
ignored.  Remember that a letter like
       $x$                   % $ ... $  and  \( ... \)  are equivalent
is a formula when it denotes a mathematical
symbol, and it should be typed as one.

\section{Displayed Text}

Text is displayed by indenting it from the left
margin.  Quotations are commonly displayed.  There
are short quotations
\begin{quote}
   This is a short quotation.  It consists of a
   single paragraph of text.  See how it is formatted.
\end{quote}
and longer ones.
\begin{quotation}
   This is a longer quotation.  It consists of two
   paragraphs of text, neither of which are
   particularly interesting.

   This is the second paragraph of the quotation.  It
   is just as dull as the first paragraph.
\end{quotation}
Another frequently-displayed structure is a list.
The following is an example of an \emph{itemized}
list.
\begin{itemize}
   \item This is the first item of an itemized list.
         Each item in the list is marked with a ``tick''.
         You don't have to worry about what kind of tick
         mark is used.

   \item This is the second item of the list.  It
         contains another list nested inside it.  The inner
         list is an \emph{enumerated} list.
         \begin{enumerate}
            \item This is the first item of an enumerated
                  list that is nested within the itemized list.

            \item This is the second item of the inner list.
                  \LaTeX\ allows you to nest lists deeper than
                  you really should.
         \end{enumerate}
         This is the rest of the second item of the outer
         list.  It is no more interesting than any other
         part of the item.
   \item This is the third item of the list.
\end{itemize}
You can even display poetry.
\begin{verse}
   There is an environment
    for verse \\             % The \\ command separates lines
   Whose features some poets % within a stanza.
   will curse.

                             % One or more blank lines separate stanzas.

   For instead of making\\
   Them do \emph{all} line breaking, \\
   It allows them to put too many words on a line when they'd rather be
   forced to be terse.
\end{verse}

Mathematical formulas may also be displayed.  A
displayed formula
is
one-line long; multiline
formulas require special formatting instructions.
   \[  \ip{\Gamma}{\psi'} = x'' + y^{2} + z_{i}^{n}\]
Don't start a paragraph with a displayed equation,
nor make one a paragraph by itself.

\end{document}               % End of document.

\section{Bench Huge Section 126}
% synthetic bench section 126
% This is a sample LaTeX input file.  (Version of 12 August 2004.)
%
% A '%' character causes TeX to ignore all remaining text on the line,
% and is used for comments like this one.

\documentclass{article}      % Specifies the document class

                             % The preamble begins here.
\title{An Example Document}  % Declares the document's title.
\author{Leslie Lamport}      % Declares the author's name.
\date{January 21, 1994}      % Deleting this command produces today's date.

\newcommand{\ip}[2]{(#1, #2)}
                             % Defines \ip{arg1}{arg2} to mean
                             % (arg1, arg2).

%\newcommand{\ip}[2]{\langle #1 | #2\rangle}
                             % This is an alternative definition of
                             % \ip that is commented out.

\begin{document}             % End of preamble and beginning of text.

\maketitle                   % Produces the title.

This is an example input file.  Comparing it with
the output it generates can show you how to
produce a simple document of your own.

\section{Ordinary Text}      % Produces section heading.  Lower-level
                             % sections are begun with similar
                             % \subsection and \subsubsection commands.

The ends  of words and sentences are marked
  by   spaces. It  doesn't matter how many
spaces    you type; one is as good as 100.  The
end of   a line counts as a space.

One   or more   blank lines denote the  end
of  a paragraph.

Since any number of consecutive spaces are treated
like a single one, the formatting of the input
file makes no difference to
      \LaTeX,                % The \LaTeX command generates the LaTeX logo.
but it makes a difference to you.  When you use
\LaTeX, making your input file as easy to read
as possible will be a great help as you write
your document and when you change it.  This sample
file shows how you can add comments to your own input
file.

Because printing is different from typewriting,
there are a number of things that you have to do
differently when preparing an input file than if
you were just typing the document directly.
Quotation marks like
       ``this''
have to be handled specially, as do quotes within
quotes:
       ``\,`this'            % \, separates the double and single quote.
        is what I just
        wrote, not  `that'\,''.

Dashes come in three sizes: an
       intra-word
dash, a medium dash for number ranges like
       1--2,
and a punctuation
       dash---like
this.

A sentence-ending space should be larger than the
space between words within a sentence.  You
sometimes have to type special commands in
conjunction with punctuation characters to get
this right, as in the following sentence.
       Gnats, gnus, etc.\ all  % `\ ' makes an inter-word space.
       begin with G\@.         % \@ marks end-of-sentence punctuation.
You should check the spaces after periods when
reading your output to make sure you haven't
forgotten any special cases.  Generating an
ellipsis
       \ldots\               % `\ ' is needed after `\ldots' because TeX
                             % ignores spaces after command names like \ldots
                             % made from \ + letters.
                             %
                             % Note how a `%' character causes TeX to ignore
                             % the end of the input line, so these blank lines
                             % do not start a new paragraph.
                             %
with the right spacing around the periods requires
a special command.

\LaTeX\ interprets some common characters as
commands, so you must type special commands to
generate them.  These characters include the
following:
       \$ \& \% \# \{ and \}.

In printing, text is usually emphasized with an
       \emph{italic}
type style.

\begin{em}
   A long segment of text can also be emphasized
   in this way.  Text within such a segment can be
   given \emph{additional} emphasis.
\end{em}

It is sometimes necessary to prevent \LaTeX\ from
breaking a line where it might otherwise do so.
This may be at a space, as between the ``Mr.''\ and
``Jones'' in
       ``Mr.~Jones'',        % ~ produces an unbreakable interword space.
or within a word---especially when the word is a
symbol like
       \mbox{\emph{itemnum}}
that makes little sense when hyphenated across
lines.

Footnotes\footnote{This is an example of a footnote.}
pose no problem.

\LaTeX\ is good at typesetting mathematical formulas
like
       \( x-3y + z = 7 \)
or
       \( a_{1} > x^{2n} + y^{2n} > x' \)
or
       \( \ip{A}{B} = \sum_{i} a_{i} b_{i} \).
The spaces you type in a formula are
ignored.  Remember that a letter like
       $x$                   % $ ... $  and  \( ... \)  are equivalent
is a formula when it denotes a mathematical
symbol, and it should be typed as one.

\section{Displayed Text}

Text is displayed by indenting it from the left
margin.  Quotations are commonly displayed.  There
are short quotations
\begin{quote}
   This is a short quotation.  It consists of a
   single paragraph of text.  See how it is formatted.
\end{quote}
and longer ones.
\begin{quotation}
   This is a longer quotation.  It consists of two
   paragraphs of text, neither of which are
   particularly interesting.

   This is the second paragraph of the quotation.  It
   is just as dull as the first paragraph.
\end{quotation}
Another frequently-displayed structure is a list.
The following is an example of an \emph{itemized}
list.
\begin{itemize}
   \item This is the first item of an itemized list.
         Each item in the list is marked with a ``tick''.
         You don't have to worry about what kind of tick
         mark is used.

   \item This is the second item of the list.  It
         contains another list nested inside it.  The inner
         list is an \emph{enumerated} list.
         \begin{enumerate}
            \item This is the first item of an enumerated
                  list that is nested within the itemized list.

            \item This is the second item of the inner list.
                  \LaTeX\ allows you to nest lists deeper than
                  you really should.
         \end{enumerate}
         This is the rest of the second item of the outer
         list.  It is no more interesting than any other
         part of the item.
   \item This is the third item of the list.
\end{itemize}
You can even display poetry.
\begin{verse}
   There is an environment
    for verse \\             % The \\ command separates lines
   Whose features some poets % within a stanza.
   will curse.

                             % One or more blank lines separate stanzas.

   For instead of making\\
   Them do \emph{all} line breaking, \\
   It allows them to put too many words on a line when they'd rather be
   forced to be terse.
\end{verse}

Mathematical formulas may also be displayed.  A
displayed formula
is
one-line long; multiline
formulas require special formatting instructions.
   \[  \ip{\Gamma}{\psi'} = x'' + y^{2} + z_{i}^{n}\]
Don't start a paragraph with a displayed equation,
nor make one a paragraph by itself.

\end{document}               % End of document.

\section{Bench Huge Section 127}
% synthetic bench section 127
% This is a sample LaTeX input file.  (Version of 12 August 2004.)
%
% A '%' character causes TeX to ignore all remaining text on the line,
% and is used for comments like this one.

\documentclass{article}      % Specifies the document class

                             % The preamble begins here.
\title{An Example Document}  % Declares the document's title.
\author{Leslie Lamport}      % Declares the author's name.
\date{January 21, 1994}      % Deleting this command produces today's date.

\newcommand{\ip}[2]{(#1, #2)}
                             % Defines \ip{arg1}{arg2} to mean
                             % (arg1, arg2).

%\newcommand{\ip}[2]{\langle #1 | #2\rangle}
                             % This is an alternative definition of
                             % \ip that is commented out.

\begin{document}             % End of preamble and beginning of text.

\maketitle                   % Produces the title.

This is an example input file.  Comparing it with
the output it generates can show you how to
produce a simple document of your own.

\section{Ordinary Text}      % Produces section heading.  Lower-level
                             % sections are begun with similar
                             % \subsection and \subsubsection commands.

The ends  of words and sentences are marked
  by   spaces. It  doesn't matter how many
spaces    you type; one is as good as 100.  The
end of   a line counts as a space.

One   or more   blank lines denote the  end
of  a paragraph.

Since any number of consecutive spaces are treated
like a single one, the formatting of the input
file makes no difference to
      \LaTeX,                % The \LaTeX command generates the LaTeX logo.
but it makes a difference to you.  When you use
\LaTeX, making your input file as easy to read
as possible will be a great help as you write
your document and when you change it.  This sample
file shows how you can add comments to your own input
file.

Because printing is different from typewriting,
there are a number of things that you have to do
differently when preparing an input file than if
you were just typing the document directly.
Quotation marks like
       ``this''
have to be handled specially, as do quotes within
quotes:
       ``\,`this'            % \, separates the double and single quote.
        is what I just
        wrote, not  `that'\,''.

Dashes come in three sizes: an
       intra-word
dash, a medium dash for number ranges like
       1--2,
and a punctuation
       dash---like
this.

A sentence-ending space should be larger than the
space between words within a sentence.  You
sometimes have to type special commands in
conjunction with punctuation characters to get
this right, as in the following sentence.
       Gnats, gnus, etc.\ all  % `\ ' makes an inter-word space.
       begin with G\@.         % \@ marks end-of-sentence punctuation.
You should check the spaces after periods when
reading your output to make sure you haven't
forgotten any special cases.  Generating an
ellipsis
       \ldots\               % `\ ' is needed after `\ldots' because TeX
                             % ignores spaces after command names like \ldots
                             % made from \ + letters.
                             %
                             % Note how a `%' character causes TeX to ignore
                             % the end of the input line, so these blank lines
                             % do not start a new paragraph.
                             %
with the right spacing around the periods requires
a special command.

\LaTeX\ interprets some common characters as
commands, so you must type special commands to
generate them.  These characters include the
following:
       \$ \& \% \# \{ and \}.

In printing, text is usually emphasized with an
       \emph{italic}
type style.

\begin{em}
   A long segment of text can also be emphasized
   in this way.  Text within such a segment can be
   given \emph{additional} emphasis.
\end{em}

It is sometimes necessary to prevent \LaTeX\ from
breaking a line where it might otherwise do so.
This may be at a space, as between the ``Mr.''\ and
``Jones'' in
       ``Mr.~Jones'',        % ~ produces an unbreakable interword space.
or within a word---especially when the word is a
symbol like
       \mbox{\emph{itemnum}}
that makes little sense when hyphenated across
lines.

Footnotes\footnote{This is an example of a footnote.}
pose no problem.

\LaTeX\ is good at typesetting mathematical formulas
like
       \( x-3y + z = 7 \)
or
       \( a_{1} > x^{2n} + y^{2n} > x' \)
or
       \( \ip{A}{B} = \sum_{i} a_{i} b_{i} \).
The spaces you type in a formula are
ignored.  Remember that a letter like
       $x$                   % $ ... $  and  \( ... \)  are equivalent
is a formula when it denotes a mathematical
symbol, and it should be typed as one.

\section{Displayed Text}

Text is displayed by indenting it from the left
margin.  Quotations are commonly displayed.  There
are short quotations
\begin{quote}
   This is a short quotation.  It consists of a
   single paragraph of text.  See how it is formatted.
\end{quote}
and longer ones.
\begin{quotation}
   This is a longer quotation.  It consists of two
   paragraphs of text, neither of which are
   particularly interesting.

   This is the second paragraph of the quotation.  It
   is just as dull as the first paragraph.
\end{quotation}
Another frequently-displayed structure is a list.
The following is an example of an \emph{itemized}
list.
\begin{itemize}
   \item This is the first item of an itemized list.
         Each item in the list is marked with a ``tick''.
         You don't have to worry about what kind of tick
         mark is used.

   \item This is the second item of the list.  It
         contains another list nested inside it.  The inner
         list is an \emph{enumerated} list.
         \begin{enumerate}
            \item This is the first item of an enumerated
                  list that is nested within the itemized list.

            \item This is the second item of the inner list.
                  \LaTeX\ allows you to nest lists deeper than
                  you really should.
         \end{enumerate}
         This is the rest of the second item of the outer
         list.  It is no more interesting than any other
         part of the item.
   \item This is the third item of the list.
\end{itemize}
You can even display poetry.
\begin{verse}
   There is an environment
    for verse \\             % The \\ command separates lines
   Whose features some poets % within a stanza.
   will curse.

                             % One or more blank lines separate stanzas.

   For instead of making\\
   Them do \emph{all} line breaking, \\
   It allows them to put too many words on a line when they'd rather be
   forced to be terse.
\end{verse}

Mathematical formulas may also be displayed.  A
displayed formula
is
one-line long; multiline
formulas require special formatting instructions.
   \[  \ip{\Gamma}{\psi'} = x'' + y^{2} + z_{i}^{n}\]
Don't start a paragraph with a displayed equation,
nor make one a paragraph by itself.

\end{document}               % End of document.

\section{Bench Huge Section 128}
% synthetic bench section 128
% This is a sample LaTeX input file.  (Version of 12 August 2004.)
%
% A '%' character causes TeX to ignore all remaining text on the line,
% and is used for comments like this one.

\documentclass{article}      % Specifies the document class

                             % The preamble begins here.
\title{An Example Document}  % Declares the document's title.
\author{Leslie Lamport}      % Declares the author's name.
\date{January 21, 1994}      % Deleting this command produces today's date.

\newcommand{\ip}[2]{(#1, #2)}
                             % Defines \ip{arg1}{arg2} to mean
                             % (arg1, arg2).

%\newcommand{\ip}[2]{\langle #1 | #2\rangle}
                             % This is an alternative definition of
                             % \ip that is commented out.

\begin{document}             % End of preamble and beginning of text.

\maketitle                   % Produces the title.

This is an example input file.  Comparing it with
the output it generates can show you how to
produce a simple document of your own.

\section{Ordinary Text}      % Produces section heading.  Lower-level
                             % sections are begun with similar
                             % \subsection and \subsubsection commands.

The ends  of words and sentences are marked
  by   spaces. It  doesn't matter how many
spaces    you type; one is as good as 100.  The
end of   a line counts as a space.

One   or more   blank lines denote the  end
of  a paragraph.

Since any number of consecutive spaces are treated
like a single one, the formatting of the input
file makes no difference to
      \LaTeX,                % The \LaTeX command generates the LaTeX logo.
but it makes a difference to you.  When you use
\LaTeX, making your input file as easy to read
as possible will be a great help as you write
your document and when you change it.  This sample
file shows how you can add comments to your own input
file.

Because printing is different from typewriting,
there are a number of things that you have to do
differently when preparing an input file than if
you were just typing the document directly.
Quotation marks like
       ``this''
have to be handled specially, as do quotes within
quotes:
       ``\,`this'            % \, separates the double and single quote.
        is what I just
        wrote, not  `that'\,''.

Dashes come in three sizes: an
       intra-word
dash, a medium dash for number ranges like
       1--2,
and a punctuation
       dash---like
this.

A sentence-ending space should be larger than the
space between words within a sentence.  You
sometimes have to type special commands in
conjunction with punctuation characters to get
this right, as in the following sentence.
       Gnats, gnus, etc.\ all  % `\ ' makes an inter-word space.
       begin with G\@.         % \@ marks end-of-sentence punctuation.
You should check the spaces after periods when
reading your output to make sure you haven't
forgotten any special cases.  Generating an
ellipsis
       \ldots\               % `\ ' is needed after `\ldots' because TeX
                             % ignores spaces after command names like \ldots
                             % made from \ + letters.
                             %
                             % Note how a `%' character causes TeX to ignore
                             % the end of the input line, so these blank lines
                             % do not start a new paragraph.
                             %
with the right spacing around the periods requires
a special command.

\LaTeX\ interprets some common characters as
commands, so you must type special commands to
generate them.  These characters include the
following:
       \$ \& \% \# \{ and \}.

In printing, text is usually emphasized with an
       \emph{italic}
type style.

\begin{em}
   A long segment of text can also be emphasized
   in this way.  Text within such a segment can be
   given \emph{additional} emphasis.
\end{em}

It is sometimes necessary to prevent \LaTeX\ from
breaking a line where it might otherwise do so.
This may be at a space, as between the ``Mr.''\ and
``Jones'' in
       ``Mr.~Jones'',        % ~ produces an unbreakable interword space.
or within a word---especially when the word is a
symbol like
       \mbox{\emph{itemnum}}
that makes little sense when hyphenated across
lines.

Footnotes\footnote{This is an example of a footnote.}
pose no problem.

\LaTeX\ is good at typesetting mathematical formulas
like
       \( x-3y + z = 7 \)
or
       \( a_{1} > x^{2n} + y^{2n} > x' \)
or
       \( \ip{A}{B} = \sum_{i} a_{i} b_{i} \).
The spaces you type in a formula are
ignored.  Remember that a letter like
       $x$                   % $ ... $  and  \( ... \)  are equivalent
is a formula when it denotes a mathematical
symbol, and it should be typed as one.

\section{Displayed Text}

Text is displayed by indenting it from the left
margin.  Quotations are commonly displayed.  There
are short quotations
\begin{quote}
   This is a short quotation.  It consists of a
   single paragraph of text.  See how it is formatted.
\end{quote}
and longer ones.
\begin{quotation}
   This is a longer quotation.  It consists of two
   paragraphs of text, neither of which are
   particularly interesting.

   This is the second paragraph of the quotation.  It
   is just as dull as the first paragraph.
\end{quotation}
Another frequently-displayed structure is a list.
The following is an example of an \emph{itemized}
list.
\begin{itemize}
   \item This is the first item of an itemized list.
         Each item in the list is marked with a ``tick''.
         You don't have to worry about what kind of tick
         mark is used.

   \item This is the second item of the list.  It
         contains another list nested inside it.  The inner
         list is an \emph{enumerated} list.
         \begin{enumerate}
            \item This is the first item of an enumerated
                  list that is nested within the itemized list.

            \item This is the second item of the inner list.
                  \LaTeX\ allows you to nest lists deeper than
                  you really should.
         \end{enumerate}
         This is the rest of the second item of the outer
         list.  It is no more interesting than any other
         part of the item.
   \item This is the third item of the list.
\end{itemize}
You can even display poetry.
\begin{verse}
   There is an environment
    for verse \\             % The \\ command separates lines
   Whose features some poets % within a stanza.
   will curse.

                             % One or more blank lines separate stanzas.

   For instead of making\\
   Them do \emph{all} line breaking, \\
   It allows them to put too many words on a line when they'd rather be
   forced to be terse.
\end{verse}

Mathematical formulas may also be displayed.  A
displayed formula
is
one-line long; multiline
formulas require special formatting instructions.
   \[  \ip{\Gamma}{\psi'} = x'' + y^{2} + z_{i}^{n}\]
Don't start a paragraph with a displayed equation,
nor make one a paragraph by itself.

\end{document}               % End of document.

\section{Bench Huge Section 129}
% synthetic bench section 129
% This is a sample LaTeX input file.  (Version of 12 August 2004.)
%
% A '%' character causes TeX to ignore all remaining text on the line,
% and is used for comments like this one.

\documentclass{article}      % Specifies the document class

                             % The preamble begins here.
\title{An Example Document}  % Declares the document's title.
\author{Leslie Lamport}      % Declares the author's name.
\date{January 21, 1994}      % Deleting this command produces today's date.

\newcommand{\ip}[2]{(#1, #2)}
                             % Defines \ip{arg1}{arg2} to mean
                             % (arg1, arg2).

%\newcommand{\ip}[2]{\langle #1 | #2\rangle}
                             % This is an alternative definition of
                             % \ip that is commented out.

\begin{document}             % End of preamble and beginning of text.

\maketitle                   % Produces the title.

This is an example input file.  Comparing it with
the output it generates can show you how to
produce a simple document of your own.

\section{Ordinary Text}      % Produces section heading.  Lower-level
                             % sections are begun with similar
                             % \subsection and \subsubsection commands.

The ends  of words and sentences are marked
  by   spaces. It  doesn't matter how many
spaces    you type; one is as good as 100.  The
end of   a line counts as a space.

One   or more   blank lines denote the  end
of  a paragraph.

Since any number of consecutive spaces are treated
like a single one, the formatting of the input
file makes no difference to
      \LaTeX,                % The \LaTeX command generates the LaTeX logo.
but it makes a difference to you.  When you use
\LaTeX, making your input file as easy to read
as possible will be a great help as you write
your document and when you change it.  This sample
file shows how you can add comments to your own input
file.

Because printing is different from typewriting,
there are a number of things that you have to do
differently when preparing an input file than if
you were just typing the document directly.
Quotation marks like
       ``this''
have to be handled specially, as do quotes within
quotes:
       ``\,`this'            % \, separates the double and single quote.
        is what I just
        wrote, not  `that'\,''.

Dashes come in three sizes: an
       intra-word
dash, a medium dash for number ranges like
       1--2,
and a punctuation
       dash---like
this.

A sentence-ending space should be larger than the
space between words within a sentence.  You
sometimes have to type special commands in
conjunction with punctuation characters to get
this right, as in the following sentence.
       Gnats, gnus, etc.\ all  % `\ ' makes an inter-word space.
       begin with G\@.         % \@ marks end-of-sentence punctuation.
You should check the spaces after periods when
reading your output to make sure you haven't
forgotten any special cases.  Generating an
ellipsis
       \ldots\               % `\ ' is needed after `\ldots' because TeX
                             % ignores spaces after command names like \ldots
                             % made from \ + letters.
                             %
                             % Note how a `%' character causes TeX to ignore
                             % the end of the input line, so these blank lines
                             % do not start a new paragraph.
                             %
with the right spacing around the periods requires
a special command.

\LaTeX\ interprets some common characters as
commands, so you must type special commands to
generate them.  These characters include the
following:
       \$ \& \% \# \{ and \}.

In printing, text is usually emphasized with an
       \emph{italic}
type style.

\begin{em}
   A long segment of text can also be emphasized
   in this way.  Text within such a segment can be
   given \emph{additional} emphasis.
\end{em}

It is sometimes necessary to prevent \LaTeX\ from
breaking a line where it might otherwise do so.
This may be at a space, as between the ``Mr.''\ and
``Jones'' in
       ``Mr.~Jones'',        % ~ produces an unbreakable interword space.
or within a word---especially when the word is a
symbol like
       \mbox{\emph{itemnum}}
that makes little sense when hyphenated across
lines.

Footnotes\footnote{This is an example of a footnote.}
pose no problem.

\LaTeX\ is good at typesetting mathematical formulas
like
       \( x-3y + z = 7 \)
or
       \( a_{1} > x^{2n} + y^{2n} > x' \)
or
       \( \ip{A}{B} = \sum_{i} a_{i} b_{i} \).
The spaces you type in a formula are
ignored.  Remember that a letter like
       $x$                   % $ ... $  and  \( ... \)  are equivalent
is a formula when it denotes a mathematical
symbol, and it should be typed as one.

\section{Displayed Text}

Text is displayed by indenting it from the left
margin.  Quotations are commonly displayed.  There
are short quotations
\begin{quote}
   This is a short quotation.  It consists of a
   single paragraph of text.  See how it is formatted.
\end{quote}
and longer ones.
\begin{quotation}
   This is a longer quotation.  It consists of two
   paragraphs of text, neither of which are
   particularly interesting.

   This is the second paragraph of the quotation.  It
   is just as dull as the first paragraph.
\end{quotation}
Another frequently-displayed structure is a list.
The following is an example of an \emph{itemized}
list.
\begin{itemize}
   \item This is the first item of an itemized list.
         Each item in the list is marked with a ``tick''.
         You don't have to worry about what kind of tick
         mark is used.

   \item This is the second item of the list.  It
         contains another list nested inside it.  The inner
         list is an \emph{enumerated} list.
         \begin{enumerate}
            \item This is the first item of an enumerated
                  list that is nested within the itemized list.

            \item This is the second item of the inner list.
                  \LaTeX\ allows you to nest lists deeper than
                  you really should.
         \end{enumerate}
         This is the rest of the second item of the outer
         list.  It is no more interesting than any other
         part of the item.
   \item This is the third item of the list.
\end{itemize}
You can even display poetry.
\begin{verse}
   There is an environment
    for verse \\             % The \\ command separates lines
   Whose features some poets % within a stanza.
   will curse.

                             % One or more blank lines separate stanzas.

   For instead of making\\
   Them do \emph{all} line breaking, \\
   It allows them to put too many words on a line when they'd rather be
   forced to be terse.
\end{verse}

Mathematical formulas may also be displayed.  A
displayed formula
is
one-line long; multiline
formulas require special formatting instructions.
   \[  \ip{\Gamma}{\psi'} = x'' + y^{2} + z_{i}^{n}\]
Don't start a paragraph with a displayed equation,
nor make one a paragraph by itself.

\end{document}               % End of document.

\section{Bench Huge Section 130}
% synthetic bench section 130
% This is a sample LaTeX input file.  (Version of 12 August 2004.)
%
% A '%' character causes TeX to ignore all remaining text on the line,
% and is used for comments like this one.

\documentclass{article}      % Specifies the document class

                             % The preamble begins here.
\title{An Example Document}  % Declares the document's title.
\author{Leslie Lamport}      % Declares the author's name.
\date{January 21, 1994}      % Deleting this command produces today's date.

\newcommand{\ip}[2]{(#1, #2)}
                             % Defines \ip{arg1}{arg2} to mean
                             % (arg1, arg2).

%\newcommand{\ip}[2]{\langle #1 | #2\rangle}
                             % This is an alternative definition of
                             % \ip that is commented out.

\begin{document}             % End of preamble and beginning of text.

\maketitle                   % Produces the title.

This is an example input file.  Comparing it with
the output it generates can show you how to
produce a simple document of your own.

\section{Ordinary Text}      % Produces section heading.  Lower-level
                             % sections are begun with similar
                             % \subsection and \subsubsection commands.

The ends  of words and sentences are marked
  by   spaces. It  doesn't matter how many
spaces    you type; one is as good as 100.  The
end of   a line counts as a space.

One   or more   blank lines denote the  end
of  a paragraph.

Since any number of consecutive spaces are treated
like a single one, the formatting of the input
file makes no difference to
      \LaTeX,                % The \LaTeX command generates the LaTeX logo.
but it makes a difference to you.  When you use
\LaTeX, making your input file as easy to read
as possible will be a great help as you write
your document and when you change it.  This sample
file shows how you can add comments to your own input
file.

Because printing is different from typewriting,
there are a number of things that you have to do
differently when preparing an input file than if
you were just typing the document directly.
Quotation marks like
       ``this''
have to be handled specially, as do quotes within
quotes:
       ``\,`this'            % \, separates the double and single quote.
        is what I just
        wrote, not  `that'\,''.

Dashes come in three sizes: an
       intra-word
dash, a medium dash for number ranges like
       1--2,
and a punctuation
       dash---like
this.

A sentence-ending space should be larger than the
space between words within a sentence.  You
sometimes have to type special commands in
conjunction with punctuation characters to get
this right, as in the following sentence.
       Gnats, gnus, etc.\ all  % `\ ' makes an inter-word space.
       begin with G\@.         % \@ marks end-of-sentence punctuation.
You should check the spaces after periods when
reading your output to make sure you haven't
forgotten any special cases.  Generating an
ellipsis
       \ldots\               % `\ ' is needed after `\ldots' because TeX
                             % ignores spaces after command names like \ldots
                             % made from \ + letters.
                             %
                             % Note how a `%' character causes TeX to ignore
                             % the end of the input line, so these blank lines
                             % do not start a new paragraph.
                             %
with the right spacing around the periods requires
a special command.

\LaTeX\ interprets some common characters as
commands, so you must type special commands to
generate them.  These characters include the
following:
       \$ \& \% \# \{ and \}.

In printing, text is usually emphasized with an
       \emph{italic}
type style.

\begin{em}
   A long segment of text can also be emphasized
   in this way.  Text within such a segment can be
   given \emph{additional} emphasis.
\end{em}

It is sometimes necessary to prevent \LaTeX\ from
breaking a line where it might otherwise do so.
This may be at a space, as between the ``Mr.''\ and
``Jones'' in
       ``Mr.~Jones'',        % ~ produces an unbreakable interword space.
or within a word---especially when the word is a
symbol like
       \mbox{\emph{itemnum}}
that makes little sense when hyphenated across
lines.

Footnotes\footnote{This is an example of a footnote.}
pose no problem.

\LaTeX\ is good at typesetting mathematical formulas
like
       \( x-3y + z = 7 \)
or
       \( a_{1} > x^{2n} + y^{2n} > x' \)
or
       \( \ip{A}{B} = \sum_{i} a_{i} b_{i} \).
The spaces you type in a formula are
ignored.  Remember that a letter like
       $x$                   % $ ... $  and  \( ... \)  are equivalent
is a formula when it denotes a mathematical
symbol, and it should be typed as one.

\section{Displayed Text}

Text is displayed by indenting it from the left
margin.  Quotations are commonly displayed.  There
are short quotations
\begin{quote}
   This is a short quotation.  It consists of a
   single paragraph of text.  See how it is formatted.
\end{quote}
and longer ones.
\begin{quotation}
   This is a longer quotation.  It consists of two
   paragraphs of text, neither of which are
   particularly interesting.

   This is the second paragraph of the quotation.  It
   is just as dull as the first paragraph.
\end{quotation}
Another frequently-displayed structure is a list.
The following is an example of an \emph{itemized}
list.
\begin{itemize}
   \item This is the first item of an itemized list.
         Each item in the list is marked with a ``tick''.
         You don't have to worry about what kind of tick
         mark is used.

   \item This is the second item of the list.  It
         contains another list nested inside it.  The inner
         list is an \emph{enumerated} list.
         \begin{enumerate}
            \item This is the first item of an enumerated
                  list that is nested within the itemized list.

            \item This is the second item of the inner list.
                  \LaTeX\ allows you to nest lists deeper than
                  you really should.
         \end{enumerate}
         This is the rest of the second item of the outer
         list.  It is no more interesting than any other
         part of the item.
   \item This is the third item of the list.
\end{itemize}
You can even display poetry.
\begin{verse}
   There is an environment
    for verse \\             % The \\ command separates lines
   Whose features some poets % within a stanza.
   will curse.

                             % One or more blank lines separate stanzas.

   For instead of making\\
   Them do \emph{all} line breaking, \\
   It allows them to put too many words on a line when they'd rather be
   forced to be terse.
\end{verse}

Mathematical formulas may also be displayed.  A
displayed formula
is
one-line long; multiline
formulas require special formatting instructions.
   \[  \ip{\Gamma}{\psi'} = x'' + y^{2} + z_{i}^{n}\]
Don't start a paragraph with a displayed equation,
nor make one a paragraph by itself.

\end{document}               % End of document.

\section{Bench Huge Section 131}
% synthetic bench section 131
% This is a sample LaTeX input file.  (Version of 12 August 2004.)
%
% A '%' character causes TeX to ignore all remaining text on the line,
% and is used for comments like this one.

\documentclass{article}      % Specifies the document class

                             % The preamble begins here.
\title{An Example Document}  % Declares the document's title.
\author{Leslie Lamport}      % Declares the author's name.
\date{January 21, 1994}      % Deleting this command produces today's date.

\newcommand{\ip}[2]{(#1, #2)}
                             % Defines \ip{arg1}{arg2} to mean
                             % (arg1, arg2).

%\newcommand{\ip}[2]{\langle #1 | #2\rangle}
                             % This is an alternative definition of
                             % \ip that is commented out.

\begin{document}             % End of preamble and beginning of text.

\maketitle                   % Produces the title.

This is an example input file.  Comparing it with
the output it generates can show you how to
produce a simple document of your own.

\section{Ordinary Text}      % Produces section heading.  Lower-level
                             % sections are begun with similar
                             % \subsection and \subsubsection commands.

The ends  of words and sentences are marked
  by   spaces. It  doesn't matter how many
spaces    you type; one is as good as 100.  The
end of   a line counts as a space.

One   or more   blank lines denote the  end
of  a paragraph.

Since any number of consecutive spaces are treated
like a single one, the formatting of the input
file makes no difference to
      \LaTeX,                % The \LaTeX command generates the LaTeX logo.
but it makes a difference to you.  When you use
\LaTeX, making your input file as easy to read
as possible will be a great help as you write
your document and when you change it.  This sample
file shows how you can add comments to your own input
file.

Because printing is different from typewriting,
there are a number of things that you have to do
differently when preparing an input file than if
you were just typing the document directly.
Quotation marks like
       ``this''
have to be handled specially, as do quotes within
quotes:
       ``\,`this'            % \, separates the double and single quote.
        is what I just
        wrote, not  `that'\,''.

Dashes come in three sizes: an
       intra-word
dash, a medium dash for number ranges like
       1--2,
and a punctuation
       dash---like
this.

A sentence-ending space should be larger than the
space between words within a sentence.  You
sometimes have to type special commands in
conjunction with punctuation characters to get
this right, as in the following sentence.
       Gnats, gnus, etc.\ all  % `\ ' makes an inter-word space.
       begin with G\@.         % \@ marks end-of-sentence punctuation.
You should check the spaces after periods when
reading your output to make sure you haven't
forgotten any special cases.  Generating an
ellipsis
       \ldots\               % `\ ' is needed after `\ldots' because TeX
                             % ignores spaces after command names like \ldots
                             % made from \ + letters.
                             %
                             % Note how a `%' character causes TeX to ignore
                             % the end of the input line, so these blank lines
                             % do not start a new paragraph.
                             %
with the right spacing around the periods requires
a special command.

\LaTeX\ interprets some common characters as
commands, so you must type special commands to
generate them.  These characters include the
following:
       \$ \& \% \# \{ and \}.

In printing, text is usually emphasized with an
       \emph{italic}
type style.

\begin{em}
   A long segment of text can also be emphasized
   in this way.  Text within such a segment can be
   given \emph{additional} emphasis.
\end{em}

It is sometimes necessary to prevent \LaTeX\ from
breaking a line where it might otherwise do so.
This may be at a space, as between the ``Mr.''\ and
``Jones'' in
       ``Mr.~Jones'',        % ~ produces an unbreakable interword space.
or within a word---especially when the word is a
symbol like
       \mbox{\emph{itemnum}}
that makes little sense when hyphenated across
lines.

Footnotes\footnote{This is an example of a footnote.}
pose no problem.

\LaTeX\ is good at typesetting mathematical formulas
like
       \( x-3y + z = 7 \)
or
       \( a_{1} > x^{2n} + y^{2n} > x' \)
or
       \( \ip{A}{B} = \sum_{i} a_{i} b_{i} \).
The spaces you type in a formula are
ignored.  Remember that a letter like
       $x$                   % $ ... $  and  \( ... \)  are equivalent
is a formula when it denotes a mathematical
symbol, and it should be typed as one.

\section{Displayed Text}

Text is displayed by indenting it from the left
margin.  Quotations are commonly displayed.  There
are short quotations
\begin{quote}
   This is a short quotation.  It consists of a
   single paragraph of text.  See how it is formatted.
\end{quote}
and longer ones.
\begin{quotation}
   This is a longer quotation.  It consists of two
   paragraphs of text, neither of which are
   particularly interesting.

   This is the second paragraph of the quotation.  It
   is just as dull as the first paragraph.
\end{quotation}
Another frequently-displayed structure is a list.
The following is an example of an \emph{itemized}
list.
\begin{itemize}
   \item This is the first item of an itemized list.
         Each item in the list is marked with a ``tick''.
         You don't have to worry about what kind of tick
         mark is used.

   \item This is the second item of the list.  It
         contains another list nested inside it.  The inner
         list is an \emph{enumerated} list.
         \begin{enumerate}
            \item This is the first item of an enumerated
                  list that is nested within the itemized list.

            \item This is the second item of the inner list.
                  \LaTeX\ allows you to nest lists deeper than
                  you really should.
         \end{enumerate}
         This is the rest of the second item of the outer
         list.  It is no more interesting than any other
         part of the item.
   \item This is the third item of the list.
\end{itemize}
You can even display poetry.
\begin{verse}
   There is an environment
    for verse \\             % The \\ command separates lines
   Whose features some poets % within a stanza.
   will curse.

                             % One or more blank lines separate stanzas.

   For instead of making\\
   Them do \emph{all} line breaking, \\
   It allows them to put too many words on a line when they'd rather be
   forced to be terse.
\end{verse}

Mathematical formulas may also be displayed.  A
displayed formula
is
one-line long; multiline
formulas require special formatting instructions.
   \[  \ip{\Gamma}{\psi'} = x'' + y^{2} + z_{i}^{n}\]
Don't start a paragraph with a displayed equation,
nor make one a paragraph by itself.

\end{document}               % End of document.

\section{Bench Huge Section 132}
% synthetic bench section 132
% This is a sample LaTeX input file.  (Version of 12 August 2004.)
%
% A '%' character causes TeX to ignore all remaining text on the line,
% and is used for comments like this one.

\documentclass{article}      % Specifies the document class

                             % The preamble begins here.
\title{An Example Document}  % Declares the document's title.
\author{Leslie Lamport}      % Declares the author's name.
\date{January 21, 1994}      % Deleting this command produces today's date.

\newcommand{\ip}[2]{(#1, #2)}
                             % Defines \ip{arg1}{arg2} to mean
                             % (arg1, arg2).

%\newcommand{\ip}[2]{\langle #1 | #2\rangle}
                             % This is an alternative definition of
                             % \ip that is commented out.

\begin{document}             % End of preamble and beginning of text.

\maketitle                   % Produces the title.

This is an example input file.  Comparing it with
the output it generates can show you how to
produce a simple document of your own.

\section{Ordinary Text}      % Produces section heading.  Lower-level
                             % sections are begun with similar
                             % \subsection and \subsubsection commands.

The ends  of words and sentences are marked
  by   spaces. It  doesn't matter how many
spaces    you type; one is as good as 100.  The
end of   a line counts as a space.

One   or more   blank lines denote the  end
of  a paragraph.

Since any number of consecutive spaces are treated
like a single one, the formatting of the input
file makes no difference to
      \LaTeX,                % The \LaTeX command generates the LaTeX logo.
but it makes a difference to you.  When you use
\LaTeX, making your input file as easy to read
as possible will be a great help as you write
your document and when you change it.  This sample
file shows how you can add comments to your own input
file.

Because printing is different from typewriting,
there are a number of things that you have to do
differently when preparing an input file than if
you were just typing the document directly.
Quotation marks like
       ``this''
have to be handled specially, as do quotes within
quotes:
       ``\,`this'            % \, separates the double and single quote.
        is what I just
        wrote, not  `that'\,''.

Dashes come in three sizes: an
       intra-word
dash, a medium dash for number ranges like
       1--2,
and a punctuation
       dash---like
this.

A sentence-ending space should be larger than the
space between words within a sentence.  You
sometimes have to type special commands in
conjunction with punctuation characters to get
this right, as in the following sentence.
       Gnats, gnus, etc.\ all  % `\ ' makes an inter-word space.
       begin with G\@.         % \@ marks end-of-sentence punctuation.
You should check the spaces after periods when
reading your output to make sure you haven't
forgotten any special cases.  Generating an
ellipsis
       \ldots\               % `\ ' is needed after `\ldots' because TeX
                             % ignores spaces after command names like \ldots
                             % made from \ + letters.
                             %
                             % Note how a `%' character causes TeX to ignore
                             % the end of the input line, so these blank lines
                             % do not start a new paragraph.
                             %
with the right spacing around the periods requires
a special command.

\LaTeX\ interprets some common characters as
commands, so you must type special commands to
generate them.  These characters include the
following:
       \$ \& \% \# \{ and \}.

In printing, text is usually emphasized with an
       \emph{italic}
type style.

\begin{em}
   A long segment of text can also be emphasized
   in this way.  Text within such a segment can be
   given \emph{additional} emphasis.
\end{em}

It is sometimes necessary to prevent \LaTeX\ from
breaking a line where it might otherwise do so.
This may be at a space, as between the ``Mr.''\ and
``Jones'' in
       ``Mr.~Jones'',        % ~ produces an unbreakable interword space.
or within a word---especially when the word is a
symbol like
       \mbox{\emph{itemnum}}
that makes little sense when hyphenated across
lines.

Footnotes\footnote{This is an example of a footnote.}
pose no problem.

\LaTeX\ is good at typesetting mathematical formulas
like
       \( x-3y + z = 7 \)
or
       \( a_{1} > x^{2n} + y^{2n} > x' \)
or
       \( \ip{A}{B} = \sum_{i} a_{i} b_{i} \).
The spaces you type in a formula are
ignored.  Remember that a letter like
       $x$                   % $ ... $  and  \( ... \)  are equivalent
is a formula when it denotes a mathematical
symbol, and it should be typed as one.

\section{Displayed Text}

Text is displayed by indenting it from the left
margin.  Quotations are commonly displayed.  There
are short quotations
\begin{quote}
   This is a short quotation.  It consists of a
   single paragraph of text.  See how it is formatted.
\end{quote}
and longer ones.
\begin{quotation}
   This is a longer quotation.  It consists of two
   paragraphs of text, neither of which are
   particularly interesting.

   This is the second paragraph of the quotation.  It
   is just as dull as the first paragraph.
\end{quotation}
Another frequently-displayed structure is a list.
The following is an example of an \emph{itemized}
list.
\begin{itemize}
   \item This is the first item of an itemized list.
         Each item in the list is marked with a ``tick''.
         You don't have to worry about what kind of tick
         mark is used.

   \item This is the second item of the list.  It
         contains another list nested inside it.  The inner
         list is an \emph{enumerated} list.
         \begin{enumerate}
            \item This is the first item of an enumerated
                  list that is nested within the itemized list.

            \item This is the second item of the inner list.
                  \LaTeX\ allows you to nest lists deeper than
                  you really should.
         \end{enumerate}
         This is the rest of the second item of the outer
         list.  It is no more interesting than any other
         part of the item.
   \item This is the third item of the list.
\end{itemize}
You can even display poetry.
\begin{verse}
   There is an environment
    for verse \\             % The \\ command separates lines
   Whose features some poets % within a stanza.
   will curse.

                             % One or more blank lines separate stanzas.

   For instead of making\\
   Them do \emph{all} line breaking, \\
   It allows them to put too many words on a line when they'd rather be
   forced to be terse.
\end{verse}

Mathematical formulas may also be displayed.  A
displayed formula
is
one-line long; multiline
formulas require special formatting instructions.
   \[  \ip{\Gamma}{\psi'} = x'' + y^{2} + z_{i}^{n}\]
Don't start a paragraph with a displayed equation,
nor make one a paragraph by itself.

\end{document}               % End of document.

\section{Bench Huge Section 133}
% synthetic bench section 133
% This is a sample LaTeX input file.  (Version of 12 August 2004.)
%
% A '%' character causes TeX to ignore all remaining text on the line,
% and is used for comments like this one.

\documentclass{article}      % Specifies the document class

                             % The preamble begins here.
\title{An Example Document}  % Declares the document's title.
\author{Leslie Lamport}      % Declares the author's name.
\date{January 21, 1994}      % Deleting this command produces today's date.

\newcommand{\ip}[2]{(#1, #2)}
                             % Defines \ip{arg1}{arg2} to mean
                             % (arg1, arg2).

%\newcommand{\ip}[2]{\langle #1 | #2\rangle}
                             % This is an alternative definition of
                             % \ip that is commented out.

\begin{document}             % End of preamble and beginning of text.

\maketitle                   % Produces the title.

This is an example input file.  Comparing it with
the output it generates can show you how to
produce a simple document of your own.

\section{Ordinary Text}      % Produces section heading.  Lower-level
                             % sections are begun with similar
                             % \subsection and \subsubsection commands.

The ends  of words and sentences are marked
  by   spaces. It  doesn't matter how many
spaces    you type; one is as good as 100.  The
end of   a line counts as a space.

One   or more   blank lines denote the  end
of  a paragraph.

Since any number of consecutive spaces are treated
like a single one, the formatting of the input
file makes no difference to
      \LaTeX,                % The \LaTeX command generates the LaTeX logo.
but it makes a difference to you.  When you use
\LaTeX, making your input file as easy to read
as possible will be a great help as you write
your document and when you change it.  This sample
file shows how you can add comments to your own input
file.

Because printing is different from typewriting,
there are a number of things that you have to do
differently when preparing an input file than if
you were just typing the document directly.
Quotation marks like
       ``this''
have to be handled specially, as do quotes within
quotes:
       ``\,`this'            % \, separates the double and single quote.
        is what I just
        wrote, not  `that'\,''.

Dashes come in three sizes: an
       intra-word
dash, a medium dash for number ranges like
       1--2,
and a punctuation
       dash---like
this.

A sentence-ending space should be larger than the
space between words within a sentence.  You
sometimes have to type special commands in
conjunction with punctuation characters to get
this right, as in the following sentence.
       Gnats, gnus, etc.\ all  % `\ ' makes an inter-word space.
       begin with G\@.         % \@ marks end-of-sentence punctuation.
You should check the spaces after periods when
reading your output to make sure you haven't
forgotten any special cases.  Generating an
ellipsis
       \ldots\               % `\ ' is needed after `\ldots' because TeX
                             % ignores spaces after command names like \ldots
                             % made from \ + letters.
                             %
                             % Note how a `%' character causes TeX to ignore
                             % the end of the input line, so these blank lines
                             % do not start a new paragraph.
                             %
with the right spacing around the periods requires
a special command.

\LaTeX\ interprets some common characters as
commands, so you must type special commands to
generate them.  These characters include the
following:
       \$ \& \% \# \{ and \}.

In printing, text is usually emphasized with an
       \emph{italic}
type style.

\begin{em}
   A long segment of text can also be emphasized
   in this way.  Text within such a segment can be
   given \emph{additional} emphasis.
\end{em}

It is sometimes necessary to prevent \LaTeX\ from
breaking a line where it might otherwise do so.
This may be at a space, as between the ``Mr.''\ and
``Jones'' in
       ``Mr.~Jones'',        % ~ produces an unbreakable interword space.
or within a word---especially when the word is a
symbol like
       \mbox{\emph{itemnum}}
that makes little sense when hyphenated across
lines.

Footnotes\footnote{This is an example of a footnote.}
pose no problem.

\LaTeX\ is good at typesetting mathematical formulas
like
       \( x-3y + z = 7 \)
or
       \( a_{1} > x^{2n} + y^{2n} > x' \)
or
       \( \ip{A}{B} = \sum_{i} a_{i} b_{i} \).
The spaces you type in a formula are
ignored.  Remember that a letter like
       $x$                   % $ ... $  and  \( ... \)  are equivalent
is a formula when it denotes a mathematical
symbol, and it should be typed as one.

\section{Displayed Text}

Text is displayed by indenting it from the left
margin.  Quotations are commonly displayed.  There
are short quotations
\begin{quote}
   This is a short quotation.  It consists of a
   single paragraph of text.  See how it is formatted.
\end{quote}
and longer ones.
\begin{quotation}
   This is a longer quotation.  It consists of two
   paragraphs of text, neither of which are
   particularly interesting.

   This is the second paragraph of the quotation.  It
   is just as dull as the first paragraph.
\end{quotation}
Another frequently-displayed structure is a list.
The following is an example of an \emph{itemized}
list.
\begin{itemize}
   \item This is the first item of an itemized list.
         Each item in the list is marked with a ``tick''.
         You don't have to worry about what kind of tick
         mark is used.

   \item This is the second item of the list.  It
         contains another list nested inside it.  The inner
         list is an \emph{enumerated} list.
         \begin{enumerate}
            \item This is the first item of an enumerated
                  list that is nested within the itemized list.

            \item This is the second item of the inner list.
                  \LaTeX\ allows you to nest lists deeper than
                  you really should.
         \end{enumerate}
         This is the rest of the second item of the outer
         list.  It is no more interesting than any other
         part of the item.
   \item This is the third item of the list.
\end{itemize}
You can even display poetry.
\begin{verse}
   There is an environment
    for verse \\             % The \\ command separates lines
   Whose features some poets % within a stanza.
   will curse.

                             % One or more blank lines separate stanzas.

   For instead of making\\
   Them do \emph{all} line breaking, \\
   It allows them to put too many words on a line when they'd rather be
   forced to be terse.
\end{verse}

Mathematical formulas may also be displayed.  A
displayed formula
is
one-line long; multiline
formulas require special formatting instructions.
   \[  \ip{\Gamma}{\psi'} = x'' + y^{2} + z_{i}^{n}\]
Don't start a paragraph with a displayed equation,
nor make one a paragraph by itself.

\end{document}               % End of document.

\section{Bench Huge Section 134}
% synthetic bench section 134
% This is a sample LaTeX input file.  (Version of 12 August 2004.)
%
% A '%' character causes TeX to ignore all remaining text on the line,
% and is used for comments like this one.

\documentclass{article}      % Specifies the document class

                             % The preamble begins here.
\title{An Example Document}  % Declares the document's title.
\author{Leslie Lamport}      % Declares the author's name.
\date{January 21, 1994}      % Deleting this command produces today's date.

\newcommand{\ip}[2]{(#1, #2)}
                             % Defines \ip{arg1}{arg2} to mean
                             % (arg1, arg2).

%\newcommand{\ip}[2]{\langle #1 | #2\rangle}
                             % This is an alternative definition of
                             % \ip that is commented out.

\begin{document}             % End of preamble and beginning of text.

\maketitle                   % Produces the title.

This is an example input file.  Comparing it with
the output it generates can show you how to
produce a simple document of your own.

\section{Ordinary Text}      % Produces section heading.  Lower-level
                             % sections are begun with similar
                             % \subsection and \subsubsection commands.

The ends  of words and sentences are marked
  by   spaces. It  doesn't matter how many
spaces    you type; one is as good as 100.  The
end of   a line counts as a space.

One   or more   blank lines denote the  end
of  a paragraph.

Since any number of consecutive spaces are treated
like a single one, the formatting of the input
file makes no difference to
      \LaTeX,                % The \LaTeX command generates the LaTeX logo.
but it makes a difference to you.  When you use
\LaTeX, making your input file as easy to read
as possible will be a great help as you write
your document and when you change it.  This sample
file shows how you can add comments to your own input
file.

Because printing is different from typewriting,
there are a number of things that you have to do
differently when preparing an input file than if
you were just typing the document directly.
Quotation marks like
       ``this''
have to be handled specially, as do quotes within
quotes:
       ``\,`this'            % \, separates the double and single quote.
        is what I just
        wrote, not  `that'\,''.

Dashes come in three sizes: an
       intra-word
dash, a medium dash for number ranges like
       1--2,
and a punctuation
       dash---like
this.

A sentence-ending space should be larger than the
space between words within a sentence.  You
sometimes have to type special commands in
conjunction with punctuation characters to get
this right, as in the following sentence.
       Gnats, gnus, etc.\ all  % `\ ' makes an inter-word space.
       begin with G\@.         % \@ marks end-of-sentence punctuation.
You should check the spaces after periods when
reading your output to make sure you haven't
forgotten any special cases.  Generating an
ellipsis
       \ldots\               % `\ ' is needed after `\ldots' because TeX
                             % ignores spaces after command names like \ldots
                             % made from \ + letters.
                             %
                             % Note how a `%' character causes TeX to ignore
                             % the end of the input line, so these blank lines
                             % do not start a new paragraph.
                             %
with the right spacing around the periods requires
a special command.

\LaTeX\ interprets some common characters as
commands, so you must type special commands to
generate them.  These characters include the
following:
       \$ \& \% \# \{ and \}.

In printing, text is usually emphasized with an
       \emph{italic}
type style.

\begin{em}
   A long segment of text can also be emphasized
   in this way.  Text within such a segment can be
   given \emph{additional} emphasis.
\end{em}

It is sometimes necessary to prevent \LaTeX\ from
breaking a line where it might otherwise do so.
This may be at a space, as between the ``Mr.''\ and
``Jones'' in
       ``Mr.~Jones'',        % ~ produces an unbreakable interword space.
or within a word---especially when the word is a
symbol like
       \mbox{\emph{itemnum}}
that makes little sense when hyphenated across
lines.

Footnotes\footnote{This is an example of a footnote.}
pose no problem.

\LaTeX\ is good at typesetting mathematical formulas
like
       \( x-3y + z = 7 \)
or
       \( a_{1} > x^{2n} + y^{2n} > x' \)
or
       \( \ip{A}{B} = \sum_{i} a_{i} b_{i} \).
The spaces you type in a formula are
ignored.  Remember that a letter like
       $x$                   % $ ... $  and  \( ... \)  are equivalent
is a formula when it denotes a mathematical
symbol, and it should be typed as one.

\section{Displayed Text}

Text is displayed by indenting it from the left
margin.  Quotations are commonly displayed.  There
are short quotations
\begin{quote}
   This is a short quotation.  It consists of a
   single paragraph of text.  See how it is formatted.
\end{quote}
and longer ones.
\begin{quotation}
   This is a longer quotation.  It consists of two
   paragraphs of text, neither of which are
   particularly interesting.

   This is the second paragraph of the quotation.  It
   is just as dull as the first paragraph.
\end{quotation}
Another frequently-displayed structure is a list.
The following is an example of an \emph{itemized}
list.
\begin{itemize}
   \item This is the first item of an itemized list.
         Each item in the list is marked with a ``tick''.
         You don't have to worry about what kind of tick
         mark is used.

   \item This is the second item of the list.  It
         contains another list nested inside it.  The inner
         list is an \emph{enumerated} list.
         \begin{enumerate}
            \item This is the first item of an enumerated
                  list that is nested within the itemized list.

            \item This is the second item of the inner list.
                  \LaTeX\ allows you to nest lists deeper than
                  you really should.
         \end{enumerate}
         This is the rest of the second item of the outer
         list.  It is no more interesting than any other
         part of the item.
   \item This is the third item of the list.
\end{itemize}
You can even display poetry.
\begin{verse}
   There is an environment
    for verse \\             % The \\ command separates lines
   Whose features some poets % within a stanza.
   will curse.

                             % One or more blank lines separate stanzas.

   For instead of making\\
   Them do \emph{all} line breaking, \\
   It allows them to put too many words on a line when they'd rather be
   forced to be terse.
\end{verse}

Mathematical formulas may also be displayed.  A
displayed formula
is
one-line long; multiline
formulas require special formatting instructions.
   \[  \ip{\Gamma}{\psi'} = x'' + y^{2} + z_{i}^{n}\]
Don't start a paragraph with a displayed equation,
nor make one a paragraph by itself.

\end{document}               % End of document.

\section{Bench Huge Section 135}
% synthetic bench section 135
% This is a sample LaTeX input file.  (Version of 12 August 2004.)
%
% A '%' character causes TeX to ignore all remaining text on the line,
% and is used for comments like this one.

\documentclass{article}      % Specifies the document class

                             % The preamble begins here.
\title{An Example Document}  % Declares the document's title.
\author{Leslie Lamport}      % Declares the author's name.
\date{January 21, 1994}      % Deleting this command produces today's date.

\newcommand{\ip}[2]{(#1, #2)}
                             % Defines \ip{arg1}{arg2} to mean
                             % (arg1, arg2).

%\newcommand{\ip}[2]{\langle #1 | #2\rangle}
                             % This is an alternative definition of
                             % \ip that is commented out.

\begin{document}             % End of preamble and beginning of text.

\maketitle                   % Produces the title.

This is an example input file.  Comparing it with
the output it generates can show you how to
produce a simple document of your own.

\section{Ordinary Text}      % Produces section heading.  Lower-level
                             % sections are begun with similar
                             % \subsection and \subsubsection commands.

The ends  of words and sentences are marked
  by   spaces. It  doesn't matter how many
spaces    you type; one is as good as 100.  The
end of   a line counts as a space.

One   or more   blank lines denote the  end
of  a paragraph.

Since any number of consecutive spaces are treated
like a single one, the formatting of the input
file makes no difference to
      \LaTeX,                % The \LaTeX command generates the LaTeX logo.
but it makes a difference to you.  When you use
\LaTeX, making your input file as easy to read
as possible will be a great help as you write
your document and when you change it.  This sample
file shows how you can add comments to your own input
file.

Because printing is different from typewriting,
there are a number of things that you have to do
differently when preparing an input file than if
you were just typing the document directly.
Quotation marks like
       ``this''
have to be handled specially, as do quotes within
quotes:
       ``\,`this'            % \, separates the double and single quote.
        is what I just
        wrote, not  `that'\,''.

Dashes come in three sizes: an
       intra-word
dash, a medium dash for number ranges like
       1--2,
and a punctuation
       dash---like
this.

A sentence-ending space should be larger than the
space between words within a sentence.  You
sometimes have to type special commands in
conjunction with punctuation characters to get
this right, as in the following sentence.
       Gnats, gnus, etc.\ all  % `\ ' makes an inter-word space.
       begin with G\@.         % \@ marks end-of-sentence punctuation.
You should check the spaces after periods when
reading your output to make sure you haven't
forgotten any special cases.  Generating an
ellipsis
       \ldots\               % `\ ' is needed after `\ldots' because TeX
                             % ignores spaces after command names like \ldots
                             % made from \ + letters.
                             %
                             % Note how a `%' character causes TeX to ignore
                             % the end of the input line, so these blank lines
                             % do not start a new paragraph.
                             %
with the right spacing around the periods requires
a special command.

\LaTeX\ interprets some common characters as
commands, so you must type special commands to
generate them.  These characters include the
following:
       \$ \& \% \# \{ and \}.

In printing, text is usually emphasized with an
       \emph{italic}
type style.

\begin{em}
   A long segment of text can also be emphasized
   in this way.  Text within such a segment can be
   given \emph{additional} emphasis.
\end{em}

It is sometimes necessary to prevent \LaTeX\ from
breaking a line where it might otherwise do so.
This may be at a space, as between the ``Mr.''\ and
``Jones'' in
       ``Mr.~Jones'',        % ~ produces an unbreakable interword space.
or within a word---especially when the word is a
symbol like
       \mbox{\emph{itemnum}}
that makes little sense when hyphenated across
lines.

Footnotes\footnote{This is an example of a footnote.}
pose no problem.

\LaTeX\ is good at typesetting mathematical formulas
like
       \( x-3y + z = 7 \)
or
       \( a_{1} > x^{2n} + y^{2n} > x' \)
or
       \( \ip{A}{B} = \sum_{i} a_{i} b_{i} \).
The spaces you type in a formula are
ignored.  Remember that a letter like
       $x$                   % $ ... $  and  \( ... \)  are equivalent
is a formula when it denotes a mathematical
symbol, and it should be typed as one.

\section{Displayed Text}

Text is displayed by indenting it from the left
margin.  Quotations are commonly displayed.  There
are short quotations
\begin{quote}
   This is a short quotation.  It consists of a
   single paragraph of text.  See how it is formatted.
\end{quote}
and longer ones.
\begin{quotation}
   This is a longer quotation.  It consists of two
   paragraphs of text, neither of which are
   particularly interesting.

   This is the second paragraph of the quotation.  It
   is just as dull as the first paragraph.
\end{quotation}
Another frequently-displayed structure is a list.
The following is an example of an \emph{itemized}
list.
\begin{itemize}
   \item This is the first item of an itemized list.
         Each item in the list is marked with a ``tick''.
         You don't have to worry about what kind of tick
         mark is used.

   \item This is the second item of the list.  It
         contains another list nested inside it.  The inner
         list is an \emph{enumerated} list.
         \begin{enumerate}
            \item This is the first item of an enumerated
                  list that is nested within the itemized list.

            \item This is the second item of the inner list.
                  \LaTeX\ allows you to nest lists deeper than
                  you really should.
         \end{enumerate}
         This is the rest of the second item of the outer
         list.  It is no more interesting than any other
         part of the item.
   \item This is the third item of the list.
\end{itemize}
You can even display poetry.
\begin{verse}
   There is an environment
    for verse \\             % The \\ command separates lines
   Whose features some poets % within a stanza.
   will curse.

                             % One or more blank lines separate stanzas.

   For instead of making\\
   Them do \emph{all} line breaking, \\
   It allows them to put too many words on a line when they'd rather be
   forced to be terse.
\end{verse}

Mathematical formulas may also be displayed.  A
displayed formula
is
one-line long; multiline
formulas require special formatting instructions.
   \[  \ip{\Gamma}{\psi'} = x'' + y^{2} + z_{i}^{n}\]
Don't start a paragraph with a displayed equation,
nor make one a paragraph by itself.

\end{document}               % End of document.

\section{Bench Huge Section 136}
% synthetic bench section 136
% This is a sample LaTeX input file.  (Version of 12 August 2004.)
%
% A '%' character causes TeX to ignore all remaining text on the line,
% and is used for comments like this one.

\documentclass{article}      % Specifies the document class

                             % The preamble begins here.
\title{An Example Document}  % Declares the document's title.
\author{Leslie Lamport}      % Declares the author's name.
\date{January 21, 1994}      % Deleting this command produces today's date.

\newcommand{\ip}[2]{(#1, #2)}
                             % Defines \ip{arg1}{arg2} to mean
                             % (arg1, arg2).

%\newcommand{\ip}[2]{\langle #1 | #2\rangle}
                             % This is an alternative definition of
                             % \ip that is commented out.

\begin{document}             % End of preamble and beginning of text.

\maketitle                   % Produces the title.

This is an example input file.  Comparing it with
the output it generates can show you how to
produce a simple document of your own.

\section{Ordinary Text}      % Produces section heading.  Lower-level
                             % sections are begun with similar
                             % \subsection and \subsubsection commands.

The ends  of words and sentences are marked
  by   spaces. It  doesn't matter how many
spaces    you type; one is as good as 100.  The
end of   a line counts as a space.

One   or more   blank lines denote the  end
of  a paragraph.

Since any number of consecutive spaces are treated
like a single one, the formatting of the input
file makes no difference to
      \LaTeX,                % The \LaTeX command generates the LaTeX logo.
but it makes a difference to you.  When you use
\LaTeX, making your input file as easy to read
as possible will be a great help as you write
your document and when you change it.  This sample
file shows how you can add comments to your own input
file.

Because printing is different from typewriting,
there are a number of things that you have to do
differently when preparing an input file than if
you were just typing the document directly.
Quotation marks like
       ``this''
have to be handled specially, as do quotes within
quotes:
       ``\,`this'            % \, separates the double and single quote.
        is what I just
        wrote, not  `that'\,''.

Dashes come in three sizes: an
       intra-word
dash, a medium dash for number ranges like
       1--2,
and a punctuation
       dash---like
this.

A sentence-ending space should be larger than the
space between words within a sentence.  You
sometimes have to type special commands in
conjunction with punctuation characters to get
this right, as in the following sentence.
       Gnats, gnus, etc.\ all  % `\ ' makes an inter-word space.
       begin with G\@.         % \@ marks end-of-sentence punctuation.
You should check the spaces after periods when
reading your output to make sure you haven't
forgotten any special cases.  Generating an
ellipsis
       \ldots\               % `\ ' is needed after `\ldots' because TeX
                             % ignores spaces after command names like \ldots
                             % made from \ + letters.
                             %
                             % Note how a `%' character causes TeX to ignore
                             % the end of the input line, so these blank lines
                             % do not start a new paragraph.
                             %
with the right spacing around the periods requires
a special command.

\LaTeX\ interprets some common characters as
commands, so you must type special commands to
generate them.  These characters include the
following:
       \$ \& \% \# \{ and \}.

In printing, text is usually emphasized with an
       \emph{italic}
type style.

\begin{em}
   A long segment of text can also be emphasized
   in this way.  Text within such a segment can be
   given \emph{additional} emphasis.
\end{em}

It is sometimes necessary to prevent \LaTeX\ from
breaking a line where it might otherwise do so.
This may be at a space, as between the ``Mr.''\ and
``Jones'' in
       ``Mr.~Jones'',        % ~ produces an unbreakable interword space.
or within a word---especially when the word is a
symbol like
       \mbox{\emph{itemnum}}
that makes little sense when hyphenated across
lines.

Footnotes\footnote{This is an example of a footnote.}
pose no problem.

\LaTeX\ is good at typesetting mathematical formulas
like
       \( x-3y + z = 7 \)
or
       \( a_{1} > x^{2n} + y^{2n} > x' \)
or
       \( \ip{A}{B} = \sum_{i} a_{i} b_{i} \).
The spaces you type in a formula are
ignored.  Remember that a letter like
       $x$                   % $ ... $  and  \( ... \)  are equivalent
is a formula when it denotes a mathematical
symbol, and it should be typed as one.

\section{Displayed Text}

Text is displayed by indenting it from the left
margin.  Quotations are commonly displayed.  There
are short quotations
\begin{quote}
   This is a short quotation.  It consists of a
   single paragraph of text.  See how it is formatted.
\end{quote}
and longer ones.
\begin{quotation}
   This is a longer quotation.  It consists of two
   paragraphs of text, neither of which are
   particularly interesting.

   This is the second paragraph of the quotation.  It
   is just as dull as the first paragraph.
\end{quotation}
Another frequently-displayed structure is a list.
The following is an example of an \emph{itemized}
list.
\begin{itemize}
   \item This is the first item of an itemized list.
         Each item in the list is marked with a ``tick''.
         You don't have to worry about what kind of tick
         mark is used.

   \item This is the second item of the list.  It
         contains another list nested inside it.  The inner
         list is an \emph{enumerated} list.
         \begin{enumerate}
            \item This is the first item of an enumerated
                  list that is nested within the itemized list.

            \item This is the second item of the inner list.
                  \LaTeX\ allows you to nest lists deeper than
                  you really should.
         \end{enumerate}
         This is the rest of the second item of the outer
         list.  It is no more interesting than any other
         part of the item.
   \item This is the third item of the list.
\end{itemize}
You can even display poetry.
\begin{verse}
   There is an environment
    for verse \\             % The \\ command separates lines
   Whose features some poets % within a stanza.
   will curse.

                             % One or more blank lines separate stanzas.

   For instead of making\\
   Them do \emph{all} line breaking, \\
   It allows them to put too many words on a line when they'd rather be
   forced to be terse.
\end{verse}

Mathematical formulas may also be displayed.  A
displayed formula
is
one-line long; multiline
formulas require special formatting instructions.
   \[  \ip{\Gamma}{\psi'} = x'' + y^{2} + z_{i}^{n}\]
Don't start a paragraph with a displayed equation,
nor make one a paragraph by itself.

\end{document}               % End of document.

\section{Bench Huge Section 137}
% synthetic bench section 137
% This is a sample LaTeX input file.  (Version of 12 August 2004.)
%
% A '%' character causes TeX to ignore all remaining text on the line,
% and is used for comments like this one.

\documentclass{article}      % Specifies the document class

                             % The preamble begins here.
\title{An Example Document}  % Declares the document's title.
\author{Leslie Lamport}      % Declares the author's name.
\date{January 21, 1994}      % Deleting this command produces today's date.

\newcommand{\ip}[2]{(#1, #2)}
                             % Defines \ip{arg1}{arg2} to mean
                             % (arg1, arg2).

%\newcommand{\ip}[2]{\langle #1 | #2\rangle}
                             % This is an alternative definition of
                             % \ip that is commented out.

\begin{document}             % End of preamble and beginning of text.

\maketitle                   % Produces the title.

This is an example input file.  Comparing it with
the output it generates can show you how to
produce a simple document of your own.

\section{Ordinary Text}      % Produces section heading.  Lower-level
                             % sections are begun with similar
                             % \subsection and \subsubsection commands.

The ends  of words and sentences are marked
  by   spaces. It  doesn't matter how many
spaces    you type; one is as good as 100.  The
end of   a line counts as a space.

One   or more   blank lines denote the  end
of  a paragraph.

Since any number of consecutive spaces are treated
like a single one, the formatting of the input
file makes no difference to
      \LaTeX,                % The \LaTeX command generates the LaTeX logo.
but it makes a difference to you.  When you use
\LaTeX, making your input file as easy to read
as possible will be a great help as you write
your document and when you change it.  This sample
file shows how you can add comments to your own input
file.

Because printing is different from typewriting,
there are a number of things that you have to do
differently when preparing an input file than if
you were just typing the document directly.
Quotation marks like
       ``this''
have to be handled specially, as do quotes within
quotes:
       ``\,`this'            % \, separates the double and single quote.
        is what I just
        wrote, not  `that'\,''.

Dashes come in three sizes: an
       intra-word
dash, a medium dash for number ranges like
       1--2,
and a punctuation
       dash---like
this.

A sentence-ending space should be larger than the
space between words within a sentence.  You
sometimes have to type special commands in
conjunction with punctuation characters to get
this right, as in the following sentence.
       Gnats, gnus, etc.\ all  % `\ ' makes an inter-word space.
       begin with G\@.         % \@ marks end-of-sentence punctuation.
You should check the spaces after periods when
reading your output to make sure you haven't
forgotten any special cases.  Generating an
ellipsis
       \ldots\               % `\ ' is needed after `\ldots' because TeX
                             % ignores spaces after command names like \ldots
                             % made from \ + letters.
                             %
                             % Note how a `%' character causes TeX to ignore
                             % the end of the input line, so these blank lines
                             % do not start a new paragraph.
                             %
with the right spacing around the periods requires
a special command.

\LaTeX\ interprets some common characters as
commands, so you must type special commands to
generate them.  These characters include the
following:
       \$ \& \% \# \{ and \}.

In printing, text is usually emphasized with an
       \emph{italic}
type style.

\begin{em}
   A long segment of text can also be emphasized
   in this way.  Text within such a segment can be
   given \emph{additional} emphasis.
\end{em}

It is sometimes necessary to prevent \LaTeX\ from
breaking a line where it might otherwise do so.
This may be at a space, as between the ``Mr.''\ and
``Jones'' in
       ``Mr.~Jones'',        % ~ produces an unbreakable interword space.
or within a word---especially when the word is a
symbol like
       \mbox{\emph{itemnum}}
that makes little sense when hyphenated across
lines.

Footnotes\footnote{This is an example of a footnote.}
pose no problem.

\LaTeX\ is good at typesetting mathematical formulas
like
       \( x-3y + z = 7 \)
or
       \( a_{1} > x^{2n} + y^{2n} > x' \)
or
       \( \ip{A}{B} = \sum_{i} a_{i} b_{i} \).
The spaces you type in a formula are
ignored.  Remember that a letter like
       $x$                   % $ ... $  and  \( ... \)  are equivalent
is a formula when it denotes a mathematical
symbol, and it should be typed as one.

\section{Displayed Text}

Text is displayed by indenting it from the left
margin.  Quotations are commonly displayed.  There
are short quotations
\begin{quote}
   This is a short quotation.  It consists of a
   single paragraph of text.  See how it is formatted.
\end{quote}
and longer ones.
\begin{quotation}
   This is a longer quotation.  It consists of two
   paragraphs of text, neither of which are
   particularly interesting.

   This is the second paragraph of the quotation.  It
   is just as dull as the first paragraph.
\end{quotation}
Another frequently-displayed structure is a list.
The following is an example of an \emph{itemized}
list.
\begin{itemize}
   \item This is the first item of an itemized list.
         Each item in the list is marked with a ``tick''.
         You don't have to worry about what kind of tick
         mark is used.

   \item This is the second item of the list.  It
         contains another list nested inside it.  The inner
         list is an \emph{enumerated} list.
         \begin{enumerate}
            \item This is the first item of an enumerated
                  list that is nested within the itemized list.

            \item This is the second item of the inner list.
                  \LaTeX\ allows you to nest lists deeper than
                  you really should.
         \end{enumerate}
         This is the rest of the second item of the outer
         list.  It is no more interesting than any other
         part of the item.
   \item This is the third item of the list.
\end{itemize}
You can even display poetry.
\begin{verse}
   There is an environment
    for verse \\             % The \\ command separates lines
   Whose features some poets % within a stanza.
   will curse.

                             % One or more blank lines separate stanzas.

   For instead of making\\
   Them do \emph{all} line breaking, \\
   It allows them to put too many words on a line when they'd rather be
   forced to be terse.
\end{verse}

Mathematical formulas may also be displayed.  A
displayed formula
is
one-line long; multiline
formulas require special formatting instructions.
   \[  \ip{\Gamma}{\psi'} = x'' + y^{2} + z_{i}^{n}\]
Don't start a paragraph with a displayed equation,
nor make one a paragraph by itself.

\end{document}               % End of document.

\section{Bench Huge Section 138}
% synthetic bench section 138
% This is a sample LaTeX input file.  (Version of 12 August 2004.)
%
% A '%' character causes TeX to ignore all remaining text on the line,
% and is used for comments like this one.

\documentclass{article}      % Specifies the document class

                             % The preamble begins here.
\title{An Example Document}  % Declares the document's title.
\author{Leslie Lamport}      % Declares the author's name.
\date{January 21, 1994}      % Deleting this command produces today's date.

\newcommand{\ip}[2]{(#1, #2)}
                             % Defines \ip{arg1}{arg2} to mean
                             % (arg1, arg2).

%\newcommand{\ip}[2]{\langle #1 | #2\rangle}
                             % This is an alternative definition of
                             % \ip that is commented out.

\begin{document}             % End of preamble and beginning of text.

\maketitle                   % Produces the title.

This is an example input file.  Comparing it with
the output it generates can show you how to
produce a simple document of your own.

\section{Ordinary Text}      % Produces section heading.  Lower-level
                             % sections are begun with similar
                             % \subsection and \subsubsection commands.

The ends  of words and sentences are marked
  by   spaces. It  doesn't matter how many
spaces    you type; one is as good as 100.  The
end of   a line counts as a space.

One   or more   blank lines denote the  end
of  a paragraph.

Since any number of consecutive spaces are treated
like a single one, the formatting of the input
file makes no difference to
      \LaTeX,                % The \LaTeX command generates the LaTeX logo.
but it makes a difference to you.  When you use
\LaTeX, making your input file as easy to read
as possible will be a great help as you write
your document and when you change it.  This sample
file shows how you can add comments to your own input
file.

Because printing is different from typewriting,
there are a number of things that you have to do
differently when preparing an input file than if
you were just typing the document directly.
Quotation marks like
       ``this''
have to be handled specially, as do quotes within
quotes:
       ``\,`this'            % \, separates the double and single quote.
        is what I just
        wrote, not  `that'\,''.

Dashes come in three sizes: an
       intra-word
dash, a medium dash for number ranges like
       1--2,
and a punctuation
       dash---like
this.

A sentence-ending space should be larger than the
space between words within a sentence.  You
sometimes have to type special commands in
conjunction with punctuation characters to get
this right, as in the following sentence.
       Gnats, gnus, etc.\ all  % `\ ' makes an inter-word space.
       begin with G\@.         % \@ marks end-of-sentence punctuation.
You should check the spaces after periods when
reading your output to make sure you haven't
forgotten any special cases.  Generating an
ellipsis
       \ldots\               % `\ ' is needed after `\ldots' because TeX
                             % ignores spaces after command names like \ldots
                             % made from \ + letters.
                             %
                             % Note how a `%' character causes TeX to ignore
                             % the end of the input line, so these blank lines
                             % do not start a new paragraph.
                             %
with the right spacing around the periods requires
a special command.

\LaTeX\ interprets some common characters as
commands, so you must type special commands to
generate them.  These characters include the
following:
       \$ \& \% \# \{ and \}.

In printing, text is usually emphasized with an
       \emph{italic}
type style.

\begin{em}
   A long segment of text can also be emphasized
   in this way.  Text within such a segment can be
   given \emph{additional} emphasis.
\end{em}

It is sometimes necessary to prevent \LaTeX\ from
breaking a line where it might otherwise do so.
This may be at a space, as between the ``Mr.''\ and
``Jones'' in
       ``Mr.~Jones'',        % ~ produces an unbreakable interword space.
or within a word---especially when the word is a
symbol like
       \mbox{\emph{itemnum}}
that makes little sense when hyphenated across
lines.

Footnotes\footnote{This is an example of a footnote.}
pose no problem.

\LaTeX\ is good at typesetting mathematical formulas
like
       \( x-3y + z = 7 \)
or
       \( a_{1} > x^{2n} + y^{2n} > x' \)
or
       \( \ip{A}{B} = \sum_{i} a_{i} b_{i} \).
The spaces you type in a formula are
ignored.  Remember that a letter like
       $x$                   % $ ... $  and  \( ... \)  are equivalent
is a formula when it denotes a mathematical
symbol, and it should be typed as one.

\section{Displayed Text}

Text is displayed by indenting it from the left
margin.  Quotations are commonly displayed.  There
are short quotations
\begin{quote}
   This is a short quotation.  It consists of a
   single paragraph of text.  See how it is formatted.
\end{quote}
and longer ones.
\begin{quotation}
   This is a longer quotation.  It consists of two
   paragraphs of text, neither of which are
   particularly interesting.

   This is the second paragraph of the quotation.  It
   is just as dull as the first paragraph.
\end{quotation}
Another frequently-displayed structure is a list.
The following is an example of an \emph{itemized}
list.
\begin{itemize}
   \item This is the first item of an itemized list.
         Each item in the list is marked with a ``tick''.
         You don't have to worry about what kind of tick
         mark is used.

   \item This is the second item of the list.  It
         contains another list nested inside it.  The inner
         list is an \emph{enumerated} list.
         \begin{enumerate}
            \item This is the first item of an enumerated
                  list that is nested within the itemized list.

            \item This is the second item of the inner list.
                  \LaTeX\ allows you to nest lists deeper than
                  you really should.
         \end{enumerate}
         This is the rest of the second item of the outer
         list.  It is no more interesting than any other
         part of the item.
   \item This is the third item of the list.
\end{itemize}
You can even display poetry.
\begin{verse}
   There is an environment
    for verse \\             % The \\ command separates lines
   Whose features some poets % within a stanza.
   will curse.

                             % One or more blank lines separate stanzas.

   For instead of making\\
   Them do \emph{all} line breaking, \\
   It allows them to put too many words on a line when they'd rather be
   forced to be terse.
\end{verse}

Mathematical formulas may also be displayed.  A
displayed formula
is
one-line long; multiline
formulas require special formatting instructions.
   \[  \ip{\Gamma}{\psi'} = x'' + y^{2} + z_{i}^{n}\]
Don't start a paragraph with a displayed equation,
nor make one a paragraph by itself.

\end{document}               % End of document.

\section{Bench Huge Section 139}
% synthetic bench section 139
% This is a sample LaTeX input file.  (Version of 12 August 2004.)
%
% A '%' character causes TeX to ignore all remaining text on the line,
% and is used for comments like this one.

\documentclass{article}      % Specifies the document class

                             % The preamble begins here.
\title{An Example Document}  % Declares the document's title.
\author{Leslie Lamport}      % Declares the author's name.
\date{January 21, 1994}      % Deleting this command produces today's date.

\newcommand{\ip}[2]{(#1, #2)}
                             % Defines \ip{arg1}{arg2} to mean
                             % (arg1, arg2).

%\newcommand{\ip}[2]{\langle #1 | #2\rangle}
                             % This is an alternative definition of
                             % \ip that is commented out.

\begin{document}             % End of preamble and beginning of text.

\maketitle                   % Produces the title.

This is an example input file.  Comparing it with
the output it generates can show you how to
produce a simple document of your own.

\section{Ordinary Text}      % Produces section heading.  Lower-level
                             % sections are begun with similar
                             % \subsection and \subsubsection commands.

The ends  of words and sentences are marked
  by   spaces. It  doesn't matter how many
spaces    you type; one is as good as 100.  The
end of   a line counts as a space.

One   or more   blank lines denote the  end
of  a paragraph.

Since any number of consecutive spaces are treated
like a single one, the formatting of the input
file makes no difference to
      \LaTeX,                % The \LaTeX command generates the LaTeX logo.
but it makes a difference to you.  When you use
\LaTeX, making your input file as easy to read
as possible will be a great help as you write
your document and when you change it.  This sample
file shows how you can add comments to your own input
file.

Because printing is different from typewriting,
there are a number of things that you have to do
differently when preparing an input file than if
you were just typing the document directly.
Quotation marks like
       ``this''
have to be handled specially, as do quotes within
quotes:
       ``\,`this'            % \, separates the double and single quote.
        is what I just
        wrote, not  `that'\,''.

Dashes come in three sizes: an
       intra-word
dash, a medium dash for number ranges like
       1--2,
and a punctuation
       dash---like
this.

A sentence-ending space should be larger than the
space between words within a sentence.  You
sometimes have to type special commands in
conjunction with punctuation characters to get
this right, as in the following sentence.
       Gnats, gnus, etc.\ all  % `\ ' makes an inter-word space.
       begin with G\@.         % \@ marks end-of-sentence punctuation.
You should check the spaces after periods when
reading your output to make sure you haven't
forgotten any special cases.  Generating an
ellipsis
       \ldots\               % `\ ' is needed after `\ldots' because TeX
                             % ignores spaces after command names like \ldots
                             % made from \ + letters.
                             %
                             % Note how a `%' character causes TeX to ignore
                             % the end of the input line, so these blank lines
                             % do not start a new paragraph.
                             %
with the right spacing around the periods requires
a special command.

\LaTeX\ interprets some common characters as
commands, so you must type special commands to
generate them.  These characters include the
following:
       \$ \& \% \# \{ and \}.

In printing, text is usually emphasized with an
       \emph{italic}
type style.

\begin{em}
   A long segment of text can also be emphasized
   in this way.  Text within such a segment can be
   given \emph{additional} emphasis.
\end{em}

It is sometimes necessary to prevent \LaTeX\ from
breaking a line where it might otherwise do so.
This may be at a space, as between the ``Mr.''\ and
``Jones'' in
       ``Mr.~Jones'',        % ~ produces an unbreakable interword space.
or within a word---especially when the word is a
symbol like
       \mbox{\emph{itemnum}}
that makes little sense when hyphenated across
lines.

Footnotes\footnote{This is an example of a footnote.}
pose no problem.

\LaTeX\ is good at typesetting mathematical formulas
like
       \( x-3y + z = 7 \)
or
       \( a_{1} > x^{2n} + y^{2n} > x' \)
or
       \( \ip{A}{B} = \sum_{i} a_{i} b_{i} \).
The spaces you type in a formula are
ignored.  Remember that a letter like
       $x$                   % $ ... $  and  \( ... \)  are equivalent
is a formula when it denotes a mathematical
symbol, and it should be typed as one.

\section{Displayed Text}

Text is displayed by indenting it from the left
margin.  Quotations are commonly displayed.  There
are short quotations
\begin{quote}
   This is a short quotation.  It consists of a
   single paragraph of text.  See how it is formatted.
\end{quote}
and longer ones.
\begin{quotation}
   This is a longer quotation.  It consists of two
   paragraphs of text, neither of which are
   particularly interesting.

   This is the second paragraph of the quotation.  It
   is just as dull as the first paragraph.
\end{quotation}
Another frequently-displayed structure is a list.
The following is an example of an \emph{itemized}
list.
\begin{itemize}
   \item This is the first item of an itemized list.
         Each item in the list is marked with a ``tick''.
         You don't have to worry about what kind of tick
         mark is used.

   \item This is the second item of the list.  It
         contains another list nested inside it.  The inner
         list is an \emph{enumerated} list.
         \begin{enumerate}
            \item This is the first item of an enumerated
                  list that is nested within the itemized list.

            \item This is the second item of the inner list.
                  \LaTeX\ allows you to nest lists deeper than
                  you really should.
         \end{enumerate}
         This is the rest of the second item of the outer
         list.  It is no more interesting than any other
         part of the item.
   \item This is the third item of the list.
\end{itemize}
You can even display poetry.
\begin{verse}
   There is an environment
    for verse \\             % The \\ command separates lines
   Whose features some poets % within a stanza.
   will curse.

                             % One or more blank lines separate stanzas.

   For instead of making\\
   Them do \emph{all} line breaking, \\
   It allows them to put too many words on a line when they'd rather be
   forced to be terse.
\end{verse}

Mathematical formulas may also be displayed.  A
displayed formula
is
one-line long; multiline
formulas require special formatting instructions.
   \[  \ip{\Gamma}{\psi'} = x'' + y^{2} + z_{i}^{n}\]
Don't start a paragraph with a displayed equation,
nor make one a paragraph by itself.

\end{document}               % End of document.

\section{Bench Huge Section 140}
% synthetic bench section 140
% This is a sample LaTeX input file.  (Version of 12 August 2004.)
%
% A '%' character causes TeX to ignore all remaining text on the line,
% and is used for comments like this one.

\documentclass{article}      % Specifies the document class

                             % The preamble begins here.
\title{An Example Document}  % Declares the document's title.
\author{Leslie Lamport}      % Declares the author's name.
\date{January 21, 1994}      % Deleting this command produces today's date.

\newcommand{\ip}[2]{(#1, #2)}
                             % Defines \ip{arg1}{arg2} to mean
                             % (arg1, arg2).

%\newcommand{\ip}[2]{\langle #1 | #2\rangle}
                             % This is an alternative definition of
                             % \ip that is commented out.

\begin{document}             % End of preamble and beginning of text.

\maketitle                   % Produces the title.

This is an example input file.  Comparing it with
the output it generates can show you how to
produce a simple document of your own.

\section{Ordinary Text}      % Produces section heading.  Lower-level
                             % sections are begun with similar
                             % \subsection and \subsubsection commands.

The ends  of words and sentences are marked
  by   spaces. It  doesn't matter how many
spaces    you type; one is as good as 100.  The
end of   a line counts as a space.

One   or more   blank lines denote the  end
of  a paragraph.

Since any number of consecutive spaces are treated
like a single one, the formatting of the input
file makes no difference to
      \LaTeX,                % The \LaTeX command generates the LaTeX logo.
but it makes a difference to you.  When you use
\LaTeX, making your input file as easy to read
as possible will be a great help as you write
your document and when you change it.  This sample
file shows how you can add comments to your own input
file.

Because printing is different from typewriting,
there are a number of things that you have to do
differently when preparing an input file than if
you were just typing the document directly.
Quotation marks like
       ``this''
have to be handled specially, as do quotes within
quotes:
       ``\,`this'            % \, separates the double and single quote.
        is what I just
        wrote, not  `that'\,''.

Dashes come in three sizes: an
       intra-word
dash, a medium dash for number ranges like
       1--2,
and a punctuation
       dash---like
this.

A sentence-ending space should be larger than the
space between words within a sentence.  You
sometimes have to type special commands in
conjunction with punctuation characters to get
this right, as in the following sentence.
       Gnats, gnus, etc.\ all  % `\ ' makes an inter-word space.
       begin with G\@.         % \@ marks end-of-sentence punctuation.
You should check the spaces after periods when
reading your output to make sure you haven't
forgotten any special cases.  Generating an
ellipsis
       \ldots\               % `\ ' is needed after `\ldots' because TeX
                             % ignores spaces after command names like \ldots
                             % made from \ + letters.
                             %
                             % Note how a `%' character causes TeX to ignore
                             % the end of the input line, so these blank lines
                             % do not start a new paragraph.
                             %
with the right spacing around the periods requires
a special command.

\LaTeX\ interprets some common characters as
commands, so you must type special commands to
generate them.  These characters include the
following:
       \$ \& \% \# \{ and \}.

In printing, text is usually emphasized with an
       \emph{italic}
type style.

\begin{em}
   A long segment of text can also be emphasized
   in this way.  Text within such a segment can be
   given \emph{additional} emphasis.
\end{em}

It is sometimes necessary to prevent \LaTeX\ from
breaking a line where it might otherwise do so.
This may be at a space, as between the ``Mr.''\ and
``Jones'' in
       ``Mr.~Jones'',        % ~ produces an unbreakable interword space.
or within a word---especially when the word is a
symbol like
       \mbox{\emph{itemnum}}
that makes little sense when hyphenated across
lines.

Footnotes\footnote{This is an example of a footnote.}
pose no problem.

\LaTeX\ is good at typesetting mathematical formulas
like
       \( x-3y + z = 7 \)
or
       \( a_{1} > x^{2n} + y^{2n} > x' \)
or
       \( \ip{A}{B} = \sum_{i} a_{i} b_{i} \).
The spaces you type in a formula are
ignored.  Remember that a letter like
       $x$                   % $ ... $  and  \( ... \)  are equivalent
is a formula when it denotes a mathematical
symbol, and it should be typed as one.

\section{Displayed Text}

Text is displayed by indenting it from the left
margin.  Quotations are commonly displayed.  There
are short quotations
\begin{quote}
   This is a short quotation.  It consists of a
   single paragraph of text.  See how it is formatted.
\end{quote}
and longer ones.
\begin{quotation}
   This is a longer quotation.  It consists of two
   paragraphs of text, neither of which are
   particularly interesting.

   This is the second paragraph of the quotation.  It
   is just as dull as the first paragraph.
\end{quotation}
Another frequently-displayed structure is a list.
The following is an example of an \emph{itemized}
list.
\begin{itemize}
   \item This is the first item of an itemized list.
         Each item in the list is marked with a ``tick''.
         You don't have to worry about what kind of tick
         mark is used.

   \item This is the second item of the list.  It
         contains another list nested inside it.  The inner
         list is an \emph{enumerated} list.
         \begin{enumerate}
            \item This is the first item of an enumerated
                  list that is nested within the itemized list.

            \item This is the second item of the inner list.
                  \LaTeX\ allows you to nest lists deeper than
                  you really should.
         \end{enumerate}
         This is the rest of the second item of the outer
         list.  It is no more interesting than any other
         part of the item.
   \item This is the third item of the list.
\end{itemize}
You can even display poetry.
\begin{verse}
   There is an environment
    for verse \\             % The \\ command separates lines
   Whose features some poets % within a stanza.
   will curse.

                             % One or more blank lines separate stanzas.

   For instead of making\\
   Them do \emph{all} line breaking, \\
   It allows them to put too many words on a line when they'd rather be
   forced to be terse.
\end{verse}

Mathematical formulas may also be displayed.  A
displayed formula
is
one-line long; multiline
formulas require special formatting instructions.
   \[  \ip{\Gamma}{\psi'} = x'' + y^{2} + z_{i}^{n}\]
Don't start a paragraph with a displayed equation,
nor make one a paragraph by itself.

\end{document}               % End of document.

\section{Bench Huge Section 141}
% synthetic bench section 141
% This is a sample LaTeX input file.  (Version of 12 August 2004.)
%
% A '%' character causes TeX to ignore all remaining text on the line,
% and is used for comments like this one.

\documentclass{article}      % Specifies the document class

                             % The preamble begins here.
\title{An Example Document}  % Declares the document's title.
\author{Leslie Lamport}      % Declares the author's name.
\date{January 21, 1994}      % Deleting this command produces today's date.

\newcommand{\ip}[2]{(#1, #2)}
                             % Defines \ip{arg1}{arg2} to mean
                             % (arg1, arg2).

%\newcommand{\ip}[2]{\langle #1 | #2\rangle}
                             % This is an alternative definition of
                             % \ip that is commented out.

\begin{document}             % End of preamble and beginning of text.

\maketitle                   % Produces the title.

This is an example input file.  Comparing it with
the output it generates can show you how to
produce a simple document of your own.

\section{Ordinary Text}      % Produces section heading.  Lower-level
                             % sections are begun with similar
                             % \subsection and \subsubsection commands.

The ends  of words and sentences are marked
  by   spaces. It  doesn't matter how many
spaces    you type; one is as good as 100.  The
end of   a line counts as a space.

One   or more   blank lines denote the  end
of  a paragraph.

Since any number of consecutive spaces are treated
like a single one, the formatting of the input
file makes no difference to
      \LaTeX,                % The \LaTeX command generates the LaTeX logo.
but it makes a difference to you.  When you use
\LaTeX, making your input file as easy to read
as possible will be a great help as you write
your document and when you change it.  This sample
file shows how you can add comments to your own input
file.

Because printing is different from typewriting,
there are a number of things that you have to do
differently when preparing an input file than if
you were just typing the document directly.
Quotation marks like
       ``this''
have to be handled specially, as do quotes within
quotes:
       ``\,`this'            % \, separates the double and single quote.
        is what I just
        wrote, not  `that'\,''.

Dashes come in three sizes: an
       intra-word
dash, a medium dash for number ranges like
       1--2,
and a punctuation
       dash---like
this.

A sentence-ending space should be larger than the
space between words within a sentence.  You
sometimes have to type special commands in
conjunction with punctuation characters to get
this right, as in the following sentence.
       Gnats, gnus, etc.\ all  % `\ ' makes an inter-word space.
       begin with G\@.         % \@ marks end-of-sentence punctuation.
You should check the spaces after periods when
reading your output to make sure you haven't
forgotten any special cases.  Generating an
ellipsis
       \ldots\               % `\ ' is needed after `\ldots' because TeX
                             % ignores spaces after command names like \ldots
                             % made from \ + letters.
                             %
                             % Note how a `%' character causes TeX to ignore
                             % the end of the input line, so these blank lines
                             % do not start a new paragraph.
                             %
with the right spacing around the periods requires
a special command.

\LaTeX\ interprets some common characters as
commands, so you must type special commands to
generate them.  These characters include the
following:
       \$ \& \% \# \{ and \}.

In printing, text is usually emphasized with an
       \emph{italic}
type style.

\begin{em}
   A long segment of text can also be emphasized
   in this way.  Text within such a segment can be
   given \emph{additional} emphasis.
\end{em}

It is sometimes necessary to prevent \LaTeX\ from
breaking a line where it might otherwise do so.
This may be at a space, as between the ``Mr.''\ and
``Jones'' in
       ``Mr.~Jones'',        % ~ produces an unbreakable interword space.
or within a word---especially when the word is a
symbol like
       \mbox{\emph{itemnum}}
that makes little sense when hyphenated across
lines.

Footnotes\footnote{This is an example of a footnote.}
pose no problem.

\LaTeX\ is good at typesetting mathematical formulas
like
       \( x-3y + z = 7 \)
or
       \( a_{1} > x^{2n} + y^{2n} > x' \)
or
       \( \ip{A}{B} = \sum_{i} a_{i} b_{i} \).
The spaces you type in a formula are
ignored.  Remember that a letter like
       $x$                   % $ ... $  and  \( ... \)  are equivalent
is a formula when it denotes a mathematical
symbol, and it should be typed as one.

\section{Displayed Text}

Text is displayed by indenting it from the left
margin.  Quotations are commonly displayed.  There
are short quotations
\begin{quote}
   This is a short quotation.  It consists of a
   single paragraph of text.  See how it is formatted.
\end{quote}
and longer ones.
\begin{quotation}
   This is a longer quotation.  It consists of two
   paragraphs of text, neither of which are
   particularly interesting.

   This is the second paragraph of the quotation.  It
   is just as dull as the first paragraph.
\end{quotation}
Another frequently-displayed structure is a list.
The following is an example of an \emph{itemized}
list.
\begin{itemize}
   \item This is the first item of an itemized list.
         Each item in the list is marked with a ``tick''.
         You don't have to worry about what kind of tick
         mark is used.

   \item This is the second item of the list.  It
         contains another list nested inside it.  The inner
         list is an \emph{enumerated} list.
         \begin{enumerate}
            \item This is the first item of an enumerated
                  list that is nested within the itemized list.

            \item This is the second item of the inner list.
                  \LaTeX\ allows you to nest lists deeper than
                  you really should.
         \end{enumerate}
         This is the rest of the second item of the outer
         list.  It is no more interesting than any other
         part of the item.
   \item This is the third item of the list.
\end{itemize}
You can even display poetry.
\begin{verse}
   There is an environment
    for verse \\             % The \\ command separates lines
   Whose features some poets % within a stanza.
   will curse.

                             % One or more blank lines separate stanzas.

   For instead of making\\
   Them do \emph{all} line breaking, \\
   It allows them to put too many words on a line when they'd rather be
   forced to be terse.
\end{verse}

Mathematical formulas may also be displayed.  A
displayed formula
is
one-line long; multiline
formulas require special formatting instructions.
   \[  \ip{\Gamma}{\psi'} = x'' + y^{2} + z_{i}^{n}\]
Don't start a paragraph with a displayed equation,
nor make one a paragraph by itself.

\end{document}               % End of document.

\section{Bench Huge Section 142}
% synthetic bench section 142
% This is a sample LaTeX input file.  (Version of 12 August 2004.)
%
% A '%' character causes TeX to ignore all remaining text on the line,
% and is used for comments like this one.

\documentclass{article}      % Specifies the document class

                             % The preamble begins here.
\title{An Example Document}  % Declares the document's title.
\author{Leslie Lamport}      % Declares the author's name.
\date{January 21, 1994}      % Deleting this command produces today's date.

\newcommand{\ip}[2]{(#1, #2)}
                             % Defines \ip{arg1}{arg2} to mean
                             % (arg1, arg2).

%\newcommand{\ip}[2]{\langle #1 | #2\rangle}
                             % This is an alternative definition of
                             % \ip that is commented out.

\begin{document}             % End of preamble and beginning of text.

\maketitle                   % Produces the title.

This is an example input file.  Comparing it with
the output it generates can show you how to
produce a simple document of your own.

\section{Ordinary Text}      % Produces section heading.  Lower-level
                             % sections are begun with similar
                             % \subsection and \subsubsection commands.

The ends  of words and sentences are marked
  by   spaces. It  doesn't matter how many
spaces    you type; one is as good as 100.  The
end of   a line counts as a space.

One   or more   blank lines denote the  end
of  a paragraph.

Since any number of consecutive spaces are treated
like a single one, the formatting of the input
file makes no difference to
      \LaTeX,                % The \LaTeX command generates the LaTeX logo.
but it makes a difference to you.  When you use
\LaTeX, making your input file as easy to read
as possible will be a great help as you write
your document and when you change it.  This sample
file shows how you can add comments to your own input
file.

Because printing is different from typewriting,
there are a number of things that you have to do
differently when preparing an input file than if
you were just typing the document directly.
Quotation marks like
       ``this''
have to be handled specially, as do quotes within
quotes:
       ``\,`this'            % \, separates the double and single quote.
        is what I just
        wrote, not  `that'\,''.

Dashes come in three sizes: an
       intra-word
dash, a medium dash for number ranges like
       1--2,
and a punctuation
       dash---like
this.

A sentence-ending space should be larger than the
space between words within a sentence.  You
sometimes have to type special commands in
conjunction with punctuation characters to get
this right, as in the following sentence.
       Gnats, gnus, etc.\ all  % `\ ' makes an inter-word space.
       begin with G\@.         % \@ marks end-of-sentence punctuation.
You should check the spaces after periods when
reading your output to make sure you haven't
forgotten any special cases.  Generating an
ellipsis
       \ldots\               % `\ ' is needed after `\ldots' because TeX
                             % ignores spaces after command names like \ldots
                             % made from \ + letters.
                             %
                             % Note how a `%' character causes TeX to ignore
                             % the end of the input line, so these blank lines
                             % do not start a new paragraph.
                             %
with the right spacing around the periods requires
a special command.

\LaTeX\ interprets some common characters as
commands, so you must type special commands to
generate them.  These characters include the
following:
       \$ \& \% \# \{ and \}.

In printing, text is usually emphasized with an
       \emph{italic}
type style.

\begin{em}
   A long segment of text can also be emphasized
   in this way.  Text within such a segment can be
   given \emph{additional} emphasis.
\end{em}

It is sometimes necessary to prevent \LaTeX\ from
breaking a line where it might otherwise do so.
This may be at a space, as between the ``Mr.''\ and
``Jones'' in
       ``Mr.~Jones'',        % ~ produces an unbreakable interword space.
or within a word---especially when the word is a
symbol like
       \mbox{\emph{itemnum}}
that makes little sense when hyphenated across
lines.

Footnotes\footnote{This is an example of a footnote.}
pose no problem.

\LaTeX\ is good at typesetting mathematical formulas
like
       \( x-3y + z = 7 \)
or
       \( a_{1} > x^{2n} + y^{2n} > x' \)
or
       \( \ip{A}{B} = \sum_{i} a_{i} b_{i} \).
The spaces you type in a formula are
ignored.  Remember that a letter like
       $x$                   % $ ... $  and  \( ... \)  are equivalent
is a formula when it denotes a mathematical
symbol, and it should be typed as one.

\section{Displayed Text}

Text is displayed by indenting it from the left
margin.  Quotations are commonly displayed.  There
are short quotations
\begin{quote}
   This is a short quotation.  It consists of a
   single paragraph of text.  See how it is formatted.
\end{quote}
and longer ones.
\begin{quotation}
   This is a longer quotation.  It consists of two
   paragraphs of text, neither of which are
   particularly interesting.

   This is the second paragraph of the quotation.  It
   is just as dull as the first paragraph.
\end{quotation}
Another frequently-displayed structure is a list.
The following is an example of an \emph{itemized}
list.
\begin{itemize}
   \item This is the first item of an itemized list.
         Each item in the list is marked with a ``tick''.
         You don't have to worry about what kind of tick
         mark is used.

   \item This is the second item of the list.  It
         contains another list nested inside it.  The inner
         list is an \emph{enumerated} list.
         \begin{enumerate}
            \item This is the first item of an enumerated
                  list that is nested within the itemized list.

            \item This is the second item of the inner list.
                  \LaTeX\ allows you to nest lists deeper than
                  you really should.
         \end{enumerate}
         This is the rest of the second item of the outer
         list.  It is no more interesting than any other
         part of the item.
   \item This is the third item of the list.
\end{itemize}
You can even display poetry.
\begin{verse}
   There is an environment
    for verse \\             % The \\ command separates lines
   Whose features some poets % within a stanza.
   will curse.

                             % One or more blank lines separate stanzas.

   For instead of making\\
   Them do \emph{all} line breaking, \\
   It allows them to put too many words on a line when they'd rather be
   forced to be terse.
\end{verse}

Mathematical formulas may also be displayed.  A
displayed formula
is
one-line long; multiline
formulas require special formatting instructions.
   \[  \ip{\Gamma}{\psi'} = x'' + y^{2} + z_{i}^{n}\]
Don't start a paragraph with a displayed equation,
nor make one a paragraph by itself.

\end{document}               % End of document.

\section{Bench Huge Section 143}
% synthetic bench section 143
% This is a sample LaTeX input file.  (Version of 12 August 2004.)
%
% A '%' character causes TeX to ignore all remaining text on the line,
% and is used for comments like this one.

\documentclass{article}      % Specifies the document class

                             % The preamble begins here.
\title{An Example Document}  % Declares the document's title.
\author{Leslie Lamport}      % Declares the author's name.
\date{January 21, 1994}      % Deleting this command produces today's date.

\newcommand{\ip}[2]{(#1, #2)}
                             % Defines \ip{arg1}{arg2} to mean
                             % (arg1, arg2).

%\newcommand{\ip}[2]{\langle #1 | #2\rangle}
                             % This is an alternative definition of
                             % \ip that is commented out.

\begin{document}             % End of preamble and beginning of text.

\maketitle                   % Produces the title.

This is an example input file.  Comparing it with
the output it generates can show you how to
produce a simple document of your own.

\section{Ordinary Text}      % Produces section heading.  Lower-level
                             % sections are begun with similar
                             % \subsection and \subsubsection commands.

The ends  of words and sentences are marked
  by   spaces. It  doesn't matter how many
spaces    you type; one is as good as 100.  The
end of   a line counts as a space.

One   or more   blank lines denote the  end
of  a paragraph.

Since any number of consecutive spaces are treated
like a single one, the formatting of the input
file makes no difference to
      \LaTeX,                % The \LaTeX command generates the LaTeX logo.
but it makes a difference to you.  When you use
\LaTeX, making your input file as easy to read
as possible will be a great help as you write
your document and when you change it.  This sample
file shows how you can add comments to your own input
file.

Because printing is different from typewriting,
there are a number of things that you have to do
differently when preparing an input file than if
you were just typing the document directly.
Quotation marks like
       ``this''
have to be handled specially, as do quotes within
quotes:
       ``\,`this'            % \, separates the double and single quote.
        is what I just
        wrote, not  `that'\,''.

Dashes come in three sizes: an
       intra-word
dash, a medium dash for number ranges like
       1--2,
and a punctuation
       dash---like
this.

A sentence-ending space should be larger than the
space between words within a sentence.  You
sometimes have to type special commands in
conjunction with punctuation characters to get
this right, as in the following sentence.
       Gnats, gnus, etc.\ all  % `\ ' makes an inter-word space.
       begin with G\@.         % \@ marks end-of-sentence punctuation.
You should check the spaces after periods when
reading your output to make sure you haven't
forgotten any special cases.  Generating an
ellipsis
       \ldots\               % `\ ' is needed after `\ldots' because TeX
                             % ignores spaces after command names like \ldots
                             % made from \ + letters.
                             %
                             % Note how a `%' character causes TeX to ignore
                             % the end of the input line, so these blank lines
                             % do not start a new paragraph.
                             %
with the right spacing around the periods requires
a special command.

\LaTeX\ interprets some common characters as
commands, so you must type special commands to
generate them.  These characters include the
following:
       \$ \& \% \# \{ and \}.

In printing, text is usually emphasized with an
       \emph{italic}
type style.

\begin{em}
   A long segment of text can also be emphasized
   in this way.  Text within such a segment can be
   given \emph{additional} emphasis.
\end{em}

It is sometimes necessary to prevent \LaTeX\ from
breaking a line where it might otherwise do so.
This may be at a space, as between the ``Mr.''\ and
``Jones'' in
       ``Mr.~Jones'',        % ~ produces an unbreakable interword space.
or within a word---especially when the word is a
symbol like
       \mbox{\emph{itemnum}}
that makes little sense when hyphenated across
lines.

Footnotes\footnote{This is an example of a footnote.}
pose no problem.

\LaTeX\ is good at typesetting mathematical formulas
like
       \( x-3y + z = 7 \)
or
       \( a_{1} > x^{2n} + y^{2n} > x' \)
or
       \( \ip{A}{B} = \sum_{i} a_{i} b_{i} \).
The spaces you type in a formula are
ignored.  Remember that a letter like
       $x$                   % $ ... $  and  \( ... \)  are equivalent
is a formula when it denotes a mathematical
symbol, and it should be typed as one.

\section{Displayed Text}

Text is displayed by indenting it from the left
margin.  Quotations are commonly displayed.  There
are short quotations
\begin{quote}
   This is a short quotation.  It consists of a
   single paragraph of text.  See how it is formatted.
\end{quote}
and longer ones.
\begin{quotation}
   This is a longer quotation.  It consists of two
   paragraphs of text, neither of which are
   particularly interesting.

   This is the second paragraph of the quotation.  It
   is just as dull as the first paragraph.
\end{quotation}
Another frequently-displayed structure is a list.
The following is an example of an \emph{itemized}
list.
\begin{itemize}
   \item This is the first item of an itemized list.
         Each item in the list is marked with a ``tick''.
         You don't have to worry about what kind of tick
         mark is used.

   \item This is the second item of the list.  It
         contains another list nested inside it.  The inner
         list is an \emph{enumerated} list.
         \begin{enumerate}
            \item This is the first item of an enumerated
                  list that is nested within the itemized list.

            \item This is the second item of the inner list.
                  \LaTeX\ allows you to nest lists deeper than
                  you really should.
         \end{enumerate}
         This is the rest of the second item of the outer
         list.  It is no more interesting than any other
         part of the item.
   \item This is the third item of the list.
\end{itemize}
You can even display poetry.
\begin{verse}
   There is an environment
    for verse \\             % The \\ command separates lines
   Whose features some poets % within a stanza.
   will curse.

                             % One or more blank lines separate stanzas.

   For instead of making\\
   Them do \emph{all} line breaking, \\
   It allows them to put too many words on a line when they'd rather be
   forced to be terse.
\end{verse}

Mathematical formulas may also be displayed.  A
displayed formula
is
one-line long; multiline
formulas require special formatting instructions.
   \[  \ip{\Gamma}{\psi'} = x'' + y^{2} + z_{i}^{n}\]
Don't start a paragraph with a displayed equation,
nor make one a paragraph by itself.

\end{document}               % End of document.

\section{Bench Huge Section 144}
% synthetic bench section 144
% This is a sample LaTeX input file.  (Version of 12 August 2004.)
%
% A '%' character causes TeX to ignore all remaining text on the line,
% and is used for comments like this one.

\documentclass{article}      % Specifies the document class

                             % The preamble begins here.
\title{An Example Document}  % Declares the document's title.
\author{Leslie Lamport}      % Declares the author's name.
\date{January 21, 1994}      % Deleting this command produces today's date.

\newcommand{\ip}[2]{(#1, #2)}
                             % Defines \ip{arg1}{arg2} to mean
                             % (arg1, arg2).

%\newcommand{\ip}[2]{\langle #1 | #2\rangle}
                             % This is an alternative definition of
                             % \ip that is commented out.

\begin{document}             % End of preamble and beginning of text.

\maketitle                   % Produces the title.

This is an example input file.  Comparing it with
the output it generates can show you how to
produce a simple document of your own.

\section{Ordinary Text}      % Produces section heading.  Lower-level
                             % sections are begun with similar
                             % \subsection and \subsubsection commands.

The ends  of words and sentences are marked
  by   spaces. It  doesn't matter how many
spaces    you type; one is as good as 100.  The
end of   a line counts as a space.

One   or more   blank lines denote the  end
of  a paragraph.

Since any number of consecutive spaces are treated
like a single one, the formatting of the input
file makes no difference to
      \LaTeX,                % The \LaTeX command generates the LaTeX logo.
but it makes a difference to you.  When you use
\LaTeX, making your input file as easy to read
as possible will be a great help as you write
your document and when you change it.  This sample
file shows how you can add comments to your own input
file.

Because printing is different from typewriting,
there are a number of things that you have to do
differently when preparing an input file than if
you were just typing the document directly.
Quotation marks like
       ``this''
have to be handled specially, as do quotes within
quotes:
       ``\,`this'            % \, separates the double and single quote.
        is what I just
        wrote, not  `that'\,''.

Dashes come in three sizes: an
       intra-word
dash, a medium dash for number ranges like
       1--2,
and a punctuation
       dash---like
this.

A sentence-ending space should be larger than the
space between words within a sentence.  You
sometimes have to type special commands in
conjunction with punctuation characters to get
this right, as in the following sentence.
       Gnats, gnus, etc.\ all  % `\ ' makes an inter-word space.
       begin with G\@.         % \@ marks end-of-sentence punctuation.
You should check the spaces after periods when
reading your output to make sure you haven't
forgotten any special cases.  Generating an
ellipsis
       \ldots\               % `\ ' is needed after `\ldots' because TeX
                             % ignores spaces after command names like \ldots
                             % made from \ + letters.
                             %
                             % Note how a `%' character causes TeX to ignore
                             % the end of the input line, so these blank lines
                             % do not start a new paragraph.
                             %
with the right spacing around the periods requires
a special command.

\LaTeX\ interprets some common characters as
commands, so you must type special commands to
generate them.  These characters include the
following:
       \$ \& \% \# \{ and \}.

In printing, text is usually emphasized with an
       \emph{italic}
type style.

\begin{em}
   A long segment of text can also be emphasized
   in this way.  Text within such a segment can be
   given \emph{additional} emphasis.
\end{em}

It is sometimes necessary to prevent \LaTeX\ from
breaking a line where it might otherwise do so.
This may be at a space, as between the ``Mr.''\ and
``Jones'' in
       ``Mr.~Jones'',        % ~ produces an unbreakable interword space.
or within a word---especially when the word is a
symbol like
       \mbox{\emph{itemnum}}
that makes little sense when hyphenated across
lines.

Footnotes\footnote{This is an example of a footnote.}
pose no problem.

\LaTeX\ is good at typesetting mathematical formulas
like
       \( x-3y + z = 7 \)
or
       \( a_{1} > x^{2n} + y^{2n} > x' \)
or
       \( \ip{A}{B} = \sum_{i} a_{i} b_{i} \).
The spaces you type in a formula are
ignored.  Remember that a letter like
       $x$                   % $ ... $  and  \( ... \)  are equivalent
is a formula when it denotes a mathematical
symbol, and it should be typed as one.

\section{Displayed Text}

Text is displayed by indenting it from the left
margin.  Quotations are commonly displayed.  There
are short quotations
\begin{quote}
   This is a short quotation.  It consists of a
   single paragraph of text.  See how it is formatted.
\end{quote}
and longer ones.
\begin{quotation}
   This is a longer quotation.  It consists of two
   paragraphs of text, neither of which are
   particularly interesting.

   This is the second paragraph of the quotation.  It
   is just as dull as the first paragraph.
\end{quotation}
Another frequently-displayed structure is a list.
The following is an example of an \emph{itemized}
list.
\begin{itemize}
   \item This is the first item of an itemized list.
         Each item in the list is marked with a ``tick''.
         You don't have to worry about what kind of tick
         mark is used.

   \item This is the second item of the list.  It
         contains another list nested inside it.  The inner
         list is an \emph{enumerated} list.
         \begin{enumerate}
            \item This is the first item of an enumerated
                  list that is nested within the itemized list.

            \item This is the second item of the inner list.
                  \LaTeX\ allows you to nest lists deeper than
                  you really should.
         \end{enumerate}
         This is the rest of the second item of the outer
         list.  It is no more interesting than any other
         part of the item.
   \item This is the third item of the list.
\end{itemize}
You can even display poetry.
\begin{verse}
   There is an environment
    for verse \\             % The \\ command separates lines
   Whose features some poets % within a stanza.
   will curse.

                             % One or more blank lines separate stanzas.

   For instead of making\\
   Them do \emph{all} line breaking, \\
   It allows them to put too many words on a line when they'd rather be
   forced to be terse.
\end{verse}

Mathematical formulas may also be displayed.  A
displayed formula
is
one-line long; multiline
formulas require special formatting instructions.
   \[  \ip{\Gamma}{\psi'} = x'' + y^{2} + z_{i}^{n}\]
Don't start a paragraph with a displayed equation,
nor make one a paragraph by itself.

\end{document}               % End of document.

\section{Bench Huge Section 145}
% synthetic bench section 145
% This is a sample LaTeX input file.  (Version of 12 August 2004.)
%
% A '%' character causes TeX to ignore all remaining text on the line,
% and is used for comments like this one.

\documentclass{article}      % Specifies the document class

                             % The preamble begins here.
\title{An Example Document}  % Declares the document's title.
\author{Leslie Lamport}      % Declares the author's name.
\date{January 21, 1994}      % Deleting this command produces today's date.

\newcommand{\ip}[2]{(#1, #2)}
                             % Defines \ip{arg1}{arg2} to mean
                             % (arg1, arg2).

%\newcommand{\ip}[2]{\langle #1 | #2\rangle}
                             % This is an alternative definition of
                             % \ip that is commented out.

\begin{document}             % End of preamble and beginning of text.

\maketitle                   % Produces the title.

This is an example input file.  Comparing it with
the output it generates can show you how to
produce a simple document of your own.

\section{Ordinary Text}      % Produces section heading.  Lower-level
                             % sections are begun with similar
                             % \subsection and \subsubsection commands.

The ends  of words and sentences are marked
  by   spaces. It  doesn't matter how many
spaces    you type; one is as good as 100.  The
end of   a line counts as a space.

One   or more   blank lines denote the  end
of  a paragraph.

Since any number of consecutive spaces are treated
like a single one, the formatting of the input
file makes no difference to
      \LaTeX,                % The \LaTeX command generates the LaTeX logo.
but it makes a difference to you.  When you use
\LaTeX, making your input file as easy to read
as possible will be a great help as you write
your document and when you change it.  This sample
file shows how you can add comments to your own input
file.

Because printing is different from typewriting,
there are a number of things that you have to do
differently when preparing an input file than if
you were just typing the document directly.
Quotation marks like
       ``this''
have to be handled specially, as do quotes within
quotes:
       ``\,`this'            % \, separates the double and single quote.
        is what I just
        wrote, not  `that'\,''.

Dashes come in three sizes: an
       intra-word
dash, a medium dash for number ranges like
       1--2,
and a punctuation
       dash---like
this.

A sentence-ending space should be larger than the
space between words within a sentence.  You
sometimes have to type special commands in
conjunction with punctuation characters to get
this right, as in the following sentence.
       Gnats, gnus, etc.\ all  % `\ ' makes an inter-word space.
       begin with G\@.         % \@ marks end-of-sentence punctuation.
You should check the spaces after periods when
reading your output to make sure you haven't
forgotten any special cases.  Generating an
ellipsis
       \ldots\               % `\ ' is needed after `\ldots' because TeX
                             % ignores spaces after command names like \ldots
                             % made from \ + letters.
                             %
                             % Note how a `%' character causes TeX to ignore
                             % the end of the input line, so these blank lines
                             % do not start a new paragraph.
                             %
with the right spacing around the periods requires
a special command.

\LaTeX\ interprets some common characters as
commands, so you must type special commands to
generate them.  These characters include the
following:
       \$ \& \% \# \{ and \}.

In printing, text is usually emphasized with an
       \emph{italic}
type style.

\begin{em}
   A long segment of text can also be emphasized
   in this way.  Text within such a segment can be
   given \emph{additional} emphasis.
\end{em}

It is sometimes necessary to prevent \LaTeX\ from
breaking a line where it might otherwise do so.
This may be at a space, as between the ``Mr.''\ and
``Jones'' in
       ``Mr.~Jones'',        % ~ produces an unbreakable interword space.
or within a word---especially when the word is a
symbol like
       \mbox{\emph{itemnum}}
that makes little sense when hyphenated across
lines.

Footnotes\footnote{This is an example of a footnote.}
pose no problem.

\LaTeX\ is good at typesetting mathematical formulas
like
       \( x-3y + z = 7 \)
or
       \( a_{1} > x^{2n} + y^{2n} > x' \)
or
       \( \ip{A}{B} = \sum_{i} a_{i} b_{i} \).
The spaces you type in a formula are
ignored.  Remember that a letter like
       $x$                   % $ ... $  and  \( ... \)  are equivalent
is a formula when it denotes a mathematical
symbol, and it should be typed as one.

\section{Displayed Text}

Text is displayed by indenting it from the left
margin.  Quotations are commonly displayed.  There
are short quotations
\begin{quote}
   This is a short quotation.  It consists of a
   single paragraph of text.  See how it is formatted.
\end{quote}
and longer ones.
\begin{quotation}
   This is a longer quotation.  It consists of two
   paragraphs of text, neither of which are
   particularly interesting.

   This is the second paragraph of the quotation.  It
   is just as dull as the first paragraph.
\end{quotation}
Another frequently-displayed structure is a list.
The following is an example of an \emph{itemized}
list.
\begin{itemize}
   \item This is the first item of an itemized list.
         Each item in the list is marked with a ``tick''.
         You don't have to worry about what kind of tick
         mark is used.

   \item This is the second item of the list.  It
         contains another list nested inside it.  The inner
         list is an \emph{enumerated} list.
         \begin{enumerate}
            \item This is the first item of an enumerated
                  list that is nested within the itemized list.

            \item This is the second item of the inner list.
                  \LaTeX\ allows you to nest lists deeper than
                  you really should.
         \end{enumerate}
         This is the rest of the second item of the outer
         list.  It is no more interesting than any other
         part of the item.
   \item This is the third item of the list.
\end{itemize}
You can even display poetry.
\begin{verse}
   There is an environment
    for verse \\             % The \\ command separates lines
   Whose features some poets % within a stanza.
   will curse.

                             % One or more blank lines separate stanzas.

   For instead of making\\
   Them do \emph{all} line breaking, \\
   It allows them to put too many words on a line when they'd rather be
   forced to be terse.
\end{verse}

Mathematical formulas may also be displayed.  A
displayed formula
is
one-line long; multiline
formulas require special formatting instructions.
   \[  \ip{\Gamma}{\psi'} = x'' + y^{2} + z_{i}^{n}\]
Don't start a paragraph with a displayed equation,
nor make one a paragraph by itself.

\end{document}               % End of document.

\section{Bench Huge Section 146}
% synthetic bench section 146
% This is a sample LaTeX input file.  (Version of 12 August 2004.)
%
% A '%' character causes TeX to ignore all remaining text on the line,
% and is used for comments like this one.

\documentclass{article}      % Specifies the document class

                             % The preamble begins here.
\title{An Example Document}  % Declares the document's title.
\author{Leslie Lamport}      % Declares the author's name.
\date{January 21, 1994}      % Deleting this command produces today's date.

\newcommand{\ip}[2]{(#1, #2)}
                             % Defines \ip{arg1}{arg2} to mean
                             % (arg1, arg2).

%\newcommand{\ip}[2]{\langle #1 | #2\rangle}
                             % This is an alternative definition of
                             % \ip that is commented out.

\begin{document}             % End of preamble and beginning of text.

\maketitle                   % Produces the title.

This is an example input file.  Comparing it with
the output it generates can show you how to
produce a simple document of your own.

\section{Ordinary Text}      % Produces section heading.  Lower-level
                             % sections are begun with similar
                             % \subsection and \subsubsection commands.

The ends  of words and sentences are marked
  by   spaces. It  doesn't matter how many
spaces    you type; one is as good as 100.  The
end of   a line counts as a space.

One   or more   blank lines denote the  end
of  a paragraph.

Since any number of consecutive spaces are treated
like a single one, the formatting of the input
file makes no difference to
      \LaTeX,                % The \LaTeX command generates the LaTeX logo.
but it makes a difference to you.  When you use
\LaTeX, making your input file as easy to read
as possible will be a great help as you write
your document and when you change it.  This sample
file shows how you can add comments to your own input
file.

Because printing is different from typewriting,
there are a number of things that you have to do
differently when preparing an input file than if
you were just typing the document directly.
Quotation marks like
       ``this''
have to be handled specially, as do quotes within
quotes:
       ``\,`this'            % \, separates the double and single quote.
        is what I just
        wrote, not  `that'\,''.

Dashes come in three sizes: an
       intra-word
dash, a medium dash for number ranges like
       1--2,
and a punctuation
       dash---like
this.

A sentence-ending space should be larger than the
space between words within a sentence.  You
sometimes have to type special commands in
conjunction with punctuation characters to get
this right, as in the following sentence.
       Gnats, gnus, etc.\ all  % `\ ' makes an inter-word space.
       begin with G\@.         % \@ marks end-of-sentence punctuation.
You should check the spaces after periods when
reading your output to make sure you haven't
forgotten any special cases.  Generating an
ellipsis
       \ldots\               % `\ ' is needed after `\ldots' because TeX
                             % ignores spaces after command names like \ldots
                             % made from \ + letters.
                             %
                             % Note how a `%' character causes TeX to ignore
                             % the end of the input line, so these blank lines
                             % do not start a new paragraph.
                             %
with the right spacing around the periods requires
a special command.

\LaTeX\ interprets some common characters as
commands, so you must type special commands to
generate them.  These characters include the
following:
       \$ \& \% \# \{ and \}.

In printing, text is usually emphasized with an
       \emph{italic}
type style.

\begin{em}
   A long segment of text can also be emphasized
   in this way.  Text within such a segment can be
   given \emph{additional} emphasis.
\end{em}

It is sometimes necessary to prevent \LaTeX\ from
breaking a line where it might otherwise do so.
This may be at a space, as between the ``Mr.''\ and
``Jones'' in
       ``Mr.~Jones'',        % ~ produces an unbreakable interword space.
or within a word---especially when the word is a
symbol like
       \mbox{\emph{itemnum}}
that makes little sense when hyphenated across
lines.

Footnotes\footnote{This is an example of a footnote.}
pose no problem.

\LaTeX\ is good at typesetting mathematical formulas
like
       \( x-3y + z = 7 \)
or
       \( a_{1} > x^{2n} + y^{2n} > x' \)
or
       \( \ip{A}{B} = \sum_{i} a_{i} b_{i} \).
The spaces you type in a formula are
ignored.  Remember that a letter like
       $x$                   % $ ... $  and  \( ... \)  are equivalent
is a formula when it denotes a mathematical
symbol, and it should be typed as one.

\section{Displayed Text}

Text is displayed by indenting it from the left
margin.  Quotations are commonly displayed.  There
are short quotations
\begin{quote}
   This is a short quotation.  It consists of a
   single paragraph of text.  See how it is formatted.
\end{quote}
and longer ones.
\begin{quotation}
   This is a longer quotation.  It consists of two
   paragraphs of text, neither of which are
   particularly interesting.

   This is the second paragraph of the quotation.  It
   is just as dull as the first paragraph.
\end{quotation}
Another frequently-displayed structure is a list.
The following is an example of an \emph{itemized}
list.
\begin{itemize}
   \item This is the first item of an itemized list.
         Each item in the list is marked with a ``tick''.
         You don't have to worry about what kind of tick
         mark is used.

   \item This is the second item of the list.  It
         contains another list nested inside it.  The inner
         list is an \emph{enumerated} list.
         \begin{enumerate}
            \item This is the first item of an enumerated
                  list that is nested within the itemized list.

            \item This is the second item of the inner list.
                  \LaTeX\ allows you to nest lists deeper than
                  you really should.
         \end{enumerate}
         This is the rest of the second item of the outer
         list.  It is no more interesting than any other
         part of the item.
   \item This is the third item of the list.
\end{itemize}
You can even display poetry.
\begin{verse}
   There is an environment
    for verse \\             % The \\ command separates lines
   Whose features some poets % within a stanza.
   will curse.

                             % One or more blank lines separate stanzas.

   For instead of making\\
   Them do \emph{all} line breaking, \\
   It allows them to put too many words on a line when they'd rather be
   forced to be terse.
\end{verse}

Mathematical formulas may also be displayed.  A
displayed formula
is
one-line long; multiline
formulas require special formatting instructions.
   \[  \ip{\Gamma}{\psi'} = x'' + y^{2} + z_{i}^{n}\]
Don't start a paragraph with a displayed equation,
nor make one a paragraph by itself.

\end{document}               % End of document.

\section{Bench Huge Section 147}
% synthetic bench section 147
% This is a sample LaTeX input file.  (Version of 12 August 2004.)
%
% A '%' character causes TeX to ignore all remaining text on the line,
% and is used for comments like this one.

\documentclass{article}      % Specifies the document class

                             % The preamble begins here.
\title{An Example Document}  % Declares the document's title.
\author{Leslie Lamport}      % Declares the author's name.
\date{January 21, 1994}      % Deleting this command produces today's date.

\newcommand{\ip}[2]{(#1, #2)}
                             % Defines \ip{arg1}{arg2} to mean
                             % (arg1, arg2).

%\newcommand{\ip}[2]{\langle #1 | #2\rangle}
                             % This is an alternative definition of
                             % \ip that is commented out.

\begin{document}             % End of preamble and beginning of text.

\maketitle                   % Produces the title.

This is an example input file.  Comparing it with
the output it generates can show you how to
produce a simple document of your own.

\section{Ordinary Text}      % Produces section heading.  Lower-level
                             % sections are begun with similar
                             % \subsection and \subsubsection commands.

The ends  of words and sentences are marked
  by   spaces. It  doesn't matter how many
spaces    you type; one is as good as 100.  The
end of   a line counts as a space.

One   or more   blank lines denote the  end
of  a paragraph.

Since any number of consecutive spaces are treated
like a single one, the formatting of the input
file makes no difference to
      \LaTeX,                % The \LaTeX command generates the LaTeX logo.
but it makes a difference to you.  When you use
\LaTeX, making your input file as easy to read
as possible will be a great help as you write
your document and when you change it.  This sample
file shows how you can add comments to your own input
file.

Because printing is different from typewriting,
there are a number of things that you have to do
differently when preparing an input file than if
you were just typing the document directly.
Quotation marks like
       ``this''
have to be handled specially, as do quotes within
quotes:
       ``\,`this'            % \, separates the double and single quote.
        is what I just
        wrote, not  `that'\,''.

Dashes come in three sizes: an
       intra-word
dash, a medium dash for number ranges like
       1--2,
and a punctuation
       dash---like
this.

A sentence-ending space should be larger than the
space between words within a sentence.  You
sometimes have to type special commands in
conjunction with punctuation characters to get
this right, as in the following sentence.
       Gnats, gnus, etc.\ all  % `\ ' makes an inter-word space.
       begin with G\@.         % \@ marks end-of-sentence punctuation.
You should check the spaces after periods when
reading your output to make sure you haven't
forgotten any special cases.  Generating an
ellipsis
       \ldots\               % `\ ' is needed after `\ldots' because TeX
                             % ignores spaces after command names like \ldots
                             % made from \ + letters.
                             %
                             % Note how a `%' character causes TeX to ignore
                             % the end of the input line, so these blank lines
                             % do not start a new paragraph.
                             %
with the right spacing around the periods requires
a special command.

\LaTeX\ interprets some common characters as
commands, so you must type special commands to
generate them.  These characters include the
following:
       \$ \& \% \# \{ and \}.

In printing, text is usually emphasized with an
       \emph{italic}
type style.

\begin{em}
   A long segment of text can also be emphasized
   in this way.  Text within such a segment can be
   given \emph{additional} emphasis.
\end{em}

It is sometimes necessary to prevent \LaTeX\ from
breaking a line where it might otherwise do so.
This may be at a space, as between the ``Mr.''\ and
``Jones'' in
       ``Mr.~Jones'',        % ~ produces an unbreakable interword space.
or within a word---especially when the word is a
symbol like
       \mbox{\emph{itemnum}}
that makes little sense when hyphenated across
lines.

Footnotes\footnote{This is an example of a footnote.}
pose no problem.

\LaTeX\ is good at typesetting mathematical formulas
like
       \( x-3y + z = 7 \)
or
       \( a_{1} > x^{2n} + y^{2n} > x' \)
or
       \( \ip{A}{B} = \sum_{i} a_{i} b_{i} \).
The spaces you type in a formula are
ignored.  Remember that a letter like
       $x$                   % $ ... $  and  \( ... \)  are equivalent
is a formula when it denotes a mathematical
symbol, and it should be typed as one.

\section{Displayed Text}

Text is displayed by indenting it from the left
margin.  Quotations are commonly displayed.  There
are short quotations
\begin{quote}
   This is a short quotation.  It consists of a
   single paragraph of text.  See how it is formatted.
\end{quote}
and longer ones.
\begin{quotation}
   This is a longer quotation.  It consists of two
   paragraphs of text, neither of which are
   particularly interesting.

   This is the second paragraph of the quotation.  It
   is just as dull as the first paragraph.
\end{quotation}
Another frequently-displayed structure is a list.
The following is an example of an \emph{itemized}
list.
\begin{itemize}
   \item This is the first item of an itemized list.
         Each item in the list is marked with a ``tick''.
         You don't have to worry about what kind of tick
         mark is used.

   \item This is the second item of the list.  It
         contains another list nested inside it.  The inner
         list is an \emph{enumerated} list.
         \begin{enumerate}
            \item This is the first item of an enumerated
                  list that is nested within the itemized list.

            \item This is the second item of the inner list.
                  \LaTeX\ allows you to nest lists deeper than
                  you really should.
         \end{enumerate}
         This is the rest of the second item of the outer
         list.  It is no more interesting than any other
         part of the item.
   \item This is the third item of the list.
\end{itemize}
You can even display poetry.
\begin{verse}
   There is an environment
    for verse \\             % The \\ command separates lines
   Whose features some poets % within a stanza.
   will curse.

                             % One or more blank lines separate stanzas.

   For instead of making\\
   Them do \emph{all} line breaking, \\
   It allows them to put too many words on a line when they'd rather be
   forced to be terse.
\end{verse}

Mathematical formulas may also be displayed.  A
displayed formula
is
one-line long; multiline
formulas require special formatting instructions.
   \[  \ip{\Gamma}{\psi'} = x'' + y^{2} + z_{i}^{n}\]
Don't start a paragraph with a displayed equation,
nor make one a paragraph by itself.

\end{document}               % End of document.

\section{Bench Huge Section 148}
% synthetic bench section 148
% This is a sample LaTeX input file.  (Version of 12 August 2004.)
%
% A '%' character causes TeX to ignore all remaining text on the line,
% and is used for comments like this one.

\documentclass{article}      % Specifies the document class

                             % The preamble begins here.
\title{An Example Document}  % Declares the document's title.
\author{Leslie Lamport}      % Declares the author's name.
\date{January 21, 1994}      % Deleting this command produces today's date.

\newcommand{\ip}[2]{(#1, #2)}
                             % Defines \ip{arg1}{arg2} to mean
                             % (arg1, arg2).

%\newcommand{\ip}[2]{\langle #1 | #2\rangle}
                             % This is an alternative definition of
                             % \ip that is commented out.

\begin{document}             % End of preamble and beginning of text.

\maketitle                   % Produces the title.

This is an example input file.  Comparing it with
the output it generates can show you how to
produce a simple document of your own.

\section{Ordinary Text}      % Produces section heading.  Lower-level
                             % sections are begun with similar
                             % \subsection and \subsubsection commands.

The ends  of words and sentences are marked
  by   spaces. It  doesn't matter how many
spaces    you type; one is as good as 100.  The
end of   a line counts as a space.

One   or more   blank lines denote the  end
of  a paragraph.

Since any number of consecutive spaces are treated
like a single one, the formatting of the input
file makes no difference to
      \LaTeX,                % The \LaTeX command generates the LaTeX logo.
but it makes a difference to you.  When you use
\LaTeX, making your input file as easy to read
as possible will be a great help as you write
your document and when you change it.  This sample
file shows how you can add comments to your own input
file.

Because printing is different from typewriting,
there are a number of things that you have to do
differently when preparing an input file than if
you were just typing the document directly.
Quotation marks like
       ``this''
have to be handled specially, as do quotes within
quotes:
       ``\,`this'            % \, separates the double and single quote.
        is what I just
        wrote, not  `that'\,''.

Dashes come in three sizes: an
       intra-word
dash, a medium dash for number ranges like
       1--2,
and a punctuation
       dash---like
this.

A sentence-ending space should be larger than the
space between words within a sentence.  You
sometimes have to type special commands in
conjunction with punctuation characters to get
this right, as in the following sentence.
       Gnats, gnus, etc.\ all  % `\ ' makes an inter-word space.
       begin with G\@.         % \@ marks end-of-sentence punctuation.
You should check the spaces after periods when
reading your output to make sure you haven't
forgotten any special cases.  Generating an
ellipsis
       \ldots\               % `\ ' is needed after `\ldots' because TeX
                             % ignores spaces after command names like \ldots
                             % made from \ + letters.
                             %
                             % Note how a `%' character causes TeX to ignore
                             % the end of the input line, so these blank lines
                             % do not start a new paragraph.
                             %
with the right spacing around the periods requires
a special command.

\LaTeX\ interprets some common characters as
commands, so you must type special commands to
generate them.  These characters include the
following:
       \$ \& \% \# \{ and \}.

In printing, text is usually emphasized with an
       \emph{italic}
type style.

\begin{em}
   A long segment of text can also be emphasized
   in this way.  Text within such a segment can be
   given \emph{additional} emphasis.
\end{em}

It is sometimes necessary to prevent \LaTeX\ from
breaking a line where it might otherwise do so.
This may be at a space, as between the ``Mr.''\ and
``Jones'' in
       ``Mr.~Jones'',        % ~ produces an unbreakable interword space.
or within a word---especially when the word is a
symbol like
       \mbox{\emph{itemnum}}
that makes little sense when hyphenated across
lines.

Footnotes\footnote{This is an example of a footnote.}
pose no problem.

\LaTeX\ is good at typesetting mathematical formulas
like
       \( x-3y + z = 7 \)
or
       \( a_{1} > x^{2n} + y^{2n} > x' \)
or
       \( \ip{A}{B} = \sum_{i} a_{i} b_{i} \).
The spaces you type in a formula are
ignored.  Remember that a letter like
       $x$                   % $ ... $  and  \( ... \)  are equivalent
is a formula when it denotes a mathematical
symbol, and it should be typed as one.

\section{Displayed Text}

Text is displayed by indenting it from the left
margin.  Quotations are commonly displayed.  There
are short quotations
\begin{quote}
   This is a short quotation.  It consists of a
   single paragraph of text.  See how it is formatted.
\end{quote}
and longer ones.
\begin{quotation}
   This is a longer quotation.  It consists of two
   paragraphs of text, neither of which are
   particularly interesting.

   This is the second paragraph of the quotation.  It
   is just as dull as the first paragraph.
\end{quotation}
Another frequently-displayed structure is a list.
The following is an example of an \emph{itemized}
list.
\begin{itemize}
   \item This is the first item of an itemized list.
         Each item in the list is marked with a ``tick''.
         You don't have to worry about what kind of tick
         mark is used.

   \item This is the second item of the list.  It
         contains another list nested inside it.  The inner
         list is an \emph{enumerated} list.
         \begin{enumerate}
            \item This is the first item of an enumerated
                  list that is nested within the itemized list.

            \item This is the second item of the inner list.
                  \LaTeX\ allows you to nest lists deeper than
                  you really should.
         \end{enumerate}
         This is the rest of the second item of the outer
         list.  It is no more interesting than any other
         part of the item.
   \item This is the third item of the list.
\end{itemize}
You can even display poetry.
\begin{verse}
   There is an environment
    for verse \\             % The \\ command separates lines
   Whose features some poets % within a stanza.
   will curse.

                             % One or more blank lines separate stanzas.

   For instead of making\\
   Them do \emph{all} line breaking, \\
   It allows them to put too many words on a line when they'd rather be
   forced to be terse.
\end{verse}

Mathematical formulas may also be displayed.  A
displayed formula
is
one-line long; multiline
formulas require special formatting instructions.
   \[  \ip{\Gamma}{\psi'} = x'' + y^{2} + z_{i}^{n}\]
Don't start a paragraph with a displayed equation,
nor make one a paragraph by itself.

\end{document}               % End of document.

\section{Bench Huge Section 149}
% synthetic bench section 149
% This is a sample LaTeX input file.  (Version of 12 August 2004.)
%
% A '%' character causes TeX to ignore all remaining text on the line,
% and is used for comments like this one.

\documentclass{article}      % Specifies the document class

                             % The preamble begins here.
\title{An Example Document}  % Declares the document's title.
\author{Leslie Lamport}      % Declares the author's name.
\date{January 21, 1994}      % Deleting this command produces today's date.

\newcommand{\ip}[2]{(#1, #2)}
                             % Defines \ip{arg1}{arg2} to mean
                             % (arg1, arg2).

%\newcommand{\ip}[2]{\langle #1 | #2\rangle}
                             % This is an alternative definition of
                             % \ip that is commented out.

\begin{document}             % End of preamble and beginning of text.

\maketitle                   % Produces the title.

This is an example input file.  Comparing it with
the output it generates can show you how to
produce a simple document of your own.

\section{Ordinary Text}      % Produces section heading.  Lower-level
                             % sections are begun with similar
                             % \subsection and \subsubsection commands.

The ends  of words and sentences are marked
  by   spaces. It  doesn't matter how many
spaces    you type; one is as good as 100.  The
end of   a line counts as a space.

One   or more   blank lines denote the  end
of  a paragraph.

Since any number of consecutive spaces are treated
like a single one, the formatting of the input
file makes no difference to
      \LaTeX,                % The \LaTeX command generates the LaTeX logo.
but it makes a difference to you.  When you use
\LaTeX, making your input file as easy to read
as possible will be a great help as you write
your document and when you change it.  This sample
file shows how you can add comments to your own input
file.

Because printing is different from typewriting,
there are a number of things that you have to do
differently when preparing an input file than if
you were just typing the document directly.
Quotation marks like
       ``this''
have to be handled specially, as do quotes within
quotes:
       ``\,`this'            % \, separates the double and single quote.
        is what I just
        wrote, not  `that'\,''.

Dashes come in three sizes: an
       intra-word
dash, a medium dash for number ranges like
       1--2,
and a punctuation
       dash---like
this.

A sentence-ending space should be larger than the
space between words within a sentence.  You
sometimes have to type special commands in
conjunction with punctuation characters to get
this right, as in the following sentence.
       Gnats, gnus, etc.\ all  % `\ ' makes an inter-word space.
       begin with G\@.         % \@ marks end-of-sentence punctuation.
You should check the spaces after periods when
reading your output to make sure you haven't
forgotten any special cases.  Generating an
ellipsis
       \ldots\               % `\ ' is needed after `\ldots' because TeX
                             % ignores spaces after command names like \ldots
                             % made from \ + letters.
                             %
                             % Note how a `%' character causes TeX to ignore
                             % the end of the input line, so these blank lines
                             % do not start a new paragraph.
                             %
with the right spacing around the periods requires
a special command.

\LaTeX\ interprets some common characters as
commands, so you must type special commands to
generate them.  These characters include the
following:
       \$ \& \% \# \{ and \}.

In printing, text is usually emphasized with an
       \emph{italic}
type style.

\begin{em}
   A long segment of text can also be emphasized
   in this way.  Text within such a segment can be
   given \emph{additional} emphasis.
\end{em}

It is sometimes necessary to prevent \LaTeX\ from
breaking a line where it might otherwise do so.
This may be at a space, as between the ``Mr.''\ and
``Jones'' in
       ``Mr.~Jones'',        % ~ produces an unbreakable interword space.
or within a word---especially when the word is a
symbol like
       \mbox{\emph{itemnum}}
that makes little sense when hyphenated across
lines.

Footnotes\footnote{This is an example of a footnote.}
pose no problem.

\LaTeX\ is good at typesetting mathematical formulas
like
       \( x-3y + z = 7 \)
or
       \( a_{1} > x^{2n} + y^{2n} > x' \)
or
       \( \ip{A}{B} = \sum_{i} a_{i} b_{i} \).
The spaces you type in a formula are
ignored.  Remember that a letter like
       $x$                   % $ ... $  and  \( ... \)  are equivalent
is a formula when it denotes a mathematical
symbol, and it should be typed as one.

\section{Displayed Text}

Text is displayed by indenting it from the left
margin.  Quotations are commonly displayed.  There
are short quotations
\begin{quote}
   This is a short quotation.  It consists of a
   single paragraph of text.  See how it is formatted.
\end{quote}
and longer ones.
\begin{quotation}
   This is a longer quotation.  It consists of two
   paragraphs of text, neither of which are
   particularly interesting.

   This is the second paragraph of the quotation.  It
   is just as dull as the first paragraph.
\end{quotation}
Another frequently-displayed structure is a list.
The following is an example of an \emph{itemized}
list.
\begin{itemize}
   \item This is the first item of an itemized list.
         Each item in the list is marked with a ``tick''.
         You don't have to worry about what kind of tick
         mark is used.

   \item This is the second item of the list.  It
         contains another list nested inside it.  The inner
         list is an \emph{enumerated} list.
         \begin{enumerate}
            \item This is the first item of an enumerated
                  list that is nested within the itemized list.

            \item This is the second item of the inner list.
                  \LaTeX\ allows you to nest lists deeper than
                  you really should.
         \end{enumerate}
         This is the rest of the second item of the outer
         list.  It is no more interesting than any other
         part of the item.
   \item This is the third item of the list.
\end{itemize}
You can even display poetry.
\begin{verse}
   There is an environment
    for verse \\             % The \\ command separates lines
   Whose features some poets % within a stanza.
   will curse.

                             % One or more blank lines separate stanzas.

   For instead of making\\
   Them do \emph{all} line breaking, \\
   It allows them to put too many words on a line when they'd rather be
   forced to be terse.
\end{verse}

Mathematical formulas may also be displayed.  A
displayed formula
is
one-line long; multiline
formulas require special formatting instructions.
   \[  \ip{\Gamma}{\psi'} = x'' + y^{2} + z_{i}^{n}\]
Don't start a paragraph with a displayed equation,
nor make one a paragraph by itself.

\end{document}               % End of document.

\section{Bench Huge Section 150}
% synthetic bench section 150
% This is a sample LaTeX input file.  (Version of 12 August 2004.)
%
% A '%' character causes TeX to ignore all remaining text on the line,
% and is used for comments like this one.

\documentclass{article}      % Specifies the document class

                             % The preamble begins here.
\title{An Example Document}  % Declares the document's title.
\author{Leslie Lamport}      % Declares the author's name.
\date{January 21, 1994}      % Deleting this command produces today's date.

\newcommand{\ip}[2]{(#1, #2)}
                             % Defines \ip{arg1}{arg2} to mean
                             % (arg1, arg2).

%\newcommand{\ip}[2]{\langle #1 | #2\rangle}
                             % This is an alternative definition of
                             % \ip that is commented out.

\begin{document}             % End of preamble and beginning of text.

\maketitle                   % Produces the title.

This is an example input file.  Comparing it with
the output it generates can show you how to
produce a simple document of your own.

\section{Ordinary Text}      % Produces section heading.  Lower-level
                             % sections are begun with similar
                             % \subsection and \subsubsection commands.

The ends  of words and sentences are marked
  by   spaces. It  doesn't matter how many
spaces    you type; one is as good as 100.  The
end of   a line counts as a space.

One   or more   blank lines denote the  end
of  a paragraph.

Since any number of consecutive spaces are treated
like a single one, the formatting of the input
file makes no difference to
      \LaTeX,                % The \LaTeX command generates the LaTeX logo.
but it makes a difference to you.  When you use
\LaTeX, making your input file as easy to read
as possible will be a great help as you write
your document and when you change it.  This sample
file shows how you can add comments to your own input
file.

Because printing is different from typewriting,
there are a number of things that you have to do
differently when preparing an input file than if
you were just typing the document directly.
Quotation marks like
       ``this''
have to be handled specially, as do quotes within
quotes:
       ``\,`this'            % \, separates the double and single quote.
        is what I just
        wrote, not  `that'\,''.

Dashes come in three sizes: an
       intra-word
dash, a medium dash for number ranges like
       1--2,
and a punctuation
       dash---like
this.

A sentence-ending space should be larger than the
space between words within a sentence.  You
sometimes have to type special commands in
conjunction with punctuation characters to get
this right, as in the following sentence.
       Gnats, gnus, etc.\ all  % `\ ' makes an inter-word space.
       begin with G\@.         % \@ marks end-of-sentence punctuation.
You should check the spaces after periods when
reading your output to make sure you haven't
forgotten any special cases.  Generating an
ellipsis
       \ldots\               % `\ ' is needed after `\ldots' because TeX
                             % ignores spaces after command names like \ldots
                             % made from \ + letters.
                             %
                             % Note how a `%' character causes TeX to ignore
                             % the end of the input line, so these blank lines
                             % do not start a new paragraph.
                             %
with the right spacing around the periods requires
a special command.

\LaTeX\ interprets some common characters as
commands, so you must type special commands to
generate them.  These characters include the
following:
       \$ \& \% \# \{ and \}.

In printing, text is usually emphasized with an
       \emph{italic}
type style.

\begin{em}
   A long segment of text can also be emphasized
   in this way.  Text within such a segment can be
   given \emph{additional} emphasis.
\end{em}

It is sometimes necessary to prevent \LaTeX\ from
breaking a line where it might otherwise do so.
This may be at a space, as between the ``Mr.''\ and
``Jones'' in
       ``Mr.~Jones'',        % ~ produces an unbreakable interword space.
or within a word---especially when the word is a
symbol like
       \mbox{\emph{itemnum}}
that makes little sense when hyphenated across
lines.

Footnotes\footnote{This is an example of a footnote.}
pose no problem.

\LaTeX\ is good at typesetting mathematical formulas
like
       \( x-3y + z = 7 \)
or
       \( a_{1} > x^{2n} + y^{2n} > x' \)
or
       \( \ip{A}{B} = \sum_{i} a_{i} b_{i} \).
The spaces you type in a formula are
ignored.  Remember that a letter like
       $x$                   % $ ... $  and  \( ... \)  are equivalent
is a formula when it denotes a mathematical
symbol, and it should be typed as one.

\section{Displayed Text}

Text is displayed by indenting it from the left
margin.  Quotations are commonly displayed.  There
are short quotations
\begin{quote}
   This is a short quotation.  It consists of a
   single paragraph of text.  See how it is formatted.
\end{quote}
and longer ones.
\begin{quotation}
   This is a longer quotation.  It consists of two
   paragraphs of text, neither of which are
   particularly interesting.

   This is the second paragraph of the quotation.  It
   is just as dull as the first paragraph.
\end{quotation}
Another frequently-displayed structure is a list.
The following is an example of an \emph{itemized}
list.
\begin{itemize}
   \item This is the first item of an itemized list.
         Each item in the list is marked with a ``tick''.
         You don't have to worry about what kind of tick
         mark is used.

   \item This is the second item of the list.  It
         contains another list nested inside it.  The inner
         list is an \emph{enumerated} list.
         \begin{enumerate}
            \item This is the first item of an enumerated
                  list that is nested within the itemized list.

            \item This is the second item of the inner list.
                  \LaTeX\ allows you to nest lists deeper than
                  you really should.
         \end{enumerate}
         This is the rest of the second item of the outer
         list.  It is no more interesting than any other
         part of the item.
   \item This is the third item of the list.
\end{itemize}
You can even display poetry.
\begin{verse}
   There is an environment
    for verse \\             % The \\ command separates lines
   Whose features some poets % within a stanza.
   will curse.

                             % One or more blank lines separate stanzas.

   For instead of making\\
   Them do \emph{all} line breaking, \\
   It allows them to put too many words on a line when they'd rather be
   forced to be terse.
\end{verse}

Mathematical formulas may also be displayed.  A
displayed formula
is
one-line long; multiline
formulas require special formatting instructions.
   \[  \ip{\Gamma}{\psi'} = x'' + y^{2} + z_{i}^{n}\]
Don't start a paragraph with a displayed equation,
nor make one a paragraph by itself.

\end{document}               % End of document.

\section{Bench Huge Section 151}
% synthetic bench section 151
% This is a sample LaTeX input file.  (Version of 12 August 2004.)
%
% A '%' character causes TeX to ignore all remaining text on the line,
% and is used for comments like this one.

\documentclass{article}      % Specifies the document class

                             % The preamble begins here.
\title{An Example Document}  % Declares the document's title.
\author{Leslie Lamport}      % Declares the author's name.
\date{January 21, 1994}      % Deleting this command produces today's date.

\newcommand{\ip}[2]{(#1, #2)}
                             % Defines \ip{arg1}{arg2} to mean
                             % (arg1, arg2).

%\newcommand{\ip}[2]{\langle #1 | #2\rangle}
                             % This is an alternative definition of
                             % \ip that is commented out.

\begin{document}             % End of preamble and beginning of text.

\maketitle                   % Produces the title.

This is an example input file.  Comparing it with
the output it generates can show you how to
produce a simple document of your own.

\section{Ordinary Text}      % Produces section heading.  Lower-level
                             % sections are begun with similar
                             % \subsection and \subsubsection commands.

The ends  of words and sentences are marked
  by   spaces. It  doesn't matter how many
spaces    you type; one is as good as 100.  The
end of   a line counts as a space.

One   or more   blank lines denote the  end
of  a paragraph.

Since any number of consecutive spaces are treated
like a single one, the formatting of the input
file makes no difference to
      \LaTeX,                % The \LaTeX command generates the LaTeX logo.
but it makes a difference to you.  When you use
\LaTeX, making your input file as easy to read
as possible will be a great help as you write
your document and when you change it.  This sample
file shows how you can add comments to your own input
file.

Because printing is different from typewriting,
there are a number of things that you have to do
differently when preparing an input file than if
you were just typing the document directly.
Quotation marks like
       ``this''
have to be handled specially, as do quotes within
quotes:
       ``\,`this'            % \, separates the double and single quote.
        is what I just
        wrote, not  `that'\,''.

Dashes come in three sizes: an
       intra-word
dash, a medium dash for number ranges like
       1--2,
and a punctuation
       dash---like
this.

A sentence-ending space should be larger than the
space between words within a sentence.  You
sometimes have to type special commands in
conjunction with punctuation characters to get
this right, as in the following sentence.
       Gnats, gnus, etc.\ all  % `\ ' makes an inter-word space.
       begin with G\@.         % \@ marks end-of-sentence punctuation.
You should check the spaces after periods when
reading your output to make sure you haven't
forgotten any special cases.  Generating an
ellipsis
       \ldots\               % `\ ' is needed after `\ldots' because TeX
                             % ignores spaces after command names like \ldots
                             % made from \ + letters.
                             %
                             % Note how a `%' character causes TeX to ignore
                             % the end of the input line, so these blank lines
                             % do not start a new paragraph.
                             %
with the right spacing around the periods requires
a special command.

\LaTeX\ interprets some common characters as
commands, so you must type special commands to
generate them.  These characters include the
following:
       \$ \& \% \# \{ and \}.

In printing, text is usually emphasized with an
       \emph{italic}
type style.

\begin{em}
   A long segment of text can also be emphasized
   in this way.  Text within such a segment can be
   given \emph{additional} emphasis.
\end{em}

It is sometimes necessary to prevent \LaTeX\ from
breaking a line where it might otherwise do so.
This may be at a space, as between the ``Mr.''\ and
``Jones'' in
       ``Mr.~Jones'',        % ~ produces an unbreakable interword space.
or within a word---especially when the word is a
symbol like
       \mbox{\emph{itemnum}}
that makes little sense when hyphenated across
lines.

Footnotes\footnote{This is an example of a footnote.}
pose no problem.

\LaTeX\ is good at typesetting mathematical formulas
like
       \( x-3y + z = 7 \)
or
       \( a_{1} > x^{2n} + y^{2n} > x' \)
or
       \( \ip{A}{B} = \sum_{i} a_{i} b_{i} \).
The spaces you type in a formula are
ignored.  Remember that a letter like
       $x$                   % $ ... $  and  \( ... \)  are equivalent
is a formula when it denotes a mathematical
symbol, and it should be typed as one.

\section{Displayed Text}

Text is displayed by indenting it from the left
margin.  Quotations are commonly displayed.  There
are short quotations
\begin{quote}
   This is a short quotation.  It consists of a
   single paragraph of text.  See how it is formatted.
\end{quote}
and longer ones.
\begin{quotation}
   This is a longer quotation.  It consists of two
   paragraphs of text, neither of which are
   particularly interesting.

   This is the second paragraph of the quotation.  It
   is just as dull as the first paragraph.
\end{quotation}
Another frequently-displayed structure is a list.
The following is an example of an \emph{itemized}
list.
\begin{itemize}
   \item This is the first item of an itemized list.
         Each item in the list is marked with a ``tick''.
         You don't have to worry about what kind of tick
         mark is used.

   \item This is the second item of the list.  It
         contains another list nested inside it.  The inner
         list is an \emph{enumerated} list.
         \begin{enumerate}
            \item This is the first item of an enumerated
                  list that is nested within the itemized list.

            \item This is the second item of the inner list.
                  \LaTeX\ allows you to nest lists deeper than
                  you really should.
         \end{enumerate}
         This is the rest of the second item of the outer
         list.  It is no more interesting than any other
         part of the item.
   \item This is the third item of the list.
\end{itemize}
You can even display poetry.
\begin{verse}
   There is an environment
    for verse \\             % The \\ command separates lines
   Whose features some poets % within a stanza.
   will curse.

                             % One or more blank lines separate stanzas.

   For instead of making\\
   Them do \emph{all} line breaking, \\
   It allows them to put too many words on a line when they'd rather be
   forced to be terse.
\end{verse}

Mathematical formulas may also be displayed.  A
displayed formula
is
one-line long; multiline
formulas require special formatting instructions.
   \[  \ip{\Gamma}{\psi'} = x'' + y^{2} + z_{i}^{n}\]
Don't start a paragraph with a displayed equation,
nor make one a paragraph by itself.

\end{document}               % End of document.

\section{Bench Huge Section 152}
% synthetic bench section 152
% This is a sample LaTeX input file.  (Version of 12 August 2004.)
%
% A '%' character causes TeX to ignore all remaining text on the line,
% and is used for comments like this one.

\documentclass{article}      % Specifies the document class

                             % The preamble begins here.
\title{An Example Document}  % Declares the document's title.
\author{Leslie Lamport}      % Declares the author's name.
\date{January 21, 1994}      % Deleting this command produces today's date.

\newcommand{\ip}[2]{(#1, #2)}
                             % Defines \ip{arg1}{arg2} to mean
                             % (arg1, arg2).

%\newcommand{\ip}[2]{\langle #1 | #2\rangle}
                             % This is an alternative definition of
                             % \ip that is commented out.

\begin{document}             % End of preamble and beginning of text.

\maketitle                   % Produces the title.

This is an example input file.  Comparing it with
the output it generates can show you how to
produce a simple document of your own.

\section{Ordinary Text}      % Produces section heading.  Lower-level
                             % sections are begun with similar
                             % \subsection and \subsubsection commands.

The ends  of words and sentences are marked
  by   spaces. It  doesn't matter how many
spaces    you type; one is as good as 100.  The
end of   a line counts as a space.

One   or more   blank lines denote the  end
of  a paragraph.

Since any number of consecutive spaces are treated
like a single one, the formatting of the input
file makes no difference to
      \LaTeX,                % The \LaTeX command generates the LaTeX logo.
but it makes a difference to you.  When you use
\LaTeX, making your input file as easy to read
as possible will be a great help as you write
your document and when you change it.  This sample
file shows how you can add comments to your own input
file.

Because printing is different from typewriting,
there are a number of things that you have to do
differently when preparing an input file than if
you were just typing the document directly.
Quotation marks like
       ``this''
have to be handled specially, as do quotes within
quotes:
       ``\,`this'            % \, separates the double and single quote.
        is what I just
        wrote, not  `that'\,''.

Dashes come in three sizes: an
       intra-word
dash, a medium dash for number ranges like
       1--2,
and a punctuation
       dash---like
this.

A sentence-ending space should be larger than the
space between words within a sentence.  You
sometimes have to type special commands in
conjunction with punctuation characters to get
this right, as in the following sentence.
       Gnats, gnus, etc.\ all  % `\ ' makes an inter-word space.
       begin with G\@.         % \@ marks end-of-sentence punctuation.
You should check the spaces after periods when
reading your output to make sure you haven't
forgotten any special cases.  Generating an
ellipsis
       \ldots\               % `\ ' is needed after `\ldots' because TeX
                             % ignores spaces after command names like \ldots
                             % made from \ + letters.
                             %
                             % Note how a `%' character causes TeX to ignore
                             % the end of the input line, so these blank lines
                             % do not start a new paragraph.
                             %
with the right spacing around the periods requires
a special command.

\LaTeX\ interprets some common characters as
commands, so you must type special commands to
generate them.  These characters include the
following:
       \$ \& \% \# \{ and \}.

In printing, text is usually emphasized with an
       \emph{italic}
type style.

\begin{em}
   A long segment of text can also be emphasized
   in this way.  Text within such a segment can be
   given \emph{additional} emphasis.
\end{em}

It is sometimes necessary to prevent \LaTeX\ from
breaking a line where it might otherwise do so.
This may be at a space, as between the ``Mr.''\ and
``Jones'' in
       ``Mr.~Jones'',        % ~ produces an unbreakable interword space.
or within a word---especially when the word is a
symbol like
       \mbox{\emph{itemnum}}
that makes little sense when hyphenated across
lines.

Footnotes\footnote{This is an example of a footnote.}
pose no problem.

\LaTeX\ is good at typesetting mathematical formulas
like
       \( x-3y + z = 7 \)
or
       \( a_{1} > x^{2n} + y^{2n} > x' \)
or
       \( \ip{A}{B} = \sum_{i} a_{i} b_{i} \).
The spaces you type in a formula are
ignored.  Remember that a letter like
       $x$                   % $ ... $  and  \( ... \)  are equivalent
is a formula when it denotes a mathematical
symbol, and it should be typed as one.

\section{Displayed Text}

Text is displayed by indenting it from the left
margin.  Quotations are commonly displayed.  There
are short quotations
\begin{quote}
   This is a short quotation.  It consists of a
   single paragraph of text.  See how it is formatted.
\end{quote}
and longer ones.
\begin{quotation}
   This is a longer quotation.  It consists of two
   paragraphs of text, neither of which are
   particularly interesting.

   This is the second paragraph of the quotation.  It
   is just as dull as the first paragraph.
\end{quotation}
Another frequently-displayed structure is a list.
The following is an example of an \emph{itemized}
list.
\begin{itemize}
   \item This is the first item of an itemized list.
         Each item in the list is marked with a ``tick''.
         You don't have to worry about what kind of tick
         mark is used.

   \item This is the second item of the list.  It
         contains another list nested inside it.  The inner
         list is an \emph{enumerated} list.
         \begin{enumerate}
            \item This is the first item of an enumerated
                  list that is nested within the itemized list.

            \item This is the second item of the inner list.
                  \LaTeX\ allows you to nest lists deeper than
                  you really should.
         \end{enumerate}
         This is the rest of the second item of the outer
         list.  It is no more interesting than any other
         part of the item.
   \item This is the third item of the list.
\end{itemize}
You can even display poetry.
\begin{verse}
   There is an environment
    for verse \\             % The \\ command separates lines
   Whose features some poets % within a stanza.
   will curse.

                             % One or more blank lines separate stanzas.

   For instead of making\\
   Them do \emph{all} line breaking, \\
   It allows them to put too many words on a line when they'd rather be
   forced to be terse.
\end{verse}

Mathematical formulas may also be displayed.  A
displayed formula
is
one-line long; multiline
formulas require special formatting instructions.
   \[  \ip{\Gamma}{\psi'} = x'' + y^{2} + z_{i}^{n}\]
Don't start a paragraph with a displayed equation,
nor make one a paragraph by itself.

\end{document}               % End of document.

\section{Bench Huge Section 153}
% synthetic bench section 153
% This is a sample LaTeX input file.  (Version of 12 August 2004.)
%
% A '%' character causes TeX to ignore all remaining text on the line,
% and is used for comments like this one.

\documentclass{article}      % Specifies the document class

                             % The preamble begins here.
\title{An Example Document}  % Declares the document's title.
\author{Leslie Lamport}      % Declares the author's name.
\date{January 21, 1994}      % Deleting this command produces today's date.

\newcommand{\ip}[2]{(#1, #2)}
                             % Defines \ip{arg1}{arg2} to mean
                             % (arg1, arg2).

%\newcommand{\ip}[2]{\langle #1 | #2\rangle}
                             % This is an alternative definition of
                             % \ip that is commented out.

\begin{document}             % End of preamble and beginning of text.

\maketitle                   % Produces the title.

This is an example input file.  Comparing it with
the output it generates can show you how to
produce a simple document of your own.

\section{Ordinary Text}      % Produces section heading.  Lower-level
                             % sections are begun with similar
                             % \subsection and \subsubsection commands.

The ends  of words and sentences are marked
  by   spaces. It  doesn't matter how many
spaces    you type; one is as good as 100.  The
end of   a line counts as a space.

One   or more   blank lines denote the  end
of  a paragraph.

Since any number of consecutive spaces are treated
like a single one, the formatting of the input
file makes no difference to
      \LaTeX,                % The \LaTeX command generates the LaTeX logo.
but it makes a difference to you.  When you use
\LaTeX, making your input file as easy to read
as possible will be a great help as you write
your document and when you change it.  This sample
file shows how you can add comments to your own input
file.

Because printing is different from typewriting,
there are a number of things that you have to do
differently when preparing an input file than if
you were just typing the document directly.
Quotation marks like
       ``this''
have to be handled specially, as do quotes within
quotes:
       ``\,`this'            % \, separates the double and single quote.
        is what I just
        wrote, not  `that'\,''.

Dashes come in three sizes: an
       intra-word
dash, a medium dash for number ranges like
       1--2,
and a punctuation
       dash---like
this.

A sentence-ending space should be larger than the
space between words within a sentence.  You
sometimes have to type special commands in
conjunction with punctuation characters to get
this right, as in the following sentence.
       Gnats, gnus, etc.\ all  % `\ ' makes an inter-word space.
       begin with G\@.         % \@ marks end-of-sentence punctuation.
You should check the spaces after periods when
reading your output to make sure you haven't
forgotten any special cases.  Generating an
ellipsis
       \ldots\               % `\ ' is needed after `\ldots' because TeX
                             % ignores spaces after command names like \ldots
                             % made from \ + letters.
                             %
                             % Note how a `%' character causes TeX to ignore
                             % the end of the input line, so these blank lines
                             % do not start a new paragraph.
                             %
with the right spacing around the periods requires
a special command.

\LaTeX\ interprets some common characters as
commands, so you must type special commands to
generate them.  These characters include the
following:
       \$ \& \% \# \{ and \}.

In printing, text is usually emphasized with an
       \emph{italic}
type style.

\begin{em}
   A long segment of text can also be emphasized
   in this way.  Text within such a segment can be
   given \emph{additional} emphasis.
\end{em}

It is sometimes necessary to prevent \LaTeX\ from
breaking a line where it might otherwise do so.
This may be at a space, as between the ``Mr.''\ and
``Jones'' in
       ``Mr.~Jones'',        % ~ produces an unbreakable interword space.
or within a word---especially when the word is a
symbol like
       \mbox{\emph{itemnum}}
that makes little sense when hyphenated across
lines.

Footnotes\footnote{This is an example of a footnote.}
pose no problem.

\LaTeX\ is good at typesetting mathematical formulas
like
       \( x-3y + z = 7 \)
or
       \( a_{1} > x^{2n} + y^{2n} > x' \)
or
       \( \ip{A}{B} = \sum_{i} a_{i} b_{i} \).
The spaces you type in a formula are
ignored.  Remember that a letter like
       $x$                   % $ ... $  and  \( ... \)  are equivalent
is a formula when it denotes a mathematical
symbol, and it should be typed as one.

\section{Displayed Text}

Text is displayed by indenting it from the left
margin.  Quotations are commonly displayed.  There
are short quotations
\begin{quote}
   This is a short quotation.  It consists of a
   single paragraph of text.  See how it is formatted.
\end{quote}
and longer ones.
\begin{quotation}
   This is a longer quotation.  It consists of two
   paragraphs of text, neither of which are
   particularly interesting.

   This is the second paragraph of the quotation.  It
   is just as dull as the first paragraph.
\end{quotation}
Another frequently-displayed structure is a list.
The following is an example of an \emph{itemized}
list.
\begin{itemize}
   \item This is the first item of an itemized list.
         Each item in the list is marked with a ``tick''.
         You don't have to worry about what kind of tick
         mark is used.

   \item This is the second item of the list.  It
         contains another list nested inside it.  The inner
         list is an \emph{enumerated} list.
         \begin{enumerate}
            \item This is the first item of an enumerated
                  list that is nested within the itemized list.

            \item This is the second item of the inner list.
                  \LaTeX\ allows you to nest lists deeper than
                  you really should.
         \end{enumerate}
         This is the rest of the second item of the outer
         list.  It is no more interesting than any other
         part of the item.
   \item This is the third item of the list.
\end{itemize}
You can even display poetry.
\begin{verse}
   There is an environment
    for verse \\             % The \\ command separates lines
   Whose features some poets % within a stanza.
   will curse.

                             % One or more blank lines separate stanzas.

   For instead of making\\
   Them do \emph{all} line breaking, \\
   It allows them to put too many words on a line when they'd rather be
   forced to be terse.
\end{verse}

Mathematical formulas may also be displayed.  A
displayed formula
is
one-line long; multiline
formulas require special formatting instructions.
   \[  \ip{\Gamma}{\psi'} = x'' + y^{2} + z_{i}^{n}\]
Don't start a paragraph with a displayed equation,
nor make one a paragraph by itself.

\end{document}               % End of document.

\section{Bench Huge Section 154}
% synthetic bench section 154
% This is a sample LaTeX input file.  (Version of 12 August 2004.)
%
% A '%' character causes TeX to ignore all remaining text on the line,
% and is used for comments like this one.

\documentclass{article}      % Specifies the document class

                             % The preamble begins here.
\title{An Example Document}  % Declares the document's title.
\author{Leslie Lamport}      % Declares the author's name.
\date{January 21, 1994}      % Deleting this command produces today's date.

\newcommand{\ip}[2]{(#1, #2)}
                             % Defines \ip{arg1}{arg2} to mean
                             % (arg1, arg2).

%\newcommand{\ip}[2]{\langle #1 | #2\rangle}
                             % This is an alternative definition of
                             % \ip that is commented out.

\begin{document}             % End of preamble and beginning of text.

\maketitle                   % Produces the title.

This is an example input file.  Comparing it with
the output it generates can show you how to
produce a simple document of your own.

\section{Ordinary Text}      % Produces section heading.  Lower-level
                             % sections are begun with similar
                             % \subsection and \subsubsection commands.

The ends  of words and sentences are marked
  by   spaces. It  doesn't matter how many
spaces    you type; one is as good as 100.  The
end of   a line counts as a space.

One   or more   blank lines denote the  end
of  a paragraph.

Since any number of consecutive spaces are treated
like a single one, the formatting of the input
file makes no difference to
      \LaTeX,                % The \LaTeX command generates the LaTeX logo.
but it makes a difference to you.  When you use
\LaTeX, making your input file as easy to read
as possible will be a great help as you write
your document and when you change it.  This sample
file shows how you can add comments to your own input
file.

Because printing is different from typewriting,
there are a number of things that you have to do
differently when preparing an input file than if
you were just typing the document directly.
Quotation marks like
       ``this''
have to be handled specially, as do quotes within
quotes:
       ``\,`this'            % \, separates the double and single quote.
        is what I just
        wrote, not  `that'\,''.

Dashes come in three sizes: an
       intra-word
dash, a medium dash for number ranges like
       1--2,
and a punctuation
       dash---like
this.

A sentence-ending space should be larger than the
space between words within a sentence.  You
sometimes have to type special commands in
conjunction with punctuation characters to get
this right, as in the following sentence.
       Gnats, gnus, etc.\ all  % `\ ' makes an inter-word space.
       begin with G\@.         % \@ marks end-of-sentence punctuation.
You should check the spaces after periods when
reading your output to make sure you haven't
forgotten any special cases.  Generating an
ellipsis
       \ldots\               % `\ ' is needed after `\ldots' because TeX
                             % ignores spaces after command names like \ldots
                             % made from \ + letters.
                             %
                             % Note how a `%' character causes TeX to ignore
                             % the end of the input line, so these blank lines
                             % do not start a new paragraph.
                             %
with the right spacing around the periods requires
a special command.

\LaTeX\ interprets some common characters as
commands, so you must type special commands to
generate them.  These characters include the
following:
       \$ \& \% \# \{ and \}.

In printing, text is usually emphasized with an
       \emph{italic}
type style.

\begin{em}
   A long segment of text can also be emphasized
   in this way.  Text within such a segment can be
   given \emph{additional} emphasis.
\end{em}

It is sometimes necessary to prevent \LaTeX\ from
breaking a line where it might otherwise do so.
This may be at a space, as between the ``Mr.''\ and
``Jones'' in
       ``Mr.~Jones'',        % ~ produces an unbreakable interword space.
or within a word---especially when the word is a
symbol like
       \mbox{\emph{itemnum}}
that makes little sense when hyphenated across
lines.

Footnotes\footnote{This is an example of a footnote.}
pose no problem.

\LaTeX\ is good at typesetting mathematical formulas
like
       \( x-3y + z = 7 \)
or
       \( a_{1} > x^{2n} + y^{2n} > x' \)
or
       \( \ip{A}{B} = \sum_{i} a_{i} b_{i} \).
The spaces you type in a formula are
ignored.  Remember that a letter like
       $x$                   % $ ... $  and  \( ... \)  are equivalent
is a formula when it denotes a mathematical
symbol, and it should be typed as one.

\section{Displayed Text}

Text is displayed by indenting it from the left
margin.  Quotations are commonly displayed.  There
are short quotations
\begin{quote}
   This is a short quotation.  It consists of a
   single paragraph of text.  See how it is formatted.
\end{quote}
and longer ones.
\begin{quotation}
   This is a longer quotation.  It consists of two
   paragraphs of text, neither of which are
   particularly interesting.

   This is the second paragraph of the quotation.  It
   is just as dull as the first paragraph.
\end{quotation}
Another frequently-displayed structure is a list.
The following is an example of an \emph{itemized}
list.
\begin{itemize}
   \item This is the first item of an itemized list.
         Each item in the list is marked with a ``tick''.
         You don't have to worry about what kind of tick
         mark is used.

   \item This is the second item of the list.  It
         contains another list nested inside it.  The inner
         list is an \emph{enumerated} list.
         \begin{enumerate}
            \item This is the first item of an enumerated
                  list that is nested within the itemized list.

            \item This is the second item of the inner list.
                  \LaTeX\ allows you to nest lists deeper than
                  you really should.
         \end{enumerate}
         This is the rest of the second item of the outer
         list.  It is no more interesting than any other
         part of the item.
   \item This is the third item of the list.
\end{itemize}
You can even display poetry.
\begin{verse}
   There is an environment
    for verse \\             % The \\ command separates lines
   Whose features some poets % within a stanza.
   will curse.

                             % One or more blank lines separate stanzas.

   For instead of making\\
   Them do \emph{all} line breaking, \\
   It allows them to put too many words on a line when they'd rather be
   forced to be terse.
\end{verse}

Mathematical formulas may also be displayed.  A
displayed formula
is
one-line long; multiline
formulas require special formatting instructions.
   \[  \ip{\Gamma}{\psi'} = x'' + y^{2} + z_{i}^{n}\]
Don't start a paragraph with a displayed equation,
nor make one a paragraph by itself.

\end{document}               % End of document.

\section{Bench Huge Section 155}
% synthetic bench section 155
% This is a sample LaTeX input file.  (Version of 12 August 2004.)
%
% A '%' character causes TeX to ignore all remaining text on the line,
% and is used for comments like this one.

\documentclass{article}      % Specifies the document class

                             % The preamble begins here.
\title{An Example Document}  % Declares the document's title.
\author{Leslie Lamport}      % Declares the author's name.
\date{January 21, 1994}      % Deleting this command produces today's date.

\newcommand{\ip}[2]{(#1, #2)}
                             % Defines \ip{arg1}{arg2} to mean
                             % (arg1, arg2).

%\newcommand{\ip}[2]{\langle #1 | #2\rangle}
                             % This is an alternative definition of
                             % \ip that is commented out.

\begin{document}             % End of preamble and beginning of text.

\maketitle                   % Produces the title.

This is an example input file.  Comparing it with
the output it generates can show you how to
produce a simple document of your own.

\section{Ordinary Text}      % Produces section heading.  Lower-level
                             % sections are begun with similar
                             % \subsection and \subsubsection commands.

The ends  of words and sentences are marked
  by   spaces. It  doesn't matter how many
spaces    you type; one is as good as 100.  The
end of   a line counts as a space.

One   or more   blank lines denote the  end
of  a paragraph.

Since any number of consecutive spaces are treated
like a single one, the formatting of the input
file makes no difference to
      \LaTeX,                % The \LaTeX command generates the LaTeX logo.
but it makes a difference to you.  When you use
\LaTeX, making your input file as easy to read
as possible will be a great help as you write
your document and when you change it.  This sample
file shows how you can add comments to your own input
file.

Because printing is different from typewriting,
there are a number of things that you have to do
differently when preparing an input file than if
you were just typing the document directly.
Quotation marks like
       ``this''
have to be handled specially, as do quotes within
quotes:
       ``\,`this'            % \, separates the double and single quote.
        is what I just
        wrote, not  `that'\,''.

Dashes come in three sizes: an
       intra-word
dash, a medium dash for number ranges like
       1--2,
and a punctuation
       dash---like
this.

A sentence-ending space should be larger than the
space between words within a sentence.  You
sometimes have to type special commands in
conjunction with punctuation characters to get
this right, as in the following sentence.
       Gnats, gnus, etc.\ all  % `\ ' makes an inter-word space.
       begin with G\@.         % \@ marks end-of-sentence punctuation.
You should check the spaces after periods when
reading your output to make sure you haven't
forgotten any special cases.  Generating an
ellipsis
       \ldots\               % `\ ' is needed after `\ldots' because TeX
                             % ignores spaces after command names like \ldots
                             % made from \ + letters.
                             %
                             % Note how a `%' character causes TeX to ignore
                             % the end of the input line, so these blank lines
                             % do not start a new paragraph.
                             %
with the right spacing around the periods requires
a special command.

\LaTeX\ interprets some common characters as
commands, so you must type special commands to
generate them.  These characters include the
following:
       \$ \& \% \# \{ and \}.

In printing, text is usually emphasized with an
       \emph{italic}
type style.

\begin{em}
   A long segment of text can also be emphasized
   in this way.  Text within such a segment can be
   given \emph{additional} emphasis.
\end{em}

It is sometimes necessary to prevent \LaTeX\ from
breaking a line where it might otherwise do so.
This may be at a space, as between the ``Mr.''\ and
``Jones'' in
       ``Mr.~Jones'',        % ~ produces an unbreakable interword space.
or within a word---especially when the word is a
symbol like
       \mbox{\emph{itemnum}}
that makes little sense when hyphenated across
lines.

Footnotes\footnote{This is an example of a footnote.}
pose no problem.

\LaTeX\ is good at typesetting mathematical formulas
like
       \( x-3y + z = 7 \)
or
       \( a_{1} > x^{2n} + y^{2n} > x' \)
or
       \( \ip{A}{B} = \sum_{i} a_{i} b_{i} \).
The spaces you type in a formula are
ignored.  Remember that a letter like
       $x$                   % $ ... $  and  \( ... \)  are equivalent
is a formula when it denotes a mathematical
symbol, and it should be typed as one.

\section{Displayed Text}

Text is displayed by indenting it from the left
margin.  Quotations are commonly displayed.  There
are short quotations
\begin{quote}
   This is a short quotation.  It consists of a
   single paragraph of text.  See how it is formatted.
\end{quote}
and longer ones.
\begin{quotation}
   This is a longer quotation.  It consists of two
   paragraphs of text, neither of which are
   particularly interesting.

   This is the second paragraph of the quotation.  It
   is just as dull as the first paragraph.
\end{quotation}
Another frequently-displayed structure is a list.
The following is an example of an \emph{itemized}
list.
\begin{itemize}
   \item This is the first item of an itemized list.
         Each item in the list is marked with a ``tick''.
         You don't have to worry about what kind of tick
         mark is used.

   \item This is the second item of the list.  It
         contains another list nested inside it.  The inner
         list is an \emph{enumerated} list.
         \begin{enumerate}
            \item This is the first item of an enumerated
                  list that is nested within the itemized list.

            \item This is the second item of the inner list.
                  \LaTeX\ allows you to nest lists deeper than
                  you really should.
         \end{enumerate}
         This is the rest of the second item of the outer
         list.  It is no more interesting than any other
         part of the item.
   \item This is the third item of the list.
\end{itemize}
You can even display poetry.
\begin{verse}
   There is an environment
    for verse \\             % The \\ command separates lines
   Whose features some poets % within a stanza.
   will curse.

                             % One or more blank lines separate stanzas.

   For instead of making\\
   Them do \emph{all} line breaking, \\
   It allows them to put too many words on a line when they'd rather be
   forced to be terse.
\end{verse}

Mathematical formulas may also be displayed.  A
displayed formula
is
one-line long; multiline
formulas require special formatting instructions.
   \[  \ip{\Gamma}{\psi'} = x'' + y^{2} + z_{i}^{n}\]
Don't start a paragraph with a displayed equation,
nor make one a paragraph by itself.

\end{document}               % End of document.

\section{Bench Huge Section 156}
% synthetic bench section 156
% This is a sample LaTeX input file.  (Version of 12 August 2004.)
%
% A '%' character causes TeX to ignore all remaining text on the line,
% and is used for comments like this one.

\documentclass{article}      % Specifies the document class

                             % The preamble begins here.
\title{An Example Document}  % Declares the document's title.
\author{Leslie Lamport}      % Declares the author's name.
\date{January 21, 1994}      % Deleting this command produces today's date.

\newcommand{\ip}[2]{(#1, #2)}
                             % Defines \ip{arg1}{arg2} to mean
                             % (arg1, arg2).

%\newcommand{\ip}[2]{\langle #1 | #2\rangle}
                             % This is an alternative definition of
                             % \ip that is commented out.

\begin{document}             % End of preamble and beginning of text.

\maketitle                   % Produces the title.

This is an example input file.  Comparing it with
the output it generates can show you how to
produce a simple document of your own.

\section{Ordinary Text}      % Produces section heading.  Lower-level
                             % sections are begun with similar
                             % \subsection and \subsubsection commands.

The ends  of words and sentences are marked
  by   spaces. It  doesn't matter how many
spaces    you type; one is as good as 100.  The
end of   a line counts as a space.

One   or more   blank lines denote the  end
of  a paragraph.

Since any number of consecutive spaces are treated
like a single one, the formatting of the input
file makes no difference to
      \LaTeX,                % The \LaTeX command generates the LaTeX logo.
but it makes a difference to you.  When you use
\LaTeX, making your input file as easy to read
as possible will be a great help as you write
your document and when you change it.  This sample
file shows how you can add comments to your own input
file.

Because printing is different from typewriting,
there are a number of things that you have to do
differently when preparing an input file than if
you were just typing the document directly.
Quotation marks like
       ``this''
have to be handled specially, as do quotes within
quotes:
       ``\,`this'            % \, separates the double and single quote.
        is what I just
        wrote, not  `that'\,''.

Dashes come in three sizes: an
       intra-word
dash, a medium dash for number ranges like
       1--2,
and a punctuation
       dash---like
this.

A sentence-ending space should be larger than the
space between words within a sentence.  You
sometimes have to type special commands in
conjunction with punctuation characters to get
this right, as in the following sentence.
       Gnats, gnus, etc.\ all  % `\ ' makes an inter-word space.
       begin with G\@.         % \@ marks end-of-sentence punctuation.
You should check the spaces after periods when
reading your output to make sure you haven't
forgotten any special cases.  Generating an
ellipsis
       \ldots\               % `\ ' is needed after `\ldots' because TeX
                             % ignores spaces after command names like \ldots
                             % made from \ + letters.
                             %
                             % Note how a `%' character causes TeX to ignore
                             % the end of the input line, so these blank lines
                             % do not start a new paragraph.
                             %
with the right spacing around the periods requires
a special command.

\LaTeX\ interprets some common characters as
commands, so you must type special commands to
generate them.  These characters include the
following:
       \$ \& \% \# \{ and \}.

In printing, text is usually emphasized with an
       \emph{italic}
type style.

\begin{em}
   A long segment of text can also be emphasized
   in this way.  Text within such a segment can be
   given \emph{additional} emphasis.
\end{em}

It is sometimes necessary to prevent \LaTeX\ from
breaking a line where it might otherwise do so.
This may be at a space, as between the ``Mr.''\ and
``Jones'' in
       ``Mr.~Jones'',        % ~ produces an unbreakable interword space.
or within a word---especially when the word is a
symbol like
       \mbox{\emph{itemnum}}
that makes little sense when hyphenated across
lines.

Footnotes\footnote{This is an example of a footnote.}
pose no problem.

\LaTeX\ is good at typesetting mathematical formulas
like
       \( x-3y + z = 7 \)
or
       \( a_{1} > x^{2n} + y^{2n} > x' \)
or
       \( \ip{A}{B} = \sum_{i} a_{i} b_{i} \).
The spaces you type in a formula are
ignored.  Remember that a letter like
       $x$                   % $ ... $  and  \( ... \)  are equivalent
is a formula when it denotes a mathematical
symbol, and it should be typed as one.

\section{Displayed Text}

Text is displayed by indenting it from the left
margin.  Quotations are commonly displayed.  There
are short quotations
\begin{quote}
   This is a short quotation.  It consists of a
   single paragraph of text.  See how it is formatted.
\end{quote}
and longer ones.
\begin{quotation}
   This is a longer quotation.  It consists of two
   paragraphs of text, neither of which are
   particularly interesting.

   This is the second paragraph of the quotation.  It
   is just as dull as the first paragraph.
\end{quotation}
Another frequently-displayed structure is a list.
The following is an example of an \emph{itemized}
list.
\begin{itemize}
   \item This is the first item of an itemized list.
         Each item in the list is marked with a ``tick''.
         You don't have to worry about what kind of tick
         mark is used.

   \item This is the second item of the list.  It
         contains another list nested inside it.  The inner
         list is an \emph{enumerated} list.
         \begin{enumerate}
            \item This is the first item of an enumerated
                  list that is nested within the itemized list.

            \item This is the second item of the inner list.
                  \LaTeX\ allows you to nest lists deeper than
                  you really should.
         \end{enumerate}
         This is the rest of the second item of the outer
         list.  It is no more interesting than any other
         part of the item.
   \item This is the third item of the list.
\end{itemize}
You can even display poetry.
\begin{verse}
   There is an environment
    for verse \\             % The \\ command separates lines
   Whose features some poets % within a stanza.
   will curse.

                             % One or more blank lines separate stanzas.

   For instead of making\\
   Them do \emph{all} line breaking, \\
   It allows them to put too many words on a line when they'd rather be
   forced to be terse.
\end{verse}

Mathematical formulas may also be displayed.  A
displayed formula
is
one-line long; multiline
formulas require special formatting instructions.
   \[  \ip{\Gamma}{\psi'} = x'' + y^{2} + z_{i}^{n}\]
Don't start a paragraph with a displayed equation,
nor make one a paragraph by itself.

\end{document}               % End of document.

\section{Bench Huge Section 157}
% synthetic bench section 157
% This is a sample LaTeX input file.  (Version of 12 August 2004.)
%
% A '%' character causes TeX to ignore all remaining text on the line,
% and is used for comments like this one.

\documentclass{article}      % Specifies the document class

                             % The preamble begins here.
\title{An Example Document}  % Declares the document's title.
\author{Leslie Lamport}      % Declares the author's name.
\date{January 21, 1994}      % Deleting this command produces today's date.

\newcommand{\ip}[2]{(#1, #2)}
                             % Defines \ip{arg1}{arg2} to mean
                             % (arg1, arg2).

%\newcommand{\ip}[2]{\langle #1 | #2\rangle}
                             % This is an alternative definition of
                             % \ip that is commented out.

\begin{document}             % End of preamble and beginning of text.

\maketitle                   % Produces the title.

This is an example input file.  Comparing it with
the output it generates can show you how to
produce a simple document of your own.

\section{Ordinary Text}      % Produces section heading.  Lower-level
                             % sections are begun with similar
                             % \subsection and \subsubsection commands.

The ends  of words and sentences are marked
  by   spaces. It  doesn't matter how many
spaces    you type; one is as good as 100.  The
end of   a line counts as a space.

One   or more   blank lines denote the  end
of  a paragraph.

Since any number of consecutive spaces are treated
like a single one, the formatting of the input
file makes no difference to
      \LaTeX,                % The \LaTeX command generates the LaTeX logo.
but it makes a difference to you.  When you use
\LaTeX, making your input file as easy to read
as possible will be a great help as you write
your document and when you change it.  This sample
file shows how you can add comments to your own input
file.

Because printing is different from typewriting,
there are a number of things that you have to do
differently when preparing an input file than if
you were just typing the document directly.
Quotation marks like
       ``this''
have to be handled specially, as do quotes within
quotes:
       ``\,`this'            % \, separates the double and single quote.
        is what I just
        wrote, not  `that'\,''.

Dashes come in three sizes: an
       intra-word
dash, a medium dash for number ranges like
       1--2,
and a punctuation
       dash---like
this.

A sentence-ending space should be larger than the
space between words within a sentence.  You
sometimes have to type special commands in
conjunction with punctuation characters to get
this right, as in the following sentence.
       Gnats, gnus, etc.\ all  % `\ ' makes an inter-word space.
       begin with G\@.         % \@ marks end-of-sentence punctuation.
You should check the spaces after periods when
reading your output to make sure you haven't
forgotten any special cases.  Generating an
ellipsis
       \ldots\               % `\ ' is needed after `\ldots' because TeX
                             % ignores spaces after command names like \ldots
                             % made from \ + letters.
                             %
                             % Note how a `%' character causes TeX to ignore
                             % the end of the input line, so these blank lines
                             % do not start a new paragraph.
                             %
with the right spacing around the periods requires
a special command.

\LaTeX\ interprets some common characters as
commands, so you must type special commands to
generate them.  These characters include the
following:
       \$ \& \% \# \{ and \}.

In printing, text is usually emphasized with an
       \emph{italic}
type style.

\begin{em}
   A long segment of text can also be emphasized
   in this way.  Text within such a segment can be
   given \emph{additional} emphasis.
\end{em}

It is sometimes necessary to prevent \LaTeX\ from
breaking a line where it might otherwise do so.
This may be at a space, as between the ``Mr.''\ and
``Jones'' in
       ``Mr.~Jones'',        % ~ produces an unbreakable interword space.
or within a word---especially when the word is a
symbol like
       \mbox{\emph{itemnum}}
that makes little sense when hyphenated across
lines.

Footnotes\footnote{This is an example of a footnote.}
pose no problem.

\LaTeX\ is good at typesetting mathematical formulas
like
       \( x-3y + z = 7 \)
or
       \( a_{1} > x^{2n} + y^{2n} > x' \)
or
       \( \ip{A}{B} = \sum_{i} a_{i} b_{i} \).
The spaces you type in a formula are
ignored.  Remember that a letter like
       $x$                   % $ ... $  and  \( ... \)  are equivalent
is a formula when it denotes a mathematical
symbol, and it should be typed as one.

\section{Displayed Text}

Text is displayed by indenting it from the left
margin.  Quotations are commonly displayed.  There
are short quotations
\begin{quote}
   This is a short quotation.  It consists of a
   single paragraph of text.  See how it is formatted.
\end{quote}
and longer ones.
\begin{quotation}
   This is a longer quotation.  It consists of two
   paragraphs of text, neither of which are
   particularly interesting.

   This is the second paragraph of the quotation.  It
   is just as dull as the first paragraph.
\end{quotation}
Another frequently-displayed structure is a list.
The following is an example of an \emph{itemized}
list.
\begin{itemize}
   \item This is the first item of an itemized list.
         Each item in the list is marked with a ``tick''.
         You don't have to worry about what kind of tick
         mark is used.

   \item This is the second item of the list.  It
         contains another list nested inside it.  The inner
         list is an \emph{enumerated} list.
         \begin{enumerate}
            \item This is the first item of an enumerated
                  list that is nested within the itemized list.

            \item This is the second item of the inner list.
                  \LaTeX\ allows you to nest lists deeper than
                  you really should.
         \end{enumerate}
         This is the rest of the second item of the outer
         list.  It is no more interesting than any other
         part of the item.
   \item This is the third item of the list.
\end{itemize}
You can even display poetry.
\begin{verse}
   There is an environment
    for verse \\             % The \\ command separates lines
   Whose features some poets % within a stanza.
   will curse.

                             % One or more blank lines separate stanzas.

   For instead of making\\
   Them do \emph{all} line breaking, \\
   It allows them to put too many words on a line when they'd rather be
   forced to be terse.
\end{verse}

Mathematical formulas may also be displayed.  A
displayed formula
is
one-line long; multiline
formulas require special formatting instructions.
   \[  \ip{\Gamma}{\psi'} = x'' + y^{2} + z_{i}^{n}\]
Don't start a paragraph with a displayed equation,
nor make one a paragraph by itself.

\end{document}               % End of document.

\section{Bench Huge Section 158}
% synthetic bench section 158
% This is a sample LaTeX input file.  (Version of 12 August 2004.)
%
% A '%' character causes TeX to ignore all remaining text on the line,
% and is used for comments like this one.

\documentclass{article}      % Specifies the document class

                             % The preamble begins here.
\title{An Example Document}  % Declares the document's title.
\author{Leslie Lamport}      % Declares the author's name.
\date{January 21, 1994}      % Deleting this command produces today's date.

\newcommand{\ip}[2]{(#1, #2)}
                             % Defines \ip{arg1}{arg2} to mean
                             % (arg1, arg2).

%\newcommand{\ip}[2]{\langle #1 | #2\rangle}
                             % This is an alternative definition of
                             % \ip that is commented out.

\begin{document}             % End of preamble and beginning of text.

\maketitle                   % Produces the title.

This is an example input file.  Comparing it with
the output it generates can show you how to
produce a simple document of your own.

\section{Ordinary Text}      % Produces section heading.  Lower-level
                             % sections are begun with similar
                             % \subsection and \subsubsection commands.

The ends  of words and sentences are marked
  by   spaces. It  doesn't matter how many
spaces    you type; one is as good as 100.  The
end of   a line counts as a space.

One   or more   blank lines denote the  end
of  a paragraph.

Since any number of consecutive spaces are treated
like a single one, the formatting of the input
file makes no difference to
      \LaTeX,                % The \LaTeX command generates the LaTeX logo.
but it makes a difference to you.  When you use
\LaTeX, making your input file as easy to read
as possible will be a great help as you write
your document and when you change it.  This sample
file shows how you can add comments to your own input
file.

Because printing is different from typewriting,
there are a number of things that you have to do
differently when preparing an input file than if
you were just typing the document directly.
Quotation marks like
       ``this''
have to be handled specially, as do quotes within
quotes:
       ``\,`this'            % \, separates the double and single quote.
        is what I just
        wrote, not  `that'\,''.

Dashes come in three sizes: an
       intra-word
dash, a medium dash for number ranges like
       1--2,
and a punctuation
       dash---like
this.

A sentence-ending space should be larger than the
space between words within a sentence.  You
sometimes have to type special commands in
conjunction with punctuation characters to get
this right, as in the following sentence.
       Gnats, gnus, etc.\ all  % `\ ' makes an inter-word space.
       begin with G\@.         % \@ marks end-of-sentence punctuation.
You should check the spaces after periods when
reading your output to make sure you haven't
forgotten any special cases.  Generating an
ellipsis
       \ldots\               % `\ ' is needed after `\ldots' because TeX
                             % ignores spaces after command names like \ldots
                             % made from \ + letters.
                             %
                             % Note how a `%' character causes TeX to ignore
                             % the end of the input line, so these blank lines
                             % do not start a new paragraph.
                             %
with the right spacing around the periods requires
a special command.

\LaTeX\ interprets some common characters as
commands, so you must type special commands to
generate them.  These characters include the
following:
       \$ \& \% \# \{ and \}.

In printing, text is usually emphasized with an
       \emph{italic}
type style.

\begin{em}
   A long segment of text can also be emphasized
   in this way.  Text within such a segment can be
   given \emph{additional} emphasis.
\end{em}

It is sometimes necessary to prevent \LaTeX\ from
breaking a line where it might otherwise do so.
This may be at a space, as between the ``Mr.''\ and
``Jones'' in
       ``Mr.~Jones'',        % ~ produces an unbreakable interword space.
or within a word---especially when the word is a
symbol like
       \mbox{\emph{itemnum}}
that makes little sense when hyphenated across
lines.

Footnotes\footnote{This is an example of a footnote.}
pose no problem.

\LaTeX\ is good at typesetting mathematical formulas
like
       \( x-3y + z = 7 \)
or
       \( a_{1} > x^{2n} + y^{2n} > x' \)
or
       \( \ip{A}{B} = \sum_{i} a_{i} b_{i} \).
The spaces you type in a formula are
ignored.  Remember that a letter like
       $x$                   % $ ... $  and  \( ... \)  are equivalent
is a formula when it denotes a mathematical
symbol, and it should be typed as one.

\section{Displayed Text}

Text is displayed by indenting it from the left
margin.  Quotations are commonly displayed.  There
are short quotations
\begin{quote}
   This is a short quotation.  It consists of a
   single paragraph of text.  See how it is formatted.
\end{quote}
and longer ones.
\begin{quotation}
   This is a longer quotation.  It consists of two
   paragraphs of text, neither of which are
   particularly interesting.

   This is the second paragraph of the quotation.  It
   is just as dull as the first paragraph.
\end{quotation}
Another frequently-displayed structure is a list.
The following is an example of an \emph{itemized}
list.
\begin{itemize}
   \item This is the first item of an itemized list.
         Each item in the list is marked with a ``tick''.
         You don't have to worry about what kind of tick
         mark is used.

   \item This is the second item of the list.  It
         contains another list nested inside it.  The inner
         list is an \emph{enumerated} list.
         \begin{enumerate}
            \item This is the first item of an enumerated
                  list that is nested within the itemized list.

            \item This is the second item of the inner list.
                  \LaTeX\ allows you to nest lists deeper than
                  you really should.
         \end{enumerate}
         This is the rest of the second item of the outer
         list.  It is no more interesting than any other
         part of the item.
   \item This is the third item of the list.
\end{itemize}
You can even display poetry.
\begin{verse}
   There is an environment
    for verse \\             % The \\ command separates lines
   Whose features some poets % within a stanza.
   will curse.

                             % One or more blank lines separate stanzas.

   For instead of making\\
   Them do \emph{all} line breaking, \\
   It allows them to put too many words on a line when they'd rather be
   forced to be terse.
\end{verse}

Mathematical formulas may also be displayed.  A
displayed formula
is
one-line long; multiline
formulas require special formatting instructions.
   \[  \ip{\Gamma}{\psi'} = x'' + y^{2} + z_{i}^{n}\]
Don't start a paragraph with a displayed equation,
nor make one a paragraph by itself.

\end{document}               % End of document.

\section{Bench Huge Section 159}
% synthetic bench section 159
% This is a sample LaTeX input file.  (Version of 12 August 2004.)
%
% A '%' character causes TeX to ignore all remaining text on the line,
% and is used for comments like this one.

\documentclass{article}      % Specifies the document class

                             % The preamble begins here.
\title{An Example Document}  % Declares the document's title.
\author{Leslie Lamport}      % Declares the author's name.
\date{January 21, 1994}      % Deleting this command produces today's date.

\newcommand{\ip}[2]{(#1, #2)}
                             % Defines \ip{arg1}{arg2} to mean
                             % (arg1, arg2).

%\newcommand{\ip}[2]{\langle #1 | #2\rangle}
                             % This is an alternative definition of
                             % \ip that is commented out.

\begin{document}             % End of preamble and beginning of text.

\maketitle                   % Produces the title.

This is an example input file.  Comparing it with
the output it generates can show you how to
produce a simple document of your own.

\section{Ordinary Text}      % Produces section heading.  Lower-level
                             % sections are begun with similar
                             % \subsection and \subsubsection commands.

The ends  of words and sentences are marked
  by   spaces. It  doesn't matter how many
spaces    you type; one is as good as 100.  The
end of   a line counts as a space.

One   or more   blank lines denote the  end
of  a paragraph.

Since any number of consecutive spaces are treated
like a single one, the formatting of the input
file makes no difference to
      \LaTeX,                % The \LaTeX command generates the LaTeX logo.
but it makes a difference to you.  When you use
\LaTeX, making your input file as easy to read
as possible will be a great help as you write
your document and when you change it.  This sample
file shows how you can add comments to your own input
file.

Because printing is different from typewriting,
there are a number of things that you have to do
differently when preparing an input file than if
you were just typing the document directly.
Quotation marks like
       ``this''
have to be handled specially, as do quotes within
quotes:
       ``\,`this'            % \, separates the double and single quote.
        is what I just
        wrote, not  `that'\,''.

Dashes come in three sizes: an
       intra-word
dash, a medium dash for number ranges like
       1--2,
and a punctuation
       dash---like
this.

A sentence-ending space should be larger than the
space between words within a sentence.  You
sometimes have to type special commands in
conjunction with punctuation characters to get
this right, as in the following sentence.
       Gnats, gnus, etc.\ all  % `\ ' makes an inter-word space.
       begin with G\@.         % \@ marks end-of-sentence punctuation.
You should check the spaces after periods when
reading your output to make sure you haven't
forgotten any special cases.  Generating an
ellipsis
       \ldots\               % `\ ' is needed after `\ldots' because TeX
                             % ignores spaces after command names like \ldots
                             % made from \ + letters.
                             %
                             % Note how a `%' character causes TeX to ignore
                             % the end of the input line, so these blank lines
                             % do not start a new paragraph.
                             %
with the right spacing around the periods requires
a special command.

\LaTeX\ interprets some common characters as
commands, so you must type special commands to
generate them.  These characters include the
following:
       \$ \& \% \# \{ and \}.

In printing, text is usually emphasized with an
       \emph{italic}
type style.

\begin{em}
   A long segment of text can also be emphasized
   in this way.  Text within such a segment can be
   given \emph{additional} emphasis.
\end{em}

It is sometimes necessary to prevent \LaTeX\ from
breaking a line where it might otherwise do so.
This may be at a space, as between the ``Mr.''\ and
``Jones'' in
       ``Mr.~Jones'',        % ~ produces an unbreakable interword space.
or within a word---especially when the word is a
symbol like
       \mbox{\emph{itemnum}}
that makes little sense when hyphenated across
lines.

Footnotes\footnote{This is an example of a footnote.}
pose no problem.

\LaTeX\ is good at typesetting mathematical formulas
like
       \( x-3y + z = 7 \)
or
       \( a_{1} > x^{2n} + y^{2n} > x' \)
or
       \( \ip{A}{B} = \sum_{i} a_{i} b_{i} \).
The spaces you type in a formula are
ignored.  Remember that a letter like
       $x$                   % $ ... $  and  \( ... \)  are equivalent
is a formula when it denotes a mathematical
symbol, and it should be typed as one.

\section{Displayed Text}

Text is displayed by indenting it from the left
margin.  Quotations are commonly displayed.  There
are short quotations
\begin{quote}
   This is a short quotation.  It consists of a
   single paragraph of text.  See how it is formatted.
\end{quote}
and longer ones.
\begin{quotation}
   This is a longer quotation.  It consists of two
   paragraphs of text, neither of which are
   particularly interesting.

   This is the second paragraph of the quotation.  It
   is just as dull as the first paragraph.
\end{quotation}
Another frequently-displayed structure is a list.
The following is an example of an \emph{itemized}
list.
\begin{itemize}
   \item This is the first item of an itemized list.
         Each item in the list is marked with a ``tick''.
         You don't have to worry about what kind of tick
         mark is used.

   \item This is the second item of the list.  It
         contains another list nested inside it.  The inner
         list is an \emph{enumerated} list.
         \begin{enumerate}
            \item This is the first item of an enumerated
                  list that is nested within the itemized list.

            \item This is the second item of the inner list.
                  \LaTeX\ allows you to nest lists deeper than
                  you really should.
         \end{enumerate}
         This is the rest of the second item of the outer
         list.  It is no more interesting than any other
         part of the item.
   \item This is the third item of the list.
\end{itemize}
You can even display poetry.
\begin{verse}
   There is an environment
    for verse \\             % The \\ command separates lines
   Whose features some poets % within a stanza.
   will curse.

                             % One or more blank lines separate stanzas.

   For instead of making\\
   Them do \emph{all} line breaking, \\
   It allows them to put too many words on a line when they'd rather be
   forced to be terse.
\end{verse}

Mathematical formulas may also be displayed.  A
displayed formula
is
one-line long; multiline
formulas require special formatting instructions.
   \[  \ip{\Gamma}{\psi'} = x'' + y^{2} + z_{i}^{n}\]
Don't start a paragraph with a displayed equation,
nor make one a paragraph by itself.

\end{document}               % End of document.

\section{Bench Huge Section 160}
% synthetic bench section 160
% This is a sample LaTeX input file.  (Version of 12 August 2004.)
%
% A '%' character causes TeX to ignore all remaining text on the line,
% and is used for comments like this one.

\documentclass{article}      % Specifies the document class

                             % The preamble begins here.
\title{An Example Document}  % Declares the document's title.
\author{Leslie Lamport}      % Declares the author's name.
\date{January 21, 1994}      % Deleting this command produces today's date.

\newcommand{\ip}[2]{(#1, #2)}
                             % Defines \ip{arg1}{arg2} to mean
                             % (arg1, arg2).

%\newcommand{\ip}[2]{\langle #1 | #2\rangle}
                             % This is an alternative definition of
                             % \ip that is commented out.

\begin{document}             % End of preamble and beginning of text.

\maketitle                   % Produces the title.

This is an example input file.  Comparing it with
the output it generates can show you how to
produce a simple document of your own.

\section{Ordinary Text}      % Produces section heading.  Lower-level
                             % sections are begun with similar
                             % \subsection and \subsubsection commands.

The ends  of words and sentences are marked
  by   spaces. It  doesn't matter how many
spaces    you type; one is as good as 100.  The
end of   a line counts as a space.

One   or more   blank lines denote the  end
of  a paragraph.

Since any number of consecutive spaces are treated
like a single one, the formatting of the input
file makes no difference to
      \LaTeX,                % The \LaTeX command generates the LaTeX logo.
but it makes a difference to you.  When you use
\LaTeX, making your input file as easy to read
as possible will be a great help as you write
your document and when you change it.  This sample
file shows how you can add comments to your own input
file.

Because printing is different from typewriting,
there are a number of things that you have to do
differently when preparing an input file than if
you were just typing the document directly.
Quotation marks like
       ``this''
have to be handled specially, as do quotes within
quotes:
       ``\,`this'            % \, separates the double and single quote.
        is what I just
        wrote, not  `that'\,''.

Dashes come in three sizes: an
       intra-word
dash, a medium dash for number ranges like
       1--2,
and a punctuation
       dash---like
this.

A sentence-ending space should be larger than the
space between words within a sentence.  You
sometimes have to type special commands in
conjunction with punctuation characters to get
this right, as in the following sentence.
       Gnats, gnus, etc.\ all  % `\ ' makes an inter-word space.
       begin with G\@.         % \@ marks end-of-sentence punctuation.
You should check the spaces after periods when
reading your output to make sure you haven't
forgotten any special cases.  Generating an
ellipsis
       \ldots\               % `\ ' is needed after `\ldots' because TeX
                             % ignores spaces after command names like \ldots
                             % made from \ + letters.
                             %
                             % Note how a `%' character causes TeX to ignore
                             % the end of the input line, so these blank lines
                             % do not start a new paragraph.
                             %
with the right spacing around the periods requires
a special command.

\LaTeX\ interprets some common characters as
commands, so you must type special commands to
generate them.  These characters include the
following:
       \$ \& \% \# \{ and \}.

In printing, text is usually emphasized with an
       \emph{italic}
type style.

\begin{em}
   A long segment of text can also be emphasized
   in this way.  Text within such a segment can be
   given \emph{additional} emphasis.
\end{em}

It is sometimes necessary to prevent \LaTeX\ from
breaking a line where it might otherwise do so.
This may be at a space, as between the ``Mr.''\ and
``Jones'' in
       ``Mr.~Jones'',        % ~ produces an unbreakable interword space.
or within a word---especially when the word is a
symbol like
       \mbox{\emph{itemnum}}
that makes little sense when hyphenated across
lines.

Footnotes\footnote{This is an example of a footnote.}
pose no problem.

\LaTeX\ is good at typesetting mathematical formulas
like
       \( x-3y + z = 7 \)
or
       \( a_{1} > x^{2n} + y^{2n} > x' \)
or
       \( \ip{A}{B} = \sum_{i} a_{i} b_{i} \).
The spaces you type in a formula are
ignored.  Remember that a letter like
       $x$                   % $ ... $  and  \( ... \)  are equivalent
is a formula when it denotes a mathematical
symbol, and it should be typed as one.

\section{Displayed Text}

Text is displayed by indenting it from the left
margin.  Quotations are commonly displayed.  There
are short quotations
\begin{quote}
   This is a short quotation.  It consists of a
   single paragraph of text.  See how it is formatted.
\end{quote}
and longer ones.
\begin{quotation}
   This is a longer quotation.  It consists of two
   paragraphs of text, neither of which are
   particularly interesting.

   This is the second paragraph of the quotation.  It
   is just as dull as the first paragraph.
\end{quotation}
Another frequently-displayed structure is a list.
The following is an example of an \emph{itemized}
list.
\begin{itemize}
   \item This is the first item of an itemized list.
         Each item in the list is marked with a ``tick''.
         You don't have to worry about what kind of tick
         mark is used.

   \item This is the second item of the list.  It
         contains another list nested inside it.  The inner
         list is an \emph{enumerated} list.
         \begin{enumerate}
            \item This is the first item of an enumerated
                  list that is nested within the itemized list.

            \item This is the second item of the inner list.
                  \LaTeX\ allows you to nest lists deeper than
                  you really should.
         \end{enumerate}
         This is the rest of the second item of the outer
         list.  It is no more interesting than any other
         part of the item.
   \item This is the third item of the list.
\end{itemize}
You can even display poetry.
\begin{verse}
   There is an environment
    for verse \\             % The \\ command separates lines
   Whose features some poets % within a stanza.
   will curse.

                             % One or more blank lines separate stanzas.

   For instead of making\\
   Them do \emph{all} line breaking, \\
   It allows them to put too many words on a line when they'd rather be
   forced to be terse.
\end{verse}

Mathematical formulas may also be displayed.  A
displayed formula
is
one-line long; multiline
formulas require special formatting instructions.
   \[  \ip{\Gamma}{\psi'} = x'' + y^{2} + z_{i}^{n}\]
Don't start a paragraph with a displayed equation,
nor make one a paragraph by itself.

\end{document}               % End of document.

\section{Bench Huge Section 161}
% synthetic bench section 161
% This is a sample LaTeX input file.  (Version of 12 August 2004.)
%
% A '%' character causes TeX to ignore all remaining text on the line,
% and is used for comments like this one.

\documentclass{article}      % Specifies the document class

                             % The preamble begins here.
\title{An Example Document}  % Declares the document's title.
\author{Leslie Lamport}      % Declares the author's name.
\date{January 21, 1994}      % Deleting this command produces today's date.

\newcommand{\ip}[2]{(#1, #2)}
                             % Defines \ip{arg1}{arg2} to mean
                             % (arg1, arg2).

%\newcommand{\ip}[2]{\langle #1 | #2\rangle}
                             % This is an alternative definition of
                             % \ip that is commented out.

\begin{document}             % End of preamble and beginning of text.

\maketitle                   % Produces the title.

This is an example input file.  Comparing it with
the output it generates can show you how to
produce a simple document of your own.

\section{Ordinary Text}      % Produces section heading.  Lower-level
                             % sections are begun with similar
                             % \subsection and \subsubsection commands.

The ends  of words and sentences are marked
  by   spaces. It  doesn't matter how many
spaces    you type; one is as good as 100.  The
end of   a line counts as a space.

One   or more   blank lines denote the  end
of  a paragraph.

Since any number of consecutive spaces are treated
like a single one, the formatting of the input
file makes no difference to
      \LaTeX,                % The \LaTeX command generates the LaTeX logo.
but it makes a difference to you.  When you use
\LaTeX, making your input file as easy to read
as possible will be a great help as you write
your document and when you change it.  This sample
file shows how you can add comments to your own input
file.

Because printing is different from typewriting,
there are a number of things that you have to do
differently when preparing an input file than if
you were just typing the document directly.
Quotation marks like
       ``this''
have to be handled specially, as do quotes within
quotes:
       ``\,`this'            % \, separates the double and single quote.
        is what I just
        wrote, not  `that'\,''.

Dashes come in three sizes: an
       intra-word
dash, a medium dash for number ranges like
       1--2,
and a punctuation
       dash---like
this.

A sentence-ending space should be larger than the
space between words within a sentence.  You
sometimes have to type special commands in
conjunction with punctuation characters to get
this right, as in the following sentence.
       Gnats, gnus, etc.\ all  % `\ ' makes an inter-word space.
       begin with G\@.         % \@ marks end-of-sentence punctuation.
You should check the spaces after periods when
reading your output to make sure you haven't
forgotten any special cases.  Generating an
ellipsis
       \ldots\               % `\ ' is needed after `\ldots' because TeX
                             % ignores spaces after command names like \ldots
                             % made from \ + letters.
                             %
                             % Note how a `%' character causes TeX to ignore
                             % the end of the input line, so these blank lines
                             % do not start a new paragraph.
                             %
with the right spacing around the periods requires
a special command.

\LaTeX\ interprets some common characters as
commands, so you must type special commands to
generate them.  These characters include the
following:
       \$ \& \% \# \{ and \}.

In printing, text is usually emphasized with an
       \emph{italic}
type style.

\begin{em}
   A long segment of text can also be emphasized
   in this way.  Text within such a segment can be
   given \emph{additional} emphasis.
\end{em}

It is sometimes necessary to prevent \LaTeX\ from
breaking a line where it might otherwise do so.
This may be at a space, as between the ``Mr.''\ and
``Jones'' in
       ``Mr.~Jones'',        % ~ produces an unbreakable interword space.
or within a word---especially when the word is a
symbol like
       \mbox{\emph{itemnum}}
that makes little sense when hyphenated across
lines.

Footnotes\footnote{This is an example of a footnote.}
pose no problem.

\LaTeX\ is good at typesetting mathematical formulas
like
       \( x-3y + z = 7 \)
or
       \( a_{1} > x^{2n} + y^{2n} > x' \)
or
       \( \ip{A}{B} = \sum_{i} a_{i} b_{i} \).
The spaces you type in a formula are
ignored.  Remember that a letter like
       $x$                   % $ ... $  and  \( ... \)  are equivalent
is a formula when it denotes a mathematical
symbol, and it should be typed as one.

\section{Displayed Text}

Text is displayed by indenting it from the left
margin.  Quotations are commonly displayed.  There
are short quotations
\begin{quote}
   This is a short quotation.  It consists of a
   single paragraph of text.  See how it is formatted.
\end{quote}
and longer ones.
\begin{quotation}
   This is a longer quotation.  It consists of two
   paragraphs of text, neither of which are
   particularly interesting.

   This is the second paragraph of the quotation.  It
   is just as dull as the first paragraph.
\end{quotation}
Another frequently-displayed structure is a list.
The following is an example of an \emph{itemized}
list.
\begin{itemize}
   \item This is the first item of an itemized list.
         Each item in the list is marked with a ``tick''.
         You don't have to worry about what kind of tick
         mark is used.

   \item This is the second item of the list.  It
         contains another list nested inside it.  The inner
         list is an \emph{enumerated} list.
         \begin{enumerate}
            \item This is the first item of an enumerated
                  list that is nested within the itemized list.

            \item This is the second item of the inner list.
                  \LaTeX\ allows you to nest lists deeper than
                  you really should.
         \end{enumerate}
         This is the rest of the second item of the outer
         list.  It is no more interesting than any other
         part of the item.
   \item This is the third item of the list.
\end{itemize}
You can even display poetry.
\begin{verse}
   There is an environment
    for verse \\             % The \\ command separates lines
   Whose features some poets % within a stanza.
   will curse.

                             % One or more blank lines separate stanzas.

   For instead of making\\
   Them do \emph{all} line breaking, \\
   It allows them to put too many words on a line when they'd rather be
   forced to be terse.
\end{verse}

Mathematical formulas may also be displayed.  A
displayed formula
is
one-line long; multiline
formulas require special formatting instructions.
   \[  \ip{\Gamma}{\psi'} = x'' + y^{2} + z_{i}^{n}\]
Don't start a paragraph with a displayed equation,
nor make one a paragraph by itself.

\end{document}               % End of document.

\section{Bench Huge Section 162}
% synthetic bench section 162
% This is a sample LaTeX input file.  (Version of 12 August 2004.)
%
% A '%' character causes TeX to ignore all remaining text on the line,
% and is used for comments like this one.

\documentclass{article}      % Specifies the document class

                             % The preamble begins here.
\title{An Example Document}  % Declares the document's title.
\author{Leslie Lamport}      % Declares the author's name.
\date{January 21, 1994}      % Deleting this command produces today's date.

\newcommand{\ip}[2]{(#1, #2)}
                             % Defines \ip{arg1}{arg2} to mean
                             % (arg1, arg2).

%\newcommand{\ip}[2]{\langle #1 | #2\rangle}
                             % This is an alternative definition of
                             % \ip that is commented out.

\begin{document}             % End of preamble and beginning of text.

\maketitle                   % Produces the title.

This is an example input file.  Comparing it with
the output it generates can show you how to
produce a simple document of your own.

\section{Ordinary Text}      % Produces section heading.  Lower-level
                             % sections are begun with similar
                             % \subsection and \subsubsection commands.

The ends  of words and sentences are marked
  by   spaces. It  doesn't matter how many
spaces    you type; one is as good as 100.  The
end of   a line counts as a space.

One   or more   blank lines denote the  end
of  a paragraph.

Since any number of consecutive spaces are treated
like a single one, the formatting of the input
file makes no difference to
      \LaTeX,                % The \LaTeX command generates the LaTeX logo.
but it makes a difference to you.  When you use
\LaTeX, making your input file as easy to read
as possible will be a great help as you write
your document and when you change it.  This sample
file shows how you can add comments to your own input
file.

Because printing is different from typewriting,
there are a number of things that you have to do
differently when preparing an input file than if
you were just typing the document directly.
Quotation marks like
       ``this''
have to be handled specially, as do quotes within
quotes:
       ``\,`this'            % \, separates the double and single quote.
        is what I just
        wrote, not  `that'\,''.

Dashes come in three sizes: an
       intra-word
dash, a medium dash for number ranges like
       1--2,
and a punctuation
       dash---like
this.

A sentence-ending space should be larger than the
space between words within a sentence.  You
sometimes have to type special commands in
conjunction with punctuation characters to get
this right, as in the following sentence.
       Gnats, gnus, etc.\ all  % `\ ' makes an inter-word space.
       begin with G\@.         % \@ marks end-of-sentence punctuation.
You should check the spaces after periods when
reading your output to make sure you haven't
forgotten any special cases.  Generating an
ellipsis
       \ldots\               % `\ ' is needed after `\ldots' because TeX
                             % ignores spaces after command names like \ldots
                             % made from \ + letters.
                             %
                             % Note how a `%' character causes TeX to ignore
                             % the end of the input line, so these blank lines
                             % do not start a new paragraph.
                             %
with the right spacing around the periods requires
a special command.

\LaTeX\ interprets some common characters as
commands, so you must type special commands to
generate them.  These characters include the
following:
       \$ \& \% \# \{ and \}.

In printing, text is usually emphasized with an
       \emph{italic}
type style.

\begin{em}
   A long segment of text can also be emphasized
   in this way.  Text within such a segment can be
   given \emph{additional} emphasis.
\end{em}

It is sometimes necessary to prevent \LaTeX\ from
breaking a line where it might otherwise do so.
This may be at a space, as between the ``Mr.''\ and
``Jones'' in
       ``Mr.~Jones'',        % ~ produces an unbreakable interword space.
or within a word---especially when the word is a
symbol like
       \mbox{\emph{itemnum}}
that makes little sense when hyphenated across
lines.

Footnotes\footnote{This is an example of a footnote.}
pose no problem.

\LaTeX\ is good at typesetting mathematical formulas
like
       \( x-3y + z = 7 \)
or
       \( a_{1} > x^{2n} + y^{2n} > x' \)
or
       \( \ip{A}{B} = \sum_{i} a_{i} b_{i} \).
The spaces you type in a formula are
ignored.  Remember that a letter like
       $x$                   % $ ... $  and  \( ... \)  are equivalent
is a formula when it denotes a mathematical
symbol, and it should be typed as one.

\section{Displayed Text}

Text is displayed by indenting it from the left
margin.  Quotations are commonly displayed.  There
are short quotations
\begin{quote}
   This is a short quotation.  It consists of a
   single paragraph of text.  See how it is formatted.
\end{quote}
and longer ones.
\begin{quotation}
   This is a longer quotation.  It consists of two
   paragraphs of text, neither of which are
   particularly interesting.

   This is the second paragraph of the quotation.  It
   is just as dull as the first paragraph.
\end{quotation}
Another frequently-displayed structure is a list.
The following is an example of an \emph{itemized}
list.
\begin{itemize}
   \item This is the first item of an itemized list.
         Each item in the list is marked with a ``tick''.
         You don't have to worry about what kind of tick
         mark is used.

   \item This is the second item of the list.  It
         contains another list nested inside it.  The inner
         list is an \emph{enumerated} list.
         \begin{enumerate}
            \item This is the first item of an enumerated
                  list that is nested within the itemized list.

            \item This is the second item of the inner list.
                  \LaTeX\ allows you to nest lists deeper than
                  you really should.
         \end{enumerate}
         This is the rest of the second item of the outer
         list.  It is no more interesting than any other
         part of the item.
   \item This is the third item of the list.
\end{itemize}
You can even display poetry.
\begin{verse}
   There is an environment
    for verse \\             % The \\ command separates lines
   Whose features some poets % within a stanza.
   will curse.

                             % One or more blank lines separate stanzas.

   For instead of making\\
   Them do \emph{all} line breaking, \\
   It allows them to put too many words on a line when they'd rather be
   forced to be terse.
\end{verse}

Mathematical formulas may also be displayed.  A
displayed formula
is
one-line long; multiline
formulas require special formatting instructions.
   \[  \ip{\Gamma}{\psi'} = x'' + y^{2} + z_{i}^{n}\]
Don't start a paragraph with a displayed equation,
nor make one a paragraph by itself.

\end{document}               % End of document.

\section{Bench Huge Section 163}
% synthetic bench section 163
% This is a sample LaTeX input file.  (Version of 12 August 2004.)
%
% A '%' character causes TeX to ignore all remaining text on the line,
% and is used for comments like this one.

\documentclass{article}      % Specifies the document class

                             % The preamble begins here.
\title{An Example Document}  % Declares the document's title.
\author{Leslie Lamport}      % Declares the author's name.
\date{January 21, 1994}      % Deleting this command produces today's date.

\newcommand{\ip}[2]{(#1, #2)}
                             % Defines \ip{arg1}{arg2} to mean
                             % (arg1, arg2).

%\newcommand{\ip}[2]{\langle #1 | #2\rangle}
                             % This is an alternative definition of
                             % \ip that is commented out.

\begin{document}             % End of preamble and beginning of text.

\maketitle                   % Produces the title.

This is an example input file.  Comparing it with
the output it generates can show you how to
produce a simple document of your own.

\section{Ordinary Text}      % Produces section heading.  Lower-level
                             % sections are begun with similar
                             % \subsection and \subsubsection commands.

The ends  of words and sentences are marked
  by   spaces. It  doesn't matter how many
spaces    you type; one is as good as 100.  The
end of   a line counts as a space.

One   or more   blank lines denote the  end
of  a paragraph.

Since any number of consecutive spaces are treated
like a single one, the formatting of the input
file makes no difference to
      \LaTeX,                % The \LaTeX command generates the LaTeX logo.
but it makes a difference to you.  When you use
\LaTeX, making your input file as easy to read
as possible will be a great help as you write
your document and when you change it.  This sample
file shows how you can add comments to your own input
file.

Because printing is different from typewriting,
there are a number of things that you have to do
differently when preparing an input file than if
you were just typing the document directly.
Quotation marks like
       ``this''
have to be handled specially, as do quotes within
quotes:
       ``\,`this'            % \, separates the double and single quote.
        is what I just
        wrote, not  `that'\,''.

Dashes come in three sizes: an
       intra-word
dash, a medium dash for number ranges like
       1--2,
and a punctuation
       dash---like
this.

A sentence-ending space should be larger than the
space between words within a sentence.  You
sometimes have to type special commands in
conjunction with punctuation characters to get
this right, as in the following sentence.
       Gnats, gnus, etc.\ all  % `\ ' makes an inter-word space.
       begin with G\@.         % \@ marks end-of-sentence punctuation.
You should check the spaces after periods when
reading your output to make sure you haven't
forgotten any special cases.  Generating an
ellipsis
       \ldots\               % `\ ' is needed after `\ldots' because TeX
                             % ignores spaces after command names like \ldots
                             % made from \ + letters.
                             %
                             % Note how a `%' character causes TeX to ignore
                             % the end of the input line, so these blank lines
                             % do not start a new paragraph.
                             %
with the right spacing around the periods requires
a special command.

\LaTeX\ interprets some common characters as
commands, so you must type special commands to
generate them.  These characters include the
following:
       \$ \& \% \# \{ and \}.

In printing, text is usually emphasized with an
       \emph{italic}
type style.

\begin{em}
   A long segment of text can also be emphasized
   in this way.  Text within such a segment can be
   given \emph{additional} emphasis.
\end{em}

It is sometimes necessary to prevent \LaTeX\ from
breaking a line where it might otherwise do so.
This may be at a space, as between the ``Mr.''\ and
``Jones'' in
       ``Mr.~Jones'',        % ~ produces an unbreakable interword space.
or within a word---especially when the word is a
symbol like
       \mbox{\emph{itemnum}}
that makes little sense when hyphenated across
lines.

Footnotes\footnote{This is an example of a footnote.}
pose no problem.

\LaTeX\ is good at typesetting mathematical formulas
like
       \( x-3y + z = 7 \)
or
       \( a_{1} > x^{2n} + y^{2n} > x' \)
or
       \( \ip{A}{B} = \sum_{i} a_{i} b_{i} \).
The spaces you type in a formula are
ignored.  Remember that a letter like
       $x$                   % $ ... $  and  \( ... \)  are equivalent
is a formula when it denotes a mathematical
symbol, and it should be typed as one.

\section{Displayed Text}

Text is displayed by indenting it from the left
margin.  Quotations are commonly displayed.  There
are short quotations
\begin{quote}
   This is a short quotation.  It consists of a
   single paragraph of text.  See how it is formatted.
\end{quote}
and longer ones.
\begin{quotation}
   This is a longer quotation.  It consists of two
   paragraphs of text, neither of which are
   particularly interesting.

   This is the second paragraph of the quotation.  It
   is just as dull as the first paragraph.
\end{quotation}
Another frequently-displayed structure is a list.
The following is an example of an \emph{itemized}
list.
\begin{itemize}
   \item This is the first item of an itemized list.
         Each item in the list is marked with a ``tick''.
         You don't have to worry about what kind of tick
         mark is used.

   \item This is the second item of the list.  It
         contains another list nested inside it.  The inner
         list is an \emph{enumerated} list.
         \begin{enumerate}
            \item This is the first item of an enumerated
                  list that is nested within the itemized list.

            \item This is the second item of the inner list.
                  \LaTeX\ allows you to nest lists deeper than
                  you really should.
         \end{enumerate}
         This is the rest of the second item of the outer
         list.  It is no more interesting than any other
         part of the item.
   \item This is the third item of the list.
\end{itemize}
You can even display poetry.
\begin{verse}
   There is an environment
    for verse \\             % The \\ command separates lines
   Whose features some poets % within a stanza.
   will curse.

                             % One or more blank lines separate stanzas.

   For instead of making\\
   Them do \emph{all} line breaking, \\
   It allows them to put too many words on a line when they'd rather be
   forced to be terse.
\end{verse}

Mathematical formulas may also be displayed.  A
displayed formula
is
one-line long; multiline
formulas require special formatting instructions.
   \[  \ip{\Gamma}{\psi'} = x'' + y^{2} + z_{i}^{n}\]
Don't start a paragraph with a displayed equation,
nor make one a paragraph by itself.

\end{document}               % End of document.

\section{Bench Huge Section 164}
% synthetic bench section 164
% This is a sample LaTeX input file.  (Version of 12 August 2004.)
%
% A '%' character causes TeX to ignore all remaining text on the line,
% and is used for comments like this one.

\documentclass{article}      % Specifies the document class

                             % The preamble begins here.
\title{An Example Document}  % Declares the document's title.
\author{Leslie Lamport}      % Declares the author's name.
\date{January 21, 1994}      % Deleting this command produces today's date.

\newcommand{\ip}[2]{(#1, #2)}
                             % Defines \ip{arg1}{arg2} to mean
                             % (arg1, arg2).

%\newcommand{\ip}[2]{\langle #1 | #2\rangle}
                             % This is an alternative definition of
                             % \ip that is commented out.

\begin{document}             % End of preamble and beginning of text.

\maketitle                   % Produces the title.

This is an example input file.  Comparing it with
the output it generates can show you how to
produce a simple document of your own.

\section{Ordinary Text}      % Produces section heading.  Lower-level
                             % sections are begun with similar
                             % \subsection and \subsubsection commands.

The ends  of words and sentences are marked
  by   spaces. It  doesn't matter how many
spaces    you type; one is as good as 100.  The
end of   a line counts as a space.

One   or more   blank lines denote the  end
of  a paragraph.

Since any number of consecutive spaces are treated
like a single one, the formatting of the input
file makes no difference to
      \LaTeX,                % The \LaTeX command generates the LaTeX logo.
but it makes a difference to you.  When you use
\LaTeX, making your input file as easy to read
as possible will be a great help as you write
your document and when you change it.  This sample
file shows how you can add comments to your own input
file.

Because printing is different from typewriting,
there are a number of things that you have to do
differently when preparing an input file than if
you were just typing the document directly.
Quotation marks like
       ``this''
have to be handled specially, as do quotes within
quotes:
       ``\,`this'            % \, separates the double and single quote.
        is what I just
        wrote, not  `that'\,''.

Dashes come in three sizes: an
       intra-word
dash, a medium dash for number ranges like
       1--2,
and a punctuation
       dash---like
this.

A sentence-ending space should be larger than the
space between words within a sentence.  You
sometimes have to type special commands in
conjunction with punctuation characters to get
this right, as in the following sentence.
       Gnats, gnus, etc.\ all  % `\ ' makes an inter-word space.
       begin with G\@.         % \@ marks end-of-sentence punctuation.
You should check the spaces after periods when
reading your output to make sure you haven't
forgotten any special cases.  Generating an
ellipsis
       \ldots\               % `\ ' is needed after `\ldots' because TeX
                             % ignores spaces after command names like \ldots
                             % made from \ + letters.
                             %
                             % Note how a `%' character causes TeX to ignore
                             % the end of the input line, so these blank lines
                             % do not start a new paragraph.
                             %
with the right spacing around the periods requires
a special command.

\LaTeX\ interprets some common characters as
commands, so you must type special commands to
generate them.  These characters include the
following:
       \$ \& \% \# \{ and \}.

In printing, text is usually emphasized with an
       \emph{italic}
type style.

\begin{em}
   A long segment of text can also be emphasized
   in this way.  Text within such a segment can be
   given \emph{additional} emphasis.
\end{em}

It is sometimes necessary to prevent \LaTeX\ from
breaking a line where it might otherwise do so.
This may be at a space, as between the ``Mr.''\ and
``Jones'' in
       ``Mr.~Jones'',        % ~ produces an unbreakable interword space.
or within a word---especially when the word is a
symbol like
       \mbox{\emph{itemnum}}
that makes little sense when hyphenated across
lines.

Footnotes\footnote{This is an example of a footnote.}
pose no problem.

\LaTeX\ is good at typesetting mathematical formulas
like
       \( x-3y + z = 7 \)
or
       \( a_{1} > x^{2n} + y^{2n} > x' \)
or
       \( \ip{A}{B} = \sum_{i} a_{i} b_{i} \).
The spaces you type in a formula are
ignored.  Remember that a letter like
       $x$                   % $ ... $  and  \( ... \)  are equivalent
is a formula when it denotes a mathematical
symbol, and it should be typed as one.

\section{Displayed Text}

Text is displayed by indenting it from the left
margin.  Quotations are commonly displayed.  There
are short quotations
\begin{quote}
   This is a short quotation.  It consists of a
   single paragraph of text.  See how it is formatted.
\end{quote}
and longer ones.
\begin{quotation}
   This is a longer quotation.  It consists of two
   paragraphs of text, neither of which are
   particularly interesting.

   This is the second paragraph of the quotation.  It
   is just as dull as the first paragraph.
\end{quotation}
Another frequently-displayed structure is a list.
The following is an example of an \emph{itemized}
list.
\begin{itemize}
   \item This is the first item of an itemized list.
         Each item in the list is marked with a ``tick''.
         You don't have to worry about what kind of tick
         mark is used.

   \item This is the second item of the list.  It
         contains another list nested inside it.  The inner
         list is an \emph{enumerated} list.
         \begin{enumerate}
            \item This is the first item of an enumerated
                  list that is nested within the itemized list.

            \item This is the second item of the inner list.
                  \LaTeX\ allows you to nest lists deeper than
                  you really should.
         \end{enumerate}
         This is the rest of the second item of the outer
         list.  It is no more interesting than any other
         part of the item.
   \item This is the third item of the list.
\end{itemize}
You can even display poetry.
\begin{verse}
   There is an environment
    for verse \\             % The \\ command separates lines
   Whose features some poets % within a stanza.
   will curse.

                             % One or more blank lines separate stanzas.

   For instead of making\\
   Them do \emph{all} line breaking, \\
   It allows them to put too many words on a line when they'd rather be
   forced to be terse.
\end{verse}

Mathematical formulas may also be displayed.  A
displayed formula
is
one-line long; multiline
formulas require special formatting instructions.
   \[  \ip{\Gamma}{\psi'} = x'' + y^{2} + z_{i}^{n}\]
Don't start a paragraph with a displayed equation,
nor make one a paragraph by itself.

\end{document}               % End of document.

\section{Bench Huge Section 165}
% synthetic bench section 165
% This is a sample LaTeX input file.  (Version of 12 August 2004.)
%
% A '%' character causes TeX to ignore all remaining text on the line,
% and is used for comments like this one.

\documentclass{article}      % Specifies the document class

                             % The preamble begins here.
\title{An Example Document}  % Declares the document's title.
\author{Leslie Lamport}      % Declares the author's name.
\date{January 21, 1994}      % Deleting this command produces today's date.

\newcommand{\ip}[2]{(#1, #2)}
                             % Defines \ip{arg1}{arg2} to mean
                             % (arg1, arg2).

%\newcommand{\ip}[2]{\langle #1 | #2\rangle}
                             % This is an alternative definition of
                             % \ip that is commented out.

\begin{document}             % End of preamble and beginning of text.

\maketitle                   % Produces the title.

This is an example input file.  Comparing it with
the output it generates can show you how to
produce a simple document of your own.

\section{Ordinary Text}      % Produces section heading.  Lower-level
                             % sections are begun with similar
                             % \subsection and \subsubsection commands.

The ends  of words and sentences are marked
  by   spaces. It  doesn't matter how many
spaces    you type; one is as good as 100.  The
end of   a line counts as a space.

One   or more   blank lines denote the  end
of  a paragraph.

Since any number of consecutive spaces are treated
like a single one, the formatting of the input
file makes no difference to
      \LaTeX,                % The \LaTeX command generates the LaTeX logo.
but it makes a difference to you.  When you use
\LaTeX, making your input file as easy to read
as possible will be a great help as you write
your document and when you change it.  This sample
file shows how you can add comments to your own input
file.

Because printing is different from typewriting,
there are a number of things that you have to do
differently when preparing an input file than if
you were just typing the document directly.
Quotation marks like
       ``this''
have to be handled specially, as do quotes within
quotes:
       ``\,`this'            % \, separates the double and single quote.
        is what I just
        wrote, not  `that'\,''.

Dashes come in three sizes: an
       intra-word
dash, a medium dash for number ranges like
       1--2,
and a punctuation
       dash---like
this.

A sentence-ending space should be larger than the
space between words within a sentence.  You
sometimes have to type special commands in
conjunction with punctuation characters to get
this right, as in the following sentence.
       Gnats, gnus, etc.\ all  % `\ ' makes an inter-word space.
       begin with G\@.         % \@ marks end-of-sentence punctuation.
You should check the spaces after periods when
reading your output to make sure you haven't
forgotten any special cases.  Generating an
ellipsis
       \ldots\               % `\ ' is needed after `\ldots' because TeX
                             % ignores spaces after command names like \ldots
                             % made from \ + letters.
                             %
                             % Note how a `%' character causes TeX to ignore
                             % the end of the input line, so these blank lines
                             % do not start a new paragraph.
                             %
with the right spacing around the periods requires
a special command.

\LaTeX\ interprets some common characters as
commands, so you must type special commands to
generate them.  These characters include the
following:
       \$ \& \% \# \{ and \}.

In printing, text is usually emphasized with an
       \emph{italic}
type style.

\begin{em}
   A long segment of text can also be emphasized
   in this way.  Text within such a segment can be
   given \emph{additional} emphasis.
\end{em}

It is sometimes necessary to prevent \LaTeX\ from
breaking a line where it might otherwise do so.
This may be at a space, as between the ``Mr.''\ and
``Jones'' in
       ``Mr.~Jones'',        % ~ produces an unbreakable interword space.
or within a word---especially when the word is a
symbol like
       \mbox{\emph{itemnum}}
that makes little sense when hyphenated across
lines.

Footnotes\footnote{This is an example of a footnote.}
pose no problem.

\LaTeX\ is good at typesetting mathematical formulas
like
       \( x-3y + z = 7 \)
or
       \( a_{1} > x^{2n} + y^{2n} > x' \)
or
       \( \ip{A}{B} = \sum_{i} a_{i} b_{i} \).
The spaces you type in a formula are
ignored.  Remember that a letter like
       $x$                   % $ ... $  and  \( ... \)  are equivalent
is a formula when it denotes a mathematical
symbol, and it should be typed as one.

\section{Displayed Text}

Text is displayed by indenting it from the left
margin.  Quotations are commonly displayed.  There
are short quotations
\begin{quote}
   This is a short quotation.  It consists of a
   single paragraph of text.  See how it is formatted.
\end{quote}
and longer ones.
\begin{quotation}
   This is a longer quotation.  It consists of two
   paragraphs of text, neither of which are
   particularly interesting.

   This is the second paragraph of the quotation.  It
   is just as dull as the first paragraph.
\end{quotation}
Another frequently-displayed structure is a list.
The following is an example of an \emph{itemized}
list.
\begin{itemize}
   \item This is the first item of an itemized list.
         Each item in the list is marked with a ``tick''.
         You don't have to worry about what kind of tick
         mark is used.

   \item This is the second item of the list.  It
         contains another list nested inside it.  The inner
         list is an \emph{enumerated} list.
         \begin{enumerate}
            \item This is the first item of an enumerated
                  list that is nested within the itemized list.

            \item This is the second item of the inner list.
                  \LaTeX\ allows you to nest lists deeper than
                  you really should.
         \end{enumerate}
         This is the rest of the second item of the outer
         list.  It is no more interesting than any other
         part of the item.
   \item This is the third item of the list.
\end{itemize}
You can even display poetry.
\begin{verse}
   There is an environment
    for verse \\             % The \\ command separates lines
   Whose features some poets % within a stanza.
   will curse.

                             % One or more blank lines separate stanzas.

   For instead of making\\
   Them do \emph{all} line breaking, \\
   It allows them to put too many words on a line when they'd rather be
   forced to be terse.
\end{verse}

Mathematical formulas may also be displayed.  A
displayed formula
is
one-line long; multiline
formulas require special formatting instructions.
   \[  \ip{\Gamma}{\psi'} = x'' + y^{2} + z_{i}^{n}\]
Don't start a paragraph with a displayed equation,
nor make one a paragraph by itself.

\end{document}               % End of document.

\section{Bench Huge Section 166}
% synthetic bench section 166
% This is a sample LaTeX input file.  (Version of 12 August 2004.)
%
% A '%' character causes TeX to ignore all remaining text on the line,
% and is used for comments like this one.

\documentclass{article}      % Specifies the document class

                             % The preamble begins here.
\title{An Example Document}  % Declares the document's title.
\author{Leslie Lamport}      % Declares the author's name.
\date{January 21, 1994}      % Deleting this command produces today's date.

\newcommand{\ip}[2]{(#1, #2)}
                             % Defines \ip{arg1}{arg2} to mean
                             % (arg1, arg2).

%\newcommand{\ip}[2]{\langle #1 | #2\rangle}
                             % This is an alternative definition of
                             % \ip that is commented out.

\begin{document}             % End of preamble and beginning of text.

\maketitle                   % Produces the title.

This is an example input file.  Comparing it with
the output it generates can show you how to
produce a simple document of your own.

\section{Ordinary Text}      % Produces section heading.  Lower-level
                             % sections are begun with similar
                             % \subsection and \subsubsection commands.

The ends  of words and sentences are marked
  by   spaces. It  doesn't matter how many
spaces    you type; one is as good as 100.  The
end of   a line counts as a space.

One   or more   blank lines denote the  end
of  a paragraph.

Since any number of consecutive spaces are treated
like a single one, the formatting of the input
file makes no difference to
      \LaTeX,                % The \LaTeX command generates the LaTeX logo.
but it makes a difference to you.  When you use
\LaTeX, making your input file as easy to read
as possible will be a great help as you write
your document and when you change it.  This sample
file shows how you can add comments to your own input
file.

Because printing is different from typewriting,
there are a number of things that you have to do
differently when preparing an input file than if
you were just typing the document directly.
Quotation marks like
       ``this''
have to be handled specially, as do quotes within
quotes:
       ``\,`this'            % \, separates the double and single quote.
        is what I just
        wrote, not  `that'\,''.

Dashes come in three sizes: an
       intra-word
dash, a medium dash for number ranges like
       1--2,
and a punctuation
       dash---like
this.

A sentence-ending space should be larger than the
space between words within a sentence.  You
sometimes have to type special commands in
conjunction with punctuation characters to get
this right, as in the following sentence.
       Gnats, gnus, etc.\ all  % `\ ' makes an inter-word space.
       begin with G\@.         % \@ marks end-of-sentence punctuation.
You should check the spaces after periods when
reading your output to make sure you haven't
forgotten any special cases.  Generating an
ellipsis
       \ldots\               % `\ ' is needed after `\ldots' because TeX
                             % ignores spaces after command names like \ldots
                             % made from \ + letters.
                             %
                             % Note how a `%' character causes TeX to ignore
                             % the end of the input line, so these blank lines
                             % do not start a new paragraph.
                             %
with the right spacing around the periods requires
a special command.

\LaTeX\ interprets some common characters as
commands, so you must type special commands to
generate them.  These characters include the
following:
       \$ \& \% \# \{ and \}.

In printing, text is usually emphasized with an
       \emph{italic}
type style.

\begin{em}
   A long segment of text can also be emphasized
   in this way.  Text within such a segment can be
   given \emph{additional} emphasis.
\end{em}

It is sometimes necessary to prevent \LaTeX\ from
breaking a line where it might otherwise do so.
This may be at a space, as between the ``Mr.''\ and
``Jones'' in
       ``Mr.~Jones'',        % ~ produces an unbreakable interword space.
or within a word---especially when the word is a
symbol like
       \mbox{\emph{itemnum}}
that makes little sense when hyphenated across
lines.

Footnotes\footnote{This is an example of a footnote.}
pose no problem.

\LaTeX\ is good at typesetting mathematical formulas
like
       \( x-3y + z = 7 \)
or
       \( a_{1} > x^{2n} + y^{2n} > x' \)
or
       \( \ip{A}{B} = \sum_{i} a_{i} b_{i} \).
The spaces you type in a formula are
ignored.  Remember that a letter like
       $x$                   % $ ... $  and  \( ... \)  are equivalent
is a formula when it denotes a mathematical
symbol, and it should be typed as one.

\section{Displayed Text}

Text is displayed by indenting it from the left
margin.  Quotations are commonly displayed.  There
are short quotations
\begin{quote}
   This is a short quotation.  It consists of a
   single paragraph of text.  See how it is formatted.
\end{quote}
and longer ones.
\begin{quotation}
   This is a longer quotation.  It consists of two
   paragraphs of text, neither of which are
   particularly interesting.

   This is the second paragraph of the quotation.  It
   is just as dull as the first paragraph.
\end{quotation}
Another frequently-displayed structure is a list.
The following is an example of an \emph{itemized}
list.
\begin{itemize}
   \item This is the first item of an itemized list.
         Each item in the list is marked with a ``tick''.
         You don't have to worry about what kind of tick
         mark is used.

   \item This is the second item of the list.  It
         contains another list nested inside it.  The inner
         list is an \emph{enumerated} list.
         \begin{enumerate}
            \item This is the first item of an enumerated
                  list that is nested within the itemized list.

            \item This is the second item of the inner list.
                  \LaTeX\ allows you to nest lists deeper than
                  you really should.
         \end{enumerate}
         This is the rest of the second item of the outer
         list.  It is no more interesting than any other
         part of the item.
   \item This is the third item of the list.
\end{itemize}
You can even display poetry.
\begin{verse}
   There is an environment
    for verse \\             % The \\ command separates lines
   Whose features some poets % within a stanza.
   will curse.

                             % One or more blank lines separate stanzas.

   For instead of making\\
   Them do \emph{all} line breaking, \\
   It allows them to put too many words on a line when they'd rather be
   forced to be terse.
\end{verse}

Mathematical formulas may also be displayed.  A
displayed formula
is
one-line long; multiline
formulas require special formatting instructions.
   \[  \ip{\Gamma}{\psi'} = x'' + y^{2} + z_{i}^{n}\]
Don't start a paragraph with a displayed equation,
nor make one a paragraph by itself.

\end{document}               % End of document.

\section{Bench Huge Section 167}
% synthetic bench section 167
% This is a sample LaTeX input file.  (Version of 12 August 2004.)
%
% A '%' character causes TeX to ignore all remaining text on the line,
% and is used for comments like this one.

\documentclass{article}      % Specifies the document class

                             % The preamble begins here.
\title{An Example Document}  % Declares the document's title.
\author{Leslie Lamport}      % Declares the author's name.
\date{January 21, 1994}      % Deleting this command produces today's date.

\newcommand{\ip}[2]{(#1, #2)}
                             % Defines \ip{arg1}{arg2} to mean
                             % (arg1, arg2).

%\newcommand{\ip}[2]{\langle #1 | #2\rangle}
                             % This is an alternative definition of
                             % \ip that is commented out.

\begin{document}             % End of preamble and beginning of text.

\maketitle                   % Produces the title.

This is an example input file.  Comparing it with
the output it generates can show you how to
produce a simple document of your own.

\section{Ordinary Text}      % Produces section heading.  Lower-level
                             % sections are begun with similar
                             % \subsection and \subsubsection commands.

The ends  of words and sentences are marked
  by   spaces. It  doesn't matter how many
spaces    you type; one is as good as 100.  The
end of   a line counts as a space.

One   or more   blank lines denote the  end
of  a paragraph.

Since any number of consecutive spaces are treated
like a single one, the formatting of the input
file makes no difference to
      \LaTeX,                % The \LaTeX command generates the LaTeX logo.
but it makes a difference to you.  When you use
\LaTeX, making your input file as easy to read
as possible will be a great help as you write
your document and when you change it.  This sample
file shows how you can add comments to your own input
file.

Because printing is different from typewriting,
there are a number of things that you have to do
differently when preparing an input file than if
you were just typing the document directly.
Quotation marks like
       ``this''
have to be handled specially, as do quotes within
quotes:
       ``\,`this'            % \, separates the double and single quote.
        is what I just
        wrote, not  `that'\,''.

Dashes come in three sizes: an
       intra-word
dash, a medium dash for number ranges like
       1--2,
and a punctuation
       dash---like
this.

A sentence-ending space should be larger than the
space between words within a sentence.  You
sometimes have to type special commands in
conjunction with punctuation characters to get
this right, as in the following sentence.
       Gnats, gnus, etc.\ all  % `\ ' makes an inter-word space.
       begin with G\@.         % \@ marks end-of-sentence punctuation.
You should check the spaces after periods when
reading your output to make sure you haven't
forgotten any special cases.  Generating an
ellipsis
       \ldots\               % `\ ' is needed after `\ldots' because TeX
                             % ignores spaces after command names like \ldots
                             % made from \ + letters.
                             %
                             % Note how a `%' character causes TeX to ignore
                             % the end of the input line, so these blank lines
                             % do not start a new paragraph.
                             %
with the right spacing around the periods requires
a special command.

\LaTeX\ interprets some common characters as
commands, so you must type special commands to
generate them.  These characters include the
following:
       \$ \& \% \# \{ and \}.

In printing, text is usually emphasized with an
       \emph{italic}
type style.

\begin{em}
   A long segment of text can also be emphasized
   in this way.  Text within such a segment can be
   given \emph{additional} emphasis.
\end{em}

It is sometimes necessary to prevent \LaTeX\ from
breaking a line where it might otherwise do so.
This may be at a space, as between the ``Mr.''\ and
``Jones'' in
       ``Mr.~Jones'',        % ~ produces an unbreakable interword space.
or within a word---especially when the word is a
symbol like
       \mbox{\emph{itemnum}}
that makes little sense when hyphenated across
lines.

Footnotes\footnote{This is an example of a footnote.}
pose no problem.

\LaTeX\ is good at typesetting mathematical formulas
like
       \( x-3y + z = 7 \)
or
       \( a_{1} > x^{2n} + y^{2n} > x' \)
or
       \( \ip{A}{B} = \sum_{i} a_{i} b_{i} \).
The spaces you type in a formula are
ignored.  Remember that a letter like
       $x$                   % $ ... $  and  \( ... \)  are equivalent
is a formula when it denotes a mathematical
symbol, and it should be typed as one.

\section{Displayed Text}

Text is displayed by indenting it from the left
margin.  Quotations are commonly displayed.  There
are short quotations
\begin{quote}
   This is a short quotation.  It consists of a
   single paragraph of text.  See how it is formatted.
\end{quote}
and longer ones.
\begin{quotation}
   This is a longer quotation.  It consists of two
   paragraphs of text, neither of which are
   particularly interesting.

   This is the second paragraph of the quotation.  It
   is just as dull as the first paragraph.
\end{quotation}
Another frequently-displayed structure is a list.
The following is an example of an \emph{itemized}
list.
\begin{itemize}
   \item This is the first item of an itemized list.
         Each item in the list is marked with a ``tick''.
         You don't have to worry about what kind of tick
         mark is used.

   \item This is the second item of the list.  It
         contains another list nested inside it.  The inner
         list is an \emph{enumerated} list.
         \begin{enumerate}
            \item This is the first item of an enumerated
                  list that is nested within the itemized list.

            \item This is the second item of the inner list.
                  \LaTeX\ allows you to nest lists deeper than
                  you really should.
         \end{enumerate}
         This is the rest of the second item of the outer
         list.  It is no more interesting than any other
         part of the item.
   \item This is the third item of the list.
\end{itemize}
You can even display poetry.
\begin{verse}
   There is an environment
    for verse \\             % The \\ command separates lines
   Whose features some poets % within a stanza.
   will curse.

                             % One or more blank lines separate stanzas.

   For instead of making\\
   Them do \emph{all} line breaking, \\
   It allows them to put too many words on a line when they'd rather be
   forced to be terse.
\end{verse}

Mathematical formulas may also be displayed.  A
displayed formula
is
one-line long; multiline
formulas require special formatting instructions.
   \[  \ip{\Gamma}{\psi'} = x'' + y^{2} + z_{i}^{n}\]
Don't start a paragraph with a displayed equation,
nor make one a paragraph by itself.

\end{document}               % End of document.

\section{Bench Huge Section 168}
% synthetic bench section 168
% This is a sample LaTeX input file.  (Version of 12 August 2004.)
%
% A '%' character causes TeX to ignore all remaining text on the line,
% and is used for comments like this one.

\documentclass{article}      % Specifies the document class

                             % The preamble begins here.
\title{An Example Document}  % Declares the document's title.
\author{Leslie Lamport}      % Declares the author's name.
\date{January 21, 1994}      % Deleting this command produces today's date.

\newcommand{\ip}[2]{(#1, #2)}
                             % Defines \ip{arg1}{arg2} to mean
                             % (arg1, arg2).

%\newcommand{\ip}[2]{\langle #1 | #2\rangle}
                             % This is an alternative definition of
                             % \ip that is commented out.

\begin{document}             % End of preamble and beginning of text.

\maketitle                   % Produces the title.

This is an example input file.  Comparing it with
the output it generates can show you how to
produce a simple document of your own.

\section{Ordinary Text}      % Produces section heading.  Lower-level
                             % sections are begun with similar
                             % \subsection and \subsubsection commands.

The ends  of words and sentences are marked
  by   spaces. It  doesn't matter how many
spaces    you type; one is as good as 100.  The
end of   a line counts as a space.

One   or more   blank lines denote the  end
of  a paragraph.

Since any number of consecutive spaces are treated
like a single one, the formatting of the input
file makes no difference to
      \LaTeX,                % The \LaTeX command generates the LaTeX logo.
but it makes a difference to you.  When you use
\LaTeX, making your input file as easy to read
as possible will be a great help as you write
your document and when you change it.  This sample
file shows how you can add comments to your own input
file.

Because printing is different from typewriting,
there are a number of things that you have to do
differently when preparing an input file than if
you were just typing the document directly.
Quotation marks like
       ``this''
have to be handled specially, as do quotes within
quotes:
       ``\,`this'            % \, separates the double and single quote.
        is what I just
        wrote, not  `that'\,''.

Dashes come in three sizes: an
       intra-word
dash, a medium dash for number ranges like
       1--2,
and a punctuation
       dash---like
this.

A sentence-ending space should be larger than the
space between words within a sentence.  You
sometimes have to type special commands in
conjunction with punctuation characters to get
this right, as in the following sentence.
       Gnats, gnus, etc.\ all  % `\ ' makes an inter-word space.
       begin with G\@.         % \@ marks end-of-sentence punctuation.
You should check the spaces after periods when
reading your output to make sure you haven't
forgotten any special cases.  Generating an
ellipsis
       \ldots\               % `\ ' is needed after `\ldots' because TeX
                             % ignores spaces after command names like \ldots
                             % made from \ + letters.
                             %
                             % Note how a `%' character causes TeX to ignore
                             % the end of the input line, so these blank lines
                             % do not start a new paragraph.
                             %
with the right spacing around the periods requires
a special command.

\LaTeX\ interprets some common characters as
commands, so you must type special commands to
generate them.  These characters include the
following:
       \$ \& \% \# \{ and \}.

In printing, text is usually emphasized with an
       \emph{italic}
type style.

\begin{em}
   A long segment of text can also be emphasized
   in this way.  Text within such a segment can be
   given \emph{additional} emphasis.
\end{em}

It is sometimes necessary to prevent \LaTeX\ from
breaking a line where it might otherwise do so.
This may be at a space, as between the ``Mr.''\ and
``Jones'' in
       ``Mr.~Jones'',        % ~ produces an unbreakable interword space.
or within a word---especially when the word is a
symbol like
       \mbox{\emph{itemnum}}
that makes little sense when hyphenated across
lines.

Footnotes\footnote{This is an example of a footnote.}
pose no problem.

\LaTeX\ is good at typesetting mathematical formulas
like
       \( x-3y + z = 7 \)
or
       \( a_{1} > x^{2n} + y^{2n} > x' \)
or
       \( \ip{A}{B} = \sum_{i} a_{i} b_{i} \).
The spaces you type in a formula are
ignored.  Remember that a letter like
       $x$                   % $ ... $  and  \( ... \)  are equivalent
is a formula when it denotes a mathematical
symbol, and it should be typed as one.

\section{Displayed Text}

Text is displayed by indenting it from the left
margin.  Quotations are commonly displayed.  There
are short quotations
\begin{quote}
   This is a short quotation.  It consists of a
   single paragraph of text.  See how it is formatted.
\end{quote}
and longer ones.
\begin{quotation}
   This is a longer quotation.  It consists of two
   paragraphs of text, neither of which are
   particularly interesting.

   This is the second paragraph of the quotation.  It
   is just as dull as the first paragraph.
\end{quotation}
Another frequently-displayed structure is a list.
The following is an example of an \emph{itemized}
list.
\begin{itemize}
   \item This is the first item of an itemized list.
         Each item in the list is marked with a ``tick''.
         You don't have to worry about what kind of tick
         mark is used.

   \item This is the second item of the list.  It
         contains another list nested inside it.  The inner
         list is an \emph{enumerated} list.
         \begin{enumerate}
            \item This is the first item of an enumerated
                  list that is nested within the itemized list.

            \item This is the second item of the inner list.
                  \LaTeX\ allows you to nest lists deeper than
                  you really should.
         \end{enumerate}
         This is the rest of the second item of the outer
         list.  It is no more interesting than any other
         part of the item.
   \item This is the third item of the list.
\end{itemize}
You can even display poetry.
\begin{verse}
   There is an environment
    for verse \\             % The \\ command separates lines
   Whose features some poets % within a stanza.
   will curse.

                             % One or more blank lines separate stanzas.

   For instead of making\\
   Them do \emph{all} line breaking, \\
   It allows them to put too many words on a line when they'd rather be
   forced to be terse.
\end{verse}

Mathematical formulas may also be displayed.  A
displayed formula
is
one-line long; multiline
formulas require special formatting instructions.
   \[  \ip{\Gamma}{\psi'} = x'' + y^{2} + z_{i}^{n}\]
Don't start a paragraph with a displayed equation,
nor make one a paragraph by itself.

\end{document}               % End of document.

\section{Bench Huge Section 169}
% synthetic bench section 169
% This is a sample LaTeX input file.  (Version of 12 August 2004.)
%
% A '%' character causes TeX to ignore all remaining text on the line,
% and is used for comments like this one.

\documentclass{article}      % Specifies the document class

                             % The preamble begins here.
\title{An Example Document}  % Declares the document's title.
\author{Leslie Lamport}      % Declares the author's name.
\date{January 21, 1994}      % Deleting this command produces today's date.

\newcommand{\ip}[2]{(#1, #2)}
                             % Defines \ip{arg1}{arg2} to mean
                             % (arg1, arg2).

%\newcommand{\ip}[2]{\langle #1 | #2\rangle}
                             % This is an alternative definition of
                             % \ip that is commented out.

\begin{document}             % End of preamble and beginning of text.

\maketitle                   % Produces the title.

This is an example input file.  Comparing it with
the output it generates can show you how to
produce a simple document of your own.

\section{Ordinary Text}      % Produces section heading.  Lower-level
                             % sections are begun with similar
                             % \subsection and \subsubsection commands.

The ends  of words and sentences are marked
  by   spaces. It  doesn't matter how many
spaces    you type; one is as good as 100.  The
end of   a line counts as a space.

One   or more   blank lines denote the  end
of  a paragraph.

Since any number of consecutive spaces are treated
like a single one, the formatting of the input
file makes no difference to
      \LaTeX,                % The \LaTeX command generates the LaTeX logo.
but it makes a difference to you.  When you use
\LaTeX, making your input file as easy to read
as possible will be a great help as you write
your document and when you change it.  This sample
file shows how you can add comments to your own input
file.

Because printing is different from typewriting,
there are a number of things that you have to do
differently when preparing an input file than if
you were just typing the document directly.
Quotation marks like
       ``this''
have to be handled specially, as do quotes within
quotes:
       ``\,`this'            % \, separates the double and single quote.
        is what I just
        wrote, not  `that'\,''.

Dashes come in three sizes: an
       intra-word
dash, a medium dash for number ranges like
       1--2,
and a punctuation
       dash---like
this.

A sentence-ending space should be larger than the
space between words within a sentence.  You
sometimes have to type special commands in
conjunction with punctuation characters to get
this right, as in the following sentence.
       Gnats, gnus, etc.\ all  % `\ ' makes an inter-word space.
       begin with G\@.         % \@ marks end-of-sentence punctuation.
You should check the spaces after periods when
reading your output to make sure you haven't
forgotten any special cases.  Generating an
ellipsis
       \ldots\               % `\ ' is needed after `\ldots' because TeX
                             % ignores spaces after command names like \ldots
                             % made from \ + letters.
                             %
                             % Note how a `%' character causes TeX to ignore
                             % the end of the input line, so these blank lines
                             % do not start a new paragraph.
                             %
with the right spacing around the periods requires
a special command.

\LaTeX\ interprets some common characters as
commands, so you must type special commands to
generate them.  These characters include the
following:
       \$ \& \% \# \{ and \}.

In printing, text is usually emphasized with an
       \emph{italic}
type style.

\begin{em}
   A long segment of text can also be emphasized
   in this way.  Text within such a segment can be
   given \emph{additional} emphasis.
\end{em}

It is sometimes necessary to prevent \LaTeX\ from
breaking a line where it might otherwise do so.
This may be at a space, as between the ``Mr.''\ and
``Jones'' in
       ``Mr.~Jones'',        % ~ produces an unbreakable interword space.
or within a word---especially when the word is a
symbol like
       \mbox{\emph{itemnum}}
that makes little sense when hyphenated across
lines.

Footnotes\footnote{This is an example of a footnote.}
pose no problem.

\LaTeX\ is good at typesetting mathematical formulas
like
       \( x-3y + z = 7 \)
or
       \( a_{1} > x^{2n} + y^{2n} > x' \)
or
       \( \ip{A}{B} = \sum_{i} a_{i} b_{i} \).
The spaces you type in a formula are
ignored.  Remember that a letter like
       $x$                   % $ ... $  and  \( ... \)  are equivalent
is a formula when it denotes a mathematical
symbol, and it should be typed as one.

\section{Displayed Text}

Text is displayed by indenting it from the left
margin.  Quotations are commonly displayed.  There
are short quotations
\begin{quote}
   This is a short quotation.  It consists of a
   single paragraph of text.  See how it is formatted.
\end{quote}
and longer ones.
\begin{quotation}
   This is a longer quotation.  It consists of two
   paragraphs of text, neither of which are
   particularly interesting.

   This is the second paragraph of the quotation.  It
   is just as dull as the first paragraph.
\end{quotation}
Another frequently-displayed structure is a list.
The following is an example of an \emph{itemized}
list.
\begin{itemize}
   \item This is the first item of an itemized list.
         Each item in the list is marked with a ``tick''.
         You don't have to worry about what kind of tick
         mark is used.

   \item This is the second item of the list.  It
         contains another list nested inside it.  The inner
         list is an \emph{enumerated} list.
         \begin{enumerate}
            \item This is the first item of an enumerated
                  list that is nested within the itemized list.

            \item This is the second item of the inner list.
                  \LaTeX\ allows you to nest lists deeper than
                  you really should.
         \end{enumerate}
         This is the rest of the second item of the outer
         list.  It is no more interesting than any other
         part of the item.
   \item This is the third item of the list.
\end{itemize}
You can even display poetry.
\begin{verse}
   There is an environment
    for verse \\             % The \\ command separates lines
   Whose features some poets % within a stanza.
   will curse.

                             % One or more blank lines separate stanzas.

   For instead of making\\
   Them do \emph{all} line breaking, \\
   It allows them to put too many words on a line when they'd rather be
   forced to be terse.
\end{verse}

Mathematical formulas may also be displayed.  A
displayed formula
is
one-line long; multiline
formulas require special formatting instructions.
   \[  \ip{\Gamma}{\psi'} = x'' + y^{2} + z_{i}^{n}\]
Don't start a paragraph with a displayed equation,
nor make one a paragraph by itself.

\end{document}               % End of document.

\section{Bench Huge Section 170}
% synthetic bench section 170
% This is a sample LaTeX input file.  (Version of 12 August 2004.)
%
% A '%' character causes TeX to ignore all remaining text on the line,
% and is used for comments like this one.

\documentclass{article}      % Specifies the document class

                             % The preamble begins here.
\title{An Example Document}  % Declares the document's title.
\author{Leslie Lamport}      % Declares the author's name.
\date{January 21, 1994}      % Deleting this command produces today's date.

\newcommand{\ip}[2]{(#1, #2)}
                             % Defines \ip{arg1}{arg2} to mean
                             % (arg1, arg2).

%\newcommand{\ip}[2]{\langle #1 | #2\rangle}
                             % This is an alternative definition of
                             % \ip that is commented out.

\begin{document}             % End of preamble and beginning of text.

\maketitle                   % Produces the title.

This is an example input file.  Comparing it with
the output it generates can show you how to
produce a simple document of your own.

\section{Ordinary Text}      % Produces section heading.  Lower-level
                             % sections are begun with similar
                             % \subsection and \subsubsection commands.

The ends  of words and sentences are marked
  by   spaces. It  doesn't matter how many
spaces    you type; one is as good as 100.  The
end of   a line counts as a space.

One   or more   blank lines denote the  end
of  a paragraph.

Since any number of consecutive spaces are treated
like a single one, the formatting of the input
file makes no difference to
      \LaTeX,                % The \LaTeX command generates the LaTeX logo.
but it makes a difference to you.  When you use
\LaTeX, making your input file as easy to read
as possible will be a great help as you write
your document and when you change it.  This sample
file shows how you can add comments to your own input
file.

Because printing is different from typewriting,
there are a number of things that you have to do
differently when preparing an input file than if
you were just typing the document directly.
Quotation marks like
       ``this''
have to be handled specially, as do quotes within
quotes:
       ``\,`this'            % \, separates the double and single quote.
        is what I just
        wrote, not  `that'\,''.

Dashes come in three sizes: an
       intra-word
dash, a medium dash for number ranges like
       1--2,
and a punctuation
       dash---like
this.

A sentence-ending space should be larger than the
space between words within a sentence.  You
sometimes have to type special commands in
conjunction with punctuation characters to get
this right, as in the following sentence.
       Gnats, gnus, etc.\ all  % `\ ' makes an inter-word space.
       begin with G\@.         % \@ marks end-of-sentence punctuation.
You should check the spaces after periods when
reading your output to make sure you haven't
forgotten any special cases.  Generating an
ellipsis
       \ldots\               % `\ ' is needed after `\ldots' because TeX
                             % ignores spaces after command names like \ldots
                             % made from \ + letters.
                             %
                             % Note how a `%' character causes TeX to ignore
                             % the end of the input line, so these blank lines
                             % do not start a new paragraph.
                             %
with the right spacing around the periods requires
a special command.

\LaTeX\ interprets some common characters as
commands, so you must type special commands to
generate them.  These characters include the
following:
       \$ \& \% \# \{ and \}.

In printing, text is usually emphasized with an
       \emph{italic}
type style.

\begin{em}
   A long segment of text can also be emphasized
   in this way.  Text within such a segment can be
   given \emph{additional} emphasis.
\end{em}

It is sometimes necessary to prevent \LaTeX\ from
breaking a line where it might otherwise do so.
This may be at a space, as between the ``Mr.''\ and
``Jones'' in
       ``Mr.~Jones'',        % ~ produces an unbreakable interword space.
or within a word---especially when the word is a
symbol like
       \mbox{\emph{itemnum}}
that makes little sense when hyphenated across
lines.

Footnotes\footnote{This is an example of a footnote.}
pose no problem.

\LaTeX\ is good at typesetting mathematical formulas
like
       \( x-3y + z = 7 \)
or
       \( a_{1} > x^{2n} + y^{2n} > x' \)
or
       \( \ip{A}{B} = \sum_{i} a_{i} b_{i} \).
The spaces you type in a formula are
ignored.  Remember that a letter like
       $x$                   % $ ... $  and  \( ... \)  are equivalent
is a formula when it denotes a mathematical
symbol, and it should be typed as one.

\section{Displayed Text}

Text is displayed by indenting it from the left
margin.  Quotations are commonly displayed.  There
are short quotations
\begin{quote}
   This is a short quotation.  It consists of a
   single paragraph of text.  See how it is formatted.
\end{quote}
and longer ones.
\begin{quotation}
   This is a longer quotation.  It consists of two
   paragraphs of text, neither of which are
   particularly interesting.

   This is the second paragraph of the quotation.  It
   is just as dull as the first paragraph.
\end{quotation}
Another frequently-displayed structure is a list.
The following is an example of an \emph{itemized}
list.
\begin{itemize}
   \item This is the first item of an itemized list.
         Each item in the list is marked with a ``tick''.
         You don't have to worry about what kind of tick
         mark is used.

   \item This is the second item of the list.  It
         contains another list nested inside it.  The inner
         list is an \emph{enumerated} list.
         \begin{enumerate}
            \item This is the first item of an enumerated
                  list that is nested within the itemized list.

            \item This is the second item of the inner list.
                  \LaTeX\ allows you to nest lists deeper than
                  you really should.
         \end{enumerate}
         This is the rest of the second item of the outer
         list.  It is no more interesting than any other
         part of the item.
   \item This is the third item of the list.
\end{itemize}
You can even display poetry.
\begin{verse}
   There is an environment
    for verse \\             % The \\ command separates lines
   Whose features some poets % within a stanza.
   will curse.

                             % One or more blank lines separate stanzas.

   For instead of making\\
   Them do \emph{all} line breaking, \\
   It allows them to put too many words on a line when they'd rather be
   forced to be terse.
\end{verse}

Mathematical formulas may also be displayed.  A
displayed formula
is
one-line long; multiline
formulas require special formatting instructions.
   \[  \ip{\Gamma}{\psi'} = x'' + y^{2} + z_{i}^{n}\]
Don't start a paragraph with a displayed equation,
nor make one a paragraph by itself.

\end{document}               % End of document.

\section{Bench Huge Section 171}
% synthetic bench section 171
% This is a sample LaTeX input file.  (Version of 12 August 2004.)
%
% A '%' character causes TeX to ignore all remaining text on the line,
% and is used for comments like this one.

\documentclass{article}      % Specifies the document class

                             % The preamble begins here.
\title{An Example Document}  % Declares the document's title.
\author{Leslie Lamport}      % Declares the author's name.
\date{January 21, 1994}      % Deleting this command produces today's date.

\newcommand{\ip}[2]{(#1, #2)}
                             % Defines \ip{arg1}{arg2} to mean
                             % (arg1, arg2).

%\newcommand{\ip}[2]{\langle #1 | #2\rangle}
                             % This is an alternative definition of
                             % \ip that is commented out.

\begin{document}             % End of preamble and beginning of text.

\maketitle                   % Produces the title.

This is an example input file.  Comparing it with
the output it generates can show you how to
produce a simple document of your own.

\section{Ordinary Text}      % Produces section heading.  Lower-level
                             % sections are begun with similar
                             % \subsection and \subsubsection commands.

The ends  of words and sentences are marked
  by   spaces. It  doesn't matter how many
spaces    you type; one is as good as 100.  The
end of   a line counts as a space.

One   or more   blank lines denote the  end
of  a paragraph.

Since any number of consecutive spaces are treated
like a single one, the formatting of the input
file makes no difference to
      \LaTeX,                % The \LaTeX command generates the LaTeX logo.
but it makes a difference to you.  When you use
\LaTeX, making your input file as easy to read
as possible will be a great help as you write
your document and when you change it.  This sample
file shows how you can add comments to your own input
file.

Because printing is different from typewriting,
there are a number of things that you have to do
differently when preparing an input file than if
you were just typing the document directly.
Quotation marks like
       ``this''
have to be handled specially, as do quotes within
quotes:
       ``\,`this'            % \, separates the double and single quote.
        is what I just
        wrote, not  `that'\,''.

Dashes come in three sizes: an
       intra-word
dash, a medium dash for number ranges like
       1--2,
and a punctuation
       dash---like
this.

A sentence-ending space should be larger than the
space between words within a sentence.  You
sometimes have to type special commands in
conjunction with punctuation characters to get
this right, as in the following sentence.
       Gnats, gnus, etc.\ all  % `\ ' makes an inter-word space.
       begin with G\@.         % \@ marks end-of-sentence punctuation.
You should check the spaces after periods when
reading your output to make sure you haven't
forgotten any special cases.  Generating an
ellipsis
       \ldots\               % `\ ' is needed after `\ldots' because TeX
                             % ignores spaces after command names like \ldots
                             % made from \ + letters.
                             %
                             % Note how a `%' character causes TeX to ignore
                             % the end of the input line, so these blank lines
                             % do not start a new paragraph.
                             %
with the right spacing around the periods requires
a special command.

\LaTeX\ interprets some common characters as
commands, so you must type special commands to
generate them.  These characters include the
following:
       \$ \& \% \# \{ and \}.

In printing, text is usually emphasized with an
       \emph{italic}
type style.

\begin{em}
   A long segment of text can also be emphasized
   in this way.  Text within such a segment can be
   given \emph{additional} emphasis.
\end{em}

It is sometimes necessary to prevent \LaTeX\ from
breaking a line where it might otherwise do so.
This may be at a space, as between the ``Mr.''\ and
``Jones'' in
       ``Mr.~Jones'',        % ~ produces an unbreakable interword space.
or within a word---especially when the word is a
symbol like
       \mbox{\emph{itemnum}}
that makes little sense when hyphenated across
lines.

Footnotes\footnote{This is an example of a footnote.}
pose no problem.

\LaTeX\ is good at typesetting mathematical formulas
like
       \( x-3y + z = 7 \)
or
       \( a_{1} > x^{2n} + y^{2n} > x' \)
or
       \( \ip{A}{B} = \sum_{i} a_{i} b_{i} \).
The spaces you type in a formula are
ignored.  Remember that a letter like
       $x$                   % $ ... $  and  \( ... \)  are equivalent
is a formula when it denotes a mathematical
symbol, and it should be typed as one.

\section{Displayed Text}

Text is displayed by indenting it from the left
margin.  Quotations are commonly displayed.  There
are short quotations
\begin{quote}
   This is a short quotation.  It consists of a
   single paragraph of text.  See how it is formatted.
\end{quote}
and longer ones.
\begin{quotation}
   This is a longer quotation.  It consists of two
   paragraphs of text, neither of which are
   particularly interesting.

   This is the second paragraph of the quotation.  It
   is just as dull as the first paragraph.
\end{quotation}
Another frequently-displayed structure is a list.
The following is an example of an \emph{itemized}
list.
\begin{itemize}
   \item This is the first item of an itemized list.
         Each item in the list is marked with a ``tick''.
         You don't have to worry about what kind of tick
         mark is used.

   \item This is the second item of the list.  It
         contains another list nested inside it.  The inner
         list is an \emph{enumerated} list.
         \begin{enumerate}
            \item This is the first item of an enumerated
                  list that is nested within the itemized list.

            \item This is the second item of the inner list.
                  \LaTeX\ allows you to nest lists deeper than
                  you really should.
         \end{enumerate}
         This is the rest of the second item of the outer
         list.  It is no more interesting than any other
         part of the item.
   \item This is the third item of the list.
\end{itemize}
You can even display poetry.
\begin{verse}
   There is an environment
    for verse \\             % The \\ command separates lines
   Whose features some poets % within a stanza.
   will curse.

                             % One or more blank lines separate stanzas.

   For instead of making\\
   Them do \emph{all} line breaking, \\
   It allows them to put too many words on a line when they'd rather be
   forced to be terse.
\end{verse}

Mathematical formulas may also be displayed.  A
displayed formula
is
one-line long; multiline
formulas require special formatting instructions.
   \[  \ip{\Gamma}{\psi'} = x'' + y^{2} + z_{i}^{n}\]
Don't start a paragraph with a displayed equation,
nor make one a paragraph by itself.

\end{document}               % End of document.

\section{Bench Huge Section 172}
% synthetic bench section 172
% This is a sample LaTeX input file.  (Version of 12 August 2004.)
%
% A '%' character causes TeX to ignore all remaining text on the line,
% and is used for comments like this one.

\documentclass{article}      % Specifies the document class

                             % The preamble begins here.
\title{An Example Document}  % Declares the document's title.
\author{Leslie Lamport}      % Declares the author's name.
\date{January 21, 1994}      % Deleting this command produces today's date.

\newcommand{\ip}[2]{(#1, #2)}
                             % Defines \ip{arg1}{arg2} to mean
                             % (arg1, arg2).

%\newcommand{\ip}[2]{\langle #1 | #2\rangle}
                             % This is an alternative definition of
                             % \ip that is commented out.

\begin{document}             % End of preamble and beginning of text.

\maketitle                   % Produces the title.

This is an example input file.  Comparing it with
the output it generates can show you how to
produce a simple document of your own.

\section{Ordinary Text}      % Produces section heading.  Lower-level
                             % sections are begun with similar
                             % \subsection and \subsubsection commands.

The ends  of words and sentences are marked
  by   spaces. It  doesn't matter how many
spaces    you type; one is as good as 100.  The
end of   a line counts as a space.

One   or more   blank lines denote the  end
of  a paragraph.

Since any number of consecutive spaces are treated
like a single one, the formatting of the input
file makes no difference to
      \LaTeX,                % The \LaTeX command generates the LaTeX logo.
but it makes a difference to you.  When you use
\LaTeX, making your input file as easy to read
as possible will be a great help as you write
your document and when you change it.  This sample
file shows how you can add comments to your own input
file.

Because printing is different from typewriting,
there are a number of things that you have to do
differently when preparing an input file than if
you were just typing the document directly.
Quotation marks like
       ``this''
have to be handled specially, as do quotes within
quotes:
       ``\,`this'            % \, separates the double and single quote.
        is what I just
        wrote, not  `that'\,''.

Dashes come in three sizes: an
       intra-word
dash, a medium dash for number ranges like
       1--2,
and a punctuation
       dash---like
this.

A sentence-ending space should be larger than the
space between words within a sentence.  You
sometimes have to type special commands in
conjunction with punctuation characters to get
this right, as in the following sentence.
       Gnats, gnus, etc.\ all  % `\ ' makes an inter-word space.
       begin with G\@.         % \@ marks end-of-sentence punctuation.
You should check the spaces after periods when
reading your output to make sure you haven't
forgotten any special cases.  Generating an
ellipsis
       \ldots\               % `\ ' is needed after `\ldots' because TeX
                             % ignores spaces after command names like \ldots
                             % made from \ + letters.
                             %
                             % Note how a `%' character causes TeX to ignore
                             % the end of the input line, so these blank lines
                             % do not start a new paragraph.
                             %
with the right spacing around the periods requires
a special command.

\LaTeX\ interprets some common characters as
commands, so you must type special commands to
generate them.  These characters include the
following:
       \$ \& \% \# \{ and \}.

In printing, text is usually emphasized with an
       \emph{italic}
type style.

\begin{em}
   A long segment of text can also be emphasized
   in this way.  Text within such a segment can be
   given \emph{additional} emphasis.
\end{em}

It is sometimes necessary to prevent \LaTeX\ from
breaking a line where it might otherwise do so.
This may be at a space, as between the ``Mr.''\ and
``Jones'' in
       ``Mr.~Jones'',        % ~ produces an unbreakable interword space.
or within a word---especially when the word is a
symbol like
       \mbox{\emph{itemnum}}
that makes little sense when hyphenated across
lines.

Footnotes\footnote{This is an example of a footnote.}
pose no problem.

\LaTeX\ is good at typesetting mathematical formulas
like
       \( x-3y + z = 7 \)
or
       \( a_{1} > x^{2n} + y^{2n} > x' \)
or
       \( \ip{A}{B} = \sum_{i} a_{i} b_{i} \).
The spaces you type in a formula are
ignored.  Remember that a letter like
       $x$                   % $ ... $  and  \( ... \)  are equivalent
is a formula when it denotes a mathematical
symbol, and it should be typed as one.

\section{Displayed Text}

Text is displayed by indenting it from the left
margin.  Quotations are commonly displayed.  There
are short quotations
\begin{quote}
   This is a short quotation.  It consists of a
   single paragraph of text.  See how it is formatted.
\end{quote}
and longer ones.
\begin{quotation}
   This is a longer quotation.  It consists of two
   paragraphs of text, neither of which are
   particularly interesting.

   This is the second paragraph of the quotation.  It
   is just as dull as the first paragraph.
\end{quotation}
Another frequently-displayed structure is a list.
The following is an example of an \emph{itemized}
list.
\begin{itemize}
   \item This is the first item of an itemized list.
         Each item in the list is marked with a ``tick''.
         You don't have to worry about what kind of tick
         mark is used.

   \item This is the second item of the list.  It
         contains another list nested inside it.  The inner
         list is an \emph{enumerated} list.
         \begin{enumerate}
            \item This is the first item of an enumerated
                  list that is nested within the itemized list.

            \item This is the second item of the inner list.
                  \LaTeX\ allows you to nest lists deeper than
                  you really should.
         \end{enumerate}
         This is the rest of the second item of the outer
         list.  It is no more interesting than any other
         part of the item.
   \item This is the third item of the list.
\end{itemize}
You can even display poetry.
\begin{verse}
   There is an environment
    for verse \\             % The \\ command separates lines
   Whose features some poets % within a stanza.
   will curse.

                             % One or more blank lines separate stanzas.

   For instead of making\\
   Them do \emph{all} line breaking, \\
   It allows them to put too many words on a line when they'd rather be
   forced to be terse.
\end{verse}

Mathematical formulas may also be displayed.  A
displayed formula
is
one-line long; multiline
formulas require special formatting instructions.
   \[  \ip{\Gamma}{\psi'} = x'' + y^{2} + z_{i}^{n}\]
Don't start a paragraph with a displayed equation,
nor make one a paragraph by itself.

\end{document}               % End of document.

\section{Bench Huge Section 173}
% synthetic bench section 173
% This is a sample LaTeX input file.  (Version of 12 August 2004.)
%
% A '%' character causes TeX to ignore all remaining text on the line,
% and is used for comments like this one.

\documentclass{article}      % Specifies the document class

                             % The preamble begins here.
\title{An Example Document}  % Declares the document's title.
\author{Leslie Lamport}      % Declares the author's name.
\date{January 21, 1994}      % Deleting this command produces today's date.

\newcommand{\ip}[2]{(#1, #2)}
                             % Defines \ip{arg1}{arg2} to mean
                             % (arg1, arg2).

%\newcommand{\ip}[2]{\langle #1 | #2\rangle}
                             % This is an alternative definition of
                             % \ip that is commented out.

\begin{document}             % End of preamble and beginning of text.

\maketitle                   % Produces the title.

This is an example input file.  Comparing it with
the output it generates can show you how to
produce a simple document of your own.

\section{Ordinary Text}      % Produces section heading.  Lower-level
                             % sections are begun with similar
                             % \subsection and \subsubsection commands.

The ends  of words and sentences are marked
  by   spaces. It  doesn't matter how many
spaces    you type; one is as good as 100.  The
end of   a line counts as a space.

One   or more   blank lines denote the  end
of  a paragraph.

Since any number of consecutive spaces are treated
like a single one, the formatting of the input
file makes no difference to
      \LaTeX,                % The \LaTeX command generates the LaTeX logo.
but it makes a difference to you.  When you use
\LaTeX, making your input file as easy to read
as possible will be a great help as you write
your document and when you change it.  This sample
file shows how you can add comments to your own input
file.

Because printing is different from typewriting,
there are a number of things that you have to do
differently when preparing an input file than if
you were just typing the document directly.
Quotation marks like
       ``this''
have to be handled specially, as do quotes within
quotes:
       ``\,`this'            % \, separates the double and single quote.
        is what I just
        wrote, not  `that'\,''.

Dashes come in three sizes: an
       intra-word
dash, a medium dash for number ranges like
       1--2,
and a punctuation
       dash---like
this.

A sentence-ending space should be larger than the
space between words within a sentence.  You
sometimes have to type special commands in
conjunction with punctuation characters to get
this right, as in the following sentence.
       Gnats, gnus, etc.\ all  % `\ ' makes an inter-word space.
       begin with G\@.         % \@ marks end-of-sentence punctuation.
You should check the spaces after periods when
reading your output to make sure you haven't
forgotten any special cases.  Generating an
ellipsis
       \ldots\               % `\ ' is needed after `\ldots' because TeX
                             % ignores spaces after command names like \ldots
                             % made from \ + letters.
                             %
                             % Note how a `%' character causes TeX to ignore
                             % the end of the input line, so these blank lines
                             % do not start a new paragraph.
                             %
with the right spacing around the periods requires
a special command.

\LaTeX\ interprets some common characters as
commands, so you must type special commands to
generate them.  These characters include the
following:
       \$ \& \% \# \{ and \}.

In printing, text is usually emphasized with an
       \emph{italic}
type style.

\begin{em}
   A long segment of text can also be emphasized
   in this way.  Text within such a segment can be
   given \emph{additional} emphasis.
\end{em}

It is sometimes necessary to prevent \LaTeX\ from
breaking a line where it might otherwise do so.
This may be at a space, as between the ``Mr.''\ and
``Jones'' in
       ``Mr.~Jones'',        % ~ produces an unbreakable interword space.
or within a word---especially when the word is a
symbol like
       \mbox{\emph{itemnum}}
that makes little sense when hyphenated across
lines.

Footnotes\footnote{This is an example of a footnote.}
pose no problem.

\LaTeX\ is good at typesetting mathematical formulas
like
       \( x-3y + z = 7 \)
or
       \( a_{1} > x^{2n} + y^{2n} > x' \)
or
       \( \ip{A}{B} = \sum_{i} a_{i} b_{i} \).
The spaces you type in a formula are
ignored.  Remember that a letter like
       $x$                   % $ ... $  and  \( ... \)  are equivalent
is a formula when it denotes a mathematical
symbol, and it should be typed as one.

\section{Displayed Text}

Text is displayed by indenting it from the left
margin.  Quotations are commonly displayed.  There
are short quotations
\begin{quote}
   This is a short quotation.  It consists of a
   single paragraph of text.  See how it is formatted.
\end{quote}
and longer ones.
\begin{quotation}
   This is a longer quotation.  It consists of two
   paragraphs of text, neither of which are
   particularly interesting.

   This is the second paragraph of the quotation.  It
   is just as dull as the first paragraph.
\end{quotation}
Another frequently-displayed structure is a list.
The following is an example of an \emph{itemized}
list.
\begin{itemize}
   \item This is the first item of an itemized list.
         Each item in the list is marked with a ``tick''.
         You don't have to worry about what kind of tick
         mark is used.

   \item This is the second item of the list.  It
         contains another list nested inside it.  The inner
         list is an \emph{enumerated} list.
         \begin{enumerate}
            \item This is the first item of an enumerated
                  list that is nested within the itemized list.

            \item This is the second item of the inner list.
                  \LaTeX\ allows you to nest lists deeper than
                  you really should.
         \end{enumerate}
         This is the rest of the second item of the outer
         list.  It is no more interesting than any other
         part of the item.
   \item This is the third item of the list.
\end{itemize}
You can even display poetry.
\begin{verse}
   There is an environment
    for verse \\             % The \\ command separates lines
   Whose features some poets % within a stanza.
   will curse.

                             % One or more blank lines separate stanzas.

   For instead of making\\
   Them do \emph{all} line breaking, \\
   It allows them to put too many words on a line when they'd rather be
   forced to be terse.
\end{verse}

Mathematical formulas may also be displayed.  A
displayed formula
is
one-line long; multiline
formulas require special formatting instructions.
   \[  \ip{\Gamma}{\psi'} = x'' + y^{2} + z_{i}^{n}\]
Don't start a paragraph with a displayed equation,
nor make one a paragraph by itself.

\end{document}               % End of document.

\section{Bench Huge Section 174}
% synthetic bench section 174
% This is a sample LaTeX input file.  (Version of 12 August 2004.)
%
% A '%' character causes TeX to ignore all remaining text on the line,
% and is used for comments like this one.

\documentclass{article}      % Specifies the document class

                             % The preamble begins here.
\title{An Example Document}  % Declares the document's title.
\author{Leslie Lamport}      % Declares the author's name.
\date{January 21, 1994}      % Deleting this command produces today's date.

\newcommand{\ip}[2]{(#1, #2)}
                             % Defines \ip{arg1}{arg2} to mean
                             % (arg1, arg2).

%\newcommand{\ip}[2]{\langle #1 | #2\rangle}
                             % This is an alternative definition of
                             % \ip that is commented out.

\begin{document}             % End of preamble and beginning of text.

\maketitle                   % Produces the title.

This is an example input file.  Comparing it with
the output it generates can show you how to
produce a simple document of your own.

\section{Ordinary Text}      % Produces section heading.  Lower-level
                             % sections are begun with similar
                             % \subsection and \subsubsection commands.

The ends  of words and sentences are marked
  by   spaces. It  doesn't matter how many
spaces    you type; one is as good as 100.  The
end of   a line counts as a space.

One   or more   blank lines denote the  end
of  a paragraph.

Since any number of consecutive spaces are treated
like a single one, the formatting of the input
file makes no difference to
      \LaTeX,                % The \LaTeX command generates the LaTeX logo.
but it makes a difference to you.  When you use
\LaTeX, making your input file as easy to read
as possible will be a great help as you write
your document and when you change it.  This sample
file shows how you can add comments to your own input
file.

Because printing is different from typewriting,
there are a number of things that you have to do
differently when preparing an input file than if
you were just typing the document directly.
Quotation marks like
       ``this''
have to be handled specially, as do quotes within
quotes:
       ``\,`this'            % \, separates the double and single quote.
        is what I just
        wrote, not  `that'\,''.

Dashes come in three sizes: an
       intra-word
dash, a medium dash for number ranges like
       1--2,
and a punctuation
       dash---like
this.

A sentence-ending space should be larger than the
space between words within a sentence.  You
sometimes have to type special commands in
conjunction with punctuation characters to get
this right, as in the following sentence.
       Gnats, gnus, etc.\ all  % `\ ' makes an inter-word space.
       begin with G\@.         % \@ marks end-of-sentence punctuation.
You should check the spaces after periods when
reading your output to make sure you haven't
forgotten any special cases.  Generating an
ellipsis
       \ldots\               % `\ ' is needed after `\ldots' because TeX
                             % ignores spaces after command names like \ldots
                             % made from \ + letters.
                             %
                             % Note how a `%' character causes TeX to ignore
                             % the end of the input line, so these blank lines
                             % do not start a new paragraph.
                             %
with the right spacing around the periods requires
a special command.

\LaTeX\ interprets some common characters as
commands, so you must type special commands to
generate them.  These characters include the
following:
       \$ \& \% \# \{ and \}.

In printing, text is usually emphasized with an
       \emph{italic}
type style.

\begin{em}
   A long segment of text can also be emphasized
   in this way.  Text within such a segment can be
   given \emph{additional} emphasis.
\end{em}

It is sometimes necessary to prevent \LaTeX\ from
breaking a line where it might otherwise do so.
This may be at a space, as between the ``Mr.''\ and
``Jones'' in
       ``Mr.~Jones'',        % ~ produces an unbreakable interword space.
or within a word---especially when the word is a
symbol like
       \mbox{\emph{itemnum}}
that makes little sense when hyphenated across
lines.

Footnotes\footnote{This is an example of a footnote.}
pose no problem.

\LaTeX\ is good at typesetting mathematical formulas
like
       \( x-3y + z = 7 \)
or
       \( a_{1} > x^{2n} + y^{2n} > x' \)
or
       \( \ip{A}{B} = \sum_{i} a_{i} b_{i} \).
The spaces you type in a formula are
ignored.  Remember that a letter like
       $x$                   % $ ... $  and  \( ... \)  are equivalent
is a formula when it denotes a mathematical
symbol, and it should be typed as one.

\section{Displayed Text}

Text is displayed by indenting it from the left
margin.  Quotations are commonly displayed.  There
are short quotations
\begin{quote}
   This is a short quotation.  It consists of a
   single paragraph of text.  See how it is formatted.
\end{quote}
and longer ones.
\begin{quotation}
   This is a longer quotation.  It consists of two
   paragraphs of text, neither of which are
   particularly interesting.

   This is the second paragraph of the quotation.  It
   is just as dull as the first paragraph.
\end{quotation}
Another frequently-displayed structure is a list.
The following is an example of an \emph{itemized}
list.
\begin{itemize}
   \item This is the first item of an itemized list.
         Each item in the list is marked with a ``tick''.
         You don't have to worry about what kind of tick
         mark is used.

   \item This is the second item of the list.  It
         contains another list nested inside it.  The inner
         list is an \emph{enumerated} list.
         \begin{enumerate}
            \item This is the first item of an enumerated
                  list that is nested within the itemized list.

            \item This is the second item of the inner list.
                  \LaTeX\ allows you to nest lists deeper than
                  you really should.
         \end{enumerate}
         This is the rest of the second item of the outer
         list.  It is no more interesting than any other
         part of the item.
   \item This is the third item of the list.
\end{itemize}
You can even display poetry.
\begin{verse}
   There is an environment
    for verse \\             % The \\ command separates lines
   Whose features some poets % within a stanza.
   will curse.

                             % One or more blank lines separate stanzas.

   For instead of making\\
   Them do \emph{all} line breaking, \\
   It allows them to put too many words on a line when they'd rather be
   forced to be terse.
\end{verse}

Mathematical formulas may also be displayed.  A
displayed formula
is
one-line long; multiline
formulas require special formatting instructions.
   \[  \ip{\Gamma}{\psi'} = x'' + y^{2} + z_{i}^{n}\]
Don't start a paragraph with a displayed equation,
nor make one a paragraph by itself.

\end{document}               % End of document.

\section{Bench Huge Section 175}
% synthetic bench section 175
% This is a sample LaTeX input file.  (Version of 12 August 2004.)
%
% A '%' character causes TeX to ignore all remaining text on the line,
% and is used for comments like this one.

\documentclass{article}      % Specifies the document class

                             % The preamble begins here.
\title{An Example Document}  % Declares the document's title.
\author{Leslie Lamport}      % Declares the author's name.
\date{January 21, 1994}      % Deleting this command produces today's date.

\newcommand{\ip}[2]{(#1, #2)}
                             % Defines \ip{arg1}{arg2} to mean
                             % (arg1, arg2).

%\newcommand{\ip}[2]{\langle #1 | #2\rangle}
                             % This is an alternative definition of
                             % \ip that is commented out.

\begin{document}             % End of preamble and beginning of text.

\maketitle                   % Produces the title.

This is an example input file.  Comparing it with
the output it generates can show you how to
produce a simple document of your own.

\section{Ordinary Text}      % Produces section heading.  Lower-level
                             % sections are begun with similar
                             % \subsection and \subsubsection commands.

The ends  of words and sentences are marked
  by   spaces. It  doesn't matter how many
spaces    you type; one is as good as 100.  The
end of   a line counts as a space.

One   or more   blank lines denote the  end
of  a paragraph.

Since any number of consecutive spaces are treated
like a single one, the formatting of the input
file makes no difference to
      \LaTeX,                % The \LaTeX command generates the LaTeX logo.
but it makes a difference to you.  When you use
\LaTeX, making your input file as easy to read
as possible will be a great help as you write
your document and when you change it.  This sample
file shows how you can add comments to your own input
file.

Because printing is different from typewriting,
there are a number of things that you have to do
differently when preparing an input file than if
you were just typing the document directly.
Quotation marks like
       ``this''
have to be handled specially, as do quotes within
quotes:
       ``\,`this'            % \, separates the double and single quote.
        is what I just
        wrote, not  `that'\,''.

Dashes come in three sizes: an
       intra-word
dash, a medium dash for number ranges like
       1--2,
and a punctuation
       dash---like
this.

A sentence-ending space should be larger than the
space between words within a sentence.  You
sometimes have to type special commands in
conjunction with punctuation characters to get
this right, as in the following sentence.
       Gnats, gnus, etc.\ all  % `\ ' makes an inter-word space.
       begin with G\@.         % \@ marks end-of-sentence punctuation.
You should check the spaces after periods when
reading your output to make sure you haven't
forgotten any special cases.  Generating an
ellipsis
       \ldots\               % `\ ' is needed after `\ldots' because TeX
                             % ignores spaces after command names like \ldots
                             % made from \ + letters.
                             %
                             % Note how a `%' character causes TeX to ignore
                             % the end of the input line, so these blank lines
                             % do not start a new paragraph.
                             %
with the right spacing around the periods requires
a special command.

\LaTeX\ interprets some common characters as
commands, so you must type special commands to
generate them.  These characters include the
following:
       \$ \& \% \# \{ and \}.

In printing, text is usually emphasized with an
       \emph{italic}
type style.

\begin{em}
   A long segment of text can also be emphasized
   in this way.  Text within such a segment can be
   given \emph{additional} emphasis.
\end{em}

It is sometimes necessary to prevent \LaTeX\ from
breaking a line where it might otherwise do so.
This may be at a space, as between the ``Mr.''\ and
``Jones'' in
       ``Mr.~Jones'',        % ~ produces an unbreakable interword space.
or within a word---especially when the word is a
symbol like
       \mbox{\emph{itemnum}}
that makes little sense when hyphenated across
lines.

Footnotes\footnote{This is an example of a footnote.}
pose no problem.

\LaTeX\ is good at typesetting mathematical formulas
like
       \( x-3y + z = 7 \)
or
       \( a_{1} > x^{2n} + y^{2n} > x' \)
or
       \( \ip{A}{B} = \sum_{i} a_{i} b_{i} \).
The spaces you type in a formula are
ignored.  Remember that a letter like
       $x$                   % $ ... $  and  \( ... \)  are equivalent
is a formula when it denotes a mathematical
symbol, and it should be typed as one.

\section{Displayed Text}

Text is displayed by indenting it from the left
margin.  Quotations are commonly displayed.  There
are short quotations
\begin{quote}
   This is a short quotation.  It consists of a
   single paragraph of text.  See how it is formatted.
\end{quote}
and longer ones.
\begin{quotation}
   This is a longer quotation.  It consists of two
   paragraphs of text, neither of which are
   particularly interesting.

   This is the second paragraph of the quotation.  It
   is just as dull as the first paragraph.
\end{quotation}
Another frequently-displayed structure is a list.
The following is an example of an \emph{itemized}
list.
\begin{itemize}
   \item This is the first item of an itemized list.
         Each item in the list is marked with a ``tick''.
         You don't have to worry about what kind of tick
         mark is used.

   \item This is the second item of the list.  It
         contains another list nested inside it.  The inner
         list is an \emph{enumerated} list.
         \begin{enumerate}
            \item This is the first item of an enumerated
                  list that is nested within the itemized list.

            \item This is the second item of the inner list.
                  \LaTeX\ allows you to nest lists deeper than
                  you really should.
         \end{enumerate}
         This is the rest of the second item of the outer
         list.  It is no more interesting than any other
         part of the item.
   \item This is the third item of the list.
\end{itemize}
You can even display poetry.
\begin{verse}
   There is an environment
    for verse \\             % The \\ command separates lines
   Whose features some poets % within a stanza.
   will curse.

                             % One or more blank lines separate stanzas.

   For instead of making\\
   Them do \emph{all} line breaking, \\
   It allows them to put too many words on a line when they'd rather be
   forced to be terse.
\end{verse}

Mathematical formulas may also be displayed.  A
displayed formula
is
one-line long; multiline
formulas require special formatting instructions.
   \[  \ip{\Gamma}{\psi'} = x'' + y^{2} + z_{i}^{n}\]
Don't start a paragraph with a displayed equation,
nor make one a paragraph by itself.

\end{document}               % End of document.

\section{Bench Huge Section 176}
% synthetic bench section 176
% This is a sample LaTeX input file.  (Version of 12 August 2004.)
%
% A '%' character causes TeX to ignore all remaining text on the line,
% and is used for comments like this one.

\documentclass{article}      % Specifies the document class

                             % The preamble begins here.
\title{An Example Document}  % Declares the document's title.
\author{Leslie Lamport}      % Declares the author's name.
\date{January 21, 1994}      % Deleting this command produces today's date.

\newcommand{\ip}[2]{(#1, #2)}
                             % Defines \ip{arg1}{arg2} to mean
                             % (arg1, arg2).

%\newcommand{\ip}[2]{\langle #1 | #2\rangle}
                             % This is an alternative definition of
                             % \ip that is commented out.

\begin{document}             % End of preamble and beginning of text.

\maketitle                   % Produces the title.

This is an example input file.  Comparing it with
the output it generates can show you how to
produce a simple document of your own.

\section{Ordinary Text}      % Produces section heading.  Lower-level
                             % sections are begun with similar
                             % \subsection and \subsubsection commands.

The ends  of words and sentences are marked
  by   spaces. It  doesn't matter how many
spaces    you type; one is as good as 100.  The
end of   a line counts as a space.

One   or more   blank lines denote the  end
of  a paragraph.

Since any number of consecutive spaces are treated
like a single one, the formatting of the input
file makes no difference to
      \LaTeX,                % The \LaTeX command generates the LaTeX logo.
but it makes a difference to you.  When you use
\LaTeX, making your input file as easy to read
as possible will be a great help as you write
your document and when you change it.  This sample
file shows how you can add comments to your own input
file.

Because printing is different from typewriting,
there are a number of things that you have to do
differently when preparing an input file than if
you were just typing the document directly.
Quotation marks like
       ``this''
have to be handled specially, as do quotes within
quotes:
       ``\,`this'            % \, separates the double and single quote.
        is what I just
        wrote, not  `that'\,''.

Dashes come in three sizes: an
       intra-word
dash, a medium dash for number ranges like
       1--2,
and a punctuation
       dash---like
this.

A sentence-ending space should be larger than the
space between words within a sentence.  You
sometimes have to type special commands in
conjunction with punctuation characters to get
this right, as in the following sentence.
       Gnats, gnus, etc.\ all  % `\ ' makes an inter-word space.
       begin with G\@.         % \@ marks end-of-sentence punctuation.
You should check the spaces after periods when
reading your output to make sure you haven't
forgotten any special cases.  Generating an
ellipsis
       \ldots\               % `\ ' is needed after `\ldots' because TeX
                             % ignores spaces after command names like \ldots
                             % made from \ + letters.
                             %
                             % Note how a `%' character causes TeX to ignore
                             % the end of the input line, so these blank lines
                             % do not start a new paragraph.
                             %
with the right spacing around the periods requires
a special command.

\LaTeX\ interprets some common characters as
commands, so you must type special commands to
generate them.  These characters include the
following:
       \$ \& \% \# \{ and \}.

In printing, text is usually emphasized with an
       \emph{italic}
type style.

\begin{em}
   A long segment of text can also be emphasized
   in this way.  Text within such a segment can be
   given \emph{additional} emphasis.
\end{em}

It is sometimes necessary to prevent \LaTeX\ from
breaking a line where it might otherwise do so.
This may be at a space, as between the ``Mr.''\ and
``Jones'' in
       ``Mr.~Jones'',        % ~ produces an unbreakable interword space.
or within a word---especially when the word is a
symbol like
       \mbox{\emph{itemnum}}
that makes little sense when hyphenated across
lines.

Footnotes\footnote{This is an example of a footnote.}
pose no problem.

\LaTeX\ is good at typesetting mathematical formulas
like
       \( x-3y + z = 7 \)
or
       \( a_{1} > x^{2n} + y^{2n} > x' \)
or
       \( \ip{A}{B} = \sum_{i} a_{i} b_{i} \).
The spaces you type in a formula are
ignored.  Remember that a letter like
       $x$                   % $ ... $  and  \( ... \)  are equivalent
is a formula when it denotes a mathematical
symbol, and it should be typed as one.

\section{Displayed Text}

Text is displayed by indenting it from the left
margin.  Quotations are commonly displayed.  There
are short quotations
\begin{quote}
   This is a short quotation.  It consists of a
   single paragraph of text.  See how it is formatted.
\end{quote}
and longer ones.
\begin{quotation}
   This is a longer quotation.  It consists of two
   paragraphs of text, neither of which are
   particularly interesting.

   This is the second paragraph of the quotation.  It
   is just as dull as the first paragraph.
\end{quotation}
Another frequently-displayed structure is a list.
The following is an example of an \emph{itemized}
list.
\begin{itemize}
   \item This is the first item of an itemized list.
         Each item in the list is marked with a ``tick''.
         You don't have to worry about what kind of tick
         mark is used.

   \item This is the second item of the list.  It
         contains another list nested inside it.  The inner
         list is an \emph{enumerated} list.
         \begin{enumerate}
            \item This is the first item of an enumerated
                  list that is nested within the itemized list.

            \item This is the second item of the inner list.
                  \LaTeX\ allows you to nest lists deeper than
                  you really should.
         \end{enumerate}
         This is the rest of the second item of the outer
         list.  It is no more interesting than any other
         part of the item.
   \item This is the third item of the list.
\end{itemize}
You can even display poetry.
\begin{verse}
   There is an environment
    for verse \\             % The \\ command separates lines
   Whose features some poets % within a stanza.
   will curse.

                             % One or more blank lines separate stanzas.

   For instead of making\\
   Them do \emph{all} line breaking, \\
   It allows them to put too many words on a line when they'd rather be
   forced to be terse.
\end{verse}

Mathematical formulas may also be displayed.  A
displayed formula
is
one-line long; multiline
formulas require special formatting instructions.
   \[  \ip{\Gamma}{\psi'} = x'' + y^{2} + z_{i}^{n}\]
Don't start a paragraph with a displayed equation,
nor make one a paragraph by itself.

\end{document}               % End of document.

\section{Bench Huge Section 177}
% synthetic bench section 177
% This is a sample LaTeX input file.  (Version of 12 August 2004.)
%
% A '%' character causes TeX to ignore all remaining text on the line,
% and is used for comments like this one.

\documentclass{article}      % Specifies the document class

                             % The preamble begins here.
\title{An Example Document}  % Declares the document's title.
\author{Leslie Lamport}      % Declares the author's name.
\date{January 21, 1994}      % Deleting this command produces today's date.

\newcommand{\ip}[2]{(#1, #2)}
                             % Defines \ip{arg1}{arg2} to mean
                             % (arg1, arg2).

%\newcommand{\ip}[2]{\langle #1 | #2\rangle}
                             % This is an alternative definition of
                             % \ip that is commented out.

\begin{document}             % End of preamble and beginning of text.

\maketitle                   % Produces the title.

This is an example input file.  Comparing it with
the output it generates can show you how to
produce a simple document of your own.

\section{Ordinary Text}      % Produces section heading.  Lower-level
                             % sections are begun with similar
                             % \subsection and \subsubsection commands.

The ends  of words and sentences are marked
  by   spaces. It  doesn't matter how many
spaces    you type; one is as good as 100.  The
end of   a line counts as a space.

One   or more   blank lines denote the  end
of  a paragraph.

Since any number of consecutive spaces are treated
like a single one, the formatting of the input
file makes no difference to
      \LaTeX,                % The \LaTeX command generates the LaTeX logo.
but it makes a difference to you.  When you use
\LaTeX, making your input file as easy to read
as possible will be a great help as you write
your document and when you change it.  This sample
file shows how you can add comments to your own input
file.

Because printing is different from typewriting,
there are a number of things that you have to do
differently when preparing an input file than if
you were just typing the document directly.
Quotation marks like
       ``this''
have to be handled specially, as do quotes within
quotes:
       ``\,`this'            % \, separates the double and single quote.
        is what I just
        wrote, not  `that'\,''.

Dashes come in three sizes: an
       intra-word
dash, a medium dash for number ranges like
       1--2,
and a punctuation
       dash---like
this.

A sentence-ending space should be larger than the
space between words within a sentence.  You
sometimes have to type special commands in
conjunction with punctuation characters to get
this right, as in the following sentence.
       Gnats, gnus, etc.\ all  % `\ ' makes an inter-word space.
       begin with G\@.         % \@ marks end-of-sentence punctuation.
You should check the spaces after periods when
reading your output to make sure you haven't
forgotten any special cases.  Generating an
ellipsis
       \ldots\               % `\ ' is needed after `\ldots' because TeX
                             % ignores spaces after command names like \ldots
                             % made from \ + letters.
                             %
                             % Note how a `%' character causes TeX to ignore
                             % the end of the input line, so these blank lines
                             % do not start a new paragraph.
                             %
with the right spacing around the periods requires
a special command.

\LaTeX\ interprets some common characters as
commands, so you must type special commands to
generate them.  These characters include the
following:
       \$ \& \% \# \{ and \}.

In printing, text is usually emphasized with an
       \emph{italic}
type style.

\begin{em}
   A long segment of text can also be emphasized
   in this way.  Text within such a segment can be
   given \emph{additional} emphasis.
\end{em}

It is sometimes necessary to prevent \LaTeX\ from
breaking a line where it might otherwise do so.
This may be at a space, as between the ``Mr.''\ and
``Jones'' in
       ``Mr.~Jones'',        % ~ produces an unbreakable interword space.
or within a word---especially when the word is a
symbol like
       \mbox{\emph{itemnum}}
that makes little sense when hyphenated across
lines.

Footnotes\footnote{This is an example of a footnote.}
pose no problem.

\LaTeX\ is good at typesetting mathematical formulas
like
       \( x-3y + z = 7 \)
or
       \( a_{1} > x^{2n} + y^{2n} > x' \)
or
       \( \ip{A}{B} = \sum_{i} a_{i} b_{i} \).
The spaces you type in a formula are
ignored.  Remember that a letter like
       $x$                   % $ ... $  and  \( ... \)  are equivalent
is a formula when it denotes a mathematical
symbol, and it should be typed as one.

\section{Displayed Text}

Text is displayed by indenting it from the left
margin.  Quotations are commonly displayed.  There
are short quotations
\begin{quote}
   This is a short quotation.  It consists of a
   single paragraph of text.  See how it is formatted.
\end{quote}
and longer ones.
\begin{quotation}
   This is a longer quotation.  It consists of two
   paragraphs of text, neither of which are
   particularly interesting.

   This is the second paragraph of the quotation.  It
   is just as dull as the first paragraph.
\end{quotation}
Another frequently-displayed structure is a list.
The following is an example of an \emph{itemized}
list.
\begin{itemize}
   \item This is the first item of an itemized list.
         Each item in the list is marked with a ``tick''.
         You don't have to worry about what kind of tick
         mark is used.

   \item This is the second item of the list.  It
         contains another list nested inside it.  The inner
         list is an \emph{enumerated} list.
         \begin{enumerate}
            \item This is the first item of an enumerated
                  list that is nested within the itemized list.

            \item This is the second item of the inner list.
                  \LaTeX\ allows you to nest lists deeper than
                  you really should.
         \end{enumerate}
         This is the rest of the second item of the outer
         list.  It is no more interesting than any other
         part of the item.
   \item This is the third item of the list.
\end{itemize}
You can even display poetry.
\begin{verse}
   There is an environment
    for verse \\             % The \\ command separates lines
   Whose features some poets % within a stanza.
   will curse.

                             % One or more blank lines separate stanzas.

   For instead of making\\
   Them do \emph{all} line breaking, \\
   It allows them to put too many words on a line when they'd rather be
   forced to be terse.
\end{verse}

Mathematical formulas may also be displayed.  A
displayed formula
is
one-line long; multiline
formulas require special formatting instructions.
   \[  \ip{\Gamma}{\psi'} = x'' + y^{2} + z_{i}^{n}\]
Don't start a paragraph with a displayed equation,
nor make one a paragraph by itself.

\end{document}               % End of document.

\section{Bench Huge Section 178}
% synthetic bench section 178
% This is a sample LaTeX input file.  (Version of 12 August 2004.)
%
% A '%' character causes TeX to ignore all remaining text on the line,
% and is used for comments like this one.

\documentclass{article}      % Specifies the document class

                             % The preamble begins here.
\title{An Example Document}  % Declares the document's title.
\author{Leslie Lamport}      % Declares the author's name.
\date{January 21, 1994}      % Deleting this command produces today's date.

\newcommand{\ip}[2]{(#1, #2)}
                             % Defines \ip{arg1}{arg2} to mean
                             % (arg1, arg2).

%\newcommand{\ip}[2]{\langle #1 | #2\rangle}
                             % This is an alternative definition of
                             % \ip that is commented out.

\begin{document}             % End of preamble and beginning of text.

\maketitle                   % Produces the title.

This is an example input file.  Comparing it with
the output it generates can show you how to
produce a simple document of your own.

\section{Ordinary Text}      % Produces section heading.  Lower-level
                             % sections are begun with similar
                             % \subsection and \subsubsection commands.

The ends  of words and sentences are marked
  by   spaces. It  doesn't matter how many
spaces    you type; one is as good as 100.  The
end of   a line counts as a space.

One   or more   blank lines denote the  end
of  a paragraph.

Since any number of consecutive spaces are treated
like a single one, the formatting of the input
file makes no difference to
      \LaTeX,                % The \LaTeX command generates the LaTeX logo.
but it makes a difference to you.  When you use
\LaTeX, making your input file as easy to read
as possible will be a great help as you write
your document and when you change it.  This sample
file shows how you can add comments to your own input
file.

Because printing is different from typewriting,
there are a number of things that you have to do
differently when preparing an input file than if
you were just typing the document directly.
Quotation marks like
       ``this''
have to be handled specially, as do quotes within
quotes:
       ``\,`this'            % \, separates the double and single quote.
        is what I just
        wrote, not  `that'\,''.

Dashes come in three sizes: an
       intra-word
dash, a medium dash for number ranges like
       1--2,
and a punctuation
       dash---like
this.

A sentence-ending space should be larger than the
space between words within a sentence.  You
sometimes have to type special commands in
conjunction with punctuation characters to get
this right, as in the following sentence.
       Gnats, gnus, etc.\ all  % `\ ' makes an inter-word space.
       begin with G\@.         % \@ marks end-of-sentence punctuation.
You should check the spaces after periods when
reading your output to make sure you haven't
forgotten any special cases.  Generating an
ellipsis
       \ldots\               % `\ ' is needed after `\ldots' because TeX
                             % ignores spaces after command names like \ldots
                             % made from \ + letters.
                             %
                             % Note how a `%' character causes TeX to ignore
                             % the end of the input line, so these blank lines
                             % do not start a new paragraph.
                             %
with the right spacing around the periods requires
a special command.

\LaTeX\ interprets some common characters as
commands, so you must type special commands to
generate them.  These characters include the
following:
       \$ \& \% \# \{ and \}.

In printing, text is usually emphasized with an
       \emph{italic}
type style.

\begin{em}
   A long segment of text can also be emphasized
   in this way.  Text within such a segment can be
   given \emph{additional} emphasis.
\end{em}

It is sometimes necessary to prevent \LaTeX\ from
breaking a line where it might otherwise do so.
This may be at a space, as between the ``Mr.''\ and
``Jones'' in
       ``Mr.~Jones'',        % ~ produces an unbreakable interword space.
or within a word---especially when the word is a
symbol like
       \mbox{\emph{itemnum}}
that makes little sense when hyphenated across
lines.

Footnotes\footnote{This is an example of a footnote.}
pose no problem.

\LaTeX\ is good at typesetting mathematical formulas
like
       \( x-3y + z = 7 \)
or
       \( a_{1} > x^{2n} + y^{2n} > x' \)
or
       \( \ip{A}{B} = \sum_{i} a_{i} b_{i} \).
The spaces you type in a formula are
ignored.  Remember that a letter like
       $x$                   % $ ... $  and  \( ... \)  are equivalent
is a formula when it denotes a mathematical
symbol, and it should be typed as one.

\section{Displayed Text}

Text is displayed by indenting it from the left
margin.  Quotations are commonly displayed.  There
are short quotations
\begin{quote}
   This is a short quotation.  It consists of a
   single paragraph of text.  See how it is formatted.
\end{quote}
and longer ones.
\begin{quotation}
   This is a longer quotation.  It consists of two
   paragraphs of text, neither of which are
   particularly interesting.

   This is the second paragraph of the quotation.  It
   is just as dull as the first paragraph.
\end{quotation}
Another frequently-displayed structure is a list.
The following is an example of an \emph{itemized}
list.
\begin{itemize}
   \item This is the first item of an itemized list.
         Each item in the list is marked with a ``tick''.
         You don't have to worry about what kind of tick
         mark is used.

   \item This is the second item of the list.  It
         contains another list nested inside it.  The inner
         list is an \emph{enumerated} list.
         \begin{enumerate}
            \item This is the first item of an enumerated
                  list that is nested within the itemized list.

            \item This is the second item of the inner list.
                  \LaTeX\ allows you to nest lists deeper than
                  you really should.
         \end{enumerate}
         This is the rest of the second item of the outer
         list.  It is no more interesting than any other
         part of the item.
   \item This is the third item of the list.
\end{itemize}
You can even display poetry.
\begin{verse}
   There is an environment
    for verse \\             % The \\ command separates lines
   Whose features some poets % within a stanza.
   will curse.

                             % One or more blank lines separate stanzas.

   For instead of making\\
   Them do \emph{all} line breaking, \\
   It allows them to put too many words on a line when they'd rather be
   forced to be terse.
\end{verse}

Mathematical formulas may also be displayed.  A
displayed formula
is
one-line long; multiline
formulas require special formatting instructions.
   \[  \ip{\Gamma}{\psi'} = x'' + y^{2} + z_{i}^{n}\]
Don't start a paragraph with a displayed equation,
nor make one a paragraph by itself.

\end{document}               % End of document.

\section{Bench Huge Section 179}
% synthetic bench section 179
% This is a sample LaTeX input file.  (Version of 12 August 2004.)
%
% A '%' character causes TeX to ignore all remaining text on the line,
% and is used for comments like this one.

\documentclass{article}      % Specifies the document class

                             % The preamble begins here.
\title{An Example Document}  % Declares the document's title.
\author{Leslie Lamport}      % Declares the author's name.
\date{January 21, 1994}      % Deleting this command produces today's date.

\newcommand{\ip}[2]{(#1, #2)}
                             % Defines \ip{arg1}{arg2} to mean
                             % (arg1, arg2).

%\newcommand{\ip}[2]{\langle #1 | #2\rangle}
                             % This is an alternative definition of
                             % \ip that is commented out.

\begin{document}             % End of preamble and beginning of text.

\maketitle                   % Produces the title.

This is an example input file.  Comparing it with
the output it generates can show you how to
produce a simple document of your own.

\section{Ordinary Text}      % Produces section heading.  Lower-level
                             % sections are begun with similar
                             % \subsection and \subsubsection commands.

The ends  of words and sentences are marked
  by   spaces. It  doesn't matter how many
spaces    you type; one is as good as 100.  The
end of   a line counts as a space.

One   or more   blank lines denote the  end
of  a paragraph.

Since any number of consecutive spaces are treated
like a single one, the formatting of the input
file makes no difference to
      \LaTeX,                % The \LaTeX command generates the LaTeX logo.
but it makes a difference to you.  When you use
\LaTeX, making your input file as easy to read
as possible will be a great help as you write
your document and when you change it.  This sample
file shows how you can add comments to your own input
file.

Because printing is different from typewriting,
there are a number of things that you have to do
differently when preparing an input file than if
you were just typing the document directly.
Quotation marks like
       ``this''
have to be handled specially, as do quotes within
quotes:
       ``\,`this'            % \, separates the double and single quote.
        is what I just
        wrote, not  `that'\,''.

Dashes come in three sizes: an
       intra-word
dash, a medium dash for number ranges like
       1--2,
and a punctuation
       dash---like
this.

A sentence-ending space should be larger than the
space between words within a sentence.  You
sometimes have to type special commands in
conjunction with punctuation characters to get
this right, as in the following sentence.
       Gnats, gnus, etc.\ all  % `\ ' makes an inter-word space.
       begin with G\@.         % \@ marks end-of-sentence punctuation.
You should check the spaces after periods when
reading your output to make sure you haven't
forgotten any special cases.  Generating an
ellipsis
       \ldots\               % `\ ' is needed after `\ldots' because TeX
                             % ignores spaces after command names like \ldots
                             % made from \ + letters.
                             %
                             % Note how a `%' character causes TeX to ignore
                             % the end of the input line, so these blank lines
                             % do not start a new paragraph.
                             %
with the right spacing around the periods requires
a special command.

\LaTeX\ interprets some common characters as
commands, so you must type special commands to
generate them.  These characters include the
following:
       \$ \& \% \# \{ and \}.

In printing, text is usually emphasized with an
       \emph{italic}
type style.

\begin{em}
   A long segment of text can also be emphasized
   in this way.  Text within such a segment can be
   given \emph{additional} emphasis.
\end{em}

It is sometimes necessary to prevent \LaTeX\ from
breaking a line where it might otherwise do so.
This may be at a space, as between the ``Mr.''\ and
``Jones'' in
       ``Mr.~Jones'',        % ~ produces an unbreakable interword space.
or within a word---especially when the word is a
symbol like
       \mbox{\emph{itemnum}}
that makes little sense when hyphenated across
lines.

Footnotes\footnote{This is an example of a footnote.}
pose no problem.

\LaTeX\ is good at typesetting mathematical formulas
like
       \( x-3y + z = 7 \)
or
       \( a_{1} > x^{2n} + y^{2n} > x' \)
or
       \( \ip{A}{B} = \sum_{i} a_{i} b_{i} \).
The spaces you type in a formula are
ignored.  Remember that a letter like
       $x$                   % $ ... $  and  \( ... \)  are equivalent
is a formula when it denotes a mathematical
symbol, and it should be typed as one.

\section{Displayed Text}

Text is displayed by indenting it from the left
margin.  Quotations are commonly displayed.  There
are short quotations
\begin{quote}
   This is a short quotation.  It consists of a
   single paragraph of text.  See how it is formatted.
\end{quote}
and longer ones.
\begin{quotation}
   This is a longer quotation.  It consists of two
   paragraphs of text, neither of which are
   particularly interesting.

   This is the second paragraph of the quotation.  It
   is just as dull as the first paragraph.
\end{quotation}
Another frequently-displayed structure is a list.
The following is an example of an \emph{itemized}
list.
\begin{itemize}
   \item This is the first item of an itemized list.
         Each item in the list is marked with a ``tick''.
         You don't have to worry about what kind of tick
         mark is used.

   \item This is the second item of the list.  It
         contains another list nested inside it.  The inner
         list is an \emph{enumerated} list.
         \begin{enumerate}
            \item This is the first item of an enumerated
                  list that is nested within the itemized list.

            \item This is the second item of the inner list.
                  \LaTeX\ allows you to nest lists deeper than
                  you really should.
         \end{enumerate}
         This is the rest of the second item of the outer
         list.  It is no more interesting than any other
         part of the item.
   \item This is the third item of the list.
\end{itemize}
You can even display poetry.
\begin{verse}
   There is an environment
    for verse \\             % The \\ command separates lines
   Whose features some poets % within a stanza.
   will curse.

                             % One or more blank lines separate stanzas.

   For instead of making\\
   Them do \emph{all} line breaking, \\
   It allows them to put too many words on a line when they'd rather be
   forced to be terse.
\end{verse}

Mathematical formulas may also be displayed.  A
displayed formula
is
one-line long; multiline
formulas require special formatting instructions.
   \[  \ip{\Gamma}{\psi'} = x'' + y^{2} + z_{i}^{n}\]
Don't start a paragraph with a displayed equation,
nor make one a paragraph by itself.

\end{document}               % End of document.

\section{Bench Huge Section 180}
% synthetic bench section 180
% This is a sample LaTeX input file.  (Version of 12 August 2004.)
%
% A '%' character causes TeX to ignore all remaining text on the line,
% and is used for comments like this one.

\documentclass{article}      % Specifies the document class

                             % The preamble begins here.
\title{An Example Document}  % Declares the document's title.
\author{Leslie Lamport}      % Declares the author's name.
\date{January 21, 1994}      % Deleting this command produces today's date.

\newcommand{\ip}[2]{(#1, #2)}
                             % Defines \ip{arg1}{arg2} to mean
                             % (arg1, arg2).

%\newcommand{\ip}[2]{\langle #1 | #2\rangle}
                             % This is an alternative definition of
                             % \ip that is commented out.

\begin{document}             % End of preamble and beginning of text.

\maketitle                   % Produces the title.

This is an example input file.  Comparing it with
the output it generates can show you how to
produce a simple document of your own.

\section{Ordinary Text}      % Produces section heading.  Lower-level
                             % sections are begun with similar
                             % \subsection and \subsubsection commands.

The ends  of words and sentences are marked
  by   spaces. It  doesn't matter how many
spaces    you type; one is as good as 100.  The
end of   a line counts as a space.

One   or more   blank lines denote the  end
of  a paragraph.

Since any number of consecutive spaces are treated
like a single one, the formatting of the input
file makes no difference to
      \LaTeX,                % The \LaTeX command generates the LaTeX logo.
but it makes a difference to you.  When you use
\LaTeX, making your input file as easy to read
as possible will be a great help as you write
your document and when you change it.  This sample
file shows how you can add comments to your own input
file.

Because printing is different from typewriting,
there are a number of things that you have to do
differently when preparing an input file than if
you were just typing the document directly.
Quotation marks like
       ``this''
have to be handled specially, as do quotes within
quotes:
       ``\,`this'            % \, separates the double and single quote.
        is what I just
        wrote, not  `that'\,''.

Dashes come in three sizes: an
       intra-word
dash, a medium dash for number ranges like
       1--2,
and a punctuation
       dash---like
this.

A sentence-ending space should be larger than the
space between words within a sentence.  You
sometimes have to type special commands in
conjunction with punctuation characters to get
this right, as in the following sentence.
       Gnats, gnus, etc.\ all  % `\ ' makes an inter-word space.
       begin with G\@.         % \@ marks end-of-sentence punctuation.
You should check the spaces after periods when
reading your output to make sure you haven't
forgotten any special cases.  Generating an
ellipsis
       \ldots\               % `\ ' is needed after `\ldots' because TeX
                             % ignores spaces after command names like \ldots
                             % made from \ + letters.
                             %
                             % Note how a `%' character causes TeX to ignore
                             % the end of the input line, so these blank lines
                             % do not start a new paragraph.
                             %
with the right spacing around the periods requires
a special command.

\LaTeX\ interprets some common characters as
commands, so you must type special commands to
generate them.  These characters include the
following:
       \$ \& \% \# \{ and \}.

In printing, text is usually emphasized with an
       \emph{italic}
type style.

\begin{em}
   A long segment of text can also be emphasized
   in this way.  Text within such a segment can be
   given \emph{additional} emphasis.
\end{em}

It is sometimes necessary to prevent \LaTeX\ from
breaking a line where it might otherwise do so.
This may be at a space, as between the ``Mr.''\ and
``Jones'' in
       ``Mr.~Jones'',        % ~ produces an unbreakable interword space.
or within a word---especially when the word is a
symbol like
       \mbox{\emph{itemnum}}
that makes little sense when hyphenated across
lines.

Footnotes\footnote{This is an example of a footnote.}
pose no problem.

\LaTeX\ is good at typesetting mathematical formulas
like
       \( x-3y + z = 7 \)
or
       \( a_{1} > x^{2n} + y^{2n} > x' \)
or
       \( \ip{A}{B} = \sum_{i} a_{i} b_{i} \).
The spaces you type in a formula are
ignored.  Remember that a letter like
       $x$                   % $ ... $  and  \( ... \)  are equivalent
is a formula when it denotes a mathematical
symbol, and it should be typed as one.

\section{Displayed Text}

Text is displayed by indenting it from the left
margin.  Quotations are commonly displayed.  There
are short quotations
\begin{quote}
   This is a short quotation.  It consists of a
   single paragraph of text.  See how it is formatted.
\end{quote}
and longer ones.
\begin{quotation}
   This is a longer quotation.  It consists of two
   paragraphs of text, neither of which are
   particularly interesting.

   This is the second paragraph of the quotation.  It
   is just as dull as the first paragraph.
\end{quotation}
Another frequently-displayed structure is a list.
The following is an example of an \emph{itemized}
list.
\begin{itemize}
   \item This is the first item of an itemized list.
         Each item in the list is marked with a ``tick''.
         You don't have to worry about what kind of tick
         mark is used.

   \item This is the second item of the list.  It
         contains another list nested inside it.  The inner
         list is an \emph{enumerated} list.
         \begin{enumerate}
            \item This is the first item of an enumerated
                  list that is nested within the itemized list.

            \item This is the second item of the inner list.
                  \LaTeX\ allows you to nest lists deeper than
                  you really should.
         \end{enumerate}
         This is the rest of the second item of the outer
         list.  It is no more interesting than any other
         part of the item.
   \item This is the third item of the list.
\end{itemize}
You can even display poetry.
\begin{verse}
   There is an environment
    for verse \\             % The \\ command separates lines
   Whose features some poets % within a stanza.
   will curse.

                             % One or more blank lines separate stanzas.

   For instead of making\\
   Them do \emph{all} line breaking, \\
   It allows them to put too many words on a line when they'd rather be
   forced to be terse.
\end{verse}

Mathematical formulas may also be displayed.  A
displayed formula
is
one-line long; multiline
formulas require special formatting instructions.
   \[  \ip{\Gamma}{\psi'} = x'' + y^{2} + z_{i}^{n}\]
Don't start a paragraph with a displayed equation,
nor make one a paragraph by itself.

\end{document}               % End of document.

\section{Bench Huge Section 181}
% synthetic bench section 181
% This is a sample LaTeX input file.  (Version of 12 August 2004.)
%
% A '%' character causes TeX to ignore all remaining text on the line,
% and is used for comments like this one.

\documentclass{article}      % Specifies the document class

                             % The preamble begins here.
\title{An Example Document}  % Declares the document's title.
\author{Leslie Lamport}      % Declares the author's name.
\date{January 21, 1994}      % Deleting this command produces today's date.

\newcommand{\ip}[2]{(#1, #2)}
                             % Defines \ip{arg1}{arg2} to mean
                             % (arg1, arg2).

%\newcommand{\ip}[2]{\langle #1 | #2\rangle}
                             % This is an alternative definition of
                             % \ip that is commented out.

\begin{document}             % End of preamble and beginning of text.

\maketitle                   % Produces the title.

This is an example input file.  Comparing it with
the output it generates can show you how to
produce a simple document of your own.

\section{Ordinary Text}      % Produces section heading.  Lower-level
                             % sections are begun with similar
                             % \subsection and \subsubsection commands.

The ends  of words and sentences are marked
  by   spaces. It  doesn't matter how many
spaces    you type; one is as good as 100.  The
end of   a line counts as a space.

One   or more   blank lines denote the  end
of  a paragraph.

Since any number of consecutive spaces are treated
like a single one, the formatting of the input
file makes no difference to
      \LaTeX,                % The \LaTeX command generates the LaTeX logo.
but it makes a difference to you.  When you use
\LaTeX, making your input file as easy to read
as possible will be a great help as you write
your document and when you change it.  This sample
file shows how you can add comments to your own input
file.

Because printing is different from typewriting,
there are a number of things that you have to do
differently when preparing an input file than if
you were just typing the document directly.
Quotation marks like
       ``this''
have to be handled specially, as do quotes within
quotes:
       ``\,`this'            % \, separates the double and single quote.
        is what I just
        wrote, not  `that'\,''.

Dashes come in three sizes: an
       intra-word
dash, a medium dash for number ranges like
       1--2,
and a punctuation
       dash---like
this.

A sentence-ending space should be larger than the
space between words within a sentence.  You
sometimes have to type special commands in
conjunction with punctuation characters to get
this right, as in the following sentence.
       Gnats, gnus, etc.\ all  % `\ ' makes an inter-word space.
       begin with G\@.         % \@ marks end-of-sentence punctuation.
You should check the spaces after periods when
reading your output to make sure you haven't
forgotten any special cases.  Generating an
ellipsis
       \ldots\               % `\ ' is needed after `\ldots' because TeX
                             % ignores spaces after command names like \ldots
                             % made from \ + letters.
                             %
                             % Note how a `%' character causes TeX to ignore
                             % the end of the input line, so these blank lines
                             % do not start a new paragraph.
                             %
with the right spacing around the periods requires
a special command.

\LaTeX\ interprets some common characters as
commands, so you must type special commands to
generate them.  These characters include the
following:
       \$ \& \% \# \{ and \}.

In printing, text is usually emphasized with an
       \emph{italic}
type style.

\begin{em}
   A long segment of text can also be emphasized
   in this way.  Text within such a segment can be
   given \emph{additional} emphasis.
\end{em}

It is sometimes necessary to prevent \LaTeX\ from
breaking a line where it might otherwise do so.
This may be at a space, as between the ``Mr.''\ and
``Jones'' in
       ``Mr.~Jones'',        % ~ produces an unbreakable interword space.
or within a word---especially when the word is a
symbol like
       \mbox{\emph{itemnum}}
that makes little sense when hyphenated across
lines.

Footnotes\footnote{This is an example of a footnote.}
pose no problem.

\LaTeX\ is good at typesetting mathematical formulas
like
       \( x-3y + z = 7 \)
or
       \( a_{1} > x^{2n} + y^{2n} > x' \)
or
       \( \ip{A}{B} = \sum_{i} a_{i} b_{i} \).
The spaces you type in a formula are
ignored.  Remember that a letter like
       $x$                   % $ ... $  and  \( ... \)  are equivalent
is a formula when it denotes a mathematical
symbol, and it should be typed as one.

\section{Displayed Text}

Text is displayed by indenting it from the left
margin.  Quotations are commonly displayed.  There
are short quotations
\begin{quote}
   This is a short quotation.  It consists of a
   single paragraph of text.  See how it is formatted.
\end{quote}
and longer ones.
\begin{quotation}
   This is a longer quotation.  It consists of two
   paragraphs of text, neither of which are
   particularly interesting.

   This is the second paragraph of the quotation.  It
   is just as dull as the first paragraph.
\end{quotation}
Another frequently-displayed structure is a list.
The following is an example of an \emph{itemized}
list.
\begin{itemize}
   \item This is the first item of an itemized list.
         Each item in the list is marked with a ``tick''.
         You don't have to worry about what kind of tick
         mark is used.

   \item This is the second item of the list.  It
         contains another list nested inside it.  The inner
         list is an \emph{enumerated} list.
         \begin{enumerate}
            \item This is the first item of an enumerated
                  list that is nested within the itemized list.

            \item This is the second item of the inner list.
                  \LaTeX\ allows you to nest lists deeper than
                  you really should.
         \end{enumerate}
         This is the rest of the second item of the outer
         list.  It is no more interesting than any other
         part of the item.
   \item This is the third item of the list.
\end{itemize}
You can even display poetry.
\begin{verse}
   There is an environment
    for verse \\             % The \\ command separates lines
   Whose features some poets % within a stanza.
   will curse.

                             % One or more blank lines separate stanzas.

   For instead of making\\
   Them do \emph{all} line breaking, \\
   It allows them to put too many words on a line when they'd rather be
   forced to be terse.
\end{verse}

Mathematical formulas may also be displayed.  A
displayed formula
is
one-line long; multiline
formulas require special formatting instructions.
   \[  \ip{\Gamma}{\psi'} = x'' + y^{2} + z_{i}^{n}\]
Don't start a paragraph with a displayed equation,
nor make one a paragraph by itself.

\end{document}               % End of document.

\section{Bench Huge Section 182}
% synthetic bench section 182
% This is a sample LaTeX input file.  (Version of 12 August 2004.)
%
% A '%' character causes TeX to ignore all remaining text on the line,
% and is used for comments like this one.

\documentclass{article}      % Specifies the document class

                             % The preamble begins here.
\title{An Example Document}  % Declares the document's title.
\author{Leslie Lamport}      % Declares the author's name.
\date{January 21, 1994}      % Deleting this command produces today's date.

\newcommand{\ip}[2]{(#1, #2)}
                             % Defines \ip{arg1}{arg2} to mean
                             % (arg1, arg2).

%\newcommand{\ip}[2]{\langle #1 | #2\rangle}
                             % This is an alternative definition of
                             % \ip that is commented out.

\begin{document}             % End of preamble and beginning of text.

\maketitle                   % Produces the title.

This is an example input file.  Comparing it with
the output it generates can show you how to
produce a simple document of your own.

\section{Ordinary Text}      % Produces section heading.  Lower-level
                             % sections are begun with similar
                             % \subsection and \subsubsection commands.

The ends  of words and sentences are marked
  by   spaces. It  doesn't matter how many
spaces    you type; one is as good as 100.  The
end of   a line counts as a space.

One   or more   blank lines denote the  end
of  a paragraph.

Since any number of consecutive spaces are treated
like a single one, the formatting of the input
file makes no difference to
      \LaTeX,                % The \LaTeX command generates the LaTeX logo.
but it makes a difference to you.  When you use
\LaTeX, making your input file as easy to read
as possible will be a great help as you write
your document and when you change it.  This sample
file shows how you can add comments to your own input
file.

Because printing is different from typewriting,
there are a number of things that you have to do
differently when preparing an input file than if
you were just typing the document directly.
Quotation marks like
       ``this''
have to be handled specially, as do quotes within
quotes:
       ``\,`this'            % \, separates the double and single quote.
        is what I just
        wrote, not  `that'\,''.

Dashes come in three sizes: an
       intra-word
dash, a medium dash for number ranges like
       1--2,
and a punctuation
       dash---like
this.

A sentence-ending space should be larger than the
space between words within a sentence.  You
sometimes have to type special commands in
conjunction with punctuation characters to get
this right, as in the following sentence.
       Gnats, gnus, etc.\ all  % `\ ' makes an inter-word space.
       begin with G\@.         % \@ marks end-of-sentence punctuation.
You should check the spaces after periods when
reading your output to make sure you haven't
forgotten any special cases.  Generating an
ellipsis
       \ldots\               % `\ ' is needed after `\ldots' because TeX
                             % ignores spaces after command names like \ldots
                             % made from \ + letters.
                             %
                             % Note how a `%' character causes TeX to ignore
                             % the end of the input line, so these blank lines
                             % do not start a new paragraph.
                             %
with the right spacing around the periods requires
a special command.

\LaTeX\ interprets some common characters as
commands, so you must type special commands to
generate them.  These characters include the
following:
       \$ \& \% \# \{ and \}.

In printing, text is usually emphasized with an
       \emph{italic}
type style.

\begin{em}
   A long segment of text can also be emphasized
   in this way.  Text within such a segment can be
   given \emph{additional} emphasis.
\end{em}

It is sometimes necessary to prevent \LaTeX\ from
breaking a line where it might otherwise do so.
This may be at a space, as between the ``Mr.''\ and
``Jones'' in
       ``Mr.~Jones'',        % ~ produces an unbreakable interword space.
or within a word---especially when the word is a
symbol like
       \mbox{\emph{itemnum}}
that makes little sense when hyphenated across
lines.

Footnotes\footnote{This is an example of a footnote.}
pose no problem.

\LaTeX\ is good at typesetting mathematical formulas
like
       \( x-3y + z = 7 \)
or
       \( a_{1} > x^{2n} + y^{2n} > x' \)
or
       \( \ip{A}{B} = \sum_{i} a_{i} b_{i} \).
The spaces you type in a formula are
ignored.  Remember that a letter like
       $x$                   % $ ... $  and  \( ... \)  are equivalent
is a formula when it denotes a mathematical
symbol, and it should be typed as one.

\section{Displayed Text}

Text is displayed by indenting it from the left
margin.  Quotations are commonly displayed.  There
are short quotations
\begin{quote}
   This is a short quotation.  It consists of a
   single paragraph of text.  See how it is formatted.
\end{quote}
and longer ones.
\begin{quotation}
   This is a longer quotation.  It consists of two
   paragraphs of text, neither of which are
   particularly interesting.

   This is the second paragraph of the quotation.  It
   is just as dull as the first paragraph.
\end{quotation}
Another frequently-displayed structure is a list.
The following is an example of an \emph{itemized}
list.
\begin{itemize}
   \item This is the first item of an itemized list.
         Each item in the list is marked with a ``tick''.
         You don't have to worry about what kind of tick
         mark is used.

   \item This is the second item of the list.  It
         contains another list nested inside it.  The inner
         list is an \emph{enumerated} list.
         \begin{enumerate}
            \item This is the first item of an enumerated
                  list that is nested within the itemized list.

            \item This is the second item of the inner list.
                  \LaTeX\ allows you to nest lists deeper than
                  you really should.
         \end{enumerate}
         This is the rest of the second item of the outer
         list.  It is no more interesting than any other
         part of the item.
   \item This is the third item of the list.
\end{itemize}
You can even display poetry.
\begin{verse}
   There is an environment
    for verse \\             % The \\ command separates lines
   Whose features some poets % within a stanza.
   will curse.

                             % One or more blank lines separate stanzas.

   For instead of making\\
   Them do \emph{all} line breaking, \\
   It allows them to put too many words on a line when they'd rather be
   forced to be terse.
\end{verse}

Mathematical formulas may also be displayed.  A
displayed formula
is
one-line long; multiline
formulas require special formatting instructions.
   \[  \ip{\Gamma}{\psi'} = x'' + y^{2} + z_{i}^{n}\]
Don't start a paragraph with a displayed equation,
nor make one a paragraph by itself.

\end{document}               % End of document.

\section{Bench Huge Section 183}
% synthetic bench section 183
% This is a sample LaTeX input file.  (Version of 12 August 2004.)
%
% A '%' character causes TeX to ignore all remaining text on the line,
% and is used for comments like this one.

\documentclass{article}      % Specifies the document class

                             % The preamble begins here.
\title{An Example Document}  % Declares the document's title.
\author{Leslie Lamport}      % Declares the author's name.
\date{January 21, 1994}      % Deleting this command produces today's date.

\newcommand{\ip}[2]{(#1, #2)}
                             % Defines \ip{arg1}{arg2} to mean
                             % (arg1, arg2).

%\newcommand{\ip}[2]{\langle #1 | #2\rangle}
                             % This is an alternative definition of
                             % \ip that is commented out.

\begin{document}             % End of preamble and beginning of text.

\maketitle                   % Produces the title.

This is an example input file.  Comparing it with
the output it generates can show you how to
produce a simple document of your own.

\section{Ordinary Text}      % Produces section heading.  Lower-level
                             % sections are begun with similar
                             % \subsection and \subsubsection commands.

The ends  of words and sentences are marked
  by   spaces. It  doesn't matter how many
spaces    you type; one is as good as 100.  The
end of   a line counts as a space.

One   or more   blank lines denote the  end
of  a paragraph.

Since any number of consecutive spaces are treated
like a single one, the formatting of the input
file makes no difference to
      \LaTeX,                % The \LaTeX command generates the LaTeX logo.
but it makes a difference to you.  When you use
\LaTeX, making your input file as easy to read
as possible will be a great help as you write
your document and when you change it.  This sample
file shows how you can add comments to your own input
file.

Because printing is different from typewriting,
there are a number of things that you have to do
differently when preparing an input file than if
you were just typing the document directly.
Quotation marks like
       ``this''
have to be handled specially, as do quotes within
quotes:
       ``\,`this'            % \, separates the double and single quote.
        is what I just
        wrote, not  `that'\,''.

Dashes come in three sizes: an
       intra-word
dash, a medium dash for number ranges like
       1--2,
and a punctuation
       dash---like
this.

A sentence-ending space should be larger than the
space between words within a sentence.  You
sometimes have to type special commands in
conjunction with punctuation characters to get
this right, as in the following sentence.
       Gnats, gnus, etc.\ all  % `\ ' makes an inter-word space.
       begin with G\@.         % \@ marks end-of-sentence punctuation.
You should check the spaces after periods when
reading your output to make sure you haven't
forgotten any special cases.  Generating an
ellipsis
       \ldots\               % `\ ' is needed after `\ldots' because TeX
                             % ignores spaces after command names like \ldots
                             % made from \ + letters.
                             %
                             % Note how a `%' character causes TeX to ignore
                             % the end of the input line, so these blank lines
                             % do not start a new paragraph.
                             %
with the right spacing around the periods requires
a special command.

\LaTeX\ interprets some common characters as
commands, so you must type special commands to
generate them.  These characters include the
following:
       \$ \& \% \# \{ and \}.

In printing, text is usually emphasized with an
       \emph{italic}
type style.

\begin{em}
   A long segment of text can also be emphasized
   in this way.  Text within such a segment can be
   given \emph{additional} emphasis.
\end{em}

It is sometimes necessary to prevent \LaTeX\ from
breaking a line where it might otherwise do so.
This may be at a space, as between the ``Mr.''\ and
``Jones'' in
       ``Mr.~Jones'',        % ~ produces an unbreakable interword space.
or within a word---especially when the word is a
symbol like
       \mbox{\emph{itemnum}}
that makes little sense when hyphenated across
lines.

Footnotes\footnote{This is an example of a footnote.}
pose no problem.

\LaTeX\ is good at typesetting mathematical formulas
like
       \( x-3y + z = 7 \)
or
       \( a_{1} > x^{2n} + y^{2n} > x' \)
or
       \( \ip{A}{B} = \sum_{i} a_{i} b_{i} \).
The spaces you type in a formula are
ignored.  Remember that a letter like
       $x$                   % $ ... $  and  \( ... \)  are equivalent
is a formula when it denotes a mathematical
symbol, and it should be typed as one.

\section{Displayed Text}

Text is displayed by indenting it from the left
margin.  Quotations are commonly displayed.  There
are short quotations
\begin{quote}
   This is a short quotation.  It consists of a
   single paragraph of text.  See how it is formatted.
\end{quote}
and longer ones.
\begin{quotation}
   This is a longer quotation.  It consists of two
   paragraphs of text, neither of which are
   particularly interesting.

   This is the second paragraph of the quotation.  It
   is just as dull as the first paragraph.
\end{quotation}
Another frequently-displayed structure is a list.
The following is an example of an \emph{itemized}
list.
\begin{itemize}
   \item This is the first item of an itemized list.
         Each item in the list is marked with a ``tick''.
         You don't have to worry about what kind of tick
         mark is used.

   \item This is the second item of the list.  It
         contains another list nested inside it.  The inner
         list is an \emph{enumerated} list.
         \begin{enumerate}
            \item This is the first item of an enumerated
                  list that is nested within the itemized list.

            \item This is the second item of the inner list.
                  \LaTeX\ allows you to nest lists deeper than
                  you really should.
         \end{enumerate}
         This is the rest of the second item of the outer
         list.  It is no more interesting than any other
         part of the item.
   \item This is the third item of the list.
\end{itemize}
You can even display poetry.
\begin{verse}
   There is an environment
    for verse \\             % The \\ command separates lines
   Whose features some poets % within a stanza.
   will curse.

                             % One or more blank lines separate stanzas.

   For instead of making\\
   Them do \emph{all} line breaking, \\
   It allows them to put too many words on a line when they'd rather be
   forced to be terse.
\end{verse}

Mathematical formulas may also be displayed.  A
displayed formula
is
one-line long; multiline
formulas require special formatting instructions.
   \[  \ip{\Gamma}{\psi'} = x'' + y^{2} + z_{i}^{n}\]
Don't start a paragraph with a displayed equation,
nor make one a paragraph by itself.

\end{document}               % End of document.

\section{Bench Huge Section 184}
% synthetic bench section 184
% This is a sample LaTeX input file.  (Version of 12 August 2004.)
%
% A '%' character causes TeX to ignore all remaining text on the line,
% and is used for comments like this one.

\documentclass{article}      % Specifies the document class

                             % The preamble begins here.
\title{An Example Document}  % Declares the document's title.
\author{Leslie Lamport}      % Declares the author's name.
\date{January 21, 1994}      % Deleting this command produces today's date.

\newcommand{\ip}[2]{(#1, #2)}
                             % Defines \ip{arg1}{arg2} to mean
                             % (arg1, arg2).

%\newcommand{\ip}[2]{\langle #1 | #2\rangle}
                             % This is an alternative definition of
                             % \ip that is commented out.

\begin{document}             % End of preamble and beginning of text.

\maketitle                   % Produces the title.

This is an example input file.  Comparing it with
the output it generates can show you how to
produce a simple document of your own.

\section{Ordinary Text}      % Produces section heading.  Lower-level
                             % sections are begun with similar
                             % \subsection and \subsubsection commands.

The ends  of words and sentences are marked
  by   spaces. It  doesn't matter how many
spaces    you type; one is as good as 100.  The
end of   a line counts as a space.

One   or more   blank lines denote the  end
of  a paragraph.

Since any number of consecutive spaces are treated
like a single one, the formatting of the input
file makes no difference to
      \LaTeX,                % The \LaTeX command generates the LaTeX logo.
but it makes a difference to you.  When you use
\LaTeX, making your input file as easy to read
as possible will be a great help as you write
your document and when you change it.  This sample
file shows how you can add comments to your own input
file.

Because printing is different from typewriting,
there are a number of things that you have to do
differently when preparing an input file than if
you were just typing the document directly.
Quotation marks like
       ``this''
have to be handled specially, as do quotes within
quotes:
       ``\,`this'            % \, separates the double and single quote.
        is what I just
        wrote, not  `that'\,''.

Dashes come in three sizes: an
       intra-word
dash, a medium dash for number ranges like
       1--2,
and a punctuation
       dash---like
this.

A sentence-ending space should be larger than the
space between words within a sentence.  You
sometimes have to type special commands in
conjunction with punctuation characters to get
this right, as in the following sentence.
       Gnats, gnus, etc.\ all  % `\ ' makes an inter-word space.
       begin with G\@.         % \@ marks end-of-sentence punctuation.
You should check the spaces after periods when
reading your output to make sure you haven't
forgotten any special cases.  Generating an
ellipsis
       \ldots\               % `\ ' is needed after `\ldots' because TeX
                             % ignores spaces after command names like \ldots
                             % made from \ + letters.
                             %
                             % Note how a `%' character causes TeX to ignore
                             % the end of the input line, so these blank lines
                             % do not start a new paragraph.
                             %
with the right spacing around the periods requires
a special command.

\LaTeX\ interprets some common characters as
commands, so you must type special commands to
generate them.  These characters include the
following:
       \$ \& \% \# \{ and \}.

In printing, text is usually emphasized with an
       \emph{italic}
type style.

\begin{em}
   A long segment of text can also be emphasized
   in this way.  Text within such a segment can be
   given \emph{additional} emphasis.
\end{em}

It is sometimes necessary to prevent \LaTeX\ from
breaking a line where it might otherwise do so.
This may be at a space, as between the ``Mr.''\ and
``Jones'' in
       ``Mr.~Jones'',        % ~ produces an unbreakable interword space.
or within a word---especially when the word is a
symbol like
       \mbox{\emph{itemnum}}
that makes little sense when hyphenated across
lines.

Footnotes\footnote{This is an example of a footnote.}
pose no problem.

\LaTeX\ is good at typesetting mathematical formulas
like
       \( x-3y + z = 7 \)
or
       \( a_{1} > x^{2n} + y^{2n} > x' \)
or
       \( \ip{A}{B} = \sum_{i} a_{i} b_{i} \).
The spaces you type in a formula are
ignored.  Remember that a letter like
       $x$                   % $ ... $  and  \( ... \)  are equivalent
is a formula when it denotes a mathematical
symbol, and it should be typed as one.

\section{Displayed Text}

Text is displayed by indenting it from the left
margin.  Quotations are commonly displayed.  There
are short quotations
\begin{quote}
   This is a short quotation.  It consists of a
   single paragraph of text.  See how it is formatted.
\end{quote}
and longer ones.
\begin{quotation}
   This is a longer quotation.  It consists of two
   paragraphs of text, neither of which are
   particularly interesting.

   This is the second paragraph of the quotation.  It
   is just as dull as the first paragraph.
\end{quotation}
Another frequently-displayed structure is a list.
The following is an example of an \emph{itemized}
list.
\begin{itemize}
   \item This is the first item of an itemized list.
         Each item in the list is marked with a ``tick''.
         You don't have to worry about what kind of tick
         mark is used.

   \item This is the second item of the list.  It
         contains another list nested inside it.  The inner
         list is an \emph{enumerated} list.
         \begin{enumerate}
            \item This is the first item of an enumerated
                  list that is nested within the itemized list.

            \item This is the second item of the inner list.
                  \LaTeX\ allows you to nest lists deeper than
                  you really should.
         \end{enumerate}
         This is the rest of the second item of the outer
         list.  It is no more interesting than any other
         part of the item.
   \item This is the third item of the list.
\end{itemize}
You can even display poetry.
\begin{verse}
   There is an environment
    for verse \\             % The \\ command separates lines
   Whose features some poets % within a stanza.
   will curse.

                             % One or more blank lines separate stanzas.

   For instead of making\\
   Them do \emph{all} line breaking, \\
   It allows them to put too many words on a line when they'd rather be
   forced to be terse.
\end{verse}

Mathematical formulas may also be displayed.  A
displayed formula
is
one-line long; multiline
formulas require special formatting instructions.
   \[  \ip{\Gamma}{\psi'} = x'' + y^{2} + z_{i}^{n}\]
Don't start a paragraph with a displayed equation,
nor make one a paragraph by itself.

\end{document}               % End of document.

\section{Bench Huge Section 185}
% synthetic bench section 185
% This is a sample LaTeX input file.  (Version of 12 August 2004.)
%
% A '%' character causes TeX to ignore all remaining text on the line,
% and is used for comments like this one.

\documentclass{article}      % Specifies the document class

                             % The preamble begins here.
\title{An Example Document}  % Declares the document's title.
\author{Leslie Lamport}      % Declares the author's name.
\date{January 21, 1994}      % Deleting this command produces today's date.

\newcommand{\ip}[2]{(#1, #2)}
                             % Defines \ip{arg1}{arg2} to mean
                             % (arg1, arg2).

%\newcommand{\ip}[2]{\langle #1 | #2\rangle}
                             % This is an alternative definition of
                             % \ip that is commented out.

\begin{document}             % End of preamble and beginning of text.

\maketitle                   % Produces the title.

This is an example input file.  Comparing it with
the output it generates can show you how to
produce a simple document of your own.

\section{Ordinary Text}      % Produces section heading.  Lower-level
                             % sections are begun with similar
                             % \subsection and \subsubsection commands.

The ends  of words and sentences are marked
  by   spaces. It  doesn't matter how many
spaces    you type; one is as good as 100.  The
end of   a line counts as a space.

One   or more   blank lines denote the  end
of  a paragraph.

Since any number of consecutive spaces are treated
like a single one, the formatting of the input
file makes no difference to
      \LaTeX,                % The \LaTeX command generates the LaTeX logo.
but it makes a difference to you.  When you use
\LaTeX, making your input file as easy to read
as possible will be a great help as you write
your document and when you change it.  This sample
file shows how you can add comments to your own input
file.

Because printing is different from typewriting,
there are a number of things that you have to do
differently when preparing an input file than if
you were just typing the document directly.
Quotation marks like
       ``this''
have to be handled specially, as do quotes within
quotes:
       ``\,`this'            % \, separates the double and single quote.
        is what I just
        wrote, not  `that'\,''.

Dashes come in three sizes: an
       intra-word
dash, a medium dash for number ranges like
       1--2,
and a punctuation
       dash---like
this.

A sentence-ending space should be larger than the
space between words within a sentence.  You
sometimes have to type special commands in
conjunction with punctuation characters to get
this right, as in the following sentence.
       Gnats, gnus, etc.\ all  % `\ ' makes an inter-word space.
       begin with G\@.         % \@ marks end-of-sentence punctuation.
You should check the spaces after periods when
reading your output to make sure you haven't
forgotten any special cases.  Generating an
ellipsis
       \ldots\               % `\ ' is needed after `\ldots' because TeX
                             % ignores spaces after command names like \ldots
                             % made from \ + letters.
                             %
                             % Note how a `%' character causes TeX to ignore
                             % the end of the input line, so these blank lines
                             % do not start a new paragraph.
                             %
with the right spacing around the periods requires
a special command.

\LaTeX\ interprets some common characters as
commands, so you must type special commands to
generate them.  These characters include the
following:
       \$ \& \% \# \{ and \}.

In printing, text is usually emphasized with an
       \emph{italic}
type style.

\begin{em}
   A long segment of text can also be emphasized
   in this way.  Text within such a segment can be
   given \emph{additional} emphasis.
\end{em}

It is sometimes necessary to prevent \LaTeX\ from
breaking a line where it might otherwise do so.
This may be at a space, as between the ``Mr.''\ and
``Jones'' in
       ``Mr.~Jones'',        % ~ produces an unbreakable interword space.
or within a word---especially when the word is a
symbol like
       \mbox{\emph{itemnum}}
that makes little sense when hyphenated across
lines.

Footnotes\footnote{This is an example of a footnote.}
pose no problem.

\LaTeX\ is good at typesetting mathematical formulas
like
       \( x-3y + z = 7 \)
or
       \( a_{1} > x^{2n} + y^{2n} > x' \)
or
       \( \ip{A}{B} = \sum_{i} a_{i} b_{i} \).
The spaces you type in a formula are
ignored.  Remember that a letter like
       $x$                   % $ ... $  and  \( ... \)  are equivalent
is a formula when it denotes a mathematical
symbol, and it should be typed as one.

\section{Displayed Text}

Text is displayed by indenting it from the left
margin.  Quotations are commonly displayed.  There
are short quotations
\begin{quote}
   This is a short quotation.  It consists of a
   single paragraph of text.  See how it is formatted.
\end{quote}
and longer ones.
\begin{quotation}
   This is a longer quotation.  It consists of two
   paragraphs of text, neither of which are
   particularly interesting.

   This is the second paragraph of the quotation.  It
   is just as dull as the first paragraph.
\end{quotation}
Another frequently-displayed structure is a list.
The following is an example of an \emph{itemized}
list.
\begin{itemize}
   \item This is the first item of an itemized list.
         Each item in the list is marked with a ``tick''.
         You don't have to worry about what kind of tick
         mark is used.

   \item This is the second item of the list.  It
         contains another list nested inside it.  The inner
         list is an \emph{enumerated} list.
         \begin{enumerate}
            \item This is the first item of an enumerated
                  list that is nested within the itemized list.

            \item This is the second item of the inner list.
                  \LaTeX\ allows you to nest lists deeper than
                  you really should.
         \end{enumerate}
         This is the rest of the second item of the outer
         list.  It is no more interesting than any other
         part of the item.
   \item This is the third item of the list.
\end{itemize}
You can even display poetry.
\begin{verse}
   There is an environment
    for verse \\             % The \\ command separates lines
   Whose features some poets % within a stanza.
   will curse.

                             % One or more blank lines separate stanzas.

   For instead of making\\
   Them do \emph{all} line breaking, \\
   It allows them to put too many words on a line when they'd rather be
   forced to be terse.
\end{verse}

Mathematical formulas may also be displayed.  A
displayed formula
is
one-line long; multiline
formulas require special formatting instructions.
   \[  \ip{\Gamma}{\psi'} = x'' + y^{2} + z_{i}^{n}\]
Don't start a paragraph with a displayed equation,
nor make one a paragraph by itself.

\end{document}               % End of document.

\section{Bench Huge Section 186}
% synthetic bench section 186
% This is a sample LaTeX input file.  (Version of 12 August 2004.)
%
% A '%' character causes TeX to ignore all remaining text on the line,
% and is used for comments like this one.

\documentclass{article}      % Specifies the document class

                             % The preamble begins here.
\title{An Example Document}  % Declares the document's title.
\author{Leslie Lamport}      % Declares the author's name.
\date{January 21, 1994}      % Deleting this command produces today's date.

\newcommand{\ip}[2]{(#1, #2)}
                             % Defines \ip{arg1}{arg2} to mean
                             % (arg1, arg2).

%\newcommand{\ip}[2]{\langle #1 | #2\rangle}
                             % This is an alternative definition of
                             % \ip that is commented out.

\begin{document}             % End of preamble and beginning of text.

\maketitle                   % Produces the title.

This is an example input file.  Comparing it with
the output it generates can show you how to
produce a simple document of your own.

\section{Ordinary Text}      % Produces section heading.  Lower-level
                             % sections are begun with similar
                             % \subsection and \subsubsection commands.

The ends  of words and sentences are marked
  by   spaces. It  doesn't matter how many
spaces    you type; one is as good as 100.  The
end of   a line counts as a space.

One   or more   blank lines denote the  end
of  a paragraph.

Since any number of consecutive spaces are treated
like a single one, the formatting of the input
file makes no difference to
      \LaTeX,                % The \LaTeX command generates the LaTeX logo.
but it makes a difference to you.  When you use
\LaTeX, making your input file as easy to read
as possible will be a great help as you write
your document and when you change it.  This sample
file shows how you can add comments to your own input
file.

Because printing is different from typewriting,
there are a number of things that you have to do
differently when preparing an input file than if
you were just typing the document directly.
Quotation marks like
       ``this''
have to be handled specially, as do quotes within
quotes:
       ``\,`this'            % \, separates the double and single quote.
        is what I just
        wrote, not  `that'\,''.

Dashes come in three sizes: an
       intra-word
dash, a medium dash for number ranges like
       1--2,
and a punctuation
       dash---like
this.

A sentence-ending space should be larger than the
space between words within a sentence.  You
sometimes have to type special commands in
conjunction with punctuation characters to get
this right, as in the following sentence.
       Gnats, gnus, etc.\ all  % `\ ' makes an inter-word space.
       begin with G\@.         % \@ marks end-of-sentence punctuation.
You should check the spaces after periods when
reading your output to make sure you haven't
forgotten any special cases.  Generating an
ellipsis
       \ldots\               % `\ ' is needed after `\ldots' because TeX
                             % ignores spaces after command names like \ldots
                             % made from \ + letters.
                             %
                             % Note how a `%' character causes TeX to ignore
                             % the end of the input line, so these blank lines
                             % do not start a new paragraph.
                             %
with the right spacing around the periods requires
a special command.

\LaTeX\ interprets some common characters as
commands, so you must type special commands to
generate them.  These characters include the
following:
       \$ \& \% \# \{ and \}.

In printing, text is usually emphasized with an
       \emph{italic}
type style.

\begin{em}
   A long segment of text can also be emphasized
   in this way.  Text within such a segment can be
   given \emph{additional} emphasis.
\end{em}

It is sometimes necessary to prevent \LaTeX\ from
breaking a line where it might otherwise do so.
This may be at a space, as between the ``Mr.''\ and
``Jones'' in
       ``Mr.~Jones'',        % ~ produces an unbreakable interword space.
or within a word---especially when the word is a
symbol like
       \mbox{\emph{itemnum}}
that makes little sense when hyphenated across
lines.

Footnotes\footnote{This is an example of a footnote.}
pose no problem.

\LaTeX\ is good at typesetting mathematical formulas
like
       \( x-3y + z = 7 \)
or
       \( a_{1} > x^{2n} + y^{2n} > x' \)
or
       \( \ip{A}{B} = \sum_{i} a_{i} b_{i} \).
The spaces you type in a formula are
ignored.  Remember that a letter like
       $x$                   % $ ... $  and  \( ... \)  are equivalent
is a formula when it denotes a mathematical
symbol, and it should be typed as one.

\section{Displayed Text}

Text is displayed by indenting it from the left
margin.  Quotations are commonly displayed.  There
are short quotations
\begin{quote}
   This is a short quotation.  It consists of a
   single paragraph of text.  See how it is formatted.
\end{quote}
and longer ones.
\begin{quotation}
   This is a longer quotation.  It consists of two
   paragraphs of text, neither of which are
   particularly interesting.

   This is the second paragraph of the quotation.  It
   is just as dull as the first paragraph.
\end{quotation}
Another frequently-displayed structure is a list.
The following is an example of an \emph{itemized}
list.
\begin{itemize}
   \item This is the first item of an itemized list.
         Each item in the list is marked with a ``tick''.
         You don't have to worry about what kind of tick
         mark is used.

   \item This is the second item of the list.  It
         contains another list nested inside it.  The inner
         list is an \emph{enumerated} list.
         \begin{enumerate}
            \item This is the first item of an enumerated
                  list that is nested within the itemized list.

            \item This is the second item of the inner list.
                  \LaTeX\ allows you to nest lists deeper than
                  you really should.
         \end{enumerate}
         This is the rest of the second item of the outer
         list.  It is no more interesting than any other
         part of the item.
   \item This is the third item of the list.
\end{itemize}
You can even display poetry.
\begin{verse}
   There is an environment
    for verse \\             % The \\ command separates lines
   Whose features some poets % within a stanza.
   will curse.

                             % One or more blank lines separate stanzas.

   For instead of making\\
   Them do \emph{all} line breaking, \\
   It allows them to put too many words on a line when they'd rather be
   forced to be terse.
\end{verse}

Mathematical formulas may also be displayed.  A
displayed formula
is
one-line long; multiline
formulas require special formatting instructions.
   \[  \ip{\Gamma}{\psi'} = x'' + y^{2} + z_{i}^{n}\]
Don't start a paragraph with a displayed equation,
nor make one a paragraph by itself.

\end{document}               % End of document.

\section{Bench Huge Section 187}
% synthetic bench section 187
% This is a sample LaTeX input file.  (Version of 12 August 2004.)
%
% A '%' character causes TeX to ignore all remaining text on the line,
% and is used for comments like this one.

\documentclass{article}      % Specifies the document class

                             % The preamble begins here.
\title{An Example Document}  % Declares the document's title.
\author{Leslie Lamport}      % Declares the author's name.
\date{January 21, 1994}      % Deleting this command produces today's date.

\newcommand{\ip}[2]{(#1, #2)}
                             % Defines \ip{arg1}{arg2} to mean
                             % (arg1, arg2).

%\newcommand{\ip}[2]{\langle #1 | #2\rangle}
                             % This is an alternative definition of
                             % \ip that is commented out.

\begin{document}             % End of preamble and beginning of text.

\maketitle                   % Produces the title.

This is an example input file.  Comparing it with
the output it generates can show you how to
produce a simple document of your own.

\section{Ordinary Text}      % Produces section heading.  Lower-level
                             % sections are begun with similar
                             % \subsection and \subsubsection commands.

The ends  of words and sentences are marked
  by   spaces. It  doesn't matter how many
spaces    you type; one is as good as 100.  The
end of   a line counts as a space.

One   or more   blank lines denote the  end
of  a paragraph.

Since any number of consecutive spaces are treated
like a single one, the formatting of the input
file makes no difference to
      \LaTeX,                % The \LaTeX command generates the LaTeX logo.
but it makes a difference to you.  When you use
\LaTeX, making your input file as easy to read
as possible will be a great help as you write
your document and when you change it.  This sample
file shows how you can add comments to your own input
file.

Because printing is different from typewriting,
there are a number of things that you have to do
differently when preparing an input file than if
you were just typing the document directly.
Quotation marks like
       ``this''
have to be handled specially, as do quotes within
quotes:
       ``\,`this'            % \, separates the double and single quote.
        is what I just
        wrote, not  `that'\,''.

Dashes come in three sizes: an
       intra-word
dash, a medium dash for number ranges like
       1--2,
and a punctuation
       dash---like
this.

A sentence-ending space should be larger than the
space between words within a sentence.  You
sometimes have to type special commands in
conjunction with punctuation characters to get
this right, as in the following sentence.
       Gnats, gnus, etc.\ all  % `\ ' makes an inter-word space.
       begin with G\@.         % \@ marks end-of-sentence punctuation.
You should check the spaces after periods when
reading your output to make sure you haven't
forgotten any special cases.  Generating an
ellipsis
       \ldots\               % `\ ' is needed after `\ldots' because TeX
                             % ignores spaces after command names like \ldots
                             % made from \ + letters.
                             %
                             % Note how a `%' character causes TeX to ignore
                             % the end of the input line, so these blank lines
                             % do not start a new paragraph.
                             %
with the right spacing around the periods requires
a special command.

\LaTeX\ interprets some common characters as
commands, so you must type special commands to
generate them.  These characters include the
following:
       \$ \& \% \# \{ and \}.

In printing, text is usually emphasized with an
       \emph{italic}
type style.

\begin{em}
   A long segment of text can also be emphasized
   in this way.  Text within such a segment can be
   given \emph{additional} emphasis.
\end{em}

It is sometimes necessary to prevent \LaTeX\ from
breaking a line where it might otherwise do so.
This may be at a space, as between the ``Mr.''\ and
``Jones'' in
       ``Mr.~Jones'',        % ~ produces an unbreakable interword space.
or within a word---especially when the word is a
symbol like
       \mbox{\emph{itemnum}}
that makes little sense when hyphenated across
lines.

Footnotes\footnote{This is an example of a footnote.}
pose no problem.

\LaTeX\ is good at typesetting mathematical formulas
like
       \( x-3y + z = 7 \)
or
       \( a_{1} > x^{2n} + y^{2n} > x' \)
or
       \( \ip{A}{B} = \sum_{i} a_{i} b_{i} \).
The spaces you type in a formula are
ignored.  Remember that a letter like
       $x$                   % $ ... $  and  \( ... \)  are equivalent
is a formula when it denotes a mathematical
symbol, and it should be typed as one.

\section{Displayed Text}

Text is displayed by indenting it from the left
margin.  Quotations are commonly displayed.  There
are short quotations
\begin{quote}
   This is a short quotation.  It consists of a
   single paragraph of text.  See how it is formatted.
\end{quote}
and longer ones.
\begin{quotation}
   This is a longer quotation.  It consists of two
   paragraphs of text, neither of which are
   particularly interesting.

   This is the second paragraph of the quotation.  It
   is just as dull as the first paragraph.
\end{quotation}
Another frequently-displayed structure is a list.
The following is an example of an \emph{itemized}
list.
\begin{itemize}
   \item This is the first item of an itemized list.
         Each item in the list is marked with a ``tick''.
         You don't have to worry about what kind of tick
         mark is used.

   \item This is the second item of the list.  It
         contains another list nested inside it.  The inner
         list is an \emph{enumerated} list.
         \begin{enumerate}
            \item This is the first item of an enumerated
                  list that is nested within the itemized list.

            \item This is the second item of the inner list.
                  \LaTeX\ allows you to nest lists deeper than
                  you really should.
         \end{enumerate}
         This is the rest of the second item of the outer
         list.  It is no more interesting than any other
         part of the item.
   \item This is the third item of the list.
\end{itemize}
You can even display poetry.
\begin{verse}
   There is an environment
    for verse \\             % The \\ command separates lines
   Whose features some poets % within a stanza.
   will curse.

                             % One or more blank lines separate stanzas.

   For instead of making\\
   Them do \emph{all} line breaking, \\
   It allows them to put too many words on a line when they'd rather be
   forced to be terse.
\end{verse}

Mathematical formulas may also be displayed.  A
displayed formula
is
one-line long; multiline
formulas require special formatting instructions.
   \[  \ip{\Gamma}{\psi'} = x'' + y^{2} + z_{i}^{n}\]
Don't start a paragraph with a displayed equation,
nor make one a paragraph by itself.

\end{document}               % End of document.

\section{Bench Huge Section 188}
% synthetic bench section 188
% This is a sample LaTeX input file.  (Version of 12 August 2004.)
%
% A '%' character causes TeX to ignore all remaining text on the line,
% and is used for comments like this one.

\documentclass{article}      % Specifies the document class

                             % The preamble begins here.
\title{An Example Document}  % Declares the document's title.
\author{Leslie Lamport}      % Declares the author's name.
\date{January 21, 1994}      % Deleting this command produces today's date.

\newcommand{\ip}[2]{(#1, #2)}
                             % Defines \ip{arg1}{arg2} to mean
                             % (arg1, arg2).

%\newcommand{\ip}[2]{\langle #1 | #2\rangle}
                             % This is an alternative definition of
                             % \ip that is commented out.

\begin{document}             % End of preamble and beginning of text.

\maketitle                   % Produces the title.

This is an example input file.  Comparing it with
the output it generates can show you how to
produce a simple document of your own.

\section{Ordinary Text}      % Produces section heading.  Lower-level
                             % sections are begun with similar
                             % \subsection and \subsubsection commands.

The ends  of words and sentences are marked
  by   spaces. It  doesn't matter how many
spaces    you type; one is as good as 100.  The
end of   a line counts as a space.

One   or more   blank lines denote the  end
of  a paragraph.

Since any number of consecutive spaces are treated
like a single one, the formatting of the input
file makes no difference to
      \LaTeX,                % The \LaTeX command generates the LaTeX logo.
but it makes a difference to you.  When you use
\LaTeX, making your input file as easy to read
as possible will be a great help as you write
your document and when you change it.  This sample
file shows how you can add comments to your own input
file.

Because printing is different from typewriting,
there are a number of things that you have to do
differently when preparing an input file than if
you were just typing the document directly.
Quotation marks like
       ``this''
have to be handled specially, as do quotes within
quotes:
       ``\,`this'            % \, separates the double and single quote.
        is what I just
        wrote, not  `that'\,''.

Dashes come in three sizes: an
       intra-word
dash, a medium dash for number ranges like
       1--2,
and a punctuation
       dash---like
this.

A sentence-ending space should be larger than the
space between words within a sentence.  You
sometimes have to type special commands in
conjunction with punctuation characters to get
this right, as in the following sentence.
       Gnats, gnus, etc.\ all  % `\ ' makes an inter-word space.
       begin with G\@.         % \@ marks end-of-sentence punctuation.
You should check the spaces after periods when
reading your output to make sure you haven't
forgotten any special cases.  Generating an
ellipsis
       \ldots\               % `\ ' is needed after `\ldots' because TeX
                             % ignores spaces after command names like \ldots
                             % made from \ + letters.
                             %
                             % Note how a `%' character causes TeX to ignore
                             % the end of the input line, so these blank lines
                             % do not start a new paragraph.
                             %
with the right spacing around the periods requires
a special command.

\LaTeX\ interprets some common characters as
commands, so you must type special commands to
generate them.  These characters include the
following:
       \$ \& \% \# \{ and \}.

In printing, text is usually emphasized with an
       \emph{italic}
type style.

\begin{em}
   A long segment of text can also be emphasized
   in this way.  Text within such a segment can be
   given \emph{additional} emphasis.
\end{em}

It is sometimes necessary to prevent \LaTeX\ from
breaking a line where it might otherwise do so.
This may be at a space, as between the ``Mr.''\ and
``Jones'' in
       ``Mr.~Jones'',        % ~ produces an unbreakable interword space.
or within a word---especially when the word is a
symbol like
       \mbox{\emph{itemnum}}
that makes little sense when hyphenated across
lines.

Footnotes\footnote{This is an example of a footnote.}
pose no problem.

\LaTeX\ is good at typesetting mathematical formulas
like
       \( x-3y + z = 7 \)
or
       \( a_{1} > x^{2n} + y^{2n} > x' \)
or
       \( \ip{A}{B} = \sum_{i} a_{i} b_{i} \).
The spaces you type in a formula are
ignored.  Remember that a letter like
       $x$                   % $ ... $  and  \( ... \)  are equivalent
is a formula when it denotes a mathematical
symbol, and it should be typed as one.

\section{Displayed Text}

Text is displayed by indenting it from the left
margin.  Quotations are commonly displayed.  There
are short quotations
\begin{quote}
   This is a short quotation.  It consists of a
   single paragraph of text.  See how it is formatted.
\end{quote}
and longer ones.
\begin{quotation}
   This is a longer quotation.  It consists of two
   paragraphs of text, neither of which are
   particularly interesting.

   This is the second paragraph of the quotation.  It
   is just as dull as the first paragraph.
\end{quotation}
Another frequently-displayed structure is a list.
The following is an example of an \emph{itemized}
list.
\begin{itemize}
   \item This is the first item of an itemized list.
         Each item in the list is marked with a ``tick''.
         You don't have to worry about what kind of tick
         mark is used.

   \item This is the second item of the list.  It
         contains another list nested inside it.  The inner
         list is an \emph{enumerated} list.
         \begin{enumerate}
            \item This is the first item of an enumerated
                  list that is nested within the itemized list.

            \item This is the second item of the inner list.
                  \LaTeX\ allows you to nest lists deeper than
                  you really should.
         \end{enumerate}
         This is the rest of the second item of the outer
         list.  It is no more interesting than any other
         part of the item.
   \item This is the third item of the list.
\end{itemize}
You can even display poetry.
\begin{verse}
   There is an environment
    for verse \\             % The \\ command separates lines
   Whose features some poets % within a stanza.
   will curse.

                             % One or more blank lines separate stanzas.

   For instead of making\\
   Them do \emph{all} line breaking, \\
   It allows them to put too many words on a line when they'd rather be
   forced to be terse.
\end{verse}

Mathematical formulas may also be displayed.  A
displayed formula
is
one-line long; multiline
formulas require special formatting instructions.
   \[  \ip{\Gamma}{\psi'} = x'' + y^{2} + z_{i}^{n}\]
Don't start a paragraph with a displayed equation,
nor make one a paragraph by itself.

\end{document}               % End of document.

\section{Bench Huge Section 189}
% synthetic bench section 189
% This is a sample LaTeX input file.  (Version of 12 August 2004.)
%
% A '%' character causes TeX to ignore all remaining text on the line,
% and is used for comments like this one.

\documentclass{article}      % Specifies the document class

                             % The preamble begins here.
\title{An Example Document}  % Declares the document's title.
\author{Leslie Lamport}      % Declares the author's name.
\date{January 21, 1994}      % Deleting this command produces today's date.

\newcommand{\ip}[2]{(#1, #2)}
                             % Defines \ip{arg1}{arg2} to mean
                             % (arg1, arg2).

%\newcommand{\ip}[2]{\langle #1 | #2\rangle}
                             % This is an alternative definition of
                             % \ip that is commented out.

\begin{document}             % End of preamble and beginning of text.

\maketitle                   % Produces the title.

This is an example input file.  Comparing it with
the output it generates can show you how to
produce a simple document of your own.

\section{Ordinary Text}      % Produces section heading.  Lower-level
                             % sections are begun with similar
                             % \subsection and \subsubsection commands.

The ends  of words and sentences are marked
  by   spaces. It  doesn't matter how many
spaces    you type; one is as good as 100.  The
end of   a line counts as a space.

One   or more   blank lines denote the  end
of  a paragraph.

Since any number of consecutive spaces are treated
like a single one, the formatting of the input
file makes no difference to
      \LaTeX,                % The \LaTeX command generates the LaTeX logo.
but it makes a difference to you.  When you use
\LaTeX, making your input file as easy to read
as possible will be a great help as you write
your document and when you change it.  This sample
file shows how you can add comments to your own input
file.

Because printing is different from typewriting,
there are a number of things that you have to do
differently when preparing an input file than if
you were just typing the document directly.
Quotation marks like
       ``this''
have to be handled specially, as do quotes within
quotes:
       ``\,`this'            % \, separates the double and single quote.
        is what I just
        wrote, not  `that'\,''.

Dashes come in three sizes: an
       intra-word
dash, a medium dash for number ranges like
       1--2,
and a punctuation
       dash---like
this.

A sentence-ending space should be larger than the
space between words within a sentence.  You
sometimes have to type special commands in
conjunction with punctuation characters to get
this right, as in the following sentence.
       Gnats, gnus, etc.\ all  % `\ ' makes an inter-word space.
       begin with G\@.         % \@ marks end-of-sentence punctuation.
You should check the spaces after periods when
reading your output to make sure you haven't
forgotten any special cases.  Generating an
ellipsis
       \ldots\               % `\ ' is needed after `\ldots' because TeX
                             % ignores spaces after command names like \ldots
                             % made from \ + letters.
                             %
                             % Note how a `%' character causes TeX to ignore
                             % the end of the input line, so these blank lines
                             % do not start a new paragraph.
                             %
with the right spacing around the periods requires
a special command.

\LaTeX\ interprets some common characters as
commands, so you must type special commands to
generate them.  These characters include the
following:
       \$ \& \% \# \{ and \}.

In printing, text is usually emphasized with an
       \emph{italic}
type style.

\begin{em}
   A long segment of text can also be emphasized
   in this way.  Text within such a segment can be
   given \emph{additional} emphasis.
\end{em}

It is sometimes necessary to prevent \LaTeX\ from
breaking a line where it might otherwise do so.
This may be at a space, as between the ``Mr.''\ and
``Jones'' in
       ``Mr.~Jones'',        % ~ produces an unbreakable interword space.
or within a word---especially when the word is a
symbol like
       \mbox{\emph{itemnum}}
that makes little sense when hyphenated across
lines.

Footnotes\footnote{This is an example of a footnote.}
pose no problem.

\LaTeX\ is good at typesetting mathematical formulas
like
       \( x-3y + z = 7 \)
or
       \( a_{1} > x^{2n} + y^{2n} > x' \)
or
       \( \ip{A}{B} = \sum_{i} a_{i} b_{i} \).
The spaces you type in a formula are
ignored.  Remember that a letter like
       $x$                   % $ ... $  and  \( ... \)  are equivalent
is a formula when it denotes a mathematical
symbol, and it should be typed as one.

\section{Displayed Text}

Text is displayed by indenting it from the left
margin.  Quotations are commonly displayed.  There
are short quotations
\begin{quote}
   This is a short quotation.  It consists of a
   single paragraph of text.  See how it is formatted.
\end{quote}
and longer ones.
\begin{quotation}
   This is a longer quotation.  It consists of two
   paragraphs of text, neither of which are
   particularly interesting.

   This is the second paragraph of the quotation.  It
   is just as dull as the first paragraph.
\end{quotation}
Another frequently-displayed structure is a list.
The following is an example of an \emph{itemized}
list.
\begin{itemize}
   \item This is the first item of an itemized list.
         Each item in the list is marked with a ``tick''.
         You don't have to worry about what kind of tick
         mark is used.

   \item This is the second item of the list.  It
         contains another list nested inside it.  The inner
         list is an \emph{enumerated} list.
         \begin{enumerate}
            \item This is the first item of an enumerated
                  list that is nested within the itemized list.

            \item This is the second item of the inner list.
                  \LaTeX\ allows you to nest lists deeper than
                  you really should.
         \end{enumerate}
         This is the rest of the second item of the outer
         list.  It is no more interesting than any other
         part of the item.
   \item This is the third item of the list.
\end{itemize}
You can even display poetry.
\begin{verse}
   There is an environment
    for verse \\             % The \\ command separates lines
   Whose features some poets % within a stanza.
   will curse.

                             % One or more blank lines separate stanzas.

   For instead of making\\
   Them do \emph{all} line breaking, \\
   It allows them to put too many words on a line when they'd rather be
   forced to be terse.
\end{verse}

Mathematical formulas may also be displayed.  A
displayed formula
is
one-line long; multiline
formulas require special formatting instructions.
   \[  \ip{\Gamma}{\psi'} = x'' + y^{2} + z_{i}^{n}\]
Don't start a paragraph with a displayed equation,
nor make one a paragraph by itself.

\end{document}               % End of document.

\section{Bench Huge Section 190}
% synthetic bench section 190
% This is a sample LaTeX input file.  (Version of 12 August 2004.)
%
% A '%' character causes TeX to ignore all remaining text on the line,
% and is used for comments like this one.

\documentclass{article}      % Specifies the document class

                             % The preamble begins here.
\title{An Example Document}  % Declares the document's title.
\author{Leslie Lamport}      % Declares the author's name.
\date{January 21, 1994}      % Deleting this command produces today's date.

\newcommand{\ip}[2]{(#1, #2)}
                             % Defines \ip{arg1}{arg2} to mean
                             % (arg1, arg2).

%\newcommand{\ip}[2]{\langle #1 | #2\rangle}
                             % This is an alternative definition of
                             % \ip that is commented out.

\begin{document}             % End of preamble and beginning of text.

\maketitle                   % Produces the title.

This is an example input file.  Comparing it with
the output it generates can show you how to
produce a simple document of your own.

\section{Ordinary Text}      % Produces section heading.  Lower-level
                             % sections are begun with similar
                             % \subsection and \subsubsection commands.

The ends  of words and sentences are marked
  by   spaces. It  doesn't matter how many
spaces    you type; one is as good as 100.  The
end of   a line counts as a space.

One   or more   blank lines denote the  end
of  a paragraph.

Since any number of consecutive spaces are treated
like a single one, the formatting of the input
file makes no difference to
      \LaTeX,                % The \LaTeX command generates the LaTeX logo.
but it makes a difference to you.  When you use
\LaTeX, making your input file as easy to read
as possible will be a great help as you write
your document and when you change it.  This sample
file shows how you can add comments to your own input
file.

Because printing is different from typewriting,
there are a number of things that you have to do
differently when preparing an input file than if
you were just typing the document directly.
Quotation marks like
       ``this''
have to be handled specially, as do quotes within
quotes:
       ``\,`this'            % \, separates the double and single quote.
        is what I just
        wrote, not  `that'\,''.

Dashes come in three sizes: an
       intra-word
dash, a medium dash for number ranges like
       1--2,
and a punctuation
       dash---like
this.

A sentence-ending space should be larger than the
space between words within a sentence.  You
sometimes have to type special commands in
conjunction with punctuation characters to get
this right, as in the following sentence.
       Gnats, gnus, etc.\ all  % `\ ' makes an inter-word space.
       begin with G\@.         % \@ marks end-of-sentence punctuation.
You should check the spaces after periods when
reading your output to make sure you haven't
forgotten any special cases.  Generating an
ellipsis
       \ldots\               % `\ ' is needed after `\ldots' because TeX
                             % ignores spaces after command names like \ldots
                             % made from \ + letters.
                             %
                             % Note how a `%' character causes TeX to ignore
                             % the end of the input line, so these blank lines
                             % do not start a new paragraph.
                             %
with the right spacing around the periods requires
a special command.

\LaTeX\ interprets some common characters as
commands, so you must type special commands to
generate them.  These characters include the
following:
       \$ \& \% \# \{ and \}.

In printing, text is usually emphasized with an
       \emph{italic}
type style.

\begin{em}
   A long segment of text can also be emphasized
   in this way.  Text within such a segment can be
   given \emph{additional} emphasis.
\end{em}

It is sometimes necessary to prevent \LaTeX\ from
breaking a line where it might otherwise do so.
This may be at a space, as between the ``Mr.''\ and
``Jones'' in
       ``Mr.~Jones'',        % ~ produces an unbreakable interword space.
or within a word---especially when the word is a
symbol like
       \mbox{\emph{itemnum}}
that makes little sense when hyphenated across
lines.

Footnotes\footnote{This is an example of a footnote.}
pose no problem.

\LaTeX\ is good at typesetting mathematical formulas
like
       \( x-3y + z = 7 \)
or
       \( a_{1} > x^{2n} + y^{2n} > x' \)
or
       \( \ip{A}{B} = \sum_{i} a_{i} b_{i} \).
The spaces you type in a formula are
ignored.  Remember that a letter like
       $x$                   % $ ... $  and  \( ... \)  are equivalent
is a formula when it denotes a mathematical
symbol, and it should be typed as one.

\section{Displayed Text}

Text is displayed by indenting it from the left
margin.  Quotations are commonly displayed.  There
are short quotations
\begin{quote}
   This is a short quotation.  It consists of a
   single paragraph of text.  See how it is formatted.
\end{quote}
and longer ones.
\begin{quotation}
   This is a longer quotation.  It consists of two
   paragraphs of text, neither of which are
   particularly interesting.

   This is the second paragraph of the quotation.  It
   is just as dull as the first paragraph.
\end{quotation}
Another frequently-displayed structure is a list.
The following is an example of an \emph{itemized}
list.
\begin{itemize}
   \item This is the first item of an itemized list.
         Each item in the list is marked with a ``tick''.
         You don't have to worry about what kind of tick
         mark is used.

   \item This is the second item of the list.  It
         contains another list nested inside it.  The inner
         list is an \emph{enumerated} list.
         \begin{enumerate}
            \item This is the first item of an enumerated
                  list that is nested within the itemized list.

            \item This is the second item of the inner list.
                  \LaTeX\ allows you to nest lists deeper than
                  you really should.
         \end{enumerate}
         This is the rest of the second item of the outer
         list.  It is no more interesting than any other
         part of the item.
   \item This is the third item of the list.
\end{itemize}
You can even display poetry.
\begin{verse}
   There is an environment
    for verse \\             % The \\ command separates lines
   Whose features some poets % within a stanza.
   will curse.

                             % One or more blank lines separate stanzas.

   For instead of making\\
   Them do \emph{all} line breaking, \\
   It allows them to put too many words on a line when they'd rather be
   forced to be terse.
\end{verse}

Mathematical formulas may also be displayed.  A
displayed formula
is
one-line long; multiline
formulas require special formatting instructions.
   \[  \ip{\Gamma}{\psi'} = x'' + y^{2} + z_{i}^{n}\]
Don't start a paragraph with a displayed equation,
nor make one a paragraph by itself.

\end{document}               % End of document.

\section{Bench Huge Section 191}
% synthetic bench section 191
% This is a sample LaTeX input file.  (Version of 12 August 2004.)
%
% A '%' character causes TeX to ignore all remaining text on the line,
% and is used for comments like this one.

\documentclass{article}      % Specifies the document class

                             % The preamble begins here.
\title{An Example Document}  % Declares the document's title.
\author{Leslie Lamport}      % Declares the author's name.
\date{January 21, 1994}      % Deleting this command produces today's date.

\newcommand{\ip}[2]{(#1, #2)}
                             % Defines \ip{arg1}{arg2} to mean
                             % (arg1, arg2).

%\newcommand{\ip}[2]{\langle #1 | #2\rangle}
                             % This is an alternative definition of
                             % \ip that is commented out.

\begin{document}             % End of preamble and beginning of text.

\maketitle                   % Produces the title.

This is an example input file.  Comparing it with
the output it generates can show you how to
produce a simple document of your own.

\section{Ordinary Text}      % Produces section heading.  Lower-level
                             % sections are begun with similar
                             % \subsection and \subsubsection commands.

The ends  of words and sentences are marked
  by   spaces. It  doesn't matter how many
spaces    you type; one is as good as 100.  The
end of   a line counts as a space.

One   or more   blank lines denote the  end
of  a paragraph.

Since any number of consecutive spaces are treated
like a single one, the formatting of the input
file makes no difference to
      \LaTeX,                % The \LaTeX command generates the LaTeX logo.
but it makes a difference to you.  When you use
\LaTeX, making your input file as easy to read
as possible will be a great help as you write
your document and when you change it.  This sample
file shows how you can add comments to your own input
file.

Because printing is different from typewriting,
there are a number of things that you have to do
differently when preparing an input file than if
you were just typing the document directly.
Quotation marks like
       ``this''
have to be handled specially, as do quotes within
quotes:
       ``\,`this'            % \, separates the double and single quote.
        is what I just
        wrote, not  `that'\,''.

Dashes come in three sizes: an
       intra-word
dash, a medium dash for number ranges like
       1--2,
and a punctuation
       dash---like
this.

A sentence-ending space should be larger than the
space between words within a sentence.  You
sometimes have to type special commands in
conjunction with punctuation characters to get
this right, as in the following sentence.
       Gnats, gnus, etc.\ all  % `\ ' makes an inter-word space.
       begin with G\@.         % \@ marks end-of-sentence punctuation.
You should check the spaces after periods when
reading your output to make sure you haven't
forgotten any special cases.  Generating an
ellipsis
       \ldots\               % `\ ' is needed after `\ldots' because TeX
                             % ignores spaces after command names like \ldots
                             % made from \ + letters.
                             %
                             % Note how a `%' character causes TeX to ignore
                             % the end of the input line, so these blank lines
                             % do not start a new paragraph.
                             %
with the right spacing around the periods requires
a special command.

\LaTeX\ interprets some common characters as
commands, so you must type special commands to
generate them.  These characters include the
following:
       \$ \& \% \# \{ and \}.

In printing, text is usually emphasized with an
       \emph{italic}
type style.

\begin{em}
   A long segment of text can also be emphasized
   in this way.  Text within such a segment can be
   given \emph{additional} emphasis.
\end{em}

It is sometimes necessary to prevent \LaTeX\ from
breaking a line where it might otherwise do so.
This may be at a space, as between the ``Mr.''\ and
``Jones'' in
       ``Mr.~Jones'',        % ~ produces an unbreakable interword space.
or within a word---especially when the word is a
symbol like
       \mbox{\emph{itemnum}}
that makes little sense when hyphenated across
lines.

Footnotes\footnote{This is an example of a footnote.}
pose no problem.

\LaTeX\ is good at typesetting mathematical formulas
like
       \( x-3y + z = 7 \)
or
       \( a_{1} > x^{2n} + y^{2n} > x' \)
or
       \( \ip{A}{B} = \sum_{i} a_{i} b_{i} \).
The spaces you type in a formula are
ignored.  Remember that a letter like
       $x$                   % $ ... $  and  \( ... \)  are equivalent
is a formula when it denotes a mathematical
symbol, and it should be typed as one.

\section{Displayed Text}

Text is displayed by indenting it from the left
margin.  Quotations are commonly displayed.  There
are short quotations
\begin{quote}
   This is a short quotation.  It consists of a
   single paragraph of text.  See how it is formatted.
\end{quote}
and longer ones.
\begin{quotation}
   This is a longer quotation.  It consists of two
   paragraphs of text, neither of which are
   particularly interesting.

   This is the second paragraph of the quotation.  It
   is just as dull as the first paragraph.
\end{quotation}
Another frequently-displayed structure is a list.
The following is an example of an \emph{itemized}
list.
\begin{itemize}
   \item This is the first item of an itemized list.
         Each item in the list is marked with a ``tick''.
         You don't have to worry about what kind of tick
         mark is used.

   \item This is the second item of the list.  It
         contains another list nested inside it.  The inner
         list is an \emph{enumerated} list.
         \begin{enumerate}
            \item This is the first item of an enumerated
                  list that is nested within the itemized list.

            \item This is the second item of the inner list.
                  \LaTeX\ allows you to nest lists deeper than
                  you really should.
         \end{enumerate}
         This is the rest of the second item of the outer
         list.  It is no more interesting than any other
         part of the item.
   \item This is the third item of the list.
\end{itemize}
You can even display poetry.
\begin{verse}
   There is an environment
    for verse \\             % The \\ command separates lines
   Whose features some poets % within a stanza.
   will curse.

                             % One or more blank lines separate stanzas.

   For instead of making\\
   Them do \emph{all} line breaking, \\
   It allows them to put too many words on a line when they'd rather be
   forced to be terse.
\end{verse}

Mathematical formulas may also be displayed.  A
displayed formula
is
one-line long; multiline
formulas require special formatting instructions.
   \[  \ip{\Gamma}{\psi'} = x'' + y^{2} + z_{i}^{n}\]
Don't start a paragraph with a displayed equation,
nor make one a paragraph by itself.

\end{document}               % End of document.

\section{Bench Huge Section 192}
% synthetic bench section 192
% This is a sample LaTeX input file.  (Version of 12 August 2004.)
%
% A '%' character causes TeX to ignore all remaining text on the line,
% and is used for comments like this one.

\documentclass{article}      % Specifies the document class

                             % The preamble begins here.
\title{An Example Document}  % Declares the document's title.
\author{Leslie Lamport}      % Declares the author's name.
\date{January 21, 1994}      % Deleting this command produces today's date.

\newcommand{\ip}[2]{(#1, #2)}
                             % Defines \ip{arg1}{arg2} to mean
                             % (arg1, arg2).

%\newcommand{\ip}[2]{\langle #1 | #2\rangle}
                             % This is an alternative definition of
                             % \ip that is commented out.

\begin{document}             % End of preamble and beginning of text.

\maketitle                   % Produces the title.

This is an example input file.  Comparing it with
the output it generates can show you how to
produce a simple document of your own.

\section{Ordinary Text}      % Produces section heading.  Lower-level
                             % sections are begun with similar
                             % \subsection and \subsubsection commands.

The ends  of words and sentences are marked
  by   spaces. It  doesn't matter how many
spaces    you type; one is as good as 100.  The
end of   a line counts as a space.

One   or more   blank lines denote the  end
of  a paragraph.

Since any number of consecutive spaces are treated
like a single one, the formatting of the input
file makes no difference to
      \LaTeX,                % The \LaTeX command generates the LaTeX logo.
but it makes a difference to you.  When you use
\LaTeX, making your input file as easy to read
as possible will be a great help as you write
your document and when you change it.  This sample
file shows how you can add comments to your own input
file.

Because printing is different from typewriting,
there are a number of things that you have to do
differently when preparing an input file than if
you were just typing the document directly.
Quotation marks like
       ``this''
have to be handled specially, as do quotes within
quotes:
       ``\,`this'            % \, separates the double and single quote.
        is what I just
        wrote, not  `that'\,''.

Dashes come in three sizes: an
       intra-word
dash, a medium dash for number ranges like
       1--2,
and a punctuation
       dash---like
this.

A sentence-ending space should be larger than the
space between words within a sentence.  You
sometimes have to type special commands in
conjunction with punctuation characters to get
this right, as in the following sentence.
       Gnats, gnus, etc.\ all  % `\ ' makes an inter-word space.
       begin with G\@.         % \@ marks end-of-sentence punctuation.
You should check the spaces after periods when
reading your output to make sure you haven't
forgotten any special cases.  Generating an
ellipsis
       \ldots\               % `\ ' is needed after `\ldots' because TeX
                             % ignores spaces after command names like \ldots
                             % made from \ + letters.
                             %
                             % Note how a `%' character causes TeX to ignore
                             % the end of the input line, so these blank lines
                             % do not start a new paragraph.
                             %
with the right spacing around the periods requires
a special command.

\LaTeX\ interprets some common characters as
commands, so you must type special commands to
generate them.  These characters include the
following:
       \$ \& \% \# \{ and \}.

In printing, text is usually emphasized with an
       \emph{italic}
type style.

\begin{em}
   A long segment of text can also be emphasized
   in this way.  Text within such a segment can be
   given \emph{additional} emphasis.
\end{em}

It is sometimes necessary to prevent \LaTeX\ from
breaking a line where it might otherwise do so.
This may be at a space, as between the ``Mr.''\ and
``Jones'' in
       ``Mr.~Jones'',        % ~ produces an unbreakable interword space.
or within a word---especially when the word is a
symbol like
       \mbox{\emph{itemnum}}
that makes little sense when hyphenated across
lines.

Footnotes\footnote{This is an example of a footnote.}
pose no problem.

\LaTeX\ is good at typesetting mathematical formulas
like
       \( x-3y + z = 7 \)
or
       \( a_{1} > x^{2n} + y^{2n} > x' \)
or
       \( \ip{A}{B} = \sum_{i} a_{i} b_{i} \).
The spaces you type in a formula are
ignored.  Remember that a letter like
       $x$                   % $ ... $  and  \( ... \)  are equivalent
is a formula when it denotes a mathematical
symbol, and it should be typed as one.

\section{Displayed Text}

Text is displayed by indenting it from the left
margin.  Quotations are commonly displayed.  There
are short quotations
\begin{quote}
   This is a short quotation.  It consists of a
   single paragraph of text.  See how it is formatted.
\end{quote}
and longer ones.
\begin{quotation}
   This is a longer quotation.  It consists of two
   paragraphs of text, neither of which are
   particularly interesting.

   This is the second paragraph of the quotation.  It
   is just as dull as the first paragraph.
\end{quotation}
Another frequently-displayed structure is a list.
The following is an example of an \emph{itemized}
list.
\begin{itemize}
   \item This is the first item of an itemized list.
         Each item in the list is marked with a ``tick''.
         You don't have to worry about what kind of tick
         mark is used.

   \item This is the second item of the list.  It
         contains another list nested inside it.  The inner
         list is an \emph{enumerated} list.
         \begin{enumerate}
            \item This is the first item of an enumerated
                  list that is nested within the itemized list.

            \item This is the second item of the inner list.
                  \LaTeX\ allows you to nest lists deeper than
                  you really should.
         \end{enumerate}
         This is the rest of the second item of the outer
         list.  It is no more interesting than any other
         part of the item.
   \item This is the third item of the list.
\end{itemize}
You can even display poetry.
\begin{verse}
   There is an environment
    for verse \\             % The \\ command separates lines
   Whose features some poets % within a stanza.
   will curse.

                             % One or more blank lines separate stanzas.

   For instead of making\\
   Them do \emph{all} line breaking, \\
   It allows them to put too many words on a line when they'd rather be
   forced to be terse.
\end{verse}

Mathematical formulas may also be displayed.  A
displayed formula
is
one-line long; multiline
formulas require special formatting instructions.
   \[  \ip{\Gamma}{\psi'} = x'' + y^{2} + z_{i}^{n}\]
Don't start a paragraph with a displayed equation,
nor make one a paragraph by itself.

\end{document}               % End of document.

\section{Bench Huge Section 193}
% synthetic bench section 193
% This is a sample LaTeX input file.  (Version of 12 August 2004.)
%
% A '%' character causes TeX to ignore all remaining text on the line,
% and is used for comments like this one.

\documentclass{article}      % Specifies the document class

                             % The preamble begins here.
\title{An Example Document}  % Declares the document's title.
\author{Leslie Lamport}      % Declares the author's name.
\date{January 21, 1994}      % Deleting this command produces today's date.

\newcommand{\ip}[2]{(#1, #2)}
                             % Defines \ip{arg1}{arg2} to mean
                             % (arg1, arg2).

%\newcommand{\ip}[2]{\langle #1 | #2\rangle}
                             % This is an alternative definition of
                             % \ip that is commented out.

\begin{document}             % End of preamble and beginning of text.

\maketitle                   % Produces the title.

This is an example input file.  Comparing it with
the output it generates can show you how to
produce a simple document of your own.

\section{Ordinary Text}      % Produces section heading.  Lower-level
                             % sections are begun with similar
                             % \subsection and \subsubsection commands.

The ends  of words and sentences are marked
  by   spaces. It  doesn't matter how many
spaces    you type; one is as good as 100.  The
end of   a line counts as a space.

One   or more   blank lines denote the  end
of  a paragraph.

Since any number of consecutive spaces are treated
like a single one, the formatting of the input
file makes no difference to
      \LaTeX,                % The \LaTeX command generates the LaTeX logo.
but it makes a difference to you.  When you use
\LaTeX, making your input file as easy to read
as possible will be a great help as you write
your document and when you change it.  This sample
file shows how you can add comments to your own input
file.

Because printing is different from typewriting,
there are a number of things that you have to do
differently when preparing an input file than if
you were just typing the document directly.
Quotation marks like
       ``this''
have to be handled specially, as do quotes within
quotes:
       ``\,`this'            % \, separates the double and single quote.
        is what I just
        wrote, not  `that'\,''.

Dashes come in three sizes: an
       intra-word
dash, a medium dash for number ranges like
       1--2,
and a punctuation
       dash---like
this.

A sentence-ending space should be larger than the
space between words within a sentence.  You
sometimes have to type special commands in
conjunction with punctuation characters to get
this right, as in the following sentence.
       Gnats, gnus, etc.\ all  % `\ ' makes an inter-word space.
       begin with G\@.         % \@ marks end-of-sentence punctuation.
You should check the spaces after periods when
reading your output to make sure you haven't
forgotten any special cases.  Generating an
ellipsis
       \ldots\               % `\ ' is needed after `\ldots' because TeX
                             % ignores spaces after command names like \ldots
                             % made from \ + letters.
                             %
                             % Note how a `%' character causes TeX to ignore
                             % the end of the input line, so these blank lines
                             % do not start a new paragraph.
                             %
with the right spacing around the periods requires
a special command.

\LaTeX\ interprets some common characters as
commands, so you must type special commands to
generate them.  These characters include the
following:
       \$ \& \% \# \{ and \}.

In printing, text is usually emphasized with an
       \emph{italic}
type style.

\begin{em}
   A long segment of text can also be emphasized
   in this way.  Text within such a segment can be
   given \emph{additional} emphasis.
\end{em}

It is sometimes necessary to prevent \LaTeX\ from
breaking a line where it might otherwise do so.
This may be at a space, as between the ``Mr.''\ and
``Jones'' in
       ``Mr.~Jones'',        % ~ produces an unbreakable interword space.
or within a word---especially when the word is a
symbol like
       \mbox{\emph{itemnum}}
that makes little sense when hyphenated across
lines.

Footnotes\footnote{This is an example of a footnote.}
pose no problem.

\LaTeX\ is good at typesetting mathematical formulas
like
       \( x-3y + z = 7 \)
or
       \( a_{1} > x^{2n} + y^{2n} > x' \)
or
       \( \ip{A}{B} = \sum_{i} a_{i} b_{i} \).
The spaces you type in a formula are
ignored.  Remember that a letter like
       $x$                   % $ ... $  and  \( ... \)  are equivalent
is a formula when it denotes a mathematical
symbol, and it should be typed as one.

\section{Displayed Text}

Text is displayed by indenting it from the left
margin.  Quotations are commonly displayed.  There
are short quotations
\begin{quote}
   This is a short quotation.  It consists of a
   single paragraph of text.  See how it is formatted.
\end{quote}
and longer ones.
\begin{quotation}
   This is a longer quotation.  It consists of two
   paragraphs of text, neither of which are
   particularly interesting.

   This is the second paragraph of the quotation.  It
   is just as dull as the first paragraph.
\end{quotation}
Another frequently-displayed structure is a list.
The following is an example of an \emph{itemized}
list.
\begin{itemize}
   \item This is the first item of an itemized list.
         Each item in the list is marked with a ``tick''.
         You don't have to worry about what kind of tick
         mark is used.

   \item This is the second item of the list.  It
         contains another list nested inside it.  The inner
         list is an \emph{enumerated} list.
         \begin{enumerate}
            \item This is the first item of an enumerated
                  list that is nested within the itemized list.

            \item This is the second item of the inner list.
                  \LaTeX\ allows you to nest lists deeper than
                  you really should.
         \end{enumerate}
         This is the rest of the second item of the outer
         list.  It is no more interesting than any other
         part of the item.
   \item This is the third item of the list.
\end{itemize}
You can even display poetry.
\begin{verse}
   There is an environment
    for verse \\             % The \\ command separates lines
   Whose features some poets % within a stanza.
   will curse.

                             % One or more blank lines separate stanzas.

   For instead of making\\
   Them do \emph{all} line breaking, \\
   It allows them to put too many words on a line when they'd rather be
   forced to be terse.
\end{verse}

Mathematical formulas may also be displayed.  A
displayed formula
is
one-line long; multiline
formulas require special formatting instructions.
   \[  \ip{\Gamma}{\psi'} = x'' + y^{2} + z_{i}^{n}\]
Don't start a paragraph with a displayed equation,
nor make one a paragraph by itself.

\end{document}               % End of document.

\section{Bench Huge Section 194}
% synthetic bench section 194
% This is a sample LaTeX input file.  (Version of 12 August 2004.)
%
% A '%' character causes TeX to ignore all remaining text on the line,
% and is used for comments like this one.

\documentclass{article}      % Specifies the document class

                             % The preamble begins here.
\title{An Example Document}  % Declares the document's title.
\author{Leslie Lamport}      % Declares the author's name.
\date{January 21, 1994}      % Deleting this command produces today's date.

\newcommand{\ip}[2]{(#1, #2)}
                             % Defines \ip{arg1}{arg2} to mean
                             % (arg1, arg2).

%\newcommand{\ip}[2]{\langle #1 | #2\rangle}
                             % This is an alternative definition of
                             % \ip that is commented out.

\begin{document}             % End of preamble and beginning of text.

\maketitle                   % Produces the title.

This is an example input file.  Comparing it with
the output it generates can show you how to
produce a simple document of your own.

\section{Ordinary Text}      % Produces section heading.  Lower-level
                             % sections are begun with similar
                             % \subsection and \subsubsection commands.

The ends  of words and sentences are marked
  by   spaces. It  doesn't matter how many
spaces    you type; one is as good as 100.  The
end of   a line counts as a space.

One   or more   blank lines denote the  end
of  a paragraph.

Since any number of consecutive spaces are treated
like a single one, the formatting of the input
file makes no difference to
      \LaTeX,                % The \LaTeX command generates the LaTeX logo.
but it makes a difference to you.  When you use
\LaTeX, making your input file as easy to read
as possible will be a great help as you write
your document and when you change it.  This sample
file shows how you can add comments to your own input
file.

Because printing is different from typewriting,
there are a number of things that you have to do
differently when preparing an input file than if
you were just typing the document directly.
Quotation marks like
       ``this''
have to be handled specially, as do quotes within
quotes:
       ``\,`this'            % \, separates the double and single quote.
        is what I just
        wrote, not  `that'\,''.

Dashes come in three sizes: an
       intra-word
dash, a medium dash for number ranges like
       1--2,
and a punctuation
       dash---like
this.

A sentence-ending space should be larger than the
space between words within a sentence.  You
sometimes have to type special commands in
conjunction with punctuation characters to get
this right, as in the following sentence.
       Gnats, gnus, etc.\ all  % `\ ' makes an inter-word space.
       begin with G\@.         % \@ marks end-of-sentence punctuation.
You should check the spaces after periods when
reading your output to make sure you haven't
forgotten any special cases.  Generating an
ellipsis
       \ldots\               % `\ ' is needed after `\ldots' because TeX
                             % ignores spaces after command names like \ldots
                             % made from \ + letters.
                             %
                             % Note how a `%' character causes TeX to ignore
                             % the end of the input line, so these blank lines
                             % do not start a new paragraph.
                             %
with the right spacing around the periods requires
a special command.

\LaTeX\ interprets some common characters as
commands, so you must type special commands to
generate them.  These characters include the
following:
       \$ \& \% \# \{ and \}.

In printing, text is usually emphasized with an
       \emph{italic}
type style.

\begin{em}
   A long segment of text can also be emphasized
   in this way.  Text within such a segment can be
   given \emph{additional} emphasis.
\end{em}

It is sometimes necessary to prevent \LaTeX\ from
breaking a line where it might otherwise do so.
This may be at a space, as between the ``Mr.''\ and
``Jones'' in
       ``Mr.~Jones'',        % ~ produces an unbreakable interword space.
or within a word---especially when the word is a
symbol like
       \mbox{\emph{itemnum}}
that makes little sense when hyphenated across
lines.

Footnotes\footnote{This is an example of a footnote.}
pose no problem.

\LaTeX\ is good at typesetting mathematical formulas
like
       \( x-3y + z = 7 \)
or
       \( a_{1} > x^{2n} + y^{2n} > x' \)
or
       \( \ip{A}{B} = \sum_{i} a_{i} b_{i} \).
The spaces you type in a formula are
ignored.  Remember that a letter like
       $x$                   % $ ... $  and  \( ... \)  are equivalent
is a formula when it denotes a mathematical
symbol, and it should be typed as one.

\section{Displayed Text}

Text is displayed by indenting it from the left
margin.  Quotations are commonly displayed.  There
are short quotations
\begin{quote}
   This is a short quotation.  It consists of a
   single paragraph of text.  See how it is formatted.
\end{quote}
and longer ones.
\begin{quotation}
   This is a longer quotation.  It consists of two
   paragraphs of text, neither of which are
   particularly interesting.

   This is the second paragraph of the quotation.  It
   is just as dull as the first paragraph.
\end{quotation}
Another frequently-displayed structure is a list.
The following is an example of an \emph{itemized}
list.
\begin{itemize}
   \item This is the first item of an itemized list.
         Each item in the list is marked with a ``tick''.
         You don't have to worry about what kind of tick
         mark is used.

   \item This is the second item of the list.  It
         contains another list nested inside it.  The inner
         list is an \emph{enumerated} list.
         \begin{enumerate}
            \item This is the first item of an enumerated
                  list that is nested within the itemized list.

            \item This is the second item of the inner list.
                  \LaTeX\ allows you to nest lists deeper than
                  you really should.
         \end{enumerate}
         This is the rest of the second item of the outer
         list.  It is no more interesting than any other
         part of the item.
   \item This is the third item of the list.
\end{itemize}
You can even display poetry.
\begin{verse}
   There is an environment
    for verse \\             % The \\ command separates lines
   Whose features some poets % within a stanza.
   will curse.

                             % One or more blank lines separate stanzas.

   For instead of making\\
   Them do \emph{all} line breaking, \\
   It allows them to put too many words on a line when they'd rather be
   forced to be terse.
\end{verse}

Mathematical formulas may also be displayed.  A
displayed formula
is
one-line long; multiline
formulas require special formatting instructions.
   \[  \ip{\Gamma}{\psi'} = x'' + y^{2} + z_{i}^{n}\]
Don't start a paragraph with a displayed equation,
nor make one a paragraph by itself.

\end{document}               % End of document.

\section{Bench Huge Section 195}
% synthetic bench section 195
% This is a sample LaTeX input file.  (Version of 12 August 2004.)
%
% A '%' character causes TeX to ignore all remaining text on the line,
% and is used for comments like this one.

\documentclass{article}      % Specifies the document class

                             % The preamble begins here.
\title{An Example Document}  % Declares the document's title.
\author{Leslie Lamport}      % Declares the author's name.
\date{January 21, 1994}      % Deleting this command produces today's date.

\newcommand{\ip}[2]{(#1, #2)}
                             % Defines \ip{arg1}{arg2} to mean
                             % (arg1, arg2).

%\newcommand{\ip}[2]{\langle #1 | #2\rangle}
                             % This is an alternative definition of
                             % \ip that is commented out.

\begin{document}             % End of preamble and beginning of text.

\maketitle                   % Produces the title.

This is an example input file.  Comparing it with
the output it generates can show you how to
produce a simple document of your own.

\section{Ordinary Text}      % Produces section heading.  Lower-level
                             % sections are begun with similar
                             % \subsection and \subsubsection commands.

The ends  of words and sentences are marked
  by   spaces. It  doesn't matter how many
spaces    you type; one is as good as 100.  The
end of   a line counts as a space.

One   or more   blank lines denote the  end
of  a paragraph.

Since any number of consecutive spaces are treated
like a single one, the formatting of the input
file makes no difference to
      \LaTeX,                % The \LaTeX command generates the LaTeX logo.
but it makes a difference to you.  When you use
\LaTeX, making your input file as easy to read
as possible will be a great help as you write
your document and when you change it.  This sample
file shows how you can add comments to your own input
file.

Because printing is different from typewriting,
there are a number of things that you have to do
differently when preparing an input file than if
you were just typing the document directly.
Quotation marks like
       ``this''
have to be handled specially, as do quotes within
quotes:
       ``\,`this'            % \, separates the double and single quote.
        is what I just
        wrote, not  `that'\,''.

Dashes come in three sizes: an
       intra-word
dash, a medium dash for number ranges like
       1--2,
and a punctuation
       dash---like
this.

A sentence-ending space should be larger than the
space between words within a sentence.  You
sometimes have to type special commands in
conjunction with punctuation characters to get
this right, as in the following sentence.
       Gnats, gnus, etc.\ all  % `\ ' makes an inter-word space.
       begin with G\@.         % \@ marks end-of-sentence punctuation.
You should check the spaces after periods when
reading your output to make sure you haven't
forgotten any special cases.  Generating an
ellipsis
       \ldots\               % `\ ' is needed after `\ldots' because TeX
                             % ignores spaces after command names like \ldots
                             % made from \ + letters.
                             %
                             % Note how a `%' character causes TeX to ignore
                             % the end of the input line, so these blank lines
                             % do not start a new paragraph.
                             %
with the right spacing around the periods requires
a special command.

\LaTeX\ interprets some common characters as
commands, so you must type special commands to
generate them.  These characters include the
following:
       \$ \& \% \# \{ and \}.

In printing, text is usually emphasized with an
       \emph{italic}
type style.

\begin{em}
   A long segment of text can also be emphasized
   in this way.  Text within such a segment can be
   given \emph{additional} emphasis.
\end{em}

It is sometimes necessary to prevent \LaTeX\ from
breaking a line where it might otherwise do so.
This may be at a space, as between the ``Mr.''\ and
``Jones'' in
       ``Mr.~Jones'',        % ~ produces an unbreakable interword space.
or within a word---especially when the word is a
symbol like
       \mbox{\emph{itemnum}}
that makes little sense when hyphenated across
lines.

Footnotes\footnote{This is an example of a footnote.}
pose no problem.

\LaTeX\ is good at typesetting mathematical formulas
like
       \( x-3y + z = 7 \)
or
       \( a_{1} > x^{2n} + y^{2n} > x' \)
or
       \( \ip{A}{B} = \sum_{i} a_{i} b_{i} \).
The spaces you type in a formula are
ignored.  Remember that a letter like
       $x$                   % $ ... $  and  \( ... \)  are equivalent
is a formula when it denotes a mathematical
symbol, and it should be typed as one.

\section{Displayed Text}

Text is displayed by indenting it from the left
margin.  Quotations are commonly displayed.  There
are short quotations
\begin{quote}
   This is a short quotation.  It consists of a
   single paragraph of text.  See how it is formatted.
\end{quote}
and longer ones.
\begin{quotation}
   This is a longer quotation.  It consists of two
   paragraphs of text, neither of which are
   particularly interesting.

   This is the second paragraph of the quotation.  It
   is just as dull as the first paragraph.
\end{quotation}
Another frequently-displayed structure is a list.
The following is an example of an \emph{itemized}
list.
\begin{itemize}
   \item This is the first item of an itemized list.
         Each item in the list is marked with a ``tick''.
         You don't have to worry about what kind of tick
         mark is used.

   \item This is the second item of the list.  It
         contains another list nested inside it.  The inner
         list is an \emph{enumerated} list.
         \begin{enumerate}
            \item This is the first item of an enumerated
                  list that is nested within the itemized list.

            \item This is the second item of the inner list.
                  \LaTeX\ allows you to nest lists deeper than
                  you really should.
         \end{enumerate}
         This is the rest of the second item of the outer
         list.  It is no more interesting than any other
         part of the item.
   \item This is the third item of the list.
\end{itemize}
You can even display poetry.
\begin{verse}
   There is an environment
    for verse \\             % The \\ command separates lines
   Whose features some poets % within a stanza.
   will curse.

                             % One or more blank lines separate stanzas.

   For instead of making\\
   Them do \emph{all} line breaking, \\
   It allows them to put too many words on a line when they'd rather be
   forced to be terse.
\end{verse}

Mathematical formulas may also be displayed.  A
displayed formula
is
one-line long; multiline
formulas require special formatting instructions.
   \[  \ip{\Gamma}{\psi'} = x'' + y^{2} + z_{i}^{n}\]
Don't start a paragraph with a displayed equation,
nor make one a paragraph by itself.

\end{document}               % End of document.

\section{Bench Huge Section 196}
% synthetic bench section 196
% This is a sample LaTeX input file.  (Version of 12 August 2004.)
%
% A '%' character causes TeX to ignore all remaining text on the line,
% and is used for comments like this one.

\documentclass{article}      % Specifies the document class

                             % The preamble begins here.
\title{An Example Document}  % Declares the document's title.
\author{Leslie Lamport}      % Declares the author's name.
\date{January 21, 1994}      % Deleting this command produces today's date.

\newcommand{\ip}[2]{(#1, #2)}
                             % Defines \ip{arg1}{arg2} to mean
                             % (arg1, arg2).

%\newcommand{\ip}[2]{\langle #1 | #2\rangle}
                             % This is an alternative definition of
                             % \ip that is commented out.

\begin{document}             % End of preamble and beginning of text.

\maketitle                   % Produces the title.

This is an example input file.  Comparing it with
the output it generates can show you how to
produce a simple document of your own.

\section{Ordinary Text}      % Produces section heading.  Lower-level
                             % sections are begun with similar
                             % \subsection and \subsubsection commands.

The ends  of words and sentences are marked
  by   spaces. It  doesn't matter how many
spaces    you type; one is as good as 100.  The
end of   a line counts as a space.

One   or more   blank lines denote the  end
of  a paragraph.

Since any number of consecutive spaces are treated
like a single one, the formatting of the input
file makes no difference to
      \LaTeX,                % The \LaTeX command generates the LaTeX logo.
but it makes a difference to you.  When you use
\LaTeX, making your input file as easy to read
as possible will be a great help as you write
your document and when you change it.  This sample
file shows how you can add comments to your own input
file.

Because printing is different from typewriting,
there are a number of things that you have to do
differently when preparing an input file than if
you were just typing the document directly.
Quotation marks like
       ``this''
have to be handled specially, as do quotes within
quotes:
       ``\,`this'            % \, separates the double and single quote.
        is what I just
        wrote, not  `that'\,''.

Dashes come in three sizes: an
       intra-word
dash, a medium dash for number ranges like
       1--2,
and a punctuation
       dash---like
this.

A sentence-ending space should be larger than the
space between words within a sentence.  You
sometimes have to type special commands in
conjunction with punctuation characters to get
this right, as in the following sentence.
       Gnats, gnus, etc.\ all  % `\ ' makes an inter-word space.
       begin with G\@.         % \@ marks end-of-sentence punctuation.
You should check the spaces after periods when
reading your output to make sure you haven't
forgotten any special cases.  Generating an
ellipsis
       \ldots\               % `\ ' is needed after `\ldots' because TeX
                             % ignores spaces after command names like \ldots
                             % made from \ + letters.
                             %
                             % Note how a `%' character causes TeX to ignore
                             % the end of the input line, so these blank lines
                             % do not start a new paragraph.
                             %
with the right spacing around the periods requires
a special command.

\LaTeX\ interprets some common characters as
commands, so you must type special commands to
generate them.  These characters include the
following:
       \$ \& \% \# \{ and \}.

In printing, text is usually emphasized with an
       \emph{italic}
type style.

\begin{em}
   A long segment of text can also be emphasized
   in this way.  Text within such a segment can be
   given \emph{additional} emphasis.
\end{em}

It is sometimes necessary to prevent \LaTeX\ from
breaking a line where it might otherwise do so.
This may be at a space, as between the ``Mr.''\ and
``Jones'' in
       ``Mr.~Jones'',        % ~ produces an unbreakable interword space.
or within a word---especially when the word is a
symbol like
       \mbox{\emph{itemnum}}
that makes little sense when hyphenated across
lines.

Footnotes\footnote{This is an example of a footnote.}
pose no problem.

\LaTeX\ is good at typesetting mathematical formulas
like
       \( x-3y + z = 7 \)
or
       \( a_{1} > x^{2n} + y^{2n} > x' \)
or
       \( \ip{A}{B} = \sum_{i} a_{i} b_{i} \).
The spaces you type in a formula are
ignored.  Remember that a letter like
       $x$                   % $ ... $  and  \( ... \)  are equivalent
is a formula when it denotes a mathematical
symbol, and it should be typed as one.

\section{Displayed Text}

Text is displayed by indenting it from the left
margin.  Quotations are commonly displayed.  There
are short quotations
\begin{quote}
   This is a short quotation.  It consists of a
   single paragraph of text.  See how it is formatted.
\end{quote}
and longer ones.
\begin{quotation}
   This is a longer quotation.  It consists of two
   paragraphs of text, neither of which are
   particularly interesting.

   This is the second paragraph of the quotation.  It
   is just as dull as the first paragraph.
\end{quotation}
Another frequently-displayed structure is a list.
The following is an example of an \emph{itemized}
list.
\begin{itemize}
   \item This is the first item of an itemized list.
         Each item in the list is marked with a ``tick''.
         You don't have to worry about what kind of tick
         mark is used.

   \item This is the second item of the list.  It
         contains another list nested inside it.  The inner
         list is an \emph{enumerated} list.
         \begin{enumerate}
            \item This is the first item of an enumerated
                  list that is nested within the itemized list.

            \item This is the second item of the inner list.
                  \LaTeX\ allows you to nest lists deeper than
                  you really should.
         \end{enumerate}
         This is the rest of the second item of the outer
         list.  It is no more interesting than any other
         part of the item.
   \item This is the third item of the list.
\end{itemize}
You can even display poetry.
\begin{verse}
   There is an environment
    for verse \\             % The \\ command separates lines
   Whose features some poets % within a stanza.
   will curse.

                             % One or more blank lines separate stanzas.

   For instead of making\\
   Them do \emph{all} line breaking, \\
   It allows them to put too many words on a line when they'd rather be
   forced to be terse.
\end{verse}

Mathematical formulas may also be displayed.  A
displayed formula
is
one-line long; multiline
formulas require special formatting instructions.
   \[  \ip{\Gamma}{\psi'} = x'' + y^{2} + z_{i}^{n}\]
Don't start a paragraph with a displayed equation,
nor make one a paragraph by itself.

\end{document}               % End of document.

\section{Bench Huge Section 197}
% synthetic bench section 197
% This is a sample LaTeX input file.  (Version of 12 August 2004.)
%
% A '%' character causes TeX to ignore all remaining text on the line,
% and is used for comments like this one.

\documentclass{article}      % Specifies the document class

                             % The preamble begins here.
\title{An Example Document}  % Declares the document's title.
\author{Leslie Lamport}      % Declares the author's name.
\date{January 21, 1994}      % Deleting this command produces today's date.

\newcommand{\ip}[2]{(#1, #2)}
                             % Defines \ip{arg1}{arg2} to mean
                             % (arg1, arg2).

%\newcommand{\ip}[2]{\langle #1 | #2\rangle}
                             % This is an alternative definition of
                             % \ip that is commented out.

\begin{document}             % End of preamble and beginning of text.

\maketitle                   % Produces the title.

This is an example input file.  Comparing it with
the output it generates can show you how to
produce a simple document of your own.

\section{Ordinary Text}      % Produces section heading.  Lower-level
                             % sections are begun with similar
                             % \subsection and \subsubsection commands.

The ends  of words and sentences are marked
  by   spaces. It  doesn't matter how many
spaces    you type; one is as good as 100.  The
end of   a line counts as a space.

One   or more   blank lines denote the  end
of  a paragraph.

Since any number of consecutive spaces are treated
like a single one, the formatting of the input
file makes no difference to
      \LaTeX,                % The \LaTeX command generates the LaTeX logo.
but it makes a difference to you.  When you use
\LaTeX, making your input file as easy to read
as possible will be a great help as you write
your document and when you change it.  This sample
file shows how you can add comments to your own input
file.

Because printing is different from typewriting,
there are a number of things that you have to do
differently when preparing an input file than if
you were just typing the document directly.
Quotation marks like
       ``this''
have to be handled specially, as do quotes within
quotes:
       ``\,`this'            % \, separates the double and single quote.
        is what I just
        wrote, not  `that'\,''.

Dashes come in three sizes: an
       intra-word
dash, a medium dash for number ranges like
       1--2,
and a punctuation
       dash---like
this.

A sentence-ending space should be larger than the
space between words within a sentence.  You
sometimes have to type special commands in
conjunction with punctuation characters to get
this right, as in the following sentence.
       Gnats, gnus, etc.\ all  % `\ ' makes an inter-word space.
       begin with G\@.         % \@ marks end-of-sentence punctuation.
You should check the spaces after periods when
reading your output to make sure you haven't
forgotten any special cases.  Generating an
ellipsis
       \ldots\               % `\ ' is needed after `\ldots' because TeX
                             % ignores spaces after command names like \ldots
                             % made from \ + letters.
                             %
                             % Note how a `%' character causes TeX to ignore
                             % the end of the input line, so these blank lines
                             % do not start a new paragraph.
                             %
with the right spacing around the periods requires
a special command.

\LaTeX\ interprets some common characters as
commands, so you must type special commands to
generate them.  These characters include the
following:
       \$ \& \% \# \{ and \}.

In printing, text is usually emphasized with an
       \emph{italic}
type style.

\begin{em}
   A long segment of text can also be emphasized
   in this way.  Text within such a segment can be
   given \emph{additional} emphasis.
\end{em}

It is sometimes necessary to prevent \LaTeX\ from
breaking a line where it might otherwise do so.
This may be at a space, as between the ``Mr.''\ and
``Jones'' in
       ``Mr.~Jones'',        % ~ produces an unbreakable interword space.
or within a word---especially when the word is a
symbol like
       \mbox{\emph{itemnum}}
that makes little sense when hyphenated across
lines.

Footnotes\footnote{This is an example of a footnote.}
pose no problem.

\LaTeX\ is good at typesetting mathematical formulas
like
       \( x-3y + z = 7 \)
or
       \( a_{1} > x^{2n} + y^{2n} > x' \)
or
       \( \ip{A}{B} = \sum_{i} a_{i} b_{i} \).
The spaces you type in a formula are
ignored.  Remember that a letter like
       $x$                   % $ ... $  and  \( ... \)  are equivalent
is a formula when it denotes a mathematical
symbol, and it should be typed as one.

\section{Displayed Text}

Text is displayed by indenting it from the left
margin.  Quotations are commonly displayed.  There
are short quotations
\begin{quote}
   This is a short quotation.  It consists of a
   single paragraph of text.  See how it is formatted.
\end{quote}
and longer ones.
\begin{quotation}
   This is a longer quotation.  It consists of two
   paragraphs of text, neither of which are
   particularly interesting.

   This is the second paragraph of the quotation.  It
   is just as dull as the first paragraph.
\end{quotation}
Another frequently-displayed structure is a list.
The following is an example of an \emph{itemized}
list.
\begin{itemize}
   \item This is the first item of an itemized list.
         Each item in the list is marked with a ``tick''.
         You don't have to worry about what kind of tick
         mark is used.

   \item This is the second item of the list.  It
         contains another list nested inside it.  The inner
         list is an \emph{enumerated} list.
         \begin{enumerate}
            \item This is the first item of an enumerated
                  list that is nested within the itemized list.

            \item This is the second item of the inner list.
                  \LaTeX\ allows you to nest lists deeper than
                  you really should.
         \end{enumerate}
         This is the rest of the second item of the outer
         list.  It is no more interesting than any other
         part of the item.
   \item This is the third item of the list.
\end{itemize}
You can even display poetry.
\begin{verse}
   There is an environment
    for verse \\             % The \\ command separates lines
   Whose features some poets % within a stanza.
   will curse.

                             % One or more blank lines separate stanzas.

   For instead of making\\
   Them do \emph{all} line breaking, \\
   It allows them to put too many words on a line when they'd rather be
   forced to be terse.
\end{verse}

Mathematical formulas may also be displayed.  A
displayed formula
is
one-line long; multiline
formulas require special formatting instructions.
   \[  \ip{\Gamma}{\psi'} = x'' + y^{2} + z_{i}^{n}\]
Don't start a paragraph with a displayed equation,
nor make one a paragraph by itself.

\end{document}               % End of document.

\section{Bench Huge Section 198}
% synthetic bench section 198
% This is a sample LaTeX input file.  (Version of 12 August 2004.)
%
% A '%' character causes TeX to ignore all remaining text on the line,
% and is used for comments like this one.

\documentclass{article}      % Specifies the document class

                             % The preamble begins here.
\title{An Example Document}  % Declares the document's title.
\author{Leslie Lamport}      % Declares the author's name.
\date{January 21, 1994}      % Deleting this command produces today's date.

\newcommand{\ip}[2]{(#1, #2)}
                             % Defines \ip{arg1}{arg2} to mean
                             % (arg1, arg2).

%\newcommand{\ip}[2]{\langle #1 | #2\rangle}
                             % This is an alternative definition of
                             % \ip that is commented out.

\begin{document}             % End of preamble and beginning of text.

\maketitle                   % Produces the title.

This is an example input file.  Comparing it with
the output it generates can show you how to
produce a simple document of your own.

\section{Ordinary Text}      % Produces section heading.  Lower-level
                             % sections are begun with similar
                             % \subsection and \subsubsection commands.

The ends  of words and sentences are marked
  by   spaces. It  doesn't matter how many
spaces    you type; one is as good as 100.  The
end of   a line counts as a space.

One   or more   blank lines denote the  end
of  a paragraph.

Since any number of consecutive spaces are treated
like a single one, the formatting of the input
file makes no difference to
      \LaTeX,                % The \LaTeX command generates the LaTeX logo.
but it makes a difference to you.  When you use
\LaTeX, making your input file as easy to read
as possible will be a great help as you write
your document and when you change it.  This sample
file shows how you can add comments to your own input
file.

Because printing is different from typewriting,
there are a number of things that you have to do
differently when preparing an input file than if
you were just typing the document directly.
Quotation marks like
       ``this''
have to be handled specially, as do quotes within
quotes:
       ``\,`this'            % \, separates the double and single quote.
        is what I just
        wrote, not  `that'\,''.

Dashes come in three sizes: an
       intra-word
dash, a medium dash for number ranges like
       1--2,
and a punctuation
       dash---like
this.

A sentence-ending space should be larger than the
space between words within a sentence.  You
sometimes have to type special commands in
conjunction with punctuation characters to get
this right, as in the following sentence.
       Gnats, gnus, etc.\ all  % `\ ' makes an inter-word space.
       begin with G\@.         % \@ marks end-of-sentence punctuation.
You should check the spaces after periods when
reading your output to make sure you haven't
forgotten any special cases.  Generating an
ellipsis
       \ldots\               % `\ ' is needed after `\ldots' because TeX
                             % ignores spaces after command names like \ldots
                             % made from \ + letters.
                             %
                             % Note how a `%' character causes TeX to ignore
                             % the end of the input line, so these blank lines
                             % do not start a new paragraph.
                             %
with the right spacing around the periods requires
a special command.

\LaTeX\ interprets some common characters as
commands, so you must type special commands to
generate them.  These characters include the
following:
       \$ \& \% \# \{ and \}.

In printing, text is usually emphasized with an
       \emph{italic}
type style.

\begin{em}
   A long segment of text can also be emphasized
   in this way.  Text within such a segment can be
   given \emph{additional} emphasis.
\end{em}

It is sometimes necessary to prevent \LaTeX\ from
breaking a line where it might otherwise do so.
This may be at a space, as between the ``Mr.''\ and
``Jones'' in
       ``Mr.~Jones'',        % ~ produces an unbreakable interword space.
or within a word---especially when the word is a
symbol like
       \mbox{\emph{itemnum}}
that makes little sense when hyphenated across
lines.

Footnotes\footnote{This is an example of a footnote.}
pose no problem.

\LaTeX\ is good at typesetting mathematical formulas
like
       \( x-3y + z = 7 \)
or
       \( a_{1} > x^{2n} + y^{2n} > x' \)
or
       \( \ip{A}{B} = \sum_{i} a_{i} b_{i} \).
The spaces you type in a formula are
ignored.  Remember that a letter like
       $x$                   % $ ... $  and  \( ... \)  are equivalent
is a formula when it denotes a mathematical
symbol, and it should be typed as one.

\section{Displayed Text}

Text is displayed by indenting it from the left
margin.  Quotations are commonly displayed.  There
are short quotations
\begin{quote}
   This is a short quotation.  It consists of a
   single paragraph of text.  See how it is formatted.
\end{quote}
and longer ones.
\begin{quotation}
   This is a longer quotation.  It consists of two
   paragraphs of text, neither of which are
   particularly interesting.

   This is the second paragraph of the quotation.  It
   is just as dull as the first paragraph.
\end{quotation}
Another frequently-displayed structure is a list.
The following is an example of an \emph{itemized}
list.
\begin{itemize}
   \item This is the first item of an itemized list.
         Each item in the list is marked with a ``tick''.
         You don't have to worry about what kind of tick
         mark is used.

   \item This is the second item of the list.  It
         contains another list nested inside it.  The inner
         list is an \emph{enumerated} list.
         \begin{enumerate}
            \item This is the first item of an enumerated
                  list that is nested within the itemized list.

            \item This is the second item of the inner list.
                  \LaTeX\ allows you to nest lists deeper than
                  you really should.
         \end{enumerate}
         This is the rest of the second item of the outer
         list.  It is no more interesting than any other
         part of the item.
   \item This is the third item of the list.
\end{itemize}
You can even display poetry.
\begin{verse}
   There is an environment
    for verse \\             % The \\ command separates lines
   Whose features some poets % within a stanza.
   will curse.

                             % One or more blank lines separate stanzas.

   For instead of making\\
   Them do \emph{all} line breaking, \\
   It allows them to put too many words on a line when they'd rather be
   forced to be terse.
\end{verse}

Mathematical formulas may also be displayed.  A
displayed formula
is
one-line long; multiline
formulas require special formatting instructions.
   \[  \ip{\Gamma}{\psi'} = x'' + y^{2} + z_{i}^{n}\]
Don't start a paragraph with a displayed equation,
nor make one a paragraph by itself.

\end{document}               % End of document.

\section{Bench Huge Section 199}
% synthetic bench section 199
% This is a sample LaTeX input file.  (Version of 12 August 2004.)
%
% A '%' character causes TeX to ignore all remaining text on the line,
% and is used for comments like this one.

\documentclass{article}      % Specifies the document class

                             % The preamble begins here.
\title{An Example Document}  % Declares the document's title.
\author{Leslie Lamport}      % Declares the author's name.
\date{January 21, 1994}      % Deleting this command produces today's date.

\newcommand{\ip}[2]{(#1, #2)}
                             % Defines \ip{arg1}{arg2} to mean
                             % (arg1, arg2).

%\newcommand{\ip}[2]{\langle #1 | #2\rangle}
                             % This is an alternative definition of
                             % \ip that is commented out.

\begin{document}             % End of preamble and beginning of text.

\maketitle                   % Produces the title.

This is an example input file.  Comparing it with
the output it generates can show you how to
produce a simple document of your own.

\section{Ordinary Text}      % Produces section heading.  Lower-level
                             % sections are begun with similar
                             % \subsection and \subsubsection commands.

The ends  of words and sentences are marked
  by   spaces. It  doesn't matter how many
spaces    you type; one is as good as 100.  The
end of   a line counts as a space.

One   or more   blank lines denote the  end
of  a paragraph.

Since any number of consecutive spaces are treated
like a single one, the formatting of the input
file makes no difference to
      \LaTeX,                % The \LaTeX command generates the LaTeX logo.
but it makes a difference to you.  When you use
\LaTeX, making your input file as easy to read
as possible will be a great help as you write
your document and when you change it.  This sample
file shows how you can add comments to your own input
file.

Because printing is different from typewriting,
there are a number of things that you have to do
differently when preparing an input file than if
you were just typing the document directly.
Quotation marks like
       ``this''
have to be handled specially, as do quotes within
quotes:
       ``\,`this'            % \, separates the double and single quote.
        is what I just
        wrote, not  `that'\,''.

Dashes come in three sizes: an
       intra-word
dash, a medium dash for number ranges like
       1--2,
and a punctuation
       dash---like
this.

A sentence-ending space should be larger than the
space between words within a sentence.  You
sometimes have to type special commands in
conjunction with punctuation characters to get
this right, as in the following sentence.
       Gnats, gnus, etc.\ all  % `\ ' makes an inter-word space.
       begin with G\@.         % \@ marks end-of-sentence punctuation.
You should check the spaces after periods when
reading your output to make sure you haven't
forgotten any special cases.  Generating an
ellipsis
       \ldots\               % `\ ' is needed after `\ldots' because TeX
                             % ignores spaces after command names like \ldots
                             % made from \ + letters.
                             %
                             % Note how a `%' character causes TeX to ignore
                             % the end of the input line, so these blank lines
                             % do not start a new paragraph.
                             %
with the right spacing around the periods requires
a special command.

\LaTeX\ interprets some common characters as
commands, so you must type special commands to
generate them.  These characters include the
following:
       \$ \& \% \# \{ and \}.

In printing, text is usually emphasized with an
       \emph{italic}
type style.

\begin{em}
   A long segment of text can also be emphasized
   in this way.  Text within such a segment can be
   given \emph{additional} emphasis.
\end{em}

It is sometimes necessary to prevent \LaTeX\ from
breaking a line where it might otherwise do so.
This may be at a space, as between the ``Mr.''\ and
``Jones'' in
       ``Mr.~Jones'',        % ~ produces an unbreakable interword space.
or within a word---especially when the word is a
symbol like
       \mbox{\emph{itemnum}}
that makes little sense when hyphenated across
lines.

Footnotes\footnote{This is an example of a footnote.}
pose no problem.

\LaTeX\ is good at typesetting mathematical formulas
like
       \( x-3y + z = 7 \)
or
       \( a_{1} > x^{2n} + y^{2n} > x' \)
or
       \( \ip{A}{B} = \sum_{i} a_{i} b_{i} \).
The spaces you type in a formula are
ignored.  Remember that a letter like
       $x$                   % $ ... $  and  \( ... \)  are equivalent
is a formula when it denotes a mathematical
symbol, and it should be typed as one.

\section{Displayed Text}

Text is displayed by indenting it from the left
margin.  Quotations are commonly displayed.  There
are short quotations
\begin{quote}
   This is a short quotation.  It consists of a
   single paragraph of text.  See how it is formatted.
\end{quote}
and longer ones.
\begin{quotation}
   This is a longer quotation.  It consists of two
   paragraphs of text, neither of which are
   particularly interesting.

   This is the second paragraph of the quotation.  It
   is just as dull as the first paragraph.
\end{quotation}
Another frequently-displayed structure is a list.
The following is an example of an \emph{itemized}
list.
\begin{itemize}
   \item This is the first item of an itemized list.
         Each item in the list is marked with a ``tick''.
         You don't have to worry about what kind of tick
         mark is used.

   \item This is the second item of the list.  It
         contains another list nested inside it.  The inner
         list is an \emph{enumerated} list.
         \begin{enumerate}
            \item This is the first item of an enumerated
                  list that is nested within the itemized list.

            \item This is the second item of the inner list.
                  \LaTeX\ allows you to nest lists deeper than
                  you really should.
         \end{enumerate}
         This is the rest of the second item of the outer
         list.  It is no more interesting than any other
         part of the item.
   \item This is the third item of the list.
\end{itemize}
You can even display poetry.
\begin{verse}
   There is an environment
    for verse \\             % The \\ command separates lines
   Whose features some poets % within a stanza.
   will curse.

                             % One or more blank lines separate stanzas.

   For instead of making\\
   Them do \emph{all} line breaking, \\
   It allows them to put too many words on a line when they'd rather be
   forced to be terse.
\end{verse}

Mathematical formulas may also be displayed.  A
displayed formula
is
one-line long; multiline
formulas require special formatting instructions.
   \[  \ip{\Gamma}{\psi'} = x'' + y^{2} + z_{i}^{n}\]
Don't start a paragraph with a displayed equation,
nor make one a paragraph by itself.

\end{document}               % End of document.

\section{Bench Huge Section 200}
% synthetic bench section 200
% This is a sample LaTeX input file.  (Version of 12 August 2004.)
%
% A '%' character causes TeX to ignore all remaining text on the line,
% and is used for comments like this one.

\documentclass{article}      % Specifies the document class

                             % The preamble begins here.
\title{An Example Document}  % Declares the document's title.
\author{Leslie Lamport}      % Declares the author's name.
\date{January 21, 1994}      % Deleting this command produces today's date.

\newcommand{\ip}[2]{(#1, #2)}
                             % Defines \ip{arg1}{arg2} to mean
                             % (arg1, arg2).

%\newcommand{\ip}[2]{\langle #1 | #2\rangle}
                             % This is an alternative definition of
                             % \ip that is commented out.

\begin{document}             % End of preamble and beginning of text.

\maketitle                   % Produces the title.

This is an example input file.  Comparing it with
the output it generates can show you how to
produce a simple document of your own.

\section{Ordinary Text}      % Produces section heading.  Lower-level
                             % sections are begun with similar
                             % \subsection and \subsubsection commands.

The ends  of words and sentences are marked
  by   spaces. It  doesn't matter how many
spaces    you type; one is as good as 100.  The
end of   a line counts as a space.

One   or more   blank lines denote the  end
of  a paragraph.

Since any number of consecutive spaces are treated
like a single one, the formatting of the input
file makes no difference to
      \LaTeX,                % The \LaTeX command generates the LaTeX logo.
but it makes a difference to you.  When you use
\LaTeX, making your input file as easy to read
as possible will be a great help as you write
your document and when you change it.  This sample
file shows how you can add comments to your own input
file.

Because printing is different from typewriting,
there are a number of things that you have to do
differently when preparing an input file than if
you were just typing the document directly.
Quotation marks like
       ``this''
have to be handled specially, as do quotes within
quotes:
       ``\,`this'            % \, separates the double and single quote.
        is what I just
        wrote, not  `that'\,''.

Dashes come in three sizes: an
       intra-word
dash, a medium dash for number ranges like
       1--2,
and a punctuation
       dash---like
this.

A sentence-ending space should be larger than the
space between words within a sentence.  You
sometimes have to type special commands in
conjunction with punctuation characters to get
this right, as in the following sentence.
       Gnats, gnus, etc.\ all  % `\ ' makes an inter-word space.
       begin with G\@.         % \@ marks end-of-sentence punctuation.
You should check the spaces after periods when
reading your output to make sure you haven't
forgotten any special cases.  Generating an
ellipsis
       \ldots\               % `\ ' is needed after `\ldots' because TeX
                             % ignores spaces after command names like \ldots
                             % made from \ + letters.
                             %
                             % Note how a `%' character causes TeX to ignore
                             % the end of the input line, so these blank lines
                             % do not start a new paragraph.
                             %
with the right spacing around the periods requires
a special command.

\LaTeX\ interprets some common characters as
commands, so you must type special commands to
generate them.  These characters include the
following:
       \$ \& \% \# \{ and \}.

In printing, text is usually emphasized with an
       \emph{italic}
type style.

\begin{em}
   A long segment of text can also be emphasized
   in this way.  Text within such a segment can be
   given \emph{additional} emphasis.
\end{em}

It is sometimes necessary to prevent \LaTeX\ from
breaking a line where it might otherwise do so.
This may be at a space, as between the ``Mr.''\ and
``Jones'' in
       ``Mr.~Jones'',        % ~ produces an unbreakable interword space.
or within a word---especially when the word is a
symbol like
       \mbox{\emph{itemnum}}
that makes little sense when hyphenated across
lines.

Footnotes\footnote{This is an example of a footnote.}
pose no problem.

\LaTeX\ is good at typesetting mathematical formulas
like
       \( x-3y + z = 7 \)
or
       \( a_{1} > x^{2n} + y^{2n} > x' \)
or
       \( \ip{A}{B} = \sum_{i} a_{i} b_{i} \).
The spaces you type in a formula are
ignored.  Remember that a letter like
       $x$                   % $ ... $  and  \( ... \)  are equivalent
is a formula when it denotes a mathematical
symbol, and it should be typed as one.

\section{Displayed Text}

Text is displayed by indenting it from the left
margin.  Quotations are commonly displayed.  There
are short quotations
\begin{quote}
   This is a short quotation.  It consists of a
   single paragraph of text.  See how it is formatted.
\end{quote}
and longer ones.
\begin{quotation}
   This is a longer quotation.  It consists of two
   paragraphs of text, neither of which are
   particularly interesting.

   This is the second paragraph of the quotation.  It
   is just as dull as the first paragraph.
\end{quotation}
Another frequently-displayed structure is a list.
The following is an example of an \emph{itemized}
list.
\begin{itemize}
   \item This is the first item of an itemized list.
         Each item in the list is marked with a ``tick''.
         You don't have to worry about what kind of tick
         mark is used.

   \item This is the second item of the list.  It
         contains another list nested inside it.  The inner
         list is an \emph{enumerated} list.
         \begin{enumerate}
            \item This is the first item of an enumerated
                  list that is nested within the itemized list.

            \item This is the second item of the inner list.
                  \LaTeX\ allows you to nest lists deeper than
                  you really should.
         \end{enumerate}
         This is the rest of the second item of the outer
         list.  It is no more interesting than any other
         part of the item.
   \item This is the third item of the list.
\end{itemize}
You can even display poetry.
\begin{verse}
   There is an environment
    for verse \\             % The \\ command separates lines
   Whose features some poets % within a stanza.
   will curse.

                             % One or more blank lines separate stanzas.

   For instead of making\\
   Them do \emph{all} line breaking, \\
   It allows them to put too many words on a line when they'd rather be
   forced to be terse.
\end{verse}

Mathematical formulas may also be displayed.  A
displayed formula
is
one-line long; multiline
formulas require special formatting instructions.
   \[  \ip{\Gamma}{\psi'} = x'' + y^{2} + z_{i}^{n}\]
Don't start a paragraph with a displayed equation,
nor make one a paragraph by itself.

\end{document}               % End of document.

\section{Bench Huge Section 201}
% synthetic bench section 201
% This is a sample LaTeX input file.  (Version of 12 August 2004.)
%
% A '%' character causes TeX to ignore all remaining text on the line,
% and is used for comments like this one.

\documentclass{article}      % Specifies the document class

                             % The preamble begins here.
\title{An Example Document}  % Declares the document's title.
\author{Leslie Lamport}      % Declares the author's name.
\date{January 21, 1994}      % Deleting this command produces today's date.

\newcommand{\ip}[2]{(#1, #2)}
                             % Defines \ip{arg1}{arg2} to mean
                             % (arg1, arg2).

%\newcommand{\ip}[2]{\langle #1 | #2\rangle}
                             % This is an alternative definition of
                             % \ip that is commented out.

\begin{document}             % End of preamble and beginning of text.

\maketitle                   % Produces the title.

This is an example input file.  Comparing it with
the output it generates can show you how to
produce a simple document of your own.

\section{Ordinary Text}      % Produces section heading.  Lower-level
                             % sections are begun with similar
                             % \subsection and \subsubsection commands.

The ends  of words and sentences are marked
  by   spaces. It  doesn't matter how many
spaces    you type; one is as good as 100.  The
end of   a line counts as a space.

One   or more   blank lines denote the  end
of  a paragraph.

Since any number of consecutive spaces are treated
like a single one, the formatting of the input
file makes no difference to
      \LaTeX,                % The \LaTeX command generates the LaTeX logo.
but it makes a difference to you.  When you use
\LaTeX, making your input file as easy to read
as possible will be a great help as you write
your document and when you change it.  This sample
file shows how you can add comments to your own input
file.

Because printing is different from typewriting,
there are a number of things that you have to do
differently when preparing an input file than if
you were just typing the document directly.
Quotation marks like
       ``this''
have to be handled specially, as do quotes within
quotes:
       ``\,`this'            % \, separates the double and single quote.
        is what I just
        wrote, not  `that'\,''.

Dashes come in three sizes: an
       intra-word
dash, a medium dash for number ranges like
       1--2,
and a punctuation
       dash---like
this.

A sentence-ending space should be larger than the
space between words within a sentence.  You
sometimes have to type special commands in
conjunction with punctuation characters to get
this right, as in the following sentence.
       Gnats, gnus, etc.\ all  % `\ ' makes an inter-word space.
       begin with G\@.         % \@ marks end-of-sentence punctuation.
You should check the spaces after periods when
reading your output to make sure you haven't
forgotten any special cases.  Generating an
ellipsis
       \ldots\               % `\ ' is needed after `\ldots' because TeX
                             % ignores spaces after command names like \ldots
                             % made from \ + letters.
                             %
                             % Note how a `%' character causes TeX to ignore
                             % the end of the input line, so these blank lines
                             % do not start a new paragraph.
                             %
with the right spacing around the periods requires
a special command.

\LaTeX\ interprets some common characters as
commands, so you must type special commands to
generate them.  These characters include the
following:
       \$ \& \% \# \{ and \}.

In printing, text is usually emphasized with an
       \emph{italic}
type style.

\begin{em}
   A long segment of text can also be emphasized
   in this way.  Text within such a segment can be
   given \emph{additional} emphasis.
\end{em}

It is sometimes necessary to prevent \LaTeX\ from
breaking a line where it might otherwise do so.
This may be at a space, as between the ``Mr.''\ and
``Jones'' in
       ``Mr.~Jones'',        % ~ produces an unbreakable interword space.
or within a word---especially when the word is a
symbol like
       \mbox{\emph{itemnum}}
that makes little sense when hyphenated across
lines.

Footnotes\footnote{This is an example of a footnote.}
pose no problem.

\LaTeX\ is good at typesetting mathematical formulas
like
       \( x-3y + z = 7 \)
or
       \( a_{1} > x^{2n} + y^{2n} > x' \)
or
       \( \ip{A}{B} = \sum_{i} a_{i} b_{i} \).
The spaces you type in a formula are
ignored.  Remember that a letter like
       $x$                   % $ ... $  and  \( ... \)  are equivalent
is a formula when it denotes a mathematical
symbol, and it should be typed as one.

\section{Displayed Text}

Text is displayed by indenting it from the left
margin.  Quotations are commonly displayed.  There
are short quotations
\begin{quote}
   This is a short quotation.  It consists of a
   single paragraph of text.  See how it is formatted.
\end{quote}
and longer ones.
\begin{quotation}
   This is a longer quotation.  It consists of two
   paragraphs of text, neither of which are
   particularly interesting.

   This is the second paragraph of the quotation.  It
   is just as dull as the first paragraph.
\end{quotation}
Another frequently-displayed structure is a list.
The following is an example of an \emph{itemized}
list.
\begin{itemize}
   \item This is the first item of an itemized list.
         Each item in the list is marked with a ``tick''.
         You don't have to worry about what kind of tick
         mark is used.

   \item This is the second item of the list.  It
         contains another list nested inside it.  The inner
         list is an \emph{enumerated} list.
         \begin{enumerate}
            \item This is the first item of an enumerated
                  list that is nested within the itemized list.

            \item This is the second item of the inner list.
                  \LaTeX\ allows you to nest lists deeper than
                  you really should.
         \end{enumerate}
         This is the rest of the second item of the outer
         list.  It is no more interesting than any other
         part of the item.
   \item This is the third item of the list.
\end{itemize}
You can even display poetry.
\begin{verse}
   There is an environment
    for verse \\             % The \\ command separates lines
   Whose features some poets % within a stanza.
   will curse.

                             % One or more blank lines separate stanzas.

   For instead of making\\
   Them do \emph{all} line breaking, \\
   It allows them to put too many words on a line when they'd rather be
   forced to be terse.
\end{verse}

Mathematical formulas may also be displayed.  A
displayed formula
is
one-line long; multiline
formulas require special formatting instructions.
   \[  \ip{\Gamma}{\psi'} = x'' + y^{2} + z_{i}^{n}\]
Don't start a paragraph with a displayed equation,
nor make one a paragraph by itself.

\end{document}               % End of document.

\section{Bench Huge Section 202}
% synthetic bench section 202
% This is a sample LaTeX input file.  (Version of 12 August 2004.)
%
% A '%' character causes TeX to ignore all remaining text on the line,
% and is used for comments like this one.

\documentclass{article}      % Specifies the document class

                             % The preamble begins here.
\title{An Example Document}  % Declares the document's title.
\author{Leslie Lamport}      % Declares the author's name.
\date{January 21, 1994}      % Deleting this command produces today's date.

\newcommand{\ip}[2]{(#1, #2)}
                             % Defines \ip{arg1}{arg2} to mean
                             % (arg1, arg2).

%\newcommand{\ip}[2]{\langle #1 | #2\rangle}
                             % This is an alternative definition of
                             % \ip that is commented out.

\begin{document}             % End of preamble and beginning of text.

\maketitle                   % Produces the title.

This is an example input file.  Comparing it with
the output it generates can show you how to
produce a simple document of your own.

\section{Ordinary Text}      % Produces section heading.  Lower-level
                             % sections are begun with similar
                             % \subsection and \subsubsection commands.

The ends  of words and sentences are marked
  by   spaces. It  doesn't matter how many
spaces    you type; one is as good as 100.  The
end of   a line counts as a space.

One   or more   blank lines denote the  end
of  a paragraph.

Since any number of consecutive spaces are treated
like a single one, the formatting of the input
file makes no difference to
      \LaTeX,                % The \LaTeX command generates the LaTeX logo.
but it makes a difference to you.  When you use
\LaTeX, making your input file as easy to read
as possible will be a great help as you write
your document and when you change it.  This sample
file shows how you can add comments to your own input
file.

Because printing is different from typewriting,
there are a number of things that you have to do
differently when preparing an input file than if
you were just typing the document directly.
Quotation marks like
       ``this''
have to be handled specially, as do quotes within
quotes:
       ``\,`this'            % \, separates the double and single quote.
        is what I just
        wrote, not  `that'\,''.

Dashes come in three sizes: an
       intra-word
dash, a medium dash for number ranges like
       1--2,
and a punctuation
       dash---like
this.

A sentence-ending space should be larger than the
space between words within a sentence.  You
sometimes have to type special commands in
conjunction with punctuation characters to get
this right, as in the following sentence.
       Gnats, gnus, etc.\ all  % `\ ' makes an inter-word space.
       begin with G\@.         % \@ marks end-of-sentence punctuation.
You should check the spaces after periods when
reading your output to make sure you haven't
forgotten any special cases.  Generating an
ellipsis
       \ldots\               % `\ ' is needed after `\ldots' because TeX
                             % ignores spaces after command names like \ldots
                             % made from \ + letters.
                             %
                             % Note how a `%' character causes TeX to ignore
                             % the end of the input line, so these blank lines
                             % do not start a new paragraph.
                             %
with the right spacing around the periods requires
a special command.

\LaTeX\ interprets some common characters as
commands, so you must type special commands to
generate them.  These characters include the
following:
       \$ \& \% \# \{ and \}.

In printing, text is usually emphasized with an
       \emph{italic}
type style.

\begin{em}
   A long segment of text can also be emphasized
   in this way.  Text within such a segment can be
   given \emph{additional} emphasis.
\end{em}

It is sometimes necessary to prevent \LaTeX\ from
breaking a line where it might otherwise do so.
This may be at a space, as between the ``Mr.''\ and
``Jones'' in
       ``Mr.~Jones'',        % ~ produces an unbreakable interword space.
or within a word---especially when the word is a
symbol like
       \mbox{\emph{itemnum}}
that makes little sense when hyphenated across
lines.

Footnotes\footnote{This is an example of a footnote.}
pose no problem.

\LaTeX\ is good at typesetting mathematical formulas
like
       \( x-3y + z = 7 \)
or
       \( a_{1} > x^{2n} + y^{2n} > x' \)
or
       \( \ip{A}{B} = \sum_{i} a_{i} b_{i} \).
The spaces you type in a formula are
ignored.  Remember that a letter like
       $x$                   % $ ... $  and  \( ... \)  are equivalent
is a formula when it denotes a mathematical
symbol, and it should be typed as one.

\section{Displayed Text}

Text is displayed by indenting it from the left
margin.  Quotations are commonly displayed.  There
are short quotations
\begin{quote}
   This is a short quotation.  It consists of a
   single paragraph of text.  See how it is formatted.
\end{quote}
and longer ones.
\begin{quotation}
   This is a longer quotation.  It consists of two
   paragraphs of text, neither of which are
   particularly interesting.

   This is the second paragraph of the quotation.  It
   is just as dull as the first paragraph.
\end{quotation}
Another frequently-displayed structure is a list.
The following is an example of an \emph{itemized}
list.
\begin{itemize}
   \item This is the first item of an itemized list.
         Each item in the list is marked with a ``tick''.
         You don't have to worry about what kind of tick
         mark is used.

   \item This is the second item of the list.  It
         contains another list nested inside it.  The inner
         list is an \emph{enumerated} list.
         \begin{enumerate}
            \item This is the first item of an enumerated
                  list that is nested within the itemized list.

            \item This is the second item of the inner list.
                  \LaTeX\ allows you to nest lists deeper than
                  you really should.
         \end{enumerate}
         This is the rest of the second item of the outer
         list.  It is no more interesting than any other
         part of the item.
   \item This is the third item of the list.
\end{itemize}
You can even display poetry.
\begin{verse}
   There is an environment
    for verse \\             % The \\ command separates lines
   Whose features some poets % within a stanza.
   will curse.

                             % One or more blank lines separate stanzas.

   For instead of making\\
   Them do \emph{all} line breaking, \\
   It allows them to put too many words on a line when they'd rather be
   forced to be terse.
\end{verse}

Mathematical formulas may also be displayed.  A
displayed formula
is
one-line long; multiline
formulas require special formatting instructions.
   \[  \ip{\Gamma}{\psi'} = x'' + y^{2} + z_{i}^{n}\]
Don't start a paragraph with a displayed equation,
nor make one a paragraph by itself.

\end{document}               % End of document.

\section{Bench Huge Section 203}
% synthetic bench section 203
% This is a sample LaTeX input file.  (Version of 12 August 2004.)
%
% A '%' character causes TeX to ignore all remaining text on the line,
% and is used for comments like this one.

\documentclass{article}      % Specifies the document class

                             % The preamble begins here.
\title{An Example Document}  % Declares the document's title.
\author{Leslie Lamport}      % Declares the author's name.
\date{January 21, 1994}      % Deleting this command produces today's date.

\newcommand{\ip}[2]{(#1, #2)}
                             % Defines \ip{arg1}{arg2} to mean
                             % (arg1, arg2).

%\newcommand{\ip}[2]{\langle #1 | #2\rangle}
                             % This is an alternative definition of
                             % \ip that is commented out.

\begin{document}             % End of preamble and beginning of text.

\maketitle                   % Produces the title.

This is an example input file.  Comparing it with
the output it generates can show you how to
produce a simple document of your own.

\section{Ordinary Text}      % Produces section heading.  Lower-level
                             % sections are begun with similar
                             % \subsection and \subsubsection commands.

The ends  of words and sentences are marked
  by   spaces. It  doesn't matter how many
spaces    you type; one is as good as 100.  The
end of   a line counts as a space.

One   or more   blank lines denote the  end
of  a paragraph.

Since any number of consecutive spaces are treated
like a single one, the formatting of the input
file makes no difference to
      \LaTeX,                % The \LaTeX command generates the LaTeX logo.
but it makes a difference to you.  When you use
\LaTeX, making your input file as easy to read
as possible will be a great help as you write
your document and when you change it.  This sample
file shows how you can add comments to your own input
file.

Because printing is different from typewriting,
there are a number of things that you have to do
differently when preparing an input file than if
you were just typing the document directly.
Quotation marks like
       ``this''
have to be handled specially, as do quotes within
quotes:
       ``\,`this'            % \, separates the double and single quote.
        is what I just
        wrote, not  `that'\,''.

Dashes come in three sizes: an
       intra-word
dash, a medium dash for number ranges like
       1--2,
and a punctuation
       dash---like
this.

A sentence-ending space should be larger than the
space between words within a sentence.  You
sometimes have to type special commands in
conjunction with punctuation characters to get
this right, as in the following sentence.
       Gnats, gnus, etc.\ all  % `\ ' makes an inter-word space.
       begin with G\@.         % \@ marks end-of-sentence punctuation.
You should check the spaces after periods when
reading your output to make sure you haven't
forgotten any special cases.  Generating an
ellipsis
       \ldots\               % `\ ' is needed after `\ldots' because TeX
                             % ignores spaces after command names like \ldots
                             % made from \ + letters.
                             %
                             % Note how a `%' character causes TeX to ignore
                             % the end of the input line, so these blank lines
                             % do not start a new paragraph.
                             %
with the right spacing around the periods requires
a special command.

\LaTeX\ interprets some common characters as
commands, so you must type special commands to
generate them.  These characters include the
following:
       \$ \& \% \# \{ and \}.

In printing, text is usually emphasized with an
       \emph{italic}
type style.

\begin{em}
   A long segment of text can also be emphasized
   in this way.  Text within such a segment can be
   given \emph{additional} emphasis.
\end{em}

It is sometimes necessary to prevent \LaTeX\ from
breaking a line where it might otherwise do so.
This may be at a space, as between the ``Mr.''\ and
``Jones'' in
       ``Mr.~Jones'',        % ~ produces an unbreakable interword space.
or within a word---especially when the word is a
symbol like
       \mbox{\emph{itemnum}}
that makes little sense when hyphenated across
lines.

Footnotes\footnote{This is an example of a footnote.}
pose no problem.

\LaTeX\ is good at typesetting mathematical formulas
like
       \( x-3y + z = 7 \)
or
       \( a_{1} > x^{2n} + y^{2n} > x' \)
or
       \( \ip{A}{B} = \sum_{i} a_{i} b_{i} \).
The spaces you type in a formula are
ignored.  Remember that a letter like
       $x$                   % $ ... $  and  \( ... \)  are equivalent
is a formula when it denotes a mathematical
symbol, and it should be typed as one.

\section{Displayed Text}

Text is displayed by indenting it from the left
margin.  Quotations are commonly displayed.  There
are short quotations
\begin{quote}
   This is a short quotation.  It consists of a
   single paragraph of text.  See how it is formatted.
\end{quote}
and longer ones.
\begin{quotation}
   This is a longer quotation.  It consists of two
   paragraphs of text, neither of which are
   particularly interesting.

   This is the second paragraph of the quotation.  It
   is just as dull as the first paragraph.
\end{quotation}
Another frequently-displayed structure is a list.
The following is an example of an \emph{itemized}
list.
\begin{itemize}
   \item This is the first item of an itemized list.
         Each item in the list is marked with a ``tick''.
         You don't have to worry about what kind of tick
         mark is used.

   \item This is the second item of the list.  It
         contains another list nested inside it.  The inner
         list is an \emph{enumerated} list.
         \begin{enumerate}
            \item This is the first item of an enumerated
                  list that is nested within the itemized list.

            \item This is the second item of the inner list.
                  \LaTeX\ allows you to nest lists deeper than
                  you really should.
         \end{enumerate}
         This is the rest of the second item of the outer
         list.  It is no more interesting than any other
         part of the item.
   \item This is the third item of the list.
\end{itemize}
You can even display poetry.
\begin{verse}
   There is an environment
    for verse \\             % The \\ command separates lines
   Whose features some poets % within a stanza.
   will curse.

                             % One or more blank lines separate stanzas.

   For instead of making\\
   Them do \emph{all} line breaking, \\
   It allows them to put too many words on a line when they'd rather be
   forced to be terse.
\end{verse}

Mathematical formulas may also be displayed.  A
displayed formula
is
one-line long; multiline
formulas require special formatting instructions.
   \[  \ip{\Gamma}{\psi'} = x'' + y^{2} + z_{i}^{n}\]
Don't start a paragraph with a displayed equation,
nor make one a paragraph by itself.

\end{document}               % End of document.

\section{Bench Huge Section 204}
% synthetic bench section 204
% This is a sample LaTeX input file.  (Version of 12 August 2004.)
%
% A '%' character causes TeX to ignore all remaining text on the line,
% and is used for comments like this one.

\documentclass{article}      % Specifies the document class

                             % The preamble begins here.
\title{An Example Document}  % Declares the document's title.
\author{Leslie Lamport}      % Declares the author's name.
\date{January 21, 1994}      % Deleting this command produces today's date.

\newcommand{\ip}[2]{(#1, #2)}
                             % Defines \ip{arg1}{arg2} to mean
                             % (arg1, arg2).

%\newcommand{\ip}[2]{\langle #1 | #2\rangle}
                             % This is an alternative definition of
                             % \ip that is commented out.

\begin{document}             % End of preamble and beginning of text.

\maketitle                   % Produces the title.

This is an example input file.  Comparing it with
the output it generates can show you how to
produce a simple document of your own.

\section{Ordinary Text}      % Produces section heading.  Lower-level
                             % sections are begun with similar
                             % \subsection and \subsubsection commands.

The ends  of words and sentences are marked
  by   spaces. It  doesn't matter how many
spaces    you type; one is as good as 100.  The
end of   a line counts as a space.

One   or more   blank lines denote the  end
of  a paragraph.

Since any number of consecutive spaces are treated
like a single one, the formatting of the input
file makes no difference to
      \LaTeX,                % The \LaTeX command generates the LaTeX logo.
but it makes a difference to you.  When you use
\LaTeX, making your input file as easy to read
as possible will be a great help as you write
your document and when you change it.  This sample
file shows how you can add comments to your own input
file.

Because printing is different from typewriting,
there are a number of things that you have to do
differently when preparing an input file than if
you were just typing the document directly.
Quotation marks like
       ``this''
have to be handled specially, as do quotes within
quotes:
       ``\,`this'            % \, separates the double and single quote.
        is what I just
        wrote, not  `that'\,''.

Dashes come in three sizes: an
       intra-word
dash, a medium dash for number ranges like
       1--2,
and a punctuation
       dash---like
this.

A sentence-ending space should be larger than the
space between words within a sentence.  You
sometimes have to type special commands in
conjunction with punctuation characters to get
this right, as in the following sentence.
       Gnats, gnus, etc.\ all  % `\ ' makes an inter-word space.
       begin with G\@.         % \@ marks end-of-sentence punctuation.
You should check the spaces after periods when
reading your output to make sure you haven't
forgotten any special cases.  Generating an
ellipsis
       \ldots\               % `\ ' is needed after `\ldots' because TeX
                             % ignores spaces after command names like \ldots
                             % made from \ + letters.
                             %
                             % Note how a `%' character causes TeX to ignore
                             % the end of the input line, so these blank lines
                             % do not start a new paragraph.
                             %
with the right spacing around the periods requires
a special command.

\LaTeX\ interprets some common characters as
commands, so you must type special commands to
generate them.  These characters include the
following:
       \$ \& \% \# \{ and \}.

In printing, text is usually emphasized with an
       \emph{italic}
type style.

\begin{em}
   A long segment of text can also be emphasized
   in this way.  Text within such a segment can be
   given \emph{additional} emphasis.
\end{em}

It is sometimes necessary to prevent \LaTeX\ from
breaking a line where it might otherwise do so.
This may be at a space, as between the ``Mr.''\ and
``Jones'' in
       ``Mr.~Jones'',        % ~ produces an unbreakable interword space.
or within a word---especially when the word is a
symbol like
       \mbox{\emph{itemnum}}
that makes little sense when hyphenated across
lines.

Footnotes\footnote{This is an example of a footnote.}
pose no problem.

\LaTeX\ is good at typesetting mathematical formulas
like
       \( x-3y + z = 7 \)
or
       \( a_{1} > x^{2n} + y^{2n} > x' \)
or
       \( \ip{A}{B} = \sum_{i} a_{i} b_{i} \).
The spaces you type in a formula are
ignored.  Remember that a letter like
       $x$                   % $ ... $  and  \( ... \)  are equivalent
is a formula when it denotes a mathematical
symbol, and it should be typed as one.

\section{Displayed Text}

Text is displayed by indenting it from the left
margin.  Quotations are commonly displayed.  There
are short quotations
\begin{quote}
   This is a short quotation.  It consists of a
   single paragraph of text.  See how it is formatted.
\end{quote}
and longer ones.
\begin{quotation}
   This is a longer quotation.  It consists of two
   paragraphs of text, neither of which are
   particularly interesting.

   This is the second paragraph of the quotation.  It
   is just as dull as the first paragraph.
\end{quotation}
Another frequently-displayed structure is a list.
The following is an example of an \emph{itemized}
list.
\begin{itemize}
   \item This is the first item of an itemized list.
         Each item in the list is marked with a ``tick''.
         You don't have to worry about what kind of tick
         mark is used.

   \item This is the second item of the list.  It
         contains another list nested inside it.  The inner
         list is an \emph{enumerated} list.
         \begin{enumerate}
            \item This is the first item of an enumerated
                  list that is nested within the itemized list.

            \item This is the second item of the inner list.
                  \LaTeX\ allows you to nest lists deeper than
                  you really should.
         \end{enumerate}
         This is the rest of the second item of the outer
         list.  It is no more interesting than any other
         part of the item.
   \item This is the third item of the list.
\end{itemize}
You can even display poetry.
\begin{verse}
   There is an environment
    for verse \\             % The \\ command separates lines
   Whose features some poets % within a stanza.
   will curse.

                             % One or more blank lines separate stanzas.

   For instead of making\\
   Them do \emph{all} line breaking, \\
   It allows them to put too many words on a line when they'd rather be
   forced to be terse.
\end{verse}

Mathematical formulas may also be displayed.  A
displayed formula
is
one-line long; multiline
formulas require special formatting instructions.
   \[  \ip{\Gamma}{\psi'} = x'' + y^{2} + z_{i}^{n}\]
Don't start a paragraph with a displayed equation,
nor make one a paragraph by itself.

\end{document}               % End of document.

\section{Bench Huge Section 205}
% synthetic bench section 205
% This is a sample LaTeX input file.  (Version of 12 August 2004.)
%
% A '%' character causes TeX to ignore all remaining text on the line,
% and is used for comments like this one.

\documentclass{article}      % Specifies the document class

                             % The preamble begins here.
\title{An Example Document}  % Declares the document's title.
\author{Leslie Lamport}      % Declares the author's name.
\date{January 21, 1994}      % Deleting this command produces today's date.

\newcommand{\ip}[2]{(#1, #2)}
                             % Defines \ip{arg1}{arg2} to mean
                             % (arg1, arg2).

%\newcommand{\ip}[2]{\langle #1 | #2\rangle}
                             % This is an alternative definition of
                             % \ip that is commented out.

\begin{document}             % End of preamble and beginning of text.

\maketitle                   % Produces the title.

This is an example input file.  Comparing it with
the output it generates can show you how to
produce a simple document of your own.

\section{Ordinary Text}      % Produces section heading.  Lower-level
                             % sections are begun with similar
                             % \subsection and \subsubsection commands.

The ends  of words and sentences are marked
  by   spaces. It  doesn't matter how many
spaces    you type; one is as good as 100.  The
end of   a line counts as a space.

One   or more   blank lines denote the  end
of  a paragraph.

Since any number of consecutive spaces are treated
like a single one, the formatting of the input
file makes no difference to
      \LaTeX,                % The \LaTeX command generates the LaTeX logo.
but it makes a difference to you.  When you use
\LaTeX, making your input file as easy to read
as possible will be a great help as you write
your document and when you change it.  This sample
file shows how you can add comments to your own input
file.

Because printing is different from typewriting,
there are a number of things that you have to do
differently when preparing an input file than if
you were just typing the document directly.
Quotation marks like
       ``this''
have to be handled specially, as do quotes within
quotes:
       ``\,`this'            % \, separates the double and single quote.
        is what I just
        wrote, not  `that'\,''.

Dashes come in three sizes: an
       intra-word
dash, a medium dash for number ranges like
       1--2,
and a punctuation
       dash---like
this.

A sentence-ending space should be larger than the
space between words within a sentence.  You
sometimes have to type special commands in
conjunction with punctuation characters to get
this right, as in the following sentence.
       Gnats, gnus, etc.\ all  % `\ ' makes an inter-word space.
       begin with G\@.         % \@ marks end-of-sentence punctuation.
You should check the spaces after periods when
reading your output to make sure you haven't
forgotten any special cases.  Generating an
ellipsis
       \ldots\               % `\ ' is needed after `\ldots' because TeX
                             % ignores spaces after command names like \ldots
                             % made from \ + letters.
                             %
                             % Note how a `%' character causes TeX to ignore
                             % the end of the input line, so these blank lines
                             % do not start a new paragraph.
                             %
with the right spacing around the periods requires
a special command.

\LaTeX\ interprets some common characters as
commands, so you must type special commands to
generate them.  These characters include the
following:
       \$ \& \% \# \{ and \}.

In printing, text is usually emphasized with an
       \emph{italic}
type style.

\begin{em}
   A long segment of text can also be emphasized
   in this way.  Text within such a segment can be
   given \emph{additional} emphasis.
\end{em}

It is sometimes necessary to prevent \LaTeX\ from
breaking a line where it might otherwise do so.
This may be at a space, as between the ``Mr.''\ and
``Jones'' in
       ``Mr.~Jones'',        % ~ produces an unbreakable interword space.
or within a word---especially when the word is a
symbol like
       \mbox{\emph{itemnum}}
that makes little sense when hyphenated across
lines.

Footnotes\footnote{This is an example of a footnote.}
pose no problem.

\LaTeX\ is good at typesetting mathematical formulas
like
       \( x-3y + z = 7 \)
or
       \( a_{1} > x^{2n} + y^{2n} > x' \)
or
       \( \ip{A}{B} = \sum_{i} a_{i} b_{i} \).
The spaces you type in a formula are
ignored.  Remember that a letter like
       $x$                   % $ ... $  and  \( ... \)  are equivalent
is a formula when it denotes a mathematical
symbol, and it should be typed as one.

\section{Displayed Text}

Text is displayed by indenting it from the left
margin.  Quotations are commonly displayed.  There
are short quotations
\begin{quote}
   This is a short quotation.  It consists of a
   single paragraph of text.  See how it is formatted.
\end{quote}
and longer ones.
\begin{quotation}
   This is a longer quotation.  It consists of two
   paragraphs of text, neither of which are
   particularly interesting.

   This is the second paragraph of the quotation.  It
   is just as dull as the first paragraph.
\end{quotation}
Another frequently-displayed structure is a list.
The following is an example of an \emph{itemized}
list.
\begin{itemize}
   \item This is the first item of an itemized list.
         Each item in the list is marked with a ``tick''.
         You don't have to worry about what kind of tick
         mark is used.

   \item This is the second item of the list.  It
         contains another list nested inside it.  The inner
         list is an \emph{enumerated} list.
         \begin{enumerate}
            \item This is the first item of an enumerated
                  list that is nested within the itemized list.

            \item This is the second item of the inner list.
                  \LaTeX\ allows you to nest lists deeper than
                  you really should.
         \end{enumerate}
         This is the rest of the second item of the outer
         list.  It is no more interesting than any other
         part of the item.
   \item This is the third item of the list.
\end{itemize}
You can even display poetry.
\begin{verse}
   There is an environment
    for verse \\             % The \\ command separates lines
   Whose features some poets % within a stanza.
   will curse.

                             % One or more blank lines separate stanzas.

   For instead of making\\
   Them do \emph{all} line breaking, \\
   It allows them to put too many words on a line when they'd rather be
   forced to be terse.
\end{verse}

Mathematical formulas may also be displayed.  A
displayed formula
is
one-line long; multiline
formulas require special formatting instructions.
   \[  \ip{\Gamma}{\psi'} = x'' + y^{2} + z_{i}^{n}\]
Don't start a paragraph with a displayed equation,
nor make one a paragraph by itself.

\end{document}               % End of document.

\section{Bench Huge Section 206}
% synthetic bench section 206
% This is a sample LaTeX input file.  (Version of 12 August 2004.)
%
% A '%' character causes TeX to ignore all remaining text on the line,
% and is used for comments like this one.

\documentclass{article}      % Specifies the document class

                             % The preamble begins here.
\title{An Example Document}  % Declares the document's title.
\author{Leslie Lamport}      % Declares the author's name.
\date{January 21, 1994}      % Deleting this command produces today's date.

\newcommand{\ip}[2]{(#1, #2)}
                             % Defines \ip{arg1}{arg2} to mean
                             % (arg1, arg2).

%\newcommand{\ip}[2]{\langle #1 | #2\rangle}
                             % This is an alternative definition of
                             % \ip that is commented out.

\begin{document}             % End of preamble and beginning of text.

\maketitle                   % Produces the title.

This is an example input file.  Comparing it with
the output it generates can show you how to
produce a simple document of your own.

\section{Ordinary Text}      % Produces section heading.  Lower-level
                             % sections are begun with similar
                             % \subsection and \subsubsection commands.

The ends  of words and sentences are marked
  by   spaces. It  doesn't matter how many
spaces    you type; one is as good as 100.  The
end of   a line counts as a space.

One   or more   blank lines denote the  end
of  a paragraph.

Since any number of consecutive spaces are treated
like a single one, the formatting of the input
file makes no difference to
      \LaTeX,                % The \LaTeX command generates the LaTeX logo.
but it makes a difference to you.  When you use
\LaTeX, making your input file as easy to read
as possible will be a great help as you write
your document and when you change it.  This sample
file shows how you can add comments to your own input
file.

Because printing is different from typewriting,
there are a number of things that you have to do
differently when preparing an input file than if
you were just typing the document directly.
Quotation marks like
       ``this''
have to be handled specially, as do quotes within
quotes:
       ``\,`this'            % \, separates the double and single quote.
        is what I just
        wrote, not  `that'\,''.

Dashes come in three sizes: an
       intra-word
dash, a medium dash for number ranges like
       1--2,
and a punctuation
       dash---like
this.

A sentence-ending space should be larger than the
space between words within a sentence.  You
sometimes have to type special commands in
conjunction with punctuation characters to get
this right, as in the following sentence.
       Gnats, gnus, etc.\ all  % `\ ' makes an inter-word space.
       begin with G\@.         % \@ marks end-of-sentence punctuation.
You should check the spaces after periods when
reading your output to make sure you haven't
forgotten any special cases.  Generating an
ellipsis
       \ldots\               % `\ ' is needed after `\ldots' because TeX
                             % ignores spaces after command names like \ldots
                             % made from \ + letters.
                             %
                             % Note how a `%' character causes TeX to ignore
                             % the end of the input line, so these blank lines
                             % do not start a new paragraph.
                             %
with the right spacing around the periods requires
a special command.

\LaTeX\ interprets some common characters as
commands, so you must type special commands to
generate them.  These characters include the
following:
       \$ \& \% \# \{ and \}.

In printing, text is usually emphasized with an
       \emph{italic}
type style.

\begin{em}
   A long segment of text can also be emphasized
   in this way.  Text within such a segment can be
   given \emph{additional} emphasis.
\end{em}

It is sometimes necessary to prevent \LaTeX\ from
breaking a line where it might otherwise do so.
This may be at a space, as between the ``Mr.''\ and
``Jones'' in
       ``Mr.~Jones'',        % ~ produces an unbreakable interword space.
or within a word---especially when the word is a
symbol like
       \mbox{\emph{itemnum}}
that makes little sense when hyphenated across
lines.

Footnotes\footnote{This is an example of a footnote.}
pose no problem.

\LaTeX\ is good at typesetting mathematical formulas
like
       \( x-3y + z = 7 \)
or
       \( a_{1} > x^{2n} + y^{2n} > x' \)
or
       \( \ip{A}{B} = \sum_{i} a_{i} b_{i} \).
The spaces you type in a formula are
ignored.  Remember that a letter like
       $x$                   % $ ... $  and  \( ... \)  are equivalent
is a formula when it denotes a mathematical
symbol, and it should be typed as one.

\section{Displayed Text}

Text is displayed by indenting it from the left
margin.  Quotations are commonly displayed.  There
are short quotations
\begin{quote}
   This is a short quotation.  It consists of a
   single paragraph of text.  See how it is formatted.
\end{quote}
and longer ones.
\begin{quotation}
   This is a longer quotation.  It consists of two
   paragraphs of text, neither of which are
   particularly interesting.

   This is the second paragraph of the quotation.  It
   is just as dull as the first paragraph.
\end{quotation}
Another frequently-displayed structure is a list.
The following is an example of an \emph{itemized}
list.
\begin{itemize}
   \item This is the first item of an itemized list.
         Each item in the list is marked with a ``tick''.
         You don't have to worry about what kind of tick
         mark is used.

   \item This is the second item of the list.  It
         contains another list nested inside it.  The inner
         list is an \emph{enumerated} list.
         \begin{enumerate}
            \item This is the first item of an enumerated
                  list that is nested within the itemized list.

            \item This is the second item of the inner list.
                  \LaTeX\ allows you to nest lists deeper than
                  you really should.
         \end{enumerate}
         This is the rest of the second item of the outer
         list.  It is no more interesting than any other
         part of the item.
   \item This is the third item of the list.
\end{itemize}
You can even display poetry.
\begin{verse}
   There is an environment
    for verse \\             % The \\ command separates lines
   Whose features some poets % within a stanza.
   will curse.

                             % One or more blank lines separate stanzas.

   For instead of making\\
   Them do \emph{all} line breaking, \\
   It allows them to put too many words on a line when they'd rather be
   forced to be terse.
\end{verse}

Mathematical formulas may also be displayed.  A
displayed formula
is
one-line long; multiline
formulas require special formatting instructions.
   \[  \ip{\Gamma}{\psi'} = x'' + y^{2} + z_{i}^{n}\]
Don't start a paragraph with a displayed equation,
nor make one a paragraph by itself.

\end{document}               % End of document.

\section{Bench Huge Section 207}
% synthetic bench section 207
% This is a sample LaTeX input file.  (Version of 12 August 2004.)
%
% A '%' character causes TeX to ignore all remaining text on the line,
% and is used for comments like this one.

\documentclass{article}      % Specifies the document class

                             % The preamble begins here.
\title{An Example Document}  % Declares the document's title.
\author{Leslie Lamport}      % Declares the author's name.
\date{January 21, 1994}      % Deleting this command produces today's date.

\newcommand{\ip}[2]{(#1, #2)}
                             % Defines \ip{arg1}{arg2} to mean
                             % (arg1, arg2).

%\newcommand{\ip}[2]{\langle #1 | #2\rangle}
                             % This is an alternative definition of
                             % \ip that is commented out.

\begin{document}             % End of preamble and beginning of text.

\maketitle                   % Produces the title.

This is an example input file.  Comparing it with
the output it generates can show you how to
produce a simple document of your own.

\section{Ordinary Text}      % Produces section heading.  Lower-level
                             % sections are begun with similar
                             % \subsection and \subsubsection commands.

The ends  of words and sentences are marked
  by   spaces. It  doesn't matter how many
spaces    you type; one is as good as 100.  The
end of   a line counts as a space.

One   or more   blank lines denote the  end
of  a paragraph.

Since any number of consecutive spaces are treated
like a single one, the formatting of the input
file makes no difference to
      \LaTeX,                % The \LaTeX command generates the LaTeX logo.
but it makes a difference to you.  When you use
\LaTeX, making your input file as easy to read
as possible will be a great help as you write
your document and when you change it.  This sample
file shows how you can add comments to your own input
file.

Because printing is different from typewriting,
there are a number of things that you have to do
differently when preparing an input file than if
you were just typing the document directly.
Quotation marks like
       ``this''
have to be handled specially, as do quotes within
quotes:
       ``\,`this'            % \, separates the double and single quote.
        is what I just
        wrote, not  `that'\,''.

Dashes come in three sizes: an
       intra-word
dash, a medium dash for number ranges like
       1--2,
and a punctuation
       dash---like
this.

A sentence-ending space should be larger than the
space between words within a sentence.  You
sometimes have to type special commands in
conjunction with punctuation characters to get
this right, as in the following sentence.
       Gnats, gnus, etc.\ all  % `\ ' makes an inter-word space.
       begin with G\@.         % \@ marks end-of-sentence punctuation.
You should check the spaces after periods when
reading your output to make sure you haven't
forgotten any special cases.  Generating an
ellipsis
       \ldots\               % `\ ' is needed after `\ldots' because TeX
                             % ignores spaces after command names like \ldots
                             % made from \ + letters.
                             %
                             % Note how a `%' character causes TeX to ignore
                             % the end of the input line, so these blank lines
                             % do not start a new paragraph.
                             %
with the right spacing around the periods requires
a special command.

\LaTeX\ interprets some common characters as
commands, so you must type special commands to
generate them.  These characters include the
following:
       \$ \& \% \# \{ and \}.

In printing, text is usually emphasized with an
       \emph{italic}
type style.

\begin{em}
   A long segment of text can also be emphasized
   in this way.  Text within such a segment can be
   given \emph{additional} emphasis.
\end{em}

It is sometimes necessary to prevent \LaTeX\ from
breaking a line where it might otherwise do so.
This may be at a space, as between the ``Mr.''\ and
``Jones'' in
       ``Mr.~Jones'',        % ~ produces an unbreakable interword space.
or within a word---especially when the word is a
symbol like
       \mbox{\emph{itemnum}}
that makes little sense when hyphenated across
lines.

Footnotes\footnote{This is an example of a footnote.}
pose no problem.

\LaTeX\ is good at typesetting mathematical formulas
like
       \( x-3y + z = 7 \)
or
       \( a_{1} > x^{2n} + y^{2n} > x' \)
or
       \( \ip{A}{B} = \sum_{i} a_{i} b_{i} \).
The spaces you type in a formula are
ignored.  Remember that a letter like
       $x$                   % $ ... $  and  \( ... \)  are equivalent
is a formula when it denotes a mathematical
symbol, and it should be typed as one.

\section{Displayed Text}

Text is displayed by indenting it from the left
margin.  Quotations are commonly displayed.  There
are short quotations
\begin{quote}
   This is a short quotation.  It consists of a
   single paragraph of text.  See how it is formatted.
\end{quote}
and longer ones.
\begin{quotation}
   This is a longer quotation.  It consists of two
   paragraphs of text, neither of which are
   particularly interesting.

   This is the second paragraph of the quotation.  It
   is just as dull as the first paragraph.
\end{quotation}
Another frequently-displayed structure is a list.
The following is an example of an \emph{itemized}
list.
\begin{itemize}
   \item This is the first item of an itemized list.
         Each item in the list is marked with a ``tick''.
         You don't have to worry about what kind of tick
         mark is used.

   \item This is the second item of the list.  It
         contains another list nested inside it.  The inner
         list is an \emph{enumerated} list.
         \begin{enumerate}
            \item This is the first item of an enumerated
                  list that is nested within the itemized list.

            \item This is the second item of the inner list.
                  \LaTeX\ allows you to nest lists deeper than
                  you really should.
         \end{enumerate}
         This is the rest of the second item of the outer
         list.  It is no more interesting than any other
         part of the item.
   \item This is the third item of the list.
\end{itemize}
You can even display poetry.
\begin{verse}
   There is an environment
    for verse \\             % The \\ command separates lines
   Whose features some poets % within a stanza.
   will curse.

                             % One or more blank lines separate stanzas.

   For instead of making\\
   Them do \emph{all} line breaking, \\
   It allows them to put too many words on a line when they'd rather be
   forced to be terse.
\end{verse}

Mathematical formulas may also be displayed.  A
displayed formula
is
one-line long; multiline
formulas require special formatting instructions.
   \[  \ip{\Gamma}{\psi'} = x'' + y^{2} + z_{i}^{n}\]
Don't start a paragraph with a displayed equation,
nor make one a paragraph by itself.

\end{document}               % End of document.

\section{Bench Huge Section 208}
% synthetic bench section 208
% This is a sample LaTeX input file.  (Version of 12 August 2004.)
%
% A '%' character causes TeX to ignore all remaining text on the line,
% and is used for comments like this one.

\documentclass{article}      % Specifies the document class

                             % The preamble begins here.
\title{An Example Document}  % Declares the document's title.
\author{Leslie Lamport}      % Declares the author's name.
\date{January 21, 1994}      % Deleting this command produces today's date.

\newcommand{\ip}[2]{(#1, #2)}
                             % Defines \ip{arg1}{arg2} to mean
                             % (arg1, arg2).

%\newcommand{\ip}[2]{\langle #1 | #2\rangle}
                             % This is an alternative definition of
                             % \ip that is commented out.

\begin{document}             % End of preamble and beginning of text.

\maketitle                   % Produces the title.

This is an example input file.  Comparing it with
the output it generates can show you how to
produce a simple document of your own.

\section{Ordinary Text}      % Produces section heading.  Lower-level
                             % sections are begun with similar
                             % \subsection and \subsubsection commands.

The ends  of words and sentences are marked
  by   spaces. It  doesn't matter how many
spaces    you type; one is as good as 100.  The
end of   a line counts as a space.

One   or more   blank lines denote the  end
of  a paragraph.

Since any number of consecutive spaces are treated
like a single one, the formatting of the input
file makes no difference to
      \LaTeX,                % The \LaTeX command generates the LaTeX logo.
but it makes a difference to you.  When you use
\LaTeX, making your input file as easy to read
as possible will be a great help as you write
your document and when you change it.  This sample
file shows how you can add comments to your own input
file.

Because printing is different from typewriting,
there are a number of things that you have to do
differently when preparing an input file than if
you were just typing the document directly.
Quotation marks like
       ``this''
have to be handled specially, as do quotes within
quotes:
       ``\,`this'            % \, separates the double and single quote.
        is what I just
        wrote, not  `that'\,''.

Dashes come in three sizes: an
       intra-word
dash, a medium dash for number ranges like
       1--2,
and a punctuation
       dash---like
this.

A sentence-ending space should be larger than the
space between words within a sentence.  You
sometimes have to type special commands in
conjunction with punctuation characters to get
this right, as in the following sentence.
       Gnats, gnus, etc.\ all  % `\ ' makes an inter-word space.
       begin with G\@.         % \@ marks end-of-sentence punctuation.
You should check the spaces after periods when
reading your output to make sure you haven't
forgotten any special cases.  Generating an
ellipsis
       \ldots\               % `\ ' is needed after `\ldots' because TeX
                             % ignores spaces after command names like \ldots
                             % made from \ + letters.
                             %
                             % Note how a `%' character causes TeX to ignore
                             % the end of the input line, so these blank lines
                             % do not start a new paragraph.
                             %
with the right spacing around the periods requires
a special command.

\LaTeX\ interprets some common characters as
commands, so you must type special commands to
generate them.  These characters include the
following:
       \$ \& \% \# \{ and \}.

In printing, text is usually emphasized with an
       \emph{italic}
type style.

\begin{em}
   A long segment of text can also be emphasized
   in this way.  Text within such a segment can be
   given \emph{additional} emphasis.
\end{em}

It is sometimes necessary to prevent \LaTeX\ from
breaking a line where it might otherwise do so.
This may be at a space, as between the ``Mr.''\ and
``Jones'' in
       ``Mr.~Jones'',        % ~ produces an unbreakable interword space.
or within a word---especially when the word is a
symbol like
       \mbox{\emph{itemnum}}
that makes little sense when hyphenated across
lines.

Footnotes\footnote{This is an example of a footnote.}
pose no problem.

\LaTeX\ is good at typesetting mathematical formulas
like
       \( x-3y + z = 7 \)
or
       \( a_{1} > x^{2n} + y^{2n} > x' \)
or
       \( \ip{A}{B} = \sum_{i} a_{i} b_{i} \).
The spaces you type in a formula are
ignored.  Remember that a letter like
       $x$                   % $ ... $  and  \( ... \)  are equivalent
is a formula when it denotes a mathematical
symbol, and it should be typed as one.

\section{Displayed Text}

Text is displayed by indenting it from the left
margin.  Quotations are commonly displayed.  There
are short quotations
\begin{quote}
   This is a short quotation.  It consists of a
   single paragraph of text.  See how it is formatted.
\end{quote}
and longer ones.
\begin{quotation}
   This is a longer quotation.  It consists of two
   paragraphs of text, neither of which are
   particularly interesting.

   This is the second paragraph of the quotation.  It
   is just as dull as the first paragraph.
\end{quotation}
Another frequently-displayed structure is a list.
The following is an example of an \emph{itemized}
list.
\begin{itemize}
   \item This is the first item of an itemized list.
         Each item in the list is marked with a ``tick''.
         You don't have to worry about what kind of tick
         mark is used.

   \item This is the second item of the list.  It
         contains another list nested inside it.  The inner
         list is an \emph{enumerated} list.
         \begin{enumerate}
            \item This is the first item of an enumerated
                  list that is nested within the itemized list.

            \item This is the second item of the inner list.
                  \LaTeX\ allows you to nest lists deeper than
                  you really should.
         \end{enumerate}
         This is the rest of the second item of the outer
         list.  It is no more interesting than any other
         part of the item.
   \item This is the third item of the list.
\end{itemize}
You can even display poetry.
\begin{verse}
   There is an environment
    for verse \\             % The \\ command separates lines
   Whose features some poets % within a stanza.
   will curse.

                             % One or more blank lines separate stanzas.

   For instead of making\\
   Them do \emph{all} line breaking, \\
   It allows them to put too many words on a line when they'd rather be
   forced to be terse.
\end{verse}

Mathematical formulas may also be displayed.  A
displayed formula
is
one-line long; multiline
formulas require special formatting instructions.
   \[  \ip{\Gamma}{\psi'} = x'' + y^{2} + z_{i}^{n}\]
Don't start a paragraph with a displayed equation,
nor make one a paragraph by itself.

\end{document}               % End of document.

\section{Bench Huge Section 209}
% synthetic bench section 209
% This is a sample LaTeX input file.  (Version of 12 August 2004.)
%
% A '%' character causes TeX to ignore all remaining text on the line,
% and is used for comments like this one.

\documentclass{article}      % Specifies the document class

                             % The preamble begins here.
\title{An Example Document}  % Declares the document's title.
\author{Leslie Lamport}      % Declares the author's name.
\date{January 21, 1994}      % Deleting this command produces today's date.

\newcommand{\ip}[2]{(#1, #2)}
                             % Defines \ip{arg1}{arg2} to mean
                             % (arg1, arg2).

%\newcommand{\ip}[2]{\langle #1 | #2\rangle}
                             % This is an alternative definition of
                             % \ip that is commented out.

\begin{document}             % End of preamble and beginning of text.

\maketitle                   % Produces the title.

This is an example input file.  Comparing it with
the output it generates can show you how to
produce a simple document of your own.

\section{Ordinary Text}      % Produces section heading.  Lower-level
                             % sections are begun with similar
                             % \subsection and \subsubsection commands.

The ends  of words and sentences are marked
  by   spaces. It  doesn't matter how many
spaces    you type; one is as good as 100.  The
end of   a line counts as a space.

One   or more   blank lines denote the  end
of  a paragraph.

Since any number of consecutive spaces are treated
like a single one, the formatting of the input
file makes no difference to
      \LaTeX,                % The \LaTeX command generates the LaTeX logo.
but it makes a difference to you.  When you use
\LaTeX, making your input file as easy to read
as possible will be a great help as you write
your document and when you change it.  This sample
file shows how you can add comments to your own input
file.

Because printing is different from typewriting,
there are a number of things that you have to do
differently when preparing an input file than if
you were just typing the document directly.
Quotation marks like
       ``this''
have to be handled specially, as do quotes within
quotes:
       ``\,`this'            % \, separates the double and single quote.
        is what I just
        wrote, not  `that'\,''.

Dashes come in three sizes: an
       intra-word
dash, a medium dash for number ranges like
       1--2,
and a punctuation
       dash---like
this.

A sentence-ending space should be larger than the
space between words within a sentence.  You
sometimes have to type special commands in
conjunction with punctuation characters to get
this right, as in the following sentence.
       Gnats, gnus, etc.\ all  % `\ ' makes an inter-word space.
       begin with G\@.         % \@ marks end-of-sentence punctuation.
You should check the spaces after periods when
reading your output to make sure you haven't
forgotten any special cases.  Generating an
ellipsis
       \ldots\               % `\ ' is needed after `\ldots' because TeX
                             % ignores spaces after command names like \ldots
                             % made from \ + letters.
                             %
                             % Note how a `%' character causes TeX to ignore
                             % the end of the input line, so these blank lines
                             % do not start a new paragraph.
                             %
with the right spacing around the periods requires
a special command.

\LaTeX\ interprets some common characters as
commands, so you must type special commands to
generate them.  These characters include the
following:
       \$ \& \% \# \{ and \}.

In printing, text is usually emphasized with an
       \emph{italic}
type style.

\begin{em}
   A long segment of text can also be emphasized
   in this way.  Text within such a segment can be
   given \emph{additional} emphasis.
\end{em}

It is sometimes necessary to prevent \LaTeX\ from
breaking a line where it might otherwise do so.
This may be at a space, as between the ``Mr.''\ and
``Jones'' in
       ``Mr.~Jones'',        % ~ produces an unbreakable interword space.
or within a word---especially when the word is a
symbol like
       \mbox{\emph{itemnum}}
that makes little sense when hyphenated across
lines.

Footnotes\footnote{This is an example of a footnote.}
pose no problem.

\LaTeX\ is good at typesetting mathematical formulas
like
       \( x-3y + z = 7 \)
or
       \( a_{1} > x^{2n} + y^{2n} > x' \)
or
       \( \ip{A}{B} = \sum_{i} a_{i} b_{i} \).
The spaces you type in a formula are
ignored.  Remember that a letter like
       $x$                   % $ ... $  and  \( ... \)  are equivalent
is a formula when it denotes a mathematical
symbol, and it should be typed as one.

\section{Displayed Text}

Text is displayed by indenting it from the left
margin.  Quotations are commonly displayed.  There
are short quotations
\begin{quote}
   This is a short quotation.  It consists of a
   single paragraph of text.  See how it is formatted.
\end{quote}
and longer ones.
\begin{quotation}
   This is a longer quotation.  It consists of two
   paragraphs of text, neither of which are
   particularly interesting.

   This is the second paragraph of the quotation.  It
   is just as dull as the first paragraph.
\end{quotation}
Another frequently-displayed structure is a list.
The following is an example of an \emph{itemized}
list.
\begin{itemize}
   \item This is the first item of an itemized list.
         Each item in the list is marked with a ``tick''.
         You don't have to worry about what kind of tick
         mark is used.

   \item This is the second item of the list.  It
         contains another list nested inside it.  The inner
         list is an \emph{enumerated} list.
         \begin{enumerate}
            \item This is the first item of an enumerated
                  list that is nested within the itemized list.

            \item This is the second item of the inner list.
                  \LaTeX\ allows you to nest lists deeper than
                  you really should.
         \end{enumerate}
         This is the rest of the second item of the outer
         list.  It is no more interesting than any other
         part of the item.
   \item This is the third item of the list.
\end{itemize}
You can even display poetry.
\begin{verse}
   There is an environment
    for verse \\             % The \\ command separates lines
   Whose features some poets % within a stanza.
   will curse.

                             % One or more blank lines separate stanzas.

   For instead of making\\
   Them do \emph{all} line breaking, \\
   It allows them to put too many words on a line when they'd rather be
   forced to be terse.
\end{verse}

Mathematical formulas may also be displayed.  A
displayed formula
is
one-line long; multiline
formulas require special formatting instructions.
   \[  \ip{\Gamma}{\psi'} = x'' + y^{2} + z_{i}^{n}\]
Don't start a paragraph with a displayed equation,
nor make one a paragraph by itself.

\end{document}               % End of document.

\section{Bench Huge Section 210}
% synthetic bench section 210
% This is a sample LaTeX input file.  (Version of 12 August 2004.)
%
% A '%' character causes TeX to ignore all remaining text on the line,
% and is used for comments like this one.

\documentclass{article}      % Specifies the document class

                             % The preamble begins here.
\title{An Example Document}  % Declares the document's title.
\author{Leslie Lamport}      % Declares the author's name.
\date{January 21, 1994}      % Deleting this command produces today's date.

\newcommand{\ip}[2]{(#1, #2)}
                             % Defines \ip{arg1}{arg2} to mean
                             % (arg1, arg2).

%\newcommand{\ip}[2]{\langle #1 | #2\rangle}
                             % This is an alternative definition of
                             % \ip that is commented out.

\begin{document}             % End of preamble and beginning of text.

\maketitle                   % Produces the title.

This is an example input file.  Comparing it with
the output it generates can show you how to
produce a simple document of your own.

\section{Ordinary Text}      % Produces section heading.  Lower-level
                             % sections are begun with similar
                             % \subsection and \subsubsection commands.

The ends  of words and sentences are marked
  by   spaces. It  doesn't matter how many
spaces    you type; one is as good as 100.  The
end of   a line counts as a space.

One   or more   blank lines denote the  end
of  a paragraph.

Since any number of consecutive spaces are treated
like a single one, the formatting of the input
file makes no difference to
      \LaTeX,                % The \LaTeX command generates the LaTeX logo.
but it makes a difference to you.  When you use
\LaTeX, making your input file as easy to read
as possible will be a great help as you write
your document and when you change it.  This sample
file shows how you can add comments to your own input
file.

Because printing is different from typewriting,
there are a number of things that you have to do
differently when preparing an input file than if
you were just typing the document directly.
Quotation marks like
       ``this''
have to be handled specially, as do quotes within
quotes:
       ``\,`this'            % \, separates the double and single quote.
        is what I just
        wrote, not  `that'\,''.

Dashes come in three sizes: an
       intra-word
dash, a medium dash for number ranges like
       1--2,
and a punctuation
       dash---like
this.

A sentence-ending space should be larger than the
space between words within a sentence.  You
sometimes have to type special commands in
conjunction with punctuation characters to get
this right, as in the following sentence.
       Gnats, gnus, etc.\ all  % `\ ' makes an inter-word space.
       begin with G\@.         % \@ marks end-of-sentence punctuation.
You should check the spaces after periods when
reading your output to make sure you haven't
forgotten any special cases.  Generating an
ellipsis
       \ldots\               % `\ ' is needed after `\ldots' because TeX
                             % ignores spaces after command names like \ldots
                             % made from \ + letters.
                             %
                             % Note how a `%' character causes TeX to ignore
                             % the end of the input line, so these blank lines
                             % do not start a new paragraph.
                             %
with the right spacing around the periods requires
a special command.

\LaTeX\ interprets some common characters as
commands, so you must type special commands to
generate them.  These characters include the
following:
       \$ \& \% \# \{ and \}.

In printing, text is usually emphasized with an
       \emph{italic}
type style.

\begin{em}
   A long segment of text can also be emphasized
   in this way.  Text within such a segment can be
   given \emph{additional} emphasis.
\end{em}

It is sometimes necessary to prevent \LaTeX\ from
breaking a line where it might otherwise do so.
This may be at a space, as between the ``Mr.''\ and
``Jones'' in
       ``Mr.~Jones'',        % ~ produces an unbreakable interword space.
or within a word---especially when the word is a
symbol like
       \mbox{\emph{itemnum}}
that makes little sense when hyphenated across
lines.

Footnotes\footnote{This is an example of a footnote.}
pose no problem.

\LaTeX\ is good at typesetting mathematical formulas
like
       \( x-3y + z = 7 \)
or
       \( a_{1} > x^{2n} + y^{2n} > x' \)
or
       \( \ip{A}{B} = \sum_{i} a_{i} b_{i} \).
The spaces you type in a formula are
ignored.  Remember that a letter like
       $x$                   % $ ... $  and  \( ... \)  are equivalent
is a formula when it denotes a mathematical
symbol, and it should be typed as one.

\section{Displayed Text}

Text is displayed by indenting it from the left
margin.  Quotations are commonly displayed.  There
are short quotations
\begin{quote}
   This is a short quotation.  It consists of a
   single paragraph of text.  See how it is formatted.
\end{quote}
and longer ones.
\begin{quotation}
   This is a longer quotation.  It consists of two
   paragraphs of text, neither of which are
   particularly interesting.

   This is the second paragraph of the quotation.  It
   is just as dull as the first paragraph.
\end{quotation}
Another frequently-displayed structure is a list.
The following is an example of an \emph{itemized}
list.
\begin{itemize}
   \item This is the first item of an itemized list.
         Each item in the list is marked with a ``tick''.
         You don't have to worry about what kind of tick
         mark is used.

   \item This is the second item of the list.  It
         contains another list nested inside it.  The inner
         list is an \emph{enumerated} list.
         \begin{enumerate}
            \item This is the first item of an enumerated
                  list that is nested within the itemized list.

            \item This is the second item of the inner list.
                  \LaTeX\ allows you to nest lists deeper than
                  you really should.
         \end{enumerate}
         This is the rest of the second item of the outer
         list.  It is no more interesting than any other
         part of the item.
   \item This is the third item of the list.
\end{itemize}
You can even display poetry.
\begin{verse}
   There is an environment
    for verse \\             % The \\ command separates lines
   Whose features some poets % within a stanza.
   will curse.

                             % One or more blank lines separate stanzas.

   For instead of making\\
   Them do \emph{all} line breaking, \\
   It allows them to put too many words on a line when they'd rather be
   forced to be terse.
\end{verse}

Mathematical formulas may also be displayed.  A
displayed formula
is
one-line long; multiline
formulas require special formatting instructions.
   \[  \ip{\Gamma}{\psi'} = x'' + y^{2} + z_{i}^{n}\]
Don't start a paragraph with a displayed equation,
nor make one a paragraph by itself.

\end{document}               % End of document.

\section{Bench Huge Section 211}
% synthetic bench section 211
% This is a sample LaTeX input file.  (Version of 12 August 2004.)
%
% A '%' character causes TeX to ignore all remaining text on the line,
% and is used for comments like this one.

\documentclass{article}      % Specifies the document class

                             % The preamble begins here.
\title{An Example Document}  % Declares the document's title.
\author{Leslie Lamport}      % Declares the author's name.
\date{January 21, 1994}      % Deleting this command produces today's date.

\newcommand{\ip}[2]{(#1, #2)}
                             % Defines \ip{arg1}{arg2} to mean
                             % (arg1, arg2).

%\newcommand{\ip}[2]{\langle #1 | #2\rangle}
                             % This is an alternative definition of
                             % \ip that is commented out.

\begin{document}             % End of preamble and beginning of text.

\maketitle                   % Produces the title.

This is an example input file.  Comparing it with
the output it generates can show you how to
produce a simple document of your own.

\section{Ordinary Text}      % Produces section heading.  Lower-level
                             % sections are begun with similar
                             % \subsection and \subsubsection commands.

The ends  of words and sentences are marked
  by   spaces. It  doesn't matter how many
spaces    you type; one is as good as 100.  The
end of   a line counts as a space.

One   or more   blank lines denote the  end
of  a paragraph.

Since any number of consecutive spaces are treated
like a single one, the formatting of the input
file makes no difference to
      \LaTeX,                % The \LaTeX command generates the LaTeX logo.
but it makes a difference to you.  When you use
\LaTeX, making your input file as easy to read
as possible will be a great help as you write
your document and when you change it.  This sample
file shows how you can add comments to your own input
file.

Because printing is different from typewriting,
there are a number of things that you have to do
differently when preparing an input file than if
you were just typing the document directly.
Quotation marks like
       ``this''
have to be handled specially, as do quotes within
quotes:
       ``\,`this'            % \, separates the double and single quote.
        is what I just
        wrote, not  `that'\,''.

Dashes come in three sizes: an
       intra-word
dash, a medium dash for number ranges like
       1--2,
and a punctuation
       dash---like
this.

A sentence-ending space should be larger than the
space between words within a sentence.  You
sometimes have to type special commands in
conjunction with punctuation characters to get
this right, as in the following sentence.
       Gnats, gnus, etc.\ all  % `\ ' makes an inter-word space.
       begin with G\@.         % \@ marks end-of-sentence punctuation.
You should check the spaces after periods when
reading your output to make sure you haven't
forgotten any special cases.  Generating an
ellipsis
       \ldots\               % `\ ' is needed after `\ldots' because TeX
                             % ignores spaces after command names like \ldots
                             % made from \ + letters.
                             %
                             % Note how a `%' character causes TeX to ignore
                             % the end of the input line, so these blank lines
                             % do not start a new paragraph.
                             %
with the right spacing around the periods requires
a special command.

\LaTeX\ interprets some common characters as
commands, so you must type special commands to
generate them.  These characters include the
following:
       \$ \& \% \# \{ and \}.

In printing, text is usually emphasized with an
       \emph{italic}
type style.

\begin{em}
   A long segment of text can also be emphasized
   in this way.  Text within such a segment can be
   given \emph{additional} emphasis.
\end{em}

It is sometimes necessary to prevent \LaTeX\ from
breaking a line where it might otherwise do so.
This may be at a space, as between the ``Mr.''\ and
``Jones'' in
       ``Mr.~Jones'',        % ~ produces an unbreakable interword space.
or within a word---especially when the word is a
symbol like
       \mbox{\emph{itemnum}}
that makes little sense when hyphenated across
lines.

Footnotes\footnote{This is an example of a footnote.}
pose no problem.

\LaTeX\ is good at typesetting mathematical formulas
like
       \( x-3y + z = 7 \)
or
       \( a_{1} > x^{2n} + y^{2n} > x' \)
or
       \( \ip{A}{B} = \sum_{i} a_{i} b_{i} \).
The spaces you type in a formula are
ignored.  Remember that a letter like
       $x$                   % $ ... $  and  \( ... \)  are equivalent
is a formula when it denotes a mathematical
symbol, and it should be typed as one.

\section{Displayed Text}

Text is displayed by indenting it from the left
margin.  Quotations are commonly displayed.  There
are short quotations
\begin{quote}
   This is a short quotation.  It consists of a
   single paragraph of text.  See how it is formatted.
\end{quote}
and longer ones.
\begin{quotation}
   This is a longer quotation.  It consists of two
   paragraphs of text, neither of which are
   particularly interesting.

   This is the second paragraph of the quotation.  It
   is just as dull as the first paragraph.
\end{quotation}
Another frequently-displayed structure is a list.
The following is an example of an \emph{itemized}
list.
\begin{itemize}
   \item This is the first item of an itemized list.
         Each item in the list is marked with a ``tick''.
         You don't have to worry about what kind of tick
         mark is used.

   \item This is the second item of the list.  It
         contains another list nested inside it.  The inner
         list is an \emph{enumerated} list.
         \begin{enumerate}
            \item This is the first item of an enumerated
                  list that is nested within the itemized list.

            \item This is the second item of the inner list.
                  \LaTeX\ allows you to nest lists deeper than
                  you really should.
         \end{enumerate}
         This is the rest of the second item of the outer
         list.  It is no more interesting than any other
         part of the item.
   \item This is the third item of the list.
\end{itemize}
You can even display poetry.
\begin{verse}
   There is an environment
    for verse \\             % The \\ command separates lines
   Whose features some poets % within a stanza.
   will curse.

                             % One or more blank lines separate stanzas.

   For instead of making\\
   Them do \emph{all} line breaking, \\
   It allows them to put too many words on a line when they'd rather be
   forced to be terse.
\end{verse}

Mathematical formulas may also be displayed.  A
displayed formula
is
one-line long; multiline
formulas require special formatting instructions.
   \[  \ip{\Gamma}{\psi'} = x'' + y^{2} + z_{i}^{n}\]
Don't start a paragraph with a displayed equation,
nor make one a paragraph by itself.

\end{document}               % End of document.

\section{Bench Huge Section 212}
% synthetic bench section 212
% This is a sample LaTeX input file.  (Version of 12 August 2004.)
%
% A '%' character causes TeX to ignore all remaining text on the line,
% and is used for comments like this one.

\documentclass{article}      % Specifies the document class

                             % The preamble begins here.
\title{An Example Document}  % Declares the document's title.
\author{Leslie Lamport}      % Declares the author's name.
\date{January 21, 1994}      % Deleting this command produces today's date.

\newcommand{\ip}[2]{(#1, #2)}
                             % Defines \ip{arg1}{arg2} to mean
                             % (arg1, arg2).

%\newcommand{\ip}[2]{\langle #1 | #2\rangle}
                             % This is an alternative definition of
                             % \ip that is commented out.

\begin{document}             % End of preamble and beginning of text.

\maketitle                   % Produces the title.

This is an example input file.  Comparing it with
the output it generates can show you how to
produce a simple document of your own.

\section{Ordinary Text}      % Produces section heading.  Lower-level
                             % sections are begun with similar
                             % \subsection and \subsubsection commands.

The ends  of words and sentences are marked
  by   spaces. It  doesn't matter how many
spaces    you type; one is as good as 100.  The
end of   a line counts as a space.

One   or more   blank lines denote the  end
of  a paragraph.

Since any number of consecutive spaces are treated
like a single one, the formatting of the input
file makes no difference to
      \LaTeX,                % The \LaTeX command generates the LaTeX logo.
but it makes a difference to you.  When you use
\LaTeX, making your input file as easy to read
as possible will be a great help as you write
your document and when you change it.  This sample
file shows how you can add comments to your own input
file.

Because printing is different from typewriting,
there are a number of things that you have to do
differently when preparing an input file than if
you were just typing the document directly.
Quotation marks like
       ``this''
have to be handled specially, as do quotes within
quotes:
       ``\,`this'            % \, separates the double and single quote.
        is what I just
        wrote, not  `that'\,''.

Dashes come in three sizes: an
       intra-word
dash, a medium dash for number ranges like
       1--2,
and a punctuation
       dash---like
this.

A sentence-ending space should be larger than the
space between words within a sentence.  You
sometimes have to type special commands in
conjunction with punctuation characters to get
this right, as in the following sentence.
       Gnats, gnus, etc.\ all  % `\ ' makes an inter-word space.
       begin with G\@.         % \@ marks end-of-sentence punctuation.
You should check the spaces after periods when
reading your output to make sure you haven't
forgotten any special cases.  Generating an
ellipsis
       \ldots\               % `\ ' is needed after `\ldots' because TeX
                             % ignores spaces after command names like \ldots
                             % made from \ + letters.
                             %
                             % Note how a `%' character causes TeX to ignore
                             % the end of the input line, so these blank lines
                             % do not start a new paragraph.
                             %
with the right spacing around the periods requires
a special command.

\LaTeX\ interprets some common characters as
commands, so you must type special commands to
generate them.  These characters include the
following:
       \$ \& \% \# \{ and \}.

In printing, text is usually emphasized with an
       \emph{italic}
type style.

\begin{em}
   A long segment of text can also be emphasized
   in this way.  Text within such a segment can be
   given \emph{additional} emphasis.
\end{em}

It is sometimes necessary to prevent \LaTeX\ from
breaking a line where it might otherwise do so.
This may be at a space, as between the ``Mr.''\ and
``Jones'' in
       ``Mr.~Jones'',        % ~ produces an unbreakable interword space.
or within a word---especially when the word is a
symbol like
       \mbox{\emph{itemnum}}
that makes little sense when hyphenated across
lines.

Footnotes\footnote{This is an example of a footnote.}
pose no problem.

\LaTeX\ is good at typesetting mathematical formulas
like
       \( x-3y + z = 7 \)
or
       \( a_{1} > x^{2n} + y^{2n} > x' \)
or
       \( \ip{A}{B} = \sum_{i} a_{i} b_{i} \).
The spaces you type in a formula are
ignored.  Remember that a letter like
       $x$                   % $ ... $  and  \( ... \)  are equivalent
is a formula when it denotes a mathematical
symbol, and it should be typed as one.

\section{Displayed Text}

Text is displayed by indenting it from the left
margin.  Quotations are commonly displayed.  There
are short quotations
\begin{quote}
   This is a short quotation.  It consists of a
   single paragraph of text.  See how it is formatted.
\end{quote}
and longer ones.
\begin{quotation}
   This is a longer quotation.  It consists of two
   paragraphs of text, neither of which are
   particularly interesting.

   This is the second paragraph of the quotation.  It
   is just as dull as the first paragraph.
\end{quotation}
Another frequently-displayed structure is a list.
The following is an example of an \emph{itemized}
list.
\begin{itemize}
   \item This is the first item of an itemized list.
         Each item in the list is marked with a ``tick''.
         You don't have to worry about what kind of tick
         mark is used.

   \item This is the second item of the list.  It
         contains another list nested inside it.  The inner
         list is an \emph{enumerated} list.
         \begin{enumerate}
            \item This is the first item of an enumerated
                  list that is nested within the itemized list.

            \item This is the second item of the inner list.
                  \LaTeX\ allows you to nest lists deeper than
                  you really should.
         \end{enumerate}
         This is the rest of the second item of the outer
         list.  It is no more interesting than any other
         part of the item.
   \item This is the third item of the list.
\end{itemize}
You can even display poetry.
\begin{verse}
   There is an environment
    for verse \\             % The \\ command separates lines
   Whose features some poets % within a stanza.
   will curse.

                             % One or more blank lines separate stanzas.

   For instead of making\\
   Them do \emph{all} line breaking, \\
   It allows them to put too many words on a line when they'd rather be
   forced to be terse.
\end{verse}

Mathematical formulas may also be displayed.  A
displayed formula
is
one-line long; multiline
formulas require special formatting instructions.
   \[  \ip{\Gamma}{\psi'} = x'' + y^{2} + z_{i}^{n}\]
Don't start a paragraph with a displayed equation,
nor make one a paragraph by itself.

\end{document}               % End of document.

\section{Bench Huge Section 213}
% synthetic bench section 213
% This is a sample LaTeX input file.  (Version of 12 August 2004.)
%
% A '%' character causes TeX to ignore all remaining text on the line,
% and is used for comments like this one.

\documentclass{article}      % Specifies the document class

                             % The preamble begins here.
\title{An Example Document}  % Declares the document's title.
\author{Leslie Lamport}      % Declares the author's name.
\date{January 21, 1994}      % Deleting this command produces today's date.

\newcommand{\ip}[2]{(#1, #2)}
                             % Defines \ip{arg1}{arg2} to mean
                             % (arg1, arg2).

%\newcommand{\ip}[2]{\langle #1 | #2\rangle}
                             % This is an alternative definition of
                             % \ip that is commented out.

\begin{document}             % End of preamble and beginning of text.

\maketitle                   % Produces the title.

This is an example input file.  Comparing it with
the output it generates can show you how to
produce a simple document of your own.

\section{Ordinary Text}      % Produces section heading.  Lower-level
                             % sections are begun with similar
                             % \subsection and \subsubsection commands.

The ends  of words and sentences are marked
  by   spaces. It  doesn't matter how many
spaces    you type; one is as good as 100.  The
end of   a line counts as a space.

One   or more   blank lines denote the  end
of  a paragraph.

Since any number of consecutive spaces are treated
like a single one, the formatting of the input
file makes no difference to
      \LaTeX,                % The \LaTeX command generates the LaTeX logo.
but it makes a difference to you.  When you use
\LaTeX, making your input file as easy to read
as possible will be a great help as you write
your document and when you change it.  This sample
file shows how you can add comments to your own input
file.

Because printing is different from typewriting,
there are a number of things that you have to do
differently when preparing an input file than if
you were just typing the document directly.
Quotation marks like
       ``this''
have to be handled specially, as do quotes within
quotes:
       ``\,`this'            % \, separates the double and single quote.
        is what I just
        wrote, not  `that'\,''.

Dashes come in three sizes: an
       intra-word
dash, a medium dash for number ranges like
       1--2,
and a punctuation
       dash---like
this.

A sentence-ending space should be larger than the
space between words within a sentence.  You
sometimes have to type special commands in
conjunction with punctuation characters to get
this right, as in the following sentence.
       Gnats, gnus, etc.\ all  % `\ ' makes an inter-word space.
       begin with G\@.         % \@ marks end-of-sentence punctuation.
You should check the spaces after periods when
reading your output to make sure you haven't
forgotten any special cases.  Generating an
ellipsis
       \ldots\               % `\ ' is needed after `\ldots' because TeX
                             % ignores spaces after command names like \ldots
                             % made from \ + letters.
                             %
                             % Note how a `%' character causes TeX to ignore
                             % the end of the input line, so these blank lines
                             % do not start a new paragraph.
                             %
with the right spacing around the periods requires
a special command.

\LaTeX\ interprets some common characters as
commands, so you must type special commands to
generate them.  These characters include the
following:
       \$ \& \% \# \{ and \}.

In printing, text is usually emphasized with an
       \emph{italic}
type style.

\begin{em}
   A long segment of text can also be emphasized
   in this way.  Text within such a segment can be
   given \emph{additional} emphasis.
\end{em}

It is sometimes necessary to prevent \LaTeX\ from
breaking a line where it might otherwise do so.
This may be at a space, as between the ``Mr.''\ and
``Jones'' in
       ``Mr.~Jones'',        % ~ produces an unbreakable interword space.
or within a word---especially when the word is a
symbol like
       \mbox{\emph{itemnum}}
that makes little sense when hyphenated across
lines.

Footnotes\footnote{This is an example of a footnote.}
pose no problem.

\LaTeX\ is good at typesetting mathematical formulas
like
       \( x-3y + z = 7 \)
or
       \( a_{1} > x^{2n} + y^{2n} > x' \)
or
       \( \ip{A}{B} = \sum_{i} a_{i} b_{i} \).
The spaces you type in a formula are
ignored.  Remember that a letter like
       $x$                   % $ ... $  and  \( ... \)  are equivalent
is a formula when it denotes a mathematical
symbol, and it should be typed as one.

\section{Displayed Text}

Text is displayed by indenting it from the left
margin.  Quotations are commonly displayed.  There
are short quotations
\begin{quote}
   This is a short quotation.  It consists of a
   single paragraph of text.  See how it is formatted.
\end{quote}
and longer ones.
\begin{quotation}
   This is a longer quotation.  It consists of two
   paragraphs of text, neither of which are
   particularly interesting.

   This is the second paragraph of the quotation.  It
   is just as dull as the first paragraph.
\end{quotation}
Another frequently-displayed structure is a list.
The following is an example of an \emph{itemized}
list.
\begin{itemize}
   \item This is the first item of an itemized list.
         Each item in the list is marked with a ``tick''.
         You don't have to worry about what kind of tick
         mark is used.

   \item This is the second item of the list.  It
         contains another list nested inside it.  The inner
         list is an \emph{enumerated} list.
         \begin{enumerate}
            \item This is the first item of an enumerated
                  list that is nested within the itemized list.

            \item This is the second item of the inner list.
                  \LaTeX\ allows you to nest lists deeper than
                  you really should.
         \end{enumerate}
         This is the rest of the second item of the outer
         list.  It is no more interesting than any other
         part of the item.
   \item This is the third item of the list.
\end{itemize}
You can even display poetry.
\begin{verse}
   There is an environment
    for verse \\             % The \\ command separates lines
   Whose features some poets % within a stanza.
   will curse.

                             % One or more blank lines separate stanzas.

   For instead of making\\
   Them do \emph{all} line breaking, \\
   It allows them to put too many words on a line when they'd rather be
   forced to be terse.
\end{verse}

Mathematical formulas may also be displayed.  A
displayed formula
is
one-line long; multiline
formulas require special formatting instructions.
   \[  \ip{\Gamma}{\psi'} = x'' + y^{2} + z_{i}^{n}\]
Don't start a paragraph with a displayed equation,
nor make one a paragraph by itself.

\end{document}               % End of document.

\section{Bench Huge Section 214}
% synthetic bench section 214
% This is a sample LaTeX input file.  (Version of 12 August 2004.)
%
% A '%' character causes TeX to ignore all remaining text on the line,
% and is used for comments like this one.

\documentclass{article}      % Specifies the document class

                             % The preamble begins here.
\title{An Example Document}  % Declares the document's title.
\author{Leslie Lamport}      % Declares the author's name.
\date{January 21, 1994}      % Deleting this command produces today's date.

\newcommand{\ip}[2]{(#1, #2)}
                             % Defines \ip{arg1}{arg2} to mean
                             % (arg1, arg2).

%\newcommand{\ip}[2]{\langle #1 | #2\rangle}
                             % This is an alternative definition of
                             % \ip that is commented out.

\begin{document}             % End of preamble and beginning of text.

\maketitle                   % Produces the title.

This is an example input file.  Comparing it with
the output it generates can show you how to
produce a simple document of your own.

\section{Ordinary Text}      % Produces section heading.  Lower-level
                             % sections are begun with similar
                             % \subsection and \subsubsection commands.

The ends  of words and sentences are marked
  by   spaces. It  doesn't matter how many
spaces    you type; one is as good as 100.  The
end of   a line counts as a space.

One   or more   blank lines denote the  end
of  a paragraph.

Since any number of consecutive spaces are treated
like a single one, the formatting of the input
file makes no difference to
      \LaTeX,                % The \LaTeX command generates the LaTeX logo.
but it makes a difference to you.  When you use
\LaTeX, making your input file as easy to read
as possible will be a great help as you write
your document and when you change it.  This sample
file shows how you can add comments to your own input
file.

Because printing is different from typewriting,
there are a number of things that you have to do
differently when preparing an input file than if
you were just typing the document directly.
Quotation marks like
       ``this''
have to be handled specially, as do quotes within
quotes:
       ``\,`this'            % \, separates the double and single quote.
        is what I just
        wrote, not  `that'\,''.

Dashes come in three sizes: an
       intra-word
dash, a medium dash for number ranges like
       1--2,
and a punctuation
       dash---like
this.

A sentence-ending space should be larger than the
space between words within a sentence.  You
sometimes have to type special commands in
conjunction with punctuation characters to get
this right, as in the following sentence.
       Gnats, gnus, etc.\ all  % `\ ' makes an inter-word space.
       begin with G\@.         % \@ marks end-of-sentence punctuation.
You should check the spaces after periods when
reading your output to make sure you haven't
forgotten any special cases.  Generating an
ellipsis
       \ldots\               % `\ ' is needed after `\ldots' because TeX
                             % ignores spaces after command names like \ldots
                             % made from \ + letters.
                             %
                             % Note how a `%' character causes TeX to ignore
                             % the end of the input line, so these blank lines
                             % do not start a new paragraph.
                             %
with the right spacing around the periods requires
a special command.

\LaTeX\ interprets some common characters as
commands, so you must type special commands to
generate them.  These characters include the
following:
       \$ \& \% \# \{ and \}.

In printing, text is usually emphasized with an
       \emph{italic}
type style.

\begin{em}
   A long segment of text can also be emphasized
   in this way.  Text within such a segment can be
   given \emph{additional} emphasis.
\end{em}

It is sometimes necessary to prevent \LaTeX\ from
breaking a line where it might otherwise do so.
This may be at a space, as between the ``Mr.''\ and
``Jones'' in
       ``Mr.~Jones'',        % ~ produces an unbreakable interword space.
or within a word---especially when the word is a
symbol like
       \mbox{\emph{itemnum}}
that makes little sense when hyphenated across
lines.

Footnotes\footnote{This is an example of a footnote.}
pose no problem.

\LaTeX\ is good at typesetting mathematical formulas
like
       \( x-3y + z = 7 \)
or
       \( a_{1} > x^{2n} + y^{2n} > x' \)
or
       \( \ip{A}{B} = \sum_{i} a_{i} b_{i} \).
The spaces you type in a formula are
ignored.  Remember that a letter like
       $x$                   % $ ... $  and  \( ... \)  are equivalent
is a formula when it denotes a mathematical
symbol, and it should be typed as one.

\section{Displayed Text}

Text is displayed by indenting it from the left
margin.  Quotations are commonly displayed.  There
are short quotations
\begin{quote}
   This is a short quotation.  It consists of a
   single paragraph of text.  See how it is formatted.
\end{quote}
and longer ones.
\begin{quotation}
   This is a longer quotation.  It consists of two
   paragraphs of text, neither of which are
   particularly interesting.

   This is the second paragraph of the quotation.  It
   is just as dull as the first paragraph.
\end{quotation}
Another frequently-displayed structure is a list.
The following is an example of an \emph{itemized}
list.
\begin{itemize}
   \item This is the first item of an itemized list.
         Each item in the list is marked with a ``tick''.
         You don't have to worry about what kind of tick
         mark is used.

   \item This is the second item of the list.  It
         contains another list nested inside it.  The inner
         list is an \emph{enumerated} list.
         \begin{enumerate}
            \item This is the first item of an enumerated
                  list that is nested within the itemized list.

            \item This is the second item of the inner list.
                  \LaTeX\ allows you to nest lists deeper than
                  you really should.
         \end{enumerate}
         This is the rest of the second item of the outer
         list.  It is no more interesting than any other
         part of the item.
   \item This is the third item of the list.
\end{itemize}
You can even display poetry.
\begin{verse}
   There is an environment
    for verse \\             % The \\ command separates lines
   Whose features some poets % within a stanza.
   will curse.

                             % One or more blank lines separate stanzas.

   For instead of making\\
   Them do \emph{all} line breaking, \\
   It allows them to put too many words on a line when they'd rather be
   forced to be terse.
\end{verse}

Mathematical formulas may also be displayed.  A
displayed formula
is
one-line long; multiline
formulas require special formatting instructions.
   \[  \ip{\Gamma}{\psi'} = x'' + y^{2} + z_{i}^{n}\]
Don't start a paragraph with a displayed equation,
nor make one a paragraph by itself.

\end{document}               % End of document.

\section{Bench Huge Section 215}
% synthetic bench section 215
% This is a sample LaTeX input file.  (Version of 12 August 2004.)
%
% A '%' character causes TeX to ignore all remaining text on the line,
% and is used for comments like this one.

\documentclass{article}      % Specifies the document class

                             % The preamble begins here.
\title{An Example Document}  % Declares the document's title.
\author{Leslie Lamport}      % Declares the author's name.
\date{January 21, 1994}      % Deleting this command produces today's date.

\newcommand{\ip}[2]{(#1, #2)}
                             % Defines \ip{arg1}{arg2} to mean
                             % (arg1, arg2).

%\newcommand{\ip}[2]{\langle #1 | #2\rangle}
                             % This is an alternative definition of
                             % \ip that is commented out.

\begin{document}             % End of preamble and beginning of text.

\maketitle                   % Produces the title.

This is an example input file.  Comparing it with
the output it generates can show you how to
produce a simple document of your own.

\section{Ordinary Text}      % Produces section heading.  Lower-level
                             % sections are begun with similar
                             % \subsection and \subsubsection commands.

The ends  of words and sentences are marked
  by   spaces. It  doesn't matter how many
spaces    you type; one is as good as 100.  The
end of   a line counts as a space.

One   or more   blank lines denote the  end
of  a paragraph.

Since any number of consecutive spaces are treated
like a single one, the formatting of the input
file makes no difference to
      \LaTeX,                % The \LaTeX command generates the LaTeX logo.
but it makes a difference to you.  When you use
\LaTeX, making your input file as easy to read
as possible will be a great help as you write
your document and when you change it.  This sample
file shows how you can add comments to your own input
file.

Because printing is different from typewriting,
there are a number of things that you have to do
differently when preparing an input file than if
you were just typing the document directly.
Quotation marks like
       ``this''
have to be handled specially, as do quotes within
quotes:
       ``\,`this'            % \, separates the double and single quote.
        is what I just
        wrote, not  `that'\,''.

Dashes come in three sizes: an
       intra-word
dash, a medium dash for number ranges like
       1--2,
and a punctuation
       dash---like
this.

A sentence-ending space should be larger than the
space between words within a sentence.  You
sometimes have to type special commands in
conjunction with punctuation characters to get
this right, as in the following sentence.
       Gnats, gnus, etc.\ all  % `\ ' makes an inter-word space.
       begin with G\@.         % \@ marks end-of-sentence punctuation.
You should check the spaces after periods when
reading your output to make sure you haven't
forgotten any special cases.  Generating an
ellipsis
       \ldots\               % `\ ' is needed after `\ldots' because TeX
                             % ignores spaces after command names like \ldots
                             % made from \ + letters.
                             %
                             % Note how a `%' character causes TeX to ignore
                             % the end of the input line, so these blank lines
                             % do not start a new paragraph.
                             %
with the right spacing around the periods requires
a special command.

\LaTeX\ interprets some common characters as
commands, so you must type special commands to
generate them.  These characters include the
following:
       \$ \& \% \# \{ and \}.

In printing, text is usually emphasized with an
       \emph{italic}
type style.

\begin{em}
   A long segment of text can also be emphasized
   in this way.  Text within such a segment can be
   given \emph{additional} emphasis.
\end{em}

It is sometimes necessary to prevent \LaTeX\ from
breaking a line where it might otherwise do so.
This may be at a space, as between the ``Mr.''\ and
``Jones'' in
       ``Mr.~Jones'',        % ~ produces an unbreakable interword space.
or within a word---especially when the word is a
symbol like
       \mbox{\emph{itemnum}}
that makes little sense when hyphenated across
lines.

Footnotes\footnote{This is an example of a footnote.}
pose no problem.

\LaTeX\ is good at typesetting mathematical formulas
like
       \( x-3y + z = 7 \)
or
       \( a_{1} > x^{2n} + y^{2n} > x' \)
or
       \( \ip{A}{B} = \sum_{i} a_{i} b_{i} \).
The spaces you type in a formula are
ignored.  Remember that a letter like
       $x$                   % $ ... $  and  \( ... \)  are equivalent
is a formula when it denotes a mathematical
symbol, and it should be typed as one.

\section{Displayed Text}

Text is displayed by indenting it from the left
margin.  Quotations are commonly displayed.  There
are short quotations
\begin{quote}
   This is a short quotation.  It consists of a
   single paragraph of text.  See how it is formatted.
\end{quote}
and longer ones.
\begin{quotation}
   This is a longer quotation.  It consists of two
   paragraphs of text, neither of which are
   particularly interesting.

   This is the second paragraph of the quotation.  It
   is just as dull as the first paragraph.
\end{quotation}
Another frequently-displayed structure is a list.
The following is an example of an \emph{itemized}
list.
\begin{itemize}
   \item This is the first item of an itemized list.
         Each item in the list is marked with a ``tick''.
         You don't have to worry about what kind of tick
         mark is used.

   \item This is the second item of the list.  It
         contains another list nested inside it.  The inner
         list is an \emph{enumerated} list.
         \begin{enumerate}
            \item This is the first item of an enumerated
                  list that is nested within the itemized list.

            \item This is the second item of the inner list.
                  \LaTeX\ allows you to nest lists deeper than
                  you really should.
         \end{enumerate}
         This is the rest of the second item of the outer
         list.  It is no more interesting than any other
         part of the item.
   \item This is the third item of the list.
\end{itemize}
You can even display poetry.
\begin{verse}
   There is an environment
    for verse \\             % The \\ command separates lines
   Whose features some poets % within a stanza.
   will curse.

                             % One or more blank lines separate stanzas.

   For instead of making\\
   Them do \emph{all} line breaking, \\
   It allows them to put too many words on a line when they'd rather be
   forced to be terse.
\end{verse}

Mathematical formulas may also be displayed.  A
displayed formula
is
one-line long; multiline
formulas require special formatting instructions.
   \[  \ip{\Gamma}{\psi'} = x'' + y^{2} + z_{i}^{n}\]
Don't start a paragraph with a displayed equation,
nor make one a paragraph by itself.

\end{document}               % End of document.

\section{Bench Huge Section 216}
% synthetic bench section 216
% This is a sample LaTeX input file.  (Version of 12 August 2004.)
%
% A '%' character causes TeX to ignore all remaining text on the line,
% and is used for comments like this one.

\documentclass{article}      % Specifies the document class

                             % The preamble begins here.
\title{An Example Document}  % Declares the document's title.
\author{Leslie Lamport}      % Declares the author's name.
\date{January 21, 1994}      % Deleting this command produces today's date.

\newcommand{\ip}[2]{(#1, #2)}
                             % Defines \ip{arg1}{arg2} to mean
                             % (arg1, arg2).

%\newcommand{\ip}[2]{\langle #1 | #2\rangle}
                             % This is an alternative definition of
                             % \ip that is commented out.

\begin{document}             % End of preamble and beginning of text.

\maketitle                   % Produces the title.

This is an example input file.  Comparing it with
the output it generates can show you how to
produce a simple document of your own.

\section{Ordinary Text}      % Produces section heading.  Lower-level
                             % sections are begun with similar
                             % \subsection and \subsubsection commands.

The ends  of words and sentences are marked
  by   spaces. It  doesn't matter how many
spaces    you type; one is as good as 100.  The
end of   a line counts as a space.

One   or more   blank lines denote the  end
of  a paragraph.

Since any number of consecutive spaces are treated
like a single one, the formatting of the input
file makes no difference to
      \LaTeX,                % The \LaTeX command generates the LaTeX logo.
but it makes a difference to you.  When you use
\LaTeX, making your input file as easy to read
as possible will be a great help as you write
your document and when you change it.  This sample
file shows how you can add comments to your own input
file.

Because printing is different from typewriting,
there are a number of things that you have to do
differently when preparing an input file than if
you were just typing the document directly.
Quotation marks like
       ``this''
have to be handled specially, as do quotes within
quotes:
       ``\,`this'            % \, separates the double and single quote.
        is what I just
        wrote, not  `that'\,''.

Dashes come in three sizes: an
       intra-word
dash, a medium dash for number ranges like
       1--2,
and a punctuation
       dash---like
this.

A sentence-ending space should be larger than the
space between words within a sentence.  You
sometimes have to type special commands in
conjunction with punctuation characters to get
this right, as in the following sentence.
       Gnats, gnus, etc.\ all  % `\ ' makes an inter-word space.
       begin with G\@.         % \@ marks end-of-sentence punctuation.
You should check the spaces after periods when
reading your output to make sure you haven't
forgotten any special cases.  Generating an
ellipsis
       \ldots\               % `\ ' is needed after `\ldots' because TeX
                             % ignores spaces after command names like \ldots
                             % made from \ + letters.
                             %
                             % Note how a `%' character causes TeX to ignore
                             % the end of the input line, so these blank lines
                             % do not start a new paragraph.
                             %
with the right spacing around the periods requires
a special command.

\LaTeX\ interprets some common characters as
commands, so you must type special commands to
generate them.  These characters include the
following:
       \$ \& \% \# \{ and \}.

In printing, text is usually emphasized with an
       \emph{italic}
type style.

\begin{em}
   A long segment of text can also be emphasized
   in this way.  Text within such a segment can be
   given \emph{additional} emphasis.
\end{em}

It is sometimes necessary to prevent \LaTeX\ from
breaking a line where it might otherwise do so.
This may be at a space, as between the ``Mr.''\ and
``Jones'' in
       ``Mr.~Jones'',        % ~ produces an unbreakable interword space.
or within a word---especially when the word is a
symbol like
       \mbox{\emph{itemnum}}
that makes little sense when hyphenated across
lines.

Footnotes\footnote{This is an example of a footnote.}
pose no problem.

\LaTeX\ is good at typesetting mathematical formulas
like
       \( x-3y + z = 7 \)
or
       \( a_{1} > x^{2n} + y^{2n} > x' \)
or
       \( \ip{A}{B} = \sum_{i} a_{i} b_{i} \).
The spaces you type in a formula are
ignored.  Remember that a letter like
       $x$                   % $ ... $  and  \( ... \)  are equivalent
is a formula when it denotes a mathematical
symbol, and it should be typed as one.

\section{Displayed Text}

Text is displayed by indenting it from the left
margin.  Quotations are commonly displayed.  There
are short quotations
\begin{quote}
   This is a short quotation.  It consists of a
   single paragraph of text.  See how it is formatted.
\end{quote}
and longer ones.
\begin{quotation}
   This is a longer quotation.  It consists of two
   paragraphs of text, neither of which are
   particularly interesting.

   This is the second paragraph of the quotation.  It
   is just as dull as the first paragraph.
\end{quotation}
Another frequently-displayed structure is a list.
The following is an example of an \emph{itemized}
list.
\begin{itemize}
   \item This is the first item of an itemized list.
         Each item in the list is marked with a ``tick''.
         You don't have to worry about what kind of tick
         mark is used.

   \item This is the second item of the list.  It
         contains another list nested inside it.  The inner
         list is an \emph{enumerated} list.
         \begin{enumerate}
            \item This is the first item of an enumerated
                  list that is nested within the itemized list.

            \item This is the second item of the inner list.
                  \LaTeX\ allows you to nest lists deeper than
                  you really should.
         \end{enumerate}
         This is the rest of the second item of the outer
         list.  It is no more interesting than any other
         part of the item.
   \item This is the third item of the list.
\end{itemize}
You can even display poetry.
\begin{verse}
   There is an environment
    for verse \\             % The \\ command separates lines
   Whose features some poets % within a stanza.
   will curse.

                             % One or more blank lines separate stanzas.

   For instead of making\\
   Them do \emph{all} line breaking, \\
   It allows them to put too many words on a line when they'd rather be
   forced to be terse.
\end{verse}

Mathematical formulas may also be displayed.  A
displayed formula
is
one-line long; multiline
formulas require special formatting instructions.
   \[  \ip{\Gamma}{\psi'} = x'' + y^{2} + z_{i}^{n}\]
Don't start a paragraph with a displayed equation,
nor make one a paragraph by itself.

\end{document}               % End of document.

\section{Bench Huge Section 217}
% synthetic bench section 217
% This is a sample LaTeX input file.  (Version of 12 August 2004.)
%
% A '%' character causes TeX to ignore all remaining text on the line,
% and is used for comments like this one.

\documentclass{article}      % Specifies the document class

                             % The preamble begins here.
\title{An Example Document}  % Declares the document's title.
\author{Leslie Lamport}      % Declares the author's name.
\date{January 21, 1994}      % Deleting this command produces today's date.

\newcommand{\ip}[2]{(#1, #2)}
                             % Defines \ip{arg1}{arg2} to mean
                             % (arg1, arg2).

%\newcommand{\ip}[2]{\langle #1 | #2\rangle}
                             % This is an alternative definition of
                             % \ip that is commented out.

\begin{document}             % End of preamble and beginning of text.

\maketitle                   % Produces the title.

This is an example input file.  Comparing it with
the output it generates can show you how to
produce a simple document of your own.

\section{Ordinary Text}      % Produces section heading.  Lower-level
                             % sections are begun with similar
                             % \subsection and \subsubsection commands.

The ends  of words and sentences are marked
  by   spaces. It  doesn't matter how many
spaces    you type; one is as good as 100.  The
end of   a line counts as a space.

One   or more   blank lines denote the  end
of  a paragraph.

Since any number of consecutive spaces are treated
like a single one, the formatting of the input
file makes no difference to
      \LaTeX,                % The \LaTeX command generates the LaTeX logo.
but it makes a difference to you.  When you use
\LaTeX, making your input file as easy to read
as possible will be a great help as you write
your document and when you change it.  This sample
file shows how you can add comments to your own input
file.

Because printing is different from typewriting,
there are a number of things that you have to do
differently when preparing an input file than if
you were just typing the document directly.
Quotation marks like
       ``this''
have to be handled specially, as do quotes within
quotes:
       ``\,`this'            % \, separates the double and single quote.
        is what I just
        wrote, not  `that'\,''.

Dashes come in three sizes: an
       intra-word
dash, a medium dash for number ranges like
       1--2,
and a punctuation
       dash---like
this.

A sentence-ending space should be larger than the
space between words within a sentence.  You
sometimes have to type special commands in
conjunction with punctuation characters to get
this right, as in the following sentence.
       Gnats, gnus, etc.\ all  % `\ ' makes an inter-word space.
       begin with G\@.         % \@ marks end-of-sentence punctuation.
You should check the spaces after periods when
reading your output to make sure you haven't
forgotten any special cases.  Generating an
ellipsis
       \ldots\               % `\ ' is needed after `\ldots' because TeX
                             % ignores spaces after command names like \ldots
                             % made from \ + letters.
                             %
                             % Note how a `%' character causes TeX to ignore
                             % the end of the input line, so these blank lines
                             % do not start a new paragraph.
                             %
with the right spacing around the periods requires
a special command.

\LaTeX\ interprets some common characters as
commands, so you must type special commands to
generate them.  These characters include the
following:
       \$ \& \% \# \{ and \}.

In printing, text is usually emphasized with an
       \emph{italic}
type style.

\begin{em}
   A long segment of text can also be emphasized
   in this way.  Text within such a segment can be
   given \emph{additional} emphasis.
\end{em}

It is sometimes necessary to prevent \LaTeX\ from
breaking a line where it might otherwise do so.
This may be at a space, as between the ``Mr.''\ and
``Jones'' in
       ``Mr.~Jones'',        % ~ produces an unbreakable interword space.
or within a word---especially when the word is a
symbol like
       \mbox{\emph{itemnum}}
that makes little sense when hyphenated across
lines.

Footnotes\footnote{This is an example of a footnote.}
pose no problem.

\LaTeX\ is good at typesetting mathematical formulas
like
       \( x-3y + z = 7 \)
or
       \( a_{1} > x^{2n} + y^{2n} > x' \)
or
       \( \ip{A}{B} = \sum_{i} a_{i} b_{i} \).
The spaces you type in a formula are
ignored.  Remember that a letter like
       $x$                   % $ ... $  and  \( ... \)  are equivalent
is a formula when it denotes a mathematical
symbol, and it should be typed as one.

\section{Displayed Text}

Text is displayed by indenting it from the left
margin.  Quotations are commonly displayed.  There
are short quotations
\begin{quote}
   This is a short quotation.  It consists of a
   single paragraph of text.  See how it is formatted.
\end{quote}
and longer ones.
\begin{quotation}
   This is a longer quotation.  It consists of two
   paragraphs of text, neither of which are
   particularly interesting.

   This is the second paragraph of the quotation.  It
   is just as dull as the first paragraph.
\end{quotation}
Another frequently-displayed structure is a list.
The following is an example of an \emph{itemized}
list.
\begin{itemize}
   \item This is the first item of an itemized list.
         Each item in the list is marked with a ``tick''.
         You don't have to worry about what kind of tick
         mark is used.

   \item This is the second item of the list.  It
         contains another list nested inside it.  The inner
         list is an \emph{enumerated} list.
         \begin{enumerate}
            \item This is the first item of an enumerated
                  list that is nested within the itemized list.

            \item This is the second item of the inner list.
                  \LaTeX\ allows you to nest lists deeper than
                  you really should.
         \end{enumerate}
         This is the rest of the second item of the outer
         list.  It is no more interesting than any other
         part of the item.
   \item This is the third item of the list.
\end{itemize}
You can even display poetry.
\begin{verse}
   There is an environment
    for verse \\             % The \\ command separates lines
   Whose features some poets % within a stanza.
   will curse.

                             % One or more blank lines separate stanzas.

   For instead of making\\
   Them do \emph{all} line breaking, \\
   It allows them to put too many words on a line when they'd rather be
   forced to be terse.
\end{verse}

Mathematical formulas may also be displayed.  A
displayed formula
is
one-line long; multiline
formulas require special formatting instructions.
   \[  \ip{\Gamma}{\psi'} = x'' + y^{2} + z_{i}^{n}\]
Don't start a paragraph with a displayed equation,
nor make one a paragraph by itself.

\end{document}               % End of document.

\section{Bench Huge Section 218}
% synthetic bench section 218
% This is a sample LaTeX input file.  (Version of 12 August 2004.)
%
% A '%' character causes TeX to ignore all remaining text on the line,
% and is used for comments like this one.

\documentclass{article}      % Specifies the document class

                             % The preamble begins here.
\title{An Example Document}  % Declares the document's title.
\author{Leslie Lamport}      % Declares the author's name.
\date{January 21, 1994}      % Deleting this command produces today's date.

\newcommand{\ip}[2]{(#1, #2)}
                             % Defines \ip{arg1}{arg2} to mean
                             % (arg1, arg2).

%\newcommand{\ip}[2]{\langle #1 | #2\rangle}
                             % This is an alternative definition of
                             % \ip that is commented out.

\begin{document}             % End of preamble and beginning of text.

\maketitle                   % Produces the title.

This is an example input file.  Comparing it with
the output it generates can show you how to
produce a simple document of your own.

\section{Ordinary Text}      % Produces section heading.  Lower-level
                             % sections are begun with similar
                             % \subsection and \subsubsection commands.

The ends  of words and sentences are marked
  by   spaces. It  doesn't matter how many
spaces    you type; one is as good as 100.  The
end of   a line counts as a space.

One   or more   blank lines denote the  end
of  a paragraph.

Since any number of consecutive spaces are treated
like a single one, the formatting of the input
file makes no difference to
      \LaTeX,                % The \LaTeX command generates the LaTeX logo.
but it makes a difference to you.  When you use
\LaTeX, making your input file as easy to read
as possible will be a great help as you write
your document and when you change it.  This sample
file shows how you can add comments to your own input
file.

Because printing is different from typewriting,
there are a number of things that you have to do
differently when preparing an input file than if
you were just typing the document directly.
Quotation marks like
       ``this''
have to be handled specially, as do quotes within
quotes:
       ``\,`this'            % \, separates the double and single quote.
        is what I just
        wrote, not  `that'\,''.

Dashes come in three sizes: an
       intra-word
dash, a medium dash for number ranges like
       1--2,
and a punctuation
       dash---like
this.

A sentence-ending space should be larger than the
space between words within a sentence.  You
sometimes have to type special commands in
conjunction with punctuation characters to get
this right, as in the following sentence.
       Gnats, gnus, etc.\ all  % `\ ' makes an inter-word space.
       begin with G\@.         % \@ marks end-of-sentence punctuation.
You should check the spaces after periods when
reading your output to make sure you haven't
forgotten any special cases.  Generating an
ellipsis
       \ldots\               % `\ ' is needed after `\ldots' because TeX
                             % ignores spaces after command names like \ldots
                             % made from \ + letters.
                             %
                             % Note how a `%' character causes TeX to ignore
                             % the end of the input line, so these blank lines
                             % do not start a new paragraph.
                             %
with the right spacing around the periods requires
a special command.

\LaTeX\ interprets some common characters as
commands, so you must type special commands to
generate them.  These characters include the
following:
       \$ \& \% \# \{ and \}.

In printing, text is usually emphasized with an
       \emph{italic}
type style.

\begin{em}
   A long segment of text can also be emphasized
   in this way.  Text within such a segment can be
   given \emph{additional} emphasis.
\end{em}

It is sometimes necessary to prevent \LaTeX\ from
breaking a line where it might otherwise do so.
This may be at a space, as between the ``Mr.''\ and
``Jones'' in
       ``Mr.~Jones'',        % ~ produces an unbreakable interword space.
or within a word---especially when the word is a
symbol like
       \mbox{\emph{itemnum}}
that makes little sense when hyphenated across
lines.

Footnotes\footnote{This is an example of a footnote.}
pose no problem.

\LaTeX\ is good at typesetting mathematical formulas
like
       \( x-3y + z = 7 \)
or
       \( a_{1} > x^{2n} + y^{2n} > x' \)
or
       \( \ip{A}{B} = \sum_{i} a_{i} b_{i} \).
The spaces you type in a formula are
ignored.  Remember that a letter like
       $x$                   % $ ... $  and  \( ... \)  are equivalent
is a formula when it denotes a mathematical
symbol, and it should be typed as one.

\section{Displayed Text}

Text is displayed by indenting it from the left
margin.  Quotations are commonly displayed.  There
are short quotations
\begin{quote}
   This is a short quotation.  It consists of a
   single paragraph of text.  See how it is formatted.
\end{quote}
and longer ones.
\begin{quotation}
   This is a longer quotation.  It consists of two
   paragraphs of text, neither of which are
   particularly interesting.

   This is the second paragraph of the quotation.  It
   is just as dull as the first paragraph.
\end{quotation}
Another frequently-displayed structure is a list.
The following is an example of an \emph{itemized}
list.
\begin{itemize}
   \item This is the first item of an itemized list.
         Each item in the list is marked with a ``tick''.
         You don't have to worry about what kind of tick
         mark is used.

   \item This is the second item of the list.  It
         contains another list nested inside it.  The inner
         list is an \emph{enumerated} list.
         \begin{enumerate}
            \item This is the first item of an enumerated
                  list that is nested within the itemized list.

            \item This is the second item of the inner list.
                  \LaTeX\ allows you to nest lists deeper than
                  you really should.
         \end{enumerate}
         This is the rest of the second item of the outer
         list.  It is no more interesting than any other
         part of the item.
   \item This is the third item of the list.
\end{itemize}
You can even display poetry.
\begin{verse}
   There is an environment
    for verse \\             % The \\ command separates lines
   Whose features some poets % within a stanza.
   will curse.

                             % One or more blank lines separate stanzas.

   For instead of making\\
   Them do \emph{all} line breaking, \\
   It allows them to put too many words on a line when they'd rather be
   forced to be terse.
\end{verse}

Mathematical formulas may also be displayed.  A
displayed formula
is
one-line long; multiline
formulas require special formatting instructions.
   \[  \ip{\Gamma}{\psi'} = x'' + y^{2} + z_{i}^{n}\]
Don't start a paragraph with a displayed equation,
nor make one a paragraph by itself.

\end{document}               % End of document.

\section{Bench Huge Section 219}
% synthetic bench section 219
% This is a sample LaTeX input file.  (Version of 12 August 2004.)
%
% A '%' character causes TeX to ignore all remaining text on the line,
% and is used for comments like this one.

\documentclass{article}      % Specifies the document class

                             % The preamble begins here.
\title{An Example Document}  % Declares the document's title.
\author{Leslie Lamport}      % Declares the author's name.
\date{January 21, 1994}      % Deleting this command produces today's date.

\newcommand{\ip}[2]{(#1, #2)}
                             % Defines \ip{arg1}{arg2} to mean
                             % (arg1, arg2).

%\newcommand{\ip}[2]{\langle #1 | #2\rangle}
                             % This is an alternative definition of
                             % \ip that is commented out.

\begin{document}             % End of preamble and beginning of text.

\maketitle                   % Produces the title.

This is an example input file.  Comparing it with
the output it generates can show you how to
produce a simple document of your own.

\section{Ordinary Text}      % Produces section heading.  Lower-level
                             % sections are begun with similar
                             % \subsection and \subsubsection commands.

The ends  of words and sentences are marked
  by   spaces. It  doesn't matter how many
spaces    you type; one is as good as 100.  The
end of   a line counts as a space.

One   or more   blank lines denote the  end
of  a paragraph.

Since any number of consecutive spaces are treated
like a single one, the formatting of the input
file makes no difference to
      \LaTeX,                % The \LaTeX command generates the LaTeX logo.
but it makes a difference to you.  When you use
\LaTeX, making your input file as easy to read
as possible will be a great help as you write
your document and when you change it.  This sample
file shows how you can add comments to your own input
file.

Because printing is different from typewriting,
there are a number of things that you have to do
differently when preparing an input file than if
you were just typing the document directly.
Quotation marks like
       ``this''
have to be handled specially, as do quotes within
quotes:
       ``\,`this'            % \, separates the double and single quote.
        is what I just
        wrote, not  `that'\,''.

Dashes come in three sizes: an
       intra-word
dash, a medium dash for number ranges like
       1--2,
and a punctuation
       dash---like
this.

A sentence-ending space should be larger than the
space between words within a sentence.  You
sometimes have to type special commands in
conjunction with punctuation characters to get
this right, as in the following sentence.
       Gnats, gnus, etc.\ all  % `\ ' makes an inter-word space.
       begin with G\@.         % \@ marks end-of-sentence punctuation.
You should check the spaces after periods when
reading your output to make sure you haven't
forgotten any special cases.  Generating an
ellipsis
       \ldots\               % `\ ' is needed after `\ldots' because TeX
                             % ignores spaces after command names like \ldots
                             % made from \ + letters.
                             %
                             % Note how a `%' character causes TeX to ignore
                             % the end of the input line, so these blank lines
                             % do not start a new paragraph.
                             %
with the right spacing around the periods requires
a special command.

\LaTeX\ interprets some common characters as
commands, so you must type special commands to
generate them.  These characters include the
following:
       \$ \& \% \# \{ and \}.

In printing, text is usually emphasized with an
       \emph{italic}
type style.

\begin{em}
   A long segment of text can also be emphasized
   in this way.  Text within such a segment can be
   given \emph{additional} emphasis.
\end{em}

It is sometimes necessary to prevent \LaTeX\ from
breaking a line where it might otherwise do so.
This may be at a space, as between the ``Mr.''\ and
``Jones'' in
       ``Mr.~Jones'',        % ~ produces an unbreakable interword space.
or within a word---especially when the word is a
symbol like
       \mbox{\emph{itemnum}}
that makes little sense when hyphenated across
lines.

Footnotes\footnote{This is an example of a footnote.}
pose no problem.

\LaTeX\ is good at typesetting mathematical formulas
like
       \( x-3y + z = 7 \)
or
       \( a_{1} > x^{2n} + y^{2n} > x' \)
or
       \( \ip{A}{B} = \sum_{i} a_{i} b_{i} \).
The spaces you type in a formula are
ignored.  Remember that a letter like
       $x$                   % $ ... $  and  \( ... \)  are equivalent
is a formula when it denotes a mathematical
symbol, and it should be typed as one.

\section{Displayed Text}

Text is displayed by indenting it from the left
margin.  Quotations are commonly displayed.  There
are short quotations
\begin{quote}
   This is a short quotation.  It consists of a
   single paragraph of text.  See how it is formatted.
\end{quote}
and longer ones.
\begin{quotation}
   This is a longer quotation.  It consists of two
   paragraphs of text, neither of which are
   particularly interesting.

   This is the second paragraph of the quotation.  It
   is just as dull as the first paragraph.
\end{quotation}
Another frequently-displayed structure is a list.
The following is an example of an \emph{itemized}
list.
\begin{itemize}
   \item This is the first item of an itemized list.
         Each item in the list is marked with a ``tick''.
         You don't have to worry about what kind of tick
         mark is used.

   \item This is the second item of the list.  It
         contains another list nested inside it.  The inner
         list is an \emph{enumerated} list.
         \begin{enumerate}
            \item This is the first item of an enumerated
                  list that is nested within the itemized list.

            \item This is the second item of the inner list.
                  \LaTeX\ allows you to nest lists deeper than
                  you really should.
         \end{enumerate}
         This is the rest of the second item of the outer
         list.  It is no more interesting than any other
         part of the item.
   \item This is the third item of the list.
\end{itemize}
You can even display poetry.
\begin{verse}
   There is an environment
    for verse \\             % The \\ command separates lines
   Whose features some poets % within a stanza.
   will curse.

                             % One or more blank lines separate stanzas.

   For instead of making\\
   Them do \emph{all} line breaking, \\
   It allows them to put too many words on a line when they'd rather be
   forced to be terse.
\end{verse}

Mathematical formulas may also be displayed.  A
displayed formula
is
one-line long; multiline
formulas require special formatting instructions.
   \[  \ip{\Gamma}{\psi'} = x'' + y^{2} + z_{i}^{n}\]
Don't start a paragraph with a displayed equation,
nor make one a paragraph by itself.

\end{document}               % End of document.

\section{Bench Huge Section 220}
% synthetic bench section 220
% This is a sample LaTeX input file.  (Version of 12 August 2004.)
%
% A '%' character causes TeX to ignore all remaining text on the line,
% and is used for comments like this one.

\documentclass{article}      % Specifies the document class

                             % The preamble begins here.
\title{An Example Document}  % Declares the document's title.
\author{Leslie Lamport}      % Declares the author's name.
\date{January 21, 1994}      % Deleting this command produces today's date.

\newcommand{\ip}[2]{(#1, #2)}
                             % Defines \ip{arg1}{arg2} to mean
                             % (arg1, arg2).

%\newcommand{\ip}[2]{\langle #1 | #2\rangle}
                             % This is an alternative definition of
                             % \ip that is commented out.

\begin{document}             % End of preamble and beginning of text.

\maketitle                   % Produces the title.

This is an example input file.  Comparing it with
the output it generates can show you how to
produce a simple document of your own.

\section{Ordinary Text}      % Produces section heading.  Lower-level
                             % sections are begun with similar
                             % \subsection and \subsubsection commands.

The ends  of words and sentences are marked
  by   spaces. It  doesn't matter how many
spaces    you type; one is as good as 100.  The
end of   a line counts as a space.

One   or more   blank lines denote the  end
of  a paragraph.

Since any number of consecutive spaces are treated
like a single one, the formatting of the input
file makes no difference to
      \LaTeX,                % The \LaTeX command generates the LaTeX logo.
but it makes a difference to you.  When you use
\LaTeX, making your input file as easy to read
as possible will be a great help as you write
your document and when you change it.  This sample
file shows how you can add comments to your own input
file.

Because printing is different from typewriting,
there are a number of things that you have to do
differently when preparing an input file than if
you were just typing the document directly.
Quotation marks like
       ``this''
have to be handled specially, as do quotes within
quotes:
       ``\,`this'            % \, separates the double and single quote.
        is what I just
        wrote, not  `that'\,''.

Dashes come in three sizes: an
       intra-word
dash, a medium dash for number ranges like
       1--2,
and a punctuation
       dash---like
this.

A sentence-ending space should be larger than the
space between words within a sentence.  You
sometimes have to type special commands in
conjunction with punctuation characters to get
this right, as in the following sentence.
       Gnats, gnus, etc.\ all  % `\ ' makes an inter-word space.
       begin with G\@.         % \@ marks end-of-sentence punctuation.
You should check the spaces after periods when
reading your output to make sure you haven't
forgotten any special cases.  Generating an
ellipsis
       \ldots\               % `\ ' is needed after `\ldots' because TeX
                             % ignores spaces after command names like \ldots
                             % made from \ + letters.
                             %
                             % Note how a `%' character causes TeX to ignore
                             % the end of the input line, so these blank lines
                             % do not start a new paragraph.
                             %
with the right spacing around the periods requires
a special command.

\LaTeX\ interprets some common characters as
commands, so you must type special commands to
generate them.  These characters include the
following:
       \$ \& \% \# \{ and \}.

In printing, text is usually emphasized with an
       \emph{italic}
type style.

\begin{em}
   A long segment of text can also be emphasized
   in this way.  Text within such a segment can be
   given \emph{additional} emphasis.
\end{em}

It is sometimes necessary to prevent \LaTeX\ from
breaking a line where it might otherwise do so.
This may be at a space, as between the ``Mr.''\ and
``Jones'' in
       ``Mr.~Jones'',        % ~ produces an unbreakable interword space.
or within a word---especially when the word is a
symbol like
       \mbox{\emph{itemnum}}
that makes little sense when hyphenated across
lines.

Footnotes\footnote{This is an example of a footnote.}
pose no problem.

\LaTeX\ is good at typesetting mathematical formulas
like
       \( x-3y + z = 7 \)
or
       \( a_{1} > x^{2n} + y^{2n} > x' \)
or
       \( \ip{A}{B} = \sum_{i} a_{i} b_{i} \).
The spaces you type in a formula are
ignored.  Remember that a letter like
       $x$                   % $ ... $  and  \( ... \)  are equivalent
is a formula when it denotes a mathematical
symbol, and it should be typed as one.

\section{Displayed Text}

Text is displayed by indenting it from the left
margin.  Quotations are commonly displayed.  There
are short quotations
\begin{quote}
   This is a short quotation.  It consists of a
   single paragraph of text.  See how it is formatted.
\end{quote}
and longer ones.
\begin{quotation}
   This is a longer quotation.  It consists of two
   paragraphs of text, neither of which are
   particularly interesting.

   This is the second paragraph of the quotation.  It
   is just as dull as the first paragraph.
\end{quotation}
Another frequently-displayed structure is a list.
The following is an example of an \emph{itemized}
list.
\begin{itemize}
   \item This is the first item of an itemized list.
         Each item in the list is marked with a ``tick''.
         You don't have to worry about what kind of tick
         mark is used.

   \item This is the second item of the list.  It
         contains another list nested inside it.  The inner
         list is an \emph{enumerated} list.
         \begin{enumerate}
            \item This is the first item of an enumerated
                  list that is nested within the itemized list.

            \item This is the second item of the inner list.
                  \LaTeX\ allows you to nest lists deeper than
                  you really should.
         \end{enumerate}
         This is the rest of the second item of the outer
         list.  It is no more interesting than any other
         part of the item.
   \item This is the third item of the list.
\end{itemize}
You can even display poetry.
\begin{verse}
   There is an environment
    for verse \\             % The \\ command separates lines
   Whose features some poets % within a stanza.
   will curse.

                             % One or more blank lines separate stanzas.

   For instead of making\\
   Them do \emph{all} line breaking, \\
   It allows them to put too many words on a line when they'd rather be
   forced to be terse.
\end{verse}

Mathematical formulas may also be displayed.  A
displayed formula
is
one-line long; multiline
formulas require special formatting instructions.
   \[  \ip{\Gamma}{\psi'} = x'' + y^{2} + z_{i}^{n}\]
Don't start a paragraph with a displayed equation,
nor make one a paragraph by itself.

\end{document}               % End of document.

\section{Bench Huge Section 221}
% synthetic bench section 221
% This is a sample LaTeX input file.  (Version of 12 August 2004.)
%
% A '%' character causes TeX to ignore all remaining text on the line,
% and is used for comments like this one.

\documentclass{article}      % Specifies the document class

                             % The preamble begins here.
\title{An Example Document}  % Declares the document's title.
\author{Leslie Lamport}      % Declares the author's name.
\date{January 21, 1994}      % Deleting this command produces today's date.

\newcommand{\ip}[2]{(#1, #2)}
                             % Defines \ip{arg1}{arg2} to mean
                             % (arg1, arg2).

%\newcommand{\ip}[2]{\langle #1 | #2\rangle}
                             % This is an alternative definition of
                             % \ip that is commented out.

\begin{document}             % End of preamble and beginning of text.

\maketitle                   % Produces the title.

This is an example input file.  Comparing it with
the output it generates can show you how to
produce a simple document of your own.

\section{Ordinary Text}      % Produces section heading.  Lower-level
                             % sections are begun with similar
                             % \subsection and \subsubsection commands.

The ends  of words and sentences are marked
  by   spaces. It  doesn't matter how many
spaces    you type; one is as good as 100.  The
end of   a line counts as a space.

One   or more   blank lines denote the  end
of  a paragraph.

Since any number of consecutive spaces are treated
like a single one, the formatting of the input
file makes no difference to
      \LaTeX,                % The \LaTeX command generates the LaTeX logo.
but it makes a difference to you.  When you use
\LaTeX, making your input file as easy to read
as possible will be a great help as you write
your document and when you change it.  This sample
file shows how you can add comments to your own input
file.

Because printing is different from typewriting,
there are a number of things that you have to do
differently when preparing an input file than if
you were just typing the document directly.
Quotation marks like
       ``this''
have to be handled specially, as do quotes within
quotes:
       ``\,`this'            % \, separates the double and single quote.
        is what I just
        wrote, not  `that'\,''.

Dashes come in three sizes: an
       intra-word
dash, a medium dash for number ranges like
       1--2,
and a punctuation
       dash---like
this.

A sentence-ending space should be larger than the
space between words within a sentence.  You
sometimes have to type special commands in
conjunction with punctuation characters to get
this right, as in the following sentence.
       Gnats, gnus, etc.\ all  % `\ ' makes an inter-word space.
       begin with G\@.         % \@ marks end-of-sentence punctuation.
You should check the spaces after periods when
reading your output to make sure you haven't
forgotten any special cases.  Generating an
ellipsis
       \ldots\               % `\ ' is needed after `\ldots' because TeX
                             % ignores spaces after command names like \ldots
                             % made from \ + letters.
                             %
                             % Note how a `%' character causes TeX to ignore
                             % the end of the input line, so these blank lines
                             % do not start a new paragraph.
                             %
with the right spacing around the periods requires
a special command.

\LaTeX\ interprets some common characters as
commands, so you must type special commands to
generate them.  These characters include the
following:
       \$ \& \% \# \{ and \}.

In printing, text is usually emphasized with an
       \emph{italic}
type style.

\begin{em}
   A long segment of text can also be emphasized
   in this way.  Text within such a segment can be
   given \emph{additional} emphasis.
\end{em}

It is sometimes necessary to prevent \LaTeX\ from
breaking a line where it might otherwise do so.
This may be at a space, as between the ``Mr.''\ and
``Jones'' in
       ``Mr.~Jones'',        % ~ produces an unbreakable interword space.
or within a word---especially when the word is a
symbol like
       \mbox{\emph{itemnum}}
that makes little sense when hyphenated across
lines.

Footnotes\footnote{This is an example of a footnote.}
pose no problem.

\LaTeX\ is good at typesetting mathematical formulas
like
       \( x-3y + z = 7 \)
or
       \( a_{1} > x^{2n} + y^{2n} > x' \)
or
       \( \ip{A}{B} = \sum_{i} a_{i} b_{i} \).
The spaces you type in a formula are
ignored.  Remember that a letter like
       $x$                   % $ ... $  and  \( ... \)  are equivalent
is a formula when it denotes a mathematical
symbol, and it should be typed as one.

\section{Displayed Text}

Text is displayed by indenting it from the left
margin.  Quotations are commonly displayed.  There
are short quotations
\begin{quote}
   This is a short quotation.  It consists of a
   single paragraph of text.  See how it is formatted.
\end{quote}
and longer ones.
\begin{quotation}
   This is a longer quotation.  It consists of two
   paragraphs of text, neither of which are
   particularly interesting.

   This is the second paragraph of the quotation.  It
   is just as dull as the first paragraph.
\end{quotation}
Another frequently-displayed structure is a list.
The following is an example of an \emph{itemized}
list.
\begin{itemize}
   \item This is the first item of an itemized list.
         Each item in the list is marked with a ``tick''.
         You don't have to worry about what kind of tick
         mark is used.

   \item This is the second item of the list.  It
         contains another list nested inside it.  The inner
         list is an \emph{enumerated} list.
         \begin{enumerate}
            \item This is the first item of an enumerated
                  list that is nested within the itemized list.

            \item This is the second item of the inner list.
                  \LaTeX\ allows you to nest lists deeper than
                  you really should.
         \end{enumerate}
         This is the rest of the second item of the outer
         list.  It is no more interesting than any other
         part of the item.
   \item This is the third item of the list.
\end{itemize}
You can even display poetry.
\begin{verse}
   There is an environment
    for verse \\             % The \\ command separates lines
   Whose features some poets % within a stanza.
   will curse.

                             % One or more blank lines separate stanzas.

   For instead of making\\
   Them do \emph{all} line breaking, \\
   It allows them to put too many words on a line when they'd rather be
   forced to be terse.
\end{verse}

Mathematical formulas may also be displayed.  A
displayed formula
is
one-line long; multiline
formulas require special formatting instructions.
   \[  \ip{\Gamma}{\psi'} = x'' + y^{2} + z_{i}^{n}\]
Don't start a paragraph with a displayed equation,
nor make one a paragraph by itself.

\end{document}               % End of document.

\section{Bench Huge Section 222}
% synthetic bench section 222
% This is a sample LaTeX input file.  (Version of 12 August 2004.)
%
% A '%' character causes TeX to ignore all remaining text on the line,
% and is used for comments like this one.

\documentclass{article}      % Specifies the document class

                             % The preamble begins here.
\title{An Example Document}  % Declares the document's title.
\author{Leslie Lamport}      % Declares the author's name.
\date{January 21, 1994}      % Deleting this command produces today's date.

\newcommand{\ip}[2]{(#1, #2)}
                             % Defines \ip{arg1}{arg2} to mean
                             % (arg1, arg2).

%\newcommand{\ip}[2]{\langle #1 | #2\rangle}
                             % This is an alternative definition of
                             % \ip that is commented out.

\begin{document}             % End of preamble and beginning of text.

\maketitle                   % Produces the title.

This is an example input file.  Comparing it with
the output it generates can show you how to
produce a simple document of your own.

\section{Ordinary Text}      % Produces section heading.  Lower-level
                             % sections are begun with similar
                             % \subsection and \subsubsection commands.

The ends  of words and sentences are marked
  by   spaces. It  doesn't matter how many
spaces    you type; one is as good as 100.  The
end of   a line counts as a space.

One   or more   blank lines denote the  end
of  a paragraph.

Since any number of consecutive spaces are treated
like a single one, the formatting of the input
file makes no difference to
      \LaTeX,                % The \LaTeX command generates the LaTeX logo.
but it makes a difference to you.  When you use
\LaTeX, making your input file as easy to read
as possible will be a great help as you write
your document and when you change it.  This sample
file shows how you can add comments to your own input
file.

Because printing is different from typewriting,
there are a number of things that you have to do
differently when preparing an input file than if
you were just typing the document directly.
Quotation marks like
       ``this''
have to be handled specially, as do quotes within
quotes:
       ``\,`this'            % \, separates the double and single quote.
        is what I just
        wrote, not  `that'\,''.

Dashes come in three sizes: an
       intra-word
dash, a medium dash for number ranges like
       1--2,
and a punctuation
       dash---like
this.

A sentence-ending space should be larger than the
space between words within a sentence.  You
sometimes have to type special commands in
conjunction with punctuation characters to get
this right, as in the following sentence.
       Gnats, gnus, etc.\ all  % `\ ' makes an inter-word space.
       begin with G\@.         % \@ marks end-of-sentence punctuation.
You should check the spaces after periods when
reading your output to make sure you haven't
forgotten any special cases.  Generating an
ellipsis
       \ldots\               % `\ ' is needed after `\ldots' because TeX
                             % ignores spaces after command names like \ldots
                             % made from \ + letters.
                             %
                             % Note how a `%' character causes TeX to ignore
                             % the end of the input line, so these blank lines
                             % do not start a new paragraph.
                             %
with the right spacing around the periods requires
a special command.

\LaTeX\ interprets some common characters as
commands, so you must type special commands to
generate them.  These characters include the
following:
       \$ \& \% \# \{ and \}.

In printing, text is usually emphasized with an
       \emph{italic}
type style.

\begin{em}
   A long segment of text can also be emphasized
   in this way.  Text within such a segment can be
   given \emph{additional} emphasis.
\end{em}

It is sometimes necessary to prevent \LaTeX\ from
breaking a line where it might otherwise do so.
This may be at a space, as between the ``Mr.''\ and
``Jones'' in
       ``Mr.~Jones'',        % ~ produces an unbreakable interword space.
or within a word---especially when the word is a
symbol like
       \mbox{\emph{itemnum}}
that makes little sense when hyphenated across
lines.

Footnotes\footnote{This is an example of a footnote.}
pose no problem.

\LaTeX\ is good at typesetting mathematical formulas
like
       \( x-3y + z = 7 \)
or
       \( a_{1} > x^{2n} + y^{2n} > x' \)
or
       \( \ip{A}{B} = \sum_{i} a_{i} b_{i} \).
The spaces you type in a formula are
ignored.  Remember that a letter like
       $x$                   % $ ... $  and  \( ... \)  are equivalent
is a formula when it denotes a mathematical
symbol, and it should be typed as one.

\section{Displayed Text}

Text is displayed by indenting it from the left
margin.  Quotations are commonly displayed.  There
are short quotations
\begin{quote}
   This is a short quotation.  It consists of a
   single paragraph of text.  See how it is formatted.
\end{quote}
and longer ones.
\begin{quotation}
   This is a longer quotation.  It consists of two
   paragraphs of text, neither of which are
   particularly interesting.

   This is the second paragraph of the quotation.  It
   is just as dull as the first paragraph.
\end{quotation}
Another frequently-displayed structure is a list.
The following is an example of an \emph{itemized}
list.
\begin{itemize}
   \item This is the first item of an itemized list.
         Each item in the list is marked with a ``tick''.
         You don't have to worry about what kind of tick
         mark is used.

   \item This is the second item of the list.  It
         contains another list nested inside it.  The inner
         list is an \emph{enumerated} list.
         \begin{enumerate}
            \item This is the first item of an enumerated
                  list that is nested within the itemized list.

            \item This is the second item of the inner list.
                  \LaTeX\ allows you to nest lists deeper than
                  you really should.
         \end{enumerate}
         This is the rest of the second item of the outer
         list.  It is no more interesting than any other
         part of the item.
   \item This is the third item of the list.
\end{itemize}
You can even display poetry.
\begin{verse}
   There is an environment
    for verse \\             % The \\ command separates lines
   Whose features some poets % within a stanza.
   will curse.

                             % One or more blank lines separate stanzas.

   For instead of making\\
   Them do \emph{all} line breaking, \\
   It allows them to put too many words on a line when they'd rather be
   forced to be terse.
\end{verse}

Mathematical formulas may also be displayed.  A
displayed formula
is
one-line long; multiline
formulas require special formatting instructions.
   \[  \ip{\Gamma}{\psi'} = x'' + y^{2} + z_{i}^{n}\]
Don't start a paragraph with a displayed equation,
nor make one a paragraph by itself.

\end{document}               % End of document.

\section{Bench Huge Section 223}
% synthetic bench section 223
% This is a sample LaTeX input file.  (Version of 12 August 2004.)
%
% A '%' character causes TeX to ignore all remaining text on the line,
% and is used for comments like this one.

\documentclass{article}      % Specifies the document class

                             % The preamble begins here.
\title{An Example Document}  % Declares the document's title.
\author{Leslie Lamport}      % Declares the author's name.
\date{January 21, 1994}      % Deleting this command produces today's date.

\newcommand{\ip}[2]{(#1, #2)}
                             % Defines \ip{arg1}{arg2} to mean
                             % (arg1, arg2).

%\newcommand{\ip}[2]{\langle #1 | #2\rangle}
                             % This is an alternative definition of
                             % \ip that is commented out.

\begin{document}             % End of preamble and beginning of text.

\maketitle                   % Produces the title.

This is an example input file.  Comparing it with
the output it generates can show you how to
produce a simple document of your own.

\section{Ordinary Text}      % Produces section heading.  Lower-level
                             % sections are begun with similar
                             % \subsection and \subsubsection commands.

The ends  of words and sentences are marked
  by   spaces. It  doesn't matter how many
spaces    you type; one is as good as 100.  The
end of   a line counts as a space.

One   or more   blank lines denote the  end
of  a paragraph.

Since any number of consecutive spaces are treated
like a single one, the formatting of the input
file makes no difference to
      \LaTeX,                % The \LaTeX command generates the LaTeX logo.
but it makes a difference to you.  When you use
\LaTeX, making your input file as easy to read
as possible will be a great help as you write
your document and when you change it.  This sample
file shows how you can add comments to your own input
file.

Because printing is different from typewriting,
there are a number of things that you have to do
differently when preparing an input file than if
you were just typing the document directly.
Quotation marks like
       ``this''
have to be handled specially, as do quotes within
quotes:
       ``\,`this'            % \, separates the double and single quote.
        is what I just
        wrote, not  `that'\,''.

Dashes come in three sizes: an
       intra-word
dash, a medium dash for number ranges like
       1--2,
and a punctuation
       dash---like
this.

A sentence-ending space should be larger than the
space between words within a sentence.  You
sometimes have to type special commands in
conjunction with punctuation characters to get
this right, as in the following sentence.
       Gnats, gnus, etc.\ all  % `\ ' makes an inter-word space.
       begin with G\@.         % \@ marks end-of-sentence punctuation.
You should check the spaces after periods when
reading your output to make sure you haven't
forgotten any special cases.  Generating an
ellipsis
       \ldots\               % `\ ' is needed after `\ldots' because TeX
                             % ignores spaces after command names like \ldots
                             % made from \ + letters.
                             %
                             % Note how a `%' character causes TeX to ignore
                             % the end of the input line, so these blank lines
                             % do not start a new paragraph.
                             %
with the right spacing around the periods requires
a special command.

\LaTeX\ interprets some common characters as
commands, so you must type special commands to
generate them.  These characters include the
following:
       \$ \& \% \# \{ and \}.

In printing, text is usually emphasized with an
       \emph{italic}
type style.

\begin{em}
   A long segment of text can also be emphasized
   in this way.  Text within such a segment can be
   given \emph{additional} emphasis.
\end{em}

It is sometimes necessary to prevent \LaTeX\ from
breaking a line where it might otherwise do so.
This may be at a space, as between the ``Mr.''\ and
``Jones'' in
       ``Mr.~Jones'',        % ~ produces an unbreakable interword space.
or within a word---especially when the word is a
symbol like
       \mbox{\emph{itemnum}}
that makes little sense when hyphenated across
lines.

Footnotes\footnote{This is an example of a footnote.}
pose no problem.

\LaTeX\ is good at typesetting mathematical formulas
like
       \( x-3y + z = 7 \)
or
       \( a_{1} > x^{2n} + y^{2n} > x' \)
or
       \( \ip{A}{B} = \sum_{i} a_{i} b_{i} \).
The spaces you type in a formula are
ignored.  Remember that a letter like
       $x$                   % $ ... $  and  \( ... \)  are equivalent
is a formula when it denotes a mathematical
symbol, and it should be typed as one.

\section{Displayed Text}

Text is displayed by indenting it from the left
margin.  Quotations are commonly displayed.  There
are short quotations
\begin{quote}
   This is a short quotation.  It consists of a
   single paragraph of text.  See how it is formatted.
\end{quote}
and longer ones.
\begin{quotation}
   This is a longer quotation.  It consists of two
   paragraphs of text, neither of which are
   particularly interesting.

   This is the second paragraph of the quotation.  It
   is just as dull as the first paragraph.
\end{quotation}
Another frequently-displayed structure is a list.
The following is an example of an \emph{itemized}
list.
\begin{itemize}
   \item This is the first item of an itemized list.
         Each item in the list is marked with a ``tick''.
         You don't have to worry about what kind of tick
         mark is used.

   \item This is the second item of the list.  It
         contains another list nested inside it.  The inner
         list is an \emph{enumerated} list.
         \begin{enumerate}
            \item This is the first item of an enumerated
                  list that is nested within the itemized list.

            \item This is the second item of the inner list.
                  \LaTeX\ allows you to nest lists deeper than
                  you really should.
         \end{enumerate}
         This is the rest of the second item of the outer
         list.  It is no more interesting than any other
         part of the item.
   \item This is the third item of the list.
\end{itemize}
You can even display poetry.
\begin{verse}
   There is an environment
    for verse \\             % The \\ command separates lines
   Whose features some poets % within a stanza.
   will curse.

                             % One or more blank lines separate stanzas.

   For instead of making\\
   Them do \emph{all} line breaking, \\
   It allows them to put too many words on a line when they'd rather be
   forced to be terse.
\end{verse}

Mathematical formulas may also be displayed.  A
displayed formula
is
one-line long; multiline
formulas require special formatting instructions.
   \[  \ip{\Gamma}{\psi'} = x'' + y^{2} + z_{i}^{n}\]
Don't start a paragraph with a displayed equation,
nor make one a paragraph by itself.

\end{document}               % End of document.

\section{Bench Huge Section 224}
% synthetic bench section 224
% This is a sample LaTeX input file.  (Version of 12 August 2004.)
%
% A '%' character causes TeX to ignore all remaining text on the line,
% and is used for comments like this one.

\documentclass{article}      % Specifies the document class

                             % The preamble begins here.
\title{An Example Document}  % Declares the document's title.
\author{Leslie Lamport}      % Declares the author's name.
\date{January 21, 1994}      % Deleting this command produces today's date.

\newcommand{\ip}[2]{(#1, #2)}
                             % Defines \ip{arg1}{arg2} to mean
                             % (arg1, arg2).

%\newcommand{\ip}[2]{\langle #1 | #2\rangle}
                             % This is an alternative definition of
                             % \ip that is commented out.

\begin{document}             % End of preamble and beginning of text.

\maketitle                   % Produces the title.

This is an example input file.  Comparing it with
the output it generates can show you how to
produce a simple document of your own.

\section{Ordinary Text}      % Produces section heading.  Lower-level
                             % sections are begun with similar
                             % \subsection and \subsubsection commands.

The ends  of words and sentences are marked
  by   spaces. It  doesn't matter how many
spaces    you type; one is as good as 100.  The
end of   a line counts as a space.

One   or more   blank lines denote the  end
of  a paragraph.

Since any number of consecutive spaces are treated
like a single one, the formatting of the input
file makes no difference to
      \LaTeX,                % The \LaTeX command generates the LaTeX logo.
but it makes a difference to you.  When you use
\LaTeX, making your input file as easy to read
as possible will be a great help as you write
your document and when you change it.  This sample
file shows how you can add comments to your own input
file.

Because printing is different from typewriting,
there are a number of things that you have to do
differently when preparing an input file than if
you were just typing the document directly.
Quotation marks like
       ``this''
have to be handled specially, as do quotes within
quotes:
       ``\,`this'            % \, separates the double and single quote.
        is what I just
        wrote, not  `that'\,''.

Dashes come in three sizes: an
       intra-word
dash, a medium dash for number ranges like
       1--2,
and a punctuation
       dash---like
this.

A sentence-ending space should be larger than the
space between words within a sentence.  You
sometimes have to type special commands in
conjunction with punctuation characters to get
this right, as in the following sentence.
       Gnats, gnus, etc.\ all  % `\ ' makes an inter-word space.
       begin with G\@.         % \@ marks end-of-sentence punctuation.
You should check the spaces after periods when
reading your output to make sure you haven't
forgotten any special cases.  Generating an
ellipsis
       \ldots\               % `\ ' is needed after `\ldots' because TeX
                             % ignores spaces after command names like \ldots
                             % made from \ + letters.
                             %
                             % Note how a `%' character causes TeX to ignore
                             % the end of the input line, so these blank lines
                             % do not start a new paragraph.
                             %
with the right spacing around the periods requires
a special command.

\LaTeX\ interprets some common characters as
commands, so you must type special commands to
generate them.  These characters include the
following:
       \$ \& \% \# \{ and \}.

In printing, text is usually emphasized with an
       \emph{italic}
type style.

\begin{em}
   A long segment of text can also be emphasized
   in this way.  Text within such a segment can be
   given \emph{additional} emphasis.
\end{em}

It is sometimes necessary to prevent \LaTeX\ from
breaking a line where it might otherwise do so.
This may be at a space, as between the ``Mr.''\ and
``Jones'' in
       ``Mr.~Jones'',        % ~ produces an unbreakable interword space.
or within a word---especially when the word is a
symbol like
       \mbox{\emph{itemnum}}
that makes little sense when hyphenated across
lines.

Footnotes\footnote{This is an example of a footnote.}
pose no problem.

\LaTeX\ is good at typesetting mathematical formulas
like
       \( x-3y + z = 7 \)
or
       \( a_{1} > x^{2n} + y^{2n} > x' \)
or
       \( \ip{A}{B} = \sum_{i} a_{i} b_{i} \).
The spaces you type in a formula are
ignored.  Remember that a letter like
       $x$                   % $ ... $  and  \( ... \)  are equivalent
is a formula when it denotes a mathematical
symbol, and it should be typed as one.

\section{Displayed Text}

Text is displayed by indenting it from the left
margin.  Quotations are commonly displayed.  There
are short quotations
\begin{quote}
   This is a short quotation.  It consists of a
   single paragraph of text.  See how it is formatted.
\end{quote}
and longer ones.
\begin{quotation}
   This is a longer quotation.  It consists of two
   paragraphs of text, neither of which are
   particularly interesting.

   This is the second paragraph of the quotation.  It
   is just as dull as the first paragraph.
\end{quotation}
Another frequently-displayed structure is a list.
The following is an example of an \emph{itemized}
list.
\begin{itemize}
   \item This is the first item of an itemized list.
         Each item in the list is marked with a ``tick''.
         You don't have to worry about what kind of tick
         mark is used.

   \item This is the second item of the list.  It
         contains another list nested inside it.  The inner
         list is an \emph{enumerated} list.
         \begin{enumerate}
            \item This is the first item of an enumerated
                  list that is nested within the itemized list.

            \item This is the second item of the inner list.
                  \LaTeX\ allows you to nest lists deeper than
                  you really should.
         \end{enumerate}
         This is the rest of the second item of the outer
         list.  It is no more interesting than any other
         part of the item.
   \item This is the third item of the list.
\end{itemize}
You can even display poetry.
\begin{verse}
   There is an environment
    for verse \\             % The \\ command separates lines
   Whose features some poets % within a stanza.
   will curse.

                             % One or more blank lines separate stanzas.

   For instead of making\\
   Them do \emph{all} line breaking, \\
   It allows them to put too many words on a line when they'd rather be
   forced to be terse.
\end{verse}

Mathematical formulas may also be displayed.  A
displayed formula
is
one-line long; multiline
formulas require special formatting instructions.
   \[  \ip{\Gamma}{\psi'} = x'' + y^{2} + z_{i}^{n}\]
Don't start a paragraph with a displayed equation,
nor make one a paragraph by itself.

\end{document}               % End of document.

\section{Bench Huge Section 225}
% synthetic bench section 225
% This is a sample LaTeX input file.  (Version of 12 August 2004.)
%
% A '%' character causes TeX to ignore all remaining text on the line,
% and is used for comments like this one.

\documentclass{article}      % Specifies the document class

                             % The preamble begins here.
\title{An Example Document}  % Declares the document's title.
\author{Leslie Lamport}      % Declares the author's name.
\date{January 21, 1994}      % Deleting this command produces today's date.

\newcommand{\ip}[2]{(#1, #2)}
                             % Defines \ip{arg1}{arg2} to mean
                             % (arg1, arg2).

%\newcommand{\ip}[2]{\langle #1 | #2\rangle}
                             % This is an alternative definition of
                             % \ip that is commented out.

\begin{document}             % End of preamble and beginning of text.

\maketitle                   % Produces the title.

This is an example input file.  Comparing it with
the output it generates can show you how to
produce a simple document of your own.

\section{Ordinary Text}      % Produces section heading.  Lower-level
                             % sections are begun with similar
                             % \subsection and \subsubsection commands.

The ends  of words and sentences are marked
  by   spaces. It  doesn't matter how many
spaces    you type; one is as good as 100.  The
end of   a line counts as a space.

One   or more   blank lines denote the  end
of  a paragraph.

Since any number of consecutive spaces are treated
like a single one, the formatting of the input
file makes no difference to
      \LaTeX,                % The \LaTeX command generates the LaTeX logo.
but it makes a difference to you.  When you use
\LaTeX, making your input file as easy to read
as possible will be a great help as you write
your document and when you change it.  This sample
file shows how you can add comments to your own input
file.

Because printing is different from typewriting,
there are a number of things that you have to do
differently when preparing an input file than if
you were just typing the document directly.
Quotation marks like
       ``this''
have to be handled specially, as do quotes within
quotes:
       ``\,`this'            % \, separates the double and single quote.
        is what I just
        wrote, not  `that'\,''.

Dashes come in three sizes: an
       intra-word
dash, a medium dash for number ranges like
       1--2,
and a punctuation
       dash---like
this.

A sentence-ending space should be larger than the
space between words within a sentence.  You
sometimes have to type special commands in
conjunction with punctuation characters to get
this right, as in the following sentence.
       Gnats, gnus, etc.\ all  % `\ ' makes an inter-word space.
       begin with G\@.         % \@ marks end-of-sentence punctuation.
You should check the spaces after periods when
reading your output to make sure you haven't
forgotten any special cases.  Generating an
ellipsis
       \ldots\               % `\ ' is needed after `\ldots' because TeX
                             % ignores spaces after command names like \ldots
                             % made from \ + letters.
                             %
                             % Note how a `%' character causes TeX to ignore
                             % the end of the input line, so these blank lines
                             % do not start a new paragraph.
                             %
with the right spacing around the periods requires
a special command.

\LaTeX\ interprets some common characters as
commands, so you must type special commands to
generate them.  These characters include the
following:
       \$ \& \% \# \{ and \}.

In printing, text is usually emphasized with an
       \emph{italic}
type style.

\begin{em}
   A long segment of text can also be emphasized
   in this way.  Text within such a segment can be
   given \emph{additional} emphasis.
\end{em}

It is sometimes necessary to prevent \LaTeX\ from
breaking a line where it might otherwise do so.
This may be at a space, as between the ``Mr.''\ and
``Jones'' in
       ``Mr.~Jones'',        % ~ produces an unbreakable interword space.
or within a word---especially when the word is a
symbol like
       \mbox{\emph{itemnum}}
that makes little sense when hyphenated across
lines.

Footnotes\footnote{This is an example of a footnote.}
pose no problem.

\LaTeX\ is good at typesetting mathematical formulas
like
       \( x-3y + z = 7 \)
or
       \( a_{1} > x^{2n} + y^{2n} > x' \)
or
       \( \ip{A}{B} = \sum_{i} a_{i} b_{i} \).
The spaces you type in a formula are
ignored.  Remember that a letter like
       $x$                   % $ ... $  and  \( ... \)  are equivalent
is a formula when it denotes a mathematical
symbol, and it should be typed as one.

\section{Displayed Text}

Text is displayed by indenting it from the left
margin.  Quotations are commonly displayed.  There
are short quotations
\begin{quote}
   This is a short quotation.  It consists of a
   single paragraph of text.  See how it is formatted.
\end{quote}
and longer ones.
\begin{quotation}
   This is a longer quotation.  It consists of two
   paragraphs of text, neither of which are
   particularly interesting.

   This is the second paragraph of the quotation.  It
   is just as dull as the first paragraph.
\end{quotation}
Another frequently-displayed structure is a list.
The following is an example of an \emph{itemized}
list.
\begin{itemize}
   \item This is the first item of an itemized list.
         Each item in the list is marked with a ``tick''.
         You don't have to worry about what kind of tick
         mark is used.

   \item This is the second item of the list.  It
         contains another list nested inside it.  The inner
         list is an \emph{enumerated} list.
         \begin{enumerate}
            \item This is the first item of an enumerated
                  list that is nested within the itemized list.

            \item This is the second item of the inner list.
                  \LaTeX\ allows you to nest lists deeper than
                  you really should.
         \end{enumerate}
         This is the rest of the second item of the outer
         list.  It is no more interesting than any other
         part of the item.
   \item This is the third item of the list.
\end{itemize}
You can even display poetry.
\begin{verse}
   There is an environment
    for verse \\             % The \\ command separates lines
   Whose features some poets % within a stanza.
   will curse.

                             % One or more blank lines separate stanzas.

   For instead of making\\
   Them do \emph{all} line breaking, \\
   It allows them to put too many words on a line when they'd rather be
   forced to be terse.
\end{verse}

Mathematical formulas may also be displayed.  A
displayed formula
is
one-line long; multiline
formulas require special formatting instructions.
   \[  \ip{\Gamma}{\psi'} = x'' + y^{2} + z_{i}^{n}\]
Don't start a paragraph with a displayed equation,
nor make one a paragraph by itself.

\end{document}               % End of document.

\section{Bench Huge Section 226}
% synthetic bench section 226
% This is a sample LaTeX input file.  (Version of 12 August 2004.)
%
% A '%' character causes TeX to ignore all remaining text on the line,
% and is used for comments like this one.

\documentclass{article}      % Specifies the document class

                             % The preamble begins here.
\title{An Example Document}  % Declares the document's title.
\author{Leslie Lamport}      % Declares the author's name.
\date{January 21, 1994}      % Deleting this command produces today's date.

\newcommand{\ip}[2]{(#1, #2)}
                             % Defines \ip{arg1}{arg2} to mean
                             % (arg1, arg2).

%\newcommand{\ip}[2]{\langle #1 | #2\rangle}
                             % This is an alternative definition of
                             % \ip that is commented out.

\begin{document}             % End of preamble and beginning of text.

\maketitle                   % Produces the title.

This is an example input file.  Comparing it with
the output it generates can show you how to
produce a simple document of your own.

\section{Ordinary Text}      % Produces section heading.  Lower-level
                             % sections are begun with similar
                             % \subsection and \subsubsection commands.

The ends  of words and sentences are marked
  by   spaces. It  doesn't matter how many
spaces    you type; one is as good as 100.  The
end of   a line counts as a space.

One   or more   blank lines denote the  end
of  a paragraph.

Since any number of consecutive spaces are treated
like a single one, the formatting of the input
file makes no difference to
      \LaTeX,                % The \LaTeX command generates the LaTeX logo.
but it makes a difference to you.  When you use
\LaTeX, making your input file as easy to read
as possible will be a great help as you write
your document and when you change it.  This sample
file shows how you can add comments to your own input
file.

Because printing is different from typewriting,
there are a number of things that you have to do
differently when preparing an input file than if
you were just typing the document directly.
Quotation marks like
       ``this''
have to be handled specially, as do quotes within
quotes:
       ``\,`this'            % \, separates the double and single quote.
        is what I just
        wrote, not  `that'\,''.

Dashes come in three sizes: an
       intra-word
dash, a medium dash for number ranges like
       1--2,
and a punctuation
       dash---like
this.

A sentence-ending space should be larger than the
space between words within a sentence.  You
sometimes have to type special commands in
conjunction with punctuation characters to get
this right, as in the following sentence.
       Gnats, gnus, etc.\ all  % `\ ' makes an inter-word space.
       begin with G\@.         % \@ marks end-of-sentence punctuation.
You should check the spaces after periods when
reading your output to make sure you haven't
forgotten any special cases.  Generating an
ellipsis
       \ldots\               % `\ ' is needed after `\ldots' because TeX
                             % ignores spaces after command names like \ldots
                             % made from \ + letters.
                             %
                             % Note how a `%' character causes TeX to ignore
                             % the end of the input line, so these blank lines
                             % do not start a new paragraph.
                             %
with the right spacing around the periods requires
a special command.

\LaTeX\ interprets some common characters as
commands, so you must type special commands to
generate them.  These characters include the
following:
       \$ \& \% \# \{ and \}.

In printing, text is usually emphasized with an
       \emph{italic}
type style.

\begin{em}
   A long segment of text can also be emphasized
   in this way.  Text within such a segment can be
   given \emph{additional} emphasis.
\end{em}

It is sometimes necessary to prevent \LaTeX\ from
breaking a line where it might otherwise do so.
This may be at a space, as between the ``Mr.''\ and
``Jones'' in
       ``Mr.~Jones'',        % ~ produces an unbreakable interword space.
or within a word---especially when the word is a
symbol like
       \mbox{\emph{itemnum}}
that makes little sense when hyphenated across
lines.

Footnotes\footnote{This is an example of a footnote.}
pose no problem.

\LaTeX\ is good at typesetting mathematical formulas
like
       \( x-3y + z = 7 \)
or
       \( a_{1} > x^{2n} + y^{2n} > x' \)
or
       \( \ip{A}{B} = \sum_{i} a_{i} b_{i} \).
The spaces you type in a formula are
ignored.  Remember that a letter like
       $x$                   % $ ... $  and  \( ... \)  are equivalent
is a formula when it denotes a mathematical
symbol, and it should be typed as one.

\section{Displayed Text}

Text is displayed by indenting it from the left
margin.  Quotations are commonly displayed.  There
are short quotations
\begin{quote}
   This is a short quotation.  It consists of a
   single paragraph of text.  See how it is formatted.
\end{quote}
and longer ones.
\begin{quotation}
   This is a longer quotation.  It consists of two
   paragraphs of text, neither of which are
   particularly interesting.

   This is the second paragraph of the quotation.  It
   is just as dull as the first paragraph.
\end{quotation}
Another frequently-displayed structure is a list.
The following is an example of an \emph{itemized}
list.
\begin{itemize}
   \item This is the first item of an itemized list.
         Each item in the list is marked with a ``tick''.
         You don't have to worry about what kind of tick
         mark is used.

   \item This is the second item of the list.  It
         contains another list nested inside it.  The inner
         list is an \emph{enumerated} list.
         \begin{enumerate}
            \item This is the first item of an enumerated
                  list that is nested within the itemized list.

            \item This is the second item of the inner list.
                  \LaTeX\ allows you to nest lists deeper than
                  you really should.
         \end{enumerate}
         This is the rest of the second item of the outer
         list.  It is no more interesting than any other
         part of the item.
   \item This is the third item of the list.
\end{itemize}
You can even display poetry.
\begin{verse}
   There is an environment
    for verse \\             % The \\ command separates lines
   Whose features some poets % within a stanza.
   will curse.

                             % One or more blank lines separate stanzas.

   For instead of making\\
   Them do \emph{all} line breaking, \\
   It allows them to put too many words on a line when they'd rather be
   forced to be terse.
\end{verse}

Mathematical formulas may also be displayed.  A
displayed formula
is
one-line long; multiline
formulas require special formatting instructions.
   \[  \ip{\Gamma}{\psi'} = x'' + y^{2} + z_{i}^{n}\]
Don't start a paragraph with a displayed equation,
nor make one a paragraph by itself.

\end{document}               % End of document.

\section{Bench Huge Section 227}
% synthetic bench section 227
% This is a sample LaTeX input file.  (Version of 12 August 2004.)
%
% A '%' character causes TeX to ignore all remaining text on the line,
% and is used for comments like this one.

\documentclass{article}      % Specifies the document class

                             % The preamble begins here.
\title{An Example Document}  % Declares the document's title.
\author{Leslie Lamport}      % Declares the author's name.
\date{January 21, 1994}      % Deleting this command produces today's date.

\newcommand{\ip}[2]{(#1, #2)}
                             % Defines \ip{arg1}{arg2} to mean
                             % (arg1, arg2).

%\newcommand{\ip}[2]{\langle #1 | #2\rangle}
                             % This is an alternative definition of
                             % \ip that is commented out.

\begin{document}             % End of preamble and beginning of text.

\maketitle                   % Produces the title.

This is an example input file.  Comparing it with
the output it generates can show you how to
produce a simple document of your own.

\section{Ordinary Text}      % Produces section heading.  Lower-level
                             % sections are begun with similar
                             % \subsection and \subsubsection commands.

The ends  of words and sentences are marked
  by   spaces. It  doesn't matter how many
spaces    you type; one is as good as 100.  The
end of   a line counts as a space.

One   or more   blank lines denote the  end
of  a paragraph.

Since any number of consecutive spaces are treated
like a single one, the formatting of the input
file makes no difference to
      \LaTeX,                % The \LaTeX command generates the LaTeX logo.
but it makes a difference to you.  When you use
\LaTeX, making your input file as easy to read
as possible will be a great help as you write
your document and when you change it.  This sample
file shows how you can add comments to your own input
file.

Because printing is different from typewriting,
there are a number of things that you have to do
differently when preparing an input file than if
you were just typing the document directly.
Quotation marks like
       ``this''
have to be handled specially, as do quotes within
quotes:
       ``\,`this'            % \, separates the double and single quote.
        is what I just
        wrote, not  `that'\,''.

Dashes come in three sizes: an
       intra-word
dash, a medium dash for number ranges like
       1--2,
and a punctuation
       dash---like
this.

A sentence-ending space should be larger than the
space between words within a sentence.  You
sometimes have to type special commands in
conjunction with punctuation characters to get
this right, as in the following sentence.
       Gnats, gnus, etc.\ all  % `\ ' makes an inter-word space.
       begin with G\@.         % \@ marks end-of-sentence punctuation.
You should check the spaces after periods when
reading your output to make sure you haven't
forgotten any special cases.  Generating an
ellipsis
       \ldots\               % `\ ' is needed after `\ldots' because TeX
                             % ignores spaces after command names like \ldots
                             % made from \ + letters.
                             %
                             % Note how a `%' character causes TeX to ignore
                             % the end of the input line, so these blank lines
                             % do not start a new paragraph.
                             %
with the right spacing around the periods requires
a special command.

\LaTeX\ interprets some common characters as
commands, so you must type special commands to
generate them.  These characters include the
following:
       \$ \& \% \# \{ and \}.

In printing, text is usually emphasized with an
       \emph{italic}
type style.

\begin{em}
   A long segment of text can also be emphasized
   in this way.  Text within such a segment can be
   given \emph{additional} emphasis.
\end{em}

It is sometimes necessary to prevent \LaTeX\ from
breaking a line where it might otherwise do so.
This may be at a space, as between the ``Mr.''\ and
``Jones'' in
       ``Mr.~Jones'',        % ~ produces an unbreakable interword space.
or within a word---especially when the word is a
symbol like
       \mbox{\emph{itemnum}}
that makes little sense when hyphenated across
lines.

Footnotes\footnote{This is an example of a footnote.}
pose no problem.

\LaTeX\ is good at typesetting mathematical formulas
like
       \( x-3y + z = 7 \)
or
       \( a_{1} > x^{2n} + y^{2n} > x' \)
or
       \( \ip{A}{B} = \sum_{i} a_{i} b_{i} \).
The spaces you type in a formula are
ignored.  Remember that a letter like
       $x$                   % $ ... $  and  \( ... \)  are equivalent
is a formula when it denotes a mathematical
symbol, and it should be typed as one.

\section{Displayed Text}

Text is displayed by indenting it from the left
margin.  Quotations are commonly displayed.  There
are short quotations
\begin{quote}
   This is a short quotation.  It consists of a
   single paragraph of text.  See how it is formatted.
\end{quote}
and longer ones.
\begin{quotation}
   This is a longer quotation.  It consists of two
   paragraphs of text, neither of which are
   particularly interesting.

   This is the second paragraph of the quotation.  It
   is just as dull as the first paragraph.
\end{quotation}
Another frequently-displayed structure is a list.
The following is an example of an \emph{itemized}
list.
\begin{itemize}
   \item This is the first item of an itemized list.
         Each item in the list is marked with a ``tick''.
         You don't have to worry about what kind of tick
         mark is used.

   \item This is the second item of the list.  It
         contains another list nested inside it.  The inner
         list is an \emph{enumerated} list.
         \begin{enumerate}
            \item This is the first item of an enumerated
                  list that is nested within the itemized list.

            \item This is the second item of the inner list.
                  \LaTeX\ allows you to nest lists deeper than
                  you really should.
         \end{enumerate}
         This is the rest of the second item of the outer
         list.  It is no more interesting than any other
         part of the item.
   \item This is the third item of the list.
\end{itemize}
You can even display poetry.
\begin{verse}
   There is an environment
    for verse \\             % The \\ command separates lines
   Whose features some poets % within a stanza.
   will curse.

                             % One or more blank lines separate stanzas.

   For instead of making\\
   Them do \emph{all} line breaking, \\
   It allows them to put too many words on a line when they'd rather be
   forced to be terse.
\end{verse}

Mathematical formulas may also be displayed.  A
displayed formula
is
one-line long; multiline
formulas require special formatting instructions.
   \[  \ip{\Gamma}{\psi'} = x'' + y^{2} + z_{i}^{n}\]
Don't start a paragraph with a displayed equation,
nor make one a paragraph by itself.

\end{document}               % End of document.

\section{Bench Huge Section 228}
% synthetic bench section 228
% This is a sample LaTeX input file.  (Version of 12 August 2004.)
%
% A '%' character causes TeX to ignore all remaining text on the line,
% and is used for comments like this one.

\documentclass{article}      % Specifies the document class

                             % The preamble begins here.
\title{An Example Document}  % Declares the document's title.
\author{Leslie Lamport}      % Declares the author's name.
\date{January 21, 1994}      % Deleting this command produces today's date.

\newcommand{\ip}[2]{(#1, #2)}
                             % Defines \ip{arg1}{arg2} to mean
                             % (arg1, arg2).

%\newcommand{\ip}[2]{\langle #1 | #2\rangle}
                             % This is an alternative definition of
                             % \ip that is commented out.

\begin{document}             % End of preamble and beginning of text.

\maketitle                   % Produces the title.

This is an example input file.  Comparing it with
the output it generates can show you how to
produce a simple document of your own.

\section{Ordinary Text}      % Produces section heading.  Lower-level
                             % sections are begun with similar
                             % \subsection and \subsubsection commands.

The ends  of words and sentences are marked
  by   spaces. It  doesn't matter how many
spaces    you type; one is as good as 100.  The
end of   a line counts as a space.

One   or more   blank lines denote the  end
of  a paragraph.

Since any number of consecutive spaces are treated
like a single one, the formatting of the input
file makes no difference to
      \LaTeX,                % The \LaTeX command generates the LaTeX logo.
but it makes a difference to you.  When you use
\LaTeX, making your input file as easy to read
as possible will be a great help as you write
your document and when you change it.  This sample
file shows how you can add comments to your own input
file.

Because printing is different from typewriting,
there are a number of things that you have to do
differently when preparing an input file than if
you were just typing the document directly.
Quotation marks like
       ``this''
have to be handled specially, as do quotes within
quotes:
       ``\,`this'            % \, separates the double and single quote.
        is what I just
        wrote, not  `that'\,''.

Dashes come in three sizes: an
       intra-word
dash, a medium dash for number ranges like
       1--2,
and a punctuation
       dash---like
this.

A sentence-ending space should be larger than the
space between words within a sentence.  You
sometimes have to type special commands in
conjunction with punctuation characters to get
this right, as in the following sentence.
       Gnats, gnus, etc.\ all  % `\ ' makes an inter-word space.
       begin with G\@.         % \@ marks end-of-sentence punctuation.
You should check the spaces after periods when
reading your output to make sure you haven't
forgotten any special cases.  Generating an
ellipsis
       \ldots\               % `\ ' is needed after `\ldots' because TeX
                             % ignores spaces after command names like \ldots
                             % made from \ + letters.
                             %
                             % Note how a `%' character causes TeX to ignore
                             % the end of the input line, so these blank lines
                             % do not start a new paragraph.
                             %
with the right spacing around the periods requires
a special command.

\LaTeX\ interprets some common characters as
commands, so you must type special commands to
generate them.  These characters include the
following:
       \$ \& \% \# \{ and \}.

In printing, text is usually emphasized with an
       \emph{italic}
type style.

\begin{em}
   A long segment of text can also be emphasized
   in this way.  Text within such a segment can be
   given \emph{additional} emphasis.
\end{em}

It is sometimes necessary to prevent \LaTeX\ from
breaking a line where it might otherwise do so.
This may be at a space, as between the ``Mr.''\ and
``Jones'' in
       ``Mr.~Jones'',        % ~ produces an unbreakable interword space.
or within a word---especially when the word is a
symbol like
       \mbox{\emph{itemnum}}
that makes little sense when hyphenated across
lines.

Footnotes\footnote{This is an example of a footnote.}
pose no problem.

\LaTeX\ is good at typesetting mathematical formulas
like
       \( x-3y + z = 7 \)
or
       \( a_{1} > x^{2n} + y^{2n} > x' \)
or
       \( \ip{A}{B} = \sum_{i} a_{i} b_{i} \).
The spaces you type in a formula are
ignored.  Remember that a letter like
       $x$                   % $ ... $  and  \( ... \)  are equivalent
is a formula when it denotes a mathematical
symbol, and it should be typed as one.

\section{Displayed Text}

Text is displayed by indenting it from the left
margin.  Quotations are commonly displayed.  There
are short quotations
\begin{quote}
   This is a short quotation.  It consists of a
   single paragraph of text.  See how it is formatted.
\end{quote}
and longer ones.
\begin{quotation}
   This is a longer quotation.  It consists of two
   paragraphs of text, neither of which are
   particularly interesting.

   This is the second paragraph of the quotation.  It
   is just as dull as the first paragraph.
\end{quotation}
Another frequently-displayed structure is a list.
The following is an example of an \emph{itemized}
list.
\begin{itemize}
   \item This is the first item of an itemized list.
         Each item in the list is marked with a ``tick''.
         You don't have to worry about what kind of tick
         mark is used.

   \item This is the second item of the list.  It
         contains another list nested inside it.  The inner
         list is an \emph{enumerated} list.
         \begin{enumerate}
            \item This is the first item of an enumerated
                  list that is nested within the itemized list.

            \item This is the second item of the inner list.
                  \LaTeX\ allows you to nest lists deeper than
                  you really should.
         \end{enumerate}
         This is the rest of the second item of the outer
         list.  It is no more interesting than any other
         part of the item.
   \item This is the third item of the list.
\end{itemize}
You can even display poetry.
\begin{verse}
   There is an environment
    for verse \\             % The \\ command separates lines
   Whose features some poets % within a stanza.
   will curse.

                             % One or more blank lines separate stanzas.

   For instead of making\\
   Them do \emph{all} line breaking, \\
   It allows them to put too many words on a line when they'd rather be
   forced to be terse.
\end{verse}

Mathematical formulas may also be displayed.  A
displayed formula
is
one-line long; multiline
formulas require special formatting instructions.
   \[  \ip{\Gamma}{\psi'} = x'' + y^{2} + z_{i}^{n}\]
Don't start a paragraph with a displayed equation,
nor make one a paragraph by itself.

\end{document}               % End of document.

\section{Bench Huge Section 229}
% synthetic bench section 229
% This is a sample LaTeX input file.  (Version of 12 August 2004.)
%
% A '%' character causes TeX to ignore all remaining text on the line,
% and is used for comments like this one.

\documentclass{article}      % Specifies the document class

                             % The preamble begins here.
\title{An Example Document}  % Declares the document's title.
\author{Leslie Lamport}      % Declares the author's name.
\date{January 21, 1994}      % Deleting this command produces today's date.

\newcommand{\ip}[2]{(#1, #2)}
                             % Defines \ip{arg1}{arg2} to mean
                             % (arg1, arg2).

%\newcommand{\ip}[2]{\langle #1 | #2\rangle}
                             % This is an alternative definition of
                             % \ip that is commented out.

\begin{document}             % End of preamble and beginning of text.

\maketitle                   % Produces the title.

This is an example input file.  Comparing it with
the output it generates can show you how to
produce a simple document of your own.

\section{Ordinary Text}      % Produces section heading.  Lower-level
                             % sections are begun with similar
                             % \subsection and \subsubsection commands.

The ends  of words and sentences are marked
  by   spaces. It  doesn't matter how many
spaces    you type; one is as good as 100.  The
end of   a line counts as a space.

One   or more   blank lines denote the  end
of  a paragraph.

Since any number of consecutive spaces are treated
like a single one, the formatting of the input
file makes no difference to
      \LaTeX,                % The \LaTeX command generates the LaTeX logo.
but it makes a difference to you.  When you use
\LaTeX, making your input file as easy to read
as possible will be a great help as you write
your document and when you change it.  This sample
file shows how you can add comments to your own input
file.

Because printing is different from typewriting,
there are a number of things that you have to do
differently when preparing an input file than if
you were just typing the document directly.
Quotation marks like
       ``this''
have to be handled specially, as do quotes within
quotes:
       ``\,`this'            % \, separates the double and single quote.
        is what I just
        wrote, not  `that'\,''.

Dashes come in three sizes: an
       intra-word
dash, a medium dash for number ranges like
       1--2,
and a punctuation
       dash---like
this.

A sentence-ending space should be larger than the
space between words within a sentence.  You
sometimes have to type special commands in
conjunction with punctuation characters to get
this right, as in the following sentence.
       Gnats, gnus, etc.\ all  % `\ ' makes an inter-word space.
       begin with G\@.         % \@ marks end-of-sentence punctuation.
You should check the spaces after periods when
reading your output to make sure you haven't
forgotten any special cases.  Generating an
ellipsis
       \ldots\               % `\ ' is needed after `\ldots' because TeX
                             % ignores spaces after command names like \ldots
                             % made from \ + letters.
                             %
                             % Note how a `%' character causes TeX to ignore
                             % the end of the input line, so these blank lines
                             % do not start a new paragraph.
                             %
with the right spacing around the periods requires
a special command.

\LaTeX\ interprets some common characters as
commands, so you must type special commands to
generate them.  These characters include the
following:
       \$ \& \% \# \{ and \}.

In printing, text is usually emphasized with an
       \emph{italic}
type style.

\begin{em}
   A long segment of text can also be emphasized
   in this way.  Text within such a segment can be
   given \emph{additional} emphasis.
\end{em}

It is sometimes necessary to prevent \LaTeX\ from
breaking a line where it might otherwise do so.
This may be at a space, as between the ``Mr.''\ and
``Jones'' in
       ``Mr.~Jones'',        % ~ produces an unbreakable interword space.
or within a word---especially when the word is a
symbol like
       \mbox{\emph{itemnum}}
that makes little sense when hyphenated across
lines.

Footnotes\footnote{This is an example of a footnote.}
pose no problem.

\LaTeX\ is good at typesetting mathematical formulas
like
       \( x-3y + z = 7 \)
or
       \( a_{1} > x^{2n} + y^{2n} > x' \)
or
       \( \ip{A}{B} = \sum_{i} a_{i} b_{i} \).
The spaces you type in a formula are
ignored.  Remember that a letter like
       $x$                   % $ ... $  and  \( ... \)  are equivalent
is a formula when it denotes a mathematical
symbol, and it should be typed as one.

\section{Displayed Text}

Text is displayed by indenting it from the left
margin.  Quotations are commonly displayed.  There
are short quotations
\begin{quote}
   This is a short quotation.  It consists of a
   single paragraph of text.  See how it is formatted.
\end{quote}
and longer ones.
\begin{quotation}
   This is a longer quotation.  It consists of two
   paragraphs of text, neither of which are
   particularly interesting.

   This is the second paragraph of the quotation.  It
   is just as dull as the first paragraph.
\end{quotation}
Another frequently-displayed structure is a list.
The following is an example of an \emph{itemized}
list.
\begin{itemize}
   \item This is the first item of an itemized list.
         Each item in the list is marked with a ``tick''.
         You don't have to worry about what kind of tick
         mark is used.

   \item This is the second item of the list.  It
         contains another list nested inside it.  The inner
         list is an \emph{enumerated} list.
         \begin{enumerate}
            \item This is the first item of an enumerated
                  list that is nested within the itemized list.

            \item This is the second item of the inner list.
                  \LaTeX\ allows you to nest lists deeper than
                  you really should.
         \end{enumerate}
         This is the rest of the second item of the outer
         list.  It is no more interesting than any other
         part of the item.
   \item This is the third item of the list.
\end{itemize}
You can even display poetry.
\begin{verse}
   There is an environment
    for verse \\             % The \\ command separates lines
   Whose features some poets % within a stanza.
   will curse.

                             % One or more blank lines separate stanzas.

   For instead of making\\
   Them do \emph{all} line breaking, \\
   It allows them to put too many words on a line when they'd rather be
   forced to be terse.
\end{verse}

Mathematical formulas may also be displayed.  A
displayed formula
is
one-line long; multiline
formulas require special formatting instructions.
   \[  \ip{\Gamma}{\psi'} = x'' + y^{2} + z_{i}^{n}\]
Don't start a paragraph with a displayed equation,
nor make one a paragraph by itself.

\end{document}               % End of document.

\section{Bench Huge Section 230}
% synthetic bench section 230
% This is a sample LaTeX input file.  (Version of 12 August 2004.)
%
% A '%' character causes TeX to ignore all remaining text on the line,
% and is used for comments like this one.

\documentclass{article}      % Specifies the document class

                             % The preamble begins here.
\title{An Example Document}  % Declares the document's title.
\author{Leslie Lamport}      % Declares the author's name.
\date{January 21, 1994}      % Deleting this command produces today's date.

\newcommand{\ip}[2]{(#1, #2)}
                             % Defines \ip{arg1}{arg2} to mean
                             % (arg1, arg2).

%\newcommand{\ip}[2]{\langle #1 | #2\rangle}
                             % This is an alternative definition of
                             % \ip that is commented out.

\begin{document}             % End of preamble and beginning of text.

\maketitle                   % Produces the title.

This is an example input file.  Comparing it with
the output it generates can show you how to
produce a simple document of your own.

\section{Ordinary Text}      % Produces section heading.  Lower-level
                             % sections are begun with similar
                             % \subsection and \subsubsection commands.

The ends  of words and sentences are marked
  by   spaces. It  doesn't matter how many
spaces    you type; one is as good as 100.  The
end of   a line counts as a space.

One   or more   blank lines denote the  end
of  a paragraph.

Since any number of consecutive spaces are treated
like a single one, the formatting of the input
file makes no difference to
      \LaTeX,                % The \LaTeX command generates the LaTeX logo.
but it makes a difference to you.  When you use
\LaTeX, making your input file as easy to read
as possible will be a great help as you write
your document and when you change it.  This sample
file shows how you can add comments to your own input
file.

Because printing is different from typewriting,
there are a number of things that you have to do
differently when preparing an input file than if
you were just typing the document directly.
Quotation marks like
       ``this''
have to be handled specially, as do quotes within
quotes:
       ``\,`this'            % \, separates the double and single quote.
        is what I just
        wrote, not  `that'\,''.

Dashes come in three sizes: an
       intra-word
dash, a medium dash for number ranges like
       1--2,
and a punctuation
       dash---like
this.

A sentence-ending space should be larger than the
space between words within a sentence.  You
sometimes have to type special commands in
conjunction with punctuation characters to get
this right, as in the following sentence.
       Gnats, gnus, etc.\ all  % `\ ' makes an inter-word space.
       begin with G\@.         % \@ marks end-of-sentence punctuation.
You should check the spaces after periods when
reading your output to make sure you haven't
forgotten any special cases.  Generating an
ellipsis
       \ldots\               % `\ ' is needed after `\ldots' because TeX
                             % ignores spaces after command names like \ldots
                             % made from \ + letters.
                             %
                             % Note how a `%' character causes TeX to ignore
                             % the end of the input line, so these blank lines
                             % do not start a new paragraph.
                             %
with the right spacing around the periods requires
a special command.

\LaTeX\ interprets some common characters as
commands, so you must type special commands to
generate them.  These characters include the
following:
       \$ \& \% \# \{ and \}.

In printing, text is usually emphasized with an
       \emph{italic}
type style.

\begin{em}
   A long segment of text can also be emphasized
   in this way.  Text within such a segment can be
   given \emph{additional} emphasis.
\end{em}

It is sometimes necessary to prevent \LaTeX\ from
breaking a line where it might otherwise do so.
This may be at a space, as between the ``Mr.''\ and
``Jones'' in
       ``Mr.~Jones'',        % ~ produces an unbreakable interword space.
or within a word---especially when the word is a
symbol like
       \mbox{\emph{itemnum}}
that makes little sense when hyphenated across
lines.

Footnotes\footnote{This is an example of a footnote.}
pose no problem.

\LaTeX\ is good at typesetting mathematical formulas
like
       \( x-3y + z = 7 \)
or
       \( a_{1} > x^{2n} + y^{2n} > x' \)
or
       \( \ip{A}{B} = \sum_{i} a_{i} b_{i} \).
The spaces you type in a formula are
ignored.  Remember that a letter like
       $x$                   % $ ... $  and  \( ... \)  are equivalent
is a formula when it denotes a mathematical
symbol, and it should be typed as one.

\section{Displayed Text}

Text is displayed by indenting it from the left
margin.  Quotations are commonly displayed.  There
are short quotations
\begin{quote}
   This is a short quotation.  It consists of a
   single paragraph of text.  See how it is formatted.
\end{quote}
and longer ones.
\begin{quotation}
   This is a longer quotation.  It consists of two
   paragraphs of text, neither of which are
   particularly interesting.

   This is the second paragraph of the quotation.  It
   is just as dull as the first paragraph.
\end{quotation}
Another frequently-displayed structure is a list.
The following is an example of an \emph{itemized}
list.
\begin{itemize}
   \item This is the first item of an itemized list.
         Each item in the list is marked with a ``tick''.
         You don't have to worry about what kind of tick
         mark is used.

   \item This is the second item of the list.  It
         contains another list nested inside it.  The inner
         list is an \emph{enumerated} list.
         \begin{enumerate}
            \item This is the first item of an enumerated
                  list that is nested within the itemized list.

            \item This is the second item of the inner list.
                  \LaTeX\ allows you to nest lists deeper than
                  you really should.
         \end{enumerate}
         This is the rest of the second item of the outer
         list.  It is no more interesting than any other
         part of the item.
   \item This is the third item of the list.
\end{itemize}
You can even display poetry.
\begin{verse}
   There is an environment
    for verse \\             % The \\ command separates lines
   Whose features some poets % within a stanza.
   will curse.

                             % One or more blank lines separate stanzas.

   For instead of making\\
   Them do \emph{all} line breaking, \\
   It allows them to put too many words on a line when they'd rather be
   forced to be terse.
\end{verse}

Mathematical formulas may also be displayed.  A
displayed formula
is
one-line long; multiline
formulas require special formatting instructions.
   \[  \ip{\Gamma}{\psi'} = x'' + y^{2} + z_{i}^{n}\]
Don't start a paragraph with a displayed equation,
nor make one a paragraph by itself.

\end{document}               % End of document.

\section{Bench Huge Section 231}
% synthetic bench section 231
% This is a sample LaTeX input file.  (Version of 12 August 2004.)
%
% A '%' character causes TeX to ignore all remaining text on the line,
% and is used for comments like this one.

\documentclass{article}      % Specifies the document class

                             % The preamble begins here.
\title{An Example Document}  % Declares the document's title.
\author{Leslie Lamport}      % Declares the author's name.
\date{January 21, 1994}      % Deleting this command produces today's date.

\newcommand{\ip}[2]{(#1, #2)}
                             % Defines \ip{arg1}{arg2} to mean
                             % (arg1, arg2).

%\newcommand{\ip}[2]{\langle #1 | #2\rangle}
                             % This is an alternative definition of
                             % \ip that is commented out.

\begin{document}             % End of preamble and beginning of text.

\maketitle                   % Produces the title.

This is an example input file.  Comparing it with
the output it generates can show you how to
produce a simple document of your own.

\section{Ordinary Text}      % Produces section heading.  Lower-level
                             % sections are begun with similar
                             % \subsection and \subsubsection commands.

The ends  of words and sentences are marked
  by   spaces. It  doesn't matter how many
spaces    you type; one is as good as 100.  The
end of   a line counts as a space.

One   or more   blank lines denote the  end
of  a paragraph.

Since any number of consecutive spaces are treated
like a single one, the formatting of the input
file makes no difference to
      \LaTeX,                % The \LaTeX command generates the LaTeX logo.
but it makes a difference to you.  When you use
\LaTeX, making your input file as easy to read
as possible will be a great help as you write
your document and when you change it.  This sample
file shows how you can add comments to your own input
file.

Because printing is different from typewriting,
there are a number of things that you have to do
differently when preparing an input file than if
you were just typing the document directly.
Quotation marks like
       ``this''
have to be handled specially, as do quotes within
quotes:
       ``\,`this'            % \, separates the double and single quote.
        is what I just
        wrote, not  `that'\,''.

Dashes come in three sizes: an
       intra-word
dash, a medium dash for number ranges like
       1--2,
and a punctuation
       dash---like
this.

A sentence-ending space should be larger than the
space between words within a sentence.  You
sometimes have to type special commands in
conjunction with punctuation characters to get
this right, as in the following sentence.
       Gnats, gnus, etc.\ all  % `\ ' makes an inter-word space.
       begin with G\@.         % \@ marks end-of-sentence punctuation.
You should check the spaces after periods when
reading your output to make sure you haven't
forgotten any special cases.  Generating an
ellipsis
       \ldots\               % `\ ' is needed after `\ldots' because TeX
                             % ignores spaces after command names like \ldots
                             % made from \ + letters.
                             %
                             % Note how a `%' character causes TeX to ignore
                             % the end of the input line, so these blank lines
                             % do not start a new paragraph.
                             %
with the right spacing around the periods requires
a special command.

\LaTeX\ interprets some common characters as
commands, so you must type special commands to
generate them.  These characters include the
following:
       \$ \& \% \# \{ and \}.

In printing, text is usually emphasized with an
       \emph{italic}
type style.

\begin{em}
   A long segment of text can also be emphasized
   in this way.  Text within such a segment can be
   given \emph{additional} emphasis.
\end{em}

It is sometimes necessary to prevent \LaTeX\ from
breaking a line where it might otherwise do so.
This may be at a space, as between the ``Mr.''\ and
``Jones'' in
       ``Mr.~Jones'',        % ~ produces an unbreakable interword space.
or within a word---especially when the word is a
symbol like
       \mbox{\emph{itemnum}}
that makes little sense when hyphenated across
lines.

Footnotes\footnote{This is an example of a footnote.}
pose no problem.

\LaTeX\ is good at typesetting mathematical formulas
like
       \( x-3y + z = 7 \)
or
       \( a_{1} > x^{2n} + y^{2n} > x' \)
or
       \( \ip{A}{B} = \sum_{i} a_{i} b_{i} \).
The spaces you type in a formula are
ignored.  Remember that a letter like
       $x$                   % $ ... $  and  \( ... \)  are equivalent
is a formula when it denotes a mathematical
symbol, and it should be typed as one.

\section{Displayed Text}

Text is displayed by indenting it from the left
margin.  Quotations are commonly displayed.  There
are short quotations
\begin{quote}
   This is a short quotation.  It consists of a
   single paragraph of text.  See how it is formatted.
\end{quote}
and longer ones.
\begin{quotation}
   This is a longer quotation.  It consists of two
   paragraphs of text, neither of which are
   particularly interesting.

   This is the second paragraph of the quotation.  It
   is just as dull as the first paragraph.
\end{quotation}
Another frequently-displayed structure is a list.
The following is an example of an \emph{itemized}
list.
\begin{itemize}
   \item This is the first item of an itemized list.
         Each item in the list is marked with a ``tick''.
         You don't have to worry about what kind of tick
         mark is used.

   \item This is the second item of the list.  It
         contains another list nested inside it.  The inner
         list is an \emph{enumerated} list.
         \begin{enumerate}
            \item This is the first item of an enumerated
                  list that is nested within the itemized list.

            \item This is the second item of the inner list.
                  \LaTeX\ allows you to nest lists deeper than
                  you really should.
         \end{enumerate}
         This is the rest of the second item of the outer
         list.  It is no more interesting than any other
         part of the item.
   \item This is the third item of the list.
\end{itemize}
You can even display poetry.
\begin{verse}
   There is an environment
    for verse \\             % The \\ command separates lines
   Whose features some poets % within a stanza.
   will curse.

                             % One or more blank lines separate stanzas.

   For instead of making\\
   Them do \emph{all} line breaking, \\
   It allows them to put too many words on a line when they'd rather be
   forced to be terse.
\end{verse}

Mathematical formulas may also be displayed.  A
displayed formula
is
one-line long; multiline
formulas require special formatting instructions.
   \[  \ip{\Gamma}{\psi'} = x'' + y^{2} + z_{i}^{n}\]
Don't start a paragraph with a displayed equation,
nor make one a paragraph by itself.

\end{document}               % End of document.

\section{Bench Huge Section 232}
% synthetic bench section 232
% This is a sample LaTeX input file.  (Version of 12 August 2004.)
%
% A '%' character causes TeX to ignore all remaining text on the line,
% and is used for comments like this one.

\documentclass{article}      % Specifies the document class

                             % The preamble begins here.
\title{An Example Document}  % Declares the document's title.
\author{Leslie Lamport}      % Declares the author's name.
\date{January 21, 1994}      % Deleting this command produces today's date.

\newcommand{\ip}[2]{(#1, #2)}
                             % Defines \ip{arg1}{arg2} to mean
                             % (arg1, arg2).

%\newcommand{\ip}[2]{\langle #1 | #2\rangle}
                             % This is an alternative definition of
                             % \ip that is commented out.

\begin{document}             % End of preamble and beginning of text.

\maketitle                   % Produces the title.

This is an example input file.  Comparing it with
the output it generates can show you how to
produce a simple document of your own.

\section{Ordinary Text}      % Produces section heading.  Lower-level
                             % sections are begun with similar
                             % \subsection and \subsubsection commands.

The ends  of words and sentences are marked
  by   spaces. It  doesn't matter how many
spaces    you type; one is as good as 100.  The
end of   a line counts as a space.

One   or more   blank lines denote the  end
of  a paragraph.

Since any number of consecutive spaces are treated
like a single one, the formatting of the input
file makes no difference to
      \LaTeX,                % The \LaTeX command generates the LaTeX logo.
but it makes a difference to you.  When you use
\LaTeX, making your input file as easy to read
as possible will be a great help as you write
your document and when you change it.  This sample
file shows how you can add comments to your own input
file.

Because printing is different from typewriting,
there are a number of things that you have to do
differently when preparing an input file than if
you were just typing the document directly.
Quotation marks like
       ``this''
have to be handled specially, as do quotes within
quotes:
       ``\,`this'            % \, separates the double and single quote.
        is what I just
        wrote, not  `that'\,''.

Dashes come in three sizes: an
       intra-word
dash, a medium dash for number ranges like
       1--2,
and a punctuation
       dash---like
this.

A sentence-ending space should be larger than the
space between words within a sentence.  You
sometimes have to type special commands in
conjunction with punctuation characters to get
this right, as in the following sentence.
       Gnats, gnus, etc.\ all  % `\ ' makes an inter-word space.
       begin with G\@.         % \@ marks end-of-sentence punctuation.
You should check the spaces after periods when
reading your output to make sure you haven't
forgotten any special cases.  Generating an
ellipsis
       \ldots\               % `\ ' is needed after `\ldots' because TeX
                             % ignores spaces after command names like \ldots
                             % made from \ + letters.
                             %
                             % Note how a `%' character causes TeX to ignore
                             % the end of the input line, so these blank lines
                             % do not start a new paragraph.
                             %
with the right spacing around the periods requires
a special command.

\LaTeX\ interprets some common characters as
commands, so you must type special commands to
generate them.  These characters include the
following:
       \$ \& \% \# \{ and \}.

In printing, text is usually emphasized with an
       \emph{italic}
type style.

\begin{em}
   A long segment of text can also be emphasized
   in this way.  Text within such a segment can be
   given \emph{additional} emphasis.
\end{em}

It is sometimes necessary to prevent \LaTeX\ from
breaking a line where it might otherwise do so.
This may be at a space, as between the ``Mr.''\ and
``Jones'' in
       ``Mr.~Jones'',        % ~ produces an unbreakable interword space.
or within a word---especially when the word is a
symbol like
       \mbox{\emph{itemnum}}
that makes little sense when hyphenated across
lines.

Footnotes\footnote{This is an example of a footnote.}
pose no problem.

\LaTeX\ is good at typesetting mathematical formulas
like
       \( x-3y + z = 7 \)
or
       \( a_{1} > x^{2n} + y^{2n} > x' \)
or
       \( \ip{A}{B} = \sum_{i} a_{i} b_{i} \).
The spaces you type in a formula are
ignored.  Remember that a letter like
       $x$                   % $ ... $  and  \( ... \)  are equivalent
is a formula when it denotes a mathematical
symbol, and it should be typed as one.

\section{Displayed Text}

Text is displayed by indenting it from the left
margin.  Quotations are commonly displayed.  There
are short quotations
\begin{quote}
   This is a short quotation.  It consists of a
   single paragraph of text.  See how it is formatted.
\end{quote}
and longer ones.
\begin{quotation}
   This is a longer quotation.  It consists of two
   paragraphs of text, neither of which are
   particularly interesting.

   This is the second paragraph of the quotation.  It
   is just as dull as the first paragraph.
\end{quotation}
Another frequently-displayed structure is a list.
The following is an example of an \emph{itemized}
list.
\begin{itemize}
   \item This is the first item of an itemized list.
         Each item in the list is marked with a ``tick''.
         You don't have to worry about what kind of tick
         mark is used.

   \item This is the second item of the list.  It
         contains another list nested inside it.  The inner
         list is an \emph{enumerated} list.
         \begin{enumerate}
            \item This is the first item of an enumerated
                  list that is nested within the itemized list.

            \item This is the second item of the inner list.
                  \LaTeX\ allows you to nest lists deeper than
                  you really should.
         \end{enumerate}
         This is the rest of the second item of the outer
         list.  It is no more interesting than any other
         part of the item.
   \item This is the third item of the list.
\end{itemize}
You can even display poetry.
\begin{verse}
   There is an environment
    for verse \\             % The \\ command separates lines
   Whose features some poets % within a stanza.
   will curse.

                             % One or more blank lines separate stanzas.

   For instead of making\\
   Them do \emph{all} line breaking, \\
   It allows them to put too many words on a line when they'd rather be
   forced to be terse.
\end{verse}

Mathematical formulas may also be displayed.  A
displayed formula
is
one-line long; multiline
formulas require special formatting instructions.
   \[  \ip{\Gamma}{\psi'} = x'' + y^{2} + z_{i}^{n}\]
Don't start a paragraph with a displayed equation,
nor make one a paragraph by itself.

\end{document}               % End of document.

\section{Bench Huge Section 233}
% synthetic bench section 233
% This is a sample LaTeX input file.  (Version of 12 August 2004.)
%
% A '%' character causes TeX to ignore all remaining text on the line,
% and is used for comments like this one.

\documentclass{article}      % Specifies the document class

                             % The preamble begins here.
\title{An Example Document}  % Declares the document's title.
\author{Leslie Lamport}      % Declares the author's name.
\date{January 21, 1994}      % Deleting this command produces today's date.

\newcommand{\ip}[2]{(#1, #2)}
                             % Defines \ip{arg1}{arg2} to mean
                             % (arg1, arg2).

%\newcommand{\ip}[2]{\langle #1 | #2\rangle}
                             % This is an alternative definition of
                             % \ip that is commented out.

\begin{document}             % End of preamble and beginning of text.

\maketitle                   % Produces the title.

This is an example input file.  Comparing it with
the output it generates can show you how to
produce a simple document of your own.

\section{Ordinary Text}      % Produces section heading.  Lower-level
                             % sections are begun with similar
                             % \subsection and \subsubsection commands.

The ends  of words and sentences are marked
  by   spaces. It  doesn't matter how many
spaces    you type; one is as good as 100.  The
end of   a line counts as a space.

One   or more   blank lines denote the  end
of  a paragraph.

Since any number of consecutive spaces are treated
like a single one, the formatting of the input
file makes no difference to
      \LaTeX,                % The \LaTeX command generates the LaTeX logo.
but it makes a difference to you.  When you use
\LaTeX, making your input file as easy to read
as possible will be a great help as you write
your document and when you change it.  This sample
file shows how you can add comments to your own input
file.

Because printing is different from typewriting,
there are a number of things that you have to do
differently when preparing an input file than if
you were just typing the document directly.
Quotation marks like
       ``this''
have to be handled specially, as do quotes within
quotes:
       ``\,`this'            % \, separates the double and single quote.
        is what I just
        wrote, not  `that'\,''.

Dashes come in three sizes: an
       intra-word
dash, a medium dash for number ranges like
       1--2,
and a punctuation
       dash---like
this.

A sentence-ending space should be larger than the
space between words within a sentence.  You
sometimes have to type special commands in
conjunction with punctuation characters to get
this right, as in the following sentence.
       Gnats, gnus, etc.\ all  % `\ ' makes an inter-word space.
       begin with G\@.         % \@ marks end-of-sentence punctuation.
You should check the spaces after periods when
reading your output to make sure you haven't
forgotten any special cases.  Generating an
ellipsis
       \ldots\               % `\ ' is needed after `\ldots' because TeX
                             % ignores spaces after command names like \ldots
                             % made from \ + letters.
                             %
                             % Note how a `%' character causes TeX to ignore
                             % the end of the input line, so these blank lines
                             % do not start a new paragraph.
                             %
with the right spacing around the periods requires
a special command.

\LaTeX\ interprets some common characters as
commands, so you must type special commands to
generate them.  These characters include the
following:
       \$ \& \% \# \{ and \}.

In printing, text is usually emphasized with an
       \emph{italic}
type style.

\begin{em}
   A long segment of text can also be emphasized
   in this way.  Text within such a segment can be
   given \emph{additional} emphasis.
\end{em}

It is sometimes necessary to prevent \LaTeX\ from
breaking a line where it might otherwise do so.
This may be at a space, as between the ``Mr.''\ and
``Jones'' in
       ``Mr.~Jones'',        % ~ produces an unbreakable interword space.
or within a word---especially when the word is a
symbol like
       \mbox{\emph{itemnum}}
that makes little sense when hyphenated across
lines.

Footnotes\footnote{This is an example of a footnote.}
pose no problem.

\LaTeX\ is good at typesetting mathematical formulas
like
       \( x-3y + z = 7 \)
or
       \( a_{1} > x^{2n} + y^{2n} > x' \)
or
       \( \ip{A}{B} = \sum_{i} a_{i} b_{i} \).
The spaces you type in a formula are
ignored.  Remember that a letter like
       $x$                   % $ ... $  and  \( ... \)  are equivalent
is a formula when it denotes a mathematical
symbol, and it should be typed as one.

\section{Displayed Text}

Text is displayed by indenting it from the left
margin.  Quotations are commonly displayed.  There
are short quotations
\begin{quote}
   This is a short quotation.  It consists of a
   single paragraph of text.  See how it is formatted.
\end{quote}
and longer ones.
\begin{quotation}
   This is a longer quotation.  It consists of two
   paragraphs of text, neither of which are
   particularly interesting.

   This is the second paragraph of the quotation.  It
   is just as dull as the first paragraph.
\end{quotation}
Another frequently-displayed structure is a list.
The following is an example of an \emph{itemized}
list.
\begin{itemize}
   \item This is the first item of an itemized list.
         Each item in the list is marked with a ``tick''.
         You don't have to worry about what kind of tick
         mark is used.

   \item This is the second item of the list.  It
         contains another list nested inside it.  The inner
         list is an \emph{enumerated} list.
         \begin{enumerate}
            \item This is the first item of an enumerated
                  list that is nested within the itemized list.

            \item This is the second item of the inner list.
                  \LaTeX\ allows you to nest lists deeper than
                  you really should.
         \end{enumerate}
         This is the rest of the second item of the outer
         list.  It is no more interesting than any other
         part of the item.
   \item This is the third item of the list.
\end{itemize}
You can even display poetry.
\begin{verse}
   There is an environment
    for verse \\             % The \\ command separates lines
   Whose features some poets % within a stanza.
   will curse.

                             % One or more blank lines separate stanzas.

   For instead of making\\
   Them do \emph{all} line breaking, \\
   It allows them to put too many words on a line when they'd rather be
   forced to be terse.
\end{verse}

Mathematical formulas may also be displayed.  A
displayed formula
is
one-line long; multiline
formulas require special formatting instructions.
   \[  \ip{\Gamma}{\psi'} = x'' + y^{2} + z_{i}^{n}\]
Don't start a paragraph with a displayed equation,
nor make one a paragraph by itself.

\end{document}               % End of document.

\section{Bench Huge Section 234}
% synthetic bench section 234
% This is a sample LaTeX input file.  (Version of 12 August 2004.)
%
% A '%' character causes TeX to ignore all remaining text on the line,
% and is used for comments like this one.

\documentclass{article}      % Specifies the document class

                             % The preamble begins here.
\title{An Example Document}  % Declares the document's title.
\author{Leslie Lamport}      % Declares the author's name.
\date{January 21, 1994}      % Deleting this command produces today's date.

\newcommand{\ip}[2]{(#1, #2)}
                             % Defines \ip{arg1}{arg2} to mean
                             % (arg1, arg2).

%\newcommand{\ip}[2]{\langle #1 | #2\rangle}
                             % This is an alternative definition of
                             % \ip that is commented out.

\begin{document}             % End of preamble and beginning of text.

\maketitle                   % Produces the title.

This is an example input file.  Comparing it with
the output it generates can show you how to
produce a simple document of your own.

\section{Ordinary Text}      % Produces section heading.  Lower-level
                             % sections are begun with similar
                             % \subsection and \subsubsection commands.

The ends  of words and sentences are marked
  by   spaces. It  doesn't matter how many
spaces    you type; one is as good as 100.  The
end of   a line counts as a space.

One   or more   blank lines denote the  end
of  a paragraph.

Since any number of consecutive spaces are treated
like a single one, the formatting of the input
file makes no difference to
      \LaTeX,                % The \LaTeX command generates the LaTeX logo.
but it makes a difference to you.  When you use
\LaTeX, making your input file as easy to read
as possible will be a great help as you write
your document and when you change it.  This sample
file shows how you can add comments to your own input
file.

Because printing is different from typewriting,
there are a number of things that you have to do
differently when preparing an input file than if
you were just typing the document directly.
Quotation marks like
       ``this''
have to be handled specially, as do quotes within
quotes:
       ``\,`this'            % \, separates the double and single quote.
        is what I just
        wrote, not  `that'\,''.

Dashes come in three sizes: an
       intra-word
dash, a medium dash for number ranges like
       1--2,
and a punctuation
       dash---like
this.

A sentence-ending space should be larger than the
space between words within a sentence.  You
sometimes have to type special commands in
conjunction with punctuation characters to get
this right, as in the following sentence.
       Gnats, gnus, etc.\ all  % `\ ' makes an inter-word space.
       begin with G\@.         % \@ marks end-of-sentence punctuation.
You should check the spaces after periods when
reading your output to make sure you haven't
forgotten any special cases.  Generating an
ellipsis
       \ldots\               % `\ ' is needed after `\ldots' because TeX
                             % ignores spaces after command names like \ldots
                             % made from \ + letters.
                             %
                             % Note how a `%' character causes TeX to ignore
                             % the end of the input line, so these blank lines
                             % do not start a new paragraph.
                             %
with the right spacing around the periods requires
a special command.

\LaTeX\ interprets some common characters as
commands, so you must type special commands to
generate them.  These characters include the
following:
       \$ \& \% \# \{ and \}.

In printing, text is usually emphasized with an
       \emph{italic}
type style.

\begin{em}
   A long segment of text can also be emphasized
   in this way.  Text within such a segment can be
   given \emph{additional} emphasis.
\end{em}

It is sometimes necessary to prevent \LaTeX\ from
breaking a line where it might otherwise do so.
This may be at a space, as between the ``Mr.''\ and
``Jones'' in
       ``Mr.~Jones'',        % ~ produces an unbreakable interword space.
or within a word---especially when the word is a
symbol like
       \mbox{\emph{itemnum}}
that makes little sense when hyphenated across
lines.

Footnotes\footnote{This is an example of a footnote.}
pose no problem.

\LaTeX\ is good at typesetting mathematical formulas
like
       \( x-3y + z = 7 \)
or
       \( a_{1} > x^{2n} + y^{2n} > x' \)
or
       \( \ip{A}{B} = \sum_{i} a_{i} b_{i} \).
The spaces you type in a formula are
ignored.  Remember that a letter like
       $x$                   % $ ... $  and  \( ... \)  are equivalent
is a formula when it denotes a mathematical
symbol, and it should be typed as one.

\section{Displayed Text}

Text is displayed by indenting it from the left
margin.  Quotations are commonly displayed.  There
are short quotations
\begin{quote}
   This is a short quotation.  It consists of a
   single paragraph of text.  See how it is formatted.
\end{quote}
and longer ones.
\begin{quotation}
   This is a longer quotation.  It consists of two
   paragraphs of text, neither of which are
   particularly interesting.

   This is the second paragraph of the quotation.  It
   is just as dull as the first paragraph.
\end{quotation}
Another frequently-displayed structure is a list.
The following is an example of an \emph{itemized}
list.
\begin{itemize}
   \item This is the first item of an itemized list.
         Each item in the list is marked with a ``tick''.
         You don't have to worry about what kind of tick
         mark is used.

   \item This is the second item of the list.  It
         contains another list nested inside it.  The inner
         list is an \emph{enumerated} list.
         \begin{enumerate}
            \item This is the first item of an enumerated
                  list that is nested within the itemized list.

            \item This is the second item of the inner list.
                  \LaTeX\ allows you to nest lists deeper than
                  you really should.
         \end{enumerate}
         This is the rest of the second item of the outer
         list.  It is no more interesting than any other
         part of the item.
   \item This is the third item of the list.
\end{itemize}
You can even display poetry.
\begin{verse}
   There is an environment
    for verse \\             % The \\ command separates lines
   Whose features some poets % within a stanza.
   will curse.

                             % One or more blank lines separate stanzas.

   For instead of making\\
   Them do \emph{all} line breaking, \\
   It allows them to put too many words on a line when they'd rather be
   forced to be terse.
\end{verse}

Mathematical formulas may also be displayed.  A
displayed formula
is
one-line long; multiline
formulas require special formatting instructions.
   \[  \ip{\Gamma}{\psi'} = x'' + y^{2} + z_{i}^{n}\]
Don't start a paragraph with a displayed equation,
nor make one a paragraph by itself.

\end{document}               % End of document.

\section{Bench Huge Section 235}
% synthetic bench section 235
% This is a sample LaTeX input file.  (Version of 12 August 2004.)
%
% A '%' character causes TeX to ignore all remaining text on the line,
% and is used for comments like this one.

\documentclass{article}      % Specifies the document class

                             % The preamble begins here.
\title{An Example Document}  % Declares the document's title.
\author{Leslie Lamport}      % Declares the author's name.
\date{January 21, 1994}      % Deleting this command produces today's date.

\newcommand{\ip}[2]{(#1, #2)}
                             % Defines \ip{arg1}{arg2} to mean
                             % (arg1, arg2).

%\newcommand{\ip}[2]{\langle #1 | #2\rangle}
                             % This is an alternative definition of
                             % \ip that is commented out.

\begin{document}             % End of preamble and beginning of text.

\maketitle                   % Produces the title.

This is an example input file.  Comparing it with
the output it generates can show you how to
produce a simple document of your own.

\section{Ordinary Text}      % Produces section heading.  Lower-level
                             % sections are begun with similar
                             % \subsection and \subsubsection commands.

The ends  of words and sentences are marked
  by   spaces. It  doesn't matter how many
spaces    you type; one is as good as 100.  The
end of   a line counts as a space.

One   or more   blank lines denote the  end
of  a paragraph.

Since any number of consecutive spaces are treated
like a single one, the formatting of the input
file makes no difference to
      \LaTeX,                % The \LaTeX command generates the LaTeX logo.
but it makes a difference to you.  When you use
\LaTeX, making your input file as easy to read
as possible will be a great help as you write
your document and when you change it.  This sample
file shows how you can add comments to your own input
file.

Because printing is different from typewriting,
there are a number of things that you have to do
differently when preparing an input file than if
you were just typing the document directly.
Quotation marks like
       ``this''
have to be handled specially, as do quotes within
quotes:
       ``\,`this'            % \, separates the double and single quote.
        is what I just
        wrote, not  `that'\,''.

Dashes come in three sizes: an
       intra-word
dash, a medium dash for number ranges like
       1--2,
and a punctuation
       dash---like
this.

A sentence-ending space should be larger than the
space between words within a sentence.  You
sometimes have to type special commands in
conjunction with punctuation characters to get
this right, as in the following sentence.
       Gnats, gnus, etc.\ all  % `\ ' makes an inter-word space.
       begin with G\@.         % \@ marks end-of-sentence punctuation.
You should check the spaces after periods when
reading your output to make sure you haven't
forgotten any special cases.  Generating an
ellipsis
       \ldots\               % `\ ' is needed after `\ldots' because TeX
                             % ignores spaces after command names like \ldots
                             % made from \ + letters.
                             %
                             % Note how a `%' character causes TeX to ignore
                             % the end of the input line, so these blank lines
                             % do not start a new paragraph.
                             %
with the right spacing around the periods requires
a special command.

\LaTeX\ interprets some common characters as
commands, so you must type special commands to
generate them.  These characters include the
following:
       \$ \& \% \# \{ and \}.

In printing, text is usually emphasized with an
       \emph{italic}
type style.

\begin{em}
   A long segment of text can also be emphasized
   in this way.  Text within such a segment can be
   given \emph{additional} emphasis.
\end{em}

It is sometimes necessary to prevent \LaTeX\ from
breaking a line where it might otherwise do so.
This may be at a space, as between the ``Mr.''\ and
``Jones'' in
       ``Mr.~Jones'',        % ~ produces an unbreakable interword space.
or within a word---especially when the word is a
symbol like
       \mbox{\emph{itemnum}}
that makes little sense when hyphenated across
lines.

Footnotes\footnote{This is an example of a footnote.}
pose no problem.

\LaTeX\ is good at typesetting mathematical formulas
like
       \( x-3y + z = 7 \)
or
       \( a_{1} > x^{2n} + y^{2n} > x' \)
or
       \( \ip{A}{B} = \sum_{i} a_{i} b_{i} \).
The spaces you type in a formula are
ignored.  Remember that a letter like
       $x$                   % $ ... $  and  \( ... \)  are equivalent
is a formula when it denotes a mathematical
symbol, and it should be typed as one.

\section{Displayed Text}

Text is displayed by indenting it from the left
margin.  Quotations are commonly displayed.  There
are short quotations
\begin{quote}
   This is a short quotation.  It consists of a
   single paragraph of text.  See how it is formatted.
\end{quote}
and longer ones.
\begin{quotation}
   This is a longer quotation.  It consists of two
   paragraphs of text, neither of which are
   particularly interesting.

   This is the second paragraph of the quotation.  It
   is just as dull as the first paragraph.
\end{quotation}
Another frequently-displayed structure is a list.
The following is an example of an \emph{itemized}
list.
\begin{itemize}
   \item This is the first item of an itemized list.
         Each item in the list is marked with a ``tick''.
         You don't have to worry about what kind of tick
         mark is used.

   \item This is the second item of the list.  It
         contains another list nested inside it.  The inner
         list is an \emph{enumerated} list.
         \begin{enumerate}
            \item This is the first item of an enumerated
                  list that is nested within the itemized list.

            \item This is the second item of the inner list.
                  \LaTeX\ allows you to nest lists deeper than
                  you really should.
         \end{enumerate}
         This is the rest of the second item of the outer
         list.  It is no more interesting than any other
         part of the item.
   \item This is the third item of the list.
\end{itemize}
You can even display poetry.
\begin{verse}
   There is an environment
    for verse \\             % The \\ command separates lines
   Whose features some poets % within a stanza.
   will curse.

                             % One or more blank lines separate stanzas.

   For instead of making\\
   Them do \emph{all} line breaking, \\
   It allows them to put too many words on a line when they'd rather be
   forced to be terse.
\end{verse}

Mathematical formulas may also be displayed.  A
displayed formula
is
one-line long; multiline
formulas require special formatting instructions.
   \[  \ip{\Gamma}{\psi'} = x'' + y^{2} + z_{i}^{n}\]
Don't start a paragraph with a displayed equation,
nor make one a paragraph by itself.

\end{document}               % End of document.

\section{Bench Huge Section 236}
% synthetic bench section 236
% This is a sample LaTeX input file.  (Version of 12 August 2004.)
%
% A '%' character causes TeX to ignore all remaining text on the line,
% and is used for comments like this one.

\documentclass{article}      % Specifies the document class

                             % The preamble begins here.
\title{An Example Document}  % Declares the document's title.
\author{Leslie Lamport}      % Declares the author's name.
\date{January 21, 1994}      % Deleting this command produces today's date.

\newcommand{\ip}[2]{(#1, #2)}
                             % Defines \ip{arg1}{arg2} to mean
                             % (arg1, arg2).

%\newcommand{\ip}[2]{\langle #1 | #2\rangle}
                             % This is an alternative definition of
                             % \ip that is commented out.

\begin{document}             % End of preamble and beginning of text.

\maketitle                   % Produces the title.

This is an example input file.  Comparing it with
the output it generates can show you how to
produce a simple document of your own.

\section{Ordinary Text}      % Produces section heading.  Lower-level
                             % sections are begun with similar
                             % \subsection and \subsubsection commands.

The ends  of words and sentences are marked
  by   spaces. It  doesn't matter how many
spaces    you type; one is as good as 100.  The
end of   a line counts as a space.

One   or more   blank lines denote the  end
of  a paragraph.

Since any number of consecutive spaces are treated
like a single one, the formatting of the input
file makes no difference to
      \LaTeX,                % The \LaTeX command generates the LaTeX logo.
but it makes a difference to you.  When you use
\LaTeX, making your input file as easy to read
as possible will be a great help as you write
your document and when you change it.  This sample
file shows how you can add comments to your own input
file.

Because printing is different from typewriting,
there are a number of things that you have to do
differently when preparing an input file than if
you were just typing the document directly.
Quotation marks like
       ``this''
have to be handled specially, as do quotes within
quotes:
       ``\,`this'            % \, separates the double and single quote.
        is what I just
        wrote, not  `that'\,''.

Dashes come in three sizes: an
       intra-word
dash, a medium dash for number ranges like
       1--2,
and a punctuation
       dash---like
this.

A sentence-ending space should be larger than the
space between words within a sentence.  You
sometimes have to type special commands in
conjunction with punctuation characters to get
this right, as in the following sentence.
       Gnats, gnus, etc.\ all  % `\ ' makes an inter-word space.
       begin with G\@.         % \@ marks end-of-sentence punctuation.
You should check the spaces after periods when
reading your output to make sure you haven't
forgotten any special cases.  Generating an
ellipsis
       \ldots\               % `\ ' is needed after `\ldots' because TeX
                             % ignores spaces after command names like \ldots
                             % made from \ + letters.
                             %
                             % Note how a `%' character causes TeX to ignore
                             % the end of the input line, so these blank lines
                             % do not start a new paragraph.
                             %
with the right spacing around the periods requires
a special command.

\LaTeX\ interprets some common characters as
commands, so you must type special commands to
generate them.  These characters include the
following:
       \$ \& \% \# \{ and \}.

In printing, text is usually emphasized with an
       \emph{italic}
type style.

\begin{em}
   A long segment of text can also be emphasized
   in this way.  Text within such a segment can be
   given \emph{additional} emphasis.
\end{em}

It is sometimes necessary to prevent \LaTeX\ from
breaking a line where it might otherwise do so.
This may be at a space, as between the ``Mr.''\ and
``Jones'' in
       ``Mr.~Jones'',        % ~ produces an unbreakable interword space.
or within a word---especially when the word is a
symbol like
       \mbox{\emph{itemnum}}
that makes little sense when hyphenated across
lines.

Footnotes\footnote{This is an example of a footnote.}
pose no problem.

\LaTeX\ is good at typesetting mathematical formulas
like
       \( x-3y + z = 7 \)
or
       \( a_{1} > x^{2n} + y^{2n} > x' \)
or
       \( \ip{A}{B} = \sum_{i} a_{i} b_{i} \).
The spaces you type in a formula are
ignored.  Remember that a letter like
       $x$                   % $ ... $  and  \( ... \)  are equivalent
is a formula when it denotes a mathematical
symbol, and it should be typed as one.

\section{Displayed Text}

Text is displayed by indenting it from the left
margin.  Quotations are commonly displayed.  There
are short quotations
\begin{quote}
   This is a short quotation.  It consists of a
   single paragraph of text.  See how it is formatted.
\end{quote}
and longer ones.
\begin{quotation}
   This is a longer quotation.  It consists of two
   paragraphs of text, neither of which are
   particularly interesting.

   This is the second paragraph of the quotation.  It
   is just as dull as the first paragraph.
\end{quotation}
Another frequently-displayed structure is a list.
The following is an example of an \emph{itemized}
list.
\begin{itemize}
   \item This is the first item of an itemized list.
         Each item in the list is marked with a ``tick''.
         You don't have to worry about what kind of tick
         mark is used.

   \item This is the second item of the list.  It
         contains another list nested inside it.  The inner
         list is an \emph{enumerated} list.
         \begin{enumerate}
            \item This is the first item of an enumerated
                  list that is nested within the itemized list.

            \item This is the second item of the inner list.
                  \LaTeX\ allows you to nest lists deeper than
                  you really should.
         \end{enumerate}
         This is the rest of the second item of the outer
         list.  It is no more interesting than any other
         part of the item.
   \item This is the third item of the list.
\end{itemize}
You can even display poetry.
\begin{verse}
   There is an environment
    for verse \\             % The \\ command separates lines
   Whose features some poets % within a stanza.
   will curse.

                             % One or more blank lines separate stanzas.

   For instead of making\\
   Them do \emph{all} line breaking, \\
   It allows them to put too many words on a line when they'd rather be
   forced to be terse.
\end{verse}

Mathematical formulas may also be displayed.  A
displayed formula
is
one-line long; multiline
formulas require special formatting instructions.
   \[  \ip{\Gamma}{\psi'} = x'' + y^{2} + z_{i}^{n}\]
Don't start a paragraph with a displayed equation,
nor make one a paragraph by itself.

\end{document}               % End of document.

\section{Bench Huge Section 237}
% synthetic bench section 237
% This is a sample LaTeX input file.  (Version of 12 August 2004.)
%
% A '%' character causes TeX to ignore all remaining text on the line,
% and is used for comments like this one.

\documentclass{article}      % Specifies the document class

                             % The preamble begins here.
\title{An Example Document}  % Declares the document's title.
\author{Leslie Lamport}      % Declares the author's name.
\date{January 21, 1994}      % Deleting this command produces today's date.

\newcommand{\ip}[2]{(#1, #2)}
                             % Defines \ip{arg1}{arg2} to mean
                             % (arg1, arg2).

%\newcommand{\ip}[2]{\langle #1 | #2\rangle}
                             % This is an alternative definition of
                             % \ip that is commented out.

\begin{document}             % End of preamble and beginning of text.

\maketitle                   % Produces the title.

This is an example input file.  Comparing it with
the output it generates can show you how to
produce a simple document of your own.

\section{Ordinary Text}      % Produces section heading.  Lower-level
                             % sections are begun with similar
                             % \subsection and \subsubsection commands.

The ends  of words and sentences are marked
  by   spaces. It  doesn't matter how many
spaces    you type; one is as good as 100.  The
end of   a line counts as a space.

One   or more   blank lines denote the  end
of  a paragraph.

Since any number of consecutive spaces are treated
like a single one, the formatting of the input
file makes no difference to
      \LaTeX,                % The \LaTeX command generates the LaTeX logo.
but it makes a difference to you.  When you use
\LaTeX, making your input file as easy to read
as possible will be a great help as you write
your document and when you change it.  This sample
file shows how you can add comments to your own input
file.

Because printing is different from typewriting,
there are a number of things that you have to do
differently when preparing an input file than if
you were just typing the document directly.
Quotation marks like
       ``this''
have to be handled specially, as do quotes within
quotes:
       ``\,`this'            % \, separates the double and single quote.
        is what I just
        wrote, not  `that'\,''.

Dashes come in three sizes: an
       intra-word
dash, a medium dash for number ranges like
       1--2,
and a punctuation
       dash---like
this.

A sentence-ending space should be larger than the
space between words within a sentence.  You
sometimes have to type special commands in
conjunction with punctuation characters to get
this right, as in the following sentence.
       Gnats, gnus, etc.\ all  % `\ ' makes an inter-word space.
       begin with G\@.         % \@ marks end-of-sentence punctuation.
You should check the spaces after periods when
reading your output to make sure you haven't
forgotten any special cases.  Generating an
ellipsis
       \ldots\               % `\ ' is needed after `\ldots' because TeX
                             % ignores spaces after command names like \ldots
                             % made from \ + letters.
                             %
                             % Note how a `%' character causes TeX to ignore
                             % the end of the input line, so these blank lines
                             % do not start a new paragraph.
                             %
with the right spacing around the periods requires
a special command.

\LaTeX\ interprets some common characters as
commands, so you must type special commands to
generate them.  These characters include the
following:
       \$ \& \% \# \{ and \}.

In printing, text is usually emphasized with an
       \emph{italic}
type style.

\begin{em}
   A long segment of text can also be emphasized
   in this way.  Text within such a segment can be
   given \emph{additional} emphasis.
\end{em}

It is sometimes necessary to prevent \LaTeX\ from
breaking a line where it might otherwise do so.
This may be at a space, as between the ``Mr.''\ and
``Jones'' in
       ``Mr.~Jones'',        % ~ produces an unbreakable interword space.
or within a word---especially when the word is a
symbol like
       \mbox{\emph{itemnum}}
that makes little sense when hyphenated across
lines.

Footnotes\footnote{This is an example of a footnote.}
pose no problem.

\LaTeX\ is good at typesetting mathematical formulas
like
       \( x-3y + z = 7 \)
or
       \( a_{1} > x^{2n} + y^{2n} > x' \)
or
       \( \ip{A}{B} = \sum_{i} a_{i} b_{i} \).
The spaces you type in a formula are
ignored.  Remember that a letter like
       $x$                   % $ ... $  and  \( ... \)  are equivalent
is a formula when it denotes a mathematical
symbol, and it should be typed as one.

\section{Displayed Text}

Text is displayed by indenting it from the left
margin.  Quotations are commonly displayed.  There
are short quotations
\begin{quote}
   This is a short quotation.  It consists of a
   single paragraph of text.  See how it is formatted.
\end{quote}
and longer ones.
\begin{quotation}
   This is a longer quotation.  It consists of two
   paragraphs of text, neither of which are
   particularly interesting.

   This is the second paragraph of the quotation.  It
   is just as dull as the first paragraph.
\end{quotation}
Another frequently-displayed structure is a list.
The following is an example of an \emph{itemized}
list.
\begin{itemize}
   \item This is the first item of an itemized list.
         Each item in the list is marked with a ``tick''.
         You don't have to worry about what kind of tick
         mark is used.

   \item This is the second item of the list.  It
         contains another list nested inside it.  The inner
         list is an \emph{enumerated} list.
         \begin{enumerate}
            \item This is the first item of an enumerated
                  list that is nested within the itemized list.

            \item This is the second item of the inner list.
                  \LaTeX\ allows you to nest lists deeper than
                  you really should.
         \end{enumerate}
         This is the rest of the second item of the outer
         list.  It is no more interesting than any other
         part of the item.
   \item This is the third item of the list.
\end{itemize}
You can even display poetry.
\begin{verse}
   There is an environment
    for verse \\             % The \\ command separates lines
   Whose features some poets % within a stanza.
   will curse.

                             % One or more blank lines separate stanzas.

   For instead of making\\
   Them do \emph{all} line breaking, \\
   It allows them to put too many words on a line when they'd rather be
   forced to be terse.
\end{verse}

Mathematical formulas may also be displayed.  A
displayed formula
is
one-line long; multiline
formulas require special formatting instructions.
   \[  \ip{\Gamma}{\psi'} = x'' + y^{2} + z_{i}^{n}\]
Don't start a paragraph with a displayed equation,
nor make one a paragraph by itself.

\end{document}               % End of document.

\section{Bench Huge Section 238}
% synthetic bench section 238
% This is a sample LaTeX input file.  (Version of 12 August 2004.)
%
% A '%' character causes TeX to ignore all remaining text on the line,
% and is used for comments like this one.

\documentclass{article}      % Specifies the document class

                             % The preamble begins here.
\title{An Example Document}  % Declares the document's title.
\author{Leslie Lamport}      % Declares the author's name.
\date{January 21, 1994}      % Deleting this command produces today's date.

\newcommand{\ip}[2]{(#1, #2)}
                             % Defines \ip{arg1}{arg2} to mean
                             % (arg1, arg2).

%\newcommand{\ip}[2]{\langle #1 | #2\rangle}
                             % This is an alternative definition of
                             % \ip that is commented out.

\begin{document}             % End of preamble and beginning of text.

\maketitle                   % Produces the title.

This is an example input file.  Comparing it with
the output it generates can show you how to
produce a simple document of your own.

\section{Ordinary Text}      % Produces section heading.  Lower-level
                             % sections are begun with similar
                             % \subsection and \subsubsection commands.

The ends  of words and sentences are marked
  by   spaces. It  doesn't matter how many
spaces    you type; one is as good as 100.  The
end of   a line counts as a space.

One   or more   blank lines denote the  end
of  a paragraph.

Since any number of consecutive spaces are treated
like a single one, the formatting of the input
file makes no difference to
      \LaTeX,                % The \LaTeX command generates the LaTeX logo.
but it makes a difference to you.  When you use
\LaTeX, making your input file as easy to read
as possible will be a great help as you write
your document and when you change it.  This sample
file shows how you can add comments to your own input
file.

Because printing is different from typewriting,
there are a number of things that you have to do
differently when preparing an input file than if
you were just typing the document directly.
Quotation marks like
       ``this''
have to be handled specially, as do quotes within
quotes:
       ``\,`this'            % \, separates the double and single quote.
        is what I just
        wrote, not  `that'\,''.

Dashes come in three sizes: an
       intra-word
dash, a medium dash for number ranges like
       1--2,
and a punctuation
       dash---like
this.

A sentence-ending space should be larger than the
space between words within a sentence.  You
sometimes have to type special commands in
conjunction with punctuation characters to get
this right, as in the following sentence.
       Gnats, gnus, etc.\ all  % `\ ' makes an inter-word space.
       begin with G\@.         % \@ marks end-of-sentence punctuation.
You should check the spaces after periods when
reading your output to make sure you haven't
forgotten any special cases.  Generating an
ellipsis
       \ldots\               % `\ ' is needed after `\ldots' because TeX
                             % ignores spaces after command names like \ldots
                             % made from \ + letters.
                             %
                             % Note how a `%' character causes TeX to ignore
                             % the end of the input line, so these blank lines
                             % do not start a new paragraph.
                             %
with the right spacing around the periods requires
a special command.

\LaTeX\ interprets some common characters as
commands, so you must type special commands to
generate them.  These characters include the
following:
       \$ \& \% \# \{ and \}.

In printing, text is usually emphasized with an
       \emph{italic}
type style.

\begin{em}
   A long segment of text can also be emphasized
   in this way.  Text within such a segment can be
   given \emph{additional} emphasis.
\end{em}

It is sometimes necessary to prevent \LaTeX\ from
breaking a line where it might otherwise do so.
This may be at a space, as between the ``Mr.''\ and
``Jones'' in
       ``Mr.~Jones'',        % ~ produces an unbreakable interword space.
or within a word---especially when the word is a
symbol like
       \mbox{\emph{itemnum}}
that makes little sense when hyphenated across
lines.

Footnotes\footnote{This is an example of a footnote.}
pose no problem.

\LaTeX\ is good at typesetting mathematical formulas
like
       \( x-3y + z = 7 \)
or
       \( a_{1} > x^{2n} + y^{2n} > x' \)
or
       \( \ip{A}{B} = \sum_{i} a_{i} b_{i} \).
The spaces you type in a formula are
ignored.  Remember that a letter like
       $x$                   % $ ... $  and  \( ... \)  are equivalent
is a formula when it denotes a mathematical
symbol, and it should be typed as one.

\section{Displayed Text}

Text is displayed by indenting it from the left
margin.  Quotations are commonly displayed.  There
are short quotations
\begin{quote}
   This is a short quotation.  It consists of a
   single paragraph of text.  See how it is formatted.
\end{quote}
and longer ones.
\begin{quotation}
   This is a longer quotation.  It consists of two
   paragraphs of text, neither of which are
   particularly interesting.

   This is the second paragraph of the quotation.  It
   is just as dull as the first paragraph.
\end{quotation}
Another frequently-displayed structure is a list.
The following is an example of an \emph{itemized}
list.
\begin{itemize}
   \item This is the first item of an itemized list.
         Each item in the list is marked with a ``tick''.
         You don't have to worry about what kind of tick
         mark is used.

   \item This is the second item of the list.  It
         contains another list nested inside it.  The inner
         list is an \emph{enumerated} list.
         \begin{enumerate}
            \item This is the first item of an enumerated
                  list that is nested within the itemized list.

            \item This is the second item of the inner list.
                  \LaTeX\ allows you to nest lists deeper than
                  you really should.
         \end{enumerate}
         This is the rest of the second item of the outer
         list.  It is no more interesting than any other
         part of the item.
   \item This is the third item of the list.
\end{itemize}
You can even display poetry.
\begin{verse}
   There is an environment
    for verse \\             % The \\ command separates lines
   Whose features some poets % within a stanza.
   will curse.

                             % One or more blank lines separate stanzas.

   For instead of making\\
   Them do \emph{all} line breaking, \\
   It allows them to put too many words on a line when they'd rather be
   forced to be terse.
\end{verse}

Mathematical formulas may also be displayed.  A
displayed formula
is
one-line long; multiline
formulas require special formatting instructions.
   \[  \ip{\Gamma}{\psi'} = x'' + y^{2} + z_{i}^{n}\]
Don't start a paragraph with a displayed equation,
nor make one a paragraph by itself.

\end{document}               % End of document.

\section{Bench Huge Section 239}
% synthetic bench section 239
% This is a sample LaTeX input file.  (Version of 12 August 2004.)
%
% A '%' character causes TeX to ignore all remaining text on the line,
% and is used for comments like this one.

\documentclass{article}      % Specifies the document class

                             % The preamble begins here.
\title{An Example Document}  % Declares the document's title.
\author{Leslie Lamport}      % Declares the author's name.
\date{January 21, 1994}      % Deleting this command produces today's date.

\newcommand{\ip}[2]{(#1, #2)}
                             % Defines \ip{arg1}{arg2} to mean
                             % (arg1, arg2).

%\newcommand{\ip}[2]{\langle #1 | #2\rangle}
                             % This is an alternative definition of
                             % \ip that is commented out.

\begin{document}             % End of preamble and beginning of text.

\maketitle                   % Produces the title.

This is an example input file.  Comparing it with
the output it generates can show you how to
produce a simple document of your own.

\section{Ordinary Text}      % Produces section heading.  Lower-level
                             % sections are begun with similar
                             % \subsection and \subsubsection commands.

The ends  of words and sentences are marked
  by   spaces. It  doesn't matter how many
spaces    you type; one is as good as 100.  The
end of   a line counts as a space.

One   or more   blank lines denote the  end
of  a paragraph.

Since any number of consecutive spaces are treated
like a single one, the formatting of the input
file makes no difference to
      \LaTeX,                % The \LaTeX command generates the LaTeX logo.
but it makes a difference to you.  When you use
\LaTeX, making your input file as easy to read
as possible will be a great help as you write
your document and when you change it.  This sample
file shows how you can add comments to your own input
file.

Because printing is different from typewriting,
there are a number of things that you have to do
differently when preparing an input file than if
you were just typing the document directly.
Quotation marks like
       ``this''
have to be handled specially, as do quotes within
quotes:
       ``\,`this'            % \, separates the double and single quote.
        is what I just
        wrote, not  `that'\,''.

Dashes come in three sizes: an
       intra-word
dash, a medium dash for number ranges like
       1--2,
and a punctuation
       dash---like
this.

A sentence-ending space should be larger than the
space between words within a sentence.  You
sometimes have to type special commands in
conjunction with punctuation characters to get
this right, as in the following sentence.
       Gnats, gnus, etc.\ all  % `\ ' makes an inter-word space.
       begin with G\@.         % \@ marks end-of-sentence punctuation.
You should check the spaces after periods when
reading your output to make sure you haven't
forgotten any special cases.  Generating an
ellipsis
       \ldots\               % `\ ' is needed after `\ldots' because TeX
                             % ignores spaces after command names like \ldots
                             % made from \ + letters.
                             %
                             % Note how a `%' character causes TeX to ignore
                             % the end of the input line, so these blank lines
                             % do not start a new paragraph.
                             %
with the right spacing around the periods requires
a special command.

\LaTeX\ interprets some common characters as
commands, so you must type special commands to
generate them.  These characters include the
following:
       \$ \& \% \# \{ and \}.

In printing, text is usually emphasized with an
       \emph{italic}
type style.

\begin{em}
   A long segment of text can also be emphasized
   in this way.  Text within such a segment can be
   given \emph{additional} emphasis.
\end{em}

It is sometimes necessary to prevent \LaTeX\ from
breaking a line where it might otherwise do so.
This may be at a space, as between the ``Mr.''\ and
``Jones'' in
       ``Mr.~Jones'',        % ~ produces an unbreakable interword space.
or within a word---especially when the word is a
symbol like
       \mbox{\emph{itemnum}}
that makes little sense when hyphenated across
lines.

Footnotes\footnote{This is an example of a footnote.}
pose no problem.

\LaTeX\ is good at typesetting mathematical formulas
like
       \( x-3y + z = 7 \)
or
       \( a_{1} > x^{2n} + y^{2n} > x' \)
or
       \( \ip{A}{B} = \sum_{i} a_{i} b_{i} \).
The spaces you type in a formula are
ignored.  Remember that a letter like
       $x$                   % $ ... $  and  \( ... \)  are equivalent
is a formula when it denotes a mathematical
symbol, and it should be typed as one.

\section{Displayed Text}

Text is displayed by indenting it from the left
margin.  Quotations are commonly displayed.  There
are short quotations
\begin{quote}
   This is a short quotation.  It consists of a
   single paragraph of text.  See how it is formatted.
\end{quote}
and longer ones.
\begin{quotation}
   This is a longer quotation.  It consists of two
   paragraphs of text, neither of which are
   particularly interesting.

   This is the second paragraph of the quotation.  It
   is just as dull as the first paragraph.
\end{quotation}
Another frequently-displayed structure is a list.
The following is an example of an \emph{itemized}
list.
\begin{itemize}
   \item This is the first item of an itemized list.
         Each item in the list is marked with a ``tick''.
         You don't have to worry about what kind of tick
         mark is used.

   \item This is the second item of the list.  It
         contains another list nested inside it.  The inner
         list is an \emph{enumerated} list.
         \begin{enumerate}
            \item This is the first item of an enumerated
                  list that is nested within the itemized list.

            \item This is the second item of the inner list.
                  \LaTeX\ allows you to nest lists deeper than
                  you really should.
         \end{enumerate}
         This is the rest of the second item of the outer
         list.  It is no more interesting than any other
         part of the item.
   \item This is the third item of the list.
\end{itemize}
You can even display poetry.
\begin{verse}
   There is an environment
    for verse \\             % The \\ command separates lines
   Whose features some poets % within a stanza.
   will curse.

                             % One or more blank lines separate stanzas.

   For instead of making\\
   Them do \emph{all} line breaking, \\
   It allows them to put too many words on a line when they'd rather be
   forced to be terse.
\end{verse}

Mathematical formulas may also be displayed.  A
displayed formula
is
one-line long; multiline
formulas require special formatting instructions.
   \[  \ip{\Gamma}{\psi'} = x'' + y^{2} + z_{i}^{n}\]
Don't start a paragraph with a displayed equation,
nor make one a paragraph by itself.

\end{document}               % End of document.

\section{Bench Huge Section 240}
% synthetic bench section 240
% This is a sample LaTeX input file.  (Version of 12 August 2004.)
%
% A '%' character causes TeX to ignore all remaining text on the line,
% and is used for comments like this one.

\documentclass{article}      % Specifies the document class

                             % The preamble begins here.
\title{An Example Document}  % Declares the document's title.
\author{Leslie Lamport}      % Declares the author's name.
\date{January 21, 1994}      % Deleting this command produces today's date.

\newcommand{\ip}[2]{(#1, #2)}
                             % Defines \ip{arg1}{arg2} to mean
                             % (arg1, arg2).

%\newcommand{\ip}[2]{\langle #1 | #2\rangle}
                             % This is an alternative definition of
                             % \ip that is commented out.

\begin{document}             % End of preamble and beginning of text.

\maketitle                   % Produces the title.

This is an example input file.  Comparing it with
the output it generates can show you how to
produce a simple document of your own.

\section{Ordinary Text}      % Produces section heading.  Lower-level
                             % sections are begun with similar
                             % \subsection and \subsubsection commands.

The ends  of words and sentences are marked
  by   spaces. It  doesn't matter how many
spaces    you type; one is as good as 100.  The
end of   a line counts as a space.

One   or more   blank lines denote the  end
of  a paragraph.

Since any number of consecutive spaces are treated
like a single one, the formatting of the input
file makes no difference to
      \LaTeX,                % The \LaTeX command generates the LaTeX logo.
but it makes a difference to you.  When you use
\LaTeX, making your input file as easy to read
as possible will be a great help as you write
your document and when you change it.  This sample
file shows how you can add comments to your own input
file.

Because printing is different from typewriting,
there are a number of things that you have to do
differently when preparing an input file than if
you were just typing the document directly.
Quotation marks like
       ``this''
have to be handled specially, as do quotes within
quotes:
       ``\,`this'            % \, separates the double and single quote.
        is what I just
        wrote, not  `that'\,''.

Dashes come in three sizes: an
       intra-word
dash, a medium dash for number ranges like
       1--2,
and a punctuation
       dash---like
this.

A sentence-ending space should be larger than the
space between words within a sentence.  You
sometimes have to type special commands in
conjunction with punctuation characters to get
this right, as in the following sentence.
       Gnats, gnus, etc.\ all  % `\ ' makes an inter-word space.
       begin with G\@.         % \@ marks end-of-sentence punctuation.
You should check the spaces after periods when
reading your output to make sure you haven't
forgotten any special cases.  Generating an
ellipsis
       \ldots\               % `\ ' is needed after `\ldots' because TeX
                             % ignores spaces after command names like \ldots
                             % made from \ + letters.
                             %
                             % Note how a `%' character causes TeX to ignore
                             % the end of the input line, so these blank lines
                             % do not start a new paragraph.
                             %
with the right spacing around the periods requires
a special command.

\LaTeX\ interprets some common characters as
commands, so you must type special commands to
generate them.  These characters include the
following:
       \$ \& \% \# \{ and \}.

In printing, text is usually emphasized with an
       \emph{italic}
type style.

\begin{em}
   A long segment of text can also be emphasized
   in this way.  Text within such a segment can be
   given \emph{additional} emphasis.
\end{em}

It is sometimes necessary to prevent \LaTeX\ from
breaking a line where it might otherwise do so.
This may be at a space, as between the ``Mr.''\ and
``Jones'' in
       ``Mr.~Jones'',        % ~ produces an unbreakable interword space.
or within a word---especially when the word is a
symbol like
       \mbox{\emph{itemnum}}
that makes little sense when hyphenated across
lines.

Footnotes\footnote{This is an example of a footnote.}
pose no problem.

\LaTeX\ is good at typesetting mathematical formulas
like
       \( x-3y + z = 7 \)
or
       \( a_{1} > x^{2n} + y^{2n} > x' \)
or
       \( \ip{A}{B} = \sum_{i} a_{i} b_{i} \).
The spaces you type in a formula are
ignored.  Remember that a letter like
       $x$                   % $ ... $  and  \( ... \)  are equivalent
is a formula when it denotes a mathematical
symbol, and it should be typed as one.

\section{Displayed Text}

Text is displayed by indenting it from the left
margin.  Quotations are commonly displayed.  There
are short quotations
\begin{quote}
   This is a short quotation.  It consists of a
   single paragraph of text.  See how it is formatted.
\end{quote}
and longer ones.
\begin{quotation}
   This is a longer quotation.  It consists of two
   paragraphs of text, neither of which are
   particularly interesting.

   This is the second paragraph of the quotation.  It
   is just as dull as the first paragraph.
\end{quotation}
Another frequently-displayed structure is a list.
The following is an example of an \emph{itemized}
list.
\begin{itemize}
   \item This is the first item of an itemized list.
         Each item in the list is marked with a ``tick''.
         You don't have to worry about what kind of tick
         mark is used.

   \item This is the second item of the list.  It
         contains another list nested inside it.  The inner
         list is an \emph{enumerated} list.
         \begin{enumerate}
            \item This is the first item of an enumerated
                  list that is nested within the itemized list.

            \item This is the second item of the inner list.
                  \LaTeX\ allows you to nest lists deeper than
                  you really should.
         \end{enumerate}
         This is the rest of the second item of the outer
         list.  It is no more interesting than any other
         part of the item.
   \item This is the third item of the list.
\end{itemize}
You can even display poetry.
\begin{verse}
   There is an environment
    for verse \\             % The \\ command separates lines
   Whose features some poets % within a stanza.
   will curse.

                             % One or more blank lines separate stanzas.

   For instead of making\\
   Them do \emph{all} line breaking, \\
   It allows them to put too many words on a line when they'd rather be
   forced to be terse.
\end{verse}

Mathematical formulas may also be displayed.  A
displayed formula
is
one-line long; multiline
formulas require special formatting instructions.
   \[  \ip{\Gamma}{\psi'} = x'' + y^{2} + z_{i}^{n}\]
Don't start a paragraph with a displayed equation,
nor make one a paragraph by itself.

\end{document}               % End of document.

\section{Bench Huge Section 241}
% synthetic bench section 241
% This is a sample LaTeX input file.  (Version of 12 August 2004.)
%
% A '%' character causes TeX to ignore all remaining text on the line,
% and is used for comments like this one.

\documentclass{article}      % Specifies the document class

                             % The preamble begins here.
\title{An Example Document}  % Declares the document's title.
\author{Leslie Lamport}      % Declares the author's name.
\date{January 21, 1994}      % Deleting this command produces today's date.

\newcommand{\ip}[2]{(#1, #2)}
                             % Defines \ip{arg1}{arg2} to mean
                             % (arg1, arg2).

%\newcommand{\ip}[2]{\langle #1 | #2\rangle}
                             % This is an alternative definition of
                             % \ip that is commented out.

\begin{document}             % End of preamble and beginning of text.

\maketitle                   % Produces the title.

This is an example input file.  Comparing it with
the output it generates can show you how to
produce a simple document of your own.

\section{Ordinary Text}      % Produces section heading.  Lower-level
                             % sections are begun with similar
                             % \subsection and \subsubsection commands.

The ends  of words and sentences are marked
  by   spaces. It  doesn't matter how many
spaces    you type; one is as good as 100.  The
end of   a line counts as a space.

One   or more   blank lines denote the  end
of  a paragraph.

Since any number of consecutive spaces are treated
like a single one, the formatting of the input
file makes no difference to
      \LaTeX,                % The \LaTeX command generates the LaTeX logo.
but it makes a difference to you.  When you use
\LaTeX, making your input file as easy to read
as possible will be a great help as you write
your document and when you change it.  This sample
file shows how you can add comments to your own input
file.

Because printing is different from typewriting,
there are a number of things that you have to do
differently when preparing an input file than if
you were just typing the document directly.
Quotation marks like
       ``this''
have to be handled specially, as do quotes within
quotes:
       ``\,`this'            % \, separates the double and single quote.
        is what I just
        wrote, not  `that'\,''.

Dashes come in three sizes: an
       intra-word
dash, a medium dash for number ranges like
       1--2,
and a punctuation
       dash---like
this.

A sentence-ending space should be larger than the
space between words within a sentence.  You
sometimes have to type special commands in
conjunction with punctuation characters to get
this right, as in the following sentence.
       Gnats, gnus, etc.\ all  % `\ ' makes an inter-word space.
       begin with G\@.         % \@ marks end-of-sentence punctuation.
You should check the spaces after periods when
reading your output to make sure you haven't
forgotten any special cases.  Generating an
ellipsis
       \ldots\               % `\ ' is needed after `\ldots' because TeX
                             % ignores spaces after command names like \ldots
                             % made from \ + letters.
                             %
                             % Note how a `%' character causes TeX to ignore
                             % the end of the input line, so these blank lines
                             % do not start a new paragraph.
                             %
with the right spacing around the periods requires
a special command.

\LaTeX\ interprets some common characters as
commands, so you must type special commands to
generate them.  These characters include the
following:
       \$ \& \% \# \{ and \}.

In printing, text is usually emphasized with an
       \emph{italic}
type style.

\begin{em}
   A long segment of text can also be emphasized
   in this way.  Text within such a segment can be
   given \emph{additional} emphasis.
\end{em}

It is sometimes necessary to prevent \LaTeX\ from
breaking a line where it might otherwise do so.
This may be at a space, as between the ``Mr.''\ and
``Jones'' in
       ``Mr.~Jones'',        % ~ produces an unbreakable interword space.
or within a word---especially when the word is a
symbol like
       \mbox{\emph{itemnum}}
that makes little sense when hyphenated across
lines.

Footnotes\footnote{This is an example of a footnote.}
pose no problem.

\LaTeX\ is good at typesetting mathematical formulas
like
       \( x-3y + z = 7 \)
or
       \( a_{1} > x^{2n} + y^{2n} > x' \)
or
       \( \ip{A}{B} = \sum_{i} a_{i} b_{i} \).
The spaces you type in a formula are
ignored.  Remember that a letter like
       $x$                   % $ ... $  and  \( ... \)  are equivalent
is a formula when it denotes a mathematical
symbol, and it should be typed as one.

\section{Displayed Text}

Text is displayed by indenting it from the left
margin.  Quotations are commonly displayed.  There
are short quotations
\begin{quote}
   This is a short quotation.  It consists of a
   single paragraph of text.  See how it is formatted.
\end{quote}
and longer ones.
\begin{quotation}
   This is a longer quotation.  It consists of two
   paragraphs of text, neither of which are
   particularly interesting.

   This is the second paragraph of the quotation.  It
   is just as dull as the first paragraph.
\end{quotation}
Another frequently-displayed structure is a list.
The following is an example of an \emph{itemized}
list.
\begin{itemize}
   \item This is the first item of an itemized list.
         Each item in the list is marked with a ``tick''.
         You don't have to worry about what kind of tick
         mark is used.

   \item This is the second item of the list.  It
         contains another list nested inside it.  The inner
         list is an \emph{enumerated} list.
         \begin{enumerate}
            \item This is the first item of an enumerated
                  list that is nested within the itemized list.

            \item This is the second item of the inner list.
                  \LaTeX\ allows you to nest lists deeper than
                  you really should.
         \end{enumerate}
         This is the rest of the second item of the outer
         list.  It is no more interesting than any other
         part of the item.
   \item This is the third item of the list.
\end{itemize}
You can even display poetry.
\begin{verse}
   There is an environment
    for verse \\             % The \\ command separates lines
   Whose features some poets % within a stanza.
   will curse.

                             % One or more blank lines separate stanzas.

   For instead of making\\
   Them do \emph{all} line breaking, \\
   It allows them to put too many words on a line when they'd rather be
   forced to be terse.
\end{verse}

Mathematical formulas may also be displayed.  A
displayed formula
is
one-line long; multiline
formulas require special formatting instructions.
   \[  \ip{\Gamma}{\psi'} = x'' + y^{2} + z_{i}^{n}\]
Don't start a paragraph with a displayed equation,
nor make one a paragraph by itself.

\end{document}               % End of document.

\section{Bench Huge Section 242}
% synthetic bench section 242
% This is a sample LaTeX input file.  (Version of 12 August 2004.)
%
% A '%' character causes TeX to ignore all remaining text on the line,
% and is used for comments like this one.

\documentclass{article}      % Specifies the document class

                             % The preamble begins here.
\title{An Example Document}  % Declares the document's title.
\author{Leslie Lamport}      % Declares the author's name.
\date{January 21, 1994}      % Deleting this command produces today's date.

\newcommand{\ip}[2]{(#1, #2)}
                             % Defines \ip{arg1}{arg2} to mean
                             % (arg1, arg2).

%\newcommand{\ip}[2]{\langle #1 | #2\rangle}
                             % This is an alternative definition of
                             % \ip that is commented out.

\begin{document}             % End of preamble and beginning of text.

\maketitle                   % Produces the title.

This is an example input file.  Comparing it with
the output it generates can show you how to
produce a simple document of your own.

\section{Ordinary Text}      % Produces section heading.  Lower-level
                             % sections are begun with similar
                             % \subsection and \subsubsection commands.

The ends  of words and sentences are marked
  by   spaces. It  doesn't matter how many
spaces    you type; one is as good as 100.  The
end of   a line counts as a space.

One   or more   blank lines denote the  end
of  a paragraph.

Since any number of consecutive spaces are treated
like a single one, the formatting of the input
file makes no difference to
      \LaTeX,                % The \LaTeX command generates the LaTeX logo.
but it makes a difference to you.  When you use
\LaTeX, making your input file as easy to read
as possible will be a great help as you write
your document and when you change it.  This sample
file shows how you can add comments to your own input
file.

Because printing is different from typewriting,
there are a number of things that you have to do
differently when preparing an input file than if
you were just typing the document directly.
Quotation marks like
       ``this''
have to be handled specially, as do quotes within
quotes:
       ``\,`this'            % \, separates the double and single quote.
        is what I just
        wrote, not  `that'\,''.

Dashes come in three sizes: an
       intra-word
dash, a medium dash for number ranges like
       1--2,
and a punctuation
       dash---like
this.

A sentence-ending space should be larger than the
space between words within a sentence.  You
sometimes have to type special commands in
conjunction with punctuation characters to get
this right, as in the following sentence.
       Gnats, gnus, etc.\ all  % `\ ' makes an inter-word space.
       begin with G\@.         % \@ marks end-of-sentence punctuation.
You should check the spaces after periods when
reading your output to make sure you haven't
forgotten any special cases.  Generating an
ellipsis
       \ldots\               % `\ ' is needed after `\ldots' because TeX
                             % ignores spaces after command names like \ldots
                             % made from \ + letters.
                             %
                             % Note how a `%' character causes TeX to ignore
                             % the end of the input line, so these blank lines
                             % do not start a new paragraph.
                             %
with the right spacing around the periods requires
a special command.

\LaTeX\ interprets some common characters as
commands, so you must type special commands to
generate them.  These characters include the
following:
       \$ \& \% \# \{ and \}.

In printing, text is usually emphasized with an
       \emph{italic}
type style.

\begin{em}
   A long segment of text can also be emphasized
   in this way.  Text within such a segment can be
   given \emph{additional} emphasis.
\end{em}

It is sometimes necessary to prevent \LaTeX\ from
breaking a line where it might otherwise do so.
This may be at a space, as between the ``Mr.''\ and
``Jones'' in
       ``Mr.~Jones'',        % ~ produces an unbreakable interword space.
or within a word---especially when the word is a
symbol like
       \mbox{\emph{itemnum}}
that makes little sense when hyphenated across
lines.

Footnotes\footnote{This is an example of a footnote.}
pose no problem.

\LaTeX\ is good at typesetting mathematical formulas
like
       \( x-3y + z = 7 \)
or
       \( a_{1} > x^{2n} + y^{2n} > x' \)
or
       \( \ip{A}{B} = \sum_{i} a_{i} b_{i} \).
The spaces you type in a formula are
ignored.  Remember that a letter like
       $x$                   % $ ... $  and  \( ... \)  are equivalent
is a formula when it denotes a mathematical
symbol, and it should be typed as one.

\section{Displayed Text}

Text is displayed by indenting it from the left
margin.  Quotations are commonly displayed.  There
are short quotations
\begin{quote}
   This is a short quotation.  It consists of a
   single paragraph of text.  See how it is formatted.
\end{quote}
and longer ones.
\begin{quotation}
   This is a longer quotation.  It consists of two
   paragraphs of text, neither of which are
   particularly interesting.

   This is the second paragraph of the quotation.  It
   is just as dull as the first paragraph.
\end{quotation}
Another frequently-displayed structure is a list.
The following is an example of an \emph{itemized}
list.
\begin{itemize}
   \item This is the first item of an itemized list.
         Each item in the list is marked with a ``tick''.
         You don't have to worry about what kind of tick
         mark is used.

   \item This is the second item of the list.  It
         contains another list nested inside it.  The inner
         list is an \emph{enumerated} list.
         \begin{enumerate}
            \item This is the first item of an enumerated
                  list that is nested within the itemized list.

            \item This is the second item of the inner list.
                  \LaTeX\ allows you to nest lists deeper than
                  you really should.
         \end{enumerate}
         This is the rest of the second item of the outer
         list.  It is no more interesting than any other
         part of the item.
   \item This is the third item of the list.
\end{itemize}
You can even display poetry.
\begin{verse}
   There is an environment
    for verse \\             % The \\ command separates lines
   Whose features some poets % within a stanza.
   will curse.

                             % One or more blank lines separate stanzas.

   For instead of making\\
   Them do \emph{all} line breaking, \\
   It allows them to put too many words on a line when they'd rather be
   forced to be terse.
\end{verse}

Mathematical formulas may also be displayed.  A
displayed formula
is
one-line long; multiline
formulas require special formatting instructions.
   \[  \ip{\Gamma}{\psi'} = x'' + y^{2} + z_{i}^{n}\]
Don't start a paragraph with a displayed equation,
nor make one a paragraph by itself.

\end{document}               % End of document.

\section{Bench Huge Section 243}
% synthetic bench section 243
% This is a sample LaTeX input file.  (Version of 12 August 2004.)
%
% A '%' character causes TeX to ignore all remaining text on the line,
% and is used for comments like this one.

\documentclass{article}      % Specifies the document class

                             % The preamble begins here.
\title{An Example Document}  % Declares the document's title.
\author{Leslie Lamport}      % Declares the author's name.
\date{January 21, 1994}      % Deleting this command produces today's date.

\newcommand{\ip}[2]{(#1, #2)}
                             % Defines \ip{arg1}{arg2} to mean
                             % (arg1, arg2).

%\newcommand{\ip}[2]{\langle #1 | #2\rangle}
                             % This is an alternative definition of
                             % \ip that is commented out.

\begin{document}             % End of preamble and beginning of text.

\maketitle                   % Produces the title.

This is an example input file.  Comparing it with
the output it generates can show you how to
produce a simple document of your own.

\section{Ordinary Text}      % Produces section heading.  Lower-level
                             % sections are begun with similar
                             % \subsection and \subsubsection commands.

The ends  of words and sentences are marked
  by   spaces. It  doesn't matter how many
spaces    you type; one is as good as 100.  The
end of   a line counts as a space.

One   or more   blank lines denote the  end
of  a paragraph.

Since any number of consecutive spaces are treated
like a single one, the formatting of the input
file makes no difference to
      \LaTeX,                % The \LaTeX command generates the LaTeX logo.
but it makes a difference to you.  When you use
\LaTeX, making your input file as easy to read
as possible will be a great help as you write
your document and when you change it.  This sample
file shows how you can add comments to your own input
file.

Because printing is different from typewriting,
there are a number of things that you have to do
differently when preparing an input file than if
you were just typing the document directly.
Quotation marks like
       ``this''
have to be handled specially, as do quotes within
quotes:
       ``\,`this'            % \, separates the double and single quote.
        is what I just
        wrote, not  `that'\,''.

Dashes come in three sizes: an
       intra-word
dash, a medium dash for number ranges like
       1--2,
and a punctuation
       dash---like
this.

A sentence-ending space should be larger than the
space between words within a sentence.  You
sometimes have to type special commands in
conjunction with punctuation characters to get
this right, as in the following sentence.
       Gnats, gnus, etc.\ all  % `\ ' makes an inter-word space.
       begin with G\@.         % \@ marks end-of-sentence punctuation.
You should check the spaces after periods when
reading your output to make sure you haven't
forgotten any special cases.  Generating an
ellipsis
       \ldots\               % `\ ' is needed after `\ldots' because TeX
                             % ignores spaces after command names like \ldots
                             % made from \ + letters.
                             %
                             % Note how a `%' character causes TeX to ignore
                             % the end of the input line, so these blank lines
                             % do not start a new paragraph.
                             %
with the right spacing around the periods requires
a special command.

\LaTeX\ interprets some common characters as
commands, so you must type special commands to
generate them.  These characters include the
following:
       \$ \& \% \# \{ and \}.

In printing, text is usually emphasized with an
       \emph{italic}
type style.

\begin{em}
   A long segment of text can also be emphasized
   in this way.  Text within such a segment can be
   given \emph{additional} emphasis.
\end{em}

It is sometimes necessary to prevent \LaTeX\ from
breaking a line where it might otherwise do so.
This may be at a space, as between the ``Mr.''\ and
``Jones'' in
       ``Mr.~Jones'',        % ~ produces an unbreakable interword space.
or within a word---especially when the word is a
symbol like
       \mbox{\emph{itemnum}}
that makes little sense when hyphenated across
lines.

Footnotes\footnote{This is an example of a footnote.}
pose no problem.

\LaTeX\ is good at typesetting mathematical formulas
like
       \( x-3y + z = 7 \)
or
       \( a_{1} > x^{2n} + y^{2n} > x' \)
or
       \( \ip{A}{B} = \sum_{i} a_{i} b_{i} \).
The spaces you type in a formula are
ignored.  Remember that a letter like
       $x$                   % $ ... $  and  \( ... \)  are equivalent
is a formula when it denotes a mathematical
symbol, and it should be typed as one.

\section{Displayed Text}

Text is displayed by indenting it from the left
margin.  Quotations are commonly displayed.  There
are short quotations
\begin{quote}
   This is a short quotation.  It consists of a
   single paragraph of text.  See how it is formatted.
\end{quote}
and longer ones.
\begin{quotation}
   This is a longer quotation.  It consists of two
   paragraphs of text, neither of which are
   particularly interesting.

   This is the second paragraph of the quotation.  It
   is just as dull as the first paragraph.
\end{quotation}
Another frequently-displayed structure is a list.
The following is an example of an \emph{itemized}
list.
\begin{itemize}
   \item This is the first item of an itemized list.
         Each item in the list is marked with a ``tick''.
         You don't have to worry about what kind of tick
         mark is used.

   \item This is the second item of the list.  It
         contains another list nested inside it.  The inner
         list is an \emph{enumerated} list.
         \begin{enumerate}
            \item This is the first item of an enumerated
                  list that is nested within the itemized list.

            \item This is the second item of the inner list.
                  \LaTeX\ allows you to nest lists deeper than
                  you really should.
         \end{enumerate}
         This is the rest of the second item of the outer
         list.  It is no more interesting than any other
         part of the item.
   \item This is the third item of the list.
\end{itemize}
You can even display poetry.
\begin{verse}
   There is an environment
    for verse \\             % The \\ command separates lines
   Whose features some poets % within a stanza.
   will curse.

                             % One or more blank lines separate stanzas.

   For instead of making\\
   Them do \emph{all} line breaking, \\
   It allows them to put too many words on a line when they'd rather be
   forced to be terse.
\end{verse}

Mathematical formulas may also be displayed.  A
displayed formula
is
one-line long; multiline
formulas require special formatting instructions.
   \[  \ip{\Gamma}{\psi'} = x'' + y^{2} + z_{i}^{n}\]
Don't start a paragraph with a displayed equation,
nor make one a paragraph by itself.

\end{document}               % End of document.

\section{Bench Huge Section 244}
% synthetic bench section 244
% This is a sample LaTeX input file.  (Version of 12 August 2004.)
%
% A '%' character causes TeX to ignore all remaining text on the line,
% and is used for comments like this one.

\documentclass{article}      % Specifies the document class

                             % The preamble begins here.
\title{An Example Document}  % Declares the document's title.
\author{Leslie Lamport}      % Declares the author's name.
\date{January 21, 1994}      % Deleting this command produces today's date.

\newcommand{\ip}[2]{(#1, #2)}
                             % Defines \ip{arg1}{arg2} to mean
                             % (arg1, arg2).

%\newcommand{\ip}[2]{\langle #1 | #2\rangle}
                             % This is an alternative definition of
                             % \ip that is commented out.

\begin{document}             % End of preamble and beginning of text.

\maketitle                   % Produces the title.

This is an example input file.  Comparing it with
the output it generates can show you how to
produce a simple document of your own.

\section{Ordinary Text}      % Produces section heading.  Lower-level
                             % sections are begun with similar
                             % \subsection and \subsubsection commands.

The ends  of words and sentences are marked
  by   spaces. It  doesn't matter how many
spaces    you type; one is as good as 100.  The
end of   a line counts as a space.

One   or more   blank lines denote the  end
of  a paragraph.

Since any number of consecutive spaces are treated
like a single one, the formatting of the input
file makes no difference to
      \LaTeX,                % The \LaTeX command generates the LaTeX logo.
but it makes a difference to you.  When you use
\LaTeX, making your input file as easy to read
as possible will be a great help as you write
your document and when you change it.  This sample
file shows how you can add comments to your own input
file.

Because printing is different from typewriting,
there are a number of things that you have to do
differently when preparing an input file than if
you were just typing the document directly.
Quotation marks like
       ``this''
have to be handled specially, as do quotes within
quotes:
       ``\,`this'            % \, separates the double and single quote.
        is what I just
        wrote, not  `that'\,''.

Dashes come in three sizes: an
       intra-word
dash, a medium dash for number ranges like
       1--2,
and a punctuation
       dash---like
this.

A sentence-ending space should be larger than the
space between words within a sentence.  You
sometimes have to type special commands in
conjunction with punctuation characters to get
this right, as in the following sentence.
       Gnats, gnus, etc.\ all  % `\ ' makes an inter-word space.
       begin with G\@.         % \@ marks end-of-sentence punctuation.
You should check the spaces after periods when
reading your output to make sure you haven't
forgotten any special cases.  Generating an
ellipsis
       \ldots\               % `\ ' is needed after `\ldots' because TeX
                             % ignores spaces after command names like \ldots
                             % made from \ + letters.
                             %
                             % Note how a `%' character causes TeX to ignore
                             % the end of the input line, so these blank lines
                             % do not start a new paragraph.
                             %
with the right spacing around the periods requires
a special command.

\LaTeX\ interprets some common characters as
commands, so you must type special commands to
generate them.  These characters include the
following:
       \$ \& \% \# \{ and \}.

In printing, text is usually emphasized with an
       \emph{italic}
type style.

\begin{em}
   A long segment of text can also be emphasized
   in this way.  Text within such a segment can be
   given \emph{additional} emphasis.
\end{em}

It is sometimes necessary to prevent \LaTeX\ from
breaking a line where it might otherwise do so.
This may be at a space, as between the ``Mr.''\ and
``Jones'' in
       ``Mr.~Jones'',        % ~ produces an unbreakable interword space.
or within a word---especially when the word is a
symbol like
       \mbox{\emph{itemnum}}
that makes little sense when hyphenated across
lines.

Footnotes\footnote{This is an example of a footnote.}
pose no problem.

\LaTeX\ is good at typesetting mathematical formulas
like
       \( x-3y + z = 7 \)
or
       \( a_{1} > x^{2n} + y^{2n} > x' \)
or
       \( \ip{A}{B} = \sum_{i} a_{i} b_{i} \).
The spaces you type in a formula are
ignored.  Remember that a letter like
       $x$                   % $ ... $  and  \( ... \)  are equivalent
is a formula when it denotes a mathematical
symbol, and it should be typed as one.

\section{Displayed Text}

Text is displayed by indenting it from the left
margin.  Quotations are commonly displayed.  There
are short quotations
\begin{quote}
   This is a short quotation.  It consists of a
   single paragraph of text.  See how it is formatted.
\end{quote}
and longer ones.
\begin{quotation}
   This is a longer quotation.  It consists of two
   paragraphs of text, neither of which are
   particularly interesting.

   This is the second paragraph of the quotation.  It
   is just as dull as the first paragraph.
\end{quotation}
Another frequently-displayed structure is a list.
The following is an example of an \emph{itemized}
list.
\begin{itemize}
   \item This is the first item of an itemized list.
         Each item in the list is marked with a ``tick''.
         You don't have to worry about what kind of tick
         mark is used.

   \item This is the second item of the list.  It
         contains another list nested inside it.  The inner
         list is an \emph{enumerated} list.
         \begin{enumerate}
            \item This is the first item of an enumerated
                  list that is nested within the itemized list.

            \item This is the second item of the inner list.
                  \LaTeX\ allows you to nest lists deeper than
                  you really should.
         \end{enumerate}
         This is the rest of the second item of the outer
         list.  It is no more interesting than any other
         part of the item.
   \item This is the third item of the list.
\end{itemize}
You can even display poetry.
\begin{verse}
   There is an environment
    for verse \\             % The \\ command separates lines
   Whose features some poets % within a stanza.
   will curse.

                             % One or more blank lines separate stanzas.

   For instead of making\\
   Them do \emph{all} line breaking, \\
   It allows them to put too many words on a line when they'd rather be
   forced to be terse.
\end{verse}

Mathematical formulas may also be displayed.  A
displayed formula
is
one-line long; multiline
formulas require special formatting instructions.
   \[  \ip{\Gamma}{\psi'} = x'' + y^{2} + z_{i}^{n}\]
Don't start a paragraph with a displayed equation,
nor make one a paragraph by itself.

\end{document}               % End of document.

\section{Bench Huge Section 245}
% synthetic bench section 245
% This is a sample LaTeX input file.  (Version of 12 August 2004.)
%
% A '%' character causes TeX to ignore all remaining text on the line,
% and is used for comments like this one.

\documentclass{article}      % Specifies the document class

                             % The preamble begins here.
\title{An Example Document}  % Declares the document's title.
\author{Leslie Lamport}      % Declares the author's name.
\date{January 21, 1994}      % Deleting this command produces today's date.

\newcommand{\ip}[2]{(#1, #2)}
                             % Defines \ip{arg1}{arg2} to mean
                             % (arg1, arg2).

%\newcommand{\ip}[2]{\langle #1 | #2\rangle}
                             % This is an alternative definition of
                             % \ip that is commented out.

\begin{document}             % End of preamble and beginning of text.

\maketitle                   % Produces the title.

This is an example input file.  Comparing it with
the output it generates can show you how to
produce a simple document of your own.

\section{Ordinary Text}      % Produces section heading.  Lower-level
                             % sections are begun with similar
                             % \subsection and \subsubsection commands.

The ends  of words and sentences are marked
  by   spaces. It  doesn't matter how many
spaces    you type; one is as good as 100.  The
end of   a line counts as a space.

One   or more   blank lines denote the  end
of  a paragraph.

Since any number of consecutive spaces are treated
like a single one, the formatting of the input
file makes no difference to
      \LaTeX,                % The \LaTeX command generates the LaTeX logo.
but it makes a difference to you.  When you use
\LaTeX, making your input file as easy to read
as possible will be a great help as you write
your document and when you change it.  This sample
file shows how you can add comments to your own input
file.

Because printing is different from typewriting,
there are a number of things that you have to do
differently when preparing an input file than if
you were just typing the document directly.
Quotation marks like
       ``this''
have to be handled specially, as do quotes within
quotes:
       ``\,`this'            % \, separates the double and single quote.
        is what I just
        wrote, not  `that'\,''.

Dashes come in three sizes: an
       intra-word
dash, a medium dash for number ranges like
       1--2,
and a punctuation
       dash---like
this.

A sentence-ending space should be larger than the
space between words within a sentence.  You
sometimes have to type special commands in
conjunction with punctuation characters to get
this right, as in the following sentence.
       Gnats, gnus, etc.\ all  % `\ ' makes an inter-word space.
       begin with G\@.         % \@ marks end-of-sentence punctuation.
You should check the spaces after periods when
reading your output to make sure you haven't
forgotten any special cases.  Generating an
ellipsis
       \ldots\               % `\ ' is needed after `\ldots' because TeX
                             % ignores spaces after command names like \ldots
                             % made from \ + letters.
                             %
                             % Note how a `%' character causes TeX to ignore
                             % the end of the input line, so these blank lines
                             % do not start a new paragraph.
                             %
with the right spacing around the periods requires
a special command.

\LaTeX\ interprets some common characters as
commands, so you must type special commands to
generate them.  These characters include the
following:
       \$ \& \% \# \{ and \}.

In printing, text is usually emphasized with an
       \emph{italic}
type style.

\begin{em}
   A long segment of text can also be emphasized
   in this way.  Text within such a segment can be
   given \emph{additional} emphasis.
\end{em}

It is sometimes necessary to prevent \LaTeX\ from
breaking a line where it might otherwise do so.
This may be at a space, as between the ``Mr.''\ and
``Jones'' in
       ``Mr.~Jones'',        % ~ produces an unbreakable interword space.
or within a word---especially when the word is a
symbol like
       \mbox{\emph{itemnum}}
that makes little sense when hyphenated across
lines.

Footnotes\footnote{This is an example of a footnote.}
pose no problem.

\LaTeX\ is good at typesetting mathematical formulas
like
       \( x-3y + z = 7 \)
or
       \( a_{1} > x^{2n} + y^{2n} > x' \)
or
       \( \ip{A}{B} = \sum_{i} a_{i} b_{i} \).
The spaces you type in a formula are
ignored.  Remember that a letter like
       $x$                   % $ ... $  and  \( ... \)  are equivalent
is a formula when it denotes a mathematical
symbol, and it should be typed as one.

\section{Displayed Text}

Text is displayed by indenting it from the left
margin.  Quotations are commonly displayed.  There
are short quotations
\begin{quote}
   This is a short quotation.  It consists of a
   single paragraph of text.  See how it is formatted.
\end{quote}
and longer ones.
\begin{quotation}
   This is a longer quotation.  It consists of two
   paragraphs of text, neither of which are
   particularly interesting.

   This is the second paragraph of the quotation.  It
   is just as dull as the first paragraph.
\end{quotation}
Another frequently-displayed structure is a list.
The following is an example of an \emph{itemized}
list.
\begin{itemize}
   \item This is the first item of an itemized list.
         Each item in the list is marked with a ``tick''.
         You don't have to worry about what kind of tick
         mark is used.

   \item This is the second item of the list.  It
         contains another list nested inside it.  The inner
         list is an \emph{enumerated} list.
         \begin{enumerate}
            \item This is the first item of an enumerated
                  list that is nested within the itemized list.

            \item This is the second item of the inner list.
                  \LaTeX\ allows you to nest lists deeper than
                  you really should.
         \end{enumerate}
         This is the rest of the second item of the outer
         list.  It is no more interesting than any other
         part of the item.
   \item This is the third item of the list.
\end{itemize}
You can even display poetry.
\begin{verse}
   There is an environment
    for verse \\             % The \\ command separates lines
   Whose features some poets % within a stanza.
   will curse.

                             % One or more blank lines separate stanzas.

   For instead of making\\
   Them do \emph{all} line breaking, \\
   It allows them to put too many words on a line when they'd rather be
   forced to be terse.
\end{verse}

Mathematical formulas may also be displayed.  A
displayed formula
is
one-line long; multiline
formulas require special formatting instructions.
   \[  \ip{\Gamma}{\psi'} = x'' + y^{2} + z_{i}^{n}\]
Don't start a paragraph with a displayed equation,
nor make one a paragraph by itself.

\end{document}               % End of document.

\section{Bench Huge Section 246}
% synthetic bench section 246
% This is a sample LaTeX input file.  (Version of 12 August 2004.)
%
% A '%' character causes TeX to ignore all remaining text on the line,
% and is used for comments like this one.

\documentclass{article}      % Specifies the document class

                             % The preamble begins here.
\title{An Example Document}  % Declares the document's title.
\author{Leslie Lamport}      % Declares the author's name.
\date{January 21, 1994}      % Deleting this command produces today's date.

\newcommand{\ip}[2]{(#1, #2)}
                             % Defines \ip{arg1}{arg2} to mean
                             % (arg1, arg2).

%\newcommand{\ip}[2]{\langle #1 | #2\rangle}
                             % This is an alternative definition of
                             % \ip that is commented out.

\begin{document}             % End of preamble and beginning of text.

\maketitle                   % Produces the title.

This is an example input file.  Comparing it with
the output it generates can show you how to
produce a simple document of your own.

\section{Ordinary Text}      % Produces section heading.  Lower-level
                             % sections are begun with similar
                             % \subsection and \subsubsection commands.

The ends  of words and sentences are marked
  by   spaces. It  doesn't matter how many
spaces    you type; one is as good as 100.  The
end of   a line counts as a space.

One   or more   blank lines denote the  end
of  a paragraph.

Since any number of consecutive spaces are treated
like a single one, the formatting of the input
file makes no difference to
      \LaTeX,                % The \LaTeX command generates the LaTeX logo.
but it makes a difference to you.  When you use
\LaTeX, making your input file as easy to read
as possible will be a great help as you write
your document and when you change it.  This sample
file shows how you can add comments to your own input
file.

Because printing is different from typewriting,
there are a number of things that you have to do
differently when preparing an input file than if
you were just typing the document directly.
Quotation marks like
       ``this''
have to be handled specially, as do quotes within
quotes:
       ``\,`this'            % \, separates the double and single quote.
        is what I just
        wrote, not  `that'\,''.

Dashes come in three sizes: an
       intra-word
dash, a medium dash for number ranges like
       1--2,
and a punctuation
       dash---like
this.

A sentence-ending space should be larger than the
space between words within a sentence.  You
sometimes have to type special commands in
conjunction with punctuation characters to get
this right, as in the following sentence.
       Gnats, gnus, etc.\ all  % `\ ' makes an inter-word space.
       begin with G\@.         % \@ marks end-of-sentence punctuation.
You should check the spaces after periods when
reading your output to make sure you haven't
forgotten any special cases.  Generating an
ellipsis
       \ldots\               % `\ ' is needed after `\ldots' because TeX
                             % ignores spaces after command names like \ldots
                             % made from \ + letters.
                             %
                             % Note how a `%' character causes TeX to ignore
                             % the end of the input line, so these blank lines
                             % do not start a new paragraph.
                             %
with the right spacing around the periods requires
a special command.

\LaTeX\ interprets some common characters as
commands, so you must type special commands to
generate them.  These characters include the
following:
       \$ \& \% \# \{ and \}.

In printing, text is usually emphasized with an
       \emph{italic}
type style.

\begin{em}
   A long segment of text can also be emphasized
   in this way.  Text within such a segment can be
   given \emph{additional} emphasis.
\end{em}

It is sometimes necessary to prevent \LaTeX\ from
breaking a line where it might otherwise do so.
This may be at a space, as between the ``Mr.''\ and
``Jones'' in
       ``Mr.~Jones'',        % ~ produces an unbreakable interword space.
or within a word---especially when the word is a
symbol like
       \mbox{\emph{itemnum}}
that makes little sense when hyphenated across
lines.

Footnotes\footnote{This is an example of a footnote.}
pose no problem.

\LaTeX\ is good at typesetting mathematical formulas
like
       \( x-3y + z = 7 \)
or
       \( a_{1} > x^{2n} + y^{2n} > x' \)
or
       \( \ip{A}{B} = \sum_{i} a_{i} b_{i} \).
The spaces you type in a formula are
ignored.  Remember that a letter like
       $x$                   % $ ... $  and  \( ... \)  are equivalent
is a formula when it denotes a mathematical
symbol, and it should be typed as one.

\section{Displayed Text}

Text is displayed by indenting it from the left
margin.  Quotations are commonly displayed.  There
are short quotations
\begin{quote}
   This is a short quotation.  It consists of a
   single paragraph of text.  See how it is formatted.
\end{quote}
and longer ones.
\begin{quotation}
   This is a longer quotation.  It consists of two
   paragraphs of text, neither of which are
   particularly interesting.

   This is the second paragraph of the quotation.  It
   is just as dull as the first paragraph.
\end{quotation}
Another frequently-displayed structure is a list.
The following is an example of an \emph{itemized}
list.
\begin{itemize}
   \item This is the first item of an itemized list.
         Each item in the list is marked with a ``tick''.
         You don't have to worry about what kind of tick
         mark is used.

   \item This is the second item of the list.  It
         contains another list nested inside it.  The inner
         list is an \emph{enumerated} list.
         \begin{enumerate}
            \item This is the first item of an enumerated
                  list that is nested within the itemized list.

            \item This is the second item of the inner list.
                  \LaTeX\ allows you to nest lists deeper than
                  you really should.
         \end{enumerate}
         This is the rest of the second item of the outer
         list.  It is no more interesting than any other
         part of the item.
   \item This is the third item of the list.
\end{itemize}
You can even display poetry.
\begin{verse}
   There is an environment
    for verse \\             % The \\ command separates lines
   Whose features some poets % within a stanza.
   will curse.

                             % One or more blank lines separate stanzas.

   For instead of making\\
   Them do \emph{all} line breaking, \\
   It allows them to put too many words on a line when they'd rather be
   forced to be terse.
\end{verse}

Mathematical formulas may also be displayed.  A
displayed formula
is
one-line long; multiline
formulas require special formatting instructions.
   \[  \ip{\Gamma}{\psi'} = x'' + y^{2} + z_{i}^{n}\]
Don't start a paragraph with a displayed equation,
nor make one a paragraph by itself.

\end{document}               % End of document.

\section{Bench Huge Section 247}
% synthetic bench section 247
% This is a sample LaTeX input file.  (Version of 12 August 2004.)
%
% A '%' character causes TeX to ignore all remaining text on the line,
% and is used for comments like this one.

\documentclass{article}      % Specifies the document class

                             % The preamble begins here.
\title{An Example Document}  % Declares the document's title.
\author{Leslie Lamport}      % Declares the author's name.
\date{January 21, 1994}      % Deleting this command produces today's date.

\newcommand{\ip}[2]{(#1, #2)}
                             % Defines \ip{arg1}{arg2} to mean
                             % (arg1, arg2).

%\newcommand{\ip}[2]{\langle #1 | #2\rangle}
                             % This is an alternative definition of
                             % \ip that is commented out.

\begin{document}             % End of preamble and beginning of text.

\maketitle                   % Produces the title.

This is an example input file.  Comparing it with
the output it generates can show you how to
produce a simple document of your own.

\section{Ordinary Text}      % Produces section heading.  Lower-level
                             % sections are begun with similar
                             % \subsection and \subsubsection commands.

The ends  of words and sentences are marked
  by   spaces. It  doesn't matter how many
spaces    you type; one is as good as 100.  The
end of   a line counts as a space.

One   or more   blank lines denote the  end
of  a paragraph.

Since any number of consecutive spaces are treated
like a single one, the formatting of the input
file makes no difference to
      \LaTeX,                % The \LaTeX command generates the LaTeX logo.
but it makes a difference to you.  When you use
\LaTeX, making your input file as easy to read
as possible will be a great help as you write
your document and when you change it.  This sample
file shows how you can add comments to your own input
file.

Because printing is different from typewriting,
there are a number of things that you have to do
differently when preparing an input file than if
you were just typing the document directly.
Quotation marks like
       ``this''
have to be handled specially, as do quotes within
quotes:
       ``\,`this'            % \, separates the double and single quote.
        is what I just
        wrote, not  `that'\,''.

Dashes come in three sizes: an
       intra-word
dash, a medium dash for number ranges like
       1--2,
and a punctuation
       dash---like
this.

A sentence-ending space should be larger than the
space between words within a sentence.  You
sometimes have to type special commands in
conjunction with punctuation characters to get
this right, as in the following sentence.
       Gnats, gnus, etc.\ all  % `\ ' makes an inter-word space.
       begin with G\@.         % \@ marks end-of-sentence punctuation.
You should check the spaces after periods when
reading your output to make sure you haven't
forgotten any special cases.  Generating an
ellipsis
       \ldots\               % `\ ' is needed after `\ldots' because TeX
                             % ignores spaces after command names like \ldots
                             % made from \ + letters.
                             %
                             % Note how a `%' character causes TeX to ignore
                             % the end of the input line, so these blank lines
                             % do not start a new paragraph.
                             %
with the right spacing around the periods requires
a special command.

\LaTeX\ interprets some common characters as
commands, so you must type special commands to
generate them.  These characters include the
following:
       \$ \& \% \# \{ and \}.

In printing, text is usually emphasized with an
       \emph{italic}
type style.

\begin{em}
   A long segment of text can also be emphasized
   in this way.  Text within such a segment can be
   given \emph{additional} emphasis.
\end{em}

It is sometimes necessary to prevent \LaTeX\ from
breaking a line where it might otherwise do so.
This may be at a space, as between the ``Mr.''\ and
``Jones'' in
       ``Mr.~Jones'',        % ~ produces an unbreakable interword space.
or within a word---especially when the word is a
symbol like
       \mbox{\emph{itemnum}}
that makes little sense when hyphenated across
lines.

Footnotes\footnote{This is an example of a footnote.}
pose no problem.

\LaTeX\ is good at typesetting mathematical formulas
like
       \( x-3y + z = 7 \)
or
       \( a_{1} > x^{2n} + y^{2n} > x' \)
or
       \( \ip{A}{B} = \sum_{i} a_{i} b_{i} \).
The spaces you type in a formula are
ignored.  Remember that a letter like
       $x$                   % $ ... $  and  \( ... \)  are equivalent
is a formula when it denotes a mathematical
symbol, and it should be typed as one.

\section{Displayed Text}

Text is displayed by indenting it from the left
margin.  Quotations are commonly displayed.  There
are short quotations
\begin{quote}
   This is a short quotation.  It consists of a
   single paragraph of text.  See how it is formatted.
\end{quote}
and longer ones.
\begin{quotation}
   This is a longer quotation.  It consists of two
   paragraphs of text, neither of which are
   particularly interesting.

   This is the second paragraph of the quotation.  It
   is just as dull as the first paragraph.
\end{quotation}
Another frequently-displayed structure is a list.
The following is an example of an \emph{itemized}
list.
\begin{itemize}
   \item This is the first item of an itemized list.
         Each item in the list is marked with a ``tick''.
         You don't have to worry about what kind of tick
         mark is used.

   \item This is the second item of the list.  It
         contains another list nested inside it.  The inner
         list is an \emph{enumerated} list.
         \begin{enumerate}
            \item This is the first item of an enumerated
                  list that is nested within the itemized list.

            \item This is the second item of the inner list.
                  \LaTeX\ allows you to nest lists deeper than
                  you really should.
         \end{enumerate}
         This is the rest of the second item of the outer
         list.  It is no more interesting than any other
         part of the item.
   \item This is the third item of the list.
\end{itemize}
You can even display poetry.
\begin{verse}
   There is an environment
    for verse \\             % The \\ command separates lines
   Whose features some poets % within a stanza.
   will curse.

                             % One or more blank lines separate stanzas.

   For instead of making\\
   Them do \emph{all} line breaking, \\
   It allows them to put too many words on a line when they'd rather be
   forced to be terse.
\end{verse}

Mathematical formulas may also be displayed.  A
displayed formula
is
one-line long; multiline
formulas require special formatting instructions.
   \[  \ip{\Gamma}{\psi'} = x'' + y^{2} + z_{i}^{n}\]
Don't start a paragraph with a displayed equation,
nor make one a paragraph by itself.

\end{document}               % End of document.

\section{Bench Huge Section 248}
% synthetic bench section 248
% This is a sample LaTeX input file.  (Version of 12 August 2004.)
%
% A '%' character causes TeX to ignore all remaining text on the line,
% and is used for comments like this one.

\documentclass{article}      % Specifies the document class

                             % The preamble begins here.
\title{An Example Document}  % Declares the document's title.
\author{Leslie Lamport}      % Declares the author's name.
\date{January 21, 1994}      % Deleting this command produces today's date.

\newcommand{\ip}[2]{(#1, #2)}
                             % Defines \ip{arg1}{arg2} to mean
                             % (arg1, arg2).

%\newcommand{\ip}[2]{\langle #1 | #2\rangle}
                             % This is an alternative definition of
                             % \ip that is commented out.

\begin{document}             % End of preamble and beginning of text.

\maketitle                   % Produces the title.

This is an example input file.  Comparing it with
the output it generates can show you how to
produce a simple document of your own.

\section{Ordinary Text}      % Produces section heading.  Lower-level
                             % sections are begun with similar
                             % \subsection and \subsubsection commands.

The ends  of words and sentences are marked
  by   spaces. It  doesn't matter how many
spaces    you type; one is as good as 100.  The
end of   a line counts as a space.

One   or more   blank lines denote the  end
of  a paragraph.

Since any number of consecutive spaces are treated
like a single one, the formatting of the input
file makes no difference to
      \LaTeX,                % The \LaTeX command generates the LaTeX logo.
but it makes a difference to you.  When you use
\LaTeX, making your input file as easy to read
as possible will be a great help as you write
your document and when you change it.  This sample
file shows how you can add comments to your own input
file.

Because printing is different from typewriting,
there are a number of things that you have to do
differently when preparing an input file than if
you were just typing the document directly.
Quotation marks like
       ``this''
have to be handled specially, as do quotes within
quotes:
       ``\,`this'            % \, separates the double and single quote.
        is what I just
        wrote, not  `that'\,''.

Dashes come in three sizes: an
       intra-word
dash, a medium dash for number ranges like
       1--2,
and a punctuation
       dash---like
this.

A sentence-ending space should be larger than the
space between words within a sentence.  You
sometimes have to type special commands in
conjunction with punctuation characters to get
this right, as in the following sentence.
       Gnats, gnus, etc.\ all  % `\ ' makes an inter-word space.
       begin with G\@.         % \@ marks end-of-sentence punctuation.
You should check the spaces after periods when
reading your output to make sure you haven't
forgotten any special cases.  Generating an
ellipsis
       \ldots\               % `\ ' is needed after `\ldots' because TeX
                             % ignores spaces after command names like \ldots
                             % made from \ + letters.
                             %
                             % Note how a `%' character causes TeX to ignore
                             % the end of the input line, so these blank lines
                             % do not start a new paragraph.
                             %
with the right spacing around the periods requires
a special command.

\LaTeX\ interprets some common characters as
commands, so you must type special commands to
generate them.  These characters include the
following:
       \$ \& \% \# \{ and \}.

In printing, text is usually emphasized with an
       \emph{italic}
type style.

\begin{em}
   A long segment of text can also be emphasized
   in this way.  Text within such a segment can be
   given \emph{additional} emphasis.
\end{em}

It is sometimes necessary to prevent \LaTeX\ from
breaking a line where it might otherwise do so.
This may be at a space, as between the ``Mr.''\ and
``Jones'' in
       ``Mr.~Jones'',        % ~ produces an unbreakable interword space.
or within a word---especially when the word is a
symbol like
       \mbox{\emph{itemnum}}
that makes little sense when hyphenated across
lines.

Footnotes\footnote{This is an example of a footnote.}
pose no problem.

\LaTeX\ is good at typesetting mathematical formulas
like
       \( x-3y + z = 7 \)
or
       \( a_{1} > x^{2n} + y^{2n} > x' \)
or
       \( \ip{A}{B} = \sum_{i} a_{i} b_{i} \).
The spaces you type in a formula are
ignored.  Remember that a letter like
       $x$                   % $ ... $  and  \( ... \)  are equivalent
is a formula when it denotes a mathematical
symbol, and it should be typed as one.

\section{Displayed Text}

Text is displayed by indenting it from the left
margin.  Quotations are commonly displayed.  There
are short quotations
\begin{quote}
   This is a short quotation.  It consists of a
   single paragraph of text.  See how it is formatted.
\end{quote}
and longer ones.
\begin{quotation}
   This is a longer quotation.  It consists of two
   paragraphs of text, neither of which are
   particularly interesting.

   This is the second paragraph of the quotation.  It
   is just as dull as the first paragraph.
\end{quotation}
Another frequently-displayed structure is a list.
The following is an example of an \emph{itemized}
list.
\begin{itemize}
   \item This is the first item of an itemized list.
         Each item in the list is marked with a ``tick''.
         You don't have to worry about what kind of tick
         mark is used.

   \item This is the second item of the list.  It
         contains another list nested inside it.  The inner
         list is an \emph{enumerated} list.
         \begin{enumerate}
            \item This is the first item of an enumerated
                  list that is nested within the itemized list.

            \item This is the second item of the inner list.
                  \LaTeX\ allows you to nest lists deeper than
                  you really should.
         \end{enumerate}
         This is the rest of the second item of the outer
         list.  It is no more interesting than any other
         part of the item.
   \item This is the third item of the list.
\end{itemize}
You can even display poetry.
\begin{verse}
   There is an environment
    for verse \\             % The \\ command separates lines
   Whose features some poets % within a stanza.
   will curse.

                             % One or more blank lines separate stanzas.

   For instead of making\\
   Them do \emph{all} line breaking, \\
   It allows them to put too many words on a line when they'd rather be
   forced to be terse.
\end{verse}

Mathematical formulas may also be displayed.  A
displayed formula
is
one-line long; multiline
formulas require special formatting instructions.
   \[  \ip{\Gamma}{\psi'} = x'' + y^{2} + z_{i}^{n}\]
Don't start a paragraph with a displayed equation,
nor make one a paragraph by itself.

\end{document}               % End of document.

\section{Bench Huge Section 249}
% synthetic bench section 249
% This is a sample LaTeX input file.  (Version of 12 August 2004.)
%
% A '%' character causes TeX to ignore all remaining text on the line,
% and is used for comments like this one.

\documentclass{article}      % Specifies the document class

                             % The preamble begins here.
\title{An Example Document}  % Declares the document's title.
\author{Leslie Lamport}      % Declares the author's name.
\date{January 21, 1994}      % Deleting this command produces today's date.

\newcommand{\ip}[2]{(#1, #2)}
                             % Defines \ip{arg1}{arg2} to mean
                             % (arg1, arg2).

%\newcommand{\ip}[2]{\langle #1 | #2\rangle}
                             % This is an alternative definition of
                             % \ip that is commented out.

\begin{document}             % End of preamble and beginning of text.

\maketitle                   % Produces the title.

This is an example input file.  Comparing it with
the output it generates can show you how to
produce a simple document of your own.

\section{Ordinary Text}      % Produces section heading.  Lower-level
                             % sections are begun with similar
                             % \subsection and \subsubsection commands.

The ends  of words and sentences are marked
  by   spaces. It  doesn't matter how many
spaces    you type; one is as good as 100.  The
end of   a line counts as a space.

One   or more   blank lines denote the  end
of  a paragraph.

Since any number of consecutive spaces are treated
like a single one, the formatting of the input
file makes no difference to
      \LaTeX,                % The \LaTeX command generates the LaTeX logo.
but it makes a difference to you.  When you use
\LaTeX, making your input file as easy to read
as possible will be a great help as you write
your document and when you change it.  This sample
file shows how you can add comments to your own input
file.

Because printing is different from typewriting,
there are a number of things that you have to do
differently when preparing an input file than if
you were just typing the document directly.
Quotation marks like
       ``this''
have to be handled specially, as do quotes within
quotes:
       ``\,`this'            % \, separates the double and single quote.
        is what I just
        wrote, not  `that'\,''.

Dashes come in three sizes: an
       intra-word
dash, a medium dash for number ranges like
       1--2,
and a punctuation
       dash---like
this.

A sentence-ending space should be larger than the
space between words within a sentence.  You
sometimes have to type special commands in
conjunction with punctuation characters to get
this right, as in the following sentence.
       Gnats, gnus, etc.\ all  % `\ ' makes an inter-word space.
       begin with G\@.         % \@ marks end-of-sentence punctuation.
You should check the spaces after periods when
reading your output to make sure you haven't
forgotten any special cases.  Generating an
ellipsis
       \ldots\               % `\ ' is needed after `\ldots' because TeX
                             % ignores spaces after command names like \ldots
                             % made from \ + letters.
                             %
                             % Note how a `%' character causes TeX to ignore
                             % the end of the input line, so these blank lines
                             % do not start a new paragraph.
                             %
with the right spacing around the periods requires
a special command.

\LaTeX\ interprets some common characters as
commands, so you must type special commands to
generate them.  These characters include the
following:
       \$ \& \% \# \{ and \}.

In printing, text is usually emphasized with an
       \emph{italic}
type style.

\begin{em}
   A long segment of text can also be emphasized
   in this way.  Text within such a segment can be
   given \emph{additional} emphasis.
\end{em}

It is sometimes necessary to prevent \LaTeX\ from
breaking a line where it might otherwise do so.
This may be at a space, as between the ``Mr.''\ and
``Jones'' in
       ``Mr.~Jones'',        % ~ produces an unbreakable interword space.
or within a word---especially when the word is a
symbol like
       \mbox{\emph{itemnum}}
that makes little sense when hyphenated across
lines.

Footnotes\footnote{This is an example of a footnote.}
pose no problem.

\LaTeX\ is good at typesetting mathematical formulas
like
       \( x-3y + z = 7 \)
or
       \( a_{1} > x^{2n} + y^{2n} > x' \)
or
       \( \ip{A}{B} = \sum_{i} a_{i} b_{i} \).
The spaces you type in a formula are
ignored.  Remember that a letter like
       $x$                   % $ ... $  and  \( ... \)  are equivalent
is a formula when it denotes a mathematical
symbol, and it should be typed as one.

\section{Displayed Text}

Text is displayed by indenting it from the left
margin.  Quotations are commonly displayed.  There
are short quotations
\begin{quote}
   This is a short quotation.  It consists of a
   single paragraph of text.  See how it is formatted.
\end{quote}
and longer ones.
\begin{quotation}
   This is a longer quotation.  It consists of two
   paragraphs of text, neither of which are
   particularly interesting.

   This is the second paragraph of the quotation.  It
   is just as dull as the first paragraph.
\end{quotation}
Another frequently-displayed structure is a list.
The following is an example of an \emph{itemized}
list.
\begin{itemize}
   \item This is the first item of an itemized list.
         Each item in the list is marked with a ``tick''.
         You don't have to worry about what kind of tick
         mark is used.

   \item This is the second item of the list.  It
         contains another list nested inside it.  The inner
         list is an \emph{enumerated} list.
         \begin{enumerate}
            \item This is the first item of an enumerated
                  list that is nested within the itemized list.

            \item This is the second item of the inner list.
                  \LaTeX\ allows you to nest lists deeper than
                  you really should.
         \end{enumerate}
         This is the rest of the second item of the outer
         list.  It is no more interesting than any other
         part of the item.
   \item This is the third item of the list.
\end{itemize}
You can even display poetry.
\begin{verse}
   There is an environment
    for verse \\             % The \\ command separates lines
   Whose features some poets % within a stanza.
   will curse.

                             % One or more blank lines separate stanzas.

   For instead of making\\
   Them do \emph{all} line breaking, \\
   It allows them to put too many words on a line when they'd rather be
   forced to be terse.
\end{verse}

Mathematical formulas may also be displayed.  A
displayed formula
is
one-line long; multiline
formulas require special formatting instructions.
   \[  \ip{\Gamma}{\psi'} = x'' + y^{2} + z_{i}^{n}\]
Don't start a paragraph with a displayed equation,
nor make one a paragraph by itself.

\end{document}               % End of document.

\section{Bench Huge Section 250}
% synthetic bench section 250
% This is a sample LaTeX input file.  (Version of 12 August 2004.)
%
% A '%' character causes TeX to ignore all remaining text on the line,
% and is used for comments like this one.

\documentclass{article}      % Specifies the document class

                             % The preamble begins here.
\title{An Example Document}  % Declares the document's title.
\author{Leslie Lamport}      % Declares the author's name.
\date{January 21, 1994}      % Deleting this command produces today's date.

\newcommand{\ip}[2]{(#1, #2)}
                             % Defines \ip{arg1}{arg2} to mean
                             % (arg1, arg2).

%\newcommand{\ip}[2]{\langle #1 | #2\rangle}
                             % This is an alternative definition of
                             % \ip that is commented out.

\begin{document}             % End of preamble and beginning of text.

\maketitle                   % Produces the title.

This is an example input file.  Comparing it with
the output it generates can show you how to
produce a simple document of your own.

\section{Ordinary Text}      % Produces section heading.  Lower-level
                             % sections are begun with similar
                             % \subsection and \subsubsection commands.

The ends  of words and sentences are marked
  by   spaces. It  doesn't matter how many
spaces    you type; one is as good as 100.  The
end of   a line counts as a space.

One   or more   blank lines denote the  end
of  a paragraph.

Since any number of consecutive spaces are treated
like a single one, the formatting of the input
file makes no difference to
      \LaTeX,                % The \LaTeX command generates the LaTeX logo.
but it makes a difference to you.  When you use
\LaTeX, making your input file as easy to read
as possible will be a great help as you write
your document and when you change it.  This sample
file shows how you can add comments to your own input
file.

Because printing is different from typewriting,
there are a number of things that you have to do
differently when preparing an input file than if
you were just typing the document directly.
Quotation marks like
       ``this''
have to be handled specially, as do quotes within
quotes:
       ``\,`this'            % \, separates the double and single quote.
        is what I just
        wrote, not  `that'\,''.

Dashes come in three sizes: an
       intra-word
dash, a medium dash for number ranges like
       1--2,
and a punctuation
       dash---like
this.

A sentence-ending space should be larger than the
space between words within a sentence.  You
sometimes have to type special commands in
conjunction with punctuation characters to get
this right, as in the following sentence.
       Gnats, gnus, etc.\ all  % `\ ' makes an inter-word space.
       begin with G\@.         % \@ marks end-of-sentence punctuation.
You should check the spaces after periods when
reading your output to make sure you haven't
forgotten any special cases.  Generating an
ellipsis
       \ldots\               % `\ ' is needed after `\ldots' because TeX
                             % ignores spaces after command names like \ldots
                             % made from \ + letters.
                             %
                             % Note how a `%' character causes TeX to ignore
                             % the end of the input line, so these blank lines
                             % do not start a new paragraph.
                             %
with the right spacing around the periods requires
a special command.

\LaTeX\ interprets some common characters as
commands, so you must type special commands to
generate them.  These characters include the
following:
       \$ \& \% \# \{ and \}.

In printing, text is usually emphasized with an
       \emph{italic}
type style.

\begin{em}
   A long segment of text can also be emphasized
   in this way.  Text within such a segment can be
   given \emph{additional} emphasis.
\end{em}

It is sometimes necessary to prevent \LaTeX\ from
breaking a line where it might otherwise do so.
This may be at a space, as between the ``Mr.''\ and
``Jones'' in
       ``Mr.~Jones'',        % ~ produces an unbreakable interword space.
or within a word---especially when the word is a
symbol like
       \mbox{\emph{itemnum}}
that makes little sense when hyphenated across
lines.

Footnotes\footnote{This is an example of a footnote.}
pose no problem.

\LaTeX\ is good at typesetting mathematical formulas
like
       \( x-3y + z = 7 \)
or
       \( a_{1} > x^{2n} + y^{2n} > x' \)
or
       \( \ip{A}{B} = \sum_{i} a_{i} b_{i} \).
The spaces you type in a formula are
ignored.  Remember that a letter like
       $x$                   % $ ... $  and  \( ... \)  are equivalent
is a formula when it denotes a mathematical
symbol, and it should be typed as one.

\section{Displayed Text}

Text is displayed by indenting it from the left
margin.  Quotations are commonly displayed.  There
are short quotations
\begin{quote}
   This is a short quotation.  It consists of a
   single paragraph of text.  See how it is formatted.
\end{quote}
and longer ones.
\begin{quotation}
   This is a longer quotation.  It consists of two
   paragraphs of text, neither of which are
   particularly interesting.

   This is the second paragraph of the quotation.  It
   is just as dull as the first paragraph.
\end{quotation}
Another frequently-displayed structure is a list.
The following is an example of an \emph{itemized}
list.
\begin{itemize}
   \item This is the first item of an itemized list.
         Each item in the list is marked with a ``tick''.
         You don't have to worry about what kind of tick
         mark is used.

   \item This is the second item of the list.  It
         contains another list nested inside it.  The inner
         list is an \emph{enumerated} list.
         \begin{enumerate}
            \item This is the first item of an enumerated
                  list that is nested within the itemized list.

            \item This is the second item of the inner list.
                  \LaTeX\ allows you to nest lists deeper than
                  you really should.
         \end{enumerate}
         This is the rest of the second item of the outer
         list.  It is no more interesting than any other
         part of the item.
   \item This is the third item of the list.
\end{itemize}
You can even display poetry.
\begin{verse}
   There is an environment
    for verse \\             % The \\ command separates lines
   Whose features some poets % within a stanza.
   will curse.

                             % One or more blank lines separate stanzas.

   For instead of making\\
   Them do \emph{all} line breaking, \\
   It allows them to put too many words on a line when they'd rather be
   forced to be terse.
\end{verse}

Mathematical formulas may also be displayed.  A
displayed formula
is
one-line long; multiline
formulas require special formatting instructions.
   \[  \ip{\Gamma}{\psi'} = x'' + y^{2} + z_{i}^{n}\]
Don't start a paragraph with a displayed equation,
nor make one a paragraph by itself.

\end{document}               % End of document.

\section{Bench Huge Section 251}
% synthetic bench section 251
% This is a sample LaTeX input file.  (Version of 12 August 2004.)
%
% A '%' character causes TeX to ignore all remaining text on the line,
% and is used for comments like this one.

\documentclass{article}      % Specifies the document class

                             % The preamble begins here.
\title{An Example Document}  % Declares the document's title.
\author{Leslie Lamport}      % Declares the author's name.
\date{January 21, 1994}      % Deleting this command produces today's date.

\newcommand{\ip}[2]{(#1, #2)}
                             % Defines \ip{arg1}{arg2} to mean
                             % (arg1, arg2).

%\newcommand{\ip}[2]{\langle #1 | #2\rangle}
                             % This is an alternative definition of
                             % \ip that is commented out.

\begin{document}             % End of preamble and beginning of text.

\maketitle                   % Produces the title.

This is an example input file.  Comparing it with
the output it generates can show you how to
produce a simple document of your own.

\section{Ordinary Text}      % Produces section heading.  Lower-level
                             % sections are begun with similar
                             % \subsection and \subsubsection commands.

The ends  of words and sentences are marked
  by   spaces. It  doesn't matter how many
spaces    you type; one is as good as 100.  The
end of   a line counts as a space.

One   or more   blank lines denote the  end
of  a paragraph.

Since any number of consecutive spaces are treated
like a single one, the formatting of the input
file makes no difference to
      \LaTeX,                % The \LaTeX command generates the LaTeX logo.
but it makes a difference to you.  When you use
\LaTeX, making your input file as easy to read
as possible will be a great help as you write
your document and when you change it.  This sample
file shows how you can add comments to your own input
file.

Because printing is different from typewriting,
there are a number of things that you have to do
differently when preparing an input file than if
you were just typing the document directly.
Quotation marks like
       ``this''
have to be handled specially, as do quotes within
quotes:
       ``\,`this'            % \, separates the double and single quote.
        is what I just
        wrote, not  `that'\,''.

Dashes come in three sizes: an
       intra-word
dash, a medium dash for number ranges like
       1--2,
and a punctuation
       dash---like
this.

A sentence-ending space should be larger than the
space between words within a sentence.  You
sometimes have to type special commands in
conjunction with punctuation characters to get
this right, as in the following sentence.
       Gnats, gnus, etc.\ all  % `\ ' makes an inter-word space.
       begin with G\@.         % \@ marks end-of-sentence punctuation.
You should check the spaces after periods when
reading your output to make sure you haven't
forgotten any special cases.  Generating an
ellipsis
       \ldots\               % `\ ' is needed after `\ldots' because TeX
                             % ignores spaces after command names like \ldots
                             % made from \ + letters.
                             %
                             % Note how a `%' character causes TeX to ignore
                             % the end of the input line, so these blank lines
                             % do not start a new paragraph.
                             %
with the right spacing around the periods requires
a special command.

\LaTeX\ interprets some common characters as
commands, so you must type special commands to
generate them.  These characters include the
following:
       \$ \& \% \# \{ and \}.

In printing, text is usually emphasized with an
       \emph{italic}
type style.

\begin{em}
   A long segment of text can also be emphasized
   in this way.  Text within such a segment can be
   given \emph{additional} emphasis.
\end{em}

It is sometimes necessary to prevent \LaTeX\ from
breaking a line where it might otherwise do so.
This may be at a space, as between the ``Mr.''\ and
``Jones'' in
       ``Mr.~Jones'',        % ~ produces an unbreakable interword space.
or within a word---especially when the word is a
symbol like
       \mbox{\emph{itemnum}}
that makes little sense when hyphenated across
lines.

Footnotes\footnote{This is an example of a footnote.}
pose no problem.

\LaTeX\ is good at typesetting mathematical formulas
like
       \( x-3y + z = 7 \)
or
       \( a_{1} > x^{2n} + y^{2n} > x' \)
or
       \( \ip{A}{B} = \sum_{i} a_{i} b_{i} \).
The spaces you type in a formula are
ignored.  Remember that a letter like
       $x$                   % $ ... $  and  \( ... \)  are equivalent
is a formula when it denotes a mathematical
symbol, and it should be typed as one.

\section{Displayed Text}

Text is displayed by indenting it from the left
margin.  Quotations are commonly displayed.  There
are short quotations
\begin{quote}
   This is a short quotation.  It consists of a
   single paragraph of text.  See how it is formatted.
\end{quote}
and longer ones.
\begin{quotation}
   This is a longer quotation.  It consists of two
   paragraphs of text, neither of which are
   particularly interesting.

   This is the second paragraph of the quotation.  It
   is just as dull as the first paragraph.
\end{quotation}
Another frequently-displayed structure is a list.
The following is an example of an \emph{itemized}
list.
\begin{itemize}
   \item This is the first item of an itemized list.
         Each item in the list is marked with a ``tick''.
         You don't have to worry about what kind of tick
         mark is used.

   \item This is the second item of the list.  It
         contains another list nested inside it.  The inner
         list is an \emph{enumerated} list.
         \begin{enumerate}
            \item This is the first item of an enumerated
                  list that is nested within the itemized list.

            \item This is the second item of the inner list.
                  \LaTeX\ allows you to nest lists deeper than
                  you really should.
         \end{enumerate}
         This is the rest of the second item of the outer
         list.  It is no more interesting than any other
         part of the item.
   \item This is the third item of the list.
\end{itemize}
You can even display poetry.
\begin{verse}
   There is an environment
    for verse \\             % The \\ command separates lines
   Whose features some poets % within a stanza.
   will curse.

                             % One or more blank lines separate stanzas.

   For instead of making\\
   Them do \emph{all} line breaking, \\
   It allows them to put too many words on a line when they'd rather be
   forced to be terse.
\end{verse}

Mathematical formulas may also be displayed.  A
displayed formula
is
one-line long; multiline
formulas require special formatting instructions.
   \[  \ip{\Gamma}{\psi'} = x'' + y^{2} + z_{i}^{n}\]
Don't start a paragraph with a displayed equation,
nor make one a paragraph by itself.

\end{document}               % End of document.

\section{Bench Huge Section 252}
% synthetic bench section 252
% This is a sample LaTeX input file.  (Version of 12 August 2004.)
%
% A '%' character causes TeX to ignore all remaining text on the line,
% and is used for comments like this one.

\documentclass{article}      % Specifies the document class

                             % The preamble begins here.
\title{An Example Document}  % Declares the document's title.
\author{Leslie Lamport}      % Declares the author's name.
\date{January 21, 1994}      % Deleting this command produces today's date.

\newcommand{\ip}[2]{(#1, #2)}
                             % Defines \ip{arg1}{arg2} to mean
                             % (arg1, arg2).

%\newcommand{\ip}[2]{\langle #1 | #2\rangle}
                             % This is an alternative definition of
                             % \ip that is commented out.

\begin{document}             % End of preamble and beginning of text.

\maketitle                   % Produces the title.

This is an example input file.  Comparing it with
the output it generates can show you how to
produce a simple document of your own.

\section{Ordinary Text}      % Produces section heading.  Lower-level
                             % sections are begun with similar
                             % \subsection and \subsubsection commands.

The ends  of words and sentences are marked
  by   spaces. It  doesn't matter how many
spaces    you type; one is as good as 100.  The
end of   a line counts as a space.

One   or more   blank lines denote the  end
of  a paragraph.

Since any number of consecutive spaces are treated
like a single one, the formatting of the input
file makes no difference to
      \LaTeX,                % The \LaTeX command generates the LaTeX logo.
but it makes a difference to you.  When you use
\LaTeX, making your input file as easy to read
as possible will be a great help as you write
your document and when you change it.  This sample
file shows how you can add comments to your own input
file.

Because printing is different from typewriting,
there are a number of things that you have to do
differently when preparing an input file than if
you were just typing the document directly.
Quotation marks like
       ``this''
have to be handled specially, as do quotes within
quotes:
       ``\,`this'            % \, separates the double and single quote.
        is what I just
        wrote, not  `that'\,''.

Dashes come in three sizes: an
       intra-word
dash, a medium dash for number ranges like
       1--2,
and a punctuation
       dash---like
this.

A sentence-ending space should be larger than the
space between words within a sentence.  You
sometimes have to type special commands in
conjunction with punctuation characters to get
this right, as in the following sentence.
       Gnats, gnus, etc.\ all  % `\ ' makes an inter-word space.
       begin with G\@.         % \@ marks end-of-sentence punctuation.
You should check the spaces after periods when
reading your output to make sure you haven't
forgotten any special cases.  Generating an
ellipsis
       \ldots\               % `\ ' is needed after `\ldots' because TeX
                             % ignores spaces after command names like \ldots
                             % made from \ + letters.
                             %
                             % Note how a `%' character causes TeX to ignore
                             % the end of the input line, so these blank lines
                             % do not start a new paragraph.
                             %
with the right spacing around the periods requires
a special command.

\LaTeX\ interprets some common characters as
commands, so you must type special commands to
generate them.  These characters include the
following:
       \$ \& \% \# \{ and \}.

In printing, text is usually emphasized with an
       \emph{italic}
type style.

\begin{em}
   A long segment of text can also be emphasized
   in this way.  Text within such a segment can be
   given \emph{additional} emphasis.
\end{em}

It is sometimes necessary to prevent \LaTeX\ from
breaking a line where it might otherwise do so.
This may be at a space, as between the ``Mr.''\ and
``Jones'' in
       ``Mr.~Jones'',        % ~ produces an unbreakable interword space.
or within a word---especially when the word is a
symbol like
       \mbox{\emph{itemnum}}
that makes little sense when hyphenated across
lines.

Footnotes\footnote{This is an example of a footnote.}
pose no problem.

\LaTeX\ is good at typesetting mathematical formulas
like
       \( x-3y + z = 7 \)
or
       \( a_{1} > x^{2n} + y^{2n} > x' \)
or
       \( \ip{A}{B} = \sum_{i} a_{i} b_{i} \).
The spaces you type in a formula are
ignored.  Remember that a letter like
       $x$                   % $ ... $  and  \( ... \)  are equivalent
is a formula when it denotes a mathematical
symbol, and it should be typed as one.

\section{Displayed Text}

Text is displayed by indenting it from the left
margin.  Quotations are commonly displayed.  There
are short quotations
\begin{quote}
   This is a short quotation.  It consists of a
   single paragraph of text.  See how it is formatted.
\end{quote}
and longer ones.
\begin{quotation}
   This is a longer quotation.  It consists of two
   paragraphs of text, neither of which are
   particularly interesting.

   This is the second paragraph of the quotation.  It
   is just as dull as the first paragraph.
\end{quotation}
Another frequently-displayed structure is a list.
The following is an example of an \emph{itemized}
list.
\begin{itemize}
   \item This is the first item of an itemized list.
         Each item in the list is marked with a ``tick''.
         You don't have to worry about what kind of tick
         mark is used.

   \item This is the second item of the list.  It
         contains another list nested inside it.  The inner
         list is an \emph{enumerated} list.
         \begin{enumerate}
            \item This is the first item of an enumerated
                  list that is nested within the itemized list.

            \item This is the second item of the inner list.
                  \LaTeX\ allows you to nest lists deeper than
                  you really should.
         \end{enumerate}
         This is the rest of the second item of the outer
         list.  It is no more interesting than any other
         part of the item.
   \item This is the third item of the list.
\end{itemize}
You can even display poetry.
\begin{verse}
   There is an environment
    for verse \\             % The \\ command separates lines
   Whose features some poets % within a stanza.
   will curse.

                             % One or more blank lines separate stanzas.

   For instead of making\\
   Them do \emph{all} line breaking, \\
   It allows them to put too many words on a line when they'd rather be
   forced to be terse.
\end{verse}

Mathematical formulas may also be displayed.  A
displayed formula
is
one-line long; multiline
formulas require special formatting instructions.
   \[  \ip{\Gamma}{\psi'} = x'' + y^{2} + z_{i}^{n}\]
Don't start a paragraph with a displayed equation,
nor make one a paragraph by itself.

\end{document}               % End of document.

\section{Bench Huge Section 253}
% synthetic bench section 253
% This is a sample LaTeX input file.  (Version of 12 August 2004.)
%
% A '%' character causes TeX to ignore all remaining text on the line,
% and is used for comments like this one.

\documentclass{article}      % Specifies the document class

                             % The preamble begins here.
\title{An Example Document}  % Declares the document's title.
\author{Leslie Lamport}      % Declares the author's name.
\date{January 21, 1994}      % Deleting this command produces today's date.

\newcommand{\ip}[2]{(#1, #2)}
                             % Defines \ip{arg1}{arg2} to mean
                             % (arg1, arg2).

%\newcommand{\ip}[2]{\langle #1 | #2\rangle}
                             % This is an alternative definition of
                             % \ip that is commented out.

\begin{document}             % End of preamble and beginning of text.

\maketitle                   % Produces the title.

This is an example input file.  Comparing it with
the output it generates can show you how to
produce a simple document of your own.

\section{Ordinary Text}      % Produces section heading.  Lower-level
                             % sections are begun with similar
                             % \subsection and \subsubsection commands.

The ends  of words and sentences are marked
  by   spaces. It  doesn't matter how many
spaces    you type; one is as good as 100.  The
end of   a line counts as a space.

One   or more   blank lines denote the  end
of  a paragraph.

Since any number of consecutive spaces are treated
like a single one, the formatting of the input
file makes no difference to
      \LaTeX,                % The \LaTeX command generates the LaTeX logo.
but it makes a difference to you.  When you use
\LaTeX, making your input file as easy to read
as possible will be a great help as you write
your document and when you change it.  This sample
file shows how you can add comments to your own input
file.

Because printing is different from typewriting,
there are a number of things that you have to do
differently when preparing an input file than if
you were just typing the document directly.
Quotation marks like
       ``this''
have to be handled specially, as do quotes within
quotes:
       ``\,`this'            % \, separates the double and single quote.
        is what I just
        wrote, not  `that'\,''.

Dashes come in three sizes: an
       intra-word
dash, a medium dash for number ranges like
       1--2,
and a punctuation
       dash---like
this.

A sentence-ending space should be larger than the
space between words within a sentence.  You
sometimes have to type special commands in
conjunction with punctuation characters to get
this right, as in the following sentence.
       Gnats, gnus, etc.\ all  % `\ ' makes an inter-word space.
       begin with G\@.         % \@ marks end-of-sentence punctuation.
You should check the spaces after periods when
reading your output to make sure you haven't
forgotten any special cases.  Generating an
ellipsis
       \ldots\               % `\ ' is needed after `\ldots' because TeX
                             % ignores spaces after command names like \ldots
                             % made from \ + letters.
                             %
                             % Note how a `%' character causes TeX to ignore
                             % the end of the input line, so these blank lines
                             % do not start a new paragraph.
                             %
with the right spacing around the periods requires
a special command.

\LaTeX\ interprets some common characters as
commands, so you must type special commands to
generate them.  These characters include the
following:
       \$ \& \% \# \{ and \}.

In printing, text is usually emphasized with an
       \emph{italic}
type style.

\begin{em}
   A long segment of text can also be emphasized
   in this way.  Text within such a segment can be
   given \emph{additional} emphasis.
\end{em}

It is sometimes necessary to prevent \LaTeX\ from
breaking a line where it might otherwise do so.
This may be at a space, as between the ``Mr.''\ and
``Jones'' in
       ``Mr.~Jones'',        % ~ produces an unbreakable interword space.
or within a word---especially when the word is a
symbol like
       \mbox{\emph{itemnum}}
that makes little sense when hyphenated across
lines.

Footnotes\footnote{This is an example of a footnote.}
pose no problem.

\LaTeX\ is good at typesetting mathematical formulas
like
       \( x-3y + z = 7 \)
or
       \( a_{1} > x^{2n} + y^{2n} > x' \)
or
       \( \ip{A}{B} = \sum_{i} a_{i} b_{i} \).
The spaces you type in a formula are
ignored.  Remember that a letter like
       $x$                   % $ ... $  and  \( ... \)  are equivalent
is a formula when it denotes a mathematical
symbol, and it should be typed as one.

\section{Displayed Text}

Text is displayed by indenting it from the left
margin.  Quotations are commonly displayed.  There
are short quotations
\begin{quote}
   This is a short quotation.  It consists of a
   single paragraph of text.  See how it is formatted.
\end{quote}
and longer ones.
\begin{quotation}
   This is a longer quotation.  It consists of two
   paragraphs of text, neither of which are
   particularly interesting.

   This is the second paragraph of the quotation.  It
   is just as dull as the first paragraph.
\end{quotation}
Another frequently-displayed structure is a list.
The following is an example of an \emph{itemized}
list.
\begin{itemize}
   \item This is the first item of an itemized list.
         Each item in the list is marked with a ``tick''.
         You don't have to worry about what kind of tick
         mark is used.

   \item This is the second item of the list.  It
         contains another list nested inside it.  The inner
         list is an \emph{enumerated} list.
         \begin{enumerate}
            \item This is the first item of an enumerated
                  list that is nested within the itemized list.

            \item This is the second item of the inner list.
                  \LaTeX\ allows you to nest lists deeper than
                  you really should.
         \end{enumerate}
         This is the rest of the second item of the outer
         list.  It is no more interesting than any other
         part of the item.
   \item This is the third item of the list.
\end{itemize}
You can even display poetry.
\begin{verse}
   There is an environment
    for verse \\             % The \\ command separates lines
   Whose features some poets % within a stanza.
   will curse.

                             % One or more blank lines separate stanzas.

   For instead of making\\
   Them do \emph{all} line breaking, \\
   It allows them to put too many words on a line when they'd rather be
   forced to be terse.
\end{verse}

Mathematical formulas may also be displayed.  A
displayed formula
is
one-line long; multiline
formulas require special formatting instructions.
   \[  \ip{\Gamma}{\psi'} = x'' + y^{2} + z_{i}^{n}\]
Don't start a paragraph with a displayed equation,
nor make one a paragraph by itself.

\end{document}               % End of document.

\section{Bench Huge Section 254}
% synthetic bench section 254
% This is a sample LaTeX input file.  (Version of 12 August 2004.)
%
% A '%' character causes TeX to ignore all remaining text on the line,
% and is used for comments like this one.

\documentclass{article}      % Specifies the document class

                             % The preamble begins here.
\title{An Example Document}  % Declares the document's title.
\author{Leslie Lamport}      % Declares the author's name.
\date{January 21, 1994}      % Deleting this command produces today's date.

\newcommand{\ip}[2]{(#1, #2)}
                             % Defines \ip{arg1}{arg2} to mean
                             % (arg1, arg2).

%\newcommand{\ip}[2]{\langle #1 | #2\rangle}
                             % This is an alternative definition of
                             % \ip that is commented out.

\begin{document}             % End of preamble and beginning of text.

\maketitle                   % Produces the title.

This is an example input file.  Comparing it with
the output it generates can show you how to
produce a simple document of your own.

\section{Ordinary Text}      % Produces section heading.  Lower-level
                             % sections are begun with similar
                             % \subsection and \subsubsection commands.

The ends  of words and sentences are marked
  by   spaces. It  doesn't matter how many
spaces    you type; one is as good as 100.  The
end of   a line counts as a space.

One   or more   blank lines denote the  end
of  a paragraph.

Since any number of consecutive spaces are treated
like a single one, the formatting of the input
file makes no difference to
      \LaTeX,                % The \LaTeX command generates the LaTeX logo.
but it makes a difference to you.  When you use
\LaTeX, making your input file as easy to read
as possible will be a great help as you write
your document and when you change it.  This sample
file shows how you can add comments to your own input
file.

Because printing is different from typewriting,
there are a number of things that you have to do
differently when preparing an input file than if
you were just typing the document directly.
Quotation marks like
       ``this''
have to be handled specially, as do quotes within
quotes:
       ``\,`this'            % \, separates the double and single quote.
        is what I just
        wrote, not  `that'\,''.

Dashes come in three sizes: an
       intra-word
dash, a medium dash for number ranges like
       1--2,
and a punctuation
       dash---like
this.

A sentence-ending space should be larger than the
space between words within a sentence.  You
sometimes have to type special commands in
conjunction with punctuation characters to get
this right, as in the following sentence.
       Gnats, gnus, etc.\ all  % `\ ' makes an inter-word space.
       begin with G\@.         % \@ marks end-of-sentence punctuation.
You should check the spaces after periods when
reading your output to make sure you haven't
forgotten any special cases.  Generating an
ellipsis
       \ldots\               % `\ ' is needed after `\ldots' because TeX
                             % ignores spaces after command names like \ldots
                             % made from \ + letters.
                             %
                             % Note how a `%' character causes TeX to ignore
                             % the end of the input line, so these blank lines
                             % do not start a new paragraph.
                             %
with the right spacing around the periods requires
a special command.

\LaTeX\ interprets some common characters as
commands, so you must type special commands to
generate them.  These characters include the
following:
       \$ \& \% \# \{ and \}.

In printing, text is usually emphasized with an
       \emph{italic}
type style.

\begin{em}
   A long segment of text can also be emphasized
   in this way.  Text within such a segment can be
   given \emph{additional} emphasis.
\end{em}

It is sometimes necessary to prevent \LaTeX\ from
breaking a line where it might otherwise do so.
This may be at a space, as between the ``Mr.''\ and
``Jones'' in
       ``Mr.~Jones'',        % ~ produces an unbreakable interword space.
or within a word---especially when the word is a
symbol like
       \mbox{\emph{itemnum}}
that makes little sense when hyphenated across
lines.

Footnotes\footnote{This is an example of a footnote.}
pose no problem.

\LaTeX\ is good at typesetting mathematical formulas
like
       \( x-3y + z = 7 \)
or
       \( a_{1} > x^{2n} + y^{2n} > x' \)
or
       \( \ip{A}{B} = \sum_{i} a_{i} b_{i} \).
The spaces you type in a formula are
ignored.  Remember that a letter like
       $x$                   % $ ... $  and  \( ... \)  are equivalent
is a formula when it denotes a mathematical
symbol, and it should be typed as one.

\section{Displayed Text}

Text is displayed by indenting it from the left
margin.  Quotations are commonly displayed.  There
are short quotations
\begin{quote}
   This is a short quotation.  It consists of a
   single paragraph of text.  See how it is formatted.
\end{quote}
and longer ones.
\begin{quotation}
   This is a longer quotation.  It consists of two
   paragraphs of text, neither of which are
   particularly interesting.

   This is the second paragraph of the quotation.  It
   is just as dull as the first paragraph.
\end{quotation}
Another frequently-displayed structure is a list.
The following is an example of an \emph{itemized}
list.
\begin{itemize}
   \item This is the first item of an itemized list.
         Each item in the list is marked with a ``tick''.
         You don't have to worry about what kind of tick
         mark is used.

   \item This is the second item of the list.  It
         contains another list nested inside it.  The inner
         list is an \emph{enumerated} list.
         \begin{enumerate}
            \item This is the first item of an enumerated
                  list that is nested within the itemized list.

            \item This is the second item of the inner list.
                  \LaTeX\ allows you to nest lists deeper than
                  you really should.
         \end{enumerate}
         This is the rest of the second item of the outer
         list.  It is no more interesting than any other
         part of the item.
   \item This is the third item of the list.
\end{itemize}
You can even display poetry.
\begin{verse}
   There is an environment
    for verse \\             % The \\ command separates lines
   Whose features some poets % within a stanza.
   will curse.

                             % One or more blank lines separate stanzas.

   For instead of making\\
   Them do \emph{all} line breaking, \\
   It allows them to put too many words on a line when they'd rather be
   forced to be terse.
\end{verse}

Mathematical formulas may also be displayed.  A
displayed formula
is
one-line long; multiline
formulas require special formatting instructions.
   \[  \ip{\Gamma}{\psi'} = x'' + y^{2} + z_{i}^{n}\]
Don't start a paragraph with a displayed equation,
nor make one a paragraph by itself.

\end{document}               % End of document.

\section{Bench Huge Section 255}
% synthetic bench section 255
% This is a sample LaTeX input file.  (Version of 12 August 2004.)
%
% A '%' character causes TeX to ignore all remaining text on the line,
% and is used for comments like this one.

\documentclass{article}      % Specifies the document class

                             % The preamble begins here.
\title{An Example Document}  % Declares the document's title.
\author{Leslie Lamport}      % Declares the author's name.
\date{January 21, 1994}      % Deleting this command produces today's date.

\newcommand{\ip}[2]{(#1, #2)}
                             % Defines \ip{arg1}{arg2} to mean
                             % (arg1, arg2).

%\newcommand{\ip}[2]{\langle #1 | #2\rangle}
                             % This is an alternative definition of
                             % \ip that is commented out.

\begin{document}             % End of preamble and beginning of text.

\maketitle                   % Produces the title.

This is an example input file.  Comparing it with
the output it generates can show you how to
produce a simple document of your own.

\section{Ordinary Text}      % Produces section heading.  Lower-level
                             % sections are begun with similar
                             % \subsection and \subsubsection commands.

The ends  of words and sentences are marked
  by   spaces. It  doesn't matter how many
spaces    you type; one is as good as 100.  The
end of   a line counts as a space.

One   or more   blank lines denote the  end
of  a paragraph.

Since any number of consecutive spaces are treated
like a single one, the formatting of the input
file makes no difference to
      \LaTeX,                % The \LaTeX command generates the LaTeX logo.
but it makes a difference to you.  When you use
\LaTeX, making your input file as easy to read
as possible will be a great help as you write
your document and when you change it.  This sample
file shows how you can add comments to your own input
file.

Because printing is different from typewriting,
there are a number of things that you have to do
differently when preparing an input file than if
you were just typing the document directly.
Quotation marks like
       ``this''
have to be handled specially, as do quotes within
quotes:
       ``\,`this'            % \, separates the double and single quote.
        is what I just
        wrote, not  `that'\,''.

Dashes come in three sizes: an
       intra-word
dash, a medium dash for number ranges like
       1--2,
and a punctuation
       dash---like
this.

A sentence-ending space should be larger than the
space between words within a sentence.  You
sometimes have to type special commands in
conjunction with punctuation characters to get
this right, as in the following sentence.
       Gnats, gnus, etc.\ all  % `\ ' makes an inter-word space.
       begin with G\@.         % \@ marks end-of-sentence punctuation.
You should check the spaces after periods when
reading your output to make sure you haven't
forgotten any special cases.  Generating an
ellipsis
       \ldots\               % `\ ' is needed after `\ldots' because TeX
                             % ignores spaces after command names like \ldots
                             % made from \ + letters.
                             %
                             % Note how a `%' character causes TeX to ignore
                             % the end of the input line, so these blank lines
                             % do not start a new paragraph.
                             %
with the right spacing around the periods requires
a special command.

\LaTeX\ interprets some common characters as
commands, so you must type special commands to
generate them.  These characters include the
following:
       \$ \& \% \# \{ and \}.

In printing, text is usually emphasized with an
       \emph{italic}
type style.

\begin{em}
   A long segment of text can also be emphasized
   in this way.  Text within such a segment can be
   given \emph{additional} emphasis.
\end{em}

It is sometimes necessary to prevent \LaTeX\ from
breaking a line where it might otherwise do so.
This may be at a space, as between the ``Mr.''\ and
``Jones'' in
       ``Mr.~Jones'',        % ~ produces an unbreakable interword space.
or within a word---especially when the word is a
symbol like
       \mbox{\emph{itemnum}}
that makes little sense when hyphenated across
lines.

Footnotes\footnote{This is an example of a footnote.}
pose no problem.

\LaTeX\ is good at typesetting mathematical formulas
like
       \( x-3y + z = 7 \)
or
       \( a_{1} > x^{2n} + y^{2n} > x' \)
or
       \( \ip{A}{B} = \sum_{i} a_{i} b_{i} \).
The spaces you type in a formula are
ignored.  Remember that a letter like
       $x$                   % $ ... $  and  \( ... \)  are equivalent
is a formula when it denotes a mathematical
symbol, and it should be typed as one.

\section{Displayed Text}

Text is displayed by indenting it from the left
margin.  Quotations are commonly displayed.  There
are short quotations
\begin{quote}
   This is a short quotation.  It consists of a
   single paragraph of text.  See how it is formatted.
\end{quote}
and longer ones.
\begin{quotation}
   This is a longer quotation.  It consists of two
   paragraphs of text, neither of which are
   particularly interesting.

   This is the second paragraph of the quotation.  It
   is just as dull as the first paragraph.
\end{quotation}
Another frequently-displayed structure is a list.
The following is an example of an \emph{itemized}
list.
\begin{itemize}
   \item This is the first item of an itemized list.
         Each item in the list is marked with a ``tick''.
         You don't have to worry about what kind of tick
         mark is used.

   \item This is the second item of the list.  It
         contains another list nested inside it.  The inner
         list is an \emph{enumerated} list.
         \begin{enumerate}
            \item This is the first item of an enumerated
                  list that is nested within the itemized list.

            \item This is the second item of the inner list.
                  \LaTeX\ allows you to nest lists deeper than
                  you really should.
         \end{enumerate}
         This is the rest of the second item of the outer
         list.  It is no more interesting than any other
         part of the item.
   \item This is the third item of the list.
\end{itemize}
You can even display poetry.
\begin{verse}
   There is an environment
    for verse \\             % The \\ command separates lines
   Whose features some poets % within a stanza.
   will curse.

                             % One or more blank lines separate stanzas.

   For instead of making\\
   Them do \emph{all} line breaking, \\
   It allows them to put too many words on a line when they'd rather be
   forced to be terse.
\end{verse}

Mathematical formulas may also be displayed.  A
displayed formula
is
one-line long; multiline
formulas require special formatting instructions.
   \[  \ip{\Gamma}{\psi'} = x'' + y^{2} + z_{i}^{n}\]
Don't start a paragraph with a displayed equation,
nor make one a paragraph by itself.

\end{document}               % End of document.

\section{Bench Huge Section 256}
% synthetic bench section 256
% This is a sample LaTeX input file.  (Version of 12 August 2004.)
%
% A '%' character causes TeX to ignore all remaining text on the line,
% and is used for comments like this one.

\documentclass{article}      % Specifies the document class

                             % The preamble begins here.
\title{An Example Document}  % Declares the document's title.
\author{Leslie Lamport}      % Declares the author's name.
\date{January 21, 1994}      % Deleting this command produces today's date.

\newcommand{\ip}[2]{(#1, #2)}
                             % Defines \ip{arg1}{arg2} to mean
                             % (arg1, arg2).

%\newcommand{\ip}[2]{\langle #1 | #2\rangle}
                             % This is an alternative definition of
                             % \ip that is commented out.

\begin{document}             % End of preamble and beginning of text.

\maketitle                   % Produces the title.

This is an example input file.  Comparing it with
the output it generates can show you how to
produce a simple document of your own.

\section{Ordinary Text}      % Produces section heading.  Lower-level
                             % sections are begun with similar
                             % \subsection and \subsubsection commands.

The ends  of words and sentences are marked
  by   spaces. It  doesn't matter how many
spaces    you type; one is as good as 100.  The
end of   a line counts as a space.

One   or more   blank lines denote the  end
of  a paragraph.

Since any number of consecutive spaces are treated
like a single one, the formatting of the input
file makes no difference to
      \LaTeX,                % The \LaTeX command generates the LaTeX logo.
but it makes a difference to you.  When you use
\LaTeX, making your input file as easy to read
as possible will be a great help as you write
your document and when you change it.  This sample
file shows how you can add comments to your own input
file.

Because printing is different from typewriting,
there are a number of things that you have to do
differently when preparing an input file than if
you were just typing the document directly.
Quotation marks like
       ``this''
have to be handled specially, as do quotes within
quotes:
       ``\,`this'            % \, separates the double and single quote.
        is what I just
        wrote, not  `that'\,''.

Dashes come in three sizes: an
       intra-word
dash, a medium dash for number ranges like
       1--2,
and a punctuation
       dash---like
this.

A sentence-ending space should be larger than the
space between words within a sentence.  You
sometimes have to type special commands in
conjunction with punctuation characters to get
this right, as in the following sentence.
       Gnats, gnus, etc.\ all  % `\ ' makes an inter-word space.
       begin with G\@.         % \@ marks end-of-sentence punctuation.
You should check the spaces after periods when
reading your output to make sure you haven't
forgotten any special cases.  Generating an
ellipsis
       \ldots\               % `\ ' is needed after `\ldots' because TeX
                             % ignores spaces after command names like \ldots
                             % made from \ + letters.
                             %
                             % Note how a `%' character causes TeX to ignore
                             % the end of the input line, so these blank lines
                             % do not start a new paragraph.
                             %
with the right spacing around the periods requires
a special command.

\LaTeX\ interprets some common characters as
commands, so you must type special commands to
generate them.  These characters include the
following:
       \$ \& \% \# \{ and \}.

In printing, text is usually emphasized with an
       \emph{italic}
type style.

\begin{em}
   A long segment of text can also be emphasized
   in this way.  Text within such a segment can be
   given \emph{additional} emphasis.
\end{em}

It is sometimes necessary to prevent \LaTeX\ from
breaking a line where it might otherwise do so.
This may be at a space, as between the ``Mr.''\ and
``Jones'' in
       ``Mr.~Jones'',        % ~ produces an unbreakable interword space.
or within a word---especially when the word is a
symbol like
       \mbox{\emph{itemnum}}
that makes little sense when hyphenated across
lines.

Footnotes\footnote{This is an example of a footnote.}
pose no problem.

\LaTeX\ is good at typesetting mathematical formulas
like
       \( x-3y + z = 7 \)
or
       \( a_{1} > x^{2n} + y^{2n} > x' \)
or
       \( \ip{A}{B} = \sum_{i} a_{i} b_{i} \).
The spaces you type in a formula are
ignored.  Remember that a letter like
       $x$                   % $ ... $  and  \( ... \)  are equivalent
is a formula when it denotes a mathematical
symbol, and it should be typed as one.

\section{Displayed Text}

Text is displayed by indenting it from the left
margin.  Quotations are commonly displayed.  There
are short quotations
\begin{quote}
   This is a short quotation.  It consists of a
   single paragraph of text.  See how it is formatted.
\end{quote}
and longer ones.
\begin{quotation}
   This is a longer quotation.  It consists of two
   paragraphs of text, neither of which are
   particularly interesting.

   This is the second paragraph of the quotation.  It
   is just as dull as the first paragraph.
\end{quotation}
Another frequently-displayed structure is a list.
The following is an example of an \emph{itemized}
list.
\begin{itemize}
   \item This is the first item of an itemized list.
         Each item in the list is marked with a ``tick''.
         You don't have to worry about what kind of tick
         mark is used.

   \item This is the second item of the list.  It
         contains another list nested inside it.  The inner
         list is an \emph{enumerated} list.
         \begin{enumerate}
            \item This is the first item of an enumerated
                  list that is nested within the itemized list.

            \item This is the second item of the inner list.
                  \LaTeX\ allows you to nest lists deeper than
                  you really should.
         \end{enumerate}
         This is the rest of the second item of the outer
         list.  It is no more interesting than any other
         part of the item.
   \item This is the third item of the list.
\end{itemize}
You can even display poetry.
\begin{verse}
   There is an environment
    for verse \\             % The \\ command separates lines
   Whose features some poets % within a stanza.
   will curse.

                             % One or more blank lines separate stanzas.

   For instead of making\\
   Them do \emph{all} line breaking, \\
   It allows them to put too many words on a line when they'd rather be
   forced to be terse.
\end{verse}

Mathematical formulas may also be displayed.  A
displayed formula
is
one-line long; multiline
formulas require special formatting instructions.
   \[  \ip{\Gamma}{\psi'} = x'' + y^{2} + z_{i}^{n}\]
Don't start a paragraph with a displayed equation,
nor make one a paragraph by itself.

\end{document}               % End of document.

\section{Bench Huge Section 257}
% synthetic bench section 257
% This is a sample LaTeX input file.  (Version of 12 August 2004.)
%
% A '%' character causes TeX to ignore all remaining text on the line,
% and is used for comments like this one.

\documentclass{article}      % Specifies the document class

                             % The preamble begins here.
\title{An Example Document}  % Declares the document's title.
\author{Leslie Lamport}      % Declares the author's name.
\date{January 21, 1994}      % Deleting this command produces today's date.

\newcommand{\ip}[2]{(#1, #2)}
                             % Defines \ip{arg1}{arg2} to mean
                             % (arg1, arg2).

%\newcommand{\ip}[2]{\langle #1 | #2\rangle}
                             % This is an alternative definition of
                             % \ip that is commented out.

\begin{document}             % End of preamble and beginning of text.

\maketitle                   % Produces the title.

This is an example input file.  Comparing it with
the output it generates can show you how to
produce a simple document of your own.

\section{Ordinary Text}      % Produces section heading.  Lower-level
                             % sections are begun with similar
                             % \subsection and \subsubsection commands.

The ends  of words and sentences are marked
  by   spaces. It  doesn't matter how many
spaces    you type; one is as good as 100.  The
end of   a line counts as a space.

One   or more   blank lines denote the  end
of  a paragraph.

Since any number of consecutive spaces are treated
like a single one, the formatting of the input
file makes no difference to
      \LaTeX,                % The \LaTeX command generates the LaTeX logo.
but it makes a difference to you.  When you use
\LaTeX, making your input file as easy to read
as possible will be a great help as you write
your document and when you change it.  This sample
file shows how you can add comments to your own input
file.

Because printing is different from typewriting,
there are a number of things that you have to do
differently when preparing an input file than if
you were just typing the document directly.
Quotation marks like
       ``this''
have to be handled specially, as do quotes within
quotes:
       ``\,`this'            % \, separates the double and single quote.
        is what I just
        wrote, not  `that'\,''.

Dashes come in three sizes: an
       intra-word
dash, a medium dash for number ranges like
       1--2,
and a punctuation
       dash---like
this.

A sentence-ending space should be larger than the
space between words within a sentence.  You
sometimes have to type special commands in
conjunction with punctuation characters to get
this right, as in the following sentence.
       Gnats, gnus, etc.\ all  % `\ ' makes an inter-word space.
       begin with G\@.         % \@ marks end-of-sentence punctuation.
You should check the spaces after periods when
reading your output to make sure you haven't
forgotten any special cases.  Generating an
ellipsis
       \ldots\               % `\ ' is needed after `\ldots' because TeX
                             % ignores spaces after command names like \ldots
                             % made from \ + letters.
                             %
                             % Note how a `%' character causes TeX to ignore
                             % the end of the input line, so these blank lines
                             % do not start a new paragraph.
                             %
with the right spacing around the periods requires
a special command.

\LaTeX\ interprets some common characters as
commands, so you must type special commands to
generate them.  These characters include the
following:
       \$ \& \% \# \{ and \}.

In printing, text is usually emphasized with an
       \emph{italic}
type style.

\begin{em}
   A long segment of text can also be emphasized
   in this way.  Text within such a segment can be
   given \emph{additional} emphasis.
\end{em}

It is sometimes necessary to prevent \LaTeX\ from
breaking a line where it might otherwise do so.
This may be at a space, as between the ``Mr.''\ and
``Jones'' in
       ``Mr.~Jones'',        % ~ produces an unbreakable interword space.
or within a word---especially when the word is a
symbol like
       \mbox{\emph{itemnum}}
that makes little sense when hyphenated across
lines.

Footnotes\footnote{This is an example of a footnote.}
pose no problem.

\LaTeX\ is good at typesetting mathematical formulas
like
       \( x-3y + z = 7 \)
or
       \( a_{1} > x^{2n} + y^{2n} > x' \)
or
       \( \ip{A}{B} = \sum_{i} a_{i} b_{i} \).
The spaces you type in a formula are
ignored.  Remember that a letter like
       $x$                   % $ ... $  and  \( ... \)  are equivalent
is a formula when it denotes a mathematical
symbol, and it should be typed as one.

\section{Displayed Text}

Text is displayed by indenting it from the left
margin.  Quotations are commonly displayed.  There
are short quotations
\begin{quote}
   This is a short quotation.  It consists of a
   single paragraph of text.  See how it is formatted.
\end{quote}
and longer ones.
\begin{quotation}
   This is a longer quotation.  It consists of two
   paragraphs of text, neither of which are
   particularly interesting.

   This is the second paragraph of the quotation.  It
   is just as dull as the first paragraph.
\end{quotation}
Another frequently-displayed structure is a list.
The following is an example of an \emph{itemized}
list.
\begin{itemize}
   \item This is the first item of an itemized list.
         Each item in the list is marked with a ``tick''.
         You don't have to worry about what kind of tick
         mark is used.

   \item This is the second item of the list.  It
         contains another list nested inside it.  The inner
         list is an \emph{enumerated} list.
         \begin{enumerate}
            \item This is the first item of an enumerated
                  list that is nested within the itemized list.

            \item This is the second item of the inner list.
                  \LaTeX\ allows you to nest lists deeper than
                  you really should.
         \end{enumerate}
         This is the rest of the second item of the outer
         list.  It is no more interesting than any other
         part of the item.
   \item This is the third item of the list.
\end{itemize}
You can even display poetry.
\begin{verse}
   There is an environment
    for verse \\             % The \\ command separates lines
   Whose features some poets % within a stanza.
   will curse.

                             % One or more blank lines separate stanzas.

   For instead of making\\
   Them do \emph{all} line breaking, \\
   It allows them to put too many words on a line when they'd rather be
   forced to be terse.
\end{verse}

Mathematical formulas may also be displayed.  A
displayed formula
is
one-line long; multiline
formulas require special formatting instructions.
   \[  \ip{\Gamma}{\psi'} = x'' + y^{2} + z_{i}^{n}\]
Don't start a paragraph with a displayed equation,
nor make one a paragraph by itself.

\end{document}               % End of document.

\section{Bench Huge Section 258}
% synthetic bench section 258
% This is a sample LaTeX input file.  (Version of 12 August 2004.)
%
% A '%' character causes TeX to ignore all remaining text on the line,
% and is used for comments like this one.

\documentclass{article}      % Specifies the document class

                             % The preamble begins here.
\title{An Example Document}  % Declares the document's title.
\author{Leslie Lamport}      % Declares the author's name.
\date{January 21, 1994}      % Deleting this command produces today's date.

\newcommand{\ip}[2]{(#1, #2)}
                             % Defines \ip{arg1}{arg2} to mean
                             % (arg1, arg2).

%\newcommand{\ip}[2]{\langle #1 | #2\rangle}
                             % This is an alternative definition of
                             % \ip that is commented out.

\begin{document}             % End of preamble and beginning of text.

\maketitle                   % Produces the title.

This is an example input file.  Comparing it with
the output it generates can show you how to
produce a simple document of your own.

\section{Ordinary Text}      % Produces section heading.  Lower-level
                             % sections are begun with similar
                             % \subsection and \subsubsection commands.

The ends  of words and sentences are marked
  by   spaces. It  doesn't matter how many
spaces    you type; one is as good as 100.  The
end of   a line counts as a space.

One   or more   blank lines denote the  end
of  a paragraph.

Since any number of consecutive spaces are treated
like a single one, the formatting of the input
file makes no difference to
      \LaTeX,                % The \LaTeX command generates the LaTeX logo.
but it makes a difference to you.  When you use
\LaTeX, making your input file as easy to read
as possible will be a great help as you write
your document and when you change it.  This sample
file shows how you can add comments to your own input
file.

Because printing is different from typewriting,
there are a number of things that you have to do
differently when preparing an input file than if
you were just typing the document directly.
Quotation marks like
       ``this''
have to be handled specially, as do quotes within
quotes:
       ``\,`this'            % \, separates the double and single quote.
        is what I just
        wrote, not  `that'\,''.

Dashes come in three sizes: an
       intra-word
dash, a medium dash for number ranges like
       1--2,
and a punctuation
       dash---like
this.

A sentence-ending space should be larger than the
space between words within a sentence.  You
sometimes have to type special commands in
conjunction with punctuation characters to get
this right, as in the following sentence.
       Gnats, gnus, etc.\ all  % `\ ' makes an inter-word space.
       begin with G\@.         % \@ marks end-of-sentence punctuation.
You should check the spaces after periods when
reading your output to make sure you haven't
forgotten any special cases.  Generating an
ellipsis
       \ldots\               % `\ ' is needed after `\ldots' because TeX
                             % ignores spaces after command names like \ldots
                             % made from \ + letters.
                             %
                             % Note how a `%' character causes TeX to ignore
                             % the end of the input line, so these blank lines
                             % do not start a new paragraph.
                             %
with the right spacing around the periods requires
a special command.

\LaTeX\ interprets some common characters as
commands, so you must type special commands to
generate them.  These characters include the
following:
       \$ \& \% \# \{ and \}.

In printing, text is usually emphasized with an
       \emph{italic}
type style.

\begin{em}
   A long segment of text can also be emphasized
   in this way.  Text within such a segment can be
   given \emph{additional} emphasis.
\end{em}

It is sometimes necessary to prevent \LaTeX\ from
breaking a line where it might otherwise do so.
This may be at a space, as between the ``Mr.''\ and
``Jones'' in
       ``Mr.~Jones'',        % ~ produces an unbreakable interword space.
or within a word---especially when the word is a
symbol like
       \mbox{\emph{itemnum}}
that makes little sense when hyphenated across
lines.

Footnotes\footnote{This is an example of a footnote.}
pose no problem.

\LaTeX\ is good at typesetting mathematical formulas
like
       \( x-3y + z = 7 \)
or
       \( a_{1} > x^{2n} + y^{2n} > x' \)
or
       \( \ip{A}{B} = \sum_{i} a_{i} b_{i} \).
The spaces you type in a formula are
ignored.  Remember that a letter like
       $x$                   % $ ... $  and  \( ... \)  are equivalent
is a formula when it denotes a mathematical
symbol, and it should be typed as one.

\section{Displayed Text}

Text is displayed by indenting it from the left
margin.  Quotations are commonly displayed.  There
are short quotations
\begin{quote}
   This is a short quotation.  It consists of a
   single paragraph of text.  See how it is formatted.
\end{quote}
and longer ones.
\begin{quotation}
   This is a longer quotation.  It consists of two
   paragraphs of text, neither of which are
   particularly interesting.

   This is the second paragraph of the quotation.  It
   is just as dull as the first paragraph.
\end{quotation}
Another frequently-displayed structure is a list.
The following is an example of an \emph{itemized}
list.
\begin{itemize}
   \item This is the first item of an itemized list.
         Each item in the list is marked with a ``tick''.
         You don't have to worry about what kind of tick
         mark is used.

   \item This is the second item of the list.  It
         contains another list nested inside it.  The inner
         list is an \emph{enumerated} list.
         \begin{enumerate}
            \item This is the first item of an enumerated
                  list that is nested within the itemized list.

            \item This is the second item of the inner list.
                  \LaTeX\ allows you to nest lists deeper than
                  you really should.
         \end{enumerate}
         This is the rest of the second item of the outer
         list.  It is no more interesting than any other
         part of the item.
   \item This is the third item of the list.
\end{itemize}
You can even display poetry.
\begin{verse}
   There is an environment
    for verse \\             % The \\ command separates lines
   Whose features some poets % within a stanza.
   will curse.

                             % One or more blank lines separate stanzas.

   For instead of making\\
   Them do \emph{all} line breaking, \\
   It allows them to put too many words on a line when they'd rather be
   forced to be terse.
\end{verse}

Mathematical formulas may also be displayed.  A
displayed formula
is
one-line long; multiline
formulas require special formatting instructions.
   \[  \ip{\Gamma}{\psi'} = x'' + y^{2} + z_{i}^{n}\]
Don't start a paragraph with a displayed equation,
nor make one a paragraph by itself.

\end{document}               % End of document.

\section{Bench Huge Section 259}
% synthetic bench section 259
% This is a sample LaTeX input file.  (Version of 12 August 2004.)
%
% A '%' character causes TeX to ignore all remaining text on the line,
% and is used for comments like this one.

\documentclass{article}      % Specifies the document class

                             % The preamble begins here.
\title{An Example Document}  % Declares the document's title.
\author{Leslie Lamport}      % Declares the author's name.
\date{January 21, 1994}      % Deleting this command produces today's date.

\newcommand{\ip}[2]{(#1, #2)}
                             % Defines \ip{arg1}{arg2} to mean
                             % (arg1, arg2).

%\newcommand{\ip}[2]{\langle #1 | #2\rangle}
                             % This is an alternative definition of
                             % \ip that is commented out.

\begin{document}             % End of preamble and beginning of text.

\maketitle                   % Produces the title.

This is an example input file.  Comparing it with
the output it generates can show you how to
produce a simple document of your own.

\section{Ordinary Text}      % Produces section heading.  Lower-level
                             % sections are begun with similar
                             % \subsection and \subsubsection commands.

The ends  of words and sentences are marked
  by   spaces. It  doesn't matter how many
spaces    you type; one is as good as 100.  The
end of   a line counts as a space.

One   or more   blank lines denote the  end
of  a paragraph.

Since any number of consecutive spaces are treated
like a single one, the formatting of the input
file makes no difference to
      \LaTeX,                % The \LaTeX command generates the LaTeX logo.
but it makes a difference to you.  When you use
\LaTeX, making your input file as easy to read
as possible will be a great help as you write
your document and when you change it.  This sample
file shows how you can add comments to your own input
file.

Because printing is different from typewriting,
there are a number of things that you have to do
differently when preparing an input file than if
you were just typing the document directly.
Quotation marks like
       ``this''
have to be handled specially, as do quotes within
quotes:
       ``\,`this'            % \, separates the double and single quote.
        is what I just
        wrote, not  `that'\,''.

Dashes come in three sizes: an
       intra-word
dash, a medium dash for number ranges like
       1--2,
and a punctuation
       dash---like
this.

A sentence-ending space should be larger than the
space between words within a sentence.  You
sometimes have to type special commands in
conjunction with punctuation characters to get
this right, as in the following sentence.
       Gnats, gnus, etc.\ all  % `\ ' makes an inter-word space.
       begin with G\@.         % \@ marks end-of-sentence punctuation.
You should check the spaces after periods when
reading your output to make sure you haven't
forgotten any special cases.  Generating an
ellipsis
       \ldots\               % `\ ' is needed after `\ldots' because TeX
                             % ignores spaces after command names like \ldots
                             % made from \ + letters.
                             %
                             % Note how a `%' character causes TeX to ignore
                             % the end of the input line, so these blank lines
                             % do not start a new paragraph.
                             %
with the right spacing around the periods requires
a special command.

\LaTeX\ interprets some common characters as
commands, so you must type special commands to
generate them.  These characters include the
following:
       \$ \& \% \# \{ and \}.

In printing, text is usually emphasized with an
       \emph{italic}
type style.

\begin{em}
   A long segment of text can also be emphasized
   in this way.  Text within such a segment can be
   given \emph{additional} emphasis.
\end{em}

It is sometimes necessary to prevent \LaTeX\ from
breaking a line where it might otherwise do so.
This may be at a space, as between the ``Mr.''\ and
``Jones'' in
       ``Mr.~Jones'',        % ~ produces an unbreakable interword space.
or within a word---especially when the word is a
symbol like
       \mbox{\emph{itemnum}}
that makes little sense when hyphenated across
lines.

Footnotes\footnote{This is an example of a footnote.}
pose no problem.

\LaTeX\ is good at typesetting mathematical formulas
like
       \( x-3y + z = 7 \)
or
       \( a_{1} > x^{2n} + y^{2n} > x' \)
or
       \( \ip{A}{B} = \sum_{i} a_{i} b_{i} \).
The spaces you type in a formula are
ignored.  Remember that a letter like
       $x$                   % $ ... $  and  \( ... \)  are equivalent
is a formula when it denotes a mathematical
symbol, and it should be typed as one.

\section{Displayed Text}

Text is displayed by indenting it from the left
margin.  Quotations are commonly displayed.  There
are short quotations
\begin{quote}
   This is a short quotation.  It consists of a
   single paragraph of text.  See how it is formatted.
\end{quote}
and longer ones.
\begin{quotation}
   This is a longer quotation.  It consists of two
   paragraphs of text, neither of which are
   particularly interesting.

   This is the second paragraph of the quotation.  It
   is just as dull as the first paragraph.
\end{quotation}
Another frequently-displayed structure is a list.
The following is an example of an \emph{itemized}
list.
\begin{itemize}
   \item This is the first item of an itemized list.
         Each item in the list is marked with a ``tick''.
         You don't have to worry about what kind of tick
         mark is used.

   \item This is the second item of the list.  It
         contains another list nested inside it.  The inner
         list is an \emph{enumerated} list.
         \begin{enumerate}
            \item This is the first item of an enumerated
                  list that is nested within the itemized list.

            \item This is the second item of the inner list.
                  \LaTeX\ allows you to nest lists deeper than
                  you really should.
         \end{enumerate}
         This is the rest of the second item of the outer
         list.  It is no more interesting than any other
         part of the item.
   \item This is the third item of the list.
\end{itemize}
You can even display poetry.
\begin{verse}
   There is an environment
    for verse \\             % The \\ command separates lines
   Whose features some poets % within a stanza.
   will curse.

                             % One or more blank lines separate stanzas.

   For instead of making\\
   Them do \emph{all} line breaking, \\
   It allows them to put too many words on a line when they'd rather be
   forced to be terse.
\end{verse}

Mathematical formulas may also be displayed.  A
displayed formula
is
one-line long; multiline
formulas require special formatting instructions.
   \[  \ip{\Gamma}{\psi'} = x'' + y^{2} + z_{i}^{n}\]
Don't start a paragraph with a displayed equation,
nor make one a paragraph by itself.

\end{document}               % End of document.

\section{Bench Huge Section 260}
% synthetic bench section 260
% This is a sample LaTeX input file.  (Version of 12 August 2004.)
%
% A '%' character causes TeX to ignore all remaining text on the line,
% and is used for comments like this one.

\documentclass{article}      % Specifies the document class

                             % The preamble begins here.
\title{An Example Document}  % Declares the document's title.
\author{Leslie Lamport}      % Declares the author's name.
\date{January 21, 1994}      % Deleting this command produces today's date.

\newcommand{\ip}[2]{(#1, #2)}
                             % Defines \ip{arg1}{arg2} to mean
                             % (arg1, arg2).

%\newcommand{\ip}[2]{\langle #1 | #2\rangle}
                             % This is an alternative definition of
                             % \ip that is commented out.

\begin{document}             % End of preamble and beginning of text.

\maketitle                   % Produces the title.

This is an example input file.  Comparing it with
the output it generates can show you how to
produce a simple document of your own.

\section{Ordinary Text}      % Produces section heading.  Lower-level
                             % sections are begun with similar
                             % \subsection and \subsubsection commands.

The ends  of words and sentences are marked
  by   spaces. It  doesn't matter how many
spaces    you type; one is as good as 100.  The
end of   a line counts as a space.

One   or more   blank lines denote the  end
of  a paragraph.

Since any number of consecutive spaces are treated
like a single one, the formatting of the input
file makes no difference to
      \LaTeX,                % The \LaTeX command generates the LaTeX logo.
but it makes a difference to you.  When you use
\LaTeX, making your input file as easy to read
as possible will be a great help as you write
your document and when you change it.  This sample
file shows how you can add comments to your own input
file.

Because printing is different from typewriting,
there are a number of things that you have to do
differently when preparing an input file than if
you were just typing the document directly.
Quotation marks like
       ``this''
have to be handled specially, as do quotes within
quotes:
       ``\,`this'            % \, separates the double and single quote.
        is what I just
        wrote, not  `that'\,''.

Dashes come in three sizes: an
       intra-word
dash, a medium dash for number ranges like
       1--2,
and a punctuation
       dash---like
this.

A sentence-ending space should be larger than the
space between words within a sentence.  You
sometimes have to type special commands in
conjunction with punctuation characters to get
this right, as in the following sentence.
       Gnats, gnus, etc.\ all  % `\ ' makes an inter-word space.
       begin with G\@.         % \@ marks end-of-sentence punctuation.
You should check the spaces after periods when
reading your output to make sure you haven't
forgotten any special cases.  Generating an
ellipsis
       \ldots\               % `\ ' is needed after `\ldots' because TeX
                             % ignores spaces after command names like \ldots
                             % made from \ + letters.
                             %
                             % Note how a `%' character causes TeX to ignore
                             % the end of the input line, so these blank lines
                             % do not start a new paragraph.
                             %
with the right spacing around the periods requires
a special command.

\LaTeX\ interprets some common characters as
commands, so you must type special commands to
generate them.  These characters include the
following:
       \$ \& \% \# \{ and \}.

In printing, text is usually emphasized with an
       \emph{italic}
type style.

\begin{em}
   A long segment of text can also be emphasized
   in this way.  Text within such a segment can be
   given \emph{additional} emphasis.
\end{em}

It is sometimes necessary to prevent \LaTeX\ from
breaking a line where it might otherwise do so.
This may be at a space, as between the ``Mr.''\ and
``Jones'' in
       ``Mr.~Jones'',        % ~ produces an unbreakable interword space.
or within a word---especially when the word is a
symbol like
       \mbox{\emph{itemnum}}
that makes little sense when hyphenated across
lines.

Footnotes\footnote{This is an example of a footnote.}
pose no problem.

\LaTeX\ is good at typesetting mathematical formulas
like
       \( x-3y + z = 7 \)
or
       \( a_{1} > x^{2n} + y^{2n} > x' \)
or
       \( \ip{A}{B} = \sum_{i} a_{i} b_{i} \).
The spaces you type in a formula are
ignored.  Remember that a letter like
       $x$                   % $ ... $  and  \( ... \)  are equivalent
is a formula when it denotes a mathematical
symbol, and it should be typed as one.

\section{Displayed Text}

Text is displayed by indenting it from the left
margin.  Quotations are commonly displayed.  There
are short quotations
\begin{quote}
   This is a short quotation.  It consists of a
   single paragraph of text.  See how it is formatted.
\end{quote}
and longer ones.
\begin{quotation}
   This is a longer quotation.  It consists of two
   paragraphs of text, neither of which are
   particularly interesting.

   This is the second paragraph of the quotation.  It
   is just as dull as the first paragraph.
\end{quotation}
Another frequently-displayed structure is a list.
The following is an example of an \emph{itemized}
list.
\begin{itemize}
   \item This is the first item of an itemized list.
         Each item in the list is marked with a ``tick''.
         You don't have to worry about what kind of tick
         mark is used.

   \item This is the second item of the list.  It
         contains another list nested inside it.  The inner
         list is an \emph{enumerated} list.
         \begin{enumerate}
            \item This is the first item of an enumerated
                  list that is nested within the itemized list.

            \item This is the second item of the inner list.
                  \LaTeX\ allows you to nest lists deeper than
                  you really should.
         \end{enumerate}
         This is the rest of the second item of the outer
         list.  It is no more interesting than any other
         part of the item.
   \item This is the third item of the list.
\end{itemize}
You can even display poetry.
\begin{verse}
   There is an environment
    for verse \\             % The \\ command separates lines
   Whose features some poets % within a stanza.
   will curse.

                             % One or more blank lines separate stanzas.

   For instead of making\\
   Them do \emph{all} line breaking, \\
   It allows them to put too many words on a line when they'd rather be
   forced to be terse.
\end{verse}

Mathematical formulas may also be displayed.  A
displayed formula
is
one-line long; multiline
formulas require special formatting instructions.
   \[  \ip{\Gamma}{\psi'} = x'' + y^{2} + z_{i}^{n}\]
Don't start a paragraph with a displayed equation,
nor make one a paragraph by itself.

\end{document}               % End of document.

\section{Bench Huge Section 261}
% synthetic bench section 261
% This is a sample LaTeX input file.  (Version of 12 August 2004.)
%
% A '%' character causes TeX to ignore all remaining text on the line,
% and is used for comments like this one.

\documentclass{article}      % Specifies the document class

                             % The preamble begins here.
\title{An Example Document}  % Declares the document's title.
\author{Leslie Lamport}      % Declares the author's name.
\date{January 21, 1994}      % Deleting this command produces today's date.

\newcommand{\ip}[2]{(#1, #2)}
                             % Defines \ip{arg1}{arg2} to mean
                             % (arg1, arg2).

%\newcommand{\ip}[2]{\langle #1 | #2\rangle}
                             % This is an alternative definition of
                             % \ip that is commented out.

\begin{document}             % End of preamble and beginning of text.

\maketitle                   % Produces the title.

This is an example input file.  Comparing it with
the output it generates can show you how to
produce a simple document of your own.

\section{Ordinary Text}      % Produces section heading.  Lower-level
                             % sections are begun with similar
                             % \subsection and \subsubsection commands.

The ends  of words and sentences are marked
  by   spaces. It  doesn't matter how many
spaces    you type; one is as good as 100.  The
end of   a line counts as a space.

One   or more   blank lines denote the  end
of  a paragraph.

Since any number of consecutive spaces are treated
like a single one, the formatting of the input
file makes no difference to
      \LaTeX,                % The \LaTeX command generates the LaTeX logo.
but it makes a difference to you.  When you use
\LaTeX, making your input file as easy to read
as possible will be a great help as you write
your document and when you change it.  This sample
file shows how you can add comments to your own input
file.

Because printing is different from typewriting,
there are a number of things that you have to do
differently when preparing an input file than if
you were just typing the document directly.
Quotation marks like
       ``this''
have to be handled specially, as do quotes within
quotes:
       ``\,`this'            % \, separates the double and single quote.
        is what I just
        wrote, not  `that'\,''.

Dashes come in three sizes: an
       intra-word
dash, a medium dash for number ranges like
       1--2,
and a punctuation
       dash---like
this.

A sentence-ending space should be larger than the
space between words within a sentence.  You
sometimes have to type special commands in
conjunction with punctuation characters to get
this right, as in the following sentence.
       Gnats, gnus, etc.\ all  % `\ ' makes an inter-word space.
       begin with G\@.         % \@ marks end-of-sentence punctuation.
You should check the spaces after periods when
reading your output to make sure you haven't
forgotten any special cases.  Generating an
ellipsis
       \ldots\               % `\ ' is needed after `\ldots' because TeX
                             % ignores spaces after command names like \ldots
                             % made from \ + letters.
                             %
                             % Note how a `%' character causes TeX to ignore
                             % the end of the input line, so these blank lines
                             % do not start a new paragraph.
                             %
with the right spacing around the periods requires
a special command.

\LaTeX\ interprets some common characters as
commands, so you must type special commands to
generate them.  These characters include the
following:
       \$ \& \% \# \{ and \}.

In printing, text is usually emphasized with an
       \emph{italic}
type style.

\begin{em}
   A long segment of text can also be emphasized
   in this way.  Text within such a segment can be
   given \emph{additional} emphasis.
\end{em}

It is sometimes necessary to prevent \LaTeX\ from
breaking a line where it might otherwise do so.
This may be at a space, as between the ``Mr.''\ and
``Jones'' in
       ``Mr.~Jones'',        % ~ produces an unbreakable interword space.
or within a word---especially when the word is a
symbol like
       \mbox{\emph{itemnum}}
that makes little sense when hyphenated across
lines.

Footnotes\footnote{This is an example of a footnote.}
pose no problem.

\LaTeX\ is good at typesetting mathematical formulas
like
       \( x-3y + z = 7 \)
or
       \( a_{1} > x^{2n} + y^{2n} > x' \)
or
       \( \ip{A}{B} = \sum_{i} a_{i} b_{i} \).
The spaces you type in a formula are
ignored.  Remember that a letter like
       $x$                   % $ ... $  and  \( ... \)  are equivalent
is a formula when it denotes a mathematical
symbol, and it should be typed as one.

\section{Displayed Text}

Text is displayed by indenting it from the left
margin.  Quotations are commonly displayed.  There
are short quotations
\begin{quote}
   This is a short quotation.  It consists of a
   single paragraph of text.  See how it is formatted.
\end{quote}
and longer ones.
\begin{quotation}
   This is a longer quotation.  It consists of two
   paragraphs of text, neither of which are
   particularly interesting.

   This is the second paragraph of the quotation.  It
   is just as dull as the first paragraph.
\end{quotation}
Another frequently-displayed structure is a list.
The following is an example of an \emph{itemized}
list.
\begin{itemize}
   \item This is the first item of an itemized list.
         Each item in the list is marked with a ``tick''.
         You don't have to worry about what kind of tick
         mark is used.

   \item This is the second item of the list.  It
         contains another list nested inside it.  The inner
         list is an \emph{enumerated} list.
         \begin{enumerate}
            \item This is the first item of an enumerated
                  list that is nested within the itemized list.

            \item This is the second item of the inner list.
                  \LaTeX\ allows you to nest lists deeper than
                  you really should.
         \end{enumerate}
         This is the rest of the second item of the outer
         list.  It is no more interesting than any other
         part of the item.
   \item This is the third item of the list.
\end{itemize}
You can even display poetry.
\begin{verse}
   There is an environment
    for verse \\             % The \\ command separates lines
   Whose features some poets % within a stanza.
   will curse.

                             % One or more blank lines separate stanzas.

   For instead of making\\
   Them do \emph{all} line breaking, \\
   It allows them to put too many words on a line when they'd rather be
   forced to be terse.
\end{verse}

Mathematical formulas may also be displayed.  A
displayed formula
is
one-line long; multiline
formulas require special formatting instructions.
   \[  \ip{\Gamma}{\psi'} = x'' + y^{2} + z_{i}^{n}\]
Don't start a paragraph with a displayed equation,
nor make one a paragraph by itself.

\end{document}               % End of document.

\section{Bench Huge Section 262}
% synthetic bench section 262
% This is a sample LaTeX input file.  (Version of 12 August 2004.)
%
% A '%' character causes TeX to ignore all remaining text on the line,
% and is used for comments like this one.

\documentclass{article}      % Specifies the document class

                             % The preamble begins here.
\title{An Example Document}  % Declares the document's title.
\author{Leslie Lamport}      % Declares the author's name.
\date{January 21, 1994}      % Deleting this command produces today's date.

\newcommand{\ip}[2]{(#1, #2)}
                             % Defines \ip{arg1}{arg2} to mean
                             % (arg1, arg2).

%\newcommand{\ip}[2]{\langle #1 | #2\rangle}
                             % This is an alternative definition of
                             % \ip that is commented out.

\begin{document}             % End of preamble and beginning of text.

\maketitle                   % Produces the title.

This is an example input file.  Comparing it with
the output it generates can show you how to
produce a simple document of your own.

\section{Ordinary Text}      % Produces section heading.  Lower-level
                             % sections are begun with similar
                             % \subsection and \subsubsection commands.

The ends  of words and sentences are marked
  by   spaces. It  doesn't matter how many
spaces    you type; one is as good as 100.  The
end of   a line counts as a space.

One   or more   blank lines denote the  end
of  a paragraph.

Since any number of consecutive spaces are treated
like a single one, the formatting of the input
file makes no difference to
      \LaTeX,                % The \LaTeX command generates the LaTeX logo.
but it makes a difference to you.  When you use
\LaTeX, making your input file as easy to read
as possible will be a great help as you write
your document and when you change it.  This sample
file shows how you can add comments to your own input
file.

Because printing is different from typewriting,
there are a number of things that you have to do
differently when preparing an input file than if
you were just typing the document directly.
Quotation marks like
       ``this''
have to be handled specially, as do quotes within
quotes:
       ``\,`this'            % \, separates the double and single quote.
        is what I just
        wrote, not  `that'\,''.

Dashes come in three sizes: an
       intra-word
dash, a medium dash for number ranges like
       1--2,
and a punctuation
       dash---like
this.

A sentence-ending space should be larger than the
space between words within a sentence.  You
sometimes have to type special commands in
conjunction with punctuation characters to get
this right, as in the following sentence.
       Gnats, gnus, etc.\ all  % `\ ' makes an inter-word space.
       begin with G\@.         % \@ marks end-of-sentence punctuation.
You should check the spaces after periods when
reading your output to make sure you haven't
forgotten any special cases.  Generating an
ellipsis
       \ldots\               % `\ ' is needed after `\ldots' because TeX
                             % ignores spaces after command names like \ldots
                             % made from \ + letters.
                             %
                             % Note how a `%' character causes TeX to ignore
                             % the end of the input line, so these blank lines
                             % do not start a new paragraph.
                             %
with the right spacing around the periods requires
a special command.

\LaTeX\ interprets some common characters as
commands, so you must type special commands to
generate them.  These characters include the
following:
       \$ \& \% \# \{ and \}.

In printing, text is usually emphasized with an
       \emph{italic}
type style.

\begin{em}
   A long segment of text can also be emphasized
   in this way.  Text within such a segment can be
   given \emph{additional} emphasis.
\end{em}

It is sometimes necessary to prevent \LaTeX\ from
breaking a line where it might otherwise do so.
This may be at a space, as between the ``Mr.''\ and
``Jones'' in
       ``Mr.~Jones'',        % ~ produces an unbreakable interword space.
or within a word---especially when the word is a
symbol like
       \mbox{\emph{itemnum}}
that makes little sense when hyphenated across
lines.

Footnotes\footnote{This is an example of a footnote.}
pose no problem.

\LaTeX\ is good at typesetting mathematical formulas
like
       \( x-3y + z = 7 \)
or
       \( a_{1} > x^{2n} + y^{2n} > x' \)
or
       \( \ip{A}{B} = \sum_{i} a_{i} b_{i} \).
The spaces you type in a formula are
ignored.  Remember that a letter like
       $x$                   % $ ... $  and  \( ... \)  are equivalent
is a formula when it denotes a mathematical
symbol, and it should be typed as one.

\section{Displayed Text}

Text is displayed by indenting it from the left
margin.  Quotations are commonly displayed.  There
are short quotations
\begin{quote}
   This is a short quotation.  It consists of a
   single paragraph of text.  See how it is formatted.
\end{quote}
and longer ones.
\begin{quotation}
   This is a longer quotation.  It consists of two
   paragraphs of text, neither of which are
   particularly interesting.

   This is the second paragraph of the quotation.  It
   is just as dull as the first paragraph.
\end{quotation}
Another frequently-displayed structure is a list.
The following is an example of an \emph{itemized}
list.
\begin{itemize}
   \item This is the first item of an itemized list.
         Each item in the list is marked with a ``tick''.
         You don't have to worry about what kind of tick
         mark is used.

   \item This is the second item of the list.  It
         contains another list nested inside it.  The inner
         list is an \emph{enumerated} list.
         \begin{enumerate}
            \item This is the first item of an enumerated
                  list that is nested within the itemized list.

            \item This is the second item of the inner list.
                  \LaTeX\ allows you to nest lists deeper than
                  you really should.
         \end{enumerate}
         This is the rest of the second item of the outer
         list.  It is no more interesting than any other
         part of the item.
   \item This is the third item of the list.
\end{itemize}
You can even display poetry.
\begin{verse}
   There is an environment
    for verse \\             % The \\ command separates lines
   Whose features some poets % within a stanza.
   will curse.

                             % One or more blank lines separate stanzas.

   For instead of making\\
   Them do \emph{all} line breaking, \\
   It allows them to put too many words on a line when they'd rather be
   forced to be terse.
\end{verse}

Mathematical formulas may also be displayed.  A
displayed formula
is
one-line long; multiline
formulas require special formatting instructions.
   \[  \ip{\Gamma}{\psi'} = x'' + y^{2} + z_{i}^{n}\]
Don't start a paragraph with a displayed equation,
nor make one a paragraph by itself.

\end{document}               % End of document.

\section{Bench Huge Section 263}
% synthetic bench section 263
% This is a sample LaTeX input file.  (Version of 12 August 2004.)
%
% A '%' character causes TeX to ignore all remaining text on the line,
% and is used for comments like this one.

\documentclass{article}      % Specifies the document class

                             % The preamble begins here.
\title{An Example Document}  % Declares the document's title.
\author{Leslie Lamport}      % Declares the author's name.
\date{January 21, 1994}      % Deleting this command produces today's date.

\newcommand{\ip}[2]{(#1, #2)}
                             % Defines \ip{arg1}{arg2} to mean
                             % (arg1, arg2).

%\newcommand{\ip}[2]{\langle #1 | #2\rangle}
                             % This is an alternative definition of
                             % \ip that is commented out.

\begin{document}             % End of preamble and beginning of text.

\maketitle                   % Produces the title.

This is an example input file.  Comparing it with
the output it generates can show you how to
produce a simple document of your own.

\section{Ordinary Text}      % Produces section heading.  Lower-level
                             % sections are begun with similar
                             % \subsection and \subsubsection commands.

The ends  of words and sentences are marked
  by   spaces. It  doesn't matter how many
spaces    you type; one is as good as 100.  The
end of   a line counts as a space.

One   or more   blank lines denote the  end
of  a paragraph.

Since any number of consecutive spaces are treated
like a single one, the formatting of the input
file makes no difference to
      \LaTeX,                % The \LaTeX command generates the LaTeX logo.
but it makes a difference to you.  When you use
\LaTeX, making your input file as easy to read
as possible will be a great help as you write
your document and when you change it.  This sample
file shows how you can add comments to your own input
file.

Because printing is different from typewriting,
there are a number of things that you have to do
differently when preparing an input file than if
you were just typing the document directly.
Quotation marks like
       ``this''
have to be handled specially, as do quotes within
quotes:
       ``\,`this'            % \, separates the double and single quote.
        is what I just
        wrote, not  `that'\,''.

Dashes come in three sizes: an
       intra-word
dash, a medium dash for number ranges like
       1--2,
and a punctuation
       dash---like
this.

A sentence-ending space should be larger than the
space between words within a sentence.  You
sometimes have to type special commands in
conjunction with punctuation characters to get
this right, as in the following sentence.
       Gnats, gnus, etc.\ all  % `\ ' makes an inter-word space.
       begin with G\@.         % \@ marks end-of-sentence punctuation.
You should check the spaces after periods when
reading your output to make sure you haven't
forgotten any special cases.  Generating an
ellipsis
       \ldots\               % `\ ' is needed after `\ldots' because TeX
                             % ignores spaces after command names like \ldots
                             % made from \ + letters.
                             %
                             % Note how a `%' character causes TeX to ignore
                             % the end of the input line, so these blank lines
                             % do not start a new paragraph.
                             %
with the right spacing around the periods requires
a special command.

\LaTeX\ interprets some common characters as
commands, so you must type special commands to
generate them.  These characters include the
following:
       \$ \& \% \# \{ and \}.

In printing, text is usually emphasized with an
       \emph{italic}
type style.

\begin{em}
   A long segment of text can also be emphasized
   in this way.  Text within such a segment can be
   given \emph{additional} emphasis.
\end{em}

It is sometimes necessary to prevent \LaTeX\ from
breaking a line where it might otherwise do so.
This may be at a space, as between the ``Mr.''\ and
``Jones'' in
       ``Mr.~Jones'',        % ~ produces an unbreakable interword space.
or within a word---especially when the word is a
symbol like
       \mbox{\emph{itemnum}}
that makes little sense when hyphenated across
lines.

Footnotes\footnote{This is an example of a footnote.}
pose no problem.

\LaTeX\ is good at typesetting mathematical formulas
like
       \( x-3y + z = 7 \)
or
       \( a_{1} > x^{2n} + y^{2n} > x' \)
or
       \( \ip{A}{B} = \sum_{i} a_{i} b_{i} \).
The spaces you type in a formula are
ignored.  Remember that a letter like
       $x$                   % $ ... $  and  \( ... \)  are equivalent
is a formula when it denotes a mathematical
symbol, and it should be typed as one.

\section{Displayed Text}

Text is displayed by indenting it from the left
margin.  Quotations are commonly displayed.  There
are short quotations
\begin{quote}
   This is a short quotation.  It consists of a
   single paragraph of text.  See how it is formatted.
\end{quote}
and longer ones.
\begin{quotation}
   This is a longer quotation.  It consists of two
   paragraphs of text, neither of which are
   particularly interesting.

   This is the second paragraph of the quotation.  It
   is just as dull as the first paragraph.
\end{quotation}
Another frequently-displayed structure is a list.
The following is an example of an \emph{itemized}
list.
\begin{itemize}
   \item This is the first item of an itemized list.
         Each item in the list is marked with a ``tick''.
         You don't have to worry about what kind of tick
         mark is used.

   \item This is the second item of the list.  It
         contains another list nested inside it.  The inner
         list is an \emph{enumerated} list.
         \begin{enumerate}
            \item This is the first item of an enumerated
                  list that is nested within the itemized list.

            \item This is the second item of the inner list.
                  \LaTeX\ allows you to nest lists deeper than
                  you really should.
         \end{enumerate}
         This is the rest of the second item of the outer
         list.  It is no more interesting than any other
         part of the item.
   \item This is the third item of the list.
\end{itemize}
You can even display poetry.
\begin{verse}
   There is an environment
    for verse \\             % The \\ command separates lines
   Whose features some poets % within a stanza.
   will curse.

                             % One or more blank lines separate stanzas.

   For instead of making\\
   Them do \emph{all} line breaking, \\
   It allows them to put too many words on a line when they'd rather be
   forced to be terse.
\end{verse}

Mathematical formulas may also be displayed.  A
displayed formula
is
one-line long; multiline
formulas require special formatting instructions.
   \[  \ip{\Gamma}{\psi'} = x'' + y^{2} + z_{i}^{n}\]
Don't start a paragraph with a displayed equation,
nor make one a paragraph by itself.

\end{document}               % End of document.

\section{Bench Huge Section 264}
% synthetic bench section 264
% This is a sample LaTeX input file.  (Version of 12 August 2004.)
%
% A '%' character causes TeX to ignore all remaining text on the line,
% and is used for comments like this one.

\documentclass{article}      % Specifies the document class

                             % The preamble begins here.
\title{An Example Document}  % Declares the document's title.
\author{Leslie Lamport}      % Declares the author's name.
\date{January 21, 1994}      % Deleting this command produces today's date.

\newcommand{\ip}[2]{(#1, #2)}
                             % Defines \ip{arg1}{arg2} to mean
                             % (arg1, arg2).

%\newcommand{\ip}[2]{\langle #1 | #2\rangle}
                             % This is an alternative definition of
                             % \ip that is commented out.

\begin{document}             % End of preamble and beginning of text.

\maketitle                   % Produces the title.

This is an example input file.  Comparing it with
the output it generates can show you how to
produce a simple document of your own.

\section{Ordinary Text}      % Produces section heading.  Lower-level
                             % sections are begun with similar
                             % \subsection and \subsubsection commands.

The ends  of words and sentences are marked
  by   spaces. It  doesn't matter how many
spaces    you type; one is as good as 100.  The
end of   a line counts as a space.

One   or more   blank lines denote the  end
of  a paragraph.

Since any number of consecutive spaces are treated
like a single one, the formatting of the input
file makes no difference to
      \LaTeX,                % The \LaTeX command generates the LaTeX logo.
but it makes a difference to you.  When you use
\LaTeX, making your input file as easy to read
as possible will be a great help as you write
your document and when you change it.  This sample
file shows how you can add comments to your own input
file.

Because printing is different from typewriting,
there are a number of things that you have to do
differently when preparing an input file than if
you were just typing the document directly.
Quotation marks like
       ``this''
have to be handled specially, as do quotes within
quotes:
       ``\,`this'            % \, separates the double and single quote.
        is what I just
        wrote, not  `that'\,''.

Dashes come in three sizes: an
       intra-word
dash, a medium dash for number ranges like
       1--2,
and a punctuation
       dash---like
this.

A sentence-ending space should be larger than the
space between words within a sentence.  You
sometimes have to type special commands in
conjunction with punctuation characters to get
this right, as in the following sentence.
       Gnats, gnus, etc.\ all  % `\ ' makes an inter-word space.
       begin with G\@.         % \@ marks end-of-sentence punctuation.
You should check the spaces after periods when
reading your output to make sure you haven't
forgotten any special cases.  Generating an
ellipsis
       \ldots\               % `\ ' is needed after `\ldots' because TeX
                             % ignores spaces after command names like \ldots
                             % made from \ + letters.
                             %
                             % Note how a `%' character causes TeX to ignore
                             % the end of the input line, so these blank lines
                             % do not start a new paragraph.
                             %
with the right spacing around the periods requires
a special command.

\LaTeX\ interprets some common characters as
commands, so you must type special commands to
generate them.  These characters include the
following:
       \$ \& \% \# \{ and \}.

In printing, text is usually emphasized with an
       \emph{italic}
type style.

\begin{em}
   A long segment of text can also be emphasized
   in this way.  Text within such a segment can be
   given \emph{additional} emphasis.
\end{em}

It is sometimes necessary to prevent \LaTeX\ from
breaking a line where it might otherwise do so.
This may be at a space, as between the ``Mr.''\ and
``Jones'' in
       ``Mr.~Jones'',        % ~ produces an unbreakable interword space.
or within a word---especially when the word is a
symbol like
       \mbox{\emph{itemnum}}
that makes little sense when hyphenated across
lines.

Footnotes\footnote{This is an example of a footnote.}
pose no problem.

\LaTeX\ is good at typesetting mathematical formulas
like
       \( x-3y + z = 7 \)
or
       \( a_{1} > x^{2n} + y^{2n} > x' \)
or
       \( \ip{A}{B} = \sum_{i} a_{i} b_{i} \).
The spaces you type in a formula are
ignored.  Remember that a letter like
       $x$                   % $ ... $  and  \( ... \)  are equivalent
is a formula when it denotes a mathematical
symbol, and it should be typed as one.

\section{Displayed Text}

Text is displayed by indenting it from the left
margin.  Quotations are commonly displayed.  There
are short quotations
\begin{quote}
   This is a short quotation.  It consists of a
   single paragraph of text.  See how it is formatted.
\end{quote}
and longer ones.
\begin{quotation}
   This is a longer quotation.  It consists of two
   paragraphs of text, neither of which are
   particularly interesting.

   This is the second paragraph of the quotation.  It
   is just as dull as the first paragraph.
\end{quotation}
Another frequently-displayed structure is a list.
The following is an example of an \emph{itemized}
list.
\begin{itemize}
   \item This is the first item of an itemized list.
         Each item in the list is marked with a ``tick''.
         You don't have to worry about what kind of tick
         mark is used.

   \item This is the second item of the list.  It
         contains another list nested inside it.  The inner
         list is an \emph{enumerated} list.
         \begin{enumerate}
            \item This is the first item of an enumerated
                  list that is nested within the itemized list.

            \item This is the second item of the inner list.
                  \LaTeX\ allows you to nest lists deeper than
                  you really should.
         \end{enumerate}
         This is the rest of the second item of the outer
         list.  It is no more interesting than any other
         part of the item.
   \item This is the third item of the list.
\end{itemize}
You can even display poetry.
\begin{verse}
   There is an environment
    for verse \\             % The \\ command separates lines
   Whose features some poets % within a stanza.
   will curse.

                             % One or more blank lines separate stanzas.

   For instead of making\\
   Them do \emph{all} line breaking, \\
   It allows them to put too many words on a line when they'd rather be
   forced to be terse.
\end{verse}

Mathematical formulas may also be displayed.  A
displayed formula
is
one-line long; multiline
formulas require special formatting instructions.
   \[  \ip{\Gamma}{\psi'} = x'' + y^{2} + z_{i}^{n}\]
Don't start a paragraph with a displayed equation,
nor make one a paragraph by itself.

\end{document}               % End of document.

\section{Bench Huge Section 265}
% synthetic bench section 265
% This is a sample LaTeX input file.  (Version of 12 August 2004.)
%
% A '%' character causes TeX to ignore all remaining text on the line,
% and is used for comments like this one.

\documentclass{article}      % Specifies the document class

                             % The preamble begins here.
\title{An Example Document}  % Declares the document's title.
\author{Leslie Lamport}      % Declares the author's name.
\date{January 21, 1994}      % Deleting this command produces today's date.

\newcommand{\ip}[2]{(#1, #2)}
                             % Defines \ip{arg1}{arg2} to mean
                             % (arg1, arg2).

%\newcommand{\ip}[2]{\langle #1 | #2\rangle}
                             % This is an alternative definition of
                             % \ip that is commented out.

\begin{document}             % End of preamble and beginning of text.

\maketitle                   % Produces the title.

This is an example input file.  Comparing it with
the output it generates can show you how to
produce a simple document of your own.

\section{Ordinary Text}      % Produces section heading.  Lower-level
                             % sections are begun with similar
                             % \subsection and \subsubsection commands.

The ends  of words and sentences are marked
  by   spaces. It  doesn't matter how many
spaces    you type; one is as good as 100.  The
end of   a line counts as a space.

One   or more   blank lines denote the  end
of  a paragraph.

Since any number of consecutive spaces are treated
like a single one, the formatting of the input
file makes no difference to
      \LaTeX,                % The \LaTeX command generates the LaTeX logo.
but it makes a difference to you.  When you use
\LaTeX, making your input file as easy to read
as possible will be a great help as you write
your document and when you change it.  This sample
file shows how you can add comments to your own input
file.

Because printing is different from typewriting,
there are a number of things that you have to do
differently when preparing an input file than if
you were just typing the document directly.
Quotation marks like
       ``this''
have to be handled specially, as do quotes within
quotes:
       ``\,`this'            % \, separates the double and single quote.
        is what I just
        wrote, not  `that'\,''.

Dashes come in three sizes: an
       intra-word
dash, a medium dash for number ranges like
       1--2,
and a punctuation
       dash---like
this.

A sentence-ending space should be larger than the
space between words within a sentence.  You
sometimes have to type special commands in
conjunction with punctuation characters to get
this right, as in the following sentence.
       Gnats, gnus, etc.\ all  % `\ ' makes an inter-word space.
       begin with G\@.         % \@ marks end-of-sentence punctuation.
You should check the spaces after periods when
reading your output to make sure you haven't
forgotten any special cases.  Generating an
ellipsis
       \ldots\               % `\ ' is needed after `\ldots' because TeX
                             % ignores spaces after command names like \ldots
                             % made from \ + letters.
                             %
                             % Note how a `%' character causes TeX to ignore
                             % the end of the input line, so these blank lines
                             % do not start a new paragraph.
                             %
with the right spacing around the periods requires
a special command.

\LaTeX\ interprets some common characters as
commands, so you must type special commands to
generate them.  These characters include the
following:
       \$ \& \% \# \{ and \}.

In printing, text is usually emphasized with an
       \emph{italic}
type style.

\begin{em}
   A long segment of text can also be emphasized
   in this way.  Text within such a segment can be
   given \emph{additional} emphasis.
\end{em}

It is sometimes necessary to prevent \LaTeX\ from
breaking a line where it might otherwise do so.
This may be at a space, as between the ``Mr.''\ and
``Jones'' in
       ``Mr.~Jones'',        % ~ produces an unbreakable interword space.
or within a word---especially when the word is a
symbol like
       \mbox{\emph{itemnum}}
that makes little sense when hyphenated across
lines.

Footnotes\footnote{This is an example of a footnote.}
pose no problem.

\LaTeX\ is good at typesetting mathematical formulas
like
       \( x-3y + z = 7 \)
or
       \( a_{1} > x^{2n} + y^{2n} > x' \)
or
       \( \ip{A}{B} = \sum_{i} a_{i} b_{i} \).
The spaces you type in a formula are
ignored.  Remember that a letter like
       $x$                   % $ ... $  and  \( ... \)  are equivalent
is a formula when it denotes a mathematical
symbol, and it should be typed as one.

\section{Displayed Text}

Text is displayed by indenting it from the left
margin.  Quotations are commonly displayed.  There
are short quotations
\begin{quote}
   This is a short quotation.  It consists of a
   single paragraph of text.  See how it is formatted.
\end{quote}
and longer ones.
\begin{quotation}
   This is a longer quotation.  It consists of two
   paragraphs of text, neither of which are
   particularly interesting.

   This is the second paragraph of the quotation.  It
   is just as dull as the first paragraph.
\end{quotation}
Another frequently-displayed structure is a list.
The following is an example of an \emph{itemized}
list.
\begin{itemize}
   \item This is the first item of an itemized list.
         Each item in the list is marked with a ``tick''.
         You don't have to worry about what kind of tick
         mark is used.

   \item This is the second item of the list.  It
         contains another list nested inside it.  The inner
         list is an \emph{enumerated} list.
         \begin{enumerate}
            \item This is the first item of an enumerated
                  list that is nested within the itemized list.

            \item This is the second item of the inner list.
                  \LaTeX\ allows you to nest lists deeper than
                  you really should.
         \end{enumerate}
         This is the rest of the second item of the outer
         list.  It is no more interesting than any other
         part of the item.
   \item This is the third item of the list.
\end{itemize}
You can even display poetry.
\begin{verse}
   There is an environment
    for verse \\             % The \\ command separates lines
   Whose features some poets % within a stanza.
   will curse.

                             % One or more blank lines separate stanzas.

   For instead of making\\
   Them do \emph{all} line breaking, \\
   It allows them to put too many words on a line when they'd rather be
   forced to be terse.
\end{verse}

Mathematical formulas may also be displayed.  A
displayed formula
is
one-line long; multiline
formulas require special formatting instructions.
   \[  \ip{\Gamma}{\psi'} = x'' + y^{2} + z_{i}^{n}\]
Don't start a paragraph with a displayed equation,
nor make one a paragraph by itself.

\end{document}               % End of document.

\section{Bench Huge Section 266}
% synthetic bench section 266
% This is a sample LaTeX input file.  (Version of 12 August 2004.)
%
% A '%' character causes TeX to ignore all remaining text on the line,
% and is used for comments like this one.

\documentclass{article}      % Specifies the document class

                             % The preamble begins here.
\title{An Example Document}  % Declares the document's title.
\author{Leslie Lamport}      % Declares the author's name.
\date{January 21, 1994}      % Deleting this command produces today's date.

\newcommand{\ip}[2]{(#1, #2)}
                             % Defines \ip{arg1}{arg2} to mean
                             % (arg1, arg2).

%\newcommand{\ip}[2]{\langle #1 | #2\rangle}
                             % This is an alternative definition of
                             % \ip that is commented out.

\begin{document}             % End of preamble and beginning of text.

\maketitle                   % Produces the title.

This is an example input file.  Comparing it with
the output it generates can show you how to
produce a simple document of your own.

\section{Ordinary Text}      % Produces section heading.  Lower-level
                             % sections are begun with similar
                             % \subsection and \subsubsection commands.

The ends  of words and sentences are marked
  by   spaces. It  doesn't matter how many
spaces    you type; one is as good as 100.  The
end of   a line counts as a space.

One   or more   blank lines denote the  end
of  a paragraph.

Since any number of consecutive spaces are treated
like a single one, the formatting of the input
file makes no difference to
      \LaTeX,                % The \LaTeX command generates the LaTeX logo.
but it makes a difference to you.  When you use
\LaTeX, making your input file as easy to read
as possible will be a great help as you write
your document and when you change it.  This sample
file shows how you can add comments to your own input
file.

Because printing is different from typewriting,
there are a number of things that you have to do
differently when preparing an input file than if
you were just typing the document directly.
Quotation marks like
       ``this''
have to be handled specially, as do quotes within
quotes:
       ``\,`this'            % \, separates the double and single quote.
        is what I just
        wrote, not  `that'\,''.

Dashes come in three sizes: an
       intra-word
dash, a medium dash for number ranges like
       1--2,
and a punctuation
       dash---like
this.

A sentence-ending space should be larger than the
space between words within a sentence.  You
sometimes have to type special commands in
conjunction with punctuation characters to get
this right, as in the following sentence.
       Gnats, gnus, etc.\ all  % `\ ' makes an inter-word space.
       begin with G\@.         % \@ marks end-of-sentence punctuation.
You should check the spaces after periods when
reading your output to make sure you haven't
forgotten any special cases.  Generating an
ellipsis
       \ldots\               % `\ ' is needed after `\ldots' because TeX
                             % ignores spaces after command names like \ldots
                             % made from \ + letters.
                             %
                             % Note how a `%' character causes TeX to ignore
                             % the end of the input line, so these blank lines
                             % do not start a new paragraph.
                             %
with the right spacing around the periods requires
a special command.

\LaTeX\ interprets some common characters as
commands, so you must type special commands to
generate them.  These characters include the
following:
       \$ \& \% \# \{ and \}.

In printing, text is usually emphasized with an
       \emph{italic}
type style.

\begin{em}
   A long segment of text can also be emphasized
   in this way.  Text within such a segment can be
   given \emph{additional} emphasis.
\end{em}

It is sometimes necessary to prevent \LaTeX\ from
breaking a line where it might otherwise do so.
This may be at a space, as between the ``Mr.''\ and
``Jones'' in
       ``Mr.~Jones'',        % ~ produces an unbreakable interword space.
or within a word---especially when the word is a
symbol like
       \mbox{\emph{itemnum}}
that makes little sense when hyphenated across
lines.

Footnotes\footnote{This is an example of a footnote.}
pose no problem.

\LaTeX\ is good at typesetting mathematical formulas
like
       \( x-3y + z = 7 \)
or
       \( a_{1} > x^{2n} + y^{2n} > x' \)
or
       \( \ip{A}{B} = \sum_{i} a_{i} b_{i} \).
The spaces you type in a formula are
ignored.  Remember that a letter like
       $x$                   % $ ... $  and  \( ... \)  are equivalent
is a formula when it denotes a mathematical
symbol, and it should be typed as one.

\section{Displayed Text}

Text is displayed by indenting it from the left
margin.  Quotations are commonly displayed.  There
are short quotations
\begin{quote}
   This is a short quotation.  It consists of a
   single paragraph of text.  See how it is formatted.
\end{quote}
and longer ones.
\begin{quotation}
   This is a longer quotation.  It consists of two
   paragraphs of text, neither of which are
   particularly interesting.

   This is the second paragraph of the quotation.  It
   is just as dull as the first paragraph.
\end{quotation}
Another frequently-displayed structure is a list.
The following is an example of an \emph{itemized}
list.
\begin{itemize}
   \item This is the first item of an itemized list.
         Each item in the list is marked with a ``tick''.
         You don't have to worry about what kind of tick
         mark is used.

   \item This is the second item of the list.  It
         contains another list nested inside it.  The inner
         list is an \emph{enumerated} list.
         \begin{enumerate}
            \item This is the first item of an enumerated
                  list that is nested within the itemized list.

            \item This is the second item of the inner list.
                  \LaTeX\ allows you to nest lists deeper than
                  you really should.
         \end{enumerate}
         This is the rest of the second item of the outer
         list.  It is no more interesting than any other
         part of the item.
   \item This is the third item of the list.
\end{itemize}
You can even display poetry.
\begin{verse}
   There is an environment
    for verse \\             % The \\ command separates lines
   Whose features some poets % within a stanza.
   will curse.

                             % One or more blank lines separate stanzas.

   For instead of making\\
   Them do \emph{all} line breaking, \\
   It allows them to put too many words on a line when they'd rather be
   forced to be terse.
\end{verse}

Mathematical formulas may also be displayed.  A
displayed formula
is
one-line long; multiline
formulas require special formatting instructions.
   \[  \ip{\Gamma}{\psi'} = x'' + y^{2} + z_{i}^{n}\]
Don't start a paragraph with a displayed equation,
nor make one a paragraph by itself.

\end{document}               % End of document.

\section{Bench Huge Section 267}
% synthetic bench section 267
% This is a sample LaTeX input file.  (Version of 12 August 2004.)
%
% A '%' character causes TeX to ignore all remaining text on the line,
% and is used for comments like this one.

\documentclass{article}      % Specifies the document class

                             % The preamble begins here.
\title{An Example Document}  % Declares the document's title.
\author{Leslie Lamport}      % Declares the author's name.
\date{January 21, 1994}      % Deleting this command produces today's date.

\newcommand{\ip}[2]{(#1, #2)}
                             % Defines \ip{arg1}{arg2} to mean
                             % (arg1, arg2).

%\newcommand{\ip}[2]{\langle #1 | #2\rangle}
                             % This is an alternative definition of
                             % \ip that is commented out.

\begin{document}             % End of preamble and beginning of text.

\maketitle                   % Produces the title.

This is an example input file.  Comparing it with
the output it generates can show you how to
produce a simple document of your own.

\section{Ordinary Text}      % Produces section heading.  Lower-level
                             % sections are begun with similar
                             % \subsection and \subsubsection commands.

The ends  of words and sentences are marked
  by   spaces. It  doesn't matter how many
spaces    you type; one is as good as 100.  The
end of   a line counts as a space.

One   or more   blank lines denote the  end
of  a paragraph.

Since any number of consecutive spaces are treated
like a single one, the formatting of the input
file makes no difference to
      \LaTeX,                % The \LaTeX command generates the LaTeX logo.
but it makes a difference to you.  When you use
\LaTeX, making your input file as easy to read
as possible will be a great help as you write
your document and when you change it.  This sample
file shows how you can add comments to your own input
file.

Because printing is different from typewriting,
there are a number of things that you have to do
differently when preparing an input file than if
you were just typing the document directly.
Quotation marks like
       ``this''
have to be handled specially, as do quotes within
quotes:
       ``\,`this'            % \, separates the double and single quote.
        is what I just
        wrote, not  `that'\,''.

Dashes come in three sizes: an
       intra-word
dash, a medium dash for number ranges like
       1--2,
and a punctuation
       dash---like
this.

A sentence-ending space should be larger than the
space between words within a sentence.  You
sometimes have to type special commands in
conjunction with punctuation characters to get
this right, as in the following sentence.
       Gnats, gnus, etc.\ all  % `\ ' makes an inter-word space.
       begin with G\@.         % \@ marks end-of-sentence punctuation.
You should check the spaces after periods when
reading your output to make sure you haven't
forgotten any special cases.  Generating an
ellipsis
       \ldots\               % `\ ' is needed after `\ldots' because TeX
                             % ignores spaces after command names like \ldots
                             % made from \ + letters.
                             %
                             % Note how a `%' character causes TeX to ignore
                             % the end of the input line, so these blank lines
                             % do not start a new paragraph.
                             %
with the right spacing around the periods requires
a special command.

\LaTeX\ interprets some common characters as
commands, so you must type special commands to
generate them.  These characters include the
following:
       \$ \& \% \# \{ and \}.

In printing, text is usually emphasized with an
       \emph{italic}
type style.

\begin{em}
   A long segment of text can also be emphasized
   in this way.  Text within such a segment can be
   given \emph{additional} emphasis.
\end{em}

It is sometimes necessary to prevent \LaTeX\ from
breaking a line where it might otherwise do so.
This may be at a space, as between the ``Mr.''\ and
``Jones'' in
       ``Mr.~Jones'',        % ~ produces an unbreakable interword space.
or within a word---especially when the word is a
symbol like
       \mbox{\emph{itemnum}}
that makes little sense when hyphenated across
lines.

Footnotes\footnote{This is an example of a footnote.}
pose no problem.

\LaTeX\ is good at typesetting mathematical formulas
like
       \( x-3y + z = 7 \)
or
       \( a_{1} > x^{2n} + y^{2n} > x' \)
or
       \( \ip{A}{B} = \sum_{i} a_{i} b_{i} \).
The spaces you type in a formula are
ignored.  Remember that a letter like
       $x$                   % $ ... $  and  \( ... \)  are equivalent
is a formula when it denotes a mathematical
symbol, and it should be typed as one.

\section{Displayed Text}

Text is displayed by indenting it from the left
margin.  Quotations are commonly displayed.  There
are short quotations
\begin{quote}
   This is a short quotation.  It consists of a
   single paragraph of text.  See how it is formatted.
\end{quote}
and longer ones.
\begin{quotation}
   This is a longer quotation.  It consists of two
   paragraphs of text, neither of which are
   particularly interesting.

   This is the second paragraph of the quotation.  It
   is just as dull as the first paragraph.
\end{quotation}
Another frequently-displayed structure is a list.
The following is an example of an \emph{itemized}
list.
\begin{itemize}
   \item This is the first item of an itemized list.
         Each item in the list is marked with a ``tick''.
         You don't have to worry about what kind of tick
         mark is used.

   \item This is the second item of the list.  It
         contains another list nested inside it.  The inner
         list is an \emph{enumerated} list.
         \begin{enumerate}
            \item This is the first item of an enumerated
                  list that is nested within the itemized list.

            \item This is the second item of the inner list.
                  \LaTeX\ allows you to nest lists deeper than
                  you really should.
         \end{enumerate}
         This is the rest of the second item of the outer
         list.  It is no more interesting than any other
         part of the item.
   \item This is the third item of the list.
\end{itemize}
You can even display poetry.
\begin{verse}
   There is an environment
    for verse \\             % The \\ command separates lines
   Whose features some poets % within a stanza.
   will curse.

                             % One or more blank lines separate stanzas.

   For instead of making\\
   Them do \emph{all} line breaking, \\
   It allows them to put too many words on a line when they'd rather be
   forced to be terse.
\end{verse}

Mathematical formulas may also be displayed.  A
displayed formula
is
one-line long; multiline
formulas require special formatting instructions.
   \[  \ip{\Gamma}{\psi'} = x'' + y^{2} + z_{i}^{n}\]
Don't start a paragraph with a displayed equation,
nor make one a paragraph by itself.

\end{document}               % End of document.

\section{Bench Huge Section 268}
% synthetic bench section 268
% This is a sample LaTeX input file.  (Version of 12 August 2004.)
%
% A '%' character causes TeX to ignore all remaining text on the line,
% and is used for comments like this one.

\documentclass{article}      % Specifies the document class

                             % The preamble begins here.
\title{An Example Document}  % Declares the document's title.
\author{Leslie Lamport}      % Declares the author's name.
\date{January 21, 1994}      % Deleting this command produces today's date.

\newcommand{\ip}[2]{(#1, #2)}
                             % Defines \ip{arg1}{arg2} to mean
                             % (arg1, arg2).

%\newcommand{\ip}[2]{\langle #1 | #2\rangle}
                             % This is an alternative definition of
                             % \ip that is commented out.

\begin{document}             % End of preamble and beginning of text.

\maketitle                   % Produces the title.

This is an example input file.  Comparing it with
the output it generates can show you how to
produce a simple document of your own.

\section{Ordinary Text}      % Produces section heading.  Lower-level
                             % sections are begun with similar
                             % \subsection and \subsubsection commands.

The ends  of words and sentences are marked
  by   spaces. It  doesn't matter how many
spaces    you type; one is as good as 100.  The
end of   a line counts as a space.

One   or more   blank lines denote the  end
of  a paragraph.

Since any number of consecutive spaces are treated
like a single one, the formatting of the input
file makes no difference to
      \LaTeX,                % The \LaTeX command generates the LaTeX logo.
but it makes a difference to you.  When you use
\LaTeX, making your input file as easy to read
as possible will be a great help as you write
your document and when you change it.  This sample
file shows how you can add comments to your own input
file.

Because printing is different from typewriting,
there are a number of things that you have to do
differently when preparing an input file than if
you were just typing the document directly.
Quotation marks like
       ``this''
have to be handled specially, as do quotes within
quotes:
       ``\,`this'            % \, separates the double and single quote.
        is what I just
        wrote, not  `that'\,''.

Dashes come in three sizes: an
       intra-word
dash, a medium dash for number ranges like
       1--2,
and a punctuation
       dash---like
this.

A sentence-ending space should be larger than the
space between words within a sentence.  You
sometimes have to type special commands in
conjunction with punctuation characters to get
this right, as in the following sentence.
       Gnats, gnus, etc.\ all  % `\ ' makes an inter-word space.
       begin with G\@.         % \@ marks end-of-sentence punctuation.
You should check the spaces after periods when
reading your output to make sure you haven't
forgotten any special cases.  Generating an
ellipsis
       \ldots\               % `\ ' is needed after `\ldots' because TeX
                             % ignores spaces after command names like \ldots
                             % made from \ + letters.
                             %
                             % Note how a `%' character causes TeX to ignore
                             % the end of the input line, so these blank lines
                             % do not start a new paragraph.
                             %
with the right spacing around the periods requires
a special command.

\LaTeX\ interprets some common characters as
commands, so you must type special commands to
generate them.  These characters include the
following:
       \$ \& \% \# \{ and \}.

In printing, text is usually emphasized with an
       \emph{italic}
type style.

\begin{em}
   A long segment of text can also be emphasized
   in this way.  Text within such a segment can be
   given \emph{additional} emphasis.
\end{em}

It is sometimes necessary to prevent \LaTeX\ from
breaking a line where it might otherwise do so.
This may be at a space, as between the ``Mr.''\ and
``Jones'' in
       ``Mr.~Jones'',        % ~ produces an unbreakable interword space.
or within a word---especially when the word is a
symbol like
       \mbox{\emph{itemnum}}
that makes little sense when hyphenated across
lines.

Footnotes\footnote{This is an example of a footnote.}
pose no problem.

\LaTeX\ is good at typesetting mathematical formulas
like
       \( x-3y + z = 7 \)
or
       \( a_{1} > x^{2n} + y^{2n} > x' \)
or
       \( \ip{A}{B} = \sum_{i} a_{i} b_{i} \).
The spaces you type in a formula are
ignored.  Remember that a letter like
       $x$                   % $ ... $  and  \( ... \)  are equivalent
is a formula when it denotes a mathematical
symbol, and it should be typed as one.

\section{Displayed Text}

Text is displayed by indenting it from the left
margin.  Quotations are commonly displayed.  There
are short quotations
\begin{quote}
   This is a short quotation.  It consists of a
   single paragraph of text.  See how it is formatted.
\end{quote}
and longer ones.
\begin{quotation}
   This is a longer quotation.  It consists of two
   paragraphs of text, neither of which are
   particularly interesting.

   This is the second paragraph of the quotation.  It
   is just as dull as the first paragraph.
\end{quotation}
Another frequently-displayed structure is a list.
The following is an example of an \emph{itemized}
list.
\begin{itemize}
   \item This is the first item of an itemized list.
         Each item in the list is marked with a ``tick''.
         You don't have to worry about what kind of tick
         mark is used.

   \item This is the second item of the list.  It
         contains another list nested inside it.  The inner
         list is an \emph{enumerated} list.
         \begin{enumerate}
            \item This is the first item of an enumerated
                  list that is nested within the itemized list.

            \item This is the second item of the inner list.
                  \LaTeX\ allows you to nest lists deeper than
                  you really should.
         \end{enumerate}
         This is the rest of the second item of the outer
         list.  It is no more interesting than any other
         part of the item.
   \item This is the third item of the list.
\end{itemize}
You can even display poetry.
\begin{verse}
   There is an environment
    for verse \\             % The \\ command separates lines
   Whose features some poets % within a stanza.
   will curse.

                             % One or more blank lines separate stanzas.

   For instead of making\\
   Them do \emph{all} line breaking, \\
   It allows them to put too many words on a line when they'd rather be
   forced to be terse.
\end{verse}

Mathematical formulas may also be displayed.  A
displayed formula
is
one-line long; multiline
formulas require special formatting instructions.
   \[  \ip{\Gamma}{\psi'} = x'' + y^{2} + z_{i}^{n}\]
Don't start a paragraph with a displayed equation,
nor make one a paragraph by itself.

\end{document}               % End of document.

\section{Bench Huge Section 269}
% synthetic bench section 269
% This is a sample LaTeX input file.  (Version of 12 August 2004.)
%
% A '%' character causes TeX to ignore all remaining text on the line,
% and is used for comments like this one.

\documentclass{article}      % Specifies the document class

                             % The preamble begins here.
\title{An Example Document}  % Declares the document's title.
\author{Leslie Lamport}      % Declares the author's name.
\date{January 21, 1994}      % Deleting this command produces today's date.

\newcommand{\ip}[2]{(#1, #2)}
                             % Defines \ip{arg1}{arg2} to mean
                             % (arg1, arg2).

%\newcommand{\ip}[2]{\langle #1 | #2\rangle}
                             % This is an alternative definition of
                             % \ip that is commented out.

\begin{document}             % End of preamble and beginning of text.

\maketitle                   % Produces the title.

This is an example input file.  Comparing it with
the output it generates can show you how to
produce a simple document of your own.

\section{Ordinary Text}      % Produces section heading.  Lower-level
                             % sections are begun with similar
                             % \subsection and \subsubsection commands.

The ends  of words and sentences are marked
  by   spaces. It  doesn't matter how many
spaces    you type; one is as good as 100.  The
end of   a line counts as a space.

One   or more   blank lines denote the  end
of  a paragraph.

Since any number of consecutive spaces are treated
like a single one, the formatting of the input
file makes no difference to
      \LaTeX,                % The \LaTeX command generates the LaTeX logo.
but it makes a difference to you.  When you use
\LaTeX, making your input file as easy to read
as possible will be a great help as you write
your document and when you change it.  This sample
file shows how you can add comments to your own input
file.

Because printing is different from typewriting,
there are a number of things that you have to do
differently when preparing an input file than if
you were just typing the document directly.
Quotation marks like
       ``this''
have to be handled specially, as do quotes within
quotes:
       ``\,`this'            % \, separates the double and single quote.
        is what I just
        wrote, not  `that'\,''.

Dashes come in three sizes: an
       intra-word
dash, a medium dash for number ranges like
       1--2,
and a punctuation
       dash---like
this.

A sentence-ending space should be larger than the
space between words within a sentence.  You
sometimes have to type special commands in
conjunction with punctuation characters to get
this right, as in the following sentence.
       Gnats, gnus, etc.\ all  % `\ ' makes an inter-word space.
       begin with G\@.         % \@ marks end-of-sentence punctuation.
You should check the spaces after periods when
reading your output to make sure you haven't
forgotten any special cases.  Generating an
ellipsis
       \ldots\               % `\ ' is needed after `\ldots' because TeX
                             % ignores spaces after command names like \ldots
                             % made from \ + letters.
                             %
                             % Note how a `%' character causes TeX to ignore
                             % the end of the input line, so these blank lines
                             % do not start a new paragraph.
                             %
with the right spacing around the periods requires
a special command.

\LaTeX\ interprets some common characters as
commands, so you must type special commands to
generate them.  These characters include the
following:
       \$ \& \% \# \{ and \}.

In printing, text is usually emphasized with an
       \emph{italic}
type style.

\begin{em}
   A long segment of text can also be emphasized
   in this way.  Text within such a segment can be
   given \emph{additional} emphasis.
\end{em}

It is sometimes necessary to prevent \LaTeX\ from
breaking a line where it might otherwise do so.
This may be at a space, as between the ``Mr.''\ and
``Jones'' in
       ``Mr.~Jones'',        % ~ produces an unbreakable interword space.
or within a word---especially when the word is a
symbol like
       \mbox{\emph{itemnum}}
that makes little sense when hyphenated across
lines.

Footnotes\footnote{This is an example of a footnote.}
pose no problem.

\LaTeX\ is good at typesetting mathematical formulas
like
       \( x-3y + z = 7 \)
or
       \( a_{1} > x^{2n} + y^{2n} > x' \)
or
       \( \ip{A}{B} = \sum_{i} a_{i} b_{i} \).
The spaces you type in a formula are
ignored.  Remember that a letter like
       $x$                   % $ ... $  and  \( ... \)  are equivalent
is a formula when it denotes a mathematical
symbol, and it should be typed as one.

\section{Displayed Text}

Text is displayed by indenting it from the left
margin.  Quotations are commonly displayed.  There
are short quotations
\begin{quote}
   This is a short quotation.  It consists of a
   single paragraph of text.  See how it is formatted.
\end{quote}
and longer ones.
\begin{quotation}
   This is a longer quotation.  It consists of two
   paragraphs of text, neither of which are
   particularly interesting.

   This is the second paragraph of the quotation.  It
   is just as dull as the first paragraph.
\end{quotation}
Another frequently-displayed structure is a list.
The following is an example of an \emph{itemized}
list.
\begin{itemize}
   \item This is the first item of an itemized list.
         Each item in the list is marked with a ``tick''.
         You don't have to worry about what kind of tick
         mark is used.

   \item This is the second item of the list.  It
         contains another list nested inside it.  The inner
         list is an \emph{enumerated} list.
         \begin{enumerate}
            \item This is the first item of an enumerated
                  list that is nested within the itemized list.

            \item This is the second item of the inner list.
                  \LaTeX\ allows you to nest lists deeper than
                  you really should.
         \end{enumerate}
         This is the rest of the second item of the outer
         list.  It is no more interesting than any other
         part of the item.
   \item This is the third item of the list.
\end{itemize}
You can even display poetry.
\begin{verse}
   There is an environment
    for verse \\             % The \\ command separates lines
   Whose features some poets % within a stanza.
   will curse.

                             % One or more blank lines separate stanzas.

   For instead of making\\
   Them do \emph{all} line breaking, \\
   It allows them to put too many words on a line when they'd rather be
   forced to be terse.
\end{verse}

Mathematical formulas may also be displayed.  A
displayed formula
is
one-line long; multiline
formulas require special formatting instructions.
   \[  \ip{\Gamma}{\psi'} = x'' + y^{2} + z_{i}^{n}\]
Don't start a paragraph with a displayed equation,
nor make one a paragraph by itself.

\end{document}               % End of document.

\section{Bench Huge Section 270}
% synthetic bench section 270
% This is a sample LaTeX input file.  (Version of 12 August 2004.)
%
% A '%' character causes TeX to ignore all remaining text on the line,
% and is used for comments like this one.

\documentclass{article}      % Specifies the document class

                             % The preamble begins here.
\title{An Example Document}  % Declares the document's title.
\author{Leslie Lamport}      % Declares the author's name.
\date{January 21, 1994}      % Deleting this command produces today's date.

\newcommand{\ip}[2]{(#1, #2)}
                             % Defines \ip{arg1}{arg2} to mean
                             % (arg1, arg2).

%\newcommand{\ip}[2]{\langle #1 | #2\rangle}
                             % This is an alternative definition of
                             % \ip that is commented out.

\begin{document}             % End of preamble and beginning of text.

\maketitle                   % Produces the title.

This is an example input file.  Comparing it with
the output it generates can show you how to
produce a simple document of your own.

\section{Ordinary Text}      % Produces section heading.  Lower-level
                             % sections are begun with similar
                             % \subsection and \subsubsection commands.

The ends  of words and sentences are marked
  by   spaces. It  doesn't matter how many
spaces    you type; one is as good as 100.  The
end of   a line counts as a space.

One   or more   blank lines denote the  end
of  a paragraph.

Since any number of consecutive spaces are treated
like a single one, the formatting of the input
file makes no difference to
      \LaTeX,                % The \LaTeX command generates the LaTeX logo.
but it makes a difference to you.  When you use
\LaTeX, making your input file as easy to read
as possible will be a great help as you write
your document and when you change it.  This sample
file shows how you can add comments to your own input
file.

Because printing is different from typewriting,
there are a number of things that you have to do
differently when preparing an input file than if
you were just typing the document directly.
Quotation marks like
       ``this''
have to be handled specially, as do quotes within
quotes:
       ``\,`this'            % \, separates the double and single quote.
        is what I just
        wrote, not  `that'\,''.

Dashes come in three sizes: an
       intra-word
dash, a medium dash for number ranges like
       1--2,
and a punctuation
       dash---like
this.

A sentence-ending space should be larger than the
space between words within a sentence.  You
sometimes have to type special commands in
conjunction with punctuation characters to get
this right, as in the following sentence.
       Gnats, gnus, etc.\ all  % `\ ' makes an inter-word space.
       begin with G\@.         % \@ marks end-of-sentence punctuation.
You should check the spaces after periods when
reading your output to make sure you haven't
forgotten any special cases.  Generating an
ellipsis
       \ldots\               % `\ ' is needed after `\ldots' because TeX
                             % ignores spaces after command names like \ldots
                             % made from \ + letters.
                             %
                             % Note how a `%' character causes TeX to ignore
                             % the end of the input line, so these blank lines
                             % do not start a new paragraph.
                             %
with the right spacing around the periods requires
a special command.

\LaTeX\ interprets some common characters as
commands, so you must type special commands to
generate them.  These characters include the
following:
       \$ \& \% \# \{ and \}.

In printing, text is usually emphasized with an
       \emph{italic}
type style.

\begin{em}
   A long segment of text can also be emphasized
   in this way.  Text within such a segment can be
   given \emph{additional} emphasis.
\end{em}

It is sometimes necessary to prevent \LaTeX\ from
breaking a line where it might otherwise do so.
This may be at a space, as between the ``Mr.''\ and
``Jones'' in
       ``Mr.~Jones'',        % ~ produces an unbreakable interword space.
or within a word---especially when the word is a
symbol like
       \mbox{\emph{itemnum}}
that makes little sense when hyphenated across
lines.

Footnotes\footnote{This is an example of a footnote.}
pose no problem.

\LaTeX\ is good at typesetting mathematical formulas
like
       \( x-3y + z = 7 \)
or
       \( a_{1} > x^{2n} + y^{2n} > x' \)
or
       \( \ip{A}{B} = \sum_{i} a_{i} b_{i} \).
The spaces you type in a formula are
ignored.  Remember that a letter like
       $x$                   % $ ... $  and  \( ... \)  are equivalent
is a formula when it denotes a mathematical
symbol, and it should be typed as one.

\section{Displayed Text}

Text is displayed by indenting it from the left
margin.  Quotations are commonly displayed.  There
are short quotations
\begin{quote}
   This is a short quotation.  It consists of a
   single paragraph of text.  See how it is formatted.
\end{quote}
and longer ones.
\begin{quotation}
   This is a longer quotation.  It consists of two
   paragraphs of text, neither of which are
   particularly interesting.

   This is the second paragraph of the quotation.  It
   is just as dull as the first paragraph.
\end{quotation}
Another frequently-displayed structure is a list.
The following is an example of an \emph{itemized}
list.
\begin{itemize}
   \item This is the first item of an itemized list.
         Each item in the list is marked with a ``tick''.
         You don't have to worry about what kind of tick
         mark is used.

   \item This is the second item of the list.  It
         contains another list nested inside it.  The inner
         list is an \emph{enumerated} list.
         \begin{enumerate}
            \item This is the first item of an enumerated
                  list that is nested within the itemized list.

            \item This is the second item of the inner list.
                  \LaTeX\ allows you to nest lists deeper than
                  you really should.
         \end{enumerate}
         This is the rest of the second item of the outer
         list.  It is no more interesting than any other
         part of the item.
   \item This is the third item of the list.
\end{itemize}
You can even display poetry.
\begin{verse}
   There is an environment
    for verse \\             % The \\ command separates lines
   Whose features some poets % within a stanza.
   will curse.

                             % One or more blank lines separate stanzas.

   For instead of making\\
   Them do \emph{all} line breaking, \\
   It allows them to put too many words on a line when they'd rather be
   forced to be terse.
\end{verse}

Mathematical formulas may also be displayed.  A
displayed formula
is
one-line long; multiline
formulas require special formatting instructions.
   \[  \ip{\Gamma}{\psi'} = x'' + y^{2} + z_{i}^{n}\]
Don't start a paragraph with a displayed equation,
nor make one a paragraph by itself.

\end{document}               % End of document.

\section{Bench Huge Section 271}
% synthetic bench section 271
% This is a sample LaTeX input file.  (Version of 12 August 2004.)
%
% A '%' character causes TeX to ignore all remaining text on the line,
% and is used for comments like this one.

\documentclass{article}      % Specifies the document class

                             % The preamble begins here.
\title{An Example Document}  % Declares the document's title.
\author{Leslie Lamport}      % Declares the author's name.
\date{January 21, 1994}      % Deleting this command produces today's date.

\newcommand{\ip}[2]{(#1, #2)}
                             % Defines \ip{arg1}{arg2} to mean
                             % (arg1, arg2).

%\newcommand{\ip}[2]{\langle #1 | #2\rangle}
                             % This is an alternative definition of
                             % \ip that is commented out.

\begin{document}             % End of preamble and beginning of text.

\maketitle                   % Produces the title.

This is an example input file.  Comparing it with
the output it generates can show you how to
produce a simple document of your own.

\section{Ordinary Text}      % Produces section heading.  Lower-level
                             % sections are begun with similar
                             % \subsection and \subsubsection commands.

The ends  of words and sentences are marked
  by   spaces. It  doesn't matter how many
spaces    you type; one is as good as 100.  The
end of   a line counts as a space.

One   or more   blank lines denote the  end
of  a paragraph.

Since any number of consecutive spaces are treated
like a single one, the formatting of the input
file makes no difference to
      \LaTeX,                % The \LaTeX command generates the LaTeX logo.
but it makes a difference to you.  When you use
\LaTeX, making your input file as easy to read
as possible will be a great help as you write
your document and when you change it.  This sample
file shows how you can add comments to your own input
file.

Because printing is different from typewriting,
there are a number of things that you have to do
differently when preparing an input file than if
you were just typing the document directly.
Quotation marks like
       ``this''
have to be handled specially, as do quotes within
quotes:
       ``\,`this'            % \, separates the double and single quote.
        is what I just
        wrote, not  `that'\,''.

Dashes come in three sizes: an
       intra-word
dash, a medium dash for number ranges like
       1--2,
and a punctuation
       dash---like
this.

A sentence-ending space should be larger than the
space between words within a sentence.  You
sometimes have to type special commands in
conjunction with punctuation characters to get
this right, as in the following sentence.
       Gnats, gnus, etc.\ all  % `\ ' makes an inter-word space.
       begin with G\@.         % \@ marks end-of-sentence punctuation.
You should check the spaces after periods when
reading your output to make sure you haven't
forgotten any special cases.  Generating an
ellipsis
       \ldots\               % `\ ' is needed after `\ldots' because TeX
                             % ignores spaces after command names like \ldots
                             % made from \ + letters.
                             %
                             % Note how a `%' character causes TeX to ignore
                             % the end of the input line, so these blank lines
                             % do not start a new paragraph.
                             %
with the right spacing around the periods requires
a special command.

\LaTeX\ interprets some common characters as
commands, so you must type special commands to
generate them.  These characters include the
following:
       \$ \& \% \# \{ and \}.

In printing, text is usually emphasized with an
       \emph{italic}
type style.

\begin{em}
   A long segment of text can also be emphasized
   in this way.  Text within such a segment can be
   given \emph{additional} emphasis.
\end{em}

It is sometimes necessary to prevent \LaTeX\ from
breaking a line where it might otherwise do so.
This may be at a space, as between the ``Mr.''\ and
``Jones'' in
       ``Mr.~Jones'',        % ~ produces an unbreakable interword space.
or within a word---especially when the word is a
symbol like
       \mbox{\emph{itemnum}}
that makes little sense when hyphenated across
lines.

Footnotes\footnote{This is an example of a footnote.}
pose no problem.

\LaTeX\ is good at typesetting mathematical formulas
like
       \( x-3y + z = 7 \)
or
       \( a_{1} > x^{2n} + y^{2n} > x' \)
or
       \( \ip{A}{B} = \sum_{i} a_{i} b_{i} \).
The spaces you type in a formula are
ignored.  Remember that a letter like
       $x$                   % $ ... $  and  \( ... \)  are equivalent
is a formula when it denotes a mathematical
symbol, and it should be typed as one.

\section{Displayed Text}

Text is displayed by indenting it from the left
margin.  Quotations are commonly displayed.  There
are short quotations
\begin{quote}
   This is a short quotation.  It consists of a
   single paragraph of text.  See how it is formatted.
\end{quote}
and longer ones.
\begin{quotation}
   This is a longer quotation.  It consists of two
   paragraphs of text, neither of which are
   particularly interesting.

   This is the second paragraph of the quotation.  It
   is just as dull as the first paragraph.
\end{quotation}
Another frequently-displayed structure is a list.
The following is an example of an \emph{itemized}
list.
\begin{itemize}
   \item This is the first item of an itemized list.
         Each item in the list is marked with a ``tick''.
         You don't have to worry about what kind of tick
         mark is used.

   \item This is the second item of the list.  It
         contains another list nested inside it.  The inner
         list is an \emph{enumerated} list.
         \begin{enumerate}
            \item This is the first item of an enumerated
                  list that is nested within the itemized list.

            \item This is the second item of the inner list.
                  \LaTeX\ allows you to nest lists deeper than
                  you really should.
         \end{enumerate}
         This is the rest of the second item of the outer
         list.  It is no more interesting than any other
         part of the item.
   \item This is the third item of the list.
\end{itemize}
You can even display poetry.
\begin{verse}
   There is an environment
    for verse \\             % The \\ command separates lines
   Whose features some poets % within a stanza.
   will curse.

                             % One or more blank lines separate stanzas.

   For instead of making\\
   Them do \emph{all} line breaking, \\
   It allows them to put too many words on a line when they'd rather be
   forced to be terse.
\end{verse}

Mathematical formulas may also be displayed.  A
displayed formula
is
one-line long; multiline
formulas require special formatting instructions.
   \[  \ip{\Gamma}{\psi'} = x'' + y^{2} + z_{i}^{n}\]
Don't start a paragraph with a displayed equation,
nor make one a paragraph by itself.

\end{document}               % End of document.

\section{Bench Huge Section 272}
% synthetic bench section 272
% This is a sample LaTeX input file.  (Version of 12 August 2004.)
%
% A '%' character causes TeX to ignore all remaining text on the line,
% and is used for comments like this one.

\documentclass{article}      % Specifies the document class

                             % The preamble begins here.
\title{An Example Document}  % Declares the document's title.
\author{Leslie Lamport}      % Declares the author's name.
\date{January 21, 1994}      % Deleting this command produces today's date.

\newcommand{\ip}[2]{(#1, #2)}
                             % Defines \ip{arg1}{arg2} to mean
                             % (arg1, arg2).

%\newcommand{\ip}[2]{\langle #1 | #2\rangle}
                             % This is an alternative definition of
                             % \ip that is commented out.

\begin{document}             % End of preamble and beginning of text.

\maketitle                   % Produces the title.

This is an example input file.  Comparing it with
the output it generates can show you how to
produce a simple document of your own.

\section{Ordinary Text}      % Produces section heading.  Lower-level
                             % sections are begun with similar
                             % \subsection and \subsubsection commands.

The ends  of words and sentences are marked
  by   spaces. It  doesn't matter how many
spaces    you type; one is as good as 100.  The
end of   a line counts as a space.

One   or more   blank lines denote the  end
of  a paragraph.

Since any number of consecutive spaces are treated
like a single one, the formatting of the input
file makes no difference to
      \LaTeX,                % The \LaTeX command generates the LaTeX logo.
but it makes a difference to you.  When you use
\LaTeX, making your input file as easy to read
as possible will be a great help as you write
your document and when you change it.  This sample
file shows how you can add comments to your own input
file.

Because printing is different from typewriting,
there are a number of things that you have to do
differently when preparing an input file than if
you were just typing the document directly.
Quotation marks like
       ``this''
have to be handled specially, as do quotes within
quotes:
       ``\,`this'            % \, separates the double and single quote.
        is what I just
        wrote, not  `that'\,''.

Dashes come in three sizes: an
       intra-word
dash, a medium dash for number ranges like
       1--2,
and a punctuation
       dash---like
this.

A sentence-ending space should be larger than the
space between words within a sentence.  You
sometimes have to type special commands in
conjunction with punctuation characters to get
this right, as in the following sentence.
       Gnats, gnus, etc.\ all  % `\ ' makes an inter-word space.
       begin with G\@.         % \@ marks end-of-sentence punctuation.
You should check the spaces after periods when
reading your output to make sure you haven't
forgotten any special cases.  Generating an
ellipsis
       \ldots\               % `\ ' is needed after `\ldots' because TeX
                             % ignores spaces after command names like \ldots
                             % made from \ + letters.
                             %
                             % Note how a `%' character causes TeX to ignore
                             % the end of the input line, so these blank lines
                             % do not start a new paragraph.
                             %
with the right spacing around the periods requires
a special command.

\LaTeX\ interprets some common characters as
commands, so you must type special commands to
generate them.  These characters include the
following:
       \$ \& \% \# \{ and \}.

In printing, text is usually emphasized with an
       \emph{italic}
type style.

\begin{em}
   A long segment of text can also be emphasized
   in this way.  Text within such a segment can be
   given \emph{additional} emphasis.
\end{em}

It is sometimes necessary to prevent \LaTeX\ from
breaking a line where it might otherwise do so.
This may be at a space, as between the ``Mr.''\ and
``Jones'' in
       ``Mr.~Jones'',        % ~ produces an unbreakable interword space.
or within a word---especially when the word is a
symbol like
       \mbox{\emph{itemnum}}
that makes little sense when hyphenated across
lines.

Footnotes\footnote{This is an example of a footnote.}
pose no problem.

\LaTeX\ is good at typesetting mathematical formulas
like
       \( x-3y + z = 7 \)
or
       \( a_{1} > x^{2n} + y^{2n} > x' \)
or
       \( \ip{A}{B} = \sum_{i} a_{i} b_{i} \).
The spaces you type in a formula are
ignored.  Remember that a letter like
       $x$                   % $ ... $  and  \( ... \)  are equivalent
is a formula when it denotes a mathematical
symbol, and it should be typed as one.

\section{Displayed Text}

Text is displayed by indenting it from the left
margin.  Quotations are commonly displayed.  There
are short quotations
\begin{quote}
   This is a short quotation.  It consists of a
   single paragraph of text.  See how it is formatted.
\end{quote}
and longer ones.
\begin{quotation}
   This is a longer quotation.  It consists of two
   paragraphs of text, neither of which are
   particularly interesting.

   This is the second paragraph of the quotation.  It
   is just as dull as the first paragraph.
\end{quotation}
Another frequently-displayed structure is a list.
The following is an example of an \emph{itemized}
list.
\begin{itemize}
   \item This is the first item of an itemized list.
         Each item in the list is marked with a ``tick''.
         You don't have to worry about what kind of tick
         mark is used.

   \item This is the second item of the list.  It
         contains another list nested inside it.  The inner
         list is an \emph{enumerated} list.
         \begin{enumerate}
            \item This is the first item of an enumerated
                  list that is nested within the itemized list.

            \item This is the second item of the inner list.
                  \LaTeX\ allows you to nest lists deeper than
                  you really should.
         \end{enumerate}
         This is the rest of the second item of the outer
         list.  It is no more interesting than any other
         part of the item.
   \item This is the third item of the list.
\end{itemize}
You can even display poetry.
\begin{verse}
   There is an environment
    for verse \\             % The \\ command separates lines
   Whose features some poets % within a stanza.
   will curse.

                             % One or more blank lines separate stanzas.

   For instead of making\\
   Them do \emph{all} line breaking, \\
   It allows them to put too many words on a line when they'd rather be
   forced to be terse.
\end{verse}

Mathematical formulas may also be displayed.  A
displayed formula
is
one-line long; multiline
formulas require special formatting instructions.
   \[  \ip{\Gamma}{\psi'} = x'' + y^{2} + z_{i}^{n}\]
Don't start a paragraph with a displayed equation,
nor make one a paragraph by itself.

\end{document}               % End of document.

\section{Bench Huge Section 273}
% synthetic bench section 273
% This is a sample LaTeX input file.  (Version of 12 August 2004.)
%
% A '%' character causes TeX to ignore all remaining text on the line,
% and is used for comments like this one.

\documentclass{article}      % Specifies the document class

                             % The preamble begins here.
\title{An Example Document}  % Declares the document's title.
\author{Leslie Lamport}      % Declares the author's name.
\date{January 21, 1994}      % Deleting this command produces today's date.

\newcommand{\ip}[2]{(#1, #2)}
                             % Defines \ip{arg1}{arg2} to mean
                             % (arg1, arg2).

%\newcommand{\ip}[2]{\langle #1 | #2\rangle}
                             % This is an alternative definition of
                             % \ip that is commented out.

\begin{document}             % End of preamble and beginning of text.

\maketitle                   % Produces the title.

This is an example input file.  Comparing it with
the output it generates can show you how to
produce a simple document of your own.

\section{Ordinary Text}      % Produces section heading.  Lower-level
                             % sections are begun with similar
                             % \subsection and \subsubsection commands.

The ends  of words and sentences are marked
  by   spaces. It  doesn't matter how many
spaces    you type; one is as good as 100.  The
end of   a line counts as a space.

One   or more   blank lines denote the  end
of  a paragraph.

Since any number of consecutive spaces are treated
like a single one, the formatting of the input
file makes no difference to
      \LaTeX,                % The \LaTeX command generates the LaTeX logo.
but it makes a difference to you.  When you use
\LaTeX, making your input file as easy to read
as possible will be a great help as you write
your document and when you change it.  This sample
file shows how you can add comments to your own input
file.

Because printing is different from typewriting,
there are a number of things that you have to do
differently when preparing an input file than if
you were just typing the document directly.
Quotation marks like
       ``this''
have to be handled specially, as do quotes within
quotes:
       ``\,`this'            % \, separates the double and single quote.
        is what I just
        wrote, not  `that'\,''.

Dashes come in three sizes: an
       intra-word
dash, a medium dash for number ranges like
       1--2,
and a punctuation
       dash---like
this.

A sentence-ending space should be larger than the
space between words within a sentence.  You
sometimes have to type special commands in
conjunction with punctuation characters to get
this right, as in the following sentence.
       Gnats, gnus, etc.\ all  % `\ ' makes an inter-word space.
       begin with G\@.         % \@ marks end-of-sentence punctuation.
You should check the spaces after periods when
reading your output to make sure you haven't
forgotten any special cases.  Generating an
ellipsis
       \ldots\               % `\ ' is needed after `\ldots' because TeX
                             % ignores spaces after command names like \ldots
                             % made from \ + letters.
                             %
                             % Note how a `%' character causes TeX to ignore
                             % the end of the input line, so these blank lines
                             % do not start a new paragraph.
                             %
with the right spacing around the periods requires
a special command.

\LaTeX\ interprets some common characters as
commands, so you must type special commands to
generate them.  These characters include the
following:
       \$ \& \% \# \{ and \}.

In printing, text is usually emphasized with an
       \emph{italic}
type style.

\begin{em}
   A long segment of text can also be emphasized
   in this way.  Text within such a segment can be
   given \emph{additional} emphasis.
\end{em}

It is sometimes necessary to prevent \LaTeX\ from
breaking a line where it might otherwise do so.
This may be at a space, as between the ``Mr.''\ and
``Jones'' in
       ``Mr.~Jones'',        % ~ produces an unbreakable interword space.
or within a word---especially when the word is a
symbol like
       \mbox{\emph{itemnum}}
that makes little sense when hyphenated across
lines.

Footnotes\footnote{This is an example of a footnote.}
pose no problem.

\LaTeX\ is good at typesetting mathematical formulas
like
       \( x-3y + z = 7 \)
or
       \( a_{1} > x^{2n} + y^{2n} > x' \)
or
       \( \ip{A}{B} = \sum_{i} a_{i} b_{i} \).
The spaces you type in a formula are
ignored.  Remember that a letter like
       $x$                   % $ ... $  and  \( ... \)  are equivalent
is a formula when it denotes a mathematical
symbol, and it should be typed as one.

\section{Displayed Text}

Text is displayed by indenting it from the left
margin.  Quotations are commonly displayed.  There
are short quotations
\begin{quote}
   This is a short quotation.  It consists of a
   single paragraph of text.  See how it is formatted.
\end{quote}
and longer ones.
\begin{quotation}
   This is a longer quotation.  It consists of two
   paragraphs of text, neither of which are
   particularly interesting.

   This is the second paragraph of the quotation.  It
   is just as dull as the first paragraph.
\end{quotation}
Another frequently-displayed structure is a list.
The following is an example of an \emph{itemized}
list.
\begin{itemize}
   \item This is the first item of an itemized list.
         Each item in the list is marked with a ``tick''.
         You don't have to worry about what kind of tick
         mark is used.

   \item This is the second item of the list.  It
         contains another list nested inside it.  The inner
         list is an \emph{enumerated} list.
         \begin{enumerate}
            \item This is the first item of an enumerated
                  list that is nested within the itemized list.

            \item This is the second item of the inner list.
                  \LaTeX\ allows you to nest lists deeper than
                  you really should.
         \end{enumerate}
         This is the rest of the second item of the outer
         list.  It is no more interesting than any other
         part of the item.
   \item This is the third item of the list.
\end{itemize}
You can even display poetry.
\begin{verse}
   There is an environment
    for verse \\             % The \\ command separates lines
   Whose features some poets % within a stanza.
   will curse.

                             % One or more blank lines separate stanzas.

   For instead of making\\
   Them do \emph{all} line breaking, \\
   It allows them to put too many words on a line when they'd rather be
   forced to be terse.
\end{verse}

Mathematical formulas may also be displayed.  A
displayed formula
is
one-line long; multiline
formulas require special formatting instructions.
   \[  \ip{\Gamma}{\psi'} = x'' + y^{2} + z_{i}^{n}\]
Don't start a paragraph with a displayed equation,
nor make one a paragraph by itself.

\end{document}               % End of document.

\section{Bench Huge Section 274}
% synthetic bench section 274
% This is a sample LaTeX input file.  (Version of 12 August 2004.)
%
% A '%' character causes TeX to ignore all remaining text on the line,
% and is used for comments like this one.

\documentclass{article}      % Specifies the document class

                             % The preamble begins here.
\title{An Example Document}  % Declares the document's title.
\author{Leslie Lamport}      % Declares the author's name.
\date{January 21, 1994}      % Deleting this command produces today's date.

\newcommand{\ip}[2]{(#1, #2)}
                             % Defines \ip{arg1}{arg2} to mean
                             % (arg1, arg2).

%\newcommand{\ip}[2]{\langle #1 | #2\rangle}
                             % This is an alternative definition of
                             % \ip that is commented out.

\begin{document}             % End of preamble and beginning of text.

\maketitle                   % Produces the title.

This is an example input file.  Comparing it with
the output it generates can show you how to
produce a simple document of your own.

\section{Ordinary Text}      % Produces section heading.  Lower-level
                             % sections are begun with similar
                             % \subsection and \subsubsection commands.

The ends  of words and sentences are marked
  by   spaces. It  doesn't matter how many
spaces    you type; one is as good as 100.  The
end of   a line counts as a space.

One   or more   blank lines denote the  end
of  a paragraph.

Since any number of consecutive spaces are treated
like a single one, the formatting of the input
file makes no difference to
      \LaTeX,                % The \LaTeX command generates the LaTeX logo.
but it makes a difference to you.  When you use
\LaTeX, making your input file as easy to read
as possible will be a great help as you write
your document and when you change it.  This sample
file shows how you can add comments to your own input
file.

Because printing is different from typewriting,
there are a number of things that you have to do
differently when preparing an input file than if
you were just typing the document directly.
Quotation marks like
       ``this''
have to be handled specially, as do quotes within
quotes:
       ``\,`this'            % \, separates the double and single quote.
        is what I just
        wrote, not  `that'\,''.

Dashes come in three sizes: an
       intra-word
dash, a medium dash for number ranges like
       1--2,
and a punctuation
       dash---like
this.

A sentence-ending space should be larger than the
space between words within a sentence.  You
sometimes have to type special commands in
conjunction with punctuation characters to get
this right, as in the following sentence.
       Gnats, gnus, etc.\ all  % `\ ' makes an inter-word space.
       begin with G\@.         % \@ marks end-of-sentence punctuation.
You should check the spaces after periods when
reading your output to make sure you haven't
forgotten any special cases.  Generating an
ellipsis
       \ldots\               % `\ ' is needed after `\ldots' because TeX
                             % ignores spaces after command names like \ldots
                             % made from \ + letters.
                             %
                             % Note how a `%' character causes TeX to ignore
                             % the end of the input line, so these blank lines
                             % do not start a new paragraph.
                             %
with the right spacing around the periods requires
a special command.

\LaTeX\ interprets some common characters as
commands, so you must type special commands to
generate them.  These characters include the
following:
       \$ \& \% \# \{ and \}.

In printing, text is usually emphasized with an
       \emph{italic}
type style.

\begin{em}
   A long segment of text can also be emphasized
   in this way.  Text within such a segment can be
   given \emph{additional} emphasis.
\end{em}

It is sometimes necessary to prevent \LaTeX\ from
breaking a line where it might otherwise do so.
This may be at a space, as between the ``Mr.''\ and
``Jones'' in
       ``Mr.~Jones'',        % ~ produces an unbreakable interword space.
or within a word---especially when the word is a
symbol like
       \mbox{\emph{itemnum}}
that makes little sense when hyphenated across
lines.

Footnotes\footnote{This is an example of a footnote.}
pose no problem.

\LaTeX\ is good at typesetting mathematical formulas
like
       \( x-3y + z = 7 \)
or
       \( a_{1} > x^{2n} + y^{2n} > x' \)
or
       \( \ip{A}{B} = \sum_{i} a_{i} b_{i} \).
The spaces you type in a formula are
ignored.  Remember that a letter like
       $x$                   % $ ... $  and  \( ... \)  are equivalent
is a formula when it denotes a mathematical
symbol, and it should be typed as one.

\section{Displayed Text}

Text is displayed by indenting it from the left
margin.  Quotations are commonly displayed.  There
are short quotations
\begin{quote}
   This is a short quotation.  It consists of a
   single paragraph of text.  See how it is formatted.
\end{quote}
and longer ones.
\begin{quotation}
   This is a longer quotation.  It consists of two
   paragraphs of text, neither of which are
   particularly interesting.

   This is the second paragraph of the quotation.  It
   is just as dull as the first paragraph.
\end{quotation}
Another frequently-displayed structure is a list.
The following is an example of an \emph{itemized}
list.
\begin{itemize}
   \item This is the first item of an itemized list.
         Each item in the list is marked with a ``tick''.
         You don't have to worry about what kind of tick
         mark is used.

   \item This is the second item of the list.  It
         contains another list nested inside it.  The inner
         list is an \emph{enumerated} list.
         \begin{enumerate}
            \item This is the first item of an enumerated
                  list that is nested within the itemized list.

            \item This is the second item of the inner list.
                  \LaTeX\ allows you to nest lists deeper than
                  you really should.
         \end{enumerate}
         This is the rest of the second item of the outer
         list.  It is no more interesting than any other
         part of the item.
   \item This is the third item of the list.
\end{itemize}
You can even display poetry.
\begin{verse}
   There is an environment
    for verse \\             % The \\ command separates lines
   Whose features some poets % within a stanza.
   will curse.

                             % One or more blank lines separate stanzas.

   For instead of making\\
   Them do \emph{all} line breaking, \\
   It allows them to put too many words on a line when they'd rather be
   forced to be terse.
\end{verse}

Mathematical formulas may also be displayed.  A
displayed formula
is
one-line long; multiline
formulas require special formatting instructions.
   \[  \ip{\Gamma}{\psi'} = x'' + y^{2} + z_{i}^{n}\]
Don't start a paragraph with a displayed equation,
nor make one a paragraph by itself.

\end{document}               % End of document.

\section{Bench Huge Section 275}
% synthetic bench section 275
% This is a sample LaTeX input file.  (Version of 12 August 2004.)
%
% A '%' character causes TeX to ignore all remaining text on the line,
% and is used for comments like this one.

\documentclass{article}      % Specifies the document class

                             % The preamble begins here.
\title{An Example Document}  % Declares the document's title.
\author{Leslie Lamport}      % Declares the author's name.
\date{January 21, 1994}      % Deleting this command produces today's date.

\newcommand{\ip}[2]{(#1, #2)}
                             % Defines \ip{arg1}{arg2} to mean
                             % (arg1, arg2).

%\newcommand{\ip}[2]{\langle #1 | #2\rangle}
                             % This is an alternative definition of
                             % \ip that is commented out.

\begin{document}             % End of preamble and beginning of text.

\maketitle                   % Produces the title.

This is an example input file.  Comparing it with
the output it generates can show you how to
produce a simple document of your own.

\section{Ordinary Text}      % Produces section heading.  Lower-level
                             % sections are begun with similar
                             % \subsection and \subsubsection commands.

The ends  of words and sentences are marked
  by   spaces. It  doesn't matter how many
spaces    you type; one is as good as 100.  The
end of   a line counts as a space.

One   or more   blank lines denote the  end
of  a paragraph.

Since any number of consecutive spaces are treated
like a single one, the formatting of the input
file makes no difference to
      \LaTeX,                % The \LaTeX command generates the LaTeX logo.
but it makes a difference to you.  When you use
\LaTeX, making your input file as easy to read
as possible will be a great help as you write
your document and when you change it.  This sample
file shows how you can add comments to your own input
file.

Because printing is different from typewriting,
there are a number of things that you have to do
differently when preparing an input file than if
you were just typing the document directly.
Quotation marks like
       ``this''
have to be handled specially, as do quotes within
quotes:
       ``\,`this'            % \, separates the double and single quote.
        is what I just
        wrote, not  `that'\,''.

Dashes come in three sizes: an
       intra-word
dash, a medium dash for number ranges like
       1--2,
and a punctuation
       dash---like
this.

A sentence-ending space should be larger than the
space between words within a sentence.  You
sometimes have to type special commands in
conjunction with punctuation characters to get
this right, as in the following sentence.
       Gnats, gnus, etc.\ all  % `\ ' makes an inter-word space.
       begin with G\@.         % \@ marks end-of-sentence punctuation.
You should check the spaces after periods when
reading your output to make sure you haven't
forgotten any special cases.  Generating an
ellipsis
       \ldots\               % `\ ' is needed after `\ldots' because TeX
                             % ignores spaces after command names like \ldots
                             % made from \ + letters.
                             %
                             % Note how a `%' character causes TeX to ignore
                             % the end of the input line, so these blank lines
                             % do not start a new paragraph.
                             %
with the right spacing around the periods requires
a special command.

\LaTeX\ interprets some common characters as
commands, so you must type special commands to
generate them.  These characters include the
following:
       \$ \& \% \# \{ and \}.

In printing, text is usually emphasized with an
       \emph{italic}
type style.

\begin{em}
   A long segment of text can also be emphasized
   in this way.  Text within such a segment can be
   given \emph{additional} emphasis.
\end{em}

It is sometimes necessary to prevent \LaTeX\ from
breaking a line where it might otherwise do so.
This may be at a space, as between the ``Mr.''\ and
``Jones'' in
       ``Mr.~Jones'',        % ~ produces an unbreakable interword space.
or within a word---especially when the word is a
symbol like
       \mbox{\emph{itemnum}}
that makes little sense when hyphenated across
lines.

Footnotes\footnote{This is an example of a footnote.}
pose no problem.

\LaTeX\ is good at typesetting mathematical formulas
like
       \( x-3y + z = 7 \)
or
       \( a_{1} > x^{2n} + y^{2n} > x' \)
or
       \( \ip{A}{B} = \sum_{i} a_{i} b_{i} \).
The spaces you type in a formula are
ignored.  Remember that a letter like
       $x$                   % $ ... $  and  \( ... \)  are equivalent
is a formula when it denotes a mathematical
symbol, and it should be typed as one.

\section{Displayed Text}

Text is displayed by indenting it from the left
margin.  Quotations are commonly displayed.  There
are short quotations
\begin{quote}
   This is a short quotation.  It consists of a
   single paragraph of text.  See how it is formatted.
\end{quote}
and longer ones.
\begin{quotation}
   This is a longer quotation.  It consists of two
   paragraphs of text, neither of which are
   particularly interesting.

   This is the second paragraph of the quotation.  It
   is just as dull as the first paragraph.
\end{quotation}
Another frequently-displayed structure is a list.
The following is an example of an \emph{itemized}
list.
\begin{itemize}
   \item This is the first item of an itemized list.
         Each item in the list is marked with a ``tick''.
         You don't have to worry about what kind of tick
         mark is used.

   \item This is the second item of the list.  It
         contains another list nested inside it.  The inner
         list is an \emph{enumerated} list.
         \begin{enumerate}
            \item This is the first item of an enumerated
                  list that is nested within the itemized list.

            \item This is the second item of the inner list.
                  \LaTeX\ allows you to nest lists deeper than
                  you really should.
         \end{enumerate}
         This is the rest of the second item of the outer
         list.  It is no more interesting than any other
         part of the item.
   \item This is the third item of the list.
\end{itemize}
You can even display poetry.
\begin{verse}
   There is an environment
    for verse \\             % The \\ command separates lines
   Whose features some poets % within a stanza.
   will curse.

                             % One or more blank lines separate stanzas.

   For instead of making\\
   Them do \emph{all} line breaking, \\
   It allows them to put too many words on a line when they'd rather be
   forced to be terse.
\end{verse}

Mathematical formulas may also be displayed.  A
displayed formula
is
one-line long; multiline
formulas require special formatting instructions.
   \[  \ip{\Gamma}{\psi'} = x'' + y^{2} + z_{i}^{n}\]
Don't start a paragraph with a displayed equation,
nor make one a paragraph by itself.

\end{document}               % End of document.

\section{Bench Huge Section 276}
% synthetic bench section 276
% This is a sample LaTeX input file.  (Version of 12 August 2004.)
%
% A '%' character causes TeX to ignore all remaining text on the line,
% and is used for comments like this one.

\documentclass{article}      % Specifies the document class

                             % The preamble begins here.
\title{An Example Document}  % Declares the document's title.
\author{Leslie Lamport}      % Declares the author's name.
\date{January 21, 1994}      % Deleting this command produces today's date.

\newcommand{\ip}[2]{(#1, #2)}
                             % Defines \ip{arg1}{arg2} to mean
                             % (arg1, arg2).

%\newcommand{\ip}[2]{\langle #1 | #2\rangle}
                             % This is an alternative definition of
                             % \ip that is commented out.

\begin{document}             % End of preamble and beginning of text.

\maketitle                   % Produces the title.

This is an example input file.  Comparing it with
the output it generates can show you how to
produce a simple document of your own.

\section{Ordinary Text}      % Produces section heading.  Lower-level
                             % sections are begun with similar
                             % \subsection and \subsubsection commands.

The ends  of words and sentences are marked
  by   spaces. It  doesn't matter how many
spaces    you type; one is as good as 100.  The
end of   a line counts as a space.

One   or more   blank lines denote the  end
of  a paragraph.

Since any number of consecutive spaces are treated
like a single one, the formatting of the input
file makes no difference to
      \LaTeX,                % The \LaTeX command generates the LaTeX logo.
but it makes a difference to you.  When you use
\LaTeX, making your input file as easy to read
as possible will be a great help as you write
your document and when you change it.  This sample
file shows how you can add comments to your own input
file.

Because printing is different from typewriting,
there are a number of things that you have to do
differently when preparing an input file than if
you were just typing the document directly.
Quotation marks like
       ``this''
have to be handled specially, as do quotes within
quotes:
       ``\,`this'            % \, separates the double and single quote.
        is what I just
        wrote, not  `that'\,''.

Dashes come in three sizes: an
       intra-word
dash, a medium dash for number ranges like
       1--2,
and a punctuation
       dash---like
this.

A sentence-ending space should be larger than the
space between words within a sentence.  You
sometimes have to type special commands in
conjunction with punctuation characters to get
this right, as in the following sentence.
       Gnats, gnus, etc.\ all  % `\ ' makes an inter-word space.
       begin with G\@.         % \@ marks end-of-sentence punctuation.
You should check the spaces after periods when
reading your output to make sure you haven't
forgotten any special cases.  Generating an
ellipsis
       \ldots\               % `\ ' is needed after `\ldots' because TeX
                             % ignores spaces after command names like \ldots
                             % made from \ + letters.
                             %
                             % Note how a `%' character causes TeX to ignore
                             % the end of the input line, so these blank lines
                             % do not start a new paragraph.
                             %
with the right spacing around the periods requires
a special command.

\LaTeX\ interprets some common characters as
commands, so you must type special commands to
generate them.  These characters include the
following:
       \$ \& \% \# \{ and \}.

In printing, text is usually emphasized with an
       \emph{italic}
type style.

\begin{em}
   A long segment of text can also be emphasized
   in this way.  Text within such a segment can be
   given \emph{additional} emphasis.
\end{em}

It is sometimes necessary to prevent \LaTeX\ from
breaking a line where it might otherwise do so.
This may be at a space, as between the ``Mr.''\ and
``Jones'' in
       ``Mr.~Jones'',        % ~ produces an unbreakable interword space.
or within a word---especially when the word is a
symbol like
       \mbox{\emph{itemnum}}
that makes little sense when hyphenated across
lines.

Footnotes\footnote{This is an example of a footnote.}
pose no problem.

\LaTeX\ is good at typesetting mathematical formulas
like
       \( x-3y + z = 7 \)
or
       \( a_{1} > x^{2n} + y^{2n} > x' \)
or
       \( \ip{A}{B} = \sum_{i} a_{i} b_{i} \).
The spaces you type in a formula are
ignored.  Remember that a letter like
       $x$                   % $ ... $  and  \( ... \)  are equivalent
is a formula when it denotes a mathematical
symbol, and it should be typed as one.

\section{Displayed Text}

Text is displayed by indenting it from the left
margin.  Quotations are commonly displayed.  There
are short quotations
\begin{quote}
   This is a short quotation.  It consists of a
   single paragraph of text.  See how it is formatted.
\end{quote}
and longer ones.
\begin{quotation}
   This is a longer quotation.  It consists of two
   paragraphs of text, neither of which are
   particularly interesting.

   This is the second paragraph of the quotation.  It
   is just as dull as the first paragraph.
\end{quotation}
Another frequently-displayed structure is a list.
The following is an example of an \emph{itemized}
list.
\begin{itemize}
   \item This is the first item of an itemized list.
         Each item in the list is marked with a ``tick''.
         You don't have to worry about what kind of tick
         mark is used.

   \item This is the second item of the list.  It
         contains another list nested inside it.  The inner
         list is an \emph{enumerated} list.
         \begin{enumerate}
            \item This is the first item of an enumerated
                  list that is nested within the itemized list.

            \item This is the second item of the inner list.
                  \LaTeX\ allows you to nest lists deeper than
                  you really should.
         \end{enumerate}
         This is the rest of the second item of the outer
         list.  It is no more interesting than any other
         part of the item.
   \item This is the third item of the list.
\end{itemize}
You can even display poetry.
\begin{verse}
   There is an environment
    for verse \\             % The \\ command separates lines
   Whose features some poets % within a stanza.
   will curse.

                             % One or more blank lines separate stanzas.

   For instead of making\\
   Them do \emph{all} line breaking, \\
   It allows them to put too many words on a line when they'd rather be
   forced to be terse.
\end{verse}

Mathematical formulas may also be displayed.  A
displayed formula
is
one-line long; multiline
formulas require special formatting instructions.
   \[  \ip{\Gamma}{\psi'} = x'' + y^{2} + z_{i}^{n}\]
Don't start a paragraph with a displayed equation,
nor make one a paragraph by itself.

\end{document}               % End of document.

\section{Bench Huge Section 277}
% synthetic bench section 277
% This is a sample LaTeX input file.  (Version of 12 August 2004.)
%
% A '%' character causes TeX to ignore all remaining text on the line,
% and is used for comments like this one.

\documentclass{article}      % Specifies the document class

                             % The preamble begins here.
\title{An Example Document}  % Declares the document's title.
\author{Leslie Lamport}      % Declares the author's name.
\date{January 21, 1994}      % Deleting this command produces today's date.

\newcommand{\ip}[2]{(#1, #2)}
                             % Defines \ip{arg1}{arg2} to mean
                             % (arg1, arg2).

%\newcommand{\ip}[2]{\langle #1 | #2\rangle}
                             % This is an alternative definition of
                             % \ip that is commented out.

\begin{document}             % End of preamble and beginning of text.

\maketitle                   % Produces the title.

This is an example input file.  Comparing it with
the output it generates can show you how to
produce a simple document of your own.

\section{Ordinary Text}      % Produces section heading.  Lower-level
                             % sections are begun with similar
                             % \subsection and \subsubsection commands.

The ends  of words and sentences are marked
  by   spaces. It  doesn't matter how many
spaces    you type; one is as good as 100.  The
end of   a line counts as a space.

One   or more   blank lines denote the  end
of  a paragraph.

Since any number of consecutive spaces are treated
like a single one, the formatting of the input
file makes no difference to
      \LaTeX,                % The \LaTeX command generates the LaTeX logo.
but it makes a difference to you.  When you use
\LaTeX, making your input file as easy to read
as possible will be a great help as you write
your document and when you change it.  This sample
file shows how you can add comments to your own input
file.

Because printing is different from typewriting,
there are a number of things that you have to do
differently when preparing an input file than if
you were just typing the document directly.
Quotation marks like
       ``this''
have to be handled specially, as do quotes within
quotes:
       ``\,`this'            % \, separates the double and single quote.
        is what I just
        wrote, not  `that'\,''.

Dashes come in three sizes: an
       intra-word
dash, a medium dash for number ranges like
       1--2,
and a punctuation
       dash---like
this.

A sentence-ending space should be larger than the
space between words within a sentence.  You
sometimes have to type special commands in
conjunction with punctuation characters to get
this right, as in the following sentence.
       Gnats, gnus, etc.\ all  % `\ ' makes an inter-word space.
       begin with G\@.         % \@ marks end-of-sentence punctuation.
You should check the spaces after periods when
reading your output to make sure you haven't
forgotten any special cases.  Generating an
ellipsis
       \ldots\               % `\ ' is needed after `\ldots' because TeX
                             % ignores spaces after command names like \ldots
                             % made from \ + letters.
                             %
                             % Note how a `%' character causes TeX to ignore
                             % the end of the input line, so these blank lines
                             % do not start a new paragraph.
                             %
with the right spacing around the periods requires
a special command.

\LaTeX\ interprets some common characters as
commands, so you must type special commands to
generate them.  These characters include the
following:
       \$ \& \% \# \{ and \}.

In printing, text is usually emphasized with an
       \emph{italic}
type style.

\begin{em}
   A long segment of text can also be emphasized
   in this way.  Text within such a segment can be
   given \emph{additional} emphasis.
\end{em}

It is sometimes necessary to prevent \LaTeX\ from
breaking a line where it might otherwise do so.
This may be at a space, as between the ``Mr.''\ and
``Jones'' in
       ``Mr.~Jones'',        % ~ produces an unbreakable interword space.
or within a word---especially when the word is a
symbol like
       \mbox{\emph{itemnum}}
that makes little sense when hyphenated across
lines.

Footnotes\footnote{This is an example of a footnote.}
pose no problem.

\LaTeX\ is good at typesetting mathematical formulas
like
       \( x-3y + z = 7 \)
or
       \( a_{1} > x^{2n} + y^{2n} > x' \)
or
       \( \ip{A}{B} = \sum_{i} a_{i} b_{i} \).
The spaces you type in a formula are
ignored.  Remember that a letter like
       $x$                   % $ ... $  and  \( ... \)  are equivalent
is a formula when it denotes a mathematical
symbol, and it should be typed as one.

\section{Displayed Text}

Text is displayed by indenting it from the left
margin.  Quotations are commonly displayed.  There
are short quotations
\begin{quote}
   This is a short quotation.  It consists of a
   single paragraph of text.  See how it is formatted.
\end{quote}
and longer ones.
\begin{quotation}
   This is a longer quotation.  It consists of two
   paragraphs of text, neither of which are
   particularly interesting.

   This is the second paragraph of the quotation.  It
   is just as dull as the first paragraph.
\end{quotation}
Another frequently-displayed structure is a list.
The following is an example of an \emph{itemized}
list.
\begin{itemize}
   \item This is the first item of an itemized list.
         Each item in the list is marked with a ``tick''.
         You don't have to worry about what kind of tick
         mark is used.

   \item This is the second item of the list.  It
         contains another list nested inside it.  The inner
         list is an \emph{enumerated} list.
         \begin{enumerate}
            \item This is the first item of an enumerated
                  list that is nested within the itemized list.

            \item This is the second item of the inner list.
                  \LaTeX\ allows you to nest lists deeper than
                  you really should.
         \end{enumerate}
         This is the rest of the second item of the outer
         list.  It is no more interesting than any other
         part of the item.
   \item This is the third item of the list.
\end{itemize}
You can even display poetry.
\begin{verse}
   There is an environment
    for verse \\             % The \\ command separates lines
   Whose features some poets % within a stanza.
   will curse.

                             % One or more blank lines separate stanzas.

   For instead of making\\
   Them do \emph{all} line breaking, \\
   It allows them to put too many words on a line when they'd rather be
   forced to be terse.
\end{verse}

Mathematical formulas may also be displayed.  A
displayed formula
is
one-line long; multiline
formulas require special formatting instructions.
   \[  \ip{\Gamma}{\psi'} = x'' + y^{2} + z_{i}^{n}\]
Don't start a paragraph with a displayed equation,
nor make one a paragraph by itself.

\end{document}               % End of document.

\section{Bench Huge Section 278}
% synthetic bench section 278
% This is a sample LaTeX input file.  (Version of 12 August 2004.)
%
% A '%' character causes TeX to ignore all remaining text on the line,
% and is used for comments like this one.

\documentclass{article}      % Specifies the document class

                             % The preamble begins here.
\title{An Example Document}  % Declares the document's title.
\author{Leslie Lamport}      % Declares the author's name.
\date{January 21, 1994}      % Deleting this command produces today's date.

\newcommand{\ip}[2]{(#1, #2)}
                             % Defines \ip{arg1}{arg2} to mean
                             % (arg1, arg2).

%\newcommand{\ip}[2]{\langle #1 | #2\rangle}
                             % This is an alternative definition of
                             % \ip that is commented out.

\begin{document}             % End of preamble and beginning of text.

\maketitle                   % Produces the title.

This is an example input file.  Comparing it with
the output it generates can show you how to
produce a simple document of your own.

\section{Ordinary Text}      % Produces section heading.  Lower-level
                             % sections are begun with similar
                             % \subsection and \subsubsection commands.

The ends  of words and sentences are marked
  by   spaces. It  doesn't matter how many
spaces    you type; one is as good as 100.  The
end of   a line counts as a space.

One   or more   blank lines denote the  end
of  a paragraph.

Since any number of consecutive spaces are treated
like a single one, the formatting of the input
file makes no difference to
      \LaTeX,                % The \LaTeX command generates the LaTeX logo.
but it makes a difference to you.  When you use
\LaTeX, making your input file as easy to read
as possible will be a great help as you write
your document and when you change it.  This sample
file shows how you can add comments to your own input
file.

Because printing is different from typewriting,
there are a number of things that you have to do
differently when preparing an input file than if
you were just typing the document directly.
Quotation marks like
       ``this''
have to be handled specially, as do quotes within
quotes:
       ``\,`this'            % \, separates the double and single quote.
        is what I just
        wrote, not  `that'\,''.

Dashes come in three sizes: an
       intra-word
dash, a medium dash for number ranges like
       1--2,
and a punctuation
       dash---like
this.

A sentence-ending space should be larger than the
space between words within a sentence.  You
sometimes have to type special commands in
conjunction with punctuation characters to get
this right, as in the following sentence.
       Gnats, gnus, etc.\ all  % `\ ' makes an inter-word space.
       begin with G\@.         % \@ marks end-of-sentence punctuation.
You should check the spaces after periods when
reading your output to make sure you haven't
forgotten any special cases.  Generating an
ellipsis
       \ldots\               % `\ ' is needed after `\ldots' because TeX
                             % ignores spaces after command names like \ldots
                             % made from \ + letters.
                             %
                             % Note how a `%' character causes TeX to ignore
                             % the end of the input line, so these blank lines
                             % do not start a new paragraph.
                             %
with the right spacing around the periods requires
a special command.

\LaTeX\ interprets some common characters as
commands, so you must type special commands to
generate them.  These characters include the
following:
       \$ \& \% \# \{ and \}.

In printing, text is usually emphasized with an
       \emph{italic}
type style.

\begin{em}
   A long segment of text can also be emphasized
   in this way.  Text within such a segment can be
   given \emph{additional} emphasis.
\end{em}

It is sometimes necessary to prevent \LaTeX\ from
breaking a line where it might otherwise do so.
This may be at a space, as between the ``Mr.''\ and
``Jones'' in
       ``Mr.~Jones'',        % ~ produces an unbreakable interword space.
or within a word---especially when the word is a
symbol like
       \mbox{\emph{itemnum}}
that makes little sense when hyphenated across
lines.

Footnotes\footnote{This is an example of a footnote.}
pose no problem.

\LaTeX\ is good at typesetting mathematical formulas
like
       \( x-3y + z = 7 \)
or
       \( a_{1} > x^{2n} + y^{2n} > x' \)
or
       \( \ip{A}{B} = \sum_{i} a_{i} b_{i} \).
The spaces you type in a formula are
ignored.  Remember that a letter like
       $x$                   % $ ... $  and  \( ... \)  are equivalent
is a formula when it denotes a mathematical
symbol, and it should be typed as one.

\section{Displayed Text}

Text is displayed by indenting it from the left
margin.  Quotations are commonly displayed.  There
are short quotations
\begin{quote}
   This is a short quotation.  It consists of a
   single paragraph of text.  See how it is formatted.
\end{quote}
and longer ones.
\begin{quotation}
   This is a longer quotation.  It consists of two
   paragraphs of text, neither of which are
   particularly interesting.

   This is the second paragraph of the quotation.  It
   is just as dull as the first paragraph.
\end{quotation}
Another frequently-displayed structure is a list.
The following is an example of an \emph{itemized}
list.
\begin{itemize}
   \item This is the first item of an itemized list.
         Each item in the list is marked with a ``tick''.
         You don't have to worry about what kind of tick
         mark is used.

   \item This is the second item of the list.  It
         contains another list nested inside it.  The inner
         list is an \emph{enumerated} list.
         \begin{enumerate}
            \item This is the first item of an enumerated
                  list that is nested within the itemized list.

            \item This is the second item of the inner list.
                  \LaTeX\ allows you to nest lists deeper than
                  you really should.
         \end{enumerate}
         This is the rest of the second item of the outer
         list.  It is no more interesting than any other
         part of the item.
   \item This is the third item of the list.
\end{itemize}
You can even display poetry.
\begin{verse}
   There is an environment
    for verse \\             % The \\ command separates lines
   Whose features some poets % within a stanza.
   will curse.

                             % One or more blank lines separate stanzas.

   For instead of making\\
   Them do \emph{all} line breaking, \\
   It allows them to put too many words on a line when they'd rather be
   forced to be terse.
\end{verse}

Mathematical formulas may also be displayed.  A
displayed formula
is
one-line long; multiline
formulas require special formatting instructions.
   \[  \ip{\Gamma}{\psi'} = x'' + y^{2} + z_{i}^{n}\]
Don't start a paragraph with a displayed equation,
nor make one a paragraph by itself.

\end{document}               % End of document.

\section{Bench Huge Section 279}
% synthetic bench section 279
% This is a sample LaTeX input file.  (Version of 12 August 2004.)
%
% A '%' character causes TeX to ignore all remaining text on the line,
% and is used for comments like this one.

\documentclass{article}      % Specifies the document class

                             % The preamble begins here.
\title{An Example Document}  % Declares the document's title.
\author{Leslie Lamport}      % Declares the author's name.
\date{January 21, 1994}      % Deleting this command produces today's date.

\newcommand{\ip}[2]{(#1, #2)}
                             % Defines \ip{arg1}{arg2} to mean
                             % (arg1, arg2).

%\newcommand{\ip}[2]{\langle #1 | #2\rangle}
                             % This is an alternative definition of
                             % \ip that is commented out.

\begin{document}             % End of preamble and beginning of text.

\maketitle                   % Produces the title.

This is an example input file.  Comparing it with
the output it generates can show you how to
produce a simple document of your own.

\section{Ordinary Text}      % Produces section heading.  Lower-level
                             % sections are begun with similar
                             % \subsection and \subsubsection commands.

The ends  of words and sentences are marked
  by   spaces. It  doesn't matter how many
spaces    you type; one is as good as 100.  The
end of   a line counts as a space.

One   or more   blank lines denote the  end
of  a paragraph.

Since any number of consecutive spaces are treated
like a single one, the formatting of the input
file makes no difference to
      \LaTeX,                % The \LaTeX command generates the LaTeX logo.
but it makes a difference to you.  When you use
\LaTeX, making your input file as easy to read
as possible will be a great help as you write
your document and when you change it.  This sample
file shows how you can add comments to your own input
file.

Because printing is different from typewriting,
there are a number of things that you have to do
differently when preparing an input file than if
you were just typing the document directly.
Quotation marks like
       ``this''
have to be handled specially, as do quotes within
quotes:
       ``\,`this'            % \, separates the double and single quote.
        is what I just
        wrote, not  `that'\,''.

Dashes come in three sizes: an
       intra-word
dash, a medium dash for number ranges like
       1--2,
and a punctuation
       dash---like
this.

A sentence-ending space should be larger than the
space between words within a sentence.  You
sometimes have to type special commands in
conjunction with punctuation characters to get
this right, as in the following sentence.
       Gnats, gnus, etc.\ all  % `\ ' makes an inter-word space.
       begin with G\@.         % \@ marks end-of-sentence punctuation.
You should check the spaces after periods when
reading your output to make sure you haven't
forgotten any special cases.  Generating an
ellipsis
       \ldots\               % `\ ' is needed after `\ldots' because TeX
                             % ignores spaces after command names like \ldots
                             % made from \ + letters.
                             %
                             % Note how a `%' character causes TeX to ignore
                             % the end of the input line, so these blank lines
                             % do not start a new paragraph.
                             %
with the right spacing around the periods requires
a special command.

\LaTeX\ interprets some common characters as
commands, so you must type special commands to
generate them.  These characters include the
following:
       \$ \& \% \# \{ and \}.

In printing, text is usually emphasized with an
       \emph{italic}
type style.

\begin{em}
   A long segment of text can also be emphasized
   in this way.  Text within such a segment can be
   given \emph{additional} emphasis.
\end{em}

It is sometimes necessary to prevent \LaTeX\ from
breaking a line where it might otherwise do so.
This may be at a space, as between the ``Mr.''\ and
``Jones'' in
       ``Mr.~Jones'',        % ~ produces an unbreakable interword space.
or within a word---especially when the word is a
symbol like
       \mbox{\emph{itemnum}}
that makes little sense when hyphenated across
lines.

Footnotes\footnote{This is an example of a footnote.}
pose no problem.

\LaTeX\ is good at typesetting mathematical formulas
like
       \( x-3y + z = 7 \)
or
       \( a_{1} > x^{2n} + y^{2n} > x' \)
or
       \( \ip{A}{B} = \sum_{i} a_{i} b_{i} \).
The spaces you type in a formula are
ignored.  Remember that a letter like
       $x$                   % $ ... $  and  \( ... \)  are equivalent
is a formula when it denotes a mathematical
symbol, and it should be typed as one.

\section{Displayed Text}

Text is displayed by indenting it from the left
margin.  Quotations are commonly displayed.  There
are short quotations
\begin{quote}
   This is a short quotation.  It consists of a
   single paragraph of text.  See how it is formatted.
\end{quote}
and longer ones.
\begin{quotation}
   This is a longer quotation.  It consists of two
   paragraphs of text, neither of which are
   particularly interesting.

   This is the second paragraph of the quotation.  It
   is just as dull as the first paragraph.
\end{quotation}
Another frequently-displayed structure is a list.
The following is an example of an \emph{itemized}
list.
\begin{itemize}
   \item This is the first item of an itemized list.
         Each item in the list is marked with a ``tick''.
         You don't have to worry about what kind of tick
         mark is used.

   \item This is the second item of the list.  It
         contains another list nested inside it.  The inner
         list is an \emph{enumerated} list.
         \begin{enumerate}
            \item This is the first item of an enumerated
                  list that is nested within the itemized list.

            \item This is the second item of the inner list.
                  \LaTeX\ allows you to nest lists deeper than
                  you really should.
         \end{enumerate}
         This is the rest of the second item of the outer
         list.  It is no more interesting than any other
         part of the item.
   \item This is the third item of the list.
\end{itemize}
You can even display poetry.
\begin{verse}
   There is an environment
    for verse \\             % The \\ command separates lines
   Whose features some poets % within a stanza.
   will curse.

                             % One or more blank lines separate stanzas.

   For instead of making\\
   Them do \emph{all} line breaking, \\
   It allows them to put too many words on a line when they'd rather be
   forced to be terse.
\end{verse}

Mathematical formulas may also be displayed.  A
displayed formula
is
one-line long; multiline
formulas require special formatting instructions.
   \[  \ip{\Gamma}{\psi'} = x'' + y^{2} + z_{i}^{n}\]
Don't start a paragraph with a displayed equation,
nor make one a paragraph by itself.

\end{document}               % End of document.

\section{Bench Huge Section 280}
% synthetic bench section 280
% This is a sample LaTeX input file.  (Version of 12 August 2004.)
%
% A '%' character causes TeX to ignore all remaining text on the line,
% and is used for comments like this one.

\documentclass{article}      % Specifies the document class

                             % The preamble begins here.
\title{An Example Document}  % Declares the document's title.
\author{Leslie Lamport}      % Declares the author's name.
\date{January 21, 1994}      % Deleting this command produces today's date.

\newcommand{\ip}[2]{(#1, #2)}
                             % Defines \ip{arg1}{arg2} to mean
                             % (arg1, arg2).

%\newcommand{\ip}[2]{\langle #1 | #2\rangle}
                             % This is an alternative definition of
                             % \ip that is commented out.

\begin{document}             % End of preamble and beginning of text.

\maketitle                   % Produces the title.

This is an example input file.  Comparing it with
the output it generates can show you how to
produce a simple document of your own.

\section{Ordinary Text}      % Produces section heading.  Lower-level
                             % sections are begun with similar
                             % \subsection and \subsubsection commands.

The ends  of words and sentences are marked
  by   spaces. It  doesn't matter how many
spaces    you type; one is as good as 100.  The
end of   a line counts as a space.

One   or more   blank lines denote the  end
of  a paragraph.

Since any number of consecutive spaces are treated
like a single one, the formatting of the input
file makes no difference to
      \LaTeX,                % The \LaTeX command generates the LaTeX logo.
but it makes a difference to you.  When you use
\LaTeX, making your input file as easy to read
as possible will be a great help as you write
your document and when you change it.  This sample
file shows how you can add comments to your own input
file.

Because printing is different from typewriting,
there are a number of things that you have to do
differently when preparing an input file than if
you were just typing the document directly.
Quotation marks like
       ``this''
have to be handled specially, as do quotes within
quotes:
       ``\,`this'            % \, separates the double and single quote.
        is what I just
        wrote, not  `that'\,''.

Dashes come in three sizes: an
       intra-word
dash, a medium dash for number ranges like
       1--2,
and a punctuation
       dash---like
this.

A sentence-ending space should be larger than the
space between words within a sentence.  You
sometimes have to type special commands in
conjunction with punctuation characters to get
this right, as in the following sentence.
       Gnats, gnus, etc.\ all  % `\ ' makes an inter-word space.
       begin with G\@.         % \@ marks end-of-sentence punctuation.
You should check the spaces after periods when
reading your output to make sure you haven't
forgotten any special cases.  Generating an
ellipsis
       \ldots\               % `\ ' is needed after `\ldots' because TeX
                             % ignores spaces after command names like \ldots
                             % made from \ + letters.
                             %
                             % Note how a `%' character causes TeX to ignore
                             % the end of the input line, so these blank lines
                             % do not start a new paragraph.
                             %
with the right spacing around the periods requires
a special command.

\LaTeX\ interprets some common characters as
commands, so you must type special commands to
generate them.  These characters include the
following:
       \$ \& \% \# \{ and \}.

In printing, text is usually emphasized with an
       \emph{italic}
type style.

\begin{em}
   A long segment of text can also be emphasized
   in this way.  Text within such a segment can be
   given \emph{additional} emphasis.
\end{em}

It is sometimes necessary to prevent \LaTeX\ from
breaking a line where it might otherwise do so.
This may be at a space, as between the ``Mr.''\ and
``Jones'' in
       ``Mr.~Jones'',        % ~ produces an unbreakable interword space.
or within a word---especially when the word is a
symbol like
       \mbox{\emph{itemnum}}
that makes little sense when hyphenated across
lines.

Footnotes\footnote{This is an example of a footnote.}
pose no problem.

\LaTeX\ is good at typesetting mathematical formulas
like
       \( x-3y + z = 7 \)
or
       \( a_{1} > x^{2n} + y^{2n} > x' \)
or
       \( \ip{A}{B} = \sum_{i} a_{i} b_{i} \).
The spaces you type in a formula are
ignored.  Remember that a letter like
       $x$                   % $ ... $  and  \( ... \)  are equivalent
is a formula when it denotes a mathematical
symbol, and it should be typed as one.

\section{Displayed Text}

Text is displayed by indenting it from the left
margin.  Quotations are commonly displayed.  There
are short quotations
\begin{quote}
   This is a short quotation.  It consists of a
   single paragraph of text.  See how it is formatted.
\end{quote}
and longer ones.
\begin{quotation}
   This is a longer quotation.  It consists of two
   paragraphs of text, neither of which are
   particularly interesting.

   This is the second paragraph of the quotation.  It
   is just as dull as the first paragraph.
\end{quotation}
Another frequently-displayed structure is a list.
The following is an example of an \emph{itemized}
list.
\begin{itemize}
   \item This is the first item of an itemized list.
         Each item in the list is marked with a ``tick''.
         You don't have to worry about what kind of tick
         mark is used.

   \item This is the second item of the list.  It
         contains another list nested inside it.  The inner
         list is an \emph{enumerated} list.
         \begin{enumerate}
            \item This is the first item of an enumerated
                  list that is nested within the itemized list.

            \item This is the second item of the inner list.
                  \LaTeX\ allows you to nest lists deeper than
                  you really should.
         \end{enumerate}
         This is the rest of the second item of the outer
         list.  It is no more interesting than any other
         part of the item.
   \item This is the third item of the list.
\end{itemize}
You can even display poetry.
\begin{verse}
   There is an environment
    for verse \\             % The \\ command separates lines
   Whose features some poets % within a stanza.
   will curse.

                             % One or more blank lines separate stanzas.

   For instead of making\\
   Them do \emph{all} line breaking, \\
   It allows them to put too many words on a line when they'd rather be
   forced to be terse.
\end{verse}

Mathematical formulas may also be displayed.  A
displayed formula
is
one-line long; multiline
formulas require special formatting instructions.
   \[  \ip{\Gamma}{\psi'} = x'' + y^{2} + z_{i}^{n}\]
Don't start a paragraph with a displayed equation,
nor make one a paragraph by itself.

\end{document}               % End of document.

\section{Bench Huge Section 281}
% synthetic bench section 281
% This is a sample LaTeX input file.  (Version of 12 August 2004.)
%
% A '%' character causes TeX to ignore all remaining text on the line,
% and is used for comments like this one.

\documentclass{article}      % Specifies the document class

                             % The preamble begins here.
\title{An Example Document}  % Declares the document's title.
\author{Leslie Lamport}      % Declares the author's name.
\date{January 21, 1994}      % Deleting this command produces today's date.

\newcommand{\ip}[2]{(#1, #2)}
                             % Defines \ip{arg1}{arg2} to mean
                             % (arg1, arg2).

%\newcommand{\ip}[2]{\langle #1 | #2\rangle}
                             % This is an alternative definition of
                             % \ip that is commented out.

\begin{document}             % End of preamble and beginning of text.

\maketitle                   % Produces the title.

This is an example input file.  Comparing it with
the output it generates can show you how to
produce a simple document of your own.

\section{Ordinary Text}      % Produces section heading.  Lower-level
                             % sections are begun with similar
                             % \subsection and \subsubsection commands.

The ends  of words and sentences are marked
  by   spaces. It  doesn't matter how many
spaces    you type; one is as good as 100.  The
end of   a line counts as a space.

One   or more   blank lines denote the  end
of  a paragraph.

Since any number of consecutive spaces are treated
like a single one, the formatting of the input
file makes no difference to
      \LaTeX,                % The \LaTeX command generates the LaTeX logo.
but it makes a difference to you.  When you use
\LaTeX, making your input file as easy to read
as possible will be a great help as you write
your document and when you change it.  This sample
file shows how you can add comments to your own input
file.

Because printing is different from typewriting,
there are a number of things that you have to do
differently when preparing an input file than if
you were just typing the document directly.
Quotation marks like
       ``this''
have to be handled specially, as do quotes within
quotes:
       ``\,`this'            % \, separates the double and single quote.
        is what I just
        wrote, not  `that'\,''.

Dashes come in three sizes: an
       intra-word
dash, a medium dash for number ranges like
       1--2,
and a punctuation
       dash---like
this.

A sentence-ending space should be larger than the
space between words within a sentence.  You
sometimes have to type special commands in
conjunction with punctuation characters to get
this right, as in the following sentence.
       Gnats, gnus, etc.\ all  % `\ ' makes an inter-word space.
       begin with G\@.         % \@ marks end-of-sentence punctuation.
You should check the spaces after periods when
reading your output to make sure you haven't
forgotten any special cases.  Generating an
ellipsis
       \ldots\               % `\ ' is needed after `\ldots' because TeX
                             % ignores spaces after command names like \ldots
                             % made from \ + letters.
                             %
                             % Note how a `%' character causes TeX to ignore
                             % the end of the input line, so these blank lines
                             % do not start a new paragraph.
                             %
with the right spacing around the periods requires
a special command.

\LaTeX\ interprets some common characters as
commands, so you must type special commands to
generate them.  These characters include the
following:
       \$ \& \% \# \{ and \}.

In printing, text is usually emphasized with an
       \emph{italic}
type style.

\begin{em}
   A long segment of text can also be emphasized
   in this way.  Text within such a segment can be
   given \emph{additional} emphasis.
\end{em}

It is sometimes necessary to prevent \LaTeX\ from
breaking a line where it might otherwise do so.
This may be at a space, as between the ``Mr.''\ and
``Jones'' in
       ``Mr.~Jones'',        % ~ produces an unbreakable interword space.
or within a word---especially when the word is a
symbol like
       \mbox{\emph{itemnum}}
that makes little sense when hyphenated across
lines.

Footnotes\footnote{This is an example of a footnote.}
pose no problem.

\LaTeX\ is good at typesetting mathematical formulas
like
       \( x-3y + z = 7 \)
or
       \( a_{1} > x^{2n} + y^{2n} > x' \)
or
       \( \ip{A}{B} = \sum_{i} a_{i} b_{i} \).
The spaces you type in a formula are
ignored.  Remember that a letter like
       $x$                   % $ ... $  and  \( ... \)  are equivalent
is a formula when it denotes a mathematical
symbol, and it should be typed as one.

\section{Displayed Text}

Text is displayed by indenting it from the left
margin.  Quotations are commonly displayed.  There
are short quotations
\begin{quote}
   This is a short quotation.  It consists of a
   single paragraph of text.  See how it is formatted.
\end{quote}
and longer ones.
\begin{quotation}
   This is a longer quotation.  It consists of two
   paragraphs of text, neither of which are
   particularly interesting.

   This is the second paragraph of the quotation.  It
   is just as dull as the first paragraph.
\end{quotation}
Another frequently-displayed structure is a list.
The following is an example of an \emph{itemized}
list.
\begin{itemize}
   \item This is the first item of an itemized list.
         Each item in the list is marked with a ``tick''.
         You don't have to worry about what kind of tick
         mark is used.

   \item This is the second item of the list.  It
         contains another list nested inside it.  The inner
         list is an \emph{enumerated} list.
         \begin{enumerate}
            \item This is the first item of an enumerated
                  list that is nested within the itemized list.

            \item This is the second item of the inner list.
                  \LaTeX\ allows you to nest lists deeper than
                  you really should.
         \end{enumerate}
         This is the rest of the second item of the outer
         list.  It is no more interesting than any other
         part of the item.
   \item This is the third item of the list.
\end{itemize}
You can even display poetry.
\begin{verse}
   There is an environment
    for verse \\             % The \\ command separates lines
   Whose features some poets % within a stanza.
   will curse.

                             % One or more blank lines separate stanzas.

   For instead of making\\
   Them do \emph{all} line breaking, \\
   It allows them to put too many words on a line when they'd rather be
   forced to be terse.
\end{verse}

Mathematical formulas may also be displayed.  A
displayed formula
is
one-line long; multiline
formulas require special formatting instructions.
   \[  \ip{\Gamma}{\psi'} = x'' + y^{2} + z_{i}^{n}\]
Don't start a paragraph with a displayed equation,
nor make one a paragraph by itself.

\end{document}               % End of document.

\section{Bench Huge Section 282}
% synthetic bench section 282
% This is a sample LaTeX input file.  (Version of 12 August 2004.)
%
% A '%' character causes TeX to ignore all remaining text on the line,
% and is used for comments like this one.

\documentclass{article}      % Specifies the document class

                             % The preamble begins here.
\title{An Example Document}  % Declares the document's title.
\author{Leslie Lamport}      % Declares the author's name.
\date{January 21, 1994}      % Deleting this command produces today's date.

\newcommand{\ip}[2]{(#1, #2)}
                             % Defines \ip{arg1}{arg2} to mean
                             % (arg1, arg2).

%\newcommand{\ip}[2]{\langle #1 | #2\rangle}
                             % This is an alternative definition of
                             % \ip that is commented out.

\begin{document}             % End of preamble and beginning of text.

\maketitle                   % Produces the title.

This is an example input file.  Comparing it with
the output it generates can show you how to
produce a simple document of your own.

\section{Ordinary Text}      % Produces section heading.  Lower-level
                             % sections are begun with similar
                             % \subsection and \subsubsection commands.

The ends  of words and sentences are marked
  by   spaces. It  doesn't matter how many
spaces    you type; one is as good as 100.  The
end of   a line counts as a space.

One   or more   blank lines denote the  end
of  a paragraph.

Since any number of consecutive spaces are treated
like a single one, the formatting of the input
file makes no difference to
      \LaTeX,                % The \LaTeX command generates the LaTeX logo.
but it makes a difference to you.  When you use
\LaTeX, making your input file as easy to read
as possible will be a great help as you write
your document and when you change it.  This sample
file shows how you can add comments to your own input
file.

Because printing is different from typewriting,
there are a number of things that you have to do
differently when preparing an input file than if
you were just typing the document directly.
Quotation marks like
       ``this''
have to be handled specially, as do quotes within
quotes:
       ``\,`this'            % \, separates the double and single quote.
        is what I just
        wrote, not  `that'\,''.

Dashes come in three sizes: an
       intra-word
dash, a medium dash for number ranges like
       1--2,
and a punctuation
       dash---like
this.

A sentence-ending space should be larger than the
space between words within a sentence.  You
sometimes have to type special commands in
conjunction with punctuation characters to get
this right, as in the following sentence.
       Gnats, gnus, etc.\ all  % `\ ' makes an inter-word space.
       begin with G\@.         % \@ marks end-of-sentence punctuation.
You should check the spaces after periods when
reading your output to make sure you haven't
forgotten any special cases.  Generating an
ellipsis
       \ldots\               % `\ ' is needed after `\ldots' because TeX
                             % ignores spaces after command names like \ldots
                             % made from \ + letters.
                             %
                             % Note how a `%' character causes TeX to ignore
                             % the end of the input line, so these blank lines
                             % do not start a new paragraph.
                             %
with the right spacing around the periods requires
a special command.

\LaTeX\ interprets some common characters as
commands, so you must type special commands to
generate them.  These characters include the
following:
       \$ \& \% \# \{ and \}.

In printing, text is usually emphasized with an
       \emph{italic}
type style.

\begin{em}
   A long segment of text can also be emphasized
   in this way.  Text within such a segment can be
   given \emph{additional} emphasis.
\end{em}

It is sometimes necessary to prevent \LaTeX\ from
breaking a line where it might otherwise do so.
This may be at a space, as between the ``Mr.''\ and
``Jones'' in
       ``Mr.~Jones'',        % ~ produces an unbreakable interword space.
or within a word---especially when the word is a
symbol like
       \mbox{\emph{itemnum}}
that makes little sense when hyphenated across
lines.

Footnotes\footnote{This is an example of a footnote.}
pose no problem.

\LaTeX\ is good at typesetting mathematical formulas
like
       \( x-3y + z = 7 \)
or
       \( a_{1} > x^{2n} + y^{2n} > x' \)
or
       \( \ip{A}{B} = \sum_{i} a_{i} b_{i} \).
The spaces you type in a formula are
ignored.  Remember that a letter like
       $x$                   % $ ... $  and  \( ... \)  are equivalent
is a formula when it denotes a mathematical
symbol, and it should be typed as one.

\section{Displayed Text}

Text is displayed by indenting it from the left
margin.  Quotations are commonly displayed.  There
are short quotations
\begin{quote}
   This is a short quotation.  It consists of a
   single paragraph of text.  See how it is formatted.
\end{quote}
and longer ones.
\begin{quotation}
   This is a longer quotation.  It consists of two
   paragraphs of text, neither of which are
   particularly interesting.

   This is the second paragraph of the quotation.  It
   is just as dull as the first paragraph.
\end{quotation}
Another frequently-displayed structure is a list.
The following is an example of an \emph{itemized}
list.
\begin{itemize}
   \item This is the first item of an itemized list.
         Each item in the list is marked with a ``tick''.
         You don't have to worry about what kind of tick
         mark is used.

   \item This is the second item of the list.  It
         contains another list nested inside it.  The inner
         list is an \emph{enumerated} list.
         \begin{enumerate}
            \item This is the first item of an enumerated
                  list that is nested within the itemized list.

            \item This is the second item of the inner list.
                  \LaTeX\ allows you to nest lists deeper than
                  you really should.
         \end{enumerate}
         This is the rest of the second item of the outer
         list.  It is no more interesting than any other
         part of the item.
   \item This is the third item of the list.
\end{itemize}
You can even display poetry.
\begin{verse}
   There is an environment
    for verse \\             % The \\ command separates lines
   Whose features some poets % within a stanza.
   will curse.

                             % One or more blank lines separate stanzas.

   For instead of making\\
   Them do \emph{all} line breaking, \\
   It allows them to put too many words on a line when they'd rather be
   forced to be terse.
\end{verse}

Mathematical formulas may also be displayed.  A
displayed formula
is
one-line long; multiline
formulas require special formatting instructions.
   \[  \ip{\Gamma}{\psi'} = x'' + y^{2} + z_{i}^{n}\]
Don't start a paragraph with a displayed equation,
nor make one a paragraph by itself.

\end{document}               % End of document.

\section{Bench Huge Section 283}
% synthetic bench section 283
% This is a sample LaTeX input file.  (Version of 12 August 2004.)
%
% A '%' character causes TeX to ignore all remaining text on the line,
% and is used for comments like this one.

\documentclass{article}      % Specifies the document class

                             % The preamble begins here.
\title{An Example Document}  % Declares the document's title.
\author{Leslie Lamport}      % Declares the author's name.
\date{January 21, 1994}      % Deleting this command produces today's date.

\newcommand{\ip}[2]{(#1, #2)}
                             % Defines \ip{arg1}{arg2} to mean
                             % (arg1, arg2).

%\newcommand{\ip}[2]{\langle #1 | #2\rangle}
                             % This is an alternative definition of
                             % \ip that is commented out.

\begin{document}             % End of preamble and beginning of text.

\maketitle                   % Produces the title.

This is an example input file.  Comparing it with
the output it generates can show you how to
produce a simple document of your own.

\section{Ordinary Text}      % Produces section heading.  Lower-level
                             % sections are begun with similar
                             % \subsection and \subsubsection commands.

The ends  of words and sentences are marked
  by   spaces. It  doesn't matter how many
spaces    you type; one is as good as 100.  The
end of   a line counts as a space.

One   or more   blank lines denote the  end
of  a paragraph.

Since any number of consecutive spaces are treated
like a single one, the formatting of the input
file makes no difference to
      \LaTeX,                % The \LaTeX command generates the LaTeX logo.
but it makes a difference to you.  When you use
\LaTeX, making your input file as easy to read
as possible will be a great help as you write
your document and when you change it.  This sample
file shows how you can add comments to your own input
file.

Because printing is different from typewriting,
there are a number of things that you have to do
differently when preparing an input file than if
you were just typing the document directly.
Quotation marks like
       ``this''
have to be handled specially, as do quotes within
quotes:
       ``\,`this'            % \, separates the double and single quote.
        is what I just
        wrote, not  `that'\,''.

Dashes come in three sizes: an
       intra-word
dash, a medium dash for number ranges like
       1--2,
and a punctuation
       dash---like
this.

A sentence-ending space should be larger than the
space between words within a sentence.  You
sometimes have to type special commands in
conjunction with punctuation characters to get
this right, as in the following sentence.
       Gnats, gnus, etc.\ all  % `\ ' makes an inter-word space.
       begin with G\@.         % \@ marks end-of-sentence punctuation.
You should check the spaces after periods when
reading your output to make sure you haven't
forgotten any special cases.  Generating an
ellipsis
       \ldots\               % `\ ' is needed after `\ldots' because TeX
                             % ignores spaces after command names like \ldots
                             % made from \ + letters.
                             %
                             % Note how a `%' character causes TeX to ignore
                             % the end of the input line, so these blank lines
                             % do not start a new paragraph.
                             %
with the right spacing around the periods requires
a special command.

\LaTeX\ interprets some common characters as
commands, so you must type special commands to
generate them.  These characters include the
following:
       \$ \& \% \# \{ and \}.

In printing, text is usually emphasized with an
       \emph{italic}
type style.

\begin{em}
   A long segment of text can also be emphasized
   in this way.  Text within such a segment can be
   given \emph{additional} emphasis.
\end{em}

It is sometimes necessary to prevent \LaTeX\ from
breaking a line where it might otherwise do so.
This may be at a space, as between the ``Mr.''\ and
``Jones'' in
       ``Mr.~Jones'',        % ~ produces an unbreakable interword space.
or within a word---especially when the word is a
symbol like
       \mbox{\emph{itemnum}}
that makes little sense when hyphenated across
lines.

Footnotes\footnote{This is an example of a footnote.}
pose no problem.

\LaTeX\ is good at typesetting mathematical formulas
like
       \( x-3y + z = 7 \)
or
       \( a_{1} > x^{2n} + y^{2n} > x' \)
or
       \( \ip{A}{B} = \sum_{i} a_{i} b_{i} \).
The spaces you type in a formula are
ignored.  Remember that a letter like
       $x$                   % $ ... $  and  \( ... \)  are equivalent
is a formula when it denotes a mathematical
symbol, and it should be typed as one.

\section{Displayed Text}

Text is displayed by indenting it from the left
margin.  Quotations are commonly displayed.  There
are short quotations
\begin{quote}
   This is a short quotation.  It consists of a
   single paragraph of text.  See how it is formatted.
\end{quote}
and longer ones.
\begin{quotation}
   This is a longer quotation.  It consists of two
   paragraphs of text, neither of which are
   particularly interesting.

   This is the second paragraph of the quotation.  It
   is just as dull as the first paragraph.
\end{quotation}
Another frequently-displayed structure is a list.
The following is an example of an \emph{itemized}
list.
\begin{itemize}
   \item This is the first item of an itemized list.
         Each item in the list is marked with a ``tick''.
         You don't have to worry about what kind of tick
         mark is used.

   \item This is the second item of the list.  It
         contains another list nested inside it.  The inner
         list is an \emph{enumerated} list.
         \begin{enumerate}
            \item This is the first item of an enumerated
                  list that is nested within the itemized list.

            \item This is the second item of the inner list.
                  \LaTeX\ allows you to nest lists deeper than
                  you really should.
         \end{enumerate}
         This is the rest of the second item of the outer
         list.  It is no more interesting than any other
         part of the item.
   \item This is the third item of the list.
\end{itemize}
You can even display poetry.
\begin{verse}
   There is an environment
    for verse \\             % The \\ command separates lines
   Whose features some poets % within a stanza.
   will curse.

                             % One or more blank lines separate stanzas.

   For instead of making\\
   Them do \emph{all} line breaking, \\
   It allows them to put too many words on a line when they'd rather be
   forced to be terse.
\end{verse}

Mathematical formulas may also be displayed.  A
displayed formula
is
one-line long; multiline
formulas require special formatting instructions.
   \[  \ip{\Gamma}{\psi'} = x'' + y^{2} + z_{i}^{n}\]
Don't start a paragraph with a displayed equation,
nor make one a paragraph by itself.

\end{document}               % End of document.

\section{Bench Huge Section 284}
% synthetic bench section 284
% This is a sample LaTeX input file.  (Version of 12 August 2004.)
%
% A '%' character causes TeX to ignore all remaining text on the line,
% and is used for comments like this one.

\documentclass{article}      % Specifies the document class

                             % The preamble begins here.
\title{An Example Document}  % Declares the document's title.
\author{Leslie Lamport}      % Declares the author's name.
\date{January 21, 1994}      % Deleting this command produces today's date.

\newcommand{\ip}[2]{(#1, #2)}
                             % Defines \ip{arg1}{arg2} to mean
                             % (arg1, arg2).

%\newcommand{\ip}[2]{\langle #1 | #2\rangle}
                             % This is an alternative definition of
                             % \ip that is commented out.

\begin{document}             % End of preamble and beginning of text.

\maketitle                   % Produces the title.

This is an example input file.  Comparing it with
the output it generates can show you how to
produce a simple document of your own.

\section{Ordinary Text}      % Produces section heading.  Lower-level
                             % sections are begun with similar
                             % \subsection and \subsubsection commands.

The ends  of words and sentences are marked
  by   spaces. It  doesn't matter how many
spaces    you type; one is as good as 100.  The
end of   a line counts as a space.

One   or more   blank lines denote the  end
of  a paragraph.

Since any number of consecutive spaces are treated
like a single one, the formatting of the input
file makes no difference to
      \LaTeX,                % The \LaTeX command generates the LaTeX logo.
but it makes a difference to you.  When you use
\LaTeX, making your input file as easy to read
as possible will be a great help as you write
your document and when you change it.  This sample
file shows how you can add comments to your own input
file.

Because printing is different from typewriting,
there are a number of things that you have to do
differently when preparing an input file than if
you were just typing the document directly.
Quotation marks like
       ``this''
have to be handled specially, as do quotes within
quotes:
       ``\,`this'            % \, separates the double and single quote.
        is what I just
        wrote, not  `that'\,''.

Dashes come in three sizes: an
       intra-word
dash, a medium dash for number ranges like
       1--2,
and a punctuation
       dash---like
this.

A sentence-ending space should be larger than the
space between words within a sentence.  You
sometimes have to type special commands in
conjunction with punctuation characters to get
this right, as in the following sentence.
       Gnats, gnus, etc.\ all  % `\ ' makes an inter-word space.
       begin with G\@.         % \@ marks end-of-sentence punctuation.
You should check the spaces after periods when
reading your output to make sure you haven't
forgotten any special cases.  Generating an
ellipsis
       \ldots\               % `\ ' is needed after `\ldots' because TeX
                             % ignores spaces after command names like \ldots
                             % made from \ + letters.
                             %
                             % Note how a `%' character causes TeX to ignore
                             % the end of the input line, so these blank lines
                             % do not start a new paragraph.
                             %
with the right spacing around the periods requires
a special command.

\LaTeX\ interprets some common characters as
commands, so you must type special commands to
generate them.  These characters include the
following:
       \$ \& \% \# \{ and \}.

In printing, text is usually emphasized with an
       \emph{italic}
type style.

\begin{em}
   A long segment of text can also be emphasized
   in this way.  Text within such a segment can be
   given \emph{additional} emphasis.
\end{em}

It is sometimes necessary to prevent \LaTeX\ from
breaking a line where it might otherwise do so.
This may be at a space, as between the ``Mr.''\ and
``Jones'' in
       ``Mr.~Jones'',        % ~ produces an unbreakable interword space.
or within a word---especially when the word is a
symbol like
       \mbox{\emph{itemnum}}
that makes little sense when hyphenated across
lines.

Footnotes\footnote{This is an example of a footnote.}
pose no problem.

\LaTeX\ is good at typesetting mathematical formulas
like
       \( x-3y + z = 7 \)
or
       \( a_{1} > x^{2n} + y^{2n} > x' \)
or
       \( \ip{A}{B} = \sum_{i} a_{i} b_{i} \).
The spaces you type in a formula are
ignored.  Remember that a letter like
       $x$                   % $ ... $  and  \( ... \)  are equivalent
is a formula when it denotes a mathematical
symbol, and it should be typed as one.

\section{Displayed Text}

Text is displayed by indenting it from the left
margin.  Quotations are commonly displayed.  There
are short quotations
\begin{quote}
   This is a short quotation.  It consists of a
   single paragraph of text.  See how it is formatted.
\end{quote}
and longer ones.
\begin{quotation}
   This is a longer quotation.  It consists of two
   paragraphs of text, neither of which are
   particularly interesting.

   This is the second paragraph of the quotation.  It
   is just as dull as the first paragraph.
\end{quotation}
Another frequently-displayed structure is a list.
The following is an example of an \emph{itemized}
list.
\begin{itemize}
   \item This is the first item of an itemized list.
         Each item in the list is marked with a ``tick''.
         You don't have to worry about what kind of tick
         mark is used.

   \item This is the second item of the list.  It
         contains another list nested inside it.  The inner
         list is an \emph{enumerated} list.
         \begin{enumerate}
            \item This is the first item of an enumerated
                  list that is nested within the itemized list.

            \item This is the second item of the inner list.
                  \LaTeX\ allows you to nest lists deeper than
                  you really should.
         \end{enumerate}
         This is the rest of the second item of the outer
         list.  It is no more interesting than any other
         part of the item.
   \item This is the third item of the list.
\end{itemize}
You can even display poetry.
\begin{verse}
   There is an environment
    for verse \\             % The \\ command separates lines
   Whose features some poets % within a stanza.
   will curse.

                             % One or more blank lines separate stanzas.

   For instead of making\\
   Them do \emph{all} line breaking, \\
   It allows them to put too many words on a line when they'd rather be
   forced to be terse.
\end{verse}

Mathematical formulas may also be displayed.  A
displayed formula
is
one-line long; multiline
formulas require special formatting instructions.
   \[  \ip{\Gamma}{\psi'} = x'' + y^{2} + z_{i}^{n}\]
Don't start a paragraph with a displayed equation,
nor make one a paragraph by itself.

\end{document}               % End of document.

\section{Bench Huge Section 285}
% synthetic bench section 285
% This is a sample LaTeX input file.  (Version of 12 August 2004.)
%
% A '%' character causes TeX to ignore all remaining text on the line,
% and is used for comments like this one.

\documentclass{article}      % Specifies the document class

                             % The preamble begins here.
\title{An Example Document}  % Declares the document's title.
\author{Leslie Lamport}      % Declares the author's name.
\date{January 21, 1994}      % Deleting this command produces today's date.

\newcommand{\ip}[2]{(#1, #2)}
                             % Defines \ip{arg1}{arg2} to mean
                             % (arg1, arg2).

%\newcommand{\ip}[2]{\langle #1 | #2\rangle}
                             % This is an alternative definition of
                             % \ip that is commented out.

\begin{document}             % End of preamble and beginning of text.

\maketitle                   % Produces the title.

This is an example input file.  Comparing it with
the output it generates can show you how to
produce a simple document of your own.

\section{Ordinary Text}      % Produces section heading.  Lower-level
                             % sections are begun with similar
                             % \subsection and \subsubsection commands.

The ends  of words and sentences are marked
  by   spaces. It  doesn't matter how many
spaces    you type; one is as good as 100.  The
end of   a line counts as a space.

One   or more   blank lines denote the  end
of  a paragraph.

Since any number of consecutive spaces are treated
like a single one, the formatting of the input
file makes no difference to
      \LaTeX,                % The \LaTeX command generates the LaTeX logo.
but it makes a difference to you.  When you use
\LaTeX, making your input file as easy to read
as possible will be a great help as you write
your document and when you change it.  This sample
file shows how you can add comments to your own input
file.

Because printing is different from typewriting,
there are a number of things that you have to do
differently when preparing an input file than if
you were just typing the document directly.
Quotation marks like
       ``this''
have to be handled specially, as do quotes within
quotes:
       ``\,`this'            % \, separates the double and single quote.
        is what I just
        wrote, not  `that'\,''.

Dashes come in three sizes: an
       intra-word
dash, a medium dash for number ranges like
       1--2,
and a punctuation
       dash---like
this.

A sentence-ending space should be larger than the
space between words within a sentence.  You
sometimes have to type special commands in
conjunction with punctuation characters to get
this right, as in the following sentence.
       Gnats, gnus, etc.\ all  % `\ ' makes an inter-word space.
       begin with G\@.         % \@ marks end-of-sentence punctuation.
You should check the spaces after periods when
reading your output to make sure you haven't
forgotten any special cases.  Generating an
ellipsis
       \ldots\               % `\ ' is needed after `\ldots' because TeX
                             % ignores spaces after command names like \ldots
                             % made from \ + letters.
                             %
                             % Note how a `%' character causes TeX to ignore
                             % the end of the input line, so these blank lines
                             % do not start a new paragraph.
                             %
with the right spacing around the periods requires
a special command.

\LaTeX\ interprets some common characters as
commands, so you must type special commands to
generate them.  These characters include the
following:
       \$ \& \% \# \{ and \}.

In printing, text is usually emphasized with an
       \emph{italic}
type style.

\begin{em}
   A long segment of text can also be emphasized
   in this way.  Text within such a segment can be
   given \emph{additional} emphasis.
\end{em}

It is sometimes necessary to prevent \LaTeX\ from
breaking a line where it might otherwise do so.
This may be at a space, as between the ``Mr.''\ and
``Jones'' in
       ``Mr.~Jones'',        % ~ produces an unbreakable interword space.
or within a word---especially when the word is a
symbol like
       \mbox{\emph{itemnum}}
that makes little sense when hyphenated across
lines.

Footnotes\footnote{This is an example of a footnote.}
pose no problem.

\LaTeX\ is good at typesetting mathematical formulas
like
       \( x-3y + z = 7 \)
or
       \( a_{1} > x^{2n} + y^{2n} > x' \)
or
       \( \ip{A}{B} = \sum_{i} a_{i} b_{i} \).
The spaces you type in a formula are
ignored.  Remember that a letter like
       $x$                   % $ ... $  and  \( ... \)  are equivalent
is a formula when it denotes a mathematical
symbol, and it should be typed as one.

\section{Displayed Text}

Text is displayed by indenting it from the left
margin.  Quotations are commonly displayed.  There
are short quotations
\begin{quote}
   This is a short quotation.  It consists of a
   single paragraph of text.  See how it is formatted.
\end{quote}
and longer ones.
\begin{quotation}
   This is a longer quotation.  It consists of two
   paragraphs of text, neither of which are
   particularly interesting.

   This is the second paragraph of the quotation.  It
   is just as dull as the first paragraph.
\end{quotation}
Another frequently-displayed structure is a list.
The following is an example of an \emph{itemized}
list.
\begin{itemize}
   \item This is the first item of an itemized list.
         Each item in the list is marked with a ``tick''.
         You don't have to worry about what kind of tick
         mark is used.

   \item This is the second item of the list.  It
         contains another list nested inside it.  The inner
         list is an \emph{enumerated} list.
         \begin{enumerate}
            \item This is the first item of an enumerated
                  list that is nested within the itemized list.

            \item This is the second item of the inner list.
                  \LaTeX\ allows you to nest lists deeper than
                  you really should.
         \end{enumerate}
         This is the rest of the second item of the outer
         list.  It is no more interesting than any other
         part of the item.
   \item This is the third item of the list.
\end{itemize}
You can even display poetry.
\begin{verse}
   There is an environment
    for verse \\             % The \\ command separates lines
   Whose features some poets % within a stanza.
   will curse.

                             % One or more blank lines separate stanzas.

   For instead of making\\
   Them do \emph{all} line breaking, \\
   It allows them to put too many words on a line when they'd rather be
   forced to be terse.
\end{verse}

Mathematical formulas may also be displayed.  A
displayed formula
is
one-line long; multiline
formulas require special formatting instructions.
   \[  \ip{\Gamma}{\psi'} = x'' + y^{2} + z_{i}^{n}\]
Don't start a paragraph with a displayed equation,
nor make one a paragraph by itself.

\end{document}               % End of document.

\section{Bench Huge Section 286}
% synthetic bench section 286
% This is a sample LaTeX input file.  (Version of 12 August 2004.)
%
% A '%' character causes TeX to ignore all remaining text on the line,
% and is used for comments like this one.

\documentclass{article}      % Specifies the document class

                             % The preamble begins here.
\title{An Example Document}  % Declares the document's title.
\author{Leslie Lamport}      % Declares the author's name.
\date{January 21, 1994}      % Deleting this command produces today's date.

\newcommand{\ip}[2]{(#1, #2)}
                             % Defines \ip{arg1}{arg2} to mean
                             % (arg1, arg2).

%\newcommand{\ip}[2]{\langle #1 | #2\rangle}
                             % This is an alternative definition of
                             % \ip that is commented out.

\begin{document}             % End of preamble and beginning of text.

\maketitle                   % Produces the title.

This is an example input file.  Comparing it with
the output it generates can show you how to
produce a simple document of your own.

\section{Ordinary Text}      % Produces section heading.  Lower-level
                             % sections are begun with similar
                             % \subsection and \subsubsection commands.

The ends  of words and sentences are marked
  by   spaces. It  doesn't matter how many
spaces    you type; one is as good as 100.  The
end of   a line counts as a space.

One   or more   blank lines denote the  end
of  a paragraph.

Since any number of consecutive spaces are treated
like a single one, the formatting of the input
file makes no difference to
      \LaTeX,                % The \LaTeX command generates the LaTeX logo.
but it makes a difference to you.  When you use
\LaTeX, making your input file as easy to read
as possible will be a great help as you write
your document and when you change it.  This sample
file shows how you can add comments to your own input
file.

Because printing is different from typewriting,
there are a number of things that you have to do
differently when preparing an input file than if
you were just typing the document directly.
Quotation marks like
       ``this''
have to be handled specially, as do quotes within
quotes:
       ``\,`this'            % \, separates the double and single quote.
        is what I just
        wrote, not  `that'\,''.

Dashes come in three sizes: an
       intra-word
dash, a medium dash for number ranges like
       1--2,
and a punctuation
       dash---like
this.

A sentence-ending space should be larger than the
space between words within a sentence.  You
sometimes have to type special commands in
conjunction with punctuation characters to get
this right, as in the following sentence.
       Gnats, gnus, etc.\ all  % `\ ' makes an inter-word space.
       begin with G\@.         % \@ marks end-of-sentence punctuation.
You should check the spaces after periods when
reading your output to make sure you haven't
forgotten any special cases.  Generating an
ellipsis
       \ldots\               % `\ ' is needed after `\ldots' because TeX
                             % ignores spaces after command names like \ldots
                             % made from \ + letters.
                             %
                             % Note how a `%' character causes TeX to ignore
                             % the end of the input line, so these blank lines
                             % do not start a new paragraph.
                             %
with the right spacing around the periods requires
a special command.

\LaTeX\ interprets some common characters as
commands, so you must type special commands to
generate them.  These characters include the
following:
       \$ \& \% \# \{ and \}.

In printing, text is usually emphasized with an
       \emph{italic}
type style.

\begin{em}
   A long segment of text can also be emphasized
   in this way.  Text within such a segment can be
   given \emph{additional} emphasis.
\end{em}

It is sometimes necessary to prevent \LaTeX\ from
breaking a line where it might otherwise do so.
This may be at a space, as between the ``Mr.''\ and
``Jones'' in
       ``Mr.~Jones'',        % ~ produces an unbreakable interword space.
or within a word---especially when the word is a
symbol like
       \mbox{\emph{itemnum}}
that makes little sense when hyphenated across
lines.

Footnotes\footnote{This is an example of a footnote.}
pose no problem.

\LaTeX\ is good at typesetting mathematical formulas
like
       \( x-3y + z = 7 \)
or
       \( a_{1} > x^{2n} + y^{2n} > x' \)
or
       \( \ip{A}{B} = \sum_{i} a_{i} b_{i} \).
The spaces you type in a formula are
ignored.  Remember that a letter like
       $x$                   % $ ... $  and  \( ... \)  are equivalent
is a formula when it denotes a mathematical
symbol, and it should be typed as one.

\section{Displayed Text}

Text is displayed by indenting it from the left
margin.  Quotations are commonly displayed.  There
are short quotations
\begin{quote}
   This is a short quotation.  It consists of a
   single paragraph of text.  See how it is formatted.
\end{quote}
and longer ones.
\begin{quotation}
   This is a longer quotation.  It consists of two
   paragraphs of text, neither of which are
   particularly interesting.

   This is the second paragraph of the quotation.  It
   is just as dull as the first paragraph.
\end{quotation}
Another frequently-displayed structure is a list.
The following is an example of an \emph{itemized}
list.
\begin{itemize}
   \item This is the first item of an itemized list.
         Each item in the list is marked with a ``tick''.
         You don't have to worry about what kind of tick
         mark is used.

   \item This is the second item of the list.  It
         contains another list nested inside it.  The inner
         list is an \emph{enumerated} list.
         \begin{enumerate}
            \item This is the first item of an enumerated
                  list that is nested within the itemized list.

            \item This is the second item of the inner list.
                  \LaTeX\ allows you to nest lists deeper than
                  you really should.
         \end{enumerate}
         This is the rest of the second item of the outer
         list.  It is no more interesting than any other
         part of the item.
   \item This is the third item of the list.
\end{itemize}
You can even display poetry.
\begin{verse}
   There is an environment
    for verse \\             % The \\ command separates lines
   Whose features some poets % within a stanza.
   will curse.

                             % One or more blank lines separate stanzas.

   For instead of making\\
   Them do \emph{all} line breaking, \\
   It allows them to put too many words on a line when they'd rather be
   forced to be terse.
\end{verse}

Mathematical formulas may also be displayed.  A
displayed formula
is
one-line long; multiline
formulas require special formatting instructions.
   \[  \ip{\Gamma}{\psi'} = x'' + y^{2} + z_{i}^{n}\]
Don't start a paragraph with a displayed equation,
nor make one a paragraph by itself.

\end{document}               % End of document.

\section{Bench Huge Section 287}
% synthetic bench section 287
% This is a sample LaTeX input file.  (Version of 12 August 2004.)
%
% A '%' character causes TeX to ignore all remaining text on the line,
% and is used for comments like this one.

\documentclass{article}      % Specifies the document class

                             % The preamble begins here.
\title{An Example Document}  % Declares the document's title.
\author{Leslie Lamport}      % Declares the author's name.
\date{January 21, 1994}      % Deleting this command produces today's date.

\newcommand{\ip}[2]{(#1, #2)}
                             % Defines \ip{arg1}{arg2} to mean
                             % (arg1, arg2).

%\newcommand{\ip}[2]{\langle #1 | #2\rangle}
                             % This is an alternative definition of
                             % \ip that is commented out.

\begin{document}             % End of preamble and beginning of text.

\maketitle                   % Produces the title.

This is an example input file.  Comparing it with
the output it generates can show you how to
produce a simple document of your own.

\section{Ordinary Text}      % Produces section heading.  Lower-level
                             % sections are begun with similar
                             % \subsection and \subsubsection commands.

The ends  of words and sentences are marked
  by   spaces. It  doesn't matter how many
spaces    you type; one is as good as 100.  The
end of   a line counts as a space.

One   or more   blank lines denote the  end
of  a paragraph.

Since any number of consecutive spaces are treated
like a single one, the formatting of the input
file makes no difference to
      \LaTeX,                % The \LaTeX command generates the LaTeX logo.
but it makes a difference to you.  When you use
\LaTeX, making your input file as easy to read
as possible will be a great help as you write
your document and when you change it.  This sample
file shows how you can add comments to your own input
file.

Because printing is different from typewriting,
there are a number of things that you have to do
differently when preparing an input file than if
you were just typing the document directly.
Quotation marks like
       ``this''
have to be handled specially, as do quotes within
quotes:
       ``\,`this'            % \, separates the double and single quote.
        is what I just
        wrote, not  `that'\,''.

Dashes come in three sizes: an
       intra-word
dash, a medium dash for number ranges like
       1--2,
and a punctuation
       dash---like
this.

A sentence-ending space should be larger than the
space between words within a sentence.  You
sometimes have to type special commands in
conjunction with punctuation characters to get
this right, as in the following sentence.
       Gnats, gnus, etc.\ all  % `\ ' makes an inter-word space.
       begin with G\@.         % \@ marks end-of-sentence punctuation.
You should check the spaces after periods when
reading your output to make sure you haven't
forgotten any special cases.  Generating an
ellipsis
       \ldots\               % `\ ' is needed after `\ldots' because TeX
                             % ignores spaces after command names like \ldots
                             % made from \ + letters.
                             %
                             % Note how a `%' character causes TeX to ignore
                             % the end of the input line, so these blank lines
                             % do not start a new paragraph.
                             %
with the right spacing around the periods requires
a special command.

\LaTeX\ interprets some common characters as
commands, so you must type special commands to
generate them.  These characters include the
following:
       \$ \& \% \# \{ and \}.

In printing, text is usually emphasized with an
       \emph{italic}
type style.

\begin{em}
   A long segment of text can also be emphasized
   in this way.  Text within such a segment can be
   given \emph{additional} emphasis.
\end{em}

It is sometimes necessary to prevent \LaTeX\ from
breaking a line where it might otherwise do so.
This may be at a space, as between the ``Mr.''\ and
``Jones'' in
       ``Mr.~Jones'',        % ~ produces an unbreakable interword space.
or within a word---especially when the word is a
symbol like
       \mbox{\emph{itemnum}}
that makes little sense when hyphenated across
lines.

Footnotes\footnote{This is an example of a footnote.}
pose no problem.

\LaTeX\ is good at typesetting mathematical formulas
like
       \( x-3y + z = 7 \)
or
       \( a_{1} > x^{2n} + y^{2n} > x' \)
or
       \( \ip{A}{B} = \sum_{i} a_{i} b_{i} \).
The spaces you type in a formula are
ignored.  Remember that a letter like
       $x$                   % $ ... $  and  \( ... \)  are equivalent
is a formula when it denotes a mathematical
symbol, and it should be typed as one.

\section{Displayed Text}

Text is displayed by indenting it from the left
margin.  Quotations are commonly displayed.  There
are short quotations
\begin{quote}
   This is a short quotation.  It consists of a
   single paragraph of text.  See how it is formatted.
\end{quote}
and longer ones.
\begin{quotation}
   This is a longer quotation.  It consists of two
   paragraphs of text, neither of which are
   particularly interesting.

   This is the second paragraph of the quotation.  It
   is just as dull as the first paragraph.
\end{quotation}
Another frequently-displayed structure is a list.
The following is an example of an \emph{itemized}
list.
\begin{itemize}
   \item This is the first item of an itemized list.
         Each item in the list is marked with a ``tick''.
         You don't have to worry about what kind of tick
         mark is used.

   \item This is the second item of the list.  It
         contains another list nested inside it.  The inner
         list is an \emph{enumerated} list.
         \begin{enumerate}
            \item This is the first item of an enumerated
                  list that is nested within the itemized list.

            \item This is the second item of the inner list.
                  \LaTeX\ allows you to nest lists deeper than
                  you really should.
         \end{enumerate}
         This is the rest of the second item of the outer
         list.  It is no more interesting than any other
         part of the item.
   \item This is the third item of the list.
\end{itemize}
You can even display poetry.
\begin{verse}
   There is an environment
    for verse \\             % The \\ command separates lines
   Whose features some poets % within a stanza.
   will curse.

                             % One or more blank lines separate stanzas.

   For instead of making\\
   Them do \emph{all} line breaking, \\
   It allows them to put too many words on a line when they'd rather be
   forced to be terse.
\end{verse}

Mathematical formulas may also be displayed.  A
displayed formula
is
one-line long; multiline
formulas require special formatting instructions.
   \[  \ip{\Gamma}{\psi'} = x'' + y^{2} + z_{i}^{n}\]
Don't start a paragraph with a displayed equation,
nor make one a paragraph by itself.

\end{document}               % End of document.

\section{Bench Huge Section 288}
% synthetic bench section 288
% This is a sample LaTeX input file.  (Version of 12 August 2004.)
%
% A '%' character causes TeX to ignore all remaining text on the line,
% and is used for comments like this one.

\documentclass{article}      % Specifies the document class

                             % The preamble begins here.
\title{An Example Document}  % Declares the document's title.
\author{Leslie Lamport}      % Declares the author's name.
\date{January 21, 1994}      % Deleting this command produces today's date.

\newcommand{\ip}[2]{(#1, #2)}
                             % Defines \ip{arg1}{arg2} to mean
                             % (arg1, arg2).

%\newcommand{\ip}[2]{\langle #1 | #2\rangle}
                             % This is an alternative definition of
                             % \ip that is commented out.

\begin{document}             % End of preamble and beginning of text.

\maketitle                   % Produces the title.

This is an example input file.  Comparing it with
the output it generates can show you how to
produce a simple document of your own.

\section{Ordinary Text}      % Produces section heading.  Lower-level
                             % sections are begun with similar
                             % \subsection and \subsubsection commands.

The ends  of words and sentences are marked
  by   spaces. It  doesn't matter how many
spaces    you type; one is as good as 100.  The
end of   a line counts as a space.

One   or more   blank lines denote the  end
of  a paragraph.

Since any number of consecutive spaces are treated
like a single one, the formatting of the input
file makes no difference to
      \LaTeX,                % The \LaTeX command generates the LaTeX logo.
but it makes a difference to you.  When you use
\LaTeX, making your input file as easy to read
as possible will be a great help as you write
your document and when you change it.  This sample
file shows how you can add comments to your own input
file.

Because printing is different from typewriting,
there are a number of things that you have to do
differently when preparing an input file than if
you were just typing the document directly.
Quotation marks like
       ``this''
have to be handled specially, as do quotes within
quotes:
       ``\,`this'            % \, separates the double and single quote.
        is what I just
        wrote, not  `that'\,''.

Dashes come in three sizes: an
       intra-word
dash, a medium dash for number ranges like
       1--2,
and a punctuation
       dash---like
this.

A sentence-ending space should be larger than the
space between words within a sentence.  You
sometimes have to type special commands in
conjunction with punctuation characters to get
this right, as in the following sentence.
       Gnats, gnus, etc.\ all  % `\ ' makes an inter-word space.
       begin with G\@.         % \@ marks end-of-sentence punctuation.
You should check the spaces after periods when
reading your output to make sure you haven't
forgotten any special cases.  Generating an
ellipsis
       \ldots\               % `\ ' is needed after `\ldots' because TeX
                             % ignores spaces after command names like \ldots
                             % made from \ + letters.
                             %
                             % Note how a `%' character causes TeX to ignore
                             % the end of the input line, so these blank lines
                             % do not start a new paragraph.
                             %
with the right spacing around the periods requires
a special command.

\LaTeX\ interprets some common characters as
commands, so you must type special commands to
generate them.  These characters include the
following:
       \$ \& \% \# \{ and \}.

In printing, text is usually emphasized with an
       \emph{italic}
type style.

\begin{em}
   A long segment of text can also be emphasized
   in this way.  Text within such a segment can be
   given \emph{additional} emphasis.
\end{em}

It is sometimes necessary to prevent \LaTeX\ from
breaking a line where it might otherwise do so.
This may be at a space, as between the ``Mr.''\ and
``Jones'' in
       ``Mr.~Jones'',        % ~ produces an unbreakable interword space.
or within a word---especially when the word is a
symbol like
       \mbox{\emph{itemnum}}
that makes little sense when hyphenated across
lines.

Footnotes\footnote{This is an example of a footnote.}
pose no problem.

\LaTeX\ is good at typesetting mathematical formulas
like
       \( x-3y + z = 7 \)
or
       \( a_{1} > x^{2n} + y^{2n} > x' \)
or
       \( \ip{A}{B} = \sum_{i} a_{i} b_{i} \).
The spaces you type in a formula are
ignored.  Remember that a letter like
       $x$                   % $ ... $  and  \( ... \)  are equivalent
is a formula when it denotes a mathematical
symbol, and it should be typed as one.

\section{Displayed Text}

Text is displayed by indenting it from the left
margin.  Quotations are commonly displayed.  There
are short quotations
\begin{quote}
   This is a short quotation.  It consists of a
   single paragraph of text.  See how it is formatted.
\end{quote}
and longer ones.
\begin{quotation}
   This is a longer quotation.  It consists of two
   paragraphs of text, neither of which are
   particularly interesting.

   This is the second paragraph of the quotation.  It
   is just as dull as the first paragraph.
\end{quotation}
Another frequently-displayed structure is a list.
The following is an example of an \emph{itemized}
list.
\begin{itemize}
   \item This is the first item of an itemized list.
         Each item in the list is marked with a ``tick''.
         You don't have to worry about what kind of tick
         mark is used.

   \item This is the second item of the list.  It
         contains another list nested inside it.  The inner
         list is an \emph{enumerated} list.
         \begin{enumerate}
            \item This is the first item of an enumerated
                  list that is nested within the itemized list.

            \item This is the second item of the inner list.
                  \LaTeX\ allows you to nest lists deeper than
                  you really should.
         \end{enumerate}
         This is the rest of the second item of the outer
         list.  It is no more interesting than any other
         part of the item.
   \item This is the third item of the list.
\end{itemize}
You can even display poetry.
\begin{verse}
   There is an environment
    for verse \\             % The \\ command separates lines
   Whose features some poets % within a stanza.
   will curse.

                             % One or more blank lines separate stanzas.

   For instead of making\\
   Them do \emph{all} line breaking, \\
   It allows them to put too many words on a line when they'd rather be
   forced to be terse.
\end{verse}

Mathematical formulas may also be displayed.  A
displayed formula
is
one-line long; multiline
formulas require special formatting instructions.
   \[  \ip{\Gamma}{\psi'} = x'' + y^{2} + z_{i}^{n}\]
Don't start a paragraph with a displayed equation,
nor make one a paragraph by itself.

\end{document}               % End of document.

\section{Bench Huge Section 289}
% synthetic bench section 289
% This is a sample LaTeX input file.  (Version of 12 August 2004.)
%
% A '%' character causes TeX to ignore all remaining text on the line,
% and is used for comments like this one.

\documentclass{article}      % Specifies the document class

                             % The preamble begins here.
\title{An Example Document}  % Declares the document's title.
\author{Leslie Lamport}      % Declares the author's name.
\date{January 21, 1994}      % Deleting this command produces today's date.

\newcommand{\ip}[2]{(#1, #2)}
                             % Defines \ip{arg1}{arg2} to mean
                             % (arg1, arg2).

%\newcommand{\ip}[2]{\langle #1 | #2\rangle}
                             % This is an alternative definition of
                             % \ip that is commented out.

\begin{document}             % End of preamble and beginning of text.

\maketitle                   % Produces the title.

This is an example input file.  Comparing it with
the output it generates can show you how to
produce a simple document of your own.

\section{Ordinary Text}      % Produces section heading.  Lower-level
                             % sections are begun with similar
                             % \subsection and \subsubsection commands.

The ends  of words and sentences are marked
  by   spaces. It  doesn't matter how many
spaces    you type; one is as good as 100.  The
end of   a line counts as a space.

One   or more   blank lines denote the  end
of  a paragraph.

Since any number of consecutive spaces are treated
like a single one, the formatting of the input
file makes no difference to
      \LaTeX,                % The \LaTeX command generates the LaTeX logo.
but it makes a difference to you.  When you use
\LaTeX, making your input file as easy to read
as possible will be a great help as you write
your document and when you change it.  This sample
file shows how you can add comments to your own input
file.

Because printing is different from typewriting,
there are a number of things that you have to do
differently when preparing an input file than if
you were just typing the document directly.
Quotation marks like
       ``this''
have to be handled specially, as do quotes within
quotes:
       ``\,`this'            % \, separates the double and single quote.
        is what I just
        wrote, not  `that'\,''.

Dashes come in three sizes: an
       intra-word
dash, a medium dash for number ranges like
       1--2,
and a punctuation
       dash---like
this.

A sentence-ending space should be larger than the
space between words within a sentence.  You
sometimes have to type special commands in
conjunction with punctuation characters to get
this right, as in the following sentence.
       Gnats, gnus, etc.\ all  % `\ ' makes an inter-word space.
       begin with G\@.         % \@ marks end-of-sentence punctuation.
You should check the spaces after periods when
reading your output to make sure you haven't
forgotten any special cases.  Generating an
ellipsis
       \ldots\               % `\ ' is needed after `\ldots' because TeX
                             % ignores spaces after command names like \ldots
                             % made from \ + letters.
                             %
                             % Note how a `%' character causes TeX to ignore
                             % the end of the input line, so these blank lines
                             % do not start a new paragraph.
                             %
with the right spacing around the periods requires
a special command.

\LaTeX\ interprets some common characters as
commands, so you must type special commands to
generate them.  These characters include the
following:
       \$ \& \% \# \{ and \}.

In printing, text is usually emphasized with an
       \emph{italic}
type style.

\begin{em}
   A long segment of text can also be emphasized
   in this way.  Text within such a segment can be
   given \emph{additional} emphasis.
\end{em}

It is sometimes necessary to prevent \LaTeX\ from
breaking a line where it might otherwise do so.
This may be at a space, as between the ``Mr.''\ and
``Jones'' in
       ``Mr.~Jones'',        % ~ produces an unbreakable interword space.
or within a word---especially when the word is a
symbol like
       \mbox{\emph{itemnum}}
that makes little sense when hyphenated across
lines.

Footnotes\footnote{This is an example of a footnote.}
pose no problem.

\LaTeX\ is good at typesetting mathematical formulas
like
       \( x-3y + z = 7 \)
or
       \( a_{1} > x^{2n} + y^{2n} > x' \)
or
       \( \ip{A}{B} = \sum_{i} a_{i} b_{i} \).
The spaces you type in a formula are
ignored.  Remember that a letter like
       $x$                   % $ ... $  and  \( ... \)  are equivalent
is a formula when it denotes a mathematical
symbol, and it should be typed as one.

\section{Displayed Text}

Text is displayed by indenting it from the left
margin.  Quotations are commonly displayed.  There
are short quotations
\begin{quote}
   This is a short quotation.  It consists of a
   single paragraph of text.  See how it is formatted.
\end{quote}
and longer ones.
\begin{quotation}
   This is a longer quotation.  It consists of two
   paragraphs of text, neither of which are
   particularly interesting.

   This is the second paragraph of the quotation.  It
   is just as dull as the first paragraph.
\end{quotation}
Another frequently-displayed structure is a list.
The following is an example of an \emph{itemized}
list.
\begin{itemize}
   \item This is the first item of an itemized list.
         Each item in the list is marked with a ``tick''.
         You don't have to worry about what kind of tick
         mark is used.

   \item This is the second item of the list.  It
         contains another list nested inside it.  The inner
         list is an \emph{enumerated} list.
         \begin{enumerate}
            \item This is the first item of an enumerated
                  list that is nested within the itemized list.

            \item This is the second item of the inner list.
                  \LaTeX\ allows you to nest lists deeper than
                  you really should.
         \end{enumerate}
         This is the rest of the second item of the outer
         list.  It is no more interesting than any other
         part of the item.
   \item This is the third item of the list.
\end{itemize}
You can even display poetry.
\begin{verse}
   There is an environment
    for verse \\             % The \\ command separates lines
   Whose features some poets % within a stanza.
   will curse.

                             % One or more blank lines separate stanzas.

   For instead of making\\
   Them do \emph{all} line breaking, \\
   It allows them to put too many words on a line when they'd rather be
   forced to be terse.
\end{verse}

Mathematical formulas may also be displayed.  A
displayed formula
is
one-line long; multiline
formulas require special formatting instructions.
   \[  \ip{\Gamma}{\psi'} = x'' + y^{2} + z_{i}^{n}\]
Don't start a paragraph with a displayed equation,
nor make one a paragraph by itself.

\end{document}               % End of document.

\section{Bench Huge Section 290}
% synthetic bench section 290
% This is a sample LaTeX input file.  (Version of 12 August 2004.)
%
% A '%' character causes TeX to ignore all remaining text on the line,
% and is used for comments like this one.

\documentclass{article}      % Specifies the document class

                             % The preamble begins here.
\title{An Example Document}  % Declares the document's title.
\author{Leslie Lamport}      % Declares the author's name.
\date{January 21, 1994}      % Deleting this command produces today's date.

\newcommand{\ip}[2]{(#1, #2)}
                             % Defines \ip{arg1}{arg2} to mean
                             % (arg1, arg2).

%\newcommand{\ip}[2]{\langle #1 | #2\rangle}
                             % This is an alternative definition of
                             % \ip that is commented out.

\begin{document}             % End of preamble and beginning of text.

\maketitle                   % Produces the title.

This is an example input file.  Comparing it with
the output it generates can show you how to
produce a simple document of your own.

\section{Ordinary Text}      % Produces section heading.  Lower-level
                             % sections are begun with similar
                             % \subsection and \subsubsection commands.

The ends  of words and sentences are marked
  by   spaces. It  doesn't matter how many
spaces    you type; one is as good as 100.  The
end of   a line counts as a space.

One   or more   blank lines denote the  end
of  a paragraph.

Since any number of consecutive spaces are treated
like a single one, the formatting of the input
file makes no difference to
      \LaTeX,                % The \LaTeX command generates the LaTeX logo.
but it makes a difference to you.  When you use
\LaTeX, making your input file as easy to read
as possible will be a great help as you write
your document and when you change it.  This sample
file shows how you can add comments to your own input
file.

Because printing is different from typewriting,
there are a number of things that you have to do
differently when preparing an input file than if
you were just typing the document directly.
Quotation marks like
       ``this''
have to be handled specially, as do quotes within
quotes:
       ``\,`this'            % \, separates the double and single quote.
        is what I just
        wrote, not  `that'\,''.

Dashes come in three sizes: an
       intra-word
dash, a medium dash for number ranges like
       1--2,
and a punctuation
       dash---like
this.

A sentence-ending space should be larger than the
space between words within a sentence.  You
sometimes have to type special commands in
conjunction with punctuation characters to get
this right, as in the following sentence.
       Gnats, gnus, etc.\ all  % `\ ' makes an inter-word space.
       begin with G\@.         % \@ marks end-of-sentence punctuation.
You should check the spaces after periods when
reading your output to make sure you haven't
forgotten any special cases.  Generating an
ellipsis
       \ldots\               % `\ ' is needed after `\ldots' because TeX
                             % ignores spaces after command names like \ldots
                             % made from \ + letters.
                             %
                             % Note how a `%' character causes TeX to ignore
                             % the end of the input line, so these blank lines
                             % do not start a new paragraph.
                             %
with the right spacing around the periods requires
a special command.

\LaTeX\ interprets some common characters as
commands, so you must type special commands to
generate them.  These characters include the
following:
       \$ \& \% \# \{ and \}.

In printing, text is usually emphasized with an
       \emph{italic}
type style.

\begin{em}
   A long segment of text can also be emphasized
   in this way.  Text within such a segment can be
   given \emph{additional} emphasis.
\end{em}

It is sometimes necessary to prevent \LaTeX\ from
breaking a line where it might otherwise do so.
This may be at a space, as between the ``Mr.''\ and
``Jones'' in
       ``Mr.~Jones'',        % ~ produces an unbreakable interword space.
or within a word---especially when the word is a
symbol like
       \mbox{\emph{itemnum}}
that makes little sense when hyphenated across
lines.

Footnotes\footnote{This is an example of a footnote.}
pose no problem.

\LaTeX\ is good at typesetting mathematical formulas
like
       \( x-3y + z = 7 \)
or
       \( a_{1} > x^{2n} + y^{2n} > x' \)
or
       \( \ip{A}{B} = \sum_{i} a_{i} b_{i} \).
The spaces you type in a formula are
ignored.  Remember that a letter like
       $x$                   % $ ... $  and  \( ... \)  are equivalent
is a formula when it denotes a mathematical
symbol, and it should be typed as one.

\section{Displayed Text}

Text is displayed by indenting it from the left
margin.  Quotations are commonly displayed.  There
are short quotations
\begin{quote}
   This is a short quotation.  It consists of a
   single paragraph of text.  See how it is formatted.
\end{quote}
and longer ones.
\begin{quotation}
   This is a longer quotation.  It consists of two
   paragraphs of text, neither of which are
   particularly interesting.

   This is the second paragraph of the quotation.  It
   is just as dull as the first paragraph.
\end{quotation}
Another frequently-displayed structure is a list.
The following is an example of an \emph{itemized}
list.
\begin{itemize}
   \item This is the first item of an itemized list.
         Each item in the list is marked with a ``tick''.
         You don't have to worry about what kind of tick
         mark is used.

   \item This is the second item of the list.  It
         contains another list nested inside it.  The inner
         list is an \emph{enumerated} list.
         \begin{enumerate}
            \item This is the first item of an enumerated
                  list that is nested within the itemized list.

            \item This is the second item of the inner list.
                  \LaTeX\ allows you to nest lists deeper than
                  you really should.
         \end{enumerate}
         This is the rest of the second item of the outer
         list.  It is no more interesting than any other
         part of the item.
   \item This is the third item of the list.
\end{itemize}
You can even display poetry.
\begin{verse}
   There is an environment
    for verse \\             % The \\ command separates lines
   Whose features some poets % within a stanza.
   will curse.

                             % One or more blank lines separate stanzas.

   For instead of making\\
   Them do \emph{all} line breaking, \\
   It allows them to put too many words on a line when they'd rather be
   forced to be terse.
\end{verse}

Mathematical formulas may also be displayed.  A
displayed formula
is
one-line long; multiline
formulas require special formatting instructions.
   \[  \ip{\Gamma}{\psi'} = x'' + y^{2} + z_{i}^{n}\]
Don't start a paragraph with a displayed equation,
nor make one a paragraph by itself.

\end{document}               % End of document.

\section{Bench Huge Section 291}
% synthetic bench section 291
% This is a sample LaTeX input file.  (Version of 12 August 2004.)
%
% A '%' character causes TeX to ignore all remaining text on the line,
% and is used for comments like this one.

\documentclass{article}      % Specifies the document class

                             % The preamble begins here.
\title{An Example Document}  % Declares the document's title.
\author{Leslie Lamport}      % Declares the author's name.
\date{January 21, 1994}      % Deleting this command produces today's date.

\newcommand{\ip}[2]{(#1, #2)}
                             % Defines \ip{arg1}{arg2} to mean
                             % (arg1, arg2).

%\newcommand{\ip}[2]{\langle #1 | #2\rangle}
                             % This is an alternative definition of
                             % \ip that is commented out.

\begin{document}             % End of preamble and beginning of text.

\maketitle                   % Produces the title.

This is an example input file.  Comparing it with
the output it generates can show you how to
produce a simple document of your own.

\section{Ordinary Text}      % Produces section heading.  Lower-level
                             % sections are begun with similar
                             % \subsection and \subsubsection commands.

The ends  of words and sentences are marked
  by   spaces. It  doesn't matter how many
spaces    you type; one is as good as 100.  The
end of   a line counts as a space.

One   or more   blank lines denote the  end
of  a paragraph.

Since any number of consecutive spaces are treated
like a single one, the formatting of the input
file makes no difference to
      \LaTeX,                % The \LaTeX command generates the LaTeX logo.
but it makes a difference to you.  When you use
\LaTeX, making your input file as easy to read
as possible will be a great help as you write
your document and when you change it.  This sample
file shows how you can add comments to your own input
file.

Because printing is different from typewriting,
there are a number of things that you have to do
differently when preparing an input file than if
you were just typing the document directly.
Quotation marks like
       ``this''
have to be handled specially, as do quotes within
quotes:
       ``\,`this'            % \, separates the double and single quote.
        is what I just
        wrote, not  `that'\,''.

Dashes come in three sizes: an
       intra-word
dash, a medium dash for number ranges like
       1--2,
and a punctuation
       dash---like
this.

A sentence-ending space should be larger than the
space between words within a sentence.  You
sometimes have to type special commands in
conjunction with punctuation characters to get
this right, as in the following sentence.
       Gnats, gnus, etc.\ all  % `\ ' makes an inter-word space.
       begin with G\@.         % \@ marks end-of-sentence punctuation.
You should check the spaces after periods when
reading your output to make sure you haven't
forgotten any special cases.  Generating an
ellipsis
       \ldots\               % `\ ' is needed after `\ldots' because TeX
                             % ignores spaces after command names like \ldots
                             % made from \ + letters.
                             %
                             % Note how a `%' character causes TeX to ignore
                             % the end of the input line, so these blank lines
                             % do not start a new paragraph.
                             %
with the right spacing around the periods requires
a special command.

\LaTeX\ interprets some common characters as
commands, so you must type special commands to
generate them.  These characters include the
following:
       \$ \& \% \# \{ and \}.

In printing, text is usually emphasized with an
       \emph{italic}
type style.

\begin{em}
   A long segment of text can also be emphasized
   in this way.  Text within such a segment can be
   given \emph{additional} emphasis.
\end{em}

It is sometimes necessary to prevent \LaTeX\ from
breaking a line where it might otherwise do so.
This may be at a space, as between the ``Mr.''\ and
``Jones'' in
       ``Mr.~Jones'',        % ~ produces an unbreakable interword space.
or within a word---especially when the word is a
symbol like
       \mbox{\emph{itemnum}}
that makes little sense when hyphenated across
lines.

Footnotes\footnote{This is an example of a footnote.}
pose no problem.

\LaTeX\ is good at typesetting mathematical formulas
like
       \( x-3y + z = 7 \)
or
       \( a_{1} > x^{2n} + y^{2n} > x' \)
or
       \( \ip{A}{B} = \sum_{i} a_{i} b_{i} \).
The spaces you type in a formula are
ignored.  Remember that a letter like
       $x$                   % $ ... $  and  \( ... \)  are equivalent
is a formula when it denotes a mathematical
symbol, and it should be typed as one.

\section{Displayed Text}

Text is displayed by indenting it from the left
margin.  Quotations are commonly displayed.  There
are short quotations
\begin{quote}
   This is a short quotation.  It consists of a
   single paragraph of text.  See how it is formatted.
\end{quote}
and longer ones.
\begin{quotation}
   This is a longer quotation.  It consists of two
   paragraphs of text, neither of which are
   particularly interesting.

   This is the second paragraph of the quotation.  It
   is just as dull as the first paragraph.
\end{quotation}
Another frequently-displayed structure is a list.
The following is an example of an \emph{itemized}
list.
\begin{itemize}
   \item This is the first item of an itemized list.
         Each item in the list is marked with a ``tick''.
         You don't have to worry about what kind of tick
         mark is used.

   \item This is the second item of the list.  It
         contains another list nested inside it.  The inner
         list is an \emph{enumerated} list.
         \begin{enumerate}
            \item This is the first item of an enumerated
                  list that is nested within the itemized list.

            \item This is the second item of the inner list.
                  \LaTeX\ allows you to nest lists deeper than
                  you really should.
         \end{enumerate}
         This is the rest of the second item of the outer
         list.  It is no more interesting than any other
         part of the item.
   \item This is the third item of the list.
\end{itemize}
You can even display poetry.
\begin{verse}
   There is an environment
    for verse \\             % The \\ command separates lines
   Whose features some poets % within a stanza.
   will curse.

                             % One or more blank lines separate stanzas.

   For instead of making\\
   Them do \emph{all} line breaking, \\
   It allows them to put too many words on a line when they'd rather be
   forced to be terse.
\end{verse}

Mathematical formulas may also be displayed.  A
displayed formula
is
one-line long; multiline
formulas require special formatting instructions.
   \[  \ip{\Gamma}{\psi'} = x'' + y^{2} + z_{i}^{n}\]
Don't start a paragraph with a displayed equation,
nor make one a paragraph by itself.

\end{document}               % End of document.

\section{Bench Huge Section 292}
% synthetic bench section 292
% This is a sample LaTeX input file.  (Version of 12 August 2004.)
%
% A '%' character causes TeX to ignore all remaining text on the line,
% and is used for comments like this one.

\documentclass{article}      % Specifies the document class

                             % The preamble begins here.
\title{An Example Document}  % Declares the document's title.
\author{Leslie Lamport}      % Declares the author's name.
\date{January 21, 1994}      % Deleting this command produces today's date.

\newcommand{\ip}[2]{(#1, #2)}
                             % Defines \ip{arg1}{arg2} to mean
                             % (arg1, arg2).

%\newcommand{\ip}[2]{\langle #1 | #2\rangle}
                             % This is an alternative definition of
                             % \ip that is commented out.

\begin{document}             % End of preamble and beginning of text.

\maketitle                   % Produces the title.

This is an example input file.  Comparing it with
the output it generates can show you how to
produce a simple document of your own.

\section{Ordinary Text}      % Produces section heading.  Lower-level
                             % sections are begun with similar
                             % \subsection and \subsubsection commands.

The ends  of words and sentences are marked
  by   spaces. It  doesn't matter how many
spaces    you type; one is as good as 100.  The
end of   a line counts as a space.

One   or more   blank lines denote the  end
of  a paragraph.

Since any number of consecutive spaces are treated
like a single one, the formatting of the input
file makes no difference to
      \LaTeX,                % The \LaTeX command generates the LaTeX logo.
but it makes a difference to you.  When you use
\LaTeX, making your input file as easy to read
as possible will be a great help as you write
your document and when you change it.  This sample
file shows how you can add comments to your own input
file.

Because printing is different from typewriting,
there are a number of things that you have to do
differently when preparing an input file than if
you were just typing the document directly.
Quotation marks like
       ``this''
have to be handled specially, as do quotes within
quotes:
       ``\,`this'            % \, separates the double and single quote.
        is what I just
        wrote, not  `that'\,''.

Dashes come in three sizes: an
       intra-word
dash, a medium dash for number ranges like
       1--2,
and a punctuation
       dash---like
this.

A sentence-ending space should be larger than the
space between words within a sentence.  You
sometimes have to type special commands in
conjunction with punctuation characters to get
this right, as in the following sentence.
       Gnats, gnus, etc.\ all  % `\ ' makes an inter-word space.
       begin with G\@.         % \@ marks end-of-sentence punctuation.
You should check the spaces after periods when
reading your output to make sure you haven't
forgotten any special cases.  Generating an
ellipsis
       \ldots\               % `\ ' is needed after `\ldots' because TeX
                             % ignores spaces after command names like \ldots
                             % made from \ + letters.
                             %
                             % Note how a `%' character causes TeX to ignore
                             % the end of the input line, so these blank lines
                             % do not start a new paragraph.
                             %
with the right spacing around the periods requires
a special command.

\LaTeX\ interprets some common characters as
commands, so you must type special commands to
generate them.  These characters include the
following:
       \$ \& \% \# \{ and \}.

In printing, text is usually emphasized with an
       \emph{italic}
type style.

\begin{em}
   A long segment of text can also be emphasized
   in this way.  Text within such a segment can be
   given \emph{additional} emphasis.
\end{em}

It is sometimes necessary to prevent \LaTeX\ from
breaking a line where it might otherwise do so.
This may be at a space, as between the ``Mr.''\ and
``Jones'' in
       ``Mr.~Jones'',        % ~ produces an unbreakable interword space.
or within a word---especially when the word is a
symbol like
       \mbox{\emph{itemnum}}
that makes little sense when hyphenated across
lines.

Footnotes\footnote{This is an example of a footnote.}
pose no problem.

\LaTeX\ is good at typesetting mathematical formulas
like
       \( x-3y + z = 7 \)
or
       \( a_{1} > x^{2n} + y^{2n} > x' \)
or
       \( \ip{A}{B} = \sum_{i} a_{i} b_{i} \).
The spaces you type in a formula are
ignored.  Remember that a letter like
       $x$                   % $ ... $  and  \( ... \)  are equivalent
is a formula when it denotes a mathematical
symbol, and it should be typed as one.

\section{Displayed Text}

Text is displayed by indenting it from the left
margin.  Quotations are commonly displayed.  There
are short quotations
\begin{quote}
   This is a short quotation.  It consists of a
   single paragraph of text.  See how it is formatted.
\end{quote}
and longer ones.
\begin{quotation}
   This is a longer quotation.  It consists of two
   paragraphs of text, neither of which are
   particularly interesting.

   This is the second paragraph of the quotation.  It
   is just as dull as the first paragraph.
\end{quotation}
Another frequently-displayed structure is a list.
The following is an example of an \emph{itemized}
list.
\begin{itemize}
   \item This is the first item of an itemized list.
         Each item in the list is marked with a ``tick''.
         You don't have to worry about what kind of tick
         mark is used.

   \item This is the second item of the list.  It
         contains another list nested inside it.  The inner
         list is an \emph{enumerated} list.
         \begin{enumerate}
            \item This is the first item of an enumerated
                  list that is nested within the itemized list.

            \item This is the second item of the inner list.
                  \LaTeX\ allows you to nest lists deeper than
                  you really should.
         \end{enumerate}
         This is the rest of the second item of the outer
         list.  It is no more interesting than any other
         part of the item.
   \item This is the third item of the list.
\end{itemize}
You can even display poetry.
\begin{verse}
   There is an environment
    for verse \\             % The \\ command separates lines
   Whose features some poets % within a stanza.
   will curse.

                             % One or more blank lines separate stanzas.

   For instead of making\\
   Them do \emph{all} line breaking, \\
   It allows them to put too many words on a line when they'd rather be
   forced to be terse.
\end{verse}

Mathematical formulas may also be displayed.  A
displayed formula
is
one-line long; multiline
formulas require special formatting instructions.
   \[  \ip{\Gamma}{\psi'} = x'' + y^{2} + z_{i}^{n}\]
Don't start a paragraph with a displayed equation,
nor make one a paragraph by itself.

\end{document}               % End of document.

\section{Bench Huge Section 293}
% synthetic bench section 293
% This is a sample LaTeX input file.  (Version of 12 August 2004.)
%
% A '%' character causes TeX to ignore all remaining text on the line,
% and is used for comments like this one.

\documentclass{article}      % Specifies the document class

                             % The preamble begins here.
\title{An Example Document}  % Declares the document's title.
\author{Leslie Lamport}      % Declares the author's name.
\date{January 21, 1994}      % Deleting this command produces today's date.

\newcommand{\ip}[2]{(#1, #2)}
                             % Defines \ip{arg1}{arg2} to mean
                             % (arg1, arg2).

%\newcommand{\ip}[2]{\langle #1 | #2\rangle}
                             % This is an alternative definition of
                             % \ip that is commented out.

\begin{document}             % End of preamble and beginning of text.

\maketitle                   % Produces the title.

This is an example input file.  Comparing it with
the output it generates can show you how to
produce a simple document of your own.

\section{Ordinary Text}      % Produces section heading.  Lower-level
                             % sections are begun with similar
                             % \subsection and \subsubsection commands.

The ends  of words and sentences are marked
  by   spaces. It  doesn't matter how many
spaces    you type; one is as good as 100.  The
end of   a line counts as a space.

One   or more   blank lines denote the  end
of  a paragraph.

Since any number of consecutive spaces are treated
like a single one, the formatting of the input
file makes no difference to
      \LaTeX,                % The \LaTeX command generates the LaTeX logo.
but it makes a difference to you.  When you use
\LaTeX, making your input file as easy to read
as possible will be a great help as you write
your document and when you change it.  This sample
file shows how you can add comments to your own input
file.

Because printing is different from typewriting,
there are a number of things that you have to do
differently when preparing an input file than if
you were just typing the document directly.
Quotation marks like
       ``this''
have to be handled specially, as do quotes within
quotes:
       ``\,`this'            % \, separates the double and single quote.
        is what I just
        wrote, not  `that'\,''.

Dashes come in three sizes: an
       intra-word
dash, a medium dash for number ranges like
       1--2,
and a punctuation
       dash---like
this.

A sentence-ending space should be larger than the
space between words within a sentence.  You
sometimes have to type special commands in
conjunction with punctuation characters to get
this right, as in the following sentence.
       Gnats, gnus, etc.\ all  % `\ ' makes an inter-word space.
       begin with G\@.         % \@ marks end-of-sentence punctuation.
You should check the spaces after periods when
reading your output to make sure you haven't
forgotten any special cases.  Generating an
ellipsis
       \ldots\               % `\ ' is needed after `\ldots' because TeX
                             % ignores spaces after command names like \ldots
                             % made from \ + letters.
                             %
                             % Note how a `%' character causes TeX to ignore
                             % the end of the input line, so these blank lines
                             % do not start a new paragraph.
                             %
with the right spacing around the periods requires
a special command.

\LaTeX\ interprets some common characters as
commands, so you must type special commands to
generate them.  These characters include the
following:
       \$ \& \% \# \{ and \}.

In printing, text is usually emphasized with an
       \emph{italic}
type style.

\begin{em}
   A long segment of text can also be emphasized
   in this way.  Text within such a segment can be
   given \emph{additional} emphasis.
\end{em}

It is sometimes necessary to prevent \LaTeX\ from
breaking a line where it might otherwise do so.
This may be at a space, as between the ``Mr.''\ and
``Jones'' in
       ``Mr.~Jones'',        % ~ produces an unbreakable interword space.
or within a word---especially when the word is a
symbol like
       \mbox{\emph{itemnum}}
that makes little sense when hyphenated across
lines.

Footnotes\footnote{This is an example of a footnote.}
pose no problem.

\LaTeX\ is good at typesetting mathematical formulas
like
       \( x-3y + z = 7 \)
or
       \( a_{1} > x^{2n} + y^{2n} > x' \)
or
       \( \ip{A}{B} = \sum_{i} a_{i} b_{i} \).
The spaces you type in a formula are
ignored.  Remember that a letter like
       $x$                   % $ ... $  and  \( ... \)  are equivalent
is a formula when it denotes a mathematical
symbol, and it should be typed as one.

\section{Displayed Text}

Text is displayed by indenting it from the left
margin.  Quotations are commonly displayed.  There
are short quotations
\begin{quote}
   This is a short quotation.  It consists of a
   single paragraph of text.  See how it is formatted.
\end{quote}
and longer ones.
\begin{quotation}
   This is a longer quotation.  It consists of two
   paragraphs of text, neither of which are
   particularly interesting.

   This is the second paragraph of the quotation.  It
   is just as dull as the first paragraph.
\end{quotation}
Another frequently-displayed structure is a list.
The following is an example of an \emph{itemized}
list.
\begin{itemize}
   \item This is the first item of an itemized list.
         Each item in the list is marked with a ``tick''.
         You don't have to worry about what kind of tick
         mark is used.

   \item This is the second item of the list.  It
         contains another list nested inside it.  The inner
         list is an \emph{enumerated} list.
         \begin{enumerate}
            \item This is the first item of an enumerated
                  list that is nested within the itemized list.

            \item This is the second item of the inner list.
                  \LaTeX\ allows you to nest lists deeper than
                  you really should.
         \end{enumerate}
         This is the rest of the second item of the outer
         list.  It is no more interesting than any other
         part of the item.
   \item This is the third item of the list.
\end{itemize}
You can even display poetry.
\begin{verse}
   There is an environment
    for verse \\             % The \\ command separates lines
   Whose features some poets % within a stanza.
   will curse.

                             % One or more blank lines separate stanzas.

   For instead of making\\
   Them do \emph{all} line breaking, \\
   It allows them to put too many words on a line when they'd rather be
   forced to be terse.
\end{verse}

Mathematical formulas may also be displayed.  A
displayed formula
is
one-line long; multiline
formulas require special formatting instructions.
   \[  \ip{\Gamma}{\psi'} = x'' + y^{2} + z_{i}^{n}\]
Don't start a paragraph with a displayed equation,
nor make one a paragraph by itself.

\end{document}               % End of document.

\section{Bench Huge Section 294}
% synthetic bench section 294
% This is a sample LaTeX input file.  (Version of 12 August 2004.)
%
% A '%' character causes TeX to ignore all remaining text on the line,
% and is used for comments like this one.

\documentclass{article}      % Specifies the document class

                             % The preamble begins here.
\title{An Example Document}  % Declares the document's title.
\author{Leslie Lamport}      % Declares the author's name.
\date{January 21, 1994}      % Deleting this command produces today's date.

\newcommand{\ip}[2]{(#1, #2)}
                             % Defines \ip{arg1}{arg2} to mean
                             % (arg1, arg2).

%\newcommand{\ip}[2]{\langle #1 | #2\rangle}
                             % This is an alternative definition of
                             % \ip that is commented out.

\begin{document}             % End of preamble and beginning of text.

\maketitle                   % Produces the title.

This is an example input file.  Comparing it with
the output it generates can show you how to
produce a simple document of your own.

\section{Ordinary Text}      % Produces section heading.  Lower-level
                             % sections are begun with similar
                             % \subsection and \subsubsection commands.

The ends  of words and sentences are marked
  by   spaces. It  doesn't matter how many
spaces    you type; one is as good as 100.  The
end of   a line counts as a space.

One   or more   blank lines denote the  end
of  a paragraph.

Since any number of consecutive spaces are treated
like a single one, the formatting of the input
file makes no difference to
      \LaTeX,                % The \LaTeX command generates the LaTeX logo.
but it makes a difference to you.  When you use
\LaTeX, making your input file as easy to read
as possible will be a great help as you write
your document and when you change it.  This sample
file shows how you can add comments to your own input
file.

Because printing is different from typewriting,
there are a number of things that you have to do
differently when preparing an input file than if
you were just typing the document directly.
Quotation marks like
       ``this''
have to be handled specially, as do quotes within
quotes:
       ``\,`this'            % \, separates the double and single quote.
        is what I just
        wrote, not  `that'\,''.

Dashes come in three sizes: an
       intra-word
dash, a medium dash for number ranges like
       1--2,
and a punctuation
       dash---like
this.

A sentence-ending space should be larger than the
space between words within a sentence.  You
sometimes have to type special commands in
conjunction with punctuation characters to get
this right, as in the following sentence.
       Gnats, gnus, etc.\ all  % `\ ' makes an inter-word space.
       begin with G\@.         % \@ marks end-of-sentence punctuation.
You should check the spaces after periods when
reading your output to make sure you haven't
forgotten any special cases.  Generating an
ellipsis
       \ldots\               % `\ ' is needed after `\ldots' because TeX
                             % ignores spaces after command names like \ldots
                             % made from \ + letters.
                             %
                             % Note how a `%' character causes TeX to ignore
                             % the end of the input line, so these blank lines
                             % do not start a new paragraph.
                             %
with the right spacing around the periods requires
a special command.

\LaTeX\ interprets some common characters as
commands, so you must type special commands to
generate them.  These characters include the
following:
       \$ \& \% \# \{ and \}.

In printing, text is usually emphasized with an
       \emph{italic}
type style.

\begin{em}
   A long segment of text can also be emphasized
   in this way.  Text within such a segment can be
   given \emph{additional} emphasis.
\end{em}

It is sometimes necessary to prevent \LaTeX\ from
breaking a line where it might otherwise do so.
This may be at a space, as between the ``Mr.''\ and
``Jones'' in
       ``Mr.~Jones'',        % ~ produces an unbreakable interword space.
or within a word---especially when the word is a
symbol like
       \mbox{\emph{itemnum}}
that makes little sense when hyphenated across
lines.

Footnotes\footnote{This is an example of a footnote.}
pose no problem.

\LaTeX\ is good at typesetting mathematical formulas
like
       \( x-3y + z = 7 \)
or
       \( a_{1} > x^{2n} + y^{2n} > x' \)
or
       \( \ip{A}{B} = \sum_{i} a_{i} b_{i} \).
The spaces you type in a formula are
ignored.  Remember that a letter like
       $x$                   % $ ... $  and  \( ... \)  are equivalent
is a formula when it denotes a mathematical
symbol, and it should be typed as one.

\section{Displayed Text}

Text is displayed by indenting it from the left
margin.  Quotations are commonly displayed.  There
are short quotations
\begin{quote}
   This is a short quotation.  It consists of a
   single paragraph of text.  See how it is formatted.
\end{quote}
and longer ones.
\begin{quotation}
   This is a longer quotation.  It consists of two
   paragraphs of text, neither of which are
   particularly interesting.

   This is the second paragraph of the quotation.  It
   is just as dull as the first paragraph.
\end{quotation}
Another frequently-displayed structure is a list.
The following is an example of an \emph{itemized}
list.
\begin{itemize}
   \item This is the first item of an itemized list.
         Each item in the list is marked with a ``tick''.
         You don't have to worry about what kind of tick
         mark is used.

   \item This is the second item of the list.  It
         contains another list nested inside it.  The inner
         list is an \emph{enumerated} list.
         \begin{enumerate}
            \item This is the first item of an enumerated
                  list that is nested within the itemized list.

            \item This is the second item of the inner list.
                  \LaTeX\ allows you to nest lists deeper than
                  you really should.
         \end{enumerate}
         This is the rest of the second item of the outer
         list.  It is no more interesting than any other
         part of the item.
   \item This is the third item of the list.
\end{itemize}
You can even display poetry.
\begin{verse}
   There is an environment
    for verse \\             % The \\ command separates lines
   Whose features some poets % within a stanza.
   will curse.

                             % One or more blank lines separate stanzas.

   For instead of making\\
   Them do \emph{all} line breaking, \\
   It allows them to put too many words on a line when they'd rather be
   forced to be terse.
\end{verse}

Mathematical formulas may also be displayed.  A
displayed formula
is
one-line long; multiline
formulas require special formatting instructions.
   \[  \ip{\Gamma}{\psi'} = x'' + y^{2} + z_{i}^{n}\]
Don't start a paragraph with a displayed equation,
nor make one a paragraph by itself.

\end{document}               % End of document.

\section{Bench Huge Section 295}
% synthetic bench section 295
% This is a sample LaTeX input file.  (Version of 12 August 2004.)
%
% A '%' character causes TeX to ignore all remaining text on the line,
% and is used for comments like this one.

\documentclass{article}      % Specifies the document class

                             % The preamble begins here.
\title{An Example Document}  % Declares the document's title.
\author{Leslie Lamport}      % Declares the author's name.
\date{January 21, 1994}      % Deleting this command produces today's date.

\newcommand{\ip}[2]{(#1, #2)}
                             % Defines \ip{arg1}{arg2} to mean
                             % (arg1, arg2).

%\newcommand{\ip}[2]{\langle #1 | #2\rangle}
                             % This is an alternative definition of
                             % \ip that is commented out.

\begin{document}             % End of preamble and beginning of text.

\maketitle                   % Produces the title.

This is an example input file.  Comparing it with
the output it generates can show you how to
produce a simple document of your own.

\section{Ordinary Text}      % Produces section heading.  Lower-level
                             % sections are begun with similar
                             % \subsection and \subsubsection commands.

The ends  of words and sentences are marked
  by   spaces. It  doesn't matter how many
spaces    you type; one is as good as 100.  The
end of   a line counts as a space.

One   or more   blank lines denote the  end
of  a paragraph.

Since any number of consecutive spaces are treated
like a single one, the formatting of the input
file makes no difference to
      \LaTeX,                % The \LaTeX command generates the LaTeX logo.
but it makes a difference to you.  When you use
\LaTeX, making your input file as easy to read
as possible will be a great help as you write
your document and when you change it.  This sample
file shows how you can add comments to your own input
file.

Because printing is different from typewriting,
there are a number of things that you have to do
differently when preparing an input file than if
you were just typing the document directly.
Quotation marks like
       ``this''
have to be handled specially, as do quotes within
quotes:
       ``\,`this'            % \, separates the double and single quote.
        is what I just
        wrote, not  `that'\,''.

Dashes come in three sizes: an
       intra-word
dash, a medium dash for number ranges like
       1--2,
and a punctuation
       dash---like
this.

A sentence-ending space should be larger than the
space between words within a sentence.  You
sometimes have to type special commands in
conjunction with punctuation characters to get
this right, as in the following sentence.
       Gnats, gnus, etc.\ all  % `\ ' makes an inter-word space.
       begin with G\@.         % \@ marks end-of-sentence punctuation.
You should check the spaces after periods when
reading your output to make sure you haven't
forgotten any special cases.  Generating an
ellipsis
       \ldots\               % `\ ' is needed after `\ldots' because TeX
                             % ignores spaces after command names like \ldots
                             % made from \ + letters.
                             %
                             % Note how a `%' character causes TeX to ignore
                             % the end of the input line, so these blank lines
                             % do not start a new paragraph.
                             %
with the right spacing around the periods requires
a special command.

\LaTeX\ interprets some common characters as
commands, so you must type special commands to
generate them.  These characters include the
following:
       \$ \& \% \# \{ and \}.

In printing, text is usually emphasized with an
       \emph{italic}
type style.

\begin{em}
   A long segment of text can also be emphasized
   in this way.  Text within such a segment can be
   given \emph{additional} emphasis.
\end{em}

It is sometimes necessary to prevent \LaTeX\ from
breaking a line where it might otherwise do so.
This may be at a space, as between the ``Mr.''\ and
``Jones'' in
       ``Mr.~Jones'',        % ~ produces an unbreakable interword space.
or within a word---especially when the word is a
symbol like
       \mbox{\emph{itemnum}}
that makes little sense when hyphenated across
lines.

Footnotes\footnote{This is an example of a footnote.}
pose no problem.

\LaTeX\ is good at typesetting mathematical formulas
like
       \( x-3y + z = 7 \)
or
       \( a_{1} > x^{2n} + y^{2n} > x' \)
or
       \( \ip{A}{B} = \sum_{i} a_{i} b_{i} \).
The spaces you type in a formula are
ignored.  Remember that a letter like
       $x$                   % $ ... $  and  \( ... \)  are equivalent
is a formula when it denotes a mathematical
symbol, and it should be typed as one.

\section{Displayed Text}

Text is displayed by indenting it from the left
margin.  Quotations are commonly displayed.  There
are short quotations
\begin{quote}
   This is a short quotation.  It consists of a
   single paragraph of text.  See how it is formatted.
\end{quote}
and longer ones.
\begin{quotation}
   This is a longer quotation.  It consists of two
   paragraphs of text, neither of which are
   particularly interesting.

   This is the second paragraph of the quotation.  It
   is just as dull as the first paragraph.
\end{quotation}
Another frequently-displayed structure is a list.
The following is an example of an \emph{itemized}
list.
\begin{itemize}
   \item This is the first item of an itemized list.
         Each item in the list is marked with a ``tick''.
         You don't have to worry about what kind of tick
         mark is used.

   \item This is the second item of the list.  It
         contains another list nested inside it.  The inner
         list is an \emph{enumerated} list.
         \begin{enumerate}
            \item This is the first item of an enumerated
                  list that is nested within the itemized list.

            \item This is the second item of the inner list.
                  \LaTeX\ allows you to nest lists deeper than
                  you really should.
         \end{enumerate}
         This is the rest of the second item of the outer
         list.  It is no more interesting than any other
         part of the item.
   \item This is the third item of the list.
\end{itemize}
You can even display poetry.
\begin{verse}
   There is an environment
    for verse \\             % The \\ command separates lines
   Whose features some poets % within a stanza.
   will curse.

                             % One or more blank lines separate stanzas.

   For instead of making\\
   Them do \emph{all} line breaking, \\
   It allows them to put too many words on a line when they'd rather be
   forced to be terse.
\end{verse}

Mathematical formulas may also be displayed.  A
displayed formula
is
one-line long; multiline
formulas require special formatting instructions.
   \[  \ip{\Gamma}{\psi'} = x'' + y^{2} + z_{i}^{n}\]
Don't start a paragraph with a displayed equation,
nor make one a paragraph by itself.

\end{document}               % End of document.

\section{Bench Huge Section 296}
% synthetic bench section 296
% This is a sample LaTeX input file.  (Version of 12 August 2004.)
%
% A '%' character causes TeX to ignore all remaining text on the line,
% and is used for comments like this one.

\documentclass{article}      % Specifies the document class

                             % The preamble begins here.
\title{An Example Document}  % Declares the document's title.
\author{Leslie Lamport}      % Declares the author's name.
\date{January 21, 1994}      % Deleting this command produces today's date.

\newcommand{\ip}[2]{(#1, #2)}
                             % Defines \ip{arg1}{arg2} to mean
                             % (arg1, arg2).

%\newcommand{\ip}[2]{\langle #1 | #2\rangle}
                             % This is an alternative definition of
                             % \ip that is commented out.

\begin{document}             % End of preamble and beginning of text.

\maketitle                   % Produces the title.

This is an example input file.  Comparing it with
the output it generates can show you how to
produce a simple document of your own.

\section{Ordinary Text}      % Produces section heading.  Lower-level
                             % sections are begun with similar
                             % \subsection and \subsubsection commands.

The ends  of words and sentences are marked
  by   spaces. It  doesn't matter how many
spaces    you type; one is as good as 100.  The
end of   a line counts as a space.

One   or more   blank lines denote the  end
of  a paragraph.

Since any number of consecutive spaces are treated
like a single one, the formatting of the input
file makes no difference to
      \LaTeX,                % The \LaTeX command generates the LaTeX logo.
but it makes a difference to you.  When you use
\LaTeX, making your input file as easy to read
as possible will be a great help as you write
your document and when you change it.  This sample
file shows how you can add comments to your own input
file.

Because printing is different from typewriting,
there are a number of things that you have to do
differently when preparing an input file than if
you were just typing the document directly.
Quotation marks like
       ``this''
have to be handled specially, as do quotes within
quotes:
       ``\,`this'            % \, separates the double and single quote.
        is what I just
        wrote, not  `that'\,''.

Dashes come in three sizes: an
       intra-word
dash, a medium dash for number ranges like
       1--2,
and a punctuation
       dash---like
this.

A sentence-ending space should be larger than the
space between words within a sentence.  You
sometimes have to type special commands in
conjunction with punctuation characters to get
this right, as in the following sentence.
       Gnats, gnus, etc.\ all  % `\ ' makes an inter-word space.
       begin with G\@.         % \@ marks end-of-sentence punctuation.
You should check the spaces after periods when
reading your output to make sure you haven't
forgotten any special cases.  Generating an
ellipsis
       \ldots\               % `\ ' is needed after `\ldots' because TeX
                             % ignores spaces after command names like \ldots
                             % made from \ + letters.
                             %
                             % Note how a `%' character causes TeX to ignore
                             % the end of the input line, so these blank lines
                             % do not start a new paragraph.
                             %
with the right spacing around the periods requires
a special command.

\LaTeX\ interprets some common characters as
commands, so you must type special commands to
generate them.  These characters include the
following:
       \$ \& \% \# \{ and \}.

In printing, text is usually emphasized with an
       \emph{italic}
type style.

\begin{em}
   A long segment of text can also be emphasized
   in this way.  Text within such a segment can be
   given \emph{additional} emphasis.
\end{em}

It is sometimes necessary to prevent \LaTeX\ from
breaking a line where it might otherwise do so.
This may be at a space, as between the ``Mr.''\ and
``Jones'' in
       ``Mr.~Jones'',        % ~ produces an unbreakable interword space.
or within a word---especially when the word is a
symbol like
       \mbox{\emph{itemnum}}
that makes little sense when hyphenated across
lines.

Footnotes\footnote{This is an example of a footnote.}
pose no problem.

\LaTeX\ is good at typesetting mathematical formulas
like
       \( x-3y + z = 7 \)
or
       \( a_{1} > x^{2n} + y^{2n} > x' \)
or
       \( \ip{A}{B} = \sum_{i} a_{i} b_{i} \).
The spaces you type in a formula are
ignored.  Remember that a letter like
       $x$                   % $ ... $  and  \( ... \)  are equivalent
is a formula when it denotes a mathematical
symbol, and it should be typed as one.

\section{Displayed Text}

Text is displayed by indenting it from the left
margin.  Quotations are commonly displayed.  There
are short quotations
\begin{quote}
   This is a short quotation.  It consists of a
   single paragraph of text.  See how it is formatted.
\end{quote}
and longer ones.
\begin{quotation}
   This is a longer quotation.  It consists of two
   paragraphs of text, neither of which are
   particularly interesting.

   This is the second paragraph of the quotation.  It
   is just as dull as the first paragraph.
\end{quotation}
Another frequently-displayed structure is a list.
The following is an example of an \emph{itemized}
list.
\begin{itemize}
   \item This is the first item of an itemized list.
         Each item in the list is marked with a ``tick''.
         You don't have to worry about what kind of tick
         mark is used.

   \item This is the second item of the list.  It
         contains another list nested inside it.  The inner
         list is an \emph{enumerated} list.
         \begin{enumerate}
            \item This is the first item of an enumerated
                  list that is nested within the itemized list.

            \item This is the second item of the inner list.
                  \LaTeX\ allows you to nest lists deeper than
                  you really should.
         \end{enumerate}
         This is the rest of the second item of the outer
         list.  It is no more interesting than any other
         part of the item.
   \item This is the third item of the list.
\end{itemize}
You can even display poetry.
\begin{verse}
   There is an environment
    for verse \\             % The \\ command separates lines
   Whose features some poets % within a stanza.
   will curse.

                             % One or more blank lines separate stanzas.

   For instead of making\\
   Them do \emph{all} line breaking, \\
   It allows them to put too many words on a line when they'd rather be
   forced to be terse.
\end{verse}

Mathematical formulas may also be displayed.  A
displayed formula
is
one-line long; multiline
formulas require special formatting instructions.
   \[  \ip{\Gamma}{\psi'} = x'' + y^{2} + z_{i}^{n}\]
Don't start a paragraph with a displayed equation,
nor make one a paragraph by itself.

\end{document}               % End of document.

\section{Bench Huge Section 297}
% synthetic bench section 297
% This is a sample LaTeX input file.  (Version of 12 August 2004.)
%
% A '%' character causes TeX to ignore all remaining text on the line,
% and is used for comments like this one.

\documentclass{article}      % Specifies the document class

                             % The preamble begins here.
\title{An Example Document}  % Declares the document's title.
\author{Leslie Lamport}      % Declares the author's name.
\date{January 21, 1994}      % Deleting this command produces today's date.

\newcommand{\ip}[2]{(#1, #2)}
                             % Defines \ip{arg1}{arg2} to mean
                             % (arg1, arg2).

%\newcommand{\ip}[2]{\langle #1 | #2\rangle}
                             % This is an alternative definition of
                             % \ip that is commented out.

\begin{document}             % End of preamble and beginning of text.

\maketitle                   % Produces the title.

This is an example input file.  Comparing it with
the output it generates can show you how to
produce a simple document of your own.

\section{Ordinary Text}      % Produces section heading.  Lower-level
                             % sections are begun with similar
                             % \subsection and \subsubsection commands.

The ends  of words and sentences are marked
  by   spaces. It  doesn't matter how many
spaces    you type; one is as good as 100.  The
end of   a line counts as a space.

One   or more   blank lines denote the  end
of  a paragraph.

Since any number of consecutive spaces are treated
like a single one, the formatting of the input
file makes no difference to
      \LaTeX,                % The \LaTeX command generates the LaTeX logo.
but it makes a difference to you.  When you use
\LaTeX, making your input file as easy to read
as possible will be a great help as you write
your document and when you change it.  This sample
file shows how you can add comments to your own input
file.

Because printing is different from typewriting,
there are a number of things that you have to do
differently when preparing an input file than if
you were just typing the document directly.
Quotation marks like
       ``this''
have to be handled specially, as do quotes within
quotes:
       ``\,`this'            % \, separates the double and single quote.
        is what I just
        wrote, not  `that'\,''.

Dashes come in three sizes: an
       intra-word
dash, a medium dash for number ranges like
       1--2,
and a punctuation
       dash---like
this.

A sentence-ending space should be larger than the
space between words within a sentence.  You
sometimes have to type special commands in
conjunction with punctuation characters to get
this right, as in the following sentence.
       Gnats, gnus, etc.\ all  % `\ ' makes an inter-word space.
       begin with G\@.         % \@ marks end-of-sentence punctuation.
You should check the spaces after periods when
reading your output to make sure you haven't
forgotten any special cases.  Generating an
ellipsis
       \ldots\               % `\ ' is needed after `\ldots' because TeX
                             % ignores spaces after command names like \ldots
                             % made from \ + letters.
                             %
                             % Note how a `%' character causes TeX to ignore
                             % the end of the input line, so these blank lines
                             % do not start a new paragraph.
                             %
with the right spacing around the periods requires
a special command.

\LaTeX\ interprets some common characters as
commands, so you must type special commands to
generate them.  These characters include the
following:
       \$ \& \% \# \{ and \}.

In printing, text is usually emphasized with an
       \emph{italic}
type style.

\begin{em}
   A long segment of text can also be emphasized
   in this way.  Text within such a segment can be
   given \emph{additional} emphasis.
\end{em}

It is sometimes necessary to prevent \LaTeX\ from
breaking a line where it might otherwise do so.
This may be at a space, as between the ``Mr.''\ and
``Jones'' in
       ``Mr.~Jones'',        % ~ produces an unbreakable interword space.
or within a word---especially when the word is a
symbol like
       \mbox{\emph{itemnum}}
that makes little sense when hyphenated across
lines.

Footnotes\footnote{This is an example of a footnote.}
pose no problem.

\LaTeX\ is good at typesetting mathematical formulas
like
       \( x-3y + z = 7 \)
or
       \( a_{1} > x^{2n} + y^{2n} > x' \)
or
       \( \ip{A}{B} = \sum_{i} a_{i} b_{i} \).
The spaces you type in a formula are
ignored.  Remember that a letter like
       $x$                   % $ ... $  and  \( ... \)  are equivalent
is a formula when it denotes a mathematical
symbol, and it should be typed as one.

\section{Displayed Text}

Text is displayed by indenting it from the left
margin.  Quotations are commonly displayed.  There
are short quotations
\begin{quote}
   This is a short quotation.  It consists of a
   single paragraph of text.  See how it is formatted.
\end{quote}
and longer ones.
\begin{quotation}
   This is a longer quotation.  It consists of two
   paragraphs of text, neither of which are
   particularly interesting.

   This is the second paragraph of the quotation.  It
   is just as dull as the first paragraph.
\end{quotation}
Another frequently-displayed structure is a list.
The following is an example of an \emph{itemized}
list.
\begin{itemize}
   \item This is the first item of an itemized list.
         Each item in the list is marked with a ``tick''.
         You don't have to worry about what kind of tick
         mark is used.

   \item This is the second item of the list.  It
         contains another list nested inside it.  The inner
         list is an \emph{enumerated} list.
         \begin{enumerate}
            \item This is the first item of an enumerated
                  list that is nested within the itemized list.

            \item This is the second item of the inner list.
                  \LaTeX\ allows you to nest lists deeper than
                  you really should.
         \end{enumerate}
         This is the rest of the second item of the outer
         list.  It is no more interesting than any other
         part of the item.
   \item This is the third item of the list.
\end{itemize}
You can even display poetry.
\begin{verse}
   There is an environment
    for verse \\             % The \\ command separates lines
   Whose features some poets % within a stanza.
   will curse.

                             % One or more blank lines separate stanzas.

   For instead of making\\
   Them do \emph{all} line breaking, \\
   It allows them to put too many words on a line when they'd rather be
   forced to be terse.
\end{verse}

Mathematical formulas may also be displayed.  A
displayed formula
is
one-line long; multiline
formulas require special formatting instructions.
   \[  \ip{\Gamma}{\psi'} = x'' + y^{2} + z_{i}^{n}\]
Don't start a paragraph with a displayed equation,
nor make one a paragraph by itself.

\end{document}               % End of document.

\section{Bench Huge Section 298}
% synthetic bench section 298
% This is a sample LaTeX input file.  (Version of 12 August 2004.)
%
% A '%' character causes TeX to ignore all remaining text on the line,
% and is used for comments like this one.

\documentclass{article}      % Specifies the document class

                             % The preamble begins here.
\title{An Example Document}  % Declares the document's title.
\author{Leslie Lamport}      % Declares the author's name.
\date{January 21, 1994}      % Deleting this command produces today's date.

\newcommand{\ip}[2]{(#1, #2)}
                             % Defines \ip{arg1}{arg2} to mean
                             % (arg1, arg2).

%\newcommand{\ip}[2]{\langle #1 | #2\rangle}
                             % This is an alternative definition of
                             % \ip that is commented out.

\begin{document}             % End of preamble and beginning of text.

\maketitle                   % Produces the title.

This is an example input file.  Comparing it with
the output it generates can show you how to
produce a simple document of your own.

\section{Ordinary Text}      % Produces section heading.  Lower-level
                             % sections are begun with similar
                             % \subsection and \subsubsection commands.

The ends  of words and sentences are marked
  by   spaces. It  doesn't matter how many
spaces    you type; one is as good as 100.  The
end of   a line counts as a space.

One   or more   blank lines denote the  end
of  a paragraph.

Since any number of consecutive spaces are treated
like a single one, the formatting of the input
file makes no difference to
      \LaTeX,                % The \LaTeX command generates the LaTeX logo.
but it makes a difference to you.  When you use
\LaTeX, making your input file as easy to read
as possible will be a great help as you write
your document and when you change it.  This sample
file shows how you can add comments to your own input
file.

Because printing is different from typewriting,
there are a number of things that you have to do
differently when preparing an input file than if
you were just typing the document directly.
Quotation marks like
       ``this''
have to be handled specially, as do quotes within
quotes:
       ``\,`this'            % \, separates the double and single quote.
        is what I just
        wrote, not  `that'\,''.

Dashes come in three sizes: an
       intra-word
dash, a medium dash for number ranges like
       1--2,
and a punctuation
       dash---like
this.

A sentence-ending space should be larger than the
space between words within a sentence.  You
sometimes have to type special commands in
conjunction with punctuation characters to get
this right, as in the following sentence.
       Gnats, gnus, etc.\ all  % `\ ' makes an inter-word space.
       begin with G\@.         % \@ marks end-of-sentence punctuation.
You should check the spaces after periods when
reading your output to make sure you haven't
forgotten any special cases.  Generating an
ellipsis
       \ldots\               % `\ ' is needed after `\ldots' because TeX
                             % ignores spaces after command names like \ldots
                             % made from \ + letters.
                             %
                             % Note how a `%' character causes TeX to ignore
                             % the end of the input line, so these blank lines
                             % do not start a new paragraph.
                             %
with the right spacing around the periods requires
a special command.

\LaTeX\ interprets some common characters as
commands, so you must type special commands to
generate them.  These characters include the
following:
       \$ \& \% \# \{ and \}.

In printing, text is usually emphasized with an
       \emph{italic}
type style.

\begin{em}
   A long segment of text can also be emphasized
   in this way.  Text within such a segment can be
   given \emph{additional} emphasis.
\end{em}

It is sometimes necessary to prevent \LaTeX\ from
breaking a line where it might otherwise do so.
This may be at a space, as between the ``Mr.''\ and
``Jones'' in
       ``Mr.~Jones'',        % ~ produces an unbreakable interword space.
or within a word---especially when the word is a
symbol like
       \mbox{\emph{itemnum}}
that makes little sense when hyphenated across
lines.

Footnotes\footnote{This is an example of a footnote.}
pose no problem.

\LaTeX\ is good at typesetting mathematical formulas
like
       \( x-3y + z = 7 \)
or
       \( a_{1} > x^{2n} + y^{2n} > x' \)
or
       \( \ip{A}{B} = \sum_{i} a_{i} b_{i} \).
The spaces you type in a formula are
ignored.  Remember that a letter like
       $x$                   % $ ... $  and  \( ... \)  are equivalent
is a formula when it denotes a mathematical
symbol, and it should be typed as one.

\section{Displayed Text}

Text is displayed by indenting it from the left
margin.  Quotations are commonly displayed.  There
are short quotations
\begin{quote}
   This is a short quotation.  It consists of a
   single paragraph of text.  See how it is formatted.
\end{quote}
and longer ones.
\begin{quotation}
   This is a longer quotation.  It consists of two
   paragraphs of text, neither of which are
   particularly interesting.

   This is the second paragraph of the quotation.  It
   is just as dull as the first paragraph.
\end{quotation}
Another frequently-displayed structure is a list.
The following is an example of an \emph{itemized}
list.
\begin{itemize}
   \item This is the first item of an itemized list.
         Each item in the list is marked with a ``tick''.
         You don't have to worry about what kind of tick
         mark is used.

   \item This is the second item of the list.  It
         contains another list nested inside it.  The inner
         list is an \emph{enumerated} list.
         \begin{enumerate}
            \item This is the first item of an enumerated
                  list that is nested within the itemized list.

            \item This is the second item of the inner list.
                  \LaTeX\ allows you to nest lists deeper than
                  you really should.
         \end{enumerate}
         This is the rest of the second item of the outer
         list.  It is no more interesting than any other
         part of the item.
   \item This is the third item of the list.
\end{itemize}
You can even display poetry.
\begin{verse}
   There is an environment
    for verse \\             % The \\ command separates lines
   Whose features some poets % within a stanza.
   will curse.

                             % One or more blank lines separate stanzas.

   For instead of making\\
   Them do \emph{all} line breaking, \\
   It allows them to put too many words on a line when they'd rather be
   forced to be terse.
\end{verse}

Mathematical formulas may also be displayed.  A
displayed formula
is
one-line long; multiline
formulas require special formatting instructions.
   \[  \ip{\Gamma}{\psi'} = x'' + y^{2} + z_{i}^{n}\]
Don't start a paragraph with a displayed equation,
nor make one a paragraph by itself.

\end{document}               % End of document.

\section{Bench Huge Section 299}
% synthetic bench section 299
% This is a sample LaTeX input file.  (Version of 12 August 2004.)
%
% A '%' character causes TeX to ignore all remaining text on the line,
% and is used for comments like this one.

\documentclass{article}      % Specifies the document class

                             % The preamble begins here.
\title{An Example Document}  % Declares the document's title.
\author{Leslie Lamport}      % Declares the author's name.
\date{January 21, 1994}      % Deleting this command produces today's date.

\newcommand{\ip}[2]{(#1, #2)}
                             % Defines \ip{arg1}{arg2} to mean
                             % (arg1, arg2).

%\newcommand{\ip}[2]{\langle #1 | #2\rangle}
                             % This is an alternative definition of
                             % \ip that is commented out.

\begin{document}             % End of preamble and beginning of text.

\maketitle                   % Produces the title.

This is an example input file.  Comparing it with
the output it generates can show you how to
produce a simple document of your own.

\section{Ordinary Text}      % Produces section heading.  Lower-level
                             % sections are begun with similar
                             % \subsection and \subsubsection commands.

The ends  of words and sentences are marked
  by   spaces. It  doesn't matter how many
spaces    you type; one is as good as 100.  The
end of   a line counts as a space.

One   or more   blank lines denote the  end
of  a paragraph.

Since any number of consecutive spaces are treated
like a single one, the formatting of the input
file makes no difference to
      \LaTeX,                % The \LaTeX command generates the LaTeX logo.
but it makes a difference to you.  When you use
\LaTeX, making your input file as easy to read
as possible will be a great help as you write
your document and when you change it.  This sample
file shows how you can add comments to your own input
file.

Because printing is different from typewriting,
there are a number of things that you have to do
differently when preparing an input file than if
you were just typing the document directly.
Quotation marks like
       ``this''
have to be handled specially, as do quotes within
quotes:
       ``\,`this'            % \, separates the double and single quote.
        is what I just
        wrote, not  `that'\,''.

Dashes come in three sizes: an
       intra-word
dash, a medium dash for number ranges like
       1--2,
and a punctuation
       dash---like
this.

A sentence-ending space should be larger than the
space between words within a sentence.  You
sometimes have to type special commands in
conjunction with punctuation characters to get
this right, as in the following sentence.
       Gnats, gnus, etc.\ all  % `\ ' makes an inter-word space.
       begin with G\@.         % \@ marks end-of-sentence punctuation.
You should check the spaces after periods when
reading your output to make sure you haven't
forgotten any special cases.  Generating an
ellipsis
       \ldots\               % `\ ' is needed after `\ldots' because TeX
                             % ignores spaces after command names like \ldots
                             % made from \ + letters.
                             %
                             % Note how a `%' character causes TeX to ignore
                             % the end of the input line, so these blank lines
                             % do not start a new paragraph.
                             %
with the right spacing around the periods requires
a special command.

\LaTeX\ interprets some common characters as
commands, so you must type special commands to
generate them.  These characters include the
following:
       \$ \& \% \# \{ and \}.

In printing, text is usually emphasized with an
       \emph{italic}
type style.

\begin{em}
   A long segment of text can also be emphasized
   in this way.  Text within such a segment can be
   given \emph{additional} emphasis.
\end{em}

It is sometimes necessary to prevent \LaTeX\ from
breaking a line where it might otherwise do so.
This may be at a space, as between the ``Mr.''\ and
``Jones'' in
       ``Mr.~Jones'',        % ~ produces an unbreakable interword space.
or within a word---especially when the word is a
symbol like
       \mbox{\emph{itemnum}}
that makes little sense when hyphenated across
lines.

Footnotes\footnote{This is an example of a footnote.}
pose no problem.

\LaTeX\ is good at typesetting mathematical formulas
like
       \( x-3y + z = 7 \)
or
       \( a_{1} > x^{2n} + y^{2n} > x' \)
or
       \( \ip{A}{B} = \sum_{i} a_{i} b_{i} \).
The spaces you type in a formula are
ignored.  Remember that a letter like
       $x$                   % $ ... $  and  \( ... \)  are equivalent
is a formula when it denotes a mathematical
symbol, and it should be typed as one.

\section{Displayed Text}

Text is displayed by indenting it from the left
margin.  Quotations are commonly displayed.  There
are short quotations
\begin{quote}
   This is a short quotation.  It consists of a
   single paragraph of text.  See how it is formatted.
\end{quote}
and longer ones.
\begin{quotation}
   This is a longer quotation.  It consists of two
   paragraphs of text, neither of which are
   particularly interesting.

   This is the second paragraph of the quotation.  It
   is just as dull as the first paragraph.
\end{quotation}
Another frequently-displayed structure is a list.
The following is an example of an \emph{itemized}
list.
\begin{itemize}
   \item This is the first item of an itemized list.
         Each item in the list is marked with a ``tick''.
         You don't have to worry about what kind of tick
         mark is used.

   \item This is the second item of the list.  It
         contains another list nested inside it.  The inner
         list is an \emph{enumerated} list.
         \begin{enumerate}
            \item This is the first item of an enumerated
                  list that is nested within the itemized list.

            \item This is the second item of the inner list.
                  \LaTeX\ allows you to nest lists deeper than
                  you really should.
         \end{enumerate}
         This is the rest of the second item of the outer
         list.  It is no more interesting than any other
         part of the item.
   \item This is the third item of the list.
\end{itemize}
You can even display poetry.
\begin{verse}
   There is an environment
    for verse \\             % The \\ command separates lines
   Whose features some poets % within a stanza.
   will curse.

                             % One or more blank lines separate stanzas.

   For instead of making\\
   Them do \emph{all} line breaking, \\
   It allows them to put too many words on a line when they'd rather be
   forced to be terse.
\end{verse}

Mathematical formulas may also be displayed.  A
displayed formula
is
one-line long; multiline
formulas require special formatting instructions.
   \[  \ip{\Gamma}{\psi'} = x'' + y^{2} + z_{i}^{n}\]
Don't start a paragraph with a displayed equation,
nor make one a paragraph by itself.

\end{document}               % End of document.

\section{Bench Huge Section 300}
% synthetic bench section 300
% This is a sample LaTeX input file.  (Version of 12 August 2004.)
%
% A '%' character causes TeX to ignore all remaining text on the line,
% and is used for comments like this one.

\documentclass{article}      % Specifies the document class

                             % The preamble begins here.
\title{An Example Document}  % Declares the document's title.
\author{Leslie Lamport}      % Declares the author's name.
\date{January 21, 1994}      % Deleting this command produces today's date.

\newcommand{\ip}[2]{(#1, #2)}
                             % Defines \ip{arg1}{arg2} to mean
                             % (arg1, arg2).

%\newcommand{\ip}[2]{\langle #1 | #2\rangle}
                             % This is an alternative definition of
                             % \ip that is commented out.

\begin{document}             % End of preamble and beginning of text.

\maketitle                   % Produces the title.

This is an example input file.  Comparing it with
the output it generates can show you how to
produce a simple document of your own.

\section{Ordinary Text}      % Produces section heading.  Lower-level
                             % sections are begun with similar
                             % \subsection and \subsubsection commands.

The ends  of words and sentences are marked
  by   spaces. It  doesn't matter how many
spaces    you type; one is as good as 100.  The
end of   a line counts as a space.

One   or more   blank lines denote the  end
of  a paragraph.

Since any number of consecutive spaces are treated
like a single one, the formatting of the input
file makes no difference to
      \LaTeX,                % The \LaTeX command generates the LaTeX logo.
but it makes a difference to you.  When you use
\LaTeX, making your input file as easy to read
as possible will be a great help as you write
your document and when you change it.  This sample
file shows how you can add comments to your own input
file.

Because printing is different from typewriting,
there are a number of things that you have to do
differently when preparing an input file than if
you were just typing the document directly.
Quotation marks like
       ``this''
have to be handled specially, as do quotes within
quotes:
       ``\,`this'            % \, separates the double and single quote.
        is what I just
        wrote, not  `that'\,''.

Dashes come in three sizes: an
       intra-word
dash, a medium dash for number ranges like
       1--2,
and a punctuation
       dash---like
this.

A sentence-ending space should be larger than the
space between words within a sentence.  You
sometimes have to type special commands in
conjunction with punctuation characters to get
this right, as in the following sentence.
       Gnats, gnus, etc.\ all  % `\ ' makes an inter-word space.
       begin with G\@.         % \@ marks end-of-sentence punctuation.
You should check the spaces after periods when
reading your output to make sure you haven't
forgotten any special cases.  Generating an
ellipsis
       \ldots\               % `\ ' is needed after `\ldots' because TeX
                             % ignores spaces after command names like \ldots
                             % made from \ + letters.
                             %
                             % Note how a `%' character causes TeX to ignore
                             % the end of the input line, so these blank lines
                             % do not start a new paragraph.
                             %
with the right spacing around the periods requires
a special command.

\LaTeX\ interprets some common characters as
commands, so you must type special commands to
generate them.  These characters include the
following:
       \$ \& \% \# \{ and \}.

In printing, text is usually emphasized with an
       \emph{italic}
type style.

\begin{em}
   A long segment of text can also be emphasized
   in this way.  Text within such a segment can be
   given \emph{additional} emphasis.
\end{em}

It is sometimes necessary to prevent \LaTeX\ from
breaking a line where it might otherwise do so.
This may be at a space, as between the ``Mr.''\ and
``Jones'' in
       ``Mr.~Jones'',        % ~ produces an unbreakable interword space.
or within a word---especially when the word is a
symbol like
       \mbox{\emph{itemnum}}
that makes little sense when hyphenated across
lines.

Footnotes\footnote{This is an example of a footnote.}
pose no problem.

\LaTeX\ is good at typesetting mathematical formulas
like
       \( x-3y + z = 7 \)
or
       \( a_{1} > x^{2n} + y^{2n} > x' \)
or
       \( \ip{A}{B} = \sum_{i} a_{i} b_{i} \).
The spaces you type in a formula are
ignored.  Remember that a letter like
       $x$                   % $ ... $  and  \( ... \)  are equivalent
is a formula when it denotes a mathematical
symbol, and it should be typed as one.

\section{Displayed Text}

Text is displayed by indenting it from the left
margin.  Quotations are commonly displayed.  There
are short quotations
\begin{quote}
   This is a short quotation.  It consists of a
   single paragraph of text.  See how it is formatted.
\end{quote}
and longer ones.
\begin{quotation}
   This is a longer quotation.  It consists of two
   paragraphs of text, neither of which are
   particularly interesting.

   This is the second paragraph of the quotation.  It
   is just as dull as the first paragraph.
\end{quotation}
Another frequently-displayed structure is a list.
The following is an example of an \emph{itemized}
list.
\begin{itemize}
   \item This is the first item of an itemized list.
         Each item in the list is marked with a ``tick''.
         You don't have to worry about what kind of tick
         mark is used.

   \item This is the second item of the list.  It
         contains another list nested inside it.  The inner
         list is an \emph{enumerated} list.
         \begin{enumerate}
            \item This is the first item of an enumerated
                  list that is nested within the itemized list.

            \item This is the second item of the inner list.
                  \LaTeX\ allows you to nest lists deeper than
                  you really should.
         \end{enumerate}
         This is the rest of the second item of the outer
         list.  It is no more interesting than any other
         part of the item.
   \item This is the third item of the list.
\end{itemize}
You can even display poetry.
\begin{verse}
   There is an environment
    for verse \\             % The \\ command separates lines
   Whose features some poets % within a stanza.
   will curse.

                             % One or more blank lines separate stanzas.

   For instead of making\\
   Them do \emph{all} line breaking, \\
   It allows them to put too many words on a line when they'd rather be
   forced to be terse.
\end{verse}

Mathematical formulas may also be displayed.  A
displayed formula
is
one-line long; multiline
formulas require special formatting instructions.
   \[  \ip{\Gamma}{\psi'} = x'' + y^{2} + z_{i}^{n}\]
Don't start a paragraph with a displayed equation,
nor make one a paragraph by itself.

\end{document}               % End of document.

\section{Bench Huge Section 301}
% synthetic bench section 301
% This is a sample LaTeX input file.  (Version of 12 August 2004.)
%
% A '%' character causes TeX to ignore all remaining text on the line,
% and is used for comments like this one.

\documentclass{article}      % Specifies the document class

                             % The preamble begins here.
\title{An Example Document}  % Declares the document's title.
\author{Leslie Lamport}      % Declares the author's name.
\date{January 21, 1994}      % Deleting this command produces today's date.

\newcommand{\ip}[2]{(#1, #2)}
                             % Defines \ip{arg1}{arg2} to mean
                             % (arg1, arg2).

%\newcommand{\ip}[2]{\langle #1 | #2\rangle}
                             % This is an alternative definition of
                             % \ip that is commented out.

\begin{document}             % End of preamble and beginning of text.

\maketitle                   % Produces the title.

This is an example input file.  Comparing it with
the output it generates can show you how to
produce a simple document of your own.

\section{Ordinary Text}      % Produces section heading.  Lower-level
                             % sections are begun with similar
                             % \subsection and \subsubsection commands.

The ends  of words and sentences are marked
  by   spaces. It  doesn't matter how many
spaces    you type; one is as good as 100.  The
end of   a line counts as a space.

One   or more   blank lines denote the  end
of  a paragraph.

Since any number of consecutive spaces are treated
like a single one, the formatting of the input
file makes no difference to
      \LaTeX,                % The \LaTeX command generates the LaTeX logo.
but it makes a difference to you.  When you use
\LaTeX, making your input file as easy to read
as possible will be a great help as you write
your document and when you change it.  This sample
file shows how you can add comments to your own input
file.

Because printing is different from typewriting,
there are a number of things that you have to do
differently when preparing an input file than if
you were just typing the document directly.
Quotation marks like
       ``this''
have to be handled specially, as do quotes within
quotes:
       ``\,`this'            % \, separates the double and single quote.
        is what I just
        wrote, not  `that'\,''.

Dashes come in three sizes: an
       intra-word
dash, a medium dash for number ranges like
       1--2,
and a punctuation
       dash---like
this.

A sentence-ending space should be larger than the
space between words within a sentence.  You
sometimes have to type special commands in
conjunction with punctuation characters to get
this right, as in the following sentence.
       Gnats, gnus, etc.\ all  % `\ ' makes an inter-word space.
       begin with G\@.         % \@ marks end-of-sentence punctuation.
You should check the spaces after periods when
reading your output to make sure you haven't
forgotten any special cases.  Generating an
ellipsis
       \ldots\               % `\ ' is needed after `\ldots' because TeX
                             % ignores spaces after command names like \ldots
                             % made from \ + letters.
                             %
                             % Note how a `%' character causes TeX to ignore
                             % the end of the input line, so these blank lines
                             % do not start a new paragraph.
                             %
with the right spacing around the periods requires
a special command.

\LaTeX\ interprets some common characters as
commands, so you must type special commands to
generate them.  These characters include the
following:
       \$ \& \% \# \{ and \}.

In printing, text is usually emphasized with an
       \emph{italic}
type style.

\begin{em}
   A long segment of text can also be emphasized
   in this way.  Text within such a segment can be
   given \emph{additional} emphasis.
\end{em}

It is sometimes necessary to prevent \LaTeX\ from
breaking a line where it might otherwise do so.
This may be at a space, as between the ``Mr.''\ and
``Jones'' in
       ``Mr.~Jones'',        % ~ produces an unbreakable interword space.
or within a word---especially when the word is a
symbol like
       \mbox{\emph{itemnum}}
that makes little sense when hyphenated across
lines.

Footnotes\footnote{This is an example of a footnote.}
pose no problem.

\LaTeX\ is good at typesetting mathematical formulas
like
       \( x-3y + z = 7 \)
or
       \( a_{1} > x^{2n} + y^{2n} > x' \)
or
       \( \ip{A}{B} = \sum_{i} a_{i} b_{i} \).
The spaces you type in a formula are
ignored.  Remember that a letter like
       $x$                   % $ ... $  and  \( ... \)  are equivalent
is a formula when it denotes a mathematical
symbol, and it should be typed as one.

\section{Displayed Text}

Text is displayed by indenting it from the left
margin.  Quotations are commonly displayed.  There
are short quotations
\begin{quote}
   This is a short quotation.  It consists of a
   single paragraph of text.  See how it is formatted.
\end{quote}
and longer ones.
\begin{quotation}
   This is a longer quotation.  It consists of two
   paragraphs of text, neither of which are
   particularly interesting.

   This is the second paragraph of the quotation.  It
   is just as dull as the first paragraph.
\end{quotation}
Another frequently-displayed structure is a list.
The following is an example of an \emph{itemized}
list.
\begin{itemize}
   \item This is the first item of an itemized list.
         Each item in the list is marked with a ``tick''.
         You don't have to worry about what kind of tick
         mark is used.

   \item This is the second item of the list.  It
         contains another list nested inside it.  The inner
         list is an \emph{enumerated} list.
         \begin{enumerate}
            \item This is the first item of an enumerated
                  list that is nested within the itemized list.

            \item This is the second item of the inner list.
                  \LaTeX\ allows you to nest lists deeper than
                  you really should.
         \end{enumerate}
         This is the rest of the second item of the outer
         list.  It is no more interesting than any other
         part of the item.
   \item This is the third item of the list.
\end{itemize}
You can even display poetry.
\begin{verse}
   There is an environment
    for verse \\             % The \\ command separates lines
   Whose features some poets % within a stanza.
   will curse.

                             % One or more blank lines separate stanzas.

   For instead of making\\
   Them do \emph{all} line breaking, \\
   It allows them to put too many words on a line when they'd rather be
   forced to be terse.
\end{verse}

Mathematical formulas may also be displayed.  A
displayed formula
is
one-line long; multiline
formulas require special formatting instructions.
   \[  \ip{\Gamma}{\psi'} = x'' + y^{2} + z_{i}^{n}\]
Don't start a paragraph with a displayed equation,
nor make one a paragraph by itself.

\end{document}               % End of document.

\section{Bench Huge Section 302}
% synthetic bench section 302
% This is a sample LaTeX input file.  (Version of 12 August 2004.)
%
% A '%' character causes TeX to ignore all remaining text on the line,
% and is used for comments like this one.

\documentclass{article}      % Specifies the document class

                             % The preamble begins here.
\title{An Example Document}  % Declares the document's title.
\author{Leslie Lamport}      % Declares the author's name.
\date{January 21, 1994}      % Deleting this command produces today's date.

\newcommand{\ip}[2]{(#1, #2)}
                             % Defines \ip{arg1}{arg2} to mean
                             % (arg1, arg2).

%\newcommand{\ip}[2]{\langle #1 | #2\rangle}
                             % This is an alternative definition of
                             % \ip that is commented out.

\begin{document}             % End of preamble and beginning of text.

\maketitle                   % Produces the title.

This is an example input file.  Comparing it with
the output it generates can show you how to
produce a simple document of your own.

\section{Ordinary Text}      % Produces section heading.  Lower-level
                             % sections are begun with similar
                             % \subsection and \subsubsection commands.

The ends  of words and sentences are marked
  by   spaces. It  doesn't matter how many
spaces    you type; one is as good as 100.  The
end of   a line counts as a space.

One   or more   blank lines denote the  end
of  a paragraph.

Since any number of consecutive spaces are treated
like a single one, the formatting of the input
file makes no difference to
      \LaTeX,                % The \LaTeX command generates the LaTeX logo.
but it makes a difference to you.  When you use
\LaTeX, making your input file as easy to read
as possible will be a great help as you write
your document and when you change it.  This sample
file shows how you can add comments to your own input
file.

Because printing is different from typewriting,
there are a number of things that you have to do
differently when preparing an input file than if
you were just typing the document directly.
Quotation marks like
       ``this''
have to be handled specially, as do quotes within
quotes:
       ``\,`this'            % \, separates the double and single quote.
        is what I just
        wrote, not  `that'\,''.

Dashes come in three sizes: an
       intra-word
dash, a medium dash for number ranges like
       1--2,
and a punctuation
       dash---like
this.

A sentence-ending space should be larger than the
space between words within a sentence.  You
sometimes have to type special commands in
conjunction with punctuation characters to get
this right, as in the following sentence.
       Gnats, gnus, etc.\ all  % `\ ' makes an inter-word space.
       begin with G\@.         % \@ marks end-of-sentence punctuation.
You should check the spaces after periods when
reading your output to make sure you haven't
forgotten any special cases.  Generating an
ellipsis
       \ldots\               % `\ ' is needed after `\ldots' because TeX
                             % ignores spaces after command names like \ldots
                             % made from \ + letters.
                             %
                             % Note how a `%' character causes TeX to ignore
                             % the end of the input line, so these blank lines
                             % do not start a new paragraph.
                             %
with the right spacing around the periods requires
a special command.

\LaTeX\ interprets some common characters as
commands, so you must type special commands to
generate them.  These characters include the
following:
       \$ \& \% \# \{ and \}.

In printing, text is usually emphasized with an
       \emph{italic}
type style.

\begin{em}
   A long segment of text can also be emphasized
   in this way.  Text within such a segment can be
   given \emph{additional} emphasis.
\end{em}

It is sometimes necessary to prevent \LaTeX\ from
breaking a line where it might otherwise do so.
This may be at a space, as between the ``Mr.''\ and
``Jones'' in
       ``Mr.~Jones'',        % ~ produces an unbreakable interword space.
or within a word---especially when the word is a
symbol like
       \mbox{\emph{itemnum}}
that makes little sense when hyphenated across
lines.

Footnotes\footnote{This is an example of a footnote.}
pose no problem.

\LaTeX\ is good at typesetting mathematical formulas
like
       \( x-3y + z = 7 \)
or
       \( a_{1} > x^{2n} + y^{2n} > x' \)
or
       \( \ip{A}{B} = \sum_{i} a_{i} b_{i} \).
The spaces you type in a formula are
ignored.  Remember that a letter like
       $x$                   % $ ... $  and  \( ... \)  are equivalent
is a formula when it denotes a mathematical
symbol, and it should be typed as one.

\section{Displayed Text}

Text is displayed by indenting it from the left
margin.  Quotations are commonly displayed.  There
are short quotations
\begin{quote}
   This is a short quotation.  It consists of a
   single paragraph of text.  See how it is formatted.
\end{quote}
and longer ones.
\begin{quotation}
   This is a longer quotation.  It consists of two
   paragraphs of text, neither of which are
   particularly interesting.

   This is the second paragraph of the quotation.  It
   is just as dull as the first paragraph.
\end{quotation}
Another frequently-displayed structure is a list.
The following is an example of an \emph{itemized}
list.
\begin{itemize}
   \item This is the first item of an itemized list.
         Each item in the list is marked with a ``tick''.
         You don't have to worry about what kind of tick
         mark is used.

   \item This is the second item of the list.  It
         contains another list nested inside it.  The inner
         list is an \emph{enumerated} list.
         \begin{enumerate}
            \item This is the first item of an enumerated
                  list that is nested within the itemized list.

            \item This is the second item of the inner list.
                  \LaTeX\ allows you to nest lists deeper than
                  you really should.
         \end{enumerate}
         This is the rest of the second item of the outer
         list.  It is no more interesting than any other
         part of the item.
   \item This is the third item of the list.
\end{itemize}
You can even display poetry.
\begin{verse}
   There is an environment
    for verse \\             % The \\ command separates lines
   Whose features some poets % within a stanza.
   will curse.

                             % One or more blank lines separate stanzas.

   For instead of making\\
   Them do \emph{all} line breaking, \\
   It allows them to put too many words on a line when they'd rather be
   forced to be terse.
\end{verse}

Mathematical formulas may also be displayed.  A
displayed formula
is
one-line long; multiline
formulas require special formatting instructions.
   \[  \ip{\Gamma}{\psi'} = x'' + y^{2} + z_{i}^{n}\]
Don't start a paragraph with a displayed equation,
nor make one a paragraph by itself.

\end{document}               % End of document.

\section{Bench Huge Section 303}
% synthetic bench section 303
% This is a sample LaTeX input file.  (Version of 12 August 2004.)
%
% A '%' character causes TeX to ignore all remaining text on the line,
% and is used for comments like this one.

\documentclass{article}      % Specifies the document class

                             % The preamble begins here.
\title{An Example Document}  % Declares the document's title.
\author{Leslie Lamport}      % Declares the author's name.
\date{January 21, 1994}      % Deleting this command produces today's date.

\newcommand{\ip}[2]{(#1, #2)}
                             % Defines \ip{arg1}{arg2} to mean
                             % (arg1, arg2).

%\newcommand{\ip}[2]{\langle #1 | #2\rangle}
                             % This is an alternative definition of
                             % \ip that is commented out.

\begin{document}             % End of preamble and beginning of text.

\maketitle                   % Produces the title.

This is an example input file.  Comparing it with
the output it generates can show you how to
produce a simple document of your own.

\section{Ordinary Text}      % Produces section heading.  Lower-level
                             % sections are begun with similar
                             % \subsection and \subsubsection commands.

The ends  of words and sentences are marked
  by   spaces. It  doesn't matter how many
spaces    you type; one is as good as 100.  The
end of   a line counts as a space.

One   or more   blank lines denote the  end
of  a paragraph.

Since any number of consecutive spaces are treated
like a single one, the formatting of the input
file makes no difference to
      \LaTeX,                % The \LaTeX command generates the LaTeX logo.
but it makes a difference to you.  When you use
\LaTeX, making your input file as easy to read
as possible will be a great help as you write
your document and when you change it.  This sample
file shows how you can add comments to your own input
file.

Because printing is different from typewriting,
there are a number of things that you have to do
differently when preparing an input file than if
you were just typing the document directly.
Quotation marks like
       ``this''
have to be handled specially, as do quotes within
quotes:
       ``\,`this'            % \, separates the double and single quote.
        is what I just
        wrote, not  `that'\,''.

Dashes come in three sizes: an
       intra-word
dash, a medium dash for number ranges like
       1--2,
and a punctuation
       dash---like
this.

A sentence-ending space should be larger than the
space between words within a sentence.  You
sometimes have to type special commands in
conjunction with punctuation characters to get
this right, as in the following sentence.
       Gnats, gnus, etc.\ all  % `\ ' makes an inter-word space.
       begin with G\@.         % \@ marks end-of-sentence punctuation.
You should check the spaces after periods when
reading your output to make sure you haven't
forgotten any special cases.  Generating an
ellipsis
       \ldots\               % `\ ' is needed after `\ldots' because TeX
                             % ignores spaces after command names like \ldots
                             % made from \ + letters.
                             %
                             % Note how a `%' character causes TeX to ignore
                             % the end of the input line, so these blank lines
                             % do not start a new paragraph.
                             %
with the right spacing around the periods requires
a special command.

\LaTeX\ interprets some common characters as
commands, so you must type special commands to
generate them.  These characters include the
following:
       \$ \& \% \# \{ and \}.

In printing, text is usually emphasized with an
       \emph{italic}
type style.

\begin{em}
   A long segment of text can also be emphasized
   in this way.  Text within such a segment can be
   given \emph{additional} emphasis.
\end{em}

It is sometimes necessary to prevent \LaTeX\ from
breaking a line where it might otherwise do so.
This may be at a space, as between the ``Mr.''\ and
``Jones'' in
       ``Mr.~Jones'',        % ~ produces an unbreakable interword space.
or within a word---especially when the word is a
symbol like
       \mbox{\emph{itemnum}}
that makes little sense when hyphenated across
lines.

Footnotes\footnote{This is an example of a footnote.}
pose no problem.

\LaTeX\ is good at typesetting mathematical formulas
like
       \( x-3y + z = 7 \)
or
       \( a_{1} > x^{2n} + y^{2n} > x' \)
or
       \( \ip{A}{B} = \sum_{i} a_{i} b_{i} \).
The spaces you type in a formula are
ignored.  Remember that a letter like
       $x$                   % $ ... $  and  \( ... \)  are equivalent
is a formula when it denotes a mathematical
symbol, and it should be typed as one.

\section{Displayed Text}

Text is displayed by indenting it from the left
margin.  Quotations are commonly displayed.  There
are short quotations
\begin{quote}
   This is a short quotation.  It consists of a
   single paragraph of text.  See how it is formatted.
\end{quote}
and longer ones.
\begin{quotation}
   This is a longer quotation.  It consists of two
   paragraphs of text, neither of which are
   particularly interesting.

   This is the second paragraph of the quotation.  It
   is just as dull as the first paragraph.
\end{quotation}
Another frequently-displayed structure is a list.
The following is an example of an \emph{itemized}
list.
\begin{itemize}
   \item This is the first item of an itemized list.
         Each item in the list is marked with a ``tick''.
         You don't have to worry about what kind of tick
         mark is used.

   \item This is the second item of the list.  It
         contains another list nested inside it.  The inner
         list is an \emph{enumerated} list.
         \begin{enumerate}
            \item This is the first item of an enumerated
                  list that is nested within the itemized list.

            \item This is the second item of the inner list.
                  \LaTeX\ allows you to nest lists deeper than
                  you really should.
         \end{enumerate}
         This is the rest of the second item of the outer
         list.  It is no more interesting than any other
         part of the item.
   \item This is the third item of the list.
\end{itemize}
You can even display poetry.
\begin{verse}
   There is an environment
    for verse \\             % The \\ command separates lines
   Whose features some poets % within a stanza.
   will curse.

                             % One or more blank lines separate stanzas.

   For instead of making\\
   Them do \emph{all} line breaking, \\
   It allows them to put too many words on a line when they'd rather be
   forced to be terse.
\end{verse}

Mathematical formulas may also be displayed.  A
displayed formula
is
one-line long; multiline
formulas require special formatting instructions.
   \[  \ip{\Gamma}{\psi'} = x'' + y^{2} + z_{i}^{n}\]
Don't start a paragraph with a displayed equation,
nor make one a paragraph by itself.

\end{document}               % End of document.

\section{Bench Huge Section 304}
% synthetic bench section 304
% This is a sample LaTeX input file.  (Version of 12 August 2004.)
%
% A '%' character causes TeX to ignore all remaining text on the line,
% and is used for comments like this one.

\documentclass{article}      % Specifies the document class

                             % The preamble begins here.
\title{An Example Document}  % Declares the document's title.
\author{Leslie Lamport}      % Declares the author's name.
\date{January 21, 1994}      % Deleting this command produces today's date.

\newcommand{\ip}[2]{(#1, #2)}
                             % Defines \ip{arg1}{arg2} to mean
                             % (arg1, arg2).

%\newcommand{\ip}[2]{\langle #1 | #2\rangle}
                             % This is an alternative definition of
                             % \ip that is commented out.

\begin{document}             % End of preamble and beginning of text.

\maketitle                   % Produces the title.

This is an example input file.  Comparing it with
the output it generates can show you how to
produce a simple document of your own.

\section{Ordinary Text}      % Produces section heading.  Lower-level
                             % sections are begun with similar
                             % \subsection and \subsubsection commands.

The ends  of words and sentences are marked
  by   spaces. It  doesn't matter how many
spaces    you type; one is as good as 100.  The
end of   a line counts as a space.

One   or more   blank lines denote the  end
of  a paragraph.

Since any number of consecutive spaces are treated
like a single one, the formatting of the input
file makes no difference to
      \LaTeX,                % The \LaTeX command generates the LaTeX logo.
but it makes a difference to you.  When you use
\LaTeX, making your input file as easy to read
as possible will be a great help as you write
your document and when you change it.  This sample
file shows how you can add comments to your own input
file.

Because printing is different from typewriting,
there are a number of things that you have to do
differently when preparing an input file than if
you were just typing the document directly.
Quotation marks like
       ``this''
have to be handled specially, as do quotes within
quotes:
       ``\,`this'            % \, separates the double and single quote.
        is what I just
        wrote, not  `that'\,''.

Dashes come in three sizes: an
       intra-word
dash, a medium dash for number ranges like
       1--2,
and a punctuation
       dash---like
this.

A sentence-ending space should be larger than the
space between words within a sentence.  You
sometimes have to type special commands in
conjunction with punctuation characters to get
this right, as in the following sentence.
       Gnats, gnus, etc.\ all  % `\ ' makes an inter-word space.
       begin with G\@.         % \@ marks end-of-sentence punctuation.
You should check the spaces after periods when
reading your output to make sure you haven't
forgotten any special cases.  Generating an
ellipsis
       \ldots\               % `\ ' is needed after `\ldots' because TeX
                             % ignores spaces after command names like \ldots
                             % made from \ + letters.
                             %
                             % Note how a `%' character causes TeX to ignore
                             % the end of the input line, so these blank lines
                             % do not start a new paragraph.
                             %
with the right spacing around the periods requires
a special command.

\LaTeX\ interprets some common characters as
commands, so you must type special commands to
generate them.  These characters include the
following:
       \$ \& \% \# \{ and \}.

In printing, text is usually emphasized with an
       \emph{italic}
type style.

\begin{em}
   A long segment of text can also be emphasized
   in this way.  Text within such a segment can be
   given \emph{additional} emphasis.
\end{em}

It is sometimes necessary to prevent \LaTeX\ from
breaking a line where it might otherwise do so.
This may be at a space, as between the ``Mr.''\ and
``Jones'' in
       ``Mr.~Jones'',        % ~ produces an unbreakable interword space.
or within a word---especially when the word is a
symbol like
       \mbox{\emph{itemnum}}
that makes little sense when hyphenated across
lines.

Footnotes\footnote{This is an example of a footnote.}
pose no problem.

\LaTeX\ is good at typesetting mathematical formulas
like
       \( x-3y + z = 7 \)
or
       \( a_{1} > x^{2n} + y^{2n} > x' \)
or
       \( \ip{A}{B} = \sum_{i} a_{i} b_{i} \).
The spaces you type in a formula are
ignored.  Remember that a letter like
       $x$                   % $ ... $  and  \( ... \)  are equivalent
is a formula when it denotes a mathematical
symbol, and it should be typed as one.

\section{Displayed Text}

Text is displayed by indenting it from the left
margin.  Quotations are commonly displayed.  There
are short quotations
\begin{quote}
   This is a short quotation.  It consists of a
   single paragraph of text.  See how it is formatted.
\end{quote}
and longer ones.
\begin{quotation}
   This is a longer quotation.  It consists of two
   paragraphs of text, neither of which are
   particularly interesting.

   This is the second paragraph of the quotation.  It
   is just as dull as the first paragraph.
\end{quotation}
Another frequently-displayed structure is a list.
The following is an example of an \emph{itemized}
list.
\begin{itemize}
   \item This is the first item of an itemized list.
         Each item in the list is marked with a ``tick''.
         You don't have to worry about what kind of tick
         mark is used.

   \item This is the second item of the list.  It
         contains another list nested inside it.  The inner
         list is an \emph{enumerated} list.
         \begin{enumerate}
            \item This is the first item of an enumerated
                  list that is nested within the itemized list.

            \item This is the second item of the inner list.
                  \LaTeX\ allows you to nest lists deeper than
                  you really should.
         \end{enumerate}
         This is the rest of the second item of the outer
         list.  It is no more interesting than any other
         part of the item.
   \item This is the third item of the list.
\end{itemize}
You can even display poetry.
\begin{verse}
   There is an environment
    for verse \\             % The \\ command separates lines
   Whose features some poets % within a stanza.
   will curse.

                             % One or more blank lines separate stanzas.

   For instead of making\\
   Them do \emph{all} line breaking, \\
   It allows them to put too many words on a line when they'd rather be
   forced to be terse.
\end{verse}

Mathematical formulas may also be displayed.  A
displayed formula
is
one-line long; multiline
formulas require special formatting instructions.
   \[  \ip{\Gamma}{\psi'} = x'' + y^{2} + z_{i}^{n}\]
Don't start a paragraph with a displayed equation,
nor make one a paragraph by itself.

\end{document}               % End of document.

\section{Bench Huge Section 305}
% synthetic bench section 305
% This is a sample LaTeX input file.  (Version of 12 August 2004.)
%
% A '%' character causes TeX to ignore all remaining text on the line,
% and is used for comments like this one.

\documentclass{article}      % Specifies the document class

                             % The preamble begins here.
\title{An Example Document}  % Declares the document's title.
\author{Leslie Lamport}      % Declares the author's name.
\date{January 21, 1994}      % Deleting this command produces today's date.

\newcommand{\ip}[2]{(#1, #2)}
                             % Defines \ip{arg1}{arg2} to mean
                             % (arg1, arg2).

%\newcommand{\ip}[2]{\langle #1 | #2\rangle}
                             % This is an alternative definition of
                             % \ip that is commented out.

\begin{document}             % End of preamble and beginning of text.

\maketitle                   % Produces the title.

This is an example input file.  Comparing it with
the output it generates can show you how to
produce a simple document of your own.

\section{Ordinary Text}      % Produces section heading.  Lower-level
                             % sections are begun with similar
                             % \subsection and \subsubsection commands.

The ends  of words and sentences are marked
  by   spaces. It  doesn't matter how many
spaces    you type; one is as good as 100.  The
end of   a line counts as a space.

One   or more   blank lines denote the  end
of  a paragraph.

Since any number of consecutive spaces are treated
like a single one, the formatting of the input
file makes no difference to
      \LaTeX,                % The \LaTeX command generates the LaTeX logo.
but it makes a difference to you.  When you use
\LaTeX, making your input file as easy to read
as possible will be a great help as you write
your document and when you change it.  This sample
file shows how you can add comments to your own input
file.

Because printing is different from typewriting,
there are a number of things that you have to do
differently when preparing an input file than if
you were just typing the document directly.
Quotation marks like
       ``this''
have to be handled specially, as do quotes within
quotes:
       ``\,`this'            % \, separates the double and single quote.
        is what I just
        wrote, not  `that'\,''.

Dashes come in three sizes: an
       intra-word
dash, a medium dash for number ranges like
       1--2,
and a punctuation
       dash---like
this.

A sentence-ending space should be larger than the
space between words within a sentence.  You
sometimes have to type special commands in
conjunction with punctuation characters to get
this right, as in the following sentence.
       Gnats, gnus, etc.\ all  % `\ ' makes an inter-word space.
       begin with G\@.         % \@ marks end-of-sentence punctuation.
You should check the spaces after periods when
reading your output to make sure you haven't
forgotten any special cases.  Generating an
ellipsis
       \ldots\               % `\ ' is needed after `\ldots' because TeX
                             % ignores spaces after command names like \ldots
                             % made from \ + letters.
                             %
                             % Note how a `%' character causes TeX to ignore
                             % the end of the input line, so these blank lines
                             % do not start a new paragraph.
                             %
with the right spacing around the periods requires
a special command.

\LaTeX\ interprets some common characters as
commands, so you must type special commands to
generate them.  These characters include the
following:
       \$ \& \% \# \{ and \}.

In printing, text is usually emphasized with an
       \emph{italic}
type style.

\begin{em}
   A long segment of text can also be emphasized
   in this way.  Text within such a segment can be
   given \emph{additional} emphasis.
\end{em}

It is sometimes necessary to prevent \LaTeX\ from
breaking a line where it might otherwise do so.
This may be at a space, as between the ``Mr.''\ and
``Jones'' in
       ``Mr.~Jones'',        % ~ produces an unbreakable interword space.
or within a word---especially when the word is a
symbol like
       \mbox{\emph{itemnum}}
that makes little sense when hyphenated across
lines.

Footnotes\footnote{This is an example of a footnote.}
pose no problem.

\LaTeX\ is good at typesetting mathematical formulas
like
       \( x-3y + z = 7 \)
or
       \( a_{1} > x^{2n} + y^{2n} > x' \)
or
       \( \ip{A}{B} = \sum_{i} a_{i} b_{i} \).
The spaces you type in a formula are
ignored.  Remember that a letter like
       $x$                   % $ ... $  and  \( ... \)  are equivalent
is a formula when it denotes a mathematical
symbol, and it should be typed as one.

\section{Displayed Text}

Text is displayed by indenting it from the left
margin.  Quotations are commonly displayed.  There
are short quotations
\begin{quote}
   This is a short quotation.  It consists of a
   single paragraph of text.  See how it is formatted.
\end{quote}
and longer ones.
\begin{quotation}
   This is a longer quotation.  It consists of two
   paragraphs of text, neither of which are
   particularly interesting.

   This is the second paragraph of the quotation.  It
   is just as dull as the first paragraph.
\end{quotation}
Another frequently-displayed structure is a list.
The following is an example of an \emph{itemized}
list.
\begin{itemize}
   \item This is the first item of an itemized list.
         Each item in the list is marked with a ``tick''.
         You don't have to worry about what kind of tick
         mark is used.

   \item This is the second item of the list.  It
         contains another list nested inside it.  The inner
         list is an \emph{enumerated} list.
         \begin{enumerate}
            \item This is the first item of an enumerated
                  list that is nested within the itemized list.

            \item This is the second item of the inner list.
                  \LaTeX\ allows you to nest lists deeper than
                  you really should.
         \end{enumerate}
         This is the rest of the second item of the outer
         list.  It is no more interesting than any other
         part of the item.
   \item This is the third item of the list.
\end{itemize}
You can even display poetry.
\begin{verse}
   There is an environment
    for verse \\             % The \\ command separates lines
   Whose features some poets % within a stanza.
   will curse.

                             % One or more blank lines separate stanzas.

   For instead of making\\
   Them do \emph{all} line breaking, \\
   It allows them to put too many words on a line when they'd rather be
   forced to be terse.
\end{verse}

Mathematical formulas may also be displayed.  A
displayed formula
is
one-line long; multiline
formulas require special formatting instructions.
   \[  \ip{\Gamma}{\psi'} = x'' + y^{2} + z_{i}^{n}\]
Don't start a paragraph with a displayed equation,
nor make one a paragraph by itself.

\end{document}               % End of document.

\section{Bench Huge Section 306}
% synthetic bench section 306
% This is a sample LaTeX input file.  (Version of 12 August 2004.)
%
% A '%' character causes TeX to ignore all remaining text on the line,
% and is used for comments like this one.

\documentclass{article}      % Specifies the document class

                             % The preamble begins here.
\title{An Example Document}  % Declares the document's title.
\author{Leslie Lamport}      % Declares the author's name.
\date{January 21, 1994}      % Deleting this command produces today's date.

\newcommand{\ip}[2]{(#1, #2)}
                             % Defines \ip{arg1}{arg2} to mean
                             % (arg1, arg2).

%\newcommand{\ip}[2]{\langle #1 | #2\rangle}
                             % This is an alternative definition of
                             % \ip that is commented out.

\begin{document}             % End of preamble and beginning of text.

\maketitle                   % Produces the title.

This is an example input file.  Comparing it with
the output it generates can show you how to
produce a simple document of your own.

\section{Ordinary Text}      % Produces section heading.  Lower-level
                             % sections are begun with similar
                             % \subsection and \subsubsection commands.

The ends  of words and sentences are marked
  by   spaces. It  doesn't matter how many
spaces    you type; one is as good as 100.  The
end of   a line counts as a space.

One   or more   blank lines denote the  end
of  a paragraph.

Since any number of consecutive spaces are treated
like a single one, the formatting of the input
file makes no difference to
      \LaTeX,                % The \LaTeX command generates the LaTeX logo.
but it makes a difference to you.  When you use
\LaTeX, making your input file as easy to read
as possible will be a great help as you write
your document and when you change it.  This sample
file shows how you can add comments to your own input
file.

Because printing is different from typewriting,
there are a number of things that you have to do
differently when preparing an input file than if
you were just typing the document directly.
Quotation marks like
       ``this''
have to be handled specially, as do quotes within
quotes:
       ``\,`this'            % \, separates the double and single quote.
        is what I just
        wrote, not  `that'\,''.

Dashes come in three sizes: an
       intra-word
dash, a medium dash for number ranges like
       1--2,
and a punctuation
       dash---like
this.

A sentence-ending space should be larger than the
space between words within a sentence.  You
sometimes have to type special commands in
conjunction with punctuation characters to get
this right, as in the following sentence.
       Gnats, gnus, etc.\ all  % `\ ' makes an inter-word space.
       begin with G\@.         % \@ marks end-of-sentence punctuation.
You should check the spaces after periods when
reading your output to make sure you haven't
forgotten any special cases.  Generating an
ellipsis
       \ldots\               % `\ ' is needed after `\ldots' because TeX
                             % ignores spaces after command names like \ldots
                             % made from \ + letters.
                             %
                             % Note how a `%' character causes TeX to ignore
                             % the end of the input line, so these blank lines
                             % do not start a new paragraph.
                             %
with the right spacing around the periods requires
a special command.

\LaTeX\ interprets some common characters as
commands, so you must type special commands to
generate them.  These characters include the
following:
       \$ \& \% \# \{ and \}.

In printing, text is usually emphasized with an
       \emph{italic}
type style.

\begin{em}
   A long segment of text can also be emphasized
   in this way.  Text within such a segment can be
   given \emph{additional} emphasis.
\end{em}

It is sometimes necessary to prevent \LaTeX\ from
breaking a line where it might otherwise do so.
This may be at a space, as between the ``Mr.''\ and
``Jones'' in
       ``Mr.~Jones'',        % ~ produces an unbreakable interword space.
or within a word---especially when the word is a
symbol like
       \mbox{\emph{itemnum}}
that makes little sense when hyphenated across
lines.

Footnotes\footnote{This is an example of a footnote.}
pose no problem.

\LaTeX\ is good at typesetting mathematical formulas
like
       \( x-3y + z = 7 \)
or
       \( a_{1} > x^{2n} + y^{2n} > x' \)
or
       \( \ip{A}{B} = \sum_{i} a_{i} b_{i} \).
The spaces you type in a formula are
ignored.  Remember that a letter like
       $x$                   % $ ... $  and  \( ... \)  are equivalent
is a formula when it denotes a mathematical
symbol, and it should be typed as one.

\section{Displayed Text}

Text is displayed by indenting it from the left
margin.  Quotations are commonly displayed.  There
are short quotations
\begin{quote}
   This is a short quotation.  It consists of a
   single paragraph of text.  See how it is formatted.
\end{quote}
and longer ones.
\begin{quotation}
   This is a longer quotation.  It consists of two
   paragraphs of text, neither of which are
   particularly interesting.

   This is the second paragraph of the quotation.  It
   is just as dull as the first paragraph.
\end{quotation}
Another frequently-displayed structure is a list.
The following is an example of an \emph{itemized}
list.
\begin{itemize}
   \item This is the first item of an itemized list.
         Each item in the list is marked with a ``tick''.
         You don't have to worry about what kind of tick
         mark is used.

   \item This is the second item of the list.  It
         contains another list nested inside it.  The inner
         list is an \emph{enumerated} list.
         \begin{enumerate}
            \item This is the first item of an enumerated
                  list that is nested within the itemized list.

            \item This is the second item of the inner list.
                  \LaTeX\ allows you to nest lists deeper than
                  you really should.
         \end{enumerate}
         This is the rest of the second item of the outer
         list.  It is no more interesting than any other
         part of the item.
   \item This is the third item of the list.
\end{itemize}
You can even display poetry.
\begin{verse}
   There is an environment
    for verse \\             % The \\ command separates lines
   Whose features some poets % within a stanza.
   will curse.

                             % One or more blank lines separate stanzas.

   For instead of making\\
   Them do \emph{all} line breaking, \\
   It allows them to put too many words on a line when they'd rather be
   forced to be terse.
\end{verse}

Mathematical formulas may also be displayed.  A
displayed formula
is
one-line long; multiline
formulas require special formatting instructions.
   \[  \ip{\Gamma}{\psi'} = x'' + y^{2} + z_{i}^{n}\]
Don't start a paragraph with a displayed equation,
nor make one a paragraph by itself.

\end{document}               % End of document.

\section{Bench Huge Section 307}
% synthetic bench section 307
% This is a sample LaTeX input file.  (Version of 12 August 2004.)
%
% A '%' character causes TeX to ignore all remaining text on the line,
% and is used for comments like this one.

\documentclass{article}      % Specifies the document class

                             % The preamble begins here.
\title{An Example Document}  % Declares the document's title.
\author{Leslie Lamport}      % Declares the author's name.
\date{January 21, 1994}      % Deleting this command produces today's date.

\newcommand{\ip}[2]{(#1, #2)}
                             % Defines \ip{arg1}{arg2} to mean
                             % (arg1, arg2).

%\newcommand{\ip}[2]{\langle #1 | #2\rangle}
                             % This is an alternative definition of
                             % \ip that is commented out.

\begin{document}             % End of preamble and beginning of text.

\maketitle                   % Produces the title.

This is an example input file.  Comparing it with
the output it generates can show you how to
produce a simple document of your own.

\section{Ordinary Text}      % Produces section heading.  Lower-level
                             % sections are begun with similar
                             % \subsection and \subsubsection commands.

The ends  of words and sentences are marked
  by   spaces. It  doesn't matter how many
spaces    you type; one is as good as 100.  The
end of   a line counts as a space.

One   or more   blank lines denote the  end
of  a paragraph.

Since any number of consecutive spaces are treated
like a single one, the formatting of the input
file makes no difference to
      \LaTeX,                % The \LaTeX command generates the LaTeX logo.
but it makes a difference to you.  When you use
\LaTeX, making your input file as easy to read
as possible will be a great help as you write
your document and when you change it.  This sample
file shows how you can add comments to your own input
file.

Because printing is different from typewriting,
there are a number of things that you have to do
differently when preparing an input file than if
you were just typing the document directly.
Quotation marks like
       ``this''
have to be handled specially, as do quotes within
quotes:
       ``\,`this'            % \, separates the double and single quote.
        is what I just
        wrote, not  `that'\,''.

Dashes come in three sizes: an
       intra-word
dash, a medium dash for number ranges like
       1--2,
and a punctuation
       dash---like
this.

A sentence-ending space should be larger than the
space between words within a sentence.  You
sometimes have to type special commands in
conjunction with punctuation characters to get
this right, as in the following sentence.
       Gnats, gnus, etc.\ all  % `\ ' makes an inter-word space.
       begin with G\@.         % \@ marks end-of-sentence punctuation.
You should check the spaces after periods when
reading your output to make sure you haven't
forgotten any special cases.  Generating an
ellipsis
       \ldots\               % `\ ' is needed after `\ldots' because TeX
                             % ignores spaces after command names like \ldots
                             % made from \ + letters.
                             %
                             % Note how a `%' character causes TeX to ignore
                             % the end of the input line, so these blank lines
                             % do not start a new paragraph.
                             %
with the right spacing around the periods requires
a special command.

\LaTeX\ interprets some common characters as
commands, so you must type special commands to
generate them.  These characters include the
following:
       \$ \& \% \# \{ and \}.

In printing, text is usually emphasized with an
       \emph{italic}
type style.

\begin{em}
   A long segment of text can also be emphasized
   in this way.  Text within such a segment can be
   given \emph{additional} emphasis.
\end{em}

It is sometimes necessary to prevent \LaTeX\ from
breaking a line where it might otherwise do so.
This may be at a space, as between the ``Mr.''\ and
``Jones'' in
       ``Mr.~Jones'',        % ~ produces an unbreakable interword space.
or within a word---especially when the word is a
symbol like
       \mbox{\emph{itemnum}}
that makes little sense when hyphenated across
lines.

Footnotes\footnote{This is an example of a footnote.}
pose no problem.

\LaTeX\ is good at typesetting mathematical formulas
like
       \( x-3y + z = 7 \)
or
       \( a_{1} > x^{2n} + y^{2n} > x' \)
or
       \( \ip{A}{B} = \sum_{i} a_{i} b_{i} \).
The spaces you type in a formula are
ignored.  Remember that a letter like
       $x$                   % $ ... $  and  \( ... \)  are equivalent
is a formula when it denotes a mathematical
symbol, and it should be typed as one.

\section{Displayed Text}

Text is displayed by indenting it from the left
margin.  Quotations are commonly displayed.  There
are short quotations
\begin{quote}
   This is a short quotation.  It consists of a
   single paragraph of text.  See how it is formatted.
\end{quote}
and longer ones.
\begin{quotation}
   This is a longer quotation.  It consists of two
   paragraphs of text, neither of which are
   particularly interesting.

   This is the second paragraph of the quotation.  It
   is just as dull as the first paragraph.
\end{quotation}
Another frequently-displayed structure is a list.
The following is an example of an \emph{itemized}
list.
\begin{itemize}
   \item This is the first item of an itemized list.
         Each item in the list is marked with a ``tick''.
         You don't have to worry about what kind of tick
         mark is used.

   \item This is the second item of the list.  It
         contains another list nested inside it.  The inner
         list is an \emph{enumerated} list.
         \begin{enumerate}
            \item This is the first item of an enumerated
                  list that is nested within the itemized list.

            \item This is the second item of the inner list.
                  \LaTeX\ allows you to nest lists deeper than
                  you really should.
         \end{enumerate}
         This is the rest of the second item of the outer
         list.  It is no more interesting than any other
         part of the item.
   \item This is the third item of the list.
\end{itemize}
You can even display poetry.
\begin{verse}
   There is an environment
    for verse \\             % The \\ command separates lines
   Whose features some poets % within a stanza.
   will curse.

                             % One or more blank lines separate stanzas.

   For instead of making\\
   Them do \emph{all} line breaking, \\
   It allows them to put too many words on a line when they'd rather be
   forced to be terse.
\end{verse}

Mathematical formulas may also be displayed.  A
displayed formula
is
one-line long; multiline
formulas require special formatting instructions.
   \[  \ip{\Gamma}{\psi'} = x'' + y^{2} + z_{i}^{n}\]
Don't start a paragraph with a displayed equation,
nor make one a paragraph by itself.

\end{document}               % End of document.

\section{Bench Huge Section 308}
% synthetic bench section 308
% This is a sample LaTeX input file.  (Version of 12 August 2004.)
%
% A '%' character causes TeX to ignore all remaining text on the line,
% and is used for comments like this one.

\documentclass{article}      % Specifies the document class

                             % The preamble begins here.
\title{An Example Document}  % Declares the document's title.
\author{Leslie Lamport}      % Declares the author's name.
\date{January 21, 1994}      % Deleting this command produces today's date.

\newcommand{\ip}[2]{(#1, #2)}
                             % Defines \ip{arg1}{arg2} to mean
                             % (arg1, arg2).

%\newcommand{\ip}[2]{\langle #1 | #2\rangle}
                             % This is an alternative definition of
                             % \ip that is commented out.

\begin{document}             % End of preamble and beginning of text.

\maketitle                   % Produces the title.

This is an example input file.  Comparing it with
the output it generates can show you how to
produce a simple document of your own.

\section{Ordinary Text}      % Produces section heading.  Lower-level
                             % sections are begun with similar
                             % \subsection and \subsubsection commands.

The ends  of words and sentences are marked
  by   spaces. It  doesn't matter how many
spaces    you type; one is as good as 100.  The
end of   a line counts as a space.

One   or more   blank lines denote the  end
of  a paragraph.

Since any number of consecutive spaces are treated
like a single one, the formatting of the input
file makes no difference to
      \LaTeX,                % The \LaTeX command generates the LaTeX logo.
but it makes a difference to you.  When you use
\LaTeX, making your input file as easy to read
as possible will be a great help as you write
your document and when you change it.  This sample
file shows how you can add comments to your own input
file.

Because printing is different from typewriting,
there are a number of things that you have to do
differently when preparing an input file than if
you were just typing the document directly.
Quotation marks like
       ``this''
have to be handled specially, as do quotes within
quotes:
       ``\,`this'            % \, separates the double and single quote.
        is what I just
        wrote, not  `that'\,''.

Dashes come in three sizes: an
       intra-word
dash, a medium dash for number ranges like
       1--2,
and a punctuation
       dash---like
this.

A sentence-ending space should be larger than the
space between words within a sentence.  You
sometimes have to type special commands in
conjunction with punctuation characters to get
this right, as in the following sentence.
       Gnats, gnus, etc.\ all  % `\ ' makes an inter-word space.
       begin with G\@.         % \@ marks end-of-sentence punctuation.
You should check the spaces after periods when
reading your output to make sure you haven't
forgotten any special cases.  Generating an
ellipsis
       \ldots\               % `\ ' is needed after `\ldots' because TeX
                             % ignores spaces after command names like \ldots
                             % made from \ + letters.
                             %
                             % Note how a `%' character causes TeX to ignore
                             % the end of the input line, so these blank lines
                             % do not start a new paragraph.
                             %
with the right spacing around the periods requires
a special command.

\LaTeX\ interprets some common characters as
commands, so you must type special commands to
generate them.  These characters include the
following:
       \$ \& \% \# \{ and \}.

In printing, text is usually emphasized with an
       \emph{italic}
type style.

\begin{em}
   A long segment of text can also be emphasized
   in this way.  Text within such a segment can be
   given \emph{additional} emphasis.
\end{em}

It is sometimes necessary to prevent \LaTeX\ from
breaking a line where it might otherwise do so.
This may be at a space, as between the ``Mr.''\ and
``Jones'' in
       ``Mr.~Jones'',        % ~ produces an unbreakable interword space.
or within a word---especially when the word is a
symbol like
       \mbox{\emph{itemnum}}
that makes little sense when hyphenated across
lines.

Footnotes\footnote{This is an example of a footnote.}
pose no problem.

\LaTeX\ is good at typesetting mathematical formulas
like
       \( x-3y + z = 7 \)
or
       \( a_{1} > x^{2n} + y^{2n} > x' \)
or
       \( \ip{A}{B} = \sum_{i} a_{i} b_{i} \).
The spaces you type in a formula are
ignored.  Remember that a letter like
       $x$                   % $ ... $  and  \( ... \)  are equivalent
is a formula when it denotes a mathematical
symbol, and it should be typed as one.

\section{Displayed Text}

Text is displayed by indenting it from the left
margin.  Quotations are commonly displayed.  There
are short quotations
\begin{quote}
   This is a short quotation.  It consists of a
   single paragraph of text.  See how it is formatted.
\end{quote}
and longer ones.
\begin{quotation}
   This is a longer quotation.  It consists of two
   paragraphs of text, neither of which are
   particularly interesting.

   This is the second paragraph of the quotation.  It
   is just as dull as the first paragraph.
\end{quotation}
Another frequently-displayed structure is a list.
The following is an example of an \emph{itemized}
list.
\begin{itemize}
   \item This is the first item of an itemized list.
         Each item in the list is marked with a ``tick''.
         You don't have to worry about what kind of tick
         mark is used.

   \item This is the second item of the list.  It
         contains another list nested inside it.  The inner
         list is an \emph{enumerated} list.
         \begin{enumerate}
            \item This is the first item of an enumerated
                  list that is nested within the itemized list.

            \item This is the second item of the inner list.
                  \LaTeX\ allows you to nest lists deeper than
                  you really should.
         \end{enumerate}
         This is the rest of the second item of the outer
         list.  It is no more interesting than any other
         part of the item.
   \item This is the third item of the list.
\end{itemize}
You can even display poetry.
\begin{verse}
   There is an environment
    for verse \\             % The \\ command separates lines
   Whose features some poets % within a stanza.
   will curse.

                             % One or more blank lines separate stanzas.

   For instead of making\\
   Them do \emph{all} line breaking, \\
   It allows them to put too many words on a line when they'd rather be
   forced to be terse.
\end{verse}

Mathematical formulas may also be displayed.  A
displayed formula
is
one-line long; multiline
formulas require special formatting instructions.
   \[  \ip{\Gamma}{\psi'} = x'' + y^{2} + z_{i}^{n}\]
Don't start a paragraph with a displayed equation,
nor make one a paragraph by itself.

\end{document}               % End of document.

\section{Bench Huge Section 309}
% synthetic bench section 309
% This is a sample LaTeX input file.  (Version of 12 August 2004.)
%
% A '%' character causes TeX to ignore all remaining text on the line,
% and is used for comments like this one.

\documentclass{article}      % Specifies the document class

                             % The preamble begins here.
\title{An Example Document}  % Declares the document's title.
\author{Leslie Lamport}      % Declares the author's name.
\date{January 21, 1994}      % Deleting this command produces today's date.

\newcommand{\ip}[2]{(#1, #2)}
                             % Defines \ip{arg1}{arg2} to mean
                             % (arg1, arg2).

%\newcommand{\ip}[2]{\langle #1 | #2\rangle}
                             % This is an alternative definition of
                             % \ip that is commented out.

\begin{document}             % End of preamble and beginning of text.

\maketitle                   % Produces the title.

This is an example input file.  Comparing it with
the output it generates can show you how to
produce a simple document of your own.

\section{Ordinary Text}      % Produces section heading.  Lower-level
                             % sections are begun with similar
                             % \subsection and \subsubsection commands.

The ends  of words and sentences are marked
  by   spaces. It  doesn't matter how many
spaces    you type; one is as good as 100.  The
end of   a line counts as a space.

One   or more   blank lines denote the  end
of  a paragraph.

Since any number of consecutive spaces are treated
like a single one, the formatting of the input
file makes no difference to
      \LaTeX,                % The \LaTeX command generates the LaTeX logo.
but it makes a difference to you.  When you use
\LaTeX, making your input file as easy to read
as possible will be a great help as you write
your document and when you change it.  This sample
file shows how you can add comments to your own input
file.

Because printing is different from typewriting,
there are a number of things that you have to do
differently when preparing an input file than if
you were just typing the document directly.
Quotation marks like
       ``this''
have to be handled specially, as do quotes within
quotes:
       ``\,`this'            % \, separates the double and single quote.
        is what I just
        wrote, not  `that'\,''.

Dashes come in three sizes: an
       intra-word
dash, a medium dash for number ranges like
       1--2,
and a punctuation
       dash---like
this.

A sentence-ending space should be larger than the
space between words within a sentence.  You
sometimes have to type special commands in
conjunction with punctuation characters to get
this right, as in the following sentence.
       Gnats, gnus, etc.\ all  % `\ ' makes an inter-word space.
       begin with G\@.         % \@ marks end-of-sentence punctuation.
You should check the spaces after periods when
reading your output to make sure you haven't
forgotten any special cases.  Generating an
ellipsis
       \ldots\               % `\ ' is needed after `\ldots' because TeX
                             % ignores spaces after command names like \ldots
                             % made from \ + letters.
                             %
                             % Note how a `%' character causes TeX to ignore
                             % the end of the input line, so these blank lines
                             % do not start a new paragraph.
                             %
with the right spacing around the periods requires
a special command.

\LaTeX\ interprets some common characters as
commands, so you must type special commands to
generate them.  These characters include the
following:
       \$ \& \% \# \{ and \}.

In printing, text is usually emphasized with an
       \emph{italic}
type style.

\begin{em}
   A long segment of text can also be emphasized
   in this way.  Text within such a segment can be
   given \emph{additional} emphasis.
\end{em}

It is sometimes necessary to prevent \LaTeX\ from
breaking a line where it might otherwise do so.
This may be at a space, as between the ``Mr.''\ and
``Jones'' in
       ``Mr.~Jones'',        % ~ produces an unbreakable interword space.
or within a word---especially when the word is a
symbol like
       \mbox{\emph{itemnum}}
that makes little sense when hyphenated across
lines.

Footnotes\footnote{This is an example of a footnote.}
pose no problem.

\LaTeX\ is good at typesetting mathematical formulas
like
       \( x-3y + z = 7 \)
or
       \( a_{1} > x^{2n} + y^{2n} > x' \)
or
       \( \ip{A}{B} = \sum_{i} a_{i} b_{i} \).
The spaces you type in a formula are
ignored.  Remember that a letter like
       $x$                   % $ ... $  and  \( ... \)  are equivalent
is a formula when it denotes a mathematical
symbol, and it should be typed as one.

\section{Displayed Text}

Text is displayed by indenting it from the left
margin.  Quotations are commonly displayed.  There
are short quotations
\begin{quote}
   This is a short quotation.  It consists of a
   single paragraph of text.  See how it is formatted.
\end{quote}
and longer ones.
\begin{quotation}
   This is a longer quotation.  It consists of two
   paragraphs of text, neither of which are
   particularly interesting.

   This is the second paragraph of the quotation.  It
   is just as dull as the first paragraph.
\end{quotation}
Another frequently-displayed structure is a list.
The following is an example of an \emph{itemized}
list.
\begin{itemize}
   \item This is the first item of an itemized list.
         Each item in the list is marked with a ``tick''.
         You don't have to worry about what kind of tick
         mark is used.

   \item This is the second item of the list.  It
         contains another list nested inside it.  The inner
         list is an \emph{enumerated} list.
         \begin{enumerate}
            \item This is the first item of an enumerated
                  list that is nested within the itemized list.

            \item This is the second item of the inner list.
                  \LaTeX\ allows you to nest lists deeper than
                  you really should.
         \end{enumerate}
         This is the rest of the second item of the outer
         list.  It is no more interesting than any other
         part of the item.
   \item This is the third item of the list.
\end{itemize}
You can even display poetry.
\begin{verse}
   There is an environment
    for verse \\             % The \\ command separates lines
   Whose features some poets % within a stanza.
   will curse.

                             % One or more blank lines separate stanzas.

   For instead of making\\
   Them do \emph{all} line breaking, \\
   It allows them to put too many words on a line when they'd rather be
   forced to be terse.
\end{verse}

Mathematical formulas may also be displayed.  A
displayed formula
is
one-line long; multiline
formulas require special formatting instructions.
   \[  \ip{\Gamma}{\psi'} = x'' + y^{2} + z_{i}^{n}\]
Don't start a paragraph with a displayed equation,
nor make one a paragraph by itself.

\end{document}               % End of document.

\section{Bench Huge Section 310}
% synthetic bench section 310
% This is a sample LaTeX input file.  (Version of 12 August 2004.)
%
% A '%' character causes TeX to ignore all remaining text on the line,
% and is used for comments like this one.

\documentclass{article}      % Specifies the document class

                             % The preamble begins here.
\title{An Example Document}  % Declares the document's title.
\author{Leslie Lamport}      % Declares the author's name.
\date{January 21, 1994}      % Deleting this command produces today's date.

\newcommand{\ip}[2]{(#1, #2)}
                             % Defines \ip{arg1}{arg2} to mean
                             % (arg1, arg2).

%\newcommand{\ip}[2]{\langle #1 | #2\rangle}
                             % This is an alternative definition of
                             % \ip that is commented out.

\begin{document}             % End of preamble and beginning of text.

\maketitle                   % Produces the title.

This is an example input file.  Comparing it with
the output it generates can show you how to
produce a simple document of your own.

\section{Ordinary Text}      % Produces section heading.  Lower-level
                             % sections are begun with similar
                             % \subsection and \subsubsection commands.

The ends  of words and sentences are marked
  by   spaces. It  doesn't matter how many
spaces    you type; one is as good as 100.  The
end of   a line counts as a space.

One   or more   blank lines denote the  end
of  a paragraph.

Since any number of consecutive spaces are treated
like a single one, the formatting of the input
file makes no difference to
      \LaTeX,                % The \LaTeX command generates the LaTeX logo.
but it makes a difference to you.  When you use
\LaTeX, making your input file as easy to read
as possible will be a great help as you write
your document and when you change it.  This sample
file shows how you can add comments to your own input
file.

Because printing is different from typewriting,
there are a number of things that you have to do
differently when preparing an input file than if
you were just typing the document directly.
Quotation marks like
       ``this''
have to be handled specially, as do quotes within
quotes:
       ``\,`this'            % \, separates the double and single quote.
        is what I just
        wrote, not  `that'\,''.

Dashes come in three sizes: an
       intra-word
dash, a medium dash for number ranges like
       1--2,
and a punctuation
       dash---like
this.

A sentence-ending space should be larger than the
space between words within a sentence.  You
sometimes have to type special commands in
conjunction with punctuation characters to get
this right, as in the following sentence.
       Gnats, gnus, etc.\ all  % `\ ' makes an inter-word space.
       begin with G\@.         % \@ marks end-of-sentence punctuation.
You should check the spaces after periods when
reading your output to make sure you haven't
forgotten any special cases.  Generating an
ellipsis
       \ldots\               % `\ ' is needed after `\ldots' because TeX
                             % ignores spaces after command names like \ldots
                             % made from \ + letters.
                             %
                             % Note how a `%' character causes TeX to ignore
                             % the end of the input line, so these blank lines
                             % do not start a new paragraph.
                             %
with the right spacing around the periods requires
a special command.

\LaTeX\ interprets some common characters as
commands, so you must type special commands to
generate them.  These characters include the
following:
       \$ \& \% \# \{ and \}.

In printing, text is usually emphasized with an
       \emph{italic}
type style.

\begin{em}
   A long segment of text can also be emphasized
   in this way.  Text within such a segment can be
   given \emph{additional} emphasis.
\end{em}

It is sometimes necessary to prevent \LaTeX\ from
breaking a line where it might otherwise do so.
This may be at a space, as between the ``Mr.''\ and
``Jones'' in
       ``Mr.~Jones'',        % ~ produces an unbreakable interword space.
or within a word---especially when the word is a
symbol like
       \mbox{\emph{itemnum}}
that makes little sense when hyphenated across
lines.

Footnotes\footnote{This is an example of a footnote.}
pose no problem.

\LaTeX\ is good at typesetting mathematical formulas
like
       \( x-3y + z = 7 \)
or
       \( a_{1} > x^{2n} + y^{2n} > x' \)
or
       \( \ip{A}{B} = \sum_{i} a_{i} b_{i} \).
The spaces you type in a formula are
ignored.  Remember that a letter like
       $x$                   % $ ... $  and  \( ... \)  are equivalent
is a formula when it denotes a mathematical
symbol, and it should be typed as one.

\section{Displayed Text}

Text is displayed by indenting it from the left
margin.  Quotations are commonly displayed.  There
are short quotations
\begin{quote}
   This is a short quotation.  It consists of a
   single paragraph of text.  See how it is formatted.
\end{quote}
and longer ones.
\begin{quotation}
   This is a longer quotation.  It consists of two
   paragraphs of text, neither of which are
   particularly interesting.

   This is the second paragraph of the quotation.  It
   is just as dull as the first paragraph.
\end{quotation}
Another frequently-displayed structure is a list.
The following is an example of an \emph{itemized}
list.
\begin{itemize}
   \item This is the first item of an itemized list.
         Each item in the list is marked with a ``tick''.
         You don't have to worry about what kind of tick
         mark is used.

   \item This is the second item of the list.  It
         contains another list nested inside it.  The inner
         list is an \emph{enumerated} list.
         \begin{enumerate}
            \item This is the first item of an enumerated
                  list that is nested within the itemized list.

            \item This is the second item of the inner list.
                  \LaTeX\ allows you to nest lists deeper than
                  you really should.
         \end{enumerate}
         This is the rest of the second item of the outer
         list.  It is no more interesting than any other
         part of the item.
   \item This is the third item of the list.
\end{itemize}
You can even display poetry.
\begin{verse}
   There is an environment
    for verse \\             % The \\ command separates lines
   Whose features some poets % within a stanza.
   will curse.

                             % One or more blank lines separate stanzas.

   For instead of making\\
   Them do \emph{all} line breaking, \\
   It allows them to put too many words on a line when they'd rather be
   forced to be terse.
\end{verse}

Mathematical formulas may also be displayed.  A
displayed formula
is
one-line long; multiline
formulas require special formatting instructions.
   \[  \ip{\Gamma}{\psi'} = x'' + y^{2} + z_{i}^{n}\]
Don't start a paragraph with a displayed equation,
nor make one a paragraph by itself.

\end{document}               % End of document.

\section{Bench Huge Section 311}
% synthetic bench section 311
% This is a sample LaTeX input file.  (Version of 12 August 2004.)
%
% A '%' character causes TeX to ignore all remaining text on the line,
% and is used for comments like this one.

\documentclass{article}      % Specifies the document class

                             % The preamble begins here.
\title{An Example Document}  % Declares the document's title.
\author{Leslie Lamport}      % Declares the author's name.
\date{January 21, 1994}      % Deleting this command produces today's date.

\newcommand{\ip}[2]{(#1, #2)}
                             % Defines \ip{arg1}{arg2} to mean
                             % (arg1, arg2).

%\newcommand{\ip}[2]{\langle #1 | #2\rangle}
                             % This is an alternative definition of
                             % \ip that is commented out.

\begin{document}             % End of preamble and beginning of text.

\maketitle                   % Produces the title.

This is an example input file.  Comparing it with
the output it generates can show you how to
produce a simple document of your own.

\section{Ordinary Text}      % Produces section heading.  Lower-level
                             % sections are begun with similar
                             % \subsection and \subsubsection commands.

The ends  of words and sentences are marked
  by   spaces. It  doesn't matter how many
spaces    you type; one is as good as 100.  The
end of   a line counts as a space.

One   or more   blank lines denote the  end
of  a paragraph.

Since any number of consecutive spaces are treated
like a single one, the formatting of the input
file makes no difference to
      \LaTeX,                % The \LaTeX command generates the LaTeX logo.
but it makes a difference to you.  When you use
\LaTeX, making your input file as easy to read
as possible will be a great help as you write
your document and when you change it.  This sample
file shows how you can add comments to your own input
file.

Because printing is different from typewriting,
there are a number of things that you have to do
differently when preparing an input file than if
you were just typing the document directly.
Quotation marks like
       ``this''
have to be handled specially, as do quotes within
quotes:
       ``\,`this'            % \, separates the double and single quote.
        is what I just
        wrote, not  `that'\,''.

Dashes come in three sizes: an
       intra-word
dash, a medium dash for number ranges like
       1--2,
and a punctuation
       dash---like
this.

A sentence-ending space should be larger than the
space between words within a sentence.  You
sometimes have to type special commands in
conjunction with punctuation characters to get
this right, as in the following sentence.
       Gnats, gnus, etc.\ all  % `\ ' makes an inter-word space.
       begin with G\@.         % \@ marks end-of-sentence punctuation.
You should check the spaces after periods when
reading your output to make sure you haven't
forgotten any special cases.  Generating an
ellipsis
       \ldots\               % `\ ' is needed after `\ldots' because TeX
                             % ignores spaces after command names like \ldots
                             % made from \ + letters.
                             %
                             % Note how a `%' character causes TeX to ignore
                             % the end of the input line, so these blank lines
                             % do not start a new paragraph.
                             %
with the right spacing around the periods requires
a special command.

\LaTeX\ interprets some common characters as
commands, so you must type special commands to
generate them.  These characters include the
following:
       \$ \& \% \# \{ and \}.

In printing, text is usually emphasized with an
       \emph{italic}
type style.

\begin{em}
   A long segment of text can also be emphasized
   in this way.  Text within such a segment can be
   given \emph{additional} emphasis.
\end{em}

It is sometimes necessary to prevent \LaTeX\ from
breaking a line where it might otherwise do so.
This may be at a space, as between the ``Mr.''\ and
``Jones'' in
       ``Mr.~Jones'',        % ~ produces an unbreakable interword space.
or within a word---especially when the word is a
symbol like
       \mbox{\emph{itemnum}}
that makes little sense when hyphenated across
lines.

Footnotes\footnote{This is an example of a footnote.}
pose no problem.

\LaTeX\ is good at typesetting mathematical formulas
like
       \( x-3y + z = 7 \)
or
       \( a_{1} > x^{2n} + y^{2n} > x' \)
or
       \( \ip{A}{B} = \sum_{i} a_{i} b_{i} \).
The spaces you type in a formula are
ignored.  Remember that a letter like
       $x$                   % $ ... $  and  \( ... \)  are equivalent
is a formula when it denotes a mathematical
symbol, and it should be typed as one.

\section{Displayed Text}

Text is displayed by indenting it from the left
margin.  Quotations are commonly displayed.  There
are short quotations
\begin{quote}
   This is a short quotation.  It consists of a
   single paragraph of text.  See how it is formatted.
\end{quote}
and longer ones.
\begin{quotation}
   This is a longer quotation.  It consists of two
   paragraphs of text, neither of which are
   particularly interesting.

   This is the second paragraph of the quotation.  It
   is just as dull as the first paragraph.
\end{quotation}
Another frequently-displayed structure is a list.
The following is an example of an \emph{itemized}
list.
\begin{itemize}
   \item This is the first item of an itemized list.
         Each item in the list is marked with a ``tick''.
         You don't have to worry about what kind of tick
         mark is used.

   \item This is the second item of the list.  It
         contains another list nested inside it.  The inner
         list is an \emph{enumerated} list.
         \begin{enumerate}
            \item This is the first item of an enumerated
                  list that is nested within the itemized list.

            \item This is the second item of the inner list.
                  \LaTeX\ allows you to nest lists deeper than
                  you really should.
         \end{enumerate}
         This is the rest of the second item of the outer
         list.  It is no more interesting than any other
         part of the item.
   \item This is the third item of the list.
\end{itemize}
You can even display poetry.
\begin{verse}
   There is an environment
    for verse \\             % The \\ command separates lines
   Whose features some poets % within a stanza.
   will curse.

                             % One or more blank lines separate stanzas.

   For instead of making\\
   Them do \emph{all} line breaking, \\
   It allows them to put too many words on a line when they'd rather be
   forced to be terse.
\end{verse}

Mathematical formulas may also be displayed.  A
displayed formula
is
one-line long; multiline
formulas require special formatting instructions.
   \[  \ip{\Gamma}{\psi'} = x'' + y^{2} + z_{i}^{n}\]
Don't start a paragraph with a displayed equation,
nor make one a paragraph by itself.

\end{document}               % End of document.

\section{Bench Huge Section 312}
% synthetic bench section 312
% This is a sample LaTeX input file.  (Version of 12 August 2004.)
%
% A '%' character causes TeX to ignore all remaining text on the line,
% and is used for comments like this one.

\documentclass{article}      % Specifies the document class

                             % The preamble begins here.
\title{An Example Document}  % Declares the document's title.
\author{Leslie Lamport}      % Declares the author's name.
\date{January 21, 1994}      % Deleting this command produces today's date.

\newcommand{\ip}[2]{(#1, #2)}
                             % Defines \ip{arg1}{arg2} to mean
                             % (arg1, arg2).

%\newcommand{\ip}[2]{\langle #1 | #2\rangle}
                             % This is an alternative definition of
                             % \ip that is commented out.

\begin{document}             % End of preamble and beginning of text.

\maketitle                   % Produces the title.

This is an example input file.  Comparing it with
the output it generates can show you how to
produce a simple document of your own.

\section{Ordinary Text}      % Produces section heading.  Lower-level
                             % sections are begun with similar
                             % \subsection and \subsubsection commands.

The ends  of words and sentences are marked
  by   spaces. It  doesn't matter how many
spaces    you type; one is as good as 100.  The
end of   a line counts as a space.

One   or more   blank lines denote the  end
of  a paragraph.

Since any number of consecutive spaces are treated
like a single one, the formatting of the input
file makes no difference to
      \LaTeX,                % The \LaTeX command generates the LaTeX logo.
but it makes a difference to you.  When you use
\LaTeX, making your input file as easy to read
as possible will be a great help as you write
your document and when you change it.  This sample
file shows how you can add comments to your own input
file.

Because printing is different from typewriting,
there are a number of things that you have to do
differently when preparing an input file than if
you were just typing the document directly.
Quotation marks like
       ``this''
have to be handled specially, as do quotes within
quotes:
       ``\,`this'            % \, separates the double and single quote.
        is what I just
        wrote, not  `that'\,''.

Dashes come in three sizes: an
       intra-word
dash, a medium dash for number ranges like
       1--2,
and a punctuation
       dash---like
this.

A sentence-ending space should be larger than the
space between words within a sentence.  You
sometimes have to type special commands in
conjunction with punctuation characters to get
this right, as in the following sentence.
       Gnats, gnus, etc.\ all  % `\ ' makes an inter-word space.
       begin with G\@.         % \@ marks end-of-sentence punctuation.
You should check the spaces after periods when
reading your output to make sure you haven't
forgotten any special cases.  Generating an
ellipsis
       \ldots\               % `\ ' is needed after `\ldots' because TeX
                             % ignores spaces after command names like \ldots
                             % made from \ + letters.
                             %
                             % Note how a `%' character causes TeX to ignore
                             % the end of the input line, so these blank lines
                             % do not start a new paragraph.
                             %
with the right spacing around the periods requires
a special command.

\LaTeX\ interprets some common characters as
commands, so you must type special commands to
generate them.  These characters include the
following:
       \$ \& \% \# \{ and \}.

In printing, text is usually emphasized with an
       \emph{italic}
type style.

\begin{em}
   A long segment of text can also be emphasized
   in this way.  Text within such a segment can be
   given \emph{additional} emphasis.
\end{em}

It is sometimes necessary to prevent \LaTeX\ from
breaking a line where it might otherwise do so.
This may be at a space, as between the ``Mr.''\ and
``Jones'' in
       ``Mr.~Jones'',        % ~ produces an unbreakable interword space.
or within a word---especially when the word is a
symbol like
       \mbox{\emph{itemnum}}
that makes little sense when hyphenated across
lines.

Footnotes\footnote{This is an example of a footnote.}
pose no problem.

\LaTeX\ is good at typesetting mathematical formulas
like
       \( x-3y + z = 7 \)
or
       \( a_{1} > x^{2n} + y^{2n} > x' \)
or
       \( \ip{A}{B} = \sum_{i} a_{i} b_{i} \).
The spaces you type in a formula are
ignored.  Remember that a letter like
       $x$                   % $ ... $  and  \( ... \)  are equivalent
is a formula when it denotes a mathematical
symbol, and it should be typed as one.

\section{Displayed Text}

Text is displayed by indenting it from the left
margin.  Quotations are commonly displayed.  There
are short quotations
\begin{quote}
   This is a short quotation.  It consists of a
   single paragraph of text.  See how it is formatted.
\end{quote}
and longer ones.
\begin{quotation}
   This is a longer quotation.  It consists of two
   paragraphs of text, neither of which are
   particularly interesting.

   This is the second paragraph of the quotation.  It
   is just as dull as the first paragraph.
\end{quotation}
Another frequently-displayed structure is a list.
The following is an example of an \emph{itemized}
list.
\begin{itemize}
   \item This is the first item of an itemized list.
         Each item in the list is marked with a ``tick''.
         You don't have to worry about what kind of tick
         mark is used.

   \item This is the second item of the list.  It
         contains another list nested inside it.  The inner
         list is an \emph{enumerated} list.
         \begin{enumerate}
            \item This is the first item of an enumerated
                  list that is nested within the itemized list.

            \item This is the second item of the inner list.
                  \LaTeX\ allows you to nest lists deeper than
                  you really should.
         \end{enumerate}
         This is the rest of the second item of the outer
         list.  It is no more interesting than any other
         part of the item.
   \item This is the third item of the list.
\end{itemize}
You can even display poetry.
\begin{verse}
   There is an environment
    for verse \\             % The \\ command separates lines
   Whose features some poets % within a stanza.
   will curse.

                             % One or more blank lines separate stanzas.

   For instead of making\\
   Them do \emph{all} line breaking, \\
   It allows them to put too many words on a line when they'd rather be
   forced to be terse.
\end{verse}

Mathematical formulas may also be displayed.  A
displayed formula
is
one-line long; multiline
formulas require special formatting instructions.
   \[  \ip{\Gamma}{\psi'} = x'' + y^{2} + z_{i}^{n}\]
Don't start a paragraph with a displayed equation,
nor make one a paragraph by itself.

\end{document}               % End of document.

\section{Bench Huge Section 313}
% synthetic bench section 313
% This is a sample LaTeX input file.  (Version of 12 August 2004.)
%
% A '%' character causes TeX to ignore all remaining text on the line,
% and is used for comments like this one.

\documentclass{article}      % Specifies the document class

                             % The preamble begins here.
\title{An Example Document}  % Declares the document's title.
\author{Leslie Lamport}      % Declares the author's name.
\date{January 21, 1994}      % Deleting this command produces today's date.

\newcommand{\ip}[2]{(#1, #2)}
                             % Defines \ip{arg1}{arg2} to mean
                             % (arg1, arg2).

%\newcommand{\ip}[2]{\langle #1 | #2\rangle}
                             % This is an alternative definition of
                             % \ip that is commented out.

\begin{document}             % End of preamble and beginning of text.

\maketitle                   % Produces the title.

This is an example input file.  Comparing it with
the output it generates can show you how to
produce a simple document of your own.

\section{Ordinary Text}      % Produces section heading.  Lower-level
                             % sections are begun with similar
                             % \subsection and \subsubsection commands.

The ends  of words and sentences are marked
  by   spaces. It  doesn't matter how many
spaces    you type; one is as good as 100.  The
end of   a line counts as a space.

One   or more   blank lines denote the  end
of  a paragraph.

Since any number of consecutive spaces are treated
like a single one, the formatting of the input
file makes no difference to
      \LaTeX,                % The \LaTeX command generates the LaTeX logo.
but it makes a difference to you.  When you use
\LaTeX, making your input file as easy to read
as possible will be a great help as you write
your document and when you change it.  This sample
file shows how you can add comments to your own input
file.

Because printing is different from typewriting,
there are a number of things that you have to do
differently when preparing an input file than if
you were just typing the document directly.
Quotation marks like
       ``this''
have to be handled specially, as do quotes within
quotes:
       ``\,`this'            % \, separates the double and single quote.
        is what I just
        wrote, not  `that'\,''.

Dashes come in three sizes: an
       intra-word
dash, a medium dash for number ranges like
       1--2,
and a punctuation
       dash---like
this.

A sentence-ending space should be larger than the
space between words within a sentence.  You
sometimes have to type special commands in
conjunction with punctuation characters to get
this right, as in the following sentence.
       Gnats, gnus, etc.\ all  % `\ ' makes an inter-word space.
       begin with G\@.         % \@ marks end-of-sentence punctuation.
You should check the spaces after periods when
reading your output to make sure you haven't
forgotten any special cases.  Generating an
ellipsis
       \ldots\               % `\ ' is needed after `\ldots' because TeX
                             % ignores spaces after command names like \ldots
                             % made from \ + letters.
                             %
                             % Note how a `%' character causes TeX to ignore
                             % the end of the input line, so these blank lines
                             % do not start a new paragraph.
                             %
with the right spacing around the periods requires
a special command.

\LaTeX\ interprets some common characters as
commands, so you must type special commands to
generate them.  These characters include the
following:
       \$ \& \% \# \{ and \}.

In printing, text is usually emphasized with an
       \emph{italic}
type style.

\begin{em}
   A long segment of text can also be emphasized
   in this way.  Text within such a segment can be
   given \emph{additional} emphasis.
\end{em}

It is sometimes necessary to prevent \LaTeX\ from
breaking a line where it might otherwise do so.
This may be at a space, as between the ``Mr.''\ and
``Jones'' in
       ``Mr.~Jones'',        % ~ produces an unbreakable interword space.
or within a word---especially when the word is a
symbol like
       \mbox{\emph{itemnum}}
that makes little sense when hyphenated across
lines.

Footnotes\footnote{This is an example of a footnote.}
pose no problem.

\LaTeX\ is good at typesetting mathematical formulas
like
       \( x-3y + z = 7 \)
or
       \( a_{1} > x^{2n} + y^{2n} > x' \)
or
       \( \ip{A}{B} = \sum_{i} a_{i} b_{i} \).
The spaces you type in a formula are
ignored.  Remember that a letter like
       $x$                   % $ ... $  and  \( ... \)  are equivalent
is a formula when it denotes a mathematical
symbol, and it should be typed as one.

\section{Displayed Text}

Text is displayed by indenting it from the left
margin.  Quotations are commonly displayed.  There
are short quotations
\begin{quote}
   This is a short quotation.  It consists of a
   single paragraph of text.  See how it is formatted.
\end{quote}
and longer ones.
\begin{quotation}
   This is a longer quotation.  It consists of two
   paragraphs of text, neither of which are
   particularly interesting.

   This is the second paragraph of the quotation.  It
   is just as dull as the first paragraph.
\end{quotation}
Another frequently-displayed structure is a list.
The following is an example of an \emph{itemized}
list.
\begin{itemize}
   \item This is the first item of an itemized list.
         Each item in the list is marked with a ``tick''.
         You don't have to worry about what kind of tick
         mark is used.

   \item This is the second item of the list.  It
         contains another list nested inside it.  The inner
         list is an \emph{enumerated} list.
         \begin{enumerate}
            \item This is the first item of an enumerated
                  list that is nested within the itemized list.

            \item This is the second item of the inner list.
                  \LaTeX\ allows you to nest lists deeper than
                  you really should.
         \end{enumerate}
         This is the rest of the second item of the outer
         list.  It is no more interesting than any other
         part of the item.
   \item This is the third item of the list.
\end{itemize}
You can even display poetry.
\begin{verse}
   There is an environment
    for verse \\             % The \\ command separates lines
   Whose features some poets % within a stanza.
   will curse.

                             % One or more blank lines separate stanzas.

   For instead of making\\
   Them do \emph{all} line breaking, \\
   It allows them to put too many words on a line when they'd rather be
   forced to be terse.
\end{verse}

Mathematical formulas may also be displayed.  A
displayed formula
is
one-line long; multiline
formulas require special formatting instructions.
   \[  \ip{\Gamma}{\psi'} = x'' + y^{2} + z_{i}^{n}\]
Don't start a paragraph with a displayed equation,
nor make one a paragraph by itself.

\end{document}               % End of document.

\section{Bench Huge Section 314}
% synthetic bench section 314
% This is a sample LaTeX input file.  (Version of 12 August 2004.)
%
% A '%' character causes TeX to ignore all remaining text on the line,
% and is used for comments like this one.

\documentclass{article}      % Specifies the document class

                             % The preamble begins here.
\title{An Example Document}  % Declares the document's title.
\author{Leslie Lamport}      % Declares the author's name.
\date{January 21, 1994}      % Deleting this command produces today's date.

\newcommand{\ip}[2]{(#1, #2)}
                             % Defines \ip{arg1}{arg2} to mean
                             % (arg1, arg2).

%\newcommand{\ip}[2]{\langle #1 | #2\rangle}
                             % This is an alternative definition of
                             % \ip that is commented out.

\begin{document}             % End of preamble and beginning of text.

\maketitle                   % Produces the title.

This is an example input file.  Comparing it with
the output it generates can show you how to
produce a simple document of your own.

\section{Ordinary Text}      % Produces section heading.  Lower-level
                             % sections are begun with similar
                             % \subsection and \subsubsection commands.

The ends  of words and sentences are marked
  by   spaces. It  doesn't matter how many
spaces    you type; one is as good as 100.  The
end of   a line counts as a space.

One   or more   blank lines denote the  end
of  a paragraph.

Since any number of consecutive spaces are treated
like a single one, the formatting of the input
file makes no difference to
      \LaTeX,                % The \LaTeX command generates the LaTeX logo.
but it makes a difference to you.  When you use
\LaTeX, making your input file as easy to read
as possible will be a great help as you write
your document and when you change it.  This sample
file shows how you can add comments to your own input
file.

Because printing is different from typewriting,
there are a number of things that you have to do
differently when preparing an input file than if
you were just typing the document directly.
Quotation marks like
       ``this''
have to be handled specially, as do quotes within
quotes:
       ``\,`this'            % \, separates the double and single quote.
        is what I just
        wrote, not  `that'\,''.

Dashes come in three sizes: an
       intra-word
dash, a medium dash for number ranges like
       1--2,
and a punctuation
       dash---like
this.

A sentence-ending space should be larger than the
space between words within a sentence.  You
sometimes have to type special commands in
conjunction with punctuation characters to get
this right, as in the following sentence.
       Gnats, gnus, etc.\ all  % `\ ' makes an inter-word space.
       begin with G\@.         % \@ marks end-of-sentence punctuation.
You should check the spaces after periods when
reading your output to make sure you haven't
forgotten any special cases.  Generating an
ellipsis
       \ldots\               % `\ ' is needed after `\ldots' because TeX
                             % ignores spaces after command names like \ldots
                             % made from \ + letters.
                             %
                             % Note how a `%' character causes TeX to ignore
                             % the end of the input line, so these blank lines
                             % do not start a new paragraph.
                             %
with the right spacing around the periods requires
a special command.

\LaTeX\ interprets some common characters as
commands, so you must type special commands to
generate them.  These characters include the
following:
       \$ \& \% \# \{ and \}.

In printing, text is usually emphasized with an
       \emph{italic}
type style.

\begin{em}
   A long segment of text can also be emphasized
   in this way.  Text within such a segment can be
   given \emph{additional} emphasis.
\end{em}

It is sometimes necessary to prevent \LaTeX\ from
breaking a line where it might otherwise do so.
This may be at a space, as between the ``Mr.''\ and
``Jones'' in
       ``Mr.~Jones'',        % ~ produces an unbreakable interword space.
or within a word---especially when the word is a
symbol like
       \mbox{\emph{itemnum}}
that makes little sense when hyphenated across
lines.

Footnotes\footnote{This is an example of a footnote.}
pose no problem.

\LaTeX\ is good at typesetting mathematical formulas
like
       \( x-3y + z = 7 \)
or
       \( a_{1} > x^{2n} + y^{2n} > x' \)
or
       \( \ip{A}{B} = \sum_{i} a_{i} b_{i} \).
The spaces you type in a formula are
ignored.  Remember that a letter like
       $x$                   % $ ... $  and  \( ... \)  are equivalent
is a formula when it denotes a mathematical
symbol, and it should be typed as one.

\section{Displayed Text}

Text is displayed by indenting it from the left
margin.  Quotations are commonly displayed.  There
are short quotations
\begin{quote}
   This is a short quotation.  It consists of a
   single paragraph of text.  See how it is formatted.
\end{quote}
and longer ones.
\begin{quotation}
   This is a longer quotation.  It consists of two
   paragraphs of text, neither of which are
   particularly interesting.

   This is the second paragraph of the quotation.  It
   is just as dull as the first paragraph.
\end{quotation}
Another frequently-displayed structure is a list.
The following is an example of an \emph{itemized}
list.
\begin{itemize}
   \item This is the first item of an itemized list.
         Each item in the list is marked with a ``tick''.
         You don't have to worry about what kind of tick
         mark is used.

   \item This is the second item of the list.  It
         contains another list nested inside it.  The inner
         list is an \emph{enumerated} list.
         \begin{enumerate}
            \item This is the first item of an enumerated
                  list that is nested within the itemized list.

            \item This is the second item of the inner list.
                  \LaTeX\ allows you to nest lists deeper than
                  you really should.
         \end{enumerate}
         This is the rest of the second item of the outer
         list.  It is no more interesting than any other
         part of the item.
   \item This is the third item of the list.
\end{itemize}
You can even display poetry.
\begin{verse}
   There is an environment
    for verse \\             % The \\ command separates lines
   Whose features some poets % within a stanza.
   will curse.

                             % One or more blank lines separate stanzas.

   For instead of making\\
   Them do \emph{all} line breaking, \\
   It allows them to put too many words on a line when they'd rather be
   forced to be terse.
\end{verse}

Mathematical formulas may also be displayed.  A
displayed formula
is
one-line long; multiline
formulas require special formatting instructions.
   \[  \ip{\Gamma}{\psi'} = x'' + y^{2} + z_{i}^{n}\]
Don't start a paragraph with a displayed equation,
nor make one a paragraph by itself.

\end{document}               % End of document.

\section{Bench Huge Section 315}
% synthetic bench section 315
% This is a sample LaTeX input file.  (Version of 12 August 2004.)
%
% A '%' character causes TeX to ignore all remaining text on the line,
% and is used for comments like this one.

\documentclass{article}      % Specifies the document class

                             % The preamble begins here.
\title{An Example Document}  % Declares the document's title.
\author{Leslie Lamport}      % Declares the author's name.
\date{January 21, 1994}      % Deleting this command produces today's date.

\newcommand{\ip}[2]{(#1, #2)}
                             % Defines \ip{arg1}{arg2} to mean
                             % (arg1, arg2).

%\newcommand{\ip}[2]{\langle #1 | #2\rangle}
                             % This is an alternative definition of
                             % \ip that is commented out.

\begin{document}             % End of preamble and beginning of text.

\maketitle                   % Produces the title.

This is an example input file.  Comparing it with
the output it generates can show you how to
produce a simple document of your own.

\section{Ordinary Text}      % Produces section heading.  Lower-level
                             % sections are begun with similar
                             % \subsection and \subsubsection commands.

The ends  of words and sentences are marked
  by   spaces. It  doesn't matter how many
spaces    you type; one is as good as 100.  The
end of   a line counts as a space.

One   or more   blank lines denote the  end
of  a paragraph.

Since any number of consecutive spaces are treated
like a single one, the formatting of the input
file makes no difference to
      \LaTeX,                % The \LaTeX command generates the LaTeX logo.
but it makes a difference to you.  When you use
\LaTeX, making your input file as easy to read
as possible will be a great help as you write
your document and when you change it.  This sample
file shows how you can add comments to your own input
file.

Because printing is different from typewriting,
there are a number of things that you have to do
differently when preparing an input file than if
you were just typing the document directly.
Quotation marks like
       ``this''
have to be handled specially, as do quotes within
quotes:
       ``\,`this'            % \, separates the double and single quote.
        is what I just
        wrote, not  `that'\,''.

Dashes come in three sizes: an
       intra-word
dash, a medium dash for number ranges like
       1--2,
and a punctuation
       dash---like
this.

A sentence-ending space should be larger than the
space between words within a sentence.  You
sometimes have to type special commands in
conjunction with punctuation characters to get
this right, as in the following sentence.
       Gnats, gnus, etc.\ all  % `\ ' makes an inter-word space.
       begin with G\@.         % \@ marks end-of-sentence punctuation.
You should check the spaces after periods when
reading your output to make sure you haven't
forgotten any special cases.  Generating an
ellipsis
       \ldots\               % `\ ' is needed after `\ldots' because TeX
                             % ignores spaces after command names like \ldots
                             % made from \ + letters.
                             %
                             % Note how a `%' character causes TeX to ignore
                             % the end of the input line, so these blank lines
                             % do not start a new paragraph.
                             %
with the right spacing around the periods requires
a special command.

\LaTeX\ interprets some common characters as
commands, so you must type special commands to
generate them.  These characters include the
following:
       \$ \& \% \# \{ and \}.

In printing, text is usually emphasized with an
       \emph{italic}
type style.

\begin{em}
   A long segment of text can also be emphasized
   in this way.  Text within such a segment can be
   given \emph{additional} emphasis.
\end{em}

It is sometimes necessary to prevent \LaTeX\ from
breaking a line where it might otherwise do so.
This may be at a space, as between the ``Mr.''\ and
``Jones'' in
       ``Mr.~Jones'',        % ~ produces an unbreakable interword space.
or within a word---especially when the word is a
symbol like
       \mbox{\emph{itemnum}}
that makes little sense when hyphenated across
lines.

Footnotes\footnote{This is an example of a footnote.}
pose no problem.

\LaTeX\ is good at typesetting mathematical formulas
like
       \( x-3y + z = 7 \)
or
       \( a_{1} > x^{2n} + y^{2n} > x' \)
or
       \( \ip{A}{B} = \sum_{i} a_{i} b_{i} \).
The spaces you type in a formula are
ignored.  Remember that a letter like
       $x$                   % $ ... $  and  \( ... \)  are equivalent
is a formula when it denotes a mathematical
symbol, and it should be typed as one.

\section{Displayed Text}

Text is displayed by indenting it from the left
margin.  Quotations are commonly displayed.  There
are short quotations
\begin{quote}
   This is a short quotation.  It consists of a
   single paragraph of text.  See how it is formatted.
\end{quote}
and longer ones.
\begin{quotation}
   This is a longer quotation.  It consists of two
   paragraphs of text, neither of which are
   particularly interesting.

   This is the second paragraph of the quotation.  It
   is just as dull as the first paragraph.
\end{quotation}
Another frequently-displayed structure is a list.
The following is an example of an \emph{itemized}
list.
\begin{itemize}
   \item This is the first item of an itemized list.
         Each item in the list is marked with a ``tick''.
         You don't have to worry about what kind of tick
         mark is used.

   \item This is the second item of the list.  It
         contains another list nested inside it.  The inner
         list is an \emph{enumerated} list.
         \begin{enumerate}
            \item This is the first item of an enumerated
                  list that is nested within the itemized list.

            \item This is the second item of the inner list.
                  \LaTeX\ allows you to nest lists deeper than
                  you really should.
         \end{enumerate}
         This is the rest of the second item of the outer
         list.  It is no more interesting than any other
         part of the item.
   \item This is the third item of the list.
\end{itemize}
You can even display poetry.
\begin{verse}
   There is an environment
    for verse \\             % The \\ command separates lines
   Whose features some poets % within a stanza.
   will curse.

                             % One or more blank lines separate stanzas.

   For instead of making\\
   Them do \emph{all} line breaking, \\
   It allows them to put too many words on a line when they'd rather be
   forced to be terse.
\end{verse}

Mathematical formulas may also be displayed.  A
displayed formula
is
one-line long; multiline
formulas require special formatting instructions.
   \[  \ip{\Gamma}{\psi'} = x'' + y^{2} + z_{i}^{n}\]
Don't start a paragraph with a displayed equation,
nor make one a paragraph by itself.

\end{document}               % End of document.

\section{Bench Huge Section 316}
% synthetic bench section 316
% This is a sample LaTeX input file.  (Version of 12 August 2004.)
%
% A '%' character causes TeX to ignore all remaining text on the line,
% and is used for comments like this one.

\documentclass{article}      % Specifies the document class

                             % The preamble begins here.
\title{An Example Document}  % Declares the document's title.
\author{Leslie Lamport}      % Declares the author's name.
\date{January 21, 1994}      % Deleting this command produces today's date.

\newcommand{\ip}[2]{(#1, #2)}
                             % Defines \ip{arg1}{arg2} to mean
                             % (arg1, arg2).

%\newcommand{\ip}[2]{\langle #1 | #2\rangle}
                             % This is an alternative definition of
                             % \ip that is commented out.

\begin{document}             % End of preamble and beginning of text.

\maketitle                   % Produces the title.

This is an example input file.  Comparing it with
the output it generates can show you how to
produce a simple document of your own.

\section{Ordinary Text}      % Produces section heading.  Lower-level
                             % sections are begun with similar
                             % \subsection and \subsubsection commands.

The ends  of words and sentences are marked
  by   spaces. It  doesn't matter how many
spaces    you type; one is as good as 100.  The
end of   a line counts as a space.

One   or more   blank lines denote the  end
of  a paragraph.

Since any number of consecutive spaces are treated
like a single one, the formatting of the input
file makes no difference to
      \LaTeX,                % The \LaTeX command generates the LaTeX logo.
but it makes a difference to you.  When you use
\LaTeX, making your input file as easy to read
as possible will be a great help as you write
your document and when you change it.  This sample
file shows how you can add comments to your own input
file.

Because printing is different from typewriting,
there are a number of things that you have to do
differently when preparing an input file than if
you were just typing the document directly.
Quotation marks like
       ``this''
have to be handled specially, as do quotes within
quotes:
       ``\,`this'            % \, separates the double and single quote.
        is what I just
        wrote, not  `that'\,''.

Dashes come in three sizes: an
       intra-word
dash, a medium dash for number ranges like
       1--2,
and a punctuation
       dash---like
this.

A sentence-ending space should be larger than the
space between words within a sentence.  You
sometimes have to type special commands in
conjunction with punctuation characters to get
this right, as in the following sentence.
       Gnats, gnus, etc.\ all  % `\ ' makes an inter-word space.
       begin with G\@.         % \@ marks end-of-sentence punctuation.
You should check the spaces after periods when
reading your output to make sure you haven't
forgotten any special cases.  Generating an
ellipsis
       \ldots\               % `\ ' is needed after `\ldots' because TeX
                             % ignores spaces after command names like \ldots
                             % made from \ + letters.
                             %
                             % Note how a `%' character causes TeX to ignore
                             % the end of the input line, so these blank lines
                             % do not start a new paragraph.
                             %
with the right spacing around the periods requires
a special command.

\LaTeX\ interprets some common characters as
commands, so you must type special commands to
generate them.  These characters include the
following:
       \$ \& \% \# \{ and \}.

In printing, text is usually emphasized with an
       \emph{italic}
type style.

\begin{em}
   A long segment of text can also be emphasized
   in this way.  Text within such a segment can be
   given \emph{additional} emphasis.
\end{em}

It is sometimes necessary to prevent \LaTeX\ from
breaking a line where it might otherwise do so.
This may be at a space, as between the ``Mr.''\ and
``Jones'' in
       ``Mr.~Jones'',        % ~ produces an unbreakable interword space.
or within a word---especially when the word is a
symbol like
       \mbox{\emph{itemnum}}
that makes little sense when hyphenated across
lines.

Footnotes\footnote{This is an example of a footnote.}
pose no problem.

\LaTeX\ is good at typesetting mathematical formulas
like
       \( x-3y + z = 7 \)
or
       \( a_{1} > x^{2n} + y^{2n} > x' \)
or
       \( \ip{A}{B} = \sum_{i} a_{i} b_{i} \).
The spaces you type in a formula are
ignored.  Remember that a letter like
       $x$                   % $ ... $  and  \( ... \)  are equivalent
is a formula when it denotes a mathematical
symbol, and it should be typed as one.

\section{Displayed Text}

Text is displayed by indenting it from the left
margin.  Quotations are commonly displayed.  There
are short quotations
\begin{quote}
   This is a short quotation.  It consists of a
   single paragraph of text.  See how it is formatted.
\end{quote}
and longer ones.
\begin{quotation}
   This is a longer quotation.  It consists of two
   paragraphs of text, neither of which are
   particularly interesting.

   This is the second paragraph of the quotation.  It
   is just as dull as the first paragraph.
\end{quotation}
Another frequently-displayed structure is a list.
The following is an example of an \emph{itemized}
list.
\begin{itemize}
   \item This is the first item of an itemized list.
         Each item in the list is marked with a ``tick''.
         You don't have to worry about what kind of tick
         mark is used.

   \item This is the second item of the list.  It
         contains another list nested inside it.  The inner
         list is an \emph{enumerated} list.
         \begin{enumerate}
            \item This is the first item of an enumerated
                  list that is nested within the itemized list.

            \item This is the second item of the inner list.
                  \LaTeX\ allows you to nest lists deeper than
                  you really should.
         \end{enumerate}
         This is the rest of the second item of the outer
         list.  It is no more interesting than any other
         part of the item.
   \item This is the third item of the list.
\end{itemize}
You can even display poetry.
\begin{verse}
   There is an environment
    for verse \\             % The \\ command separates lines
   Whose features some poets % within a stanza.
   will curse.

                             % One or more blank lines separate stanzas.

   For instead of making\\
   Them do \emph{all} line breaking, \\
   It allows them to put too many words on a line when they'd rather be
   forced to be terse.
\end{verse}

Mathematical formulas may also be displayed.  A
displayed formula
is
one-line long; multiline
formulas require special formatting instructions.
   \[  \ip{\Gamma}{\psi'} = x'' + y^{2} + z_{i}^{n}\]
Don't start a paragraph with a displayed equation,
nor make one a paragraph by itself.

\end{document}               % End of document.

\section{Bench Huge Section 317}
% synthetic bench section 317
% This is a sample LaTeX input file.  (Version of 12 August 2004.)
%
% A '%' character causes TeX to ignore all remaining text on the line,
% and is used for comments like this one.

\documentclass{article}      % Specifies the document class

                             % The preamble begins here.
\title{An Example Document}  % Declares the document's title.
\author{Leslie Lamport}      % Declares the author's name.
\date{January 21, 1994}      % Deleting this command produces today's date.

\newcommand{\ip}[2]{(#1, #2)}
                             % Defines \ip{arg1}{arg2} to mean
                             % (arg1, arg2).

%\newcommand{\ip}[2]{\langle #1 | #2\rangle}
                             % This is an alternative definition of
                             % \ip that is commented out.

\begin{document}             % End of preamble and beginning of text.

\maketitle                   % Produces the title.

This is an example input file.  Comparing it with
the output it generates can show you how to
produce a simple document of your own.

\section{Ordinary Text}      % Produces section heading.  Lower-level
                             % sections are begun with similar
                             % \subsection and \subsubsection commands.

The ends  of words and sentences are marked
  by   spaces. It  doesn't matter how many
spaces    you type; one is as good as 100.  The
end of   a line counts as a space.

One   or more   blank lines denote the  end
of  a paragraph.

Since any number of consecutive spaces are treated
like a single one, the formatting of the input
file makes no difference to
      \LaTeX,                % The \LaTeX command generates the LaTeX logo.
but it makes a difference to you.  When you use
\LaTeX, making your input file as easy to read
as possible will be a great help as you write
your document and when you change it.  This sample
file shows how you can add comments to your own input
file.

Because printing is different from typewriting,
there are a number of things that you have to do
differently when preparing an input file than if
you were just typing the document directly.
Quotation marks like
       ``this''
have to be handled specially, as do quotes within
quotes:
       ``\,`this'            % \, separates the double and single quote.
        is what I just
        wrote, not  `that'\,''.

Dashes come in three sizes: an
       intra-word
dash, a medium dash for number ranges like
       1--2,
and a punctuation
       dash---like
this.

A sentence-ending space should be larger than the
space between words within a sentence.  You
sometimes have to type special commands in
conjunction with punctuation characters to get
this right, as in the following sentence.
       Gnats, gnus, etc.\ all  % `\ ' makes an inter-word space.
       begin with G\@.         % \@ marks end-of-sentence punctuation.
You should check the spaces after periods when
reading your output to make sure you haven't
forgotten any special cases.  Generating an
ellipsis
       \ldots\               % `\ ' is needed after `\ldots' because TeX
                             % ignores spaces after command names like \ldots
                             % made from \ + letters.
                             %
                             % Note how a `%' character causes TeX to ignore
                             % the end of the input line, so these blank lines
                             % do not start a new paragraph.
                             %
with the right spacing around the periods requires
a special command.

\LaTeX\ interprets some common characters as
commands, so you must type special commands to
generate them.  These characters include the
following:
       \$ \& \% \# \{ and \}.

In printing, text is usually emphasized with an
       \emph{italic}
type style.

\begin{em}
   A long segment of text can also be emphasized
   in this way.  Text within such a segment can be
   given \emph{additional} emphasis.
\end{em}

It is sometimes necessary to prevent \LaTeX\ from
breaking a line where it might otherwise do so.
This may be at a space, as between the ``Mr.''\ and
``Jones'' in
       ``Mr.~Jones'',        % ~ produces an unbreakable interword space.
or within a word---especially when the word is a
symbol like
       \mbox{\emph{itemnum}}
that makes little sense when hyphenated across
lines.

Footnotes\footnote{This is an example of a footnote.}
pose no problem.

\LaTeX\ is good at typesetting mathematical formulas
like
       \( x-3y + z = 7 \)
or
       \( a_{1} > x^{2n} + y^{2n} > x' \)
or
       \( \ip{A}{B} = \sum_{i} a_{i} b_{i} \).
The spaces you type in a formula are
ignored.  Remember that a letter like
       $x$                   % $ ... $  and  \( ... \)  are equivalent
is a formula when it denotes a mathematical
symbol, and it should be typed as one.

\section{Displayed Text}

Text is displayed by indenting it from the left
margin.  Quotations are commonly displayed.  There
are short quotations
\begin{quote}
   This is a short quotation.  It consists of a
   single paragraph of text.  See how it is formatted.
\end{quote}
and longer ones.
\begin{quotation}
   This is a longer quotation.  It consists of two
   paragraphs of text, neither of which are
   particularly interesting.

   This is the second paragraph of the quotation.  It
   is just as dull as the first paragraph.
\end{quotation}
Another frequently-displayed structure is a list.
The following is an example of an \emph{itemized}
list.
\begin{itemize}
   \item This is the first item of an itemized list.
         Each item in the list is marked with a ``tick''.
         You don't have to worry about what kind of tick
         mark is used.

   \item This is the second item of the list.  It
         contains another list nested inside it.  The inner
         list is an \emph{enumerated} list.
         \begin{enumerate}
            \item This is the first item of an enumerated
                  list that is nested within the itemized list.

            \item This is the second item of the inner list.
                  \LaTeX\ allows you to nest lists deeper than
                  you really should.
         \end{enumerate}
         This is the rest of the second item of the outer
         list.  It is no more interesting than any other
         part of the item.
   \item This is the third item of the list.
\end{itemize}
You can even display poetry.
\begin{verse}
   There is an environment
    for verse \\             % The \\ command separates lines
   Whose features some poets % within a stanza.
   will curse.

                             % One or more blank lines separate stanzas.

   For instead of making\\
   Them do \emph{all} line breaking, \\
   It allows them to put too many words on a line when they'd rather be
   forced to be terse.
\end{verse}

Mathematical formulas may also be displayed.  A
displayed formula
is
one-line long; multiline
formulas require special formatting instructions.
   \[  \ip{\Gamma}{\psi'} = x'' + y^{2} + z_{i}^{n}\]
Don't start a paragraph with a displayed equation,
nor make one a paragraph by itself.

\end{document}               % End of document.

\section{Bench Huge Section 318}
% synthetic bench section 318
% This is a sample LaTeX input file.  (Version of 12 August 2004.)
%
% A '%' character causes TeX to ignore all remaining text on the line,
% and is used for comments like this one.

\documentclass{article}      % Specifies the document class

                             % The preamble begins here.
\title{An Example Document}  % Declares the document's title.
\author{Leslie Lamport}      % Declares the author's name.
\date{January 21, 1994}      % Deleting this command produces today's date.

\newcommand{\ip}[2]{(#1, #2)}
                             % Defines \ip{arg1}{arg2} to mean
                             % (arg1, arg2).

%\newcommand{\ip}[2]{\langle #1 | #2\rangle}
                             % This is an alternative definition of
                             % \ip that is commented out.

\begin{document}             % End of preamble and beginning of text.

\maketitle                   % Produces the title.

This is an example input file.  Comparing it with
the output it generates can show you how to
produce a simple document of your own.

\section{Ordinary Text}      % Produces section heading.  Lower-level
                             % sections are begun with similar
                             % \subsection and \subsubsection commands.

The ends  of words and sentences are marked
  by   spaces. It  doesn't matter how many
spaces    you type; one is as good as 100.  The
end of   a line counts as a space.

One   or more   blank lines denote the  end
of  a paragraph.

Since any number of consecutive spaces are treated
like a single one, the formatting of the input
file makes no difference to
      \LaTeX,                % The \LaTeX command generates the LaTeX logo.
but it makes a difference to you.  When you use
\LaTeX, making your input file as easy to read
as possible will be a great help as you write
your document and when you change it.  This sample
file shows how you can add comments to your own input
file.

Because printing is different from typewriting,
there are a number of things that you have to do
differently when preparing an input file than if
you were just typing the document directly.
Quotation marks like
       ``this''
have to be handled specially, as do quotes within
quotes:
       ``\,`this'            % \, separates the double and single quote.
        is what I just
        wrote, not  `that'\,''.

Dashes come in three sizes: an
       intra-word
dash, a medium dash for number ranges like
       1--2,
and a punctuation
       dash---like
this.

A sentence-ending space should be larger than the
space between words within a sentence.  You
sometimes have to type special commands in
conjunction with punctuation characters to get
this right, as in the following sentence.
       Gnats, gnus, etc.\ all  % `\ ' makes an inter-word space.
       begin with G\@.         % \@ marks end-of-sentence punctuation.
You should check the spaces after periods when
reading your output to make sure you haven't
forgotten any special cases.  Generating an
ellipsis
       \ldots\               % `\ ' is needed after `\ldots' because TeX
                             % ignores spaces after command names like \ldots
                             % made from \ + letters.
                             %
                             % Note how a `%' character causes TeX to ignore
                             % the end of the input line, so these blank lines
                             % do not start a new paragraph.
                             %
with the right spacing around the periods requires
a special command.

\LaTeX\ interprets some common characters as
commands, so you must type special commands to
generate them.  These characters include the
following:
       \$ \& \% \# \{ and \}.

In printing, text is usually emphasized with an
       \emph{italic}
type style.

\begin{em}
   A long segment of text can also be emphasized
   in this way.  Text within such a segment can be
   given \emph{additional} emphasis.
\end{em}

It is sometimes necessary to prevent \LaTeX\ from
breaking a line where it might otherwise do so.
This may be at a space, as between the ``Mr.''\ and
``Jones'' in
       ``Mr.~Jones'',        % ~ produces an unbreakable interword space.
or within a word---especially when the word is a
symbol like
       \mbox{\emph{itemnum}}
that makes little sense when hyphenated across
lines.

Footnotes\footnote{This is an example of a footnote.}
pose no problem.

\LaTeX\ is good at typesetting mathematical formulas
like
       \( x-3y + z = 7 \)
or
       \( a_{1} > x^{2n} + y^{2n} > x' \)
or
       \( \ip{A}{B} = \sum_{i} a_{i} b_{i} \).
The spaces you type in a formula are
ignored.  Remember that a letter like
       $x$                   % $ ... $  and  \( ... \)  are equivalent
is a formula when it denotes a mathematical
symbol, and it should be typed as one.

\section{Displayed Text}

Text is displayed by indenting it from the left
margin.  Quotations are commonly displayed.  There
are short quotations
\begin{quote}
   This is a short quotation.  It consists of a
   single paragraph of text.  See how it is formatted.
\end{quote}
and longer ones.
\begin{quotation}
   This is a longer quotation.  It consists of two
   paragraphs of text, neither of which are
   particularly interesting.

   This is the second paragraph of the quotation.  It
   is just as dull as the first paragraph.
\end{quotation}
Another frequently-displayed structure is a list.
The following is an example of an \emph{itemized}
list.
\begin{itemize}
   \item This is the first item of an itemized list.
         Each item in the list is marked with a ``tick''.
         You don't have to worry about what kind of tick
         mark is used.

   \item This is the second item of the list.  It
         contains another list nested inside it.  The inner
         list is an \emph{enumerated} list.
         \begin{enumerate}
            \item This is the first item of an enumerated
                  list that is nested within the itemized list.

            \item This is the second item of the inner list.
                  \LaTeX\ allows you to nest lists deeper than
                  you really should.
         \end{enumerate}
         This is the rest of the second item of the outer
         list.  It is no more interesting than any other
         part of the item.
   \item This is the third item of the list.
\end{itemize}
You can even display poetry.
\begin{verse}
   There is an environment
    for verse \\             % The \\ command separates lines
   Whose features some poets % within a stanza.
   will curse.

                             % One or more blank lines separate stanzas.

   For instead of making\\
   Them do \emph{all} line breaking, \\
   It allows them to put too many words on a line when they'd rather be
   forced to be terse.
\end{verse}

Mathematical formulas may also be displayed.  A
displayed formula
is
one-line long; multiline
formulas require special formatting instructions.
   \[  \ip{\Gamma}{\psi'} = x'' + y^{2} + z_{i}^{n}\]
Don't start a paragraph with a displayed equation,
nor make one a paragraph by itself.

\end{document}               % End of document.

\section{Bench Huge Section 319}
% synthetic bench section 319
% This is a sample LaTeX input file.  (Version of 12 August 2004.)
%
% A '%' character causes TeX to ignore all remaining text on the line,
% and is used for comments like this one.

\documentclass{article}      % Specifies the document class

                             % The preamble begins here.
\title{An Example Document}  % Declares the document's title.
\author{Leslie Lamport}      % Declares the author's name.
\date{January 21, 1994}      % Deleting this command produces today's date.

\newcommand{\ip}[2]{(#1, #2)}
                             % Defines \ip{arg1}{arg2} to mean
                             % (arg1, arg2).

%\newcommand{\ip}[2]{\langle #1 | #2\rangle}
                             % This is an alternative definition of
                             % \ip that is commented out.

\begin{document}             % End of preamble and beginning of text.

\maketitle                   % Produces the title.

This is an example input file.  Comparing it with
the output it generates can show you how to
produce a simple document of your own.

\section{Ordinary Text}      % Produces section heading.  Lower-level
                             % sections are begun with similar
                             % \subsection and \subsubsection commands.

The ends  of words and sentences are marked
  by   spaces. It  doesn't matter how many
spaces    you type; one is as good as 100.  The
end of   a line counts as a space.

One   or more   blank lines denote the  end
of  a paragraph.

Since any number of consecutive spaces are treated
like a single one, the formatting of the input
file makes no difference to
      \LaTeX,                % The \LaTeX command generates the LaTeX logo.
but it makes a difference to you.  When you use
\LaTeX, making your input file as easy to read
as possible will be a great help as you write
your document and when you change it.  This sample
file shows how you can add comments to your own input
file.

Because printing is different from typewriting,
there are a number of things that you have to do
differently when preparing an input file than if
you were just typing the document directly.
Quotation marks like
       ``this''
have to be handled specially, as do quotes within
quotes:
       ``\,`this'            % \, separates the double and single quote.
        is what I just
        wrote, not  `that'\,''.

Dashes come in three sizes: an
       intra-word
dash, a medium dash for number ranges like
       1--2,
and a punctuation
       dash---like
this.

A sentence-ending space should be larger than the
space between words within a sentence.  You
sometimes have to type special commands in
conjunction with punctuation characters to get
this right, as in the following sentence.
       Gnats, gnus, etc.\ all  % `\ ' makes an inter-word space.
       begin with G\@.         % \@ marks end-of-sentence punctuation.
You should check the spaces after periods when
reading your output to make sure you haven't
forgotten any special cases.  Generating an
ellipsis
       \ldots\               % `\ ' is needed after `\ldots' because TeX
                             % ignores spaces after command names like \ldots
                             % made from \ + letters.
                             %
                             % Note how a `%' character causes TeX to ignore
                             % the end of the input line, so these blank lines
                             % do not start a new paragraph.
                             %
with the right spacing around the periods requires
a special command.

\LaTeX\ interprets some common characters as
commands, so you must type special commands to
generate them.  These characters include the
following:
       \$ \& \% \# \{ and \}.

In printing, text is usually emphasized with an
       \emph{italic}
type style.

\begin{em}
   A long segment of text can also be emphasized
   in this way.  Text within such a segment can be
   given \emph{additional} emphasis.
\end{em}

It is sometimes necessary to prevent \LaTeX\ from
breaking a line where it might otherwise do so.
This may be at a space, as between the ``Mr.''\ and
``Jones'' in
       ``Mr.~Jones'',        % ~ produces an unbreakable interword space.
or within a word---especially when the word is a
symbol like
       \mbox{\emph{itemnum}}
that makes little sense when hyphenated across
lines.

Footnotes\footnote{This is an example of a footnote.}
pose no problem.

\LaTeX\ is good at typesetting mathematical formulas
like
       \( x-3y + z = 7 \)
or
       \( a_{1} > x^{2n} + y^{2n} > x' \)
or
       \( \ip{A}{B} = \sum_{i} a_{i} b_{i} \).
The spaces you type in a formula are
ignored.  Remember that a letter like
       $x$                   % $ ... $  and  \( ... \)  are equivalent
is a formula when it denotes a mathematical
symbol, and it should be typed as one.

\section{Displayed Text}

Text is displayed by indenting it from the left
margin.  Quotations are commonly displayed.  There
are short quotations
\begin{quote}
   This is a short quotation.  It consists of a
   single paragraph of text.  See how it is formatted.
\end{quote}
and longer ones.
\begin{quotation}
   This is a longer quotation.  It consists of two
   paragraphs of text, neither of which are
   particularly interesting.

   This is the second paragraph of the quotation.  It
   is just as dull as the first paragraph.
\end{quotation}
Another frequently-displayed structure is a list.
The following is an example of an \emph{itemized}
list.
\begin{itemize}
   \item This is the first item of an itemized list.
         Each item in the list is marked with a ``tick''.
         You don't have to worry about what kind of tick
         mark is used.

   \item This is the second item of the list.  It
         contains another list nested inside it.  The inner
         list is an \emph{enumerated} list.
         \begin{enumerate}
            \item This is the first item of an enumerated
                  list that is nested within the itemized list.

            \item This is the second item of the inner list.
                  \LaTeX\ allows you to nest lists deeper than
                  you really should.
         \end{enumerate}
         This is the rest of the second item of the outer
         list.  It is no more interesting than any other
         part of the item.
   \item This is the third item of the list.
\end{itemize}
You can even display poetry.
\begin{verse}
   There is an environment
    for verse \\             % The \\ command separates lines
   Whose features some poets % within a stanza.
   will curse.

                             % One or more blank lines separate stanzas.

   For instead of making\\
   Them do \emph{all} line breaking, \\
   It allows them to put too many words on a line when they'd rather be
   forced to be terse.
\end{verse}

Mathematical formulas may also be displayed.  A
displayed formula
is
one-line long; multiline
formulas require special formatting instructions.
   \[  \ip{\Gamma}{\psi'} = x'' + y^{2} + z_{i}^{n}\]
Don't start a paragraph with a displayed equation,
nor make one a paragraph by itself.

\end{document}               % End of document.

\section{Bench Huge Section 320}
% synthetic bench section 320
% This is a sample LaTeX input file.  (Version of 12 August 2004.)
%
% A '%' character causes TeX to ignore all remaining text on the line,
% and is used for comments like this one.

\documentclass{article}      % Specifies the document class

                             % The preamble begins here.
\title{An Example Document}  % Declares the document's title.
\author{Leslie Lamport}      % Declares the author's name.
\date{January 21, 1994}      % Deleting this command produces today's date.

\newcommand{\ip}[2]{(#1, #2)}
                             % Defines \ip{arg1}{arg2} to mean
                             % (arg1, arg2).

%\newcommand{\ip}[2]{\langle #1 | #2\rangle}
                             % This is an alternative definition of
                             % \ip that is commented out.

\begin{document}             % End of preamble and beginning of text.

\maketitle                   % Produces the title.

This is an example input file.  Comparing it with
the output it generates can show you how to
produce a simple document of your own.

\section{Ordinary Text}      % Produces section heading.  Lower-level
                             % sections are begun with similar
                             % \subsection and \subsubsection commands.

The ends  of words and sentences are marked
  by   spaces. It  doesn't matter how many
spaces    you type; one is as good as 100.  The
end of   a line counts as a space.

One   or more   blank lines denote the  end
of  a paragraph.

Since any number of consecutive spaces are treated
like a single one, the formatting of the input
file makes no difference to
      \LaTeX,                % The \LaTeX command generates the LaTeX logo.
but it makes a difference to you.  When you use
\LaTeX, making your input file as easy to read
as possible will be a great help as you write
your document and when you change it.  This sample
file shows how you can add comments to your own input
file.

Because printing is different from typewriting,
there are a number of things that you have to do
differently when preparing an input file than if
you were just typing the document directly.
Quotation marks like
       ``this''
have to be handled specially, as do quotes within
quotes:
       ``\,`this'            % \, separates the double and single quote.
        is what I just
        wrote, not  `that'\,''.

Dashes come in three sizes: an
       intra-word
dash, a medium dash for number ranges like
       1--2,
and a punctuation
       dash---like
this.

A sentence-ending space should be larger than the
space between words within a sentence.  You
sometimes have to type special commands in
conjunction with punctuation characters to get
this right, as in the following sentence.
       Gnats, gnus, etc.\ all  % `\ ' makes an inter-word space.
       begin with G\@.         % \@ marks end-of-sentence punctuation.
You should check the spaces after periods when
reading your output to make sure you haven't
forgotten any special cases.  Generating an
ellipsis
       \ldots\               % `\ ' is needed after `\ldots' because TeX
                             % ignores spaces after command names like \ldots
                             % made from \ + letters.
                             %
                             % Note how a `%' character causes TeX to ignore
                             % the end of the input line, so these blank lines
                             % do not start a new paragraph.
                             %
with the right spacing around the periods requires
a special command.

\LaTeX\ interprets some common characters as
commands, so you must type special commands to
generate them.  These characters include the
following:
       \$ \& \% \# \{ and \}.

In printing, text is usually emphasized with an
       \emph{italic}
type style.

\begin{em}
   A long segment of text can also be emphasized
   in this way.  Text within such a segment can be
   given \emph{additional} emphasis.
\end{em}

It is sometimes necessary to prevent \LaTeX\ from
breaking a line where it might otherwise do so.
This may be at a space, as between the ``Mr.''\ and
``Jones'' in
       ``Mr.~Jones'',        % ~ produces an unbreakable interword space.
or within a word---especially when the word is a
symbol like
       \mbox{\emph{itemnum}}
that makes little sense when hyphenated across
lines.

Footnotes\footnote{This is an example of a footnote.}
pose no problem.

\LaTeX\ is good at typesetting mathematical formulas
like
       \( x-3y + z = 7 \)
or
       \( a_{1} > x^{2n} + y^{2n} > x' \)
or
       \( \ip{A}{B} = \sum_{i} a_{i} b_{i} \).
The spaces you type in a formula are
ignored.  Remember that a letter like
       $x$                   % $ ... $  and  \( ... \)  are equivalent
is a formula when it denotes a mathematical
symbol, and it should be typed as one.

\section{Displayed Text}

Text is displayed by indenting it from the left
margin.  Quotations are commonly displayed.  There
are short quotations
\begin{quote}
   This is a short quotation.  It consists of a
   single paragraph of text.  See how it is formatted.
\end{quote}
and longer ones.
\begin{quotation}
   This is a longer quotation.  It consists of two
   paragraphs of text, neither of which are
   particularly interesting.

   This is the second paragraph of the quotation.  It
   is just as dull as the first paragraph.
\end{quotation}
Another frequently-displayed structure is a list.
The following is an example of an \emph{itemized}
list.
\begin{itemize}
   \item This is the first item of an itemized list.
         Each item in the list is marked with a ``tick''.
         You don't have to worry about what kind of tick
         mark is used.

   \item This is the second item of the list.  It
         contains another list nested inside it.  The inner
         list is an \emph{enumerated} list.
         \begin{enumerate}
            \item This is the first item of an enumerated
                  list that is nested within the itemized list.

            \item This is the second item of the inner list.
                  \LaTeX\ allows you to nest lists deeper than
                  you really should.
         \end{enumerate}
         This is the rest of the second item of the outer
         list.  It is no more interesting than any other
         part of the item.
   \item This is the third item of the list.
\end{itemize}
You can even display poetry.
\begin{verse}
   There is an environment
    for verse \\             % The \\ command separates lines
   Whose features some poets % within a stanza.
   will curse.

                             % One or more blank lines separate stanzas.

   For instead of making\\
   Them do \emph{all} line breaking, \\
   It allows them to put too many words on a line when they'd rather be
   forced to be terse.
\end{verse}

Mathematical formulas may also be displayed.  A
displayed formula
is
one-line long; multiline
formulas require special formatting instructions.
   \[  \ip{\Gamma}{\psi'} = x'' + y^{2} + z_{i}^{n}\]
Don't start a paragraph with a displayed equation,
nor make one a paragraph by itself.

\end{document}               % End of document.

\section{Bench Huge Section 321}
% synthetic bench section 321
% This is a sample LaTeX input file.  (Version of 12 August 2004.)
%
% A '%' character causes TeX to ignore all remaining text on the line,
% and is used for comments like this one.

\documentclass{article}      % Specifies the document class

                             % The preamble begins here.
\title{An Example Document}  % Declares the document's title.
\author{Leslie Lamport}      % Declares the author's name.
\date{January 21, 1994}      % Deleting this command produces today's date.

\newcommand{\ip}[2]{(#1, #2)}
                             % Defines \ip{arg1}{arg2} to mean
                             % (arg1, arg2).

%\newcommand{\ip}[2]{\langle #1 | #2\rangle}
                             % This is an alternative definition of
                             % \ip that is commented out.

\begin{document}             % End of preamble and beginning of text.

\maketitle                   % Produces the title.

This is an example input file.  Comparing it with
the output it generates can show you how to
produce a simple document of your own.

\section{Ordinary Text}      % Produces section heading.  Lower-level
                             % sections are begun with similar
                             % \subsection and \subsubsection commands.

The ends  of words and sentences are marked
  by   spaces. It  doesn't matter how many
spaces    you type; one is as good as 100.  The
end of   a line counts as a space.

One   or more   blank lines denote the  end
of  a paragraph.

Since any number of consecutive spaces are treated
like a single one, the formatting of the input
file makes no difference to
      \LaTeX,                % The \LaTeX command generates the LaTeX logo.
but it makes a difference to you.  When you use
\LaTeX, making your input file as easy to read
as possible will be a great help as you write
your document and when you change it.  This sample
file shows how you can add comments to your own input
file.

Because printing is different from typewriting,
there are a number of things that you have to do
differently when preparing an input file than if
you were just typing the document directly.
Quotation marks like
       ``this''
have to be handled specially, as do quotes within
quotes:
       ``\,`this'            % \, separates the double and single quote.
        is what I just
        wrote, not  `that'\,''.

Dashes come in three sizes: an
       intra-word
dash, a medium dash for number ranges like
       1--2,
and a punctuation
       dash---like
this.

A sentence-ending space should be larger than the
space between words within a sentence.  You
sometimes have to type special commands in
conjunction with punctuation characters to get
this right, as in the following sentence.
       Gnats, gnus, etc.\ all  % `\ ' makes an inter-word space.
       begin with G\@.         % \@ marks end-of-sentence punctuation.
You should check the spaces after periods when
reading your output to make sure you haven't
forgotten any special cases.  Generating an
ellipsis
       \ldots\               % `\ ' is needed after `\ldots' because TeX
                             % ignores spaces after command names like \ldots
                             % made from \ + letters.
                             %
                             % Note how a `%' character causes TeX to ignore
                             % the end of the input line, so these blank lines
                             % do not start a new paragraph.
                             %
with the right spacing around the periods requires
a special command.

\LaTeX\ interprets some common characters as
commands, so you must type special commands to
generate them.  These characters include the
following:
       \$ \& \% \# \{ and \}.

In printing, text is usually emphasized with an
       \emph{italic}
type style.

\begin{em}
   A long segment of text can also be emphasized
   in this way.  Text within such a segment can be
   given \emph{additional} emphasis.
\end{em}

It is sometimes necessary to prevent \LaTeX\ from
breaking a line where it might otherwise do so.
This may be at a space, as between the ``Mr.''\ and
``Jones'' in
       ``Mr.~Jones'',        % ~ produces an unbreakable interword space.
or within a word---especially when the word is a
symbol like
       \mbox{\emph{itemnum}}
that makes little sense when hyphenated across
lines.

Footnotes\footnote{This is an example of a footnote.}
pose no problem.

\LaTeX\ is good at typesetting mathematical formulas
like
       \( x-3y + z = 7 \)
or
       \( a_{1} > x^{2n} + y^{2n} > x' \)
or
       \( \ip{A}{B} = \sum_{i} a_{i} b_{i} \).
The spaces you type in a formula are
ignored.  Remember that a letter like
       $x$                   % $ ... $  and  \( ... \)  are equivalent
is a formula when it denotes a mathematical
symbol, and it should be typed as one.

\section{Displayed Text}

Text is displayed by indenting it from the left
margin.  Quotations are commonly displayed.  There
are short quotations
\begin{quote}
   This is a short quotation.  It consists of a
   single paragraph of text.  See how it is formatted.
\end{quote}
and longer ones.
\begin{quotation}
   This is a longer quotation.  It consists of two
   paragraphs of text, neither of which are
   particularly interesting.

   This is the second paragraph of the quotation.  It
   is just as dull as the first paragraph.
\end{quotation}
Another frequently-displayed structure is a list.
The following is an example of an \emph{itemized}
list.
\begin{itemize}
   \item This is the first item of an itemized list.
         Each item in the list is marked with a ``tick''.
         You don't have to worry about what kind of tick
         mark is used.

   \item This is the second item of the list.  It
         contains another list nested inside it.  The inner
         list is an \emph{enumerated} list.
         \begin{enumerate}
            \item This is the first item of an enumerated
                  list that is nested within the itemized list.

            \item This is the second item of the inner list.
                  \LaTeX\ allows you to nest lists deeper than
                  you really should.
         \end{enumerate}
         This is the rest of the second item of the outer
         list.  It is no more interesting than any other
         part of the item.
   \item This is the third item of the list.
\end{itemize}
You can even display poetry.
\begin{verse}
   There is an environment
    for verse \\             % The \\ command separates lines
   Whose features some poets % within a stanza.
   will curse.

                             % One or more blank lines separate stanzas.

   For instead of making\\
   Them do \emph{all} line breaking, \\
   It allows them to put too many words on a line when they'd rather be
   forced to be terse.
\end{verse}

Mathematical formulas may also be displayed.  A
displayed formula
is
one-line long; multiline
formulas require special formatting instructions.
   \[  \ip{\Gamma}{\psi'} = x'' + y^{2} + z_{i}^{n}\]
Don't start a paragraph with a displayed equation,
nor make one a paragraph by itself.

\end{document}               % End of document.

\section{Bench Huge Section 322}
% synthetic bench section 322
% This is a sample LaTeX input file.  (Version of 12 August 2004.)
%
% A '%' character causes TeX to ignore all remaining text on the line,
% and is used for comments like this one.

\documentclass{article}      % Specifies the document class

                             % The preamble begins here.
\title{An Example Document}  % Declares the document's title.
\author{Leslie Lamport}      % Declares the author's name.
\date{January 21, 1994}      % Deleting this command produces today's date.

\newcommand{\ip}[2]{(#1, #2)}
                             % Defines \ip{arg1}{arg2} to mean
                             % (arg1, arg2).

%\newcommand{\ip}[2]{\langle #1 | #2\rangle}
                             % This is an alternative definition of
                             % \ip that is commented out.

\begin{document}             % End of preamble and beginning of text.

\maketitle                   % Produces the title.

This is an example input file.  Comparing it with
the output it generates can show you how to
produce a simple document of your own.

\section{Ordinary Text}      % Produces section heading.  Lower-level
                             % sections are begun with similar
                             % \subsection and \subsubsection commands.

The ends  of words and sentences are marked
  by   spaces. It  doesn't matter how many
spaces    you type; one is as good as 100.  The
end of   a line counts as a space.

One   or more   blank lines denote the  end
of  a paragraph.

Since any number of consecutive spaces are treated
like a single one, the formatting of the input
file makes no difference to
      \LaTeX,                % The \LaTeX command generates the LaTeX logo.
but it makes a difference to you.  When you use
\LaTeX, making your input file as easy to read
as possible will be a great help as you write
your document and when you change it.  This sample
file shows how you can add comments to your own input
file.

Because printing is different from typewriting,
there are a number of things that you have to do
differently when preparing an input file than if
you were just typing the document directly.
Quotation marks like
       ``this''
have to be handled specially, as do quotes within
quotes:
       ``\,`this'            % \, separates the double and single quote.
        is what I just
        wrote, not  `that'\,''.

Dashes come in three sizes: an
       intra-word
dash, a medium dash for number ranges like
       1--2,
and a punctuation
       dash---like
this.

A sentence-ending space should be larger than the
space between words within a sentence.  You
sometimes have to type special commands in
conjunction with punctuation characters to get
this right, as in the following sentence.
       Gnats, gnus, etc.\ all  % `\ ' makes an inter-word space.
       begin with G\@.         % \@ marks end-of-sentence punctuation.
You should check the spaces after periods when
reading your output to make sure you haven't
forgotten any special cases.  Generating an
ellipsis
       \ldots\               % `\ ' is needed after `\ldots' because TeX
                             % ignores spaces after command names like \ldots
                             % made from \ + letters.
                             %
                             % Note how a `%' character causes TeX to ignore
                             % the end of the input line, so these blank lines
                             % do not start a new paragraph.
                             %
with the right spacing around the periods requires
a special command.

\LaTeX\ interprets some common characters as
commands, so you must type special commands to
generate them.  These characters include the
following:
       \$ \& \% \# \{ and \}.

In printing, text is usually emphasized with an
       \emph{italic}
type style.

\begin{em}
   A long segment of text can also be emphasized
   in this way.  Text within such a segment can be
   given \emph{additional} emphasis.
\end{em}

It is sometimes necessary to prevent \LaTeX\ from
breaking a line where it might otherwise do so.
This may be at a space, as between the ``Mr.''\ and
``Jones'' in
       ``Mr.~Jones'',        % ~ produces an unbreakable interword space.
or within a word---especially when the word is a
symbol like
       \mbox{\emph{itemnum}}
that makes little sense when hyphenated across
lines.

Footnotes\footnote{This is an example of a footnote.}
pose no problem.

\LaTeX\ is good at typesetting mathematical formulas
like
       \( x-3y + z = 7 \)
or
       \( a_{1} > x^{2n} + y^{2n} > x' \)
or
       \( \ip{A}{B} = \sum_{i} a_{i} b_{i} \).
The spaces you type in a formula are
ignored.  Remember that a letter like
       $x$                   % $ ... $  and  \( ... \)  are equivalent
is a formula when it denotes a mathematical
symbol, and it should be typed as one.

\section{Displayed Text}

Text is displayed by indenting it from the left
margin.  Quotations are commonly displayed.  There
are short quotations
\begin{quote}
   This is a short quotation.  It consists of a
   single paragraph of text.  See how it is formatted.
\end{quote}
and longer ones.
\begin{quotation}
   This is a longer quotation.  It consists of two
   paragraphs of text, neither of which are
   particularly interesting.

   This is the second paragraph of the quotation.  It
   is just as dull as the first paragraph.
\end{quotation}
Another frequently-displayed structure is a list.
The following is an example of an \emph{itemized}
list.
\begin{itemize}
   \item This is the first item of an itemized list.
         Each item in the list is marked with a ``tick''.
         You don't have to worry about what kind of tick
         mark is used.

   \item This is the second item of the list.  It
         contains another list nested inside it.  The inner
         list is an \emph{enumerated} list.
         \begin{enumerate}
            \item This is the first item of an enumerated
                  list that is nested within the itemized list.

            \item This is the second item of the inner list.
                  \LaTeX\ allows you to nest lists deeper than
                  you really should.
         \end{enumerate}
         This is the rest of the second item of the outer
         list.  It is no more interesting than any other
         part of the item.
   \item This is the third item of the list.
\end{itemize}
You can even display poetry.
\begin{verse}
   There is an environment
    for verse \\             % The \\ command separates lines
   Whose features some poets % within a stanza.
   will curse.

                             % One or more blank lines separate stanzas.

   For instead of making\\
   Them do \emph{all} line breaking, \\
   It allows them to put too many words on a line when they'd rather be
   forced to be terse.
\end{verse}

Mathematical formulas may also be displayed.  A
displayed formula
is
one-line long; multiline
formulas require special formatting instructions.
   \[  \ip{\Gamma}{\psi'} = x'' + y^{2} + z_{i}^{n}\]
Don't start a paragraph with a displayed equation,
nor make one a paragraph by itself.

\end{document}               % End of document.

\section{Bench Huge Section 323}
% synthetic bench section 323
% This is a sample LaTeX input file.  (Version of 12 August 2004.)
%
% A '%' character causes TeX to ignore all remaining text on the line,
% and is used for comments like this one.

\documentclass{article}      % Specifies the document class

                             % The preamble begins here.
\title{An Example Document}  % Declares the document's title.
\author{Leslie Lamport}      % Declares the author's name.
\date{January 21, 1994}      % Deleting this command produces today's date.

\newcommand{\ip}[2]{(#1, #2)}
                             % Defines \ip{arg1}{arg2} to mean
                             % (arg1, arg2).

%\newcommand{\ip}[2]{\langle #1 | #2\rangle}
                             % This is an alternative definition of
                             % \ip that is commented out.

\begin{document}             % End of preamble and beginning of text.

\maketitle                   % Produces the title.

This is an example input file.  Comparing it with
the output it generates can show you how to
produce a simple document of your own.

\section{Ordinary Text}      % Produces section heading.  Lower-level
                             % sections are begun with similar
                             % \subsection and \subsubsection commands.

The ends  of words and sentences are marked
  by   spaces. It  doesn't matter how many
spaces    you type; one is as good as 100.  The
end of   a line counts as a space.

One   or more   blank lines denote the  end
of  a paragraph.

Since any number of consecutive spaces are treated
like a single one, the formatting of the input
file makes no difference to
      \LaTeX,                % The \LaTeX command generates the LaTeX logo.
but it makes a difference to you.  When you use
\LaTeX, making your input file as easy to read
as possible will be a great help as you write
your document and when you change it.  This sample
file shows how you can add comments to your own input
file.

Because printing is different from typewriting,
there are a number of things that you have to do
differently when preparing an input file than if
you were just typing the document directly.
Quotation marks like
       ``this''
have to be handled specially, as do quotes within
quotes:
       ``\,`this'            % \, separates the double and single quote.
        is what I just
        wrote, not  `that'\,''.

Dashes come in three sizes: an
       intra-word
dash, a medium dash for number ranges like
       1--2,
and a punctuation
       dash---like
this.

A sentence-ending space should be larger than the
space between words within a sentence.  You
sometimes have to type special commands in
conjunction with punctuation characters to get
this right, as in the following sentence.
       Gnats, gnus, etc.\ all  % `\ ' makes an inter-word space.
       begin with G\@.         % \@ marks end-of-sentence punctuation.
You should check the spaces after periods when
reading your output to make sure you haven't
forgotten any special cases.  Generating an
ellipsis
       \ldots\               % `\ ' is needed after `\ldots' because TeX
                             % ignores spaces after command names like \ldots
                             % made from \ + letters.
                             %
                             % Note how a `%' character causes TeX to ignore
                             % the end of the input line, so these blank lines
                             % do not start a new paragraph.
                             %
with the right spacing around the periods requires
a special command.

\LaTeX\ interprets some common characters as
commands, so you must type special commands to
generate them.  These characters include the
following:
       \$ \& \% \# \{ and \}.

In printing, text is usually emphasized with an
       \emph{italic}
type style.

\begin{em}
   A long segment of text can also be emphasized
   in this way.  Text within such a segment can be
   given \emph{additional} emphasis.
\end{em}

It is sometimes necessary to prevent \LaTeX\ from
breaking a line where it might otherwise do so.
This may be at a space, as between the ``Mr.''\ and
``Jones'' in
       ``Mr.~Jones'',        % ~ produces an unbreakable interword space.
or within a word---especially when the word is a
symbol like
       \mbox{\emph{itemnum}}
that makes little sense when hyphenated across
lines.

Footnotes\footnote{This is an example of a footnote.}
pose no problem.

\LaTeX\ is good at typesetting mathematical formulas
like
       \( x-3y + z = 7 \)
or
       \( a_{1} > x^{2n} + y^{2n} > x' \)
or
       \( \ip{A}{B} = \sum_{i} a_{i} b_{i} \).
The spaces you type in a formula are
ignored.  Remember that a letter like
       $x$                   % $ ... $  and  \( ... \)  are equivalent
is a formula when it denotes a mathematical
symbol, and it should be typed as one.

\section{Displayed Text}

Text is displayed by indenting it from the left
margin.  Quotations are commonly displayed.  There
are short quotations
\begin{quote}
   This is a short quotation.  It consists of a
   single paragraph of text.  See how it is formatted.
\end{quote}
and longer ones.
\begin{quotation}
   This is a longer quotation.  It consists of two
   paragraphs of text, neither of which are
   particularly interesting.

   This is the second paragraph of the quotation.  It
   is just as dull as the first paragraph.
\end{quotation}
Another frequently-displayed structure is a list.
The following is an example of an \emph{itemized}
list.
\begin{itemize}
   \item This is the first item of an itemized list.
         Each item in the list is marked with a ``tick''.
         You don't have to worry about what kind of tick
         mark is used.

   \item This is the second item of the list.  It
         contains another list nested inside it.  The inner
         list is an \emph{enumerated} list.
         \begin{enumerate}
            \item This is the first item of an enumerated
                  list that is nested within the itemized list.

            \item This is the second item of the inner list.
                  \LaTeX\ allows you to nest lists deeper than
                  you really should.
         \end{enumerate}
         This is the rest of the second item of the outer
         list.  It is no more interesting than any other
         part of the item.
   \item This is the third item of the list.
\end{itemize}
You can even display poetry.
\begin{verse}
   There is an environment
    for verse \\             % The \\ command separates lines
   Whose features some poets % within a stanza.
   will curse.

                             % One or more blank lines separate stanzas.

   For instead of making\\
   Them do \emph{all} line breaking, \\
   It allows them to put too many words on a line when they'd rather be
   forced to be terse.
\end{verse}

Mathematical formulas may also be displayed.  A
displayed formula
is
one-line long; multiline
formulas require special formatting instructions.
   \[  \ip{\Gamma}{\psi'} = x'' + y^{2} + z_{i}^{n}\]
Don't start a paragraph with a displayed equation,
nor make one a paragraph by itself.

\end{document}               % End of document.

\section{Bench Huge Section 324}
% synthetic bench section 324
% This is a sample LaTeX input file.  (Version of 12 August 2004.)
%
% A '%' character causes TeX to ignore all remaining text on the line,
% and is used for comments like this one.

\documentclass{article}      % Specifies the document class

                             % The preamble begins here.
\title{An Example Document}  % Declares the document's title.
\author{Leslie Lamport}      % Declares the author's name.
\date{January 21, 1994}      % Deleting this command produces today's date.

\newcommand{\ip}[2]{(#1, #2)}
                             % Defines \ip{arg1}{arg2} to mean
                             % (arg1, arg2).

%\newcommand{\ip}[2]{\langle #1 | #2\rangle}
                             % This is an alternative definition of
                             % \ip that is commented out.

\begin{document}             % End of preamble and beginning of text.

\maketitle                   % Produces the title.

This is an example input file.  Comparing it with
the output it generates can show you how to
produce a simple document of your own.

\section{Ordinary Text}      % Produces section heading.  Lower-level
                             % sections are begun with similar
                             % \subsection and \subsubsection commands.

The ends  of words and sentences are marked
  by   spaces. It  doesn't matter how many
spaces    you type; one is as good as 100.  The
end of   a line counts as a space.

One   or more   blank lines denote the  end
of  a paragraph.

Since any number of consecutive spaces are treated
like a single one, the formatting of the input
file makes no difference to
      \LaTeX,                % The \LaTeX command generates the LaTeX logo.
but it makes a difference to you.  When you use
\LaTeX, making your input file as easy to read
as possible will be a great help as you write
your document and when you change it.  This sample
file shows how you can add comments to your own input
file.

Because printing is different from typewriting,
there are a number of things that you have to do
differently when preparing an input file than if
you were just typing the document directly.
Quotation marks like
       ``this''
have to be handled specially, as do quotes within
quotes:
       ``\,`this'            % \, separates the double and single quote.
        is what I just
        wrote, not  `that'\,''.

Dashes come in three sizes: an
       intra-word
dash, a medium dash for number ranges like
       1--2,
and a punctuation
       dash---like
this.

A sentence-ending space should be larger than the
space between words within a sentence.  You
sometimes have to type special commands in
conjunction with punctuation characters to get
this right, as in the following sentence.
       Gnats, gnus, etc.\ all  % `\ ' makes an inter-word space.
       begin with G\@.         % \@ marks end-of-sentence punctuation.
You should check the spaces after periods when
reading your output to make sure you haven't
forgotten any special cases.  Generating an
ellipsis
       \ldots\               % `\ ' is needed after `\ldots' because TeX
                             % ignores spaces after command names like \ldots
                             % made from \ + letters.
                             %
                             % Note how a `%' character causes TeX to ignore
                             % the end of the input line, so these blank lines
                             % do not start a new paragraph.
                             %
with the right spacing around the periods requires
a special command.

\LaTeX\ interprets some common characters as
commands, so you must type special commands to
generate them.  These characters include the
following:
       \$ \& \% \# \{ and \}.

In printing, text is usually emphasized with an
       \emph{italic}
type style.

\begin{em}
   A long segment of text can also be emphasized
   in this way.  Text within such a segment can be
   given \emph{additional} emphasis.
\end{em}

It is sometimes necessary to prevent \LaTeX\ from
breaking a line where it might otherwise do so.
This may be at a space, as between the ``Mr.''\ and
``Jones'' in
       ``Mr.~Jones'',        % ~ produces an unbreakable interword space.
or within a word---especially when the word is a
symbol like
       \mbox{\emph{itemnum}}
that makes little sense when hyphenated across
lines.

Footnotes\footnote{This is an example of a footnote.}
pose no problem.

\LaTeX\ is good at typesetting mathematical formulas
like
       \( x-3y + z = 7 \)
or
       \( a_{1} > x^{2n} + y^{2n} > x' \)
or
       \( \ip{A}{B} = \sum_{i} a_{i} b_{i} \).
The spaces you type in a formula are
ignored.  Remember that a letter like
       $x$                   % $ ... $  and  \( ... \)  are equivalent
is a formula when it denotes a mathematical
symbol, and it should be typed as one.

\section{Displayed Text}

Text is displayed by indenting it from the left
margin.  Quotations are commonly displayed.  There
are short quotations
\begin{quote}
   This is a short quotation.  It consists of a
   single paragraph of text.  See how it is formatted.
\end{quote}
and longer ones.
\begin{quotation}
   This is a longer quotation.  It consists of two
   paragraphs of text, neither of which are
   particularly interesting.

   This is the second paragraph of the quotation.  It
   is just as dull as the first paragraph.
\end{quotation}
Another frequently-displayed structure is a list.
The following is an example of an \emph{itemized}
list.
\begin{itemize}
   \item This is the first item of an itemized list.
         Each item in the list is marked with a ``tick''.
         You don't have to worry about what kind of tick
         mark is used.

   \item This is the second item of the list.  It
         contains another list nested inside it.  The inner
         list is an \emph{enumerated} list.
         \begin{enumerate}
            \item This is the first item of an enumerated
                  list that is nested within the itemized list.

            \item This is the second item of the inner list.
                  \LaTeX\ allows you to nest lists deeper than
                  you really should.
         \end{enumerate}
         This is the rest of the second item of the outer
         list.  It is no more interesting than any other
         part of the item.
   \item This is the third item of the list.
\end{itemize}
You can even display poetry.
\begin{verse}
   There is an environment
    for verse \\             % The \\ command separates lines
   Whose features some poets % within a stanza.
   will curse.

                             % One or more blank lines separate stanzas.

   For instead of making\\
   Them do \emph{all} line breaking, \\
   It allows them to put too many words on a line when they'd rather be
   forced to be terse.
\end{verse}

Mathematical formulas may also be displayed.  A
displayed formula
is
one-line long; multiline
formulas require special formatting instructions.
   \[  \ip{\Gamma}{\psi'} = x'' + y^{2} + z_{i}^{n}\]
Don't start a paragraph with a displayed equation,
nor make one a paragraph by itself.

\end{document}               % End of document.

\section{Bench Huge Section 325}
% synthetic bench section 325
% This is a sample LaTeX input file.  (Version of 12 August 2004.)
%
% A '%' character causes TeX to ignore all remaining text on the line,
% and is used for comments like this one.

\documentclass{article}      % Specifies the document class

                             % The preamble begins here.
\title{An Example Document}  % Declares the document's title.
\author{Leslie Lamport}      % Declares the author's name.
\date{January 21, 1994}      % Deleting this command produces today's date.

\newcommand{\ip}[2]{(#1, #2)}
                             % Defines \ip{arg1}{arg2} to mean
                             % (arg1, arg2).

%\newcommand{\ip}[2]{\langle #1 | #2\rangle}
                             % This is an alternative definition of
                             % \ip that is commented out.

\begin{document}             % End of preamble and beginning of text.

\maketitle                   % Produces the title.

This is an example input file.  Comparing it with
the output it generates can show you how to
produce a simple document of your own.

\section{Ordinary Text}      % Produces section heading.  Lower-level
                             % sections are begun with similar
                             % \subsection and \subsubsection commands.

The ends  of words and sentences are marked
  by   spaces. It  doesn't matter how many
spaces    you type; one is as good as 100.  The
end of   a line counts as a space.

One   or more   blank lines denote the  end
of  a paragraph.

Since any number of consecutive spaces are treated
like a single one, the formatting of the input
file makes no difference to
      \LaTeX,                % The \LaTeX command generates the LaTeX logo.
but it makes a difference to you.  When you use
\LaTeX, making your input file as easy to read
as possible will be a great help as you write
your document and when you change it.  This sample
file shows how you can add comments to your own input
file.

Because printing is different from typewriting,
there are a number of things that you have to do
differently when preparing an input file than if
you were just typing the document directly.
Quotation marks like
       ``this''
have to be handled specially, as do quotes within
quotes:
       ``\,`this'            % \, separates the double and single quote.
        is what I just
        wrote, not  `that'\,''.

Dashes come in three sizes: an
       intra-word
dash, a medium dash for number ranges like
       1--2,
and a punctuation
       dash---like
this.

A sentence-ending space should be larger than the
space between words within a sentence.  You
sometimes have to type special commands in
conjunction with punctuation characters to get
this right, as in the following sentence.
       Gnats, gnus, etc.\ all  % `\ ' makes an inter-word space.
       begin with G\@.         % \@ marks end-of-sentence punctuation.
You should check the spaces after periods when
reading your output to make sure you haven't
forgotten any special cases.  Generating an
ellipsis
       \ldots\               % `\ ' is needed after `\ldots' because TeX
                             % ignores spaces after command names like \ldots
                             % made from \ + letters.
                             %
                             % Note how a `%' character causes TeX to ignore
                             % the end of the input line, so these blank lines
                             % do not start a new paragraph.
                             %
with the right spacing around the periods requires
a special command.

\LaTeX\ interprets some common characters as
commands, so you must type special commands to
generate them.  These characters include the
following:
       \$ \& \% \# \{ and \}.

In printing, text is usually emphasized with an
       \emph{italic}
type style.

\begin{em}
   A long segment of text can also be emphasized
   in this way.  Text within such a segment can be
   given \emph{additional} emphasis.
\end{em}

It is sometimes necessary to prevent \LaTeX\ from
breaking a line where it might otherwise do so.
This may be at a space, as between the ``Mr.''\ and
``Jones'' in
       ``Mr.~Jones'',        % ~ produces an unbreakable interword space.
or within a word---especially when the word is a
symbol like
       \mbox{\emph{itemnum}}
that makes little sense when hyphenated across
lines.

Footnotes\footnote{This is an example of a footnote.}
pose no problem.

\LaTeX\ is good at typesetting mathematical formulas
like
       \( x-3y + z = 7 \)
or
       \( a_{1} > x^{2n} + y^{2n} > x' \)
or
       \( \ip{A}{B} = \sum_{i} a_{i} b_{i} \).
The spaces you type in a formula are
ignored.  Remember that a letter like
       $x$                   % $ ... $  and  \( ... \)  are equivalent
is a formula when it denotes a mathematical
symbol, and it should be typed as one.

\section{Displayed Text}

Text is displayed by indenting it from the left
margin.  Quotations are commonly displayed.  There
are short quotations
\begin{quote}
   This is a short quotation.  It consists of a
   single paragraph of text.  See how it is formatted.
\end{quote}
and longer ones.
\begin{quotation}
   This is a longer quotation.  It consists of two
   paragraphs of text, neither of which are
   particularly interesting.

   This is the second paragraph of the quotation.  It
   is just as dull as the first paragraph.
\end{quotation}
Another frequently-displayed structure is a list.
The following is an example of an \emph{itemized}
list.
\begin{itemize}
   \item This is the first item of an itemized list.
         Each item in the list is marked with a ``tick''.
         You don't have to worry about what kind of tick
         mark is used.

   \item This is the second item of the list.  It
         contains another list nested inside it.  The inner
         list is an \emph{enumerated} list.
         \begin{enumerate}
            \item This is the first item of an enumerated
                  list that is nested within the itemized list.

            \item This is the second item of the inner list.
                  \LaTeX\ allows you to nest lists deeper than
                  you really should.
         \end{enumerate}
         This is the rest of the second item of the outer
         list.  It is no more interesting than any other
         part of the item.
   \item This is the third item of the list.
\end{itemize}
You can even display poetry.
\begin{verse}
   There is an environment
    for verse \\             % The \\ command separates lines
   Whose features some poets % within a stanza.
   will curse.

                             % One or more blank lines separate stanzas.

   For instead of making\\
   Them do \emph{all} line breaking, \\
   It allows them to put too many words on a line when they'd rather be
   forced to be terse.
\end{verse}

Mathematical formulas may also be displayed.  A
displayed formula
is
one-line long; multiline
formulas require special formatting instructions.
   \[  \ip{\Gamma}{\psi'} = x'' + y^{2} + z_{i}^{n}\]
Don't start a paragraph with a displayed equation,
nor make one a paragraph by itself.

\end{document}               % End of document.

\section{Bench Huge Section 326}
% synthetic bench section 326
% This is a sample LaTeX input file.  (Version of 12 August 2004.)
%
% A '%' character causes TeX to ignore all remaining text on the line,
% and is used for comments like this one.

\documentclass{article}      % Specifies the document class

                             % The preamble begins here.
\title{An Example Document}  % Declares the document's title.
\author{Leslie Lamport}      % Declares the author's name.
\date{January 21, 1994}      % Deleting this command produces today's date.

\newcommand{\ip}[2]{(#1, #2)}
                             % Defines \ip{arg1}{arg2} to mean
                             % (arg1, arg2).

%\newcommand{\ip}[2]{\langle #1 | #2\rangle}
                             % This is an alternative definition of
                             % \ip that is commented out.

\begin{document}             % End of preamble and beginning of text.

\maketitle                   % Produces the title.

This is an example input file.  Comparing it with
the output it generates can show you how to
produce a simple document of your own.

\section{Ordinary Text}      % Produces section heading.  Lower-level
                             % sections are begun with similar
                             % \subsection and \subsubsection commands.

The ends  of words and sentences are marked
  by   spaces. It  doesn't matter how many
spaces    you type; one is as good as 100.  The
end of   a line counts as a space.

One   or more   blank lines denote the  end
of  a paragraph.

Since any number of consecutive spaces are treated
like a single one, the formatting of the input
file makes no difference to
      \LaTeX,                % The \LaTeX command generates the LaTeX logo.
but it makes a difference to you.  When you use
\LaTeX, making your input file as easy to read
as possible will be a great help as you write
your document and when you change it.  This sample
file shows how you can add comments to your own input
file.

Because printing is different from typewriting,
there are a number of things that you have to do
differently when preparing an input file than if
you were just typing the document directly.
Quotation marks like
       ``this''
have to be handled specially, as do quotes within
quotes:
       ``\,`this'            % \, separates the double and single quote.
        is what I just
        wrote, not  `that'\,''.

Dashes come in three sizes: an
       intra-word
dash, a medium dash for number ranges like
       1--2,
and a punctuation
       dash---like
this.

A sentence-ending space should be larger than the
space between words within a sentence.  You
sometimes have to type special commands in
conjunction with punctuation characters to get
this right, as in the following sentence.
       Gnats, gnus, etc.\ all  % `\ ' makes an inter-word space.
       begin with G\@.         % \@ marks end-of-sentence punctuation.
You should check the spaces after periods when
reading your output to make sure you haven't
forgotten any special cases.  Generating an
ellipsis
       \ldots\               % `\ ' is needed after `\ldots' because TeX
                             % ignores spaces after command names like \ldots
                             % made from \ + letters.
                             %
                             % Note how a `%' character causes TeX to ignore
                             % the end of the input line, so these blank lines
                             % do not start a new paragraph.
                             %
with the right spacing around the periods requires
a special command.

\LaTeX\ interprets some common characters as
commands, so you must type special commands to
generate them.  These characters include the
following:
       \$ \& \% \# \{ and \}.

In printing, text is usually emphasized with an
       \emph{italic}
type style.

\begin{em}
   A long segment of text can also be emphasized
   in this way.  Text within such a segment can be
   given \emph{additional} emphasis.
\end{em}

It is sometimes necessary to prevent \LaTeX\ from
breaking a line where it might otherwise do so.
This may be at a space, as between the ``Mr.''\ and
``Jones'' in
       ``Mr.~Jones'',        % ~ produces an unbreakable interword space.
or within a word---especially when the word is a
symbol like
       \mbox{\emph{itemnum}}
that makes little sense when hyphenated across
lines.

Footnotes\footnote{This is an example of a footnote.}
pose no problem.

\LaTeX\ is good at typesetting mathematical formulas
like
       \( x-3y + z = 7 \)
or
       \( a_{1} > x^{2n} + y^{2n} > x' \)
or
       \( \ip{A}{B} = \sum_{i} a_{i} b_{i} \).
The spaces you type in a formula are
ignored.  Remember that a letter like
       $x$                   % $ ... $  and  \( ... \)  are equivalent
is a formula when it denotes a mathematical
symbol, and it should be typed as one.

\section{Displayed Text}

Text is displayed by indenting it from the left
margin.  Quotations are commonly displayed.  There
are short quotations
\begin{quote}
   This is a short quotation.  It consists of a
   single paragraph of text.  See how it is formatted.
\end{quote}
and longer ones.
\begin{quotation}
   This is a longer quotation.  It consists of two
   paragraphs of text, neither of which are
   particularly interesting.

   This is the second paragraph of the quotation.  It
   is just as dull as the first paragraph.
\end{quotation}
Another frequently-displayed structure is a list.
The following is an example of an \emph{itemized}
list.
\begin{itemize}
   \item This is the first item of an itemized list.
         Each item in the list is marked with a ``tick''.
         You don't have to worry about what kind of tick
         mark is used.

   \item This is the second item of the list.  It
         contains another list nested inside it.  The inner
         list is an \emph{enumerated} list.
         \begin{enumerate}
            \item This is the first item of an enumerated
                  list that is nested within the itemized list.

            \item This is the second item of the inner list.
                  \LaTeX\ allows you to nest lists deeper than
                  you really should.
         \end{enumerate}
         This is the rest of the second item of the outer
         list.  It is no more interesting than any other
         part of the item.
   \item This is the third item of the list.
\end{itemize}
You can even display poetry.
\begin{verse}
   There is an environment
    for verse \\             % The \\ command separates lines
   Whose features some poets % within a stanza.
   will curse.

                             % One or more blank lines separate stanzas.

   For instead of making\\
   Them do \emph{all} line breaking, \\
   It allows them to put too many words on a line when they'd rather be
   forced to be terse.
\end{verse}

Mathematical formulas may also be displayed.  A
displayed formula
is
one-line long; multiline
formulas require special formatting instructions.
   \[  \ip{\Gamma}{\psi'} = x'' + y^{2} + z_{i}^{n}\]
Don't start a paragraph with a displayed equation,
nor make one a paragraph by itself.

\end{document}               % End of document.

\section{Bench Huge Section 327}
% synthetic bench section 327
% This is a sample LaTeX input file.  (Version of 12 August 2004.)
%
% A '%' character causes TeX to ignore all remaining text on the line,
% and is used for comments like this one.

\documentclass{article}      % Specifies the document class

                             % The preamble begins here.
\title{An Example Document}  % Declares the document's title.
\author{Leslie Lamport}      % Declares the author's name.
\date{January 21, 1994}      % Deleting this command produces today's date.

\newcommand{\ip}[2]{(#1, #2)}
                             % Defines \ip{arg1}{arg2} to mean
                             % (arg1, arg2).

%\newcommand{\ip}[2]{\langle #1 | #2\rangle}
                             % This is an alternative definition of
                             % \ip that is commented out.

\begin{document}             % End of preamble and beginning of text.

\maketitle                   % Produces the title.

This is an example input file.  Comparing it with
the output it generates can show you how to
produce a simple document of your own.

\section{Ordinary Text}      % Produces section heading.  Lower-level
                             % sections are begun with similar
                             % \subsection and \subsubsection commands.

The ends  of words and sentences are marked
  by   spaces. It  doesn't matter how many
spaces    you type; one is as good as 100.  The
end of   a line counts as a space.

One   or more   blank lines denote the  end
of  a paragraph.

Since any number of consecutive spaces are treated
like a single one, the formatting of the input
file makes no difference to
      \LaTeX,                % The \LaTeX command generates the LaTeX logo.
but it makes a difference to you.  When you use
\LaTeX, making your input file as easy to read
as possible will be a great help as you write
your document and when you change it.  This sample
file shows how you can add comments to your own input
file.

Because printing is different from typewriting,
there are a number of things that you have to do
differently when preparing an input file than if
you were just typing the document directly.
Quotation marks like
       ``this''
have to be handled specially, as do quotes within
quotes:
       ``\,`this'            % \, separates the double and single quote.
        is what I just
        wrote, not  `that'\,''.

Dashes come in three sizes: an
       intra-word
dash, a medium dash for number ranges like
       1--2,
and a punctuation
       dash---like
this.

A sentence-ending space should be larger than the
space between words within a sentence.  You
sometimes have to type special commands in
conjunction with punctuation characters to get
this right, as in the following sentence.
       Gnats, gnus, etc.\ all  % `\ ' makes an inter-word space.
       begin with G\@.         % \@ marks end-of-sentence punctuation.
You should check the spaces after periods when
reading your output to make sure you haven't
forgotten any special cases.  Generating an
ellipsis
       \ldots\               % `\ ' is needed after `\ldots' because TeX
                             % ignores spaces after command names like \ldots
                             % made from \ + letters.
                             %
                             % Note how a `%' character causes TeX to ignore
                             % the end of the input line, so these blank lines
                             % do not start a new paragraph.
                             %
with the right spacing around the periods requires
a special command.

\LaTeX\ interprets some common characters as
commands, so you must type special commands to
generate them.  These characters include the
following:
       \$ \& \% \# \{ and \}.

In printing, text is usually emphasized with an
       \emph{italic}
type style.

\begin{em}
   A long segment of text can also be emphasized
   in this way.  Text within such a segment can be
   given \emph{additional} emphasis.
\end{em}

It is sometimes necessary to prevent \LaTeX\ from
breaking a line where it might otherwise do so.
This may be at a space, as between the ``Mr.''\ and
``Jones'' in
       ``Mr.~Jones'',        % ~ produces an unbreakable interword space.
or within a word---especially when the word is a
symbol like
       \mbox{\emph{itemnum}}
that makes little sense when hyphenated across
lines.

Footnotes\footnote{This is an example of a footnote.}
pose no problem.

\LaTeX\ is good at typesetting mathematical formulas
like
       \( x-3y + z = 7 \)
or
       \( a_{1} > x^{2n} + y^{2n} > x' \)
or
       \( \ip{A}{B} = \sum_{i} a_{i} b_{i} \).
The spaces you type in a formula are
ignored.  Remember that a letter like
       $x$                   % $ ... $  and  \( ... \)  are equivalent
is a formula when it denotes a mathematical
symbol, and it should be typed as one.

\section{Displayed Text}

Text is displayed by indenting it from the left
margin.  Quotations are commonly displayed.  There
are short quotations
\begin{quote}
   This is a short quotation.  It consists of a
   single paragraph of text.  See how it is formatted.
\end{quote}
and longer ones.
\begin{quotation}
   This is a longer quotation.  It consists of two
   paragraphs of text, neither of which are
   particularly interesting.

   This is the second paragraph of the quotation.  It
   is just as dull as the first paragraph.
\end{quotation}
Another frequently-displayed structure is a list.
The following is an example of an \emph{itemized}
list.
\begin{itemize}
   \item This is the first item of an itemized list.
         Each item in the list is marked with a ``tick''.
         You don't have to worry about what kind of tick
         mark is used.

   \item This is the second item of the list.  It
         contains another list nested inside it.  The inner
         list is an \emph{enumerated} list.
         \begin{enumerate}
            \item This is the first item of an enumerated
                  list that is nested within the itemized list.

            \item This is the second item of the inner list.
                  \LaTeX\ allows you to nest lists deeper than
                  you really should.
         \end{enumerate}
         This is the rest of the second item of the outer
         list.  It is no more interesting than any other
         part of the item.
   \item This is the third item of the list.
\end{itemize}
You can even display poetry.
\begin{verse}
   There is an environment
    for verse \\             % The \\ command separates lines
   Whose features some poets % within a stanza.
   will curse.

                             % One or more blank lines separate stanzas.

   For instead of making\\
   Them do \emph{all} line breaking, \\
   It allows them to put too many words on a line when they'd rather be
   forced to be terse.
\end{verse}

Mathematical formulas may also be displayed.  A
displayed formula
is
one-line long; multiline
formulas require special formatting instructions.
   \[  \ip{\Gamma}{\psi'} = x'' + y^{2} + z_{i}^{n}\]
Don't start a paragraph with a displayed equation,
nor make one a paragraph by itself.

\end{document}               % End of document.

\section{Bench Huge Section 328}
% synthetic bench section 328
% This is a sample LaTeX input file.  (Version of 12 August 2004.)
%
% A '%' character causes TeX to ignore all remaining text on the line,
% and is used for comments like this one.

\documentclass{article}      % Specifies the document class

                             % The preamble begins here.
\title{An Example Document}  % Declares the document's title.
\author{Leslie Lamport}      % Declares the author's name.
\date{January 21, 1994}      % Deleting this command produces today's date.

\newcommand{\ip}[2]{(#1, #2)}
                             % Defines \ip{arg1}{arg2} to mean
                             % (arg1, arg2).

%\newcommand{\ip}[2]{\langle #1 | #2\rangle}
                             % This is an alternative definition of
                             % \ip that is commented out.

\begin{document}             % End of preamble and beginning of text.

\maketitle                   % Produces the title.

This is an example input file.  Comparing it with
the output it generates can show you how to
produce a simple document of your own.

\section{Ordinary Text}      % Produces section heading.  Lower-level
                             % sections are begun with similar
                             % \subsection and \subsubsection commands.

The ends  of words and sentences are marked
  by   spaces. It  doesn't matter how many
spaces    you type; one is as good as 100.  The
end of   a line counts as a space.

One   or more   blank lines denote the  end
of  a paragraph.

Since any number of consecutive spaces are treated
like a single one, the formatting of the input
file makes no difference to
      \LaTeX,                % The \LaTeX command generates the LaTeX logo.
but it makes a difference to you.  When you use
\LaTeX, making your input file as easy to read
as possible will be a great help as you write
your document and when you change it.  This sample
file shows how you can add comments to your own input
file.

Because printing is different from typewriting,
there are a number of things that you have to do
differently when preparing an input file than if
you were just typing the document directly.
Quotation marks like
       ``this''
have to be handled specially, as do quotes within
quotes:
       ``\,`this'            % \, separates the double and single quote.
        is what I just
        wrote, not  `that'\,''.

Dashes come in three sizes: an
       intra-word
dash, a medium dash for number ranges like
       1--2,
and a punctuation
       dash---like
this.

A sentence-ending space should be larger than the
space between words within a sentence.  You
sometimes have to type special commands in
conjunction with punctuation characters to get
this right, as in the following sentence.
       Gnats, gnus, etc.\ all  % `\ ' makes an inter-word space.
       begin with G\@.         % \@ marks end-of-sentence punctuation.
You should check the spaces after periods when
reading your output to make sure you haven't
forgotten any special cases.  Generating an
ellipsis
       \ldots\               % `\ ' is needed after `\ldots' because TeX
                             % ignores spaces after command names like \ldots
                             % made from \ + letters.
                             %
                             % Note how a `%' character causes TeX to ignore
                             % the end of the input line, so these blank lines
                             % do not start a new paragraph.
                             %
with the right spacing around the periods requires
a special command.

\LaTeX\ interprets some common characters as
commands, so you must type special commands to
generate them.  These characters include the
following:
       \$ \& \% \# \{ and \}.

In printing, text is usually emphasized with an
       \emph{italic}
type style.

\begin{em}
   A long segment of text can also be emphasized
   in this way.  Text within such a segment can be
   given \emph{additional} emphasis.
\end{em}

It is sometimes necessary to prevent \LaTeX\ from
breaking a line where it might otherwise do so.
This may be at a space, as between the ``Mr.''\ and
``Jones'' in
       ``Mr.~Jones'',        % ~ produces an unbreakable interword space.
or within a word---especially when the word is a
symbol like
       \mbox{\emph{itemnum}}
that makes little sense when hyphenated across
lines.

Footnotes\footnote{This is an example of a footnote.}
pose no problem.

\LaTeX\ is good at typesetting mathematical formulas
like
       \( x-3y + z = 7 \)
or
       \( a_{1} > x^{2n} + y^{2n} > x' \)
or
       \( \ip{A}{B} = \sum_{i} a_{i} b_{i} \).
The spaces you type in a formula are
ignored.  Remember that a letter like
       $x$                   % $ ... $  and  \( ... \)  are equivalent
is a formula when it denotes a mathematical
symbol, and it should be typed as one.

\section{Displayed Text}

Text is displayed by indenting it from the left
margin.  Quotations are commonly displayed.  There
are short quotations
\begin{quote}
   This is a short quotation.  It consists of a
   single paragraph of text.  See how it is formatted.
\end{quote}
and longer ones.
\begin{quotation}
   This is a longer quotation.  It consists of two
   paragraphs of text, neither of which are
   particularly interesting.

   This is the second paragraph of the quotation.  It
   is just as dull as the first paragraph.
\end{quotation}
Another frequently-displayed structure is a list.
The following is an example of an \emph{itemized}
list.
\begin{itemize}
   \item This is the first item of an itemized list.
         Each item in the list is marked with a ``tick''.
         You don't have to worry about what kind of tick
         mark is used.

   \item This is the second item of the list.  It
         contains another list nested inside it.  The inner
         list is an \emph{enumerated} list.
         \begin{enumerate}
            \item This is the first item of an enumerated
                  list that is nested within the itemized list.

            \item This is the second item of the inner list.
                  \LaTeX\ allows you to nest lists deeper than
                  you really should.
         \end{enumerate}
         This is the rest of the second item of the outer
         list.  It is no more interesting than any other
         part of the item.
   \item This is the third item of the list.
\end{itemize}
You can even display poetry.
\begin{verse}
   There is an environment
    for verse \\             % The \\ command separates lines
   Whose features some poets % within a stanza.
   will curse.

                             % One or more blank lines separate stanzas.

   For instead of making\\
   Them do \emph{all} line breaking, \\
   It allows them to put too many words on a line when they'd rather be
   forced to be terse.
\end{verse}

Mathematical formulas may also be displayed.  A
displayed formula
is
one-line long; multiline
formulas require special formatting instructions.
   \[  \ip{\Gamma}{\psi'} = x'' + y^{2} + z_{i}^{n}\]
Don't start a paragraph with a displayed equation,
nor make one a paragraph by itself.

\end{document}               % End of document.

\section{Bench Huge Section 329}
% synthetic bench section 329
% This is a sample LaTeX input file.  (Version of 12 August 2004.)
%
% A '%' character causes TeX to ignore all remaining text on the line,
% and is used for comments like this one.

\documentclass{article}      % Specifies the document class

                             % The preamble begins here.
\title{An Example Document}  % Declares the document's title.
\author{Leslie Lamport}      % Declares the author's name.
\date{January 21, 1994}      % Deleting this command produces today's date.

\newcommand{\ip}[2]{(#1, #2)}
                             % Defines \ip{arg1}{arg2} to mean
                             % (arg1, arg2).

%\newcommand{\ip}[2]{\langle #1 | #2\rangle}
                             % This is an alternative definition of
                             % \ip that is commented out.

\begin{document}             % End of preamble and beginning of text.

\maketitle                   % Produces the title.

This is an example input file.  Comparing it with
the output it generates can show you how to
produce a simple document of your own.

\section{Ordinary Text}      % Produces section heading.  Lower-level
                             % sections are begun with similar
                             % \subsection and \subsubsection commands.

The ends  of words and sentences are marked
  by   spaces. It  doesn't matter how many
spaces    you type; one is as good as 100.  The
end of   a line counts as a space.

One   or more   blank lines denote the  end
of  a paragraph.

Since any number of consecutive spaces are treated
like a single one, the formatting of the input
file makes no difference to
      \LaTeX,                % The \LaTeX command generates the LaTeX logo.
but it makes a difference to you.  When you use
\LaTeX, making your input file as easy to read
as possible will be a great help as you write
your document and when you change it.  This sample
file shows how you can add comments to your own input
file.

Because printing is different from typewriting,
there are a number of things that you have to do
differently when preparing an input file than if
you were just typing the document directly.
Quotation marks like
       ``this''
have to be handled specially, as do quotes within
quotes:
       ``\,`this'            % \, separates the double and single quote.
        is what I just
        wrote, not  `that'\,''.

Dashes come in three sizes: an
       intra-word
dash, a medium dash for number ranges like
       1--2,
and a punctuation
       dash---like
this.

A sentence-ending space should be larger than the
space between words within a sentence.  You
sometimes have to type special commands in
conjunction with punctuation characters to get
this right, as in the following sentence.
       Gnats, gnus, etc.\ all  % `\ ' makes an inter-word space.
       begin with G\@.         % \@ marks end-of-sentence punctuation.
You should check the spaces after periods when
reading your output to make sure you haven't
forgotten any special cases.  Generating an
ellipsis
       \ldots\               % `\ ' is needed after `\ldots' because TeX
                             % ignores spaces after command names like \ldots
                             % made from \ + letters.
                             %
                             % Note how a `%' character causes TeX to ignore
                             % the end of the input line, so these blank lines
                             % do not start a new paragraph.
                             %
with the right spacing around the periods requires
a special command.

\LaTeX\ interprets some common characters as
commands, so you must type special commands to
generate them.  These characters include the
following:
       \$ \& \% \# \{ and \}.

In printing, text is usually emphasized with an
       \emph{italic}
type style.

\begin{em}
   A long segment of text can also be emphasized
   in this way.  Text within such a segment can be
   given \emph{additional} emphasis.
\end{em}

It is sometimes necessary to prevent \LaTeX\ from
breaking a line where it might otherwise do so.
This may be at a space, as between the ``Mr.''\ and
``Jones'' in
       ``Mr.~Jones'',        % ~ produces an unbreakable interword space.
or within a word---especially when the word is a
symbol like
       \mbox{\emph{itemnum}}
that makes little sense when hyphenated across
lines.

Footnotes\footnote{This is an example of a footnote.}
pose no problem.

\LaTeX\ is good at typesetting mathematical formulas
like
       \( x-3y + z = 7 \)
or
       \( a_{1} > x^{2n} + y^{2n} > x' \)
or
       \( \ip{A}{B} = \sum_{i} a_{i} b_{i} \).
The spaces you type in a formula are
ignored.  Remember that a letter like
       $x$                   % $ ... $  and  \( ... \)  are equivalent
is a formula when it denotes a mathematical
symbol, and it should be typed as one.

\section{Displayed Text}

Text is displayed by indenting it from the left
margin.  Quotations are commonly displayed.  There
are short quotations
\begin{quote}
   This is a short quotation.  It consists of a
   single paragraph of text.  See how it is formatted.
\end{quote}
and longer ones.
\begin{quotation}
   This is a longer quotation.  It consists of two
   paragraphs of text, neither of which are
   particularly interesting.

   This is the second paragraph of the quotation.  It
   is just as dull as the first paragraph.
\end{quotation}
Another frequently-displayed structure is a list.
The following is an example of an \emph{itemized}
list.
\begin{itemize}
   \item This is the first item of an itemized list.
         Each item in the list is marked with a ``tick''.
         You don't have to worry about what kind of tick
         mark is used.

   \item This is the second item of the list.  It
         contains another list nested inside it.  The inner
         list is an \emph{enumerated} list.
         \begin{enumerate}
            \item This is the first item of an enumerated
                  list that is nested within the itemized list.

            \item This is the second item of the inner list.
                  \LaTeX\ allows you to nest lists deeper than
                  you really should.
         \end{enumerate}
         This is the rest of the second item of the outer
         list.  It is no more interesting than any other
         part of the item.
   \item This is the third item of the list.
\end{itemize}
You can even display poetry.
\begin{verse}
   There is an environment
    for verse \\             % The \\ command separates lines
   Whose features some poets % within a stanza.
   will curse.

                             % One or more blank lines separate stanzas.

   For instead of making\\
   Them do \emph{all} line breaking, \\
   It allows them to put too many words on a line when they'd rather be
   forced to be terse.
\end{verse}

Mathematical formulas may also be displayed.  A
displayed formula
is
one-line long; multiline
formulas require special formatting instructions.
   \[  \ip{\Gamma}{\psi'} = x'' + y^{2} + z_{i}^{n}\]
Don't start a paragraph with a displayed equation,
nor make one a paragraph by itself.

\end{document}               % End of document.

\section{Bench Huge Section 330}
% synthetic bench section 330
% This is a sample LaTeX input file.  (Version of 12 August 2004.)
%
% A '%' character causes TeX to ignore all remaining text on the line,
% and is used for comments like this one.

\documentclass{article}      % Specifies the document class

                             % The preamble begins here.
\title{An Example Document}  % Declares the document's title.
\author{Leslie Lamport}      % Declares the author's name.
\date{January 21, 1994}      % Deleting this command produces today's date.

\newcommand{\ip}[2]{(#1, #2)}
                             % Defines \ip{arg1}{arg2} to mean
                             % (arg1, arg2).

%\newcommand{\ip}[2]{\langle #1 | #2\rangle}
                             % This is an alternative definition of
                             % \ip that is commented out.

\begin{document}             % End of preamble and beginning of text.

\maketitle                   % Produces the title.

This is an example input file.  Comparing it with
the output it generates can show you how to
produce a simple document of your own.

\section{Ordinary Text}      % Produces section heading.  Lower-level
                             % sections are begun with similar
                             % \subsection and \subsubsection commands.

The ends  of words and sentences are marked
  by   spaces. It  doesn't matter how many
spaces    you type; one is as good as 100.  The
end of   a line counts as a space.

One   or more   blank lines denote the  end
of  a paragraph.

Since any number of consecutive spaces are treated
like a single one, the formatting of the input
file makes no difference to
      \LaTeX,                % The \LaTeX command generates the LaTeX logo.
but it makes a difference to you.  When you use
\LaTeX, making your input file as easy to read
as possible will be a great help as you write
your document and when you change it.  This sample
file shows how you can add comments to your own input
file.

Because printing is different from typewriting,
there are a number of things that you have to do
differently when preparing an input file than if
you were just typing the document directly.
Quotation marks like
       ``this''
have to be handled specially, as do quotes within
quotes:
       ``\,`this'            % \, separates the double and single quote.
        is what I just
        wrote, not  `that'\,''.

Dashes come in three sizes: an
       intra-word
dash, a medium dash for number ranges like
       1--2,
and a punctuation
       dash---like
this.

A sentence-ending space should be larger than the
space between words within a sentence.  You
sometimes have to type special commands in
conjunction with punctuation characters to get
this right, as in the following sentence.
       Gnats, gnus, etc.\ all  % `\ ' makes an inter-word space.
       begin with G\@.         % \@ marks end-of-sentence punctuation.
You should check the spaces after periods when
reading your output to make sure you haven't
forgotten any special cases.  Generating an
ellipsis
       \ldots\               % `\ ' is needed after `\ldots' because TeX
                             % ignores spaces after command names like \ldots
                             % made from \ + letters.
                             %
                             % Note how a `%' character causes TeX to ignore
                             % the end of the input line, so these blank lines
                             % do not start a new paragraph.
                             %
with the right spacing around the periods requires
a special command.

\LaTeX\ interprets some common characters as
commands, so you must type special commands to
generate them.  These characters include the
following:
       \$ \& \% \# \{ and \}.

In printing, text is usually emphasized with an
       \emph{italic}
type style.

\begin{em}
   A long segment of text can also be emphasized
   in this way.  Text within such a segment can be
   given \emph{additional} emphasis.
\end{em}

It is sometimes necessary to prevent \LaTeX\ from
breaking a line where it might otherwise do so.
This may be at a space, as between the ``Mr.''\ and
``Jones'' in
       ``Mr.~Jones'',        % ~ produces an unbreakable interword space.
or within a word---especially when the word is a
symbol like
       \mbox{\emph{itemnum}}
that makes little sense when hyphenated across
lines.

Footnotes\footnote{This is an example of a footnote.}
pose no problem.

\LaTeX\ is good at typesetting mathematical formulas
like
       \( x-3y + z = 7 \)
or
       \( a_{1} > x^{2n} + y^{2n} > x' \)
or
       \( \ip{A}{B} = \sum_{i} a_{i} b_{i} \).
The spaces you type in a formula are
ignored.  Remember that a letter like
       $x$                   % $ ... $  and  \( ... \)  are equivalent
is a formula when it denotes a mathematical
symbol, and it should be typed as one.

\section{Displayed Text}

Text is displayed by indenting it from the left
margin.  Quotations are commonly displayed.  There
are short quotations
\begin{quote}
   This is a short quotation.  It consists of a
   single paragraph of text.  See how it is formatted.
\end{quote}
and longer ones.
\begin{quotation}
   This is a longer quotation.  It consists of two
   paragraphs of text, neither of which are
   particularly interesting.

   This is the second paragraph of the quotation.  It
   is just as dull as the first paragraph.
\end{quotation}
Another frequently-displayed structure is a list.
The following is an example of an \emph{itemized}
list.
\begin{itemize}
   \item This is the first item of an itemized list.
         Each item in the list is marked with a ``tick''.
         You don't have to worry about what kind of tick
         mark is used.

   \item This is the second item of the list.  It
         contains another list nested inside it.  The inner
         list is an \emph{enumerated} list.
         \begin{enumerate}
            \item This is the first item of an enumerated
                  list that is nested within the itemized list.

            \item This is the second item of the inner list.
                  \LaTeX\ allows you to nest lists deeper than
                  you really should.
         \end{enumerate}
         This is the rest of the second item of the outer
         list.  It is no more interesting than any other
         part of the item.
   \item This is the third item of the list.
\end{itemize}
You can even display poetry.
\begin{verse}
   There is an environment
    for verse \\             % The \\ command separates lines
   Whose features some poets % within a stanza.
   will curse.

                             % One or more blank lines separate stanzas.

   For instead of making\\
   Them do \emph{all} line breaking, \\
   It allows them to put too many words on a line when they'd rather be
   forced to be terse.
\end{verse}

Mathematical formulas may also be displayed.  A
displayed formula
is
one-line long; multiline
formulas require special formatting instructions.
   \[  \ip{\Gamma}{\psi'} = x'' + y^{2} + z_{i}^{n}\]
Don't start a paragraph with a displayed equation,
nor make one a paragraph by itself.

\end{document}               % End of document.

\section{Bench Huge Section 331}
% synthetic bench section 331
% This is a sample LaTeX input file.  (Version of 12 August 2004.)
%
% A '%' character causes TeX to ignore all remaining text on the line,
% and is used for comments like this one.

\documentclass{article}      % Specifies the document class

                             % The preamble begins here.
\title{An Example Document}  % Declares the document's title.
\author{Leslie Lamport}      % Declares the author's name.
\date{January 21, 1994}      % Deleting this command produces today's date.

\newcommand{\ip}[2]{(#1, #2)}
                             % Defines \ip{arg1}{arg2} to mean
                             % (arg1, arg2).

%\newcommand{\ip}[2]{\langle #1 | #2\rangle}
                             % This is an alternative definition of
                             % \ip that is commented out.

\begin{document}             % End of preamble and beginning of text.

\maketitle                   % Produces the title.

This is an example input file.  Comparing it with
the output it generates can show you how to
produce a simple document of your own.

\section{Ordinary Text}      % Produces section heading.  Lower-level
                             % sections are begun with similar
                             % \subsection and \subsubsection commands.

The ends  of words and sentences are marked
  by   spaces. It  doesn't matter how many
spaces    you type; one is as good as 100.  The
end of   a line counts as a space.

One   or more   blank lines denote the  end
of  a paragraph.

Since any number of consecutive spaces are treated
like a single one, the formatting of the input
file makes no difference to
      \LaTeX,                % The \LaTeX command generates the LaTeX logo.
but it makes a difference to you.  When you use
\LaTeX, making your input file as easy to read
as possible will be a great help as you write
your document and when you change it.  This sample
file shows how you can add comments to your own input
file.

Because printing is different from typewriting,
there are a number of things that you have to do
differently when preparing an input file than if
you were just typing the document directly.
Quotation marks like
       ``this''
have to be handled specially, as do quotes within
quotes:
       ``\,`this'            % \, separates the double and single quote.
        is what I just
        wrote, not  `that'\,''.

Dashes come in three sizes: an
       intra-word
dash, a medium dash for number ranges like
       1--2,
and a punctuation
       dash---like
this.

A sentence-ending space should be larger than the
space between words within a sentence.  You
sometimes have to type special commands in
conjunction with punctuation characters to get
this right, as in the following sentence.
       Gnats, gnus, etc.\ all  % `\ ' makes an inter-word space.
       begin with G\@.         % \@ marks end-of-sentence punctuation.
You should check the spaces after periods when
reading your output to make sure you haven't
forgotten any special cases.  Generating an
ellipsis
       \ldots\               % `\ ' is needed after `\ldots' because TeX
                             % ignores spaces after command names like \ldots
                             % made from \ + letters.
                             %
                             % Note how a `%' character causes TeX to ignore
                             % the end of the input line, so these blank lines
                             % do not start a new paragraph.
                             %
with the right spacing around the periods requires
a special command.

\LaTeX\ interprets some common characters as
commands, so you must type special commands to
generate them.  These characters include the
following:
       \$ \& \% \# \{ and \}.

In printing, text is usually emphasized with an
       \emph{italic}
type style.

\begin{em}
   A long segment of text can also be emphasized
   in this way.  Text within such a segment can be
   given \emph{additional} emphasis.
\end{em}

It is sometimes necessary to prevent \LaTeX\ from
breaking a line where it might otherwise do so.
This may be at a space, as between the ``Mr.''\ and
``Jones'' in
       ``Mr.~Jones'',        % ~ produces an unbreakable interword space.
or within a word---especially when the word is a
symbol like
       \mbox{\emph{itemnum}}
that makes little sense when hyphenated across
lines.

Footnotes\footnote{This is an example of a footnote.}
pose no problem.

\LaTeX\ is good at typesetting mathematical formulas
like
       \( x-3y + z = 7 \)
or
       \( a_{1} > x^{2n} + y^{2n} > x' \)
or
       \( \ip{A}{B} = \sum_{i} a_{i} b_{i} \).
The spaces you type in a formula are
ignored.  Remember that a letter like
       $x$                   % $ ... $  and  \( ... \)  are equivalent
is a formula when it denotes a mathematical
symbol, and it should be typed as one.

\section{Displayed Text}

Text is displayed by indenting it from the left
margin.  Quotations are commonly displayed.  There
are short quotations
\begin{quote}
   This is a short quotation.  It consists of a
   single paragraph of text.  See how it is formatted.
\end{quote}
and longer ones.
\begin{quotation}
   This is a longer quotation.  It consists of two
   paragraphs of text, neither of which are
   particularly interesting.

   This is the second paragraph of the quotation.  It
   is just as dull as the first paragraph.
\end{quotation}
Another frequently-displayed structure is a list.
The following is an example of an \emph{itemized}
list.
\begin{itemize}
   \item This is the first item of an itemized list.
         Each item in the list is marked with a ``tick''.
         You don't have to worry about what kind of tick
         mark is used.

   \item This is the second item of the list.  It
         contains another list nested inside it.  The inner
         list is an \emph{enumerated} list.
         \begin{enumerate}
            \item This is the first item of an enumerated
                  list that is nested within the itemized list.

            \item This is the second item of the inner list.
                  \LaTeX\ allows you to nest lists deeper than
                  you really should.
         \end{enumerate}
         This is the rest of the second item of the outer
         list.  It is no more interesting than any other
         part of the item.
   \item This is the third item of the list.
\end{itemize}
You can even display poetry.
\begin{verse}
   There is an environment
    for verse \\             % The \\ command separates lines
   Whose features some poets % within a stanza.
   will curse.

                             % One or more blank lines separate stanzas.

   For instead of making\\
   Them do \emph{all} line breaking, \\
   It allows them to put too many words on a line when they'd rather be
   forced to be terse.
\end{verse}

Mathematical formulas may also be displayed.  A
displayed formula
is
one-line long; multiline
formulas require special formatting instructions.
   \[  \ip{\Gamma}{\psi'} = x'' + y^{2} + z_{i}^{n}\]
Don't start a paragraph with a displayed equation,
nor make one a paragraph by itself.

\end{document}               % End of document.

\section{Bench Huge Section 332}
% synthetic bench section 332
% This is a sample LaTeX input file.  (Version of 12 August 2004.)
%
% A '%' character causes TeX to ignore all remaining text on the line,
% and is used for comments like this one.

\documentclass{article}      % Specifies the document class

                             % The preamble begins here.
\title{An Example Document}  % Declares the document's title.
\author{Leslie Lamport}      % Declares the author's name.
\date{January 21, 1994}      % Deleting this command produces today's date.

\newcommand{\ip}[2]{(#1, #2)}
                             % Defines \ip{arg1}{arg2} to mean
                             % (arg1, arg2).

%\newcommand{\ip}[2]{\langle #1 | #2\rangle}
                             % This is an alternative definition of
                             % \ip that is commented out.

\begin{document}             % End of preamble and beginning of text.

\maketitle                   % Produces the title.

This is an example input file.  Comparing it with
the output it generates can show you how to
produce a simple document of your own.

\section{Ordinary Text}      % Produces section heading.  Lower-level
                             % sections are begun with similar
                             % \subsection and \subsubsection commands.

The ends  of words and sentences are marked
  by   spaces. It  doesn't matter how many
spaces    you type; one is as good as 100.  The
end of   a line counts as a space.

One   or more   blank lines denote the  end
of  a paragraph.

Since any number of consecutive spaces are treated
like a single one, the formatting of the input
file makes no difference to
      \LaTeX,                % The \LaTeX command generates the LaTeX logo.
but it makes a difference to you.  When you use
\LaTeX, making your input file as easy to read
as possible will be a great help as you write
your document and when you change it.  This sample
file shows how you can add comments to your own input
file.

Because printing is different from typewriting,
there are a number of things that you have to do
differently when preparing an input file than if
you were just typing the document directly.
Quotation marks like
       ``this''
have to be handled specially, as do quotes within
quotes:
       ``\,`this'            % \, separates the double and single quote.
        is what I just
        wrote, not  `that'\,''.

Dashes come in three sizes: an
       intra-word
dash, a medium dash for number ranges like
       1--2,
and a punctuation
       dash---like
this.

A sentence-ending space should be larger than the
space between words within a sentence.  You
sometimes have to type special commands in
conjunction with punctuation characters to get
this right, as in the following sentence.
       Gnats, gnus, etc.\ all  % `\ ' makes an inter-word space.
       begin with G\@.         % \@ marks end-of-sentence punctuation.
You should check the spaces after periods when
reading your output to make sure you haven't
forgotten any special cases.  Generating an
ellipsis
       \ldots\               % `\ ' is needed after `\ldots' because TeX
                             % ignores spaces after command names like \ldots
                             % made from \ + letters.
                             %
                             % Note how a `%' character causes TeX to ignore
                             % the end of the input line, so these blank lines
                             % do not start a new paragraph.
                             %
with the right spacing around the periods requires
a special command.

\LaTeX\ interprets some common characters as
commands, so you must type special commands to
generate them.  These characters include the
following:
       \$ \& \% \# \{ and \}.

In printing, text is usually emphasized with an
       \emph{italic}
type style.

\begin{em}
   A long segment of text can also be emphasized
   in this way.  Text within such a segment can be
   given \emph{additional} emphasis.
\end{em}

It is sometimes necessary to prevent \LaTeX\ from
breaking a line where it might otherwise do so.
This may be at a space, as between the ``Mr.''\ and
``Jones'' in
       ``Mr.~Jones'',        % ~ produces an unbreakable interword space.
or within a word---especially when the word is a
symbol like
       \mbox{\emph{itemnum}}
that makes little sense when hyphenated across
lines.

Footnotes\footnote{This is an example of a footnote.}
pose no problem.

\LaTeX\ is good at typesetting mathematical formulas
like
       \( x-3y + z = 7 \)
or
       \( a_{1} > x^{2n} + y^{2n} > x' \)
or
       \( \ip{A}{B} = \sum_{i} a_{i} b_{i} \).
The spaces you type in a formula are
ignored.  Remember that a letter like
       $x$                   % $ ... $  and  \( ... \)  are equivalent
is a formula when it denotes a mathematical
symbol, and it should be typed as one.

\section{Displayed Text}

Text is displayed by indenting it from the left
margin.  Quotations are commonly displayed.  There
are short quotations
\begin{quote}
   This is a short quotation.  It consists of a
   single paragraph of text.  See how it is formatted.
\end{quote}
and longer ones.
\begin{quotation}
   This is a longer quotation.  It consists of two
   paragraphs of text, neither of which are
   particularly interesting.

   This is the second paragraph of the quotation.  It
   is just as dull as the first paragraph.
\end{quotation}
Another frequently-displayed structure is a list.
The following is an example of an \emph{itemized}
list.
\begin{itemize}
   \item This is the first item of an itemized list.
         Each item in the list is marked with a ``tick''.
         You don't have to worry about what kind of tick
         mark is used.

   \item This is the second item of the list.  It
         contains another list nested inside it.  The inner
         list is an \emph{enumerated} list.
         \begin{enumerate}
            \item This is the first item of an enumerated
                  list that is nested within the itemized list.

            \item This is the second item of the inner list.
                  \LaTeX\ allows you to nest lists deeper than
                  you really should.
         \end{enumerate}
         This is the rest of the second item of the outer
         list.  It is no more interesting than any other
         part of the item.
   \item This is the third item of the list.
\end{itemize}
You can even display poetry.
\begin{verse}
   There is an environment
    for verse \\             % The \\ command separates lines
   Whose features some poets % within a stanza.
   will curse.

                             % One or more blank lines separate stanzas.

   For instead of making\\
   Them do \emph{all} line breaking, \\
   It allows them to put too many words on a line when they'd rather be
   forced to be terse.
\end{verse}

Mathematical formulas may also be displayed.  A
displayed formula
is
one-line long; multiline
formulas require special formatting instructions.
   \[  \ip{\Gamma}{\psi'} = x'' + y^{2} + z_{i}^{n}\]
Don't start a paragraph with a displayed equation,
nor make one a paragraph by itself.

\end{document}               % End of document.

\section{Bench Huge Section 333}
% synthetic bench section 333
% This is a sample LaTeX input file.  (Version of 12 August 2004.)
%
% A '%' character causes TeX to ignore all remaining text on the line,
% and is used for comments like this one.

\documentclass{article}      % Specifies the document class

                             % The preamble begins here.
\title{An Example Document}  % Declares the document's title.
\author{Leslie Lamport}      % Declares the author's name.
\date{January 21, 1994}      % Deleting this command produces today's date.

\newcommand{\ip}[2]{(#1, #2)}
                             % Defines \ip{arg1}{arg2} to mean
                             % (arg1, arg2).

%\newcommand{\ip}[2]{\langle #1 | #2\rangle}
                             % This is an alternative definition of
                             % \ip that is commented out.

\begin{document}             % End of preamble and beginning of text.

\maketitle                   % Produces the title.

This is an example input file.  Comparing it with
the output it generates can show you how to
produce a simple document of your own.

\section{Ordinary Text}      % Produces section heading.  Lower-level
                             % sections are begun with similar
                             % \subsection and \subsubsection commands.

The ends  of words and sentences are marked
  by   spaces. It  doesn't matter how many
spaces    you type; one is as good as 100.  The
end of   a line counts as a space.

One   or more   blank lines denote the  end
of  a paragraph.

Since any number of consecutive spaces are treated
like a single one, the formatting of the input
file makes no difference to
      \LaTeX,                % The \LaTeX command generates the LaTeX logo.
but it makes a difference to you.  When you use
\LaTeX, making your input file as easy to read
as possible will be a great help as you write
your document and when you change it.  This sample
file shows how you can add comments to your own input
file.

Because printing is different from typewriting,
there are a number of things that you have to do
differently when preparing an input file than if
you were just typing the document directly.
Quotation marks like
       ``this''
have to be handled specially, as do quotes within
quotes:
       ``\,`this'            % \, separates the double and single quote.
        is what I just
        wrote, not  `that'\,''.

Dashes come in three sizes: an
       intra-word
dash, a medium dash for number ranges like
       1--2,
and a punctuation
       dash---like
this.

A sentence-ending space should be larger than the
space between words within a sentence.  You
sometimes have to type special commands in
conjunction with punctuation characters to get
this right, as in the following sentence.
       Gnats, gnus, etc.\ all  % `\ ' makes an inter-word space.
       begin with G\@.         % \@ marks end-of-sentence punctuation.
You should check the spaces after periods when
reading your output to make sure you haven't
forgotten any special cases.  Generating an
ellipsis
       \ldots\               % `\ ' is needed after `\ldots' because TeX
                             % ignores spaces after command names like \ldots
                             % made from \ + letters.
                             %
                             % Note how a `%' character causes TeX to ignore
                             % the end of the input line, so these blank lines
                             % do not start a new paragraph.
                             %
with the right spacing around the periods requires
a special command.

\LaTeX\ interprets some common characters as
commands, so you must type special commands to
generate them.  These characters include the
following:
       \$ \& \% \# \{ and \}.

In printing, text is usually emphasized with an
       \emph{italic}
type style.

\begin{em}
   A long segment of text can also be emphasized
   in this way.  Text within such a segment can be
   given \emph{additional} emphasis.
\end{em}

It is sometimes necessary to prevent \LaTeX\ from
breaking a line where it might otherwise do so.
This may be at a space, as between the ``Mr.''\ and
``Jones'' in
       ``Mr.~Jones'',        % ~ produces an unbreakable interword space.
or within a word---especially when the word is a
symbol like
       \mbox{\emph{itemnum}}
that makes little sense when hyphenated across
lines.

Footnotes\footnote{This is an example of a footnote.}
pose no problem.

\LaTeX\ is good at typesetting mathematical formulas
like
       \( x-3y + z = 7 \)
or
       \( a_{1} > x^{2n} + y^{2n} > x' \)
or
       \( \ip{A}{B} = \sum_{i} a_{i} b_{i} \).
The spaces you type in a formula are
ignored.  Remember that a letter like
       $x$                   % $ ... $  and  \( ... \)  are equivalent
is a formula when it denotes a mathematical
symbol, and it should be typed as one.

\section{Displayed Text}

Text is displayed by indenting it from the left
margin.  Quotations are commonly displayed.  There
are short quotations
\begin{quote}
   This is a short quotation.  It consists of a
   single paragraph of text.  See how it is formatted.
\end{quote}
and longer ones.
\begin{quotation}
   This is a longer quotation.  It consists of two
   paragraphs of text, neither of which are
   particularly interesting.

   This is the second paragraph of the quotation.  It
   is just as dull as the first paragraph.
\end{quotation}
Another frequently-displayed structure is a list.
The following is an example of an \emph{itemized}
list.
\begin{itemize}
   \item This is the first item of an itemized list.
         Each item in the list is marked with a ``tick''.
         You don't have to worry about what kind of tick
         mark is used.

   \item This is the second item of the list.  It
         contains another list nested inside it.  The inner
         list is an \emph{enumerated} list.
         \begin{enumerate}
            \item This is the first item of an enumerated
                  list that is nested within the itemized list.

            \item This is the second item of the inner list.
                  \LaTeX\ allows you to nest lists deeper than
                  you really should.
         \end{enumerate}
         This is the rest of the second item of the outer
         list.  It is no more interesting than any other
         part of the item.
   \item This is the third item of the list.
\end{itemize}
You can even display poetry.
\begin{verse}
   There is an environment
    for verse \\             % The \\ command separates lines
   Whose features some poets % within a stanza.
   will curse.

                             % One or more blank lines separate stanzas.

   For instead of making\\
   Them do \emph{all} line breaking, \\
   It allows them to put too many words on a line when they'd rather be
   forced to be terse.
\end{verse}

Mathematical formulas may also be displayed.  A
displayed formula
is
one-line long; multiline
formulas require special formatting instructions.
   \[  \ip{\Gamma}{\psi'} = x'' + y^{2} + z_{i}^{n}\]
Don't start a paragraph with a displayed equation,
nor make one a paragraph by itself.

\end{document}               % End of document.

\section{Bench Huge Section 334}
% synthetic bench section 334
% This is a sample LaTeX input file.  (Version of 12 August 2004.)
%
% A '%' character causes TeX to ignore all remaining text on the line,
% and is used for comments like this one.

\documentclass{article}      % Specifies the document class

                             % The preamble begins here.
\title{An Example Document}  % Declares the document's title.
\author{Leslie Lamport}      % Declares the author's name.
\date{January 21, 1994}      % Deleting this command produces today's date.

\newcommand{\ip}[2]{(#1, #2)}
                             % Defines \ip{arg1}{arg2} to mean
                             % (arg1, arg2).

%\newcommand{\ip}[2]{\langle #1 | #2\rangle}
                             % This is an alternative definition of
                             % \ip that is commented out.

\begin{document}             % End of preamble and beginning of text.

\maketitle                   % Produces the title.

This is an example input file.  Comparing it with
the output it generates can show you how to
produce a simple document of your own.

\section{Ordinary Text}      % Produces section heading.  Lower-level
                             % sections are begun with similar
                             % \subsection and \subsubsection commands.

The ends  of words and sentences are marked
  by   spaces. It  doesn't matter how many
spaces    you type; one is as good as 100.  The
end of   a line counts as a space.

One   or more   blank lines denote the  end
of  a paragraph.

Since any number of consecutive spaces are treated
like a single one, the formatting of the input
file makes no difference to
      \LaTeX,                % The \LaTeX command generates the LaTeX logo.
but it makes a difference to you.  When you use
\LaTeX, making your input file as easy to read
as possible will be a great help as you write
your document and when you change it.  This sample
file shows how you can add comments to your own input
file.

Because printing is different from typewriting,
there are a number of things that you have to do
differently when preparing an input file than if
you were just typing the document directly.
Quotation marks like
       ``this''
have to be handled specially, as do quotes within
quotes:
       ``\,`this'            % \, separates the double and single quote.
        is what I just
        wrote, not  `that'\,''.

Dashes come in three sizes: an
       intra-word
dash, a medium dash for number ranges like
       1--2,
and a punctuation
       dash---like
this.

A sentence-ending space should be larger than the
space between words within a sentence.  You
sometimes have to type special commands in
conjunction with punctuation characters to get
this right, as in the following sentence.
       Gnats, gnus, etc.\ all  % `\ ' makes an inter-word space.
       begin with G\@.         % \@ marks end-of-sentence punctuation.
You should check the spaces after periods when
reading your output to make sure you haven't
forgotten any special cases.  Generating an
ellipsis
       \ldots\               % `\ ' is needed after `\ldots' because TeX
                             % ignores spaces after command names like \ldots
                             % made from \ + letters.
                             %
                             % Note how a `%' character causes TeX to ignore
                             % the end of the input line, so these blank lines
                             % do not start a new paragraph.
                             %
with the right spacing around the periods requires
a special command.

\LaTeX\ interprets some common characters as
commands, so you must type special commands to
generate them.  These characters include the
following:
       \$ \& \% \# \{ and \}.

In printing, text is usually emphasized with an
       \emph{italic}
type style.

\begin{em}
   A long segment of text can also be emphasized
   in this way.  Text within such a segment can be
   given \emph{additional} emphasis.
\end{em}

It is sometimes necessary to prevent \LaTeX\ from
breaking a line where it might otherwise do so.
This may be at a space, as between the ``Mr.''\ and
``Jones'' in
       ``Mr.~Jones'',        % ~ produces an unbreakable interword space.
or within a word---especially when the word is a
symbol like
       \mbox{\emph{itemnum}}
that makes little sense when hyphenated across
lines.

Footnotes\footnote{This is an example of a footnote.}
pose no problem.

\LaTeX\ is good at typesetting mathematical formulas
like
       \( x-3y + z = 7 \)
or
       \( a_{1} > x^{2n} + y^{2n} > x' \)
or
       \( \ip{A}{B} = \sum_{i} a_{i} b_{i} \).
The spaces you type in a formula are
ignored.  Remember that a letter like
       $x$                   % $ ... $  and  \( ... \)  are equivalent
is a formula when it denotes a mathematical
symbol, and it should be typed as one.

\section{Displayed Text}

Text is displayed by indenting it from the left
margin.  Quotations are commonly displayed.  There
are short quotations
\begin{quote}
   This is a short quotation.  It consists of a
   single paragraph of text.  See how it is formatted.
\end{quote}
and longer ones.
\begin{quotation}
   This is a longer quotation.  It consists of two
   paragraphs of text, neither of which are
   particularly interesting.

   This is the second paragraph of the quotation.  It
   is just as dull as the first paragraph.
\end{quotation}
Another frequently-displayed structure is a list.
The following is an example of an \emph{itemized}
list.
\begin{itemize}
   \item This is the first item of an itemized list.
         Each item in the list is marked with a ``tick''.
         You don't have to worry about what kind of tick
         mark is used.

   \item This is the second item of the list.  It
         contains another list nested inside it.  The inner
         list is an \emph{enumerated} list.
         \begin{enumerate}
            \item This is the first item of an enumerated
                  list that is nested within the itemized list.

            \item This is the second item of the inner list.
                  \LaTeX\ allows you to nest lists deeper than
                  you really should.
         \end{enumerate}
         This is the rest of the second item of the outer
         list.  It is no more interesting than any other
         part of the item.
   \item This is the third item of the list.
\end{itemize}
You can even display poetry.
\begin{verse}
   There is an environment
    for verse \\             % The \\ command separates lines
   Whose features some poets % within a stanza.
   will curse.

                             % One or more blank lines separate stanzas.

   For instead of making\\
   Them do \emph{all} line breaking, \\
   It allows them to put too many words on a line when they'd rather be
   forced to be terse.
\end{verse}

Mathematical formulas may also be displayed.  A
displayed formula
is
one-line long; multiline
formulas require special formatting instructions.
   \[  \ip{\Gamma}{\psi'} = x'' + y^{2} + z_{i}^{n}\]
Don't start a paragraph with a displayed equation,
nor make one a paragraph by itself.

\end{document}               % End of document.

\section{Bench Huge Section 335}
% synthetic bench section 335
% This is a sample LaTeX input file.  (Version of 12 August 2004.)
%
% A '%' character causes TeX to ignore all remaining text on the line,
% and is used for comments like this one.

\documentclass{article}      % Specifies the document class

                             % The preamble begins here.
\title{An Example Document}  % Declares the document's title.
\author{Leslie Lamport}      % Declares the author's name.
\date{January 21, 1994}      % Deleting this command produces today's date.

\newcommand{\ip}[2]{(#1, #2)}
                             % Defines \ip{arg1}{arg2} to mean
                             % (arg1, arg2).

%\newcommand{\ip}[2]{\langle #1 | #2\rangle}
                             % This is an alternative definition of
                             % \ip that is commented out.

\begin{document}             % End of preamble and beginning of text.

\maketitle                   % Produces the title.

This is an example input file.  Comparing it with
the output it generates can show you how to
produce a simple document of your own.

\section{Ordinary Text}      % Produces section heading.  Lower-level
                             % sections are begun with similar
                             % \subsection and \subsubsection commands.

The ends  of words and sentences are marked
  by   spaces. It  doesn't matter how many
spaces    you type; one is as good as 100.  The
end of   a line counts as a space.

One   or more   blank lines denote the  end
of  a paragraph.

Since any number of consecutive spaces are treated
like a single one, the formatting of the input
file makes no difference to
      \LaTeX,                % The \LaTeX command generates the LaTeX logo.
but it makes a difference to you.  When you use
\LaTeX, making your input file as easy to read
as possible will be a great help as you write
your document and when you change it.  This sample
file shows how you can add comments to your own input
file.

Because printing is different from typewriting,
there are a number of things that you have to do
differently when preparing an input file than if
you were just typing the document directly.
Quotation marks like
       ``this''
have to be handled specially, as do quotes within
quotes:
       ``\,`this'            % \, separates the double and single quote.
        is what I just
        wrote, not  `that'\,''.

Dashes come in three sizes: an
       intra-word
dash, a medium dash for number ranges like
       1--2,
and a punctuation
       dash---like
this.

A sentence-ending space should be larger than the
space between words within a sentence.  You
sometimes have to type special commands in
conjunction with punctuation characters to get
this right, as in the following sentence.
       Gnats, gnus, etc.\ all  % `\ ' makes an inter-word space.
       begin with G\@.         % \@ marks end-of-sentence punctuation.
You should check the spaces after periods when
reading your output to make sure you haven't
forgotten any special cases.  Generating an
ellipsis
       \ldots\               % `\ ' is needed after `\ldots' because TeX
                             % ignores spaces after command names like \ldots
                             % made from \ + letters.
                             %
                             % Note how a `%' character causes TeX to ignore
                             % the end of the input line, so these blank lines
                             % do not start a new paragraph.
                             %
with the right spacing around the periods requires
a special command.

\LaTeX\ interprets some common characters as
commands, so you must type special commands to
generate them.  These characters include the
following:
       \$ \& \% \# \{ and \}.

In printing, text is usually emphasized with an
       \emph{italic}
type style.

\begin{em}
   A long segment of text can also be emphasized
   in this way.  Text within such a segment can be
   given \emph{additional} emphasis.
\end{em}

It is sometimes necessary to prevent \LaTeX\ from
breaking a line where it might otherwise do so.
This may be at a space, as between the ``Mr.''\ and
``Jones'' in
       ``Mr.~Jones'',        % ~ produces an unbreakable interword space.
or within a word---especially when the word is a
symbol like
       \mbox{\emph{itemnum}}
that makes little sense when hyphenated across
lines.

Footnotes\footnote{This is an example of a footnote.}
pose no problem.

\LaTeX\ is good at typesetting mathematical formulas
like
       \( x-3y + z = 7 \)
or
       \( a_{1} > x^{2n} + y^{2n} > x' \)
or
       \( \ip{A}{B} = \sum_{i} a_{i} b_{i} \).
The spaces you type in a formula are
ignored.  Remember that a letter like
       $x$                   % $ ... $  and  \( ... \)  are equivalent
is a formula when it denotes a mathematical
symbol, and it should be typed as one.

\section{Displayed Text}

Text is displayed by indenting it from the left
margin.  Quotations are commonly displayed.  There
are short quotations
\begin{quote}
   This is a short quotation.  It consists of a
   single paragraph of text.  See how it is formatted.
\end{quote}
and longer ones.
\begin{quotation}
   This is a longer quotation.  It consists of two
   paragraphs of text, neither of which are
   particularly interesting.

   This is the second paragraph of the quotation.  It
   is just as dull as the first paragraph.
\end{quotation}
Another frequently-displayed structure is a list.
The following is an example of an \emph{itemized}
list.
\begin{itemize}
   \item This is the first item of an itemized list.
         Each item in the list is marked with a ``tick''.
         You don't have to worry about what kind of tick
         mark is used.

   \item This is the second item of the list.  It
         contains another list nested inside it.  The inner
         list is an \emph{enumerated} list.
         \begin{enumerate}
            \item This is the first item of an enumerated
                  list that is nested within the itemized list.

            \item This is the second item of the inner list.
                  \LaTeX\ allows you to nest lists deeper than
                  you really should.
         \end{enumerate}
         This is the rest of the second item of the outer
         list.  It is no more interesting than any other
         part of the item.
   \item This is the third item of the list.
\end{itemize}
You can even display poetry.
\begin{verse}
   There is an environment
    for verse \\             % The \\ command separates lines
   Whose features some poets % within a stanza.
   will curse.

                             % One or more blank lines separate stanzas.

   For instead of making\\
   Them do \emph{all} line breaking, \\
   It allows them to put too many words on a line when they'd rather be
   forced to be terse.
\end{verse}

Mathematical formulas may also be displayed.  A
displayed formula
is
one-line long; multiline
formulas require special formatting instructions.
   \[  \ip{\Gamma}{\psi'} = x'' + y^{2} + z_{i}^{n}\]
Don't start a paragraph with a displayed equation,
nor make one a paragraph by itself.

\end{document}               % End of document.

\section{Bench Huge Section 336}
% synthetic bench section 336
% This is a sample LaTeX input file.  (Version of 12 August 2004.)
%
% A '%' character causes TeX to ignore all remaining text on the line,
% and is used for comments like this one.

\documentclass{article}      % Specifies the document class

                             % The preamble begins here.
\title{An Example Document}  % Declares the document's title.
\author{Leslie Lamport}      % Declares the author's name.
\date{January 21, 1994}      % Deleting this command produces today's date.

\newcommand{\ip}[2]{(#1, #2)}
                             % Defines \ip{arg1}{arg2} to mean
                             % (arg1, arg2).

%\newcommand{\ip}[2]{\langle #1 | #2\rangle}
                             % This is an alternative definition of
                             % \ip that is commented out.

\begin{document}             % End of preamble and beginning of text.

\maketitle                   % Produces the title.

This is an example input file.  Comparing it with
the output it generates can show you how to
produce a simple document of your own.

\section{Ordinary Text}      % Produces section heading.  Lower-level
                             % sections are begun with similar
                             % \subsection and \subsubsection commands.

The ends  of words and sentences are marked
  by   spaces. It  doesn't matter how many
spaces    you type; one is as good as 100.  The
end of   a line counts as a space.

One   or more   blank lines denote the  end
of  a paragraph.

Since any number of consecutive spaces are treated
like a single one, the formatting of the input
file makes no difference to
      \LaTeX,                % The \LaTeX command generates the LaTeX logo.
but it makes a difference to you.  When you use
\LaTeX, making your input file as easy to read
as possible will be a great help as you write
your document and when you change it.  This sample
file shows how you can add comments to your own input
file.

Because printing is different from typewriting,
there are a number of things that you have to do
differently when preparing an input file than if
you were just typing the document directly.
Quotation marks like
       ``this''
have to be handled specially, as do quotes within
quotes:
       ``\,`this'            % \, separates the double and single quote.
        is what I just
        wrote, not  `that'\,''.

Dashes come in three sizes: an
       intra-word
dash, a medium dash for number ranges like
       1--2,
and a punctuation
       dash---like
this.

A sentence-ending space should be larger than the
space between words within a sentence.  You
sometimes have to type special commands in
conjunction with punctuation characters to get
this right, as in the following sentence.
       Gnats, gnus, etc.\ all  % `\ ' makes an inter-word space.
       begin with G\@.         % \@ marks end-of-sentence punctuation.
You should check the spaces after periods when
reading your output to make sure you haven't
forgotten any special cases.  Generating an
ellipsis
       \ldots\               % `\ ' is needed after `\ldots' because TeX
                             % ignores spaces after command names like \ldots
                             % made from \ + letters.
                             %
                             % Note how a `%' character causes TeX to ignore
                             % the end of the input line, so these blank lines
                             % do not start a new paragraph.
                             %
with the right spacing around the periods requires
a special command.

\LaTeX\ interprets some common characters as
commands, so you must type special commands to
generate them.  These characters include the
following:
       \$ \& \% \# \{ and \}.

In printing, text is usually emphasized with an
       \emph{italic}
type style.

\begin{em}
   A long segment of text can also be emphasized
   in this way.  Text within such a segment can be
   given \emph{additional} emphasis.
\end{em}

It is sometimes necessary to prevent \LaTeX\ from
breaking a line where it might otherwise do so.
This may be at a space, as between the ``Mr.''\ and
``Jones'' in
       ``Mr.~Jones'',        % ~ produces an unbreakable interword space.
or within a word---especially when the word is a
symbol like
       \mbox{\emph{itemnum}}
that makes little sense when hyphenated across
lines.

Footnotes\footnote{This is an example of a footnote.}
pose no problem.

\LaTeX\ is good at typesetting mathematical formulas
like
       \( x-3y + z = 7 \)
or
       \( a_{1} > x^{2n} + y^{2n} > x' \)
or
       \( \ip{A}{B} = \sum_{i} a_{i} b_{i} \).
The spaces you type in a formula are
ignored.  Remember that a letter like
       $x$                   % $ ... $  and  \( ... \)  are equivalent
is a formula when it denotes a mathematical
symbol, and it should be typed as one.

\section{Displayed Text}

Text is displayed by indenting it from the left
margin.  Quotations are commonly displayed.  There
are short quotations
\begin{quote}
   This is a short quotation.  It consists of a
   single paragraph of text.  See how it is formatted.
\end{quote}
and longer ones.
\begin{quotation}
   This is a longer quotation.  It consists of two
   paragraphs of text, neither of which are
   particularly interesting.

   This is the second paragraph of the quotation.  It
   is just as dull as the first paragraph.
\end{quotation}
Another frequently-displayed structure is a list.
The following is an example of an \emph{itemized}
list.
\begin{itemize}
   \item This is the first item of an itemized list.
         Each item in the list is marked with a ``tick''.
         You don't have to worry about what kind of tick
         mark is used.

   \item This is the second item of the list.  It
         contains another list nested inside it.  The inner
         list is an \emph{enumerated} list.
         \begin{enumerate}
            \item This is the first item of an enumerated
                  list that is nested within the itemized list.

            \item This is the second item of the inner list.
                  \LaTeX\ allows you to nest lists deeper than
                  you really should.
         \end{enumerate}
         This is the rest of the second item of the outer
         list.  It is no more interesting than any other
         part of the item.
   \item This is the third item of the list.
\end{itemize}
You can even display poetry.
\begin{verse}
   There is an environment
    for verse \\             % The \\ command separates lines
   Whose features some poets % within a stanza.
   will curse.

                             % One or more blank lines separate stanzas.

   For instead of making\\
   Them do \emph{all} line breaking, \\
   It allows them to put too many words on a line when they'd rather be
   forced to be terse.
\end{verse}

Mathematical formulas may also be displayed.  A
displayed formula
is
one-line long; multiline
formulas require special formatting instructions.
   \[  \ip{\Gamma}{\psi'} = x'' + y^{2} + z_{i}^{n}\]
Don't start a paragraph with a displayed equation,
nor make one a paragraph by itself.

\end{document}               % End of document.

\section{Bench Huge Section 337}
% synthetic bench section 337
% This is a sample LaTeX input file.  (Version of 12 August 2004.)
%
% A '%' character causes TeX to ignore all remaining text on the line,
% and is used for comments like this one.

\documentclass{article}      % Specifies the document class

                             % The preamble begins here.
\title{An Example Document}  % Declares the document's title.
\author{Leslie Lamport}      % Declares the author's name.
\date{January 21, 1994}      % Deleting this command produces today's date.

\newcommand{\ip}[2]{(#1, #2)}
                             % Defines \ip{arg1}{arg2} to mean
                             % (arg1, arg2).

%\newcommand{\ip}[2]{\langle #1 | #2\rangle}
                             % This is an alternative definition of
                             % \ip that is commented out.

\begin{document}             % End of preamble and beginning of text.

\maketitle                   % Produces the title.

This is an example input file.  Comparing it with
the output it generates can show you how to
produce a simple document of your own.

\section{Ordinary Text}      % Produces section heading.  Lower-level
                             % sections are begun with similar
                             % \subsection and \subsubsection commands.

The ends  of words and sentences are marked
  by   spaces. It  doesn't matter how many
spaces    you type; one is as good as 100.  The
end of   a line counts as a space.

One   or more   blank lines denote the  end
of  a paragraph.

Since any number of consecutive spaces are treated
like a single one, the formatting of the input
file makes no difference to
      \LaTeX,                % The \LaTeX command generates the LaTeX logo.
but it makes a difference to you.  When you use
\LaTeX, making your input file as easy to read
as possible will be a great help as you write
your document and when you change it.  This sample
file shows how you can add comments to your own input
file.

Because printing is different from typewriting,
there are a number of things that you have to do
differently when preparing an input file than if
you were just typing the document directly.
Quotation marks like
       ``this''
have to be handled specially, as do quotes within
quotes:
       ``\,`this'            % \, separates the double and single quote.
        is what I just
        wrote, not  `that'\,''.

Dashes come in three sizes: an
       intra-word
dash, a medium dash for number ranges like
       1--2,
and a punctuation
       dash---like
this.

A sentence-ending space should be larger than the
space between words within a sentence.  You
sometimes have to type special commands in
conjunction with punctuation characters to get
this right, as in the following sentence.
       Gnats, gnus, etc.\ all  % `\ ' makes an inter-word space.
       begin with G\@.         % \@ marks end-of-sentence punctuation.
You should check the spaces after periods when
reading your output to make sure you haven't
forgotten any special cases.  Generating an
ellipsis
       \ldots\               % `\ ' is needed after `\ldots' because TeX
                             % ignores spaces after command names like \ldots
                             % made from \ + letters.
                             %
                             % Note how a `%' character causes TeX to ignore
                             % the end of the input line, so these blank lines
                             % do not start a new paragraph.
                             %
with the right spacing around the periods requires
a special command.

\LaTeX\ interprets some common characters as
commands, so you must type special commands to
generate them.  These characters include the
following:
       \$ \& \% \# \{ and \}.

In printing, text is usually emphasized with an
       \emph{italic}
type style.

\begin{em}
   A long segment of text can also be emphasized
   in this way.  Text within such a segment can be
   given \emph{additional} emphasis.
\end{em}

It is sometimes necessary to prevent \LaTeX\ from
breaking a line where it might otherwise do so.
This may be at a space, as between the ``Mr.''\ and
``Jones'' in
       ``Mr.~Jones'',        % ~ produces an unbreakable interword space.
or within a word---especially when the word is a
symbol like
       \mbox{\emph{itemnum}}
that makes little sense when hyphenated across
lines.

Footnotes\footnote{This is an example of a footnote.}
pose no problem.

\LaTeX\ is good at typesetting mathematical formulas
like
       \( x-3y + z = 7 \)
or
       \( a_{1} > x^{2n} + y^{2n} > x' \)
or
       \( \ip{A}{B} = \sum_{i} a_{i} b_{i} \).
The spaces you type in a formula are
ignored.  Remember that a letter like
       $x$                   % $ ... $  and  \( ... \)  are equivalent
is a formula when it denotes a mathematical
symbol, and it should be typed as one.

\section{Displayed Text}

Text is displayed by indenting it from the left
margin.  Quotations are commonly displayed.  There
are short quotations
\begin{quote}
   This is a short quotation.  It consists of a
   single paragraph of text.  See how it is formatted.
\end{quote}
and longer ones.
\begin{quotation}
   This is a longer quotation.  It consists of two
   paragraphs of text, neither of which are
   particularly interesting.

   This is the second paragraph of the quotation.  It
   is just as dull as the first paragraph.
\end{quotation}
Another frequently-displayed structure is a list.
The following is an example of an \emph{itemized}
list.
\begin{itemize}
   \item This is the first item of an itemized list.
         Each item in the list is marked with a ``tick''.
         You don't have to worry about what kind of tick
         mark is used.

   \item This is the second item of the list.  It
         contains another list nested inside it.  The inner
         list is an \emph{enumerated} list.
         \begin{enumerate}
            \item This is the first item of an enumerated
                  list that is nested within the itemized list.

            \item This is the second item of the inner list.
                  \LaTeX\ allows you to nest lists deeper than
                  you really should.
         \end{enumerate}
         This is the rest of the second item of the outer
         list.  It is no more interesting than any other
         part of the item.
   \item This is the third item of the list.
\end{itemize}
You can even display poetry.
\begin{verse}
   There is an environment
    for verse \\             % The \\ command separates lines
   Whose features some poets % within a stanza.
   will curse.

                             % One or more blank lines separate stanzas.

   For instead of making\\
   Them do \emph{all} line breaking, \\
   It allows them to put too many words on a line when they'd rather be
   forced to be terse.
\end{verse}

Mathematical formulas may also be displayed.  A
displayed formula
is
one-line long; multiline
formulas require special formatting instructions.
   \[  \ip{\Gamma}{\psi'} = x'' + y^{2} + z_{i}^{n}\]
Don't start a paragraph with a displayed equation,
nor make one a paragraph by itself.

\end{document}               % End of document.

\section{Bench Huge Section 338}
% synthetic bench section 338
% This is a sample LaTeX input file.  (Version of 12 August 2004.)
%
% A '%' character causes TeX to ignore all remaining text on the line,
% and is used for comments like this one.

\documentclass{article}      % Specifies the document class

                             % The preamble begins here.
\title{An Example Document}  % Declares the document's title.
\author{Leslie Lamport}      % Declares the author's name.
\date{January 21, 1994}      % Deleting this command produces today's date.

\newcommand{\ip}[2]{(#1, #2)}
                             % Defines \ip{arg1}{arg2} to mean
                             % (arg1, arg2).

%\newcommand{\ip}[2]{\langle #1 | #2\rangle}
                             % This is an alternative definition of
                             % \ip that is commented out.

\begin{document}             % End of preamble and beginning of text.

\maketitle                   % Produces the title.

This is an example input file.  Comparing it with
the output it generates can show you how to
produce a simple document of your own.

\section{Ordinary Text}      % Produces section heading.  Lower-level
                             % sections are begun with similar
                             % \subsection and \subsubsection commands.

The ends  of words and sentences are marked
  by   spaces. It  doesn't matter how many
spaces    you type; one is as good as 100.  The
end of   a line counts as a space.

One   or more   blank lines denote the  end
of  a paragraph.

Since any number of consecutive spaces are treated
like a single one, the formatting of the input
file makes no difference to
      \LaTeX,                % The \LaTeX command generates the LaTeX logo.
but it makes a difference to you.  When you use
\LaTeX, making your input file as easy to read
as possible will be a great help as you write
your document and when you change it.  This sample
file shows how you can add comments to your own input
file.

Because printing is different from typewriting,
there are a number of things that you have to do
differently when preparing an input file than if
you were just typing the document directly.
Quotation marks like
       ``this''
have to be handled specially, as do quotes within
quotes:
       ``\,`this'            % \, separates the double and single quote.
        is what I just
        wrote, not  `that'\,''.

Dashes come in three sizes: an
       intra-word
dash, a medium dash for number ranges like
       1--2,
and a punctuation
       dash---like
this.

A sentence-ending space should be larger than the
space between words within a sentence.  You
sometimes have to type special commands in
conjunction with punctuation characters to get
this right, as in the following sentence.
       Gnats, gnus, etc.\ all  % `\ ' makes an inter-word space.
       begin with G\@.         % \@ marks end-of-sentence punctuation.
You should check the spaces after periods when
reading your output to make sure you haven't
forgotten any special cases.  Generating an
ellipsis
       \ldots\               % `\ ' is needed after `\ldots' because TeX
                             % ignores spaces after command names like \ldots
                             % made from \ + letters.
                             %
                             % Note how a `%' character causes TeX to ignore
                             % the end of the input line, so these blank lines
                             % do not start a new paragraph.
                             %
with the right spacing around the periods requires
a special command.

\LaTeX\ interprets some common characters as
commands, so you must type special commands to
generate them.  These characters include the
following:
       \$ \& \% \# \{ and \}.

In printing, text is usually emphasized with an
       \emph{italic}
type style.

\begin{em}
   A long segment of text can also be emphasized
   in this way.  Text within such a segment can be
   given \emph{additional} emphasis.
\end{em}

It is sometimes necessary to prevent \LaTeX\ from
breaking a line where it might otherwise do so.
This may be at a space, as between the ``Mr.''\ and
``Jones'' in
       ``Mr.~Jones'',        % ~ produces an unbreakable interword space.
or within a word---especially when the word is a
symbol like
       \mbox{\emph{itemnum}}
that makes little sense when hyphenated across
lines.

Footnotes\footnote{This is an example of a footnote.}
pose no problem.

\LaTeX\ is good at typesetting mathematical formulas
like
       \( x-3y + z = 7 \)
or
       \( a_{1} > x^{2n} + y^{2n} > x' \)
or
       \( \ip{A}{B} = \sum_{i} a_{i} b_{i} \).
The spaces you type in a formula are
ignored.  Remember that a letter like
       $x$                   % $ ... $  and  \( ... \)  are equivalent
is a formula when it denotes a mathematical
symbol, and it should be typed as one.

\section{Displayed Text}

Text is displayed by indenting it from the left
margin.  Quotations are commonly displayed.  There
are short quotations
\begin{quote}
   This is a short quotation.  It consists of a
   single paragraph of text.  See how it is formatted.
\end{quote}
and longer ones.
\begin{quotation}
   This is a longer quotation.  It consists of two
   paragraphs of text, neither of which are
   particularly interesting.

   This is the second paragraph of the quotation.  It
   is just as dull as the first paragraph.
\end{quotation}
Another frequently-displayed structure is a list.
The following is an example of an \emph{itemized}
list.
\begin{itemize}
   \item This is the first item of an itemized list.
         Each item in the list is marked with a ``tick''.
         You don't have to worry about what kind of tick
         mark is used.

   \item This is the second item of the list.  It
         contains another list nested inside it.  The inner
         list is an \emph{enumerated} list.
         \begin{enumerate}
            \item This is the first item of an enumerated
                  list that is nested within the itemized list.

            \item This is the second item of the inner list.
                  \LaTeX\ allows you to nest lists deeper than
                  you really should.
         \end{enumerate}
         This is the rest of the second item of the outer
         list.  It is no more interesting than any other
         part of the item.
   \item This is the third item of the list.
\end{itemize}
You can even display poetry.
\begin{verse}
   There is an environment
    for verse \\             % The \\ command separates lines
   Whose features some poets % within a stanza.
   will curse.

                             % One or more blank lines separate stanzas.

   For instead of making\\
   Them do \emph{all} line breaking, \\
   It allows them to put too many words on a line when they'd rather be
   forced to be terse.
\end{verse}

Mathematical formulas may also be displayed.  A
displayed formula
is
one-line long; multiline
formulas require special formatting instructions.
   \[  \ip{\Gamma}{\psi'} = x'' + y^{2} + z_{i}^{n}\]
Don't start a paragraph with a displayed equation,
nor make one a paragraph by itself.

\end{document}               % End of document.

\section{Bench Huge Section 339}
% synthetic bench section 339
% This is a sample LaTeX input file.  (Version of 12 August 2004.)
%
% A '%' character causes TeX to ignore all remaining text on the line,
% and is used for comments like this one.

\documentclass{article}      % Specifies the document class

                             % The preamble begins here.
\title{An Example Document}  % Declares the document's title.
\author{Leslie Lamport}      % Declares the author's name.
\date{January 21, 1994}      % Deleting this command produces today's date.

\newcommand{\ip}[2]{(#1, #2)}
                             % Defines \ip{arg1}{arg2} to mean
                             % (arg1, arg2).

%\newcommand{\ip}[2]{\langle #1 | #2\rangle}
                             % This is an alternative definition of
                             % \ip that is commented out.

\begin{document}             % End of preamble and beginning of text.

\maketitle                   % Produces the title.

This is an example input file.  Comparing it with
the output it generates can show you how to
produce a simple document of your own.

\section{Ordinary Text}      % Produces section heading.  Lower-level
                             % sections are begun with similar
                             % \subsection and \subsubsection commands.

The ends  of words and sentences are marked
  by   spaces. It  doesn't matter how many
spaces    you type; one is as good as 100.  The
end of   a line counts as a space.

One   or more   blank lines denote the  end
of  a paragraph.

Since any number of consecutive spaces are treated
like a single one, the formatting of the input
file makes no difference to
      \LaTeX,                % The \LaTeX command generates the LaTeX logo.
but it makes a difference to you.  When you use
\LaTeX, making your input file as easy to read
as possible will be a great help as you write
your document and when you change it.  This sample
file shows how you can add comments to your own input
file.

Because printing is different from typewriting,
there are a number of things that you have to do
differently when preparing an input file than if
you were just typing the document directly.
Quotation marks like
       ``this''
have to be handled specially, as do quotes within
quotes:
       ``\,`this'            % \, separates the double and single quote.
        is what I just
        wrote, not  `that'\,''.

Dashes come in three sizes: an
       intra-word
dash, a medium dash for number ranges like
       1--2,
and a punctuation
       dash---like
this.

A sentence-ending space should be larger than the
space between words within a sentence.  You
sometimes have to type special commands in
conjunction with punctuation characters to get
this right, as in the following sentence.
       Gnats, gnus, etc.\ all  % `\ ' makes an inter-word space.
       begin with G\@.         % \@ marks end-of-sentence punctuation.
You should check the spaces after periods when
reading your output to make sure you haven't
forgotten any special cases.  Generating an
ellipsis
       \ldots\               % `\ ' is needed after `\ldots' because TeX
                             % ignores spaces after command names like \ldots
                             % made from \ + letters.
                             %
                             % Note how a `%' character causes TeX to ignore
                             % the end of the input line, so these blank lines
                             % do not start a new paragraph.
                             %
with the right spacing around the periods requires
a special command.

\LaTeX\ interprets some common characters as
commands, so you must type special commands to
generate them.  These characters include the
following:
       \$ \& \% \# \{ and \}.

In printing, text is usually emphasized with an
       \emph{italic}
type style.

\begin{em}
   A long segment of text can also be emphasized
   in this way.  Text within such a segment can be
   given \emph{additional} emphasis.
\end{em}

It is sometimes necessary to prevent \LaTeX\ from
breaking a line where it might otherwise do so.
This may be at a space, as between the ``Mr.''\ and
``Jones'' in
       ``Mr.~Jones'',        % ~ produces an unbreakable interword space.
or within a word---especially when the word is a
symbol like
       \mbox{\emph{itemnum}}
that makes little sense when hyphenated across
lines.

Footnotes\footnote{This is an example of a footnote.}
pose no problem.

\LaTeX\ is good at typesetting mathematical formulas
like
       \( x-3y + z = 7 \)
or
       \( a_{1} > x^{2n} + y^{2n} > x' \)
or
       \( \ip{A}{B} = \sum_{i} a_{i} b_{i} \).
The spaces you type in a formula are
ignored.  Remember that a letter like
       $x$                   % $ ... $  and  \( ... \)  are equivalent
is a formula when it denotes a mathematical
symbol, and it should be typed as one.

\section{Displayed Text}

Text is displayed by indenting it from the left
margin.  Quotations are commonly displayed.  There
are short quotations
\begin{quote}
   This is a short quotation.  It consists of a
   single paragraph of text.  See how it is formatted.
\end{quote}
and longer ones.
\begin{quotation}
   This is a longer quotation.  It consists of two
   paragraphs of text, neither of which are
   particularly interesting.

   This is the second paragraph of the quotation.  It
   is just as dull as the first paragraph.
\end{quotation}
Another frequently-displayed structure is a list.
The following is an example of an \emph{itemized}
list.
\begin{itemize}
   \item This is the first item of an itemized list.
         Each item in the list is marked with a ``tick''.
         You don't have to worry about what kind of tick
         mark is used.

   \item This is the second item of the list.  It
         contains another list nested inside it.  The inner
         list is an \emph{enumerated} list.
         \begin{enumerate}
            \item This is the first item of an enumerated
                  list that is nested within the itemized list.

            \item This is the second item of the inner list.
                  \LaTeX\ allows you to nest lists deeper than
                  you really should.
         \end{enumerate}
         This is the rest of the second item of the outer
         list.  It is no more interesting than any other
         part of the item.
   \item This is the third item of the list.
\end{itemize}
You can even display poetry.
\begin{verse}
   There is an environment
    for verse \\             % The \\ command separates lines
   Whose features some poets % within a stanza.
   will curse.

                             % One or more blank lines separate stanzas.

   For instead of making\\
   Them do \emph{all} line breaking, \\
   It allows them to put too many words on a line when they'd rather be
   forced to be terse.
\end{verse}

Mathematical formulas may also be displayed.  A
displayed formula
is
one-line long; multiline
formulas require special formatting instructions.
   \[  \ip{\Gamma}{\psi'} = x'' + y^{2} + z_{i}^{n}\]
Don't start a paragraph with a displayed equation,
nor make one a paragraph by itself.

\end{document}               % End of document.

\section{Bench Huge Section 340}
% synthetic bench section 340
% This is a sample LaTeX input file.  (Version of 12 August 2004.)
%
% A '%' character causes TeX to ignore all remaining text on the line,
% and is used for comments like this one.

\documentclass{article}      % Specifies the document class

                             % The preamble begins here.
\title{An Example Document}  % Declares the document's title.
\author{Leslie Lamport}      % Declares the author's name.
\date{January 21, 1994}      % Deleting this command produces today's date.

\newcommand{\ip}[2]{(#1, #2)}
                             % Defines \ip{arg1}{arg2} to mean
                             % (arg1, arg2).

%\newcommand{\ip}[2]{\langle #1 | #2\rangle}
                             % This is an alternative definition of
                             % \ip that is commented out.

\begin{document}             % End of preamble and beginning of text.

\maketitle                   % Produces the title.

This is an example input file.  Comparing it with
the output it generates can show you how to
produce a simple document of your own.

\section{Ordinary Text}      % Produces section heading.  Lower-level
                             % sections are begun with similar
                             % \subsection and \subsubsection commands.

The ends  of words and sentences are marked
  by   spaces. It  doesn't matter how many
spaces    you type; one is as good as 100.  The
end of   a line counts as a space.

One   or more   blank lines denote the  end
of  a paragraph.

Since any number of consecutive spaces are treated
like a single one, the formatting of the input
file makes no difference to
      \LaTeX,                % The \LaTeX command generates the LaTeX logo.
but it makes a difference to you.  When you use
\LaTeX, making your input file as easy to read
as possible will be a great help as you write
your document and when you change it.  This sample
file shows how you can add comments to your own input
file.

Because printing is different from typewriting,
there are a number of things that you have to do
differently when preparing an input file than if
you were just typing the document directly.
Quotation marks like
       ``this''
have to be handled specially, as do quotes within
quotes:
       ``\,`this'            % \, separates the double and single quote.
        is what I just
        wrote, not  `that'\,''.

Dashes come in three sizes: an
       intra-word
dash, a medium dash for number ranges like
       1--2,
and a punctuation
       dash---like
this.

A sentence-ending space should be larger than the
space between words within a sentence.  You
sometimes have to type special commands in
conjunction with punctuation characters to get
this right, as in the following sentence.
       Gnats, gnus, etc.\ all  % `\ ' makes an inter-word space.
       begin with G\@.         % \@ marks end-of-sentence punctuation.
You should check the spaces after periods when
reading your output to make sure you haven't
forgotten any special cases.  Generating an
ellipsis
       \ldots\               % `\ ' is needed after `\ldots' because TeX
                             % ignores spaces after command names like \ldots
                             % made from \ + letters.
                             %
                             % Note how a `%' character causes TeX to ignore
                             % the end of the input line, so these blank lines
                             % do not start a new paragraph.
                             %
with the right spacing around the periods requires
a special command.

\LaTeX\ interprets some common characters as
commands, so you must type special commands to
generate them.  These characters include the
following:
       \$ \& \% \# \{ and \}.

In printing, text is usually emphasized with an
       \emph{italic}
type style.

\begin{em}
   A long segment of text can also be emphasized
   in this way.  Text within such a segment can be
   given \emph{additional} emphasis.
\end{em}

It is sometimes necessary to prevent \LaTeX\ from
breaking a line where it might otherwise do so.
This may be at a space, as between the ``Mr.''\ and
``Jones'' in
       ``Mr.~Jones'',        % ~ produces an unbreakable interword space.
or within a word---especially when the word is a
symbol like
       \mbox{\emph{itemnum}}
that makes little sense when hyphenated across
lines.

Footnotes\footnote{This is an example of a footnote.}
pose no problem.

\LaTeX\ is good at typesetting mathematical formulas
like
       \( x-3y + z = 7 \)
or
       \( a_{1} > x^{2n} + y^{2n} > x' \)
or
       \( \ip{A}{B} = \sum_{i} a_{i} b_{i} \).
The spaces you type in a formula are
ignored.  Remember that a letter like
       $x$                   % $ ... $  and  \( ... \)  are equivalent
is a formula when it denotes a mathematical
symbol, and it should be typed as one.

\section{Displayed Text}

Text is displayed by indenting it from the left
margin.  Quotations are commonly displayed.  There
are short quotations
\begin{quote}
   This is a short quotation.  It consists of a
   single paragraph of text.  See how it is formatted.
\end{quote}
and longer ones.
\begin{quotation}
   This is a longer quotation.  It consists of two
   paragraphs of text, neither of which are
   particularly interesting.

   This is the second paragraph of the quotation.  It
   is just as dull as the first paragraph.
\end{quotation}
Another frequently-displayed structure is a list.
The following is an example of an \emph{itemized}
list.
\begin{itemize}
   \item This is the first item of an itemized list.
         Each item in the list is marked with a ``tick''.
         You don't have to worry about what kind of tick
         mark is used.

   \item This is the second item of the list.  It
         contains another list nested inside it.  The inner
         list is an \emph{enumerated} list.
         \begin{enumerate}
            \item This is the first item of an enumerated
                  list that is nested within the itemized list.

            \item This is the second item of the inner list.
                  \LaTeX\ allows you to nest lists deeper than
                  you really should.
         \end{enumerate}
         This is the rest of the second item of the outer
         list.  It is no more interesting than any other
         part of the item.
   \item This is the third item of the list.
\end{itemize}
You can even display poetry.
\begin{verse}
   There is an environment
    for verse \\             % The \\ command separates lines
   Whose features some poets % within a stanza.
   will curse.

                             % One or more blank lines separate stanzas.

   For instead of making\\
   Them do \emph{all} line breaking, \\
   It allows them to put too many words on a line when they'd rather be
   forced to be terse.
\end{verse}

Mathematical formulas may also be displayed.  A
displayed formula
is
one-line long; multiline
formulas require special formatting instructions.
   \[  \ip{\Gamma}{\psi'} = x'' + y^{2} + z_{i}^{n}\]
Don't start a paragraph with a displayed equation,
nor make one a paragraph by itself.

\end{document}               % End of document.

\section{Bench Huge Section 341}
% synthetic bench section 341
% This is a sample LaTeX input file.  (Version of 12 August 2004.)
%
% A '%' character causes TeX to ignore all remaining text on the line,
% and is used for comments like this one.

\documentclass{article}      % Specifies the document class

                             % The preamble begins here.
\title{An Example Document}  % Declares the document's title.
\author{Leslie Lamport}      % Declares the author's name.
\date{January 21, 1994}      % Deleting this command produces today's date.

\newcommand{\ip}[2]{(#1, #2)}
                             % Defines \ip{arg1}{arg2} to mean
                             % (arg1, arg2).

%\newcommand{\ip}[2]{\langle #1 | #2\rangle}
                             % This is an alternative definition of
                             % \ip that is commented out.

\begin{document}             % End of preamble and beginning of text.

\maketitle                   % Produces the title.

This is an example input file.  Comparing it with
the output it generates can show you how to
produce a simple document of your own.

\section{Ordinary Text}      % Produces section heading.  Lower-level
                             % sections are begun with similar
                             % \subsection and \subsubsection commands.

The ends  of words and sentences are marked
  by   spaces. It  doesn't matter how many
spaces    you type; one is as good as 100.  The
end of   a line counts as a space.

One   or more   blank lines denote the  end
of  a paragraph.

Since any number of consecutive spaces are treated
like a single one, the formatting of the input
file makes no difference to
      \LaTeX,                % The \LaTeX command generates the LaTeX logo.
but it makes a difference to you.  When you use
\LaTeX, making your input file as easy to read
as possible will be a great help as you write
your document and when you change it.  This sample
file shows how you can add comments to your own input
file.

Because printing is different from typewriting,
there are a number of things that you have to do
differently when preparing an input file than if
you were just typing the document directly.
Quotation marks like
       ``this''
have to be handled specially, as do quotes within
quotes:
       ``\,`this'            % \, separates the double and single quote.
        is what I just
        wrote, not  `that'\,''.

Dashes come in three sizes: an
       intra-word
dash, a medium dash for number ranges like
       1--2,
and a punctuation
       dash---like
this.

A sentence-ending space should be larger than the
space between words within a sentence.  You
sometimes have to type special commands in
conjunction with punctuation characters to get
this right, as in the following sentence.
       Gnats, gnus, etc.\ all  % `\ ' makes an inter-word space.
       begin with G\@.         % \@ marks end-of-sentence punctuation.
You should check the spaces after periods when
reading your output to make sure you haven't
forgotten any special cases.  Generating an
ellipsis
       \ldots\               % `\ ' is needed after `\ldots' because TeX
                             % ignores spaces after command names like \ldots
                             % made from \ + letters.
                             %
                             % Note how a `%' character causes TeX to ignore
                             % the end of the input line, so these blank lines
                             % do not start a new paragraph.
                             %
with the right spacing around the periods requires
a special command.

\LaTeX\ interprets some common characters as
commands, so you must type special commands to
generate them.  These characters include the
following:
       \$ \& \% \# \{ and \}.

In printing, text is usually emphasized with an
       \emph{italic}
type style.

\begin{em}
   A long segment of text can also be emphasized
   in this way.  Text within such a segment can be
   given \emph{additional} emphasis.
\end{em}

It is sometimes necessary to prevent \LaTeX\ from
breaking a line where it might otherwise do so.
This may be at a space, as between the ``Mr.''\ and
``Jones'' in
       ``Mr.~Jones'',        % ~ produces an unbreakable interword space.
or within a word---especially when the word is a
symbol like
       \mbox{\emph{itemnum}}
that makes little sense when hyphenated across
lines.

Footnotes\footnote{This is an example of a footnote.}
pose no problem.

\LaTeX\ is good at typesetting mathematical formulas
like
       \( x-3y + z = 7 \)
or
       \( a_{1} > x^{2n} + y^{2n} > x' \)
or
       \( \ip{A}{B} = \sum_{i} a_{i} b_{i} \).
The spaces you type in a formula are
ignored.  Remember that a letter like
       $x$                   % $ ... $  and  \( ... \)  are equivalent
is a formula when it denotes a mathematical
symbol, and it should be typed as one.

\section{Displayed Text}

Text is displayed by indenting it from the left
margin.  Quotations are commonly displayed.  There
are short quotations
\begin{quote}
   This is a short quotation.  It consists of a
   single paragraph of text.  See how it is formatted.
\end{quote}
and longer ones.
\begin{quotation}
   This is a longer quotation.  It consists of two
   paragraphs of text, neither of which are
   particularly interesting.

   This is the second paragraph of the quotation.  It
   is just as dull as the first paragraph.
\end{quotation}
Another frequently-displayed structure is a list.
The following is an example of an \emph{itemized}
list.
\begin{itemize}
   \item This is the first item of an itemized list.
         Each item in the list is marked with a ``tick''.
         You don't have to worry about what kind of tick
         mark is used.

   \item This is the second item of the list.  It
         contains another list nested inside it.  The inner
         list is an \emph{enumerated} list.
         \begin{enumerate}
            \item This is the first item of an enumerated
                  list that is nested within the itemized list.

            \item This is the second item of the inner list.
                  \LaTeX\ allows you to nest lists deeper than
                  you really should.
         \end{enumerate}
         This is the rest of the second item of the outer
         list.  It is no more interesting than any other
         part of the item.
   \item This is the third item of the list.
\end{itemize}
You can even display poetry.
\begin{verse}
   There is an environment
    for verse \\             % The \\ command separates lines
   Whose features some poets % within a stanza.
   will curse.

                             % One or more blank lines separate stanzas.

   For instead of making\\
   Them do \emph{all} line breaking, \\
   It allows them to put too many words on a line when they'd rather be
   forced to be terse.
\end{verse}

Mathematical formulas may also be displayed.  A
displayed formula
is
one-line long; multiline
formulas require special formatting instructions.
   \[  \ip{\Gamma}{\psi'} = x'' + y^{2} + z_{i}^{n}\]
Don't start a paragraph with a displayed equation,
nor make one a paragraph by itself.

\end{document}               % End of document.

\section{Bench Huge Section 342}
% synthetic bench section 342
% This is a sample LaTeX input file.  (Version of 12 August 2004.)
%
% A '%' character causes TeX to ignore all remaining text on the line,
% and is used for comments like this one.

\documentclass{article}      % Specifies the document class

                             % The preamble begins here.
\title{An Example Document}  % Declares the document's title.
\author{Leslie Lamport}      % Declares the author's name.
\date{January 21, 1994}      % Deleting this command produces today's date.

\newcommand{\ip}[2]{(#1, #2)}
                             % Defines \ip{arg1}{arg2} to mean
                             % (arg1, arg2).

%\newcommand{\ip}[2]{\langle #1 | #2\rangle}
                             % This is an alternative definition of
                             % \ip that is commented out.

\begin{document}             % End of preamble and beginning of text.

\maketitle                   % Produces the title.

This is an example input file.  Comparing it with
the output it generates can show you how to
produce a simple document of your own.

\section{Ordinary Text}      % Produces section heading.  Lower-level
                             % sections are begun with similar
                             % \subsection and \subsubsection commands.

The ends  of words and sentences are marked
  by   spaces. It  doesn't matter how many
spaces    you type; one is as good as 100.  The
end of   a line counts as a space.

One   or more   blank lines denote the  end
of  a paragraph.

Since any number of consecutive spaces are treated
like a single one, the formatting of the input
file makes no difference to
      \LaTeX,                % The \LaTeX command generates the LaTeX logo.
but it makes a difference to you.  When you use
\LaTeX, making your input file as easy to read
as possible will be a great help as you write
your document and when you change it.  This sample
file shows how you can add comments to your own input
file.

Because printing is different from typewriting,
there are a number of things that you have to do
differently when preparing an input file than if
you were just typing the document directly.
Quotation marks like
       ``this''
have to be handled specially, as do quotes within
quotes:
       ``\,`this'            % \, separates the double and single quote.
        is what I just
        wrote, not  `that'\,''.

Dashes come in three sizes: an
       intra-word
dash, a medium dash for number ranges like
       1--2,
and a punctuation
       dash---like
this.

A sentence-ending space should be larger than the
space between words within a sentence.  You
sometimes have to type special commands in
conjunction with punctuation characters to get
this right, as in the following sentence.
       Gnats, gnus, etc.\ all  % `\ ' makes an inter-word space.
       begin with G\@.         % \@ marks end-of-sentence punctuation.
You should check the spaces after periods when
reading your output to make sure you haven't
forgotten any special cases.  Generating an
ellipsis
       \ldots\               % `\ ' is needed after `\ldots' because TeX
                             % ignores spaces after command names like \ldots
                             % made from \ + letters.
                             %
                             % Note how a `%' character causes TeX to ignore
                             % the end of the input line, so these blank lines
                             % do not start a new paragraph.
                             %
with the right spacing around the periods requires
a special command.

\LaTeX\ interprets some common characters as
commands, so you must type special commands to
generate them.  These characters include the
following:
       \$ \& \% \# \{ and \}.

In printing, text is usually emphasized with an
       \emph{italic}
type style.

\begin{em}
   A long segment of text can also be emphasized
   in this way.  Text within such a segment can be
   given \emph{additional} emphasis.
\end{em}

It is sometimes necessary to prevent \LaTeX\ from
breaking a line where it might otherwise do so.
This may be at a space, as between the ``Mr.''\ and
``Jones'' in
       ``Mr.~Jones'',        % ~ produces an unbreakable interword space.
or within a word---especially when the word is a
symbol like
       \mbox{\emph{itemnum}}
that makes little sense when hyphenated across
lines.

Footnotes\footnote{This is an example of a footnote.}
pose no problem.

\LaTeX\ is good at typesetting mathematical formulas
like
       \( x-3y + z = 7 \)
or
       \( a_{1} > x^{2n} + y^{2n} > x' \)
or
       \( \ip{A}{B} = \sum_{i} a_{i} b_{i} \).
The spaces you type in a formula are
ignored.  Remember that a letter like
       $x$                   % $ ... $  and  \( ... \)  are equivalent
is a formula when it denotes a mathematical
symbol, and it should be typed as one.

\section{Displayed Text}

Text is displayed by indenting it from the left
margin.  Quotations are commonly displayed.  There
are short quotations
\begin{quote}
   This is a short quotation.  It consists of a
   single paragraph of text.  See how it is formatted.
\end{quote}
and longer ones.
\begin{quotation}
   This is a longer quotation.  It consists of two
   paragraphs of text, neither of which are
   particularly interesting.

   This is the second paragraph of the quotation.  It
   is just as dull as the first paragraph.
\end{quotation}
Another frequently-displayed structure is a list.
The following is an example of an \emph{itemized}
list.
\begin{itemize}
   \item This is the first item of an itemized list.
         Each item in the list is marked with a ``tick''.
         You don't have to worry about what kind of tick
         mark is used.

   \item This is the second item of the list.  It
         contains another list nested inside it.  The inner
         list is an \emph{enumerated} list.
         \begin{enumerate}
            \item This is the first item of an enumerated
                  list that is nested within the itemized list.

            \item This is the second item of the inner list.
                  \LaTeX\ allows you to nest lists deeper than
                  you really should.
         \end{enumerate}
         This is the rest of the second item of the outer
         list.  It is no more interesting than any other
         part of the item.
   \item This is the third item of the list.
\end{itemize}
You can even display poetry.
\begin{verse}
   There is an environment
    for verse \\             % The \\ command separates lines
   Whose features some poets % within a stanza.
   will curse.

                             % One or more blank lines separate stanzas.

   For instead of making\\
   Them do \emph{all} line breaking, \\
   It allows them to put too many words on a line when they'd rather be
   forced to be terse.
\end{verse}

Mathematical formulas may also be displayed.  A
displayed formula
is
one-line long; multiline
formulas require special formatting instructions.
   \[  \ip{\Gamma}{\psi'} = x'' + y^{2} + z_{i}^{n}\]
Don't start a paragraph with a displayed equation,
nor make one a paragraph by itself.

\end{document}               % End of document.

\section{Bench Huge Section 343}
% synthetic bench section 343
% This is a sample LaTeX input file.  (Version of 12 August 2004.)
%
% A '%' character causes TeX to ignore all remaining text on the line,
% and is used for comments like this one.

\documentclass{article}      % Specifies the document class

                             % The preamble begins here.
\title{An Example Document}  % Declares the document's title.
\author{Leslie Lamport}      % Declares the author's name.
\date{January 21, 1994}      % Deleting this command produces today's date.

\newcommand{\ip}[2]{(#1, #2)}
                             % Defines \ip{arg1}{arg2} to mean
                             % (arg1, arg2).

%\newcommand{\ip}[2]{\langle #1 | #2\rangle}
                             % This is an alternative definition of
                             % \ip that is commented out.

\begin{document}             % End of preamble and beginning of text.

\maketitle                   % Produces the title.

This is an example input file.  Comparing it with
the output it generates can show you how to
produce a simple document of your own.

\section{Ordinary Text}      % Produces section heading.  Lower-level
                             % sections are begun with similar
                             % \subsection and \subsubsection commands.

The ends  of words and sentences are marked
  by   spaces. It  doesn't matter how many
spaces    you type; one is as good as 100.  The
end of   a line counts as a space.

One   or more   blank lines denote the  end
of  a paragraph.

Since any number of consecutive spaces are treated
like a single one, the formatting of the input
file makes no difference to
      \LaTeX,                % The \LaTeX command generates the LaTeX logo.
but it makes a difference to you.  When you use
\LaTeX, making your input file as easy to read
as possible will be a great help as you write
your document and when you change it.  This sample
file shows how you can add comments to your own input
file.

Because printing is different from typewriting,
there are a number of things that you have to do
differently when preparing an input file than if
you were just typing the document directly.
Quotation marks like
       ``this''
have to be handled specially, as do quotes within
quotes:
       ``\,`this'            % \, separates the double and single quote.
        is what I just
        wrote, not  `that'\,''.

Dashes come in three sizes: an
       intra-word
dash, a medium dash for number ranges like
       1--2,
and a punctuation
       dash---like
this.

A sentence-ending space should be larger than the
space between words within a sentence.  You
sometimes have to type special commands in
conjunction with punctuation characters to get
this right, as in the following sentence.
       Gnats, gnus, etc.\ all  % `\ ' makes an inter-word space.
       begin with G\@.         % \@ marks end-of-sentence punctuation.
You should check the spaces after periods when
reading your output to make sure you haven't
forgotten any special cases.  Generating an
ellipsis
       \ldots\               % `\ ' is needed after `\ldots' because TeX
                             % ignores spaces after command names like \ldots
                             % made from \ + letters.
                             %
                             % Note how a `%' character causes TeX to ignore
                             % the end of the input line, so these blank lines
                             % do not start a new paragraph.
                             %
with the right spacing around the periods requires
a special command.

\LaTeX\ interprets some common characters as
commands, so you must type special commands to
generate them.  These characters include the
following:
       \$ \& \% \# \{ and \}.

In printing, text is usually emphasized with an
       \emph{italic}
type style.

\begin{em}
   A long segment of text can also be emphasized
   in this way.  Text within such a segment can be
   given \emph{additional} emphasis.
\end{em}

It is sometimes necessary to prevent \LaTeX\ from
breaking a line where it might otherwise do so.
This may be at a space, as between the ``Mr.''\ and
``Jones'' in
       ``Mr.~Jones'',        % ~ produces an unbreakable interword space.
or within a word---especially when the word is a
symbol like
       \mbox{\emph{itemnum}}
that makes little sense when hyphenated across
lines.

Footnotes\footnote{This is an example of a footnote.}
pose no problem.

\LaTeX\ is good at typesetting mathematical formulas
like
       \( x-3y + z = 7 \)
or
       \( a_{1} > x^{2n} + y^{2n} > x' \)
or
       \( \ip{A}{B} = \sum_{i} a_{i} b_{i} \).
The spaces you type in a formula are
ignored.  Remember that a letter like
       $x$                   % $ ... $  and  \( ... \)  are equivalent
is a formula when it denotes a mathematical
symbol, and it should be typed as one.

\section{Displayed Text}

Text is displayed by indenting it from the left
margin.  Quotations are commonly displayed.  There
are short quotations
\begin{quote}
   This is a short quotation.  It consists of a
   single paragraph of text.  See how it is formatted.
\end{quote}
and longer ones.
\begin{quotation}
   This is a longer quotation.  It consists of two
   paragraphs of text, neither of which are
   particularly interesting.

   This is the second paragraph of the quotation.  It
   is just as dull as the first paragraph.
\end{quotation}
Another frequently-displayed structure is a list.
The following is an example of an \emph{itemized}
list.
\begin{itemize}
   \item This is the first item of an itemized list.
         Each item in the list is marked with a ``tick''.
         You don't have to worry about what kind of tick
         mark is used.

   \item This is the second item of the list.  It
         contains another list nested inside it.  The inner
         list is an \emph{enumerated} list.
         \begin{enumerate}
            \item This is the first item of an enumerated
                  list that is nested within the itemized list.

            \item This is the second item of the inner list.
                  \LaTeX\ allows you to nest lists deeper than
                  you really should.
         \end{enumerate}
         This is the rest of the second item of the outer
         list.  It is no more interesting than any other
         part of the item.
   \item This is the third item of the list.
\end{itemize}
You can even display poetry.
\begin{verse}
   There is an environment
    for verse \\             % The \\ command separates lines
   Whose features some poets % within a stanza.
   will curse.

                             % One or more blank lines separate stanzas.

   For instead of making\\
   Them do \emph{all} line breaking, \\
   It allows them to put too many words on a line when they'd rather be
   forced to be terse.
\end{verse}

Mathematical formulas may also be displayed.  A
displayed formula
is
one-line long; multiline
formulas require special formatting instructions.
   \[  \ip{\Gamma}{\psi'} = x'' + y^{2} + z_{i}^{n}\]
Don't start a paragraph with a displayed equation,
nor make one a paragraph by itself.

\end{document}               % End of document.

\section{Bench Huge Section 344}
% synthetic bench section 344
% This is a sample LaTeX input file.  (Version of 12 August 2004.)
%
% A '%' character causes TeX to ignore all remaining text on the line,
% and is used for comments like this one.

\documentclass{article}      % Specifies the document class

                             % The preamble begins here.
\title{An Example Document}  % Declares the document's title.
\author{Leslie Lamport}      % Declares the author's name.
\date{January 21, 1994}      % Deleting this command produces today's date.

\newcommand{\ip}[2]{(#1, #2)}
                             % Defines \ip{arg1}{arg2} to mean
                             % (arg1, arg2).

%\newcommand{\ip}[2]{\langle #1 | #2\rangle}
                             % This is an alternative definition of
                             % \ip that is commented out.

\begin{document}             % End of preamble and beginning of text.

\maketitle                   % Produces the title.

This is an example input file.  Comparing it with
the output it generates can show you how to
produce a simple document of your own.

\section{Ordinary Text}      % Produces section heading.  Lower-level
                             % sections are begun with similar
                             % \subsection and \subsubsection commands.

The ends  of words and sentences are marked
  by   spaces. It  doesn't matter how many
spaces    you type; one is as good as 100.  The
end of   a line counts as a space.

One   or more   blank lines denote the  end
of  a paragraph.

Since any number of consecutive spaces are treated
like a single one, the formatting of the input
file makes no difference to
      \LaTeX,                % The \LaTeX command generates the LaTeX logo.
but it makes a difference to you.  When you use
\LaTeX, making your input file as easy to read
as possible will be a great help as you write
your document and when you change it.  This sample
file shows how you can add comments to your own input
file.

Because printing is different from typewriting,
there are a number of things that you have to do
differently when preparing an input file than if
you were just typing the document directly.
Quotation marks like
       ``this''
have to be handled specially, as do quotes within
quotes:
       ``\,`this'            % \, separates the double and single quote.
        is what I just
        wrote, not  `that'\,''.

Dashes come in three sizes: an
       intra-word
dash, a medium dash for number ranges like
       1--2,
and a punctuation
       dash---like
this.

A sentence-ending space should be larger than the
space between words within a sentence.  You
sometimes have to type special commands in
conjunction with punctuation characters to get
this right, as in the following sentence.
       Gnats, gnus, etc.\ all  % `\ ' makes an inter-word space.
       begin with G\@.         % \@ marks end-of-sentence punctuation.
You should check the spaces after periods when
reading your output to make sure you haven't
forgotten any special cases.  Generating an
ellipsis
       \ldots\               % `\ ' is needed after `\ldots' because TeX
                             % ignores spaces after command names like \ldots
                             % made from \ + letters.
                             %
                             % Note how a `%' character causes TeX to ignore
                             % the end of the input line, so these blank lines
                             % do not start a new paragraph.
                             %
with the right spacing around the periods requires
a special command.

\LaTeX\ interprets some common characters as
commands, so you must type special commands to
generate them.  These characters include the
following:
       \$ \& \% \# \{ and \}.

In printing, text is usually emphasized with an
       \emph{italic}
type style.

\begin{em}
   A long segment of text can also be emphasized
   in this way.  Text within such a segment can be
   given \emph{additional} emphasis.
\end{em}

It is sometimes necessary to prevent \LaTeX\ from
breaking a line where it might otherwise do so.
This may be at a space, as between the ``Mr.''\ and
``Jones'' in
       ``Mr.~Jones'',        % ~ produces an unbreakable interword space.
or within a word---especially when the word is a
symbol like
       \mbox{\emph{itemnum}}
that makes little sense when hyphenated across
lines.

Footnotes\footnote{This is an example of a footnote.}
pose no problem.

\LaTeX\ is good at typesetting mathematical formulas
like
       \( x-3y + z = 7 \)
or
       \( a_{1} > x^{2n} + y^{2n} > x' \)
or
       \( \ip{A}{B} = \sum_{i} a_{i} b_{i} \).
The spaces you type in a formula are
ignored.  Remember that a letter like
       $x$                   % $ ... $  and  \( ... \)  are equivalent
is a formula when it denotes a mathematical
symbol, and it should be typed as one.

\section{Displayed Text}

Text is displayed by indenting it from the left
margin.  Quotations are commonly displayed.  There
are short quotations
\begin{quote}
   This is a short quotation.  It consists of a
   single paragraph of text.  See how it is formatted.
\end{quote}
and longer ones.
\begin{quotation}
   This is a longer quotation.  It consists of two
   paragraphs of text, neither of which are
   particularly interesting.

   This is the second paragraph of the quotation.  It
   is just as dull as the first paragraph.
\end{quotation}
Another frequently-displayed structure is a list.
The following is an example of an \emph{itemized}
list.
\begin{itemize}
   \item This is the first item of an itemized list.
         Each item in the list is marked with a ``tick''.
         You don't have to worry about what kind of tick
         mark is used.

   \item This is the second item of the list.  It
         contains another list nested inside it.  The inner
         list is an \emph{enumerated} list.
         \begin{enumerate}
            \item This is the first item of an enumerated
                  list that is nested within the itemized list.

            \item This is the second item of the inner list.
                  \LaTeX\ allows you to nest lists deeper than
                  you really should.
         \end{enumerate}
         This is the rest of the second item of the outer
         list.  It is no more interesting than any other
         part of the item.
   \item This is the third item of the list.
\end{itemize}
You can even display poetry.
\begin{verse}
   There is an environment
    for verse \\             % The \\ command separates lines
   Whose features some poets % within a stanza.
   will curse.

                             % One or more blank lines separate stanzas.

   For instead of making\\
   Them do \emph{all} line breaking, \\
   It allows them to put too many words on a line when they'd rather be
   forced to be terse.
\end{verse}

Mathematical formulas may also be displayed.  A
displayed formula
is
one-line long; multiline
formulas require special formatting instructions.
   \[  \ip{\Gamma}{\psi'} = x'' + y^{2} + z_{i}^{n}\]
Don't start a paragraph with a displayed equation,
nor make one a paragraph by itself.

\end{document}               % End of document.

\section{Bench Huge Section 345}
% synthetic bench section 345
% This is a sample LaTeX input file.  (Version of 12 August 2004.)
%
% A '%' character causes TeX to ignore all remaining text on the line,
% and is used for comments like this one.

\documentclass{article}      % Specifies the document class

                             % The preamble begins here.
\title{An Example Document}  % Declares the document's title.
\author{Leslie Lamport}      % Declares the author's name.
\date{January 21, 1994}      % Deleting this command produces today's date.

\newcommand{\ip}[2]{(#1, #2)}
                             % Defines \ip{arg1}{arg2} to mean
                             % (arg1, arg2).

%\newcommand{\ip}[2]{\langle #1 | #2\rangle}
                             % This is an alternative definition of
                             % \ip that is commented out.

\begin{document}             % End of preamble and beginning of text.

\maketitle                   % Produces the title.

This is an example input file.  Comparing it with
the output it generates can show you how to
produce a simple document of your own.

\section{Ordinary Text}      % Produces section heading.  Lower-level
                             % sections are begun with similar
                             % \subsection and \subsubsection commands.

The ends  of words and sentences are marked
  by   spaces. It  doesn't matter how many
spaces    you type; one is as good as 100.  The
end of   a line counts as a space.

One   or more   blank lines denote the  end
of  a paragraph.

Since any number of consecutive spaces are treated
like a single one, the formatting of the input
file makes no difference to
      \LaTeX,                % The \LaTeX command generates the LaTeX logo.
but it makes a difference to you.  When you use
\LaTeX, making your input file as easy to read
as possible will be a great help as you write
your document and when you change it.  This sample
file shows how you can add comments to your own input
file.

Because printing is different from typewriting,
there are a number of things that you have to do
differently when preparing an input file than if
you were just typing the document directly.
Quotation marks like
       ``this''
have to be handled specially, as do quotes within
quotes:
       ``\,`this'            % \, separates the double and single quote.
        is what I just
        wrote, not  `that'\,''.

Dashes come in three sizes: an
       intra-word
dash, a medium dash for number ranges like
       1--2,
and a punctuation
       dash---like
this.

A sentence-ending space should be larger than the
space between words within a sentence.  You
sometimes have to type special commands in
conjunction with punctuation characters to get
this right, as in the following sentence.
       Gnats, gnus, etc.\ all  % `\ ' makes an inter-word space.
       begin with G\@.         % \@ marks end-of-sentence punctuation.
You should check the spaces after periods when
reading your output to make sure you haven't
forgotten any special cases.  Generating an
ellipsis
       \ldots\               % `\ ' is needed after `\ldots' because TeX
                             % ignores spaces after command names like \ldots
                             % made from \ + letters.
                             %
                             % Note how a `%' character causes TeX to ignore
                             % the end of the input line, so these blank lines
                             % do not start a new paragraph.
                             %
with the right spacing around the periods requires
a special command.

\LaTeX\ interprets some common characters as
commands, so you must type special commands to
generate them.  These characters include the
following:
       \$ \& \% \# \{ and \}.

In printing, text is usually emphasized with an
       \emph{italic}
type style.

\begin{em}
   A long segment of text can also be emphasized
   in this way.  Text within such a segment can be
   given \emph{additional} emphasis.
\end{em}

It is sometimes necessary to prevent \LaTeX\ from
breaking a line where it might otherwise do so.
This may be at a space, as between the ``Mr.''\ and
``Jones'' in
       ``Mr.~Jones'',        % ~ produces an unbreakable interword space.
or within a word---especially when the word is a
symbol like
       \mbox{\emph{itemnum}}
that makes little sense when hyphenated across
lines.

Footnotes\footnote{This is an example of a footnote.}
pose no problem.

\LaTeX\ is good at typesetting mathematical formulas
like
       \( x-3y + z = 7 \)
or
       \( a_{1} > x^{2n} + y^{2n} > x' \)
or
       \( \ip{A}{B} = \sum_{i} a_{i} b_{i} \).
The spaces you type in a formula are
ignored.  Remember that a letter like
       $x$                   % $ ... $  and  \( ... \)  are equivalent
is a formula when it denotes a mathematical
symbol, and it should be typed as one.

\section{Displayed Text}

Text is displayed by indenting it from the left
margin.  Quotations are commonly displayed.  There
are short quotations
\begin{quote}
   This is a short quotation.  It consists of a
   single paragraph of text.  See how it is formatted.
\end{quote}
and longer ones.
\begin{quotation}
   This is a longer quotation.  It consists of two
   paragraphs of text, neither of which are
   particularly interesting.

   This is the second paragraph of the quotation.  It
   is just as dull as the first paragraph.
\end{quotation}
Another frequently-displayed structure is a list.
The following is an example of an \emph{itemized}
list.
\begin{itemize}
   \item This is the first item of an itemized list.
         Each item in the list is marked with a ``tick''.
         You don't have to worry about what kind of tick
         mark is used.

   \item This is the second item of the list.  It
         contains another list nested inside it.  The inner
         list is an \emph{enumerated} list.
         \begin{enumerate}
            \item This is the first item of an enumerated
                  list that is nested within the itemized list.

            \item This is the second item of the inner list.
                  \LaTeX\ allows you to nest lists deeper than
                  you really should.
         \end{enumerate}
         This is the rest of the second item of the outer
         list.  It is no more interesting than any other
         part of the item.
   \item This is the third item of the list.
\end{itemize}
You can even display poetry.
\begin{verse}
   There is an environment
    for verse \\             % The \\ command separates lines
   Whose features some poets % within a stanza.
   will curse.

                             % One or more blank lines separate stanzas.

   For instead of making\\
   Them do \emph{all} line breaking, \\
   It allows them to put too many words on a line when they'd rather be
   forced to be terse.
\end{verse}

Mathematical formulas may also be displayed.  A
displayed formula
is
one-line long; multiline
formulas require special formatting instructions.
   \[  \ip{\Gamma}{\psi'} = x'' + y^{2} + z_{i}^{n}\]
Don't start a paragraph with a displayed equation,
nor make one a paragraph by itself.

\end{document}               % End of document.

\section{Bench Huge Section 346}
% synthetic bench section 346
% This is a sample LaTeX input file.  (Version of 12 August 2004.)
%
% A '%' character causes TeX to ignore all remaining text on the line,
% and is used for comments like this one.

\documentclass{article}      % Specifies the document class

                             % The preamble begins here.
\title{An Example Document}  % Declares the document's title.
\author{Leslie Lamport}      % Declares the author's name.
\date{January 21, 1994}      % Deleting this command produces today's date.

\newcommand{\ip}[2]{(#1, #2)}
                             % Defines \ip{arg1}{arg2} to mean
                             % (arg1, arg2).

%\newcommand{\ip}[2]{\langle #1 | #2\rangle}
                             % This is an alternative definition of
                             % \ip that is commented out.

\begin{document}             % End of preamble and beginning of text.

\maketitle                   % Produces the title.

This is an example input file.  Comparing it with
the output it generates can show you how to
produce a simple document of your own.

\section{Ordinary Text}      % Produces section heading.  Lower-level
                             % sections are begun with similar
                             % \subsection and \subsubsection commands.

The ends  of words and sentences are marked
  by   spaces. It  doesn't matter how many
spaces    you type; one is as good as 100.  The
end of   a line counts as a space.

One   or more   blank lines denote the  end
of  a paragraph.

Since any number of consecutive spaces are treated
like a single one, the formatting of the input
file makes no difference to
      \LaTeX,                % The \LaTeX command generates the LaTeX logo.
but it makes a difference to you.  When you use
\LaTeX, making your input file as easy to read
as possible will be a great help as you write
your document and when you change it.  This sample
file shows how you can add comments to your own input
file.

Because printing is different from typewriting,
there are a number of things that you have to do
differently when preparing an input file than if
you were just typing the document directly.
Quotation marks like
       ``this''
have to be handled specially, as do quotes within
quotes:
       ``\,`this'            % \, separates the double and single quote.
        is what I just
        wrote, not  `that'\,''.

Dashes come in three sizes: an
       intra-word
dash, a medium dash for number ranges like
       1--2,
and a punctuation
       dash---like
this.

A sentence-ending space should be larger than the
space between words within a sentence.  You
sometimes have to type special commands in
conjunction with punctuation characters to get
this right, as in the following sentence.
       Gnats, gnus, etc.\ all  % `\ ' makes an inter-word space.
       begin with G\@.         % \@ marks end-of-sentence punctuation.
You should check the spaces after periods when
reading your output to make sure you haven't
forgotten any special cases.  Generating an
ellipsis
       \ldots\               % `\ ' is needed after `\ldots' because TeX
                             % ignores spaces after command names like \ldots
                             % made from \ + letters.
                             %
                             % Note how a `%' character causes TeX to ignore
                             % the end of the input line, so these blank lines
                             % do not start a new paragraph.
                             %
with the right spacing around the periods requires
a special command.

\LaTeX\ interprets some common characters as
commands, so you must type special commands to
generate them.  These characters include the
following:
       \$ \& \% \# \{ and \}.

In printing, text is usually emphasized with an
       \emph{italic}
type style.

\begin{em}
   A long segment of text can also be emphasized
   in this way.  Text within such a segment can be
   given \emph{additional} emphasis.
\end{em}

It is sometimes necessary to prevent \LaTeX\ from
breaking a line where it might otherwise do so.
This may be at a space, as between the ``Mr.''\ and
``Jones'' in
       ``Mr.~Jones'',        % ~ produces an unbreakable interword space.
or within a word---especially when the word is a
symbol like
       \mbox{\emph{itemnum}}
that makes little sense when hyphenated across
lines.

Footnotes\footnote{This is an example of a footnote.}
pose no problem.

\LaTeX\ is good at typesetting mathematical formulas
like
       \( x-3y + z = 7 \)
or
       \( a_{1} > x^{2n} + y^{2n} > x' \)
or
       \( \ip{A}{B} = \sum_{i} a_{i} b_{i} \).
The spaces you type in a formula are
ignored.  Remember that a letter like
       $x$                   % $ ... $  and  \( ... \)  are equivalent
is a formula when it denotes a mathematical
symbol, and it should be typed as one.

\section{Displayed Text}

Text is displayed by indenting it from the left
margin.  Quotations are commonly displayed.  There
are short quotations
\begin{quote}
   This is a short quotation.  It consists of a
   single paragraph of text.  See how it is formatted.
\end{quote}
and longer ones.
\begin{quotation}
   This is a longer quotation.  It consists of two
   paragraphs of text, neither of which are
   particularly interesting.

   This is the second paragraph of the quotation.  It
   is just as dull as the first paragraph.
\end{quotation}
Another frequently-displayed structure is a list.
The following is an example of an \emph{itemized}
list.
\begin{itemize}
   \item This is the first item of an itemized list.
         Each item in the list is marked with a ``tick''.
         You don't have to worry about what kind of tick
         mark is used.

   \item This is the second item of the list.  It
         contains another list nested inside it.  The inner
         list is an \emph{enumerated} list.
         \begin{enumerate}
            \item This is the first item of an enumerated
                  list that is nested within the itemized list.

            \item This is the second item of the inner list.
                  \LaTeX\ allows you to nest lists deeper than
                  you really should.
         \end{enumerate}
         This is the rest of the second item of the outer
         list.  It is no more interesting than any other
         part of the item.
   \item This is the third item of the list.
\end{itemize}
You can even display poetry.
\begin{verse}
   There is an environment
    for verse \\             % The \\ command separates lines
   Whose features some poets % within a stanza.
   will curse.

                             % One or more blank lines separate stanzas.

   For instead of making\\
   Them do \emph{all} line breaking, \\
   It allows them to put too many words on a line when they'd rather be
   forced to be terse.
\end{verse}

Mathematical formulas may also be displayed.  A
displayed formula
is
one-line long; multiline
formulas require special formatting instructions.
   \[  \ip{\Gamma}{\psi'} = x'' + y^{2} + z_{i}^{n}\]
Don't start a paragraph with a displayed equation,
nor make one a paragraph by itself.

\end{document}               % End of document.

\section{Bench Huge Section 347}
% synthetic bench section 347
% This is a sample LaTeX input file.  (Version of 12 August 2004.)
%
% A '%' character causes TeX to ignore all remaining text on the line,
% and is used for comments like this one.

\documentclass{article}      % Specifies the document class

                             % The preamble begins here.
\title{An Example Document}  % Declares the document's title.
\author{Leslie Lamport}      % Declares the author's name.
\date{January 21, 1994}      % Deleting this command produces today's date.

\newcommand{\ip}[2]{(#1, #2)}
                             % Defines \ip{arg1}{arg2} to mean
                             % (arg1, arg2).

%\newcommand{\ip}[2]{\langle #1 | #2\rangle}
                             % This is an alternative definition of
                             % \ip that is commented out.

\begin{document}             % End of preamble and beginning of text.

\maketitle                   % Produces the title.

This is an example input file.  Comparing it with
the output it generates can show you how to
produce a simple document of your own.

\section{Ordinary Text}      % Produces section heading.  Lower-level
                             % sections are begun with similar
                             % \subsection and \subsubsection commands.

The ends  of words and sentences are marked
  by   spaces. It  doesn't matter how many
spaces    you type; one is as good as 100.  The
end of   a line counts as a space.

One   or more   blank lines denote the  end
of  a paragraph.

Since any number of consecutive spaces are treated
like a single one, the formatting of the input
file makes no difference to
      \LaTeX,                % The \LaTeX command generates the LaTeX logo.
but it makes a difference to you.  When you use
\LaTeX, making your input file as easy to read
as possible will be a great help as you write
your document and when you change it.  This sample
file shows how you can add comments to your own input
file.

Because printing is different from typewriting,
there are a number of things that you have to do
differently when preparing an input file than if
you were just typing the document directly.
Quotation marks like
       ``this''
have to be handled specially, as do quotes within
quotes:
       ``\,`this'            % \, separates the double and single quote.
        is what I just
        wrote, not  `that'\,''.

Dashes come in three sizes: an
       intra-word
dash, a medium dash for number ranges like
       1--2,
and a punctuation
       dash---like
this.

A sentence-ending space should be larger than the
space between words within a sentence.  You
sometimes have to type special commands in
conjunction with punctuation characters to get
this right, as in the following sentence.
       Gnats, gnus, etc.\ all  % `\ ' makes an inter-word space.
       begin with G\@.         % \@ marks end-of-sentence punctuation.
You should check the spaces after periods when
reading your output to make sure you haven't
forgotten any special cases.  Generating an
ellipsis
       \ldots\               % `\ ' is needed after `\ldots' because TeX
                             % ignores spaces after command names like \ldots
                             % made from \ + letters.
                             %
                             % Note how a `%' character causes TeX to ignore
                             % the end of the input line, so these blank lines
                             % do not start a new paragraph.
                             %
with the right spacing around the periods requires
a special command.

\LaTeX\ interprets some common characters as
commands, so you must type special commands to
generate them.  These characters include the
following:
       \$ \& \% \# \{ and \}.

In printing, text is usually emphasized with an
       \emph{italic}
type style.

\begin{em}
   A long segment of text can also be emphasized
   in this way.  Text within such a segment can be
   given \emph{additional} emphasis.
\end{em}

It is sometimes necessary to prevent \LaTeX\ from
breaking a line where it might otherwise do so.
This may be at a space, as between the ``Mr.''\ and
``Jones'' in
       ``Mr.~Jones'',        % ~ produces an unbreakable interword space.
or within a word---especially when the word is a
symbol like
       \mbox{\emph{itemnum}}
that makes little sense when hyphenated across
lines.

Footnotes\footnote{This is an example of a footnote.}
pose no problem.

\LaTeX\ is good at typesetting mathematical formulas
like
       \( x-3y + z = 7 \)
or
       \( a_{1} > x^{2n} + y^{2n} > x' \)
or
       \( \ip{A}{B} = \sum_{i} a_{i} b_{i} \).
The spaces you type in a formula are
ignored.  Remember that a letter like
       $x$                   % $ ... $  and  \( ... \)  are equivalent
is a formula when it denotes a mathematical
symbol, and it should be typed as one.

\section{Displayed Text}

Text is displayed by indenting it from the left
margin.  Quotations are commonly displayed.  There
are short quotations
\begin{quote}
   This is a short quotation.  It consists of a
   single paragraph of text.  See how it is formatted.
\end{quote}
and longer ones.
\begin{quotation}
   This is a longer quotation.  It consists of two
   paragraphs of text, neither of which are
   particularly interesting.

   This is the second paragraph of the quotation.  It
   is just as dull as the first paragraph.
\end{quotation}
Another frequently-displayed structure is a list.
The following is an example of an \emph{itemized}
list.
\begin{itemize}
   \item This is the first item of an itemized list.
         Each item in the list is marked with a ``tick''.
         You don't have to worry about what kind of tick
         mark is used.

   \item This is the second item of the list.  It
         contains another list nested inside it.  The inner
         list is an \emph{enumerated} list.
         \begin{enumerate}
            \item This is the first item of an enumerated
                  list that is nested within the itemized list.

            \item This is the second item of the inner list.
                  \LaTeX\ allows you to nest lists deeper than
                  you really should.
         \end{enumerate}
         This is the rest of the second item of the outer
         list.  It is no more interesting than any other
         part of the item.
   \item This is the third item of the list.
\end{itemize}
You can even display poetry.
\begin{verse}
   There is an environment
    for verse \\             % The \\ command separates lines
   Whose features some poets % within a stanza.
   will curse.

                             % One or more blank lines separate stanzas.

   For instead of making\\
   Them do \emph{all} line breaking, \\
   It allows them to put too many words on a line when they'd rather be
   forced to be terse.
\end{verse}

Mathematical formulas may also be displayed.  A
displayed formula
is
one-line long; multiline
formulas require special formatting instructions.
   \[  \ip{\Gamma}{\psi'} = x'' + y^{2} + z_{i}^{n}\]
Don't start a paragraph with a displayed equation,
nor make one a paragraph by itself.

\end{document}               % End of document.

\section{Bench Huge Section 348}
% synthetic bench section 348
% This is a sample LaTeX input file.  (Version of 12 August 2004.)
%
% A '%' character causes TeX to ignore all remaining text on the line,
% and is used for comments like this one.

\documentclass{article}      % Specifies the document class

                             % The preamble begins here.
\title{An Example Document}  % Declares the document's title.
\author{Leslie Lamport}      % Declares the author's name.
\date{January 21, 1994}      % Deleting this command produces today's date.

\newcommand{\ip}[2]{(#1, #2)}
                             % Defines \ip{arg1}{arg2} to mean
                             % (arg1, arg2).

%\newcommand{\ip}[2]{\langle #1 | #2\rangle}
                             % This is an alternative definition of
                             % \ip that is commented out.

\begin{document}             % End of preamble and beginning of text.

\maketitle                   % Produces the title.

This is an example input file.  Comparing it with
the output it generates can show you how to
produce a simple document of your own.

\section{Ordinary Text}      % Produces section heading.  Lower-level
                             % sections are begun with similar
                             % \subsection and \subsubsection commands.

The ends  of words and sentences are marked
  by   spaces. It  doesn't matter how many
spaces    you type; one is as good as 100.  The
end of   a line counts as a space.

One   or more   blank lines denote the  end
of  a paragraph.

Since any number of consecutive spaces are treated
like a single one, the formatting of the input
file makes no difference to
      \LaTeX,                % The \LaTeX command generates the LaTeX logo.
but it makes a difference to you.  When you use
\LaTeX, making your input file as easy to read
as possible will be a great help as you write
your document and when you change it.  This sample
file shows how you can add comments to your own input
file.

Because printing is different from typewriting,
there are a number of things that you have to do
differently when preparing an input file than if
you were just typing the document directly.
Quotation marks like
       ``this''
have to be handled specially, as do quotes within
quotes:
       ``\,`this'            % \, separates the double and single quote.
        is what I just
        wrote, not  `that'\,''.

Dashes come in three sizes: an
       intra-word
dash, a medium dash for number ranges like
       1--2,
and a punctuation
       dash---like
this.

A sentence-ending space should be larger than the
space between words within a sentence.  You
sometimes have to type special commands in
conjunction with punctuation characters to get
this right, as in the following sentence.
       Gnats, gnus, etc.\ all  % `\ ' makes an inter-word space.
       begin with G\@.         % \@ marks end-of-sentence punctuation.
You should check the spaces after periods when
reading your output to make sure you haven't
forgotten any special cases.  Generating an
ellipsis
       \ldots\               % `\ ' is needed after `\ldots' because TeX
                             % ignores spaces after command names like \ldots
                             % made from \ + letters.
                             %
                             % Note how a `%' character causes TeX to ignore
                             % the end of the input line, so these blank lines
                             % do not start a new paragraph.
                             %
with the right spacing around the periods requires
a special command.

\LaTeX\ interprets some common characters as
commands, so you must type special commands to
generate them.  These characters include the
following:
       \$ \& \% \# \{ and \}.

In printing, text is usually emphasized with an
       \emph{italic}
type style.

\begin{em}
   A long segment of text can also be emphasized
   in this way.  Text within such a segment can be
   given \emph{additional} emphasis.
\end{em}

It is sometimes necessary to prevent \LaTeX\ from
breaking a line where it might otherwise do so.
This may be at a space, as between the ``Mr.''\ and
``Jones'' in
       ``Mr.~Jones'',        % ~ produces an unbreakable interword space.
or within a word---especially when the word is a
symbol like
       \mbox{\emph{itemnum}}
that makes little sense when hyphenated across
lines.

Footnotes\footnote{This is an example of a footnote.}
pose no problem.

\LaTeX\ is good at typesetting mathematical formulas
like
       \( x-3y + z = 7 \)
or
       \( a_{1} > x^{2n} + y^{2n} > x' \)
or
       \( \ip{A}{B} = \sum_{i} a_{i} b_{i} \).
The spaces you type in a formula are
ignored.  Remember that a letter like
       $x$                   % $ ... $  and  \( ... \)  are equivalent
is a formula when it denotes a mathematical
symbol, and it should be typed as one.

\section{Displayed Text}

Text is displayed by indenting it from the left
margin.  Quotations are commonly displayed.  There
are short quotations
\begin{quote}
   This is a short quotation.  It consists of a
   single paragraph of text.  See how it is formatted.
\end{quote}
and longer ones.
\begin{quotation}
   This is a longer quotation.  It consists of two
   paragraphs of text, neither of which are
   particularly interesting.

   This is the second paragraph of the quotation.  It
   is just as dull as the first paragraph.
\end{quotation}
Another frequently-displayed structure is a list.
The following is an example of an \emph{itemized}
list.
\begin{itemize}
   \item This is the first item of an itemized list.
         Each item in the list is marked with a ``tick''.
         You don't have to worry about what kind of tick
         mark is used.

   \item This is the second item of the list.  It
         contains another list nested inside it.  The inner
         list is an \emph{enumerated} list.
         \begin{enumerate}
            \item This is the first item of an enumerated
                  list that is nested within the itemized list.

            \item This is the second item of the inner list.
                  \LaTeX\ allows you to nest lists deeper than
                  you really should.
         \end{enumerate}
         This is the rest of the second item of the outer
         list.  It is no more interesting than any other
         part of the item.
   \item This is the third item of the list.
\end{itemize}
You can even display poetry.
\begin{verse}
   There is an environment
    for verse \\             % The \\ command separates lines
   Whose features some poets % within a stanza.
   will curse.

                             % One or more blank lines separate stanzas.

   For instead of making\\
   Them do \emph{all} line breaking, \\
   It allows them to put too many words on a line when they'd rather be
   forced to be terse.
\end{verse}

Mathematical formulas may also be displayed.  A
displayed formula
is
one-line long; multiline
formulas require special formatting instructions.
   \[  \ip{\Gamma}{\psi'} = x'' + y^{2} + z_{i}^{n}\]
Don't start a paragraph with a displayed equation,
nor make one a paragraph by itself.

\end{document}               % End of document.

\section{Bench Huge Section 349}
% synthetic bench section 349
% This is a sample LaTeX input file.  (Version of 12 August 2004.)
%
% A '%' character causes TeX to ignore all remaining text on the line,
% and is used for comments like this one.

\documentclass{article}      % Specifies the document class

                             % The preamble begins here.
\title{An Example Document}  % Declares the document's title.
\author{Leslie Lamport}      % Declares the author's name.
\date{January 21, 1994}      % Deleting this command produces today's date.

\newcommand{\ip}[2]{(#1, #2)}
                             % Defines \ip{arg1}{arg2} to mean
                             % (arg1, arg2).

%\newcommand{\ip}[2]{\langle #1 | #2\rangle}
                             % This is an alternative definition of
                             % \ip that is commented out.

\begin{document}             % End of preamble and beginning of text.

\maketitle                   % Produces the title.

This is an example input file.  Comparing it with
the output it generates can show you how to
produce a simple document of your own.

\section{Ordinary Text}      % Produces section heading.  Lower-level
                             % sections are begun with similar
                             % \subsection and \subsubsection commands.

The ends  of words and sentences are marked
  by   spaces. It  doesn't matter how many
spaces    you type; one is as good as 100.  The
end of   a line counts as a space.

One   or more   blank lines denote the  end
of  a paragraph.

Since any number of consecutive spaces are treated
like a single one, the formatting of the input
file makes no difference to
      \LaTeX,                % The \LaTeX command generates the LaTeX logo.
but it makes a difference to you.  When you use
\LaTeX, making your input file as easy to read
as possible will be a great help as you write
your document and when you change it.  This sample
file shows how you can add comments to your own input
file.

Because printing is different from typewriting,
there are a number of things that you have to do
differently when preparing an input file than if
you were just typing the document directly.
Quotation marks like
       ``this''
have to be handled specially, as do quotes within
quotes:
       ``\,`this'            % \, separates the double and single quote.
        is what I just
        wrote, not  `that'\,''.

Dashes come in three sizes: an
       intra-word
dash, a medium dash for number ranges like
       1--2,
and a punctuation
       dash---like
this.

A sentence-ending space should be larger than the
space between words within a sentence.  You
sometimes have to type special commands in
conjunction with punctuation characters to get
this right, as in the following sentence.
       Gnats, gnus, etc.\ all  % `\ ' makes an inter-word space.
       begin with G\@.         % \@ marks end-of-sentence punctuation.
You should check the spaces after periods when
reading your output to make sure you haven't
forgotten any special cases.  Generating an
ellipsis
       \ldots\               % `\ ' is needed after `\ldots' because TeX
                             % ignores spaces after command names like \ldots
                             % made from \ + letters.
                             %
                             % Note how a `%' character causes TeX to ignore
                             % the end of the input line, so these blank lines
                             % do not start a new paragraph.
                             %
with the right spacing around the periods requires
a special command.

\LaTeX\ interprets some common characters as
commands, so you must type special commands to
generate them.  These characters include the
following:
       \$ \& \% \# \{ and \}.

In printing, text is usually emphasized with an
       \emph{italic}
type style.

\begin{em}
   A long segment of text can also be emphasized
   in this way.  Text within such a segment can be
   given \emph{additional} emphasis.
\end{em}

It is sometimes necessary to prevent \LaTeX\ from
breaking a line where it might otherwise do so.
This may be at a space, as between the ``Mr.''\ and
``Jones'' in
       ``Mr.~Jones'',        % ~ produces an unbreakable interword space.
or within a word---especially when the word is a
symbol like
       \mbox{\emph{itemnum}}
that makes little sense when hyphenated across
lines.

Footnotes\footnote{This is an example of a footnote.}
pose no problem.

\LaTeX\ is good at typesetting mathematical formulas
like
       \( x-3y + z = 7 \)
or
       \( a_{1} > x^{2n} + y^{2n} > x' \)
or
       \( \ip{A}{B} = \sum_{i} a_{i} b_{i} \).
The spaces you type in a formula are
ignored.  Remember that a letter like
       $x$                   % $ ... $  and  \( ... \)  are equivalent
is a formula when it denotes a mathematical
symbol, and it should be typed as one.

\section{Displayed Text}

Text is displayed by indenting it from the left
margin.  Quotations are commonly displayed.  There
are short quotations
\begin{quote}
   This is a short quotation.  It consists of a
   single paragraph of text.  See how it is formatted.
\end{quote}
and longer ones.
\begin{quotation}
   This is a longer quotation.  It consists of two
   paragraphs of text, neither of which are
   particularly interesting.

   This is the second paragraph of the quotation.  It
   is just as dull as the first paragraph.
\end{quotation}
Another frequently-displayed structure is a list.
The following is an example of an \emph{itemized}
list.
\begin{itemize}
   \item This is the first item of an itemized list.
         Each item in the list is marked with a ``tick''.
         You don't have to worry about what kind of tick
         mark is used.

   \item This is the second item of the list.  It
         contains another list nested inside it.  The inner
         list is an \emph{enumerated} list.
         \begin{enumerate}
            \item This is the first item of an enumerated
                  list that is nested within the itemized list.

            \item This is the second item of the inner list.
                  \LaTeX\ allows you to nest lists deeper than
                  you really should.
         \end{enumerate}
         This is the rest of the second item of the outer
         list.  It is no more interesting than any other
         part of the item.
   \item This is the third item of the list.
\end{itemize}
You can even display poetry.
\begin{verse}
   There is an environment
    for verse \\             % The \\ command separates lines
   Whose features some poets % within a stanza.
   will curse.

                             % One or more blank lines separate stanzas.

   For instead of making\\
   Them do \emph{all} line breaking, \\
   It allows them to put too many words on a line when they'd rather be
   forced to be terse.
\end{verse}

Mathematical formulas may also be displayed.  A
displayed formula
is
one-line long; multiline
formulas require special formatting instructions.
   \[  \ip{\Gamma}{\psi'} = x'' + y^{2} + z_{i}^{n}\]
Don't start a paragraph with a displayed equation,
nor make one a paragraph by itself.

\end{document}               % End of document.

\section{Bench Huge Section 350}
% synthetic bench section 350
% This is a sample LaTeX input file.  (Version of 12 August 2004.)
%
% A '%' character causes TeX to ignore all remaining text on the line,
% and is used for comments like this one.

\documentclass{article}      % Specifies the document class

                             % The preamble begins here.
\title{An Example Document}  % Declares the document's title.
\author{Leslie Lamport}      % Declares the author's name.
\date{January 21, 1994}      % Deleting this command produces today's date.

\newcommand{\ip}[2]{(#1, #2)}
                             % Defines \ip{arg1}{arg2} to mean
                             % (arg1, arg2).

%\newcommand{\ip}[2]{\langle #1 | #2\rangle}
                             % This is an alternative definition of
                             % \ip that is commented out.

\begin{document}             % End of preamble and beginning of text.

\maketitle                   % Produces the title.

This is an example input file.  Comparing it with
the output it generates can show you how to
produce a simple document of your own.

\section{Ordinary Text}      % Produces section heading.  Lower-level
                             % sections are begun with similar
                             % \subsection and \subsubsection commands.

The ends  of words and sentences are marked
  by   spaces. It  doesn't matter how many
spaces    you type; one is as good as 100.  The
end of   a line counts as a space.

One   or more   blank lines denote the  end
of  a paragraph.

Since any number of consecutive spaces are treated
like a single one, the formatting of the input
file makes no difference to
      \LaTeX,                % The \LaTeX command generates the LaTeX logo.
but it makes a difference to you.  When you use
\LaTeX, making your input file as easy to read
as possible will be a great help as you write
your document and when you change it.  This sample
file shows how you can add comments to your own input
file.

Because printing is different from typewriting,
there are a number of things that you have to do
differently when preparing an input file than if
you were just typing the document directly.
Quotation marks like
       ``this''
have to be handled specially, as do quotes within
quotes:
       ``\,`this'            % \, separates the double and single quote.
        is what I just
        wrote, not  `that'\,''.

Dashes come in three sizes: an
       intra-word
dash, a medium dash for number ranges like
       1--2,
and a punctuation
       dash---like
this.

A sentence-ending space should be larger than the
space between words within a sentence.  You
sometimes have to type special commands in
conjunction with punctuation characters to get
this right, as in the following sentence.
       Gnats, gnus, etc.\ all  % `\ ' makes an inter-word space.
       begin with G\@.         % \@ marks end-of-sentence punctuation.
You should check the spaces after periods when
reading your output to make sure you haven't
forgotten any special cases.  Generating an
ellipsis
       \ldots\               % `\ ' is needed after `\ldots' because TeX
                             % ignores spaces after command names like \ldots
                             % made from \ + letters.
                             %
                             % Note how a `%' character causes TeX to ignore
                             % the end of the input line, so these blank lines
                             % do not start a new paragraph.
                             %
with the right spacing around the periods requires
a special command.

\LaTeX\ interprets some common characters as
commands, so you must type special commands to
generate them.  These characters include the
following:
       \$ \& \% \# \{ and \}.

In printing, text is usually emphasized with an
       \emph{italic}
type style.

\begin{em}
   A long segment of text can also be emphasized
   in this way.  Text within such a segment can be
   given \emph{additional} emphasis.
\end{em}

It is sometimes necessary to prevent \LaTeX\ from
breaking a line where it might otherwise do so.
This may be at a space, as between the ``Mr.''\ and
``Jones'' in
       ``Mr.~Jones'',        % ~ produces an unbreakable interword space.
or within a word---especially when the word is a
symbol like
       \mbox{\emph{itemnum}}
that makes little sense when hyphenated across
lines.

Footnotes\footnote{This is an example of a footnote.}
pose no problem.

\LaTeX\ is good at typesetting mathematical formulas
like
       \( x-3y + z = 7 \)
or
       \( a_{1} > x^{2n} + y^{2n} > x' \)
or
       \( \ip{A}{B} = \sum_{i} a_{i} b_{i} \).
The spaces you type in a formula are
ignored.  Remember that a letter like
       $x$                   % $ ... $  and  \( ... \)  are equivalent
is a formula when it denotes a mathematical
symbol, and it should be typed as one.

\section{Displayed Text}

Text is displayed by indenting it from the left
margin.  Quotations are commonly displayed.  There
are short quotations
\begin{quote}
   This is a short quotation.  It consists of a
   single paragraph of text.  See how it is formatted.
\end{quote}
and longer ones.
\begin{quotation}
   This is a longer quotation.  It consists of two
   paragraphs of text, neither of which are
   particularly interesting.

   This is the second paragraph of the quotation.  It
   is just as dull as the first paragraph.
\end{quotation}
Another frequently-displayed structure is a list.
The following is an example of an \emph{itemized}
list.
\begin{itemize}
   \item This is the first item of an itemized list.
         Each item in the list is marked with a ``tick''.
         You don't have to worry about what kind of tick
         mark is used.

   \item This is the second item of the list.  It
         contains another list nested inside it.  The inner
         list is an \emph{enumerated} list.
         \begin{enumerate}
            \item This is the first item of an enumerated
                  list that is nested within the itemized list.

            \item This is the second item of the inner list.
                  \LaTeX\ allows you to nest lists deeper than
                  you really should.
         \end{enumerate}
         This is the rest of the second item of the outer
         list.  It is no more interesting than any other
         part of the item.
   \item This is the third item of the list.
\end{itemize}
You can even display poetry.
\begin{verse}
   There is an environment
    for verse \\             % The \\ command separates lines
   Whose features some poets % within a stanza.
   will curse.

                             % One or more blank lines separate stanzas.

   For instead of making\\
   Them do \emph{all} line breaking, \\
   It allows them to put too many words on a line when they'd rather be
   forced to be terse.
\end{verse}

Mathematical formulas may also be displayed.  A
displayed formula
is
one-line long; multiline
formulas require special formatting instructions.
   \[  \ip{\Gamma}{\psi'} = x'' + y^{2} + z_{i}^{n}\]
Don't start a paragraph with a displayed equation,
nor make one a paragraph by itself.

\end{document}               % End of document.

\section{Bench Huge Section 351}
% synthetic bench section 351
% This is a sample LaTeX input file.  (Version of 12 August 2004.)
%
% A '%' character causes TeX to ignore all remaining text on the line,
% and is used for comments like this one.

\documentclass{article}      % Specifies the document class

                             % The preamble begins here.
\title{An Example Document}  % Declares the document's title.
\author{Leslie Lamport}      % Declares the author's name.
\date{January 21, 1994}      % Deleting this command produces today's date.

\newcommand{\ip}[2]{(#1, #2)}
                             % Defines \ip{arg1}{arg2} to mean
                             % (arg1, arg2).

%\newcommand{\ip}[2]{\langle #1 | #2\rangle}
                             % This is an alternative definition of
                             % \ip that is commented out.

\begin{document}             % End of preamble and beginning of text.

\maketitle                   % Produces the title.

This is an example input file.  Comparing it with
the output it generates can show you how to
produce a simple document of your own.

\section{Ordinary Text}      % Produces section heading.  Lower-level
                             % sections are begun with similar
                             % \subsection and \subsubsection commands.

The ends  of words and sentences are marked
  by   spaces. It  doesn't matter how many
spaces    you type; one is as good as 100.  The
end of   a line counts as a space.

One   or more   blank lines denote the  end
of  a paragraph.

Since any number of consecutive spaces are treated
like a single one, the formatting of the input
file makes no difference to
      \LaTeX,                % The \LaTeX command generates the LaTeX logo.
but it makes a difference to you.  When you use
\LaTeX, making your input file as easy to read
as possible will be a great help as you write
your document and when you change it.  This sample
file shows how you can add comments to your own input
file.

Because printing is different from typewriting,
there are a number of things that you have to do
differently when preparing an input file than if
you were just typing the document directly.
Quotation marks like
       ``this''
have to be handled specially, as do quotes within
quotes:
       ``\,`this'            % \, separates the double and single quote.
        is what I just
        wrote, not  `that'\,''.

Dashes come in three sizes: an
       intra-word
dash, a medium dash for number ranges like
       1--2,
and a punctuation
       dash---like
this.

A sentence-ending space should be larger than the
space between words within a sentence.  You
sometimes have to type special commands in
conjunction with punctuation characters to get
this right, as in the following sentence.
       Gnats, gnus, etc.\ all  % `\ ' makes an inter-word space.
       begin with G\@.         % \@ marks end-of-sentence punctuation.
You should check the spaces after periods when
reading your output to make sure you haven't
forgotten any special cases.  Generating an
ellipsis
       \ldots\               % `\ ' is needed after `\ldots' because TeX
                             % ignores spaces after command names like \ldots
                             % made from \ + letters.
                             %
                             % Note how a `%' character causes TeX to ignore
                             % the end of the input line, so these blank lines
                             % do not start a new paragraph.
                             %
with the right spacing around the periods requires
a special command.

\LaTeX\ interprets some common characters as
commands, so you must type special commands to
generate them.  These characters include the
following:
       \$ \& \% \# \{ and \}.

In printing, text is usually emphasized with an
       \emph{italic}
type style.

\begin{em}
   A long segment of text can also be emphasized
   in this way.  Text within such a segment can be
   given \emph{additional} emphasis.
\end{em}

It is sometimes necessary to prevent \LaTeX\ from
breaking a line where it might otherwise do so.
This may be at a space, as between the ``Mr.''\ and
``Jones'' in
       ``Mr.~Jones'',        % ~ produces an unbreakable interword space.
or within a word---especially when the word is a
symbol like
       \mbox{\emph{itemnum}}
that makes little sense when hyphenated across
lines.

Footnotes\footnote{This is an example of a footnote.}
pose no problem.

\LaTeX\ is good at typesetting mathematical formulas
like
       \( x-3y + z = 7 \)
or
       \( a_{1} > x^{2n} + y^{2n} > x' \)
or
       \( \ip{A}{B} = \sum_{i} a_{i} b_{i} \).
The spaces you type in a formula are
ignored.  Remember that a letter like
       $x$                   % $ ... $  and  \( ... \)  are equivalent
is a formula when it denotes a mathematical
symbol, and it should be typed as one.

\section{Displayed Text}

Text is displayed by indenting it from the left
margin.  Quotations are commonly displayed.  There
are short quotations
\begin{quote}
   This is a short quotation.  It consists of a
   single paragraph of text.  See how it is formatted.
\end{quote}
and longer ones.
\begin{quotation}
   This is a longer quotation.  It consists of two
   paragraphs of text, neither of which are
   particularly interesting.

   This is the second paragraph of the quotation.  It
   is just as dull as the first paragraph.
\end{quotation}
Another frequently-displayed structure is a list.
The following is an example of an \emph{itemized}
list.
\begin{itemize}
   \item This is the first item of an itemized list.
         Each item in the list is marked with a ``tick''.
         You don't have to worry about what kind of tick
         mark is used.

   \item This is the second item of the list.  It
         contains another list nested inside it.  The inner
         list is an \emph{enumerated} list.
         \begin{enumerate}
            \item This is the first item of an enumerated
                  list that is nested within the itemized list.

            \item This is the second item of the inner list.
                  \LaTeX\ allows you to nest lists deeper than
                  you really should.
         \end{enumerate}
         This is the rest of the second item of the outer
         list.  It is no more interesting than any other
         part of the item.
   \item This is the third item of the list.
\end{itemize}
You can even display poetry.
\begin{verse}
   There is an environment
    for verse \\             % The \\ command separates lines
   Whose features some poets % within a stanza.
   will curse.

                             % One or more blank lines separate stanzas.

   For instead of making\\
   Them do \emph{all} line breaking, \\
   It allows them to put too many words on a line when they'd rather be
   forced to be terse.
\end{verse}

Mathematical formulas may also be displayed.  A
displayed formula
is
one-line long; multiline
formulas require special formatting instructions.
   \[  \ip{\Gamma}{\psi'} = x'' + y^{2} + z_{i}^{n}\]
Don't start a paragraph with a displayed equation,
nor make one a paragraph by itself.

\end{document}               % End of document.

\section{Bench Huge Section 352}
% synthetic bench section 352
% This is a sample LaTeX input file.  (Version of 12 August 2004.)
%
% A '%' character causes TeX to ignore all remaining text on the line,
% and is used for comments like this one.

\documentclass{article}      % Specifies the document class

                             % The preamble begins here.
\title{An Example Document}  % Declares the document's title.
\author{Leslie Lamport}      % Declares the author's name.
\date{January 21, 1994}      % Deleting this command produces today's date.

\newcommand{\ip}[2]{(#1, #2)}
                             % Defines \ip{arg1}{arg2} to mean
                             % (arg1, arg2).

%\newcommand{\ip}[2]{\langle #1 | #2\rangle}
                             % This is an alternative definition of
                             % \ip that is commented out.

\begin{document}             % End of preamble and beginning of text.

\maketitle                   % Produces the title.

This is an example input file.  Comparing it with
the output it generates can show you how to
produce a simple document of your own.

\section{Ordinary Text}      % Produces section heading.  Lower-level
                             % sections are begun with similar
                             % \subsection and \subsubsection commands.

The ends  of words and sentences are marked
  by   spaces. It  doesn't matter how many
spaces    you type; one is as good as 100.  The
end of   a line counts as a space.

One   or more   blank lines denote the  end
of  a paragraph.

Since any number of consecutive spaces are treated
like a single one, the formatting of the input
file makes no difference to
      \LaTeX,                % The \LaTeX command generates the LaTeX logo.
but it makes a difference to you.  When you use
\LaTeX, making your input file as easy to read
as possible will be a great help as you write
your document and when you change it.  This sample
file shows how you can add comments to your own input
file.

Because printing is different from typewriting,
there are a number of things that you have to do
differently when preparing an input file than if
you were just typing the document directly.
Quotation marks like
       ``this''
have to be handled specially, as do quotes within
quotes:
       ``\,`this'            % \, separates the double and single quote.
        is what I just
        wrote, not  `that'\,''.

Dashes come in three sizes: an
       intra-word
dash, a medium dash for number ranges like
       1--2,
and a punctuation
       dash---like
this.

A sentence-ending space should be larger than the
space between words within a sentence.  You
sometimes have to type special commands in
conjunction with punctuation characters to get
this right, as in the following sentence.
       Gnats, gnus, etc.\ all  % `\ ' makes an inter-word space.
       begin with G\@.         % \@ marks end-of-sentence punctuation.
You should check the spaces after periods when
reading your output to make sure you haven't
forgotten any special cases.  Generating an
ellipsis
       \ldots\               % `\ ' is needed after `\ldots' because TeX
                             % ignores spaces after command names like \ldots
                             % made from \ + letters.
                             %
                             % Note how a `%' character causes TeX to ignore
                             % the end of the input line, so these blank lines
                             % do not start a new paragraph.
                             %
with the right spacing around the periods requires
a special command.

\LaTeX\ interprets some common characters as
commands, so you must type special commands to
generate them.  These characters include the
following:
       \$ \& \% \# \{ and \}.

In printing, text is usually emphasized with an
       \emph{italic}
type style.

\begin{em}
   A long segment of text can also be emphasized
   in this way.  Text within such a segment can be
   given \emph{additional} emphasis.
\end{em}

It is sometimes necessary to prevent \LaTeX\ from
breaking a line where it might otherwise do so.
This may be at a space, as between the ``Mr.''\ and
``Jones'' in
       ``Mr.~Jones'',        % ~ produces an unbreakable interword space.
or within a word---especially when the word is a
symbol like
       \mbox{\emph{itemnum}}
that makes little sense when hyphenated across
lines.

Footnotes\footnote{This is an example of a footnote.}
pose no problem.

\LaTeX\ is good at typesetting mathematical formulas
like
       \( x-3y + z = 7 \)
or
       \( a_{1} > x^{2n} + y^{2n} > x' \)
or
       \( \ip{A}{B} = \sum_{i} a_{i} b_{i} \).
The spaces you type in a formula are
ignored.  Remember that a letter like
       $x$                   % $ ... $  and  \( ... \)  are equivalent
is a formula when it denotes a mathematical
symbol, and it should be typed as one.

\section{Displayed Text}

Text is displayed by indenting it from the left
margin.  Quotations are commonly displayed.  There
are short quotations
\begin{quote}
   This is a short quotation.  It consists of a
   single paragraph of text.  See how it is formatted.
\end{quote}
and longer ones.
\begin{quotation}
   This is a longer quotation.  It consists of two
   paragraphs of text, neither of which are
   particularly interesting.

   This is the second paragraph of the quotation.  It
   is just as dull as the first paragraph.
\end{quotation}
Another frequently-displayed structure is a list.
The following is an example of an \emph{itemized}
list.
\begin{itemize}
   \item This is the first item of an itemized list.
         Each item in the list is marked with a ``tick''.
         You don't have to worry about what kind of tick
         mark is used.

   \item This is the second item of the list.  It
         contains another list nested inside it.  The inner
         list is an \emph{enumerated} list.
         \begin{enumerate}
            \item This is the first item of an enumerated
                  list that is nested within the itemized list.

            \item This is the second item of the inner list.
                  \LaTeX\ allows you to nest lists deeper than
                  you really should.
         \end{enumerate}
         This is the rest of the second item of the outer
         list.  It is no more interesting than any other
         part of the item.
   \item This is the third item of the list.
\end{itemize}
You can even display poetry.
\begin{verse}
   There is an environment
    for verse \\             % The \\ command separates lines
   Whose features some poets % within a stanza.
   will curse.

                             % One or more blank lines separate stanzas.

   For instead of making\\
   Them do \emph{all} line breaking, \\
   It allows them to put too many words on a line when they'd rather be
   forced to be terse.
\end{verse}

Mathematical formulas may also be displayed.  A
displayed formula
is
one-line long; multiline
formulas require special formatting instructions.
   \[  \ip{\Gamma}{\psi'} = x'' + y^{2} + z_{i}^{n}\]
Don't start a paragraph with a displayed equation,
nor make one a paragraph by itself.

\end{document}               % End of document.

\section{Bench Huge Section 353}
% synthetic bench section 353
% This is a sample LaTeX input file.  (Version of 12 August 2004.)
%
% A '%' character causes TeX to ignore all remaining text on the line,
% and is used for comments like this one.

\documentclass{article}      % Specifies the document class

                             % The preamble begins here.
\title{An Example Document}  % Declares the document's title.
\author{Leslie Lamport}      % Declares the author's name.
\date{January 21, 1994}      % Deleting this command produces today's date.

\newcommand{\ip}[2]{(#1, #2)}
                             % Defines \ip{arg1}{arg2} to mean
                             % (arg1, arg2).

%\newcommand{\ip}[2]{\langle #1 | #2\rangle}
                             % This is an alternative definition of
                             % \ip that is commented out.

\begin{document}             % End of preamble and beginning of text.

\maketitle                   % Produces the title.

This is an example input file.  Comparing it with
the output it generates can show you how to
produce a simple document of your own.

\section{Ordinary Text}      % Produces section heading.  Lower-level
                             % sections are begun with similar
                             % \subsection and \subsubsection commands.

The ends  of words and sentences are marked
  by   spaces. It  doesn't matter how many
spaces    you type; one is as good as 100.  The
end of   a line counts as a space.

One   or more   blank lines denote the  end
of  a paragraph.

Since any number of consecutive spaces are treated
like a single one, the formatting of the input
file makes no difference to
      \LaTeX,                % The \LaTeX command generates the LaTeX logo.
but it makes a difference to you.  When you use
\LaTeX, making your input file as easy to read
as possible will be a great help as you write
your document and when you change it.  This sample
file shows how you can add comments to your own input
file.

Because printing is different from typewriting,
there are a number of things that you have to do
differently when preparing an input file than if
you were just typing the document directly.
Quotation marks like
       ``this''
have to be handled specially, as do quotes within
quotes:
       ``\,`this'            % \, separates the double and single quote.
        is what I just
        wrote, not  `that'\,''.

Dashes come in three sizes: an
       intra-word
dash, a medium dash for number ranges like
       1--2,
and a punctuation
       dash---like
this.

A sentence-ending space should be larger than the
space between words within a sentence.  You
sometimes have to type special commands in
conjunction with punctuation characters to get
this right, as in the following sentence.
       Gnats, gnus, etc.\ all  % `\ ' makes an inter-word space.
       begin with G\@.         % \@ marks end-of-sentence punctuation.
You should check the spaces after periods when
reading your output to make sure you haven't
forgotten any special cases.  Generating an
ellipsis
       \ldots\               % `\ ' is needed after `\ldots' because TeX
                             % ignores spaces after command names like \ldots
                             % made from \ + letters.
                             %
                             % Note how a `%' character causes TeX to ignore
                             % the end of the input line, so these blank lines
                             % do not start a new paragraph.
                             %
with the right spacing around the periods requires
a special command.

\LaTeX\ interprets some common characters as
commands, so you must type special commands to
generate them.  These characters include the
following:
       \$ \& \% \# \{ and \}.

In printing, text is usually emphasized with an
       \emph{italic}
type style.

\begin{em}
   A long segment of text can also be emphasized
   in this way.  Text within such a segment can be
   given \emph{additional} emphasis.
\end{em}

It is sometimes necessary to prevent \LaTeX\ from
breaking a line where it might otherwise do so.
This may be at a space, as between the ``Mr.''\ and
``Jones'' in
       ``Mr.~Jones'',        % ~ produces an unbreakable interword space.
or within a word---especially when the word is a
symbol like
       \mbox{\emph{itemnum}}
that makes little sense when hyphenated across
lines.

Footnotes\footnote{This is an example of a footnote.}
pose no problem.

\LaTeX\ is good at typesetting mathematical formulas
like
       \( x-3y + z = 7 \)
or
       \( a_{1} > x^{2n} + y^{2n} > x' \)
or
       \( \ip{A}{B} = \sum_{i} a_{i} b_{i} \).
The spaces you type in a formula are
ignored.  Remember that a letter like
       $x$                   % $ ... $  and  \( ... \)  are equivalent
is a formula when it denotes a mathematical
symbol, and it should be typed as one.

\section{Displayed Text}

Text is displayed by indenting it from the left
margin.  Quotations are commonly displayed.  There
are short quotations
\begin{quote}
   This is a short quotation.  It consists of a
   single paragraph of text.  See how it is formatted.
\end{quote}
and longer ones.
\begin{quotation}
   This is a longer quotation.  It consists of two
   paragraphs of text, neither of which are
   particularly interesting.

   This is the second paragraph of the quotation.  It
   is just as dull as the first paragraph.
\end{quotation}
Another frequently-displayed structure is a list.
The following is an example of an \emph{itemized}
list.
\begin{itemize}
   \item This is the first item of an itemized list.
         Each item in the list is marked with a ``tick''.
         You don't have to worry about what kind of tick
         mark is used.

   \item This is the second item of the list.  It
         contains another list nested inside it.  The inner
         list is an \emph{enumerated} list.
         \begin{enumerate}
            \item This is the first item of an enumerated
                  list that is nested within the itemized list.

            \item This is the second item of the inner list.
                  \LaTeX\ allows you to nest lists deeper than
                  you really should.
         \end{enumerate}
         This is the rest of the second item of the outer
         list.  It is no more interesting than any other
         part of the item.
   \item This is the third item of the list.
\end{itemize}
You can even display poetry.
\begin{verse}
   There is an environment
    for verse \\             % The \\ command separates lines
   Whose features some poets % within a stanza.
   will curse.

                             % One or more blank lines separate stanzas.

   For instead of making\\
   Them do \emph{all} line breaking, \\
   It allows them to put too many words on a line when they'd rather be
   forced to be terse.
\end{verse}

Mathematical formulas may also be displayed.  A
displayed formula
is
one-line long; multiline
formulas require special formatting instructions.
   \[  \ip{\Gamma}{\psi'} = x'' + y^{2} + z_{i}^{n}\]
Don't start a paragraph with a displayed equation,
nor make one a paragraph by itself.

\end{document}               % End of document.

\section{Bench Huge Section 354}
% synthetic bench section 354
% This is a sample LaTeX input file.  (Version of 12 August 2004.)
%
% A '%' character causes TeX to ignore all remaining text on the line,
% and is used for comments like this one.

\documentclass{article}      % Specifies the document class

                             % The preamble begins here.
\title{An Example Document}  % Declares the document's title.
\author{Leslie Lamport}      % Declares the author's name.
\date{January 21, 1994}      % Deleting this command produces today's date.

\newcommand{\ip}[2]{(#1, #2)}
                             % Defines \ip{arg1}{arg2} to mean
                             % (arg1, arg2).

%\newcommand{\ip}[2]{\langle #1 | #2\rangle}
                             % This is an alternative definition of
                             % \ip that is commented out.

\begin{document}             % End of preamble and beginning of text.

\maketitle                   % Produces the title.

This is an example input file.  Comparing it with
the output it generates can show you how to
produce a simple document of your own.

\section{Ordinary Text}      % Produces section heading.  Lower-level
                             % sections are begun with similar
                             % \subsection and \subsubsection commands.

The ends  of words and sentences are marked
  by   spaces. It  doesn't matter how many
spaces    you type; one is as good as 100.  The
end of   a line counts as a space.

One   or more   blank lines denote the  end
of  a paragraph.

Since any number of consecutive spaces are treated
like a single one, the formatting of the input
file makes no difference to
      \LaTeX,                % The \LaTeX command generates the LaTeX logo.
but it makes a difference to you.  When you use
\LaTeX, making your input file as easy to read
as possible will be a great help as you write
your document and when you change it.  This sample
file shows how you can add comments to your own input
file.

Because printing is different from typewriting,
there are a number of things that you have to do
differently when preparing an input file than if
you were just typing the document directly.
Quotation marks like
       ``this''
have to be handled specially, as do quotes within
quotes:
       ``\,`this'            % \, separates the double and single quote.
        is what I just
        wrote, not  `that'\,''.

Dashes come in three sizes: an
       intra-word
dash, a medium dash for number ranges like
       1--2,
and a punctuation
       dash---like
this.

A sentence-ending space should be larger than the
space between words within a sentence.  You
sometimes have to type special commands in
conjunction with punctuation characters to get
this right, as in the following sentence.
       Gnats, gnus, etc.\ all  % `\ ' makes an inter-word space.
       begin with G\@.         % \@ marks end-of-sentence punctuation.
You should check the spaces after periods when
reading your output to make sure you haven't
forgotten any special cases.  Generating an
ellipsis
       \ldots\               % `\ ' is needed after `\ldots' because TeX
                             % ignores spaces after command names like \ldots
                             % made from \ + letters.
                             %
                             % Note how a `%' character causes TeX to ignore
                             % the end of the input line, so these blank lines
                             % do not start a new paragraph.
                             %
with the right spacing around the periods requires
a special command.

\LaTeX\ interprets some common characters as
commands, so you must type special commands to
generate them.  These characters include the
following:
       \$ \& \% \# \{ and \}.

In printing, text is usually emphasized with an
       \emph{italic}
type style.

\begin{em}
   A long segment of text can also be emphasized
   in this way.  Text within such a segment can be
   given \emph{additional} emphasis.
\end{em}

It is sometimes necessary to prevent \LaTeX\ from
breaking a line where it might otherwise do so.
This may be at a space, as between the ``Mr.''\ and
``Jones'' in
       ``Mr.~Jones'',        % ~ produces an unbreakable interword space.
or within a word---especially when the word is a
symbol like
       \mbox{\emph{itemnum}}
that makes little sense when hyphenated across
lines.

Footnotes\footnote{This is an example of a footnote.}
pose no problem.

\LaTeX\ is good at typesetting mathematical formulas
like
       \( x-3y + z = 7 \)
or
       \( a_{1} > x^{2n} + y^{2n} > x' \)
or
       \( \ip{A}{B} = \sum_{i} a_{i} b_{i} \).
The spaces you type in a formula are
ignored.  Remember that a letter like
       $x$                   % $ ... $  and  \( ... \)  are equivalent
is a formula when it denotes a mathematical
symbol, and it should be typed as one.

\section{Displayed Text}

Text is displayed by indenting it from the left
margin.  Quotations are commonly displayed.  There
are short quotations
\begin{quote}
   This is a short quotation.  It consists of a
   single paragraph of text.  See how it is formatted.
\end{quote}
and longer ones.
\begin{quotation}
   This is a longer quotation.  It consists of two
   paragraphs of text, neither of which are
   particularly interesting.

   This is the second paragraph of the quotation.  It
   is just as dull as the first paragraph.
\end{quotation}
Another frequently-displayed structure is a list.
The following is an example of an \emph{itemized}
list.
\begin{itemize}
   \item This is the first item of an itemized list.
         Each item in the list is marked with a ``tick''.
         You don't have to worry about what kind of tick
         mark is used.

   \item This is the second item of the list.  It
         contains another list nested inside it.  The inner
         list is an \emph{enumerated} list.
         \begin{enumerate}
            \item This is the first item of an enumerated
                  list that is nested within the itemized list.

            \item This is the second item of the inner list.
                  \LaTeX\ allows you to nest lists deeper than
                  you really should.
         \end{enumerate}
         This is the rest of the second item of the outer
         list.  It is no more interesting than any other
         part of the item.
   \item This is the third item of the list.
\end{itemize}
You can even display poetry.
\begin{verse}
   There is an environment
    for verse \\             % The \\ command separates lines
   Whose features some poets % within a stanza.
   will curse.

                             % One or more blank lines separate stanzas.

   For instead of making\\
   Them do \emph{all} line breaking, \\
   It allows them to put too many words on a line when they'd rather be
   forced to be terse.
\end{verse}

Mathematical formulas may also be displayed.  A
displayed formula
is
one-line long; multiline
formulas require special formatting instructions.
   \[  \ip{\Gamma}{\psi'} = x'' + y^{2} + z_{i}^{n}\]
Don't start a paragraph with a displayed equation,
nor make one a paragraph by itself.

\end{document}               % End of document.

\section{Bench Huge Section 355}
% synthetic bench section 355
% This is a sample LaTeX input file.  (Version of 12 August 2004.)
%
% A '%' character causes TeX to ignore all remaining text on the line,
% and is used for comments like this one.

\documentclass{article}      % Specifies the document class

                             % The preamble begins here.
\title{An Example Document}  % Declares the document's title.
\author{Leslie Lamport}      % Declares the author's name.
\date{January 21, 1994}      % Deleting this command produces today's date.

\newcommand{\ip}[2]{(#1, #2)}
                             % Defines \ip{arg1}{arg2} to mean
                             % (arg1, arg2).

%\newcommand{\ip}[2]{\langle #1 | #2\rangle}
                             % This is an alternative definition of
                             % \ip that is commented out.

\begin{document}             % End of preamble and beginning of text.

\maketitle                   % Produces the title.

This is an example input file.  Comparing it with
the output it generates can show you how to
produce a simple document of your own.

\section{Ordinary Text}      % Produces section heading.  Lower-level
                             % sections are begun with similar
                             % \subsection and \subsubsection commands.

The ends  of words and sentences are marked
  by   spaces. It  doesn't matter how many
spaces    you type; one is as good as 100.  The
end of   a line counts as a space.

One   or more   blank lines denote the  end
of  a paragraph.

Since any number of consecutive spaces are treated
like a single one, the formatting of the input
file makes no difference to
      \LaTeX,                % The \LaTeX command generates the LaTeX logo.
but it makes a difference to you.  When you use
\LaTeX, making your input file as easy to read
as possible will be a great help as you write
your document and when you change it.  This sample
file shows how you can add comments to your own input
file.

Because printing is different from typewriting,
there are a number of things that you have to do
differently when preparing an input file than if
you were just typing the document directly.
Quotation marks like
       ``this''
have to be handled specially, as do quotes within
quotes:
       ``\,`this'            % \, separates the double and single quote.
        is what I just
        wrote, not  `that'\,''.

Dashes come in three sizes: an
       intra-word
dash, a medium dash for number ranges like
       1--2,
and a punctuation
       dash---like
this.

A sentence-ending space should be larger than the
space between words within a sentence.  You
sometimes have to type special commands in
conjunction with punctuation characters to get
this right, as in the following sentence.
       Gnats, gnus, etc.\ all  % `\ ' makes an inter-word space.
       begin with G\@.         % \@ marks end-of-sentence punctuation.
You should check the spaces after periods when
reading your output to make sure you haven't
forgotten any special cases.  Generating an
ellipsis
       \ldots\               % `\ ' is needed after `\ldots' because TeX
                             % ignores spaces after command names like \ldots
                             % made from \ + letters.
                             %
                             % Note how a `%' character causes TeX to ignore
                             % the end of the input line, so these blank lines
                             % do not start a new paragraph.
                             %
with the right spacing around the periods requires
a special command.

\LaTeX\ interprets some common characters as
commands, so you must type special commands to
generate them.  These characters include the
following:
       \$ \& \% \# \{ and \}.

In printing, text is usually emphasized with an
       \emph{italic}
type style.

\begin{em}
   A long segment of text can also be emphasized
   in this way.  Text within such a segment can be
   given \emph{additional} emphasis.
\end{em}

It is sometimes necessary to prevent \LaTeX\ from
breaking a line where it might otherwise do so.
This may be at a space, as between the ``Mr.''\ and
``Jones'' in
       ``Mr.~Jones'',        % ~ produces an unbreakable interword space.
or within a word---especially when the word is a
symbol like
       \mbox{\emph{itemnum}}
that makes little sense when hyphenated across
lines.

Footnotes\footnote{This is an example of a footnote.}
pose no problem.

\LaTeX\ is good at typesetting mathematical formulas
like
       \( x-3y + z = 7 \)
or
       \( a_{1} > x^{2n} + y^{2n} > x' \)
or
       \( \ip{A}{B} = \sum_{i} a_{i} b_{i} \).
The spaces you type in a formula are
ignored.  Remember that a letter like
       $x$                   % $ ... $  and  \( ... \)  are equivalent
is a formula when it denotes a mathematical
symbol, and it should be typed as one.

\section{Displayed Text}

Text is displayed by indenting it from the left
margin.  Quotations are commonly displayed.  There
are short quotations
\begin{quote}
   This is a short quotation.  It consists of a
   single paragraph of text.  See how it is formatted.
\end{quote}
and longer ones.
\begin{quotation}
   This is a longer quotation.  It consists of two
   paragraphs of text, neither of which are
   particularly interesting.

   This is the second paragraph of the quotation.  It
   is just as dull as the first paragraph.
\end{quotation}
Another frequently-displayed structure is a list.
The following is an example of an \emph{itemized}
list.
\begin{itemize}
   \item This is the first item of an itemized list.
         Each item in the list is marked with a ``tick''.
         You don't have to worry about what kind of tick
         mark is used.

   \item This is the second item of the list.  It
         contains another list nested inside it.  The inner
         list is an \emph{enumerated} list.
         \begin{enumerate}
            \item This is the first item of an enumerated
                  list that is nested within the itemized list.

            \item This is the second item of the inner list.
                  \LaTeX\ allows you to nest lists deeper than
                  you really should.
         \end{enumerate}
         This is the rest of the second item of the outer
         list.  It is no more interesting than any other
         part of the item.
   \item This is the third item of the list.
\end{itemize}
You can even display poetry.
\begin{verse}
   There is an environment
    for verse \\             % The \\ command separates lines
   Whose features some poets % within a stanza.
   will curse.

                             % One or more blank lines separate stanzas.

   For instead of making\\
   Them do \emph{all} line breaking, \\
   It allows them to put too many words on a line when they'd rather be
   forced to be terse.
\end{verse}

Mathematical formulas may also be displayed.  A
displayed formula
is
one-line long; multiline
formulas require special formatting instructions.
   \[  \ip{\Gamma}{\psi'} = x'' + y^{2} + z_{i}^{n}\]
Don't start a paragraph with a displayed equation,
nor make one a paragraph by itself.

\end{document}               % End of document.

\section{Bench Huge Section 356}
% synthetic bench section 356
% This is a sample LaTeX input file.  (Version of 12 August 2004.)
%
% A '%' character causes TeX to ignore all remaining text on the line,
% and is used for comments like this one.

\documentclass{article}      % Specifies the document class

                             % The preamble begins here.
\title{An Example Document}  % Declares the document's title.
\author{Leslie Lamport}      % Declares the author's name.
\date{January 21, 1994}      % Deleting this command produces today's date.

\newcommand{\ip}[2]{(#1, #2)}
                             % Defines \ip{arg1}{arg2} to mean
                             % (arg1, arg2).

%\newcommand{\ip}[2]{\langle #1 | #2\rangle}
                             % This is an alternative definition of
                             % \ip that is commented out.

\begin{document}             % End of preamble and beginning of text.

\maketitle                   % Produces the title.

This is an example input file.  Comparing it with
the output it generates can show you how to
produce a simple document of your own.

\section{Ordinary Text}      % Produces section heading.  Lower-level
                             % sections are begun with similar
                             % \subsection and \subsubsection commands.

The ends  of words and sentences are marked
  by   spaces. It  doesn't matter how many
spaces    you type; one is as good as 100.  The
end of   a line counts as a space.

One   or more   blank lines denote the  end
of  a paragraph.

Since any number of consecutive spaces are treated
like a single one, the formatting of the input
file makes no difference to
      \LaTeX,                % The \LaTeX command generates the LaTeX logo.
but it makes a difference to you.  When you use
\LaTeX, making your input file as easy to read
as possible will be a great help as you write
your document and when you change it.  This sample
file shows how you can add comments to your own input
file.

Because printing is different from typewriting,
there are a number of things that you have to do
differently when preparing an input file than if
you were just typing the document directly.
Quotation marks like
       ``this''
have to be handled specially, as do quotes within
quotes:
       ``\,`this'            % \, separates the double and single quote.
        is what I just
        wrote, not  `that'\,''.

Dashes come in three sizes: an
       intra-word
dash, a medium dash for number ranges like
       1--2,
and a punctuation
       dash---like
this.

A sentence-ending space should be larger than the
space between words within a sentence.  You
sometimes have to type special commands in
conjunction with punctuation characters to get
this right, as in the following sentence.
       Gnats, gnus, etc.\ all  % `\ ' makes an inter-word space.
       begin with G\@.         % \@ marks end-of-sentence punctuation.
You should check the spaces after periods when
reading your output to make sure you haven't
forgotten any special cases.  Generating an
ellipsis
       \ldots\               % `\ ' is needed after `\ldots' because TeX
                             % ignores spaces after command names like \ldots
                             % made from \ + letters.
                             %
                             % Note how a `%' character causes TeX to ignore
                             % the end of the input line, so these blank lines
                             % do not start a new paragraph.
                             %
with the right spacing around the periods requires
a special command.

\LaTeX\ interprets some common characters as
commands, so you must type special commands to
generate them.  These characters include the
following:
       \$ \& \% \# \{ and \}.

In printing, text is usually emphasized with an
       \emph{italic}
type style.

\begin{em}
   A long segment of text can also be emphasized
   in this way.  Text within such a segment can be
   given \emph{additional} emphasis.
\end{em}

It is sometimes necessary to prevent \LaTeX\ from
breaking a line where it might otherwise do so.
This may be at a space, as between the ``Mr.''\ and
``Jones'' in
       ``Mr.~Jones'',        % ~ produces an unbreakable interword space.
or within a word---especially when the word is a
symbol like
       \mbox{\emph{itemnum}}
that makes little sense when hyphenated across
lines.

Footnotes\footnote{This is an example of a footnote.}
pose no problem.

\LaTeX\ is good at typesetting mathematical formulas
like
       \( x-3y + z = 7 \)
or
       \( a_{1} > x^{2n} + y^{2n} > x' \)
or
       \( \ip{A}{B} = \sum_{i} a_{i} b_{i} \).
The spaces you type in a formula are
ignored.  Remember that a letter like
       $x$                   % $ ... $  and  \( ... \)  are equivalent
is a formula when it denotes a mathematical
symbol, and it should be typed as one.

\section{Displayed Text}

Text is displayed by indenting it from the left
margin.  Quotations are commonly displayed.  There
are short quotations
\begin{quote}
   This is a short quotation.  It consists of a
   single paragraph of text.  See how it is formatted.
\end{quote}
and longer ones.
\begin{quotation}
   This is a longer quotation.  It consists of two
   paragraphs of text, neither of which are
   particularly interesting.

   This is the second paragraph of the quotation.  It
   is just as dull as the first paragraph.
\end{quotation}
Another frequently-displayed structure is a list.
The following is an example of an \emph{itemized}
list.
\begin{itemize}
   \item This is the first item of an itemized list.
         Each item in the list is marked with a ``tick''.
         You don't have to worry about what kind of tick
         mark is used.

   \item This is the second item of the list.  It
         contains another list nested inside it.  The inner
         list is an \emph{enumerated} list.
         \begin{enumerate}
            \item This is the first item of an enumerated
                  list that is nested within the itemized list.

            \item This is the second item of the inner list.
                  \LaTeX\ allows you to nest lists deeper than
                  you really should.
         \end{enumerate}
         This is the rest of the second item of the outer
         list.  It is no more interesting than any other
         part of the item.
   \item This is the third item of the list.
\end{itemize}
You can even display poetry.
\begin{verse}
   There is an environment
    for verse \\             % The \\ command separates lines
   Whose features some poets % within a stanza.
   will curse.

                             % One or more blank lines separate stanzas.

   For instead of making\\
   Them do \emph{all} line breaking, \\
   It allows them to put too many words on a line when they'd rather be
   forced to be terse.
\end{verse}

Mathematical formulas may also be displayed.  A
displayed formula
is
one-line long; multiline
formulas require special formatting instructions.
   \[  \ip{\Gamma}{\psi'} = x'' + y^{2} + z_{i}^{n}\]
Don't start a paragraph with a displayed equation,
nor make one a paragraph by itself.

\end{document}               % End of document.

\section{Bench Huge Section 357}
% synthetic bench section 357
% This is a sample LaTeX input file.  (Version of 12 August 2004.)
%
% A '%' character causes TeX to ignore all remaining text on the line,
% and is used for comments like this one.

\documentclass{article}      % Specifies the document class

                             % The preamble begins here.
\title{An Example Document}  % Declares the document's title.
\author{Leslie Lamport}      % Declares the author's name.
\date{January 21, 1994}      % Deleting this command produces today's date.

\newcommand{\ip}[2]{(#1, #2)}
                             % Defines \ip{arg1}{arg2} to mean
                             % (arg1, arg2).

%\newcommand{\ip}[2]{\langle #1 | #2\rangle}
                             % This is an alternative definition of
                             % \ip that is commented out.

\begin{document}             % End of preamble and beginning of text.

\maketitle                   % Produces the title.

This is an example input file.  Comparing it with
the output it generates can show you how to
produce a simple document of your own.

\section{Ordinary Text}      % Produces section heading.  Lower-level
                             % sections are begun with similar
                             % \subsection and \subsubsection commands.

The ends  of words and sentences are marked
  by   spaces. It  doesn't matter how many
spaces    you type; one is as good as 100.  The
end of   a line counts as a space.

One   or more   blank lines denote the  end
of  a paragraph.

Since any number of consecutive spaces are treated
like a single one, the formatting of the input
file makes no difference to
      \LaTeX,                % The \LaTeX command generates the LaTeX logo.
but it makes a difference to you.  When you use
\LaTeX, making your input file as easy to read
as possible will be a great help as you write
your document and when you change it.  This sample
file shows how you can add comments to your own input
file.

Because printing is different from typewriting,
there are a number of things that you have to do
differently when preparing an input file than if
you were just typing the document directly.
Quotation marks like
       ``this''
have to be handled specially, as do quotes within
quotes:
       ``\,`this'            % \, separates the double and single quote.
        is what I just
        wrote, not  `that'\,''.

Dashes come in three sizes: an
       intra-word
dash, a medium dash for number ranges like
       1--2,
and a punctuation
       dash---like
this.

A sentence-ending space should be larger than the
space between words within a sentence.  You
sometimes have to type special commands in
conjunction with punctuation characters to get
this right, as in the following sentence.
       Gnats, gnus, etc.\ all  % `\ ' makes an inter-word space.
       begin with G\@.         % \@ marks end-of-sentence punctuation.
You should check the spaces after periods when
reading your output to make sure you haven't
forgotten any special cases.  Generating an
ellipsis
       \ldots\               % `\ ' is needed after `\ldots' because TeX
                             % ignores spaces after command names like \ldots
                             % made from \ + letters.
                             %
                             % Note how a `%' character causes TeX to ignore
                             % the end of the input line, so these blank lines
                             % do not start a new paragraph.
                             %
with the right spacing around the periods requires
a special command.

\LaTeX\ interprets some common characters as
commands, so you must type special commands to
generate them.  These characters include the
following:
       \$ \& \% \# \{ and \}.

In printing, text is usually emphasized with an
       \emph{italic}
type style.

\begin{em}
   A long segment of text can also be emphasized
   in this way.  Text within such a segment can be
   given \emph{additional} emphasis.
\end{em}

It is sometimes necessary to prevent \LaTeX\ from
breaking a line where it might otherwise do so.
This may be at a space, as between the ``Mr.''\ and
``Jones'' in
       ``Mr.~Jones'',        % ~ produces an unbreakable interword space.
or within a word---especially when the word is a
symbol like
       \mbox{\emph{itemnum}}
that makes little sense when hyphenated across
lines.

Footnotes\footnote{This is an example of a footnote.}
pose no problem.

\LaTeX\ is good at typesetting mathematical formulas
like
       \( x-3y + z = 7 \)
or
       \( a_{1} > x^{2n} + y^{2n} > x' \)
or
       \( \ip{A}{B} = \sum_{i} a_{i} b_{i} \).
The spaces you type in a formula are
ignored.  Remember that a letter like
       $x$                   % $ ... $  and  \( ... \)  are equivalent
is a formula when it denotes a mathematical
symbol, and it should be typed as one.

\section{Displayed Text}

Text is displayed by indenting it from the left
margin.  Quotations are commonly displayed.  There
are short quotations
\begin{quote}
   This is a short quotation.  It consists of a
   single paragraph of text.  See how it is formatted.
\end{quote}
and longer ones.
\begin{quotation}
   This is a longer quotation.  It consists of two
   paragraphs of text, neither of which are
   particularly interesting.

   This is the second paragraph of the quotation.  It
   is just as dull as the first paragraph.
\end{quotation}
Another frequently-displayed structure is a list.
The following is an example of an \emph{itemized}
list.
\begin{itemize}
   \item This is the first item of an itemized list.
         Each item in the list is marked with a ``tick''.
         You don't have to worry about what kind of tick
         mark is used.

   \item This is the second item of the list.  It
         contains another list nested inside it.  The inner
         list is an \emph{enumerated} list.
         \begin{enumerate}
            \item This is the first item of an enumerated
                  list that is nested within the itemized list.

            \item This is the second item of the inner list.
                  \LaTeX\ allows you to nest lists deeper than
                  you really should.
         \end{enumerate}
         This is the rest of the second item of the outer
         list.  It is no more interesting than any other
         part of the item.
   \item This is the third item of the list.
\end{itemize}
You can even display poetry.
\begin{verse}
   There is an environment
    for verse \\             % The \\ command separates lines
   Whose features some poets % within a stanza.
   will curse.

                             % One or more blank lines separate stanzas.

   For instead of making\\
   Them do \emph{all} line breaking, \\
   It allows them to put too many words on a line when they'd rather be
   forced to be terse.
\end{verse}

Mathematical formulas may also be displayed.  A
displayed formula
is
one-line long; multiline
formulas require special formatting instructions.
   \[  \ip{\Gamma}{\psi'} = x'' + y^{2} + z_{i}^{n}\]
Don't start a paragraph with a displayed equation,
nor make one a paragraph by itself.

\end{document}               % End of document.

\section{Bench Huge Section 358}
% synthetic bench section 358
% This is a sample LaTeX input file.  (Version of 12 August 2004.)
%
% A '%' character causes TeX to ignore all remaining text on the line,
% and is used for comments like this one.

\documentclass{article}      % Specifies the document class

                             % The preamble begins here.
\title{An Example Document}  % Declares the document's title.
\author{Leslie Lamport}      % Declares the author's name.
\date{January 21, 1994}      % Deleting this command produces today's date.

\newcommand{\ip}[2]{(#1, #2)}
                             % Defines \ip{arg1}{arg2} to mean
                             % (arg1, arg2).

%\newcommand{\ip}[2]{\langle #1 | #2\rangle}
                             % This is an alternative definition of
                             % \ip that is commented out.

\begin{document}             % End of preamble and beginning of text.

\maketitle                   % Produces the title.

This is an example input file.  Comparing it with
the output it generates can show you how to
produce a simple document of your own.

\section{Ordinary Text}      % Produces section heading.  Lower-level
                             % sections are begun with similar
                             % \subsection and \subsubsection commands.

The ends  of words and sentences are marked
  by   spaces. It  doesn't matter how many
spaces    you type; one is as good as 100.  The
end of   a line counts as a space.

One   or more   blank lines denote the  end
of  a paragraph.

Since any number of consecutive spaces are treated
like a single one, the formatting of the input
file makes no difference to
      \LaTeX,                % The \LaTeX command generates the LaTeX logo.
but it makes a difference to you.  When you use
\LaTeX, making your input file as easy to read
as possible will be a great help as you write
your document and when you change it.  This sample
file shows how you can add comments to your own input
file.

Because printing is different from typewriting,
there are a number of things that you have to do
differently when preparing an input file than if
you were just typing the document directly.
Quotation marks like
       ``this''
have to be handled specially, as do quotes within
quotes:
       ``\,`this'            % \, separates the double and single quote.
        is what I just
        wrote, not  `that'\,''.

Dashes come in three sizes: an
       intra-word
dash, a medium dash for number ranges like
       1--2,
and a punctuation
       dash---like
this.

A sentence-ending space should be larger than the
space between words within a sentence.  You
sometimes have to type special commands in
conjunction with punctuation characters to get
this right, as in the following sentence.
       Gnats, gnus, etc.\ all  % `\ ' makes an inter-word space.
       begin with G\@.         % \@ marks end-of-sentence punctuation.
You should check the spaces after periods when
reading your output to make sure you haven't
forgotten any special cases.  Generating an
ellipsis
       \ldots\               % `\ ' is needed after `\ldots' because TeX
                             % ignores spaces after command names like \ldots
                             % made from \ + letters.
                             %
                             % Note how a `%' character causes TeX to ignore
                             % the end of the input line, so these blank lines
                             % do not start a new paragraph.
                             %
with the right spacing around the periods requires
a special command.

\LaTeX\ interprets some common characters as
commands, so you must type special commands to
generate them.  These characters include the
following:
       \$ \& \% \# \{ and \}.

In printing, text is usually emphasized with an
       \emph{italic}
type style.

\begin{em}
   A long segment of text can also be emphasized
   in this way.  Text within such a segment can be
   given \emph{additional} emphasis.
\end{em}

It is sometimes necessary to prevent \LaTeX\ from
breaking a line where it might otherwise do so.
This may be at a space, as between the ``Mr.''\ and
``Jones'' in
       ``Mr.~Jones'',        % ~ produces an unbreakable interword space.
or within a word---especially when the word is a
symbol like
       \mbox{\emph{itemnum}}
that makes little sense when hyphenated across
lines.

Footnotes\footnote{This is an example of a footnote.}
pose no problem.

\LaTeX\ is good at typesetting mathematical formulas
like
       \( x-3y + z = 7 \)
or
       \( a_{1} > x^{2n} + y^{2n} > x' \)
or
       \( \ip{A}{B} = \sum_{i} a_{i} b_{i} \).
The spaces you type in a formula are
ignored.  Remember that a letter like
       $x$                   % $ ... $  and  \( ... \)  are equivalent
is a formula when it denotes a mathematical
symbol, and it should be typed as one.

\section{Displayed Text}

Text is displayed by indenting it from the left
margin.  Quotations are commonly displayed.  There
are short quotations
\begin{quote}
   This is a short quotation.  It consists of a
   single paragraph of text.  See how it is formatted.
\end{quote}
and longer ones.
\begin{quotation}
   This is a longer quotation.  It consists of two
   paragraphs of text, neither of which are
   particularly interesting.

   This is the second paragraph of the quotation.  It
   is just as dull as the first paragraph.
\end{quotation}
Another frequently-displayed structure is a list.
The following is an example of an \emph{itemized}
list.
\begin{itemize}
   \item This is the first item of an itemized list.
         Each item in the list is marked with a ``tick''.
         You don't have to worry about what kind of tick
         mark is used.

   \item This is the second item of the list.  It
         contains another list nested inside it.  The inner
         list is an \emph{enumerated} list.
         \begin{enumerate}
            \item This is the first item of an enumerated
                  list that is nested within the itemized list.

            \item This is the second item of the inner list.
                  \LaTeX\ allows you to nest lists deeper than
                  you really should.
         \end{enumerate}
         This is the rest of the second item of the outer
         list.  It is no more interesting than any other
         part of the item.
   \item This is the third item of the list.
\end{itemize}
You can even display poetry.
\begin{verse}
   There is an environment
    for verse \\             % The \\ command separates lines
   Whose features some poets % within a stanza.
   will curse.

                             % One or more blank lines separate stanzas.

   For instead of making\\
   Them do \emph{all} line breaking, \\
   It allows them to put too many words on a line when they'd rather be
   forced to be terse.
\end{verse}

Mathematical formulas may also be displayed.  A
displayed formula
is
one-line long; multiline
formulas require special formatting instructions.
   \[  \ip{\Gamma}{\psi'} = x'' + y^{2} + z_{i}^{n}\]
Don't start a paragraph with a displayed equation,
nor make one a paragraph by itself.

\end{document}               % End of document.

\section{Bench Huge Section 359}
% synthetic bench section 359
% This is a sample LaTeX input file.  (Version of 12 August 2004.)
%
% A '%' character causes TeX to ignore all remaining text on the line,
% and is used for comments like this one.

\documentclass{article}      % Specifies the document class

                             % The preamble begins here.
\title{An Example Document}  % Declares the document's title.
\author{Leslie Lamport}      % Declares the author's name.
\date{January 21, 1994}      % Deleting this command produces today's date.

\newcommand{\ip}[2]{(#1, #2)}
                             % Defines \ip{arg1}{arg2} to mean
                             % (arg1, arg2).

%\newcommand{\ip}[2]{\langle #1 | #2\rangle}
                             % This is an alternative definition of
                             % \ip that is commented out.

\begin{document}             % End of preamble and beginning of text.

\maketitle                   % Produces the title.

This is an example input file.  Comparing it with
the output it generates can show you how to
produce a simple document of your own.

\section{Ordinary Text}      % Produces section heading.  Lower-level
                             % sections are begun with similar
                             % \subsection and \subsubsection commands.

The ends  of words and sentences are marked
  by   spaces. It  doesn't matter how many
spaces    you type; one is as good as 100.  The
end of   a line counts as a space.

One   or more   blank lines denote the  end
of  a paragraph.

Since any number of consecutive spaces are treated
like a single one, the formatting of the input
file makes no difference to
      \LaTeX,                % The \LaTeX command generates the LaTeX logo.
but it makes a difference to you.  When you use
\LaTeX, making your input file as easy to read
as possible will be a great help as you write
your document and when you change it.  This sample
file shows how you can add comments to your own input
file.

Because printing is different from typewriting,
there are a number of things that you have to do
differently when preparing an input file than if
you were just typing the document directly.
Quotation marks like
       ``this''
have to be handled specially, as do quotes within
quotes:
       ``\,`this'            % \, separates the double and single quote.
        is what I just
        wrote, not  `that'\,''.

Dashes come in three sizes: an
       intra-word
dash, a medium dash for number ranges like
       1--2,
and a punctuation
       dash---like
this.

A sentence-ending space should be larger than the
space between words within a sentence.  You
sometimes have to type special commands in
conjunction with punctuation characters to get
this right, as in the following sentence.
       Gnats, gnus, etc.\ all  % `\ ' makes an inter-word space.
       begin with G\@.         % \@ marks end-of-sentence punctuation.
You should check the spaces after periods when
reading your output to make sure you haven't
forgotten any special cases.  Generating an
ellipsis
       \ldots\               % `\ ' is needed after `\ldots' because TeX
                             % ignores spaces after command names like \ldots
                             % made from \ + letters.
                             %
                             % Note how a `%' character causes TeX to ignore
                             % the end of the input line, so these blank lines
                             % do not start a new paragraph.
                             %
with the right spacing around the periods requires
a special command.

\LaTeX\ interprets some common characters as
commands, so you must type special commands to
generate them.  These characters include the
following:
       \$ \& \% \# \{ and \}.

In printing, text is usually emphasized with an
       \emph{italic}
type style.

\begin{em}
   A long segment of text can also be emphasized
   in this way.  Text within such a segment can be
   given \emph{additional} emphasis.
\end{em}

It is sometimes necessary to prevent \LaTeX\ from
breaking a line where it might otherwise do so.
This may be at a space, as between the ``Mr.''\ and
``Jones'' in
       ``Mr.~Jones'',        % ~ produces an unbreakable interword space.
or within a word---especially when the word is a
symbol like
       \mbox{\emph{itemnum}}
that makes little sense when hyphenated across
lines.

Footnotes\footnote{This is an example of a footnote.}
pose no problem.

\LaTeX\ is good at typesetting mathematical formulas
like
       \( x-3y + z = 7 \)
or
       \( a_{1} > x^{2n} + y^{2n} > x' \)
or
       \( \ip{A}{B} = \sum_{i} a_{i} b_{i} \).
The spaces you type in a formula are
ignored.  Remember that a letter like
       $x$                   % $ ... $  and  \( ... \)  are equivalent
is a formula when it denotes a mathematical
symbol, and it should be typed as one.

\section{Displayed Text}

Text is displayed by indenting it from the left
margin.  Quotations are commonly displayed.  There
are short quotations
\begin{quote}
   This is a short quotation.  It consists of a
   single paragraph of text.  See how it is formatted.
\end{quote}
and longer ones.
\begin{quotation}
   This is a longer quotation.  It consists of two
   paragraphs of text, neither of which are
   particularly interesting.

   This is the second paragraph of the quotation.  It
   is just as dull as the first paragraph.
\end{quotation}
Another frequently-displayed structure is a list.
The following is an example of an \emph{itemized}
list.
\begin{itemize}
   \item This is the first item of an itemized list.
         Each item in the list is marked with a ``tick''.
         You don't have to worry about what kind of tick
         mark is used.

   \item This is the second item of the list.  It
         contains another list nested inside it.  The inner
         list is an \emph{enumerated} list.
         \begin{enumerate}
            \item This is the first item of an enumerated
                  list that is nested within the itemized list.

            \item This is the second item of the inner list.
                  \LaTeX\ allows you to nest lists deeper than
                  you really should.
         \end{enumerate}
         This is the rest of the second item of the outer
         list.  It is no more interesting than any other
         part of the item.
   \item This is the third item of the list.
\end{itemize}
You can even display poetry.
\begin{verse}
   There is an environment
    for verse \\             % The \\ command separates lines
   Whose features some poets % within a stanza.
   will curse.

                             % One or more blank lines separate stanzas.

   For instead of making\\
   Them do \emph{all} line breaking, \\
   It allows them to put too many words on a line when they'd rather be
   forced to be terse.
\end{verse}

Mathematical formulas may also be displayed.  A
displayed formula
is
one-line long; multiline
formulas require special formatting instructions.
   \[  \ip{\Gamma}{\psi'} = x'' + y^{2} + z_{i}^{n}\]
Don't start a paragraph with a displayed equation,
nor make one a paragraph by itself.

\end{document}               % End of document.

\section{Bench Huge Section 360}
% synthetic bench section 360
% This is a sample LaTeX input file.  (Version of 12 August 2004.)
%
% A '%' character causes TeX to ignore all remaining text on the line,
% and is used for comments like this one.

\documentclass{article}      % Specifies the document class

                             % The preamble begins here.
\title{An Example Document}  % Declares the document's title.
\author{Leslie Lamport}      % Declares the author's name.
\date{January 21, 1994}      % Deleting this command produces today's date.

\newcommand{\ip}[2]{(#1, #2)}
                             % Defines \ip{arg1}{arg2} to mean
                             % (arg1, arg2).

%\newcommand{\ip}[2]{\langle #1 | #2\rangle}
                             % This is an alternative definition of
                             % \ip that is commented out.

\begin{document}             % End of preamble and beginning of text.

\maketitle                   % Produces the title.

This is an example input file.  Comparing it with
the output it generates can show you how to
produce a simple document of your own.

\section{Ordinary Text}      % Produces section heading.  Lower-level
                             % sections are begun with similar
                             % \subsection and \subsubsection commands.

The ends  of words and sentences are marked
  by   spaces. It  doesn't matter how many
spaces    you type; one is as good as 100.  The
end of   a line counts as a space.

One   or more   blank lines denote the  end
of  a paragraph.

Since any number of consecutive spaces are treated
like a single one, the formatting of the input
file makes no difference to
      \LaTeX,                % The \LaTeX command generates the LaTeX logo.
but it makes a difference to you.  When you use
\LaTeX, making your input file as easy to read
as possible will be a great help as you write
your document and when you change it.  This sample
file shows how you can add comments to your own input
file.

Because printing is different from typewriting,
there are a number of things that you have to do
differently when preparing an input file than if
you were just typing the document directly.
Quotation marks like
       ``this''
have to be handled specially, as do quotes within
quotes:
       ``\,`this'            % \, separates the double and single quote.
        is what I just
        wrote, not  `that'\,''.

Dashes come in three sizes: an
       intra-word
dash, a medium dash for number ranges like
       1--2,
and a punctuation
       dash---like
this.

A sentence-ending space should be larger than the
space between words within a sentence.  You
sometimes have to type special commands in
conjunction with punctuation characters to get
this right, as in the following sentence.
       Gnats, gnus, etc.\ all  % `\ ' makes an inter-word space.
       begin with G\@.         % \@ marks end-of-sentence punctuation.
You should check the spaces after periods when
reading your output to make sure you haven't
forgotten any special cases.  Generating an
ellipsis
       \ldots\               % `\ ' is needed after `\ldots' because TeX
                             % ignores spaces after command names like \ldots
                             % made from \ + letters.
                             %
                             % Note how a `%' character causes TeX to ignore
                             % the end of the input line, so these blank lines
                             % do not start a new paragraph.
                             %
with the right spacing around the periods requires
a special command.

\LaTeX\ interprets some common characters as
commands, so you must type special commands to
generate them.  These characters include the
following:
       \$ \& \% \# \{ and \}.

In printing, text is usually emphasized with an
       \emph{italic}
type style.

\begin{em}
   A long segment of text can also be emphasized
   in this way.  Text within such a segment can be
   given \emph{additional} emphasis.
\end{em}

It is sometimes necessary to prevent \LaTeX\ from
breaking a line where it might otherwise do so.
This may be at a space, as between the ``Mr.''\ and
``Jones'' in
       ``Mr.~Jones'',        % ~ produces an unbreakable interword space.
or within a word---especially when the word is a
symbol like
       \mbox{\emph{itemnum}}
that makes little sense when hyphenated across
lines.

Footnotes\footnote{This is an example of a footnote.}
pose no problem.

\LaTeX\ is good at typesetting mathematical formulas
like
       \( x-3y + z = 7 \)
or
       \( a_{1} > x^{2n} + y^{2n} > x' \)
or
       \( \ip{A}{B} = \sum_{i} a_{i} b_{i} \).
The spaces you type in a formula are
ignored.  Remember that a letter like
       $x$                   % $ ... $  and  \( ... \)  are equivalent
is a formula when it denotes a mathematical
symbol, and it should be typed as one.

\section{Displayed Text}

Text is displayed by indenting it from the left
margin.  Quotations are commonly displayed.  There
are short quotations
\begin{quote}
   This is a short quotation.  It consists of a
   single paragraph of text.  See how it is formatted.
\end{quote}
and longer ones.
\begin{quotation}
   This is a longer quotation.  It consists of two
   paragraphs of text, neither of which are
   particularly interesting.

   This is the second paragraph of the quotation.  It
   is just as dull as the first paragraph.
\end{quotation}
Another frequently-displayed structure is a list.
The following is an example of an \emph{itemized}
list.
\begin{itemize}
   \item This is the first item of an itemized list.
         Each item in the list is marked with a ``tick''.
         You don't have to worry about what kind of tick
         mark is used.

   \item This is the second item of the list.  It
         contains another list nested inside it.  The inner
         list is an \emph{enumerated} list.
         \begin{enumerate}
            \item This is the first item of an enumerated
                  list that is nested within the itemized list.

            \item This is the second item of the inner list.
                  \LaTeX\ allows you to nest lists deeper than
                  you really should.
         \end{enumerate}
         This is the rest of the second item of the outer
         list.  It is no more interesting than any other
         part of the item.
   \item This is the third item of the list.
\end{itemize}
You can even display poetry.
\begin{verse}
   There is an environment
    for verse \\             % The \\ command separates lines
   Whose features some poets % within a stanza.
   will curse.

                             % One or more blank lines separate stanzas.

   For instead of making\\
   Them do \emph{all} line breaking, \\
   It allows them to put too many words on a line when they'd rather be
   forced to be terse.
\end{verse}

Mathematical formulas may also be displayed.  A
displayed formula
is
one-line long; multiline
formulas require special formatting instructions.
   \[  \ip{\Gamma}{\psi'} = x'' + y^{2} + z_{i}^{n}\]
Don't start a paragraph with a displayed equation,
nor make one a paragraph by itself.

\end{document}               % End of document.

\section{Bench Huge Section 361}
% synthetic bench section 361
% This is a sample LaTeX input file.  (Version of 12 August 2004.)
%
% A '%' character causes TeX to ignore all remaining text on the line,
% and is used for comments like this one.

\documentclass{article}      % Specifies the document class

                             % The preamble begins here.
\title{An Example Document}  % Declares the document's title.
\author{Leslie Lamport}      % Declares the author's name.
\date{January 21, 1994}      % Deleting this command produces today's date.

\newcommand{\ip}[2]{(#1, #2)}
                             % Defines \ip{arg1}{arg2} to mean
                             % (arg1, arg2).

%\newcommand{\ip}[2]{\langle #1 | #2\rangle}
                             % This is an alternative definition of
                             % \ip that is commented out.

\begin{document}             % End of preamble and beginning of text.

\maketitle                   % Produces the title.

This is an example input file.  Comparing it with
the output it generates can show you how to
produce a simple document of your own.

\section{Ordinary Text}      % Produces section heading.  Lower-level
                             % sections are begun with similar
                             % \subsection and \subsubsection commands.

The ends  of words and sentences are marked
  by   spaces. It  doesn't matter how many
spaces    you type; one is as good as 100.  The
end of   a line counts as a space.

One   or more   blank lines denote the  end
of  a paragraph.

Since any number of consecutive spaces are treated
like a single one, the formatting of the input
file makes no difference to
      \LaTeX,                % The \LaTeX command generates the LaTeX logo.
but it makes a difference to you.  When you use
\LaTeX, making your input file as easy to read
as possible will be a great help as you write
your document and when you change it.  This sample
file shows how you can add comments to your own input
file.

Because printing is different from typewriting,
there are a number of things that you have to do
differently when preparing an input file than if
you were just typing the document directly.
Quotation marks like
       ``this''
have to be handled specially, as do quotes within
quotes:
       ``\,`this'            % \, separates the double and single quote.
        is what I just
        wrote, not  `that'\,''.

Dashes come in three sizes: an
       intra-word
dash, a medium dash for number ranges like
       1--2,
and a punctuation
       dash---like
this.

A sentence-ending space should be larger than the
space between words within a sentence.  You
sometimes have to type special commands in
conjunction with punctuation characters to get
this right, as in the following sentence.
       Gnats, gnus, etc.\ all  % `\ ' makes an inter-word space.
       begin with G\@.         % \@ marks end-of-sentence punctuation.
You should check the spaces after periods when
reading your output to make sure you haven't
forgotten any special cases.  Generating an
ellipsis
       \ldots\               % `\ ' is needed after `\ldots' because TeX
                             % ignores spaces after command names like \ldots
                             % made from \ + letters.
                             %
                             % Note how a `%' character causes TeX to ignore
                             % the end of the input line, so these blank lines
                             % do not start a new paragraph.
                             %
with the right spacing around the periods requires
a special command.

\LaTeX\ interprets some common characters as
commands, so you must type special commands to
generate them.  These characters include the
following:
       \$ \& \% \# \{ and \}.

In printing, text is usually emphasized with an
       \emph{italic}
type style.

\begin{em}
   A long segment of text can also be emphasized
   in this way.  Text within such a segment can be
   given \emph{additional} emphasis.
\end{em}

It is sometimes necessary to prevent \LaTeX\ from
breaking a line where it might otherwise do so.
This may be at a space, as between the ``Mr.''\ and
``Jones'' in
       ``Mr.~Jones'',        % ~ produces an unbreakable interword space.
or within a word---especially when the word is a
symbol like
       \mbox{\emph{itemnum}}
that makes little sense when hyphenated across
lines.

Footnotes\footnote{This is an example of a footnote.}
pose no problem.

\LaTeX\ is good at typesetting mathematical formulas
like
       \( x-3y + z = 7 \)
or
       \( a_{1} > x^{2n} + y^{2n} > x' \)
or
       \( \ip{A}{B} = \sum_{i} a_{i} b_{i} \).
The spaces you type in a formula are
ignored.  Remember that a letter like
       $x$                   % $ ... $  and  \( ... \)  are equivalent
is a formula when it denotes a mathematical
symbol, and it should be typed as one.

\section{Displayed Text}

Text is displayed by indenting it from the left
margin.  Quotations are commonly displayed.  There
are short quotations
\begin{quote}
   This is a short quotation.  It consists of a
   single paragraph of text.  See how it is formatted.
\end{quote}
and longer ones.
\begin{quotation}
   This is a longer quotation.  It consists of two
   paragraphs of text, neither of which are
   particularly interesting.

   This is the second paragraph of the quotation.  It
   is just as dull as the first paragraph.
\end{quotation}
Another frequently-displayed structure is a list.
The following is an example of an \emph{itemized}
list.
\begin{itemize}
   \item This is the first item of an itemized list.
         Each item in the list is marked with a ``tick''.
         You don't have to worry about what kind of tick
         mark is used.

   \item This is the second item of the list.  It
         contains another list nested inside it.  The inner
         list is an \emph{enumerated} list.
         \begin{enumerate}
            \item This is the first item of an enumerated
                  list that is nested within the itemized list.

            \item This is the second item of the inner list.
                  \LaTeX\ allows you to nest lists deeper than
                  you really should.
         \end{enumerate}
         This is the rest of the second item of the outer
         list.  It is no more interesting than any other
         part of the item.
   \item This is the third item of the list.
\end{itemize}
You can even display poetry.
\begin{verse}
   There is an environment
    for verse \\             % The \\ command separates lines
   Whose features some poets % within a stanza.
   will curse.

                             % One or more blank lines separate stanzas.

   For instead of making\\
   Them do \emph{all} line breaking, \\
   It allows them to put too many words on a line when they'd rather be
   forced to be terse.
\end{verse}

Mathematical formulas may also be displayed.  A
displayed formula
is
one-line long; multiline
formulas require special formatting instructions.
   \[  \ip{\Gamma}{\psi'} = x'' + y^{2} + z_{i}^{n}\]
Don't start a paragraph with a displayed equation,
nor make one a paragraph by itself.

\end{document}               % End of document.

\section{Bench Huge Section 362}
% synthetic bench section 362
% This is a sample LaTeX input file.  (Version of 12 August 2004.)
%
% A '%' character causes TeX to ignore all remaining text on the line,
% and is used for comments like this one.

\documentclass{article}      % Specifies the document class

                             % The preamble begins here.
\title{An Example Document}  % Declares the document's title.
\author{Leslie Lamport}      % Declares the author's name.
\date{January 21, 1994}      % Deleting this command produces today's date.

\newcommand{\ip}[2]{(#1, #2)}
                             % Defines \ip{arg1}{arg2} to mean
                             % (arg1, arg2).

%\newcommand{\ip}[2]{\langle #1 | #2\rangle}
                             % This is an alternative definition of
                             % \ip that is commented out.

\begin{document}             % End of preamble and beginning of text.

\maketitle                   % Produces the title.

This is an example input file.  Comparing it with
the output it generates can show you how to
produce a simple document of your own.

\section{Ordinary Text}      % Produces section heading.  Lower-level
                             % sections are begun with similar
                             % \subsection and \subsubsection commands.

The ends  of words and sentences are marked
  by   spaces. It  doesn't matter how many
spaces    you type; one is as good as 100.  The
end of   a line counts as a space.

One   or more   blank lines denote the  end
of  a paragraph.

Since any number of consecutive spaces are treated
like a single one, the formatting of the input
file makes no difference to
      \LaTeX,                % The \LaTeX command generates the LaTeX logo.
but it makes a difference to you.  When you use
\LaTeX, making your input file as easy to read
as possible will be a great help as you write
your document and when you change it.  This sample
file shows how you can add comments to your own input
file.

Because printing is different from typewriting,
there are a number of things that you have to do
differently when preparing an input file than if
you were just typing the document directly.
Quotation marks like
       ``this''
have to be handled specially, as do quotes within
quotes:
       ``\,`this'            % \, separates the double and single quote.
        is what I just
        wrote, not  `that'\,''.

Dashes come in three sizes: an
       intra-word
dash, a medium dash for number ranges like
       1--2,
and a punctuation
       dash---like
this.

A sentence-ending space should be larger than the
space between words within a sentence.  You
sometimes have to type special commands in
conjunction with punctuation characters to get
this right, as in the following sentence.
       Gnats, gnus, etc.\ all  % `\ ' makes an inter-word space.
       begin with G\@.         % \@ marks end-of-sentence punctuation.
You should check the spaces after periods when
reading your output to make sure you haven't
forgotten any special cases.  Generating an
ellipsis
       \ldots\               % `\ ' is needed after `\ldots' because TeX
                             % ignores spaces after command names like \ldots
                             % made from \ + letters.
                             %
                             % Note how a `%' character causes TeX to ignore
                             % the end of the input line, so these blank lines
                             % do not start a new paragraph.
                             %
with the right spacing around the periods requires
a special command.

\LaTeX\ interprets some common characters as
commands, so you must type special commands to
generate them.  These characters include the
following:
       \$ \& \% \# \{ and \}.

In printing, text is usually emphasized with an
       \emph{italic}
type style.

\begin{em}
   A long segment of text can also be emphasized
   in this way.  Text within such a segment can be
   given \emph{additional} emphasis.
\end{em}

It is sometimes necessary to prevent \LaTeX\ from
breaking a line where it might otherwise do so.
This may be at a space, as between the ``Mr.''\ and
``Jones'' in
       ``Mr.~Jones'',        % ~ produces an unbreakable interword space.
or within a word---especially when the word is a
symbol like
       \mbox{\emph{itemnum}}
that makes little sense when hyphenated across
lines.

Footnotes\footnote{This is an example of a footnote.}
pose no problem.

\LaTeX\ is good at typesetting mathematical formulas
like
       \( x-3y + z = 7 \)
or
       \( a_{1} > x^{2n} + y^{2n} > x' \)
or
       \( \ip{A}{B} = \sum_{i} a_{i} b_{i} \).
The spaces you type in a formula are
ignored.  Remember that a letter like
       $x$                   % $ ... $  and  \( ... \)  are equivalent
is a formula when it denotes a mathematical
symbol, and it should be typed as one.

\section{Displayed Text}

Text is displayed by indenting it from the left
margin.  Quotations are commonly displayed.  There
are short quotations
\begin{quote}
   This is a short quotation.  It consists of a
   single paragraph of text.  See how it is formatted.
\end{quote}
and longer ones.
\begin{quotation}
   This is a longer quotation.  It consists of two
   paragraphs of text, neither of which are
   particularly interesting.

   This is the second paragraph of the quotation.  It
   is just as dull as the first paragraph.
\end{quotation}
Another frequently-displayed structure is a list.
The following is an example of an \emph{itemized}
list.
\begin{itemize}
   \item This is the first item of an itemized list.
         Each item in the list is marked with a ``tick''.
         You don't have to worry about what kind of tick
         mark is used.

   \item This is the second item of the list.  It
         contains another list nested inside it.  The inner
         list is an \emph{enumerated} list.
         \begin{enumerate}
            \item This is the first item of an enumerated
                  list that is nested within the itemized list.

            \item This is the second item of the inner list.
                  \LaTeX\ allows you to nest lists deeper than
                  you really should.
         \end{enumerate}
         This is the rest of the second item of the outer
         list.  It is no more interesting than any other
         part of the item.
   \item This is the third item of the list.
\end{itemize}
You can even display poetry.
\begin{verse}
   There is an environment
    for verse \\             % The \\ command separates lines
   Whose features some poets % within a stanza.
   will curse.

                             % One or more blank lines separate stanzas.

   For instead of making\\
   Them do \emph{all} line breaking, \\
   It allows them to put too many words on a line when they'd rather be
   forced to be terse.
\end{verse}

Mathematical formulas may also be displayed.  A
displayed formula
is
one-line long; multiline
formulas require special formatting instructions.
   \[  \ip{\Gamma}{\psi'} = x'' + y^{2} + z_{i}^{n}\]
Don't start a paragraph with a displayed equation,
nor make one a paragraph by itself.

\end{document}               % End of document.

\section{Bench Huge Section 363}
% synthetic bench section 363
% This is a sample LaTeX input file.  (Version of 12 August 2004.)
%
% A '%' character causes TeX to ignore all remaining text on the line,
% and is used for comments like this one.

\documentclass{article}      % Specifies the document class

                             % The preamble begins here.
\title{An Example Document}  % Declares the document's title.
\author{Leslie Lamport}      % Declares the author's name.
\date{January 21, 1994}      % Deleting this command produces today's date.

\newcommand{\ip}[2]{(#1, #2)}
                             % Defines \ip{arg1}{arg2} to mean
                             % (arg1, arg2).

%\newcommand{\ip}[2]{\langle #1 | #2\rangle}
                             % This is an alternative definition of
                             % \ip that is commented out.

\begin{document}             % End of preamble and beginning of text.

\maketitle                   % Produces the title.

This is an example input file.  Comparing it with
the output it generates can show you how to
produce a simple document of your own.

\section{Ordinary Text}      % Produces section heading.  Lower-level
                             % sections are begun with similar
                             % \subsection and \subsubsection commands.

The ends  of words and sentences are marked
  by   spaces. It  doesn't matter how many
spaces    you type; one is as good as 100.  The
end of   a line counts as a space.

One   or more   blank lines denote the  end
of  a paragraph.

Since any number of consecutive spaces are treated
like a single one, the formatting of the input
file makes no difference to
      \LaTeX,                % The \LaTeX command generates the LaTeX logo.
but it makes a difference to you.  When you use
\LaTeX, making your input file as easy to read
as possible will be a great help as you write
your document and when you change it.  This sample
file shows how you can add comments to your own input
file.

Because printing is different from typewriting,
there are a number of things that you have to do
differently when preparing an input file than if
you were just typing the document directly.
Quotation marks like
       ``this''
have to be handled specially, as do quotes within
quotes:
       ``\,`this'            % \, separates the double and single quote.
        is what I just
        wrote, not  `that'\,''.

Dashes come in three sizes: an
       intra-word
dash, a medium dash for number ranges like
       1--2,
and a punctuation
       dash---like
this.

A sentence-ending space should be larger than the
space between words within a sentence.  You
sometimes have to type special commands in
conjunction with punctuation characters to get
this right, as in the following sentence.
       Gnats, gnus, etc.\ all  % `\ ' makes an inter-word space.
       begin with G\@.         % \@ marks end-of-sentence punctuation.
You should check the spaces after periods when
reading your output to make sure you haven't
forgotten any special cases.  Generating an
ellipsis
       \ldots\               % `\ ' is needed after `\ldots' because TeX
                             % ignores spaces after command names like \ldots
                             % made from \ + letters.
                             %
                             % Note how a `%' character causes TeX to ignore
                             % the end of the input line, so these blank lines
                             % do not start a new paragraph.
                             %
with the right spacing around the periods requires
a special command.

\LaTeX\ interprets some common characters as
commands, so you must type special commands to
generate them.  These characters include the
following:
       \$ \& \% \# \{ and \}.

In printing, text is usually emphasized with an
       \emph{italic}
type style.

\begin{em}
   A long segment of text can also be emphasized
   in this way.  Text within such a segment can be
   given \emph{additional} emphasis.
\end{em}

It is sometimes necessary to prevent \LaTeX\ from
breaking a line where it might otherwise do so.
This may be at a space, as between the ``Mr.''\ and
``Jones'' in
       ``Mr.~Jones'',        % ~ produces an unbreakable interword space.
or within a word---especially when the word is a
symbol like
       \mbox{\emph{itemnum}}
that makes little sense when hyphenated across
lines.

Footnotes\footnote{This is an example of a footnote.}
pose no problem.

\LaTeX\ is good at typesetting mathematical formulas
like
       \( x-3y + z = 7 \)
or
       \( a_{1} > x^{2n} + y^{2n} > x' \)
or
       \( \ip{A}{B} = \sum_{i} a_{i} b_{i} \).
The spaces you type in a formula are
ignored.  Remember that a letter like
       $x$                   % $ ... $  and  \( ... \)  are equivalent
is a formula when it denotes a mathematical
symbol, and it should be typed as one.

\section{Displayed Text}

Text is displayed by indenting it from the left
margin.  Quotations are commonly displayed.  There
are short quotations
\begin{quote}
   This is a short quotation.  It consists of a
   single paragraph of text.  See how it is formatted.
\end{quote}
and longer ones.
\begin{quotation}
   This is a longer quotation.  It consists of two
   paragraphs of text, neither of which are
   particularly interesting.

   This is the second paragraph of the quotation.  It
   is just as dull as the first paragraph.
\end{quotation}
Another frequently-displayed structure is a list.
The following is an example of an \emph{itemized}
list.
\begin{itemize}
   \item This is the first item of an itemized list.
         Each item in the list is marked with a ``tick''.
         You don't have to worry about what kind of tick
         mark is used.

   \item This is the second item of the list.  It
         contains another list nested inside it.  The inner
         list is an \emph{enumerated} list.
         \begin{enumerate}
            \item This is the first item of an enumerated
                  list that is nested within the itemized list.

            \item This is the second item of the inner list.
                  \LaTeX\ allows you to nest lists deeper than
                  you really should.
         \end{enumerate}
         This is the rest of the second item of the outer
         list.  It is no more interesting than any other
         part of the item.
   \item This is the third item of the list.
\end{itemize}
You can even display poetry.
\begin{verse}
   There is an environment
    for verse \\             % The \\ command separates lines
   Whose features some poets % within a stanza.
   will curse.

                             % One or more blank lines separate stanzas.

   For instead of making\\
   Them do \emph{all} line breaking, \\
   It allows them to put too many words on a line when they'd rather be
   forced to be terse.
\end{verse}

Mathematical formulas may also be displayed.  A
displayed formula
is
one-line long; multiline
formulas require special formatting instructions.
   \[  \ip{\Gamma}{\psi'} = x'' + y^{2} + z_{i}^{n}\]
Don't start a paragraph with a displayed equation,
nor make one a paragraph by itself.

\end{document}               % End of document.

\section{Bench Huge Section 364}
% synthetic bench section 364
% This is a sample LaTeX input file.  (Version of 12 August 2004.)
%
% A '%' character causes TeX to ignore all remaining text on the line,
% and is used for comments like this one.

\documentclass{article}      % Specifies the document class

                             % The preamble begins here.
\title{An Example Document}  % Declares the document's title.
\author{Leslie Lamport}      % Declares the author's name.
\date{January 21, 1994}      % Deleting this command produces today's date.

\newcommand{\ip}[2]{(#1, #2)}
                             % Defines \ip{arg1}{arg2} to mean
                             % (arg1, arg2).

%\newcommand{\ip}[2]{\langle #1 | #2\rangle}
                             % This is an alternative definition of
                             % \ip that is commented out.

\begin{document}             % End of preamble and beginning of text.

\maketitle                   % Produces the title.

This is an example input file.  Comparing it with
the output it generates can show you how to
produce a simple document of your own.

\section{Ordinary Text}      % Produces section heading.  Lower-level
                             % sections are begun with similar
                             % \subsection and \subsubsection commands.

The ends  of words and sentences are marked
  by   spaces. It  doesn't matter how many
spaces    you type; one is as good as 100.  The
end of   a line counts as a space.

One   or more   blank lines denote the  end
of  a paragraph.

Since any number of consecutive spaces are treated
like a single one, the formatting of the input
file makes no difference to
      \LaTeX,                % The \LaTeX command generates the LaTeX logo.
but it makes a difference to you.  When you use
\LaTeX, making your input file as easy to read
as possible will be a great help as you write
your document and when you change it.  This sample
file shows how you can add comments to your own input
file.

Because printing is different from typewriting,
there are a number of things that you have to do
differently when preparing an input file than if
you were just typing the document directly.
Quotation marks like
       ``this''
have to be handled specially, as do quotes within
quotes:
       ``\,`this'            % \, separates the double and single quote.
        is what I just
        wrote, not  `that'\,''.

Dashes come in three sizes: an
       intra-word
dash, a medium dash for number ranges like
       1--2,
and a punctuation
       dash---like
this.

A sentence-ending space should be larger than the
space between words within a sentence.  You
sometimes have to type special commands in
conjunction with punctuation characters to get
this right, as in the following sentence.
       Gnats, gnus, etc.\ all  % `\ ' makes an inter-word space.
       begin with G\@.         % \@ marks end-of-sentence punctuation.
You should check the spaces after periods when
reading your output to make sure you haven't
forgotten any special cases.  Generating an
ellipsis
       \ldots\               % `\ ' is needed after `\ldots' because TeX
                             % ignores spaces after command names like \ldots
                             % made from \ + letters.
                             %
                             % Note how a `%' character causes TeX to ignore
                             % the end of the input line, so these blank lines
                             % do not start a new paragraph.
                             %
with the right spacing around the periods requires
a special command.

\LaTeX\ interprets some common characters as
commands, so you must type special commands to
generate them.  These characters include the
following:
       \$ \& \% \# \{ and \}.

In printing, text is usually emphasized with an
       \emph{italic}
type style.

\begin{em}
   A long segment of text can also be emphasized
   in this way.  Text within such a segment can be
   given \emph{additional} emphasis.
\end{em}

It is sometimes necessary to prevent \LaTeX\ from
breaking a line where it might otherwise do so.
This may be at a space, as between the ``Mr.''\ and
``Jones'' in
       ``Mr.~Jones'',        % ~ produces an unbreakable interword space.
or within a word---especially when the word is a
symbol like
       \mbox{\emph{itemnum}}
that makes little sense when hyphenated across
lines.

Footnotes\footnote{This is an example of a footnote.}
pose no problem.

\LaTeX\ is good at typesetting mathematical formulas
like
       \( x-3y + z = 7 \)
or
       \( a_{1} > x^{2n} + y^{2n} > x' \)
or
       \( \ip{A}{B} = \sum_{i} a_{i} b_{i} \).
The spaces you type in a formula are
ignored.  Remember that a letter like
       $x$                   % $ ... $  and  \( ... \)  are equivalent
is a formula when it denotes a mathematical
symbol, and it should be typed as one.

\section{Displayed Text}

Text is displayed by indenting it from the left
margin.  Quotations are commonly displayed.  There
are short quotations
\begin{quote}
   This is a short quotation.  It consists of a
   single paragraph of text.  See how it is formatted.
\end{quote}
and longer ones.
\begin{quotation}
   This is a longer quotation.  It consists of two
   paragraphs of text, neither of which are
   particularly interesting.

   This is the second paragraph of the quotation.  It
   is just as dull as the first paragraph.
\end{quotation}
Another frequently-displayed structure is a list.
The following is an example of an \emph{itemized}
list.
\begin{itemize}
   \item This is the first item of an itemized list.
         Each item in the list is marked with a ``tick''.
         You don't have to worry about what kind of tick
         mark is used.

   \item This is the second item of the list.  It
         contains another list nested inside it.  The inner
         list is an \emph{enumerated} list.
         \begin{enumerate}
            \item This is the first item of an enumerated
                  list that is nested within the itemized list.

            \item This is the second item of the inner list.
                  \LaTeX\ allows you to nest lists deeper than
                  you really should.
         \end{enumerate}
         This is the rest of the second item of the outer
         list.  It is no more interesting than any other
         part of the item.
   \item This is the third item of the list.
\end{itemize}
You can even display poetry.
\begin{verse}
   There is an environment
    for verse \\             % The \\ command separates lines
   Whose features some poets % within a stanza.
   will curse.

                             % One or more blank lines separate stanzas.

   For instead of making\\
   Them do \emph{all} line breaking, \\
   It allows them to put too many words on a line when they'd rather be
   forced to be terse.
\end{verse}

Mathematical formulas may also be displayed.  A
displayed formula
is
one-line long; multiline
formulas require special formatting instructions.
   \[  \ip{\Gamma}{\psi'} = x'' + y^{2} + z_{i}^{n}\]
Don't start a paragraph with a displayed equation,
nor make one a paragraph by itself.

\end{document}               % End of document.

\section{Bench Huge Section 365}
% synthetic bench section 365
% This is a sample LaTeX input file.  (Version of 12 August 2004.)
%
% A '%' character causes TeX to ignore all remaining text on the line,
% and is used for comments like this one.

\documentclass{article}      % Specifies the document class

                             % The preamble begins here.
\title{An Example Document}  % Declares the document's title.
\author{Leslie Lamport}      % Declares the author's name.
\date{January 21, 1994}      % Deleting this command produces today's date.

\newcommand{\ip}[2]{(#1, #2)}
                             % Defines \ip{arg1}{arg2} to mean
                             % (arg1, arg2).

%\newcommand{\ip}[2]{\langle #1 | #2\rangle}
                             % This is an alternative definition of
                             % \ip that is commented out.

\begin{document}             % End of preamble and beginning of text.

\maketitle                   % Produces the title.

This is an example input file.  Comparing it with
the output it generates can show you how to
produce a simple document of your own.

\section{Ordinary Text}      % Produces section heading.  Lower-level
                             % sections are begun with similar
                             % \subsection and \subsubsection commands.

The ends  of words and sentences are marked
  by   spaces. It  doesn't matter how many
spaces    you type; one is as good as 100.  The
end of   a line counts as a space.

One   or more   blank lines denote the  end
of  a paragraph.

Since any number of consecutive spaces are treated
like a single one, the formatting of the input
file makes no difference to
      \LaTeX,                % The \LaTeX command generates the LaTeX logo.
but it makes a difference to you.  When you use
\LaTeX, making your input file as easy to read
as possible will be a great help as you write
your document and when you change it.  This sample
file shows how you can add comments to your own input
file.

Because printing is different from typewriting,
there are a number of things that you have to do
differently when preparing an input file than if
you were just typing the document directly.
Quotation marks like
       ``this''
have to be handled specially, as do quotes within
quotes:
       ``\,`this'            % \, separates the double and single quote.
        is what I just
        wrote, not  `that'\,''.

Dashes come in three sizes: an
       intra-word
dash, a medium dash for number ranges like
       1--2,
and a punctuation
       dash---like
this.

A sentence-ending space should be larger than the
space between words within a sentence.  You
sometimes have to type special commands in
conjunction with punctuation characters to get
this right, as in the following sentence.
       Gnats, gnus, etc.\ all  % `\ ' makes an inter-word space.
       begin with G\@.         % \@ marks end-of-sentence punctuation.
You should check the spaces after periods when
reading your output to make sure you haven't
forgotten any special cases.  Generating an
ellipsis
       \ldots\               % `\ ' is needed after `\ldots' because TeX
                             % ignores spaces after command names like \ldots
                             % made from \ + letters.
                             %
                             % Note how a `%' character causes TeX to ignore
                             % the end of the input line, so these blank lines
                             % do not start a new paragraph.
                             %
with the right spacing around the periods requires
a special command.

\LaTeX\ interprets some common characters as
commands, so you must type special commands to
generate them.  These characters include the
following:
       \$ \& \% \# \{ and \}.

In printing, text is usually emphasized with an
       \emph{italic}
type style.

\begin{em}
   A long segment of text can also be emphasized
   in this way.  Text within such a segment can be
   given \emph{additional} emphasis.
\end{em}

It is sometimes necessary to prevent \LaTeX\ from
breaking a line where it might otherwise do so.
This may be at a space, as between the ``Mr.''\ and
``Jones'' in
       ``Mr.~Jones'',        % ~ produces an unbreakable interword space.
or within a word---especially when the word is a
symbol like
       \mbox{\emph{itemnum}}
that makes little sense when hyphenated across
lines.

Footnotes\footnote{This is an example of a footnote.}
pose no problem.

\LaTeX\ is good at typesetting mathematical formulas
like
       \( x-3y + z = 7 \)
or
       \( a_{1} > x^{2n} + y^{2n} > x' \)
or
       \( \ip{A}{B} = \sum_{i} a_{i} b_{i} \).
The spaces you type in a formula are
ignored.  Remember that a letter like
       $x$                   % $ ... $  and  \( ... \)  are equivalent
is a formula when it denotes a mathematical
symbol, and it should be typed as one.

\section{Displayed Text}

Text is displayed by indenting it from the left
margin.  Quotations are commonly displayed.  There
are short quotations
\begin{quote}
   This is a short quotation.  It consists of a
   single paragraph of text.  See how it is formatted.
\end{quote}
and longer ones.
\begin{quotation}
   This is a longer quotation.  It consists of two
   paragraphs of text, neither of which are
   particularly interesting.

   This is the second paragraph of the quotation.  It
   is just as dull as the first paragraph.
\end{quotation}
Another frequently-displayed structure is a list.
The following is an example of an \emph{itemized}
list.
\begin{itemize}
   \item This is the first item of an itemized list.
         Each item in the list is marked with a ``tick''.
         You don't have to worry about what kind of tick
         mark is used.

   \item This is the second item of the list.  It
         contains another list nested inside it.  The inner
         list is an \emph{enumerated} list.
         \begin{enumerate}
            \item This is the first item of an enumerated
                  list that is nested within the itemized list.

            \item This is the second item of the inner list.
                  \LaTeX\ allows you to nest lists deeper than
                  you really should.
         \end{enumerate}
         This is the rest of the second item of the outer
         list.  It is no more interesting than any other
         part of the item.
   \item This is the third item of the list.
\end{itemize}
You can even display poetry.
\begin{verse}
   There is an environment
    for verse \\             % The \\ command separates lines
   Whose features some poets % within a stanza.
   will curse.

                             % One or more blank lines separate stanzas.

   For instead of making\\
   Them do \emph{all} line breaking, \\
   It allows them to put too many words on a line when they'd rather be
   forced to be terse.
\end{verse}

Mathematical formulas may also be displayed.  A
displayed formula
is
one-line long; multiline
formulas require special formatting instructions.
   \[  \ip{\Gamma}{\psi'} = x'' + y^{2} + z_{i}^{n}\]
Don't start a paragraph with a displayed equation,
nor make one a paragraph by itself.

\end{document}               % End of document.

\section{Bench Huge Section 366}
% synthetic bench section 366
% This is a sample LaTeX input file.  (Version of 12 August 2004.)
%
% A '%' character causes TeX to ignore all remaining text on the line,
% and is used for comments like this one.

\documentclass{article}      % Specifies the document class

                             % The preamble begins here.
\title{An Example Document}  % Declares the document's title.
\author{Leslie Lamport}      % Declares the author's name.
\date{January 21, 1994}      % Deleting this command produces today's date.

\newcommand{\ip}[2]{(#1, #2)}
                             % Defines \ip{arg1}{arg2} to mean
                             % (arg1, arg2).

%\newcommand{\ip}[2]{\langle #1 | #2\rangle}
                             % This is an alternative definition of
                             % \ip that is commented out.

\begin{document}             % End of preamble and beginning of text.

\maketitle                   % Produces the title.

This is an example input file.  Comparing it with
the output it generates can show you how to
produce a simple document of your own.

\section{Ordinary Text}      % Produces section heading.  Lower-level
                             % sections are begun with similar
                             % \subsection and \subsubsection commands.

The ends  of words and sentences are marked
  by   spaces. It  doesn't matter how many
spaces    you type; one is as good as 100.  The
end of   a line counts as a space.

One   or more   blank lines denote the  end
of  a paragraph.

Since any number of consecutive spaces are treated
like a single one, the formatting of the input
file makes no difference to
      \LaTeX,                % The \LaTeX command generates the LaTeX logo.
but it makes a difference to you.  When you use
\LaTeX, making your input file as easy to read
as possible will be a great help as you write
your document and when you change it.  This sample
file shows how you can add comments to your own input
file.

Because printing is different from typewriting,
there are a number of things that you have to do
differently when preparing an input file than if
you were just typing the document directly.
Quotation marks like
       ``this''
have to be handled specially, as do quotes within
quotes:
       ``\,`this'            % \, separates the double and single quote.
        is what I just
        wrote, not  `that'\,''.

Dashes come in three sizes: an
       intra-word
dash, a medium dash for number ranges like
       1--2,
and a punctuation
       dash---like
this.

A sentence-ending space should be larger than the
space between words within a sentence.  You
sometimes have to type special commands in
conjunction with punctuation characters to get
this right, as in the following sentence.
       Gnats, gnus, etc.\ all  % `\ ' makes an inter-word space.
       begin with G\@.         % \@ marks end-of-sentence punctuation.
You should check the spaces after periods when
reading your output to make sure you haven't
forgotten any special cases.  Generating an
ellipsis
       \ldots\               % `\ ' is needed after `\ldots' because TeX
                             % ignores spaces after command names like \ldots
                             % made from \ + letters.
                             %
                             % Note how a `%' character causes TeX to ignore
                             % the end of the input line, so these blank lines
                             % do not start a new paragraph.
                             %
with the right spacing around the periods requires
a special command.

\LaTeX\ interprets some common characters as
commands, so you must type special commands to
generate them.  These characters include the
following:
       \$ \& \% \# \{ and \}.

In printing, text is usually emphasized with an
       \emph{italic}
type style.

\begin{em}
   A long segment of text can also be emphasized
   in this way.  Text within such a segment can be
   given \emph{additional} emphasis.
\end{em}

It is sometimes necessary to prevent \LaTeX\ from
breaking a line where it might otherwise do so.
This may be at a space, as between the ``Mr.''\ and
``Jones'' in
       ``Mr.~Jones'',        % ~ produces an unbreakable interword space.
or within a word---especially when the word is a
symbol like
       \mbox{\emph{itemnum}}
that makes little sense when hyphenated across
lines.

Footnotes\footnote{This is an example of a footnote.}
pose no problem.

\LaTeX\ is good at typesetting mathematical formulas
like
       \( x-3y + z = 7 \)
or
       \( a_{1} > x^{2n} + y^{2n} > x' \)
or
       \( \ip{A}{B} = \sum_{i} a_{i} b_{i} \).
The spaces you type in a formula are
ignored.  Remember that a letter like
       $x$                   % $ ... $  and  \( ... \)  are equivalent
is a formula when it denotes a mathematical
symbol, and it should be typed as one.

\section{Displayed Text}

Text is displayed by indenting it from the left
margin.  Quotations are commonly displayed.  There
are short quotations
\begin{quote}
   This is a short quotation.  It consists of a
   single paragraph of text.  See how it is formatted.
\end{quote}
and longer ones.
\begin{quotation}
   This is a longer quotation.  It consists of two
   paragraphs of text, neither of which are
   particularly interesting.

   This is the second paragraph of the quotation.  It
   is just as dull as the first paragraph.
\end{quotation}
Another frequently-displayed structure is a list.
The following is an example of an \emph{itemized}
list.
\begin{itemize}
   \item This is the first item of an itemized list.
         Each item in the list is marked with a ``tick''.
         You don't have to worry about what kind of tick
         mark is used.

   \item This is the second item of the list.  It
         contains another list nested inside it.  The inner
         list is an \emph{enumerated} list.
         \begin{enumerate}
            \item This is the first item of an enumerated
                  list that is nested within the itemized list.

            \item This is the second item of the inner list.
                  \LaTeX\ allows you to nest lists deeper than
                  you really should.
         \end{enumerate}
         This is the rest of the second item of the outer
         list.  It is no more interesting than any other
         part of the item.
   \item This is the third item of the list.
\end{itemize}
You can even display poetry.
\begin{verse}
   There is an environment
    for verse \\             % The \\ command separates lines
   Whose features some poets % within a stanza.
   will curse.

                             % One or more blank lines separate stanzas.

   For instead of making\\
   Them do \emph{all} line breaking, \\
   It allows them to put too many words on a line when they'd rather be
   forced to be terse.
\end{verse}

Mathematical formulas may also be displayed.  A
displayed formula
is
one-line long; multiline
formulas require special formatting instructions.
   \[  \ip{\Gamma}{\psi'} = x'' + y^{2} + z_{i}^{n}\]
Don't start a paragraph with a displayed equation,
nor make one a paragraph by itself.

\end{document}               % End of document.

\section{Bench Huge Section 367}
% synthetic bench section 367
% This is a sample LaTeX input file.  (Version of 12 August 2004.)
%
% A '%' character causes TeX to ignore all remaining text on the line,
% and is used for comments like this one.

\documentclass{article}      % Specifies the document class

                             % The preamble begins here.
\title{An Example Document}  % Declares the document's title.
\author{Leslie Lamport}      % Declares the author's name.
\date{January 21, 1994}      % Deleting this command produces today's date.

\newcommand{\ip}[2]{(#1, #2)}
                             % Defines \ip{arg1}{arg2} to mean
                             % (arg1, arg2).

%\newcommand{\ip}[2]{\langle #1 | #2\rangle}
                             % This is an alternative definition of
                             % \ip that is commented out.

\begin{document}             % End of preamble and beginning of text.

\maketitle                   % Produces the title.

This is an example input file.  Comparing it with
the output it generates can show you how to
produce a simple document of your own.

\section{Ordinary Text}      % Produces section heading.  Lower-level
                             % sections are begun with similar
                             % \subsection and \subsubsection commands.

The ends  of words and sentences are marked
  by   spaces. It  doesn't matter how many
spaces    you type; one is as good as 100.  The
end of   a line counts as a space.

One   or more   blank lines denote the  end
of  a paragraph.

Since any number of consecutive spaces are treated
like a single one, the formatting of the input
file makes no difference to
      \LaTeX,                % The \LaTeX command generates the LaTeX logo.
but it makes a difference to you.  When you use
\LaTeX, making your input file as easy to read
as possible will be a great help as you write
your document and when you change it.  This sample
file shows how you can add comments to your own input
file.

Because printing is different from typewriting,
there are a number of things that you have to do
differently when preparing an input file than if
you were just typing the document directly.
Quotation marks like
       ``this''
have to be handled specially, as do quotes within
quotes:
       ``\,`this'            % \, separates the double and single quote.
        is what I just
        wrote, not  `that'\,''.

Dashes come in three sizes: an
       intra-word
dash, a medium dash for number ranges like
       1--2,
and a punctuation
       dash---like
this.

A sentence-ending space should be larger than the
space between words within a sentence.  You
sometimes have to type special commands in
conjunction with punctuation characters to get
this right, as in the following sentence.
       Gnats, gnus, etc.\ all  % `\ ' makes an inter-word space.
       begin with G\@.         % \@ marks end-of-sentence punctuation.
You should check the spaces after periods when
reading your output to make sure you haven't
forgotten any special cases.  Generating an
ellipsis
       \ldots\               % `\ ' is needed after `\ldots' because TeX
                             % ignores spaces after command names like \ldots
                             % made from \ + letters.
                             %
                             % Note how a `%' character causes TeX to ignore
                             % the end of the input line, so these blank lines
                             % do not start a new paragraph.
                             %
with the right spacing around the periods requires
a special command.

\LaTeX\ interprets some common characters as
commands, so you must type special commands to
generate them.  These characters include the
following:
       \$ \& \% \# \{ and \}.

In printing, text is usually emphasized with an
       \emph{italic}
type style.

\begin{em}
   A long segment of text can also be emphasized
   in this way.  Text within such a segment can be
   given \emph{additional} emphasis.
\end{em}

It is sometimes necessary to prevent \LaTeX\ from
breaking a line where it might otherwise do so.
This may be at a space, as between the ``Mr.''\ and
``Jones'' in
       ``Mr.~Jones'',        % ~ produces an unbreakable interword space.
or within a word---especially when the word is a
symbol like
       \mbox{\emph{itemnum}}
that makes little sense when hyphenated across
lines.

Footnotes\footnote{This is an example of a footnote.}
pose no problem.

\LaTeX\ is good at typesetting mathematical formulas
like
       \( x-3y + z = 7 \)
or
       \( a_{1} > x^{2n} + y^{2n} > x' \)
or
       \( \ip{A}{B} = \sum_{i} a_{i} b_{i} \).
The spaces you type in a formula are
ignored.  Remember that a letter like
       $x$                   % $ ... $  and  \( ... \)  are equivalent
is a formula when it denotes a mathematical
symbol, and it should be typed as one.

\section{Displayed Text}

Text is displayed by indenting it from the left
margin.  Quotations are commonly displayed.  There
are short quotations
\begin{quote}
   This is a short quotation.  It consists of a
   single paragraph of text.  See how it is formatted.
\end{quote}
and longer ones.
\begin{quotation}
   This is a longer quotation.  It consists of two
   paragraphs of text, neither of which are
   particularly interesting.

   This is the second paragraph of the quotation.  It
   is just as dull as the first paragraph.
\end{quotation}
Another frequently-displayed structure is a list.
The following is an example of an \emph{itemized}
list.
\begin{itemize}
   \item This is the first item of an itemized list.
         Each item in the list is marked with a ``tick''.
         You don't have to worry about what kind of tick
         mark is used.

   \item This is the second item of the list.  It
         contains another list nested inside it.  The inner
         list is an \emph{enumerated} list.
         \begin{enumerate}
            \item This is the first item of an enumerated
                  list that is nested within the itemized list.

            \item This is the second item of the inner list.
                  \LaTeX\ allows you to nest lists deeper than
                  you really should.
         \end{enumerate}
         This is the rest of the second item of the outer
         list.  It is no more interesting than any other
         part of the item.
   \item This is the third item of the list.
\end{itemize}
You can even display poetry.
\begin{verse}
   There is an environment
    for verse \\             % The \\ command separates lines
   Whose features some poets % within a stanza.
   will curse.

                             % One or more blank lines separate stanzas.

   For instead of making\\
   Them do \emph{all} line breaking, \\
   It allows them to put too many words on a line when they'd rather be
   forced to be terse.
\end{verse}

Mathematical formulas may also be displayed.  A
displayed formula
is
one-line long; multiline
formulas require special formatting instructions.
   \[  \ip{\Gamma}{\psi'} = x'' + y^{2} + z_{i}^{n}\]
Don't start a paragraph with a displayed equation,
nor make one a paragraph by itself.

\end{document}               % End of document.

\section{Bench Huge Section 368}
% synthetic bench section 368
% This is a sample LaTeX input file.  (Version of 12 August 2004.)
%
% A '%' character causes TeX to ignore all remaining text on the line,
% and is used for comments like this one.

\documentclass{article}      % Specifies the document class

                             % The preamble begins here.
\title{An Example Document}  % Declares the document's title.
\author{Leslie Lamport}      % Declares the author's name.
\date{January 21, 1994}      % Deleting this command produces today's date.

\newcommand{\ip}[2]{(#1, #2)}
                             % Defines \ip{arg1}{arg2} to mean
                             % (arg1, arg2).

%\newcommand{\ip}[2]{\langle #1 | #2\rangle}
                             % This is an alternative definition of
                             % \ip that is commented out.

\begin{document}             % End of preamble and beginning of text.

\maketitle                   % Produces the title.

This is an example input file.  Comparing it with
the output it generates can show you how to
produce a simple document of your own.

\section{Ordinary Text}      % Produces section heading.  Lower-level
                             % sections are begun with similar
                             % \subsection and \subsubsection commands.

The ends  of words and sentences are marked
  by   spaces. It  doesn't matter how many
spaces    you type; one is as good as 100.  The
end of   a line counts as a space.

One   or more   blank lines denote the  end
of  a paragraph.

Since any number of consecutive spaces are treated
like a single one, the formatting of the input
file makes no difference to
      \LaTeX,                % The \LaTeX command generates the LaTeX logo.
but it makes a difference to you.  When you use
\LaTeX, making your input file as easy to read
as possible will be a great help as you write
your document and when you change it.  This sample
file shows how you can add comments to your own input
file.

Because printing is different from typewriting,
there are a number of things that you have to do
differently when preparing an input file than if
you were just typing the document directly.
Quotation marks like
       ``this''
have to be handled specially, as do quotes within
quotes:
       ``\,`this'            % \, separates the double and single quote.
        is what I just
        wrote, not  `that'\,''.

Dashes come in three sizes: an
       intra-word
dash, a medium dash for number ranges like
       1--2,
and a punctuation
       dash---like
this.

A sentence-ending space should be larger than the
space between words within a sentence.  You
sometimes have to type special commands in
conjunction with punctuation characters to get
this right, as in the following sentence.
       Gnats, gnus, etc.\ all  % `\ ' makes an inter-word space.
       begin with G\@.         % \@ marks end-of-sentence punctuation.
You should check the spaces after periods when
reading your output to make sure you haven't
forgotten any special cases.  Generating an
ellipsis
       \ldots\               % `\ ' is needed after `\ldots' because TeX
                             % ignores spaces after command names like \ldots
                             % made from \ + letters.
                             %
                             % Note how a `%' character causes TeX to ignore
                             % the end of the input line, so these blank lines
                             % do not start a new paragraph.
                             %
with the right spacing around the periods requires
a special command.

\LaTeX\ interprets some common characters as
commands, so you must type special commands to
generate them.  These characters include the
following:
       \$ \& \% \# \{ and \}.

In printing, text is usually emphasized with an
       \emph{italic}
type style.

\begin{em}
   A long segment of text can also be emphasized
   in this way.  Text within such a segment can be
   given \emph{additional} emphasis.
\end{em}

It is sometimes necessary to prevent \LaTeX\ from
breaking a line where it might otherwise do so.
This may be at a space, as between the ``Mr.''\ and
``Jones'' in
       ``Mr.~Jones'',        % ~ produces an unbreakable interword space.
or within a word---especially when the word is a
symbol like
       \mbox{\emph{itemnum}}
that makes little sense when hyphenated across
lines.

Footnotes\footnote{This is an example of a footnote.}
pose no problem.

\LaTeX\ is good at typesetting mathematical formulas
like
       \( x-3y + z = 7 \)
or
       \( a_{1} > x^{2n} + y^{2n} > x' \)
or
       \( \ip{A}{B} = \sum_{i} a_{i} b_{i} \).
The spaces you type in a formula are
ignored.  Remember that a letter like
       $x$                   % $ ... $  and  \( ... \)  are equivalent
is a formula when it denotes a mathematical
symbol, and it should be typed as one.

\section{Displayed Text}

Text is displayed by indenting it from the left
margin.  Quotations are commonly displayed.  There
are short quotations
\begin{quote}
   This is a short quotation.  It consists of a
   single paragraph of text.  See how it is formatted.
\end{quote}
and longer ones.
\begin{quotation}
   This is a longer quotation.  It consists of two
   paragraphs of text, neither of which are
   particularly interesting.

   This is the second paragraph of the quotation.  It
   is just as dull as the first paragraph.
\end{quotation}
Another frequently-displayed structure is a list.
The following is an example of an \emph{itemized}
list.
\begin{itemize}
   \item This is the first item of an itemized list.
         Each item in the list is marked with a ``tick''.
         You don't have to worry about what kind of tick
         mark is used.

   \item This is the second item of the list.  It
         contains another list nested inside it.  The inner
         list is an \emph{enumerated} list.
         \begin{enumerate}
            \item This is the first item of an enumerated
                  list that is nested within the itemized list.

            \item This is the second item of the inner list.
                  \LaTeX\ allows you to nest lists deeper than
                  you really should.
         \end{enumerate}
         This is the rest of the second item of the outer
         list.  It is no more interesting than any other
         part of the item.
   \item This is the third item of the list.
\end{itemize}
You can even display poetry.
\begin{verse}
   There is an environment
    for verse \\             % The \\ command separates lines
   Whose features some poets % within a stanza.
   will curse.

                             % One or more blank lines separate stanzas.

   For instead of making\\
   Them do \emph{all} line breaking, \\
   It allows them to put too many words on a line when they'd rather be
   forced to be terse.
\end{verse}

Mathematical formulas may also be displayed.  A
displayed formula
is
one-line long; multiline
formulas require special formatting instructions.
   \[  \ip{\Gamma}{\psi'} = x'' + y^{2} + z_{i}^{n}\]
Don't start a paragraph with a displayed equation,
nor make one a paragraph by itself.

\end{document}               % End of document.

\section{Bench Huge Section 369}
% synthetic bench section 369
% This is a sample LaTeX input file.  (Version of 12 August 2004.)
%
% A '%' character causes TeX to ignore all remaining text on the line,
% and is used for comments like this one.

\documentclass{article}      % Specifies the document class

                             % The preamble begins here.
\title{An Example Document}  % Declares the document's title.
\author{Leslie Lamport}      % Declares the author's name.
\date{January 21, 1994}      % Deleting this command produces today's date.

\newcommand{\ip}[2]{(#1, #2)}
                             % Defines \ip{arg1}{arg2} to mean
                             % (arg1, arg2).

%\newcommand{\ip}[2]{\langle #1 | #2\rangle}
                             % This is an alternative definition of
                             % \ip that is commented out.

\begin{document}             % End of preamble and beginning of text.

\maketitle                   % Produces the title.

This is an example input file.  Comparing it with
the output it generates can show you how to
produce a simple document of your own.

\section{Ordinary Text}      % Produces section heading.  Lower-level
                             % sections are begun with similar
                             % \subsection and \subsubsection commands.

The ends  of words and sentences are marked
  by   spaces. It  doesn't matter how many
spaces    you type; one is as good as 100.  The
end of   a line counts as a space.

One   or more   blank lines denote the  end
of  a paragraph.

Since any number of consecutive spaces are treated
like a single one, the formatting of the input
file makes no difference to
      \LaTeX,                % The \LaTeX command generates the LaTeX logo.
but it makes a difference to you.  When you use
\LaTeX, making your input file as easy to read
as possible will be a great help as you write
your document and when you change it.  This sample
file shows how you can add comments to your own input
file.

Because printing is different from typewriting,
there are a number of things that you have to do
differently when preparing an input file than if
you were just typing the document directly.
Quotation marks like
       ``this''
have to be handled specially, as do quotes within
quotes:
       ``\,`this'            % \, separates the double and single quote.
        is what I just
        wrote, not  `that'\,''.

Dashes come in three sizes: an
       intra-word
dash, a medium dash for number ranges like
       1--2,
and a punctuation
       dash---like
this.

A sentence-ending space should be larger than the
space between words within a sentence.  You
sometimes have to type special commands in
conjunction with punctuation characters to get
this right, as in the following sentence.
       Gnats, gnus, etc.\ all  % `\ ' makes an inter-word space.
       begin with G\@.         % \@ marks end-of-sentence punctuation.
You should check the spaces after periods when
reading your output to make sure you haven't
forgotten any special cases.  Generating an
ellipsis
       \ldots\               % `\ ' is needed after `\ldots' because TeX
                             % ignores spaces after command names like \ldots
                             % made from \ + letters.
                             %
                             % Note how a `%' character causes TeX to ignore
                             % the end of the input line, so these blank lines
                             % do not start a new paragraph.
                             %
with the right spacing around the periods requires
a special command.

\LaTeX\ interprets some common characters as
commands, so you must type special commands to
generate them.  These characters include the
following:
       \$ \& \% \# \{ and \}.

In printing, text is usually emphasized with an
       \emph{italic}
type style.

\begin{em}
   A long segment of text can also be emphasized
   in this way.  Text within such a segment can be
   given \emph{additional} emphasis.
\end{em}

It is sometimes necessary to prevent \LaTeX\ from
breaking a line where it might otherwise do so.
This may be at a space, as between the ``Mr.''\ and
``Jones'' in
       ``Mr.~Jones'',        % ~ produces an unbreakable interword space.
or within a word---especially when the word is a
symbol like
       \mbox{\emph{itemnum}}
that makes little sense when hyphenated across
lines.

Footnotes\footnote{This is an example of a footnote.}
pose no problem.

\LaTeX\ is good at typesetting mathematical formulas
like
       \( x-3y + z = 7 \)
or
       \( a_{1} > x^{2n} + y^{2n} > x' \)
or
       \( \ip{A}{B} = \sum_{i} a_{i} b_{i} \).
The spaces you type in a formula are
ignored.  Remember that a letter like
       $x$                   % $ ... $  and  \( ... \)  are equivalent
is a formula when it denotes a mathematical
symbol, and it should be typed as one.

\section{Displayed Text}

Text is displayed by indenting it from the left
margin.  Quotations are commonly displayed.  There
are short quotations
\begin{quote}
   This is a short quotation.  It consists of a
   single paragraph of text.  See how it is formatted.
\end{quote}
and longer ones.
\begin{quotation}
   This is a longer quotation.  It consists of two
   paragraphs of text, neither of which are
   particularly interesting.

   This is the second paragraph of the quotation.  It
   is just as dull as the first paragraph.
\end{quotation}
Another frequently-displayed structure is a list.
The following is an example of an \emph{itemized}
list.
\begin{itemize}
   \item This is the first item of an itemized list.
         Each item in the list is marked with a ``tick''.
         You don't have to worry about what kind of tick
         mark is used.

   \item This is the second item of the list.  It
         contains another list nested inside it.  The inner
         list is an \emph{enumerated} list.
         \begin{enumerate}
            \item This is the first item of an enumerated
                  list that is nested within the itemized list.

            \item This is the second item of the inner list.
                  \LaTeX\ allows you to nest lists deeper than
                  you really should.
         \end{enumerate}
         This is the rest of the second item of the outer
         list.  It is no more interesting than any other
         part of the item.
   \item This is the third item of the list.
\end{itemize}
You can even display poetry.
\begin{verse}
   There is an environment
    for verse \\             % The \\ command separates lines
   Whose features some poets % within a stanza.
   will curse.

                             % One or more blank lines separate stanzas.

   For instead of making\\
   Them do \emph{all} line breaking, \\
   It allows them to put too many words on a line when they'd rather be
   forced to be terse.
\end{verse}

Mathematical formulas may also be displayed.  A
displayed formula
is
one-line long; multiline
formulas require special formatting instructions.
   \[  \ip{\Gamma}{\psi'} = x'' + y^{2} + z_{i}^{n}\]
Don't start a paragraph with a displayed equation,
nor make one a paragraph by itself.

\end{document}               % End of document.

\section{Bench Huge Section 370}
% synthetic bench section 370
% This is a sample LaTeX input file.  (Version of 12 August 2004.)
%
% A '%' character causes TeX to ignore all remaining text on the line,
% and is used for comments like this one.

\documentclass{article}      % Specifies the document class

                             % The preamble begins here.
\title{An Example Document}  % Declares the document's title.
\author{Leslie Lamport}      % Declares the author's name.
\date{January 21, 1994}      % Deleting this command produces today's date.

\newcommand{\ip}[2]{(#1, #2)}
                             % Defines \ip{arg1}{arg2} to mean
                             % (arg1, arg2).

%\newcommand{\ip}[2]{\langle #1 | #2\rangle}
                             % This is an alternative definition of
                             % \ip that is commented out.

\begin{document}             % End of preamble and beginning of text.

\maketitle                   % Produces the title.

This is an example input file.  Comparing it with
the output it generates can show you how to
produce a simple document of your own.

\section{Ordinary Text}      % Produces section heading.  Lower-level
                             % sections are begun with similar
                             % \subsection and \subsubsection commands.

The ends  of words and sentences are marked
  by   spaces. It  doesn't matter how many
spaces    you type; one is as good as 100.  The
end of   a line counts as a space.

One   or more   blank lines denote the  end
of  a paragraph.

Since any number of consecutive spaces are treated
like a single one, the formatting of the input
file makes no difference to
      \LaTeX,                % The \LaTeX command generates the LaTeX logo.
but it makes a difference to you.  When you use
\LaTeX, making your input file as easy to read
as possible will be a great help as you write
your document and when you change it.  This sample
file shows how you can add comments to your own input
file.

Because printing is different from typewriting,
there are a number of things that you have to do
differently when preparing an input file than if
you were just typing the document directly.
Quotation marks like
       ``this''
have to be handled specially, as do quotes within
quotes:
       ``\,`this'            % \, separates the double and single quote.
        is what I just
        wrote, not  `that'\,''.

Dashes come in three sizes: an
       intra-word
dash, a medium dash for number ranges like
       1--2,
and a punctuation
       dash---like
this.

A sentence-ending space should be larger than the
space between words within a sentence.  You
sometimes have to type special commands in
conjunction with punctuation characters to get
this right, as in the following sentence.
       Gnats, gnus, etc.\ all  % `\ ' makes an inter-word space.
       begin with G\@.         % \@ marks end-of-sentence punctuation.
You should check the spaces after periods when
reading your output to make sure you haven't
forgotten any special cases.  Generating an
ellipsis
       \ldots\               % `\ ' is needed after `\ldots' because TeX
                             % ignores spaces after command names like \ldots
                             % made from \ + letters.
                             %
                             % Note how a `%' character causes TeX to ignore
                             % the end of the input line, so these blank lines
                             % do not start a new paragraph.
                             %
with the right spacing around the periods requires
a special command.

\LaTeX\ interprets some common characters as
commands, so you must type special commands to
generate them.  These characters include the
following:
       \$ \& \% \# \{ and \}.

In printing, text is usually emphasized with an
       \emph{italic}
type style.

\begin{em}
   A long segment of text can also be emphasized
   in this way.  Text within such a segment can be
   given \emph{additional} emphasis.
\end{em}

It is sometimes necessary to prevent \LaTeX\ from
breaking a line where it might otherwise do so.
This may be at a space, as between the ``Mr.''\ and
``Jones'' in
       ``Mr.~Jones'',        % ~ produces an unbreakable interword space.
or within a word---especially when the word is a
symbol like
       \mbox{\emph{itemnum}}
that makes little sense when hyphenated across
lines.

Footnotes\footnote{This is an example of a footnote.}
pose no problem.

\LaTeX\ is good at typesetting mathematical formulas
like
       \( x-3y + z = 7 \)
or
       \( a_{1} > x^{2n} + y^{2n} > x' \)
or
       \( \ip{A}{B} = \sum_{i} a_{i} b_{i} \).
The spaces you type in a formula are
ignored.  Remember that a letter like
       $x$                   % $ ... $  and  \( ... \)  are equivalent
is a formula when it denotes a mathematical
symbol, and it should be typed as one.

\section{Displayed Text}

Text is displayed by indenting it from the left
margin.  Quotations are commonly displayed.  There
are short quotations
\begin{quote}
   This is a short quotation.  It consists of a
   single paragraph of text.  See how it is formatted.
\end{quote}
and longer ones.
\begin{quotation}
   This is a longer quotation.  It consists of two
   paragraphs of text, neither of which are
   particularly interesting.

   This is the second paragraph of the quotation.  It
   is just as dull as the first paragraph.
\end{quotation}
Another frequently-displayed structure is a list.
The following is an example of an \emph{itemized}
list.
\begin{itemize}
   \item This is the first item of an itemized list.
         Each item in the list is marked with a ``tick''.
         You don't have to worry about what kind of tick
         mark is used.

   \item This is the second item of the list.  It
         contains another list nested inside it.  The inner
         list is an \emph{enumerated} list.
         \begin{enumerate}
            \item This is the first item of an enumerated
                  list that is nested within the itemized list.

            \item This is the second item of the inner list.
                  \LaTeX\ allows you to nest lists deeper than
                  you really should.
         \end{enumerate}
         This is the rest of the second item of the outer
         list.  It is no more interesting than any other
         part of the item.
   \item This is the third item of the list.
\end{itemize}
You can even display poetry.
\begin{verse}
   There is an environment
    for verse \\             % The \\ command separates lines
   Whose features some poets % within a stanza.
   will curse.

                             % One or more blank lines separate stanzas.

   For instead of making\\
   Them do \emph{all} line breaking, \\
   It allows them to put too many words on a line when they'd rather be
   forced to be terse.
\end{verse}

Mathematical formulas may also be displayed.  A
displayed formula
is
one-line long; multiline
formulas require special formatting instructions.
   \[  \ip{\Gamma}{\psi'} = x'' + y^{2} + z_{i}^{n}\]
Don't start a paragraph with a displayed equation,
nor make one a paragraph by itself.

\end{document}               % End of document.

\section{Bench Huge Section 371}
% synthetic bench section 371
% This is a sample LaTeX input file.  (Version of 12 August 2004.)
%
% A '%' character causes TeX to ignore all remaining text on the line,
% and is used for comments like this one.

\documentclass{article}      % Specifies the document class

                             % The preamble begins here.
\title{An Example Document}  % Declares the document's title.
\author{Leslie Lamport}      % Declares the author's name.
\date{January 21, 1994}      % Deleting this command produces today's date.

\newcommand{\ip}[2]{(#1, #2)}
                             % Defines \ip{arg1}{arg2} to mean
                             % (arg1, arg2).

%\newcommand{\ip}[2]{\langle #1 | #2\rangle}
                             % This is an alternative definition of
                             % \ip that is commented out.

\begin{document}             % End of preamble and beginning of text.

\maketitle                   % Produces the title.

This is an example input file.  Comparing it with
the output it generates can show you how to
produce a simple document of your own.

\section{Ordinary Text}      % Produces section heading.  Lower-level
                             % sections are begun with similar
                             % \subsection and \subsubsection commands.

The ends  of words and sentences are marked
  by   spaces. It  doesn't matter how many
spaces    you type; one is as good as 100.  The
end of   a line counts as a space.

One   or more   blank lines denote the  end
of  a paragraph.

Since any number of consecutive spaces are treated
like a single one, the formatting of the input
file makes no difference to
      \LaTeX,                % The \LaTeX command generates the LaTeX logo.
but it makes a difference to you.  When you use
\LaTeX, making your input file as easy to read
as possible will be a great help as you write
your document and when you change it.  This sample
file shows how you can add comments to your own input
file.

Because printing is different from typewriting,
there are a number of things that you have to do
differently when preparing an input file than if
you were just typing the document directly.
Quotation marks like
       ``this''
have to be handled specially, as do quotes within
quotes:
       ``\,`this'            % \, separates the double and single quote.
        is what I just
        wrote, not  `that'\,''.

Dashes come in three sizes: an
       intra-word
dash, a medium dash for number ranges like
       1--2,
and a punctuation
       dash---like
this.

A sentence-ending space should be larger than the
space between words within a sentence.  You
sometimes have to type special commands in
conjunction with punctuation characters to get
this right, as in the following sentence.
       Gnats, gnus, etc.\ all  % `\ ' makes an inter-word space.
       begin with G\@.         % \@ marks end-of-sentence punctuation.
You should check the spaces after periods when
reading your output to make sure you haven't
forgotten any special cases.  Generating an
ellipsis
       \ldots\               % `\ ' is needed after `\ldots' because TeX
                             % ignores spaces after command names like \ldots
                             % made from \ + letters.
                             %
                             % Note how a `%' character causes TeX to ignore
                             % the end of the input line, so these blank lines
                             % do not start a new paragraph.
                             %
with the right spacing around the periods requires
a special command.

\LaTeX\ interprets some common characters as
commands, so you must type special commands to
generate them.  These characters include the
following:
       \$ \& \% \# \{ and \}.

In printing, text is usually emphasized with an
       \emph{italic}
type style.

\begin{em}
   A long segment of text can also be emphasized
   in this way.  Text within such a segment can be
   given \emph{additional} emphasis.
\end{em}

It is sometimes necessary to prevent \LaTeX\ from
breaking a line where it might otherwise do so.
This may be at a space, as between the ``Mr.''\ and
``Jones'' in
       ``Mr.~Jones'',        % ~ produces an unbreakable interword space.
or within a word---especially when the word is a
symbol like
       \mbox{\emph{itemnum}}
that makes little sense when hyphenated across
lines.

Footnotes\footnote{This is an example of a footnote.}
pose no problem.

\LaTeX\ is good at typesetting mathematical formulas
like
       \( x-3y + z = 7 \)
or
       \( a_{1} > x^{2n} + y^{2n} > x' \)
or
       \( \ip{A}{B} = \sum_{i} a_{i} b_{i} \).
The spaces you type in a formula are
ignored.  Remember that a letter like
       $x$                   % $ ... $  and  \( ... \)  are equivalent
is a formula when it denotes a mathematical
symbol, and it should be typed as one.

\section{Displayed Text}

Text is displayed by indenting it from the left
margin.  Quotations are commonly displayed.  There
are short quotations
\begin{quote}
   This is a short quotation.  It consists of a
   single paragraph of text.  See how it is formatted.
\end{quote}
and longer ones.
\begin{quotation}
   This is a longer quotation.  It consists of two
   paragraphs of text, neither of which are
   particularly interesting.

   This is the second paragraph of the quotation.  It
   is just as dull as the first paragraph.
\end{quotation}
Another frequently-displayed structure is a list.
The following is an example of an \emph{itemized}
list.
\begin{itemize}
   \item This is the first item of an itemized list.
         Each item in the list is marked with a ``tick''.
         You don't have to worry about what kind of tick
         mark is used.

   \item This is the second item of the list.  It
         contains another list nested inside it.  The inner
         list is an \emph{enumerated} list.
         \begin{enumerate}
            \item This is the first item of an enumerated
                  list that is nested within the itemized list.

            \item This is the second item of the inner list.
                  \LaTeX\ allows you to nest lists deeper than
                  you really should.
         \end{enumerate}
         This is the rest of the second item of the outer
         list.  It is no more interesting than any other
         part of the item.
   \item This is the third item of the list.
\end{itemize}
You can even display poetry.
\begin{verse}
   There is an environment
    for verse \\             % The \\ command separates lines
   Whose features some poets % within a stanza.
   will curse.

                             % One or more blank lines separate stanzas.

   For instead of making\\
   Them do \emph{all} line breaking, \\
   It allows them to put too many words on a line when they'd rather be
   forced to be terse.
\end{verse}

Mathematical formulas may also be displayed.  A
displayed formula
is
one-line long; multiline
formulas require special formatting instructions.
   \[  \ip{\Gamma}{\psi'} = x'' + y^{2} + z_{i}^{n}\]
Don't start a paragraph with a displayed equation,
nor make one a paragraph by itself.

\end{document}               % End of document.

\section{Bench Huge Section 372}
% synthetic bench section 372
% This is a sample LaTeX input file.  (Version of 12 August 2004.)
%
% A '%' character causes TeX to ignore all remaining text on the line,
% and is used for comments like this one.

\documentclass{article}      % Specifies the document class

                             % The preamble begins here.
\title{An Example Document}  % Declares the document's title.
\author{Leslie Lamport}      % Declares the author's name.
\date{January 21, 1994}      % Deleting this command produces today's date.

\newcommand{\ip}[2]{(#1, #2)}
                             % Defines \ip{arg1}{arg2} to mean
                             % (arg1, arg2).

%\newcommand{\ip}[2]{\langle #1 | #2\rangle}
                             % This is an alternative definition of
                             % \ip that is commented out.

\begin{document}             % End of preamble and beginning of text.

\maketitle                   % Produces the title.

This is an example input file.  Comparing it with
the output it generates can show you how to
produce a simple document of your own.

\section{Ordinary Text}      % Produces section heading.  Lower-level
                             % sections are begun with similar
                             % \subsection and \subsubsection commands.

The ends  of words and sentences are marked
  by   spaces. It  doesn't matter how many
spaces    you type; one is as good as 100.  The
end of   a line counts as a space.

One   or more   blank lines denote the  end
of  a paragraph.

Since any number of consecutive spaces are treated
like a single one, the formatting of the input
file makes no difference to
      \LaTeX,                % The \LaTeX command generates the LaTeX logo.
but it makes a difference to you.  When you use
\LaTeX, making your input file as easy to read
as possible will be a great help as you write
your document and when you change it.  This sample
file shows how you can add comments to your own input
file.

Because printing is different from typewriting,
there are a number of things that you have to do
differently when preparing an input file than if
you were just typing the document directly.
Quotation marks like
       ``this''
have to be handled specially, as do quotes within
quotes:
       ``\,`this'            % \, separates the double and single quote.
        is what I just
        wrote, not  `that'\,''.

Dashes come in three sizes: an
       intra-word
dash, a medium dash for number ranges like
       1--2,
and a punctuation
       dash---like
this.

A sentence-ending space should be larger than the
space between words within a sentence.  You
sometimes have to type special commands in
conjunction with punctuation characters to get
this right, as in the following sentence.
       Gnats, gnus, etc.\ all  % `\ ' makes an inter-word space.
       begin with G\@.         % \@ marks end-of-sentence punctuation.
You should check the spaces after periods when
reading your output to make sure you haven't
forgotten any special cases.  Generating an
ellipsis
       \ldots\               % `\ ' is needed after `\ldots' because TeX
                             % ignores spaces after command names like \ldots
                             % made from \ + letters.
                             %
                             % Note how a `%' character causes TeX to ignore
                             % the end of the input line, so these blank lines
                             % do not start a new paragraph.
                             %
with the right spacing around the periods requires
a special command.

\LaTeX\ interprets some common characters as
commands, so you must type special commands to
generate them.  These characters include the
following:
       \$ \& \% \# \{ and \}.

In printing, text is usually emphasized with an
       \emph{italic}
type style.

\begin{em}
   A long segment of text can also be emphasized
   in this way.  Text within such a segment can be
   given \emph{additional} emphasis.
\end{em}

It is sometimes necessary to prevent \LaTeX\ from
breaking a line where it might otherwise do so.
This may be at a space, as between the ``Mr.''\ and
``Jones'' in
       ``Mr.~Jones'',        % ~ produces an unbreakable interword space.
or within a word---especially when the word is a
symbol like
       \mbox{\emph{itemnum}}
that makes little sense when hyphenated across
lines.

Footnotes\footnote{This is an example of a footnote.}
pose no problem.

\LaTeX\ is good at typesetting mathematical formulas
like
       \( x-3y + z = 7 \)
or
       \( a_{1} > x^{2n} + y^{2n} > x' \)
or
       \( \ip{A}{B} = \sum_{i} a_{i} b_{i} \).
The spaces you type in a formula are
ignored.  Remember that a letter like
       $x$                   % $ ... $  and  \( ... \)  are equivalent
is a formula when it denotes a mathematical
symbol, and it should be typed as one.

\section{Displayed Text}

Text is displayed by indenting it from the left
margin.  Quotations are commonly displayed.  There
are short quotations
\begin{quote}
   This is a short quotation.  It consists of a
   single paragraph of text.  See how it is formatted.
\end{quote}
and longer ones.
\begin{quotation}
   This is a longer quotation.  It consists of two
   paragraphs of text, neither of which are
   particularly interesting.

   This is the second paragraph of the quotation.  It
   is just as dull as the first paragraph.
\end{quotation}
Another frequently-displayed structure is a list.
The following is an example of an \emph{itemized}
list.
\begin{itemize}
   \item This is the first item of an itemized list.
         Each item in the list is marked with a ``tick''.
         You don't have to worry about what kind of tick
         mark is used.

   \item This is the second item of the list.  It
         contains another list nested inside it.  The inner
         list is an \emph{enumerated} list.
         \begin{enumerate}
            \item This is the first item of an enumerated
                  list that is nested within the itemized list.

            \item This is the second item of the inner list.
                  \LaTeX\ allows you to nest lists deeper than
                  you really should.
         \end{enumerate}
         This is the rest of the second item of the outer
         list.  It is no more interesting than any other
         part of the item.
   \item This is the third item of the list.
\end{itemize}
You can even display poetry.
\begin{verse}
   There is an environment
    for verse \\             % The \\ command separates lines
   Whose features some poets % within a stanza.
   will curse.

                             % One or more blank lines separate stanzas.

   For instead of making\\
   Them do \emph{all} line breaking, \\
   It allows them to put too many words on a line when they'd rather be
   forced to be terse.
\end{verse}

Mathematical formulas may also be displayed.  A
displayed formula
is
one-line long; multiline
formulas require special formatting instructions.
   \[  \ip{\Gamma}{\psi'} = x'' + y^{2} + z_{i}^{n}\]
Don't start a paragraph with a displayed equation,
nor make one a paragraph by itself.

\end{document}               % End of document.

\section{Bench Huge Section 373}
% synthetic bench section 373
% This is a sample LaTeX input file.  (Version of 12 August 2004.)
%
% A '%' character causes TeX to ignore all remaining text on the line,
% and is used for comments like this one.

\documentclass{article}      % Specifies the document class

                             % The preamble begins here.
\title{An Example Document}  % Declares the document's title.
\author{Leslie Lamport}      % Declares the author's name.
\date{January 21, 1994}      % Deleting this command produces today's date.

\newcommand{\ip}[2]{(#1, #2)}
                             % Defines \ip{arg1}{arg2} to mean
                             % (arg1, arg2).

%\newcommand{\ip}[2]{\langle #1 | #2\rangle}
                             % This is an alternative definition of
                             % \ip that is commented out.

\begin{document}             % End of preamble and beginning of text.

\maketitle                   % Produces the title.

This is an example input file.  Comparing it with
the output it generates can show you how to
produce a simple document of your own.

\section{Ordinary Text}      % Produces section heading.  Lower-level
                             % sections are begun with similar
                             % \subsection and \subsubsection commands.

The ends  of words and sentences are marked
  by   spaces. It  doesn't matter how many
spaces    you type; one is as good as 100.  The
end of   a line counts as a space.

One   or more   blank lines denote the  end
of  a paragraph.

Since any number of consecutive spaces are treated
like a single one, the formatting of the input
file makes no difference to
      \LaTeX,                % The \LaTeX command generates the LaTeX logo.
but it makes a difference to you.  When you use
\LaTeX, making your input file as easy to read
as possible will be a great help as you write
your document and when you change it.  This sample
file shows how you can add comments to your own input
file.

Because printing is different from typewriting,
there are a number of things that you have to do
differently when preparing an input file than if
you were just typing the document directly.
Quotation marks like
       ``this''
have to be handled specially, as do quotes within
quotes:
       ``\,`this'            % \, separates the double and single quote.
        is what I just
        wrote, not  `that'\,''.

Dashes come in three sizes: an
       intra-word
dash, a medium dash for number ranges like
       1--2,
and a punctuation
       dash---like
this.

A sentence-ending space should be larger than the
space between words within a sentence.  You
sometimes have to type special commands in
conjunction with punctuation characters to get
this right, as in the following sentence.
       Gnats, gnus, etc.\ all  % `\ ' makes an inter-word space.
       begin with G\@.         % \@ marks end-of-sentence punctuation.
You should check the spaces after periods when
reading your output to make sure you haven't
forgotten any special cases.  Generating an
ellipsis
       \ldots\               % `\ ' is needed after `\ldots' because TeX
                             % ignores spaces after command names like \ldots
                             % made from \ + letters.
                             %
                             % Note how a `%' character causes TeX to ignore
                             % the end of the input line, so these blank lines
                             % do not start a new paragraph.
                             %
with the right spacing around the periods requires
a special command.

\LaTeX\ interprets some common characters as
commands, so you must type special commands to
generate them.  These characters include the
following:
       \$ \& \% \# \{ and \}.

In printing, text is usually emphasized with an
       \emph{italic}
type style.

\begin{em}
   A long segment of text can also be emphasized
   in this way.  Text within such a segment can be
   given \emph{additional} emphasis.
\end{em}

It is sometimes necessary to prevent \LaTeX\ from
breaking a line where it might otherwise do so.
This may be at a space, as between the ``Mr.''\ and
``Jones'' in
       ``Mr.~Jones'',        % ~ produces an unbreakable interword space.
or within a word---especially when the word is a
symbol like
       \mbox{\emph{itemnum}}
that makes little sense when hyphenated across
lines.

Footnotes\footnote{This is an example of a footnote.}
pose no problem.

\LaTeX\ is good at typesetting mathematical formulas
like
       \( x-3y + z = 7 \)
or
       \( a_{1} > x^{2n} + y^{2n} > x' \)
or
       \( \ip{A}{B} = \sum_{i} a_{i} b_{i} \).
The spaces you type in a formula are
ignored.  Remember that a letter like
       $x$                   % $ ... $  and  \( ... \)  are equivalent
is a formula when it denotes a mathematical
symbol, and it should be typed as one.

\section{Displayed Text}

Text is displayed by indenting it from the left
margin.  Quotations are commonly displayed.  There
are short quotations
\begin{quote}
   This is a short quotation.  It consists of a
   single paragraph of text.  See how it is formatted.
\end{quote}
and longer ones.
\begin{quotation}
   This is a longer quotation.  It consists of two
   paragraphs of text, neither of which are
   particularly interesting.

   This is the second paragraph of the quotation.  It
   is just as dull as the first paragraph.
\end{quotation}
Another frequently-displayed structure is a list.
The following is an example of an \emph{itemized}
list.
\begin{itemize}
   \item This is the first item of an itemized list.
         Each item in the list is marked with a ``tick''.
         You don't have to worry about what kind of tick
         mark is used.

   \item This is the second item of the list.  It
         contains another list nested inside it.  The inner
         list is an \emph{enumerated} list.
         \begin{enumerate}
            \item This is the first item of an enumerated
                  list that is nested within the itemized list.

            \item This is the second item of the inner list.
                  \LaTeX\ allows you to nest lists deeper than
                  you really should.
         \end{enumerate}
         This is the rest of the second item of the outer
         list.  It is no more interesting than any other
         part of the item.
   \item This is the third item of the list.
\end{itemize}
You can even display poetry.
\begin{verse}
   There is an environment
    for verse \\             % The \\ command separates lines
   Whose features some poets % within a stanza.
   will curse.

                             % One or more blank lines separate stanzas.

   For instead of making\\
   Them do \emph{all} line breaking, \\
   It allows them to put too many words on a line when they'd rather be
   forced to be terse.
\end{verse}

Mathematical formulas may also be displayed.  A
displayed formula
is
one-line long; multiline
formulas require special formatting instructions.
   \[  \ip{\Gamma}{\psi'} = x'' + y^{2} + z_{i}^{n}\]
Don't start a paragraph with a displayed equation,
nor make one a paragraph by itself.

\end{document}               % End of document.

\section{Bench Huge Section 374}
% synthetic bench section 374
% This is a sample LaTeX input file.  (Version of 12 August 2004.)
%
% A '%' character causes TeX to ignore all remaining text on the line,
% and is used for comments like this one.

\documentclass{article}      % Specifies the document class

                             % The preamble begins here.
\title{An Example Document}  % Declares the document's title.
\author{Leslie Lamport}      % Declares the author's name.
\date{January 21, 1994}      % Deleting this command produces today's date.

\newcommand{\ip}[2]{(#1, #2)}
                             % Defines \ip{arg1}{arg2} to mean
                             % (arg1, arg2).

%\newcommand{\ip}[2]{\langle #1 | #2\rangle}
                             % This is an alternative definition of
                             % \ip that is commented out.

\begin{document}             % End of preamble and beginning of text.

\maketitle                   % Produces the title.

This is an example input file.  Comparing it with
the output it generates can show you how to
produce a simple document of your own.

\section{Ordinary Text}      % Produces section heading.  Lower-level
                             % sections are begun with similar
                             % \subsection and \subsubsection commands.

The ends  of words and sentences are marked
  by   spaces. It  doesn't matter how many
spaces    you type; one is as good as 100.  The
end of   a line counts as a space.

One   or more   blank lines denote the  end
of  a paragraph.

Since any number of consecutive spaces are treated
like a single one, the formatting of the input
file makes no difference to
      \LaTeX,                % The \LaTeX command generates the LaTeX logo.
but it makes a difference to you.  When you use
\LaTeX, making your input file as easy to read
as possible will be a great help as you write
your document and when you change it.  This sample
file shows how you can add comments to your own input
file.

Because printing is different from typewriting,
there are a number of things that you have to do
differently when preparing an input file than if
you were just typing the document directly.
Quotation marks like
       ``this''
have to be handled specially, as do quotes within
quotes:
       ``\,`this'            % \, separates the double and single quote.
        is what I just
        wrote, not  `that'\,''.

Dashes come in three sizes: an
       intra-word
dash, a medium dash for number ranges like
       1--2,
and a punctuation
       dash---like
this.

A sentence-ending space should be larger than the
space between words within a sentence.  You
sometimes have to type special commands in
conjunction with punctuation characters to get
this right, as in the following sentence.
       Gnats, gnus, etc.\ all  % `\ ' makes an inter-word space.
       begin with G\@.         % \@ marks end-of-sentence punctuation.
You should check the spaces after periods when
reading your output to make sure you haven't
forgotten any special cases.  Generating an
ellipsis
       \ldots\               % `\ ' is needed after `\ldots' because TeX
                             % ignores spaces after command names like \ldots
                             % made from \ + letters.
                             %
                             % Note how a `%' character causes TeX to ignore
                             % the end of the input line, so these blank lines
                             % do not start a new paragraph.
                             %
with the right spacing around the periods requires
a special command.

\LaTeX\ interprets some common characters as
commands, so you must type special commands to
generate them.  These characters include the
following:
       \$ \& \% \# \{ and \}.

In printing, text is usually emphasized with an
       \emph{italic}
type style.

\begin{em}
   A long segment of text can also be emphasized
   in this way.  Text within such a segment can be
   given \emph{additional} emphasis.
\end{em}

It is sometimes necessary to prevent \LaTeX\ from
breaking a line where it might otherwise do so.
This may be at a space, as between the ``Mr.''\ and
``Jones'' in
       ``Mr.~Jones'',        % ~ produces an unbreakable interword space.
or within a word---especially when the word is a
symbol like
       \mbox{\emph{itemnum}}
that makes little sense when hyphenated across
lines.

Footnotes\footnote{This is an example of a footnote.}
pose no problem.

\LaTeX\ is good at typesetting mathematical formulas
like
       \( x-3y + z = 7 \)
or
       \( a_{1} > x^{2n} + y^{2n} > x' \)
or
       \( \ip{A}{B} = \sum_{i} a_{i} b_{i} \).
The spaces you type in a formula are
ignored.  Remember that a letter like
       $x$                   % $ ... $  and  \( ... \)  are equivalent
is a formula when it denotes a mathematical
symbol, and it should be typed as one.

\section{Displayed Text}

Text is displayed by indenting it from the left
margin.  Quotations are commonly displayed.  There
are short quotations
\begin{quote}
   This is a short quotation.  It consists of a
   single paragraph of text.  See how it is formatted.
\end{quote}
and longer ones.
\begin{quotation}
   This is a longer quotation.  It consists of two
   paragraphs of text, neither of which are
   particularly interesting.

   This is the second paragraph of the quotation.  It
   is just as dull as the first paragraph.
\end{quotation}
Another frequently-displayed structure is a list.
The following is an example of an \emph{itemized}
list.
\begin{itemize}
   \item This is the first item of an itemized list.
         Each item in the list is marked with a ``tick''.
         You don't have to worry about what kind of tick
         mark is used.

   \item This is the second item of the list.  It
         contains another list nested inside it.  The inner
         list is an \emph{enumerated} list.
         \begin{enumerate}
            \item This is the first item of an enumerated
                  list that is nested within the itemized list.

            \item This is the second item of the inner list.
                  \LaTeX\ allows you to nest lists deeper than
                  you really should.
         \end{enumerate}
         This is the rest of the second item of the outer
         list.  It is no more interesting than any other
         part of the item.
   \item This is the third item of the list.
\end{itemize}
You can even display poetry.
\begin{verse}
   There is an environment
    for verse \\             % The \\ command separates lines
   Whose features some poets % within a stanza.
   will curse.

                             % One or more blank lines separate stanzas.

   For instead of making\\
   Them do \emph{all} line breaking, \\
   It allows them to put too many words on a line when they'd rather be
   forced to be terse.
\end{verse}

Mathematical formulas may also be displayed.  A
displayed formula
is
one-line long; multiline
formulas require special formatting instructions.
   \[  \ip{\Gamma}{\psi'} = x'' + y^{2} + z_{i}^{n}\]
Don't start a paragraph with a displayed equation,
nor make one a paragraph by itself.

\end{document}               % End of document.

\section{Bench Huge Section 375}
% synthetic bench section 375
% This is a sample LaTeX input file.  (Version of 12 August 2004.)
%
% A '%' character causes TeX to ignore all remaining text on the line,
% and is used for comments like this one.

\documentclass{article}      % Specifies the document class

                             % The preamble begins here.
\title{An Example Document}  % Declares the document's title.
\author{Leslie Lamport}      % Declares the author's name.
\date{January 21, 1994}      % Deleting this command produces today's date.

\newcommand{\ip}[2]{(#1, #2)}
                             % Defines \ip{arg1}{arg2} to mean
                             % (arg1, arg2).

%\newcommand{\ip}[2]{\langle #1 | #2\rangle}
                             % This is an alternative definition of
                             % \ip that is commented out.

\begin{document}             % End of preamble and beginning of text.

\maketitle                   % Produces the title.

This is an example input file.  Comparing it with
the output it generates can show you how to
produce a simple document of your own.

\section{Ordinary Text}      % Produces section heading.  Lower-level
                             % sections are begun with similar
                             % \subsection and \subsubsection commands.

The ends  of words and sentences are marked
  by   spaces. It  doesn't matter how many
spaces    you type; one is as good as 100.  The
end of   a line counts as a space.

One   or more   blank lines denote the  end
of  a paragraph.

Since any number of consecutive spaces are treated
like a single one, the formatting of the input
file makes no difference to
      \LaTeX,                % The \LaTeX command generates the LaTeX logo.
but it makes a difference to you.  When you use
\LaTeX, making your input file as easy to read
as possible will be a great help as you write
your document and when you change it.  This sample
file shows how you can add comments to your own input
file.

Because printing is different from typewriting,
there are a number of things that you have to do
differently when preparing an input file than if
you were just typing the document directly.
Quotation marks like
       ``this''
have to be handled specially, as do quotes within
quotes:
       ``\,`this'            % \, separates the double and single quote.
        is what I just
        wrote, not  `that'\,''.

Dashes come in three sizes: an
       intra-word
dash, a medium dash for number ranges like
       1--2,
and a punctuation
       dash---like
this.

A sentence-ending space should be larger than the
space between words within a sentence.  You
sometimes have to type special commands in
conjunction with punctuation characters to get
this right, as in the following sentence.
       Gnats, gnus, etc.\ all  % `\ ' makes an inter-word space.
       begin with G\@.         % \@ marks end-of-sentence punctuation.
You should check the spaces after periods when
reading your output to make sure you haven't
forgotten any special cases.  Generating an
ellipsis
       \ldots\               % `\ ' is needed after `\ldots' because TeX
                             % ignores spaces after command names like \ldots
                             % made from \ + letters.
                             %
                             % Note how a `%' character causes TeX to ignore
                             % the end of the input line, so these blank lines
                             % do not start a new paragraph.
                             %
with the right spacing around the periods requires
a special command.

\LaTeX\ interprets some common characters as
commands, so you must type special commands to
generate them.  These characters include the
following:
       \$ \& \% \# \{ and \}.

In printing, text is usually emphasized with an
       \emph{italic}
type style.

\begin{em}
   A long segment of text can also be emphasized
   in this way.  Text within such a segment can be
   given \emph{additional} emphasis.
\end{em}

It is sometimes necessary to prevent \LaTeX\ from
breaking a line where it might otherwise do so.
This may be at a space, as between the ``Mr.''\ and
``Jones'' in
       ``Mr.~Jones'',        % ~ produces an unbreakable interword space.
or within a word---especially when the word is a
symbol like
       \mbox{\emph{itemnum}}
that makes little sense when hyphenated across
lines.

Footnotes\footnote{This is an example of a footnote.}
pose no problem.

\LaTeX\ is good at typesetting mathematical formulas
like
       \( x-3y + z = 7 \)
or
       \( a_{1} > x^{2n} + y^{2n} > x' \)
or
       \( \ip{A}{B} = \sum_{i} a_{i} b_{i} \).
The spaces you type in a formula are
ignored.  Remember that a letter like
       $x$                   % $ ... $  and  \( ... \)  are equivalent
is a formula when it denotes a mathematical
symbol, and it should be typed as one.

\section{Displayed Text}

Text is displayed by indenting it from the left
margin.  Quotations are commonly displayed.  There
are short quotations
\begin{quote}
   This is a short quotation.  It consists of a
   single paragraph of text.  See how it is formatted.
\end{quote}
and longer ones.
\begin{quotation}
   This is a longer quotation.  It consists of two
   paragraphs of text, neither of which are
   particularly interesting.

   This is the second paragraph of the quotation.  It
   is just as dull as the first paragraph.
\end{quotation}
Another frequently-displayed structure is a list.
The following is an example of an \emph{itemized}
list.
\begin{itemize}
   \item This is the first item of an itemized list.
         Each item in the list is marked with a ``tick''.
         You don't have to worry about what kind of tick
         mark is used.

   \item This is the second item of the list.  It
         contains another list nested inside it.  The inner
         list is an \emph{enumerated} list.
         \begin{enumerate}
            \item This is the first item of an enumerated
                  list that is nested within the itemized list.

            \item This is the second item of the inner list.
                  \LaTeX\ allows you to nest lists deeper than
                  you really should.
         \end{enumerate}
         This is the rest of the second item of the outer
         list.  It is no more interesting than any other
         part of the item.
   \item This is the third item of the list.
\end{itemize}
You can even display poetry.
\begin{verse}
   There is an environment
    for verse \\             % The \\ command separates lines
   Whose features some poets % within a stanza.
   will curse.

                             % One or more blank lines separate stanzas.

   For instead of making\\
   Them do \emph{all} line breaking, \\
   It allows them to put too many words on a line when they'd rather be
   forced to be terse.
\end{verse}

Mathematical formulas may also be displayed.  A
displayed formula
is
one-line long; multiline
formulas require special formatting instructions.
   \[  \ip{\Gamma}{\psi'} = x'' + y^{2} + z_{i}^{n}\]
Don't start a paragraph with a displayed equation,
nor make one a paragraph by itself.

\end{document}               % End of document.

\section{Bench Huge Section 376}
% synthetic bench section 376
% This is a sample LaTeX input file.  (Version of 12 August 2004.)
%
% A '%' character causes TeX to ignore all remaining text on the line,
% and is used for comments like this one.

\documentclass{article}      % Specifies the document class

                             % The preamble begins here.
\title{An Example Document}  % Declares the document's title.
\author{Leslie Lamport}      % Declares the author's name.
\date{January 21, 1994}      % Deleting this command produces today's date.

\newcommand{\ip}[2]{(#1, #2)}
                             % Defines \ip{arg1}{arg2} to mean
                             % (arg1, arg2).

%\newcommand{\ip}[2]{\langle #1 | #2\rangle}
                             % This is an alternative definition of
                             % \ip that is commented out.

\begin{document}             % End of preamble and beginning of text.

\maketitle                   % Produces the title.

This is an example input file.  Comparing it with
the output it generates can show you how to
produce a simple document of your own.

\section{Ordinary Text}      % Produces section heading.  Lower-level
                             % sections are begun with similar
                             % \subsection and \subsubsection commands.

The ends  of words and sentences are marked
  by   spaces. It  doesn't matter how many
spaces    you type; one is as good as 100.  The
end of   a line counts as a space.

One   or more   blank lines denote the  end
of  a paragraph.

Since any number of consecutive spaces are treated
like a single one, the formatting of the input
file makes no difference to
      \LaTeX,                % The \LaTeX command generates the LaTeX logo.
but it makes a difference to you.  When you use
\LaTeX, making your input file as easy to read
as possible will be a great help as you write
your document and when you change it.  This sample
file shows how you can add comments to your own input
file.

Because printing is different from typewriting,
there are a number of things that you have to do
differently when preparing an input file than if
you were just typing the document directly.
Quotation marks like
       ``this''
have to be handled specially, as do quotes within
quotes:
       ``\,`this'            % \, separates the double and single quote.
        is what I just
        wrote, not  `that'\,''.

Dashes come in three sizes: an
       intra-word
dash, a medium dash for number ranges like
       1--2,
and a punctuation
       dash---like
this.

A sentence-ending space should be larger than the
space between words within a sentence.  You
sometimes have to type special commands in
conjunction with punctuation characters to get
this right, as in the following sentence.
       Gnats, gnus, etc.\ all  % `\ ' makes an inter-word space.
       begin with G\@.         % \@ marks end-of-sentence punctuation.
You should check the spaces after periods when
reading your output to make sure you haven't
forgotten any special cases.  Generating an
ellipsis
       \ldots\               % `\ ' is needed after `\ldots' because TeX
                             % ignores spaces after command names like \ldots
                             % made from \ + letters.
                             %
                             % Note how a `%' character causes TeX to ignore
                             % the end of the input line, so these blank lines
                             % do not start a new paragraph.
                             %
with the right spacing around the periods requires
a special command.

\LaTeX\ interprets some common characters as
commands, so you must type special commands to
generate them.  These characters include the
following:
       \$ \& \% \# \{ and \}.

In printing, text is usually emphasized with an
       \emph{italic}
type style.

\begin{em}
   A long segment of text can also be emphasized
   in this way.  Text within such a segment can be
   given \emph{additional} emphasis.
\end{em}

It is sometimes necessary to prevent \LaTeX\ from
breaking a line where it might otherwise do so.
This may be at a space, as between the ``Mr.''\ and
``Jones'' in
       ``Mr.~Jones'',        % ~ produces an unbreakable interword space.
or within a word---especially when the word is a
symbol like
       \mbox{\emph{itemnum}}
that makes little sense when hyphenated across
lines.

Footnotes\footnote{This is an example of a footnote.}
pose no problem.

\LaTeX\ is good at typesetting mathematical formulas
like
       \( x-3y + z = 7 \)
or
       \( a_{1} > x^{2n} + y^{2n} > x' \)
or
       \( \ip{A}{B} = \sum_{i} a_{i} b_{i} \).
The spaces you type in a formula are
ignored.  Remember that a letter like
       $x$                   % $ ... $  and  \( ... \)  are equivalent
is a formula when it denotes a mathematical
symbol, and it should be typed as one.

\section{Displayed Text}

Text is displayed by indenting it from the left
margin.  Quotations are commonly displayed.  There
are short quotations
\begin{quote}
   This is a short quotation.  It consists of a
   single paragraph of text.  See how it is formatted.
\end{quote}
and longer ones.
\begin{quotation}
   This is a longer quotation.  It consists of two
   paragraphs of text, neither of which are
   particularly interesting.

   This is the second paragraph of the quotation.  It
   is just as dull as the first paragraph.
\end{quotation}
Another frequently-displayed structure is a list.
The following is an example of an \emph{itemized}
list.
\begin{itemize}
   \item This is the first item of an itemized list.
         Each item in the list is marked with a ``tick''.
         You don't have to worry about what kind of tick
         mark is used.

   \item This is the second item of the list.  It
         contains another list nested inside it.  The inner
         list is an \emph{enumerated} list.
         \begin{enumerate}
            \item This is the first item of an enumerated
                  list that is nested within the itemized list.

            \item This is the second item of the inner list.
                  \LaTeX\ allows you to nest lists deeper than
                  you really should.
         \end{enumerate}
         This is the rest of the second item of the outer
         list.  It is no more interesting than any other
         part of the item.
   \item This is the third item of the list.
\end{itemize}
You can even display poetry.
\begin{verse}
   There is an environment
    for verse \\             % The \\ command separates lines
   Whose features some poets % within a stanza.
   will curse.

                             % One or more blank lines separate stanzas.

   For instead of making\\
   Them do \emph{all} line breaking, \\
   It allows them to put too many words on a line when they'd rather be
   forced to be terse.
\end{verse}

Mathematical formulas may also be displayed.  A
displayed formula
is
one-line long; multiline
formulas require special formatting instructions.
   \[  \ip{\Gamma}{\psi'} = x'' + y^{2} + z_{i}^{n}\]
Don't start a paragraph with a displayed equation,
nor make one a paragraph by itself.

\end{document}               % End of document.

\section{Bench Huge Section 377}
% synthetic bench section 377
% This is a sample LaTeX input file.  (Version of 12 August 2004.)
%
% A '%' character causes TeX to ignore all remaining text on the line,
% and is used for comments like this one.

\documentclass{article}      % Specifies the document class

                             % The preamble begins here.
\title{An Example Document}  % Declares the document's title.
\author{Leslie Lamport}      % Declares the author's name.
\date{January 21, 1994}      % Deleting this command produces today's date.

\newcommand{\ip}[2]{(#1, #2)}
                             % Defines \ip{arg1}{arg2} to mean
                             % (arg1, arg2).

%\newcommand{\ip}[2]{\langle #1 | #2\rangle}
                             % This is an alternative definition of
                             % \ip that is commented out.

\begin{document}             % End of preamble and beginning of text.

\maketitle                   % Produces the title.

This is an example input file.  Comparing it with
the output it generates can show you how to
produce a simple document of your own.

\section{Ordinary Text}      % Produces section heading.  Lower-level
                             % sections are begun with similar
                             % \subsection and \subsubsection commands.

The ends  of words and sentences are marked
  by   spaces. It  doesn't matter how many
spaces    you type; one is as good as 100.  The
end of   a line counts as a space.

One   or more   blank lines denote the  end
of  a paragraph.

Since any number of consecutive spaces are treated
like a single one, the formatting of the input
file makes no difference to
      \LaTeX,                % The \LaTeX command generates the LaTeX logo.
but it makes a difference to you.  When you use
\LaTeX, making your input file as easy to read
as possible will be a great help as you write
your document and when you change it.  This sample
file shows how you can add comments to your own input
file.

Because printing is different from typewriting,
there are a number of things that you have to do
differently when preparing an input file than if
you were just typing the document directly.
Quotation marks like
       ``this''
have to be handled specially, as do quotes within
quotes:
       ``\,`this'            % \, separates the double and single quote.
        is what I just
        wrote, not  `that'\,''.

Dashes come in three sizes: an
       intra-word
dash, a medium dash for number ranges like
       1--2,
and a punctuation
       dash---like
this.

A sentence-ending space should be larger than the
space between words within a sentence.  You
sometimes have to type special commands in
conjunction with punctuation characters to get
this right, as in the following sentence.
       Gnats, gnus, etc.\ all  % `\ ' makes an inter-word space.
       begin with G\@.         % \@ marks end-of-sentence punctuation.
You should check the spaces after periods when
reading your output to make sure you haven't
forgotten any special cases.  Generating an
ellipsis
       \ldots\               % `\ ' is needed after `\ldots' because TeX
                             % ignores spaces after command names like \ldots
                             % made from \ + letters.
                             %
                             % Note how a `%' character causes TeX to ignore
                             % the end of the input line, so these blank lines
                             % do not start a new paragraph.
                             %
with the right spacing around the periods requires
a special command.

\LaTeX\ interprets some common characters as
commands, so you must type special commands to
generate them.  These characters include the
following:
       \$ \& \% \# \{ and \}.

In printing, text is usually emphasized with an
       \emph{italic}
type style.

\begin{em}
   A long segment of text can also be emphasized
   in this way.  Text within such a segment can be
   given \emph{additional} emphasis.
\end{em}

It is sometimes necessary to prevent \LaTeX\ from
breaking a line where it might otherwise do so.
This may be at a space, as between the ``Mr.''\ and
``Jones'' in
       ``Mr.~Jones'',        % ~ produces an unbreakable interword space.
or within a word---especially when the word is a
symbol like
       \mbox{\emph{itemnum}}
that makes little sense when hyphenated across
lines.

Footnotes\footnote{This is an example of a footnote.}
pose no problem.

\LaTeX\ is good at typesetting mathematical formulas
like
       \( x-3y + z = 7 \)
or
       \( a_{1} > x^{2n} + y^{2n} > x' \)
or
       \( \ip{A}{B} = \sum_{i} a_{i} b_{i} \).
The spaces you type in a formula are
ignored.  Remember that a letter like
       $x$                   % $ ... $  and  \( ... \)  are equivalent
is a formula when it denotes a mathematical
symbol, and it should be typed as one.

\section{Displayed Text}

Text is displayed by indenting it from the left
margin.  Quotations are commonly displayed.  There
are short quotations
\begin{quote}
   This is a short quotation.  It consists of a
   single paragraph of text.  See how it is formatted.
\end{quote}
and longer ones.
\begin{quotation}
   This is a longer quotation.  It consists of two
   paragraphs of text, neither of which are
   particularly interesting.

   This is the second paragraph of the quotation.  It
   is just as dull as the first paragraph.
\end{quotation}
Another frequently-displayed structure is a list.
The following is an example of an \emph{itemized}
list.
\begin{itemize}
   \item This is the first item of an itemized list.
         Each item in the list is marked with a ``tick''.
         You don't have to worry about what kind of tick
         mark is used.

   \item This is the second item of the list.  It
         contains another list nested inside it.  The inner
         list is an \emph{enumerated} list.
         \begin{enumerate}
            \item This is the first item of an enumerated
                  list that is nested within the itemized list.

            \item This is the second item of the inner list.
                  \LaTeX\ allows you to nest lists deeper than
                  you really should.
         \end{enumerate}
         This is the rest of the second item of the outer
         list.  It is no more interesting than any other
         part of the item.
   \item This is the third item of the list.
\end{itemize}
You can even display poetry.
\begin{verse}
   There is an environment
    for verse \\             % The \\ command separates lines
   Whose features some poets % within a stanza.
   will curse.

                             % One or more blank lines separate stanzas.

   For instead of making\\
   Them do \emph{all} line breaking, \\
   It allows them to put too many words on a line when they'd rather be
   forced to be terse.
\end{verse}

Mathematical formulas may also be displayed.  A
displayed formula
is
one-line long; multiline
formulas require special formatting instructions.
   \[  \ip{\Gamma}{\psi'} = x'' + y^{2} + z_{i}^{n}\]
Don't start a paragraph with a displayed equation,
nor make one a paragraph by itself.

\end{document}               % End of document.

\section{Bench Huge Section 378}
% synthetic bench section 378
% This is a sample LaTeX input file.  (Version of 12 August 2004.)
%
% A '%' character causes TeX to ignore all remaining text on the line,
% and is used for comments like this one.

\documentclass{article}      % Specifies the document class

                             % The preamble begins here.
\title{An Example Document}  % Declares the document's title.
\author{Leslie Lamport}      % Declares the author's name.
\date{January 21, 1994}      % Deleting this command produces today's date.

\newcommand{\ip}[2]{(#1, #2)}
                             % Defines \ip{arg1}{arg2} to mean
                             % (arg1, arg2).

%\newcommand{\ip}[2]{\langle #1 | #2\rangle}
                             % This is an alternative definition of
                             % \ip that is commented out.

\begin{document}             % End of preamble and beginning of text.

\maketitle                   % Produces the title.

This is an example input file.  Comparing it with
the output it generates can show you how to
produce a simple document of your own.

\section{Ordinary Text}      % Produces section heading.  Lower-level
                             % sections are begun with similar
                             % \subsection and \subsubsection commands.

The ends  of words and sentences are marked
  by   spaces. It  doesn't matter how many
spaces    you type; one is as good as 100.  The
end of   a line counts as a space.

One   or more   blank lines denote the  end
of  a paragraph.

Since any number of consecutive spaces are treated
like a single one, the formatting of the input
file makes no difference to
      \LaTeX,                % The \LaTeX command generates the LaTeX logo.
but it makes a difference to you.  When you use
\LaTeX, making your input file as easy to read
as possible will be a great help as you write
your document and when you change it.  This sample
file shows how you can add comments to your own input
file.

Because printing is different from typewriting,
there are a number of things that you have to do
differently when preparing an input file than if
you were just typing the document directly.
Quotation marks like
       ``this''
have to be handled specially, as do quotes within
quotes:
       ``\,`this'            % \, separates the double and single quote.
        is what I just
        wrote, not  `that'\,''.

Dashes come in three sizes: an
       intra-word
dash, a medium dash for number ranges like
       1--2,
and a punctuation
       dash---like
this.

A sentence-ending space should be larger than the
space between words within a sentence.  You
sometimes have to type special commands in
conjunction with punctuation characters to get
this right, as in the following sentence.
       Gnats, gnus, etc.\ all  % `\ ' makes an inter-word space.
       begin with G\@.         % \@ marks end-of-sentence punctuation.
You should check the spaces after periods when
reading your output to make sure you haven't
forgotten any special cases.  Generating an
ellipsis
       \ldots\               % `\ ' is needed after `\ldots' because TeX
                             % ignores spaces after command names like \ldots
                             % made from \ + letters.
                             %
                             % Note how a `%' character causes TeX to ignore
                             % the end of the input line, so these blank lines
                             % do not start a new paragraph.
                             %
with the right spacing around the periods requires
a special command.

\LaTeX\ interprets some common characters as
commands, so you must type special commands to
generate them.  These characters include the
following:
       \$ \& \% \# \{ and \}.

In printing, text is usually emphasized with an
       \emph{italic}
type style.

\begin{em}
   A long segment of text can also be emphasized
   in this way.  Text within such a segment can be
   given \emph{additional} emphasis.
\end{em}

It is sometimes necessary to prevent \LaTeX\ from
breaking a line where it might otherwise do so.
This may be at a space, as between the ``Mr.''\ and
``Jones'' in
       ``Mr.~Jones'',        % ~ produces an unbreakable interword space.
or within a word---especially when the word is a
symbol like
       \mbox{\emph{itemnum}}
that makes little sense when hyphenated across
lines.

Footnotes\footnote{This is an example of a footnote.}
pose no problem.

\LaTeX\ is good at typesetting mathematical formulas
like
       \( x-3y + z = 7 \)
or
       \( a_{1} > x^{2n} + y^{2n} > x' \)
or
       \( \ip{A}{B} = \sum_{i} a_{i} b_{i} \).
The spaces you type in a formula are
ignored.  Remember that a letter like
       $x$                   % $ ... $  and  \( ... \)  are equivalent
is a formula when it denotes a mathematical
symbol, and it should be typed as one.

\section{Displayed Text}

Text is displayed by indenting it from the left
margin.  Quotations are commonly displayed.  There
are short quotations
\begin{quote}
   This is a short quotation.  It consists of a
   single paragraph of text.  See how it is formatted.
\end{quote}
and longer ones.
\begin{quotation}
   This is a longer quotation.  It consists of two
   paragraphs of text, neither of which are
   particularly interesting.

   This is the second paragraph of the quotation.  It
   is just as dull as the first paragraph.
\end{quotation}
Another frequently-displayed structure is a list.
The following is an example of an \emph{itemized}
list.
\begin{itemize}
   \item This is the first item of an itemized list.
         Each item in the list is marked with a ``tick''.
         You don't have to worry about what kind of tick
         mark is used.

   \item This is the second item of the list.  It
         contains another list nested inside it.  The inner
         list is an \emph{enumerated} list.
         \begin{enumerate}
            \item This is the first item of an enumerated
                  list that is nested within the itemized list.

            \item This is the second item of the inner list.
                  \LaTeX\ allows you to nest lists deeper than
                  you really should.
         \end{enumerate}
         This is the rest of the second item of the outer
         list.  It is no more interesting than any other
         part of the item.
   \item This is the third item of the list.
\end{itemize}
You can even display poetry.
\begin{verse}
   There is an environment
    for verse \\             % The \\ command separates lines
   Whose features some poets % within a stanza.
   will curse.

                             % One or more blank lines separate stanzas.

   For instead of making\\
   Them do \emph{all} line breaking, \\
   It allows them to put too many words on a line when they'd rather be
   forced to be terse.
\end{verse}

Mathematical formulas may also be displayed.  A
displayed formula
is
one-line long; multiline
formulas require special formatting instructions.
   \[  \ip{\Gamma}{\psi'} = x'' + y^{2} + z_{i}^{n}\]
Don't start a paragraph with a displayed equation,
nor make one a paragraph by itself.

\end{document}               % End of document.

\section{Bench Huge Section 379}
% synthetic bench section 379
% This is a sample LaTeX input file.  (Version of 12 August 2004.)
%
% A '%' character causes TeX to ignore all remaining text on the line,
% and is used for comments like this one.

\documentclass{article}      % Specifies the document class

                             % The preamble begins here.
\title{An Example Document}  % Declares the document's title.
\author{Leslie Lamport}      % Declares the author's name.
\date{January 21, 1994}      % Deleting this command produces today's date.

\newcommand{\ip}[2]{(#1, #2)}
                             % Defines \ip{arg1}{arg2} to mean
                             % (arg1, arg2).

%\newcommand{\ip}[2]{\langle #1 | #2\rangle}
                             % This is an alternative definition of
                             % \ip that is commented out.

\begin{document}             % End of preamble and beginning of text.

\maketitle                   % Produces the title.

This is an example input file.  Comparing it with
the output it generates can show you how to
produce a simple document of your own.

\section{Ordinary Text}      % Produces section heading.  Lower-level
                             % sections are begun with similar
                             % \subsection and \subsubsection commands.

The ends  of words and sentences are marked
  by   spaces. It  doesn't matter how many
spaces    you type; one is as good as 100.  The
end of   a line counts as a space.

One   or more   blank lines denote the  end
of  a paragraph.

Since any number of consecutive spaces are treated
like a single one, the formatting of the input
file makes no difference to
      \LaTeX,                % The \LaTeX command generates the LaTeX logo.
but it makes a difference to you.  When you use
\LaTeX, making your input file as easy to read
as possible will be a great help as you write
your document and when you change it.  This sample
file shows how you can add comments to your own input
file.

Because printing is different from typewriting,
there are a number of things that you have to do
differently when preparing an input file than if
you were just typing the document directly.
Quotation marks like
       ``this''
have to be handled specially, as do quotes within
quotes:
       ``\,`this'            % \, separates the double and single quote.
        is what I just
        wrote, not  `that'\,''.

Dashes come in three sizes: an
       intra-word
dash, a medium dash for number ranges like
       1--2,
and a punctuation
       dash---like
this.

A sentence-ending space should be larger than the
space between words within a sentence.  You
sometimes have to type special commands in
conjunction with punctuation characters to get
this right, as in the following sentence.
       Gnats, gnus, etc.\ all  % `\ ' makes an inter-word space.
       begin with G\@.         % \@ marks end-of-sentence punctuation.
You should check the spaces after periods when
reading your output to make sure you haven't
forgotten any special cases.  Generating an
ellipsis
       \ldots\               % `\ ' is needed after `\ldots' because TeX
                             % ignores spaces after command names like \ldots
                             % made from \ + letters.
                             %
                             % Note how a `%' character causes TeX to ignore
                             % the end of the input line, so these blank lines
                             % do not start a new paragraph.
                             %
with the right spacing around the periods requires
a special command.

\LaTeX\ interprets some common characters as
commands, so you must type special commands to
generate them.  These characters include the
following:
       \$ \& \% \# \{ and \}.

In printing, text is usually emphasized with an
       \emph{italic}
type style.

\begin{em}
   A long segment of text can also be emphasized
   in this way.  Text within such a segment can be
   given \emph{additional} emphasis.
\end{em}

It is sometimes necessary to prevent \LaTeX\ from
breaking a line where it might otherwise do so.
This may be at a space, as between the ``Mr.''\ and
``Jones'' in
       ``Mr.~Jones'',        % ~ produces an unbreakable interword space.
or within a word---especially when the word is a
symbol like
       \mbox{\emph{itemnum}}
that makes little sense when hyphenated across
lines.

Footnotes\footnote{This is an example of a footnote.}
pose no problem.

\LaTeX\ is good at typesetting mathematical formulas
like
       \( x-3y + z = 7 \)
or
       \( a_{1} > x^{2n} + y^{2n} > x' \)
or
       \( \ip{A}{B} = \sum_{i} a_{i} b_{i} \).
The spaces you type in a formula are
ignored.  Remember that a letter like
       $x$                   % $ ... $  and  \( ... \)  are equivalent
is a formula when it denotes a mathematical
symbol, and it should be typed as one.

\section{Displayed Text}

Text is displayed by indenting it from the left
margin.  Quotations are commonly displayed.  There
are short quotations
\begin{quote}
   This is a short quotation.  It consists of a
   single paragraph of text.  See how it is formatted.
\end{quote}
and longer ones.
\begin{quotation}
   This is a longer quotation.  It consists of two
   paragraphs of text, neither of which are
   particularly interesting.

   This is the second paragraph of the quotation.  It
   is just as dull as the first paragraph.
\end{quotation}
Another frequently-displayed structure is a list.
The following is an example of an \emph{itemized}
list.
\begin{itemize}
   \item This is the first item of an itemized list.
         Each item in the list is marked with a ``tick''.
         You don't have to worry about what kind of tick
         mark is used.

   \item This is the second item of the list.  It
         contains another list nested inside it.  The inner
         list is an \emph{enumerated} list.
         \begin{enumerate}
            \item This is the first item of an enumerated
                  list that is nested within the itemized list.

            \item This is the second item of the inner list.
                  \LaTeX\ allows you to nest lists deeper than
                  you really should.
         \end{enumerate}
         This is the rest of the second item of the outer
         list.  It is no more interesting than any other
         part of the item.
   \item This is the third item of the list.
\end{itemize}
You can even display poetry.
\begin{verse}
   There is an environment
    for verse \\             % The \\ command separates lines
   Whose features some poets % within a stanza.
   will curse.

                             % One or more blank lines separate stanzas.

   For instead of making\\
   Them do \emph{all} line breaking, \\
   It allows them to put too many words on a line when they'd rather be
   forced to be terse.
\end{verse}

Mathematical formulas may also be displayed.  A
displayed formula
is
one-line long; multiline
formulas require special formatting instructions.
   \[  \ip{\Gamma}{\psi'} = x'' + y^{2} + z_{i}^{n}\]
Don't start a paragraph with a displayed equation,
nor make one a paragraph by itself.

\end{document}               % End of document.

\section{Bench Huge Section 380}
% synthetic bench section 380
% This is a sample LaTeX input file.  (Version of 12 August 2004.)
%
% A '%' character causes TeX to ignore all remaining text on the line,
% and is used for comments like this one.

\documentclass{article}      % Specifies the document class

                             % The preamble begins here.
\title{An Example Document}  % Declares the document's title.
\author{Leslie Lamport}      % Declares the author's name.
\date{January 21, 1994}      % Deleting this command produces today's date.

\newcommand{\ip}[2]{(#1, #2)}
                             % Defines \ip{arg1}{arg2} to mean
                             % (arg1, arg2).

%\newcommand{\ip}[2]{\langle #1 | #2\rangle}
                             % This is an alternative definition of
                             % \ip that is commented out.

\begin{document}             % End of preamble and beginning of text.

\maketitle                   % Produces the title.

This is an example input file.  Comparing it with
the output it generates can show you how to
produce a simple document of your own.

\section{Ordinary Text}      % Produces section heading.  Lower-level
                             % sections are begun with similar
                             % \subsection and \subsubsection commands.

The ends  of words and sentences are marked
  by   spaces. It  doesn't matter how many
spaces    you type; one is as good as 100.  The
end of   a line counts as a space.

One   or more   blank lines denote the  end
of  a paragraph.

Since any number of consecutive spaces are treated
like a single one, the formatting of the input
file makes no difference to
      \LaTeX,                % The \LaTeX command generates the LaTeX logo.
but it makes a difference to you.  When you use
\LaTeX, making your input file as easy to read
as possible will be a great help as you write
your document and when you change it.  This sample
file shows how you can add comments to your own input
file.

Because printing is different from typewriting,
there are a number of things that you have to do
differently when preparing an input file than if
you were just typing the document directly.
Quotation marks like
       ``this''
have to be handled specially, as do quotes within
quotes:
       ``\,`this'            % \, separates the double and single quote.
        is what I just
        wrote, not  `that'\,''.

Dashes come in three sizes: an
       intra-word
dash, a medium dash for number ranges like
       1--2,
and a punctuation
       dash---like
this.

A sentence-ending space should be larger than the
space between words within a sentence.  You
sometimes have to type special commands in
conjunction with punctuation characters to get
this right, as in the following sentence.
       Gnats, gnus, etc.\ all  % `\ ' makes an inter-word space.
       begin with G\@.         % \@ marks end-of-sentence punctuation.
You should check the spaces after periods when
reading your output to make sure you haven't
forgotten any special cases.  Generating an
ellipsis
       \ldots\               % `\ ' is needed after `\ldots' because TeX
                             % ignores spaces after command names like \ldots
                             % made from \ + letters.
                             %
                             % Note how a `%' character causes TeX to ignore
                             % the end of the input line, so these blank lines
                             % do not start a new paragraph.
                             %
with the right spacing around the periods requires
a special command.

\LaTeX\ interprets some common characters as
commands, so you must type special commands to
generate them.  These characters include the
following:
       \$ \& \% \# \{ and \}.

In printing, text is usually emphasized with an
       \emph{italic}
type style.

\begin{em}
   A long segment of text can also be emphasized
   in this way.  Text within such a segment can be
   given \emph{additional} emphasis.
\end{em}

It is sometimes necessary to prevent \LaTeX\ from
breaking a line where it might otherwise do so.
This may be at a space, as between the ``Mr.''\ and
``Jones'' in
       ``Mr.~Jones'',        % ~ produces an unbreakable interword space.
or within a word---especially when the word is a
symbol like
       \mbox{\emph{itemnum}}
that makes little sense when hyphenated across
lines.

Footnotes\footnote{This is an example of a footnote.}
pose no problem.

\LaTeX\ is good at typesetting mathematical formulas
like
       \( x-3y + z = 7 \)
or
       \( a_{1} > x^{2n} + y^{2n} > x' \)
or
       \( \ip{A}{B} = \sum_{i} a_{i} b_{i} \).
The spaces you type in a formula are
ignored.  Remember that a letter like
       $x$                   % $ ... $  and  \( ... \)  are equivalent
is a formula when it denotes a mathematical
symbol, and it should be typed as one.

\section{Displayed Text}

Text is displayed by indenting it from the left
margin.  Quotations are commonly displayed.  There
are short quotations
\begin{quote}
   This is a short quotation.  It consists of a
   single paragraph of text.  See how it is formatted.
\end{quote}
and longer ones.
\begin{quotation}
   This is a longer quotation.  It consists of two
   paragraphs of text, neither of which are
   particularly interesting.

   This is the second paragraph of the quotation.  It
   is just as dull as the first paragraph.
\end{quotation}
Another frequently-displayed structure is a list.
The following is an example of an \emph{itemized}
list.
\begin{itemize}
   \item This is the first item of an itemized list.
         Each item in the list is marked with a ``tick''.
         You don't have to worry about what kind of tick
         mark is used.

   \item This is the second item of the list.  It
         contains another list nested inside it.  The inner
         list is an \emph{enumerated} list.
         \begin{enumerate}
            \item This is the first item of an enumerated
                  list that is nested within the itemized list.

            \item This is the second item of the inner list.
                  \LaTeX\ allows you to nest lists deeper than
                  you really should.
         \end{enumerate}
         This is the rest of the second item of the outer
         list.  It is no more interesting than any other
         part of the item.
   \item This is the third item of the list.
\end{itemize}
You can even display poetry.
\begin{verse}
   There is an environment
    for verse \\             % The \\ command separates lines
   Whose features some poets % within a stanza.
   will curse.

                             % One or more blank lines separate stanzas.

   For instead of making\\
   Them do \emph{all} line breaking, \\
   It allows them to put too many words on a line when they'd rather be
   forced to be terse.
\end{verse}

Mathematical formulas may also be displayed.  A
displayed formula
is
one-line long; multiline
formulas require special formatting instructions.
   \[  \ip{\Gamma}{\psi'} = x'' + y^{2} + z_{i}^{n}\]
Don't start a paragraph with a displayed equation,
nor make one a paragraph by itself.

\end{document}               % End of document.

\section{Bench Huge Section 381}
% synthetic bench section 381
% This is a sample LaTeX input file.  (Version of 12 August 2004.)
%
% A '%' character causes TeX to ignore all remaining text on the line,
% and is used for comments like this one.

\documentclass{article}      % Specifies the document class

                             % The preamble begins here.
\title{An Example Document}  % Declares the document's title.
\author{Leslie Lamport}      % Declares the author's name.
\date{January 21, 1994}      % Deleting this command produces today's date.

\newcommand{\ip}[2]{(#1, #2)}
                             % Defines \ip{arg1}{arg2} to mean
                             % (arg1, arg2).

%\newcommand{\ip}[2]{\langle #1 | #2\rangle}
                             % This is an alternative definition of
                             % \ip that is commented out.

\begin{document}             % End of preamble and beginning of text.

\maketitle                   % Produces the title.

This is an example input file.  Comparing it with
the output it generates can show you how to
produce a simple document of your own.

\section{Ordinary Text}      % Produces section heading.  Lower-level
                             % sections are begun with similar
                             % \subsection and \subsubsection commands.

The ends  of words and sentences are marked
  by   spaces. It  doesn't matter how many
spaces    you type; one is as good as 100.  The
end of   a line counts as a space.

One   or more   blank lines denote the  end
of  a paragraph.

Since any number of consecutive spaces are treated
like a single one, the formatting of the input
file makes no difference to
      \LaTeX,                % The \LaTeX command generates the LaTeX logo.
but it makes a difference to you.  When you use
\LaTeX, making your input file as easy to read
as possible will be a great help as you write
your document and when you change it.  This sample
file shows how you can add comments to your own input
file.

Because printing is different from typewriting,
there are a number of things that you have to do
differently when preparing an input file than if
you were just typing the document directly.
Quotation marks like
       ``this''
have to be handled specially, as do quotes within
quotes:
       ``\,`this'            % \, separates the double and single quote.
        is what I just
        wrote, not  `that'\,''.

Dashes come in three sizes: an
       intra-word
dash, a medium dash for number ranges like
       1--2,
and a punctuation
       dash---like
this.

A sentence-ending space should be larger than the
space between words within a sentence.  You
sometimes have to type special commands in
conjunction with punctuation characters to get
this right, as in the following sentence.
       Gnats, gnus, etc.\ all  % `\ ' makes an inter-word space.
       begin with G\@.         % \@ marks end-of-sentence punctuation.
You should check the spaces after periods when
reading your output to make sure you haven't
forgotten any special cases.  Generating an
ellipsis
       \ldots\               % `\ ' is needed after `\ldots' because TeX
                             % ignores spaces after command names like \ldots
                             % made from \ + letters.
                             %
                             % Note how a `%' character causes TeX to ignore
                             % the end of the input line, so these blank lines
                             % do not start a new paragraph.
                             %
with the right spacing around the periods requires
a special command.

\LaTeX\ interprets some common characters as
commands, so you must type special commands to
generate them.  These characters include the
following:
       \$ \& \% \# \{ and \}.

In printing, text is usually emphasized with an
       \emph{italic}
type style.

\begin{em}
   A long segment of text can also be emphasized
   in this way.  Text within such a segment can be
   given \emph{additional} emphasis.
\end{em}

It is sometimes necessary to prevent \LaTeX\ from
breaking a line where it might otherwise do so.
This may be at a space, as between the ``Mr.''\ and
``Jones'' in
       ``Mr.~Jones'',        % ~ produces an unbreakable interword space.
or within a word---especially when the word is a
symbol like
       \mbox{\emph{itemnum}}
that makes little sense when hyphenated across
lines.

Footnotes\footnote{This is an example of a footnote.}
pose no problem.

\LaTeX\ is good at typesetting mathematical formulas
like
       \( x-3y + z = 7 \)
or
       \( a_{1} > x^{2n} + y^{2n} > x' \)
or
       \( \ip{A}{B} = \sum_{i} a_{i} b_{i} \).
The spaces you type in a formula are
ignored.  Remember that a letter like
       $x$                   % $ ... $  and  \( ... \)  are equivalent
is a formula when it denotes a mathematical
symbol, and it should be typed as one.

\section{Displayed Text}

Text is displayed by indenting it from the left
margin.  Quotations are commonly displayed.  There
are short quotations
\begin{quote}
   This is a short quotation.  It consists of a
   single paragraph of text.  See how it is formatted.
\end{quote}
and longer ones.
\begin{quotation}
   This is a longer quotation.  It consists of two
   paragraphs of text, neither of which are
   particularly interesting.

   This is the second paragraph of the quotation.  It
   is just as dull as the first paragraph.
\end{quotation}
Another frequently-displayed structure is a list.
The following is an example of an \emph{itemized}
list.
\begin{itemize}
   \item This is the first item of an itemized list.
         Each item in the list is marked with a ``tick''.
         You don't have to worry about what kind of tick
         mark is used.

   \item This is the second item of the list.  It
         contains another list nested inside it.  The inner
         list is an \emph{enumerated} list.
         \begin{enumerate}
            \item This is the first item of an enumerated
                  list that is nested within the itemized list.

            \item This is the second item of the inner list.
                  \LaTeX\ allows you to nest lists deeper than
                  you really should.
         \end{enumerate}
         This is the rest of the second item of the outer
         list.  It is no more interesting than any other
         part of the item.
   \item This is the third item of the list.
\end{itemize}
You can even display poetry.
\begin{verse}
   There is an environment
    for verse \\             % The \\ command separates lines
   Whose features some poets % within a stanza.
   will curse.

                             % One or more blank lines separate stanzas.

   For instead of making\\
   Them do \emph{all} line breaking, \\
   It allows them to put too many words on a line when they'd rather be
   forced to be terse.
\end{verse}

Mathematical formulas may also be displayed.  A
displayed formula
is
one-line long; multiline
formulas require special formatting instructions.
   \[  \ip{\Gamma}{\psi'} = x'' + y^{2} + z_{i}^{n}\]
Don't start a paragraph with a displayed equation,
nor make one a paragraph by itself.

\end{document}               % End of document.

\section{Bench Huge Section 382}
% synthetic bench section 382
% This is a sample LaTeX input file.  (Version of 12 August 2004.)
%
% A '%' character causes TeX to ignore all remaining text on the line,
% and is used for comments like this one.

\documentclass{article}      % Specifies the document class

                             % The preamble begins here.
\title{An Example Document}  % Declares the document's title.
\author{Leslie Lamport}      % Declares the author's name.
\date{January 21, 1994}      % Deleting this command produces today's date.

\newcommand{\ip}[2]{(#1, #2)}
                             % Defines \ip{arg1}{arg2} to mean
                             % (arg1, arg2).

%\newcommand{\ip}[2]{\langle #1 | #2\rangle}
                             % This is an alternative definition of
                             % \ip that is commented out.

\begin{document}             % End of preamble and beginning of text.

\maketitle                   % Produces the title.

This is an example input file.  Comparing it with
the output it generates can show you how to
produce a simple document of your own.

\section{Ordinary Text}      % Produces section heading.  Lower-level
                             % sections are begun with similar
                             % \subsection and \subsubsection commands.

The ends  of words and sentences are marked
  by   spaces. It  doesn't matter how many
spaces    you type; one is as good as 100.  The
end of   a line counts as a space.

One   or more   blank lines denote the  end
of  a paragraph.

Since any number of consecutive spaces are treated
like a single one, the formatting of the input
file makes no difference to
      \LaTeX,                % The \LaTeX command generates the LaTeX logo.
but it makes a difference to you.  When you use
\LaTeX, making your input file as easy to read
as possible will be a great help as you write
your document and when you change it.  This sample
file shows how you can add comments to your own input
file.

Because printing is different from typewriting,
there are a number of things that you have to do
differently when preparing an input file than if
you were just typing the document directly.
Quotation marks like
       ``this''
have to be handled specially, as do quotes within
quotes:
       ``\,`this'            % \, separates the double and single quote.
        is what I just
        wrote, not  `that'\,''.

Dashes come in three sizes: an
       intra-word
dash, a medium dash for number ranges like
       1--2,
and a punctuation
       dash---like
this.

A sentence-ending space should be larger than the
space between words within a sentence.  You
sometimes have to type special commands in
conjunction with punctuation characters to get
this right, as in the following sentence.
       Gnats, gnus, etc.\ all  % `\ ' makes an inter-word space.
       begin with G\@.         % \@ marks end-of-sentence punctuation.
You should check the spaces after periods when
reading your output to make sure you haven't
forgotten any special cases.  Generating an
ellipsis
       \ldots\               % `\ ' is needed after `\ldots' because TeX
                             % ignores spaces after command names like \ldots
                             % made from \ + letters.
                             %
                             % Note how a `%' character causes TeX to ignore
                             % the end of the input line, so these blank lines
                             % do not start a new paragraph.
                             %
with the right spacing around the periods requires
a special command.

\LaTeX\ interprets some common characters as
commands, so you must type special commands to
generate them.  These characters include the
following:
       \$ \& \% \# \{ and \}.

In printing, text is usually emphasized with an
       \emph{italic}
type style.

\begin{em}
   A long segment of text can also be emphasized
   in this way.  Text within such a segment can be
   given \emph{additional} emphasis.
\end{em}

It is sometimes necessary to prevent \LaTeX\ from
breaking a line where it might otherwise do so.
This may be at a space, as between the ``Mr.''\ and
``Jones'' in
       ``Mr.~Jones'',        % ~ produces an unbreakable interword space.
or within a word---especially when the word is a
symbol like
       \mbox{\emph{itemnum}}
that makes little sense when hyphenated across
lines.

Footnotes\footnote{This is an example of a footnote.}
pose no problem.

\LaTeX\ is good at typesetting mathematical formulas
like
       \( x-3y + z = 7 \)
or
       \( a_{1} > x^{2n} + y^{2n} > x' \)
or
       \( \ip{A}{B} = \sum_{i} a_{i} b_{i} \).
The spaces you type in a formula are
ignored.  Remember that a letter like
       $x$                   % $ ... $  and  \( ... \)  are equivalent
is a formula when it denotes a mathematical
symbol, and it should be typed as one.

\section{Displayed Text}

Text is displayed by indenting it from the left
margin.  Quotations are commonly displayed.  There
are short quotations
\begin{quote}
   This is a short quotation.  It consists of a
   single paragraph of text.  See how it is formatted.
\end{quote}
and longer ones.
\begin{quotation}
   This is a longer quotation.  It consists of two
   paragraphs of text, neither of which are
   particularly interesting.

   This is the second paragraph of the quotation.  It
   is just as dull as the first paragraph.
\end{quotation}
Another frequently-displayed structure is a list.
The following is an example of an \emph{itemized}
list.
\begin{itemize}
   \item This is the first item of an itemized list.
         Each item in the list is marked with a ``tick''.
         You don't have to worry about what kind of tick
         mark is used.

   \item This is the second item of the list.  It
         contains another list nested inside it.  The inner
         list is an \emph{enumerated} list.
         \begin{enumerate}
            \item This is the first item of an enumerated
                  list that is nested within the itemized list.

            \item This is the second item of the inner list.
                  \LaTeX\ allows you to nest lists deeper than
                  you really should.
         \end{enumerate}
         This is the rest of the second item of the outer
         list.  It is no more interesting than any other
         part of the item.
   \item This is the third item of the list.
\end{itemize}
You can even display poetry.
\begin{verse}
   There is an environment
    for verse \\             % The \\ command separates lines
   Whose features some poets % within a stanza.
   will curse.

                             % One or more blank lines separate stanzas.

   For instead of making\\
   Them do \emph{all} line breaking, \\
   It allows them to put too many words on a line when they'd rather be
   forced to be terse.
\end{verse}

Mathematical formulas may also be displayed.  A
displayed formula
is
one-line long; multiline
formulas require special formatting instructions.
   \[  \ip{\Gamma}{\psi'} = x'' + y^{2} + z_{i}^{n}\]
Don't start a paragraph with a displayed equation,
nor make one a paragraph by itself.

\end{document}               % End of document.

\section{Bench Huge Section 383}
% synthetic bench section 383
% This is a sample LaTeX input file.  (Version of 12 August 2004.)
%
% A '%' character causes TeX to ignore all remaining text on the line,
% and is used for comments like this one.

\documentclass{article}      % Specifies the document class

                             % The preamble begins here.
\title{An Example Document}  % Declares the document's title.
\author{Leslie Lamport}      % Declares the author's name.
\date{January 21, 1994}      % Deleting this command produces today's date.

\newcommand{\ip}[2]{(#1, #2)}
                             % Defines \ip{arg1}{arg2} to mean
                             % (arg1, arg2).

%\newcommand{\ip}[2]{\langle #1 | #2\rangle}
                             % This is an alternative definition of
                             % \ip that is commented out.

\begin{document}             % End of preamble and beginning of text.

\maketitle                   % Produces the title.

This is an example input file.  Comparing it with
the output it generates can show you how to
produce a simple document of your own.

\section{Ordinary Text}      % Produces section heading.  Lower-level
                             % sections are begun with similar
                             % \subsection and \subsubsection commands.

The ends  of words and sentences are marked
  by   spaces. It  doesn't matter how many
spaces    you type; one is as good as 100.  The
end of   a line counts as a space.

One   or more   blank lines denote the  end
of  a paragraph.

Since any number of consecutive spaces are treated
like a single one, the formatting of the input
file makes no difference to
      \LaTeX,                % The \LaTeX command generates the LaTeX logo.
but it makes a difference to you.  When you use
\LaTeX, making your input file as easy to read
as possible will be a great help as you write
your document and when you change it.  This sample
file shows how you can add comments to your own input
file.

Because printing is different from typewriting,
there are a number of things that you have to do
differently when preparing an input file than if
you were just typing the document directly.
Quotation marks like
       ``this''
have to be handled specially, as do quotes within
quotes:
       ``\,`this'            % \, separates the double and single quote.
        is what I just
        wrote, not  `that'\,''.

Dashes come in three sizes: an
       intra-word
dash, a medium dash for number ranges like
       1--2,
and a punctuation
       dash---like
this.

A sentence-ending space should be larger than the
space between words within a sentence.  You
sometimes have to type special commands in
conjunction with punctuation characters to get
this right, as in the following sentence.
       Gnats, gnus, etc.\ all  % `\ ' makes an inter-word space.
       begin with G\@.         % \@ marks end-of-sentence punctuation.
You should check the spaces after periods when
reading your output to make sure you haven't
forgotten any special cases.  Generating an
ellipsis
       \ldots\               % `\ ' is needed after `\ldots' because TeX
                             % ignores spaces after command names like \ldots
                             % made from \ + letters.
                             %
                             % Note how a `%' character causes TeX to ignore
                             % the end of the input line, so these blank lines
                             % do not start a new paragraph.
                             %
with the right spacing around the periods requires
a special command.

\LaTeX\ interprets some common characters as
commands, so you must type special commands to
generate them.  These characters include the
following:
       \$ \& \% \# \{ and \}.

In printing, text is usually emphasized with an
       \emph{italic}
type style.

\begin{em}
   A long segment of text can also be emphasized
   in this way.  Text within such a segment can be
   given \emph{additional} emphasis.
\end{em}

It is sometimes necessary to prevent \LaTeX\ from
breaking a line where it might otherwise do so.
This may be at a space, as between the ``Mr.''\ and
``Jones'' in
       ``Mr.~Jones'',        % ~ produces an unbreakable interword space.
or within a word---especially when the word is a
symbol like
       \mbox{\emph{itemnum}}
that makes little sense when hyphenated across
lines.

Footnotes\footnote{This is an example of a footnote.}
pose no problem.

\LaTeX\ is good at typesetting mathematical formulas
like
       \( x-3y + z = 7 \)
or
       \( a_{1} > x^{2n} + y^{2n} > x' \)
or
       \( \ip{A}{B} = \sum_{i} a_{i} b_{i} \).
The spaces you type in a formula are
ignored.  Remember that a letter like
       $x$                   % $ ... $  and  \( ... \)  are equivalent
is a formula when it denotes a mathematical
symbol, and it should be typed as one.

\section{Displayed Text}

Text is displayed by indenting it from the left
margin.  Quotations are commonly displayed.  There
are short quotations
\begin{quote}
   This is a short quotation.  It consists of a
   single paragraph of text.  See how it is formatted.
\end{quote}
and longer ones.
\begin{quotation}
   This is a longer quotation.  It consists of two
   paragraphs of text, neither of which are
   particularly interesting.

   This is the second paragraph of the quotation.  It
   is just as dull as the first paragraph.
\end{quotation}
Another frequently-displayed structure is a list.
The following is an example of an \emph{itemized}
list.
\begin{itemize}
   \item This is the first item of an itemized list.
         Each item in the list is marked with a ``tick''.
         You don't have to worry about what kind of tick
         mark is used.

   \item This is the second item of the list.  It
         contains another list nested inside it.  The inner
         list is an \emph{enumerated} list.
         \begin{enumerate}
            \item This is the first item of an enumerated
                  list that is nested within the itemized list.

            \item This is the second item of the inner list.
                  \LaTeX\ allows you to nest lists deeper than
                  you really should.
         \end{enumerate}
         This is the rest of the second item of the outer
         list.  It is no more interesting than any other
         part of the item.
   \item This is the third item of the list.
\end{itemize}
You can even display poetry.
\begin{verse}
   There is an environment
    for verse \\             % The \\ command separates lines
   Whose features some poets % within a stanza.
   will curse.

                             % One or more blank lines separate stanzas.

   For instead of making\\
   Them do \emph{all} line breaking, \\
   It allows them to put too many words on a line when they'd rather be
   forced to be terse.
\end{verse}

Mathematical formulas may also be displayed.  A
displayed formula
is
one-line long; multiline
formulas require special formatting instructions.
   \[  \ip{\Gamma}{\psi'} = x'' + y^{2} + z_{i}^{n}\]
Don't start a paragraph with a displayed equation,
nor make one a paragraph by itself.

\end{document}               % End of document.

\section{Bench Huge Section 384}
% synthetic bench section 384
% This is a sample LaTeX input file.  (Version of 12 August 2004.)
%
% A '%' character causes TeX to ignore all remaining text on the line,
% and is used for comments like this one.

\documentclass{article}      % Specifies the document class

                             % The preamble begins here.
\title{An Example Document}  % Declares the document's title.
\author{Leslie Lamport}      % Declares the author's name.
\date{January 21, 1994}      % Deleting this command produces today's date.

\newcommand{\ip}[2]{(#1, #2)}
                             % Defines \ip{arg1}{arg2} to mean
                             % (arg1, arg2).

%\newcommand{\ip}[2]{\langle #1 | #2\rangle}
                             % This is an alternative definition of
                             % \ip that is commented out.

\begin{document}             % End of preamble and beginning of text.

\maketitle                   % Produces the title.

This is an example input file.  Comparing it with
the output it generates can show you how to
produce a simple document of your own.

\section{Ordinary Text}      % Produces section heading.  Lower-level
                             % sections are begun with similar
                             % \subsection and \subsubsection commands.

The ends  of words and sentences are marked
  by   spaces. It  doesn't matter how many
spaces    you type; one is as good as 100.  The
end of   a line counts as a space.

One   or more   blank lines denote the  end
of  a paragraph.

Since any number of consecutive spaces are treated
like a single one, the formatting of the input
file makes no difference to
      \LaTeX,                % The \LaTeX command generates the LaTeX logo.
but it makes a difference to you.  When you use
\LaTeX, making your input file as easy to read
as possible will be a great help as you write
your document and when you change it.  This sample
file shows how you can add comments to your own input
file.

Because printing is different from typewriting,
there are a number of things that you have to do
differently when preparing an input file than if
you were just typing the document directly.
Quotation marks like
       ``this''
have to be handled specially, as do quotes within
quotes:
       ``\,`this'            % \, separates the double and single quote.
        is what I just
        wrote, not  `that'\,''.

Dashes come in three sizes: an
       intra-word
dash, a medium dash for number ranges like
       1--2,
and a punctuation
       dash---like
this.

A sentence-ending space should be larger than the
space between words within a sentence.  You
sometimes have to type special commands in
conjunction with punctuation characters to get
this right, as in the following sentence.
       Gnats, gnus, etc.\ all  % `\ ' makes an inter-word space.
       begin with G\@.         % \@ marks end-of-sentence punctuation.
You should check the spaces after periods when
reading your output to make sure you haven't
forgotten any special cases.  Generating an
ellipsis
       \ldots\               % `\ ' is needed after `\ldots' because TeX
                             % ignores spaces after command names like \ldots
                             % made from \ + letters.
                             %
                             % Note how a `%' character causes TeX to ignore
                             % the end of the input line, so these blank lines
                             % do not start a new paragraph.
                             %
with the right spacing around the periods requires
a special command.

\LaTeX\ interprets some common characters as
commands, so you must type special commands to
generate them.  These characters include the
following:
       \$ \& \% \# \{ and \}.

In printing, text is usually emphasized with an
       \emph{italic}
type style.

\begin{em}
   A long segment of text can also be emphasized
   in this way.  Text within such a segment can be
   given \emph{additional} emphasis.
\end{em}

It is sometimes necessary to prevent \LaTeX\ from
breaking a line where it might otherwise do so.
This may be at a space, as between the ``Mr.''\ and
``Jones'' in
       ``Mr.~Jones'',        % ~ produces an unbreakable interword space.
or within a word---especially when the word is a
symbol like
       \mbox{\emph{itemnum}}
that makes little sense when hyphenated across
lines.

Footnotes\footnote{This is an example of a footnote.}
pose no problem.

\LaTeX\ is good at typesetting mathematical formulas
like
       \( x-3y + z = 7 \)
or
       \( a_{1} > x^{2n} + y^{2n} > x' \)
or
       \( \ip{A}{B} = \sum_{i} a_{i} b_{i} \).
The spaces you type in a formula are
ignored.  Remember that a letter like
       $x$                   % $ ... $  and  \( ... \)  are equivalent
is a formula when it denotes a mathematical
symbol, and it should be typed as one.

\section{Displayed Text}

Text is displayed by indenting it from the left
margin.  Quotations are commonly displayed.  There
are short quotations
\begin{quote}
   This is a short quotation.  It consists of a
   single paragraph of text.  See how it is formatted.
\end{quote}
and longer ones.
\begin{quotation}
   This is a longer quotation.  It consists of two
   paragraphs of text, neither of which are
   particularly interesting.

   This is the second paragraph of the quotation.  It
   is just as dull as the first paragraph.
\end{quotation}
Another frequently-displayed structure is a list.
The following is an example of an \emph{itemized}
list.
\begin{itemize}
   \item This is the first item of an itemized list.
         Each item in the list is marked with a ``tick''.
         You don't have to worry about what kind of tick
         mark is used.

   \item This is the second item of the list.  It
         contains another list nested inside it.  The inner
         list is an \emph{enumerated} list.
         \begin{enumerate}
            \item This is the first item of an enumerated
                  list that is nested within the itemized list.

            \item This is the second item of the inner list.
                  \LaTeX\ allows you to nest lists deeper than
                  you really should.
         \end{enumerate}
         This is the rest of the second item of the outer
         list.  It is no more interesting than any other
         part of the item.
   \item This is the third item of the list.
\end{itemize}
You can even display poetry.
\begin{verse}
   There is an environment
    for verse \\             % The \\ command separates lines
   Whose features some poets % within a stanza.
   will curse.

                             % One or more blank lines separate stanzas.

   For instead of making\\
   Them do \emph{all} line breaking, \\
   It allows them to put too many words on a line when they'd rather be
   forced to be terse.
\end{verse}

Mathematical formulas may also be displayed.  A
displayed formula
is
one-line long; multiline
formulas require special formatting instructions.
   \[  \ip{\Gamma}{\psi'} = x'' + y^{2} + z_{i}^{n}\]
Don't start a paragraph with a displayed equation,
nor make one a paragraph by itself.

\end{document}               % End of document.

\section{Bench Huge Section 385}
% synthetic bench section 385
% This is a sample LaTeX input file.  (Version of 12 August 2004.)
%
% A '%' character causes TeX to ignore all remaining text on the line,
% and is used for comments like this one.

\documentclass{article}      % Specifies the document class

                             % The preamble begins here.
\title{An Example Document}  % Declares the document's title.
\author{Leslie Lamport}      % Declares the author's name.
\date{January 21, 1994}      % Deleting this command produces today's date.

\newcommand{\ip}[2]{(#1, #2)}
                             % Defines \ip{arg1}{arg2} to mean
                             % (arg1, arg2).

%\newcommand{\ip}[2]{\langle #1 | #2\rangle}
                             % This is an alternative definition of
                             % \ip that is commented out.

\begin{document}             % End of preamble and beginning of text.

\maketitle                   % Produces the title.

This is an example input file.  Comparing it with
the output it generates can show you how to
produce a simple document of your own.

\section{Ordinary Text}      % Produces section heading.  Lower-level
                             % sections are begun with similar
                             % \subsection and \subsubsection commands.

The ends  of words and sentences are marked
  by   spaces. It  doesn't matter how many
spaces    you type; one is as good as 100.  The
end of   a line counts as a space.

One   or more   blank lines denote the  end
of  a paragraph.

Since any number of consecutive spaces are treated
like a single one, the formatting of the input
file makes no difference to
      \LaTeX,                % The \LaTeX command generates the LaTeX logo.
but it makes a difference to you.  When you use
\LaTeX, making your input file as easy to read
as possible will be a great help as you write
your document and when you change it.  This sample
file shows how you can add comments to your own input
file.

Because printing is different from typewriting,
there are a number of things that you have to do
differently when preparing an input file than if
you were just typing the document directly.
Quotation marks like
       ``this''
have to be handled specially, as do quotes within
quotes:
       ``\,`this'            % \, separates the double and single quote.
        is what I just
        wrote, not  `that'\,''.

Dashes come in three sizes: an
       intra-word
dash, a medium dash for number ranges like
       1--2,
and a punctuation
       dash---like
this.

A sentence-ending space should be larger than the
space between words within a sentence.  You
sometimes have to type special commands in
conjunction with punctuation characters to get
this right, as in the following sentence.
       Gnats, gnus, etc.\ all  % `\ ' makes an inter-word space.
       begin with G\@.         % \@ marks end-of-sentence punctuation.
You should check the spaces after periods when
reading your output to make sure you haven't
forgotten any special cases.  Generating an
ellipsis
       \ldots\               % `\ ' is needed after `\ldots' because TeX
                             % ignores spaces after command names like \ldots
                             % made from \ + letters.
                             %
                             % Note how a `%' character causes TeX to ignore
                             % the end of the input line, so these blank lines
                             % do not start a new paragraph.
                             %
with the right spacing around the periods requires
a special command.

\LaTeX\ interprets some common characters as
commands, so you must type special commands to
generate them.  These characters include the
following:
       \$ \& \% \# \{ and \}.

In printing, text is usually emphasized with an
       \emph{italic}
type style.

\begin{em}
   A long segment of text can also be emphasized
   in this way.  Text within such a segment can be
   given \emph{additional} emphasis.
\end{em}

It is sometimes necessary to prevent \LaTeX\ from
breaking a line where it might otherwise do so.
This may be at a space, as between the ``Mr.''\ and
``Jones'' in
       ``Mr.~Jones'',        % ~ produces an unbreakable interword space.
or within a word---especially when the word is a
symbol like
       \mbox{\emph{itemnum}}
that makes little sense when hyphenated across
lines.

Footnotes\footnote{This is an example of a footnote.}
pose no problem.

\LaTeX\ is good at typesetting mathematical formulas
like
       \( x-3y + z = 7 \)
or
       \( a_{1} > x^{2n} + y^{2n} > x' \)
or
       \( \ip{A}{B} = \sum_{i} a_{i} b_{i} \).
The spaces you type in a formula are
ignored.  Remember that a letter like
       $x$                   % $ ... $  and  \( ... \)  are equivalent
is a formula when it denotes a mathematical
symbol, and it should be typed as one.

\section{Displayed Text}

Text is displayed by indenting it from the left
margin.  Quotations are commonly displayed.  There
are short quotations
\begin{quote}
   This is a short quotation.  It consists of a
   single paragraph of text.  See how it is formatted.
\end{quote}
and longer ones.
\begin{quotation}
   This is a longer quotation.  It consists of two
   paragraphs of text, neither of which are
   particularly interesting.

   This is the second paragraph of the quotation.  It
   is just as dull as the first paragraph.
\end{quotation}
Another frequently-displayed structure is a list.
The following is an example of an \emph{itemized}
list.
\begin{itemize}
   \item This is the first item of an itemized list.
         Each item in the list is marked with a ``tick''.
         You don't have to worry about what kind of tick
         mark is used.

   \item This is the second item of the list.  It
         contains another list nested inside it.  The inner
         list is an \emph{enumerated} list.
         \begin{enumerate}
            \item This is the first item of an enumerated
                  list that is nested within the itemized list.

            \item This is the second item of the inner list.
                  \LaTeX\ allows you to nest lists deeper than
                  you really should.
         \end{enumerate}
         This is the rest of the second item of the outer
         list.  It is no more interesting than any other
         part of the item.
   \item This is the third item of the list.
\end{itemize}
You can even display poetry.
\begin{verse}
   There is an environment
    for verse \\             % The \\ command separates lines
   Whose features some poets % within a stanza.
   will curse.

                             % One or more blank lines separate stanzas.

   For instead of making\\
   Them do \emph{all} line breaking, \\
   It allows them to put too many words on a line when they'd rather be
   forced to be terse.
\end{verse}

Mathematical formulas may also be displayed.  A
displayed formula
is
one-line long; multiline
formulas require special formatting instructions.
   \[  \ip{\Gamma}{\psi'} = x'' + y^{2} + z_{i}^{n}\]
Don't start a paragraph with a displayed equation,
nor make one a paragraph by itself.

\end{document}               % End of document.

\section{Bench Huge Section 386}
% synthetic bench section 386
% This is a sample LaTeX input file.  (Version of 12 August 2004.)
%
% A '%' character causes TeX to ignore all remaining text on the line,
% and is used for comments like this one.

\documentclass{article}      % Specifies the document class

                             % The preamble begins here.
\title{An Example Document}  % Declares the document's title.
\author{Leslie Lamport}      % Declares the author's name.
\date{January 21, 1994}      % Deleting this command produces today's date.

\newcommand{\ip}[2]{(#1, #2)}
                             % Defines \ip{arg1}{arg2} to mean
                             % (arg1, arg2).

%\newcommand{\ip}[2]{\langle #1 | #2\rangle}
                             % This is an alternative definition of
                             % \ip that is commented out.

\begin{document}             % End of preamble and beginning of text.

\maketitle                   % Produces the title.

This is an example input file.  Comparing it with
the output it generates can show you how to
produce a simple document of your own.

\section{Ordinary Text}      % Produces section heading.  Lower-level
                             % sections are begun with similar
                             % \subsection and \subsubsection commands.

The ends  of words and sentences are marked
  by   spaces. It  doesn't matter how many
spaces    you type; one is as good as 100.  The
end of   a line counts as a space.

One   or more   blank lines denote the  end
of  a paragraph.

Since any number of consecutive spaces are treated
like a single one, the formatting of the input
file makes no difference to
      \LaTeX,                % The \LaTeX command generates the LaTeX logo.
but it makes a difference to you.  When you use
\LaTeX, making your input file as easy to read
as possible will be a great help as you write
your document and when you change it.  This sample
file shows how you can add comments to your own input
file.

Because printing is different from typewriting,
there are a number of things that you have to do
differently when preparing an input file than if
you were just typing the document directly.
Quotation marks like
       ``this''
have to be handled specially, as do quotes within
quotes:
       ``\,`this'            % \, separates the double and single quote.
        is what I just
        wrote, not  `that'\,''.

Dashes come in three sizes: an
       intra-word
dash, a medium dash for number ranges like
       1--2,
and a punctuation
       dash---like
this.

A sentence-ending space should be larger than the
space between words within a sentence.  You
sometimes have to type special commands in
conjunction with punctuation characters to get
this right, as in the following sentence.
       Gnats, gnus, etc.\ all  % `\ ' makes an inter-word space.
       begin with G\@.         % \@ marks end-of-sentence punctuation.
You should check the spaces after periods when
reading your output to make sure you haven't
forgotten any special cases.  Generating an
ellipsis
       \ldots\               % `\ ' is needed after `\ldots' because TeX
                             % ignores spaces after command names like \ldots
                             % made from \ + letters.
                             %
                             % Note how a `%' character causes TeX to ignore
                             % the end of the input line, so these blank lines
                             % do not start a new paragraph.
                             %
with the right spacing around the periods requires
a special command.

\LaTeX\ interprets some common characters as
commands, so you must type special commands to
generate them.  These characters include the
following:
       \$ \& \% \# \{ and \}.

In printing, text is usually emphasized with an
       \emph{italic}
type style.

\begin{em}
   A long segment of text can also be emphasized
   in this way.  Text within such a segment can be
   given \emph{additional} emphasis.
\end{em}

It is sometimes necessary to prevent \LaTeX\ from
breaking a line where it might otherwise do so.
This may be at a space, as between the ``Mr.''\ and
``Jones'' in
       ``Mr.~Jones'',        % ~ produces an unbreakable interword space.
or within a word---especially when the word is a
symbol like
       \mbox{\emph{itemnum}}
that makes little sense when hyphenated across
lines.

Footnotes\footnote{This is an example of a footnote.}
pose no problem.

\LaTeX\ is good at typesetting mathematical formulas
like
       \( x-3y + z = 7 \)
or
       \( a_{1} > x^{2n} + y^{2n} > x' \)
or
       \( \ip{A}{B} = \sum_{i} a_{i} b_{i} \).
The spaces you type in a formula are
ignored.  Remember that a letter like
       $x$                   % $ ... $  and  \( ... \)  are equivalent
is a formula when it denotes a mathematical
symbol, and it should be typed as one.

\section{Displayed Text}

Text is displayed by indenting it from the left
margin.  Quotations are commonly displayed.  There
are short quotations
\begin{quote}
   This is a short quotation.  It consists of a
   single paragraph of text.  See how it is formatted.
\end{quote}
and longer ones.
\begin{quotation}
   This is a longer quotation.  It consists of two
   paragraphs of text, neither of which are
   particularly interesting.

   This is the second paragraph of the quotation.  It
   is just as dull as the first paragraph.
\end{quotation}
Another frequently-displayed structure is a list.
The following is an example of an \emph{itemized}
list.
\begin{itemize}
   \item This is the first item of an itemized list.
         Each item in the list is marked with a ``tick''.
         You don't have to worry about what kind of tick
         mark is used.

   \item This is the second item of the list.  It
         contains another list nested inside it.  The inner
         list is an \emph{enumerated} list.
         \begin{enumerate}
            \item This is the first item of an enumerated
                  list that is nested within the itemized list.

            \item This is the second item of the inner list.
                  \LaTeX\ allows you to nest lists deeper than
                  you really should.
         \end{enumerate}
         This is the rest of the second item of the outer
         list.  It is no more interesting than any other
         part of the item.
   \item This is the third item of the list.
\end{itemize}
You can even display poetry.
\begin{verse}
   There is an environment
    for verse \\             % The \\ command separates lines
   Whose features some poets % within a stanza.
   will curse.

                             % One or more blank lines separate stanzas.

   For instead of making\\
   Them do \emph{all} line breaking, \\
   It allows them to put too many words on a line when they'd rather be
   forced to be terse.
\end{verse}

Mathematical formulas may also be displayed.  A
displayed formula
is
one-line long; multiline
formulas require special formatting instructions.
   \[  \ip{\Gamma}{\psi'} = x'' + y^{2} + z_{i}^{n}\]
Don't start a paragraph with a displayed equation,
nor make one a paragraph by itself.

\end{document}               % End of document.

\section{Bench Huge Section 387}
% synthetic bench section 387
% This is a sample LaTeX input file.  (Version of 12 August 2004.)
%
% A '%' character causes TeX to ignore all remaining text on the line,
% and is used for comments like this one.

\documentclass{article}      % Specifies the document class

                             % The preamble begins here.
\title{An Example Document}  % Declares the document's title.
\author{Leslie Lamport}      % Declares the author's name.
\date{January 21, 1994}      % Deleting this command produces today's date.

\newcommand{\ip}[2]{(#1, #2)}
                             % Defines \ip{arg1}{arg2} to mean
                             % (arg1, arg2).

%\newcommand{\ip}[2]{\langle #1 | #2\rangle}
                             % This is an alternative definition of
                             % \ip that is commented out.

\begin{document}             % End of preamble and beginning of text.

\maketitle                   % Produces the title.

This is an example input file.  Comparing it with
the output it generates can show you how to
produce a simple document of your own.

\section{Ordinary Text}      % Produces section heading.  Lower-level
                             % sections are begun with similar
                             % \subsection and \subsubsection commands.

The ends  of words and sentences are marked
  by   spaces. It  doesn't matter how many
spaces    you type; one is as good as 100.  The
end of   a line counts as a space.

One   or more   blank lines denote the  end
of  a paragraph.

Since any number of consecutive spaces are treated
like a single one, the formatting of the input
file makes no difference to
      \LaTeX,                % The \LaTeX command generates the LaTeX logo.
but it makes a difference to you.  When you use
\LaTeX, making your input file as easy to read
as possible will be a great help as you write
your document and when you change it.  This sample
file shows how you can add comments to your own input
file.

Because printing is different from typewriting,
there are a number of things that you have to do
differently when preparing an input file than if
you were just typing the document directly.
Quotation marks like
       ``this''
have to be handled specially, as do quotes within
quotes:
       ``\,`this'            % \, separates the double and single quote.
        is what I just
        wrote, not  `that'\,''.

Dashes come in three sizes: an
       intra-word
dash, a medium dash for number ranges like
       1--2,
and a punctuation
       dash---like
this.

A sentence-ending space should be larger than the
space between words within a sentence.  You
sometimes have to type special commands in
conjunction with punctuation characters to get
this right, as in the following sentence.
       Gnats, gnus, etc.\ all  % `\ ' makes an inter-word space.
       begin with G\@.         % \@ marks end-of-sentence punctuation.
You should check the spaces after periods when
reading your output to make sure you haven't
forgotten any special cases.  Generating an
ellipsis
       \ldots\               % `\ ' is needed after `\ldots' because TeX
                             % ignores spaces after command names like \ldots
                             % made from \ + letters.
                             %
                             % Note how a `%' character causes TeX to ignore
                             % the end of the input line, so these blank lines
                             % do not start a new paragraph.
                             %
with the right spacing around the periods requires
a special command.

\LaTeX\ interprets some common characters as
commands, so you must type special commands to
generate them.  These characters include the
following:
       \$ \& \% \# \{ and \}.

In printing, text is usually emphasized with an
       \emph{italic}
type style.

\begin{em}
   A long segment of text can also be emphasized
   in this way.  Text within such a segment can be
   given \emph{additional} emphasis.
\end{em}

It is sometimes necessary to prevent \LaTeX\ from
breaking a line where it might otherwise do so.
This may be at a space, as between the ``Mr.''\ and
``Jones'' in
       ``Mr.~Jones'',        % ~ produces an unbreakable interword space.
or within a word---especially when the word is a
symbol like
       \mbox{\emph{itemnum}}
that makes little sense when hyphenated across
lines.

Footnotes\footnote{This is an example of a footnote.}
pose no problem.

\LaTeX\ is good at typesetting mathematical formulas
like
       \( x-3y + z = 7 \)
or
       \( a_{1} > x^{2n} + y^{2n} > x' \)
or
       \( \ip{A}{B} = \sum_{i} a_{i} b_{i} \).
The spaces you type in a formula are
ignored.  Remember that a letter like
       $x$                   % $ ... $  and  \( ... \)  are equivalent
is a formula when it denotes a mathematical
symbol, and it should be typed as one.

\section{Displayed Text}

Text is displayed by indenting it from the left
margin.  Quotations are commonly displayed.  There
are short quotations
\begin{quote}
   This is a short quotation.  It consists of a
   single paragraph of text.  See how it is formatted.
\end{quote}
and longer ones.
\begin{quotation}
   This is a longer quotation.  It consists of two
   paragraphs of text, neither of which are
   particularly interesting.

   This is the second paragraph of the quotation.  It
   is just as dull as the first paragraph.
\end{quotation}
Another frequently-displayed structure is a list.
The following is an example of an \emph{itemized}
list.
\begin{itemize}
   \item This is the first item of an itemized list.
         Each item in the list is marked with a ``tick''.
         You don't have to worry about what kind of tick
         mark is used.

   \item This is the second item of the list.  It
         contains another list nested inside it.  The inner
         list is an \emph{enumerated} list.
         \begin{enumerate}
            \item This is the first item of an enumerated
                  list that is nested within the itemized list.

            \item This is the second item of the inner list.
                  \LaTeX\ allows you to nest lists deeper than
                  you really should.
         \end{enumerate}
         This is the rest of the second item of the outer
         list.  It is no more interesting than any other
         part of the item.
   \item This is the third item of the list.
\end{itemize}
You can even display poetry.
\begin{verse}
   There is an environment
    for verse \\             % The \\ command separates lines
   Whose features some poets % within a stanza.
   will curse.

                             % One or more blank lines separate stanzas.

   For instead of making\\
   Them do \emph{all} line breaking, \\
   It allows them to put too many words on a line when they'd rather be
   forced to be terse.
\end{verse}

Mathematical formulas may also be displayed.  A
displayed formula
is
one-line long; multiline
formulas require special formatting instructions.
   \[  \ip{\Gamma}{\psi'} = x'' + y^{2} + z_{i}^{n}\]
Don't start a paragraph with a displayed equation,
nor make one a paragraph by itself.

\end{document}               % End of document.

\section{Bench Huge Section 388}
% synthetic bench section 388
% This is a sample LaTeX input file.  (Version of 12 August 2004.)
%
% A '%' character causes TeX to ignore all remaining text on the line,
% and is used for comments like this one.

\documentclass{article}      % Specifies the document class

                             % The preamble begins here.
\title{An Example Document}  % Declares the document's title.
\author{Leslie Lamport}      % Declares the author's name.
\date{January 21, 1994}      % Deleting this command produces today's date.

\newcommand{\ip}[2]{(#1, #2)}
                             % Defines \ip{arg1}{arg2} to mean
                             % (arg1, arg2).

%\newcommand{\ip}[2]{\langle #1 | #2\rangle}
                             % This is an alternative definition of
                             % \ip that is commented out.

\begin{document}             % End of preamble and beginning of text.

\maketitle                   % Produces the title.

This is an example input file.  Comparing it with
the output it generates can show you how to
produce a simple document of your own.

\section{Ordinary Text}      % Produces section heading.  Lower-level
                             % sections are begun with similar
                             % \subsection and \subsubsection commands.

The ends  of words and sentences are marked
  by   spaces. It  doesn't matter how many
spaces    you type; one is as good as 100.  The
end of   a line counts as a space.

One   or more   blank lines denote the  end
of  a paragraph.

Since any number of consecutive spaces are treated
like a single one, the formatting of the input
file makes no difference to
      \LaTeX,                % The \LaTeX command generates the LaTeX logo.
but it makes a difference to you.  When you use
\LaTeX, making your input file as easy to read
as possible will be a great help as you write
your document and when you change it.  This sample
file shows how you can add comments to your own input
file.

Because printing is different from typewriting,
there are a number of things that you have to do
differently when preparing an input file than if
you were just typing the document directly.
Quotation marks like
       ``this''
have to be handled specially, as do quotes within
quotes:
       ``\,`this'            % \, separates the double and single quote.
        is what I just
        wrote, not  `that'\,''.

Dashes come in three sizes: an
       intra-word
dash, a medium dash for number ranges like
       1--2,
and a punctuation
       dash---like
this.

A sentence-ending space should be larger than the
space between words within a sentence.  You
sometimes have to type special commands in
conjunction with punctuation characters to get
this right, as in the following sentence.
       Gnats, gnus, etc.\ all  % `\ ' makes an inter-word space.
       begin with G\@.         % \@ marks end-of-sentence punctuation.
You should check the spaces after periods when
reading your output to make sure you haven't
forgotten any special cases.  Generating an
ellipsis
       \ldots\               % `\ ' is needed after `\ldots' because TeX
                             % ignores spaces after command names like \ldots
                             % made from \ + letters.
                             %
                             % Note how a `%' character causes TeX to ignore
                             % the end of the input line, so these blank lines
                             % do not start a new paragraph.
                             %
with the right spacing around the periods requires
a special command.

\LaTeX\ interprets some common characters as
commands, so you must type special commands to
generate them.  These characters include the
following:
       \$ \& \% \# \{ and \}.

In printing, text is usually emphasized with an
       \emph{italic}
type style.

\begin{em}
   A long segment of text can also be emphasized
   in this way.  Text within such a segment can be
   given \emph{additional} emphasis.
\end{em}

It is sometimes necessary to prevent \LaTeX\ from
breaking a line where it might otherwise do so.
This may be at a space, as between the ``Mr.''\ and
``Jones'' in
       ``Mr.~Jones'',        % ~ produces an unbreakable interword space.
or within a word---especially when the word is a
symbol like
       \mbox{\emph{itemnum}}
that makes little sense when hyphenated across
lines.

Footnotes\footnote{This is an example of a footnote.}
pose no problem.

\LaTeX\ is good at typesetting mathematical formulas
like
       \( x-3y + z = 7 \)
or
       \( a_{1} > x^{2n} + y^{2n} > x' \)
or
       \( \ip{A}{B} = \sum_{i} a_{i} b_{i} \).
The spaces you type in a formula are
ignored.  Remember that a letter like
       $x$                   % $ ... $  and  \( ... \)  are equivalent
is a formula when it denotes a mathematical
symbol, and it should be typed as one.

\section{Displayed Text}

Text is displayed by indenting it from the left
margin.  Quotations are commonly displayed.  There
are short quotations
\begin{quote}
   This is a short quotation.  It consists of a
   single paragraph of text.  See how it is formatted.
\end{quote}
and longer ones.
\begin{quotation}
   This is a longer quotation.  It consists of two
   paragraphs of text, neither of which are
   particularly interesting.

   This is the second paragraph of the quotation.  It
   is just as dull as the first paragraph.
\end{quotation}
Another frequently-displayed structure is a list.
The following is an example of an \emph{itemized}
list.
\begin{itemize}
   \item This is the first item of an itemized list.
         Each item in the list is marked with a ``tick''.
         You don't have to worry about what kind of tick
         mark is used.

   \item This is the second item of the list.  It
         contains another list nested inside it.  The inner
         list is an \emph{enumerated} list.
         \begin{enumerate}
            \item This is the first item of an enumerated
                  list that is nested within the itemized list.

            \item This is the second item of the inner list.
                  \LaTeX\ allows you to nest lists deeper than
                  you really should.
         \end{enumerate}
         This is the rest of the second item of the outer
         list.  It is no more interesting than any other
         part of the item.
   \item This is the third item of the list.
\end{itemize}
You can even display poetry.
\begin{verse}
   There is an environment
    for verse \\             % The \\ command separates lines
   Whose features some poets % within a stanza.
   will curse.

                             % One or more blank lines separate stanzas.

   For instead of making\\
   Them do \emph{all} line breaking, \\
   It allows them to put too many words on a line when they'd rather be
   forced to be terse.
\end{verse}

Mathematical formulas may also be displayed.  A
displayed formula
is
one-line long; multiline
formulas require special formatting instructions.
   \[  \ip{\Gamma}{\psi'} = x'' + y^{2} + z_{i}^{n}\]
Don't start a paragraph with a displayed equation,
nor make one a paragraph by itself.

\end{document}               % End of document.

\section{Bench Huge Section 389}
% synthetic bench section 389
% This is a sample LaTeX input file.  (Version of 12 August 2004.)
%
% A '%' character causes TeX to ignore all remaining text on the line,
% and is used for comments like this one.

\documentclass{article}      % Specifies the document class

                             % The preamble begins here.
\title{An Example Document}  % Declares the document's title.
\author{Leslie Lamport}      % Declares the author's name.
\date{January 21, 1994}      % Deleting this command produces today's date.

\newcommand{\ip}[2]{(#1, #2)}
                             % Defines \ip{arg1}{arg2} to mean
                             % (arg1, arg2).

%\newcommand{\ip}[2]{\langle #1 | #2\rangle}
                             % This is an alternative definition of
                             % \ip that is commented out.

\begin{document}             % End of preamble and beginning of text.

\maketitle                   % Produces the title.

This is an example input file.  Comparing it with
the output it generates can show you how to
produce a simple document of your own.

\section{Ordinary Text}      % Produces section heading.  Lower-level
                             % sections are begun with similar
                             % \subsection and \subsubsection commands.

The ends  of words and sentences are marked
  by   spaces. It  doesn't matter how many
spaces    you type; one is as good as 100.  The
end of   a line counts as a space.

One   or more   blank lines denote the  end
of  a paragraph.

Since any number of consecutive spaces are treated
like a single one, the formatting of the input
file makes no difference to
      \LaTeX,                % The \LaTeX command generates the LaTeX logo.
but it makes a difference to you.  When you use
\LaTeX, making your input file as easy to read
as possible will be a great help as you write
your document and when you change it.  This sample
file shows how you can add comments to your own input
file.

Because printing is different from typewriting,
there are a number of things that you have to do
differently when preparing an input file than if
you were just typing the document directly.
Quotation marks like
       ``this''
have to be handled specially, as do quotes within
quotes:
       ``\,`this'            % \, separates the double and single quote.
        is what I just
        wrote, not  `that'\,''.

Dashes come in three sizes: an
       intra-word
dash, a medium dash for number ranges like
       1--2,
and a punctuation
       dash---like
this.

A sentence-ending space should be larger than the
space between words within a sentence.  You
sometimes have to type special commands in
conjunction with punctuation characters to get
this right, as in the following sentence.
       Gnats, gnus, etc.\ all  % `\ ' makes an inter-word space.
       begin with G\@.         % \@ marks end-of-sentence punctuation.
You should check the spaces after periods when
reading your output to make sure you haven't
forgotten any special cases.  Generating an
ellipsis
       \ldots\               % `\ ' is needed after `\ldots' because TeX
                             % ignores spaces after command names like \ldots
                             % made from \ + letters.
                             %
                             % Note how a `%' character causes TeX to ignore
                             % the end of the input line, so these blank lines
                             % do not start a new paragraph.
                             %
with the right spacing around the periods requires
a special command.

\LaTeX\ interprets some common characters as
commands, so you must type special commands to
generate them.  These characters include the
following:
       \$ \& \% \# \{ and \}.

In printing, text is usually emphasized with an
       \emph{italic}
type style.

\begin{em}
   A long segment of text can also be emphasized
   in this way.  Text within such a segment can be
   given \emph{additional} emphasis.
\end{em}

It is sometimes necessary to prevent \LaTeX\ from
breaking a line where it might otherwise do so.
This may be at a space, as between the ``Mr.''\ and
``Jones'' in
       ``Mr.~Jones'',        % ~ produces an unbreakable interword space.
or within a word---especially when the word is a
symbol like
       \mbox{\emph{itemnum}}
that makes little sense when hyphenated across
lines.

Footnotes\footnote{This is an example of a footnote.}
pose no problem.

\LaTeX\ is good at typesetting mathematical formulas
like
       \( x-3y + z = 7 \)
or
       \( a_{1} > x^{2n} + y^{2n} > x' \)
or
       \( \ip{A}{B} = \sum_{i} a_{i} b_{i} \).
The spaces you type in a formula are
ignored.  Remember that a letter like
       $x$                   % $ ... $  and  \( ... \)  are equivalent
is a formula when it denotes a mathematical
symbol, and it should be typed as one.

\section{Displayed Text}

Text is displayed by indenting it from the left
margin.  Quotations are commonly displayed.  There
are short quotations
\begin{quote}
   This is a short quotation.  It consists of a
   single paragraph of text.  See how it is formatted.
\end{quote}
and longer ones.
\begin{quotation}
   This is a longer quotation.  It consists of two
   paragraphs of text, neither of which are
   particularly interesting.

   This is the second paragraph of the quotation.  It
   is just as dull as the first paragraph.
\end{quotation}
Another frequently-displayed structure is a list.
The following is an example of an \emph{itemized}
list.
\begin{itemize}
   \item This is the first item of an itemized list.
         Each item in the list is marked with a ``tick''.
         You don't have to worry about what kind of tick
         mark is used.

   \item This is the second item of the list.  It
         contains another list nested inside it.  The inner
         list is an \emph{enumerated} list.
         \begin{enumerate}
            \item This is the first item of an enumerated
                  list that is nested within the itemized list.

            \item This is the second item of the inner list.
                  \LaTeX\ allows you to nest lists deeper than
                  you really should.
         \end{enumerate}
         This is the rest of the second item of the outer
         list.  It is no more interesting than any other
         part of the item.
   \item This is the third item of the list.
\end{itemize}
You can even display poetry.
\begin{verse}
   There is an environment
    for verse \\             % The \\ command separates lines
   Whose features some poets % within a stanza.
   will curse.

                             % One or more blank lines separate stanzas.

   For instead of making\\
   Them do \emph{all} line breaking, \\
   It allows them to put too many words on a line when they'd rather be
   forced to be terse.
\end{verse}

Mathematical formulas may also be displayed.  A
displayed formula
is
one-line long; multiline
formulas require special formatting instructions.
   \[  \ip{\Gamma}{\psi'} = x'' + y^{2} + z_{i}^{n}\]
Don't start a paragraph with a displayed equation,
nor make one a paragraph by itself.

\end{document}               % End of document.

\section{Bench Huge Section 390}
% synthetic bench section 390
% This is a sample LaTeX input file.  (Version of 12 August 2004.)
%
% A '%' character causes TeX to ignore all remaining text on the line,
% and is used for comments like this one.

\documentclass{article}      % Specifies the document class

                             % The preamble begins here.
\title{An Example Document}  % Declares the document's title.
\author{Leslie Lamport}      % Declares the author's name.
\date{January 21, 1994}      % Deleting this command produces today's date.

\newcommand{\ip}[2]{(#1, #2)}
                             % Defines \ip{arg1}{arg2} to mean
                             % (arg1, arg2).

%\newcommand{\ip}[2]{\langle #1 | #2\rangle}
                             % This is an alternative definition of
                             % \ip that is commented out.

\begin{document}             % End of preamble and beginning of text.

\maketitle                   % Produces the title.

This is an example input file.  Comparing it with
the output it generates can show you how to
produce a simple document of your own.

\section{Ordinary Text}      % Produces section heading.  Lower-level
                             % sections are begun with similar
                             % \subsection and \subsubsection commands.

The ends  of words and sentences are marked
  by   spaces. It  doesn't matter how many
spaces    you type; one is as good as 100.  The
end of   a line counts as a space.

One   or more   blank lines denote the  end
of  a paragraph.

Since any number of consecutive spaces are treated
like a single one, the formatting of the input
file makes no difference to
      \LaTeX,                % The \LaTeX command generates the LaTeX logo.
but it makes a difference to you.  When you use
\LaTeX, making your input file as easy to read
as possible will be a great help as you write
your document and when you change it.  This sample
file shows how you can add comments to your own input
file.

Because printing is different from typewriting,
there are a number of things that you have to do
differently when preparing an input file than if
you were just typing the document directly.
Quotation marks like
       ``this''
have to be handled specially, as do quotes within
quotes:
       ``\,`this'            % \, separates the double and single quote.
        is what I just
        wrote, not  `that'\,''.

Dashes come in three sizes: an
       intra-word
dash, a medium dash for number ranges like
       1--2,
and a punctuation
       dash---like
this.

A sentence-ending space should be larger than the
space between words within a sentence.  You
sometimes have to type special commands in
conjunction with punctuation characters to get
this right, as in the following sentence.
       Gnats, gnus, etc.\ all  % `\ ' makes an inter-word space.
       begin with G\@.         % \@ marks end-of-sentence punctuation.
You should check the spaces after periods when
reading your output to make sure you haven't
forgotten any special cases.  Generating an
ellipsis
       \ldots\               % `\ ' is needed after `\ldots' because TeX
                             % ignores spaces after command names like \ldots
                             % made from \ + letters.
                             %
                             % Note how a `%' character causes TeX to ignore
                             % the end of the input line, so these blank lines
                             % do not start a new paragraph.
                             %
with the right spacing around the periods requires
a special command.

\LaTeX\ interprets some common characters as
commands, so you must type special commands to
generate them.  These characters include the
following:
       \$ \& \% \# \{ and \}.

In printing, text is usually emphasized with an
       \emph{italic}
type style.

\begin{em}
   A long segment of text can also be emphasized
   in this way.  Text within such a segment can be
   given \emph{additional} emphasis.
\end{em}

It is sometimes necessary to prevent \LaTeX\ from
breaking a line where it might otherwise do so.
This may be at a space, as between the ``Mr.''\ and
``Jones'' in
       ``Mr.~Jones'',        % ~ produces an unbreakable interword space.
or within a word---especially when the word is a
symbol like
       \mbox{\emph{itemnum}}
that makes little sense when hyphenated across
lines.

Footnotes\footnote{This is an example of a footnote.}
pose no problem.

\LaTeX\ is good at typesetting mathematical formulas
like
       \( x-3y + z = 7 \)
or
       \( a_{1} > x^{2n} + y^{2n} > x' \)
or
       \( \ip{A}{B} = \sum_{i} a_{i} b_{i} \).
The spaces you type in a formula are
ignored.  Remember that a letter like
       $x$                   % $ ... $  and  \( ... \)  are equivalent
is a formula when it denotes a mathematical
symbol, and it should be typed as one.

\section{Displayed Text}

Text is displayed by indenting it from the left
margin.  Quotations are commonly displayed.  There
are short quotations
\begin{quote}
   This is a short quotation.  It consists of a
   single paragraph of text.  See how it is formatted.
\end{quote}
and longer ones.
\begin{quotation}
   This is a longer quotation.  It consists of two
   paragraphs of text, neither of which are
   particularly interesting.

   This is the second paragraph of the quotation.  It
   is just as dull as the first paragraph.
\end{quotation}
Another frequently-displayed structure is a list.
The following is an example of an \emph{itemized}
list.
\begin{itemize}
   \item This is the first item of an itemized list.
         Each item in the list is marked with a ``tick''.
         You don't have to worry about what kind of tick
         mark is used.

   \item This is the second item of the list.  It
         contains another list nested inside it.  The inner
         list is an \emph{enumerated} list.
         \begin{enumerate}
            \item This is the first item of an enumerated
                  list that is nested within the itemized list.

            \item This is the second item of the inner list.
                  \LaTeX\ allows you to nest lists deeper than
                  you really should.
         \end{enumerate}
         This is the rest of the second item of the outer
         list.  It is no more interesting than any other
         part of the item.
   \item This is the third item of the list.
\end{itemize}
You can even display poetry.
\begin{verse}
   There is an environment
    for verse \\             % The \\ command separates lines
   Whose features some poets % within a stanza.
   will curse.

                             % One or more blank lines separate stanzas.

   For instead of making\\
   Them do \emph{all} line breaking, \\
   It allows them to put too many words on a line when they'd rather be
   forced to be terse.
\end{verse}

Mathematical formulas may also be displayed.  A
displayed formula
is
one-line long; multiline
formulas require special formatting instructions.
   \[  \ip{\Gamma}{\psi'} = x'' + y^{2} + z_{i}^{n}\]
Don't start a paragraph with a displayed equation,
nor make one a paragraph by itself.

\end{document}               % End of document.

\section{Bench Huge Section 391}
% synthetic bench section 391
% This is a sample LaTeX input file.  (Version of 12 August 2004.)
%
% A '%' character causes TeX to ignore all remaining text on the line,
% and is used for comments like this one.

\documentclass{article}      % Specifies the document class

                             % The preamble begins here.
\title{An Example Document}  % Declares the document's title.
\author{Leslie Lamport}      % Declares the author's name.
\date{January 21, 1994}      % Deleting this command produces today's date.

\newcommand{\ip}[2]{(#1, #2)}
                             % Defines \ip{arg1}{arg2} to mean
                             % (arg1, arg2).

%\newcommand{\ip}[2]{\langle #1 | #2\rangle}
                             % This is an alternative definition of
                             % \ip that is commented out.

\begin{document}             % End of preamble and beginning of text.

\maketitle                   % Produces the title.

This is an example input file.  Comparing it with
the output it generates can show you how to
produce a simple document of your own.

\section{Ordinary Text}      % Produces section heading.  Lower-level
                             % sections are begun with similar
                             % \subsection and \subsubsection commands.

The ends  of words and sentences are marked
  by   spaces. It  doesn't matter how many
spaces    you type; one is as good as 100.  The
end of   a line counts as a space.

One   or more   blank lines denote the  end
of  a paragraph.

Since any number of consecutive spaces are treated
like a single one, the formatting of the input
file makes no difference to
      \LaTeX,                % The \LaTeX command generates the LaTeX logo.
but it makes a difference to you.  When you use
\LaTeX, making your input file as easy to read
as possible will be a great help as you write
your document and when you change it.  This sample
file shows how you can add comments to your own input
file.

Because printing is different from typewriting,
there are a number of things that you have to do
differently when preparing an input file than if
you were just typing the document directly.
Quotation marks like
       ``this''
have to be handled specially, as do quotes within
quotes:
       ``\,`this'            % \, separates the double and single quote.
        is what I just
        wrote, not  `that'\,''.

Dashes come in three sizes: an
       intra-word
dash, a medium dash for number ranges like
       1--2,
and a punctuation
       dash---like
this.

A sentence-ending space should be larger than the
space between words within a sentence.  You
sometimes have to type special commands in
conjunction with punctuation characters to get
this right, as in the following sentence.
       Gnats, gnus, etc.\ all  % `\ ' makes an inter-word space.
       begin with G\@.         % \@ marks end-of-sentence punctuation.
You should check the spaces after periods when
reading your output to make sure you haven't
forgotten any special cases.  Generating an
ellipsis
       \ldots\               % `\ ' is needed after `\ldots' because TeX
                             % ignores spaces after command names like \ldots
                             % made from \ + letters.
                             %
                             % Note how a `%' character causes TeX to ignore
                             % the end of the input line, so these blank lines
                             % do not start a new paragraph.
                             %
with the right spacing around the periods requires
a special command.

\LaTeX\ interprets some common characters as
commands, so you must type special commands to
generate them.  These characters include the
following:
       \$ \& \% \# \{ and \}.

In printing, text is usually emphasized with an
       \emph{italic}
type style.

\begin{em}
   A long segment of text can also be emphasized
   in this way.  Text within such a segment can be
   given \emph{additional} emphasis.
\end{em}

It is sometimes necessary to prevent \LaTeX\ from
breaking a line where it might otherwise do so.
This may be at a space, as between the ``Mr.''\ and
``Jones'' in
       ``Mr.~Jones'',        % ~ produces an unbreakable interword space.
or within a word---especially when the word is a
symbol like
       \mbox{\emph{itemnum}}
that makes little sense when hyphenated across
lines.

Footnotes\footnote{This is an example of a footnote.}
pose no problem.

\LaTeX\ is good at typesetting mathematical formulas
like
       \( x-3y + z = 7 \)
or
       \( a_{1} > x^{2n} + y^{2n} > x' \)
or
       \( \ip{A}{B} = \sum_{i} a_{i} b_{i} \).
The spaces you type in a formula are
ignored.  Remember that a letter like
       $x$                   % $ ... $  and  \( ... \)  are equivalent
is a formula when it denotes a mathematical
symbol, and it should be typed as one.

\section{Displayed Text}

Text is displayed by indenting it from the left
margin.  Quotations are commonly displayed.  There
are short quotations
\begin{quote}
   This is a short quotation.  It consists of a
   single paragraph of text.  See how it is formatted.
\end{quote}
and longer ones.
\begin{quotation}
   This is a longer quotation.  It consists of two
   paragraphs of text, neither of which are
   particularly interesting.

   This is the second paragraph of the quotation.  It
   is just as dull as the first paragraph.
\end{quotation}
Another frequently-displayed structure is a list.
The following is an example of an \emph{itemized}
list.
\begin{itemize}
   \item This is the first item of an itemized list.
         Each item in the list is marked with a ``tick''.
         You don't have to worry about what kind of tick
         mark is used.

   \item This is the second item of the list.  It
         contains another list nested inside it.  The inner
         list is an \emph{enumerated} list.
         \begin{enumerate}
            \item This is the first item of an enumerated
                  list that is nested within the itemized list.

            \item This is the second item of the inner list.
                  \LaTeX\ allows you to nest lists deeper than
                  you really should.
         \end{enumerate}
         This is the rest of the second item of the outer
         list.  It is no more interesting than any other
         part of the item.
   \item This is the third item of the list.
\end{itemize}
You can even display poetry.
\begin{verse}
   There is an environment
    for verse \\             % The \\ command separates lines
   Whose features some poets % within a stanza.
   will curse.

                             % One or more blank lines separate stanzas.

   For instead of making\\
   Them do \emph{all} line breaking, \\
   It allows them to put too many words on a line when they'd rather be
   forced to be terse.
\end{verse}

Mathematical formulas may also be displayed.  A
displayed formula
is
one-line long; multiline
formulas require special formatting instructions.
   \[  \ip{\Gamma}{\psi'} = x'' + y^{2} + z_{i}^{n}\]
Don't start a paragraph with a displayed equation,
nor make one a paragraph by itself.

\end{document}               % End of document.

\section{Bench Huge Section 392}
% synthetic bench section 392
% This is a sample LaTeX input file.  (Version of 12 August 2004.)
%
% A '%' character causes TeX to ignore all remaining text on the line,
% and is used for comments like this one.

\documentclass{article}      % Specifies the document class

                             % The preamble begins here.
\title{An Example Document}  % Declares the document's title.
\author{Leslie Lamport}      % Declares the author's name.
\date{January 21, 1994}      % Deleting this command produces today's date.

\newcommand{\ip}[2]{(#1, #2)}
                             % Defines \ip{arg1}{arg2} to mean
                             % (arg1, arg2).

%\newcommand{\ip}[2]{\langle #1 | #2\rangle}
                             % This is an alternative definition of
                             % \ip that is commented out.

\begin{document}             % End of preamble and beginning of text.

\maketitle                   % Produces the title.

This is an example input file.  Comparing it with
the output it generates can show you how to
produce a simple document of your own.

\section{Ordinary Text}      % Produces section heading.  Lower-level
                             % sections are begun with similar
                             % \subsection and \subsubsection commands.

The ends  of words and sentences are marked
  by   spaces. It  doesn't matter how many
spaces    you type; one is as good as 100.  The
end of   a line counts as a space.

One   or more   blank lines denote the  end
of  a paragraph.

Since any number of consecutive spaces are treated
like a single one, the formatting of the input
file makes no difference to
      \LaTeX,                % The \LaTeX command generates the LaTeX logo.
but it makes a difference to you.  When you use
\LaTeX, making your input file as easy to read
as possible will be a great help as you write
your document and when you change it.  This sample
file shows how you can add comments to your own input
file.

Because printing is different from typewriting,
there are a number of things that you have to do
differently when preparing an input file than if
you were just typing the document directly.
Quotation marks like
       ``this''
have to be handled specially, as do quotes within
quotes:
       ``\,`this'            % \, separates the double and single quote.
        is what I just
        wrote, not  `that'\,''.

Dashes come in three sizes: an
       intra-word
dash, a medium dash for number ranges like
       1--2,
and a punctuation
       dash---like
this.

A sentence-ending space should be larger than the
space between words within a sentence.  You
sometimes have to type special commands in
conjunction with punctuation characters to get
this right, as in the following sentence.
       Gnats, gnus, etc.\ all  % `\ ' makes an inter-word space.
       begin with G\@.         % \@ marks end-of-sentence punctuation.
You should check the spaces after periods when
reading your output to make sure you haven't
forgotten any special cases.  Generating an
ellipsis
       \ldots\               % `\ ' is needed after `\ldots' because TeX
                             % ignores spaces after command names like \ldots
                             % made from \ + letters.
                             %
                             % Note how a `%' character causes TeX to ignore
                             % the end of the input line, so these blank lines
                             % do not start a new paragraph.
                             %
with the right spacing around the periods requires
a special command.

\LaTeX\ interprets some common characters as
commands, so you must type special commands to
generate them.  These characters include the
following:
       \$ \& \% \# \{ and \}.

In printing, text is usually emphasized with an
       \emph{italic}
type style.

\begin{em}
   A long segment of text can also be emphasized
   in this way.  Text within such a segment can be
   given \emph{additional} emphasis.
\end{em}

It is sometimes necessary to prevent \LaTeX\ from
breaking a line where it might otherwise do so.
This may be at a space, as between the ``Mr.''\ and
``Jones'' in
       ``Mr.~Jones'',        % ~ produces an unbreakable interword space.
or within a word---especially when the word is a
symbol like
       \mbox{\emph{itemnum}}
that makes little sense when hyphenated across
lines.

Footnotes\footnote{This is an example of a footnote.}
pose no problem.

\LaTeX\ is good at typesetting mathematical formulas
like
       \( x-3y + z = 7 \)
or
       \( a_{1} > x^{2n} + y^{2n} > x' \)
or
       \( \ip{A}{B} = \sum_{i} a_{i} b_{i} \).
The spaces you type in a formula are
ignored.  Remember that a letter like
       $x$                   % $ ... $  and  \( ... \)  are equivalent
is a formula when it denotes a mathematical
symbol, and it should be typed as one.

\section{Displayed Text}

Text is displayed by indenting it from the left
margin.  Quotations are commonly displayed.  There
are short quotations
\begin{quote}
   This is a short quotation.  It consists of a
   single paragraph of text.  See how it is formatted.
\end{quote}
and longer ones.
\begin{quotation}
   This is a longer quotation.  It consists of two
   paragraphs of text, neither of which are
   particularly interesting.

   This is the second paragraph of the quotation.  It
   is just as dull as the first paragraph.
\end{quotation}
Another frequently-displayed structure is a list.
The following is an example of an \emph{itemized}
list.
\begin{itemize}
   \item This is the first item of an itemized list.
         Each item in the list is marked with a ``tick''.
         You don't have to worry about what kind of tick
         mark is used.

   \item This is the second item of the list.  It
         contains another list nested inside it.  The inner
         list is an \emph{enumerated} list.
         \begin{enumerate}
            \item This is the first item of an enumerated
                  list that is nested within the itemized list.

            \item This is the second item of the inner list.
                  \LaTeX\ allows you to nest lists deeper than
                  you really should.
         \end{enumerate}
         This is the rest of the second item of the outer
         list.  It is no more interesting than any other
         part of the item.
   \item This is the third item of the list.
\end{itemize}
You can even display poetry.
\begin{verse}
   There is an environment
    for verse \\             % The \\ command separates lines
   Whose features some poets % within a stanza.
   will curse.

                             % One or more blank lines separate stanzas.

   For instead of making\\
   Them do \emph{all} line breaking, \\
   It allows them to put too many words on a line when they'd rather be
   forced to be terse.
\end{verse}

Mathematical formulas may also be displayed.  A
displayed formula
is
one-line long; multiline
formulas require special formatting instructions.
   \[  \ip{\Gamma}{\psi'} = x'' + y^{2} + z_{i}^{n}\]
Don't start a paragraph with a displayed equation,
nor make one a paragraph by itself.

\end{document}               % End of document.

\section{Bench Huge Section 393}
% synthetic bench section 393
% This is a sample LaTeX input file.  (Version of 12 August 2004.)
%
% A '%' character causes TeX to ignore all remaining text on the line,
% and is used for comments like this one.

\documentclass{article}      % Specifies the document class

                             % The preamble begins here.
\title{An Example Document}  % Declares the document's title.
\author{Leslie Lamport}      % Declares the author's name.
\date{January 21, 1994}      % Deleting this command produces today's date.

\newcommand{\ip}[2]{(#1, #2)}
                             % Defines \ip{arg1}{arg2} to mean
                             % (arg1, arg2).

%\newcommand{\ip}[2]{\langle #1 | #2\rangle}
                             % This is an alternative definition of
                             % \ip that is commented out.

\begin{document}             % End of preamble and beginning of text.

\maketitle                   % Produces the title.

This is an example input file.  Comparing it with
the output it generates can show you how to
produce a simple document of your own.

\section{Ordinary Text}      % Produces section heading.  Lower-level
                             % sections are begun with similar
                             % \subsection and \subsubsection commands.

The ends  of words and sentences are marked
  by   spaces. It  doesn't matter how many
spaces    you type; one is as good as 100.  The
end of   a line counts as a space.

One   or more   blank lines denote the  end
of  a paragraph.

Since any number of consecutive spaces are treated
like a single one, the formatting of the input
file makes no difference to
      \LaTeX,                % The \LaTeX command generates the LaTeX logo.
but it makes a difference to you.  When you use
\LaTeX, making your input file as easy to read
as possible will be a great help as you write
your document and when you change it.  This sample
file shows how you can add comments to your own input
file.

Because printing is different from typewriting,
there are a number of things that you have to do
differently when preparing an input file than if
you were just typing the document directly.
Quotation marks like
       ``this''
have to be handled specially, as do quotes within
quotes:
       ``\,`this'            % \, separates the double and single quote.
        is what I just
        wrote, not  `that'\,''.

Dashes come in three sizes: an
       intra-word
dash, a medium dash for number ranges like
       1--2,
and a punctuation
       dash---like
this.

A sentence-ending space should be larger than the
space between words within a sentence.  You
sometimes have to type special commands in
conjunction with punctuation characters to get
this right, as in the following sentence.
       Gnats, gnus, etc.\ all  % `\ ' makes an inter-word space.
       begin with G\@.         % \@ marks end-of-sentence punctuation.
You should check the spaces after periods when
reading your output to make sure you haven't
forgotten any special cases.  Generating an
ellipsis
       \ldots\               % `\ ' is needed after `\ldots' because TeX
                             % ignores spaces after command names like \ldots
                             % made from \ + letters.
                             %
                             % Note how a `%' character causes TeX to ignore
                             % the end of the input line, so these blank lines
                             % do not start a new paragraph.
                             %
with the right spacing around the periods requires
a special command.

\LaTeX\ interprets some common characters as
commands, so you must type special commands to
generate them.  These characters include the
following:
       \$ \& \% \# \{ and \}.

In printing, text is usually emphasized with an
       \emph{italic}
type style.

\begin{em}
   A long segment of text can also be emphasized
   in this way.  Text within such a segment can be
   given \emph{additional} emphasis.
\end{em}

It is sometimes necessary to prevent \LaTeX\ from
breaking a line where it might otherwise do so.
This may be at a space, as between the ``Mr.''\ and
``Jones'' in
       ``Mr.~Jones'',        % ~ produces an unbreakable interword space.
or within a word---especially when the word is a
symbol like
       \mbox{\emph{itemnum}}
that makes little sense when hyphenated across
lines.

Footnotes\footnote{This is an example of a footnote.}
pose no problem.

\LaTeX\ is good at typesetting mathematical formulas
like
       \( x-3y + z = 7 \)
or
       \( a_{1} > x^{2n} + y^{2n} > x' \)
or
       \( \ip{A}{B} = \sum_{i} a_{i} b_{i} \).
The spaces you type in a formula are
ignored.  Remember that a letter like
       $x$                   % $ ... $  and  \( ... \)  are equivalent
is a formula when it denotes a mathematical
symbol, and it should be typed as one.

\section{Displayed Text}

Text is displayed by indenting it from the left
margin.  Quotations are commonly displayed.  There
are short quotations
\begin{quote}
   This is a short quotation.  It consists of a
   single paragraph of text.  See how it is formatted.
\end{quote}
and longer ones.
\begin{quotation}
   This is a longer quotation.  It consists of two
   paragraphs of text, neither of which are
   particularly interesting.

   This is the second paragraph of the quotation.  It
   is just as dull as the first paragraph.
\end{quotation}
Another frequently-displayed structure is a list.
The following is an example of an \emph{itemized}
list.
\begin{itemize}
   \item This is the first item of an itemized list.
         Each item in the list is marked with a ``tick''.
         You don't have to worry about what kind of tick
         mark is used.

   \item This is the second item of the list.  It
         contains another list nested inside it.  The inner
         list is an \emph{enumerated} list.
         \begin{enumerate}
            \item This is the first item of an enumerated
                  list that is nested within the itemized list.

            \item This is the second item of the inner list.
                  \LaTeX\ allows you to nest lists deeper than
                  you really should.
         \end{enumerate}
         This is the rest of the second item of the outer
         list.  It is no more interesting than any other
         part of the item.
   \item This is the third item of the list.
\end{itemize}
You can even display poetry.
\begin{verse}
   There is an environment
    for verse \\             % The \\ command separates lines
   Whose features some poets % within a stanza.
   will curse.

                             % One or more blank lines separate stanzas.

   For instead of making\\
   Them do \emph{all} line breaking, \\
   It allows them to put too many words on a line when they'd rather be
   forced to be terse.
\end{verse}

Mathematical formulas may also be displayed.  A
displayed formula
is
one-line long; multiline
formulas require special formatting instructions.
   \[  \ip{\Gamma}{\psi'} = x'' + y^{2} + z_{i}^{n}\]
Don't start a paragraph with a displayed equation,
nor make one a paragraph by itself.

\end{document}               % End of document.

\section{Bench Huge Section 394}
% synthetic bench section 394
% This is a sample LaTeX input file.  (Version of 12 August 2004.)
%
% A '%' character causes TeX to ignore all remaining text on the line,
% and is used for comments like this one.

\documentclass{article}      % Specifies the document class

                             % The preamble begins here.
\title{An Example Document}  % Declares the document's title.
\author{Leslie Lamport}      % Declares the author's name.
\date{January 21, 1994}      % Deleting this command produces today's date.

\newcommand{\ip}[2]{(#1, #2)}
                             % Defines \ip{arg1}{arg2} to mean
                             % (arg1, arg2).

%\newcommand{\ip}[2]{\langle #1 | #2\rangle}
                             % This is an alternative definition of
                             % \ip that is commented out.

\begin{document}             % End of preamble and beginning of text.

\maketitle                   % Produces the title.

This is an example input file.  Comparing it with
the output it generates can show you how to
produce a simple document of your own.

\section{Ordinary Text}      % Produces section heading.  Lower-level
                             % sections are begun with similar
                             % \subsection and \subsubsection commands.

The ends  of words and sentences are marked
  by   spaces. It  doesn't matter how many
spaces    you type; one is as good as 100.  The
end of   a line counts as a space.

One   or more   blank lines denote the  end
of  a paragraph.

Since any number of consecutive spaces are treated
like a single one, the formatting of the input
file makes no difference to
      \LaTeX,                % The \LaTeX command generates the LaTeX logo.
but it makes a difference to you.  When you use
\LaTeX, making your input file as easy to read
as possible will be a great help as you write
your document and when you change it.  This sample
file shows how you can add comments to your own input
file.

Because printing is different from typewriting,
there are a number of things that you have to do
differently when preparing an input file than if
you were just typing the document directly.
Quotation marks like
       ``this''
have to be handled specially, as do quotes within
quotes:
       ``\,`this'            % \, separates the double and single quote.
        is what I just
        wrote, not  `that'\,''.

Dashes come in three sizes: an
       intra-word
dash, a medium dash for number ranges like
       1--2,
and a punctuation
       dash---like
this.

A sentence-ending space should be larger than the
space between words within a sentence.  You
sometimes have to type special commands in
conjunction with punctuation characters to get
this right, as in the following sentence.
       Gnats, gnus, etc.\ all  % `\ ' makes an inter-word space.
       begin with G\@.         % \@ marks end-of-sentence punctuation.
You should check the spaces after periods when
reading your output to make sure you haven't
forgotten any special cases.  Generating an
ellipsis
       \ldots\               % `\ ' is needed after `\ldots' because TeX
                             % ignores spaces after command names like \ldots
                             % made from \ + letters.
                             %
                             % Note how a `%' character causes TeX to ignore
                             % the end of the input line, so these blank lines
                             % do not start a new paragraph.
                             %
with the right spacing around the periods requires
a special command.

\LaTeX\ interprets some common characters as
commands, so you must type special commands to
generate them.  These characters include the
following:
       \$ \& \% \# \{ and \}.

In printing, text is usually emphasized with an
       \emph{italic}
type style.

\begin{em}
   A long segment of text can also be emphasized
   in this way.  Text within such a segment can be
   given \emph{additional} emphasis.
\end{em}

It is sometimes necessary to prevent \LaTeX\ from
breaking a line where it might otherwise do so.
This may be at a space, as between the ``Mr.''\ and
``Jones'' in
       ``Mr.~Jones'',        % ~ produces an unbreakable interword space.
or within a word---especially when the word is a
symbol like
       \mbox{\emph{itemnum}}
that makes little sense when hyphenated across
lines.

Footnotes\footnote{This is an example of a footnote.}
pose no problem.

\LaTeX\ is good at typesetting mathematical formulas
like
       \( x-3y + z = 7 \)
or
       \( a_{1} > x^{2n} + y^{2n} > x' \)
or
       \( \ip{A}{B} = \sum_{i} a_{i} b_{i} \).
The spaces you type in a formula are
ignored.  Remember that a letter like
       $x$                   % $ ... $  and  \( ... \)  are equivalent
is a formula when it denotes a mathematical
symbol, and it should be typed as one.

\section{Displayed Text}

Text is displayed by indenting it from the left
margin.  Quotations are commonly displayed.  There
are short quotations
\begin{quote}
   This is a short quotation.  It consists of a
   single paragraph of text.  See how it is formatted.
\end{quote}
and longer ones.
\begin{quotation}
   This is a longer quotation.  It consists of two
   paragraphs of text, neither of which are
   particularly interesting.

   This is the second paragraph of the quotation.  It
   is just as dull as the first paragraph.
\end{quotation}
Another frequently-displayed structure is a list.
The following is an example of an \emph{itemized}
list.
\begin{itemize}
   \item This is the first item of an itemized list.
         Each item in the list is marked with a ``tick''.
         You don't have to worry about what kind of tick
         mark is used.

   \item This is the second item of the list.  It
         contains another list nested inside it.  The inner
         list is an \emph{enumerated} list.
         \begin{enumerate}
            \item This is the first item of an enumerated
                  list that is nested within the itemized list.

            \item This is the second item of the inner list.
                  \LaTeX\ allows you to nest lists deeper than
                  you really should.
         \end{enumerate}
         This is the rest of the second item of the outer
         list.  It is no more interesting than any other
         part of the item.
   \item This is the third item of the list.
\end{itemize}
You can even display poetry.
\begin{verse}
   There is an environment
    for verse \\             % The \\ command separates lines
   Whose features some poets % within a stanza.
   will curse.

                             % One or more blank lines separate stanzas.

   For instead of making\\
   Them do \emph{all} line breaking, \\
   It allows them to put too many words on a line when they'd rather be
   forced to be terse.
\end{verse}

Mathematical formulas may also be displayed.  A
displayed formula
is
one-line long; multiline
formulas require special formatting instructions.
   \[  \ip{\Gamma}{\psi'} = x'' + y^{2} + z_{i}^{n}\]
Don't start a paragraph with a displayed equation,
nor make one a paragraph by itself.

\end{document}               % End of document.

\section{Bench Huge Section 395}
% synthetic bench section 395
% This is a sample LaTeX input file.  (Version of 12 August 2004.)
%
% A '%' character causes TeX to ignore all remaining text on the line,
% and is used for comments like this one.

\documentclass{article}      % Specifies the document class

                             % The preamble begins here.
\title{An Example Document}  % Declares the document's title.
\author{Leslie Lamport}      % Declares the author's name.
\date{January 21, 1994}      % Deleting this command produces today's date.

\newcommand{\ip}[2]{(#1, #2)}
                             % Defines \ip{arg1}{arg2} to mean
                             % (arg1, arg2).

%\newcommand{\ip}[2]{\langle #1 | #2\rangle}
                             % This is an alternative definition of
                             % \ip that is commented out.

\begin{document}             % End of preamble and beginning of text.

\maketitle                   % Produces the title.

This is an example input file.  Comparing it with
the output it generates can show you how to
produce a simple document of your own.

\section{Ordinary Text}      % Produces section heading.  Lower-level
                             % sections are begun with similar
                             % \subsection and \subsubsection commands.

The ends  of words and sentences are marked
  by   spaces. It  doesn't matter how many
spaces    you type; one is as good as 100.  The
end of   a line counts as a space.

One   or more   blank lines denote the  end
of  a paragraph.

Since any number of consecutive spaces are treated
like a single one, the formatting of the input
file makes no difference to
      \LaTeX,                % The \LaTeX command generates the LaTeX logo.
but it makes a difference to you.  When you use
\LaTeX, making your input file as easy to read
as possible will be a great help as you write
your document and when you change it.  This sample
file shows how you can add comments to your own input
file.

Because printing is different from typewriting,
there are a number of things that you have to do
differently when preparing an input file than if
you were just typing the document directly.
Quotation marks like
       ``this''
have to be handled specially, as do quotes within
quotes:
       ``\,`this'            % \, separates the double and single quote.
        is what I just
        wrote, not  `that'\,''.

Dashes come in three sizes: an
       intra-word
dash, a medium dash for number ranges like
       1--2,
and a punctuation
       dash---like
this.

A sentence-ending space should be larger than the
space between words within a sentence.  You
sometimes have to type special commands in
conjunction with punctuation characters to get
this right, as in the following sentence.
       Gnats, gnus, etc.\ all  % `\ ' makes an inter-word space.
       begin with G\@.         % \@ marks end-of-sentence punctuation.
You should check the spaces after periods when
reading your output to make sure you haven't
forgotten any special cases.  Generating an
ellipsis
       \ldots\               % `\ ' is needed after `\ldots' because TeX
                             % ignores spaces after command names like \ldots
                             % made from \ + letters.
                             %
                             % Note how a `%' character causes TeX to ignore
                             % the end of the input line, so these blank lines
                             % do not start a new paragraph.
                             %
with the right spacing around the periods requires
a special command.

\LaTeX\ interprets some common characters as
commands, so you must type special commands to
generate them.  These characters include the
following:
       \$ \& \% \# \{ and \}.

In printing, text is usually emphasized with an
       \emph{italic}
type style.

\begin{em}
   A long segment of text can also be emphasized
   in this way.  Text within such a segment can be
   given \emph{additional} emphasis.
\end{em}

It is sometimes necessary to prevent \LaTeX\ from
breaking a line where it might otherwise do so.
This may be at a space, as between the ``Mr.''\ and
``Jones'' in
       ``Mr.~Jones'',        % ~ produces an unbreakable interword space.
or within a word---especially when the word is a
symbol like
       \mbox{\emph{itemnum}}
that makes little sense when hyphenated across
lines.

Footnotes\footnote{This is an example of a footnote.}
pose no problem.

\LaTeX\ is good at typesetting mathematical formulas
like
       \( x-3y + z = 7 \)
or
       \( a_{1} > x^{2n} + y^{2n} > x' \)
or
       \( \ip{A}{B} = \sum_{i} a_{i} b_{i} \).
The spaces you type in a formula are
ignored.  Remember that a letter like
       $x$                   % $ ... $  and  \( ... \)  are equivalent
is a formula when it denotes a mathematical
symbol, and it should be typed as one.

\section{Displayed Text}

Text is displayed by indenting it from the left
margin.  Quotations are commonly displayed.  There
are short quotations
\begin{quote}
   This is a short quotation.  It consists of a
   single paragraph of text.  See how it is formatted.
\end{quote}
and longer ones.
\begin{quotation}
   This is a longer quotation.  It consists of two
   paragraphs of text, neither of which are
   particularly interesting.

   This is the second paragraph of the quotation.  It
   is just as dull as the first paragraph.
\end{quotation}
Another frequently-displayed structure is a list.
The following is an example of an \emph{itemized}
list.
\begin{itemize}
   \item This is the first item of an itemized list.
         Each item in the list is marked with a ``tick''.
         You don't have to worry about what kind of tick
         mark is used.

   \item This is the second item of the list.  It
         contains another list nested inside it.  The inner
         list is an \emph{enumerated} list.
         \begin{enumerate}
            \item This is the first item of an enumerated
                  list that is nested within the itemized list.

            \item This is the second item of the inner list.
                  \LaTeX\ allows you to nest lists deeper than
                  you really should.
         \end{enumerate}
         This is the rest of the second item of the outer
         list.  It is no more interesting than any other
         part of the item.
   \item This is the third item of the list.
\end{itemize}
You can even display poetry.
\begin{verse}
   There is an environment
    for verse \\             % The \\ command separates lines
   Whose features some poets % within a stanza.
   will curse.

                             % One or more blank lines separate stanzas.

   For instead of making\\
   Them do \emph{all} line breaking, \\
   It allows them to put too many words on a line when they'd rather be
   forced to be terse.
\end{verse}

Mathematical formulas may also be displayed.  A
displayed formula
is
one-line long; multiline
formulas require special formatting instructions.
   \[  \ip{\Gamma}{\psi'} = x'' + y^{2} + z_{i}^{n}\]
Don't start a paragraph with a displayed equation,
nor make one a paragraph by itself.

\end{document}               % End of document.

\section{Bench Huge Section 396}
% synthetic bench section 396
% This is a sample LaTeX input file.  (Version of 12 August 2004.)
%
% A '%' character causes TeX to ignore all remaining text on the line,
% and is used for comments like this one.

\documentclass{article}      % Specifies the document class

                             % The preamble begins here.
\title{An Example Document}  % Declares the document's title.
\author{Leslie Lamport}      % Declares the author's name.
\date{January 21, 1994}      % Deleting this command produces today's date.

\newcommand{\ip}[2]{(#1, #2)}
                             % Defines \ip{arg1}{arg2} to mean
                             % (arg1, arg2).

%\newcommand{\ip}[2]{\langle #1 | #2\rangle}
                             % This is an alternative definition of
                             % \ip that is commented out.

\begin{document}             % End of preamble and beginning of text.

\maketitle                   % Produces the title.

This is an example input file.  Comparing it with
the output it generates can show you how to
produce a simple document of your own.

\section{Ordinary Text}      % Produces section heading.  Lower-level
                             % sections are begun with similar
                             % \subsection and \subsubsection commands.

The ends  of words and sentences are marked
  by   spaces. It  doesn't matter how many
spaces    you type; one is as good as 100.  The
end of   a line counts as a space.

One   or more   blank lines denote the  end
of  a paragraph.

Since any number of consecutive spaces are treated
like a single one, the formatting of the input
file makes no difference to
      \LaTeX,                % The \LaTeX command generates the LaTeX logo.
but it makes a difference to you.  When you use
\LaTeX, making your input file as easy to read
as possible will be a great help as you write
your document and when you change it.  This sample
file shows how you can add comments to your own input
file.

Because printing is different from typewriting,
there are a number of things that you have to do
differently when preparing an input file than if
you were just typing the document directly.
Quotation marks like
       ``this''
have to be handled specially, as do quotes within
quotes:
       ``\,`this'            % \, separates the double and single quote.
        is what I just
        wrote, not  `that'\,''.

Dashes come in three sizes: an
       intra-word
dash, a medium dash for number ranges like
       1--2,
and a punctuation
       dash---like
this.

A sentence-ending space should be larger than the
space between words within a sentence.  You
sometimes have to type special commands in
conjunction with punctuation characters to get
this right, as in the following sentence.
       Gnats, gnus, etc.\ all  % `\ ' makes an inter-word space.
       begin with G\@.         % \@ marks end-of-sentence punctuation.
You should check the spaces after periods when
reading your output to make sure you haven't
forgotten any special cases.  Generating an
ellipsis
       \ldots\               % `\ ' is needed after `\ldots' because TeX
                             % ignores spaces after command names like \ldots
                             % made from \ + letters.
                             %
                             % Note how a `%' character causes TeX to ignore
                             % the end of the input line, so these blank lines
                             % do not start a new paragraph.
                             %
with the right spacing around the periods requires
a special command.

\LaTeX\ interprets some common characters as
commands, so you must type special commands to
generate them.  These characters include the
following:
       \$ \& \% \# \{ and \}.

In printing, text is usually emphasized with an
       \emph{italic}
type style.

\begin{em}
   A long segment of text can also be emphasized
   in this way.  Text within such a segment can be
   given \emph{additional} emphasis.
\end{em}

It is sometimes necessary to prevent \LaTeX\ from
breaking a line where it might otherwise do so.
This may be at a space, as between the ``Mr.''\ and
``Jones'' in
       ``Mr.~Jones'',        % ~ produces an unbreakable interword space.
or within a word---especially when the word is a
symbol like
       \mbox{\emph{itemnum}}
that makes little sense when hyphenated across
lines.

Footnotes\footnote{This is an example of a footnote.}
pose no problem.

\LaTeX\ is good at typesetting mathematical formulas
like
       \( x-3y + z = 7 \)
or
       \( a_{1} > x^{2n} + y^{2n} > x' \)
or
       \( \ip{A}{B} = \sum_{i} a_{i} b_{i} \).
The spaces you type in a formula are
ignored.  Remember that a letter like
       $x$                   % $ ... $  and  \( ... \)  are equivalent
is a formula when it denotes a mathematical
symbol, and it should be typed as one.

\section{Displayed Text}

Text is displayed by indenting it from the left
margin.  Quotations are commonly displayed.  There
are short quotations
\begin{quote}
   This is a short quotation.  It consists of a
   single paragraph of text.  See how it is formatted.
\end{quote}
and longer ones.
\begin{quotation}
   This is a longer quotation.  It consists of two
   paragraphs of text, neither of which are
   particularly interesting.

   This is the second paragraph of the quotation.  It
   is just as dull as the first paragraph.
\end{quotation}
Another frequently-displayed structure is a list.
The following is an example of an \emph{itemized}
list.
\begin{itemize}
   \item This is the first item of an itemized list.
         Each item in the list is marked with a ``tick''.
         You don't have to worry about what kind of tick
         mark is used.

   \item This is the second item of the list.  It
         contains another list nested inside it.  The inner
         list is an \emph{enumerated} list.
         \begin{enumerate}
            \item This is the first item of an enumerated
                  list that is nested within the itemized list.

            \item This is the second item of the inner list.
                  \LaTeX\ allows you to nest lists deeper than
                  you really should.
         \end{enumerate}
         This is the rest of the second item of the outer
         list.  It is no more interesting than any other
         part of the item.
   \item This is the third item of the list.
\end{itemize}
You can even display poetry.
\begin{verse}
   There is an environment
    for verse \\             % The \\ command separates lines
   Whose features some poets % within a stanza.
   will curse.

                             % One or more blank lines separate stanzas.

   For instead of making\\
   Them do \emph{all} line breaking, \\
   It allows them to put too many words on a line when they'd rather be
   forced to be terse.
\end{verse}

Mathematical formulas may also be displayed.  A
displayed formula
is
one-line long; multiline
formulas require special formatting instructions.
   \[  \ip{\Gamma}{\psi'} = x'' + y^{2} + z_{i}^{n}\]
Don't start a paragraph with a displayed equation,
nor make one a paragraph by itself.

\end{document}               % End of document.

\section{Bench Huge Section 397}
% synthetic bench section 397
% This is a sample LaTeX input file.  (Version of 12 August 2004.)
%
% A '%' character causes TeX to ignore all remaining text on the line,
% and is used for comments like this one.

\documentclass{article}      % Specifies the document class

                             % The preamble begins here.
\title{An Example Document}  % Declares the document's title.
\author{Leslie Lamport}      % Declares the author's name.
\date{January 21, 1994}      % Deleting this command produces today's date.

\newcommand{\ip}[2]{(#1, #2)}
                             % Defines \ip{arg1}{arg2} to mean
                             % (arg1, arg2).

%\newcommand{\ip}[2]{\langle #1 | #2\rangle}
                             % This is an alternative definition of
                             % \ip that is commented out.

\begin{document}             % End of preamble and beginning of text.

\maketitle                   % Produces the title.

This is an example input file.  Comparing it with
the output it generates can show you how to
produce a simple document of your own.

\section{Ordinary Text}      % Produces section heading.  Lower-level
                             % sections are begun with similar
                             % \subsection and \subsubsection commands.

The ends  of words and sentences are marked
  by   spaces. It  doesn't matter how many
spaces    you type; one is as good as 100.  The
end of   a line counts as a space.

One   or more   blank lines denote the  end
of  a paragraph.

Since any number of consecutive spaces are treated
like a single one, the formatting of the input
file makes no difference to
      \LaTeX,                % The \LaTeX command generates the LaTeX logo.
but it makes a difference to you.  When you use
\LaTeX, making your input file as easy to read
as possible will be a great help as you write
your document and when you change it.  This sample
file shows how you can add comments to your own input
file.

Because printing is different from typewriting,
there are a number of things that you have to do
differently when preparing an input file than if
you were just typing the document directly.
Quotation marks like
       ``this''
have to be handled specially, as do quotes within
quotes:
       ``\,`this'            % \, separates the double and single quote.
        is what I just
        wrote, not  `that'\,''.

Dashes come in three sizes: an
       intra-word
dash, a medium dash for number ranges like
       1--2,
and a punctuation
       dash---like
this.

A sentence-ending space should be larger than the
space between words within a sentence.  You
sometimes have to type special commands in
conjunction with punctuation characters to get
this right, as in the following sentence.
       Gnats, gnus, etc.\ all  % `\ ' makes an inter-word space.
       begin with G\@.         % \@ marks end-of-sentence punctuation.
You should check the spaces after periods when
reading your output to make sure you haven't
forgotten any special cases.  Generating an
ellipsis
       \ldots\               % `\ ' is needed after `\ldots' because TeX
                             % ignores spaces after command names like \ldots
                             % made from \ + letters.
                             %
                             % Note how a `%' character causes TeX to ignore
                             % the end of the input line, so these blank lines
                             % do not start a new paragraph.
                             %
with the right spacing around the periods requires
a special command.

\LaTeX\ interprets some common characters as
commands, so you must type special commands to
generate them.  These characters include the
following:
       \$ \& \% \# \{ and \}.

In printing, text is usually emphasized with an
       \emph{italic}
type style.

\begin{em}
   A long segment of text can also be emphasized
   in this way.  Text within such a segment can be
   given \emph{additional} emphasis.
\end{em}

It is sometimes necessary to prevent \LaTeX\ from
breaking a line where it might otherwise do so.
This may be at a space, as between the ``Mr.''\ and
``Jones'' in
       ``Mr.~Jones'',        % ~ produces an unbreakable interword space.
or within a word---especially when the word is a
symbol like
       \mbox{\emph{itemnum}}
that makes little sense when hyphenated across
lines.

Footnotes\footnote{This is an example of a footnote.}
pose no problem.

\LaTeX\ is good at typesetting mathematical formulas
like
       \( x-3y + z = 7 \)
or
       \( a_{1} > x^{2n} + y^{2n} > x' \)
or
       \( \ip{A}{B} = \sum_{i} a_{i} b_{i} \).
The spaces you type in a formula are
ignored.  Remember that a letter like
       $x$                   % $ ... $  and  \( ... \)  are equivalent
is a formula when it denotes a mathematical
symbol, and it should be typed as one.

\section{Displayed Text}

Text is displayed by indenting it from the left
margin.  Quotations are commonly displayed.  There
are short quotations
\begin{quote}
   This is a short quotation.  It consists of a
   single paragraph of text.  See how it is formatted.
\end{quote}
and longer ones.
\begin{quotation}
   This is a longer quotation.  It consists of two
   paragraphs of text, neither of which are
   particularly interesting.

   This is the second paragraph of the quotation.  It
   is just as dull as the first paragraph.
\end{quotation}
Another frequently-displayed structure is a list.
The following is an example of an \emph{itemized}
list.
\begin{itemize}
   \item This is the first item of an itemized list.
         Each item in the list is marked with a ``tick''.
         You don't have to worry about what kind of tick
         mark is used.

   \item This is the second item of the list.  It
         contains another list nested inside it.  The inner
         list is an \emph{enumerated} list.
         \begin{enumerate}
            \item This is the first item of an enumerated
                  list that is nested within the itemized list.

            \item This is the second item of the inner list.
                  \LaTeX\ allows you to nest lists deeper than
                  you really should.
         \end{enumerate}
         This is the rest of the second item of the outer
         list.  It is no more interesting than any other
         part of the item.
   \item This is the third item of the list.
\end{itemize}
You can even display poetry.
\begin{verse}
   There is an environment
    for verse \\             % The \\ command separates lines
   Whose features some poets % within a stanza.
   will curse.

                             % One or more blank lines separate stanzas.

   For instead of making\\
   Them do \emph{all} line breaking, \\
   It allows them to put too many words on a line when they'd rather be
   forced to be terse.
\end{verse}

Mathematical formulas may also be displayed.  A
displayed formula
is
one-line long; multiline
formulas require special formatting instructions.
   \[  \ip{\Gamma}{\psi'} = x'' + y^{2} + z_{i}^{n}\]
Don't start a paragraph with a displayed equation,
nor make one a paragraph by itself.

\end{document}               % End of document.

\section{Bench Huge Section 398}
% synthetic bench section 398
% This is a sample LaTeX input file.  (Version of 12 August 2004.)
%
% A '%' character causes TeX to ignore all remaining text on the line,
% and is used for comments like this one.

\documentclass{article}      % Specifies the document class

                             % The preamble begins here.
\title{An Example Document}  % Declares the document's title.
\author{Leslie Lamport}      % Declares the author's name.
\date{January 21, 1994}      % Deleting this command produces today's date.

\newcommand{\ip}[2]{(#1, #2)}
                             % Defines \ip{arg1}{arg2} to mean
                             % (arg1, arg2).

%\newcommand{\ip}[2]{\langle #1 | #2\rangle}
                             % This is an alternative definition of
                             % \ip that is commented out.

\begin{document}             % End of preamble and beginning of text.

\maketitle                   % Produces the title.

This is an example input file.  Comparing it with
the output it generates can show you how to
produce a simple document of your own.

\section{Ordinary Text}      % Produces section heading.  Lower-level
                             % sections are begun with similar
                             % \subsection and \subsubsection commands.

The ends  of words and sentences are marked
  by   spaces. It  doesn't matter how many
spaces    you type; one is as good as 100.  The
end of   a line counts as a space.

One   or more   blank lines denote the  end
of  a paragraph.

Since any number of consecutive spaces are treated
like a single one, the formatting of the input
file makes no difference to
      \LaTeX,                % The \LaTeX command generates the LaTeX logo.
but it makes a difference to you.  When you use
\LaTeX, making your input file as easy to read
as possible will be a great help as you write
your document and when you change it.  This sample
file shows how you can add comments to your own input
file.

Because printing is different from typewriting,
there are a number of things that you have to do
differently when preparing an input file than if
you were just typing the document directly.
Quotation marks like
       ``this''
have to be handled specially, as do quotes within
quotes:
       ``\,`this'            % \, separates the double and single quote.
        is what I just
        wrote, not  `that'\,''.

Dashes come in three sizes: an
       intra-word
dash, a medium dash for number ranges like
       1--2,
and a punctuation
       dash---like
this.

A sentence-ending space should be larger than the
space between words within a sentence.  You
sometimes have to type special commands in
conjunction with punctuation characters to get
this right, as in the following sentence.
       Gnats, gnus, etc.\ all  % `\ ' makes an inter-word space.
       begin with G\@.         % \@ marks end-of-sentence punctuation.
You should check the spaces after periods when
reading your output to make sure you haven't
forgotten any special cases.  Generating an
ellipsis
       \ldots\               % `\ ' is needed after `\ldots' because TeX
                             % ignores spaces after command names like \ldots
                             % made from \ + letters.
                             %
                             % Note how a `%' character causes TeX to ignore
                             % the end of the input line, so these blank lines
                             % do not start a new paragraph.
                             %
with the right spacing around the periods requires
a special command.

\LaTeX\ interprets some common characters as
commands, so you must type special commands to
generate them.  These characters include the
following:
       \$ \& \% \# \{ and \}.

In printing, text is usually emphasized with an
       \emph{italic}
type style.

\begin{em}
   A long segment of text can also be emphasized
   in this way.  Text within such a segment can be
   given \emph{additional} emphasis.
\end{em}

It is sometimes necessary to prevent \LaTeX\ from
breaking a line where it might otherwise do so.
This may be at a space, as between the ``Mr.''\ and
``Jones'' in
       ``Mr.~Jones'',        % ~ produces an unbreakable interword space.
or within a word---especially when the word is a
symbol like
       \mbox{\emph{itemnum}}
that makes little sense when hyphenated across
lines.

Footnotes\footnote{This is an example of a footnote.}
pose no problem.

\LaTeX\ is good at typesetting mathematical formulas
like
       \( x-3y + z = 7 \)
or
       \( a_{1} > x^{2n} + y^{2n} > x' \)
or
       \( \ip{A}{B} = \sum_{i} a_{i} b_{i} \).
The spaces you type in a formula are
ignored.  Remember that a letter like
       $x$                   % $ ... $  and  \( ... \)  are equivalent
is a formula when it denotes a mathematical
symbol, and it should be typed as one.

\section{Displayed Text}

Text is displayed by indenting it from the left
margin.  Quotations are commonly displayed.  There
are short quotations
\begin{quote}
   This is a short quotation.  It consists of a
   single paragraph of text.  See how it is formatted.
\end{quote}
and longer ones.
\begin{quotation}
   This is a longer quotation.  It consists of two
   paragraphs of text, neither of which are
   particularly interesting.

   This is the second paragraph of the quotation.  It
   is just as dull as the first paragraph.
\end{quotation}
Another frequently-displayed structure is a list.
The following is an example of an \emph{itemized}
list.
\begin{itemize}
   \item This is the first item of an itemized list.
         Each item in the list is marked with a ``tick''.
         You don't have to worry about what kind of tick
         mark is used.

   \item This is the second item of the list.  It
         contains another list nested inside it.  The inner
         list is an \emph{enumerated} list.
         \begin{enumerate}
            \item This is the first item of an enumerated
                  list that is nested within the itemized list.

            \item This is the second item of the inner list.
                  \LaTeX\ allows you to nest lists deeper than
                  you really should.
         \end{enumerate}
         This is the rest of the second item of the outer
         list.  It is no more interesting than any other
         part of the item.
   \item This is the third item of the list.
\end{itemize}
You can even display poetry.
\begin{verse}
   There is an environment
    for verse \\             % The \\ command separates lines
   Whose features some poets % within a stanza.
   will curse.

                             % One or more blank lines separate stanzas.

   For instead of making\\
   Them do \emph{all} line breaking, \\
   It allows them to put too many words on a line when they'd rather be
   forced to be terse.
\end{verse}

Mathematical formulas may also be displayed.  A
displayed formula
is
one-line long; multiline
formulas require special formatting instructions.
   \[  \ip{\Gamma}{\psi'} = x'' + y^{2} + z_{i}^{n}\]
Don't start a paragraph with a displayed equation,
nor make one a paragraph by itself.

\end{document}               % End of document.

\section{Bench Huge Section 399}
% synthetic bench section 399
% This is a sample LaTeX input file.  (Version of 12 August 2004.)
%
% A '%' character causes TeX to ignore all remaining text on the line,
% and is used for comments like this one.

\documentclass{article}      % Specifies the document class

                             % The preamble begins here.
\title{An Example Document}  % Declares the document's title.
\author{Leslie Lamport}      % Declares the author's name.
\date{January 21, 1994}      % Deleting this command produces today's date.

\newcommand{\ip}[2]{(#1, #2)}
                             % Defines \ip{arg1}{arg2} to mean
                             % (arg1, arg2).

%\newcommand{\ip}[2]{\langle #1 | #2\rangle}
                             % This is an alternative definition of
                             % \ip that is commented out.

\begin{document}             % End of preamble and beginning of text.

\maketitle                   % Produces the title.

This is an example input file.  Comparing it with
the output it generates can show you how to
produce a simple document of your own.

\section{Ordinary Text}      % Produces section heading.  Lower-level
                             % sections are begun with similar
                             % \subsection and \subsubsection commands.

The ends  of words and sentences are marked
  by   spaces. It  doesn't matter how many
spaces    you type; one is as good as 100.  The
end of   a line counts as a space.

One   or more   blank lines denote the  end
of  a paragraph.

Since any number of consecutive spaces are treated
like a single one, the formatting of the input
file makes no difference to
      \LaTeX,                % The \LaTeX command generates the LaTeX logo.
but it makes a difference to you.  When you use
\LaTeX, making your input file as easy to read
as possible will be a great help as you write
your document and when you change it.  This sample
file shows how you can add comments to your own input
file.

Because printing is different from typewriting,
there are a number of things that you have to do
differently when preparing an input file than if
you were just typing the document directly.
Quotation marks like
       ``this''
have to be handled specially, as do quotes within
quotes:
       ``\,`this'            % \, separates the double and single quote.
        is what I just
        wrote, not  `that'\,''.

Dashes come in three sizes: an
       intra-word
dash, a medium dash for number ranges like
       1--2,
and a punctuation
       dash---like
this.

A sentence-ending space should be larger than the
space between words within a sentence.  You
sometimes have to type special commands in
conjunction with punctuation characters to get
this right, as in the following sentence.
       Gnats, gnus, etc.\ all  % `\ ' makes an inter-word space.
       begin with G\@.         % \@ marks end-of-sentence punctuation.
You should check the spaces after periods when
reading your output to make sure you haven't
forgotten any special cases.  Generating an
ellipsis
       \ldots\               % `\ ' is needed after `\ldots' because TeX
                             % ignores spaces after command names like \ldots
                             % made from \ + letters.
                             %
                             % Note how a `%' character causes TeX to ignore
                             % the end of the input line, so these blank lines
                             % do not start a new paragraph.
                             %
with the right spacing around the periods requires
a special command.

\LaTeX\ interprets some common characters as
commands, so you must type special commands to
generate them.  These characters include the
following:
       \$ \& \% \# \{ and \}.

In printing, text is usually emphasized with an
       \emph{italic}
type style.

\begin{em}
   A long segment of text can also be emphasized
   in this way.  Text within such a segment can be
   given \emph{additional} emphasis.
\end{em}

It is sometimes necessary to prevent \LaTeX\ from
breaking a line where it might otherwise do so.
This may be at a space, as between the ``Mr.''\ and
``Jones'' in
       ``Mr.~Jones'',        % ~ produces an unbreakable interword space.
or within a word---especially when the word is a
symbol like
       \mbox{\emph{itemnum}}
that makes little sense when hyphenated across
lines.

Footnotes\footnote{This is an example of a footnote.}
pose no problem.

\LaTeX\ is good at typesetting mathematical formulas
like
       \( x-3y + z = 7 \)
or
       \( a_{1} > x^{2n} + y^{2n} > x' \)
or
       \( \ip{A}{B} = \sum_{i} a_{i} b_{i} \).
The spaces you type in a formula are
ignored.  Remember that a letter like
       $x$                   % $ ... $  and  \( ... \)  are equivalent
is a formula when it denotes a mathematical
symbol, and it should be typed as one.

\section{Displayed Text}

Text is displayed by indenting it from the left
margin.  Quotations are commonly displayed.  There
are short quotations
\begin{quote}
   This is a short quotation.  It consists of a
   single paragraph of text.  See how it is formatted.
\end{quote}
and longer ones.
\begin{quotation}
   This is a longer quotation.  It consists of two
   paragraphs of text, neither of which are
   particularly interesting.

   This is the second paragraph of the quotation.  It
   is just as dull as the first paragraph.
\end{quotation}
Another frequently-displayed structure is a list.
The following is an example of an \emph{itemized}
list.
\begin{itemize}
   \item This is the first item of an itemized list.
         Each item in the list is marked with a ``tick''.
         You don't have to worry about what kind of tick
         mark is used.

   \item This is the second item of the list.  It
         contains another list nested inside it.  The inner
         list is an \emph{enumerated} list.
         \begin{enumerate}
            \item This is the first item of an enumerated
                  list that is nested within the itemized list.

            \item This is the second item of the inner list.
                  \LaTeX\ allows you to nest lists deeper than
                  you really should.
         \end{enumerate}
         This is the rest of the second item of the outer
         list.  It is no more interesting than any other
         part of the item.
   \item This is the third item of the list.
\end{itemize}
You can even display poetry.
\begin{verse}
   There is an environment
    for verse \\             % The \\ command separates lines
   Whose features some poets % within a stanza.
   will curse.

                             % One or more blank lines separate stanzas.

   For instead of making\\
   Them do \emph{all} line breaking, \\
   It allows them to put too many words on a line when they'd rather be
   forced to be terse.
\end{verse}

Mathematical formulas may also be displayed.  A
displayed formula
is
one-line long; multiline
formulas require special formatting instructions.
   \[  \ip{\Gamma}{\psi'} = x'' + y^{2} + z_{i}^{n}\]
Don't start a paragraph with a displayed equation,
nor make one a paragraph by itself.

\end{document}               % End of document.

\section{Bench Huge Section 400}
% synthetic bench section 400
% This is a sample LaTeX input file.  (Version of 12 August 2004.)
%
% A '%' character causes TeX to ignore all remaining text on the line,
% and is used for comments like this one.

\documentclass{article}      % Specifies the document class

                             % The preamble begins here.
\title{An Example Document}  % Declares the document's title.
\author{Leslie Lamport}      % Declares the author's name.
\date{January 21, 1994}      % Deleting this command produces today's date.

\newcommand{\ip}[2]{(#1, #2)}
                             % Defines \ip{arg1}{arg2} to mean
                             % (arg1, arg2).

%\newcommand{\ip}[2]{\langle #1 | #2\rangle}
                             % This is an alternative definition of
                             % \ip that is commented out.

\begin{document}             % End of preamble and beginning of text.

\maketitle                   % Produces the title.

This is an example input file.  Comparing it with
the output it generates can show you how to
produce a simple document of your own.

\section{Ordinary Text}      % Produces section heading.  Lower-level
                             % sections are begun with similar
                             % \subsection and \subsubsection commands.

The ends  of words and sentences are marked
  by   spaces. It  doesn't matter how many
spaces    you type; one is as good as 100.  The
end of   a line counts as a space.

One   or more   blank lines denote the  end
of  a paragraph.

Since any number of consecutive spaces are treated
like a single one, the formatting of the input
file makes no difference to
      \LaTeX,                % The \LaTeX command generates the LaTeX logo.
but it makes a difference to you.  When you use
\LaTeX, making your input file as easy to read
as possible will be a great help as you write
your document and when you change it.  This sample
file shows how you can add comments to your own input
file.

Because printing is different from typewriting,
there are a number of things that you have to do
differently when preparing an input file than if
you were just typing the document directly.
Quotation marks like
       ``this''
have to be handled specially, as do quotes within
quotes:
       ``\,`this'            % \, separates the double and single quote.
        is what I just
        wrote, not  `that'\,''.

Dashes come in three sizes: an
       intra-word
dash, a medium dash for number ranges like
       1--2,
and a punctuation
       dash---like
this.

A sentence-ending space should be larger than the
space between words within a sentence.  You
sometimes have to type special commands in
conjunction with punctuation characters to get
this right, as in the following sentence.
       Gnats, gnus, etc.\ all  % `\ ' makes an inter-word space.
       begin with G\@.         % \@ marks end-of-sentence punctuation.
You should check the spaces after periods when
reading your output to make sure you haven't
forgotten any special cases.  Generating an
ellipsis
       \ldots\               % `\ ' is needed after `\ldots' because TeX
                             % ignores spaces after command names like \ldots
                             % made from \ + letters.
                             %
                             % Note how a `%' character causes TeX to ignore
                             % the end of the input line, so these blank lines
                             % do not start a new paragraph.
                             %
with the right spacing around the periods requires
a special command.

\LaTeX\ interprets some common characters as
commands, so you must type special commands to
generate them.  These characters include the
following:
       \$ \& \% \# \{ and \}.

In printing, text is usually emphasized with an
       \emph{italic}
type style.

\begin{em}
   A long segment of text can also be emphasized
   in this way.  Text within such a segment can be
   given \emph{additional} emphasis.
\end{em}

It is sometimes necessary to prevent \LaTeX\ from
breaking a line where it might otherwise do so.
This may be at a space, as between the ``Mr.''\ and
``Jones'' in
       ``Mr.~Jones'',        % ~ produces an unbreakable interword space.
or within a word---especially when the word is a
symbol like
       \mbox{\emph{itemnum}}
that makes little sense when hyphenated across
lines.

Footnotes\footnote{This is an example of a footnote.}
pose no problem.

\LaTeX\ is good at typesetting mathematical formulas
like
       \( x-3y + z = 7 \)
or
       \( a_{1} > x^{2n} + y^{2n} > x' \)
or
       \( \ip{A}{B} = \sum_{i} a_{i} b_{i} \).
The spaces you type in a formula are
ignored.  Remember that a letter like
       $x$                   % $ ... $  and  \( ... \)  are equivalent
is a formula when it denotes a mathematical
symbol, and it should be typed as one.

\section{Displayed Text}

Text is displayed by indenting it from the left
margin.  Quotations are commonly displayed.  There
are short quotations
\begin{quote}
   This is a short quotation.  It consists of a
   single paragraph of text.  See how it is formatted.
\end{quote}
and longer ones.
\begin{quotation}
   This is a longer quotation.  It consists of two
   paragraphs of text, neither of which are
   particularly interesting.

   This is the second paragraph of the quotation.  It
   is just as dull as the first paragraph.
\end{quotation}
Another frequently-displayed structure is a list.
The following is an example of an \emph{itemized}
list.
\begin{itemize}
   \item This is the first item of an itemized list.
         Each item in the list is marked with a ``tick''.
         You don't have to worry about what kind of tick
         mark is used.

   \item This is the second item of the list.  It
         contains another list nested inside it.  The inner
         list is an \emph{enumerated} list.
         \begin{enumerate}
            \item This is the first item of an enumerated
                  list that is nested within the itemized list.

            \item This is the second item of the inner list.
                  \LaTeX\ allows you to nest lists deeper than
                  you really should.
         \end{enumerate}
         This is the rest of the second item of the outer
         list.  It is no more interesting than any other
         part of the item.
   \item This is the third item of the list.
\end{itemize}
You can even display poetry.
\begin{verse}
   There is an environment
    for verse \\             % The \\ command separates lines
   Whose features some poets % within a stanza.
   will curse.

                             % One or more blank lines separate stanzas.

   For instead of making\\
   Them do \emph{all} line breaking, \\
   It allows them to put too many words on a line when they'd rather be
   forced to be terse.
\end{verse}

Mathematical formulas may also be displayed.  A
displayed formula
is
one-line long; multiline
formulas require special formatting instructions.
   \[  \ip{\Gamma}{\psi'} = x'' + y^{2} + z_{i}^{n}\]
Don't start a paragraph with a displayed equation,
nor make one a paragraph by itself.

\end{document}               % End of document.

\section{Bench Huge Section 401}
% synthetic bench section 401
% This is a sample LaTeX input file.  (Version of 12 August 2004.)
%
% A '%' character causes TeX to ignore all remaining text on the line,
% and is used for comments like this one.

\documentclass{article}      % Specifies the document class

                             % The preamble begins here.
\title{An Example Document}  % Declares the document's title.
\author{Leslie Lamport}      % Declares the author's name.
\date{January 21, 1994}      % Deleting this command produces today's date.

\newcommand{\ip}[2]{(#1, #2)}
                             % Defines \ip{arg1}{arg2} to mean
                             % (arg1, arg2).

%\newcommand{\ip}[2]{\langle #1 | #2\rangle}
                             % This is an alternative definition of
                             % \ip that is commented out.

\begin{document}             % End of preamble and beginning of text.

\maketitle                   % Produces the title.

This is an example input file.  Comparing it with
the output it generates can show you how to
produce a simple document of your own.

\section{Ordinary Text}      % Produces section heading.  Lower-level
                             % sections are begun with similar
                             % \subsection and \subsubsection commands.

The ends  of words and sentences are marked
  by   spaces. It  doesn't matter how many
spaces    you type; one is as good as 100.  The
end of   a line counts as a space.

One   or more   blank lines denote the  end
of  a paragraph.

Since any number of consecutive spaces are treated
like a single one, the formatting of the input
file makes no difference to
      \LaTeX,                % The \LaTeX command generates the LaTeX logo.
but it makes a difference to you.  When you use
\LaTeX, making your input file as easy to read
as possible will be a great help as you write
your document and when you change it.  This sample
file shows how you can add comments to your own input
file.

Because printing is different from typewriting,
there are a number of things that you have to do
differently when preparing an input file than if
you were just typing the document directly.
Quotation marks like
       ``this''
have to be handled specially, as do quotes within
quotes:
       ``\,`this'            % \, separates the double and single quote.
        is what I just
        wrote, not  `that'\,''.

Dashes come in three sizes: an
       intra-word
dash, a medium dash for number ranges like
       1--2,
and a punctuation
       dash---like
this.

A sentence-ending space should be larger than the
space between words within a sentence.  You
sometimes have to type special commands in
conjunction with punctuation characters to get
this right, as in the following sentence.
       Gnats, gnus, etc.\ all  % `\ ' makes an inter-word space.
       begin with G\@.         % \@ marks end-of-sentence punctuation.
You should check the spaces after periods when
reading your output to make sure you haven't
forgotten any special cases.  Generating an
ellipsis
       \ldots\               % `\ ' is needed after `\ldots' because TeX
                             % ignores spaces after command names like \ldots
                             % made from \ + letters.
                             %
                             % Note how a `%' character causes TeX to ignore
                             % the end of the input line, so these blank lines
                             % do not start a new paragraph.
                             %
with the right spacing around the periods requires
a special command.

\LaTeX\ interprets some common characters as
commands, so you must type special commands to
generate them.  These characters include the
following:
       \$ \& \% \# \{ and \}.

In printing, text is usually emphasized with an
       \emph{italic}
type style.

\begin{em}
   A long segment of text can also be emphasized
   in this way.  Text within such a segment can be
   given \emph{additional} emphasis.
\end{em}

It is sometimes necessary to prevent \LaTeX\ from
breaking a line where it might otherwise do so.
This may be at a space, as between the ``Mr.''\ and
``Jones'' in
       ``Mr.~Jones'',        % ~ produces an unbreakable interword space.
or within a word---especially when the word is a
symbol like
       \mbox{\emph{itemnum}}
that makes little sense when hyphenated across
lines.

Footnotes\footnote{This is an example of a footnote.}
pose no problem.

\LaTeX\ is good at typesetting mathematical formulas
like
       \( x-3y + z = 7 \)
or
       \( a_{1} > x^{2n} + y^{2n} > x' \)
or
       \( \ip{A}{B} = \sum_{i} a_{i} b_{i} \).
The spaces you type in a formula are
ignored.  Remember that a letter like
       $x$                   % $ ... $  and  \( ... \)  are equivalent
is a formula when it denotes a mathematical
symbol, and it should be typed as one.

\section{Displayed Text}

Text is displayed by indenting it from the left
margin.  Quotations are commonly displayed.  There
are short quotations
\begin{quote}
   This is a short quotation.  It consists of a
   single paragraph of text.  See how it is formatted.
\end{quote}
and longer ones.
\begin{quotation}
   This is a longer quotation.  It consists of two
   paragraphs of text, neither of which are
   particularly interesting.

   This is the second paragraph of the quotation.  It
   is just as dull as the first paragraph.
\end{quotation}
Another frequently-displayed structure is a list.
The following is an example of an \emph{itemized}
list.
\begin{itemize}
   \item This is the first item of an itemized list.
         Each item in the list is marked with a ``tick''.
         You don't have to worry about what kind of tick
         mark is used.

   \item This is the second item of the list.  It
         contains another list nested inside it.  The inner
         list is an \emph{enumerated} list.
         \begin{enumerate}
            \item This is the first item of an enumerated
                  list that is nested within the itemized list.

            \item This is the second item of the inner list.
                  \LaTeX\ allows you to nest lists deeper than
                  you really should.
         \end{enumerate}
         This is the rest of the second item of the outer
         list.  It is no more interesting than any other
         part of the item.
   \item This is the third item of the list.
\end{itemize}
You can even display poetry.
\begin{verse}
   There is an environment
    for verse \\             % The \\ command separates lines
   Whose features some poets % within a stanza.
   will curse.

                             % One or more blank lines separate stanzas.

   For instead of making\\
   Them do \emph{all} line breaking, \\
   It allows them to put too many words on a line when they'd rather be
   forced to be terse.
\end{verse}

Mathematical formulas may also be displayed.  A
displayed formula
is
one-line long; multiline
formulas require special formatting instructions.
   \[  \ip{\Gamma}{\psi'} = x'' + y^{2} + z_{i}^{n}\]
Don't start a paragraph with a displayed equation,
nor make one a paragraph by itself.

\end{document}               % End of document.

\section{Bench Huge Section 402}
% synthetic bench section 402
% This is a sample LaTeX input file.  (Version of 12 August 2004.)
%
% A '%' character causes TeX to ignore all remaining text on the line,
% and is used for comments like this one.

\documentclass{article}      % Specifies the document class

                             % The preamble begins here.
\title{An Example Document}  % Declares the document's title.
\author{Leslie Lamport}      % Declares the author's name.
\date{January 21, 1994}      % Deleting this command produces today's date.

\newcommand{\ip}[2]{(#1, #2)}
                             % Defines \ip{arg1}{arg2} to mean
                             % (arg1, arg2).

%\newcommand{\ip}[2]{\langle #1 | #2\rangle}
                             % This is an alternative definition of
                             % \ip that is commented out.

\begin{document}             % End of preamble and beginning of text.

\maketitle                   % Produces the title.

This is an example input file.  Comparing it with
the output it generates can show you how to
produce a simple document of your own.

\section{Ordinary Text}      % Produces section heading.  Lower-level
                             % sections are begun with similar
                             % \subsection and \subsubsection commands.

The ends  of words and sentences are marked
  by   spaces. It  doesn't matter how many
spaces    you type; one is as good as 100.  The
end of   a line counts as a space.

One   or more   blank lines denote the  end
of  a paragraph.

Since any number of consecutive spaces are treated
like a single one, the formatting of the input
file makes no difference to
      \LaTeX,                % The \LaTeX command generates the LaTeX logo.
but it makes a difference to you.  When you use
\LaTeX, making your input file as easy to read
as possible will be a great help as you write
your document and when you change it.  This sample
file shows how you can add comments to your own input
file.

Because printing is different from typewriting,
there are a number of things that you have to do
differently when preparing an input file than if
you were just typing the document directly.
Quotation marks like
       ``this''
have to be handled specially, as do quotes within
quotes:
       ``\,`this'            % \, separates the double and single quote.
        is what I just
        wrote, not  `that'\,''.

Dashes come in three sizes: an
       intra-word
dash, a medium dash for number ranges like
       1--2,
and a punctuation
       dash---like
this.

A sentence-ending space should be larger than the
space between words within a sentence.  You
sometimes have to type special commands in
conjunction with punctuation characters to get
this right, as in the following sentence.
       Gnats, gnus, etc.\ all  % `\ ' makes an inter-word space.
       begin with G\@.         % \@ marks end-of-sentence punctuation.
You should check the spaces after periods when
reading your output to make sure you haven't
forgotten any special cases.  Generating an
ellipsis
       \ldots\               % `\ ' is needed after `\ldots' because TeX
                             % ignores spaces after command names like \ldots
                             % made from \ + letters.
                             %
                             % Note how a `%' character causes TeX to ignore
                             % the end of the input line, so these blank lines
                             % do not start a new paragraph.
                             %
with the right spacing around the periods requires
a special command.

\LaTeX\ interprets some common characters as
commands, so you must type special commands to
generate them.  These characters include the
following:
       \$ \& \% \# \{ and \}.

In printing, text is usually emphasized with an
       \emph{italic}
type style.

\begin{em}
   A long segment of text can also be emphasized
   in this way.  Text within such a segment can be
   given \emph{additional} emphasis.
\end{em}

It is sometimes necessary to prevent \LaTeX\ from
breaking a line where it might otherwise do so.
This may be at a space, as between the ``Mr.''\ and
``Jones'' in
       ``Mr.~Jones'',        % ~ produces an unbreakable interword space.
or within a word---especially when the word is a
symbol like
       \mbox{\emph{itemnum}}
that makes little sense when hyphenated across
lines.

Footnotes\footnote{This is an example of a footnote.}
pose no problem.

\LaTeX\ is good at typesetting mathematical formulas
like
       \( x-3y + z = 7 \)
or
       \( a_{1} > x^{2n} + y^{2n} > x' \)
or
       \( \ip{A}{B} = \sum_{i} a_{i} b_{i} \).
The spaces you type in a formula are
ignored.  Remember that a letter like
       $x$                   % $ ... $  and  \( ... \)  are equivalent
is a formula when it denotes a mathematical
symbol, and it should be typed as one.

\section{Displayed Text}

Text is displayed by indenting it from the left
margin.  Quotations are commonly displayed.  There
are short quotations
\begin{quote}
   This is a short quotation.  It consists of a
   single paragraph of text.  See how it is formatted.
\end{quote}
and longer ones.
\begin{quotation}
   This is a longer quotation.  It consists of two
   paragraphs of text, neither of which are
   particularly interesting.

   This is the second paragraph of the quotation.  It
   is just as dull as the first paragraph.
\end{quotation}
Another frequently-displayed structure is a list.
The following is an example of an \emph{itemized}
list.
\begin{itemize}
   \item This is the first item of an itemized list.
         Each item in the list is marked with a ``tick''.
         You don't have to worry about what kind of tick
         mark is used.

   \item This is the second item of the list.  It
         contains another list nested inside it.  The inner
         list is an \emph{enumerated} list.
         \begin{enumerate}
            \item This is the first item of an enumerated
                  list that is nested within the itemized list.

            \item This is the second item of the inner list.
                  \LaTeX\ allows you to nest lists deeper than
                  you really should.
         \end{enumerate}
         This is the rest of the second item of the outer
         list.  It is no more interesting than any other
         part of the item.
   \item This is the third item of the list.
\end{itemize}
You can even display poetry.
\begin{verse}
   There is an environment
    for verse \\             % The \\ command separates lines
   Whose features some poets % within a stanza.
   will curse.

                             % One or more blank lines separate stanzas.

   For instead of making\\
   Them do \emph{all} line breaking, \\
   It allows them to put too many words on a line when they'd rather be
   forced to be terse.
\end{verse}

Mathematical formulas may also be displayed.  A
displayed formula
is
one-line long; multiline
formulas require special formatting instructions.
   \[  \ip{\Gamma}{\psi'} = x'' + y^{2} + z_{i}^{n}\]
Don't start a paragraph with a displayed equation,
nor make one a paragraph by itself.

\end{document}               % End of document.

\section{Bench Huge Section 403}
% synthetic bench section 403
% This is a sample LaTeX input file.  (Version of 12 August 2004.)
%
% A '%' character causes TeX to ignore all remaining text on the line,
% and is used for comments like this one.

\documentclass{article}      % Specifies the document class

                             % The preamble begins here.
\title{An Example Document}  % Declares the document's title.
\author{Leslie Lamport}      % Declares the author's name.
\date{January 21, 1994}      % Deleting this command produces today's date.

\newcommand{\ip}[2]{(#1, #2)}
                             % Defines \ip{arg1}{arg2} to mean
                             % (arg1, arg2).

%\newcommand{\ip}[2]{\langle #1 | #2\rangle}
                             % This is an alternative definition of
                             % \ip that is commented out.

\begin{document}             % End of preamble and beginning of text.

\maketitle                   % Produces the title.

This is an example input file.  Comparing it with
the output it generates can show you how to
produce a simple document of your own.

\section{Ordinary Text}      % Produces section heading.  Lower-level
                             % sections are begun with similar
                             % \subsection and \subsubsection commands.

The ends  of words and sentences are marked
  by   spaces. It  doesn't matter how many
spaces    you type; one is as good as 100.  The
end of   a line counts as a space.

One   or more   blank lines denote the  end
of  a paragraph.

Since any number of consecutive spaces are treated
like a single one, the formatting of the input
file makes no difference to
      \LaTeX,                % The \LaTeX command generates the LaTeX logo.
but it makes a difference to you.  When you use
\LaTeX, making your input file as easy to read
as possible will be a great help as you write
your document and when you change it.  This sample
file shows how you can add comments to your own input
file.

Because printing is different from typewriting,
there are a number of things that you have to do
differently when preparing an input file than if
you were just typing the document directly.
Quotation marks like
       ``this''
have to be handled specially, as do quotes within
quotes:
       ``\,`this'            % \, separates the double and single quote.
        is what I just
        wrote, not  `that'\,''.

Dashes come in three sizes: an
       intra-word
dash, a medium dash for number ranges like
       1--2,
and a punctuation
       dash---like
this.

A sentence-ending space should be larger than the
space between words within a sentence.  You
sometimes have to type special commands in
conjunction with punctuation characters to get
this right, as in the following sentence.
       Gnats, gnus, etc.\ all  % `\ ' makes an inter-word space.
       begin with G\@.         % \@ marks end-of-sentence punctuation.
You should check the spaces after periods when
reading your output to make sure you haven't
forgotten any special cases.  Generating an
ellipsis
       \ldots\               % `\ ' is needed after `\ldots' because TeX
                             % ignores spaces after command names like \ldots
                             % made from \ + letters.
                             %
                             % Note how a `%' character causes TeX to ignore
                             % the end of the input line, so these blank lines
                             % do not start a new paragraph.
                             %
with the right spacing around the periods requires
a special command.

\LaTeX\ interprets some common characters as
commands, so you must type special commands to
generate them.  These characters include the
following:
       \$ \& \% \# \{ and \}.

In printing, text is usually emphasized with an
       \emph{italic}
type style.

\begin{em}
   A long segment of text can also be emphasized
   in this way.  Text within such a segment can be
   given \emph{additional} emphasis.
\end{em}

It is sometimes necessary to prevent \LaTeX\ from
breaking a line where it might otherwise do so.
This may be at a space, as between the ``Mr.''\ and
``Jones'' in
       ``Mr.~Jones'',        % ~ produces an unbreakable interword space.
or within a word---especially when the word is a
symbol like
       \mbox{\emph{itemnum}}
that makes little sense when hyphenated across
lines.

Footnotes\footnote{This is an example of a footnote.}
pose no problem.

\LaTeX\ is good at typesetting mathematical formulas
like
       \( x-3y + z = 7 \)
or
       \( a_{1} > x^{2n} + y^{2n} > x' \)
or
       \( \ip{A}{B} = \sum_{i} a_{i} b_{i} \).
The spaces you type in a formula are
ignored.  Remember that a letter like
       $x$                   % $ ... $  and  \( ... \)  are equivalent
is a formula when it denotes a mathematical
symbol, and it should be typed as one.

\section{Displayed Text}

Text is displayed by indenting it from the left
margin.  Quotations are commonly displayed.  There
are short quotations
\begin{quote}
   This is a short quotation.  It consists of a
   single paragraph of text.  See how it is formatted.
\end{quote}
and longer ones.
\begin{quotation}
   This is a longer quotation.  It consists of two
   paragraphs of text, neither of which are
   particularly interesting.

   This is the second paragraph of the quotation.  It
   is just as dull as the first paragraph.
\end{quotation}
Another frequently-displayed structure is a list.
The following is an example of an \emph{itemized}
list.
\begin{itemize}
   \item This is the first item of an itemized list.
         Each item in the list is marked with a ``tick''.
         You don't have to worry about what kind of tick
         mark is used.

   \item This is the second item of the list.  It
         contains another list nested inside it.  The inner
         list is an \emph{enumerated} list.
         \begin{enumerate}
            \item This is the first item of an enumerated
                  list that is nested within the itemized list.

            \item This is the second item of the inner list.
                  \LaTeX\ allows you to nest lists deeper than
                  you really should.
         \end{enumerate}
         This is the rest of the second item of the outer
         list.  It is no more interesting than any other
         part of the item.
   \item This is the third item of the list.
\end{itemize}
You can even display poetry.
\begin{verse}
   There is an environment
    for verse \\             % The \\ command separates lines
   Whose features some poets % within a stanza.
   will curse.

                             % One or more blank lines separate stanzas.

   For instead of making\\
   Them do \emph{all} line breaking, \\
   It allows them to put too many words on a line when they'd rather be
   forced to be terse.
\end{verse}

Mathematical formulas may also be displayed.  A
displayed formula
is
one-line long; multiline
formulas require special formatting instructions.
   \[  \ip{\Gamma}{\psi'} = x'' + y^{2} + z_{i}^{n}\]
Don't start a paragraph with a displayed equation,
nor make one a paragraph by itself.

\end{document}               % End of document.

\section{Bench Huge Section 404}
% synthetic bench section 404
% This is a sample LaTeX input file.  (Version of 12 August 2004.)
%
% A '%' character causes TeX to ignore all remaining text on the line,
% and is used for comments like this one.

\documentclass{article}      % Specifies the document class

                             % The preamble begins here.
\title{An Example Document}  % Declares the document's title.
\author{Leslie Lamport}      % Declares the author's name.
\date{January 21, 1994}      % Deleting this command produces today's date.

\newcommand{\ip}[2]{(#1, #2)}
                             % Defines \ip{arg1}{arg2} to mean
                             % (arg1, arg2).

%\newcommand{\ip}[2]{\langle #1 | #2\rangle}
                             % This is an alternative definition of
                             % \ip that is commented out.

\begin{document}             % End of preamble and beginning of text.

\maketitle                   % Produces the title.

This is an example input file.  Comparing it with
the output it generates can show you how to
produce a simple document of your own.

\section{Ordinary Text}      % Produces section heading.  Lower-level
                             % sections are begun with similar
                             % \subsection and \subsubsection commands.

The ends  of words and sentences are marked
  by   spaces. It  doesn't matter how many
spaces    you type; one is as good as 100.  The
end of   a line counts as a space.

One   or more   blank lines denote the  end
of  a paragraph.

Since any number of consecutive spaces are treated
like a single one, the formatting of the input
file makes no difference to
      \LaTeX,                % The \LaTeX command generates the LaTeX logo.
but it makes a difference to you.  When you use
\LaTeX, making your input file as easy to read
as possible will be a great help as you write
your document and when you change it.  This sample
file shows how you can add comments to your own input
file.

Because printing is different from typewriting,
there are a number of things that you have to do
differently when preparing an input file than if
you were just typing the document directly.
Quotation marks like
       ``this''
have to be handled specially, as do quotes within
quotes:
       ``\,`this'            % \, separates the double and single quote.
        is what I just
        wrote, not  `that'\,''.

Dashes come in three sizes: an
       intra-word
dash, a medium dash for number ranges like
       1--2,
and a punctuation
       dash---like
this.

A sentence-ending space should be larger than the
space between words within a sentence.  You
sometimes have to type special commands in
conjunction with punctuation characters to get
this right, as in the following sentence.
       Gnats, gnus, etc.\ all  % `\ ' makes an inter-word space.
       begin with G\@.         % \@ marks end-of-sentence punctuation.
You should check the spaces after periods when
reading your output to make sure you haven't
forgotten any special cases.  Generating an
ellipsis
       \ldots\               % `\ ' is needed after `\ldots' because TeX
                             % ignores spaces after command names like \ldots
                             % made from \ + letters.
                             %
                             % Note how a `%' character causes TeX to ignore
                             % the end of the input line, so these blank lines
                             % do not start a new paragraph.
                             %
with the right spacing around the periods requires
a special command.

\LaTeX\ interprets some common characters as
commands, so you must type special commands to
generate them.  These characters include the
following:
       \$ \& \% \# \{ and \}.

In printing, text is usually emphasized with an
       \emph{italic}
type style.

\begin{em}
   A long segment of text can also be emphasized
   in this way.  Text within such a segment can be
   given \emph{additional} emphasis.
\end{em}

It is sometimes necessary to prevent \LaTeX\ from
breaking a line where it might otherwise do so.
This may be at a space, as between the ``Mr.''\ and
``Jones'' in
       ``Mr.~Jones'',        % ~ produces an unbreakable interword space.
or within a word---especially when the word is a
symbol like
       \mbox{\emph{itemnum}}
that makes little sense when hyphenated across
lines.

Footnotes\footnote{This is an example of a footnote.}
pose no problem.

\LaTeX\ is good at typesetting mathematical formulas
like
       \( x-3y + z = 7 \)
or
       \( a_{1} > x^{2n} + y^{2n} > x' \)
or
       \( \ip{A}{B} = \sum_{i} a_{i} b_{i} \).
The spaces you type in a formula are
ignored.  Remember that a letter like
       $x$                   % $ ... $  and  \( ... \)  are equivalent
is a formula when it denotes a mathematical
symbol, and it should be typed as one.

\section{Displayed Text}

Text is displayed by indenting it from the left
margin.  Quotations are commonly displayed.  There
are short quotations
\begin{quote}
   This is a short quotation.  It consists of a
   single paragraph of text.  See how it is formatted.
\end{quote}
and longer ones.
\begin{quotation}
   This is a longer quotation.  It consists of two
   paragraphs of text, neither of which are
   particularly interesting.

   This is the second paragraph of the quotation.  It
   is just as dull as the first paragraph.
\end{quotation}
Another frequently-displayed structure is a list.
The following is an example of an \emph{itemized}
list.
\begin{itemize}
   \item This is the first item of an itemized list.
         Each item in the list is marked with a ``tick''.
         You don't have to worry about what kind of tick
         mark is used.

   \item This is the second item of the list.  It
         contains another list nested inside it.  The inner
         list is an \emph{enumerated} list.
         \begin{enumerate}
            \item This is the first item of an enumerated
                  list that is nested within the itemized list.

            \item This is the second item of the inner list.
                  \LaTeX\ allows you to nest lists deeper than
                  you really should.
         \end{enumerate}
         This is the rest of the second item of the outer
         list.  It is no more interesting than any other
         part of the item.
   \item This is the third item of the list.
\end{itemize}
You can even display poetry.
\begin{verse}
   There is an environment
    for verse \\             % The \\ command separates lines
   Whose features some poets % within a stanza.
   will curse.

                             % One or more blank lines separate stanzas.

   For instead of making\\
   Them do \emph{all} line breaking, \\
   It allows them to put too many words on a line when they'd rather be
   forced to be terse.
\end{verse}

Mathematical formulas may also be displayed.  A
displayed formula
is
one-line long; multiline
formulas require special formatting instructions.
   \[  \ip{\Gamma}{\psi'} = x'' + y^{2} + z_{i}^{n}\]
Don't start a paragraph with a displayed equation,
nor make one a paragraph by itself.

\end{document}               % End of document.

\section{Bench Huge Section 405}
% synthetic bench section 405
% This is a sample LaTeX input file.  (Version of 12 August 2004.)
%
% A '%' character causes TeX to ignore all remaining text on the line,
% and is used for comments like this one.

\documentclass{article}      % Specifies the document class

                             % The preamble begins here.
\title{An Example Document}  % Declares the document's title.
\author{Leslie Lamport}      % Declares the author's name.
\date{January 21, 1994}      % Deleting this command produces today's date.

\newcommand{\ip}[2]{(#1, #2)}
                             % Defines \ip{arg1}{arg2} to mean
                             % (arg1, arg2).

%\newcommand{\ip}[2]{\langle #1 | #2\rangle}
                             % This is an alternative definition of
                             % \ip that is commented out.

\begin{document}             % End of preamble and beginning of text.

\maketitle                   % Produces the title.

This is an example input file.  Comparing it with
the output it generates can show you how to
produce a simple document of your own.

\section{Ordinary Text}      % Produces section heading.  Lower-level
                             % sections are begun with similar
                             % \subsection and \subsubsection commands.

The ends  of words and sentences are marked
  by   spaces. It  doesn't matter how many
spaces    you type; one is as good as 100.  The
end of   a line counts as a space.

One   or more   blank lines denote the  end
of  a paragraph.

Since any number of consecutive spaces are treated
like a single one, the formatting of the input
file makes no difference to
      \LaTeX,                % The \LaTeX command generates the LaTeX logo.
but it makes a difference to you.  When you use
\LaTeX, making your input file as easy to read
as possible will be a great help as you write
your document and when you change it.  This sample
file shows how you can add comments to your own input
file.

Because printing is different from typewriting,
there are a number of things that you have to do
differently when preparing an input file than if
you were just typing the document directly.
Quotation marks like
       ``this''
have to be handled specially, as do quotes within
quotes:
       ``\,`this'            % \, separates the double and single quote.
        is what I just
        wrote, not  `that'\,''.

Dashes come in three sizes: an
       intra-word
dash, a medium dash for number ranges like
       1--2,
and a punctuation
       dash---like
this.

A sentence-ending space should be larger than the
space between words within a sentence.  You
sometimes have to type special commands in
conjunction with punctuation characters to get
this right, as in the following sentence.
       Gnats, gnus, etc.\ all  % `\ ' makes an inter-word space.
       begin with G\@.         % \@ marks end-of-sentence punctuation.
You should check the spaces after periods when
reading your output to make sure you haven't
forgotten any special cases.  Generating an
ellipsis
       \ldots\               % `\ ' is needed after `\ldots' because TeX
                             % ignores spaces after command names like \ldots
                             % made from \ + letters.
                             %
                             % Note how a `%' character causes TeX to ignore
                             % the end of the input line, so these blank lines
                             % do not start a new paragraph.
                             %
with the right spacing around the periods requires
a special command.

\LaTeX\ interprets some common characters as
commands, so you must type special commands to
generate them.  These characters include the
following:
       \$ \& \% \# \{ and \}.

In printing, text is usually emphasized with an
       \emph{italic}
type style.

\begin{em}
   A long segment of text can also be emphasized
   in this way.  Text within such a segment can be
   given \emph{additional} emphasis.
\end{em}

It is sometimes necessary to prevent \LaTeX\ from
breaking a line where it might otherwise do so.
This may be at a space, as between the ``Mr.''\ and
``Jones'' in
       ``Mr.~Jones'',        % ~ produces an unbreakable interword space.
or within a word---especially when the word is a
symbol like
       \mbox{\emph{itemnum}}
that makes little sense when hyphenated across
lines.

Footnotes\footnote{This is an example of a footnote.}
pose no problem.

\LaTeX\ is good at typesetting mathematical formulas
like
       \( x-3y + z = 7 \)
or
       \( a_{1} > x^{2n} + y^{2n} > x' \)
or
       \( \ip{A}{B} = \sum_{i} a_{i} b_{i} \).
The spaces you type in a formula are
ignored.  Remember that a letter like
       $x$                   % $ ... $  and  \( ... \)  are equivalent
is a formula when it denotes a mathematical
symbol, and it should be typed as one.

\section{Displayed Text}

Text is displayed by indenting it from the left
margin.  Quotations are commonly displayed.  There
are short quotations
\begin{quote}
   This is a short quotation.  It consists of a
   single paragraph of text.  See how it is formatted.
\end{quote}
and longer ones.
\begin{quotation}
   This is a longer quotation.  It consists of two
   paragraphs of text, neither of which are
   particularly interesting.

   This is the second paragraph of the quotation.  It
   is just as dull as the first paragraph.
\end{quotation}
Another frequently-displayed structure is a list.
The following is an example of an \emph{itemized}
list.
\begin{itemize}
   \item This is the first item of an itemized list.
         Each item in the list is marked with a ``tick''.
         You don't have to worry about what kind of tick
         mark is used.

   \item This is the second item of the list.  It
         contains another list nested inside it.  The inner
         list is an \emph{enumerated} list.
         \begin{enumerate}
            \item This is the first item of an enumerated
                  list that is nested within the itemized list.

            \item This is the second item of the inner list.
                  \LaTeX\ allows you to nest lists deeper than
                  you really should.
         \end{enumerate}
         This is the rest of the second item of the outer
         list.  It is no more interesting than any other
         part of the item.
   \item This is the third item of the list.
\end{itemize}
You can even display poetry.
\begin{verse}
   There is an environment
    for verse \\             % The \\ command separates lines
   Whose features some poets % within a stanza.
   will curse.

                             % One or more blank lines separate stanzas.

   For instead of making\\
   Them do \emph{all} line breaking, \\
   It allows them to put too many words on a line when they'd rather be
   forced to be terse.
\end{verse}

Mathematical formulas may also be displayed.  A
displayed formula
is
one-line long; multiline
formulas require special formatting instructions.
   \[  \ip{\Gamma}{\psi'} = x'' + y^{2} + z_{i}^{n}\]
Don't start a paragraph with a displayed equation,
nor make one a paragraph by itself.

\end{document}               % End of document.

\section{Bench Huge Section 406}
% synthetic bench section 406
% This is a sample LaTeX input file.  (Version of 12 August 2004.)
%
% A '%' character causes TeX to ignore all remaining text on the line,
% and is used for comments like this one.

\documentclass{article}      % Specifies the document class

                             % The preamble begins here.
\title{An Example Document}  % Declares the document's title.
\author{Leslie Lamport}      % Declares the author's name.
\date{January 21, 1994}      % Deleting this command produces today's date.

\newcommand{\ip}[2]{(#1, #2)}
                             % Defines \ip{arg1}{arg2} to mean
                             % (arg1, arg2).

%\newcommand{\ip}[2]{\langle #1 | #2\rangle}
                             % This is an alternative definition of
                             % \ip that is commented out.

\begin{document}             % End of preamble and beginning of text.

\maketitle                   % Produces the title.

This is an example input file.  Comparing it with
the output it generates can show you how to
produce a simple document of your own.

\section{Ordinary Text}      % Produces section heading.  Lower-level
                             % sections are begun with similar
                             % \subsection and \subsubsection commands.

The ends  of words and sentences are marked
  by   spaces. It  doesn't matter how many
spaces    you type; one is as good as 100.  The
end of   a line counts as a space.

One   or more   blank lines denote the  end
of  a paragraph.

Since any number of consecutive spaces are treated
like a single one, the formatting of the input
file makes no difference to
      \LaTeX,                % The \LaTeX command generates the LaTeX logo.
but it makes a difference to you.  When you use
\LaTeX, making your input file as easy to read
as possible will be a great help as you write
your document and when you change it.  This sample
file shows how you can add comments to your own input
file.

Because printing is different from typewriting,
there are a number of things that you have to do
differently when preparing an input file than if
you were just typing the document directly.
Quotation marks like
       ``this''
have to be handled specially, as do quotes within
quotes:
       ``\,`this'            % \, separates the double and single quote.
        is what I just
        wrote, not  `that'\,''.

Dashes come in three sizes: an
       intra-word
dash, a medium dash for number ranges like
       1--2,
and a punctuation
       dash---like
this.

A sentence-ending space should be larger than the
space between words within a sentence.  You
sometimes have to type special commands in
conjunction with punctuation characters to get
this right, as in the following sentence.
       Gnats, gnus, etc.\ all  % `\ ' makes an inter-word space.
       begin with G\@.         % \@ marks end-of-sentence punctuation.
You should check the spaces after periods when
reading your output to make sure you haven't
forgotten any special cases.  Generating an
ellipsis
       \ldots\               % `\ ' is needed after `\ldots' because TeX
                             % ignores spaces after command names like \ldots
                             % made from \ + letters.
                             %
                             % Note how a `%' character causes TeX to ignore
                             % the end of the input line, so these blank lines
                             % do not start a new paragraph.
                             %
with the right spacing around the periods requires
a special command.

\LaTeX\ interprets some common characters as
commands, so you must type special commands to
generate them.  These characters include the
following:
       \$ \& \% \# \{ and \}.

In printing, text is usually emphasized with an
       \emph{italic}
type style.

\begin{em}
   A long segment of text can also be emphasized
   in this way.  Text within such a segment can be
   given \emph{additional} emphasis.
\end{em}

It is sometimes necessary to prevent \LaTeX\ from
breaking a line where it might otherwise do so.
This may be at a space, as between the ``Mr.''\ and
``Jones'' in
       ``Mr.~Jones'',        % ~ produces an unbreakable interword space.
or within a word---especially when the word is a
symbol like
       \mbox{\emph{itemnum}}
that makes little sense when hyphenated across
lines.

Footnotes\footnote{This is an example of a footnote.}
pose no problem.

\LaTeX\ is good at typesetting mathematical formulas
like
       \( x-3y + z = 7 \)
or
       \( a_{1} > x^{2n} + y^{2n} > x' \)
or
       \( \ip{A}{B} = \sum_{i} a_{i} b_{i} \).
The spaces you type in a formula are
ignored.  Remember that a letter like
       $x$                   % $ ... $  and  \( ... \)  are equivalent
is a formula when it denotes a mathematical
symbol, and it should be typed as one.

\section{Displayed Text}

Text is displayed by indenting it from the left
margin.  Quotations are commonly displayed.  There
are short quotations
\begin{quote}
   This is a short quotation.  It consists of a
   single paragraph of text.  See how it is formatted.
\end{quote}
and longer ones.
\begin{quotation}
   This is a longer quotation.  It consists of two
   paragraphs of text, neither of which are
   particularly interesting.

   This is the second paragraph of the quotation.  It
   is just as dull as the first paragraph.
\end{quotation}
Another frequently-displayed structure is a list.
The following is an example of an \emph{itemized}
list.
\begin{itemize}
   \item This is the first item of an itemized list.
         Each item in the list is marked with a ``tick''.
         You don't have to worry about what kind of tick
         mark is used.

   \item This is the second item of the list.  It
         contains another list nested inside it.  The inner
         list is an \emph{enumerated} list.
         \begin{enumerate}
            \item This is the first item of an enumerated
                  list that is nested within the itemized list.

            \item This is the second item of the inner list.
                  \LaTeX\ allows you to nest lists deeper than
                  you really should.
         \end{enumerate}
         This is the rest of the second item of the outer
         list.  It is no more interesting than any other
         part of the item.
   \item This is the third item of the list.
\end{itemize}
You can even display poetry.
\begin{verse}
   There is an environment
    for verse \\             % The \\ command separates lines
   Whose features some poets % within a stanza.
   will curse.

                             % One or more blank lines separate stanzas.

   For instead of making\\
   Them do \emph{all} line breaking, \\
   It allows them to put too many words on a line when they'd rather be
   forced to be terse.
\end{verse}

Mathematical formulas may also be displayed.  A
displayed formula
is
one-line long; multiline
formulas require special formatting instructions.
   \[  \ip{\Gamma}{\psi'} = x'' + y^{2} + z_{i}^{n}\]
Don't start a paragraph with a displayed equation,
nor make one a paragraph by itself.

\end{document}               % End of document.

\section{Bench Huge Section 407}
% synthetic bench section 407
% This is a sample LaTeX input file.  (Version of 12 August 2004.)
%
% A '%' character causes TeX to ignore all remaining text on the line,
% and is used for comments like this one.

\documentclass{article}      % Specifies the document class

                             % The preamble begins here.
\title{An Example Document}  % Declares the document's title.
\author{Leslie Lamport}      % Declares the author's name.
\date{January 21, 1994}      % Deleting this command produces today's date.

\newcommand{\ip}[2]{(#1, #2)}
                             % Defines \ip{arg1}{arg2} to mean
                             % (arg1, arg2).

%\newcommand{\ip}[2]{\langle #1 | #2\rangle}
                             % This is an alternative definition of
                             % \ip that is commented out.

\begin{document}             % End of preamble and beginning of text.

\maketitle                   % Produces the title.

This is an example input file.  Comparing it with
the output it generates can show you how to
produce a simple document of your own.

\section{Ordinary Text}      % Produces section heading.  Lower-level
                             % sections are begun with similar
                             % \subsection and \subsubsection commands.

The ends  of words and sentences are marked
  by   spaces. It  doesn't matter how many
spaces    you type; one is as good as 100.  The
end of   a line counts as a space.

One   or more   blank lines denote the  end
of  a paragraph.

Since any number of consecutive spaces are treated
like a single one, the formatting of the input
file makes no difference to
      \LaTeX,                % The \LaTeX command generates the LaTeX logo.
but it makes a difference to you.  When you use
\LaTeX, making your input file as easy to read
as possible will be a great help as you write
your document and when you change it.  This sample
file shows how you can add comments to your own input
file.

Because printing is different from typewriting,
there are a number of things that you have to do
differently when preparing an input file than if
you were just typing the document directly.
Quotation marks like
       ``this''
have to be handled specially, as do quotes within
quotes:
       ``\,`this'            % \, separates the double and single quote.
        is what I just
        wrote, not  `that'\,''.

Dashes come in three sizes: an
       intra-word
dash, a medium dash for number ranges like
       1--2,
and a punctuation
       dash---like
this.

A sentence-ending space should be larger than the
space between words within a sentence.  You
sometimes have to type special commands in
conjunction with punctuation characters to get
this right, as in the following sentence.
       Gnats, gnus, etc.\ all  % `\ ' makes an inter-word space.
       begin with G\@.         % \@ marks end-of-sentence punctuation.
You should check the spaces after periods when
reading your output to make sure you haven't
forgotten any special cases.  Generating an
ellipsis
       \ldots\               % `\ ' is needed after `\ldots' because TeX
                             % ignores spaces after command names like \ldots
                             % made from \ + letters.
                             %
                             % Note how a `%' character causes TeX to ignore
                             % the end of the input line, so these blank lines
                             % do not start a new paragraph.
                             %
with the right spacing around the periods requires
a special command.

\LaTeX\ interprets some common characters as
commands, so you must type special commands to
generate them.  These characters include the
following:
       \$ \& \% \# \{ and \}.

In printing, text is usually emphasized with an
       \emph{italic}
type style.

\begin{em}
   A long segment of text can also be emphasized
   in this way.  Text within such a segment can be
   given \emph{additional} emphasis.
\end{em}

It is sometimes necessary to prevent \LaTeX\ from
breaking a line where it might otherwise do so.
This may be at a space, as between the ``Mr.''\ and
``Jones'' in
       ``Mr.~Jones'',        % ~ produces an unbreakable interword space.
or within a word---especially when the word is a
symbol like
       \mbox{\emph{itemnum}}
that makes little sense when hyphenated across
lines.

Footnotes\footnote{This is an example of a footnote.}
pose no problem.

\LaTeX\ is good at typesetting mathematical formulas
like
       \( x-3y + z = 7 \)
or
       \( a_{1} > x^{2n} + y^{2n} > x' \)
or
       \( \ip{A}{B} = \sum_{i} a_{i} b_{i} \).
The spaces you type in a formula are
ignored.  Remember that a letter like
       $x$                   % $ ... $  and  \( ... \)  are equivalent
is a formula when it denotes a mathematical
symbol, and it should be typed as one.

\section{Displayed Text}

Text is displayed by indenting it from the left
margin.  Quotations are commonly displayed.  There
are short quotations
\begin{quote}
   This is a short quotation.  It consists of a
   single paragraph of text.  See how it is formatted.
\end{quote}
and longer ones.
\begin{quotation}
   This is a longer quotation.  It consists of two
   paragraphs of text, neither of which are
   particularly interesting.

   This is the second paragraph of the quotation.  It
   is just as dull as the first paragraph.
\end{quotation}
Another frequently-displayed structure is a list.
The following is an example of an \emph{itemized}
list.
\begin{itemize}
   \item This is the first item of an itemized list.
         Each item in the list is marked with a ``tick''.
         You don't have to worry about what kind of tick
         mark is used.

   \item This is the second item of the list.  It
         contains another list nested inside it.  The inner
         list is an \emph{enumerated} list.
         \begin{enumerate}
            \item This is the first item of an enumerated
                  list that is nested within the itemized list.

            \item This is the second item of the inner list.
                  \LaTeX\ allows you to nest lists deeper than
                  you really should.
         \end{enumerate}
         This is the rest of the second item of the outer
         list.  It is no more interesting than any other
         part of the item.
   \item This is the third item of the list.
\end{itemize}
You can even display poetry.
\begin{verse}
   There is an environment
    for verse \\             % The \\ command separates lines
   Whose features some poets % within a stanza.
   will curse.

                             % One or more blank lines separate stanzas.

   For instead of making\\
   Them do \emph{all} line breaking, \\
   It allows them to put too many words on a line when they'd rather be
   forced to be terse.
\end{verse}

Mathematical formulas may also be displayed.  A
displayed formula
is
one-line long; multiline
formulas require special formatting instructions.
   \[  \ip{\Gamma}{\psi'} = x'' + y^{2} + z_{i}^{n}\]
Don't start a paragraph with a displayed equation,
nor make one a paragraph by itself.

\end{document}               % End of document.

\section{Bench Huge Section 408}
% synthetic bench section 408
% This is a sample LaTeX input file.  (Version of 12 August 2004.)
%
% A '%' character causes TeX to ignore all remaining text on the line,
% and is used for comments like this one.

\documentclass{article}      % Specifies the document class

                             % The preamble begins here.
\title{An Example Document}  % Declares the document's title.
\author{Leslie Lamport}      % Declares the author's name.
\date{January 21, 1994}      % Deleting this command produces today's date.

\newcommand{\ip}[2]{(#1, #2)}
                             % Defines \ip{arg1}{arg2} to mean
                             % (arg1, arg2).

%\newcommand{\ip}[2]{\langle #1 | #2\rangle}
                             % This is an alternative definition of
                             % \ip that is commented out.

\begin{document}             % End of preamble and beginning of text.

\maketitle                   % Produces the title.

This is an example input file.  Comparing it with
the output it generates can show you how to
produce a simple document of your own.

\section{Ordinary Text}      % Produces section heading.  Lower-level
                             % sections are begun with similar
                             % \subsection and \subsubsection commands.

The ends  of words and sentences are marked
  by   spaces. It  doesn't matter how many
spaces    you type; one is as good as 100.  The
end of   a line counts as a space.

One   or more   blank lines denote the  end
of  a paragraph.

Since any number of consecutive spaces are treated
like a single one, the formatting of the input
file makes no difference to
      \LaTeX,                % The \LaTeX command generates the LaTeX logo.
but it makes a difference to you.  When you use
\LaTeX, making your input file as easy to read
as possible will be a great help as you write
your document and when you change it.  This sample
file shows how you can add comments to your own input
file.

Because printing is different from typewriting,
there are a number of things that you have to do
differently when preparing an input file than if
you were just typing the document directly.
Quotation marks like
       ``this''
have to be handled specially, as do quotes within
quotes:
       ``\,`this'            % \, separates the double and single quote.
        is what I just
        wrote, not  `that'\,''.

Dashes come in three sizes: an
       intra-word
dash, a medium dash for number ranges like
       1--2,
and a punctuation
       dash---like
this.

A sentence-ending space should be larger than the
space between words within a sentence.  You
sometimes have to type special commands in
conjunction with punctuation characters to get
this right, as in the following sentence.
       Gnats, gnus, etc.\ all  % `\ ' makes an inter-word space.
       begin with G\@.         % \@ marks end-of-sentence punctuation.
You should check the spaces after periods when
reading your output to make sure you haven't
forgotten any special cases.  Generating an
ellipsis
       \ldots\               % `\ ' is needed after `\ldots' because TeX
                             % ignores spaces after command names like \ldots
                             % made from \ + letters.
                             %
                             % Note how a `%' character causes TeX to ignore
                             % the end of the input line, so these blank lines
                             % do not start a new paragraph.
                             %
with the right spacing around the periods requires
a special command.

\LaTeX\ interprets some common characters as
commands, so you must type special commands to
generate them.  These characters include the
following:
       \$ \& \% \# \{ and \}.

In printing, text is usually emphasized with an
       \emph{italic}
type style.

\begin{em}
   A long segment of text can also be emphasized
   in this way.  Text within such a segment can be
   given \emph{additional} emphasis.
\end{em}

It is sometimes necessary to prevent \LaTeX\ from
breaking a line where it might otherwise do so.
This may be at a space, as between the ``Mr.''\ and
``Jones'' in
       ``Mr.~Jones'',        % ~ produces an unbreakable interword space.
or within a word---especially when the word is a
symbol like
       \mbox{\emph{itemnum}}
that makes little sense when hyphenated across
lines.

Footnotes\footnote{This is an example of a footnote.}
pose no problem.

\LaTeX\ is good at typesetting mathematical formulas
like
       \( x-3y + z = 7 \)
or
       \( a_{1} > x^{2n} + y^{2n} > x' \)
or
       \( \ip{A}{B} = \sum_{i} a_{i} b_{i} \).
The spaces you type in a formula are
ignored.  Remember that a letter like
       $x$                   % $ ... $  and  \( ... \)  are equivalent
is a formula when it denotes a mathematical
symbol, and it should be typed as one.

\section{Displayed Text}

Text is displayed by indenting it from the left
margin.  Quotations are commonly displayed.  There
are short quotations
\begin{quote}
   This is a short quotation.  It consists of a
   single paragraph of text.  See how it is formatted.
\end{quote}
and longer ones.
\begin{quotation}
   This is a longer quotation.  It consists of two
   paragraphs of text, neither of which are
   particularly interesting.

   This is the second paragraph of the quotation.  It
   is just as dull as the first paragraph.
\end{quotation}
Another frequently-displayed structure is a list.
The following is an example of an \emph{itemized}
list.
\begin{itemize}
   \item This is the first item of an itemized list.
         Each item in the list is marked with a ``tick''.
         You don't have to worry about what kind of tick
         mark is used.

   \item This is the second item of the list.  It
         contains another list nested inside it.  The inner
         list is an \emph{enumerated} list.
         \begin{enumerate}
            \item This is the first item of an enumerated
                  list that is nested within the itemized list.

            \item This is the second item of the inner list.
                  \LaTeX\ allows you to nest lists deeper than
                  you really should.
         \end{enumerate}
         This is the rest of the second item of the outer
         list.  It is no more interesting than any other
         part of the item.
   \item This is the third item of the list.
\end{itemize}
You can even display poetry.
\begin{verse}
   There is an environment
    for verse \\             % The \\ command separates lines
   Whose features some poets % within a stanza.
   will curse.

                             % One or more blank lines separate stanzas.

   For instead of making\\
   Them do \emph{all} line breaking, \\
   It allows them to put too many words on a line when they'd rather be
   forced to be terse.
\end{verse}

Mathematical formulas may also be displayed.  A
displayed formula
is
one-line long; multiline
formulas require special formatting instructions.
   \[  \ip{\Gamma}{\psi'} = x'' + y^{2} + z_{i}^{n}\]
Don't start a paragraph with a displayed equation,
nor make one a paragraph by itself.

\end{document}               % End of document.

\section{Bench Huge Section 409}
% synthetic bench section 409
% This is a sample LaTeX input file.  (Version of 12 August 2004.)
%
% A '%' character causes TeX to ignore all remaining text on the line,
% and is used for comments like this one.

\documentclass{article}      % Specifies the document class

                             % The preamble begins here.
\title{An Example Document}  % Declares the document's title.
\author{Leslie Lamport}      % Declares the author's name.
\date{January 21, 1994}      % Deleting this command produces today's date.

\newcommand{\ip}[2]{(#1, #2)}
                             % Defines \ip{arg1}{arg2} to mean
                             % (arg1, arg2).

%\newcommand{\ip}[2]{\langle #1 | #2\rangle}
                             % This is an alternative definition of
                             % \ip that is commented out.

\begin{document}             % End of preamble and beginning of text.

\maketitle                   % Produces the title.

This is an example input file.  Comparing it with
the output it generates can show you how to
produce a simple document of your own.

\section{Ordinary Text}      % Produces section heading.  Lower-level
                             % sections are begun with similar
                             % \subsection and \subsubsection commands.

The ends  of words and sentences are marked
  by   spaces. It  doesn't matter how many
spaces    you type; one is as good as 100.  The
end of   a line counts as a space.

One   or more   blank lines denote the  end
of  a paragraph.

Since any number of consecutive spaces are treated
like a single one, the formatting of the input
file makes no difference to
      \LaTeX,                % The \LaTeX command generates the LaTeX logo.
but it makes a difference to you.  When you use
\LaTeX, making your input file as easy to read
as possible will be a great help as you write
your document and when you change it.  This sample
file shows how you can add comments to your own input
file.

Because printing is different from typewriting,
there are a number of things that you have to do
differently when preparing an input file than if
you were just typing the document directly.
Quotation marks like
       ``this''
have to be handled specially, as do quotes within
quotes:
       ``\,`this'            % \, separates the double and single quote.
        is what I just
        wrote, not  `that'\,''.

Dashes come in three sizes: an
       intra-word
dash, a medium dash for number ranges like
       1--2,
and a punctuation
       dash---like
this.

A sentence-ending space should be larger than the
space between words within a sentence.  You
sometimes have to type special commands in
conjunction with punctuation characters to get
this right, as in the following sentence.
       Gnats, gnus, etc.\ all  % `\ ' makes an inter-word space.
       begin with G\@.         % \@ marks end-of-sentence punctuation.
You should check the spaces after periods when
reading your output to make sure you haven't
forgotten any special cases.  Generating an
ellipsis
       \ldots\               % `\ ' is needed after `\ldots' because TeX
                             % ignores spaces after command names like \ldots
                             % made from \ + letters.
                             %
                             % Note how a `%' character causes TeX to ignore
                             % the end of the input line, so these blank lines
                             % do not start a new paragraph.
                             %
with the right spacing around the periods requires
a special command.

\LaTeX\ interprets some common characters as
commands, so you must type special commands to
generate them.  These characters include the
following:
       \$ \& \% \# \{ and \}.

In printing, text is usually emphasized with an
       \emph{italic}
type style.

\begin{em}
   A long segment of text can also be emphasized
   in this way.  Text within such a segment can be
   given \emph{additional} emphasis.
\end{em}

It is sometimes necessary to prevent \LaTeX\ from
breaking a line where it might otherwise do so.
This may be at a space, as between the ``Mr.''\ and
``Jones'' in
       ``Mr.~Jones'',        % ~ produces an unbreakable interword space.
or within a word---especially when the word is a
symbol like
       \mbox{\emph{itemnum}}
that makes little sense when hyphenated across
lines.

Footnotes\footnote{This is an example of a footnote.}
pose no problem.

\LaTeX\ is good at typesetting mathematical formulas
like
       \( x-3y + z = 7 \)
or
       \( a_{1} > x^{2n} + y^{2n} > x' \)
or
       \( \ip{A}{B} = \sum_{i} a_{i} b_{i} \).
The spaces you type in a formula are
ignored.  Remember that a letter like
       $x$                   % $ ... $  and  \( ... \)  are equivalent
is a formula when it denotes a mathematical
symbol, and it should be typed as one.

\section{Displayed Text}

Text is displayed by indenting it from the left
margin.  Quotations are commonly displayed.  There
are short quotations
\begin{quote}
   This is a short quotation.  It consists of a
   single paragraph of text.  See how it is formatted.
\end{quote}
and longer ones.
\begin{quotation}
   This is a longer quotation.  It consists of two
   paragraphs of text, neither of which are
   particularly interesting.

   This is the second paragraph of the quotation.  It
   is just as dull as the first paragraph.
\end{quotation}
Another frequently-displayed structure is a list.
The following is an example of an \emph{itemized}
list.
\begin{itemize}
   \item This is the first item of an itemized list.
         Each item in the list is marked with a ``tick''.
         You don't have to worry about what kind of tick
         mark is used.

   \item This is the second item of the list.  It
         contains another list nested inside it.  The inner
         list is an \emph{enumerated} list.
         \begin{enumerate}
            \item This is the first item of an enumerated
                  list that is nested within the itemized list.

            \item This is the second item of the inner list.
                  \LaTeX\ allows you to nest lists deeper than
                  you really should.
         \end{enumerate}
         This is the rest of the second item of the outer
         list.  It is no more interesting than any other
         part of the item.
   \item This is the third item of the list.
\end{itemize}
You can even display poetry.
\begin{verse}
   There is an environment
    for verse \\             % The \\ command separates lines
   Whose features some poets % within a stanza.
   will curse.

                             % One or more blank lines separate stanzas.

   For instead of making\\
   Them do \emph{all} line breaking, \\
   It allows them to put too many words on a line when they'd rather be
   forced to be terse.
\end{verse}

Mathematical formulas may also be displayed.  A
displayed formula
is
one-line long; multiline
formulas require special formatting instructions.
   \[  \ip{\Gamma}{\psi'} = x'' + y^{2} + z_{i}^{n}\]
Don't start a paragraph with a displayed equation,
nor make one a paragraph by itself.

\end{document}               % End of document.

\section{Bench Huge Section 410}
% synthetic bench section 410
% This is a sample LaTeX input file.  (Version of 12 August 2004.)
%
% A '%' character causes TeX to ignore all remaining text on the line,
% and is used for comments like this one.

\documentclass{article}      % Specifies the document class

                             % The preamble begins here.
\title{An Example Document}  % Declares the document's title.
\author{Leslie Lamport}      % Declares the author's name.
\date{January 21, 1994}      % Deleting this command produces today's date.

\newcommand{\ip}[2]{(#1, #2)}
                             % Defines \ip{arg1}{arg2} to mean
                             % (arg1, arg2).

%\newcommand{\ip}[2]{\langle #1 | #2\rangle}
                             % This is an alternative definition of
                             % \ip that is commented out.

\begin{document}             % End of preamble and beginning of text.

\maketitle                   % Produces the title.

This is an example input file.  Comparing it with
the output it generates can show you how to
produce a simple document of your own.

\section{Ordinary Text}      % Produces section heading.  Lower-level
                             % sections are begun with similar
                             % \subsection and \subsubsection commands.

The ends  of words and sentences are marked
  by   spaces. It  doesn't matter how many
spaces    you type; one is as good as 100.  The
end of   a line counts as a space.

One   or more   blank lines denote the  end
of  a paragraph.

Since any number of consecutive spaces are treated
like a single one, the formatting of the input
file makes no difference to
      \LaTeX,                % The \LaTeX command generates the LaTeX logo.
but it makes a difference to you.  When you use
\LaTeX, making your input file as easy to read
as possible will be a great help as you write
your document and when you change it.  This sample
file shows how you can add comments to your own input
file.

Because printing is different from typewriting,
there are a number of things that you have to do
differently when preparing an input file than if
you were just typing the document directly.
Quotation marks like
       ``this''
have to be handled specially, as do quotes within
quotes:
       ``\,`this'            % \, separates the double and single quote.
        is what I just
        wrote, not  `that'\,''.

Dashes come in three sizes: an
       intra-word
dash, a medium dash for number ranges like
       1--2,
and a punctuation
       dash---like
this.

A sentence-ending space should be larger than the
space between words within a sentence.  You
sometimes have to type special commands in
conjunction with punctuation characters to get
this right, as in the following sentence.
       Gnats, gnus, etc.\ all  % `\ ' makes an inter-word space.
       begin with G\@.         % \@ marks end-of-sentence punctuation.
You should check the spaces after periods when
reading your output to make sure you haven't
forgotten any special cases.  Generating an
ellipsis
       \ldots\               % `\ ' is needed after `\ldots' because TeX
                             % ignores spaces after command names like \ldots
                             % made from \ + letters.
                             %
                             % Note how a `%' character causes TeX to ignore
                             % the end of the input line, so these blank lines
                             % do not start a new paragraph.
                             %
with the right spacing around the periods requires
a special command.

\LaTeX\ interprets some common characters as
commands, so you must type special commands to
generate them.  These characters include the
following:
       \$ \& \% \# \{ and \}.

In printing, text is usually emphasized with an
       \emph{italic}
type style.

\begin{em}
   A long segment of text can also be emphasized
   in this way.  Text within such a segment can be
   given \emph{additional} emphasis.
\end{em}

It is sometimes necessary to prevent \LaTeX\ from
breaking a line where it might otherwise do so.
This may be at a space, as between the ``Mr.''\ and
``Jones'' in
       ``Mr.~Jones'',        % ~ produces an unbreakable interword space.
or within a word---especially when the word is a
symbol like
       \mbox{\emph{itemnum}}
that makes little sense when hyphenated across
lines.

Footnotes\footnote{This is an example of a footnote.}
pose no problem.

\LaTeX\ is good at typesetting mathematical formulas
like
       \( x-3y + z = 7 \)
or
       \( a_{1} > x^{2n} + y^{2n} > x' \)
or
       \( \ip{A}{B} = \sum_{i} a_{i} b_{i} \).
The spaces you type in a formula are
ignored.  Remember that a letter like
       $x$                   % $ ... $  and  \( ... \)  are equivalent
is a formula when it denotes a mathematical
symbol, and it should be typed as one.

\section{Displayed Text}

Text is displayed by indenting it from the left
margin.  Quotations are commonly displayed.  There
are short quotations
\begin{quote}
   This is a short quotation.  It consists of a
   single paragraph of text.  See how it is formatted.
\end{quote}
and longer ones.
\begin{quotation}
   This is a longer quotation.  It consists of two
   paragraphs of text, neither of which are
   particularly interesting.

   This is the second paragraph of the quotation.  It
   is just as dull as the first paragraph.
\end{quotation}
Another frequently-displayed structure is a list.
The following is an example of an \emph{itemized}
list.
\begin{itemize}
   \item This is the first item of an itemized list.
         Each item in the list is marked with a ``tick''.
         You don't have to worry about what kind of tick
         mark is used.

   \item This is the second item of the list.  It
         contains another list nested inside it.  The inner
         list is an \emph{enumerated} list.
         \begin{enumerate}
            \item This is the first item of an enumerated
                  list that is nested within the itemized list.

            \item This is the second item of the inner list.
                  \LaTeX\ allows you to nest lists deeper than
                  you really should.
         \end{enumerate}
         This is the rest of the second item of the outer
         list.  It is no more interesting than any other
         part of the item.
   \item This is the third item of the list.
\end{itemize}
You can even display poetry.
\begin{verse}
   There is an environment
    for verse \\             % The \\ command separates lines
   Whose features some poets % within a stanza.
   will curse.

                             % One or more blank lines separate stanzas.

   For instead of making\\
   Them do \emph{all} line breaking, \\
   It allows them to put too many words on a line when they'd rather be
   forced to be terse.
\end{verse}

Mathematical formulas may also be displayed.  A
displayed formula
is
one-line long; multiline
formulas require special formatting instructions.
   \[  \ip{\Gamma}{\psi'} = x'' + y^{2} + z_{i}^{n}\]
Don't start a paragraph with a displayed equation,
nor make one a paragraph by itself.

\end{document}               % End of document.

\section{Bench Huge Section 411}
% synthetic bench section 411
% This is a sample LaTeX input file.  (Version of 12 August 2004.)
%
% A '%' character causes TeX to ignore all remaining text on the line,
% and is used for comments like this one.

\documentclass{article}      % Specifies the document class

                             % The preamble begins here.
\title{An Example Document}  % Declares the document's title.
\author{Leslie Lamport}      % Declares the author's name.
\date{January 21, 1994}      % Deleting this command produces today's date.

\newcommand{\ip}[2]{(#1, #2)}
                             % Defines \ip{arg1}{arg2} to mean
                             % (arg1, arg2).

%\newcommand{\ip}[2]{\langle #1 | #2\rangle}
                             % This is an alternative definition of
                             % \ip that is commented out.

\begin{document}             % End of preamble and beginning of text.

\maketitle                   % Produces the title.

This is an example input file.  Comparing it with
the output it generates can show you how to
produce a simple document of your own.

\section{Ordinary Text}      % Produces section heading.  Lower-level
                             % sections are begun with similar
                             % \subsection and \subsubsection commands.

The ends  of words and sentences are marked
  by   spaces. It  doesn't matter how many
spaces    you type; one is as good as 100.  The
end of   a line counts as a space.

One   or more   blank lines denote the  end
of  a paragraph.

Since any number of consecutive spaces are treated
like a single one, the formatting of the input
file makes no difference to
      \LaTeX,                % The \LaTeX command generates the LaTeX logo.
but it makes a difference to you.  When you use
\LaTeX, making your input file as easy to read
as possible will be a great help as you write
your document and when you change it.  This sample
file shows how you can add comments to your own input
file.

Because printing is different from typewriting,
there are a number of things that you have to do
differently when preparing an input file than if
you were just typing the document directly.
Quotation marks like
       ``this''
have to be handled specially, as do quotes within
quotes:
       ``\,`this'            % \, separates the double and single quote.
        is what I just
        wrote, not  `that'\,''.

Dashes come in three sizes: an
       intra-word
dash, a medium dash for number ranges like
       1--2,
and a punctuation
       dash---like
this.

A sentence-ending space should be larger than the
space between words within a sentence.  You
sometimes have to type special commands in
conjunction with punctuation characters to get
this right, as in the following sentence.
       Gnats, gnus, etc.\ all  % `\ ' makes an inter-word space.
       begin with G\@.         % \@ marks end-of-sentence punctuation.
You should check the spaces after periods when
reading your output to make sure you haven't
forgotten any special cases.  Generating an
ellipsis
       \ldots\               % `\ ' is needed after `\ldots' because TeX
                             % ignores spaces after command names like \ldots
                             % made from \ + letters.
                             %
                             % Note how a `%' character causes TeX to ignore
                             % the end of the input line, so these blank lines
                             % do not start a new paragraph.
                             %
with the right spacing around the periods requires
a special command.

\LaTeX\ interprets some common characters as
commands, so you must type special commands to
generate them.  These characters include the
following:
       \$ \& \% \# \{ and \}.

In printing, text is usually emphasized with an
       \emph{italic}
type style.

\begin{em}
   A long segment of text can also be emphasized
   in this way.  Text within such a segment can be
   given \emph{additional} emphasis.
\end{em}

It is sometimes necessary to prevent \LaTeX\ from
breaking a line where it might otherwise do so.
This may be at a space, as between the ``Mr.''\ and
``Jones'' in
       ``Mr.~Jones'',        % ~ produces an unbreakable interword space.
or within a word---especially when the word is a
symbol like
       \mbox{\emph{itemnum}}
that makes little sense when hyphenated across
lines.

Footnotes\footnote{This is an example of a footnote.}
pose no problem.

\LaTeX\ is good at typesetting mathematical formulas
like
       \( x-3y + z = 7 \)
or
       \( a_{1} > x^{2n} + y^{2n} > x' \)
or
       \( \ip{A}{B} = \sum_{i} a_{i} b_{i} \).
The spaces you type in a formula are
ignored.  Remember that a letter like
       $x$                   % $ ... $  and  \( ... \)  are equivalent
is a formula when it denotes a mathematical
symbol, and it should be typed as one.

\section{Displayed Text}

Text is displayed by indenting it from the left
margin.  Quotations are commonly displayed.  There
are short quotations
\begin{quote}
   This is a short quotation.  It consists of a
   single paragraph of text.  See how it is formatted.
\end{quote}
and longer ones.
\begin{quotation}
   This is a longer quotation.  It consists of two
   paragraphs of text, neither of which are
   particularly interesting.

   This is the second paragraph of the quotation.  It
   is just as dull as the first paragraph.
\end{quotation}
Another frequently-displayed structure is a list.
The following is an example of an \emph{itemized}
list.
\begin{itemize}
   \item This is the first item of an itemized list.
         Each item in the list is marked with a ``tick''.
         You don't have to worry about what kind of tick
         mark is used.

   \item This is the second item of the list.  It
         contains another list nested inside it.  The inner
         list is an \emph{enumerated} list.
         \begin{enumerate}
            \item This is the first item of an enumerated
                  list that is nested within the itemized list.

            \item This is the second item of the inner list.
                  \LaTeX\ allows you to nest lists deeper than
                  you really should.
         \end{enumerate}
         This is the rest of the second item of the outer
         list.  It is no more interesting than any other
         part of the item.
   \item This is the third item of the list.
\end{itemize}
You can even display poetry.
\begin{verse}
   There is an environment
    for verse \\             % The \\ command separates lines
   Whose features some poets % within a stanza.
   will curse.

                             % One or more blank lines separate stanzas.

   For instead of making\\
   Them do \emph{all} line breaking, \\
   It allows them to put too many words on a line when they'd rather be
   forced to be terse.
\end{verse}

Mathematical formulas may also be displayed.  A
displayed formula
is
one-line long; multiline
formulas require special formatting instructions.
   \[  \ip{\Gamma}{\psi'} = x'' + y^{2} + z_{i}^{n}\]
Don't start a paragraph with a displayed equation,
nor make one a paragraph by itself.

\end{document}               % End of document.

\section{Bench Huge Section 412}
% synthetic bench section 412
% This is a sample LaTeX input file.  (Version of 12 August 2004.)
%
% A '%' character causes TeX to ignore all remaining text on the line,
% and is used for comments like this one.

\documentclass{article}      % Specifies the document class

                             % The preamble begins here.
\title{An Example Document}  % Declares the document's title.
\author{Leslie Lamport}      % Declares the author's name.
\date{January 21, 1994}      % Deleting this command produces today's date.

\newcommand{\ip}[2]{(#1, #2)}
                             % Defines \ip{arg1}{arg2} to mean
                             % (arg1, arg2).

%\newcommand{\ip}[2]{\langle #1 | #2\rangle}
                             % This is an alternative definition of
                             % \ip that is commented out.

\begin{document}             % End of preamble and beginning of text.

\maketitle                   % Produces the title.

This is an example input file.  Comparing it with
the output it generates can show you how to
produce a simple document of your own.

\section{Ordinary Text}      % Produces section heading.  Lower-level
                             % sections are begun with similar
                             % \subsection and \subsubsection commands.

The ends  of words and sentences are marked
  by   spaces. It  doesn't matter how many
spaces    you type; one is as good as 100.  The
end of   a line counts as a space.

One   or more   blank lines denote the  end
of  a paragraph.

Since any number of consecutive spaces are treated
like a single one, the formatting of the input
file makes no difference to
      \LaTeX,                % The \LaTeX command generates the LaTeX logo.
but it makes a difference to you.  When you use
\LaTeX, making your input file as easy to read
as possible will be a great help as you write
your document and when you change it.  This sample
file shows how you can add comments to your own input
file.

Because printing is different from typewriting,
there are a number of things that you have to do
differently when preparing an input file than if
you were just typing the document directly.
Quotation marks like
       ``this''
have to be handled specially, as do quotes within
quotes:
       ``\,`this'            % \, separates the double and single quote.
        is what I just
        wrote, not  `that'\,''.

Dashes come in three sizes: an
       intra-word
dash, a medium dash for number ranges like
       1--2,
and a punctuation
       dash---like
this.

A sentence-ending space should be larger than the
space between words within a sentence.  You
sometimes have to type special commands in
conjunction with punctuation characters to get
this right, as in the following sentence.
       Gnats, gnus, etc.\ all  % `\ ' makes an inter-word space.
       begin with G\@.         % \@ marks end-of-sentence punctuation.
You should check the spaces after periods when
reading your output to make sure you haven't
forgotten any special cases.  Generating an
ellipsis
       \ldots\               % `\ ' is needed after `\ldots' because TeX
                             % ignores spaces after command names like \ldots
                             % made from \ + letters.
                             %
                             % Note how a `%' character causes TeX to ignore
                             % the end of the input line, so these blank lines
                             % do not start a new paragraph.
                             %
with the right spacing around the periods requires
a special command.

\LaTeX\ interprets some common characters as
commands, so you must type special commands to
generate them.  These characters include the
following:
       \$ \& \% \# \{ and \}.

In printing, text is usually emphasized with an
       \emph{italic}
type style.

\begin{em}
   A long segment of text can also be emphasized
   in this way.  Text within such a segment can be
   given \emph{additional} emphasis.
\end{em}

It is sometimes necessary to prevent \LaTeX\ from
breaking a line where it might otherwise do so.
This may be at a space, as between the ``Mr.''\ and
``Jones'' in
       ``Mr.~Jones'',        % ~ produces an unbreakable interword space.
or within a word---especially when the word is a
symbol like
       \mbox{\emph{itemnum}}
that makes little sense when hyphenated across
lines.

Footnotes\footnote{This is an example of a footnote.}
pose no problem.

\LaTeX\ is good at typesetting mathematical formulas
like
       \( x-3y + z = 7 \)
or
       \( a_{1} > x^{2n} + y^{2n} > x' \)
or
       \( \ip{A}{B} = \sum_{i} a_{i} b_{i} \).
The spaces you type in a formula are
ignored.  Remember that a letter like
       $x$                   % $ ... $  and  \( ... \)  are equivalent
is a formula when it denotes a mathematical
symbol, and it should be typed as one.

\section{Displayed Text}

Text is displayed by indenting it from the left
margin.  Quotations are commonly displayed.  There
are short quotations
\begin{quote}
   This is a short quotation.  It consists of a
   single paragraph of text.  See how it is formatted.
\end{quote}
and longer ones.
\begin{quotation}
   This is a longer quotation.  It consists of two
   paragraphs of text, neither of which are
   particularly interesting.

   This is the second paragraph of the quotation.  It
   is just as dull as the first paragraph.
\end{quotation}
Another frequently-displayed structure is a list.
The following is an example of an \emph{itemized}
list.
\begin{itemize}
   \item This is the first item of an itemized list.
         Each item in the list is marked with a ``tick''.
         You don't have to worry about what kind of tick
         mark is used.

   \item This is the second item of the list.  It
         contains another list nested inside it.  The inner
         list is an \emph{enumerated} list.
         \begin{enumerate}
            \item This is the first item of an enumerated
                  list that is nested within the itemized list.

            \item This is the second item of the inner list.
                  \LaTeX\ allows you to nest lists deeper than
                  you really should.
         \end{enumerate}
         This is the rest of the second item of the outer
         list.  It is no more interesting than any other
         part of the item.
   \item This is the third item of the list.
\end{itemize}
You can even display poetry.
\begin{verse}
   There is an environment
    for verse \\             % The \\ command separates lines
   Whose features some poets % within a stanza.
   will curse.

                             % One or more blank lines separate stanzas.

   For instead of making\\
   Them do \emph{all} line breaking, \\
   It allows them to put too many words on a line when they'd rather be
   forced to be terse.
\end{verse}

Mathematical formulas may also be displayed.  A
displayed formula
is
one-line long; multiline
formulas require special formatting instructions.
   \[  \ip{\Gamma}{\psi'} = x'' + y^{2} + z_{i}^{n}\]
Don't start a paragraph with a displayed equation,
nor make one a paragraph by itself.

\end{document}               % End of document.

\section{Bench Huge Section 413}
% synthetic bench section 413
% This is a sample LaTeX input file.  (Version of 12 August 2004.)
%
% A '%' character causes TeX to ignore all remaining text on the line,
% and is used for comments like this one.

\documentclass{article}      % Specifies the document class

                             % The preamble begins here.
\title{An Example Document}  % Declares the document's title.
\author{Leslie Lamport}      % Declares the author's name.
\date{January 21, 1994}      % Deleting this command produces today's date.

\newcommand{\ip}[2]{(#1, #2)}
                             % Defines \ip{arg1}{arg2} to mean
                             % (arg1, arg2).

%\newcommand{\ip}[2]{\langle #1 | #2\rangle}
                             % This is an alternative definition of
                             % \ip that is commented out.

\begin{document}             % End of preamble and beginning of text.

\maketitle                   % Produces the title.

This is an example input file.  Comparing it with
the output it generates can show you how to
produce a simple document of your own.

\section{Ordinary Text}      % Produces section heading.  Lower-level
                             % sections are begun with similar
                             % \subsection and \subsubsection commands.

The ends  of words and sentences are marked
  by   spaces. It  doesn't matter how many
spaces    you type; one is as good as 100.  The
end of   a line counts as a space.

One   or more   blank lines denote the  end
of  a paragraph.

Since any number of consecutive spaces are treated
like a single one, the formatting of the input
file makes no difference to
      \LaTeX,                % The \LaTeX command generates the LaTeX logo.
but it makes a difference to you.  When you use
\LaTeX, making your input file as easy to read
as possible will be a great help as you write
your document and when you change it.  This sample
file shows how you can add comments to your own input
file.

Because printing is different from typewriting,
there are a number of things that you have to do
differently when preparing an input file than if
you were just typing the document directly.
Quotation marks like
       ``this''
have to be handled specially, as do quotes within
quotes:
       ``\,`this'            % \, separates the double and single quote.
        is what I just
        wrote, not  `that'\,''.

Dashes come in three sizes: an
       intra-word
dash, a medium dash for number ranges like
       1--2,
and a punctuation
       dash---like
this.

A sentence-ending space should be larger than the
space between words within a sentence.  You
sometimes have to type special commands in
conjunction with punctuation characters to get
this right, as in the following sentence.
       Gnats, gnus, etc.\ all  % `\ ' makes an inter-word space.
       begin with G\@.         % \@ marks end-of-sentence punctuation.
You should check the spaces after periods when
reading your output to make sure you haven't
forgotten any special cases.  Generating an
ellipsis
       \ldots\               % `\ ' is needed after `\ldots' because TeX
                             % ignores spaces after command names like \ldots
                             % made from \ + letters.
                             %
                             % Note how a `%' character causes TeX to ignore
                             % the end of the input line, so these blank lines
                             % do not start a new paragraph.
                             %
with the right spacing around the periods requires
a special command.

\LaTeX\ interprets some common characters as
commands, so you must type special commands to
generate them.  These characters include the
following:
       \$ \& \% \# \{ and \}.

In printing, text is usually emphasized with an
       \emph{italic}
type style.

\begin{em}
   A long segment of text can also be emphasized
   in this way.  Text within such a segment can be
   given \emph{additional} emphasis.
\end{em}

It is sometimes necessary to prevent \LaTeX\ from
breaking a line where it might otherwise do so.
This may be at a space, as between the ``Mr.''\ and
``Jones'' in
       ``Mr.~Jones'',        % ~ produces an unbreakable interword space.
or within a word---especially when the word is a
symbol like
       \mbox{\emph{itemnum}}
that makes little sense when hyphenated across
lines.

Footnotes\footnote{This is an example of a footnote.}
pose no problem.

\LaTeX\ is good at typesetting mathematical formulas
like
       \( x-3y + z = 7 \)
or
       \( a_{1} > x^{2n} + y^{2n} > x' \)
or
       \( \ip{A}{B} = \sum_{i} a_{i} b_{i} \).
The spaces you type in a formula are
ignored.  Remember that a letter like
       $x$                   % $ ... $  and  \( ... \)  are equivalent
is a formula when it denotes a mathematical
symbol, and it should be typed as one.

\section{Displayed Text}

Text is displayed by indenting it from the left
margin.  Quotations are commonly displayed.  There
are short quotations
\begin{quote}
   This is a short quotation.  It consists of a
   single paragraph of text.  See how it is formatted.
\end{quote}
and longer ones.
\begin{quotation}
   This is a longer quotation.  It consists of two
   paragraphs of text, neither of which are
   particularly interesting.

   This is the second paragraph of the quotation.  It
   is just as dull as the first paragraph.
\end{quotation}
Another frequently-displayed structure is a list.
The following is an example of an \emph{itemized}
list.
\begin{itemize}
   \item This is the first item of an itemized list.
         Each item in the list is marked with a ``tick''.
         You don't have to worry about what kind of tick
         mark is used.

   \item This is the second item of the list.  It
         contains another list nested inside it.  The inner
         list is an \emph{enumerated} list.
         \begin{enumerate}
            \item This is the first item of an enumerated
                  list that is nested within the itemized list.

            \item This is the second item of the inner list.
                  \LaTeX\ allows you to nest lists deeper than
                  you really should.
         \end{enumerate}
         This is the rest of the second item of the outer
         list.  It is no more interesting than any other
         part of the item.
   \item This is the third item of the list.
\end{itemize}
You can even display poetry.
\begin{verse}
   There is an environment
    for verse \\             % The \\ command separates lines
   Whose features some poets % within a stanza.
   will curse.

                             % One or more blank lines separate stanzas.

   For instead of making\\
   Them do \emph{all} line breaking, \\
   It allows them to put too many words on a line when they'd rather be
   forced to be terse.
\end{verse}

Mathematical formulas may also be displayed.  A
displayed formula
is
one-line long; multiline
formulas require special formatting instructions.
   \[  \ip{\Gamma}{\psi'} = x'' + y^{2} + z_{i}^{n}\]
Don't start a paragraph with a displayed equation,
nor make one a paragraph by itself.

\end{document}               % End of document.

\section{Bench Huge Section 414}
% synthetic bench section 414
% This is a sample LaTeX input file.  (Version of 12 August 2004.)
%
% A '%' character causes TeX to ignore all remaining text on the line,
% and is used for comments like this one.

\documentclass{article}      % Specifies the document class

                             % The preamble begins here.
\title{An Example Document}  % Declares the document's title.
\author{Leslie Lamport}      % Declares the author's name.
\date{January 21, 1994}      % Deleting this command produces today's date.

\newcommand{\ip}[2]{(#1, #2)}
                             % Defines \ip{arg1}{arg2} to mean
                             % (arg1, arg2).

%\newcommand{\ip}[2]{\langle #1 | #2\rangle}
                             % This is an alternative definition of
                             % \ip that is commented out.

\begin{document}             % End of preamble and beginning of text.

\maketitle                   % Produces the title.

This is an example input file.  Comparing it with
the output it generates can show you how to
produce a simple document of your own.

\section{Ordinary Text}      % Produces section heading.  Lower-level
                             % sections are begun with similar
                             % \subsection and \subsubsection commands.

The ends  of words and sentences are marked
  by   spaces. It  doesn't matter how many
spaces    you type; one is as good as 100.  The
end of   a line counts as a space.

One   or more   blank lines denote the  end
of  a paragraph.

Since any number of consecutive spaces are treated
like a single one, the formatting of the input
file makes no difference to
      \LaTeX,                % The \LaTeX command generates the LaTeX logo.
but it makes a difference to you.  When you use
\LaTeX, making your input file as easy to read
as possible will be a great help as you write
your document and when you change it.  This sample
file shows how you can add comments to your own input
file.

Because printing is different from typewriting,
there are a number of things that you have to do
differently when preparing an input file than if
you were just typing the document directly.
Quotation marks like
       ``this''
have to be handled specially, as do quotes within
quotes:
       ``\,`this'            % \, separates the double and single quote.
        is what I just
        wrote, not  `that'\,''.

Dashes come in three sizes: an
       intra-word
dash, a medium dash for number ranges like
       1--2,
and a punctuation
       dash---like
this.

A sentence-ending space should be larger than the
space between words within a sentence.  You
sometimes have to type special commands in
conjunction with punctuation characters to get
this right, as in the following sentence.
       Gnats, gnus, etc.\ all  % `\ ' makes an inter-word space.
       begin with G\@.         % \@ marks end-of-sentence punctuation.
You should check the spaces after periods when
reading your output to make sure you haven't
forgotten any special cases.  Generating an
ellipsis
       \ldots\               % `\ ' is needed after `\ldots' because TeX
                             % ignores spaces after command names like \ldots
                             % made from \ + letters.
                             %
                             % Note how a `%' character causes TeX to ignore
                             % the end of the input line, so these blank lines
                             % do not start a new paragraph.
                             %
with the right spacing around the periods requires
a special command.

\LaTeX\ interprets some common characters as
commands, so you must type special commands to
generate them.  These characters include the
following:
       \$ \& \% \# \{ and \}.

In printing, text is usually emphasized with an
       \emph{italic}
type style.

\begin{em}
   A long segment of text can also be emphasized
   in this way.  Text within such a segment can be
   given \emph{additional} emphasis.
\end{em}

It is sometimes necessary to prevent \LaTeX\ from
breaking a line where it might otherwise do so.
This may be at a space, as between the ``Mr.''\ and
``Jones'' in
       ``Mr.~Jones'',        % ~ produces an unbreakable interword space.
or within a word---especially when the word is a
symbol like
       \mbox{\emph{itemnum}}
that makes little sense when hyphenated across
lines.

Footnotes\footnote{This is an example of a footnote.}
pose no problem.

\LaTeX\ is good at typesetting mathematical formulas
like
       \( x-3y + z = 7 \)
or
       \( a_{1} > x^{2n} + y^{2n} > x' \)
or
       \( \ip{A}{B} = \sum_{i} a_{i} b_{i} \).
The spaces you type in a formula are
ignored.  Remember that a letter like
       $x$                   % $ ... $  and  \( ... \)  are equivalent
is a formula when it denotes a mathematical
symbol, and it should be typed as one.

\section{Displayed Text}

Text is displayed by indenting it from the left
margin.  Quotations are commonly displayed.  There
are short quotations
\begin{quote}
   This is a short quotation.  It consists of a
   single paragraph of text.  See how it is formatted.
\end{quote}
and longer ones.
\begin{quotation}
   This is a longer quotation.  It consists of two
   paragraphs of text, neither of which are
   particularly interesting.

   This is the second paragraph of the quotation.  It
   is just as dull as the first paragraph.
\end{quotation}
Another frequently-displayed structure is a list.
The following is an example of an \emph{itemized}
list.
\begin{itemize}
   \item This is the first item of an itemized list.
         Each item in the list is marked with a ``tick''.
         You don't have to worry about what kind of tick
         mark is used.

   \item This is the second item of the list.  It
         contains another list nested inside it.  The inner
         list is an \emph{enumerated} list.
         \begin{enumerate}
            \item This is the first item of an enumerated
                  list that is nested within the itemized list.

            \item This is the second item of the inner list.
                  \LaTeX\ allows you to nest lists deeper than
                  you really should.
         \end{enumerate}
         This is the rest of the second item of the outer
         list.  It is no more interesting than any other
         part of the item.
   \item This is the third item of the list.
\end{itemize}
You can even display poetry.
\begin{verse}
   There is an environment
    for verse \\             % The \\ command separates lines
   Whose features some poets % within a stanza.
   will curse.

                             % One or more blank lines separate stanzas.

   For instead of making\\
   Them do \emph{all} line breaking, \\
   It allows them to put too many words on a line when they'd rather be
   forced to be terse.
\end{verse}

Mathematical formulas may also be displayed.  A
displayed formula
is
one-line long; multiline
formulas require special formatting instructions.
   \[  \ip{\Gamma}{\psi'} = x'' + y^{2} + z_{i}^{n}\]
Don't start a paragraph with a displayed equation,
nor make one a paragraph by itself.

\end{document}               % End of document.

\section{Bench Huge Section 415}
% synthetic bench section 415
% This is a sample LaTeX input file.  (Version of 12 August 2004.)
%
% A '%' character causes TeX to ignore all remaining text on the line,
% and is used for comments like this one.

\documentclass{article}      % Specifies the document class

                             % The preamble begins here.
\title{An Example Document}  % Declares the document's title.
\author{Leslie Lamport}      % Declares the author's name.
\date{January 21, 1994}      % Deleting this command produces today's date.

\newcommand{\ip}[2]{(#1, #2)}
                             % Defines \ip{arg1}{arg2} to mean
                             % (arg1, arg2).

%\newcommand{\ip}[2]{\langle #1 | #2\rangle}
                             % This is an alternative definition of
                             % \ip that is commented out.

\begin{document}             % End of preamble and beginning of text.

\maketitle                   % Produces the title.

This is an example input file.  Comparing it with
the output it generates can show you how to
produce a simple document of your own.

\section{Ordinary Text}      % Produces section heading.  Lower-level
                             % sections are begun with similar
                             % \subsection and \subsubsection commands.

The ends  of words and sentences are marked
  by   spaces. It  doesn't matter how many
spaces    you type; one is as good as 100.  The
end of   a line counts as a space.

One   or more   blank lines denote the  end
of  a paragraph.

Since any number of consecutive spaces are treated
like a single one, the formatting of the input
file makes no difference to
      \LaTeX,                % The \LaTeX command generates the LaTeX logo.
but it makes a difference to you.  When you use
\LaTeX, making your input file as easy to read
as possible will be a great help as you write
your document and when you change it.  This sample
file shows how you can add comments to your own input
file.

Because printing is different from typewriting,
there are a number of things that you have to do
differently when preparing an input file than if
you were just typing the document directly.
Quotation marks like
       ``this''
have to be handled specially, as do quotes within
quotes:
       ``\,`this'            % \, separates the double and single quote.
        is what I just
        wrote, not  `that'\,''.

Dashes come in three sizes: an
       intra-word
dash, a medium dash for number ranges like
       1--2,
and a punctuation
       dash---like
this.

A sentence-ending space should be larger than the
space between words within a sentence.  You
sometimes have to type special commands in
conjunction with punctuation characters to get
this right, as in the following sentence.
       Gnats, gnus, etc.\ all  % `\ ' makes an inter-word space.
       begin with G\@.         % \@ marks end-of-sentence punctuation.
You should check the spaces after periods when
reading your output to make sure you haven't
forgotten any special cases.  Generating an
ellipsis
       \ldots\               % `\ ' is needed after `\ldots' because TeX
                             % ignores spaces after command names like \ldots
                             % made from \ + letters.
                             %
                             % Note how a `%' character causes TeX to ignore
                             % the end of the input line, so these blank lines
                             % do not start a new paragraph.
                             %
with the right spacing around the periods requires
a special command.

\LaTeX\ interprets some common characters as
commands, so you must type special commands to
generate them.  These characters include the
following:
       \$ \& \% \# \{ and \}.

In printing, text is usually emphasized with an
       \emph{italic}
type style.

\begin{em}
   A long segment of text can also be emphasized
   in this way.  Text within such a segment can be
   given \emph{additional} emphasis.
\end{em}

It is sometimes necessary to prevent \LaTeX\ from
breaking a line where it might otherwise do so.
This may be at a space, as between the ``Mr.''\ and
``Jones'' in
       ``Mr.~Jones'',        % ~ produces an unbreakable interword space.
or within a word---especially when the word is a
symbol like
       \mbox{\emph{itemnum}}
that makes little sense when hyphenated across
lines.

Footnotes\footnote{This is an example of a footnote.}
pose no problem.

\LaTeX\ is good at typesetting mathematical formulas
like
       \( x-3y + z = 7 \)
or
       \( a_{1} > x^{2n} + y^{2n} > x' \)
or
       \( \ip{A}{B} = \sum_{i} a_{i} b_{i} \).
The spaces you type in a formula are
ignored.  Remember that a letter like
       $x$                   % $ ... $  and  \( ... \)  are equivalent
is a formula when it denotes a mathematical
symbol, and it should be typed as one.

\section{Displayed Text}

Text is displayed by indenting it from the left
margin.  Quotations are commonly displayed.  There
are short quotations
\begin{quote}
   This is a short quotation.  It consists of a
   single paragraph of text.  See how it is formatted.
\end{quote}
and longer ones.
\begin{quotation}
   This is a longer quotation.  It consists of two
   paragraphs of text, neither of which are
   particularly interesting.

   This is the second paragraph of the quotation.  It
   is just as dull as the first paragraph.
\end{quotation}
Another frequently-displayed structure is a list.
The following is an example of an \emph{itemized}
list.
\begin{itemize}
   \item This is the first item of an itemized list.
         Each item in the list is marked with a ``tick''.
         You don't have to worry about what kind of tick
         mark is used.

   \item This is the second item of the list.  It
         contains another list nested inside it.  The inner
         list is an \emph{enumerated} list.
         \begin{enumerate}
            \item This is the first item of an enumerated
                  list that is nested within the itemized list.

            \item This is the second item of the inner list.
                  \LaTeX\ allows you to nest lists deeper than
                  you really should.
         \end{enumerate}
         This is the rest of the second item of the outer
         list.  It is no more interesting than any other
         part of the item.
   \item This is the third item of the list.
\end{itemize}
You can even display poetry.
\begin{verse}
   There is an environment
    for verse \\             % The \\ command separates lines
   Whose features some poets % within a stanza.
   will curse.

                             % One or more blank lines separate stanzas.

   For instead of making\\
   Them do \emph{all} line breaking, \\
   It allows them to put too many words on a line when they'd rather be
   forced to be terse.
\end{verse}

Mathematical formulas may also be displayed.  A
displayed formula
is
one-line long; multiline
formulas require special formatting instructions.
   \[  \ip{\Gamma}{\psi'} = x'' + y^{2} + z_{i}^{n}\]
Don't start a paragraph with a displayed equation,
nor make one a paragraph by itself.

\end{document}               % End of document.

\section{Bench Huge Section 416}
% synthetic bench section 416
% This is a sample LaTeX input file.  (Version of 12 August 2004.)
%
% A '%' character causes TeX to ignore all remaining text on the line,
% and is used for comments like this one.

\documentclass{article}      % Specifies the document class

                             % The preamble begins here.
\title{An Example Document}  % Declares the document's title.
\author{Leslie Lamport}      % Declares the author's name.
\date{January 21, 1994}      % Deleting this command produces today's date.

\newcommand{\ip}[2]{(#1, #2)}
                             % Defines \ip{arg1}{arg2} to mean
                             % (arg1, arg2).

%\newcommand{\ip}[2]{\langle #1 | #2\rangle}
                             % This is an alternative definition of
                             % \ip that is commented out.

\begin{document}             % End of preamble and beginning of text.

\maketitle                   % Produces the title.

This is an example input file.  Comparing it with
the output it generates can show you how to
produce a simple document of your own.

\section{Ordinary Text}      % Produces section heading.  Lower-level
                             % sections are begun with similar
                             % \subsection and \subsubsection commands.

The ends  of words and sentences are marked
  by   spaces. It  doesn't matter how many
spaces    you type; one is as good as 100.  The
end of   a line counts as a space.

One   or more   blank lines denote the  end
of  a paragraph.

Since any number of consecutive spaces are treated
like a single one, the formatting of the input
file makes no difference to
      \LaTeX,                % The \LaTeX command generates the LaTeX logo.
but it makes a difference to you.  When you use
\LaTeX, making your input file as easy to read
as possible will be a great help as you write
your document and when you change it.  This sample
file shows how you can add comments to your own input
file.

Because printing is different from typewriting,
there are a number of things that you have to do
differently when preparing an input file than if
you were just typing the document directly.
Quotation marks like
       ``this''
have to be handled specially, as do quotes within
quotes:
       ``\,`this'            % \, separates the double and single quote.
        is what I just
        wrote, not  `that'\,''.

Dashes come in three sizes: an
       intra-word
dash, a medium dash for number ranges like
       1--2,
and a punctuation
       dash---like
this.

A sentence-ending space should be larger than the
space between words within a sentence.  You
sometimes have to type special commands in
conjunction with punctuation characters to get
this right, as in the following sentence.
       Gnats, gnus, etc.\ all  % `\ ' makes an inter-word space.
       begin with G\@.         % \@ marks end-of-sentence punctuation.
You should check the spaces after periods when
reading your output to make sure you haven't
forgotten any special cases.  Generating an
ellipsis
       \ldots\               % `\ ' is needed after `\ldots' because TeX
                             % ignores spaces after command names like \ldots
                             % made from \ + letters.
                             %
                             % Note how a `%' character causes TeX to ignore
                             % the end of the input line, so these blank lines
                             % do not start a new paragraph.
                             %
with the right spacing around the periods requires
a special command.

\LaTeX\ interprets some common characters as
commands, so you must type special commands to
generate them.  These characters include the
following:
       \$ \& \% \# \{ and \}.

In printing, text is usually emphasized with an
       \emph{italic}
type style.

\begin{em}
   A long segment of text can also be emphasized
   in this way.  Text within such a segment can be
   given \emph{additional} emphasis.
\end{em}

It is sometimes necessary to prevent \LaTeX\ from
breaking a line where it might otherwise do so.
This may be at a space, as between the ``Mr.''\ and
``Jones'' in
       ``Mr.~Jones'',        % ~ produces an unbreakable interword space.
or within a word---especially when the word is a
symbol like
       \mbox{\emph{itemnum}}
that makes little sense when hyphenated across
lines.

Footnotes\footnote{This is an example of a footnote.}
pose no problem.

\LaTeX\ is good at typesetting mathematical formulas
like
       \( x-3y + z = 7 \)
or
       \( a_{1} > x^{2n} + y^{2n} > x' \)
or
       \( \ip{A}{B} = \sum_{i} a_{i} b_{i} \).
The spaces you type in a formula are
ignored.  Remember that a letter like
       $x$                   % $ ... $  and  \( ... \)  are equivalent
is a formula when it denotes a mathematical
symbol, and it should be typed as one.

\section{Displayed Text}

Text is displayed by indenting it from the left
margin.  Quotations are commonly displayed.  There
are short quotations
\begin{quote}
   This is a short quotation.  It consists of a
   single paragraph of text.  See how it is formatted.
\end{quote}
and longer ones.
\begin{quotation}
   This is a longer quotation.  It consists of two
   paragraphs of text, neither of which are
   particularly interesting.

   This is the second paragraph of the quotation.  It
   is just as dull as the first paragraph.
\end{quotation}
Another frequently-displayed structure is a list.
The following is an example of an \emph{itemized}
list.
\begin{itemize}
   \item This is the first item of an itemized list.
         Each item in the list is marked with a ``tick''.
         You don't have to worry about what kind of tick
         mark is used.

   \item This is the second item of the list.  It
         contains another list nested inside it.  The inner
         list is an \emph{enumerated} list.
         \begin{enumerate}
            \item This is the first item of an enumerated
                  list that is nested within the itemized list.

            \item This is the second item of the inner list.
                  \LaTeX\ allows you to nest lists deeper than
                  you really should.
         \end{enumerate}
         This is the rest of the second item of the outer
         list.  It is no more interesting than any other
         part of the item.
   \item This is the third item of the list.
\end{itemize}
You can even display poetry.
\begin{verse}
   There is an environment
    for verse \\             % The \\ command separates lines
   Whose features some poets % within a stanza.
   will curse.

                             % One or more blank lines separate stanzas.

   For instead of making\\
   Them do \emph{all} line breaking, \\
   It allows them to put too many words on a line when they'd rather be
   forced to be terse.
\end{verse}

Mathematical formulas may also be displayed.  A
displayed formula
is
one-line long; multiline
formulas require special formatting instructions.
   \[  \ip{\Gamma}{\psi'} = x'' + y^{2} + z_{i}^{n}\]
Don't start a paragraph with a displayed equation,
nor make one a paragraph by itself.

\end{document}               % End of document.

\section{Bench Huge Section 417}
% synthetic bench section 417
% This is a sample LaTeX input file.  (Version of 12 August 2004.)
%
% A '%' character causes TeX to ignore all remaining text on the line,
% and is used for comments like this one.

\documentclass{article}      % Specifies the document class

                             % The preamble begins here.
\title{An Example Document}  % Declares the document's title.
\author{Leslie Lamport}      % Declares the author's name.
\date{January 21, 1994}      % Deleting this command produces today's date.

\newcommand{\ip}[2]{(#1, #2)}
                             % Defines \ip{arg1}{arg2} to mean
                             % (arg1, arg2).

%\newcommand{\ip}[2]{\langle #1 | #2\rangle}
                             % This is an alternative definition of
                             % \ip that is commented out.

\begin{document}             % End of preamble and beginning of text.

\maketitle                   % Produces the title.

This is an example input file.  Comparing it with
the output it generates can show you how to
produce a simple document of your own.

\section{Ordinary Text}      % Produces section heading.  Lower-level
                             % sections are begun with similar
                             % \subsection and \subsubsection commands.

The ends  of words and sentences are marked
  by   spaces. It  doesn't matter how many
spaces    you type; one is as good as 100.  The
end of   a line counts as a space.

One   or more   blank lines denote the  end
of  a paragraph.

Since any number of consecutive spaces are treated
like a single one, the formatting of the input
file makes no difference to
      \LaTeX,                % The \LaTeX command generates the LaTeX logo.
but it makes a difference to you.  When you use
\LaTeX, making your input file as easy to read
as possible will be a great help as you write
your document and when you change it.  This sample
file shows how you can add comments to your own input
file.

Because printing is different from typewriting,
there are a number of things that you have to do
differently when preparing an input file than if
you were just typing the document directly.
Quotation marks like
       ``this''
have to be handled specially, as do quotes within
quotes:
       ``\,`this'            % \, separates the double and single quote.
        is what I just
        wrote, not  `that'\,''.

Dashes come in three sizes: an
       intra-word
dash, a medium dash for number ranges like
       1--2,
and a punctuation
       dash---like
this.

A sentence-ending space should be larger than the
space between words within a sentence.  You
sometimes have to type special commands in
conjunction with punctuation characters to get
this right, as in the following sentence.
       Gnats, gnus, etc.\ all  % `\ ' makes an inter-word space.
       begin with G\@.         % \@ marks end-of-sentence punctuation.
You should check the spaces after periods when
reading your output to make sure you haven't
forgotten any special cases.  Generating an
ellipsis
       \ldots\               % `\ ' is needed after `\ldots' because TeX
                             % ignores spaces after command names like \ldots
                             % made from \ + letters.
                             %
                             % Note how a `%' character causes TeX to ignore
                             % the end of the input line, so these blank lines
                             % do not start a new paragraph.
                             %
with the right spacing around the periods requires
a special command.

\LaTeX\ interprets some common characters as
commands, so you must type special commands to
generate them.  These characters include the
following:
       \$ \& \% \# \{ and \}.

In printing, text is usually emphasized with an
       \emph{italic}
type style.

\begin{em}
   A long segment of text can also be emphasized
   in this way.  Text within such a segment can be
   given \emph{additional} emphasis.
\end{em}

It is sometimes necessary to prevent \LaTeX\ from
breaking a line where it might otherwise do so.
This may be at a space, as between the ``Mr.''\ and
``Jones'' in
       ``Mr.~Jones'',        % ~ produces an unbreakable interword space.
or within a word---especially when the word is a
symbol like
       \mbox{\emph{itemnum}}
that makes little sense when hyphenated across
lines.

Footnotes\footnote{This is an example of a footnote.}
pose no problem.

\LaTeX\ is good at typesetting mathematical formulas
like
       \( x-3y + z = 7 \)
or
       \( a_{1} > x^{2n} + y^{2n} > x' \)
or
       \( \ip{A}{B} = \sum_{i} a_{i} b_{i} \).
The spaces you type in a formula are
ignored.  Remember that a letter like
       $x$                   % $ ... $  and  \( ... \)  are equivalent
is a formula when it denotes a mathematical
symbol, and it should be typed as one.

\section{Displayed Text}

Text is displayed by indenting it from the left
margin.  Quotations are commonly displayed.  There
are short quotations
\begin{quote}
   This is a short quotation.  It consists of a
   single paragraph of text.  See how it is formatted.
\end{quote}
and longer ones.
\begin{quotation}
   This is a longer quotation.  It consists of two
   paragraphs of text, neither of which are
   particularly interesting.

   This is the second paragraph of the quotation.  It
   is just as dull as the first paragraph.
\end{quotation}
Another frequently-displayed structure is a list.
The following is an example of an \emph{itemized}
list.
\begin{itemize}
   \item This is the first item of an itemized list.
         Each item in the list is marked with a ``tick''.
         You don't have to worry about what kind of tick
         mark is used.

   \item This is the second item of the list.  It
         contains another list nested inside it.  The inner
         list is an \emph{enumerated} list.
         \begin{enumerate}
            \item This is the first item of an enumerated
                  list that is nested within the itemized list.

            \item This is the second item of the inner list.
                  \LaTeX\ allows you to nest lists deeper than
                  you really should.
         \end{enumerate}
         This is the rest of the second item of the outer
         list.  It is no more interesting than any other
         part of the item.
   \item This is the third item of the list.
\end{itemize}
You can even display poetry.
\begin{verse}
   There is an environment
    for verse \\             % The \\ command separates lines
   Whose features some poets % within a stanza.
   will curse.

                             % One or more blank lines separate stanzas.

   For instead of making\\
   Them do \emph{all} line breaking, \\
   It allows them to put too many words on a line when they'd rather be
   forced to be terse.
\end{verse}

Mathematical formulas may also be displayed.  A
displayed formula
is
one-line long; multiline
formulas require special formatting instructions.
   \[  \ip{\Gamma}{\psi'} = x'' + y^{2} + z_{i}^{n}\]
Don't start a paragraph with a displayed equation,
nor make one a paragraph by itself.

\end{document}               % End of document.

\section{Bench Huge Section 418}
% synthetic bench section 418
% This is a sample LaTeX input file.  (Version of 12 August 2004.)
%
% A '%' character causes TeX to ignore all remaining text on the line,
% and is used for comments like this one.

\documentclass{article}      % Specifies the document class

                             % The preamble begins here.
\title{An Example Document}  % Declares the document's title.
\author{Leslie Lamport}      % Declares the author's name.
\date{January 21, 1994}      % Deleting this command produces today's date.

\newcommand{\ip}[2]{(#1, #2)}
                             % Defines \ip{arg1}{arg2} to mean
                             % (arg1, arg2).

%\newcommand{\ip}[2]{\langle #1 | #2\rangle}
                             % This is an alternative definition of
                             % \ip that is commented out.

\begin{document}             % End of preamble and beginning of text.

\maketitle                   % Produces the title.

This is an example input file.  Comparing it with
the output it generates can show you how to
produce a simple document of your own.

\section{Ordinary Text}      % Produces section heading.  Lower-level
                             % sections are begun with similar
                             % \subsection and \subsubsection commands.

The ends  of words and sentences are marked
  by   spaces. It  doesn't matter how many
spaces    you type; one is as good as 100.  The
end of   a line counts as a space.

One   or more   blank lines denote the  end
of  a paragraph.

Since any number of consecutive spaces are treated
like a single one, the formatting of the input
file makes no difference to
      \LaTeX,                % The \LaTeX command generates the LaTeX logo.
but it makes a difference to you.  When you use
\LaTeX, making your input file as easy to read
as possible will be a great help as you write
your document and when you change it.  This sample
file shows how you can add comments to your own input
file.

Because printing is different from typewriting,
there are a number of things that you have to do
differently when preparing an input file than if
you were just typing the document directly.
Quotation marks like
       ``this''
have to be handled specially, as do quotes within
quotes:
       ``\,`this'            % \, separates the double and single quote.
        is what I just
        wrote, not  `that'\,''.

Dashes come in three sizes: an
       intra-word
dash, a medium dash for number ranges like
       1--2,
and a punctuation
       dash---like
this.

A sentence-ending space should be larger than the
space between words within a sentence.  You
sometimes have to type special commands in
conjunction with punctuation characters to get
this right, as in the following sentence.
       Gnats, gnus, etc.\ all  % `\ ' makes an inter-word space.
       begin with G\@.         % \@ marks end-of-sentence punctuation.
You should check the spaces after periods when
reading your output to make sure you haven't
forgotten any special cases.  Generating an
ellipsis
       \ldots\               % `\ ' is needed after `\ldots' because TeX
                             % ignores spaces after command names like \ldots
                             % made from \ + letters.
                             %
                             % Note how a `%' character causes TeX to ignore
                             % the end of the input line, so these blank lines
                             % do not start a new paragraph.
                             %
with the right spacing around the periods requires
a special command.

\LaTeX\ interprets some common characters as
commands, so you must type special commands to
generate them.  These characters include the
following:
       \$ \& \% \# \{ and \}.

In printing, text is usually emphasized with an
       \emph{italic}
type style.

\begin{em}
   A long segment of text can also be emphasized
   in this way.  Text within such a segment can be
   given \emph{additional} emphasis.
\end{em}

It is sometimes necessary to prevent \LaTeX\ from
breaking a line where it might otherwise do so.
This may be at a space, as between the ``Mr.''\ and
``Jones'' in
       ``Mr.~Jones'',        % ~ produces an unbreakable interword space.
or within a word---especially when the word is a
symbol like
       \mbox{\emph{itemnum}}
that makes little sense when hyphenated across
lines.

Footnotes\footnote{This is an example of a footnote.}
pose no problem.

\LaTeX\ is good at typesetting mathematical formulas
like
       \( x-3y + z = 7 \)
or
       \( a_{1} > x^{2n} + y^{2n} > x' \)
or
       \( \ip{A}{B} = \sum_{i} a_{i} b_{i} \).
The spaces you type in a formula are
ignored.  Remember that a letter like
       $x$                   % $ ... $  and  \( ... \)  are equivalent
is a formula when it denotes a mathematical
symbol, and it should be typed as one.

\section{Displayed Text}

Text is displayed by indenting it from the left
margin.  Quotations are commonly displayed.  There
are short quotations
\begin{quote}
   This is a short quotation.  It consists of a
   single paragraph of text.  See how it is formatted.
\end{quote}
and longer ones.
\begin{quotation}
   This is a longer quotation.  It consists of two
   paragraphs of text, neither of which are
   particularly interesting.

   This is the second paragraph of the quotation.  It
   is just as dull as the first paragraph.
\end{quotation}
Another frequently-displayed structure is a list.
The following is an example of an \emph{itemized}
list.
\begin{itemize}
   \item This is the first item of an itemized list.
         Each item in the list is marked with a ``tick''.
         You don't have to worry about what kind of tick
         mark is used.

   \item This is the second item of the list.  It
         contains another list nested inside it.  The inner
         list is an \emph{enumerated} list.
         \begin{enumerate}
            \item This is the first item of an enumerated
                  list that is nested within the itemized list.

            \item This is the second item of the inner list.
                  \LaTeX\ allows you to nest lists deeper than
                  you really should.
         \end{enumerate}
         This is the rest of the second item of the outer
         list.  It is no more interesting than any other
         part of the item.
   \item This is the third item of the list.
\end{itemize}
You can even display poetry.
\begin{verse}
   There is an environment
    for verse \\             % The \\ command separates lines
   Whose features some poets % within a stanza.
   will curse.

                             % One or more blank lines separate stanzas.

   For instead of making\\
   Them do \emph{all} line breaking, \\
   It allows them to put too many words on a line when they'd rather be
   forced to be terse.
\end{verse}

Mathematical formulas may also be displayed.  A
displayed formula
is
one-line long; multiline
formulas require special formatting instructions.
   \[  \ip{\Gamma}{\psi'} = x'' + y^{2} + z_{i}^{n}\]
Don't start a paragraph with a displayed equation,
nor make one a paragraph by itself.

\end{document}               % End of document.

\section{Bench Huge Section 419}
% synthetic bench section 419
% This is a sample LaTeX input file.  (Version of 12 August 2004.)
%
% A '%' character causes TeX to ignore all remaining text on the line,
% and is used for comments like this one.

\documentclass{article}      % Specifies the document class

                             % The preamble begins here.
\title{An Example Document}  % Declares the document's title.
\author{Leslie Lamport}      % Declares the author's name.
\date{January 21, 1994}      % Deleting this command produces today's date.

\newcommand{\ip}[2]{(#1, #2)}
                             % Defines \ip{arg1}{arg2} to mean
                             % (arg1, arg2).

%\newcommand{\ip}[2]{\langle #1 | #2\rangle}
                             % This is an alternative definition of
                             % \ip that is commented out.

\begin{document}             % End of preamble and beginning of text.

\maketitle                   % Produces the title.

This is an example input file.  Comparing it with
the output it generates can show you how to
produce a simple document of your own.

\section{Ordinary Text}      % Produces section heading.  Lower-level
                             % sections are begun with similar
                             % \subsection and \subsubsection commands.

The ends  of words and sentences are marked
  by   spaces. It  doesn't matter how many
spaces    you type; one is as good as 100.  The
end of   a line counts as a space.

One   or more   blank lines denote the  end
of  a paragraph.

Since any number of consecutive spaces are treated
like a single one, the formatting of the input
file makes no difference to
      \LaTeX,                % The \LaTeX command generates the LaTeX logo.
but it makes a difference to you.  When you use
\LaTeX, making your input file as easy to read
as possible will be a great help as you write
your document and when you change it.  This sample
file shows how you can add comments to your own input
file.

Because printing is different from typewriting,
there are a number of things that you have to do
differently when preparing an input file than if
you were just typing the document directly.
Quotation marks like
       ``this''
have to be handled specially, as do quotes within
quotes:
       ``\,`this'            % \, separates the double and single quote.
        is what I just
        wrote, not  `that'\,''.

Dashes come in three sizes: an
       intra-word
dash, a medium dash for number ranges like
       1--2,
and a punctuation
       dash---like
this.

A sentence-ending space should be larger than the
space between words within a sentence.  You
sometimes have to type special commands in
conjunction with punctuation characters to get
this right, as in the following sentence.
       Gnats, gnus, etc.\ all  % `\ ' makes an inter-word space.
       begin with G\@.         % \@ marks end-of-sentence punctuation.
You should check the spaces after periods when
reading your output to make sure you haven't
forgotten any special cases.  Generating an
ellipsis
       \ldots\               % `\ ' is needed after `\ldots' because TeX
                             % ignores spaces after command names like \ldots
                             % made from \ + letters.
                             %
                             % Note how a `%' character causes TeX to ignore
                             % the end of the input line, so these blank lines
                             % do not start a new paragraph.
                             %
with the right spacing around the periods requires
a special command.

\LaTeX\ interprets some common characters as
commands, so you must type special commands to
generate them.  These characters include the
following:
       \$ \& \% \# \{ and \}.

In printing, text is usually emphasized with an
       \emph{italic}
type style.

\begin{em}
   A long segment of text can also be emphasized
   in this way.  Text within such a segment can be
   given \emph{additional} emphasis.
\end{em}

It is sometimes necessary to prevent \LaTeX\ from
breaking a line where it might otherwise do so.
This may be at a space, as between the ``Mr.''\ and
``Jones'' in
       ``Mr.~Jones'',        % ~ produces an unbreakable interword space.
or within a word---especially when the word is a
symbol like
       \mbox{\emph{itemnum}}
that makes little sense when hyphenated across
lines.

Footnotes\footnote{This is an example of a footnote.}
pose no problem.

\LaTeX\ is good at typesetting mathematical formulas
like
       \( x-3y + z = 7 \)
or
       \( a_{1} > x^{2n} + y^{2n} > x' \)
or
       \( \ip{A}{B} = \sum_{i} a_{i} b_{i} \).
The spaces you type in a formula are
ignored.  Remember that a letter like
       $x$                   % $ ... $  and  \( ... \)  are equivalent
is a formula when it denotes a mathematical
symbol, and it should be typed as one.

\section{Displayed Text}

Text is displayed by indenting it from the left
margin.  Quotations are commonly displayed.  There
are short quotations
\begin{quote}
   This is a short quotation.  It consists of a
   single paragraph of text.  See how it is formatted.
\end{quote}
and longer ones.
\begin{quotation}
   This is a longer quotation.  It consists of two
   paragraphs of text, neither of which are
   particularly interesting.

   This is the second paragraph of the quotation.  It
   is just as dull as the first paragraph.
\end{quotation}
Another frequently-displayed structure is a list.
The following is an example of an \emph{itemized}
list.
\begin{itemize}
   \item This is the first item of an itemized list.
         Each item in the list is marked with a ``tick''.
         You don't have to worry about what kind of tick
         mark is used.

   \item This is the second item of the list.  It
         contains another list nested inside it.  The inner
         list is an \emph{enumerated} list.
         \begin{enumerate}
            \item This is the first item of an enumerated
                  list that is nested within the itemized list.

            \item This is the second item of the inner list.
                  \LaTeX\ allows you to nest lists deeper than
                  you really should.
         \end{enumerate}
         This is the rest of the second item of the outer
         list.  It is no more interesting than any other
         part of the item.
   \item This is the third item of the list.
\end{itemize}
You can even display poetry.
\begin{verse}
   There is an environment
    for verse \\             % The \\ command separates lines
   Whose features some poets % within a stanza.
   will curse.

                             % One or more blank lines separate stanzas.

   For instead of making\\
   Them do \emph{all} line breaking, \\
   It allows them to put too many words on a line when they'd rather be
   forced to be terse.
\end{verse}

Mathematical formulas may also be displayed.  A
displayed formula
is
one-line long; multiline
formulas require special formatting instructions.
   \[  \ip{\Gamma}{\psi'} = x'' + y^{2} + z_{i}^{n}\]
Don't start a paragraph with a displayed equation,
nor make one a paragraph by itself.

\end{document}               % End of document.

\section{Bench Huge Section 420}
% synthetic bench section 420
% This is a sample LaTeX input file.  (Version of 12 August 2004.)
%
% A '%' character causes TeX to ignore all remaining text on the line,
% and is used for comments like this one.

\documentclass{article}      % Specifies the document class

                             % The preamble begins here.
\title{An Example Document}  % Declares the document's title.
\author{Leslie Lamport}      % Declares the author's name.
\date{January 21, 1994}      % Deleting this command produces today's date.

\newcommand{\ip}[2]{(#1, #2)}
                             % Defines \ip{arg1}{arg2} to mean
                             % (arg1, arg2).

%\newcommand{\ip}[2]{\langle #1 | #2\rangle}
                             % This is an alternative definition of
                             % \ip that is commented out.

\begin{document}             % End of preamble and beginning of text.

\maketitle                   % Produces the title.

This is an example input file.  Comparing it with
the output it generates can show you how to
produce a simple document of your own.

\section{Ordinary Text}      % Produces section heading.  Lower-level
                             % sections are begun with similar
                             % \subsection and \subsubsection commands.

The ends  of words and sentences are marked
  by   spaces. It  doesn't matter how many
spaces    you type; one is as good as 100.  The
end of   a line counts as a space.

One   or more   blank lines denote the  end
of  a paragraph.

Since any number of consecutive spaces are treated
like a single one, the formatting of the input
file makes no difference to
      \LaTeX,                % The \LaTeX command generates the LaTeX logo.
but it makes a difference to you.  When you use
\LaTeX, making your input file as easy to read
as possible will be a great help as you write
your document and when you change it.  This sample
file shows how you can add comments to your own input
file.

Because printing is different from typewriting,
there are a number of things that you have to do
differently when preparing an input file than if
you were just typing the document directly.
Quotation marks like
       ``this''
have to be handled specially, as do quotes within
quotes:
       ``\,`this'            % \, separates the double and single quote.
        is what I just
        wrote, not  `that'\,''.

Dashes come in three sizes: an
       intra-word
dash, a medium dash for number ranges like
       1--2,
and a punctuation
       dash---like
this.

A sentence-ending space should be larger than the
space between words within a sentence.  You
sometimes have to type special commands in
conjunction with punctuation characters to get
this right, as in the following sentence.
       Gnats, gnus, etc.\ all  % `\ ' makes an inter-word space.
       begin with G\@.         % \@ marks end-of-sentence punctuation.
You should check the spaces after periods when
reading your output to make sure you haven't
forgotten any special cases.  Generating an
ellipsis
       \ldots\               % `\ ' is needed after `\ldots' because TeX
                             % ignores spaces after command names like \ldots
                             % made from \ + letters.
                             %
                             % Note how a `%' character causes TeX to ignore
                             % the end of the input line, so these blank lines
                             % do not start a new paragraph.
                             %
with the right spacing around the periods requires
a special command.

\LaTeX\ interprets some common characters as
commands, so you must type special commands to
generate them.  These characters include the
following:
       \$ \& \% \# \{ and \}.

In printing, text is usually emphasized with an
       \emph{italic}
type style.

\begin{em}
   A long segment of text can also be emphasized
   in this way.  Text within such a segment can be
   given \emph{additional} emphasis.
\end{em}

It is sometimes necessary to prevent \LaTeX\ from
breaking a line where it might otherwise do so.
This may be at a space, as between the ``Mr.''\ and
``Jones'' in
       ``Mr.~Jones'',        % ~ produces an unbreakable interword space.
or within a word---especially when the word is a
symbol like
       \mbox{\emph{itemnum}}
that makes little sense when hyphenated across
lines.

Footnotes\footnote{This is an example of a footnote.}
pose no problem.

\LaTeX\ is good at typesetting mathematical formulas
like
       \( x-3y + z = 7 \)
or
       \( a_{1} > x^{2n} + y^{2n} > x' \)
or
       \( \ip{A}{B} = \sum_{i} a_{i} b_{i} \).
The spaces you type in a formula are
ignored.  Remember that a letter like
       $x$                   % $ ... $  and  \( ... \)  are equivalent
is a formula when it denotes a mathematical
symbol, and it should be typed as one.

\section{Displayed Text}

Text is displayed by indenting it from the left
margin.  Quotations are commonly displayed.  There
are short quotations
\begin{quote}
   This is a short quotation.  It consists of a
   single paragraph of text.  See how it is formatted.
\end{quote}
and longer ones.
\begin{quotation}
   This is a longer quotation.  It consists of two
   paragraphs of text, neither of which are
   particularly interesting.

   This is the second paragraph of the quotation.  It
   is just as dull as the first paragraph.
\end{quotation}
Another frequently-displayed structure is a list.
The following is an example of an \emph{itemized}
list.
\begin{itemize}
   \item This is the first item of an itemized list.
         Each item in the list is marked with a ``tick''.
         You don't have to worry about what kind of tick
         mark is used.

   \item This is the second item of the list.  It
         contains another list nested inside it.  The inner
         list is an \emph{enumerated} list.
         \begin{enumerate}
            \item This is the first item of an enumerated
                  list that is nested within the itemized list.

            \item This is the second item of the inner list.
                  \LaTeX\ allows you to nest lists deeper than
                  you really should.
         \end{enumerate}
         This is the rest of the second item of the outer
         list.  It is no more interesting than any other
         part of the item.
   \item This is the third item of the list.
\end{itemize}
You can even display poetry.
\begin{verse}
   There is an environment
    for verse \\             % The \\ command separates lines
   Whose features some poets % within a stanza.
   will curse.

                             % One or more blank lines separate stanzas.

   For instead of making\\
   Them do \emph{all} line breaking, \\
   It allows them to put too many words on a line when they'd rather be
   forced to be terse.
\end{verse}

Mathematical formulas may also be displayed.  A
displayed formula
is
one-line long; multiline
formulas require special formatting instructions.
   \[  \ip{\Gamma}{\psi'} = x'' + y^{2} + z_{i}^{n}\]
Don't start a paragraph with a displayed equation,
nor make one a paragraph by itself.

\end{document}               % End of document.

\section{Bench Huge Section 421}
% synthetic bench section 421
% This is a sample LaTeX input file.  (Version of 12 August 2004.)
%
% A '%' character causes TeX to ignore all remaining text on the line,
% and is used for comments like this one.

\documentclass{article}      % Specifies the document class

                             % The preamble begins here.
\title{An Example Document}  % Declares the document's title.
\author{Leslie Lamport}      % Declares the author's name.
\date{January 21, 1994}      % Deleting this command produces today's date.

\newcommand{\ip}[2]{(#1, #2)}
                             % Defines \ip{arg1}{arg2} to mean
                             % (arg1, arg2).

%\newcommand{\ip}[2]{\langle #1 | #2\rangle}
                             % This is an alternative definition of
                             % \ip that is commented out.

\begin{document}             % End of preamble and beginning of text.

\maketitle                   % Produces the title.

This is an example input file.  Comparing it with
the output it generates can show you how to
produce a simple document of your own.

\section{Ordinary Text}      % Produces section heading.  Lower-level
                             % sections are begun with similar
                             % \subsection and \subsubsection commands.

The ends  of words and sentences are marked
  by   spaces. It  doesn't matter how many
spaces    you type; one is as good as 100.  The
end of   a line counts as a space.

One   or more   blank lines denote the  end
of  a paragraph.

Since any number of consecutive spaces are treated
like a single one, the formatting of the input
file makes no difference to
      \LaTeX,                % The \LaTeX command generates the LaTeX logo.
but it makes a difference to you.  When you use
\LaTeX, making your input file as easy to read
as possible will be a great help as you write
your document and when you change it.  This sample
file shows how you can add comments to your own input
file.

Because printing is different from typewriting,
there are a number of things that you have to do
differently when preparing an input file than if
you were just typing the document directly.
Quotation marks like
       ``this''
have to be handled specially, as do quotes within
quotes:
       ``\,`this'            % \, separates the double and single quote.
        is what I just
        wrote, not  `that'\,''.

Dashes come in three sizes: an
       intra-word
dash, a medium dash for number ranges like
       1--2,
and a punctuation
       dash---like
this.

A sentence-ending space should be larger than the
space between words within a sentence.  You
sometimes have to type special commands in
conjunction with punctuation characters to get
this right, as in the following sentence.
       Gnats, gnus, etc.\ all  % `\ ' makes an inter-word space.
       begin with G\@.         % \@ marks end-of-sentence punctuation.
You should check the spaces after periods when
reading your output to make sure you haven't
forgotten any special cases.  Generating an
ellipsis
       \ldots\               % `\ ' is needed after `\ldots' because TeX
                             % ignores spaces after command names like \ldots
                             % made from \ + letters.
                             %
                             % Note how a `%' character causes TeX to ignore
                             % the end of the input line, so these blank lines
                             % do not start a new paragraph.
                             %
with the right spacing around the periods requires
a special command.

\LaTeX\ interprets some common characters as
commands, so you must type special commands to
generate them.  These characters include the
following:
       \$ \& \% \# \{ and \}.

In printing, text is usually emphasized with an
       \emph{italic}
type style.

\begin{em}
   A long segment of text can also be emphasized
   in this way.  Text within such a segment can be
   given \emph{additional} emphasis.
\end{em}

It is sometimes necessary to prevent \LaTeX\ from
breaking a line where it might otherwise do so.
This may be at a space, as between the ``Mr.''\ and
``Jones'' in
       ``Mr.~Jones'',        % ~ produces an unbreakable interword space.
or within a word---especially when the word is a
symbol like
       \mbox{\emph{itemnum}}
that makes little sense when hyphenated across
lines.

Footnotes\footnote{This is an example of a footnote.}
pose no problem.

\LaTeX\ is good at typesetting mathematical formulas
like
       \( x-3y + z = 7 \)
or
       \( a_{1} > x^{2n} + y^{2n} > x' \)
or
       \( \ip{A}{B} = \sum_{i} a_{i} b_{i} \).
The spaces you type in a formula are
ignored.  Remember that a letter like
       $x$                   % $ ... $  and  \( ... \)  are equivalent
is a formula when it denotes a mathematical
symbol, and it should be typed as one.

\section{Displayed Text}

Text is displayed by indenting it from the left
margin.  Quotations are commonly displayed.  There
are short quotations
\begin{quote}
   This is a short quotation.  It consists of a
   single paragraph of text.  See how it is formatted.
\end{quote}
and longer ones.
\begin{quotation}
   This is a longer quotation.  It consists of two
   paragraphs of text, neither of which are
   particularly interesting.

   This is the second paragraph of the quotation.  It
   is just as dull as the first paragraph.
\end{quotation}
Another frequently-displayed structure is a list.
The following is an example of an \emph{itemized}
list.
\begin{itemize}
   \item This is the first item of an itemized list.
         Each item in the list is marked with a ``tick''.
         You don't have to worry about what kind of tick
         mark is used.

   \item This is the second item of the list.  It
         contains another list nested inside it.  The inner
         list is an \emph{enumerated} list.
         \begin{enumerate}
            \item This is the first item of an enumerated
                  list that is nested within the itemized list.

            \item This is the second item of the inner list.
                  \LaTeX\ allows you to nest lists deeper than
                  you really should.
         \end{enumerate}
         This is the rest of the second item of the outer
         list.  It is no more interesting than any other
         part of the item.
   \item This is the third item of the list.
\end{itemize}
You can even display poetry.
\begin{verse}
   There is an environment
    for verse \\             % The \\ command separates lines
   Whose features some poets % within a stanza.
   will curse.

                             % One or more blank lines separate stanzas.

   For instead of making\\
   Them do \emph{all} line breaking, \\
   It allows them to put too many words on a line when they'd rather be
   forced to be terse.
\end{verse}

Mathematical formulas may also be displayed.  A
displayed formula
is
one-line long; multiline
formulas require special formatting instructions.
   \[  \ip{\Gamma}{\psi'} = x'' + y^{2} + z_{i}^{n}\]
Don't start a paragraph with a displayed equation,
nor make one a paragraph by itself.

\end{document}               % End of document.

\section{Bench Huge Section 422}
% synthetic bench section 422
% This is a sample LaTeX input file.  (Version of 12 August 2004.)
%
% A '%' character causes TeX to ignore all remaining text on the line,
% and is used for comments like this one.

\documentclass{article}      % Specifies the document class

                             % The preamble begins here.
\title{An Example Document}  % Declares the document's title.
\author{Leslie Lamport}      % Declares the author's name.
\date{January 21, 1994}      % Deleting this command produces today's date.

\newcommand{\ip}[2]{(#1, #2)}
                             % Defines \ip{arg1}{arg2} to mean
                             % (arg1, arg2).

%\newcommand{\ip}[2]{\langle #1 | #2\rangle}
                             % This is an alternative definition of
                             % \ip that is commented out.

\begin{document}             % End of preamble and beginning of text.

\maketitle                   % Produces the title.

This is an example input file.  Comparing it with
the output it generates can show you how to
produce a simple document of your own.

\section{Ordinary Text}      % Produces section heading.  Lower-level
                             % sections are begun with similar
                             % \subsection and \subsubsection commands.

The ends  of words and sentences are marked
  by   spaces. It  doesn't matter how many
spaces    you type; one is as good as 100.  The
end of   a line counts as a space.

One   or more   blank lines denote the  end
of  a paragraph.

Since any number of consecutive spaces are treated
like a single one, the formatting of the input
file makes no difference to
      \LaTeX,                % The \LaTeX command generates the LaTeX logo.
but it makes a difference to you.  When you use
\LaTeX, making your input file as easy to read
as possible will be a great help as you write
your document and when you change it.  This sample
file shows how you can add comments to your own input
file.

Because printing is different from typewriting,
there are a number of things that you have to do
differently when preparing an input file than if
you were just typing the document directly.
Quotation marks like
       ``this''
have to be handled specially, as do quotes within
quotes:
       ``\,`this'            % \, separates the double and single quote.
        is what I just
        wrote, not  `that'\,''.

Dashes come in three sizes: an
       intra-word
dash, a medium dash for number ranges like
       1--2,
and a punctuation
       dash---like
this.

A sentence-ending space should be larger than the
space between words within a sentence.  You
sometimes have to type special commands in
conjunction with punctuation characters to get
this right, as in the following sentence.
       Gnats, gnus, etc.\ all  % `\ ' makes an inter-word space.
       begin with G\@.         % \@ marks end-of-sentence punctuation.
You should check the spaces after periods when
reading your output to make sure you haven't
forgotten any special cases.  Generating an
ellipsis
       \ldots\               % `\ ' is needed after `\ldots' because TeX
                             % ignores spaces after command names like \ldots
                             % made from \ + letters.
                             %
                             % Note how a `%' character causes TeX to ignore
                             % the end of the input line, so these blank lines
                             % do not start a new paragraph.
                             %
with the right spacing around the periods requires
a special command.

\LaTeX\ interprets some common characters as
commands, so you must type special commands to
generate them.  These characters include the
following:
       \$ \& \% \# \{ and \}.

In printing, text is usually emphasized with an
       \emph{italic}
type style.

\begin{em}
   A long segment of text can also be emphasized
   in this way.  Text within such a segment can be
   given \emph{additional} emphasis.
\end{em}

It is sometimes necessary to prevent \LaTeX\ from
breaking a line where it might otherwise do so.
This may be at a space, as between the ``Mr.''\ and
``Jones'' in
       ``Mr.~Jones'',        % ~ produces an unbreakable interword space.
or within a word---especially when the word is a
symbol like
       \mbox{\emph{itemnum}}
that makes little sense when hyphenated across
lines.

Footnotes\footnote{This is an example of a footnote.}
pose no problem.

\LaTeX\ is good at typesetting mathematical formulas
like
       \( x-3y + z = 7 \)
or
       \( a_{1} > x^{2n} + y^{2n} > x' \)
or
       \( \ip{A}{B} = \sum_{i} a_{i} b_{i} \).
The spaces you type in a formula are
ignored.  Remember that a letter like
       $x$                   % $ ... $  and  \( ... \)  are equivalent
is a formula when it denotes a mathematical
symbol, and it should be typed as one.

\section{Displayed Text}

Text is displayed by indenting it from the left
margin.  Quotations are commonly displayed.  There
are short quotations
\begin{quote}
   This is a short quotation.  It consists of a
   single paragraph of text.  See how it is formatted.
\end{quote}
and longer ones.
\begin{quotation}
   This is a longer quotation.  It consists of two
   paragraphs of text, neither of which are
   particularly interesting.

   This is the second paragraph of the quotation.  It
   is just as dull as the first paragraph.
\end{quotation}
Another frequently-displayed structure is a list.
The following is an example of an \emph{itemized}
list.
\begin{itemize}
   \item This is the first item of an itemized list.
         Each item in the list is marked with a ``tick''.
         You don't have to worry about what kind of tick
         mark is used.

   \item This is the second item of the list.  It
         contains another list nested inside it.  The inner
         list is an \emph{enumerated} list.
         \begin{enumerate}
            \item This is the first item of an enumerated
                  list that is nested within the itemized list.

            \item This is the second item of the inner list.
                  \LaTeX\ allows you to nest lists deeper than
                  you really should.
         \end{enumerate}
         This is the rest of the second item of the outer
         list.  It is no more interesting than any other
         part of the item.
   \item This is the third item of the list.
\end{itemize}
You can even display poetry.
\begin{verse}
   There is an environment
    for verse \\             % The \\ command separates lines
   Whose features some poets % within a stanza.
   will curse.

                             % One or more blank lines separate stanzas.

   For instead of making\\
   Them do \emph{all} line breaking, \\
   It allows them to put too many words on a line when they'd rather be
   forced to be terse.
\end{verse}

Mathematical formulas may also be displayed.  A
displayed formula
is
one-line long; multiline
formulas require special formatting instructions.
   \[  \ip{\Gamma}{\psi'} = x'' + y^{2} + z_{i}^{n}\]
Don't start a paragraph with a displayed equation,
nor make one a paragraph by itself.

\end{document}               % End of document.

\section{Bench Huge Section 423}
% synthetic bench section 423
% This is a sample LaTeX input file.  (Version of 12 August 2004.)
%
% A '%' character causes TeX to ignore all remaining text on the line,
% and is used for comments like this one.

\documentclass{article}      % Specifies the document class

                             % The preamble begins here.
\title{An Example Document}  % Declares the document's title.
\author{Leslie Lamport}      % Declares the author's name.
\date{January 21, 1994}      % Deleting this command produces today's date.

\newcommand{\ip}[2]{(#1, #2)}
                             % Defines \ip{arg1}{arg2} to mean
                             % (arg1, arg2).

%\newcommand{\ip}[2]{\langle #1 | #2\rangle}
                             % This is an alternative definition of
                             % \ip that is commented out.

\begin{document}             % End of preamble and beginning of text.

\maketitle                   % Produces the title.

This is an example input file.  Comparing it with
the output it generates can show you how to
produce a simple document of your own.

\section{Ordinary Text}      % Produces section heading.  Lower-level
                             % sections are begun with similar
                             % \subsection and \subsubsection commands.

The ends  of words and sentences are marked
  by   spaces. It  doesn't matter how many
spaces    you type; one is as good as 100.  The
end of   a line counts as a space.

One   or more   blank lines denote the  end
of  a paragraph.

Since any number of consecutive spaces are treated
like a single one, the formatting of the input
file makes no difference to
      \LaTeX,                % The \LaTeX command generates the LaTeX logo.
but it makes a difference to you.  When you use
\LaTeX, making your input file as easy to read
as possible will be a great help as you write
your document and when you change it.  This sample
file shows how you can add comments to your own input
file.

Because printing is different from typewriting,
there are a number of things that you have to do
differently when preparing an input file than if
you were just typing the document directly.
Quotation marks like
       ``this''
have to be handled specially, as do quotes within
quotes:
       ``\,`this'            % \, separates the double and single quote.
        is what I just
        wrote, not  `that'\,''.

Dashes come in three sizes: an
       intra-word
dash, a medium dash for number ranges like
       1--2,
and a punctuation
       dash---like
this.

A sentence-ending space should be larger than the
space between words within a sentence.  You
sometimes have to type special commands in
conjunction with punctuation characters to get
this right, as in the following sentence.
       Gnats, gnus, etc.\ all  % `\ ' makes an inter-word space.
       begin with G\@.         % \@ marks end-of-sentence punctuation.
You should check the spaces after periods when
reading your output to make sure you haven't
forgotten any special cases.  Generating an
ellipsis
       \ldots\               % `\ ' is needed after `\ldots' because TeX
                             % ignores spaces after command names like \ldots
                             % made from \ + letters.
                             %
                             % Note how a `%' character causes TeX to ignore
                             % the end of the input line, so these blank lines
                             % do not start a new paragraph.
                             %
with the right spacing around the periods requires
a special command.

\LaTeX\ interprets some common characters as
commands, so you must type special commands to
generate them.  These characters include the
following:
       \$ \& \% \# \{ and \}.

In printing, text is usually emphasized with an
       \emph{italic}
type style.

\begin{em}
   A long segment of text can also be emphasized
   in this way.  Text within such a segment can be
   given \emph{additional} emphasis.
\end{em}

It is sometimes necessary to prevent \LaTeX\ from
breaking a line where it might otherwise do so.
This may be at a space, as between the ``Mr.''\ and
``Jones'' in
       ``Mr.~Jones'',        % ~ produces an unbreakable interword space.
or within a word---especially when the word is a
symbol like
       \mbox{\emph{itemnum}}
that makes little sense when hyphenated across
lines.

Footnotes\footnote{This is an example of a footnote.}
pose no problem.

\LaTeX\ is good at typesetting mathematical formulas
like
       \( x-3y + z = 7 \)
or
       \( a_{1} > x^{2n} + y^{2n} > x' \)
or
       \( \ip{A}{B} = \sum_{i} a_{i} b_{i} \).
The spaces you type in a formula are
ignored.  Remember that a letter like
       $x$                   % $ ... $  and  \( ... \)  are equivalent
is a formula when it denotes a mathematical
symbol, and it should be typed as one.

\section{Displayed Text}

Text is displayed by indenting it from the left
margin.  Quotations are commonly displayed.  There
are short quotations
\begin{quote}
   This is a short quotation.  It consists of a
   single paragraph of text.  See how it is formatted.
\end{quote}
and longer ones.
\begin{quotation}
   This is a longer quotation.  It consists of two
   paragraphs of text, neither of which are
   particularly interesting.

   This is the second paragraph of the quotation.  It
   is just as dull as the first paragraph.
\end{quotation}
Another frequently-displayed structure is a list.
The following is an example of an \emph{itemized}
list.
\begin{itemize}
   \item This is the first item of an itemized list.
         Each item in the list is marked with a ``tick''.
         You don't have to worry about what kind of tick
         mark is used.

   \item This is the second item of the list.  It
         contains another list nested inside it.  The inner
         list is an \emph{enumerated} list.
         \begin{enumerate}
            \item This is the first item of an enumerated
                  list that is nested within the itemized list.

            \item This is the second item of the inner list.
                  \LaTeX\ allows you to nest lists deeper than
                  you really should.
         \end{enumerate}
         This is the rest of the second item of the outer
         list.  It is no more interesting than any other
         part of the item.
   \item This is the third item of the list.
\end{itemize}
You can even display poetry.
\begin{verse}
   There is an environment
    for verse \\             % The \\ command separates lines
   Whose features some poets % within a stanza.
   will curse.

                             % One or more blank lines separate stanzas.

   For instead of making\\
   Them do \emph{all} line breaking, \\
   It allows them to put too many words on a line when they'd rather be
   forced to be terse.
\end{verse}

Mathematical formulas may also be displayed.  A
displayed formula
is
one-line long; multiline
formulas require special formatting instructions.
   \[  \ip{\Gamma}{\psi'} = x'' + y^{2} + z_{i}^{n}\]
Don't start a paragraph with a displayed equation,
nor make one a paragraph by itself.

\end{document}               % End of document.

\section{Bench Huge Section 424}
% synthetic bench section 424
% This is a sample LaTeX input file.  (Version of 12 August 2004.)
%
% A '%' character causes TeX to ignore all remaining text on the line,
% and is used for comments like this one.

\documentclass{article}      % Specifies the document class

                             % The preamble begins here.
\title{An Example Document}  % Declares the document's title.
\author{Leslie Lamport}      % Declares the author's name.
\date{January 21, 1994}      % Deleting this command produces today's date.

\newcommand{\ip}[2]{(#1, #2)}
                             % Defines \ip{arg1}{arg2} to mean
                             % (arg1, arg2).

%\newcommand{\ip}[2]{\langle #1 | #2\rangle}
                             % This is an alternative definition of
                             % \ip that is commented out.

\begin{document}             % End of preamble and beginning of text.

\maketitle                   % Produces the title.

This is an example input file.  Comparing it with
the output it generates can show you how to
produce a simple document of your own.

\section{Ordinary Text}      % Produces section heading.  Lower-level
                             % sections are begun with similar
                             % \subsection and \subsubsection commands.

The ends  of words and sentences are marked
  by   spaces. It  doesn't matter how many
spaces    you type; one is as good as 100.  The
end of   a line counts as a space.

One   or more   blank lines denote the  end
of  a paragraph.

Since any number of consecutive spaces are treated
like a single one, the formatting of the input
file makes no difference to
      \LaTeX,                % The \LaTeX command generates the LaTeX logo.
but it makes a difference to you.  When you use
\LaTeX, making your input file as easy to read
as possible will be a great help as you write
your document and when you change it.  This sample
file shows how you can add comments to your own input
file.

Because printing is different from typewriting,
there are a number of things that you have to do
differently when preparing an input file than if
you were just typing the document directly.
Quotation marks like
       ``this''
have to be handled specially, as do quotes within
quotes:
       ``\,`this'            % \, separates the double and single quote.
        is what I just
        wrote, not  `that'\,''.

Dashes come in three sizes: an
       intra-word
dash, a medium dash for number ranges like
       1--2,
and a punctuation
       dash---like
this.

A sentence-ending space should be larger than the
space between words within a sentence.  You
sometimes have to type special commands in
conjunction with punctuation characters to get
this right, as in the following sentence.
       Gnats, gnus, etc.\ all  % `\ ' makes an inter-word space.
       begin with G\@.         % \@ marks end-of-sentence punctuation.
You should check the spaces after periods when
reading your output to make sure you haven't
forgotten any special cases.  Generating an
ellipsis
       \ldots\               % `\ ' is needed after `\ldots' because TeX
                             % ignores spaces after command names like \ldots
                             % made from \ + letters.
                             %
                             % Note how a `%' character causes TeX to ignore
                             % the end of the input line, so these blank lines
                             % do not start a new paragraph.
                             %
with the right spacing around the periods requires
a special command.

\LaTeX\ interprets some common characters as
commands, so you must type special commands to
generate them.  These characters include the
following:
       \$ \& \% \# \{ and \}.

In printing, text is usually emphasized with an
       \emph{italic}
type style.

\begin{em}
   A long segment of text can also be emphasized
   in this way.  Text within such a segment can be
   given \emph{additional} emphasis.
\end{em}

It is sometimes necessary to prevent \LaTeX\ from
breaking a line where it might otherwise do so.
This may be at a space, as between the ``Mr.''\ and
``Jones'' in
       ``Mr.~Jones'',        % ~ produces an unbreakable interword space.
or within a word---especially when the word is a
symbol like
       \mbox{\emph{itemnum}}
that makes little sense when hyphenated across
lines.

Footnotes\footnote{This is an example of a footnote.}
pose no problem.

\LaTeX\ is good at typesetting mathematical formulas
like
       \( x-3y + z = 7 \)
or
       \( a_{1} > x^{2n} + y^{2n} > x' \)
or
       \( \ip{A}{B} = \sum_{i} a_{i} b_{i} \).
The spaces you type in a formula are
ignored.  Remember that a letter like
       $x$                   % $ ... $  and  \( ... \)  are equivalent
is a formula when it denotes a mathematical
symbol, and it should be typed as one.

\section{Displayed Text}

Text is displayed by indenting it from the left
margin.  Quotations are commonly displayed.  There
are short quotations
\begin{quote}
   This is a short quotation.  It consists of a
   single paragraph of text.  See how it is formatted.
\end{quote}
and longer ones.
\begin{quotation}
   This is a longer quotation.  It consists of two
   paragraphs of text, neither of which are
   particularly interesting.

   This is the second paragraph of the quotation.  It
   is just as dull as the first paragraph.
\end{quotation}
Another frequently-displayed structure is a list.
The following is an example of an \emph{itemized}
list.
\begin{itemize}
   \item This is the first item of an itemized list.
         Each item in the list is marked with a ``tick''.
         You don't have to worry about what kind of tick
         mark is used.

   \item This is the second item of the list.  It
         contains another list nested inside it.  The inner
         list is an \emph{enumerated} list.
         \begin{enumerate}
            \item This is the first item of an enumerated
                  list that is nested within the itemized list.

            \item This is the second item of the inner list.
                  \LaTeX\ allows you to nest lists deeper than
                  you really should.
         \end{enumerate}
         This is the rest of the second item of the outer
         list.  It is no more interesting than any other
         part of the item.
   \item This is the third item of the list.
\end{itemize}
You can even display poetry.
\begin{verse}
   There is an environment
    for verse \\             % The \\ command separates lines
   Whose features some poets % within a stanza.
   will curse.

                             % One or more blank lines separate stanzas.

   For instead of making\\
   Them do \emph{all} line breaking, \\
   It allows them to put too many words on a line when they'd rather be
   forced to be terse.
\end{verse}

Mathematical formulas may also be displayed.  A
displayed formula
is
one-line long; multiline
formulas require special formatting instructions.
   \[  \ip{\Gamma}{\psi'} = x'' + y^{2} + z_{i}^{n}\]
Don't start a paragraph with a displayed equation,
nor make one a paragraph by itself.

\end{document}               % End of document.

\section{Bench Huge Section 425}
% synthetic bench section 425
% This is a sample LaTeX input file.  (Version of 12 August 2004.)
%
% A '%' character causes TeX to ignore all remaining text on the line,
% and is used for comments like this one.

\documentclass{article}      % Specifies the document class

                             % The preamble begins here.
\title{An Example Document}  % Declares the document's title.
\author{Leslie Lamport}      % Declares the author's name.
\date{January 21, 1994}      % Deleting this command produces today's date.

\newcommand{\ip}[2]{(#1, #2)}
                             % Defines \ip{arg1}{arg2} to mean
                             % (arg1, arg2).

%\newcommand{\ip}[2]{\langle #1 | #2\rangle}
                             % This is an alternative definition of
                             % \ip that is commented out.

\begin{document}             % End of preamble and beginning of text.

\maketitle                   % Produces the title.

This is an example input file.  Comparing it with
the output it generates can show you how to
produce a simple document of your own.

\section{Ordinary Text}      % Produces section heading.  Lower-level
                             % sections are begun with similar
                             % \subsection and \subsubsection commands.

The ends  of words and sentences are marked
  by   spaces. It  doesn't matter how many
spaces    you type; one is as good as 100.  The
end of   a line counts as a space.

One   or more   blank lines denote the  end
of  a paragraph.

Since any number of consecutive spaces are treated
like a single one, the formatting of the input
file makes no difference to
      \LaTeX,                % The \LaTeX command generates the LaTeX logo.
but it makes a difference to you.  When you use
\LaTeX, making your input file as easy to read
as possible will be a great help as you write
your document and when you change it.  This sample
file shows how you can add comments to your own input
file.

Because printing is different from typewriting,
there are a number of things that you have to do
differently when preparing an input file than if
you were just typing the document directly.
Quotation marks like
       ``this''
have to be handled specially, as do quotes within
quotes:
       ``\,`this'            % \, separates the double and single quote.
        is what I just
        wrote, not  `that'\,''.

Dashes come in three sizes: an
       intra-word
dash, a medium dash for number ranges like
       1--2,
and a punctuation
       dash---like
this.

A sentence-ending space should be larger than the
space between words within a sentence.  You
sometimes have to type special commands in
conjunction with punctuation characters to get
this right, as in the following sentence.
       Gnats, gnus, etc.\ all  % `\ ' makes an inter-word space.
       begin with G\@.         % \@ marks end-of-sentence punctuation.
You should check the spaces after periods when
reading your output to make sure you haven't
forgotten any special cases.  Generating an
ellipsis
       \ldots\               % `\ ' is needed after `\ldots' because TeX
                             % ignores spaces after command names like \ldots
                             % made from \ + letters.
                             %
                             % Note how a `%' character causes TeX to ignore
                             % the end of the input line, so these blank lines
                             % do not start a new paragraph.
                             %
with the right spacing around the periods requires
a special command.

\LaTeX\ interprets some common characters as
commands, so you must type special commands to
generate them.  These characters include the
following:
       \$ \& \% \# \{ and \}.

In printing, text is usually emphasized with an
       \emph{italic}
type style.

\begin{em}
   A long segment of text can also be emphasized
   in this way.  Text within such a segment can be
   given \emph{additional} emphasis.
\end{em}

It is sometimes necessary to prevent \LaTeX\ from
breaking a line where it might otherwise do so.
This may be at a space, as between the ``Mr.''\ and
``Jones'' in
       ``Mr.~Jones'',        % ~ produces an unbreakable interword space.
or within a word---especially when the word is a
symbol like
       \mbox{\emph{itemnum}}
that makes little sense when hyphenated across
lines.

Footnotes\footnote{This is an example of a footnote.}
pose no problem.

\LaTeX\ is good at typesetting mathematical formulas
like
       \( x-3y + z = 7 \)
or
       \( a_{1} > x^{2n} + y^{2n} > x' \)
or
       \( \ip{A}{B} = \sum_{i} a_{i} b_{i} \).
The spaces you type in a formula are
ignored.  Remember that a letter like
       $x$                   % $ ... $  and  \( ... \)  are equivalent
is a formula when it denotes a mathematical
symbol, and it should be typed as one.

\section{Displayed Text}

Text is displayed by indenting it from the left
margin.  Quotations are commonly displayed.  There
are short quotations
\begin{quote}
   This is a short quotation.  It consists of a
   single paragraph of text.  See how it is formatted.
\end{quote}
and longer ones.
\begin{quotation}
   This is a longer quotation.  It consists of two
   paragraphs of text, neither of which are
   particularly interesting.

   This is the second paragraph of the quotation.  It
   is just as dull as the first paragraph.
\end{quotation}
Another frequently-displayed structure is a list.
The following is an example of an \emph{itemized}
list.
\begin{itemize}
   \item This is the first item of an itemized list.
         Each item in the list is marked with a ``tick''.
         You don't have to worry about what kind of tick
         mark is used.

   \item This is the second item of the list.  It
         contains another list nested inside it.  The inner
         list is an \emph{enumerated} list.
         \begin{enumerate}
            \item This is the first item of an enumerated
                  list that is nested within the itemized list.

            \item This is the second item of the inner list.
                  \LaTeX\ allows you to nest lists deeper than
                  you really should.
         \end{enumerate}
         This is the rest of the second item of the outer
         list.  It is no more interesting than any other
         part of the item.
   \item This is the third item of the list.
\end{itemize}
You can even display poetry.
\begin{verse}
   There is an environment
    for verse \\             % The \\ command separates lines
   Whose features some poets % within a stanza.
   will curse.

                             % One or more blank lines separate stanzas.

   For instead of making\\
   Them do \emph{all} line breaking, \\
   It allows them to put too many words on a line when they'd rather be
   forced to be terse.
\end{verse}

Mathematical formulas may also be displayed.  A
displayed formula
is
one-line long; multiline
formulas require special formatting instructions.
   \[  \ip{\Gamma}{\psi'} = x'' + y^{2} + z_{i}^{n}\]
Don't start a paragraph with a displayed equation,
nor make one a paragraph by itself.

\end{document}               % End of document.

\section{Bench Huge Section 426}
% synthetic bench section 426
% This is a sample LaTeX input file.  (Version of 12 August 2004.)
%
% A '%' character causes TeX to ignore all remaining text on the line,
% and is used for comments like this one.

\documentclass{article}      % Specifies the document class

                             % The preamble begins here.
\title{An Example Document}  % Declares the document's title.
\author{Leslie Lamport}      % Declares the author's name.
\date{January 21, 1994}      % Deleting this command produces today's date.

\newcommand{\ip}[2]{(#1, #2)}
                             % Defines \ip{arg1}{arg2} to mean
                             % (arg1, arg2).

%\newcommand{\ip}[2]{\langle #1 | #2\rangle}
                             % This is an alternative definition of
                             % \ip that is commented out.

\begin{document}             % End of preamble and beginning of text.

\maketitle                   % Produces the title.

This is an example input file.  Comparing it with
the output it generates can show you how to
produce a simple document of your own.

\section{Ordinary Text}      % Produces section heading.  Lower-level
                             % sections are begun with similar
                             % \subsection and \subsubsection commands.

The ends  of words and sentences are marked
  by   spaces. It  doesn't matter how many
spaces    you type; one is as good as 100.  The
end of   a line counts as a space.

One   or more   blank lines denote the  end
of  a paragraph.

Since any number of consecutive spaces are treated
like a single one, the formatting of the input
file makes no difference to
      \LaTeX,                % The \LaTeX command generates the LaTeX logo.
but it makes a difference to you.  When you use
\LaTeX, making your input file as easy to read
as possible will be a great help as you write
your document and when you change it.  This sample
file shows how you can add comments to your own input
file.

Because printing is different from typewriting,
there are a number of things that you have to do
differently when preparing an input file than if
you were just typing the document directly.
Quotation marks like
       ``this''
have to be handled specially, as do quotes within
quotes:
       ``\,`this'            % \, separates the double and single quote.
        is what I just
        wrote, not  `that'\,''.

Dashes come in three sizes: an
       intra-word
dash, a medium dash for number ranges like
       1--2,
and a punctuation
       dash---like
this.

A sentence-ending space should be larger than the
space between words within a sentence.  You
sometimes have to type special commands in
conjunction with punctuation characters to get
this right, as in the following sentence.
       Gnats, gnus, etc.\ all  % `\ ' makes an inter-word space.
       begin with G\@.         % \@ marks end-of-sentence punctuation.
You should check the spaces after periods when
reading your output to make sure you haven't
forgotten any special cases.  Generating an
ellipsis
       \ldots\               % `\ ' is needed after `\ldots' because TeX
                             % ignores spaces after command names like \ldots
                             % made from \ + letters.
                             %
                             % Note how a `%' character causes TeX to ignore
                             % the end of the input line, so these blank lines
                             % do not start a new paragraph.
                             %
with the right spacing around the periods requires
a special command.

\LaTeX\ interprets some common characters as
commands, so you must type special commands to
generate them.  These characters include the
following:
       \$ \& \% \# \{ and \}.

In printing, text is usually emphasized with an
       \emph{italic}
type style.

\begin{em}
   A long segment of text can also be emphasized
   in this way.  Text within such a segment can be
   given \emph{additional} emphasis.
\end{em}

It is sometimes necessary to prevent \LaTeX\ from
breaking a line where it might otherwise do so.
This may be at a space, as between the ``Mr.''\ and
``Jones'' in
       ``Mr.~Jones'',        % ~ produces an unbreakable interword space.
or within a word---especially when the word is a
symbol like
       \mbox{\emph{itemnum}}
that makes little sense when hyphenated across
lines.

Footnotes\footnote{This is an example of a footnote.}
pose no problem.

\LaTeX\ is good at typesetting mathematical formulas
like
       \( x-3y + z = 7 \)
or
       \( a_{1} > x^{2n} + y^{2n} > x' \)
or
       \( \ip{A}{B} = \sum_{i} a_{i} b_{i} \).
The spaces you type in a formula are
ignored.  Remember that a letter like
       $x$                   % $ ... $  and  \( ... \)  are equivalent
is a formula when it denotes a mathematical
symbol, and it should be typed as one.

\section{Displayed Text}

Text is displayed by indenting it from the left
margin.  Quotations are commonly displayed.  There
are short quotations
\begin{quote}
   This is a short quotation.  It consists of a
   single paragraph of text.  See how it is formatted.
\end{quote}
and longer ones.
\begin{quotation}
   This is a longer quotation.  It consists of two
   paragraphs of text, neither of which are
   particularly interesting.

   This is the second paragraph of the quotation.  It
   is just as dull as the first paragraph.
\end{quotation}
Another frequently-displayed structure is a list.
The following is an example of an \emph{itemized}
list.
\begin{itemize}
   \item This is the first item of an itemized list.
         Each item in the list is marked with a ``tick''.
         You don't have to worry about what kind of tick
         mark is used.

   \item This is the second item of the list.  It
         contains another list nested inside it.  The inner
         list is an \emph{enumerated} list.
         \begin{enumerate}
            \item This is the first item of an enumerated
                  list that is nested within the itemized list.

            \item This is the second item of the inner list.
                  \LaTeX\ allows you to nest lists deeper than
                  you really should.
         \end{enumerate}
         This is the rest of the second item of the outer
         list.  It is no more interesting than any other
         part of the item.
   \item This is the third item of the list.
\end{itemize}
You can even display poetry.
\begin{verse}
   There is an environment
    for verse \\             % The \\ command separates lines
   Whose features some poets % within a stanza.
   will curse.

                             % One or more blank lines separate stanzas.

   For instead of making\\
   Them do \emph{all} line breaking, \\
   It allows them to put too many words on a line when they'd rather be
   forced to be terse.
\end{verse}

Mathematical formulas may also be displayed.  A
displayed formula
is
one-line long; multiline
formulas require special formatting instructions.
   \[  \ip{\Gamma}{\psi'} = x'' + y^{2} + z_{i}^{n}\]
Don't start a paragraph with a displayed equation,
nor make one a paragraph by itself.

\end{document}               % End of document.

\section{Bench Huge Section 427}
% synthetic bench section 427
% This is a sample LaTeX input file.  (Version of 12 August 2004.)
%
% A '%' character causes TeX to ignore all remaining text on the line,
% and is used for comments like this one.

\documentclass{article}      % Specifies the document class

                             % The preamble begins here.
\title{An Example Document}  % Declares the document's title.
\author{Leslie Lamport}      % Declares the author's name.
\date{January 21, 1994}      % Deleting this command produces today's date.

\newcommand{\ip}[2]{(#1, #2)}
                             % Defines \ip{arg1}{arg2} to mean
                             % (arg1, arg2).

%\newcommand{\ip}[2]{\langle #1 | #2\rangle}
                             % This is an alternative definition of
                             % \ip that is commented out.

\begin{document}             % End of preamble and beginning of text.

\maketitle                   % Produces the title.

This is an example input file.  Comparing it with
the output it generates can show you how to
produce a simple document of your own.

\section{Ordinary Text}      % Produces section heading.  Lower-level
                             % sections are begun with similar
                             % \subsection and \subsubsection commands.

The ends  of words and sentences are marked
  by   spaces. It  doesn't matter how many
spaces    you type; one is as good as 100.  The
end of   a line counts as a space.

One   or more   blank lines denote the  end
of  a paragraph.

Since any number of consecutive spaces are treated
like a single one, the formatting of the input
file makes no difference to
      \LaTeX,                % The \LaTeX command generates the LaTeX logo.
but it makes a difference to you.  When you use
\LaTeX, making your input file as easy to read
as possible will be a great help as you write
your document and when you change it.  This sample
file shows how you can add comments to your own input
file.

Because printing is different from typewriting,
there are a number of things that you have to do
differently when preparing an input file than if
you were just typing the document directly.
Quotation marks like
       ``this''
have to be handled specially, as do quotes within
quotes:
       ``\,`this'            % \, separates the double and single quote.
        is what I just
        wrote, not  `that'\,''.

Dashes come in three sizes: an
       intra-word
dash, a medium dash for number ranges like
       1--2,
and a punctuation
       dash---like
this.

A sentence-ending space should be larger than the
space between words within a sentence.  You
sometimes have to type special commands in
conjunction with punctuation characters to get
this right, as in the following sentence.
       Gnats, gnus, etc.\ all  % `\ ' makes an inter-word space.
       begin with G\@.         % \@ marks end-of-sentence punctuation.
You should check the spaces after periods when
reading your output to make sure you haven't
forgotten any special cases.  Generating an
ellipsis
       \ldots\               % `\ ' is needed after `\ldots' because TeX
                             % ignores spaces after command names like \ldots
                             % made from \ + letters.
                             %
                             % Note how a `%' character causes TeX to ignore
                             % the end of the input line, so these blank lines
                             % do not start a new paragraph.
                             %
with the right spacing around the periods requires
a special command.

\LaTeX\ interprets some common characters as
commands, so you must type special commands to
generate them.  These characters include the
following:
       \$ \& \% \# \{ and \}.

In printing, text is usually emphasized with an
       \emph{italic}
type style.

\begin{em}
   A long segment of text can also be emphasized
   in this way.  Text within such a segment can be
   given \emph{additional} emphasis.
\end{em}

It is sometimes necessary to prevent \LaTeX\ from
breaking a line where it might otherwise do so.
This may be at a space, as between the ``Mr.''\ and
``Jones'' in
       ``Mr.~Jones'',        % ~ produces an unbreakable interword space.
or within a word---especially when the word is a
symbol like
       \mbox{\emph{itemnum}}
that makes little sense when hyphenated across
lines.

Footnotes\footnote{This is an example of a footnote.}
pose no problem.

\LaTeX\ is good at typesetting mathematical formulas
like
       \( x-3y + z = 7 \)
or
       \( a_{1} > x^{2n} + y^{2n} > x' \)
or
       \( \ip{A}{B} = \sum_{i} a_{i} b_{i} \).
The spaces you type in a formula are
ignored.  Remember that a letter like
       $x$                   % $ ... $  and  \( ... \)  are equivalent
is a formula when it denotes a mathematical
symbol, and it should be typed as one.

\section{Displayed Text}

Text is displayed by indenting it from the left
margin.  Quotations are commonly displayed.  There
are short quotations
\begin{quote}
   This is a short quotation.  It consists of a
   single paragraph of text.  See how it is formatted.
\end{quote}
and longer ones.
\begin{quotation}
   This is a longer quotation.  It consists of two
   paragraphs of text, neither of which are
   particularly interesting.

   This is the second paragraph of the quotation.  It
   is just as dull as the first paragraph.
\end{quotation}
Another frequently-displayed structure is a list.
The following is an example of an \emph{itemized}
list.
\begin{itemize}
   \item This is the first item of an itemized list.
         Each item in the list is marked with a ``tick''.
         You don't have to worry about what kind of tick
         mark is used.

   \item This is the second item of the list.  It
         contains another list nested inside it.  The inner
         list is an \emph{enumerated} list.
         \begin{enumerate}
            \item This is the first item of an enumerated
                  list that is nested within the itemized list.

            \item This is the second item of the inner list.
                  \LaTeX\ allows you to nest lists deeper than
                  you really should.
         \end{enumerate}
         This is the rest of the second item of the outer
         list.  It is no more interesting than any other
         part of the item.
   \item This is the third item of the list.
\end{itemize}
You can even display poetry.
\begin{verse}
   There is an environment
    for verse \\             % The \\ command separates lines
   Whose features some poets % within a stanza.
   will curse.

                             % One or more blank lines separate stanzas.

   For instead of making\\
   Them do \emph{all} line breaking, \\
   It allows them to put too many words on a line when they'd rather be
   forced to be terse.
\end{verse}

Mathematical formulas may also be displayed.  A
displayed formula
is
one-line long; multiline
formulas require special formatting instructions.
   \[  \ip{\Gamma}{\psi'} = x'' + y^{2} + z_{i}^{n}\]
Don't start a paragraph with a displayed equation,
nor make one a paragraph by itself.

\end{document}               % End of document.

\section{Bench Huge Section 428}
% synthetic bench section 428
% This is a sample LaTeX input file.  (Version of 12 August 2004.)
%
% A '%' character causes TeX to ignore all remaining text on the line,
% and is used for comments like this one.

\documentclass{article}      % Specifies the document class

                             % The preamble begins here.
\title{An Example Document}  % Declares the document's title.
\author{Leslie Lamport}      % Declares the author's name.
\date{January 21, 1994}      % Deleting this command produces today's date.

\newcommand{\ip}[2]{(#1, #2)}
                             % Defines \ip{arg1}{arg2} to mean
                             % (arg1, arg2).

%\newcommand{\ip}[2]{\langle #1 | #2\rangle}
                             % This is an alternative definition of
                             % \ip that is commented out.

\begin{document}             % End of preamble and beginning of text.

\maketitle                   % Produces the title.

This is an example input file.  Comparing it with
the output it generates can show you how to
produce a simple document of your own.

\section{Ordinary Text}      % Produces section heading.  Lower-level
                             % sections are begun with similar
                             % \subsection and \subsubsection commands.

The ends  of words and sentences are marked
  by   spaces. It  doesn't matter how many
spaces    you type; one is as good as 100.  The
end of   a line counts as a space.

One   or more   blank lines denote the  end
of  a paragraph.

Since any number of consecutive spaces are treated
like a single one, the formatting of the input
file makes no difference to
      \LaTeX,                % The \LaTeX command generates the LaTeX logo.
but it makes a difference to you.  When you use
\LaTeX, making your input file as easy to read
as possible will be a great help as you write
your document and when you change it.  This sample
file shows how you can add comments to your own input
file.

Because printing is different from typewriting,
there are a number of things that you have to do
differently when preparing an input file than if
you were just typing the document directly.
Quotation marks like
       ``this''
have to be handled specially, as do quotes within
quotes:
       ``\,`this'            % \, separates the double and single quote.
        is what I just
        wrote, not  `that'\,''.

Dashes come in three sizes: an
       intra-word
dash, a medium dash for number ranges like
       1--2,
and a punctuation
       dash---like
this.

A sentence-ending space should be larger than the
space between words within a sentence.  You
sometimes have to type special commands in
conjunction with punctuation characters to get
this right, as in the following sentence.
       Gnats, gnus, etc.\ all  % `\ ' makes an inter-word space.
       begin with G\@.         % \@ marks end-of-sentence punctuation.
You should check the spaces after periods when
reading your output to make sure you haven't
forgotten any special cases.  Generating an
ellipsis
       \ldots\               % `\ ' is needed after `\ldots' because TeX
                             % ignores spaces after command names like \ldots
                             % made from \ + letters.
                             %
                             % Note how a `%' character causes TeX to ignore
                             % the end of the input line, so these blank lines
                             % do not start a new paragraph.
                             %
with the right spacing around the periods requires
a special command.

\LaTeX\ interprets some common characters as
commands, so you must type special commands to
generate them.  These characters include the
following:
       \$ \& \% \# \{ and \}.

In printing, text is usually emphasized with an
       \emph{italic}
type style.

\begin{em}
   A long segment of text can also be emphasized
   in this way.  Text within such a segment can be
   given \emph{additional} emphasis.
\end{em}

It is sometimes necessary to prevent \LaTeX\ from
breaking a line where it might otherwise do so.
This may be at a space, as between the ``Mr.''\ and
``Jones'' in
       ``Mr.~Jones'',        % ~ produces an unbreakable interword space.
or within a word---especially when the word is a
symbol like
       \mbox{\emph{itemnum}}
that makes little sense when hyphenated across
lines.

Footnotes\footnote{This is an example of a footnote.}
pose no problem.

\LaTeX\ is good at typesetting mathematical formulas
like
       \( x-3y + z = 7 \)
or
       \( a_{1} > x^{2n} + y^{2n} > x' \)
or
       \( \ip{A}{B} = \sum_{i} a_{i} b_{i} \).
The spaces you type in a formula are
ignored.  Remember that a letter like
       $x$                   % $ ... $  and  \( ... \)  are equivalent
is a formula when it denotes a mathematical
symbol, and it should be typed as one.

\section{Displayed Text}

Text is displayed by indenting it from the left
margin.  Quotations are commonly displayed.  There
are short quotations
\begin{quote}
   This is a short quotation.  It consists of a
   single paragraph of text.  See how it is formatted.
\end{quote}
and longer ones.
\begin{quotation}
   This is a longer quotation.  It consists of two
   paragraphs of text, neither of which are
   particularly interesting.

   This is the second paragraph of the quotation.  It
   is just as dull as the first paragraph.
\end{quotation}
Another frequently-displayed structure is a list.
The following is an example of an \emph{itemized}
list.
\begin{itemize}
   \item This is the first item of an itemized list.
         Each item in the list is marked with a ``tick''.
         You don't have to worry about what kind of tick
         mark is used.

   \item This is the second item of the list.  It
         contains another list nested inside it.  The inner
         list is an \emph{enumerated} list.
         \begin{enumerate}
            \item This is the first item of an enumerated
                  list that is nested within the itemized list.

            \item This is the second item of the inner list.
                  \LaTeX\ allows you to nest lists deeper than
                  you really should.
         \end{enumerate}
         This is the rest of the second item of the outer
         list.  It is no more interesting than any other
         part of the item.
   \item This is the third item of the list.
\end{itemize}
You can even display poetry.
\begin{verse}
   There is an environment
    for verse \\             % The \\ command separates lines
   Whose features some poets % within a stanza.
   will curse.

                             % One or more blank lines separate stanzas.

   For instead of making\\
   Them do \emph{all} line breaking, \\
   It allows them to put too many words on a line when they'd rather be
   forced to be terse.
\end{verse}

Mathematical formulas may also be displayed.  A
displayed formula
is
one-line long; multiline
formulas require special formatting instructions.
   \[  \ip{\Gamma}{\psi'} = x'' + y^{2} + z_{i}^{n}\]
Don't start a paragraph with a displayed equation,
nor make one a paragraph by itself.

\end{document}               % End of document.

\section{Bench Huge Section 429}
% synthetic bench section 429
% This is a sample LaTeX input file.  (Version of 12 August 2004.)
%
% A '%' character causes TeX to ignore all remaining text on the line,
% and is used for comments like this one.

\documentclass{article}      % Specifies the document class

                             % The preamble begins here.
\title{An Example Document}  % Declares the document's title.
\author{Leslie Lamport}      % Declares the author's name.
\date{January 21, 1994}      % Deleting this command produces today's date.

\newcommand{\ip}[2]{(#1, #2)}
                             % Defines \ip{arg1}{arg2} to mean
                             % (arg1, arg2).

%\newcommand{\ip}[2]{\langle #1 | #2\rangle}
                             % This is an alternative definition of
                             % \ip that is commented out.

\begin{document}             % End of preamble and beginning of text.

\maketitle                   % Produces the title.

This is an example input file.  Comparing it with
the output it generates can show you how to
produce a simple document of your own.

\section{Ordinary Text}      % Produces section heading.  Lower-level
                             % sections are begun with similar
                             % \subsection and \subsubsection commands.

The ends  of words and sentences are marked
  by   spaces. It  doesn't matter how many
spaces    you type; one is as good as 100.  The
end of   a line counts as a space.

One   or more   blank lines denote the  end
of  a paragraph.

Since any number of consecutive spaces are treated
like a single one, the formatting of the input
file makes no difference to
      \LaTeX,                % The \LaTeX command generates the LaTeX logo.
but it makes a difference to you.  When you use
\LaTeX, making your input file as easy to read
as possible will be a great help as you write
your document and when you change it.  This sample
file shows how you can add comments to your own input
file.

Because printing is different from typewriting,
there are a number of things that you have to do
differently when preparing an input file than if
you were just typing the document directly.
Quotation marks like
       ``this''
have to be handled specially, as do quotes within
quotes:
       ``\,`this'            % \, separates the double and single quote.
        is what I just
        wrote, not  `that'\,''.

Dashes come in three sizes: an
       intra-word
dash, a medium dash for number ranges like
       1--2,
and a punctuation
       dash---like
this.

A sentence-ending space should be larger than the
space between words within a sentence.  You
sometimes have to type special commands in
conjunction with punctuation characters to get
this right, as in the following sentence.
       Gnats, gnus, etc.\ all  % `\ ' makes an inter-word space.
       begin with G\@.         % \@ marks end-of-sentence punctuation.
You should check the spaces after periods when
reading your output to make sure you haven't
forgotten any special cases.  Generating an
ellipsis
       \ldots\               % `\ ' is needed after `\ldots' because TeX
                             % ignores spaces after command names like \ldots
                             % made from \ + letters.
                             %
                             % Note how a `%' character causes TeX to ignore
                             % the end of the input line, so these blank lines
                             % do not start a new paragraph.
                             %
with the right spacing around the periods requires
a special command.

\LaTeX\ interprets some common characters as
commands, so you must type special commands to
generate them.  These characters include the
following:
       \$ \& \% \# \{ and \}.

In printing, text is usually emphasized with an
       \emph{italic}
type style.

\begin{em}
   A long segment of text can also be emphasized
   in this way.  Text within such a segment can be
   given \emph{additional} emphasis.
\end{em}

It is sometimes necessary to prevent \LaTeX\ from
breaking a line where it might otherwise do so.
This may be at a space, as between the ``Mr.''\ and
``Jones'' in
       ``Mr.~Jones'',        % ~ produces an unbreakable interword space.
or within a word---especially when the word is a
symbol like
       \mbox{\emph{itemnum}}
that makes little sense when hyphenated across
lines.

Footnotes\footnote{This is an example of a footnote.}
pose no problem.

\LaTeX\ is good at typesetting mathematical formulas
like
       \( x-3y + z = 7 \)
or
       \( a_{1} > x^{2n} + y^{2n} > x' \)
or
       \( \ip{A}{B} = \sum_{i} a_{i} b_{i} \).
The spaces you type in a formula are
ignored.  Remember that a letter like
       $x$                   % $ ... $  and  \( ... \)  are equivalent
is a formula when it denotes a mathematical
symbol, and it should be typed as one.

\section{Displayed Text}

Text is displayed by indenting it from the left
margin.  Quotations are commonly displayed.  There
are short quotations
\begin{quote}
   This is a short quotation.  It consists of a
   single paragraph of text.  See how it is formatted.
\end{quote}
and longer ones.
\begin{quotation}
   This is a longer quotation.  It consists of two
   paragraphs of text, neither of which are
   particularly interesting.

   This is the second paragraph of the quotation.  It
   is just as dull as the first paragraph.
\end{quotation}
Another frequently-displayed structure is a list.
The following is an example of an \emph{itemized}
list.
\begin{itemize}
   \item This is the first item of an itemized list.
         Each item in the list is marked with a ``tick''.
         You don't have to worry about what kind of tick
         mark is used.

   \item This is the second item of the list.  It
         contains another list nested inside it.  The inner
         list is an \emph{enumerated} list.
         \begin{enumerate}
            \item This is the first item of an enumerated
                  list that is nested within the itemized list.

            \item This is the second item of the inner list.
                  \LaTeX\ allows you to nest lists deeper than
                  you really should.
         \end{enumerate}
         This is the rest of the second item of the outer
         list.  It is no more interesting than any other
         part of the item.
   \item This is the third item of the list.
\end{itemize}
You can even display poetry.
\begin{verse}
   There is an environment
    for verse \\             % The \\ command separates lines
   Whose features some poets % within a stanza.
   will curse.

                             % One or more blank lines separate stanzas.

   For instead of making\\
   Them do \emph{all} line breaking, \\
   It allows them to put too many words on a line when they'd rather be
   forced to be terse.
\end{verse}

Mathematical formulas may also be displayed.  A
displayed formula
is
one-line long; multiline
formulas require special formatting instructions.
   \[  \ip{\Gamma}{\psi'} = x'' + y^{2} + z_{i}^{n}\]
Don't start a paragraph with a displayed equation,
nor make one a paragraph by itself.

\end{document}               % End of document.

\section{Bench Huge Section 430}
% synthetic bench section 430
% This is a sample LaTeX input file.  (Version of 12 August 2004.)
%
% A '%' character causes TeX to ignore all remaining text on the line,
% and is used for comments like this one.

\documentclass{article}      % Specifies the document class

                             % The preamble begins here.
\title{An Example Document}  % Declares the document's title.
\author{Leslie Lamport}      % Declares the author's name.
\date{January 21, 1994}      % Deleting this command produces today's date.

\newcommand{\ip}[2]{(#1, #2)}
                             % Defines \ip{arg1}{arg2} to mean
                             % (arg1, arg2).

%\newcommand{\ip}[2]{\langle #1 | #2\rangle}
                             % This is an alternative definition of
                             % \ip that is commented out.

\begin{document}             % End of preamble and beginning of text.

\maketitle                   % Produces the title.

This is an example input file.  Comparing it with
the output it generates can show you how to
produce a simple document of your own.

\section{Ordinary Text}      % Produces section heading.  Lower-level
                             % sections are begun with similar
                             % \subsection and \subsubsection commands.

The ends  of words and sentences are marked
  by   spaces. It  doesn't matter how many
spaces    you type; one is as good as 100.  The
end of   a line counts as a space.

One   or more   blank lines denote the  end
of  a paragraph.

Since any number of consecutive spaces are treated
like a single one, the formatting of the input
file makes no difference to
      \LaTeX,                % The \LaTeX command generates the LaTeX logo.
but it makes a difference to you.  When you use
\LaTeX, making your input file as easy to read
as possible will be a great help as you write
your document and when you change it.  This sample
file shows how you can add comments to your own input
file.

Because printing is different from typewriting,
there are a number of things that you have to do
differently when preparing an input file than if
you were just typing the document directly.
Quotation marks like
       ``this''
have to be handled specially, as do quotes within
quotes:
       ``\,`this'            % \, separates the double and single quote.
        is what I just
        wrote, not  `that'\,''.

Dashes come in three sizes: an
       intra-word
dash, a medium dash for number ranges like
       1--2,
and a punctuation
       dash---like
this.

A sentence-ending space should be larger than the
space between words within a sentence.  You
sometimes have to type special commands in
conjunction with punctuation characters to get
this right, as in the following sentence.
       Gnats, gnus, etc.\ all  % `\ ' makes an inter-word space.
       begin with G\@.         % \@ marks end-of-sentence punctuation.
You should check the spaces after periods when
reading your output to make sure you haven't
forgotten any special cases.  Generating an
ellipsis
       \ldots\               % `\ ' is needed after `\ldots' because TeX
                             % ignores spaces after command names like \ldots
                             % made from \ + letters.
                             %
                             % Note how a `%' character causes TeX to ignore
                             % the end of the input line, so these blank lines
                             % do not start a new paragraph.
                             %
with the right spacing around the periods requires
a special command.

\LaTeX\ interprets some common characters as
commands, so you must type special commands to
generate them.  These characters include the
following:
       \$ \& \% \# \{ and \}.

In printing, text is usually emphasized with an
       \emph{italic}
type style.

\begin{em}
   A long segment of text can also be emphasized
   in this way.  Text within such a segment can be
   given \emph{additional} emphasis.
\end{em}

It is sometimes necessary to prevent \LaTeX\ from
breaking a line where it might otherwise do so.
This may be at a space, as between the ``Mr.''\ and
``Jones'' in
       ``Mr.~Jones'',        % ~ produces an unbreakable interword space.
or within a word---especially when the word is a
symbol like
       \mbox{\emph{itemnum}}
that makes little sense when hyphenated across
lines.

Footnotes\footnote{This is an example of a footnote.}
pose no problem.

\LaTeX\ is good at typesetting mathematical formulas
like
       \( x-3y + z = 7 \)
or
       \( a_{1} > x^{2n} + y^{2n} > x' \)
or
       \( \ip{A}{B} = \sum_{i} a_{i} b_{i} \).
The spaces you type in a formula are
ignored.  Remember that a letter like
       $x$                   % $ ... $  and  \( ... \)  are equivalent
is a formula when it denotes a mathematical
symbol, and it should be typed as one.

\section{Displayed Text}

Text is displayed by indenting it from the left
margin.  Quotations are commonly displayed.  There
are short quotations
\begin{quote}
   This is a short quotation.  It consists of a
   single paragraph of text.  See how it is formatted.
\end{quote}
and longer ones.
\begin{quotation}
   This is a longer quotation.  It consists of two
   paragraphs of text, neither of which are
   particularly interesting.

   This is the second paragraph of the quotation.  It
   is just as dull as the first paragraph.
\end{quotation}
Another frequently-displayed structure is a list.
The following is an example of an \emph{itemized}
list.
\begin{itemize}
   \item This is the first item of an itemized list.
         Each item in the list is marked with a ``tick''.
         You don't have to worry about what kind of tick
         mark is used.

   \item This is the second item of the list.  It
         contains another list nested inside it.  The inner
         list is an \emph{enumerated} list.
         \begin{enumerate}
            \item This is the first item of an enumerated
                  list that is nested within the itemized list.

            \item This is the second item of the inner list.
                  \LaTeX\ allows you to nest lists deeper than
                  you really should.
         \end{enumerate}
         This is the rest of the second item of the outer
         list.  It is no more interesting than any other
         part of the item.
   \item This is the third item of the list.
\end{itemize}
You can even display poetry.
\begin{verse}
   There is an environment
    for verse \\             % The \\ command separates lines
   Whose features some poets % within a stanza.
   will curse.

                             % One or more blank lines separate stanzas.

   For instead of making\\
   Them do \emph{all} line breaking, \\
   It allows them to put too many words on a line when they'd rather be
   forced to be terse.
\end{verse}

Mathematical formulas may also be displayed.  A
displayed formula
is
one-line long; multiline
formulas require special formatting instructions.
   \[  \ip{\Gamma}{\psi'} = x'' + y^{2} + z_{i}^{n}\]
Don't start a paragraph with a displayed equation,
nor make one a paragraph by itself.

\end{document}               % End of document.

\section{Bench Huge Section 431}
% synthetic bench section 431
% This is a sample LaTeX input file.  (Version of 12 August 2004.)
%
% A '%' character causes TeX to ignore all remaining text on the line,
% and is used for comments like this one.

\documentclass{article}      % Specifies the document class

                             % The preamble begins here.
\title{An Example Document}  % Declares the document's title.
\author{Leslie Lamport}      % Declares the author's name.
\date{January 21, 1994}      % Deleting this command produces today's date.

\newcommand{\ip}[2]{(#1, #2)}
                             % Defines \ip{arg1}{arg2} to mean
                             % (arg1, arg2).

%\newcommand{\ip}[2]{\langle #1 | #2\rangle}
                             % This is an alternative definition of
                             % \ip that is commented out.

\begin{document}             % End of preamble and beginning of text.

\maketitle                   % Produces the title.

This is an example input file.  Comparing it with
the output it generates can show you how to
produce a simple document of your own.

\section{Ordinary Text}      % Produces section heading.  Lower-level
                             % sections are begun with similar
                             % \subsection and \subsubsection commands.

The ends  of words and sentences are marked
  by   spaces. It  doesn't matter how many
spaces    you type; one is as good as 100.  The
end of   a line counts as a space.

One   or more   blank lines denote the  end
of  a paragraph.

Since any number of consecutive spaces are treated
like a single one, the formatting of the input
file makes no difference to
      \LaTeX,                % The \LaTeX command generates the LaTeX logo.
but it makes a difference to you.  When you use
\LaTeX, making your input file as easy to read
as possible will be a great help as you write
your document and when you change it.  This sample
file shows how you can add comments to your own input
file.

Because printing is different from typewriting,
there are a number of things that you have to do
differently when preparing an input file than if
you were just typing the document directly.
Quotation marks like
       ``this''
have to be handled specially, as do quotes within
quotes:
       ``\,`this'            % \, separates the double and single quote.
        is what I just
        wrote, not  `that'\,''.

Dashes come in three sizes: an
       intra-word
dash, a medium dash for number ranges like
       1--2,
and a punctuation
       dash---like
this.

A sentence-ending space should be larger than the
space between words within a sentence.  You
sometimes have to type special commands in
conjunction with punctuation characters to get
this right, as in the following sentence.
       Gnats, gnus, etc.\ all  % `\ ' makes an inter-word space.
       begin with G\@.         % \@ marks end-of-sentence punctuation.
You should check the spaces after periods when
reading your output to make sure you haven't
forgotten any special cases.  Generating an
ellipsis
       \ldots\               % `\ ' is needed after `\ldots' because TeX
                             % ignores spaces after command names like \ldots
                             % made from \ + letters.
                             %
                             % Note how a `%' character causes TeX to ignore
                             % the end of the input line, so these blank lines
                             % do not start a new paragraph.
                             %
with the right spacing around the periods requires
a special command.

\LaTeX\ interprets some common characters as
commands, so you must type special commands to
generate them.  These characters include the
following:
       \$ \& \% \# \{ and \}.

In printing, text is usually emphasized with an
       \emph{italic}
type style.

\begin{em}
   A long segment of text can also be emphasized
   in this way.  Text within such a segment can be
   given \emph{additional} emphasis.
\end{em}

It is sometimes necessary to prevent \LaTeX\ from
breaking a line where it might otherwise do so.
This may be at a space, as between the ``Mr.''\ and
``Jones'' in
       ``Mr.~Jones'',        % ~ produces an unbreakable interword space.
or within a word---especially when the word is a
symbol like
       \mbox{\emph{itemnum}}
that makes little sense when hyphenated across
lines.

Footnotes\footnote{This is an example of a footnote.}
pose no problem.

\LaTeX\ is good at typesetting mathematical formulas
like
       \( x-3y + z = 7 \)
or
       \( a_{1} > x^{2n} + y^{2n} > x' \)
or
       \( \ip{A}{B} = \sum_{i} a_{i} b_{i} \).
The spaces you type in a formula are
ignored.  Remember that a letter like
       $x$                   % $ ... $  and  \( ... \)  are equivalent
is a formula when it denotes a mathematical
symbol, and it should be typed as one.

\section{Displayed Text}

Text is displayed by indenting it from the left
margin.  Quotations are commonly displayed.  There
are short quotations
\begin{quote}
   This is a short quotation.  It consists of a
   single paragraph of text.  See how it is formatted.
\end{quote}
and longer ones.
\begin{quotation}
   This is a longer quotation.  It consists of two
   paragraphs of text, neither of which are
   particularly interesting.

   This is the second paragraph of the quotation.  It
   is just as dull as the first paragraph.
\end{quotation}
Another frequently-displayed structure is a list.
The following is an example of an \emph{itemized}
list.
\begin{itemize}
   \item This is the first item of an itemized list.
         Each item in the list is marked with a ``tick''.
         You don't have to worry about what kind of tick
         mark is used.

   \item This is the second item of the list.  It
         contains another list nested inside it.  The inner
         list is an \emph{enumerated} list.
         \begin{enumerate}
            \item This is the first item of an enumerated
                  list that is nested within the itemized list.

            \item This is the second item of the inner list.
                  \LaTeX\ allows you to nest lists deeper than
                  you really should.
         \end{enumerate}
         This is the rest of the second item of the outer
         list.  It is no more interesting than any other
         part of the item.
   \item This is the third item of the list.
\end{itemize}
You can even display poetry.
\begin{verse}
   There is an environment
    for verse \\             % The \\ command separates lines
   Whose features some poets % within a stanza.
   will curse.

                             % One or more blank lines separate stanzas.

   For instead of making\\
   Them do \emph{all} line breaking, \\
   It allows them to put too many words on a line when they'd rather be
   forced to be terse.
\end{verse}

Mathematical formulas may also be displayed.  A
displayed formula
is
one-line long; multiline
formulas require special formatting instructions.
   \[  \ip{\Gamma}{\psi'} = x'' + y^{2} + z_{i}^{n}\]
Don't start a paragraph with a displayed equation,
nor make one a paragraph by itself.

\end{document}               % End of document.

\section{Bench Huge Section 432}
% synthetic bench section 432
% This is a sample LaTeX input file.  (Version of 12 August 2004.)
%
% A '%' character causes TeX to ignore all remaining text on the line,
% and is used for comments like this one.

\documentclass{article}      % Specifies the document class

                             % The preamble begins here.
\title{An Example Document}  % Declares the document's title.
\author{Leslie Lamport}      % Declares the author's name.
\date{January 21, 1994}      % Deleting this command produces today's date.

\newcommand{\ip}[2]{(#1, #2)}
                             % Defines \ip{arg1}{arg2} to mean
                             % (arg1, arg2).

%\newcommand{\ip}[2]{\langle #1 | #2\rangle}
                             % This is an alternative definition of
                             % \ip that is commented out.

\begin{document}             % End of preamble and beginning of text.

\maketitle                   % Produces the title.

This is an example input file.  Comparing it with
the output it generates can show you how to
produce a simple document of your own.

\section{Ordinary Text}      % Produces section heading.  Lower-level
                             % sections are begun with similar
                             % \subsection and \subsubsection commands.

The ends  of words and sentences are marked
  by   spaces. It  doesn't matter how many
spaces    you type; one is as good as 100.  The
end of   a line counts as a space.

One   or more   blank lines denote the  end
of  a paragraph.

Since any number of consecutive spaces are treated
like a single one, the formatting of the input
file makes no difference to
      \LaTeX,                % The \LaTeX command generates the LaTeX logo.
but it makes a difference to you.  When you use
\LaTeX, making your input file as easy to read
as possible will be a great help as you write
your document and when you change it.  This sample
file shows how you can add comments to your own input
file.

Because printing is different from typewriting,
there are a number of things that you have to do
differently when preparing an input file than if
you were just typing the document directly.
Quotation marks like
       ``this''
have to be handled specially, as do quotes within
quotes:
       ``\,`this'            % \, separates the double and single quote.
        is what I just
        wrote, not  `that'\,''.

Dashes come in three sizes: an
       intra-word
dash, a medium dash for number ranges like
       1--2,
and a punctuation
       dash---like
this.

A sentence-ending space should be larger than the
space between words within a sentence.  You
sometimes have to type special commands in
conjunction with punctuation characters to get
this right, as in the following sentence.
       Gnats, gnus, etc.\ all  % `\ ' makes an inter-word space.
       begin with G\@.         % \@ marks end-of-sentence punctuation.
You should check the spaces after periods when
reading your output to make sure you haven't
forgotten any special cases.  Generating an
ellipsis
       \ldots\               % `\ ' is needed after `\ldots' because TeX
                             % ignores spaces after command names like \ldots
                             % made from \ + letters.
                             %
                             % Note how a `%' character causes TeX to ignore
                             % the end of the input line, so these blank lines
                             % do not start a new paragraph.
                             %
with the right spacing around the periods requires
a special command.

\LaTeX\ interprets some common characters as
commands, so you must type special commands to
generate them.  These characters include the
following:
       \$ \& \% \# \{ and \}.

In printing, text is usually emphasized with an
       \emph{italic}
type style.

\begin{em}
   A long segment of text can also be emphasized
   in this way.  Text within such a segment can be
   given \emph{additional} emphasis.
\end{em}

It is sometimes necessary to prevent \LaTeX\ from
breaking a line where it might otherwise do so.
This may be at a space, as between the ``Mr.''\ and
``Jones'' in
       ``Mr.~Jones'',        % ~ produces an unbreakable interword space.
or within a word---especially when the word is a
symbol like
       \mbox{\emph{itemnum}}
that makes little sense when hyphenated across
lines.

Footnotes\footnote{This is an example of a footnote.}
pose no problem.

\LaTeX\ is good at typesetting mathematical formulas
like
       \( x-3y + z = 7 \)
or
       \( a_{1} > x^{2n} + y^{2n} > x' \)
or
       \( \ip{A}{B} = \sum_{i} a_{i} b_{i} \).
The spaces you type in a formula are
ignored.  Remember that a letter like
       $x$                   % $ ... $  and  \( ... \)  are equivalent
is a formula when it denotes a mathematical
symbol, and it should be typed as one.

\section{Displayed Text}

Text is displayed by indenting it from the left
margin.  Quotations are commonly displayed.  There
are short quotations
\begin{quote}
   This is a short quotation.  It consists of a
   single paragraph of text.  See how it is formatted.
\end{quote}
and longer ones.
\begin{quotation}
   This is a longer quotation.  It consists of two
   paragraphs of text, neither of which are
   particularly interesting.

   This is the second paragraph of the quotation.  It
   is just as dull as the first paragraph.
\end{quotation}
Another frequently-displayed structure is a list.
The following is an example of an \emph{itemized}
list.
\begin{itemize}
   \item This is the first item of an itemized list.
         Each item in the list is marked with a ``tick''.
         You don't have to worry about what kind of tick
         mark is used.

   \item This is the second item of the list.  It
         contains another list nested inside it.  The inner
         list is an \emph{enumerated} list.
         \begin{enumerate}
            \item This is the first item of an enumerated
                  list that is nested within the itemized list.

            \item This is the second item of the inner list.
                  \LaTeX\ allows you to nest lists deeper than
                  you really should.
         \end{enumerate}
         This is the rest of the second item of the outer
         list.  It is no more interesting than any other
         part of the item.
   \item This is the third item of the list.
\end{itemize}
You can even display poetry.
\begin{verse}
   There is an environment
    for verse \\             % The \\ command separates lines
   Whose features some poets % within a stanza.
   will curse.

                             % One or more blank lines separate stanzas.

   For instead of making\\
   Them do \emph{all} line breaking, \\
   It allows them to put too many words on a line when they'd rather be
   forced to be terse.
\end{verse}

Mathematical formulas may also be displayed.  A
displayed formula
is
one-line long; multiline
formulas require special formatting instructions.
   \[  \ip{\Gamma}{\psi'} = x'' + y^{2} + z_{i}^{n}\]
Don't start a paragraph with a displayed equation,
nor make one a paragraph by itself.

\end{document}               % End of document.

\section{Bench Huge Section 433}
% synthetic bench section 433
% This is a sample LaTeX input file.  (Version of 12 August 2004.)
%
% A '%' character causes TeX to ignore all remaining text on the line,
% and is used for comments like this one.

\documentclass{article}      % Specifies the document class

                             % The preamble begins here.
\title{An Example Document}  % Declares the document's title.
\author{Leslie Lamport}      % Declares the author's name.
\date{January 21, 1994}      % Deleting this command produces today's date.

\newcommand{\ip}[2]{(#1, #2)}
                             % Defines \ip{arg1}{arg2} to mean
                             % (arg1, arg2).

%\newcommand{\ip}[2]{\langle #1 | #2\rangle}
                             % This is an alternative definition of
                             % \ip that is commented out.

\begin{document}             % End of preamble and beginning of text.

\maketitle                   % Produces the title.

This is an example input file.  Comparing it with
the output it generates can show you how to
produce a simple document of your own.

\section{Ordinary Text}      % Produces section heading.  Lower-level
                             % sections are begun with similar
                             % \subsection and \subsubsection commands.

The ends  of words and sentences are marked
  by   spaces. It  doesn't matter how many
spaces    you type; one is as good as 100.  The
end of   a line counts as a space.

One   or more   blank lines denote the  end
of  a paragraph.

Since any number of consecutive spaces are treated
like a single one, the formatting of the input
file makes no difference to
      \LaTeX,                % The \LaTeX command generates the LaTeX logo.
but it makes a difference to you.  When you use
\LaTeX, making your input file as easy to read
as possible will be a great help as you write
your document and when you change it.  This sample
file shows how you can add comments to your own input
file.

Because printing is different from typewriting,
there are a number of things that you have to do
differently when preparing an input file than if
you were just typing the document directly.
Quotation marks like
       ``this''
have to be handled specially, as do quotes within
quotes:
       ``\,`this'            % \, separates the double and single quote.
        is what I just
        wrote, not  `that'\,''.

Dashes come in three sizes: an
       intra-word
dash, a medium dash for number ranges like
       1--2,
and a punctuation
       dash---like
this.

A sentence-ending space should be larger than the
space between words within a sentence.  You
sometimes have to type special commands in
conjunction with punctuation characters to get
this right, as in the following sentence.
       Gnats, gnus, etc.\ all  % `\ ' makes an inter-word space.
       begin with G\@.         % \@ marks end-of-sentence punctuation.
You should check the spaces after periods when
reading your output to make sure you haven't
forgotten any special cases.  Generating an
ellipsis
       \ldots\               % `\ ' is needed after `\ldots' because TeX
                             % ignores spaces after command names like \ldots
                             % made from \ + letters.
                             %
                             % Note how a `%' character causes TeX to ignore
                             % the end of the input line, so these blank lines
                             % do not start a new paragraph.
                             %
with the right spacing around the periods requires
a special command.

\LaTeX\ interprets some common characters as
commands, so you must type special commands to
generate them.  These characters include the
following:
       \$ \& \% \# \{ and \}.

In printing, text is usually emphasized with an
       \emph{italic}
type style.

\begin{em}
   A long segment of text can also be emphasized
   in this way.  Text within such a segment can be
   given \emph{additional} emphasis.
\end{em}

It is sometimes necessary to prevent \LaTeX\ from
breaking a line where it might otherwise do so.
This may be at a space, as between the ``Mr.''\ and
``Jones'' in
       ``Mr.~Jones'',        % ~ produces an unbreakable interword space.
or within a word---especially when the word is a
symbol like
       \mbox{\emph{itemnum}}
that makes little sense when hyphenated across
lines.

Footnotes\footnote{This is an example of a footnote.}
pose no problem.

\LaTeX\ is good at typesetting mathematical formulas
like
       \( x-3y + z = 7 \)
or
       \( a_{1} > x^{2n} + y^{2n} > x' \)
or
       \( \ip{A}{B} = \sum_{i} a_{i} b_{i} \).
The spaces you type in a formula are
ignored.  Remember that a letter like
       $x$                   % $ ... $  and  \( ... \)  are equivalent
is a formula when it denotes a mathematical
symbol, and it should be typed as one.

\section{Displayed Text}

Text is displayed by indenting it from the left
margin.  Quotations are commonly displayed.  There
are short quotations
\begin{quote}
   This is a short quotation.  It consists of a
   single paragraph of text.  See how it is formatted.
\end{quote}
and longer ones.
\begin{quotation}
   This is a longer quotation.  It consists of two
   paragraphs of text, neither of which are
   particularly interesting.

   This is the second paragraph of the quotation.  It
   is just as dull as the first paragraph.
\end{quotation}
Another frequently-displayed structure is a list.
The following is an example of an \emph{itemized}
list.
\begin{itemize}
   \item This is the first item of an itemized list.
         Each item in the list is marked with a ``tick''.
         You don't have to worry about what kind of tick
         mark is used.

   \item This is the second item of the list.  It
         contains another list nested inside it.  The inner
         list is an \emph{enumerated} list.
         \begin{enumerate}
            \item This is the first item of an enumerated
                  list that is nested within the itemized list.

            \item This is the second item of the inner list.
                  \LaTeX\ allows you to nest lists deeper than
                  you really should.
         \end{enumerate}
         This is the rest of the second item of the outer
         list.  It is no more interesting than any other
         part of the item.
   \item This is the third item of the list.
\end{itemize}
You can even display poetry.
\begin{verse}
   There is an environment
    for verse \\             % The \\ command separates lines
   Whose features some poets % within a stanza.
   will curse.

                             % One or more blank lines separate stanzas.

   For instead of making\\
   Them do \emph{all} line breaking, \\
   It allows them to put too many words on a line when they'd rather be
   forced to be terse.
\end{verse}

Mathematical formulas may also be displayed.  A
displayed formula
is
one-line long; multiline
formulas require special formatting instructions.
   \[  \ip{\Gamma}{\psi'} = x'' + y^{2} + z_{i}^{n}\]
Don't start a paragraph with a displayed equation,
nor make one a paragraph by itself.

\end{document}               % End of document.

\section{Bench Huge Section 434}
% synthetic bench section 434
% This is a sample LaTeX input file.  (Version of 12 August 2004.)
%
% A '%' character causes TeX to ignore all remaining text on the line,
% and is used for comments like this one.

\documentclass{article}      % Specifies the document class

                             % The preamble begins here.
\title{An Example Document}  % Declares the document's title.
\author{Leslie Lamport}      % Declares the author's name.
\date{January 21, 1994}      % Deleting this command produces today's date.

\newcommand{\ip}[2]{(#1, #2)}
                             % Defines \ip{arg1}{arg2} to mean
                             % (arg1, arg2).

%\newcommand{\ip}[2]{\langle #1 | #2\rangle}
                             % This is an alternative definition of
                             % \ip that is commented out.

\begin{document}             % End of preamble and beginning of text.

\maketitle                   % Produces the title.

This is an example input file.  Comparing it with
the output it generates can show you how to
produce a simple document of your own.

\section{Ordinary Text}      % Produces section heading.  Lower-level
                             % sections are begun with similar
                             % \subsection and \subsubsection commands.

The ends  of words and sentences are marked
  by   spaces. It  doesn't matter how many
spaces    you type; one is as good as 100.  The
end of   a line counts as a space.

One   or more   blank lines denote the  end
of  a paragraph.

Since any number of consecutive spaces are treated
like a single one, the formatting of the input
file makes no difference to
      \LaTeX,                % The \LaTeX command generates the LaTeX logo.
but it makes a difference to you.  When you use
\LaTeX, making your input file as easy to read
as possible will be a great help as you write
your document and when you change it.  This sample
file shows how you can add comments to your own input
file.

Because printing is different from typewriting,
there are a number of things that you have to do
differently when preparing an input file than if
you were just typing the document directly.
Quotation marks like
       ``this''
have to be handled specially, as do quotes within
quotes:
       ``\,`this'            % \, separates the double and single quote.
        is what I just
        wrote, not  `that'\,''.

Dashes come in three sizes: an
       intra-word
dash, a medium dash for number ranges like
       1--2,
and a punctuation
       dash---like
this.

A sentence-ending space should be larger than the
space between words within a sentence.  You
sometimes have to type special commands in
conjunction with punctuation characters to get
this right, as in the following sentence.
       Gnats, gnus, etc.\ all  % `\ ' makes an inter-word space.
       begin with G\@.         % \@ marks end-of-sentence punctuation.
You should check the spaces after periods when
reading your output to make sure you haven't
forgotten any special cases.  Generating an
ellipsis
       \ldots\               % `\ ' is needed after `\ldots' because TeX
                             % ignores spaces after command names like \ldots
                             % made from \ + letters.
                             %
                             % Note how a `%' character causes TeX to ignore
                             % the end of the input line, so these blank lines
                             % do not start a new paragraph.
                             %
with the right spacing around the periods requires
a special command.

\LaTeX\ interprets some common characters as
commands, so you must type special commands to
generate them.  These characters include the
following:
       \$ \& \% \# \{ and \}.

In printing, text is usually emphasized with an
       \emph{italic}
type style.

\begin{em}
   A long segment of text can also be emphasized
   in this way.  Text within such a segment can be
   given \emph{additional} emphasis.
\end{em}

It is sometimes necessary to prevent \LaTeX\ from
breaking a line where it might otherwise do so.
This may be at a space, as between the ``Mr.''\ and
``Jones'' in
       ``Mr.~Jones'',        % ~ produces an unbreakable interword space.
or within a word---especially when the word is a
symbol like
       \mbox{\emph{itemnum}}
that makes little sense when hyphenated across
lines.

Footnotes\footnote{This is an example of a footnote.}
pose no problem.

\LaTeX\ is good at typesetting mathematical formulas
like
       \( x-3y + z = 7 \)
or
       \( a_{1} > x^{2n} + y^{2n} > x' \)
or
       \( \ip{A}{B} = \sum_{i} a_{i} b_{i} \).
The spaces you type in a formula are
ignored.  Remember that a letter like
       $x$                   % $ ... $  and  \( ... \)  are equivalent
is a formula when it denotes a mathematical
symbol, and it should be typed as one.

\section{Displayed Text}

Text is displayed by indenting it from the left
margin.  Quotations are commonly displayed.  There
are short quotations
\begin{quote}
   This is a short quotation.  It consists of a
   single paragraph of text.  See how it is formatted.
\end{quote}
and longer ones.
\begin{quotation}
   This is a longer quotation.  It consists of two
   paragraphs of text, neither of which are
   particularly interesting.

   This is the second paragraph of the quotation.  It
   is just as dull as the first paragraph.
\end{quotation}
Another frequently-displayed structure is a list.
The following is an example of an \emph{itemized}
list.
\begin{itemize}
   \item This is the first item of an itemized list.
         Each item in the list is marked with a ``tick''.
         You don't have to worry about what kind of tick
         mark is used.

   \item This is the second item of the list.  It
         contains another list nested inside it.  The inner
         list is an \emph{enumerated} list.
         \begin{enumerate}
            \item This is the first item of an enumerated
                  list that is nested within the itemized list.

            \item This is the second item of the inner list.
                  \LaTeX\ allows you to nest lists deeper than
                  you really should.
         \end{enumerate}
         This is the rest of the second item of the outer
         list.  It is no more interesting than any other
         part of the item.
   \item This is the third item of the list.
\end{itemize}
You can even display poetry.
\begin{verse}
   There is an environment
    for verse \\             % The \\ command separates lines
   Whose features some poets % within a stanza.
   will curse.

                             % One or more blank lines separate stanzas.

   For instead of making\\
   Them do \emph{all} line breaking, \\
   It allows them to put too many words on a line when they'd rather be
   forced to be terse.
\end{verse}

Mathematical formulas may also be displayed.  A
displayed formula
is
one-line long; multiline
formulas require special formatting instructions.
   \[  \ip{\Gamma}{\psi'} = x'' + y^{2} + z_{i}^{n}\]
Don't start a paragraph with a displayed equation,
nor make one a paragraph by itself.

\end{document}               % End of document.

\section{Bench Huge Section 435}
% synthetic bench section 435
% This is a sample LaTeX input file.  (Version of 12 August 2004.)
%
% A '%' character causes TeX to ignore all remaining text on the line,
% and is used for comments like this one.

\documentclass{article}      % Specifies the document class

                             % The preamble begins here.
\title{An Example Document}  % Declares the document's title.
\author{Leslie Lamport}      % Declares the author's name.
\date{January 21, 1994}      % Deleting this command produces today's date.

\newcommand{\ip}[2]{(#1, #2)}
                             % Defines \ip{arg1}{arg2} to mean
                             % (arg1, arg2).

%\newcommand{\ip}[2]{\langle #1 | #2\rangle}
                             % This is an alternative definition of
                             % \ip that is commented out.

\begin{document}             % End of preamble and beginning of text.

\maketitle                   % Produces the title.

This is an example input file.  Comparing it with
the output it generates can show you how to
produce a simple document of your own.

\section{Ordinary Text}      % Produces section heading.  Lower-level
                             % sections are begun with similar
                             % \subsection and \subsubsection commands.

The ends  of words and sentences are marked
  by   spaces. It  doesn't matter how many
spaces    you type; one is as good as 100.  The
end of   a line counts as a space.

One   or more   blank lines denote the  end
of  a paragraph.

Since any number of consecutive spaces are treated
like a single one, the formatting of the input
file makes no difference to
      \LaTeX,                % The \LaTeX command generates the LaTeX logo.
but it makes a difference to you.  When you use
\LaTeX, making your input file as easy to read
as possible will be a great help as you write
your document and when you change it.  This sample
file shows how you can add comments to your own input
file.

Because printing is different from typewriting,
there are a number of things that you have to do
differently when preparing an input file than if
you were just typing the document directly.
Quotation marks like
       ``this''
have to be handled specially, as do quotes within
quotes:
       ``\,`this'            % \, separates the double and single quote.
        is what I just
        wrote, not  `that'\,''.

Dashes come in three sizes: an
       intra-word
dash, a medium dash for number ranges like
       1--2,
and a punctuation
       dash---like
this.

A sentence-ending space should be larger than the
space between words within a sentence.  You
sometimes have to type special commands in
conjunction with punctuation characters to get
this right, as in the following sentence.
       Gnats, gnus, etc.\ all  % `\ ' makes an inter-word space.
       begin with G\@.         % \@ marks end-of-sentence punctuation.
You should check the spaces after periods when
reading your output to make sure you haven't
forgotten any special cases.  Generating an
ellipsis
       \ldots\               % `\ ' is needed after `\ldots' because TeX
                             % ignores spaces after command names like \ldots
                             % made from \ + letters.
                             %
                             % Note how a `%' character causes TeX to ignore
                             % the end of the input line, so these blank lines
                             % do not start a new paragraph.
                             %
with the right spacing around the periods requires
a special command.

\LaTeX\ interprets some common characters as
commands, so you must type special commands to
generate them.  These characters include the
following:
       \$ \& \% \# \{ and \}.

In printing, text is usually emphasized with an
       \emph{italic}
type style.

\begin{em}
   A long segment of text can also be emphasized
   in this way.  Text within such a segment can be
   given \emph{additional} emphasis.
\end{em}

It is sometimes necessary to prevent \LaTeX\ from
breaking a line where it might otherwise do so.
This may be at a space, as between the ``Mr.''\ and
``Jones'' in
       ``Mr.~Jones'',        % ~ produces an unbreakable interword space.
or within a word---especially when the word is a
symbol like
       \mbox{\emph{itemnum}}
that makes little sense when hyphenated across
lines.

Footnotes\footnote{This is an example of a footnote.}
pose no problem.

\LaTeX\ is good at typesetting mathematical formulas
like
       \( x-3y + z = 7 \)
or
       \( a_{1} > x^{2n} + y^{2n} > x' \)
or
       \( \ip{A}{B} = \sum_{i} a_{i} b_{i} \).
The spaces you type in a formula are
ignored.  Remember that a letter like
       $x$                   % $ ... $  and  \( ... \)  are equivalent
is a formula when it denotes a mathematical
symbol, and it should be typed as one.

\section{Displayed Text}

Text is displayed by indenting it from the left
margin.  Quotations are commonly displayed.  There
are short quotations
\begin{quote}
   This is a short quotation.  It consists of a
   single paragraph of text.  See how it is formatted.
\end{quote}
and longer ones.
\begin{quotation}
   This is a longer quotation.  It consists of two
   paragraphs of text, neither of which are
   particularly interesting.

   This is the second paragraph of the quotation.  It
   is just as dull as the first paragraph.
\end{quotation}
Another frequently-displayed structure is a list.
The following is an example of an \emph{itemized}
list.
\begin{itemize}
   \item This is the first item of an itemized list.
         Each item in the list is marked with a ``tick''.
         You don't have to worry about what kind of tick
         mark is used.

   \item This is the second item of the list.  It
         contains another list nested inside it.  The inner
         list is an \emph{enumerated} list.
         \begin{enumerate}
            \item This is the first item of an enumerated
                  list that is nested within the itemized list.

            \item This is the second item of the inner list.
                  \LaTeX\ allows you to nest lists deeper than
                  you really should.
         \end{enumerate}
         This is the rest of the second item of the outer
         list.  It is no more interesting than any other
         part of the item.
   \item This is the third item of the list.
\end{itemize}
You can even display poetry.
\begin{verse}
   There is an environment
    for verse \\             % The \\ command separates lines
   Whose features some poets % within a stanza.
   will curse.

                             % One or more blank lines separate stanzas.

   For instead of making\\
   Them do \emph{all} line breaking, \\
   It allows them to put too many words on a line when they'd rather be
   forced to be terse.
\end{verse}

Mathematical formulas may also be displayed.  A
displayed formula
is
one-line long; multiline
formulas require special formatting instructions.
   \[  \ip{\Gamma}{\psi'} = x'' + y^{2} + z_{i}^{n}\]
Don't start a paragraph with a displayed equation,
nor make one a paragraph by itself.

\end{document}               % End of document.

\section{Bench Huge Section 436}
% synthetic bench section 436
% This is a sample LaTeX input file.  (Version of 12 August 2004.)
%
% A '%' character causes TeX to ignore all remaining text on the line,
% and is used for comments like this one.

\documentclass{article}      % Specifies the document class

                             % The preamble begins here.
\title{An Example Document}  % Declares the document's title.
\author{Leslie Lamport}      % Declares the author's name.
\date{January 21, 1994}      % Deleting this command produces today's date.

\newcommand{\ip}[2]{(#1, #2)}
                             % Defines \ip{arg1}{arg2} to mean
                             % (arg1, arg2).

%\newcommand{\ip}[2]{\langle #1 | #2\rangle}
                             % This is an alternative definition of
                             % \ip that is commented out.

\begin{document}             % End of preamble and beginning of text.

\maketitle                   % Produces the title.

This is an example input file.  Comparing it with
the output it generates can show you how to
produce a simple document of your own.

\section{Ordinary Text}      % Produces section heading.  Lower-level
                             % sections are begun with similar
                             % \subsection and \subsubsection commands.

The ends  of words and sentences are marked
  by   spaces. It  doesn't matter how many
spaces    you type; one is as good as 100.  The
end of   a line counts as a space.

One   or more   blank lines denote the  end
of  a paragraph.

Since any number of consecutive spaces are treated
like a single one, the formatting of the input
file makes no difference to
      \LaTeX,                % The \LaTeX command generates the LaTeX logo.
but it makes a difference to you.  When you use
\LaTeX, making your input file as easy to read
as possible will be a great help as you write
your document and when you change it.  This sample
file shows how you can add comments to your own input
file.

Because printing is different from typewriting,
there are a number of things that you have to do
differently when preparing an input file than if
you were just typing the document directly.
Quotation marks like
       ``this''
have to be handled specially, as do quotes within
quotes:
       ``\,`this'            % \, separates the double and single quote.
        is what I just
        wrote, not  `that'\,''.

Dashes come in three sizes: an
       intra-word
dash, a medium dash for number ranges like
       1--2,
and a punctuation
       dash---like
this.

A sentence-ending space should be larger than the
space between words within a sentence.  You
sometimes have to type special commands in
conjunction with punctuation characters to get
this right, as in the following sentence.
       Gnats, gnus, etc.\ all  % `\ ' makes an inter-word space.
       begin with G\@.         % \@ marks end-of-sentence punctuation.
You should check the spaces after periods when
reading your output to make sure you haven't
forgotten any special cases.  Generating an
ellipsis
       \ldots\               % `\ ' is needed after `\ldots' because TeX
                             % ignores spaces after command names like \ldots
                             % made from \ + letters.
                             %
                             % Note how a `%' character causes TeX to ignore
                             % the end of the input line, so these blank lines
                             % do not start a new paragraph.
                             %
with the right spacing around the periods requires
a special command.

\LaTeX\ interprets some common characters as
commands, so you must type special commands to
generate them.  These characters include the
following:
       \$ \& \% \# \{ and \}.

In printing, text is usually emphasized with an
       \emph{italic}
type style.

\begin{em}
   A long segment of text can also be emphasized
   in this way.  Text within such a segment can be
   given \emph{additional} emphasis.
\end{em}

It is sometimes necessary to prevent \LaTeX\ from
breaking a line where it might otherwise do so.
This may be at a space, as between the ``Mr.''\ and
``Jones'' in
       ``Mr.~Jones'',        % ~ produces an unbreakable interword space.
or within a word---especially when the word is a
symbol like
       \mbox{\emph{itemnum}}
that makes little sense when hyphenated across
lines.

Footnotes\footnote{This is an example of a footnote.}
pose no problem.

\LaTeX\ is good at typesetting mathematical formulas
like
       \( x-3y + z = 7 \)
or
       \( a_{1} > x^{2n} + y^{2n} > x' \)
or
       \( \ip{A}{B} = \sum_{i} a_{i} b_{i} \).
The spaces you type in a formula are
ignored.  Remember that a letter like
       $x$                   % $ ... $  and  \( ... \)  are equivalent
is a formula when it denotes a mathematical
symbol, and it should be typed as one.

\section{Displayed Text}

Text is displayed by indenting it from the left
margin.  Quotations are commonly displayed.  There
are short quotations
\begin{quote}
   This is a short quotation.  It consists of a
   single paragraph of text.  See how it is formatted.
\end{quote}
and longer ones.
\begin{quotation}
   This is a longer quotation.  It consists of two
   paragraphs of text, neither of which are
   particularly interesting.

   This is the second paragraph of the quotation.  It
   is just as dull as the first paragraph.
\end{quotation}
Another frequently-displayed structure is a list.
The following is an example of an \emph{itemized}
list.
\begin{itemize}
   \item This is the first item of an itemized list.
         Each item in the list is marked with a ``tick''.
         You don't have to worry about what kind of tick
         mark is used.

   \item This is the second item of the list.  It
         contains another list nested inside it.  The inner
         list is an \emph{enumerated} list.
         \begin{enumerate}
            \item This is the first item of an enumerated
                  list that is nested within the itemized list.

            \item This is the second item of the inner list.
                  \LaTeX\ allows you to nest lists deeper than
                  you really should.
         \end{enumerate}
         This is the rest of the second item of the outer
         list.  It is no more interesting than any other
         part of the item.
   \item This is the third item of the list.
\end{itemize}
You can even display poetry.
\begin{verse}
   There is an environment
    for verse \\             % The \\ command separates lines
   Whose features some poets % within a stanza.
   will curse.

                             % One or more blank lines separate stanzas.

   For instead of making\\
   Them do \emph{all} line breaking, \\
   It allows them to put too many words on a line when they'd rather be
   forced to be terse.
\end{verse}

Mathematical formulas may also be displayed.  A
displayed formula
is
one-line long; multiline
formulas require special formatting instructions.
   \[  \ip{\Gamma}{\psi'} = x'' + y^{2} + z_{i}^{n}\]
Don't start a paragraph with a displayed equation,
nor make one a paragraph by itself.

\end{document}               % End of document.

\section{Bench Huge Section 437}
% synthetic bench section 437
% This is a sample LaTeX input file.  (Version of 12 August 2004.)
%
% A '%' character causes TeX to ignore all remaining text on the line,
% and is used for comments like this one.

\documentclass{article}      % Specifies the document class

                             % The preamble begins here.
\title{An Example Document}  % Declares the document's title.
\author{Leslie Lamport}      % Declares the author's name.
\date{January 21, 1994}      % Deleting this command produces today's date.

\newcommand{\ip}[2]{(#1, #2)}
                             % Defines \ip{arg1}{arg2} to mean
                             % (arg1, arg2).

%\newcommand{\ip}[2]{\langle #1 | #2\rangle}
                             % This is an alternative definition of
                             % \ip that is commented out.

\begin{document}             % End of preamble and beginning of text.

\maketitle                   % Produces the title.

This is an example input file.  Comparing it with
the output it generates can show you how to
produce a simple document of your own.

\section{Ordinary Text}      % Produces section heading.  Lower-level
                             % sections are begun with similar
                             % \subsection and \subsubsection commands.

The ends  of words and sentences are marked
  by   spaces. It  doesn't matter how many
spaces    you type; one is as good as 100.  The
end of   a line counts as a space.

One   or more   blank lines denote the  end
of  a paragraph.

Since any number of consecutive spaces are treated
like a single one, the formatting of the input
file makes no difference to
      \LaTeX,                % The \LaTeX command generates the LaTeX logo.
but it makes a difference to you.  When you use
\LaTeX, making your input file as easy to read
as possible will be a great help as you write
your document and when you change it.  This sample
file shows how you can add comments to your own input
file.

Because printing is different from typewriting,
there are a number of things that you have to do
differently when preparing an input file than if
you were just typing the document directly.
Quotation marks like
       ``this''
have to be handled specially, as do quotes within
quotes:
       ``\,`this'            % \, separates the double and single quote.
        is what I just
        wrote, not  `that'\,''.

Dashes come in three sizes: an
       intra-word
dash, a medium dash for number ranges like
       1--2,
and a punctuation
       dash---like
this.

A sentence-ending space should be larger than the
space between words within a sentence.  You
sometimes have to type special commands in
conjunction with punctuation characters to get
this right, as in the following sentence.
       Gnats, gnus, etc.\ all  % `\ ' makes an inter-word space.
       begin with G\@.         % \@ marks end-of-sentence punctuation.
You should check the spaces after periods when
reading your output to make sure you haven't
forgotten any special cases.  Generating an
ellipsis
       \ldots\               % `\ ' is needed after `\ldots' because TeX
                             % ignores spaces after command names like \ldots
                             % made from \ + letters.
                             %
                             % Note how a `%' character causes TeX to ignore
                             % the end of the input line, so these blank lines
                             % do not start a new paragraph.
                             %
with the right spacing around the periods requires
a special command.

\LaTeX\ interprets some common characters as
commands, so you must type special commands to
generate them.  These characters include the
following:
       \$ \& \% \# \{ and \}.

In printing, text is usually emphasized with an
       \emph{italic}
type style.

\begin{em}
   A long segment of text can also be emphasized
   in this way.  Text within such a segment can be
   given \emph{additional} emphasis.
\end{em}

It is sometimes necessary to prevent \LaTeX\ from
breaking a line where it might otherwise do so.
This may be at a space, as between the ``Mr.''\ and
``Jones'' in
       ``Mr.~Jones'',        % ~ produces an unbreakable interword space.
or within a word---especially when the word is a
symbol like
       \mbox{\emph{itemnum}}
that makes little sense when hyphenated across
lines.

Footnotes\footnote{This is an example of a footnote.}
pose no problem.

\LaTeX\ is good at typesetting mathematical formulas
like
       \( x-3y + z = 7 \)
or
       \( a_{1} > x^{2n} + y^{2n} > x' \)
or
       \( \ip{A}{B} = \sum_{i} a_{i} b_{i} \).
The spaces you type in a formula are
ignored.  Remember that a letter like
       $x$                   % $ ... $  and  \( ... \)  are equivalent
is a formula when it denotes a mathematical
symbol, and it should be typed as one.

\section{Displayed Text}

Text is displayed by indenting it from the left
margin.  Quotations are commonly displayed.  There
are short quotations
\begin{quote}
   This is a short quotation.  It consists of a
   single paragraph of text.  See how it is formatted.
\end{quote}
and longer ones.
\begin{quotation}
   This is a longer quotation.  It consists of two
   paragraphs of text, neither of which are
   particularly interesting.

   This is the second paragraph of the quotation.  It
   is just as dull as the first paragraph.
\end{quotation}
Another frequently-displayed structure is a list.
The following is an example of an \emph{itemized}
list.
\begin{itemize}
   \item This is the first item of an itemized list.
         Each item in the list is marked with a ``tick''.
         You don't have to worry about what kind of tick
         mark is used.

   \item This is the second item of the list.  It
         contains another list nested inside it.  The inner
         list is an \emph{enumerated} list.
         \begin{enumerate}
            \item This is the first item of an enumerated
                  list that is nested within the itemized list.

            \item This is the second item of the inner list.
                  \LaTeX\ allows you to nest lists deeper than
                  you really should.
         \end{enumerate}
         This is the rest of the second item of the outer
         list.  It is no more interesting than any other
         part of the item.
   \item This is the third item of the list.
\end{itemize}
You can even display poetry.
\begin{verse}
   There is an environment
    for verse \\             % The \\ command separates lines
   Whose features some poets % within a stanza.
   will curse.

                             % One or more blank lines separate stanzas.

   For instead of making\\
   Them do \emph{all} line breaking, \\
   It allows them to put too many words on a line when they'd rather be
   forced to be terse.
\end{verse}

Mathematical formulas may also be displayed.  A
displayed formula
is
one-line long; multiline
formulas require special formatting instructions.
   \[  \ip{\Gamma}{\psi'} = x'' + y^{2} + z_{i}^{n}\]
Don't start a paragraph with a displayed equation,
nor make one a paragraph by itself.

\end{document}               % End of document.

\section{Bench Huge Section 438}
% synthetic bench section 438
% This is a sample LaTeX input file.  (Version of 12 August 2004.)
%
% A '%' character causes TeX to ignore all remaining text on the line,
% and is used for comments like this one.

\documentclass{article}      % Specifies the document class

                             % The preamble begins here.
\title{An Example Document}  % Declares the document's title.
\author{Leslie Lamport}      % Declares the author's name.
\date{January 21, 1994}      % Deleting this command produces today's date.

\newcommand{\ip}[2]{(#1, #2)}
                             % Defines \ip{arg1}{arg2} to mean
                             % (arg1, arg2).

%\newcommand{\ip}[2]{\langle #1 | #2\rangle}
                             % This is an alternative definition of
                             % \ip that is commented out.

\begin{document}             % End of preamble and beginning of text.

\maketitle                   % Produces the title.

This is an example input file.  Comparing it with
the output it generates can show you how to
produce a simple document of your own.

\section{Ordinary Text}      % Produces section heading.  Lower-level
                             % sections are begun with similar
                             % \subsection and \subsubsection commands.

The ends  of words and sentences are marked
  by   spaces. It  doesn't matter how many
spaces    you type; one is as good as 100.  The
end of   a line counts as a space.

One   or more   blank lines denote the  end
of  a paragraph.

Since any number of consecutive spaces are treated
like a single one, the formatting of the input
file makes no difference to
      \LaTeX,                % The \LaTeX command generates the LaTeX logo.
but it makes a difference to you.  When you use
\LaTeX, making your input file as easy to read
as possible will be a great help as you write
your document and when you change it.  This sample
file shows how you can add comments to your own input
file.

Because printing is different from typewriting,
there are a number of things that you have to do
differently when preparing an input file than if
you were just typing the document directly.
Quotation marks like
       ``this''
have to be handled specially, as do quotes within
quotes:
       ``\,`this'            % \, separates the double and single quote.
        is what I just
        wrote, not  `that'\,''.

Dashes come in three sizes: an
       intra-word
dash, a medium dash for number ranges like
       1--2,
and a punctuation
       dash---like
this.

A sentence-ending space should be larger than the
space between words within a sentence.  You
sometimes have to type special commands in
conjunction with punctuation characters to get
this right, as in the following sentence.
       Gnats, gnus, etc.\ all  % `\ ' makes an inter-word space.
       begin with G\@.         % \@ marks end-of-sentence punctuation.
You should check the spaces after periods when
reading your output to make sure you haven't
forgotten any special cases.  Generating an
ellipsis
       \ldots\               % `\ ' is needed after `\ldots' because TeX
                             % ignores spaces after command names like \ldots
                             % made from \ + letters.
                             %
                             % Note how a `%' character causes TeX to ignore
                             % the end of the input line, so these blank lines
                             % do not start a new paragraph.
                             %
with the right spacing around the periods requires
a special command.

\LaTeX\ interprets some common characters as
commands, so you must type special commands to
generate them.  These characters include the
following:
       \$ \& \% \# \{ and \}.

In printing, text is usually emphasized with an
       \emph{italic}
type style.

\begin{em}
   A long segment of text can also be emphasized
   in this way.  Text within such a segment can be
   given \emph{additional} emphasis.
\end{em}

It is sometimes necessary to prevent \LaTeX\ from
breaking a line where it might otherwise do so.
This may be at a space, as between the ``Mr.''\ and
``Jones'' in
       ``Mr.~Jones'',        % ~ produces an unbreakable interword space.
or within a word---especially when the word is a
symbol like
       \mbox{\emph{itemnum}}
that makes little sense when hyphenated across
lines.

Footnotes\footnote{This is an example of a footnote.}
pose no problem.

\LaTeX\ is good at typesetting mathematical formulas
like
       \( x-3y + z = 7 \)
or
       \( a_{1} > x^{2n} + y^{2n} > x' \)
or
       \( \ip{A}{B} = \sum_{i} a_{i} b_{i} \).
The spaces you type in a formula are
ignored.  Remember that a letter like
       $x$                   % $ ... $  and  \( ... \)  are equivalent
is a formula when it denotes a mathematical
symbol, and it should be typed as one.

\section{Displayed Text}

Text is displayed by indenting it from the left
margin.  Quotations are commonly displayed.  There
are short quotations
\begin{quote}
   This is a short quotation.  It consists of a
   single paragraph of text.  See how it is formatted.
\end{quote}
and longer ones.
\begin{quotation}
   This is a longer quotation.  It consists of two
   paragraphs of text, neither of which are
   particularly interesting.

   This is the second paragraph of the quotation.  It
   is just as dull as the first paragraph.
\end{quotation}
Another frequently-displayed structure is a list.
The following is an example of an \emph{itemized}
list.
\begin{itemize}
   \item This is the first item of an itemized list.
         Each item in the list is marked with a ``tick''.
         You don't have to worry about what kind of tick
         mark is used.

   \item This is the second item of the list.  It
         contains another list nested inside it.  The inner
         list is an \emph{enumerated} list.
         \begin{enumerate}
            \item This is the first item of an enumerated
                  list that is nested within the itemized list.

            \item This is the second item of the inner list.
                  \LaTeX\ allows you to nest lists deeper than
                  you really should.
         \end{enumerate}
         This is the rest of the second item of the outer
         list.  It is no more interesting than any other
         part of the item.
   \item This is the third item of the list.
\end{itemize}
You can even display poetry.
\begin{verse}
   There is an environment
    for verse \\             % The \\ command separates lines
   Whose features some poets % within a stanza.
   will curse.

                             % One or more blank lines separate stanzas.

   For instead of making\\
   Them do \emph{all} line breaking, \\
   It allows them to put too many words on a line when they'd rather be
   forced to be terse.
\end{verse}

Mathematical formulas may also be displayed.  A
displayed formula
is
one-line long; multiline
formulas require special formatting instructions.
   \[  \ip{\Gamma}{\psi'} = x'' + y^{2} + z_{i}^{n}\]
Don't start a paragraph with a displayed equation,
nor make one a paragraph by itself.

\end{document}               % End of document.

\section{Bench Huge Section 439}
% synthetic bench section 439
% This is a sample LaTeX input file.  (Version of 12 August 2004.)
%
% A '%' character causes TeX to ignore all remaining text on the line,
% and is used for comments like this one.

\documentclass{article}      % Specifies the document class

                             % The preamble begins here.
\title{An Example Document}  % Declares the document's title.
\author{Leslie Lamport}      % Declares the author's name.
\date{January 21, 1994}      % Deleting this command produces today's date.

\newcommand{\ip}[2]{(#1, #2)}
                             % Defines \ip{arg1}{arg2} to mean
                             % (arg1, arg2).

%\newcommand{\ip}[2]{\langle #1 | #2\rangle}
                             % This is an alternative definition of
                             % \ip that is commented out.

\begin{document}             % End of preamble and beginning of text.

\maketitle                   % Produces the title.

This is an example input file.  Comparing it with
the output it generates can show you how to
produce a simple document of your own.

\section{Ordinary Text}      % Produces section heading.  Lower-level
                             % sections are begun with similar
                             % \subsection and \subsubsection commands.

The ends  of words and sentences are marked
  by   spaces. It  doesn't matter how many
spaces    you type; one is as good as 100.  The
end of   a line counts as a space.

One   or more   blank lines denote the  end
of  a paragraph.

Since any number of consecutive spaces are treated
like a single one, the formatting of the input
file makes no difference to
      \LaTeX,                % The \LaTeX command generates the LaTeX logo.
but it makes a difference to you.  When you use
\LaTeX, making your input file as easy to read
as possible will be a great help as you write
your document and when you change it.  This sample
file shows how you can add comments to your own input
file.

Because printing is different from typewriting,
there are a number of things that you have to do
differently when preparing an input file than if
you were just typing the document directly.
Quotation marks like
       ``this''
have to be handled specially, as do quotes within
quotes:
       ``\,`this'            % \, separates the double and single quote.
        is what I just
        wrote, not  `that'\,''.

Dashes come in three sizes: an
       intra-word
dash, a medium dash for number ranges like
       1--2,
and a punctuation
       dash---like
this.

A sentence-ending space should be larger than the
space between words within a sentence.  You
sometimes have to type special commands in
conjunction with punctuation characters to get
this right, as in the following sentence.
       Gnats, gnus, etc.\ all  % `\ ' makes an inter-word space.
       begin with G\@.         % \@ marks end-of-sentence punctuation.
You should check the spaces after periods when
reading your output to make sure you haven't
forgotten any special cases.  Generating an
ellipsis
       \ldots\               % `\ ' is needed after `\ldots' because TeX
                             % ignores spaces after command names like \ldots
                             % made from \ + letters.
                             %
                             % Note how a `%' character causes TeX to ignore
                             % the end of the input line, so these blank lines
                             % do not start a new paragraph.
                             %
with the right spacing around the periods requires
a special command.

\LaTeX\ interprets some common characters as
commands, so you must type special commands to
generate them.  These characters include the
following:
       \$ \& \% \# \{ and \}.

In printing, text is usually emphasized with an
       \emph{italic}
type style.

\begin{em}
   A long segment of text can also be emphasized
   in this way.  Text within such a segment can be
   given \emph{additional} emphasis.
\end{em}

It is sometimes necessary to prevent \LaTeX\ from
breaking a line where it might otherwise do so.
This may be at a space, as between the ``Mr.''\ and
``Jones'' in
       ``Mr.~Jones'',        % ~ produces an unbreakable interword space.
or within a word---especially when the word is a
symbol like
       \mbox{\emph{itemnum}}
that makes little sense when hyphenated across
lines.

Footnotes\footnote{This is an example of a footnote.}
pose no problem.

\LaTeX\ is good at typesetting mathematical formulas
like
       \( x-3y + z = 7 \)
or
       \( a_{1} > x^{2n} + y^{2n} > x' \)
or
       \( \ip{A}{B} = \sum_{i} a_{i} b_{i} \).
The spaces you type in a formula are
ignored.  Remember that a letter like
       $x$                   % $ ... $  and  \( ... \)  are equivalent
is a formula when it denotes a mathematical
symbol, and it should be typed as one.

\section{Displayed Text}

Text is displayed by indenting it from the left
margin.  Quotations are commonly displayed.  There
are short quotations
\begin{quote}
   This is a short quotation.  It consists of a
   single paragraph of text.  See how it is formatted.
\end{quote}
and longer ones.
\begin{quotation}
   This is a longer quotation.  It consists of two
   paragraphs of text, neither of which are
   particularly interesting.

   This is the second paragraph of the quotation.  It
   is just as dull as the first paragraph.
\end{quotation}
Another frequently-displayed structure is a list.
The following is an example of an \emph{itemized}
list.
\begin{itemize}
   \item This is the first item of an itemized list.
         Each item in the list is marked with a ``tick''.
         You don't have to worry about what kind of tick
         mark is used.

   \item This is the second item of the list.  It
         contains another list nested inside it.  The inner
         list is an \emph{enumerated} list.
         \begin{enumerate}
            \item This is the first item of an enumerated
                  list that is nested within the itemized list.

            \item This is the second item of the inner list.
                  \LaTeX\ allows you to nest lists deeper than
                  you really should.
         \end{enumerate}
         This is the rest of the second item of the outer
         list.  It is no more interesting than any other
         part of the item.
   \item This is the third item of the list.
\end{itemize}
You can even display poetry.
\begin{verse}
   There is an environment
    for verse \\             % The \\ command separates lines
   Whose features some poets % within a stanza.
   will curse.

                             % One or more blank lines separate stanzas.

   For instead of making\\
   Them do \emph{all} line breaking, \\
   It allows them to put too many words on a line when they'd rather be
   forced to be terse.
\end{verse}

Mathematical formulas may also be displayed.  A
displayed formula
is
one-line long; multiline
formulas require special formatting instructions.
   \[  \ip{\Gamma}{\psi'} = x'' + y^{2} + z_{i}^{n}\]
Don't start a paragraph with a displayed equation,
nor make one a paragraph by itself.

\end{document}               % End of document.

\section{Bench Huge Section 440}
% synthetic bench section 440
% This is a sample LaTeX input file.  (Version of 12 August 2004.)
%
% A '%' character causes TeX to ignore all remaining text on the line,
% and is used for comments like this one.

\documentclass{article}      % Specifies the document class

                             % The preamble begins here.
\title{An Example Document}  % Declares the document's title.
\author{Leslie Lamport}      % Declares the author's name.
\date{January 21, 1994}      % Deleting this command produces today's date.

\newcommand{\ip}[2]{(#1, #2)}
                             % Defines \ip{arg1}{arg2} to mean
                             % (arg1, arg2).

%\newcommand{\ip}[2]{\langle #1 | #2\rangle}
                             % This is an alternative definition of
                             % \ip that is commented out.

\begin{document}             % End of preamble and beginning of text.

\maketitle                   % Produces the title.

This is an example input file.  Comparing it with
the output it generates can show you how to
produce a simple document of your own.

\section{Ordinary Text}      % Produces section heading.  Lower-level
                             % sections are begun with similar
                             % \subsection and \subsubsection commands.

The ends  of words and sentences are marked
  by   spaces. It  doesn't matter how many
spaces    you type; one is as good as 100.  The
end of   a line counts as a space.

One   or more   blank lines denote the  end
of  a paragraph.

Since any number of consecutive spaces are treated
like a single one, the formatting of the input
file makes no difference to
      \LaTeX,                % The \LaTeX command generates the LaTeX logo.
but it makes a difference to you.  When you use
\LaTeX, making your input file as easy to read
as possible will be a great help as you write
your document and when you change it.  This sample
file shows how you can add comments to your own input
file.

Because printing is different from typewriting,
there are a number of things that you have to do
differently when preparing an input file than if
you were just typing the document directly.
Quotation marks like
       ``this''
have to be handled specially, as do quotes within
quotes:
       ``\,`this'            % \, separates the double and single quote.
        is what I just
        wrote, not  `that'\,''.

Dashes come in three sizes: an
       intra-word
dash, a medium dash for number ranges like
       1--2,
and a punctuation
       dash---like
this.

A sentence-ending space should be larger than the
space between words within a sentence.  You
sometimes have to type special commands in
conjunction with punctuation characters to get
this right, as in the following sentence.
       Gnats, gnus, etc.\ all  % `\ ' makes an inter-word space.
       begin with G\@.         % \@ marks end-of-sentence punctuation.
You should check the spaces after periods when
reading your output to make sure you haven't
forgotten any special cases.  Generating an
ellipsis
       \ldots\               % `\ ' is needed after `\ldots' because TeX
                             % ignores spaces after command names like \ldots
                             % made from \ + letters.
                             %
                             % Note how a `%' character causes TeX to ignore
                             % the end of the input line, so these blank lines
                             % do not start a new paragraph.
                             %
with the right spacing around the periods requires
a special command.

\LaTeX\ interprets some common characters as
commands, so you must type special commands to
generate them.  These characters include the
following:
       \$ \& \% \# \{ and \}.

In printing, text is usually emphasized with an
       \emph{italic}
type style.

\begin{em}
   A long segment of text can also be emphasized
   in this way.  Text within such a segment can be
   given \emph{additional} emphasis.
\end{em}

It is sometimes necessary to prevent \LaTeX\ from
breaking a line where it might otherwise do so.
This may be at a space, as between the ``Mr.''\ and
``Jones'' in
       ``Mr.~Jones'',        % ~ produces an unbreakable interword space.
or within a word---especially when the word is a
symbol like
       \mbox{\emph{itemnum}}
that makes little sense when hyphenated across
lines.

Footnotes\footnote{This is an example of a footnote.}
pose no problem.

\LaTeX\ is good at typesetting mathematical formulas
like
       \( x-3y + z = 7 \)
or
       \( a_{1} > x^{2n} + y^{2n} > x' \)
or
       \( \ip{A}{B} = \sum_{i} a_{i} b_{i} \).
The spaces you type in a formula are
ignored.  Remember that a letter like
       $x$                   % $ ... $  and  \( ... \)  are equivalent
is a formula when it denotes a mathematical
symbol, and it should be typed as one.

\section{Displayed Text}

Text is displayed by indenting it from the left
margin.  Quotations are commonly displayed.  There
are short quotations
\begin{quote}
   This is a short quotation.  It consists of a
   single paragraph of text.  See how it is formatted.
\end{quote}
and longer ones.
\begin{quotation}
   This is a longer quotation.  It consists of two
   paragraphs of text, neither of which are
   particularly interesting.

   This is the second paragraph of the quotation.  It
   is just as dull as the first paragraph.
\end{quotation}
Another frequently-displayed structure is a list.
The following is an example of an \emph{itemized}
list.
\begin{itemize}
   \item This is the first item of an itemized list.
         Each item in the list is marked with a ``tick''.
         You don't have to worry about what kind of tick
         mark is used.

   \item This is the second item of the list.  It
         contains another list nested inside it.  The inner
         list is an \emph{enumerated} list.
         \begin{enumerate}
            \item This is the first item of an enumerated
                  list that is nested within the itemized list.

            \item This is the second item of the inner list.
                  \LaTeX\ allows you to nest lists deeper than
                  you really should.
         \end{enumerate}
         This is the rest of the second item of the outer
         list.  It is no more interesting than any other
         part of the item.
   \item This is the third item of the list.
\end{itemize}
You can even display poetry.
\begin{verse}
   There is an environment
    for verse \\             % The \\ command separates lines
   Whose features some poets % within a stanza.
   will curse.

                             % One or more blank lines separate stanzas.

   For instead of making\\
   Them do \emph{all} line breaking, \\
   It allows them to put too many words on a line when they'd rather be
   forced to be terse.
\end{verse}

Mathematical formulas may also be displayed.  A
displayed formula
is
one-line long; multiline
formulas require special formatting instructions.
   \[  \ip{\Gamma}{\psi'} = x'' + y^{2} + z_{i}^{n}\]
Don't start a paragraph with a displayed equation,
nor make one a paragraph by itself.

\end{document}               % End of document.

\section{Bench Huge Section 441}
% synthetic bench section 441
% This is a sample LaTeX input file.  (Version of 12 August 2004.)
%
% A '%' character causes TeX to ignore all remaining text on the line,
% and is used for comments like this one.

\documentclass{article}      % Specifies the document class

                             % The preamble begins here.
\title{An Example Document}  % Declares the document's title.
\author{Leslie Lamport}      % Declares the author's name.
\date{January 21, 1994}      % Deleting this command produces today's date.

\newcommand{\ip}[2]{(#1, #2)}
                             % Defines \ip{arg1}{arg2} to mean
                             % (arg1, arg2).

%\newcommand{\ip}[2]{\langle #1 | #2\rangle}
                             % This is an alternative definition of
                             % \ip that is commented out.

\begin{document}             % End of preamble and beginning of text.

\maketitle                   % Produces the title.

This is an example input file.  Comparing it with
the output it generates can show you how to
produce a simple document of your own.

\section{Ordinary Text}      % Produces section heading.  Lower-level
                             % sections are begun with similar
                             % \subsection and \subsubsection commands.

The ends  of words and sentences are marked
  by   spaces. It  doesn't matter how many
spaces    you type; one is as good as 100.  The
end of   a line counts as a space.

One   or more   blank lines denote the  end
of  a paragraph.

Since any number of consecutive spaces are treated
like a single one, the formatting of the input
file makes no difference to
      \LaTeX,                % The \LaTeX command generates the LaTeX logo.
but it makes a difference to you.  When you use
\LaTeX, making your input file as easy to read
as possible will be a great help as you write
your document and when you change it.  This sample
file shows how you can add comments to your own input
file.

Because printing is different from typewriting,
there are a number of things that you have to do
differently when preparing an input file than if
you were just typing the document directly.
Quotation marks like
       ``this''
have to be handled specially, as do quotes within
quotes:
       ``\,`this'            % \, separates the double and single quote.
        is what I just
        wrote, not  `that'\,''.

Dashes come in three sizes: an
       intra-word
dash, a medium dash for number ranges like
       1--2,
and a punctuation
       dash---like
this.

A sentence-ending space should be larger than the
space between words within a sentence.  You
sometimes have to type special commands in
conjunction with punctuation characters to get
this right, as in the following sentence.
       Gnats, gnus, etc.\ all  % `\ ' makes an inter-word space.
       begin with G\@.         % \@ marks end-of-sentence punctuation.
You should check the spaces after periods when
reading your output to make sure you haven't
forgotten any special cases.  Generating an
ellipsis
       \ldots\               % `\ ' is needed after `\ldots' because TeX
                             % ignores spaces after command names like \ldots
                             % made from \ + letters.
                             %
                             % Note how a `%' character causes TeX to ignore
                             % the end of the input line, so these blank lines
                             % do not start a new paragraph.
                             %
with the right spacing around the periods requires
a special command.

\LaTeX\ interprets some common characters as
commands, so you must type special commands to
generate them.  These characters include the
following:
       \$ \& \% \# \{ and \}.

In printing, text is usually emphasized with an
       \emph{italic}
type style.

\begin{em}
   A long segment of text can also be emphasized
   in this way.  Text within such a segment can be
   given \emph{additional} emphasis.
\end{em}

It is sometimes necessary to prevent \LaTeX\ from
breaking a line where it might otherwise do so.
This may be at a space, as between the ``Mr.''\ and
``Jones'' in
       ``Mr.~Jones'',        % ~ produces an unbreakable interword space.
or within a word---especially when the word is a
symbol like
       \mbox{\emph{itemnum}}
that makes little sense when hyphenated across
lines.

Footnotes\footnote{This is an example of a footnote.}
pose no problem.

\LaTeX\ is good at typesetting mathematical formulas
like
       \( x-3y + z = 7 \)
or
       \( a_{1} > x^{2n} + y^{2n} > x' \)
or
       \( \ip{A}{B} = \sum_{i} a_{i} b_{i} \).
The spaces you type in a formula are
ignored.  Remember that a letter like
       $x$                   % $ ... $  and  \( ... \)  are equivalent
is a formula when it denotes a mathematical
symbol, and it should be typed as one.

\section{Displayed Text}

Text is displayed by indenting it from the left
margin.  Quotations are commonly displayed.  There
are short quotations
\begin{quote}
   This is a short quotation.  It consists of a
   single paragraph of text.  See how it is formatted.
\end{quote}
and longer ones.
\begin{quotation}
   This is a longer quotation.  It consists of two
   paragraphs of text, neither of which are
   particularly interesting.

   This is the second paragraph of the quotation.  It
   is just as dull as the first paragraph.
\end{quotation}
Another frequently-displayed structure is a list.
The following is an example of an \emph{itemized}
list.
\begin{itemize}
   \item This is the first item of an itemized list.
         Each item in the list is marked with a ``tick''.
         You don't have to worry about what kind of tick
         mark is used.

   \item This is the second item of the list.  It
         contains another list nested inside it.  The inner
         list is an \emph{enumerated} list.
         \begin{enumerate}
            \item This is the first item of an enumerated
                  list that is nested within the itemized list.

            \item This is the second item of the inner list.
                  \LaTeX\ allows you to nest lists deeper than
                  you really should.
         \end{enumerate}
         This is the rest of the second item of the outer
         list.  It is no more interesting than any other
         part of the item.
   \item This is the third item of the list.
\end{itemize}
You can even display poetry.
\begin{verse}
   There is an environment
    for verse \\             % The \\ command separates lines
   Whose features some poets % within a stanza.
   will curse.

                             % One or more blank lines separate stanzas.

   For instead of making\\
   Them do \emph{all} line breaking, \\
   It allows them to put too many words on a line when they'd rather be
   forced to be terse.
\end{verse}

Mathematical formulas may also be displayed.  A
displayed formula
is
one-line long; multiline
formulas require special formatting instructions.
   \[  \ip{\Gamma}{\psi'} = x'' + y^{2} + z_{i}^{n}\]
Don't start a paragraph with a displayed equation,
nor make one a paragraph by itself.

\end{document}               % End of document.

\section{Bench Huge Section 442}
% synthetic bench section 442
% This is a sample LaTeX input file.  (Version of 12 August 2004.)
%
% A '%' character causes TeX to ignore all remaining text on the line,
% and is used for comments like this one.

\documentclass{article}      % Specifies the document class

                             % The preamble begins here.
\title{An Example Document}  % Declares the document's title.
\author{Leslie Lamport}      % Declares the author's name.
\date{January 21, 1994}      % Deleting this command produces today's date.

\newcommand{\ip}[2]{(#1, #2)}
                             % Defines \ip{arg1}{arg2} to mean
                             % (arg1, arg2).

%\newcommand{\ip}[2]{\langle #1 | #2\rangle}
                             % This is an alternative definition of
                             % \ip that is commented out.

\begin{document}             % End of preamble and beginning of text.

\maketitle                   % Produces the title.

This is an example input file.  Comparing it with
the output it generates can show you how to
produce a simple document of your own.

\section{Ordinary Text}      % Produces section heading.  Lower-level
                             % sections are begun with similar
                             % \subsection and \subsubsection commands.

The ends  of words and sentences are marked
  by   spaces. It  doesn't matter how many
spaces    you type; one is as good as 100.  The
end of   a line counts as a space.

One   or more   blank lines denote the  end
of  a paragraph.

Since any number of consecutive spaces are treated
like a single one, the formatting of the input
file makes no difference to
      \LaTeX,                % The \LaTeX command generates the LaTeX logo.
but it makes a difference to you.  When you use
\LaTeX, making your input file as easy to read
as possible will be a great help as you write
your document and when you change it.  This sample
file shows how you can add comments to your own input
file.

Because printing is different from typewriting,
there are a number of things that you have to do
differently when preparing an input file than if
you were just typing the document directly.
Quotation marks like
       ``this''
have to be handled specially, as do quotes within
quotes:
       ``\,`this'            % \, separates the double and single quote.
        is what I just
        wrote, not  `that'\,''.

Dashes come in three sizes: an
       intra-word
dash, a medium dash for number ranges like
       1--2,
and a punctuation
       dash---like
this.

A sentence-ending space should be larger than the
space between words within a sentence.  You
sometimes have to type special commands in
conjunction with punctuation characters to get
this right, as in the following sentence.
       Gnats, gnus, etc.\ all  % `\ ' makes an inter-word space.
       begin with G\@.         % \@ marks end-of-sentence punctuation.
You should check the spaces after periods when
reading your output to make sure you haven't
forgotten any special cases.  Generating an
ellipsis
       \ldots\               % `\ ' is needed after `\ldots' because TeX
                             % ignores spaces after command names like \ldots
                             % made from \ + letters.
                             %
                             % Note how a `%' character causes TeX to ignore
                             % the end of the input line, so these blank lines
                             % do not start a new paragraph.
                             %
with the right spacing around the periods requires
a special command.

\LaTeX\ interprets some common characters as
commands, so you must type special commands to
generate them.  These characters include the
following:
       \$ \& \% \# \{ and \}.

In printing, text is usually emphasized with an
       \emph{italic}
type style.

\begin{em}
   A long segment of text can also be emphasized
   in this way.  Text within such a segment can be
   given \emph{additional} emphasis.
\end{em}

It is sometimes necessary to prevent \LaTeX\ from
breaking a line where it might otherwise do so.
This may be at a space, as between the ``Mr.''\ and
``Jones'' in
       ``Mr.~Jones'',        % ~ produces an unbreakable interword space.
or within a word---especially when the word is a
symbol like
       \mbox{\emph{itemnum}}
that makes little sense when hyphenated across
lines.

Footnotes\footnote{This is an example of a footnote.}
pose no problem.

\LaTeX\ is good at typesetting mathematical formulas
like
       \( x-3y + z = 7 \)
or
       \( a_{1} > x^{2n} + y^{2n} > x' \)
or
       \( \ip{A}{B} = \sum_{i} a_{i} b_{i} \).
The spaces you type in a formula are
ignored.  Remember that a letter like
       $x$                   % $ ... $  and  \( ... \)  are equivalent
is a formula when it denotes a mathematical
symbol, and it should be typed as one.

\section{Displayed Text}

Text is displayed by indenting it from the left
margin.  Quotations are commonly displayed.  There
are short quotations
\begin{quote}
   This is a short quotation.  It consists of a
   single paragraph of text.  See how it is formatted.
\end{quote}
and longer ones.
\begin{quotation}
   This is a longer quotation.  It consists of two
   paragraphs of text, neither of which are
   particularly interesting.

   This is the second paragraph of the quotation.  It
   is just as dull as the first paragraph.
\end{quotation}
Another frequently-displayed structure is a list.
The following is an example of an \emph{itemized}
list.
\begin{itemize}
   \item This is the first item of an itemized list.
         Each item in the list is marked with a ``tick''.
         You don't have to worry about what kind of tick
         mark is used.

   \item This is the second item of the list.  It
         contains another list nested inside it.  The inner
         list is an \emph{enumerated} list.
         \begin{enumerate}
            \item This is the first item of an enumerated
                  list that is nested within the itemized list.

            \item This is the second item of the inner list.
                  \LaTeX\ allows you to nest lists deeper than
                  you really should.
         \end{enumerate}
         This is the rest of the second item of the outer
         list.  It is no more interesting than any other
         part of the item.
   \item This is the third item of the list.
\end{itemize}
You can even display poetry.
\begin{verse}
   There is an environment
    for verse \\             % The \\ command separates lines
   Whose features some poets % within a stanza.
   will curse.

                             % One or more blank lines separate stanzas.

   For instead of making\\
   Them do \emph{all} line breaking, \\
   It allows them to put too many words on a line when they'd rather be
   forced to be terse.
\end{verse}

Mathematical formulas may also be displayed.  A
displayed formula
is
one-line long; multiline
formulas require special formatting instructions.
   \[  \ip{\Gamma}{\psi'} = x'' + y^{2} + z_{i}^{n}\]
Don't start a paragraph with a displayed equation,
nor make one a paragraph by itself.

\end{document}               % End of document.

\section{Bench Huge Section 443}
% synthetic bench section 443
% This is a sample LaTeX input file.  (Version of 12 August 2004.)
%
% A '%' character causes TeX to ignore all remaining text on the line,
% and is used for comments like this one.

\documentclass{article}      % Specifies the document class

                             % The preamble begins here.
\title{An Example Document}  % Declares the document's title.
\author{Leslie Lamport}      % Declares the author's name.
\date{January 21, 1994}      % Deleting this command produces today's date.

\newcommand{\ip}[2]{(#1, #2)}
                             % Defines \ip{arg1}{arg2} to mean
                             % (arg1, arg2).

%\newcommand{\ip}[2]{\langle #1 | #2\rangle}
                             % This is an alternative definition of
                             % \ip that is commented out.

\begin{document}             % End of preamble and beginning of text.

\maketitle                   % Produces the title.

This is an example input file.  Comparing it with
the output it generates can show you how to
produce a simple document of your own.

\section{Ordinary Text}      % Produces section heading.  Lower-level
                             % sections are begun with similar
                             % \subsection and \subsubsection commands.

The ends  of words and sentences are marked
  by   spaces. It  doesn't matter how many
spaces    you type; one is as good as 100.  The
end of   a line counts as a space.

One   or more   blank lines denote the  end
of  a paragraph.

Since any number of consecutive spaces are treated
like a single one, the formatting of the input
file makes no difference to
      \LaTeX,                % The \LaTeX command generates the LaTeX logo.
but it makes a difference to you.  When you use
\LaTeX, making your input file as easy to read
as possible will be a great help as you write
your document and when you change it.  This sample
file shows how you can add comments to your own input
file.

Because printing is different from typewriting,
there are a number of things that you have to do
differently when preparing an input file than if
you were just typing the document directly.
Quotation marks like
       ``this''
have to be handled specially, as do quotes within
quotes:
       ``\,`this'            % \, separates the double and single quote.
        is what I just
        wrote, not  `that'\,''.

Dashes come in three sizes: an
       intra-word
dash, a medium dash for number ranges like
       1--2,
and a punctuation
       dash---like
this.

A sentence-ending space should be larger than the
space between words within a sentence.  You
sometimes have to type special commands in
conjunction with punctuation characters to get
this right, as in the following sentence.
       Gnats, gnus, etc.\ all  % `\ ' makes an inter-word space.
       begin with G\@.         % \@ marks end-of-sentence punctuation.
You should check the spaces after periods when
reading your output to make sure you haven't
forgotten any special cases.  Generating an
ellipsis
       \ldots\               % `\ ' is needed after `\ldots' because TeX
                             % ignores spaces after command names like \ldots
                             % made from \ + letters.
                             %
                             % Note how a `%' character causes TeX to ignore
                             % the end of the input line, so these blank lines
                             % do not start a new paragraph.
                             %
with the right spacing around the periods requires
a special command.

\LaTeX\ interprets some common characters as
commands, so you must type special commands to
generate them.  These characters include the
following:
       \$ \& \% \# \{ and \}.

In printing, text is usually emphasized with an
       \emph{italic}
type style.

\begin{em}
   A long segment of text can also be emphasized
   in this way.  Text within such a segment can be
   given \emph{additional} emphasis.
\end{em}

It is sometimes necessary to prevent \LaTeX\ from
breaking a line where it might otherwise do so.
This may be at a space, as between the ``Mr.''\ and
``Jones'' in
       ``Mr.~Jones'',        % ~ produces an unbreakable interword space.
or within a word---especially when the word is a
symbol like
       \mbox{\emph{itemnum}}
that makes little sense when hyphenated across
lines.

Footnotes\footnote{This is an example of a footnote.}
pose no problem.

\LaTeX\ is good at typesetting mathematical formulas
like
       \( x-3y + z = 7 \)
or
       \( a_{1} > x^{2n} + y^{2n} > x' \)
or
       \( \ip{A}{B} = \sum_{i} a_{i} b_{i} \).
The spaces you type in a formula are
ignored.  Remember that a letter like
       $x$                   % $ ... $  and  \( ... \)  are equivalent
is a formula when it denotes a mathematical
symbol, and it should be typed as one.

\section{Displayed Text}

Text is displayed by indenting it from the left
margin.  Quotations are commonly displayed.  There
are short quotations
\begin{quote}
   This is a short quotation.  It consists of a
   single paragraph of text.  See how it is formatted.
\end{quote}
and longer ones.
\begin{quotation}
   This is a longer quotation.  It consists of two
   paragraphs of text, neither of which are
   particularly interesting.

   This is the second paragraph of the quotation.  It
   is just as dull as the first paragraph.
\end{quotation}
Another frequently-displayed structure is a list.
The following is an example of an \emph{itemized}
list.
\begin{itemize}
   \item This is the first item of an itemized list.
         Each item in the list is marked with a ``tick''.
         You don't have to worry about what kind of tick
         mark is used.

   \item This is the second item of the list.  It
         contains another list nested inside it.  The inner
         list is an \emph{enumerated} list.
         \begin{enumerate}
            \item This is the first item of an enumerated
                  list that is nested within the itemized list.

            \item This is the second item of the inner list.
                  \LaTeX\ allows you to nest lists deeper than
                  you really should.
         \end{enumerate}
         This is the rest of the second item of the outer
         list.  It is no more interesting than any other
         part of the item.
   \item This is the third item of the list.
\end{itemize}
You can even display poetry.
\begin{verse}
   There is an environment
    for verse \\             % The \\ command separates lines
   Whose features some poets % within a stanza.
   will curse.

                             % One or more blank lines separate stanzas.

   For instead of making\\
   Them do \emph{all} line breaking, \\
   It allows them to put too many words on a line when they'd rather be
   forced to be terse.
\end{verse}

Mathematical formulas may also be displayed.  A
displayed formula
is
one-line long; multiline
formulas require special formatting instructions.
   \[  \ip{\Gamma}{\psi'} = x'' + y^{2} + z_{i}^{n}\]
Don't start a paragraph with a displayed equation,
nor make one a paragraph by itself.

\end{document}               % End of document.

\section{Bench Huge Section 444}
% synthetic bench section 444
% This is a sample LaTeX input file.  (Version of 12 August 2004.)
%
% A '%' character causes TeX to ignore all remaining text on the line,
% and is used for comments like this one.

\documentclass{article}      % Specifies the document class

                             % The preamble begins here.
\title{An Example Document}  % Declares the document's title.
\author{Leslie Lamport}      % Declares the author's name.
\date{January 21, 1994}      % Deleting this command produces today's date.

\newcommand{\ip}[2]{(#1, #2)}
                             % Defines \ip{arg1}{arg2} to mean
                             % (arg1, arg2).

%\newcommand{\ip}[2]{\langle #1 | #2\rangle}
                             % This is an alternative definition of
                             % \ip that is commented out.

\begin{document}             % End of preamble and beginning of text.

\maketitle                   % Produces the title.

This is an example input file.  Comparing it with
the output it generates can show you how to
produce a simple document of your own.

\section{Ordinary Text}      % Produces section heading.  Lower-level
                             % sections are begun with similar
                             % \subsection and \subsubsection commands.

The ends  of words and sentences are marked
  by   spaces. It  doesn't matter how many
spaces    you type; one is as good as 100.  The
end of   a line counts as a space.

One   or more   blank lines denote the  end
of  a paragraph.

Since any number of consecutive spaces are treated
like a single one, the formatting of the input
file makes no difference to
      \LaTeX,                % The \LaTeX command generates the LaTeX logo.
but it makes a difference to you.  When you use
\LaTeX, making your input file as easy to read
as possible will be a great help as you write
your document and when you change it.  This sample
file shows how you can add comments to your own input
file.

Because printing is different from typewriting,
there are a number of things that you have to do
differently when preparing an input file than if
you were just typing the document directly.
Quotation marks like
       ``this''
have to be handled specially, as do quotes within
quotes:
       ``\,`this'            % \, separates the double and single quote.
        is what I just
        wrote, not  `that'\,''.

Dashes come in three sizes: an
       intra-word
dash, a medium dash for number ranges like
       1--2,
and a punctuation
       dash---like
this.

A sentence-ending space should be larger than the
space between words within a sentence.  You
sometimes have to type special commands in
conjunction with punctuation characters to get
this right, as in the following sentence.
       Gnats, gnus, etc.\ all  % `\ ' makes an inter-word space.
       begin with G\@.         % \@ marks end-of-sentence punctuation.
You should check the spaces after periods when
reading your output to make sure you haven't
forgotten any special cases.  Generating an
ellipsis
       \ldots\               % `\ ' is needed after `\ldots' because TeX
                             % ignores spaces after command names like \ldots
                             % made from \ + letters.
                             %
                             % Note how a `%' character causes TeX to ignore
                             % the end of the input line, so these blank lines
                             % do not start a new paragraph.
                             %
with the right spacing around the periods requires
a special command.

\LaTeX\ interprets some common characters as
commands, so you must type special commands to
generate them.  These characters include the
following:
       \$ \& \% \# \{ and \}.

In printing, text is usually emphasized with an
       \emph{italic}
type style.

\begin{em}
   A long segment of text can also be emphasized
   in this way.  Text within such a segment can be
   given \emph{additional} emphasis.
\end{em}

It is sometimes necessary to prevent \LaTeX\ from
breaking a line where it might otherwise do so.
This may be at a space, as between the ``Mr.''\ and
``Jones'' in
       ``Mr.~Jones'',        % ~ produces an unbreakable interword space.
or within a word---especially when the word is a
symbol like
       \mbox{\emph{itemnum}}
that makes little sense when hyphenated across
lines.

Footnotes\footnote{This is an example of a footnote.}
pose no problem.

\LaTeX\ is good at typesetting mathematical formulas
like
       \( x-3y + z = 7 \)
or
       \( a_{1} > x^{2n} + y^{2n} > x' \)
or
       \( \ip{A}{B} = \sum_{i} a_{i} b_{i} \).
The spaces you type in a formula are
ignored.  Remember that a letter like
       $x$                   % $ ... $  and  \( ... \)  are equivalent
is a formula when it denotes a mathematical
symbol, and it should be typed as one.

\section{Displayed Text}

Text is displayed by indenting it from the left
margin.  Quotations are commonly displayed.  There
are short quotations
\begin{quote}
   This is a short quotation.  It consists of a
   single paragraph of text.  See how it is formatted.
\end{quote}
and longer ones.
\begin{quotation}
   This is a longer quotation.  It consists of two
   paragraphs of text, neither of which are
   particularly interesting.

   This is the second paragraph of the quotation.  It
   is just as dull as the first paragraph.
\end{quotation}
Another frequently-displayed structure is a list.
The following is an example of an \emph{itemized}
list.
\begin{itemize}
   \item This is the first item of an itemized list.
         Each item in the list is marked with a ``tick''.
         You don't have to worry about what kind of tick
         mark is used.

   \item This is the second item of the list.  It
         contains another list nested inside it.  The inner
         list is an \emph{enumerated} list.
         \begin{enumerate}
            \item This is the first item of an enumerated
                  list that is nested within the itemized list.

            \item This is the second item of the inner list.
                  \LaTeX\ allows you to nest lists deeper than
                  you really should.
         \end{enumerate}
         This is the rest of the second item of the outer
         list.  It is no more interesting than any other
         part of the item.
   \item This is the third item of the list.
\end{itemize}
You can even display poetry.
\begin{verse}
   There is an environment
    for verse \\             % The \\ command separates lines
   Whose features some poets % within a stanza.
   will curse.

                             % One or more blank lines separate stanzas.

   For instead of making\\
   Them do \emph{all} line breaking, \\
   It allows them to put too many words on a line when they'd rather be
   forced to be terse.
\end{verse}

Mathematical formulas may also be displayed.  A
displayed formula
is
one-line long; multiline
formulas require special formatting instructions.
   \[  \ip{\Gamma}{\psi'} = x'' + y^{2} + z_{i}^{n}\]
Don't start a paragraph with a displayed equation,
nor make one a paragraph by itself.

\end{document}               % End of document.

\section{Bench Huge Section 445}
% synthetic bench section 445
% This is a sample LaTeX input file.  (Version of 12 August 2004.)
%
% A '%' character causes TeX to ignore all remaining text on the line,
% and is used for comments like this one.

\documentclass{article}      % Specifies the document class

                             % The preamble begins here.
\title{An Example Document}  % Declares the document's title.
\author{Leslie Lamport}      % Declares the author's name.
\date{January 21, 1994}      % Deleting this command produces today's date.

\newcommand{\ip}[2]{(#1, #2)}
                             % Defines \ip{arg1}{arg2} to mean
                             % (arg1, arg2).

%\newcommand{\ip}[2]{\langle #1 | #2\rangle}
                             % This is an alternative definition of
                             % \ip that is commented out.

\begin{document}             % End of preamble and beginning of text.

\maketitle                   % Produces the title.

This is an example input file.  Comparing it with
the output it generates can show you how to
produce a simple document of your own.

\section{Ordinary Text}      % Produces section heading.  Lower-level
                             % sections are begun with similar
                             % \subsection and \subsubsection commands.

The ends  of words and sentences are marked
  by   spaces. It  doesn't matter how many
spaces    you type; one is as good as 100.  The
end of   a line counts as a space.

One   or more   blank lines denote the  end
of  a paragraph.

Since any number of consecutive spaces are treated
like a single one, the formatting of the input
file makes no difference to
      \LaTeX,                % The \LaTeX command generates the LaTeX logo.
but it makes a difference to you.  When you use
\LaTeX, making your input file as easy to read
as possible will be a great help as you write
your document and when you change it.  This sample
file shows how you can add comments to your own input
file.

Because printing is different from typewriting,
there are a number of things that you have to do
differently when preparing an input file than if
you were just typing the document directly.
Quotation marks like
       ``this''
have to be handled specially, as do quotes within
quotes:
       ``\,`this'            % \, separates the double and single quote.
        is what I just
        wrote, not  `that'\,''.

Dashes come in three sizes: an
       intra-word
dash, a medium dash for number ranges like
       1--2,
and a punctuation
       dash---like
this.

A sentence-ending space should be larger than the
space between words within a sentence.  You
sometimes have to type special commands in
conjunction with punctuation characters to get
this right, as in the following sentence.
       Gnats, gnus, etc.\ all  % `\ ' makes an inter-word space.
       begin with G\@.         % \@ marks end-of-sentence punctuation.
You should check the spaces after periods when
reading your output to make sure you haven't
forgotten any special cases.  Generating an
ellipsis
       \ldots\               % `\ ' is needed after `\ldots' because TeX
                             % ignores spaces after command names like \ldots
                             % made from \ + letters.
                             %
                             % Note how a `%' character causes TeX to ignore
                             % the end of the input line, so these blank lines
                             % do not start a new paragraph.
                             %
with the right spacing around the periods requires
a special command.

\LaTeX\ interprets some common characters as
commands, so you must type special commands to
generate them.  These characters include the
following:
       \$ \& \% \# \{ and \}.

In printing, text is usually emphasized with an
       \emph{italic}
type style.

\begin{em}
   A long segment of text can also be emphasized
   in this way.  Text within such a segment can be
   given \emph{additional} emphasis.
\end{em}

It is sometimes necessary to prevent \LaTeX\ from
breaking a line where it might otherwise do so.
This may be at a space, as between the ``Mr.''\ and
``Jones'' in
       ``Mr.~Jones'',        % ~ produces an unbreakable interword space.
or within a word---especially when the word is a
symbol like
       \mbox{\emph{itemnum}}
that makes little sense when hyphenated across
lines.

Footnotes\footnote{This is an example of a footnote.}
pose no problem.

\LaTeX\ is good at typesetting mathematical formulas
like
       \( x-3y + z = 7 \)
or
       \( a_{1} > x^{2n} + y^{2n} > x' \)
or
       \( \ip{A}{B} = \sum_{i} a_{i} b_{i} \).
The spaces you type in a formula are
ignored.  Remember that a letter like
       $x$                   % $ ... $  and  \( ... \)  are equivalent
is a formula when it denotes a mathematical
symbol, and it should be typed as one.

\section{Displayed Text}

Text is displayed by indenting it from the left
margin.  Quotations are commonly displayed.  There
are short quotations
\begin{quote}
   This is a short quotation.  It consists of a
   single paragraph of text.  See how it is formatted.
\end{quote}
and longer ones.
\begin{quotation}
   This is a longer quotation.  It consists of two
   paragraphs of text, neither of which are
   particularly interesting.

   This is the second paragraph of the quotation.  It
   is just as dull as the first paragraph.
\end{quotation}
Another frequently-displayed structure is a list.
The following is an example of an \emph{itemized}
list.
\begin{itemize}
   \item This is the first item of an itemized list.
         Each item in the list is marked with a ``tick''.
         You don't have to worry about what kind of tick
         mark is used.

   \item This is the second item of the list.  It
         contains another list nested inside it.  The inner
         list is an \emph{enumerated} list.
         \begin{enumerate}
            \item This is the first item of an enumerated
                  list that is nested within the itemized list.

            \item This is the second item of the inner list.
                  \LaTeX\ allows you to nest lists deeper than
                  you really should.
         \end{enumerate}
         This is the rest of the second item of the outer
         list.  It is no more interesting than any other
         part of the item.
   \item This is the third item of the list.
\end{itemize}
You can even display poetry.
\begin{verse}
   There is an environment
    for verse \\             % The \\ command separates lines
   Whose features some poets % within a stanza.
   will curse.

                             % One or more blank lines separate stanzas.

   For instead of making\\
   Them do \emph{all} line breaking, \\
   It allows them to put too many words on a line when they'd rather be
   forced to be terse.
\end{verse}

Mathematical formulas may also be displayed.  A
displayed formula
is
one-line long; multiline
formulas require special formatting instructions.
   \[  \ip{\Gamma}{\psi'} = x'' + y^{2} + z_{i}^{n}\]
Don't start a paragraph with a displayed equation,
nor make one a paragraph by itself.

\end{document}               % End of document.

\section{Bench Huge Section 446}
% synthetic bench section 446
% This is a sample LaTeX input file.  (Version of 12 August 2004.)
%
% A '%' character causes TeX to ignore all remaining text on the line,
% and is used for comments like this one.

\documentclass{article}      % Specifies the document class

                             % The preamble begins here.
\title{An Example Document}  % Declares the document's title.
\author{Leslie Lamport}      % Declares the author's name.
\date{January 21, 1994}      % Deleting this command produces today's date.

\newcommand{\ip}[2]{(#1, #2)}
                             % Defines \ip{arg1}{arg2} to mean
                             % (arg1, arg2).

%\newcommand{\ip}[2]{\langle #1 | #2\rangle}
                             % This is an alternative definition of
                             % \ip that is commented out.

\begin{document}             % End of preamble and beginning of text.

\maketitle                   % Produces the title.

This is an example input file.  Comparing it with
the output it generates can show you how to
produce a simple document of your own.

\section{Ordinary Text}      % Produces section heading.  Lower-level
                             % sections are begun with similar
                             % \subsection and \subsubsection commands.

The ends  of words and sentences are marked
  by   spaces. It  doesn't matter how many
spaces    you type; one is as good as 100.  The
end of   a line counts as a space.

One   or more   blank lines denote the  end
of  a paragraph.

Since any number of consecutive spaces are treated
like a single one, the formatting of the input
file makes no difference to
      \LaTeX,                % The \LaTeX command generates the LaTeX logo.
but it makes a difference to you.  When you use
\LaTeX, making your input file as easy to read
as possible will be a great help as you write
your document and when you change it.  This sample
file shows how you can add comments to your own input
file.

Because printing is different from typewriting,
there are a number of things that you have to do
differently when preparing an input file than if
you were just typing the document directly.
Quotation marks like
       ``this''
have to be handled specially, as do quotes within
quotes:
       ``\,`this'            % \, separates the double and single quote.
        is what I just
        wrote, not  `that'\,''.

Dashes come in three sizes: an
       intra-word
dash, a medium dash for number ranges like
       1--2,
and a punctuation
       dash---like
this.

A sentence-ending space should be larger than the
space between words within a sentence.  You
sometimes have to type special commands in
conjunction with punctuation characters to get
this right, as in the following sentence.
       Gnats, gnus, etc.\ all  % `\ ' makes an inter-word space.
       begin with G\@.         % \@ marks end-of-sentence punctuation.
You should check the spaces after periods when
reading your output to make sure you haven't
forgotten any special cases.  Generating an
ellipsis
       \ldots\               % `\ ' is needed after `\ldots' because TeX
                             % ignores spaces after command names like \ldots
                             % made from \ + letters.
                             %
                             % Note how a `%' character causes TeX to ignore
                             % the end of the input line, so these blank lines
                             % do not start a new paragraph.
                             %
with the right spacing around the periods requires
a special command.

\LaTeX\ interprets some common characters as
commands, so you must type special commands to
generate them.  These characters include the
following:
       \$ \& \% \# \{ and \}.

In printing, text is usually emphasized with an
       \emph{italic}
type style.

\begin{em}
   A long segment of text can also be emphasized
   in this way.  Text within such a segment can be
   given \emph{additional} emphasis.
\end{em}

It is sometimes necessary to prevent \LaTeX\ from
breaking a line where it might otherwise do so.
This may be at a space, as between the ``Mr.''\ and
``Jones'' in
       ``Mr.~Jones'',        % ~ produces an unbreakable interword space.
or within a word---especially when the word is a
symbol like
       \mbox{\emph{itemnum}}
that makes little sense when hyphenated across
lines.

Footnotes\footnote{This is an example of a footnote.}
pose no problem.

\LaTeX\ is good at typesetting mathematical formulas
like
       \( x-3y + z = 7 \)
or
       \( a_{1} > x^{2n} + y^{2n} > x' \)
or
       \( \ip{A}{B} = \sum_{i} a_{i} b_{i} \).
The spaces you type in a formula are
ignored.  Remember that a letter like
       $x$                   % $ ... $  and  \( ... \)  are equivalent
is a formula when it denotes a mathematical
symbol, and it should be typed as one.

\section{Displayed Text}

Text is displayed by indenting it from the left
margin.  Quotations are commonly displayed.  There
are short quotations
\begin{quote}
   This is a short quotation.  It consists of a
   single paragraph of text.  See how it is formatted.
\end{quote}
and longer ones.
\begin{quotation}
   This is a longer quotation.  It consists of two
   paragraphs of text, neither of which are
   particularly interesting.

   This is the second paragraph of the quotation.  It
   is just as dull as the first paragraph.
\end{quotation}
Another frequently-displayed structure is a list.
The following is an example of an \emph{itemized}
list.
\begin{itemize}
   \item This is the first item of an itemized list.
         Each item in the list is marked with a ``tick''.
         You don't have to worry about what kind of tick
         mark is used.

   \item This is the second item of the list.  It
         contains another list nested inside it.  The inner
         list is an \emph{enumerated} list.
         \begin{enumerate}
            \item This is the first item of an enumerated
                  list that is nested within the itemized list.

            \item This is the second item of the inner list.
                  \LaTeX\ allows you to nest lists deeper than
                  you really should.
         \end{enumerate}
         This is the rest of the second item of the outer
         list.  It is no more interesting than any other
         part of the item.
   \item This is the third item of the list.
\end{itemize}
You can even display poetry.
\begin{verse}
   There is an environment
    for verse \\             % The \\ command separates lines
   Whose features some poets % within a stanza.
   will curse.

                             % One or more blank lines separate stanzas.

   For instead of making\\
   Them do \emph{all} line breaking, \\
   It allows them to put too many words on a line when they'd rather be
   forced to be terse.
\end{verse}

Mathematical formulas may also be displayed.  A
displayed formula
is
one-line long; multiline
formulas require special formatting instructions.
   \[  \ip{\Gamma}{\psi'} = x'' + y^{2} + z_{i}^{n}\]
Don't start a paragraph with a displayed equation,
nor make one a paragraph by itself.

\end{document}               % End of document.

\section{Bench Huge Section 447}
% synthetic bench section 447
% This is a sample LaTeX input file.  (Version of 12 August 2004.)
%
% A '%' character causes TeX to ignore all remaining text on the line,
% and is used for comments like this one.

\documentclass{article}      % Specifies the document class

                             % The preamble begins here.
\title{An Example Document}  % Declares the document's title.
\author{Leslie Lamport}      % Declares the author's name.
\date{January 21, 1994}      % Deleting this command produces today's date.

\newcommand{\ip}[2]{(#1, #2)}
                             % Defines \ip{arg1}{arg2} to mean
                             % (arg1, arg2).

%\newcommand{\ip}[2]{\langle #1 | #2\rangle}
                             % This is an alternative definition of
                             % \ip that is commented out.

\begin{document}             % End of preamble and beginning of text.

\maketitle                   % Produces the title.

This is an example input file.  Comparing it with
the output it generates can show you how to
produce a simple document of your own.

\section{Ordinary Text}      % Produces section heading.  Lower-level
                             % sections are begun with similar
                             % \subsection and \subsubsection commands.

The ends  of words and sentences are marked
  by   spaces. It  doesn't matter how many
spaces    you type; one is as good as 100.  The
end of   a line counts as a space.

One   or more   blank lines denote the  end
of  a paragraph.

Since any number of consecutive spaces are treated
like a single one, the formatting of the input
file makes no difference to
      \LaTeX,                % The \LaTeX command generates the LaTeX logo.
but it makes a difference to you.  When you use
\LaTeX, making your input file as easy to read
as possible will be a great help as you write
your document and when you change it.  This sample
file shows how you can add comments to your own input
file.

Because printing is different from typewriting,
there are a number of things that you have to do
differently when preparing an input file than if
you were just typing the document directly.
Quotation marks like
       ``this''
have to be handled specially, as do quotes within
quotes:
       ``\,`this'            % \, separates the double and single quote.
        is what I just
        wrote, not  `that'\,''.

Dashes come in three sizes: an
       intra-word
dash, a medium dash for number ranges like
       1--2,
and a punctuation
       dash---like
this.

A sentence-ending space should be larger than the
space between words within a sentence.  You
sometimes have to type special commands in
conjunction with punctuation characters to get
this right, as in the following sentence.
       Gnats, gnus, etc.\ all  % `\ ' makes an inter-word space.
       begin with G\@.         % \@ marks end-of-sentence punctuation.
You should check the spaces after periods when
reading your output to make sure you haven't
forgotten any special cases.  Generating an
ellipsis
       \ldots\               % `\ ' is needed after `\ldots' because TeX
                             % ignores spaces after command names like \ldots
                             % made from \ + letters.
                             %
                             % Note how a `%' character causes TeX to ignore
                             % the end of the input line, so these blank lines
                             % do not start a new paragraph.
                             %
with the right spacing around the periods requires
a special command.

\LaTeX\ interprets some common characters as
commands, so you must type special commands to
generate them.  These characters include the
following:
       \$ \& \% \# \{ and \}.

In printing, text is usually emphasized with an
       \emph{italic}
type style.

\begin{em}
   A long segment of text can also be emphasized
   in this way.  Text within such a segment can be
   given \emph{additional} emphasis.
\end{em}

It is sometimes necessary to prevent \LaTeX\ from
breaking a line where it might otherwise do so.
This may be at a space, as between the ``Mr.''\ and
``Jones'' in
       ``Mr.~Jones'',        % ~ produces an unbreakable interword space.
or within a word---especially when the word is a
symbol like
       \mbox{\emph{itemnum}}
that makes little sense when hyphenated across
lines.

Footnotes\footnote{This is an example of a footnote.}
pose no problem.

\LaTeX\ is good at typesetting mathematical formulas
like
       \( x-3y + z = 7 \)
or
       \( a_{1} > x^{2n} + y^{2n} > x' \)
or
       \( \ip{A}{B} = \sum_{i} a_{i} b_{i} \).
The spaces you type in a formula are
ignored.  Remember that a letter like
       $x$                   % $ ... $  and  \( ... \)  are equivalent
is a formula when it denotes a mathematical
symbol, and it should be typed as one.

\section{Displayed Text}

Text is displayed by indenting it from the left
margin.  Quotations are commonly displayed.  There
are short quotations
\begin{quote}
   This is a short quotation.  It consists of a
   single paragraph of text.  See how it is formatted.
\end{quote}
and longer ones.
\begin{quotation}
   This is a longer quotation.  It consists of two
   paragraphs of text, neither of which are
   particularly interesting.

   This is the second paragraph of the quotation.  It
   is just as dull as the first paragraph.
\end{quotation}
Another frequently-displayed structure is a list.
The following is an example of an \emph{itemized}
list.
\begin{itemize}
   \item This is the first item of an itemized list.
         Each item in the list is marked with a ``tick''.
         You don't have to worry about what kind of tick
         mark is used.

   \item This is the second item of the list.  It
         contains another list nested inside it.  The inner
         list is an \emph{enumerated} list.
         \begin{enumerate}
            \item This is the first item of an enumerated
                  list that is nested within the itemized list.

            \item This is the second item of the inner list.
                  \LaTeX\ allows you to nest lists deeper than
                  you really should.
         \end{enumerate}
         This is the rest of the second item of the outer
         list.  It is no more interesting than any other
         part of the item.
   \item This is the third item of the list.
\end{itemize}
You can even display poetry.
\begin{verse}
   There is an environment
    for verse \\             % The \\ command separates lines
   Whose features some poets % within a stanza.
   will curse.

                             % One or more blank lines separate stanzas.

   For instead of making\\
   Them do \emph{all} line breaking, \\
   It allows them to put too many words on a line when they'd rather be
   forced to be terse.
\end{verse}

Mathematical formulas may also be displayed.  A
displayed formula
is
one-line long; multiline
formulas require special formatting instructions.
   \[  \ip{\Gamma}{\psi'} = x'' + y^{2} + z_{i}^{n}\]
Don't start a paragraph with a displayed equation,
nor make one a paragraph by itself.

\end{document}               % End of document.

\section{Bench Huge Section 448}
% synthetic bench section 448
% This is a sample LaTeX input file.  (Version of 12 August 2004.)
%
% A '%' character causes TeX to ignore all remaining text on the line,
% and is used for comments like this one.

\documentclass{article}      % Specifies the document class

                             % The preamble begins here.
\title{An Example Document}  % Declares the document's title.
\author{Leslie Lamport}      % Declares the author's name.
\date{January 21, 1994}      % Deleting this command produces today's date.

\newcommand{\ip}[2]{(#1, #2)}
                             % Defines \ip{arg1}{arg2} to mean
                             % (arg1, arg2).

%\newcommand{\ip}[2]{\langle #1 | #2\rangle}
                             % This is an alternative definition of
                             % \ip that is commented out.

\begin{document}             % End of preamble and beginning of text.

\maketitle                   % Produces the title.

This is an example input file.  Comparing it with
the output it generates can show you how to
produce a simple document of your own.

\section{Ordinary Text}      % Produces section heading.  Lower-level
                             % sections are begun with similar
                             % \subsection and \subsubsection commands.

The ends  of words and sentences are marked
  by   spaces. It  doesn't matter how many
spaces    you type; one is as good as 100.  The
end of   a line counts as a space.

One   or more   blank lines denote the  end
of  a paragraph.

Since any number of consecutive spaces are treated
like a single one, the formatting of the input
file makes no difference to
      \LaTeX,                % The \LaTeX command generates the LaTeX logo.
but it makes a difference to you.  When you use
\LaTeX, making your input file as easy to read
as possible will be a great help as you write
your document and when you change it.  This sample
file shows how you can add comments to your own input
file.

Because printing is different from typewriting,
there are a number of things that you have to do
differently when preparing an input file than if
you were just typing the document directly.
Quotation marks like
       ``this''
have to be handled specially, as do quotes within
quotes:
       ``\,`this'            % \, separates the double and single quote.
        is what I just
        wrote, not  `that'\,''.

Dashes come in three sizes: an
       intra-word
dash, a medium dash for number ranges like
       1--2,
and a punctuation
       dash---like
this.

A sentence-ending space should be larger than the
space between words within a sentence.  You
sometimes have to type special commands in
conjunction with punctuation characters to get
this right, as in the following sentence.
       Gnats, gnus, etc.\ all  % `\ ' makes an inter-word space.
       begin with G\@.         % \@ marks end-of-sentence punctuation.
You should check the spaces after periods when
reading your output to make sure you haven't
forgotten any special cases.  Generating an
ellipsis
       \ldots\               % `\ ' is needed after `\ldots' because TeX
                             % ignores spaces after command names like \ldots
                             % made from \ + letters.
                             %
                             % Note how a `%' character causes TeX to ignore
                             % the end of the input line, so these blank lines
                             % do not start a new paragraph.
                             %
with the right spacing around the periods requires
a special command.

\LaTeX\ interprets some common characters as
commands, so you must type special commands to
generate them.  These characters include the
following:
       \$ \& \% \# \{ and \}.

In printing, text is usually emphasized with an
       \emph{italic}
type style.

\begin{em}
   A long segment of text can also be emphasized
   in this way.  Text within such a segment can be
   given \emph{additional} emphasis.
\end{em}

It is sometimes necessary to prevent \LaTeX\ from
breaking a line where it might otherwise do so.
This may be at a space, as between the ``Mr.''\ and
``Jones'' in
       ``Mr.~Jones'',        % ~ produces an unbreakable interword space.
or within a word---especially when the word is a
symbol like
       \mbox{\emph{itemnum}}
that makes little sense when hyphenated across
lines.

Footnotes\footnote{This is an example of a footnote.}
pose no problem.

\LaTeX\ is good at typesetting mathematical formulas
like
       \( x-3y + z = 7 \)
or
       \( a_{1} > x^{2n} + y^{2n} > x' \)
or
       \( \ip{A}{B} = \sum_{i} a_{i} b_{i} \).
The spaces you type in a formula are
ignored.  Remember that a letter like
       $x$                   % $ ... $  and  \( ... \)  are equivalent
is a formula when it denotes a mathematical
symbol, and it should be typed as one.

\section{Displayed Text}

Text is displayed by indenting it from the left
margin.  Quotations are commonly displayed.  There
are short quotations
\begin{quote}
   This is a short quotation.  It consists of a
   single paragraph of text.  See how it is formatted.
\end{quote}
and longer ones.
\begin{quotation}
   This is a longer quotation.  It consists of two
   paragraphs of text, neither of which are
   particularly interesting.

   This is the second paragraph of the quotation.  It
   is just as dull as the first paragraph.
\end{quotation}
Another frequently-displayed structure is a list.
The following is an example of an \emph{itemized}
list.
\begin{itemize}
   \item This is the first item of an itemized list.
         Each item in the list is marked with a ``tick''.
         You don't have to worry about what kind of tick
         mark is used.

   \item This is the second item of the list.  It
         contains another list nested inside it.  The inner
         list is an \emph{enumerated} list.
         \begin{enumerate}
            \item This is the first item of an enumerated
                  list that is nested within the itemized list.

            \item This is the second item of the inner list.
                  \LaTeX\ allows you to nest lists deeper than
                  you really should.
         \end{enumerate}
         This is the rest of the second item of the outer
         list.  It is no more interesting than any other
         part of the item.
   \item This is the third item of the list.
\end{itemize}
You can even display poetry.
\begin{verse}
   There is an environment
    for verse \\             % The \\ command separates lines
   Whose features some poets % within a stanza.
   will curse.

                             % One or more blank lines separate stanzas.

   For instead of making\\
   Them do \emph{all} line breaking, \\
   It allows them to put too many words on a line when they'd rather be
   forced to be terse.
\end{verse}

Mathematical formulas may also be displayed.  A
displayed formula
is
one-line long; multiline
formulas require special formatting instructions.
   \[  \ip{\Gamma}{\psi'} = x'' + y^{2} + z_{i}^{n}\]
Don't start a paragraph with a displayed equation,
nor make one a paragraph by itself.

\end{document}               % End of document.

\section{Bench Huge Section 449}
% synthetic bench section 449
% This is a sample LaTeX input file.  (Version of 12 August 2004.)
%
% A '%' character causes TeX to ignore all remaining text on the line,
% and is used for comments like this one.

\documentclass{article}      % Specifies the document class

                             % The preamble begins here.
\title{An Example Document}  % Declares the document's title.
\author{Leslie Lamport}      % Declares the author's name.
\date{January 21, 1994}      % Deleting this command produces today's date.

\newcommand{\ip}[2]{(#1, #2)}
                             % Defines \ip{arg1}{arg2} to mean
                             % (arg1, arg2).

%\newcommand{\ip}[2]{\langle #1 | #2\rangle}
                             % This is an alternative definition of
                             % \ip that is commented out.

\begin{document}             % End of preamble and beginning of text.

\maketitle                   % Produces the title.

This is an example input file.  Comparing it with
the output it generates can show you how to
produce a simple document of your own.

\section{Ordinary Text}      % Produces section heading.  Lower-level
                             % sections are begun with similar
                             % \subsection and \subsubsection commands.

The ends  of words and sentences are marked
  by   spaces. It  doesn't matter how many
spaces    you type; one is as good as 100.  The
end of   a line counts as a space.

One   or more   blank lines denote the  end
of  a paragraph.

Since any number of consecutive spaces are treated
like a single one, the formatting of the input
file makes no difference to
      \LaTeX,                % The \LaTeX command generates the LaTeX logo.
but it makes a difference to you.  When you use
\LaTeX, making your input file as easy to read
as possible will be a great help as you write
your document and when you change it.  This sample
file shows how you can add comments to your own input
file.

Because printing is different from typewriting,
there are a number of things that you have to do
differently when preparing an input file than if
you were just typing the document directly.
Quotation marks like
       ``this''
have to be handled specially, as do quotes within
quotes:
       ``\,`this'            % \, separates the double and single quote.
        is what I just
        wrote, not  `that'\,''.

Dashes come in three sizes: an
       intra-word
dash, a medium dash for number ranges like
       1--2,
and a punctuation
       dash---like
this.

A sentence-ending space should be larger than the
space between words within a sentence.  You
sometimes have to type special commands in
conjunction with punctuation characters to get
this right, as in the following sentence.
       Gnats, gnus, etc.\ all  % `\ ' makes an inter-word space.
       begin with G\@.         % \@ marks end-of-sentence punctuation.
You should check the spaces after periods when
reading your output to make sure you haven't
forgotten any special cases.  Generating an
ellipsis
       \ldots\               % `\ ' is needed after `\ldots' because TeX
                             % ignores spaces after command names like \ldots
                             % made from \ + letters.
                             %
                             % Note how a `%' character causes TeX to ignore
                             % the end of the input line, so these blank lines
                             % do not start a new paragraph.
                             %
with the right spacing around the periods requires
a special command.

\LaTeX\ interprets some common characters as
commands, so you must type special commands to
generate them.  These characters include the
following:
       \$ \& \% \# \{ and \}.

In printing, text is usually emphasized with an
       \emph{italic}
type style.

\begin{em}
   A long segment of text can also be emphasized
   in this way.  Text within such a segment can be
   given \emph{additional} emphasis.
\end{em}

It is sometimes necessary to prevent \LaTeX\ from
breaking a line where it might otherwise do so.
This may be at a space, as between the ``Mr.''\ and
``Jones'' in
       ``Mr.~Jones'',        % ~ produces an unbreakable interword space.
or within a word---especially when the word is a
symbol like
       \mbox{\emph{itemnum}}
that makes little sense when hyphenated across
lines.

Footnotes\footnote{This is an example of a footnote.}
pose no problem.

\LaTeX\ is good at typesetting mathematical formulas
like
       \( x-3y + z = 7 \)
or
       \( a_{1} > x^{2n} + y^{2n} > x' \)
or
       \( \ip{A}{B} = \sum_{i} a_{i} b_{i} \).
The spaces you type in a formula are
ignored.  Remember that a letter like
       $x$                   % $ ... $  and  \( ... \)  are equivalent
is a formula when it denotes a mathematical
symbol, and it should be typed as one.

\section{Displayed Text}

Text is displayed by indenting it from the left
margin.  Quotations are commonly displayed.  There
are short quotations
\begin{quote}
   This is a short quotation.  It consists of a
   single paragraph of text.  See how it is formatted.
\end{quote}
and longer ones.
\begin{quotation}
   This is a longer quotation.  It consists of two
   paragraphs of text, neither of which are
   particularly interesting.

   This is the second paragraph of the quotation.  It
   is just as dull as the first paragraph.
\end{quotation}
Another frequently-displayed structure is a list.
The following is an example of an \emph{itemized}
list.
\begin{itemize}
   \item This is the first item of an itemized list.
         Each item in the list is marked with a ``tick''.
         You don't have to worry about what kind of tick
         mark is used.

   \item This is the second item of the list.  It
         contains another list nested inside it.  The inner
         list is an \emph{enumerated} list.
         \begin{enumerate}
            \item This is the first item of an enumerated
                  list that is nested within the itemized list.

            \item This is the second item of the inner list.
                  \LaTeX\ allows you to nest lists deeper than
                  you really should.
         \end{enumerate}
         This is the rest of the second item of the outer
         list.  It is no more interesting than any other
         part of the item.
   \item This is the third item of the list.
\end{itemize}
You can even display poetry.
\begin{verse}
   There is an environment
    for verse \\             % The \\ command separates lines
   Whose features some poets % within a stanza.
   will curse.

                             % One or more blank lines separate stanzas.

   For instead of making\\
   Them do \emph{all} line breaking, \\
   It allows them to put too many words on a line when they'd rather be
   forced to be terse.
\end{verse}

Mathematical formulas may also be displayed.  A
displayed formula
is
one-line long; multiline
formulas require special formatting instructions.
   \[  \ip{\Gamma}{\psi'} = x'' + y^{2} + z_{i}^{n}\]
Don't start a paragraph with a displayed equation,
nor make one a paragraph by itself.

\end{document}               % End of document.

\section{Bench Huge Section 450}
% synthetic bench section 450
% This is a sample LaTeX input file.  (Version of 12 August 2004.)
%
% A '%' character causes TeX to ignore all remaining text on the line,
% and is used for comments like this one.

\documentclass{article}      % Specifies the document class

                             % The preamble begins here.
\title{An Example Document}  % Declares the document's title.
\author{Leslie Lamport}      % Declares the author's name.
\date{January 21, 1994}      % Deleting this command produces today's date.

\newcommand{\ip}[2]{(#1, #2)}
                             % Defines \ip{arg1}{arg2} to mean
                             % (arg1, arg2).

%\newcommand{\ip}[2]{\langle #1 | #2\rangle}
                             % This is an alternative definition of
                             % \ip that is commented out.

\begin{document}             % End of preamble and beginning of text.

\maketitle                   % Produces the title.

This is an example input file.  Comparing it with
the output it generates can show you how to
produce a simple document of your own.

\section{Ordinary Text}      % Produces section heading.  Lower-level
                             % sections are begun with similar
                             % \subsection and \subsubsection commands.

The ends  of words and sentences are marked
  by   spaces. It  doesn't matter how many
spaces    you type; one is as good as 100.  The
end of   a line counts as a space.

One   or more   blank lines denote the  end
of  a paragraph.

Since any number of consecutive spaces are treated
like a single one, the formatting of the input
file makes no difference to
      \LaTeX,                % The \LaTeX command generates the LaTeX logo.
but it makes a difference to you.  When you use
\LaTeX, making your input file as easy to read
as possible will be a great help as you write
your document and when you change it.  This sample
file shows how you can add comments to your own input
file.

Because printing is different from typewriting,
there are a number of things that you have to do
differently when preparing an input file than if
you were just typing the document directly.
Quotation marks like
       ``this''
have to be handled specially, as do quotes within
quotes:
       ``\,`this'            % \, separates the double and single quote.
        is what I just
        wrote, not  `that'\,''.

Dashes come in three sizes: an
       intra-word
dash, a medium dash for number ranges like
       1--2,
and a punctuation
       dash---like
this.

A sentence-ending space should be larger than the
space between words within a sentence.  You
sometimes have to type special commands in
conjunction with punctuation characters to get
this right, as in the following sentence.
       Gnats, gnus, etc.\ all  % `\ ' makes an inter-word space.
       begin with G\@.         % \@ marks end-of-sentence punctuation.
You should check the spaces after periods when
reading your output to make sure you haven't
forgotten any special cases.  Generating an
ellipsis
       \ldots\               % `\ ' is needed after `\ldots' because TeX
                             % ignores spaces after command names like \ldots
                             % made from \ + letters.
                             %
                             % Note how a `%' character causes TeX to ignore
                             % the end of the input line, so these blank lines
                             % do not start a new paragraph.
                             %
with the right spacing around the periods requires
a special command.

\LaTeX\ interprets some common characters as
commands, so you must type special commands to
generate them.  These characters include the
following:
       \$ \& \% \# \{ and \}.

In printing, text is usually emphasized with an
       \emph{italic}
type style.

\begin{em}
   A long segment of text can also be emphasized
   in this way.  Text within such a segment can be
   given \emph{additional} emphasis.
\end{em}

It is sometimes necessary to prevent \LaTeX\ from
breaking a line where it might otherwise do so.
This may be at a space, as between the ``Mr.''\ and
``Jones'' in
       ``Mr.~Jones'',        % ~ produces an unbreakable interword space.
or within a word---especially when the word is a
symbol like
       \mbox{\emph{itemnum}}
that makes little sense when hyphenated across
lines.

Footnotes\footnote{This is an example of a footnote.}
pose no problem.

\LaTeX\ is good at typesetting mathematical formulas
like
       \( x-3y + z = 7 \)
or
       \( a_{1} > x^{2n} + y^{2n} > x' \)
or
       \( \ip{A}{B} = \sum_{i} a_{i} b_{i} \).
The spaces you type in a formula are
ignored.  Remember that a letter like
       $x$                   % $ ... $  and  \( ... \)  are equivalent
is a formula when it denotes a mathematical
symbol, and it should be typed as one.

\section{Displayed Text}

Text is displayed by indenting it from the left
margin.  Quotations are commonly displayed.  There
are short quotations
\begin{quote}
   This is a short quotation.  It consists of a
   single paragraph of text.  See how it is formatted.
\end{quote}
and longer ones.
\begin{quotation}
   This is a longer quotation.  It consists of two
   paragraphs of text, neither of which are
   particularly interesting.

   This is the second paragraph of the quotation.  It
   is just as dull as the first paragraph.
\end{quotation}
Another frequently-displayed structure is a list.
The following is an example of an \emph{itemized}
list.
\begin{itemize}
   \item This is the first item of an itemized list.
         Each item in the list is marked with a ``tick''.
         You don't have to worry about what kind of tick
         mark is used.

   \item This is the second item of the list.  It
         contains another list nested inside it.  The inner
         list is an \emph{enumerated} list.
         \begin{enumerate}
            \item This is the first item of an enumerated
                  list that is nested within the itemized list.

            \item This is the second item of the inner list.
                  \LaTeX\ allows you to nest lists deeper than
                  you really should.
         \end{enumerate}
         This is the rest of the second item of the outer
         list.  It is no more interesting than any other
         part of the item.
   \item This is the third item of the list.
\end{itemize}
You can even display poetry.
\begin{verse}
   There is an environment
    for verse \\             % The \\ command separates lines
   Whose features some poets % within a stanza.
   will curse.

                             % One or more blank lines separate stanzas.

   For instead of making\\
   Them do \emph{all} line breaking, \\
   It allows them to put too many words on a line when they'd rather be
   forced to be terse.
\end{verse}

Mathematical formulas may also be displayed.  A
displayed formula
is
one-line long; multiline
formulas require special formatting instructions.
   \[  \ip{\Gamma}{\psi'} = x'' + y^{2} + z_{i}^{n}\]
Don't start a paragraph with a displayed equation,
nor make one a paragraph by itself.

\end{document}               % End of document.

\section{Bench Huge Section 451}
% synthetic bench section 451
% This is a sample LaTeX input file.  (Version of 12 August 2004.)
%
% A '%' character causes TeX to ignore all remaining text on the line,
% and is used for comments like this one.

\documentclass{article}      % Specifies the document class

                             % The preamble begins here.
\title{An Example Document}  % Declares the document's title.
\author{Leslie Lamport}      % Declares the author's name.
\date{January 21, 1994}      % Deleting this command produces today's date.

\newcommand{\ip}[2]{(#1, #2)}
                             % Defines \ip{arg1}{arg2} to mean
                             % (arg1, arg2).

%\newcommand{\ip}[2]{\langle #1 | #2\rangle}
                             % This is an alternative definition of
                             % \ip that is commented out.

\begin{document}             % End of preamble and beginning of text.

\maketitle                   % Produces the title.

This is an example input file.  Comparing it with
the output it generates can show you how to
produce a simple document of your own.

\section{Ordinary Text}      % Produces section heading.  Lower-level
                             % sections are begun with similar
                             % \subsection and \subsubsection commands.

The ends  of words and sentences are marked
  by   spaces. It  doesn't matter how many
spaces    you type; one is as good as 100.  The
end of   a line counts as a space.

One   or more   blank lines denote the  end
of  a paragraph.

Since any number of consecutive spaces are treated
like a single one, the formatting of the input
file makes no difference to
      \LaTeX,                % The \LaTeX command generates the LaTeX logo.
but it makes a difference to you.  When you use
\LaTeX, making your input file as easy to read
as possible will be a great help as you write
your document and when you change it.  This sample
file shows how you can add comments to your own input
file.

Because printing is different from typewriting,
there are a number of things that you have to do
differently when preparing an input file than if
you were just typing the document directly.
Quotation marks like
       ``this''
have to be handled specially, as do quotes within
quotes:
       ``\,`this'            % \, separates the double and single quote.
        is what I just
        wrote, not  `that'\,''.

Dashes come in three sizes: an
       intra-word
dash, a medium dash for number ranges like
       1--2,
and a punctuation
       dash---like
this.

A sentence-ending space should be larger than the
space between words within a sentence.  You
sometimes have to type special commands in
conjunction with punctuation characters to get
this right, as in the following sentence.
       Gnats, gnus, etc.\ all  % `\ ' makes an inter-word space.
       begin with G\@.         % \@ marks end-of-sentence punctuation.
You should check the spaces after periods when
reading your output to make sure you haven't
forgotten any special cases.  Generating an
ellipsis
       \ldots\               % `\ ' is needed after `\ldots' because TeX
                             % ignores spaces after command names like \ldots
                             % made from \ + letters.
                             %
                             % Note how a `%' character causes TeX to ignore
                             % the end of the input line, so these blank lines
                             % do not start a new paragraph.
                             %
with the right spacing around the periods requires
a special command.

\LaTeX\ interprets some common characters as
commands, so you must type special commands to
generate them.  These characters include the
following:
       \$ \& \% \# \{ and \}.

In printing, text is usually emphasized with an
       \emph{italic}
type style.

\begin{em}
   A long segment of text can also be emphasized
   in this way.  Text within such a segment can be
   given \emph{additional} emphasis.
\end{em}

It is sometimes necessary to prevent \LaTeX\ from
breaking a line where it might otherwise do so.
This may be at a space, as between the ``Mr.''\ and
``Jones'' in
       ``Mr.~Jones'',        % ~ produces an unbreakable interword space.
or within a word---especially when the word is a
symbol like
       \mbox{\emph{itemnum}}
that makes little sense when hyphenated across
lines.

Footnotes\footnote{This is an example of a footnote.}
pose no problem.

\LaTeX\ is good at typesetting mathematical formulas
like
       \( x-3y + z = 7 \)
or
       \( a_{1} > x^{2n} + y^{2n} > x' \)
or
       \( \ip{A}{B} = \sum_{i} a_{i} b_{i} \).
The spaces you type in a formula are
ignored.  Remember that a letter like
       $x$                   % $ ... $  and  \( ... \)  are equivalent
is a formula when it denotes a mathematical
symbol, and it should be typed as one.

\section{Displayed Text}

Text is displayed by indenting it from the left
margin.  Quotations are commonly displayed.  There
are short quotations
\begin{quote}
   This is a short quotation.  It consists of a
   single paragraph of text.  See how it is formatted.
\end{quote}
and longer ones.
\begin{quotation}
   This is a longer quotation.  It consists of two
   paragraphs of text, neither of which are
   particularly interesting.

   This is the second paragraph of the quotation.  It
   is just as dull as the first paragraph.
\end{quotation}
Another frequently-displayed structure is a list.
The following is an example of an \emph{itemized}
list.
\begin{itemize}
   \item This is the first item of an itemized list.
         Each item in the list is marked with a ``tick''.
         You don't have to worry about what kind of tick
         mark is used.

   \item This is the second item of the list.  It
         contains another list nested inside it.  The inner
         list is an \emph{enumerated} list.
         \begin{enumerate}
            \item This is the first item of an enumerated
                  list that is nested within the itemized list.

            \item This is the second item of the inner list.
                  \LaTeX\ allows you to nest lists deeper than
                  you really should.
         \end{enumerate}
         This is the rest of the second item of the outer
         list.  It is no more interesting than any other
         part of the item.
   \item This is the third item of the list.
\end{itemize}
You can even display poetry.
\begin{verse}
   There is an environment
    for verse \\             % The \\ command separates lines
   Whose features some poets % within a stanza.
   will curse.

                             % One or more blank lines separate stanzas.

   For instead of making\\
   Them do \emph{all} line breaking, \\
   It allows them to put too many words on a line when they'd rather be
   forced to be terse.
\end{verse}

Mathematical formulas may also be displayed.  A
displayed formula
is
one-line long; multiline
formulas require special formatting instructions.
   \[  \ip{\Gamma}{\psi'} = x'' + y^{2} + z_{i}^{n}\]
Don't start a paragraph with a displayed equation,
nor make one a paragraph by itself.

\end{document}               % End of document.

\section{Bench Huge Section 452}
% synthetic bench section 452
% This is a sample LaTeX input file.  (Version of 12 August 2004.)
%
% A '%' character causes TeX to ignore all remaining text on the line,
% and is used for comments like this one.

\documentclass{article}      % Specifies the document class

                             % The preamble begins here.
\title{An Example Document}  % Declares the document's title.
\author{Leslie Lamport}      % Declares the author's name.
\date{January 21, 1994}      % Deleting this command produces today's date.

\newcommand{\ip}[2]{(#1, #2)}
                             % Defines \ip{arg1}{arg2} to mean
                             % (arg1, arg2).

%\newcommand{\ip}[2]{\langle #1 | #2\rangle}
                             % This is an alternative definition of
                             % \ip that is commented out.

\begin{document}             % End of preamble and beginning of text.

\maketitle                   % Produces the title.

This is an example input file.  Comparing it with
the output it generates can show you how to
produce a simple document of your own.

\section{Ordinary Text}      % Produces section heading.  Lower-level
                             % sections are begun with similar
                             % \subsection and \subsubsection commands.

The ends  of words and sentences are marked
  by   spaces. It  doesn't matter how many
spaces    you type; one is as good as 100.  The
end of   a line counts as a space.

One   or more   blank lines denote the  end
of  a paragraph.

Since any number of consecutive spaces are treated
like a single one, the formatting of the input
file makes no difference to
      \LaTeX,                % The \LaTeX command generates the LaTeX logo.
but it makes a difference to you.  When you use
\LaTeX, making your input file as easy to read
as possible will be a great help as you write
your document and when you change it.  This sample
file shows how you can add comments to your own input
file.

Because printing is different from typewriting,
there are a number of things that you have to do
differently when preparing an input file than if
you were just typing the document directly.
Quotation marks like
       ``this''
have to be handled specially, as do quotes within
quotes:
       ``\,`this'            % \, separates the double and single quote.
        is what I just
        wrote, not  `that'\,''.

Dashes come in three sizes: an
       intra-word
dash, a medium dash for number ranges like
       1--2,
and a punctuation
       dash---like
this.

A sentence-ending space should be larger than the
space between words within a sentence.  You
sometimes have to type special commands in
conjunction with punctuation characters to get
this right, as in the following sentence.
       Gnats, gnus, etc.\ all  % `\ ' makes an inter-word space.
       begin with G\@.         % \@ marks end-of-sentence punctuation.
You should check the spaces after periods when
reading your output to make sure you haven't
forgotten any special cases.  Generating an
ellipsis
       \ldots\               % `\ ' is needed after `\ldots' because TeX
                             % ignores spaces after command names like \ldots
                             % made from \ + letters.
                             %
                             % Note how a `%' character causes TeX to ignore
                             % the end of the input line, so these blank lines
                             % do not start a new paragraph.
                             %
with the right spacing around the periods requires
a special command.

\LaTeX\ interprets some common characters as
commands, so you must type special commands to
generate them.  These characters include the
following:
       \$ \& \% \# \{ and \}.

In printing, text is usually emphasized with an
       \emph{italic}
type style.

\begin{em}
   A long segment of text can also be emphasized
   in this way.  Text within such a segment can be
   given \emph{additional} emphasis.
\end{em}

It is sometimes necessary to prevent \LaTeX\ from
breaking a line where it might otherwise do so.
This may be at a space, as between the ``Mr.''\ and
``Jones'' in
       ``Mr.~Jones'',        % ~ produces an unbreakable interword space.
or within a word---especially when the word is a
symbol like
       \mbox{\emph{itemnum}}
that makes little sense when hyphenated across
lines.

Footnotes\footnote{This is an example of a footnote.}
pose no problem.

\LaTeX\ is good at typesetting mathematical formulas
like
       \( x-3y + z = 7 \)
or
       \( a_{1} > x^{2n} + y^{2n} > x' \)
or
       \( \ip{A}{B} = \sum_{i} a_{i} b_{i} \).
The spaces you type in a formula are
ignored.  Remember that a letter like
       $x$                   % $ ... $  and  \( ... \)  are equivalent
is a formula when it denotes a mathematical
symbol, and it should be typed as one.

\section{Displayed Text}

Text is displayed by indenting it from the left
margin.  Quotations are commonly displayed.  There
are short quotations
\begin{quote}
   This is a short quotation.  It consists of a
   single paragraph of text.  See how it is formatted.
\end{quote}
and longer ones.
\begin{quotation}
   This is a longer quotation.  It consists of two
   paragraphs of text, neither of which are
   particularly interesting.

   This is the second paragraph of the quotation.  It
   is just as dull as the first paragraph.
\end{quotation}
Another frequently-displayed structure is a list.
The following is an example of an \emph{itemized}
list.
\begin{itemize}
   \item This is the first item of an itemized list.
         Each item in the list is marked with a ``tick''.
         You don't have to worry about what kind of tick
         mark is used.

   \item This is the second item of the list.  It
         contains another list nested inside it.  The inner
         list is an \emph{enumerated} list.
         \begin{enumerate}
            \item This is the first item of an enumerated
                  list that is nested within the itemized list.

            \item This is the second item of the inner list.
                  \LaTeX\ allows you to nest lists deeper than
                  you really should.
         \end{enumerate}
         This is the rest of the second item of the outer
         list.  It is no more interesting than any other
         part of the item.
   \item This is the third item of the list.
\end{itemize}
You can even display poetry.
\begin{verse}
   There is an environment
    for verse \\             % The \\ command separates lines
   Whose features some poets % within a stanza.
   will curse.

                             % One or more blank lines separate stanzas.

   For instead of making\\
   Them do \emph{all} line breaking, \\
   It allows them to put too many words on a line when they'd rather be
   forced to be terse.
\end{verse}

Mathematical formulas may also be displayed.  A
displayed formula
is
one-line long; multiline
formulas require special formatting instructions.
   \[  \ip{\Gamma}{\psi'} = x'' + y^{2} + z_{i}^{n}\]
Don't start a paragraph with a displayed equation,
nor make one a paragraph by itself.

\end{document}               % End of document.

\section{Bench Huge Section 453}
% synthetic bench section 453
% This is a sample LaTeX input file.  (Version of 12 August 2004.)
%
% A '%' character causes TeX to ignore all remaining text on the line,
% and is used for comments like this one.

\documentclass{article}      % Specifies the document class

                             % The preamble begins here.
\title{An Example Document}  % Declares the document's title.
\author{Leslie Lamport}      % Declares the author's name.
\date{January 21, 1994}      % Deleting this command produces today's date.

\newcommand{\ip}[2]{(#1, #2)}
                             % Defines \ip{arg1}{arg2} to mean
                             % (arg1, arg2).

%\newcommand{\ip}[2]{\langle #1 | #2\rangle}
                             % This is an alternative definition of
                             % \ip that is commented out.

\begin{document}             % End of preamble and beginning of text.

\maketitle                   % Produces the title.

This is an example input file.  Comparing it with
the output it generates can show you how to
produce a simple document of your own.

\section{Ordinary Text}      % Produces section heading.  Lower-level
                             % sections are begun with similar
                             % \subsection and \subsubsection commands.

The ends  of words and sentences are marked
  by   spaces. It  doesn't matter how many
spaces    you type; one is as good as 100.  The
end of   a line counts as a space.

One   or more   blank lines denote the  end
of  a paragraph.

Since any number of consecutive spaces are treated
like a single one, the formatting of the input
file makes no difference to
      \LaTeX,                % The \LaTeX command generates the LaTeX logo.
but it makes a difference to you.  When you use
\LaTeX, making your input file as easy to read
as possible will be a great help as you write
your document and when you change it.  This sample
file shows how you can add comments to your own input
file.

Because printing is different from typewriting,
there are a number of things that you have to do
differently when preparing an input file than if
you were just typing the document directly.
Quotation marks like
       ``this''
have to be handled specially, as do quotes within
quotes:
       ``\,`this'            % \, separates the double and single quote.
        is what I just
        wrote, not  `that'\,''.

Dashes come in three sizes: an
       intra-word
dash, a medium dash for number ranges like
       1--2,
and a punctuation
       dash---like
this.

A sentence-ending space should be larger than the
space between words within a sentence.  You
sometimes have to type special commands in
conjunction with punctuation characters to get
this right, as in the following sentence.
       Gnats, gnus, etc.\ all  % `\ ' makes an inter-word space.
       begin with G\@.         % \@ marks end-of-sentence punctuation.
You should check the spaces after periods when
reading your output to make sure you haven't
forgotten any special cases.  Generating an
ellipsis
       \ldots\               % `\ ' is needed after `\ldots' because TeX
                             % ignores spaces after command names like \ldots
                             % made from \ + letters.
                             %
                             % Note how a `%' character causes TeX to ignore
                             % the end of the input line, so these blank lines
                             % do not start a new paragraph.
                             %
with the right spacing around the periods requires
a special command.

\LaTeX\ interprets some common characters as
commands, so you must type special commands to
generate them.  These characters include the
following:
       \$ \& \% \# \{ and \}.

In printing, text is usually emphasized with an
       \emph{italic}
type style.

\begin{em}
   A long segment of text can also be emphasized
   in this way.  Text within such a segment can be
   given \emph{additional} emphasis.
\end{em}

It is sometimes necessary to prevent \LaTeX\ from
breaking a line where it might otherwise do so.
This may be at a space, as between the ``Mr.''\ and
``Jones'' in
       ``Mr.~Jones'',        % ~ produces an unbreakable interword space.
or within a word---especially when the word is a
symbol like
       \mbox{\emph{itemnum}}
that makes little sense when hyphenated across
lines.

Footnotes\footnote{This is an example of a footnote.}
pose no problem.

\LaTeX\ is good at typesetting mathematical formulas
like
       \( x-3y + z = 7 \)
or
       \( a_{1} > x^{2n} + y^{2n} > x' \)
or
       \( \ip{A}{B} = \sum_{i} a_{i} b_{i} \).
The spaces you type in a formula are
ignored.  Remember that a letter like
       $x$                   % $ ... $  and  \( ... \)  are equivalent
is a formula when it denotes a mathematical
symbol, and it should be typed as one.

\section{Displayed Text}

Text is displayed by indenting it from the left
margin.  Quotations are commonly displayed.  There
are short quotations
\begin{quote}
   This is a short quotation.  It consists of a
   single paragraph of text.  See how it is formatted.
\end{quote}
and longer ones.
\begin{quotation}
   This is a longer quotation.  It consists of two
   paragraphs of text, neither of which are
   particularly interesting.

   This is the second paragraph of the quotation.  It
   is just as dull as the first paragraph.
\end{quotation}
Another frequently-displayed structure is a list.
The following is an example of an \emph{itemized}
list.
\begin{itemize}
   \item This is the first item of an itemized list.
         Each item in the list is marked with a ``tick''.
         You don't have to worry about what kind of tick
         mark is used.

   \item This is the second item of the list.  It
         contains another list nested inside it.  The inner
         list is an \emph{enumerated} list.
         \begin{enumerate}
            \item This is the first item of an enumerated
                  list that is nested within the itemized list.

            \item This is the second item of the inner list.
                  \LaTeX\ allows you to nest lists deeper than
                  you really should.
         \end{enumerate}
         This is the rest of the second item of the outer
         list.  It is no more interesting than any other
         part of the item.
   \item This is the third item of the list.
\end{itemize}
You can even display poetry.
\begin{verse}
   There is an environment
    for verse \\             % The \\ command separates lines
   Whose features some poets % within a stanza.
   will curse.

                             % One or more blank lines separate stanzas.

   For instead of making\\
   Them do \emph{all} line breaking, \\
   It allows them to put too many words on a line when they'd rather be
   forced to be terse.
\end{verse}

Mathematical formulas may also be displayed.  A
displayed formula
is
one-line long; multiline
formulas require special formatting instructions.
   \[  \ip{\Gamma}{\psi'} = x'' + y^{2} + z_{i}^{n}\]
Don't start a paragraph with a displayed equation,
nor make one a paragraph by itself.

\end{document}               % End of document.

\section{Bench Huge Section 454}
% synthetic bench section 454
% This is a sample LaTeX input file.  (Version of 12 August 2004.)
%
% A '%' character causes TeX to ignore all remaining text on the line,
% and is used for comments like this one.

\documentclass{article}      % Specifies the document class

                             % The preamble begins here.
\title{An Example Document}  % Declares the document's title.
\author{Leslie Lamport}      % Declares the author's name.
\date{January 21, 1994}      % Deleting this command produces today's date.

\newcommand{\ip}[2]{(#1, #2)}
                             % Defines \ip{arg1}{arg2} to mean
                             % (arg1, arg2).

%\newcommand{\ip}[2]{\langle #1 | #2\rangle}
                             % This is an alternative definition of
                             % \ip that is commented out.

\begin{document}             % End of preamble and beginning of text.

\maketitle                   % Produces the title.

This is an example input file.  Comparing it with
the output it generates can show you how to
produce a simple document of your own.

\section{Ordinary Text}      % Produces section heading.  Lower-level
                             % sections are begun with similar
                             % \subsection and \subsubsection commands.

The ends  of words and sentences are marked
  by   spaces. It  doesn't matter how many
spaces    you type; one is as good as 100.  The
end of   a line counts as a space.

One   or more   blank lines denote the  end
of  a paragraph.

Since any number of consecutive spaces are treated
like a single one, the formatting of the input
file makes no difference to
      \LaTeX,                % The \LaTeX command generates the LaTeX logo.
but it makes a difference to you.  When you use
\LaTeX, making your input file as easy to read
as possible will be a great help as you write
your document and when you change it.  This sample
file shows how you can add comments to your own input
file.

Because printing is different from typewriting,
there are a number of things that you have to do
differently when preparing an input file than if
you were just typing the document directly.
Quotation marks like
       ``this''
have to be handled specially, as do quotes within
quotes:
       ``\,`this'            % \, separates the double and single quote.
        is what I just
        wrote, not  `that'\,''.

Dashes come in three sizes: an
       intra-word
dash, a medium dash for number ranges like
       1--2,
and a punctuation
       dash---like
this.

A sentence-ending space should be larger than the
space between words within a sentence.  You
sometimes have to type special commands in
conjunction with punctuation characters to get
this right, as in the following sentence.
       Gnats, gnus, etc.\ all  % `\ ' makes an inter-word space.
       begin with G\@.         % \@ marks end-of-sentence punctuation.
You should check the spaces after periods when
reading your output to make sure you haven't
forgotten any special cases.  Generating an
ellipsis
       \ldots\               % `\ ' is needed after `\ldots' because TeX
                             % ignores spaces after command names like \ldots
                             % made from \ + letters.
                             %
                             % Note how a `%' character causes TeX to ignore
                             % the end of the input line, so these blank lines
                             % do not start a new paragraph.
                             %
with the right spacing around the periods requires
a special command.

\LaTeX\ interprets some common characters as
commands, so you must type special commands to
generate them.  These characters include the
following:
       \$ \& \% \# \{ and \}.

In printing, text is usually emphasized with an
       \emph{italic}
type style.

\begin{em}
   A long segment of text can also be emphasized
   in this way.  Text within such a segment can be
   given \emph{additional} emphasis.
\end{em}

It is sometimes necessary to prevent \LaTeX\ from
breaking a line where it might otherwise do so.
This may be at a space, as between the ``Mr.''\ and
``Jones'' in
       ``Mr.~Jones'',        % ~ produces an unbreakable interword space.
or within a word---especially when the word is a
symbol like
       \mbox{\emph{itemnum}}
that makes little sense when hyphenated across
lines.

Footnotes\footnote{This is an example of a footnote.}
pose no problem.

\LaTeX\ is good at typesetting mathematical formulas
like
       \( x-3y + z = 7 \)
or
       \( a_{1} > x^{2n} + y^{2n} > x' \)
or
       \( \ip{A}{B} = \sum_{i} a_{i} b_{i} \).
The spaces you type in a formula are
ignored.  Remember that a letter like
       $x$                   % $ ... $  and  \( ... \)  are equivalent
is a formula when it denotes a mathematical
symbol, and it should be typed as one.

\section{Displayed Text}

Text is displayed by indenting it from the left
margin.  Quotations are commonly displayed.  There
are short quotations
\begin{quote}
   This is a short quotation.  It consists of a
   single paragraph of text.  See how it is formatted.
\end{quote}
and longer ones.
\begin{quotation}
   This is a longer quotation.  It consists of two
   paragraphs of text, neither of which are
   particularly interesting.

   This is the second paragraph of the quotation.  It
   is just as dull as the first paragraph.
\end{quotation}
Another frequently-displayed structure is a list.
The following is an example of an \emph{itemized}
list.
\begin{itemize}
   \item This is the first item of an itemized list.
         Each item in the list is marked with a ``tick''.
         You don't have to worry about what kind of tick
         mark is used.

   \item This is the second item of the list.  It
         contains another list nested inside it.  The inner
         list is an \emph{enumerated} list.
         \begin{enumerate}
            \item This is the first item of an enumerated
                  list that is nested within the itemized list.

            \item This is the second item of the inner list.
                  \LaTeX\ allows you to nest lists deeper than
                  you really should.
         \end{enumerate}
         This is the rest of the second item of the outer
         list.  It is no more interesting than any other
         part of the item.
   \item This is the third item of the list.
\end{itemize}
You can even display poetry.
\begin{verse}
   There is an environment
    for verse \\             % The \\ command separates lines
   Whose features some poets % within a stanza.
   will curse.

                             % One or more blank lines separate stanzas.

   For instead of making\\
   Them do \emph{all} line breaking, \\
   It allows them to put too many words on a line when they'd rather be
   forced to be terse.
\end{verse}

Mathematical formulas may also be displayed.  A
displayed formula
is
one-line long; multiline
formulas require special formatting instructions.
   \[  \ip{\Gamma}{\psi'} = x'' + y^{2} + z_{i}^{n}\]
Don't start a paragraph with a displayed equation,
nor make one a paragraph by itself.

\end{document}               % End of document.

\section{Bench Huge Section 455}
% synthetic bench section 455
% This is a sample LaTeX input file.  (Version of 12 August 2004.)
%
% A '%' character causes TeX to ignore all remaining text on the line,
% and is used for comments like this one.

\documentclass{article}      % Specifies the document class

                             % The preamble begins here.
\title{An Example Document}  % Declares the document's title.
\author{Leslie Lamport}      % Declares the author's name.
\date{January 21, 1994}      % Deleting this command produces today's date.

\newcommand{\ip}[2]{(#1, #2)}
                             % Defines \ip{arg1}{arg2} to mean
                             % (arg1, arg2).

%\newcommand{\ip}[2]{\langle #1 | #2\rangle}
                             % This is an alternative definition of
                             % \ip that is commented out.

\begin{document}             % End of preamble and beginning of text.

\maketitle                   % Produces the title.

This is an example input file.  Comparing it with
the output it generates can show you how to
produce a simple document of your own.

\section{Ordinary Text}      % Produces section heading.  Lower-level
                             % sections are begun with similar
                             % \subsection and \subsubsection commands.

The ends  of words and sentences are marked
  by   spaces. It  doesn't matter how many
spaces    you type; one is as good as 100.  The
end of   a line counts as a space.

One   or more   blank lines denote the  end
of  a paragraph.

Since any number of consecutive spaces are treated
like a single one, the formatting of the input
file makes no difference to
      \LaTeX,                % The \LaTeX command generates the LaTeX logo.
but it makes a difference to you.  When you use
\LaTeX, making your input file as easy to read
as possible will be a great help as you write
your document and when you change it.  This sample
file shows how you can add comments to your own input
file.

Because printing is different from typewriting,
there are a number of things that you have to do
differently when preparing an input file than if
you were just typing the document directly.
Quotation marks like
       ``this''
have to be handled specially, as do quotes within
quotes:
       ``\,`this'            % \, separates the double and single quote.
        is what I just
        wrote, not  `that'\,''.

Dashes come in three sizes: an
       intra-word
dash, a medium dash for number ranges like
       1--2,
and a punctuation
       dash---like
this.

A sentence-ending space should be larger than the
space between words within a sentence.  You
sometimes have to type special commands in
conjunction with punctuation characters to get
this right, as in the following sentence.
       Gnats, gnus, etc.\ all  % `\ ' makes an inter-word space.
       begin with G\@.         % \@ marks end-of-sentence punctuation.
You should check the spaces after periods when
reading your output to make sure you haven't
forgotten any special cases.  Generating an
ellipsis
       \ldots\               % `\ ' is needed after `\ldots' because TeX
                             % ignores spaces after command names like \ldots
                             % made from \ + letters.
                             %
                             % Note how a `%' character causes TeX to ignore
                             % the end of the input line, so these blank lines
                             % do not start a new paragraph.
                             %
with the right spacing around the periods requires
a special command.

\LaTeX\ interprets some common characters as
commands, so you must type special commands to
generate them.  These characters include the
following:
       \$ \& \% \# \{ and \}.

In printing, text is usually emphasized with an
       \emph{italic}
type style.

\begin{em}
   A long segment of text can also be emphasized
   in this way.  Text within such a segment can be
   given \emph{additional} emphasis.
\end{em}

It is sometimes necessary to prevent \LaTeX\ from
breaking a line where it might otherwise do so.
This may be at a space, as between the ``Mr.''\ and
``Jones'' in
       ``Mr.~Jones'',        % ~ produces an unbreakable interword space.
or within a word---especially when the word is a
symbol like
       \mbox{\emph{itemnum}}
that makes little sense when hyphenated across
lines.

Footnotes\footnote{This is an example of a footnote.}
pose no problem.

\LaTeX\ is good at typesetting mathematical formulas
like
       \( x-3y + z = 7 \)
or
       \( a_{1} > x^{2n} + y^{2n} > x' \)
or
       \( \ip{A}{B} = \sum_{i} a_{i} b_{i} \).
The spaces you type in a formula are
ignored.  Remember that a letter like
       $x$                   % $ ... $  and  \( ... \)  are equivalent
is a formula when it denotes a mathematical
symbol, and it should be typed as one.

\section{Displayed Text}

Text is displayed by indenting it from the left
margin.  Quotations are commonly displayed.  There
are short quotations
\begin{quote}
   This is a short quotation.  It consists of a
   single paragraph of text.  See how it is formatted.
\end{quote}
and longer ones.
\begin{quotation}
   This is a longer quotation.  It consists of two
   paragraphs of text, neither of which are
   particularly interesting.

   This is the second paragraph of the quotation.  It
   is just as dull as the first paragraph.
\end{quotation}
Another frequently-displayed structure is a list.
The following is an example of an \emph{itemized}
list.
\begin{itemize}
   \item This is the first item of an itemized list.
         Each item in the list is marked with a ``tick''.
         You don't have to worry about what kind of tick
         mark is used.

   \item This is the second item of the list.  It
         contains another list nested inside it.  The inner
         list is an \emph{enumerated} list.
         \begin{enumerate}
            \item This is the first item of an enumerated
                  list that is nested within the itemized list.

            \item This is the second item of the inner list.
                  \LaTeX\ allows you to nest lists deeper than
                  you really should.
         \end{enumerate}
         This is the rest of the second item of the outer
         list.  It is no more interesting than any other
         part of the item.
   \item This is the third item of the list.
\end{itemize}
You can even display poetry.
\begin{verse}
   There is an environment
    for verse \\             % The \\ command separates lines
   Whose features some poets % within a stanza.
   will curse.

                             % One or more blank lines separate stanzas.

   For instead of making\\
   Them do \emph{all} line breaking, \\
   It allows them to put too many words on a line when they'd rather be
   forced to be terse.
\end{verse}

Mathematical formulas may also be displayed.  A
displayed formula
is
one-line long; multiline
formulas require special formatting instructions.
   \[  \ip{\Gamma}{\psi'} = x'' + y^{2} + z_{i}^{n}\]
Don't start a paragraph with a displayed equation,
nor make one a paragraph by itself.

\end{document}               % End of document.

\section{Bench Huge Section 456}
% synthetic bench section 456
% This is a sample LaTeX input file.  (Version of 12 August 2004.)
%
% A '%' character causes TeX to ignore all remaining text on the line,
% and is used for comments like this one.

\documentclass{article}      % Specifies the document class

                             % The preamble begins here.
\title{An Example Document}  % Declares the document's title.
\author{Leslie Lamport}      % Declares the author's name.
\date{January 21, 1994}      % Deleting this command produces today's date.

\newcommand{\ip}[2]{(#1, #2)}
                             % Defines \ip{arg1}{arg2} to mean
                             % (arg1, arg2).

%\newcommand{\ip}[2]{\langle #1 | #2\rangle}
                             % This is an alternative definition of
                             % \ip that is commented out.

\begin{document}             % End of preamble and beginning of text.

\maketitle                   % Produces the title.

This is an example input file.  Comparing it with
the output it generates can show you how to
produce a simple document of your own.

\section{Ordinary Text}      % Produces section heading.  Lower-level
                             % sections are begun with similar
                             % \subsection and \subsubsection commands.

The ends  of words and sentences are marked
  by   spaces. It  doesn't matter how many
spaces    you type; one is as good as 100.  The
end of   a line counts as a space.

One   or more   blank lines denote the  end
of  a paragraph.

Since any number of consecutive spaces are treated
like a single one, the formatting of the input
file makes no difference to
      \LaTeX,                % The \LaTeX command generates the LaTeX logo.
but it makes a difference to you.  When you use
\LaTeX, making your input file as easy to read
as possible will be a great help as you write
your document and when you change it.  This sample
file shows how you can add comments to your own input
file.

Because printing is different from typewriting,
there are a number of things that you have to do
differently when preparing an input file than if
you were just typing the document directly.
Quotation marks like
       ``this''
have to be handled specially, as do quotes within
quotes:
       ``\,`this'            % \, separates the double and single quote.
        is what I just
        wrote, not  `that'\,''.

Dashes come in three sizes: an
       intra-word
dash, a medium dash for number ranges like
       1--2,
and a punctuation
       dash---like
this.

A sentence-ending space should be larger than the
space between words within a sentence.  You
sometimes have to type special commands in
conjunction with punctuation characters to get
this right, as in the following sentence.
       Gnats, gnus, etc.\ all  % `\ ' makes an inter-word space.
       begin with G\@.         % \@ marks end-of-sentence punctuation.
You should check the spaces after periods when
reading your output to make sure you haven't
forgotten any special cases.  Generating an
ellipsis
       \ldots\               % `\ ' is needed after `\ldots' because TeX
                             % ignores spaces after command names like \ldots
                             % made from \ + letters.
                             %
                             % Note how a `%' character causes TeX to ignore
                             % the end of the input line, so these blank lines
                             % do not start a new paragraph.
                             %
with the right spacing around the periods requires
a special command.

\LaTeX\ interprets some common characters as
commands, so you must type special commands to
generate them.  These characters include the
following:
       \$ \& \% \# \{ and \}.

In printing, text is usually emphasized with an
       \emph{italic}
type style.

\begin{em}
   A long segment of text can also be emphasized
   in this way.  Text within such a segment can be
   given \emph{additional} emphasis.
\end{em}

It is sometimes necessary to prevent \LaTeX\ from
breaking a line where it might otherwise do so.
This may be at a space, as between the ``Mr.''\ and
``Jones'' in
       ``Mr.~Jones'',        % ~ produces an unbreakable interword space.
or within a word---especially when the word is a
symbol like
       \mbox{\emph{itemnum}}
that makes little sense when hyphenated across
lines.

Footnotes\footnote{This is an example of a footnote.}
pose no problem.

\LaTeX\ is good at typesetting mathematical formulas
like
       \( x-3y + z = 7 \)
or
       \( a_{1} > x^{2n} + y^{2n} > x' \)
or
       \( \ip{A}{B} = \sum_{i} a_{i} b_{i} \).
The spaces you type in a formula are
ignored.  Remember that a letter like
       $x$                   % $ ... $  and  \( ... \)  are equivalent
is a formula when it denotes a mathematical
symbol, and it should be typed as one.

\section{Displayed Text}

Text is displayed by indenting it from the left
margin.  Quotations are commonly displayed.  There
are short quotations
\begin{quote}
   This is a short quotation.  It consists of a
   single paragraph of text.  See how it is formatted.
\end{quote}
and longer ones.
\begin{quotation}
   This is a longer quotation.  It consists of two
   paragraphs of text, neither of which are
   particularly interesting.

   This is the second paragraph of the quotation.  It
   is just as dull as the first paragraph.
\end{quotation}
Another frequently-displayed structure is a list.
The following is an example of an \emph{itemized}
list.
\begin{itemize}
   \item This is the first item of an itemized list.
         Each item in the list is marked with a ``tick''.
         You don't have to worry about what kind of tick
         mark is used.

   \item This is the second item of the list.  It
         contains another list nested inside it.  The inner
         list is an \emph{enumerated} list.
         \begin{enumerate}
            \item This is the first item of an enumerated
                  list that is nested within the itemized list.

            \item This is the second item of the inner list.
                  \LaTeX\ allows you to nest lists deeper than
                  you really should.
         \end{enumerate}
         This is the rest of the second item of the outer
         list.  It is no more interesting than any other
         part of the item.
   \item This is the third item of the list.
\end{itemize}
You can even display poetry.
\begin{verse}
   There is an environment
    for verse \\             % The \\ command separates lines
   Whose features some poets % within a stanza.
   will curse.

                             % One or more blank lines separate stanzas.

   For instead of making\\
   Them do \emph{all} line breaking, \\
   It allows them to put too many words on a line when they'd rather be
   forced to be terse.
\end{verse}

Mathematical formulas may also be displayed.  A
displayed formula
is
one-line long; multiline
formulas require special formatting instructions.
   \[  \ip{\Gamma}{\psi'} = x'' + y^{2} + z_{i}^{n}\]
Don't start a paragraph with a displayed equation,
nor make one a paragraph by itself.

\end{document}               % End of document.

\section{Bench Huge Section 457}
% synthetic bench section 457
% This is a sample LaTeX input file.  (Version of 12 August 2004.)
%
% A '%' character causes TeX to ignore all remaining text on the line,
% and is used for comments like this one.

\documentclass{article}      % Specifies the document class

                             % The preamble begins here.
\title{An Example Document}  % Declares the document's title.
\author{Leslie Lamport}      % Declares the author's name.
\date{January 21, 1994}      % Deleting this command produces today's date.

\newcommand{\ip}[2]{(#1, #2)}
                             % Defines \ip{arg1}{arg2} to mean
                             % (arg1, arg2).

%\newcommand{\ip}[2]{\langle #1 | #2\rangle}
                             % This is an alternative definition of
                             % \ip that is commented out.

\begin{document}             % End of preamble and beginning of text.

\maketitle                   % Produces the title.

This is an example input file.  Comparing it with
the output it generates can show you how to
produce a simple document of your own.

\section{Ordinary Text}      % Produces section heading.  Lower-level
                             % sections are begun with similar
                             % \subsection and \subsubsection commands.

The ends  of words and sentences are marked
  by   spaces. It  doesn't matter how many
spaces    you type; one is as good as 100.  The
end of   a line counts as a space.

One   or more   blank lines denote the  end
of  a paragraph.

Since any number of consecutive spaces are treated
like a single one, the formatting of the input
file makes no difference to
      \LaTeX,                % The \LaTeX command generates the LaTeX logo.
but it makes a difference to you.  When you use
\LaTeX, making your input file as easy to read
as possible will be a great help as you write
your document and when you change it.  This sample
file shows how you can add comments to your own input
file.

Because printing is different from typewriting,
there are a number of things that you have to do
differently when preparing an input file than if
you were just typing the document directly.
Quotation marks like
       ``this''
have to be handled specially, as do quotes within
quotes:
       ``\,`this'            % \, separates the double and single quote.
        is what I just
        wrote, not  `that'\,''.

Dashes come in three sizes: an
       intra-word
dash, a medium dash for number ranges like
       1--2,
and a punctuation
       dash---like
this.

A sentence-ending space should be larger than the
space between words within a sentence.  You
sometimes have to type special commands in
conjunction with punctuation characters to get
this right, as in the following sentence.
       Gnats, gnus, etc.\ all  % `\ ' makes an inter-word space.
       begin with G\@.         % \@ marks end-of-sentence punctuation.
You should check the spaces after periods when
reading your output to make sure you haven't
forgotten any special cases.  Generating an
ellipsis
       \ldots\               % `\ ' is needed after `\ldots' because TeX
                             % ignores spaces after command names like \ldots
                             % made from \ + letters.
                             %
                             % Note how a `%' character causes TeX to ignore
                             % the end of the input line, so these blank lines
                             % do not start a new paragraph.
                             %
with the right spacing around the periods requires
a special command.

\LaTeX\ interprets some common characters as
commands, so you must type special commands to
generate them.  These characters include the
following:
       \$ \& \% \# \{ and \}.

In printing, text is usually emphasized with an
       \emph{italic}
type style.

\begin{em}
   A long segment of text can also be emphasized
   in this way.  Text within such a segment can be
   given \emph{additional} emphasis.
\end{em}

It is sometimes necessary to prevent \LaTeX\ from
breaking a line where it might otherwise do so.
This may be at a space, as between the ``Mr.''\ and
``Jones'' in
       ``Mr.~Jones'',        % ~ produces an unbreakable interword space.
or within a word---especially when the word is a
symbol like
       \mbox{\emph{itemnum}}
that makes little sense when hyphenated across
lines.

Footnotes\footnote{This is an example of a footnote.}
pose no problem.

\LaTeX\ is good at typesetting mathematical formulas
like
       \( x-3y + z = 7 \)
or
       \( a_{1} > x^{2n} + y^{2n} > x' \)
or
       \( \ip{A}{B} = \sum_{i} a_{i} b_{i} \).
The spaces you type in a formula are
ignored.  Remember that a letter like
       $x$                   % $ ... $  and  \( ... \)  are equivalent
is a formula when it denotes a mathematical
symbol, and it should be typed as one.

\section{Displayed Text}

Text is displayed by indenting it from the left
margin.  Quotations are commonly displayed.  There
are short quotations
\begin{quote}
   This is a short quotation.  It consists of a
   single paragraph of text.  See how it is formatted.
\end{quote}
and longer ones.
\begin{quotation}
   This is a longer quotation.  It consists of two
   paragraphs of text, neither of which are
   particularly interesting.

   This is the second paragraph of the quotation.  It
   is just as dull as the first paragraph.
\end{quotation}
Another frequently-displayed structure is a list.
The following is an example of an \emph{itemized}
list.
\begin{itemize}
   \item This is the first item of an itemized list.
         Each item in the list is marked with a ``tick''.
         You don't have to worry about what kind of tick
         mark is used.

   \item This is the second item of the list.  It
         contains another list nested inside it.  The inner
         list is an \emph{enumerated} list.
         \begin{enumerate}
            \item This is the first item of an enumerated
                  list that is nested within the itemized list.

            \item This is the second item of the inner list.
                  \LaTeX\ allows you to nest lists deeper than
                  you really should.
         \end{enumerate}
         This is the rest of the second item of the outer
         list.  It is no more interesting than any other
         part of the item.
   \item This is the third item of the list.
\end{itemize}
You can even display poetry.
\begin{verse}
   There is an environment
    for verse \\             % The \\ command separates lines
   Whose features some poets % within a stanza.
   will curse.

                             % One or more blank lines separate stanzas.

   For instead of making\\
   Them do \emph{all} line breaking, \\
   It allows them to put too many words on a line when they'd rather be
   forced to be terse.
\end{verse}

Mathematical formulas may also be displayed.  A
displayed formula
is
one-line long; multiline
formulas require special formatting instructions.
   \[  \ip{\Gamma}{\psi'} = x'' + y^{2} + z_{i}^{n}\]
Don't start a paragraph with a displayed equation,
nor make one a paragraph by itself.

\end{document}               % End of document.

\section{Bench Huge Section 458}
% synthetic bench section 458
% This is a sample LaTeX input file.  (Version of 12 August 2004.)
%
% A '%' character causes TeX to ignore all remaining text on the line,
% and is used for comments like this one.

\documentclass{article}      % Specifies the document class

                             % The preamble begins here.
\title{An Example Document}  % Declares the document's title.
\author{Leslie Lamport}      % Declares the author's name.
\date{January 21, 1994}      % Deleting this command produces today's date.

\newcommand{\ip}[2]{(#1, #2)}
                             % Defines \ip{arg1}{arg2} to mean
                             % (arg1, arg2).

%\newcommand{\ip}[2]{\langle #1 | #2\rangle}
                             % This is an alternative definition of
                             % \ip that is commented out.

\begin{document}             % End of preamble and beginning of text.

\maketitle                   % Produces the title.

This is an example input file.  Comparing it with
the output it generates can show you how to
produce a simple document of your own.

\section{Ordinary Text}      % Produces section heading.  Lower-level
                             % sections are begun with similar
                             % \subsection and \subsubsection commands.

The ends  of words and sentences are marked
  by   spaces. It  doesn't matter how many
spaces    you type; one is as good as 100.  The
end of   a line counts as a space.

One   or more   blank lines denote the  end
of  a paragraph.

Since any number of consecutive spaces are treated
like a single one, the formatting of the input
file makes no difference to
      \LaTeX,                % The \LaTeX command generates the LaTeX logo.
but it makes a difference to you.  When you use
\LaTeX, making your input file as easy to read
as possible will be a great help as you write
your document and when you change it.  This sample
file shows how you can add comments to your own input
file.

Because printing is different from typewriting,
there are a number of things that you have to do
differently when preparing an input file than if
you were just typing the document directly.
Quotation marks like
       ``this''
have to be handled specially, as do quotes within
quotes:
       ``\,`this'            % \, separates the double and single quote.
        is what I just
        wrote, not  `that'\,''.

Dashes come in three sizes: an
       intra-word
dash, a medium dash for number ranges like
       1--2,
and a punctuation
       dash---like
this.

A sentence-ending space should be larger than the
space between words within a sentence.  You
sometimes have to type special commands in
conjunction with punctuation characters to get
this right, as in the following sentence.
       Gnats, gnus, etc.\ all  % `\ ' makes an inter-word space.
       begin with G\@.         % \@ marks end-of-sentence punctuation.
You should check the spaces after periods when
reading your output to make sure you haven't
forgotten any special cases.  Generating an
ellipsis
       \ldots\               % `\ ' is needed after `\ldots' because TeX
                             % ignores spaces after command names like \ldots
                             % made from \ + letters.
                             %
                             % Note how a `%' character causes TeX to ignore
                             % the end of the input line, so these blank lines
                             % do not start a new paragraph.
                             %
with the right spacing around the periods requires
a special command.

\LaTeX\ interprets some common characters as
commands, so you must type special commands to
generate them.  These characters include the
following:
       \$ \& \% \# \{ and \}.

In printing, text is usually emphasized with an
       \emph{italic}
type style.

\begin{em}
   A long segment of text can also be emphasized
   in this way.  Text within such a segment can be
   given \emph{additional} emphasis.
\end{em}

It is sometimes necessary to prevent \LaTeX\ from
breaking a line where it might otherwise do so.
This may be at a space, as between the ``Mr.''\ and
``Jones'' in
       ``Mr.~Jones'',        % ~ produces an unbreakable interword space.
or within a word---especially when the word is a
symbol like
       \mbox{\emph{itemnum}}
that makes little sense when hyphenated across
lines.

Footnotes\footnote{This is an example of a footnote.}
pose no problem.

\LaTeX\ is good at typesetting mathematical formulas
like
       \( x-3y + z = 7 \)
or
       \( a_{1} > x^{2n} + y^{2n} > x' \)
or
       \( \ip{A}{B} = \sum_{i} a_{i} b_{i} \).
The spaces you type in a formula are
ignored.  Remember that a letter like
       $x$                   % $ ... $  and  \( ... \)  are equivalent
is a formula when it denotes a mathematical
symbol, and it should be typed as one.

\section{Displayed Text}

Text is displayed by indenting it from the left
margin.  Quotations are commonly displayed.  There
are short quotations
\begin{quote}
   This is a short quotation.  It consists of a
   single paragraph of text.  See how it is formatted.
\end{quote}
and longer ones.
\begin{quotation}
   This is a longer quotation.  It consists of two
   paragraphs of text, neither of which are
   particularly interesting.

   This is the second paragraph of the quotation.  It
   is just as dull as the first paragraph.
\end{quotation}
Another frequently-displayed structure is a list.
The following is an example of an \emph{itemized}
list.
\begin{itemize}
   \item This is the first item of an itemized list.
         Each item in the list is marked with a ``tick''.
         You don't have to worry about what kind of tick
         mark is used.

   \item This is the second item of the list.  It
         contains another list nested inside it.  The inner
         list is an \emph{enumerated} list.
         \begin{enumerate}
            \item This is the first item of an enumerated
                  list that is nested within the itemized list.

            \item This is the second item of the inner list.
                  \LaTeX\ allows you to nest lists deeper than
                  you really should.
         \end{enumerate}
         This is the rest of the second item of the outer
         list.  It is no more interesting than any other
         part of the item.
   \item This is the third item of the list.
\end{itemize}
You can even display poetry.
\begin{verse}
   There is an environment
    for verse \\             % The \\ command separates lines
   Whose features some poets % within a stanza.
   will curse.

                             % One or more blank lines separate stanzas.

   For instead of making\\
   Them do \emph{all} line breaking, \\
   It allows them to put too many words on a line when they'd rather be
   forced to be terse.
\end{verse}

Mathematical formulas may also be displayed.  A
displayed formula
is
one-line long; multiline
formulas require special formatting instructions.
   \[  \ip{\Gamma}{\psi'} = x'' + y^{2} + z_{i}^{n}\]
Don't start a paragraph with a displayed equation,
nor make one a paragraph by itself.

\end{document}               % End of document.

\section{Bench Huge Section 459}
% synthetic bench section 459
% This is a sample LaTeX input file.  (Version of 12 August 2004.)
%
% A '%' character causes TeX to ignore all remaining text on the line,
% and is used for comments like this one.

\documentclass{article}      % Specifies the document class

                             % The preamble begins here.
\title{An Example Document}  % Declares the document's title.
\author{Leslie Lamport}      % Declares the author's name.
\date{January 21, 1994}      % Deleting this command produces today's date.

\newcommand{\ip}[2]{(#1, #2)}
                             % Defines \ip{arg1}{arg2} to mean
                             % (arg1, arg2).

%\newcommand{\ip}[2]{\langle #1 | #2\rangle}
                             % This is an alternative definition of
                             % \ip that is commented out.

\begin{document}             % End of preamble and beginning of text.

\maketitle                   % Produces the title.

This is an example input file.  Comparing it with
the output it generates can show you how to
produce a simple document of your own.

\section{Ordinary Text}      % Produces section heading.  Lower-level
                             % sections are begun with similar
                             % \subsection and \subsubsection commands.

The ends  of words and sentences are marked
  by   spaces. It  doesn't matter how many
spaces    you type; one is as good as 100.  The
end of   a line counts as a space.

One   or more   blank lines denote the  end
of  a paragraph.

Since any number of consecutive spaces are treated
like a single one, the formatting of the input
file makes no difference to
      \LaTeX,                % The \LaTeX command generates the LaTeX logo.
but it makes a difference to you.  When you use
\LaTeX, making your input file as easy to read
as possible will be a great help as you write
your document and when you change it.  This sample
file shows how you can add comments to your own input
file.

Because printing is different from typewriting,
there are a number of things that you have to do
differently when preparing an input file than if
you were just typing the document directly.
Quotation marks like
       ``this''
have to be handled specially, as do quotes within
quotes:
       ``\,`this'            % \, separates the double and single quote.
        is what I just
        wrote, not  `that'\,''.

Dashes come in three sizes: an
       intra-word
dash, a medium dash for number ranges like
       1--2,
and a punctuation
       dash---like
this.

A sentence-ending space should be larger than the
space between words within a sentence.  You
sometimes have to type special commands in
conjunction with punctuation characters to get
this right, as in the following sentence.
       Gnats, gnus, etc.\ all  % `\ ' makes an inter-word space.
       begin with G\@.         % \@ marks end-of-sentence punctuation.
You should check the spaces after periods when
reading your output to make sure you haven't
forgotten any special cases.  Generating an
ellipsis
       \ldots\               % `\ ' is needed after `\ldots' because TeX
                             % ignores spaces after command names like \ldots
                             % made from \ + letters.
                             %
                             % Note how a `%' character causes TeX to ignore
                             % the end of the input line, so these blank lines
                             % do not start a new paragraph.
                             %
with the right spacing around the periods requires
a special command.

\LaTeX\ interprets some common characters as
commands, so you must type special commands to
generate them.  These characters include the
following:
       \$ \& \% \# \{ and \}.

In printing, text is usually emphasized with an
       \emph{italic}
type style.

\begin{em}
   A long segment of text can also be emphasized
   in this way.  Text within such a segment can be
   given \emph{additional} emphasis.
\end{em}

It is sometimes necessary to prevent \LaTeX\ from
breaking a line where it might otherwise do so.
This may be at a space, as between the ``Mr.''\ and
``Jones'' in
       ``Mr.~Jones'',        % ~ produces an unbreakable interword space.
or within a word---especially when the word is a
symbol like
       \mbox{\emph{itemnum}}
that makes little sense when hyphenated across
lines.

Footnotes\footnote{This is an example of a footnote.}
pose no problem.

\LaTeX\ is good at typesetting mathematical formulas
like
       \( x-3y + z = 7 \)
or
       \( a_{1} > x^{2n} + y^{2n} > x' \)
or
       \( \ip{A}{B} = \sum_{i} a_{i} b_{i} \).
The spaces you type in a formula are
ignored.  Remember that a letter like
       $x$                   % $ ... $  and  \( ... \)  are equivalent
is a formula when it denotes a mathematical
symbol, and it should be typed as one.

\section{Displayed Text}

Text is displayed by indenting it from the left
margin.  Quotations are commonly displayed.  There
are short quotations
\begin{quote}
   This is a short quotation.  It consists of a
   single paragraph of text.  See how it is formatted.
\end{quote}
and longer ones.
\begin{quotation}
   This is a longer quotation.  It consists of two
   paragraphs of text, neither of which are
   particularly interesting.

   This is the second paragraph of the quotation.  It
   is just as dull as the first paragraph.
\end{quotation}
Another frequently-displayed structure is a list.
The following is an example of an \emph{itemized}
list.
\begin{itemize}
   \item This is the first item of an itemized list.
         Each item in the list is marked with a ``tick''.
         You don't have to worry about what kind of tick
         mark is used.

   \item This is the second item of the list.  It
         contains another list nested inside it.  The inner
         list is an \emph{enumerated} list.
         \begin{enumerate}
            \item This is the first item of an enumerated
                  list that is nested within the itemized list.

            \item This is the second item of the inner list.
                  \LaTeX\ allows you to nest lists deeper than
                  you really should.
         \end{enumerate}
         This is the rest of the second item of the outer
         list.  It is no more interesting than any other
         part of the item.
   \item This is the third item of the list.
\end{itemize}
You can even display poetry.
\begin{verse}
   There is an environment
    for verse \\             % The \\ command separates lines
   Whose features some poets % within a stanza.
   will curse.

                             % One or more blank lines separate stanzas.

   For instead of making\\
   Them do \emph{all} line breaking, \\
   It allows them to put too many words on a line when they'd rather be
   forced to be terse.
\end{verse}

Mathematical formulas may also be displayed.  A
displayed formula
is
one-line long; multiline
formulas require special formatting instructions.
   \[  \ip{\Gamma}{\psi'} = x'' + y^{2} + z_{i}^{n}\]
Don't start a paragraph with a displayed equation,
nor make one a paragraph by itself.

\end{document}               % End of document.

\section{Bench Huge Section 460}
% synthetic bench section 460
% This is a sample LaTeX input file.  (Version of 12 August 2004.)
%
% A '%' character causes TeX to ignore all remaining text on the line,
% and is used for comments like this one.

\documentclass{article}      % Specifies the document class

                             % The preamble begins here.
\title{An Example Document}  % Declares the document's title.
\author{Leslie Lamport}      % Declares the author's name.
\date{January 21, 1994}      % Deleting this command produces today's date.

\newcommand{\ip}[2]{(#1, #2)}
                             % Defines \ip{arg1}{arg2} to mean
                             % (arg1, arg2).

%\newcommand{\ip}[2]{\langle #1 | #2\rangle}
                             % This is an alternative definition of
                             % \ip that is commented out.

\begin{document}             % End of preamble and beginning of text.

\maketitle                   % Produces the title.

This is an example input file.  Comparing it with
the output it generates can show you how to
produce a simple document of your own.

\section{Ordinary Text}      % Produces section heading.  Lower-level
                             % sections are begun with similar
                             % \subsection and \subsubsection commands.

The ends  of words and sentences are marked
  by   spaces. It  doesn't matter how many
spaces    you type; one is as good as 100.  The
end of   a line counts as a space.

One   or more   blank lines denote the  end
of  a paragraph.

Since any number of consecutive spaces are treated
like a single one, the formatting of the input
file makes no difference to
      \LaTeX,                % The \LaTeX command generates the LaTeX logo.
but it makes a difference to you.  When you use
\LaTeX, making your input file as easy to read
as possible will be a great help as you write
your document and when you change it.  This sample
file shows how you can add comments to your own input
file.

Because printing is different from typewriting,
there are a number of things that you have to do
differently when preparing an input file than if
you were just typing the document directly.
Quotation marks like
       ``this''
have to be handled specially, as do quotes within
quotes:
       ``\,`this'            % \, separates the double and single quote.
        is what I just
        wrote, not  `that'\,''.

Dashes come in three sizes: an
       intra-word
dash, a medium dash for number ranges like
       1--2,
and a punctuation
       dash---like
this.

A sentence-ending space should be larger than the
space between words within a sentence.  You
sometimes have to type special commands in
conjunction with punctuation characters to get
this right, as in the following sentence.
       Gnats, gnus, etc.\ all  % `\ ' makes an inter-word space.
       begin with G\@.         % \@ marks end-of-sentence punctuation.
You should check the spaces after periods when
reading your output to make sure you haven't
forgotten any special cases.  Generating an
ellipsis
       \ldots\               % `\ ' is needed after `\ldots' because TeX
                             % ignores spaces after command names like \ldots
                             % made from \ + letters.
                             %
                             % Note how a `%' character causes TeX to ignore
                             % the end of the input line, so these blank lines
                             % do not start a new paragraph.
                             %
with the right spacing around the periods requires
a special command.

\LaTeX\ interprets some common characters as
commands, so you must type special commands to
generate them.  These characters include the
following:
       \$ \& \% \# \{ and \}.

In printing, text is usually emphasized with an
       \emph{italic}
type style.

\begin{em}
   A long segment of text can also be emphasized
   in this way.  Text within such a segment can be
   given \emph{additional} emphasis.
\end{em}

It is sometimes necessary to prevent \LaTeX\ from
breaking a line where it might otherwise do so.
This may be at a space, as between the ``Mr.''\ and
``Jones'' in
       ``Mr.~Jones'',        % ~ produces an unbreakable interword space.
or within a word---especially when the word is a
symbol like
       \mbox{\emph{itemnum}}
that makes little sense when hyphenated across
lines.

Footnotes\footnote{This is an example of a footnote.}
pose no problem.

\LaTeX\ is good at typesetting mathematical formulas
like
       \( x-3y + z = 7 \)
or
       \( a_{1} > x^{2n} + y^{2n} > x' \)
or
       \( \ip{A}{B} = \sum_{i} a_{i} b_{i} \).
The spaces you type in a formula are
ignored.  Remember that a letter like
       $x$                   % $ ... $  and  \( ... \)  are equivalent
is a formula when it denotes a mathematical
symbol, and it should be typed as one.

\section{Displayed Text}

Text is displayed by indenting it from the left
margin.  Quotations are commonly displayed.  There
are short quotations
\begin{quote}
   This is a short quotation.  It consists of a
   single paragraph of text.  See how it is formatted.
\end{quote}
and longer ones.
\begin{quotation}
   This is a longer quotation.  It consists of two
   paragraphs of text, neither of which are
   particularly interesting.

   This is the second paragraph of the quotation.  It
   is just as dull as the first paragraph.
\end{quotation}
Another frequently-displayed structure is a list.
The following is an example of an \emph{itemized}
list.
\begin{itemize}
   \item This is the first item of an itemized list.
         Each item in the list is marked with a ``tick''.
         You don't have to worry about what kind of tick
         mark is used.

   \item This is the second item of the list.  It
         contains another list nested inside it.  The inner
         list is an \emph{enumerated} list.
         \begin{enumerate}
            \item This is the first item of an enumerated
                  list that is nested within the itemized list.

            \item This is the second item of the inner list.
                  \LaTeX\ allows you to nest lists deeper than
                  you really should.
         \end{enumerate}
         This is the rest of the second item of the outer
         list.  It is no more interesting than any other
         part of the item.
   \item This is the third item of the list.
\end{itemize}
You can even display poetry.
\begin{verse}
   There is an environment
    for verse \\             % The \\ command separates lines
   Whose features some poets % within a stanza.
   will curse.

                             % One or more blank lines separate stanzas.

   For instead of making\\
   Them do \emph{all} line breaking, \\
   It allows them to put too many words on a line when they'd rather be
   forced to be terse.
\end{verse}

Mathematical formulas may also be displayed.  A
displayed formula
is
one-line long; multiline
formulas require special formatting instructions.
   \[  \ip{\Gamma}{\psi'} = x'' + y^{2} + z_{i}^{n}\]
Don't start a paragraph with a displayed equation,
nor make one a paragraph by itself.

\end{document}               % End of document.

\section{Bench Huge Section 461}
% synthetic bench section 461
% This is a sample LaTeX input file.  (Version of 12 August 2004.)
%
% A '%' character causes TeX to ignore all remaining text on the line,
% and is used for comments like this one.

\documentclass{article}      % Specifies the document class

                             % The preamble begins here.
\title{An Example Document}  % Declares the document's title.
\author{Leslie Lamport}      % Declares the author's name.
\date{January 21, 1994}      % Deleting this command produces today's date.

\newcommand{\ip}[2]{(#1, #2)}
                             % Defines \ip{arg1}{arg2} to mean
                             % (arg1, arg2).

%\newcommand{\ip}[2]{\langle #1 | #2\rangle}
                             % This is an alternative definition of
                             % \ip that is commented out.

\begin{document}             % End of preamble and beginning of text.

\maketitle                   % Produces the title.

This is an example input file.  Comparing it with
the output it generates can show you how to
produce a simple document of your own.

\section{Ordinary Text}      % Produces section heading.  Lower-level
                             % sections are begun with similar
                             % \subsection and \subsubsection commands.

The ends  of words and sentences are marked
  by   spaces. It  doesn't matter how many
spaces    you type; one is as good as 100.  The
end of   a line counts as a space.

One   or more   blank lines denote the  end
of  a paragraph.

Since any number of consecutive spaces are treated
like a single one, the formatting of the input
file makes no difference to
      \LaTeX,                % The \LaTeX command generates the LaTeX logo.
but it makes a difference to you.  When you use
\LaTeX, making your input file as easy to read
as possible will be a great help as you write
your document and when you change it.  This sample
file shows how you can add comments to your own input
file.

Because printing is different from typewriting,
there are a number of things that you have to do
differently when preparing an input file than if
you were just typing the document directly.
Quotation marks like
       ``this''
have to be handled specially, as do quotes within
quotes:
       ``\,`this'            % \, separates the double and single quote.
        is what I just
        wrote, not  `that'\,''.

Dashes come in three sizes: an
       intra-word
dash, a medium dash for number ranges like
       1--2,
and a punctuation
       dash---like
this.

A sentence-ending space should be larger than the
space between words within a sentence.  You
sometimes have to type special commands in
conjunction with punctuation characters to get
this right, as in the following sentence.
       Gnats, gnus, etc.\ all  % `\ ' makes an inter-word space.
       begin with G\@.         % \@ marks end-of-sentence punctuation.
You should check the spaces after periods when
reading your output to make sure you haven't
forgotten any special cases.  Generating an
ellipsis
       \ldots\               % `\ ' is needed after `\ldots' because TeX
                             % ignores spaces after command names like \ldots
                             % made from \ + letters.
                             %
                             % Note how a `%' character causes TeX to ignore
                             % the end of the input line, so these blank lines
                             % do not start a new paragraph.
                             %
with the right spacing around the periods requires
a special command.

\LaTeX\ interprets some common characters as
commands, so you must type special commands to
generate them.  These characters include the
following:
       \$ \& \% \# \{ and \}.

In printing, text is usually emphasized with an
       \emph{italic}
type style.

\begin{em}
   A long segment of text can also be emphasized
   in this way.  Text within such a segment can be
   given \emph{additional} emphasis.
\end{em}

It is sometimes necessary to prevent \LaTeX\ from
breaking a line where it might otherwise do so.
This may be at a space, as between the ``Mr.''\ and
``Jones'' in
       ``Mr.~Jones'',        % ~ produces an unbreakable interword space.
or within a word---especially when the word is a
symbol like
       \mbox{\emph{itemnum}}
that makes little sense when hyphenated across
lines.

Footnotes\footnote{This is an example of a footnote.}
pose no problem.

\LaTeX\ is good at typesetting mathematical formulas
like
       \( x-3y + z = 7 \)
or
       \( a_{1} > x^{2n} + y^{2n} > x' \)
or
       \( \ip{A}{B} = \sum_{i} a_{i} b_{i} \).
The spaces you type in a formula are
ignored.  Remember that a letter like
       $x$                   % $ ... $  and  \( ... \)  are equivalent
is a formula when it denotes a mathematical
symbol, and it should be typed as one.

\section{Displayed Text}

Text is displayed by indenting it from the left
margin.  Quotations are commonly displayed.  There
are short quotations
\begin{quote}
   This is a short quotation.  It consists of a
   single paragraph of text.  See how it is formatted.
\end{quote}
and longer ones.
\begin{quotation}
   This is a longer quotation.  It consists of two
   paragraphs of text, neither of which are
   particularly interesting.

   This is the second paragraph of the quotation.  It
   is just as dull as the first paragraph.
\end{quotation}
Another frequently-displayed structure is a list.
The following is an example of an \emph{itemized}
list.
\begin{itemize}
   \item This is the first item of an itemized list.
         Each item in the list is marked with a ``tick''.
         You don't have to worry about what kind of tick
         mark is used.

   \item This is the second item of the list.  It
         contains another list nested inside it.  The inner
         list is an \emph{enumerated} list.
         \begin{enumerate}
            \item This is the first item of an enumerated
                  list that is nested within the itemized list.

            \item This is the second item of the inner list.
                  \LaTeX\ allows you to nest lists deeper than
                  you really should.
         \end{enumerate}
         This is the rest of the second item of the outer
         list.  It is no more interesting than any other
         part of the item.
   \item This is the third item of the list.
\end{itemize}
You can even display poetry.
\begin{verse}
   There is an environment
    for verse \\             % The \\ command separates lines
   Whose features some poets % within a stanza.
   will curse.

                             % One or more blank lines separate stanzas.

   For instead of making\\
   Them do \emph{all} line breaking, \\
   It allows them to put too many words on a line when they'd rather be
   forced to be terse.
\end{verse}

Mathematical formulas may also be displayed.  A
displayed formula
is
one-line long; multiline
formulas require special formatting instructions.
   \[  \ip{\Gamma}{\psi'} = x'' + y^{2} + z_{i}^{n}\]
Don't start a paragraph with a displayed equation,
nor make one a paragraph by itself.

\end{document}               % End of document.

\section{Bench Huge Section 462}
% synthetic bench section 462
% This is a sample LaTeX input file.  (Version of 12 August 2004.)
%
% A '%' character causes TeX to ignore all remaining text on the line,
% and is used for comments like this one.

\documentclass{article}      % Specifies the document class

                             % The preamble begins here.
\title{An Example Document}  % Declares the document's title.
\author{Leslie Lamport}      % Declares the author's name.
\date{January 21, 1994}      % Deleting this command produces today's date.

\newcommand{\ip}[2]{(#1, #2)}
                             % Defines \ip{arg1}{arg2} to mean
                             % (arg1, arg2).

%\newcommand{\ip}[2]{\langle #1 | #2\rangle}
                             % This is an alternative definition of
                             % \ip that is commented out.

\begin{document}             % End of preamble and beginning of text.

\maketitle                   % Produces the title.

This is an example input file.  Comparing it with
the output it generates can show you how to
produce a simple document of your own.

\section{Ordinary Text}      % Produces section heading.  Lower-level
                             % sections are begun with similar
                             % \subsection and \subsubsection commands.

The ends  of words and sentences are marked
  by   spaces. It  doesn't matter how many
spaces    you type; one is as good as 100.  The
end of   a line counts as a space.

One   or more   blank lines denote the  end
of  a paragraph.

Since any number of consecutive spaces are treated
like a single one, the formatting of the input
file makes no difference to
      \LaTeX,                % The \LaTeX command generates the LaTeX logo.
but it makes a difference to you.  When you use
\LaTeX, making your input file as easy to read
as possible will be a great help as you write
your document and when you change it.  This sample
file shows how you can add comments to your own input
file.

Because printing is different from typewriting,
there are a number of things that you have to do
differently when preparing an input file than if
you were just typing the document directly.
Quotation marks like
       ``this''
have to be handled specially, as do quotes within
quotes:
       ``\,`this'            % \, separates the double and single quote.
        is what I just
        wrote, not  `that'\,''.

Dashes come in three sizes: an
       intra-word
dash, a medium dash for number ranges like
       1--2,
and a punctuation
       dash---like
this.

A sentence-ending space should be larger than the
space between words within a sentence.  You
sometimes have to type special commands in
conjunction with punctuation characters to get
this right, as in the following sentence.
       Gnats, gnus, etc.\ all  % `\ ' makes an inter-word space.
       begin with G\@.         % \@ marks end-of-sentence punctuation.
You should check the spaces after periods when
reading your output to make sure you haven't
forgotten any special cases.  Generating an
ellipsis
       \ldots\               % `\ ' is needed after `\ldots' because TeX
                             % ignores spaces after command names like \ldots
                             % made from \ + letters.
                             %
                             % Note how a `%' character causes TeX to ignore
                             % the end of the input line, so these blank lines
                             % do not start a new paragraph.
                             %
with the right spacing around the periods requires
a special command.

\LaTeX\ interprets some common characters as
commands, so you must type special commands to
generate them.  These characters include the
following:
       \$ \& \% \# \{ and \}.

In printing, text is usually emphasized with an
       \emph{italic}
type style.

\begin{em}
   A long segment of text can also be emphasized
   in this way.  Text within such a segment can be
   given \emph{additional} emphasis.
\end{em}

It is sometimes necessary to prevent \LaTeX\ from
breaking a line where it might otherwise do so.
This may be at a space, as between the ``Mr.''\ and
``Jones'' in
       ``Mr.~Jones'',        % ~ produces an unbreakable interword space.
or within a word---especially when the word is a
symbol like
       \mbox{\emph{itemnum}}
that makes little sense when hyphenated across
lines.

Footnotes\footnote{This is an example of a footnote.}
pose no problem.

\LaTeX\ is good at typesetting mathematical formulas
like
       \( x-3y + z = 7 \)
or
       \( a_{1} > x^{2n} + y^{2n} > x' \)
or
       \( \ip{A}{B} = \sum_{i} a_{i} b_{i} \).
The spaces you type in a formula are
ignored.  Remember that a letter like
       $x$                   % $ ... $  and  \( ... \)  are equivalent
is a formula when it denotes a mathematical
symbol, and it should be typed as one.

\section{Displayed Text}

Text is displayed by indenting it from the left
margin.  Quotations are commonly displayed.  There
are short quotations
\begin{quote}
   This is a short quotation.  It consists of a
   single paragraph of text.  See how it is formatted.
\end{quote}
and longer ones.
\begin{quotation}
   This is a longer quotation.  It consists of two
   paragraphs of text, neither of which are
   particularly interesting.

   This is the second paragraph of the quotation.  It
   is just as dull as the first paragraph.
\end{quotation}
Another frequently-displayed structure is a list.
The following is an example of an \emph{itemized}
list.
\begin{itemize}
   \item This is the first item of an itemized list.
         Each item in the list is marked with a ``tick''.
         You don't have to worry about what kind of tick
         mark is used.

   \item This is the second item of the list.  It
         contains another list nested inside it.  The inner
         list is an \emph{enumerated} list.
         \begin{enumerate}
            \item This is the first item of an enumerated
                  list that is nested within the itemized list.

            \item This is the second item of the inner list.
                  \LaTeX\ allows you to nest lists deeper than
                  you really should.
         \end{enumerate}
         This is the rest of the second item of the outer
         list.  It is no more interesting than any other
         part of the item.
   \item This is the third item of the list.
\end{itemize}
You can even display poetry.
\begin{verse}
   There is an environment
    for verse \\             % The \\ command separates lines
   Whose features some poets % within a stanza.
   will curse.

                             % One or more blank lines separate stanzas.

   For instead of making\\
   Them do \emph{all} line breaking, \\
   It allows them to put too many words on a line when they'd rather be
   forced to be terse.
\end{verse}

Mathematical formulas may also be displayed.  A
displayed formula
is
one-line long; multiline
formulas require special formatting instructions.
   \[  \ip{\Gamma}{\psi'} = x'' + y^{2} + z_{i}^{n}\]
Don't start a paragraph with a displayed equation,
nor make one a paragraph by itself.

\end{document}               % End of document.

\section{Bench Huge Section 463}
% synthetic bench section 463
% This is a sample LaTeX input file.  (Version of 12 August 2004.)
%
% A '%' character causes TeX to ignore all remaining text on the line,
% and is used for comments like this one.

\documentclass{article}      % Specifies the document class

                             % The preamble begins here.
\title{An Example Document}  % Declares the document's title.
\author{Leslie Lamport}      % Declares the author's name.
\date{January 21, 1994}      % Deleting this command produces today's date.

\newcommand{\ip}[2]{(#1, #2)}
                             % Defines \ip{arg1}{arg2} to mean
                             % (arg1, arg2).

%\newcommand{\ip}[2]{\langle #1 | #2\rangle}
                             % This is an alternative definition of
                             % \ip that is commented out.

\begin{document}             % End of preamble and beginning of text.

\maketitle                   % Produces the title.

This is an example input file.  Comparing it with
the output it generates can show you how to
produce a simple document of your own.

\section{Ordinary Text}      % Produces section heading.  Lower-level
                             % sections are begun with similar
                             % \subsection and \subsubsection commands.

The ends  of words and sentences are marked
  by   spaces. It  doesn't matter how many
spaces    you type; one is as good as 100.  The
end of   a line counts as a space.

One   or more   blank lines denote the  end
of  a paragraph.

Since any number of consecutive spaces are treated
like a single one, the formatting of the input
file makes no difference to
      \LaTeX,                % The \LaTeX command generates the LaTeX logo.
but it makes a difference to you.  When you use
\LaTeX, making your input file as easy to read
as possible will be a great help as you write
your document and when you change it.  This sample
file shows how you can add comments to your own input
file.

Because printing is different from typewriting,
there are a number of things that you have to do
differently when preparing an input file than if
you were just typing the document directly.
Quotation marks like
       ``this''
have to be handled specially, as do quotes within
quotes:
       ``\,`this'            % \, separates the double and single quote.
        is what I just
        wrote, not  `that'\,''.

Dashes come in three sizes: an
       intra-word
dash, a medium dash for number ranges like
       1--2,
and a punctuation
       dash---like
this.

A sentence-ending space should be larger than the
space between words within a sentence.  You
sometimes have to type special commands in
conjunction with punctuation characters to get
this right, as in the following sentence.
       Gnats, gnus, etc.\ all  % `\ ' makes an inter-word space.
       begin with G\@.         % \@ marks end-of-sentence punctuation.
You should check the spaces after periods when
reading your output to make sure you haven't
forgotten any special cases.  Generating an
ellipsis
       \ldots\               % `\ ' is needed after `\ldots' because TeX
                             % ignores spaces after command names like \ldots
                             % made from \ + letters.
                             %
                             % Note how a `%' character causes TeX to ignore
                             % the end of the input line, so these blank lines
                             % do not start a new paragraph.
                             %
with the right spacing around the periods requires
a special command.

\LaTeX\ interprets some common characters as
commands, so you must type special commands to
generate them.  These characters include the
following:
       \$ \& \% \# \{ and \}.

In printing, text is usually emphasized with an
       \emph{italic}
type style.

\begin{em}
   A long segment of text can also be emphasized
   in this way.  Text within such a segment can be
   given \emph{additional} emphasis.
\end{em}

It is sometimes necessary to prevent \LaTeX\ from
breaking a line where it might otherwise do so.
This may be at a space, as between the ``Mr.''\ and
``Jones'' in
       ``Mr.~Jones'',        % ~ produces an unbreakable interword space.
or within a word---especially when the word is a
symbol like
       \mbox{\emph{itemnum}}
that makes little sense when hyphenated across
lines.

Footnotes\footnote{This is an example of a footnote.}
pose no problem.

\LaTeX\ is good at typesetting mathematical formulas
like
       \( x-3y + z = 7 \)
or
       \( a_{1} > x^{2n} + y^{2n} > x' \)
or
       \( \ip{A}{B} = \sum_{i} a_{i} b_{i} \).
The spaces you type in a formula are
ignored.  Remember that a letter like
       $x$                   % $ ... $  and  \( ... \)  are equivalent
is a formula when it denotes a mathematical
symbol, and it should be typed as one.

\section{Displayed Text}

Text is displayed by indenting it from the left
margin.  Quotations are commonly displayed.  There
are short quotations
\begin{quote}
   This is a short quotation.  It consists of a
   single paragraph of text.  See how it is formatted.
\end{quote}
and longer ones.
\begin{quotation}
   This is a longer quotation.  It consists of two
   paragraphs of text, neither of which are
   particularly interesting.

   This is the second paragraph of the quotation.  It
   is just as dull as the first paragraph.
\end{quotation}
Another frequently-displayed structure is a list.
The following is an example of an \emph{itemized}
list.
\begin{itemize}
   \item This is the first item of an itemized list.
         Each item in the list is marked with a ``tick''.
         You don't have to worry about what kind of tick
         mark is used.

   \item This is the second item of the list.  It
         contains another list nested inside it.  The inner
         list is an \emph{enumerated} list.
         \begin{enumerate}
            \item This is the first item of an enumerated
                  list that is nested within the itemized list.

            \item This is the second item of the inner list.
                  \LaTeX\ allows you to nest lists deeper than
                  you really should.
         \end{enumerate}
         This is the rest of the second item of the outer
         list.  It is no more interesting than any other
         part of the item.
   \item This is the third item of the list.
\end{itemize}
You can even display poetry.
\begin{verse}
   There is an environment
    for verse \\             % The \\ command separates lines
   Whose features some poets % within a stanza.
   will curse.

                             % One or more blank lines separate stanzas.

   For instead of making\\
   Them do \emph{all} line breaking, \\
   It allows them to put too many words on a line when they'd rather be
   forced to be terse.
\end{verse}

Mathematical formulas may also be displayed.  A
displayed formula
is
one-line long; multiline
formulas require special formatting instructions.
   \[  \ip{\Gamma}{\psi'} = x'' + y^{2} + z_{i}^{n}\]
Don't start a paragraph with a displayed equation,
nor make one a paragraph by itself.

\end{document}               % End of document.

\section{Bench Huge Section 464}
% synthetic bench section 464
% This is a sample LaTeX input file.  (Version of 12 August 2004.)
%
% A '%' character causes TeX to ignore all remaining text on the line,
% and is used for comments like this one.

\documentclass{article}      % Specifies the document class

                             % The preamble begins here.
\title{An Example Document}  % Declares the document's title.
\author{Leslie Lamport}      % Declares the author's name.
\date{January 21, 1994}      % Deleting this command produces today's date.

\newcommand{\ip}[2]{(#1, #2)}
                             % Defines \ip{arg1}{arg2} to mean
                             % (arg1, arg2).

%\newcommand{\ip}[2]{\langle #1 | #2\rangle}
                             % This is an alternative definition of
                             % \ip that is commented out.

\begin{document}             % End of preamble and beginning of text.

\maketitle                   % Produces the title.

This is an example input file.  Comparing it with
the output it generates can show you how to
produce a simple document of your own.

\section{Ordinary Text}      % Produces section heading.  Lower-level
                             % sections are begun with similar
                             % \subsection and \subsubsection commands.

The ends  of words and sentences are marked
  by   spaces. It  doesn't matter how many
spaces    you type; one is as good as 100.  The
end of   a line counts as a space.

One   or more   blank lines denote the  end
of  a paragraph.

Since any number of consecutive spaces are treated
like a single one, the formatting of the input
file makes no difference to
      \LaTeX,                % The \LaTeX command generates the LaTeX logo.
but it makes a difference to you.  When you use
\LaTeX, making your input file as easy to read
as possible will be a great help as you write
your document and when you change it.  This sample
file shows how you can add comments to your own input
file.

Because printing is different from typewriting,
there are a number of things that you have to do
differently when preparing an input file than if
you were just typing the document directly.
Quotation marks like
       ``this''
have to be handled specially, as do quotes within
quotes:
       ``\,`this'            % \, separates the double and single quote.
        is what I just
        wrote, not  `that'\,''.

Dashes come in three sizes: an
       intra-word
dash, a medium dash for number ranges like
       1--2,
and a punctuation
       dash---like
this.

A sentence-ending space should be larger than the
space between words within a sentence.  You
sometimes have to type special commands in
conjunction with punctuation characters to get
this right, as in the following sentence.
       Gnats, gnus, etc.\ all  % `\ ' makes an inter-word space.
       begin with G\@.         % \@ marks end-of-sentence punctuation.
You should check the spaces after periods when
reading your output to make sure you haven't
forgotten any special cases.  Generating an
ellipsis
       \ldots\               % `\ ' is needed after `\ldots' because TeX
                             % ignores spaces after command names like \ldots
                             % made from \ + letters.
                             %
                             % Note how a `%' character causes TeX to ignore
                             % the end of the input line, so these blank lines
                             % do not start a new paragraph.
                             %
with the right spacing around the periods requires
a special command.

\LaTeX\ interprets some common characters as
commands, so you must type special commands to
generate them.  These characters include the
following:
       \$ \& \% \# \{ and \}.

In printing, text is usually emphasized with an
       \emph{italic}
type style.

\begin{em}
   A long segment of text can also be emphasized
   in this way.  Text within such a segment can be
   given \emph{additional} emphasis.
\end{em}

It is sometimes necessary to prevent \LaTeX\ from
breaking a line where it might otherwise do so.
This may be at a space, as between the ``Mr.''\ and
``Jones'' in
       ``Mr.~Jones'',        % ~ produces an unbreakable interword space.
or within a word---especially when the word is a
symbol like
       \mbox{\emph{itemnum}}
that makes little sense when hyphenated across
lines.

Footnotes\footnote{This is an example of a footnote.}
pose no problem.

\LaTeX\ is good at typesetting mathematical formulas
like
       \( x-3y + z = 7 \)
or
       \( a_{1} > x^{2n} + y^{2n} > x' \)
or
       \( \ip{A}{B} = \sum_{i} a_{i} b_{i} \).
The spaces you type in a formula are
ignored.  Remember that a letter like
       $x$                   % $ ... $  and  \( ... \)  are equivalent
is a formula when it denotes a mathematical
symbol, and it should be typed as one.

\section{Displayed Text}

Text is displayed by indenting it from the left
margin.  Quotations are commonly displayed.  There
are short quotations
\begin{quote}
   This is a short quotation.  It consists of a
   single paragraph of text.  See how it is formatted.
\end{quote}
and longer ones.
\begin{quotation}
   This is a longer quotation.  It consists of two
   paragraphs of text, neither of which are
   particularly interesting.

   This is the second paragraph of the quotation.  It
   is just as dull as the first paragraph.
\end{quotation}
Another frequently-displayed structure is a list.
The following is an example of an \emph{itemized}
list.
\begin{itemize}
   \item This is the first item of an itemized list.
         Each item in the list is marked with a ``tick''.
         You don't have to worry about what kind of tick
         mark is used.

   \item This is the second item of the list.  It
         contains another list nested inside it.  The inner
         list is an \emph{enumerated} list.
         \begin{enumerate}
            \item This is the first item of an enumerated
                  list that is nested within the itemized list.

            \item This is the second item of the inner list.
                  \LaTeX\ allows you to nest lists deeper than
                  you really should.
         \end{enumerate}
         This is the rest of the second item of the outer
         list.  It is no more interesting than any other
         part of the item.
   \item This is the third item of the list.
\end{itemize}
You can even display poetry.
\begin{verse}
   There is an environment
    for verse \\             % The \\ command separates lines
   Whose features some poets % within a stanza.
   will curse.

                             % One or more blank lines separate stanzas.

   For instead of making\\
   Them do \emph{all} line breaking, \\
   It allows them to put too many words on a line when they'd rather be
   forced to be terse.
\end{verse}

Mathematical formulas may also be displayed.  A
displayed formula
is
one-line long; multiline
formulas require special formatting instructions.
   \[  \ip{\Gamma}{\psi'} = x'' + y^{2} + z_{i}^{n}\]
Don't start a paragraph with a displayed equation,
nor make one a paragraph by itself.

\end{document}               % End of document.

\section{Bench Huge Section 465}
% synthetic bench section 465
% This is a sample LaTeX input file.  (Version of 12 August 2004.)
%
% A '%' character causes TeX to ignore all remaining text on the line,
% and is used for comments like this one.

\documentclass{article}      % Specifies the document class

                             % The preamble begins here.
\title{An Example Document}  % Declares the document's title.
\author{Leslie Lamport}      % Declares the author's name.
\date{January 21, 1994}      % Deleting this command produces today's date.

\newcommand{\ip}[2]{(#1, #2)}
                             % Defines \ip{arg1}{arg2} to mean
                             % (arg1, arg2).

%\newcommand{\ip}[2]{\langle #1 | #2\rangle}
                             % This is an alternative definition of
                             % \ip that is commented out.

\begin{document}             % End of preamble and beginning of text.

\maketitle                   % Produces the title.

This is an example input file.  Comparing it with
the output it generates can show you how to
produce a simple document of your own.

\section{Ordinary Text}      % Produces section heading.  Lower-level
                             % sections are begun with similar
                             % \subsection and \subsubsection commands.

The ends  of words and sentences are marked
  by   spaces. It  doesn't matter how many
spaces    you type; one is as good as 100.  The
end of   a line counts as a space.

One   or more   blank lines denote the  end
of  a paragraph.

Since any number of consecutive spaces are treated
like a single one, the formatting of the input
file makes no difference to
      \LaTeX,                % The \LaTeX command generates the LaTeX logo.
but it makes a difference to you.  When you use
\LaTeX, making your input file as easy to read
as possible will be a great help as you write
your document and when you change it.  This sample
file shows how you can add comments to your own input
file.

Because printing is different from typewriting,
there are a number of things that you have to do
differently when preparing an input file than if
you were just typing the document directly.
Quotation marks like
       ``this''
have to be handled specially, as do quotes within
quotes:
       ``\,`this'            % \, separates the double and single quote.
        is what I just
        wrote, not  `that'\,''.

Dashes come in three sizes: an
       intra-word
dash, a medium dash for number ranges like
       1--2,
and a punctuation
       dash---like
this.

A sentence-ending space should be larger than the
space between words within a sentence.  You
sometimes have to type special commands in
conjunction with punctuation characters to get
this right, as in the following sentence.
       Gnats, gnus, etc.\ all  % `\ ' makes an inter-word space.
       begin with G\@.         % \@ marks end-of-sentence punctuation.
You should check the spaces after periods when
reading your output to make sure you haven't
forgotten any special cases.  Generating an
ellipsis
       \ldots\               % `\ ' is needed after `\ldots' because TeX
                             % ignores spaces after command names like \ldots
                             % made from \ + letters.
                             %
                             % Note how a `%' character causes TeX to ignore
                             % the end of the input line, so these blank lines
                             % do not start a new paragraph.
                             %
with the right spacing around the periods requires
a special command.

\LaTeX\ interprets some common characters as
commands, so you must type special commands to
generate them.  These characters include the
following:
       \$ \& \% \# \{ and \}.

In printing, text is usually emphasized with an
       \emph{italic}
type style.

\begin{em}
   A long segment of text can also be emphasized
   in this way.  Text within such a segment can be
   given \emph{additional} emphasis.
\end{em}

It is sometimes necessary to prevent \LaTeX\ from
breaking a line where it might otherwise do so.
This may be at a space, as between the ``Mr.''\ and
``Jones'' in
       ``Mr.~Jones'',        % ~ produces an unbreakable interword space.
or within a word---especially when the word is a
symbol like
       \mbox{\emph{itemnum}}
that makes little sense when hyphenated across
lines.

Footnotes\footnote{This is an example of a footnote.}
pose no problem.

\LaTeX\ is good at typesetting mathematical formulas
like
       \( x-3y + z = 7 \)
or
       \( a_{1} > x^{2n} + y^{2n} > x' \)
or
       \( \ip{A}{B} = \sum_{i} a_{i} b_{i} \).
The spaces you type in a formula are
ignored.  Remember that a letter like
       $x$                   % $ ... $  and  \( ... \)  are equivalent
is a formula when it denotes a mathematical
symbol, and it should be typed as one.

\section{Displayed Text}

Text is displayed by indenting it from the left
margin.  Quotations are commonly displayed.  There
are short quotations
\begin{quote}
   This is a short quotation.  It consists of a
   single paragraph of text.  See how it is formatted.
\end{quote}
and longer ones.
\begin{quotation}
   This is a longer quotation.  It consists of two
   paragraphs of text, neither of which are
   particularly interesting.

   This is the second paragraph of the quotation.  It
   is just as dull as the first paragraph.
\end{quotation}
Another frequently-displayed structure is a list.
The following is an example of an \emph{itemized}
list.
\begin{itemize}
   \item This is the first item of an itemized list.
         Each item in the list is marked with a ``tick''.
         You don't have to worry about what kind of tick
         mark is used.

   \item This is the second item of the list.  It
         contains another list nested inside it.  The inner
         list is an \emph{enumerated} list.
         \begin{enumerate}
            \item This is the first item of an enumerated
                  list that is nested within the itemized list.

            \item This is the second item of the inner list.
                  \LaTeX\ allows you to nest lists deeper than
                  you really should.
         \end{enumerate}
         This is the rest of the second item of the outer
         list.  It is no more interesting than any other
         part of the item.
   \item This is the third item of the list.
\end{itemize}
You can even display poetry.
\begin{verse}
   There is an environment
    for verse \\             % The \\ command separates lines
   Whose features some poets % within a stanza.
   will curse.

                             % One or more blank lines separate stanzas.

   For instead of making\\
   Them do \emph{all} line breaking, \\
   It allows them to put too many words on a line when they'd rather be
   forced to be terse.
\end{verse}

Mathematical formulas may also be displayed.  A
displayed formula
is
one-line long; multiline
formulas require special formatting instructions.
   \[  \ip{\Gamma}{\psi'} = x'' + y^{2} + z_{i}^{n}\]
Don't start a paragraph with a displayed equation,
nor make one a paragraph by itself.

\end{document}               % End of document.

\section{Bench Huge Section 466}
% synthetic bench section 466
% This is a sample LaTeX input file.  (Version of 12 August 2004.)
%
% A '%' character causes TeX to ignore all remaining text on the line,
% and is used for comments like this one.

\documentclass{article}      % Specifies the document class

                             % The preamble begins here.
\title{An Example Document}  % Declares the document's title.
\author{Leslie Lamport}      % Declares the author's name.
\date{January 21, 1994}      % Deleting this command produces today's date.

\newcommand{\ip}[2]{(#1, #2)}
                             % Defines \ip{arg1}{arg2} to mean
                             % (arg1, arg2).

%\newcommand{\ip}[2]{\langle #1 | #2\rangle}
                             % This is an alternative definition of
                             % \ip that is commented out.

\begin{document}             % End of preamble and beginning of text.

\maketitle                   % Produces the title.

This is an example input file.  Comparing it with
the output it generates can show you how to
produce a simple document of your own.

\section{Ordinary Text}      % Produces section heading.  Lower-level
                             % sections are begun with similar
                             % \subsection and \subsubsection commands.

The ends  of words and sentences are marked
  by   spaces. It  doesn't matter how many
spaces    you type; one is as good as 100.  The
end of   a line counts as a space.

One   or more   blank lines denote the  end
of  a paragraph.

Since any number of consecutive spaces are treated
like a single one, the formatting of the input
file makes no difference to
      \LaTeX,                % The \LaTeX command generates the LaTeX logo.
but it makes a difference to you.  When you use
\LaTeX, making your input file as easy to read
as possible will be a great help as you write
your document and when you change it.  This sample
file shows how you can add comments to your own input
file.

Because printing is different from typewriting,
there are a number of things that you have to do
differently when preparing an input file than if
you were just typing the document directly.
Quotation marks like
       ``this''
have to be handled specially, as do quotes within
quotes:
       ``\,`this'            % \, separates the double and single quote.
        is what I just
        wrote, not  `that'\,''.

Dashes come in three sizes: an
       intra-word
dash, a medium dash for number ranges like
       1--2,
and a punctuation
       dash---like
this.

A sentence-ending space should be larger than the
space between words within a sentence.  You
sometimes have to type special commands in
conjunction with punctuation characters to get
this right, as in the following sentence.
       Gnats, gnus, etc.\ all  % `\ ' makes an inter-word space.
       begin with G\@.         % \@ marks end-of-sentence punctuation.
You should check the spaces after periods when
reading your output to make sure you haven't
forgotten any special cases.  Generating an
ellipsis
       \ldots\               % `\ ' is needed after `\ldots' because TeX
                             % ignores spaces after command names like \ldots
                             % made from \ + letters.
                             %
                             % Note how a `%' character causes TeX to ignore
                             % the end of the input line, so these blank lines
                             % do not start a new paragraph.
                             %
with the right spacing around the periods requires
a special command.

\LaTeX\ interprets some common characters as
commands, so you must type special commands to
generate them.  These characters include the
following:
       \$ \& \% \# \{ and \}.

In printing, text is usually emphasized with an
       \emph{italic}
type style.

\begin{em}
   A long segment of text can also be emphasized
   in this way.  Text within such a segment can be
   given \emph{additional} emphasis.
\end{em}

It is sometimes necessary to prevent \LaTeX\ from
breaking a line where it might otherwise do so.
This may be at a space, as between the ``Mr.''\ and
``Jones'' in
       ``Mr.~Jones'',        % ~ produces an unbreakable interword space.
or within a word---especially when the word is a
symbol like
       \mbox{\emph{itemnum}}
that makes little sense when hyphenated across
lines.

Footnotes\footnote{This is an example of a footnote.}
pose no problem.

\LaTeX\ is good at typesetting mathematical formulas
like
       \( x-3y + z = 7 \)
or
       \( a_{1} > x^{2n} + y^{2n} > x' \)
or
       \( \ip{A}{B} = \sum_{i} a_{i} b_{i} \).
The spaces you type in a formula are
ignored.  Remember that a letter like
       $x$                   % $ ... $  and  \( ... \)  are equivalent
is a formula when it denotes a mathematical
symbol, and it should be typed as one.

\section{Displayed Text}

Text is displayed by indenting it from the left
margin.  Quotations are commonly displayed.  There
are short quotations
\begin{quote}
   This is a short quotation.  It consists of a
   single paragraph of text.  See how it is formatted.
\end{quote}
and longer ones.
\begin{quotation}
   This is a longer quotation.  It consists of two
   paragraphs of text, neither of which are
   particularly interesting.

   This is the second paragraph of the quotation.  It
   is just as dull as the first paragraph.
\end{quotation}
Another frequently-displayed structure is a list.
The following is an example of an \emph{itemized}
list.
\begin{itemize}
   \item This is the first item of an itemized list.
         Each item in the list is marked with a ``tick''.
         You don't have to worry about what kind of tick
         mark is used.

   \item This is the second item of the list.  It
         contains another list nested inside it.  The inner
         list is an \emph{enumerated} list.
         \begin{enumerate}
            \item This is the first item of an enumerated
                  list that is nested within the itemized list.

            \item This is the second item of the inner list.
                  \LaTeX\ allows you to nest lists deeper than
                  you really should.
         \end{enumerate}
         This is the rest of the second item of the outer
         list.  It is no more interesting than any other
         part of the item.
   \item This is the third item of the list.
\end{itemize}
You can even display poetry.
\begin{verse}
   There is an environment
    for verse \\             % The \\ command separates lines
   Whose features some poets % within a stanza.
   will curse.

                             % One or more blank lines separate stanzas.

   For instead of making\\
   Them do \emph{all} line breaking, \\
   It allows them to put too many words on a line when they'd rather be
   forced to be terse.
\end{verse}

Mathematical formulas may also be displayed.  A
displayed formula
is
one-line long; multiline
formulas require special formatting instructions.
   \[  \ip{\Gamma}{\psi'} = x'' + y^{2} + z_{i}^{n}\]
Don't start a paragraph with a displayed equation,
nor make one a paragraph by itself.

\end{document}               % End of document.

\section{Bench Huge Section 467}
% synthetic bench section 467
% This is a sample LaTeX input file.  (Version of 12 August 2004.)
%
% A '%' character causes TeX to ignore all remaining text on the line,
% and is used for comments like this one.

\documentclass{article}      % Specifies the document class

                             % The preamble begins here.
\title{An Example Document}  % Declares the document's title.
\author{Leslie Lamport}      % Declares the author's name.
\date{January 21, 1994}      % Deleting this command produces today's date.

\newcommand{\ip}[2]{(#1, #2)}
                             % Defines \ip{arg1}{arg2} to mean
                             % (arg1, arg2).

%\newcommand{\ip}[2]{\langle #1 | #2\rangle}
                             % This is an alternative definition of
                             % \ip that is commented out.

\begin{document}             % End of preamble and beginning of text.

\maketitle                   % Produces the title.

This is an example input file.  Comparing it with
the output it generates can show you how to
produce a simple document of your own.

\section{Ordinary Text}      % Produces section heading.  Lower-level
                             % sections are begun with similar
                             % \subsection and \subsubsection commands.

The ends  of words and sentences are marked
  by   spaces. It  doesn't matter how many
spaces    you type; one is as good as 100.  The
end of   a line counts as a space.

One   or more   blank lines denote the  end
of  a paragraph.

Since any number of consecutive spaces are treated
like a single one, the formatting of the input
file makes no difference to
      \LaTeX,                % The \LaTeX command generates the LaTeX logo.
but it makes a difference to you.  When you use
\LaTeX, making your input file as easy to read
as possible will be a great help as you write
your document and when you change it.  This sample
file shows how you can add comments to your own input
file.

Because printing is different from typewriting,
there are a number of things that you have to do
differently when preparing an input file than if
you were just typing the document directly.
Quotation marks like
       ``this''
have to be handled specially, as do quotes within
quotes:
       ``\,`this'            % \, separates the double and single quote.
        is what I just
        wrote, not  `that'\,''.

Dashes come in three sizes: an
       intra-word
dash, a medium dash for number ranges like
       1--2,
and a punctuation
       dash---like
this.

A sentence-ending space should be larger than the
space between words within a sentence.  You
sometimes have to type special commands in
conjunction with punctuation characters to get
this right, as in the following sentence.
       Gnats, gnus, etc.\ all  % `\ ' makes an inter-word space.
       begin with G\@.         % \@ marks end-of-sentence punctuation.
You should check the spaces after periods when
reading your output to make sure you haven't
forgotten any special cases.  Generating an
ellipsis
       \ldots\               % `\ ' is needed after `\ldots' because TeX
                             % ignores spaces after command names like \ldots
                             % made from \ + letters.
                             %
                             % Note how a `%' character causes TeX to ignore
                             % the end of the input line, so these blank lines
                             % do not start a new paragraph.
                             %
with the right spacing around the periods requires
a special command.

\LaTeX\ interprets some common characters as
commands, so you must type special commands to
generate them.  These characters include the
following:
       \$ \& \% \# \{ and \}.

In printing, text is usually emphasized with an
       \emph{italic}
type style.

\begin{em}
   A long segment of text can also be emphasized
   in this way.  Text within such a segment can be
   given \emph{additional} emphasis.
\end{em}

It is sometimes necessary to prevent \LaTeX\ from
breaking a line where it might otherwise do so.
This may be at a space, as between the ``Mr.''\ and
``Jones'' in
       ``Mr.~Jones'',        % ~ produces an unbreakable interword space.
or within a word---especially when the word is a
symbol like
       \mbox{\emph{itemnum}}
that makes little sense when hyphenated across
lines.

Footnotes\footnote{This is an example of a footnote.}
pose no problem.

\LaTeX\ is good at typesetting mathematical formulas
like
       \( x-3y + z = 7 \)
or
       \( a_{1} > x^{2n} + y^{2n} > x' \)
or
       \( \ip{A}{B} = \sum_{i} a_{i} b_{i} \).
The spaces you type in a formula are
ignored.  Remember that a letter like
       $x$                   % $ ... $  and  \( ... \)  are equivalent
is a formula when it denotes a mathematical
symbol, and it should be typed as one.

\section{Displayed Text}

Text is displayed by indenting it from the left
margin.  Quotations are commonly displayed.  There
are short quotations
\begin{quote}
   This is a short quotation.  It consists of a
   single paragraph of text.  See how it is formatted.
\end{quote}
and longer ones.
\begin{quotation}
   This is a longer quotation.  It consists of two
   paragraphs of text, neither of which are
   particularly interesting.

   This is the second paragraph of the quotation.  It
   is just as dull as the first paragraph.
\end{quotation}
Another frequently-displayed structure is a list.
The following is an example of an \emph{itemized}
list.
\begin{itemize}
   \item This is the first item of an itemized list.
         Each item in the list is marked with a ``tick''.
         You don't have to worry about what kind of tick
         mark is used.

   \item This is the second item of the list.  It
         contains another list nested inside it.  The inner
         list is an \emph{enumerated} list.
         \begin{enumerate}
            \item This is the first item of an enumerated
                  list that is nested within the itemized list.

            \item This is the second item of the inner list.
                  \LaTeX\ allows you to nest lists deeper than
                  you really should.
         \end{enumerate}
         This is the rest of the second item of the outer
         list.  It is no more interesting than any other
         part of the item.
   \item This is the third item of the list.
\end{itemize}
You can even display poetry.
\begin{verse}
   There is an environment
    for verse \\             % The \\ command separates lines
   Whose features some poets % within a stanza.
   will curse.

                             % One or more blank lines separate stanzas.

   For instead of making\\
   Them do \emph{all} line breaking, \\
   It allows them to put too many words on a line when they'd rather be
   forced to be terse.
\end{verse}

Mathematical formulas may also be displayed.  A
displayed formula
is
one-line long; multiline
formulas require special formatting instructions.
   \[  \ip{\Gamma}{\psi'} = x'' + y^{2} + z_{i}^{n}\]
Don't start a paragraph with a displayed equation,
nor make one a paragraph by itself.

\end{document}               % End of document.

\section{Bench Huge Section 468}
% synthetic bench section 468
% This is a sample LaTeX input file.  (Version of 12 August 2004.)
%
% A '%' character causes TeX to ignore all remaining text on the line,
% and is used for comments like this one.

\documentclass{article}      % Specifies the document class

                             % The preamble begins here.
\title{An Example Document}  % Declares the document's title.
\author{Leslie Lamport}      % Declares the author's name.
\date{January 21, 1994}      % Deleting this command produces today's date.

\newcommand{\ip}[2]{(#1, #2)}
                             % Defines \ip{arg1}{arg2} to mean
                             % (arg1, arg2).

%\newcommand{\ip}[2]{\langle #1 | #2\rangle}
                             % This is an alternative definition of
                             % \ip that is commented out.

\begin{document}             % End of preamble and beginning of text.

\maketitle                   % Produces the title.

This is an example input file.  Comparing it with
the output it generates can show you how to
produce a simple document of your own.

\section{Ordinary Text}      % Produces section heading.  Lower-level
                             % sections are begun with similar
                             % \subsection and \subsubsection commands.

The ends  of words and sentences are marked
  by   spaces. It  doesn't matter how many
spaces    you type; one is as good as 100.  The
end of   a line counts as a space.

One   or more   blank lines denote the  end
of  a paragraph.

Since any number of consecutive spaces are treated
like a single one, the formatting of the input
file makes no difference to
      \LaTeX,                % The \LaTeX command generates the LaTeX logo.
but it makes a difference to you.  When you use
\LaTeX, making your input file as easy to read
as possible will be a great help as you write
your document and when you change it.  This sample
file shows how you can add comments to your own input
file.

Because printing is different from typewriting,
there are a number of things that you have to do
differently when preparing an input file than if
you were just typing the document directly.
Quotation marks like
       ``this''
have to be handled specially, as do quotes within
quotes:
       ``\,`this'            % \, separates the double and single quote.
        is what I just
        wrote, not  `that'\,''.

Dashes come in three sizes: an
       intra-word
dash, a medium dash for number ranges like
       1--2,
and a punctuation
       dash---like
this.

A sentence-ending space should be larger than the
space between words within a sentence.  You
sometimes have to type special commands in
conjunction with punctuation characters to get
this right, as in the following sentence.
       Gnats, gnus, etc.\ all  % `\ ' makes an inter-word space.
       begin with G\@.         % \@ marks end-of-sentence punctuation.
You should check the spaces after periods when
reading your output to make sure you haven't
forgotten any special cases.  Generating an
ellipsis
       \ldots\               % `\ ' is needed after `\ldots' because TeX
                             % ignores spaces after command names like \ldots
                             % made from \ + letters.
                             %
                             % Note how a `%' character causes TeX to ignore
                             % the end of the input line, so these blank lines
                             % do not start a new paragraph.
                             %
with the right spacing around the periods requires
a special command.

\LaTeX\ interprets some common characters as
commands, so you must type special commands to
generate them.  These characters include the
following:
       \$ \& \% \# \{ and \}.

In printing, text is usually emphasized with an
       \emph{italic}
type style.

\begin{em}
   A long segment of text can also be emphasized
   in this way.  Text within such a segment can be
   given \emph{additional} emphasis.
\end{em}

It is sometimes necessary to prevent \LaTeX\ from
breaking a line where it might otherwise do so.
This may be at a space, as between the ``Mr.''\ and
``Jones'' in
       ``Mr.~Jones'',        % ~ produces an unbreakable interword space.
or within a word---especially when the word is a
symbol like
       \mbox{\emph{itemnum}}
that makes little sense when hyphenated across
lines.

Footnotes\footnote{This is an example of a footnote.}
pose no problem.

\LaTeX\ is good at typesetting mathematical formulas
like
       \( x-3y + z = 7 \)
or
       \( a_{1} > x^{2n} + y^{2n} > x' \)
or
       \( \ip{A}{B} = \sum_{i} a_{i} b_{i} \).
The spaces you type in a formula are
ignored.  Remember that a letter like
       $x$                   % $ ... $  and  \( ... \)  are equivalent
is a formula when it denotes a mathematical
symbol, and it should be typed as one.

\section{Displayed Text}

Text is displayed by indenting it from the left
margin.  Quotations are commonly displayed.  There
are short quotations
\begin{quote}
   This is a short quotation.  It consists of a
   single paragraph of text.  See how it is formatted.
\end{quote}
and longer ones.
\begin{quotation}
   This is a longer quotation.  It consists of two
   paragraphs of text, neither of which are
   particularly interesting.

   This is the second paragraph of the quotation.  It
   is just as dull as the first paragraph.
\end{quotation}
Another frequently-displayed structure is a list.
The following is an example of an \emph{itemized}
list.
\begin{itemize}
   \item This is the first item of an itemized list.
         Each item in the list is marked with a ``tick''.
         You don't have to worry about what kind of tick
         mark is used.

   \item This is the second item of the list.  It
         contains another list nested inside it.  The inner
         list is an \emph{enumerated} list.
         \begin{enumerate}
            \item This is the first item of an enumerated
                  list that is nested within the itemized list.

            \item This is the second item of the inner list.
                  \LaTeX\ allows you to nest lists deeper than
                  you really should.
         \end{enumerate}
         This is the rest of the second item of the outer
         list.  It is no more interesting than any other
         part of the item.
   \item This is the third item of the list.
\end{itemize}
You can even display poetry.
\begin{verse}
   There is an environment
    for verse \\             % The \\ command separates lines
   Whose features some poets % within a stanza.
   will curse.

                             % One or more blank lines separate stanzas.

   For instead of making\\
   Them do \emph{all} line breaking, \\
   It allows them to put too many words on a line when they'd rather be
   forced to be terse.
\end{verse}

Mathematical formulas may also be displayed.  A
displayed formula
is
one-line long; multiline
formulas require special formatting instructions.
   \[  \ip{\Gamma}{\psi'} = x'' + y^{2} + z_{i}^{n}\]
Don't start a paragraph with a displayed equation,
nor make one a paragraph by itself.

\end{document}               % End of document.

\section{Bench Huge Section 469}
% synthetic bench section 469
% This is a sample LaTeX input file.  (Version of 12 August 2004.)
%
% A '%' character causes TeX to ignore all remaining text on the line,
% and is used for comments like this one.

\documentclass{article}      % Specifies the document class

                             % The preamble begins here.
\title{An Example Document}  % Declares the document's title.
\author{Leslie Lamport}      % Declares the author's name.
\date{January 21, 1994}      % Deleting this command produces today's date.

\newcommand{\ip}[2]{(#1, #2)}
                             % Defines \ip{arg1}{arg2} to mean
                             % (arg1, arg2).

%\newcommand{\ip}[2]{\langle #1 | #2\rangle}
                             % This is an alternative definition of
                             % \ip that is commented out.

\begin{document}             % End of preamble and beginning of text.

\maketitle                   % Produces the title.

This is an example input file.  Comparing it with
the output it generates can show you how to
produce a simple document of your own.

\section{Ordinary Text}      % Produces section heading.  Lower-level
                             % sections are begun with similar
                             % \subsection and \subsubsection commands.

The ends  of words and sentences are marked
  by   spaces. It  doesn't matter how many
spaces    you type; one is as good as 100.  The
end of   a line counts as a space.

One   or more   blank lines denote the  end
of  a paragraph.

Since any number of consecutive spaces are treated
like a single one, the formatting of the input
file makes no difference to
      \LaTeX,                % The \LaTeX command generates the LaTeX logo.
but it makes a difference to you.  When you use
\LaTeX, making your input file as easy to read
as possible will be a great help as you write
your document and when you change it.  This sample
file shows how you can add comments to your own input
file.

Because printing is different from typewriting,
there are a number of things that you have to do
differently when preparing an input file than if
you were just typing the document directly.
Quotation marks like
       ``this''
have to be handled specially, as do quotes within
quotes:
       ``\,`this'            % \, separates the double and single quote.
        is what I just
        wrote, not  `that'\,''.

Dashes come in three sizes: an
       intra-word
dash, a medium dash for number ranges like
       1--2,
and a punctuation
       dash---like
this.

A sentence-ending space should be larger than the
space between words within a sentence.  You
sometimes have to type special commands in
conjunction with punctuation characters to get
this right, as in the following sentence.
       Gnats, gnus, etc.\ all  % `\ ' makes an inter-word space.
       begin with G\@.         % \@ marks end-of-sentence punctuation.
You should check the spaces after periods when
reading your output to make sure you haven't
forgotten any special cases.  Generating an
ellipsis
       \ldots\               % `\ ' is needed after `\ldots' because TeX
                             % ignores spaces after command names like \ldots
                             % made from \ + letters.
                             %
                             % Note how a `%' character causes TeX to ignore
                             % the end of the input line, so these blank lines
                             % do not start a new paragraph.
                             %
with the right spacing around the periods requires
a special command.

\LaTeX\ interprets some common characters as
commands, so you must type special commands to
generate them.  These characters include the
following:
       \$ \& \% \# \{ and \}.

In printing, text is usually emphasized with an
       \emph{italic}
type style.

\begin{em}
   A long segment of text can also be emphasized
   in this way.  Text within such a segment can be
   given \emph{additional} emphasis.
\end{em}

It is sometimes necessary to prevent \LaTeX\ from
breaking a line where it might otherwise do so.
This may be at a space, as between the ``Mr.''\ and
``Jones'' in
       ``Mr.~Jones'',        % ~ produces an unbreakable interword space.
or within a word---especially when the word is a
symbol like
       \mbox{\emph{itemnum}}
that makes little sense when hyphenated across
lines.

Footnotes\footnote{This is an example of a footnote.}
pose no problem.

\LaTeX\ is good at typesetting mathematical formulas
like
       \( x-3y + z = 7 \)
or
       \( a_{1} > x^{2n} + y^{2n} > x' \)
or
       \( \ip{A}{B} = \sum_{i} a_{i} b_{i} \).
The spaces you type in a formula are
ignored.  Remember that a letter like
       $x$                   % $ ... $  and  \( ... \)  are equivalent
is a formula when it denotes a mathematical
symbol, and it should be typed as one.

\section{Displayed Text}

Text is displayed by indenting it from the left
margin.  Quotations are commonly displayed.  There
are short quotations
\begin{quote}
   This is a short quotation.  It consists of a
   single paragraph of text.  See how it is formatted.
\end{quote}
and longer ones.
\begin{quotation}
   This is a longer quotation.  It consists of two
   paragraphs of text, neither of which are
   particularly interesting.

   This is the second paragraph of the quotation.  It
   is just as dull as the first paragraph.
\end{quotation}
Another frequently-displayed structure is a list.
The following is an example of an \emph{itemized}
list.
\begin{itemize}
   \item This is the first item of an itemized list.
         Each item in the list is marked with a ``tick''.
         You don't have to worry about what kind of tick
         mark is used.

   \item This is the second item of the list.  It
         contains another list nested inside it.  The inner
         list is an \emph{enumerated} list.
         \begin{enumerate}
            \item This is the first item of an enumerated
                  list that is nested within the itemized list.

            \item This is the second item of the inner list.
                  \LaTeX\ allows you to nest lists deeper than
                  you really should.
         \end{enumerate}
         This is the rest of the second item of the outer
         list.  It is no more interesting than any other
         part of the item.
   \item This is the third item of the list.
\end{itemize}
You can even display poetry.
\begin{verse}
   There is an environment
    for verse \\             % The \\ command separates lines
   Whose features some poets % within a stanza.
   will curse.

                             % One or more blank lines separate stanzas.

   For instead of making\\
   Them do \emph{all} line breaking, \\
   It allows them to put too many words on a line when they'd rather be
   forced to be terse.
\end{verse}

Mathematical formulas may also be displayed.  A
displayed formula
is
one-line long; multiline
formulas require special formatting instructions.
   \[  \ip{\Gamma}{\psi'} = x'' + y^{2} + z_{i}^{n}\]
Don't start a paragraph with a displayed equation,
nor make one a paragraph by itself.

\end{document}               % End of document.

\section{Bench Huge Section 470}
% synthetic bench section 470
% This is a sample LaTeX input file.  (Version of 12 August 2004.)
%
% A '%' character causes TeX to ignore all remaining text on the line,
% and is used for comments like this one.

\documentclass{article}      % Specifies the document class

                             % The preamble begins here.
\title{An Example Document}  % Declares the document's title.
\author{Leslie Lamport}      % Declares the author's name.
\date{January 21, 1994}      % Deleting this command produces today's date.

\newcommand{\ip}[2]{(#1, #2)}
                             % Defines \ip{arg1}{arg2} to mean
                             % (arg1, arg2).

%\newcommand{\ip}[2]{\langle #1 | #2\rangle}
                             % This is an alternative definition of
                             % \ip that is commented out.

\begin{document}             % End of preamble and beginning of text.

\maketitle                   % Produces the title.

This is an example input file.  Comparing it with
the output it generates can show you how to
produce a simple document of your own.

\section{Ordinary Text}      % Produces section heading.  Lower-level
                             % sections are begun with similar
                             % \subsection and \subsubsection commands.

The ends  of words and sentences are marked
  by   spaces. It  doesn't matter how many
spaces    you type; one is as good as 100.  The
end of   a line counts as a space.

One   or more   blank lines denote the  end
of  a paragraph.

Since any number of consecutive spaces are treated
like a single one, the formatting of the input
file makes no difference to
      \LaTeX,                % The \LaTeX command generates the LaTeX logo.
but it makes a difference to you.  When you use
\LaTeX, making your input file as easy to read
as possible will be a great help as you write
your document and when you change it.  This sample
file shows how you can add comments to your own input
file.

Because printing is different from typewriting,
there are a number of things that you have to do
differently when preparing an input file than if
you were just typing the document directly.
Quotation marks like
       ``this''
have to be handled specially, as do quotes within
quotes:
       ``\,`this'            % \, separates the double and single quote.
        is what I just
        wrote, not  `that'\,''.

Dashes come in three sizes: an
       intra-word
dash, a medium dash for number ranges like
       1--2,
and a punctuation
       dash---like
this.

A sentence-ending space should be larger than the
space between words within a sentence.  You
sometimes have to type special commands in
conjunction with punctuation characters to get
this right, as in the following sentence.
       Gnats, gnus, etc.\ all  % `\ ' makes an inter-word space.
       begin with G\@.         % \@ marks end-of-sentence punctuation.
You should check the spaces after periods when
reading your output to make sure you haven't
forgotten any special cases.  Generating an
ellipsis
       \ldots\               % `\ ' is needed after `\ldots' because TeX
                             % ignores spaces after command names like \ldots
                             % made from \ + letters.
                             %
                             % Note how a `%' character causes TeX to ignore
                             % the end of the input line, so these blank lines
                             % do not start a new paragraph.
                             %
with the right spacing around the periods requires
a special command.

\LaTeX\ interprets some common characters as
commands, so you must type special commands to
generate them.  These characters include the
following:
       \$ \& \% \# \{ and \}.

In printing, text is usually emphasized with an
       \emph{italic}
type style.

\begin{em}
   A long segment of text can also be emphasized
   in this way.  Text within such a segment can be
   given \emph{additional} emphasis.
\end{em}

It is sometimes necessary to prevent \LaTeX\ from
breaking a line where it might otherwise do so.
This may be at a space, as between the ``Mr.''\ and
``Jones'' in
       ``Mr.~Jones'',        % ~ produces an unbreakable interword space.
or within a word---especially when the word is a
symbol like
       \mbox{\emph{itemnum}}
that makes little sense when hyphenated across
lines.

Footnotes\footnote{This is an example of a footnote.}
pose no problem.

\LaTeX\ is good at typesetting mathematical formulas
like
       \( x-3y + z = 7 \)
or
       \( a_{1} > x^{2n} + y^{2n} > x' \)
or
       \( \ip{A}{B} = \sum_{i} a_{i} b_{i} \).
The spaces you type in a formula are
ignored.  Remember that a letter like
       $x$                   % $ ... $  and  \( ... \)  are equivalent
is a formula when it denotes a mathematical
symbol, and it should be typed as one.

\section{Displayed Text}

Text is displayed by indenting it from the left
margin.  Quotations are commonly displayed.  There
are short quotations
\begin{quote}
   This is a short quotation.  It consists of a
   single paragraph of text.  See how it is formatted.
\end{quote}
and longer ones.
\begin{quotation}
   This is a longer quotation.  It consists of two
   paragraphs of text, neither of which are
   particularly interesting.

   This is the second paragraph of the quotation.  It
   is just as dull as the first paragraph.
\end{quotation}
Another frequently-displayed structure is a list.
The following is an example of an \emph{itemized}
list.
\begin{itemize}
   \item This is the first item of an itemized list.
         Each item in the list is marked with a ``tick''.
         You don't have to worry about what kind of tick
         mark is used.

   \item This is the second item of the list.  It
         contains another list nested inside it.  The inner
         list is an \emph{enumerated} list.
         \begin{enumerate}
            \item This is the first item of an enumerated
                  list that is nested within the itemized list.

            \item This is the second item of the inner list.
                  \LaTeX\ allows you to nest lists deeper than
                  you really should.
         \end{enumerate}
         This is the rest of the second item of the outer
         list.  It is no more interesting than any other
         part of the item.
   \item This is the third item of the list.
\end{itemize}
You can even display poetry.
\begin{verse}
   There is an environment
    for verse \\             % The \\ command separates lines
   Whose features some poets % within a stanza.
   will curse.

                             % One or more blank lines separate stanzas.

   For instead of making\\
   Them do \emph{all} line breaking, \\
   It allows them to put too many words on a line when they'd rather be
   forced to be terse.
\end{verse}

Mathematical formulas may also be displayed.  A
displayed formula
is
one-line long; multiline
formulas require special formatting instructions.
   \[  \ip{\Gamma}{\psi'} = x'' + y^{2} + z_{i}^{n}\]
Don't start a paragraph with a displayed equation,
nor make one a paragraph by itself.

\end{document}               % End of document.

\section{Bench Huge Section 471}
% synthetic bench section 471
% This is a sample LaTeX input file.  (Version of 12 August 2004.)
%
% A '%' character causes TeX to ignore all remaining text on the line,
% and is used for comments like this one.

\documentclass{article}      % Specifies the document class

                             % The preamble begins here.
\title{An Example Document}  % Declares the document's title.
\author{Leslie Lamport}      % Declares the author's name.
\date{January 21, 1994}      % Deleting this command produces today's date.

\newcommand{\ip}[2]{(#1, #2)}
                             % Defines \ip{arg1}{arg2} to mean
                             % (arg1, arg2).

%\newcommand{\ip}[2]{\langle #1 | #2\rangle}
                             % This is an alternative definition of
                             % \ip that is commented out.

\begin{document}             % End of preamble and beginning of text.

\maketitle                   % Produces the title.

This is an example input file.  Comparing it with
the output it generates can show you how to
produce a simple document of your own.

\section{Ordinary Text}      % Produces section heading.  Lower-level
                             % sections are begun with similar
                             % \subsection and \subsubsection commands.

The ends  of words and sentences are marked
  by   spaces. It  doesn't matter how many
spaces    you type; one is as good as 100.  The
end of   a line counts as a space.

One   or more   blank lines denote the  end
of  a paragraph.

Since any number of consecutive spaces are treated
like a single one, the formatting of the input
file makes no difference to
      \LaTeX,                % The \LaTeX command generates the LaTeX logo.
but it makes a difference to you.  When you use
\LaTeX, making your input file as easy to read
as possible will be a great help as you write
your document and when you change it.  This sample
file shows how you can add comments to your own input
file.

Because printing is different from typewriting,
there are a number of things that you have to do
differently when preparing an input file than if
you were just typing the document directly.
Quotation marks like
       ``this''
have to be handled specially, as do quotes within
quotes:
       ``\,`this'            % \, separates the double and single quote.
        is what I just
        wrote, not  `that'\,''.

Dashes come in three sizes: an
       intra-word
dash, a medium dash for number ranges like
       1--2,
and a punctuation
       dash---like
this.

A sentence-ending space should be larger than the
space between words within a sentence.  You
sometimes have to type special commands in
conjunction with punctuation characters to get
this right, as in the following sentence.
       Gnats, gnus, etc.\ all  % `\ ' makes an inter-word space.
       begin with G\@.         % \@ marks end-of-sentence punctuation.
You should check the spaces after periods when
reading your output to make sure you haven't
forgotten any special cases.  Generating an
ellipsis
       \ldots\               % `\ ' is needed after `\ldots' because TeX
                             % ignores spaces after command names like \ldots
                             % made from \ + letters.
                             %
                             % Note how a `%' character causes TeX to ignore
                             % the end of the input line, so these blank lines
                             % do not start a new paragraph.
                             %
with the right spacing around the periods requires
a special command.

\LaTeX\ interprets some common characters as
commands, so you must type special commands to
generate them.  These characters include the
following:
       \$ \& \% \# \{ and \}.

In printing, text is usually emphasized with an
       \emph{italic}
type style.

\begin{em}
   A long segment of text can also be emphasized
   in this way.  Text within such a segment can be
   given \emph{additional} emphasis.
\end{em}

It is sometimes necessary to prevent \LaTeX\ from
breaking a line where it might otherwise do so.
This may be at a space, as between the ``Mr.''\ and
``Jones'' in
       ``Mr.~Jones'',        % ~ produces an unbreakable interword space.
or within a word---especially when the word is a
symbol like
       \mbox{\emph{itemnum}}
that makes little sense when hyphenated across
lines.

Footnotes\footnote{This is an example of a footnote.}
pose no problem.

\LaTeX\ is good at typesetting mathematical formulas
like
       \( x-3y + z = 7 \)
or
       \( a_{1} > x^{2n} + y^{2n} > x' \)
or
       \( \ip{A}{B} = \sum_{i} a_{i} b_{i} \).
The spaces you type in a formula are
ignored.  Remember that a letter like
       $x$                   % $ ... $  and  \( ... \)  are equivalent
is a formula when it denotes a mathematical
symbol, and it should be typed as one.

\section{Displayed Text}

Text is displayed by indenting it from the left
margin.  Quotations are commonly displayed.  There
are short quotations
\begin{quote}
   This is a short quotation.  It consists of a
   single paragraph of text.  See how it is formatted.
\end{quote}
and longer ones.
\begin{quotation}
   This is a longer quotation.  It consists of two
   paragraphs of text, neither of which are
   particularly interesting.

   This is the second paragraph of the quotation.  It
   is just as dull as the first paragraph.
\end{quotation}
Another frequently-displayed structure is a list.
The following is an example of an \emph{itemized}
list.
\begin{itemize}
   \item This is the first item of an itemized list.
         Each item in the list is marked with a ``tick''.
         You don't have to worry about what kind of tick
         mark is used.

   \item This is the second item of the list.  It
         contains another list nested inside it.  The inner
         list is an \emph{enumerated} list.
         \begin{enumerate}
            \item This is the first item of an enumerated
                  list that is nested within the itemized list.

            \item This is the second item of the inner list.
                  \LaTeX\ allows you to nest lists deeper than
                  you really should.
         \end{enumerate}
         This is the rest of the second item of the outer
         list.  It is no more interesting than any other
         part of the item.
   \item This is the third item of the list.
\end{itemize}
You can even display poetry.
\begin{verse}
   There is an environment
    for verse \\             % The \\ command separates lines
   Whose features some poets % within a stanza.
   will curse.

                             % One or more blank lines separate stanzas.

   For instead of making\\
   Them do \emph{all} line breaking, \\
   It allows them to put too many words on a line when they'd rather be
   forced to be terse.
\end{verse}

Mathematical formulas may also be displayed.  A
displayed formula
is
one-line long; multiline
formulas require special formatting instructions.
   \[  \ip{\Gamma}{\psi'} = x'' + y^{2} + z_{i}^{n}\]
Don't start a paragraph with a displayed equation,
nor make one a paragraph by itself.

\end{document}               % End of document.

\section{Bench Huge Section 472}
% synthetic bench section 472
% This is a sample LaTeX input file.  (Version of 12 August 2004.)
%
% A '%' character causes TeX to ignore all remaining text on the line,
% and is used for comments like this one.

\documentclass{article}      % Specifies the document class

                             % The preamble begins here.
\title{An Example Document}  % Declares the document's title.
\author{Leslie Lamport}      % Declares the author's name.
\date{January 21, 1994}      % Deleting this command produces today's date.

\newcommand{\ip}[2]{(#1, #2)}
                             % Defines \ip{arg1}{arg2} to mean
                             % (arg1, arg2).

%\newcommand{\ip}[2]{\langle #1 | #2\rangle}
                             % This is an alternative definition of
                             % \ip that is commented out.

\begin{document}             % End of preamble and beginning of text.

\maketitle                   % Produces the title.

This is an example input file.  Comparing it with
the output it generates can show you how to
produce a simple document of your own.

\section{Ordinary Text}      % Produces section heading.  Lower-level
                             % sections are begun with similar
                             % \subsection and \subsubsection commands.

The ends  of words and sentences are marked
  by   spaces. It  doesn't matter how many
spaces    you type; one is as good as 100.  The
end of   a line counts as a space.

One   or more   blank lines denote the  end
of  a paragraph.

Since any number of consecutive spaces are treated
like a single one, the formatting of the input
file makes no difference to
      \LaTeX,                % The \LaTeX command generates the LaTeX logo.
but it makes a difference to you.  When you use
\LaTeX, making your input file as easy to read
as possible will be a great help as you write
your document and when you change it.  This sample
file shows how you can add comments to your own input
file.

Because printing is different from typewriting,
there are a number of things that you have to do
differently when preparing an input file than if
you were just typing the document directly.
Quotation marks like
       ``this''
have to be handled specially, as do quotes within
quotes:
       ``\,`this'            % \, separates the double and single quote.
        is what I just
        wrote, not  `that'\,''.

Dashes come in three sizes: an
       intra-word
dash, a medium dash for number ranges like
       1--2,
and a punctuation
       dash---like
this.

A sentence-ending space should be larger than the
space between words within a sentence.  You
sometimes have to type special commands in
conjunction with punctuation characters to get
this right, as in the following sentence.
       Gnats, gnus, etc.\ all  % `\ ' makes an inter-word space.
       begin with G\@.         % \@ marks end-of-sentence punctuation.
You should check the spaces after periods when
reading your output to make sure you haven't
forgotten any special cases.  Generating an
ellipsis
       \ldots\               % `\ ' is needed after `\ldots' because TeX
                             % ignores spaces after command names like \ldots
                             % made from \ + letters.
                             %
                             % Note how a `%' character causes TeX to ignore
                             % the end of the input line, so these blank lines
                             % do not start a new paragraph.
                             %
with the right spacing around the periods requires
a special command.

\LaTeX\ interprets some common characters as
commands, so you must type special commands to
generate them.  These characters include the
following:
       \$ \& \% \# \{ and \}.

In printing, text is usually emphasized with an
       \emph{italic}
type style.

\begin{em}
   A long segment of text can also be emphasized
   in this way.  Text within such a segment can be
   given \emph{additional} emphasis.
\end{em}

It is sometimes necessary to prevent \LaTeX\ from
breaking a line where it might otherwise do so.
This may be at a space, as between the ``Mr.''\ and
``Jones'' in
       ``Mr.~Jones'',        % ~ produces an unbreakable interword space.
or within a word---especially when the word is a
symbol like
       \mbox{\emph{itemnum}}
that makes little sense when hyphenated across
lines.

Footnotes\footnote{This is an example of a footnote.}
pose no problem.

\LaTeX\ is good at typesetting mathematical formulas
like
       \( x-3y + z = 7 \)
or
       \( a_{1} > x^{2n} + y^{2n} > x' \)
or
       \( \ip{A}{B} = \sum_{i} a_{i} b_{i} \).
The spaces you type in a formula are
ignored.  Remember that a letter like
       $x$                   % $ ... $  and  \( ... \)  are equivalent
is a formula when it denotes a mathematical
symbol, and it should be typed as one.

\section{Displayed Text}

Text is displayed by indenting it from the left
margin.  Quotations are commonly displayed.  There
are short quotations
\begin{quote}
   This is a short quotation.  It consists of a
   single paragraph of text.  See how it is formatted.
\end{quote}
and longer ones.
\begin{quotation}
   This is a longer quotation.  It consists of two
   paragraphs of text, neither of which are
   particularly interesting.

   This is the second paragraph of the quotation.  It
   is just as dull as the first paragraph.
\end{quotation}
Another frequently-displayed structure is a list.
The following is an example of an \emph{itemized}
list.
\begin{itemize}
   \item This is the first item of an itemized list.
         Each item in the list is marked with a ``tick''.
         You don't have to worry about what kind of tick
         mark is used.

   \item This is the second item of the list.  It
         contains another list nested inside it.  The inner
         list is an \emph{enumerated} list.
         \begin{enumerate}
            \item This is the first item of an enumerated
                  list that is nested within the itemized list.

            \item This is the second item of the inner list.
                  \LaTeX\ allows you to nest lists deeper than
                  you really should.
         \end{enumerate}
         This is the rest of the second item of the outer
         list.  It is no more interesting than any other
         part of the item.
   \item This is the third item of the list.
\end{itemize}
You can even display poetry.
\begin{verse}
   There is an environment
    for verse \\             % The \\ command separates lines
   Whose features some poets % within a stanza.
   will curse.

                             % One or more blank lines separate stanzas.

   For instead of making\\
   Them do \emph{all} line breaking, \\
   It allows them to put too many words on a line when they'd rather be
   forced to be terse.
\end{verse}

Mathematical formulas may also be displayed.  A
displayed formula
is
one-line long; multiline
formulas require special formatting instructions.
   \[  \ip{\Gamma}{\psi'} = x'' + y^{2} + z_{i}^{n}\]
Don't start a paragraph with a displayed equation,
nor make one a paragraph by itself.

\end{document}               % End of document.

\section{Bench Huge Section 473}
% synthetic bench section 473
% This is a sample LaTeX input file.  (Version of 12 August 2004.)
%
% A '%' character causes TeX to ignore all remaining text on the line,
% and is used for comments like this one.

\documentclass{article}      % Specifies the document class

                             % The preamble begins here.
\title{An Example Document}  % Declares the document's title.
\author{Leslie Lamport}      % Declares the author's name.
\date{January 21, 1994}      % Deleting this command produces today's date.

\newcommand{\ip}[2]{(#1, #2)}
                             % Defines \ip{arg1}{arg2} to mean
                             % (arg1, arg2).

%\newcommand{\ip}[2]{\langle #1 | #2\rangle}
                             % This is an alternative definition of
                             % \ip that is commented out.

\begin{document}             % End of preamble and beginning of text.

\maketitle                   % Produces the title.

This is an example input file.  Comparing it with
the output it generates can show you how to
produce a simple document of your own.

\section{Ordinary Text}      % Produces section heading.  Lower-level
                             % sections are begun with similar
                             % \subsection and \subsubsection commands.

The ends  of words and sentences are marked
  by   spaces. It  doesn't matter how many
spaces    you type; one is as good as 100.  The
end of   a line counts as a space.

One   or more   blank lines denote the  end
of  a paragraph.

Since any number of consecutive spaces are treated
like a single one, the formatting of the input
file makes no difference to
      \LaTeX,                % The \LaTeX command generates the LaTeX logo.
but it makes a difference to you.  When you use
\LaTeX, making your input file as easy to read
as possible will be a great help as you write
your document and when you change it.  This sample
file shows how you can add comments to your own input
file.

Because printing is different from typewriting,
there are a number of things that you have to do
differently when preparing an input file than if
you were just typing the document directly.
Quotation marks like
       ``this''
have to be handled specially, as do quotes within
quotes:
       ``\,`this'            % \, separates the double and single quote.
        is what I just
        wrote, not  `that'\,''.

Dashes come in three sizes: an
       intra-word
dash, a medium dash for number ranges like
       1--2,
and a punctuation
       dash---like
this.

A sentence-ending space should be larger than the
space between words within a sentence.  You
sometimes have to type special commands in
conjunction with punctuation characters to get
this right, as in the following sentence.
       Gnats, gnus, etc.\ all  % `\ ' makes an inter-word space.
       begin with G\@.         % \@ marks end-of-sentence punctuation.
You should check the spaces after periods when
reading your output to make sure you haven't
forgotten any special cases.  Generating an
ellipsis
       \ldots\               % `\ ' is needed after `\ldots' because TeX
                             % ignores spaces after command names like \ldots
                             % made from \ + letters.
                             %
                             % Note how a `%' character causes TeX to ignore
                             % the end of the input line, so these blank lines
                             % do not start a new paragraph.
                             %
with the right spacing around the periods requires
a special command.

\LaTeX\ interprets some common characters as
commands, so you must type special commands to
generate them.  These characters include the
following:
       \$ \& \% \# \{ and \}.

In printing, text is usually emphasized with an
       \emph{italic}
type style.

\begin{em}
   A long segment of text can also be emphasized
   in this way.  Text within such a segment can be
   given \emph{additional} emphasis.
\end{em}

It is sometimes necessary to prevent \LaTeX\ from
breaking a line where it might otherwise do so.
This may be at a space, as between the ``Mr.''\ and
``Jones'' in
       ``Mr.~Jones'',        % ~ produces an unbreakable interword space.
or within a word---especially when the word is a
symbol like
       \mbox{\emph{itemnum}}
that makes little sense when hyphenated across
lines.

Footnotes\footnote{This is an example of a footnote.}
pose no problem.

\LaTeX\ is good at typesetting mathematical formulas
like
       \( x-3y + z = 7 \)
or
       \( a_{1} > x^{2n} + y^{2n} > x' \)
or
       \( \ip{A}{B} = \sum_{i} a_{i} b_{i} \).
The spaces you type in a formula are
ignored.  Remember that a letter like
       $x$                   % $ ... $  and  \( ... \)  are equivalent
is a formula when it denotes a mathematical
symbol, and it should be typed as one.

\section{Displayed Text}

Text is displayed by indenting it from the left
margin.  Quotations are commonly displayed.  There
are short quotations
\begin{quote}
   This is a short quotation.  It consists of a
   single paragraph of text.  See how it is formatted.
\end{quote}
and longer ones.
\begin{quotation}
   This is a longer quotation.  It consists of two
   paragraphs of text, neither of which are
   particularly interesting.

   This is the second paragraph of the quotation.  It
   is just as dull as the first paragraph.
\end{quotation}
Another frequently-displayed structure is a list.
The following is an example of an \emph{itemized}
list.
\begin{itemize}
   \item This is the first item of an itemized list.
         Each item in the list is marked with a ``tick''.
         You don't have to worry about what kind of tick
         mark is used.

   \item This is the second item of the list.  It
         contains another list nested inside it.  The inner
         list is an \emph{enumerated} list.
         \begin{enumerate}
            \item This is the first item of an enumerated
                  list that is nested within the itemized list.

            \item This is the second item of the inner list.
                  \LaTeX\ allows you to nest lists deeper than
                  you really should.
         \end{enumerate}
         This is the rest of the second item of the outer
         list.  It is no more interesting than any other
         part of the item.
   \item This is the third item of the list.
\end{itemize}
You can even display poetry.
\begin{verse}
   There is an environment
    for verse \\             % The \\ command separates lines
   Whose features some poets % within a stanza.
   will curse.

                             % One or more blank lines separate stanzas.

   For instead of making\\
   Them do \emph{all} line breaking, \\
   It allows them to put too many words on a line when they'd rather be
   forced to be terse.
\end{verse}

Mathematical formulas may also be displayed.  A
displayed formula
is
one-line long; multiline
formulas require special formatting instructions.
   \[  \ip{\Gamma}{\psi'} = x'' + y^{2} + z_{i}^{n}\]
Don't start a paragraph with a displayed equation,
nor make one a paragraph by itself.

\end{document}               % End of document.

\section{Bench Huge Section 474}
% synthetic bench section 474
% This is a sample LaTeX input file.  (Version of 12 August 2004.)
%
% A '%' character causes TeX to ignore all remaining text on the line,
% and is used for comments like this one.

\documentclass{article}      % Specifies the document class

                             % The preamble begins here.
\title{An Example Document}  % Declares the document's title.
\author{Leslie Lamport}      % Declares the author's name.
\date{January 21, 1994}      % Deleting this command produces today's date.

\newcommand{\ip}[2]{(#1, #2)}
                             % Defines \ip{arg1}{arg2} to mean
                             % (arg1, arg2).

%\newcommand{\ip}[2]{\langle #1 | #2\rangle}
                             % This is an alternative definition of
                             % \ip that is commented out.

\begin{document}             % End of preamble and beginning of text.

\maketitle                   % Produces the title.

This is an example input file.  Comparing it with
the output it generates can show you how to
produce a simple document of your own.

\section{Ordinary Text}      % Produces section heading.  Lower-level
                             % sections are begun with similar
                             % \subsection and \subsubsection commands.

The ends  of words and sentences are marked
  by   spaces. It  doesn't matter how many
spaces    you type; one is as good as 100.  The
end of   a line counts as a space.

One   or more   blank lines denote the  end
of  a paragraph.

Since any number of consecutive spaces are treated
like a single one, the formatting of the input
file makes no difference to
      \LaTeX,                % The \LaTeX command generates the LaTeX logo.
but it makes a difference to you.  When you use
\LaTeX, making your input file as easy to read
as possible will be a great help as you write
your document and when you change it.  This sample
file shows how you can add comments to your own input
file.

Because printing is different from typewriting,
there are a number of things that you have to do
differently when preparing an input file than if
you were just typing the document directly.
Quotation marks like
       ``this''
have to be handled specially, as do quotes within
quotes:
       ``\,`this'            % \, separates the double and single quote.
        is what I just
        wrote, not  `that'\,''.

Dashes come in three sizes: an
       intra-word
dash, a medium dash for number ranges like
       1--2,
and a punctuation
       dash---like
this.

A sentence-ending space should be larger than the
space between words within a sentence.  You
sometimes have to type special commands in
conjunction with punctuation characters to get
this right, as in the following sentence.
       Gnats, gnus, etc.\ all  % `\ ' makes an inter-word space.
       begin with G\@.         % \@ marks end-of-sentence punctuation.
You should check the spaces after periods when
reading your output to make sure you haven't
forgotten any special cases.  Generating an
ellipsis
       \ldots\               % `\ ' is needed after `\ldots' because TeX
                             % ignores spaces after command names like \ldots
                             % made from \ + letters.
                             %
                             % Note how a `%' character causes TeX to ignore
                             % the end of the input line, so these blank lines
                             % do not start a new paragraph.
                             %
with the right spacing around the periods requires
a special command.

\LaTeX\ interprets some common characters as
commands, so you must type special commands to
generate them.  These characters include the
following:
       \$ \& \% \# \{ and \}.

In printing, text is usually emphasized with an
       \emph{italic}
type style.

\begin{em}
   A long segment of text can also be emphasized
   in this way.  Text within such a segment can be
   given \emph{additional} emphasis.
\end{em}

It is sometimes necessary to prevent \LaTeX\ from
breaking a line where it might otherwise do so.
This may be at a space, as between the ``Mr.''\ and
``Jones'' in
       ``Mr.~Jones'',        % ~ produces an unbreakable interword space.
or within a word---especially when the word is a
symbol like
       \mbox{\emph{itemnum}}
that makes little sense when hyphenated across
lines.

Footnotes\footnote{This is an example of a footnote.}
pose no problem.

\LaTeX\ is good at typesetting mathematical formulas
like
       \( x-3y + z = 7 \)
or
       \( a_{1} > x^{2n} + y^{2n} > x' \)
or
       \( \ip{A}{B} = \sum_{i} a_{i} b_{i} \).
The spaces you type in a formula are
ignored.  Remember that a letter like
       $x$                   % $ ... $  and  \( ... \)  are equivalent
is a formula when it denotes a mathematical
symbol, and it should be typed as one.

\section{Displayed Text}

Text is displayed by indenting it from the left
margin.  Quotations are commonly displayed.  There
are short quotations
\begin{quote}
   This is a short quotation.  It consists of a
   single paragraph of text.  See how it is formatted.
\end{quote}
and longer ones.
\begin{quotation}
   This is a longer quotation.  It consists of two
   paragraphs of text, neither of which are
   particularly interesting.

   This is the second paragraph of the quotation.  It
   is just as dull as the first paragraph.
\end{quotation}
Another frequently-displayed structure is a list.
The following is an example of an \emph{itemized}
list.
\begin{itemize}
   \item This is the first item of an itemized list.
         Each item in the list is marked with a ``tick''.
         You don't have to worry about what kind of tick
         mark is used.

   \item This is the second item of the list.  It
         contains another list nested inside it.  The inner
         list is an \emph{enumerated} list.
         \begin{enumerate}
            \item This is the first item of an enumerated
                  list that is nested within the itemized list.

            \item This is the second item of the inner list.
                  \LaTeX\ allows you to nest lists deeper than
                  you really should.
         \end{enumerate}
         This is the rest of the second item of the outer
         list.  It is no more interesting than any other
         part of the item.
   \item This is the third item of the list.
\end{itemize}
You can even display poetry.
\begin{verse}
   There is an environment
    for verse \\             % The \\ command separates lines
   Whose features some poets % within a stanza.
   will curse.

                             % One or more blank lines separate stanzas.

   For instead of making\\
   Them do \emph{all} line breaking, \\
   It allows them to put too many words on a line when they'd rather be
   forced to be terse.
\end{verse}

Mathematical formulas may also be displayed.  A
displayed formula
is
one-line long; multiline
formulas require special formatting instructions.
   \[  \ip{\Gamma}{\psi'} = x'' + y^{2} + z_{i}^{n}\]
Don't start a paragraph with a displayed equation,
nor make one a paragraph by itself.

\end{document}               % End of document.

\section{Bench Huge Section 475}
% synthetic bench section 475
% This is a sample LaTeX input file.  (Version of 12 August 2004.)
%
% A '%' character causes TeX to ignore all remaining text on the line,
% and is used for comments like this one.

\documentclass{article}      % Specifies the document class

                             % The preamble begins here.
\title{An Example Document}  % Declares the document's title.
\author{Leslie Lamport}      % Declares the author's name.
\date{January 21, 1994}      % Deleting this command produces today's date.

\newcommand{\ip}[2]{(#1, #2)}
                             % Defines \ip{arg1}{arg2} to mean
                             % (arg1, arg2).

%\newcommand{\ip}[2]{\langle #1 | #2\rangle}
                             % This is an alternative definition of
                             % \ip that is commented out.

\begin{document}             % End of preamble and beginning of text.

\maketitle                   % Produces the title.

This is an example input file.  Comparing it with
the output it generates can show you how to
produce a simple document of your own.

\section{Ordinary Text}      % Produces section heading.  Lower-level
                             % sections are begun with similar
                             % \subsection and \subsubsection commands.

The ends  of words and sentences are marked
  by   spaces. It  doesn't matter how many
spaces    you type; one is as good as 100.  The
end of   a line counts as a space.

One   or more   blank lines denote the  end
of  a paragraph.

Since any number of consecutive spaces are treated
like a single one, the formatting of the input
file makes no difference to
      \LaTeX,                % The \LaTeX command generates the LaTeX logo.
but it makes a difference to you.  When you use
\LaTeX, making your input file as easy to read
as possible will be a great help as you write
your document and when you change it.  This sample
file shows how you can add comments to your own input
file.

Because printing is different from typewriting,
there are a number of things that you have to do
differently when preparing an input file than if
you were just typing the document directly.
Quotation marks like
       ``this''
have to be handled specially, as do quotes within
quotes:
       ``\,`this'            % \, separates the double and single quote.
        is what I just
        wrote, not  `that'\,''.

Dashes come in three sizes: an
       intra-word
dash, a medium dash for number ranges like
       1--2,
and a punctuation
       dash---like
this.

A sentence-ending space should be larger than the
space between words within a sentence.  You
sometimes have to type special commands in
conjunction with punctuation characters to get
this right, as in the following sentence.
       Gnats, gnus, etc.\ all  % `\ ' makes an inter-word space.
       begin with G\@.         % \@ marks end-of-sentence punctuation.
You should check the spaces after periods when
reading your output to make sure you haven't
forgotten any special cases.  Generating an
ellipsis
       \ldots\               % `\ ' is needed after `\ldots' because TeX
                             % ignores spaces after command names like \ldots
                             % made from \ + letters.
                             %
                             % Note how a `%' character causes TeX to ignore
                             % the end of the input line, so these blank lines
                             % do not start a new paragraph.
                             %
with the right spacing around the periods requires
a special command.

\LaTeX\ interprets some common characters as
commands, so you must type special commands to
generate them.  These characters include the
following:
       \$ \& \% \# \{ and \}.

In printing, text is usually emphasized with an
       \emph{italic}
type style.

\begin{em}
   A long segment of text can also be emphasized
   in this way.  Text within such a segment can be
   given \emph{additional} emphasis.
\end{em}

It is sometimes necessary to prevent \LaTeX\ from
breaking a line where it might otherwise do so.
This may be at a space, as between the ``Mr.''\ and
``Jones'' in
       ``Mr.~Jones'',        % ~ produces an unbreakable interword space.
or within a word---especially when the word is a
symbol like
       \mbox{\emph{itemnum}}
that makes little sense when hyphenated across
lines.

Footnotes\footnote{This is an example of a footnote.}
pose no problem.

\LaTeX\ is good at typesetting mathematical formulas
like
       \( x-3y + z = 7 \)
or
       \( a_{1} > x^{2n} + y^{2n} > x' \)
or
       \( \ip{A}{B} = \sum_{i} a_{i} b_{i} \).
The spaces you type in a formula are
ignored.  Remember that a letter like
       $x$                   % $ ... $  and  \( ... \)  are equivalent
is a formula when it denotes a mathematical
symbol, and it should be typed as one.

\section{Displayed Text}

Text is displayed by indenting it from the left
margin.  Quotations are commonly displayed.  There
are short quotations
\begin{quote}
   This is a short quotation.  It consists of a
   single paragraph of text.  See how it is formatted.
\end{quote}
and longer ones.
\begin{quotation}
   This is a longer quotation.  It consists of two
   paragraphs of text, neither of which are
   particularly interesting.

   This is the second paragraph of the quotation.  It
   is just as dull as the first paragraph.
\end{quotation}
Another frequently-displayed structure is a list.
The following is an example of an \emph{itemized}
list.
\begin{itemize}
   \item This is the first item of an itemized list.
         Each item in the list is marked with a ``tick''.
         You don't have to worry about what kind of tick
         mark is used.

   \item This is the second item of the list.  It
         contains another list nested inside it.  The inner
         list is an \emph{enumerated} list.
         \begin{enumerate}
            \item This is the first item of an enumerated
                  list that is nested within the itemized list.

            \item This is the second item of the inner list.
                  \LaTeX\ allows you to nest lists deeper than
                  you really should.
         \end{enumerate}
         This is the rest of the second item of the outer
         list.  It is no more interesting than any other
         part of the item.
   \item This is the third item of the list.
\end{itemize}
You can even display poetry.
\begin{verse}
   There is an environment
    for verse \\             % The \\ command separates lines
   Whose features some poets % within a stanza.
   will curse.

                             % One or more blank lines separate stanzas.

   For instead of making\\
   Them do \emph{all} line breaking, \\
   It allows them to put too many words on a line when they'd rather be
   forced to be terse.
\end{verse}

Mathematical formulas may also be displayed.  A
displayed formula
is
one-line long; multiline
formulas require special formatting instructions.
   \[  \ip{\Gamma}{\psi'} = x'' + y^{2} + z_{i}^{n}\]
Don't start a paragraph with a displayed equation,
nor make one a paragraph by itself.

\end{document}               % End of document.

\section{Bench Huge Section 476}
% synthetic bench section 476
% This is a sample LaTeX input file.  (Version of 12 August 2004.)
%
% A '%' character causes TeX to ignore all remaining text on the line,
% and is used for comments like this one.

\documentclass{article}      % Specifies the document class

                             % The preamble begins here.
\title{An Example Document}  % Declares the document's title.
\author{Leslie Lamport}      % Declares the author's name.
\date{January 21, 1994}      % Deleting this command produces today's date.

\newcommand{\ip}[2]{(#1, #2)}
                             % Defines \ip{arg1}{arg2} to mean
                             % (arg1, arg2).

%\newcommand{\ip}[2]{\langle #1 | #2\rangle}
                             % This is an alternative definition of
                             % \ip that is commented out.

\begin{document}             % End of preamble and beginning of text.

\maketitle                   % Produces the title.

This is an example input file.  Comparing it with
the output it generates can show you how to
produce a simple document of your own.

\section{Ordinary Text}      % Produces section heading.  Lower-level
                             % sections are begun with similar
                             % \subsection and \subsubsection commands.

The ends  of words and sentences are marked
  by   spaces. It  doesn't matter how many
spaces    you type; one is as good as 100.  The
end of   a line counts as a space.

One   or more   blank lines denote the  end
of  a paragraph.

Since any number of consecutive spaces are treated
like a single one, the formatting of the input
file makes no difference to
      \LaTeX,                % The \LaTeX command generates the LaTeX logo.
but it makes a difference to you.  When you use
\LaTeX, making your input file as easy to read
as possible will be a great help as you write
your document and when you change it.  This sample
file shows how you can add comments to your own input
file.

Because printing is different from typewriting,
there are a number of things that you have to do
differently when preparing an input file than if
you were just typing the document directly.
Quotation marks like
       ``this''
have to be handled specially, as do quotes within
quotes:
       ``\,`this'            % \, separates the double and single quote.
        is what I just
        wrote, not  `that'\,''.

Dashes come in three sizes: an
       intra-word
dash, a medium dash for number ranges like
       1--2,
and a punctuation
       dash---like
this.

A sentence-ending space should be larger than the
space between words within a sentence.  You
sometimes have to type special commands in
conjunction with punctuation characters to get
this right, as in the following sentence.
       Gnats, gnus, etc.\ all  % `\ ' makes an inter-word space.
       begin with G\@.         % \@ marks end-of-sentence punctuation.
You should check the spaces after periods when
reading your output to make sure you haven't
forgotten any special cases.  Generating an
ellipsis
       \ldots\               % `\ ' is needed after `\ldots' because TeX
                             % ignores spaces after command names like \ldots
                             % made from \ + letters.
                             %
                             % Note how a `%' character causes TeX to ignore
                             % the end of the input line, so these blank lines
                             % do not start a new paragraph.
                             %
with the right spacing around the periods requires
a special command.

\LaTeX\ interprets some common characters as
commands, so you must type special commands to
generate them.  These characters include the
following:
       \$ \& \% \# \{ and \}.

In printing, text is usually emphasized with an
       \emph{italic}
type style.

\begin{em}
   A long segment of text can also be emphasized
   in this way.  Text within such a segment can be
   given \emph{additional} emphasis.
\end{em}

It is sometimes necessary to prevent \LaTeX\ from
breaking a line where it might otherwise do so.
This may be at a space, as between the ``Mr.''\ and
``Jones'' in
       ``Mr.~Jones'',        % ~ produces an unbreakable interword space.
or within a word---especially when the word is a
symbol like
       \mbox{\emph{itemnum}}
that makes little sense when hyphenated across
lines.

Footnotes\footnote{This is an example of a footnote.}
pose no problem.

\LaTeX\ is good at typesetting mathematical formulas
like
       \( x-3y + z = 7 \)
or
       \( a_{1} > x^{2n} + y^{2n} > x' \)
or
       \( \ip{A}{B} = \sum_{i} a_{i} b_{i} \).
The spaces you type in a formula are
ignored.  Remember that a letter like
       $x$                   % $ ... $  and  \( ... \)  are equivalent
is a formula when it denotes a mathematical
symbol, and it should be typed as one.

\section{Displayed Text}

Text is displayed by indenting it from the left
margin.  Quotations are commonly displayed.  There
are short quotations
\begin{quote}
   This is a short quotation.  It consists of a
   single paragraph of text.  See how it is formatted.
\end{quote}
and longer ones.
\begin{quotation}
   This is a longer quotation.  It consists of two
   paragraphs of text, neither of which are
   particularly interesting.

   This is the second paragraph of the quotation.  It
   is just as dull as the first paragraph.
\end{quotation}
Another frequently-displayed structure is a list.
The following is an example of an \emph{itemized}
list.
\begin{itemize}
   \item This is the first item of an itemized list.
         Each item in the list is marked with a ``tick''.
         You don't have to worry about what kind of tick
         mark is used.

   \item This is the second item of the list.  It
         contains another list nested inside it.  The inner
         list is an \emph{enumerated} list.
         \begin{enumerate}
            \item This is the first item of an enumerated
                  list that is nested within the itemized list.

            \item This is the second item of the inner list.
                  \LaTeX\ allows you to nest lists deeper than
                  you really should.
         \end{enumerate}
         This is the rest of the second item of the outer
         list.  It is no more interesting than any other
         part of the item.
   \item This is the third item of the list.
\end{itemize}
You can even display poetry.
\begin{verse}
   There is an environment
    for verse \\             % The \\ command separates lines
   Whose features some poets % within a stanza.
   will curse.

                             % One or more blank lines separate stanzas.

   For instead of making\\
   Them do \emph{all} line breaking, \\
   It allows them to put too many words on a line when they'd rather be
   forced to be terse.
\end{verse}

Mathematical formulas may also be displayed.  A
displayed formula
is
one-line long; multiline
formulas require special formatting instructions.
   \[  \ip{\Gamma}{\psi'} = x'' + y^{2} + z_{i}^{n}\]
Don't start a paragraph with a displayed equation,
nor make one a paragraph by itself.

\end{document}               % End of document.

\section{Bench Huge Section 477}
% synthetic bench section 477
% This is a sample LaTeX input file.  (Version of 12 August 2004.)
%
% A '%' character causes TeX to ignore all remaining text on the line,
% and is used for comments like this one.

\documentclass{article}      % Specifies the document class

                             % The preamble begins here.
\title{An Example Document}  % Declares the document's title.
\author{Leslie Lamport}      % Declares the author's name.
\date{January 21, 1994}      % Deleting this command produces today's date.

\newcommand{\ip}[2]{(#1, #2)}
                             % Defines \ip{arg1}{arg2} to mean
                             % (arg1, arg2).

%\newcommand{\ip}[2]{\langle #1 | #2\rangle}
                             % This is an alternative definition of
                             % \ip that is commented out.

\begin{document}             % End of preamble and beginning of text.

\maketitle                   % Produces the title.

This is an example input file.  Comparing it with
the output it generates can show you how to
produce a simple document of your own.

\section{Ordinary Text}      % Produces section heading.  Lower-level
                             % sections are begun with similar
                             % \subsection and \subsubsection commands.

The ends  of words and sentences are marked
  by   spaces. It  doesn't matter how many
spaces    you type; one is as good as 100.  The
end of   a line counts as a space.

One   or more   blank lines denote the  end
of  a paragraph.

Since any number of consecutive spaces are treated
like a single one, the formatting of the input
file makes no difference to
      \LaTeX,                % The \LaTeX command generates the LaTeX logo.
but it makes a difference to you.  When you use
\LaTeX, making your input file as easy to read
as possible will be a great help as you write
your document and when you change it.  This sample
file shows how you can add comments to your own input
file.

Because printing is different from typewriting,
there are a number of things that you have to do
differently when preparing an input file than if
you were just typing the document directly.
Quotation marks like
       ``this''
have to be handled specially, as do quotes within
quotes:
       ``\,`this'            % \, separates the double and single quote.
        is what I just
        wrote, not  `that'\,''.

Dashes come in three sizes: an
       intra-word
dash, a medium dash for number ranges like
       1--2,
and a punctuation
       dash---like
this.

A sentence-ending space should be larger than the
space between words within a sentence.  You
sometimes have to type special commands in
conjunction with punctuation characters to get
this right, as in the following sentence.
       Gnats, gnus, etc.\ all  % `\ ' makes an inter-word space.
       begin with G\@.         % \@ marks end-of-sentence punctuation.
You should check the spaces after periods when
reading your output to make sure you haven't
forgotten any special cases.  Generating an
ellipsis
       \ldots\               % `\ ' is needed after `\ldots' because TeX
                             % ignores spaces after command names like \ldots
                             % made from \ + letters.
                             %
                             % Note how a `%' character causes TeX to ignore
                             % the end of the input line, so these blank lines
                             % do not start a new paragraph.
                             %
with the right spacing around the periods requires
a special command.

\LaTeX\ interprets some common characters as
commands, so you must type special commands to
generate them.  These characters include the
following:
       \$ \& \% \# \{ and \}.

In printing, text is usually emphasized with an
       \emph{italic}
type style.

\begin{em}
   A long segment of text can also be emphasized
   in this way.  Text within such a segment can be
   given \emph{additional} emphasis.
\end{em}

It is sometimes necessary to prevent \LaTeX\ from
breaking a line where it might otherwise do so.
This may be at a space, as between the ``Mr.''\ and
``Jones'' in
       ``Mr.~Jones'',        % ~ produces an unbreakable interword space.
or within a word---especially when the word is a
symbol like
       \mbox{\emph{itemnum}}
that makes little sense when hyphenated across
lines.

Footnotes\footnote{This is an example of a footnote.}
pose no problem.

\LaTeX\ is good at typesetting mathematical formulas
like
       \( x-3y + z = 7 \)
or
       \( a_{1} > x^{2n} + y^{2n} > x' \)
or
       \( \ip{A}{B} = \sum_{i} a_{i} b_{i} \).
The spaces you type in a formula are
ignored.  Remember that a letter like
       $x$                   % $ ... $  and  \( ... \)  are equivalent
is a formula when it denotes a mathematical
symbol, and it should be typed as one.

\section{Displayed Text}

Text is displayed by indenting it from the left
margin.  Quotations are commonly displayed.  There
are short quotations
\begin{quote}
   This is a short quotation.  It consists of a
   single paragraph of text.  See how it is formatted.
\end{quote}
and longer ones.
\begin{quotation}
   This is a longer quotation.  It consists of two
   paragraphs of text, neither of which are
   particularly interesting.

   This is the second paragraph of the quotation.  It
   is just as dull as the first paragraph.
\end{quotation}
Another frequently-displayed structure is a list.
The following is an example of an \emph{itemized}
list.
\begin{itemize}
   \item This is the first item of an itemized list.
         Each item in the list is marked with a ``tick''.
         You don't have to worry about what kind of tick
         mark is used.

   \item This is the second item of the list.  It
         contains another list nested inside it.  The inner
         list is an \emph{enumerated} list.
         \begin{enumerate}
            \item This is the first item of an enumerated
                  list that is nested within the itemized list.

            \item This is the second item of the inner list.
                  \LaTeX\ allows you to nest lists deeper than
                  you really should.
         \end{enumerate}
         This is the rest of the second item of the outer
         list.  It is no more interesting than any other
         part of the item.
   \item This is the third item of the list.
\end{itemize}
You can even display poetry.
\begin{verse}
   There is an environment
    for verse \\             % The \\ command separates lines
   Whose features some poets % within a stanza.
   will curse.

                             % One or more blank lines separate stanzas.

   For instead of making\\
   Them do \emph{all} line breaking, \\
   It allows them to put too many words on a line when they'd rather be
   forced to be terse.
\end{verse}

Mathematical formulas may also be displayed.  A
displayed formula
is
one-line long; multiline
formulas require special formatting instructions.
   \[  \ip{\Gamma}{\psi'} = x'' + y^{2} + z_{i}^{n}\]
Don't start a paragraph with a displayed equation,
nor make one a paragraph by itself.

\end{document}               % End of document.

\section{Bench Huge Section 478}
% synthetic bench section 478
% This is a sample LaTeX input file.  (Version of 12 August 2004.)
%
% A '%' character causes TeX to ignore all remaining text on the line,
% and is used for comments like this one.

\documentclass{article}      % Specifies the document class

                             % The preamble begins here.
\title{An Example Document}  % Declares the document's title.
\author{Leslie Lamport}      % Declares the author's name.
\date{January 21, 1994}      % Deleting this command produces today's date.

\newcommand{\ip}[2]{(#1, #2)}
                             % Defines \ip{arg1}{arg2} to mean
                             % (arg1, arg2).

%\newcommand{\ip}[2]{\langle #1 | #2\rangle}
                             % This is an alternative definition of
                             % \ip that is commented out.

\begin{document}             % End of preamble and beginning of text.

\maketitle                   % Produces the title.

This is an example input file.  Comparing it with
the output it generates can show you how to
produce a simple document of your own.

\section{Ordinary Text}      % Produces section heading.  Lower-level
                             % sections are begun with similar
                             % \subsection and \subsubsection commands.

The ends  of words and sentences are marked
  by   spaces. It  doesn't matter how many
spaces    you type; one is as good as 100.  The
end of   a line counts as a space.

One   or more   blank lines denote the  end
of  a paragraph.

Since any number of consecutive spaces are treated
like a single one, the formatting of the input
file makes no difference to
      \LaTeX,                % The \LaTeX command generates the LaTeX logo.
but it makes a difference to you.  When you use
\LaTeX, making your input file as easy to read
as possible will be a great help as you write
your document and when you change it.  This sample
file shows how you can add comments to your own input
file.

Because printing is different from typewriting,
there are a number of things that you have to do
differently when preparing an input file than if
you were just typing the document directly.
Quotation marks like
       ``this''
have to be handled specially, as do quotes within
quotes:
       ``\,`this'            % \, separates the double and single quote.
        is what I just
        wrote, not  `that'\,''.

Dashes come in three sizes: an
       intra-word
dash, a medium dash for number ranges like
       1--2,
and a punctuation
       dash---like
this.

A sentence-ending space should be larger than the
space between words within a sentence.  You
sometimes have to type special commands in
conjunction with punctuation characters to get
this right, as in the following sentence.
       Gnats, gnus, etc.\ all  % `\ ' makes an inter-word space.
       begin with G\@.         % \@ marks end-of-sentence punctuation.
You should check the spaces after periods when
reading your output to make sure you haven't
forgotten any special cases.  Generating an
ellipsis
       \ldots\               % `\ ' is needed after `\ldots' because TeX
                             % ignores spaces after command names like \ldots
                             % made from \ + letters.
                             %
                             % Note how a `%' character causes TeX to ignore
                             % the end of the input line, so these blank lines
                             % do not start a new paragraph.
                             %
with the right spacing around the periods requires
a special command.

\LaTeX\ interprets some common characters as
commands, so you must type special commands to
generate them.  These characters include the
following:
       \$ \& \% \# \{ and \}.

In printing, text is usually emphasized with an
       \emph{italic}
type style.

\begin{em}
   A long segment of text can also be emphasized
   in this way.  Text within such a segment can be
   given \emph{additional} emphasis.
\end{em}

It is sometimes necessary to prevent \LaTeX\ from
breaking a line where it might otherwise do so.
This may be at a space, as between the ``Mr.''\ and
``Jones'' in
       ``Mr.~Jones'',        % ~ produces an unbreakable interword space.
or within a word---especially when the word is a
symbol like
       \mbox{\emph{itemnum}}
that makes little sense when hyphenated across
lines.

Footnotes\footnote{This is an example of a footnote.}
pose no problem.

\LaTeX\ is good at typesetting mathematical formulas
like
       \( x-3y + z = 7 \)
or
       \( a_{1} > x^{2n} + y^{2n} > x' \)
or
       \( \ip{A}{B} = \sum_{i} a_{i} b_{i} \).
The spaces you type in a formula are
ignored.  Remember that a letter like
       $x$                   % $ ... $  and  \( ... \)  are equivalent
is a formula when it denotes a mathematical
symbol, and it should be typed as one.

\section{Displayed Text}

Text is displayed by indenting it from the left
margin.  Quotations are commonly displayed.  There
are short quotations
\begin{quote}
   This is a short quotation.  It consists of a
   single paragraph of text.  See how it is formatted.
\end{quote}
and longer ones.
\begin{quotation}
   This is a longer quotation.  It consists of two
   paragraphs of text, neither of which are
   particularly interesting.

   This is the second paragraph of the quotation.  It
   is just as dull as the first paragraph.
\end{quotation}
Another frequently-displayed structure is a list.
The following is an example of an \emph{itemized}
list.
\begin{itemize}
   \item This is the first item of an itemized list.
         Each item in the list is marked with a ``tick''.
         You don't have to worry about what kind of tick
         mark is used.

   \item This is the second item of the list.  It
         contains another list nested inside it.  The inner
         list is an \emph{enumerated} list.
         \begin{enumerate}
            \item This is the first item of an enumerated
                  list that is nested within the itemized list.

            \item This is the second item of the inner list.
                  \LaTeX\ allows you to nest lists deeper than
                  you really should.
         \end{enumerate}
         This is the rest of the second item of the outer
         list.  It is no more interesting than any other
         part of the item.
   \item This is the third item of the list.
\end{itemize}
You can even display poetry.
\begin{verse}
   There is an environment
    for verse \\             % The \\ command separates lines
   Whose features some poets % within a stanza.
   will curse.

                             % One or more blank lines separate stanzas.

   For instead of making\\
   Them do \emph{all} line breaking, \\
   It allows them to put too many words on a line when they'd rather be
   forced to be terse.
\end{verse}

Mathematical formulas may also be displayed.  A
displayed formula
is
one-line long; multiline
formulas require special formatting instructions.
   \[  \ip{\Gamma}{\psi'} = x'' + y^{2} + z_{i}^{n}\]
Don't start a paragraph with a displayed equation,
nor make one a paragraph by itself.

\end{document}               % End of document.

\section{Bench Huge Section 479}
% synthetic bench section 479
% This is a sample LaTeX input file.  (Version of 12 August 2004.)
%
% A '%' character causes TeX to ignore all remaining text on the line,
% and is used for comments like this one.

\documentclass{article}      % Specifies the document class

                             % The preamble begins here.
\title{An Example Document}  % Declares the document's title.
\author{Leslie Lamport}      % Declares the author's name.
\date{January 21, 1994}      % Deleting this command produces today's date.

\newcommand{\ip}[2]{(#1, #2)}
                             % Defines \ip{arg1}{arg2} to mean
                             % (arg1, arg2).

%\newcommand{\ip}[2]{\langle #1 | #2\rangle}
                             % This is an alternative definition of
                             % \ip that is commented out.

\begin{document}             % End of preamble and beginning of text.

\maketitle                   % Produces the title.

This is an example input file.  Comparing it with
the output it generates can show you how to
produce a simple document of your own.

\section{Ordinary Text}      % Produces section heading.  Lower-level
                             % sections are begun with similar
                             % \subsection and \subsubsection commands.

The ends  of words and sentences are marked
  by   spaces. It  doesn't matter how many
spaces    you type; one is as good as 100.  The
end of   a line counts as a space.

One   or more   blank lines denote the  end
of  a paragraph.

Since any number of consecutive spaces are treated
like a single one, the formatting of the input
file makes no difference to
      \LaTeX,                % The \LaTeX command generates the LaTeX logo.
but it makes a difference to you.  When you use
\LaTeX, making your input file as easy to read
as possible will be a great help as you write
your document and when you change it.  This sample
file shows how you can add comments to your own input
file.

Because printing is different from typewriting,
there are a number of things that you have to do
differently when preparing an input file than if
you were just typing the document directly.
Quotation marks like
       ``this''
have to be handled specially, as do quotes within
quotes:
       ``\,`this'            % \, separates the double and single quote.
        is what I just
        wrote, not  `that'\,''.

Dashes come in three sizes: an
       intra-word
dash, a medium dash for number ranges like
       1--2,
and a punctuation
       dash---like
this.

A sentence-ending space should be larger than the
space between words within a sentence.  You
sometimes have to type special commands in
conjunction with punctuation characters to get
this right, as in the following sentence.
       Gnats, gnus, etc.\ all  % `\ ' makes an inter-word space.
       begin with G\@.         % \@ marks end-of-sentence punctuation.
You should check the spaces after periods when
reading your output to make sure you haven't
forgotten any special cases.  Generating an
ellipsis
       \ldots\               % `\ ' is needed after `\ldots' because TeX
                             % ignores spaces after command names like \ldots
                             % made from \ + letters.
                             %
                             % Note how a `%' character causes TeX to ignore
                             % the end of the input line, so these blank lines
                             % do not start a new paragraph.
                             %
with the right spacing around the periods requires
a special command.

\LaTeX\ interprets some common characters as
commands, so you must type special commands to
generate them.  These characters include the
following:
       \$ \& \% \# \{ and \}.

In printing, text is usually emphasized with an
       \emph{italic}
type style.

\begin{em}
   A long segment of text can also be emphasized
   in this way.  Text within such a segment can be
   given \emph{additional} emphasis.
\end{em}

It is sometimes necessary to prevent \LaTeX\ from
breaking a line where it might otherwise do so.
This may be at a space, as between the ``Mr.''\ and
``Jones'' in
       ``Mr.~Jones'',        % ~ produces an unbreakable interword space.
or within a word---especially when the word is a
symbol like
       \mbox{\emph{itemnum}}
that makes little sense when hyphenated across
lines.

Footnotes\footnote{This is an example of a footnote.}
pose no problem.

\LaTeX\ is good at typesetting mathematical formulas
like
       \( x-3y + z = 7 \)
or
       \( a_{1} > x^{2n} + y^{2n} > x' \)
or
       \( \ip{A}{B} = \sum_{i} a_{i} b_{i} \).
The spaces you type in a formula are
ignored.  Remember that a letter like
       $x$                   % $ ... $  and  \( ... \)  are equivalent
is a formula when it denotes a mathematical
symbol, and it should be typed as one.

\section{Displayed Text}

Text is displayed by indenting it from the left
margin.  Quotations are commonly displayed.  There
are short quotations
\begin{quote}
   This is a short quotation.  It consists of a
   single paragraph of text.  See how it is formatted.
\end{quote}
and longer ones.
\begin{quotation}
   This is a longer quotation.  It consists of two
   paragraphs of text, neither of which are
   particularly interesting.

   This is the second paragraph of the quotation.  It
   is just as dull as the first paragraph.
\end{quotation}
Another frequently-displayed structure is a list.
The following is an example of an \emph{itemized}
list.
\begin{itemize}
   \item This is the first item of an itemized list.
         Each item in the list is marked with a ``tick''.
         You don't have to worry about what kind of tick
         mark is used.

   \item This is the second item of the list.  It
         contains another list nested inside it.  The inner
         list is an \emph{enumerated} list.
         \begin{enumerate}
            \item This is the first item of an enumerated
                  list that is nested within the itemized list.

            \item This is the second item of the inner list.
                  \LaTeX\ allows you to nest lists deeper than
                  you really should.
         \end{enumerate}
         This is the rest of the second item of the outer
         list.  It is no more interesting than any other
         part of the item.
   \item This is the third item of the list.
\end{itemize}
You can even display poetry.
\begin{verse}
   There is an environment
    for verse \\             % The \\ command separates lines
   Whose features some poets % within a stanza.
   will curse.

                             % One or more blank lines separate stanzas.

   For instead of making\\
   Them do \emph{all} line breaking, \\
   It allows them to put too many words on a line when they'd rather be
   forced to be terse.
\end{verse}

Mathematical formulas may also be displayed.  A
displayed formula
is
one-line long; multiline
formulas require special formatting instructions.
   \[  \ip{\Gamma}{\psi'} = x'' + y^{2} + z_{i}^{n}\]
Don't start a paragraph with a displayed equation,
nor make one a paragraph by itself.

\end{document}               % End of document.

\section{Bench Huge Section 480}
% synthetic bench section 480
% This is a sample LaTeX input file.  (Version of 12 August 2004.)
%
% A '%' character causes TeX to ignore all remaining text on the line,
% and is used for comments like this one.

\documentclass{article}      % Specifies the document class

                             % The preamble begins here.
\title{An Example Document}  % Declares the document's title.
\author{Leslie Lamport}      % Declares the author's name.
\date{January 21, 1994}      % Deleting this command produces today's date.

\newcommand{\ip}[2]{(#1, #2)}
                             % Defines \ip{arg1}{arg2} to mean
                             % (arg1, arg2).

%\newcommand{\ip}[2]{\langle #1 | #2\rangle}
                             % This is an alternative definition of
                             % \ip that is commented out.

\begin{document}             % End of preamble and beginning of text.

\maketitle                   % Produces the title.

This is an example input file.  Comparing it with
the output it generates can show you how to
produce a simple document of your own.

\section{Ordinary Text}      % Produces section heading.  Lower-level
                             % sections are begun with similar
                             % \subsection and \subsubsection commands.

The ends  of words and sentences are marked
  by   spaces. It  doesn't matter how many
spaces    you type; one is as good as 100.  The
end of   a line counts as a space.

One   or more   blank lines denote the  end
of  a paragraph.

Since any number of consecutive spaces are treated
like a single one, the formatting of the input
file makes no difference to
      \LaTeX,                % The \LaTeX command generates the LaTeX logo.
but it makes a difference to you.  When you use
\LaTeX, making your input file as easy to read
as possible will be a great help as you write
your document and when you change it.  This sample
file shows how you can add comments to your own input
file.

Because printing is different from typewriting,
there are a number of things that you have to do
differently when preparing an input file than if
you were just typing the document directly.
Quotation marks like
       ``this''
have to be handled specially, as do quotes within
quotes:
       ``\,`this'            % \, separates the double and single quote.
        is what I just
        wrote, not  `that'\,''.

Dashes come in three sizes: an
       intra-word
dash, a medium dash for number ranges like
       1--2,
and a punctuation
       dash---like
this.

A sentence-ending space should be larger than the
space between words within a sentence.  You
sometimes have to type special commands in
conjunction with punctuation characters to get
this right, as in the following sentence.
       Gnats, gnus, etc.\ all  % `\ ' makes an inter-word space.
       begin with G\@.         % \@ marks end-of-sentence punctuation.
You should check the spaces after periods when
reading your output to make sure you haven't
forgotten any special cases.  Generating an
ellipsis
       \ldots\               % `\ ' is needed after `\ldots' because TeX
                             % ignores spaces after command names like \ldots
                             % made from \ + letters.
                             %
                             % Note how a `%' character causes TeX to ignore
                             % the end of the input line, so these blank lines
                             % do not start a new paragraph.
                             %
with the right spacing around the periods requires
a special command.

\LaTeX\ interprets some common characters as
commands, so you must type special commands to
generate them.  These characters include the
following:
       \$ \& \% \# \{ and \}.

In printing, text is usually emphasized with an
       \emph{italic}
type style.

\begin{em}
   A long segment of text can also be emphasized
   in this way.  Text within such a segment can be
   given \emph{additional} emphasis.
\end{em}

It is sometimes necessary to prevent \LaTeX\ from
breaking a line where it might otherwise do so.
This may be at a space, as between the ``Mr.''\ and
``Jones'' in
       ``Mr.~Jones'',        % ~ produces an unbreakable interword space.
or within a word---especially when the word is a
symbol like
       \mbox{\emph{itemnum}}
that makes little sense when hyphenated across
lines.

Footnotes\footnote{This is an example of a footnote.}
pose no problem.

\LaTeX\ is good at typesetting mathematical formulas
like
       \( x-3y + z = 7 \)
or
       \( a_{1} > x^{2n} + y^{2n} > x' \)
or
       \( \ip{A}{B} = \sum_{i} a_{i} b_{i} \).
The spaces you type in a formula are
ignored.  Remember that a letter like
       $x$                   % $ ... $  and  \( ... \)  are equivalent
is a formula when it denotes a mathematical
symbol, and it should be typed as one.

\section{Displayed Text}

Text is displayed by indenting it from the left
margin.  Quotations are commonly displayed.  There
are short quotations
\begin{quote}
   This is a short quotation.  It consists of a
   single paragraph of text.  See how it is formatted.
\end{quote}
and longer ones.
\begin{quotation}
   This is a longer quotation.  It consists of two
   paragraphs of text, neither of which are
   particularly interesting.

   This is the second paragraph of the quotation.  It
   is just as dull as the first paragraph.
\end{quotation}
Another frequently-displayed structure is a list.
The following is an example of an \emph{itemized}
list.
\begin{itemize}
   \item This is the first item of an itemized list.
         Each item in the list is marked with a ``tick''.
         You don't have to worry about what kind of tick
         mark is used.

   \item This is the second item of the list.  It
         contains another list nested inside it.  The inner
         list is an \emph{enumerated} list.
         \begin{enumerate}
            \item This is the first item of an enumerated
                  list that is nested within the itemized list.

            \item This is the second item of the inner list.
                  \LaTeX\ allows you to nest lists deeper than
                  you really should.
         \end{enumerate}
         This is the rest of the second item of the outer
         list.  It is no more interesting than any other
         part of the item.
   \item This is the third item of the list.
\end{itemize}
You can even display poetry.
\begin{verse}
   There is an environment
    for verse \\             % The \\ command separates lines
   Whose features some poets % within a stanza.
   will curse.

                             % One or more blank lines separate stanzas.

   For instead of making\\
   Them do \emph{all} line breaking, \\
   It allows them to put too many words on a line when they'd rather be
   forced to be terse.
\end{verse}

Mathematical formulas may also be displayed.  A
displayed formula
is
one-line long; multiline
formulas require special formatting instructions.
   \[  \ip{\Gamma}{\psi'} = x'' + y^{2} + z_{i}^{n}\]
Don't start a paragraph with a displayed equation,
nor make one a paragraph by itself.

\end{document}               % End of document.

\section{Bench Huge Section 481}
% synthetic bench section 481
% This is a sample LaTeX input file.  (Version of 12 August 2004.)
%
% A '%' character causes TeX to ignore all remaining text on the line,
% and is used for comments like this one.

\documentclass{article}      % Specifies the document class

                             % The preamble begins here.
\title{An Example Document}  % Declares the document's title.
\author{Leslie Lamport}      % Declares the author's name.
\date{January 21, 1994}      % Deleting this command produces today's date.

\newcommand{\ip}[2]{(#1, #2)}
                             % Defines \ip{arg1}{arg2} to mean
                             % (arg1, arg2).

%\newcommand{\ip}[2]{\langle #1 | #2\rangle}
                             % This is an alternative definition of
                             % \ip that is commented out.

\begin{document}             % End of preamble and beginning of text.

\maketitle                   % Produces the title.

This is an example input file.  Comparing it with
the output it generates can show you how to
produce a simple document of your own.

\section{Ordinary Text}      % Produces section heading.  Lower-level
                             % sections are begun with similar
                             % \subsection and \subsubsection commands.

The ends  of words and sentences are marked
  by   spaces. It  doesn't matter how many
spaces    you type; one is as good as 100.  The
end of   a line counts as a space.

One   or more   blank lines denote the  end
of  a paragraph.

Since any number of consecutive spaces are treated
like a single one, the formatting of the input
file makes no difference to
      \LaTeX,                % The \LaTeX command generates the LaTeX logo.
but it makes a difference to you.  When you use
\LaTeX, making your input file as easy to read
as possible will be a great help as you write
your document and when you change it.  This sample
file shows how you can add comments to your own input
file.

Because printing is different from typewriting,
there are a number of things that you have to do
differently when preparing an input file than if
you were just typing the document directly.
Quotation marks like
       ``this''
have to be handled specially, as do quotes within
quotes:
       ``\,`this'            % \, separates the double and single quote.
        is what I just
        wrote, not  `that'\,''.

Dashes come in three sizes: an
       intra-word
dash, a medium dash for number ranges like
       1--2,
and a punctuation
       dash---like
this.

A sentence-ending space should be larger than the
space between words within a sentence.  You
sometimes have to type special commands in
conjunction with punctuation characters to get
this right, as in the following sentence.
       Gnats, gnus, etc.\ all  % `\ ' makes an inter-word space.
       begin with G\@.         % \@ marks end-of-sentence punctuation.
You should check the spaces after periods when
reading your output to make sure you haven't
forgotten any special cases.  Generating an
ellipsis
       \ldots\               % `\ ' is needed after `\ldots' because TeX
                             % ignores spaces after command names like \ldots
                             % made from \ + letters.
                             %
                             % Note how a `%' character causes TeX to ignore
                             % the end of the input line, so these blank lines
                             % do not start a new paragraph.
                             %
with the right spacing around the periods requires
a special command.

\LaTeX\ interprets some common characters as
commands, so you must type special commands to
generate them.  These characters include the
following:
       \$ \& \% \# \{ and \}.

In printing, text is usually emphasized with an
       \emph{italic}
type style.

\begin{em}
   A long segment of text can also be emphasized
   in this way.  Text within such a segment can be
   given \emph{additional} emphasis.
\end{em}

It is sometimes necessary to prevent \LaTeX\ from
breaking a line where it might otherwise do so.
This may be at a space, as between the ``Mr.''\ and
``Jones'' in
       ``Mr.~Jones'',        % ~ produces an unbreakable interword space.
or within a word---especially when the word is a
symbol like
       \mbox{\emph{itemnum}}
that makes little sense when hyphenated across
lines.

Footnotes\footnote{This is an example of a footnote.}
pose no problem.

\LaTeX\ is good at typesetting mathematical formulas
like
       \( x-3y + z = 7 \)
or
       \( a_{1} > x^{2n} + y^{2n} > x' \)
or
       \( \ip{A}{B} = \sum_{i} a_{i} b_{i} \).
The spaces you type in a formula are
ignored.  Remember that a letter like
       $x$                   % $ ... $  and  \( ... \)  are equivalent
is a formula when it denotes a mathematical
symbol, and it should be typed as one.

\section{Displayed Text}

Text is displayed by indenting it from the left
margin.  Quotations are commonly displayed.  There
are short quotations
\begin{quote}
   This is a short quotation.  It consists of a
   single paragraph of text.  See how it is formatted.
\end{quote}
and longer ones.
\begin{quotation}
   This is a longer quotation.  It consists of two
   paragraphs of text, neither of which are
   particularly interesting.

   This is the second paragraph of the quotation.  It
   is just as dull as the first paragraph.
\end{quotation}
Another frequently-displayed structure is a list.
The following is an example of an \emph{itemized}
list.
\begin{itemize}
   \item This is the first item of an itemized list.
         Each item in the list is marked with a ``tick''.
         You don't have to worry about what kind of tick
         mark is used.

   \item This is the second item of the list.  It
         contains another list nested inside it.  The inner
         list is an \emph{enumerated} list.
         \begin{enumerate}
            \item This is the first item of an enumerated
                  list that is nested within the itemized list.

            \item This is the second item of the inner list.
                  \LaTeX\ allows you to nest lists deeper than
                  you really should.
         \end{enumerate}
         This is the rest of the second item of the outer
         list.  It is no more interesting than any other
         part of the item.
   \item This is the third item of the list.
\end{itemize}
You can even display poetry.
\begin{verse}
   There is an environment
    for verse \\             % The \\ command separates lines
   Whose features some poets % within a stanza.
   will curse.

                             % One or more blank lines separate stanzas.

   For instead of making\\
   Them do \emph{all} line breaking, \\
   It allows them to put too many words on a line when they'd rather be
   forced to be terse.
\end{verse}

Mathematical formulas may also be displayed.  A
displayed formula
is
one-line long; multiline
formulas require special formatting instructions.
   \[  \ip{\Gamma}{\psi'} = x'' + y^{2} + z_{i}^{n}\]
Don't start a paragraph with a displayed equation,
nor make one a paragraph by itself.

\end{document}               % End of document.

\section{Bench Huge Section 482}
% synthetic bench section 482
% This is a sample LaTeX input file.  (Version of 12 August 2004.)
%
% A '%' character causes TeX to ignore all remaining text on the line,
% and is used for comments like this one.

\documentclass{article}      % Specifies the document class

                             % The preamble begins here.
\title{An Example Document}  % Declares the document's title.
\author{Leslie Lamport}      % Declares the author's name.
\date{January 21, 1994}      % Deleting this command produces today's date.

\newcommand{\ip}[2]{(#1, #2)}
                             % Defines \ip{arg1}{arg2} to mean
                             % (arg1, arg2).

%\newcommand{\ip}[2]{\langle #1 | #2\rangle}
                             % This is an alternative definition of
                             % \ip that is commented out.

\begin{document}             % End of preamble and beginning of text.

\maketitle                   % Produces the title.

This is an example input file.  Comparing it with
the output it generates can show you how to
produce a simple document of your own.

\section{Ordinary Text}      % Produces section heading.  Lower-level
                             % sections are begun with similar
                             % \subsection and \subsubsection commands.

The ends  of words and sentences are marked
  by   spaces. It  doesn't matter how many
spaces    you type; one is as good as 100.  The
end of   a line counts as a space.

One   or more   blank lines denote the  end
of  a paragraph.

Since any number of consecutive spaces are treated
like a single one, the formatting of the input
file makes no difference to
      \LaTeX,                % The \LaTeX command generates the LaTeX logo.
but it makes a difference to you.  When you use
\LaTeX, making your input file as easy to read
as possible will be a great help as you write
your document and when you change it.  This sample
file shows how you can add comments to your own input
file.

Because printing is different from typewriting,
there are a number of things that you have to do
differently when preparing an input file than if
you were just typing the document directly.
Quotation marks like
       ``this''
have to be handled specially, as do quotes within
quotes:
       ``\,`this'            % \, separates the double and single quote.
        is what I just
        wrote, not  `that'\,''.

Dashes come in three sizes: an
       intra-word
dash, a medium dash for number ranges like
       1--2,
and a punctuation
       dash---like
this.

A sentence-ending space should be larger than the
space between words within a sentence.  You
sometimes have to type special commands in
conjunction with punctuation characters to get
this right, as in the following sentence.
       Gnats, gnus, etc.\ all  % `\ ' makes an inter-word space.
       begin with G\@.         % \@ marks end-of-sentence punctuation.
You should check the spaces after periods when
reading your output to make sure you haven't
forgotten any special cases.  Generating an
ellipsis
       \ldots\               % `\ ' is needed after `\ldots' because TeX
                             % ignores spaces after command names like \ldots
                             % made from \ + letters.
                             %
                             % Note how a `%' character causes TeX to ignore
                             % the end of the input line, so these blank lines
                             % do not start a new paragraph.
                             %
with the right spacing around the periods requires
a special command.

\LaTeX\ interprets some common characters as
commands, so you must type special commands to
generate them.  These characters include the
following:
       \$ \& \% \# \{ and \}.

In printing, text is usually emphasized with an
       \emph{italic}
type style.

\begin{em}
   A long segment of text can also be emphasized
   in this way.  Text within such a segment can be
   given \emph{additional} emphasis.
\end{em}

It is sometimes necessary to prevent \LaTeX\ from
breaking a line where it might otherwise do so.
This may be at a space, as between the ``Mr.''\ and
``Jones'' in
       ``Mr.~Jones'',        % ~ produces an unbreakable interword space.
or within a word---especially when the word is a
symbol like
       \mbox{\emph{itemnum}}
that makes little sense when hyphenated across
lines.

Footnotes\footnote{This is an example of a footnote.}
pose no problem.

\LaTeX\ is good at typesetting mathematical formulas
like
       \( x-3y + z = 7 \)
or
       \( a_{1} > x^{2n} + y^{2n} > x' \)
or
       \( \ip{A}{B} = \sum_{i} a_{i} b_{i} \).
The spaces you type in a formula are
ignored.  Remember that a letter like
       $x$                   % $ ... $  and  \( ... \)  are equivalent
is a formula when it denotes a mathematical
symbol, and it should be typed as one.

\section{Displayed Text}

Text is displayed by indenting it from the left
margin.  Quotations are commonly displayed.  There
are short quotations
\begin{quote}
   This is a short quotation.  It consists of a
   single paragraph of text.  See how it is formatted.
\end{quote}
and longer ones.
\begin{quotation}
   This is a longer quotation.  It consists of two
   paragraphs of text, neither of which are
   particularly interesting.

   This is the second paragraph of the quotation.  It
   is just as dull as the first paragraph.
\end{quotation}
Another frequently-displayed structure is a list.
The following is an example of an \emph{itemized}
list.
\begin{itemize}
   \item This is the first item of an itemized list.
         Each item in the list is marked with a ``tick''.
         You don't have to worry about what kind of tick
         mark is used.

   \item This is the second item of the list.  It
         contains another list nested inside it.  The inner
         list is an \emph{enumerated} list.
         \begin{enumerate}
            \item This is the first item of an enumerated
                  list that is nested within the itemized list.

            \item This is the second item of the inner list.
                  \LaTeX\ allows you to nest lists deeper than
                  you really should.
         \end{enumerate}
         This is the rest of the second item of the outer
         list.  It is no more interesting than any other
         part of the item.
   \item This is the third item of the list.
\end{itemize}
You can even display poetry.
\begin{verse}
   There is an environment
    for verse \\             % The \\ command separates lines
   Whose features some poets % within a stanza.
   will curse.

                             % One or more blank lines separate stanzas.

   For instead of making\\
   Them do \emph{all} line breaking, \\
   It allows them to put too many words on a line when they'd rather be
   forced to be terse.
\end{verse}

Mathematical formulas may also be displayed.  A
displayed formula
is
one-line long; multiline
formulas require special formatting instructions.
   \[  \ip{\Gamma}{\psi'} = x'' + y^{2} + z_{i}^{n}\]
Don't start a paragraph with a displayed equation,
nor make one a paragraph by itself.

\end{document}               % End of document.

\section{Bench Huge Section 483}
% synthetic bench section 483
% This is a sample LaTeX input file.  (Version of 12 August 2004.)
%
% A '%' character causes TeX to ignore all remaining text on the line,
% and is used for comments like this one.

\documentclass{article}      % Specifies the document class

                             % The preamble begins here.
\title{An Example Document}  % Declares the document's title.
\author{Leslie Lamport}      % Declares the author's name.
\date{January 21, 1994}      % Deleting this command produces today's date.

\newcommand{\ip}[2]{(#1, #2)}
                             % Defines \ip{arg1}{arg2} to mean
                             % (arg1, arg2).

%\newcommand{\ip}[2]{\langle #1 | #2\rangle}
                             % This is an alternative definition of
                             % \ip that is commented out.

\begin{document}             % End of preamble and beginning of text.

\maketitle                   % Produces the title.

This is an example input file.  Comparing it with
the output it generates can show you how to
produce a simple document of your own.

\section{Ordinary Text}      % Produces section heading.  Lower-level
                             % sections are begun with similar
                             % \subsection and \subsubsection commands.

The ends  of words and sentences are marked
  by   spaces. It  doesn't matter how many
spaces    you type; one is as good as 100.  The
end of   a line counts as a space.

One   or more   blank lines denote the  end
of  a paragraph.

Since any number of consecutive spaces are treated
like a single one, the formatting of the input
file makes no difference to
      \LaTeX,                % The \LaTeX command generates the LaTeX logo.
but it makes a difference to you.  When you use
\LaTeX, making your input file as easy to read
as possible will be a great help as you write
your document and when you change it.  This sample
file shows how you can add comments to your own input
file.

Because printing is different from typewriting,
there are a number of things that you have to do
differently when preparing an input file than if
you were just typing the document directly.
Quotation marks like
       ``this''
have to be handled specially, as do quotes within
quotes:
       ``\,`this'            % \, separates the double and single quote.
        is what I just
        wrote, not  `that'\,''.

Dashes come in three sizes: an
       intra-word
dash, a medium dash for number ranges like
       1--2,
and a punctuation
       dash---like
this.

A sentence-ending space should be larger than the
space between words within a sentence.  You
sometimes have to type special commands in
conjunction with punctuation characters to get
this right, as in the following sentence.
       Gnats, gnus, etc.\ all  % `\ ' makes an inter-word space.
       begin with G\@.         % \@ marks end-of-sentence punctuation.
You should check the spaces after periods when
reading your output to make sure you haven't
forgotten any special cases.  Generating an
ellipsis
       \ldots\               % `\ ' is needed after `\ldots' because TeX
                             % ignores spaces after command names like \ldots
                             % made from \ + letters.
                             %
                             % Note how a `%' character causes TeX to ignore
                             % the end of the input line, so these blank lines
                             % do not start a new paragraph.
                             %
with the right spacing around the periods requires
a special command.

\LaTeX\ interprets some common characters as
commands, so you must type special commands to
generate them.  These characters include the
following:
       \$ \& \% \# \{ and \}.

In printing, text is usually emphasized with an
       \emph{italic}
type style.

\begin{em}
   A long segment of text can also be emphasized
   in this way.  Text within such a segment can be
   given \emph{additional} emphasis.
\end{em}

It is sometimes necessary to prevent \LaTeX\ from
breaking a line where it might otherwise do so.
This may be at a space, as between the ``Mr.''\ and
``Jones'' in
       ``Mr.~Jones'',        % ~ produces an unbreakable interword space.
or within a word---especially when the word is a
symbol like
       \mbox{\emph{itemnum}}
that makes little sense when hyphenated across
lines.

Footnotes\footnote{This is an example of a footnote.}
pose no problem.

\LaTeX\ is good at typesetting mathematical formulas
like
       \( x-3y + z = 7 \)
or
       \( a_{1} > x^{2n} + y^{2n} > x' \)
or
       \( \ip{A}{B} = \sum_{i} a_{i} b_{i} \).
The spaces you type in a formula are
ignored.  Remember that a letter like
       $x$                   % $ ... $  and  \( ... \)  are equivalent
is a formula when it denotes a mathematical
symbol, and it should be typed as one.

\section{Displayed Text}

Text is displayed by indenting it from the left
margin.  Quotations are commonly displayed.  There
are short quotations
\begin{quote}
   This is a short quotation.  It consists of a
   single paragraph of text.  See how it is formatted.
\end{quote}
and longer ones.
\begin{quotation}
   This is a longer quotation.  It consists of two
   paragraphs of text, neither of which are
   particularly interesting.

   This is the second paragraph of the quotation.  It
   is just as dull as the first paragraph.
\end{quotation}
Another frequently-displayed structure is a list.
The following is an example of an \emph{itemized}
list.
\begin{itemize}
   \item This is the first item of an itemized list.
         Each item in the list is marked with a ``tick''.
         You don't have to worry about what kind of tick
         mark is used.

   \item This is the second item of the list.  It
         contains another list nested inside it.  The inner
         list is an \emph{enumerated} list.
         \begin{enumerate}
            \item This is the first item of an enumerated
                  list that is nested within the itemized list.

            \item This is the second item of the inner list.
                  \LaTeX\ allows you to nest lists deeper than
                  you really should.
         \end{enumerate}
         This is the rest of the second item of the outer
         list.  It is no more interesting than any other
         part of the item.
   \item This is the third item of the list.
\end{itemize}
You can even display poetry.
\begin{verse}
   There is an environment
    for verse \\             % The \\ command separates lines
   Whose features some poets % within a stanza.
   will curse.

                             % One or more blank lines separate stanzas.

   For instead of making\\
   Them do \emph{all} line breaking, \\
   It allows them to put too many words on a line when they'd rather be
   forced to be terse.
\end{verse}

Mathematical formulas may also be displayed.  A
displayed formula
is
one-line long; multiline
formulas require special formatting instructions.
   \[  \ip{\Gamma}{\psi'} = x'' + y^{2} + z_{i}^{n}\]
Don't start a paragraph with a displayed equation,
nor make one a paragraph by itself.

\end{document}               % End of document.

\section{Bench Huge Section 484}
% synthetic bench section 484
% This is a sample LaTeX input file.  (Version of 12 August 2004.)
%
% A '%' character causes TeX to ignore all remaining text on the line,
% and is used for comments like this one.

\documentclass{article}      % Specifies the document class

                             % The preamble begins here.
\title{An Example Document}  % Declares the document's title.
\author{Leslie Lamport}      % Declares the author's name.
\date{January 21, 1994}      % Deleting this command produces today's date.

\newcommand{\ip}[2]{(#1, #2)}
                             % Defines \ip{arg1}{arg2} to mean
                             % (arg1, arg2).

%\newcommand{\ip}[2]{\langle #1 | #2\rangle}
                             % This is an alternative definition of
                             % \ip that is commented out.

\begin{document}             % End of preamble and beginning of text.

\maketitle                   % Produces the title.

This is an example input file.  Comparing it with
the output it generates can show you how to
produce a simple document of your own.

\section{Ordinary Text}      % Produces section heading.  Lower-level
                             % sections are begun with similar
                             % \subsection and \subsubsection commands.

The ends  of words and sentences are marked
  by   spaces. It  doesn't matter how many
spaces    you type; one is as good as 100.  The
end of   a line counts as a space.

One   or more   blank lines denote the  end
of  a paragraph.

Since any number of consecutive spaces are treated
like a single one, the formatting of the input
file makes no difference to
      \LaTeX,                % The \LaTeX command generates the LaTeX logo.
but it makes a difference to you.  When you use
\LaTeX, making your input file as easy to read
as possible will be a great help as you write
your document and when you change it.  This sample
file shows how you can add comments to your own input
file.

Because printing is different from typewriting,
there are a number of things that you have to do
differently when preparing an input file than if
you were just typing the document directly.
Quotation marks like
       ``this''
have to be handled specially, as do quotes within
quotes:
       ``\,`this'            % \, separates the double and single quote.
        is what I just
        wrote, not  `that'\,''.

Dashes come in three sizes: an
       intra-word
dash, a medium dash for number ranges like
       1--2,
and a punctuation
       dash---like
this.

A sentence-ending space should be larger than the
space between words within a sentence.  You
sometimes have to type special commands in
conjunction with punctuation characters to get
this right, as in the following sentence.
       Gnats, gnus, etc.\ all  % `\ ' makes an inter-word space.
       begin with G\@.         % \@ marks end-of-sentence punctuation.
You should check the spaces after periods when
reading your output to make sure you haven't
forgotten any special cases.  Generating an
ellipsis
       \ldots\               % `\ ' is needed after `\ldots' because TeX
                             % ignores spaces after command names like \ldots
                             % made from \ + letters.
                             %
                             % Note how a `%' character causes TeX to ignore
                             % the end of the input line, so these blank lines
                             % do not start a new paragraph.
                             %
with the right spacing around the periods requires
a special command.

\LaTeX\ interprets some common characters as
commands, so you must type special commands to
generate them.  These characters include the
following:
       \$ \& \% \# \{ and \}.

In printing, text is usually emphasized with an
       \emph{italic}
type style.

\begin{em}
   A long segment of text can also be emphasized
   in this way.  Text within such a segment can be
   given \emph{additional} emphasis.
\end{em}

It is sometimes necessary to prevent \LaTeX\ from
breaking a line where it might otherwise do so.
This may be at a space, as between the ``Mr.''\ and
``Jones'' in
       ``Mr.~Jones'',        % ~ produces an unbreakable interword space.
or within a word---especially when the word is a
symbol like
       \mbox{\emph{itemnum}}
that makes little sense when hyphenated across
lines.

Footnotes\footnote{This is an example of a footnote.}
pose no problem.

\LaTeX\ is good at typesetting mathematical formulas
like
       \( x-3y + z = 7 \)
or
       \( a_{1} > x^{2n} + y^{2n} > x' \)
or
       \( \ip{A}{B} = \sum_{i} a_{i} b_{i} \).
The spaces you type in a formula are
ignored.  Remember that a letter like
       $x$                   % $ ... $  and  \( ... \)  are equivalent
is a formula when it denotes a mathematical
symbol, and it should be typed as one.

\section{Displayed Text}

Text is displayed by indenting it from the left
margin.  Quotations are commonly displayed.  There
are short quotations
\begin{quote}
   This is a short quotation.  It consists of a
   single paragraph of text.  See how it is formatted.
\end{quote}
and longer ones.
\begin{quotation}
   This is a longer quotation.  It consists of two
   paragraphs of text, neither of which are
   particularly interesting.

   This is the second paragraph of the quotation.  It
   is just as dull as the first paragraph.
\end{quotation}
Another frequently-displayed structure is a list.
The following is an example of an \emph{itemized}
list.
\begin{itemize}
   \item This is the first item of an itemized list.
         Each item in the list is marked with a ``tick''.
         You don't have to worry about what kind of tick
         mark is used.

   \item This is the second item of the list.  It
         contains another list nested inside it.  The inner
         list is an \emph{enumerated} list.
         \begin{enumerate}
            \item This is the first item of an enumerated
                  list that is nested within the itemized list.

            \item This is the second item of the inner list.
                  \LaTeX\ allows you to nest lists deeper than
                  you really should.
         \end{enumerate}
         This is the rest of the second item of the outer
         list.  It is no more interesting than any other
         part of the item.
   \item This is the third item of the list.
\end{itemize}
You can even display poetry.
\begin{verse}
   There is an environment
    for verse \\             % The \\ command separates lines
   Whose features some poets % within a stanza.
   will curse.

                             % One or more blank lines separate stanzas.

   For instead of making\\
   Them do \emph{all} line breaking, \\
   It allows them to put too many words on a line when they'd rather be
   forced to be terse.
\end{verse}

Mathematical formulas may also be displayed.  A
displayed formula
is
one-line long; multiline
formulas require special formatting instructions.
   \[  \ip{\Gamma}{\psi'} = x'' + y^{2} + z_{i}^{n}\]
Don't start a paragraph with a displayed equation,
nor make one a paragraph by itself.

\end{document}               % End of document.

\section{Bench Huge Section 485}
% synthetic bench section 485
% This is a sample LaTeX input file.  (Version of 12 August 2004.)
%
% A '%' character causes TeX to ignore all remaining text on the line,
% and is used for comments like this one.

\documentclass{article}      % Specifies the document class

                             % The preamble begins here.
\title{An Example Document}  % Declares the document's title.
\author{Leslie Lamport}      % Declares the author's name.
\date{January 21, 1994}      % Deleting this command produces today's date.

\newcommand{\ip}[2]{(#1, #2)}
                             % Defines \ip{arg1}{arg2} to mean
                             % (arg1, arg2).

%\newcommand{\ip}[2]{\langle #1 | #2\rangle}
                             % This is an alternative definition of
                             % \ip that is commented out.

\begin{document}             % End of preamble and beginning of text.

\maketitle                   % Produces the title.

This is an example input file.  Comparing it with
the output it generates can show you how to
produce a simple document of your own.

\section{Ordinary Text}      % Produces section heading.  Lower-level
                             % sections are begun with similar
                             % \subsection and \subsubsection commands.

The ends  of words and sentences are marked
  by   spaces. It  doesn't matter how many
spaces    you type; one is as good as 100.  The
end of   a line counts as a space.

One   or more   blank lines denote the  end
of  a paragraph.

Since any number of consecutive spaces are treated
like a single one, the formatting of the input
file makes no difference to
      \LaTeX,                % The \LaTeX command generates the LaTeX logo.
but it makes a difference to you.  When you use
\LaTeX, making your input file as easy to read
as possible will be a great help as you write
your document and when you change it.  This sample
file shows how you can add comments to your own input
file.

Because printing is different from typewriting,
there are a number of things that you have to do
differently when preparing an input file than if
you were just typing the document directly.
Quotation marks like
       ``this''
have to be handled specially, as do quotes within
quotes:
       ``\,`this'            % \, separates the double and single quote.
        is what I just
        wrote, not  `that'\,''.

Dashes come in three sizes: an
       intra-word
dash, a medium dash for number ranges like
       1--2,
and a punctuation
       dash---like
this.

A sentence-ending space should be larger than the
space between words within a sentence.  You
sometimes have to type special commands in
conjunction with punctuation characters to get
this right, as in the following sentence.
       Gnats, gnus, etc.\ all  % `\ ' makes an inter-word space.
       begin with G\@.         % \@ marks end-of-sentence punctuation.
You should check the spaces after periods when
reading your output to make sure you haven't
forgotten any special cases.  Generating an
ellipsis
       \ldots\               % `\ ' is needed after `\ldots' because TeX
                             % ignores spaces after command names like \ldots
                             % made from \ + letters.
                             %
                             % Note how a `%' character causes TeX to ignore
                             % the end of the input line, so these blank lines
                             % do not start a new paragraph.
                             %
with the right spacing around the periods requires
a special command.

\LaTeX\ interprets some common characters as
commands, so you must type special commands to
generate them.  These characters include the
following:
       \$ \& \% \# \{ and \}.

In printing, text is usually emphasized with an
       \emph{italic}
type style.

\begin{em}
   A long segment of text can also be emphasized
   in this way.  Text within such a segment can be
   given \emph{additional} emphasis.
\end{em}

It is sometimes necessary to prevent \LaTeX\ from
breaking a line where it might otherwise do so.
This may be at a space, as between the ``Mr.''\ and
``Jones'' in
       ``Mr.~Jones'',        % ~ produces an unbreakable interword space.
or within a word---especially when the word is a
symbol like
       \mbox{\emph{itemnum}}
that makes little sense when hyphenated across
lines.

Footnotes\footnote{This is an example of a footnote.}
pose no problem.

\LaTeX\ is good at typesetting mathematical formulas
like
       \( x-3y + z = 7 \)
or
       \( a_{1} > x^{2n} + y^{2n} > x' \)
or
       \( \ip{A}{B} = \sum_{i} a_{i} b_{i} \).
The spaces you type in a formula are
ignored.  Remember that a letter like
       $x$                   % $ ... $  and  \( ... \)  are equivalent
is a formula when it denotes a mathematical
symbol, and it should be typed as one.

\section{Displayed Text}

Text is displayed by indenting it from the left
margin.  Quotations are commonly displayed.  There
are short quotations
\begin{quote}
   This is a short quotation.  It consists of a
   single paragraph of text.  See how it is formatted.
\end{quote}
and longer ones.
\begin{quotation}
   This is a longer quotation.  It consists of two
   paragraphs of text, neither of which are
   particularly interesting.

   This is the second paragraph of the quotation.  It
   is just as dull as the first paragraph.
\end{quotation}
Another frequently-displayed structure is a list.
The following is an example of an \emph{itemized}
list.
\begin{itemize}
   \item This is the first item of an itemized list.
         Each item in the list is marked with a ``tick''.
         You don't have to worry about what kind of tick
         mark is used.

   \item This is the second item of the list.  It
         contains another list nested inside it.  The inner
         list is an \emph{enumerated} list.
         \begin{enumerate}
            \item This is the first item of an enumerated
                  list that is nested within the itemized list.

            \item This is the second item of the inner list.
                  \LaTeX\ allows you to nest lists deeper than
                  you really should.
         \end{enumerate}
         This is the rest of the second item of the outer
         list.  It is no more interesting than any other
         part of the item.
   \item This is the third item of the list.
\end{itemize}
You can even display poetry.
\begin{verse}
   There is an environment
    for verse \\             % The \\ command separates lines
   Whose features some poets % within a stanza.
   will curse.

                             % One or more blank lines separate stanzas.

   For instead of making\\
   Them do \emph{all} line breaking, \\
   It allows them to put too many words on a line when they'd rather be
   forced to be terse.
\end{verse}

Mathematical formulas may also be displayed.  A
displayed formula
is
one-line long; multiline
formulas require special formatting instructions.
   \[  \ip{\Gamma}{\psi'} = x'' + y^{2} + z_{i}^{n}\]
Don't start a paragraph with a displayed equation,
nor make one a paragraph by itself.

\end{document}               % End of document.

\section{Bench Huge Section 486}
% synthetic bench section 486
% This is a sample LaTeX input file.  (Version of 12 August 2004.)
%
% A '%' character causes TeX to ignore all remaining text on the line,
% and is used for comments like this one.

\documentclass{article}      % Specifies the document class

                             % The preamble begins here.
\title{An Example Document}  % Declares the document's title.
\author{Leslie Lamport}      % Declares the author's name.
\date{January 21, 1994}      % Deleting this command produces today's date.

\newcommand{\ip}[2]{(#1, #2)}
                             % Defines \ip{arg1}{arg2} to mean
                             % (arg1, arg2).

%\newcommand{\ip}[2]{\langle #1 | #2\rangle}
                             % This is an alternative definition of
                             % \ip that is commented out.

\begin{document}             % End of preamble and beginning of text.

\maketitle                   % Produces the title.

This is an example input file.  Comparing it with
the output it generates can show you how to
produce a simple document of your own.

\section{Ordinary Text}      % Produces section heading.  Lower-level
                             % sections are begun with similar
                             % \subsection and \subsubsection commands.

The ends  of words and sentences are marked
  by   spaces. It  doesn't matter how many
spaces    you type; one is as good as 100.  The
end of   a line counts as a space.

One   or more   blank lines denote the  end
of  a paragraph.

Since any number of consecutive spaces are treated
like a single one, the formatting of the input
file makes no difference to
      \LaTeX,                % The \LaTeX command generates the LaTeX logo.
but it makes a difference to you.  When you use
\LaTeX, making your input file as easy to read
as possible will be a great help as you write
your document and when you change it.  This sample
file shows how you can add comments to your own input
file.

Because printing is different from typewriting,
there are a number of things that you have to do
differently when preparing an input file than if
you were just typing the document directly.
Quotation marks like
       ``this''
have to be handled specially, as do quotes within
quotes:
       ``\,`this'            % \, separates the double and single quote.
        is what I just
        wrote, not  `that'\,''.

Dashes come in three sizes: an
       intra-word
dash, a medium dash for number ranges like
       1--2,
and a punctuation
       dash---like
this.

A sentence-ending space should be larger than the
space between words within a sentence.  You
sometimes have to type special commands in
conjunction with punctuation characters to get
this right, as in the following sentence.
       Gnats, gnus, etc.\ all  % `\ ' makes an inter-word space.
       begin with G\@.         % \@ marks end-of-sentence punctuation.
You should check the spaces after periods when
reading your output to make sure you haven't
forgotten any special cases.  Generating an
ellipsis
       \ldots\               % `\ ' is needed after `\ldots' because TeX
                             % ignores spaces after command names like \ldots
                             % made from \ + letters.
                             %
                             % Note how a `%' character causes TeX to ignore
                             % the end of the input line, so these blank lines
                             % do not start a new paragraph.
                             %
with the right spacing around the periods requires
a special command.

\LaTeX\ interprets some common characters as
commands, so you must type special commands to
generate them.  These characters include the
following:
       \$ \& \% \# \{ and \}.

In printing, text is usually emphasized with an
       \emph{italic}
type style.

\begin{em}
   A long segment of text can also be emphasized
   in this way.  Text within such a segment can be
   given \emph{additional} emphasis.
\end{em}

It is sometimes necessary to prevent \LaTeX\ from
breaking a line where it might otherwise do so.
This may be at a space, as between the ``Mr.''\ and
``Jones'' in
       ``Mr.~Jones'',        % ~ produces an unbreakable interword space.
or within a word---especially when the word is a
symbol like
       \mbox{\emph{itemnum}}
that makes little sense when hyphenated across
lines.

Footnotes\footnote{This is an example of a footnote.}
pose no problem.

\LaTeX\ is good at typesetting mathematical formulas
like
       \( x-3y + z = 7 \)
or
       \( a_{1} > x^{2n} + y^{2n} > x' \)
or
       \( \ip{A}{B} = \sum_{i} a_{i} b_{i} \).
The spaces you type in a formula are
ignored.  Remember that a letter like
       $x$                   % $ ... $  and  \( ... \)  are equivalent
is a formula when it denotes a mathematical
symbol, and it should be typed as one.

\section{Displayed Text}

Text is displayed by indenting it from the left
margin.  Quotations are commonly displayed.  There
are short quotations
\begin{quote}
   This is a short quotation.  It consists of a
   single paragraph of text.  See how it is formatted.
\end{quote}
and longer ones.
\begin{quotation}
   This is a longer quotation.  It consists of two
   paragraphs of text, neither of which are
   particularly interesting.

   This is the second paragraph of the quotation.  It
   is just as dull as the first paragraph.
\end{quotation}
Another frequently-displayed structure is a list.
The following is an example of an \emph{itemized}
list.
\begin{itemize}
   \item This is the first item of an itemized list.
         Each item in the list is marked with a ``tick''.
         You don't have to worry about what kind of tick
         mark is used.

   \item This is the second item of the list.  It
         contains another list nested inside it.  The inner
         list is an \emph{enumerated} list.
         \begin{enumerate}
            \item This is the first item of an enumerated
                  list that is nested within the itemized list.

            \item This is the second item of the inner list.
                  \LaTeX\ allows you to nest lists deeper than
                  you really should.
         \end{enumerate}
         This is the rest of the second item of the outer
         list.  It is no more interesting than any other
         part of the item.
   \item This is the third item of the list.
\end{itemize}
You can even display poetry.
\begin{verse}
   There is an environment
    for verse \\             % The \\ command separates lines
   Whose features some poets % within a stanza.
   will curse.

                             % One or more blank lines separate stanzas.

   For instead of making\\
   Them do \emph{all} line breaking, \\
   It allows them to put too many words on a line when they'd rather be
   forced to be terse.
\end{verse}

Mathematical formulas may also be displayed.  A
displayed formula
is
one-line long; multiline
formulas require special formatting instructions.
   \[  \ip{\Gamma}{\psi'} = x'' + y^{2} + z_{i}^{n}\]
Don't start a paragraph with a displayed equation,
nor make one a paragraph by itself.

\end{document}               % End of document.

\section{Bench Huge Section 487}
% synthetic bench section 487
% This is a sample LaTeX input file.  (Version of 12 August 2004.)
%
% A '%' character causes TeX to ignore all remaining text on the line,
% and is used for comments like this one.

\documentclass{article}      % Specifies the document class

                             % The preamble begins here.
\title{An Example Document}  % Declares the document's title.
\author{Leslie Lamport}      % Declares the author's name.
\date{January 21, 1994}      % Deleting this command produces today's date.

\newcommand{\ip}[2]{(#1, #2)}
                             % Defines \ip{arg1}{arg2} to mean
                             % (arg1, arg2).

%\newcommand{\ip}[2]{\langle #1 | #2\rangle}
                             % This is an alternative definition of
                             % \ip that is commented out.

\begin{document}             % End of preamble and beginning of text.

\maketitle                   % Produces the title.

This is an example input file.  Comparing it with
the output it generates can show you how to
produce a simple document of your own.

\section{Ordinary Text}      % Produces section heading.  Lower-level
                             % sections are begun with similar
                             % \subsection and \subsubsection commands.

The ends  of words and sentences are marked
  by   spaces. It  doesn't matter how many
spaces    you type; one is as good as 100.  The
end of   a line counts as a space.

One   or more   blank lines denote the  end
of  a paragraph.

Since any number of consecutive spaces are treated
like a single one, the formatting of the input
file makes no difference to
      \LaTeX,                % The \LaTeX command generates the LaTeX logo.
but it makes a difference to you.  When you use
\LaTeX, making your input file as easy to read
as possible will be a great help as you write
your document and when you change it.  This sample
file shows how you can add comments to your own input
file.

Because printing is different from typewriting,
there are a number of things that you have to do
differently when preparing an input file than if
you were just typing the document directly.
Quotation marks like
       ``this''
have to be handled specially, as do quotes within
quotes:
       ``\,`this'            % \, separates the double and single quote.
        is what I just
        wrote, not  `that'\,''.

Dashes come in three sizes: an
       intra-word
dash, a medium dash for number ranges like
       1--2,
and a punctuation
       dash---like
this.

A sentence-ending space should be larger than the
space between words within a sentence.  You
sometimes have to type special commands in
conjunction with punctuation characters to get
this right, as in the following sentence.
       Gnats, gnus, etc.\ all  % `\ ' makes an inter-word space.
       begin with G\@.         % \@ marks end-of-sentence punctuation.
You should check the spaces after periods when
reading your output to make sure you haven't
forgotten any special cases.  Generating an
ellipsis
       \ldots\               % `\ ' is needed after `\ldots' because TeX
                             % ignores spaces after command names like \ldots
                             % made from \ + letters.
                             %
                             % Note how a `%' character causes TeX to ignore
                             % the end of the input line, so these blank lines
                             % do not start a new paragraph.
                             %
with the right spacing around the periods requires
a special command.

\LaTeX\ interprets some common characters as
commands, so you must type special commands to
generate them.  These characters include the
following:
       \$ \& \% \# \{ and \}.

In printing, text is usually emphasized with an
       \emph{italic}
type style.

\begin{em}
   A long segment of text can also be emphasized
   in this way.  Text within such a segment can be
   given \emph{additional} emphasis.
\end{em}

It is sometimes necessary to prevent \LaTeX\ from
breaking a line where it might otherwise do so.
This may be at a space, as between the ``Mr.''\ and
``Jones'' in
       ``Mr.~Jones'',        % ~ produces an unbreakable interword space.
or within a word---especially when the word is a
symbol like
       \mbox{\emph{itemnum}}
that makes little sense when hyphenated across
lines.

Footnotes\footnote{This is an example of a footnote.}
pose no problem.

\LaTeX\ is good at typesetting mathematical formulas
like
       \( x-3y + z = 7 \)
or
       \( a_{1} > x^{2n} + y^{2n} > x' \)
or
       \( \ip{A}{B} = \sum_{i} a_{i} b_{i} \).
The spaces you type in a formula are
ignored.  Remember that a letter like
       $x$                   % $ ... $  and  \( ... \)  are equivalent
is a formula when it denotes a mathematical
symbol, and it should be typed as one.

\section{Displayed Text}

Text is displayed by indenting it from the left
margin.  Quotations are commonly displayed.  There
are short quotations
\begin{quote}
   This is a short quotation.  It consists of a
   single paragraph of text.  See how it is formatted.
\end{quote}
and longer ones.
\begin{quotation}
   This is a longer quotation.  It consists of two
   paragraphs of text, neither of which are
   particularly interesting.

   This is the second paragraph of the quotation.  It
   is just as dull as the first paragraph.
\end{quotation}
Another frequently-displayed structure is a list.
The following is an example of an \emph{itemized}
list.
\begin{itemize}
   \item This is the first item of an itemized list.
         Each item in the list is marked with a ``tick''.
         You don't have to worry about what kind of tick
         mark is used.

   \item This is the second item of the list.  It
         contains another list nested inside it.  The inner
         list is an \emph{enumerated} list.
         \begin{enumerate}
            \item This is the first item of an enumerated
                  list that is nested within the itemized list.

            \item This is the second item of the inner list.
                  \LaTeX\ allows you to nest lists deeper than
                  you really should.
         \end{enumerate}
         This is the rest of the second item of the outer
         list.  It is no more interesting than any other
         part of the item.
   \item This is the third item of the list.
\end{itemize}
You can even display poetry.
\begin{verse}
   There is an environment
    for verse \\             % The \\ command separates lines
   Whose features some poets % within a stanza.
   will curse.

                             % One or more blank lines separate stanzas.

   For instead of making\\
   Them do \emph{all} line breaking, \\
   It allows them to put too many words on a line when they'd rather be
   forced to be terse.
\end{verse}

Mathematical formulas may also be displayed.  A
displayed formula
is
one-line long; multiline
formulas require special formatting instructions.
   \[  \ip{\Gamma}{\psi'} = x'' + y^{2} + z_{i}^{n}\]
Don't start a paragraph with a displayed equation,
nor make one a paragraph by itself.

\end{document}               % End of document.

\section{Bench Huge Section 488}
% synthetic bench section 488
% This is a sample LaTeX input file.  (Version of 12 August 2004.)
%
% A '%' character causes TeX to ignore all remaining text on the line,
% and is used for comments like this one.

\documentclass{article}      % Specifies the document class

                             % The preamble begins here.
\title{An Example Document}  % Declares the document's title.
\author{Leslie Lamport}      % Declares the author's name.
\date{January 21, 1994}      % Deleting this command produces today's date.

\newcommand{\ip}[2]{(#1, #2)}
                             % Defines \ip{arg1}{arg2} to mean
                             % (arg1, arg2).

%\newcommand{\ip}[2]{\langle #1 | #2\rangle}
                             % This is an alternative definition of
                             % \ip that is commented out.

\begin{document}             % End of preamble and beginning of text.

\maketitle                   % Produces the title.

This is an example input file.  Comparing it with
the output it generates can show you how to
produce a simple document of your own.

\section{Ordinary Text}      % Produces section heading.  Lower-level
                             % sections are begun with similar
                             % \subsection and \subsubsection commands.

The ends  of words and sentences are marked
  by   spaces. It  doesn't matter how many
spaces    you type; one is as good as 100.  The
end of   a line counts as a space.

One   or more   blank lines denote the  end
of  a paragraph.

Since any number of consecutive spaces are treated
like a single one, the formatting of the input
file makes no difference to
      \LaTeX,                % The \LaTeX command generates the LaTeX logo.
but it makes a difference to you.  When you use
\LaTeX, making your input file as easy to read
as possible will be a great help as you write
your document and when you change it.  This sample
file shows how you can add comments to your own input
file.

Because printing is different from typewriting,
there are a number of things that you have to do
differently when preparing an input file than if
you were just typing the document directly.
Quotation marks like
       ``this''
have to be handled specially, as do quotes within
quotes:
       ``\,`this'            % \, separates the double and single quote.
        is what I just
        wrote, not  `that'\,''.

Dashes come in three sizes: an
       intra-word
dash, a medium dash for number ranges like
       1--2,
and a punctuation
       dash---like
this.

A sentence-ending space should be larger than the
space between words within a sentence.  You
sometimes have to type special commands in
conjunction with punctuation characters to get
this right, as in the following sentence.
       Gnats, gnus, etc.\ all  % `\ ' makes an inter-word space.
       begin with G\@.         % \@ marks end-of-sentence punctuation.
You should check the spaces after periods when
reading your output to make sure you haven't
forgotten any special cases.  Generating an
ellipsis
       \ldots\               % `\ ' is needed after `\ldots' because TeX
                             % ignores spaces after command names like \ldots
                             % made from \ + letters.
                             %
                             % Note how a `%' character causes TeX to ignore
                             % the end of the input line, so these blank lines
                             % do not start a new paragraph.
                             %
with the right spacing around the periods requires
a special command.

\LaTeX\ interprets some common characters as
commands, so you must type special commands to
generate them.  These characters include the
following:
       \$ \& \% \# \{ and \}.

In printing, text is usually emphasized with an
       \emph{italic}
type style.

\begin{em}
   A long segment of text can also be emphasized
   in this way.  Text within such a segment can be
   given \emph{additional} emphasis.
\end{em}

It is sometimes necessary to prevent \LaTeX\ from
breaking a line where it might otherwise do so.
This may be at a space, as between the ``Mr.''\ and
``Jones'' in
       ``Mr.~Jones'',        % ~ produces an unbreakable interword space.
or within a word---especially when the word is a
symbol like
       \mbox{\emph{itemnum}}
that makes little sense when hyphenated across
lines.

Footnotes\footnote{This is an example of a footnote.}
pose no problem.

\LaTeX\ is good at typesetting mathematical formulas
like
       \( x-3y + z = 7 \)
or
       \( a_{1} > x^{2n} + y^{2n} > x' \)
or
       \( \ip{A}{B} = \sum_{i} a_{i} b_{i} \).
The spaces you type in a formula are
ignored.  Remember that a letter like
       $x$                   % $ ... $  and  \( ... \)  are equivalent
is a formula when it denotes a mathematical
symbol, and it should be typed as one.

\section{Displayed Text}

Text is displayed by indenting it from the left
margin.  Quotations are commonly displayed.  There
are short quotations
\begin{quote}
   This is a short quotation.  It consists of a
   single paragraph of text.  See how it is formatted.
\end{quote}
and longer ones.
\begin{quotation}
   This is a longer quotation.  It consists of two
   paragraphs of text, neither of which are
   particularly interesting.

   This is the second paragraph of the quotation.  It
   is just as dull as the first paragraph.
\end{quotation}
Another frequently-displayed structure is a list.
The following is an example of an \emph{itemized}
list.
\begin{itemize}
   \item This is the first item of an itemized list.
         Each item in the list is marked with a ``tick''.
         You don't have to worry about what kind of tick
         mark is used.

   \item This is the second item of the list.  It
         contains another list nested inside it.  The inner
         list is an \emph{enumerated} list.
         \begin{enumerate}
            \item This is the first item of an enumerated
                  list that is nested within the itemized list.

            \item This is the second item of the inner list.
                  \LaTeX\ allows you to nest lists deeper than
                  you really should.
         \end{enumerate}
         This is the rest of the second item of the outer
         list.  It is no more interesting than any other
         part of the item.
   \item This is the third item of the list.
\end{itemize}
You can even display poetry.
\begin{verse}
   There is an environment
    for verse \\             % The \\ command separates lines
   Whose features some poets % within a stanza.
   will curse.

                             % One or more blank lines separate stanzas.

   For instead of making\\
   Them do \emph{all} line breaking, \\
   It allows them to put too many words on a line when they'd rather be
   forced to be terse.
\end{verse}

Mathematical formulas may also be displayed.  A
displayed formula
is
one-line long; multiline
formulas require special formatting instructions.
   \[  \ip{\Gamma}{\psi'} = x'' + y^{2} + z_{i}^{n}\]
Don't start a paragraph with a displayed equation,
nor make one a paragraph by itself.

\end{document}               % End of document.

\section{Bench Huge Section 489}
% synthetic bench section 489
% This is a sample LaTeX input file.  (Version of 12 August 2004.)
%
% A '%' character causes TeX to ignore all remaining text on the line,
% and is used for comments like this one.

\documentclass{article}      % Specifies the document class

                             % The preamble begins here.
\title{An Example Document}  % Declares the document's title.
\author{Leslie Lamport}      % Declares the author's name.
\date{January 21, 1994}      % Deleting this command produces today's date.

\newcommand{\ip}[2]{(#1, #2)}
                             % Defines \ip{arg1}{arg2} to mean
                             % (arg1, arg2).

%\newcommand{\ip}[2]{\langle #1 | #2\rangle}
                             % This is an alternative definition of
                             % \ip that is commented out.

\begin{document}             % End of preamble and beginning of text.

\maketitle                   % Produces the title.

This is an example input file.  Comparing it with
the output it generates can show you how to
produce a simple document of your own.

\section{Ordinary Text}      % Produces section heading.  Lower-level
                             % sections are begun with similar
                             % \subsection and \subsubsection commands.

The ends  of words and sentences are marked
  by   spaces. It  doesn't matter how many
spaces    you type; one is as good as 100.  The
end of   a line counts as a space.

One   or more   blank lines denote the  end
of  a paragraph.

Since any number of consecutive spaces are treated
like a single one, the formatting of the input
file makes no difference to
      \LaTeX,                % The \LaTeX command generates the LaTeX logo.
but it makes a difference to you.  When you use
\LaTeX, making your input file as easy to read
as possible will be a great help as you write
your document and when you change it.  This sample
file shows how you can add comments to your own input
file.

Because printing is different from typewriting,
there are a number of things that you have to do
differently when preparing an input file than if
you were just typing the document directly.
Quotation marks like
       ``this''
have to be handled specially, as do quotes within
quotes:
       ``\,`this'            % \, separates the double and single quote.
        is what I just
        wrote, not  `that'\,''.

Dashes come in three sizes: an
       intra-word
dash, a medium dash for number ranges like
       1--2,
and a punctuation
       dash---like
this.

A sentence-ending space should be larger than the
space between words within a sentence.  You
sometimes have to type special commands in
conjunction with punctuation characters to get
this right, as in the following sentence.
       Gnats, gnus, etc.\ all  % `\ ' makes an inter-word space.
       begin with G\@.         % \@ marks end-of-sentence punctuation.
You should check the spaces after periods when
reading your output to make sure you haven't
forgotten any special cases.  Generating an
ellipsis
       \ldots\               % `\ ' is needed after `\ldots' because TeX
                             % ignores spaces after command names like \ldots
                             % made from \ + letters.
                             %
                             % Note how a `%' character causes TeX to ignore
                             % the end of the input line, so these blank lines
                             % do not start a new paragraph.
                             %
with the right spacing around the periods requires
a special command.

\LaTeX\ interprets some common characters as
commands, so you must type special commands to
generate them.  These characters include the
following:
       \$ \& \% \# \{ and \}.

In printing, text is usually emphasized with an
       \emph{italic}
type style.

\begin{em}
   A long segment of text can also be emphasized
   in this way.  Text within such a segment can be
   given \emph{additional} emphasis.
\end{em}

It is sometimes necessary to prevent \LaTeX\ from
breaking a line where it might otherwise do so.
This may be at a space, as between the ``Mr.''\ and
``Jones'' in
       ``Mr.~Jones'',        % ~ produces an unbreakable interword space.
or within a word---especially when the word is a
symbol like
       \mbox{\emph{itemnum}}
that makes little sense when hyphenated across
lines.

Footnotes\footnote{This is an example of a footnote.}
pose no problem.

\LaTeX\ is good at typesetting mathematical formulas
like
       \( x-3y + z = 7 \)
or
       \( a_{1} > x^{2n} + y^{2n} > x' \)
or
       \( \ip{A}{B} = \sum_{i} a_{i} b_{i} \).
The spaces you type in a formula are
ignored.  Remember that a letter like
       $x$                   % $ ... $  and  \( ... \)  are equivalent
is a formula when it denotes a mathematical
symbol, and it should be typed as one.

\section{Displayed Text}

Text is displayed by indenting it from the left
margin.  Quotations are commonly displayed.  There
are short quotations
\begin{quote}
   This is a short quotation.  It consists of a
   single paragraph of text.  See how it is formatted.
\end{quote}
and longer ones.
\begin{quotation}
   This is a longer quotation.  It consists of two
   paragraphs of text, neither of which are
   particularly interesting.

   This is the second paragraph of the quotation.  It
   is just as dull as the first paragraph.
\end{quotation}
Another frequently-displayed structure is a list.
The following is an example of an \emph{itemized}
list.
\begin{itemize}
   \item This is the first item of an itemized list.
         Each item in the list is marked with a ``tick''.
         You don't have to worry about what kind of tick
         mark is used.

   \item This is the second item of the list.  It
         contains another list nested inside it.  The inner
         list is an \emph{enumerated} list.
         \begin{enumerate}
            \item This is the first item of an enumerated
                  list that is nested within the itemized list.

            \item This is the second item of the inner list.
                  \LaTeX\ allows you to nest lists deeper than
                  you really should.
         \end{enumerate}
         This is the rest of the second item of the outer
         list.  It is no more interesting than any other
         part of the item.
   \item This is the third item of the list.
\end{itemize}
You can even display poetry.
\begin{verse}
   There is an environment
    for verse \\             % The \\ command separates lines
   Whose features some poets % within a stanza.
   will curse.

                             % One or more blank lines separate stanzas.

   For instead of making\\
   Them do \emph{all} line breaking, \\
   It allows them to put too many words on a line when they'd rather be
   forced to be terse.
\end{verse}

Mathematical formulas may also be displayed.  A
displayed formula
is
one-line long; multiline
formulas require special formatting instructions.
   \[  \ip{\Gamma}{\psi'} = x'' + y^{2} + z_{i}^{n}\]
Don't start a paragraph with a displayed equation,
nor make one a paragraph by itself.

\end{document}               % End of document.

\section{Bench Huge Section 490}
% synthetic bench section 490
% This is a sample LaTeX input file.  (Version of 12 August 2004.)
%
% A '%' character causes TeX to ignore all remaining text on the line,
% and is used for comments like this one.

\documentclass{article}      % Specifies the document class

                             % The preamble begins here.
\title{An Example Document}  % Declares the document's title.
\author{Leslie Lamport}      % Declares the author's name.
\date{January 21, 1994}      % Deleting this command produces today's date.

\newcommand{\ip}[2]{(#1, #2)}
                             % Defines \ip{arg1}{arg2} to mean
                             % (arg1, arg2).

%\newcommand{\ip}[2]{\langle #1 | #2\rangle}
                             % This is an alternative definition of
                             % \ip that is commented out.

\begin{document}             % End of preamble and beginning of text.

\maketitle                   % Produces the title.

This is an example input file.  Comparing it with
the output it generates can show you how to
produce a simple document of your own.

\section{Ordinary Text}      % Produces section heading.  Lower-level
                             % sections are begun with similar
                             % \subsection and \subsubsection commands.

The ends  of words and sentences are marked
  by   spaces. It  doesn't matter how many
spaces    you type; one is as good as 100.  The
end of   a line counts as a space.

One   or more   blank lines denote the  end
of  a paragraph.

Since any number of consecutive spaces are treated
like a single one, the formatting of the input
file makes no difference to
      \LaTeX,                % The \LaTeX command generates the LaTeX logo.
but it makes a difference to you.  When you use
\LaTeX, making your input file as easy to read
as possible will be a great help as you write
your document and when you change it.  This sample
file shows how you can add comments to your own input
file.

Because printing is different from typewriting,
there are a number of things that you have to do
differently when preparing an input file than if
you were just typing the document directly.
Quotation marks like
       ``this''
have to be handled specially, as do quotes within
quotes:
       ``\,`this'            % \, separates the double and single quote.
        is what I just
        wrote, not  `that'\,''.

Dashes come in three sizes: an
       intra-word
dash, a medium dash for number ranges like
       1--2,
and a punctuation
       dash---like
this.

A sentence-ending space should be larger than the
space between words within a sentence.  You
sometimes have to type special commands in
conjunction with punctuation characters to get
this right, as in the following sentence.
       Gnats, gnus, etc.\ all  % `\ ' makes an inter-word space.
       begin with G\@.         % \@ marks end-of-sentence punctuation.
You should check the spaces after periods when
reading your output to make sure you haven't
forgotten any special cases.  Generating an
ellipsis
       \ldots\               % `\ ' is needed after `\ldots' because TeX
                             % ignores spaces after command names like \ldots
                             % made from \ + letters.
                             %
                             % Note how a `%' character causes TeX to ignore
                             % the end of the input line, so these blank lines
                             % do not start a new paragraph.
                             %
with the right spacing around the periods requires
a special command.

\LaTeX\ interprets some common characters as
commands, so you must type special commands to
generate them.  These characters include the
following:
       \$ \& \% \# \{ and \}.

In printing, text is usually emphasized with an
       \emph{italic}
type style.

\begin{em}
   A long segment of text can also be emphasized
   in this way.  Text within such a segment can be
   given \emph{additional} emphasis.
\end{em}

It is sometimes necessary to prevent \LaTeX\ from
breaking a line where it might otherwise do so.
This may be at a space, as between the ``Mr.''\ and
``Jones'' in
       ``Mr.~Jones'',        % ~ produces an unbreakable interword space.
or within a word---especially when the word is a
symbol like
       \mbox{\emph{itemnum}}
that makes little sense when hyphenated across
lines.

Footnotes\footnote{This is an example of a footnote.}
pose no problem.

\LaTeX\ is good at typesetting mathematical formulas
like
       \( x-3y + z = 7 \)
or
       \( a_{1} > x^{2n} + y^{2n} > x' \)
or
       \( \ip{A}{B} = \sum_{i} a_{i} b_{i} \).
The spaces you type in a formula are
ignored.  Remember that a letter like
       $x$                   % $ ... $  and  \( ... \)  are equivalent
is a formula when it denotes a mathematical
symbol, and it should be typed as one.

\section{Displayed Text}

Text is displayed by indenting it from the left
margin.  Quotations are commonly displayed.  There
are short quotations
\begin{quote}
   This is a short quotation.  It consists of a
   single paragraph of text.  See how it is formatted.
\end{quote}
and longer ones.
\begin{quotation}
   This is a longer quotation.  It consists of two
   paragraphs of text, neither of which are
   particularly interesting.

   This is the second paragraph of the quotation.  It
   is just as dull as the first paragraph.
\end{quotation}
Another frequently-displayed structure is a list.
The following is an example of an \emph{itemized}
list.
\begin{itemize}
   \item This is the first item of an itemized list.
         Each item in the list is marked with a ``tick''.
         You don't have to worry about what kind of tick
         mark is used.

   \item This is the second item of the list.  It
         contains another list nested inside it.  The inner
         list is an \emph{enumerated} list.
         \begin{enumerate}
            \item This is the first item of an enumerated
                  list that is nested within the itemized list.

            \item This is the second item of the inner list.
                  \LaTeX\ allows you to nest lists deeper than
                  you really should.
         \end{enumerate}
         This is the rest of the second item of the outer
         list.  It is no more interesting than any other
         part of the item.
   \item This is the third item of the list.
\end{itemize}
You can even display poetry.
\begin{verse}
   There is an environment
    for verse \\             % The \\ command separates lines
   Whose features some poets % within a stanza.
   will curse.

                             % One or more blank lines separate stanzas.

   For instead of making\\
   Them do \emph{all} line breaking, \\
   It allows them to put too many words on a line when they'd rather be
   forced to be terse.
\end{verse}

Mathematical formulas may also be displayed.  A
displayed formula
is
one-line long; multiline
formulas require special formatting instructions.
   \[  \ip{\Gamma}{\psi'} = x'' + y^{2} + z_{i}^{n}\]
Don't start a paragraph with a displayed equation,
nor make one a paragraph by itself.

\end{document}               % End of document.

\section{Bench Huge Section 491}
% synthetic bench section 491
% This is a sample LaTeX input file.  (Version of 12 August 2004.)
%
% A '%' character causes TeX to ignore all remaining text on the line,
% and is used for comments like this one.

\documentclass{article}      % Specifies the document class

                             % The preamble begins here.
\title{An Example Document}  % Declares the document's title.
\author{Leslie Lamport}      % Declares the author's name.
\date{January 21, 1994}      % Deleting this command produces today's date.

\newcommand{\ip}[2]{(#1, #2)}
                             % Defines \ip{arg1}{arg2} to mean
                             % (arg1, arg2).

%\newcommand{\ip}[2]{\langle #1 | #2\rangle}
                             % This is an alternative definition of
                             % \ip that is commented out.

\begin{document}             % End of preamble and beginning of text.

\maketitle                   % Produces the title.

This is an example input file.  Comparing it with
the output it generates can show you how to
produce a simple document of your own.

\section{Ordinary Text}      % Produces section heading.  Lower-level
                             % sections are begun with similar
                             % \subsection and \subsubsection commands.

The ends  of words and sentences are marked
  by   spaces. It  doesn't matter how many
spaces    you type; one is as good as 100.  The
end of   a line counts as a space.

One   or more   blank lines denote the  end
of  a paragraph.

Since any number of consecutive spaces are treated
like a single one, the formatting of the input
file makes no difference to
      \LaTeX,                % The \LaTeX command generates the LaTeX logo.
but it makes a difference to you.  When you use
\LaTeX, making your input file as easy to read
as possible will be a great help as you write
your document and when you change it.  This sample
file shows how you can add comments to your own input
file.

Because printing is different from typewriting,
there are a number of things that you have to do
differently when preparing an input file than if
you were just typing the document directly.
Quotation marks like
       ``this''
have to be handled specially, as do quotes within
quotes:
       ``\,`this'            % \, separates the double and single quote.
        is what I just
        wrote, not  `that'\,''.

Dashes come in three sizes: an
       intra-word
dash, a medium dash for number ranges like
       1--2,
and a punctuation
       dash---like
this.

A sentence-ending space should be larger than the
space between words within a sentence.  You
sometimes have to type special commands in
conjunction with punctuation characters to get
this right, as in the following sentence.
       Gnats, gnus, etc.\ all  % `\ ' makes an inter-word space.
       begin with G\@.         % \@ marks end-of-sentence punctuation.
You should check the spaces after periods when
reading your output to make sure you haven't
forgotten any special cases.  Generating an
ellipsis
       \ldots\               % `\ ' is needed after `\ldots' because TeX
                             % ignores spaces after command names like \ldots
                             % made from \ + letters.
                             %
                             % Note how a `%' character causes TeX to ignore
                             % the end of the input line, so these blank lines
                             % do not start a new paragraph.
                             %
with the right spacing around the periods requires
a special command.

\LaTeX\ interprets some common characters as
commands, so you must type special commands to
generate them.  These characters include the
following:
       \$ \& \% \# \{ and \}.

In printing, text is usually emphasized with an
       \emph{italic}
type style.

\begin{em}
   A long segment of text can also be emphasized
   in this way.  Text within such a segment can be
   given \emph{additional} emphasis.
\end{em}

It is sometimes necessary to prevent \LaTeX\ from
breaking a line where it might otherwise do so.
This may be at a space, as between the ``Mr.''\ and
``Jones'' in
       ``Mr.~Jones'',        % ~ produces an unbreakable interword space.
or within a word---especially when the word is a
symbol like
       \mbox{\emph{itemnum}}
that makes little sense when hyphenated across
lines.

Footnotes\footnote{This is an example of a footnote.}
pose no problem.

\LaTeX\ is good at typesetting mathematical formulas
like
       \( x-3y + z = 7 \)
or
       \( a_{1} > x^{2n} + y^{2n} > x' \)
or
       \( \ip{A}{B} = \sum_{i} a_{i} b_{i} \).
The spaces you type in a formula are
ignored.  Remember that a letter like
       $x$                   % $ ... $  and  \( ... \)  are equivalent
is a formula when it denotes a mathematical
symbol, and it should be typed as one.

\section{Displayed Text}

Text is displayed by indenting it from the left
margin.  Quotations are commonly displayed.  There
are short quotations
\begin{quote}
   This is a short quotation.  It consists of a
   single paragraph of text.  See how it is formatted.
\end{quote}
and longer ones.
\begin{quotation}
   This is a longer quotation.  It consists of two
   paragraphs of text, neither of which are
   particularly interesting.

   This is the second paragraph of the quotation.  It
   is just as dull as the first paragraph.
\end{quotation}
Another frequently-displayed structure is a list.
The following is an example of an \emph{itemized}
list.
\begin{itemize}
   \item This is the first item of an itemized list.
         Each item in the list is marked with a ``tick''.
         You don't have to worry about what kind of tick
         mark is used.

   \item This is the second item of the list.  It
         contains another list nested inside it.  The inner
         list is an \emph{enumerated} list.
         \begin{enumerate}
            \item This is the first item of an enumerated
                  list that is nested within the itemized list.

            \item This is the second item of the inner list.
                  \LaTeX\ allows you to nest lists deeper than
                  you really should.
         \end{enumerate}
         This is the rest of the second item of the outer
         list.  It is no more interesting than any other
         part of the item.
   \item This is the third item of the list.
\end{itemize}
You can even display poetry.
\begin{verse}
   There is an environment
    for verse \\             % The \\ command separates lines
   Whose features some poets % within a stanza.
   will curse.

                             % One or more blank lines separate stanzas.

   For instead of making\\
   Them do \emph{all} line breaking, \\
   It allows them to put too many words on a line when they'd rather be
   forced to be terse.
\end{verse}

Mathematical formulas may also be displayed.  A
displayed formula
is
one-line long; multiline
formulas require special formatting instructions.
   \[  \ip{\Gamma}{\psi'} = x'' + y^{2} + z_{i}^{n}\]
Don't start a paragraph with a displayed equation,
nor make one a paragraph by itself.

\end{document}               % End of document.

\section{Bench Huge Section 492}
% synthetic bench section 492
% This is a sample LaTeX input file.  (Version of 12 August 2004.)
%
% A '%' character causes TeX to ignore all remaining text on the line,
% and is used for comments like this one.

\documentclass{article}      % Specifies the document class

                             % The preamble begins here.
\title{An Example Document}  % Declares the document's title.
\author{Leslie Lamport}      % Declares the author's name.
\date{January 21, 1994}      % Deleting this command produces today's date.

\newcommand{\ip}[2]{(#1, #2)}
                             % Defines \ip{arg1}{arg2} to mean
                             % (arg1, arg2).

%\newcommand{\ip}[2]{\langle #1 | #2\rangle}
                             % This is an alternative definition of
                             % \ip that is commented out.

\begin{document}             % End of preamble and beginning of text.

\maketitle                   % Produces the title.

This is an example input file.  Comparing it with
the output it generates can show you how to
produce a simple document of your own.

\section{Ordinary Text}      % Produces section heading.  Lower-level
                             % sections are begun with similar
                             % \subsection and \subsubsection commands.

The ends  of words and sentences are marked
  by   spaces. It  doesn't matter how many
spaces    you type; one is as good as 100.  The
end of   a line counts as a space.

One   or more   blank lines denote the  end
of  a paragraph.

Since any number of consecutive spaces are treated
like a single one, the formatting of the input
file makes no difference to
      \LaTeX,                % The \LaTeX command generates the LaTeX logo.
but it makes a difference to you.  When you use
\LaTeX, making your input file as easy to read
as possible will be a great help as you write
your document and when you change it.  This sample
file shows how you can add comments to your own input
file.

Because printing is different from typewriting,
there are a number of things that you have to do
differently when preparing an input file than if
you were just typing the document directly.
Quotation marks like
       ``this''
have to be handled specially, as do quotes within
quotes:
       ``\,`this'            % \, separates the double and single quote.
        is what I just
        wrote, not  `that'\,''.

Dashes come in three sizes: an
       intra-word
dash, a medium dash for number ranges like
       1--2,
and a punctuation
       dash---like
this.

A sentence-ending space should be larger than the
space between words within a sentence.  You
sometimes have to type special commands in
conjunction with punctuation characters to get
this right, as in the following sentence.
       Gnats, gnus, etc.\ all  % `\ ' makes an inter-word space.
       begin with G\@.         % \@ marks end-of-sentence punctuation.
You should check the spaces after periods when
reading your output to make sure you haven't
forgotten any special cases.  Generating an
ellipsis
       \ldots\               % `\ ' is needed after `\ldots' because TeX
                             % ignores spaces after command names like \ldots
                             % made from \ + letters.
                             %
                             % Note how a `%' character causes TeX to ignore
                             % the end of the input line, so these blank lines
                             % do not start a new paragraph.
                             %
with the right spacing around the periods requires
a special command.

\LaTeX\ interprets some common characters as
commands, so you must type special commands to
generate them.  These characters include the
following:
       \$ \& \% \# \{ and \}.

In printing, text is usually emphasized with an
       \emph{italic}
type style.

\begin{em}
   A long segment of text can also be emphasized
   in this way.  Text within such a segment can be
   given \emph{additional} emphasis.
\end{em}

It is sometimes necessary to prevent \LaTeX\ from
breaking a line where it might otherwise do so.
This may be at a space, as between the ``Mr.''\ and
``Jones'' in
       ``Mr.~Jones'',        % ~ produces an unbreakable interword space.
or within a word---especially when the word is a
symbol like
       \mbox{\emph{itemnum}}
that makes little sense when hyphenated across
lines.

Footnotes\footnote{This is an example of a footnote.}
pose no problem.

\LaTeX\ is good at typesetting mathematical formulas
like
       \( x-3y + z = 7 \)
or
       \( a_{1} > x^{2n} + y^{2n} > x' \)
or
       \( \ip{A}{B} = \sum_{i} a_{i} b_{i} \).
The spaces you type in a formula are
ignored.  Remember that a letter like
       $x$                   % $ ... $  and  \( ... \)  are equivalent
is a formula when it denotes a mathematical
symbol, and it should be typed as one.

\section{Displayed Text}

Text is displayed by indenting it from the left
margin.  Quotations are commonly displayed.  There
are short quotations
\begin{quote}
   This is a short quotation.  It consists of a
   single paragraph of text.  See how it is formatted.
\end{quote}
and longer ones.
\begin{quotation}
   This is a longer quotation.  It consists of two
   paragraphs of text, neither of which are
   particularly interesting.

   This is the second paragraph of the quotation.  It
   is just as dull as the first paragraph.
\end{quotation}
Another frequently-displayed structure is a list.
The following is an example of an \emph{itemized}
list.
\begin{itemize}
   \item This is the first item of an itemized list.
         Each item in the list is marked with a ``tick''.
         You don't have to worry about what kind of tick
         mark is used.

   \item This is the second item of the list.  It
         contains another list nested inside it.  The inner
         list is an \emph{enumerated} list.
         \begin{enumerate}
            \item This is the first item of an enumerated
                  list that is nested within the itemized list.

            \item This is the second item of the inner list.
                  \LaTeX\ allows you to nest lists deeper than
                  you really should.
         \end{enumerate}
         This is the rest of the second item of the outer
         list.  It is no more interesting than any other
         part of the item.
   \item This is the third item of the list.
\end{itemize}
You can even display poetry.
\begin{verse}
   There is an environment
    for verse \\             % The \\ command separates lines
   Whose features some poets % within a stanza.
   will curse.

                             % One or more blank lines separate stanzas.

   For instead of making\\
   Them do \emph{all} line breaking, \\
   It allows them to put too many words on a line when they'd rather be
   forced to be terse.
\end{verse}

Mathematical formulas may also be displayed.  A
displayed formula
is
one-line long; multiline
formulas require special formatting instructions.
   \[  \ip{\Gamma}{\psi'} = x'' + y^{2} + z_{i}^{n}\]
Don't start a paragraph with a displayed equation,
nor make one a paragraph by itself.

\end{document}               % End of document.

\section{Bench Huge Section 493}
% synthetic bench section 493
% This is a sample LaTeX input file.  (Version of 12 August 2004.)
%
% A '%' character causes TeX to ignore all remaining text on the line,
% and is used for comments like this one.

\documentclass{article}      % Specifies the document class

                             % The preamble begins here.
\title{An Example Document}  % Declares the document's title.
\author{Leslie Lamport}      % Declares the author's name.
\date{January 21, 1994}      % Deleting this command produces today's date.

\newcommand{\ip}[2]{(#1, #2)}
                             % Defines \ip{arg1}{arg2} to mean
                             % (arg1, arg2).

%\newcommand{\ip}[2]{\langle #1 | #2\rangle}
                             % This is an alternative definition of
                             % \ip that is commented out.

\begin{document}             % End of preamble and beginning of text.

\maketitle                   % Produces the title.

This is an example input file.  Comparing it with
the output it generates can show you how to
produce a simple document of your own.

\section{Ordinary Text}      % Produces section heading.  Lower-level
                             % sections are begun with similar
                             % \subsection and \subsubsection commands.

The ends  of words and sentences are marked
  by   spaces. It  doesn't matter how many
spaces    you type; one is as good as 100.  The
end of   a line counts as a space.

One   or more   blank lines denote the  end
of  a paragraph.

Since any number of consecutive spaces are treated
like a single one, the formatting of the input
file makes no difference to
      \LaTeX,                % The \LaTeX command generates the LaTeX logo.
but it makes a difference to you.  When you use
\LaTeX, making your input file as easy to read
as possible will be a great help as you write
your document and when you change it.  This sample
file shows how you can add comments to your own input
file.

Because printing is different from typewriting,
there are a number of things that you have to do
differently when preparing an input file than if
you were just typing the document directly.
Quotation marks like
       ``this''
have to be handled specially, as do quotes within
quotes:
       ``\,`this'            % \, separates the double and single quote.
        is what I just
        wrote, not  `that'\,''.

Dashes come in three sizes: an
       intra-word
dash, a medium dash for number ranges like
       1--2,
and a punctuation
       dash---like
this.

A sentence-ending space should be larger than the
space between words within a sentence.  You
sometimes have to type special commands in
conjunction with punctuation characters to get
this right, as in the following sentence.
       Gnats, gnus, etc.\ all  % `\ ' makes an inter-word space.
       begin with G\@.         % \@ marks end-of-sentence punctuation.
You should check the spaces after periods when
reading your output to make sure you haven't
forgotten any special cases.  Generating an
ellipsis
       \ldots\               % `\ ' is needed after `\ldots' because TeX
                             % ignores spaces after command names like \ldots
                             % made from \ + letters.
                             %
                             % Note how a `%' character causes TeX to ignore
                             % the end of the input line, so these blank lines
                             % do not start a new paragraph.
                             %
with the right spacing around the periods requires
a special command.

\LaTeX\ interprets some common characters as
commands, so you must type special commands to
generate them.  These characters include the
following:
       \$ \& \% \# \{ and \}.

In printing, text is usually emphasized with an
       \emph{italic}
type style.

\begin{em}
   A long segment of text can also be emphasized
   in this way.  Text within such a segment can be
   given \emph{additional} emphasis.
\end{em}

It is sometimes necessary to prevent \LaTeX\ from
breaking a line where it might otherwise do so.
This may be at a space, as between the ``Mr.''\ and
``Jones'' in
       ``Mr.~Jones'',        % ~ produces an unbreakable interword space.
or within a word---especially when the word is a
symbol like
       \mbox{\emph{itemnum}}
that makes little sense when hyphenated across
lines.

Footnotes\footnote{This is an example of a footnote.}
pose no problem.

\LaTeX\ is good at typesetting mathematical formulas
like
       \( x-3y + z = 7 \)
or
       \( a_{1} > x^{2n} + y^{2n} > x' \)
or
       \( \ip{A}{B} = \sum_{i} a_{i} b_{i} \).
The spaces you type in a formula are
ignored.  Remember that a letter like
       $x$                   % $ ... $  and  \( ... \)  are equivalent
is a formula when it denotes a mathematical
symbol, and it should be typed as one.

\section{Displayed Text}

Text is displayed by indenting it from the left
margin.  Quotations are commonly displayed.  There
are short quotations
\begin{quote}
   This is a short quotation.  It consists of a
   single paragraph of text.  See how it is formatted.
\end{quote}
and longer ones.
\begin{quotation}
   This is a longer quotation.  It consists of two
   paragraphs of text, neither of which are
   particularly interesting.

   This is the second paragraph of the quotation.  It
   is just as dull as the first paragraph.
\end{quotation}
Another frequently-displayed structure is a list.
The following is an example of an \emph{itemized}
list.
\begin{itemize}
   \item This is the first item of an itemized list.
         Each item in the list is marked with a ``tick''.
         You don't have to worry about what kind of tick
         mark is used.

   \item This is the second item of the list.  It
         contains another list nested inside it.  The inner
         list is an \emph{enumerated} list.
         \begin{enumerate}
            \item This is the first item of an enumerated
                  list that is nested within the itemized list.

            \item This is the second item of the inner list.
                  \LaTeX\ allows you to nest lists deeper than
                  you really should.
         \end{enumerate}
         This is the rest of the second item of the outer
         list.  It is no more interesting than any other
         part of the item.
   \item This is the third item of the list.
\end{itemize}
You can even display poetry.
\begin{verse}
   There is an environment
    for verse \\             % The \\ command separates lines
   Whose features some poets % within a stanza.
   will curse.

                             % One or more blank lines separate stanzas.

   For instead of making\\
   Them do \emph{all} line breaking, \\
   It allows them to put too many words on a line when they'd rather be
   forced to be terse.
\end{verse}

Mathematical formulas may also be displayed.  A
displayed formula
is
one-line long; multiline
formulas require special formatting instructions.
   \[  \ip{\Gamma}{\psi'} = x'' + y^{2} + z_{i}^{n}\]
Don't start a paragraph with a displayed equation,
nor make one a paragraph by itself.

\end{document}               % End of document.

\section{Bench Huge Section 494}
% synthetic bench section 494
% This is a sample LaTeX input file.  (Version of 12 August 2004.)
%
% A '%' character causes TeX to ignore all remaining text on the line,
% and is used for comments like this one.

\documentclass{article}      % Specifies the document class

                             % The preamble begins here.
\title{An Example Document}  % Declares the document's title.
\author{Leslie Lamport}      % Declares the author's name.
\date{January 21, 1994}      % Deleting this command produces today's date.

\newcommand{\ip}[2]{(#1, #2)}
                             % Defines \ip{arg1}{arg2} to mean
                             % (arg1, arg2).

%\newcommand{\ip}[2]{\langle #1 | #2\rangle}
                             % This is an alternative definition of
                             % \ip that is commented out.

\begin{document}             % End of preamble and beginning of text.

\maketitle                   % Produces the title.

This is an example input file.  Comparing it with
the output it generates can show you how to
produce a simple document of your own.

\section{Ordinary Text}      % Produces section heading.  Lower-level
                             % sections are begun with similar
                             % \subsection and \subsubsection commands.

The ends  of words and sentences are marked
  by   spaces. It  doesn't matter how many
spaces    you type; one is as good as 100.  The
end of   a line counts as a space.

One   or more   blank lines denote the  end
of  a paragraph.

Since any number of consecutive spaces are treated
like a single one, the formatting of the input
file makes no difference to
      \LaTeX,                % The \LaTeX command generates the LaTeX logo.
but it makes a difference to you.  When you use
\LaTeX, making your input file as easy to read
as possible will be a great help as you write
your document and when you change it.  This sample
file shows how you can add comments to your own input
file.

Because printing is different from typewriting,
there are a number of things that you have to do
differently when preparing an input file than if
you were just typing the document directly.
Quotation marks like
       ``this''
have to be handled specially, as do quotes within
quotes:
       ``\,`this'            % \, separates the double and single quote.
        is what I just
        wrote, not  `that'\,''.

Dashes come in three sizes: an
       intra-word
dash, a medium dash for number ranges like
       1--2,
and a punctuation
       dash---like
this.

A sentence-ending space should be larger than the
space between words within a sentence.  You
sometimes have to type special commands in
conjunction with punctuation characters to get
this right, as in the following sentence.
       Gnats, gnus, etc.\ all  % `\ ' makes an inter-word space.
       begin with G\@.         % \@ marks end-of-sentence punctuation.
You should check the spaces after periods when
reading your output to make sure you haven't
forgotten any special cases.  Generating an
ellipsis
       \ldots\               % `\ ' is needed after `\ldots' because TeX
                             % ignores spaces after command names like \ldots
                             % made from \ + letters.
                             %
                             % Note how a `%' character causes TeX to ignore
                             % the end of the input line, so these blank lines
                             % do not start a new paragraph.
                             %
with the right spacing around the periods requires
a special command.

\LaTeX\ interprets some common characters as
commands, so you must type special commands to
generate them.  These characters include the
following:
       \$ \& \% \# \{ and \}.

In printing, text is usually emphasized with an
       \emph{italic}
type style.

\begin{em}
   A long segment of text can also be emphasized
   in this way.  Text within such a segment can be
   given \emph{additional} emphasis.
\end{em}

It is sometimes necessary to prevent \LaTeX\ from
breaking a line where it might otherwise do so.
This may be at a space, as between the ``Mr.''\ and
``Jones'' in
       ``Mr.~Jones'',        % ~ produces an unbreakable interword space.
or within a word---especially when the word is a
symbol like
       \mbox{\emph{itemnum}}
that makes little sense when hyphenated across
lines.

Footnotes\footnote{This is an example of a footnote.}
pose no problem.

\LaTeX\ is good at typesetting mathematical formulas
like
       \( x-3y + z = 7 \)
or
       \( a_{1} > x^{2n} + y^{2n} > x' \)
or
       \( \ip{A}{B} = \sum_{i} a_{i} b_{i} \).
The spaces you type in a formula are
ignored.  Remember that a letter like
       $x$                   % $ ... $  and  \( ... \)  are equivalent
is a formula when it denotes a mathematical
symbol, and it should be typed as one.

\section{Displayed Text}

Text is displayed by indenting it from the left
margin.  Quotations are commonly displayed.  There
are short quotations
\begin{quote}
   This is a short quotation.  It consists of a
   single paragraph of text.  See how it is formatted.
\end{quote}
and longer ones.
\begin{quotation}
   This is a longer quotation.  It consists of two
   paragraphs of text, neither of which are
   particularly interesting.

   This is the second paragraph of the quotation.  It
   is just as dull as the first paragraph.
\end{quotation}
Another frequently-displayed structure is a list.
The following is an example of an \emph{itemized}
list.
\begin{itemize}
   \item This is the first item of an itemized list.
         Each item in the list is marked with a ``tick''.
         You don't have to worry about what kind of tick
         mark is used.

   \item This is the second item of the list.  It
         contains another list nested inside it.  The inner
         list is an \emph{enumerated} list.
         \begin{enumerate}
            \item This is the first item of an enumerated
                  list that is nested within the itemized list.

            \item This is the second item of the inner list.
                  \LaTeX\ allows you to nest lists deeper than
                  you really should.
         \end{enumerate}
         This is the rest of the second item of the outer
         list.  It is no more interesting than any other
         part of the item.
   \item This is the third item of the list.
\end{itemize}
You can even display poetry.
\begin{verse}
   There is an environment
    for verse \\             % The \\ command separates lines
   Whose features some poets % within a stanza.
   will curse.

                             % One or more blank lines separate stanzas.

   For instead of making\\
   Them do \emph{all} line breaking, \\
   It allows them to put too many words on a line when they'd rather be
   forced to be terse.
\end{verse}

Mathematical formulas may also be displayed.  A
displayed formula
is
one-line long; multiline
formulas require special formatting instructions.
   \[  \ip{\Gamma}{\psi'} = x'' + y^{2} + z_{i}^{n}\]
Don't start a paragraph with a displayed equation,
nor make one a paragraph by itself.

\end{document}               % End of document.

\section{Bench Huge Section 495}
% synthetic bench section 495
% This is a sample LaTeX input file.  (Version of 12 August 2004.)
%
% A '%' character causes TeX to ignore all remaining text on the line,
% and is used for comments like this one.

\documentclass{article}      % Specifies the document class

                             % The preamble begins here.
\title{An Example Document}  % Declares the document's title.
\author{Leslie Lamport}      % Declares the author's name.
\date{January 21, 1994}      % Deleting this command produces today's date.

\newcommand{\ip}[2]{(#1, #2)}
                             % Defines \ip{arg1}{arg2} to mean
                             % (arg1, arg2).

%\newcommand{\ip}[2]{\langle #1 | #2\rangle}
                             % This is an alternative definition of
                             % \ip that is commented out.

\begin{document}             % End of preamble and beginning of text.

\maketitle                   % Produces the title.

This is an example input file.  Comparing it with
the output it generates can show you how to
produce a simple document of your own.

\section{Ordinary Text}      % Produces section heading.  Lower-level
                             % sections are begun with similar
                             % \subsection and \subsubsection commands.

The ends  of words and sentences are marked
  by   spaces. It  doesn't matter how many
spaces    you type; one is as good as 100.  The
end of   a line counts as a space.

One   or more   blank lines denote the  end
of  a paragraph.

Since any number of consecutive spaces are treated
like a single one, the formatting of the input
file makes no difference to
      \LaTeX,                % The \LaTeX command generates the LaTeX logo.
but it makes a difference to you.  When you use
\LaTeX, making your input file as easy to read
as possible will be a great help as you write
your document and when you change it.  This sample
file shows how you can add comments to your own input
file.

Because printing is different from typewriting,
there are a number of things that you have to do
differently when preparing an input file than if
you were just typing the document directly.
Quotation marks like
       ``this''
have to be handled specially, as do quotes within
quotes:
       ``\,`this'            % \, separates the double and single quote.
        is what I just
        wrote, not  `that'\,''.

Dashes come in three sizes: an
       intra-word
dash, a medium dash for number ranges like
       1--2,
and a punctuation
       dash---like
this.

A sentence-ending space should be larger than the
space between words within a sentence.  You
sometimes have to type special commands in
conjunction with punctuation characters to get
this right, as in the following sentence.
       Gnats, gnus, etc.\ all  % `\ ' makes an inter-word space.
       begin with G\@.         % \@ marks end-of-sentence punctuation.
You should check the spaces after periods when
reading your output to make sure you haven't
forgotten any special cases.  Generating an
ellipsis
       \ldots\               % `\ ' is needed after `\ldots' because TeX
                             % ignores spaces after command names like \ldots
                             % made from \ + letters.
                             %
                             % Note how a `%' character causes TeX to ignore
                             % the end of the input line, so these blank lines
                             % do not start a new paragraph.
                             %
with the right spacing around the periods requires
a special command.

\LaTeX\ interprets some common characters as
commands, so you must type special commands to
generate them.  These characters include the
following:
       \$ \& \% \# \{ and \}.

In printing, text is usually emphasized with an
       \emph{italic}
type style.

\begin{em}
   A long segment of text can also be emphasized
   in this way.  Text within such a segment can be
   given \emph{additional} emphasis.
\end{em}

It is sometimes necessary to prevent \LaTeX\ from
breaking a line where it might otherwise do so.
This may be at a space, as between the ``Mr.''\ and
``Jones'' in
       ``Mr.~Jones'',        % ~ produces an unbreakable interword space.
or within a word---especially when the word is a
symbol like
       \mbox{\emph{itemnum}}
that makes little sense when hyphenated across
lines.

Footnotes\footnote{This is an example of a footnote.}
pose no problem.

\LaTeX\ is good at typesetting mathematical formulas
like
       \( x-3y + z = 7 \)
or
       \( a_{1} > x^{2n} + y^{2n} > x' \)
or
       \( \ip{A}{B} = \sum_{i} a_{i} b_{i} \).
The spaces you type in a formula are
ignored.  Remember that a letter like
       $x$                   % $ ... $  and  \( ... \)  are equivalent
is a formula when it denotes a mathematical
symbol, and it should be typed as one.

\section{Displayed Text}

Text is displayed by indenting it from the left
margin.  Quotations are commonly displayed.  There
are short quotations
\begin{quote}
   This is a short quotation.  It consists of a
   single paragraph of text.  See how it is formatted.
\end{quote}
and longer ones.
\begin{quotation}
   This is a longer quotation.  It consists of two
   paragraphs of text, neither of which are
   particularly interesting.

   This is the second paragraph of the quotation.  It
   is just as dull as the first paragraph.
\end{quotation}
Another frequently-displayed structure is a list.
The following is an example of an \emph{itemized}
list.
\begin{itemize}
   \item This is the first item of an itemized list.
         Each item in the list is marked with a ``tick''.
         You don't have to worry about what kind of tick
         mark is used.

   \item This is the second item of the list.  It
         contains another list nested inside it.  The inner
         list is an \emph{enumerated} list.
         \begin{enumerate}
            \item This is the first item of an enumerated
                  list that is nested within the itemized list.

            \item This is the second item of the inner list.
                  \LaTeX\ allows you to nest lists deeper than
                  you really should.
         \end{enumerate}
         This is the rest of the second item of the outer
         list.  It is no more interesting than any other
         part of the item.
   \item This is the third item of the list.
\end{itemize}
You can even display poetry.
\begin{verse}
   There is an environment
    for verse \\             % The \\ command separates lines
   Whose features some poets % within a stanza.
   will curse.

                             % One or more blank lines separate stanzas.

   For instead of making\\
   Them do \emph{all} line breaking, \\
   It allows them to put too many words on a line when they'd rather be
   forced to be terse.
\end{verse}

Mathematical formulas may also be displayed.  A
displayed formula
is
one-line long; multiline
formulas require special formatting instructions.
   \[  \ip{\Gamma}{\psi'} = x'' + y^{2} + z_{i}^{n}\]
Don't start a paragraph with a displayed equation,
nor make one a paragraph by itself.

\end{document}               % End of document.

\section{Bench Huge Section 496}
% synthetic bench section 496
% This is a sample LaTeX input file.  (Version of 12 August 2004.)
%
% A '%' character causes TeX to ignore all remaining text on the line,
% and is used for comments like this one.

\documentclass{article}      % Specifies the document class

                             % The preamble begins here.
\title{An Example Document}  % Declares the document's title.
\author{Leslie Lamport}      % Declares the author's name.
\date{January 21, 1994}      % Deleting this command produces today's date.

\newcommand{\ip}[2]{(#1, #2)}
                             % Defines \ip{arg1}{arg2} to mean
                             % (arg1, arg2).

%\newcommand{\ip}[2]{\langle #1 | #2\rangle}
                             % This is an alternative definition of
                             % \ip that is commented out.

\begin{document}             % End of preamble and beginning of text.

\maketitle                   % Produces the title.

This is an example input file.  Comparing it with
the output it generates can show you how to
produce a simple document of your own.

\section{Ordinary Text}      % Produces section heading.  Lower-level
                             % sections are begun with similar
                             % \subsection and \subsubsection commands.

The ends  of words and sentences are marked
  by   spaces. It  doesn't matter how many
spaces    you type; one is as good as 100.  The
end of   a line counts as a space.

One   or more   blank lines denote the  end
of  a paragraph.

Since any number of consecutive spaces are treated
like a single one, the formatting of the input
file makes no difference to
      \LaTeX,                % The \LaTeX command generates the LaTeX logo.
but it makes a difference to you.  When you use
\LaTeX, making your input file as easy to read
as possible will be a great help as you write
your document and when you change it.  This sample
file shows how you can add comments to your own input
file.

Because printing is different from typewriting,
there are a number of things that you have to do
differently when preparing an input file than if
you were just typing the document directly.
Quotation marks like
       ``this''
have to be handled specially, as do quotes within
quotes:
       ``\,`this'            % \, separates the double and single quote.
        is what I just
        wrote, not  `that'\,''.

Dashes come in three sizes: an
       intra-word
dash, a medium dash for number ranges like
       1--2,
and a punctuation
       dash---like
this.

A sentence-ending space should be larger than the
space between words within a sentence.  You
sometimes have to type special commands in
conjunction with punctuation characters to get
this right, as in the following sentence.
       Gnats, gnus, etc.\ all  % `\ ' makes an inter-word space.
       begin with G\@.         % \@ marks end-of-sentence punctuation.
You should check the spaces after periods when
reading your output to make sure you haven't
forgotten any special cases.  Generating an
ellipsis
       \ldots\               % `\ ' is needed after `\ldots' because TeX
                             % ignores spaces after command names like \ldots
                             % made from \ + letters.
                             %
                             % Note how a `%' character causes TeX to ignore
                             % the end of the input line, so these blank lines
                             % do not start a new paragraph.
                             %
with the right spacing around the periods requires
a special command.

\LaTeX\ interprets some common characters as
commands, so you must type special commands to
generate them.  These characters include the
following:
       \$ \& \% \# \{ and \}.

In printing, text is usually emphasized with an
       \emph{italic}
type style.

\begin{em}
   A long segment of text can also be emphasized
   in this way.  Text within such a segment can be
   given \emph{additional} emphasis.
\end{em}

It is sometimes necessary to prevent \LaTeX\ from
breaking a line where it might otherwise do so.
This may be at a space, as between the ``Mr.''\ and
``Jones'' in
       ``Mr.~Jones'',        % ~ produces an unbreakable interword space.
or within a word---especially when the word is a
symbol like
       \mbox{\emph{itemnum}}
that makes little sense when hyphenated across
lines.

Footnotes\footnote{This is an example of a footnote.}
pose no problem.

\LaTeX\ is good at typesetting mathematical formulas
like
       \( x-3y + z = 7 \)
or
       \( a_{1} > x^{2n} + y^{2n} > x' \)
or
       \( \ip{A}{B} = \sum_{i} a_{i} b_{i} \).
The spaces you type in a formula are
ignored.  Remember that a letter like
       $x$                   % $ ... $  and  \( ... \)  are equivalent
is a formula when it denotes a mathematical
symbol, and it should be typed as one.

\section{Displayed Text}

Text is displayed by indenting it from the left
margin.  Quotations are commonly displayed.  There
are short quotations
\begin{quote}
   This is a short quotation.  It consists of a
   single paragraph of text.  See how it is formatted.
\end{quote}
and longer ones.
\begin{quotation}
   This is a longer quotation.  It consists of two
   paragraphs of text, neither of which are
   particularly interesting.

   This is the second paragraph of the quotation.  It
   is just as dull as the first paragraph.
\end{quotation}
Another frequently-displayed structure is a list.
The following is an example of an \emph{itemized}
list.
\begin{itemize}
   \item This is the first item of an itemized list.
         Each item in the list is marked with a ``tick''.
         You don't have to worry about what kind of tick
         mark is used.

   \item This is the second item of the list.  It
         contains another list nested inside it.  The inner
         list is an \emph{enumerated} list.
         \begin{enumerate}
            \item This is the first item of an enumerated
                  list that is nested within the itemized list.

            \item This is the second item of the inner list.
                  \LaTeX\ allows you to nest lists deeper than
                  you really should.
         \end{enumerate}
         This is the rest of the second item of the outer
         list.  It is no more interesting than any other
         part of the item.
   \item This is the third item of the list.
\end{itemize}
You can even display poetry.
\begin{verse}
   There is an environment
    for verse \\             % The \\ command separates lines
   Whose features some poets % within a stanza.
   will curse.

                             % One or more blank lines separate stanzas.

   For instead of making\\
   Them do \emph{all} line breaking, \\
   It allows them to put too many words on a line when they'd rather be
   forced to be terse.
\end{verse}

Mathematical formulas may also be displayed.  A
displayed formula
is
one-line long; multiline
formulas require special formatting instructions.
   \[  \ip{\Gamma}{\psi'} = x'' + y^{2} + z_{i}^{n}\]
Don't start a paragraph with a displayed equation,
nor make one a paragraph by itself.

\end{document}               % End of document.

\section{Bench Huge Section 497}
% synthetic bench section 497
% This is a sample LaTeX input file.  (Version of 12 August 2004.)
%
% A '%' character causes TeX to ignore all remaining text on the line,
% and is used for comments like this one.

\documentclass{article}      % Specifies the document class

                             % The preamble begins here.
\title{An Example Document}  % Declares the document's title.
\author{Leslie Lamport}      % Declares the author's name.
\date{January 21, 1994}      % Deleting this command produces today's date.

\newcommand{\ip}[2]{(#1, #2)}
                             % Defines \ip{arg1}{arg2} to mean
                             % (arg1, arg2).

%\newcommand{\ip}[2]{\langle #1 | #2\rangle}
                             % This is an alternative definition of
                             % \ip that is commented out.

\begin{document}             % End of preamble and beginning of text.

\maketitle                   % Produces the title.

This is an example input file.  Comparing it with
the output it generates can show you how to
produce a simple document of your own.

\section{Ordinary Text}      % Produces section heading.  Lower-level
                             % sections are begun with similar
                             % \subsection and \subsubsection commands.

The ends  of words and sentences are marked
  by   spaces. It  doesn't matter how many
spaces    you type; one is as good as 100.  The
end of   a line counts as a space.

One   or more   blank lines denote the  end
of  a paragraph.

Since any number of consecutive spaces are treated
like a single one, the formatting of the input
file makes no difference to
      \LaTeX,                % The \LaTeX command generates the LaTeX logo.
but it makes a difference to you.  When you use
\LaTeX, making your input file as easy to read
as possible will be a great help as you write
your document and when you change it.  This sample
file shows how you can add comments to your own input
file.

Because printing is different from typewriting,
there are a number of things that you have to do
differently when preparing an input file than if
you were just typing the document directly.
Quotation marks like
       ``this''
have to be handled specially, as do quotes within
quotes:
       ``\,`this'            % \, separates the double and single quote.
        is what I just
        wrote, not  `that'\,''.

Dashes come in three sizes: an
       intra-word
dash, a medium dash for number ranges like
       1--2,
and a punctuation
       dash---like
this.

A sentence-ending space should be larger than the
space between words within a sentence.  You
sometimes have to type special commands in
conjunction with punctuation characters to get
this right, as in the following sentence.
       Gnats, gnus, etc.\ all  % `\ ' makes an inter-word space.
       begin with G\@.         % \@ marks end-of-sentence punctuation.
You should check the spaces after periods when
reading your output to make sure you haven't
forgotten any special cases.  Generating an
ellipsis
       \ldots\               % `\ ' is needed after `\ldots' because TeX
                             % ignores spaces after command names like \ldots
                             % made from \ + letters.
                             %
                             % Note how a `%' character causes TeX to ignore
                             % the end of the input line, so these blank lines
                             % do not start a new paragraph.
                             %
with the right spacing around the periods requires
a special command.

\LaTeX\ interprets some common characters as
commands, so you must type special commands to
generate them.  These characters include the
following:
       \$ \& \% \# \{ and \}.

In printing, text is usually emphasized with an
       \emph{italic}
type style.

\begin{em}
   A long segment of text can also be emphasized
   in this way.  Text within such a segment can be
   given \emph{additional} emphasis.
\end{em}

It is sometimes necessary to prevent \LaTeX\ from
breaking a line where it might otherwise do so.
This may be at a space, as between the ``Mr.''\ and
``Jones'' in
       ``Mr.~Jones'',        % ~ produces an unbreakable interword space.
or within a word---especially when the word is a
symbol like
       \mbox{\emph{itemnum}}
that makes little sense when hyphenated across
lines.

Footnotes\footnote{This is an example of a footnote.}
pose no problem.

\LaTeX\ is good at typesetting mathematical formulas
like
       \( x-3y + z = 7 \)
or
       \( a_{1} > x^{2n} + y^{2n} > x' \)
or
       \( \ip{A}{B} = \sum_{i} a_{i} b_{i} \).
The spaces you type in a formula are
ignored.  Remember that a letter like
       $x$                   % $ ... $  and  \( ... \)  are equivalent
is a formula when it denotes a mathematical
symbol, and it should be typed as one.

\section{Displayed Text}

Text is displayed by indenting it from the left
margin.  Quotations are commonly displayed.  There
are short quotations
\begin{quote}
   This is a short quotation.  It consists of a
   single paragraph of text.  See how it is formatted.
\end{quote}
and longer ones.
\begin{quotation}
   This is a longer quotation.  It consists of two
   paragraphs of text, neither of which are
   particularly interesting.

   This is the second paragraph of the quotation.  It
   is just as dull as the first paragraph.
\end{quotation}
Another frequently-displayed structure is a list.
The following is an example of an \emph{itemized}
list.
\begin{itemize}
   \item This is the first item of an itemized list.
         Each item in the list is marked with a ``tick''.
         You don't have to worry about what kind of tick
         mark is used.

   \item This is the second item of the list.  It
         contains another list nested inside it.  The inner
         list is an \emph{enumerated} list.
         \begin{enumerate}
            \item This is the first item of an enumerated
                  list that is nested within the itemized list.

            \item This is the second item of the inner list.
                  \LaTeX\ allows you to nest lists deeper than
                  you really should.
         \end{enumerate}
         This is the rest of the second item of the outer
         list.  It is no more interesting than any other
         part of the item.
   \item This is the third item of the list.
\end{itemize}
You can even display poetry.
\begin{verse}
   There is an environment
    for verse \\             % The \\ command separates lines
   Whose features some poets % within a stanza.
   will curse.

                             % One or more blank lines separate stanzas.

   For instead of making\\
   Them do \emph{all} line breaking, \\
   It allows them to put too many words on a line when they'd rather be
   forced to be terse.
\end{verse}

Mathematical formulas may also be displayed.  A
displayed formula
is
one-line long; multiline
formulas require special formatting instructions.
   \[  \ip{\Gamma}{\psi'} = x'' + y^{2} + z_{i}^{n}\]
Don't start a paragraph with a displayed equation,
nor make one a paragraph by itself.

\end{document}               % End of document.

\section{Bench Huge Section 498}
% synthetic bench section 498
% This is a sample LaTeX input file.  (Version of 12 August 2004.)
%
% A '%' character causes TeX to ignore all remaining text on the line,
% and is used for comments like this one.

\documentclass{article}      % Specifies the document class

                             % The preamble begins here.
\title{An Example Document}  % Declares the document's title.
\author{Leslie Lamport}      % Declares the author's name.
\date{January 21, 1994}      % Deleting this command produces today's date.

\newcommand{\ip}[2]{(#1, #2)}
                             % Defines \ip{arg1}{arg2} to mean
                             % (arg1, arg2).

%\newcommand{\ip}[2]{\langle #1 | #2\rangle}
                             % This is an alternative definition of
                             % \ip that is commented out.

\begin{document}             % End of preamble and beginning of text.

\maketitle                   % Produces the title.

This is an example input file.  Comparing it with
the output it generates can show you how to
produce a simple document of your own.

\section{Ordinary Text}      % Produces section heading.  Lower-level
                             % sections are begun with similar
                             % \subsection and \subsubsection commands.

The ends  of words and sentences are marked
  by   spaces. It  doesn't matter how many
spaces    you type; one is as good as 100.  The
end of   a line counts as a space.

One   or more   blank lines denote the  end
of  a paragraph.

Since any number of consecutive spaces are treated
like a single one, the formatting of the input
file makes no difference to
      \LaTeX,                % The \LaTeX command generates the LaTeX logo.
but it makes a difference to you.  When you use
\LaTeX, making your input file as easy to read
as possible will be a great help as you write
your document and when you change it.  This sample
file shows how you can add comments to your own input
file.

Because printing is different from typewriting,
there are a number of things that you have to do
differently when preparing an input file than if
you were just typing the document directly.
Quotation marks like
       ``this''
have to be handled specially, as do quotes within
quotes:
       ``\,`this'            % \, separates the double and single quote.
        is what I just
        wrote, not  `that'\,''.

Dashes come in three sizes: an
       intra-word
dash, a medium dash for number ranges like
       1--2,
and a punctuation
       dash---like
this.

A sentence-ending space should be larger than the
space between words within a sentence.  You
sometimes have to type special commands in
conjunction with punctuation characters to get
this right, as in the following sentence.
       Gnats, gnus, etc.\ all  % `\ ' makes an inter-word space.
       begin with G\@.         % \@ marks end-of-sentence punctuation.
You should check the spaces after periods when
reading your output to make sure you haven't
forgotten any special cases.  Generating an
ellipsis
       \ldots\               % `\ ' is needed after `\ldots' because TeX
                             % ignores spaces after command names like \ldots
                             % made from \ + letters.
                             %
                             % Note how a `%' character causes TeX to ignore
                             % the end of the input line, so these blank lines
                             % do not start a new paragraph.
                             %
with the right spacing around the periods requires
a special command.

\LaTeX\ interprets some common characters as
commands, so you must type special commands to
generate them.  These characters include the
following:
       \$ \& \% \# \{ and \}.

In printing, text is usually emphasized with an
       \emph{italic}
type style.

\begin{em}
   A long segment of text can also be emphasized
   in this way.  Text within such a segment can be
   given \emph{additional} emphasis.
\end{em}

It is sometimes necessary to prevent \LaTeX\ from
breaking a line where it might otherwise do so.
This may be at a space, as between the ``Mr.''\ and
``Jones'' in
       ``Mr.~Jones'',        % ~ produces an unbreakable interword space.
or within a word---especially when the word is a
symbol like
       \mbox{\emph{itemnum}}
that makes little sense when hyphenated across
lines.

Footnotes\footnote{This is an example of a footnote.}
pose no problem.

\LaTeX\ is good at typesetting mathematical formulas
like
       \( x-3y + z = 7 \)
or
       \( a_{1} > x^{2n} + y^{2n} > x' \)
or
       \( \ip{A}{B} = \sum_{i} a_{i} b_{i} \).
The spaces you type in a formula are
ignored.  Remember that a letter like
       $x$                   % $ ... $  and  \( ... \)  are equivalent
is a formula when it denotes a mathematical
symbol, and it should be typed as one.

\section{Displayed Text}

Text is displayed by indenting it from the left
margin.  Quotations are commonly displayed.  There
are short quotations
\begin{quote}
   This is a short quotation.  It consists of a
   single paragraph of text.  See how it is formatted.
\end{quote}
and longer ones.
\begin{quotation}
   This is a longer quotation.  It consists of two
   paragraphs of text, neither of which are
   particularly interesting.

   This is the second paragraph of the quotation.  It
   is just as dull as the first paragraph.
\end{quotation}
Another frequently-displayed structure is a list.
The following is an example of an \emph{itemized}
list.
\begin{itemize}
   \item This is the first item of an itemized list.
         Each item in the list is marked with a ``tick''.
         You don't have to worry about what kind of tick
         mark is used.

   \item This is the second item of the list.  It
         contains another list nested inside it.  The inner
         list is an \emph{enumerated} list.
         \begin{enumerate}
            \item This is the first item of an enumerated
                  list that is nested within the itemized list.

            \item This is the second item of the inner list.
                  \LaTeX\ allows you to nest lists deeper than
                  you really should.
         \end{enumerate}
         This is the rest of the second item of the outer
         list.  It is no more interesting than any other
         part of the item.
   \item This is the third item of the list.
\end{itemize}
You can even display poetry.
\begin{verse}
   There is an environment
    for verse \\             % The \\ command separates lines
   Whose features some poets % within a stanza.
   will curse.

                             % One or more blank lines separate stanzas.

   For instead of making\\
   Them do \emph{all} line breaking, \\
   It allows them to put too many words on a line when they'd rather be
   forced to be terse.
\end{verse}

Mathematical formulas may also be displayed.  A
displayed formula
is
one-line long; multiline
formulas require special formatting instructions.
   \[  \ip{\Gamma}{\psi'} = x'' + y^{2} + z_{i}^{n}\]
Don't start a paragraph with a displayed equation,
nor make one a paragraph by itself.

\end{document}               % End of document.

\section{Bench Huge Section 499}
% synthetic bench section 499
% This is a sample LaTeX input file.  (Version of 12 August 2004.)
%
% A '%' character causes TeX to ignore all remaining text on the line,
% and is used for comments like this one.

\documentclass{article}      % Specifies the document class

                             % The preamble begins here.
\title{An Example Document}  % Declares the document's title.
\author{Leslie Lamport}      % Declares the author's name.
\date{January 21, 1994}      % Deleting this command produces today's date.

\newcommand{\ip}[2]{(#1, #2)}
                             % Defines \ip{arg1}{arg2} to mean
                             % (arg1, arg2).

%\newcommand{\ip}[2]{\langle #1 | #2\rangle}
                             % This is an alternative definition of
                             % \ip that is commented out.

\begin{document}             % End of preamble and beginning of text.

\maketitle                   % Produces the title.

This is an example input file.  Comparing it with
the output it generates can show you how to
produce a simple document of your own.

\section{Ordinary Text}      % Produces section heading.  Lower-level
                             % sections are begun with similar
                             % \subsection and \subsubsection commands.

The ends  of words and sentences are marked
  by   spaces. It  doesn't matter how many
spaces    you type; one is as good as 100.  The
end of   a line counts as a space.

One   or more   blank lines denote the  end
of  a paragraph.

Since any number of consecutive spaces are treated
like a single one, the formatting of the input
file makes no difference to
      \LaTeX,                % The \LaTeX command generates the LaTeX logo.
but it makes a difference to you.  When you use
\LaTeX, making your input file as easy to read
as possible will be a great help as you write
your document and when you change it.  This sample
file shows how you can add comments to your own input
file.

Because printing is different from typewriting,
there are a number of things that you have to do
differently when preparing an input file than if
you were just typing the document directly.
Quotation marks like
       ``this''
have to be handled specially, as do quotes within
quotes:
       ``\,`this'            % \, separates the double and single quote.
        is what I just
        wrote, not  `that'\,''.

Dashes come in three sizes: an
       intra-word
dash, a medium dash for number ranges like
       1--2,
and a punctuation
       dash---like
this.

A sentence-ending space should be larger than the
space between words within a sentence.  You
sometimes have to type special commands in
conjunction with punctuation characters to get
this right, as in the following sentence.
       Gnats, gnus, etc.\ all  % `\ ' makes an inter-word space.
       begin with G\@.         % \@ marks end-of-sentence punctuation.
You should check the spaces after periods when
reading your output to make sure you haven't
forgotten any special cases.  Generating an
ellipsis
       \ldots\               % `\ ' is needed after `\ldots' because TeX
                             % ignores spaces after command names like \ldots
                             % made from \ + letters.
                             %
                             % Note how a `%' character causes TeX to ignore
                             % the end of the input line, so these blank lines
                             % do not start a new paragraph.
                             %
with the right spacing around the periods requires
a special command.

\LaTeX\ interprets some common characters as
commands, so you must type special commands to
generate them.  These characters include the
following:
       \$ \& \% \# \{ and \}.

In printing, text is usually emphasized with an
       \emph{italic}
type style.

\begin{em}
   A long segment of text can also be emphasized
   in this way.  Text within such a segment can be
   given \emph{additional} emphasis.
\end{em}

It is sometimes necessary to prevent \LaTeX\ from
breaking a line where it might otherwise do so.
This may be at a space, as between the ``Mr.''\ and
``Jones'' in
       ``Mr.~Jones'',        % ~ produces an unbreakable interword space.
or within a word---especially when the word is a
symbol like
       \mbox{\emph{itemnum}}
that makes little sense when hyphenated across
lines.

Footnotes\footnote{This is an example of a footnote.}
pose no problem.

\LaTeX\ is good at typesetting mathematical formulas
like
       \( x-3y + z = 7 \)
or
       \( a_{1} > x^{2n} + y^{2n} > x' \)
or
       \( \ip{A}{B} = \sum_{i} a_{i} b_{i} \).
The spaces you type in a formula are
ignored.  Remember that a letter like
       $x$                   % $ ... $  and  \( ... \)  are equivalent
is a formula when it denotes a mathematical
symbol, and it should be typed as one.

\section{Displayed Text}

Text is displayed by indenting it from the left
margin.  Quotations are commonly displayed.  There
are short quotations
\begin{quote}
   This is a short quotation.  It consists of a
   single paragraph of text.  See how it is formatted.
\end{quote}
and longer ones.
\begin{quotation}
   This is a longer quotation.  It consists of two
   paragraphs of text, neither of which are
   particularly interesting.

   This is the second paragraph of the quotation.  It
   is just as dull as the first paragraph.
\end{quotation}
Another frequently-displayed structure is a list.
The following is an example of an \emph{itemized}
list.
\begin{itemize}
   \item This is the first item of an itemized list.
         Each item in the list is marked with a ``tick''.
         You don't have to worry about what kind of tick
         mark is used.

   \item This is the second item of the list.  It
         contains another list nested inside it.  The inner
         list is an \emph{enumerated} list.
         \begin{enumerate}
            \item This is the first item of an enumerated
                  list that is nested within the itemized list.

            \item This is the second item of the inner list.
                  \LaTeX\ allows you to nest lists deeper than
                  you really should.
         \end{enumerate}
         This is the rest of the second item of the outer
         list.  It is no more interesting than any other
         part of the item.
   \item This is the third item of the list.
\end{itemize}
You can even display poetry.
\begin{verse}
   There is an environment
    for verse \\             % The \\ command separates lines
   Whose features some poets % within a stanza.
   will curse.

                             % One or more blank lines separate stanzas.

   For instead of making\\
   Them do \emph{all} line breaking, \\
   It allows them to put too many words on a line when they'd rather be
   forced to be terse.
\end{verse}

Mathematical formulas may also be displayed.  A
displayed formula
is
one-line long; multiline
formulas require special formatting instructions.
   \[  \ip{\Gamma}{\psi'} = x'' + y^{2} + z_{i}^{n}\]
Don't start a paragraph with a displayed equation,
nor make one a paragraph by itself.

\end{document}               % End of document.

\section{Bench Huge Section 500}
% synthetic bench section 500
% This is a sample LaTeX input file.  (Version of 12 August 2004.)
%
% A '%' character causes TeX to ignore all remaining text on the line,
% and is used for comments like this one.

\documentclass{article}      % Specifies the document class

                             % The preamble begins here.
\title{An Example Document}  % Declares the document's title.
\author{Leslie Lamport}      % Declares the author's name.
\date{January 21, 1994}      % Deleting this command produces today's date.

\newcommand{\ip}[2]{(#1, #2)}
                             % Defines \ip{arg1}{arg2} to mean
                             % (arg1, arg2).

%\newcommand{\ip}[2]{\langle #1 | #2\rangle}
                             % This is an alternative definition of
                             % \ip that is commented out.

\begin{document}             % End of preamble and beginning of text.

\maketitle                   % Produces the title.

This is an example input file.  Comparing it with
the output it generates can show you how to
produce a simple document of your own.

\section{Ordinary Text}      % Produces section heading.  Lower-level
                             % sections are begun with similar
                             % \subsection and \subsubsection commands.

The ends  of words and sentences are marked
  by   spaces. It  doesn't matter how many
spaces    you type; one is as good as 100.  The
end of   a line counts as a space.

One   or more   blank lines denote the  end
of  a paragraph.

Since any number of consecutive spaces are treated
like a single one, the formatting of the input
file makes no difference to
      \LaTeX,                % The \LaTeX command generates the LaTeX logo.
but it makes a difference to you.  When you use
\LaTeX, making your input file as easy to read
as possible will be a great help as you write
your document and when you change it.  This sample
file shows how you can add comments to your own input
file.

Because printing is different from typewriting,
there are a number of things that you have to do
differently when preparing an input file than if
you were just typing the document directly.
Quotation marks like
       ``this''
have to be handled specially, as do quotes within
quotes:
       ``\,`this'            % \, separates the double and single quote.
        is what I just
        wrote, not  `that'\,''.

Dashes come in three sizes: an
       intra-word
dash, a medium dash for number ranges like
       1--2,
and a punctuation
       dash---like
this.

A sentence-ending space should be larger than the
space between words within a sentence.  You
sometimes have to type special commands in
conjunction with punctuation characters to get
this right, as in the following sentence.
       Gnats, gnus, etc.\ all  % `\ ' makes an inter-word space.
       begin with G\@.         % \@ marks end-of-sentence punctuation.
You should check the spaces after periods when
reading your output to make sure you haven't
forgotten any special cases.  Generating an
ellipsis
       \ldots\               % `\ ' is needed after `\ldots' because TeX
                             % ignores spaces after command names like \ldots
                             % made from \ + letters.
                             %
                             % Note how a `%' character causes TeX to ignore
                             % the end of the input line, so these blank lines
                             % do not start a new paragraph.
                             %
with the right spacing around the periods requires
a special command.

\LaTeX\ interprets some common characters as
commands, so you must type special commands to
generate them.  These characters include the
following:
       \$ \& \% \# \{ and \}.

In printing, text is usually emphasized with an
       \emph{italic}
type style.

\begin{em}
   A long segment of text can also be emphasized
   in this way.  Text within such a segment can be
   given \emph{additional} emphasis.
\end{em}

It is sometimes necessary to prevent \LaTeX\ from
breaking a line where it might otherwise do so.
This may be at a space, as between the ``Mr.''\ and
``Jones'' in
       ``Mr.~Jones'',        % ~ produces an unbreakable interword space.
or within a word---especially when the word is a
symbol like
       \mbox{\emph{itemnum}}
that makes little sense when hyphenated across
lines.

Footnotes\footnote{This is an example of a footnote.}
pose no problem.

\LaTeX\ is good at typesetting mathematical formulas
like
       \( x-3y + z = 7 \)
or
       \( a_{1} > x^{2n} + y^{2n} > x' \)
or
       \( \ip{A}{B} = \sum_{i} a_{i} b_{i} \).
The spaces you type in a formula are
ignored.  Remember that a letter like
       $x$                   % $ ... $  and  \( ... \)  are equivalent
is a formula when it denotes a mathematical
symbol, and it should be typed as one.

\section{Displayed Text}

Text is displayed by indenting it from the left
margin.  Quotations are commonly displayed.  There
are short quotations
\begin{quote}
   This is a short quotation.  It consists of a
   single paragraph of text.  See how it is formatted.
\end{quote}
and longer ones.
\begin{quotation}
   This is a longer quotation.  It consists of two
   paragraphs of text, neither of which are
   particularly interesting.

   This is the second paragraph of the quotation.  It
   is just as dull as the first paragraph.
\end{quotation}
Another frequently-displayed structure is a list.
The following is an example of an \emph{itemized}
list.
\begin{itemize}
   \item This is the first item of an itemized list.
         Each item in the list is marked with a ``tick''.
         You don't have to worry about what kind of tick
         mark is used.

   \item This is the second item of the list.  It
         contains another list nested inside it.  The inner
         list is an \emph{enumerated} list.
         \begin{enumerate}
            \item This is the first item of an enumerated
                  list that is nested within the itemized list.

            \item This is the second item of the inner list.
                  \LaTeX\ allows you to nest lists deeper than
                  you really should.
         \end{enumerate}
         This is the rest of the second item of the outer
         list.  It is no more interesting than any other
         part of the item.
   \item This is the third item of the list.
\end{itemize}
You can even display poetry.
\begin{verse}
   There is an environment
    for verse \\             % The \\ command separates lines
   Whose features some poets % within a stanza.
   will curse.

                             % One or more blank lines separate stanzas.

   For instead of making\\
   Them do \emph{all} line breaking, \\
   It allows them to put too many words on a line when they'd rather be
   forced to be terse.
\end{verse}

Mathematical formulas may also be displayed.  A
displayed formula
is
one-line long; multiline
formulas require special formatting instructions.
   \[  \ip{\Gamma}{\psi'} = x'' + y^{2} + z_{i}^{n}\]
Don't start a paragraph with a displayed equation,
nor make one a paragraph by itself.

\end{document}               % End of document.

\section{Bench Huge Section 501}
% synthetic bench section 501
% This is a sample LaTeX input file.  (Version of 12 August 2004.)
%
% A '%' character causes TeX to ignore all remaining text on the line,
% and is used for comments like this one.

\documentclass{article}      % Specifies the document class

                             % The preamble begins here.
\title{An Example Document}  % Declares the document's title.
\author{Leslie Lamport}      % Declares the author's name.
\date{January 21, 1994}      % Deleting this command produces today's date.

\newcommand{\ip}[2]{(#1, #2)}
                             % Defines \ip{arg1}{arg2} to mean
                             % (arg1, arg2).

%\newcommand{\ip}[2]{\langle #1 | #2\rangle}
                             % This is an alternative definition of
                             % \ip that is commented out.

\begin{document}             % End of preamble and beginning of text.

\maketitle                   % Produces the title.

This is an example input file.  Comparing it with
the output it generates can show you how to
produce a simple document of your own.

\section{Ordinary Text}      % Produces section heading.  Lower-level
                             % sections are begun with similar
                             % \subsection and \subsubsection commands.

The ends  of words and sentences are marked
  by   spaces. It  doesn't matter how many
spaces    you type; one is as good as 100.  The
end of   a line counts as a space.

One   or more   blank lines denote the  end
of  a paragraph.

Since any number of consecutive spaces are treated
like a single one, the formatting of the input
file makes no difference to
      \LaTeX,                % The \LaTeX command generates the LaTeX logo.
but it makes a difference to you.  When you use
\LaTeX, making your input file as easy to read
as possible will be a great help as you write
your document and when you change it.  This sample
file shows how you can add comments to your own input
file.

Because printing is different from typewriting,
there are a number of things that you have to do
differently when preparing an input file than if
you were just typing the document directly.
Quotation marks like
       ``this''
have to be handled specially, as do quotes within
quotes:
       ``\,`this'            % \, separates the double and single quote.
        is what I just
        wrote, not  `that'\,''.

Dashes come in three sizes: an
       intra-word
dash, a medium dash for number ranges like
       1--2,
and a punctuation
       dash---like
this.

A sentence-ending space should be larger than the
space between words within a sentence.  You
sometimes have to type special commands in
conjunction with punctuation characters to get
this right, as in the following sentence.
       Gnats, gnus, etc.\ all  % `\ ' makes an inter-word space.
       begin with G\@.         % \@ marks end-of-sentence punctuation.
You should check the spaces after periods when
reading your output to make sure you haven't
forgotten any special cases.  Generating an
ellipsis
       \ldots\               % `\ ' is needed after `\ldots' because TeX
                             % ignores spaces after command names like \ldots
                             % made from \ + letters.
                             %
                             % Note how a `%' character causes TeX to ignore
                             % the end of the input line, so these blank lines
                             % do not start a new paragraph.
                             %
with the right spacing around the periods requires
a special command.

\LaTeX\ interprets some common characters as
commands, so you must type special commands to
generate them.  These characters include the
following:
       \$ \& \% \# \{ and \}.

In printing, text is usually emphasized with an
       \emph{italic}
type style.

\begin{em}
   A long segment of text can also be emphasized
   in this way.  Text within such a segment can be
   given \emph{additional} emphasis.
\end{em}

It is sometimes necessary to prevent \LaTeX\ from
breaking a line where it might otherwise do so.
This may be at a space, as between the ``Mr.''\ and
``Jones'' in
       ``Mr.~Jones'',        % ~ produces an unbreakable interword space.
or within a word---especially when the word is a
symbol like
       \mbox{\emph{itemnum}}
that makes little sense when hyphenated across
lines.

Footnotes\footnote{This is an example of a footnote.}
pose no problem.

\LaTeX\ is good at typesetting mathematical formulas
like
       \( x-3y + z = 7 \)
or
       \( a_{1} > x^{2n} + y^{2n} > x' \)
or
       \( \ip{A}{B} = \sum_{i} a_{i} b_{i} \).
The spaces you type in a formula are
ignored.  Remember that a letter like
       $x$                   % $ ... $  and  \( ... \)  are equivalent
is a formula when it denotes a mathematical
symbol, and it should be typed as one.

\section{Displayed Text}

Text is displayed by indenting it from the left
margin.  Quotations are commonly displayed.  There
are short quotations
\begin{quote}
   This is a short quotation.  It consists of a
   single paragraph of text.  See how it is formatted.
\end{quote}
and longer ones.
\begin{quotation}
   This is a longer quotation.  It consists of two
   paragraphs of text, neither of which are
   particularly interesting.

   This is the second paragraph of the quotation.  It
   is just as dull as the first paragraph.
\end{quotation}
Another frequently-displayed structure is a list.
The following is an example of an \emph{itemized}
list.
\begin{itemize}
   \item This is the first item of an itemized list.
         Each item in the list is marked with a ``tick''.
         You don't have to worry about what kind of tick
         mark is used.

   \item This is the second item of the list.  It
         contains another list nested inside it.  The inner
         list is an \emph{enumerated} list.
         \begin{enumerate}
            \item This is the first item of an enumerated
                  list that is nested within the itemized list.

            \item This is the second item of the inner list.
                  \LaTeX\ allows you to nest lists deeper than
                  you really should.
         \end{enumerate}
         This is the rest of the second item of the outer
         list.  It is no more interesting than any other
         part of the item.
   \item This is the third item of the list.
\end{itemize}
You can even display poetry.
\begin{verse}
   There is an environment
    for verse \\             % The \\ command separates lines
   Whose features some poets % within a stanza.
   will curse.

                             % One or more blank lines separate stanzas.

   For instead of making\\
   Them do \emph{all} line breaking, \\
   It allows them to put too many words on a line when they'd rather be
   forced to be terse.
\end{verse}

Mathematical formulas may also be displayed.  A
displayed formula
is
one-line long; multiline
formulas require special formatting instructions.
   \[  \ip{\Gamma}{\psi'} = x'' + y^{2} + z_{i}^{n}\]
Don't start a paragraph with a displayed equation,
nor make one a paragraph by itself.

\end{document}               % End of document.

\section{Bench Huge Section 502}
% synthetic bench section 502
% This is a sample LaTeX input file.  (Version of 12 August 2004.)
%
% A '%' character causes TeX to ignore all remaining text on the line,
% and is used for comments like this one.

\documentclass{article}      % Specifies the document class

                             % The preamble begins here.
\title{An Example Document}  % Declares the document's title.
\author{Leslie Lamport}      % Declares the author's name.
\date{January 21, 1994}      % Deleting this command produces today's date.

\newcommand{\ip}[2]{(#1, #2)}
                             % Defines \ip{arg1}{arg2} to mean
                             % (arg1, arg2).

%\newcommand{\ip}[2]{\langle #1 | #2\rangle}
                             % This is an alternative definition of
                             % \ip that is commented out.

\begin{document}             % End of preamble and beginning of text.

\maketitle                   % Produces the title.

This is an example input file.  Comparing it with
the output it generates can show you how to
produce a simple document of your own.

\section{Ordinary Text}      % Produces section heading.  Lower-level
                             % sections are begun with similar
                             % \subsection and \subsubsection commands.

The ends  of words and sentences are marked
  by   spaces. It  doesn't matter how many
spaces    you type; one is as good as 100.  The
end of   a line counts as a space.

One   or more   blank lines denote the  end
of  a paragraph.

Since any number of consecutive spaces are treated
like a single one, the formatting of the input
file makes no difference to
      \LaTeX,                % The \LaTeX command generates the LaTeX logo.
but it makes a difference to you.  When you use
\LaTeX, making your input file as easy to read
as possible will be a great help as you write
your document and when you change it.  This sample
file shows how you can add comments to your own input
file.

Because printing is different from typewriting,
there are a number of things that you have to do
differently when preparing an input file than if
you were just typing the document directly.
Quotation marks like
       ``this''
have to be handled specially, as do quotes within
quotes:
       ``\,`this'            % \, separates the double and single quote.
        is what I just
        wrote, not  `that'\,''.

Dashes come in three sizes: an
       intra-word
dash, a medium dash for number ranges like
       1--2,
and a punctuation
       dash---like
this.

A sentence-ending space should be larger than the
space between words within a sentence.  You
sometimes have to type special commands in
conjunction with punctuation characters to get
this right, as in the following sentence.
       Gnats, gnus, etc.\ all  % `\ ' makes an inter-word space.
       begin with G\@.         % \@ marks end-of-sentence punctuation.
You should check the spaces after periods when
reading your output to make sure you haven't
forgotten any special cases.  Generating an
ellipsis
       \ldots\               % `\ ' is needed after `\ldots' because TeX
                             % ignores spaces after command names like \ldots
                             % made from \ + letters.
                             %
                             % Note how a `%' character causes TeX to ignore
                             % the end of the input line, so these blank lines
                             % do not start a new paragraph.
                             %
with the right spacing around the periods requires
a special command.

\LaTeX\ interprets some common characters as
commands, so you must type special commands to
generate them.  These characters include the
following:
       \$ \& \% \# \{ and \}.

In printing, text is usually emphasized with an
       \emph{italic}
type style.

\begin{em}
   A long segment of text can also be emphasized
   in this way.  Text within such a segment can be
   given \emph{additional} emphasis.
\end{em}

It is sometimes necessary to prevent \LaTeX\ from
breaking a line where it might otherwise do so.
This may be at a space, as between the ``Mr.''\ and
``Jones'' in
       ``Mr.~Jones'',        % ~ produces an unbreakable interword space.
or within a word---especially when the word is a
symbol like
       \mbox{\emph{itemnum}}
that makes little sense when hyphenated across
lines.

Footnotes\footnote{This is an example of a footnote.}
pose no problem.

\LaTeX\ is good at typesetting mathematical formulas
like
       \( x-3y + z = 7 \)
or
       \( a_{1} > x^{2n} + y^{2n} > x' \)
or
       \( \ip{A}{B} = \sum_{i} a_{i} b_{i} \).
The spaces you type in a formula are
ignored.  Remember that a letter like
       $x$                   % $ ... $  and  \( ... \)  are equivalent
is a formula when it denotes a mathematical
symbol, and it should be typed as one.

\section{Displayed Text}

Text is displayed by indenting it from the left
margin.  Quotations are commonly displayed.  There
are short quotations
\begin{quote}
   This is a short quotation.  It consists of a
   single paragraph of text.  See how it is formatted.
\end{quote}
and longer ones.
\begin{quotation}
   This is a longer quotation.  It consists of two
   paragraphs of text, neither of which are
   particularly interesting.

   This is the second paragraph of the quotation.  It
   is just as dull as the first paragraph.
\end{quotation}
Another frequently-displayed structure is a list.
The following is an example of an \emph{itemized}
list.
\begin{itemize}
   \item This is the first item of an itemized list.
         Each item in the list is marked with a ``tick''.
         You don't have to worry about what kind of tick
         mark is used.

   \item This is the second item of the list.  It
         contains another list nested inside it.  The inner
         list is an \emph{enumerated} list.
         \begin{enumerate}
            \item This is the first item of an enumerated
                  list that is nested within the itemized list.

            \item This is the second item of the inner list.
                  \LaTeX\ allows you to nest lists deeper than
                  you really should.
         \end{enumerate}
         This is the rest of the second item of the outer
         list.  It is no more interesting than any other
         part of the item.
   \item This is the third item of the list.
\end{itemize}
You can even display poetry.
\begin{verse}
   There is an environment
    for verse \\             % The \\ command separates lines
   Whose features some poets % within a stanza.
   will curse.

                             % One or more blank lines separate stanzas.

   For instead of making\\
   Them do \emph{all} line breaking, \\
   It allows them to put too many words on a line when they'd rather be
   forced to be terse.
\end{verse}

Mathematical formulas may also be displayed.  A
displayed formula
is
one-line long; multiline
formulas require special formatting instructions.
   \[  \ip{\Gamma}{\psi'} = x'' + y^{2} + z_{i}^{n}\]
Don't start a paragraph with a displayed equation,
nor make one a paragraph by itself.

\end{document}               % End of document.

\section{Bench Huge Section 503}
% synthetic bench section 503
% This is a sample LaTeX input file.  (Version of 12 August 2004.)
%
% A '%' character causes TeX to ignore all remaining text on the line,
% and is used for comments like this one.

\documentclass{article}      % Specifies the document class

                             % The preamble begins here.
\title{An Example Document}  % Declares the document's title.
\author{Leslie Lamport}      % Declares the author's name.
\date{January 21, 1994}      % Deleting this command produces today's date.

\newcommand{\ip}[2]{(#1, #2)}
                             % Defines \ip{arg1}{arg2} to mean
                             % (arg1, arg2).

%\newcommand{\ip}[2]{\langle #1 | #2\rangle}
                             % This is an alternative definition of
                             % \ip that is commented out.

\begin{document}             % End of preamble and beginning of text.

\maketitle                   % Produces the title.

This is an example input file.  Comparing it with
the output it generates can show you how to
produce a simple document of your own.

\section{Ordinary Text}      % Produces section heading.  Lower-level
                             % sections are begun with similar
                             % \subsection and \subsubsection commands.

The ends  of words and sentences are marked
  by   spaces. It  doesn't matter how many
spaces    you type; one is as good as 100.  The
end of   a line counts as a space.

One   or more   blank lines denote the  end
of  a paragraph.

Since any number of consecutive spaces are treated
like a single one, the formatting of the input
file makes no difference to
      \LaTeX,                % The \LaTeX command generates the LaTeX logo.
but it makes a difference to you.  When you use
\LaTeX, making your input file as easy to read
as possible will be a great help as you write
your document and when you change it.  This sample
file shows how you can add comments to your own input
file.

Because printing is different from typewriting,
there are a number of things that you have to do
differently when preparing an input file than if
you were just typing the document directly.
Quotation marks like
       ``this''
have to be handled specially, as do quotes within
quotes:
       ``\,`this'            % \, separates the double and single quote.
        is what I just
        wrote, not  `that'\,''.

Dashes come in three sizes: an
       intra-word
dash, a medium dash for number ranges like
       1--2,
and a punctuation
       dash---like
this.

A sentence-ending space should be larger than the
space between words within a sentence.  You
sometimes have to type special commands in
conjunction with punctuation characters to get
this right, as in the following sentence.
       Gnats, gnus, etc.\ all  % `\ ' makes an inter-word space.
       begin with G\@.         % \@ marks end-of-sentence punctuation.
You should check the spaces after periods when
reading your output to make sure you haven't
forgotten any special cases.  Generating an
ellipsis
       \ldots\               % `\ ' is needed after `\ldots' because TeX
                             % ignores spaces after command names like \ldots
                             % made from \ + letters.
                             %
                             % Note how a `%' character causes TeX to ignore
                             % the end of the input line, so these blank lines
                             % do not start a new paragraph.
                             %
with the right spacing around the periods requires
a special command.

\LaTeX\ interprets some common characters as
commands, so you must type special commands to
generate them.  These characters include the
following:
       \$ \& \% \# \{ and \}.

In printing, text is usually emphasized with an
       \emph{italic}
type style.

\begin{em}
   A long segment of text can also be emphasized
   in this way.  Text within such a segment can be
   given \emph{additional} emphasis.
\end{em}

It is sometimes necessary to prevent \LaTeX\ from
breaking a line where it might otherwise do so.
This may be at a space, as between the ``Mr.''\ and
``Jones'' in
       ``Mr.~Jones'',        % ~ produces an unbreakable interword space.
or within a word---especially when the word is a
symbol like
       \mbox{\emph{itemnum}}
that makes little sense when hyphenated across
lines.

Footnotes\footnote{This is an example of a footnote.}
pose no problem.

\LaTeX\ is good at typesetting mathematical formulas
like
       \( x-3y + z = 7 \)
or
       \( a_{1} > x^{2n} + y^{2n} > x' \)
or
       \( \ip{A}{B} = \sum_{i} a_{i} b_{i} \).
The spaces you type in a formula are
ignored.  Remember that a letter like
       $x$                   % $ ... $  and  \( ... \)  are equivalent
is a formula when it denotes a mathematical
symbol, and it should be typed as one.

\section{Displayed Text}

Text is displayed by indenting it from the left
margin.  Quotations are commonly displayed.  There
are short quotations
\begin{quote}
   This is a short quotation.  It consists of a
   single paragraph of text.  See how it is formatted.
\end{quote}
and longer ones.
\begin{quotation}
   This is a longer quotation.  It consists of two
   paragraphs of text, neither of which are
   particularly interesting.

   This is the second paragraph of the quotation.  It
   is just as dull as the first paragraph.
\end{quotation}
Another frequently-displayed structure is a list.
The following is an example of an \emph{itemized}
list.
\begin{itemize}
   \item This is the first item of an itemized list.
         Each item in the list is marked with a ``tick''.
         You don't have to worry about what kind of tick
         mark is used.

   \item This is the second item of the list.  It
         contains another list nested inside it.  The inner
         list is an \emph{enumerated} list.
         \begin{enumerate}
            \item This is the first item of an enumerated
                  list that is nested within the itemized list.

            \item This is the second item of the inner list.
                  \LaTeX\ allows you to nest lists deeper than
                  you really should.
         \end{enumerate}
         This is the rest of the second item of the outer
         list.  It is no more interesting than any other
         part of the item.
   \item This is the third item of the list.
\end{itemize}
You can even display poetry.
\begin{verse}
   There is an environment
    for verse \\             % The \\ command separates lines
   Whose features some poets % within a stanza.
   will curse.

                             % One or more blank lines separate stanzas.

   For instead of making\\
   Them do \emph{all} line breaking, \\
   It allows them to put too many words on a line when they'd rather be
   forced to be terse.
\end{verse}

Mathematical formulas may also be displayed.  A
displayed formula
is
one-line long; multiline
formulas require special formatting instructions.
   \[  \ip{\Gamma}{\psi'} = x'' + y^{2} + z_{i}^{n}\]
Don't start a paragraph with a displayed equation,
nor make one a paragraph by itself.

\end{document}               % End of document.

\section{Bench Huge Section 504}
% synthetic bench section 504
% This is a sample LaTeX input file.  (Version of 12 August 2004.)
%
% A '%' character causes TeX to ignore all remaining text on the line,
% and is used for comments like this one.

\documentclass{article}      % Specifies the document class

                             % The preamble begins here.
\title{An Example Document}  % Declares the document's title.
\author{Leslie Lamport}      % Declares the author's name.
\date{January 21, 1994}      % Deleting this command produces today's date.

\newcommand{\ip}[2]{(#1, #2)}
                             % Defines \ip{arg1}{arg2} to mean
                             % (arg1, arg2).

%\newcommand{\ip}[2]{\langle #1 | #2\rangle}
                             % This is an alternative definition of
                             % \ip that is commented out.

\begin{document}             % End of preamble and beginning of text.

\maketitle                   % Produces the title.

This is an example input file.  Comparing it with
the output it generates can show you how to
produce a simple document of your own.

\section{Ordinary Text}      % Produces section heading.  Lower-level
                             % sections are begun with similar
                             % \subsection and \subsubsection commands.

The ends  of words and sentences are marked
  by   spaces. It  doesn't matter how many
spaces    you type; one is as good as 100.  The
end of   a line counts as a space.

One   or more   blank lines denote the  end
of  a paragraph.

Since any number of consecutive spaces are treated
like a single one, the formatting of the input
file makes no difference to
      \LaTeX,                % The \LaTeX command generates the LaTeX logo.
but it makes a difference to you.  When you use
\LaTeX, making your input file as easy to read
as possible will be a great help as you write
your document and when you change it.  This sample
file shows how you can add comments to your own input
file.

Because printing is different from typewriting,
there are a number of things that you have to do
differently when preparing an input file than if
you were just typing the document directly.
Quotation marks like
       ``this''
have to be handled specially, as do quotes within
quotes:
       ``\,`this'            % \, separates the double and single quote.
        is what I just
        wrote, not  `that'\,''.

Dashes come in three sizes: an
       intra-word
dash, a medium dash for number ranges like
       1--2,
and a punctuation
       dash---like
this.

A sentence-ending space should be larger than the
space between words within a sentence.  You
sometimes have to type special commands in
conjunction with punctuation characters to get
this right, as in the following sentence.
       Gnats, gnus, etc.\ all  % `\ ' makes an inter-word space.
       begin with G\@.         % \@ marks end-of-sentence punctuation.
You should check the spaces after periods when
reading your output to make sure you haven't
forgotten any special cases.  Generating an
ellipsis
       \ldots\               % `\ ' is needed after `\ldots' because TeX
                             % ignores spaces after command names like \ldots
                             % made from \ + letters.
                             %
                             % Note how a `%' character causes TeX to ignore
                             % the end of the input line, so these blank lines
                             % do not start a new paragraph.
                             %
with the right spacing around the periods requires
a special command.

\LaTeX\ interprets some common characters as
commands, so you must type special commands to
generate them.  These characters include the
following:
       \$ \& \% \# \{ and \}.

In printing, text is usually emphasized with an
       \emph{italic}
type style.

\begin{em}
   A long segment of text can also be emphasized
   in this way.  Text within such a segment can be
   given \emph{additional} emphasis.
\end{em}

It is sometimes necessary to prevent \LaTeX\ from
breaking a line where it might otherwise do so.
This may be at a space, as between the ``Mr.''\ and
``Jones'' in
       ``Mr.~Jones'',        % ~ produces an unbreakable interword space.
or within a word---especially when the word is a
symbol like
       \mbox{\emph{itemnum}}
that makes little sense when hyphenated across
lines.

Footnotes\footnote{This is an example of a footnote.}
pose no problem.

\LaTeX\ is good at typesetting mathematical formulas
like
       \( x-3y + z = 7 \)
or
       \( a_{1} > x^{2n} + y^{2n} > x' \)
or
       \( \ip{A}{B} = \sum_{i} a_{i} b_{i} \).
The spaces you type in a formula are
ignored.  Remember that a letter like
       $x$                   % $ ... $  and  \( ... \)  are equivalent
is a formula when it denotes a mathematical
symbol, and it should be typed as one.

\section{Displayed Text}

Text is displayed by indenting it from the left
margin.  Quotations are commonly displayed.  There
are short quotations
\begin{quote}
   This is a short quotation.  It consists of a
   single paragraph of text.  See how it is formatted.
\end{quote}
and longer ones.
\begin{quotation}
   This is a longer quotation.  It consists of two
   paragraphs of text, neither of which are
   particularly interesting.

   This is the second paragraph of the quotation.  It
   is just as dull as the first paragraph.
\end{quotation}
Another frequently-displayed structure is a list.
The following is an example of an \emph{itemized}
list.
\begin{itemize}
   \item This is the first item of an itemized list.
         Each item in the list is marked with a ``tick''.
         You don't have to worry about what kind of tick
         mark is used.

   \item This is the second item of the list.  It
         contains another list nested inside it.  The inner
         list is an \emph{enumerated} list.
         \begin{enumerate}
            \item This is the first item of an enumerated
                  list that is nested within the itemized list.

            \item This is the second item of the inner list.
                  \LaTeX\ allows you to nest lists deeper than
                  you really should.
         \end{enumerate}
         This is the rest of the second item of the outer
         list.  It is no more interesting than any other
         part of the item.
   \item This is the third item of the list.
\end{itemize}
You can even display poetry.
\begin{verse}
   There is an environment
    for verse \\             % The \\ command separates lines
   Whose features some poets % within a stanza.
   will curse.

                             % One or more blank lines separate stanzas.

   For instead of making\\
   Them do \emph{all} line breaking, \\
   It allows them to put too many words on a line when they'd rather be
   forced to be terse.
\end{verse}

Mathematical formulas may also be displayed.  A
displayed formula
is
one-line long; multiline
formulas require special formatting instructions.
   \[  \ip{\Gamma}{\psi'} = x'' + y^{2} + z_{i}^{n}\]
Don't start a paragraph with a displayed equation,
nor make one a paragraph by itself.

\end{document}               % End of document.

\section{Bench Huge Section 505}
% synthetic bench section 505
% This is a sample LaTeX input file.  (Version of 12 August 2004.)
%
% A '%' character causes TeX to ignore all remaining text on the line,
% and is used for comments like this one.

\documentclass{article}      % Specifies the document class

                             % The preamble begins here.
\title{An Example Document}  % Declares the document's title.
\author{Leslie Lamport}      % Declares the author's name.
\date{January 21, 1994}      % Deleting this command produces today's date.

\newcommand{\ip}[2]{(#1, #2)}
                             % Defines \ip{arg1}{arg2} to mean
                             % (arg1, arg2).

%\newcommand{\ip}[2]{\langle #1 | #2\rangle}
                             % This is an alternative definition of
                             % \ip that is commented out.

\begin{document}             % End of preamble and beginning of text.

\maketitle                   % Produces the title.

This is an example input file.  Comparing it with
the output it generates can show you how to
produce a simple document of your own.

\section{Ordinary Text}      % Produces section heading.  Lower-level
                             % sections are begun with similar
                             % \subsection and \subsubsection commands.

The ends  of words and sentences are marked
  by   spaces. It  doesn't matter how many
spaces    you type; one is as good as 100.  The
end of   a line counts as a space.

One   or more   blank lines denote the  end
of  a paragraph.

Since any number of consecutive spaces are treated
like a single one, the formatting of the input
file makes no difference to
      \LaTeX,                % The \LaTeX command generates the LaTeX logo.
but it makes a difference to you.  When you use
\LaTeX, making your input file as easy to read
as possible will be a great help as you write
your document and when you change it.  This sample
file shows how you can add comments to your own input
file.

Because printing is different from typewriting,
there are a number of things that you have to do
differently when preparing an input file than if
you were just typing the document directly.
Quotation marks like
       ``this''
have to be handled specially, as do quotes within
quotes:
       ``\,`this'            % \, separates the double and single quote.
        is what I just
        wrote, not  `that'\,''.

Dashes come in three sizes: an
       intra-word
dash, a medium dash for number ranges like
       1--2,
and a punctuation
       dash---like
this.

A sentence-ending space should be larger than the
space between words within a sentence.  You
sometimes have to type special commands in
conjunction with punctuation characters to get
this right, as in the following sentence.
       Gnats, gnus, etc.\ all  % `\ ' makes an inter-word space.
       begin with G\@.         % \@ marks end-of-sentence punctuation.
You should check the spaces after periods when
reading your output to make sure you haven't
forgotten any special cases.  Generating an
ellipsis
       \ldots\               % `\ ' is needed after `\ldots' because TeX
                             % ignores spaces after command names like \ldots
                             % made from \ + letters.
                             %
                             % Note how a `%' character causes TeX to ignore
                             % the end of the input line, so these blank lines
                             % do not start a new paragraph.
                             %
with the right spacing around the periods requires
a special command.

\LaTeX\ interprets some common characters as
commands, so you must type special commands to
generate them.  These characters include the
following:
       \$ \& \% \# \{ and \}.

In printing, text is usually emphasized with an
       \emph{italic}
type style.

\begin{em}
   A long segment of text can also be emphasized
   in this way.  Text within such a segment can be
   given \emph{additional} emphasis.
\end{em}

It is sometimes necessary to prevent \LaTeX\ from
breaking a line where it might otherwise do so.
This may be at a space, as between the ``Mr.''\ and
``Jones'' in
       ``Mr.~Jones'',        % ~ produces an unbreakable interword space.
or within a word---especially when the word is a
symbol like
       \mbox{\emph{itemnum}}
that makes little sense when hyphenated across
lines.

Footnotes\footnote{This is an example of a footnote.}
pose no problem.

\LaTeX\ is good at typesetting mathematical formulas
like
       \( x-3y + z = 7 \)
or
       \( a_{1} > x^{2n} + y^{2n} > x' \)
or
       \( \ip{A}{B} = \sum_{i} a_{i} b_{i} \).
The spaces you type in a formula are
ignored.  Remember that a letter like
       $x$                   % $ ... $  and  \( ... \)  are equivalent
is a formula when it denotes a mathematical
symbol, and it should be typed as one.

\section{Displayed Text}

Text is displayed by indenting it from the left
margin.  Quotations are commonly displayed.  There
are short quotations
\begin{quote}
   This is a short quotation.  It consists of a
   single paragraph of text.  See how it is formatted.
\end{quote}
and longer ones.
\begin{quotation}
   This is a longer quotation.  It consists of two
   paragraphs of text, neither of which are
   particularly interesting.

   This is the second paragraph of the quotation.  It
   is just as dull as the first paragraph.
\end{quotation}
Another frequently-displayed structure is a list.
The following is an example of an \emph{itemized}
list.
\begin{itemize}
   \item This is the first item of an itemized list.
         Each item in the list is marked with a ``tick''.
         You don't have to worry about what kind of tick
         mark is used.

   \item This is the second item of the list.  It
         contains another list nested inside it.  The inner
         list is an \emph{enumerated} list.
         \begin{enumerate}
            \item This is the first item of an enumerated
                  list that is nested within the itemized list.

            \item This is the second item of the inner list.
                  \LaTeX\ allows you to nest lists deeper than
                  you really should.
         \end{enumerate}
         This is the rest of the second item of the outer
         list.  It is no more interesting than any other
         part of the item.
   \item This is the third item of the list.
\end{itemize}
You can even display poetry.
\begin{verse}
   There is an environment
    for verse \\             % The \\ command separates lines
   Whose features some poets % within a stanza.
   will curse.

                             % One or more blank lines separate stanzas.

   For instead of making\\
   Them do \emph{all} line breaking, \\
   It allows them to put too many words on a line when they'd rather be
   forced to be terse.
\end{verse}

Mathematical formulas may also be displayed.  A
displayed formula
is
one-line long; multiline
formulas require special formatting instructions.
   \[  \ip{\Gamma}{\psi'} = x'' + y^{2} + z_{i}^{n}\]
Don't start a paragraph with a displayed equation,
nor make one a paragraph by itself.

\end{document}               % End of document.

\section{Bench Huge Section 506}
% synthetic bench section 506
% This is a sample LaTeX input file.  (Version of 12 August 2004.)
%
% A '%' character causes TeX to ignore all remaining text on the line,
% and is used for comments like this one.

\documentclass{article}      % Specifies the document class

                             % The preamble begins here.
\title{An Example Document}  % Declares the document's title.
\author{Leslie Lamport}      % Declares the author's name.
\date{January 21, 1994}      % Deleting this command produces today's date.

\newcommand{\ip}[2]{(#1, #2)}
                             % Defines \ip{arg1}{arg2} to mean
                             % (arg1, arg2).

%\newcommand{\ip}[2]{\langle #1 | #2\rangle}
                             % This is an alternative definition of
                             % \ip that is commented out.

\begin{document}             % End of preamble and beginning of text.

\maketitle                   % Produces the title.

This is an example input file.  Comparing it with
the output it generates can show you how to
produce a simple document of your own.

\section{Ordinary Text}      % Produces section heading.  Lower-level
                             % sections are begun with similar
                             % \subsection and \subsubsection commands.

The ends  of words and sentences are marked
  by   spaces. It  doesn't matter how many
spaces    you type; one is as good as 100.  The
end of   a line counts as a space.

One   or more   blank lines denote the  end
of  a paragraph.

Since any number of consecutive spaces are treated
like a single one, the formatting of the input
file makes no difference to
      \LaTeX,                % The \LaTeX command generates the LaTeX logo.
but it makes a difference to you.  When you use
\LaTeX, making your input file as easy to read
as possible will be a great help as you write
your document and when you change it.  This sample
file shows how you can add comments to your own input
file.

Because printing is different from typewriting,
there are a number of things that you have to do
differently when preparing an input file than if
you were just typing the document directly.
Quotation marks like
       ``this''
have to be handled specially, as do quotes within
quotes:
       ``\,`this'            % \, separates the double and single quote.
        is what I just
        wrote, not  `that'\,''.

Dashes come in three sizes: an
       intra-word
dash, a medium dash for number ranges like
       1--2,
and a punctuation
       dash---like
this.

A sentence-ending space should be larger than the
space between words within a sentence.  You
sometimes have to type special commands in
conjunction with punctuation characters to get
this right, as in the following sentence.
       Gnats, gnus, etc.\ all  % `\ ' makes an inter-word space.
       begin with G\@.         % \@ marks end-of-sentence punctuation.
You should check the spaces after periods when
reading your output to make sure you haven't
forgotten any special cases.  Generating an
ellipsis
       \ldots\               % `\ ' is needed after `\ldots' because TeX
                             % ignores spaces after command names like \ldots
                             % made from \ + letters.
                             %
                             % Note how a `%' character causes TeX to ignore
                             % the end of the input line, so these blank lines
                             % do not start a new paragraph.
                             %
with the right spacing around the periods requires
a special command.

\LaTeX\ interprets some common characters as
commands, so you must type special commands to
generate them.  These characters include the
following:
       \$ \& \% \# \{ and \}.

In printing, text is usually emphasized with an
       \emph{italic}
type style.

\begin{em}
   A long segment of text can also be emphasized
   in this way.  Text within such a segment can be
   given \emph{additional} emphasis.
\end{em}

It is sometimes necessary to prevent \LaTeX\ from
breaking a line where it might otherwise do so.
This may be at a space, as between the ``Mr.''\ and
``Jones'' in
       ``Mr.~Jones'',        % ~ produces an unbreakable interword space.
or within a word---especially when the word is a
symbol like
       \mbox{\emph{itemnum}}
that makes little sense when hyphenated across
lines.

Footnotes\footnote{This is an example of a footnote.}
pose no problem.

\LaTeX\ is good at typesetting mathematical formulas
like
       \( x-3y + z = 7 \)
or
       \( a_{1} > x^{2n} + y^{2n} > x' \)
or
       \( \ip{A}{B} = \sum_{i} a_{i} b_{i} \).
The spaces you type in a formula are
ignored.  Remember that a letter like
       $x$                   % $ ... $  and  \( ... \)  are equivalent
is a formula when it denotes a mathematical
symbol, and it should be typed as one.

\section{Displayed Text}

Text is displayed by indenting it from the left
margin.  Quotations are commonly displayed.  There
are short quotations
\begin{quote}
   This is a short quotation.  It consists of a
   single paragraph of text.  See how it is formatted.
\end{quote}
and longer ones.
\begin{quotation}
   This is a longer quotation.  It consists of two
   paragraphs of text, neither of which are
   particularly interesting.

   This is the second paragraph of the quotation.  It
   is just as dull as the first paragraph.
\end{quotation}
Another frequently-displayed structure is a list.
The following is an example of an \emph{itemized}
list.
\begin{itemize}
   \item This is the first item of an itemized list.
         Each item in the list is marked with a ``tick''.
         You don't have to worry about what kind of tick
         mark is used.

   \item This is the second item of the list.  It
         contains another list nested inside it.  The inner
         list is an \emph{enumerated} list.
         \begin{enumerate}
            \item This is the first item of an enumerated
                  list that is nested within the itemized list.

            \item This is the second item of the inner list.
                  \LaTeX\ allows you to nest lists deeper than
                  you really should.
         \end{enumerate}
         This is the rest of the second item of the outer
         list.  It is no more interesting than any other
         part of the item.
   \item This is the third item of the list.
\end{itemize}
You can even display poetry.
\begin{verse}
   There is an environment
    for verse \\             % The \\ command separates lines
   Whose features some poets % within a stanza.
   will curse.

                             % One or more blank lines separate stanzas.

   For instead of making\\
   Them do \emph{all} line breaking, \\
   It allows them to put too many words on a line when they'd rather be
   forced to be terse.
\end{verse}

Mathematical formulas may also be displayed.  A
displayed formula
is
one-line long; multiline
formulas require special formatting instructions.
   \[  \ip{\Gamma}{\psi'} = x'' + y^{2} + z_{i}^{n}\]
Don't start a paragraph with a displayed equation,
nor make one a paragraph by itself.

\end{document}               % End of document.

\section{Bench Huge Section 507}
% synthetic bench section 507
% This is a sample LaTeX input file.  (Version of 12 August 2004.)
%
% A '%' character causes TeX to ignore all remaining text on the line,
% and is used for comments like this one.

\documentclass{article}      % Specifies the document class

                             % The preamble begins here.
\title{An Example Document}  % Declares the document's title.
\author{Leslie Lamport}      % Declares the author's name.
\date{January 21, 1994}      % Deleting this command produces today's date.

\newcommand{\ip}[2]{(#1, #2)}
                             % Defines \ip{arg1}{arg2} to mean
                             % (arg1, arg2).

%\newcommand{\ip}[2]{\langle #1 | #2\rangle}
                             % This is an alternative definition of
                             % \ip that is commented out.

\begin{document}             % End of preamble and beginning of text.

\maketitle                   % Produces the title.

This is an example input file.  Comparing it with
the output it generates can show you how to
produce a simple document of your own.

\section{Ordinary Text}      % Produces section heading.  Lower-level
                             % sections are begun with similar
                             % \subsection and \subsubsection commands.

The ends  of words and sentences are marked
  by   spaces. It  doesn't matter how many
spaces    you type; one is as good as 100.  The
end of   a line counts as a space.

One   or more   blank lines denote the  end
of  a paragraph.

Since any number of consecutive spaces are treated
like a single one, the formatting of the input
file makes no difference to
      \LaTeX,                % The \LaTeX command generates the LaTeX logo.
but it makes a difference to you.  When you use
\LaTeX, making your input file as easy to read
as possible will be a great help as you write
your document and when you change it.  This sample
file shows how you can add comments to your own input
file.

Because printing is different from typewriting,
there are a number of things that you have to do
differently when preparing an input file than if
you were just typing the document directly.
Quotation marks like
       ``this''
have to be handled specially, as do quotes within
quotes:
       ``\,`this'            % \, separates the double and single quote.
        is what I just
        wrote, not  `that'\,''.

Dashes come in three sizes: an
       intra-word
dash, a medium dash for number ranges like
       1--2,
and a punctuation
       dash---like
this.

A sentence-ending space should be larger than the
space between words within a sentence.  You
sometimes have to type special commands in
conjunction with punctuation characters to get
this right, as in the following sentence.
       Gnats, gnus, etc.\ all  % `\ ' makes an inter-word space.
       begin with G\@.         % \@ marks end-of-sentence punctuation.
You should check the spaces after periods when
reading your output to make sure you haven't
forgotten any special cases.  Generating an
ellipsis
       \ldots\               % `\ ' is needed after `\ldots' because TeX
                             % ignores spaces after command names like \ldots
                             % made from \ + letters.
                             %
                             % Note how a `%' character causes TeX to ignore
                             % the end of the input line, so these blank lines
                             % do not start a new paragraph.
                             %
with the right spacing around the periods requires
a special command.

\LaTeX\ interprets some common characters as
commands, so you must type special commands to
generate them.  These characters include the
following:
       \$ \& \% \# \{ and \}.

In printing, text is usually emphasized with an
       \emph{italic}
type style.

\begin{em}
   A long segment of text can also be emphasized
   in this way.  Text within such a segment can be
   given \emph{additional} emphasis.
\end{em}

It is sometimes necessary to prevent \LaTeX\ from
breaking a line where it might otherwise do so.
This may be at a space, as between the ``Mr.''\ and
``Jones'' in
       ``Mr.~Jones'',        % ~ produces an unbreakable interword space.
or within a word---especially when the word is a
symbol like
       \mbox{\emph{itemnum}}
that makes little sense when hyphenated across
lines.

Footnotes\footnote{This is an example of a footnote.}
pose no problem.

\LaTeX\ is good at typesetting mathematical formulas
like
       \( x-3y + z = 7 \)
or
       \( a_{1} > x^{2n} + y^{2n} > x' \)
or
       \( \ip{A}{B} = \sum_{i} a_{i} b_{i} \).
The spaces you type in a formula are
ignored.  Remember that a letter like
       $x$                   % $ ... $  and  \( ... \)  are equivalent
is a formula when it denotes a mathematical
symbol, and it should be typed as one.

\section{Displayed Text}

Text is displayed by indenting it from the left
margin.  Quotations are commonly displayed.  There
are short quotations
\begin{quote}
   This is a short quotation.  It consists of a
   single paragraph of text.  See how it is formatted.
\end{quote}
and longer ones.
\begin{quotation}
   This is a longer quotation.  It consists of two
   paragraphs of text, neither of which are
   particularly interesting.

   This is the second paragraph of the quotation.  It
   is just as dull as the first paragraph.
\end{quotation}
Another frequently-displayed structure is a list.
The following is an example of an \emph{itemized}
list.
\begin{itemize}
   \item This is the first item of an itemized list.
         Each item in the list is marked with a ``tick''.
         You don't have to worry about what kind of tick
         mark is used.

   \item This is the second item of the list.  It
         contains another list nested inside it.  The inner
         list is an \emph{enumerated} list.
         \begin{enumerate}
            \item This is the first item of an enumerated
                  list that is nested within the itemized list.

            \item This is the second item of the inner list.
                  \LaTeX\ allows you to nest lists deeper than
                  you really should.
         \end{enumerate}
         This is the rest of the second item of the outer
         list.  It is no more interesting than any other
         part of the item.
   \item This is the third item of the list.
\end{itemize}
You can even display poetry.
\begin{verse}
   There is an environment
    for verse \\             % The \\ command separates lines
   Whose features some poets % within a stanza.
   will curse.

                             % One or more blank lines separate stanzas.

   For instead of making\\
   Them do \emph{all} line breaking, \\
   It allows them to put too many words on a line when they'd rather be
   forced to be terse.
\end{verse}

Mathematical formulas may also be displayed.  A
displayed formula
is
one-line long; multiline
formulas require special formatting instructions.
   \[  \ip{\Gamma}{\psi'} = x'' + y^{2} + z_{i}^{n}\]
Don't start a paragraph with a displayed equation,
nor make one a paragraph by itself.

\end{document}               % End of document.

\section{Bench Huge Section 508}
% synthetic bench section 508
% This is a sample LaTeX input file.  (Version of 12 August 2004.)
%
% A '%' character causes TeX to ignore all remaining text on the line,
% and is used for comments like this one.

\documentclass{article}      % Specifies the document class

                             % The preamble begins here.
\title{An Example Document}  % Declares the document's title.
\author{Leslie Lamport}      % Declares the author's name.
\date{January 21, 1994}      % Deleting this command produces today's date.

\newcommand{\ip}[2]{(#1, #2)}
                             % Defines \ip{arg1}{arg2} to mean
                             % (arg1, arg2).

%\newcommand{\ip}[2]{\langle #1 | #2\rangle}
                             % This is an alternative definition of
                             % \ip that is commented out.

\begin{document}             % End of preamble and beginning of text.

\maketitle                   % Produces the title.

This is an example input file.  Comparing it with
the output it generates can show you how to
produce a simple document of your own.

\section{Ordinary Text}      % Produces section heading.  Lower-level
                             % sections are begun with similar
                             % \subsection and \subsubsection commands.

The ends  of words and sentences are marked
  by   spaces. It  doesn't matter how many
spaces    you type; one is as good as 100.  The
end of   a line counts as a space.

One   or more   blank lines denote the  end
of  a paragraph.

Since any number of consecutive spaces are treated
like a single one, the formatting of the input
file makes no difference to
      \LaTeX,                % The \LaTeX command generates the LaTeX logo.
but it makes a difference to you.  When you use
\LaTeX, making your input file as easy to read
as possible will be a great help as you write
your document and when you change it.  This sample
file shows how you can add comments to your own input
file.

Because printing is different from typewriting,
there are a number of things that you have to do
differently when preparing an input file than if
you were just typing the document directly.
Quotation marks like
       ``this''
have to be handled specially, as do quotes within
quotes:
       ``\,`this'            % \, separates the double and single quote.
        is what I just
        wrote, not  `that'\,''.

Dashes come in three sizes: an
       intra-word
dash, a medium dash for number ranges like
       1--2,
and a punctuation
       dash---like
this.

A sentence-ending space should be larger than the
space between words within a sentence.  You
sometimes have to type special commands in
conjunction with punctuation characters to get
this right, as in the following sentence.
       Gnats, gnus, etc.\ all  % `\ ' makes an inter-word space.
       begin with G\@.         % \@ marks end-of-sentence punctuation.
You should check the spaces after periods when
reading your output to make sure you haven't
forgotten any special cases.  Generating an
ellipsis
       \ldots\               % `\ ' is needed after `\ldots' because TeX
                             % ignores spaces after command names like \ldots
                             % made from \ + letters.
                             %
                             % Note how a `%' character causes TeX to ignore
                             % the end of the input line, so these blank lines
                             % do not start a new paragraph.
                             %
with the right spacing around the periods requires
a special command.

\LaTeX\ interprets some common characters as
commands, so you must type special commands to
generate them.  These characters include the
following:
       \$ \& \% \# \{ and \}.

In printing, text is usually emphasized with an
       \emph{italic}
type style.

\begin{em}
   A long segment of text can also be emphasized
   in this way.  Text within such a segment can be
   given \emph{additional} emphasis.
\end{em}

It is sometimes necessary to prevent \LaTeX\ from
breaking a line where it might otherwise do so.
This may be at a space, as between the ``Mr.''\ and
``Jones'' in
       ``Mr.~Jones'',        % ~ produces an unbreakable interword space.
or within a word---especially when the word is a
symbol like
       \mbox{\emph{itemnum}}
that makes little sense when hyphenated across
lines.

Footnotes\footnote{This is an example of a footnote.}
pose no problem.

\LaTeX\ is good at typesetting mathematical formulas
like
       \( x-3y + z = 7 \)
or
       \( a_{1} > x^{2n} + y^{2n} > x' \)
or
       \( \ip{A}{B} = \sum_{i} a_{i} b_{i} \).
The spaces you type in a formula are
ignored.  Remember that a letter like
       $x$                   % $ ... $  and  \( ... \)  are equivalent
is a formula when it denotes a mathematical
symbol, and it should be typed as one.

\section{Displayed Text}

Text is displayed by indenting it from the left
margin.  Quotations are commonly displayed.  There
are short quotations
\begin{quote}
   This is a short quotation.  It consists of a
   single paragraph of text.  See how it is formatted.
\end{quote}
and longer ones.
\begin{quotation}
   This is a longer quotation.  It consists of two
   paragraphs of text, neither of which are
   particularly interesting.

   This is the second paragraph of the quotation.  It
   is just as dull as the first paragraph.
\end{quotation}
Another frequently-displayed structure is a list.
The following is an example of an \emph{itemized}
list.
\begin{itemize}
   \item This is the first item of an itemized list.
         Each item in the list is marked with a ``tick''.
         You don't have to worry about what kind of tick
         mark is used.

   \item This is the second item of the list.  It
         contains another list nested inside it.  The inner
         list is an \emph{enumerated} list.
         \begin{enumerate}
            \item This is the first item of an enumerated
                  list that is nested within the itemized list.

            \item This is the second item of the inner list.
                  \LaTeX\ allows you to nest lists deeper than
                  you really should.
         \end{enumerate}
         This is the rest of the second item of the outer
         list.  It is no more interesting than any other
         part of the item.
   \item This is the third item of the list.
\end{itemize}
You can even display poetry.
\begin{verse}
   There is an environment
    for verse \\             % The \\ command separates lines
   Whose features some poets % within a stanza.
   will curse.

                             % One or more blank lines separate stanzas.

   For instead of making\\
   Them do \emph{all} line breaking, \\
   It allows them to put too many words on a line when they'd rather be
   forced to be terse.
\end{verse}

Mathematical formulas may also be displayed.  A
displayed formula
is
one-line long; multiline
formulas require special formatting instructions.
   \[  \ip{\Gamma}{\psi'} = x'' + y^{2} + z_{i}^{n}\]
Don't start a paragraph with a displayed equation,
nor make one a paragraph by itself.

\end{document}               % End of document.

\section{Bench Huge Section 509}
% synthetic bench section 509
% This is a sample LaTeX input file.  (Version of 12 August 2004.)
%
% A '%' character causes TeX to ignore all remaining text on the line,
% and is used for comments like this one.

\documentclass{article}      % Specifies the document class

                             % The preamble begins here.
\title{An Example Document}  % Declares the document's title.
\author{Leslie Lamport}      % Declares the author's name.
\date{January 21, 1994}      % Deleting this command produces today's date.

\newcommand{\ip}[2]{(#1, #2)}
                             % Defines \ip{arg1}{arg2} to mean
                             % (arg1, arg2).

%\newcommand{\ip}[2]{\langle #1 | #2\rangle}
                             % This is an alternative definition of
                             % \ip that is commented out.

\begin{document}             % End of preamble and beginning of text.

\maketitle                   % Produces the title.

This is an example input file.  Comparing it with
the output it generates can show you how to
produce a simple document of your own.

\section{Ordinary Text}      % Produces section heading.  Lower-level
                             % sections are begun with similar
                             % \subsection and \subsubsection commands.

The ends  of words and sentences are marked
  by   spaces. It  doesn't matter how many
spaces    you type; one is as good as 100.  The
end of   a line counts as a space.

One   or more   blank lines denote the  end
of  a paragraph.

Since any number of consecutive spaces are treated
like a single one, the formatting of the input
file makes no difference to
      \LaTeX,                % The \LaTeX command generates the LaTeX logo.
but it makes a difference to you.  When you use
\LaTeX, making your input file as easy to read
as possible will be a great help as you write
your document and when you change it.  This sample
file shows how you can add comments to your own input
file.

Because printing is different from typewriting,
there are a number of things that you have to do
differently when preparing an input file than if
you were just typing the document directly.
Quotation marks like
       ``this''
have to be handled specially, as do quotes within
quotes:
       ``\,`this'            % \, separates the double and single quote.
        is what I just
        wrote, not  `that'\,''.

Dashes come in three sizes: an
       intra-word
dash, a medium dash for number ranges like
       1--2,
and a punctuation
       dash---like
this.

A sentence-ending space should be larger than the
space between words within a sentence.  You
sometimes have to type special commands in
conjunction with punctuation characters to get
this right, as in the following sentence.
       Gnats, gnus, etc.\ all  % `\ ' makes an inter-word space.
       begin with G\@.         % \@ marks end-of-sentence punctuation.
You should check the spaces after periods when
reading your output to make sure you haven't
forgotten any special cases.  Generating an
ellipsis
       \ldots\               % `\ ' is needed after `\ldots' because TeX
                             % ignores spaces after command names like \ldots
                             % made from \ + letters.
                             %
                             % Note how a `%' character causes TeX to ignore
                             % the end of the input line, so these blank lines
                             % do not start a new paragraph.
                             %
with the right spacing around the periods requires
a special command.

\LaTeX\ interprets some common characters as
commands, so you must type special commands to
generate them.  These characters include the
following:
       \$ \& \% \# \{ and \}.

In printing, text is usually emphasized with an
       \emph{italic}
type style.

\begin{em}
   A long segment of text can also be emphasized
   in this way.  Text within such a segment can be
   given \emph{additional} emphasis.
\end{em}

It is sometimes necessary to prevent \LaTeX\ from
breaking a line where it might otherwise do so.
This may be at a space, as between the ``Mr.''\ and
``Jones'' in
       ``Mr.~Jones'',        % ~ produces an unbreakable interword space.
or within a word---especially when the word is a
symbol like
       \mbox{\emph{itemnum}}
that makes little sense when hyphenated across
lines.

Footnotes\footnote{This is an example of a footnote.}
pose no problem.

\LaTeX\ is good at typesetting mathematical formulas
like
       \( x-3y + z = 7 \)
or
       \( a_{1} > x^{2n} + y^{2n} > x' \)
or
       \( \ip{A}{B} = \sum_{i} a_{i} b_{i} \).
The spaces you type in a formula are
ignored.  Remember that a letter like
       $x$                   % $ ... $  and  \( ... \)  are equivalent
is a formula when it denotes a mathematical
symbol, and it should be typed as one.

\section{Displayed Text}

Text is displayed by indenting it from the left
margin.  Quotations are commonly displayed.  There
are short quotations
\begin{quote}
   This is a short quotation.  It consists of a
   single paragraph of text.  See how it is formatted.
\end{quote}
and longer ones.
\begin{quotation}
   This is a longer quotation.  It consists of two
   paragraphs of text, neither of which are
   particularly interesting.

   This is the second paragraph of the quotation.  It
   is just as dull as the first paragraph.
\end{quotation}
Another frequently-displayed structure is a list.
The following is an example of an \emph{itemized}
list.
\begin{itemize}
   \item This is the first item of an itemized list.
         Each item in the list is marked with a ``tick''.
         You don't have to worry about what kind of tick
         mark is used.

   \item This is the second item of the list.  It
         contains another list nested inside it.  The inner
         list is an \emph{enumerated} list.
         \begin{enumerate}
            \item This is the first item of an enumerated
                  list that is nested within the itemized list.

            \item This is the second item of the inner list.
                  \LaTeX\ allows you to nest lists deeper than
                  you really should.
         \end{enumerate}
         This is the rest of the second item of the outer
         list.  It is no more interesting than any other
         part of the item.
   \item This is the third item of the list.
\end{itemize}
You can even display poetry.
\begin{verse}
   There is an environment
    for verse \\             % The \\ command separates lines
   Whose features some poets % within a stanza.
   will curse.

                             % One or more blank lines separate stanzas.

   For instead of making\\
   Them do \emph{all} line breaking, \\
   It allows them to put too many words on a line when they'd rather be
   forced to be terse.
\end{verse}

Mathematical formulas may also be displayed.  A
displayed formula
is
one-line long; multiline
formulas require special formatting instructions.
   \[  \ip{\Gamma}{\psi'} = x'' + y^{2} + z_{i}^{n}\]
Don't start a paragraph with a displayed equation,
nor make one a paragraph by itself.

\end{document}               % End of document.

\section{Bench Huge Section 510}
% synthetic bench section 510
% This is a sample LaTeX input file.  (Version of 12 August 2004.)
%
% A '%' character causes TeX to ignore all remaining text on the line,
% and is used for comments like this one.

\documentclass{article}      % Specifies the document class

                             % The preamble begins here.
\title{An Example Document}  % Declares the document's title.
\author{Leslie Lamport}      % Declares the author's name.
\date{January 21, 1994}      % Deleting this command produces today's date.

\newcommand{\ip}[2]{(#1, #2)}
                             % Defines \ip{arg1}{arg2} to mean
                             % (arg1, arg2).

%\newcommand{\ip}[2]{\langle #1 | #2\rangle}
                             % This is an alternative definition of
                             % \ip that is commented out.

\begin{document}             % End of preamble and beginning of text.

\maketitle                   % Produces the title.

This is an example input file.  Comparing it with
the output it generates can show you how to
produce a simple document of your own.

\section{Ordinary Text}      % Produces section heading.  Lower-level
                             % sections are begun with similar
                             % \subsection and \subsubsection commands.

The ends  of words and sentences are marked
  by   spaces. It  doesn't matter how many
spaces    you type; one is as good as 100.  The
end of   a line counts as a space.

One   or more   blank lines denote the  end
of  a paragraph.

Since any number of consecutive spaces are treated
like a single one, the formatting of the input
file makes no difference to
      \LaTeX,                % The \LaTeX command generates the LaTeX logo.
but it makes a difference to you.  When you use
\LaTeX, making your input file as easy to read
as possible will be a great help as you write
your document and when you change it.  This sample
file shows how you can add comments to your own input
file.

Because printing is different from typewriting,
there are a number of things that you have to do
differently when preparing an input file than if
you were just typing the document directly.
Quotation marks like
       ``this''
have to be handled specially, as do quotes within
quotes:
       ``\,`this'            % \, separates the double and single quote.
        is what I just
        wrote, not  `that'\,''.

Dashes come in three sizes: an
       intra-word
dash, a medium dash for number ranges like
       1--2,
and a punctuation
       dash---like
this.

A sentence-ending space should be larger than the
space between words within a sentence.  You
sometimes have to type special commands in
conjunction with punctuation characters to get
this right, as in the following sentence.
       Gnats, gnus, etc.\ all  % `\ ' makes an inter-word space.
       begin with G\@.         % \@ marks end-of-sentence punctuation.
You should check the spaces after periods when
reading your output to make sure you haven't
forgotten any special cases.  Generating an
ellipsis
       \ldots\               % `\ ' is needed after `\ldots' because TeX
                             % ignores spaces after command names like \ldots
                             % made from \ + letters.
                             %
                             % Note how a `%' character causes TeX to ignore
                             % the end of the input line, so these blank lines
                             % do not start a new paragraph.
                             %
with the right spacing around the periods requires
a special command.

\LaTeX\ interprets some common characters as
commands, so you must type special commands to
generate them.  These characters include the
following:
       \$ \& \% \# \{ and \}.

In printing, text is usually emphasized with an
       \emph{italic}
type style.

\begin{em}
   A long segment of text can also be emphasized
   in this way.  Text within such a segment can be
   given \emph{additional} emphasis.
\end{em}

It is sometimes necessary to prevent \LaTeX\ from
breaking a line where it might otherwise do so.
This may be at a space, as between the ``Mr.''\ and
``Jones'' in
       ``Mr.~Jones'',        % ~ produces an unbreakable interword space.
or within a word---especially when the word is a
symbol like
       \mbox{\emph{itemnum}}
that makes little sense when hyphenated across
lines.

Footnotes\footnote{This is an example of a footnote.}
pose no problem.

\LaTeX\ is good at typesetting mathematical formulas
like
       \( x-3y + z = 7 \)
or
       \( a_{1} > x^{2n} + y^{2n} > x' \)
or
       \( \ip{A}{B} = \sum_{i} a_{i} b_{i} \).
The spaces you type in a formula are
ignored.  Remember that a letter like
       $x$                   % $ ... $  and  \( ... \)  are equivalent
is a formula when it denotes a mathematical
symbol, and it should be typed as one.

\section{Displayed Text}

Text is displayed by indenting it from the left
margin.  Quotations are commonly displayed.  There
are short quotations
\begin{quote}
   This is a short quotation.  It consists of a
   single paragraph of text.  See how it is formatted.
\end{quote}
and longer ones.
\begin{quotation}
   This is a longer quotation.  It consists of two
   paragraphs of text, neither of which are
   particularly interesting.

   This is the second paragraph of the quotation.  It
   is just as dull as the first paragraph.
\end{quotation}
Another frequently-displayed structure is a list.
The following is an example of an \emph{itemized}
list.
\begin{itemize}
   \item This is the first item of an itemized list.
         Each item in the list is marked with a ``tick''.
         You don't have to worry about what kind of tick
         mark is used.

   \item This is the second item of the list.  It
         contains another list nested inside it.  The inner
         list is an \emph{enumerated} list.
         \begin{enumerate}
            \item This is the first item of an enumerated
                  list that is nested within the itemized list.

            \item This is the second item of the inner list.
                  \LaTeX\ allows you to nest lists deeper than
                  you really should.
         \end{enumerate}
         This is the rest of the second item of the outer
         list.  It is no more interesting than any other
         part of the item.
   \item This is the third item of the list.
\end{itemize}
You can even display poetry.
\begin{verse}
   There is an environment
    for verse \\             % The \\ command separates lines
   Whose features some poets % within a stanza.
   will curse.

                             % One or more blank lines separate stanzas.

   For instead of making\\
   Them do \emph{all} line breaking, \\
   It allows them to put too many words on a line when they'd rather be
   forced to be terse.
\end{verse}

Mathematical formulas may also be displayed.  A
displayed formula
is
one-line long; multiline
formulas require special formatting instructions.
   \[  \ip{\Gamma}{\psi'} = x'' + y^{2} + z_{i}^{n}\]
Don't start a paragraph with a displayed equation,
nor make one a paragraph by itself.

\end{document}               % End of document.

\section{Bench Huge Section 511}
% synthetic bench section 511
% This is a sample LaTeX input file.  (Version of 12 August 2004.)
%
% A '%' character causes TeX to ignore all remaining text on the line,
% and is used for comments like this one.

\documentclass{article}      % Specifies the document class

                             % The preamble begins here.
\title{An Example Document}  % Declares the document's title.
\author{Leslie Lamport}      % Declares the author's name.
\date{January 21, 1994}      % Deleting this command produces today's date.

\newcommand{\ip}[2]{(#1, #2)}
                             % Defines \ip{arg1}{arg2} to mean
                             % (arg1, arg2).

%\newcommand{\ip}[2]{\langle #1 | #2\rangle}
                             % This is an alternative definition of
                             % \ip that is commented out.

\begin{document}             % End of preamble and beginning of text.

\maketitle                   % Produces the title.

This is an example input file.  Comparing it with
the output it generates can show you how to
produce a simple document of your own.

\section{Ordinary Text}      % Produces section heading.  Lower-level
                             % sections are begun with similar
                             % \subsection and \subsubsection commands.

The ends  of words and sentences are marked
  by   spaces. It  doesn't matter how many
spaces    you type; one is as good as 100.  The
end of   a line counts as a space.

One   or more   blank lines denote the  end
of  a paragraph.

Since any number of consecutive spaces are treated
like a single one, the formatting of the input
file makes no difference to
      \LaTeX,                % The \LaTeX command generates the LaTeX logo.
but it makes a difference to you.  When you use
\LaTeX, making your input file as easy to read
as possible will be a great help as you write
your document and when you change it.  This sample
file shows how you can add comments to your own input
file.

Because printing is different from typewriting,
there are a number of things that you have to do
differently when preparing an input file than if
you were just typing the document directly.
Quotation marks like
       ``this''
have to be handled specially, as do quotes within
quotes:
       ``\,`this'            % \, separates the double and single quote.
        is what I just
        wrote, not  `that'\,''.

Dashes come in three sizes: an
       intra-word
dash, a medium dash for number ranges like
       1--2,
and a punctuation
       dash---like
this.

A sentence-ending space should be larger than the
space between words within a sentence.  You
sometimes have to type special commands in
conjunction with punctuation characters to get
this right, as in the following sentence.
       Gnats, gnus, etc.\ all  % `\ ' makes an inter-word space.
       begin with G\@.         % \@ marks end-of-sentence punctuation.
You should check the spaces after periods when
reading your output to make sure you haven't
forgotten any special cases.  Generating an
ellipsis
       \ldots\               % `\ ' is needed after `\ldots' because TeX
                             % ignores spaces after command names like \ldots
                             % made from \ + letters.
                             %
                             % Note how a `%' character causes TeX to ignore
                             % the end of the input line, so these blank lines
                             % do not start a new paragraph.
                             %
with the right spacing around the periods requires
a special command.

\LaTeX\ interprets some common characters as
commands, so you must type special commands to
generate them.  These characters include the
following:
       \$ \& \% \# \{ and \}.

In printing, text is usually emphasized with an
       \emph{italic}
type style.

\begin{em}
   A long segment of text can also be emphasized
   in this way.  Text within such a segment can be
   given \emph{additional} emphasis.
\end{em}

It is sometimes necessary to prevent \LaTeX\ from
breaking a line where it might otherwise do so.
This may be at a space, as between the ``Mr.''\ and
``Jones'' in
       ``Mr.~Jones'',        % ~ produces an unbreakable interword space.
or within a word---especially when the word is a
symbol like
       \mbox{\emph{itemnum}}
that makes little sense when hyphenated across
lines.

Footnotes\footnote{This is an example of a footnote.}
pose no problem.

\LaTeX\ is good at typesetting mathematical formulas
like
       \( x-3y + z = 7 \)
or
       \( a_{1} > x^{2n} + y^{2n} > x' \)
or
       \( \ip{A}{B} = \sum_{i} a_{i} b_{i} \).
The spaces you type in a formula are
ignored.  Remember that a letter like
       $x$                   % $ ... $  and  \( ... \)  are equivalent
is a formula when it denotes a mathematical
symbol, and it should be typed as one.

\section{Displayed Text}

Text is displayed by indenting it from the left
margin.  Quotations are commonly displayed.  There
are short quotations
\begin{quote}
   This is a short quotation.  It consists of a
   single paragraph of text.  See how it is formatted.
\end{quote}
and longer ones.
\begin{quotation}
   This is a longer quotation.  It consists of two
   paragraphs of text, neither of which are
   particularly interesting.

   This is the second paragraph of the quotation.  It
   is just as dull as the first paragraph.
\end{quotation}
Another frequently-displayed structure is a list.
The following is an example of an \emph{itemized}
list.
\begin{itemize}
   \item This is the first item of an itemized list.
         Each item in the list is marked with a ``tick''.
         You don't have to worry about what kind of tick
         mark is used.

   \item This is the second item of the list.  It
         contains another list nested inside it.  The inner
         list is an \emph{enumerated} list.
         \begin{enumerate}
            \item This is the first item of an enumerated
                  list that is nested within the itemized list.

            \item This is the second item of the inner list.
                  \LaTeX\ allows you to nest lists deeper than
                  you really should.
         \end{enumerate}
         This is the rest of the second item of the outer
         list.  It is no more interesting than any other
         part of the item.
   \item This is the third item of the list.
\end{itemize}
You can even display poetry.
\begin{verse}
   There is an environment
    for verse \\             % The \\ command separates lines
   Whose features some poets % within a stanza.
   will curse.

                             % One or more blank lines separate stanzas.

   For instead of making\\
   Them do \emph{all} line breaking, \\
   It allows them to put too many words on a line when they'd rather be
   forced to be terse.
\end{verse}

Mathematical formulas may also be displayed.  A
displayed formula
is
one-line long; multiline
formulas require special formatting instructions.
   \[  \ip{\Gamma}{\psi'} = x'' + y^{2} + z_{i}^{n}\]
Don't start a paragraph with a displayed equation,
nor make one a paragraph by itself.

\end{document}               % End of document.

\section{Bench Huge Section 512}
% synthetic bench section 512
% This is a sample LaTeX input file.  (Version of 12 August 2004.)
%
% A '%' character causes TeX to ignore all remaining text on the line,
% and is used for comments like this one.

\documentclass{article}      % Specifies the document class

                             % The preamble begins here.
\title{An Example Document}  % Declares the document's title.
\author{Leslie Lamport}      % Declares the author's name.
\date{January 21, 1994}      % Deleting this command produces today's date.

\newcommand{\ip}[2]{(#1, #2)}
                             % Defines \ip{arg1}{arg2} to mean
                             % (arg1, arg2).

%\newcommand{\ip}[2]{\langle #1 | #2\rangle}
                             % This is an alternative definition of
                             % \ip that is commented out.

\begin{document}             % End of preamble and beginning of text.

\maketitle                   % Produces the title.

This is an example input file.  Comparing it with
the output it generates can show you how to
produce a simple document of your own.

\section{Ordinary Text}      % Produces section heading.  Lower-level
                             % sections are begun with similar
                             % \subsection and \subsubsection commands.

The ends  of words and sentences are marked
  by   spaces. It  doesn't matter how many
spaces    you type; one is as good as 100.  The
end of   a line counts as a space.

One   or more   blank lines denote the  end
of  a paragraph.

Since any number of consecutive spaces are treated
like a single one, the formatting of the input
file makes no difference to
      \LaTeX,                % The \LaTeX command generates the LaTeX logo.
but it makes a difference to you.  When you use
\LaTeX, making your input file as easy to read
as possible will be a great help as you write
your document and when you change it.  This sample
file shows how you can add comments to your own input
file.

Because printing is different from typewriting,
there are a number of things that you have to do
differently when preparing an input file than if
you were just typing the document directly.
Quotation marks like
       ``this''
have to be handled specially, as do quotes within
quotes:
       ``\,`this'            % \, separates the double and single quote.
        is what I just
        wrote, not  `that'\,''.

Dashes come in three sizes: an
       intra-word
dash, a medium dash for number ranges like
       1--2,
and a punctuation
       dash---like
this.

A sentence-ending space should be larger than the
space between words within a sentence.  You
sometimes have to type special commands in
conjunction with punctuation characters to get
this right, as in the following sentence.
       Gnats, gnus, etc.\ all  % `\ ' makes an inter-word space.
       begin with G\@.         % \@ marks end-of-sentence punctuation.
You should check the spaces after periods when
reading your output to make sure you haven't
forgotten any special cases.  Generating an
ellipsis
       \ldots\               % `\ ' is needed after `\ldots' because TeX
                             % ignores spaces after command names like \ldots
                             % made from \ + letters.
                             %
                             % Note how a `%' character causes TeX to ignore
                             % the end of the input line, so these blank lines
                             % do not start a new paragraph.
                             %
with the right spacing around the periods requires
a special command.

\LaTeX\ interprets some common characters as
commands, so you must type special commands to
generate them.  These characters include the
following:
       \$ \& \% \# \{ and \}.

In printing, text is usually emphasized with an
       \emph{italic}
type style.

\begin{em}
   A long segment of text can also be emphasized
   in this way.  Text within such a segment can be
   given \emph{additional} emphasis.
\end{em}

It is sometimes necessary to prevent \LaTeX\ from
breaking a line where it might otherwise do so.
This may be at a space, as between the ``Mr.''\ and
``Jones'' in
       ``Mr.~Jones'',        % ~ produces an unbreakable interword space.
or within a word---especially when the word is a
symbol like
       \mbox{\emph{itemnum}}
that makes little sense when hyphenated across
lines.

Footnotes\footnote{This is an example of a footnote.}
pose no problem.

\LaTeX\ is good at typesetting mathematical formulas
like
       \( x-3y + z = 7 \)
or
       \( a_{1} > x^{2n} + y^{2n} > x' \)
or
       \( \ip{A}{B} = \sum_{i} a_{i} b_{i} \).
The spaces you type in a formula are
ignored.  Remember that a letter like
       $x$                   % $ ... $  and  \( ... \)  are equivalent
is a formula when it denotes a mathematical
symbol, and it should be typed as one.

\section{Displayed Text}

Text is displayed by indenting it from the left
margin.  Quotations are commonly displayed.  There
are short quotations
\begin{quote}
   This is a short quotation.  It consists of a
   single paragraph of text.  See how it is formatted.
\end{quote}
and longer ones.
\begin{quotation}
   This is a longer quotation.  It consists of two
   paragraphs of text, neither of which are
   particularly interesting.

   This is the second paragraph of the quotation.  It
   is just as dull as the first paragraph.
\end{quotation}
Another frequently-displayed structure is a list.
The following is an example of an \emph{itemized}
list.
\begin{itemize}
   \item This is the first item of an itemized list.
         Each item in the list is marked with a ``tick''.
         You don't have to worry about what kind of tick
         mark is used.

   \item This is the second item of the list.  It
         contains another list nested inside it.  The inner
         list is an \emph{enumerated} list.
         \begin{enumerate}
            \item This is the first item of an enumerated
                  list that is nested within the itemized list.

            \item This is the second item of the inner list.
                  \LaTeX\ allows you to nest lists deeper than
                  you really should.
         \end{enumerate}
         This is the rest of the second item of the outer
         list.  It is no more interesting than any other
         part of the item.
   \item This is the third item of the list.
\end{itemize}
You can even display poetry.
\begin{verse}
   There is an environment
    for verse \\             % The \\ command separates lines
   Whose features some poets % within a stanza.
   will curse.

                             % One or more blank lines separate stanzas.

   For instead of making\\
   Them do \emph{all} line breaking, \\
   It allows them to put too many words on a line when they'd rather be
   forced to be terse.
\end{verse}

Mathematical formulas may also be displayed.  A
displayed formula
is
one-line long; multiline
formulas require special formatting instructions.
   \[  \ip{\Gamma}{\psi'} = x'' + y^{2} + z_{i}^{n}\]
Don't start a paragraph with a displayed equation,
nor make one a paragraph by itself.

\end{document}               % End of document.

\section{Bench Huge Section 513}
% synthetic bench section 513
% This is a sample LaTeX input file.  (Version of 12 August 2004.)
%
% A '%' character causes TeX to ignore all remaining text on the line,
% and is used for comments like this one.

\documentclass{article}      % Specifies the document class

                             % The preamble begins here.
\title{An Example Document}  % Declares the document's title.
\author{Leslie Lamport}      % Declares the author's name.
\date{January 21, 1994}      % Deleting this command produces today's date.

\newcommand{\ip}[2]{(#1, #2)}
                             % Defines \ip{arg1}{arg2} to mean
                             % (arg1, arg2).

%\newcommand{\ip}[2]{\langle #1 | #2\rangle}
                             % This is an alternative definition of
                             % \ip that is commented out.

\begin{document}             % End of preamble and beginning of text.

\maketitle                   % Produces the title.

This is an example input file.  Comparing it with
the output it generates can show you how to
produce a simple document of your own.

\section{Ordinary Text}      % Produces section heading.  Lower-level
                             % sections are begun with similar
                             % \subsection and \subsubsection commands.

The ends  of words and sentences are marked
  by   spaces. It  doesn't matter how many
spaces    you type; one is as good as 100.  The
end of   a line counts as a space.

One   or more   blank lines denote the  end
of  a paragraph.

Since any number of consecutive spaces are treated
like a single one, the formatting of the input
file makes no difference to
      \LaTeX,                % The \LaTeX command generates the LaTeX logo.
but it makes a difference to you.  When you use
\LaTeX, making your input file as easy to read
as possible will be a great help as you write
your document and when you change it.  This sample
file shows how you can add comments to your own input
file.

Because printing is different from typewriting,
there are a number of things that you have to do
differently when preparing an input file than if
you were just typing the document directly.
Quotation marks like
       ``this''
have to be handled specially, as do quotes within
quotes:
       ``\,`this'            % \, separates the double and single quote.
        is what I just
        wrote, not  `that'\,''.

Dashes come in three sizes: an
       intra-word
dash, a medium dash for number ranges like
       1--2,
and a punctuation
       dash---like
this.

A sentence-ending space should be larger than the
space between words within a sentence.  You
sometimes have to type special commands in
conjunction with punctuation characters to get
this right, as in the following sentence.
       Gnats, gnus, etc.\ all  % `\ ' makes an inter-word space.
       begin with G\@.         % \@ marks end-of-sentence punctuation.
You should check the spaces after periods when
reading your output to make sure you haven't
forgotten any special cases.  Generating an
ellipsis
       \ldots\               % `\ ' is needed after `\ldots' because TeX
                             % ignores spaces after command names like \ldots
                             % made from \ + letters.
                             %
                             % Note how a `%' character causes TeX to ignore
                             % the end of the input line, so these blank lines
                             % do not start a new paragraph.
                             %
with the right spacing around the periods requires
a special command.

\LaTeX\ interprets some common characters as
commands, so you must type special commands to
generate them.  These characters include the
following:
       \$ \& \% \# \{ and \}.

In printing, text is usually emphasized with an
       \emph{italic}
type style.

\begin{em}
   A long segment of text can also be emphasized
   in this way.  Text within such a segment can be
   given \emph{additional} emphasis.
\end{em}

It is sometimes necessary to prevent \LaTeX\ from
breaking a line where it might otherwise do so.
This may be at a space, as between the ``Mr.''\ and
``Jones'' in
       ``Mr.~Jones'',        % ~ produces an unbreakable interword space.
or within a word---especially when the word is a
symbol like
       \mbox{\emph{itemnum}}
that makes little sense when hyphenated across
lines.

Footnotes\footnote{This is an example of a footnote.}
pose no problem.

\LaTeX\ is good at typesetting mathematical formulas
like
       \( x-3y + z = 7 \)
or
       \( a_{1} > x^{2n} + y^{2n} > x' \)
or
       \( \ip{A}{B} = \sum_{i} a_{i} b_{i} \).
The spaces you type in a formula are
ignored.  Remember that a letter like
       $x$                   % $ ... $  and  \( ... \)  are equivalent
is a formula when it denotes a mathematical
symbol, and it should be typed as one.

\section{Displayed Text}

Text is displayed by indenting it from the left
margin.  Quotations are commonly displayed.  There
are short quotations
\begin{quote}
   This is a short quotation.  It consists of a
   single paragraph of text.  See how it is formatted.
\end{quote}
and longer ones.
\begin{quotation}
   This is a longer quotation.  It consists of two
   paragraphs of text, neither of which are
   particularly interesting.

   This is the second paragraph of the quotation.  It
   is just as dull as the first paragraph.
\end{quotation}
Another frequently-displayed structure is a list.
The following is an example of an \emph{itemized}
list.
\begin{itemize}
   \item This is the first item of an itemized list.
         Each item in the list is marked with a ``tick''.
         You don't have to worry about what kind of tick
         mark is used.

   \item This is the second item of the list.  It
         contains another list nested inside it.  The inner
         list is an \emph{enumerated} list.
         \begin{enumerate}
            \item This is the first item of an enumerated
                  list that is nested within the itemized list.

            \item This is the second item of the inner list.
                  \LaTeX\ allows you to nest lists deeper than
                  you really should.
         \end{enumerate}
         This is the rest of the second item of the outer
         list.  It is no more interesting than any other
         part of the item.
   \item This is the third item of the list.
\end{itemize}
You can even display poetry.
\begin{verse}
   There is an environment
    for verse \\             % The \\ command separates lines
   Whose features some poets % within a stanza.
   will curse.

                             % One or more blank lines separate stanzas.

   For instead of making\\
   Them do \emph{all} line breaking, \\
   It allows them to put too many words on a line when they'd rather be
   forced to be terse.
\end{verse}

Mathematical formulas may also be displayed.  A
displayed formula
is
one-line long; multiline
formulas require special formatting instructions.
   \[  \ip{\Gamma}{\psi'} = x'' + y^{2} + z_{i}^{n}\]
Don't start a paragraph with a displayed equation,
nor make one a paragraph by itself.

\end{document}               % End of document.

\section{Bench Huge Section 514}
% synthetic bench section 514
% This is a sample LaTeX input file.  (Version of 12 August 2004.)
%
% A '%' character causes TeX to ignore all remaining text on the line,
% and is used for comments like this one.

\documentclass{article}      % Specifies the document class

                             % The preamble begins here.
\title{An Example Document}  % Declares the document's title.
\author{Leslie Lamport}      % Declares the author's name.
\date{January 21, 1994}      % Deleting this command produces today's date.

\newcommand{\ip}[2]{(#1, #2)}
                             % Defines \ip{arg1}{arg2} to mean
                             % (arg1, arg2).

%\newcommand{\ip}[2]{\langle #1 | #2\rangle}
                             % This is an alternative definition of
                             % \ip that is commented out.

\begin{document}             % End of preamble and beginning of text.

\maketitle                   % Produces the title.

This is an example input file.  Comparing it with
the output it generates can show you how to
produce a simple document of your own.

\section{Ordinary Text}      % Produces section heading.  Lower-level
                             % sections are begun with similar
                             % \subsection and \subsubsection commands.

The ends  of words and sentences are marked
  by   spaces. It  doesn't matter how many
spaces    you type; one is as good as 100.  The
end of   a line counts as a space.

One   or more   blank lines denote the  end
of  a paragraph.

Since any number of consecutive spaces are treated
like a single one, the formatting of the input
file makes no difference to
      \LaTeX,                % The \LaTeX command generates the LaTeX logo.
but it makes a difference to you.  When you use
\LaTeX, making your input file as easy to read
as possible will be a great help as you write
your document and when you change it.  This sample
file shows how you can add comments to your own input
file.

Because printing is different from typewriting,
there are a number of things that you have to do
differently when preparing an input file than if
you were just typing the document directly.
Quotation marks like
       ``this''
have to be handled specially, as do quotes within
quotes:
       ``\,`this'            % \, separates the double and single quote.
        is what I just
        wrote, not  `that'\,''.

Dashes come in three sizes: an
       intra-word
dash, a medium dash for number ranges like
       1--2,
and a punctuation
       dash---like
this.

A sentence-ending space should be larger than the
space between words within a sentence.  You
sometimes have to type special commands in
conjunction with punctuation characters to get
this right, as in the following sentence.
       Gnats, gnus, etc.\ all  % `\ ' makes an inter-word space.
       begin with G\@.         % \@ marks end-of-sentence punctuation.
You should check the spaces after periods when
reading your output to make sure you haven't
forgotten any special cases.  Generating an
ellipsis
       \ldots\               % `\ ' is needed after `\ldots' because TeX
                             % ignores spaces after command names like \ldots
                             % made from \ + letters.
                             %
                             % Note how a `%' character causes TeX to ignore
                             % the end of the input line, so these blank lines
                             % do not start a new paragraph.
                             %
with the right spacing around the periods requires
a special command.

\LaTeX\ interprets some common characters as
commands, so you must type special commands to
generate them.  These characters include the
following:
       \$ \& \% \# \{ and \}.

In printing, text is usually emphasized with an
       \emph{italic}
type style.

\begin{em}
   A long segment of text can also be emphasized
   in this way.  Text within such a segment can be
   given \emph{additional} emphasis.
\end{em}

It is sometimes necessary to prevent \LaTeX\ from
breaking a line where it might otherwise do so.
This may be at a space, as between the ``Mr.''\ and
``Jones'' in
       ``Mr.~Jones'',        % ~ produces an unbreakable interword space.
or within a word---especially when the word is a
symbol like
       \mbox{\emph{itemnum}}
that makes little sense when hyphenated across
lines.

Footnotes\footnote{This is an example of a footnote.}
pose no problem.

\LaTeX\ is good at typesetting mathematical formulas
like
       \( x-3y + z = 7 \)
or
       \( a_{1} > x^{2n} + y^{2n} > x' \)
or
       \( \ip{A}{B} = \sum_{i} a_{i} b_{i} \).
The spaces you type in a formula are
ignored.  Remember that a letter like
       $x$                   % $ ... $  and  \( ... \)  are equivalent
is a formula when it denotes a mathematical
symbol, and it should be typed as one.

\section{Displayed Text}

Text is displayed by indenting it from the left
margin.  Quotations are commonly displayed.  There
are short quotations
\begin{quote}
   This is a short quotation.  It consists of a
   single paragraph of text.  See how it is formatted.
\end{quote}
and longer ones.
\begin{quotation}
   This is a longer quotation.  It consists of two
   paragraphs of text, neither of which are
   particularly interesting.

   This is the second paragraph of the quotation.  It
   is just as dull as the first paragraph.
\end{quotation}
Another frequently-displayed structure is a list.
The following is an example of an \emph{itemized}
list.
\begin{itemize}
   \item This is the first item of an itemized list.
         Each item in the list is marked with a ``tick''.
         You don't have to worry about what kind of tick
         mark is used.

   \item This is the second item of the list.  It
         contains another list nested inside it.  The inner
         list is an \emph{enumerated} list.
         \begin{enumerate}
            \item This is the first item of an enumerated
                  list that is nested within the itemized list.

            \item This is the second item of the inner list.
                  \LaTeX\ allows you to nest lists deeper than
                  you really should.
         \end{enumerate}
         This is the rest of the second item of the outer
         list.  It is no more interesting than any other
         part of the item.
   \item This is the third item of the list.
\end{itemize}
You can even display poetry.
\begin{verse}
   There is an environment
    for verse \\             % The \\ command separates lines
   Whose features some poets % within a stanza.
   will curse.

                             % One or more blank lines separate stanzas.

   For instead of making\\
   Them do \emph{all} line breaking, \\
   It allows them to put too many words on a line when they'd rather be
   forced to be terse.
\end{verse}

Mathematical formulas may also be displayed.  A
displayed formula
is
one-line long; multiline
formulas require special formatting instructions.
   \[  \ip{\Gamma}{\psi'} = x'' + y^{2} + z_{i}^{n}\]
Don't start a paragraph with a displayed equation,
nor make one a paragraph by itself.

\end{document}               % End of document.

\section{Bench Huge Section 515}
% synthetic bench section 515
% This is a sample LaTeX input file.  (Version of 12 August 2004.)
%
% A '%' character causes TeX to ignore all remaining text on the line,
% and is used for comments like this one.

\documentclass{article}      % Specifies the document class

                             % The preamble begins here.
\title{An Example Document}  % Declares the document's title.
\author{Leslie Lamport}      % Declares the author's name.
\date{January 21, 1994}      % Deleting this command produces today's date.

\newcommand{\ip}[2]{(#1, #2)}
                             % Defines \ip{arg1}{arg2} to mean
                             % (arg1, arg2).

%\newcommand{\ip}[2]{\langle #1 | #2\rangle}
                             % This is an alternative definition of
                             % \ip that is commented out.

\begin{document}             % End of preamble and beginning of text.

\maketitle                   % Produces the title.

This is an example input file.  Comparing it with
the output it generates can show you how to
produce a simple document of your own.

\section{Ordinary Text}      % Produces section heading.  Lower-level
                             % sections are begun with similar
                             % \subsection and \subsubsection commands.

The ends  of words and sentences are marked
  by   spaces. It  doesn't matter how many
spaces    you type; one is as good as 100.  The
end of   a line counts as a space.

One   or more   blank lines denote the  end
of  a paragraph.

Since any number of consecutive spaces are treated
like a single one, the formatting of the input
file makes no difference to
      \LaTeX,                % The \LaTeX command generates the LaTeX logo.
but it makes a difference to you.  When you use
\LaTeX, making your input file as easy to read
as possible will be a great help as you write
your document and when you change it.  This sample
file shows how you can add comments to your own input
file.

Because printing is different from typewriting,
there are a number of things that you have to do
differently when preparing an input file than if
you were just typing the document directly.
Quotation marks like
       ``this''
have to be handled specially, as do quotes within
quotes:
       ``\,`this'            % \, separates the double and single quote.
        is what I just
        wrote, not  `that'\,''.

Dashes come in three sizes: an
       intra-word
dash, a medium dash for number ranges like
       1--2,
and a punctuation
       dash---like
this.

A sentence-ending space should be larger than the
space between words within a sentence.  You
sometimes have to type special commands in
conjunction with punctuation characters to get
this right, as in the following sentence.
       Gnats, gnus, etc.\ all  % `\ ' makes an inter-word space.
       begin with G\@.         % \@ marks end-of-sentence punctuation.
You should check the spaces after periods when
reading your output to make sure you haven't
forgotten any special cases.  Generating an
ellipsis
       \ldots\               % `\ ' is needed after `\ldots' because TeX
                             % ignores spaces after command names like \ldots
                             % made from \ + letters.
                             %
                             % Note how a `%' character causes TeX to ignore
                             % the end of the input line, so these blank lines
                             % do not start a new paragraph.
                             %
with the right spacing around the periods requires
a special command.

\LaTeX\ interprets some common characters as
commands, so you must type special commands to
generate them.  These characters include the
following:
       \$ \& \% \# \{ and \}.

In printing, text is usually emphasized with an
       \emph{italic}
type style.

\begin{em}
   A long segment of text can also be emphasized
   in this way.  Text within such a segment can be
   given \emph{additional} emphasis.
\end{em}

It is sometimes necessary to prevent \LaTeX\ from
breaking a line where it might otherwise do so.
This may be at a space, as between the ``Mr.''\ and
``Jones'' in
       ``Mr.~Jones'',        % ~ produces an unbreakable interword space.
or within a word---especially when the word is a
symbol like
       \mbox{\emph{itemnum}}
that makes little sense when hyphenated across
lines.

Footnotes\footnote{This is an example of a footnote.}
pose no problem.

\LaTeX\ is good at typesetting mathematical formulas
like
       \( x-3y + z = 7 \)
or
       \( a_{1} > x^{2n} + y^{2n} > x' \)
or
       \( \ip{A}{B} = \sum_{i} a_{i} b_{i} \).
The spaces you type in a formula are
ignored.  Remember that a letter like
       $x$                   % $ ... $  and  \( ... \)  are equivalent
is a formula when it denotes a mathematical
symbol, and it should be typed as one.

\section{Displayed Text}

Text is displayed by indenting it from the left
margin.  Quotations are commonly displayed.  There
are short quotations
\begin{quote}
   This is a short quotation.  It consists of a
   single paragraph of text.  See how it is formatted.
\end{quote}
and longer ones.
\begin{quotation}
   This is a longer quotation.  It consists of two
   paragraphs of text, neither of which are
   particularly interesting.

   This is the second paragraph of the quotation.  It
   is just as dull as the first paragraph.
\end{quotation}
Another frequently-displayed structure is a list.
The following is an example of an \emph{itemized}
list.
\begin{itemize}
   \item This is the first item of an itemized list.
         Each item in the list is marked with a ``tick''.
         You don't have to worry about what kind of tick
         mark is used.

   \item This is the second item of the list.  It
         contains another list nested inside it.  The inner
         list is an \emph{enumerated} list.
         \begin{enumerate}
            \item This is the first item of an enumerated
                  list that is nested within the itemized list.

            \item This is the second item of the inner list.
                  \LaTeX\ allows you to nest lists deeper than
                  you really should.
         \end{enumerate}
         This is the rest of the second item of the outer
         list.  It is no more interesting than any other
         part of the item.
   \item This is the third item of the list.
\end{itemize}
You can even display poetry.
\begin{verse}
   There is an environment
    for verse \\             % The \\ command separates lines
   Whose features some poets % within a stanza.
   will curse.

                             % One or more blank lines separate stanzas.

   For instead of making\\
   Them do \emph{all} line breaking, \\
   It allows them to put too many words on a line when they'd rather be
   forced to be terse.
\end{verse}

Mathematical formulas may also be displayed.  A
displayed formula
is
one-line long; multiline
formulas require special formatting instructions.
   \[  \ip{\Gamma}{\psi'} = x'' + y^{2} + z_{i}^{n}\]
Don't start a paragraph with a displayed equation,
nor make one a paragraph by itself.

\end{document}               % End of document.

\section{Bench Huge Section 516}
% synthetic bench section 516
% This is a sample LaTeX input file.  (Version of 12 August 2004.)
%
% A '%' character causes TeX to ignore all remaining text on the line,
% and is used for comments like this one.

\documentclass{article}      % Specifies the document class

                             % The preamble begins here.
\title{An Example Document}  % Declares the document's title.
\author{Leslie Lamport}      % Declares the author's name.
\date{January 21, 1994}      % Deleting this command produces today's date.

\newcommand{\ip}[2]{(#1, #2)}
                             % Defines \ip{arg1}{arg2} to mean
                             % (arg1, arg2).

%\newcommand{\ip}[2]{\langle #1 | #2\rangle}
                             % This is an alternative definition of
                             % \ip that is commented out.

\begin{document}             % End of preamble and beginning of text.

\maketitle                   % Produces the title.

This is an example input file.  Comparing it with
the output it generates can show you how to
produce a simple document of your own.

\section{Ordinary Text}      % Produces section heading.  Lower-level
                             % sections are begun with similar
                             % \subsection and \subsubsection commands.

The ends  of words and sentences are marked
  by   spaces. It  doesn't matter how many
spaces    you type; one is as good as 100.  The
end of   a line counts as a space.

One   or more   blank lines denote the  end
of  a paragraph.

Since any number of consecutive spaces are treated
like a single one, the formatting of the input
file makes no difference to
      \LaTeX,                % The \LaTeX command generates the LaTeX logo.
but it makes a difference to you.  When you use
\LaTeX, making your input file as easy to read
as possible will be a great help as you write
your document and when you change it.  This sample
file shows how you can add comments to your own input
file.

Because printing is different from typewriting,
there are a number of things that you have to do
differently when preparing an input file than if
you were just typing the document directly.
Quotation marks like
       ``this''
have to be handled specially, as do quotes within
quotes:
       ``\,`this'            % \, separates the double and single quote.
        is what I just
        wrote, not  `that'\,''.

Dashes come in three sizes: an
       intra-word
dash, a medium dash for number ranges like
       1--2,
and a punctuation
       dash---like
this.

A sentence-ending space should be larger than the
space between words within a sentence.  You
sometimes have to type special commands in
conjunction with punctuation characters to get
this right, as in the following sentence.
       Gnats, gnus, etc.\ all  % `\ ' makes an inter-word space.
       begin with G\@.         % \@ marks end-of-sentence punctuation.
You should check the spaces after periods when
reading your output to make sure you haven't
forgotten any special cases.  Generating an
ellipsis
       \ldots\               % `\ ' is needed after `\ldots' because TeX
                             % ignores spaces after command names like \ldots
                             % made from \ + letters.
                             %
                             % Note how a `%' character causes TeX to ignore
                             % the end of the input line, so these blank lines
                             % do not start a new paragraph.
                             %
with the right spacing around the periods requires
a special command.

\LaTeX\ interprets some common characters as
commands, so you must type special commands to
generate them.  These characters include the
following:
       \$ \& \% \# \{ and \}.

In printing, text is usually emphasized with an
       \emph{italic}
type style.

\begin{em}
   A long segment of text can also be emphasized
   in this way.  Text within such a segment can be
   given \emph{additional} emphasis.
\end{em}

It is sometimes necessary to prevent \LaTeX\ from
breaking a line where it might otherwise do so.
This may be at a space, as between the ``Mr.''\ and
``Jones'' in
       ``Mr.~Jones'',        % ~ produces an unbreakable interword space.
or within a word---especially when the word is a
symbol like
       \mbox{\emph{itemnum}}
that makes little sense when hyphenated across
lines.

Footnotes\footnote{This is an example of a footnote.}
pose no problem.

\LaTeX\ is good at typesetting mathematical formulas
like
       \( x-3y + z = 7 \)
or
       \( a_{1} > x^{2n} + y^{2n} > x' \)
or
       \( \ip{A}{B} = \sum_{i} a_{i} b_{i} \).
The spaces you type in a formula are
ignored.  Remember that a letter like
       $x$                   % $ ... $  and  \( ... \)  are equivalent
is a formula when it denotes a mathematical
symbol, and it should be typed as one.

\section{Displayed Text}

Text is displayed by indenting it from the left
margin.  Quotations are commonly displayed.  There
are short quotations
\begin{quote}
   This is a short quotation.  It consists of a
   single paragraph of text.  See how it is formatted.
\end{quote}
and longer ones.
\begin{quotation}
   This is a longer quotation.  It consists of two
   paragraphs of text, neither of which are
   particularly interesting.

   This is the second paragraph of the quotation.  It
   is just as dull as the first paragraph.
\end{quotation}
Another frequently-displayed structure is a list.
The following is an example of an \emph{itemized}
list.
\begin{itemize}
   \item This is the first item of an itemized list.
         Each item in the list is marked with a ``tick''.
         You don't have to worry about what kind of tick
         mark is used.

   \item This is the second item of the list.  It
         contains another list nested inside it.  The inner
         list is an \emph{enumerated} list.
         \begin{enumerate}
            \item This is the first item of an enumerated
                  list that is nested within the itemized list.

            \item This is the second item of the inner list.
                  \LaTeX\ allows you to nest lists deeper than
                  you really should.
         \end{enumerate}
         This is the rest of the second item of the outer
         list.  It is no more interesting than any other
         part of the item.
   \item This is the third item of the list.
\end{itemize}
You can even display poetry.
\begin{verse}
   There is an environment
    for verse \\             % The \\ command separates lines
   Whose features some poets % within a stanza.
   will curse.

                             % One or more blank lines separate stanzas.

   For instead of making\\
   Them do \emph{all} line breaking, \\
   It allows them to put too many words on a line when they'd rather be
   forced to be terse.
\end{verse}

Mathematical formulas may also be displayed.  A
displayed formula
is
one-line long; multiline
formulas require special formatting instructions.
   \[  \ip{\Gamma}{\psi'} = x'' + y^{2} + z_{i}^{n}\]
Don't start a paragraph with a displayed equation,
nor make one a paragraph by itself.

\end{document}               % End of document.

\section{Bench Huge Section 517}
% synthetic bench section 517
% This is a sample LaTeX input file.  (Version of 12 August 2004.)
%
% A '%' character causes TeX to ignore all remaining text on the line,
% and is used for comments like this one.

\documentclass{article}      % Specifies the document class

                             % The preamble begins here.
\title{An Example Document}  % Declares the document's title.
\author{Leslie Lamport}      % Declares the author's name.
\date{January 21, 1994}      % Deleting this command produces today's date.

\newcommand{\ip}[2]{(#1, #2)}
                             % Defines \ip{arg1}{arg2} to mean
                             % (arg1, arg2).

%\newcommand{\ip}[2]{\langle #1 | #2\rangle}
                             % This is an alternative definition of
                             % \ip that is commented out.

\begin{document}             % End of preamble and beginning of text.

\maketitle                   % Produces the title.

This is an example input file.  Comparing it with
the output it generates can show you how to
produce a simple document of your own.

\section{Ordinary Text}      % Produces section heading.  Lower-level
                             % sections are begun with similar
                             % \subsection and \subsubsection commands.

The ends  of words and sentences are marked
  by   spaces. It  doesn't matter how many
spaces    you type; one is as good as 100.  The
end of   a line counts as a space.

One   or more   blank lines denote the  end
of  a paragraph.

Since any number of consecutive spaces are treated
like a single one, the formatting of the input
file makes no difference to
      \LaTeX,                % The \LaTeX command generates the LaTeX logo.
but it makes a difference to you.  When you use
\LaTeX, making your input file as easy to read
as possible will be a great help as you write
your document and when you change it.  This sample
file shows how you can add comments to your own input
file.

Because printing is different from typewriting,
there are a number of things that you have to do
differently when preparing an input file than if
you were just typing the document directly.
Quotation marks like
       ``this''
have to be handled specially, as do quotes within
quotes:
       ``\,`this'            % \, separates the double and single quote.
        is what I just
        wrote, not  `that'\,''.

Dashes come in three sizes: an
       intra-word
dash, a medium dash for number ranges like
       1--2,
and a punctuation
       dash---like
this.

A sentence-ending space should be larger than the
space between words within a sentence.  You
sometimes have to type special commands in
conjunction with punctuation characters to get
this right, as in the following sentence.
       Gnats, gnus, etc.\ all  % `\ ' makes an inter-word space.
       begin with G\@.         % \@ marks end-of-sentence punctuation.
You should check the spaces after periods when
reading your output to make sure you haven't
forgotten any special cases.  Generating an
ellipsis
       \ldots\               % `\ ' is needed after `\ldots' because TeX
                             % ignores spaces after command names like \ldots
                             % made from \ + letters.
                             %
                             % Note how a `%' character causes TeX to ignore
                             % the end of the input line, so these blank lines
                             % do not start a new paragraph.
                             %
with the right spacing around the periods requires
a special command.

\LaTeX\ interprets some common characters as
commands, so you must type special commands to
generate them.  These characters include the
following:
       \$ \& \% \# \{ and \}.

In printing, text is usually emphasized with an
       \emph{italic}
type style.

\begin{em}
   A long segment of text can also be emphasized
   in this way.  Text within such a segment can be
   given \emph{additional} emphasis.
\end{em}

It is sometimes necessary to prevent \LaTeX\ from
breaking a line where it might otherwise do so.
This may be at a space, as between the ``Mr.''\ and
``Jones'' in
       ``Mr.~Jones'',        % ~ produces an unbreakable interword space.
or within a word---especially when the word is a
symbol like
       \mbox{\emph{itemnum}}
that makes little sense when hyphenated across
lines.

Footnotes\footnote{This is an example of a footnote.}
pose no problem.

\LaTeX\ is good at typesetting mathematical formulas
like
       \( x-3y + z = 7 \)
or
       \( a_{1} > x^{2n} + y^{2n} > x' \)
or
       \( \ip{A}{B} = \sum_{i} a_{i} b_{i} \).
The spaces you type in a formula are
ignored.  Remember that a letter like
       $x$                   % $ ... $  and  \( ... \)  are equivalent
is a formula when it denotes a mathematical
symbol, and it should be typed as one.

\section{Displayed Text}

Text is displayed by indenting it from the left
margin.  Quotations are commonly displayed.  There
are short quotations
\begin{quote}
   This is a short quotation.  It consists of a
   single paragraph of text.  See how it is formatted.
\end{quote}
and longer ones.
\begin{quotation}
   This is a longer quotation.  It consists of two
   paragraphs of text, neither of which are
   particularly interesting.

   This is the second paragraph of the quotation.  It
   is just as dull as the first paragraph.
\end{quotation}
Another frequently-displayed structure is a list.
The following is an example of an \emph{itemized}
list.
\begin{itemize}
   \item This is the first item of an itemized list.
         Each item in the list is marked with a ``tick''.
         You don't have to worry about what kind of tick
         mark is used.

   \item This is the second item of the list.  It
         contains another list nested inside it.  The inner
         list is an \emph{enumerated} list.
         \begin{enumerate}
            \item This is the first item of an enumerated
                  list that is nested within the itemized list.

            \item This is the second item of the inner list.
                  \LaTeX\ allows you to nest lists deeper than
                  you really should.
         \end{enumerate}
         This is the rest of the second item of the outer
         list.  It is no more interesting than any other
         part of the item.
   \item This is the third item of the list.
\end{itemize}
You can even display poetry.
\begin{verse}
   There is an environment
    for verse \\             % The \\ command separates lines
   Whose features some poets % within a stanza.
   will curse.

                             % One or more blank lines separate stanzas.

   For instead of making\\
   Them do \emph{all} line breaking, \\
   It allows them to put too many words on a line when they'd rather be
   forced to be terse.
\end{verse}

Mathematical formulas may also be displayed.  A
displayed formula
is
one-line long; multiline
formulas require special formatting instructions.
   \[  \ip{\Gamma}{\psi'} = x'' + y^{2} + z_{i}^{n}\]
Don't start a paragraph with a displayed equation,
nor make one a paragraph by itself.

\end{document}               % End of document.

\section{Bench Huge Section 518}
% synthetic bench section 518
% This is a sample LaTeX input file.  (Version of 12 August 2004.)
%
% A '%' character causes TeX to ignore all remaining text on the line,
% and is used for comments like this one.

\documentclass{article}      % Specifies the document class

                             % The preamble begins here.
\title{An Example Document}  % Declares the document's title.
\author{Leslie Lamport}      % Declares the author's name.
\date{January 21, 1994}      % Deleting this command produces today's date.

\newcommand{\ip}[2]{(#1, #2)}
                             % Defines \ip{arg1}{arg2} to mean
                             % (arg1, arg2).

%\newcommand{\ip}[2]{\langle #1 | #2\rangle}
                             % This is an alternative definition of
                             % \ip that is commented out.

\begin{document}             % End of preamble and beginning of text.

\maketitle                   % Produces the title.

This is an example input file.  Comparing it with
the output it generates can show you how to
produce a simple document of your own.

\section{Ordinary Text}      % Produces section heading.  Lower-level
                             % sections are begun with similar
                             % \subsection and \subsubsection commands.

The ends  of words and sentences are marked
  by   spaces. It  doesn't matter how many
spaces    you type; one is as good as 100.  The
end of   a line counts as a space.

One   or more   blank lines denote the  end
of  a paragraph.

Since any number of consecutive spaces are treated
like a single one, the formatting of the input
file makes no difference to
      \LaTeX,                % The \LaTeX command generates the LaTeX logo.
but it makes a difference to you.  When you use
\LaTeX, making your input file as easy to read
as possible will be a great help as you write
your document and when you change it.  This sample
file shows how you can add comments to your own input
file.

Because printing is different from typewriting,
there are a number of things that you have to do
differently when preparing an input file than if
you were just typing the document directly.
Quotation marks like
       ``this''
have to be handled specially, as do quotes within
quotes:
       ``\,`this'            % \, separates the double and single quote.
        is what I just
        wrote, not  `that'\,''.

Dashes come in three sizes: an
       intra-word
dash, a medium dash for number ranges like
       1--2,
and a punctuation
       dash---like
this.

A sentence-ending space should be larger than the
space between words within a sentence.  You
sometimes have to type special commands in
conjunction with punctuation characters to get
this right, as in the following sentence.
       Gnats, gnus, etc.\ all  % `\ ' makes an inter-word space.
       begin with G\@.         % \@ marks end-of-sentence punctuation.
You should check the spaces after periods when
reading your output to make sure you haven't
forgotten any special cases.  Generating an
ellipsis
       \ldots\               % `\ ' is needed after `\ldots' because TeX
                             % ignores spaces after command names like \ldots
                             % made from \ + letters.
                             %
                             % Note how a `%' character causes TeX to ignore
                             % the end of the input line, so these blank lines
                             % do not start a new paragraph.
                             %
with the right spacing around the periods requires
a special command.

\LaTeX\ interprets some common characters as
commands, so you must type special commands to
generate them.  These characters include the
following:
       \$ \& \% \# \{ and \}.

In printing, text is usually emphasized with an
       \emph{italic}
type style.

\begin{em}
   A long segment of text can also be emphasized
   in this way.  Text within such a segment can be
   given \emph{additional} emphasis.
\end{em}

It is sometimes necessary to prevent \LaTeX\ from
breaking a line where it might otherwise do so.
This may be at a space, as between the ``Mr.''\ and
``Jones'' in
       ``Mr.~Jones'',        % ~ produces an unbreakable interword space.
or within a word---especially when the word is a
symbol like
       \mbox{\emph{itemnum}}
that makes little sense when hyphenated across
lines.

Footnotes\footnote{This is an example of a footnote.}
pose no problem.

\LaTeX\ is good at typesetting mathematical formulas
like
       \( x-3y + z = 7 \)
or
       \( a_{1} > x^{2n} + y^{2n} > x' \)
or
       \( \ip{A}{B} = \sum_{i} a_{i} b_{i} \).
The spaces you type in a formula are
ignored.  Remember that a letter like
       $x$                   % $ ... $  and  \( ... \)  are equivalent
is a formula when it denotes a mathematical
symbol, and it should be typed as one.

\section{Displayed Text}

Text is displayed by indenting it from the left
margin.  Quotations are commonly displayed.  There
are short quotations
\begin{quote}
   This is a short quotation.  It consists of a
   single paragraph of text.  See how it is formatted.
\end{quote}
and longer ones.
\begin{quotation}
   This is a longer quotation.  It consists of two
   paragraphs of text, neither of which are
   particularly interesting.

   This is the second paragraph of the quotation.  It
   is just as dull as the first paragraph.
\end{quotation}
Another frequently-displayed structure is a list.
The following is an example of an \emph{itemized}
list.
\begin{itemize}
   \item This is the first item of an itemized list.
         Each item in the list is marked with a ``tick''.
         You don't have to worry about what kind of tick
         mark is used.

   \item This is the second item of the list.  It
         contains another list nested inside it.  The inner
         list is an \emph{enumerated} list.
         \begin{enumerate}
            \item This is the first item of an enumerated
                  list that is nested within the itemized list.

            \item This is the second item of the inner list.
                  \LaTeX\ allows you to nest lists deeper than
                  you really should.
         \end{enumerate}
         This is the rest of the second item of the outer
         list.  It is no more interesting than any other
         part of the item.
   \item This is the third item of the list.
\end{itemize}
You can even display poetry.
\begin{verse}
   There is an environment
    for verse \\             % The \\ command separates lines
   Whose features some poets % within a stanza.
   will curse.

                             % One or more blank lines separate stanzas.

   For instead of making\\
   Them do \emph{all} line breaking, \\
   It allows them to put too many words on a line when they'd rather be
   forced to be terse.
\end{verse}

Mathematical formulas may also be displayed.  A
displayed formula
is
one-line long; multiline
formulas require special formatting instructions.
   \[  \ip{\Gamma}{\psi'} = x'' + y^{2} + z_{i}^{n}\]
Don't start a paragraph with a displayed equation,
nor make one a paragraph by itself.

\end{document}               % End of document.

\section{Bench Huge Section 519}
% synthetic bench section 519
% This is a sample LaTeX input file.  (Version of 12 August 2004.)
%
% A '%' character causes TeX to ignore all remaining text on the line,
% and is used for comments like this one.

\documentclass{article}      % Specifies the document class

                             % The preamble begins here.
\title{An Example Document}  % Declares the document's title.
\author{Leslie Lamport}      % Declares the author's name.
\date{January 21, 1994}      % Deleting this command produces today's date.

\newcommand{\ip}[2]{(#1, #2)}
                             % Defines \ip{arg1}{arg2} to mean
                             % (arg1, arg2).

%\newcommand{\ip}[2]{\langle #1 | #2\rangle}
                             % This is an alternative definition of
                             % \ip that is commented out.

\begin{document}             % End of preamble and beginning of text.

\maketitle                   % Produces the title.

This is an example input file.  Comparing it with
the output it generates can show you how to
produce a simple document of your own.

\section{Ordinary Text}      % Produces section heading.  Lower-level
                             % sections are begun with similar
                             % \subsection and \subsubsection commands.

The ends  of words and sentences are marked
  by   spaces. It  doesn't matter how many
spaces    you type; one is as good as 100.  The
end of   a line counts as a space.

One   or more   blank lines denote the  end
of  a paragraph.

Since any number of consecutive spaces are treated
like a single one, the formatting of the input
file makes no difference to
      \LaTeX,                % The \LaTeX command generates the LaTeX logo.
but it makes a difference to you.  When you use
\LaTeX, making your input file as easy to read
as possible will be a great help as you write
your document and when you change it.  This sample
file shows how you can add comments to your own input
file.

Because printing is different from typewriting,
there are a number of things that you have to do
differently when preparing an input file than if
you were just typing the document directly.
Quotation marks like
       ``this''
have to be handled specially, as do quotes within
quotes:
       ``\,`this'            % \, separates the double and single quote.
        is what I just
        wrote, not  `that'\,''.

Dashes come in three sizes: an
       intra-word
dash, a medium dash for number ranges like
       1--2,
and a punctuation
       dash---like
this.

A sentence-ending space should be larger than the
space between words within a sentence.  You
sometimes have to type special commands in
conjunction with punctuation characters to get
this right, as in the following sentence.
       Gnats, gnus, etc.\ all  % `\ ' makes an inter-word space.
       begin with G\@.         % \@ marks end-of-sentence punctuation.
You should check the spaces after periods when
reading your output to make sure you haven't
forgotten any special cases.  Generating an
ellipsis
       \ldots\               % `\ ' is needed after `\ldots' because TeX
                             % ignores spaces after command names like \ldots
                             % made from \ + letters.
                             %
                             % Note how a `%' character causes TeX to ignore
                             % the end of the input line, so these blank lines
                             % do not start a new paragraph.
                             %
with the right spacing around the periods requires
a special command.

\LaTeX\ interprets some common characters as
commands, so you must type special commands to
generate them.  These characters include the
following:
       \$ \& \% \# \{ and \}.

In printing, text is usually emphasized with an
       \emph{italic}
type style.

\begin{em}
   A long segment of text can also be emphasized
   in this way.  Text within such a segment can be
   given \emph{additional} emphasis.
\end{em}

It is sometimes necessary to prevent \LaTeX\ from
breaking a line where it might otherwise do so.
This may be at a space, as between the ``Mr.''\ and
``Jones'' in
       ``Mr.~Jones'',        % ~ produces an unbreakable interword space.
or within a word---especially when the word is a
symbol like
       \mbox{\emph{itemnum}}
that makes little sense when hyphenated across
lines.

Footnotes\footnote{This is an example of a footnote.}
pose no problem.

\LaTeX\ is good at typesetting mathematical formulas
like
       \( x-3y + z = 7 \)
or
       \( a_{1} > x^{2n} + y^{2n} > x' \)
or
       \( \ip{A}{B} = \sum_{i} a_{i} b_{i} \).
The spaces you type in a formula are
ignored.  Remember that a letter like
       $x$                   % $ ... $  and  \( ... \)  are equivalent
is a formula when it denotes a mathematical
symbol, and it should be typed as one.

\section{Displayed Text}

Text is displayed by indenting it from the left
margin.  Quotations are commonly displayed.  There
are short quotations
\begin{quote}
   This is a short quotation.  It consists of a
   single paragraph of text.  See how it is formatted.
\end{quote}
and longer ones.
\begin{quotation}
   This is a longer quotation.  It consists of two
   paragraphs of text, neither of which are
   particularly interesting.

   This is the second paragraph of the quotation.  It
   is just as dull as the first paragraph.
\end{quotation}
Another frequently-displayed structure is a list.
The following is an example of an \emph{itemized}
list.
\begin{itemize}
   \item This is the first item of an itemized list.
         Each item in the list is marked with a ``tick''.
         You don't have to worry about what kind of tick
         mark is used.

   \item This is the second item of the list.  It
         contains another list nested inside it.  The inner
         list is an \emph{enumerated} list.
         \begin{enumerate}
            \item This is the first item of an enumerated
                  list that is nested within the itemized list.

            \item This is the second item of the inner list.
                  \LaTeX\ allows you to nest lists deeper than
                  you really should.
         \end{enumerate}
         This is the rest of the second item of the outer
         list.  It is no more interesting than any other
         part of the item.
   \item This is the third item of the list.
\end{itemize}
You can even display poetry.
\begin{verse}
   There is an environment
    for verse \\             % The \\ command separates lines
   Whose features some poets % within a stanza.
   will curse.

                             % One or more blank lines separate stanzas.

   For instead of making\\
   Them do \emph{all} line breaking, \\
   It allows them to put too many words on a line when they'd rather be
   forced to be terse.
\end{verse}

Mathematical formulas may also be displayed.  A
displayed formula
is
one-line long; multiline
formulas require special formatting instructions.
   \[  \ip{\Gamma}{\psi'} = x'' + y^{2} + z_{i}^{n}\]
Don't start a paragraph with a displayed equation,
nor make one a paragraph by itself.

\end{document}               % End of document.

\section{Bench Huge Section 520}
% synthetic bench section 520
% This is a sample LaTeX input file.  (Version of 12 August 2004.)
%
% A '%' character causes TeX to ignore all remaining text on the line,
% and is used for comments like this one.

\documentclass{article}      % Specifies the document class

                             % The preamble begins here.
\title{An Example Document}  % Declares the document's title.
\author{Leslie Lamport}      % Declares the author's name.
\date{January 21, 1994}      % Deleting this command produces today's date.

\newcommand{\ip}[2]{(#1, #2)}
                             % Defines \ip{arg1}{arg2} to mean
                             % (arg1, arg2).

%\newcommand{\ip}[2]{\langle #1 | #2\rangle}
                             % This is an alternative definition of
                             % \ip that is commented out.

\begin{document}             % End of preamble and beginning of text.

\maketitle                   % Produces the title.

This is an example input file.  Comparing it with
the output it generates can show you how to
produce a simple document of your own.

\section{Ordinary Text}      % Produces section heading.  Lower-level
                             % sections are begun with similar
                             % \subsection and \subsubsection commands.

The ends  of words and sentences are marked
  by   spaces. It  doesn't matter how many
spaces    you type; one is as good as 100.  The
end of   a line counts as a space.

One   or more   blank lines denote the  end
of  a paragraph.

Since any number of consecutive spaces are treated
like a single one, the formatting of the input
file makes no difference to
      \LaTeX,                % The \LaTeX command generates the LaTeX logo.
but it makes a difference to you.  When you use
\LaTeX, making your input file as easy to read
as possible will be a great help as you write
your document and when you change it.  This sample
file shows how you can add comments to your own input
file.

Because printing is different from typewriting,
there are a number of things that you have to do
differently when preparing an input file than if
you were just typing the document directly.
Quotation marks like
       ``this''
have to be handled specially, as do quotes within
quotes:
       ``\,`this'            % \, separates the double and single quote.
        is what I just
        wrote, not  `that'\,''.

Dashes come in three sizes: an
       intra-word
dash, a medium dash for number ranges like
       1--2,
and a punctuation
       dash---like
this.

A sentence-ending space should be larger than the
space between words within a sentence.  You
sometimes have to type special commands in
conjunction with punctuation characters to get
this right, as in the following sentence.
       Gnats, gnus, etc.\ all  % `\ ' makes an inter-word space.
       begin with G\@.         % \@ marks end-of-sentence punctuation.
You should check the spaces after periods when
reading your output to make sure you haven't
forgotten any special cases.  Generating an
ellipsis
       \ldots\               % `\ ' is needed after `\ldots' because TeX
                             % ignores spaces after command names like \ldots
                             % made from \ + letters.
                             %
                             % Note how a `%' character causes TeX to ignore
                             % the end of the input line, so these blank lines
                             % do not start a new paragraph.
                             %
with the right spacing around the periods requires
a special command.

\LaTeX\ interprets some common characters as
commands, so you must type special commands to
generate them.  These characters include the
following:
       \$ \& \% \# \{ and \}.

In printing, text is usually emphasized with an
       \emph{italic}
type style.

\begin{em}
   A long segment of text can also be emphasized
   in this way.  Text within such a segment can be
   given \emph{additional} emphasis.
\end{em}

It is sometimes necessary to prevent \LaTeX\ from
breaking a line where it might otherwise do so.
This may be at a space, as between the ``Mr.''\ and
``Jones'' in
       ``Mr.~Jones'',        % ~ produces an unbreakable interword space.
or within a word---especially when the word is a
symbol like
       \mbox{\emph{itemnum}}
that makes little sense when hyphenated across
lines.

Footnotes\footnote{This is an example of a footnote.}
pose no problem.

\LaTeX\ is good at typesetting mathematical formulas
like
       \( x-3y + z = 7 \)
or
       \( a_{1} > x^{2n} + y^{2n} > x' \)
or
       \( \ip{A}{B} = \sum_{i} a_{i} b_{i} \).
The spaces you type in a formula are
ignored.  Remember that a letter like
       $x$                   % $ ... $  and  \( ... \)  are equivalent
is a formula when it denotes a mathematical
symbol, and it should be typed as one.

\section{Displayed Text}

Text is displayed by indenting it from the left
margin.  Quotations are commonly displayed.  There
are short quotations
\begin{quote}
   This is a short quotation.  It consists of a
   single paragraph of text.  See how it is formatted.
\end{quote}
and longer ones.
\begin{quotation}
   This is a longer quotation.  It consists of two
   paragraphs of text, neither of which are
   particularly interesting.

   This is the second paragraph of the quotation.  It
   is just as dull as the first paragraph.
\end{quotation}
Another frequently-displayed structure is a list.
The following is an example of an \emph{itemized}
list.
\begin{itemize}
   \item This is the first item of an itemized list.
         Each item in the list is marked with a ``tick''.
         You don't have to worry about what kind of tick
         mark is used.

   \item This is the second item of the list.  It
         contains another list nested inside it.  The inner
         list is an \emph{enumerated} list.
         \begin{enumerate}
            \item This is the first item of an enumerated
                  list that is nested within the itemized list.

            \item This is the second item of the inner list.
                  \LaTeX\ allows you to nest lists deeper than
                  you really should.
         \end{enumerate}
         This is the rest of the second item of the outer
         list.  It is no more interesting than any other
         part of the item.
   \item This is the third item of the list.
\end{itemize}
You can even display poetry.
\begin{verse}
   There is an environment
    for verse \\             % The \\ command separates lines
   Whose features some poets % within a stanza.
   will curse.

                             % One or more blank lines separate stanzas.

   For instead of making\\
   Them do \emph{all} line breaking, \\
   It allows them to put too many words on a line when they'd rather be
   forced to be terse.
\end{verse}

Mathematical formulas may also be displayed.  A
displayed formula
is
one-line long; multiline
formulas require special formatting instructions.
   \[  \ip{\Gamma}{\psi'} = x'' + y^{2} + z_{i}^{n}\]
Don't start a paragraph with a displayed equation,
nor make one a paragraph by itself.

\end{document}               % End of document.

\section{Bench Huge Section 521}
% synthetic bench section 521
% This is a sample LaTeX input file.  (Version of 12 August 2004.)
%
% A '%' character causes TeX to ignore all remaining text on the line,
% and is used for comments like this one.

\documentclass{article}      % Specifies the document class

                             % The preamble begins here.
\title{An Example Document}  % Declares the document's title.
\author{Leslie Lamport}      % Declares the author's name.
\date{January 21, 1994}      % Deleting this command produces today's date.

\newcommand{\ip}[2]{(#1, #2)}
                             % Defines \ip{arg1}{arg2} to mean
                             % (arg1, arg2).

%\newcommand{\ip}[2]{\langle #1 | #2\rangle}
                             % This is an alternative definition of
                             % \ip that is commented out.

\begin{document}             % End of preamble and beginning of text.

\maketitle                   % Produces the title.

This is an example input file.  Comparing it with
the output it generates can show you how to
produce a simple document of your own.

\section{Ordinary Text}      % Produces section heading.  Lower-level
                             % sections are begun with similar
                             % \subsection and \subsubsection commands.

The ends  of words and sentences are marked
  by   spaces. It  doesn't matter how many
spaces    you type; one is as good as 100.  The
end of   a line counts as a space.

One   or more   blank lines denote the  end
of  a paragraph.

Since any number of consecutive spaces are treated
like a single one, the formatting of the input
file makes no difference to
      \LaTeX,                % The \LaTeX command generates the LaTeX logo.
but it makes a difference to you.  When you use
\LaTeX, making your input file as easy to read
as possible will be a great help as you write
your document and when you change it.  This sample
file shows how you can add comments to your own input
file.

Because printing is different from typewriting,
there are a number of things that you have to do
differently when preparing an input file than if
you were just typing the document directly.
Quotation marks like
       ``this''
have to be handled specially, as do quotes within
quotes:
       ``\,`this'            % \, separates the double and single quote.
        is what I just
        wrote, not  `that'\,''.

Dashes come in three sizes: an
       intra-word
dash, a medium dash for number ranges like
       1--2,
and a punctuation
       dash---like
this.

A sentence-ending space should be larger than the
space between words within a sentence.  You
sometimes have to type special commands in
conjunction with punctuation characters to get
this right, as in the following sentence.
       Gnats, gnus, etc.\ all  % `\ ' makes an inter-word space.
       begin with G\@.         % \@ marks end-of-sentence punctuation.
You should check the spaces after periods when
reading your output to make sure you haven't
forgotten any special cases.  Generating an
ellipsis
       \ldots\               % `\ ' is needed after `\ldots' because TeX
                             % ignores spaces after command names like \ldots
                             % made from \ + letters.
                             %
                             % Note how a `%' character causes TeX to ignore
                             % the end of the input line, so these blank lines
                             % do not start a new paragraph.
                             %
with the right spacing around the periods requires
a special command.

\LaTeX\ interprets some common characters as
commands, so you must type special commands to
generate them.  These characters include the
following:
       \$ \& \% \# \{ and \}.

In printing, text is usually emphasized with an
       \emph{italic}
type style.

\begin{em}
   A long segment of text can also be emphasized
   in this way.  Text within such a segment can be
   given \emph{additional} emphasis.
\end{em}

It is sometimes necessary to prevent \LaTeX\ from
breaking a line where it might otherwise do so.
This may be at a space, as between the ``Mr.''\ and
``Jones'' in
       ``Mr.~Jones'',        % ~ produces an unbreakable interword space.
or within a word---especially when the word is a
symbol like
       \mbox{\emph{itemnum}}
that makes little sense when hyphenated across
lines.

Footnotes\footnote{This is an example of a footnote.}
pose no problem.

\LaTeX\ is good at typesetting mathematical formulas
like
       \( x-3y + z = 7 \)
or
       \( a_{1} > x^{2n} + y^{2n} > x' \)
or
       \( \ip{A}{B} = \sum_{i} a_{i} b_{i} \).
The spaces you type in a formula are
ignored.  Remember that a letter like
       $x$                   % $ ... $  and  \( ... \)  are equivalent
is a formula when it denotes a mathematical
symbol, and it should be typed as one.

\section{Displayed Text}

Text is displayed by indenting it from the left
margin.  Quotations are commonly displayed.  There
are short quotations
\begin{quote}
   This is a short quotation.  It consists of a
   single paragraph of text.  See how it is formatted.
\end{quote}
and longer ones.
\begin{quotation}
   This is a longer quotation.  It consists of two
   paragraphs of text, neither of which are
   particularly interesting.

   This is the second paragraph of the quotation.  It
   is just as dull as the first paragraph.
\end{quotation}
Another frequently-displayed structure is a list.
The following is an example of an \emph{itemized}
list.
\begin{itemize}
   \item This is the first item of an itemized list.
         Each item in the list is marked with a ``tick''.
         You don't have to worry about what kind of tick
         mark is used.

   \item This is the second item of the list.  It
         contains another list nested inside it.  The inner
         list is an \emph{enumerated} list.
         \begin{enumerate}
            \item This is the first item of an enumerated
                  list that is nested within the itemized list.

            \item This is the second item of the inner list.
                  \LaTeX\ allows you to nest lists deeper than
                  you really should.
         \end{enumerate}
         This is the rest of the second item of the outer
         list.  It is no more interesting than any other
         part of the item.
   \item This is the third item of the list.
\end{itemize}
You can even display poetry.
\begin{verse}
   There is an environment
    for verse \\             % The \\ command separates lines
   Whose features some poets % within a stanza.
   will curse.

                             % One or more blank lines separate stanzas.

   For instead of making\\
   Them do \emph{all} line breaking, \\
   It allows them to put too many words on a line when they'd rather be
   forced to be terse.
\end{verse}

Mathematical formulas may also be displayed.  A
displayed formula
is
one-line long; multiline
formulas require special formatting instructions.
   \[  \ip{\Gamma}{\psi'} = x'' + y^{2} + z_{i}^{n}\]
Don't start a paragraph with a displayed equation,
nor make one a paragraph by itself.

\end{document}               % End of document.

\section{Bench Huge Section 522}
% synthetic bench section 522
% This is a sample LaTeX input file.  (Version of 12 August 2004.)
%
% A '%' character causes TeX to ignore all remaining text on the line,
% and is used for comments like this one.

\documentclass{article}      % Specifies the document class

                             % The preamble begins here.
\title{An Example Document}  % Declares the document's title.
\author{Leslie Lamport}      % Declares the author's name.
\date{January 21, 1994}      % Deleting this command produces today's date.

\newcommand{\ip}[2]{(#1, #2)}
                             % Defines \ip{arg1}{arg2} to mean
                             % (arg1, arg2).

%\newcommand{\ip}[2]{\langle #1 | #2\rangle}
                             % This is an alternative definition of
                             % \ip that is commented out.

\begin{document}             % End of preamble and beginning of text.

\maketitle                   % Produces the title.

This is an example input file.  Comparing it with
the output it generates can show you how to
produce a simple document of your own.

\section{Ordinary Text}      % Produces section heading.  Lower-level
                             % sections are begun with similar
                             % \subsection and \subsubsection commands.

The ends  of words and sentences are marked
  by   spaces. It  doesn't matter how many
spaces    you type; one is as good as 100.  The
end of   a line counts as a space.

One   or more   blank lines denote the  end
of  a paragraph.

Since any number of consecutive spaces are treated
like a single one, the formatting of the input
file makes no difference to
      \LaTeX,                % The \LaTeX command generates the LaTeX logo.
but it makes a difference to you.  When you use
\LaTeX, making your input file as easy to read
as possible will be a great help as you write
your document and when you change it.  This sample
file shows how you can add comments to your own input
file.

Because printing is different from typewriting,
there are a number of things that you have to do
differently when preparing an input file than if
you were just typing the document directly.
Quotation marks like
       ``this''
have to be handled specially, as do quotes within
quotes:
       ``\,`this'            % \, separates the double and single quote.
        is what I just
        wrote, not  `that'\,''.

Dashes come in three sizes: an
       intra-word
dash, a medium dash for number ranges like
       1--2,
and a punctuation
       dash---like
this.

A sentence-ending space should be larger than the
space between words within a sentence.  You
sometimes have to type special commands in
conjunction with punctuation characters to get
this right, as in the following sentence.
       Gnats, gnus, etc.\ all  % `\ ' makes an inter-word space.
       begin with G\@.         % \@ marks end-of-sentence punctuation.
You should check the spaces after periods when
reading your output to make sure you haven't
forgotten any special cases.  Generating an
ellipsis
       \ldots\               % `\ ' is needed after `\ldots' because TeX
                             % ignores spaces after command names like \ldots
                             % made from \ + letters.
                             %
                             % Note how a `%' character causes TeX to ignore
                             % the end of the input line, so these blank lines
                             % do not start a new paragraph.
                             %
with the right spacing around the periods requires
a special command.

\LaTeX\ interprets some common characters as
commands, so you must type special commands to
generate them.  These characters include the
following:
       \$ \& \% \# \{ and \}.

In printing, text is usually emphasized with an
       \emph{italic}
type style.

\begin{em}
   A long segment of text can also be emphasized
   in this way.  Text within such a segment can be
   given \emph{additional} emphasis.
\end{em}

It is sometimes necessary to prevent \LaTeX\ from
breaking a line where it might otherwise do so.
This may be at a space, as between the ``Mr.''\ and
``Jones'' in
       ``Mr.~Jones'',        % ~ produces an unbreakable interword space.
or within a word---especially when the word is a
symbol like
       \mbox{\emph{itemnum}}
that makes little sense when hyphenated across
lines.

Footnotes\footnote{This is an example of a footnote.}
pose no problem.

\LaTeX\ is good at typesetting mathematical formulas
like
       \( x-3y + z = 7 \)
or
       \( a_{1} > x^{2n} + y^{2n} > x' \)
or
       \( \ip{A}{B} = \sum_{i} a_{i} b_{i} \).
The spaces you type in a formula are
ignored.  Remember that a letter like
       $x$                   % $ ... $  and  \( ... \)  are equivalent
is a formula when it denotes a mathematical
symbol, and it should be typed as one.

\section{Displayed Text}

Text is displayed by indenting it from the left
margin.  Quotations are commonly displayed.  There
are short quotations
\begin{quote}
   This is a short quotation.  It consists of a
   single paragraph of text.  See how it is formatted.
\end{quote}
and longer ones.
\begin{quotation}
   This is a longer quotation.  It consists of two
   paragraphs of text, neither of which are
   particularly interesting.

   This is the second paragraph of the quotation.  It
   is just as dull as the first paragraph.
\end{quotation}
Another frequently-displayed structure is a list.
The following is an example of an \emph{itemized}
list.
\begin{itemize}
   \item This is the first item of an itemized list.
         Each item in the list is marked with a ``tick''.
         You don't have to worry about what kind of tick
         mark is used.

   \item This is the second item of the list.  It
         contains another list nested inside it.  The inner
         list is an \emph{enumerated} list.
         \begin{enumerate}
            \item This is the first item of an enumerated
                  list that is nested within the itemized list.

            \item This is the second item of the inner list.
                  \LaTeX\ allows you to nest lists deeper than
                  you really should.
         \end{enumerate}
         This is the rest of the second item of the outer
         list.  It is no more interesting than any other
         part of the item.
   \item This is the third item of the list.
\end{itemize}
You can even display poetry.
\begin{verse}
   There is an environment
    for verse \\             % The \\ command separates lines
   Whose features some poets % within a stanza.
   will curse.

                             % One or more blank lines separate stanzas.

   For instead of making\\
   Them do \emph{all} line breaking, \\
   It allows them to put too many words on a line when they'd rather be
   forced to be terse.
\end{verse}

Mathematical formulas may also be displayed.  A
displayed formula
is
one-line long; multiline
formulas require special formatting instructions.
   \[  \ip{\Gamma}{\psi'} = x'' + y^{2} + z_{i}^{n}\]
Don't start a paragraph with a displayed equation,
nor make one a paragraph by itself.

\end{document}               % End of document.

\section{Bench Huge Section 523}
% synthetic bench section 523
% This is a sample LaTeX input file.  (Version of 12 August 2004.)
%
% A '%' character causes TeX to ignore all remaining text on the line,
% and is used for comments like this one.

\documentclass{article}      % Specifies the document class

                             % The preamble begins here.
\title{An Example Document}  % Declares the document's title.
\author{Leslie Lamport}      % Declares the author's name.
\date{January 21, 1994}      % Deleting this command produces today's date.

\newcommand{\ip}[2]{(#1, #2)}
                             % Defines \ip{arg1}{arg2} to mean
                             % (arg1, arg2).

%\newcommand{\ip}[2]{\langle #1 | #2\rangle}
                             % This is an alternative definition of
                             % \ip that is commented out.

\begin{document}             % End of preamble and beginning of text.

\maketitle                   % Produces the title.

This is an example input file.  Comparing it with
the output it generates can show you how to
produce a simple document of your own.

\section{Ordinary Text}      % Produces section heading.  Lower-level
                             % sections are begun with similar
                             % \subsection and \subsubsection commands.

The ends  of words and sentences are marked
  by   spaces. It  doesn't matter how many
spaces    you type; one is as good as 100.  The
end of   a line counts as a space.

One   or more   blank lines denote the  end
of  a paragraph.

Since any number of consecutive spaces are treated
like a single one, the formatting of the input
file makes no difference to
      \LaTeX,                % The \LaTeX command generates the LaTeX logo.
but it makes a difference to you.  When you use
\LaTeX, making your input file as easy to read
as possible will be a great help as you write
your document and when you change it.  This sample
file shows how you can add comments to your own input
file.

Because printing is different from typewriting,
there are a number of things that you have to do
differently when preparing an input file than if
you were just typing the document directly.
Quotation marks like
       ``this''
have to be handled specially, as do quotes within
quotes:
       ``\,`this'            % \, separates the double and single quote.
        is what I just
        wrote, not  `that'\,''.

Dashes come in three sizes: an
       intra-word
dash, a medium dash for number ranges like
       1--2,
and a punctuation
       dash---like
this.

A sentence-ending space should be larger than the
space between words within a sentence.  You
sometimes have to type special commands in
conjunction with punctuation characters to get
this right, as in the following sentence.
       Gnats, gnus, etc.\ all  % `\ ' makes an inter-word space.
       begin with G\@.         % \@ marks end-of-sentence punctuation.
You should check the spaces after periods when
reading your output to make sure you haven't
forgotten any special cases.  Generating an
ellipsis
       \ldots\               % `\ ' is needed after `\ldots' because TeX
                             % ignores spaces after command names like \ldots
                             % made from \ + letters.
                             %
                             % Note how a `%' character causes TeX to ignore
                             % the end of the input line, so these blank lines
                             % do not start a new paragraph.
                             %
with the right spacing around the periods requires
a special command.

\LaTeX\ interprets some common characters as
commands, so you must type special commands to
generate them.  These characters include the
following:
       \$ \& \% \# \{ and \}.

In printing, text is usually emphasized with an
       \emph{italic}
type style.

\begin{em}
   A long segment of text can also be emphasized
   in this way.  Text within such a segment can be
   given \emph{additional} emphasis.
\end{em}

It is sometimes necessary to prevent \LaTeX\ from
breaking a line where it might otherwise do so.
This may be at a space, as between the ``Mr.''\ and
``Jones'' in
       ``Mr.~Jones'',        % ~ produces an unbreakable interword space.
or within a word---especially when the word is a
symbol like
       \mbox{\emph{itemnum}}
that makes little sense when hyphenated across
lines.

Footnotes\footnote{This is an example of a footnote.}
pose no problem.

\LaTeX\ is good at typesetting mathematical formulas
like
       \( x-3y + z = 7 \)
or
       \( a_{1} > x^{2n} + y^{2n} > x' \)
or
       \( \ip{A}{B} = \sum_{i} a_{i} b_{i} \).
The spaces you type in a formula are
ignored.  Remember that a letter like
       $x$                   % $ ... $  and  \( ... \)  are equivalent
is a formula when it denotes a mathematical
symbol, and it should be typed as one.

\section{Displayed Text}

Text is displayed by indenting it from the left
margin.  Quotations are commonly displayed.  There
are short quotations
\begin{quote}
   This is a short quotation.  It consists of a
   single paragraph of text.  See how it is formatted.
\end{quote}
and longer ones.
\begin{quotation}
   This is a longer quotation.  It consists of two
   paragraphs of text, neither of which are
   particularly interesting.

   This is the second paragraph of the quotation.  It
   is just as dull as the first paragraph.
\end{quotation}
Another frequently-displayed structure is a list.
The following is an example of an \emph{itemized}
list.
\begin{itemize}
   \item This is the first item of an itemized list.
         Each item in the list is marked with a ``tick''.
         You don't have to worry about what kind of tick
         mark is used.

   \item This is the second item of the list.  It
         contains another list nested inside it.  The inner
         list is an \emph{enumerated} list.
         \begin{enumerate}
            \item This is the first item of an enumerated
                  list that is nested within the itemized list.

            \item This is the second item of the inner list.
                  \LaTeX\ allows you to nest lists deeper than
                  you really should.
         \end{enumerate}
         This is the rest of the second item of the outer
         list.  It is no more interesting than any other
         part of the item.
   \item This is the third item of the list.
\end{itemize}
You can even display poetry.
\begin{verse}
   There is an environment
    for verse \\             % The \\ command separates lines
   Whose features some poets % within a stanza.
   will curse.

                             % One or more blank lines separate stanzas.

   For instead of making\\
   Them do \emph{all} line breaking, \\
   It allows them to put too many words on a line when they'd rather be
   forced to be terse.
\end{verse}

Mathematical formulas may also be displayed.  A
displayed formula
is
one-line long; multiline
formulas require special formatting instructions.
   \[  \ip{\Gamma}{\psi'} = x'' + y^{2} + z_{i}^{n}\]
Don't start a paragraph with a displayed equation,
nor make one a paragraph by itself.

\end{document}               % End of document.

\section{Bench Huge Section 524}
% synthetic bench section 524
% This is a sample LaTeX input file.  (Version of 12 August 2004.)
%
% A '%' character causes TeX to ignore all remaining text on the line,
% and is used for comments like this one.

\documentclass{article}      % Specifies the document class

                             % The preamble begins here.
\title{An Example Document}  % Declares the document's title.
\author{Leslie Lamport}      % Declares the author's name.
\date{January 21, 1994}      % Deleting this command produces today's date.

\newcommand{\ip}[2]{(#1, #2)}
                             % Defines \ip{arg1}{arg2} to mean
                             % (arg1, arg2).

%\newcommand{\ip}[2]{\langle #1 | #2\rangle}
                             % This is an alternative definition of
                             % \ip that is commented out.

\begin{document}             % End of preamble and beginning of text.

\maketitle                   % Produces the title.

This is an example input file.  Comparing it with
the output it generates can show you how to
produce a simple document of your own.

\section{Ordinary Text}      % Produces section heading.  Lower-level
                             % sections are begun with similar
                             % \subsection and \subsubsection commands.

The ends  of words and sentences are marked
  by   spaces. It  doesn't matter how many
spaces    you type; one is as good as 100.  The
end of   a line counts as a space.

One   or more   blank lines denote the  end
of  a paragraph.

Since any number of consecutive spaces are treated
like a single one, the formatting of the input
file makes no difference to
      \LaTeX,                % The \LaTeX command generates the LaTeX logo.
but it makes a difference to you.  When you use
\LaTeX, making your input file as easy to read
as possible will be a great help as you write
your document and when you change it.  This sample
file shows how you can add comments to your own input
file.

Because printing is different from typewriting,
there are a number of things that you have to do
differently when preparing an input file than if
you were just typing the document directly.
Quotation marks like
       ``this''
have to be handled specially, as do quotes within
quotes:
       ``\,`this'            % \, separates the double and single quote.
        is what I just
        wrote, not  `that'\,''.

Dashes come in three sizes: an
       intra-word
dash, a medium dash for number ranges like
       1--2,
and a punctuation
       dash---like
this.

A sentence-ending space should be larger than the
space between words within a sentence.  You
sometimes have to type special commands in
conjunction with punctuation characters to get
this right, as in the following sentence.
       Gnats, gnus, etc.\ all  % `\ ' makes an inter-word space.
       begin with G\@.         % \@ marks end-of-sentence punctuation.
You should check the spaces after periods when
reading your output to make sure you haven't
forgotten any special cases.  Generating an
ellipsis
       \ldots\               % `\ ' is needed after `\ldots' because TeX
                             % ignores spaces after command names like \ldots
                             % made from \ + letters.
                             %
                             % Note how a `%' character causes TeX to ignore
                             % the end of the input line, so these blank lines
                             % do not start a new paragraph.
                             %
with the right spacing around the periods requires
a special command.

\LaTeX\ interprets some common characters as
commands, so you must type special commands to
generate them.  These characters include the
following:
       \$ \& \% \# \{ and \}.

In printing, text is usually emphasized with an
       \emph{italic}
type style.

\begin{em}
   A long segment of text can also be emphasized
   in this way.  Text within such a segment can be
   given \emph{additional} emphasis.
\end{em}

It is sometimes necessary to prevent \LaTeX\ from
breaking a line where it might otherwise do so.
This may be at a space, as between the ``Mr.''\ and
``Jones'' in
       ``Mr.~Jones'',        % ~ produces an unbreakable interword space.
or within a word---especially when the word is a
symbol like
       \mbox{\emph{itemnum}}
that makes little sense when hyphenated across
lines.

Footnotes\footnote{This is an example of a footnote.}
pose no problem.

\LaTeX\ is good at typesetting mathematical formulas
like
       \( x-3y + z = 7 \)
or
       \( a_{1} > x^{2n} + y^{2n} > x' \)
or
       \( \ip{A}{B} = \sum_{i} a_{i} b_{i} \).
The spaces you type in a formula are
ignored.  Remember that a letter like
       $x$                   % $ ... $  and  \( ... \)  are equivalent
is a formula when it denotes a mathematical
symbol, and it should be typed as one.

\section{Displayed Text}

Text is displayed by indenting it from the left
margin.  Quotations are commonly displayed.  There
are short quotations
\begin{quote}
   This is a short quotation.  It consists of a
   single paragraph of text.  See how it is formatted.
\end{quote}
and longer ones.
\begin{quotation}
   This is a longer quotation.  It consists of two
   paragraphs of text, neither of which are
   particularly interesting.

   This is the second paragraph of the quotation.  It
   is just as dull as the first paragraph.
\end{quotation}
Another frequently-displayed structure is a list.
The following is an example of an \emph{itemized}
list.
\begin{itemize}
   \item This is the first item of an itemized list.
         Each item in the list is marked with a ``tick''.
         You don't have to worry about what kind of tick
         mark is used.

   \item This is the second item of the list.  It
         contains another list nested inside it.  The inner
         list is an \emph{enumerated} list.
         \begin{enumerate}
            \item This is the first item of an enumerated
                  list that is nested within the itemized list.

            \item This is the second item of the inner list.
                  \LaTeX\ allows you to nest lists deeper than
                  you really should.
         \end{enumerate}
         This is the rest of the second item of the outer
         list.  It is no more interesting than any other
         part of the item.
   \item This is the third item of the list.
\end{itemize}
You can even display poetry.
\begin{verse}
   There is an environment
    for verse \\             % The \\ command separates lines
   Whose features some poets % within a stanza.
   will curse.

                             % One or more blank lines separate stanzas.

   For instead of making\\
   Them do \emph{all} line breaking, \\
   It allows them to put too many words on a line when they'd rather be
   forced to be terse.
\end{verse}

Mathematical formulas may also be displayed.  A
displayed formula
is
one-line long; multiline
formulas require special formatting instructions.
   \[  \ip{\Gamma}{\psi'} = x'' + y^{2} + z_{i}^{n}\]
Don't start a paragraph with a displayed equation,
nor make one a paragraph by itself.

\end{document}               % End of document.

\section{Bench Huge Section 525}
% synthetic bench section 525
% This is a sample LaTeX input file.  (Version of 12 August 2004.)
%
% A '%' character causes TeX to ignore all remaining text on the line,
% and is used for comments like this one.

\documentclass{article}      % Specifies the document class

                             % The preamble begins here.
\title{An Example Document}  % Declares the document's title.
\author{Leslie Lamport}      % Declares the author's name.
\date{January 21, 1994}      % Deleting this command produces today's date.

\newcommand{\ip}[2]{(#1, #2)}
                             % Defines \ip{arg1}{arg2} to mean
                             % (arg1, arg2).

%\newcommand{\ip}[2]{\langle #1 | #2\rangle}
                             % This is an alternative definition of
                             % \ip that is commented out.

\begin{document}             % End of preamble and beginning of text.

\maketitle                   % Produces the title.

This is an example input file.  Comparing it with
the output it generates can show you how to
produce a simple document of your own.

\section{Ordinary Text}      % Produces section heading.  Lower-level
                             % sections are begun with similar
                             % \subsection and \subsubsection commands.

The ends  of words and sentences are marked
  by   spaces. It  doesn't matter how many
spaces    you type; one is as good as 100.  The
end of   a line counts as a space.

One   or more   blank lines denote the  end
of  a paragraph.

Since any number of consecutive spaces are treated
like a single one, the formatting of the input
file makes no difference to
      \LaTeX,                % The \LaTeX command generates the LaTeX logo.
but it makes a difference to you.  When you use
\LaTeX, making your input file as easy to read
as possible will be a great help as you write
your document and when you change it.  This sample
file shows how you can add comments to your own input
file.

Because printing is different from typewriting,
there are a number of things that you have to do
differently when preparing an input file than if
you were just typing the document directly.
Quotation marks like
       ``this''
have to be handled specially, as do quotes within
quotes:
       ``\,`this'            % \, separates the double and single quote.
        is what I just
        wrote, not  `that'\,''.

Dashes come in three sizes: an
       intra-word
dash, a medium dash for number ranges like
       1--2,
and a punctuation
       dash---like
this.

A sentence-ending space should be larger than the
space between words within a sentence.  You
sometimes have to type special commands in
conjunction with punctuation characters to get
this right, as in the following sentence.
       Gnats, gnus, etc.\ all  % `\ ' makes an inter-word space.
       begin with G\@.         % \@ marks end-of-sentence punctuation.
You should check the spaces after periods when
reading your output to make sure you haven't
forgotten any special cases.  Generating an
ellipsis
       \ldots\               % `\ ' is needed after `\ldots' because TeX
                             % ignores spaces after command names like \ldots
                             % made from \ + letters.
                             %
                             % Note how a `%' character causes TeX to ignore
                             % the end of the input line, so these blank lines
                             % do not start a new paragraph.
                             %
with the right spacing around the periods requires
a special command.

\LaTeX\ interprets some common characters as
commands, so you must type special commands to
generate them.  These characters include the
following:
       \$ \& \% \# \{ and \}.

In printing, text is usually emphasized with an
       \emph{italic}
type style.

\begin{em}
   A long segment of text can also be emphasized
   in this way.  Text within such a segment can be
   given \emph{additional} emphasis.
\end{em}

It is sometimes necessary to prevent \LaTeX\ from
breaking a line where it might otherwise do so.
This may be at a space, as between the ``Mr.''\ and
``Jones'' in
       ``Mr.~Jones'',        % ~ produces an unbreakable interword space.
or within a word---especially when the word is a
symbol like
       \mbox{\emph{itemnum}}
that makes little sense when hyphenated across
lines.

Footnotes\footnote{This is an example of a footnote.}
pose no problem.

\LaTeX\ is good at typesetting mathematical formulas
like
       \( x-3y + z = 7 \)
or
       \( a_{1} > x^{2n} + y^{2n} > x' \)
or
       \( \ip{A}{B} = \sum_{i} a_{i} b_{i} \).
The spaces you type in a formula are
ignored.  Remember that a letter like
       $x$                   % $ ... $  and  \( ... \)  are equivalent
is a formula when it denotes a mathematical
symbol, and it should be typed as one.

\section{Displayed Text}

Text is displayed by indenting it from the left
margin.  Quotations are commonly displayed.  There
are short quotations
\begin{quote}
   This is a short quotation.  It consists of a
   single paragraph of text.  See how it is formatted.
\end{quote}
and longer ones.
\begin{quotation}
   This is a longer quotation.  It consists of two
   paragraphs of text, neither of which are
   particularly interesting.

   This is the second paragraph of the quotation.  It
   is just as dull as the first paragraph.
\end{quotation}
Another frequently-displayed structure is a list.
The following is an example of an \emph{itemized}
list.
\begin{itemize}
   \item This is the first item of an itemized list.
         Each item in the list is marked with a ``tick''.
         You don't have to worry about what kind of tick
         mark is used.

   \item This is the second item of the list.  It
         contains another list nested inside it.  The inner
         list is an \emph{enumerated} list.
         \begin{enumerate}
            \item This is the first item of an enumerated
                  list that is nested within the itemized list.

            \item This is the second item of the inner list.
                  \LaTeX\ allows you to nest lists deeper than
                  you really should.
         \end{enumerate}
         This is the rest of the second item of the outer
         list.  It is no more interesting than any other
         part of the item.
   \item This is the third item of the list.
\end{itemize}
You can even display poetry.
\begin{verse}
   There is an environment
    for verse \\             % The \\ command separates lines
   Whose features some poets % within a stanza.
   will curse.

                             % One or more blank lines separate stanzas.

   For instead of making\\
   Them do \emph{all} line breaking, \\
   It allows them to put too many words on a line when they'd rather be
   forced to be terse.
\end{verse}

Mathematical formulas may also be displayed.  A
displayed formula
is
one-line long; multiline
formulas require special formatting instructions.
   \[  \ip{\Gamma}{\psi'} = x'' + y^{2} + z_{i}^{n}\]
Don't start a paragraph with a displayed equation,
nor make one a paragraph by itself.

\end{document}               % End of document.

\section{Bench Huge Section 526}
% synthetic bench section 526
% This is a sample LaTeX input file.  (Version of 12 August 2004.)
%
% A '%' character causes TeX to ignore all remaining text on the line,
% and is used for comments like this one.

\documentclass{article}      % Specifies the document class

                             % The preamble begins here.
\title{An Example Document}  % Declares the document's title.
\author{Leslie Lamport}      % Declares the author's name.
\date{January 21, 1994}      % Deleting this command produces today's date.

\newcommand{\ip}[2]{(#1, #2)}
                             % Defines \ip{arg1}{arg2} to mean
                             % (arg1, arg2).

%\newcommand{\ip}[2]{\langle #1 | #2\rangle}
                             % This is an alternative definition of
                             % \ip that is commented out.

\begin{document}             % End of preamble and beginning of text.

\maketitle                   % Produces the title.

This is an example input file.  Comparing it with
the output it generates can show you how to
produce a simple document of your own.

\section{Ordinary Text}      % Produces section heading.  Lower-level
                             % sections are begun with similar
                             % \subsection and \subsubsection commands.

The ends  of words and sentences are marked
  by   spaces. It  doesn't matter how many
spaces    you type; one is as good as 100.  The
end of   a line counts as a space.

One   or more   blank lines denote the  end
of  a paragraph.

Since any number of consecutive spaces are treated
like a single one, the formatting of the input
file makes no difference to
      \LaTeX,                % The \LaTeX command generates the LaTeX logo.
but it makes a difference to you.  When you use
\LaTeX, making your input file as easy to read
as possible will be a great help as you write
your document and when you change it.  This sample
file shows how you can add comments to your own input
file.

Because printing is different from typewriting,
there are a number of things that you have to do
differently when preparing an input file than if
you were just typing the document directly.
Quotation marks like
       ``this''
have to be handled specially, as do quotes within
quotes:
       ``\,`this'            % \, separates the double and single quote.
        is what I just
        wrote, not  `that'\,''.

Dashes come in three sizes: an
       intra-word
dash, a medium dash for number ranges like
       1--2,
and a punctuation
       dash---like
this.

A sentence-ending space should be larger than the
space between words within a sentence.  You
sometimes have to type special commands in
conjunction with punctuation characters to get
this right, as in the following sentence.
       Gnats, gnus, etc.\ all  % `\ ' makes an inter-word space.
       begin with G\@.         % \@ marks end-of-sentence punctuation.
You should check the spaces after periods when
reading your output to make sure you haven't
forgotten any special cases.  Generating an
ellipsis
       \ldots\               % `\ ' is needed after `\ldots' because TeX
                             % ignores spaces after command names like \ldots
                             % made from \ + letters.
                             %
                             % Note how a `%' character causes TeX to ignore
                             % the end of the input line, so these blank lines
                             % do not start a new paragraph.
                             %
with the right spacing around the periods requires
a special command.

\LaTeX\ interprets some common characters as
commands, so you must type special commands to
generate them.  These characters include the
following:
       \$ \& \% \# \{ and \}.

In printing, text is usually emphasized with an
       \emph{italic}
type style.

\begin{em}
   A long segment of text can also be emphasized
   in this way.  Text within such a segment can be
   given \emph{additional} emphasis.
\end{em}

It is sometimes necessary to prevent \LaTeX\ from
breaking a line where it might otherwise do so.
This may be at a space, as between the ``Mr.''\ and
``Jones'' in
       ``Mr.~Jones'',        % ~ produces an unbreakable interword space.
or within a word---especially when the word is a
symbol like
       \mbox{\emph{itemnum}}
that makes little sense when hyphenated across
lines.

Footnotes\footnote{This is an example of a footnote.}
pose no problem.

\LaTeX\ is good at typesetting mathematical formulas
like
       \( x-3y + z = 7 \)
or
       \( a_{1} > x^{2n} + y^{2n} > x' \)
or
       \( \ip{A}{B} = \sum_{i} a_{i} b_{i} \).
The spaces you type in a formula are
ignored.  Remember that a letter like
       $x$                   % $ ... $  and  \( ... \)  are equivalent
is a formula when it denotes a mathematical
symbol, and it should be typed as one.

\section{Displayed Text}

Text is displayed by indenting it from the left
margin.  Quotations are commonly displayed.  There
are short quotations
\begin{quote}
   This is a short quotation.  It consists of a
   single paragraph of text.  See how it is formatted.
\end{quote}
and longer ones.
\begin{quotation}
   This is a longer quotation.  It consists of two
   paragraphs of text, neither of which are
   particularly interesting.

   This is the second paragraph of the quotation.  It
   is just as dull as the first paragraph.
\end{quotation}
Another frequently-displayed structure is a list.
The following is an example of an \emph{itemized}
list.
\begin{itemize}
   \item This is the first item of an itemized list.
         Each item in the list is marked with a ``tick''.
         You don't have to worry about what kind of tick
         mark is used.

   \item This is the second item of the list.  It
         contains another list nested inside it.  The inner
         list is an \emph{enumerated} list.
         \begin{enumerate}
            \item This is the first item of an enumerated
                  list that is nested within the itemized list.

            \item This is the second item of the inner list.
                  \LaTeX\ allows you to nest lists deeper than
                  you really should.
         \end{enumerate}
         This is the rest of the second item of the outer
         list.  It is no more interesting than any other
         part of the item.
   \item This is the third item of the list.
\end{itemize}
You can even display poetry.
\begin{verse}
   There is an environment
    for verse \\             % The \\ command separates lines
   Whose features some poets % within a stanza.
   will curse.

                             % One or more blank lines separate stanzas.

   For instead of making\\
   Them do \emph{all} line breaking, \\
   It allows them to put too many words on a line when they'd rather be
   forced to be terse.
\end{verse}

Mathematical formulas may also be displayed.  A
displayed formula
is
one-line long; multiline
formulas require special formatting instructions.
   \[  \ip{\Gamma}{\psi'} = x'' + y^{2} + z_{i}^{n}\]
Don't start a paragraph with a displayed equation,
nor make one a paragraph by itself.

\end{document}               % End of document.

\section{Bench Huge Section 527}
% synthetic bench section 527
% This is a sample LaTeX input file.  (Version of 12 August 2004.)
%
% A '%' character causes TeX to ignore all remaining text on the line,
% and is used for comments like this one.

\documentclass{article}      % Specifies the document class

                             % The preamble begins here.
\title{An Example Document}  % Declares the document's title.
\author{Leslie Lamport}      % Declares the author's name.
\date{January 21, 1994}      % Deleting this command produces today's date.

\newcommand{\ip}[2]{(#1, #2)}
                             % Defines \ip{arg1}{arg2} to mean
                             % (arg1, arg2).

%\newcommand{\ip}[2]{\langle #1 | #2\rangle}
                             % This is an alternative definition of
                             % \ip that is commented out.

\begin{document}             % End of preamble and beginning of text.

\maketitle                   % Produces the title.

This is an example input file.  Comparing it with
the output it generates can show you how to
produce a simple document of your own.

\section{Ordinary Text}      % Produces section heading.  Lower-level
                             % sections are begun with similar
                             % \subsection and \subsubsection commands.

The ends  of words and sentences are marked
  by   spaces. It  doesn't matter how many
spaces    you type; one is as good as 100.  The
end of   a line counts as a space.

One   or more   blank lines denote the  end
of  a paragraph.

Since any number of consecutive spaces are treated
like a single one, the formatting of the input
file makes no difference to
      \LaTeX,                % The \LaTeX command generates the LaTeX logo.
but it makes a difference to you.  When you use
\LaTeX, making your input file as easy to read
as possible will be a great help as you write
your document and when you change it.  This sample
file shows how you can add comments to your own input
file.

Because printing is different from typewriting,
there are a number of things that you have to do
differently when preparing an input file than if
you were just typing the document directly.
Quotation marks like
       ``this''
have to be handled specially, as do quotes within
quotes:
       ``\,`this'            % \, separates the double and single quote.
        is what I just
        wrote, not  `that'\,''.

Dashes come in three sizes: an
       intra-word
dash, a medium dash for number ranges like
       1--2,
and a punctuation
       dash---like
this.

A sentence-ending space should be larger than the
space between words within a sentence.  You
sometimes have to type special commands in
conjunction with punctuation characters to get
this right, as in the following sentence.
       Gnats, gnus, etc.\ all  % `\ ' makes an inter-word space.
       begin with G\@.         % \@ marks end-of-sentence punctuation.
You should check the spaces after periods when
reading your output to make sure you haven't
forgotten any special cases.  Generating an
ellipsis
       \ldots\               % `\ ' is needed after `\ldots' because TeX
                             % ignores spaces after command names like \ldots
                             % made from \ + letters.
                             %
                             % Note how a `%' character causes TeX to ignore
                             % the end of the input line, so these blank lines
                             % do not start a new paragraph.
                             %
with the right spacing around the periods requires
a special command.

\LaTeX\ interprets some common characters as
commands, so you must type special commands to
generate them.  These characters include the
following:
       \$ \& \% \# \{ and \}.

In printing, text is usually emphasized with an
       \emph{italic}
type style.

\begin{em}
   A long segment of text can also be emphasized
   in this way.  Text within such a segment can be
   given \emph{additional} emphasis.
\end{em}

It is sometimes necessary to prevent \LaTeX\ from
breaking a line where it might otherwise do so.
This may be at a space, as between the ``Mr.''\ and
``Jones'' in
       ``Mr.~Jones'',        % ~ produces an unbreakable interword space.
or within a word---especially when the word is a
symbol like
       \mbox{\emph{itemnum}}
that makes little sense when hyphenated across
lines.

Footnotes\footnote{This is an example of a footnote.}
pose no problem.

\LaTeX\ is good at typesetting mathematical formulas
like
       \( x-3y + z = 7 \)
or
       \( a_{1} > x^{2n} + y^{2n} > x' \)
or
       \( \ip{A}{B} = \sum_{i} a_{i} b_{i} \).
The spaces you type in a formula are
ignored.  Remember that a letter like
       $x$                   % $ ... $  and  \( ... \)  are equivalent
is a formula when it denotes a mathematical
symbol, and it should be typed as one.

\section{Displayed Text}

Text is displayed by indenting it from the left
margin.  Quotations are commonly displayed.  There
are short quotations
\begin{quote}
   This is a short quotation.  It consists of a
   single paragraph of text.  See how it is formatted.
\end{quote}
and longer ones.
\begin{quotation}
   This is a longer quotation.  It consists of two
   paragraphs of text, neither of which are
   particularly interesting.

   This is the second paragraph of the quotation.  It
   is just as dull as the first paragraph.
\end{quotation}
Another frequently-displayed structure is a list.
The following is an example of an \emph{itemized}
list.
\begin{itemize}
   \item This is the first item of an itemized list.
         Each item in the list is marked with a ``tick''.
         You don't have to worry about what kind of tick
         mark is used.

   \item This is the second item of the list.  It
         contains another list nested inside it.  The inner
         list is an \emph{enumerated} list.
         \begin{enumerate}
            \item This is the first item of an enumerated
                  list that is nested within the itemized list.

            \item This is the second item of the inner list.
                  \LaTeX\ allows you to nest lists deeper than
                  you really should.
         \end{enumerate}
         This is the rest of the second item of the outer
         list.  It is no more interesting than any other
         part of the item.
   \item This is the third item of the list.
\end{itemize}
You can even display poetry.
\begin{verse}
   There is an environment
    for verse \\             % The \\ command separates lines
   Whose features some poets % within a stanza.
   will curse.

                             % One or more blank lines separate stanzas.

   For instead of making\\
   Them do \emph{all} line breaking, \\
   It allows them to put too many words on a line when they'd rather be
   forced to be terse.
\end{verse}

Mathematical formulas may also be displayed.  A
displayed formula
is
one-line long; multiline
formulas require special formatting instructions.
   \[  \ip{\Gamma}{\psi'} = x'' + y^{2} + z_{i}^{n}\]
Don't start a paragraph with a displayed equation,
nor make one a paragraph by itself.

\end{document}               % End of document.

\section{Bench Huge Section 528}
% synthetic bench section 528
% This is a sample LaTeX input file.  (Version of 12 August 2004.)
%
% A '%' character causes TeX to ignore all remaining text on the line,
% and is used for comments like this one.

\documentclass{article}      % Specifies the document class

                             % The preamble begins here.
\title{An Example Document}  % Declares the document's title.
\author{Leslie Lamport}      % Declares the author's name.
\date{January 21, 1994}      % Deleting this command produces today's date.

\newcommand{\ip}[2]{(#1, #2)}
                             % Defines \ip{arg1}{arg2} to mean
                             % (arg1, arg2).

%\newcommand{\ip}[2]{\langle #1 | #2\rangle}
                             % This is an alternative definition of
                             % \ip that is commented out.

\begin{document}             % End of preamble and beginning of text.

\maketitle                   % Produces the title.

This is an example input file.  Comparing it with
the output it generates can show you how to
produce a simple document of your own.

\section{Ordinary Text}      % Produces section heading.  Lower-level
                             % sections are begun with similar
                             % \subsection and \subsubsection commands.

The ends  of words and sentences are marked
  by   spaces. It  doesn't matter how many
spaces    you type; one is as good as 100.  The
end of   a line counts as a space.

One   or more   blank lines denote the  end
of  a paragraph.

Since any number of consecutive spaces are treated
like a single one, the formatting of the input
file makes no difference to
      \LaTeX,                % The \LaTeX command generates the LaTeX logo.
but it makes a difference to you.  When you use
\LaTeX, making your input file as easy to read
as possible will be a great help as you write
your document and when you change it.  This sample
file shows how you can add comments to your own input
file.

Because printing is different from typewriting,
there are a number of things that you have to do
differently when preparing an input file than if
you were just typing the document directly.
Quotation marks like
       ``this''
have to be handled specially, as do quotes within
quotes:
       ``\,`this'            % \, separates the double and single quote.
        is what I just
        wrote, not  `that'\,''.

Dashes come in three sizes: an
       intra-word
dash, a medium dash for number ranges like
       1--2,
and a punctuation
       dash---like
this.

A sentence-ending space should be larger than the
space between words within a sentence.  You
sometimes have to type special commands in
conjunction with punctuation characters to get
this right, as in the following sentence.
       Gnats, gnus, etc.\ all  % `\ ' makes an inter-word space.
       begin with G\@.         % \@ marks end-of-sentence punctuation.
You should check the spaces after periods when
reading your output to make sure you haven't
forgotten any special cases.  Generating an
ellipsis
       \ldots\               % `\ ' is needed after `\ldots' because TeX
                             % ignores spaces after command names like \ldots
                             % made from \ + letters.
                             %
                             % Note how a `%' character causes TeX to ignore
                             % the end of the input line, so these blank lines
                             % do not start a new paragraph.
                             %
with the right spacing around the periods requires
a special command.

\LaTeX\ interprets some common characters as
commands, so you must type special commands to
generate them.  These characters include the
following:
       \$ \& \% \# \{ and \}.

In printing, text is usually emphasized with an
       \emph{italic}
type style.

\begin{em}
   A long segment of text can also be emphasized
   in this way.  Text within such a segment can be
   given \emph{additional} emphasis.
\end{em}

It is sometimes necessary to prevent \LaTeX\ from
breaking a line where it might otherwise do so.
This may be at a space, as between the ``Mr.''\ and
``Jones'' in
       ``Mr.~Jones'',        % ~ produces an unbreakable interword space.
or within a word---especially when the word is a
symbol like
       \mbox{\emph{itemnum}}
that makes little sense when hyphenated across
lines.

Footnotes\footnote{This is an example of a footnote.}
pose no problem.

\LaTeX\ is good at typesetting mathematical formulas
like
       \( x-3y + z = 7 \)
or
       \( a_{1} > x^{2n} + y^{2n} > x' \)
or
       \( \ip{A}{B} = \sum_{i} a_{i} b_{i} \).
The spaces you type in a formula are
ignored.  Remember that a letter like
       $x$                   % $ ... $  and  \( ... \)  are equivalent
is a formula when it denotes a mathematical
symbol, and it should be typed as one.

\section{Displayed Text}

Text is displayed by indenting it from the left
margin.  Quotations are commonly displayed.  There
are short quotations
\begin{quote}
   This is a short quotation.  It consists of a
   single paragraph of text.  See how it is formatted.
\end{quote}
and longer ones.
\begin{quotation}
   This is a longer quotation.  It consists of two
   paragraphs of text, neither of which are
   particularly interesting.

   This is the second paragraph of the quotation.  It
   is just as dull as the first paragraph.
\end{quotation}
Another frequently-displayed structure is a list.
The following is an example of an \emph{itemized}
list.
\begin{itemize}
   \item This is the first item of an itemized list.
         Each item in the list is marked with a ``tick''.
         You don't have to worry about what kind of tick
         mark is used.

   \item This is the second item of the list.  It
         contains another list nested inside it.  The inner
         list is an \emph{enumerated} list.
         \begin{enumerate}
            \item This is the first item of an enumerated
                  list that is nested within the itemized list.

            \item This is the second item of the inner list.
                  \LaTeX\ allows you to nest lists deeper than
                  you really should.
         \end{enumerate}
         This is the rest of the second item of the outer
         list.  It is no more interesting than any other
         part of the item.
   \item This is the third item of the list.
\end{itemize}
You can even display poetry.
\begin{verse}
   There is an environment
    for verse \\             % The \\ command separates lines
   Whose features some poets % within a stanza.
   will curse.

                             % One or more blank lines separate stanzas.

   For instead of making\\
   Them do \emph{all} line breaking, \\
   It allows them to put too many words on a line when they'd rather be
   forced to be terse.
\end{verse}

Mathematical formulas may also be displayed.  A
displayed formula
is
one-line long; multiline
formulas require special formatting instructions.
   \[  \ip{\Gamma}{\psi'} = x'' + y^{2} + z_{i}^{n}\]
Don't start a paragraph with a displayed equation,
nor make one a paragraph by itself.

\end{document}               % End of document.

\section{Bench Huge Section 529}
% synthetic bench section 529
% This is a sample LaTeX input file.  (Version of 12 August 2004.)
%
% A '%' character causes TeX to ignore all remaining text on the line,
% and is used for comments like this one.

\documentclass{article}      % Specifies the document class

                             % The preamble begins here.
\title{An Example Document}  % Declares the document's title.
\author{Leslie Lamport}      % Declares the author's name.
\date{January 21, 1994}      % Deleting this command produces today's date.

\newcommand{\ip}[2]{(#1, #2)}
                             % Defines \ip{arg1}{arg2} to mean
                             % (arg1, arg2).

%\newcommand{\ip}[2]{\langle #1 | #2\rangle}
                             % This is an alternative definition of
                             % \ip that is commented out.

\begin{document}             % End of preamble and beginning of text.

\maketitle                   % Produces the title.

This is an example input file.  Comparing it with
the output it generates can show you how to
produce a simple document of your own.

\section{Ordinary Text}      % Produces section heading.  Lower-level
                             % sections are begun with similar
                             % \subsection and \subsubsection commands.

The ends  of words and sentences are marked
  by   spaces. It  doesn't matter how many
spaces    you type; one is as good as 100.  The
end of   a line counts as a space.

One   or more   blank lines denote the  end
of  a paragraph.

Since any number of consecutive spaces are treated
like a single one, the formatting of the input
file makes no difference to
      \LaTeX,                % The \LaTeX command generates the LaTeX logo.
but it makes a difference to you.  When you use
\LaTeX, making your input file as easy to read
as possible will be a great help as you write
your document and when you change it.  This sample
file shows how you can add comments to your own input
file.

Because printing is different from typewriting,
there are a number of things that you have to do
differently when preparing an input file than if
you were just typing the document directly.
Quotation marks like
       ``this''
have to be handled specially, as do quotes within
quotes:
       ``\,`this'            % \, separates the double and single quote.
        is what I just
        wrote, not  `that'\,''.

Dashes come in three sizes: an
       intra-word
dash, a medium dash for number ranges like
       1--2,
and a punctuation
       dash---like
this.

A sentence-ending space should be larger than the
space between words within a sentence.  You
sometimes have to type special commands in
conjunction with punctuation characters to get
this right, as in the following sentence.
       Gnats, gnus, etc.\ all  % `\ ' makes an inter-word space.
       begin with G\@.         % \@ marks end-of-sentence punctuation.
You should check the spaces after periods when
reading your output to make sure you haven't
forgotten any special cases.  Generating an
ellipsis
       \ldots\               % `\ ' is needed after `\ldots' because TeX
                             % ignores spaces after command names like \ldots
                             % made from \ + letters.
                             %
                             % Note how a `%' character causes TeX to ignore
                             % the end of the input line, so these blank lines
                             % do not start a new paragraph.
                             %
with the right spacing around the periods requires
a special command.

\LaTeX\ interprets some common characters as
commands, so you must type special commands to
generate them.  These characters include the
following:
       \$ \& \% \# \{ and \}.

In printing, text is usually emphasized with an
       \emph{italic}
type style.

\begin{em}
   A long segment of text can also be emphasized
   in this way.  Text within such a segment can be
   given \emph{additional} emphasis.
\end{em}

It is sometimes necessary to prevent \LaTeX\ from
breaking a line where it might otherwise do so.
This may be at a space, as between the ``Mr.''\ and
``Jones'' in
       ``Mr.~Jones'',        % ~ produces an unbreakable interword space.
or within a word---especially when the word is a
symbol like
       \mbox{\emph{itemnum}}
that makes little sense when hyphenated across
lines.

Footnotes\footnote{This is an example of a footnote.}
pose no problem.

\LaTeX\ is good at typesetting mathematical formulas
like
       \( x-3y + z = 7 \)
or
       \( a_{1} > x^{2n} + y^{2n} > x' \)
or
       \( \ip{A}{B} = \sum_{i} a_{i} b_{i} \).
The spaces you type in a formula are
ignored.  Remember that a letter like
       $x$                   % $ ... $  and  \( ... \)  are equivalent
is a formula when it denotes a mathematical
symbol, and it should be typed as one.

\section{Displayed Text}

Text is displayed by indenting it from the left
margin.  Quotations are commonly displayed.  There
are short quotations
\begin{quote}
   This is a short quotation.  It consists of a
   single paragraph of text.  See how it is formatted.
\end{quote}
and longer ones.
\begin{quotation}
   This is a longer quotation.  It consists of two
   paragraphs of text, neither of which are
   particularly interesting.

   This is the second paragraph of the quotation.  It
   is just as dull as the first paragraph.
\end{quotation}
Another frequently-displayed structure is a list.
The following is an example of an \emph{itemized}
list.
\begin{itemize}
   \item This is the first item of an itemized list.
         Each item in the list is marked with a ``tick''.
         You don't have to worry about what kind of tick
         mark is used.

   \item This is the second item of the list.  It
         contains another list nested inside it.  The inner
         list is an \emph{enumerated} list.
         \begin{enumerate}
            \item This is the first item of an enumerated
                  list that is nested within the itemized list.

            \item This is the second item of the inner list.
                  \LaTeX\ allows you to nest lists deeper than
                  you really should.
         \end{enumerate}
         This is the rest of the second item of the outer
         list.  It is no more interesting than any other
         part of the item.
   \item This is the third item of the list.
\end{itemize}
You can even display poetry.
\begin{verse}
   There is an environment
    for verse \\             % The \\ command separates lines
   Whose features some poets % within a stanza.
   will curse.

                             % One or more blank lines separate stanzas.

   For instead of making\\
   Them do \emph{all} line breaking, \\
   It allows them to put too many words on a line when they'd rather be
   forced to be terse.
\end{verse}

Mathematical formulas may also be displayed.  A
displayed formula
is
one-line long; multiline
formulas require special formatting instructions.
   \[  \ip{\Gamma}{\psi'} = x'' + y^{2} + z_{i}^{n}\]
Don't start a paragraph with a displayed equation,
nor make one a paragraph by itself.

\end{document}               % End of document.

\section{Bench Huge Section 530}
% synthetic bench section 530
% This is a sample LaTeX input file.  (Version of 12 August 2004.)
%
% A '%' character causes TeX to ignore all remaining text on the line,
% and is used for comments like this one.

\documentclass{article}      % Specifies the document class

                             % The preamble begins here.
\title{An Example Document}  % Declares the document's title.
\author{Leslie Lamport}      % Declares the author's name.
\date{January 21, 1994}      % Deleting this command produces today's date.

\newcommand{\ip}[2]{(#1, #2)}
                             % Defines \ip{arg1}{arg2} to mean
                             % (arg1, arg2).

%\newcommand{\ip}[2]{\langle #1 | #2\rangle}
                             % This is an alternative definition of
                             % \ip that is commented out.

\begin{document}             % End of preamble and beginning of text.

\maketitle                   % Produces the title.

This is an example input file.  Comparing it with
the output it generates can show you how to
produce a simple document of your own.

\section{Ordinary Text}      % Produces section heading.  Lower-level
                             % sections are begun with similar
                             % \subsection and \subsubsection commands.

The ends  of words and sentences are marked
  by   spaces. It  doesn't matter how many
spaces    you type; one is as good as 100.  The
end of   a line counts as a space.

One   or more   blank lines denote the  end
of  a paragraph.

Since any number of consecutive spaces are treated
like a single one, the formatting of the input
file makes no difference to
      \LaTeX,                % The \LaTeX command generates the LaTeX logo.
but it makes a difference to you.  When you use
\LaTeX, making your input file as easy to read
as possible will be a great help as you write
your document and when you change it.  This sample
file shows how you can add comments to your own input
file.

Because printing is different from typewriting,
there are a number of things that you have to do
differently when preparing an input file than if
you were just typing the document directly.
Quotation marks like
       ``this''
have to be handled specially, as do quotes within
quotes:
       ``\,`this'            % \, separates the double and single quote.
        is what I just
        wrote, not  `that'\,''.

Dashes come in three sizes: an
       intra-word
dash, a medium dash for number ranges like
       1--2,
and a punctuation
       dash---like
this.

A sentence-ending space should be larger than the
space between words within a sentence.  You
sometimes have to type special commands in
conjunction with punctuation characters to get
this right, as in the following sentence.
       Gnats, gnus, etc.\ all  % `\ ' makes an inter-word space.
       begin with G\@.         % \@ marks end-of-sentence punctuation.
You should check the spaces after periods when
reading your output to make sure you haven't
forgotten any special cases.  Generating an
ellipsis
       \ldots\               % `\ ' is needed after `\ldots' because TeX
                             % ignores spaces after command names like \ldots
                             % made from \ + letters.
                             %
                             % Note how a `%' character causes TeX to ignore
                             % the end of the input line, so these blank lines
                             % do not start a new paragraph.
                             %
with the right spacing around the periods requires
a special command.

\LaTeX\ interprets some common characters as
commands, so you must type special commands to
generate them.  These characters include the
following:
       \$ \& \% \# \{ and \}.

In printing, text is usually emphasized with an
       \emph{italic}
type style.

\begin{em}
   A long segment of text can also be emphasized
   in this way.  Text within such a segment can be
   given \emph{additional} emphasis.
\end{em}

It is sometimes necessary to prevent \LaTeX\ from
breaking a line where it might otherwise do so.
This may be at a space, as between the ``Mr.''\ and
``Jones'' in
       ``Mr.~Jones'',        % ~ produces an unbreakable interword space.
or within a word---especially when the word is a
symbol like
       \mbox{\emph{itemnum}}
that makes little sense when hyphenated across
lines.

Footnotes\footnote{This is an example of a footnote.}
pose no problem.

\LaTeX\ is good at typesetting mathematical formulas
like
       \( x-3y + z = 7 \)
or
       \( a_{1} > x^{2n} + y^{2n} > x' \)
or
       \( \ip{A}{B} = \sum_{i} a_{i} b_{i} \).
The spaces you type in a formula are
ignored.  Remember that a letter like
       $x$                   % $ ... $  and  \( ... \)  are equivalent
is a formula when it denotes a mathematical
symbol, and it should be typed as one.

\section{Displayed Text}

Text is displayed by indenting it from the left
margin.  Quotations are commonly displayed.  There
are short quotations
\begin{quote}
   This is a short quotation.  It consists of a
   single paragraph of text.  See how it is formatted.
\end{quote}
and longer ones.
\begin{quotation}
   This is a longer quotation.  It consists of two
   paragraphs of text, neither of which are
   particularly interesting.

   This is the second paragraph of the quotation.  It
   is just as dull as the first paragraph.
\end{quotation}
Another frequently-displayed structure is a list.
The following is an example of an \emph{itemized}
list.
\begin{itemize}
   \item This is the first item of an itemized list.
         Each item in the list is marked with a ``tick''.
         You don't have to worry about what kind of tick
         mark is used.

   \item This is the second item of the list.  It
         contains another list nested inside it.  The inner
         list is an \emph{enumerated} list.
         \begin{enumerate}
            \item This is the first item of an enumerated
                  list that is nested within the itemized list.

            \item This is the second item of the inner list.
                  \LaTeX\ allows you to nest lists deeper than
                  you really should.
         \end{enumerate}
         This is the rest of the second item of the outer
         list.  It is no more interesting than any other
         part of the item.
   \item This is the third item of the list.
\end{itemize}
You can even display poetry.
\begin{verse}
   There is an environment
    for verse \\             % The \\ command separates lines
   Whose features some poets % within a stanza.
   will curse.

                             % One or more blank lines separate stanzas.

   For instead of making\\
   Them do \emph{all} line breaking, \\
   It allows them to put too many words on a line when they'd rather be
   forced to be terse.
\end{verse}

Mathematical formulas may also be displayed.  A
displayed formula
is
one-line long; multiline
formulas require special formatting instructions.
   \[  \ip{\Gamma}{\psi'} = x'' + y^{2} + z_{i}^{n}\]
Don't start a paragraph with a displayed equation,
nor make one a paragraph by itself.

\end{document}               % End of document.

\section{Bench Huge Section 531}
% synthetic bench section 531
% This is a sample LaTeX input file.  (Version of 12 August 2004.)
%
% A '%' character causes TeX to ignore all remaining text on the line,
% and is used for comments like this one.

\documentclass{article}      % Specifies the document class

                             % The preamble begins here.
\title{An Example Document}  % Declares the document's title.
\author{Leslie Lamport}      % Declares the author's name.
\date{January 21, 1994}      % Deleting this command produces today's date.

\newcommand{\ip}[2]{(#1, #2)}
                             % Defines \ip{arg1}{arg2} to mean
                             % (arg1, arg2).

%\newcommand{\ip}[2]{\langle #1 | #2\rangle}
                             % This is an alternative definition of
                             % \ip that is commented out.

\begin{document}             % End of preamble and beginning of text.

\maketitle                   % Produces the title.

This is an example input file.  Comparing it with
the output it generates can show you how to
produce a simple document of your own.

\section{Ordinary Text}      % Produces section heading.  Lower-level
                             % sections are begun with similar
                             % \subsection and \subsubsection commands.

The ends  of words and sentences are marked
  by   spaces. It  doesn't matter how many
spaces    you type; one is as good as 100.  The
end of   a line counts as a space.

One   or more   blank lines denote the  end
of  a paragraph.

Since any number of consecutive spaces are treated
like a single one, the formatting of the input
file makes no difference to
      \LaTeX,                % The \LaTeX command generates the LaTeX logo.
but it makes a difference to you.  When you use
\LaTeX, making your input file as easy to read
as possible will be a great help as you write
your document and when you change it.  This sample
file shows how you can add comments to your own input
file.

Because printing is different from typewriting,
there are a number of things that you have to do
differently when preparing an input file than if
you were just typing the document directly.
Quotation marks like
       ``this''
have to be handled specially, as do quotes within
quotes:
       ``\,`this'            % \, separates the double and single quote.
        is what I just
        wrote, not  `that'\,''.

Dashes come in three sizes: an
       intra-word
dash, a medium dash for number ranges like
       1--2,
and a punctuation
       dash---like
this.

A sentence-ending space should be larger than the
space between words within a sentence.  You
sometimes have to type special commands in
conjunction with punctuation characters to get
this right, as in the following sentence.
       Gnats, gnus, etc.\ all  % `\ ' makes an inter-word space.
       begin with G\@.         % \@ marks end-of-sentence punctuation.
You should check the spaces after periods when
reading your output to make sure you haven't
forgotten any special cases.  Generating an
ellipsis
       \ldots\               % `\ ' is needed after `\ldots' because TeX
                             % ignores spaces after command names like \ldots
                             % made from \ + letters.
                             %
                             % Note how a `%' character causes TeX to ignore
                             % the end of the input line, so these blank lines
                             % do not start a new paragraph.
                             %
with the right spacing around the periods requires
a special command.

\LaTeX\ interprets some common characters as
commands, so you must type special commands to
generate them.  These characters include the
following:
       \$ \& \% \# \{ and \}.

In printing, text is usually emphasized with an
       \emph{italic}
type style.

\begin{em}
   A long segment of text can also be emphasized
   in this way.  Text within such a segment can be
   given \emph{additional} emphasis.
\end{em}

It is sometimes necessary to prevent \LaTeX\ from
breaking a line where it might otherwise do so.
This may be at a space, as between the ``Mr.''\ and
``Jones'' in
       ``Mr.~Jones'',        % ~ produces an unbreakable interword space.
or within a word---especially when the word is a
symbol like
       \mbox{\emph{itemnum}}
that makes little sense when hyphenated across
lines.

Footnotes\footnote{This is an example of a footnote.}
pose no problem.

\LaTeX\ is good at typesetting mathematical formulas
like
       \( x-3y + z = 7 \)
or
       \( a_{1} > x^{2n} + y^{2n} > x' \)
or
       \( \ip{A}{B} = \sum_{i} a_{i} b_{i} \).
The spaces you type in a formula are
ignored.  Remember that a letter like
       $x$                   % $ ... $  and  \( ... \)  are equivalent
is a formula when it denotes a mathematical
symbol, and it should be typed as one.

\section{Displayed Text}

Text is displayed by indenting it from the left
margin.  Quotations are commonly displayed.  There
are short quotations
\begin{quote}
   This is a short quotation.  It consists of a
   single paragraph of text.  See how it is formatted.
\end{quote}
and longer ones.
\begin{quotation}
   This is a longer quotation.  It consists of two
   paragraphs of text, neither of which are
   particularly interesting.

   This is the second paragraph of the quotation.  It
   is just as dull as the first paragraph.
\end{quotation}
Another frequently-displayed structure is a list.
The following is an example of an \emph{itemized}
list.
\begin{itemize}
   \item This is the first item of an itemized list.
         Each item in the list is marked with a ``tick''.
         You don't have to worry about what kind of tick
         mark is used.

   \item This is the second item of the list.  It
         contains another list nested inside it.  The inner
         list is an \emph{enumerated} list.
         \begin{enumerate}
            \item This is the first item of an enumerated
                  list that is nested within the itemized list.

            \item This is the second item of the inner list.
                  \LaTeX\ allows you to nest lists deeper than
                  you really should.
         \end{enumerate}
         This is the rest of the second item of the outer
         list.  It is no more interesting than any other
         part of the item.
   \item This is the third item of the list.
\end{itemize}
You can even display poetry.
\begin{verse}
   There is an environment
    for verse \\             % The \\ command separates lines
   Whose features some poets % within a stanza.
   will curse.

                             % One or more blank lines separate stanzas.

   For instead of making\\
   Them do \emph{all} line breaking, \\
   It allows them to put too many words on a line when they'd rather be
   forced to be terse.
\end{verse}

Mathematical formulas may also be displayed.  A
displayed formula
is
one-line long; multiline
formulas require special formatting instructions.
   \[  \ip{\Gamma}{\psi'} = x'' + y^{2} + z_{i}^{n}\]
Don't start a paragraph with a displayed equation,
nor make one a paragraph by itself.

\end{document}               % End of document.

\section{Bench Huge Section 532}
% synthetic bench section 532
% This is a sample LaTeX input file.  (Version of 12 August 2004.)
%
% A '%' character causes TeX to ignore all remaining text on the line,
% and is used for comments like this one.

\documentclass{article}      % Specifies the document class

                             % The preamble begins here.
\title{An Example Document}  % Declares the document's title.
\author{Leslie Lamport}      % Declares the author's name.
\date{January 21, 1994}      % Deleting this command produces today's date.

\newcommand{\ip}[2]{(#1, #2)}
                             % Defines \ip{arg1}{arg2} to mean
                             % (arg1, arg2).

%\newcommand{\ip}[2]{\langle #1 | #2\rangle}
                             % This is an alternative definition of
                             % \ip that is commented out.

\begin{document}             % End of preamble and beginning of text.

\maketitle                   % Produces the title.

This is an example input file.  Comparing it with
the output it generates can show you how to
produce a simple document of your own.

\section{Ordinary Text}      % Produces section heading.  Lower-level
                             % sections are begun with similar
                             % \subsection and \subsubsection commands.

The ends  of words and sentences are marked
  by   spaces. It  doesn't matter how many
spaces    you type; one is as good as 100.  The
end of   a line counts as a space.

One   or more   blank lines denote the  end
of  a paragraph.

Since any number of consecutive spaces are treated
like a single one, the formatting of the input
file makes no difference to
      \LaTeX,                % The \LaTeX command generates the LaTeX logo.
but it makes a difference to you.  When you use
\LaTeX, making your input file as easy to read
as possible will be a great help as you write
your document and when you change it.  This sample
file shows how you can add comments to your own input
file.

Because printing is different from typewriting,
there are a number of things that you have to do
differently when preparing an input file than if
you were just typing the document directly.
Quotation marks like
       ``this''
have to be handled specially, as do quotes within
quotes:
       ``\,`this'            % \, separates the double and single quote.
        is what I just
        wrote, not  `that'\,''.

Dashes come in three sizes: an
       intra-word
dash, a medium dash for number ranges like
       1--2,
and a punctuation
       dash---like
this.

A sentence-ending space should be larger than the
space between words within a sentence.  You
sometimes have to type special commands in
conjunction with punctuation characters to get
this right, as in the following sentence.
       Gnats, gnus, etc.\ all  % `\ ' makes an inter-word space.
       begin with G\@.         % \@ marks end-of-sentence punctuation.
You should check the spaces after periods when
reading your output to make sure you haven't
forgotten any special cases.  Generating an
ellipsis
       \ldots\               % `\ ' is needed after `\ldots' because TeX
                             % ignores spaces after command names like \ldots
                             % made from \ + letters.
                             %
                             % Note how a `%' character causes TeX to ignore
                             % the end of the input line, so these blank lines
                             % do not start a new paragraph.
                             %
with the right spacing around the periods requires
a special command.

\LaTeX\ interprets some common characters as
commands, so you must type special commands to
generate them.  These characters include the
following:
       \$ \& \% \# \{ and \}.

In printing, text is usually emphasized with an
       \emph{italic}
type style.

\begin{em}
   A long segment of text can also be emphasized
   in this way.  Text within such a segment can be
   given \emph{additional} emphasis.
\end{em}

It is sometimes necessary to prevent \LaTeX\ from
breaking a line where it might otherwise do so.
This may be at a space, as between the ``Mr.''\ and
``Jones'' in
       ``Mr.~Jones'',        % ~ produces an unbreakable interword space.
or within a word---especially when the word is a
symbol like
       \mbox{\emph{itemnum}}
that makes little sense when hyphenated across
lines.

Footnotes\footnote{This is an example of a footnote.}
pose no problem.

\LaTeX\ is good at typesetting mathematical formulas
like
       \( x-3y + z = 7 \)
or
       \( a_{1} > x^{2n} + y^{2n} > x' \)
or
       \( \ip{A}{B} = \sum_{i} a_{i} b_{i} \).
The spaces you type in a formula are
ignored.  Remember that a letter like
       $x$                   % $ ... $  and  \( ... \)  are equivalent
is a formula when it denotes a mathematical
symbol, and it should be typed as one.

\section{Displayed Text}

Text is displayed by indenting it from the left
margin.  Quotations are commonly displayed.  There
are short quotations
\begin{quote}
   This is a short quotation.  It consists of a
   single paragraph of text.  See how it is formatted.
\end{quote}
and longer ones.
\begin{quotation}
   This is a longer quotation.  It consists of two
   paragraphs of text, neither of which are
   particularly interesting.

   This is the second paragraph of the quotation.  It
   is just as dull as the first paragraph.
\end{quotation}
Another frequently-displayed structure is a list.
The following is an example of an \emph{itemized}
list.
\begin{itemize}
   \item This is the first item of an itemized list.
         Each item in the list is marked with a ``tick''.
         You don't have to worry about what kind of tick
         mark is used.

   \item This is the second item of the list.  It
         contains another list nested inside it.  The inner
         list is an \emph{enumerated} list.
         \begin{enumerate}
            \item This is the first item of an enumerated
                  list that is nested within the itemized list.

            \item This is the second item of the inner list.
                  \LaTeX\ allows you to nest lists deeper than
                  you really should.
         \end{enumerate}
         This is the rest of the second item of the outer
         list.  It is no more interesting than any other
         part of the item.
   \item This is the third item of the list.
\end{itemize}
You can even display poetry.
\begin{verse}
   There is an environment
    for verse \\             % The \\ command separates lines
   Whose features some poets % within a stanza.
   will curse.

                             % One or more blank lines separate stanzas.

   For instead of making\\
   Them do \emph{all} line breaking, \\
   It allows them to put too many words on a line when they'd rather be
   forced to be terse.
\end{verse}

Mathematical formulas may also be displayed.  A
displayed formula
is
one-line long; multiline
formulas require special formatting instructions.
   \[  \ip{\Gamma}{\psi'} = x'' + y^{2} + z_{i}^{n}\]
Don't start a paragraph with a displayed equation,
nor make one a paragraph by itself.

\end{document}               % End of document.

\section{Bench Huge Section 533}
% synthetic bench section 533
% This is a sample LaTeX input file.  (Version of 12 August 2004.)
%
% A '%' character causes TeX to ignore all remaining text on the line,
% and is used for comments like this one.

\documentclass{article}      % Specifies the document class

                             % The preamble begins here.
\title{An Example Document}  % Declares the document's title.
\author{Leslie Lamport}      % Declares the author's name.
\date{January 21, 1994}      % Deleting this command produces today's date.

\newcommand{\ip}[2]{(#1, #2)}
                             % Defines \ip{arg1}{arg2} to mean
                             % (arg1, arg2).

%\newcommand{\ip}[2]{\langle #1 | #2\rangle}
                             % This is an alternative definition of
                             % \ip that is commented out.

\begin{document}             % End of preamble and beginning of text.

\maketitle                   % Produces the title.

This is an example input file.  Comparing it with
the output it generates can show you how to
produce a simple document of your own.

\section{Ordinary Text}      % Produces section heading.  Lower-level
                             % sections are begun with similar
                             % \subsection and \subsubsection commands.

The ends  of words and sentences are marked
  by   spaces. It  doesn't matter how many
spaces    you type; one is as good as 100.  The
end of   a line counts as a space.

One   or more   blank lines denote the  end
of  a paragraph.

Since any number of consecutive spaces are treated
like a single one, the formatting of the input
file makes no difference to
      \LaTeX,                % The \LaTeX command generates the LaTeX logo.
but it makes a difference to you.  When you use
\LaTeX, making your input file as easy to read
as possible will be a great help as you write
your document and when you change it.  This sample
file shows how you can add comments to your own input
file.

Because printing is different from typewriting,
there are a number of things that you have to do
differently when preparing an input file than if
you were just typing the document directly.
Quotation marks like
       ``this''
have to be handled specially, as do quotes within
quotes:
       ``\,`this'            % \, separates the double and single quote.
        is what I just
        wrote, not  `that'\,''.

Dashes come in three sizes: an
       intra-word
dash, a medium dash for number ranges like
       1--2,
and a punctuation
       dash---like
this.

A sentence-ending space should be larger than the
space between words within a sentence.  You
sometimes have to type special commands in
conjunction with punctuation characters to get
this right, as in the following sentence.
       Gnats, gnus, etc.\ all  % `\ ' makes an inter-word space.
       begin with G\@.         % \@ marks end-of-sentence punctuation.
You should check the spaces after periods when
reading your output to make sure you haven't
forgotten any special cases.  Generating an
ellipsis
       \ldots\               % `\ ' is needed after `\ldots' because TeX
                             % ignores spaces after command names like \ldots
                             % made from \ + letters.
                             %
                             % Note how a `%' character causes TeX to ignore
                             % the end of the input line, so these blank lines
                             % do not start a new paragraph.
                             %
with the right spacing around the periods requires
a special command.

\LaTeX\ interprets some common characters as
commands, so you must type special commands to
generate them.  These characters include the
following:
       \$ \& \% \# \{ and \}.

In printing, text is usually emphasized with an
       \emph{italic}
type style.

\begin{em}
   A long segment of text can also be emphasized
   in this way.  Text within such a segment can be
   given \emph{additional} emphasis.
\end{em}

It is sometimes necessary to prevent \LaTeX\ from
breaking a line where it might otherwise do so.
This may be at a space, as between the ``Mr.''\ and
``Jones'' in
       ``Mr.~Jones'',        % ~ produces an unbreakable interword space.
or within a word---especially when the word is a
symbol like
       \mbox{\emph{itemnum}}
that makes little sense when hyphenated across
lines.

Footnotes\footnote{This is an example of a footnote.}
pose no problem.

\LaTeX\ is good at typesetting mathematical formulas
like
       \( x-3y + z = 7 \)
or
       \( a_{1} > x^{2n} + y^{2n} > x' \)
or
       \( \ip{A}{B} = \sum_{i} a_{i} b_{i} \).
The spaces you type in a formula are
ignored.  Remember that a letter like
       $x$                   % $ ... $  and  \( ... \)  are equivalent
is a formula when it denotes a mathematical
symbol, and it should be typed as one.

\section{Displayed Text}

Text is displayed by indenting it from the left
margin.  Quotations are commonly displayed.  There
are short quotations
\begin{quote}
   This is a short quotation.  It consists of a
   single paragraph of text.  See how it is formatted.
\end{quote}
and longer ones.
\begin{quotation}
   This is a longer quotation.  It consists of two
   paragraphs of text, neither of which are
   particularly interesting.

   This is the second paragraph of the quotation.  It
   is just as dull as the first paragraph.
\end{quotation}
Another frequently-displayed structure is a list.
The following is an example of an \emph{itemized}
list.
\begin{itemize}
   \item This is the first item of an itemized list.
         Each item in the list is marked with a ``tick''.
         You don't have to worry about what kind of tick
         mark is used.

   \item This is the second item of the list.  It
         contains another list nested inside it.  The inner
         list is an \emph{enumerated} list.
         \begin{enumerate}
            \item This is the first item of an enumerated
                  list that is nested within the itemized list.

            \item This is the second item of the inner list.
                  \LaTeX\ allows you to nest lists deeper than
                  you really should.
         \end{enumerate}
         This is the rest of the second item of the outer
         list.  It is no more interesting than any other
         part of the item.
   \item This is the third item of the list.
\end{itemize}
You can even display poetry.
\begin{verse}
   There is an environment
    for verse \\             % The \\ command separates lines
   Whose features some poets % within a stanza.
   will curse.

                             % One or more blank lines separate stanzas.

   For instead of making\\
   Them do \emph{all} line breaking, \\
   It allows them to put too many words on a line when they'd rather be
   forced to be terse.
\end{verse}

Mathematical formulas may also be displayed.  A
displayed formula
is
one-line long; multiline
formulas require special formatting instructions.
   \[  \ip{\Gamma}{\psi'} = x'' + y^{2} + z_{i}^{n}\]
Don't start a paragraph with a displayed equation,
nor make one a paragraph by itself.

\end{document}               % End of document.

\section{Bench Huge Section 534}
% synthetic bench section 534
% This is a sample LaTeX input file.  (Version of 12 August 2004.)
%
% A '%' character causes TeX to ignore all remaining text on the line,
% and is used for comments like this one.

\documentclass{article}      % Specifies the document class

                             % The preamble begins here.
\title{An Example Document}  % Declares the document's title.
\author{Leslie Lamport}      % Declares the author's name.
\date{January 21, 1994}      % Deleting this command produces today's date.

\newcommand{\ip}[2]{(#1, #2)}
                             % Defines \ip{arg1}{arg2} to mean
                             % (arg1, arg2).

%\newcommand{\ip}[2]{\langle #1 | #2\rangle}
                             % This is an alternative definition of
                             % \ip that is commented out.

\begin{document}             % End of preamble and beginning of text.

\maketitle                   % Produces the title.

This is an example input file.  Comparing it with
the output it generates can show you how to
produce a simple document of your own.

\section{Ordinary Text}      % Produces section heading.  Lower-level
                             % sections are begun with similar
                             % \subsection and \subsubsection commands.

The ends  of words and sentences are marked
  by   spaces. It  doesn't matter how many
spaces    you type; one is as good as 100.  The
end of   a line counts as a space.

One   or more   blank lines denote the  end
of  a paragraph.

Since any number of consecutive spaces are treated
like a single one, the formatting of the input
file makes no difference to
      \LaTeX,                % The \LaTeX command generates the LaTeX logo.
but it makes a difference to you.  When you use
\LaTeX, making your input file as easy to read
as possible will be a great help as you write
your document and when you change it.  This sample
file shows how you can add comments to your own input
file.

Because printing is different from typewriting,
there are a number of things that you have to do
differently when preparing an input file than if
you were just typing the document directly.
Quotation marks like
       ``this''
have to be handled specially, as do quotes within
quotes:
       ``\,`this'            % \, separates the double and single quote.
        is what I just
        wrote, not  `that'\,''.

Dashes come in three sizes: an
       intra-word
dash, a medium dash for number ranges like
       1--2,
and a punctuation
       dash---like
this.

A sentence-ending space should be larger than the
space between words within a sentence.  You
sometimes have to type special commands in
conjunction with punctuation characters to get
this right, as in the following sentence.
       Gnats, gnus, etc.\ all  % `\ ' makes an inter-word space.
       begin with G\@.         % \@ marks end-of-sentence punctuation.
You should check the spaces after periods when
reading your output to make sure you haven't
forgotten any special cases.  Generating an
ellipsis
       \ldots\               % `\ ' is needed after `\ldots' because TeX
                             % ignores spaces after command names like \ldots
                             % made from \ + letters.
                             %
                             % Note how a `%' character causes TeX to ignore
                             % the end of the input line, so these blank lines
                             % do not start a new paragraph.
                             %
with the right spacing around the periods requires
a special command.

\LaTeX\ interprets some common characters as
commands, so you must type special commands to
generate them.  These characters include the
following:
       \$ \& \% \# \{ and \}.

In printing, text is usually emphasized with an
       \emph{italic}
type style.

\begin{em}
   A long segment of text can also be emphasized
   in this way.  Text within such a segment can be
   given \emph{additional} emphasis.
\end{em}

It is sometimes necessary to prevent \LaTeX\ from
breaking a line where it might otherwise do so.
This may be at a space, as between the ``Mr.''\ and
``Jones'' in
       ``Mr.~Jones'',        % ~ produces an unbreakable interword space.
or within a word---especially when the word is a
symbol like
       \mbox{\emph{itemnum}}
that makes little sense when hyphenated across
lines.

Footnotes\footnote{This is an example of a footnote.}
pose no problem.

\LaTeX\ is good at typesetting mathematical formulas
like
       \( x-3y + z = 7 \)
or
       \( a_{1} > x^{2n} + y^{2n} > x' \)
or
       \( \ip{A}{B} = \sum_{i} a_{i} b_{i} \).
The spaces you type in a formula are
ignored.  Remember that a letter like
       $x$                   % $ ... $  and  \( ... \)  are equivalent
is a formula when it denotes a mathematical
symbol, and it should be typed as one.

\section{Displayed Text}

Text is displayed by indenting it from the left
margin.  Quotations are commonly displayed.  There
are short quotations
\begin{quote}
   This is a short quotation.  It consists of a
   single paragraph of text.  See how it is formatted.
\end{quote}
and longer ones.
\begin{quotation}
   This is a longer quotation.  It consists of two
   paragraphs of text, neither of which are
   particularly interesting.

   This is the second paragraph of the quotation.  It
   is just as dull as the first paragraph.
\end{quotation}
Another frequently-displayed structure is a list.
The following is an example of an \emph{itemized}
list.
\begin{itemize}
   \item This is the first item of an itemized list.
         Each item in the list is marked with a ``tick''.
         You don't have to worry about what kind of tick
         mark is used.

   \item This is the second item of the list.  It
         contains another list nested inside it.  The inner
         list is an \emph{enumerated} list.
         \begin{enumerate}
            \item This is the first item of an enumerated
                  list that is nested within the itemized list.

            \item This is the second item of the inner list.
                  \LaTeX\ allows you to nest lists deeper than
                  you really should.
         \end{enumerate}
         This is the rest of the second item of the outer
         list.  It is no more interesting than any other
         part of the item.
   \item This is the third item of the list.
\end{itemize}
You can even display poetry.
\begin{verse}
   There is an environment
    for verse \\             % The \\ command separates lines
   Whose features some poets % within a stanza.
   will curse.

                             % One or more blank lines separate stanzas.

   For instead of making\\
   Them do \emph{all} line breaking, \\
   It allows them to put too many words on a line when they'd rather be
   forced to be terse.
\end{verse}

Mathematical formulas may also be displayed.  A
displayed formula
is
one-line long; multiline
formulas require special formatting instructions.
   \[  \ip{\Gamma}{\psi'} = x'' + y^{2} + z_{i}^{n}\]
Don't start a paragraph with a displayed equation,
nor make one a paragraph by itself.

\end{document}               % End of document.

\section{Bench Huge Section 535}
% synthetic bench section 535
% This is a sample LaTeX input file.  (Version of 12 August 2004.)
%
% A '%' character causes TeX to ignore all remaining text on the line,
% and is used for comments like this one.

\documentclass{article}      % Specifies the document class

                             % The preamble begins here.
\title{An Example Document}  % Declares the document's title.
\author{Leslie Lamport}      % Declares the author's name.
\date{January 21, 1994}      % Deleting this command produces today's date.

\newcommand{\ip}[2]{(#1, #2)}
                             % Defines \ip{arg1}{arg2} to mean
                             % (arg1, arg2).

%\newcommand{\ip}[2]{\langle #1 | #2\rangle}
                             % This is an alternative definition of
                             % \ip that is commented out.

\begin{document}             % End of preamble and beginning of text.

\maketitle                   % Produces the title.

This is an example input file.  Comparing it with
the output it generates can show you how to
produce a simple document of your own.

\section{Ordinary Text}      % Produces section heading.  Lower-level
                             % sections are begun with similar
                             % \subsection and \subsubsection commands.

The ends  of words and sentences are marked
  by   spaces. It  doesn't matter how many
spaces    you type; one is as good as 100.  The
end of   a line counts as a space.

One   or more   blank lines denote the  end
of  a paragraph.

Since any number of consecutive spaces are treated
like a single one, the formatting of the input
file makes no difference to
      \LaTeX,                % The \LaTeX command generates the LaTeX logo.
but it makes a difference to you.  When you use
\LaTeX, making your input file as easy to read
as possible will be a great help as you write
your document and when you change it.  This sample
file shows how you can add comments to your own input
file.

Because printing is different from typewriting,
there are a number of things that you have to do
differently when preparing an input file than if
you were just typing the document directly.
Quotation marks like
       ``this''
have to be handled specially, as do quotes within
quotes:
       ``\,`this'            % \, separates the double and single quote.
        is what I just
        wrote, not  `that'\,''.

Dashes come in three sizes: an
       intra-word
dash, a medium dash for number ranges like
       1--2,
and a punctuation
       dash---like
this.

A sentence-ending space should be larger than the
space between words within a sentence.  You
sometimes have to type special commands in
conjunction with punctuation characters to get
this right, as in the following sentence.
       Gnats, gnus, etc.\ all  % `\ ' makes an inter-word space.
       begin with G\@.         % \@ marks end-of-sentence punctuation.
You should check the spaces after periods when
reading your output to make sure you haven't
forgotten any special cases.  Generating an
ellipsis
       \ldots\               % `\ ' is needed after `\ldots' because TeX
                             % ignores spaces after command names like \ldots
                             % made from \ + letters.
                             %
                             % Note how a `%' character causes TeX to ignore
                             % the end of the input line, so these blank lines
                             % do not start a new paragraph.
                             %
with the right spacing around the periods requires
a special command.

\LaTeX\ interprets some common characters as
commands, so you must type special commands to
generate them.  These characters include the
following:
       \$ \& \% \# \{ and \}.

In printing, text is usually emphasized with an
       \emph{italic}
type style.

\begin{em}
   A long segment of text can also be emphasized
   in this way.  Text within such a segment can be
   given \emph{additional} emphasis.
\end{em}

It is sometimes necessary to prevent \LaTeX\ from
breaking a line where it might otherwise do so.
This may be at a space, as between the ``Mr.''\ and
``Jones'' in
       ``Mr.~Jones'',        % ~ produces an unbreakable interword space.
or within a word---especially when the word is a
symbol like
       \mbox{\emph{itemnum}}
that makes little sense when hyphenated across
lines.

Footnotes\footnote{This is an example of a footnote.}
pose no problem.

\LaTeX\ is good at typesetting mathematical formulas
like
       \( x-3y + z = 7 \)
or
       \( a_{1} > x^{2n} + y^{2n} > x' \)
or
       \( \ip{A}{B} = \sum_{i} a_{i} b_{i} \).
The spaces you type in a formula are
ignored.  Remember that a letter like
       $x$                   % $ ... $  and  \( ... \)  are equivalent
is a formula when it denotes a mathematical
symbol, and it should be typed as one.

\section{Displayed Text}

Text is displayed by indenting it from the left
margin.  Quotations are commonly displayed.  There
are short quotations
\begin{quote}
   This is a short quotation.  It consists of a
   single paragraph of text.  See how it is formatted.
\end{quote}
and longer ones.
\begin{quotation}
   This is a longer quotation.  It consists of two
   paragraphs of text, neither of which are
   particularly interesting.

   This is the second paragraph of the quotation.  It
   is just as dull as the first paragraph.
\end{quotation}
Another frequently-displayed structure is a list.
The following is an example of an \emph{itemized}
list.
\begin{itemize}
   \item This is the first item of an itemized list.
         Each item in the list is marked with a ``tick''.
         You don't have to worry about what kind of tick
         mark is used.

   \item This is the second item of the list.  It
         contains another list nested inside it.  The inner
         list is an \emph{enumerated} list.
         \begin{enumerate}
            \item This is the first item of an enumerated
                  list that is nested within the itemized list.

            \item This is the second item of the inner list.
                  \LaTeX\ allows you to nest lists deeper than
                  you really should.
         \end{enumerate}
         This is the rest of the second item of the outer
         list.  It is no more interesting than any other
         part of the item.
   \item This is the third item of the list.
\end{itemize}
You can even display poetry.
\begin{verse}
   There is an environment
    for verse \\             % The \\ command separates lines
   Whose features some poets % within a stanza.
   will curse.

                             % One or more blank lines separate stanzas.

   For instead of making\\
   Them do \emph{all} line breaking, \\
   It allows them to put too many words on a line when they'd rather be
   forced to be terse.
\end{verse}

Mathematical formulas may also be displayed.  A
displayed formula
is
one-line long; multiline
formulas require special formatting instructions.
   \[  \ip{\Gamma}{\psi'} = x'' + y^{2} + z_{i}^{n}\]
Don't start a paragraph with a displayed equation,
nor make one a paragraph by itself.

\end{document}               % End of document.

\section{Bench Huge Section 536}
% synthetic bench section 536
% This is a sample LaTeX input file.  (Version of 12 August 2004.)
%
% A '%' character causes TeX to ignore all remaining text on the line,
% and is used for comments like this one.

\documentclass{article}      % Specifies the document class

                             % The preamble begins here.
\title{An Example Document}  % Declares the document's title.
\author{Leslie Lamport}      % Declares the author's name.
\date{January 21, 1994}      % Deleting this command produces today's date.

\newcommand{\ip}[2]{(#1, #2)}
                             % Defines \ip{arg1}{arg2} to mean
                             % (arg1, arg2).

%\newcommand{\ip}[2]{\langle #1 | #2\rangle}
                             % This is an alternative definition of
                             % \ip that is commented out.

\begin{document}             % End of preamble and beginning of text.

\maketitle                   % Produces the title.

This is an example input file.  Comparing it with
the output it generates can show you how to
produce a simple document of your own.

\section{Ordinary Text}      % Produces section heading.  Lower-level
                             % sections are begun with similar
                             % \subsection and \subsubsection commands.

The ends  of words and sentences are marked
  by   spaces. It  doesn't matter how many
spaces    you type; one is as good as 100.  The
end of   a line counts as a space.

One   or more   blank lines denote the  end
of  a paragraph.

Since any number of consecutive spaces are treated
like a single one, the formatting of the input
file makes no difference to
      \LaTeX,                % The \LaTeX command generates the LaTeX logo.
but it makes a difference to you.  When you use
\LaTeX, making your input file as easy to read
as possible will be a great help as you write
your document and when you change it.  This sample
file shows how you can add comments to your own input
file.

Because printing is different from typewriting,
there are a number of things that you have to do
differently when preparing an input file than if
you were just typing the document directly.
Quotation marks like
       ``this''
have to be handled specially, as do quotes within
quotes:
       ``\,`this'            % \, separates the double and single quote.
        is what I just
        wrote, not  `that'\,''.

Dashes come in three sizes: an
       intra-word
dash, a medium dash for number ranges like
       1--2,
and a punctuation
       dash---like
this.

A sentence-ending space should be larger than the
space between words within a sentence.  You
sometimes have to type special commands in
conjunction with punctuation characters to get
this right, as in the following sentence.
       Gnats, gnus, etc.\ all  % `\ ' makes an inter-word space.
       begin with G\@.         % \@ marks end-of-sentence punctuation.
You should check the spaces after periods when
reading your output to make sure you haven't
forgotten any special cases.  Generating an
ellipsis
       \ldots\               % `\ ' is needed after `\ldots' because TeX
                             % ignores spaces after command names like \ldots
                             % made from \ + letters.
                             %
                             % Note how a `%' character causes TeX to ignore
                             % the end of the input line, so these blank lines
                             % do not start a new paragraph.
                             %
with the right spacing around the periods requires
a special command.

\LaTeX\ interprets some common characters as
commands, so you must type special commands to
generate them.  These characters include the
following:
       \$ \& \% \# \{ and \}.

In printing, text is usually emphasized with an
       \emph{italic}
type style.

\begin{em}
   A long segment of text can also be emphasized
   in this way.  Text within such a segment can be
   given \emph{additional} emphasis.
\end{em}

It is sometimes necessary to prevent \LaTeX\ from
breaking a line where it might otherwise do so.
This may be at a space, as between the ``Mr.''\ and
``Jones'' in
       ``Mr.~Jones'',        % ~ produces an unbreakable interword space.
or within a word---especially when the word is a
symbol like
       \mbox{\emph{itemnum}}
that makes little sense when hyphenated across
lines.

Footnotes\footnote{This is an example of a footnote.}
pose no problem.

\LaTeX\ is good at typesetting mathematical formulas
like
       \( x-3y + z = 7 \)
or
       \( a_{1} > x^{2n} + y^{2n} > x' \)
or
       \( \ip{A}{B} = \sum_{i} a_{i} b_{i} \).
The spaces you type in a formula are
ignored.  Remember that a letter like
       $x$                   % $ ... $  and  \( ... \)  are equivalent
is a formula when it denotes a mathematical
symbol, and it should be typed as one.

\section{Displayed Text}

Text is displayed by indenting it from the left
margin.  Quotations are commonly displayed.  There
are short quotations
\begin{quote}
   This is a short quotation.  It consists of a
   single paragraph of text.  See how it is formatted.
\end{quote}
and longer ones.
\begin{quotation}
   This is a longer quotation.  It consists of two
   paragraphs of text, neither of which are
   particularly interesting.

   This is the second paragraph of the quotation.  It
   is just as dull as the first paragraph.
\end{quotation}
Another frequently-displayed structure is a list.
The following is an example of an \emph{itemized}
list.
\begin{itemize}
   \item This is the first item of an itemized list.
         Each item in the list is marked with a ``tick''.
         You don't have to worry about what kind of tick
         mark is used.

   \item This is the second item of the list.  It
         contains another list nested inside it.  The inner
         list is an \emph{enumerated} list.
         \begin{enumerate}
            \item This is the first item of an enumerated
                  list that is nested within the itemized list.

            \item This is the second item of the inner list.
                  \LaTeX\ allows you to nest lists deeper than
                  you really should.
         \end{enumerate}
         This is the rest of the second item of the outer
         list.  It is no more interesting than any other
         part of the item.
   \item This is the third item of the list.
\end{itemize}
You can even display poetry.
\begin{verse}
   There is an environment
    for verse \\             % The \\ command separates lines
   Whose features some poets % within a stanza.
   will curse.

                             % One or more blank lines separate stanzas.

   For instead of making\\
   Them do \emph{all} line breaking, \\
   It allows them to put too many words on a line when they'd rather be
   forced to be terse.
\end{verse}

Mathematical formulas may also be displayed.  A
displayed formula
is
one-line long; multiline
formulas require special formatting instructions.
   \[  \ip{\Gamma}{\psi'} = x'' + y^{2} + z_{i}^{n}\]
Don't start a paragraph with a displayed equation,
nor make one a paragraph by itself.

\end{document}               % End of document.

\section{Bench Huge Section 537}
% synthetic bench section 537
% This is a sample LaTeX input file.  (Version of 12 August 2004.)
%
% A '%' character causes TeX to ignore all remaining text on the line,
% and is used for comments like this one.

\documentclass{article}      % Specifies the document class

                             % The preamble begins here.
\title{An Example Document}  % Declares the document's title.
\author{Leslie Lamport}      % Declares the author's name.
\date{January 21, 1994}      % Deleting this command produces today's date.

\newcommand{\ip}[2]{(#1, #2)}
                             % Defines \ip{arg1}{arg2} to mean
                             % (arg1, arg2).

%\newcommand{\ip}[2]{\langle #1 | #2\rangle}
                             % This is an alternative definition of
                             % \ip that is commented out.

\begin{document}             % End of preamble and beginning of text.

\maketitle                   % Produces the title.

This is an example input file.  Comparing it with
the output it generates can show you how to
produce a simple document of your own.

\section{Ordinary Text}      % Produces section heading.  Lower-level
                             % sections are begun with similar
                             % \subsection and \subsubsection commands.

The ends  of words and sentences are marked
  by   spaces. It  doesn't matter how many
spaces    you type; one is as good as 100.  The
end of   a line counts as a space.

One   or more   blank lines denote the  end
of  a paragraph.

Since any number of consecutive spaces are treated
like a single one, the formatting of the input
file makes no difference to
      \LaTeX,                % The \LaTeX command generates the LaTeX logo.
but it makes a difference to you.  When you use
\LaTeX, making your input file as easy to read
as possible will be a great help as you write
your document and when you change it.  This sample
file shows how you can add comments to your own input
file.

Because printing is different from typewriting,
there are a number of things that you have to do
differently when preparing an input file than if
you were just typing the document directly.
Quotation marks like
       ``this''
have to be handled specially, as do quotes within
quotes:
       ``\,`this'            % \, separates the double and single quote.
        is what I just
        wrote, not  `that'\,''.

Dashes come in three sizes: an
       intra-word
dash, a medium dash for number ranges like
       1--2,
and a punctuation
       dash---like
this.

A sentence-ending space should be larger than the
space between words within a sentence.  You
sometimes have to type special commands in
conjunction with punctuation characters to get
this right, as in the following sentence.
       Gnats, gnus, etc.\ all  % `\ ' makes an inter-word space.
       begin with G\@.         % \@ marks end-of-sentence punctuation.
You should check the spaces after periods when
reading your output to make sure you haven't
forgotten any special cases.  Generating an
ellipsis
       \ldots\               % `\ ' is needed after `\ldots' because TeX
                             % ignores spaces after command names like \ldots
                             % made from \ + letters.
                             %
                             % Note how a `%' character causes TeX to ignore
                             % the end of the input line, so these blank lines
                             % do not start a new paragraph.
                             %
with the right spacing around the periods requires
a special command.

\LaTeX\ interprets some common characters as
commands, so you must type special commands to
generate them.  These characters include the
following:
       \$ \& \% \# \{ and \}.

In printing, text is usually emphasized with an
       \emph{italic}
type style.

\begin{em}
   A long segment of text can also be emphasized
   in this way.  Text within such a segment can be
   given \emph{additional} emphasis.
\end{em}

It is sometimes necessary to prevent \LaTeX\ from
breaking a line where it might otherwise do so.
This may be at a space, as between the ``Mr.''\ and
``Jones'' in
       ``Mr.~Jones'',        % ~ produces an unbreakable interword space.
or within a word---especially when the word is a
symbol like
       \mbox{\emph{itemnum}}
that makes little sense when hyphenated across
lines.

Footnotes\footnote{This is an example of a footnote.}
pose no problem.

\LaTeX\ is good at typesetting mathematical formulas
like
       \( x-3y + z = 7 \)
or
       \( a_{1} > x^{2n} + y^{2n} > x' \)
or
       \( \ip{A}{B} = \sum_{i} a_{i} b_{i} \).
The spaces you type in a formula are
ignored.  Remember that a letter like
       $x$                   % $ ... $  and  \( ... \)  are equivalent
is a formula when it denotes a mathematical
symbol, and it should be typed as one.

\section{Displayed Text}

Text is displayed by indenting it from the left
margin.  Quotations are commonly displayed.  There
are short quotations
\begin{quote}
   This is a short quotation.  It consists of a
   single paragraph of text.  See how it is formatted.
\end{quote}
and longer ones.
\begin{quotation}
   This is a longer quotation.  It consists of two
   paragraphs of text, neither of which are
   particularly interesting.

   This is the second paragraph of the quotation.  It
   is just as dull as the first paragraph.
\end{quotation}
Another frequently-displayed structure is a list.
The following is an example of an \emph{itemized}
list.
\begin{itemize}
   \item This is the first item of an itemized list.
         Each item in the list is marked with a ``tick''.
         You don't have to worry about what kind of tick
         mark is used.

   \item This is the second item of the list.  It
         contains another list nested inside it.  The inner
         list is an \emph{enumerated} list.
         \begin{enumerate}
            \item This is the first item of an enumerated
                  list that is nested within the itemized list.

            \item This is the second item of the inner list.
                  \LaTeX\ allows you to nest lists deeper than
                  you really should.
         \end{enumerate}
         This is the rest of the second item of the outer
         list.  It is no more interesting than any other
         part of the item.
   \item This is the third item of the list.
\end{itemize}
You can even display poetry.
\begin{verse}
   There is an environment
    for verse \\             % The \\ command separates lines
   Whose features some poets % within a stanza.
   will curse.

                             % One or more blank lines separate stanzas.

   For instead of making\\
   Them do \emph{all} line breaking, \\
   It allows them to put too many words on a line when they'd rather be
   forced to be terse.
\end{verse}

Mathematical formulas may also be displayed.  A
displayed formula
is
one-line long; multiline
formulas require special formatting instructions.
   \[  \ip{\Gamma}{\psi'} = x'' + y^{2} + z_{i}^{n}\]
Don't start a paragraph with a displayed equation,
nor make one a paragraph by itself.

\end{document}               % End of document.

\section{Bench Huge Section 538}
% synthetic bench section 538
% This is a sample LaTeX input file.  (Version of 12 August 2004.)
%
% A '%' character causes TeX to ignore all remaining text on the line,
% and is used for comments like this one.

\documentclass{article}      % Specifies the document class

                             % The preamble begins here.
\title{An Example Document}  % Declares the document's title.
\author{Leslie Lamport}      % Declares the author's name.
\date{January 21, 1994}      % Deleting this command produces today's date.

\newcommand{\ip}[2]{(#1, #2)}
                             % Defines \ip{arg1}{arg2} to mean
                             % (arg1, arg2).

%\newcommand{\ip}[2]{\langle #1 | #2\rangle}
                             % This is an alternative definition of
                             % \ip that is commented out.

\begin{document}             % End of preamble and beginning of text.

\maketitle                   % Produces the title.

This is an example input file.  Comparing it with
the output it generates can show you how to
produce a simple document of your own.

\section{Ordinary Text}      % Produces section heading.  Lower-level
                             % sections are begun with similar
                             % \subsection and \subsubsection commands.

The ends  of words and sentences are marked
  by   spaces. It  doesn't matter how many
spaces    you type; one is as good as 100.  The
end of   a line counts as a space.

One   or more   blank lines denote the  end
of  a paragraph.

Since any number of consecutive spaces are treated
like a single one, the formatting of the input
file makes no difference to
      \LaTeX,                % The \LaTeX command generates the LaTeX logo.
but it makes a difference to you.  When you use
\LaTeX, making your input file as easy to read
as possible will be a great help as you write
your document and when you change it.  This sample
file shows how you can add comments to your own input
file.

Because printing is different from typewriting,
there are a number of things that you have to do
differently when preparing an input file than if
you were just typing the document directly.
Quotation marks like
       ``this''
have to be handled specially, as do quotes within
quotes:
       ``\,`this'            % \, separates the double and single quote.
        is what I just
        wrote, not  `that'\,''.

Dashes come in three sizes: an
       intra-word
dash, a medium dash for number ranges like
       1--2,
and a punctuation
       dash---like
this.

A sentence-ending space should be larger than the
space between words within a sentence.  You
sometimes have to type special commands in
conjunction with punctuation characters to get
this right, as in the following sentence.
       Gnats, gnus, etc.\ all  % `\ ' makes an inter-word space.
       begin with G\@.         % \@ marks end-of-sentence punctuation.
You should check the spaces after periods when
reading your output to make sure you haven't
forgotten any special cases.  Generating an
ellipsis
       \ldots\               % `\ ' is needed after `\ldots' because TeX
                             % ignores spaces after command names like \ldots
                             % made from \ + letters.
                             %
                             % Note how a `%' character causes TeX to ignore
                             % the end of the input line, so these blank lines
                             % do not start a new paragraph.
                             %
with the right spacing around the periods requires
a special command.

\LaTeX\ interprets some common characters as
commands, so you must type special commands to
generate them.  These characters include the
following:
       \$ \& \% \# \{ and \}.

In printing, text is usually emphasized with an
       \emph{italic}
type style.

\begin{em}
   A long segment of text can also be emphasized
   in this way.  Text within such a segment can be
   given \emph{additional} emphasis.
\end{em}

It is sometimes necessary to prevent \LaTeX\ from
breaking a line where it might otherwise do so.
This may be at a space, as between the ``Mr.''\ and
``Jones'' in
       ``Mr.~Jones'',        % ~ produces an unbreakable interword space.
or within a word---especially when the word is a
symbol like
       \mbox{\emph{itemnum}}
that makes little sense when hyphenated across
lines.

Footnotes\footnote{This is an example of a footnote.}
pose no problem.

\LaTeX\ is good at typesetting mathematical formulas
like
       \( x-3y + z = 7 \)
or
       \( a_{1} > x^{2n} + y^{2n} > x' \)
or
       \( \ip{A}{B} = \sum_{i} a_{i} b_{i} \).
The spaces you type in a formula are
ignored.  Remember that a letter like
       $x$                   % $ ... $  and  \( ... \)  are equivalent
is a formula when it denotes a mathematical
symbol, and it should be typed as one.

\section{Displayed Text}

Text is displayed by indenting it from the left
margin.  Quotations are commonly displayed.  There
are short quotations
\begin{quote}
   This is a short quotation.  It consists of a
   single paragraph of text.  See how it is formatted.
\end{quote}
and longer ones.
\begin{quotation}
   This is a longer quotation.  It consists of two
   paragraphs of text, neither of which are
   particularly interesting.

   This is the second paragraph of the quotation.  It
   is just as dull as the first paragraph.
\end{quotation}
Another frequently-displayed structure is a list.
The following is an example of an \emph{itemized}
list.
\begin{itemize}
   \item This is the first item of an itemized list.
         Each item in the list is marked with a ``tick''.
         You don't have to worry about what kind of tick
         mark is used.

   \item This is the second item of the list.  It
         contains another list nested inside it.  The inner
         list is an \emph{enumerated} list.
         \begin{enumerate}
            \item This is the first item of an enumerated
                  list that is nested within the itemized list.

            \item This is the second item of the inner list.
                  \LaTeX\ allows you to nest lists deeper than
                  you really should.
         \end{enumerate}
         This is the rest of the second item of the outer
         list.  It is no more interesting than any other
         part of the item.
   \item This is the third item of the list.
\end{itemize}
You can even display poetry.
\begin{verse}
   There is an environment
    for verse \\             % The \\ command separates lines
   Whose features some poets % within a stanza.
   will curse.

                             % One or more blank lines separate stanzas.

   For instead of making\\
   Them do \emph{all} line breaking, \\
   It allows them to put too many words on a line when they'd rather be
   forced to be terse.
\end{verse}

Mathematical formulas may also be displayed.  A
displayed formula
is
one-line long; multiline
formulas require special formatting instructions.
   \[  \ip{\Gamma}{\psi'} = x'' + y^{2} + z_{i}^{n}\]
Don't start a paragraph with a displayed equation,
nor make one a paragraph by itself.

\end{document}               % End of document.

\section{Bench Huge Section 539}
% synthetic bench section 539
% This is a sample LaTeX input file.  (Version of 12 August 2004.)
%
% A '%' character causes TeX to ignore all remaining text on the line,
% and is used for comments like this one.

\documentclass{article}      % Specifies the document class

                             % The preamble begins here.
\title{An Example Document}  % Declares the document's title.
\author{Leslie Lamport}      % Declares the author's name.
\date{January 21, 1994}      % Deleting this command produces today's date.

\newcommand{\ip}[2]{(#1, #2)}
                             % Defines \ip{arg1}{arg2} to mean
                             % (arg1, arg2).

%\newcommand{\ip}[2]{\langle #1 | #2\rangle}
                             % This is an alternative definition of
                             % \ip that is commented out.

\begin{document}             % End of preamble and beginning of text.

\maketitle                   % Produces the title.

This is an example input file.  Comparing it with
the output it generates can show you how to
produce a simple document of your own.

\section{Ordinary Text}      % Produces section heading.  Lower-level
                             % sections are begun with similar
                             % \subsection and \subsubsection commands.

The ends  of words and sentences are marked
  by   spaces. It  doesn't matter how many
spaces    you type; one is as good as 100.  The
end of   a line counts as a space.

One   or more   blank lines denote the  end
of  a paragraph.

Since any number of consecutive spaces are treated
like a single one, the formatting of the input
file makes no difference to
      \LaTeX,                % The \LaTeX command generates the LaTeX logo.
but it makes a difference to you.  When you use
\LaTeX, making your input file as easy to read
as possible will be a great help as you write
your document and when you change it.  This sample
file shows how you can add comments to your own input
file.

Because printing is different from typewriting,
there are a number of things that you have to do
differently when preparing an input file than if
you were just typing the document directly.
Quotation marks like
       ``this''
have to be handled specially, as do quotes within
quotes:
       ``\,`this'            % \, separates the double and single quote.
        is what I just
        wrote, not  `that'\,''.

Dashes come in three sizes: an
       intra-word
dash, a medium dash for number ranges like
       1--2,
and a punctuation
       dash---like
this.

A sentence-ending space should be larger than the
space between words within a sentence.  You
sometimes have to type special commands in
conjunction with punctuation characters to get
this right, as in the following sentence.
       Gnats, gnus, etc.\ all  % `\ ' makes an inter-word space.
       begin with G\@.         % \@ marks end-of-sentence punctuation.
You should check the spaces after periods when
reading your output to make sure you haven't
forgotten any special cases.  Generating an
ellipsis
       \ldots\               % `\ ' is needed after `\ldots' because TeX
                             % ignores spaces after command names like \ldots
                             % made from \ + letters.
                             %
                             % Note how a `%' character causes TeX to ignore
                             % the end of the input line, so these blank lines
                             % do not start a new paragraph.
                             %
with the right spacing around the periods requires
a special command.

\LaTeX\ interprets some common characters as
commands, so you must type special commands to
generate them.  These characters include the
following:
       \$ \& \% \# \{ and \}.

In printing, text is usually emphasized with an
       \emph{italic}
type style.

\begin{em}
   A long segment of text can also be emphasized
   in this way.  Text within such a segment can be
   given \emph{additional} emphasis.
\end{em}

It is sometimes necessary to prevent \LaTeX\ from
breaking a line where it might otherwise do so.
This may be at a space, as between the ``Mr.''\ and
``Jones'' in
       ``Mr.~Jones'',        % ~ produces an unbreakable interword space.
or within a word---especially when the word is a
symbol like
       \mbox{\emph{itemnum}}
that makes little sense when hyphenated across
lines.

Footnotes\footnote{This is an example of a footnote.}
pose no problem.

\LaTeX\ is good at typesetting mathematical formulas
like
       \( x-3y + z = 7 \)
or
       \( a_{1} > x^{2n} + y^{2n} > x' \)
or
       \( \ip{A}{B} = \sum_{i} a_{i} b_{i} \).
The spaces you type in a formula are
ignored.  Remember that a letter like
       $x$                   % $ ... $  and  \( ... \)  are equivalent
is a formula when it denotes a mathematical
symbol, and it should be typed as one.

\section{Displayed Text}

Text is displayed by indenting it from the left
margin.  Quotations are commonly displayed.  There
are short quotations
\begin{quote}
   This is a short quotation.  It consists of a
   single paragraph of text.  See how it is formatted.
\end{quote}
and longer ones.
\begin{quotation}
   This is a longer quotation.  It consists of two
   paragraphs of text, neither of which are
   particularly interesting.

   This is the second paragraph of the quotation.  It
   is just as dull as the first paragraph.
\end{quotation}
Another frequently-displayed structure is a list.
The following is an example of an \emph{itemized}
list.
\begin{itemize}
   \item This is the first item of an itemized list.
         Each item in the list is marked with a ``tick''.
         You don't have to worry about what kind of tick
         mark is used.

   \item This is the second item of the list.  It
         contains another list nested inside it.  The inner
         list is an \emph{enumerated} list.
         \begin{enumerate}
            \item This is the first item of an enumerated
                  list that is nested within the itemized list.

            \item This is the second item of the inner list.
                  \LaTeX\ allows you to nest lists deeper than
                  you really should.
         \end{enumerate}
         This is the rest of the second item of the outer
         list.  It is no more interesting than any other
         part of the item.
   \item This is the third item of the list.
\end{itemize}
You can even display poetry.
\begin{verse}
   There is an environment
    for verse \\             % The \\ command separates lines
   Whose features some poets % within a stanza.
   will curse.

                             % One or more blank lines separate stanzas.

   For instead of making\\
   Them do \emph{all} line breaking, \\
   It allows them to put too many words on a line when they'd rather be
   forced to be terse.
\end{verse}

Mathematical formulas may also be displayed.  A
displayed formula
is
one-line long; multiline
formulas require special formatting instructions.
   \[  \ip{\Gamma}{\psi'} = x'' + y^{2} + z_{i}^{n}\]
Don't start a paragraph with a displayed equation,
nor make one a paragraph by itself.

\end{document}               % End of document.

\section{Bench Huge Section 540}
% synthetic bench section 540
% This is a sample LaTeX input file.  (Version of 12 August 2004.)
%
% A '%' character causes TeX to ignore all remaining text on the line,
% and is used for comments like this one.

\documentclass{article}      % Specifies the document class

                             % The preamble begins here.
\title{An Example Document}  % Declares the document's title.
\author{Leslie Lamport}      % Declares the author's name.
\date{January 21, 1994}      % Deleting this command produces today's date.

\newcommand{\ip}[2]{(#1, #2)}
                             % Defines \ip{arg1}{arg2} to mean
                             % (arg1, arg2).

%\newcommand{\ip}[2]{\langle #1 | #2\rangle}
                             % This is an alternative definition of
                             % \ip that is commented out.

\begin{document}             % End of preamble and beginning of text.

\maketitle                   % Produces the title.

This is an example input file.  Comparing it with
the output it generates can show you how to
produce a simple document of your own.

\section{Ordinary Text}      % Produces section heading.  Lower-level
                             % sections are begun with similar
                             % \subsection and \subsubsection commands.

The ends  of words and sentences are marked
  by   spaces. It  doesn't matter how many
spaces    you type; one is as good as 100.  The
end of   a line counts as a space.

One   or more   blank lines denote the  end
of  a paragraph.

Since any number of consecutive spaces are treated
like a single one, the formatting of the input
file makes no difference to
      \LaTeX,                % The \LaTeX command generates the LaTeX logo.
but it makes a difference to you.  When you use
\LaTeX, making your input file as easy to read
as possible will be a great help as you write
your document and when you change it.  This sample
file shows how you can add comments to your own input
file.

Because printing is different from typewriting,
there are a number of things that you have to do
differently when preparing an input file than if
you were just typing the document directly.
Quotation marks like
       ``this''
have to be handled specially, as do quotes within
quotes:
       ``\,`this'            % \, separates the double and single quote.
        is what I just
        wrote, not  `that'\,''.

Dashes come in three sizes: an
       intra-word
dash, a medium dash for number ranges like
       1--2,
and a punctuation
       dash---like
this.

A sentence-ending space should be larger than the
space between words within a sentence.  You
sometimes have to type special commands in
conjunction with punctuation characters to get
this right, as in the following sentence.
       Gnats, gnus, etc.\ all  % `\ ' makes an inter-word space.
       begin with G\@.         % \@ marks end-of-sentence punctuation.
You should check the spaces after periods when
reading your output to make sure you haven't
forgotten any special cases.  Generating an
ellipsis
       \ldots\               % `\ ' is needed after `\ldots' because TeX
                             % ignores spaces after command names like \ldots
                             % made from \ + letters.
                             %
                             % Note how a `%' character causes TeX to ignore
                             % the end of the input line, so these blank lines
                             % do not start a new paragraph.
                             %
with the right spacing around the periods requires
a special command.

\LaTeX\ interprets some common characters as
commands, so you must type special commands to
generate them.  These characters include the
following:
       \$ \& \% \# \{ and \}.

In printing, text is usually emphasized with an
       \emph{italic}
type style.

\begin{em}
   A long segment of text can also be emphasized
   in this way.  Text within such a segment can be
   given \emph{additional} emphasis.
\end{em}

It is sometimes necessary to prevent \LaTeX\ from
breaking a line where it might otherwise do so.
This may be at a space, as between the ``Mr.''\ and
``Jones'' in
       ``Mr.~Jones'',        % ~ produces an unbreakable interword space.
or within a word---especially when the word is a
symbol like
       \mbox{\emph{itemnum}}
that makes little sense when hyphenated across
lines.

Footnotes\footnote{This is an example of a footnote.}
pose no problem.

\LaTeX\ is good at typesetting mathematical formulas
like
       \( x-3y + z = 7 \)
or
       \( a_{1} > x^{2n} + y^{2n} > x' \)
or
       \( \ip{A}{B} = \sum_{i} a_{i} b_{i} \).
The spaces you type in a formula are
ignored.  Remember that a letter like
       $x$                   % $ ... $  and  \( ... \)  are equivalent
is a formula when it denotes a mathematical
symbol, and it should be typed as one.

\section{Displayed Text}

Text is displayed by indenting it from the left
margin.  Quotations are commonly displayed.  There
are short quotations
\begin{quote}
   This is a short quotation.  It consists of a
   single paragraph of text.  See how it is formatted.
\end{quote}
and longer ones.
\begin{quotation}
   This is a longer quotation.  It consists of two
   paragraphs of text, neither of which are
   particularly interesting.

   This is the second paragraph of the quotation.  It
   is just as dull as the first paragraph.
\end{quotation}
Another frequently-displayed structure is a list.
The following is an example of an \emph{itemized}
list.
\begin{itemize}
   \item This is the first item of an itemized list.
         Each item in the list is marked with a ``tick''.
         You don't have to worry about what kind of tick
         mark is used.

   \item This is the second item of the list.  It
         contains another list nested inside it.  The inner
         list is an \emph{enumerated} list.
         \begin{enumerate}
            \item This is the first item of an enumerated
                  list that is nested within the itemized list.

            \item This is the second item of the inner list.
                  \LaTeX\ allows you to nest lists deeper than
                  you really should.
         \end{enumerate}
         This is the rest of the second item of the outer
         list.  It is no more interesting than any other
         part of the item.
   \item This is the third item of the list.
\end{itemize}
You can even display poetry.
\begin{verse}
   There is an environment
    for verse \\             % The \\ command separates lines
   Whose features some poets % within a stanza.
   will curse.

                             % One or more blank lines separate stanzas.

   For instead of making\\
   Them do \emph{all} line breaking, \\
   It allows them to put too many words on a line when they'd rather be
   forced to be terse.
\end{verse}

Mathematical formulas may also be displayed.  A
displayed formula
is
one-line long; multiline
formulas require special formatting instructions.
   \[  \ip{\Gamma}{\psi'} = x'' + y^{2} + z_{i}^{n}\]
Don't start a paragraph with a displayed equation,
nor make one a paragraph by itself.

\end{document}               % End of document.

\section{Bench Huge Section 541}
% synthetic bench section 541
% This is a sample LaTeX input file.  (Version of 12 August 2004.)
%
% A '%' character causes TeX to ignore all remaining text on the line,
% and is used for comments like this one.

\documentclass{article}      % Specifies the document class

                             % The preamble begins here.
\title{An Example Document}  % Declares the document's title.
\author{Leslie Lamport}      % Declares the author's name.
\date{January 21, 1994}      % Deleting this command produces today's date.

\newcommand{\ip}[2]{(#1, #2)}
                             % Defines \ip{arg1}{arg2} to mean
                             % (arg1, arg2).

%\newcommand{\ip}[2]{\langle #1 | #2\rangle}
                             % This is an alternative definition of
                             % \ip that is commented out.

\begin{document}             % End of preamble and beginning of text.

\maketitle                   % Produces the title.

This is an example input file.  Comparing it with
the output it generates can show you how to
produce a simple document of your own.

\section{Ordinary Text}      % Produces section heading.  Lower-level
                             % sections are begun with similar
                             % \subsection and \subsubsection commands.

The ends  of words and sentences are marked
  by   spaces. It  doesn't matter how many
spaces    you type; one is as good as 100.  The
end of   a line counts as a space.

One   or more   blank lines denote the  end
of  a paragraph.

Since any number of consecutive spaces are treated
like a single one, the formatting of the input
file makes no difference to
      \LaTeX,                % The \LaTeX command generates the LaTeX logo.
but it makes a difference to you.  When you use
\LaTeX, making your input file as easy to read
as possible will be a great help as you write
your document and when you change it.  This sample
file shows how you can add comments to your own input
file.

Because printing is different from typewriting,
there are a number of things that you have to do
differently when preparing an input file than if
you were just typing the document directly.
Quotation marks like
       ``this''
have to be handled specially, as do quotes within
quotes:
       ``\,`this'            % \, separates the double and single quote.
        is what I just
        wrote, not  `that'\,''.

Dashes come in three sizes: an
       intra-word
dash, a medium dash for number ranges like
       1--2,
and a punctuation
       dash---like
this.

A sentence-ending space should be larger than the
space between words within a sentence.  You
sometimes have to type special commands in
conjunction with punctuation characters to get
this right, as in the following sentence.
       Gnats, gnus, etc.\ all  % `\ ' makes an inter-word space.
       begin with G\@.         % \@ marks end-of-sentence punctuation.
You should check the spaces after periods when
reading your output to make sure you haven't
forgotten any special cases.  Generating an
ellipsis
       \ldots\               % `\ ' is needed after `\ldots' because TeX
                             % ignores spaces after command names like \ldots
                             % made from \ + letters.
                             %
                             % Note how a `%' character causes TeX to ignore
                             % the end of the input line, so these blank lines
                             % do not start a new paragraph.
                             %
with the right spacing around the periods requires
a special command.

\LaTeX\ interprets some common characters as
commands, so you must type special commands to
generate them.  These characters include the
following:
       \$ \& \% \# \{ and \}.

In printing, text is usually emphasized with an
       \emph{italic}
type style.

\begin{em}
   A long segment of text can also be emphasized
   in this way.  Text within such a segment can be
   given \emph{additional} emphasis.
\end{em}

It is sometimes necessary to prevent \LaTeX\ from
breaking a line where it might otherwise do so.
This may be at a space, as between the ``Mr.''\ and
``Jones'' in
       ``Mr.~Jones'',        % ~ produces an unbreakable interword space.
or within a word---especially when the word is a
symbol like
       \mbox{\emph{itemnum}}
that makes little sense when hyphenated across
lines.

Footnotes\footnote{This is an example of a footnote.}
pose no problem.

\LaTeX\ is good at typesetting mathematical formulas
like
       \( x-3y + z = 7 \)
or
       \( a_{1} > x^{2n} + y^{2n} > x' \)
or
       \( \ip{A}{B} = \sum_{i} a_{i} b_{i} \).
The spaces you type in a formula are
ignored.  Remember that a letter like
       $x$                   % $ ... $  and  \( ... \)  are equivalent
is a formula when it denotes a mathematical
symbol, and it should be typed as one.

\section{Displayed Text}

Text is displayed by indenting it from the left
margin.  Quotations are commonly displayed.  There
are short quotations
\begin{quote}
   This is a short quotation.  It consists of a
   single paragraph of text.  See how it is formatted.
\end{quote}
and longer ones.
\begin{quotation}
   This is a longer quotation.  It consists of two
   paragraphs of text, neither of which are
   particularly interesting.

   This is the second paragraph of the quotation.  It
   is just as dull as the first paragraph.
\end{quotation}
Another frequently-displayed structure is a list.
The following is an example of an \emph{itemized}
list.
\begin{itemize}
   \item This is the first item of an itemized list.
         Each item in the list is marked with a ``tick''.
         You don't have to worry about what kind of tick
         mark is used.

   \item This is the second item of the list.  It
         contains another list nested inside it.  The inner
         list is an \emph{enumerated} list.
         \begin{enumerate}
            \item This is the first item of an enumerated
                  list that is nested within the itemized list.

            \item This is the second item of the inner list.
                  \LaTeX\ allows you to nest lists deeper than
                  you really should.
         \end{enumerate}
         This is the rest of the second item of the outer
         list.  It is no more interesting than any other
         part of the item.
   \item This is the third item of the list.
\end{itemize}
You can even display poetry.
\begin{verse}
   There is an environment
    for verse \\             % The \\ command separates lines
   Whose features some poets % within a stanza.
   will curse.

                             % One or more blank lines separate stanzas.

   For instead of making\\
   Them do \emph{all} line breaking, \\
   It allows them to put too many words on a line when they'd rather be
   forced to be terse.
\end{verse}

Mathematical formulas may also be displayed.  A
displayed formula
is
one-line long; multiline
formulas require special formatting instructions.
   \[  \ip{\Gamma}{\psi'} = x'' + y^{2} + z_{i}^{n}\]
Don't start a paragraph with a displayed equation,
nor make one a paragraph by itself.

\end{document}               % End of document.

\section{Bench Huge Section 542}
% synthetic bench section 542
% This is a sample LaTeX input file.  (Version of 12 August 2004.)
%
% A '%' character causes TeX to ignore all remaining text on the line,
% and is used for comments like this one.

\documentclass{article}      % Specifies the document class

                             % The preamble begins here.
\title{An Example Document}  % Declares the document's title.
\author{Leslie Lamport}      % Declares the author's name.
\date{January 21, 1994}      % Deleting this command produces today's date.

\newcommand{\ip}[2]{(#1, #2)}
                             % Defines \ip{arg1}{arg2} to mean
                             % (arg1, arg2).

%\newcommand{\ip}[2]{\langle #1 | #2\rangle}
                             % This is an alternative definition of
                             % \ip that is commented out.

\begin{document}             % End of preamble and beginning of text.

\maketitle                   % Produces the title.

This is an example input file.  Comparing it with
the output it generates can show you how to
produce a simple document of your own.

\section{Ordinary Text}      % Produces section heading.  Lower-level
                             % sections are begun with similar
                             % \subsection and \subsubsection commands.

The ends  of words and sentences are marked
  by   spaces. It  doesn't matter how many
spaces    you type; one is as good as 100.  The
end of   a line counts as a space.

One   or more   blank lines denote the  end
of  a paragraph.

Since any number of consecutive spaces are treated
like a single one, the formatting of the input
file makes no difference to
      \LaTeX,                % The \LaTeX command generates the LaTeX logo.
but it makes a difference to you.  When you use
\LaTeX, making your input file as easy to read
as possible will be a great help as you write
your document and when you change it.  This sample
file shows how you can add comments to your own input
file.

Because printing is different from typewriting,
there are a number of things that you have to do
differently when preparing an input file than if
you were just typing the document directly.
Quotation marks like
       ``this''
have to be handled specially, as do quotes within
quotes:
       ``\,`this'            % \, separates the double and single quote.
        is what I just
        wrote, not  `that'\,''.

Dashes come in three sizes: an
       intra-word
dash, a medium dash for number ranges like
       1--2,
and a punctuation
       dash---like
this.

A sentence-ending space should be larger than the
space between words within a sentence.  You
sometimes have to type special commands in
conjunction with punctuation characters to get
this right, as in the following sentence.
       Gnats, gnus, etc.\ all  % `\ ' makes an inter-word space.
       begin with G\@.         % \@ marks end-of-sentence punctuation.
You should check the spaces after periods when
reading your output to make sure you haven't
forgotten any special cases.  Generating an
ellipsis
       \ldots\               % `\ ' is needed after `\ldots' because TeX
                             % ignores spaces after command names like \ldots
                             % made from \ + letters.
                             %
                             % Note how a `%' character causes TeX to ignore
                             % the end of the input line, so these blank lines
                             % do not start a new paragraph.
                             %
with the right spacing around the periods requires
a special command.

\LaTeX\ interprets some common characters as
commands, so you must type special commands to
generate them.  These characters include the
following:
       \$ \& \% \# \{ and \}.

In printing, text is usually emphasized with an
       \emph{italic}
type style.

\begin{em}
   A long segment of text can also be emphasized
   in this way.  Text within such a segment can be
   given \emph{additional} emphasis.
\end{em}

It is sometimes necessary to prevent \LaTeX\ from
breaking a line where it might otherwise do so.
This may be at a space, as between the ``Mr.''\ and
``Jones'' in
       ``Mr.~Jones'',        % ~ produces an unbreakable interword space.
or within a word---especially when the word is a
symbol like
       \mbox{\emph{itemnum}}
that makes little sense when hyphenated across
lines.

Footnotes\footnote{This is an example of a footnote.}
pose no problem.

\LaTeX\ is good at typesetting mathematical formulas
like
       \( x-3y + z = 7 \)
or
       \( a_{1} > x^{2n} + y^{2n} > x' \)
or
       \( \ip{A}{B} = \sum_{i} a_{i} b_{i} \).
The spaces you type in a formula are
ignored.  Remember that a letter like
       $x$                   % $ ... $  and  \( ... \)  are equivalent
is a formula when it denotes a mathematical
symbol, and it should be typed as one.

\section{Displayed Text}

Text is displayed by indenting it from the left
margin.  Quotations are commonly displayed.  There
are short quotations
\begin{quote}
   This is a short quotation.  It consists of a
   single paragraph of text.  See how it is formatted.
\end{quote}
and longer ones.
\begin{quotation}
   This is a longer quotation.  It consists of two
   paragraphs of text, neither of which are
   particularly interesting.

   This is the second paragraph of the quotation.  It
   is just as dull as the first paragraph.
\end{quotation}
Another frequently-displayed structure is a list.
The following is an example of an \emph{itemized}
list.
\begin{itemize}
   \item This is the first item of an itemized list.
         Each item in the list is marked with a ``tick''.
         You don't have to worry about what kind of tick
         mark is used.

   \item This is the second item of the list.  It
         contains another list nested inside it.  The inner
         list is an \emph{enumerated} list.
         \begin{enumerate}
            \item This is the first item of an enumerated
                  list that is nested within the itemized list.

            \item This is the second item of the inner list.
                  \LaTeX\ allows you to nest lists deeper than
                  you really should.
         \end{enumerate}
         This is the rest of the second item of the outer
         list.  It is no more interesting than any other
         part of the item.
   \item This is the third item of the list.
\end{itemize}
You can even display poetry.
\begin{verse}
   There is an environment
    for verse \\             % The \\ command separates lines
   Whose features some poets % within a stanza.
   will curse.

                             % One or more blank lines separate stanzas.

   For instead of making\\
   Them do \emph{all} line breaking, \\
   It allows them to put too many words on a line when they'd rather be
   forced to be terse.
\end{verse}

Mathematical formulas may also be displayed.  A
displayed formula
is
one-line long; multiline
formulas require special formatting instructions.
   \[  \ip{\Gamma}{\psi'} = x'' + y^{2} + z_{i}^{n}\]
Don't start a paragraph with a displayed equation,
nor make one a paragraph by itself.

\end{document}               % End of document.

\section{Bench Huge Section 543}
% synthetic bench section 543
% This is a sample LaTeX input file.  (Version of 12 August 2004.)
%
% A '%' character causes TeX to ignore all remaining text on the line,
% and is used for comments like this one.

\documentclass{article}      % Specifies the document class

                             % The preamble begins here.
\title{An Example Document}  % Declares the document's title.
\author{Leslie Lamport}      % Declares the author's name.
\date{January 21, 1994}      % Deleting this command produces today's date.

\newcommand{\ip}[2]{(#1, #2)}
                             % Defines \ip{arg1}{arg2} to mean
                             % (arg1, arg2).

%\newcommand{\ip}[2]{\langle #1 | #2\rangle}
                             % This is an alternative definition of
                             % \ip that is commented out.

\begin{document}             % End of preamble and beginning of text.

\maketitle                   % Produces the title.

This is an example input file.  Comparing it with
the output it generates can show you how to
produce a simple document of your own.

\section{Ordinary Text}      % Produces section heading.  Lower-level
                             % sections are begun with similar
                             % \subsection and \subsubsection commands.

The ends  of words and sentences are marked
  by   spaces. It  doesn't matter how many
spaces    you type; one is as good as 100.  The
end of   a line counts as a space.

One   or more   blank lines denote the  end
of  a paragraph.

Since any number of consecutive spaces are treated
like a single one, the formatting of the input
file makes no difference to
      \LaTeX,                % The \LaTeX command generates the LaTeX logo.
but it makes a difference to you.  When you use
\LaTeX, making your input file as easy to read
as possible will be a great help as you write
your document and when you change it.  This sample
file shows how you can add comments to your own input
file.

Because printing is different from typewriting,
there are a number of things that you have to do
differently when preparing an input file than if
you were just typing the document directly.
Quotation marks like
       ``this''
have to be handled specially, as do quotes within
quotes:
       ``\,`this'            % \, separates the double and single quote.
        is what I just
        wrote, not  `that'\,''.

Dashes come in three sizes: an
       intra-word
dash, a medium dash for number ranges like
       1--2,
and a punctuation
       dash---like
this.

A sentence-ending space should be larger than the
space between words within a sentence.  You
sometimes have to type special commands in
conjunction with punctuation characters to get
this right, as in the following sentence.
       Gnats, gnus, etc.\ all  % `\ ' makes an inter-word space.
       begin with G\@.         % \@ marks end-of-sentence punctuation.
You should check the spaces after periods when
reading your output to make sure you haven't
forgotten any special cases.  Generating an
ellipsis
       \ldots\               % `\ ' is needed after `\ldots' because TeX
                             % ignores spaces after command names like \ldots
                             % made from \ + letters.
                             %
                             % Note how a `%' character causes TeX to ignore
                             % the end of the input line, so these blank lines
                             % do not start a new paragraph.
                             %
with the right spacing around the periods requires
a special command.

\LaTeX\ interprets some common characters as
commands, so you must type special commands to
generate them.  These characters include the
following:
       \$ \& \% \# \{ and \}.

In printing, text is usually emphasized with an
       \emph{italic}
type style.

\begin{em}
   A long segment of text can also be emphasized
   in this way.  Text within such a segment can be
   given \emph{additional} emphasis.
\end{em}

It is sometimes necessary to prevent \LaTeX\ from
breaking a line where it might otherwise do so.
This may be at a space, as between the ``Mr.''\ and
``Jones'' in
       ``Mr.~Jones'',        % ~ produces an unbreakable interword space.
or within a word---especially when the word is a
symbol like
       \mbox{\emph{itemnum}}
that makes little sense when hyphenated across
lines.

Footnotes\footnote{This is an example of a footnote.}
pose no problem.

\LaTeX\ is good at typesetting mathematical formulas
like
       \( x-3y + z = 7 \)
or
       \( a_{1} > x^{2n} + y^{2n} > x' \)
or
       \( \ip{A}{B} = \sum_{i} a_{i} b_{i} \).
The spaces you type in a formula are
ignored.  Remember that a letter like
       $x$                   % $ ... $  and  \( ... \)  are equivalent
is a formula when it denotes a mathematical
symbol, and it should be typed as one.

\section{Displayed Text}

Text is displayed by indenting it from the left
margin.  Quotations are commonly displayed.  There
are short quotations
\begin{quote}
   This is a short quotation.  It consists of a
   single paragraph of text.  See how it is formatted.
\end{quote}
and longer ones.
\begin{quotation}
   This is a longer quotation.  It consists of two
   paragraphs of text, neither of which are
   particularly interesting.

   This is the second paragraph of the quotation.  It
   is just as dull as the first paragraph.
\end{quotation}
Another frequently-displayed structure is a list.
The following is an example of an \emph{itemized}
list.
\begin{itemize}
   \item This is the first item of an itemized list.
         Each item in the list is marked with a ``tick''.
         You don't have to worry about what kind of tick
         mark is used.

   \item This is the second item of the list.  It
         contains another list nested inside it.  The inner
         list is an \emph{enumerated} list.
         \begin{enumerate}
            \item This is the first item of an enumerated
                  list that is nested within the itemized list.

            \item This is the second item of the inner list.
                  \LaTeX\ allows you to nest lists deeper than
                  you really should.
         \end{enumerate}
         This is the rest of the second item of the outer
         list.  It is no more interesting than any other
         part of the item.
   \item This is the third item of the list.
\end{itemize}
You can even display poetry.
\begin{verse}
   There is an environment
    for verse \\             % The \\ command separates lines
   Whose features some poets % within a stanza.
   will curse.

                             % One or more blank lines separate stanzas.

   For instead of making\\
   Them do \emph{all} line breaking, \\
   It allows them to put too many words on a line when they'd rather be
   forced to be terse.
\end{verse}

Mathematical formulas may also be displayed.  A
displayed formula
is
one-line long; multiline
formulas require special formatting instructions.
   \[  \ip{\Gamma}{\psi'} = x'' + y^{2} + z_{i}^{n}\]
Don't start a paragraph with a displayed equation,
nor make one a paragraph by itself.

\end{document}               % End of document.

\section{Bench Huge Section 544}
% synthetic bench section 544
% This is a sample LaTeX input file.  (Version of 12 August 2004.)
%
% A '%' character causes TeX to ignore all remaining text on the line,
% and is used for comments like this one.

\documentclass{article}      % Specifies the document class

                             % The preamble begins here.
\title{An Example Document}  % Declares the document's title.
\author{Leslie Lamport}      % Declares the author's name.
\date{January 21, 1994}      % Deleting this command produces today's date.

\newcommand{\ip}[2]{(#1, #2)}
                             % Defines \ip{arg1}{arg2} to mean
                             % (arg1, arg2).

%\newcommand{\ip}[2]{\langle #1 | #2\rangle}
                             % This is an alternative definition of
                             % \ip that is commented out.

\begin{document}             % End of preamble and beginning of text.

\maketitle                   % Produces the title.

This is an example input file.  Comparing it with
the output it generates can show you how to
produce a simple document of your own.

\section{Ordinary Text}      % Produces section heading.  Lower-level
                             % sections are begun with similar
                             % \subsection and \subsubsection commands.

The ends  of words and sentences are marked
  by   spaces. It  doesn't matter how many
spaces    you type; one is as good as 100.  The
end of   a line counts as a space.

One   or more   blank lines denote the  end
of  a paragraph.

Since any number of consecutive spaces are treated
like a single one, the formatting of the input
file makes no difference to
      \LaTeX,                % The \LaTeX command generates the LaTeX logo.
but it makes a difference to you.  When you use
\LaTeX, making your input file as easy to read
as possible will be a great help as you write
your document and when you change it.  This sample
file shows how you can add comments to your own input
file.

Because printing is different from typewriting,
there are a number of things that you have to do
differently when preparing an input file than if
you were just typing the document directly.
Quotation marks like
       ``this''
have to be handled specially, as do quotes within
quotes:
       ``\,`this'            % \, separates the double and single quote.
        is what I just
        wrote, not  `that'\,''.

Dashes come in three sizes: an
       intra-word
dash, a medium dash for number ranges like
       1--2,
and a punctuation
       dash---like
this.

A sentence-ending space should be larger than the
space between words within a sentence.  You
sometimes have to type special commands in
conjunction with punctuation characters to get
this right, as in the following sentence.
       Gnats, gnus, etc.\ all  % `\ ' makes an inter-word space.
       begin with G\@.         % \@ marks end-of-sentence punctuation.
You should check the spaces after periods when
reading your output to make sure you haven't
forgotten any special cases.  Generating an
ellipsis
       \ldots\               % `\ ' is needed after `\ldots' because TeX
                             % ignores spaces after command names like \ldots
                             % made from \ + letters.
                             %
                             % Note how a `%' character causes TeX to ignore
                             % the end of the input line, so these blank lines
                             % do not start a new paragraph.
                             %
with the right spacing around the periods requires
a special command.

\LaTeX\ interprets some common characters as
commands, so you must type special commands to
generate them.  These characters include the
following:
       \$ \& \% \# \{ and \}.

In printing, text is usually emphasized with an
       \emph{italic}
type style.

\begin{em}
   A long segment of text can also be emphasized
   in this way.  Text within such a segment can be
   given \emph{additional} emphasis.
\end{em}

It is sometimes necessary to prevent \LaTeX\ from
breaking a line where it might otherwise do so.
This may be at a space, as between the ``Mr.''\ and
``Jones'' in
       ``Mr.~Jones'',        % ~ produces an unbreakable interword space.
or within a word---especially when the word is a
symbol like
       \mbox{\emph{itemnum}}
that makes little sense when hyphenated across
lines.

Footnotes\footnote{This is an example of a footnote.}
pose no problem.

\LaTeX\ is good at typesetting mathematical formulas
like
       \( x-3y + z = 7 \)
or
       \( a_{1} > x^{2n} + y^{2n} > x' \)
or
       \( \ip{A}{B} = \sum_{i} a_{i} b_{i} \).
The spaces you type in a formula are
ignored.  Remember that a letter like
       $x$                   % $ ... $  and  \( ... \)  are equivalent
is a formula when it denotes a mathematical
symbol, and it should be typed as one.

\section{Displayed Text}

Text is displayed by indenting it from the left
margin.  Quotations are commonly displayed.  There
are short quotations
\begin{quote}
   This is a short quotation.  It consists of a
   single paragraph of text.  See how it is formatted.
\end{quote}
and longer ones.
\begin{quotation}
   This is a longer quotation.  It consists of two
   paragraphs of text, neither of which are
   particularly interesting.

   This is the second paragraph of the quotation.  It
   is just as dull as the first paragraph.
\end{quotation}
Another frequently-displayed structure is a list.
The following is an example of an \emph{itemized}
list.
\begin{itemize}
   \item This is the first item of an itemized list.
         Each item in the list is marked with a ``tick''.
         You don't have to worry about what kind of tick
         mark is used.

   \item This is the second item of the list.  It
         contains another list nested inside it.  The inner
         list is an \emph{enumerated} list.
         \begin{enumerate}
            \item This is the first item of an enumerated
                  list that is nested within the itemized list.

            \item This is the second item of the inner list.
                  \LaTeX\ allows you to nest lists deeper than
                  you really should.
         \end{enumerate}
         This is the rest of the second item of the outer
         list.  It is no more interesting than any other
         part of the item.
   \item This is the third item of the list.
\end{itemize}
You can even display poetry.
\begin{verse}
   There is an environment
    for verse \\             % The \\ command separates lines
   Whose features some poets % within a stanza.
   will curse.

                             % One or more blank lines separate stanzas.

   For instead of making\\
   Them do \emph{all} line breaking, \\
   It allows them to put too many words on a line when they'd rather be
   forced to be terse.
\end{verse}

Mathematical formulas may also be displayed.  A
displayed formula
is
one-line long; multiline
formulas require special formatting instructions.
   \[  \ip{\Gamma}{\psi'} = x'' + y^{2} + z_{i}^{n}\]
Don't start a paragraph with a displayed equation,
nor make one a paragraph by itself.

\end{document}               % End of document.

\section{Bench Huge Section 545}
% synthetic bench section 545
% This is a sample LaTeX input file.  (Version of 12 August 2004.)
%
% A '%' character causes TeX to ignore all remaining text on the line,
% and is used for comments like this one.

\documentclass{article}      % Specifies the document class

                             % The preamble begins here.
\title{An Example Document}  % Declares the document's title.
\author{Leslie Lamport}      % Declares the author's name.
\date{January 21, 1994}      % Deleting this command produces today's date.

\newcommand{\ip}[2]{(#1, #2)}
                             % Defines \ip{arg1}{arg2} to mean
                             % (arg1, arg2).

%\newcommand{\ip}[2]{\langle #1 | #2\rangle}
                             % This is an alternative definition of
                             % \ip that is commented out.

\begin{document}             % End of preamble and beginning of text.

\maketitle                   % Produces the title.

This is an example input file.  Comparing it with
the output it generates can show you how to
produce a simple document of your own.

\section{Ordinary Text}      % Produces section heading.  Lower-level
                             % sections are begun with similar
                             % \subsection and \subsubsection commands.

The ends  of words and sentences are marked
  by   spaces. It  doesn't matter how many
spaces    you type; one is as good as 100.  The
end of   a line counts as a space.

One   or more   blank lines denote the  end
of  a paragraph.

Since any number of consecutive spaces are treated
like a single one, the formatting of the input
file makes no difference to
      \LaTeX,                % The \LaTeX command generates the LaTeX logo.
but it makes a difference to you.  When you use
\LaTeX, making your input file as easy to read
as possible will be a great help as you write
your document and when you change it.  This sample
file shows how you can add comments to your own input
file.

Because printing is different from typewriting,
there are a number of things that you have to do
differently when preparing an input file than if
you were just typing the document directly.
Quotation marks like
       ``this''
have to be handled specially, as do quotes within
quotes:
       ``\,`this'            % \, separates the double and single quote.
        is what I just
        wrote, not  `that'\,''.

Dashes come in three sizes: an
       intra-word
dash, a medium dash for number ranges like
       1--2,
and a punctuation
       dash---like
this.

A sentence-ending space should be larger than the
space between words within a sentence.  You
sometimes have to type special commands in
conjunction with punctuation characters to get
this right, as in the following sentence.
       Gnats, gnus, etc.\ all  % `\ ' makes an inter-word space.
       begin with G\@.         % \@ marks end-of-sentence punctuation.
You should check the spaces after periods when
reading your output to make sure you haven't
forgotten any special cases.  Generating an
ellipsis
       \ldots\               % `\ ' is needed after `\ldots' because TeX
                             % ignores spaces after command names like \ldots
                             % made from \ + letters.
                             %
                             % Note how a `%' character causes TeX to ignore
                             % the end of the input line, so these blank lines
                             % do not start a new paragraph.
                             %
with the right spacing around the periods requires
a special command.

\LaTeX\ interprets some common characters as
commands, so you must type special commands to
generate them.  These characters include the
following:
       \$ \& \% \# \{ and \}.

In printing, text is usually emphasized with an
       \emph{italic}
type style.

\begin{em}
   A long segment of text can also be emphasized
   in this way.  Text within such a segment can be
   given \emph{additional} emphasis.
\end{em}

It is sometimes necessary to prevent \LaTeX\ from
breaking a line where it might otherwise do so.
This may be at a space, as between the ``Mr.''\ and
``Jones'' in
       ``Mr.~Jones'',        % ~ produces an unbreakable interword space.
or within a word---especially when the word is a
symbol like
       \mbox{\emph{itemnum}}
that makes little sense when hyphenated across
lines.

Footnotes\footnote{This is an example of a footnote.}
pose no problem.

\LaTeX\ is good at typesetting mathematical formulas
like
       \( x-3y + z = 7 \)
or
       \( a_{1} > x^{2n} + y^{2n} > x' \)
or
       \( \ip{A}{B} = \sum_{i} a_{i} b_{i} \).
The spaces you type in a formula are
ignored.  Remember that a letter like
       $x$                   % $ ... $  and  \( ... \)  are equivalent
is a formula when it denotes a mathematical
symbol, and it should be typed as one.

\section{Displayed Text}

Text is displayed by indenting it from the left
margin.  Quotations are commonly displayed.  There
are short quotations
\begin{quote}
   This is a short quotation.  It consists of a
   single paragraph of text.  See how it is formatted.
\end{quote}
and longer ones.
\begin{quotation}
   This is a longer quotation.  It consists of two
   paragraphs of text, neither of which are
   particularly interesting.

   This is the second paragraph of the quotation.  It
   is just as dull as the first paragraph.
\end{quotation}
Another frequently-displayed structure is a list.
The following is an example of an \emph{itemized}
list.
\begin{itemize}
   \item This is the first item of an itemized list.
         Each item in the list is marked with a ``tick''.
         You don't have to worry about what kind of tick
         mark is used.

   \item This is the second item of the list.  It
         contains another list nested inside it.  The inner
         list is an \emph{enumerated} list.
         \begin{enumerate}
            \item This is the first item of an enumerated
                  list that is nested within the itemized list.

            \item This is the second item of the inner list.
                  \LaTeX\ allows you to nest lists deeper than
                  you really should.
         \end{enumerate}
         This is the rest of the second item of the outer
         list.  It is no more interesting than any other
         part of the item.
   \item This is the third item of the list.
\end{itemize}
You can even display poetry.
\begin{verse}
   There is an environment
    for verse \\             % The \\ command separates lines
   Whose features some poets % within a stanza.
   will curse.

                             % One or more blank lines separate stanzas.

   For instead of making\\
   Them do \emph{all} line breaking, \\
   It allows them to put too many words on a line when they'd rather be
   forced to be terse.
\end{verse}

Mathematical formulas may also be displayed.  A
displayed formula
is
one-line long; multiline
formulas require special formatting instructions.
   \[  \ip{\Gamma}{\psi'} = x'' + y^{2} + z_{i}^{n}\]
Don't start a paragraph with a displayed equation,
nor make one a paragraph by itself.

\end{document}               % End of document.

\section{Bench Huge Section 546}
% synthetic bench section 546
% This is a sample LaTeX input file.  (Version of 12 August 2004.)
%
% A '%' character causes TeX to ignore all remaining text on the line,
% and is used for comments like this one.

\documentclass{article}      % Specifies the document class

                             % The preamble begins here.
\title{An Example Document}  % Declares the document's title.
\author{Leslie Lamport}      % Declares the author's name.
\date{January 21, 1994}      % Deleting this command produces today's date.

\newcommand{\ip}[2]{(#1, #2)}
                             % Defines \ip{arg1}{arg2} to mean
                             % (arg1, arg2).

%\newcommand{\ip}[2]{\langle #1 | #2\rangle}
                             % This is an alternative definition of
                             % \ip that is commented out.

\begin{document}             % End of preamble and beginning of text.

\maketitle                   % Produces the title.

This is an example input file.  Comparing it with
the output it generates can show you how to
produce a simple document of your own.

\section{Ordinary Text}      % Produces section heading.  Lower-level
                             % sections are begun with similar
                             % \subsection and \subsubsection commands.

The ends  of words and sentences are marked
  by   spaces. It  doesn't matter how many
spaces    you type; one is as good as 100.  The
end of   a line counts as a space.

One   or more   blank lines denote the  end
of  a paragraph.

Since any number of consecutive spaces are treated
like a single one, the formatting of the input
file makes no difference to
      \LaTeX,                % The \LaTeX command generates the LaTeX logo.
but it makes a difference to you.  When you use
\LaTeX, making your input file as easy to read
as possible will be a great help as you write
your document and when you change it.  This sample
file shows how you can add comments to your own input
file.

Because printing is different from typewriting,
there are a number of things that you have to do
differently when preparing an input file than if
you were just typing the document directly.
Quotation marks like
       ``this''
have to be handled specially, as do quotes within
quotes:
       ``\,`this'            % \, separates the double and single quote.
        is what I just
        wrote, not  `that'\,''.

Dashes come in three sizes: an
       intra-word
dash, a medium dash for number ranges like
       1--2,
and a punctuation
       dash---like
this.

A sentence-ending space should be larger than the
space between words within a sentence.  You
sometimes have to type special commands in
conjunction with punctuation characters to get
this right, as in the following sentence.
       Gnats, gnus, etc.\ all  % `\ ' makes an inter-word space.
       begin with G\@.         % \@ marks end-of-sentence punctuation.
You should check the spaces after periods when
reading your output to make sure you haven't
forgotten any special cases.  Generating an
ellipsis
       \ldots\               % `\ ' is needed after `\ldots' because TeX
                             % ignores spaces after command names like \ldots
                             % made from \ + letters.
                             %
                             % Note how a `%' character causes TeX to ignore
                             % the end of the input line, so these blank lines
                             % do not start a new paragraph.
                             %
with the right spacing around the periods requires
a special command.

\LaTeX\ interprets some common characters as
commands, so you must type special commands to
generate them.  These characters include the
following:
       \$ \& \% \# \{ and \}.

In printing, text is usually emphasized with an
       \emph{italic}
type style.

\begin{em}
   A long segment of text can also be emphasized
   in this way.  Text within such a segment can be
   given \emph{additional} emphasis.
\end{em}

It is sometimes necessary to prevent \LaTeX\ from
breaking a line where it might otherwise do so.
This may be at a space, as between the ``Mr.''\ and
``Jones'' in
       ``Mr.~Jones'',        % ~ produces an unbreakable interword space.
or within a word---especially when the word is a
symbol like
       \mbox{\emph{itemnum}}
that makes little sense when hyphenated across
lines.

Footnotes\footnote{This is an example of a footnote.}
pose no problem.

\LaTeX\ is good at typesetting mathematical formulas
like
       \( x-3y + z = 7 \)
or
       \( a_{1} > x^{2n} + y^{2n} > x' \)
or
       \( \ip{A}{B} = \sum_{i} a_{i} b_{i} \).
The spaces you type in a formula are
ignored.  Remember that a letter like
       $x$                   % $ ... $  and  \( ... \)  are equivalent
is a formula when it denotes a mathematical
symbol, and it should be typed as one.

\section{Displayed Text}

Text is displayed by indenting it from the left
margin.  Quotations are commonly displayed.  There
are short quotations
\begin{quote}
   This is a short quotation.  It consists of a
   single paragraph of text.  See how it is formatted.
\end{quote}
and longer ones.
\begin{quotation}
   This is a longer quotation.  It consists of two
   paragraphs of text, neither of which are
   particularly interesting.

   This is the second paragraph of the quotation.  It
   is just as dull as the first paragraph.
\end{quotation}
Another frequently-displayed structure is a list.
The following is an example of an \emph{itemized}
list.
\begin{itemize}
   \item This is the first item of an itemized list.
         Each item in the list is marked with a ``tick''.
         You don't have to worry about what kind of tick
         mark is used.

   \item This is the second item of the list.  It
         contains another list nested inside it.  The inner
         list is an \emph{enumerated} list.
         \begin{enumerate}
            \item This is the first item of an enumerated
                  list that is nested within the itemized list.

            \item This is the second item of the inner list.
                  \LaTeX\ allows you to nest lists deeper than
                  you really should.
         \end{enumerate}
         This is the rest of the second item of the outer
         list.  It is no more interesting than any other
         part of the item.
   \item This is the third item of the list.
\end{itemize}
You can even display poetry.
\begin{verse}
   There is an environment
    for verse \\             % The \\ command separates lines
   Whose features some poets % within a stanza.
   will curse.

                             % One or more blank lines separate stanzas.

   For instead of making\\
   Them do \emph{all} line breaking, \\
   It allows them to put too many words on a line when they'd rather be
   forced to be terse.
\end{verse}

Mathematical formulas may also be displayed.  A
displayed formula
is
one-line long; multiline
formulas require special formatting instructions.
   \[  \ip{\Gamma}{\psi'} = x'' + y^{2} + z_{i}^{n}\]
Don't start a paragraph with a displayed equation,
nor make one a paragraph by itself.

\end{document}               % End of document.

\section{Bench Huge Section 547}
% synthetic bench section 547
% This is a sample LaTeX input file.  (Version of 12 August 2004.)
%
% A '%' character causes TeX to ignore all remaining text on the line,
% and is used for comments like this one.

\documentclass{article}      % Specifies the document class

                             % The preamble begins here.
\title{An Example Document}  % Declares the document's title.
\author{Leslie Lamport}      % Declares the author's name.
\date{January 21, 1994}      % Deleting this command produces today's date.

\newcommand{\ip}[2]{(#1, #2)}
                             % Defines \ip{arg1}{arg2} to mean
                             % (arg1, arg2).

%\newcommand{\ip}[2]{\langle #1 | #2\rangle}
                             % This is an alternative definition of
                             % \ip that is commented out.

\begin{document}             % End of preamble and beginning of text.

\maketitle                   % Produces the title.

This is an example input file.  Comparing it with
the output it generates can show you how to
produce a simple document of your own.

\section{Ordinary Text}      % Produces section heading.  Lower-level
                             % sections are begun with similar
                             % \subsection and \subsubsection commands.

The ends  of words and sentences are marked
  by   spaces. It  doesn't matter how many
spaces    you type; one is as good as 100.  The
end of   a line counts as a space.

One   or more   blank lines denote the  end
of  a paragraph.

Since any number of consecutive spaces are treated
like a single one, the formatting of the input
file makes no difference to
      \LaTeX,                % The \LaTeX command generates the LaTeX logo.
but it makes a difference to you.  When you use
\LaTeX, making your input file as easy to read
as possible will be a great help as you write
your document and when you change it.  This sample
file shows how you can add comments to your own input
file.

Because printing is different from typewriting,
there are a number of things that you have to do
differently when preparing an input file than if
you were just typing the document directly.
Quotation marks like
       ``this''
have to be handled specially, as do quotes within
quotes:
       ``\,`this'            % \, separates the double and single quote.
        is what I just
        wrote, not  `that'\,''.

Dashes come in three sizes: an
       intra-word
dash, a medium dash for number ranges like
       1--2,
and a punctuation
       dash---like
this.

A sentence-ending space should be larger than the
space between words within a sentence.  You
sometimes have to type special commands in
conjunction with punctuation characters to get
this right, as in the following sentence.
       Gnats, gnus, etc.\ all  % `\ ' makes an inter-word space.
       begin with G\@.         % \@ marks end-of-sentence punctuation.
You should check the spaces after periods when
reading your output to make sure you haven't
forgotten any special cases.  Generating an
ellipsis
       \ldots\               % `\ ' is needed after `\ldots' because TeX
                             % ignores spaces after command names like \ldots
                             % made from \ + letters.
                             %
                             % Note how a `%' character causes TeX to ignore
                             % the end of the input line, so these blank lines
                             % do not start a new paragraph.
                             %
with the right spacing around the periods requires
a special command.

\LaTeX\ interprets some common characters as
commands, so you must type special commands to
generate them.  These characters include the
following:
       \$ \& \% \# \{ and \}.

In printing, text is usually emphasized with an
       \emph{italic}
type style.

\begin{em}
   A long segment of text can also be emphasized
   in this way.  Text within such a segment can be
   given \emph{additional} emphasis.
\end{em}

It is sometimes necessary to prevent \LaTeX\ from
breaking a line where it might otherwise do so.
This may be at a space, as between the ``Mr.''\ and
``Jones'' in
       ``Mr.~Jones'',        % ~ produces an unbreakable interword space.
or within a word---especially when the word is a
symbol like
       \mbox{\emph{itemnum}}
that makes little sense when hyphenated across
lines.

Footnotes\footnote{This is an example of a footnote.}
pose no problem.

\LaTeX\ is good at typesetting mathematical formulas
like
       \( x-3y + z = 7 \)
or
       \( a_{1} > x^{2n} + y^{2n} > x' \)
or
       \( \ip{A}{B} = \sum_{i} a_{i} b_{i} \).
The spaces you type in a formula are
ignored.  Remember that a letter like
       $x$                   % $ ... $  and  \( ... \)  are equivalent
is a formula when it denotes a mathematical
symbol, and it should be typed as one.

\section{Displayed Text}

Text is displayed by indenting it from the left
margin.  Quotations are commonly displayed.  There
are short quotations
\begin{quote}
   This is a short quotation.  It consists of a
   single paragraph of text.  See how it is formatted.
\end{quote}
and longer ones.
\begin{quotation}
   This is a longer quotation.  It consists of two
   paragraphs of text, neither of which are
   particularly interesting.

   This is the second paragraph of the quotation.  It
   is just as dull as the first paragraph.
\end{quotation}
Another frequently-displayed structure is a list.
The following is an example of an \emph{itemized}
list.
\begin{itemize}
   \item This is the first item of an itemized list.
         Each item in the list is marked with a ``tick''.
         You don't have to worry about what kind of tick
         mark is used.

   \item This is the second item of the list.  It
         contains another list nested inside it.  The inner
         list is an \emph{enumerated} list.
         \begin{enumerate}
            \item This is the first item of an enumerated
                  list that is nested within the itemized list.

            \item This is the second item of the inner list.
                  \LaTeX\ allows you to nest lists deeper than
                  you really should.
         \end{enumerate}
         This is the rest of the second item of the outer
         list.  It is no more interesting than any other
         part of the item.
   \item This is the third item of the list.
\end{itemize}
You can even display poetry.
\begin{verse}
   There is an environment
    for verse \\             % The \\ command separates lines
   Whose features some poets % within a stanza.
   will curse.

                             % One or more blank lines separate stanzas.

   For instead of making\\
   Them do \emph{all} line breaking, \\
   It allows them to put too many words on a line when they'd rather be
   forced to be terse.
\end{verse}

Mathematical formulas may also be displayed.  A
displayed formula
is
one-line long; multiline
formulas require special formatting instructions.
   \[  \ip{\Gamma}{\psi'} = x'' + y^{2} + z_{i}^{n}\]
Don't start a paragraph with a displayed equation,
nor make one a paragraph by itself.

\end{document}               % End of document.

\section{Bench Huge Section 548}
% synthetic bench section 548
% This is a sample LaTeX input file.  (Version of 12 August 2004.)
%
% A '%' character causes TeX to ignore all remaining text on the line,
% and is used for comments like this one.

\documentclass{article}      % Specifies the document class

                             % The preamble begins here.
\title{An Example Document}  % Declares the document's title.
\author{Leslie Lamport}      % Declares the author's name.
\date{January 21, 1994}      % Deleting this command produces today's date.

\newcommand{\ip}[2]{(#1, #2)}
                             % Defines \ip{arg1}{arg2} to mean
                             % (arg1, arg2).

%\newcommand{\ip}[2]{\langle #1 | #2\rangle}
                             % This is an alternative definition of
                             % \ip that is commented out.

\begin{document}             % End of preamble and beginning of text.

\maketitle                   % Produces the title.

This is an example input file.  Comparing it with
the output it generates can show you how to
produce a simple document of your own.

\section{Ordinary Text}      % Produces section heading.  Lower-level
                             % sections are begun with similar
                             % \subsection and \subsubsection commands.

The ends  of words and sentences are marked
  by   spaces. It  doesn't matter how many
spaces    you type; one is as good as 100.  The
end of   a line counts as a space.

One   or more   blank lines denote the  end
of  a paragraph.

Since any number of consecutive spaces are treated
like a single one, the formatting of the input
file makes no difference to
      \LaTeX,                % The \LaTeX command generates the LaTeX logo.
but it makes a difference to you.  When you use
\LaTeX, making your input file as easy to read
as possible will be a great help as you write
your document and when you change it.  This sample
file shows how you can add comments to your own input
file.

Because printing is different from typewriting,
there are a number of things that you have to do
differently when preparing an input file than if
you were just typing the document directly.
Quotation marks like
       ``this''
have to be handled specially, as do quotes within
quotes:
       ``\,`this'            % \, separates the double and single quote.
        is what I just
        wrote, not  `that'\,''.

Dashes come in three sizes: an
       intra-word
dash, a medium dash for number ranges like
       1--2,
and a punctuation
       dash---like
this.

A sentence-ending space should be larger than the
space between words within a sentence.  You
sometimes have to type special commands in
conjunction with punctuation characters to get
this right, as in the following sentence.
       Gnats, gnus, etc.\ all  % `\ ' makes an inter-word space.
       begin with G\@.         % \@ marks end-of-sentence punctuation.
You should check the spaces after periods when
reading your output to make sure you haven't
forgotten any special cases.  Generating an
ellipsis
       \ldots\               % `\ ' is needed after `\ldots' because TeX
                             % ignores spaces after command names like \ldots
                             % made from \ + letters.
                             %
                             % Note how a `%' character causes TeX to ignore
                             % the end of the input line, so these blank lines
                             % do not start a new paragraph.
                             %
with the right spacing around the periods requires
a special command.

\LaTeX\ interprets some common characters as
commands, so you must type special commands to
generate them.  These characters include the
following:
       \$ \& \% \# \{ and \}.

In printing, text is usually emphasized with an
       \emph{italic}
type style.

\begin{em}
   A long segment of text can also be emphasized
   in this way.  Text within such a segment can be
   given \emph{additional} emphasis.
\end{em}

It is sometimes necessary to prevent \LaTeX\ from
breaking a line where it might otherwise do so.
This may be at a space, as between the ``Mr.''\ and
``Jones'' in
       ``Mr.~Jones'',        % ~ produces an unbreakable interword space.
or within a word---especially when the word is a
symbol like
       \mbox{\emph{itemnum}}
that makes little sense when hyphenated across
lines.

Footnotes\footnote{This is an example of a footnote.}
pose no problem.

\LaTeX\ is good at typesetting mathematical formulas
like
       \( x-3y + z = 7 \)
or
       \( a_{1} > x^{2n} + y^{2n} > x' \)
or
       \( \ip{A}{B} = \sum_{i} a_{i} b_{i} \).
The spaces you type in a formula are
ignored.  Remember that a letter like
       $x$                   % $ ... $  and  \( ... \)  are equivalent
is a formula when it denotes a mathematical
symbol, and it should be typed as one.

\section{Displayed Text}

Text is displayed by indenting it from the left
margin.  Quotations are commonly displayed.  There
are short quotations
\begin{quote}
   This is a short quotation.  It consists of a
   single paragraph of text.  See how it is formatted.
\end{quote}
and longer ones.
\begin{quotation}
   This is a longer quotation.  It consists of two
   paragraphs of text, neither of which are
   particularly interesting.

   This is the second paragraph of the quotation.  It
   is just as dull as the first paragraph.
\end{quotation}
Another frequently-displayed structure is a list.
The following is an example of an \emph{itemized}
list.
\begin{itemize}
   \item This is the first item of an itemized list.
         Each item in the list is marked with a ``tick''.
         You don't have to worry about what kind of tick
         mark is used.

   \item This is the second item of the list.  It
         contains another list nested inside it.  The inner
         list is an \emph{enumerated} list.
         \begin{enumerate}
            \item This is the first item of an enumerated
                  list that is nested within the itemized list.

            \item This is the second item of the inner list.
                  \LaTeX\ allows you to nest lists deeper than
                  you really should.
         \end{enumerate}
         This is the rest of the second item of the outer
         list.  It is no more interesting than any other
         part of the item.
   \item This is the third item of the list.
\end{itemize}
You can even display poetry.
\begin{verse}
   There is an environment
    for verse \\             % The \\ command separates lines
   Whose features some poets % within a stanza.
   will curse.

                             % One or more blank lines separate stanzas.

   For instead of making\\
   Them do \emph{all} line breaking, \\
   It allows them to put too many words on a line when they'd rather be
   forced to be terse.
\end{verse}

Mathematical formulas may also be displayed.  A
displayed formula
is
one-line long; multiline
formulas require special formatting instructions.
   \[  \ip{\Gamma}{\psi'} = x'' + y^{2} + z_{i}^{n}\]
Don't start a paragraph with a displayed equation,
nor make one a paragraph by itself.

\end{document}               % End of document.

\section{Bench Huge Section 549}
% synthetic bench section 549
% This is a sample LaTeX input file.  (Version of 12 August 2004.)
%
% A '%' character causes TeX to ignore all remaining text on the line,
% and is used for comments like this one.

\documentclass{article}      % Specifies the document class

                             % The preamble begins here.
\title{An Example Document}  % Declares the document's title.
\author{Leslie Lamport}      % Declares the author's name.
\date{January 21, 1994}      % Deleting this command produces today's date.

\newcommand{\ip}[2]{(#1, #2)}
                             % Defines \ip{arg1}{arg2} to mean
                             % (arg1, arg2).

%\newcommand{\ip}[2]{\langle #1 | #2\rangle}
                             % This is an alternative definition of
                             % \ip that is commented out.

\begin{document}             % End of preamble and beginning of text.

\maketitle                   % Produces the title.

This is an example input file.  Comparing it with
the output it generates can show you how to
produce a simple document of your own.

\section{Ordinary Text}      % Produces section heading.  Lower-level
                             % sections are begun with similar
                             % \subsection and \subsubsection commands.

The ends  of words and sentences are marked
  by   spaces. It  doesn't matter how many
spaces    you type; one is as good as 100.  The
end of   a line counts as a space.

One   or more   blank lines denote the  end
of  a paragraph.

Since any number of consecutive spaces are treated
like a single one, the formatting of the input
file makes no difference to
      \LaTeX,                % The \LaTeX command generates the LaTeX logo.
but it makes a difference to you.  When you use
\LaTeX, making your input file as easy to read
as possible will be a great help as you write
your document and when you change it.  This sample
file shows how you can add comments to your own input
file.

Because printing is different from typewriting,
there are a number of things that you have to do
differently when preparing an input file than if
you were just typing the document directly.
Quotation marks like
       ``this''
have to be handled specially, as do quotes within
quotes:
       ``\,`this'            % \, separates the double and single quote.
        is what I just
        wrote, not  `that'\,''.

Dashes come in three sizes: an
       intra-word
dash, a medium dash for number ranges like
       1--2,
and a punctuation
       dash---like
this.

A sentence-ending space should be larger than the
space between words within a sentence.  You
sometimes have to type special commands in
conjunction with punctuation characters to get
this right, as in the following sentence.
       Gnats, gnus, etc.\ all  % `\ ' makes an inter-word space.
       begin with G\@.         % \@ marks end-of-sentence punctuation.
You should check the spaces after periods when
reading your output to make sure you haven't
forgotten any special cases.  Generating an
ellipsis
       \ldots\               % `\ ' is needed after `\ldots' because TeX
                             % ignores spaces after command names like \ldots
                             % made from \ + letters.
                             %
                             % Note how a `%' character causes TeX to ignore
                             % the end of the input line, so these blank lines
                             % do not start a new paragraph.
                             %
with the right spacing around the periods requires
a special command.

\LaTeX\ interprets some common characters as
commands, so you must type special commands to
generate them.  These characters include the
following:
       \$ \& \% \# \{ and \}.

In printing, text is usually emphasized with an
       \emph{italic}
type style.

\begin{em}
   A long segment of text can also be emphasized
   in this way.  Text within such a segment can be
   given \emph{additional} emphasis.
\end{em}

It is sometimes necessary to prevent \LaTeX\ from
breaking a line where it might otherwise do so.
This may be at a space, as between the ``Mr.''\ and
``Jones'' in
       ``Mr.~Jones'',        % ~ produces an unbreakable interword space.
or within a word---especially when the word is a
symbol like
       \mbox{\emph{itemnum}}
that makes little sense when hyphenated across
lines.

Footnotes\footnote{This is an example of a footnote.}
pose no problem.

\LaTeX\ is good at typesetting mathematical formulas
like
       \( x-3y + z = 7 \)
or
       \( a_{1} > x^{2n} + y^{2n} > x' \)
or
       \( \ip{A}{B} = \sum_{i} a_{i} b_{i} \).
The spaces you type in a formula are
ignored.  Remember that a letter like
       $x$                   % $ ... $  and  \( ... \)  are equivalent
is a formula when it denotes a mathematical
symbol, and it should be typed as one.

\section{Displayed Text}

Text is displayed by indenting it from the left
margin.  Quotations are commonly displayed.  There
are short quotations
\begin{quote}
   This is a short quotation.  It consists of a
   single paragraph of text.  See how it is formatted.
\end{quote}
and longer ones.
\begin{quotation}
   This is a longer quotation.  It consists of two
   paragraphs of text, neither of which are
   particularly interesting.

   This is the second paragraph of the quotation.  It
   is just as dull as the first paragraph.
\end{quotation}
Another frequently-displayed structure is a list.
The following is an example of an \emph{itemized}
list.
\begin{itemize}
   \item This is the first item of an itemized list.
         Each item in the list is marked with a ``tick''.
         You don't have to worry about what kind of tick
         mark is used.

   \item This is the second item of the list.  It
         contains another list nested inside it.  The inner
         list is an \emph{enumerated} list.
         \begin{enumerate}
            \item This is the first item of an enumerated
                  list that is nested within the itemized list.

            \item This is the second item of the inner list.
                  \LaTeX\ allows you to nest lists deeper than
                  you really should.
         \end{enumerate}
         This is the rest of the second item of the outer
         list.  It is no more interesting than any other
         part of the item.
   \item This is the third item of the list.
\end{itemize}
You can even display poetry.
\begin{verse}
   There is an environment
    for verse \\             % The \\ command separates lines
   Whose features some poets % within a stanza.
   will curse.

                             % One or more blank lines separate stanzas.

   For instead of making\\
   Them do \emph{all} line breaking, \\
   It allows them to put too many words on a line when they'd rather be
   forced to be terse.
\end{verse}

Mathematical formulas may also be displayed.  A
displayed formula
is
one-line long; multiline
formulas require special formatting instructions.
   \[  \ip{\Gamma}{\psi'} = x'' + y^{2} + z_{i}^{n}\]
Don't start a paragraph with a displayed equation,
nor make one a paragraph by itself.

\end{document}               % End of document.

\section{Bench Huge Section 550}
% synthetic bench section 550
% This is a sample LaTeX input file.  (Version of 12 August 2004.)
%
% A '%' character causes TeX to ignore all remaining text on the line,
% and is used for comments like this one.

\documentclass{article}      % Specifies the document class

                             % The preamble begins here.
\title{An Example Document}  % Declares the document's title.
\author{Leslie Lamport}      % Declares the author's name.
\date{January 21, 1994}      % Deleting this command produces today's date.

\newcommand{\ip}[2]{(#1, #2)}
                             % Defines \ip{arg1}{arg2} to mean
                             % (arg1, arg2).

%\newcommand{\ip}[2]{\langle #1 | #2\rangle}
                             % This is an alternative definition of
                             % \ip that is commented out.

\begin{document}             % End of preamble and beginning of text.

\maketitle                   % Produces the title.

This is an example input file.  Comparing it with
the output it generates can show you how to
produce a simple document of your own.

\section{Ordinary Text}      % Produces section heading.  Lower-level
                             % sections are begun with similar
                             % \subsection and \subsubsection commands.

The ends  of words and sentences are marked
  by   spaces. It  doesn't matter how many
spaces    you type; one is as good as 100.  The
end of   a line counts as a space.

One   or more   blank lines denote the  end
of  a paragraph.

Since any number of consecutive spaces are treated
like a single one, the formatting of the input
file makes no difference to
      \LaTeX,                % The \LaTeX command generates the LaTeX logo.
but it makes a difference to you.  When you use
\LaTeX, making your input file as easy to read
as possible will be a great help as you write
your document and when you change it.  This sample
file shows how you can add comments to your own input
file.

Because printing is different from typewriting,
there are a number of things that you have to do
differently when preparing an input file than if
you were just typing the document directly.
Quotation marks like
       ``this''
have to be handled specially, as do quotes within
quotes:
       ``\,`this'            % \, separates the double and single quote.
        is what I just
        wrote, not  `that'\,''.

Dashes come in three sizes: an
       intra-word
dash, a medium dash for number ranges like
       1--2,
and a punctuation
       dash---like
this.

A sentence-ending space should be larger than the
space between words within a sentence.  You
sometimes have to type special commands in
conjunction with punctuation characters to get
this right, as in the following sentence.
       Gnats, gnus, etc.\ all  % `\ ' makes an inter-word space.
       begin with G\@.         % \@ marks end-of-sentence punctuation.
You should check the spaces after periods when
reading your output to make sure you haven't
forgotten any special cases.  Generating an
ellipsis
       \ldots\               % `\ ' is needed after `\ldots' because TeX
                             % ignores spaces after command names like \ldots
                             % made from \ + letters.
                             %
                             % Note how a `%' character causes TeX to ignore
                             % the end of the input line, so these blank lines
                             % do not start a new paragraph.
                             %
with the right spacing around the periods requires
a special command.

\LaTeX\ interprets some common characters as
commands, so you must type special commands to
generate them.  These characters include the
following:
       \$ \& \% \# \{ and \}.

In printing, text is usually emphasized with an
       \emph{italic}
type style.

\begin{em}
   A long segment of text can also be emphasized
   in this way.  Text within such a segment can be
   given \emph{additional} emphasis.
\end{em}

It is sometimes necessary to prevent \LaTeX\ from
breaking a line where it might otherwise do so.
This may be at a space, as between the ``Mr.''\ and
``Jones'' in
       ``Mr.~Jones'',        % ~ produces an unbreakable interword space.
or within a word---especially when the word is a
symbol like
       \mbox{\emph{itemnum}}
that makes little sense when hyphenated across
lines.

Footnotes\footnote{This is an example of a footnote.}
pose no problem.

\LaTeX\ is good at typesetting mathematical formulas
like
       \( x-3y + z = 7 \)
or
       \( a_{1} > x^{2n} + y^{2n} > x' \)
or
       \( \ip{A}{B} = \sum_{i} a_{i} b_{i} \).
The spaces you type in a formula are
ignored.  Remember that a letter like
       $x$                   % $ ... $  and  \( ... \)  are equivalent
is a formula when it denotes a mathematical
symbol, and it should be typed as one.

\section{Displayed Text}

Text is displayed by indenting it from the left
margin.  Quotations are commonly displayed.  There
are short quotations
\begin{quote}
   This is a short quotation.  It consists of a
   single paragraph of text.  See how it is formatted.
\end{quote}
and longer ones.
\begin{quotation}
   This is a longer quotation.  It consists of two
   paragraphs of text, neither of which are
   particularly interesting.

   This is the second paragraph of the quotation.  It
   is just as dull as the first paragraph.
\end{quotation}
Another frequently-displayed structure is a list.
The following is an example of an \emph{itemized}
list.
\begin{itemize}
   \item This is the first item of an itemized list.
         Each item in the list is marked with a ``tick''.
         You don't have to worry about what kind of tick
         mark is used.

   \item This is the second item of the list.  It
         contains another list nested inside it.  The inner
         list is an \emph{enumerated} list.
         \begin{enumerate}
            \item This is the first item of an enumerated
                  list that is nested within the itemized list.

            \item This is the second item of the inner list.
                  \LaTeX\ allows you to nest lists deeper than
                  you really should.
         \end{enumerate}
         This is the rest of the second item of the outer
         list.  It is no more interesting than any other
         part of the item.
   \item This is the third item of the list.
\end{itemize}
You can even display poetry.
\begin{verse}
   There is an environment
    for verse \\             % The \\ command separates lines
   Whose features some poets % within a stanza.
   will curse.

                             % One or more blank lines separate stanzas.

   For instead of making\\
   Them do \emph{all} line breaking, \\
   It allows them to put too many words on a line when they'd rather be
   forced to be terse.
\end{verse}

Mathematical formulas may also be displayed.  A
displayed formula
is
one-line long; multiline
formulas require special formatting instructions.
   \[  \ip{\Gamma}{\psi'} = x'' + y^{2} + z_{i}^{n}\]
Don't start a paragraph with a displayed equation,
nor make one a paragraph by itself.

\end{document}               % End of document.

\section{Bench Huge Section 551}
% synthetic bench section 551
% This is a sample LaTeX input file.  (Version of 12 August 2004.)
%
% A '%' character causes TeX to ignore all remaining text on the line,
% and is used for comments like this one.

\documentclass{article}      % Specifies the document class

                             % The preamble begins here.
\title{An Example Document}  % Declares the document's title.
\author{Leslie Lamport}      % Declares the author's name.
\date{January 21, 1994}      % Deleting this command produces today's date.

\newcommand{\ip}[2]{(#1, #2)}
                             % Defines \ip{arg1}{arg2} to mean
                             % (arg1, arg2).

%\newcommand{\ip}[2]{\langle #1 | #2\rangle}
                             % This is an alternative definition of
                             % \ip that is commented out.

\begin{document}             % End of preamble and beginning of text.

\maketitle                   % Produces the title.

This is an example input file.  Comparing it with
the output it generates can show you how to
produce a simple document of your own.

\section{Ordinary Text}      % Produces section heading.  Lower-level
                             % sections are begun with similar
                             % \subsection and \subsubsection commands.

The ends  of words and sentences are marked
  by   spaces. It  doesn't matter how many
spaces    you type; one is as good as 100.  The
end of   a line counts as a space.

One   or more   blank lines denote the  end
of  a paragraph.

Since any number of consecutive spaces are treated
like a single one, the formatting of the input
file makes no difference to
      \LaTeX,                % The \LaTeX command generates the LaTeX logo.
but it makes a difference to you.  When you use
\LaTeX, making your input file as easy to read
as possible will be a great help as you write
your document and when you change it.  This sample
file shows how you can add comments to your own input
file.

Because printing is different from typewriting,
there are a number of things that you have to do
differently when preparing an input file than if
you were just typing the document directly.
Quotation marks like
       ``this''
have to be handled specially, as do quotes within
quotes:
       ``\,`this'            % \, separates the double and single quote.
        is what I just
        wrote, not  `that'\,''.

Dashes come in three sizes: an
       intra-word
dash, a medium dash for number ranges like
       1--2,
and a punctuation
       dash---like
this.

A sentence-ending space should be larger than the
space between words within a sentence.  You
sometimes have to type special commands in
conjunction with punctuation characters to get
this right, as in the following sentence.
       Gnats, gnus, etc.\ all  % `\ ' makes an inter-word space.
       begin with G\@.         % \@ marks end-of-sentence punctuation.
You should check the spaces after periods when
reading your output to make sure you haven't
forgotten any special cases.  Generating an
ellipsis
       \ldots\               % `\ ' is needed after `\ldots' because TeX
                             % ignores spaces after command names like \ldots
                             % made from \ + letters.
                             %
                             % Note how a `%' character causes TeX to ignore
                             % the end of the input line, so these blank lines
                             % do not start a new paragraph.
                             %
with the right spacing around the periods requires
a special command.

\LaTeX\ interprets some common characters as
commands, so you must type special commands to
generate them.  These characters include the
following:
       \$ \& \% \# \{ and \}.

In printing, text is usually emphasized with an
       \emph{italic}
type style.

\begin{em}
   A long segment of text can also be emphasized
   in this way.  Text within such a segment can be
   given \emph{additional} emphasis.
\end{em}

It is sometimes necessary to prevent \LaTeX\ from
breaking a line where it might otherwise do so.
This may be at a space, as between the ``Mr.''\ and
``Jones'' in
       ``Mr.~Jones'',        % ~ produces an unbreakable interword space.
or within a word---especially when the word is a
symbol like
       \mbox{\emph{itemnum}}
that makes little sense when hyphenated across
lines.

Footnotes\footnote{This is an example of a footnote.}
pose no problem.

\LaTeX\ is good at typesetting mathematical formulas
like
       \( x-3y + z = 7 \)
or
       \( a_{1} > x^{2n} + y^{2n} > x' \)
or
       \( \ip{A}{B} = \sum_{i} a_{i} b_{i} \).
The spaces you type in a formula are
ignored.  Remember that a letter like
       $x$                   % $ ... $  and  \( ... \)  are equivalent
is a formula when it denotes a mathematical
symbol, and it should be typed as one.

\section{Displayed Text}

Text is displayed by indenting it from the left
margin.  Quotations are commonly displayed.  There
are short quotations
\begin{quote}
   This is a short quotation.  It consists of a
   single paragraph of text.  See how it is formatted.
\end{quote}
and longer ones.
\begin{quotation}
   This is a longer quotation.  It consists of two
   paragraphs of text, neither of which are
   particularly interesting.

   This is the second paragraph of the quotation.  It
   is just as dull as the first paragraph.
\end{quotation}
Another frequently-displayed structure is a list.
The following is an example of an \emph{itemized}
list.
\begin{itemize}
   \item This is the first item of an itemized list.
         Each item in the list is marked with a ``tick''.
         You don't have to worry about what kind of tick
         mark is used.

   \item This is the second item of the list.  It
         contains another list nested inside it.  The inner
         list is an \emph{enumerated} list.
         \begin{enumerate}
            \item This is the first item of an enumerated
                  list that is nested within the itemized list.

            \item This is the second item of the inner list.
                  \LaTeX\ allows you to nest lists deeper than
                  you really should.
         \end{enumerate}
         This is the rest of the second item of the outer
         list.  It is no more interesting than any other
         part of the item.
   \item This is the third item of the list.
\end{itemize}
You can even display poetry.
\begin{verse}
   There is an environment
    for verse \\             % The \\ command separates lines
   Whose features some poets % within a stanza.
   will curse.

                             % One or more blank lines separate stanzas.

   For instead of making\\
   Them do \emph{all} line breaking, \\
   It allows them to put too many words on a line when they'd rather be
   forced to be terse.
\end{verse}

Mathematical formulas may also be displayed.  A
displayed formula
is
one-line long; multiline
formulas require special formatting instructions.
   \[  \ip{\Gamma}{\psi'} = x'' + y^{2} + z_{i}^{n}\]
Don't start a paragraph with a displayed equation,
nor make one a paragraph by itself.

\end{document}               % End of document.

\section{Bench Huge Section 552}
% synthetic bench section 552
% This is a sample LaTeX input file.  (Version of 12 August 2004.)
%
% A '%' character causes TeX to ignore all remaining text on the line,
% and is used for comments like this one.

\documentclass{article}      % Specifies the document class

                             % The preamble begins here.
\title{An Example Document}  % Declares the document's title.
\author{Leslie Lamport}      % Declares the author's name.
\date{January 21, 1994}      % Deleting this command produces today's date.

\newcommand{\ip}[2]{(#1, #2)}
                             % Defines \ip{arg1}{arg2} to mean
                             % (arg1, arg2).

%\newcommand{\ip}[2]{\langle #1 | #2\rangle}
                             % This is an alternative definition of
                             % \ip that is commented out.

\begin{document}             % End of preamble and beginning of text.

\maketitle                   % Produces the title.

This is an example input file.  Comparing it with
the output it generates can show you how to
produce a simple document of your own.

\section{Ordinary Text}      % Produces section heading.  Lower-level
                             % sections are begun with similar
                             % \subsection and \subsubsection commands.

The ends  of words and sentences are marked
  by   spaces. It  doesn't matter how many
spaces    you type; one is as good as 100.  The
end of   a line counts as a space.

One   or more   blank lines denote the  end
of  a paragraph.

Since any number of consecutive spaces are treated
like a single one, the formatting of the input
file makes no difference to
      \LaTeX,                % The \LaTeX command generates the LaTeX logo.
but it makes a difference to you.  When you use
\LaTeX, making your input file as easy to read
as possible will be a great help as you write
your document and when you change it.  This sample
file shows how you can add comments to your own input
file.

Because printing is different from typewriting,
there are a number of things that you have to do
differently when preparing an input file than if
you were just typing the document directly.
Quotation marks like
       ``this''
have to be handled specially, as do quotes within
quotes:
       ``\,`this'            % \, separates the double and single quote.
        is what I just
        wrote, not  `that'\,''.

Dashes come in three sizes: an
       intra-word
dash, a medium dash for number ranges like
       1--2,
and a punctuation
       dash---like
this.

A sentence-ending space should be larger than the
space between words within a sentence.  You
sometimes have to type special commands in
conjunction with punctuation characters to get
this right, as in the following sentence.
       Gnats, gnus, etc.\ all  % `\ ' makes an inter-word space.
       begin with G\@.         % \@ marks end-of-sentence punctuation.
You should check the spaces after periods when
reading your output to make sure you haven't
forgotten any special cases.  Generating an
ellipsis
       \ldots\               % `\ ' is needed after `\ldots' because TeX
                             % ignores spaces after command names like \ldots
                             % made from \ + letters.
                             %
                             % Note how a `%' character causes TeX to ignore
                             % the end of the input line, so these blank lines
                             % do not start a new paragraph.
                             %
with the right spacing around the periods requires
a special command.

\LaTeX\ interprets some common characters as
commands, so you must type special commands to
generate them.  These characters include the
following:
       \$ \& \% \# \{ and \}.

In printing, text is usually emphasized with an
       \emph{italic}
type style.

\begin{em}
   A long segment of text can also be emphasized
   in this way.  Text within such a segment can be
   given \emph{additional} emphasis.
\end{em}

It is sometimes necessary to prevent \LaTeX\ from
breaking a line where it might otherwise do so.
This may be at a space, as between the ``Mr.''\ and
``Jones'' in
       ``Mr.~Jones'',        % ~ produces an unbreakable interword space.
or within a word---especially when the word is a
symbol like
       \mbox{\emph{itemnum}}
that makes little sense when hyphenated across
lines.

Footnotes\footnote{This is an example of a footnote.}
pose no problem.

\LaTeX\ is good at typesetting mathematical formulas
like
       \( x-3y + z = 7 \)
or
       \( a_{1} > x^{2n} + y^{2n} > x' \)
or
       \( \ip{A}{B} = \sum_{i} a_{i} b_{i} \).
The spaces you type in a formula are
ignored.  Remember that a letter like
       $x$                   % $ ... $  and  \( ... \)  are equivalent
is a formula when it denotes a mathematical
symbol, and it should be typed as one.

\section{Displayed Text}

Text is displayed by indenting it from the left
margin.  Quotations are commonly displayed.  There
are short quotations
\begin{quote}
   This is a short quotation.  It consists of a
   single paragraph of text.  See how it is formatted.
\end{quote}
and longer ones.
\begin{quotation}
   This is a longer quotation.  It consists of two
   paragraphs of text, neither of which are
   particularly interesting.

   This is the second paragraph of the quotation.  It
   is just as dull as the first paragraph.
\end{quotation}
Another frequently-displayed structure is a list.
The following is an example of an \emph{itemized}
list.
\begin{itemize}
   \item This is the first item of an itemized list.
         Each item in the list is marked with a ``tick''.
         You don't have to worry about what kind of tick
         mark is used.

   \item This is the second item of the list.  It
         contains another list nested inside it.  The inner
         list is an \emph{enumerated} list.
         \begin{enumerate}
            \item This is the first item of an enumerated
                  list that is nested within the itemized list.

            \item This is the second item of the inner list.
                  \LaTeX\ allows you to nest lists deeper than
                  you really should.
         \end{enumerate}
         This is the rest of the second item of the outer
         list.  It is no more interesting than any other
         part of the item.
   \item This is the third item of the list.
\end{itemize}
You can even display poetry.
\begin{verse}
   There is an environment
    for verse \\             % The \\ command separates lines
   Whose features some poets % within a stanza.
   will curse.

                             % One or more blank lines separate stanzas.

   For instead of making\\
   Them do \emph{all} line breaking, \\
   It allows them to put too many words on a line when they'd rather be
   forced to be terse.
\end{verse}

Mathematical formulas may also be displayed.  A
displayed formula
is
one-line long; multiline
formulas require special formatting instructions.
   \[  \ip{\Gamma}{\psi'} = x'' + y^{2} + z_{i}^{n}\]
Don't start a paragraph with a displayed equation,
nor make one a paragraph by itself.

\end{document}               % End of document.

\section{Bench Huge Section 553}
% synthetic bench section 553
% This is a sample LaTeX input file.  (Version of 12 August 2004.)
%
% A '%' character causes TeX to ignore all remaining text on the line,
% and is used for comments like this one.

\documentclass{article}      % Specifies the document class

                             % The preamble begins here.
\title{An Example Document}  % Declares the document's title.
\author{Leslie Lamport}      % Declares the author's name.
\date{January 21, 1994}      % Deleting this command produces today's date.

\newcommand{\ip}[2]{(#1, #2)}
                             % Defines \ip{arg1}{arg2} to mean
                             % (arg1, arg2).

%\newcommand{\ip}[2]{\langle #1 | #2\rangle}
                             % This is an alternative definition of
                             % \ip that is commented out.

\begin{document}             % End of preamble and beginning of text.

\maketitle                   % Produces the title.

This is an example input file.  Comparing it with
the output it generates can show you how to
produce a simple document of your own.

\section{Ordinary Text}      % Produces section heading.  Lower-level
                             % sections are begun with similar
                             % \subsection and \subsubsection commands.

The ends  of words and sentences are marked
  by   spaces. It  doesn't matter how many
spaces    you type; one is as good as 100.  The
end of   a line counts as a space.

One   or more   blank lines denote the  end
of  a paragraph.

Since any number of consecutive spaces are treated
like a single one, the formatting of the input
file makes no difference to
      \LaTeX,                % The \LaTeX command generates the LaTeX logo.
but it makes a difference to you.  When you use
\LaTeX, making your input file as easy to read
as possible will be a great help as you write
your document and when you change it.  This sample
file shows how you can add comments to your own input
file.

Because printing is different from typewriting,
there are a number of things that you have to do
differently when preparing an input file than if
you were just typing the document directly.
Quotation marks like
       ``this''
have to be handled specially, as do quotes within
quotes:
       ``\,`this'            % \, separates the double and single quote.
        is what I just
        wrote, not  `that'\,''.

Dashes come in three sizes: an
       intra-word
dash, a medium dash for number ranges like
       1--2,
and a punctuation
       dash---like
this.

A sentence-ending space should be larger than the
space between words within a sentence.  You
sometimes have to type special commands in
conjunction with punctuation characters to get
this right, as in the following sentence.
       Gnats, gnus, etc.\ all  % `\ ' makes an inter-word space.
       begin with G\@.         % \@ marks end-of-sentence punctuation.
You should check the spaces after periods when
reading your output to make sure you haven't
forgotten any special cases.  Generating an
ellipsis
       \ldots\               % `\ ' is needed after `\ldots' because TeX
                             % ignores spaces after command names like \ldots
                             % made from \ + letters.
                             %
                             % Note how a `%' character causes TeX to ignore
                             % the end of the input line, so these blank lines
                             % do not start a new paragraph.
                             %
with the right spacing around the periods requires
a special command.

\LaTeX\ interprets some common characters as
commands, so you must type special commands to
generate them.  These characters include the
following:
       \$ \& \% \# \{ and \}.

In printing, text is usually emphasized with an
       \emph{italic}
type style.

\begin{em}
   A long segment of text can also be emphasized
   in this way.  Text within such a segment can be
   given \emph{additional} emphasis.
\end{em}

It is sometimes necessary to prevent \LaTeX\ from
breaking a line where it might otherwise do so.
This may be at a space, as between the ``Mr.''\ and
``Jones'' in
       ``Mr.~Jones'',        % ~ produces an unbreakable interword space.
or within a word---especially when the word is a
symbol like
       \mbox{\emph{itemnum}}
that makes little sense when hyphenated across
lines.

Footnotes\footnote{This is an example of a footnote.}
pose no problem.

\LaTeX\ is good at typesetting mathematical formulas
like
       \( x-3y + z = 7 \)
or
       \( a_{1} > x^{2n} + y^{2n} > x' \)
or
       \( \ip{A}{B} = \sum_{i} a_{i} b_{i} \).
The spaces you type in a formula are
ignored.  Remember that a letter like
       $x$                   % $ ... $  and  \( ... \)  are equivalent
is a formula when it denotes a mathematical
symbol, and it should be typed as one.

\section{Displayed Text}

Text is displayed by indenting it from the left
margin.  Quotations are commonly displayed.  There
are short quotations
\begin{quote}
   This is a short quotation.  It consists of a
   single paragraph of text.  See how it is formatted.
\end{quote}
and longer ones.
\begin{quotation}
   This is a longer quotation.  It consists of two
   paragraphs of text, neither of which are
   particularly interesting.

   This is the second paragraph of the quotation.  It
   is just as dull as the first paragraph.
\end{quotation}
Another frequently-displayed structure is a list.
The following is an example of an \emph{itemized}
list.
\begin{itemize}
   \item This is the first item of an itemized list.
         Each item in the list is marked with a ``tick''.
         You don't have to worry about what kind of tick
         mark is used.

   \item This is the second item of the list.  It
         contains another list nested inside it.  The inner
         list is an \emph{enumerated} list.
         \begin{enumerate}
            \item This is the first item of an enumerated
                  list that is nested within the itemized list.

            \item This is the second item of the inner list.
                  \LaTeX\ allows you to nest lists deeper than
                  you really should.
         \end{enumerate}
         This is the rest of the second item of the outer
         list.  It is no more interesting than any other
         part of the item.
   \item This is the third item of the list.
\end{itemize}
You can even display poetry.
\begin{verse}
   There is an environment
    for verse \\             % The \\ command separates lines
   Whose features some poets % within a stanza.
   will curse.

                             % One or more blank lines separate stanzas.

   For instead of making\\
   Them do \emph{all} line breaking, \\
   It allows them to put too many words on a line when they'd rather be
   forced to be terse.
\end{verse}

Mathematical formulas may also be displayed.  A
displayed formula
is
one-line long; multiline
formulas require special formatting instructions.
   \[  \ip{\Gamma}{\psi'} = x'' + y^{2} + z_{i}^{n}\]
Don't start a paragraph with a displayed equation,
nor make one a paragraph by itself.

\end{document}               % End of document.

\section{Bench Huge Section 554}
% synthetic bench section 554
% This is a sample LaTeX input file.  (Version of 12 August 2004.)
%
% A '%' character causes TeX to ignore all remaining text on the line,
% and is used for comments like this one.

\documentclass{article}      % Specifies the document class

                             % The preamble begins here.
\title{An Example Document}  % Declares the document's title.
\author{Leslie Lamport}      % Declares the author's name.
\date{January 21, 1994}      % Deleting this command produces today's date.

\newcommand{\ip}[2]{(#1, #2)}
                             % Defines \ip{arg1}{arg2} to mean
                             % (arg1, arg2).

%\newcommand{\ip}[2]{\langle #1 | #2\rangle}
                             % This is an alternative definition of
                             % \ip that is commented out.

\begin{document}             % End of preamble and beginning of text.

\maketitle                   % Produces the title.

This is an example input file.  Comparing it with
the output it generates can show you how to
produce a simple document of your own.

\section{Ordinary Text}      % Produces section heading.  Lower-level
                             % sections are begun with similar
                             % \subsection and \subsubsection commands.

The ends  of words and sentences are marked
  by   spaces. It  doesn't matter how many
spaces    you type; one is as good as 100.  The
end of   a line counts as a space.

One   or more   blank lines denote the  end
of  a paragraph.

Since any number of consecutive spaces are treated
like a single one, the formatting of the input
file makes no difference to
      \LaTeX,                % The \LaTeX command generates the LaTeX logo.
but it makes a difference to you.  When you use
\LaTeX, making your input file as easy to read
as possible will be a great help as you write
your document and when you change it.  This sample
file shows how you can add comments to your own input
file.

Because printing is different from typewriting,
there are a number of things that you have to do
differently when preparing an input file than if
you were just typing the document directly.
Quotation marks like
       ``this''
have to be handled specially, as do quotes within
quotes:
       ``\,`this'            % \, separates the double and single quote.
        is what I just
        wrote, not  `that'\,''.

Dashes come in three sizes: an
       intra-word
dash, a medium dash for number ranges like
       1--2,
and a punctuation
       dash---like
this.

A sentence-ending space should be larger than the
space between words within a sentence.  You
sometimes have to type special commands in
conjunction with punctuation characters to get
this right, as in the following sentence.
       Gnats, gnus, etc.\ all  % `\ ' makes an inter-word space.
       begin with G\@.         % \@ marks end-of-sentence punctuation.
You should check the spaces after periods when
reading your output to make sure you haven't
forgotten any special cases.  Generating an
ellipsis
       \ldots\               % `\ ' is needed after `\ldots' because TeX
                             % ignores spaces after command names like \ldots
                             % made from \ + letters.
                             %
                             % Note how a `%' character causes TeX to ignore
                             % the end of the input line, so these blank lines
                             % do not start a new paragraph.
                             %
with the right spacing around the periods requires
a special command.

\LaTeX\ interprets some common characters as
commands, so you must type special commands to
generate them.  These characters include the
following:
       \$ \& \% \# \{ and \}.

In printing, text is usually emphasized with an
       \emph{italic}
type style.

\begin{em}
   A long segment of text can also be emphasized
   in this way.  Text within such a segment can be
   given \emph{additional} emphasis.
\end{em}

It is sometimes necessary to prevent \LaTeX\ from
breaking a line where it might otherwise do so.
This may be at a space, as between the ``Mr.''\ and
``Jones'' in
       ``Mr.~Jones'',        % ~ produces an unbreakable interword space.
or within a word---especially when the word is a
symbol like
       \mbox{\emph{itemnum}}
that makes little sense when hyphenated across
lines.

Footnotes\footnote{This is an example of a footnote.}
pose no problem.

\LaTeX\ is good at typesetting mathematical formulas
like
       \( x-3y + z = 7 \)
or
       \( a_{1} > x^{2n} + y^{2n} > x' \)
or
       \( \ip{A}{B} = \sum_{i} a_{i} b_{i} \).
The spaces you type in a formula are
ignored.  Remember that a letter like
       $x$                   % $ ... $  and  \( ... \)  are equivalent
is a formula when it denotes a mathematical
symbol, and it should be typed as one.

\section{Displayed Text}

Text is displayed by indenting it from the left
margin.  Quotations are commonly displayed.  There
are short quotations
\begin{quote}
   This is a short quotation.  It consists of a
   single paragraph of text.  See how it is formatted.
\end{quote}
and longer ones.
\begin{quotation}
   This is a longer quotation.  It consists of two
   paragraphs of text, neither of which are
   particularly interesting.

   This is the second paragraph of the quotation.  It
   is just as dull as the first paragraph.
\end{quotation}
Another frequently-displayed structure is a list.
The following is an example of an \emph{itemized}
list.
\begin{itemize}
   \item This is the first item of an itemized list.
         Each item in the list is marked with a ``tick''.
         You don't have to worry about what kind of tick
         mark is used.

   \item This is the second item of the list.  It
         contains another list nested inside it.  The inner
         list is an \emph{enumerated} list.
         \begin{enumerate}
            \item This is the first item of an enumerated
                  list that is nested within the itemized list.

            \item This is the second item of the inner list.
                  \LaTeX\ allows you to nest lists deeper than
                  you really should.
         \end{enumerate}
         This is the rest of the second item of the outer
         list.  It is no more interesting than any other
         part of the item.
   \item This is the third item of the list.
\end{itemize}
You can even display poetry.
\begin{verse}
   There is an environment
    for verse \\             % The \\ command separates lines
   Whose features some poets % within a stanza.
   will curse.

                             % One or more blank lines separate stanzas.

   For instead of making\\
   Them do \emph{all} line breaking, \\
   It allows them to put too many words on a line when they'd rather be
   forced to be terse.
\end{verse}

Mathematical formulas may also be displayed.  A
displayed formula
is
one-line long; multiline
formulas require special formatting instructions.
   \[  \ip{\Gamma}{\psi'} = x'' + y^{2} + z_{i}^{n}\]
Don't start a paragraph with a displayed equation,
nor make one a paragraph by itself.

\end{document}               % End of document.

\section{Bench Huge Section 555}
% synthetic bench section 555
% This is a sample LaTeX input file.  (Version of 12 August 2004.)
%
% A '%' character causes TeX to ignore all remaining text on the line,
% and is used for comments like this one.

\documentclass{article}      % Specifies the document class

                             % The preamble begins here.
\title{An Example Document}  % Declares the document's title.
\author{Leslie Lamport}      % Declares the author's name.
\date{January 21, 1994}      % Deleting this command produces today's date.

\newcommand{\ip}[2]{(#1, #2)}
                             % Defines \ip{arg1}{arg2} to mean
                             % (arg1, arg2).

%\newcommand{\ip}[2]{\langle #1 | #2\rangle}
                             % This is an alternative definition of
                             % \ip that is commented out.

\begin{document}             % End of preamble and beginning of text.

\maketitle                   % Produces the title.

This is an example input file.  Comparing it with
the output it generates can show you how to
produce a simple document of your own.

\section{Ordinary Text}      % Produces section heading.  Lower-level
                             % sections are begun with similar
                             % \subsection and \subsubsection commands.

The ends  of words and sentences are marked
  by   spaces. It  doesn't matter how many
spaces    you type; one is as good as 100.  The
end of   a line counts as a space.

One   or more   blank lines denote the  end
of  a paragraph.

Since any number of consecutive spaces are treated
like a single one, the formatting of the input
file makes no difference to
      \LaTeX,                % The \LaTeX command generates the LaTeX logo.
but it makes a difference to you.  When you use
\LaTeX, making your input file as easy to read
as possible will be a great help as you write
your document and when you change it.  This sample
file shows how you can add comments to your own input
file.

Because printing is different from typewriting,
there are a number of things that you have to do
differently when preparing an input file than if
you were just typing the document directly.
Quotation marks like
       ``this''
have to be handled specially, as do quotes within
quotes:
       ``\,`this'            % \, separates the double and single quote.
        is what I just
        wrote, not  `that'\,''.

Dashes come in three sizes: an
       intra-word
dash, a medium dash for number ranges like
       1--2,
and a punctuation
       dash---like
this.

A sentence-ending space should be larger than the
space between words within a sentence.  You
sometimes have to type special commands in
conjunction with punctuation characters to get
this right, as in the following sentence.
       Gnats, gnus, etc.\ all  % `\ ' makes an inter-word space.
       begin with G\@.         % \@ marks end-of-sentence punctuation.
You should check the spaces after periods when
reading your output to make sure you haven't
forgotten any special cases.  Generating an
ellipsis
       \ldots\               % `\ ' is needed after `\ldots' because TeX
                             % ignores spaces after command names like \ldots
                             % made from \ + letters.
                             %
                             % Note how a `%' character causes TeX to ignore
                             % the end of the input line, so these blank lines
                             % do not start a new paragraph.
                             %
with the right spacing around the periods requires
a special command.

\LaTeX\ interprets some common characters as
commands, so you must type special commands to
generate them.  These characters include the
following:
       \$ \& \% \# \{ and \}.

In printing, text is usually emphasized with an
       \emph{italic}
type style.

\begin{em}
   A long segment of text can also be emphasized
   in this way.  Text within such a segment can be
   given \emph{additional} emphasis.
\end{em}

It is sometimes necessary to prevent \LaTeX\ from
breaking a line where it might otherwise do so.
This may be at a space, as between the ``Mr.''\ and
``Jones'' in
       ``Mr.~Jones'',        % ~ produces an unbreakable interword space.
or within a word---especially when the word is a
symbol like
       \mbox{\emph{itemnum}}
that makes little sense when hyphenated across
lines.

Footnotes\footnote{This is an example of a footnote.}
pose no problem.

\LaTeX\ is good at typesetting mathematical formulas
like
       \( x-3y + z = 7 \)
or
       \( a_{1} > x^{2n} + y^{2n} > x' \)
or
       \( \ip{A}{B} = \sum_{i} a_{i} b_{i} \).
The spaces you type in a formula are
ignored.  Remember that a letter like
       $x$                   % $ ... $  and  \( ... \)  are equivalent
is a formula when it denotes a mathematical
symbol, and it should be typed as one.

\section{Displayed Text}

Text is displayed by indenting it from the left
margin.  Quotations are commonly displayed.  There
are short quotations
\begin{quote}
   This is a short quotation.  It consists of a
   single paragraph of text.  See how it is formatted.
\end{quote}
and longer ones.
\begin{quotation}
   This is a longer quotation.  It consists of two
   paragraphs of text, neither of which are
   particularly interesting.

   This is the second paragraph of the quotation.  It
   is just as dull as the first paragraph.
\end{quotation}
Another frequently-displayed structure is a list.
The following is an example of an \emph{itemized}
list.
\begin{itemize}
   \item This is the first item of an itemized list.
         Each item in the list is marked with a ``tick''.
         You don't have to worry about what kind of tick
         mark is used.

   \item This is the second item of the list.  It
         contains another list nested inside it.  The inner
         list is an \emph{enumerated} list.
         \begin{enumerate}
            \item This is the first item of an enumerated
                  list that is nested within the itemized list.

            \item This is the second item of the inner list.
                  \LaTeX\ allows you to nest lists deeper than
                  you really should.
         \end{enumerate}
         This is the rest of the second item of the outer
         list.  It is no more interesting than any other
         part of the item.
   \item This is the third item of the list.
\end{itemize}
You can even display poetry.
\begin{verse}
   There is an environment
    for verse \\             % The \\ command separates lines
   Whose features some poets % within a stanza.
   will curse.

                             % One or more blank lines separate stanzas.

   For instead of making\\
   Them do \emph{all} line breaking, \\
   It allows them to put too many words on a line when they'd rather be
   forced to be terse.
\end{verse}

Mathematical formulas may also be displayed.  A
displayed formula
is
one-line long; multiline
formulas require special formatting instructions.
   \[  \ip{\Gamma}{\psi'} = x'' + y^{2} + z_{i}^{n}\]
Don't start a paragraph with a displayed equation,
nor make one a paragraph by itself.

\end{document}               % End of document.

\section{Bench Huge Section 556}
% synthetic bench section 556
% This is a sample LaTeX input file.  (Version of 12 August 2004.)
%
% A '%' character causes TeX to ignore all remaining text on the line,
% and is used for comments like this one.

\documentclass{article}      % Specifies the document class

                             % The preamble begins here.
\title{An Example Document}  % Declares the document's title.
\author{Leslie Lamport}      % Declares the author's name.
\date{January 21, 1994}      % Deleting this command produces today's date.

\newcommand{\ip}[2]{(#1, #2)}
                             % Defines \ip{arg1}{arg2} to mean
                             % (arg1, arg2).

%\newcommand{\ip}[2]{\langle #1 | #2\rangle}
                             % This is an alternative definition of
                             % \ip that is commented out.

\begin{document}             % End of preamble and beginning of text.

\maketitle                   % Produces the title.

This is an example input file.  Comparing it with
the output it generates can show you how to
produce a simple document of your own.

\section{Ordinary Text}      % Produces section heading.  Lower-level
                             % sections are begun with similar
                             % \subsection and \subsubsection commands.

The ends  of words and sentences are marked
  by   spaces. It  doesn't matter how many
spaces    you type; one is as good as 100.  The
end of   a line counts as a space.

One   or more   blank lines denote the  end
of  a paragraph.

Since any number of consecutive spaces are treated
like a single one, the formatting of the input
file makes no difference to
      \LaTeX,                % The \LaTeX command generates the LaTeX logo.
but it makes a difference to you.  When you use
\LaTeX, making your input file as easy to read
as possible will be a great help as you write
your document and when you change it.  This sample
file shows how you can add comments to your own input
file.

Because printing is different from typewriting,
there are a number of things that you have to do
differently when preparing an input file than if
you were just typing the document directly.
Quotation marks like
       ``this''
have to be handled specially, as do quotes within
quotes:
       ``\,`this'            % \, separates the double and single quote.
        is what I just
        wrote, not  `that'\,''.

Dashes come in three sizes: an
       intra-word
dash, a medium dash for number ranges like
       1--2,
and a punctuation
       dash---like
this.

A sentence-ending space should be larger than the
space between words within a sentence.  You
sometimes have to type special commands in
conjunction with punctuation characters to get
this right, as in the following sentence.
       Gnats, gnus, etc.\ all  % `\ ' makes an inter-word space.
       begin with G\@.         % \@ marks end-of-sentence punctuation.
You should check the spaces after periods when
reading your output to make sure you haven't
forgotten any special cases.  Generating an
ellipsis
       \ldots\               % `\ ' is needed after `\ldots' because TeX
                             % ignores spaces after command names like \ldots
                             % made from \ + letters.
                             %
                             % Note how a `%' character causes TeX to ignore
                             % the end of the input line, so these blank lines
                             % do not start a new paragraph.
                             %
with the right spacing around the periods requires
a special command.

\LaTeX\ interprets some common characters as
commands, so you must type special commands to
generate them.  These characters include the
following:
       \$ \& \% \# \{ and \}.

In printing, text is usually emphasized with an
       \emph{italic}
type style.

\begin{em}
   A long segment of text can also be emphasized
   in this way.  Text within such a segment can be
   given \emph{additional} emphasis.
\end{em}

It is sometimes necessary to prevent \LaTeX\ from
breaking a line where it might otherwise do so.
This may be at a space, as between the ``Mr.''\ and
``Jones'' in
       ``Mr.~Jones'',        % ~ produces an unbreakable interword space.
or within a word---especially when the word is a
symbol like
       \mbox{\emph{itemnum}}
that makes little sense when hyphenated across
lines.

Footnotes\footnote{This is an example of a footnote.}
pose no problem.

\LaTeX\ is good at typesetting mathematical formulas
like
       \( x-3y + z = 7 \)
or
       \( a_{1} > x^{2n} + y^{2n} > x' \)
or
       \( \ip{A}{B} = \sum_{i} a_{i} b_{i} \).
The spaces you type in a formula are
ignored.  Remember that a letter like
       $x$                   % $ ... $  and  \( ... \)  are equivalent
is a formula when it denotes a mathematical
symbol, and it should be typed as one.

\section{Displayed Text}

Text is displayed by indenting it from the left
margin.  Quotations are commonly displayed.  There
are short quotations
\begin{quote}
   This is a short quotation.  It consists of a
   single paragraph of text.  See how it is formatted.
\end{quote}
and longer ones.
\begin{quotation}
   This is a longer quotation.  It consists of two
   paragraphs of text, neither of which are
   particularly interesting.

   This is the second paragraph of the quotation.  It
   is just as dull as the first paragraph.
\end{quotation}
Another frequently-displayed structure is a list.
The following is an example of an \emph{itemized}
list.
\begin{itemize}
   \item This is the first item of an itemized list.
         Each item in the list is marked with a ``tick''.
         You don't have to worry about what kind of tick
         mark is used.

   \item This is the second item of the list.  It
         contains another list nested inside it.  The inner
         list is an \emph{enumerated} list.
         \begin{enumerate}
            \item This is the first item of an enumerated
                  list that is nested within the itemized list.

            \item This is the second item of the inner list.
                  \LaTeX\ allows you to nest lists deeper than
                  you really should.
         \end{enumerate}
         This is the rest of the second item of the outer
         list.  It is no more interesting than any other
         part of the item.
   \item This is the third item of the list.
\end{itemize}
You can even display poetry.
\begin{verse}
   There is an environment
    for verse \\             % The \\ command separates lines
   Whose features some poets % within a stanza.
   will curse.

                             % One or more blank lines separate stanzas.

   For instead of making\\
   Them do \emph{all} line breaking, \\
   It allows them to put too many words on a line when they'd rather be
   forced to be terse.
\end{verse}

Mathematical formulas may also be displayed.  A
displayed formula
is
one-line long; multiline
formulas require special formatting instructions.
   \[  \ip{\Gamma}{\psi'} = x'' + y^{2} + z_{i}^{n}\]
Don't start a paragraph with a displayed equation,
nor make one a paragraph by itself.

\end{document}               % End of document.

\section{Bench Huge Section 557}
% synthetic bench section 557
% This is a sample LaTeX input file.  (Version of 12 August 2004.)
%
% A '%' character causes TeX to ignore all remaining text on the line,
% and is used for comments like this one.

\documentclass{article}      % Specifies the document class

                             % The preamble begins here.
\title{An Example Document}  % Declares the document's title.
\author{Leslie Lamport}      % Declares the author's name.
\date{January 21, 1994}      % Deleting this command produces today's date.

\newcommand{\ip}[2]{(#1, #2)}
                             % Defines \ip{arg1}{arg2} to mean
                             % (arg1, arg2).

%\newcommand{\ip}[2]{\langle #1 | #2\rangle}
                             % This is an alternative definition of
                             % \ip that is commented out.

\begin{document}             % End of preamble and beginning of text.

\maketitle                   % Produces the title.

This is an example input file.  Comparing it with
the output it generates can show you how to
produce a simple document of your own.

\section{Ordinary Text}      % Produces section heading.  Lower-level
                             % sections are begun with similar
                             % \subsection and \subsubsection commands.

The ends  of words and sentences are marked
  by   spaces. It  doesn't matter how many
spaces    you type; one is as good as 100.  The
end of   a line counts as a space.

One   or more   blank lines denote the  end
of  a paragraph.

Since any number of consecutive spaces are treated
like a single one, the formatting of the input
file makes no difference to
      \LaTeX,                % The \LaTeX command generates the LaTeX logo.
but it makes a difference to you.  When you use
\LaTeX, making your input file as easy to read
as possible will be a great help as you write
your document and when you change it.  This sample
file shows how you can add comments to your own input
file.

Because printing is different from typewriting,
there are a number of things that you have to do
differently when preparing an input file than if
you were just typing the document directly.
Quotation marks like
       ``this''
have to be handled specially, as do quotes within
quotes:
       ``\,`this'            % \, separates the double and single quote.
        is what I just
        wrote, not  `that'\,''.

Dashes come in three sizes: an
       intra-word
dash, a medium dash for number ranges like
       1--2,
and a punctuation
       dash---like
this.

A sentence-ending space should be larger than the
space between words within a sentence.  You
sometimes have to type special commands in
conjunction with punctuation characters to get
this right, as in the following sentence.
       Gnats, gnus, etc.\ all  % `\ ' makes an inter-word space.
       begin with G\@.         % \@ marks end-of-sentence punctuation.
You should check the spaces after periods when
reading your output to make sure you haven't
forgotten any special cases.  Generating an
ellipsis
       \ldots\               % `\ ' is needed after `\ldots' because TeX
                             % ignores spaces after command names like \ldots
                             % made from \ + letters.
                             %
                             % Note how a `%' character causes TeX to ignore
                             % the end of the input line, so these blank lines
                             % do not start a new paragraph.
                             %
with the right spacing around the periods requires
a special command.

\LaTeX\ interprets some common characters as
commands, so you must type special commands to
generate them.  These characters include the
following:
       \$ \& \% \# \{ and \}.

In printing, text is usually emphasized with an
       \emph{italic}
type style.

\begin{em}
   A long segment of text can also be emphasized
   in this way.  Text within such a segment can be
   given \emph{additional} emphasis.
\end{em}

It is sometimes necessary to prevent \LaTeX\ from
breaking a line where it might otherwise do so.
This may be at a space, as between the ``Mr.''\ and
``Jones'' in
       ``Mr.~Jones'',        % ~ produces an unbreakable interword space.
or within a word---especially when the word is a
symbol like
       \mbox{\emph{itemnum}}
that makes little sense when hyphenated across
lines.

Footnotes\footnote{This is an example of a footnote.}
pose no problem.

\LaTeX\ is good at typesetting mathematical formulas
like
       \( x-3y + z = 7 \)
or
       \( a_{1} > x^{2n} + y^{2n} > x' \)
or
       \( \ip{A}{B} = \sum_{i} a_{i} b_{i} \).
The spaces you type in a formula are
ignored.  Remember that a letter like
       $x$                   % $ ... $  and  \( ... \)  are equivalent
is a formula when it denotes a mathematical
symbol, and it should be typed as one.

\section{Displayed Text}

Text is displayed by indenting it from the left
margin.  Quotations are commonly displayed.  There
are short quotations
\begin{quote}
   This is a short quotation.  It consists of a
   single paragraph of text.  See how it is formatted.
\end{quote}
and longer ones.
\begin{quotation}
   This is a longer quotation.  It consists of two
   paragraphs of text, neither of which are
   particularly interesting.

   This is the second paragraph of the quotation.  It
   is just as dull as the first paragraph.
\end{quotation}
Another frequently-displayed structure is a list.
The following is an example of an \emph{itemized}
list.
\begin{itemize}
   \item This is the first item of an itemized list.
         Each item in the list is marked with a ``tick''.
         You don't have to worry about what kind of tick
         mark is used.

   \item This is the second item of the list.  It
         contains another list nested inside it.  The inner
         list is an \emph{enumerated} list.
         \begin{enumerate}
            \item This is the first item of an enumerated
                  list that is nested within the itemized list.

            \item This is the second item of the inner list.
                  \LaTeX\ allows you to nest lists deeper than
                  you really should.
         \end{enumerate}
         This is the rest of the second item of the outer
         list.  It is no more interesting than any other
         part of the item.
   \item This is the third item of the list.
\end{itemize}
You can even display poetry.
\begin{verse}
   There is an environment
    for verse \\             % The \\ command separates lines
   Whose features some poets % within a stanza.
   will curse.

                             % One or more blank lines separate stanzas.

   For instead of making\\
   Them do \emph{all} line breaking, \\
   It allows them to put too many words on a line when they'd rather be
   forced to be terse.
\end{verse}

Mathematical formulas may also be displayed.  A
displayed formula
is
one-line long; multiline
formulas require special formatting instructions.
   \[  \ip{\Gamma}{\psi'} = x'' + y^{2} + z_{i}^{n}\]
Don't start a paragraph with a displayed equation,
nor make one a paragraph by itself.

\end{document}               % End of document.

\section{Bench Huge Section 558}
% synthetic bench section 558
% This is a sample LaTeX input file.  (Version of 12 August 2004.)
%
% A '%' character causes TeX to ignore all remaining text on the line,
% and is used for comments like this one.

\documentclass{article}      % Specifies the document class

                             % The preamble begins here.
\title{An Example Document}  % Declares the document's title.
\author{Leslie Lamport}      % Declares the author's name.
\date{January 21, 1994}      % Deleting this command produces today's date.

\newcommand{\ip}[2]{(#1, #2)}
                             % Defines \ip{arg1}{arg2} to mean
                             % (arg1, arg2).

%\newcommand{\ip}[2]{\langle #1 | #2\rangle}
                             % This is an alternative definition of
                             % \ip that is commented out.

\begin{document}             % End of preamble and beginning of text.

\maketitle                   % Produces the title.

This is an example input file.  Comparing it with
the output it generates can show you how to
produce a simple document of your own.

\section{Ordinary Text}      % Produces section heading.  Lower-level
                             % sections are begun with similar
                             % \subsection and \subsubsection commands.

The ends  of words and sentences are marked
  by   spaces. It  doesn't matter how many
spaces    you type; one is as good as 100.  The
end of   a line counts as a space.

One   or more   blank lines denote the  end
of  a paragraph.

Since any number of consecutive spaces are treated
like a single one, the formatting of the input
file makes no difference to
      \LaTeX,                % The \LaTeX command generates the LaTeX logo.
but it makes a difference to you.  When you use
\LaTeX, making your input file as easy to read
as possible will be a great help as you write
your document and when you change it.  This sample
file shows how you can add comments to your own input
file.

Because printing is different from typewriting,
there are a number of things that you have to do
differently when preparing an input file than if
you were just typing the document directly.
Quotation marks like
       ``this''
have to be handled specially, as do quotes within
quotes:
       ``\,`this'            % \, separates the double and single quote.
        is what I just
        wrote, not  `that'\,''.

Dashes come in three sizes: an
       intra-word
dash, a medium dash for number ranges like
       1--2,
and a punctuation
       dash---like
this.

A sentence-ending space should be larger than the
space between words within a sentence.  You
sometimes have to type special commands in
conjunction with punctuation characters to get
this right, as in the following sentence.
       Gnats, gnus, etc.\ all  % `\ ' makes an inter-word space.
       begin with G\@.         % \@ marks end-of-sentence punctuation.
You should check the spaces after periods when
reading your output to make sure you haven't
forgotten any special cases.  Generating an
ellipsis
       \ldots\               % `\ ' is needed after `\ldots' because TeX
                             % ignores spaces after command names like \ldots
                             % made from \ + letters.
                             %
                             % Note how a `%' character causes TeX to ignore
                             % the end of the input line, so these blank lines
                             % do not start a new paragraph.
                             %
with the right spacing around the periods requires
a special command.

\LaTeX\ interprets some common characters as
commands, so you must type special commands to
generate them.  These characters include the
following:
       \$ \& \% \# \{ and \}.

In printing, text is usually emphasized with an
       \emph{italic}
type style.

\begin{em}
   A long segment of text can also be emphasized
   in this way.  Text within such a segment can be
   given \emph{additional} emphasis.
\end{em}

It is sometimes necessary to prevent \LaTeX\ from
breaking a line where it might otherwise do so.
This may be at a space, as between the ``Mr.''\ and
``Jones'' in
       ``Mr.~Jones'',        % ~ produces an unbreakable interword space.
or within a word---especially when the word is a
symbol like
       \mbox{\emph{itemnum}}
that makes little sense when hyphenated across
lines.

Footnotes\footnote{This is an example of a footnote.}
pose no problem.

\LaTeX\ is good at typesetting mathematical formulas
like
       \( x-3y + z = 7 \)
or
       \( a_{1} > x^{2n} + y^{2n} > x' \)
or
       \( \ip{A}{B} = \sum_{i} a_{i} b_{i} \).
The spaces you type in a formula are
ignored.  Remember that a letter like
       $x$                   % $ ... $  and  \( ... \)  are equivalent
is a formula when it denotes a mathematical
symbol, and it should be typed as one.

\section{Displayed Text}

Text is displayed by indenting it from the left
margin.  Quotations are commonly displayed.  There
are short quotations
\begin{quote}
   This is a short quotation.  It consists of a
   single paragraph of text.  See how it is formatted.
\end{quote}
and longer ones.
\begin{quotation}
   This is a longer quotation.  It consists of two
   paragraphs of text, neither of which are
   particularly interesting.

   This is the second paragraph of the quotation.  It
   is just as dull as the first paragraph.
\end{quotation}
Another frequently-displayed structure is a list.
The following is an example of an \emph{itemized}
list.
\begin{itemize}
   \item This is the first item of an itemized list.
         Each item in the list is marked with a ``tick''.
         You don't have to worry about what kind of tick
         mark is used.

   \item This is the second item of the list.  It
         contains another list nested inside it.  The inner
         list is an \emph{enumerated} list.
         \begin{enumerate}
            \item This is the first item of an enumerated
                  list that is nested within the itemized list.

            \item This is the second item of the inner list.
                  \LaTeX\ allows you to nest lists deeper than
                  you really should.
         \end{enumerate}
         This is the rest of the second item of the outer
         list.  It is no more interesting than any other
         part of the item.
   \item This is the third item of the list.
\end{itemize}
You can even display poetry.
\begin{verse}
   There is an environment
    for verse \\             % The \\ command separates lines
   Whose features some poets % within a stanza.
   will curse.

                             % One or more blank lines separate stanzas.

   For instead of making\\
   Them do \emph{all} line breaking, \\
   It allows them to put too many words on a line when they'd rather be
   forced to be terse.
\end{verse}

Mathematical formulas may also be displayed.  A
displayed formula
is
one-line long; multiline
formulas require special formatting instructions.
   \[  \ip{\Gamma}{\psi'} = x'' + y^{2} + z_{i}^{n}\]
Don't start a paragraph with a displayed equation,
nor make one a paragraph by itself.

\end{document}               % End of document.

\section{Bench Huge Section 559}
% synthetic bench section 559
% This is a sample LaTeX input file.  (Version of 12 August 2004.)
%
% A '%' character causes TeX to ignore all remaining text on the line,
% and is used for comments like this one.

\documentclass{article}      % Specifies the document class

                             % The preamble begins here.
\title{An Example Document}  % Declares the document's title.
\author{Leslie Lamport}      % Declares the author's name.
\date{January 21, 1994}      % Deleting this command produces today's date.

\newcommand{\ip}[2]{(#1, #2)}
                             % Defines \ip{arg1}{arg2} to mean
                             % (arg1, arg2).

%\newcommand{\ip}[2]{\langle #1 | #2\rangle}
                             % This is an alternative definition of
                             % \ip that is commented out.

\begin{document}             % End of preamble and beginning of text.

\maketitle                   % Produces the title.

This is an example input file.  Comparing it with
the output it generates can show you how to
produce a simple document of your own.

\section{Ordinary Text}      % Produces section heading.  Lower-level
                             % sections are begun with similar
                             % \subsection and \subsubsection commands.

The ends  of words and sentences are marked
  by   spaces. It  doesn't matter how many
spaces    you type; one is as good as 100.  The
end of   a line counts as a space.

One   or more   blank lines denote the  end
of  a paragraph.

Since any number of consecutive spaces are treated
like a single one, the formatting of the input
file makes no difference to
      \LaTeX,                % The \LaTeX command generates the LaTeX logo.
but it makes a difference to you.  When you use
\LaTeX, making your input file as easy to read
as possible will be a great help as you write
your document and when you change it.  This sample
file shows how you can add comments to your own input
file.

Because printing is different from typewriting,
there are a number of things that you have to do
differently when preparing an input file than if
you were just typing the document directly.
Quotation marks like
       ``this''
have to be handled specially, as do quotes within
quotes:
       ``\,`this'            % \, separates the double and single quote.
        is what I just
        wrote, not  `that'\,''.

Dashes come in three sizes: an
       intra-word
dash, a medium dash for number ranges like
       1--2,
and a punctuation
       dash---like
this.

A sentence-ending space should be larger than the
space between words within a sentence.  You
sometimes have to type special commands in
conjunction with punctuation characters to get
this right, as in the following sentence.
       Gnats, gnus, etc.\ all  % `\ ' makes an inter-word space.
       begin with G\@.         % \@ marks end-of-sentence punctuation.
You should check the spaces after periods when
reading your output to make sure you haven't
forgotten any special cases.  Generating an
ellipsis
       \ldots\               % `\ ' is needed after `\ldots' because TeX
                             % ignores spaces after command names like \ldots
                             % made from \ + letters.
                             %
                             % Note how a `%' character causes TeX to ignore
                             % the end of the input line, so these blank lines
                             % do not start a new paragraph.
                             %
with the right spacing around the periods requires
a special command.

\LaTeX\ interprets some common characters as
commands, so you must type special commands to
generate them.  These characters include the
following:
       \$ \& \% \# \{ and \}.

In printing, text is usually emphasized with an
       \emph{italic}
type style.

\begin{em}
   A long segment of text can also be emphasized
   in this way.  Text within such a segment can be
   given \emph{additional} emphasis.
\end{em}

It is sometimes necessary to prevent \LaTeX\ from
breaking a line where it might otherwise do so.
This may be at a space, as between the ``Mr.''\ and
``Jones'' in
       ``Mr.~Jones'',        % ~ produces an unbreakable interword space.
or within a word---especially when the word is a
symbol like
       \mbox{\emph{itemnum}}
that makes little sense when hyphenated across
lines.

Footnotes\footnote{This is an example of a footnote.}
pose no problem.

\LaTeX\ is good at typesetting mathematical formulas
like
       \( x-3y + z = 7 \)
or
       \( a_{1} > x^{2n} + y^{2n} > x' \)
or
       \( \ip{A}{B} = \sum_{i} a_{i} b_{i} \).
The spaces you type in a formula are
ignored.  Remember that a letter like
       $x$                   % $ ... $  and  \( ... \)  are equivalent
is a formula when it denotes a mathematical
symbol, and it should be typed as one.

\section{Displayed Text}

Text is displayed by indenting it from the left
margin.  Quotations are commonly displayed.  There
are short quotations
\begin{quote}
   This is a short quotation.  It consists of a
   single paragraph of text.  See how it is formatted.
\end{quote}
and longer ones.
\begin{quotation}
   This is a longer quotation.  It consists of two
   paragraphs of text, neither of which are
   particularly interesting.

   This is the second paragraph of the quotation.  It
   is just as dull as the first paragraph.
\end{quotation}
Another frequently-displayed structure is a list.
The following is an example of an \emph{itemized}
list.
\begin{itemize}
   \item This is the first item of an itemized list.
         Each item in the list is marked with a ``tick''.
         You don't have to worry about what kind of tick
         mark is used.

   \item This is the second item of the list.  It
         contains another list nested inside it.  The inner
         list is an \emph{enumerated} list.
         \begin{enumerate}
            \item This is the first item of an enumerated
                  list that is nested within the itemized list.

            \item This is the second item of the inner list.
                  \LaTeX\ allows you to nest lists deeper than
                  you really should.
         \end{enumerate}
         This is the rest of the second item of the outer
         list.  It is no more interesting than any other
         part of the item.
   \item This is the third item of the list.
\end{itemize}
You can even display poetry.
\begin{verse}
   There is an environment
    for verse \\             % The \\ command separates lines
   Whose features some poets % within a stanza.
   will curse.

                             % One or more blank lines separate stanzas.

   For instead of making\\
   Them do \emph{all} line breaking, \\
   It allows them to put too many words on a line when they'd rather be
   forced to be terse.
\end{verse}

Mathematical formulas may also be displayed.  A
displayed formula
is
one-line long; multiline
formulas require special formatting instructions.
   \[  \ip{\Gamma}{\psi'} = x'' + y^{2} + z_{i}^{n}\]
Don't start a paragraph with a displayed equation,
nor make one a paragraph by itself.

\end{document}               % End of document.

\section{Bench Huge Section 560}
% synthetic bench section 560
% This is a sample LaTeX input file.  (Version of 12 August 2004.)
%
% A '%' character causes TeX to ignore all remaining text on the line,
% and is used for comments like this one.

\documentclass{article}      % Specifies the document class

                             % The preamble begins here.
\title{An Example Document}  % Declares the document's title.
\author{Leslie Lamport}      % Declares the author's name.
\date{January 21, 1994}      % Deleting this command produces today's date.

\newcommand{\ip}[2]{(#1, #2)}
                             % Defines \ip{arg1}{arg2} to mean
                             % (arg1, arg2).

%\newcommand{\ip}[2]{\langle #1 | #2\rangle}
                             % This is an alternative definition of
                             % \ip that is commented out.

\begin{document}             % End of preamble and beginning of text.

\maketitle                   % Produces the title.

This is an example input file.  Comparing it with
the output it generates can show you how to
produce a simple document of your own.

\section{Ordinary Text}      % Produces section heading.  Lower-level
                             % sections are begun with similar
                             % \subsection and \subsubsection commands.

The ends  of words and sentences are marked
  by   spaces. It  doesn't matter how many
spaces    you type; one is as good as 100.  The
end of   a line counts as a space.

One   or more   blank lines denote the  end
of  a paragraph.

Since any number of consecutive spaces are treated
like a single one, the formatting of the input
file makes no difference to
      \LaTeX,                % The \LaTeX command generates the LaTeX logo.
but it makes a difference to you.  When you use
\LaTeX, making your input file as easy to read
as possible will be a great help as you write
your document and when you change it.  This sample
file shows how you can add comments to your own input
file.

Because printing is different from typewriting,
there are a number of things that you have to do
differently when preparing an input file than if
you were just typing the document directly.
Quotation marks like
       ``this''
have to be handled specially, as do quotes within
quotes:
       ``\,`this'            % \, separates the double and single quote.
        is what I just
        wrote, not  `that'\,''.

Dashes come in three sizes: an
       intra-word
dash, a medium dash for number ranges like
       1--2,
and a punctuation
       dash---like
this.

A sentence-ending space should be larger than the
space between words within a sentence.  You
sometimes have to type special commands in
conjunction with punctuation characters to get
this right, as in the following sentence.
       Gnats, gnus, etc.\ all  % `\ ' makes an inter-word space.
       begin with G\@.         % \@ marks end-of-sentence punctuation.
You should check the spaces after periods when
reading your output to make sure you haven't
forgotten any special cases.  Generating an
ellipsis
       \ldots\               % `\ ' is needed after `\ldots' because TeX
                             % ignores spaces after command names like \ldots
                             % made from \ + letters.
                             %
                             % Note how a `%' character causes TeX to ignore
                             % the end of the input line, so these blank lines
                             % do not start a new paragraph.
                             %
with the right spacing around the periods requires
a special command.

\LaTeX\ interprets some common characters as
commands, so you must type special commands to
generate them.  These characters include the
following:
       \$ \& \% \# \{ and \}.

In printing, text is usually emphasized with an
       \emph{italic}
type style.

\begin{em}
   A long segment of text can also be emphasized
   in this way.  Text within such a segment can be
   given \emph{additional} emphasis.
\end{em}

It is sometimes necessary to prevent \LaTeX\ from
breaking a line where it might otherwise do so.
This may be at a space, as between the ``Mr.''\ and
``Jones'' in
       ``Mr.~Jones'',        % ~ produces an unbreakable interword space.
or within a word---especially when the word is a
symbol like
       \mbox{\emph{itemnum}}
that makes little sense when hyphenated across
lines.

Footnotes\footnote{This is an example of a footnote.}
pose no problem.

\LaTeX\ is good at typesetting mathematical formulas
like
       \( x-3y + z = 7 \)
or
       \( a_{1} > x^{2n} + y^{2n} > x' \)
or
       \( \ip{A}{B} = \sum_{i} a_{i} b_{i} \).
The spaces you type in a formula are
ignored.  Remember that a letter like
       $x$                   % $ ... $  and  \( ... \)  are equivalent
is a formula when it denotes a mathematical
symbol, and it should be typed as one.

\section{Displayed Text}

Text is displayed by indenting it from the left
margin.  Quotations are commonly displayed.  There
are short quotations
\begin{quote}
   This is a short quotation.  It consists of a
   single paragraph of text.  See how it is formatted.
\end{quote}
and longer ones.
\begin{quotation}
   This is a longer quotation.  It consists of two
   paragraphs of text, neither of which are
   particularly interesting.

   This is the second paragraph of the quotation.  It
   is just as dull as the first paragraph.
\end{quotation}
Another frequently-displayed structure is a list.
The following is an example of an \emph{itemized}
list.
\begin{itemize}
   \item This is the first item of an itemized list.
         Each item in the list is marked with a ``tick''.
         You don't have to worry about what kind of tick
         mark is used.

   \item This is the second item of the list.  It
         contains another list nested inside it.  The inner
         list is an \emph{enumerated} list.
         \begin{enumerate}
            \item This is the first item of an enumerated
                  list that is nested within the itemized list.

            \item This is the second item of the inner list.
                  \LaTeX\ allows you to nest lists deeper than
                  you really should.
         \end{enumerate}
         This is the rest of the second item of the outer
         list.  It is no more interesting than any other
         part of the item.
   \item This is the third item of the list.
\end{itemize}
You can even display poetry.
\begin{verse}
   There is an environment
    for verse \\             % The \\ command separates lines
   Whose features some poets % within a stanza.
   will curse.

                             % One or more blank lines separate stanzas.

   For instead of making\\
   Them do \emph{all} line breaking, \\
   It allows them to put too many words on a line when they'd rather be
   forced to be terse.
\end{verse}

Mathematical formulas may also be displayed.  A
displayed formula
is
one-line long; multiline
formulas require special formatting instructions.
   \[  \ip{\Gamma}{\psi'} = x'' + y^{2} + z_{i}^{n}\]
Don't start a paragraph with a displayed equation,
nor make one a paragraph by itself.

\end{document}               % End of document.

\section{Bench Huge Section 561}
% synthetic bench section 561
% This is a sample LaTeX input file.  (Version of 12 August 2004.)
%
% A '%' character causes TeX to ignore all remaining text on the line,
% and is used for comments like this one.

\documentclass{article}      % Specifies the document class

                             % The preamble begins here.
\title{An Example Document}  % Declares the document's title.
\author{Leslie Lamport}      % Declares the author's name.
\date{January 21, 1994}      % Deleting this command produces today's date.

\newcommand{\ip}[2]{(#1, #2)}
                             % Defines \ip{arg1}{arg2} to mean
                             % (arg1, arg2).

%\newcommand{\ip}[2]{\langle #1 | #2\rangle}
                             % This is an alternative definition of
                             % \ip that is commented out.

\begin{document}             % End of preamble and beginning of text.

\maketitle                   % Produces the title.

This is an example input file.  Comparing it with
the output it generates can show you how to
produce a simple document of your own.

\section{Ordinary Text}      % Produces section heading.  Lower-level
                             % sections are begun with similar
                             % \subsection and \subsubsection commands.

The ends  of words and sentences are marked
  by   spaces. It  doesn't matter how many
spaces    you type; one is as good as 100.  The
end of   a line counts as a space.

One   or more   blank lines denote the  end
of  a paragraph.

Since any number of consecutive spaces are treated
like a single one, the formatting of the input
file makes no difference to
      \LaTeX,                % The \LaTeX command generates the LaTeX logo.
but it makes a difference to you.  When you use
\LaTeX, making your input file as easy to read
as possible will be a great help as you write
your document and when you change it.  This sample
file shows how you can add comments to your own input
file.

Because printing is different from typewriting,
there are a number of things that you have to do
differently when preparing an input file than if
you were just typing the document directly.
Quotation marks like
       ``this''
have to be handled specially, as do quotes within
quotes:
       ``\,`this'            % \, separates the double and single quote.
        is what I just
        wrote, not  `that'\,''.

Dashes come in three sizes: an
       intra-word
dash, a medium dash for number ranges like
       1--2,
and a punctuation
       dash---like
this.

A sentence-ending space should be larger than the
space between words within a sentence.  You
sometimes have to type special commands in
conjunction with punctuation characters to get
this right, as in the following sentence.
       Gnats, gnus, etc.\ all  % `\ ' makes an inter-word space.
       begin with G\@.         % \@ marks end-of-sentence punctuation.
You should check the spaces after periods when
reading your output to make sure you haven't
forgotten any special cases.  Generating an
ellipsis
       \ldots\               % `\ ' is needed after `\ldots' because TeX
                             % ignores spaces after command names like \ldots
                             % made from \ + letters.
                             %
                             % Note how a `%' character causes TeX to ignore
                             % the end of the input line, so these blank lines
                             % do not start a new paragraph.
                             %
with the right spacing around the periods requires
a special command.

\LaTeX\ interprets some common characters as
commands, so you must type special commands to
generate them.  These characters include the
following:
       \$ \& \% \# \{ and \}.

In printing, text is usually emphasized with an
       \emph{italic}
type style.

\begin{em}
   A long segment of text can also be emphasized
   in this way.  Text within such a segment can be
   given \emph{additional} emphasis.
\end{em}

It is sometimes necessary to prevent \LaTeX\ from
breaking a line where it might otherwise do so.
This may be at a space, as between the ``Mr.''\ and
``Jones'' in
       ``Mr.~Jones'',        % ~ produces an unbreakable interword space.
or within a word---especially when the word is a
symbol like
       \mbox{\emph{itemnum}}
that makes little sense when hyphenated across
lines.

Footnotes\footnote{This is an example of a footnote.}
pose no problem.

\LaTeX\ is good at typesetting mathematical formulas
like
       \( x-3y + z = 7 \)
or
       \( a_{1} > x^{2n} + y^{2n} > x' \)
or
       \( \ip{A}{B} = \sum_{i} a_{i} b_{i} \).
The spaces you type in a formula are
ignored.  Remember that a letter like
       $x$                   % $ ... $  and  \( ... \)  are equivalent
is a formula when it denotes a mathematical
symbol, and it should be typed as one.

\section{Displayed Text}

Text is displayed by indenting it from the left
margin.  Quotations are commonly displayed.  There
are short quotations
\begin{quote}
   This is a short quotation.  It consists of a
   single paragraph of text.  See how it is formatted.
\end{quote}
and longer ones.
\begin{quotation}
   This is a longer quotation.  It consists of two
   paragraphs of text, neither of which are
   particularly interesting.

   This is the second paragraph of the quotation.  It
   is just as dull as the first paragraph.
\end{quotation}
Another frequently-displayed structure is a list.
The following is an example of an \emph{itemized}
list.
\begin{itemize}
   \item This is the first item of an itemized list.
         Each item in the list is marked with a ``tick''.
         You don't have to worry about what kind of tick
         mark is used.

   \item This is the second item of the list.  It
         contains another list nested inside it.  The inner
         list is an \emph{enumerated} list.
         \begin{enumerate}
            \item This is the first item of an enumerated
                  list that is nested within the itemized list.

            \item This is the second item of the inner list.
                  \LaTeX\ allows you to nest lists deeper than
                  you really should.
         \end{enumerate}
         This is the rest of the second item of the outer
         list.  It is no more interesting than any other
         part of the item.
   \item This is the third item of the list.
\end{itemize}
You can even display poetry.
\begin{verse}
   There is an environment
    for verse \\             % The \\ command separates lines
   Whose features some poets % within a stanza.
   will curse.

                             % One or more blank lines separate stanzas.

   For instead of making\\
   Them do \emph{all} line breaking, \\
   It allows them to put too many words on a line when they'd rather be
   forced to be terse.
\end{verse}

Mathematical formulas may also be displayed.  A
displayed formula
is
one-line long; multiline
formulas require special formatting instructions.
   \[  \ip{\Gamma}{\psi'} = x'' + y^{2} + z_{i}^{n}\]
Don't start a paragraph with a displayed equation,
nor make one a paragraph by itself.

\end{document}               % End of document.

\section{Bench Huge Section 562}
% synthetic bench section 562
% This is a sample LaTeX input file.  (Version of 12 August 2004.)
%
% A '%' character causes TeX to ignore all remaining text on the line,
% and is used for comments like this one.

\documentclass{article}      % Specifies the document class

                             % The preamble begins here.
\title{An Example Document}  % Declares the document's title.
\author{Leslie Lamport}      % Declares the author's name.
\date{January 21, 1994}      % Deleting this command produces today's date.

\newcommand{\ip}[2]{(#1, #2)}
                             % Defines \ip{arg1}{arg2} to mean
                             % (arg1, arg2).

%\newcommand{\ip}[2]{\langle #1 | #2\rangle}
                             % This is an alternative definition of
                             % \ip that is commented out.

\begin{document}             % End of preamble and beginning of text.

\maketitle                   % Produces the title.

This is an example input file.  Comparing it with
the output it generates can show you how to
produce a simple document of your own.

\section{Ordinary Text}      % Produces section heading.  Lower-level
                             % sections are begun with similar
                             % \subsection and \subsubsection commands.

The ends  of words and sentences are marked
  by   spaces. It  doesn't matter how many
spaces    you type; one is as good as 100.  The
end of   a line counts as a space.

One   or more   blank lines denote the  end
of  a paragraph.

Since any number of consecutive spaces are treated
like a single one, the formatting of the input
file makes no difference to
      \LaTeX,                % The \LaTeX command generates the LaTeX logo.
but it makes a difference to you.  When you use
\LaTeX, making your input file as easy to read
as possible will be a great help as you write
your document and when you change it.  This sample
file shows how you can add comments to your own input
file.

Because printing is different from typewriting,
there are a number of things that you have to do
differently when preparing an input file than if
you were just typing the document directly.
Quotation marks like
       ``this''
have to be handled specially, as do quotes within
quotes:
       ``\,`this'            % \, separates the double and single quote.
        is what I just
        wrote, not  `that'\,''.

Dashes come in three sizes: an
       intra-word
dash, a medium dash for number ranges like
       1--2,
and a punctuation
       dash---like
this.

A sentence-ending space should be larger than the
space between words within a sentence.  You
sometimes have to type special commands in
conjunction with punctuation characters to get
this right, as in the following sentence.
       Gnats, gnus, etc.\ all  % `\ ' makes an inter-word space.
       begin with G\@.         % \@ marks end-of-sentence punctuation.
You should check the spaces after periods when
reading your output to make sure you haven't
forgotten any special cases.  Generating an
ellipsis
       \ldots\               % `\ ' is needed after `\ldots' because TeX
                             % ignores spaces after command names like \ldots
                             % made from \ + letters.
                             %
                             % Note how a `%' character causes TeX to ignore
                             % the end of the input line, so these blank lines
                             % do not start a new paragraph.
                             %
with the right spacing around the periods requires
a special command.

\LaTeX\ interprets some common characters as
commands, so you must type special commands to
generate them.  These characters include the
following:
       \$ \& \% \# \{ and \}.

In printing, text is usually emphasized with an
       \emph{italic}
type style.

\begin{em}
   A long segment of text can also be emphasized
   in this way.  Text within such a segment can be
   given \emph{additional} emphasis.
\end{em}

It is sometimes necessary to prevent \LaTeX\ from
breaking a line where it might otherwise do so.
This may be at a space, as between the ``Mr.''\ and
``Jones'' in
       ``Mr.~Jones'',        % ~ produces an unbreakable interword space.
or within a word---especially when the word is a
symbol like
       \mbox{\emph{itemnum}}
that makes little sense when hyphenated across
lines.

Footnotes\footnote{This is an example of a footnote.}
pose no problem.

\LaTeX\ is good at typesetting mathematical formulas
like
       \( x-3y + z = 7 \)
or
       \( a_{1} > x^{2n} + y^{2n} > x' \)
or
       \( \ip{A}{B} = \sum_{i} a_{i} b_{i} \).
The spaces you type in a formula are
ignored.  Remember that a letter like
       $x$                   % $ ... $  and  \( ... \)  are equivalent
is a formula when it denotes a mathematical
symbol, and it should be typed as one.

\section{Displayed Text}

Text is displayed by indenting it from the left
margin.  Quotations are commonly displayed.  There
are short quotations
\begin{quote}
   This is a short quotation.  It consists of a
   single paragraph of text.  See how it is formatted.
\end{quote}
and longer ones.
\begin{quotation}
   This is a longer quotation.  It consists of two
   paragraphs of text, neither of which are
   particularly interesting.

   This is the second paragraph of the quotation.  It
   is just as dull as the first paragraph.
\end{quotation}
Another frequently-displayed structure is a list.
The following is an example of an \emph{itemized}
list.
\begin{itemize}
   \item This is the first item of an itemized list.
         Each item in the list is marked with a ``tick''.
         You don't have to worry about what kind of tick
         mark is used.

   \item This is the second item of the list.  It
         contains another list nested inside it.  The inner
         list is an \emph{enumerated} list.
         \begin{enumerate}
            \item This is the first item of an enumerated
                  list that is nested within the itemized list.

            \item This is the second item of the inner list.
                  \LaTeX\ allows you to nest lists deeper than
                  you really should.
         \end{enumerate}
         This is the rest of the second item of the outer
         list.  It is no more interesting than any other
         part of the item.
   \item This is the third item of the list.
\end{itemize}
You can even display poetry.
\begin{verse}
   There is an environment
    for verse \\             % The \\ command separates lines
   Whose features some poets % within a stanza.
   will curse.

                             % One or more blank lines separate stanzas.

   For instead of making\\
   Them do \emph{all} line breaking, \\
   It allows them to put too many words on a line when they'd rather be
   forced to be terse.
\end{verse}

Mathematical formulas may also be displayed.  A
displayed formula
is
one-line long; multiline
formulas require special formatting instructions.
   \[  \ip{\Gamma}{\psi'} = x'' + y^{2} + z_{i}^{n}\]
Don't start a paragraph with a displayed equation,
nor make one a paragraph by itself.

\end{document}               % End of document.

\section{Bench Huge Section 563}
% synthetic bench section 563
% This is a sample LaTeX input file.  (Version of 12 August 2004.)
%
% A '%' character causes TeX to ignore all remaining text on the line,
% and is used for comments like this one.

\documentclass{article}      % Specifies the document class

                             % The preamble begins here.
\title{An Example Document}  % Declares the document's title.
\author{Leslie Lamport}      % Declares the author's name.
\date{January 21, 1994}      % Deleting this command produces today's date.

\newcommand{\ip}[2]{(#1, #2)}
                             % Defines \ip{arg1}{arg2} to mean
                             % (arg1, arg2).

%\newcommand{\ip}[2]{\langle #1 | #2\rangle}
                             % This is an alternative definition of
                             % \ip that is commented out.

\begin{document}             % End of preamble and beginning of text.

\maketitle                   % Produces the title.

This is an example input file.  Comparing it with
the output it generates can show you how to
produce a simple document of your own.

\section{Ordinary Text}      % Produces section heading.  Lower-level
                             % sections are begun with similar
                             % \subsection and \subsubsection commands.

The ends  of words and sentences are marked
  by   spaces. It  doesn't matter how many
spaces    you type; one is as good as 100.  The
end of   a line counts as a space.

One   or more   blank lines denote the  end
of  a paragraph.

Since any number of consecutive spaces are treated
like a single one, the formatting of the input
file makes no difference to
      \LaTeX,                % The \LaTeX command generates the LaTeX logo.
but it makes a difference to you.  When you use
\LaTeX, making your input file as easy to read
as possible will be a great help as you write
your document and when you change it.  This sample
file shows how you can add comments to your own input
file.

Because printing is different from typewriting,
there are a number of things that you have to do
differently when preparing an input file than if
you were just typing the document directly.
Quotation marks like
       ``this''
have to be handled specially, as do quotes within
quotes:
       ``\,`this'            % \, separates the double and single quote.
        is what I just
        wrote, not  `that'\,''.

Dashes come in three sizes: an
       intra-word
dash, a medium dash for number ranges like
       1--2,
and a punctuation
       dash---like
this.

A sentence-ending space should be larger than the
space between words within a sentence.  You
sometimes have to type special commands in
conjunction with punctuation characters to get
this right, as in the following sentence.
       Gnats, gnus, etc.\ all  % `\ ' makes an inter-word space.
       begin with G\@.         % \@ marks end-of-sentence punctuation.
You should check the spaces after periods when
reading your output to make sure you haven't
forgotten any special cases.  Generating an
ellipsis
       \ldots\               % `\ ' is needed after `\ldots' because TeX
                             % ignores spaces after command names like \ldots
                             % made from \ + letters.
                             %
                             % Note how a `%' character causes TeX to ignore
                             % the end of the input line, so these blank lines
                             % do not start a new paragraph.
                             %
with the right spacing around the periods requires
a special command.

\LaTeX\ interprets some common characters as
commands, so you must type special commands to
generate them.  These characters include the
following:
       \$ \& \% \# \{ and \}.

In printing, text is usually emphasized with an
       \emph{italic}
type style.

\begin{em}
   A long segment of text can also be emphasized
   in this way.  Text within such a segment can be
   given \emph{additional} emphasis.
\end{em}

It is sometimes necessary to prevent \LaTeX\ from
breaking a line where it might otherwise do so.
This may be at a space, as between the ``Mr.''\ and
``Jones'' in
       ``Mr.~Jones'',        % ~ produces an unbreakable interword space.
or within a word---especially when the word is a
symbol like
       \mbox{\emph{itemnum}}
that makes little sense when hyphenated across
lines.

Footnotes\footnote{This is an example of a footnote.}
pose no problem.

\LaTeX\ is good at typesetting mathematical formulas
like
       \( x-3y + z = 7 \)
or
       \( a_{1} > x^{2n} + y^{2n} > x' \)
or
       \( \ip{A}{B} = \sum_{i} a_{i} b_{i} \).
The spaces you type in a formula are
ignored.  Remember that a letter like
       $x$                   % $ ... $  and  \( ... \)  are equivalent
is a formula when it denotes a mathematical
symbol, and it should be typed as one.

\section{Displayed Text}

Text is displayed by indenting it from the left
margin.  Quotations are commonly displayed.  There
are short quotations
\begin{quote}
   This is a short quotation.  It consists of a
   single paragraph of text.  See how it is formatted.
\end{quote}
and longer ones.
\begin{quotation}
   This is a longer quotation.  It consists of two
   paragraphs of text, neither of which are
   particularly interesting.

   This is the second paragraph of the quotation.  It
   is just as dull as the first paragraph.
\end{quotation}
Another frequently-displayed structure is a list.
The following is an example of an \emph{itemized}
list.
\begin{itemize}
   \item This is the first item of an itemized list.
         Each item in the list is marked with a ``tick''.
         You don't have to worry about what kind of tick
         mark is used.

   \item This is the second item of the list.  It
         contains another list nested inside it.  The inner
         list is an \emph{enumerated} list.
         \begin{enumerate}
            \item This is the first item of an enumerated
                  list that is nested within the itemized list.

            \item This is the second item of the inner list.
                  \LaTeX\ allows you to nest lists deeper than
                  you really should.
         \end{enumerate}
         This is the rest of the second item of the outer
         list.  It is no more interesting than any other
         part of the item.
   \item This is the third item of the list.
\end{itemize}
You can even display poetry.
\begin{verse}
   There is an environment
    for verse \\             % The \\ command separates lines
   Whose features some poets % within a stanza.
   will curse.

                             % One or more blank lines separate stanzas.

   For instead of making\\
   Them do \emph{all} line breaking, \\
   It allows them to put too many words on a line when they'd rather be
   forced to be terse.
\end{verse}

Mathematical formulas may also be displayed.  A
displayed formula
is
one-line long; multiline
formulas require special formatting instructions.
   \[  \ip{\Gamma}{\psi'} = x'' + y^{2} + z_{i}^{n}\]
Don't start a paragraph with a displayed equation,
nor make one a paragraph by itself.

\end{document}               % End of document.

\section{Bench Huge Section 564}
% synthetic bench section 564
% This is a sample LaTeX input file.  (Version of 12 August 2004.)
%
% A '%' character causes TeX to ignore all remaining text on the line,
% and is used for comments like this one.

\documentclass{article}      % Specifies the document class

                             % The preamble begins here.
\title{An Example Document}  % Declares the document's title.
\author{Leslie Lamport}      % Declares the author's name.
\date{January 21, 1994}      % Deleting this command produces today's date.

\newcommand{\ip}[2]{(#1, #2)}
                             % Defines \ip{arg1}{arg2} to mean
                             % (arg1, arg2).

%\newcommand{\ip}[2]{\langle #1 | #2\rangle}
                             % This is an alternative definition of
                             % \ip that is commented out.

\begin{document}             % End of preamble and beginning of text.

\maketitle                   % Produces the title.

This is an example input file.  Comparing it with
the output it generates can show you how to
produce a simple document of your own.

\section{Ordinary Text}      % Produces section heading.  Lower-level
                             % sections are begun with similar
                             % \subsection and \subsubsection commands.

The ends  of words and sentences are marked
  by   spaces. It  doesn't matter how many
spaces    you type; one is as good as 100.  The
end of   a line counts as a space.

One   or more   blank lines denote the  end
of  a paragraph.

Since any number of consecutive spaces are treated
like a single one, the formatting of the input
file makes no difference to
      \LaTeX,                % The \LaTeX command generates the LaTeX logo.
but it makes a difference to you.  When you use
\LaTeX, making your input file as easy to read
as possible will be a great help as you write
your document and when you change it.  This sample
file shows how you can add comments to your own input
file.

Because printing is different from typewriting,
there are a number of things that you have to do
differently when preparing an input file than if
you were just typing the document directly.
Quotation marks like
       ``this''
have to be handled specially, as do quotes within
quotes:
       ``\,`this'            % \, separates the double and single quote.
        is what I just
        wrote, not  `that'\,''.

Dashes come in three sizes: an
       intra-word
dash, a medium dash for number ranges like
       1--2,
and a punctuation
       dash---like
this.

A sentence-ending space should be larger than the
space between words within a sentence.  You
sometimes have to type special commands in
conjunction with punctuation characters to get
this right, as in the following sentence.
       Gnats, gnus, etc.\ all  % `\ ' makes an inter-word space.
       begin with G\@.         % \@ marks end-of-sentence punctuation.
You should check the spaces after periods when
reading your output to make sure you haven't
forgotten any special cases.  Generating an
ellipsis
       \ldots\               % `\ ' is needed after `\ldots' because TeX
                             % ignores spaces after command names like \ldots
                             % made from \ + letters.
                             %
                             % Note how a `%' character causes TeX to ignore
                             % the end of the input line, so these blank lines
                             % do not start a new paragraph.
                             %
with the right spacing around the periods requires
a special command.

\LaTeX\ interprets some common characters as
commands, so you must type special commands to
generate them.  These characters include the
following:
       \$ \& \% \# \{ and \}.

In printing, text is usually emphasized with an
       \emph{italic}
type style.

\begin{em}
   A long segment of text can also be emphasized
   in this way.  Text within such a segment can be
   given \emph{additional} emphasis.
\end{em}

It is sometimes necessary to prevent \LaTeX\ from
breaking a line where it might otherwise do so.
This may be at a space, as between the ``Mr.''\ and
``Jones'' in
       ``Mr.~Jones'',        % ~ produces an unbreakable interword space.
or within a word---especially when the word is a
symbol like
       \mbox{\emph{itemnum}}
that makes little sense when hyphenated across
lines.

Footnotes\footnote{This is an example of a footnote.}
pose no problem.

\LaTeX\ is good at typesetting mathematical formulas
like
       \( x-3y + z = 7 \)
or
       \( a_{1} > x^{2n} + y^{2n} > x' \)
or
       \( \ip{A}{B} = \sum_{i} a_{i} b_{i} \).
The spaces you type in a formula are
ignored.  Remember that a letter like
       $x$                   % $ ... $  and  \( ... \)  are equivalent
is a formula when it denotes a mathematical
symbol, and it should be typed as one.

\section{Displayed Text}

Text is displayed by indenting it from the left
margin.  Quotations are commonly displayed.  There
are short quotations
\begin{quote}
   This is a short quotation.  It consists of a
   single paragraph of text.  See how it is formatted.
\end{quote}
and longer ones.
\begin{quotation}
   This is a longer quotation.  It consists of two
   paragraphs of text, neither of which are
   particularly interesting.

   This is the second paragraph of the quotation.  It
   is just as dull as the first paragraph.
\end{quotation}
Another frequently-displayed structure is a list.
The following is an example of an \emph{itemized}
list.
\begin{itemize}
   \item This is the first item of an itemized list.
         Each item in the list is marked with a ``tick''.
         You don't have to worry about what kind of tick
         mark is used.

   \item This is the second item of the list.  It
         contains another list nested inside it.  The inner
         list is an \emph{enumerated} list.
         \begin{enumerate}
            \item This is the first item of an enumerated
                  list that is nested within the itemized list.

            \item This is the second item of the inner list.
                  \LaTeX\ allows you to nest lists deeper than
                  you really should.
         \end{enumerate}
         This is the rest of the second item of the outer
         list.  It is no more interesting than any other
         part of the item.
   \item This is the third item of the list.
\end{itemize}
You can even display poetry.
\begin{verse}
   There is an environment
    for verse \\             % The \\ command separates lines
   Whose features some poets % within a stanza.
   will curse.

                             % One or more blank lines separate stanzas.

   For instead of making\\
   Them do \emph{all} line breaking, \\
   It allows them to put too many words on a line when they'd rather be
   forced to be terse.
\end{verse}

Mathematical formulas may also be displayed.  A
displayed formula
is
one-line long; multiline
formulas require special formatting instructions.
   \[  \ip{\Gamma}{\psi'} = x'' + y^{2} + z_{i}^{n}\]
Don't start a paragraph with a displayed equation,
nor make one a paragraph by itself.

\end{document}               % End of document.

\section{Bench Huge Section 565}
% synthetic bench section 565
% This is a sample LaTeX input file.  (Version of 12 August 2004.)
%
% A '%' character causes TeX to ignore all remaining text on the line,
% and is used for comments like this one.

\documentclass{article}      % Specifies the document class

                             % The preamble begins here.
\title{An Example Document}  % Declares the document's title.
\author{Leslie Lamport}      % Declares the author's name.
\date{January 21, 1994}      % Deleting this command produces today's date.

\newcommand{\ip}[2]{(#1, #2)}
                             % Defines \ip{arg1}{arg2} to mean
                             % (arg1, arg2).

%\newcommand{\ip}[2]{\langle #1 | #2\rangle}
                             % This is an alternative definition of
                             % \ip that is commented out.

\begin{document}             % End of preamble and beginning of text.

\maketitle                   % Produces the title.

This is an example input file.  Comparing it with
the output it generates can show you how to
produce a simple document of your own.

\section{Ordinary Text}      % Produces section heading.  Lower-level
                             % sections are begun with similar
                             % \subsection and \subsubsection commands.

The ends  of words and sentences are marked
  by   spaces. It  doesn't matter how many
spaces    you type; one is as good as 100.  The
end of   a line counts as a space.

One   or more   blank lines denote the  end
of  a paragraph.

Since any number of consecutive spaces are treated
like a single one, the formatting of the input
file makes no difference to
      \LaTeX,                % The \LaTeX command generates the LaTeX logo.
but it makes a difference to you.  When you use
\LaTeX, making your input file as easy to read
as possible will be a great help as you write
your document and when you change it.  This sample
file shows how you can add comments to your own input
file.

Because printing is different from typewriting,
there are a number of things that you have to do
differently when preparing an input file than if
you were just typing the document directly.
Quotation marks like
       ``this''
have to be handled specially, as do quotes within
quotes:
       ``\,`this'            % \, separates the double and single quote.
        is what I just
        wrote, not  `that'\,''.

Dashes come in three sizes: an
       intra-word
dash, a medium dash for number ranges like
       1--2,
and a punctuation
       dash---like
this.

A sentence-ending space should be larger than the
space between words within a sentence.  You
sometimes have to type special commands in
conjunction with punctuation characters to get
this right, as in the following sentence.
       Gnats, gnus, etc.\ all  % `\ ' makes an inter-word space.
       begin with G\@.         % \@ marks end-of-sentence punctuation.
You should check the spaces after periods when
reading your output to make sure you haven't
forgotten any special cases.  Generating an
ellipsis
       \ldots\               % `\ ' is needed after `\ldots' because TeX
                             % ignores spaces after command names like \ldots
                             % made from \ + letters.
                             %
                             % Note how a `%' character causes TeX to ignore
                             % the end of the input line, so these blank lines
                             % do not start a new paragraph.
                             %
with the right spacing around the periods requires
a special command.

\LaTeX\ interprets some common characters as
commands, so you must type special commands to
generate them.  These characters include the
following:
       \$ \& \% \# \{ and \}.

In printing, text is usually emphasized with an
       \emph{italic}
type style.

\begin{em}
   A long segment of text can also be emphasized
   in this way.  Text within such a segment can be
   given \emph{additional} emphasis.
\end{em}

It is sometimes necessary to prevent \LaTeX\ from
breaking a line where it might otherwise do so.
This may be at a space, as between the ``Mr.''\ and
``Jones'' in
       ``Mr.~Jones'',        % ~ produces an unbreakable interword space.
or within a word---especially when the word is a
symbol like
       \mbox{\emph{itemnum}}
that makes little sense when hyphenated across
lines.

Footnotes\footnote{This is an example of a footnote.}
pose no problem.

\LaTeX\ is good at typesetting mathematical formulas
like
       \( x-3y + z = 7 \)
or
       \( a_{1} > x^{2n} + y^{2n} > x' \)
or
       \( \ip{A}{B} = \sum_{i} a_{i} b_{i} \).
The spaces you type in a formula are
ignored.  Remember that a letter like
       $x$                   % $ ... $  and  \( ... \)  are equivalent
is a formula when it denotes a mathematical
symbol, and it should be typed as one.

\section{Displayed Text}

Text is displayed by indenting it from the left
margin.  Quotations are commonly displayed.  There
are short quotations
\begin{quote}
   This is a short quotation.  It consists of a
   single paragraph of text.  See how it is formatted.
\end{quote}
and longer ones.
\begin{quotation}
   This is a longer quotation.  It consists of two
   paragraphs of text, neither of which are
   particularly interesting.

   This is the second paragraph of the quotation.  It
   is just as dull as the first paragraph.
\end{quotation}
Another frequently-displayed structure is a list.
The following is an example of an \emph{itemized}
list.
\begin{itemize}
   \item This is the first item of an itemized list.
         Each item in the list is marked with a ``tick''.
         You don't have to worry about what kind of tick
         mark is used.

   \item This is the second item of the list.  It
         contains another list nested inside it.  The inner
         list is an \emph{enumerated} list.
         \begin{enumerate}
            \item This is the first item of an enumerated
                  list that is nested within the itemized list.

            \item This is the second item of the inner list.
                  \LaTeX\ allows you to nest lists deeper than
                  you really should.
         \end{enumerate}
         This is the rest of the second item of the outer
         list.  It is no more interesting than any other
         part of the item.
   \item This is the third item of the list.
\end{itemize}
You can even display poetry.
\begin{verse}
   There is an environment
    for verse \\             % The \\ command separates lines
   Whose features some poets % within a stanza.
   will curse.

                             % One or more blank lines separate stanzas.

   For instead of making\\
   Them do \emph{all} line breaking, \\
   It allows them to put too many words on a line when they'd rather be
   forced to be terse.
\end{verse}

Mathematical formulas may also be displayed.  A
displayed formula
is
one-line long; multiline
formulas require special formatting instructions.
   \[  \ip{\Gamma}{\psi'} = x'' + y^{2} + z_{i}^{n}\]
Don't start a paragraph with a displayed equation,
nor make one a paragraph by itself.

\end{document}               % End of document.

\section{Bench Huge Section 566}
% synthetic bench section 566
% This is a sample LaTeX input file.  (Version of 12 August 2004.)
%
% A '%' character causes TeX to ignore all remaining text on the line,
% and is used for comments like this one.

\documentclass{article}      % Specifies the document class

                             % The preamble begins here.
\title{An Example Document}  % Declares the document's title.
\author{Leslie Lamport}      % Declares the author's name.
\date{January 21, 1994}      % Deleting this command produces today's date.

\newcommand{\ip}[2]{(#1, #2)}
                             % Defines \ip{arg1}{arg2} to mean
                             % (arg1, arg2).

%\newcommand{\ip}[2]{\langle #1 | #2\rangle}
                             % This is an alternative definition of
                             % \ip that is commented out.

\begin{document}             % End of preamble and beginning of text.

\maketitle                   % Produces the title.

This is an example input file.  Comparing it with
the output it generates can show you how to
produce a simple document of your own.

\section{Ordinary Text}      % Produces section heading.  Lower-level
                             % sections are begun with similar
                             % \subsection and \subsubsection commands.

The ends  of words and sentences are marked
  by   spaces. It  doesn't matter how many
spaces    you type; one is as good as 100.  The
end of   a line counts as a space.

One   or more   blank lines denote the  end
of  a paragraph.

Since any number of consecutive spaces are treated
like a single one, the formatting of the input
file makes no difference to
      \LaTeX,                % The \LaTeX command generates the LaTeX logo.
but it makes a difference to you.  When you use
\LaTeX, making your input file as easy to read
as possible will be a great help as you write
your document and when you change it.  This sample
file shows how you can add comments to your own input
file.

Because printing is different from typewriting,
there are a number of things that you have to do
differently when preparing an input file than if
you were just typing the document directly.
Quotation marks like
       ``this''
have to be handled specially, as do quotes within
quotes:
       ``\,`this'            % \, separates the double and single quote.
        is what I just
        wrote, not  `that'\,''.

Dashes come in three sizes: an
       intra-word
dash, a medium dash for number ranges like
       1--2,
and a punctuation
       dash---like
this.

A sentence-ending space should be larger than the
space between words within a sentence.  You
sometimes have to type special commands in
conjunction with punctuation characters to get
this right, as in the following sentence.
       Gnats, gnus, etc.\ all  % `\ ' makes an inter-word space.
       begin with G\@.         % \@ marks end-of-sentence punctuation.
You should check the spaces after periods when
reading your output to make sure you haven't
forgotten any special cases.  Generating an
ellipsis
       \ldots\               % `\ ' is needed after `\ldots' because TeX
                             % ignores spaces after command names like \ldots
                             % made from \ + letters.
                             %
                             % Note how a `%' character causes TeX to ignore
                             % the end of the input line, so these blank lines
                             % do not start a new paragraph.
                             %
with the right spacing around the periods requires
a special command.

\LaTeX\ interprets some common characters as
commands, so you must type special commands to
generate them.  These characters include the
following:
       \$ \& \% \# \{ and \}.

In printing, text is usually emphasized with an
       \emph{italic}
type style.

\begin{em}
   A long segment of text can also be emphasized
   in this way.  Text within such a segment can be
   given \emph{additional} emphasis.
\end{em}

It is sometimes necessary to prevent \LaTeX\ from
breaking a line where it might otherwise do so.
This may be at a space, as between the ``Mr.''\ and
``Jones'' in
       ``Mr.~Jones'',        % ~ produces an unbreakable interword space.
or within a word---especially when the word is a
symbol like
       \mbox{\emph{itemnum}}
that makes little sense when hyphenated across
lines.

Footnotes\footnote{This is an example of a footnote.}
pose no problem.

\LaTeX\ is good at typesetting mathematical formulas
like
       \( x-3y + z = 7 \)
or
       \( a_{1} > x^{2n} + y^{2n} > x' \)
or
       \( \ip{A}{B} = \sum_{i} a_{i} b_{i} \).
The spaces you type in a formula are
ignored.  Remember that a letter like
       $x$                   % $ ... $  and  \( ... \)  are equivalent
is a formula when it denotes a mathematical
symbol, and it should be typed as one.

\section{Displayed Text}

Text is displayed by indenting it from the left
margin.  Quotations are commonly displayed.  There
are short quotations
\begin{quote}
   This is a short quotation.  It consists of a
   single paragraph of text.  See how it is formatted.
\end{quote}
and longer ones.
\begin{quotation}
   This is a longer quotation.  It consists of two
   paragraphs of text, neither of which are
   particularly interesting.

   This is the second paragraph of the quotation.  It
   is just as dull as the first paragraph.
\end{quotation}
Another frequently-displayed structure is a list.
The following is an example of an \emph{itemized}
list.
\begin{itemize}
   \item This is the first item of an itemized list.
         Each item in the list is marked with a ``tick''.
         You don't have to worry about what kind of tick
         mark is used.

   \item This is the second item of the list.  It
         contains another list nested inside it.  The inner
         list is an \emph{enumerated} list.
         \begin{enumerate}
            \item This is the first item of an enumerated
                  list that is nested within the itemized list.

            \item This is the second item of the inner list.
                  \LaTeX\ allows you to nest lists deeper than
                  you really should.
         \end{enumerate}
         This is the rest of the second item of the outer
         list.  It is no more interesting than any other
         part of the item.
   \item This is the third item of the list.
\end{itemize}
You can even display poetry.
\begin{verse}
   There is an environment
    for verse \\             % The \\ command separates lines
   Whose features some poets % within a stanza.
   will curse.

                             % One or more blank lines separate stanzas.

   For instead of making\\
   Them do \emph{all} line breaking, \\
   It allows them to put too many words on a line when they'd rather be
   forced to be terse.
\end{verse}

Mathematical formulas may also be displayed.  A
displayed formula
is
one-line long; multiline
formulas require special formatting instructions.
   \[  \ip{\Gamma}{\psi'} = x'' + y^{2} + z_{i}^{n}\]
Don't start a paragraph with a displayed equation,
nor make one a paragraph by itself.

\end{document}               % End of document.

\section{Bench Huge Section 567}
% synthetic bench section 567
% This is a sample LaTeX input file.  (Version of 12 August 2004.)
%
% A '%' character causes TeX to ignore all remaining text on the line,
% and is used for comments like this one.

\documentclass{article}      % Specifies the document class

                             % The preamble begins here.
\title{An Example Document}  % Declares the document's title.
\author{Leslie Lamport}      % Declares the author's name.
\date{January 21, 1994}      % Deleting this command produces today's date.

\newcommand{\ip}[2]{(#1, #2)}
                             % Defines \ip{arg1}{arg2} to mean
                             % (arg1, arg2).

%\newcommand{\ip}[2]{\langle #1 | #2\rangle}
                             % This is an alternative definition of
                             % \ip that is commented out.

\begin{document}             % End of preamble and beginning of text.

\maketitle                   % Produces the title.

This is an example input file.  Comparing it with
the output it generates can show you how to
produce a simple document of your own.

\section{Ordinary Text}      % Produces section heading.  Lower-level
                             % sections are begun with similar
                             % \subsection and \subsubsection commands.

The ends  of words and sentences are marked
  by   spaces. It  doesn't matter how many
spaces    you type; one is as good as 100.  The
end of   a line counts as a space.

One   or more   blank lines denote the  end
of  a paragraph.

Since any number of consecutive spaces are treated
like a single one, the formatting of the input
file makes no difference to
      \LaTeX,                % The \LaTeX command generates the LaTeX logo.
but it makes a difference to you.  When you use
\LaTeX, making your input file as easy to read
as possible will be a great help as you write
your document and when you change it.  This sample
file shows how you can add comments to your own input
file.

Because printing is different from typewriting,
there are a number of things that you have to do
differently when preparing an input file than if
you were just typing the document directly.
Quotation marks like
       ``this''
have to be handled specially, as do quotes within
quotes:
       ``\,`this'            % \, separates the double and single quote.
        is what I just
        wrote, not  `that'\,''.

Dashes come in three sizes: an
       intra-word
dash, a medium dash for number ranges like
       1--2,
and a punctuation
       dash---like
this.

A sentence-ending space should be larger than the
space between words within a sentence.  You
sometimes have to type special commands in
conjunction with punctuation characters to get
this right, as in the following sentence.
       Gnats, gnus, etc.\ all  % `\ ' makes an inter-word space.
       begin with G\@.         % \@ marks end-of-sentence punctuation.
You should check the spaces after periods when
reading your output to make sure you haven't
forgotten any special cases.  Generating an
ellipsis
       \ldots\               % `\ ' is needed after `\ldots' because TeX
                             % ignores spaces after command names like \ldots
                             % made from \ + letters.
                             %
                             % Note how a `%' character causes TeX to ignore
                             % the end of the input line, so these blank lines
                             % do not start a new paragraph.
                             %
with the right spacing around the periods requires
a special command.

\LaTeX\ interprets some common characters as
commands, so you must type special commands to
generate them.  These characters include the
following:
       \$ \& \% \# \{ and \}.

In printing, text is usually emphasized with an
       \emph{italic}
type style.

\begin{em}
   A long segment of text can also be emphasized
   in this way.  Text within such a segment can be
   given \emph{additional} emphasis.
\end{em}

It is sometimes necessary to prevent \LaTeX\ from
breaking a line where it might otherwise do so.
This may be at a space, as between the ``Mr.''\ and
``Jones'' in
       ``Mr.~Jones'',        % ~ produces an unbreakable interword space.
or within a word---especially when the word is a
symbol like
       \mbox{\emph{itemnum}}
that makes little sense when hyphenated across
lines.

Footnotes\footnote{This is an example of a footnote.}
pose no problem.

\LaTeX\ is good at typesetting mathematical formulas
like
       \( x-3y + z = 7 \)
or
       \( a_{1} > x^{2n} + y^{2n} > x' \)
or
       \( \ip{A}{B} = \sum_{i} a_{i} b_{i} \).
The spaces you type in a formula are
ignored.  Remember that a letter like
       $x$                   % $ ... $  and  \( ... \)  are equivalent
is a formula when it denotes a mathematical
symbol, and it should be typed as one.

\section{Displayed Text}

Text is displayed by indenting it from the left
margin.  Quotations are commonly displayed.  There
are short quotations
\begin{quote}
   This is a short quotation.  It consists of a
   single paragraph of text.  See how it is formatted.
\end{quote}
and longer ones.
\begin{quotation}
   This is a longer quotation.  It consists of two
   paragraphs of text, neither of which are
   particularly interesting.

   This is the second paragraph of the quotation.  It
   is just as dull as the first paragraph.
\end{quotation}
Another frequently-displayed structure is a list.
The following is an example of an \emph{itemized}
list.
\begin{itemize}
   \item This is the first item of an itemized list.
         Each item in the list is marked with a ``tick''.
         You don't have to worry about what kind of tick
         mark is used.

   \item This is the second item of the list.  It
         contains another list nested inside it.  The inner
         list is an \emph{enumerated} list.
         \begin{enumerate}
            \item This is the first item of an enumerated
                  list that is nested within the itemized list.

            \item This is the second item of the inner list.
                  \LaTeX\ allows you to nest lists deeper than
                  you really should.
         \end{enumerate}
         This is the rest of the second item of the outer
         list.  It is no more interesting than any other
         part of the item.
   \item This is the third item of the list.
\end{itemize}
You can even display poetry.
\begin{verse}
   There is an environment
    for verse \\             % The \\ command separates lines
   Whose features some poets % within a stanza.
   will curse.

                             % One or more blank lines separate stanzas.

   For instead of making\\
   Them do \emph{all} line breaking, \\
   It allows them to put too many words on a line when they'd rather be
   forced to be terse.
\end{verse}

Mathematical formulas may also be displayed.  A
displayed formula
is
one-line long; multiline
formulas require special formatting instructions.
   \[  \ip{\Gamma}{\psi'} = x'' + y^{2} + z_{i}^{n}\]
Don't start a paragraph with a displayed equation,
nor make one a paragraph by itself.

\end{document}               % End of document.

\section{Bench Huge Section 568}
% synthetic bench section 568
% This is a sample LaTeX input file.  (Version of 12 August 2004.)
%
% A '%' character causes TeX to ignore all remaining text on the line,
% and is used for comments like this one.

\documentclass{article}      % Specifies the document class

                             % The preamble begins here.
\title{An Example Document}  % Declares the document's title.
\author{Leslie Lamport}      % Declares the author's name.
\date{January 21, 1994}      % Deleting this command produces today's date.

\newcommand{\ip}[2]{(#1, #2)}
                             % Defines \ip{arg1}{arg2} to mean
                             % (arg1, arg2).

%\newcommand{\ip}[2]{\langle #1 | #2\rangle}
                             % This is an alternative definition of
                             % \ip that is commented out.

\begin{document}             % End of preamble and beginning of text.

\maketitle                   % Produces the title.

This is an example input file.  Comparing it with
the output it generates can show you how to
produce a simple document of your own.

\section{Ordinary Text}      % Produces section heading.  Lower-level
                             % sections are begun with similar
                             % \subsection and \subsubsection commands.

The ends  of words and sentences are marked
  by   spaces. It  doesn't matter how many
spaces    you type; one is as good as 100.  The
end of   a line counts as a space.

One   or more   blank lines denote the  end
of  a paragraph.

Since any number of consecutive spaces are treated
like a single one, the formatting of the input
file makes no difference to
      \LaTeX,                % The \LaTeX command generates the LaTeX logo.
but it makes a difference to you.  When you use
\LaTeX, making your input file as easy to read
as possible will be a great help as you write
your document and when you change it.  This sample
file shows how you can add comments to your own input
file.

Because printing is different from typewriting,
there are a number of things that you have to do
differently when preparing an input file than if
you were just typing the document directly.
Quotation marks like
       ``this''
have to be handled specially, as do quotes within
quotes:
       ``\,`this'            % \, separates the double and single quote.
        is what I just
        wrote, not  `that'\,''.

Dashes come in three sizes: an
       intra-word
dash, a medium dash for number ranges like
       1--2,
and a punctuation
       dash---like
this.

A sentence-ending space should be larger than the
space between words within a sentence.  You
sometimes have to type special commands in
conjunction with punctuation characters to get
this right, as in the following sentence.
       Gnats, gnus, etc.\ all  % `\ ' makes an inter-word space.
       begin with G\@.         % \@ marks end-of-sentence punctuation.
You should check the spaces after periods when
reading your output to make sure you haven't
forgotten any special cases.  Generating an
ellipsis
       \ldots\               % `\ ' is needed after `\ldots' because TeX
                             % ignores spaces after command names like \ldots
                             % made from \ + letters.
                             %
                             % Note how a `%' character causes TeX to ignore
                             % the end of the input line, so these blank lines
                             % do not start a new paragraph.
                             %
with the right spacing around the periods requires
a special command.

\LaTeX\ interprets some common characters as
commands, so you must type special commands to
generate them.  These characters include the
following:
       \$ \& \% \# \{ and \}.

In printing, text is usually emphasized with an
       \emph{italic}
type style.

\begin{em}
   A long segment of text can also be emphasized
   in this way.  Text within such a segment can be
   given \emph{additional} emphasis.
\end{em}

It is sometimes necessary to prevent \LaTeX\ from
breaking a line where it might otherwise do so.
This may be at a space, as between the ``Mr.''\ and
``Jones'' in
       ``Mr.~Jones'',        % ~ produces an unbreakable interword space.
or within a word---especially when the word is a
symbol like
       \mbox{\emph{itemnum}}
that makes little sense when hyphenated across
lines.

Footnotes\footnote{This is an example of a footnote.}
pose no problem.

\LaTeX\ is good at typesetting mathematical formulas
like
       \( x-3y + z = 7 \)
or
       \( a_{1} > x^{2n} + y^{2n} > x' \)
or
       \( \ip{A}{B} = \sum_{i} a_{i} b_{i} \).
The spaces you type in a formula are
ignored.  Remember that a letter like
       $x$                   % $ ... $  and  \( ... \)  are equivalent
is a formula when it denotes a mathematical
symbol, and it should be typed as one.

\section{Displayed Text}

Text is displayed by indenting it from the left
margin.  Quotations are commonly displayed.  There
are short quotations
\begin{quote}
   This is a short quotation.  It consists of a
   single paragraph of text.  See how it is formatted.
\end{quote}
and longer ones.
\begin{quotation}
   This is a longer quotation.  It consists of two
   paragraphs of text, neither of which are
   particularly interesting.

   This is the second paragraph of the quotation.  It
   is just as dull as the first paragraph.
\end{quotation}
Another frequently-displayed structure is a list.
The following is an example of an \emph{itemized}
list.
\begin{itemize}
   \item This is the first item of an itemized list.
         Each item in the list is marked with a ``tick''.
         You don't have to worry about what kind of tick
         mark is used.

   \item This is the second item of the list.  It
         contains another list nested inside it.  The inner
         list is an \emph{enumerated} list.
         \begin{enumerate}
            \item This is the first item of an enumerated
                  list that is nested within the itemized list.

            \item This is the second item of the inner list.
                  \LaTeX\ allows you to nest lists deeper than
                  you really should.
         \end{enumerate}
         This is the rest of the second item of the outer
         list.  It is no more interesting than any other
         part of the item.
   \item This is the third item of the list.
\end{itemize}
You can even display poetry.
\begin{verse}
   There is an environment
    for verse \\             % The \\ command separates lines
   Whose features some poets % within a stanza.
   will curse.

                             % One or more blank lines separate stanzas.

   For instead of making\\
   Them do \emph{all} line breaking, \\
   It allows them to put too many words on a line when they'd rather be
   forced to be terse.
\end{verse}

Mathematical formulas may also be displayed.  A
displayed formula
is
one-line long; multiline
formulas require special formatting instructions.
   \[  \ip{\Gamma}{\psi'} = x'' + y^{2} + z_{i}^{n}\]
Don't start a paragraph with a displayed equation,
nor make one a paragraph by itself.

\end{document}               % End of document.

\section{Bench Huge Section 569}
% synthetic bench section 569
% This is a sample LaTeX input file.  (Version of 12 August 2004.)
%
% A '%' character causes TeX to ignore all remaining text on the line,
% and is used for comments like this one.

\documentclass{article}      % Specifies the document class

                             % The preamble begins here.
\title{An Example Document}  % Declares the document's title.
\author{Leslie Lamport}      % Declares the author's name.
\date{January 21, 1994}      % Deleting this command produces today's date.

\newcommand{\ip}[2]{(#1, #2)}
                             % Defines \ip{arg1}{arg2} to mean
                             % (arg1, arg2).

%\newcommand{\ip}[2]{\langle #1 | #2\rangle}
                             % This is an alternative definition of
                             % \ip that is commented out.

\begin{document}             % End of preamble and beginning of text.

\maketitle                   % Produces the title.

This is an example input file.  Comparing it with
the output it generates can show you how to
produce a simple document of your own.

\section{Ordinary Text}      % Produces section heading.  Lower-level
                             % sections are begun with similar
                             % \subsection and \subsubsection commands.

The ends  of words and sentences are marked
  by   spaces. It  doesn't matter how many
spaces    you type; one is as good as 100.  The
end of   a line counts as a space.

One   or more   blank lines denote the  end
of  a paragraph.

Since any number of consecutive spaces are treated
like a single one, the formatting of the input
file makes no difference to
      \LaTeX,                % The \LaTeX command generates the LaTeX logo.
but it makes a difference to you.  When you use
\LaTeX, making your input file as easy to read
as possible will be a great help as you write
your document and when you change it.  This sample
file shows how you can add comments to your own input
file.

Because printing is different from typewriting,
there are a number of things that you have to do
differently when preparing an input file than if
you were just typing the document directly.
Quotation marks like
       ``this''
have to be handled specially, as do quotes within
quotes:
       ``\,`this'            % \, separates the double and single quote.
        is what I just
        wrote, not  `that'\,''.

Dashes come in three sizes: an
       intra-word
dash, a medium dash for number ranges like
       1--2,
and a punctuation
       dash---like
this.

A sentence-ending space should be larger than the
space between words within a sentence.  You
sometimes have to type special commands in
conjunction with punctuation characters to get
this right, as in the following sentence.
       Gnats, gnus, etc.\ all  % `\ ' makes an inter-word space.
       begin with G\@.         % \@ marks end-of-sentence punctuation.
You should check the spaces after periods when
reading your output to make sure you haven't
forgotten any special cases.  Generating an
ellipsis
       \ldots\               % `\ ' is needed after `\ldots' because TeX
                             % ignores spaces after command names like \ldots
                             % made from \ + letters.
                             %
                             % Note how a `%' character causes TeX to ignore
                             % the end of the input line, so these blank lines
                             % do not start a new paragraph.
                             %
with the right spacing around the periods requires
a special command.

\LaTeX\ interprets some common characters as
commands, so you must type special commands to
generate them.  These characters include the
following:
       \$ \& \% \# \{ and \}.

In printing, text is usually emphasized with an
       \emph{italic}
type style.

\begin{em}
   A long segment of text can also be emphasized
   in this way.  Text within such a segment can be
   given \emph{additional} emphasis.
\end{em}

It is sometimes necessary to prevent \LaTeX\ from
breaking a line where it might otherwise do so.
This may be at a space, as between the ``Mr.''\ and
``Jones'' in
       ``Mr.~Jones'',        % ~ produces an unbreakable interword space.
or within a word---especially when the word is a
symbol like
       \mbox{\emph{itemnum}}
that makes little sense when hyphenated across
lines.

Footnotes\footnote{This is an example of a footnote.}
pose no problem.

\LaTeX\ is good at typesetting mathematical formulas
like
       \( x-3y + z = 7 \)
or
       \( a_{1} > x^{2n} + y^{2n} > x' \)
or
       \( \ip{A}{B} = \sum_{i} a_{i} b_{i} \).
The spaces you type in a formula are
ignored.  Remember that a letter like
       $x$                   % $ ... $  and  \( ... \)  are equivalent
is a formula when it denotes a mathematical
symbol, and it should be typed as one.

\section{Displayed Text}

Text is displayed by indenting it from the left
margin.  Quotations are commonly displayed.  There
are short quotations
\begin{quote}
   This is a short quotation.  It consists of a
   single paragraph of text.  See how it is formatted.
\end{quote}
and longer ones.
\begin{quotation}
   This is a longer quotation.  It consists of two
   paragraphs of text, neither of which are
   particularly interesting.

   This is the second paragraph of the quotation.  It
   is just as dull as the first paragraph.
\end{quotation}
Another frequently-displayed structure is a list.
The following is an example of an \emph{itemized}
list.
\begin{itemize}
   \item This is the first item of an itemized list.
         Each item in the list is marked with a ``tick''.
         You don't have to worry about what kind of tick
         mark is used.

   \item This is the second item of the list.  It
         contains another list nested inside it.  The inner
         list is an \emph{enumerated} list.
         \begin{enumerate}
            \item This is the first item of an enumerated
                  list that is nested within the itemized list.

            \item This is the second item of the inner list.
                  \LaTeX\ allows you to nest lists deeper than
                  you really should.
         \end{enumerate}
         This is the rest of the second item of the outer
         list.  It is no more interesting than any other
         part of the item.
   \item This is the third item of the list.
\end{itemize}
You can even display poetry.
\begin{verse}
   There is an environment
    for verse \\             % The \\ command separates lines
   Whose features some poets % within a stanza.
   will curse.

                             % One or more blank lines separate stanzas.

   For instead of making\\
   Them do \emph{all} line breaking, \\
   It allows them to put too many words on a line when they'd rather be
   forced to be terse.
\end{verse}

Mathematical formulas may also be displayed.  A
displayed formula
is
one-line long; multiline
formulas require special formatting instructions.
   \[  \ip{\Gamma}{\psi'} = x'' + y^{2} + z_{i}^{n}\]
Don't start a paragraph with a displayed equation,
nor make one a paragraph by itself.

\end{document}               % End of document.

\section{Bench Huge Section 570}
% synthetic bench section 570
% This is a sample LaTeX input file.  (Version of 12 August 2004.)
%
% A '%' character causes TeX to ignore all remaining text on the line,
% and is used for comments like this one.

\documentclass{article}      % Specifies the document class

                             % The preamble begins here.
\title{An Example Document}  % Declares the document's title.
\author{Leslie Lamport}      % Declares the author's name.
\date{January 21, 1994}      % Deleting this command produces today's date.

\newcommand{\ip}[2]{(#1, #2)}
                             % Defines \ip{arg1}{arg2} to mean
                             % (arg1, arg2).

%\newcommand{\ip}[2]{\langle #1 | #2\rangle}
                             % This is an alternative definition of
                             % \ip that is commented out.

\begin{document}             % End of preamble and beginning of text.

\maketitle                   % Produces the title.

This is an example input file.  Comparing it with
the output it generates can show you how to
produce a simple document of your own.

\section{Ordinary Text}      % Produces section heading.  Lower-level
                             % sections are begun with similar
                             % \subsection and \subsubsection commands.

The ends  of words and sentences are marked
  by   spaces. It  doesn't matter how many
spaces    you type; one is as good as 100.  The
end of   a line counts as a space.

One   or more   blank lines denote the  end
of  a paragraph.

Since any number of consecutive spaces are treated
like a single one, the formatting of the input
file makes no difference to
      \LaTeX,                % The \LaTeX command generates the LaTeX logo.
but it makes a difference to you.  When you use
\LaTeX, making your input file as easy to read
as possible will be a great help as you write
your document and when you change it.  This sample
file shows how you can add comments to your own input
file.

Because printing is different from typewriting,
there are a number of things that you have to do
differently when preparing an input file than if
you were just typing the document directly.
Quotation marks like
       ``this''
have to be handled specially, as do quotes within
quotes:
       ``\,`this'            % \, separates the double and single quote.
        is what I just
        wrote, not  `that'\,''.

Dashes come in three sizes: an
       intra-word
dash, a medium dash for number ranges like
       1--2,
and a punctuation
       dash---like
this.

A sentence-ending space should be larger than the
space between words within a sentence.  You
sometimes have to type special commands in
conjunction with punctuation characters to get
this right, as in the following sentence.
       Gnats, gnus, etc.\ all  % `\ ' makes an inter-word space.
       begin with G\@.         % \@ marks end-of-sentence punctuation.
You should check the spaces after periods when
reading your output to make sure you haven't
forgotten any special cases.  Generating an
ellipsis
       \ldots\               % `\ ' is needed after `\ldots' because TeX
                             % ignores spaces after command names like \ldots
                             % made from \ + letters.
                             %
                             % Note how a `%' character causes TeX to ignore
                             % the end of the input line, so these blank lines
                             % do not start a new paragraph.
                             %
with the right spacing around the periods requires
a special command.

\LaTeX\ interprets some common characters as
commands, so you must type special commands to
generate them.  These characters include the
following:
       \$ \& \% \# \{ and \}.

In printing, text is usually emphasized with an
       \emph{italic}
type style.

\begin{em}
   A long segment of text can also be emphasized
   in this way.  Text within such a segment can be
   given \emph{additional} emphasis.
\end{em}

It is sometimes necessary to prevent \LaTeX\ from
breaking a line where it might otherwise do so.
This may be at a space, as between the ``Mr.''\ and
``Jones'' in
       ``Mr.~Jones'',        % ~ produces an unbreakable interword space.
or within a word---especially when the word is a
symbol like
       \mbox{\emph{itemnum}}
that makes little sense when hyphenated across
lines.

Footnotes\footnote{This is an example of a footnote.}
pose no problem.

\LaTeX\ is good at typesetting mathematical formulas
like
       \( x-3y + z = 7 \)
or
       \( a_{1} > x^{2n} + y^{2n} > x' \)
or
       \( \ip{A}{B} = \sum_{i} a_{i} b_{i} \).
The spaces you type in a formula are
ignored.  Remember that a letter like
       $x$                   % $ ... $  and  \( ... \)  are equivalent
is a formula when it denotes a mathematical
symbol, and it should be typed as one.

\section{Displayed Text}

Text is displayed by indenting it from the left
margin.  Quotations are commonly displayed.  There
are short quotations
\begin{quote}
   This is a short quotation.  It consists of a
   single paragraph of text.  See how it is formatted.
\end{quote}
and longer ones.
\begin{quotation}
   This is a longer quotation.  It consists of two
   paragraphs of text, neither of which are
   particularly interesting.

   This is the second paragraph of the quotation.  It
   is just as dull as the first paragraph.
\end{quotation}
Another frequently-displayed structure is a list.
The following is an example of an \emph{itemized}
list.
\begin{itemize}
   \item This is the first item of an itemized list.
         Each item in the list is marked with a ``tick''.
         You don't have to worry about what kind of tick
         mark is used.

   \item This is the second item of the list.  It
         contains another list nested inside it.  The inner
         list is an \emph{enumerated} list.
         \begin{enumerate}
            \item This is the first item of an enumerated
                  list that is nested within the itemized list.

            \item This is the second item of the inner list.
                  \LaTeX\ allows you to nest lists deeper than
                  you really should.
         \end{enumerate}
         This is the rest of the second item of the outer
         list.  It is no more interesting than any other
         part of the item.
   \item This is the third item of the list.
\end{itemize}
You can even display poetry.
\begin{verse}
   There is an environment
    for verse \\             % The \\ command separates lines
   Whose features some poets % within a stanza.
   will curse.

                             % One or more blank lines separate stanzas.

   For instead of making\\
   Them do \emph{all} line breaking, \\
   It allows them to put too many words on a line when they'd rather be
   forced to be terse.
\end{verse}

Mathematical formulas may also be displayed.  A
displayed formula
is
one-line long; multiline
formulas require special formatting instructions.
   \[  \ip{\Gamma}{\psi'} = x'' + y^{2} + z_{i}^{n}\]
Don't start a paragraph with a displayed equation,
nor make one a paragraph by itself.

\end{document}               % End of document.

\section{Bench Huge Section 571}
% synthetic bench section 571
% This is a sample LaTeX input file.  (Version of 12 August 2004.)
%
% A '%' character causes TeX to ignore all remaining text on the line,
% and is used for comments like this one.

\documentclass{article}      % Specifies the document class

                             % The preamble begins here.
\title{An Example Document}  % Declares the document's title.
\author{Leslie Lamport}      % Declares the author's name.
\date{January 21, 1994}      % Deleting this command produces today's date.

\newcommand{\ip}[2]{(#1, #2)}
                             % Defines \ip{arg1}{arg2} to mean
                             % (arg1, arg2).

%\newcommand{\ip}[2]{\langle #1 | #2\rangle}
                             % This is an alternative definition of
                             % \ip that is commented out.

\begin{document}             % End of preamble and beginning of text.

\maketitle                   % Produces the title.

This is an example input file.  Comparing it with
the output it generates can show you how to
produce a simple document of your own.

\section{Ordinary Text}      % Produces section heading.  Lower-level
                             % sections are begun with similar
                             % \subsection and \subsubsection commands.

The ends  of words and sentences are marked
  by   spaces. It  doesn't matter how many
spaces    you type; one is as good as 100.  The
end of   a line counts as a space.

One   or more   blank lines denote the  end
of  a paragraph.

Since any number of consecutive spaces are treated
like a single one, the formatting of the input
file makes no difference to
      \LaTeX,                % The \LaTeX command generates the LaTeX logo.
but it makes a difference to you.  When you use
\LaTeX, making your input file as easy to read
as possible will be a great help as you write
your document and when you change it.  This sample
file shows how you can add comments to your own input
file.

Because printing is different from typewriting,
there are a number of things that you have to do
differently when preparing an input file than if
you were just typing the document directly.
Quotation marks like
       ``this''
have to be handled specially, as do quotes within
quotes:
       ``\,`this'            % \, separates the double and single quote.
        is what I just
        wrote, not  `that'\,''.

Dashes come in three sizes: an
       intra-word
dash, a medium dash for number ranges like
       1--2,
and a punctuation
       dash---like
this.

A sentence-ending space should be larger than the
space between words within a sentence.  You
sometimes have to type special commands in
conjunction with punctuation characters to get
this right, as in the following sentence.
       Gnats, gnus, etc.\ all  % `\ ' makes an inter-word space.
       begin with G\@.         % \@ marks end-of-sentence punctuation.
You should check the spaces after periods when
reading your output to make sure you haven't
forgotten any special cases.  Generating an
ellipsis
       \ldots\               % `\ ' is needed after `\ldots' because TeX
                             % ignores spaces after command names like \ldots
                             % made from \ + letters.
                             %
                             % Note how a `%' character causes TeX to ignore
                             % the end of the input line, so these blank lines
                             % do not start a new paragraph.
                             %
with the right spacing around the periods requires
a special command.

\LaTeX\ interprets some common characters as
commands, so you must type special commands to
generate them.  These characters include the
following:
       \$ \& \% \# \{ and \}.

In printing, text is usually emphasized with an
       \emph{italic}
type style.

\begin{em}
   A long segment of text can also be emphasized
   in this way.  Text within such a segment can be
   given \emph{additional} emphasis.
\end{em}

It is sometimes necessary to prevent \LaTeX\ from
breaking a line where it might otherwise do so.
This may be at a space, as between the ``Mr.''\ and
``Jones'' in
       ``Mr.~Jones'',        % ~ produces an unbreakable interword space.
or within a word---especially when the word is a
symbol like
       \mbox{\emph{itemnum}}
that makes little sense when hyphenated across
lines.

Footnotes\footnote{This is an example of a footnote.}
pose no problem.

\LaTeX\ is good at typesetting mathematical formulas
like
       \( x-3y + z = 7 \)
or
       \( a_{1} > x^{2n} + y^{2n} > x' \)
or
       \( \ip{A}{B} = \sum_{i} a_{i} b_{i} \).
The spaces you type in a formula are
ignored.  Remember that a letter like
       $x$                   % $ ... $  and  \( ... \)  are equivalent
is a formula when it denotes a mathematical
symbol, and it should be typed as one.

\section{Displayed Text}

Text is displayed by indenting it from the left
margin.  Quotations are commonly displayed.  There
are short quotations
\begin{quote}
   This is a short quotation.  It consists of a
   single paragraph of text.  See how it is formatted.
\end{quote}
and longer ones.
\begin{quotation}
   This is a longer quotation.  It consists of two
   paragraphs of text, neither of which are
   particularly interesting.

   This is the second paragraph of the quotation.  It
   is just as dull as the first paragraph.
\end{quotation}
Another frequently-displayed structure is a list.
The following is an example of an \emph{itemized}
list.
\begin{itemize}
   \item This is the first item of an itemized list.
         Each item in the list is marked with a ``tick''.
         You don't have to worry about what kind of tick
         mark is used.

   \item This is the second item of the list.  It
         contains another list nested inside it.  The inner
         list is an \emph{enumerated} list.
         \begin{enumerate}
            \item This is the first item of an enumerated
                  list that is nested within the itemized list.

            \item This is the second item of the inner list.
                  \LaTeX\ allows you to nest lists deeper than
                  you really should.
         \end{enumerate}
         This is the rest of the second item of the outer
         list.  It is no more interesting than any other
         part of the item.
   \item This is the third item of the list.
\end{itemize}
You can even display poetry.
\begin{verse}
   There is an environment
    for verse \\             % The \\ command separates lines
   Whose features some poets % within a stanza.
   will curse.

                             % One or more blank lines separate stanzas.

   For instead of making\\
   Them do \emph{all} line breaking, \\
   It allows them to put too many words on a line when they'd rather be
   forced to be terse.
\end{verse}

Mathematical formulas may also be displayed.  A
displayed formula
is
one-line long; multiline
formulas require special formatting instructions.
   \[  \ip{\Gamma}{\psi'} = x'' + y^{2} + z_{i}^{n}\]
Don't start a paragraph with a displayed equation,
nor make one a paragraph by itself.

\end{document}               % End of document.

\section{Bench Huge Section 572}
% synthetic bench section 572
% This is a sample LaTeX input file.  (Version of 12 August 2004.)
%
% A '%' character causes TeX to ignore all remaining text on the line,
% and is used for comments like this one.

\documentclass{article}      % Specifies the document class

                             % The preamble begins here.
\title{An Example Document}  % Declares the document's title.
\author{Leslie Lamport}      % Declares the author's name.
\date{January 21, 1994}      % Deleting this command produces today's date.

\newcommand{\ip}[2]{(#1, #2)}
                             % Defines \ip{arg1}{arg2} to mean
                             % (arg1, arg2).

%\newcommand{\ip}[2]{\langle #1 | #2\rangle}
                             % This is an alternative definition of
                             % \ip that is commented out.

\begin{document}             % End of preamble and beginning of text.

\maketitle                   % Produces the title.

This is an example input file.  Comparing it with
the output it generates can show you how to
produce a simple document of your own.

\section{Ordinary Text}      % Produces section heading.  Lower-level
                             % sections are begun with similar
                             % \subsection and \subsubsection commands.

The ends  of words and sentences are marked
  by   spaces. It  doesn't matter how many
spaces    you type; one is as good as 100.  The
end of   a line counts as a space.

One   or more   blank lines denote the  end
of  a paragraph.

Since any number of consecutive spaces are treated
like a single one, the formatting of the input
file makes no difference to
      \LaTeX,                % The \LaTeX command generates the LaTeX logo.
but it makes a difference to you.  When you use
\LaTeX, making your input file as easy to read
as possible will be a great help as you write
your document and when you change it.  This sample
file shows how you can add comments to your own input
file.

Because printing is different from typewriting,
there are a number of things that you have to do
differently when preparing an input file than if
you were just typing the document directly.
Quotation marks like
       ``this''
have to be handled specially, as do quotes within
quotes:
       ``\,`this'            % \, separates the double and single quote.
        is what I just
        wrote, not  `that'\,''.

Dashes come in three sizes: an
       intra-word
dash, a medium dash for number ranges like
       1--2,
and a punctuation
       dash---like
this.

A sentence-ending space should be larger than the
space between words within a sentence.  You
sometimes have to type special commands in
conjunction with punctuation characters to get
this right, as in the following sentence.
       Gnats, gnus, etc.\ all  % `\ ' makes an inter-word space.
       begin with G\@.         % \@ marks end-of-sentence punctuation.
You should check the spaces after periods when
reading your output to make sure you haven't
forgotten any special cases.  Generating an
ellipsis
       \ldots\               % `\ ' is needed after `\ldots' because TeX
                             % ignores spaces after command names like \ldots
                             % made from \ + letters.
                             %
                             % Note how a `%' character causes TeX to ignore
                             % the end of the input line, so these blank lines
                             % do not start a new paragraph.
                             %
with the right spacing around the periods requires
a special command.

\LaTeX\ interprets some common characters as
commands, so you must type special commands to
generate them.  These characters include the
following:
       \$ \& \% \# \{ and \}.

In printing, text is usually emphasized with an
       \emph{italic}
type style.

\begin{em}
   A long segment of text can also be emphasized
   in this way.  Text within such a segment can be
   given \emph{additional} emphasis.
\end{em}

It is sometimes necessary to prevent \LaTeX\ from
breaking a line where it might otherwise do so.
This may be at a space, as between the ``Mr.''\ and
``Jones'' in
       ``Mr.~Jones'',        % ~ produces an unbreakable interword space.
or within a word---especially when the word is a
symbol like
       \mbox{\emph{itemnum}}
that makes little sense when hyphenated across
lines.

Footnotes\footnote{This is an example of a footnote.}
pose no problem.

\LaTeX\ is good at typesetting mathematical formulas
like
       \( x-3y + z = 7 \)
or
       \( a_{1} > x^{2n} + y^{2n} > x' \)
or
       \( \ip{A}{B} = \sum_{i} a_{i} b_{i} \).
The spaces you type in a formula are
ignored.  Remember that a letter like
       $x$                   % $ ... $  and  \( ... \)  are equivalent
is a formula when it denotes a mathematical
symbol, and it should be typed as one.

\section{Displayed Text}

Text is displayed by indenting it from the left
margin.  Quotations are commonly displayed.  There
are short quotations
\begin{quote}
   This is a short quotation.  It consists of a
   single paragraph of text.  See how it is formatted.
\end{quote}
and longer ones.
\begin{quotation}
   This is a longer quotation.  It consists of two
   paragraphs of text, neither of which are
   particularly interesting.

   This is the second paragraph of the quotation.  It
   is just as dull as the first paragraph.
\end{quotation}
Another frequently-displayed structure is a list.
The following is an example of an \emph{itemized}
list.
\begin{itemize}
   \item This is the first item of an itemized list.
         Each item in the list is marked with a ``tick''.
         You don't have to worry about what kind of tick
         mark is used.

   \item This is the second item of the list.  It
         contains another list nested inside it.  The inner
         list is an \emph{enumerated} list.
         \begin{enumerate}
            \item This is the first item of an enumerated
                  list that is nested within the itemized list.

            \item This is the second item of the inner list.
                  \LaTeX\ allows you to nest lists deeper than
                  you really should.
         \end{enumerate}
         This is the rest of the second item of the outer
         list.  It is no more interesting than any other
         part of the item.
   \item This is the third item of the list.
\end{itemize}
You can even display poetry.
\begin{verse}
   There is an environment
    for verse \\             % The \\ command separates lines
   Whose features some poets % within a stanza.
   will curse.

                             % One or more blank lines separate stanzas.

   For instead of making\\
   Them do \emph{all} line breaking, \\
   It allows them to put too many words on a line when they'd rather be
   forced to be terse.
\end{verse}

Mathematical formulas may also be displayed.  A
displayed formula
is
one-line long; multiline
formulas require special formatting instructions.
   \[  \ip{\Gamma}{\psi'} = x'' + y^{2} + z_{i}^{n}\]
Don't start a paragraph with a displayed equation,
nor make one a paragraph by itself.

\end{document}               % End of document.

\section{Bench Huge Section 573}
% synthetic bench section 573
% This is a sample LaTeX input file.  (Version of 12 August 2004.)
%
% A '%' character causes TeX to ignore all remaining text on the line,
% and is used for comments like this one.

\documentclass{article}      % Specifies the document class

                             % The preamble begins here.
\title{An Example Document}  % Declares the document's title.
\author{Leslie Lamport}      % Declares the author's name.
\date{January 21, 1994}      % Deleting this command produces today's date.

\newcommand{\ip}[2]{(#1, #2)}
                             % Defines \ip{arg1}{arg2} to mean
                             % (arg1, arg2).

%\newcommand{\ip}[2]{\langle #1 | #2\rangle}
                             % This is an alternative definition of
                             % \ip that is commented out.

\begin{document}             % End of preamble and beginning of text.

\maketitle                   % Produces the title.

This is an example input file.  Comparing it with
the output it generates can show you how to
produce a simple document of your own.

\section{Ordinary Text}      % Produces section heading.  Lower-level
                             % sections are begun with similar
                             % \subsection and \subsubsection commands.

The ends  of words and sentences are marked
  by   spaces. It  doesn't matter how many
spaces    you type; one is as good as 100.  The
end of   a line counts as a space.

One   or more   blank lines denote the  end
of  a paragraph.

Since any number of consecutive spaces are treated
like a single one, the formatting of the input
file makes no difference to
      \LaTeX,                % The \LaTeX command generates the LaTeX logo.
but it makes a difference to you.  When you use
\LaTeX, making your input file as easy to read
as possible will be a great help as you write
your document and when you change it.  This sample
file shows how you can add comments to your own input
file.

Because printing is different from typewriting,
there are a number of things that you have to do
differently when preparing an input file than if
you were just typing the document directly.
Quotation marks like
       ``this''
have to be handled specially, as do quotes within
quotes:
       ``\,`this'            % \, separates the double and single quote.
        is what I just
        wrote, not  `that'\,''.

Dashes come in three sizes: an
       intra-word
dash, a medium dash for number ranges like
       1--2,
and a punctuation
       dash---like
this.

A sentence-ending space should be larger than the
space between words within a sentence.  You
sometimes have to type special commands in
conjunction with punctuation characters to get
this right, as in the following sentence.
       Gnats, gnus, etc.\ all  % `\ ' makes an inter-word space.
       begin with G\@.         % \@ marks end-of-sentence punctuation.
You should check the spaces after periods when
reading your output to make sure you haven't
forgotten any special cases.  Generating an
ellipsis
       \ldots\               % `\ ' is needed after `\ldots' because TeX
                             % ignores spaces after command names like \ldots
                             % made from \ + letters.
                             %
                             % Note how a `%' character causes TeX to ignore
                             % the end of the input line, so these blank lines
                             % do not start a new paragraph.
                             %
with the right spacing around the periods requires
a special command.

\LaTeX\ interprets some common characters as
commands, so you must type special commands to
generate them.  These characters include the
following:
       \$ \& \% \# \{ and \}.

In printing, text is usually emphasized with an
       \emph{italic}
type style.

\begin{em}
   A long segment of text can also be emphasized
   in this way.  Text within such a segment can be
   given \emph{additional} emphasis.
\end{em}

It is sometimes necessary to prevent \LaTeX\ from
breaking a line where it might otherwise do so.
This may be at a space, as between the ``Mr.''\ and
``Jones'' in
       ``Mr.~Jones'',        % ~ produces an unbreakable interword space.
or within a word---especially when the word is a
symbol like
       \mbox{\emph{itemnum}}
that makes little sense when hyphenated across
lines.

Footnotes\footnote{This is an example of a footnote.}
pose no problem.

\LaTeX\ is good at typesetting mathematical formulas
like
       \( x-3y + z = 7 \)
or
       \( a_{1} > x^{2n} + y^{2n} > x' \)
or
       \( \ip{A}{B} = \sum_{i} a_{i} b_{i} \).
The spaces you type in a formula are
ignored.  Remember that a letter like
       $x$                   % $ ... $  and  \( ... \)  are equivalent
is a formula when it denotes a mathematical
symbol, and it should be typed as one.

\section{Displayed Text}

Text is displayed by indenting it from the left
margin.  Quotations are commonly displayed.  There
are short quotations
\begin{quote}
   This is a short quotation.  It consists of a
   single paragraph of text.  See how it is formatted.
\end{quote}
and longer ones.
\begin{quotation}
   This is a longer quotation.  It consists of two
   paragraphs of text, neither of which are
   particularly interesting.

   This is the second paragraph of the quotation.  It
   is just as dull as the first paragraph.
\end{quotation}
Another frequently-displayed structure is a list.
The following is an example of an \emph{itemized}
list.
\begin{itemize}
   \item This is the first item of an itemized list.
         Each item in the list is marked with a ``tick''.
         You don't have to worry about what kind of tick
         mark is used.

   \item This is the second item of the list.  It
         contains another list nested inside it.  The inner
         list is an \emph{enumerated} list.
         \begin{enumerate}
            \item This is the first item of an enumerated
                  list that is nested within the itemized list.

            \item This is the second item of the inner list.
                  \LaTeX\ allows you to nest lists deeper than
                  you really should.
         \end{enumerate}
         This is the rest of the second item of the outer
         list.  It is no more interesting than any other
         part of the item.
   \item This is the third item of the list.
\end{itemize}
You can even display poetry.
\begin{verse}
   There is an environment
    for verse \\             % The \\ command separates lines
   Whose features some poets % within a stanza.
   will curse.

                             % One or more blank lines separate stanzas.

   For instead of making\\
   Them do \emph{all} line breaking, \\
   It allows them to put too many words on a line when they'd rather be
   forced to be terse.
\end{verse}

Mathematical formulas may also be displayed.  A
displayed formula
is
one-line long; multiline
formulas require special formatting instructions.
   \[  \ip{\Gamma}{\psi'} = x'' + y^{2} + z_{i}^{n}\]
Don't start a paragraph with a displayed equation,
nor make one a paragraph by itself.

\end{document}               % End of document.

\section{Bench Huge Section 574}
% synthetic bench section 574
% This is a sample LaTeX input file.  (Version of 12 August 2004.)
%
% A '%' character causes TeX to ignore all remaining text on the line,
% and is used for comments like this one.

\documentclass{article}      % Specifies the document class

                             % The preamble begins here.
\title{An Example Document}  % Declares the document's title.
\author{Leslie Lamport}      % Declares the author's name.
\date{January 21, 1994}      % Deleting this command produces today's date.

\newcommand{\ip}[2]{(#1, #2)}
                             % Defines \ip{arg1}{arg2} to mean
                             % (arg1, arg2).

%\newcommand{\ip}[2]{\langle #1 | #2\rangle}
                             % This is an alternative definition of
                             % \ip that is commented out.

\begin{document}             % End of preamble and beginning of text.

\maketitle                   % Produces the title.

This is an example input file.  Comparing it with
the output it generates can show you how to
produce a simple document of your own.

\section{Ordinary Text}      % Produces section heading.  Lower-level
                             % sections are begun with similar
                             % \subsection and \subsubsection commands.

The ends  of words and sentences are marked
  by   spaces. It  doesn't matter how many
spaces    you type; one is as good as 100.  The
end of   a line counts as a space.

One   or more   blank lines denote the  end
of  a paragraph.

Since any number of consecutive spaces are treated
like a single one, the formatting of the input
file makes no difference to
      \LaTeX,                % The \LaTeX command generates the LaTeX logo.
but it makes a difference to you.  When you use
\LaTeX, making your input file as easy to read
as possible will be a great help as you write
your document and when you change it.  This sample
file shows how you can add comments to your own input
file.

Because printing is different from typewriting,
there are a number of things that you have to do
differently when preparing an input file than if
you were just typing the document directly.
Quotation marks like
       ``this''
have to be handled specially, as do quotes within
quotes:
       ``\,`this'            % \, separates the double and single quote.
        is what I just
        wrote, not  `that'\,''.

Dashes come in three sizes: an
       intra-word
dash, a medium dash for number ranges like
       1--2,
and a punctuation
       dash---like
this.

A sentence-ending space should be larger than the
space between words within a sentence.  You
sometimes have to type special commands in
conjunction with punctuation characters to get
this right, as in the following sentence.
       Gnats, gnus, etc.\ all  % `\ ' makes an inter-word space.
       begin with G\@.         % \@ marks end-of-sentence punctuation.
You should check the spaces after periods when
reading your output to make sure you haven't
forgotten any special cases.  Generating an
ellipsis
       \ldots\               % `\ ' is needed after `\ldots' because TeX
                             % ignores spaces after command names like \ldots
                             % made from \ + letters.
                             %
                             % Note how a `%' character causes TeX to ignore
                             % the end of the input line, so these blank lines
                             % do not start a new paragraph.
                             %
with the right spacing around the periods requires
a special command.

\LaTeX\ interprets some common characters as
commands, so you must type special commands to
generate them.  These characters include the
following:
       \$ \& \% \# \{ and \}.

In printing, text is usually emphasized with an
       \emph{italic}
type style.

\begin{em}
   A long segment of text can also be emphasized
   in this way.  Text within such a segment can be
   given \emph{additional} emphasis.
\end{em}

It is sometimes necessary to prevent \LaTeX\ from
breaking a line where it might otherwise do so.
This may be at a space, as between the ``Mr.''\ and
``Jones'' in
       ``Mr.~Jones'',        % ~ produces an unbreakable interword space.
or within a word---especially when the word is a
symbol like
       \mbox{\emph{itemnum}}
that makes little sense when hyphenated across
lines.

Footnotes\footnote{This is an example of a footnote.}
pose no problem.

\LaTeX\ is good at typesetting mathematical formulas
like
       \( x-3y + z = 7 \)
or
       \( a_{1} > x^{2n} + y^{2n} > x' \)
or
       \( \ip{A}{B} = \sum_{i} a_{i} b_{i} \).
The spaces you type in a formula are
ignored.  Remember that a letter like
       $x$                   % $ ... $  and  \( ... \)  are equivalent
is a formula when it denotes a mathematical
symbol, and it should be typed as one.

\section{Displayed Text}

Text is displayed by indenting it from the left
margin.  Quotations are commonly displayed.  There
are short quotations
\begin{quote}
   This is a short quotation.  It consists of a
   single paragraph of text.  See how it is formatted.
\end{quote}
and longer ones.
\begin{quotation}
   This is a longer quotation.  It consists of two
   paragraphs of text, neither of which are
   particularly interesting.

   This is the second paragraph of the quotation.  It
   is just as dull as the first paragraph.
\end{quotation}
Another frequently-displayed structure is a list.
The following is an example of an \emph{itemized}
list.
\begin{itemize}
   \item This is the first item of an itemized list.
         Each item in the list is marked with a ``tick''.
         You don't have to worry about what kind of tick
         mark is used.

   \item This is the second item of the list.  It
         contains another list nested inside it.  The inner
         list is an \emph{enumerated} list.
         \begin{enumerate}
            \item This is the first item of an enumerated
                  list that is nested within the itemized list.

            \item This is the second item of the inner list.
                  \LaTeX\ allows you to nest lists deeper than
                  you really should.
         \end{enumerate}
         This is the rest of the second item of the outer
         list.  It is no more interesting than any other
         part of the item.
   \item This is the third item of the list.
\end{itemize}
You can even display poetry.
\begin{verse}
   There is an environment
    for verse \\             % The \\ command separates lines
   Whose features some poets % within a stanza.
   will curse.

                             % One or more blank lines separate stanzas.

   For instead of making\\
   Them do \emph{all} line breaking, \\
   It allows them to put too many words on a line when they'd rather be
   forced to be terse.
\end{verse}

Mathematical formulas may also be displayed.  A
displayed formula
is
one-line long; multiline
formulas require special formatting instructions.
   \[  \ip{\Gamma}{\psi'} = x'' + y^{2} + z_{i}^{n}\]
Don't start a paragraph with a displayed equation,
nor make one a paragraph by itself.

\end{document}               % End of document.

\section{Bench Huge Section 575}
% synthetic bench section 575
% This is a sample LaTeX input file.  (Version of 12 August 2004.)
%
% A '%' character causes TeX to ignore all remaining text on the line,
% and is used for comments like this one.

\documentclass{article}      % Specifies the document class

                             % The preamble begins here.
\title{An Example Document}  % Declares the document's title.
\author{Leslie Lamport}      % Declares the author's name.
\date{January 21, 1994}      % Deleting this command produces today's date.

\newcommand{\ip}[2]{(#1, #2)}
                             % Defines \ip{arg1}{arg2} to mean
                             % (arg1, arg2).

%\newcommand{\ip}[2]{\langle #1 | #2\rangle}
                             % This is an alternative definition of
                             % \ip that is commented out.

\begin{document}             % End of preamble and beginning of text.

\maketitle                   % Produces the title.

This is an example input file.  Comparing it with
the output it generates can show you how to
produce a simple document of your own.

\section{Ordinary Text}      % Produces section heading.  Lower-level
                             % sections are begun with similar
                             % \subsection and \subsubsection commands.

The ends  of words and sentences are marked
  by   spaces. It  doesn't matter how many
spaces    you type; one is as good as 100.  The
end of   a line counts as a space.

One   or more   blank lines denote the  end
of  a paragraph.

Since any number of consecutive spaces are treated
like a single one, the formatting of the input
file makes no difference to
      \LaTeX,                % The \LaTeX command generates the LaTeX logo.
but it makes a difference to you.  When you use
\LaTeX, making your input file as easy to read
as possible will be a great help as you write
your document and when you change it.  This sample
file shows how you can add comments to your own input
file.

Because printing is different from typewriting,
there are a number of things that you have to do
differently when preparing an input file than if
you were just typing the document directly.
Quotation marks like
       ``this''
have to be handled specially, as do quotes within
quotes:
       ``\,`this'            % \, separates the double and single quote.
        is what I just
        wrote, not  `that'\,''.

Dashes come in three sizes: an
       intra-word
dash, a medium dash for number ranges like
       1--2,
and a punctuation
       dash---like
this.

A sentence-ending space should be larger than the
space between words within a sentence.  You
sometimes have to type special commands in
conjunction with punctuation characters to get
this right, as in the following sentence.
       Gnats, gnus, etc.\ all  % `\ ' makes an inter-word space.
       begin with G\@.         % \@ marks end-of-sentence punctuation.
You should check the spaces after periods when
reading your output to make sure you haven't
forgotten any special cases.  Generating an
ellipsis
       \ldots\               % `\ ' is needed after `\ldots' because TeX
                             % ignores spaces after command names like \ldots
                             % made from \ + letters.
                             %
                             % Note how a `%' character causes TeX to ignore
                             % the end of the input line, so these blank lines
                             % do not start a new paragraph.
                             %
with the right spacing around the periods requires
a special command.

\LaTeX\ interprets some common characters as
commands, so you must type special commands to
generate them.  These characters include the
following:
       \$ \& \% \# \{ and \}.

In printing, text is usually emphasized with an
       \emph{italic}
type style.

\begin{em}
   A long segment of text can also be emphasized
   in this way.  Text within such a segment can be
   given \emph{additional} emphasis.
\end{em}

It is sometimes necessary to prevent \LaTeX\ from
breaking a line where it might otherwise do so.
This may be at a space, as between the ``Mr.''\ and
``Jones'' in
       ``Mr.~Jones'',        % ~ produces an unbreakable interword space.
or within a word---especially when the word is a
symbol like
       \mbox{\emph{itemnum}}
that makes little sense when hyphenated across
lines.

Footnotes\footnote{This is an example of a footnote.}
pose no problem.

\LaTeX\ is good at typesetting mathematical formulas
like
       \( x-3y + z = 7 \)
or
       \( a_{1} > x^{2n} + y^{2n} > x' \)
or
       \( \ip{A}{B} = \sum_{i} a_{i} b_{i} \).
The spaces you type in a formula are
ignored.  Remember that a letter like
       $x$                   % $ ... $  and  \( ... \)  are equivalent
is a formula when it denotes a mathematical
symbol, and it should be typed as one.

\section{Displayed Text}

Text is displayed by indenting it from the left
margin.  Quotations are commonly displayed.  There
are short quotations
\begin{quote}
   This is a short quotation.  It consists of a
   single paragraph of text.  See how it is formatted.
\end{quote}
and longer ones.
\begin{quotation}
   This is a longer quotation.  It consists of two
   paragraphs of text, neither of which are
   particularly interesting.

   This is the second paragraph of the quotation.  It
   is just as dull as the first paragraph.
\end{quotation}
Another frequently-displayed structure is a list.
The following is an example of an \emph{itemized}
list.
\begin{itemize}
   \item This is the first item of an itemized list.
         Each item in the list is marked with a ``tick''.
         You don't have to worry about what kind of tick
         mark is used.

   \item This is the second item of the list.  It
         contains another list nested inside it.  The inner
         list is an \emph{enumerated} list.
         \begin{enumerate}
            \item This is the first item of an enumerated
                  list that is nested within the itemized list.

            \item This is the second item of the inner list.
                  \LaTeX\ allows you to nest lists deeper than
                  you really should.
         \end{enumerate}
         This is the rest of the second item of the outer
         list.  It is no more interesting than any other
         part of the item.
   \item This is the third item of the list.
\end{itemize}
You can even display poetry.
\begin{verse}
   There is an environment
    for verse \\             % The \\ command separates lines
   Whose features some poets % within a stanza.
   will curse.

                             % One or more blank lines separate stanzas.

   For instead of making\\
   Them do \emph{all} line breaking, \\
   It allows them to put too many words on a line when they'd rather be
   forced to be terse.
\end{verse}

Mathematical formulas may also be displayed.  A
displayed formula
is
one-line long; multiline
formulas require special formatting instructions.
   \[  \ip{\Gamma}{\psi'} = x'' + y^{2} + z_{i}^{n}\]
Don't start a paragraph with a displayed equation,
nor make one a paragraph by itself.

\end{document}               % End of document.

\section{Bench Huge Section 576}
% synthetic bench section 576
% This is a sample LaTeX input file.  (Version of 12 August 2004.)
%
% A '%' character causes TeX to ignore all remaining text on the line,
% and is used for comments like this one.

\documentclass{article}      % Specifies the document class

                             % The preamble begins here.
\title{An Example Document}  % Declares the document's title.
\author{Leslie Lamport}      % Declares the author's name.
\date{January 21, 1994}      % Deleting this command produces today's date.

\newcommand{\ip}[2]{(#1, #2)}
                             % Defines \ip{arg1}{arg2} to mean
                             % (arg1, arg2).

%\newcommand{\ip}[2]{\langle #1 | #2\rangle}
                             % This is an alternative definition of
                             % \ip that is commented out.

\begin{document}             % End of preamble and beginning of text.

\maketitle                   % Produces the title.

This is an example input file.  Comparing it with
the output it generates can show you how to
produce a simple document of your own.

\section{Ordinary Text}      % Produces section heading.  Lower-level
                             % sections are begun with similar
                             % \subsection and \subsubsection commands.

The ends  of words and sentences are marked
  by   spaces. It  doesn't matter how many
spaces    you type; one is as good as 100.  The
end of   a line counts as a space.

One   or more   blank lines denote the  end
of  a paragraph.

Since any number of consecutive spaces are treated
like a single one, the formatting of the input
file makes no difference to
      \LaTeX,                % The \LaTeX command generates the LaTeX logo.
but it makes a difference to you.  When you use
\LaTeX, making your input file as easy to read
as possible will be a great help as you write
your document and when you change it.  This sample
file shows how you can add comments to your own input
file.

Because printing is different from typewriting,
there are a number of things that you have to do
differently when preparing an input file than if
you were just typing the document directly.
Quotation marks like
       ``this''
have to be handled specially, as do quotes within
quotes:
       ``\,`this'            % \, separates the double and single quote.
        is what I just
        wrote, not  `that'\,''.

Dashes come in three sizes: an
       intra-word
dash, a medium dash for number ranges like
       1--2,
and a punctuation
       dash---like
this.

A sentence-ending space should be larger than the
space between words within a sentence.  You
sometimes have to type special commands in
conjunction with punctuation characters to get
this right, as in the following sentence.
       Gnats, gnus, etc.\ all  % `\ ' makes an inter-word space.
       begin with G\@.         % \@ marks end-of-sentence punctuation.
You should check the spaces after periods when
reading your output to make sure you haven't
forgotten any special cases.  Generating an
ellipsis
       \ldots\               % `\ ' is needed after `\ldots' because TeX
                             % ignores spaces after command names like \ldots
                             % made from \ + letters.
                             %
                             % Note how a `%' character causes TeX to ignore
                             % the end of the input line, so these blank lines
                             % do not start a new paragraph.
                             %
with the right spacing around the periods requires
a special command.

\LaTeX\ interprets some common characters as
commands, so you must type special commands to
generate them.  These characters include the
following:
       \$ \& \% \# \{ and \}.

In printing, text is usually emphasized with an
       \emph{italic}
type style.

\begin{em}
   A long segment of text can also be emphasized
   in this way.  Text within such a segment can be
   given \emph{additional} emphasis.
\end{em}

It is sometimes necessary to prevent \LaTeX\ from
breaking a line where it might otherwise do so.
This may be at a space, as between the ``Mr.''\ and
``Jones'' in
       ``Mr.~Jones'',        % ~ produces an unbreakable interword space.
or within a word---especially when the word is a
symbol like
       \mbox{\emph{itemnum}}
that makes little sense when hyphenated across
lines.

Footnotes\footnote{This is an example of a footnote.}
pose no problem.

\LaTeX\ is good at typesetting mathematical formulas
like
       \( x-3y + z = 7 \)
or
       \( a_{1} > x^{2n} + y^{2n} > x' \)
or
       \( \ip{A}{B} = \sum_{i} a_{i} b_{i} \).
The spaces you type in a formula are
ignored.  Remember that a letter like
       $x$                   % $ ... $  and  \( ... \)  are equivalent
is a formula when it denotes a mathematical
symbol, and it should be typed as one.

\section{Displayed Text}

Text is displayed by indenting it from the left
margin.  Quotations are commonly displayed.  There
are short quotations
\begin{quote}
   This is a short quotation.  It consists of a
   single paragraph of text.  See how it is formatted.
\end{quote}
and longer ones.
\begin{quotation}
   This is a longer quotation.  It consists of two
   paragraphs of text, neither of which are
   particularly interesting.

   This is the second paragraph of the quotation.  It
   is just as dull as the first paragraph.
\end{quotation}
Another frequently-displayed structure is a list.
The following is an example of an \emph{itemized}
list.
\begin{itemize}
   \item This is the first item of an itemized list.
         Each item in the list is marked with a ``tick''.
         You don't have to worry about what kind of tick
         mark is used.

   \item This is the second item of the list.  It
         contains another list nested inside it.  The inner
         list is an \emph{enumerated} list.
         \begin{enumerate}
            \item This is the first item of an enumerated
                  list that is nested within the itemized list.

            \item This is the second item of the inner list.
                  \LaTeX\ allows you to nest lists deeper than
                  you really should.
         \end{enumerate}
         This is the rest of the second item of the outer
         list.  It is no more interesting than any other
         part of the item.
   \item This is the third item of the list.
\end{itemize}
You can even display poetry.
\begin{verse}
   There is an environment
    for verse \\             % The \\ command separates lines
   Whose features some poets % within a stanza.
   will curse.

                             % One or more blank lines separate stanzas.

   For instead of making\\
   Them do \emph{all} line breaking, \\
   It allows them to put too many words on a line when they'd rather be
   forced to be terse.
\end{verse}

Mathematical formulas may also be displayed.  A
displayed formula
is
one-line long; multiline
formulas require special formatting instructions.
   \[  \ip{\Gamma}{\psi'} = x'' + y^{2} + z_{i}^{n}\]
Don't start a paragraph with a displayed equation,
nor make one a paragraph by itself.

\end{document}               % End of document.

\section{Bench Huge Section 577}
% synthetic bench section 577
% This is a sample LaTeX input file.  (Version of 12 August 2004.)
%
% A '%' character causes TeX to ignore all remaining text on the line,
% and is used for comments like this one.

\documentclass{article}      % Specifies the document class

                             % The preamble begins here.
\title{An Example Document}  % Declares the document's title.
\author{Leslie Lamport}      % Declares the author's name.
\date{January 21, 1994}      % Deleting this command produces today's date.

\newcommand{\ip}[2]{(#1, #2)}
                             % Defines \ip{arg1}{arg2} to mean
                             % (arg1, arg2).

%\newcommand{\ip}[2]{\langle #1 | #2\rangle}
                             % This is an alternative definition of
                             % \ip that is commented out.

\begin{document}             % End of preamble and beginning of text.

\maketitle                   % Produces the title.

This is an example input file.  Comparing it with
the output it generates can show you how to
produce a simple document of your own.

\section{Ordinary Text}      % Produces section heading.  Lower-level
                             % sections are begun with similar
                             % \subsection and \subsubsection commands.

The ends  of words and sentences are marked
  by   spaces. It  doesn't matter how many
spaces    you type; one is as good as 100.  The
end of   a line counts as a space.

One   or more   blank lines denote the  end
of  a paragraph.

Since any number of consecutive spaces are treated
like a single one, the formatting of the input
file makes no difference to
      \LaTeX,                % The \LaTeX command generates the LaTeX logo.
but it makes a difference to you.  When you use
\LaTeX, making your input file as easy to read
as possible will be a great help as you write
your document and when you change it.  This sample
file shows how you can add comments to your own input
file.

Because printing is different from typewriting,
there are a number of things that you have to do
differently when preparing an input file than if
you were just typing the document directly.
Quotation marks like
       ``this''
have to be handled specially, as do quotes within
quotes:
       ``\,`this'            % \, separates the double and single quote.
        is what I just
        wrote, not  `that'\,''.

Dashes come in three sizes: an
       intra-word
dash, a medium dash for number ranges like
       1--2,
and a punctuation
       dash---like
this.

A sentence-ending space should be larger than the
space between words within a sentence.  You
sometimes have to type special commands in
conjunction with punctuation characters to get
this right, as in the following sentence.
       Gnats, gnus, etc.\ all  % `\ ' makes an inter-word space.
       begin with G\@.         % \@ marks end-of-sentence punctuation.
You should check the spaces after periods when
reading your output to make sure you haven't
forgotten any special cases.  Generating an
ellipsis
       \ldots\               % `\ ' is needed after `\ldots' because TeX
                             % ignores spaces after command names like \ldots
                             % made from \ + letters.
                             %
                             % Note how a `%' character causes TeX to ignore
                             % the end of the input line, so these blank lines
                             % do not start a new paragraph.
                             %
with the right spacing around the periods requires
a special command.

\LaTeX\ interprets some common characters as
commands, so you must type special commands to
generate them.  These characters include the
following:
       \$ \& \% \# \{ and \}.

In printing, text is usually emphasized with an
       \emph{italic}
type style.

\begin{em}
   A long segment of text can also be emphasized
   in this way.  Text within such a segment can be
   given \emph{additional} emphasis.
\end{em}

It is sometimes necessary to prevent \LaTeX\ from
breaking a line where it might otherwise do so.
This may be at a space, as between the ``Mr.''\ and
``Jones'' in
       ``Mr.~Jones'',        % ~ produces an unbreakable interword space.
or within a word---especially when the word is a
symbol like
       \mbox{\emph{itemnum}}
that makes little sense when hyphenated across
lines.

Footnotes\footnote{This is an example of a footnote.}
pose no problem.

\LaTeX\ is good at typesetting mathematical formulas
like
       \( x-3y + z = 7 \)
or
       \( a_{1} > x^{2n} + y^{2n} > x' \)
or
       \( \ip{A}{B} = \sum_{i} a_{i} b_{i} \).
The spaces you type in a formula are
ignored.  Remember that a letter like
       $x$                   % $ ... $  and  \( ... \)  are equivalent
is a formula when it denotes a mathematical
symbol, and it should be typed as one.

\section{Displayed Text}

Text is displayed by indenting it from the left
margin.  Quotations are commonly displayed.  There
are short quotations
\begin{quote}
   This is a short quotation.  It consists of a
   single paragraph of text.  See how it is formatted.
\end{quote}
and longer ones.
\begin{quotation}
   This is a longer quotation.  It consists of two
   paragraphs of text, neither of which are
   particularly interesting.

   This is the second paragraph of the quotation.  It
   is just as dull as the first paragraph.
\end{quotation}
Another frequently-displayed structure is a list.
The following is an example of an \emph{itemized}
list.
\begin{itemize}
   \item This is the first item of an itemized list.
         Each item in the list is marked with a ``tick''.
         You don't have to worry about what kind of tick
         mark is used.

   \item This is the second item of the list.  It
         contains another list nested inside it.  The inner
         list is an \emph{enumerated} list.
         \begin{enumerate}
            \item This is the first item of an enumerated
                  list that is nested within the itemized list.

            \item This is the second item of the inner list.
                  \LaTeX\ allows you to nest lists deeper than
                  you really should.
         \end{enumerate}
         This is the rest of the second item of the outer
         list.  It is no more interesting than any other
         part of the item.
   \item This is the third item of the list.
\end{itemize}
You can even display poetry.
\begin{verse}
   There is an environment
    for verse \\             % The \\ command separates lines
   Whose features some poets % within a stanza.
   will curse.

                             % One or more blank lines separate stanzas.

   For instead of making\\
   Them do \emph{all} line breaking, \\
   It allows them to put too many words on a line when they'd rather be
   forced to be terse.
\end{verse}

Mathematical formulas may also be displayed.  A
displayed formula
is
one-line long; multiline
formulas require special formatting instructions.
   \[  \ip{\Gamma}{\psi'} = x'' + y^{2} + z_{i}^{n}\]
Don't start a paragraph with a displayed equation,
nor make one a paragraph by itself.

\end{document}               % End of document.

\section{Bench Huge Section 578}
% synthetic bench section 578
% This is a sample LaTeX input file.  (Version of 12 August 2004.)
%
% A '%' character causes TeX to ignore all remaining text on the line,
% and is used for comments like this one.

\documentclass{article}      % Specifies the document class

                             % The preamble begins here.
\title{An Example Document}  % Declares the document's title.
\author{Leslie Lamport}      % Declares the author's name.
\date{January 21, 1994}      % Deleting this command produces today's date.

\newcommand{\ip}[2]{(#1, #2)}
                             % Defines \ip{arg1}{arg2} to mean
                             % (arg1, arg2).

%\newcommand{\ip}[2]{\langle #1 | #2\rangle}
                             % This is an alternative definition of
                             % \ip that is commented out.

\begin{document}             % End of preamble and beginning of text.

\maketitle                   % Produces the title.

This is an example input file.  Comparing it with
the output it generates can show you how to
produce a simple document of your own.

\section{Ordinary Text}      % Produces section heading.  Lower-level
                             % sections are begun with similar
                             % \subsection and \subsubsection commands.

The ends  of words and sentences are marked
  by   spaces. It  doesn't matter how many
spaces    you type; one is as good as 100.  The
end of   a line counts as a space.

One   or more   blank lines denote the  end
of  a paragraph.

Since any number of consecutive spaces are treated
like a single one, the formatting of the input
file makes no difference to
      \LaTeX,                % The \LaTeX command generates the LaTeX logo.
but it makes a difference to you.  When you use
\LaTeX, making your input file as easy to read
as possible will be a great help as you write
your document and when you change it.  This sample
file shows how you can add comments to your own input
file.

Because printing is different from typewriting,
there are a number of things that you have to do
differently when preparing an input file than if
you were just typing the document directly.
Quotation marks like
       ``this''
have to be handled specially, as do quotes within
quotes:
       ``\,`this'            % \, separates the double and single quote.
        is what I just
        wrote, not  `that'\,''.

Dashes come in three sizes: an
       intra-word
dash, a medium dash for number ranges like
       1--2,
and a punctuation
       dash---like
this.

A sentence-ending space should be larger than the
space between words within a sentence.  You
sometimes have to type special commands in
conjunction with punctuation characters to get
this right, as in the following sentence.
       Gnats, gnus, etc.\ all  % `\ ' makes an inter-word space.
       begin with G\@.         % \@ marks end-of-sentence punctuation.
You should check the spaces after periods when
reading your output to make sure you haven't
forgotten any special cases.  Generating an
ellipsis
       \ldots\               % `\ ' is needed after `\ldots' because TeX
                             % ignores spaces after command names like \ldots
                             % made from \ + letters.
                             %
                             % Note how a `%' character causes TeX to ignore
                             % the end of the input line, so these blank lines
                             % do not start a new paragraph.
                             %
with the right spacing around the periods requires
a special command.

\LaTeX\ interprets some common characters as
commands, so you must type special commands to
generate them.  These characters include the
following:
       \$ \& \% \# \{ and \}.

In printing, text is usually emphasized with an
       \emph{italic}
type style.

\begin{em}
   A long segment of text can also be emphasized
   in this way.  Text within such a segment can be
   given \emph{additional} emphasis.
\end{em}

It is sometimes necessary to prevent \LaTeX\ from
breaking a line where it might otherwise do so.
This may be at a space, as between the ``Mr.''\ and
``Jones'' in
       ``Mr.~Jones'',        % ~ produces an unbreakable interword space.
or within a word---especially when the word is a
symbol like
       \mbox{\emph{itemnum}}
that makes little sense when hyphenated across
lines.

Footnotes\footnote{This is an example of a footnote.}
pose no problem.

\LaTeX\ is good at typesetting mathematical formulas
like
       \( x-3y + z = 7 \)
or
       \( a_{1} > x^{2n} + y^{2n} > x' \)
or
       \( \ip{A}{B} = \sum_{i} a_{i} b_{i} \).
The spaces you type in a formula are
ignored.  Remember that a letter like
       $x$                   % $ ... $  and  \( ... \)  are equivalent
is a formula when it denotes a mathematical
symbol, and it should be typed as one.

\section{Displayed Text}

Text is displayed by indenting it from the left
margin.  Quotations are commonly displayed.  There
are short quotations
\begin{quote}
   This is a short quotation.  It consists of a
   single paragraph of text.  See how it is formatted.
\end{quote}
and longer ones.
\begin{quotation}
   This is a longer quotation.  It consists of two
   paragraphs of text, neither of which are
   particularly interesting.

   This is the second paragraph of the quotation.  It
   is just as dull as the first paragraph.
\end{quotation}
Another frequently-displayed structure is a list.
The following is an example of an \emph{itemized}
list.
\begin{itemize}
   \item This is the first item of an itemized list.
         Each item in the list is marked with a ``tick''.
         You don't have to worry about what kind of tick
         mark is used.

   \item This is the second item of the list.  It
         contains another list nested inside it.  The inner
         list is an \emph{enumerated} list.
         \begin{enumerate}
            \item This is the first item of an enumerated
                  list that is nested within the itemized list.

            \item This is the second item of the inner list.
                  \LaTeX\ allows you to nest lists deeper than
                  you really should.
         \end{enumerate}
         This is the rest of the second item of the outer
         list.  It is no more interesting than any other
         part of the item.
   \item This is the third item of the list.
\end{itemize}
You can even display poetry.
\begin{verse}
   There is an environment
    for verse \\             % The \\ command separates lines
   Whose features some poets % within a stanza.
   will curse.

                             % One or more blank lines separate stanzas.

   For instead of making\\
   Them do \emph{all} line breaking, \\
   It allows them to put too many words on a line when they'd rather be
   forced to be terse.
\end{verse}

Mathematical formulas may also be displayed.  A
displayed formula
is
one-line long; multiline
formulas require special formatting instructions.
   \[  \ip{\Gamma}{\psi'} = x'' + y^{2} + z_{i}^{n}\]
Don't start a paragraph with a displayed equation,
nor make one a paragraph by itself.

\end{document}               % End of document.

\section{Bench Huge Section 579}
% synthetic bench section 579
% This is a sample LaTeX input file.  (Version of 12 August 2004.)
%
% A '%' character causes TeX to ignore all remaining text on the line,
% and is used for comments like this one.

\documentclass{article}      % Specifies the document class

                             % The preamble begins here.
\title{An Example Document}  % Declares the document's title.
\author{Leslie Lamport}      % Declares the author's name.
\date{January 21, 1994}      % Deleting this command produces today's date.

\newcommand{\ip}[2]{(#1, #2)}
                             % Defines \ip{arg1}{arg2} to mean
                             % (arg1, arg2).

%\newcommand{\ip}[2]{\langle #1 | #2\rangle}
                             % This is an alternative definition of
                             % \ip that is commented out.

\begin{document}             % End of preamble and beginning of text.

\maketitle                   % Produces the title.

This is an example input file.  Comparing it with
the output it generates can show you how to
produce a simple document of your own.

\section{Ordinary Text}      % Produces section heading.  Lower-level
                             % sections are begun with similar
                             % \subsection and \subsubsection commands.

The ends  of words and sentences are marked
  by   spaces. It  doesn't matter how many
spaces    you type; one is as good as 100.  The
end of   a line counts as a space.

One   or more   blank lines denote the  end
of  a paragraph.

Since any number of consecutive spaces are treated
like a single one, the formatting of the input
file makes no difference to
      \LaTeX,                % The \LaTeX command generates the LaTeX logo.
but it makes a difference to you.  When you use
\LaTeX, making your input file as easy to read
as possible will be a great help as you write
your document and when you change it.  This sample
file shows how you can add comments to your own input
file.

Because printing is different from typewriting,
there are a number of things that you have to do
differently when preparing an input file than if
you were just typing the document directly.
Quotation marks like
       ``this''
have to be handled specially, as do quotes within
quotes:
       ``\,`this'            % \, separates the double and single quote.
        is what I just
        wrote, not  `that'\,''.

Dashes come in three sizes: an
       intra-word
dash, a medium dash for number ranges like
       1--2,
and a punctuation
       dash---like
this.

A sentence-ending space should be larger than the
space between words within a sentence.  You
sometimes have to type special commands in
conjunction with punctuation characters to get
this right, as in the following sentence.
       Gnats, gnus, etc.\ all  % `\ ' makes an inter-word space.
       begin with G\@.         % \@ marks end-of-sentence punctuation.
You should check the spaces after periods when
reading your output to make sure you haven't
forgotten any special cases.  Generating an
ellipsis
       \ldots\               % `\ ' is needed after `\ldots' because TeX
                             % ignores spaces after command names like \ldots
                             % made from \ + letters.
                             %
                             % Note how a `%' character causes TeX to ignore
                             % the end of the input line, so these blank lines
                             % do not start a new paragraph.
                             %
with the right spacing around the periods requires
a special command.

\LaTeX\ interprets some common characters as
commands, so you must type special commands to
generate them.  These characters include the
following:
       \$ \& \% \# \{ and \}.

In printing, text is usually emphasized with an
       \emph{italic}
type style.

\begin{em}
   A long segment of text can also be emphasized
   in this way.  Text within such a segment can be
   given \emph{additional} emphasis.
\end{em}

It is sometimes necessary to prevent \LaTeX\ from
breaking a line where it might otherwise do so.
This may be at a space, as between the ``Mr.''\ and
``Jones'' in
       ``Mr.~Jones'',        % ~ produces an unbreakable interword space.
or within a word---especially when the word is a
symbol like
       \mbox{\emph{itemnum}}
that makes little sense when hyphenated across
lines.

Footnotes\footnote{This is an example of a footnote.}
pose no problem.

\LaTeX\ is good at typesetting mathematical formulas
like
       \( x-3y + z = 7 \)
or
       \( a_{1} > x^{2n} + y^{2n} > x' \)
or
       \( \ip{A}{B} = \sum_{i} a_{i} b_{i} \).
The spaces you type in a formula are
ignored.  Remember that a letter like
       $x$                   % $ ... $  and  \( ... \)  are equivalent
is a formula when it denotes a mathematical
symbol, and it should be typed as one.

\section{Displayed Text}

Text is displayed by indenting it from the left
margin.  Quotations are commonly displayed.  There
are short quotations
\begin{quote}
   This is a short quotation.  It consists of a
   single paragraph of text.  See how it is formatted.
\end{quote}
and longer ones.
\begin{quotation}
   This is a longer quotation.  It consists of two
   paragraphs of text, neither of which are
   particularly interesting.

   This is the second paragraph of the quotation.  It
   is just as dull as the first paragraph.
\end{quotation}
Another frequently-displayed structure is a list.
The following is an example of an \emph{itemized}
list.
\begin{itemize}
   \item This is the first item of an itemized list.
         Each item in the list is marked with a ``tick''.
         You don't have to worry about what kind of tick
         mark is used.

   \item This is the second item of the list.  It
         contains another list nested inside it.  The inner
         list is an \emph{enumerated} list.
         \begin{enumerate}
            \item This is the first item of an enumerated
                  list that is nested within the itemized list.

            \item This is the second item of the inner list.
                  \LaTeX\ allows you to nest lists deeper than
                  you really should.
         \end{enumerate}
         This is the rest of the second item of the outer
         list.  It is no more interesting than any other
         part of the item.
   \item This is the third item of the list.
\end{itemize}
You can even display poetry.
\begin{verse}
   There is an environment
    for verse \\             % The \\ command separates lines
   Whose features some poets % within a stanza.
   will curse.

                             % One or more blank lines separate stanzas.

   For instead of making\\
   Them do \emph{all} line breaking, \\
   It allows them to put too many words on a line when they'd rather be
   forced to be terse.
\end{verse}

Mathematical formulas may also be displayed.  A
displayed formula
is
one-line long; multiline
formulas require special formatting instructions.
   \[  \ip{\Gamma}{\psi'} = x'' + y^{2} + z_{i}^{n}\]
Don't start a paragraph with a displayed equation,
nor make one a paragraph by itself.

\end{document}               % End of document.

\section{Bench Huge Section 580}
% synthetic bench section 580
% This is a sample LaTeX input file.  (Version of 12 August 2004.)
%
% A '%' character causes TeX to ignore all remaining text on the line,
% and is used for comments like this one.

\documentclass{article}      % Specifies the document class

                             % The preamble begins here.
\title{An Example Document}  % Declares the document's title.
\author{Leslie Lamport}      % Declares the author's name.
\date{January 21, 1994}      % Deleting this command produces today's date.

\newcommand{\ip}[2]{(#1, #2)}
                             % Defines \ip{arg1}{arg2} to mean
                             % (arg1, arg2).

%\newcommand{\ip}[2]{\langle #1 | #2\rangle}
                             % This is an alternative definition of
                             % \ip that is commented out.

\begin{document}             % End of preamble and beginning of text.

\maketitle                   % Produces the title.

This is an example input file.  Comparing it with
the output it generates can show you how to
produce a simple document of your own.

\section{Ordinary Text}      % Produces section heading.  Lower-level
                             % sections are begun with similar
                             % \subsection and \subsubsection commands.

The ends  of words and sentences are marked
  by   spaces. It  doesn't matter how many
spaces    you type; one is as good as 100.  The
end of   a line counts as a space.

One   or more   blank lines denote the  end
of  a paragraph.

Since any number of consecutive spaces are treated
like a single one, the formatting of the input
file makes no difference to
      \LaTeX,                % The \LaTeX command generates the LaTeX logo.
but it makes a difference to you.  When you use
\LaTeX, making your input file as easy to read
as possible will be a great help as you write
your document and when you change it.  This sample
file shows how you can add comments to your own input
file.

Because printing is different from typewriting,
there are a number of things that you have to do
differently when preparing an input file than if
you were just typing the document directly.
Quotation marks like
       ``this''
have to be handled specially, as do quotes within
quotes:
       ``\,`this'            % \, separates the double and single quote.
        is what I just
        wrote, not  `that'\,''.

Dashes come in three sizes: an
       intra-word
dash, a medium dash for number ranges like
       1--2,
and a punctuation
       dash---like
this.

A sentence-ending space should be larger than the
space between words within a sentence.  You
sometimes have to type special commands in
conjunction with punctuation characters to get
this right, as in the following sentence.
       Gnats, gnus, etc.\ all  % `\ ' makes an inter-word space.
       begin with G\@.         % \@ marks end-of-sentence punctuation.
You should check the spaces after periods when
reading your output to make sure you haven't
forgotten any special cases.  Generating an
ellipsis
       \ldots\               % `\ ' is needed after `\ldots' because TeX
                             % ignores spaces after command names like \ldots
                             % made from \ + letters.
                             %
                             % Note how a `%' character causes TeX to ignore
                             % the end of the input line, so these blank lines
                             % do not start a new paragraph.
                             %
with the right spacing around the periods requires
a special command.

\LaTeX\ interprets some common characters as
commands, so you must type special commands to
generate them.  These characters include the
following:
       \$ \& \% \# \{ and \}.

In printing, text is usually emphasized with an
       \emph{italic}
type style.

\begin{em}
   A long segment of text can also be emphasized
   in this way.  Text within such a segment can be
   given \emph{additional} emphasis.
\end{em}

It is sometimes necessary to prevent \LaTeX\ from
breaking a line where it might otherwise do so.
This may be at a space, as between the ``Mr.''\ and
``Jones'' in
       ``Mr.~Jones'',        % ~ produces an unbreakable interword space.
or within a word---especially when the word is a
symbol like
       \mbox{\emph{itemnum}}
that makes little sense when hyphenated across
lines.

Footnotes\footnote{This is an example of a footnote.}
pose no problem.

\LaTeX\ is good at typesetting mathematical formulas
like
       \( x-3y + z = 7 \)
or
       \( a_{1} > x^{2n} + y^{2n} > x' \)
or
       \( \ip{A}{B} = \sum_{i} a_{i} b_{i} \).
The spaces you type in a formula are
ignored.  Remember that a letter like
       $x$                   % $ ... $  and  \( ... \)  are equivalent
is a formula when it denotes a mathematical
symbol, and it should be typed as one.

\section{Displayed Text}

Text is displayed by indenting it from the left
margin.  Quotations are commonly displayed.  There
are short quotations
\begin{quote}
   This is a short quotation.  It consists of a
   single paragraph of text.  See how it is formatted.
\end{quote}
and longer ones.
\begin{quotation}
   This is a longer quotation.  It consists of two
   paragraphs of text, neither of which are
   particularly interesting.

   This is the second paragraph of the quotation.  It
   is just as dull as the first paragraph.
\end{quotation}
Another frequently-displayed structure is a list.
The following is an example of an \emph{itemized}
list.
\begin{itemize}
   \item This is the first item of an itemized list.
         Each item in the list is marked with a ``tick''.
         You don't have to worry about what kind of tick
         mark is used.

   \item This is the second item of the list.  It
         contains another list nested inside it.  The inner
         list is an \emph{enumerated} list.
         \begin{enumerate}
            \item This is the first item of an enumerated
                  list that is nested within the itemized list.

            \item This is the second item of the inner list.
                  \LaTeX\ allows you to nest lists deeper than
                  you really should.
         \end{enumerate}
         This is the rest of the second item of the outer
         list.  It is no more interesting than any other
         part of the item.
   \item This is the third item of the list.
\end{itemize}
You can even display poetry.
\begin{verse}
   There is an environment
    for verse \\             % The \\ command separates lines
   Whose features some poets % within a stanza.
   will curse.

                             % One or more blank lines separate stanzas.

   For instead of making\\
   Them do \emph{all} line breaking, \\
   It allows them to put too many words on a line when they'd rather be
   forced to be terse.
\end{verse}

Mathematical formulas may also be displayed.  A
displayed formula
is
one-line long; multiline
formulas require special formatting instructions.
   \[  \ip{\Gamma}{\psi'} = x'' + y^{2} + z_{i}^{n}\]
Don't start a paragraph with a displayed equation,
nor make one a paragraph by itself.

\end{document}               % End of document.

\section{Bench Huge Section 581}
% synthetic bench section 581
% This is a sample LaTeX input file.  (Version of 12 August 2004.)
%
% A '%' character causes TeX to ignore all remaining text on the line,
% and is used for comments like this one.

\documentclass{article}      % Specifies the document class

                             % The preamble begins here.
\title{An Example Document}  % Declares the document's title.
\author{Leslie Lamport}      % Declares the author's name.
\date{January 21, 1994}      % Deleting this command produces today's date.

\newcommand{\ip}[2]{(#1, #2)}
                             % Defines \ip{arg1}{arg2} to mean
                             % (arg1, arg2).

%\newcommand{\ip}[2]{\langle #1 | #2\rangle}
                             % This is an alternative definition of
                             % \ip that is commented out.

\begin{document}             % End of preamble and beginning of text.

\maketitle                   % Produces the title.

This is an example input file.  Comparing it with
the output it generates can show you how to
produce a simple document of your own.

\section{Ordinary Text}      % Produces section heading.  Lower-level
                             % sections are begun with similar
                             % \subsection and \subsubsection commands.

The ends  of words and sentences are marked
  by   spaces. It  doesn't matter how many
spaces    you type; one is as good as 100.  The
end of   a line counts as a space.

One   or more   blank lines denote the  end
of  a paragraph.

Since any number of consecutive spaces are treated
like a single one, the formatting of the input
file makes no difference to
      \LaTeX,                % The \LaTeX command generates the LaTeX logo.
but it makes a difference to you.  When you use
\LaTeX, making your input file as easy to read
as possible will be a great help as you write
your document and when you change it.  This sample
file shows how you can add comments to your own input
file.

Because printing is different from typewriting,
there are a number of things that you have to do
differently when preparing an input file than if
you were just typing the document directly.
Quotation marks like
       ``this''
have to be handled specially, as do quotes within
quotes:
       ``\,`this'            % \, separates the double and single quote.
        is what I just
        wrote, not  `that'\,''.

Dashes come in three sizes: an
       intra-word
dash, a medium dash for number ranges like
       1--2,
and a punctuation
       dash---like
this.

A sentence-ending space should be larger than the
space between words within a sentence.  You
sometimes have to type special commands in
conjunction with punctuation characters to get
this right, as in the following sentence.
       Gnats, gnus, etc.\ all  % `\ ' makes an inter-word space.
       begin with G\@.         % \@ marks end-of-sentence punctuation.
You should check the spaces after periods when
reading your output to make sure you haven't
forgotten any special cases.  Generating an
ellipsis
       \ldots\               % `\ ' is needed after `\ldots' because TeX
                             % ignores spaces after command names like \ldots
                             % made from \ + letters.
                             %
                             % Note how a `%' character causes TeX to ignore
                             % the end of the input line, so these blank lines
                             % do not start a new paragraph.
                             %
with the right spacing around the periods requires
a special command.

\LaTeX\ interprets some common characters as
commands, so you must type special commands to
generate them.  These characters include the
following:
       \$ \& \% \# \{ and \}.

In printing, text is usually emphasized with an
       \emph{italic}
type style.

\begin{em}
   A long segment of text can also be emphasized
   in this way.  Text within such a segment can be
   given \emph{additional} emphasis.
\end{em}

It is sometimes necessary to prevent \LaTeX\ from
breaking a line where it might otherwise do so.
This may be at a space, as between the ``Mr.''\ and
``Jones'' in
       ``Mr.~Jones'',        % ~ produces an unbreakable interword space.
or within a word---especially when the word is a
symbol like
       \mbox{\emph{itemnum}}
that makes little sense when hyphenated across
lines.

Footnotes\footnote{This is an example of a footnote.}
pose no problem.

\LaTeX\ is good at typesetting mathematical formulas
like
       \( x-3y + z = 7 \)
or
       \( a_{1} > x^{2n} + y^{2n} > x' \)
or
       \( \ip{A}{B} = \sum_{i} a_{i} b_{i} \).
The spaces you type in a formula are
ignored.  Remember that a letter like
       $x$                   % $ ... $  and  \( ... \)  are equivalent
is a formula when it denotes a mathematical
symbol, and it should be typed as one.

\section{Displayed Text}

Text is displayed by indenting it from the left
margin.  Quotations are commonly displayed.  There
are short quotations
\begin{quote}
   This is a short quotation.  It consists of a
   single paragraph of text.  See how it is formatted.
\end{quote}
and longer ones.
\begin{quotation}
   This is a longer quotation.  It consists of two
   paragraphs of text, neither of which are
   particularly interesting.

   This is the second paragraph of the quotation.  It
   is just as dull as the first paragraph.
\end{quotation}
Another frequently-displayed structure is a list.
The following is an example of an \emph{itemized}
list.
\begin{itemize}
   \item This is the first item of an itemized list.
         Each item in the list is marked with a ``tick''.
         You don't have to worry about what kind of tick
         mark is used.

   \item This is the second item of the list.  It
         contains another list nested inside it.  The inner
         list is an \emph{enumerated} list.
         \begin{enumerate}
            \item This is the first item of an enumerated
                  list that is nested within the itemized list.

            \item This is the second item of the inner list.
                  \LaTeX\ allows you to nest lists deeper than
                  you really should.
         \end{enumerate}
         This is the rest of the second item of the outer
         list.  It is no more interesting than any other
         part of the item.
   \item This is the third item of the list.
\end{itemize}
You can even display poetry.
\begin{verse}
   There is an environment
    for verse \\             % The \\ command separates lines
   Whose features some poets % within a stanza.
   will curse.

                             % One or more blank lines separate stanzas.

   For instead of making\\
   Them do \emph{all} line breaking, \\
   It allows them to put too many words on a line when they'd rather be
   forced to be terse.
\end{verse}

Mathematical formulas may also be displayed.  A
displayed formula
is
one-line long; multiline
formulas require special formatting instructions.
   \[  \ip{\Gamma}{\psi'} = x'' + y^{2} + z_{i}^{n}\]
Don't start a paragraph with a displayed equation,
nor make one a paragraph by itself.

\end{document}               % End of document.

\section{Bench Huge Section 582}
% synthetic bench section 582
% This is a sample LaTeX input file.  (Version of 12 August 2004.)
%
% A '%' character causes TeX to ignore all remaining text on the line,
% and is used for comments like this one.

\documentclass{article}      % Specifies the document class

                             % The preamble begins here.
\title{An Example Document}  % Declares the document's title.
\author{Leslie Lamport}      % Declares the author's name.
\date{January 21, 1994}      % Deleting this command produces today's date.

\newcommand{\ip}[2]{(#1, #2)}
                             % Defines \ip{arg1}{arg2} to mean
                             % (arg1, arg2).

%\newcommand{\ip}[2]{\langle #1 | #2\rangle}
                             % This is an alternative definition of
                             % \ip that is commented out.

\begin{document}             % End of preamble and beginning of text.

\maketitle                   % Produces the title.

This is an example input file.  Comparing it with
the output it generates can show you how to
produce a simple document of your own.

\section{Ordinary Text}      % Produces section heading.  Lower-level
                             % sections are begun with similar
                             % \subsection and \subsubsection commands.

The ends  of words and sentences are marked
  by   spaces. It  doesn't matter how many
spaces    you type; one is as good as 100.  The
end of   a line counts as a space.

One   or more   blank lines denote the  end
of  a paragraph.

Since any number of consecutive spaces are treated
like a single one, the formatting of the input
file makes no difference to
      \LaTeX,                % The \LaTeX command generates the LaTeX logo.
but it makes a difference to you.  When you use
\LaTeX, making your input file as easy to read
as possible will be a great help as you write
your document and when you change it.  This sample
file shows how you can add comments to your own input
file.

Because printing is different from typewriting,
there are a number of things that you have to do
differently when preparing an input file than if
you were just typing the document directly.
Quotation marks like
       ``this''
have to be handled specially, as do quotes within
quotes:
       ``\,`this'            % \, separates the double and single quote.
        is what I just
        wrote, not  `that'\,''.

Dashes come in three sizes: an
       intra-word
dash, a medium dash for number ranges like
       1--2,
and a punctuation
       dash---like
this.

A sentence-ending space should be larger than the
space between words within a sentence.  You
sometimes have to type special commands in
conjunction with punctuation characters to get
this right, as in the following sentence.
       Gnats, gnus, etc.\ all  % `\ ' makes an inter-word space.
       begin with G\@.         % \@ marks end-of-sentence punctuation.
You should check the spaces after periods when
reading your output to make sure you haven't
forgotten any special cases.  Generating an
ellipsis
       \ldots\               % `\ ' is needed after `\ldots' because TeX
                             % ignores spaces after command names like \ldots
                             % made from \ + letters.
                             %
                             % Note how a `%' character causes TeX to ignore
                             % the end of the input line, so these blank lines
                             % do not start a new paragraph.
                             %
with the right spacing around the periods requires
a special command.

\LaTeX\ interprets some common characters as
commands, so you must type special commands to
generate them.  These characters include the
following:
       \$ \& \% \# \{ and \}.

In printing, text is usually emphasized with an
       \emph{italic}
type style.

\begin{em}
   A long segment of text can also be emphasized
   in this way.  Text within such a segment can be
   given \emph{additional} emphasis.
\end{em}

It is sometimes necessary to prevent \LaTeX\ from
breaking a line where it might otherwise do so.
This may be at a space, as between the ``Mr.''\ and
``Jones'' in
       ``Mr.~Jones'',        % ~ produces an unbreakable interword space.
or within a word---especially when the word is a
symbol like
       \mbox{\emph{itemnum}}
that makes little sense when hyphenated across
lines.

Footnotes\footnote{This is an example of a footnote.}
pose no problem.

\LaTeX\ is good at typesetting mathematical formulas
like
       \( x-3y + z = 7 \)
or
       \( a_{1} > x^{2n} + y^{2n} > x' \)
or
       \( \ip{A}{B} = \sum_{i} a_{i} b_{i} \).
The spaces you type in a formula are
ignored.  Remember that a letter like
       $x$                   % $ ... $  and  \( ... \)  are equivalent
is a formula when it denotes a mathematical
symbol, and it should be typed as one.

\section{Displayed Text}

Text is displayed by indenting it from the left
margin.  Quotations are commonly displayed.  There
are short quotations
\begin{quote}
   This is a short quotation.  It consists of a
   single paragraph of text.  See how it is formatted.
\end{quote}
and longer ones.
\begin{quotation}
   This is a longer quotation.  It consists of two
   paragraphs of text, neither of which are
   particularly interesting.

   This is the second paragraph of the quotation.  It
   is just as dull as the first paragraph.
\end{quotation}
Another frequently-displayed structure is a list.
The following is an example of an \emph{itemized}
list.
\begin{itemize}
   \item This is the first item of an itemized list.
         Each item in the list is marked with a ``tick''.
         You don't have to worry about what kind of tick
         mark is used.

   \item This is the second item of the list.  It
         contains another list nested inside it.  The inner
         list is an \emph{enumerated} list.
         \begin{enumerate}
            \item This is the first item of an enumerated
                  list that is nested within the itemized list.

            \item This is the second item of the inner list.
                  \LaTeX\ allows you to nest lists deeper than
                  you really should.
         \end{enumerate}
         This is the rest of the second item of the outer
         list.  It is no more interesting than any other
         part of the item.
   \item This is the third item of the list.
\end{itemize}
You can even display poetry.
\begin{verse}
   There is an environment
    for verse \\             % The \\ command separates lines
   Whose features some poets % within a stanza.
   will curse.

                             % One or more blank lines separate stanzas.

   For instead of making\\
   Them do \emph{all} line breaking, \\
   It allows them to put too many words on a line when they'd rather be
   forced to be terse.
\end{verse}

Mathematical formulas may also be displayed.  A
displayed formula
is
one-line long; multiline
formulas require special formatting instructions.
   \[  \ip{\Gamma}{\psi'} = x'' + y^{2} + z_{i}^{n}\]
Don't start a paragraph with a displayed equation,
nor make one a paragraph by itself.

\end{document}               % End of document.

\section{Bench Huge Section 583}
% synthetic bench section 583
% This is a sample LaTeX input file.  (Version of 12 August 2004.)
%
% A '%' character causes TeX to ignore all remaining text on the line,
% and is used for comments like this one.

\documentclass{article}      % Specifies the document class

                             % The preamble begins here.
\title{An Example Document}  % Declares the document's title.
\author{Leslie Lamport}      % Declares the author's name.
\date{January 21, 1994}      % Deleting this command produces today's date.

\newcommand{\ip}[2]{(#1, #2)}
                             % Defines \ip{arg1}{arg2} to mean
                             % (arg1, arg2).

%\newcommand{\ip}[2]{\langle #1 | #2\rangle}
                             % This is an alternative definition of
                             % \ip that is commented out.

\begin{document}             % End of preamble and beginning of text.

\maketitle                   % Produces the title.

This is an example input file.  Comparing it with
the output it generates can show you how to
produce a simple document of your own.

\section{Ordinary Text}      % Produces section heading.  Lower-level
                             % sections are begun with similar
                             % \subsection and \subsubsection commands.

The ends  of words and sentences are marked
  by   spaces. It  doesn't matter how many
spaces    you type; one is as good as 100.  The
end of   a line counts as a space.

One   or more   blank lines denote the  end
of  a paragraph.

Since any number of consecutive spaces are treated
like a single one, the formatting of the input
file makes no difference to
      \LaTeX,                % The \LaTeX command generates the LaTeX logo.
but it makes a difference to you.  When you use
\LaTeX, making your input file as easy to read
as possible will be a great help as you write
your document and when you change it.  This sample
file shows how you can add comments to your own input
file.

Because printing is different from typewriting,
there are a number of things that you have to do
differently when preparing an input file than if
you were just typing the document directly.
Quotation marks like
       ``this''
have to be handled specially, as do quotes within
quotes:
       ``\,`this'            % \, separates the double and single quote.
        is what I just
        wrote, not  `that'\,''.

Dashes come in three sizes: an
       intra-word
dash, a medium dash for number ranges like
       1--2,
and a punctuation
       dash---like
this.

A sentence-ending space should be larger than the
space between words within a sentence.  You
sometimes have to type special commands in
conjunction with punctuation characters to get
this right, as in the following sentence.
       Gnats, gnus, etc.\ all  % `\ ' makes an inter-word space.
       begin with G\@.         % \@ marks end-of-sentence punctuation.
You should check the spaces after periods when
reading your output to make sure you haven't
forgotten any special cases.  Generating an
ellipsis
       \ldots\               % `\ ' is needed after `\ldots' because TeX
                             % ignores spaces after command names like \ldots
                             % made from \ + letters.
                             %
                             % Note how a `%' character causes TeX to ignore
                             % the end of the input line, so these blank lines
                             % do not start a new paragraph.
                             %
with the right spacing around the periods requires
a special command.

\LaTeX\ interprets some common characters as
commands, so you must type special commands to
generate them.  These characters include the
following:
       \$ \& \% \# \{ and \}.

In printing, text is usually emphasized with an
       \emph{italic}
type style.

\begin{em}
   A long segment of text can also be emphasized
   in this way.  Text within such a segment can be
   given \emph{additional} emphasis.
\end{em}

It is sometimes necessary to prevent \LaTeX\ from
breaking a line where it might otherwise do so.
This may be at a space, as between the ``Mr.''\ and
``Jones'' in
       ``Mr.~Jones'',        % ~ produces an unbreakable interword space.
or within a word---especially when the word is a
symbol like
       \mbox{\emph{itemnum}}
that makes little sense when hyphenated across
lines.

Footnotes\footnote{This is an example of a footnote.}
pose no problem.

\LaTeX\ is good at typesetting mathematical formulas
like
       \( x-3y + z = 7 \)
or
       \( a_{1} > x^{2n} + y^{2n} > x' \)
or
       \( \ip{A}{B} = \sum_{i} a_{i} b_{i} \).
The spaces you type in a formula are
ignored.  Remember that a letter like
       $x$                   % $ ... $  and  \( ... \)  are equivalent
is a formula when it denotes a mathematical
symbol, and it should be typed as one.

\section{Displayed Text}

Text is displayed by indenting it from the left
margin.  Quotations are commonly displayed.  There
are short quotations
\begin{quote}
   This is a short quotation.  It consists of a
   single paragraph of text.  See how it is formatted.
\end{quote}
and longer ones.
\begin{quotation}
   This is a longer quotation.  It consists of two
   paragraphs of text, neither of which are
   particularly interesting.

   This is the second paragraph of the quotation.  It
   is just as dull as the first paragraph.
\end{quotation}
Another frequently-displayed structure is a list.
The following is an example of an \emph{itemized}
list.
\begin{itemize}
   \item This is the first item of an itemized list.
         Each item in the list is marked with a ``tick''.
         You don't have to worry about what kind of tick
         mark is used.

   \item This is the second item of the list.  It
         contains another list nested inside it.  The inner
         list is an \emph{enumerated} list.
         \begin{enumerate}
            \item This is the first item of an enumerated
                  list that is nested within the itemized list.

            \item This is the second item of the inner list.
                  \LaTeX\ allows you to nest lists deeper than
                  you really should.
         \end{enumerate}
         This is the rest of the second item of the outer
         list.  It is no more interesting than any other
         part of the item.
   \item This is the third item of the list.
\end{itemize}
You can even display poetry.
\begin{verse}
   There is an environment
    for verse \\             % The \\ command separates lines
   Whose features some poets % within a stanza.
   will curse.

                             % One or more blank lines separate stanzas.

   For instead of making\\
   Them do \emph{all} line breaking, \\
   It allows them to put too many words on a line when they'd rather be
   forced to be terse.
\end{verse}

Mathematical formulas may also be displayed.  A
displayed formula
is
one-line long; multiline
formulas require special formatting instructions.
   \[  \ip{\Gamma}{\psi'} = x'' + y^{2} + z_{i}^{n}\]
Don't start a paragraph with a displayed equation,
nor make one a paragraph by itself.

\end{document}               % End of document.

\section{Bench Huge Section 584}
% synthetic bench section 584
% This is a sample LaTeX input file.  (Version of 12 August 2004.)
%
% A '%' character causes TeX to ignore all remaining text on the line,
% and is used for comments like this one.

\documentclass{article}      % Specifies the document class

                             % The preamble begins here.
\title{An Example Document}  % Declares the document's title.
\author{Leslie Lamport}      % Declares the author's name.
\date{January 21, 1994}      % Deleting this command produces today's date.

\newcommand{\ip}[2]{(#1, #2)}
                             % Defines \ip{arg1}{arg2} to mean
                             % (arg1, arg2).

%\newcommand{\ip}[2]{\langle #1 | #2\rangle}
                             % This is an alternative definition of
                             % \ip that is commented out.

\begin{document}             % End of preamble and beginning of text.

\maketitle                   % Produces the title.

This is an example input file.  Comparing it with
the output it generates can show you how to
produce a simple document of your own.

\section{Ordinary Text}      % Produces section heading.  Lower-level
                             % sections are begun with similar
                             % \subsection and \subsubsection commands.

The ends  of words and sentences are marked
  by   spaces. It  doesn't matter how many
spaces    you type; one is as good as 100.  The
end of   a line counts as a space.

One   or more   blank lines denote the  end
of  a paragraph.

Since any number of consecutive spaces are treated
like a single one, the formatting of the input
file makes no difference to
      \LaTeX,                % The \LaTeX command generates the LaTeX logo.
but it makes a difference to you.  When you use
\LaTeX, making your input file as easy to read
as possible will be a great help as you write
your document and when you change it.  This sample
file shows how you can add comments to your own input
file.

Because printing is different from typewriting,
there are a number of things that you have to do
differently when preparing an input file than if
you were just typing the document directly.
Quotation marks like
       ``this''
have to be handled specially, as do quotes within
quotes:
       ``\,`this'            % \, separates the double and single quote.
        is what I just
        wrote, not  `that'\,''.

Dashes come in three sizes: an
       intra-word
dash, a medium dash for number ranges like
       1--2,
and a punctuation
       dash---like
this.

A sentence-ending space should be larger than the
space between words within a sentence.  You
sometimes have to type special commands in
conjunction with punctuation characters to get
this right, as in the following sentence.
       Gnats, gnus, etc.\ all  % `\ ' makes an inter-word space.
       begin with G\@.         % \@ marks end-of-sentence punctuation.
You should check the spaces after periods when
reading your output to make sure you haven't
forgotten any special cases.  Generating an
ellipsis
       \ldots\               % `\ ' is needed after `\ldots' because TeX
                             % ignores spaces after command names like \ldots
                             % made from \ + letters.
                             %
                             % Note how a `%' character causes TeX to ignore
                             % the end of the input line, so these blank lines
                             % do not start a new paragraph.
                             %
with the right spacing around the periods requires
a special command.

\LaTeX\ interprets some common characters as
commands, so you must type special commands to
generate them.  These characters include the
following:
       \$ \& \% \# \{ and \}.

In printing, text is usually emphasized with an
       \emph{italic}
type style.

\begin{em}
   A long segment of text can also be emphasized
   in this way.  Text within such a segment can be
   given \emph{additional} emphasis.
\end{em}

It is sometimes necessary to prevent \LaTeX\ from
breaking a line where it might otherwise do so.
This may be at a space, as between the ``Mr.''\ and
``Jones'' in
       ``Mr.~Jones'',        % ~ produces an unbreakable interword space.
or within a word---especially when the word is a
symbol like
       \mbox{\emph{itemnum}}
that makes little sense when hyphenated across
lines.

Footnotes\footnote{This is an example of a footnote.}
pose no problem.

\LaTeX\ is good at typesetting mathematical formulas
like
       \( x-3y + z = 7 \)
or
       \( a_{1} > x^{2n} + y^{2n} > x' \)
or
       \( \ip{A}{B} = \sum_{i} a_{i} b_{i} \).
The spaces you type in a formula are
ignored.  Remember that a letter like
       $x$                   % $ ... $  and  \( ... \)  are equivalent
is a formula when it denotes a mathematical
symbol, and it should be typed as one.

\section{Displayed Text}

Text is displayed by indenting it from the left
margin.  Quotations are commonly displayed.  There
are short quotations
\begin{quote}
   This is a short quotation.  It consists of a
   single paragraph of text.  See how it is formatted.
\end{quote}
and longer ones.
\begin{quotation}
   This is a longer quotation.  It consists of two
   paragraphs of text, neither of which are
   particularly interesting.

   This is the second paragraph of the quotation.  It
   is just as dull as the first paragraph.
\end{quotation}
Another frequently-displayed structure is a list.
The following is an example of an \emph{itemized}
list.
\begin{itemize}
   \item This is the first item of an itemized list.
         Each item in the list is marked with a ``tick''.
         You don't have to worry about what kind of tick
         mark is used.

   \item This is the second item of the list.  It
         contains another list nested inside it.  The inner
         list is an \emph{enumerated} list.
         \begin{enumerate}
            \item This is the first item of an enumerated
                  list that is nested within the itemized list.

            \item This is the second item of the inner list.
                  \LaTeX\ allows you to nest lists deeper than
                  you really should.
         \end{enumerate}
         This is the rest of the second item of the outer
         list.  It is no more interesting than any other
         part of the item.
   \item This is the third item of the list.
\end{itemize}
You can even display poetry.
\begin{verse}
   There is an environment
    for verse \\             % The \\ command separates lines
   Whose features some poets % within a stanza.
   will curse.

                             % One or more blank lines separate stanzas.

   For instead of making\\
   Them do \emph{all} line breaking, \\
   It allows them to put too many words on a line when they'd rather be
   forced to be terse.
\end{verse}

Mathematical formulas may also be displayed.  A
displayed formula
is
one-line long; multiline
formulas require special formatting instructions.
   \[  \ip{\Gamma}{\psi'} = x'' + y^{2} + z_{i}^{n}\]
Don't start a paragraph with a displayed equation,
nor make one a paragraph by itself.

\end{document}               % End of document.

\section{Bench Huge Section 585}
% synthetic bench section 585
% This is a sample LaTeX input file.  (Version of 12 August 2004.)
%
% A '%' character causes TeX to ignore all remaining text on the line,
% and is used for comments like this one.

\documentclass{article}      % Specifies the document class

                             % The preamble begins here.
\title{An Example Document}  % Declares the document's title.
\author{Leslie Lamport}      % Declares the author's name.
\date{January 21, 1994}      % Deleting this command produces today's date.

\newcommand{\ip}[2]{(#1, #2)}
                             % Defines \ip{arg1}{arg2} to mean
                             % (arg1, arg2).

%\newcommand{\ip}[2]{\langle #1 | #2\rangle}
                             % This is an alternative definition of
                             % \ip that is commented out.

\begin{document}             % End of preamble and beginning of text.

\maketitle                   % Produces the title.

This is an example input file.  Comparing it with
the output it generates can show you how to
produce a simple document of your own.

\section{Ordinary Text}      % Produces section heading.  Lower-level
                             % sections are begun with similar
                             % \subsection and \subsubsection commands.

The ends  of words and sentences are marked
  by   spaces. It  doesn't matter how many
spaces    you type; one is as good as 100.  The
end of   a line counts as a space.

One   or more   blank lines denote the  end
of  a paragraph.

Since any number of consecutive spaces are treated
like a single one, the formatting of the input
file makes no difference to
      \LaTeX,                % The \LaTeX command generates the LaTeX logo.
but it makes a difference to you.  When you use
\LaTeX, making your input file as easy to read
as possible will be a great help as you write
your document and when you change it.  This sample
file shows how you can add comments to your own input
file.

Because printing is different from typewriting,
there are a number of things that you have to do
differently when preparing an input file than if
you were just typing the document directly.
Quotation marks like
       ``this''
have to be handled specially, as do quotes within
quotes:
       ``\,`this'            % \, separates the double and single quote.
        is what I just
        wrote, not  `that'\,''.

Dashes come in three sizes: an
       intra-word
dash, a medium dash for number ranges like
       1--2,
and a punctuation
       dash---like
this.

A sentence-ending space should be larger than the
space between words within a sentence.  You
sometimes have to type special commands in
conjunction with punctuation characters to get
this right, as in the following sentence.
       Gnats, gnus, etc.\ all  % `\ ' makes an inter-word space.
       begin with G\@.         % \@ marks end-of-sentence punctuation.
You should check the spaces after periods when
reading your output to make sure you haven't
forgotten any special cases.  Generating an
ellipsis
       \ldots\               % `\ ' is needed after `\ldots' because TeX
                             % ignores spaces after command names like \ldots
                             % made from \ + letters.
                             %
                             % Note how a `%' character causes TeX to ignore
                             % the end of the input line, so these blank lines
                             % do not start a new paragraph.
                             %
with the right spacing around the periods requires
a special command.

\LaTeX\ interprets some common characters as
commands, so you must type special commands to
generate them.  These characters include the
following:
       \$ \& \% \# \{ and \}.

In printing, text is usually emphasized with an
       \emph{italic}
type style.

\begin{em}
   A long segment of text can also be emphasized
   in this way.  Text within such a segment can be
   given \emph{additional} emphasis.
\end{em}

It is sometimes necessary to prevent \LaTeX\ from
breaking a line where it might otherwise do so.
This may be at a space, as between the ``Mr.''\ and
``Jones'' in
       ``Mr.~Jones'',        % ~ produces an unbreakable interword space.
or within a word---especially when the word is a
symbol like
       \mbox{\emph{itemnum}}
that makes little sense when hyphenated across
lines.

Footnotes\footnote{This is an example of a footnote.}
pose no problem.

\LaTeX\ is good at typesetting mathematical formulas
like
       \( x-3y + z = 7 \)
or
       \( a_{1} > x^{2n} + y^{2n} > x' \)
or
       \( \ip{A}{B} = \sum_{i} a_{i} b_{i} \).
The spaces you type in a formula are
ignored.  Remember that a letter like
       $x$                   % $ ... $  and  \( ... \)  are equivalent
is a formula when it denotes a mathematical
symbol, and it should be typed as one.

\section{Displayed Text}

Text is displayed by indenting it from the left
margin.  Quotations are commonly displayed.  There
are short quotations
\begin{quote}
   This is a short quotation.  It consists of a
   single paragraph of text.  See how it is formatted.
\end{quote}
and longer ones.
\begin{quotation}
   This is a longer quotation.  It consists of two
   paragraphs of text, neither of which are
   particularly interesting.

   This is the second paragraph of the quotation.  It
   is just as dull as the first paragraph.
\end{quotation}
Another frequently-displayed structure is a list.
The following is an example of an \emph{itemized}
list.
\begin{itemize}
   \item This is the first item of an itemized list.
         Each item in the list is marked with a ``tick''.
         You don't have to worry about what kind of tick
         mark is used.

   \item This is the second item of the list.  It
         contains another list nested inside it.  The inner
         list is an \emph{enumerated} list.
         \begin{enumerate}
            \item This is the first item of an enumerated
                  list that is nested within the itemized list.

            \item This is the second item of the inner list.
                  \LaTeX\ allows you to nest lists deeper than
                  you really should.
         \end{enumerate}
         This is the rest of the second item of the outer
         list.  It is no more interesting than any other
         part of the item.
   \item This is the third item of the list.
\end{itemize}
You can even display poetry.
\begin{verse}
   There is an environment
    for verse \\             % The \\ command separates lines
   Whose features some poets % within a stanza.
   will curse.

                             % One or more blank lines separate stanzas.

   For instead of making\\
   Them do \emph{all} line breaking, \\
   It allows them to put too many words on a line when they'd rather be
   forced to be terse.
\end{verse}

Mathematical formulas may also be displayed.  A
displayed formula
is
one-line long; multiline
formulas require special formatting instructions.
   \[  \ip{\Gamma}{\psi'} = x'' + y^{2} + z_{i}^{n}\]
Don't start a paragraph with a displayed equation,
nor make one a paragraph by itself.

\end{document}               % End of document.

\section{Bench Huge Section 586}
% synthetic bench section 586
% This is a sample LaTeX input file.  (Version of 12 August 2004.)
%
% A '%' character causes TeX to ignore all remaining text on the line,
% and is used for comments like this one.

\documentclass{article}      % Specifies the document class

                             % The preamble begins here.
\title{An Example Document}  % Declares the document's title.
\author{Leslie Lamport}      % Declares the author's name.
\date{January 21, 1994}      % Deleting this command produces today's date.

\newcommand{\ip}[2]{(#1, #2)}
                             % Defines \ip{arg1}{arg2} to mean
                             % (arg1, arg2).

%\newcommand{\ip}[2]{\langle #1 | #2\rangle}
                             % This is an alternative definition of
                             % \ip that is commented out.

\begin{document}             % End of preamble and beginning of text.

\maketitle                   % Produces the title.

This is an example input file.  Comparing it with
the output it generates can show you how to
produce a simple document of your own.

\section{Ordinary Text}      % Produces section heading.  Lower-level
                             % sections are begun with similar
                             % \subsection and \subsubsection commands.

The ends  of words and sentences are marked
  by   spaces. It  doesn't matter how many
spaces    you type; one is as good as 100.  The
end of   a line counts as a space.

One   or more   blank lines denote the  end
of  a paragraph.

Since any number of consecutive spaces are treated
like a single one, the formatting of the input
file makes no difference to
      \LaTeX,                % The \LaTeX command generates the LaTeX logo.
but it makes a difference to you.  When you use
\LaTeX, making your input file as easy to read
as possible will be a great help as you write
your document and when you change it.  This sample
file shows how you can add comments to your own input
file.

Because printing is different from typewriting,
there are a number of things that you have to do
differently when preparing an input file than if
you were just typing the document directly.
Quotation marks like
       ``this''
have to be handled specially, as do quotes within
quotes:
       ``\,`this'            % \, separates the double and single quote.
        is what I just
        wrote, not  `that'\,''.

Dashes come in three sizes: an
       intra-word
dash, a medium dash for number ranges like
       1--2,
and a punctuation
       dash---like
this.

A sentence-ending space should be larger than the
space between words within a sentence.  You
sometimes have to type special commands in
conjunction with punctuation characters to get
this right, as in the following sentence.
       Gnats, gnus, etc.\ all  % `\ ' makes an inter-word space.
       begin with G\@.         % \@ marks end-of-sentence punctuation.
You should check the spaces after periods when
reading your output to make sure you haven't
forgotten any special cases.  Generating an
ellipsis
       \ldots\               % `\ ' is needed after `\ldots' because TeX
                             % ignores spaces after command names like \ldots
                             % made from \ + letters.
                             %
                             % Note how a `%' character causes TeX to ignore
                             % the end of the input line, so these blank lines
                             % do not start a new paragraph.
                             %
with the right spacing around the periods requires
a special command.

\LaTeX\ interprets some common characters as
commands, so you must type special commands to
generate them.  These characters include the
following:
       \$ \& \% \# \{ and \}.

In printing, text is usually emphasized with an
       \emph{italic}
type style.

\begin{em}
   A long segment of text can also be emphasized
   in this way.  Text within such a segment can be
   given \emph{additional} emphasis.
\end{em}

It is sometimes necessary to prevent \LaTeX\ from
breaking a line where it might otherwise do so.
This may be at a space, as between the ``Mr.''\ and
``Jones'' in
       ``Mr.~Jones'',        % ~ produces an unbreakable interword space.
or within a word---especially when the word is a
symbol like
       \mbox{\emph{itemnum}}
that makes little sense when hyphenated across
lines.

Footnotes\footnote{This is an example of a footnote.}
pose no problem.

\LaTeX\ is good at typesetting mathematical formulas
like
       \( x-3y + z = 7 \)
or
       \( a_{1} > x^{2n} + y^{2n} > x' \)
or
       \( \ip{A}{B} = \sum_{i} a_{i} b_{i} \).
The spaces you type in a formula are
ignored.  Remember that a letter like
       $x$                   % $ ... $  and  \( ... \)  are equivalent
is a formula when it denotes a mathematical
symbol, and it should be typed as one.

\section{Displayed Text}

Text is displayed by indenting it from the left
margin.  Quotations are commonly displayed.  There
are short quotations
\begin{quote}
   This is a short quotation.  It consists of a
   single paragraph of text.  See how it is formatted.
\end{quote}
and longer ones.
\begin{quotation}
   This is a longer quotation.  It consists of two
   paragraphs of text, neither of which are
   particularly interesting.

   This is the second paragraph of the quotation.  It
   is just as dull as the first paragraph.
\end{quotation}
Another frequently-displayed structure is a list.
The following is an example of an \emph{itemized}
list.
\begin{itemize}
   \item This is the first item of an itemized list.
         Each item in the list is marked with a ``tick''.
         You don't have to worry about what kind of tick
         mark is used.

   \item This is the second item of the list.  It
         contains another list nested inside it.  The inner
         list is an \emph{enumerated} list.
         \begin{enumerate}
            \item This is the first item of an enumerated
                  list that is nested within the itemized list.

            \item This is the second item of the inner list.
                  \LaTeX\ allows you to nest lists deeper than
                  you really should.
         \end{enumerate}
         This is the rest of the second item of the outer
         list.  It is no more interesting than any other
         part of the item.
   \item This is the third item of the list.
\end{itemize}
You can even display poetry.
\begin{verse}
   There is an environment
    for verse \\             % The \\ command separates lines
   Whose features some poets % within a stanza.
   will curse.

                             % One or more blank lines separate stanzas.

   For instead of making\\
   Them do \emph{all} line breaking, \\
   It allows them to put too many words on a line when they'd rather be
   forced to be terse.
\end{verse}

Mathematical formulas may also be displayed.  A
displayed formula
is
one-line long; multiline
formulas require special formatting instructions.
   \[  \ip{\Gamma}{\psi'} = x'' + y^{2} + z_{i}^{n}\]
Don't start a paragraph with a displayed equation,
nor make one a paragraph by itself.

\end{document}               % End of document.

\section{Bench Huge Section 587}
% synthetic bench section 587
% This is a sample LaTeX input file.  (Version of 12 August 2004.)
%
% A '%' character causes TeX to ignore all remaining text on the line,
% and is used for comments like this one.

\documentclass{article}      % Specifies the document class

                             % The preamble begins here.
\title{An Example Document}  % Declares the document's title.
\author{Leslie Lamport}      % Declares the author's name.
\date{January 21, 1994}      % Deleting this command produces today's date.

\newcommand{\ip}[2]{(#1, #2)}
                             % Defines \ip{arg1}{arg2} to mean
                             % (arg1, arg2).

%\newcommand{\ip}[2]{\langle #1 | #2\rangle}
                             % This is an alternative definition of
                             % \ip that is commented out.

\begin{document}             % End of preamble and beginning of text.

\maketitle                   % Produces the title.

This is an example input file.  Comparing it with
the output it generates can show you how to
produce a simple document of your own.

\section{Ordinary Text}      % Produces section heading.  Lower-level
                             % sections are begun with similar
                             % \subsection and \subsubsection commands.

The ends  of words and sentences are marked
  by   spaces. It  doesn't matter how many
spaces    you type; one is as good as 100.  The
end of   a line counts as a space.

One   or more   blank lines denote the  end
of  a paragraph.

Since any number of consecutive spaces are treated
like a single one, the formatting of the input
file makes no difference to
      \LaTeX,                % The \LaTeX command generates the LaTeX logo.
but it makes a difference to you.  When you use
\LaTeX, making your input file as easy to read
as possible will be a great help as you write
your document and when you change it.  This sample
file shows how you can add comments to your own input
file.

Because printing is different from typewriting,
there are a number of things that you have to do
differently when preparing an input file than if
you were just typing the document directly.
Quotation marks like
       ``this''
have to be handled specially, as do quotes within
quotes:
       ``\,`this'            % \, separates the double and single quote.
        is what I just
        wrote, not  `that'\,''.

Dashes come in three sizes: an
       intra-word
dash, a medium dash for number ranges like
       1--2,
and a punctuation
       dash---like
this.

A sentence-ending space should be larger than the
space between words within a sentence.  You
sometimes have to type special commands in
conjunction with punctuation characters to get
this right, as in the following sentence.
       Gnats, gnus, etc.\ all  % `\ ' makes an inter-word space.
       begin with G\@.         % \@ marks end-of-sentence punctuation.
You should check the spaces after periods when
reading your output to make sure you haven't
forgotten any special cases.  Generating an
ellipsis
       \ldots\               % `\ ' is needed after `\ldots' because TeX
                             % ignores spaces after command names like \ldots
                             % made from \ + letters.
                             %
                             % Note how a `%' character causes TeX to ignore
                             % the end of the input line, so these blank lines
                             % do not start a new paragraph.
                             %
with the right spacing around the periods requires
a special command.

\LaTeX\ interprets some common characters as
commands, so you must type special commands to
generate them.  These characters include the
following:
       \$ \& \% \# \{ and \}.

In printing, text is usually emphasized with an
       \emph{italic}
type style.

\begin{em}
   A long segment of text can also be emphasized
   in this way.  Text within such a segment can be
   given \emph{additional} emphasis.
\end{em}

It is sometimes necessary to prevent \LaTeX\ from
breaking a line where it might otherwise do so.
This may be at a space, as between the ``Mr.''\ and
``Jones'' in
       ``Mr.~Jones'',        % ~ produces an unbreakable interword space.
or within a word---especially when the word is a
symbol like
       \mbox{\emph{itemnum}}
that makes little sense when hyphenated across
lines.

Footnotes\footnote{This is an example of a footnote.}
pose no problem.

\LaTeX\ is good at typesetting mathematical formulas
like
       \( x-3y + z = 7 \)
or
       \( a_{1} > x^{2n} + y^{2n} > x' \)
or
       \( \ip{A}{B} = \sum_{i} a_{i} b_{i} \).
The spaces you type in a formula are
ignored.  Remember that a letter like
       $x$                   % $ ... $  and  \( ... \)  are equivalent
is a formula when it denotes a mathematical
symbol, and it should be typed as one.

\section{Displayed Text}

Text is displayed by indenting it from the left
margin.  Quotations are commonly displayed.  There
are short quotations
\begin{quote}
   This is a short quotation.  It consists of a
   single paragraph of text.  See how it is formatted.
\end{quote}
and longer ones.
\begin{quotation}
   This is a longer quotation.  It consists of two
   paragraphs of text, neither of which are
   particularly interesting.

   This is the second paragraph of the quotation.  It
   is just as dull as the first paragraph.
\end{quotation}
Another frequently-displayed structure is a list.
The following is an example of an \emph{itemized}
list.
\begin{itemize}
   \item This is the first item of an itemized list.
         Each item in the list is marked with a ``tick''.
         You don't have to worry about what kind of tick
         mark is used.

   \item This is the second item of the list.  It
         contains another list nested inside it.  The inner
         list is an \emph{enumerated} list.
         \begin{enumerate}
            \item This is the first item of an enumerated
                  list that is nested within the itemized list.

            \item This is the second item of the inner list.
                  \LaTeX\ allows you to nest lists deeper than
                  you really should.
         \end{enumerate}
         This is the rest of the second item of the outer
         list.  It is no more interesting than any other
         part of the item.
   \item This is the third item of the list.
\end{itemize}
You can even display poetry.
\begin{verse}
   There is an environment
    for verse \\             % The \\ command separates lines
   Whose features some poets % within a stanza.
   will curse.

                             % One or more blank lines separate stanzas.

   For instead of making\\
   Them do \emph{all} line breaking, \\
   It allows them to put too many words on a line when they'd rather be
   forced to be terse.
\end{verse}

Mathematical formulas may also be displayed.  A
displayed formula
is
one-line long; multiline
formulas require special formatting instructions.
   \[  \ip{\Gamma}{\psi'} = x'' + y^{2} + z_{i}^{n}\]
Don't start a paragraph with a displayed equation,
nor make one a paragraph by itself.

\end{document}               % End of document.

\section{Bench Huge Section 588}
% synthetic bench section 588
% This is a sample LaTeX input file.  (Version of 12 August 2004.)
%
% A '%' character causes TeX to ignore all remaining text on the line,
% and is used for comments like this one.

\documentclass{article}      % Specifies the document class

                             % The preamble begins here.
\title{An Example Document}  % Declares the document's title.
\author{Leslie Lamport}      % Declares the author's name.
\date{January 21, 1994}      % Deleting this command produces today's date.

\newcommand{\ip}[2]{(#1, #2)}
                             % Defines \ip{arg1}{arg2} to mean
                             % (arg1, arg2).

%\newcommand{\ip}[2]{\langle #1 | #2\rangle}
                             % This is an alternative definition of
                             % \ip that is commented out.

\begin{document}             % End of preamble and beginning of text.

\maketitle                   % Produces the title.

This is an example input file.  Comparing it with
the output it generates can show you how to
produce a simple document of your own.

\section{Ordinary Text}      % Produces section heading.  Lower-level
                             % sections are begun with similar
                             % \subsection and \subsubsection commands.

The ends  of words and sentences are marked
  by   spaces. It  doesn't matter how many
spaces    you type; one is as good as 100.  The
end of   a line counts as a space.

One   or more   blank lines denote the  end
of  a paragraph.

Since any number of consecutive spaces are treated
like a single one, the formatting of the input
file makes no difference to
      \LaTeX,                % The \LaTeX command generates the LaTeX logo.
but it makes a difference to you.  When you use
\LaTeX, making your input file as easy to read
as possible will be a great help as you write
your document and when you change it.  This sample
file shows how you can add comments to your own input
file.

Because printing is different from typewriting,
there are a number of things that you have to do
differently when preparing an input file than if
you were just typing the document directly.
Quotation marks like
       ``this''
have to be handled specially, as do quotes within
quotes:
       ``\,`this'            % \, separates the double and single quote.
        is what I just
        wrote, not  `that'\,''.

Dashes come in three sizes: an
       intra-word
dash, a medium dash for number ranges like
       1--2,
and a punctuation
       dash---like
this.

A sentence-ending space should be larger than the
space between words within a sentence.  You
sometimes have to type special commands in
conjunction with punctuation characters to get
this right, as in the following sentence.
       Gnats, gnus, etc.\ all  % `\ ' makes an inter-word space.
       begin with G\@.         % \@ marks end-of-sentence punctuation.
You should check the spaces after periods when
reading your output to make sure you haven't
forgotten any special cases.  Generating an
ellipsis
       \ldots\               % `\ ' is needed after `\ldots' because TeX
                             % ignores spaces after command names like \ldots
                             % made from \ + letters.
                             %
                             % Note how a `%' character causes TeX to ignore
                             % the end of the input line, so these blank lines
                             % do not start a new paragraph.
                             %
with the right spacing around the periods requires
a special command.

\LaTeX\ interprets some common characters as
commands, so you must type special commands to
generate them.  These characters include the
following:
       \$ \& \% \# \{ and \}.

In printing, text is usually emphasized with an
       \emph{italic}
type style.

\begin{em}
   A long segment of text can also be emphasized
   in this way.  Text within such a segment can be
   given \emph{additional} emphasis.
\end{em}

It is sometimes necessary to prevent \LaTeX\ from
breaking a line where it might otherwise do so.
This may be at a space, as between the ``Mr.''\ and
``Jones'' in
       ``Mr.~Jones'',        % ~ produces an unbreakable interword space.
or within a word---especially when the word is a
symbol like
       \mbox{\emph{itemnum}}
that makes little sense when hyphenated across
lines.

Footnotes\footnote{This is an example of a footnote.}
pose no problem.

\LaTeX\ is good at typesetting mathematical formulas
like
       \( x-3y + z = 7 \)
or
       \( a_{1} > x^{2n} + y^{2n} > x' \)
or
       \( \ip{A}{B} = \sum_{i} a_{i} b_{i} \).
The spaces you type in a formula are
ignored.  Remember that a letter like
       $x$                   % $ ... $  and  \( ... \)  are equivalent
is a formula when it denotes a mathematical
symbol, and it should be typed as one.

\section{Displayed Text}

Text is displayed by indenting it from the left
margin.  Quotations are commonly displayed.  There
are short quotations
\begin{quote}
   This is a short quotation.  It consists of a
   single paragraph of text.  See how it is formatted.
\end{quote}
and longer ones.
\begin{quotation}
   This is a longer quotation.  It consists of two
   paragraphs of text, neither of which are
   particularly interesting.

   This is the second paragraph of the quotation.  It
   is just as dull as the first paragraph.
\end{quotation}
Another frequently-displayed structure is a list.
The following is an example of an \emph{itemized}
list.
\begin{itemize}
   \item This is the first item of an itemized list.
         Each item in the list is marked with a ``tick''.
         You don't have to worry about what kind of tick
         mark is used.

   \item This is the second item of the list.  It
         contains another list nested inside it.  The inner
         list is an \emph{enumerated} list.
         \begin{enumerate}
            \item This is the first item of an enumerated
                  list that is nested within the itemized list.

            \item This is the second item of the inner list.
                  \LaTeX\ allows you to nest lists deeper than
                  you really should.
         \end{enumerate}
         This is the rest of the second item of the outer
         list.  It is no more interesting than any other
         part of the item.
   \item This is the third item of the list.
\end{itemize}
You can even display poetry.
\begin{verse}
   There is an environment
    for verse \\             % The \\ command separates lines
   Whose features some poets % within a stanza.
   will curse.

                             % One or more blank lines separate stanzas.

   For instead of making\\
   Them do \emph{all} line breaking, \\
   It allows them to put too many words on a line when they'd rather be
   forced to be terse.
\end{verse}

Mathematical formulas may also be displayed.  A
displayed formula
is
one-line long; multiline
formulas require special formatting instructions.
   \[  \ip{\Gamma}{\psi'} = x'' + y^{2} + z_{i}^{n}\]
Don't start a paragraph with a displayed equation,
nor make one a paragraph by itself.

\end{document}               % End of document.

\section{Bench Huge Section 589}
% synthetic bench section 589
% This is a sample LaTeX input file.  (Version of 12 August 2004.)
%
% A '%' character causes TeX to ignore all remaining text on the line,
% and is used for comments like this one.

\documentclass{article}      % Specifies the document class

                             % The preamble begins here.
\title{An Example Document}  % Declares the document's title.
\author{Leslie Lamport}      % Declares the author's name.
\date{January 21, 1994}      % Deleting this command produces today's date.

\newcommand{\ip}[2]{(#1, #2)}
                             % Defines \ip{arg1}{arg2} to mean
                             % (arg1, arg2).

%\newcommand{\ip}[2]{\langle #1 | #2\rangle}
                             % This is an alternative definition of
                             % \ip that is commented out.

\begin{document}             % End of preamble and beginning of text.

\maketitle                   % Produces the title.

This is an example input file.  Comparing it with
the output it generates can show you how to
produce a simple document of your own.

\section{Ordinary Text}      % Produces section heading.  Lower-level
                             % sections are begun with similar
                             % \subsection and \subsubsection commands.

The ends  of words and sentences are marked
  by   spaces. It  doesn't matter how many
spaces    you type; one is as good as 100.  The
end of   a line counts as a space.

One   or more   blank lines denote the  end
of  a paragraph.

Since any number of consecutive spaces are treated
like a single one, the formatting of the input
file makes no difference to
      \LaTeX,                % The \LaTeX command generates the LaTeX logo.
but it makes a difference to you.  When you use
\LaTeX, making your input file as easy to read
as possible will be a great help as you write
your document and when you change it.  This sample
file shows how you can add comments to your own input
file.

Because printing is different from typewriting,
there are a number of things that you have to do
differently when preparing an input file than if
you were just typing the document directly.
Quotation marks like
       ``this''
have to be handled specially, as do quotes within
quotes:
       ``\,`this'            % \, separates the double and single quote.
        is what I just
        wrote, not  `that'\,''.

Dashes come in three sizes: an
       intra-word
dash, a medium dash for number ranges like
       1--2,
and a punctuation
       dash---like
this.

A sentence-ending space should be larger than the
space between words within a sentence.  You
sometimes have to type special commands in
conjunction with punctuation characters to get
this right, as in the following sentence.
       Gnats, gnus, etc.\ all  % `\ ' makes an inter-word space.
       begin with G\@.         % \@ marks end-of-sentence punctuation.
You should check the spaces after periods when
reading your output to make sure you haven't
forgotten any special cases.  Generating an
ellipsis
       \ldots\               % `\ ' is needed after `\ldots' because TeX
                             % ignores spaces after command names like \ldots
                             % made from \ + letters.
                             %
                             % Note how a `%' character causes TeX to ignore
                             % the end of the input line, so these blank lines
                             % do not start a new paragraph.
                             %
with the right spacing around the periods requires
a special command.

\LaTeX\ interprets some common characters as
commands, so you must type special commands to
generate them.  These characters include the
following:
       \$ \& \% \# \{ and \}.

In printing, text is usually emphasized with an
       \emph{italic}
type style.

\begin{em}
   A long segment of text can also be emphasized
   in this way.  Text within such a segment can be
   given \emph{additional} emphasis.
\end{em}

It is sometimes necessary to prevent \LaTeX\ from
breaking a line where it might otherwise do so.
This may be at a space, as between the ``Mr.''\ and
``Jones'' in
       ``Mr.~Jones'',        % ~ produces an unbreakable interword space.
or within a word---especially when the word is a
symbol like
       \mbox{\emph{itemnum}}
that makes little sense when hyphenated across
lines.

Footnotes\footnote{This is an example of a footnote.}
pose no problem.

\LaTeX\ is good at typesetting mathematical formulas
like
       \( x-3y + z = 7 \)
or
       \( a_{1} > x^{2n} + y^{2n} > x' \)
or
       \( \ip{A}{B} = \sum_{i} a_{i} b_{i} \).
The spaces you type in a formula are
ignored.  Remember that a letter like
       $x$                   % $ ... $  and  \( ... \)  are equivalent
is a formula when it denotes a mathematical
symbol, and it should be typed as one.

\section{Displayed Text}

Text is displayed by indenting it from the left
margin.  Quotations are commonly displayed.  There
are short quotations
\begin{quote}
   This is a short quotation.  It consists of a
   single paragraph of text.  See how it is formatted.
\end{quote}
and longer ones.
\begin{quotation}
   This is a longer quotation.  It consists of two
   paragraphs of text, neither of which are
   particularly interesting.

   This is the second paragraph of the quotation.  It
   is just as dull as the first paragraph.
\end{quotation}
Another frequently-displayed structure is a list.
The following is an example of an \emph{itemized}
list.
\begin{itemize}
   \item This is the first item of an itemized list.
         Each item in the list is marked with a ``tick''.
         You don't have to worry about what kind of tick
         mark is used.

   \item This is the second item of the list.  It
         contains another list nested inside it.  The inner
         list is an \emph{enumerated} list.
         \begin{enumerate}
            \item This is the first item of an enumerated
                  list that is nested within the itemized list.

            \item This is the second item of the inner list.
                  \LaTeX\ allows you to nest lists deeper than
                  you really should.
         \end{enumerate}
         This is the rest of the second item of the outer
         list.  It is no more interesting than any other
         part of the item.
   \item This is the third item of the list.
\end{itemize}
You can even display poetry.
\begin{verse}
   There is an environment
    for verse \\             % The \\ command separates lines
   Whose features some poets % within a stanza.
   will curse.

                             % One or more blank lines separate stanzas.

   For instead of making\\
   Them do \emph{all} line breaking, \\
   It allows them to put too many words on a line when they'd rather be
   forced to be terse.
\end{verse}

Mathematical formulas may also be displayed.  A
displayed formula
is
one-line long; multiline
formulas require special formatting instructions.
   \[  \ip{\Gamma}{\psi'} = x'' + y^{2} + z_{i}^{n}\]
Don't start a paragraph with a displayed equation,
nor make one a paragraph by itself.

\end{document}               % End of document.

\section{Bench Huge Section 590}
% synthetic bench section 590
% This is a sample LaTeX input file.  (Version of 12 August 2004.)
%
% A '%' character causes TeX to ignore all remaining text on the line,
% and is used for comments like this one.

\documentclass{article}      % Specifies the document class

                             % The preamble begins here.
\title{An Example Document}  % Declares the document's title.
\author{Leslie Lamport}      % Declares the author's name.
\date{January 21, 1994}      % Deleting this command produces today's date.

\newcommand{\ip}[2]{(#1, #2)}
                             % Defines \ip{arg1}{arg2} to mean
                             % (arg1, arg2).

%\newcommand{\ip}[2]{\langle #1 | #2\rangle}
                             % This is an alternative definition of
                             % \ip that is commented out.

\begin{document}             % End of preamble and beginning of text.

\maketitle                   % Produces the title.

This is an example input file.  Comparing it with
the output it generates can show you how to
produce a simple document of your own.

\section{Ordinary Text}      % Produces section heading.  Lower-level
                             % sections are begun with similar
                             % \subsection and \subsubsection commands.

The ends  of words and sentences are marked
  by   spaces. It  doesn't matter how many
spaces    you type; one is as good as 100.  The
end of   a line counts as a space.

One   or more   blank lines denote the  end
of  a paragraph.

Since any number of consecutive spaces are treated
like a single one, the formatting of the input
file makes no difference to
      \LaTeX,                % The \LaTeX command generates the LaTeX logo.
but it makes a difference to you.  When you use
\LaTeX, making your input file as easy to read
as possible will be a great help as you write
your document and when you change it.  This sample
file shows how you can add comments to your own input
file.

Because printing is different from typewriting,
there are a number of things that you have to do
differently when preparing an input file than if
you were just typing the document directly.
Quotation marks like
       ``this''
have to be handled specially, as do quotes within
quotes:
       ``\,`this'            % \, separates the double and single quote.
        is what I just
        wrote, not  `that'\,''.

Dashes come in three sizes: an
       intra-word
dash, a medium dash for number ranges like
       1--2,
and a punctuation
       dash---like
this.

A sentence-ending space should be larger than the
space between words within a sentence.  You
sometimes have to type special commands in
conjunction with punctuation characters to get
this right, as in the following sentence.
       Gnats, gnus, etc.\ all  % `\ ' makes an inter-word space.
       begin with G\@.         % \@ marks end-of-sentence punctuation.
You should check the spaces after periods when
reading your output to make sure you haven't
forgotten any special cases.  Generating an
ellipsis
       \ldots\               % `\ ' is needed after `\ldots' because TeX
                             % ignores spaces after command names like \ldots
                             % made from \ + letters.
                             %
                             % Note how a `%' character causes TeX to ignore
                             % the end of the input line, so these blank lines
                             % do not start a new paragraph.
                             %
with the right spacing around the periods requires
a special command.

\LaTeX\ interprets some common characters as
commands, so you must type special commands to
generate them.  These characters include the
following:
       \$ \& \% \# \{ and \}.

In printing, text is usually emphasized with an
       \emph{italic}
type style.

\begin{em}
   A long segment of text can also be emphasized
   in this way.  Text within such a segment can be
   given \emph{additional} emphasis.
\end{em}

It is sometimes necessary to prevent \LaTeX\ from
breaking a line where it might otherwise do so.
This may be at a space, as between the ``Mr.''\ and
``Jones'' in
       ``Mr.~Jones'',        % ~ produces an unbreakable interword space.
or within a word---especially when the word is a
symbol like
       \mbox{\emph{itemnum}}
that makes little sense when hyphenated across
lines.

Footnotes\footnote{This is an example of a footnote.}
pose no problem.

\LaTeX\ is good at typesetting mathematical formulas
like
       \( x-3y + z = 7 \)
or
       \( a_{1} > x^{2n} + y^{2n} > x' \)
or
       \( \ip{A}{B} = \sum_{i} a_{i} b_{i} \).
The spaces you type in a formula are
ignored.  Remember that a letter like
       $x$                   % $ ... $  and  \( ... \)  are equivalent
is a formula when it denotes a mathematical
symbol, and it should be typed as one.

\section{Displayed Text}

Text is displayed by indenting it from the left
margin.  Quotations are commonly displayed.  There
are short quotations
\begin{quote}
   This is a short quotation.  It consists of a
   single paragraph of text.  See how it is formatted.
\end{quote}
and longer ones.
\begin{quotation}
   This is a longer quotation.  It consists of two
   paragraphs of text, neither of which are
   particularly interesting.

   This is the second paragraph of the quotation.  It
   is just as dull as the first paragraph.
\end{quotation}
Another frequently-displayed structure is a list.
The following is an example of an \emph{itemized}
list.
\begin{itemize}
   \item This is the first item of an itemized list.
         Each item in the list is marked with a ``tick''.
         You don't have to worry about what kind of tick
         mark is used.

   \item This is the second item of the list.  It
         contains another list nested inside it.  The inner
         list is an \emph{enumerated} list.
         \begin{enumerate}
            \item This is the first item of an enumerated
                  list that is nested within the itemized list.

            \item This is the second item of the inner list.
                  \LaTeX\ allows you to nest lists deeper than
                  you really should.
         \end{enumerate}
         This is the rest of the second item of the outer
         list.  It is no more interesting than any other
         part of the item.
   \item This is the third item of the list.
\end{itemize}
You can even display poetry.
\begin{verse}
   There is an environment
    for verse \\             % The \\ command separates lines
   Whose features some poets % within a stanza.
   will curse.

                             % One or more blank lines separate stanzas.

   For instead of making\\
   Them do \emph{all} line breaking, \\
   It allows them to put too many words on a line when they'd rather be
   forced to be terse.
\end{verse}

Mathematical formulas may also be displayed.  A
displayed formula
is
one-line long; multiline
formulas require special formatting instructions.
   \[  \ip{\Gamma}{\psi'} = x'' + y^{2} + z_{i}^{n}\]
Don't start a paragraph with a displayed equation,
nor make one a paragraph by itself.

\end{document}               % End of document.

\section{Bench Huge Section 591}
% synthetic bench section 591
% This is a sample LaTeX input file.  (Version of 12 August 2004.)
%
% A '%' character causes TeX to ignore all remaining text on the line,
% and is used for comments like this one.

\documentclass{article}      % Specifies the document class

                             % The preamble begins here.
\title{An Example Document}  % Declares the document's title.
\author{Leslie Lamport}      % Declares the author's name.
\date{January 21, 1994}      % Deleting this command produces today's date.

\newcommand{\ip}[2]{(#1, #2)}
                             % Defines \ip{arg1}{arg2} to mean
                             % (arg1, arg2).

%\newcommand{\ip}[2]{\langle #1 | #2\rangle}
                             % This is an alternative definition of
                             % \ip that is commented out.

\begin{document}             % End of preamble and beginning of text.

\maketitle                   % Produces the title.

This is an example input file.  Comparing it with
the output it generates can show you how to
produce a simple document of your own.

\section{Ordinary Text}      % Produces section heading.  Lower-level
                             % sections are begun with similar
                             % \subsection and \subsubsection commands.

The ends  of words and sentences are marked
  by   spaces. It  doesn't matter how many
spaces    you type; one is as good as 100.  The
end of   a line counts as a space.

One   or more   blank lines denote the  end
of  a paragraph.

Since any number of consecutive spaces are treated
like a single one, the formatting of the input
file makes no difference to
      \LaTeX,                % The \LaTeX command generates the LaTeX logo.
but it makes a difference to you.  When you use
\LaTeX, making your input file as easy to read
as possible will be a great help as you write
your document and when you change it.  This sample
file shows how you can add comments to your own input
file.

Because printing is different from typewriting,
there are a number of things that you have to do
differently when preparing an input file than if
you were just typing the document directly.
Quotation marks like
       ``this''
have to be handled specially, as do quotes within
quotes:
       ``\,`this'            % \, separates the double and single quote.
        is what I just
        wrote, not  `that'\,''.

Dashes come in three sizes: an
       intra-word
dash, a medium dash for number ranges like
       1--2,
and a punctuation
       dash---like
this.

A sentence-ending space should be larger than the
space between words within a sentence.  You
sometimes have to type special commands in
conjunction with punctuation characters to get
this right, as in the following sentence.
       Gnats, gnus, etc.\ all  % `\ ' makes an inter-word space.
       begin with G\@.         % \@ marks end-of-sentence punctuation.
You should check the spaces after periods when
reading your output to make sure you haven't
forgotten any special cases.  Generating an
ellipsis
       \ldots\               % `\ ' is needed after `\ldots' because TeX
                             % ignores spaces after command names like \ldots
                             % made from \ + letters.
                             %
                             % Note how a `%' character causes TeX to ignore
                             % the end of the input line, so these blank lines
                             % do not start a new paragraph.
                             %
with the right spacing around the periods requires
a special command.

\LaTeX\ interprets some common characters as
commands, so you must type special commands to
generate them.  These characters include the
following:
       \$ \& \% \# \{ and \}.

In printing, text is usually emphasized with an
       \emph{italic}
type style.

\begin{em}
   A long segment of text can also be emphasized
   in this way.  Text within such a segment can be
   given \emph{additional} emphasis.
\end{em}

It is sometimes necessary to prevent \LaTeX\ from
breaking a line where it might otherwise do so.
This may be at a space, as between the ``Mr.''\ and
``Jones'' in
       ``Mr.~Jones'',        % ~ produces an unbreakable interword space.
or within a word---especially when the word is a
symbol like
       \mbox{\emph{itemnum}}
that makes little sense when hyphenated across
lines.

Footnotes\footnote{This is an example of a footnote.}
pose no problem.

\LaTeX\ is good at typesetting mathematical formulas
like
       \( x-3y + z = 7 \)
or
       \( a_{1} > x^{2n} + y^{2n} > x' \)
or
       \( \ip{A}{B} = \sum_{i} a_{i} b_{i} \).
The spaces you type in a formula are
ignored.  Remember that a letter like
       $x$                   % $ ... $  and  \( ... \)  are equivalent
is a formula when it denotes a mathematical
symbol, and it should be typed as one.

\section{Displayed Text}

Text is displayed by indenting it from the left
margin.  Quotations are commonly displayed.  There
are short quotations
\begin{quote}
   This is a short quotation.  It consists of a
   single paragraph of text.  See how it is formatted.
\end{quote}
and longer ones.
\begin{quotation}
   This is a longer quotation.  It consists of two
   paragraphs of text, neither of which are
   particularly interesting.

   This is the second paragraph of the quotation.  It
   is just as dull as the first paragraph.
\end{quotation}
Another frequently-displayed structure is a list.
The following is an example of an \emph{itemized}
list.
\begin{itemize}
   \item This is the first item of an itemized list.
         Each item in the list is marked with a ``tick''.
         You don't have to worry about what kind of tick
         mark is used.

   \item This is the second item of the list.  It
         contains another list nested inside it.  The inner
         list is an \emph{enumerated} list.
         \begin{enumerate}
            \item This is the first item of an enumerated
                  list that is nested within the itemized list.

            \item This is the second item of the inner list.
                  \LaTeX\ allows you to nest lists deeper than
                  you really should.
         \end{enumerate}
         This is the rest of the second item of the outer
         list.  It is no more interesting than any other
         part of the item.
   \item This is the third item of the list.
\end{itemize}
You can even display poetry.
\begin{verse}
   There is an environment
    for verse \\             % The \\ command separates lines
   Whose features some poets % within a stanza.
   will curse.

                             % One or more blank lines separate stanzas.

   For instead of making\\
   Them do \emph{all} line breaking, \\
   It allows them to put too many words on a line when they'd rather be
   forced to be terse.
\end{verse}

Mathematical formulas may also be displayed.  A
displayed formula
is
one-line long; multiline
formulas require special formatting instructions.
   \[  \ip{\Gamma}{\psi'} = x'' + y^{2} + z_{i}^{n}\]
Don't start a paragraph with a displayed equation,
nor make one a paragraph by itself.

\end{document}               % End of document.

\section{Bench Huge Section 592}
% synthetic bench section 592
% This is a sample LaTeX input file.  (Version of 12 August 2004.)
%
% A '%' character causes TeX to ignore all remaining text on the line,
% and is used for comments like this one.

\documentclass{article}      % Specifies the document class

                             % The preamble begins here.
\title{An Example Document}  % Declares the document's title.
\author{Leslie Lamport}      % Declares the author's name.
\date{January 21, 1994}      % Deleting this command produces today's date.

\newcommand{\ip}[2]{(#1, #2)}
                             % Defines \ip{arg1}{arg2} to mean
                             % (arg1, arg2).

%\newcommand{\ip}[2]{\langle #1 | #2\rangle}
                             % This is an alternative definition of
                             % \ip that is commented out.

\begin{document}             % End of preamble and beginning of text.

\maketitle                   % Produces the title.

This is an example input file.  Comparing it with
the output it generates can show you how to
produce a simple document of your own.

\section{Ordinary Text}      % Produces section heading.  Lower-level
                             % sections are begun with similar
                             % \subsection and \subsubsection commands.

The ends  of words and sentences are marked
  by   spaces. It  doesn't matter how many
spaces    you type; one is as good as 100.  The
end of   a line counts as a space.

One   or more   blank lines denote the  end
of  a paragraph.

Since any number of consecutive spaces are treated
like a single one, the formatting of the input
file makes no difference to
      \LaTeX,                % The \LaTeX command generates the LaTeX logo.
but it makes a difference to you.  When you use
\LaTeX, making your input file as easy to read
as possible will be a great help as you write
your document and when you change it.  This sample
file shows how you can add comments to your own input
file.

Because printing is different from typewriting,
there are a number of things that you have to do
differently when preparing an input file than if
you were just typing the document directly.
Quotation marks like
       ``this''
have to be handled specially, as do quotes within
quotes:
       ``\,`this'            % \, separates the double and single quote.
        is what I just
        wrote, not  `that'\,''.

Dashes come in three sizes: an
       intra-word
dash, a medium dash for number ranges like
       1--2,
and a punctuation
       dash---like
this.

A sentence-ending space should be larger than the
space between words within a sentence.  You
sometimes have to type special commands in
conjunction with punctuation characters to get
this right, as in the following sentence.
       Gnats, gnus, etc.\ all  % `\ ' makes an inter-word space.
       begin with G\@.         % \@ marks end-of-sentence punctuation.
You should check the spaces after periods when
reading your output to make sure you haven't
forgotten any special cases.  Generating an
ellipsis
       \ldots\               % `\ ' is needed after `\ldots' because TeX
                             % ignores spaces after command names like \ldots
                             % made from \ + letters.
                             %
                             % Note how a `%' character causes TeX to ignore
                             % the end of the input line, so these blank lines
                             % do not start a new paragraph.
                             %
with the right spacing around the periods requires
a special command.

\LaTeX\ interprets some common characters as
commands, so you must type special commands to
generate them.  These characters include the
following:
       \$ \& \% \# \{ and \}.

In printing, text is usually emphasized with an
       \emph{italic}
type style.

\begin{em}
   A long segment of text can also be emphasized
   in this way.  Text within such a segment can be
   given \emph{additional} emphasis.
\end{em}

It is sometimes necessary to prevent \LaTeX\ from
breaking a line where it might otherwise do so.
This may be at a space, as between the ``Mr.''\ and
``Jones'' in
       ``Mr.~Jones'',        % ~ produces an unbreakable interword space.
or within a word---especially when the word is a
symbol like
       \mbox{\emph{itemnum}}
that makes little sense when hyphenated across
lines.

Footnotes\footnote{This is an example of a footnote.}
pose no problem.

\LaTeX\ is good at typesetting mathematical formulas
like
       \( x-3y + z = 7 \)
or
       \( a_{1} > x^{2n} + y^{2n} > x' \)
or
       \( \ip{A}{B} = \sum_{i} a_{i} b_{i} \).
The spaces you type in a formula are
ignored.  Remember that a letter like
       $x$                   % $ ... $  and  \( ... \)  are equivalent
is a formula when it denotes a mathematical
symbol, and it should be typed as one.

\section{Displayed Text}

Text is displayed by indenting it from the left
margin.  Quotations are commonly displayed.  There
are short quotations
\begin{quote}
   This is a short quotation.  It consists of a
   single paragraph of text.  See how it is formatted.
\end{quote}
and longer ones.
\begin{quotation}
   This is a longer quotation.  It consists of two
   paragraphs of text, neither of which are
   particularly interesting.

   This is the second paragraph of the quotation.  It
   is just as dull as the first paragraph.
\end{quotation}
Another frequently-displayed structure is a list.
The following is an example of an \emph{itemized}
list.
\begin{itemize}
   \item This is the first item of an itemized list.
         Each item in the list is marked with a ``tick''.
         You don't have to worry about what kind of tick
         mark is used.

   \item This is the second item of the list.  It
         contains another list nested inside it.  The inner
         list is an \emph{enumerated} list.
         \begin{enumerate}
            \item This is the first item of an enumerated
                  list that is nested within the itemized list.

            \item This is the second item of the inner list.
                  \LaTeX\ allows you to nest lists deeper than
                  you really should.
         \end{enumerate}
         This is the rest of the second item of the outer
         list.  It is no more interesting than any other
         part of the item.
   \item This is the third item of the list.
\end{itemize}
You can even display poetry.
\begin{verse}
   There is an environment
    for verse \\             % The \\ command separates lines
   Whose features some poets % within a stanza.
   will curse.

                             % One or more blank lines separate stanzas.

   For instead of making\\
   Them do \emph{all} line breaking, \\
   It allows them to put too many words on a line when they'd rather be
   forced to be terse.
\end{verse}

Mathematical formulas may also be displayed.  A
displayed formula
is
one-line long; multiline
formulas require special formatting instructions.
   \[  \ip{\Gamma}{\psi'} = x'' + y^{2} + z_{i}^{n}\]
Don't start a paragraph with a displayed equation,
nor make one a paragraph by itself.

\end{document}               % End of document.

\section{Bench Huge Section 593}
% synthetic bench section 593
% This is a sample LaTeX input file.  (Version of 12 August 2004.)
%
% A '%' character causes TeX to ignore all remaining text on the line,
% and is used for comments like this one.

\documentclass{article}      % Specifies the document class

                             % The preamble begins here.
\title{An Example Document}  % Declares the document's title.
\author{Leslie Lamport}      % Declares the author's name.
\date{January 21, 1994}      % Deleting this command produces today's date.

\newcommand{\ip}[2]{(#1, #2)}
                             % Defines \ip{arg1}{arg2} to mean
                             % (arg1, arg2).

%\newcommand{\ip}[2]{\langle #1 | #2\rangle}
                             % This is an alternative definition of
                             % \ip that is commented out.

\begin{document}             % End of preamble and beginning of text.

\maketitle                   % Produces the title.

This is an example input file.  Comparing it with
the output it generates can show you how to
produce a simple document of your own.

\section{Ordinary Text}      % Produces section heading.  Lower-level
                             % sections are begun with similar
                             % \subsection and \subsubsection commands.

The ends  of words and sentences are marked
  by   spaces. It  doesn't matter how many
spaces    you type; one is as good as 100.  The
end of   a line counts as a space.

One   or more   blank lines denote the  end
of  a paragraph.

Since any number of consecutive spaces are treated
like a single one, the formatting of the input
file makes no difference to
      \LaTeX,                % The \LaTeX command generates the LaTeX logo.
but it makes a difference to you.  When you use
\LaTeX, making your input file as easy to read
as possible will be a great help as you write
your document and when you change it.  This sample
file shows how you can add comments to your own input
file.

Because printing is different from typewriting,
there are a number of things that you have to do
differently when preparing an input file than if
you were just typing the document directly.
Quotation marks like
       ``this''
have to be handled specially, as do quotes within
quotes:
       ``\,`this'            % \, separates the double and single quote.
        is what I just
        wrote, not  `that'\,''.

Dashes come in three sizes: an
       intra-word
dash, a medium dash for number ranges like
       1--2,
and a punctuation
       dash---like
this.

A sentence-ending space should be larger than the
space between words within a sentence.  You
sometimes have to type special commands in
conjunction with punctuation characters to get
this right, as in the following sentence.
       Gnats, gnus, etc.\ all  % `\ ' makes an inter-word space.
       begin with G\@.         % \@ marks end-of-sentence punctuation.
You should check the spaces after periods when
reading your output to make sure you haven't
forgotten any special cases.  Generating an
ellipsis
       \ldots\               % `\ ' is needed after `\ldots' because TeX
                             % ignores spaces after command names like \ldots
                             % made from \ + letters.
                             %
                             % Note how a `%' character causes TeX to ignore
                             % the end of the input line, so these blank lines
                             % do not start a new paragraph.
                             %
with the right spacing around the periods requires
a special command.

\LaTeX\ interprets some common characters as
commands, so you must type special commands to
generate them.  These characters include the
following:
       \$ \& \% \# \{ and \}.

In printing, text is usually emphasized with an
       \emph{italic}
type style.

\begin{em}
   A long segment of text can also be emphasized
   in this way.  Text within such a segment can be
   given \emph{additional} emphasis.
\end{em}

It is sometimes necessary to prevent \LaTeX\ from
breaking a line where it might otherwise do so.
This may be at a space, as between the ``Mr.''\ and
``Jones'' in
       ``Mr.~Jones'',        % ~ produces an unbreakable interword space.
or within a word---especially when the word is a
symbol like
       \mbox{\emph{itemnum}}
that makes little sense when hyphenated across
lines.

Footnotes\footnote{This is an example of a footnote.}
pose no problem.

\LaTeX\ is good at typesetting mathematical formulas
like
       \( x-3y + z = 7 \)
or
       \( a_{1} > x^{2n} + y^{2n} > x' \)
or
       \( \ip{A}{B} = \sum_{i} a_{i} b_{i} \).
The spaces you type in a formula are
ignored.  Remember that a letter like
       $x$                   % $ ... $  and  \( ... \)  are equivalent
is a formula when it denotes a mathematical
symbol, and it should be typed as one.

\section{Displayed Text}

Text is displayed by indenting it from the left
margin.  Quotations are commonly displayed.  There
are short quotations
\begin{quote}
   This is a short quotation.  It consists of a
   single paragraph of text.  See how it is formatted.
\end{quote}
and longer ones.
\begin{quotation}
   This is a longer quotation.  It consists of two
   paragraphs of text, neither of which are
   particularly interesting.

   This is the second paragraph of the quotation.  It
   is just as dull as the first paragraph.
\end{quotation}
Another frequently-displayed structure is a list.
The following is an example of an \emph{itemized}
list.
\begin{itemize}
   \item This is the first item of an itemized list.
         Each item in the list is marked with a ``tick''.
         You don't have to worry about what kind of tick
         mark is used.

   \item This is the second item of the list.  It
         contains another list nested inside it.  The inner
         list is an \emph{enumerated} list.
         \begin{enumerate}
            \item This is the first item of an enumerated
                  list that is nested within the itemized list.

            \item This is the second item of the inner list.
                  \LaTeX\ allows you to nest lists deeper than
                  you really should.
         \end{enumerate}
         This is the rest of the second item of the outer
         list.  It is no more interesting than any other
         part of the item.
   \item This is the third item of the list.
\end{itemize}
You can even display poetry.
\begin{verse}
   There is an environment
    for verse \\             % The \\ command separates lines
   Whose features some poets % within a stanza.
   will curse.

                             % One or more blank lines separate stanzas.

   For instead of making\\
   Them do \emph{all} line breaking, \\
   It allows them to put too many words on a line when they'd rather be
   forced to be terse.
\end{verse}

Mathematical formulas may also be displayed.  A
displayed formula
is
one-line long; multiline
formulas require special formatting instructions.
   \[  \ip{\Gamma}{\psi'} = x'' + y^{2} + z_{i}^{n}\]
Don't start a paragraph with a displayed equation,
nor make one a paragraph by itself.

\end{document}               % End of document.

\section{Bench Huge Section 594}
% synthetic bench section 594
% This is a sample LaTeX input file.  (Version of 12 August 2004.)
%
% A '%' character causes TeX to ignore all remaining text on the line,
% and is used for comments like this one.

\documentclass{article}      % Specifies the document class

                             % The preamble begins here.
\title{An Example Document}  % Declares the document's title.
\author{Leslie Lamport}      % Declares the author's name.
\date{January 21, 1994}      % Deleting this command produces today's date.

\newcommand{\ip}[2]{(#1, #2)}
                             % Defines \ip{arg1}{arg2} to mean
                             % (arg1, arg2).

%\newcommand{\ip}[2]{\langle #1 | #2\rangle}
                             % This is an alternative definition of
                             % \ip that is commented out.

\begin{document}             % End of preamble and beginning of text.

\maketitle                   % Produces the title.

This is an example input file.  Comparing it with
the output it generates can show you how to
produce a simple document of your own.

\section{Ordinary Text}      % Produces section heading.  Lower-level
                             % sections are begun with similar
                             % \subsection and \subsubsection commands.

The ends  of words and sentences are marked
  by   spaces. It  doesn't matter how many
spaces    you type; one is as good as 100.  The
end of   a line counts as a space.

One   or more   blank lines denote the  end
of  a paragraph.

Since any number of consecutive spaces are treated
like a single one, the formatting of the input
file makes no difference to
      \LaTeX,                % The \LaTeX command generates the LaTeX logo.
but it makes a difference to you.  When you use
\LaTeX, making your input file as easy to read
as possible will be a great help as you write
your document and when you change it.  This sample
file shows how you can add comments to your own input
file.

Because printing is different from typewriting,
there are a number of things that you have to do
differently when preparing an input file than if
you were just typing the document directly.
Quotation marks like
       ``this''
have to be handled specially, as do quotes within
quotes:
       ``\,`this'            % \, separates the double and single quote.
        is what I just
        wrote, not  `that'\,''.

Dashes come in three sizes: an
       intra-word
dash, a medium dash for number ranges like
       1--2,
and a punctuation
       dash---like
this.

A sentence-ending space should be larger than the
space between words within a sentence.  You
sometimes have to type special commands in
conjunction with punctuation characters to get
this right, as in the following sentence.
       Gnats, gnus, etc.\ all  % `\ ' makes an inter-word space.
       begin with G\@.         % \@ marks end-of-sentence punctuation.
You should check the spaces after periods when
reading your output to make sure you haven't
forgotten any special cases.  Generating an
ellipsis
       \ldots\               % `\ ' is needed after `\ldots' because TeX
                             % ignores spaces after command names like \ldots
                             % made from \ + letters.
                             %
                             % Note how a `%' character causes TeX to ignore
                             % the end of the input line, so these blank lines
                             % do not start a new paragraph.
                             %
with the right spacing around the periods requires
a special command.

\LaTeX\ interprets some common characters as
commands, so you must type special commands to
generate them.  These characters include the
following:
       \$ \& \% \# \{ and \}.

In printing, text is usually emphasized with an
       \emph{italic}
type style.

\begin{em}
   A long segment of text can also be emphasized
   in this way.  Text within such a segment can be
   given \emph{additional} emphasis.
\end{em}

It is sometimes necessary to prevent \LaTeX\ from
breaking a line where it might otherwise do so.
This may be at a space, as between the ``Mr.''\ and
``Jones'' in
       ``Mr.~Jones'',        % ~ produces an unbreakable interword space.
or within a word---especially when the word is a
symbol like
       \mbox{\emph{itemnum}}
that makes little sense when hyphenated across
lines.

Footnotes\footnote{This is an example of a footnote.}
pose no problem.

\LaTeX\ is good at typesetting mathematical formulas
like
       \( x-3y + z = 7 \)
or
       \( a_{1} > x^{2n} + y^{2n} > x' \)
or
       \( \ip{A}{B} = \sum_{i} a_{i} b_{i} \).
The spaces you type in a formula are
ignored.  Remember that a letter like
       $x$                   % $ ... $  and  \( ... \)  are equivalent
is a formula when it denotes a mathematical
symbol, and it should be typed as one.

\section{Displayed Text}

Text is displayed by indenting it from the left
margin.  Quotations are commonly displayed.  There
are short quotations
\begin{quote}
   This is a short quotation.  It consists of a
   single paragraph of text.  See how it is formatted.
\end{quote}
and longer ones.
\begin{quotation}
   This is a longer quotation.  It consists of two
   paragraphs of text, neither of which are
   particularly interesting.

   This is the second paragraph of the quotation.  It
   is just as dull as the first paragraph.
\end{quotation}
Another frequently-displayed structure is a list.
The following is an example of an \emph{itemized}
list.
\begin{itemize}
   \item This is the first item of an itemized list.
         Each item in the list is marked with a ``tick''.
         You don't have to worry about what kind of tick
         mark is used.

   \item This is the second item of the list.  It
         contains another list nested inside it.  The inner
         list is an \emph{enumerated} list.
         \begin{enumerate}
            \item This is the first item of an enumerated
                  list that is nested within the itemized list.

            \item This is the second item of the inner list.
                  \LaTeX\ allows you to nest lists deeper than
                  you really should.
         \end{enumerate}
         This is the rest of the second item of the outer
         list.  It is no more interesting than any other
         part of the item.
   \item This is the third item of the list.
\end{itemize}
You can even display poetry.
\begin{verse}
   There is an environment
    for verse \\             % The \\ command separates lines
   Whose features some poets % within a stanza.
   will curse.

                             % One or more blank lines separate stanzas.

   For instead of making\\
   Them do \emph{all} line breaking, \\
   It allows them to put too many words on a line when they'd rather be
   forced to be terse.
\end{verse}

Mathematical formulas may also be displayed.  A
displayed formula
is
one-line long; multiline
formulas require special formatting instructions.
   \[  \ip{\Gamma}{\psi'} = x'' + y^{2} + z_{i}^{n}\]
Don't start a paragraph with a displayed equation,
nor make one a paragraph by itself.

\end{document}               % End of document.

\section{Bench Huge Section 595}
% synthetic bench section 595
% This is a sample LaTeX input file.  (Version of 12 August 2004.)
%
% A '%' character causes TeX to ignore all remaining text on the line,
% and is used for comments like this one.

\documentclass{article}      % Specifies the document class

                             % The preamble begins here.
\title{An Example Document}  % Declares the document's title.
\author{Leslie Lamport}      % Declares the author's name.
\date{January 21, 1994}      % Deleting this command produces today's date.

\newcommand{\ip}[2]{(#1, #2)}
                             % Defines \ip{arg1}{arg2} to mean
                             % (arg1, arg2).

%\newcommand{\ip}[2]{\langle #1 | #2\rangle}
                             % This is an alternative definition of
                             % \ip that is commented out.

\begin{document}             % End of preamble and beginning of text.

\maketitle                   % Produces the title.

This is an example input file.  Comparing it with
the output it generates can show you how to
produce a simple document of your own.

\section{Ordinary Text}      % Produces section heading.  Lower-level
                             % sections are begun with similar
                             % \subsection and \subsubsection commands.

The ends  of words and sentences are marked
  by   spaces. It  doesn't matter how many
spaces    you type; one is as good as 100.  The
end of   a line counts as a space.

One   or more   blank lines denote the  end
of  a paragraph.

Since any number of consecutive spaces are treated
like a single one, the formatting of the input
file makes no difference to
      \LaTeX,                % The \LaTeX command generates the LaTeX logo.
but it makes a difference to you.  When you use
\LaTeX, making your input file as easy to read
as possible will be a great help as you write
your document and when you change it.  This sample
file shows how you can add comments to your own input
file.

Because printing is different from typewriting,
there are a number of things that you have to do
differently when preparing an input file than if
you were just typing the document directly.
Quotation marks like
       ``this''
have to be handled specially, as do quotes within
quotes:
       ``\,`this'            % \, separates the double and single quote.
        is what I just
        wrote, not  `that'\,''.

Dashes come in three sizes: an
       intra-word
dash, a medium dash for number ranges like
       1--2,
and a punctuation
       dash---like
this.

A sentence-ending space should be larger than the
space between words within a sentence.  You
sometimes have to type special commands in
conjunction with punctuation characters to get
this right, as in the following sentence.
       Gnats, gnus, etc.\ all  % `\ ' makes an inter-word space.
       begin with G\@.         % \@ marks end-of-sentence punctuation.
You should check the spaces after periods when
reading your output to make sure you haven't
forgotten any special cases.  Generating an
ellipsis
       \ldots\               % `\ ' is needed after `\ldots' because TeX
                             % ignores spaces after command names like \ldots
                             % made from \ + letters.
                             %
                             % Note how a `%' character causes TeX to ignore
                             % the end of the input line, so these blank lines
                             % do not start a new paragraph.
                             %
with the right spacing around the periods requires
a special command.

\LaTeX\ interprets some common characters as
commands, so you must type special commands to
generate them.  These characters include the
following:
       \$ \& \% \# \{ and \}.

In printing, text is usually emphasized with an
       \emph{italic}
type style.

\begin{em}
   A long segment of text can also be emphasized
   in this way.  Text within such a segment can be
   given \emph{additional} emphasis.
\end{em}

It is sometimes necessary to prevent \LaTeX\ from
breaking a line where it might otherwise do so.
This may be at a space, as between the ``Mr.''\ and
``Jones'' in
       ``Mr.~Jones'',        % ~ produces an unbreakable interword space.
or within a word---especially when the word is a
symbol like
       \mbox{\emph{itemnum}}
that makes little sense when hyphenated across
lines.

Footnotes\footnote{This is an example of a footnote.}
pose no problem.

\LaTeX\ is good at typesetting mathematical formulas
like
       \( x-3y + z = 7 \)
or
       \( a_{1} > x^{2n} + y^{2n} > x' \)
or
       \( \ip{A}{B} = \sum_{i} a_{i} b_{i} \).
The spaces you type in a formula are
ignored.  Remember that a letter like
       $x$                   % $ ... $  and  \( ... \)  are equivalent
is a formula when it denotes a mathematical
symbol, and it should be typed as one.

\section{Displayed Text}

Text is displayed by indenting it from the left
margin.  Quotations are commonly displayed.  There
are short quotations
\begin{quote}
   This is a short quotation.  It consists of a
   single paragraph of text.  See how it is formatted.
\end{quote}
and longer ones.
\begin{quotation}
   This is a longer quotation.  It consists of two
   paragraphs of text, neither of which are
   particularly interesting.

   This is the second paragraph of the quotation.  It
   is just as dull as the first paragraph.
\end{quotation}
Another frequently-displayed structure is a list.
The following is an example of an \emph{itemized}
list.
\begin{itemize}
   \item This is the first item of an itemized list.
         Each item in the list is marked with a ``tick''.
         You don't have to worry about what kind of tick
         mark is used.

   \item This is the second item of the list.  It
         contains another list nested inside it.  The inner
         list is an \emph{enumerated} list.
         \begin{enumerate}
            \item This is the first item of an enumerated
                  list that is nested within the itemized list.

            \item This is the second item of the inner list.
                  \LaTeX\ allows you to nest lists deeper than
                  you really should.
         \end{enumerate}
         This is the rest of the second item of the outer
         list.  It is no more interesting than any other
         part of the item.
   \item This is the third item of the list.
\end{itemize}
You can even display poetry.
\begin{verse}
   There is an environment
    for verse \\             % The \\ command separates lines
   Whose features some poets % within a stanza.
   will curse.

                             % One or more blank lines separate stanzas.

   For instead of making\\
   Them do \emph{all} line breaking, \\
   It allows them to put too many words on a line when they'd rather be
   forced to be terse.
\end{verse}

Mathematical formulas may also be displayed.  A
displayed formula
is
one-line long; multiline
formulas require special formatting instructions.
   \[  \ip{\Gamma}{\psi'} = x'' + y^{2} + z_{i}^{n}\]
Don't start a paragraph with a displayed equation,
nor make one a paragraph by itself.

\end{document}               % End of document.

\section{Bench Huge Section 596}
% synthetic bench section 596
% This is a sample LaTeX input file.  (Version of 12 August 2004.)
%
% A '%' character causes TeX to ignore all remaining text on the line,
% and is used for comments like this one.

\documentclass{article}      % Specifies the document class

                             % The preamble begins here.
\title{An Example Document}  % Declares the document's title.
\author{Leslie Lamport}      % Declares the author's name.
\date{January 21, 1994}      % Deleting this command produces today's date.

\newcommand{\ip}[2]{(#1, #2)}
                             % Defines \ip{arg1}{arg2} to mean
                             % (arg1, arg2).

%\newcommand{\ip}[2]{\langle #1 | #2\rangle}
                             % This is an alternative definition of
                             % \ip that is commented out.

\begin{document}             % End of preamble and beginning of text.

\maketitle                   % Produces the title.

This is an example input file.  Comparing it with
the output it generates can show you how to
produce a simple document of your own.

\section{Ordinary Text}      % Produces section heading.  Lower-level
                             % sections are begun with similar
                             % \subsection and \subsubsection commands.

The ends  of words and sentences are marked
  by   spaces. It  doesn't matter how many
spaces    you type; one is as good as 100.  The
end of   a line counts as a space.

One   or more   blank lines denote the  end
of  a paragraph.

Since any number of consecutive spaces are treated
like a single one, the formatting of the input
file makes no difference to
      \LaTeX,                % The \LaTeX command generates the LaTeX logo.
but it makes a difference to you.  When you use
\LaTeX, making your input file as easy to read
as possible will be a great help as you write
your document and when you change it.  This sample
file shows how you can add comments to your own input
file.

Because printing is different from typewriting,
there are a number of things that you have to do
differently when preparing an input file than if
you were just typing the document directly.
Quotation marks like
       ``this''
have to be handled specially, as do quotes within
quotes:
       ``\,`this'            % \, separates the double and single quote.
        is what I just
        wrote, not  `that'\,''.

Dashes come in three sizes: an
       intra-word
dash, a medium dash for number ranges like
       1--2,
and a punctuation
       dash---like
this.

A sentence-ending space should be larger than the
space between words within a sentence.  You
sometimes have to type special commands in
conjunction with punctuation characters to get
this right, as in the following sentence.
       Gnats, gnus, etc.\ all  % `\ ' makes an inter-word space.
       begin with G\@.         % \@ marks end-of-sentence punctuation.
You should check the spaces after periods when
reading your output to make sure you haven't
forgotten any special cases.  Generating an
ellipsis
       \ldots\               % `\ ' is needed after `\ldots' because TeX
                             % ignores spaces after command names like \ldots
                             % made from \ + letters.
                             %
                             % Note how a `%' character causes TeX to ignore
                             % the end of the input line, so these blank lines
                             % do not start a new paragraph.
                             %
with the right spacing around the periods requires
a special command.

\LaTeX\ interprets some common characters as
commands, so you must type special commands to
generate them.  These characters include the
following:
       \$ \& \% \# \{ and \}.

In printing, text is usually emphasized with an
       \emph{italic}
type style.

\begin{em}
   A long segment of text can also be emphasized
   in this way.  Text within such a segment can be
   given \emph{additional} emphasis.
\end{em}

It is sometimes necessary to prevent \LaTeX\ from
breaking a line where it might otherwise do so.
This may be at a space, as between the ``Mr.''\ and
``Jones'' in
       ``Mr.~Jones'',        % ~ produces an unbreakable interword space.
or within a word---especially when the word is a
symbol like
       \mbox{\emph{itemnum}}
that makes little sense when hyphenated across
lines.

Footnotes\footnote{This is an example of a footnote.}
pose no problem.

\LaTeX\ is good at typesetting mathematical formulas
like
       \( x-3y + z = 7 \)
or
       \( a_{1} > x^{2n} + y^{2n} > x' \)
or
       \( \ip{A}{B} = \sum_{i} a_{i} b_{i} \).
The spaces you type in a formula are
ignored.  Remember that a letter like
       $x$                   % $ ... $  and  \( ... \)  are equivalent
is a formula when it denotes a mathematical
symbol, and it should be typed as one.

\section{Displayed Text}

Text is displayed by indenting it from the left
margin.  Quotations are commonly displayed.  There
are short quotations
\begin{quote}
   This is a short quotation.  It consists of a
   single paragraph of text.  See how it is formatted.
\end{quote}
and longer ones.
\begin{quotation}
   This is a longer quotation.  It consists of two
   paragraphs of text, neither of which are
   particularly interesting.

   This is the second paragraph of the quotation.  It
   is just as dull as the first paragraph.
\end{quotation}
Another frequently-displayed structure is a list.
The following is an example of an \emph{itemized}
list.
\begin{itemize}
   \item This is the first item of an itemized list.
         Each item in the list is marked with a ``tick''.
         You don't have to worry about what kind of tick
         mark is used.

   \item This is the second item of the list.  It
         contains another list nested inside it.  The inner
         list is an \emph{enumerated} list.
         \begin{enumerate}
            \item This is the first item of an enumerated
                  list that is nested within the itemized list.

            \item This is the second item of the inner list.
                  \LaTeX\ allows you to nest lists deeper than
                  you really should.
         \end{enumerate}
         This is the rest of the second item of the outer
         list.  It is no more interesting than any other
         part of the item.
   \item This is the third item of the list.
\end{itemize}
You can even display poetry.
\begin{verse}
   There is an environment
    for verse \\             % The \\ command separates lines
   Whose features some poets % within a stanza.
   will curse.

                             % One or more blank lines separate stanzas.

   For instead of making\\
   Them do \emph{all} line breaking, \\
   It allows them to put too many words on a line when they'd rather be
   forced to be terse.
\end{verse}

Mathematical formulas may also be displayed.  A
displayed formula
is
one-line long; multiline
formulas require special formatting instructions.
   \[  \ip{\Gamma}{\psi'} = x'' + y^{2} + z_{i}^{n}\]
Don't start a paragraph with a displayed equation,
nor make one a paragraph by itself.

\end{document}               % End of document.

\section{Bench Huge Section 597}
% synthetic bench section 597
% This is a sample LaTeX input file.  (Version of 12 August 2004.)
%
% A '%' character causes TeX to ignore all remaining text on the line,
% and is used for comments like this one.

\documentclass{article}      % Specifies the document class

                             % The preamble begins here.
\title{An Example Document}  % Declares the document's title.
\author{Leslie Lamport}      % Declares the author's name.
\date{January 21, 1994}      % Deleting this command produces today's date.

\newcommand{\ip}[2]{(#1, #2)}
                             % Defines \ip{arg1}{arg2} to mean
                             % (arg1, arg2).

%\newcommand{\ip}[2]{\langle #1 | #2\rangle}
                             % This is an alternative definition of
                             % \ip that is commented out.

\begin{document}             % End of preamble and beginning of text.

\maketitle                   % Produces the title.

This is an example input file.  Comparing it with
the output it generates can show you how to
produce a simple document of your own.

\section{Ordinary Text}      % Produces section heading.  Lower-level
                             % sections are begun with similar
                             % \subsection and \subsubsection commands.

The ends  of words and sentences are marked
  by   spaces. It  doesn't matter how many
spaces    you type; one is as good as 100.  The
end of   a line counts as a space.

One   or more   blank lines denote the  end
of  a paragraph.

Since any number of consecutive spaces are treated
like a single one, the formatting of the input
file makes no difference to
      \LaTeX,                % The \LaTeX command generates the LaTeX logo.
but it makes a difference to you.  When you use
\LaTeX, making your input file as easy to read
as possible will be a great help as you write
your document and when you change it.  This sample
file shows how you can add comments to your own input
file.

Because printing is different from typewriting,
there are a number of things that you have to do
differently when preparing an input file than if
you were just typing the document directly.
Quotation marks like
       ``this''
have to be handled specially, as do quotes within
quotes:
       ``\,`this'            % \, separates the double and single quote.
        is what I just
        wrote, not  `that'\,''.

Dashes come in three sizes: an
       intra-word
dash, a medium dash for number ranges like
       1--2,
and a punctuation
       dash---like
this.

A sentence-ending space should be larger than the
space between words within a sentence.  You
sometimes have to type special commands in
conjunction with punctuation characters to get
this right, as in the following sentence.
       Gnats, gnus, etc.\ all  % `\ ' makes an inter-word space.
       begin with G\@.         % \@ marks end-of-sentence punctuation.
You should check the spaces after periods when
reading your output to make sure you haven't
forgotten any special cases.  Generating an
ellipsis
       \ldots\               % `\ ' is needed after `\ldots' because TeX
                             % ignores spaces after command names like \ldots
                             % made from \ + letters.
                             %
                             % Note how a `%' character causes TeX to ignore
                             % the end of the input line, so these blank lines
                             % do not start a new paragraph.
                             %
with the right spacing around the periods requires
a special command.

\LaTeX\ interprets some common characters as
commands, so you must type special commands to
generate them.  These characters include the
following:
       \$ \& \% \# \{ and \}.

In printing, text is usually emphasized with an
       \emph{italic}
type style.

\begin{em}
   A long segment of text can also be emphasized
   in this way.  Text within such a segment can be
   given \emph{additional} emphasis.
\end{em}

It is sometimes necessary to prevent \LaTeX\ from
breaking a line where it might otherwise do so.
This may be at a space, as between the ``Mr.''\ and
``Jones'' in
       ``Mr.~Jones'',        % ~ produces an unbreakable interword space.
or within a word---especially when the word is a
symbol like
       \mbox{\emph{itemnum}}
that makes little sense when hyphenated across
lines.

Footnotes\footnote{This is an example of a footnote.}
pose no problem.

\LaTeX\ is good at typesetting mathematical formulas
like
       \( x-3y + z = 7 \)
or
       \( a_{1} > x^{2n} + y^{2n} > x' \)
or
       \( \ip{A}{B} = \sum_{i} a_{i} b_{i} \).
The spaces you type in a formula are
ignored.  Remember that a letter like
       $x$                   % $ ... $  and  \( ... \)  are equivalent
is a formula when it denotes a mathematical
symbol, and it should be typed as one.

\section{Displayed Text}

Text is displayed by indenting it from the left
margin.  Quotations are commonly displayed.  There
are short quotations
\begin{quote}
   This is a short quotation.  It consists of a
   single paragraph of text.  See how it is formatted.
\end{quote}
and longer ones.
\begin{quotation}
   This is a longer quotation.  It consists of two
   paragraphs of text, neither of which are
   particularly interesting.

   This is the second paragraph of the quotation.  It
   is just as dull as the first paragraph.
\end{quotation}
Another frequently-displayed structure is a list.
The following is an example of an \emph{itemized}
list.
\begin{itemize}
   \item This is the first item of an itemized list.
         Each item in the list is marked with a ``tick''.
         You don't have to worry about what kind of tick
         mark is used.

   \item This is the second item of the list.  It
         contains another list nested inside it.  The inner
         list is an \emph{enumerated} list.
         \begin{enumerate}
            \item This is the first item of an enumerated
                  list that is nested within the itemized list.

            \item This is the second item of the inner list.
                  \LaTeX\ allows you to nest lists deeper than
                  you really should.
         \end{enumerate}
         This is the rest of the second item of the outer
         list.  It is no more interesting than any other
         part of the item.
   \item This is the third item of the list.
\end{itemize}
You can even display poetry.
\begin{verse}
   There is an environment
    for verse \\             % The \\ command separates lines
   Whose features some poets % within a stanza.
   will curse.

                             % One or more blank lines separate stanzas.

   For instead of making\\
   Them do \emph{all} line breaking, \\
   It allows them to put too many words on a line when they'd rather be
   forced to be terse.
\end{verse}

Mathematical formulas may also be displayed.  A
displayed formula
is
one-line long; multiline
formulas require special formatting instructions.
   \[  \ip{\Gamma}{\psi'} = x'' + y^{2} + z_{i}^{n}\]
Don't start a paragraph with a displayed equation,
nor make one a paragraph by itself.

\end{document}               % End of document.

\section{Bench Huge Section 598}
% synthetic bench section 598
% This is a sample LaTeX input file.  (Version of 12 August 2004.)
%
% A '%' character causes TeX to ignore all remaining text on the line,
% and is used for comments like this one.

\documentclass{article}      % Specifies the document class

                             % The preamble begins here.
\title{An Example Document}  % Declares the document's title.
\author{Leslie Lamport}      % Declares the author's name.
\date{January 21, 1994}      % Deleting this command produces today's date.

\newcommand{\ip}[2]{(#1, #2)}
                             % Defines \ip{arg1}{arg2} to mean
                             % (arg1, arg2).

%\newcommand{\ip}[2]{\langle #1 | #2\rangle}
                             % This is an alternative definition of
                             % \ip that is commented out.

\begin{document}             % End of preamble and beginning of text.

\maketitle                   % Produces the title.

This is an example input file.  Comparing it with
the output it generates can show you how to
produce a simple document of your own.

\section{Ordinary Text}      % Produces section heading.  Lower-level
                             % sections are begun with similar
                             % \subsection and \subsubsection commands.

The ends  of words and sentences are marked
  by   spaces. It  doesn't matter how many
spaces    you type; one is as good as 100.  The
end of   a line counts as a space.

One   or more   blank lines denote the  end
of  a paragraph.

Since any number of consecutive spaces are treated
like a single one, the formatting of the input
file makes no difference to
      \LaTeX,                % The \LaTeX command generates the LaTeX logo.
but it makes a difference to you.  When you use
\LaTeX, making your input file as easy to read
as possible will be a great help as you write
your document and when you change it.  This sample
file shows how you can add comments to your own input
file.

Because printing is different from typewriting,
there are a number of things that you have to do
differently when preparing an input file than if
you were just typing the document directly.
Quotation marks like
       ``this''
have to be handled specially, as do quotes within
quotes:
       ``\,`this'            % \, separates the double and single quote.
        is what I just
        wrote, not  `that'\,''.

Dashes come in three sizes: an
       intra-word
dash, a medium dash for number ranges like
       1--2,
and a punctuation
       dash---like
this.

A sentence-ending space should be larger than the
space between words within a sentence.  You
sometimes have to type special commands in
conjunction with punctuation characters to get
this right, as in the following sentence.
       Gnats, gnus, etc.\ all  % `\ ' makes an inter-word space.
       begin with G\@.         % \@ marks end-of-sentence punctuation.
You should check the spaces after periods when
reading your output to make sure you haven't
forgotten any special cases.  Generating an
ellipsis
       \ldots\               % `\ ' is needed after `\ldots' because TeX
                             % ignores spaces after command names like \ldots
                             % made from \ + letters.
                             %
                             % Note how a `%' character causes TeX to ignore
                             % the end of the input line, so these blank lines
                             % do not start a new paragraph.
                             %
with the right spacing around the periods requires
a special command.

\LaTeX\ interprets some common characters as
commands, so you must type special commands to
generate them.  These characters include the
following:
       \$ \& \% \# \{ and \}.

In printing, text is usually emphasized with an
       \emph{italic}
type style.

\begin{em}
   A long segment of text can also be emphasized
   in this way.  Text within such a segment can be
   given \emph{additional} emphasis.
\end{em}

It is sometimes necessary to prevent \LaTeX\ from
breaking a line where it might otherwise do so.
This may be at a space, as between the ``Mr.''\ and
``Jones'' in
       ``Mr.~Jones'',        % ~ produces an unbreakable interword space.
or within a word---especially when the word is a
symbol like
       \mbox{\emph{itemnum}}
that makes little sense when hyphenated across
lines.

Footnotes\footnote{This is an example of a footnote.}
pose no problem.

\LaTeX\ is good at typesetting mathematical formulas
like
       \( x-3y + z = 7 \)
or
       \( a_{1} > x^{2n} + y^{2n} > x' \)
or
       \( \ip{A}{B} = \sum_{i} a_{i} b_{i} \).
The spaces you type in a formula are
ignored.  Remember that a letter like
       $x$                   % $ ... $  and  \( ... \)  are equivalent
is a formula when it denotes a mathematical
symbol, and it should be typed as one.

\section{Displayed Text}

Text is displayed by indenting it from the left
margin.  Quotations are commonly displayed.  There
are short quotations
\begin{quote}
   This is a short quotation.  It consists of a
   single paragraph of text.  See how it is formatted.
\end{quote}
and longer ones.
\begin{quotation}
   This is a longer quotation.  It consists of two
   paragraphs of text, neither of which are
   particularly interesting.

   This is the second paragraph of the quotation.  It
   is just as dull as the first paragraph.
\end{quotation}
Another frequently-displayed structure is a list.
The following is an example of an \emph{itemized}
list.
\begin{itemize}
   \item This is the first item of an itemized list.
         Each item in the list is marked with a ``tick''.
         You don't have to worry about what kind of tick
         mark is used.

   \item This is the second item of the list.  It
         contains another list nested inside it.  The inner
         list is an \emph{enumerated} list.
         \begin{enumerate}
            \item This is the first item of an enumerated
                  list that is nested within the itemized list.

            \item This is the second item of the inner list.
                  \LaTeX\ allows you to nest lists deeper than
                  you really should.
         \end{enumerate}
         This is the rest of the second item of the outer
         list.  It is no more interesting than any other
         part of the item.
   \item This is the third item of the list.
\end{itemize}
You can even display poetry.
\begin{verse}
   There is an environment
    for verse \\             % The \\ command separates lines
   Whose features some poets % within a stanza.
   will curse.

                             % One or more blank lines separate stanzas.

   For instead of making\\
   Them do \emph{all} line breaking, \\
   It allows them to put too many words on a line when they'd rather be
   forced to be terse.
\end{verse}

Mathematical formulas may also be displayed.  A
displayed formula
is
one-line long; multiline
formulas require special formatting instructions.
   \[  \ip{\Gamma}{\psi'} = x'' + y^{2} + z_{i}^{n}\]
Don't start a paragraph with a displayed equation,
nor make one a paragraph by itself.

\end{document}               % End of document.

\section{Bench Huge Section 599}
% synthetic bench section 599
% This is a sample LaTeX input file.  (Version of 12 August 2004.)
%
% A '%' character causes TeX to ignore all remaining text on the line,
% and is used for comments like this one.

\documentclass{article}      % Specifies the document class

                             % The preamble begins here.
\title{An Example Document}  % Declares the document's title.
\author{Leslie Lamport}      % Declares the author's name.
\date{January 21, 1994}      % Deleting this command produces today's date.

\newcommand{\ip}[2]{(#1, #2)}
                             % Defines \ip{arg1}{arg2} to mean
                             % (arg1, arg2).

%\newcommand{\ip}[2]{\langle #1 | #2\rangle}
                             % This is an alternative definition of
                             % \ip that is commented out.

\begin{document}             % End of preamble and beginning of text.

\maketitle                   % Produces the title.

This is an example input file.  Comparing it with
the output it generates can show you how to
produce a simple document of your own.

\section{Ordinary Text}      % Produces section heading.  Lower-level
                             % sections are begun with similar
                             % \subsection and \subsubsection commands.

The ends  of words and sentences are marked
  by   spaces. It  doesn't matter how many
spaces    you type; one is as good as 100.  The
end of   a line counts as a space.

One   or more   blank lines denote the  end
of  a paragraph.

Since any number of consecutive spaces are treated
like a single one, the formatting of the input
file makes no difference to
      \LaTeX,                % The \LaTeX command generates the LaTeX logo.
but it makes a difference to you.  When you use
\LaTeX, making your input file as easy to read
as possible will be a great help as you write
your document and when you change it.  This sample
file shows how you can add comments to your own input
file.

Because printing is different from typewriting,
there are a number of things that you have to do
differently when preparing an input file than if
you were just typing the document directly.
Quotation marks like
       ``this''
have to be handled specially, as do quotes within
quotes:
       ``\,`this'            % \, separates the double and single quote.
        is what I just
        wrote, not  `that'\,''.

Dashes come in three sizes: an
       intra-word
dash, a medium dash for number ranges like
       1--2,
and a punctuation
       dash---like
this.

A sentence-ending space should be larger than the
space between words within a sentence.  You
sometimes have to type special commands in
conjunction with punctuation characters to get
this right, as in the following sentence.
       Gnats, gnus, etc.\ all  % `\ ' makes an inter-word space.
       begin with G\@.         % \@ marks end-of-sentence punctuation.
You should check the spaces after periods when
reading your output to make sure you haven't
forgotten any special cases.  Generating an
ellipsis
       \ldots\               % `\ ' is needed after `\ldots' because TeX
                             % ignores spaces after command names like \ldots
                             % made from \ + letters.
                             %
                             % Note how a `%' character causes TeX to ignore
                             % the end of the input line, so these blank lines
                             % do not start a new paragraph.
                             %
with the right spacing around the periods requires
a special command.

\LaTeX\ interprets some common characters as
commands, so you must type special commands to
generate them.  These characters include the
following:
       \$ \& \% \# \{ and \}.

In printing, text is usually emphasized with an
       \emph{italic}
type style.

\begin{em}
   A long segment of text can also be emphasized
   in this way.  Text within such a segment can be
   given \emph{additional} emphasis.
\end{em}

It is sometimes necessary to prevent \LaTeX\ from
breaking a line where it might otherwise do so.
This may be at a space, as between the ``Mr.''\ and
``Jones'' in
       ``Mr.~Jones'',        % ~ produces an unbreakable interword space.
or within a word---especially when the word is a
symbol like
       \mbox{\emph{itemnum}}
that makes little sense when hyphenated across
lines.

Footnotes\footnote{This is an example of a footnote.}
pose no problem.

\LaTeX\ is good at typesetting mathematical formulas
like
       \( x-3y + z = 7 \)
or
       \( a_{1} > x^{2n} + y^{2n} > x' \)
or
       \( \ip{A}{B} = \sum_{i} a_{i} b_{i} \).
The spaces you type in a formula are
ignored.  Remember that a letter like
       $x$                   % $ ... $  and  \( ... \)  are equivalent
is a formula when it denotes a mathematical
symbol, and it should be typed as one.

\section{Displayed Text}

Text is displayed by indenting it from the left
margin.  Quotations are commonly displayed.  There
are short quotations
\begin{quote}
   This is a short quotation.  It consists of a
   single paragraph of text.  See how it is formatted.
\end{quote}
and longer ones.
\begin{quotation}
   This is a longer quotation.  It consists of two
   paragraphs of text, neither of which are
   particularly interesting.

   This is the second paragraph of the quotation.  It
   is just as dull as the first paragraph.
\end{quotation}
Another frequently-displayed structure is a list.
The following is an example of an \emph{itemized}
list.
\begin{itemize}
   \item This is the first item of an itemized list.
         Each item in the list is marked with a ``tick''.
         You don't have to worry about what kind of tick
         mark is used.

   \item This is the second item of the list.  It
         contains another list nested inside it.  The inner
         list is an \emph{enumerated} list.
         \begin{enumerate}
            \item This is the first item of an enumerated
                  list that is nested within the itemized list.

            \item This is the second item of the inner list.
                  \LaTeX\ allows you to nest lists deeper than
                  you really should.
         \end{enumerate}
         This is the rest of the second item of the outer
         list.  It is no more interesting than any other
         part of the item.
   \item This is the third item of the list.
\end{itemize}
You can even display poetry.
\begin{verse}
   There is an environment
    for verse \\             % The \\ command separates lines
   Whose features some poets % within a stanza.
   will curse.

                             % One or more blank lines separate stanzas.

   For instead of making\\
   Them do \emph{all} line breaking, \\
   It allows them to put too many words on a line when they'd rather be
   forced to be terse.
\end{verse}

Mathematical formulas may also be displayed.  A
displayed formula
is
one-line long; multiline
formulas require special formatting instructions.
   \[  \ip{\Gamma}{\psi'} = x'' + y^{2} + z_{i}^{n}\]
Don't start a paragraph with a displayed equation,
nor make one a paragraph by itself.

\end{document}               % End of document.

\section{Bench Huge Section 600}
% synthetic bench section 600
% This is a sample LaTeX input file.  (Version of 12 August 2004.)
%
% A '%' character causes TeX to ignore all remaining text on the line,
% and is used for comments like this one.

\documentclass{article}      % Specifies the document class

                             % The preamble begins here.
\title{An Example Document}  % Declares the document's title.
\author{Leslie Lamport}      % Declares the author's name.
\date{January 21, 1994}      % Deleting this command produces today's date.

\newcommand{\ip}[2]{(#1, #2)}
                             % Defines \ip{arg1}{arg2} to mean
                             % (arg1, arg2).

%\newcommand{\ip}[2]{\langle #1 | #2\rangle}
                             % This is an alternative definition of
                             % \ip that is commented out.

\begin{document}             % End of preamble and beginning of text.

\maketitle                   % Produces the title.

This is an example input file.  Comparing it with
the output it generates can show you how to
produce a simple document of your own.

\section{Ordinary Text}      % Produces section heading.  Lower-level
                             % sections are begun with similar
                             % \subsection and \subsubsection commands.

The ends  of words and sentences are marked
  by   spaces. It  doesn't matter how many
spaces    you type; one is as good as 100.  The
end of   a line counts as a space.

One   or more   blank lines denote the  end
of  a paragraph.

Since any number of consecutive spaces are treated
like a single one, the formatting of the input
file makes no difference to
      \LaTeX,                % The \LaTeX command generates the LaTeX logo.
but it makes a difference to you.  When you use
\LaTeX, making your input file as easy to read
as possible will be a great help as you write
your document and when you change it.  This sample
file shows how you can add comments to your own input
file.

Because printing is different from typewriting,
there are a number of things that you have to do
differently when preparing an input file than if
you were just typing the document directly.
Quotation marks like
       ``this''
have to be handled specially, as do quotes within
quotes:
       ``\,`this'            % \, separates the double and single quote.
        is what I just
        wrote, not  `that'\,''.

Dashes come in three sizes: an
       intra-word
dash, a medium dash for number ranges like
       1--2,
and a punctuation
       dash---like
this.

A sentence-ending space should be larger than the
space between words within a sentence.  You
sometimes have to type special commands in
conjunction with punctuation characters to get
this right, as in the following sentence.
       Gnats, gnus, etc.\ all  % `\ ' makes an inter-word space.
       begin with G\@.         % \@ marks end-of-sentence punctuation.
You should check the spaces after periods when
reading your output to make sure you haven't
forgotten any special cases.  Generating an
ellipsis
       \ldots\               % `\ ' is needed after `\ldots' because TeX
                             % ignores spaces after command names like \ldots
                             % made from \ + letters.
                             %
                             % Note how a `%' character causes TeX to ignore
                             % the end of the input line, so these blank lines
                             % do not start a new paragraph.
                             %
with the right spacing around the periods requires
a special command.

\LaTeX\ interprets some common characters as
commands, so you must type special commands to
generate them.  These characters include the
following:
       \$ \& \% \# \{ and \}.

In printing, text is usually emphasized with an
       \emph{italic}
type style.

\begin{em}
   A long segment of text can also be emphasized
   in this way.  Text within such a segment can be
   given \emph{additional} emphasis.
\end{em}

It is sometimes necessary to prevent \LaTeX\ from
breaking a line where it might otherwise do so.
This may be at a space, as between the ``Mr.''\ and
``Jones'' in
       ``Mr.~Jones'',        % ~ produces an unbreakable interword space.
or within a word---especially when the word is a
symbol like
       \mbox{\emph{itemnum}}
that makes little sense when hyphenated across
lines.

Footnotes\footnote{This is an example of a footnote.}
pose no problem.

\LaTeX\ is good at typesetting mathematical formulas
like
       \( x-3y + z = 7 \)
or
       \( a_{1} > x^{2n} + y^{2n} > x' \)
or
       \( \ip{A}{B} = \sum_{i} a_{i} b_{i} \).
The spaces you type in a formula are
ignored.  Remember that a letter like
       $x$                   % $ ... $  and  \( ... \)  are equivalent
is a formula when it denotes a mathematical
symbol, and it should be typed as one.

\section{Displayed Text}

Text is displayed by indenting it from the left
margin.  Quotations are commonly displayed.  There
are short quotations
\begin{quote}
   This is a short quotation.  It consists of a
   single paragraph of text.  See how it is formatted.
\end{quote}
and longer ones.
\begin{quotation}
   This is a longer quotation.  It consists of two
   paragraphs of text, neither of which are
   particularly interesting.

   This is the second paragraph of the quotation.  It
   is just as dull as the first paragraph.
\end{quotation}
Another frequently-displayed structure is a list.
The following is an example of an \emph{itemized}
list.
\begin{itemize}
   \item This is the first item of an itemized list.
         Each item in the list is marked with a ``tick''.
         You don't have to worry about what kind of tick
         mark is used.

   \item This is the second item of the list.  It
         contains another list nested inside it.  The inner
         list is an \emph{enumerated} list.
         \begin{enumerate}
            \item This is the first item of an enumerated
                  list that is nested within the itemized list.

            \item This is the second item of the inner list.
                  \LaTeX\ allows you to nest lists deeper than
                  you really should.
         \end{enumerate}
         This is the rest of the second item of the outer
         list.  It is no more interesting than any other
         part of the item.
   \item This is the third item of the list.
\end{itemize}
You can even display poetry.
\begin{verse}
   There is an environment
    for verse \\             % The \\ command separates lines
   Whose features some poets % within a stanza.
   will curse.

                             % One or more blank lines separate stanzas.

   For instead of making\\
   Them do \emph{all} line breaking, \\
   It allows them to put too many words on a line when they'd rather be
   forced to be terse.
\end{verse}

Mathematical formulas may also be displayed.  A
displayed formula
is
one-line long; multiline
formulas require special formatting instructions.
   \[  \ip{\Gamma}{\psi'} = x'' + y^{2} + z_{i}^{n}\]
Don't start a paragraph with a displayed equation,
nor make one a paragraph by itself.

\end{document}               % End of document.

\section{Bench Huge Section 601}
% synthetic bench section 601
% This is a sample LaTeX input file.  (Version of 12 August 2004.)
%
% A '%' character causes TeX to ignore all remaining text on the line,
% and is used for comments like this one.

\documentclass{article}      % Specifies the document class

                             % The preamble begins here.
\title{An Example Document}  % Declares the document's title.
\author{Leslie Lamport}      % Declares the author's name.
\date{January 21, 1994}      % Deleting this command produces today's date.

\newcommand{\ip}[2]{(#1, #2)}
                             % Defines \ip{arg1}{arg2} to mean
                             % (arg1, arg2).

%\newcommand{\ip}[2]{\langle #1 | #2\rangle}
                             % This is an alternative definition of
                             % \ip that is commented out.

\begin{document}             % End of preamble and beginning of text.

\maketitle                   % Produces the title.

This is an example input file.  Comparing it with
the output it generates can show you how to
produce a simple document of your own.

\section{Ordinary Text}      % Produces section heading.  Lower-level
                             % sections are begun with similar
                             % \subsection and \subsubsection commands.

The ends  of words and sentences are marked
  by   spaces. It  doesn't matter how many
spaces    you type; one is as good as 100.  The
end of   a line counts as a space.

One   or more   blank lines denote the  end
of  a paragraph.

Since any number of consecutive spaces are treated
like a single one, the formatting of the input
file makes no difference to
      \LaTeX,                % The \LaTeX command generates the LaTeX logo.
but it makes a difference to you.  When you use
\LaTeX, making your input file as easy to read
as possible will be a great help as you write
your document and when you change it.  This sample
file shows how you can add comments to your own input
file.

Because printing is different from typewriting,
there are a number of things that you have to do
differently when preparing an input file than if
you were just typing the document directly.
Quotation marks like
       ``this''
have to be handled specially, as do quotes within
quotes:
       ``\,`this'            % \, separates the double and single quote.
        is what I just
        wrote, not  `that'\,''.

Dashes come in three sizes: an
       intra-word
dash, a medium dash for number ranges like
       1--2,
and a punctuation
       dash---like
this.

A sentence-ending space should be larger than the
space between words within a sentence.  You
sometimes have to type special commands in
conjunction with punctuation characters to get
this right, as in the following sentence.
       Gnats, gnus, etc.\ all  % `\ ' makes an inter-word space.
       begin with G\@.         % \@ marks end-of-sentence punctuation.
You should check the spaces after periods when
reading your output to make sure you haven't
forgotten any special cases.  Generating an
ellipsis
       \ldots\               % `\ ' is needed after `\ldots' because TeX
                             % ignores spaces after command names like \ldots
                             % made from \ + letters.
                             %
                             % Note how a `%' character causes TeX to ignore
                             % the end of the input line, so these blank lines
                             % do not start a new paragraph.
                             %
with the right spacing around the periods requires
a special command.

\LaTeX\ interprets some common characters as
commands, so you must type special commands to
generate them.  These characters include the
following:
       \$ \& \% \# \{ and \}.

In printing, text is usually emphasized with an
       \emph{italic}
type style.

\begin{em}
   A long segment of text can also be emphasized
   in this way.  Text within such a segment can be
   given \emph{additional} emphasis.
\end{em}

It is sometimes necessary to prevent \LaTeX\ from
breaking a line where it might otherwise do so.
This may be at a space, as between the ``Mr.''\ and
``Jones'' in
       ``Mr.~Jones'',        % ~ produces an unbreakable interword space.
or within a word---especially when the word is a
symbol like
       \mbox{\emph{itemnum}}
that makes little sense when hyphenated across
lines.

Footnotes\footnote{This is an example of a footnote.}
pose no problem.

\LaTeX\ is good at typesetting mathematical formulas
like
       \( x-3y + z = 7 \)
or
       \( a_{1} > x^{2n} + y^{2n} > x' \)
or
       \( \ip{A}{B} = \sum_{i} a_{i} b_{i} \).
The spaces you type in a formula are
ignored.  Remember that a letter like
       $x$                   % $ ... $  and  \( ... \)  are equivalent
is a formula when it denotes a mathematical
symbol, and it should be typed as one.

\section{Displayed Text}

Text is displayed by indenting it from the left
margin.  Quotations are commonly displayed.  There
are short quotations
\begin{quote}
   This is a short quotation.  It consists of a
   single paragraph of text.  See how it is formatted.
\end{quote}
and longer ones.
\begin{quotation}
   This is a longer quotation.  It consists of two
   paragraphs of text, neither of which are
   particularly interesting.

   This is the second paragraph of the quotation.  It
   is just as dull as the first paragraph.
\end{quotation}
Another frequently-displayed structure is a list.
The following is an example of an \emph{itemized}
list.
\begin{itemize}
   \item This is the first item of an itemized list.
         Each item in the list is marked with a ``tick''.
         You don't have to worry about what kind of tick
         mark is used.

   \item This is the second item of the list.  It
         contains another list nested inside it.  The inner
         list is an \emph{enumerated} list.
         \begin{enumerate}
            \item This is the first item of an enumerated
                  list that is nested within the itemized list.

            \item This is the second item of the inner list.
                  \LaTeX\ allows you to nest lists deeper than
                  you really should.
         \end{enumerate}
         This is the rest of the second item of the outer
         list.  It is no more interesting than any other
         part of the item.
   \item This is the third item of the list.
\end{itemize}
You can even display poetry.
\begin{verse}
   There is an environment
    for verse \\             % The \\ command separates lines
   Whose features some poets % within a stanza.
   will curse.

                             % One or more blank lines separate stanzas.

   For instead of making\\
   Them do \emph{all} line breaking, \\
   It allows them to put too many words on a line when they'd rather be
   forced to be terse.
\end{verse}

Mathematical formulas may also be displayed.  A
displayed formula
is
one-line long; multiline
formulas require special formatting instructions.
   \[  \ip{\Gamma}{\psi'} = x'' + y^{2} + z_{i}^{n}\]
Don't start a paragraph with a displayed equation,
nor make one a paragraph by itself.

\end{document}               % End of document.

\section{Bench Huge Section 602}
% synthetic bench section 602
% This is a sample LaTeX input file.  (Version of 12 August 2004.)
%
% A '%' character causes TeX to ignore all remaining text on the line,
% and is used for comments like this one.

\documentclass{article}      % Specifies the document class

                             % The preamble begins here.
\title{An Example Document}  % Declares the document's title.
\author{Leslie Lamport}      % Declares the author's name.
\date{January 21, 1994}      % Deleting this command produces today's date.

\newcommand{\ip}[2]{(#1, #2)}
                             % Defines \ip{arg1}{arg2} to mean
                             % (arg1, arg2).

%\newcommand{\ip}[2]{\langle #1 | #2\rangle}
                             % This is an alternative definition of
                             % \ip that is commented out.

\begin{document}             % End of preamble and beginning of text.

\maketitle                   % Produces the title.

This is an example input file.  Comparing it with
the output it generates can show you how to
produce a simple document of your own.

\section{Ordinary Text}      % Produces section heading.  Lower-level
                             % sections are begun with similar
                             % \subsection and \subsubsection commands.

The ends  of words and sentences are marked
  by   spaces. It  doesn't matter how many
spaces    you type; one is as good as 100.  The
end of   a line counts as a space.

One   or more   blank lines denote the  end
of  a paragraph.

Since any number of consecutive spaces are treated
like a single one, the formatting of the input
file makes no difference to
      \LaTeX,                % The \LaTeX command generates the LaTeX logo.
but it makes a difference to you.  When you use
\LaTeX, making your input file as easy to read
as possible will be a great help as you write
your document and when you change it.  This sample
file shows how you can add comments to your own input
file.

Because printing is different from typewriting,
there are a number of things that you have to do
differently when preparing an input file than if
you were just typing the document directly.
Quotation marks like
       ``this''
have to be handled specially, as do quotes within
quotes:
       ``\,`this'            % \, separates the double and single quote.
        is what I just
        wrote, not  `that'\,''.

Dashes come in three sizes: an
       intra-word
dash, a medium dash for number ranges like
       1--2,
and a punctuation
       dash---like
this.

A sentence-ending space should be larger than the
space between words within a sentence.  You
sometimes have to type special commands in
conjunction with punctuation characters to get
this right, as in the following sentence.
       Gnats, gnus, etc.\ all  % `\ ' makes an inter-word space.
       begin with G\@.         % \@ marks end-of-sentence punctuation.
You should check the spaces after periods when
reading your output to make sure you haven't
forgotten any special cases.  Generating an
ellipsis
       \ldots\               % `\ ' is needed after `\ldots' because TeX
                             % ignores spaces after command names like \ldots
                             % made from \ + letters.
                             %
                             % Note how a `%' character causes TeX to ignore
                             % the end of the input line, so these blank lines
                             % do not start a new paragraph.
                             %
with the right spacing around the periods requires
a special command.

\LaTeX\ interprets some common characters as
commands, so you must type special commands to
generate them.  These characters include the
following:
       \$ \& \% \# \{ and \}.

In printing, text is usually emphasized with an
       \emph{italic}
type style.

\begin{em}
   A long segment of text can also be emphasized
   in this way.  Text within such a segment can be
   given \emph{additional} emphasis.
\end{em}

It is sometimes necessary to prevent \LaTeX\ from
breaking a line where it might otherwise do so.
This may be at a space, as between the ``Mr.''\ and
``Jones'' in
       ``Mr.~Jones'',        % ~ produces an unbreakable interword space.
or within a word---especially when the word is a
symbol like
       \mbox{\emph{itemnum}}
that makes little sense when hyphenated across
lines.

Footnotes\footnote{This is an example of a footnote.}
pose no problem.

\LaTeX\ is good at typesetting mathematical formulas
like
       \( x-3y + z = 7 \)
or
       \( a_{1} > x^{2n} + y^{2n} > x' \)
or
       \( \ip{A}{B} = \sum_{i} a_{i} b_{i} \).
The spaces you type in a formula are
ignored.  Remember that a letter like
       $x$                   % $ ... $  and  \( ... \)  are equivalent
is a formula when it denotes a mathematical
symbol, and it should be typed as one.

\section{Displayed Text}

Text is displayed by indenting it from the left
margin.  Quotations are commonly displayed.  There
are short quotations
\begin{quote}
   This is a short quotation.  It consists of a
   single paragraph of text.  See how it is formatted.
\end{quote}
and longer ones.
\begin{quotation}
   This is a longer quotation.  It consists of two
   paragraphs of text, neither of which are
   particularly interesting.

   This is the second paragraph of the quotation.  It
   is just as dull as the first paragraph.
\end{quotation}
Another frequently-displayed structure is a list.
The following is an example of an \emph{itemized}
list.
\begin{itemize}
   \item This is the first item of an itemized list.
         Each item in the list is marked with a ``tick''.
         You don't have to worry about what kind of tick
         mark is used.

   \item This is the second item of the list.  It
         contains another list nested inside it.  The inner
         list is an \emph{enumerated} list.
         \begin{enumerate}
            \item This is the first item of an enumerated
                  list that is nested within the itemized list.

            \item This is the second item of the inner list.
                  \LaTeX\ allows you to nest lists deeper than
                  you really should.
         \end{enumerate}
         This is the rest of the second item of the outer
         list.  It is no more interesting than any other
         part of the item.
   \item This is the third item of the list.
\end{itemize}
You can even display poetry.
\begin{verse}
   There is an environment
    for verse \\             % The \\ command separates lines
   Whose features some poets % within a stanza.
   will curse.

                             % One or more blank lines separate stanzas.

   For instead of making\\
   Them do \emph{all} line breaking, \\
   It allows them to put too many words on a line when they'd rather be
   forced to be terse.
\end{verse}

Mathematical formulas may also be displayed.  A
displayed formula
is
one-line long; multiline
formulas require special formatting instructions.
   \[  \ip{\Gamma}{\psi'} = x'' + y^{2} + z_{i}^{n}\]
Don't start a paragraph with a displayed equation,
nor make one a paragraph by itself.

\end{document}               % End of document.

\section{Bench Huge Section 603}
% synthetic bench section 603
% This is a sample LaTeX input file.  (Version of 12 August 2004.)
%
% A '%' character causes TeX to ignore all remaining text on the line,
% and is used for comments like this one.

\documentclass{article}      % Specifies the document class

                             % The preamble begins here.
\title{An Example Document}  % Declares the document's title.
\author{Leslie Lamport}      % Declares the author's name.
\date{January 21, 1994}      % Deleting this command produces today's date.

\newcommand{\ip}[2]{(#1, #2)}
                             % Defines \ip{arg1}{arg2} to mean
                             % (arg1, arg2).

%\newcommand{\ip}[2]{\langle #1 | #2\rangle}
                             % This is an alternative definition of
                             % \ip that is commented out.

\begin{document}             % End of preamble and beginning of text.

\maketitle                   % Produces the title.

This is an example input file.  Comparing it with
the output it generates can show you how to
produce a simple document of your own.

\section{Ordinary Text}      % Produces section heading.  Lower-level
                             % sections are begun with similar
                             % \subsection and \subsubsection commands.

The ends  of words and sentences are marked
  by   spaces. It  doesn't matter how many
spaces    you type; one is as good as 100.  The
end of   a line counts as a space.

One   or more   blank lines denote the  end
of  a paragraph.

Since any number of consecutive spaces are treated
like a single one, the formatting of the input
file makes no difference to
      \LaTeX,                % The \LaTeX command generates the LaTeX logo.
but it makes a difference to you.  When you use
\LaTeX, making your input file as easy to read
as possible will be a great help as you write
your document and when you change it.  This sample
file shows how you can add comments to your own input
file.

Because printing is different from typewriting,
there are a number of things that you have to do
differently when preparing an input file than if
you were just typing the document directly.
Quotation marks like
       ``this''
have to be handled specially, as do quotes within
quotes:
       ``\,`this'            % \, separates the double and single quote.
        is what I just
        wrote, not  `that'\,''.

Dashes come in three sizes: an
       intra-word
dash, a medium dash for number ranges like
       1--2,
and a punctuation
       dash---like
this.

A sentence-ending space should be larger than the
space between words within a sentence.  You
sometimes have to type special commands in
conjunction with punctuation characters to get
this right, as in the following sentence.
       Gnats, gnus, etc.\ all  % `\ ' makes an inter-word space.
       begin with G\@.         % \@ marks end-of-sentence punctuation.
You should check the spaces after periods when
reading your output to make sure you haven't
forgotten any special cases.  Generating an
ellipsis
       \ldots\               % `\ ' is needed after `\ldots' because TeX
                             % ignores spaces after command names like \ldots
                             % made from \ + letters.
                             %
                             % Note how a `%' character causes TeX to ignore
                             % the end of the input line, so these blank lines
                             % do not start a new paragraph.
                             %
with the right spacing around the periods requires
a special command.

\LaTeX\ interprets some common characters as
commands, so you must type special commands to
generate them.  These characters include the
following:
       \$ \& \% \# \{ and \}.

In printing, text is usually emphasized with an
       \emph{italic}
type style.

\begin{em}
   A long segment of text can also be emphasized
   in this way.  Text within such a segment can be
   given \emph{additional} emphasis.
\end{em}

It is sometimes necessary to prevent \LaTeX\ from
breaking a line where it might otherwise do so.
This may be at a space, as between the ``Mr.''\ and
``Jones'' in
       ``Mr.~Jones'',        % ~ produces an unbreakable interword space.
or within a word---especially when the word is a
symbol like
       \mbox{\emph{itemnum}}
that makes little sense when hyphenated across
lines.

Footnotes\footnote{This is an example of a footnote.}
pose no problem.

\LaTeX\ is good at typesetting mathematical formulas
like
       \( x-3y + z = 7 \)
or
       \( a_{1} > x^{2n} + y^{2n} > x' \)
or
       \( \ip{A}{B} = \sum_{i} a_{i} b_{i} \).
The spaces you type in a formula are
ignored.  Remember that a letter like
       $x$                   % $ ... $  and  \( ... \)  are equivalent
is a formula when it denotes a mathematical
symbol, and it should be typed as one.

\section{Displayed Text}

Text is displayed by indenting it from the left
margin.  Quotations are commonly displayed.  There
are short quotations
\begin{quote}
   This is a short quotation.  It consists of a
   single paragraph of text.  See how it is formatted.
\end{quote}
and longer ones.
\begin{quotation}
   This is a longer quotation.  It consists of two
   paragraphs of text, neither of which are
   particularly interesting.

   This is the second paragraph of the quotation.  It
   is just as dull as the first paragraph.
\end{quotation}
Another frequently-displayed structure is a list.
The following is an example of an \emph{itemized}
list.
\begin{itemize}
   \item This is the first item of an itemized list.
         Each item in the list is marked with a ``tick''.
         You don't have to worry about what kind of tick
         mark is used.

   \item This is the second item of the list.  It
         contains another list nested inside it.  The inner
         list is an \emph{enumerated} list.
         \begin{enumerate}
            \item This is the first item of an enumerated
                  list that is nested within the itemized list.

            \item This is the second item of the inner list.
                  \LaTeX\ allows you to nest lists deeper than
                  you really should.
         \end{enumerate}
         This is the rest of the second item of the outer
         list.  It is no more interesting than any other
         part of the item.
   \item This is the third item of the list.
\end{itemize}
You can even display poetry.
\begin{verse}
   There is an environment
    for verse \\             % The \\ command separates lines
   Whose features some poets % within a stanza.
   will curse.

                             % One or more blank lines separate stanzas.

   For instead of making\\
   Them do \emph{all} line breaking, \\
   It allows them to put too many words on a line when they'd rather be
   forced to be terse.
\end{verse}

Mathematical formulas may also be displayed.  A
displayed formula
is
one-line long; multiline
formulas require special formatting instructions.
   \[  \ip{\Gamma}{\psi'} = x'' + y^{2} + z_{i}^{n}\]
Don't start a paragraph with a displayed equation,
nor make one a paragraph by itself.

\end{document}               % End of document.

\section{Bench Huge Section 604}
% synthetic bench section 604
% This is a sample LaTeX input file.  (Version of 12 August 2004.)
%
% A '%' character causes TeX to ignore all remaining text on the line,
% and is used for comments like this one.

\documentclass{article}      % Specifies the document class

                             % The preamble begins here.
\title{An Example Document}  % Declares the document's title.
\author{Leslie Lamport}      % Declares the author's name.
\date{January 21, 1994}      % Deleting this command produces today's date.

\newcommand{\ip}[2]{(#1, #2)}
                             % Defines \ip{arg1}{arg2} to mean
                             % (arg1, arg2).

%\newcommand{\ip}[2]{\langle #1 | #2\rangle}
                             % This is an alternative definition of
                             % \ip that is commented out.

\begin{document}             % End of preamble and beginning of text.

\maketitle                   % Produces the title.

This is an example input file.  Comparing it with
the output it generates can show you how to
produce a simple document of your own.

\section{Ordinary Text}      % Produces section heading.  Lower-level
                             % sections are begun with similar
                             % \subsection and \subsubsection commands.

The ends  of words and sentences are marked
  by   spaces. It  doesn't matter how many
spaces    you type; one is as good as 100.  The
end of   a line counts as a space.

One   or more   blank lines denote the  end
of  a paragraph.

Since any number of consecutive spaces are treated
like a single one, the formatting of the input
file makes no difference to
      \LaTeX,                % The \LaTeX command generates the LaTeX logo.
but it makes a difference to you.  When you use
\LaTeX, making your input file as easy to read
as possible will be a great help as you write
your document and when you change it.  This sample
file shows how you can add comments to your own input
file.

Because printing is different from typewriting,
there are a number of things that you have to do
differently when preparing an input file than if
you were just typing the document directly.
Quotation marks like
       ``this''
have to be handled specially, as do quotes within
quotes:
       ``\,`this'            % \, separates the double and single quote.
        is what I just
        wrote, not  `that'\,''.

Dashes come in three sizes: an
       intra-word
dash, a medium dash for number ranges like
       1--2,
and a punctuation
       dash---like
this.

A sentence-ending space should be larger than the
space between words within a sentence.  You
sometimes have to type special commands in
conjunction with punctuation characters to get
this right, as in the following sentence.
       Gnats, gnus, etc.\ all  % `\ ' makes an inter-word space.
       begin with G\@.         % \@ marks end-of-sentence punctuation.
You should check the spaces after periods when
reading your output to make sure you haven't
forgotten any special cases.  Generating an
ellipsis
       \ldots\               % `\ ' is needed after `\ldots' because TeX
                             % ignores spaces after command names like \ldots
                             % made from \ + letters.
                             %
                             % Note how a `%' character causes TeX to ignore
                             % the end of the input line, so these blank lines
                             % do not start a new paragraph.
                             %
with the right spacing around the periods requires
a special command.

\LaTeX\ interprets some common characters as
commands, so you must type special commands to
generate them.  These characters include the
following:
       \$ \& \% \# \{ and \}.

In printing, text is usually emphasized with an
       \emph{italic}
type style.

\begin{em}
   A long segment of text can also be emphasized
   in this way.  Text within such a segment can be
   given \emph{additional} emphasis.
\end{em}

It is sometimes necessary to prevent \LaTeX\ from
breaking a line where it might otherwise do so.
This may be at a space, as between the ``Mr.''\ and
``Jones'' in
       ``Mr.~Jones'',        % ~ produces an unbreakable interword space.
or within a word---especially when the word is a
symbol like
       \mbox{\emph{itemnum}}
that makes little sense when hyphenated across
lines.

Footnotes\footnote{This is an example of a footnote.}
pose no problem.

\LaTeX\ is good at typesetting mathematical formulas
like
       \( x-3y + z = 7 \)
or
       \( a_{1} > x^{2n} + y^{2n} > x' \)
or
       \( \ip{A}{B} = \sum_{i} a_{i} b_{i} \).
The spaces you type in a formula are
ignored.  Remember that a letter like
       $x$                   % $ ... $  and  \( ... \)  are equivalent
is a formula when it denotes a mathematical
symbol, and it should be typed as one.

\section{Displayed Text}

Text is displayed by indenting it from the left
margin.  Quotations are commonly displayed.  There
are short quotations
\begin{quote}
   This is a short quotation.  It consists of a
   single paragraph of text.  See how it is formatted.
\end{quote}
and longer ones.
\begin{quotation}
   This is a longer quotation.  It consists of two
   paragraphs of text, neither of which are
   particularly interesting.

   This is the second paragraph of the quotation.  It
   is just as dull as the first paragraph.
\end{quotation}
Another frequently-displayed structure is a list.
The following is an example of an \emph{itemized}
list.
\begin{itemize}
   \item This is the first item of an itemized list.
         Each item in the list is marked with a ``tick''.
         You don't have to worry about what kind of tick
         mark is used.

   \item This is the second item of the list.  It
         contains another list nested inside it.  The inner
         list is an \emph{enumerated} list.
         \begin{enumerate}
            \item This is the first item of an enumerated
                  list that is nested within the itemized list.

            \item This is the second item of the inner list.
                  \LaTeX\ allows you to nest lists deeper than
                  you really should.
         \end{enumerate}
         This is the rest of the second item of the outer
         list.  It is no more interesting than any other
         part of the item.
   \item This is the third item of the list.
\end{itemize}
You can even display poetry.
\begin{verse}
   There is an environment
    for verse \\             % The \\ command separates lines
   Whose features some poets % within a stanza.
   will curse.

                             % One or more blank lines separate stanzas.

   For instead of making\\
   Them do \emph{all} line breaking, \\
   It allows them to put too many words on a line when they'd rather be
   forced to be terse.
\end{verse}

Mathematical formulas may also be displayed.  A
displayed formula
is
one-line long; multiline
formulas require special formatting instructions.
   \[  \ip{\Gamma}{\psi'} = x'' + y^{2} + z_{i}^{n}\]
Don't start a paragraph with a displayed equation,
nor make one a paragraph by itself.

\end{document}               % End of document.

\section{Bench Huge Section 605}
% synthetic bench section 605
% This is a sample LaTeX input file.  (Version of 12 August 2004.)
%
% A '%' character causes TeX to ignore all remaining text on the line,
% and is used for comments like this one.

\documentclass{article}      % Specifies the document class

                             % The preamble begins here.
\title{An Example Document}  % Declares the document's title.
\author{Leslie Lamport}      % Declares the author's name.
\date{January 21, 1994}      % Deleting this command produces today's date.

\newcommand{\ip}[2]{(#1, #2)}
                             % Defines \ip{arg1}{arg2} to mean
                             % (arg1, arg2).

%\newcommand{\ip}[2]{\langle #1 | #2\rangle}
                             % This is an alternative definition of
                             % \ip that is commented out.

\begin{document}             % End of preamble and beginning of text.

\maketitle                   % Produces the title.

This is an example input file.  Comparing it with
the output it generates can show you how to
produce a simple document of your own.

\section{Ordinary Text}      % Produces section heading.  Lower-level
                             % sections are begun with similar
                             % \subsection and \subsubsection commands.

The ends  of words and sentences are marked
  by   spaces. It  doesn't matter how many
spaces    you type; one is as good as 100.  The
end of   a line counts as a space.

One   or more   blank lines denote the  end
of  a paragraph.

Since any number of consecutive spaces are treated
like a single one, the formatting of the input
file makes no difference to
      \LaTeX,                % The \LaTeX command generates the LaTeX logo.
but it makes a difference to you.  When you use
\LaTeX, making your input file as easy to read
as possible will be a great help as you write
your document and when you change it.  This sample
file shows how you can add comments to your own input
file.

Because printing is different from typewriting,
there are a number of things that you have to do
differently when preparing an input file than if
you were just typing the document directly.
Quotation marks like
       ``this''
have to be handled specially, as do quotes within
quotes:
       ``\,`this'            % \, separates the double and single quote.
        is what I just
        wrote, not  `that'\,''.

Dashes come in three sizes: an
       intra-word
dash, a medium dash for number ranges like
       1--2,
and a punctuation
       dash---like
this.

A sentence-ending space should be larger than the
space between words within a sentence.  You
sometimes have to type special commands in
conjunction with punctuation characters to get
this right, as in the following sentence.
       Gnats, gnus, etc.\ all  % `\ ' makes an inter-word space.
       begin with G\@.         % \@ marks end-of-sentence punctuation.
You should check the spaces after periods when
reading your output to make sure you haven't
forgotten any special cases.  Generating an
ellipsis
       \ldots\               % `\ ' is needed after `\ldots' because TeX
                             % ignores spaces after command names like \ldots
                             % made from \ + letters.
                             %
                             % Note how a `%' character causes TeX to ignore
                             % the end of the input line, so these blank lines
                             % do not start a new paragraph.
                             %
with the right spacing around the periods requires
a special command.

\LaTeX\ interprets some common characters as
commands, so you must type special commands to
generate them.  These characters include the
following:
       \$ \& \% \# \{ and \}.

In printing, text is usually emphasized with an
       \emph{italic}
type style.

\begin{em}
   A long segment of text can also be emphasized
   in this way.  Text within such a segment can be
   given \emph{additional} emphasis.
\end{em}

It is sometimes necessary to prevent \LaTeX\ from
breaking a line where it might otherwise do so.
This may be at a space, as between the ``Mr.''\ and
``Jones'' in
       ``Mr.~Jones'',        % ~ produces an unbreakable interword space.
or within a word---especially when the word is a
symbol like
       \mbox{\emph{itemnum}}
that makes little sense when hyphenated across
lines.

Footnotes\footnote{This is an example of a footnote.}
pose no problem.

\LaTeX\ is good at typesetting mathematical formulas
like
       \( x-3y + z = 7 \)
or
       \( a_{1} > x^{2n} + y^{2n} > x' \)
or
       \( \ip{A}{B} = \sum_{i} a_{i} b_{i} \).
The spaces you type in a formula are
ignored.  Remember that a letter like
       $x$                   % $ ... $  and  \( ... \)  are equivalent
is a formula when it denotes a mathematical
symbol, and it should be typed as one.

\section{Displayed Text}

Text is displayed by indenting it from the left
margin.  Quotations are commonly displayed.  There
are short quotations
\begin{quote}
   This is a short quotation.  It consists of a
   single paragraph of text.  See how it is formatted.
\end{quote}
and longer ones.
\begin{quotation}
   This is a longer quotation.  It consists of two
   paragraphs of text, neither of which are
   particularly interesting.

   This is the second paragraph of the quotation.  It
   is just as dull as the first paragraph.
\end{quotation}
Another frequently-displayed structure is a list.
The following is an example of an \emph{itemized}
list.
\begin{itemize}
   \item This is the first item of an itemized list.
         Each item in the list is marked with a ``tick''.
         You don't have to worry about what kind of tick
         mark is used.

   \item This is the second item of the list.  It
         contains another list nested inside it.  The inner
         list is an \emph{enumerated} list.
         \begin{enumerate}
            \item This is the first item of an enumerated
                  list that is nested within the itemized list.

            \item This is the second item of the inner list.
                  \LaTeX\ allows you to nest lists deeper than
                  you really should.
         \end{enumerate}
         This is the rest of the second item of the outer
         list.  It is no more interesting than any other
         part of the item.
   \item This is the third item of the list.
\end{itemize}
You can even display poetry.
\begin{verse}
   There is an environment
    for verse \\             % The \\ command separates lines
   Whose features some poets % within a stanza.
   will curse.

                             % One or more blank lines separate stanzas.

   For instead of making\\
   Them do \emph{all} line breaking, \\
   It allows them to put too many words on a line when they'd rather be
   forced to be terse.
\end{verse}

Mathematical formulas may also be displayed.  A
displayed formula
is
one-line long; multiline
formulas require special formatting instructions.
   \[  \ip{\Gamma}{\psi'} = x'' + y^{2} + z_{i}^{n}\]
Don't start a paragraph with a displayed equation,
nor make one a paragraph by itself.

\end{document}               % End of document.

\section{Bench Huge Section 606}
% synthetic bench section 606
% This is a sample LaTeX input file.  (Version of 12 August 2004.)
%
% A '%' character causes TeX to ignore all remaining text on the line,
% and is used for comments like this one.

\documentclass{article}      % Specifies the document class

                             % The preamble begins here.
\title{An Example Document}  % Declares the document's title.
\author{Leslie Lamport}      % Declares the author's name.
\date{January 21, 1994}      % Deleting this command produces today's date.

\newcommand{\ip}[2]{(#1, #2)}
                             % Defines \ip{arg1}{arg2} to mean
                             % (arg1, arg2).

%\newcommand{\ip}[2]{\langle #1 | #2\rangle}
                             % This is an alternative definition of
                             % \ip that is commented out.

\begin{document}             % End of preamble and beginning of text.

\maketitle                   % Produces the title.

This is an example input file.  Comparing it with
the output it generates can show you how to
produce a simple document of your own.

\section{Ordinary Text}      % Produces section heading.  Lower-level
                             % sections are begun with similar
                             % \subsection and \subsubsection commands.

The ends  of words and sentences are marked
  by   spaces. It  doesn't matter how many
spaces    you type; one is as good as 100.  The
end of   a line counts as a space.

One   or more   blank lines denote the  end
of  a paragraph.

Since any number of consecutive spaces are treated
like a single one, the formatting of the input
file makes no difference to
      \LaTeX,                % The \LaTeX command generates the LaTeX logo.
but it makes a difference to you.  When you use
\LaTeX, making your input file as easy to read
as possible will be a great help as you write
your document and when you change it.  This sample
file shows how you can add comments to your own input
file.

Because printing is different from typewriting,
there are a number of things that you have to do
differently when preparing an input file than if
you were just typing the document directly.
Quotation marks like
       ``this''
have to be handled specially, as do quotes within
quotes:
       ``\,`this'            % \, separates the double and single quote.
        is what I just
        wrote, not  `that'\,''.

Dashes come in three sizes: an
       intra-word
dash, a medium dash for number ranges like
       1--2,
and a punctuation
       dash---like
this.

A sentence-ending space should be larger than the
space between words within a sentence.  You
sometimes have to type special commands in
conjunction with punctuation characters to get
this right, as in the following sentence.
       Gnats, gnus, etc.\ all  % `\ ' makes an inter-word space.
       begin with G\@.         % \@ marks end-of-sentence punctuation.
You should check the spaces after periods when
reading your output to make sure you haven't
forgotten any special cases.  Generating an
ellipsis
       \ldots\               % `\ ' is needed after `\ldots' because TeX
                             % ignores spaces after command names like \ldots
                             % made from \ + letters.
                             %
                             % Note how a `%' character causes TeX to ignore
                             % the end of the input line, so these blank lines
                             % do not start a new paragraph.
                             %
with the right spacing around the periods requires
a special command.

\LaTeX\ interprets some common characters as
commands, so you must type special commands to
generate them.  These characters include the
following:
       \$ \& \% \# \{ and \}.

In printing, text is usually emphasized with an
       \emph{italic}
type style.

\begin{em}
   A long segment of text can also be emphasized
   in this way.  Text within such a segment can be
   given \emph{additional} emphasis.
\end{em}

It is sometimes necessary to prevent \LaTeX\ from
breaking a line where it might otherwise do so.
This may be at a space, as between the ``Mr.''\ and
``Jones'' in
       ``Mr.~Jones'',        % ~ produces an unbreakable interword space.
or within a word---especially when the word is a
symbol like
       \mbox{\emph{itemnum}}
that makes little sense when hyphenated across
lines.

Footnotes\footnote{This is an example of a footnote.}
pose no problem.

\LaTeX\ is good at typesetting mathematical formulas
like
       \( x-3y + z = 7 \)
or
       \( a_{1} > x^{2n} + y^{2n} > x' \)
or
       \( \ip{A}{B} = \sum_{i} a_{i} b_{i} \).
The spaces you type in a formula are
ignored.  Remember that a letter like
       $x$                   % $ ... $  and  \( ... \)  are equivalent
is a formula when it denotes a mathematical
symbol, and it should be typed as one.

\section{Displayed Text}

Text is displayed by indenting it from the left
margin.  Quotations are commonly displayed.  There
are short quotations
\begin{quote}
   This is a short quotation.  It consists of a
   single paragraph of text.  See how it is formatted.
\end{quote}
and longer ones.
\begin{quotation}
   This is a longer quotation.  It consists of two
   paragraphs of text, neither of which are
   particularly interesting.

   This is the second paragraph of the quotation.  It
   is just as dull as the first paragraph.
\end{quotation}
Another frequently-displayed structure is a list.
The following is an example of an \emph{itemized}
list.
\begin{itemize}
   \item This is the first item of an itemized list.
         Each item in the list is marked with a ``tick''.
         You don't have to worry about what kind of tick
         mark is used.

   \item This is the second item of the list.  It
         contains another list nested inside it.  The inner
         list is an \emph{enumerated} list.
         \begin{enumerate}
            \item This is the first item of an enumerated
                  list that is nested within the itemized list.

            \item This is the second item of the inner list.
                  \LaTeX\ allows you to nest lists deeper than
                  you really should.
         \end{enumerate}
         This is the rest of the second item of the outer
         list.  It is no more interesting than any other
         part of the item.
   \item This is the third item of the list.
\end{itemize}
You can even display poetry.
\begin{verse}
   There is an environment
    for verse \\             % The \\ command separates lines
   Whose features some poets % within a stanza.
   will curse.

                             % One or more blank lines separate stanzas.

   For instead of making\\
   Them do \emph{all} line breaking, \\
   It allows them to put too many words on a line when they'd rather be
   forced to be terse.
\end{verse}

Mathematical formulas may also be displayed.  A
displayed formula
is
one-line long; multiline
formulas require special formatting instructions.
   \[  \ip{\Gamma}{\psi'} = x'' + y^{2} + z_{i}^{n}\]
Don't start a paragraph with a displayed equation,
nor make one a paragraph by itself.

\end{document}               % End of document.

\section{Bench Huge Section 607}
% synthetic bench section 607
% This is a sample LaTeX input file.  (Version of 12 August 2004.)
%
% A '%' character causes TeX to ignore all remaining text on the line,
% and is used for comments like this one.

\documentclass{article}      % Specifies the document class

                             % The preamble begins here.
\title{An Example Document}  % Declares the document's title.
\author{Leslie Lamport}      % Declares the author's name.
\date{January 21, 1994}      % Deleting this command produces today's date.

\newcommand{\ip}[2]{(#1, #2)}
                             % Defines \ip{arg1}{arg2} to mean
                             % (arg1, arg2).

%\newcommand{\ip}[2]{\langle #1 | #2\rangle}
                             % This is an alternative definition of
                             % \ip that is commented out.

\begin{document}             % End of preamble and beginning of text.

\maketitle                   % Produces the title.

This is an example input file.  Comparing it with
the output it generates can show you how to
produce a simple document of your own.

\section{Ordinary Text}      % Produces section heading.  Lower-level
                             % sections are begun with similar
                             % \subsection and \subsubsection commands.

The ends  of words and sentences are marked
  by   spaces. It  doesn't matter how many
spaces    you type; one is as good as 100.  The
end of   a line counts as a space.

One   or more   blank lines denote the  end
of  a paragraph.

Since any number of consecutive spaces are treated
like a single one, the formatting of the input
file makes no difference to
      \LaTeX,                % The \LaTeX command generates the LaTeX logo.
but it makes a difference to you.  When you use
\LaTeX, making your input file as easy to read
as possible will be a great help as you write
your document and when you change it.  This sample
file shows how you can add comments to your own input
file.

Because printing is different from typewriting,
there are a number of things that you have to do
differently when preparing an input file than if
you were just typing the document directly.
Quotation marks like
       ``this''
have to be handled specially, as do quotes within
quotes:
       ``\,`this'            % \, separates the double and single quote.
        is what I just
        wrote, not  `that'\,''.

Dashes come in three sizes: an
       intra-word
dash, a medium dash for number ranges like
       1--2,
and a punctuation
       dash---like
this.

A sentence-ending space should be larger than the
space between words within a sentence.  You
sometimes have to type special commands in
conjunction with punctuation characters to get
this right, as in the following sentence.
       Gnats, gnus, etc.\ all  % `\ ' makes an inter-word space.
       begin with G\@.         % \@ marks end-of-sentence punctuation.
You should check the spaces after periods when
reading your output to make sure you haven't
forgotten any special cases.  Generating an
ellipsis
       \ldots\               % `\ ' is needed after `\ldots' because TeX
                             % ignores spaces after command names like \ldots
                             % made from \ + letters.
                             %
                             % Note how a `%' character causes TeX to ignore
                             % the end of the input line, so these blank lines
                             % do not start a new paragraph.
                             %
with the right spacing around the periods requires
a special command.

\LaTeX\ interprets some common characters as
commands, so you must type special commands to
generate them.  These characters include the
following:
       \$ \& \% \# \{ and \}.

In printing, text is usually emphasized with an
       \emph{italic}
type style.

\begin{em}
   A long segment of text can also be emphasized
   in this way.  Text within such a segment can be
   given \emph{additional} emphasis.
\end{em}

It is sometimes necessary to prevent \LaTeX\ from
breaking a line where it might otherwise do so.
This may be at a space, as between the ``Mr.''\ and
``Jones'' in
       ``Mr.~Jones'',        % ~ produces an unbreakable interword space.
or within a word---especially when the word is a
symbol like
       \mbox{\emph{itemnum}}
that makes little sense when hyphenated across
lines.

Footnotes\footnote{This is an example of a footnote.}
pose no problem.

\LaTeX\ is good at typesetting mathematical formulas
like
       \( x-3y + z = 7 \)
or
       \( a_{1} > x^{2n} + y^{2n} > x' \)
or
       \( \ip{A}{B} = \sum_{i} a_{i} b_{i} \).
The spaces you type in a formula are
ignored.  Remember that a letter like
       $x$                   % $ ... $  and  \( ... \)  are equivalent
is a formula when it denotes a mathematical
symbol, and it should be typed as one.

\section{Displayed Text}

Text is displayed by indenting it from the left
margin.  Quotations are commonly displayed.  There
are short quotations
\begin{quote}
   This is a short quotation.  It consists of a
   single paragraph of text.  See how it is formatted.
\end{quote}
and longer ones.
\begin{quotation}
   This is a longer quotation.  It consists of two
   paragraphs of text, neither of which are
   particularly interesting.

   This is the second paragraph of the quotation.  It
   is just as dull as the first paragraph.
\end{quotation}
Another frequently-displayed structure is a list.
The following is an example of an \emph{itemized}
list.
\begin{itemize}
   \item This is the first item of an itemized list.
         Each item in the list is marked with a ``tick''.
         You don't have to worry about what kind of tick
         mark is used.

   \item This is the second item of the list.  It
         contains another list nested inside it.  The inner
         list is an \emph{enumerated} list.
         \begin{enumerate}
            \item This is the first item of an enumerated
                  list that is nested within the itemized list.

            \item This is the second item of the inner list.
                  \LaTeX\ allows you to nest lists deeper than
                  you really should.
         \end{enumerate}
         This is the rest of the second item of the outer
         list.  It is no more interesting than any other
         part of the item.
   \item This is the third item of the list.
\end{itemize}
You can even display poetry.
\begin{verse}
   There is an environment
    for verse \\             % The \\ command separates lines
   Whose features some poets % within a stanza.
   will curse.

                             % One or more blank lines separate stanzas.

   For instead of making\\
   Them do \emph{all} line breaking, \\
   It allows them to put too many words on a line when they'd rather be
   forced to be terse.
\end{verse}

Mathematical formulas may also be displayed.  A
displayed formula
is
one-line long; multiline
formulas require special formatting instructions.
   \[  \ip{\Gamma}{\psi'} = x'' + y^{2} + z_{i}^{n}\]
Don't start a paragraph with a displayed equation,
nor make one a paragraph by itself.

\end{document}               % End of document.

\section{Bench Huge Section 608}
% synthetic bench section 608
% This is a sample LaTeX input file.  (Version of 12 August 2004.)
%
% A '%' character causes TeX to ignore all remaining text on the line,
% and is used for comments like this one.

\documentclass{article}      % Specifies the document class

                             % The preamble begins here.
\title{An Example Document}  % Declares the document's title.
\author{Leslie Lamport}      % Declares the author's name.
\date{January 21, 1994}      % Deleting this command produces today's date.

\newcommand{\ip}[2]{(#1, #2)}
                             % Defines \ip{arg1}{arg2} to mean
                             % (arg1, arg2).

%\newcommand{\ip}[2]{\langle #1 | #2\rangle}
                             % This is an alternative definition of
                             % \ip that is commented out.

\begin{document}             % End of preamble and beginning of text.

\maketitle                   % Produces the title.

This is an example input file.  Comparing it with
the output it generates can show you how to
produce a simple document of your own.

\section{Ordinary Text}      % Produces section heading.  Lower-level
                             % sections are begun with similar
                             % \subsection and \subsubsection commands.

The ends  of words and sentences are marked
  by   spaces. It  doesn't matter how many
spaces    you type; one is as good as 100.  The
end of   a line counts as a space.

One   or more   blank lines denote the  end
of  a paragraph.

Since any number of consecutive spaces are treated
like a single one, the formatting of the input
file makes no difference to
      \LaTeX,                % The \LaTeX command generates the LaTeX logo.
but it makes a difference to you.  When you use
\LaTeX, making your input file as easy to read
as possible will be a great help as you write
your document and when you change it.  This sample
file shows how you can add comments to your own input
file.

Because printing is different from typewriting,
there are a number of things that you have to do
differently when preparing an input file than if
you were just typing the document directly.
Quotation marks like
       ``this''
have to be handled specially, as do quotes within
quotes:
       ``\,`this'            % \, separates the double and single quote.
        is what I just
        wrote, not  `that'\,''.

Dashes come in three sizes: an
       intra-word
dash, a medium dash for number ranges like
       1--2,
and a punctuation
       dash---like
this.

A sentence-ending space should be larger than the
space between words within a sentence.  You
sometimes have to type special commands in
conjunction with punctuation characters to get
this right, as in the following sentence.
       Gnats, gnus, etc.\ all  % `\ ' makes an inter-word space.
       begin with G\@.         % \@ marks end-of-sentence punctuation.
You should check the spaces after periods when
reading your output to make sure you haven't
forgotten any special cases.  Generating an
ellipsis
       \ldots\               % `\ ' is needed after `\ldots' because TeX
                             % ignores spaces after command names like \ldots
                             % made from \ + letters.
                             %
                             % Note how a `%' character causes TeX to ignore
                             % the end of the input line, so these blank lines
                             % do not start a new paragraph.
                             %
with the right spacing around the periods requires
a special command.

\LaTeX\ interprets some common characters as
commands, so you must type special commands to
generate them.  These characters include the
following:
       \$ \& \% \# \{ and \}.

In printing, text is usually emphasized with an
       \emph{italic}
type style.

\begin{em}
   A long segment of text can also be emphasized
   in this way.  Text within such a segment can be
   given \emph{additional} emphasis.
\end{em}

It is sometimes necessary to prevent \LaTeX\ from
breaking a line where it might otherwise do so.
This may be at a space, as between the ``Mr.''\ and
``Jones'' in
       ``Mr.~Jones'',        % ~ produces an unbreakable interword space.
or within a word---especially when the word is a
symbol like
       \mbox{\emph{itemnum}}
that makes little sense when hyphenated across
lines.

Footnotes\footnote{This is an example of a footnote.}
pose no problem.

\LaTeX\ is good at typesetting mathematical formulas
like
       \( x-3y + z = 7 \)
or
       \( a_{1} > x^{2n} + y^{2n} > x' \)
or
       \( \ip{A}{B} = \sum_{i} a_{i} b_{i} \).
The spaces you type in a formula are
ignored.  Remember that a letter like
       $x$                   % $ ... $  and  \( ... \)  are equivalent
is a formula when it denotes a mathematical
symbol, and it should be typed as one.

\section{Displayed Text}

Text is displayed by indenting it from the left
margin.  Quotations are commonly displayed.  There
are short quotations
\begin{quote}
   This is a short quotation.  It consists of a
   single paragraph of text.  See how it is formatted.
\end{quote}
and longer ones.
\begin{quotation}
   This is a longer quotation.  It consists of two
   paragraphs of text, neither of which are
   particularly interesting.

   This is the second paragraph of the quotation.  It
   is just as dull as the first paragraph.
\end{quotation}
Another frequently-displayed structure is a list.
The following is an example of an \emph{itemized}
list.
\begin{itemize}
   \item This is the first item of an itemized list.
         Each item in the list is marked with a ``tick''.
         You don't have to worry about what kind of tick
         mark is used.

   \item This is the second item of the list.  It
         contains another list nested inside it.  The inner
         list is an \emph{enumerated} list.
         \begin{enumerate}
            \item This is the first item of an enumerated
                  list that is nested within the itemized list.

            \item This is the second item of the inner list.
                  \LaTeX\ allows you to nest lists deeper than
                  you really should.
         \end{enumerate}
         This is the rest of the second item of the outer
         list.  It is no more interesting than any other
         part of the item.
   \item This is the third item of the list.
\end{itemize}
You can even display poetry.
\begin{verse}
   There is an environment
    for verse \\             % The \\ command separates lines
   Whose features some poets % within a stanza.
   will curse.

                             % One or more blank lines separate stanzas.

   For instead of making\\
   Them do \emph{all} line breaking, \\
   It allows them to put too many words on a line when they'd rather be
   forced to be terse.
\end{verse}

Mathematical formulas may also be displayed.  A
displayed formula
is
one-line long; multiline
formulas require special formatting instructions.
   \[  \ip{\Gamma}{\psi'} = x'' + y^{2} + z_{i}^{n}\]
Don't start a paragraph with a displayed equation,
nor make one a paragraph by itself.

\end{document}               % End of document.

\section{Bench Huge Section 609}
% synthetic bench section 609
% This is a sample LaTeX input file.  (Version of 12 August 2004.)
%
% A '%' character causes TeX to ignore all remaining text on the line,
% and is used for comments like this one.

\documentclass{article}      % Specifies the document class

                             % The preamble begins here.
\title{An Example Document}  % Declares the document's title.
\author{Leslie Lamport}      % Declares the author's name.
\date{January 21, 1994}      % Deleting this command produces today's date.

\newcommand{\ip}[2]{(#1, #2)}
                             % Defines \ip{arg1}{arg2} to mean
                             % (arg1, arg2).

%\newcommand{\ip}[2]{\langle #1 | #2\rangle}
                             % This is an alternative definition of
                             % \ip that is commented out.

\begin{document}             % End of preamble and beginning of text.

\maketitle                   % Produces the title.

This is an example input file.  Comparing it with
the output it generates can show you how to
produce a simple document of your own.

\section{Ordinary Text}      % Produces section heading.  Lower-level
                             % sections are begun with similar
                             % \subsection and \subsubsection commands.

The ends  of words and sentences are marked
  by   spaces. It  doesn't matter how many
spaces    you type; one is as good as 100.  The
end of   a line counts as a space.

One   or more   blank lines denote the  end
of  a paragraph.

Since any number of consecutive spaces are treated
like a single one, the formatting of the input
file makes no difference to
      \LaTeX,                % The \LaTeX command generates the LaTeX logo.
but it makes a difference to you.  When you use
\LaTeX, making your input file as easy to read
as possible will be a great help as you write
your document and when you change it.  This sample
file shows how you can add comments to your own input
file.

Because printing is different from typewriting,
there are a number of things that you have to do
differently when preparing an input file than if
you were just typing the document directly.
Quotation marks like
       ``this''
have to be handled specially, as do quotes within
quotes:
       ``\,`this'            % \, separates the double and single quote.
        is what I just
        wrote, not  `that'\,''.

Dashes come in three sizes: an
       intra-word
dash, a medium dash for number ranges like
       1--2,
and a punctuation
       dash---like
this.

A sentence-ending space should be larger than the
space between words within a sentence.  You
sometimes have to type special commands in
conjunction with punctuation characters to get
this right, as in the following sentence.
       Gnats, gnus, etc.\ all  % `\ ' makes an inter-word space.
       begin with G\@.         % \@ marks end-of-sentence punctuation.
You should check the spaces after periods when
reading your output to make sure you haven't
forgotten any special cases.  Generating an
ellipsis
       \ldots\               % `\ ' is needed after `\ldots' because TeX
                             % ignores spaces after command names like \ldots
                             % made from \ + letters.
                             %
                             % Note how a `%' character causes TeX to ignore
                             % the end of the input line, so these blank lines
                             % do not start a new paragraph.
                             %
with the right spacing around the periods requires
a special command.

\LaTeX\ interprets some common characters as
commands, so you must type special commands to
generate them.  These characters include the
following:
       \$ \& \% \# \{ and \}.

In printing, text is usually emphasized with an
       \emph{italic}
type style.

\begin{em}
   A long segment of text can also be emphasized
   in this way.  Text within such a segment can be
   given \emph{additional} emphasis.
\end{em}

It is sometimes necessary to prevent \LaTeX\ from
breaking a line where it might otherwise do so.
This may be at a space, as between the ``Mr.''\ and
``Jones'' in
       ``Mr.~Jones'',        % ~ produces an unbreakable interword space.
or within a word---especially when the word is a
symbol like
       \mbox{\emph{itemnum}}
that makes little sense when hyphenated across
lines.

Footnotes\footnote{This is an example of a footnote.}
pose no problem.

\LaTeX\ is good at typesetting mathematical formulas
like
       \( x-3y + z = 7 \)
or
       \( a_{1} > x^{2n} + y^{2n} > x' \)
or
       \( \ip{A}{B} = \sum_{i} a_{i} b_{i} \).
The spaces you type in a formula are
ignored.  Remember that a letter like
       $x$                   % $ ... $  and  \( ... \)  are equivalent
is a formula when it denotes a mathematical
symbol, and it should be typed as one.

\section{Displayed Text}

Text is displayed by indenting it from the left
margin.  Quotations are commonly displayed.  There
are short quotations
\begin{quote}
   This is a short quotation.  It consists of a
   single paragraph of text.  See how it is formatted.
\end{quote}
and longer ones.
\begin{quotation}
   This is a longer quotation.  It consists of two
   paragraphs of text, neither of which are
   particularly interesting.

   This is the second paragraph of the quotation.  It
   is just as dull as the first paragraph.
\end{quotation}
Another frequently-displayed structure is a list.
The following is an example of an \emph{itemized}
list.
\begin{itemize}
   \item This is the first item of an itemized list.
         Each item in the list is marked with a ``tick''.
         You don't have to worry about what kind of tick
         mark is used.

   \item This is the second item of the list.  It
         contains another list nested inside it.  The inner
         list is an \emph{enumerated} list.
         \begin{enumerate}
            \item This is the first item of an enumerated
                  list that is nested within the itemized list.

            \item This is the second item of the inner list.
                  \LaTeX\ allows you to nest lists deeper than
                  you really should.
         \end{enumerate}
         This is the rest of the second item of the outer
         list.  It is no more interesting than any other
         part of the item.
   \item This is the third item of the list.
\end{itemize}
You can even display poetry.
\begin{verse}
   There is an environment
    for verse \\             % The \\ command separates lines
   Whose features some poets % within a stanza.
   will curse.

                             % One or more blank lines separate stanzas.

   For instead of making\\
   Them do \emph{all} line breaking, \\
   It allows them to put too many words on a line when they'd rather be
   forced to be terse.
\end{verse}

Mathematical formulas may also be displayed.  A
displayed formula
is
one-line long; multiline
formulas require special formatting instructions.
   \[  \ip{\Gamma}{\psi'} = x'' + y^{2} + z_{i}^{n}\]
Don't start a paragraph with a displayed equation,
nor make one a paragraph by itself.

\end{document}               % End of document.

\section{Bench Huge Section 610}
% synthetic bench section 610
% This is a sample LaTeX input file.  (Version of 12 August 2004.)
%
% A '%' character causes TeX to ignore all remaining text on the line,
% and is used for comments like this one.

\documentclass{article}      % Specifies the document class

                             % The preamble begins here.
\title{An Example Document}  % Declares the document's title.
\author{Leslie Lamport}      % Declares the author's name.
\date{January 21, 1994}      % Deleting this command produces today's date.

\newcommand{\ip}[2]{(#1, #2)}
                             % Defines \ip{arg1}{arg2} to mean
                             % (arg1, arg2).

%\newcommand{\ip}[2]{\langle #1 | #2\rangle}
                             % This is an alternative definition of
                             % \ip that is commented out.

\begin{document}             % End of preamble and beginning of text.

\maketitle                   % Produces the title.

This is an example input file.  Comparing it with
the output it generates can show you how to
produce a simple document of your own.

\section{Ordinary Text}      % Produces section heading.  Lower-level
                             % sections are begun with similar
                             % \subsection and \subsubsection commands.

The ends  of words and sentences are marked
  by   spaces. It  doesn't matter how many
spaces    you type; one is as good as 100.  The
end of   a line counts as a space.

One   or more   blank lines denote the  end
of  a paragraph.

Since any number of consecutive spaces are treated
like a single one, the formatting of the input
file makes no difference to
      \LaTeX,                % The \LaTeX command generates the LaTeX logo.
but it makes a difference to you.  When you use
\LaTeX, making your input file as easy to read
as possible will be a great help as you write
your document and when you change it.  This sample
file shows how you can add comments to your own input
file.

Because printing is different from typewriting,
there are a number of things that you have to do
differently when preparing an input file than if
you were just typing the document directly.
Quotation marks like
       ``this''
have to be handled specially, as do quotes within
quotes:
       ``\,`this'            % \, separates the double and single quote.
        is what I just
        wrote, not  `that'\,''.

Dashes come in three sizes: an
       intra-word
dash, a medium dash for number ranges like
       1--2,
and a punctuation
       dash---like
this.

A sentence-ending space should be larger than the
space between words within a sentence.  You
sometimes have to type special commands in
conjunction with punctuation characters to get
this right, as in the following sentence.
       Gnats, gnus, etc.\ all  % `\ ' makes an inter-word space.
       begin with G\@.         % \@ marks end-of-sentence punctuation.
You should check the spaces after periods when
reading your output to make sure you haven't
forgotten any special cases.  Generating an
ellipsis
       \ldots\               % `\ ' is needed after `\ldots' because TeX
                             % ignores spaces after command names like \ldots
                             % made from \ + letters.
                             %
                             % Note how a `%' character causes TeX to ignore
                             % the end of the input line, so these blank lines
                             % do not start a new paragraph.
                             %
with the right spacing around the periods requires
a special command.

\LaTeX\ interprets some common characters as
commands, so you must type special commands to
generate them.  These characters include the
following:
       \$ \& \% \# \{ and \}.

In printing, text is usually emphasized with an
       \emph{italic}
type style.

\begin{em}
   A long segment of text can also be emphasized
   in this way.  Text within such a segment can be
   given \emph{additional} emphasis.
\end{em}

It is sometimes necessary to prevent \LaTeX\ from
breaking a line where it might otherwise do so.
This may be at a space, as between the ``Mr.''\ and
``Jones'' in
       ``Mr.~Jones'',        % ~ produces an unbreakable interword space.
or within a word---especially when the word is a
symbol like
       \mbox{\emph{itemnum}}
that makes little sense when hyphenated across
lines.

Footnotes\footnote{This is an example of a footnote.}
pose no problem.

\LaTeX\ is good at typesetting mathematical formulas
like
       \( x-3y + z = 7 \)
or
       \( a_{1} > x^{2n} + y^{2n} > x' \)
or
       \( \ip{A}{B} = \sum_{i} a_{i} b_{i} \).
The spaces you type in a formula are
ignored.  Remember that a letter like
       $x$                   % $ ... $  and  \( ... \)  are equivalent
is a formula when it denotes a mathematical
symbol, and it should be typed as one.

\section{Displayed Text}

Text is displayed by indenting it from the left
margin.  Quotations are commonly displayed.  There
are short quotations
\begin{quote}
   This is a short quotation.  It consists of a
   single paragraph of text.  See how it is formatted.
\end{quote}
and longer ones.
\begin{quotation}
   This is a longer quotation.  It consists of two
   paragraphs of text, neither of which are
   particularly interesting.

   This is the second paragraph of the quotation.  It
   is just as dull as the first paragraph.
\end{quotation}
Another frequently-displayed structure is a list.
The following is an example of an \emph{itemized}
list.
\begin{itemize}
   \item This is the first item of an itemized list.
         Each item in the list is marked with a ``tick''.
         You don't have to worry about what kind of tick
         mark is used.

   \item This is the second item of the list.  It
         contains another list nested inside it.  The inner
         list is an \emph{enumerated} list.
         \begin{enumerate}
            \item This is the first item of an enumerated
                  list that is nested within the itemized list.

            \item This is the second item of the inner list.
                  \LaTeX\ allows you to nest lists deeper than
                  you really should.
         \end{enumerate}
         This is the rest of the second item of the outer
         list.  It is no more interesting than any other
         part of the item.
   \item This is the third item of the list.
\end{itemize}
You can even display poetry.
\begin{verse}
   There is an environment
    for verse \\             % The \\ command separates lines
   Whose features some poets % within a stanza.
   will curse.

                             % One or more blank lines separate stanzas.

   For instead of making\\
   Them do \emph{all} line breaking, \\
   It allows them to put too many words on a line when they'd rather be
   forced to be terse.
\end{verse}

Mathematical formulas may also be displayed.  A
displayed formula
is
one-line long; multiline
formulas require special formatting instructions.
   \[  \ip{\Gamma}{\psi'} = x'' + y^{2} + z_{i}^{n}\]
Don't start a paragraph with a displayed equation,
nor make one a paragraph by itself.

\end{document}               % End of document.

\section{Bench Huge Section 611}
% synthetic bench section 611
% This is a sample LaTeX input file.  (Version of 12 August 2004.)
%
% A '%' character causes TeX to ignore all remaining text on the line,
% and is used for comments like this one.

\documentclass{article}      % Specifies the document class

                             % The preamble begins here.
\title{An Example Document}  % Declares the document's title.
\author{Leslie Lamport}      % Declares the author's name.
\date{January 21, 1994}      % Deleting this command produces today's date.

\newcommand{\ip}[2]{(#1, #2)}
                             % Defines \ip{arg1}{arg2} to mean
                             % (arg1, arg2).

%\newcommand{\ip}[2]{\langle #1 | #2\rangle}
                             % This is an alternative definition of
                             % \ip that is commented out.

\begin{document}             % End of preamble and beginning of text.

\maketitle                   % Produces the title.

This is an example input file.  Comparing it with
the output it generates can show you how to
produce a simple document of your own.

\section{Ordinary Text}      % Produces section heading.  Lower-level
                             % sections are begun with similar
                             % \subsection and \subsubsection commands.

The ends  of words and sentences are marked
  by   spaces. It  doesn't matter how many
spaces    you type; one is as good as 100.  The
end of   a line counts as a space.

One   or more   blank lines denote the  end
of  a paragraph.

Since any number of consecutive spaces are treated
like a single one, the formatting of the input
file makes no difference to
      \LaTeX,                % The \LaTeX command generates the LaTeX logo.
but it makes a difference to you.  When you use
\LaTeX, making your input file as easy to read
as possible will be a great help as you write
your document and when you change it.  This sample
file shows how you can add comments to your own input
file.

Because printing is different from typewriting,
there are a number of things that you have to do
differently when preparing an input file than if
you were just typing the document directly.
Quotation marks like
       ``this''
have to be handled specially, as do quotes within
quotes:
       ``\,`this'            % \, separates the double and single quote.
        is what I just
        wrote, not  `that'\,''.

Dashes come in three sizes: an
       intra-word
dash, a medium dash for number ranges like
       1--2,
and a punctuation
       dash---like
this.

A sentence-ending space should be larger than the
space between words within a sentence.  You
sometimes have to type special commands in
conjunction with punctuation characters to get
this right, as in the following sentence.
       Gnats, gnus, etc.\ all  % `\ ' makes an inter-word space.
       begin with G\@.         % \@ marks end-of-sentence punctuation.
You should check the spaces after periods when
reading your output to make sure you haven't
forgotten any special cases.  Generating an
ellipsis
       \ldots\               % `\ ' is needed after `\ldots' because TeX
                             % ignores spaces after command names like \ldots
                             % made from \ + letters.
                             %
                             % Note how a `%' character causes TeX to ignore
                             % the end of the input line, so these blank lines
                             % do not start a new paragraph.
                             %
with the right spacing around the periods requires
a special command.

\LaTeX\ interprets some common characters as
commands, so you must type special commands to
generate them.  These characters include the
following:
       \$ \& \% \# \{ and \}.

In printing, text is usually emphasized with an
       \emph{italic}
type style.

\begin{em}
   A long segment of text can also be emphasized
   in this way.  Text within such a segment can be
   given \emph{additional} emphasis.
\end{em}

It is sometimes necessary to prevent \LaTeX\ from
breaking a line where it might otherwise do so.
This may be at a space, as between the ``Mr.''\ and
``Jones'' in
       ``Mr.~Jones'',        % ~ produces an unbreakable interword space.
or within a word---especially when the word is a
symbol like
       \mbox{\emph{itemnum}}
that makes little sense when hyphenated across
lines.

Footnotes\footnote{This is an example of a footnote.}
pose no problem.

\LaTeX\ is good at typesetting mathematical formulas
like
       \( x-3y + z = 7 \)
or
       \( a_{1} > x^{2n} + y^{2n} > x' \)
or
       \( \ip{A}{B} = \sum_{i} a_{i} b_{i} \).
The spaces you type in a formula are
ignored.  Remember that a letter like
       $x$                   % $ ... $  and  \( ... \)  are equivalent
is a formula when it denotes a mathematical
symbol, and it should be typed as one.

\section{Displayed Text}

Text is displayed by indenting it from the left
margin.  Quotations are commonly displayed.  There
are short quotations
\begin{quote}
   This is a short quotation.  It consists of a
   single paragraph of text.  See how it is formatted.
\end{quote}
and longer ones.
\begin{quotation}
   This is a longer quotation.  It consists of two
   paragraphs of text, neither of which are
   particularly interesting.

   This is the second paragraph of the quotation.  It
   is just as dull as the first paragraph.
\end{quotation}
Another frequently-displayed structure is a list.
The following is an example of an \emph{itemized}
list.
\begin{itemize}
   \item This is the first item of an itemized list.
         Each item in the list is marked with a ``tick''.
         You don't have to worry about what kind of tick
         mark is used.

   \item This is the second item of the list.  It
         contains another list nested inside it.  The inner
         list is an \emph{enumerated} list.
         \begin{enumerate}
            \item This is the first item of an enumerated
                  list that is nested within the itemized list.

            \item This is the second item of the inner list.
                  \LaTeX\ allows you to nest lists deeper than
                  you really should.
         \end{enumerate}
         This is the rest of the second item of the outer
         list.  It is no more interesting than any other
         part of the item.
   \item This is the third item of the list.
\end{itemize}
You can even display poetry.
\begin{verse}
   There is an environment
    for verse \\             % The \\ command separates lines
   Whose features some poets % within a stanza.
   will curse.

                             % One or more blank lines separate stanzas.

   For instead of making\\
   Them do \emph{all} line breaking, \\
   It allows them to put too many words on a line when they'd rather be
   forced to be terse.
\end{verse}

Mathematical formulas may also be displayed.  A
displayed formula
is
one-line long; multiline
formulas require special formatting instructions.
   \[  \ip{\Gamma}{\psi'} = x'' + y^{2} + z_{i}^{n}\]
Don't start a paragraph with a displayed equation,
nor make one a paragraph by itself.

\end{document}               % End of document.

\section{Bench Huge Section 612}
% synthetic bench section 612
% This is a sample LaTeX input file.  (Version of 12 August 2004.)
%
% A '%' character causes TeX to ignore all remaining text on the line,
% and is used for comments like this one.

\documentclass{article}      % Specifies the document class

                             % The preamble begins here.
\title{An Example Document}  % Declares the document's title.
\author{Leslie Lamport}      % Declares the author's name.
\date{January 21, 1994}      % Deleting this command produces today's date.

\newcommand{\ip}[2]{(#1, #2)}
                             % Defines \ip{arg1}{arg2} to mean
                             % (arg1, arg2).

%\newcommand{\ip}[2]{\langle #1 | #2\rangle}
                             % This is an alternative definition of
                             % \ip that is commented out.

\begin{document}             % End of preamble and beginning of text.

\maketitle                   % Produces the title.

This is an example input file.  Comparing it with
the output it generates can show you how to
produce a simple document of your own.

\section{Ordinary Text}      % Produces section heading.  Lower-level
                             % sections are begun with similar
                             % \subsection and \subsubsection commands.

The ends  of words and sentences are marked
  by   spaces. It  doesn't matter how many
spaces    you type; one is as good as 100.  The
end of   a line counts as a space.

One   or more   blank lines denote the  end
of  a paragraph.

Since any number of consecutive spaces are treated
like a single one, the formatting of the input
file makes no difference to
      \LaTeX,                % The \LaTeX command generates the LaTeX logo.
but it makes a difference to you.  When you use
\LaTeX, making your input file as easy to read
as possible will be a great help as you write
your document and when you change it.  This sample
file shows how you can add comments to your own input
file.

Because printing is different from typewriting,
there are a number of things that you have to do
differently when preparing an input file than if
you were just typing the document directly.
Quotation marks like
       ``this''
have to be handled specially, as do quotes within
quotes:
       ``\,`this'            % \, separates the double and single quote.
        is what I just
        wrote, not  `that'\,''.

Dashes come in three sizes: an
       intra-word
dash, a medium dash for number ranges like
       1--2,
and a punctuation
       dash---like
this.

A sentence-ending space should be larger than the
space between words within a sentence.  You
sometimes have to type special commands in
conjunction with punctuation characters to get
this right, as in the following sentence.
       Gnats, gnus, etc.\ all  % `\ ' makes an inter-word space.
       begin with G\@.         % \@ marks end-of-sentence punctuation.
You should check the spaces after periods when
reading your output to make sure you haven't
forgotten any special cases.  Generating an
ellipsis
       \ldots\               % `\ ' is needed after `\ldots' because TeX
                             % ignores spaces after command names like \ldots
                             % made from \ + letters.
                             %
                             % Note how a `%' character causes TeX to ignore
                             % the end of the input line, so these blank lines
                             % do not start a new paragraph.
                             %
with the right spacing around the periods requires
a special command.

\LaTeX\ interprets some common characters as
commands, so you must type special commands to
generate them.  These characters include the
following:
       \$ \& \% \# \{ and \}.

In printing, text is usually emphasized with an
       \emph{italic}
type style.

\begin{em}
   A long segment of text can also be emphasized
   in this way.  Text within such a segment can be
   given \emph{additional} emphasis.
\end{em}

It is sometimes necessary to prevent \LaTeX\ from
breaking a line where it might otherwise do so.
This may be at a space, as between the ``Mr.''\ and
``Jones'' in
       ``Mr.~Jones'',        % ~ produces an unbreakable interword space.
or within a word---especially when the word is a
symbol like
       \mbox{\emph{itemnum}}
that makes little sense when hyphenated across
lines.

Footnotes\footnote{This is an example of a footnote.}
pose no problem.

\LaTeX\ is good at typesetting mathematical formulas
like
       \( x-3y + z = 7 \)
or
       \( a_{1} > x^{2n} + y^{2n} > x' \)
or
       \( \ip{A}{B} = \sum_{i} a_{i} b_{i} \).
The spaces you type in a formula are
ignored.  Remember that a letter like
       $x$                   % $ ... $  and  \( ... \)  are equivalent
is a formula when it denotes a mathematical
symbol, and it should be typed as one.

\section{Displayed Text}

Text is displayed by indenting it from the left
margin.  Quotations are commonly displayed.  There
are short quotations
\begin{quote}
   This is a short quotation.  It consists of a
   single paragraph of text.  See how it is formatted.
\end{quote}
and longer ones.
\begin{quotation}
   This is a longer quotation.  It consists of two
   paragraphs of text, neither of which are
   particularly interesting.

   This is the second paragraph of the quotation.  It
   is just as dull as the first paragraph.
\end{quotation}
Another frequently-displayed structure is a list.
The following is an example of an \emph{itemized}
list.
\begin{itemize}
   \item This is the first item of an itemized list.
         Each item in the list is marked with a ``tick''.
         You don't have to worry about what kind of tick
         mark is used.

   \item This is the second item of the list.  It
         contains another list nested inside it.  The inner
         list is an \emph{enumerated} list.
         \begin{enumerate}
            \item This is the first item of an enumerated
                  list that is nested within the itemized list.

            \item This is the second item of the inner list.
                  \LaTeX\ allows you to nest lists deeper than
                  you really should.
         \end{enumerate}
         This is the rest of the second item of the outer
         list.  It is no more interesting than any other
         part of the item.
   \item This is the third item of the list.
\end{itemize}
You can even display poetry.
\begin{verse}
   There is an environment
    for verse \\             % The \\ command separates lines
   Whose features some poets % within a stanza.
   will curse.

                             % One or more blank lines separate stanzas.

   For instead of making\\
   Them do \emph{all} line breaking, \\
   It allows them to put too many words on a line when they'd rather be
   forced to be terse.
\end{verse}

Mathematical formulas may also be displayed.  A
displayed formula
is
one-line long; multiline
formulas require special formatting instructions.
   \[  \ip{\Gamma}{\psi'} = x'' + y^{2} + z_{i}^{n}\]
Don't start a paragraph with a displayed equation,
nor make one a paragraph by itself.

\end{document}               % End of document.

\section{Bench Huge Section 613}
% synthetic bench section 613
% This is a sample LaTeX input file.  (Version of 12 August 2004.)
%
% A '%' character causes TeX to ignore all remaining text on the line,
% and is used for comments like this one.

\documentclass{article}      % Specifies the document class

                             % The preamble begins here.
\title{An Example Document}  % Declares the document's title.
\author{Leslie Lamport}      % Declares the author's name.
\date{January 21, 1994}      % Deleting this command produces today's date.

\newcommand{\ip}[2]{(#1, #2)}
                             % Defines \ip{arg1}{arg2} to mean
                             % (arg1, arg2).

%\newcommand{\ip}[2]{\langle #1 | #2\rangle}
                             % This is an alternative definition of
                             % \ip that is commented out.

\begin{document}             % End of preamble and beginning of text.

\maketitle                   % Produces the title.

This is an example input file.  Comparing it with
the output it generates can show you how to
produce a simple document of your own.

\section{Ordinary Text}      % Produces section heading.  Lower-level
                             % sections are begun with similar
                             % \subsection and \subsubsection commands.

The ends  of words and sentences are marked
  by   spaces. It  doesn't matter how many
spaces    you type; one is as good as 100.  The
end of   a line counts as a space.

One   or more   blank lines denote the  end
of  a paragraph.

Since any number of consecutive spaces are treated
like a single one, the formatting of the input
file makes no difference to
      \LaTeX,                % The \LaTeX command generates the LaTeX logo.
but it makes a difference to you.  When you use
\LaTeX, making your input file as easy to read
as possible will be a great help as you write
your document and when you change it.  This sample
file shows how you can add comments to your own input
file.

Because printing is different from typewriting,
there are a number of things that you have to do
differently when preparing an input file than if
you were just typing the document directly.
Quotation marks like
       ``this''
have to be handled specially, as do quotes within
quotes:
       ``\,`this'            % \, separates the double and single quote.
        is what I just
        wrote, not  `that'\,''.

Dashes come in three sizes: an
       intra-word
dash, a medium dash for number ranges like
       1--2,
and a punctuation
       dash---like
this.

A sentence-ending space should be larger than the
space between words within a sentence.  You
sometimes have to type special commands in
conjunction with punctuation characters to get
this right, as in the following sentence.
       Gnats, gnus, etc.\ all  % `\ ' makes an inter-word space.
       begin with G\@.         % \@ marks end-of-sentence punctuation.
You should check the spaces after periods when
reading your output to make sure you haven't
forgotten any special cases.  Generating an
ellipsis
       \ldots\               % `\ ' is needed after `\ldots' because TeX
                             % ignores spaces after command names like \ldots
                             % made from \ + letters.
                             %
                             % Note how a `%' character causes TeX to ignore
                             % the end of the input line, so these blank lines
                             % do not start a new paragraph.
                             %
with the right spacing around the periods requires
a special command.

\LaTeX\ interprets some common characters as
commands, so you must type special commands to
generate them.  These characters include the
following:
       \$ \& \% \# \{ and \}.

In printing, text is usually emphasized with an
       \emph{italic}
type style.

\begin{em}
   A long segment of text can also be emphasized
   in this way.  Text within such a segment can be
   given \emph{additional} emphasis.
\end{em}

It is sometimes necessary to prevent \LaTeX\ from
breaking a line where it might otherwise do so.
This may be at a space, as between the ``Mr.''\ and
``Jones'' in
       ``Mr.~Jones'',        % ~ produces an unbreakable interword space.
or within a word---especially when the word is a
symbol like
       \mbox{\emph{itemnum}}
that makes little sense when hyphenated across
lines.

Footnotes\footnote{This is an example of a footnote.}
pose no problem.

\LaTeX\ is good at typesetting mathematical formulas
like
       \( x-3y + z = 7 \)
or
       \( a_{1} > x^{2n} + y^{2n} > x' \)
or
       \( \ip{A}{B} = \sum_{i} a_{i} b_{i} \).
The spaces you type in a formula are
ignored.  Remember that a letter like
       $x$                   % $ ... $  and  \( ... \)  are equivalent
is a formula when it denotes a mathematical
symbol, and it should be typed as one.

\section{Displayed Text}

Text is displayed by indenting it from the left
margin.  Quotations are commonly displayed.  There
are short quotations
\begin{quote}
   This is a short quotation.  It consists of a
   single paragraph of text.  See how it is formatted.
\end{quote}
and longer ones.
\begin{quotation}
   This is a longer quotation.  It consists of two
   paragraphs of text, neither of which are
   particularly interesting.

   This is the second paragraph of the quotation.  It
   is just as dull as the first paragraph.
\end{quotation}
Another frequently-displayed structure is a list.
The following is an example of an \emph{itemized}
list.
\begin{itemize}
   \item This is the first item of an itemized list.
         Each item in the list is marked with a ``tick''.
         You don't have to worry about what kind of tick
         mark is used.

   \item This is the second item of the list.  It
         contains another list nested inside it.  The inner
         list is an \emph{enumerated} list.
         \begin{enumerate}
            \item This is the first item of an enumerated
                  list that is nested within the itemized list.

            \item This is the second item of the inner list.
                  \LaTeX\ allows you to nest lists deeper than
                  you really should.
         \end{enumerate}
         This is the rest of the second item of the outer
         list.  It is no more interesting than any other
         part of the item.
   \item This is the third item of the list.
\end{itemize}
You can even display poetry.
\begin{verse}
   There is an environment
    for verse \\             % The \\ command separates lines
   Whose features some poets % within a stanza.
   will curse.

                             % One or more blank lines separate stanzas.

   For instead of making\\
   Them do \emph{all} line breaking, \\
   It allows them to put too many words on a line when they'd rather be
   forced to be terse.
\end{verse}

Mathematical formulas may also be displayed.  A
displayed formula
is
one-line long; multiline
formulas require special formatting instructions.
   \[  \ip{\Gamma}{\psi'} = x'' + y^{2} + z_{i}^{n}\]
Don't start a paragraph with a displayed equation,
nor make one a paragraph by itself.

\end{document}               % End of document.

\section{Bench Huge Section 614}
% synthetic bench section 614
% This is a sample LaTeX input file.  (Version of 12 August 2004.)
%
% A '%' character causes TeX to ignore all remaining text on the line,
% and is used for comments like this one.

\documentclass{article}      % Specifies the document class

                             % The preamble begins here.
\title{An Example Document}  % Declares the document's title.
\author{Leslie Lamport}      % Declares the author's name.
\date{January 21, 1994}      % Deleting this command produces today's date.

\newcommand{\ip}[2]{(#1, #2)}
                             % Defines \ip{arg1}{arg2} to mean
                             % (arg1, arg2).

%\newcommand{\ip}[2]{\langle #1 | #2\rangle}
                             % This is an alternative definition of
                             % \ip that is commented out.

\begin{document}             % End of preamble and beginning of text.

\maketitle                   % Produces the title.

This is an example input file.  Comparing it with
the output it generates can show you how to
produce a simple document of your own.

\section{Ordinary Text}      % Produces section heading.  Lower-level
                             % sections are begun with similar
                             % \subsection and \subsubsection commands.

The ends  of words and sentences are marked
  by   spaces. It  doesn't matter how many
spaces    you type; one is as good as 100.  The
end of   a line counts as a space.

One   or more   blank lines denote the  end
of  a paragraph.

Since any number of consecutive spaces are treated
like a single one, the formatting of the input
file makes no difference to
      \LaTeX,                % The \LaTeX command generates the LaTeX logo.
but it makes a difference to you.  When you use
\LaTeX, making your input file as easy to read
as possible will be a great help as you write
your document and when you change it.  This sample
file shows how you can add comments to your own input
file.

Because printing is different from typewriting,
there are a number of things that you have to do
differently when preparing an input file than if
you were just typing the document directly.
Quotation marks like
       ``this''
have to be handled specially, as do quotes within
quotes:
       ``\,`this'            % \, separates the double and single quote.
        is what I just
        wrote, not  `that'\,''.

Dashes come in three sizes: an
       intra-word
dash, a medium dash for number ranges like
       1--2,
and a punctuation
       dash---like
this.

A sentence-ending space should be larger than the
space between words within a sentence.  You
sometimes have to type special commands in
conjunction with punctuation characters to get
this right, as in the following sentence.
       Gnats, gnus, etc.\ all  % `\ ' makes an inter-word space.
       begin with G\@.         % \@ marks end-of-sentence punctuation.
You should check the spaces after periods when
reading your output to make sure you haven't
forgotten any special cases.  Generating an
ellipsis
       \ldots\               % `\ ' is needed after `\ldots' because TeX
                             % ignores spaces after command names like \ldots
                             % made from \ + letters.
                             %
                             % Note how a `%' character causes TeX to ignore
                             % the end of the input line, so these blank lines
                             % do not start a new paragraph.
                             %
with the right spacing around the periods requires
a special command.

\LaTeX\ interprets some common characters as
commands, so you must type special commands to
generate them.  These characters include the
following:
       \$ \& \% \# \{ and \}.

In printing, text is usually emphasized with an
       \emph{italic}
type style.

\begin{em}
   A long segment of text can also be emphasized
   in this way.  Text within such a segment can be
   given \emph{additional} emphasis.
\end{em}

It is sometimes necessary to prevent \LaTeX\ from
breaking a line where it might otherwise do so.
This may be at a space, as between the ``Mr.''\ and
``Jones'' in
       ``Mr.~Jones'',        % ~ produces an unbreakable interword space.
or within a word---especially when the word is a
symbol like
       \mbox{\emph{itemnum}}
that makes little sense when hyphenated across
lines.

Footnotes\footnote{This is an example of a footnote.}
pose no problem.

\LaTeX\ is good at typesetting mathematical formulas
like
       \( x-3y + z = 7 \)
or
       \( a_{1} > x^{2n} + y^{2n} > x' \)
or
       \( \ip{A}{B} = \sum_{i} a_{i} b_{i} \).
The spaces you type in a formula are
ignored.  Remember that a letter like
       $x$                   % $ ... $  and  \( ... \)  are equivalent
is a formula when it denotes a mathematical
symbol, and it should be typed as one.

\section{Displayed Text}

Text is displayed by indenting it from the left
margin.  Quotations are commonly displayed.  There
are short quotations
\begin{quote}
   This is a short quotation.  It consists of a
   single paragraph of text.  See how it is formatted.
\end{quote}
and longer ones.
\begin{quotation}
   This is a longer quotation.  It consists of two
   paragraphs of text, neither of which are
   particularly interesting.

   This is the second paragraph of the quotation.  It
   is just as dull as the first paragraph.
\end{quotation}
Another frequently-displayed structure is a list.
The following is an example of an \emph{itemized}
list.
\begin{itemize}
   \item This is the first item of an itemized list.
         Each item in the list is marked with a ``tick''.
         You don't have to worry about what kind of tick
         mark is used.

   \item This is the second item of the list.  It
         contains another list nested inside it.  The inner
         list is an \emph{enumerated} list.
         \begin{enumerate}
            \item This is the first item of an enumerated
                  list that is nested within the itemized list.

            \item This is the second item of the inner list.
                  \LaTeX\ allows you to nest lists deeper than
                  you really should.
         \end{enumerate}
         This is the rest of the second item of the outer
         list.  It is no more interesting than any other
         part of the item.
   \item This is the third item of the list.
\end{itemize}
You can even display poetry.
\begin{verse}
   There is an environment
    for verse \\             % The \\ command separates lines
   Whose features some poets % within a stanza.
   will curse.

                             % One or more blank lines separate stanzas.

   For instead of making\\
   Them do \emph{all} line breaking, \\
   It allows them to put too many words on a line when they'd rather be
   forced to be terse.
\end{verse}

Mathematical formulas may also be displayed.  A
displayed formula
is
one-line long; multiline
formulas require special formatting instructions.
   \[  \ip{\Gamma}{\psi'} = x'' + y^{2} + z_{i}^{n}\]
Don't start a paragraph with a displayed equation,
nor make one a paragraph by itself.

\end{document}               % End of document.

\section{Bench Huge Section 615}
% synthetic bench section 615
% This is a sample LaTeX input file.  (Version of 12 August 2004.)
%
% A '%' character causes TeX to ignore all remaining text on the line,
% and is used for comments like this one.

\documentclass{article}      % Specifies the document class

                             % The preamble begins here.
\title{An Example Document}  % Declares the document's title.
\author{Leslie Lamport}      % Declares the author's name.
\date{January 21, 1994}      % Deleting this command produces today's date.

\newcommand{\ip}[2]{(#1, #2)}
                             % Defines \ip{arg1}{arg2} to mean
                             % (arg1, arg2).

%\newcommand{\ip}[2]{\langle #1 | #2\rangle}
                             % This is an alternative definition of
                             % \ip that is commented out.

\begin{document}             % End of preamble and beginning of text.

\maketitle                   % Produces the title.

This is an example input file.  Comparing it with
the output it generates can show you how to
produce a simple document of your own.

\section{Ordinary Text}      % Produces section heading.  Lower-level
                             % sections are begun with similar
                             % \subsection and \subsubsection commands.

The ends  of words and sentences are marked
  by   spaces. It  doesn't matter how many
spaces    you type; one is as good as 100.  The
end of   a line counts as a space.

One   or more   blank lines denote the  end
of  a paragraph.

Since any number of consecutive spaces are treated
like a single one, the formatting of the input
file makes no difference to
      \LaTeX,                % The \LaTeX command generates the LaTeX logo.
but it makes a difference to you.  When you use
\LaTeX, making your input file as easy to read
as possible will be a great help as you write
your document and when you change it.  This sample
file shows how you can add comments to your own input
file.

Because printing is different from typewriting,
there are a number of things that you have to do
differently when preparing an input file than if
you were just typing the document directly.
Quotation marks like
       ``this''
have to be handled specially, as do quotes within
quotes:
       ``\,`this'            % \, separates the double and single quote.
        is what I just
        wrote, not  `that'\,''.

Dashes come in three sizes: an
       intra-word
dash, a medium dash for number ranges like
       1--2,
and a punctuation
       dash---like
this.

A sentence-ending space should be larger than the
space between words within a sentence.  You
sometimes have to type special commands in
conjunction with punctuation characters to get
this right, as in the following sentence.
       Gnats, gnus, etc.\ all  % `\ ' makes an inter-word space.
       begin with G\@.         % \@ marks end-of-sentence punctuation.
You should check the spaces after periods when
reading your output to make sure you haven't
forgotten any special cases.  Generating an
ellipsis
       \ldots\               % `\ ' is needed after `\ldots' because TeX
                             % ignores spaces after command names like \ldots
                             % made from \ + letters.
                             %
                             % Note how a `%' character causes TeX to ignore
                             % the end of the input line, so these blank lines
                             % do not start a new paragraph.
                             %
with the right spacing around the periods requires
a special command.

\LaTeX\ interprets some common characters as
commands, so you must type special commands to
generate them.  These characters include the
following:
       \$ \& \% \# \{ and \}.

In printing, text is usually emphasized with an
       \emph{italic}
type style.

\begin{em}
   A long segment of text can also be emphasized
   in this way.  Text within such a segment can be
   given \emph{additional} emphasis.
\end{em}

It is sometimes necessary to prevent \LaTeX\ from
breaking a line where it might otherwise do so.
This may be at a space, as between the ``Mr.''\ and
``Jones'' in
       ``Mr.~Jones'',        % ~ produces an unbreakable interword space.
or within a word---especially when the word is a
symbol like
       \mbox{\emph{itemnum}}
that makes little sense when hyphenated across
lines.

Footnotes\footnote{This is an example of a footnote.}
pose no problem.

\LaTeX\ is good at typesetting mathematical formulas
like
       \( x-3y + z = 7 \)
or
       \( a_{1} > x^{2n} + y^{2n} > x' \)
or
       \( \ip{A}{B} = \sum_{i} a_{i} b_{i} \).
The spaces you type in a formula are
ignored.  Remember that a letter like
       $x$                   % $ ... $  and  \( ... \)  are equivalent
is a formula when it denotes a mathematical
symbol, and it should be typed as one.

\section{Displayed Text}

Text is displayed by indenting it from the left
margin.  Quotations are commonly displayed.  There
are short quotations
\begin{quote}
   This is a short quotation.  It consists of a
   single paragraph of text.  See how it is formatted.
\end{quote}
and longer ones.
\begin{quotation}
   This is a longer quotation.  It consists of two
   paragraphs of text, neither of which are
   particularly interesting.

   This is the second paragraph of the quotation.  It
   is just as dull as the first paragraph.
\end{quotation}
Another frequently-displayed structure is a list.
The following is an example of an \emph{itemized}
list.
\begin{itemize}
   \item This is the first item of an itemized list.
         Each item in the list is marked with a ``tick''.
         You don't have to worry about what kind of tick
         mark is used.

   \item This is the second item of the list.  It
         contains another list nested inside it.  The inner
         list is an \emph{enumerated} list.
         \begin{enumerate}
            \item This is the first item of an enumerated
                  list that is nested within the itemized list.

            \item This is the second item of the inner list.
                  \LaTeX\ allows you to nest lists deeper than
                  you really should.
         \end{enumerate}
         This is the rest of the second item of the outer
         list.  It is no more interesting than any other
         part of the item.
   \item This is the third item of the list.
\end{itemize}
You can even display poetry.
\begin{verse}
   There is an environment
    for verse \\             % The \\ command separates lines
   Whose features some poets % within a stanza.
   will curse.

                             % One or more blank lines separate stanzas.

   For instead of making\\
   Them do \emph{all} line breaking, \\
   It allows them to put too many words on a line when they'd rather be
   forced to be terse.
\end{verse}

Mathematical formulas may also be displayed.  A
displayed formula
is
one-line long; multiline
formulas require special formatting instructions.
   \[  \ip{\Gamma}{\psi'} = x'' + y^{2} + z_{i}^{n}\]
Don't start a paragraph with a displayed equation,
nor make one a paragraph by itself.

\end{document}               % End of document.

\section{Bench Huge Section 616}
% synthetic bench section 616
% This is a sample LaTeX input file.  (Version of 12 August 2004.)
%
% A '%' character causes TeX to ignore all remaining text on the line,
% and is used for comments like this one.

\documentclass{article}      % Specifies the document class

                             % The preamble begins here.
\title{An Example Document}  % Declares the document's title.
\author{Leslie Lamport}      % Declares the author's name.
\date{January 21, 1994}      % Deleting this command produces today's date.

\newcommand{\ip}[2]{(#1, #2)}
                             % Defines \ip{arg1}{arg2} to mean
                             % (arg1, arg2).

%\newcommand{\ip}[2]{\langle #1 | #2\rangle}
                             % This is an alternative definition of
                             % \ip that is commented out.

\begin{document}             % End of preamble and beginning of text.

\maketitle                   % Produces the title.

This is an example input file.  Comparing it with
the output it generates can show you how to
produce a simple document of your own.

\section{Ordinary Text}      % Produces section heading.  Lower-level
                             % sections are begun with similar
                             % \subsection and \subsubsection commands.

The ends  of words and sentences are marked
  by   spaces. It  doesn't matter how many
spaces    you type; one is as good as 100.  The
end of   a line counts as a space.

One   or more   blank lines denote the  end
of  a paragraph.

Since any number of consecutive spaces are treated
like a single one, the formatting of the input
file makes no difference to
      \LaTeX,                % The \LaTeX command generates the LaTeX logo.
but it makes a difference to you.  When you use
\LaTeX, making your input file as easy to read
as possible will be a great help as you write
your document and when you change it.  This sample
file shows how you can add comments to your own input
file.

Because printing is different from typewriting,
there are a number of things that you have to do
differently when preparing an input file than if
you were just typing the document directly.
Quotation marks like
       ``this''
have to be handled specially, as do quotes within
quotes:
       ``\,`this'            % \, separates the double and single quote.
        is what I just
        wrote, not  `that'\,''.

Dashes come in three sizes: an
       intra-word
dash, a medium dash for number ranges like
       1--2,
and a punctuation
       dash---like
this.

A sentence-ending space should be larger than the
space between words within a sentence.  You
sometimes have to type special commands in
conjunction with punctuation characters to get
this right, as in the following sentence.
       Gnats, gnus, etc.\ all  % `\ ' makes an inter-word space.
       begin with G\@.         % \@ marks end-of-sentence punctuation.
You should check the spaces after periods when
reading your output to make sure you haven't
forgotten any special cases.  Generating an
ellipsis
       \ldots\               % `\ ' is needed after `\ldots' because TeX
                             % ignores spaces after command names like \ldots
                             % made from \ + letters.
                             %
                             % Note how a `%' character causes TeX to ignore
                             % the end of the input line, so these blank lines
                             % do not start a new paragraph.
                             %
with the right spacing around the periods requires
a special command.

\LaTeX\ interprets some common characters as
commands, so you must type special commands to
generate them.  These characters include the
following:
       \$ \& \% \# \{ and \}.

In printing, text is usually emphasized with an
       \emph{italic}
type style.

\begin{em}
   A long segment of text can also be emphasized
   in this way.  Text within such a segment can be
   given \emph{additional} emphasis.
\end{em}

It is sometimes necessary to prevent \LaTeX\ from
breaking a line where it might otherwise do so.
This may be at a space, as between the ``Mr.''\ and
``Jones'' in
       ``Mr.~Jones'',        % ~ produces an unbreakable interword space.
or within a word---especially when the word is a
symbol like
       \mbox{\emph{itemnum}}
that makes little sense when hyphenated across
lines.

Footnotes\footnote{This is an example of a footnote.}
pose no problem.

\LaTeX\ is good at typesetting mathematical formulas
like
       \( x-3y + z = 7 \)
or
       \( a_{1} > x^{2n} + y^{2n} > x' \)
or
       \( \ip{A}{B} = \sum_{i} a_{i} b_{i} \).
The spaces you type in a formula are
ignored.  Remember that a letter like
       $x$                   % $ ... $  and  \( ... \)  are equivalent
is a formula when it denotes a mathematical
symbol, and it should be typed as one.

\section{Displayed Text}

Text is displayed by indenting it from the left
margin.  Quotations are commonly displayed.  There
are short quotations
\begin{quote}
   This is a short quotation.  It consists of a
   single paragraph of text.  See how it is formatted.
\end{quote}
and longer ones.
\begin{quotation}
   This is a longer quotation.  It consists of two
   paragraphs of text, neither of which are
   particularly interesting.

   This is the second paragraph of the quotation.  It
   is just as dull as the first paragraph.
\end{quotation}
Another frequently-displayed structure is a list.
The following is an example of an \emph{itemized}
list.
\begin{itemize}
   \item This is the first item of an itemized list.
         Each item in the list is marked with a ``tick''.
         You don't have to worry about what kind of tick
         mark is used.

   \item This is the second item of the list.  It
         contains another list nested inside it.  The inner
         list is an \emph{enumerated} list.
         \begin{enumerate}
            \item This is the first item of an enumerated
                  list that is nested within the itemized list.

            \item This is the second item of the inner list.
                  \LaTeX\ allows you to nest lists deeper than
                  you really should.
         \end{enumerate}
         This is the rest of the second item of the outer
         list.  It is no more interesting than any other
         part of the item.
   \item This is the third item of the list.
\end{itemize}
You can even display poetry.
\begin{verse}
   There is an environment
    for verse \\             % The \\ command separates lines
   Whose features some poets % within a stanza.
   will curse.

                             % One or more blank lines separate stanzas.

   For instead of making\\
   Them do \emph{all} line breaking, \\
   It allows them to put too many words on a line when they'd rather be
   forced to be terse.
\end{verse}

Mathematical formulas may also be displayed.  A
displayed formula
is
one-line long; multiline
formulas require special formatting instructions.
   \[  \ip{\Gamma}{\psi'} = x'' + y^{2} + z_{i}^{n}\]
Don't start a paragraph with a displayed equation,
nor make one a paragraph by itself.

\end{document}               % End of document.

\section{Bench Huge Section 617}
% synthetic bench section 617
% This is a sample LaTeX input file.  (Version of 12 August 2004.)
%
% A '%' character causes TeX to ignore all remaining text on the line,
% and is used for comments like this one.

\documentclass{article}      % Specifies the document class

                             % The preamble begins here.
\title{An Example Document}  % Declares the document's title.
\author{Leslie Lamport}      % Declares the author's name.
\date{January 21, 1994}      % Deleting this command produces today's date.

\newcommand{\ip}[2]{(#1, #2)}
                             % Defines \ip{arg1}{arg2} to mean
                             % (arg1, arg2).

%\newcommand{\ip}[2]{\langle #1 | #2\rangle}
                             % This is an alternative definition of
                             % \ip that is commented out.

\begin{document}             % End of preamble and beginning of text.

\maketitle                   % Produces the title.

This is an example input file.  Comparing it with
the output it generates can show you how to
produce a simple document of your own.

\section{Ordinary Text}      % Produces section heading.  Lower-level
                             % sections are begun with similar
                             % \subsection and \subsubsection commands.

The ends  of words and sentences are marked
  by   spaces. It  doesn't matter how many
spaces    you type; one is as good as 100.  The
end of   a line counts as a space.

One   or more   blank lines denote the  end
of  a paragraph.

Since any number of consecutive spaces are treated
like a single one, the formatting of the input
file makes no difference to
      \LaTeX,                % The \LaTeX command generates the LaTeX logo.
but it makes a difference to you.  When you use
\LaTeX, making your input file as easy to read
as possible will be a great help as you write
your document and when you change it.  This sample
file shows how you can add comments to your own input
file.

Because printing is different from typewriting,
there are a number of things that you have to do
differently when preparing an input file than if
you were just typing the document directly.
Quotation marks like
       ``this''
have to be handled specially, as do quotes within
quotes:
       ``\,`this'            % \, separates the double and single quote.
        is what I just
        wrote, not  `that'\,''.

Dashes come in three sizes: an
       intra-word
dash, a medium dash for number ranges like
       1--2,
and a punctuation
       dash---like
this.

A sentence-ending space should be larger than the
space between words within a sentence.  You
sometimes have to type special commands in
conjunction with punctuation characters to get
this right, as in the following sentence.
       Gnats, gnus, etc.\ all  % `\ ' makes an inter-word space.
       begin with G\@.         % \@ marks end-of-sentence punctuation.
You should check the spaces after periods when
reading your output to make sure you haven't
forgotten any special cases.  Generating an
ellipsis
       \ldots\               % `\ ' is needed after `\ldots' because TeX
                             % ignores spaces after command names like \ldots
                             % made from \ + letters.
                             %
                             % Note how a `%' character causes TeX to ignore
                             % the end of the input line, so these blank lines
                             % do not start a new paragraph.
                             %
with the right spacing around the periods requires
a special command.

\LaTeX\ interprets some common characters as
commands, so you must type special commands to
generate them.  These characters include the
following:
       \$ \& \% \# \{ and \}.

In printing, text is usually emphasized with an
       \emph{italic}
type style.

\begin{em}
   A long segment of text can also be emphasized
   in this way.  Text within such a segment can be
   given \emph{additional} emphasis.
\end{em}

It is sometimes necessary to prevent \LaTeX\ from
breaking a line where it might otherwise do so.
This may be at a space, as between the ``Mr.''\ and
``Jones'' in
       ``Mr.~Jones'',        % ~ produces an unbreakable interword space.
or within a word---especially when the word is a
symbol like
       \mbox{\emph{itemnum}}
that makes little sense when hyphenated across
lines.

Footnotes\footnote{This is an example of a footnote.}
pose no problem.

\LaTeX\ is good at typesetting mathematical formulas
like
       \( x-3y + z = 7 \)
or
       \( a_{1} > x^{2n} + y^{2n} > x' \)
or
       \( \ip{A}{B} = \sum_{i} a_{i} b_{i} \).
The spaces you type in a formula are
ignored.  Remember that a letter like
       $x$                   % $ ... $  and  \( ... \)  are equivalent
is a formula when it denotes a mathematical
symbol, and it should be typed as one.

\section{Displayed Text}

Text is displayed by indenting it from the left
margin.  Quotations are commonly displayed.  There
are short quotations
\begin{quote}
   This is a short quotation.  It consists of a
   single paragraph of text.  See how it is formatted.
\end{quote}
and longer ones.
\begin{quotation}
   This is a longer quotation.  It consists of two
   paragraphs of text, neither of which are
   particularly interesting.

   This is the second paragraph of the quotation.  It
   is just as dull as the first paragraph.
\end{quotation}
Another frequently-displayed structure is a list.
The following is an example of an \emph{itemized}
list.
\begin{itemize}
   \item This is the first item of an itemized list.
         Each item in the list is marked with a ``tick''.
         You don't have to worry about what kind of tick
         mark is used.

   \item This is the second item of the list.  It
         contains another list nested inside it.  The inner
         list is an \emph{enumerated} list.
         \begin{enumerate}
            \item This is the first item of an enumerated
                  list that is nested within the itemized list.

            \item This is the second item of the inner list.
                  \LaTeX\ allows you to nest lists deeper than
                  you really should.
         \end{enumerate}
         This is the rest of the second item of the outer
         list.  It is no more interesting than any other
         part of the item.
   \item This is the third item of the list.
\end{itemize}
You can even display poetry.
\begin{verse}
   There is an environment
    for verse \\             % The \\ command separates lines
   Whose features some poets % within a stanza.
   will curse.

                             % One or more blank lines separate stanzas.

   For instead of making\\
   Them do \emph{all} line breaking, \\
   It allows them to put too many words on a line when they'd rather be
   forced to be terse.
\end{verse}

Mathematical formulas may also be displayed.  A
displayed formula
is
one-line long; multiline
formulas require special formatting instructions.
   \[  \ip{\Gamma}{\psi'} = x'' + y^{2} + z_{i}^{n}\]
Don't start a paragraph with a displayed equation,
nor make one a paragraph by itself.

\end{document}               % End of document.

\section{Bench Huge Section 618}
% synthetic bench section 618
% This is a sample LaTeX input file.  (Version of 12 August 2004.)
%
% A '%' character causes TeX to ignore all remaining text on the line,
% and is used for comments like this one.

\documentclass{article}      % Specifies the document class

                             % The preamble begins here.
\title{An Example Document}  % Declares the document's title.
\author{Leslie Lamport}      % Declares the author's name.
\date{January 21, 1994}      % Deleting this command produces today's date.

\newcommand{\ip}[2]{(#1, #2)}
                             % Defines \ip{arg1}{arg2} to mean
                             % (arg1, arg2).

%\newcommand{\ip}[2]{\langle #1 | #2\rangle}
                             % This is an alternative definition of
                             % \ip that is commented out.

\begin{document}             % End of preamble and beginning of text.

\maketitle                   % Produces the title.

This is an example input file.  Comparing it with
the output it generates can show you how to
produce a simple document of your own.

\section{Ordinary Text}      % Produces section heading.  Lower-level
                             % sections are begun with similar
                             % \subsection and \subsubsection commands.

The ends  of words and sentences are marked
  by   spaces. It  doesn't matter how many
spaces    you type; one is as good as 100.  The
end of   a line counts as a space.

One   or more   blank lines denote the  end
of  a paragraph.

Since any number of consecutive spaces are treated
like a single one, the formatting of the input
file makes no difference to
      \LaTeX,                % The \LaTeX command generates the LaTeX logo.
but it makes a difference to you.  When you use
\LaTeX, making your input file as easy to read
as possible will be a great help as you write
your document and when you change it.  This sample
file shows how you can add comments to your own input
file.

Because printing is different from typewriting,
there are a number of things that you have to do
differently when preparing an input file than if
you were just typing the document directly.
Quotation marks like
       ``this''
have to be handled specially, as do quotes within
quotes:
       ``\,`this'            % \, separates the double and single quote.
        is what I just
        wrote, not  `that'\,''.

Dashes come in three sizes: an
       intra-word
dash, a medium dash for number ranges like
       1--2,
and a punctuation
       dash---like
this.

A sentence-ending space should be larger than the
space between words within a sentence.  You
sometimes have to type special commands in
conjunction with punctuation characters to get
this right, as in the following sentence.
       Gnats, gnus, etc.\ all  % `\ ' makes an inter-word space.
       begin with G\@.         % \@ marks end-of-sentence punctuation.
You should check the spaces after periods when
reading your output to make sure you haven't
forgotten any special cases.  Generating an
ellipsis
       \ldots\               % `\ ' is needed after `\ldots' because TeX
                             % ignores spaces after command names like \ldots
                             % made from \ + letters.
                             %
                             % Note how a `%' character causes TeX to ignore
                             % the end of the input line, so these blank lines
                             % do not start a new paragraph.
                             %
with the right spacing around the periods requires
a special command.

\LaTeX\ interprets some common characters as
commands, so you must type special commands to
generate them.  These characters include the
following:
       \$ \& \% \# \{ and \}.

In printing, text is usually emphasized with an
       \emph{italic}
type style.

\begin{em}
   A long segment of text can also be emphasized
   in this way.  Text within such a segment can be
   given \emph{additional} emphasis.
\end{em}

It is sometimes necessary to prevent \LaTeX\ from
breaking a line where it might otherwise do so.
This may be at a space, as between the ``Mr.''\ and
``Jones'' in
       ``Mr.~Jones'',        % ~ produces an unbreakable interword space.
or within a word---especially when the word is a
symbol like
       \mbox{\emph{itemnum}}
that makes little sense when hyphenated across
lines.

Footnotes\footnote{This is an example of a footnote.}
pose no problem.

\LaTeX\ is good at typesetting mathematical formulas
like
       \( x-3y + z = 7 \)
or
       \( a_{1} > x^{2n} + y^{2n} > x' \)
or
       \( \ip{A}{B} = \sum_{i} a_{i} b_{i} \).
The spaces you type in a formula are
ignored.  Remember that a letter like
       $x$                   % $ ... $  and  \( ... \)  are equivalent
is a formula when it denotes a mathematical
symbol, and it should be typed as one.

\section{Displayed Text}

Text is displayed by indenting it from the left
margin.  Quotations are commonly displayed.  There
are short quotations
\begin{quote}
   This is a short quotation.  It consists of a
   single paragraph of text.  See how it is formatted.
\end{quote}
and longer ones.
\begin{quotation}
   This is a longer quotation.  It consists of two
   paragraphs of text, neither of which are
   particularly interesting.

   This is the second paragraph of the quotation.  It
   is just as dull as the first paragraph.
\end{quotation}
Another frequently-displayed structure is a list.
The following is an example of an \emph{itemized}
list.
\begin{itemize}
   \item This is the first item of an itemized list.
         Each item in the list is marked with a ``tick''.
         You don't have to worry about what kind of tick
         mark is used.

   \item This is the second item of the list.  It
         contains another list nested inside it.  The inner
         list is an \emph{enumerated} list.
         \begin{enumerate}
            \item This is the first item of an enumerated
                  list that is nested within the itemized list.

            \item This is the second item of the inner list.
                  \LaTeX\ allows you to nest lists deeper than
                  you really should.
         \end{enumerate}
         This is the rest of the second item of the outer
         list.  It is no more interesting than any other
         part of the item.
   \item This is the third item of the list.
\end{itemize}
You can even display poetry.
\begin{verse}
   There is an environment
    for verse \\             % The \\ command separates lines
   Whose features some poets % within a stanza.
   will curse.

                             % One or more blank lines separate stanzas.

   For instead of making\\
   Them do \emph{all} line breaking, \\
   It allows them to put too many words on a line when they'd rather be
   forced to be terse.
\end{verse}

Mathematical formulas may also be displayed.  A
displayed formula
is
one-line long; multiline
formulas require special formatting instructions.
   \[  \ip{\Gamma}{\psi'} = x'' + y^{2} + z_{i}^{n}\]
Don't start a paragraph with a displayed equation,
nor make one a paragraph by itself.

\end{document}               % End of document.

\section{Bench Huge Section 619}
% synthetic bench section 619
% This is a sample LaTeX input file.  (Version of 12 August 2004.)
%
% A '%' character causes TeX to ignore all remaining text on the line,
% and is used for comments like this one.

\documentclass{article}      % Specifies the document class

                             % The preamble begins here.
\title{An Example Document}  % Declares the document's title.
\author{Leslie Lamport}      % Declares the author's name.
\date{January 21, 1994}      % Deleting this command produces today's date.

\newcommand{\ip}[2]{(#1, #2)}
                             % Defines \ip{arg1}{arg2} to mean
                             % (arg1, arg2).

%\newcommand{\ip}[2]{\langle #1 | #2\rangle}
                             % This is an alternative definition of
                             % \ip that is commented out.

\begin{document}             % End of preamble and beginning of text.

\maketitle                   % Produces the title.

This is an example input file.  Comparing it with
the output it generates can show you how to
produce a simple document of your own.

\section{Ordinary Text}      % Produces section heading.  Lower-level
                             % sections are begun with similar
                             % \subsection and \subsubsection commands.

The ends  of words and sentences are marked
  by   spaces. It  doesn't matter how many
spaces    you type; one is as good as 100.  The
end of   a line counts as a space.

One   or more   blank lines denote the  end
of  a paragraph.

Since any number of consecutive spaces are treated
like a single one, the formatting of the input
file makes no difference to
      \LaTeX,                % The \LaTeX command generates the LaTeX logo.
but it makes a difference to you.  When you use
\LaTeX, making your input file as easy to read
as possible will be a great help as you write
your document and when you change it.  This sample
file shows how you can add comments to your own input
file.

Because printing is different from typewriting,
there are a number of things that you have to do
differently when preparing an input file than if
you were just typing the document directly.
Quotation marks like
       ``this''
have to be handled specially, as do quotes within
quotes:
       ``\,`this'            % \, separates the double and single quote.
        is what I just
        wrote, not  `that'\,''.

Dashes come in three sizes: an
       intra-word
dash, a medium dash for number ranges like
       1--2,
and a punctuation
       dash---like
this.

A sentence-ending space should be larger than the
space between words within a sentence.  You
sometimes have to type special commands in
conjunction with punctuation characters to get
this right, as in the following sentence.
       Gnats, gnus, etc.\ all  % `\ ' makes an inter-word space.
       begin with G\@.         % \@ marks end-of-sentence punctuation.
You should check the spaces after periods when
reading your output to make sure you haven't
forgotten any special cases.  Generating an
ellipsis
       \ldots\               % `\ ' is needed after `\ldots' because TeX
                             % ignores spaces after command names like \ldots
                             % made from \ + letters.
                             %
                             % Note how a `%' character causes TeX to ignore
                             % the end of the input line, so these blank lines
                             % do not start a new paragraph.
                             %
with the right spacing around the periods requires
a special command.

\LaTeX\ interprets some common characters as
commands, so you must type special commands to
generate them.  These characters include the
following:
       \$ \& \% \# \{ and \}.

In printing, text is usually emphasized with an
       \emph{italic}
type style.

\begin{em}
   A long segment of text can also be emphasized
   in this way.  Text within such a segment can be
   given \emph{additional} emphasis.
\end{em}

It is sometimes necessary to prevent \LaTeX\ from
breaking a line where it might otherwise do so.
This may be at a space, as between the ``Mr.''\ and
``Jones'' in
       ``Mr.~Jones'',        % ~ produces an unbreakable interword space.
or within a word---especially when the word is a
symbol like
       \mbox{\emph{itemnum}}
that makes little sense when hyphenated across
lines.

Footnotes\footnote{This is an example of a footnote.}
pose no problem.

\LaTeX\ is good at typesetting mathematical formulas
like
       \( x-3y + z = 7 \)
or
       \( a_{1} > x^{2n} + y^{2n} > x' \)
or
       \( \ip{A}{B} = \sum_{i} a_{i} b_{i} \).
The spaces you type in a formula are
ignored.  Remember that a letter like
       $x$                   % $ ... $  and  \( ... \)  are equivalent
is a formula when it denotes a mathematical
symbol, and it should be typed as one.

\section{Displayed Text}

Text is displayed by indenting it from the left
margin.  Quotations are commonly displayed.  There
are short quotations
\begin{quote}
   This is a short quotation.  It consists of a
   single paragraph of text.  See how it is formatted.
\end{quote}
and longer ones.
\begin{quotation}
   This is a longer quotation.  It consists of two
   paragraphs of text, neither of which are
   particularly interesting.

   This is the second paragraph of the quotation.  It
   is just as dull as the first paragraph.
\end{quotation}
Another frequently-displayed structure is a list.
The following is an example of an \emph{itemized}
list.
\begin{itemize}
   \item This is the first item of an itemized list.
         Each item in the list is marked with a ``tick''.
         You don't have to worry about what kind of tick
         mark is used.

   \item This is the second item of the list.  It
         contains another list nested inside it.  The inner
         list is an \emph{enumerated} list.
         \begin{enumerate}
            \item This is the first item of an enumerated
                  list that is nested within the itemized list.

            \item This is the second item of the inner list.
                  \LaTeX\ allows you to nest lists deeper than
                  you really should.
         \end{enumerate}
         This is the rest of the second item of the outer
         list.  It is no more interesting than any other
         part of the item.
   \item This is the third item of the list.
\end{itemize}
You can even display poetry.
\begin{verse}
   There is an environment
    for verse \\             % The \\ command separates lines
   Whose features some poets % within a stanza.
   will curse.

                             % One or more blank lines separate stanzas.

   For instead of making\\
   Them do \emph{all} line breaking, \\
   It allows them to put too many words on a line when they'd rather be
   forced to be terse.
\end{verse}

Mathematical formulas may also be displayed.  A
displayed formula
is
one-line long; multiline
formulas require special formatting instructions.
   \[  \ip{\Gamma}{\psi'} = x'' + y^{2} + z_{i}^{n}\]
Don't start a paragraph with a displayed equation,
nor make one a paragraph by itself.

\end{document}               % End of document.

\section{Bench Huge Section 620}
% synthetic bench section 620
% This is a sample LaTeX input file.  (Version of 12 August 2004.)
%
% A '%' character causes TeX to ignore all remaining text on the line,
% and is used for comments like this one.

\documentclass{article}      % Specifies the document class

                             % The preamble begins here.
\title{An Example Document}  % Declares the document's title.
\author{Leslie Lamport}      % Declares the author's name.
\date{January 21, 1994}      % Deleting this command produces today's date.

\newcommand{\ip}[2]{(#1, #2)}
                             % Defines \ip{arg1}{arg2} to mean
                             % (arg1, arg2).

%\newcommand{\ip}[2]{\langle #1 | #2\rangle}
                             % This is an alternative definition of
                             % \ip that is commented out.

\begin{document}             % End of preamble and beginning of text.

\maketitle                   % Produces the title.

This is an example input file.  Comparing it with
the output it generates can show you how to
produce a simple document of your own.

\section{Ordinary Text}      % Produces section heading.  Lower-level
                             % sections are begun with similar
                             % \subsection and \subsubsection commands.

The ends  of words and sentences are marked
  by   spaces. It  doesn't matter how many
spaces    you type; one is as good as 100.  The
end of   a line counts as a space.

One   or more   blank lines denote the  end
of  a paragraph.

Since any number of consecutive spaces are treated
like a single one, the formatting of the input
file makes no difference to
      \LaTeX,                % The \LaTeX command generates the LaTeX logo.
but it makes a difference to you.  When you use
\LaTeX, making your input file as easy to read
as possible will be a great help as you write
your document and when you change it.  This sample
file shows how you can add comments to your own input
file.

Because printing is different from typewriting,
there are a number of things that you have to do
differently when preparing an input file than if
you were just typing the document directly.
Quotation marks like
       ``this''
have to be handled specially, as do quotes within
quotes:
       ``\,`this'            % \, separates the double and single quote.
        is what I just
        wrote, not  `that'\,''.

Dashes come in three sizes: an
       intra-word
dash, a medium dash for number ranges like
       1--2,
and a punctuation
       dash---like
this.

A sentence-ending space should be larger than the
space between words within a sentence.  You
sometimes have to type special commands in
conjunction with punctuation characters to get
this right, as in the following sentence.
       Gnats, gnus, etc.\ all  % `\ ' makes an inter-word space.
       begin with G\@.         % \@ marks end-of-sentence punctuation.
You should check the spaces after periods when
reading your output to make sure you haven't
forgotten any special cases.  Generating an
ellipsis
       \ldots\               % `\ ' is needed after `\ldots' because TeX
                             % ignores spaces after command names like \ldots
                             % made from \ + letters.
                             %
                             % Note how a `%' character causes TeX to ignore
                             % the end of the input line, so these blank lines
                             % do not start a new paragraph.
                             %
with the right spacing around the periods requires
a special command.

\LaTeX\ interprets some common characters as
commands, so you must type special commands to
generate them.  These characters include the
following:
       \$ \& \% \# \{ and \}.

In printing, text is usually emphasized with an
       \emph{italic}
type style.

\begin{em}
   A long segment of text can also be emphasized
   in this way.  Text within such a segment can be
   given \emph{additional} emphasis.
\end{em}

It is sometimes necessary to prevent \LaTeX\ from
breaking a line where it might otherwise do so.
This may be at a space, as between the ``Mr.''\ and
``Jones'' in
       ``Mr.~Jones'',        % ~ produces an unbreakable interword space.
or within a word---especially when the word is a
symbol like
       \mbox{\emph{itemnum}}
that makes little sense when hyphenated across
lines.

Footnotes\footnote{This is an example of a footnote.}
pose no problem.

\LaTeX\ is good at typesetting mathematical formulas
like
       \( x-3y + z = 7 \)
or
       \( a_{1} > x^{2n} + y^{2n} > x' \)
or
       \( \ip{A}{B} = \sum_{i} a_{i} b_{i} \).
The spaces you type in a formula are
ignored.  Remember that a letter like
       $x$                   % $ ... $  and  \( ... \)  are equivalent
is a formula when it denotes a mathematical
symbol, and it should be typed as one.

\section{Displayed Text}

Text is displayed by indenting it from the left
margin.  Quotations are commonly displayed.  There
are short quotations
\begin{quote}
   This is a short quotation.  It consists of a
   single paragraph of text.  See how it is formatted.
\end{quote}
and longer ones.
\begin{quotation}
   This is a longer quotation.  It consists of two
   paragraphs of text, neither of which are
   particularly interesting.

   This is the second paragraph of the quotation.  It
   is just as dull as the first paragraph.
\end{quotation}
Another frequently-displayed structure is a list.
The following is an example of an \emph{itemized}
list.
\begin{itemize}
   \item This is the first item of an itemized list.
         Each item in the list is marked with a ``tick''.
         You don't have to worry about what kind of tick
         mark is used.

   \item This is the second item of the list.  It
         contains another list nested inside it.  The inner
         list is an \emph{enumerated} list.
         \begin{enumerate}
            \item This is the first item of an enumerated
                  list that is nested within the itemized list.

            \item This is the second item of the inner list.
                  \LaTeX\ allows you to nest lists deeper than
                  you really should.
         \end{enumerate}
         This is the rest of the second item of the outer
         list.  It is no more interesting than any other
         part of the item.
   \item This is the third item of the list.
\end{itemize}
You can even display poetry.
\begin{verse}
   There is an environment
    for verse \\             % The \\ command separates lines
   Whose features some poets % within a stanza.
   will curse.

                             % One or more blank lines separate stanzas.

   For instead of making\\
   Them do \emph{all} line breaking, \\
   It allows them to put too many words on a line when they'd rather be
   forced to be terse.
\end{verse}

Mathematical formulas may also be displayed.  A
displayed formula
is
one-line long; multiline
formulas require special formatting instructions.
   \[  \ip{\Gamma}{\psi'} = x'' + y^{2} + z_{i}^{n}\]
Don't start a paragraph with a displayed equation,
nor make one a paragraph by itself.

\end{document}               % End of document.

\section{Bench Huge Section 621}
% synthetic bench section 621
% This is a sample LaTeX input file.  (Version of 12 August 2004.)
%
% A '%' character causes TeX to ignore all remaining text on the line,
% and is used for comments like this one.

\documentclass{article}      % Specifies the document class

                             % The preamble begins here.
\title{An Example Document}  % Declares the document's title.
\author{Leslie Lamport}      % Declares the author's name.
\date{January 21, 1994}      % Deleting this command produces today's date.

\newcommand{\ip}[2]{(#1, #2)}
                             % Defines \ip{arg1}{arg2} to mean
                             % (arg1, arg2).

%\newcommand{\ip}[2]{\langle #1 | #2\rangle}
                             % This is an alternative definition of
                             % \ip that is commented out.

\begin{document}             % End of preamble and beginning of text.

\maketitle                   % Produces the title.

This is an example input file.  Comparing it with
the output it generates can show you how to
produce a simple document of your own.

\section{Ordinary Text}      % Produces section heading.  Lower-level
                             % sections are begun with similar
                             % \subsection and \subsubsection commands.

The ends  of words and sentences are marked
  by   spaces. It  doesn't matter how many
spaces    you type; one is as good as 100.  The
end of   a line counts as a space.

One   or more   blank lines denote the  end
of  a paragraph.

Since any number of consecutive spaces are treated
like a single one, the formatting of the input
file makes no difference to
      \LaTeX,                % The \LaTeX command generates the LaTeX logo.
but it makes a difference to you.  When you use
\LaTeX, making your input file as easy to read
as possible will be a great help as you write
your document and when you change it.  This sample
file shows how you can add comments to your own input
file.

Because printing is different from typewriting,
there are a number of things that you have to do
differently when preparing an input file than if
you were just typing the document directly.
Quotation marks like
       ``this''
have to be handled specially, as do quotes within
quotes:
       ``\,`this'            % \, separates the double and single quote.
        is what I just
        wrote, not  `that'\,''.

Dashes come in three sizes: an
       intra-word
dash, a medium dash for number ranges like
       1--2,
and a punctuation
       dash---like
this.

A sentence-ending space should be larger than the
space between words within a sentence.  You
sometimes have to type special commands in
conjunction with punctuation characters to get
this right, as in the following sentence.
       Gnats, gnus, etc.\ all  % `\ ' makes an inter-word space.
       begin with G\@.         % \@ marks end-of-sentence punctuation.
You should check the spaces after periods when
reading your output to make sure you haven't
forgotten any special cases.  Generating an
ellipsis
       \ldots\               % `\ ' is needed after `\ldots' because TeX
                             % ignores spaces after command names like \ldots
                             % made from \ + letters.
                             %
                             % Note how a `%' character causes TeX to ignore
                             % the end of the input line, so these blank lines
                             % do not start a new paragraph.
                             %
with the right spacing around the periods requires
a special command.

\LaTeX\ interprets some common characters as
commands, so you must type special commands to
generate them.  These characters include the
following:
       \$ \& \% \# \{ and \}.

In printing, text is usually emphasized with an
       \emph{italic}
type style.

\begin{em}
   A long segment of text can also be emphasized
   in this way.  Text within such a segment can be
   given \emph{additional} emphasis.
\end{em}

It is sometimes necessary to prevent \LaTeX\ from
breaking a line where it might otherwise do so.
This may be at a space, as between the ``Mr.''\ and
``Jones'' in
       ``Mr.~Jones'',        % ~ produces an unbreakable interword space.
or within a word---especially when the word is a
symbol like
       \mbox{\emph{itemnum}}
that makes little sense when hyphenated across
lines.

Footnotes\footnote{This is an example of a footnote.}
pose no problem.

\LaTeX\ is good at typesetting mathematical formulas
like
       \( x-3y + z = 7 \)
or
       \( a_{1} > x^{2n} + y^{2n} > x' \)
or
       \( \ip{A}{B} = \sum_{i} a_{i} b_{i} \).
The spaces you type in a formula are
ignored.  Remember that a letter like
       $x$                   % $ ... $  and  \( ... \)  are equivalent
is a formula when it denotes a mathematical
symbol, and it should be typed as one.

\section{Displayed Text}

Text is displayed by indenting it from the left
margin.  Quotations are commonly displayed.  There
are short quotations
\begin{quote}
   This is a short quotation.  It consists of a
   single paragraph of text.  See how it is formatted.
\end{quote}
and longer ones.
\begin{quotation}
   This is a longer quotation.  It consists of two
   paragraphs of text, neither of which are
   particularly interesting.

   This is the second paragraph of the quotation.  It
   is just as dull as the first paragraph.
\end{quotation}
Another frequently-displayed structure is a list.
The following is an example of an \emph{itemized}
list.
\begin{itemize}
   \item This is the first item of an itemized list.
         Each item in the list is marked with a ``tick''.
         You don't have to worry about what kind of tick
         mark is used.

   \item This is the second item of the list.  It
         contains another list nested inside it.  The inner
         list is an \emph{enumerated} list.
         \begin{enumerate}
            \item This is the first item of an enumerated
                  list that is nested within the itemized list.

            \item This is the second item of the inner list.
                  \LaTeX\ allows you to nest lists deeper than
                  you really should.
         \end{enumerate}
         This is the rest of the second item of the outer
         list.  It is no more interesting than any other
         part of the item.
   \item This is the third item of the list.
\end{itemize}
You can even display poetry.
\begin{verse}
   There is an environment
    for verse \\             % The \\ command separates lines
   Whose features some poets % within a stanza.
   will curse.

                             % One or more blank lines separate stanzas.

   For instead of making\\
   Them do \emph{all} line breaking, \\
   It allows them to put too many words on a line when they'd rather be
   forced to be terse.
\end{verse}

Mathematical formulas may also be displayed.  A
displayed formula
is
one-line long; multiline
formulas require special formatting instructions.
   \[  \ip{\Gamma}{\psi'} = x'' + y^{2} + z_{i}^{n}\]
Don't start a paragraph with a displayed equation,
nor make one a paragraph by itself.

\end{document}               % End of document.

\section{Bench Huge Section 622}
% synthetic bench section 622
% This is a sample LaTeX input file.  (Version of 12 August 2004.)
%
% A '%' character causes TeX to ignore all remaining text on the line,
% and is used for comments like this one.

\documentclass{article}      % Specifies the document class

                             % The preamble begins here.
\title{An Example Document}  % Declares the document's title.
\author{Leslie Lamport}      % Declares the author's name.
\date{January 21, 1994}      % Deleting this command produces today's date.

\newcommand{\ip}[2]{(#1, #2)}
                             % Defines \ip{arg1}{arg2} to mean
                             % (arg1, arg2).

%\newcommand{\ip}[2]{\langle #1 | #2\rangle}
                             % This is an alternative definition of
                             % \ip that is commented out.

\begin{document}             % End of preamble and beginning of text.

\maketitle                   % Produces the title.

This is an example input file.  Comparing it with
the output it generates can show you how to
produce a simple document of your own.

\section{Ordinary Text}      % Produces section heading.  Lower-level
                             % sections are begun with similar
                             % \subsection and \subsubsection commands.

The ends  of words and sentences are marked
  by   spaces. It  doesn't matter how many
spaces    you type; one is as good as 100.  The
end of   a line counts as a space.

One   or more   blank lines denote the  end
of  a paragraph.

Since any number of consecutive spaces are treated
like a single one, the formatting of the input
file makes no difference to
      \LaTeX,                % The \LaTeX command generates the LaTeX logo.
but it makes a difference to you.  When you use
\LaTeX, making your input file as easy to read
as possible will be a great help as you write
your document and when you change it.  This sample
file shows how you can add comments to your own input
file.

Because printing is different from typewriting,
there are a number of things that you have to do
differently when preparing an input file than if
you were just typing the document directly.
Quotation marks like
       ``this''
have to be handled specially, as do quotes within
quotes:
       ``\,`this'            % \, separates the double and single quote.
        is what I just
        wrote, not  `that'\,''.

Dashes come in three sizes: an
       intra-word
dash, a medium dash for number ranges like
       1--2,
and a punctuation
       dash---like
this.

A sentence-ending space should be larger than the
space between words within a sentence.  You
sometimes have to type special commands in
conjunction with punctuation characters to get
this right, as in the following sentence.
       Gnats, gnus, etc.\ all  % `\ ' makes an inter-word space.
       begin with G\@.         % \@ marks end-of-sentence punctuation.
You should check the spaces after periods when
reading your output to make sure you haven't
forgotten any special cases.  Generating an
ellipsis
       \ldots\               % `\ ' is needed after `\ldots' because TeX
                             % ignores spaces after command names like \ldots
                             % made from \ + letters.
                             %
                             % Note how a `%' character causes TeX to ignore
                             % the end of the input line, so these blank lines
                             % do not start a new paragraph.
                             %
with the right spacing around the periods requires
a special command.

\LaTeX\ interprets some common characters as
commands, so you must type special commands to
generate them.  These characters include the
following:
       \$ \& \% \# \{ and \}.

In printing, text is usually emphasized with an
       \emph{italic}
type style.

\begin{em}
   A long segment of text can also be emphasized
   in this way.  Text within such a segment can be
   given \emph{additional} emphasis.
\end{em}

It is sometimes necessary to prevent \LaTeX\ from
breaking a line where it might otherwise do so.
This may be at a space, as between the ``Mr.''\ and
``Jones'' in
       ``Mr.~Jones'',        % ~ produces an unbreakable interword space.
or within a word---especially when the word is a
symbol like
       \mbox{\emph{itemnum}}
that makes little sense when hyphenated across
lines.

Footnotes\footnote{This is an example of a footnote.}
pose no problem.

\LaTeX\ is good at typesetting mathematical formulas
like
       \( x-3y + z = 7 \)
or
       \( a_{1} > x^{2n} + y^{2n} > x' \)
or
       \( \ip{A}{B} = \sum_{i} a_{i} b_{i} \).
The spaces you type in a formula are
ignored.  Remember that a letter like
       $x$                   % $ ... $  and  \( ... \)  are equivalent
is a formula when it denotes a mathematical
symbol, and it should be typed as one.

\section{Displayed Text}

Text is displayed by indenting it from the left
margin.  Quotations are commonly displayed.  There
are short quotations
\begin{quote}
   This is a short quotation.  It consists of a
   single paragraph of text.  See how it is formatted.
\end{quote}
and longer ones.
\begin{quotation}
   This is a longer quotation.  It consists of two
   paragraphs of text, neither of which are
   particularly interesting.

   This is the second paragraph of the quotation.  It
   is just as dull as the first paragraph.
\end{quotation}
Another frequently-displayed structure is a list.
The following is an example of an \emph{itemized}
list.
\begin{itemize}
   \item This is the first item of an itemized list.
         Each item in the list is marked with a ``tick''.
         You don't have to worry about what kind of tick
         mark is used.

   \item This is the second item of the list.  It
         contains another list nested inside it.  The inner
         list is an \emph{enumerated} list.
         \begin{enumerate}
            \item This is the first item of an enumerated
                  list that is nested within the itemized list.

            \item This is the second item of the inner list.
                  \LaTeX\ allows you to nest lists deeper than
                  you really should.
         \end{enumerate}
         This is the rest of the second item of the outer
         list.  It is no more interesting than any other
         part of the item.
   \item This is the third item of the list.
\end{itemize}
You can even display poetry.
\begin{verse}
   There is an environment
    for verse \\             % The \\ command separates lines
   Whose features some poets % within a stanza.
   will curse.

                             % One or more blank lines separate stanzas.

   For instead of making\\
   Them do \emph{all} line breaking, \\
   It allows them to put too many words on a line when they'd rather be
   forced to be terse.
\end{verse}

Mathematical formulas may also be displayed.  A
displayed formula
is
one-line long; multiline
formulas require special formatting instructions.
   \[  \ip{\Gamma}{\psi'} = x'' + y^{2} + z_{i}^{n}\]
Don't start a paragraph with a displayed equation,
nor make one a paragraph by itself.

\end{document}               % End of document.

\section{Bench Huge Section 623}
% synthetic bench section 623
% This is a sample LaTeX input file.  (Version of 12 August 2004.)
%
% A '%' character causes TeX to ignore all remaining text on the line,
% and is used for comments like this one.

\documentclass{article}      % Specifies the document class

                             % The preamble begins here.
\title{An Example Document}  % Declares the document's title.
\author{Leslie Lamport}      % Declares the author's name.
\date{January 21, 1994}      % Deleting this command produces today's date.

\newcommand{\ip}[2]{(#1, #2)}
                             % Defines \ip{arg1}{arg2} to mean
                             % (arg1, arg2).

%\newcommand{\ip}[2]{\langle #1 | #2\rangle}
                             % This is an alternative definition of
                             % \ip that is commented out.

\begin{document}             % End of preamble and beginning of text.

\maketitle                   % Produces the title.

This is an example input file.  Comparing it with
the output it generates can show you how to
produce a simple document of your own.

\section{Ordinary Text}      % Produces section heading.  Lower-level
                             % sections are begun with similar
                             % \subsection and \subsubsection commands.

The ends  of words and sentences are marked
  by   spaces. It  doesn't matter how many
spaces    you type; one is as good as 100.  The
end of   a line counts as a space.

One   or more   blank lines denote the  end
of  a paragraph.

Since any number of consecutive spaces are treated
like a single one, the formatting of the input
file makes no difference to
      \LaTeX,                % The \LaTeX command generates the LaTeX logo.
but it makes a difference to you.  When you use
\LaTeX, making your input file as easy to read
as possible will be a great help as you write
your document and when you change it.  This sample
file shows how you can add comments to your own input
file.

Because printing is different from typewriting,
there are a number of things that you have to do
differently when preparing an input file than if
you were just typing the document directly.
Quotation marks like
       ``this''
have to be handled specially, as do quotes within
quotes:
       ``\,`this'            % \, separates the double and single quote.
        is what I just
        wrote, not  `that'\,''.

Dashes come in three sizes: an
       intra-word
dash, a medium dash for number ranges like
       1--2,
and a punctuation
       dash---like
this.

A sentence-ending space should be larger than the
space between words within a sentence.  You
sometimes have to type special commands in
conjunction with punctuation characters to get
this right, as in the following sentence.
       Gnats, gnus, etc.\ all  % `\ ' makes an inter-word space.
       begin with G\@.         % \@ marks end-of-sentence punctuation.
You should check the spaces after periods when
reading your output to make sure you haven't
forgotten any special cases.  Generating an
ellipsis
       \ldots\               % `\ ' is needed after `\ldots' because TeX
                             % ignores spaces after command names like \ldots
                             % made from \ + letters.
                             %
                             % Note how a `%' character causes TeX to ignore
                             % the end of the input line, so these blank lines
                             % do not start a new paragraph.
                             %
with the right spacing around the periods requires
a special command.

\LaTeX\ interprets some common characters as
commands, so you must type special commands to
generate them.  These characters include the
following:
       \$ \& \% \# \{ and \}.

In printing, text is usually emphasized with an
       \emph{italic}
type style.

\begin{em}
   A long segment of text can also be emphasized
   in this way.  Text within such a segment can be
   given \emph{additional} emphasis.
\end{em}

It is sometimes necessary to prevent \LaTeX\ from
breaking a line where it might otherwise do so.
This may be at a space, as between the ``Mr.''\ and
``Jones'' in
       ``Mr.~Jones'',        % ~ produces an unbreakable interword space.
or within a word---especially when the word is a
symbol like
       \mbox{\emph{itemnum}}
that makes little sense when hyphenated across
lines.

Footnotes\footnote{This is an example of a footnote.}
pose no problem.

\LaTeX\ is good at typesetting mathematical formulas
like
       \( x-3y + z = 7 \)
or
       \( a_{1} > x^{2n} + y^{2n} > x' \)
or
       \( \ip{A}{B} = \sum_{i} a_{i} b_{i} \).
The spaces you type in a formula are
ignored.  Remember that a letter like
       $x$                   % $ ... $  and  \( ... \)  are equivalent
is a formula when it denotes a mathematical
symbol, and it should be typed as one.

\section{Displayed Text}

Text is displayed by indenting it from the left
margin.  Quotations are commonly displayed.  There
are short quotations
\begin{quote}
   This is a short quotation.  It consists of a
   single paragraph of text.  See how it is formatted.
\end{quote}
and longer ones.
\begin{quotation}
   This is a longer quotation.  It consists of two
   paragraphs of text, neither of which are
   particularly interesting.

   This is the second paragraph of the quotation.  It
   is just as dull as the first paragraph.
\end{quotation}
Another frequently-displayed structure is a list.
The following is an example of an \emph{itemized}
list.
\begin{itemize}
   \item This is the first item of an itemized list.
         Each item in the list is marked with a ``tick''.
         You don't have to worry about what kind of tick
         mark is used.

   \item This is the second item of the list.  It
         contains another list nested inside it.  The inner
         list is an \emph{enumerated} list.
         \begin{enumerate}
            \item This is the first item of an enumerated
                  list that is nested within the itemized list.

            \item This is the second item of the inner list.
                  \LaTeX\ allows you to nest lists deeper than
                  you really should.
         \end{enumerate}
         This is the rest of the second item of the outer
         list.  It is no more interesting than any other
         part of the item.
   \item This is the third item of the list.
\end{itemize}
You can even display poetry.
\begin{verse}
   There is an environment
    for verse \\             % The \\ command separates lines
   Whose features some poets % within a stanza.
   will curse.

                             % One or more blank lines separate stanzas.

   For instead of making\\
   Them do \emph{all} line breaking, \\
   It allows them to put too many words on a line when they'd rather be
   forced to be terse.
\end{verse}

Mathematical formulas may also be displayed.  A
displayed formula
is
one-line long; multiline
formulas require special formatting instructions.
   \[  \ip{\Gamma}{\psi'} = x'' + y^{2} + z_{i}^{n}\]
Don't start a paragraph with a displayed equation,
nor make one a paragraph by itself.

\end{document}               % End of document.

\section{Bench Huge Section 624}
% synthetic bench section 624
% This is a sample LaTeX input file.  (Version of 12 August 2004.)
%
% A '%' character causes TeX to ignore all remaining text on the line,
% and is used for comments like this one.

\documentclass{article}      % Specifies the document class

                             % The preamble begins here.
\title{An Example Document}  % Declares the document's title.
\author{Leslie Lamport}      % Declares the author's name.
\date{January 21, 1994}      % Deleting this command produces today's date.

\newcommand{\ip}[2]{(#1, #2)}
                             % Defines \ip{arg1}{arg2} to mean
                             % (arg1, arg2).

%\newcommand{\ip}[2]{\langle #1 | #2\rangle}
                             % This is an alternative definition of
                             % \ip that is commented out.

\begin{document}             % End of preamble and beginning of text.

\maketitle                   % Produces the title.

This is an example input file.  Comparing it with
the output it generates can show you how to
produce a simple document of your own.

\section{Ordinary Text}      % Produces section heading.  Lower-level
                             % sections are begun with similar
                             % \subsection and \subsubsection commands.

The ends  of words and sentences are marked
  by   spaces. It  doesn't matter how many
spaces    you type; one is as good as 100.  The
end of   a line counts as a space.

One   or more   blank lines denote the  end
of  a paragraph.

Since any number of consecutive spaces are treated
like a single one, the formatting of the input
file makes no difference to
      \LaTeX,                % The \LaTeX command generates the LaTeX logo.
but it makes a difference to you.  When you use
\LaTeX, making your input file as easy to read
as possible will be a great help as you write
your document and when you change it.  This sample
file shows how you can add comments to your own input
file.

Because printing is different from typewriting,
there are a number of things that you have to do
differently when preparing an input file than if
you were just typing the document directly.
Quotation marks like
       ``this''
have to be handled specially, as do quotes within
quotes:
       ``\,`this'            % \, separates the double and single quote.
        is what I just
        wrote, not  `that'\,''.

Dashes come in three sizes: an
       intra-word
dash, a medium dash for number ranges like
       1--2,
and a punctuation
       dash---like
this.

A sentence-ending space should be larger than the
space between words within a sentence.  You
sometimes have to type special commands in
conjunction with punctuation characters to get
this right, as in the following sentence.
       Gnats, gnus, etc.\ all  % `\ ' makes an inter-word space.
       begin with G\@.         % \@ marks end-of-sentence punctuation.
You should check the spaces after periods when
reading your output to make sure you haven't
forgotten any special cases.  Generating an
ellipsis
       \ldots\               % `\ ' is needed after `\ldots' because TeX
                             % ignores spaces after command names like \ldots
                             % made from \ + letters.
                             %
                             % Note how a `%' character causes TeX to ignore
                             % the end of the input line, so these blank lines
                             % do not start a new paragraph.
                             %
with the right spacing around the periods requires
a special command.

\LaTeX\ interprets some common characters as
commands, so you must type special commands to
generate them.  These characters include the
following:
       \$ \& \% \# \{ and \}.

In printing, text is usually emphasized with an
       \emph{italic}
type style.

\begin{em}
   A long segment of text can also be emphasized
   in this way.  Text within such a segment can be
   given \emph{additional} emphasis.
\end{em}

It is sometimes necessary to prevent \LaTeX\ from
breaking a line where it might otherwise do so.
This may be at a space, as between the ``Mr.''\ and
``Jones'' in
       ``Mr.~Jones'',        % ~ produces an unbreakable interword space.
or within a word---especially when the word is a
symbol like
       \mbox{\emph{itemnum}}
that makes little sense when hyphenated across
lines.

Footnotes\footnote{This is an example of a footnote.}
pose no problem.

\LaTeX\ is good at typesetting mathematical formulas
like
       \( x-3y + z = 7 \)
or
       \( a_{1} > x^{2n} + y^{2n} > x' \)
or
       \( \ip{A}{B} = \sum_{i} a_{i} b_{i} \).
The spaces you type in a formula are
ignored.  Remember that a letter like
       $x$                   % $ ... $  and  \( ... \)  are equivalent
is a formula when it denotes a mathematical
symbol, and it should be typed as one.

\section{Displayed Text}

Text is displayed by indenting it from the left
margin.  Quotations are commonly displayed.  There
are short quotations
\begin{quote}
   This is a short quotation.  It consists of a
   single paragraph of text.  See how it is formatted.
\end{quote}
and longer ones.
\begin{quotation}
   This is a longer quotation.  It consists of two
   paragraphs of text, neither of which are
   particularly interesting.

   This is the second paragraph of the quotation.  It
   is just as dull as the first paragraph.
\end{quotation}
Another frequently-displayed structure is a list.
The following is an example of an \emph{itemized}
list.
\begin{itemize}
   \item This is the first item of an itemized list.
         Each item in the list is marked with a ``tick''.
         You don't have to worry about what kind of tick
         mark is used.

   \item This is the second item of the list.  It
         contains another list nested inside it.  The inner
         list is an \emph{enumerated} list.
         \begin{enumerate}
            \item This is the first item of an enumerated
                  list that is nested within the itemized list.

            \item This is the second item of the inner list.
                  \LaTeX\ allows you to nest lists deeper than
                  you really should.
         \end{enumerate}
         This is the rest of the second item of the outer
         list.  It is no more interesting than any other
         part of the item.
   \item This is the third item of the list.
\end{itemize}
You can even display poetry.
\begin{verse}
   There is an environment
    for verse \\             % The \\ command separates lines
   Whose features some poets % within a stanza.
   will curse.

                             % One or more blank lines separate stanzas.

   For instead of making\\
   Them do \emph{all} line breaking, \\
   It allows them to put too many words on a line when they'd rather be
   forced to be terse.
\end{verse}

Mathematical formulas may also be displayed.  A
displayed formula
is
one-line long; multiline
formulas require special formatting instructions.
   \[  \ip{\Gamma}{\psi'} = x'' + y^{2} + z_{i}^{n}\]
Don't start a paragraph with a displayed equation,
nor make one a paragraph by itself.

\end{document}               % End of document.

\section{Bench Huge Section 625}
% synthetic bench section 625
% This is a sample LaTeX input file.  (Version of 12 August 2004.)
%
% A '%' character causes TeX to ignore all remaining text on the line,
% and is used for comments like this one.

\documentclass{article}      % Specifies the document class

                             % The preamble begins here.
\title{An Example Document}  % Declares the document's title.
\author{Leslie Lamport}      % Declares the author's name.
\date{January 21, 1994}      % Deleting this command produces today's date.

\newcommand{\ip}[2]{(#1, #2)}
                             % Defines \ip{arg1}{arg2} to mean
                             % (arg1, arg2).

%\newcommand{\ip}[2]{\langle #1 | #2\rangle}
                             % This is an alternative definition of
                             % \ip that is commented out.

\begin{document}             % End of preamble and beginning of text.

\maketitle                   % Produces the title.

This is an example input file.  Comparing it with
the output it generates can show you how to
produce a simple document of your own.

\section{Ordinary Text}      % Produces section heading.  Lower-level
                             % sections are begun with similar
                             % \subsection and \subsubsection commands.

The ends  of words and sentences are marked
  by   spaces. It  doesn't matter how many
spaces    you type; one is as good as 100.  The
end of   a line counts as a space.

One   or more   blank lines denote the  end
of  a paragraph.

Since any number of consecutive spaces are treated
like a single one, the formatting of the input
file makes no difference to
      \LaTeX,                % The \LaTeX command generates the LaTeX logo.
but it makes a difference to you.  When you use
\LaTeX, making your input file as easy to read
as possible will be a great help as you write
your document and when you change it.  This sample
file shows how you can add comments to your own input
file.

Because printing is different from typewriting,
there are a number of things that you have to do
differently when preparing an input file than if
you were just typing the document directly.
Quotation marks like
       ``this''
have to be handled specially, as do quotes within
quotes:
       ``\,`this'            % \, separates the double and single quote.
        is what I just
        wrote, not  `that'\,''.

Dashes come in three sizes: an
       intra-word
dash, a medium dash for number ranges like
       1--2,
and a punctuation
       dash---like
this.

A sentence-ending space should be larger than the
space between words within a sentence.  You
sometimes have to type special commands in
conjunction with punctuation characters to get
this right, as in the following sentence.
       Gnats, gnus, etc.\ all  % `\ ' makes an inter-word space.
       begin with G\@.         % \@ marks end-of-sentence punctuation.
You should check the spaces after periods when
reading your output to make sure you haven't
forgotten any special cases.  Generating an
ellipsis
       \ldots\               % `\ ' is needed after `\ldots' because TeX
                             % ignores spaces after command names like \ldots
                             % made from \ + letters.
                             %
                             % Note how a `%' character causes TeX to ignore
                             % the end of the input line, so these blank lines
                             % do not start a new paragraph.
                             %
with the right spacing around the periods requires
a special command.

\LaTeX\ interprets some common characters as
commands, so you must type special commands to
generate them.  These characters include the
following:
       \$ \& \% \# \{ and \}.

In printing, text is usually emphasized with an
       \emph{italic}
type style.

\begin{em}
   A long segment of text can also be emphasized
   in this way.  Text within such a segment can be
   given \emph{additional} emphasis.
\end{em}

It is sometimes necessary to prevent \LaTeX\ from
breaking a line where it might otherwise do so.
This may be at a space, as between the ``Mr.''\ and
``Jones'' in
       ``Mr.~Jones'',        % ~ produces an unbreakable interword space.
or within a word---especially when the word is a
symbol like
       \mbox{\emph{itemnum}}
that makes little sense when hyphenated across
lines.

Footnotes\footnote{This is an example of a footnote.}
pose no problem.

\LaTeX\ is good at typesetting mathematical formulas
like
       \( x-3y + z = 7 \)
or
       \( a_{1} > x^{2n} + y^{2n} > x' \)
or
       \( \ip{A}{B} = \sum_{i} a_{i} b_{i} \).
The spaces you type in a formula are
ignored.  Remember that a letter like
       $x$                   % $ ... $  and  \( ... \)  are equivalent
is a formula when it denotes a mathematical
symbol, and it should be typed as one.

\section{Displayed Text}

Text is displayed by indenting it from the left
margin.  Quotations are commonly displayed.  There
are short quotations
\begin{quote}
   This is a short quotation.  It consists of a
   single paragraph of text.  See how it is formatted.
\end{quote}
and longer ones.
\begin{quotation}
   This is a longer quotation.  It consists of two
   paragraphs of text, neither of which are
   particularly interesting.

   This is the second paragraph of the quotation.  It
   is just as dull as the first paragraph.
\end{quotation}
Another frequently-displayed structure is a list.
The following is an example of an \emph{itemized}
list.
\begin{itemize}
   \item This is the first item of an itemized list.
         Each item in the list is marked with a ``tick''.
         You don't have to worry about what kind of tick
         mark is used.

   \item This is the second item of the list.  It
         contains another list nested inside it.  The inner
         list is an \emph{enumerated} list.
         \begin{enumerate}
            \item This is the first item of an enumerated
                  list that is nested within the itemized list.

            \item This is the second item of the inner list.
                  \LaTeX\ allows you to nest lists deeper than
                  you really should.
         \end{enumerate}
         This is the rest of the second item of the outer
         list.  It is no more interesting than any other
         part of the item.
   \item This is the third item of the list.
\end{itemize}
You can even display poetry.
\begin{verse}
   There is an environment
    for verse \\             % The \\ command separates lines
   Whose features some poets % within a stanza.
   will curse.

                             % One or more blank lines separate stanzas.

   For instead of making\\
   Them do \emph{all} line breaking, \\
   It allows them to put too many words on a line when they'd rather be
   forced to be terse.
\end{verse}

Mathematical formulas may also be displayed.  A
displayed formula
is
one-line long; multiline
formulas require special formatting instructions.
   \[  \ip{\Gamma}{\psi'} = x'' + y^{2} + z_{i}^{n}\]
Don't start a paragraph with a displayed equation,
nor make one a paragraph by itself.

\end{document}               % End of document.

\section{Bench Huge Section 626}
% synthetic bench section 626
% This is a sample LaTeX input file.  (Version of 12 August 2004.)
%
% A '%' character causes TeX to ignore all remaining text on the line,
% and is used for comments like this one.

\documentclass{article}      % Specifies the document class

                             % The preamble begins here.
\title{An Example Document}  % Declares the document's title.
\author{Leslie Lamport}      % Declares the author's name.
\date{January 21, 1994}      % Deleting this command produces today's date.

\newcommand{\ip}[2]{(#1, #2)}
                             % Defines \ip{arg1}{arg2} to mean
                             % (arg1, arg2).

%\newcommand{\ip}[2]{\langle #1 | #2\rangle}
                             % This is an alternative definition of
                             % \ip that is commented out.

\begin{document}             % End of preamble and beginning of text.

\maketitle                   % Produces the title.

This is an example input file.  Comparing it with
the output it generates can show you how to
produce a simple document of your own.

\section{Ordinary Text}      % Produces section heading.  Lower-level
                             % sections are begun with similar
                             % \subsection and \subsubsection commands.

The ends  of words and sentences are marked
  by   spaces. It  doesn't matter how many
spaces    you type; one is as good as 100.  The
end of   a line counts as a space.

One   or more   blank lines denote the  end
of  a paragraph.

Since any number of consecutive spaces are treated
like a single one, the formatting of the input
file makes no difference to
      \LaTeX,                % The \LaTeX command generates the LaTeX logo.
but it makes a difference to you.  When you use
\LaTeX, making your input file as easy to read
as possible will be a great help as you write
your document and when you change it.  This sample
file shows how you can add comments to your own input
file.

Because printing is different from typewriting,
there are a number of things that you have to do
differently when preparing an input file than if
you were just typing the document directly.
Quotation marks like
       ``this''
have to be handled specially, as do quotes within
quotes:
       ``\,`this'            % \, separates the double and single quote.
        is what I just
        wrote, not  `that'\,''.

Dashes come in three sizes: an
       intra-word
dash, a medium dash for number ranges like
       1--2,
and a punctuation
       dash---like
this.

A sentence-ending space should be larger than the
space between words within a sentence.  You
sometimes have to type special commands in
conjunction with punctuation characters to get
this right, as in the following sentence.
       Gnats, gnus, etc.\ all  % `\ ' makes an inter-word space.
       begin with G\@.         % \@ marks end-of-sentence punctuation.
You should check the spaces after periods when
reading your output to make sure you haven't
forgotten any special cases.  Generating an
ellipsis
       \ldots\               % `\ ' is needed after `\ldots' because TeX
                             % ignores spaces after command names like \ldots
                             % made from \ + letters.
                             %
                             % Note how a `%' character causes TeX to ignore
                             % the end of the input line, so these blank lines
                             % do not start a new paragraph.
                             %
with the right spacing around the periods requires
a special command.

\LaTeX\ interprets some common characters as
commands, so you must type special commands to
generate them.  These characters include the
following:
       \$ \& \% \# \{ and \}.

In printing, text is usually emphasized with an
       \emph{italic}
type style.

\begin{em}
   A long segment of text can also be emphasized
   in this way.  Text within such a segment can be
   given \emph{additional} emphasis.
\end{em}

It is sometimes necessary to prevent \LaTeX\ from
breaking a line where it might otherwise do so.
This may be at a space, as between the ``Mr.''\ and
``Jones'' in
       ``Mr.~Jones'',        % ~ produces an unbreakable interword space.
or within a word---especially when the word is a
symbol like
       \mbox{\emph{itemnum}}
that makes little sense when hyphenated across
lines.

Footnotes\footnote{This is an example of a footnote.}
pose no problem.

\LaTeX\ is good at typesetting mathematical formulas
like
       \( x-3y + z = 7 \)
or
       \( a_{1} > x^{2n} + y^{2n} > x' \)
or
       \( \ip{A}{B} = \sum_{i} a_{i} b_{i} \).
The spaces you type in a formula are
ignored.  Remember that a letter like
       $x$                   % $ ... $  and  \( ... \)  are equivalent
is a formula when it denotes a mathematical
symbol, and it should be typed as one.

\section{Displayed Text}

Text is displayed by indenting it from the left
margin.  Quotations are commonly displayed.  There
are short quotations
\begin{quote}
   This is a short quotation.  It consists of a
   single paragraph of text.  See how it is formatted.
\end{quote}
and longer ones.
\begin{quotation}
   This is a longer quotation.  It consists of two
   paragraphs of text, neither of which are
   particularly interesting.

   This is the second paragraph of the quotation.  It
   is just as dull as the first paragraph.
\end{quotation}
Another frequently-displayed structure is a list.
The following is an example of an \emph{itemized}
list.
\begin{itemize}
   \item This is the first item of an itemized list.
         Each item in the list is marked with a ``tick''.
         You don't have to worry about what kind of tick
         mark is used.

   \item This is the second item of the list.  It
         contains another list nested inside it.  The inner
         list is an \emph{enumerated} list.
         \begin{enumerate}
            \item This is the first item of an enumerated
                  list that is nested within the itemized list.

            \item This is the second item of the inner list.
                  \LaTeX\ allows you to nest lists deeper than
                  you really should.
         \end{enumerate}
         This is the rest of the second item of the outer
         list.  It is no more interesting than any other
         part of the item.
   \item This is the third item of the list.
\end{itemize}
You can even display poetry.
\begin{verse}
   There is an environment
    for verse \\             % The \\ command separates lines
   Whose features some poets % within a stanza.
   will curse.

                             % One or more blank lines separate stanzas.

   For instead of making\\
   Them do \emph{all} line breaking, \\
   It allows them to put too many words on a line when they'd rather be
   forced to be terse.
\end{verse}

Mathematical formulas may also be displayed.  A
displayed formula
is
one-line long; multiline
formulas require special formatting instructions.
   \[  \ip{\Gamma}{\psi'} = x'' + y^{2} + z_{i}^{n}\]
Don't start a paragraph with a displayed equation,
nor make one a paragraph by itself.

\end{document}               % End of document.

\section{Bench Huge Section 627}
% synthetic bench section 627
% This is a sample LaTeX input file.  (Version of 12 August 2004.)
%
% A '%' character causes TeX to ignore all remaining text on the line,
% and is used for comments like this one.

\documentclass{article}      % Specifies the document class

                             % The preamble begins here.
\title{An Example Document}  % Declares the document's title.
\author{Leslie Lamport}      % Declares the author's name.
\date{January 21, 1994}      % Deleting this command produces today's date.

\newcommand{\ip}[2]{(#1, #2)}
                             % Defines \ip{arg1}{arg2} to mean
                             % (arg1, arg2).

%\newcommand{\ip}[2]{\langle #1 | #2\rangle}
                             % This is an alternative definition of
                             % \ip that is commented out.

\begin{document}             % End of preamble and beginning of text.

\maketitle                   % Produces the title.

This is an example input file.  Comparing it with
the output it generates can show you how to
produce a simple document of your own.

\section{Ordinary Text}      % Produces section heading.  Lower-level
                             % sections are begun with similar
                             % \subsection and \subsubsection commands.

The ends  of words and sentences are marked
  by   spaces. It  doesn't matter how many
spaces    you type; one is as good as 100.  The
end of   a line counts as a space.

One   or more   blank lines denote the  end
of  a paragraph.

Since any number of consecutive spaces are treated
like a single one, the formatting of the input
file makes no difference to
      \LaTeX,                % The \LaTeX command generates the LaTeX logo.
but it makes a difference to you.  When you use
\LaTeX, making your input file as easy to read
as possible will be a great help as you write
your document and when you change it.  This sample
file shows how you can add comments to your own input
file.

Because printing is different from typewriting,
there are a number of things that you have to do
differently when preparing an input file than if
you were just typing the document directly.
Quotation marks like
       ``this''
have to be handled specially, as do quotes within
quotes:
       ``\,`this'            % \, separates the double and single quote.
        is what I just
        wrote, not  `that'\,''.

Dashes come in three sizes: an
       intra-word
dash, a medium dash for number ranges like
       1--2,
and a punctuation
       dash---like
this.

A sentence-ending space should be larger than the
space between words within a sentence.  You
sometimes have to type special commands in
conjunction with punctuation characters to get
this right, as in the following sentence.
       Gnats, gnus, etc.\ all  % `\ ' makes an inter-word space.
       begin with G\@.         % \@ marks end-of-sentence punctuation.
You should check the spaces after periods when
reading your output to make sure you haven't
forgotten any special cases.  Generating an
ellipsis
       \ldots\               % `\ ' is needed after `\ldots' because TeX
                             % ignores spaces after command names like \ldots
                             % made from \ + letters.
                             %
                             % Note how a `%' character causes TeX to ignore
                             % the end of the input line, so these blank lines
                             % do not start a new paragraph.
                             %
with the right spacing around the periods requires
a special command.

\LaTeX\ interprets some common characters as
commands, so you must type special commands to
generate them.  These characters include the
following:
       \$ \& \% \# \{ and \}.

In printing, text is usually emphasized with an
       \emph{italic}
type style.

\begin{em}
   A long segment of text can also be emphasized
   in this way.  Text within such a segment can be
   given \emph{additional} emphasis.
\end{em}

It is sometimes necessary to prevent \LaTeX\ from
breaking a line where it might otherwise do so.
This may be at a space, as between the ``Mr.''\ and
``Jones'' in
       ``Mr.~Jones'',        % ~ produces an unbreakable interword space.
or within a word---especially when the word is a
symbol like
       \mbox{\emph{itemnum}}
that makes little sense when hyphenated across
lines.

Footnotes\footnote{This is an example of a footnote.}
pose no problem.

\LaTeX\ is good at typesetting mathematical formulas
like
       \( x-3y + z = 7 \)
or
       \( a_{1} > x^{2n} + y^{2n} > x' \)
or
       \( \ip{A}{B} = \sum_{i} a_{i} b_{i} \).
The spaces you type in a formula are
ignored.  Remember that a letter like
       $x$                   % $ ... $  and  \( ... \)  are equivalent
is a formula when it denotes a mathematical
symbol, and it should be typed as one.

\section{Displayed Text}

Text is displayed by indenting it from the left
margin.  Quotations are commonly displayed.  There
are short quotations
\begin{quote}
   This is a short quotation.  It consists of a
   single paragraph of text.  See how it is formatted.
\end{quote}
and longer ones.
\begin{quotation}
   This is a longer quotation.  It consists of two
   paragraphs of text, neither of which are
   particularly interesting.

   This is the second paragraph of the quotation.  It
   is just as dull as the first paragraph.
\end{quotation}
Another frequently-displayed structure is a list.
The following is an example of an \emph{itemized}
list.
\begin{itemize}
   \item This is the first item of an itemized list.
         Each item in the list is marked with a ``tick''.
         You don't have to worry about what kind of tick
         mark is used.

   \item This is the second item of the list.  It
         contains another list nested inside it.  The inner
         list is an \emph{enumerated} list.
         \begin{enumerate}
            \item This is the first item of an enumerated
                  list that is nested within the itemized list.

            \item This is the second item of the inner list.
                  \LaTeX\ allows you to nest lists deeper than
                  you really should.
         \end{enumerate}
         This is the rest of the second item of the outer
         list.  It is no more interesting than any other
         part of the item.
   \item This is the third item of the list.
\end{itemize}
You can even display poetry.
\begin{verse}
   There is an environment
    for verse \\             % The \\ command separates lines
   Whose features some poets % within a stanza.
   will curse.

                             % One or more blank lines separate stanzas.

   For instead of making\\
   Them do \emph{all} line breaking, \\
   It allows them to put too many words on a line when they'd rather be
   forced to be terse.
\end{verse}

Mathematical formulas may also be displayed.  A
displayed formula
is
one-line long; multiline
formulas require special formatting instructions.
   \[  \ip{\Gamma}{\psi'} = x'' + y^{2} + z_{i}^{n}\]
Don't start a paragraph with a displayed equation,
nor make one a paragraph by itself.

\end{document}               % End of document.

\section{Bench Huge Section 628}
% synthetic bench section 628
% This is a sample LaTeX input file.  (Version of 12 August 2004.)
%
% A '%' character causes TeX to ignore all remaining text on the line,
% and is used for comments like this one.

\documentclass{article}      % Specifies the document class

                             % The preamble begins here.
\title{An Example Document}  % Declares the document's title.
\author{Leslie Lamport}      % Declares the author's name.
\date{January 21, 1994}      % Deleting this command produces today's date.

\newcommand{\ip}[2]{(#1, #2)}
                             % Defines \ip{arg1}{arg2} to mean
                             % (arg1, arg2).

%\newcommand{\ip}[2]{\langle #1 | #2\rangle}
                             % This is an alternative definition of
                             % \ip that is commented out.

\begin{document}             % End of preamble and beginning of text.

\maketitle                   % Produces the title.

This is an example input file.  Comparing it with
the output it generates can show you how to
produce a simple document of your own.

\section{Ordinary Text}      % Produces section heading.  Lower-level
                             % sections are begun with similar
                             % \subsection and \subsubsection commands.

The ends  of words and sentences are marked
  by   spaces. It  doesn't matter how many
spaces    you type; one is as good as 100.  The
end of   a line counts as a space.

One   or more   blank lines denote the  end
of  a paragraph.

Since any number of consecutive spaces are treated
like a single one, the formatting of the input
file makes no difference to
      \LaTeX,                % The \LaTeX command generates the LaTeX logo.
but it makes a difference to you.  When you use
\LaTeX, making your input file as easy to read
as possible will be a great help as you write
your document and when you change it.  This sample
file shows how you can add comments to your own input
file.

Because printing is different from typewriting,
there are a number of things that you have to do
differently when preparing an input file than if
you were just typing the document directly.
Quotation marks like
       ``this''
have to be handled specially, as do quotes within
quotes:
       ``\,`this'            % \, separates the double and single quote.
        is what I just
        wrote, not  `that'\,''.

Dashes come in three sizes: an
       intra-word
dash, a medium dash for number ranges like
       1--2,
and a punctuation
       dash---like
this.

A sentence-ending space should be larger than the
space between words within a sentence.  You
sometimes have to type special commands in
conjunction with punctuation characters to get
this right, as in the following sentence.
       Gnats, gnus, etc.\ all  % `\ ' makes an inter-word space.
       begin with G\@.         % \@ marks end-of-sentence punctuation.
You should check the spaces after periods when
reading your output to make sure you haven't
forgotten any special cases.  Generating an
ellipsis
       \ldots\               % `\ ' is needed after `\ldots' because TeX
                             % ignores spaces after command names like \ldots
                             % made from \ + letters.
                             %
                             % Note how a `%' character causes TeX to ignore
                             % the end of the input line, so these blank lines
                             % do not start a new paragraph.
                             %
with the right spacing around the periods requires
a special command.

\LaTeX\ interprets some common characters as
commands, so you must type special commands to
generate them.  These characters include the
following:
       \$ \& \% \# \{ and \}.

In printing, text is usually emphasized with an
       \emph{italic}
type style.

\begin{em}
   A long segment of text can also be emphasized
   in this way.  Text within such a segment can be
   given \emph{additional} emphasis.
\end{em}

It is sometimes necessary to prevent \LaTeX\ from
breaking a line where it might otherwise do so.
This may be at a space, as between the ``Mr.''\ and
``Jones'' in
       ``Mr.~Jones'',        % ~ produces an unbreakable interword space.
or within a word---especially when the word is a
symbol like
       \mbox{\emph{itemnum}}
that makes little sense when hyphenated across
lines.

Footnotes\footnote{This is an example of a footnote.}
pose no problem.

\LaTeX\ is good at typesetting mathematical formulas
like
       \( x-3y + z = 7 \)
or
       \( a_{1} > x^{2n} + y^{2n} > x' \)
or
       \( \ip{A}{B} = \sum_{i} a_{i} b_{i} \).
The spaces you type in a formula are
ignored.  Remember that a letter like
       $x$                   % $ ... $  and  \( ... \)  are equivalent
is a formula when it denotes a mathematical
symbol, and it should be typed as one.

\section{Displayed Text}

Text is displayed by indenting it from the left
margin.  Quotations are commonly displayed.  There
are short quotations
\begin{quote}
   This is a short quotation.  It consists of a
   single paragraph of text.  See how it is formatted.
\end{quote}
and longer ones.
\begin{quotation}
   This is a longer quotation.  It consists of two
   paragraphs of text, neither of which are
   particularly interesting.

   This is the second paragraph of the quotation.  It
   is just as dull as the first paragraph.
\end{quotation}
Another frequently-displayed structure is a list.
The following is an example of an \emph{itemized}
list.
\begin{itemize}
   \item This is the first item of an itemized list.
         Each item in the list is marked with a ``tick''.
         You don't have to worry about what kind of tick
         mark is used.

   \item This is the second item of the list.  It
         contains another list nested inside it.  The inner
         list is an \emph{enumerated} list.
         \begin{enumerate}
            \item This is the first item of an enumerated
                  list that is nested within the itemized list.

            \item This is the second item of the inner list.
                  \LaTeX\ allows you to nest lists deeper than
                  you really should.
         \end{enumerate}
         This is the rest of the second item of the outer
         list.  It is no more interesting than any other
         part of the item.
   \item This is the third item of the list.
\end{itemize}
You can even display poetry.
\begin{verse}
   There is an environment
    for verse \\             % The \\ command separates lines
   Whose features some poets % within a stanza.
   will curse.

                             % One or more blank lines separate stanzas.

   For instead of making\\
   Them do \emph{all} line breaking, \\
   It allows them to put too many words on a line when they'd rather be
   forced to be terse.
\end{verse}

Mathematical formulas may also be displayed.  A
displayed formula
is
one-line long; multiline
formulas require special formatting instructions.
   \[  \ip{\Gamma}{\psi'} = x'' + y^{2} + z_{i}^{n}\]
Don't start a paragraph with a displayed equation,
nor make one a paragraph by itself.

\end{document}               % End of document.

\section{Bench Huge Section 629}
% synthetic bench section 629
% This is a sample LaTeX input file.  (Version of 12 August 2004.)
%
% A '%' character causes TeX to ignore all remaining text on the line,
% and is used for comments like this one.

\documentclass{article}      % Specifies the document class

                             % The preamble begins here.
\title{An Example Document}  % Declares the document's title.
\author{Leslie Lamport}      % Declares the author's name.
\date{January 21, 1994}      % Deleting this command produces today's date.

\newcommand{\ip}[2]{(#1, #2)}
                             % Defines \ip{arg1}{arg2} to mean
                             % (arg1, arg2).

%\newcommand{\ip}[2]{\langle #1 | #2\rangle}
                             % This is an alternative definition of
                             % \ip that is commented out.

\begin{document}             % End of preamble and beginning of text.

\maketitle                   % Produces the title.

This is an example input file.  Comparing it with
the output it generates can show you how to
produce a simple document of your own.

\section{Ordinary Text}      % Produces section heading.  Lower-level
                             % sections are begun with similar
                             % \subsection and \subsubsection commands.

The ends  of words and sentences are marked
  by   spaces. It  doesn't matter how many
spaces    you type; one is as good as 100.  The
end of   a line counts as a space.

One   or more   blank lines denote the  end
of  a paragraph.

Since any number of consecutive spaces are treated
like a single one, the formatting of the input
file makes no difference to
      \LaTeX,                % The \LaTeX command generates the LaTeX logo.
but it makes a difference to you.  When you use
\LaTeX, making your input file as easy to read
as possible will be a great help as you write
your document and when you change it.  This sample
file shows how you can add comments to your own input
file.

Because printing is different from typewriting,
there are a number of things that you have to do
differently when preparing an input file than if
you were just typing the document directly.
Quotation marks like
       ``this''
have to be handled specially, as do quotes within
quotes:
       ``\,`this'            % \, separates the double and single quote.
        is what I just
        wrote, not  `that'\,''.

Dashes come in three sizes: an
       intra-word
dash, a medium dash for number ranges like
       1--2,
and a punctuation
       dash---like
this.

A sentence-ending space should be larger than the
space between words within a sentence.  You
sometimes have to type special commands in
conjunction with punctuation characters to get
this right, as in the following sentence.
       Gnats, gnus, etc.\ all  % `\ ' makes an inter-word space.
       begin with G\@.         % \@ marks end-of-sentence punctuation.
You should check the spaces after periods when
reading your output to make sure you haven't
forgotten any special cases.  Generating an
ellipsis
       \ldots\               % `\ ' is needed after `\ldots' because TeX
                             % ignores spaces after command names like \ldots
                             % made from \ + letters.
                             %
                             % Note how a `%' character causes TeX to ignore
                             % the end of the input line, so these blank lines
                             % do not start a new paragraph.
                             %
with the right spacing around the periods requires
a special command.

\LaTeX\ interprets some common characters as
commands, so you must type special commands to
generate them.  These characters include the
following:
       \$ \& \% \# \{ and \}.

In printing, text is usually emphasized with an
       \emph{italic}
type style.

\begin{em}
   A long segment of text can also be emphasized
   in this way.  Text within such a segment can be
   given \emph{additional} emphasis.
\end{em}

It is sometimes necessary to prevent \LaTeX\ from
breaking a line where it might otherwise do so.
This may be at a space, as between the ``Mr.''\ and
``Jones'' in
       ``Mr.~Jones'',        % ~ produces an unbreakable interword space.
or within a word---especially when the word is a
symbol like
       \mbox{\emph{itemnum}}
that makes little sense when hyphenated across
lines.

Footnotes\footnote{This is an example of a footnote.}
pose no problem.

\LaTeX\ is good at typesetting mathematical formulas
like
       \( x-3y + z = 7 \)
or
       \( a_{1} > x^{2n} + y^{2n} > x' \)
or
       \( \ip{A}{B} = \sum_{i} a_{i} b_{i} \).
The spaces you type in a formula are
ignored.  Remember that a letter like
       $x$                   % $ ... $  and  \( ... \)  are equivalent
is a formula when it denotes a mathematical
symbol, and it should be typed as one.

\section{Displayed Text}

Text is displayed by indenting it from the left
margin.  Quotations are commonly displayed.  There
are short quotations
\begin{quote}
   This is a short quotation.  It consists of a
   single paragraph of text.  See how it is formatted.
\end{quote}
and longer ones.
\begin{quotation}
   This is a longer quotation.  It consists of two
   paragraphs of text, neither of which are
   particularly interesting.

   This is the second paragraph of the quotation.  It
   is just as dull as the first paragraph.
\end{quotation}
Another frequently-displayed structure is a list.
The following is an example of an \emph{itemized}
list.
\begin{itemize}
   \item This is the first item of an itemized list.
         Each item in the list is marked with a ``tick''.
         You don't have to worry about what kind of tick
         mark is used.

   \item This is the second item of the list.  It
         contains another list nested inside it.  The inner
         list is an \emph{enumerated} list.
         \begin{enumerate}
            \item This is the first item of an enumerated
                  list that is nested within the itemized list.

            \item This is the second item of the inner list.
                  \LaTeX\ allows you to nest lists deeper than
                  you really should.
         \end{enumerate}
         This is the rest of the second item of the outer
         list.  It is no more interesting than any other
         part of the item.
   \item This is the third item of the list.
\end{itemize}
You can even display poetry.
\begin{verse}
   There is an environment
    for verse \\             % The \\ command separates lines
   Whose features some poets % within a stanza.
   will curse.

                             % One or more blank lines separate stanzas.

   For instead of making\\
   Them do \emph{all} line breaking, \\
   It allows them to put too many words on a line when they'd rather be
   forced to be terse.
\end{verse}

Mathematical formulas may also be displayed.  A
displayed formula
is
one-line long; multiline
formulas require special formatting instructions.
   \[  \ip{\Gamma}{\psi'} = x'' + y^{2} + z_{i}^{n}\]
Don't start a paragraph with a displayed equation,
nor make one a paragraph by itself.

\end{document}               % End of document.

\section{Bench Huge Section 630}
% synthetic bench section 630
% This is a sample LaTeX input file.  (Version of 12 August 2004.)
%
% A '%' character causes TeX to ignore all remaining text on the line,
% and is used for comments like this one.

\documentclass{article}      % Specifies the document class

                             % The preamble begins here.
\title{An Example Document}  % Declares the document's title.
\author{Leslie Lamport}      % Declares the author's name.
\date{January 21, 1994}      % Deleting this command produces today's date.

\newcommand{\ip}[2]{(#1, #2)}
                             % Defines \ip{arg1}{arg2} to mean
                             % (arg1, arg2).

%\newcommand{\ip}[2]{\langle #1 | #2\rangle}
                             % This is an alternative definition of
                             % \ip that is commented out.

\begin{document}             % End of preamble and beginning of text.

\maketitle                   % Produces the title.

This is an example input file.  Comparing it with
the output it generates can show you how to
produce a simple document of your own.

\section{Ordinary Text}      % Produces section heading.  Lower-level
                             % sections are begun with similar
                             % \subsection and \subsubsection commands.

The ends  of words and sentences are marked
  by   spaces. It  doesn't matter how many
spaces    you type; one is as good as 100.  The
end of   a line counts as a space.

One   or more   blank lines denote the  end
of  a paragraph.

Since any number of consecutive spaces are treated
like a single one, the formatting of the input
file makes no difference to
      \LaTeX,                % The \LaTeX command generates the LaTeX logo.
but it makes a difference to you.  When you use
\LaTeX, making your input file as easy to read
as possible will be a great help as you write
your document and when you change it.  This sample
file shows how you can add comments to your own input
file.

Because printing is different from typewriting,
there are a number of things that you have to do
differently when preparing an input file than if
you were just typing the document directly.
Quotation marks like
       ``this''
have to be handled specially, as do quotes within
quotes:
       ``\,`this'            % \, separates the double and single quote.
        is what I just
        wrote, not  `that'\,''.

Dashes come in three sizes: an
       intra-word
dash, a medium dash for number ranges like
       1--2,
and a punctuation
       dash---like
this.

A sentence-ending space should be larger than the
space between words within a sentence.  You
sometimes have to type special commands in
conjunction with punctuation characters to get
this right, as in the following sentence.
       Gnats, gnus, etc.\ all  % `\ ' makes an inter-word space.
       begin with G\@.         % \@ marks end-of-sentence punctuation.
You should check the spaces after periods when
reading your output to make sure you haven't
forgotten any special cases.  Generating an
ellipsis
       \ldots\               % `\ ' is needed after `\ldots' because TeX
                             % ignores spaces after command names like \ldots
                             % made from \ + letters.
                             %
                             % Note how a `%' character causes TeX to ignore
                             % the end of the input line, so these blank lines
                             % do not start a new paragraph.
                             %
with the right spacing around the periods requires
a special command.

\LaTeX\ interprets some common characters as
commands, so you must type special commands to
generate them.  These characters include the
following:
       \$ \& \% \# \{ and \}.

In printing, text is usually emphasized with an
       \emph{italic}
type style.

\begin{em}
   A long segment of text can also be emphasized
   in this way.  Text within such a segment can be
   given \emph{additional} emphasis.
\end{em}

It is sometimes necessary to prevent \LaTeX\ from
breaking a line where it might otherwise do so.
This may be at a space, as between the ``Mr.''\ and
``Jones'' in
       ``Mr.~Jones'',        % ~ produces an unbreakable interword space.
or within a word---especially when the word is a
symbol like
       \mbox{\emph{itemnum}}
that makes little sense when hyphenated across
lines.

Footnotes\footnote{This is an example of a footnote.}
pose no problem.

\LaTeX\ is good at typesetting mathematical formulas
like
       \( x-3y + z = 7 \)
or
       \( a_{1} > x^{2n} + y^{2n} > x' \)
or
       \( \ip{A}{B} = \sum_{i} a_{i} b_{i} \).
The spaces you type in a formula are
ignored.  Remember that a letter like
       $x$                   % $ ... $  and  \( ... \)  are equivalent
is a formula when it denotes a mathematical
symbol, and it should be typed as one.

\section{Displayed Text}

Text is displayed by indenting it from the left
margin.  Quotations are commonly displayed.  There
are short quotations
\begin{quote}
   This is a short quotation.  It consists of a
   single paragraph of text.  See how it is formatted.
\end{quote}
and longer ones.
\begin{quotation}
   This is a longer quotation.  It consists of two
   paragraphs of text, neither of which are
   particularly interesting.

   This is the second paragraph of the quotation.  It
   is just as dull as the first paragraph.
\end{quotation}
Another frequently-displayed structure is a list.
The following is an example of an \emph{itemized}
list.
\begin{itemize}
   \item This is the first item of an itemized list.
         Each item in the list is marked with a ``tick''.
         You don't have to worry about what kind of tick
         mark is used.

   \item This is the second item of the list.  It
         contains another list nested inside it.  The inner
         list is an \emph{enumerated} list.
         \begin{enumerate}
            \item This is the first item of an enumerated
                  list that is nested within the itemized list.

            \item This is the second item of the inner list.
                  \LaTeX\ allows you to nest lists deeper than
                  you really should.
         \end{enumerate}
         This is the rest of the second item of the outer
         list.  It is no more interesting than any other
         part of the item.
   \item This is the third item of the list.
\end{itemize}
You can even display poetry.
\begin{verse}
   There is an environment
    for verse \\             % The \\ command separates lines
   Whose features some poets % within a stanza.
   will curse.

                             % One or more blank lines separate stanzas.

   For instead of making\\
   Them do \emph{all} line breaking, \\
   It allows them to put too many words on a line when they'd rather be
   forced to be terse.
\end{verse}

Mathematical formulas may also be displayed.  A
displayed formula
is
one-line long; multiline
formulas require special formatting instructions.
   \[  \ip{\Gamma}{\psi'} = x'' + y^{2} + z_{i}^{n}\]
Don't start a paragraph with a displayed equation,
nor make one a paragraph by itself.

\end{document}               % End of document.

\section{Bench Huge Section 631}
% synthetic bench section 631
% This is a sample LaTeX input file.  (Version of 12 August 2004.)
%
% A '%' character causes TeX to ignore all remaining text on the line,
% and is used for comments like this one.

\documentclass{article}      % Specifies the document class

                             % The preamble begins here.
\title{An Example Document}  % Declares the document's title.
\author{Leslie Lamport}      % Declares the author's name.
\date{January 21, 1994}      % Deleting this command produces today's date.

\newcommand{\ip}[2]{(#1, #2)}
                             % Defines \ip{arg1}{arg2} to mean
                             % (arg1, arg2).

%\newcommand{\ip}[2]{\langle #1 | #2\rangle}
                             % This is an alternative definition of
                             % \ip that is commented out.

\begin{document}             % End of preamble and beginning of text.

\maketitle                   % Produces the title.

This is an example input file.  Comparing it with
the output it generates can show you how to
produce a simple document of your own.

\section{Ordinary Text}      % Produces section heading.  Lower-level
                             % sections are begun with similar
                             % \subsection and \subsubsection commands.

The ends  of words and sentences are marked
  by   spaces. It  doesn't matter how many
spaces    you type; one is as good as 100.  The
end of   a line counts as a space.

One   or more   blank lines denote the  end
of  a paragraph.

Since any number of consecutive spaces are treated
like a single one, the formatting of the input
file makes no difference to
      \LaTeX,                % The \LaTeX command generates the LaTeX logo.
but it makes a difference to you.  When you use
\LaTeX, making your input file as easy to read
as possible will be a great help as you write
your document and when you change it.  This sample
file shows how you can add comments to your own input
file.

Because printing is different from typewriting,
there are a number of things that you have to do
differently when preparing an input file than if
you were just typing the document directly.
Quotation marks like
       ``this''
have to be handled specially, as do quotes within
quotes:
       ``\,`this'            % \, separates the double and single quote.
        is what I just
        wrote, not  `that'\,''.

Dashes come in three sizes: an
       intra-word
dash, a medium dash for number ranges like
       1--2,
and a punctuation
       dash---like
this.

A sentence-ending space should be larger than the
space between words within a sentence.  You
sometimes have to type special commands in
conjunction with punctuation characters to get
this right, as in the following sentence.
       Gnats, gnus, etc.\ all  % `\ ' makes an inter-word space.
       begin with G\@.         % \@ marks end-of-sentence punctuation.
You should check the spaces after periods when
reading your output to make sure you haven't
forgotten any special cases.  Generating an
ellipsis
       \ldots\               % `\ ' is needed after `\ldots' because TeX
                             % ignores spaces after command names like \ldots
                             % made from \ + letters.
                             %
                             % Note how a `%' character causes TeX to ignore
                             % the end of the input line, so these blank lines
                             % do not start a new paragraph.
                             %
with the right spacing around the periods requires
a special command.

\LaTeX\ interprets some common characters as
commands, so you must type special commands to
generate them.  These characters include the
following:
       \$ \& \% \# \{ and \}.

In printing, text is usually emphasized with an
       \emph{italic}
type style.

\begin{em}
   A long segment of text can also be emphasized
   in this way.  Text within such a segment can be
   given \emph{additional} emphasis.
\end{em}

It is sometimes necessary to prevent \LaTeX\ from
breaking a line where it might otherwise do so.
This may be at a space, as between the ``Mr.''\ and
``Jones'' in
       ``Mr.~Jones'',        % ~ produces an unbreakable interword space.
or within a word---especially when the word is a
symbol like
       \mbox{\emph{itemnum}}
that makes little sense when hyphenated across
lines.

Footnotes\footnote{This is an example of a footnote.}
pose no problem.

\LaTeX\ is good at typesetting mathematical formulas
like
       \( x-3y + z = 7 \)
or
       \( a_{1} > x^{2n} + y^{2n} > x' \)
or
       \( \ip{A}{B} = \sum_{i} a_{i} b_{i} \).
The spaces you type in a formula are
ignored.  Remember that a letter like
       $x$                   % $ ... $  and  \( ... \)  are equivalent
is a formula when it denotes a mathematical
symbol, and it should be typed as one.

\section{Displayed Text}

Text is displayed by indenting it from the left
margin.  Quotations are commonly displayed.  There
are short quotations
\begin{quote}
   This is a short quotation.  It consists of a
   single paragraph of text.  See how it is formatted.
\end{quote}
and longer ones.
\begin{quotation}
   This is a longer quotation.  It consists of two
   paragraphs of text, neither of which are
   particularly interesting.

   This is the second paragraph of the quotation.  It
   is just as dull as the first paragraph.
\end{quotation}
Another frequently-displayed structure is a list.
The following is an example of an \emph{itemized}
list.
\begin{itemize}
   \item This is the first item of an itemized list.
         Each item in the list is marked with a ``tick''.
         You don't have to worry about what kind of tick
         mark is used.

   \item This is the second item of the list.  It
         contains another list nested inside it.  The inner
         list is an \emph{enumerated} list.
         \begin{enumerate}
            \item This is the first item of an enumerated
                  list that is nested within the itemized list.

            \item This is the second item of the inner list.
                  \LaTeX\ allows you to nest lists deeper than
                  you really should.
         \end{enumerate}
         This is the rest of the second item of the outer
         list.  It is no more interesting than any other
         part of the item.
   \item This is the third item of the list.
\end{itemize}
You can even display poetry.
\begin{verse}
   There is an environment
    for verse \\             % The \\ command separates lines
   Whose features some poets % within a stanza.
   will curse.

                             % One or more blank lines separate stanzas.

   For instead of making\\
   Them do \emph{all} line breaking, \\
   It allows them to put too many words on a line when they'd rather be
   forced to be terse.
\end{verse}

Mathematical formulas may also be displayed.  A
displayed formula
is
one-line long; multiline
formulas require special formatting instructions.
   \[  \ip{\Gamma}{\psi'} = x'' + y^{2} + z_{i}^{n}\]
Don't start a paragraph with a displayed equation,
nor make one a paragraph by itself.

\end{document}               % End of document.

\section{Bench Huge Section 632}
% synthetic bench section 632
% This is a sample LaTeX input file.  (Version of 12 August 2004.)
%
% A '%' character causes TeX to ignore all remaining text on the line,
% and is used for comments like this one.

\documentclass{article}      % Specifies the document class

                             % The preamble begins here.
\title{An Example Document}  % Declares the document's title.
\author{Leslie Lamport}      % Declares the author's name.
\date{January 21, 1994}      % Deleting this command produces today's date.

\newcommand{\ip}[2]{(#1, #2)}
                             % Defines \ip{arg1}{arg2} to mean
                             % (arg1, arg2).

%\newcommand{\ip}[2]{\langle #1 | #2\rangle}
                             % This is an alternative definition of
                             % \ip that is commented out.

\begin{document}             % End of preamble and beginning of text.

\maketitle                   % Produces the title.

This is an example input file.  Comparing it with
the output it generates can show you how to
produce a simple document of your own.

\section{Ordinary Text}      % Produces section heading.  Lower-level
                             % sections are begun with similar
                             % \subsection and \subsubsection commands.

The ends  of words and sentences are marked
  by   spaces. It  doesn't matter how many
spaces    you type; one is as good as 100.  The
end of   a line counts as a space.

One   or more   blank lines denote the  end
of  a paragraph.

Since any number of consecutive spaces are treated
like a single one, the formatting of the input
file makes no difference to
      \LaTeX,                % The \LaTeX command generates the LaTeX logo.
but it makes a difference to you.  When you use
\LaTeX, making your input file as easy to read
as possible will be a great help as you write
your document and when you change it.  This sample
file shows how you can add comments to your own input
file.

Because printing is different from typewriting,
there are a number of things that you have to do
differently when preparing an input file than if
you were just typing the document directly.
Quotation marks like
       ``this''
have to be handled specially, as do quotes within
quotes:
       ``\,`this'            % \, separates the double and single quote.
        is what I just
        wrote, not  `that'\,''.

Dashes come in three sizes: an
       intra-word
dash, a medium dash for number ranges like
       1--2,
and a punctuation
       dash---like
this.

A sentence-ending space should be larger than the
space between words within a sentence.  You
sometimes have to type special commands in
conjunction with punctuation characters to get
this right, as in the following sentence.
       Gnats, gnus, etc.\ all  % `\ ' makes an inter-word space.
       begin with G\@.         % \@ marks end-of-sentence punctuation.
You should check the spaces after periods when
reading your output to make sure you haven't
forgotten any special cases.  Generating an
ellipsis
       \ldots\               % `\ ' is needed after `\ldots' because TeX
                             % ignores spaces after command names like \ldots
                             % made from \ + letters.
                             %
                             % Note how a `%' character causes TeX to ignore
                             % the end of the input line, so these blank lines
                             % do not start a new paragraph.
                             %
with the right spacing around the periods requires
a special command.

\LaTeX\ interprets some common characters as
commands, so you must type special commands to
generate them.  These characters include the
following:
       \$ \& \% \# \{ and \}.

In printing, text is usually emphasized with an
       \emph{italic}
type style.

\begin{em}
   A long segment of text can also be emphasized
   in this way.  Text within such a segment can be
   given \emph{additional} emphasis.
\end{em}

It is sometimes necessary to prevent \LaTeX\ from
breaking a line where it might otherwise do so.
This may be at a space, as between the ``Mr.''\ and
``Jones'' in
       ``Mr.~Jones'',        % ~ produces an unbreakable interword space.
or within a word---especially when the word is a
symbol like
       \mbox{\emph{itemnum}}
that makes little sense when hyphenated across
lines.

Footnotes\footnote{This is an example of a footnote.}
pose no problem.

\LaTeX\ is good at typesetting mathematical formulas
like
       \( x-3y + z = 7 \)
or
       \( a_{1} > x^{2n} + y^{2n} > x' \)
or
       \( \ip{A}{B} = \sum_{i} a_{i} b_{i} \).
The spaces you type in a formula are
ignored.  Remember that a letter like
       $x$                   % $ ... $  and  \( ... \)  are equivalent
is a formula when it denotes a mathematical
symbol, and it should be typed as one.

\section{Displayed Text}

Text is displayed by indenting it from the left
margin.  Quotations are commonly displayed.  There
are short quotations
\begin{quote}
   This is a short quotation.  It consists of a
   single paragraph of text.  See how it is formatted.
\end{quote}
and longer ones.
\begin{quotation}
   This is a longer quotation.  It consists of two
   paragraphs of text, neither of which are
   particularly interesting.

   This is the second paragraph of the quotation.  It
   is just as dull as the first paragraph.
\end{quotation}
Another frequently-displayed structure is a list.
The following is an example of an \emph{itemized}
list.
\begin{itemize}
   \item This is the first item of an itemized list.
         Each item in the list is marked with a ``tick''.
         You don't have to worry about what kind of tick
         mark is used.

   \item This is the second item of the list.  It
         contains another list nested inside it.  The inner
         list is an \emph{enumerated} list.
         \begin{enumerate}
            \item This is the first item of an enumerated
                  list that is nested within the itemized list.

            \item This is the second item of the inner list.
                  \LaTeX\ allows you to nest lists deeper than
                  you really should.
         \end{enumerate}
         This is the rest of the second item of the outer
         list.  It is no more interesting than any other
         part of the item.
   \item This is the third item of the list.
\end{itemize}
You can even display poetry.
\begin{verse}
   There is an environment
    for verse \\             % The \\ command separates lines
   Whose features some poets % within a stanza.
   will curse.

                             % One or more blank lines separate stanzas.

   For instead of making\\
   Them do \emph{all} line breaking, \\
   It allows them to put too many words on a line when they'd rather be
   forced to be terse.
\end{verse}

Mathematical formulas may also be displayed.  A
displayed formula
is
one-line long; multiline
formulas require special formatting instructions.
   \[  \ip{\Gamma}{\psi'} = x'' + y^{2} + z_{i}^{n}\]
Don't start a paragraph with a displayed equation,
nor make one a paragraph by itself.

\end{document}               % End of document.

\section{Bench Huge Section 633}
% synthetic bench section 633
% This is a sample LaTeX input file.  (Version of 12 August 2004.)
%
% A '%' character causes TeX to ignore all remaining text on the line,
% and is used for comments like this one.

\documentclass{article}      % Specifies the document class

                             % The preamble begins here.
\title{An Example Document}  % Declares the document's title.
\author{Leslie Lamport}      % Declares the author's name.
\date{January 21, 1994}      % Deleting this command produces today's date.

\newcommand{\ip}[2]{(#1, #2)}
                             % Defines \ip{arg1}{arg2} to mean
                             % (arg1, arg2).

%\newcommand{\ip}[2]{\langle #1 | #2\rangle}
                             % This is an alternative definition of
                             % \ip that is commented out.

\begin{document}             % End of preamble and beginning of text.

\maketitle                   % Produces the title.

This is an example input file.  Comparing it with
the output it generates can show you how to
produce a simple document of your own.

\section{Ordinary Text}      % Produces section heading.  Lower-level
                             % sections are begun with similar
                             % \subsection and \subsubsection commands.

The ends  of words and sentences are marked
  by   spaces. It  doesn't matter how many
spaces    you type; one is as good as 100.  The
end of   a line counts as a space.

One   or more   blank lines denote the  end
of  a paragraph.

Since any number of consecutive spaces are treated
like a single one, the formatting of the input
file makes no difference to
      \LaTeX,                % The \LaTeX command generates the LaTeX logo.
but it makes a difference to you.  When you use
\LaTeX, making your input file as easy to read
as possible will be a great help as you write
your document and when you change it.  This sample
file shows how you can add comments to your own input
file.

Because printing is different from typewriting,
there are a number of things that you have to do
differently when preparing an input file than if
you were just typing the document directly.
Quotation marks like
       ``this''
have to be handled specially, as do quotes within
quotes:
       ``\,`this'            % \, separates the double and single quote.
        is what I just
        wrote, not  `that'\,''.

Dashes come in three sizes: an
       intra-word
dash, a medium dash for number ranges like
       1--2,
and a punctuation
       dash---like
this.

A sentence-ending space should be larger than the
space between words within a sentence.  You
sometimes have to type special commands in
conjunction with punctuation characters to get
this right, as in the following sentence.
       Gnats, gnus, etc.\ all  % `\ ' makes an inter-word space.
       begin with G\@.         % \@ marks end-of-sentence punctuation.
You should check the spaces after periods when
reading your output to make sure you haven't
forgotten any special cases.  Generating an
ellipsis
       \ldots\               % `\ ' is needed after `\ldots' because TeX
                             % ignores spaces after command names like \ldots
                             % made from \ + letters.
                             %
                             % Note how a `%' character causes TeX to ignore
                             % the end of the input line, so these blank lines
                             % do not start a new paragraph.
                             %
with the right spacing around the periods requires
a special command.

\LaTeX\ interprets some common characters as
commands, so you must type special commands to
generate them.  These characters include the
following:
       \$ \& \% \# \{ and \}.

In printing, text is usually emphasized with an
       \emph{italic}
type style.

\begin{em}
   A long segment of text can also be emphasized
   in this way.  Text within such a segment can be
   given \emph{additional} emphasis.
\end{em}

It is sometimes necessary to prevent \LaTeX\ from
breaking a line where it might otherwise do so.
This may be at a space, as between the ``Mr.''\ and
``Jones'' in
       ``Mr.~Jones'',        % ~ produces an unbreakable interword space.
or within a word---especially when the word is a
symbol like
       \mbox{\emph{itemnum}}
that makes little sense when hyphenated across
lines.

Footnotes\footnote{This is an example of a footnote.}
pose no problem.

\LaTeX\ is good at typesetting mathematical formulas
like
       \( x-3y + z = 7 \)
or
       \( a_{1} > x^{2n} + y^{2n} > x' \)
or
       \( \ip{A}{B} = \sum_{i} a_{i} b_{i} \).
The spaces you type in a formula are
ignored.  Remember that a letter like
       $x$                   % $ ... $  and  \( ... \)  are equivalent
is a formula when it denotes a mathematical
symbol, and it should be typed as one.

\section{Displayed Text}

Text is displayed by indenting it from the left
margin.  Quotations are commonly displayed.  There
are short quotations
\begin{quote}
   This is a short quotation.  It consists of a
   single paragraph of text.  See how it is formatted.
\end{quote}
and longer ones.
\begin{quotation}
   This is a longer quotation.  It consists of two
   paragraphs of text, neither of which are
   particularly interesting.

   This is the second paragraph of the quotation.  It
   is just as dull as the first paragraph.
\end{quotation}
Another frequently-displayed structure is a list.
The following is an example of an \emph{itemized}
list.
\begin{itemize}
   \item This is the first item of an itemized list.
         Each item in the list is marked with a ``tick''.
         You don't have to worry about what kind of tick
         mark is used.

   \item This is the second item of the list.  It
         contains another list nested inside it.  The inner
         list is an \emph{enumerated} list.
         \begin{enumerate}
            \item This is the first item of an enumerated
                  list that is nested within the itemized list.

            \item This is the second item of the inner list.
                  \LaTeX\ allows you to nest lists deeper than
                  you really should.
         \end{enumerate}
         This is the rest of the second item of the outer
         list.  It is no more interesting than any other
         part of the item.
   \item This is the third item of the list.
\end{itemize}
You can even display poetry.
\begin{verse}
   There is an environment
    for verse \\             % The \\ command separates lines
   Whose features some poets % within a stanza.
   will curse.

                             % One or more blank lines separate stanzas.

   For instead of making\\
   Them do \emph{all} line breaking, \\
   It allows them to put too many words on a line when they'd rather be
   forced to be terse.
\end{verse}

Mathematical formulas may also be displayed.  A
displayed formula
is
one-line long; multiline
formulas require special formatting instructions.
   \[  \ip{\Gamma}{\psi'} = x'' + y^{2} + z_{i}^{n}\]
Don't start a paragraph with a displayed equation,
nor make one a paragraph by itself.

\end{document}               % End of document.

\section{Bench Huge Section 634}
% synthetic bench section 634
% This is a sample LaTeX input file.  (Version of 12 August 2004.)
%
% A '%' character causes TeX to ignore all remaining text on the line,
% and is used for comments like this one.

\documentclass{article}      % Specifies the document class

                             % The preamble begins here.
\title{An Example Document}  % Declares the document's title.
\author{Leslie Lamport}      % Declares the author's name.
\date{January 21, 1994}      % Deleting this command produces today's date.

\newcommand{\ip}[2]{(#1, #2)}
                             % Defines \ip{arg1}{arg2} to mean
                             % (arg1, arg2).

%\newcommand{\ip}[2]{\langle #1 | #2\rangle}
                             % This is an alternative definition of
                             % \ip that is commented out.

\begin{document}             % End of preamble and beginning of text.

\maketitle                   % Produces the title.

This is an example input file.  Comparing it with
the output it generates can show you how to
produce a simple document of your own.

\section{Ordinary Text}      % Produces section heading.  Lower-level
                             % sections are begun with similar
                             % \subsection and \subsubsection commands.

The ends  of words and sentences are marked
  by   spaces. It  doesn't matter how many
spaces    you type; one is as good as 100.  The
end of   a line counts as a space.

One   or more   blank lines denote the  end
of  a paragraph.

Since any number of consecutive spaces are treated
like a single one, the formatting of the input
file makes no difference to
      \LaTeX,                % The \LaTeX command generates the LaTeX logo.
but it makes a difference to you.  When you use
\LaTeX, making your input file as easy to read
as possible will be a great help as you write
your document and when you change it.  This sample
file shows how you can add comments to your own input
file.

Because printing is different from typewriting,
there are a number of things that you have to do
differently when preparing an input file than if
you were just typing the document directly.
Quotation marks like
       ``this''
have to be handled specially, as do quotes within
quotes:
       ``\,`this'            % \, separates the double and single quote.
        is what I just
        wrote, not  `that'\,''.

Dashes come in three sizes: an
       intra-word
dash, a medium dash for number ranges like
       1--2,
and a punctuation
       dash---like
this.

A sentence-ending space should be larger than the
space between words within a sentence.  You
sometimes have to type special commands in
conjunction with punctuation characters to get
this right, as in the following sentence.
       Gnats, gnus, etc.\ all  % `\ ' makes an inter-word space.
       begin with G\@.         % \@ marks end-of-sentence punctuation.
You should check the spaces after periods when
reading your output to make sure you haven't
forgotten any special cases.  Generating an
ellipsis
       \ldots\               % `\ ' is needed after `\ldots' because TeX
                             % ignores spaces after command names like \ldots
                             % made from \ + letters.
                             %
                             % Note how a `%' character causes TeX to ignore
                             % the end of the input line, so these blank lines
                             % do not start a new paragraph.
                             %
with the right spacing around the periods requires
a special command.

\LaTeX\ interprets some common characters as
commands, so you must type special commands to
generate them.  These characters include the
following:
       \$ \& \% \# \{ and \}.

In printing, text is usually emphasized with an
       \emph{italic}
type style.

\begin{em}
   A long segment of text can also be emphasized
   in this way.  Text within such a segment can be
   given \emph{additional} emphasis.
\end{em}

It is sometimes necessary to prevent \LaTeX\ from
breaking a line where it might otherwise do so.
This may be at a space, as between the ``Mr.''\ and
``Jones'' in
       ``Mr.~Jones'',        % ~ produces an unbreakable interword space.
or within a word---especially when the word is a
symbol like
       \mbox{\emph{itemnum}}
that makes little sense when hyphenated across
lines.

Footnotes\footnote{This is an example of a footnote.}
pose no problem.

\LaTeX\ is good at typesetting mathematical formulas
like
       \( x-3y + z = 7 \)
or
       \( a_{1} > x^{2n} + y^{2n} > x' \)
or
       \( \ip{A}{B} = \sum_{i} a_{i} b_{i} \).
The spaces you type in a formula are
ignored.  Remember that a letter like
       $x$                   % $ ... $  and  \( ... \)  are equivalent
is a formula when it denotes a mathematical
symbol, and it should be typed as one.

\section{Displayed Text}

Text is displayed by indenting it from the left
margin.  Quotations are commonly displayed.  There
are short quotations
\begin{quote}
   This is a short quotation.  It consists of a
   single paragraph of text.  See how it is formatted.
\end{quote}
and longer ones.
\begin{quotation}
   This is a longer quotation.  It consists of two
   paragraphs of text, neither of which are
   particularly interesting.

   This is the second paragraph of the quotation.  It
   is just as dull as the first paragraph.
\end{quotation}
Another frequently-displayed structure is a list.
The following is an example of an \emph{itemized}
list.
\begin{itemize}
   \item This is the first item of an itemized list.
         Each item in the list is marked with a ``tick''.
         You don't have to worry about what kind of tick
         mark is used.

   \item This is the second item of the list.  It
         contains another list nested inside it.  The inner
         list is an \emph{enumerated} list.
         \begin{enumerate}
            \item This is the first item of an enumerated
                  list that is nested within the itemized list.

            \item This is the second item of the inner list.
                  \LaTeX\ allows you to nest lists deeper than
                  you really should.
         \end{enumerate}
         This is the rest of the second item of the outer
         list.  It is no more interesting than any other
         part of the item.
   \item This is the third item of the list.
\end{itemize}
You can even display poetry.
\begin{verse}
   There is an environment
    for verse \\             % The \\ command separates lines
   Whose features some poets % within a stanza.
   will curse.

                             % One or more blank lines separate stanzas.

   For instead of making\\
   Them do \emph{all} line breaking, \\
   It allows them to put too many words on a line when they'd rather be
   forced to be terse.
\end{verse}

Mathematical formulas may also be displayed.  A
displayed formula
is
one-line long; multiline
formulas require special formatting instructions.
   \[  \ip{\Gamma}{\psi'} = x'' + y^{2} + z_{i}^{n}\]
Don't start a paragraph with a displayed equation,
nor make one a paragraph by itself.

\end{document}               % End of document.

\section{Bench Huge Section 635}
% synthetic bench section 635
% This is a sample LaTeX input file.  (Version of 12 August 2004.)
%
% A '%' character causes TeX to ignore all remaining text on the line,
% and is used for comments like this one.

\documentclass{article}      % Specifies the document class

                             % The preamble begins here.
\title{An Example Document}  % Declares the document's title.
\author{Leslie Lamport}      % Declares the author's name.
\date{January 21, 1994}      % Deleting this command produces today's date.

\newcommand{\ip}[2]{(#1, #2)}
                             % Defines \ip{arg1}{arg2} to mean
                             % (arg1, arg2).

%\newcommand{\ip}[2]{\langle #1 | #2\rangle}
                             % This is an alternative definition of
                             % \ip that is commented out.

\begin{document}             % End of preamble and beginning of text.

\maketitle                   % Produces the title.

This is an example input file.  Comparing it with
the output it generates can show you how to
produce a simple document of your own.

\section{Ordinary Text}      % Produces section heading.  Lower-level
                             % sections are begun with similar
                             % \subsection and \subsubsection commands.

The ends  of words and sentences are marked
  by   spaces. It  doesn't matter how many
spaces    you type; one is as good as 100.  The
end of   a line counts as a space.

One   or more   blank lines denote the  end
of  a paragraph.

Since any number of consecutive spaces are treated
like a single one, the formatting of the input
file makes no difference to
      \LaTeX,                % The \LaTeX command generates the LaTeX logo.
but it makes a difference to you.  When you use
\LaTeX, making your input file as easy to read
as possible will be a great help as you write
your document and when you change it.  This sample
file shows how you can add comments to your own input
file.

Because printing is different from typewriting,
there are a number of things that you have to do
differently when preparing an input file than if
you were just typing the document directly.
Quotation marks like
       ``this''
have to be handled specially, as do quotes within
quotes:
       ``\,`this'            % \, separates the double and single quote.
        is what I just
        wrote, not  `that'\,''.

Dashes come in three sizes: an
       intra-word
dash, a medium dash for number ranges like
       1--2,
and a punctuation
       dash---like
this.

A sentence-ending space should be larger than the
space between words within a sentence.  You
sometimes have to type special commands in
conjunction with punctuation characters to get
this right, as in the following sentence.
       Gnats, gnus, etc.\ all  % `\ ' makes an inter-word space.
       begin with G\@.         % \@ marks end-of-sentence punctuation.
You should check the spaces after periods when
reading your output to make sure you haven't
forgotten any special cases.  Generating an
ellipsis
       \ldots\               % `\ ' is needed after `\ldots' because TeX
                             % ignores spaces after command names like \ldots
                             % made from \ + letters.
                             %
                             % Note how a `%' character causes TeX to ignore
                             % the end of the input line, so these blank lines
                             % do not start a new paragraph.
                             %
with the right spacing around the periods requires
a special command.

\LaTeX\ interprets some common characters as
commands, so you must type special commands to
generate them.  These characters include the
following:
       \$ \& \% \# \{ and \}.

In printing, text is usually emphasized with an
       \emph{italic}
type style.

\begin{em}
   A long segment of text can also be emphasized
   in this way.  Text within such a segment can be
   given \emph{additional} emphasis.
\end{em}

It is sometimes necessary to prevent \LaTeX\ from
breaking a line where it might otherwise do so.
This may be at a space, as between the ``Mr.''\ and
``Jones'' in
       ``Mr.~Jones'',        % ~ produces an unbreakable interword space.
or within a word---especially when the word is a
symbol like
       \mbox{\emph{itemnum}}
that makes little sense when hyphenated across
lines.

Footnotes\footnote{This is an example of a footnote.}
pose no problem.

\LaTeX\ is good at typesetting mathematical formulas
like
       \( x-3y + z = 7 \)
or
       \( a_{1} > x^{2n} + y^{2n} > x' \)
or
       \( \ip{A}{B} = \sum_{i} a_{i} b_{i} \).
The spaces you type in a formula are
ignored.  Remember that a letter like
       $x$                   % $ ... $  and  \( ... \)  are equivalent
is a formula when it denotes a mathematical
symbol, and it should be typed as one.

\section{Displayed Text}

Text is displayed by indenting it from the left
margin.  Quotations are commonly displayed.  There
are short quotations
\begin{quote}
   This is a short quotation.  It consists of a
   single paragraph of text.  See how it is formatted.
\end{quote}
and longer ones.
\begin{quotation}
   This is a longer quotation.  It consists of two
   paragraphs of text, neither of which are
   particularly interesting.

   This is the second paragraph of the quotation.  It
   is just as dull as the first paragraph.
\end{quotation}
Another frequently-displayed structure is a list.
The following is an example of an \emph{itemized}
list.
\begin{itemize}
   \item This is the first item of an itemized list.
         Each item in the list is marked with a ``tick''.
         You don't have to worry about what kind of tick
         mark is used.

   \item This is the second item of the list.  It
         contains another list nested inside it.  The inner
         list is an \emph{enumerated} list.
         \begin{enumerate}
            \item This is the first item of an enumerated
                  list that is nested within the itemized list.

            \item This is the second item of the inner list.
                  \LaTeX\ allows you to nest lists deeper than
                  you really should.
         \end{enumerate}
         This is the rest of the second item of the outer
         list.  It is no more interesting than any other
         part of the item.
   \item This is the third item of the list.
\end{itemize}
You can even display poetry.
\begin{verse}
   There is an environment
    for verse \\             % The \\ command separates lines
   Whose features some poets % within a stanza.
   will curse.

                             % One or more blank lines separate stanzas.

   For instead of making\\
   Them do \emph{all} line breaking, \\
   It allows them to put too many words on a line when they'd rather be
   forced to be terse.
\end{verse}

Mathematical formulas may also be displayed.  A
displayed formula
is
one-line long; multiline
formulas require special formatting instructions.
   \[  \ip{\Gamma}{\psi'} = x'' + y^{2} + z_{i}^{n}\]
Don't start a paragraph with a displayed equation,
nor make one a paragraph by itself.

\end{document}               % End of document.

\section{Bench Huge Section 636}
% synthetic bench section 636
% This is a sample LaTeX input file.  (Version of 12 August 2004.)
%
% A '%' character causes TeX to ignore all remaining text on the line,
% and is used for comments like this one.

\documentclass{article}      % Specifies the document class

                             % The preamble begins here.
\title{An Example Document}  % Declares the document's title.
\author{Leslie Lamport}      % Declares the author's name.
\date{January 21, 1994}      % Deleting this command produces today's date.

\newcommand{\ip}[2]{(#1, #2)}
                             % Defines \ip{arg1}{arg2} to mean
                             % (arg1, arg2).

%\newcommand{\ip}[2]{\langle #1 | #2\rangle}
                             % This is an alternative definition of
                             % \ip that is commented out.

\begin{document}             % End of preamble and beginning of text.

\maketitle                   % Produces the title.

This is an example input file.  Comparing it with
the output it generates can show you how to
produce a simple document of your own.

\section{Ordinary Text}      % Produces section heading.  Lower-level
                             % sections are begun with similar
                             % \subsection and \subsubsection commands.

The ends  of words and sentences are marked
  by   spaces. It  doesn't matter how many
spaces    you type; one is as good as 100.  The
end of   a line counts as a space.

One   or more   blank lines denote the  end
of  a paragraph.

Since any number of consecutive spaces are treated
like a single one, the formatting of the input
file makes no difference to
      \LaTeX,                % The \LaTeX command generates the LaTeX logo.
but it makes a difference to you.  When you use
\LaTeX, making your input file as easy to read
as possible will be a great help as you write
your document and when you change it.  This sample
file shows how you can add comments to your own input
file.

Because printing is different from typewriting,
there are a number of things that you have to do
differently when preparing an input file than if
you were just typing the document directly.
Quotation marks like
       ``this''
have to be handled specially, as do quotes within
quotes:
       ``\,`this'            % \, separates the double and single quote.
        is what I just
        wrote, not  `that'\,''.

Dashes come in three sizes: an
       intra-word
dash, a medium dash for number ranges like
       1--2,
and a punctuation
       dash---like
this.

A sentence-ending space should be larger than the
space between words within a sentence.  You
sometimes have to type special commands in
conjunction with punctuation characters to get
this right, as in the following sentence.
       Gnats, gnus, etc.\ all  % `\ ' makes an inter-word space.
       begin with G\@.         % \@ marks end-of-sentence punctuation.
You should check the spaces after periods when
reading your output to make sure you haven't
forgotten any special cases.  Generating an
ellipsis
       \ldots\               % `\ ' is needed after `\ldots' because TeX
                             % ignores spaces after command names like \ldots
                             % made from \ + letters.
                             %
                             % Note how a `%' character causes TeX to ignore
                             % the end of the input line, so these blank lines
                             % do not start a new paragraph.
                             %
with the right spacing around the periods requires
a special command.

\LaTeX\ interprets some common characters as
commands, so you must type special commands to
generate them.  These characters include the
following:
       \$ \& \% \# \{ and \}.

In printing, text is usually emphasized with an
       \emph{italic}
type style.

\begin{em}
   A long segment of text can also be emphasized
   in this way.  Text within such a segment can be
   given \emph{additional} emphasis.
\end{em}

It is sometimes necessary to prevent \LaTeX\ from
breaking a line where it might otherwise do so.
This may be at a space, as between the ``Mr.''\ and
``Jones'' in
       ``Mr.~Jones'',        % ~ produces an unbreakable interword space.
or within a word---especially when the word is a
symbol like
       \mbox{\emph{itemnum}}
that makes little sense when hyphenated across
lines.

Footnotes\footnote{This is an example of a footnote.}
pose no problem.

\LaTeX\ is good at typesetting mathematical formulas
like
       \( x-3y + z = 7 \)
or
       \( a_{1} > x^{2n} + y^{2n} > x' \)
or
       \( \ip{A}{B} = \sum_{i} a_{i} b_{i} \).
The spaces you type in a formula are
ignored.  Remember that a letter like
       $x$                   % $ ... $  and  \( ... \)  are equivalent
is a formula when it denotes a mathematical
symbol, and it should be typed as one.

\section{Displayed Text}

Text is displayed by indenting it from the left
margin.  Quotations are commonly displayed.  There
are short quotations
\begin{quote}
   This is a short quotation.  It consists of a
   single paragraph of text.  See how it is formatted.
\end{quote}
and longer ones.
\begin{quotation}
   This is a longer quotation.  It consists of two
   paragraphs of text, neither of which are
   particularly interesting.

   This is the second paragraph of the quotation.  It
   is just as dull as the first paragraph.
\end{quotation}
Another frequently-displayed structure is a list.
The following is an example of an \emph{itemized}
list.
\begin{itemize}
   \item This is the first item of an itemized list.
         Each item in the list is marked with a ``tick''.
         You don't have to worry about what kind of tick
         mark is used.

   \item This is the second item of the list.  It
         contains another list nested inside it.  The inner
         list is an \emph{enumerated} list.
         \begin{enumerate}
            \item This is the first item of an enumerated
                  list that is nested within the itemized list.

            \item This is the second item of the inner list.
                  \LaTeX\ allows you to nest lists deeper than
                  you really should.
         \end{enumerate}
         This is the rest of the second item of the outer
         list.  It is no more interesting than any other
         part of the item.
   \item This is the third item of the list.
\end{itemize}
You can even display poetry.
\begin{verse}
   There is an environment
    for verse \\             % The \\ command separates lines
   Whose features some poets % within a stanza.
   will curse.

                             % One or more blank lines separate stanzas.

   For instead of making\\
   Them do \emph{all} line breaking, \\
   It allows them to put too many words on a line when they'd rather be
   forced to be terse.
\end{verse}

Mathematical formulas may also be displayed.  A
displayed formula
is
one-line long; multiline
formulas require special formatting instructions.
   \[  \ip{\Gamma}{\psi'} = x'' + y^{2} + z_{i}^{n}\]
Don't start a paragraph with a displayed equation,
nor make one a paragraph by itself.

\end{document}               % End of document.

\section{Bench Huge Section 637}
% synthetic bench section 637
% This is a sample LaTeX input file.  (Version of 12 August 2004.)
%
% A '%' character causes TeX to ignore all remaining text on the line,
% and is used for comments like this one.

\documentclass{article}      % Specifies the document class

                             % The preamble begins here.
\title{An Example Document}  % Declares the document's title.
\author{Leslie Lamport}      % Declares the author's name.
\date{January 21, 1994}      % Deleting this command produces today's date.

\newcommand{\ip}[2]{(#1, #2)}
                             % Defines \ip{arg1}{arg2} to mean
                             % (arg1, arg2).

%\newcommand{\ip}[2]{\langle #1 | #2\rangle}
                             % This is an alternative definition of
                             % \ip that is commented out.

\begin{document}             % End of preamble and beginning of text.

\maketitle                   % Produces the title.

This is an example input file.  Comparing it with
the output it generates can show you how to
produce a simple document of your own.

\section{Ordinary Text}      % Produces section heading.  Lower-level
                             % sections are begun with similar
                             % \subsection and \subsubsection commands.

The ends  of words and sentences are marked
  by   spaces. It  doesn't matter how many
spaces    you type; one is as good as 100.  The
end of   a line counts as a space.

One   or more   blank lines denote the  end
of  a paragraph.

Since any number of consecutive spaces are treated
like a single one, the formatting of the input
file makes no difference to
      \LaTeX,                % The \LaTeX command generates the LaTeX logo.
but it makes a difference to you.  When you use
\LaTeX, making your input file as easy to read
as possible will be a great help as you write
your document and when you change it.  This sample
file shows how you can add comments to your own input
file.

Because printing is different from typewriting,
there are a number of things that you have to do
differently when preparing an input file than if
you were just typing the document directly.
Quotation marks like
       ``this''
have to be handled specially, as do quotes within
quotes:
       ``\,`this'            % \, separates the double and single quote.
        is what I just
        wrote, not  `that'\,''.

Dashes come in three sizes: an
       intra-word
dash, a medium dash for number ranges like
       1--2,
and a punctuation
       dash---like
this.

A sentence-ending space should be larger than the
space between words within a sentence.  You
sometimes have to type special commands in
conjunction with punctuation characters to get
this right, as in the following sentence.
       Gnats, gnus, etc.\ all  % `\ ' makes an inter-word space.
       begin with G\@.         % \@ marks end-of-sentence punctuation.
You should check the spaces after periods when
reading your output to make sure you haven't
forgotten any special cases.  Generating an
ellipsis
       \ldots\               % `\ ' is needed after `\ldots' because TeX
                             % ignores spaces after command names like \ldots
                             % made from \ + letters.
                             %
                             % Note how a `%' character causes TeX to ignore
                             % the end of the input line, so these blank lines
                             % do not start a new paragraph.
                             %
with the right spacing around the periods requires
a special command.

\LaTeX\ interprets some common characters as
commands, so you must type special commands to
generate them.  These characters include the
following:
       \$ \& \% \# \{ and \}.

In printing, text is usually emphasized with an
       \emph{italic}
type style.

\begin{em}
   A long segment of text can also be emphasized
   in this way.  Text within such a segment can be
   given \emph{additional} emphasis.
\end{em}

It is sometimes necessary to prevent \LaTeX\ from
breaking a line where it might otherwise do so.
This may be at a space, as between the ``Mr.''\ and
``Jones'' in
       ``Mr.~Jones'',        % ~ produces an unbreakable interword space.
or within a word---especially when the word is a
symbol like
       \mbox{\emph{itemnum}}
that makes little sense when hyphenated across
lines.

Footnotes\footnote{This is an example of a footnote.}
pose no problem.

\LaTeX\ is good at typesetting mathematical formulas
like
       \( x-3y + z = 7 \)
or
       \( a_{1} > x^{2n} + y^{2n} > x' \)
or
       \( \ip{A}{B} = \sum_{i} a_{i} b_{i} \).
The spaces you type in a formula are
ignored.  Remember that a letter like
       $x$                   % $ ... $  and  \( ... \)  are equivalent
is a formula when it denotes a mathematical
symbol, and it should be typed as one.

\section{Displayed Text}

Text is displayed by indenting it from the left
margin.  Quotations are commonly displayed.  There
are short quotations
\begin{quote}
   This is a short quotation.  It consists of a
   single paragraph of text.  See how it is formatted.
\end{quote}
and longer ones.
\begin{quotation}
   This is a longer quotation.  It consists of two
   paragraphs of text, neither of which are
   particularly interesting.

   This is the second paragraph of the quotation.  It
   is just as dull as the first paragraph.
\end{quotation}
Another frequently-displayed structure is a list.
The following is an example of an \emph{itemized}
list.
\begin{itemize}
   \item This is the first item of an itemized list.
         Each item in the list is marked with a ``tick''.
         You don't have to worry about what kind of tick
         mark is used.

   \item This is the second item of the list.  It
         contains another list nested inside it.  The inner
         list is an \emph{enumerated} list.
         \begin{enumerate}
            \item This is the first item of an enumerated
                  list that is nested within the itemized list.

            \item This is the second item of the inner list.
                  \LaTeX\ allows you to nest lists deeper than
                  you really should.
         \end{enumerate}
         This is the rest of the second item of the outer
         list.  It is no more interesting than any other
         part of the item.
   \item This is the third item of the list.
\end{itemize}
You can even display poetry.
\begin{verse}
   There is an environment
    for verse \\             % The \\ command separates lines
   Whose features some poets % within a stanza.
   will curse.

                             % One or more blank lines separate stanzas.

   For instead of making\\
   Them do \emph{all} line breaking, \\
   It allows them to put too many words on a line when they'd rather be
   forced to be terse.
\end{verse}

Mathematical formulas may also be displayed.  A
displayed formula
is
one-line long; multiline
formulas require special formatting instructions.
   \[  \ip{\Gamma}{\psi'} = x'' + y^{2} + z_{i}^{n}\]
Don't start a paragraph with a displayed equation,
nor make one a paragraph by itself.

\end{document}               % End of document.

\section{Bench Huge Section 638}
% synthetic bench section 638
% This is a sample LaTeX input file.  (Version of 12 August 2004.)
%
% A '%' character causes TeX to ignore all remaining text on the line,
% and is used for comments like this one.

\documentclass{article}      % Specifies the document class

                             % The preamble begins here.
\title{An Example Document}  % Declares the document's title.
\author{Leslie Lamport}      % Declares the author's name.
\date{January 21, 1994}      % Deleting this command produces today's date.

\newcommand{\ip}[2]{(#1, #2)}
                             % Defines \ip{arg1}{arg2} to mean
                             % (arg1, arg2).

%\newcommand{\ip}[2]{\langle #1 | #2\rangle}
                             % This is an alternative definition of
                             % \ip that is commented out.

\begin{document}             % End of preamble and beginning of text.

\maketitle                   % Produces the title.

This is an example input file.  Comparing it with
the output it generates can show you how to
produce a simple document of your own.

\section{Ordinary Text}      % Produces section heading.  Lower-level
                             % sections are begun with similar
                             % \subsection and \subsubsection commands.

The ends  of words and sentences are marked
  by   spaces. It  doesn't matter how many
spaces    you type; one is as good as 100.  The
end of   a line counts as a space.

One   or more   blank lines denote the  end
of  a paragraph.

Since any number of consecutive spaces are treated
like a single one, the formatting of the input
file makes no difference to
      \LaTeX,                % The \LaTeX command generates the LaTeX logo.
but it makes a difference to you.  When you use
\LaTeX, making your input file as easy to read
as possible will be a great help as you write
your document and when you change it.  This sample
file shows how you can add comments to your own input
file.

Because printing is different from typewriting,
there are a number of things that you have to do
differently when preparing an input file than if
you were just typing the document directly.
Quotation marks like
       ``this''
have to be handled specially, as do quotes within
quotes:
       ``\,`this'            % \, separates the double and single quote.
        is what I just
        wrote, not  `that'\,''.

Dashes come in three sizes: an
       intra-word
dash, a medium dash for number ranges like
       1--2,
and a punctuation
       dash---like
this.

A sentence-ending space should be larger than the
space between words within a sentence.  You
sometimes have to type special commands in
conjunction with punctuation characters to get
this right, as in the following sentence.
       Gnats, gnus, etc.\ all  % `\ ' makes an inter-word space.
       begin with G\@.         % \@ marks end-of-sentence punctuation.
You should check the spaces after periods when
reading your output to make sure you haven't
forgotten any special cases.  Generating an
ellipsis
       \ldots\               % `\ ' is needed after `\ldots' because TeX
                             % ignores spaces after command names like \ldots
                             % made from \ + letters.
                             %
                             % Note how a `%' character causes TeX to ignore
                             % the end of the input line, so these blank lines
                             % do not start a new paragraph.
                             %
with the right spacing around the periods requires
a special command.

\LaTeX\ interprets some common characters as
commands, so you must type special commands to
generate them.  These characters include the
following:
       \$ \& \% \# \{ and \}.

In printing, text is usually emphasized with an
       \emph{italic}
type style.

\begin{em}
   A long segment of text can also be emphasized
   in this way.  Text within such a segment can be
   given \emph{additional} emphasis.
\end{em}

It is sometimes necessary to prevent \LaTeX\ from
breaking a line where it might otherwise do so.
This may be at a space, as between the ``Mr.''\ and
``Jones'' in
       ``Mr.~Jones'',        % ~ produces an unbreakable interword space.
or within a word---especially when the word is a
symbol like
       \mbox{\emph{itemnum}}
that makes little sense when hyphenated across
lines.

Footnotes\footnote{This is an example of a footnote.}
pose no problem.

\LaTeX\ is good at typesetting mathematical formulas
like
       \( x-3y + z = 7 \)
or
       \( a_{1} > x^{2n} + y^{2n} > x' \)
or
       \( \ip{A}{B} = \sum_{i} a_{i} b_{i} \).
The spaces you type in a formula are
ignored.  Remember that a letter like
       $x$                   % $ ... $  and  \( ... \)  are equivalent
is a formula when it denotes a mathematical
symbol, and it should be typed as one.

\section{Displayed Text}

Text is displayed by indenting it from the left
margin.  Quotations are commonly displayed.  There
are short quotations
\begin{quote}
   This is a short quotation.  It consists of a
   single paragraph of text.  See how it is formatted.
\end{quote}
and longer ones.
\begin{quotation}
   This is a longer quotation.  It consists of two
   paragraphs of text, neither of which are
   particularly interesting.

   This is the second paragraph of the quotation.  It
   is just as dull as the first paragraph.
\end{quotation}
Another frequently-displayed structure is a list.
The following is an example of an \emph{itemized}
list.
\begin{itemize}
   \item This is the first item of an itemized list.
         Each item in the list is marked with a ``tick''.
         You don't have to worry about what kind of tick
         mark is used.

   \item This is the second item of the list.  It
         contains another list nested inside it.  The inner
         list is an \emph{enumerated} list.
         \begin{enumerate}
            \item This is the first item of an enumerated
                  list that is nested within the itemized list.

            \item This is the second item of the inner list.
                  \LaTeX\ allows you to nest lists deeper than
                  you really should.
         \end{enumerate}
         This is the rest of the second item of the outer
         list.  It is no more interesting than any other
         part of the item.
   \item This is the third item of the list.
\end{itemize}
You can even display poetry.
\begin{verse}
   There is an environment
    for verse \\             % The \\ command separates lines
   Whose features some poets % within a stanza.
   will curse.

                             % One or more blank lines separate stanzas.

   For instead of making\\
   Them do \emph{all} line breaking, \\
   It allows them to put too many words on a line when they'd rather be
   forced to be terse.
\end{verse}

Mathematical formulas may also be displayed.  A
displayed formula
is
one-line long; multiline
formulas require special formatting instructions.
   \[  \ip{\Gamma}{\psi'} = x'' + y^{2} + z_{i}^{n}\]
Don't start a paragraph with a displayed equation,
nor make one a paragraph by itself.

\end{document}               % End of document.

\section{Bench Huge Section 639}
% synthetic bench section 639
% This is a sample LaTeX input file.  (Version of 12 August 2004.)
%
% A '%' character causes TeX to ignore all remaining text on the line,
% and is used for comments like this one.

\documentclass{article}      % Specifies the document class

                             % The preamble begins here.
\title{An Example Document}  % Declares the document's title.
\author{Leslie Lamport}      % Declares the author's name.
\date{January 21, 1994}      % Deleting this command produces today's date.

\newcommand{\ip}[2]{(#1, #2)}
                             % Defines \ip{arg1}{arg2} to mean
                             % (arg1, arg2).

%\newcommand{\ip}[2]{\langle #1 | #2\rangle}
                             % This is an alternative definition of
                             % \ip that is commented out.

\begin{document}             % End of preamble and beginning of text.

\maketitle                   % Produces the title.

This is an example input file.  Comparing it with
the output it generates can show you how to
produce a simple document of your own.

\section{Ordinary Text}      % Produces section heading.  Lower-level
                             % sections are begun with similar
                             % \subsection and \subsubsection commands.

The ends  of words and sentences are marked
  by   spaces. It  doesn't matter how many
spaces    you type; one is as good as 100.  The
end of   a line counts as a space.

One   or more   blank lines denote the  end
of  a paragraph.

Since any number of consecutive spaces are treated
like a single one, the formatting of the input
file makes no difference to
      \LaTeX,                % The \LaTeX command generates the LaTeX logo.
but it makes a difference to you.  When you use
\LaTeX, making your input file as easy to read
as possible will be a great help as you write
your document and when you change it.  This sample
file shows how you can add comments to your own input
file.

Because printing is different from typewriting,
there are a number of things that you have to do
differently when preparing an input file than if
you were just typing the document directly.
Quotation marks like
       ``this''
have to be handled specially, as do quotes within
quotes:
       ``\,`this'            % \, separates the double and single quote.
        is what I just
        wrote, not  `that'\,''.

Dashes come in three sizes: an
       intra-word
dash, a medium dash for number ranges like
       1--2,
and a punctuation
       dash---like
this.

A sentence-ending space should be larger than the
space between words within a sentence.  You
sometimes have to type special commands in
conjunction with punctuation characters to get
this right, as in the following sentence.
       Gnats, gnus, etc.\ all  % `\ ' makes an inter-word space.
       begin with G\@.         % \@ marks end-of-sentence punctuation.
You should check the spaces after periods when
reading your output to make sure you haven't
forgotten any special cases.  Generating an
ellipsis
       \ldots\               % `\ ' is needed after `\ldots' because TeX
                             % ignores spaces after command names like \ldots
                             % made from \ + letters.
                             %
                             % Note how a `%' character causes TeX to ignore
                             % the end of the input line, so these blank lines
                             % do not start a new paragraph.
                             %
with the right spacing around the periods requires
a special command.

\LaTeX\ interprets some common characters as
commands, so you must type special commands to
generate them.  These characters include the
following:
       \$ \& \% \# \{ and \}.

In printing, text is usually emphasized with an
       \emph{italic}
type style.

\begin{em}
   A long segment of text can also be emphasized
   in this way.  Text within such a segment can be
   given \emph{additional} emphasis.
\end{em}

It is sometimes necessary to prevent \LaTeX\ from
breaking a line where it might otherwise do so.
This may be at a space, as between the ``Mr.''\ and
``Jones'' in
       ``Mr.~Jones'',        % ~ produces an unbreakable interword space.
or within a word---especially when the word is a
symbol like
       \mbox{\emph{itemnum}}
that makes little sense when hyphenated across
lines.

Footnotes\footnote{This is an example of a footnote.}
pose no problem.

\LaTeX\ is good at typesetting mathematical formulas
like
       \( x-3y + z = 7 \)
or
       \( a_{1} > x^{2n} + y^{2n} > x' \)
or
       \( \ip{A}{B} = \sum_{i} a_{i} b_{i} \).
The spaces you type in a formula are
ignored.  Remember that a letter like
       $x$                   % $ ... $  and  \( ... \)  are equivalent
is a formula when it denotes a mathematical
symbol, and it should be typed as one.

\section{Displayed Text}

Text is displayed by indenting it from the left
margin.  Quotations are commonly displayed.  There
are short quotations
\begin{quote}
   This is a short quotation.  It consists of a
   single paragraph of text.  See how it is formatted.
\end{quote}
and longer ones.
\begin{quotation}
   This is a longer quotation.  It consists of two
   paragraphs of text, neither of which are
   particularly interesting.

   This is the second paragraph of the quotation.  It
   is just as dull as the first paragraph.
\end{quotation}
Another frequently-displayed structure is a list.
The following is an example of an \emph{itemized}
list.
\begin{itemize}
   \item This is the first item of an itemized list.
         Each item in the list is marked with a ``tick''.
         You don't have to worry about what kind of tick
         mark is used.

   \item This is the second item of the list.  It
         contains another list nested inside it.  The inner
         list is an \emph{enumerated} list.
         \begin{enumerate}
            \item This is the first item of an enumerated
                  list that is nested within the itemized list.

            \item This is the second item of the inner list.
                  \LaTeX\ allows you to nest lists deeper than
                  you really should.
         \end{enumerate}
         This is the rest of the second item of the outer
         list.  It is no more interesting than any other
         part of the item.
   \item This is the third item of the list.
\end{itemize}
You can even display poetry.
\begin{verse}
   There is an environment
    for verse \\             % The \\ command separates lines
   Whose features some poets % within a stanza.
   will curse.

                             % One or more blank lines separate stanzas.

   For instead of making\\
   Them do \emph{all} line breaking, \\
   It allows them to put too many words on a line when they'd rather be
   forced to be terse.
\end{verse}

Mathematical formulas may also be displayed.  A
displayed formula
is
one-line long; multiline
formulas require special formatting instructions.
   \[  \ip{\Gamma}{\psi'} = x'' + y^{2} + z_{i}^{n}\]
Don't start a paragraph with a displayed equation,
nor make one a paragraph by itself.

\end{document}               % End of document.

\section{Bench Huge Section 640}
% synthetic bench section 640
% This is a sample LaTeX input file.  (Version of 12 August 2004.)
%
% A '%' character causes TeX to ignore all remaining text on the line,
% and is used for comments like this one.

\documentclass{article}      % Specifies the document class

                             % The preamble begins here.
\title{An Example Document}  % Declares the document's title.
\author{Leslie Lamport}      % Declares the author's name.
\date{January 21, 1994}      % Deleting this command produces today's date.

\newcommand{\ip}[2]{(#1, #2)}
                             % Defines \ip{arg1}{arg2} to mean
                             % (arg1, arg2).

%\newcommand{\ip}[2]{\langle #1 | #2\rangle}
                             % This is an alternative definition of
                             % \ip that is commented out.

\begin{document}             % End of preamble and beginning of text.

\maketitle                   % Produces the title.

This is an example input file.  Comparing it with
the output it generates can show you how to
produce a simple document of your own.

\section{Ordinary Text}      % Produces section heading.  Lower-level
                             % sections are begun with similar
                             % \subsection and \subsubsection commands.

The ends  of words and sentences are marked
  by   spaces. It  doesn't matter how many
spaces    you type; one is as good as 100.  The
end of   a line counts as a space.

One   or more   blank lines denote the  end
of  a paragraph.

Since any number of consecutive spaces are treated
like a single one, the formatting of the input
file makes no difference to
      \LaTeX,                % The \LaTeX command generates the LaTeX logo.
but it makes a difference to you.  When you use
\LaTeX, making your input file as easy to read
as possible will be a great help as you write
your document and when you change it.  This sample
file shows how you can add comments to your own input
file.

Because printing is different from typewriting,
there are a number of things that you have to do
differently when preparing an input file than if
you were just typing the document directly.
Quotation marks like
       ``this''
have to be handled specially, as do quotes within
quotes:
       ``\,`this'            % \, separates the double and single quote.
        is what I just
        wrote, not  `that'\,''.

Dashes come in three sizes: an
       intra-word
dash, a medium dash for number ranges like
       1--2,
and a punctuation
       dash---like
this.

A sentence-ending space should be larger than the
space between words within a sentence.  You
sometimes have to type special commands in
conjunction with punctuation characters to get
this right, as in the following sentence.
       Gnats, gnus, etc.\ all  % `\ ' makes an inter-word space.
       begin with G\@.         % \@ marks end-of-sentence punctuation.
You should check the spaces after periods when
reading your output to make sure you haven't
forgotten any special cases.  Generating an
ellipsis
       \ldots\               % `\ ' is needed after `\ldots' because TeX
                             % ignores spaces after command names like \ldots
                             % made from \ + letters.
                             %
                             % Note how a `%' character causes TeX to ignore
                             % the end of the input line, so these blank lines
                             % do not start a new paragraph.
                             %
with the right spacing around the periods requires
a special command.

\LaTeX\ interprets some common characters as
commands, so you must type special commands to
generate them.  These characters include the
following:
       \$ \& \% \# \{ and \}.

In printing, text is usually emphasized with an
       \emph{italic}
type style.

\begin{em}
   A long segment of text can also be emphasized
   in this way.  Text within such a segment can be
   given \emph{additional} emphasis.
\end{em}

It is sometimes necessary to prevent \LaTeX\ from
breaking a line where it might otherwise do so.
This may be at a space, as between the ``Mr.''\ and
``Jones'' in
       ``Mr.~Jones'',        % ~ produces an unbreakable interword space.
or within a word---especially when the word is a
symbol like
       \mbox{\emph{itemnum}}
that makes little sense when hyphenated across
lines.

Footnotes\footnote{This is an example of a footnote.}
pose no problem.

\LaTeX\ is good at typesetting mathematical formulas
like
       \( x-3y + z = 7 \)
or
       \( a_{1} > x^{2n} + y^{2n} > x' \)
or
       \( \ip{A}{B} = \sum_{i} a_{i} b_{i} \).
The spaces you type in a formula are
ignored.  Remember that a letter like
       $x$                   % $ ... $  and  \( ... \)  are equivalent
is a formula when it denotes a mathematical
symbol, and it should be typed as one.

\section{Displayed Text}

Text is displayed by indenting it from the left
margin.  Quotations are commonly displayed.  There
are short quotations
\begin{quote}
   This is a short quotation.  It consists of a
   single paragraph of text.  See how it is formatted.
\end{quote}
and longer ones.
\begin{quotation}
   This is a longer quotation.  It consists of two
   paragraphs of text, neither of which are
   particularly interesting.

   This is the second paragraph of the quotation.  It
   is just as dull as the first paragraph.
\end{quotation}
Another frequently-displayed structure is a list.
The following is an example of an \emph{itemized}
list.
\begin{itemize}
   \item This is the first item of an itemized list.
         Each item in the list is marked with a ``tick''.
         You don't have to worry about what kind of tick
         mark is used.

   \item This is the second item of the list.  It
         contains another list nested inside it.  The inner
         list is an \emph{enumerated} list.
         \begin{enumerate}
            \item This is the first item of an enumerated
                  list that is nested within the itemized list.

            \item This is the second item of the inner list.
                  \LaTeX\ allows you to nest lists deeper than
                  you really should.
         \end{enumerate}
         This is the rest of the second item of the outer
         list.  It is no more interesting than any other
         part of the item.
   \item This is the third item of the list.
\end{itemize}
You can even display poetry.
\begin{verse}
   There is an environment
    for verse \\             % The \\ command separates lines
   Whose features some poets % within a stanza.
   will curse.

                             % One or more blank lines separate stanzas.

   For instead of making\\
   Them do \emph{all} line breaking, \\
   It allows them to put too many words on a line when they'd rather be
   forced to be terse.
\end{verse}

Mathematical formulas may also be displayed.  A
displayed formula
is
one-line long; multiline
formulas require special formatting instructions.
   \[  \ip{\Gamma}{\psi'} = x'' + y^{2} + z_{i}^{n}\]
Don't start a paragraph with a displayed equation,
nor make one a paragraph by itself.

\end{document}               % End of document.

\section{Bench Huge Section 641}
% synthetic bench section 641
% This is a sample LaTeX input file.  (Version of 12 August 2004.)
%
% A '%' character causes TeX to ignore all remaining text on the line,
% and is used for comments like this one.

\documentclass{article}      % Specifies the document class

                             % The preamble begins here.
\title{An Example Document}  % Declares the document's title.
\author{Leslie Lamport}      % Declares the author's name.
\date{January 21, 1994}      % Deleting this command produces today's date.

\newcommand{\ip}[2]{(#1, #2)}
                             % Defines \ip{arg1}{arg2} to mean
                             % (arg1, arg2).

%\newcommand{\ip}[2]{\langle #1 | #2\rangle}
                             % This is an alternative definition of
                             % \ip that is commented out.

\begin{document}             % End of preamble and beginning of text.

\maketitle                   % Produces the title.

This is an example input file.  Comparing it with
the output it generates can show you how to
produce a simple document of your own.

\section{Ordinary Text}      % Produces section heading.  Lower-level
                             % sections are begun with similar
                             % \subsection and \subsubsection commands.

The ends  of words and sentences are marked
  by   spaces. It  doesn't matter how many
spaces    you type; one is as good as 100.  The
end of   a line counts as a space.

One   or more   blank lines denote the  end
of  a paragraph.

Since any number of consecutive spaces are treated
like a single one, the formatting of the input
file makes no difference to
      \LaTeX,                % The \LaTeX command generates the LaTeX logo.
but it makes a difference to you.  When you use
\LaTeX, making your input file as easy to read
as possible will be a great help as you write
your document and when you change it.  This sample
file shows how you can add comments to your own input
file.

Because printing is different from typewriting,
there are a number of things that you have to do
differently when preparing an input file than if
you were just typing the document directly.
Quotation marks like
       ``this''
have to be handled specially, as do quotes within
quotes:
       ``\,`this'            % \, separates the double and single quote.
        is what I just
        wrote, not  `that'\,''.

Dashes come in three sizes: an
       intra-word
dash, a medium dash for number ranges like
       1--2,
and a punctuation
       dash---like
this.

A sentence-ending space should be larger than the
space between words within a sentence.  You
sometimes have to type special commands in
conjunction with punctuation characters to get
this right, as in the following sentence.
       Gnats, gnus, etc.\ all  % `\ ' makes an inter-word space.
       begin with G\@.         % \@ marks end-of-sentence punctuation.
You should check the spaces after periods when
reading your output to make sure you haven't
forgotten any special cases.  Generating an
ellipsis
       \ldots\               % `\ ' is needed after `\ldots' because TeX
                             % ignores spaces after command names like \ldots
                             % made from \ + letters.
                             %
                             % Note how a `%' character causes TeX to ignore
                             % the end of the input line, so these blank lines
                             % do not start a new paragraph.
                             %
with the right spacing around the periods requires
a special command.

\LaTeX\ interprets some common characters as
commands, so you must type special commands to
generate them.  These characters include the
following:
       \$ \& \% \# \{ and \}.

In printing, text is usually emphasized with an
       \emph{italic}
type style.

\begin{em}
   A long segment of text can also be emphasized
   in this way.  Text within such a segment can be
   given \emph{additional} emphasis.
\end{em}

It is sometimes necessary to prevent \LaTeX\ from
breaking a line where it might otherwise do so.
This may be at a space, as between the ``Mr.''\ and
``Jones'' in
       ``Mr.~Jones'',        % ~ produces an unbreakable interword space.
or within a word---especially when the word is a
symbol like
       \mbox{\emph{itemnum}}
that makes little sense when hyphenated across
lines.

Footnotes\footnote{This is an example of a footnote.}
pose no problem.

\LaTeX\ is good at typesetting mathematical formulas
like
       \( x-3y + z = 7 \)
or
       \( a_{1} > x^{2n} + y^{2n} > x' \)
or
       \( \ip{A}{B} = \sum_{i} a_{i} b_{i} \).
The spaces you type in a formula are
ignored.  Remember that a letter like
       $x$                   % $ ... $  and  \( ... \)  are equivalent
is a formula when it denotes a mathematical
symbol, and it should be typed as one.

\section{Displayed Text}

Text is displayed by indenting it from the left
margin.  Quotations are commonly displayed.  There
are short quotations
\begin{quote}
   This is a short quotation.  It consists of a
   single paragraph of text.  See how it is formatted.
\end{quote}
and longer ones.
\begin{quotation}
   This is a longer quotation.  It consists of two
   paragraphs of text, neither of which are
   particularly interesting.

   This is the second paragraph of the quotation.  It
   is just as dull as the first paragraph.
\end{quotation}
Another frequently-displayed structure is a list.
The following is an example of an \emph{itemized}
list.
\begin{itemize}
   \item This is the first item of an itemized list.
         Each item in the list is marked with a ``tick''.
         You don't have to worry about what kind of tick
         mark is used.

   \item This is the second item of the list.  It
         contains another list nested inside it.  The inner
         list is an \emph{enumerated} list.
         \begin{enumerate}
            \item This is the first item of an enumerated
                  list that is nested within the itemized list.

            \item This is the second item of the inner list.
                  \LaTeX\ allows you to nest lists deeper than
                  you really should.
         \end{enumerate}
         This is the rest of the second item of the outer
         list.  It is no more interesting than any other
         part of the item.
   \item This is the third item of the list.
\end{itemize}
You can even display poetry.
\begin{verse}
   There is an environment
    for verse \\             % The \\ command separates lines
   Whose features some poets % within a stanza.
   will curse.

                             % One or more blank lines separate stanzas.

   For instead of making\\
   Them do \emph{all} line breaking, \\
   It allows them to put too many words on a line when they'd rather be
   forced to be terse.
\end{verse}

Mathematical formulas may also be displayed.  A
displayed formula
is
one-line long; multiline
formulas require special formatting instructions.
   \[  \ip{\Gamma}{\psi'} = x'' + y^{2} + z_{i}^{n}\]
Don't start a paragraph with a displayed equation,
nor make one a paragraph by itself.

\end{document}               % End of document.

\section{Bench Huge Section 642}
% synthetic bench section 642
% This is a sample LaTeX input file.  (Version of 12 August 2004.)
%
% A '%' character causes TeX to ignore all remaining text on the line,
% and is used for comments like this one.

\documentclass{article}      % Specifies the document class

                             % The preamble begins here.
\title{An Example Document}  % Declares the document's title.
\author{Leslie Lamport}      % Declares the author's name.
\date{January 21, 1994}      % Deleting this command produces today's date.

\newcommand{\ip}[2]{(#1, #2)}
                             % Defines \ip{arg1}{arg2} to mean
                             % (arg1, arg2).

%\newcommand{\ip}[2]{\langle #1 | #2\rangle}
                             % This is an alternative definition of
                             % \ip that is commented out.

\begin{document}             % End of preamble and beginning of text.

\maketitle                   % Produces the title.

This is an example input file.  Comparing it with
the output it generates can show you how to
produce a simple document of your own.

\section{Ordinary Text}      % Produces section heading.  Lower-level
                             % sections are begun with similar
                             % \subsection and \subsubsection commands.

The ends  of words and sentences are marked
  by   spaces. It  doesn't matter how many
spaces    you type; one is as good as 100.  The
end of   a line counts as a space.

One   or more   blank lines denote the  end
of  a paragraph.

Since any number of consecutive spaces are treated
like a single one, the formatting of the input
file makes no difference to
      \LaTeX,                % The \LaTeX command generates the LaTeX logo.
but it makes a difference to you.  When you use
\LaTeX, making your input file as easy to read
as possible will be a great help as you write
your document and when you change it.  This sample
file shows how you can add comments to your own input
file.

Because printing is different from typewriting,
there are a number of things that you have to do
differently when preparing an input file than if
you were just typing the document directly.
Quotation marks like
       ``this''
have to be handled specially, as do quotes within
quotes:
       ``\,`this'            % \, separates the double and single quote.
        is what I just
        wrote, not  `that'\,''.

Dashes come in three sizes: an
       intra-word
dash, a medium dash for number ranges like
       1--2,
and a punctuation
       dash---like
this.

A sentence-ending space should be larger than the
space between words within a sentence.  You
sometimes have to type special commands in
conjunction with punctuation characters to get
this right, as in the following sentence.
       Gnats, gnus, etc.\ all  % `\ ' makes an inter-word space.
       begin with G\@.         % \@ marks end-of-sentence punctuation.
You should check the spaces after periods when
reading your output to make sure you haven't
forgotten any special cases.  Generating an
ellipsis
       \ldots\               % `\ ' is needed after `\ldots' because TeX
                             % ignores spaces after command names like \ldots
                             % made from \ + letters.
                             %
                             % Note how a `%' character causes TeX to ignore
                             % the end of the input line, so these blank lines
                             % do not start a new paragraph.
                             %
with the right spacing around the periods requires
a special command.

\LaTeX\ interprets some common characters as
commands, so you must type special commands to
generate them.  These characters include the
following:
       \$ \& \% \# \{ and \}.

In printing, text is usually emphasized with an
       \emph{italic}
type style.

\begin{em}
   A long segment of text can also be emphasized
   in this way.  Text within such a segment can be
   given \emph{additional} emphasis.
\end{em}

It is sometimes necessary to prevent \LaTeX\ from
breaking a line where it might otherwise do so.
This may be at a space, as between the ``Mr.''\ and
``Jones'' in
       ``Mr.~Jones'',        % ~ produces an unbreakable interword space.
or within a word---especially when the word is a
symbol like
       \mbox{\emph{itemnum}}
that makes little sense when hyphenated across
lines.

Footnotes\footnote{This is an example of a footnote.}
pose no problem.

\LaTeX\ is good at typesetting mathematical formulas
like
       \( x-3y + z = 7 \)
or
       \( a_{1} > x^{2n} + y^{2n} > x' \)
or
       \( \ip{A}{B} = \sum_{i} a_{i} b_{i} \).
The spaces you type in a formula are
ignored.  Remember that a letter like
       $x$                   % $ ... $  and  \( ... \)  are equivalent
is a formula when it denotes a mathematical
symbol, and it should be typed as one.

\section{Displayed Text}

Text is displayed by indenting it from the left
margin.  Quotations are commonly displayed.  There
are short quotations
\begin{quote}
   This is a short quotation.  It consists of a
   single paragraph of text.  See how it is formatted.
\end{quote}
and longer ones.
\begin{quotation}
   This is a longer quotation.  It consists of two
   paragraphs of text, neither of which are
   particularly interesting.

   This is the second paragraph of the quotation.  It
   is just as dull as the first paragraph.
\end{quotation}
Another frequently-displayed structure is a list.
The following is an example of an \emph{itemized}
list.
\begin{itemize}
   \item This is the first item of an itemized list.
         Each item in the list is marked with a ``tick''.
         You don't have to worry about what kind of tick
         mark is used.

   \item This is the second item of the list.  It
         contains another list nested inside it.  The inner
         list is an \emph{enumerated} list.
         \begin{enumerate}
            \item This is the first item of an enumerated
                  list that is nested within the itemized list.

            \item This is the second item of the inner list.
                  \LaTeX\ allows you to nest lists deeper than
                  you really should.
         \end{enumerate}
         This is the rest of the second item of the outer
         list.  It is no more interesting than any other
         part of the item.
   \item This is the third item of the list.
\end{itemize}
You can even display poetry.
\begin{verse}
   There is an environment
    for verse \\             % The \\ command separates lines
   Whose features some poets % within a stanza.
   will curse.

                             % One or more blank lines separate stanzas.

   For instead of making\\
   Them do \emph{all} line breaking, \\
   It allows them to put too many words on a line when they'd rather be
   forced to be terse.
\end{verse}

Mathematical formulas may also be displayed.  A
displayed formula
is
one-line long; multiline
formulas require special formatting instructions.
   \[  \ip{\Gamma}{\psi'} = x'' + y^{2} + z_{i}^{n}\]
Don't start a paragraph with a displayed equation,
nor make one a paragraph by itself.

\end{document}               % End of document.

\section{Bench Huge Section 643}
% synthetic bench section 643
% This is a sample LaTeX input file.  (Version of 12 August 2004.)
%
% A '%' character causes TeX to ignore all remaining text on the line,
% and is used for comments like this one.

\documentclass{article}      % Specifies the document class

                             % The preamble begins here.
\title{An Example Document}  % Declares the document's title.
\author{Leslie Lamport}      % Declares the author's name.
\date{January 21, 1994}      % Deleting this command produces today's date.

\newcommand{\ip}[2]{(#1, #2)}
                             % Defines \ip{arg1}{arg2} to mean
                             % (arg1, arg2).

%\newcommand{\ip}[2]{\langle #1 | #2\rangle}
                             % This is an alternative definition of
                             % \ip that is commented out.

\begin{document}             % End of preamble and beginning of text.

\maketitle                   % Produces the title.

This is an example input file.  Comparing it with
the output it generates can show you how to
produce a simple document of your own.

\section{Ordinary Text}      % Produces section heading.  Lower-level
                             % sections are begun with similar
                             % \subsection and \subsubsection commands.

The ends  of words and sentences are marked
  by   spaces. It  doesn't matter how many
spaces    you type; one is as good as 100.  The
end of   a line counts as a space.

One   or more   blank lines denote the  end
of  a paragraph.

Since any number of consecutive spaces are treated
like a single one, the formatting of the input
file makes no difference to
      \LaTeX,                % The \LaTeX command generates the LaTeX logo.
but it makes a difference to you.  When you use
\LaTeX, making your input file as easy to read
as possible will be a great help as you write
your document and when you change it.  This sample
file shows how you can add comments to your own input
file.

Because printing is different from typewriting,
there are a number of things that you have to do
differently when preparing an input file than if
you were just typing the document directly.
Quotation marks like
       ``this''
have to be handled specially, as do quotes within
quotes:
       ``\,`this'            % \, separates the double and single quote.
        is what I just
        wrote, not  `that'\,''.

Dashes come in three sizes: an
       intra-word
dash, a medium dash for number ranges like
       1--2,
and a punctuation
       dash---like
this.

A sentence-ending space should be larger than the
space between words within a sentence.  You
sometimes have to type special commands in
conjunction with punctuation characters to get
this right, as in the following sentence.
       Gnats, gnus, etc.\ all  % `\ ' makes an inter-word space.
       begin with G\@.         % \@ marks end-of-sentence punctuation.
You should check the spaces after periods when
reading your output to make sure you haven't
forgotten any special cases.  Generating an
ellipsis
       \ldots\               % `\ ' is needed after `\ldots' because TeX
                             % ignores spaces after command names like \ldots
                             % made from \ + letters.
                             %
                             % Note how a `%' character causes TeX to ignore
                             % the end of the input line, so these blank lines
                             % do not start a new paragraph.
                             %
with the right spacing around the periods requires
a special command.

\LaTeX\ interprets some common characters as
commands, so you must type special commands to
generate them.  These characters include the
following:
       \$ \& \% \# \{ and \}.

In printing, text is usually emphasized with an
       \emph{italic}
type style.

\begin{em}
   A long segment of text can also be emphasized
   in this way.  Text within such a segment can be
   given \emph{additional} emphasis.
\end{em}

It is sometimes necessary to prevent \LaTeX\ from
breaking a line where it might otherwise do so.
This may be at a space, as between the ``Mr.''\ and
``Jones'' in
       ``Mr.~Jones'',        % ~ produces an unbreakable interword space.
or within a word---especially when the word is a
symbol like
       \mbox{\emph{itemnum}}
that makes little sense when hyphenated across
lines.

Footnotes\footnote{This is an example of a footnote.}
pose no problem.

\LaTeX\ is good at typesetting mathematical formulas
like
       \( x-3y + z = 7 \)
or
       \( a_{1} > x^{2n} + y^{2n} > x' \)
or
       \( \ip{A}{B} = \sum_{i} a_{i} b_{i} \).
The spaces you type in a formula are
ignored.  Remember that a letter like
       $x$                   % $ ... $  and  \( ... \)  are equivalent
is a formula when it denotes a mathematical
symbol, and it should be typed as one.

\section{Displayed Text}

Text is displayed by indenting it from the left
margin.  Quotations are commonly displayed.  There
are short quotations
\begin{quote}
   This is a short quotation.  It consists of a
   single paragraph of text.  See how it is formatted.
\end{quote}
and longer ones.
\begin{quotation}
   This is a longer quotation.  It consists of two
   paragraphs of text, neither of which are
   particularly interesting.

   This is the second paragraph of the quotation.  It
   is just as dull as the first paragraph.
\end{quotation}
Another frequently-displayed structure is a list.
The following is an example of an \emph{itemized}
list.
\begin{itemize}
   \item This is the first item of an itemized list.
         Each item in the list is marked with a ``tick''.
         You don't have to worry about what kind of tick
         mark is used.

   \item This is the second item of the list.  It
         contains another list nested inside it.  The inner
         list is an \emph{enumerated} list.
         \begin{enumerate}
            \item This is the first item of an enumerated
                  list that is nested within the itemized list.

            \item This is the second item of the inner list.
                  \LaTeX\ allows you to nest lists deeper than
                  you really should.
         \end{enumerate}
         This is the rest of the second item of the outer
         list.  It is no more interesting than any other
         part of the item.
   \item This is the third item of the list.
\end{itemize}
You can even display poetry.
\begin{verse}
   There is an environment
    for verse \\             % The \\ command separates lines
   Whose features some poets % within a stanza.
   will curse.

                             % One or more blank lines separate stanzas.

   For instead of making\\
   Them do \emph{all} line breaking, \\
   It allows them to put too many words on a line when they'd rather be
   forced to be terse.
\end{verse}

Mathematical formulas may also be displayed.  A
displayed formula
is
one-line long; multiline
formulas require special formatting instructions.
   \[  \ip{\Gamma}{\psi'} = x'' + y^{2} + z_{i}^{n}\]
Don't start a paragraph with a displayed equation,
nor make one a paragraph by itself.

\end{document}               % End of document.

\section{Bench Huge Section 644}
% synthetic bench section 644
% This is a sample LaTeX input file.  (Version of 12 August 2004.)
%
% A '%' character causes TeX to ignore all remaining text on the line,
% and is used for comments like this one.

\documentclass{article}      % Specifies the document class

                             % The preamble begins here.
\title{An Example Document}  % Declares the document's title.
\author{Leslie Lamport}      % Declares the author's name.
\date{January 21, 1994}      % Deleting this command produces today's date.

\newcommand{\ip}[2]{(#1, #2)}
                             % Defines \ip{arg1}{arg2} to mean
                             % (arg1, arg2).

%\newcommand{\ip}[2]{\langle #1 | #2\rangle}
                             % This is an alternative definition of
                             % \ip that is commented out.

\begin{document}             % End of preamble and beginning of text.

\maketitle                   % Produces the title.

This is an example input file.  Comparing it with
the output it generates can show you how to
produce a simple document of your own.

\section{Ordinary Text}      % Produces section heading.  Lower-level
                             % sections are begun with similar
                             % \subsection and \subsubsection commands.

The ends  of words and sentences are marked
  by   spaces. It  doesn't matter how many
spaces    you type; one is as good as 100.  The
end of   a line counts as a space.

One   or more   blank lines denote the  end
of  a paragraph.

Since any number of consecutive spaces are treated
like a single one, the formatting of the input
file makes no difference to
      \LaTeX,                % The \LaTeX command generates the LaTeX logo.
but it makes a difference to you.  When you use
\LaTeX, making your input file as easy to read
as possible will be a great help as you write
your document and when you change it.  This sample
file shows how you can add comments to your own input
file.

Because printing is different from typewriting,
there are a number of things that you have to do
differently when preparing an input file than if
you were just typing the document directly.
Quotation marks like
       ``this''
have to be handled specially, as do quotes within
quotes:
       ``\,`this'            % \, separates the double and single quote.
        is what I just
        wrote, not  `that'\,''.

Dashes come in three sizes: an
       intra-word
dash, a medium dash for number ranges like
       1--2,
and a punctuation
       dash---like
this.

A sentence-ending space should be larger than the
space between words within a sentence.  You
sometimes have to type special commands in
conjunction with punctuation characters to get
this right, as in the following sentence.
       Gnats, gnus, etc.\ all  % `\ ' makes an inter-word space.
       begin with G\@.         % \@ marks end-of-sentence punctuation.
You should check the spaces after periods when
reading your output to make sure you haven't
forgotten any special cases.  Generating an
ellipsis
       \ldots\               % `\ ' is needed after `\ldots' because TeX
                             % ignores spaces after command names like \ldots
                             % made from \ + letters.
                             %
                             % Note how a `%' character causes TeX to ignore
                             % the end of the input line, so these blank lines
                             % do not start a new paragraph.
                             %
with the right spacing around the periods requires
a special command.

\LaTeX\ interprets some common characters as
commands, so you must type special commands to
generate them.  These characters include the
following:
       \$ \& \% \# \{ and \}.

In printing, text is usually emphasized with an
       \emph{italic}
type style.

\begin{em}
   A long segment of text can also be emphasized
   in this way.  Text within such a segment can be
   given \emph{additional} emphasis.
\end{em}

It is sometimes necessary to prevent \LaTeX\ from
breaking a line where it might otherwise do so.
This may be at a space, as between the ``Mr.''\ and
``Jones'' in
       ``Mr.~Jones'',        % ~ produces an unbreakable interword space.
or within a word---especially when the word is a
symbol like
       \mbox{\emph{itemnum}}
that makes little sense when hyphenated across
lines.

Footnotes\footnote{This is an example of a footnote.}
pose no problem.

\LaTeX\ is good at typesetting mathematical formulas
like
       \( x-3y + z = 7 \)
or
       \( a_{1} > x^{2n} + y^{2n} > x' \)
or
       \( \ip{A}{B} = \sum_{i} a_{i} b_{i} \).
The spaces you type in a formula are
ignored.  Remember that a letter like
       $x$                   % $ ... $  and  \( ... \)  are equivalent
is a formula when it denotes a mathematical
symbol, and it should be typed as one.

\section{Displayed Text}

Text is displayed by indenting it from the left
margin.  Quotations are commonly displayed.  There
are short quotations
\begin{quote}
   This is a short quotation.  It consists of a
   single paragraph of text.  See how it is formatted.
\end{quote}
and longer ones.
\begin{quotation}
   This is a longer quotation.  It consists of two
   paragraphs of text, neither of which are
   particularly interesting.

   This is the second paragraph of the quotation.  It
   is just as dull as the first paragraph.
\end{quotation}
Another frequently-displayed structure is a list.
The following is an example of an \emph{itemized}
list.
\begin{itemize}
   \item This is the first item of an itemized list.
         Each item in the list is marked with a ``tick''.
         You don't have to worry about what kind of tick
         mark is used.

   \item This is the second item of the list.  It
         contains another list nested inside it.  The inner
         list is an \emph{enumerated} list.
         \begin{enumerate}
            \item This is the first item of an enumerated
                  list that is nested within the itemized list.

            \item This is the second item of the inner list.
                  \LaTeX\ allows you to nest lists deeper than
                  you really should.
         \end{enumerate}
         This is the rest of the second item of the outer
         list.  It is no more interesting than any other
         part of the item.
   \item This is the third item of the list.
\end{itemize}
You can even display poetry.
\begin{verse}
   There is an environment
    for verse \\             % The \\ command separates lines
   Whose features some poets % within a stanza.
   will curse.

                             % One or more blank lines separate stanzas.

   For instead of making\\
   Them do \emph{all} line breaking, \\
   It allows them to put too many words on a line when they'd rather be
   forced to be terse.
\end{verse}

Mathematical formulas may also be displayed.  A
displayed formula
is
one-line long; multiline
formulas require special formatting instructions.
   \[  \ip{\Gamma}{\psi'} = x'' + y^{2} + z_{i}^{n}\]
Don't start a paragraph with a displayed equation,
nor make one a paragraph by itself.

\end{document}               % End of document.

\section{Bench Huge Section 645}
% synthetic bench section 645
% This is a sample LaTeX input file.  (Version of 12 August 2004.)
%
% A '%' character causes TeX to ignore all remaining text on the line,
% and is used for comments like this one.

\documentclass{article}      % Specifies the document class

                             % The preamble begins here.
\title{An Example Document}  % Declares the document's title.
\author{Leslie Lamport}      % Declares the author's name.
\date{January 21, 1994}      % Deleting this command produces today's date.

\newcommand{\ip}[2]{(#1, #2)}
                             % Defines \ip{arg1}{arg2} to mean
                             % (arg1, arg2).

%\newcommand{\ip}[2]{\langle #1 | #2\rangle}
                             % This is an alternative definition of
                             % \ip that is commented out.

\begin{document}             % End of preamble and beginning of text.

\maketitle                   % Produces the title.

This is an example input file.  Comparing it with
the output it generates can show you how to
produce a simple document of your own.

\section{Ordinary Text}      % Produces section heading.  Lower-level
                             % sections are begun with similar
                             % \subsection and \subsubsection commands.

The ends  of words and sentences are marked
  by   spaces. It  doesn't matter how many
spaces    you type; one is as good as 100.  The
end of   a line counts as a space.

One   or more   blank lines denote the  end
of  a paragraph.

Since any number of consecutive spaces are treated
like a single one, the formatting of the input
file makes no difference to
      \LaTeX,                % The \LaTeX command generates the LaTeX logo.
but it makes a difference to you.  When you use
\LaTeX, making your input file as easy to read
as possible will be a great help as you write
your document and when you change it.  This sample
file shows how you can add comments to your own input
file.

Because printing is different from typewriting,
there are a number of things that you have to do
differently when preparing an input file than if
you were just typing the document directly.
Quotation marks like
       ``this''
have to be handled specially, as do quotes within
quotes:
       ``\,`this'            % \, separates the double and single quote.
        is what I just
        wrote, not  `that'\,''.

Dashes come in three sizes: an
       intra-word
dash, a medium dash for number ranges like
       1--2,
and a punctuation
       dash---like
this.

A sentence-ending space should be larger than the
space between words within a sentence.  You
sometimes have to type special commands in
conjunction with punctuation characters to get
this right, as in the following sentence.
       Gnats, gnus, etc.\ all  % `\ ' makes an inter-word space.
       begin with G\@.         % \@ marks end-of-sentence punctuation.
You should check the spaces after periods when
reading your output to make sure you haven't
forgotten any special cases.  Generating an
ellipsis
       \ldots\               % `\ ' is needed after `\ldots' because TeX
                             % ignores spaces after command names like \ldots
                             % made from \ + letters.
                             %
                             % Note how a `%' character causes TeX to ignore
                             % the end of the input line, so these blank lines
                             % do not start a new paragraph.
                             %
with the right spacing around the periods requires
a special command.

\LaTeX\ interprets some common characters as
commands, so you must type special commands to
generate them.  These characters include the
following:
       \$ \& \% \# \{ and \}.

In printing, text is usually emphasized with an
       \emph{italic}
type style.

\begin{em}
   A long segment of text can also be emphasized
   in this way.  Text within such a segment can be
   given \emph{additional} emphasis.
\end{em}

It is sometimes necessary to prevent \LaTeX\ from
breaking a line where it might otherwise do so.
This may be at a space, as between the ``Mr.''\ and
``Jones'' in
       ``Mr.~Jones'',        % ~ produces an unbreakable interword space.
or within a word---especially when the word is a
symbol like
       \mbox{\emph{itemnum}}
that makes little sense when hyphenated across
lines.

Footnotes\footnote{This is an example of a footnote.}
pose no problem.

\LaTeX\ is good at typesetting mathematical formulas
like
       \( x-3y + z = 7 \)
or
       \( a_{1} > x^{2n} + y^{2n} > x' \)
or
       \( \ip{A}{B} = \sum_{i} a_{i} b_{i} \).
The spaces you type in a formula are
ignored.  Remember that a letter like
       $x$                   % $ ... $  and  \( ... \)  are equivalent
is a formula when it denotes a mathematical
symbol, and it should be typed as one.

\section{Displayed Text}

Text is displayed by indenting it from the left
margin.  Quotations are commonly displayed.  There
are short quotations
\begin{quote}
   This is a short quotation.  It consists of a
   single paragraph of text.  See how it is formatted.
\end{quote}
and longer ones.
\begin{quotation}
   This is a longer quotation.  It consists of two
   paragraphs of text, neither of which are
   particularly interesting.

   This is the second paragraph of the quotation.  It
   is just as dull as the first paragraph.
\end{quotation}
Another frequently-displayed structure is a list.
The following is an example of an \emph{itemized}
list.
\begin{itemize}
   \item This is the first item of an itemized list.
         Each item in the list is marked with a ``tick''.
         You don't have to worry about what kind of tick
         mark is used.

   \item This is the second item of the list.  It
         contains another list nested inside it.  The inner
         list is an \emph{enumerated} list.
         \begin{enumerate}
            \item This is the first item of an enumerated
                  list that is nested within the itemized list.

            \item This is the second item of the inner list.
                  \LaTeX\ allows you to nest lists deeper than
                  you really should.
         \end{enumerate}
         This is the rest of the second item of the outer
         list.  It is no more interesting than any other
         part of the item.
   \item This is the third item of the list.
\end{itemize}
You can even display poetry.
\begin{verse}
   There is an environment
    for verse \\             % The \\ command separates lines
   Whose features some poets % within a stanza.
   will curse.

                             % One or more blank lines separate stanzas.

   For instead of making\\
   Them do \emph{all} line breaking, \\
   It allows them to put too many words on a line when they'd rather be
   forced to be terse.
\end{verse}

Mathematical formulas may also be displayed.  A
displayed formula
is
one-line long; multiline
formulas require special formatting instructions.
   \[  \ip{\Gamma}{\psi'} = x'' + y^{2} + z_{i}^{n}\]
Don't start a paragraph with a displayed equation,
nor make one a paragraph by itself.

\end{document}               % End of document.

\section{Bench Huge Section 646}
% synthetic bench section 646
% This is a sample LaTeX input file.  (Version of 12 August 2004.)
%
% A '%' character causes TeX to ignore all remaining text on the line,
% and is used for comments like this one.

\documentclass{article}      % Specifies the document class

                             % The preamble begins here.
\title{An Example Document}  % Declares the document's title.
\author{Leslie Lamport}      % Declares the author's name.
\date{January 21, 1994}      % Deleting this command produces today's date.

\newcommand{\ip}[2]{(#1, #2)}
                             % Defines \ip{arg1}{arg2} to mean
                             % (arg1, arg2).

%\newcommand{\ip}[2]{\langle #1 | #2\rangle}
                             % This is an alternative definition of
                             % \ip that is commented out.

\begin{document}             % End of preamble and beginning of text.

\maketitle                   % Produces the title.

This is an example input file.  Comparing it with
the output it generates can show you how to
produce a simple document of your own.

\section{Ordinary Text}      % Produces section heading.  Lower-level
                             % sections are begun with similar
                             % \subsection and \subsubsection commands.

The ends  of words and sentences are marked
  by   spaces. It  doesn't matter how many
spaces    you type; one is as good as 100.  The
end of   a line counts as a space.

One   or more   blank lines denote the  end
of  a paragraph.

Since any number of consecutive spaces are treated
like a single one, the formatting of the input
file makes no difference to
      \LaTeX,                % The \LaTeX command generates the LaTeX logo.
but it makes a difference to you.  When you use
\LaTeX, making your input file as easy to read
as possible will be a great help as you write
your document and when you change it.  This sample
file shows how you can add comments to your own input
file.

Because printing is different from typewriting,
there are a number of things that you have to do
differently when preparing an input file than if
you were just typing the document directly.
Quotation marks like
       ``this''
have to be handled specially, as do quotes within
quotes:
       ``\,`this'            % \, separates the double and single quote.
        is what I just
        wrote, not  `that'\,''.

Dashes come in three sizes: an
       intra-word
dash, a medium dash for number ranges like
       1--2,
and a punctuation
       dash---like
this.

A sentence-ending space should be larger than the
space between words within a sentence.  You
sometimes have to type special commands in
conjunction with punctuation characters to get
this right, as in the following sentence.
       Gnats, gnus, etc.\ all  % `\ ' makes an inter-word space.
       begin with G\@.         % \@ marks end-of-sentence punctuation.
You should check the spaces after periods when
reading your output to make sure you haven't
forgotten any special cases.  Generating an
ellipsis
       \ldots\               % `\ ' is needed after `\ldots' because TeX
                             % ignores spaces after command names like \ldots
                             % made from \ + letters.
                             %
                             % Note how a `%' character causes TeX to ignore
                             % the end of the input line, so these blank lines
                             % do not start a new paragraph.
                             %
with the right spacing around the periods requires
a special command.

\LaTeX\ interprets some common characters as
commands, so you must type special commands to
generate them.  These characters include the
following:
       \$ \& \% \# \{ and \}.

In printing, text is usually emphasized with an
       \emph{italic}
type style.

\begin{em}
   A long segment of text can also be emphasized
   in this way.  Text within such a segment can be
   given \emph{additional} emphasis.
\end{em}

It is sometimes necessary to prevent \LaTeX\ from
breaking a line where it might otherwise do so.
This may be at a space, as between the ``Mr.''\ and
``Jones'' in
       ``Mr.~Jones'',        % ~ produces an unbreakable interword space.
or within a word---especially when the word is a
symbol like
       \mbox{\emph{itemnum}}
that makes little sense when hyphenated across
lines.

Footnotes\footnote{This is an example of a footnote.}
pose no problem.

\LaTeX\ is good at typesetting mathematical formulas
like
       \( x-3y + z = 7 \)
or
       \( a_{1} > x^{2n} + y^{2n} > x' \)
or
       \( \ip{A}{B} = \sum_{i} a_{i} b_{i} \).
The spaces you type in a formula are
ignored.  Remember that a letter like
       $x$                   % $ ... $  and  \( ... \)  are equivalent
is a formula when it denotes a mathematical
symbol, and it should be typed as one.

\section{Displayed Text}

Text is displayed by indenting it from the left
margin.  Quotations are commonly displayed.  There
are short quotations
\begin{quote}
   This is a short quotation.  It consists of a
   single paragraph of text.  See how it is formatted.
\end{quote}
and longer ones.
\begin{quotation}
   This is a longer quotation.  It consists of two
   paragraphs of text, neither of which are
   particularly interesting.

   This is the second paragraph of the quotation.  It
   is just as dull as the first paragraph.
\end{quotation}
Another frequently-displayed structure is a list.
The following is an example of an \emph{itemized}
list.
\begin{itemize}
   \item This is the first item of an itemized list.
         Each item in the list is marked with a ``tick''.
         You don't have to worry about what kind of tick
         mark is used.

   \item This is the second item of the list.  It
         contains another list nested inside it.  The inner
         list is an \emph{enumerated} list.
         \begin{enumerate}
            \item This is the first item of an enumerated
                  list that is nested within the itemized list.

            \item This is the second item of the inner list.
                  \LaTeX\ allows you to nest lists deeper than
                  you really should.
         \end{enumerate}
         This is the rest of the second item of the outer
         list.  It is no more interesting than any other
         part of the item.
   \item This is the third item of the list.
\end{itemize}
You can even display poetry.
\begin{verse}
   There is an environment
    for verse \\             % The \\ command separates lines
   Whose features some poets % within a stanza.
   will curse.

                             % One or more blank lines separate stanzas.

   For instead of making\\
   Them do \emph{all} line breaking, \\
   It allows them to put too many words on a line when they'd rather be
   forced to be terse.
\end{verse}

Mathematical formulas may also be displayed.  A
displayed formula
is
one-line long; multiline
formulas require special formatting instructions.
   \[  \ip{\Gamma}{\psi'} = x'' + y^{2} + z_{i}^{n}\]
Don't start a paragraph with a displayed equation,
nor make one a paragraph by itself.

\end{document}               % End of document.

\section{Bench Huge Section 647}
% synthetic bench section 647
% This is a sample LaTeX input file.  (Version of 12 August 2004.)
%
% A '%' character causes TeX to ignore all remaining text on the line,
% and is used for comments like this one.

\documentclass{article}      % Specifies the document class

                             % The preamble begins here.
\title{An Example Document}  % Declares the document's title.
\author{Leslie Lamport}      % Declares the author's name.
\date{January 21, 1994}      % Deleting this command produces today's date.

\newcommand{\ip}[2]{(#1, #2)}
                             % Defines \ip{arg1}{arg2} to mean
                             % (arg1, arg2).

%\newcommand{\ip}[2]{\langle #1 | #2\rangle}
                             % This is an alternative definition of
                             % \ip that is commented out.

\begin{document}             % End of preamble and beginning of text.

\maketitle                   % Produces the title.

This is an example input file.  Comparing it with
the output it generates can show you how to
produce a simple document of your own.

\section{Ordinary Text}      % Produces section heading.  Lower-level
                             % sections are begun with similar
                             % \subsection and \subsubsection commands.

The ends  of words and sentences are marked
  by   spaces. It  doesn't matter how many
spaces    you type; one is as good as 100.  The
end of   a line counts as a space.

One   or more   blank lines denote the  end
of  a paragraph.

Since any number of consecutive spaces are treated
like a single one, the formatting of the input
file makes no difference to
      \LaTeX,                % The \LaTeX command generates the LaTeX logo.
but it makes a difference to you.  When you use
\LaTeX, making your input file as easy to read
as possible will be a great help as you write
your document and when you change it.  This sample
file shows how you can add comments to your own input
file.

Because printing is different from typewriting,
there are a number of things that you have to do
differently when preparing an input file than if
you were just typing the document directly.
Quotation marks like
       ``this''
have to be handled specially, as do quotes within
quotes:
       ``\,`this'            % \, separates the double and single quote.
        is what I just
        wrote, not  `that'\,''.

Dashes come in three sizes: an
       intra-word
dash, a medium dash for number ranges like
       1--2,
and a punctuation
       dash---like
this.

A sentence-ending space should be larger than the
space between words within a sentence.  You
sometimes have to type special commands in
conjunction with punctuation characters to get
this right, as in the following sentence.
       Gnats, gnus, etc.\ all  % `\ ' makes an inter-word space.
       begin with G\@.         % \@ marks end-of-sentence punctuation.
You should check the spaces after periods when
reading your output to make sure you haven't
forgotten any special cases.  Generating an
ellipsis
       \ldots\               % `\ ' is needed after `\ldots' because TeX
                             % ignores spaces after command names like \ldots
                             % made from \ + letters.
                             %
                             % Note how a `%' character causes TeX to ignore
                             % the end of the input line, so these blank lines
                             % do not start a new paragraph.
                             %
with the right spacing around the periods requires
a special command.

\LaTeX\ interprets some common characters as
commands, so you must type special commands to
generate them.  These characters include the
following:
       \$ \& \% \# \{ and \}.

In printing, text is usually emphasized with an
       \emph{italic}
type style.

\begin{em}
   A long segment of text can also be emphasized
   in this way.  Text within such a segment can be
   given \emph{additional} emphasis.
\end{em}

It is sometimes necessary to prevent \LaTeX\ from
breaking a line where it might otherwise do so.
This may be at a space, as between the ``Mr.''\ and
``Jones'' in
       ``Mr.~Jones'',        % ~ produces an unbreakable interword space.
or within a word---especially when the word is a
symbol like
       \mbox{\emph{itemnum}}
that makes little sense when hyphenated across
lines.

Footnotes\footnote{This is an example of a footnote.}
pose no problem.

\LaTeX\ is good at typesetting mathematical formulas
like
       \( x-3y + z = 7 \)
or
       \( a_{1} > x^{2n} + y^{2n} > x' \)
or
       \( \ip{A}{B} = \sum_{i} a_{i} b_{i} \).
The spaces you type in a formula are
ignored.  Remember that a letter like
       $x$                   % $ ... $  and  \( ... \)  are equivalent
is a formula when it denotes a mathematical
symbol, and it should be typed as one.

\section{Displayed Text}

Text is displayed by indenting it from the left
margin.  Quotations are commonly displayed.  There
are short quotations
\begin{quote}
   This is a short quotation.  It consists of a
   single paragraph of text.  See how it is formatted.
\end{quote}
and longer ones.
\begin{quotation}
   This is a longer quotation.  It consists of two
   paragraphs of text, neither of which are
   particularly interesting.

   This is the second paragraph of the quotation.  It
   is just as dull as the first paragraph.
\end{quotation}
Another frequently-displayed structure is a list.
The following is an example of an \emph{itemized}
list.
\begin{itemize}
   \item This is the first item of an itemized list.
         Each item in the list is marked with a ``tick''.
         You don't have to worry about what kind of tick
         mark is used.

   \item This is the second item of the list.  It
         contains another list nested inside it.  The inner
         list is an \emph{enumerated} list.
         \begin{enumerate}
            \item This is the first item of an enumerated
                  list that is nested within the itemized list.

            \item This is the second item of the inner list.
                  \LaTeX\ allows you to nest lists deeper than
                  you really should.
         \end{enumerate}
         This is the rest of the second item of the outer
         list.  It is no more interesting than any other
         part of the item.
   \item This is the third item of the list.
\end{itemize}
You can even display poetry.
\begin{verse}
   There is an environment
    for verse \\             % The \\ command separates lines
   Whose features some poets % within a stanza.
   will curse.

                             % One or more blank lines separate stanzas.

   For instead of making\\
   Them do \emph{all} line breaking, \\
   It allows them to put too many words on a line when they'd rather be
   forced to be terse.
\end{verse}

Mathematical formulas may also be displayed.  A
displayed formula
is
one-line long; multiline
formulas require special formatting instructions.
   \[  \ip{\Gamma}{\psi'} = x'' + y^{2} + z_{i}^{n}\]
Don't start a paragraph with a displayed equation,
nor make one a paragraph by itself.

\end{document}               % End of document.

\section{Bench Huge Section 648}
% synthetic bench section 648
% This is a sample LaTeX input file.  (Version of 12 August 2004.)
%
% A '%' character causes TeX to ignore all remaining text on the line,
% and is used for comments like this one.

\documentclass{article}      % Specifies the document class

                             % The preamble begins here.
\title{An Example Document}  % Declares the document's title.
\author{Leslie Lamport}      % Declares the author's name.
\date{January 21, 1994}      % Deleting this command produces today's date.

\newcommand{\ip}[2]{(#1, #2)}
                             % Defines \ip{arg1}{arg2} to mean
                             % (arg1, arg2).

%\newcommand{\ip}[2]{\langle #1 | #2\rangle}
                             % This is an alternative definition of
                             % \ip that is commented out.

\begin{document}             % End of preamble and beginning of text.

\maketitle                   % Produces the title.

This is an example input file.  Comparing it with
the output it generates can show you how to
produce a simple document of your own.

\section{Ordinary Text}      % Produces section heading.  Lower-level
                             % sections are begun with similar
                             % \subsection and \subsubsection commands.

The ends  of words and sentences are marked
  by   spaces. It  doesn't matter how many
spaces    you type; one is as good as 100.  The
end of   a line counts as a space.

One   or more   blank lines denote the  end
of  a paragraph.

Since any number of consecutive spaces are treated
like a single one, the formatting of the input
file makes no difference to
      \LaTeX,                % The \LaTeX command generates the LaTeX logo.
but it makes a difference to you.  When you use
\LaTeX, making your input file as easy to read
as possible will be a great help as you write
your document and when you change it.  This sample
file shows how you can add comments to your own input
file.

Because printing is different from typewriting,
there are a number of things that you have to do
differently when preparing an input file than if
you were just typing the document directly.
Quotation marks like
       ``this''
have to be handled specially, as do quotes within
quotes:
       ``\,`this'            % \, separates the double and single quote.
        is what I just
        wrote, not  `that'\,''.

Dashes come in three sizes: an
       intra-word
dash, a medium dash for number ranges like
       1--2,
and a punctuation
       dash---like
this.

A sentence-ending space should be larger than the
space between words within a sentence.  You
sometimes have to type special commands in
conjunction with punctuation characters to get
this right, as in the following sentence.
       Gnats, gnus, etc.\ all  % `\ ' makes an inter-word space.
       begin with G\@.         % \@ marks end-of-sentence punctuation.
You should check the spaces after periods when
reading your output to make sure you haven't
forgotten any special cases.  Generating an
ellipsis
       \ldots\               % `\ ' is needed after `\ldots' because TeX
                             % ignores spaces after command names like \ldots
                             % made from \ + letters.
                             %
                             % Note how a `%' character causes TeX to ignore
                             % the end of the input line, so these blank lines
                             % do not start a new paragraph.
                             %
with the right spacing around the periods requires
a special command.

\LaTeX\ interprets some common characters as
commands, so you must type special commands to
generate them.  These characters include the
following:
       \$ \& \% \# \{ and \}.

In printing, text is usually emphasized with an
       \emph{italic}
type style.

\begin{em}
   A long segment of text can also be emphasized
   in this way.  Text within such a segment can be
   given \emph{additional} emphasis.
\end{em}

It is sometimes necessary to prevent \LaTeX\ from
breaking a line where it might otherwise do so.
This may be at a space, as between the ``Mr.''\ and
``Jones'' in
       ``Mr.~Jones'',        % ~ produces an unbreakable interword space.
or within a word---especially when the word is a
symbol like
       \mbox{\emph{itemnum}}
that makes little sense when hyphenated across
lines.

Footnotes\footnote{This is an example of a footnote.}
pose no problem.

\LaTeX\ is good at typesetting mathematical formulas
like
       \( x-3y + z = 7 \)
or
       \( a_{1} > x^{2n} + y^{2n} > x' \)
or
       \( \ip{A}{B} = \sum_{i} a_{i} b_{i} \).
The spaces you type in a formula are
ignored.  Remember that a letter like
       $x$                   % $ ... $  and  \( ... \)  are equivalent
is a formula when it denotes a mathematical
symbol, and it should be typed as one.

\section{Displayed Text}

Text is displayed by indenting it from the left
margin.  Quotations are commonly displayed.  There
are short quotations
\begin{quote}
   This is a short quotation.  It consists of a
   single paragraph of text.  See how it is formatted.
\end{quote}
and longer ones.
\begin{quotation}
   This is a longer quotation.  It consists of two
   paragraphs of text, neither of which are
   particularly interesting.

   This is the second paragraph of the quotation.  It
   is just as dull as the first paragraph.
\end{quotation}
Another frequently-displayed structure is a list.
The following is an example of an \emph{itemized}
list.
\begin{itemize}
   \item This is the first item of an itemized list.
         Each item in the list is marked with a ``tick''.
         You don't have to worry about what kind of tick
         mark is used.

   \item This is the second item of the list.  It
         contains another list nested inside it.  The inner
         list is an \emph{enumerated} list.
         \begin{enumerate}
            \item This is the first item of an enumerated
                  list that is nested within the itemized list.

            \item This is the second item of the inner list.
                  \LaTeX\ allows you to nest lists deeper than
                  you really should.
         \end{enumerate}
         This is the rest of the second item of the outer
         list.  It is no more interesting than any other
         part of the item.
   \item This is the third item of the list.
\end{itemize}
You can even display poetry.
\begin{verse}
   There is an environment
    for verse \\             % The \\ command separates lines
   Whose features some poets % within a stanza.
   will curse.

                             % One or more blank lines separate stanzas.

   For instead of making\\
   Them do \emph{all} line breaking, \\
   It allows them to put too many words on a line when they'd rather be
   forced to be terse.
\end{verse}

Mathematical formulas may also be displayed.  A
displayed formula
is
one-line long; multiline
formulas require special formatting instructions.
   \[  \ip{\Gamma}{\psi'} = x'' + y^{2} + z_{i}^{n}\]
Don't start a paragraph with a displayed equation,
nor make one a paragraph by itself.

\end{document}               % End of document.

\section{Bench Huge Section 649}
% synthetic bench section 649
% This is a sample LaTeX input file.  (Version of 12 August 2004.)
%
% A '%' character causes TeX to ignore all remaining text on the line,
% and is used for comments like this one.

\documentclass{article}      % Specifies the document class

                             % The preamble begins here.
\title{An Example Document}  % Declares the document's title.
\author{Leslie Lamport}      % Declares the author's name.
\date{January 21, 1994}      % Deleting this command produces today's date.

\newcommand{\ip}[2]{(#1, #2)}
                             % Defines \ip{arg1}{arg2} to mean
                             % (arg1, arg2).

%\newcommand{\ip}[2]{\langle #1 | #2\rangle}
                             % This is an alternative definition of
                             % \ip that is commented out.

\begin{document}             % End of preamble and beginning of text.

\maketitle                   % Produces the title.

This is an example input file.  Comparing it with
the output it generates can show you how to
produce a simple document of your own.

\section{Ordinary Text}      % Produces section heading.  Lower-level
                             % sections are begun with similar
                             % \subsection and \subsubsection commands.

The ends  of words and sentences are marked
  by   spaces. It  doesn't matter how many
spaces    you type; one is as good as 100.  The
end of   a line counts as a space.

One   or more   blank lines denote the  end
of  a paragraph.

Since any number of consecutive spaces are treated
like a single one, the formatting of the input
file makes no difference to
      \LaTeX,                % The \LaTeX command generates the LaTeX logo.
but it makes a difference to you.  When you use
\LaTeX, making your input file as easy to read
as possible will be a great help as you write
your document and when you change it.  This sample
file shows how you can add comments to your own input
file.

Because printing is different from typewriting,
there are a number of things that you have to do
differently when preparing an input file than if
you were just typing the document directly.
Quotation marks like
       ``this''
have to be handled specially, as do quotes within
quotes:
       ``\,`this'            % \, separates the double and single quote.
        is what I just
        wrote, not  `that'\,''.

Dashes come in three sizes: an
       intra-word
dash, a medium dash for number ranges like
       1--2,
and a punctuation
       dash---like
this.

A sentence-ending space should be larger than the
space between words within a sentence.  You
sometimes have to type special commands in
conjunction with punctuation characters to get
this right, as in the following sentence.
       Gnats, gnus, etc.\ all  % `\ ' makes an inter-word space.
       begin with G\@.         % \@ marks end-of-sentence punctuation.
You should check the spaces after periods when
reading your output to make sure you haven't
forgotten any special cases.  Generating an
ellipsis
       \ldots\               % `\ ' is needed after `\ldots' because TeX
                             % ignores spaces after command names like \ldots
                             % made from \ + letters.
                             %
                             % Note how a `%' character causes TeX to ignore
                             % the end of the input line, so these blank lines
                             % do not start a new paragraph.
                             %
with the right spacing around the periods requires
a special command.

\LaTeX\ interprets some common characters as
commands, so you must type special commands to
generate them.  These characters include the
following:
       \$ \& \% \# \{ and \}.

In printing, text is usually emphasized with an
       \emph{italic}
type style.

\begin{em}
   A long segment of text can also be emphasized
   in this way.  Text within such a segment can be
   given \emph{additional} emphasis.
\end{em}

It is sometimes necessary to prevent \LaTeX\ from
breaking a line where it might otherwise do so.
This may be at a space, as between the ``Mr.''\ and
``Jones'' in
       ``Mr.~Jones'',        % ~ produces an unbreakable interword space.
or within a word---especially when the word is a
symbol like
       \mbox{\emph{itemnum}}
that makes little sense when hyphenated across
lines.

Footnotes\footnote{This is an example of a footnote.}
pose no problem.

\LaTeX\ is good at typesetting mathematical formulas
like
       \( x-3y + z = 7 \)
or
       \( a_{1} > x^{2n} + y^{2n} > x' \)
or
       \( \ip{A}{B} = \sum_{i} a_{i} b_{i} \).
The spaces you type in a formula are
ignored.  Remember that a letter like
       $x$                   % $ ... $  and  \( ... \)  are equivalent
is a formula when it denotes a mathematical
symbol, and it should be typed as one.

\section{Displayed Text}

Text is displayed by indenting it from the left
margin.  Quotations are commonly displayed.  There
are short quotations
\begin{quote}
   This is a short quotation.  It consists of a
   single paragraph of text.  See how it is formatted.
\end{quote}
and longer ones.
\begin{quotation}
   This is a longer quotation.  It consists of two
   paragraphs of text, neither of which are
   particularly interesting.

   This is the second paragraph of the quotation.  It
   is just as dull as the first paragraph.
\end{quotation}
Another frequently-displayed structure is a list.
The following is an example of an \emph{itemized}
list.
\begin{itemize}
   \item This is the first item of an itemized list.
         Each item in the list is marked with a ``tick''.
         You don't have to worry about what kind of tick
         mark is used.

   \item This is the second item of the list.  It
         contains another list nested inside it.  The inner
         list is an \emph{enumerated} list.
         \begin{enumerate}
            \item This is the first item of an enumerated
                  list that is nested within the itemized list.

            \item This is the second item of the inner list.
                  \LaTeX\ allows you to nest lists deeper than
                  you really should.
         \end{enumerate}
         This is the rest of the second item of the outer
         list.  It is no more interesting than any other
         part of the item.
   \item This is the third item of the list.
\end{itemize}
You can even display poetry.
\begin{verse}
   There is an environment
    for verse \\             % The \\ command separates lines
   Whose features some poets % within a stanza.
   will curse.

                             % One or more blank lines separate stanzas.

   For instead of making\\
   Them do \emph{all} line breaking, \\
   It allows them to put too many words on a line when they'd rather be
   forced to be terse.
\end{verse}

Mathematical formulas may also be displayed.  A
displayed formula
is
one-line long; multiline
formulas require special formatting instructions.
   \[  \ip{\Gamma}{\psi'} = x'' + y^{2} + z_{i}^{n}\]
Don't start a paragraph with a displayed equation,
nor make one a paragraph by itself.

\end{document}               % End of document.

\section{Bench Huge Section 650}
% synthetic bench section 650
% This is a sample LaTeX input file.  (Version of 12 August 2004.)
%
% A '%' character causes TeX to ignore all remaining text on the line,
% and is used for comments like this one.

\documentclass{article}      % Specifies the document class

                             % The preamble begins here.
\title{An Example Document}  % Declares the document's title.
\author{Leslie Lamport}      % Declares the author's name.
\date{January 21, 1994}      % Deleting this command produces today's date.

\newcommand{\ip}[2]{(#1, #2)}
                             % Defines \ip{arg1}{arg2} to mean
                             % (arg1, arg2).

%\newcommand{\ip}[2]{\langle #1 | #2\rangle}
                             % This is an alternative definition of
                             % \ip that is commented out.

\begin{document}             % End of preamble and beginning of text.

\maketitle                   % Produces the title.

This is an example input file.  Comparing it with
the output it generates can show you how to
produce a simple document of your own.

\section{Ordinary Text}      % Produces section heading.  Lower-level
                             % sections are begun with similar
                             % \subsection and \subsubsection commands.

The ends  of words and sentences are marked
  by   spaces. It  doesn't matter how many
spaces    you type; one is as good as 100.  The
end of   a line counts as a space.

One   or more   blank lines denote the  end
of  a paragraph.

Since any number of consecutive spaces are treated
like a single one, the formatting of the input
file makes no difference to
      \LaTeX,                % The \LaTeX command generates the LaTeX logo.
but it makes a difference to you.  When you use
\LaTeX, making your input file as easy to read
as possible will be a great help as you write
your document and when you change it.  This sample
file shows how you can add comments to your own input
file.

Because printing is different from typewriting,
there are a number of things that you have to do
differently when preparing an input file than if
you were just typing the document directly.
Quotation marks like
       ``this''
have to be handled specially, as do quotes within
quotes:
       ``\,`this'            % \, separates the double and single quote.
        is what I just
        wrote, not  `that'\,''.

Dashes come in three sizes: an
       intra-word
dash, a medium dash for number ranges like
       1--2,
and a punctuation
       dash---like
this.

A sentence-ending space should be larger than the
space between words within a sentence.  You
sometimes have to type special commands in
conjunction with punctuation characters to get
this right, as in the following sentence.
       Gnats, gnus, etc.\ all  % `\ ' makes an inter-word space.
       begin with G\@.         % \@ marks end-of-sentence punctuation.
You should check the spaces after periods when
reading your output to make sure you haven't
forgotten any special cases.  Generating an
ellipsis
       \ldots\               % `\ ' is needed after `\ldots' because TeX
                             % ignores spaces after command names like \ldots
                             % made from \ + letters.
                             %
                             % Note how a `%' character causes TeX to ignore
                             % the end of the input line, so these blank lines
                             % do not start a new paragraph.
                             %
with the right spacing around the periods requires
a special command.

\LaTeX\ interprets some common characters as
commands, so you must type special commands to
generate them.  These characters include the
following:
       \$ \& \% \# \{ and \}.

In printing, text is usually emphasized with an
       \emph{italic}
type style.

\begin{em}
   A long segment of text can also be emphasized
   in this way.  Text within such a segment can be
   given \emph{additional} emphasis.
\end{em}

It is sometimes necessary to prevent \LaTeX\ from
breaking a line where it might otherwise do so.
This may be at a space, as between the ``Mr.''\ and
``Jones'' in
       ``Mr.~Jones'',        % ~ produces an unbreakable interword space.
or within a word---especially when the word is a
symbol like
       \mbox{\emph{itemnum}}
that makes little sense when hyphenated across
lines.

Footnotes\footnote{This is an example of a footnote.}
pose no problem.

\LaTeX\ is good at typesetting mathematical formulas
like
       \( x-3y + z = 7 \)
or
       \( a_{1} > x^{2n} + y^{2n} > x' \)
or
       \( \ip{A}{B} = \sum_{i} a_{i} b_{i} \).
The spaces you type in a formula are
ignored.  Remember that a letter like
       $x$                   % $ ... $  and  \( ... \)  are equivalent
is a formula when it denotes a mathematical
symbol, and it should be typed as one.

\section{Displayed Text}

Text is displayed by indenting it from the left
margin.  Quotations are commonly displayed.  There
are short quotations
\begin{quote}
   This is a short quotation.  It consists of a
   single paragraph of text.  See how it is formatted.
\end{quote}
and longer ones.
\begin{quotation}
   This is a longer quotation.  It consists of two
   paragraphs of text, neither of which are
   particularly interesting.

   This is the second paragraph of the quotation.  It
   is just as dull as the first paragraph.
\end{quotation}
Another frequently-displayed structure is a list.
The following is an example of an \emph{itemized}
list.
\begin{itemize}
   \item This is the first item of an itemized list.
         Each item in the list is marked with a ``tick''.
         You don't have to worry about what kind of tick
         mark is used.

   \item This is the second item of the list.  It
         contains another list nested inside it.  The inner
         list is an \emph{enumerated} list.
         \begin{enumerate}
            \item This is the first item of an enumerated
                  list that is nested within the itemized list.

            \item This is the second item of the inner list.
                  \LaTeX\ allows you to nest lists deeper than
                  you really should.
         \end{enumerate}
         This is the rest of the second item of the outer
         list.  It is no more interesting than any other
         part of the item.
   \item This is the third item of the list.
\end{itemize}
You can even display poetry.
\begin{verse}
   There is an environment
    for verse \\             % The \\ command separates lines
   Whose features some poets % within a stanza.
   will curse.

                             % One or more blank lines separate stanzas.

   For instead of making\\
   Them do \emph{all} line breaking, \\
   It allows them to put too many words on a line when they'd rather be
   forced to be terse.
\end{verse}

Mathematical formulas may also be displayed.  A
displayed formula
is
one-line long; multiline
formulas require special formatting instructions.
   \[  \ip{\Gamma}{\psi'} = x'' + y^{2} + z_{i}^{n}\]
Don't start a paragraph with a displayed equation,
nor make one a paragraph by itself.

\end{document}               % End of document.

\section{Bench Huge Section 651}
% synthetic bench section 651
% This is a sample LaTeX input file.  (Version of 12 August 2004.)
%
% A '%' character causes TeX to ignore all remaining text on the line,
% and is used for comments like this one.

\documentclass{article}      % Specifies the document class

                             % The preamble begins here.
\title{An Example Document}  % Declares the document's title.
\author{Leslie Lamport}      % Declares the author's name.
\date{January 21, 1994}      % Deleting this command produces today's date.

\newcommand{\ip}[2]{(#1, #2)}
                             % Defines \ip{arg1}{arg2} to mean
                             % (arg1, arg2).

%\newcommand{\ip}[2]{\langle #1 | #2\rangle}
                             % This is an alternative definition of
                             % \ip that is commented out.

\begin{document}             % End of preamble and beginning of text.

\maketitle                   % Produces the title.

This is an example input file.  Comparing it with
the output it generates can show you how to
produce a simple document of your own.

\section{Ordinary Text}      % Produces section heading.  Lower-level
                             % sections are begun with similar
                             % \subsection and \subsubsection commands.

The ends  of words and sentences are marked
  by   spaces. It  doesn't matter how many
spaces    you type; one is as good as 100.  The
end of   a line counts as a space.

One   or more   blank lines denote the  end
of  a paragraph.

Since any number of consecutive spaces are treated
like a single one, the formatting of the input
file makes no difference to
      \LaTeX,                % The \LaTeX command generates the LaTeX logo.
but it makes a difference to you.  When you use
\LaTeX, making your input file as easy to read
as possible will be a great help as you write
your document and when you change it.  This sample
file shows how you can add comments to your own input
file.

Because printing is different from typewriting,
there are a number of things that you have to do
differently when preparing an input file than if
you were just typing the document directly.
Quotation marks like
       ``this''
have to be handled specially, as do quotes within
quotes:
       ``\,`this'            % \, separates the double and single quote.
        is what I just
        wrote, not  `that'\,''.

Dashes come in three sizes: an
       intra-word
dash, a medium dash for number ranges like
       1--2,
and a punctuation
       dash---like
this.

A sentence-ending space should be larger than the
space between words within a sentence.  You
sometimes have to type special commands in
conjunction with punctuation characters to get
this right, as in the following sentence.
       Gnats, gnus, etc.\ all  % `\ ' makes an inter-word space.
       begin with G\@.         % \@ marks end-of-sentence punctuation.
You should check the spaces after periods when
reading your output to make sure you haven't
forgotten any special cases.  Generating an
ellipsis
       \ldots\               % `\ ' is needed after `\ldots' because TeX
                             % ignores spaces after command names like \ldots
                             % made from \ + letters.
                             %
                             % Note how a `%' character causes TeX to ignore
                             % the end of the input line, so these blank lines
                             % do not start a new paragraph.
                             %
with the right spacing around the periods requires
a special command.

\LaTeX\ interprets some common characters as
commands, so you must type special commands to
generate them.  These characters include the
following:
       \$ \& \% \# \{ and \}.

In printing, text is usually emphasized with an
       \emph{italic}
type style.

\begin{em}
   A long segment of text can also be emphasized
   in this way.  Text within such a segment can be
   given \emph{additional} emphasis.
\end{em}

It is sometimes necessary to prevent \LaTeX\ from
breaking a line where it might otherwise do so.
This may be at a space, as between the ``Mr.''\ and
``Jones'' in
       ``Mr.~Jones'',        % ~ produces an unbreakable interword space.
or within a word---especially when the word is a
symbol like
       \mbox{\emph{itemnum}}
that makes little sense when hyphenated across
lines.

Footnotes\footnote{This is an example of a footnote.}
pose no problem.

\LaTeX\ is good at typesetting mathematical formulas
like
       \( x-3y + z = 7 \)
or
       \( a_{1} > x^{2n} + y^{2n} > x' \)
or
       \( \ip{A}{B} = \sum_{i} a_{i} b_{i} \).
The spaces you type in a formula are
ignored.  Remember that a letter like
       $x$                   % $ ... $  and  \( ... \)  are equivalent
is a formula when it denotes a mathematical
symbol, and it should be typed as one.

\section{Displayed Text}

Text is displayed by indenting it from the left
margin.  Quotations are commonly displayed.  There
are short quotations
\begin{quote}
   This is a short quotation.  It consists of a
   single paragraph of text.  See how it is formatted.
\end{quote}
and longer ones.
\begin{quotation}
   This is a longer quotation.  It consists of two
   paragraphs of text, neither of which are
   particularly interesting.

   This is the second paragraph of the quotation.  It
   is just as dull as the first paragraph.
\end{quotation}
Another frequently-displayed structure is a list.
The following is an example of an \emph{itemized}
list.
\begin{itemize}
   \item This is the first item of an itemized list.
         Each item in the list is marked with a ``tick''.
         You don't have to worry about what kind of tick
         mark is used.

   \item This is the second item of the list.  It
         contains another list nested inside it.  The inner
         list is an \emph{enumerated} list.
         \begin{enumerate}
            \item This is the first item of an enumerated
                  list that is nested within the itemized list.

            \item This is the second item of the inner list.
                  \LaTeX\ allows you to nest lists deeper than
                  you really should.
         \end{enumerate}
         This is the rest of the second item of the outer
         list.  It is no more interesting than any other
         part of the item.
   \item This is the third item of the list.
\end{itemize}
You can even display poetry.
\begin{verse}
   There is an environment
    for verse \\             % The \\ command separates lines
   Whose features some poets % within a stanza.
   will curse.

                             % One or more blank lines separate stanzas.

   For instead of making\\
   Them do \emph{all} line breaking, \\
   It allows them to put too many words on a line when they'd rather be
   forced to be terse.
\end{verse}

Mathematical formulas may also be displayed.  A
displayed formula
is
one-line long; multiline
formulas require special formatting instructions.
   \[  \ip{\Gamma}{\psi'} = x'' + y^{2} + z_{i}^{n}\]
Don't start a paragraph with a displayed equation,
nor make one a paragraph by itself.

\end{document}               % End of document.

\section{Bench Huge Section 652}
% synthetic bench section 652
% This is a sample LaTeX input file.  (Version of 12 August 2004.)
%
% A '%' character causes TeX to ignore all remaining text on the line,
% and is used for comments like this one.

\documentclass{article}      % Specifies the document class

                             % The preamble begins here.
\title{An Example Document}  % Declares the document's title.
\author{Leslie Lamport}      % Declares the author's name.
\date{January 21, 1994}      % Deleting this command produces today's date.

\newcommand{\ip}[2]{(#1, #2)}
                             % Defines \ip{arg1}{arg2} to mean
                             % (arg1, arg2).

%\newcommand{\ip}[2]{\langle #1 | #2\rangle}
                             % This is an alternative definition of
                             % \ip that is commented out.

\begin{document}             % End of preamble and beginning of text.

\maketitle                   % Produces the title.

This is an example input file.  Comparing it with
the output it generates can show you how to
produce a simple document of your own.

\section{Ordinary Text}      % Produces section heading.  Lower-level
                             % sections are begun with similar
                             % \subsection and \subsubsection commands.

The ends  of words and sentences are marked
  by   spaces. It  doesn't matter how many
spaces    you type; one is as good as 100.  The
end of   a line counts as a space.

One   or more   blank lines denote the  end
of  a paragraph.

Since any number of consecutive spaces are treated
like a single one, the formatting of the input
file makes no difference to
      \LaTeX,                % The \LaTeX command generates the LaTeX logo.
but it makes a difference to you.  When you use
\LaTeX, making your input file as easy to read
as possible will be a great help as you write
your document and when you change it.  This sample
file shows how you can add comments to your own input
file.

Because printing is different from typewriting,
there are a number of things that you have to do
differently when preparing an input file than if
you were just typing the document directly.
Quotation marks like
       ``this''
have to be handled specially, as do quotes within
quotes:
       ``\,`this'            % \, separates the double and single quote.
        is what I just
        wrote, not  `that'\,''.

Dashes come in three sizes: an
       intra-word
dash, a medium dash for number ranges like
       1--2,
and a punctuation
       dash---like
this.

A sentence-ending space should be larger than the
space between words within a sentence.  You
sometimes have to type special commands in
conjunction with punctuation characters to get
this right, as in the following sentence.
       Gnats, gnus, etc.\ all  % `\ ' makes an inter-word space.
       begin with G\@.         % \@ marks end-of-sentence punctuation.
You should check the spaces after periods when
reading your output to make sure you haven't
forgotten any special cases.  Generating an
ellipsis
       \ldots\               % `\ ' is needed after `\ldots' because TeX
                             % ignores spaces after command names like \ldots
                             % made from \ + letters.
                             %
                             % Note how a `%' character causes TeX to ignore
                             % the end of the input line, so these blank lines
                             % do not start a new paragraph.
                             %
with the right spacing around the periods requires
a special command.

\LaTeX\ interprets some common characters as
commands, so you must type special commands to
generate them.  These characters include the
following:
       \$ \& \% \# \{ and \}.

In printing, text is usually emphasized with an
       \emph{italic}
type style.

\begin{em}
   A long segment of text can also be emphasized
   in this way.  Text within such a segment can be
   given \emph{additional} emphasis.
\end{em}

It is sometimes necessary to prevent \LaTeX\ from
breaking a line where it might otherwise do so.
This may be at a space, as between the ``Mr.''\ and
``Jones'' in
       ``Mr.~Jones'',        % ~ produces an unbreakable interword space.
or within a word---especially when the word is a
symbol like
       \mbox{\emph{itemnum}}
that makes little sense when hyphenated across
lines.

Footnotes\footnote{This is an example of a footnote.}
pose no problem.

\LaTeX\ is good at typesetting mathematical formulas
like
       \( x-3y + z = 7 \)
or
       \( a_{1} > x^{2n} + y^{2n} > x' \)
or
       \( \ip{A}{B} = \sum_{i} a_{i} b_{i} \).
The spaces you type in a formula are
ignored.  Remember that a letter like
       $x$                   % $ ... $  and  \( ... \)  are equivalent
is a formula when it denotes a mathematical
symbol, and it should be typed as one.

\section{Displayed Text}

Text is displayed by indenting it from the left
margin.  Quotations are commonly displayed.  There
are short quotations
\begin{quote}
   This is a short quotation.  It consists of a
   single paragraph of text.  See how it is formatted.
\end{quote}
and longer ones.
\begin{quotation}
   This is a longer quotation.  It consists of two
   paragraphs of text, neither of which are
   particularly interesting.

   This is the second paragraph of the quotation.  It
   is just as dull as the first paragraph.
\end{quotation}
Another frequently-displayed structure is a list.
The following is an example of an \emph{itemized}
list.
\begin{itemize}
   \item This is the first item of an itemized list.
         Each item in the list is marked with a ``tick''.
         You don't have to worry about what kind of tick
         mark is used.

   \item This is the second item of the list.  It
         contains another list nested inside it.  The inner
         list is an \emph{enumerated} list.
         \begin{enumerate}
            \item This is the first item of an enumerated
                  list that is nested within the itemized list.

            \item This is the second item of the inner list.
                  \LaTeX\ allows you to nest lists deeper than
                  you really should.
         \end{enumerate}
         This is the rest of the second item of the outer
         list.  It is no more interesting than any other
         part of the item.
   \item This is the third item of the list.
\end{itemize}
You can even display poetry.
\begin{verse}
   There is an environment
    for verse \\             % The \\ command separates lines
   Whose features some poets % within a stanza.
   will curse.

                             % One or more blank lines separate stanzas.

   For instead of making\\
   Them do \emph{all} line breaking, \\
   It allows them to put too many words on a line when they'd rather be
   forced to be terse.
\end{verse}

Mathematical formulas may also be displayed.  A
displayed formula
is
one-line long; multiline
formulas require special formatting instructions.
   \[  \ip{\Gamma}{\psi'} = x'' + y^{2} + z_{i}^{n}\]
Don't start a paragraph with a displayed equation,
nor make one a paragraph by itself.

\end{document}               % End of document.

\section{Bench Huge Section 653}
% synthetic bench section 653
% This is a sample LaTeX input file.  (Version of 12 August 2004.)
%
% A '%' character causes TeX to ignore all remaining text on the line,
% and is used for comments like this one.

\documentclass{article}      % Specifies the document class

                             % The preamble begins here.
\title{An Example Document}  % Declares the document's title.
\author{Leslie Lamport}      % Declares the author's name.
\date{January 21, 1994}      % Deleting this command produces today's date.

\newcommand{\ip}[2]{(#1, #2)}
                             % Defines \ip{arg1}{arg2} to mean
                             % (arg1, arg2).

%\newcommand{\ip}[2]{\langle #1 | #2\rangle}
                             % This is an alternative definition of
                             % \ip that is commented out.

\begin{document}             % End of preamble and beginning of text.

\maketitle                   % Produces the title.

This is an example input file.  Comparing it with
the output it generates can show you how to
produce a simple document of your own.

\section{Ordinary Text}      % Produces section heading.  Lower-level
                             % sections are begun with similar
                             % \subsection and \subsubsection commands.

The ends  of words and sentences are marked
  by   spaces. It  doesn't matter how many
spaces    you type; one is as good as 100.  The
end of   a line counts as a space.

One   or more   blank lines denote the  end
of  a paragraph.

Since any number of consecutive spaces are treated
like a single one, the formatting of the input
file makes no difference to
      \LaTeX,                % The \LaTeX command generates the LaTeX logo.
but it makes a difference to you.  When you use
\LaTeX, making your input file as easy to read
as possible will be a great help as you write
your document and when you change it.  This sample
file shows how you can add comments to your own input
file.

Because printing is different from typewriting,
there are a number of things that you have to do
differently when preparing an input file than if
you were just typing the document directly.
Quotation marks like
       ``this''
have to be handled specially, as do quotes within
quotes:
       ``\,`this'            % \, separates the double and single quote.
        is what I just
        wrote, not  `that'\,''.

Dashes come in three sizes: an
       intra-word
dash, a medium dash for number ranges like
       1--2,
and a punctuation
       dash---like
this.

A sentence-ending space should be larger than the
space between words within a sentence.  You
sometimes have to type special commands in
conjunction with punctuation characters to get
this right, as in the following sentence.
       Gnats, gnus, etc.\ all  % `\ ' makes an inter-word space.
       begin with G\@.         % \@ marks end-of-sentence punctuation.
You should check the spaces after periods when
reading your output to make sure you haven't
forgotten any special cases.  Generating an
ellipsis
       \ldots\               % `\ ' is needed after `\ldots' because TeX
                             % ignores spaces after command names like \ldots
                             % made from \ + letters.
                             %
                             % Note how a `%' character causes TeX to ignore
                             % the end of the input line, so these blank lines
                             % do not start a new paragraph.
                             %
with the right spacing around the periods requires
a special command.

\LaTeX\ interprets some common characters as
commands, so you must type special commands to
generate them.  These characters include the
following:
       \$ \& \% \# \{ and \}.

In printing, text is usually emphasized with an
       \emph{italic}
type style.

\begin{em}
   A long segment of text can also be emphasized
   in this way.  Text within such a segment can be
   given \emph{additional} emphasis.
\end{em}

It is sometimes necessary to prevent \LaTeX\ from
breaking a line where it might otherwise do so.
This may be at a space, as between the ``Mr.''\ and
``Jones'' in
       ``Mr.~Jones'',        % ~ produces an unbreakable interword space.
or within a word---especially when the word is a
symbol like
       \mbox{\emph{itemnum}}
that makes little sense when hyphenated across
lines.

Footnotes\footnote{This is an example of a footnote.}
pose no problem.

\LaTeX\ is good at typesetting mathematical formulas
like
       \( x-3y + z = 7 \)
or
       \( a_{1} > x^{2n} + y^{2n} > x' \)
or
       \( \ip{A}{B} = \sum_{i} a_{i} b_{i} \).
The spaces you type in a formula are
ignored.  Remember that a letter like
       $x$                   % $ ... $  and  \( ... \)  are equivalent
is a formula when it denotes a mathematical
symbol, and it should be typed as one.

\section{Displayed Text}

Text is displayed by indenting it from the left
margin.  Quotations are commonly displayed.  There
are short quotations
\begin{quote}
   This is a short quotation.  It consists of a
   single paragraph of text.  See how it is formatted.
\end{quote}
and longer ones.
\begin{quotation}
   This is a longer quotation.  It consists of two
   paragraphs of text, neither of which are
   particularly interesting.

   This is the second paragraph of the quotation.  It
   is just as dull as the first paragraph.
\end{quotation}
Another frequently-displayed structure is a list.
The following is an example of an \emph{itemized}
list.
\begin{itemize}
   \item This is the first item of an itemized list.
         Each item in the list is marked with a ``tick''.
         You don't have to worry about what kind of tick
         mark is used.

   \item This is the second item of the list.  It
         contains another list nested inside it.  The inner
         list is an \emph{enumerated} list.
         \begin{enumerate}
            \item This is the first item of an enumerated
                  list that is nested within the itemized list.

            \item This is the second item of the inner list.
                  \LaTeX\ allows you to nest lists deeper than
                  you really should.
         \end{enumerate}
         This is the rest of the second item of the outer
         list.  It is no more interesting than any other
         part of the item.
   \item This is the third item of the list.
\end{itemize}
You can even display poetry.
\begin{verse}
   There is an environment
    for verse \\             % The \\ command separates lines
   Whose features some poets % within a stanza.
   will curse.

                             % One or more blank lines separate stanzas.

   For instead of making\\
   Them do \emph{all} line breaking, \\
   It allows them to put too many words on a line when they'd rather be
   forced to be terse.
\end{verse}

Mathematical formulas may also be displayed.  A
displayed formula
is
one-line long; multiline
formulas require special formatting instructions.
   \[  \ip{\Gamma}{\psi'} = x'' + y^{2} + z_{i}^{n}\]
Don't start a paragraph with a displayed equation,
nor make one a paragraph by itself.

\end{document}               % End of document.

\section{Bench Huge Section 654}
% synthetic bench section 654
% This is a sample LaTeX input file.  (Version of 12 August 2004.)
%
% A '%' character causes TeX to ignore all remaining text on the line,
% and is used for comments like this one.

\documentclass{article}      % Specifies the document class

                             % The preamble begins here.
\title{An Example Document}  % Declares the document's title.
\author{Leslie Lamport}      % Declares the author's name.
\date{January 21, 1994}      % Deleting this command produces today's date.

\newcommand{\ip}[2]{(#1, #2)}
                             % Defines \ip{arg1}{arg2} to mean
                             % (arg1, arg2).

%\newcommand{\ip}[2]{\langle #1 | #2\rangle}
                             % This is an alternative definition of
                             % \ip that is commented out.

\begin{document}             % End of preamble and beginning of text.

\maketitle                   % Produces the title.

This is an example input file.  Comparing it with
the output it generates can show you how to
produce a simple document of your own.

\section{Ordinary Text}      % Produces section heading.  Lower-level
                             % sections are begun with similar
                             % \subsection and \subsubsection commands.

The ends  of words and sentences are marked
  by   spaces. It  doesn't matter how many
spaces    you type; one is as good as 100.  The
end of   a line counts as a space.

One   or more   blank lines denote the  end
of  a paragraph.

Since any number of consecutive spaces are treated
like a single one, the formatting of the input
file makes no difference to
      \LaTeX,                % The \LaTeX command generates the LaTeX logo.
but it makes a difference to you.  When you use
\LaTeX, making your input file as easy to read
as possible will be a great help as you write
your document and when you change it.  This sample
file shows how you can add comments to your own input
file.

Because printing is different from typewriting,
there are a number of things that you have to do
differently when preparing an input file than if
you were just typing the document directly.
Quotation marks like
       ``this''
have to be handled specially, as do quotes within
quotes:
       ``\,`this'            % \, separates the double and single quote.
        is what I just
        wrote, not  `that'\,''.

Dashes come in three sizes: an
       intra-word
dash, a medium dash for number ranges like
       1--2,
and a punctuation
       dash---like
this.

A sentence-ending space should be larger than the
space between words within a sentence.  You
sometimes have to type special commands in
conjunction with punctuation characters to get
this right, as in the following sentence.
       Gnats, gnus, etc.\ all  % `\ ' makes an inter-word space.
       begin with G\@.         % \@ marks end-of-sentence punctuation.
You should check the spaces after periods when
reading your output to make sure you haven't
forgotten any special cases.  Generating an
ellipsis
       \ldots\               % `\ ' is needed after `\ldots' because TeX
                             % ignores spaces after command names like \ldots
                             % made from \ + letters.
                             %
                             % Note how a `%' character causes TeX to ignore
                             % the end of the input line, so these blank lines
                             % do not start a new paragraph.
                             %
with the right spacing around the periods requires
a special command.

\LaTeX\ interprets some common characters as
commands, so you must type special commands to
generate them.  These characters include the
following:
       \$ \& \% \# \{ and \}.

In printing, text is usually emphasized with an
       \emph{italic}
type style.

\begin{em}
   A long segment of text can also be emphasized
   in this way.  Text within such a segment can be
   given \emph{additional} emphasis.
\end{em}

It is sometimes necessary to prevent \LaTeX\ from
breaking a line where it might otherwise do so.
This may be at a space, as between the ``Mr.''\ and
``Jones'' in
       ``Mr.~Jones'',        % ~ produces an unbreakable interword space.
or within a word---especially when the word is a
symbol like
       \mbox{\emph{itemnum}}
that makes little sense when hyphenated across
lines.

Footnotes\footnote{This is an example of a footnote.}
pose no problem.

\LaTeX\ is good at typesetting mathematical formulas
like
       \( x-3y + z = 7 \)
or
       \( a_{1} > x^{2n} + y^{2n} > x' \)
or
       \( \ip{A}{B} = \sum_{i} a_{i} b_{i} \).
The spaces you type in a formula are
ignored.  Remember that a letter like
       $x$                   % $ ... $  and  \( ... \)  are equivalent
is a formula when it denotes a mathematical
symbol, and it should be typed as one.

\section{Displayed Text}

Text is displayed by indenting it from the left
margin.  Quotations are commonly displayed.  There
are short quotations
\begin{quote}
   This is a short quotation.  It consists of a
   single paragraph of text.  See how it is formatted.
\end{quote}
and longer ones.
\begin{quotation}
   This is a longer quotation.  It consists of two
   paragraphs of text, neither of which are
   particularly interesting.

   This is the second paragraph of the quotation.  It
   is just as dull as the first paragraph.
\end{quotation}
Another frequently-displayed structure is a list.
The following is an example of an \emph{itemized}
list.
\begin{itemize}
   \item This is the first item of an itemized list.
         Each item in the list is marked with a ``tick''.
         You don't have to worry about what kind of tick
         mark is used.

   \item This is the second item of the list.  It
         contains another list nested inside it.  The inner
         list is an \emph{enumerated} list.
         \begin{enumerate}
            \item This is the first item of an enumerated
                  list that is nested within the itemized list.

            \item This is the second item of the inner list.
                  \LaTeX\ allows you to nest lists deeper than
                  you really should.
         \end{enumerate}
         This is the rest of the second item of the outer
         list.  It is no more interesting than any other
         part of the item.
   \item This is the third item of the list.
\end{itemize}
You can even display poetry.
\begin{verse}
   There is an environment
    for verse \\             % The \\ command separates lines
   Whose features some poets % within a stanza.
   will curse.

                             % One or more blank lines separate stanzas.

   For instead of making\\
   Them do \emph{all} line breaking, \\
   It allows them to put too many words on a line when they'd rather be
   forced to be terse.
\end{verse}

Mathematical formulas may also be displayed.  A
displayed formula
is
one-line long; multiline
formulas require special formatting instructions.
   \[  \ip{\Gamma}{\psi'} = x'' + y^{2} + z_{i}^{n}\]
Don't start a paragraph with a displayed equation,
nor make one a paragraph by itself.

\end{document}               % End of document.

\section{Bench Huge Section 655}
% synthetic bench section 655
% This is a sample LaTeX input file.  (Version of 12 August 2004.)
%
% A '%' character causes TeX to ignore all remaining text on the line,
% and is used for comments like this one.

\documentclass{article}      % Specifies the document class

                             % The preamble begins here.
\title{An Example Document}  % Declares the document's title.
\author{Leslie Lamport}      % Declares the author's name.
\date{January 21, 1994}      % Deleting this command produces today's date.

\newcommand{\ip}[2]{(#1, #2)}
                             % Defines \ip{arg1}{arg2} to mean
                             % (arg1, arg2).

%\newcommand{\ip}[2]{\langle #1 | #2\rangle}
                             % This is an alternative definition of
                             % \ip that is commented out.

\begin{document}             % End of preamble and beginning of text.

\maketitle                   % Produces the title.

This is an example input file.  Comparing it with
the output it generates can show you how to
produce a simple document of your own.

\section{Ordinary Text}      % Produces section heading.  Lower-level
                             % sections are begun with similar
                             % \subsection and \subsubsection commands.

The ends  of words and sentences are marked
  by   spaces. It  doesn't matter how many
spaces    you type; one is as good as 100.  The
end of   a line counts as a space.

One   or more   blank lines denote the  end
of  a paragraph.

Since any number of consecutive spaces are treated
like a single one, the formatting of the input
file makes no difference to
      \LaTeX,                % The \LaTeX command generates the LaTeX logo.
but it makes a difference to you.  When you use
\LaTeX, making your input file as easy to read
as possible will be a great help as you write
your document and when you change it.  This sample
file shows how you can add comments to your own input
file.

Because printing is different from typewriting,
there are a number of things that you have to do
differently when preparing an input file than if
you were just typing the document directly.
Quotation marks like
       ``this''
have to be handled specially, as do quotes within
quotes:
       ``\,`this'            % \, separates the double and single quote.
        is what I just
        wrote, not  `that'\,''.

Dashes come in three sizes: an
       intra-word
dash, a medium dash for number ranges like
       1--2,
and a punctuation
       dash---like
this.

A sentence-ending space should be larger than the
space between words within a sentence.  You
sometimes have to type special commands in
conjunction with punctuation characters to get
this right, as in the following sentence.
       Gnats, gnus, etc.\ all  % `\ ' makes an inter-word space.
       begin with G\@.         % \@ marks end-of-sentence punctuation.
You should check the spaces after periods when
reading your output to make sure you haven't
forgotten any special cases.  Generating an
ellipsis
       \ldots\               % `\ ' is needed after `\ldots' because TeX
                             % ignores spaces after command names like \ldots
                             % made from \ + letters.
                             %
                             % Note how a `%' character causes TeX to ignore
                             % the end of the input line, so these blank lines
                             % do not start a new paragraph.
                             %
with the right spacing around the periods requires
a special command.

\LaTeX\ interprets some common characters as
commands, so you must type special commands to
generate them.  These characters include the
following:
       \$ \& \% \# \{ and \}.

In printing, text is usually emphasized with an
       \emph{italic}
type style.

\begin{em}
   A long segment of text can also be emphasized
   in this way.  Text within such a segment can be
   given \emph{additional} emphasis.
\end{em}

It is sometimes necessary to prevent \LaTeX\ from
breaking a line where it might otherwise do so.
This may be at a space, as between the ``Mr.''\ and
``Jones'' in
       ``Mr.~Jones'',        % ~ produces an unbreakable interword space.
or within a word---especially when the word is a
symbol like
       \mbox{\emph{itemnum}}
that makes little sense when hyphenated across
lines.

Footnotes\footnote{This is an example of a footnote.}
pose no problem.

\LaTeX\ is good at typesetting mathematical formulas
like
       \( x-3y + z = 7 \)
or
       \( a_{1} > x^{2n} + y^{2n} > x' \)
or
       \( \ip{A}{B} = \sum_{i} a_{i} b_{i} \).
The spaces you type in a formula are
ignored.  Remember that a letter like
       $x$                   % $ ... $  and  \( ... \)  are equivalent
is a formula when it denotes a mathematical
symbol, and it should be typed as one.

\section{Displayed Text}

Text is displayed by indenting it from the left
margin.  Quotations are commonly displayed.  There
are short quotations
\begin{quote}
   This is a short quotation.  It consists of a
   single paragraph of text.  See how it is formatted.
\end{quote}
and longer ones.
\begin{quotation}
   This is a longer quotation.  It consists of two
   paragraphs of text, neither of which are
   particularly interesting.

   This is the second paragraph of the quotation.  It
   is just as dull as the first paragraph.
\end{quotation}
Another frequently-displayed structure is a list.
The following is an example of an \emph{itemized}
list.
\begin{itemize}
   \item This is the first item of an itemized list.
         Each item in the list is marked with a ``tick''.
         You don't have to worry about what kind of tick
         mark is used.

   \item This is the second item of the list.  It
         contains another list nested inside it.  The inner
         list is an \emph{enumerated} list.
         \begin{enumerate}
            \item This is the first item of an enumerated
                  list that is nested within the itemized list.

            \item This is the second item of the inner list.
                  \LaTeX\ allows you to nest lists deeper than
                  you really should.
         \end{enumerate}
         This is the rest of the second item of the outer
         list.  It is no more interesting than any other
         part of the item.
   \item This is the third item of the list.
\end{itemize}
You can even display poetry.
\begin{verse}
   There is an environment
    for verse \\             % The \\ command separates lines
   Whose features some poets % within a stanza.
   will curse.

                             % One or more blank lines separate stanzas.

   For instead of making\\
   Them do \emph{all} line breaking, \\
   It allows them to put too many words on a line when they'd rather be
   forced to be terse.
\end{verse}

Mathematical formulas may also be displayed.  A
displayed formula
is
one-line long; multiline
formulas require special formatting instructions.
   \[  \ip{\Gamma}{\psi'} = x'' + y^{2} + z_{i}^{n}\]
Don't start a paragraph with a displayed equation,
nor make one a paragraph by itself.

\end{document}               % End of document.

\section{Bench Huge Section 656}
% synthetic bench section 656
% This is a sample LaTeX input file.  (Version of 12 August 2004.)
%
% A '%' character causes TeX to ignore all remaining text on the line,
% and is used for comments like this one.

\documentclass{article}      % Specifies the document class

                             % The preamble begins here.
\title{An Example Document}  % Declares the document's title.
\author{Leslie Lamport}      % Declares the author's name.
\date{January 21, 1994}      % Deleting this command produces today's date.

\newcommand{\ip}[2]{(#1, #2)}
                             % Defines \ip{arg1}{arg2} to mean
                             % (arg1, arg2).

%\newcommand{\ip}[2]{\langle #1 | #2\rangle}
                             % This is an alternative definition of
                             % \ip that is commented out.

\begin{document}             % End of preamble and beginning of text.

\maketitle                   % Produces the title.

This is an example input file.  Comparing it with
the output it generates can show you how to
produce a simple document of your own.

\section{Ordinary Text}      % Produces section heading.  Lower-level
                             % sections are begun with similar
                             % \subsection and \subsubsection commands.

The ends  of words and sentences are marked
  by   spaces. It  doesn't matter how many
spaces    you type; one is as good as 100.  The
end of   a line counts as a space.

One   or more   blank lines denote the  end
of  a paragraph.

Since any number of consecutive spaces are treated
like a single one, the formatting of the input
file makes no difference to
      \LaTeX,                % The \LaTeX command generates the LaTeX logo.
but it makes a difference to you.  When you use
\LaTeX, making your input file as easy to read
as possible will be a great help as you write
your document and when you change it.  This sample
file shows how you can add comments to your own input
file.

Because printing is different from typewriting,
there are a number of things that you have to do
differently when preparing an input file than if
you were just typing the document directly.
Quotation marks like
       ``this''
have to be handled specially, as do quotes within
quotes:
       ``\,`this'            % \, separates the double and single quote.
        is what I just
        wrote, not  `that'\,''.

Dashes come in three sizes: an
       intra-word
dash, a medium dash for number ranges like
       1--2,
and a punctuation
       dash---like
this.

A sentence-ending space should be larger than the
space between words within a sentence.  You
sometimes have to type special commands in
conjunction with punctuation characters to get
this right, as in the following sentence.
       Gnats, gnus, etc.\ all  % `\ ' makes an inter-word space.
       begin with G\@.         % \@ marks end-of-sentence punctuation.
You should check the spaces after periods when
reading your output to make sure you haven't
forgotten any special cases.  Generating an
ellipsis
       \ldots\               % `\ ' is needed after `\ldots' because TeX
                             % ignores spaces after command names like \ldots
                             % made from \ + letters.
                             %
                             % Note how a `%' character causes TeX to ignore
                             % the end of the input line, so these blank lines
                             % do not start a new paragraph.
                             %
with the right spacing around the periods requires
a special command.

\LaTeX\ interprets some common characters as
commands, so you must type special commands to
generate them.  These characters include the
following:
       \$ \& \% \# \{ and \}.

In printing, text is usually emphasized with an
       \emph{italic}
type style.

\begin{em}
   A long segment of text can also be emphasized
   in this way.  Text within such a segment can be
   given \emph{additional} emphasis.
\end{em}

It is sometimes necessary to prevent \LaTeX\ from
breaking a line where it might otherwise do so.
This may be at a space, as between the ``Mr.''\ and
``Jones'' in
       ``Mr.~Jones'',        % ~ produces an unbreakable interword space.
or within a word---especially when the word is a
symbol like
       \mbox{\emph{itemnum}}
that makes little sense when hyphenated across
lines.

Footnotes\footnote{This is an example of a footnote.}
pose no problem.

\LaTeX\ is good at typesetting mathematical formulas
like
       \( x-3y + z = 7 \)
or
       \( a_{1} > x^{2n} + y^{2n} > x' \)
or
       \( \ip{A}{B} = \sum_{i} a_{i} b_{i} \).
The spaces you type in a formula are
ignored.  Remember that a letter like
       $x$                   % $ ... $  and  \( ... \)  are equivalent
is a formula when it denotes a mathematical
symbol, and it should be typed as one.

\section{Displayed Text}

Text is displayed by indenting it from the left
margin.  Quotations are commonly displayed.  There
are short quotations
\begin{quote}
   This is a short quotation.  It consists of a
   single paragraph of text.  See how it is formatted.
\end{quote}
and longer ones.
\begin{quotation}
   This is a longer quotation.  It consists of two
   paragraphs of text, neither of which are
   particularly interesting.

   This is the second paragraph of the quotation.  It
   is just as dull as the first paragraph.
\end{quotation}
Another frequently-displayed structure is a list.
The following is an example of an \emph{itemized}
list.
\begin{itemize}
   \item This is the first item of an itemized list.
         Each item in the list is marked with a ``tick''.
         You don't have to worry about what kind of tick
         mark is used.

   \item This is the second item of the list.  It
         contains another list nested inside it.  The inner
         list is an \emph{enumerated} list.
         \begin{enumerate}
            \item This is the first item of an enumerated
                  list that is nested within the itemized list.

            \item This is the second item of the inner list.
                  \LaTeX\ allows you to nest lists deeper than
                  you really should.
         \end{enumerate}
         This is the rest of the second item of the outer
         list.  It is no more interesting than any other
         part of the item.
   \item This is the third item of the list.
\end{itemize}
You can even display poetry.
\begin{verse}
   There is an environment
    for verse \\             % The \\ command separates lines
   Whose features some poets % within a stanza.
   will curse.

                             % One or more blank lines separate stanzas.

   For instead of making\\
   Them do \emph{all} line breaking, \\
   It allows them to put too many words on a line when they'd rather be
   forced to be terse.
\end{verse}

Mathematical formulas may also be displayed.  A
displayed formula
is
one-line long; multiline
formulas require special formatting instructions.
   \[  \ip{\Gamma}{\psi'} = x'' + y^{2} + z_{i}^{n}\]
Don't start a paragraph with a displayed equation,
nor make one a paragraph by itself.

\end{document}               % End of document.

\section{Bench Huge Section 657}
% synthetic bench section 657
% This is a sample LaTeX input file.  (Version of 12 August 2004.)
%
% A '%' character causes TeX to ignore all remaining text on the line,
% and is used for comments like this one.

\documentclass{article}      % Specifies the document class

                             % The preamble begins here.
\title{An Example Document}  % Declares the document's title.
\author{Leslie Lamport}      % Declares the author's name.
\date{January 21, 1994}      % Deleting this command produces today's date.

\newcommand{\ip}[2]{(#1, #2)}
                             % Defines \ip{arg1}{arg2} to mean
                             % (arg1, arg2).

%\newcommand{\ip}[2]{\langle #1 | #2\rangle}
                             % This is an alternative definition of
                             % \ip that is commented out.

\begin{document}             % End of preamble and beginning of text.

\maketitle                   % Produces the title.

This is an example input file.  Comparing it with
the output it generates can show you how to
produce a simple document of your own.

\section{Ordinary Text}      % Produces section heading.  Lower-level
                             % sections are begun with similar
                             % \subsection and \subsubsection commands.

The ends  of words and sentences are marked
  by   spaces. It  doesn't matter how many
spaces    you type; one is as good as 100.  The
end of   a line counts as a space.

One   or more   blank lines denote the  end
of  a paragraph.

Since any number of consecutive spaces are treated
like a single one, the formatting of the input
file makes no difference to
      \LaTeX,                % The \LaTeX command generates the LaTeX logo.
but it makes a difference to you.  When you use
\LaTeX, making your input file as easy to read
as possible will be a great help as you write
your document and when you change it.  This sample
file shows how you can add comments to your own input
file.

Because printing is different from typewriting,
there are a number of things that you have to do
differently when preparing an input file than if
you were just typing the document directly.
Quotation marks like
       ``this''
have to be handled specially, as do quotes within
quotes:
       ``\,`this'            % \, separates the double and single quote.
        is what I just
        wrote, not  `that'\,''.

Dashes come in three sizes: an
       intra-word
dash, a medium dash for number ranges like
       1--2,
and a punctuation
       dash---like
this.

A sentence-ending space should be larger than the
space between words within a sentence.  You
sometimes have to type special commands in
conjunction with punctuation characters to get
this right, as in the following sentence.
       Gnats, gnus, etc.\ all  % `\ ' makes an inter-word space.
       begin with G\@.         % \@ marks end-of-sentence punctuation.
You should check the spaces after periods when
reading your output to make sure you haven't
forgotten any special cases.  Generating an
ellipsis
       \ldots\               % `\ ' is needed after `\ldots' because TeX
                             % ignores spaces after command names like \ldots
                             % made from \ + letters.
                             %
                             % Note how a `%' character causes TeX to ignore
                             % the end of the input line, so these blank lines
                             % do not start a new paragraph.
                             %
with the right spacing around the periods requires
a special command.

\LaTeX\ interprets some common characters as
commands, so you must type special commands to
generate them.  These characters include the
following:
       \$ \& \% \# \{ and \}.

In printing, text is usually emphasized with an
       \emph{italic}
type style.

\begin{em}
   A long segment of text can also be emphasized
   in this way.  Text within such a segment can be
   given \emph{additional} emphasis.
\end{em}

It is sometimes necessary to prevent \LaTeX\ from
breaking a line where it might otherwise do so.
This may be at a space, as between the ``Mr.''\ and
``Jones'' in
       ``Mr.~Jones'',        % ~ produces an unbreakable interword space.
or within a word---especially when the word is a
symbol like
       \mbox{\emph{itemnum}}
that makes little sense when hyphenated across
lines.

Footnotes\footnote{This is an example of a footnote.}
pose no problem.

\LaTeX\ is good at typesetting mathematical formulas
like
       \( x-3y + z = 7 \)
or
       \( a_{1} > x^{2n} + y^{2n} > x' \)
or
       \( \ip{A}{B} = \sum_{i} a_{i} b_{i} \).
The spaces you type in a formula are
ignored.  Remember that a letter like
       $x$                   % $ ... $  and  \( ... \)  are equivalent
is a formula when it denotes a mathematical
symbol, and it should be typed as one.

\section{Displayed Text}

Text is displayed by indenting it from the left
margin.  Quotations are commonly displayed.  There
are short quotations
\begin{quote}
   This is a short quotation.  It consists of a
   single paragraph of text.  See how it is formatted.
\end{quote}
and longer ones.
\begin{quotation}
   This is a longer quotation.  It consists of two
   paragraphs of text, neither of which are
   particularly interesting.

   This is the second paragraph of the quotation.  It
   is just as dull as the first paragraph.
\end{quotation}
Another frequently-displayed structure is a list.
The following is an example of an \emph{itemized}
list.
\begin{itemize}
   \item This is the first item of an itemized list.
         Each item in the list is marked with a ``tick''.
         You don't have to worry about what kind of tick
         mark is used.

   \item This is the second item of the list.  It
         contains another list nested inside it.  The inner
         list is an \emph{enumerated} list.
         \begin{enumerate}
            \item This is the first item of an enumerated
                  list that is nested within the itemized list.

            \item This is the second item of the inner list.
                  \LaTeX\ allows you to nest lists deeper than
                  you really should.
         \end{enumerate}
         This is the rest of the second item of the outer
         list.  It is no more interesting than any other
         part of the item.
   \item This is the third item of the list.
\end{itemize}
You can even display poetry.
\begin{verse}
   There is an environment
    for verse \\             % The \\ command separates lines
   Whose features some poets % within a stanza.
   will curse.

                             % One or more blank lines separate stanzas.

   For instead of making\\
   Them do \emph{all} line breaking, \\
   It allows them to put too many words on a line when they'd rather be
   forced to be terse.
\end{verse}

Mathematical formulas may also be displayed.  A
displayed formula
is
one-line long; multiline
formulas require special formatting instructions.
   \[  \ip{\Gamma}{\psi'} = x'' + y^{2} + z_{i}^{n}\]
Don't start a paragraph with a displayed equation,
nor make one a paragraph by itself.

\end{document}               % End of document.

\section{Bench Huge Section 658}
% synthetic bench section 658
% This is a sample LaTeX input file.  (Version of 12 August 2004.)
%
% A '%' character causes TeX to ignore all remaining text on the line,
% and is used for comments like this one.

\documentclass{article}      % Specifies the document class

                             % The preamble begins here.
\title{An Example Document}  % Declares the document's title.
\author{Leslie Lamport}      % Declares the author's name.
\date{January 21, 1994}      % Deleting this command produces today's date.

\newcommand{\ip}[2]{(#1, #2)}
                             % Defines \ip{arg1}{arg2} to mean
                             % (arg1, arg2).

%\newcommand{\ip}[2]{\langle #1 | #2\rangle}
                             % This is an alternative definition of
                             % \ip that is commented out.

\begin{document}             % End of preamble and beginning of text.

\maketitle                   % Produces the title.

This is an example input file.  Comparing it with
the output it generates can show you how to
produce a simple document of your own.

\section{Ordinary Text}      % Produces section heading.  Lower-level
                             % sections are begun with similar
                             % \subsection and \subsubsection commands.

The ends  of words and sentences are marked
  by   spaces. It  doesn't matter how many
spaces    you type; one is as good as 100.  The
end of   a line counts as a space.

One   or more   blank lines denote the  end
of  a paragraph.

Since any number of consecutive spaces are treated
like a single one, the formatting of the input
file makes no difference to
      \LaTeX,                % The \LaTeX command generates the LaTeX logo.
but it makes a difference to you.  When you use
\LaTeX, making your input file as easy to read
as possible will be a great help as you write
your document and when you change it.  This sample
file shows how you can add comments to your own input
file.

Because printing is different from typewriting,
there are a number of things that you have to do
differently when preparing an input file than if
you were just typing the document directly.
Quotation marks like
       ``this''
have to be handled specially, as do quotes within
quotes:
       ``\,`this'            % \, separates the double and single quote.
        is what I just
        wrote, not  `that'\,''.

Dashes come in three sizes: an
       intra-word
dash, a medium dash for number ranges like
       1--2,
and a punctuation
       dash---like
this.

A sentence-ending space should be larger than the
space between words within a sentence.  You
sometimes have to type special commands in
conjunction with punctuation characters to get
this right, as in the following sentence.
       Gnats, gnus, etc.\ all  % `\ ' makes an inter-word space.
       begin with G\@.         % \@ marks end-of-sentence punctuation.
You should check the spaces after periods when
reading your output to make sure you haven't
forgotten any special cases.  Generating an
ellipsis
       \ldots\               % `\ ' is needed after `\ldots' because TeX
                             % ignores spaces after command names like \ldots
                             % made from \ + letters.
                             %
                             % Note how a `%' character causes TeX to ignore
                             % the end of the input line, so these blank lines
                             % do not start a new paragraph.
                             %
with the right spacing around the periods requires
a special command.

\LaTeX\ interprets some common characters as
commands, so you must type special commands to
generate them.  These characters include the
following:
       \$ \& \% \# \{ and \}.

In printing, text is usually emphasized with an
       \emph{italic}
type style.

\begin{em}
   A long segment of text can also be emphasized
   in this way.  Text within such a segment can be
   given \emph{additional} emphasis.
\end{em}

It is sometimes necessary to prevent \LaTeX\ from
breaking a line where it might otherwise do so.
This may be at a space, as between the ``Mr.''\ and
``Jones'' in
       ``Mr.~Jones'',        % ~ produces an unbreakable interword space.
or within a word---especially when the word is a
symbol like
       \mbox{\emph{itemnum}}
that makes little sense when hyphenated across
lines.

Footnotes\footnote{This is an example of a footnote.}
pose no problem.

\LaTeX\ is good at typesetting mathematical formulas
like
       \( x-3y + z = 7 \)
or
       \( a_{1} > x^{2n} + y^{2n} > x' \)
or
       \( \ip{A}{B} = \sum_{i} a_{i} b_{i} \).
The spaces you type in a formula are
ignored.  Remember that a letter like
       $x$                   % $ ... $  and  \( ... \)  are equivalent
is a formula when it denotes a mathematical
symbol, and it should be typed as one.

\section{Displayed Text}

Text is displayed by indenting it from the left
margin.  Quotations are commonly displayed.  There
are short quotations
\begin{quote}
   This is a short quotation.  It consists of a
   single paragraph of text.  See how it is formatted.
\end{quote}
and longer ones.
\begin{quotation}
   This is a longer quotation.  It consists of two
   paragraphs of text, neither of which are
   particularly interesting.

   This is the second paragraph of the quotation.  It
   is just as dull as the first paragraph.
\end{quotation}
Another frequently-displayed structure is a list.
The following is an example of an \emph{itemized}
list.
\begin{itemize}
   \item This is the first item of an itemized list.
         Each item in the list is marked with a ``tick''.
         You don't have to worry about what kind of tick
         mark is used.

   \item This is the second item of the list.  It
         contains another list nested inside it.  The inner
         list is an \emph{enumerated} list.
         \begin{enumerate}
            \item This is the first item of an enumerated
                  list that is nested within the itemized list.

            \item This is the second item of the inner list.
                  \LaTeX\ allows you to nest lists deeper than
                  you really should.
         \end{enumerate}
         This is the rest of the second item of the outer
         list.  It is no more interesting than any other
         part of the item.
   \item This is the third item of the list.
\end{itemize}
You can even display poetry.
\begin{verse}
   There is an environment
    for verse \\             % The \\ command separates lines
   Whose features some poets % within a stanza.
   will curse.

                             % One or more blank lines separate stanzas.

   For instead of making\\
   Them do \emph{all} line breaking, \\
   It allows them to put too many words on a line when they'd rather be
   forced to be terse.
\end{verse}

Mathematical formulas may also be displayed.  A
displayed formula
is
one-line long; multiline
formulas require special formatting instructions.
   \[  \ip{\Gamma}{\psi'} = x'' + y^{2} + z_{i}^{n}\]
Don't start a paragraph with a displayed equation,
nor make one a paragraph by itself.

\end{document}               % End of document.

\section{Bench Huge Section 659}
% synthetic bench section 659
% This is a sample LaTeX input file.  (Version of 12 August 2004.)
%
% A '%' character causes TeX to ignore all remaining text on the line,
% and is used for comments like this one.

\documentclass{article}      % Specifies the document class

                             % The preamble begins here.
\title{An Example Document}  % Declares the document's title.
\author{Leslie Lamport}      % Declares the author's name.
\date{January 21, 1994}      % Deleting this command produces today's date.

\newcommand{\ip}[2]{(#1, #2)}
                             % Defines \ip{arg1}{arg2} to mean
                             % (arg1, arg2).

%\newcommand{\ip}[2]{\langle #1 | #2\rangle}
                             % This is an alternative definition of
                             % \ip that is commented out.

\begin{document}             % End of preamble and beginning of text.

\maketitle                   % Produces the title.

This is an example input file.  Comparing it with
the output it generates can show you how to
produce a simple document of your own.

\section{Ordinary Text}      % Produces section heading.  Lower-level
                             % sections are begun with similar
                             % \subsection and \subsubsection commands.

The ends  of words and sentences are marked
  by   spaces. It  doesn't matter how many
spaces    you type; one is as good as 100.  The
end of   a line counts as a space.

One   or more   blank lines denote the  end
of  a paragraph.

Since any number of consecutive spaces are treated
like a single one, the formatting of the input
file makes no difference to
      \LaTeX,                % The \LaTeX command generates the LaTeX logo.
but it makes a difference to you.  When you use
\LaTeX, making your input file as easy to read
as possible will be a great help as you write
your document and when you change it.  This sample
file shows how you can add comments to your own input
file.

Because printing is different from typewriting,
there are a number of things that you have to do
differently when preparing an input file than if
you were just typing the document directly.
Quotation marks like
       ``this''
have to be handled specially, as do quotes within
quotes:
       ``\,`this'            % \, separates the double and single quote.
        is what I just
        wrote, not  `that'\,''.

Dashes come in three sizes: an
       intra-word
dash, a medium dash for number ranges like
       1--2,
and a punctuation
       dash---like
this.

A sentence-ending space should be larger than the
space between words within a sentence.  You
sometimes have to type special commands in
conjunction with punctuation characters to get
this right, as in the following sentence.
       Gnats, gnus, etc.\ all  % `\ ' makes an inter-word space.
       begin with G\@.         % \@ marks end-of-sentence punctuation.
You should check the spaces after periods when
reading your output to make sure you haven't
forgotten any special cases.  Generating an
ellipsis
       \ldots\               % `\ ' is needed after `\ldots' because TeX
                             % ignores spaces after command names like \ldots
                             % made from \ + letters.
                             %
                             % Note how a `%' character causes TeX to ignore
                             % the end of the input line, so these blank lines
                             % do not start a new paragraph.
                             %
with the right spacing around the periods requires
a special command.

\LaTeX\ interprets some common characters as
commands, so you must type special commands to
generate them.  These characters include the
following:
       \$ \& \% \# \{ and \}.

In printing, text is usually emphasized with an
       \emph{italic}
type style.

\begin{em}
   A long segment of text can also be emphasized
   in this way.  Text within such a segment can be
   given \emph{additional} emphasis.
\end{em}

It is sometimes necessary to prevent \LaTeX\ from
breaking a line where it might otherwise do so.
This may be at a space, as between the ``Mr.''\ and
``Jones'' in
       ``Mr.~Jones'',        % ~ produces an unbreakable interword space.
or within a word---especially when the word is a
symbol like
       \mbox{\emph{itemnum}}
that makes little sense when hyphenated across
lines.

Footnotes\footnote{This is an example of a footnote.}
pose no problem.

\LaTeX\ is good at typesetting mathematical formulas
like
       \( x-3y + z = 7 \)
or
       \( a_{1} > x^{2n} + y^{2n} > x' \)
or
       \( \ip{A}{B} = \sum_{i} a_{i} b_{i} \).
The spaces you type in a formula are
ignored.  Remember that a letter like
       $x$                   % $ ... $  and  \( ... \)  are equivalent
is a formula when it denotes a mathematical
symbol, and it should be typed as one.

\section{Displayed Text}

Text is displayed by indenting it from the left
margin.  Quotations are commonly displayed.  There
are short quotations
\begin{quote}
   This is a short quotation.  It consists of a
   single paragraph of text.  See how it is formatted.
\end{quote}
and longer ones.
\begin{quotation}
   This is a longer quotation.  It consists of two
   paragraphs of text, neither of which are
   particularly interesting.

   This is the second paragraph of the quotation.  It
   is just as dull as the first paragraph.
\end{quotation}
Another frequently-displayed structure is a list.
The following is an example of an \emph{itemized}
list.
\begin{itemize}
   \item This is the first item of an itemized list.
         Each item in the list is marked with a ``tick''.
         You don't have to worry about what kind of tick
         mark is used.

   \item This is the second item of the list.  It
         contains another list nested inside it.  The inner
         list is an \emph{enumerated} list.
         \begin{enumerate}
            \item This is the first item of an enumerated
                  list that is nested within the itemized list.

            \item This is the second item of the inner list.
                  \LaTeX\ allows you to nest lists deeper than
                  you really should.
         \end{enumerate}
         This is the rest of the second item of the outer
         list.  It is no more interesting than any other
         part of the item.
   \item This is the third item of the list.
\end{itemize}
You can even display poetry.
\begin{verse}
   There is an environment
    for verse \\             % The \\ command separates lines
   Whose features some poets % within a stanza.
   will curse.

                             % One or more blank lines separate stanzas.

   For instead of making\\
   Them do \emph{all} line breaking, \\
   It allows them to put too many words on a line when they'd rather be
   forced to be terse.
\end{verse}

Mathematical formulas may also be displayed.  A
displayed formula
is
one-line long; multiline
formulas require special formatting instructions.
   \[  \ip{\Gamma}{\psi'} = x'' + y^{2} + z_{i}^{n}\]
Don't start a paragraph with a displayed equation,
nor make one a paragraph by itself.

\end{document}               % End of document.

\section{Bench Huge Section 660}
% synthetic bench section 660
% This is a sample LaTeX input file.  (Version of 12 August 2004.)
%
% A '%' character causes TeX to ignore all remaining text on the line,
% and is used for comments like this one.

\documentclass{article}      % Specifies the document class

                             % The preamble begins here.
\title{An Example Document}  % Declares the document's title.
\author{Leslie Lamport}      % Declares the author's name.
\date{January 21, 1994}      % Deleting this command produces today's date.

\newcommand{\ip}[2]{(#1, #2)}
                             % Defines \ip{arg1}{arg2} to mean
                             % (arg1, arg2).

%\newcommand{\ip}[2]{\langle #1 | #2\rangle}
                             % This is an alternative definition of
                             % \ip that is commented out.

\begin{document}             % End of preamble and beginning of text.

\maketitle                   % Produces the title.

This is an example input file.  Comparing it with
the output it generates can show you how to
produce a simple document of your own.

\section{Ordinary Text}      % Produces section heading.  Lower-level
                             % sections are begun with similar
                             % \subsection and \subsubsection commands.

The ends  of words and sentences are marked
  by   spaces. It  doesn't matter how many
spaces    you type; one is as good as 100.  The
end of   a line counts as a space.

One   or more   blank lines denote the  end
of  a paragraph.

Since any number of consecutive spaces are treated
like a single one, the formatting of the input
file makes no difference to
      \LaTeX,                % The \LaTeX command generates the LaTeX logo.
but it makes a difference to you.  When you use
\LaTeX, making your input file as easy to read
as possible will be a great help as you write
your document and when you change it.  This sample
file shows how you can add comments to your own input
file.

Because printing is different from typewriting,
there are a number of things that you have to do
differently when preparing an input file than if
you were just typing the document directly.
Quotation marks like
       ``this''
have to be handled specially, as do quotes within
quotes:
       ``\,`this'            % \, separates the double and single quote.
        is what I just
        wrote, not  `that'\,''.

Dashes come in three sizes: an
       intra-word
dash, a medium dash for number ranges like
       1--2,
and a punctuation
       dash---like
this.

A sentence-ending space should be larger than the
space between words within a sentence.  You
sometimes have to type special commands in
conjunction with punctuation characters to get
this right, as in the following sentence.
       Gnats, gnus, etc.\ all  % `\ ' makes an inter-word space.
       begin with G\@.         % \@ marks end-of-sentence punctuation.
You should check the spaces after periods when
reading your output to make sure you haven't
forgotten any special cases.  Generating an
ellipsis
       \ldots\               % `\ ' is needed after `\ldots' because TeX
                             % ignores spaces after command names like \ldots
                             % made from \ + letters.
                             %
                             % Note how a `%' character causes TeX to ignore
                             % the end of the input line, so these blank lines
                             % do not start a new paragraph.
                             %
with the right spacing around the periods requires
a special command.

\LaTeX\ interprets some common characters as
commands, so you must type special commands to
generate them.  These characters include the
following:
       \$ \& \% \# \{ and \}.

In printing, text is usually emphasized with an
       \emph{italic}
type style.

\begin{em}
   A long segment of text can also be emphasized
   in this way.  Text within such a segment can be
   given \emph{additional} emphasis.
\end{em}

It is sometimes necessary to prevent \LaTeX\ from
breaking a line where it might otherwise do so.
This may be at a space, as between the ``Mr.''\ and
``Jones'' in
       ``Mr.~Jones'',        % ~ produces an unbreakable interword space.
or within a word---especially when the word is a
symbol like
       \mbox{\emph{itemnum}}
that makes little sense when hyphenated across
lines.

Footnotes\footnote{This is an example of a footnote.}
pose no problem.

\LaTeX\ is good at typesetting mathematical formulas
like
       \( x-3y + z = 7 \)
or
       \( a_{1} > x^{2n} + y^{2n} > x' \)
or
       \( \ip{A}{B} = \sum_{i} a_{i} b_{i} \).
The spaces you type in a formula are
ignored.  Remember that a letter like
       $x$                   % $ ... $  and  \( ... \)  are equivalent
is a formula when it denotes a mathematical
symbol, and it should be typed as one.

\section{Displayed Text}

Text is displayed by indenting it from the left
margin.  Quotations are commonly displayed.  There
are short quotations
\begin{quote}
   This is a short quotation.  It consists of a
   single paragraph of text.  See how it is formatted.
\end{quote}
and longer ones.
\begin{quotation}
   This is a longer quotation.  It consists of two
   paragraphs of text, neither of which are
   particularly interesting.

   This is the second paragraph of the quotation.  It
   is just as dull as the first paragraph.
\end{quotation}
Another frequently-displayed structure is a list.
The following is an example of an \emph{itemized}
list.
\begin{itemize}
   \item This is the first item of an itemized list.
         Each item in the list is marked with a ``tick''.
         You don't have to worry about what kind of tick
         mark is used.

   \item This is the second item of the list.  It
         contains another list nested inside it.  The inner
         list is an \emph{enumerated} list.
         \begin{enumerate}
            \item This is the first item of an enumerated
                  list that is nested within the itemized list.

            \item This is the second item of the inner list.
                  \LaTeX\ allows you to nest lists deeper than
                  you really should.
         \end{enumerate}
         This is the rest of the second item of the outer
         list.  It is no more interesting than any other
         part of the item.
   \item This is the third item of the list.
\end{itemize}
You can even display poetry.
\begin{verse}
   There is an environment
    for verse \\             % The \\ command separates lines
   Whose features some poets % within a stanza.
   will curse.

                             % One or more blank lines separate stanzas.

   For instead of making\\
   Them do \emph{all} line breaking, \\
   It allows them to put too many words on a line when they'd rather be
   forced to be terse.
\end{verse}

Mathematical formulas may also be displayed.  A
displayed formula
is
one-line long; multiline
formulas require special formatting instructions.
   \[  \ip{\Gamma}{\psi'} = x'' + y^{2} + z_{i}^{n}\]
Don't start a paragraph with a displayed equation,
nor make one a paragraph by itself.

\end{document}               % End of document.

\section{Bench Huge Section 661}
% synthetic bench section 661
% This is a sample LaTeX input file.  (Version of 12 August 2004.)
%
% A '%' character causes TeX to ignore all remaining text on the line,
% and is used for comments like this one.

\documentclass{article}      % Specifies the document class

                             % The preamble begins here.
\title{An Example Document}  % Declares the document's title.
\author{Leslie Lamport}      % Declares the author's name.
\date{January 21, 1994}      % Deleting this command produces today's date.

\newcommand{\ip}[2]{(#1, #2)}
                             % Defines \ip{arg1}{arg2} to mean
                             % (arg1, arg2).

%\newcommand{\ip}[2]{\langle #1 | #2\rangle}
                             % This is an alternative definition of
                             % \ip that is commented out.

\begin{document}             % End of preamble and beginning of text.

\maketitle                   % Produces the title.

This is an example input file.  Comparing it with
the output it generates can show you how to
produce a simple document of your own.

\section{Ordinary Text}      % Produces section heading.  Lower-level
                             % sections are begun with similar
                             % \subsection and \subsubsection commands.

The ends  of words and sentences are marked
  by   spaces. It  doesn't matter how many
spaces    you type; one is as good as 100.  The
end of   a line counts as a space.

One   or more   blank lines denote the  end
of  a paragraph.

Since any number of consecutive spaces are treated
like a single one, the formatting of the input
file makes no difference to
      \LaTeX,                % The \LaTeX command generates the LaTeX logo.
but it makes a difference to you.  When you use
\LaTeX, making your input file as easy to read
as possible will be a great help as you write
your document and when you change it.  This sample
file shows how you can add comments to your own input
file.

Because printing is different from typewriting,
there are a number of things that you have to do
differently when preparing an input file than if
you were just typing the document directly.
Quotation marks like
       ``this''
have to be handled specially, as do quotes within
quotes:
       ``\,`this'            % \, separates the double and single quote.
        is what I just
        wrote, not  `that'\,''.

Dashes come in three sizes: an
       intra-word
dash, a medium dash for number ranges like
       1--2,
and a punctuation
       dash---like
this.

A sentence-ending space should be larger than the
space between words within a sentence.  You
sometimes have to type special commands in
conjunction with punctuation characters to get
this right, as in the following sentence.
       Gnats, gnus, etc.\ all  % `\ ' makes an inter-word space.
       begin with G\@.         % \@ marks end-of-sentence punctuation.
You should check the spaces after periods when
reading your output to make sure you haven't
forgotten any special cases.  Generating an
ellipsis
       \ldots\               % `\ ' is needed after `\ldots' because TeX
                             % ignores spaces after command names like \ldots
                             % made from \ + letters.
                             %
                             % Note how a `%' character causes TeX to ignore
                             % the end of the input line, so these blank lines
                             % do not start a new paragraph.
                             %
with the right spacing around the periods requires
a special command.

\LaTeX\ interprets some common characters as
commands, so you must type special commands to
generate them.  These characters include the
following:
       \$ \& \% \# \{ and \}.

In printing, text is usually emphasized with an
       \emph{italic}
type style.

\begin{em}
   A long segment of text can also be emphasized
   in this way.  Text within such a segment can be
   given \emph{additional} emphasis.
\end{em}

It is sometimes necessary to prevent \LaTeX\ from
breaking a line where it might otherwise do so.
This may be at a space, as between the ``Mr.''\ and
``Jones'' in
       ``Mr.~Jones'',        % ~ produces an unbreakable interword space.
or within a word---especially when the word is a
symbol like
       \mbox{\emph{itemnum}}
that makes little sense when hyphenated across
lines.

Footnotes\footnote{This is an example of a footnote.}
pose no problem.

\LaTeX\ is good at typesetting mathematical formulas
like
       \( x-3y + z = 7 \)
or
       \( a_{1} > x^{2n} + y^{2n} > x' \)
or
       \( \ip{A}{B} = \sum_{i} a_{i} b_{i} \).
The spaces you type in a formula are
ignored.  Remember that a letter like
       $x$                   % $ ... $  and  \( ... \)  are equivalent
is a formula when it denotes a mathematical
symbol, and it should be typed as one.

\section{Displayed Text}

Text is displayed by indenting it from the left
margin.  Quotations are commonly displayed.  There
are short quotations
\begin{quote}
   This is a short quotation.  It consists of a
   single paragraph of text.  See how it is formatted.
\end{quote}
and longer ones.
\begin{quotation}
   This is a longer quotation.  It consists of two
   paragraphs of text, neither of which are
   particularly interesting.

   This is the second paragraph of the quotation.  It
   is just as dull as the first paragraph.
\end{quotation}
Another frequently-displayed structure is a list.
The following is an example of an \emph{itemized}
list.
\begin{itemize}
   \item This is the first item of an itemized list.
         Each item in the list is marked with a ``tick''.
         You don't have to worry about what kind of tick
         mark is used.

   \item This is the second item of the list.  It
         contains another list nested inside it.  The inner
         list is an \emph{enumerated} list.
         \begin{enumerate}
            \item This is the first item of an enumerated
                  list that is nested within the itemized list.

            \item This is the second item of the inner list.
                  \LaTeX\ allows you to nest lists deeper than
                  you really should.
         \end{enumerate}
         This is the rest of the second item of the outer
         list.  It is no more interesting than any other
         part of the item.
   \item This is the third item of the list.
\end{itemize}
You can even display poetry.
\begin{verse}
   There is an environment
    for verse \\             % The \\ command separates lines
   Whose features some poets % within a stanza.
   will curse.

                             % One or more blank lines separate stanzas.

   For instead of making\\
   Them do \emph{all} line breaking, \\
   It allows them to put too many words on a line when they'd rather be
   forced to be terse.
\end{verse}

Mathematical formulas may also be displayed.  A
displayed formula
is
one-line long; multiline
formulas require special formatting instructions.
   \[  \ip{\Gamma}{\psi'} = x'' + y^{2} + z_{i}^{n}\]
Don't start a paragraph with a displayed equation,
nor make one a paragraph by itself.

\end{document}               % End of document.

\section{Bench Huge Section 662}
% synthetic bench section 662
% This is a sample LaTeX input file.  (Version of 12 August 2004.)
%
% A '%' character causes TeX to ignore all remaining text on the line,
% and is used for comments like this one.

\documentclass{article}      % Specifies the document class

                             % The preamble begins here.
\title{An Example Document}  % Declares the document's title.
\author{Leslie Lamport}      % Declares the author's name.
\date{January 21, 1994}      % Deleting this command produces today's date.

\newcommand{\ip}[2]{(#1, #2)}
                             % Defines \ip{arg1}{arg2} to mean
                             % (arg1, arg2).

%\newcommand{\ip}[2]{\langle #1 | #2\rangle}
                             % This is an alternative definition of
                             % \ip that is commented out.

\begin{document}             % End of preamble and beginning of text.

\maketitle                   % Produces the title.

This is an example input file.  Comparing it with
the output it generates can show you how to
produce a simple document of your own.

\section{Ordinary Text}      % Produces section heading.  Lower-level
                             % sections are begun with similar
                             % \subsection and \subsubsection commands.

The ends  of words and sentences are marked
  by   spaces. It  doesn't matter how many
spaces    you type; one is as good as 100.  The
end of   a line counts as a space.

One   or more   blank lines denote the  end
of  a paragraph.

Since any number of consecutive spaces are treated
like a single one, the formatting of the input
file makes no difference to
      \LaTeX,                % The \LaTeX command generates the LaTeX logo.
but it makes a difference to you.  When you use
\LaTeX, making your input file as easy to read
as possible will be a great help as you write
your document and when you change it.  This sample
file shows how you can add comments to your own input
file.

Because printing is different from typewriting,
there are a number of things that you have to do
differently when preparing an input file than if
you were just typing the document directly.
Quotation marks like
       ``this''
have to be handled specially, as do quotes within
quotes:
       ``\,`this'            % \, separates the double and single quote.
        is what I just
        wrote, not  `that'\,''.

Dashes come in three sizes: an
       intra-word
dash, a medium dash for number ranges like
       1--2,
and a punctuation
       dash---like
this.

A sentence-ending space should be larger than the
space between words within a sentence.  You
sometimes have to type special commands in
conjunction with punctuation characters to get
this right, as in the following sentence.
       Gnats, gnus, etc.\ all  % `\ ' makes an inter-word space.
       begin with G\@.         % \@ marks end-of-sentence punctuation.
You should check the spaces after periods when
reading your output to make sure you haven't
forgotten any special cases.  Generating an
ellipsis
       \ldots\               % `\ ' is needed after `\ldots' because TeX
                             % ignores spaces after command names like \ldots
                             % made from \ + letters.
                             %
                             % Note how a `%' character causes TeX to ignore
                             % the end of the input line, so these blank lines
                             % do not start a new paragraph.
                             %
with the right spacing around the periods requires
a special command.

\LaTeX\ interprets some common characters as
commands, so you must type special commands to
generate them.  These characters include the
following:
       \$ \& \% \# \{ and \}.

In printing, text is usually emphasized with an
       \emph{italic}
type style.

\begin{em}
   A long segment of text can also be emphasized
   in this way.  Text within such a segment can be
   given \emph{additional} emphasis.
\end{em}

It is sometimes necessary to prevent \LaTeX\ from
breaking a line where it might otherwise do so.
This may be at a space, as between the ``Mr.''\ and
``Jones'' in
       ``Mr.~Jones'',        % ~ produces an unbreakable interword space.
or within a word---especially when the word is a
symbol like
       \mbox{\emph{itemnum}}
that makes little sense when hyphenated across
lines.

Footnotes\footnote{This is an example of a footnote.}
pose no problem.

\LaTeX\ is good at typesetting mathematical formulas
like
       \( x-3y + z = 7 \)
or
       \( a_{1} > x^{2n} + y^{2n} > x' \)
or
       \( \ip{A}{B} = \sum_{i} a_{i} b_{i} \).
The spaces you type in a formula are
ignored.  Remember that a letter like
       $x$                   % $ ... $  and  \( ... \)  are equivalent
is a formula when it denotes a mathematical
symbol, and it should be typed as one.

\section{Displayed Text}

Text is displayed by indenting it from the left
margin.  Quotations are commonly displayed.  There
are short quotations
\begin{quote}
   This is a short quotation.  It consists of a
   single paragraph of text.  See how it is formatted.
\end{quote}
and longer ones.
\begin{quotation}
   This is a longer quotation.  It consists of two
   paragraphs of text, neither of which are
   particularly interesting.

   This is the second paragraph of the quotation.  It
   is just as dull as the first paragraph.
\end{quotation}
Another frequently-displayed structure is a list.
The following is an example of an \emph{itemized}
list.
\begin{itemize}
   \item This is the first item of an itemized list.
         Each item in the list is marked with a ``tick''.
         You don't have to worry about what kind of tick
         mark is used.

   \item This is the second item of the list.  It
         contains another list nested inside it.  The inner
         list is an \emph{enumerated} list.
         \begin{enumerate}
            \item This is the first item of an enumerated
                  list that is nested within the itemized list.

            \item This is the second item of the inner list.
                  \LaTeX\ allows you to nest lists deeper than
                  you really should.
         \end{enumerate}
         This is the rest of the second item of the outer
         list.  It is no more interesting than any other
         part of the item.
   \item This is the third item of the list.
\end{itemize}
You can even display poetry.
\begin{verse}
   There is an environment
    for verse \\             % The \\ command separates lines
   Whose features some poets % within a stanza.
   will curse.

                             % One or more blank lines separate stanzas.

   For instead of making\\
   Them do \emph{all} line breaking, \\
   It allows them to put too many words on a line when they'd rather be
   forced to be terse.
\end{verse}

Mathematical formulas may also be displayed.  A
displayed formula
is
one-line long; multiline
formulas require special formatting instructions.
   \[  \ip{\Gamma}{\psi'} = x'' + y^{2} + z_{i}^{n}\]
Don't start a paragraph with a displayed equation,
nor make one a paragraph by itself.

\end{document}               % End of document.

\section{Bench Huge Section 663}
% synthetic bench section 663
% This is a sample LaTeX input file.  (Version of 12 August 2004.)
%
% A '%' character causes TeX to ignore all remaining text on the line,
% and is used for comments like this one.

\documentclass{article}      % Specifies the document class

                             % The preamble begins here.
\title{An Example Document}  % Declares the document's title.
\author{Leslie Lamport}      % Declares the author's name.
\date{January 21, 1994}      % Deleting this command produces today's date.

\newcommand{\ip}[2]{(#1, #2)}
                             % Defines \ip{arg1}{arg2} to mean
                             % (arg1, arg2).

%\newcommand{\ip}[2]{\langle #1 | #2\rangle}
                             % This is an alternative definition of
                             % \ip that is commented out.

\begin{document}             % End of preamble and beginning of text.

\maketitle                   % Produces the title.

This is an example input file.  Comparing it with
the output it generates can show you how to
produce a simple document of your own.

\section{Ordinary Text}      % Produces section heading.  Lower-level
                             % sections are begun with similar
                             % \subsection and \subsubsection commands.

The ends  of words and sentences are marked
  by   spaces. It  doesn't matter how many
spaces    you type; one is as good as 100.  The
end of   a line counts as a space.

One   or more   blank lines denote the  end
of  a paragraph.

Since any number of consecutive spaces are treated
like a single one, the formatting of the input
file makes no difference to
      \LaTeX,                % The \LaTeX command generates the LaTeX logo.
but it makes a difference to you.  When you use
\LaTeX, making your input file as easy to read
as possible will be a great help as you write
your document and when you change it.  This sample
file shows how you can add comments to your own input
file.

Because printing is different from typewriting,
there are a number of things that you have to do
differently when preparing an input file than if
you were just typing the document directly.
Quotation marks like
       ``this''
have to be handled specially, as do quotes within
quotes:
       ``\,`this'            % \, separates the double and single quote.
        is what I just
        wrote, not  `that'\,''.

Dashes come in three sizes: an
       intra-word
dash, a medium dash for number ranges like
       1--2,
and a punctuation
       dash---like
this.

A sentence-ending space should be larger than the
space between words within a sentence.  You
sometimes have to type special commands in
conjunction with punctuation characters to get
this right, as in the following sentence.
       Gnats, gnus, etc.\ all  % `\ ' makes an inter-word space.
       begin with G\@.         % \@ marks end-of-sentence punctuation.
You should check the spaces after periods when
reading your output to make sure you haven't
forgotten any special cases.  Generating an
ellipsis
       \ldots\               % `\ ' is needed after `\ldots' because TeX
                             % ignores spaces after command names like \ldots
                             % made from \ + letters.
                             %
                             % Note how a `%' character causes TeX to ignore
                             % the end of the input line, so these blank lines
                             % do not start a new paragraph.
                             %
with the right spacing around the periods requires
a special command.

\LaTeX\ interprets some common characters as
commands, so you must type special commands to
generate them.  These characters include the
following:
       \$ \& \% \# \{ and \}.

In printing, text is usually emphasized with an
       \emph{italic}
type style.

\begin{em}
   A long segment of text can also be emphasized
   in this way.  Text within such a segment can be
   given \emph{additional} emphasis.
\end{em}

It is sometimes necessary to prevent \LaTeX\ from
breaking a line where it might otherwise do so.
This may be at a space, as between the ``Mr.''\ and
``Jones'' in
       ``Mr.~Jones'',        % ~ produces an unbreakable interword space.
or within a word---especially when the word is a
symbol like
       \mbox{\emph{itemnum}}
that makes little sense when hyphenated across
lines.

Footnotes\footnote{This is an example of a footnote.}
pose no problem.

\LaTeX\ is good at typesetting mathematical formulas
like
       \( x-3y + z = 7 \)
or
       \( a_{1} > x^{2n} + y^{2n} > x' \)
or
       \( \ip{A}{B} = \sum_{i} a_{i} b_{i} \).
The spaces you type in a formula are
ignored.  Remember that a letter like
       $x$                   % $ ... $  and  \( ... \)  are equivalent
is a formula when it denotes a mathematical
symbol, and it should be typed as one.

\section{Displayed Text}

Text is displayed by indenting it from the left
margin.  Quotations are commonly displayed.  There
are short quotations
\begin{quote}
   This is a short quotation.  It consists of a
   single paragraph of text.  See how it is formatted.
\end{quote}
and longer ones.
\begin{quotation}
   This is a longer quotation.  It consists of two
   paragraphs of text, neither of which are
   particularly interesting.

   This is the second paragraph of the quotation.  It
   is just as dull as the first paragraph.
\end{quotation}
Another frequently-displayed structure is a list.
The following is an example of an \emph{itemized}
list.
\begin{itemize}
   \item This is the first item of an itemized list.
         Each item in the list is marked with a ``tick''.
         You don't have to worry about what kind of tick
         mark is used.

   \item This is the second item of the list.  It
         contains another list nested inside it.  The inner
         list is an \emph{enumerated} list.
         \begin{enumerate}
            \item This is the first item of an enumerated
                  list that is nested within the itemized list.

            \item This is the second item of the inner list.
                  \LaTeX\ allows you to nest lists deeper than
                  you really should.
         \end{enumerate}
         This is the rest of the second item of the outer
         list.  It is no more interesting than any other
         part of the item.
   \item This is the third item of the list.
\end{itemize}
You can even display poetry.
\begin{verse}
   There is an environment
    for verse \\             % The \\ command separates lines
   Whose features some poets % within a stanza.
   will curse.

                             % One or more blank lines separate stanzas.

   For instead of making\\
   Them do \emph{all} line breaking, \\
   It allows them to put too many words on a line when they'd rather be
   forced to be terse.
\end{verse}

Mathematical formulas may also be displayed.  A
displayed formula
is
one-line long; multiline
formulas require special formatting instructions.
   \[  \ip{\Gamma}{\psi'} = x'' + y^{2} + z_{i}^{n}\]
Don't start a paragraph with a displayed equation,
nor make one a paragraph by itself.

\end{document}               % End of document.

\section{Bench Huge Section 664}
% synthetic bench section 664
% This is a sample LaTeX input file.  (Version of 12 August 2004.)
%
% A '%' character causes TeX to ignore all remaining text on the line,
% and is used for comments like this one.

\documentclass{article}      % Specifies the document class

                             % The preamble begins here.
\title{An Example Document}  % Declares the document's title.
\author{Leslie Lamport}      % Declares the author's name.
\date{January 21, 1994}      % Deleting this command produces today's date.

\newcommand{\ip}[2]{(#1, #2)}
                             % Defines \ip{arg1}{arg2} to mean
                             % (arg1, arg2).

%\newcommand{\ip}[2]{\langle #1 | #2\rangle}
                             % This is an alternative definition of
                             % \ip that is commented out.

\begin{document}             % End of preamble and beginning of text.

\maketitle                   % Produces the title.

This is an example input file.  Comparing it with
the output it generates can show you how to
produce a simple document of your own.

\section{Ordinary Text}      % Produces section heading.  Lower-level
                             % sections are begun with similar
                             % \subsection and \subsubsection commands.

The ends  of words and sentences are marked
  by   spaces. It  doesn't matter how many
spaces    you type; one is as good as 100.  The
end of   a line counts as a space.

One   or more   blank lines denote the  end
of  a paragraph.

Since any number of consecutive spaces are treated
like a single one, the formatting of the input
file makes no difference to
      \LaTeX,                % The \LaTeX command generates the LaTeX logo.
but it makes a difference to you.  When you use
\LaTeX, making your input file as easy to read
as possible will be a great help as you write
your document and when you change it.  This sample
file shows how you can add comments to your own input
file.

Because printing is different from typewriting,
there are a number of things that you have to do
differently when preparing an input file than if
you were just typing the document directly.
Quotation marks like
       ``this''
have to be handled specially, as do quotes within
quotes:
       ``\,`this'            % \, separates the double and single quote.
        is what I just
        wrote, not  `that'\,''.

Dashes come in three sizes: an
       intra-word
dash, a medium dash for number ranges like
       1--2,
and a punctuation
       dash---like
this.

A sentence-ending space should be larger than the
space between words within a sentence.  You
sometimes have to type special commands in
conjunction with punctuation characters to get
this right, as in the following sentence.
       Gnats, gnus, etc.\ all  % `\ ' makes an inter-word space.
       begin with G\@.         % \@ marks end-of-sentence punctuation.
You should check the spaces after periods when
reading your output to make sure you haven't
forgotten any special cases.  Generating an
ellipsis
       \ldots\               % `\ ' is needed after `\ldots' because TeX
                             % ignores spaces after command names like \ldots
                             % made from \ + letters.
                             %
                             % Note how a `%' character causes TeX to ignore
                             % the end of the input line, so these blank lines
                             % do not start a new paragraph.
                             %
with the right spacing around the periods requires
a special command.

\LaTeX\ interprets some common characters as
commands, so you must type special commands to
generate them.  These characters include the
following:
       \$ \& \% \# \{ and \}.

In printing, text is usually emphasized with an
       \emph{italic}
type style.

\begin{em}
   A long segment of text can also be emphasized
   in this way.  Text within such a segment can be
   given \emph{additional} emphasis.
\end{em}

It is sometimes necessary to prevent \LaTeX\ from
breaking a line where it might otherwise do so.
This may be at a space, as between the ``Mr.''\ and
``Jones'' in
       ``Mr.~Jones'',        % ~ produces an unbreakable interword space.
or within a word---especially when the word is a
symbol like
       \mbox{\emph{itemnum}}
that makes little sense when hyphenated across
lines.

Footnotes\footnote{This is an example of a footnote.}
pose no problem.

\LaTeX\ is good at typesetting mathematical formulas
like
       \( x-3y + z = 7 \)
or
       \( a_{1} > x^{2n} + y^{2n} > x' \)
or
       \( \ip{A}{B} = \sum_{i} a_{i} b_{i} \).
The spaces you type in a formula are
ignored.  Remember that a letter like
       $x$                   % $ ... $  and  \( ... \)  are equivalent
is a formula when it denotes a mathematical
symbol, and it should be typed as one.

\section{Displayed Text}

Text is displayed by indenting it from the left
margin.  Quotations are commonly displayed.  There
are short quotations
\begin{quote}
   This is a short quotation.  It consists of a
   single paragraph of text.  See how it is formatted.
\end{quote}
and longer ones.
\begin{quotation}
   This is a longer quotation.  It consists of two
   paragraphs of text, neither of which are
   particularly interesting.

   This is the second paragraph of the quotation.  It
   is just as dull as the first paragraph.
\end{quotation}
Another frequently-displayed structure is a list.
The following is an example of an \emph{itemized}
list.
\begin{itemize}
   \item This is the first item of an itemized list.
         Each item in the list is marked with a ``tick''.
         You don't have to worry about what kind of tick
         mark is used.

   \item This is the second item of the list.  It
         contains another list nested inside it.  The inner
         list is an \emph{enumerated} list.
         \begin{enumerate}
            \item This is the first item of an enumerated
                  list that is nested within the itemized list.

            \item This is the second item of the inner list.
                  \LaTeX\ allows you to nest lists deeper than
                  you really should.
         \end{enumerate}
         This is the rest of the second item of the outer
         list.  It is no more interesting than any other
         part of the item.
   \item This is the third item of the list.
\end{itemize}
You can even display poetry.
\begin{verse}
   There is an environment
    for verse \\             % The \\ command separates lines
   Whose features some poets % within a stanza.
   will curse.

                             % One or more blank lines separate stanzas.

   For instead of making\\
   Them do \emph{all} line breaking, \\
   It allows them to put too many words on a line when they'd rather be
   forced to be terse.
\end{verse}

Mathematical formulas may also be displayed.  A
displayed formula
is
one-line long; multiline
formulas require special formatting instructions.
   \[  \ip{\Gamma}{\psi'} = x'' + y^{2} + z_{i}^{n}\]
Don't start a paragraph with a displayed equation,
nor make one a paragraph by itself.

\end{document}               % End of document.

\section{Bench Huge Section 665}
% synthetic bench section 665
% This is a sample LaTeX input file.  (Version of 12 August 2004.)
%
% A '%' character causes TeX to ignore all remaining text on the line,
% and is used for comments like this one.

\documentclass{article}      % Specifies the document class

                             % The preamble begins here.
\title{An Example Document}  % Declares the document's title.
\author{Leslie Lamport}      % Declares the author's name.
\date{January 21, 1994}      % Deleting this command produces today's date.

\newcommand{\ip}[2]{(#1, #2)}
                             % Defines \ip{arg1}{arg2} to mean
                             % (arg1, arg2).

%\newcommand{\ip}[2]{\langle #1 | #2\rangle}
                             % This is an alternative definition of
                             % \ip that is commented out.

\begin{document}             % End of preamble and beginning of text.

\maketitle                   % Produces the title.

This is an example input file.  Comparing it with
the output it generates can show you how to
produce a simple document of your own.

\section{Ordinary Text}      % Produces section heading.  Lower-level
                             % sections are begun with similar
                             % \subsection and \subsubsection commands.

The ends  of words and sentences are marked
  by   spaces. It  doesn't matter how many
spaces    you type; one is as good as 100.  The
end of   a line counts as a space.

One   or more   blank lines denote the  end
of  a paragraph.

Since any number of consecutive spaces are treated
like a single one, the formatting of the input
file makes no difference to
      \LaTeX,                % The \LaTeX command generates the LaTeX logo.
but it makes a difference to you.  When you use
\LaTeX, making your input file as easy to read
as possible will be a great help as you write
your document and when you change it.  This sample
file shows how you can add comments to your own input
file.

Because printing is different from typewriting,
there are a number of things that you have to do
differently when preparing an input file than if
you were just typing the document directly.
Quotation marks like
       ``this''
have to be handled specially, as do quotes within
quotes:
       ``\,`this'            % \, separates the double and single quote.
        is what I just
        wrote, not  `that'\,''.

Dashes come in three sizes: an
       intra-word
dash, a medium dash for number ranges like
       1--2,
and a punctuation
       dash---like
this.

A sentence-ending space should be larger than the
space between words within a sentence.  You
sometimes have to type special commands in
conjunction with punctuation characters to get
this right, as in the following sentence.
       Gnats, gnus, etc.\ all  % `\ ' makes an inter-word space.
       begin with G\@.         % \@ marks end-of-sentence punctuation.
You should check the spaces after periods when
reading your output to make sure you haven't
forgotten any special cases.  Generating an
ellipsis
       \ldots\               % `\ ' is needed after `\ldots' because TeX
                             % ignores spaces after command names like \ldots
                             % made from \ + letters.
                             %
                             % Note how a `%' character causes TeX to ignore
                             % the end of the input line, so these blank lines
                             % do not start a new paragraph.
                             %
with the right spacing around the periods requires
a special command.

\LaTeX\ interprets some common characters as
commands, so you must type special commands to
generate them.  These characters include the
following:
       \$ \& \% \# \{ and \}.

In printing, text is usually emphasized with an
       \emph{italic}
type style.

\begin{em}
   A long segment of text can also be emphasized
   in this way.  Text within such a segment can be
   given \emph{additional} emphasis.
\end{em}

It is sometimes necessary to prevent \LaTeX\ from
breaking a line where it might otherwise do so.
This may be at a space, as between the ``Mr.''\ and
``Jones'' in
       ``Mr.~Jones'',        % ~ produces an unbreakable interword space.
or within a word---especially when the word is a
symbol like
       \mbox{\emph{itemnum}}
that makes little sense when hyphenated across
lines.

Footnotes\footnote{This is an example of a footnote.}
pose no problem.

\LaTeX\ is good at typesetting mathematical formulas
like
       \( x-3y + z = 7 \)
or
       \( a_{1} > x^{2n} + y^{2n} > x' \)
or
       \( \ip{A}{B} = \sum_{i} a_{i} b_{i} \).
The spaces you type in a formula are
ignored.  Remember that a letter like
       $x$                   % $ ... $  and  \( ... \)  are equivalent
is a formula when it denotes a mathematical
symbol, and it should be typed as one.

\section{Displayed Text}

Text is displayed by indenting it from the left
margin.  Quotations are commonly displayed.  There
are short quotations
\begin{quote}
   This is a short quotation.  It consists of a
   single paragraph of text.  See how it is formatted.
\end{quote}
and longer ones.
\begin{quotation}
   This is a longer quotation.  It consists of two
   paragraphs of text, neither of which are
   particularly interesting.

   This is the second paragraph of the quotation.  It
   is just as dull as the first paragraph.
\end{quotation}
Another frequently-displayed structure is a list.
The following is an example of an \emph{itemized}
list.
\begin{itemize}
   \item This is the first item of an itemized list.
         Each item in the list is marked with a ``tick''.
         You don't have to worry about what kind of tick
         mark is used.

   \item This is the second item of the list.  It
         contains another list nested inside it.  The inner
         list is an \emph{enumerated} list.
         \begin{enumerate}
            \item This is the first item of an enumerated
                  list that is nested within the itemized list.

            \item This is the second item of the inner list.
                  \LaTeX\ allows you to nest lists deeper than
                  you really should.
         \end{enumerate}
         This is the rest of the second item of the outer
         list.  It is no more interesting than any other
         part of the item.
   \item This is the third item of the list.
\end{itemize}
You can even display poetry.
\begin{verse}
   There is an environment
    for verse \\             % The \\ command separates lines
   Whose features some poets % within a stanza.
   will curse.

                             % One or more blank lines separate stanzas.

   For instead of making\\
   Them do \emph{all} line breaking, \\
   It allows them to put too many words on a line when they'd rather be
   forced to be terse.
\end{verse}

Mathematical formulas may also be displayed.  A
displayed formula
is
one-line long; multiline
formulas require special formatting instructions.
   \[  \ip{\Gamma}{\psi'} = x'' + y^{2} + z_{i}^{n}\]
Don't start a paragraph with a displayed equation,
nor make one a paragraph by itself.

\end{document}               % End of document.

\section{Bench Huge Section 666}
% synthetic bench section 666
% This is a sample LaTeX input file.  (Version of 12 August 2004.)
%
% A '%' character causes TeX to ignore all remaining text on the line,
% and is used for comments like this one.

\documentclass{article}      % Specifies the document class

                             % The preamble begins here.
\title{An Example Document}  % Declares the document's title.
\author{Leslie Lamport}      % Declares the author's name.
\date{January 21, 1994}      % Deleting this command produces today's date.

\newcommand{\ip}[2]{(#1, #2)}
                             % Defines \ip{arg1}{arg2} to mean
                             % (arg1, arg2).

%\newcommand{\ip}[2]{\langle #1 | #2\rangle}
                             % This is an alternative definition of
                             % \ip that is commented out.

\begin{document}             % End of preamble and beginning of text.

\maketitle                   % Produces the title.

This is an example input file.  Comparing it with
the output it generates can show you how to
produce a simple document of your own.

\section{Ordinary Text}      % Produces section heading.  Lower-level
                             % sections are begun with similar
                             % \subsection and \subsubsection commands.

The ends  of words and sentences are marked
  by   spaces. It  doesn't matter how many
spaces    you type; one is as good as 100.  The
end of   a line counts as a space.

One   or more   blank lines denote the  end
of  a paragraph.

Since any number of consecutive spaces are treated
like a single one, the formatting of the input
file makes no difference to
      \LaTeX,                % The \LaTeX command generates the LaTeX logo.
but it makes a difference to you.  When you use
\LaTeX, making your input file as easy to read
as possible will be a great help as you write
your document and when you change it.  This sample
file shows how you can add comments to your own input
file.

Because printing is different from typewriting,
there are a number of things that you have to do
differently when preparing an input file than if
you were just typing the document directly.
Quotation marks like
       ``this''
have to be handled specially, as do quotes within
quotes:
       ``\,`this'            % \, separates the double and single quote.
        is what I just
        wrote, not  `that'\,''.

Dashes come in three sizes: an
       intra-word
dash, a medium dash for number ranges like
       1--2,
and a punctuation
       dash---like
this.

A sentence-ending space should be larger than the
space between words within a sentence.  You
sometimes have to type special commands in
conjunction with punctuation characters to get
this right, as in the following sentence.
       Gnats, gnus, etc.\ all  % `\ ' makes an inter-word space.
       begin with G\@.         % \@ marks end-of-sentence punctuation.
You should check the spaces after periods when
reading your output to make sure you haven't
forgotten any special cases.  Generating an
ellipsis
       \ldots\               % `\ ' is needed after `\ldots' because TeX
                             % ignores spaces after command names like \ldots
                             % made from \ + letters.
                             %
                             % Note how a `%' character causes TeX to ignore
                             % the end of the input line, so these blank lines
                             % do not start a new paragraph.
                             %
with the right spacing around the periods requires
a special command.

\LaTeX\ interprets some common characters as
commands, so you must type special commands to
generate them.  These characters include the
following:
       \$ \& \% \# \{ and \}.

In printing, text is usually emphasized with an
       \emph{italic}
type style.

\begin{em}
   A long segment of text can also be emphasized
   in this way.  Text within such a segment can be
   given \emph{additional} emphasis.
\end{em}

It is sometimes necessary to prevent \LaTeX\ from
breaking a line where it might otherwise do so.
This may be at a space, as between the ``Mr.''\ and
``Jones'' in
       ``Mr.~Jones'',        % ~ produces an unbreakable interword space.
or within a word---especially when the word is a
symbol like
       \mbox{\emph{itemnum}}
that makes little sense when hyphenated across
lines.

Footnotes\footnote{This is an example of a footnote.}
pose no problem.

\LaTeX\ is good at typesetting mathematical formulas
like
       \( x-3y + z = 7 \)
or
       \( a_{1} > x^{2n} + y^{2n} > x' \)
or
       \( \ip{A}{B} = \sum_{i} a_{i} b_{i} \).
The spaces you type in a formula are
ignored.  Remember that a letter like
       $x$                   % $ ... $  and  \( ... \)  are equivalent
is a formula when it denotes a mathematical
symbol, and it should be typed as one.

\section{Displayed Text}

Text is displayed by indenting it from the left
margin.  Quotations are commonly displayed.  There
are short quotations
\begin{quote}
   This is a short quotation.  It consists of a
   single paragraph of text.  See how it is formatted.
\end{quote}
and longer ones.
\begin{quotation}
   This is a longer quotation.  It consists of two
   paragraphs of text, neither of which are
   particularly interesting.

   This is the second paragraph of the quotation.  It
   is just as dull as the first paragraph.
\end{quotation}
Another frequently-displayed structure is a list.
The following is an example of an \emph{itemized}
list.
\begin{itemize}
   \item This is the first item of an itemized list.
         Each item in the list is marked with a ``tick''.
         You don't have to worry about what kind of tick
         mark is used.

   \item This is the second item of the list.  It
         contains another list nested inside it.  The inner
         list is an \emph{enumerated} list.
         \begin{enumerate}
            \item This is the first item of an enumerated
                  list that is nested within the itemized list.

            \item This is the second item of the inner list.
                  \LaTeX\ allows you to nest lists deeper than
                  you really should.
         \end{enumerate}
         This is the rest of the second item of the outer
         list.  It is no more interesting than any other
         part of the item.
   \item This is the third item of the list.
\end{itemize}
You can even display poetry.
\begin{verse}
   There is an environment
    for verse \\             % The \\ command separates lines
   Whose features some poets % within a stanza.
   will curse.

                             % One or more blank lines separate stanzas.

   For instead of making\\
   Them do \emph{all} line breaking, \\
   It allows them to put too many words on a line when they'd rather be
   forced to be terse.
\end{verse}

Mathematical formulas may also be displayed.  A
displayed formula
is
one-line long; multiline
formulas require special formatting instructions.
   \[  \ip{\Gamma}{\psi'} = x'' + y^{2} + z_{i}^{n}\]
Don't start a paragraph with a displayed equation,
nor make one a paragraph by itself.

\end{document}               % End of document.

\section{Bench Huge Section 667}
% synthetic bench section 667
% This is a sample LaTeX input file.  (Version of 12 August 2004.)
%
% A '%' character causes TeX to ignore all remaining text on the line,
% and is used for comments like this one.

\documentclass{article}      % Specifies the document class

                             % The preamble begins here.
\title{An Example Document}  % Declares the document's title.
\author{Leslie Lamport}      % Declares the author's name.
\date{January 21, 1994}      % Deleting this command produces today's date.

\newcommand{\ip}[2]{(#1, #2)}
                             % Defines \ip{arg1}{arg2} to mean
                             % (arg1, arg2).

%\newcommand{\ip}[2]{\langle #1 | #2\rangle}
                             % This is an alternative definition of
                             % \ip that is commented out.

\begin{document}             % End of preamble and beginning of text.

\maketitle                   % Produces the title.

This is an example input file.  Comparing it with
the output it generates can show you how to
produce a simple document of your own.

\section{Ordinary Text}      % Produces section heading.  Lower-level
                             % sections are begun with similar
                             % \subsection and \subsubsection commands.

The ends  of words and sentences are marked
  by   spaces. It  doesn't matter how many
spaces    you type; one is as good as 100.  The
end of   a line counts as a space.

One   or more   blank lines denote the  end
of  a paragraph.

Since any number of consecutive spaces are treated
like a single one, the formatting of the input
file makes no difference to
      \LaTeX,                % The \LaTeX command generates the LaTeX logo.
but it makes a difference to you.  When you use
\LaTeX, making your input file as easy to read
as possible will be a great help as you write
your document and when you change it.  This sample
file shows how you can add comments to your own input
file.

Because printing is different from typewriting,
there are a number of things that you have to do
differently when preparing an input file than if
you were just typing the document directly.
Quotation marks like
       ``this''
have to be handled specially, as do quotes within
quotes:
       ``\,`this'            % \, separates the double and single quote.
        is what I just
        wrote, not  `that'\,''.

Dashes come in three sizes: an
       intra-word
dash, a medium dash for number ranges like
       1--2,
and a punctuation
       dash---like
this.

A sentence-ending space should be larger than the
space between words within a sentence.  You
sometimes have to type special commands in
conjunction with punctuation characters to get
this right, as in the following sentence.
       Gnats, gnus, etc.\ all  % `\ ' makes an inter-word space.
       begin with G\@.         % \@ marks end-of-sentence punctuation.
You should check the spaces after periods when
reading your output to make sure you haven't
forgotten any special cases.  Generating an
ellipsis
       \ldots\               % `\ ' is needed after `\ldots' because TeX
                             % ignores spaces after command names like \ldots
                             % made from \ + letters.
                             %
                             % Note how a `%' character causes TeX to ignore
                             % the end of the input line, so these blank lines
                             % do not start a new paragraph.
                             %
with the right spacing around the periods requires
a special command.

\LaTeX\ interprets some common characters as
commands, so you must type special commands to
generate them.  These characters include the
following:
       \$ \& \% \# \{ and \}.

In printing, text is usually emphasized with an
       \emph{italic}
type style.

\begin{em}
   A long segment of text can also be emphasized
   in this way.  Text within such a segment can be
   given \emph{additional} emphasis.
\end{em}

It is sometimes necessary to prevent \LaTeX\ from
breaking a line where it might otherwise do so.
This may be at a space, as between the ``Mr.''\ and
``Jones'' in
       ``Mr.~Jones'',        % ~ produces an unbreakable interword space.
or within a word---especially when the word is a
symbol like
       \mbox{\emph{itemnum}}
that makes little sense when hyphenated across
lines.

Footnotes\footnote{This is an example of a footnote.}
pose no problem.

\LaTeX\ is good at typesetting mathematical formulas
like
       \( x-3y + z = 7 \)
or
       \( a_{1} > x^{2n} + y^{2n} > x' \)
or
       \( \ip{A}{B} = \sum_{i} a_{i} b_{i} \).
The spaces you type in a formula are
ignored.  Remember that a letter like
       $x$                   % $ ... $  and  \( ... \)  are equivalent
is a formula when it denotes a mathematical
symbol, and it should be typed as one.

\section{Displayed Text}

Text is displayed by indenting it from the left
margin.  Quotations are commonly displayed.  There
are short quotations
\begin{quote}
   This is a short quotation.  It consists of a
   single paragraph of text.  See how it is formatted.
\end{quote}
and longer ones.
\begin{quotation}
   This is a longer quotation.  It consists of two
   paragraphs of text, neither of which are
   particularly interesting.

   This is the second paragraph of the quotation.  It
   is just as dull as the first paragraph.
\end{quotation}
Another frequently-displayed structure is a list.
The following is an example of an \emph{itemized}
list.
\begin{itemize}
   \item This is the first item of an itemized list.
         Each item in the list is marked with a ``tick''.
         You don't have to worry about what kind of tick
         mark is used.

   \item This is the second item of the list.  It
         contains another list nested inside it.  The inner
         list is an \emph{enumerated} list.
         \begin{enumerate}
            \item This is the first item of an enumerated
                  list that is nested within the itemized list.

            \item This is the second item of the inner list.
                  \LaTeX\ allows you to nest lists deeper than
                  you really should.
         \end{enumerate}
         This is the rest of the second item of the outer
         list.  It is no more interesting than any other
         part of the item.
   \item This is the third item of the list.
\end{itemize}
You can even display poetry.
\begin{verse}
   There is an environment
    for verse \\             % The \\ command separates lines
   Whose features some poets % within a stanza.
   will curse.

                             % One or more blank lines separate stanzas.

   For instead of making\\
   Them do \emph{all} line breaking, \\
   It allows them to put too many words on a line when they'd rather be
   forced to be terse.
\end{verse}

Mathematical formulas may also be displayed.  A
displayed formula
is
one-line long; multiline
formulas require special formatting instructions.
   \[  \ip{\Gamma}{\psi'} = x'' + y^{2} + z_{i}^{n}\]
Don't start a paragraph with a displayed equation,
nor make one a paragraph by itself.

\end{document}               % End of document.

\section{Bench Huge Section 668}
% synthetic bench section 668
% This is a sample LaTeX input file.  (Version of 12 August 2004.)
%
% A '%' character causes TeX to ignore all remaining text on the line,
% and is used for comments like this one.

\documentclass{article}      % Specifies the document class

                             % The preamble begins here.
\title{An Example Document}  % Declares the document's title.
\author{Leslie Lamport}      % Declares the author's name.
\date{January 21, 1994}      % Deleting this command produces today's date.

\newcommand{\ip}[2]{(#1, #2)}
                             % Defines \ip{arg1}{arg2} to mean
                             % (arg1, arg2).

%\newcommand{\ip}[2]{\langle #1 | #2\rangle}
                             % This is an alternative definition of
                             % \ip that is commented out.

\begin{document}             % End of preamble and beginning of text.

\maketitle                   % Produces the title.

This is an example input file.  Comparing it with
the output it generates can show you how to
produce a simple document of your own.

\section{Ordinary Text}      % Produces section heading.  Lower-level
                             % sections are begun with similar
                             % \subsection and \subsubsection commands.

The ends  of words and sentences are marked
  by   spaces. It  doesn't matter how many
spaces    you type; one is as good as 100.  The
end of   a line counts as a space.

One   or more   blank lines denote the  end
of  a paragraph.

Since any number of consecutive spaces are treated
like a single one, the formatting of the input
file makes no difference to
      \LaTeX,                % The \LaTeX command generates the LaTeX logo.
but it makes a difference to you.  When you use
\LaTeX, making your input file as easy to read
as possible will be a great help as you write
your document and when you change it.  This sample
file shows how you can add comments to your own input
file.

Because printing is different from typewriting,
there are a number of things that you have to do
differently when preparing an input file than if
you were just typing the document directly.
Quotation marks like
       ``this''
have to be handled specially, as do quotes within
quotes:
       ``\,`this'            % \, separates the double and single quote.
        is what I just
        wrote, not  `that'\,''.

Dashes come in three sizes: an
       intra-word
dash, a medium dash for number ranges like
       1--2,
and a punctuation
       dash---like
this.

A sentence-ending space should be larger than the
space between words within a sentence.  You
sometimes have to type special commands in
conjunction with punctuation characters to get
this right, as in the following sentence.
       Gnats, gnus, etc.\ all  % `\ ' makes an inter-word space.
       begin with G\@.         % \@ marks end-of-sentence punctuation.
You should check the spaces after periods when
reading your output to make sure you haven't
forgotten any special cases.  Generating an
ellipsis
       \ldots\               % `\ ' is needed after `\ldots' because TeX
                             % ignores spaces after command names like \ldots
                             % made from \ + letters.
                             %
                             % Note how a `%' character causes TeX to ignore
                             % the end of the input line, so these blank lines
                             % do not start a new paragraph.
                             %
with the right spacing around the periods requires
a special command.

\LaTeX\ interprets some common characters as
commands, so you must type special commands to
generate them.  These characters include the
following:
       \$ \& \% \# \{ and \}.

In printing, text is usually emphasized with an
       \emph{italic}
type style.

\begin{em}
   A long segment of text can also be emphasized
   in this way.  Text within such a segment can be
   given \emph{additional} emphasis.
\end{em}

It is sometimes necessary to prevent \LaTeX\ from
breaking a line where it might otherwise do so.
This may be at a space, as between the ``Mr.''\ and
``Jones'' in
       ``Mr.~Jones'',        % ~ produces an unbreakable interword space.
or within a word---especially when the word is a
symbol like
       \mbox{\emph{itemnum}}
that makes little sense when hyphenated across
lines.

Footnotes\footnote{This is an example of a footnote.}
pose no problem.

\LaTeX\ is good at typesetting mathematical formulas
like
       \( x-3y + z = 7 \)
or
       \( a_{1} > x^{2n} + y^{2n} > x' \)
or
       \( \ip{A}{B} = \sum_{i} a_{i} b_{i} \).
The spaces you type in a formula are
ignored.  Remember that a letter like
       $x$                   % $ ... $  and  \( ... \)  are equivalent
is a formula when it denotes a mathematical
symbol, and it should be typed as one.

\section{Displayed Text}

Text is displayed by indenting it from the left
margin.  Quotations are commonly displayed.  There
are short quotations
\begin{quote}
   This is a short quotation.  It consists of a
   single paragraph of text.  See how it is formatted.
\end{quote}
and longer ones.
\begin{quotation}
   This is a longer quotation.  It consists of two
   paragraphs of text, neither of which are
   particularly interesting.

   This is the second paragraph of the quotation.  It
   is just as dull as the first paragraph.
\end{quotation}
Another frequently-displayed structure is a list.
The following is an example of an \emph{itemized}
list.
\begin{itemize}
   \item This is the first item of an itemized list.
         Each item in the list is marked with a ``tick''.
         You don't have to worry about what kind of tick
         mark is used.

   \item This is the second item of the list.  It
         contains another list nested inside it.  The inner
         list is an \emph{enumerated} list.
         \begin{enumerate}
            \item This is the first item of an enumerated
                  list that is nested within the itemized list.

            \item This is the second item of the inner list.
                  \LaTeX\ allows you to nest lists deeper than
                  you really should.
         \end{enumerate}
         This is the rest of the second item of the outer
         list.  It is no more interesting than any other
         part of the item.
   \item This is the third item of the list.
\end{itemize}
You can even display poetry.
\begin{verse}
   There is an environment
    for verse \\             % The \\ command separates lines
   Whose features some poets % within a stanza.
   will curse.

                             % One or more blank lines separate stanzas.

   For instead of making\\
   Them do \emph{all} line breaking, \\
   It allows them to put too many words on a line when they'd rather be
   forced to be terse.
\end{verse}

Mathematical formulas may also be displayed.  A
displayed formula
is
one-line long; multiline
formulas require special formatting instructions.
   \[  \ip{\Gamma}{\psi'} = x'' + y^{2} + z_{i}^{n}\]
Don't start a paragraph with a displayed equation,
nor make one a paragraph by itself.

\end{document}               % End of document.

\section{Bench Huge Section 669}
% synthetic bench section 669
% This is a sample LaTeX input file.  (Version of 12 August 2004.)
%
% A '%' character causes TeX to ignore all remaining text on the line,
% and is used for comments like this one.

\documentclass{article}      % Specifies the document class

                             % The preamble begins here.
\title{An Example Document}  % Declares the document's title.
\author{Leslie Lamport}      % Declares the author's name.
\date{January 21, 1994}      % Deleting this command produces today's date.

\newcommand{\ip}[2]{(#1, #2)}
                             % Defines \ip{arg1}{arg2} to mean
                             % (arg1, arg2).

%\newcommand{\ip}[2]{\langle #1 | #2\rangle}
                             % This is an alternative definition of
                             % \ip that is commented out.

\begin{document}             % End of preamble and beginning of text.

\maketitle                   % Produces the title.

This is an example input file.  Comparing it with
the output it generates can show you how to
produce a simple document of your own.

\section{Ordinary Text}      % Produces section heading.  Lower-level
                             % sections are begun with similar
                             % \subsection and \subsubsection commands.

The ends  of words and sentences are marked
  by   spaces. It  doesn't matter how many
spaces    you type; one is as good as 100.  The
end of   a line counts as a space.

One   or more   blank lines denote the  end
of  a paragraph.

Since any number of consecutive spaces are treated
like a single one, the formatting of the input
file makes no difference to
      \LaTeX,                % The \LaTeX command generates the LaTeX logo.
but it makes a difference to you.  When you use
\LaTeX, making your input file as easy to read
as possible will be a great help as you write
your document and when you change it.  This sample
file shows how you can add comments to your own input
file.

Because printing is different from typewriting,
there are a number of things that you have to do
differently when preparing an input file than if
you were just typing the document directly.
Quotation marks like
       ``this''
have to be handled specially, as do quotes within
quotes:
       ``\,`this'            % \, separates the double and single quote.
        is what I just
        wrote, not  `that'\,''.

Dashes come in three sizes: an
       intra-word
dash, a medium dash for number ranges like
       1--2,
and a punctuation
       dash---like
this.

A sentence-ending space should be larger than the
space between words within a sentence.  You
sometimes have to type special commands in
conjunction with punctuation characters to get
this right, as in the following sentence.
       Gnats, gnus, etc.\ all  % `\ ' makes an inter-word space.
       begin with G\@.         % \@ marks end-of-sentence punctuation.
You should check the spaces after periods when
reading your output to make sure you haven't
forgotten any special cases.  Generating an
ellipsis
       \ldots\               % `\ ' is needed after `\ldots' because TeX
                             % ignores spaces after command names like \ldots
                             % made from \ + letters.
                             %
                             % Note how a `%' character causes TeX to ignore
                             % the end of the input line, so these blank lines
                             % do not start a new paragraph.
                             %
with the right spacing around the periods requires
a special command.

\LaTeX\ interprets some common characters as
commands, so you must type special commands to
generate them.  These characters include the
following:
       \$ \& \% \# \{ and \}.

In printing, text is usually emphasized with an
       \emph{italic}
type style.

\begin{em}
   A long segment of text can also be emphasized
   in this way.  Text within such a segment can be
   given \emph{additional} emphasis.
\end{em}

It is sometimes necessary to prevent \LaTeX\ from
breaking a line where it might otherwise do so.
This may be at a space, as between the ``Mr.''\ and
``Jones'' in
       ``Mr.~Jones'',        % ~ produces an unbreakable interword space.
or within a word---especially when the word is a
symbol like
       \mbox{\emph{itemnum}}
that makes little sense when hyphenated across
lines.

Footnotes\footnote{This is an example of a footnote.}
pose no problem.

\LaTeX\ is good at typesetting mathematical formulas
like
       \( x-3y + z = 7 \)
or
       \( a_{1} > x^{2n} + y^{2n} > x' \)
or
       \( \ip{A}{B} = \sum_{i} a_{i} b_{i} \).
The spaces you type in a formula are
ignored.  Remember that a letter like
       $x$                   % $ ... $  and  \( ... \)  are equivalent
is a formula when it denotes a mathematical
symbol, and it should be typed as one.

\section{Displayed Text}

Text is displayed by indenting it from the left
margin.  Quotations are commonly displayed.  There
are short quotations
\begin{quote}
   This is a short quotation.  It consists of a
   single paragraph of text.  See how it is formatted.
\end{quote}
and longer ones.
\begin{quotation}
   This is a longer quotation.  It consists of two
   paragraphs of text, neither of which are
   particularly interesting.

   This is the second paragraph of the quotation.  It
   is just as dull as the first paragraph.
\end{quotation}
Another frequently-displayed structure is a list.
The following is an example of an \emph{itemized}
list.
\begin{itemize}
   \item This is the first item of an itemized list.
         Each item in the list is marked with a ``tick''.
         You don't have to worry about what kind of tick
         mark is used.

   \item This is the second item of the list.  It
         contains another list nested inside it.  The inner
         list is an \emph{enumerated} list.
         \begin{enumerate}
            \item This is the first item of an enumerated
                  list that is nested within the itemized list.

            \item This is the second item of the inner list.
                  \LaTeX\ allows you to nest lists deeper than
                  you really should.
         \end{enumerate}
         This is the rest of the second item of the outer
         list.  It is no more interesting than any other
         part of the item.
   \item This is the third item of the list.
\end{itemize}
You can even display poetry.
\begin{verse}
   There is an environment
    for verse \\             % The \\ command separates lines
   Whose features some poets % within a stanza.
   will curse.

                             % One or more blank lines separate stanzas.

   For instead of making\\
   Them do \emph{all} line breaking, \\
   It allows them to put too many words on a line when they'd rather be
   forced to be terse.
\end{verse}

Mathematical formulas may also be displayed.  A
displayed formula
is
one-line long; multiline
formulas require special formatting instructions.
   \[  \ip{\Gamma}{\psi'} = x'' + y^{2} + z_{i}^{n}\]
Don't start a paragraph with a displayed equation,
nor make one a paragraph by itself.

\end{document}               % End of document.

\section{Bench Huge Section 670}
% synthetic bench section 670
% This is a sample LaTeX input file.  (Version of 12 August 2004.)
%
% A '%' character causes TeX to ignore all remaining text on the line,
% and is used for comments like this one.

\documentclass{article}      % Specifies the document class

                             % The preamble begins here.
\title{An Example Document}  % Declares the document's title.
\author{Leslie Lamport}      % Declares the author's name.
\date{January 21, 1994}      % Deleting this command produces today's date.

\newcommand{\ip}[2]{(#1, #2)}
                             % Defines \ip{arg1}{arg2} to mean
                             % (arg1, arg2).

%\newcommand{\ip}[2]{\langle #1 | #2\rangle}
                             % This is an alternative definition of
                             % \ip that is commented out.

\begin{document}             % End of preamble and beginning of text.

\maketitle                   % Produces the title.

This is an example input file.  Comparing it with
the output it generates can show you how to
produce a simple document of your own.

\section{Ordinary Text}      % Produces section heading.  Lower-level
                             % sections are begun with similar
                             % \subsection and \subsubsection commands.

The ends  of words and sentences are marked
  by   spaces. It  doesn't matter how many
spaces    you type; one is as good as 100.  The
end of   a line counts as a space.

One   or more   blank lines denote the  end
of  a paragraph.

Since any number of consecutive spaces are treated
like a single one, the formatting of the input
file makes no difference to
      \LaTeX,                % The \LaTeX command generates the LaTeX logo.
but it makes a difference to you.  When you use
\LaTeX, making your input file as easy to read
as possible will be a great help as you write
your document and when you change it.  This sample
file shows how you can add comments to your own input
file.

Because printing is different from typewriting,
there are a number of things that you have to do
differently when preparing an input file than if
you were just typing the document directly.
Quotation marks like
       ``this''
have to be handled specially, as do quotes within
quotes:
       ``\,`this'            % \, separates the double and single quote.
        is what I just
        wrote, not  `that'\,''.

Dashes come in three sizes: an
       intra-word
dash, a medium dash for number ranges like
       1--2,
and a punctuation
       dash---like
this.

A sentence-ending space should be larger than the
space between words within a sentence.  You
sometimes have to type special commands in
conjunction with punctuation characters to get
this right, as in the following sentence.
       Gnats, gnus, etc.\ all  % `\ ' makes an inter-word space.
       begin with G\@.         % \@ marks end-of-sentence punctuation.
You should check the spaces after periods when
reading your output to make sure you haven't
forgotten any special cases.  Generating an
ellipsis
       \ldots\               % `\ ' is needed after `\ldots' because TeX
                             % ignores spaces after command names like \ldots
                             % made from \ + letters.
                             %
                             % Note how a `%' character causes TeX to ignore
                             % the end of the input line, so these blank lines
                             % do not start a new paragraph.
                             %
with the right spacing around the periods requires
a special command.

\LaTeX\ interprets some common characters as
commands, so you must type special commands to
generate them.  These characters include the
following:
       \$ \& \% \# \{ and \}.

In printing, text is usually emphasized with an
       \emph{italic}
type style.

\begin{em}
   A long segment of text can also be emphasized
   in this way.  Text within such a segment can be
   given \emph{additional} emphasis.
\end{em}

It is sometimes necessary to prevent \LaTeX\ from
breaking a line where it might otherwise do so.
This may be at a space, as between the ``Mr.''\ and
``Jones'' in
       ``Mr.~Jones'',        % ~ produces an unbreakable interword space.
or within a word---especially when the word is a
symbol like
       \mbox{\emph{itemnum}}
that makes little sense when hyphenated across
lines.

Footnotes\footnote{This is an example of a footnote.}
pose no problem.

\LaTeX\ is good at typesetting mathematical formulas
like
       \( x-3y + z = 7 \)
or
       \( a_{1} > x^{2n} + y^{2n} > x' \)
or
       \( \ip{A}{B} = \sum_{i} a_{i} b_{i} \).
The spaces you type in a formula are
ignored.  Remember that a letter like
       $x$                   % $ ... $  and  \( ... \)  are equivalent
is a formula when it denotes a mathematical
symbol, and it should be typed as one.

\section{Displayed Text}

Text is displayed by indenting it from the left
margin.  Quotations are commonly displayed.  There
are short quotations
\begin{quote}
   This is a short quotation.  It consists of a
   single paragraph of text.  See how it is formatted.
\end{quote}
and longer ones.
\begin{quotation}
   This is a longer quotation.  It consists of two
   paragraphs of text, neither of which are
   particularly interesting.

   This is the second paragraph of the quotation.  It
   is just as dull as the first paragraph.
\end{quotation}
Another frequently-displayed structure is a list.
The following is an example of an \emph{itemized}
list.
\begin{itemize}
   \item This is the first item of an itemized list.
         Each item in the list is marked with a ``tick''.
         You don't have to worry about what kind of tick
         mark is used.

   \item This is the second item of the list.  It
         contains another list nested inside it.  The inner
         list is an \emph{enumerated} list.
         \begin{enumerate}
            \item This is the first item of an enumerated
                  list that is nested within the itemized list.

            \item This is the second item of the inner list.
                  \LaTeX\ allows you to nest lists deeper than
                  you really should.
         \end{enumerate}
         This is the rest of the second item of the outer
         list.  It is no more interesting than any other
         part of the item.
   \item This is the third item of the list.
\end{itemize}
You can even display poetry.
\begin{verse}
   There is an environment
    for verse \\             % The \\ command separates lines
   Whose features some poets % within a stanza.
   will curse.

                             % One or more blank lines separate stanzas.

   For instead of making\\
   Them do \emph{all} line breaking, \\
   It allows them to put too many words on a line when they'd rather be
   forced to be terse.
\end{verse}

Mathematical formulas may also be displayed.  A
displayed formula
is
one-line long; multiline
formulas require special formatting instructions.
   \[  \ip{\Gamma}{\psi'} = x'' + y^{2} + z_{i}^{n}\]
Don't start a paragraph with a displayed equation,
nor make one a paragraph by itself.

\end{document}               % End of document.

\section{Bench Huge Section 671}
% synthetic bench section 671
% This is a sample LaTeX input file.  (Version of 12 August 2004.)
%
% A '%' character causes TeX to ignore all remaining text on the line,
% and is used for comments like this one.

\documentclass{article}      % Specifies the document class

                             % The preamble begins here.
\title{An Example Document}  % Declares the document's title.
\author{Leslie Lamport}      % Declares the author's name.
\date{January 21, 1994}      % Deleting this command produces today's date.

\newcommand{\ip}[2]{(#1, #2)}
                             % Defines \ip{arg1}{arg2} to mean
                             % (arg1, arg2).

%\newcommand{\ip}[2]{\langle #1 | #2\rangle}
                             % This is an alternative definition of
                             % \ip that is commented out.

\begin{document}             % End of preamble and beginning of text.

\maketitle                   % Produces the title.

This is an example input file.  Comparing it with
the output it generates can show you how to
produce a simple document of your own.

\section{Ordinary Text}      % Produces section heading.  Lower-level
                             % sections are begun with similar
                             % \subsection and \subsubsection commands.

The ends  of words and sentences are marked
  by   spaces. It  doesn't matter how many
spaces    you type; one is as good as 100.  The
end of   a line counts as a space.

One   or more   blank lines denote the  end
of  a paragraph.

Since any number of consecutive spaces are treated
like a single one, the formatting of the input
file makes no difference to
      \LaTeX,                % The \LaTeX command generates the LaTeX logo.
but it makes a difference to you.  When you use
\LaTeX, making your input file as easy to read
as possible will be a great help as you write
your document and when you change it.  This sample
file shows how you can add comments to your own input
file.

Because printing is different from typewriting,
there are a number of things that you have to do
differently when preparing an input file than if
you were just typing the document directly.
Quotation marks like
       ``this''
have to be handled specially, as do quotes within
quotes:
       ``\,`this'            % \, separates the double and single quote.
        is what I just
        wrote, not  `that'\,''.

Dashes come in three sizes: an
       intra-word
dash, a medium dash for number ranges like
       1--2,
and a punctuation
       dash---like
this.

A sentence-ending space should be larger than the
space between words within a sentence.  You
sometimes have to type special commands in
conjunction with punctuation characters to get
this right, as in the following sentence.
       Gnats, gnus, etc.\ all  % `\ ' makes an inter-word space.
       begin with G\@.         % \@ marks end-of-sentence punctuation.
You should check the spaces after periods when
reading your output to make sure you haven't
forgotten any special cases.  Generating an
ellipsis
       \ldots\               % `\ ' is needed after `\ldots' because TeX
                             % ignores spaces after command names like \ldots
                             % made from \ + letters.
                             %
                             % Note how a `%' character causes TeX to ignore
                             % the end of the input line, so these blank lines
                             % do not start a new paragraph.
                             %
with the right spacing around the periods requires
a special command.

\LaTeX\ interprets some common characters as
commands, so you must type special commands to
generate them.  These characters include the
following:
       \$ \& \% \# \{ and \}.

In printing, text is usually emphasized with an
       \emph{italic}
type style.

\begin{em}
   A long segment of text can also be emphasized
   in this way.  Text within such a segment can be
   given \emph{additional} emphasis.
\end{em}

It is sometimes necessary to prevent \LaTeX\ from
breaking a line where it might otherwise do so.
This may be at a space, as between the ``Mr.''\ and
``Jones'' in
       ``Mr.~Jones'',        % ~ produces an unbreakable interword space.
or within a word---especially when the word is a
symbol like
       \mbox{\emph{itemnum}}
that makes little sense when hyphenated across
lines.

Footnotes\footnote{This is an example of a footnote.}
pose no problem.

\LaTeX\ is good at typesetting mathematical formulas
like
       \( x-3y + z = 7 \)
or
       \( a_{1} > x^{2n} + y^{2n} > x' \)
or
       \( \ip{A}{B} = \sum_{i} a_{i} b_{i} \).
The spaces you type in a formula are
ignored.  Remember that a letter like
       $x$                   % $ ... $  and  \( ... \)  are equivalent
is a formula when it denotes a mathematical
symbol, and it should be typed as one.

\section{Displayed Text}

Text is displayed by indenting it from the left
margin.  Quotations are commonly displayed.  There
are short quotations
\begin{quote}
   This is a short quotation.  It consists of a
   single paragraph of text.  See how it is formatted.
\end{quote}
and longer ones.
\begin{quotation}
   This is a longer quotation.  It consists of two
   paragraphs of text, neither of which are
   particularly interesting.

   This is the second paragraph of the quotation.  It
   is just as dull as the first paragraph.
\end{quotation}
Another frequently-displayed structure is a list.
The following is an example of an \emph{itemized}
list.
\begin{itemize}
   \item This is the first item of an itemized list.
         Each item in the list is marked with a ``tick''.
         You don't have to worry about what kind of tick
         mark is used.

   \item This is the second item of the list.  It
         contains another list nested inside it.  The inner
         list is an \emph{enumerated} list.
         \begin{enumerate}
            \item This is the first item of an enumerated
                  list that is nested within the itemized list.

            \item This is the second item of the inner list.
                  \LaTeX\ allows you to nest lists deeper than
                  you really should.
         \end{enumerate}
         This is the rest of the second item of the outer
         list.  It is no more interesting than any other
         part of the item.
   \item This is the third item of the list.
\end{itemize}
You can even display poetry.
\begin{verse}
   There is an environment
    for verse \\             % The \\ command separates lines
   Whose features some poets % within a stanza.
   will curse.

                             % One or more blank lines separate stanzas.

   For instead of making\\
   Them do \emph{all} line breaking, \\
   It allows them to put too many words on a line when they'd rather be
   forced to be terse.
\end{verse}

Mathematical formulas may also be displayed.  A
displayed formula
is
one-line long; multiline
formulas require special formatting instructions.
   \[  \ip{\Gamma}{\psi'} = x'' + y^{2} + z_{i}^{n}\]
Don't start a paragraph with a displayed equation,
nor make one a paragraph by itself.

\end{document}               % End of document.

\section{Bench Huge Section 672}
% synthetic bench section 672
% This is a sample LaTeX input file.  (Version of 12 August 2004.)
%
% A '%' character causes TeX to ignore all remaining text on the line,
% and is used for comments like this one.

\documentclass{article}      % Specifies the document class

                             % The preamble begins here.
\title{An Example Document}  % Declares the document's title.
\author{Leslie Lamport}      % Declares the author's name.
\date{January 21, 1994}      % Deleting this command produces today's date.

\newcommand{\ip}[2]{(#1, #2)}
                             % Defines \ip{arg1}{arg2} to mean
                             % (arg1, arg2).

%\newcommand{\ip}[2]{\langle #1 | #2\rangle}
                             % This is an alternative definition of
                             % \ip that is commented out.

\begin{document}             % End of preamble and beginning of text.

\maketitle                   % Produces the title.

This is an example input file.  Comparing it with
the output it generates can show you how to
produce a simple document of your own.

\section{Ordinary Text}      % Produces section heading.  Lower-level
                             % sections are begun with similar
                             % \subsection and \subsubsection commands.

The ends  of words and sentences are marked
  by   spaces. It  doesn't matter how many
spaces    you type; one is as good as 100.  The
end of   a line counts as a space.

One   or more   blank lines denote the  end
of  a paragraph.

Since any number of consecutive spaces are treated
like a single one, the formatting of the input
file makes no difference to
      \LaTeX,                % The \LaTeX command generates the LaTeX logo.
but it makes a difference to you.  When you use
\LaTeX, making your input file as easy to read
as possible will be a great help as you write
your document and when you change it.  This sample
file shows how you can add comments to your own input
file.

Because printing is different from typewriting,
there are a number of things that you have to do
differently when preparing an input file than if
you were just typing the document directly.
Quotation marks like
       ``this''
have to be handled specially, as do quotes within
quotes:
       ``\,`this'            % \, separates the double and single quote.
        is what I just
        wrote, not  `that'\,''.

Dashes come in three sizes: an
       intra-word
dash, a medium dash for number ranges like
       1--2,
and a punctuation
       dash---like
this.

A sentence-ending space should be larger than the
space between words within a sentence.  You
sometimes have to type special commands in
conjunction with punctuation characters to get
this right, as in the following sentence.
       Gnats, gnus, etc.\ all  % `\ ' makes an inter-word space.
       begin with G\@.         % \@ marks end-of-sentence punctuation.
You should check the spaces after periods when
reading your output to make sure you haven't
forgotten any special cases.  Generating an
ellipsis
       \ldots\               % `\ ' is needed after `\ldots' because TeX
                             % ignores spaces after command names like \ldots
                             % made from \ + letters.
                             %
                             % Note how a `%' character causes TeX to ignore
                             % the end of the input line, so these blank lines
                             % do not start a new paragraph.
                             %
with the right spacing around the periods requires
a special command.

\LaTeX\ interprets some common characters as
commands, so you must type special commands to
generate them.  These characters include the
following:
       \$ \& \% \# \{ and \}.

In printing, text is usually emphasized with an
       \emph{italic}
type style.

\begin{em}
   A long segment of text can also be emphasized
   in this way.  Text within such a segment can be
   given \emph{additional} emphasis.
\end{em}

It is sometimes necessary to prevent \LaTeX\ from
breaking a line where it might otherwise do so.
This may be at a space, as between the ``Mr.''\ and
``Jones'' in
       ``Mr.~Jones'',        % ~ produces an unbreakable interword space.
or within a word---especially when the word is a
symbol like
       \mbox{\emph{itemnum}}
that makes little sense when hyphenated across
lines.

Footnotes\footnote{This is an example of a footnote.}
pose no problem.

\LaTeX\ is good at typesetting mathematical formulas
like
       \( x-3y + z = 7 \)
or
       \( a_{1} > x^{2n} + y^{2n} > x' \)
or
       \( \ip{A}{B} = \sum_{i} a_{i} b_{i} \).
The spaces you type in a formula are
ignored.  Remember that a letter like
       $x$                   % $ ... $  and  \( ... \)  are equivalent
is a formula when it denotes a mathematical
symbol, and it should be typed as one.

\section{Displayed Text}

Text is displayed by indenting it from the left
margin.  Quotations are commonly displayed.  There
are short quotations
\begin{quote}
   This is a short quotation.  It consists of a
   single paragraph of text.  See how it is formatted.
\end{quote}
and longer ones.
\begin{quotation}
   This is a longer quotation.  It consists of two
   paragraphs of text, neither of which are
   particularly interesting.

   This is the second paragraph of the quotation.  It
   is just as dull as the first paragraph.
\end{quotation}
Another frequently-displayed structure is a list.
The following is an example of an \emph{itemized}
list.
\begin{itemize}
   \item This is the first item of an itemized list.
         Each item in the list is marked with a ``tick''.
         You don't have to worry about what kind of tick
         mark is used.

   \item This is the second item of the list.  It
         contains another list nested inside it.  The inner
         list is an \emph{enumerated} list.
         \begin{enumerate}
            \item This is the first item of an enumerated
                  list that is nested within the itemized list.

            \item This is the second item of the inner list.
                  \LaTeX\ allows you to nest lists deeper than
                  you really should.
         \end{enumerate}
         This is the rest of the second item of the outer
         list.  It is no more interesting than any other
         part of the item.
   \item This is the third item of the list.
\end{itemize}
You can even display poetry.
\begin{verse}
   There is an environment
    for verse \\             % The \\ command separates lines
   Whose features some poets % within a stanza.
   will curse.

                             % One or more blank lines separate stanzas.

   For instead of making\\
   Them do \emph{all} line breaking, \\
   It allows them to put too many words on a line when they'd rather be
   forced to be terse.
\end{verse}

Mathematical formulas may also be displayed.  A
displayed formula
is
one-line long; multiline
formulas require special formatting instructions.
   \[  \ip{\Gamma}{\psi'} = x'' + y^{2} + z_{i}^{n}\]
Don't start a paragraph with a displayed equation,
nor make one a paragraph by itself.

\end{document}               % End of document.

\section{Bench Huge Section 673}
% synthetic bench section 673
% This is a sample LaTeX input file.  (Version of 12 August 2004.)
%
% A '%' character causes TeX to ignore all remaining text on the line,
% and is used for comments like this one.

\documentclass{article}      % Specifies the document class

                             % The preamble begins here.
\title{An Example Document}  % Declares the document's title.
\author{Leslie Lamport}      % Declares the author's name.
\date{January 21, 1994}      % Deleting this command produces today's date.

\newcommand{\ip}[2]{(#1, #2)}
                             % Defines \ip{arg1}{arg2} to mean
                             % (arg1, arg2).

%\newcommand{\ip}[2]{\langle #1 | #2\rangle}
                             % This is an alternative definition of
                             % \ip that is commented out.

\begin{document}             % End of preamble and beginning of text.

\maketitle                   % Produces the title.

This is an example input file.  Comparing it with
the output it generates can show you how to
produce a simple document of your own.

\section{Ordinary Text}      % Produces section heading.  Lower-level
                             % sections are begun with similar
                             % \subsection and \subsubsection commands.

The ends  of words and sentences are marked
  by   spaces. It  doesn't matter how many
spaces    you type; one is as good as 100.  The
end of   a line counts as a space.

One   or more   blank lines denote the  end
of  a paragraph.

Since any number of consecutive spaces are treated
like a single one, the formatting of the input
file makes no difference to
      \LaTeX,                % The \LaTeX command generates the LaTeX logo.
but it makes a difference to you.  When you use
\LaTeX, making your input file as easy to read
as possible will be a great help as you write
your document and when you change it.  This sample
file shows how you can add comments to your own input
file.

Because printing is different from typewriting,
there are a number of things that you have to do
differently when preparing an input file than if
you were just typing the document directly.
Quotation marks like
       ``this''
have to be handled specially, as do quotes within
quotes:
       ``\,`this'            % \, separates the double and single quote.
        is what I just
        wrote, not  `that'\,''.

Dashes come in three sizes: an
       intra-word
dash, a medium dash for number ranges like
       1--2,
and a punctuation
       dash---like
this.

A sentence-ending space should be larger than the
space between words within a sentence.  You
sometimes have to type special commands in
conjunction with punctuation characters to get
this right, as in the following sentence.
       Gnats, gnus, etc.\ all  % `\ ' makes an inter-word space.
       begin with G\@.         % \@ marks end-of-sentence punctuation.
You should check the spaces after periods when
reading your output to make sure you haven't
forgotten any special cases.  Generating an
ellipsis
       \ldots\               % `\ ' is needed after `\ldots' because TeX
                             % ignores spaces after command names like \ldots
                             % made from \ + letters.
                             %
                             % Note how a `%' character causes TeX to ignore
                             % the end of the input line, so these blank lines
                             % do not start a new paragraph.
                             %
with the right spacing around the periods requires
a special command.

\LaTeX\ interprets some common characters as
commands, so you must type special commands to
generate them.  These characters include the
following:
       \$ \& \% \# \{ and \}.

In printing, text is usually emphasized with an
       \emph{italic}
type style.

\begin{em}
   A long segment of text can also be emphasized
   in this way.  Text within such a segment can be
   given \emph{additional} emphasis.
\end{em}

It is sometimes necessary to prevent \LaTeX\ from
breaking a line where it might otherwise do so.
This may be at a space, as between the ``Mr.''\ and
``Jones'' in
       ``Mr.~Jones'',        % ~ produces an unbreakable interword space.
or within a word---especially when the word is a
symbol like
       \mbox{\emph{itemnum}}
that makes little sense when hyphenated across
lines.

Footnotes\footnote{This is an example of a footnote.}
pose no problem.

\LaTeX\ is good at typesetting mathematical formulas
like
       \( x-3y + z = 7 \)
or
       \( a_{1} > x^{2n} + y^{2n} > x' \)
or
       \( \ip{A}{B} = \sum_{i} a_{i} b_{i} \).
The spaces you type in a formula are
ignored.  Remember that a letter like
       $x$                   % $ ... $  and  \( ... \)  are equivalent
is a formula when it denotes a mathematical
symbol, and it should be typed as one.

\section{Displayed Text}

Text is displayed by indenting it from the left
margin.  Quotations are commonly displayed.  There
are short quotations
\begin{quote}
   This is a short quotation.  It consists of a
   single paragraph of text.  See how it is formatted.
\end{quote}
and longer ones.
\begin{quotation}
   This is a longer quotation.  It consists of two
   paragraphs of text, neither of which are
   particularly interesting.

   This is the second paragraph of the quotation.  It
   is just as dull as the first paragraph.
\end{quotation}
Another frequently-displayed structure is a list.
The following is an example of an \emph{itemized}
list.
\begin{itemize}
   \item This is the first item of an itemized list.
         Each item in the list is marked with a ``tick''.
         You don't have to worry about what kind of tick
         mark is used.

   \item This is the second item of the list.  It
         contains another list nested inside it.  The inner
         list is an \emph{enumerated} list.
         \begin{enumerate}
            \item This is the first item of an enumerated
                  list that is nested within the itemized list.

            \item This is the second item of the inner list.
                  \LaTeX\ allows you to nest lists deeper than
                  you really should.
         \end{enumerate}
         This is the rest of the second item of the outer
         list.  It is no more interesting than any other
         part of the item.
   \item This is the third item of the list.
\end{itemize}
You can even display poetry.
\begin{verse}
   There is an environment
    for verse \\             % The \\ command separates lines
   Whose features some poets % within a stanza.
   will curse.

                             % One or more blank lines separate stanzas.

   For instead of making\\
   Them do \emph{all} line breaking, \\
   It allows them to put too many words on a line when they'd rather be
   forced to be terse.
\end{verse}

Mathematical formulas may also be displayed.  A
displayed formula
is
one-line long; multiline
formulas require special formatting instructions.
   \[  \ip{\Gamma}{\psi'} = x'' + y^{2} + z_{i}^{n}\]
Don't start a paragraph with a displayed equation,
nor make one a paragraph by itself.

\end{document}               % End of document.

\section{Bench Huge Section 674}
% synthetic bench section 674
% This is a sample LaTeX input file.  (Version of 12 August 2004.)
%
% A '%' character causes TeX to ignore all remaining text on the line,
% and is used for comments like this one.

\documentclass{article}      % Specifies the document class

                             % The preamble begins here.
\title{An Example Document}  % Declares the document's title.
\author{Leslie Lamport}      % Declares the author's name.
\date{January 21, 1994}      % Deleting this command produces today's date.

\newcommand{\ip}[2]{(#1, #2)}
                             % Defines \ip{arg1}{arg2} to mean
                             % (arg1, arg2).

%\newcommand{\ip}[2]{\langle #1 | #2\rangle}
                             % This is an alternative definition of
                             % \ip that is commented out.

\begin{document}             % End of preamble and beginning of text.

\maketitle                   % Produces the title.

This is an example input file.  Comparing it with
the output it generates can show you how to
produce a simple document of your own.

\section{Ordinary Text}      % Produces section heading.  Lower-level
                             % sections are begun with similar
                             % \subsection and \subsubsection commands.

The ends  of words and sentences are marked
  by   spaces. It  doesn't matter how many
spaces    you type; one is as good as 100.  The
end of   a line counts as a space.

One   or more   blank lines denote the  end
of  a paragraph.

Since any number of consecutive spaces are treated
like a single one, the formatting of the input
file makes no difference to
      \LaTeX,                % The \LaTeX command generates the LaTeX logo.
but it makes a difference to you.  When you use
\LaTeX, making your input file as easy to read
as possible will be a great help as you write
your document and when you change it.  This sample
file shows how you can add comments to your own input
file.

Because printing is different from typewriting,
there are a number of things that you have to do
differently when preparing an input file than if
you were just typing the document directly.
Quotation marks like
       ``this''
have to be handled specially, as do quotes within
quotes:
       ``\,`this'            % \, separates the double and single quote.
        is what I just
        wrote, not  `that'\,''.

Dashes come in three sizes: an
       intra-word
dash, a medium dash for number ranges like
       1--2,
and a punctuation
       dash---like
this.

A sentence-ending space should be larger than the
space between words within a sentence.  You
sometimes have to type special commands in
conjunction with punctuation characters to get
this right, as in the following sentence.
       Gnats, gnus, etc.\ all  % `\ ' makes an inter-word space.
       begin with G\@.         % \@ marks end-of-sentence punctuation.
You should check the spaces after periods when
reading your output to make sure you haven't
forgotten any special cases.  Generating an
ellipsis
       \ldots\               % `\ ' is needed after `\ldots' because TeX
                             % ignores spaces after command names like \ldots
                             % made from \ + letters.
                             %
                             % Note how a `%' character causes TeX to ignore
                             % the end of the input line, so these blank lines
                             % do not start a new paragraph.
                             %
with the right spacing around the periods requires
a special command.

\LaTeX\ interprets some common characters as
commands, so you must type special commands to
generate them.  These characters include the
following:
       \$ \& \% \# \{ and \}.

In printing, text is usually emphasized with an
       \emph{italic}
type style.

\begin{em}
   A long segment of text can also be emphasized
   in this way.  Text within such a segment can be
   given \emph{additional} emphasis.
\end{em}

It is sometimes necessary to prevent \LaTeX\ from
breaking a line where it might otherwise do so.
This may be at a space, as between the ``Mr.''\ and
``Jones'' in
       ``Mr.~Jones'',        % ~ produces an unbreakable interword space.
or within a word---especially when the word is a
symbol like
       \mbox{\emph{itemnum}}
that makes little sense when hyphenated across
lines.

Footnotes\footnote{This is an example of a footnote.}
pose no problem.

\LaTeX\ is good at typesetting mathematical formulas
like
       \( x-3y + z = 7 \)
or
       \( a_{1} > x^{2n} + y^{2n} > x' \)
or
       \( \ip{A}{B} = \sum_{i} a_{i} b_{i} \).
The spaces you type in a formula are
ignored.  Remember that a letter like
       $x$                   % $ ... $  and  \( ... \)  are equivalent
is a formula when it denotes a mathematical
symbol, and it should be typed as one.

\section{Displayed Text}

Text is displayed by indenting it from the left
margin.  Quotations are commonly displayed.  There
are short quotations
\begin{quote}
   This is a short quotation.  It consists of a
   single paragraph of text.  See how it is formatted.
\end{quote}
and longer ones.
\begin{quotation}
   This is a longer quotation.  It consists of two
   paragraphs of text, neither of which are
   particularly interesting.

   This is the second paragraph of the quotation.  It
   is just as dull as the first paragraph.
\end{quotation}
Another frequently-displayed structure is a list.
The following is an example of an \emph{itemized}
list.
\begin{itemize}
   \item This is the first item of an itemized list.
         Each item in the list is marked with a ``tick''.
         You don't have to worry about what kind of tick
         mark is used.

   \item This is the second item of the list.  It
         contains another list nested inside it.  The inner
         list is an \emph{enumerated} list.
         \begin{enumerate}
            \item This is the first item of an enumerated
                  list that is nested within the itemized list.

            \item This is the second item of the inner list.
                  \LaTeX\ allows you to nest lists deeper than
                  you really should.
         \end{enumerate}
         This is the rest of the second item of the outer
         list.  It is no more interesting than any other
         part of the item.
   \item This is the third item of the list.
\end{itemize}
You can even display poetry.
\begin{verse}
   There is an environment
    for verse \\             % The \\ command separates lines
   Whose features some poets % within a stanza.
   will curse.

                             % One or more blank lines separate stanzas.

   For instead of making\\
   Them do \emph{all} line breaking, \\
   It allows them to put too many words on a line when they'd rather be
   forced to be terse.
\end{verse}

Mathematical formulas may also be displayed.  A
displayed formula
is
one-line long; multiline
formulas require special formatting instructions.
   \[  \ip{\Gamma}{\psi'} = x'' + y^{2} + z_{i}^{n}\]
Don't start a paragraph with a displayed equation,
nor make one a paragraph by itself.

\end{document}               % End of document.

\section{Bench Huge Section 675}
% synthetic bench section 675
% This is a sample LaTeX input file.  (Version of 12 August 2004.)
%
% A '%' character causes TeX to ignore all remaining text on the line,
% and is used for comments like this one.

\documentclass{article}      % Specifies the document class

                             % The preamble begins here.
\title{An Example Document}  % Declares the document's title.
\author{Leslie Lamport}      % Declares the author's name.
\date{January 21, 1994}      % Deleting this command produces today's date.

\newcommand{\ip}[2]{(#1, #2)}
                             % Defines \ip{arg1}{arg2} to mean
                             % (arg1, arg2).

%\newcommand{\ip}[2]{\langle #1 | #2\rangle}
                             % This is an alternative definition of
                             % \ip that is commented out.

\begin{document}             % End of preamble and beginning of text.

\maketitle                   % Produces the title.

This is an example input file.  Comparing it with
the output it generates can show you how to
produce a simple document of your own.

\section{Ordinary Text}      % Produces section heading.  Lower-level
                             % sections are begun with similar
                             % \subsection and \subsubsection commands.

The ends  of words and sentences are marked
  by   spaces. It  doesn't matter how many
spaces    you type; one is as good as 100.  The
end of   a line counts as a space.

One   or more   blank lines denote the  end
of  a paragraph.

Since any number of consecutive spaces are treated
like a single one, the formatting of the input
file makes no difference to
      \LaTeX,                % The \LaTeX command generates the LaTeX logo.
but it makes a difference to you.  When you use
\LaTeX, making your input file as easy to read
as possible will be a great help as you write
your document and when you change it.  This sample
file shows how you can add comments to your own input
file.

Because printing is different from typewriting,
there are a number of things that you have to do
differently when preparing an input file than if
you were just typing the document directly.
Quotation marks like
       ``this''
have to be handled specially, as do quotes within
quotes:
       ``\,`this'            % \, separates the double and single quote.
        is what I just
        wrote, not  `that'\,''.

Dashes come in three sizes: an
       intra-word
dash, a medium dash for number ranges like
       1--2,
and a punctuation
       dash---like
this.

A sentence-ending space should be larger than the
space between words within a sentence.  You
sometimes have to type special commands in
conjunction with punctuation characters to get
this right, as in the following sentence.
       Gnats, gnus, etc.\ all  % `\ ' makes an inter-word space.
       begin with G\@.         % \@ marks end-of-sentence punctuation.
You should check the spaces after periods when
reading your output to make sure you haven't
forgotten any special cases.  Generating an
ellipsis
       \ldots\               % `\ ' is needed after `\ldots' because TeX
                             % ignores spaces after command names like \ldots
                             % made from \ + letters.
                             %
                             % Note how a `%' character causes TeX to ignore
                             % the end of the input line, so these blank lines
                             % do not start a new paragraph.
                             %
with the right spacing around the periods requires
a special command.

\LaTeX\ interprets some common characters as
commands, so you must type special commands to
generate them.  These characters include the
following:
       \$ \& \% \# \{ and \}.

In printing, text is usually emphasized with an
       \emph{italic}
type style.

\begin{em}
   A long segment of text can also be emphasized
   in this way.  Text within such a segment can be
   given \emph{additional} emphasis.
\end{em}

It is sometimes necessary to prevent \LaTeX\ from
breaking a line where it might otherwise do so.
This may be at a space, as between the ``Mr.''\ and
``Jones'' in
       ``Mr.~Jones'',        % ~ produces an unbreakable interword space.
or within a word---especially when the word is a
symbol like
       \mbox{\emph{itemnum}}
that makes little sense when hyphenated across
lines.

Footnotes\footnote{This is an example of a footnote.}
pose no problem.

\LaTeX\ is good at typesetting mathematical formulas
like
       \( x-3y + z = 7 \)
or
       \( a_{1} > x^{2n} + y^{2n} > x' \)
or
       \( \ip{A}{B} = \sum_{i} a_{i} b_{i} \).
The spaces you type in a formula are
ignored.  Remember that a letter like
       $x$                   % $ ... $  and  \( ... \)  are equivalent
is a formula when it denotes a mathematical
symbol, and it should be typed as one.

\section{Displayed Text}

Text is displayed by indenting it from the left
margin.  Quotations are commonly displayed.  There
are short quotations
\begin{quote}
   This is a short quotation.  It consists of a
   single paragraph of text.  See how it is formatted.
\end{quote}
and longer ones.
\begin{quotation}
   This is a longer quotation.  It consists of two
   paragraphs of text, neither of which are
   particularly interesting.

   This is the second paragraph of the quotation.  It
   is just as dull as the first paragraph.
\end{quotation}
Another frequently-displayed structure is a list.
The following is an example of an \emph{itemized}
list.
\begin{itemize}
   \item This is the first item of an itemized list.
         Each item in the list is marked with a ``tick''.
         You don't have to worry about what kind of tick
         mark is used.

   \item This is the second item of the list.  It
         contains another list nested inside it.  The inner
         list is an \emph{enumerated} list.
         \begin{enumerate}
            \item This is the first item of an enumerated
                  list that is nested within the itemized list.

            \item This is the second item of the inner list.
                  \LaTeX\ allows you to nest lists deeper than
                  you really should.
         \end{enumerate}
         This is the rest of the second item of the outer
         list.  It is no more interesting than any other
         part of the item.
   \item This is the third item of the list.
\end{itemize}
You can even display poetry.
\begin{verse}
   There is an environment
    for verse \\             % The \\ command separates lines
   Whose features some poets % within a stanza.
   will curse.

                             % One or more blank lines separate stanzas.

   For instead of making\\
   Them do \emph{all} line breaking, \\
   It allows them to put too many words on a line when they'd rather be
   forced to be terse.
\end{verse}

Mathematical formulas may also be displayed.  A
displayed formula
is
one-line long; multiline
formulas require special formatting instructions.
   \[  \ip{\Gamma}{\psi'} = x'' + y^{2} + z_{i}^{n}\]
Don't start a paragraph with a displayed equation,
nor make one a paragraph by itself.

\end{document}               % End of document.

\section{Bench Huge Section 676}
% synthetic bench section 676
% This is a sample LaTeX input file.  (Version of 12 August 2004.)
%
% A '%' character causes TeX to ignore all remaining text on the line,
% and is used for comments like this one.

\documentclass{article}      % Specifies the document class

                             % The preamble begins here.
\title{An Example Document}  % Declares the document's title.
\author{Leslie Lamport}      % Declares the author's name.
\date{January 21, 1994}      % Deleting this command produces today's date.

\newcommand{\ip}[2]{(#1, #2)}
                             % Defines \ip{arg1}{arg2} to mean
                             % (arg1, arg2).

%\newcommand{\ip}[2]{\langle #1 | #2\rangle}
                             % This is an alternative definition of
                             % \ip that is commented out.

\begin{document}             % End of preamble and beginning of text.

\maketitle                   % Produces the title.

This is an example input file.  Comparing it with
the output it generates can show you how to
produce a simple document of your own.

\section{Ordinary Text}      % Produces section heading.  Lower-level
                             % sections are begun with similar
                             % \subsection and \subsubsection commands.

The ends  of words and sentences are marked
  by   spaces. It  doesn't matter how many
spaces    you type; one is as good as 100.  The
end of   a line counts as a space.

One   or more   blank lines denote the  end
of  a paragraph.

Since any number of consecutive spaces are treated
like a single one, the formatting of the input
file makes no difference to
      \LaTeX,                % The \LaTeX command generates the LaTeX logo.
but it makes a difference to you.  When you use
\LaTeX, making your input file as easy to read
as possible will be a great help as you write
your document and when you change it.  This sample
file shows how you can add comments to your own input
file.

Because printing is different from typewriting,
there are a number of things that you have to do
differently when preparing an input file than if
you were just typing the document directly.
Quotation marks like
       ``this''
have to be handled specially, as do quotes within
quotes:
       ``\,`this'            % \, separates the double and single quote.
        is what I just
        wrote, not  `that'\,''.

Dashes come in three sizes: an
       intra-word
dash, a medium dash for number ranges like
       1--2,
and a punctuation
       dash---like
this.

A sentence-ending space should be larger than the
space between words within a sentence.  You
sometimes have to type special commands in
conjunction with punctuation characters to get
this right, as in the following sentence.
       Gnats, gnus, etc.\ all  % `\ ' makes an inter-word space.
       begin with G\@.         % \@ marks end-of-sentence punctuation.
You should check the spaces after periods when
reading your output to make sure you haven't
forgotten any special cases.  Generating an
ellipsis
       \ldots\               % `\ ' is needed after `\ldots' because TeX
                             % ignores spaces after command names like \ldots
                             % made from \ + letters.
                             %
                             % Note how a `%' character causes TeX to ignore
                             % the end of the input line, so these blank lines
                             % do not start a new paragraph.
                             %
with the right spacing around the periods requires
a special command.

\LaTeX\ interprets some common characters as
commands, so you must type special commands to
generate them.  These characters include the
following:
       \$ \& \% \# \{ and \}.

In printing, text is usually emphasized with an
       \emph{italic}
type style.

\begin{em}
   A long segment of text can also be emphasized
   in this way.  Text within such a segment can be
   given \emph{additional} emphasis.
\end{em}

It is sometimes necessary to prevent \LaTeX\ from
breaking a line where it might otherwise do so.
This may be at a space, as between the ``Mr.''\ and
``Jones'' in
       ``Mr.~Jones'',        % ~ produces an unbreakable interword space.
or within a word---especially when the word is a
symbol like
       \mbox{\emph{itemnum}}
that makes little sense when hyphenated across
lines.

Footnotes\footnote{This is an example of a footnote.}
pose no problem.

\LaTeX\ is good at typesetting mathematical formulas
like
       \( x-3y + z = 7 \)
or
       \( a_{1} > x^{2n} + y^{2n} > x' \)
or
       \( \ip{A}{B} = \sum_{i} a_{i} b_{i} \).
The spaces you type in a formula are
ignored.  Remember that a letter like
       $x$                   % $ ... $  and  \( ... \)  are equivalent
is a formula when it denotes a mathematical
symbol, and it should be typed as one.

\section{Displayed Text}

Text is displayed by indenting it from the left
margin.  Quotations are commonly displayed.  There
are short quotations
\begin{quote}
   This is a short quotation.  It consists of a
   single paragraph of text.  See how it is formatted.
\end{quote}
and longer ones.
\begin{quotation}
   This is a longer quotation.  It consists of two
   paragraphs of text, neither of which are
   particularly interesting.

   This is the second paragraph of the quotation.  It
   is just as dull as the first paragraph.
\end{quotation}
Another frequently-displayed structure is a list.
The following is an example of an \emph{itemized}
list.
\begin{itemize}
   \item This is the first item of an itemized list.
         Each item in the list is marked with a ``tick''.
         You don't have to worry about what kind of tick
         mark is used.

   \item This is the second item of the list.  It
         contains another list nested inside it.  The inner
         list is an \emph{enumerated} list.
         \begin{enumerate}
            \item This is the first item of an enumerated
                  list that is nested within the itemized list.

            \item This is the second item of the inner list.
                  \LaTeX\ allows you to nest lists deeper than
                  you really should.
         \end{enumerate}
         This is the rest of the second item of the outer
         list.  It is no more interesting than any other
         part of the item.
   \item This is the third item of the list.
\end{itemize}
You can even display poetry.
\begin{verse}
   There is an environment
    for verse \\             % The \\ command separates lines
   Whose features some poets % within a stanza.
   will curse.

                             % One or more blank lines separate stanzas.

   For instead of making\\
   Them do \emph{all} line breaking, \\
   It allows them to put too many words on a line when they'd rather be
   forced to be terse.
\end{verse}

Mathematical formulas may also be displayed.  A
displayed formula
is
one-line long; multiline
formulas require special formatting instructions.
   \[  \ip{\Gamma}{\psi'} = x'' + y^{2} + z_{i}^{n}\]
Don't start a paragraph with a displayed equation,
nor make one a paragraph by itself.

\end{document}               % End of document.

\section{Bench Huge Section 677}
% synthetic bench section 677
% This is a sample LaTeX input file.  (Version of 12 August 2004.)
%
% A '%' character causes TeX to ignore all remaining text on the line,
% and is used for comments like this one.

\documentclass{article}      % Specifies the document class

                             % The preamble begins here.
\title{An Example Document}  % Declares the document's title.
\author{Leslie Lamport}      % Declares the author's name.
\date{January 21, 1994}      % Deleting this command produces today's date.

\newcommand{\ip}[2]{(#1, #2)}
                             % Defines \ip{arg1}{arg2} to mean
                             % (arg1, arg2).

%\newcommand{\ip}[2]{\langle #1 | #2\rangle}
                             % This is an alternative definition of
                             % \ip that is commented out.

\begin{document}             % End of preamble and beginning of text.

\maketitle                   % Produces the title.

This is an example input file.  Comparing it with
the output it generates can show you how to
produce a simple document of your own.

\section{Ordinary Text}      % Produces section heading.  Lower-level
                             % sections are begun with similar
                             % \subsection and \subsubsection commands.

The ends  of words and sentences are marked
  by   spaces. It  doesn't matter how many
spaces    you type; one is as good as 100.  The
end of   a line counts as a space.

One   or more   blank lines denote the  end
of  a paragraph.

Since any number of consecutive spaces are treated
like a single one, the formatting of the input
file makes no difference to
      \LaTeX,                % The \LaTeX command generates the LaTeX logo.
but it makes a difference to you.  When you use
\LaTeX, making your input file as easy to read
as possible will be a great help as you write
your document and when you change it.  This sample
file shows how you can add comments to your own input
file.

Because printing is different from typewriting,
there are a number of things that you have to do
differently when preparing an input file than if
you were just typing the document directly.
Quotation marks like
       ``this''
have to be handled specially, as do quotes within
quotes:
       ``\,`this'            % \, separates the double and single quote.
        is what I just
        wrote, not  `that'\,''.

Dashes come in three sizes: an
       intra-word
dash, a medium dash for number ranges like
       1--2,
and a punctuation
       dash---like
this.

A sentence-ending space should be larger than the
space between words within a sentence.  You
sometimes have to type special commands in
conjunction with punctuation characters to get
this right, as in the following sentence.
       Gnats, gnus, etc.\ all  % `\ ' makes an inter-word space.
       begin with G\@.         % \@ marks end-of-sentence punctuation.
You should check the spaces after periods when
reading your output to make sure you haven't
forgotten any special cases.  Generating an
ellipsis
       \ldots\               % `\ ' is needed after `\ldots' because TeX
                             % ignores spaces after command names like \ldots
                             % made from \ + letters.
                             %
                             % Note how a `%' character causes TeX to ignore
                             % the end of the input line, so these blank lines
                             % do not start a new paragraph.
                             %
with the right spacing around the periods requires
a special command.

\LaTeX\ interprets some common characters as
commands, so you must type special commands to
generate them.  These characters include the
following:
       \$ \& \% \# \{ and \}.

In printing, text is usually emphasized with an
       \emph{italic}
type style.

\begin{em}
   A long segment of text can also be emphasized
   in this way.  Text within such a segment can be
   given \emph{additional} emphasis.
\end{em}

It is sometimes necessary to prevent \LaTeX\ from
breaking a line where it might otherwise do so.
This may be at a space, as between the ``Mr.''\ and
``Jones'' in
       ``Mr.~Jones'',        % ~ produces an unbreakable interword space.
or within a word---especially when the word is a
symbol like
       \mbox{\emph{itemnum}}
that makes little sense when hyphenated across
lines.

Footnotes\footnote{This is an example of a footnote.}
pose no problem.

\LaTeX\ is good at typesetting mathematical formulas
like
       \( x-3y + z = 7 \)
or
       \( a_{1} > x^{2n} + y^{2n} > x' \)
or
       \( \ip{A}{B} = \sum_{i} a_{i} b_{i} \).
The spaces you type in a formula are
ignored.  Remember that a letter like
       $x$                   % $ ... $  and  \( ... \)  are equivalent
is a formula when it denotes a mathematical
symbol, and it should be typed as one.

\section{Displayed Text}

Text is displayed by indenting it from the left
margin.  Quotations are commonly displayed.  There
are short quotations
\begin{quote}
   This is a short quotation.  It consists of a
   single paragraph of text.  See how it is formatted.
\end{quote}
and longer ones.
\begin{quotation}
   This is a longer quotation.  It consists of two
   paragraphs of text, neither of which are
   particularly interesting.

   This is the second paragraph of the quotation.  It
   is just as dull as the first paragraph.
\end{quotation}
Another frequently-displayed structure is a list.
The following is an example of an \emph{itemized}
list.
\begin{itemize}
   \item This is the first item of an itemized list.
         Each item in the list is marked with a ``tick''.
         You don't have to worry about what kind of tick
         mark is used.

   \item This is the second item of the list.  It
         contains another list nested inside it.  The inner
         list is an \emph{enumerated} list.
         \begin{enumerate}
            \item This is the first item of an enumerated
                  list that is nested within the itemized list.

            \item This is the second item of the inner list.
                  \LaTeX\ allows you to nest lists deeper than
                  you really should.
         \end{enumerate}
         This is the rest of the second item of the outer
         list.  It is no more interesting than any other
         part of the item.
   \item This is the third item of the list.
\end{itemize}
You can even display poetry.
\begin{verse}
   There is an environment
    for verse \\             % The \\ command separates lines
   Whose features some poets % within a stanza.
   will curse.

                             % One or more blank lines separate stanzas.

   For instead of making\\
   Them do \emph{all} line breaking, \\
   It allows them to put too many words on a line when they'd rather be
   forced to be terse.
\end{verse}

Mathematical formulas may also be displayed.  A
displayed formula
is
one-line long; multiline
formulas require special formatting instructions.
   \[  \ip{\Gamma}{\psi'} = x'' + y^{2} + z_{i}^{n}\]
Don't start a paragraph with a displayed equation,
nor make one a paragraph by itself.

\end{document}               % End of document.

\section{Bench Huge Section 678}
% synthetic bench section 678
% This is a sample LaTeX input file.  (Version of 12 August 2004.)
%
% A '%' character causes TeX to ignore all remaining text on the line,
% and is used for comments like this one.

\documentclass{article}      % Specifies the document class

                             % The preamble begins here.
\title{An Example Document}  % Declares the document's title.
\author{Leslie Lamport}      % Declares the author's name.
\date{January 21, 1994}      % Deleting this command produces today's date.

\newcommand{\ip}[2]{(#1, #2)}
                             % Defines \ip{arg1}{arg2} to mean
                             % (arg1, arg2).

%\newcommand{\ip}[2]{\langle #1 | #2\rangle}
                             % This is an alternative definition of
                             % \ip that is commented out.

\begin{document}             % End of preamble and beginning of text.

\maketitle                   % Produces the title.

This is an example input file.  Comparing it with
the output it generates can show you how to
produce a simple document of your own.

\section{Ordinary Text}      % Produces section heading.  Lower-level
                             % sections are begun with similar
                             % \subsection and \subsubsection commands.

The ends  of words and sentences are marked
  by   spaces. It  doesn't matter how many
spaces    you type; one is as good as 100.  The
end of   a line counts as a space.

One   or more   blank lines denote the  end
of  a paragraph.

Since any number of consecutive spaces are treated
like a single one, the formatting of the input
file makes no difference to
      \LaTeX,                % The \LaTeX command generates the LaTeX logo.
but it makes a difference to you.  When you use
\LaTeX, making your input file as easy to read
as possible will be a great help as you write
your document and when you change it.  This sample
file shows how you can add comments to your own input
file.

Because printing is different from typewriting,
there are a number of things that you have to do
differently when preparing an input file than if
you were just typing the document directly.
Quotation marks like
       ``this''
have to be handled specially, as do quotes within
quotes:
       ``\,`this'            % \, separates the double and single quote.
        is what I just
        wrote, not  `that'\,''.

Dashes come in three sizes: an
       intra-word
dash, a medium dash for number ranges like
       1--2,
and a punctuation
       dash---like
this.

A sentence-ending space should be larger than the
space between words within a sentence.  You
sometimes have to type special commands in
conjunction with punctuation characters to get
this right, as in the following sentence.
       Gnats, gnus, etc.\ all  % `\ ' makes an inter-word space.
       begin with G\@.         % \@ marks end-of-sentence punctuation.
You should check the spaces after periods when
reading your output to make sure you haven't
forgotten any special cases.  Generating an
ellipsis
       \ldots\               % `\ ' is needed after `\ldots' because TeX
                             % ignores spaces after command names like \ldots
                             % made from \ + letters.
                             %
                             % Note how a `%' character causes TeX to ignore
                             % the end of the input line, so these blank lines
                             % do not start a new paragraph.
                             %
with the right spacing around the periods requires
a special command.

\LaTeX\ interprets some common characters as
commands, so you must type special commands to
generate them.  These characters include the
following:
       \$ \& \% \# \{ and \}.

In printing, text is usually emphasized with an
       \emph{italic}
type style.

\begin{em}
   A long segment of text can also be emphasized
   in this way.  Text within such a segment can be
   given \emph{additional} emphasis.
\end{em}

It is sometimes necessary to prevent \LaTeX\ from
breaking a line where it might otherwise do so.
This may be at a space, as between the ``Mr.''\ and
``Jones'' in
       ``Mr.~Jones'',        % ~ produces an unbreakable interword space.
or within a word---especially when the word is a
symbol like
       \mbox{\emph{itemnum}}
that makes little sense when hyphenated across
lines.

Footnotes\footnote{This is an example of a footnote.}
pose no problem.

\LaTeX\ is good at typesetting mathematical formulas
like
       \( x-3y + z = 7 \)
or
       \( a_{1} > x^{2n} + y^{2n} > x' \)
or
       \( \ip{A}{B} = \sum_{i} a_{i} b_{i} \).
The spaces you type in a formula are
ignored.  Remember that a letter like
       $x$                   % $ ... $  and  \( ... \)  are equivalent
is a formula when it denotes a mathematical
symbol, and it should be typed as one.

\section{Displayed Text}

Text is displayed by indenting it from the left
margin.  Quotations are commonly displayed.  There
are short quotations
\begin{quote}
   This is a short quotation.  It consists of a
   single paragraph of text.  See how it is formatted.
\end{quote}
and longer ones.
\begin{quotation}
   This is a longer quotation.  It consists of two
   paragraphs of text, neither of which are
   particularly interesting.

   This is the second paragraph of the quotation.  It
   is just as dull as the first paragraph.
\end{quotation}
Another frequently-displayed structure is a list.
The following is an example of an \emph{itemized}
list.
\begin{itemize}
   \item This is the first item of an itemized list.
         Each item in the list is marked with a ``tick''.
         You don't have to worry about what kind of tick
         mark is used.

   \item This is the second item of the list.  It
         contains another list nested inside it.  The inner
         list is an \emph{enumerated} list.
         \begin{enumerate}
            \item This is the first item of an enumerated
                  list that is nested within the itemized list.

            \item This is the second item of the inner list.
                  \LaTeX\ allows you to nest lists deeper than
                  you really should.
         \end{enumerate}
         This is the rest of the second item of the outer
         list.  It is no more interesting than any other
         part of the item.
   \item This is the third item of the list.
\end{itemize}
You can even display poetry.
\begin{verse}
   There is an environment
    for verse \\             % The \\ command separates lines
   Whose features some poets % within a stanza.
   will curse.

                             % One or more blank lines separate stanzas.

   For instead of making\\
   Them do \emph{all} line breaking, \\
   It allows them to put too many words on a line when they'd rather be
   forced to be terse.
\end{verse}

Mathematical formulas may also be displayed.  A
displayed formula
is
one-line long; multiline
formulas require special formatting instructions.
   \[  \ip{\Gamma}{\psi'} = x'' + y^{2} + z_{i}^{n}\]
Don't start a paragraph with a displayed equation,
nor make one a paragraph by itself.

\end{document}               % End of document.

\section{Bench Huge Section 679}
% synthetic bench section 679
% This is a sample LaTeX input file.  (Version of 12 August 2004.)
%
% A '%' character causes TeX to ignore all remaining text on the line,
% and is used for comments like this one.

\documentclass{article}      % Specifies the document class

                             % The preamble begins here.
\title{An Example Document}  % Declares the document's title.
\author{Leslie Lamport}      % Declares the author's name.
\date{January 21, 1994}      % Deleting this command produces today's date.

\newcommand{\ip}[2]{(#1, #2)}
                             % Defines \ip{arg1}{arg2} to mean
                             % (arg1, arg2).

%\newcommand{\ip}[2]{\langle #1 | #2\rangle}
                             % This is an alternative definition of
                             % \ip that is commented out.

\begin{document}             % End of preamble and beginning of text.

\maketitle                   % Produces the title.

This is an example input file.  Comparing it with
the output it generates can show you how to
produce a simple document of your own.

\section{Ordinary Text}      % Produces section heading.  Lower-level
                             % sections are begun with similar
                             % \subsection and \subsubsection commands.

The ends  of words and sentences are marked
  by   spaces. It  doesn't matter how many
spaces    you type; one is as good as 100.  The
end of   a line counts as a space.

One   or more   blank lines denote the  end
of  a paragraph.

Since any number of consecutive spaces are treated
like a single one, the formatting of the input
file makes no difference to
      \LaTeX,                % The \LaTeX command generates the LaTeX logo.
but it makes a difference to you.  When you use
\LaTeX, making your input file as easy to read
as possible will be a great help as you write
your document and when you change it.  This sample
file shows how you can add comments to your own input
file.

Because printing is different from typewriting,
there are a number of things that you have to do
differently when preparing an input file than if
you were just typing the document directly.
Quotation marks like
       ``this''
have to be handled specially, as do quotes within
quotes:
       ``\,`this'            % \, separates the double and single quote.
        is what I just
        wrote, not  `that'\,''.

Dashes come in three sizes: an
       intra-word
dash, a medium dash for number ranges like
       1--2,
and a punctuation
       dash---like
this.

A sentence-ending space should be larger than the
space between words within a sentence.  You
sometimes have to type special commands in
conjunction with punctuation characters to get
this right, as in the following sentence.
       Gnats, gnus, etc.\ all  % `\ ' makes an inter-word space.
       begin with G\@.         % \@ marks end-of-sentence punctuation.
You should check the spaces after periods when
reading your output to make sure you haven't
forgotten any special cases.  Generating an
ellipsis
       \ldots\               % `\ ' is needed after `\ldots' because TeX
                             % ignores spaces after command names like \ldots
                             % made from \ + letters.
                             %
                             % Note how a `%' character causes TeX to ignore
                             % the end of the input line, so these blank lines
                             % do not start a new paragraph.
                             %
with the right spacing around the periods requires
a special command.

\LaTeX\ interprets some common characters as
commands, so you must type special commands to
generate them.  These characters include the
following:
       \$ \& \% \# \{ and \}.

In printing, text is usually emphasized with an
       \emph{italic}
type style.

\begin{em}
   A long segment of text can also be emphasized
   in this way.  Text within such a segment can be
   given \emph{additional} emphasis.
\end{em}

It is sometimes necessary to prevent \LaTeX\ from
breaking a line where it might otherwise do so.
This may be at a space, as between the ``Mr.''\ and
``Jones'' in
       ``Mr.~Jones'',        % ~ produces an unbreakable interword space.
or within a word---especially when the word is a
symbol like
       \mbox{\emph{itemnum}}
that makes little sense when hyphenated across
lines.

Footnotes\footnote{This is an example of a footnote.}
pose no problem.

\LaTeX\ is good at typesetting mathematical formulas
like
       \( x-3y + z = 7 \)
or
       \( a_{1} > x^{2n} + y^{2n} > x' \)
or
       \( \ip{A}{B} = \sum_{i} a_{i} b_{i} \).
The spaces you type in a formula are
ignored.  Remember that a letter like
       $x$                   % $ ... $  and  \( ... \)  are equivalent
is a formula when it denotes a mathematical
symbol, and it should be typed as one.

\section{Displayed Text}

Text is displayed by indenting it from the left
margin.  Quotations are commonly displayed.  There
are short quotations
\begin{quote}
   This is a short quotation.  It consists of a
   single paragraph of text.  See how it is formatted.
\end{quote}
and longer ones.
\begin{quotation}
   This is a longer quotation.  It consists of two
   paragraphs of text, neither of which are
   particularly interesting.

   This is the second paragraph of the quotation.  It
   is just as dull as the first paragraph.
\end{quotation}
Another frequently-displayed structure is a list.
The following is an example of an \emph{itemized}
list.
\begin{itemize}
   \item This is the first item of an itemized list.
         Each item in the list is marked with a ``tick''.
         You don't have to worry about what kind of tick
         mark is used.

   \item This is the second item of the list.  It
         contains another list nested inside it.  The inner
         list is an \emph{enumerated} list.
         \begin{enumerate}
            \item This is the first item of an enumerated
                  list that is nested within the itemized list.

            \item This is the second item of the inner list.
                  \LaTeX\ allows you to nest lists deeper than
                  you really should.
         \end{enumerate}
         This is the rest of the second item of the outer
         list.  It is no more interesting than any other
         part of the item.
   \item This is the third item of the list.
\end{itemize}
You can even display poetry.
\begin{verse}
   There is an environment
    for verse \\             % The \\ command separates lines
   Whose features some poets % within a stanza.
   will curse.

                             % One or more blank lines separate stanzas.

   For instead of making\\
   Them do \emph{all} line breaking, \\
   It allows them to put too many words on a line when they'd rather be
   forced to be terse.
\end{verse}

Mathematical formulas may also be displayed.  A
displayed formula
is
one-line long; multiline
formulas require special formatting instructions.
   \[  \ip{\Gamma}{\psi'} = x'' + y^{2} + z_{i}^{n}\]
Don't start a paragraph with a displayed equation,
nor make one a paragraph by itself.

\end{document}               % End of document.

\section{Bench Huge Section 680}
% synthetic bench section 680
% This is a sample LaTeX input file.  (Version of 12 August 2004.)
%
% A '%' character causes TeX to ignore all remaining text on the line,
% and is used for comments like this one.

\documentclass{article}      % Specifies the document class

                             % The preamble begins here.
\title{An Example Document}  % Declares the document's title.
\author{Leslie Lamport}      % Declares the author's name.
\date{January 21, 1994}      % Deleting this command produces today's date.

\newcommand{\ip}[2]{(#1, #2)}
                             % Defines \ip{arg1}{arg2} to mean
                             % (arg1, arg2).

%\newcommand{\ip}[2]{\langle #1 | #2\rangle}
                             % This is an alternative definition of
                             % \ip that is commented out.

\begin{document}             % End of preamble and beginning of text.

\maketitle                   % Produces the title.

This is an example input file.  Comparing it with
the output it generates can show you how to
produce a simple document of your own.

\section{Ordinary Text}      % Produces section heading.  Lower-level
                             % sections are begun with similar
                             % \subsection and \subsubsection commands.

The ends  of words and sentences are marked
  by   spaces. It  doesn't matter how many
spaces    you type; one is as good as 100.  The
end of   a line counts as a space.

One   or more   blank lines denote the  end
of  a paragraph.

Since any number of consecutive spaces are treated
like a single one, the formatting of the input
file makes no difference to
      \LaTeX,                % The \LaTeX command generates the LaTeX logo.
but it makes a difference to you.  When you use
\LaTeX, making your input file as easy to read
as possible will be a great help as you write
your document and when you change it.  This sample
file shows how you can add comments to your own input
file.

Because printing is different from typewriting,
there are a number of things that you have to do
differently when preparing an input file than if
you were just typing the document directly.
Quotation marks like
       ``this''
have to be handled specially, as do quotes within
quotes:
       ``\,`this'            % \, separates the double and single quote.
        is what I just
        wrote, not  `that'\,''.

Dashes come in three sizes: an
       intra-word
dash, a medium dash for number ranges like
       1--2,
and a punctuation
       dash---like
this.

A sentence-ending space should be larger than the
space between words within a sentence.  You
sometimes have to type special commands in
conjunction with punctuation characters to get
this right, as in the following sentence.
       Gnats, gnus, etc.\ all  % `\ ' makes an inter-word space.
       begin with G\@.         % \@ marks end-of-sentence punctuation.
You should check the spaces after periods when
reading your output to make sure you haven't
forgotten any special cases.  Generating an
ellipsis
       \ldots\               % `\ ' is needed after `\ldots' because TeX
                             % ignores spaces after command names like \ldots
                             % made from \ + letters.
                             %
                             % Note how a `%' character causes TeX to ignore
                             % the end of the input line, so these blank lines
                             % do not start a new paragraph.
                             %
with the right spacing around the periods requires
a special command.

\LaTeX\ interprets some common characters as
commands, so you must type special commands to
generate them.  These characters include the
following:
       \$ \& \% \# \{ and \}.

In printing, text is usually emphasized with an
       \emph{italic}
type style.

\begin{em}
   A long segment of text can also be emphasized
   in this way.  Text within such a segment can be
   given \emph{additional} emphasis.
\end{em}

It is sometimes necessary to prevent \LaTeX\ from
breaking a line where it might otherwise do so.
This may be at a space, as between the ``Mr.''\ and
``Jones'' in
       ``Mr.~Jones'',        % ~ produces an unbreakable interword space.
or within a word---especially when the word is a
symbol like
       \mbox{\emph{itemnum}}
that makes little sense when hyphenated across
lines.

Footnotes\footnote{This is an example of a footnote.}
pose no problem.

\LaTeX\ is good at typesetting mathematical formulas
like
       \( x-3y + z = 7 \)
or
       \( a_{1} > x^{2n} + y^{2n} > x' \)
or
       \( \ip{A}{B} = \sum_{i} a_{i} b_{i} \).
The spaces you type in a formula are
ignored.  Remember that a letter like
       $x$                   % $ ... $  and  \( ... \)  are equivalent
is a formula when it denotes a mathematical
symbol, and it should be typed as one.

\section{Displayed Text}

Text is displayed by indenting it from the left
margin.  Quotations are commonly displayed.  There
are short quotations
\begin{quote}
   This is a short quotation.  It consists of a
   single paragraph of text.  See how it is formatted.
\end{quote}
and longer ones.
\begin{quotation}
   This is a longer quotation.  It consists of two
   paragraphs of text, neither of which are
   particularly interesting.

   This is the second paragraph of the quotation.  It
   is just as dull as the first paragraph.
\end{quotation}
Another frequently-displayed structure is a list.
The following is an example of an \emph{itemized}
list.
\begin{itemize}
   \item This is the first item of an itemized list.
         Each item in the list is marked with a ``tick''.
         You don't have to worry about what kind of tick
         mark is used.

   \item This is the second item of the list.  It
         contains another list nested inside it.  The inner
         list is an \emph{enumerated} list.
         \begin{enumerate}
            \item This is the first item of an enumerated
                  list that is nested within the itemized list.

            \item This is the second item of the inner list.
                  \LaTeX\ allows you to nest lists deeper than
                  you really should.
         \end{enumerate}
         This is the rest of the second item of the outer
         list.  It is no more interesting than any other
         part of the item.
   \item This is the third item of the list.
\end{itemize}
You can even display poetry.
\begin{verse}
   There is an environment
    for verse \\             % The \\ command separates lines
   Whose features some poets % within a stanza.
   will curse.

                             % One or more blank lines separate stanzas.

   For instead of making\\
   Them do \emph{all} line breaking, \\
   It allows them to put too many words on a line when they'd rather be
   forced to be terse.
\end{verse}

Mathematical formulas may also be displayed.  A
displayed formula
is
one-line long; multiline
formulas require special formatting instructions.
   \[  \ip{\Gamma}{\psi'} = x'' + y^{2} + z_{i}^{n}\]
Don't start a paragraph with a displayed equation,
nor make one a paragraph by itself.

\end{document}               % End of document.

\section{Bench Huge Section 681}
% synthetic bench section 681
% This is a sample LaTeX input file.  (Version of 12 August 2004.)
%
% A '%' character causes TeX to ignore all remaining text on the line,
% and is used for comments like this one.

\documentclass{article}      % Specifies the document class

                             % The preamble begins here.
\title{An Example Document}  % Declares the document's title.
\author{Leslie Lamport}      % Declares the author's name.
\date{January 21, 1994}      % Deleting this command produces today's date.

\newcommand{\ip}[2]{(#1, #2)}
                             % Defines \ip{arg1}{arg2} to mean
                             % (arg1, arg2).

%\newcommand{\ip}[2]{\langle #1 | #2\rangle}
                             % This is an alternative definition of
                             % \ip that is commented out.

\begin{document}             % End of preamble and beginning of text.

\maketitle                   % Produces the title.

This is an example input file.  Comparing it with
the output it generates can show you how to
produce a simple document of your own.

\section{Ordinary Text}      % Produces section heading.  Lower-level
                             % sections are begun with similar
                             % \subsection and \subsubsection commands.

The ends  of words and sentences are marked
  by   spaces. It  doesn't matter how many
spaces    you type; one is as good as 100.  The
end of   a line counts as a space.

One   or more   blank lines denote the  end
of  a paragraph.

Since any number of consecutive spaces are treated
like a single one, the formatting of the input
file makes no difference to
      \LaTeX,                % The \LaTeX command generates the LaTeX logo.
but it makes a difference to you.  When you use
\LaTeX, making your input file as easy to read
as possible will be a great help as you write
your document and when you change it.  This sample
file shows how you can add comments to your own input
file.

Because printing is different from typewriting,
there are a number of things that you have to do
differently when preparing an input file than if
you were just typing the document directly.
Quotation marks like
       ``this''
have to be handled specially, as do quotes within
quotes:
       ``\,`this'            % \, separates the double and single quote.
        is what I just
        wrote, not  `that'\,''.

Dashes come in three sizes: an
       intra-word
dash, a medium dash for number ranges like
       1--2,
and a punctuation
       dash---like
this.

A sentence-ending space should be larger than the
space between words within a sentence.  You
sometimes have to type special commands in
conjunction with punctuation characters to get
this right, as in the following sentence.
       Gnats, gnus, etc.\ all  % `\ ' makes an inter-word space.
       begin with G\@.         % \@ marks end-of-sentence punctuation.
You should check the spaces after periods when
reading your output to make sure you haven't
forgotten any special cases.  Generating an
ellipsis
       \ldots\               % `\ ' is needed after `\ldots' because TeX
                             % ignores spaces after command names like \ldots
                             % made from \ + letters.
                             %
                             % Note how a `%' character causes TeX to ignore
                             % the end of the input line, so these blank lines
                             % do not start a new paragraph.
                             %
with the right spacing around the periods requires
a special command.

\LaTeX\ interprets some common characters as
commands, so you must type special commands to
generate them.  These characters include the
following:
       \$ \& \% \# \{ and \}.

In printing, text is usually emphasized with an
       \emph{italic}
type style.

\begin{em}
   A long segment of text can also be emphasized
   in this way.  Text within such a segment can be
   given \emph{additional} emphasis.
\end{em}

It is sometimes necessary to prevent \LaTeX\ from
breaking a line where it might otherwise do so.
This may be at a space, as between the ``Mr.''\ and
``Jones'' in
       ``Mr.~Jones'',        % ~ produces an unbreakable interword space.
or within a word---especially when the word is a
symbol like
       \mbox{\emph{itemnum}}
that makes little sense when hyphenated across
lines.

Footnotes\footnote{This is an example of a footnote.}
pose no problem.

\LaTeX\ is good at typesetting mathematical formulas
like
       \( x-3y + z = 7 \)
or
       \( a_{1} > x^{2n} + y^{2n} > x' \)
or
       \( \ip{A}{B} = \sum_{i} a_{i} b_{i} \).
The spaces you type in a formula are
ignored.  Remember that a letter like
       $x$                   % $ ... $  and  \( ... \)  are equivalent
is a formula when it denotes a mathematical
symbol, and it should be typed as one.

\section{Displayed Text}

Text is displayed by indenting it from the left
margin.  Quotations are commonly displayed.  There
are short quotations
\begin{quote}
   This is a short quotation.  It consists of a
   single paragraph of text.  See how it is formatted.
\end{quote}
and longer ones.
\begin{quotation}
   This is a longer quotation.  It consists of two
   paragraphs of text, neither of which are
   particularly interesting.

   This is the second paragraph of the quotation.  It
   is just as dull as the first paragraph.
\end{quotation}
Another frequently-displayed structure is a list.
The following is an example of an \emph{itemized}
list.
\begin{itemize}
   \item This is the first item of an itemized list.
         Each item in the list is marked with a ``tick''.
         You don't have to worry about what kind of tick
         mark is used.

   \item This is the second item of the list.  It
         contains another list nested inside it.  The inner
         list is an \emph{enumerated} list.
         \begin{enumerate}
            \item This is the first item of an enumerated
                  list that is nested within the itemized list.

            \item This is the second item of the inner list.
                  \LaTeX\ allows you to nest lists deeper than
                  you really should.
         \end{enumerate}
         This is the rest of the second item of the outer
         list.  It is no more interesting than any other
         part of the item.
   \item This is the third item of the list.
\end{itemize}
You can even display poetry.
\begin{verse}
   There is an environment
    for verse \\             % The \\ command separates lines
   Whose features some poets % within a stanza.
   will curse.

                             % One or more blank lines separate stanzas.

   For instead of making\\
   Them do \emph{all} line breaking, \\
   It allows them to put too many words on a line when they'd rather be
   forced to be terse.
\end{verse}

Mathematical formulas may also be displayed.  A
displayed formula
is
one-line long; multiline
formulas require special formatting instructions.
   \[  \ip{\Gamma}{\psi'} = x'' + y^{2} + z_{i}^{n}\]
Don't start a paragraph with a displayed equation,
nor make one a paragraph by itself.

\end{document}               % End of document.

\section{Bench Huge Section 682}
% synthetic bench section 682
% This is a sample LaTeX input file.  (Version of 12 August 2004.)
%
% A '%' character causes TeX to ignore all remaining text on the line,
% and is used for comments like this one.

\documentclass{article}      % Specifies the document class

                             % The preamble begins here.
\title{An Example Document}  % Declares the document's title.
\author{Leslie Lamport}      % Declares the author's name.
\date{January 21, 1994}      % Deleting this command produces today's date.

\newcommand{\ip}[2]{(#1, #2)}
                             % Defines \ip{arg1}{arg2} to mean
                             % (arg1, arg2).

%\newcommand{\ip}[2]{\langle #1 | #2\rangle}
                             % This is an alternative definition of
                             % \ip that is commented out.

\begin{document}             % End of preamble and beginning of text.

\maketitle                   % Produces the title.

This is an example input file.  Comparing it with
the output it generates can show you how to
produce a simple document of your own.

\section{Ordinary Text}      % Produces section heading.  Lower-level
                             % sections are begun with similar
                             % \subsection and \subsubsection commands.

The ends  of words and sentences are marked
  by   spaces. It  doesn't matter how many
spaces    you type; one is as good as 100.  The
end of   a line counts as a space.

One   or more   blank lines denote the  end
of  a paragraph.

Since any number of consecutive spaces are treated
like a single one, the formatting of the input
file makes no difference to
      \LaTeX,                % The \LaTeX command generates the LaTeX logo.
but it makes a difference to you.  When you use
\LaTeX, making your input file as easy to read
as possible will be a great help as you write
your document and when you change it.  This sample
file shows how you can add comments to your own input
file.

Because printing is different from typewriting,
there are a number of things that you have to do
differently when preparing an input file than if
you were just typing the document directly.
Quotation marks like
       ``this''
have to be handled specially, as do quotes within
quotes:
       ``\,`this'            % \, separates the double and single quote.
        is what I just
        wrote, not  `that'\,''.

Dashes come in three sizes: an
       intra-word
dash, a medium dash for number ranges like
       1--2,
and a punctuation
       dash---like
this.

A sentence-ending space should be larger than the
space between words within a sentence.  You
sometimes have to type special commands in
conjunction with punctuation characters to get
this right, as in the following sentence.
       Gnats, gnus, etc.\ all  % `\ ' makes an inter-word space.
       begin with G\@.         % \@ marks end-of-sentence punctuation.
You should check the spaces after periods when
reading your output to make sure you haven't
forgotten any special cases.  Generating an
ellipsis
       \ldots\               % `\ ' is needed after `\ldots' because TeX
                             % ignores spaces after command names like \ldots
                             % made from \ + letters.
                             %
                             % Note how a `%' character causes TeX to ignore
                             % the end of the input line, so these blank lines
                             % do not start a new paragraph.
                             %
with the right spacing around the periods requires
a special command.

\LaTeX\ interprets some common characters as
commands, so you must type special commands to
generate them.  These characters include the
following:
       \$ \& \% \# \{ and \}.

In printing, text is usually emphasized with an
       \emph{italic}
type style.

\begin{em}
   A long segment of text can also be emphasized
   in this way.  Text within such a segment can be
   given \emph{additional} emphasis.
\end{em}

It is sometimes necessary to prevent \LaTeX\ from
breaking a line where it might otherwise do so.
This may be at a space, as between the ``Mr.''\ and
``Jones'' in
       ``Mr.~Jones'',        % ~ produces an unbreakable interword space.
or within a word---especially when the word is a
symbol like
       \mbox{\emph{itemnum}}
that makes little sense when hyphenated across
lines.

Footnotes\footnote{This is an example of a footnote.}
pose no problem.

\LaTeX\ is good at typesetting mathematical formulas
like
       \( x-3y + z = 7 \)
or
       \( a_{1} > x^{2n} + y^{2n} > x' \)
or
       \( \ip{A}{B} = \sum_{i} a_{i} b_{i} \).
The spaces you type in a formula are
ignored.  Remember that a letter like
       $x$                   % $ ... $  and  \( ... \)  are equivalent
is a formula when it denotes a mathematical
symbol, and it should be typed as one.

\section{Displayed Text}

Text is displayed by indenting it from the left
margin.  Quotations are commonly displayed.  There
are short quotations
\begin{quote}
   This is a short quotation.  It consists of a
   single paragraph of text.  See how it is formatted.
\end{quote}
and longer ones.
\begin{quotation}
   This is a longer quotation.  It consists of two
   paragraphs of text, neither of which are
   particularly interesting.

   This is the second paragraph of the quotation.  It
   is just as dull as the first paragraph.
\end{quotation}
Another frequently-displayed structure is a list.
The following is an example of an \emph{itemized}
list.
\begin{itemize}
   \item This is the first item of an itemized list.
         Each item in the list is marked with a ``tick''.
         You don't have to worry about what kind of tick
         mark is used.

   \item This is the second item of the list.  It
         contains another list nested inside it.  The inner
         list is an \emph{enumerated} list.
         \begin{enumerate}
            \item This is the first item of an enumerated
                  list that is nested within the itemized list.

            \item This is the second item of the inner list.
                  \LaTeX\ allows you to nest lists deeper than
                  you really should.
         \end{enumerate}
         This is the rest of the second item of the outer
         list.  It is no more interesting than any other
         part of the item.
   \item This is the third item of the list.
\end{itemize}
You can even display poetry.
\begin{verse}
   There is an environment
    for verse \\             % The \\ command separates lines
   Whose features some poets % within a stanza.
   will curse.

                             % One or more blank lines separate stanzas.

   For instead of making\\
   Them do \emph{all} line breaking, \\
   It allows them to put too many words on a line when they'd rather be
   forced to be terse.
\end{verse}

Mathematical formulas may also be displayed.  A
displayed formula
is
one-line long; multiline
formulas require special formatting instructions.
   \[  \ip{\Gamma}{\psi'} = x'' + y^{2} + z_{i}^{n}\]
Don't start a paragraph with a displayed equation,
nor make one a paragraph by itself.

\end{document}               % End of document.

\section{Bench Huge Section 683}
% synthetic bench section 683
% This is a sample LaTeX input file.  (Version of 12 August 2004.)
%
% A '%' character causes TeX to ignore all remaining text on the line,
% and is used for comments like this one.

\documentclass{article}      % Specifies the document class

                             % The preamble begins here.
\title{An Example Document}  % Declares the document's title.
\author{Leslie Lamport}      % Declares the author's name.
\date{January 21, 1994}      % Deleting this command produces today's date.

\newcommand{\ip}[2]{(#1, #2)}
                             % Defines \ip{arg1}{arg2} to mean
                             % (arg1, arg2).

%\newcommand{\ip}[2]{\langle #1 | #2\rangle}
                             % This is an alternative definition of
                             % \ip that is commented out.

\begin{document}             % End of preamble and beginning of text.

\maketitle                   % Produces the title.

This is an example input file.  Comparing it with
the output it generates can show you how to
produce a simple document of your own.

\section{Ordinary Text}      % Produces section heading.  Lower-level
                             % sections are begun with similar
                             % \subsection and \subsubsection commands.

The ends  of words and sentences are marked
  by   spaces. It  doesn't matter how many
spaces    you type; one is as good as 100.  The
end of   a line counts as a space.

One   or more   blank lines denote the  end
of  a paragraph.

Since any number of consecutive spaces are treated
like a single one, the formatting of the input
file makes no difference to
      \LaTeX,                % The \LaTeX command generates the LaTeX logo.
but it makes a difference to you.  When you use
\LaTeX, making your input file as easy to read
as possible will be a great help as you write
your document and when you change it.  This sample
file shows how you can add comments to your own input
file.

Because printing is different from typewriting,
there are a number of things that you have to do
differently when preparing an input file than if
you were just typing the document directly.
Quotation marks like
       ``this''
have to be handled specially, as do quotes within
quotes:
       ``\,`this'            % \, separates the double and single quote.
        is what I just
        wrote, not  `that'\,''.

Dashes come in three sizes: an
       intra-word
dash, a medium dash for number ranges like
       1--2,
and a punctuation
       dash---like
this.

A sentence-ending space should be larger than the
space between words within a sentence.  You
sometimes have to type special commands in
conjunction with punctuation characters to get
this right, as in the following sentence.
       Gnats, gnus, etc.\ all  % `\ ' makes an inter-word space.
       begin with G\@.         % \@ marks end-of-sentence punctuation.
You should check the spaces after periods when
reading your output to make sure you haven't
forgotten any special cases.  Generating an
ellipsis
       \ldots\               % `\ ' is needed after `\ldots' because TeX
                             % ignores spaces after command names like \ldots
                             % made from \ + letters.
                             %
                             % Note how a `%' character causes TeX to ignore
                             % the end of the input line, so these blank lines
                             % do not start a new paragraph.
                             %
with the right spacing around the periods requires
a special command.

\LaTeX\ interprets some common characters as
commands, so you must type special commands to
generate them.  These characters include the
following:
       \$ \& \% \# \{ and \}.

In printing, text is usually emphasized with an
       \emph{italic}
type style.

\begin{em}
   A long segment of text can also be emphasized
   in this way.  Text within such a segment can be
   given \emph{additional} emphasis.
\end{em}

It is sometimes necessary to prevent \LaTeX\ from
breaking a line where it might otherwise do so.
This may be at a space, as between the ``Mr.''\ and
``Jones'' in
       ``Mr.~Jones'',        % ~ produces an unbreakable interword space.
or within a word---especially when the word is a
symbol like
       \mbox{\emph{itemnum}}
that makes little sense when hyphenated across
lines.

Footnotes\footnote{This is an example of a footnote.}
pose no problem.

\LaTeX\ is good at typesetting mathematical formulas
like
       \( x-3y + z = 7 \)
or
       \( a_{1} > x^{2n} + y^{2n} > x' \)
or
       \( \ip{A}{B} = \sum_{i} a_{i} b_{i} \).
The spaces you type in a formula are
ignored.  Remember that a letter like
       $x$                   % $ ... $  and  \( ... \)  are equivalent
is a formula when it denotes a mathematical
symbol, and it should be typed as one.

\section{Displayed Text}

Text is displayed by indenting it from the left
margin.  Quotations are commonly displayed.  There
are short quotations
\begin{quote}
   This is a short quotation.  It consists of a
   single paragraph of text.  See how it is formatted.
\end{quote}
and longer ones.
\begin{quotation}
   This is a longer quotation.  It consists of two
   paragraphs of text, neither of which are
   particularly interesting.

   This is the second paragraph of the quotation.  It
   is just as dull as the first paragraph.
\end{quotation}
Another frequently-displayed structure is a list.
The following is an example of an \emph{itemized}
list.
\begin{itemize}
   \item This is the first item of an itemized list.
         Each item in the list is marked with a ``tick''.
         You don't have to worry about what kind of tick
         mark is used.

   \item This is the second item of the list.  It
         contains another list nested inside it.  The inner
         list is an \emph{enumerated} list.
         \begin{enumerate}
            \item This is the first item of an enumerated
                  list that is nested within the itemized list.

            \item This is the second item of the inner list.
                  \LaTeX\ allows you to nest lists deeper than
                  you really should.
         \end{enumerate}
         This is the rest of the second item of the outer
         list.  It is no more interesting than any other
         part of the item.
   \item This is the third item of the list.
\end{itemize}
You can even display poetry.
\begin{verse}
   There is an environment
    for verse \\             % The \\ command separates lines
   Whose features some poets % within a stanza.
   will curse.

                             % One or more blank lines separate stanzas.

   For instead of making\\
   Them do \emph{all} line breaking, \\
   It allows them to put too many words on a line when they'd rather be
   forced to be terse.
\end{verse}

Mathematical formulas may also be displayed.  A
displayed formula
is
one-line long; multiline
formulas require special formatting instructions.
   \[  \ip{\Gamma}{\psi'} = x'' + y^{2} + z_{i}^{n}\]
Don't start a paragraph with a displayed equation,
nor make one a paragraph by itself.

\end{document}               % End of document.

\section{Bench Huge Section 684}
% synthetic bench section 684
% This is a sample LaTeX input file.  (Version of 12 August 2004.)
%
% A '%' character causes TeX to ignore all remaining text on the line,
% and is used for comments like this one.

\documentclass{article}      % Specifies the document class

                             % The preamble begins here.
\title{An Example Document}  % Declares the document's title.
\author{Leslie Lamport}      % Declares the author's name.
\date{January 21, 1994}      % Deleting this command produces today's date.

\newcommand{\ip}[2]{(#1, #2)}
                             % Defines \ip{arg1}{arg2} to mean
                             % (arg1, arg2).

%\newcommand{\ip}[2]{\langle #1 | #2\rangle}
                             % This is an alternative definition of
                             % \ip that is commented out.

\begin{document}             % End of preamble and beginning of text.

\maketitle                   % Produces the title.

This is an example input file.  Comparing it with
the output it generates can show you how to
produce a simple document of your own.

\section{Ordinary Text}      % Produces section heading.  Lower-level
                             % sections are begun with similar
                             % \subsection and \subsubsection commands.

The ends  of words and sentences are marked
  by   spaces. It  doesn't matter how many
spaces    you type; one is as good as 100.  The
end of   a line counts as a space.

One   or more   blank lines denote the  end
of  a paragraph.

Since any number of consecutive spaces are treated
like a single one, the formatting of the input
file makes no difference to
      \LaTeX,                % The \LaTeX command generates the LaTeX logo.
but it makes a difference to you.  When you use
\LaTeX, making your input file as easy to read
as possible will be a great help as you write
your document and when you change it.  This sample
file shows how you can add comments to your own input
file.

Because printing is different from typewriting,
there are a number of things that you have to do
differently when preparing an input file than if
you were just typing the document directly.
Quotation marks like
       ``this''
have to be handled specially, as do quotes within
quotes:
       ``\,`this'            % \, separates the double and single quote.
        is what I just
        wrote, not  `that'\,''.

Dashes come in three sizes: an
       intra-word
dash, a medium dash for number ranges like
       1--2,
and a punctuation
       dash---like
this.

A sentence-ending space should be larger than the
space between words within a sentence.  You
sometimes have to type special commands in
conjunction with punctuation characters to get
this right, as in the following sentence.
       Gnats, gnus, etc.\ all  % `\ ' makes an inter-word space.
       begin with G\@.         % \@ marks end-of-sentence punctuation.
You should check the spaces after periods when
reading your output to make sure you haven't
forgotten any special cases.  Generating an
ellipsis
       \ldots\               % `\ ' is needed after `\ldots' because TeX
                             % ignores spaces after command names like \ldots
                             % made from \ + letters.
                             %
                             % Note how a `%' character causes TeX to ignore
                             % the end of the input line, so these blank lines
                             % do not start a new paragraph.
                             %
with the right spacing around the periods requires
a special command.

\LaTeX\ interprets some common characters as
commands, so you must type special commands to
generate them.  These characters include the
following:
       \$ \& \% \# \{ and \}.

In printing, text is usually emphasized with an
       \emph{italic}
type style.

\begin{em}
   A long segment of text can also be emphasized
   in this way.  Text within such a segment can be
   given \emph{additional} emphasis.
\end{em}

It is sometimes necessary to prevent \LaTeX\ from
breaking a line where it might otherwise do so.
This may be at a space, as between the ``Mr.''\ and
``Jones'' in
       ``Mr.~Jones'',        % ~ produces an unbreakable interword space.
or within a word---especially when the word is a
symbol like
       \mbox{\emph{itemnum}}
that makes little sense when hyphenated across
lines.

Footnotes\footnote{This is an example of a footnote.}
pose no problem.

\LaTeX\ is good at typesetting mathematical formulas
like
       \( x-3y + z = 7 \)
or
       \( a_{1} > x^{2n} + y^{2n} > x' \)
or
       \( \ip{A}{B} = \sum_{i} a_{i} b_{i} \).
The spaces you type in a formula are
ignored.  Remember that a letter like
       $x$                   % $ ... $  and  \( ... \)  are equivalent
is a formula when it denotes a mathematical
symbol, and it should be typed as one.

\section{Displayed Text}

Text is displayed by indenting it from the left
margin.  Quotations are commonly displayed.  There
are short quotations
\begin{quote}
   This is a short quotation.  It consists of a
   single paragraph of text.  See how it is formatted.
\end{quote}
and longer ones.
\begin{quotation}
   This is a longer quotation.  It consists of two
   paragraphs of text, neither of which are
   particularly interesting.

   This is the second paragraph of the quotation.  It
   is just as dull as the first paragraph.
\end{quotation}
Another frequently-displayed structure is a list.
The following is an example of an \emph{itemized}
list.
\begin{itemize}
   \item This is the first item of an itemized list.
         Each item in the list is marked with a ``tick''.
         You don't have to worry about what kind of tick
         mark is used.

   \item This is the second item of the list.  It
         contains another list nested inside it.  The inner
         list is an \emph{enumerated} list.
         \begin{enumerate}
            \item This is the first item of an enumerated
                  list that is nested within the itemized list.

            \item This is the second item of the inner list.
                  \LaTeX\ allows you to nest lists deeper than
                  you really should.
         \end{enumerate}
         This is the rest of the second item of the outer
         list.  It is no more interesting than any other
         part of the item.
   \item This is the third item of the list.
\end{itemize}
You can even display poetry.
\begin{verse}
   There is an environment
    for verse \\             % The \\ command separates lines
   Whose features some poets % within a stanza.
   will curse.

                             % One or more blank lines separate stanzas.

   For instead of making\\
   Them do \emph{all} line breaking, \\
   It allows them to put too many words on a line when they'd rather be
   forced to be terse.
\end{verse}

Mathematical formulas may also be displayed.  A
displayed formula
is
one-line long; multiline
formulas require special formatting instructions.
   \[  \ip{\Gamma}{\psi'} = x'' + y^{2} + z_{i}^{n}\]
Don't start a paragraph with a displayed equation,
nor make one a paragraph by itself.

\end{document}               % End of document.

\section{Bench Huge Section 685}
% synthetic bench section 685
% This is a sample LaTeX input file.  (Version of 12 August 2004.)
%
% A '%' character causes TeX to ignore all remaining text on the line,
% and is used for comments like this one.

\documentclass{article}      % Specifies the document class

                             % The preamble begins here.
\title{An Example Document}  % Declares the document's title.
\author{Leslie Lamport}      % Declares the author's name.
\date{January 21, 1994}      % Deleting this command produces today's date.

\newcommand{\ip}[2]{(#1, #2)}
                             % Defines \ip{arg1}{arg2} to mean
                             % (arg1, arg2).

%\newcommand{\ip}[2]{\langle #1 | #2\rangle}
                             % This is an alternative definition of
                             % \ip that is commented out.

\begin{document}             % End of preamble and beginning of text.

\maketitle                   % Produces the title.

This is an example input file.  Comparing it with
the output it generates can show you how to
produce a simple document of your own.

\section{Ordinary Text}      % Produces section heading.  Lower-level
                             % sections are begun with similar
                             % \subsection and \subsubsection commands.

The ends  of words and sentences are marked
  by   spaces. It  doesn't matter how many
spaces    you type; one is as good as 100.  The
end of   a line counts as a space.

One   or more   blank lines denote the  end
of  a paragraph.

Since any number of consecutive spaces are treated
like a single one, the formatting of the input
file makes no difference to
      \LaTeX,                % The \LaTeX command generates the LaTeX logo.
but it makes a difference to you.  When you use
\LaTeX, making your input file as easy to read
as possible will be a great help as you write
your document and when you change it.  This sample
file shows how you can add comments to your own input
file.

Because printing is different from typewriting,
there are a number of things that you have to do
differently when preparing an input file than if
you were just typing the document directly.
Quotation marks like
       ``this''
have to be handled specially, as do quotes within
quotes:
       ``\,`this'            % \, separates the double and single quote.
        is what I just
        wrote, not  `that'\,''.

Dashes come in three sizes: an
       intra-word
dash, a medium dash for number ranges like
       1--2,
and a punctuation
       dash---like
this.

A sentence-ending space should be larger than the
space between words within a sentence.  You
sometimes have to type special commands in
conjunction with punctuation characters to get
this right, as in the following sentence.
       Gnats, gnus, etc.\ all  % `\ ' makes an inter-word space.
       begin with G\@.         % \@ marks end-of-sentence punctuation.
You should check the spaces after periods when
reading your output to make sure you haven't
forgotten any special cases.  Generating an
ellipsis
       \ldots\               % `\ ' is needed after `\ldots' because TeX
                             % ignores spaces after command names like \ldots
                             % made from \ + letters.
                             %
                             % Note how a `%' character causes TeX to ignore
                             % the end of the input line, so these blank lines
                             % do not start a new paragraph.
                             %
with the right spacing around the periods requires
a special command.

\LaTeX\ interprets some common characters as
commands, so you must type special commands to
generate them.  These characters include the
following:
       \$ \& \% \# \{ and \}.

In printing, text is usually emphasized with an
       \emph{italic}
type style.

\begin{em}
   A long segment of text can also be emphasized
   in this way.  Text within such a segment can be
   given \emph{additional} emphasis.
\end{em}

It is sometimes necessary to prevent \LaTeX\ from
breaking a line where it might otherwise do so.
This may be at a space, as between the ``Mr.''\ and
``Jones'' in
       ``Mr.~Jones'',        % ~ produces an unbreakable interword space.
or within a word---especially when the word is a
symbol like
       \mbox{\emph{itemnum}}
that makes little sense when hyphenated across
lines.

Footnotes\footnote{This is an example of a footnote.}
pose no problem.

\LaTeX\ is good at typesetting mathematical formulas
like
       \( x-3y + z = 7 \)
or
       \( a_{1} > x^{2n} + y^{2n} > x' \)
or
       \( \ip{A}{B} = \sum_{i} a_{i} b_{i} \).
The spaces you type in a formula are
ignored.  Remember that a letter like
       $x$                   % $ ... $  and  \( ... \)  are equivalent
is a formula when it denotes a mathematical
symbol, and it should be typed as one.

\section{Displayed Text}

Text is displayed by indenting it from the left
margin.  Quotations are commonly displayed.  There
are short quotations
\begin{quote}
   This is a short quotation.  It consists of a
   single paragraph of text.  See how it is formatted.
\end{quote}
and longer ones.
\begin{quotation}
   This is a longer quotation.  It consists of two
   paragraphs of text, neither of which are
   particularly interesting.

   This is the second paragraph of the quotation.  It
   is just as dull as the first paragraph.
\end{quotation}
Another frequently-displayed structure is a list.
The following is an example of an \emph{itemized}
list.
\begin{itemize}
   \item This is the first item of an itemized list.
         Each item in the list is marked with a ``tick''.
         You don't have to worry about what kind of tick
         mark is used.

   \item This is the second item of the list.  It
         contains another list nested inside it.  The inner
         list is an \emph{enumerated} list.
         \begin{enumerate}
            \item This is the first item of an enumerated
                  list that is nested within the itemized list.

            \item This is the second item of the inner list.
                  \LaTeX\ allows you to nest lists deeper than
                  you really should.
         \end{enumerate}
         This is the rest of the second item of the outer
         list.  It is no more interesting than any other
         part of the item.
   \item This is the third item of the list.
\end{itemize}
You can even display poetry.
\begin{verse}
   There is an environment
    for verse \\             % The \\ command separates lines
   Whose features some poets % within a stanza.
   will curse.

                             % One or more blank lines separate stanzas.

   For instead of making\\
   Them do \emph{all} line breaking, \\
   It allows them to put too many words on a line when they'd rather be
   forced to be terse.
\end{verse}

Mathematical formulas may also be displayed.  A
displayed formula
is
one-line long; multiline
formulas require special formatting instructions.
   \[  \ip{\Gamma}{\psi'} = x'' + y^{2} + z_{i}^{n}\]
Don't start a paragraph with a displayed equation,
nor make one a paragraph by itself.

\end{document}               % End of document.

\section{Bench Huge Section 686}
% synthetic bench section 686
% This is a sample LaTeX input file.  (Version of 12 August 2004.)
%
% A '%' character causes TeX to ignore all remaining text on the line,
% and is used for comments like this one.

\documentclass{article}      % Specifies the document class

                             % The preamble begins here.
\title{An Example Document}  % Declares the document's title.
\author{Leslie Lamport}      % Declares the author's name.
\date{January 21, 1994}      % Deleting this command produces today's date.

\newcommand{\ip}[2]{(#1, #2)}
                             % Defines \ip{arg1}{arg2} to mean
                             % (arg1, arg2).

%\newcommand{\ip}[2]{\langle #1 | #2\rangle}
                             % This is an alternative definition of
                             % \ip that is commented out.

\begin{document}             % End of preamble and beginning of text.

\maketitle                   % Produces the title.

This is an example input file.  Comparing it with
the output it generates can show you how to
produce a simple document of your own.

\section{Ordinary Text}      % Produces section heading.  Lower-level
                             % sections are begun with similar
                             % \subsection and \subsubsection commands.

The ends  of words and sentences are marked
  by   spaces. It  doesn't matter how many
spaces    you type; one is as good as 100.  The
end of   a line counts as a space.

One   or more   blank lines denote the  end
of  a paragraph.

Since any number of consecutive spaces are treated
like a single one, the formatting of the input
file makes no difference to
      \LaTeX,                % The \LaTeX command generates the LaTeX logo.
but it makes a difference to you.  When you use
\LaTeX, making your input file as easy to read
as possible will be a great help as you write
your document and when you change it.  This sample
file shows how you can add comments to your own input
file.

Because printing is different from typewriting,
there are a number of things that you have to do
differently when preparing an input file than if
you were just typing the document directly.
Quotation marks like
       ``this''
have to be handled specially, as do quotes within
quotes:
       ``\,`this'            % \, separates the double and single quote.
        is what I just
        wrote, not  `that'\,''.

Dashes come in three sizes: an
       intra-word
dash, a medium dash for number ranges like
       1--2,
and a punctuation
       dash---like
this.

A sentence-ending space should be larger than the
space between words within a sentence.  You
sometimes have to type special commands in
conjunction with punctuation characters to get
this right, as in the following sentence.
       Gnats, gnus, etc.\ all  % `\ ' makes an inter-word space.
       begin with G\@.         % \@ marks end-of-sentence punctuation.
You should check the spaces after periods when
reading your output to make sure you haven't
forgotten any special cases.  Generating an
ellipsis
       \ldots\               % `\ ' is needed after `\ldots' because TeX
                             % ignores spaces after command names like \ldots
                             % made from \ + letters.
                             %
                             % Note how a `%' character causes TeX to ignore
                             % the end of the input line, so these blank lines
                             % do not start a new paragraph.
                             %
with the right spacing around the periods requires
a special command.

\LaTeX\ interprets some common characters as
commands, so you must type special commands to
generate them.  These characters include the
following:
       \$ \& \% \# \{ and \}.

In printing, text is usually emphasized with an
       \emph{italic}
type style.

\begin{em}
   A long segment of text can also be emphasized
   in this way.  Text within such a segment can be
   given \emph{additional} emphasis.
\end{em}

It is sometimes necessary to prevent \LaTeX\ from
breaking a line where it might otherwise do so.
This may be at a space, as between the ``Mr.''\ and
``Jones'' in
       ``Mr.~Jones'',        % ~ produces an unbreakable interword space.
or within a word---especially when the word is a
symbol like
       \mbox{\emph{itemnum}}
that makes little sense when hyphenated across
lines.

Footnotes\footnote{This is an example of a footnote.}
pose no problem.

\LaTeX\ is good at typesetting mathematical formulas
like
       \( x-3y + z = 7 \)
or
       \( a_{1} > x^{2n} + y^{2n} > x' \)
or
       \( \ip{A}{B} = \sum_{i} a_{i} b_{i} \).
The spaces you type in a formula are
ignored.  Remember that a letter like
       $x$                   % $ ... $  and  \( ... \)  are equivalent
is a formula when it denotes a mathematical
symbol, and it should be typed as one.

\section{Displayed Text}

Text is displayed by indenting it from the left
margin.  Quotations are commonly displayed.  There
are short quotations
\begin{quote}
   This is a short quotation.  It consists of a
   single paragraph of text.  See how it is formatted.
\end{quote}
and longer ones.
\begin{quotation}
   This is a longer quotation.  It consists of two
   paragraphs of text, neither of which are
   particularly interesting.

   This is the second paragraph of the quotation.  It
   is just as dull as the first paragraph.
\end{quotation}
Another frequently-displayed structure is a list.
The following is an example of an \emph{itemized}
list.
\begin{itemize}
   \item This is the first item of an itemized list.
         Each item in the list is marked with a ``tick''.
         You don't have to worry about what kind of tick
         mark is used.

   \item This is the second item of the list.  It
         contains another list nested inside it.  The inner
         list is an \emph{enumerated} list.
         \begin{enumerate}
            \item This is the first item of an enumerated
                  list that is nested within the itemized list.

            \item This is the second item of the inner list.
                  \LaTeX\ allows you to nest lists deeper than
                  you really should.
         \end{enumerate}
         This is the rest of the second item of the outer
         list.  It is no more interesting than any other
         part of the item.
   \item This is the third item of the list.
\end{itemize}
You can even display poetry.
\begin{verse}
   There is an environment
    for verse \\             % The \\ command separates lines
   Whose features some poets % within a stanza.
   will curse.

                             % One or more blank lines separate stanzas.

   For instead of making\\
   Them do \emph{all} line breaking, \\
   It allows them to put too many words on a line when they'd rather be
   forced to be terse.
\end{verse}

Mathematical formulas may also be displayed.  A
displayed formula
is
one-line long; multiline
formulas require special formatting instructions.
   \[  \ip{\Gamma}{\psi'} = x'' + y^{2} + z_{i}^{n}\]
Don't start a paragraph with a displayed equation,
nor make one a paragraph by itself.

\end{document}               % End of document.

\section{Bench Huge Section 687}
% synthetic bench section 687
% This is a sample LaTeX input file.  (Version of 12 August 2004.)
%
% A '%' character causes TeX to ignore all remaining text on the line,
% and is used for comments like this one.

\documentclass{article}      % Specifies the document class

                             % The preamble begins here.
\title{An Example Document}  % Declares the document's title.
\author{Leslie Lamport}      % Declares the author's name.
\date{January 21, 1994}      % Deleting this command produces today's date.

\newcommand{\ip}[2]{(#1, #2)}
                             % Defines \ip{arg1}{arg2} to mean
                             % (arg1, arg2).

%\newcommand{\ip}[2]{\langle #1 | #2\rangle}
                             % This is an alternative definition of
                             % \ip that is commented out.

\begin{document}             % End of preamble and beginning of text.

\maketitle                   % Produces the title.

This is an example input file.  Comparing it with
the output it generates can show you how to
produce a simple document of your own.

\section{Ordinary Text}      % Produces section heading.  Lower-level
                             % sections are begun with similar
                             % \subsection and \subsubsection commands.

The ends  of words and sentences are marked
  by   spaces. It  doesn't matter how many
spaces    you type; one is as good as 100.  The
end of   a line counts as a space.

One   or more   blank lines denote the  end
of  a paragraph.

Since any number of consecutive spaces are treated
like a single one, the formatting of the input
file makes no difference to
      \LaTeX,                % The \LaTeX command generates the LaTeX logo.
but it makes a difference to you.  When you use
\LaTeX, making your input file as easy to read
as possible will be a great help as you write
your document and when you change it.  This sample
file shows how you can add comments to your own input
file.

Because printing is different from typewriting,
there are a number of things that you have to do
differently when preparing an input file than if
you were just typing the document directly.
Quotation marks like
       ``this''
have to be handled specially, as do quotes within
quotes:
       ``\,`this'            % \, separates the double and single quote.
        is what I just
        wrote, not  `that'\,''.

Dashes come in three sizes: an
       intra-word
dash, a medium dash for number ranges like
       1--2,
and a punctuation
       dash---like
this.

A sentence-ending space should be larger than the
space between words within a sentence.  You
sometimes have to type special commands in
conjunction with punctuation characters to get
this right, as in the following sentence.
       Gnats, gnus, etc.\ all  % `\ ' makes an inter-word space.
       begin with G\@.         % \@ marks end-of-sentence punctuation.
You should check the spaces after periods when
reading your output to make sure you haven't
forgotten any special cases.  Generating an
ellipsis
       \ldots\               % `\ ' is needed after `\ldots' because TeX
                             % ignores spaces after command names like \ldots
                             % made from \ + letters.
                             %
                             % Note how a `%' character causes TeX to ignore
                             % the end of the input line, so these blank lines
                             % do not start a new paragraph.
                             %
with the right spacing around the periods requires
a special command.

\LaTeX\ interprets some common characters as
commands, so you must type special commands to
generate them.  These characters include the
following:
       \$ \& \% \# \{ and \}.

In printing, text is usually emphasized with an
       \emph{italic}
type style.

\begin{em}
   A long segment of text can also be emphasized
   in this way.  Text within such a segment can be
   given \emph{additional} emphasis.
\end{em}

It is sometimes necessary to prevent \LaTeX\ from
breaking a line where it might otherwise do so.
This may be at a space, as between the ``Mr.''\ and
``Jones'' in
       ``Mr.~Jones'',        % ~ produces an unbreakable interword space.
or within a word---especially when the word is a
symbol like
       \mbox{\emph{itemnum}}
that makes little sense when hyphenated across
lines.

Footnotes\footnote{This is an example of a footnote.}
pose no problem.

\LaTeX\ is good at typesetting mathematical formulas
like
       \( x-3y + z = 7 \)
or
       \( a_{1} > x^{2n} + y^{2n} > x' \)
or
       \( \ip{A}{B} = \sum_{i} a_{i} b_{i} \).
The spaces you type in a formula are
ignored.  Remember that a letter like
       $x$                   % $ ... $  and  \( ... \)  are equivalent
is a formula when it denotes a mathematical
symbol, and it should be typed as one.

\section{Displayed Text}

Text is displayed by indenting it from the left
margin.  Quotations are commonly displayed.  There
are short quotations
\begin{quote}
   This is a short quotation.  It consists of a
   single paragraph of text.  See how it is formatted.
\end{quote}
and longer ones.
\begin{quotation}
   This is a longer quotation.  It consists of two
   paragraphs of text, neither of which are
   particularly interesting.

   This is the second paragraph of the quotation.  It
   is just as dull as the first paragraph.
\end{quotation}
Another frequently-displayed structure is a list.
The following is an example of an \emph{itemized}
list.
\begin{itemize}
   \item This is the first item of an itemized list.
         Each item in the list is marked with a ``tick''.
         You don't have to worry about what kind of tick
         mark is used.

   \item This is the second item of the list.  It
         contains another list nested inside it.  The inner
         list is an \emph{enumerated} list.
         \begin{enumerate}
            \item This is the first item of an enumerated
                  list that is nested within the itemized list.

            \item This is the second item of the inner list.
                  \LaTeX\ allows you to nest lists deeper than
                  you really should.
         \end{enumerate}
         This is the rest of the second item of the outer
         list.  It is no more interesting than any other
         part of the item.
   \item This is the third item of the list.
\end{itemize}
You can even display poetry.
\begin{verse}
   There is an environment
    for verse \\             % The \\ command separates lines
   Whose features some poets % within a stanza.
   will curse.

                             % One or more blank lines separate stanzas.

   For instead of making\\
   Them do \emph{all} line breaking, \\
   It allows them to put too many words on a line when they'd rather be
   forced to be terse.
\end{verse}

Mathematical formulas may also be displayed.  A
displayed formula
is
one-line long; multiline
formulas require special formatting instructions.
   \[  \ip{\Gamma}{\psi'} = x'' + y^{2} + z_{i}^{n}\]
Don't start a paragraph with a displayed equation,
nor make one a paragraph by itself.

\end{document}               % End of document.

\section{Bench Huge Section 688}
% synthetic bench section 688
% This is a sample LaTeX input file.  (Version of 12 August 2004.)
%
% A '%' character causes TeX to ignore all remaining text on the line,
% and is used for comments like this one.

\documentclass{article}      % Specifies the document class

                             % The preamble begins here.
\title{An Example Document}  % Declares the document's title.
\author{Leslie Lamport}      % Declares the author's name.
\date{January 21, 1994}      % Deleting this command produces today's date.

\newcommand{\ip}[2]{(#1, #2)}
                             % Defines \ip{arg1}{arg2} to mean
                             % (arg1, arg2).

%\newcommand{\ip}[2]{\langle #1 | #2\rangle}
                             % This is an alternative definition of
                             % \ip that is commented out.

\begin{document}             % End of preamble and beginning of text.

\maketitle                   % Produces the title.

This is an example input file.  Comparing it with
the output it generates can show you how to
produce a simple document of your own.

\section{Ordinary Text}      % Produces section heading.  Lower-level
                             % sections are begun with similar
                             % \subsection and \subsubsection commands.

The ends  of words and sentences are marked
  by   spaces. It  doesn't matter how many
spaces    you type; one is as good as 100.  The
end of   a line counts as a space.

One   or more   blank lines denote the  end
of  a paragraph.

Since any number of consecutive spaces are treated
like a single one, the formatting of the input
file makes no difference to
      \LaTeX,                % The \LaTeX command generates the LaTeX logo.
but it makes a difference to you.  When you use
\LaTeX, making your input file as easy to read
as possible will be a great help as you write
your document and when you change it.  This sample
file shows how you can add comments to your own input
file.

Because printing is different from typewriting,
there are a number of things that you have to do
differently when preparing an input file than if
you were just typing the document directly.
Quotation marks like
       ``this''
have to be handled specially, as do quotes within
quotes:
       ``\,`this'            % \, separates the double and single quote.
        is what I just
        wrote, not  `that'\,''.

Dashes come in three sizes: an
       intra-word
dash, a medium dash for number ranges like
       1--2,
and a punctuation
       dash---like
this.

A sentence-ending space should be larger than the
space between words within a sentence.  You
sometimes have to type special commands in
conjunction with punctuation characters to get
this right, as in the following sentence.
       Gnats, gnus, etc.\ all  % `\ ' makes an inter-word space.
       begin with G\@.         % \@ marks end-of-sentence punctuation.
You should check the spaces after periods when
reading your output to make sure you haven't
forgotten any special cases.  Generating an
ellipsis
       \ldots\               % `\ ' is needed after `\ldots' because TeX
                             % ignores spaces after command names like \ldots
                             % made from \ + letters.
                             %
                             % Note how a `%' character causes TeX to ignore
                             % the end of the input line, so these blank lines
                             % do not start a new paragraph.
                             %
with the right spacing around the periods requires
a special command.

\LaTeX\ interprets some common characters as
commands, so you must type special commands to
generate them.  These characters include the
following:
       \$ \& \% \# \{ and \}.

In printing, text is usually emphasized with an
       \emph{italic}
type style.

\begin{em}
   A long segment of text can also be emphasized
   in this way.  Text within such a segment can be
   given \emph{additional} emphasis.
\end{em}

It is sometimes necessary to prevent \LaTeX\ from
breaking a line where it might otherwise do so.
This may be at a space, as between the ``Mr.''\ and
``Jones'' in
       ``Mr.~Jones'',        % ~ produces an unbreakable interword space.
or within a word---especially when the word is a
symbol like
       \mbox{\emph{itemnum}}
that makes little sense when hyphenated across
lines.

Footnotes\footnote{This is an example of a footnote.}
pose no problem.

\LaTeX\ is good at typesetting mathematical formulas
like
       \( x-3y + z = 7 \)
or
       \( a_{1} > x^{2n} + y^{2n} > x' \)
or
       \( \ip{A}{B} = \sum_{i} a_{i} b_{i} \).
The spaces you type in a formula are
ignored.  Remember that a letter like
       $x$                   % $ ... $  and  \( ... \)  are equivalent
is a formula when it denotes a mathematical
symbol, and it should be typed as one.

\section{Displayed Text}

Text is displayed by indenting it from the left
margin.  Quotations are commonly displayed.  There
are short quotations
\begin{quote}
   This is a short quotation.  It consists of a
   single paragraph of text.  See how it is formatted.
\end{quote}
and longer ones.
\begin{quotation}
   This is a longer quotation.  It consists of two
   paragraphs of text, neither of which are
   particularly interesting.

   This is the second paragraph of the quotation.  It
   is just as dull as the first paragraph.
\end{quotation}
Another frequently-displayed structure is a list.
The following is an example of an \emph{itemized}
list.
\begin{itemize}
   \item This is the first item of an itemized list.
         Each item in the list is marked with a ``tick''.
         You don't have to worry about what kind of tick
         mark is used.

   \item This is the second item of the list.  It
         contains another list nested inside it.  The inner
         list is an \emph{enumerated} list.
         \begin{enumerate}
            \item This is the first item of an enumerated
                  list that is nested within the itemized list.

            \item This is the second item of the inner list.
                  \LaTeX\ allows you to nest lists deeper than
                  you really should.
         \end{enumerate}
         This is the rest of the second item of the outer
         list.  It is no more interesting than any other
         part of the item.
   \item This is the third item of the list.
\end{itemize}
You can even display poetry.
\begin{verse}
   There is an environment
    for verse \\             % The \\ command separates lines
   Whose features some poets % within a stanza.
   will curse.

                             % One or more blank lines separate stanzas.

   For instead of making\\
   Them do \emph{all} line breaking, \\
   It allows them to put too many words on a line when they'd rather be
   forced to be terse.
\end{verse}

Mathematical formulas may also be displayed.  A
displayed formula
is
one-line long; multiline
formulas require special formatting instructions.
   \[  \ip{\Gamma}{\psi'} = x'' + y^{2} + z_{i}^{n}\]
Don't start a paragraph with a displayed equation,
nor make one a paragraph by itself.

\end{document}               % End of document.

\section{Bench Huge Section 689}
% synthetic bench section 689
% This is a sample LaTeX input file.  (Version of 12 August 2004.)
%
% A '%' character causes TeX to ignore all remaining text on the line,
% and is used for comments like this one.

\documentclass{article}      % Specifies the document class

                             % The preamble begins here.
\title{An Example Document}  % Declares the document's title.
\author{Leslie Lamport}      % Declares the author's name.
\date{January 21, 1994}      % Deleting this command produces today's date.

\newcommand{\ip}[2]{(#1, #2)}
                             % Defines \ip{arg1}{arg2} to mean
                             % (arg1, arg2).

%\newcommand{\ip}[2]{\langle #1 | #2\rangle}
                             % This is an alternative definition of
                             % \ip that is commented out.

\begin{document}             % End of preamble and beginning of text.

\maketitle                   % Produces the title.

This is an example input file.  Comparing it with
the output it generates can show you how to
produce a simple document of your own.

\section{Ordinary Text}      % Produces section heading.  Lower-level
                             % sections are begun with similar
                             % \subsection and \subsubsection commands.

The ends  of words and sentences are marked
  by   spaces. It  doesn't matter how many
spaces    you type; one is as good as 100.  The
end of   a line counts as a space.

One   or more   blank lines denote the  end
of  a paragraph.

Since any number of consecutive spaces are treated
like a single one, the formatting of the input
file makes no difference to
      \LaTeX,                % The \LaTeX command generates the LaTeX logo.
but it makes a difference to you.  When you use
\LaTeX, making your input file as easy to read
as possible will be a great help as you write
your document and when you change it.  This sample
file shows how you can add comments to your own input
file.

Because printing is different from typewriting,
there are a number of things that you have to do
differently when preparing an input file than if
you were just typing the document directly.
Quotation marks like
       ``this''
have to be handled specially, as do quotes within
quotes:
       ``\,`this'            % \, separates the double and single quote.
        is what I just
        wrote, not  `that'\,''.

Dashes come in three sizes: an
       intra-word
dash, a medium dash for number ranges like
       1--2,
and a punctuation
       dash---like
this.

A sentence-ending space should be larger than the
space between words within a sentence.  You
sometimes have to type special commands in
conjunction with punctuation characters to get
this right, as in the following sentence.
       Gnats, gnus, etc.\ all  % `\ ' makes an inter-word space.
       begin with G\@.         % \@ marks end-of-sentence punctuation.
You should check the spaces after periods when
reading your output to make sure you haven't
forgotten any special cases.  Generating an
ellipsis
       \ldots\               % `\ ' is needed after `\ldots' because TeX
                             % ignores spaces after command names like \ldots
                             % made from \ + letters.
                             %
                             % Note how a `%' character causes TeX to ignore
                             % the end of the input line, so these blank lines
                             % do not start a new paragraph.
                             %
with the right spacing around the periods requires
a special command.

\LaTeX\ interprets some common characters as
commands, so you must type special commands to
generate them.  These characters include the
following:
       \$ \& \% \# \{ and \}.

In printing, text is usually emphasized with an
       \emph{italic}
type style.

\begin{em}
   A long segment of text can also be emphasized
   in this way.  Text within such a segment can be
   given \emph{additional} emphasis.
\end{em}

It is sometimes necessary to prevent \LaTeX\ from
breaking a line where it might otherwise do so.
This may be at a space, as between the ``Mr.''\ and
``Jones'' in
       ``Mr.~Jones'',        % ~ produces an unbreakable interword space.
or within a word---especially when the word is a
symbol like
       \mbox{\emph{itemnum}}
that makes little sense when hyphenated across
lines.

Footnotes\footnote{This is an example of a footnote.}
pose no problem.

\LaTeX\ is good at typesetting mathematical formulas
like
       \( x-3y + z = 7 \)
or
       \( a_{1} > x^{2n} + y^{2n} > x' \)
or
       \( \ip{A}{B} = \sum_{i} a_{i} b_{i} \).
The spaces you type in a formula are
ignored.  Remember that a letter like
       $x$                   % $ ... $  and  \( ... \)  are equivalent
is a formula when it denotes a mathematical
symbol, and it should be typed as one.

\section{Displayed Text}

Text is displayed by indenting it from the left
margin.  Quotations are commonly displayed.  There
are short quotations
\begin{quote}
   This is a short quotation.  It consists of a
   single paragraph of text.  See how it is formatted.
\end{quote}
and longer ones.
\begin{quotation}
   This is a longer quotation.  It consists of two
   paragraphs of text, neither of which are
   particularly interesting.

   This is the second paragraph of the quotation.  It
   is just as dull as the first paragraph.
\end{quotation}
Another frequently-displayed structure is a list.
The following is an example of an \emph{itemized}
list.
\begin{itemize}
   \item This is the first item of an itemized list.
         Each item in the list is marked with a ``tick''.
         You don't have to worry about what kind of tick
         mark is used.

   \item This is the second item of the list.  It
         contains another list nested inside it.  The inner
         list is an \emph{enumerated} list.
         \begin{enumerate}
            \item This is the first item of an enumerated
                  list that is nested within the itemized list.

            \item This is the second item of the inner list.
                  \LaTeX\ allows you to nest lists deeper than
                  you really should.
         \end{enumerate}
         This is the rest of the second item of the outer
         list.  It is no more interesting than any other
         part of the item.
   \item This is the third item of the list.
\end{itemize}
You can even display poetry.
\begin{verse}
   There is an environment
    for verse \\             % The \\ command separates lines
   Whose features some poets % within a stanza.
   will curse.

                             % One or more blank lines separate stanzas.

   For instead of making\\
   Them do \emph{all} line breaking, \\
   It allows them to put too many words on a line when they'd rather be
   forced to be terse.
\end{verse}

Mathematical formulas may also be displayed.  A
displayed formula
is
one-line long; multiline
formulas require special formatting instructions.
   \[  \ip{\Gamma}{\psi'} = x'' + y^{2} + z_{i}^{n}\]
Don't start a paragraph with a displayed equation,
nor make one a paragraph by itself.

\end{document}               % End of document.

\section{Bench Huge Section 690}
% synthetic bench section 690
% This is a sample LaTeX input file.  (Version of 12 August 2004.)
%
% A '%' character causes TeX to ignore all remaining text on the line,
% and is used for comments like this one.

\documentclass{article}      % Specifies the document class

                             % The preamble begins here.
\title{An Example Document}  % Declares the document's title.
\author{Leslie Lamport}      % Declares the author's name.
\date{January 21, 1994}      % Deleting this command produces today's date.

\newcommand{\ip}[2]{(#1, #2)}
                             % Defines \ip{arg1}{arg2} to mean
                             % (arg1, arg2).

%\newcommand{\ip}[2]{\langle #1 | #2\rangle}
                             % This is an alternative definition of
                             % \ip that is commented out.

\begin{document}             % End of preamble and beginning of text.

\maketitle                   % Produces the title.

This is an example input file.  Comparing it with
the output it generates can show you how to
produce a simple document of your own.

\section{Ordinary Text}      % Produces section heading.  Lower-level
                             % sections are begun with similar
                             % \subsection and \subsubsection commands.

The ends  of words and sentences are marked
  by   spaces. It  doesn't matter how many
spaces    you type; one is as good as 100.  The
end of   a line counts as a space.

One   or more   blank lines denote the  end
of  a paragraph.

Since any number of consecutive spaces are treated
like a single one, the formatting of the input
file makes no difference to
      \LaTeX,                % The \LaTeX command generates the LaTeX logo.
but it makes a difference to you.  When you use
\LaTeX, making your input file as easy to read
as possible will be a great help as you write
your document and when you change it.  This sample
file shows how you can add comments to your own input
file.

Because printing is different from typewriting,
there are a number of things that you have to do
differently when preparing an input file than if
you were just typing the document directly.
Quotation marks like
       ``this''
have to be handled specially, as do quotes within
quotes:
       ``\,`this'            % \, separates the double and single quote.
        is what I just
        wrote, not  `that'\,''.

Dashes come in three sizes: an
       intra-word
dash, a medium dash for number ranges like
       1--2,
and a punctuation
       dash---like
this.

A sentence-ending space should be larger than the
space between words within a sentence.  You
sometimes have to type special commands in
conjunction with punctuation characters to get
this right, as in the following sentence.
       Gnats, gnus, etc.\ all  % `\ ' makes an inter-word space.
       begin with G\@.         % \@ marks end-of-sentence punctuation.
You should check the spaces after periods when
reading your output to make sure you haven't
forgotten any special cases.  Generating an
ellipsis
       \ldots\               % `\ ' is needed after `\ldots' because TeX
                             % ignores spaces after command names like \ldots
                             % made from \ + letters.
                             %
                             % Note how a `%' character causes TeX to ignore
                             % the end of the input line, so these blank lines
                             % do not start a new paragraph.
                             %
with the right spacing around the periods requires
a special command.

\LaTeX\ interprets some common characters as
commands, so you must type special commands to
generate them.  These characters include the
following:
       \$ \& \% \# \{ and \}.

In printing, text is usually emphasized with an
       \emph{italic}
type style.

\begin{em}
   A long segment of text can also be emphasized
   in this way.  Text within such a segment can be
   given \emph{additional} emphasis.
\end{em}

It is sometimes necessary to prevent \LaTeX\ from
breaking a line where it might otherwise do so.
This may be at a space, as between the ``Mr.''\ and
``Jones'' in
       ``Mr.~Jones'',        % ~ produces an unbreakable interword space.
or within a word---especially when the word is a
symbol like
       \mbox{\emph{itemnum}}
that makes little sense when hyphenated across
lines.

Footnotes\footnote{This is an example of a footnote.}
pose no problem.

\LaTeX\ is good at typesetting mathematical formulas
like
       \( x-3y + z = 7 \)
or
       \( a_{1} > x^{2n} + y^{2n} > x' \)
or
       \( \ip{A}{B} = \sum_{i} a_{i} b_{i} \).
The spaces you type in a formula are
ignored.  Remember that a letter like
       $x$                   % $ ... $  and  \( ... \)  are equivalent
is a formula when it denotes a mathematical
symbol, and it should be typed as one.

\section{Displayed Text}

Text is displayed by indenting it from the left
margin.  Quotations are commonly displayed.  There
are short quotations
\begin{quote}
   This is a short quotation.  It consists of a
   single paragraph of text.  See how it is formatted.
\end{quote}
and longer ones.
\begin{quotation}
   This is a longer quotation.  It consists of two
   paragraphs of text, neither of which are
   particularly interesting.

   This is the second paragraph of the quotation.  It
   is just as dull as the first paragraph.
\end{quotation}
Another frequently-displayed structure is a list.
The following is an example of an \emph{itemized}
list.
\begin{itemize}
   \item This is the first item of an itemized list.
         Each item in the list is marked with a ``tick''.
         You don't have to worry about what kind of tick
         mark is used.

   \item This is the second item of the list.  It
         contains another list nested inside it.  The inner
         list is an \emph{enumerated} list.
         \begin{enumerate}
            \item This is the first item of an enumerated
                  list that is nested within the itemized list.

            \item This is the second item of the inner list.
                  \LaTeX\ allows you to nest lists deeper than
                  you really should.
         \end{enumerate}
         This is the rest of the second item of the outer
         list.  It is no more interesting than any other
         part of the item.
   \item This is the third item of the list.
\end{itemize}
You can even display poetry.
\begin{verse}
   There is an environment
    for verse \\             % The \\ command separates lines
   Whose features some poets % within a stanza.
   will curse.

                             % One or more blank lines separate stanzas.

   For instead of making\\
   Them do \emph{all} line breaking, \\
   It allows them to put too many words on a line when they'd rather be
   forced to be terse.
\end{verse}

Mathematical formulas may also be displayed.  A
displayed formula
is
one-line long; multiline
formulas require special formatting instructions.
   \[  \ip{\Gamma}{\psi'} = x'' + y^{2} + z_{i}^{n}\]
Don't start a paragraph with a displayed equation,
nor make one a paragraph by itself.

\end{document}               % End of document.

\section{Bench Huge Section 691}
% synthetic bench section 691
% This is a sample LaTeX input file.  (Version of 12 August 2004.)
%
% A '%' character causes TeX to ignore all remaining text on the line,
% and is used for comments like this one.

\documentclass{article}      % Specifies the document class

                             % The preamble begins here.
\title{An Example Document}  % Declares the document's title.
\author{Leslie Lamport}      % Declares the author's name.
\date{January 21, 1994}      % Deleting this command produces today's date.

\newcommand{\ip}[2]{(#1, #2)}
                             % Defines \ip{arg1}{arg2} to mean
                             % (arg1, arg2).

%\newcommand{\ip}[2]{\langle #1 | #2\rangle}
                             % This is an alternative definition of
                             % \ip that is commented out.

\begin{document}             % End of preamble and beginning of text.

\maketitle                   % Produces the title.

This is an example input file.  Comparing it with
the output it generates can show you how to
produce a simple document of your own.

\section{Ordinary Text}      % Produces section heading.  Lower-level
                             % sections are begun with similar
                             % \subsection and \subsubsection commands.

The ends  of words and sentences are marked
  by   spaces. It  doesn't matter how many
spaces    you type; one is as good as 100.  The
end of   a line counts as a space.

One   or more   blank lines denote the  end
of  a paragraph.

Since any number of consecutive spaces are treated
like a single one, the formatting of the input
file makes no difference to
      \LaTeX,                % The \LaTeX command generates the LaTeX logo.
but it makes a difference to you.  When you use
\LaTeX, making your input file as easy to read
as possible will be a great help as you write
your document and when you change it.  This sample
file shows how you can add comments to your own input
file.

Because printing is different from typewriting,
there are a number of things that you have to do
differently when preparing an input file than if
you were just typing the document directly.
Quotation marks like
       ``this''
have to be handled specially, as do quotes within
quotes:
       ``\,`this'            % \, separates the double and single quote.
        is what I just
        wrote, not  `that'\,''.

Dashes come in three sizes: an
       intra-word
dash, a medium dash for number ranges like
       1--2,
and a punctuation
       dash---like
this.

A sentence-ending space should be larger than the
space between words within a sentence.  You
sometimes have to type special commands in
conjunction with punctuation characters to get
this right, as in the following sentence.
       Gnats, gnus, etc.\ all  % `\ ' makes an inter-word space.
       begin with G\@.         % \@ marks end-of-sentence punctuation.
You should check the spaces after periods when
reading your output to make sure you haven't
forgotten any special cases.  Generating an
ellipsis
       \ldots\               % `\ ' is needed after `\ldots' because TeX
                             % ignores spaces after command names like \ldots
                             % made from \ + letters.
                             %
                             % Note how a `%' character causes TeX to ignore
                             % the end of the input line, so these blank lines
                             % do not start a new paragraph.
                             %
with the right spacing around the periods requires
a special command.

\LaTeX\ interprets some common characters as
commands, so you must type special commands to
generate them.  These characters include the
following:
       \$ \& \% \# \{ and \}.

In printing, text is usually emphasized with an
       \emph{italic}
type style.

\begin{em}
   A long segment of text can also be emphasized
   in this way.  Text within such a segment can be
   given \emph{additional} emphasis.
\end{em}

It is sometimes necessary to prevent \LaTeX\ from
breaking a line where it might otherwise do so.
This may be at a space, as between the ``Mr.''\ and
``Jones'' in
       ``Mr.~Jones'',        % ~ produces an unbreakable interword space.
or within a word---especially when the word is a
symbol like
       \mbox{\emph{itemnum}}
that makes little sense when hyphenated across
lines.

Footnotes\footnote{This is an example of a footnote.}
pose no problem.

\LaTeX\ is good at typesetting mathematical formulas
like
       \( x-3y + z = 7 \)
or
       \( a_{1} > x^{2n} + y^{2n} > x' \)
or
       \( \ip{A}{B} = \sum_{i} a_{i} b_{i} \).
The spaces you type in a formula are
ignored.  Remember that a letter like
       $x$                   % $ ... $  and  \( ... \)  are equivalent
is a formula when it denotes a mathematical
symbol, and it should be typed as one.

\section{Displayed Text}

Text is displayed by indenting it from the left
margin.  Quotations are commonly displayed.  There
are short quotations
\begin{quote}
   This is a short quotation.  It consists of a
   single paragraph of text.  See how it is formatted.
\end{quote}
and longer ones.
\begin{quotation}
   This is a longer quotation.  It consists of two
   paragraphs of text, neither of which are
   particularly interesting.

   This is the second paragraph of the quotation.  It
   is just as dull as the first paragraph.
\end{quotation}
Another frequently-displayed structure is a list.
The following is an example of an \emph{itemized}
list.
\begin{itemize}
   \item This is the first item of an itemized list.
         Each item in the list is marked with a ``tick''.
         You don't have to worry about what kind of tick
         mark is used.

   \item This is the second item of the list.  It
         contains another list nested inside it.  The inner
         list is an \emph{enumerated} list.
         \begin{enumerate}
            \item This is the first item of an enumerated
                  list that is nested within the itemized list.

            \item This is the second item of the inner list.
                  \LaTeX\ allows you to nest lists deeper than
                  you really should.
         \end{enumerate}
         This is the rest of the second item of the outer
         list.  It is no more interesting than any other
         part of the item.
   \item This is the third item of the list.
\end{itemize}
You can even display poetry.
\begin{verse}
   There is an environment
    for verse \\             % The \\ command separates lines
   Whose features some poets % within a stanza.
   will curse.

                             % One or more blank lines separate stanzas.

   For instead of making\\
   Them do \emph{all} line breaking, \\
   It allows them to put too many words on a line when they'd rather be
   forced to be terse.
\end{verse}

Mathematical formulas may also be displayed.  A
displayed formula
is
one-line long; multiline
formulas require special formatting instructions.
   \[  \ip{\Gamma}{\psi'} = x'' + y^{2} + z_{i}^{n}\]
Don't start a paragraph with a displayed equation,
nor make one a paragraph by itself.

\end{document}               % End of document.

\section{Bench Huge Section 692}
% synthetic bench section 692
% This is a sample LaTeX input file.  (Version of 12 August 2004.)
%
% A '%' character causes TeX to ignore all remaining text on the line,
% and is used for comments like this one.

\documentclass{article}      % Specifies the document class

                             % The preamble begins here.
\title{An Example Document}  % Declares the document's title.
\author{Leslie Lamport}      % Declares the author's name.
\date{January 21, 1994}      % Deleting this command produces today's date.

\newcommand{\ip}[2]{(#1, #2)}
                             % Defines \ip{arg1}{arg2} to mean
                             % (arg1, arg2).

%\newcommand{\ip}[2]{\langle #1 | #2\rangle}
                             % This is an alternative definition of
                             % \ip that is commented out.

\begin{document}             % End of preamble and beginning of text.

\maketitle                   % Produces the title.

This is an example input file.  Comparing it with
the output it generates can show you how to
produce a simple document of your own.

\section{Ordinary Text}      % Produces section heading.  Lower-level
                             % sections are begun with similar
                             % \subsection and \subsubsection commands.

The ends  of words and sentences are marked
  by   spaces. It  doesn't matter how many
spaces    you type; one is as good as 100.  The
end of   a line counts as a space.

One   or more   blank lines denote the  end
of  a paragraph.

Since any number of consecutive spaces are treated
like a single one, the formatting of the input
file makes no difference to
      \LaTeX,                % The \LaTeX command generates the LaTeX logo.
but it makes a difference to you.  When you use
\LaTeX, making your input file as easy to read
as possible will be a great help as you write
your document and when you change it.  This sample
file shows how you can add comments to your own input
file.

Because printing is different from typewriting,
there are a number of things that you have to do
differently when preparing an input file than if
you were just typing the document directly.
Quotation marks like
       ``this''
have to be handled specially, as do quotes within
quotes:
       ``\,`this'            % \, separates the double and single quote.
        is what I just
        wrote, not  `that'\,''.

Dashes come in three sizes: an
       intra-word
dash, a medium dash for number ranges like
       1--2,
and a punctuation
       dash---like
this.

A sentence-ending space should be larger than the
space between words within a sentence.  You
sometimes have to type special commands in
conjunction with punctuation characters to get
this right, as in the following sentence.
       Gnats, gnus, etc.\ all  % `\ ' makes an inter-word space.
       begin with G\@.         % \@ marks end-of-sentence punctuation.
You should check the spaces after periods when
reading your output to make sure you haven't
forgotten any special cases.  Generating an
ellipsis
       \ldots\               % `\ ' is needed after `\ldots' because TeX
                             % ignores spaces after command names like \ldots
                             % made from \ + letters.
                             %
                             % Note how a `%' character causes TeX to ignore
                             % the end of the input line, so these blank lines
                             % do not start a new paragraph.
                             %
with the right spacing around the periods requires
a special command.

\LaTeX\ interprets some common characters as
commands, so you must type special commands to
generate them.  These characters include the
following:
       \$ \& \% \# \{ and \}.

In printing, text is usually emphasized with an
       \emph{italic}
type style.

\begin{em}
   A long segment of text can also be emphasized
   in this way.  Text within such a segment can be
   given \emph{additional} emphasis.
\end{em}

It is sometimes necessary to prevent \LaTeX\ from
breaking a line where it might otherwise do so.
This may be at a space, as between the ``Mr.''\ and
``Jones'' in
       ``Mr.~Jones'',        % ~ produces an unbreakable interword space.
or within a word---especially when the word is a
symbol like
       \mbox{\emph{itemnum}}
that makes little sense when hyphenated across
lines.

Footnotes\footnote{This is an example of a footnote.}
pose no problem.

\LaTeX\ is good at typesetting mathematical formulas
like
       \( x-3y + z = 7 \)
or
       \( a_{1} > x^{2n} + y^{2n} > x' \)
or
       \( \ip{A}{B} = \sum_{i} a_{i} b_{i} \).
The spaces you type in a formula are
ignored.  Remember that a letter like
       $x$                   % $ ... $  and  \( ... \)  are equivalent
is a formula when it denotes a mathematical
symbol, and it should be typed as one.

\section{Displayed Text}

Text is displayed by indenting it from the left
margin.  Quotations are commonly displayed.  There
are short quotations
\begin{quote}
   This is a short quotation.  It consists of a
   single paragraph of text.  See how it is formatted.
\end{quote}
and longer ones.
\begin{quotation}
   This is a longer quotation.  It consists of two
   paragraphs of text, neither of which are
   particularly interesting.

   This is the second paragraph of the quotation.  It
   is just as dull as the first paragraph.
\end{quotation}
Another frequently-displayed structure is a list.
The following is an example of an \emph{itemized}
list.
\begin{itemize}
   \item This is the first item of an itemized list.
         Each item in the list is marked with a ``tick''.
         You don't have to worry about what kind of tick
         mark is used.

   \item This is the second item of the list.  It
         contains another list nested inside it.  The inner
         list is an \emph{enumerated} list.
         \begin{enumerate}
            \item This is the first item of an enumerated
                  list that is nested within the itemized list.

            \item This is the second item of the inner list.
                  \LaTeX\ allows you to nest lists deeper than
                  you really should.
         \end{enumerate}
         This is the rest of the second item of the outer
         list.  It is no more interesting than any other
         part of the item.
   \item This is the third item of the list.
\end{itemize}
You can even display poetry.
\begin{verse}
   There is an environment
    for verse \\             % The \\ command separates lines
   Whose features some poets % within a stanza.
   will curse.

                             % One or more blank lines separate stanzas.

   For instead of making\\
   Them do \emph{all} line breaking, \\
   It allows them to put too many words on a line when they'd rather be
   forced to be terse.
\end{verse}

Mathematical formulas may also be displayed.  A
displayed formula
is
one-line long; multiline
formulas require special formatting instructions.
   \[  \ip{\Gamma}{\psi'} = x'' + y^{2} + z_{i}^{n}\]
Don't start a paragraph with a displayed equation,
nor make one a paragraph by itself.

\end{document}               % End of document.

\section{Bench Huge Section 693}
% synthetic bench section 693
% This is a sample LaTeX input file.  (Version of 12 August 2004.)
%
% A '%' character causes TeX to ignore all remaining text on the line,
% and is used for comments like this one.

\documentclass{article}      % Specifies the document class

                             % The preamble begins here.
\title{An Example Document}  % Declares the document's title.
\author{Leslie Lamport}      % Declares the author's name.
\date{January 21, 1994}      % Deleting this command produces today's date.

\newcommand{\ip}[2]{(#1, #2)}
                             % Defines \ip{arg1}{arg2} to mean
                             % (arg1, arg2).

%\newcommand{\ip}[2]{\langle #1 | #2\rangle}
                             % This is an alternative definition of
                             % \ip that is commented out.

\begin{document}             % End of preamble and beginning of text.

\maketitle                   % Produces the title.

This is an example input file.  Comparing it with
the output it generates can show you how to
produce a simple document of your own.

\section{Ordinary Text}      % Produces section heading.  Lower-level
                             % sections are begun with similar
                             % \subsection and \subsubsection commands.

The ends  of words and sentences are marked
  by   spaces. It  doesn't matter how many
spaces    you type; one is as good as 100.  The
end of   a line counts as a space.

One   or more   blank lines denote the  end
of  a paragraph.

Since any number of consecutive spaces are treated
like a single one, the formatting of the input
file makes no difference to
      \LaTeX,                % The \LaTeX command generates the LaTeX logo.
but it makes a difference to you.  When you use
\LaTeX, making your input file as easy to read
as possible will be a great help as you write
your document and when you change it.  This sample
file shows how you can add comments to your own input
file.

Because printing is different from typewriting,
there are a number of things that you have to do
differently when preparing an input file than if
you were just typing the document directly.
Quotation marks like
       ``this''
have to be handled specially, as do quotes within
quotes:
       ``\,`this'            % \, separates the double and single quote.
        is what I just
        wrote, not  `that'\,''.

Dashes come in three sizes: an
       intra-word
dash, a medium dash for number ranges like
       1--2,
and a punctuation
       dash---like
this.

A sentence-ending space should be larger than the
space between words within a sentence.  You
sometimes have to type special commands in
conjunction with punctuation characters to get
this right, as in the following sentence.
       Gnats, gnus, etc.\ all  % `\ ' makes an inter-word space.
       begin with G\@.         % \@ marks end-of-sentence punctuation.
You should check the spaces after periods when
reading your output to make sure you haven't
forgotten any special cases.  Generating an
ellipsis
       \ldots\               % `\ ' is needed after `\ldots' because TeX
                             % ignores spaces after command names like \ldots
                             % made from \ + letters.
                             %
                             % Note how a `%' character causes TeX to ignore
                             % the end of the input line, so these blank lines
                             % do not start a new paragraph.
                             %
with the right spacing around the periods requires
a special command.

\LaTeX\ interprets some common characters as
commands, so you must type special commands to
generate them.  These characters include the
following:
       \$ \& \% \# \{ and \}.

In printing, text is usually emphasized with an
       \emph{italic}
type style.

\begin{em}
   A long segment of text can also be emphasized
   in this way.  Text within such a segment can be
   given \emph{additional} emphasis.
\end{em}

It is sometimes necessary to prevent \LaTeX\ from
breaking a line where it might otherwise do so.
This may be at a space, as between the ``Mr.''\ and
``Jones'' in
       ``Mr.~Jones'',        % ~ produces an unbreakable interword space.
or within a word---especially when the word is a
symbol like
       \mbox{\emph{itemnum}}
that makes little sense when hyphenated across
lines.

Footnotes\footnote{This is an example of a footnote.}
pose no problem.

\LaTeX\ is good at typesetting mathematical formulas
like
       \( x-3y + z = 7 \)
or
       \( a_{1} > x^{2n} + y^{2n} > x' \)
or
       \( \ip{A}{B} = \sum_{i} a_{i} b_{i} \).
The spaces you type in a formula are
ignored.  Remember that a letter like
       $x$                   % $ ... $  and  \( ... \)  are equivalent
is a formula when it denotes a mathematical
symbol, and it should be typed as one.

\section{Displayed Text}

Text is displayed by indenting it from the left
margin.  Quotations are commonly displayed.  There
are short quotations
\begin{quote}
   This is a short quotation.  It consists of a
   single paragraph of text.  See how it is formatted.
\end{quote}
and longer ones.
\begin{quotation}
   This is a longer quotation.  It consists of two
   paragraphs of text, neither of which are
   particularly interesting.

   This is the second paragraph of the quotation.  It
   is just as dull as the first paragraph.
\end{quotation}
Another frequently-displayed structure is a list.
The following is an example of an \emph{itemized}
list.
\begin{itemize}
   \item This is the first item of an itemized list.
         Each item in the list is marked with a ``tick''.
         You don't have to worry about what kind of tick
         mark is used.

   \item This is the second item of the list.  It
         contains another list nested inside it.  The inner
         list is an \emph{enumerated} list.
         \begin{enumerate}
            \item This is the first item of an enumerated
                  list that is nested within the itemized list.

            \item This is the second item of the inner list.
                  \LaTeX\ allows you to nest lists deeper than
                  you really should.
         \end{enumerate}
         This is the rest of the second item of the outer
         list.  It is no more interesting than any other
         part of the item.
   \item This is the third item of the list.
\end{itemize}
You can even display poetry.
\begin{verse}
   There is an environment
    for verse \\             % The \\ command separates lines
   Whose features some poets % within a stanza.
   will curse.

                             % One or more blank lines separate stanzas.

   For instead of making\\
   Them do \emph{all} line breaking, \\
   It allows them to put too many words on a line when they'd rather be
   forced to be terse.
\end{verse}

Mathematical formulas may also be displayed.  A
displayed formula
is
one-line long; multiline
formulas require special formatting instructions.
   \[  \ip{\Gamma}{\psi'} = x'' + y^{2} + z_{i}^{n}\]
Don't start a paragraph with a displayed equation,
nor make one a paragraph by itself.

\end{document}               % End of document.

\section{Bench Huge Section 694}
% synthetic bench section 694
% This is a sample LaTeX input file.  (Version of 12 August 2004.)
%
% A '%' character causes TeX to ignore all remaining text on the line,
% and is used for comments like this one.

\documentclass{article}      % Specifies the document class

                             % The preamble begins here.
\title{An Example Document}  % Declares the document's title.
\author{Leslie Lamport}      % Declares the author's name.
\date{January 21, 1994}      % Deleting this command produces today's date.

\newcommand{\ip}[2]{(#1, #2)}
                             % Defines \ip{arg1}{arg2} to mean
                             % (arg1, arg2).

%\newcommand{\ip}[2]{\langle #1 | #2\rangle}
                             % This is an alternative definition of
                             % \ip that is commented out.

\begin{document}             % End of preamble and beginning of text.

\maketitle                   % Produces the title.

This is an example input file.  Comparing it with
the output it generates can show you how to
produce a simple document of your own.

\section{Ordinary Text}      % Produces section heading.  Lower-level
                             % sections are begun with similar
                             % \subsection and \subsubsection commands.

The ends  of words and sentences are marked
  by   spaces. It  doesn't matter how many
spaces    you type; one is as good as 100.  The
end of   a line counts as a space.

One   or more   blank lines denote the  end
of  a paragraph.

Since any number of consecutive spaces are treated
like a single one, the formatting of the input
file makes no difference to
      \LaTeX,                % The \LaTeX command generates the LaTeX logo.
but it makes a difference to you.  When you use
\LaTeX, making your input file as easy to read
as possible will be a great help as you write
your document and when you change it.  This sample
file shows how you can add comments to your own input
file.

Because printing is different from typewriting,
there are a number of things that you have to do
differently when preparing an input file than if
you were just typing the document directly.
Quotation marks like
       ``this''
have to be handled specially, as do quotes within
quotes:
       ``\,`this'            % \, separates the double and single quote.
        is what I just
        wrote, not  `that'\,''.

Dashes come in three sizes: an
       intra-word
dash, a medium dash for number ranges like
       1--2,
and a punctuation
       dash---like
this.

A sentence-ending space should be larger than the
space between words within a sentence.  You
sometimes have to type special commands in
conjunction with punctuation characters to get
this right, as in the following sentence.
       Gnats, gnus, etc.\ all  % `\ ' makes an inter-word space.
       begin with G\@.         % \@ marks end-of-sentence punctuation.
You should check the spaces after periods when
reading your output to make sure you haven't
forgotten any special cases.  Generating an
ellipsis
       \ldots\               % `\ ' is needed after `\ldots' because TeX
                             % ignores spaces after command names like \ldots
                             % made from \ + letters.
                             %
                             % Note how a `%' character causes TeX to ignore
                             % the end of the input line, so these blank lines
                             % do not start a new paragraph.
                             %
with the right spacing around the periods requires
a special command.

\LaTeX\ interprets some common characters as
commands, so you must type special commands to
generate them.  These characters include the
following:
       \$ \& \% \# \{ and \}.

In printing, text is usually emphasized with an
       \emph{italic}
type style.

\begin{em}
   A long segment of text can also be emphasized
   in this way.  Text within such a segment can be
   given \emph{additional} emphasis.
\end{em}

It is sometimes necessary to prevent \LaTeX\ from
breaking a line where it might otherwise do so.
This may be at a space, as between the ``Mr.''\ and
``Jones'' in
       ``Mr.~Jones'',        % ~ produces an unbreakable interword space.
or within a word---especially when the word is a
symbol like
       \mbox{\emph{itemnum}}
that makes little sense when hyphenated across
lines.

Footnotes\footnote{This is an example of a footnote.}
pose no problem.

\LaTeX\ is good at typesetting mathematical formulas
like
       \( x-3y + z = 7 \)
or
       \( a_{1} > x^{2n} + y^{2n} > x' \)
or
       \( \ip{A}{B} = \sum_{i} a_{i} b_{i} \).
The spaces you type in a formula are
ignored.  Remember that a letter like
       $x$                   % $ ... $  and  \( ... \)  are equivalent
is a formula when it denotes a mathematical
symbol, and it should be typed as one.

\section{Displayed Text}

Text is displayed by indenting it from the left
margin.  Quotations are commonly displayed.  There
are short quotations
\begin{quote}
   This is a short quotation.  It consists of a
   single paragraph of text.  See how it is formatted.
\end{quote}
and longer ones.
\begin{quotation}
   This is a longer quotation.  It consists of two
   paragraphs of text, neither of which are
   particularly interesting.

   This is the second paragraph of the quotation.  It
   is just as dull as the first paragraph.
\end{quotation}
Another frequently-displayed structure is a list.
The following is an example of an \emph{itemized}
list.
\begin{itemize}
   \item This is the first item of an itemized list.
         Each item in the list is marked with a ``tick''.
         You don't have to worry about what kind of tick
         mark is used.

   \item This is the second item of the list.  It
         contains another list nested inside it.  The inner
         list is an \emph{enumerated} list.
         \begin{enumerate}
            \item This is the first item of an enumerated
                  list that is nested within the itemized list.

            \item This is the second item of the inner list.
                  \LaTeX\ allows you to nest lists deeper than
                  you really should.
         \end{enumerate}
         This is the rest of the second item of the outer
         list.  It is no more interesting than any other
         part of the item.
   \item This is the third item of the list.
\end{itemize}
You can even display poetry.
\begin{verse}
   There is an environment
    for verse \\             % The \\ command separates lines
   Whose features some poets % within a stanza.
   will curse.

                             % One or more blank lines separate stanzas.

   For instead of making\\
   Them do \emph{all} line breaking, \\
   It allows them to put too many words on a line when they'd rather be
   forced to be terse.
\end{verse}

Mathematical formulas may also be displayed.  A
displayed formula
is
one-line long; multiline
formulas require special formatting instructions.
   \[  \ip{\Gamma}{\psi'} = x'' + y^{2} + z_{i}^{n}\]
Don't start a paragraph with a displayed equation,
nor make one a paragraph by itself.

\end{document}               % End of document.

\section{Bench Huge Section 695}
% synthetic bench section 695
% This is a sample LaTeX input file.  (Version of 12 August 2004.)
%
% A '%' character causes TeX to ignore all remaining text on the line,
% and is used for comments like this one.

\documentclass{article}      % Specifies the document class

                             % The preamble begins here.
\title{An Example Document}  % Declares the document's title.
\author{Leslie Lamport}      % Declares the author's name.
\date{January 21, 1994}      % Deleting this command produces today's date.

\newcommand{\ip}[2]{(#1, #2)}
                             % Defines \ip{arg1}{arg2} to mean
                             % (arg1, arg2).

%\newcommand{\ip}[2]{\langle #1 | #2\rangle}
                             % This is an alternative definition of
                             % \ip that is commented out.

\begin{document}             % End of preamble and beginning of text.

\maketitle                   % Produces the title.

This is an example input file.  Comparing it with
the output it generates can show you how to
produce a simple document of your own.

\section{Ordinary Text}      % Produces section heading.  Lower-level
                             % sections are begun with similar
                             % \subsection and \subsubsection commands.

The ends  of words and sentences are marked
  by   spaces. It  doesn't matter how many
spaces    you type; one is as good as 100.  The
end of   a line counts as a space.

One   or more   blank lines denote the  end
of  a paragraph.

Since any number of consecutive spaces are treated
like a single one, the formatting of the input
file makes no difference to
      \LaTeX,                % The \LaTeX command generates the LaTeX logo.
but it makes a difference to you.  When you use
\LaTeX, making your input file as easy to read
as possible will be a great help as you write
your document and when you change it.  This sample
file shows how you can add comments to your own input
file.

Because printing is different from typewriting,
there are a number of things that you have to do
differently when preparing an input file than if
you were just typing the document directly.
Quotation marks like
       ``this''
have to be handled specially, as do quotes within
quotes:
       ``\,`this'            % \, separates the double and single quote.
        is what I just
        wrote, not  `that'\,''.

Dashes come in three sizes: an
       intra-word
dash, a medium dash for number ranges like
       1--2,
and a punctuation
       dash---like
this.

A sentence-ending space should be larger than the
space between words within a sentence.  You
sometimes have to type special commands in
conjunction with punctuation characters to get
this right, as in the following sentence.
       Gnats, gnus, etc.\ all  % `\ ' makes an inter-word space.
       begin with G\@.         % \@ marks end-of-sentence punctuation.
You should check the spaces after periods when
reading your output to make sure you haven't
forgotten any special cases.  Generating an
ellipsis
       \ldots\               % `\ ' is needed after `\ldots' because TeX
                             % ignores spaces after command names like \ldots
                             % made from \ + letters.
                             %
                             % Note how a `%' character causes TeX to ignore
                             % the end of the input line, so these blank lines
                             % do not start a new paragraph.
                             %
with the right spacing around the periods requires
a special command.

\LaTeX\ interprets some common characters as
commands, so you must type special commands to
generate them.  These characters include the
following:
       \$ \& \% \# \{ and \}.

In printing, text is usually emphasized with an
       \emph{italic}
type style.

\begin{em}
   A long segment of text can also be emphasized
   in this way.  Text within such a segment can be
   given \emph{additional} emphasis.
\end{em}

It is sometimes necessary to prevent \LaTeX\ from
breaking a line where it might otherwise do so.
This may be at a space, as between the ``Mr.''\ and
``Jones'' in
       ``Mr.~Jones'',        % ~ produces an unbreakable interword space.
or within a word---especially when the word is a
symbol like
       \mbox{\emph{itemnum}}
that makes little sense when hyphenated across
lines.

Footnotes\footnote{This is an example of a footnote.}
pose no problem.

\LaTeX\ is good at typesetting mathematical formulas
like
       \( x-3y + z = 7 \)
or
       \( a_{1} > x^{2n} + y^{2n} > x' \)
or
       \( \ip{A}{B} = \sum_{i} a_{i} b_{i} \).
The spaces you type in a formula are
ignored.  Remember that a letter like
       $x$                   % $ ... $  and  \( ... \)  are equivalent
is a formula when it denotes a mathematical
symbol, and it should be typed as one.

\section{Displayed Text}

Text is displayed by indenting it from the left
margin.  Quotations are commonly displayed.  There
are short quotations
\begin{quote}
   This is a short quotation.  It consists of a
   single paragraph of text.  See how it is formatted.
\end{quote}
and longer ones.
\begin{quotation}
   This is a longer quotation.  It consists of two
   paragraphs of text, neither of which are
   particularly interesting.

   This is the second paragraph of the quotation.  It
   is just as dull as the first paragraph.
\end{quotation}
Another frequently-displayed structure is a list.
The following is an example of an \emph{itemized}
list.
\begin{itemize}
   \item This is the first item of an itemized list.
         Each item in the list is marked with a ``tick''.
         You don't have to worry about what kind of tick
         mark is used.

   \item This is the second item of the list.  It
         contains another list nested inside it.  The inner
         list is an \emph{enumerated} list.
         \begin{enumerate}
            \item This is the first item of an enumerated
                  list that is nested within the itemized list.

            \item This is the second item of the inner list.
                  \LaTeX\ allows you to nest lists deeper than
                  you really should.
         \end{enumerate}
         This is the rest of the second item of the outer
         list.  It is no more interesting than any other
         part of the item.
   \item This is the third item of the list.
\end{itemize}
You can even display poetry.
\begin{verse}
   There is an environment
    for verse \\             % The \\ command separates lines
   Whose features some poets % within a stanza.
   will curse.

                             % One or more blank lines separate stanzas.

   For instead of making\\
   Them do \emph{all} line breaking, \\
   It allows them to put too many words on a line when they'd rather be
   forced to be terse.
\end{verse}

Mathematical formulas may also be displayed.  A
displayed formula
is
one-line long; multiline
formulas require special formatting instructions.
   \[  \ip{\Gamma}{\psi'} = x'' + y^{2} + z_{i}^{n}\]
Don't start a paragraph with a displayed equation,
nor make one a paragraph by itself.

\end{document}               % End of document.

\section{Bench Huge Section 696}
% synthetic bench section 696
% This is a sample LaTeX input file.  (Version of 12 August 2004.)
%
% A '%' character causes TeX to ignore all remaining text on the line,
% and is used for comments like this one.

\documentclass{article}      % Specifies the document class

                             % The preamble begins here.
\title{An Example Document}  % Declares the document's title.
\author{Leslie Lamport}      % Declares the author's name.
\date{January 21, 1994}      % Deleting this command produces today's date.

\newcommand{\ip}[2]{(#1, #2)}
                             % Defines \ip{arg1}{arg2} to mean
                             % (arg1, arg2).

%\newcommand{\ip}[2]{\langle #1 | #2\rangle}
                             % This is an alternative definition of
                             % \ip that is commented out.

\begin{document}             % End of preamble and beginning of text.

\maketitle                   % Produces the title.

This is an example input file.  Comparing it with
the output it generates can show you how to
produce a simple document of your own.

\section{Ordinary Text}      % Produces section heading.  Lower-level
                             % sections are begun with similar
                             % \subsection and \subsubsection commands.

The ends  of words and sentences are marked
  by   spaces. It  doesn't matter how many
spaces    you type; one is as good as 100.  The
end of   a line counts as a space.

One   or more   blank lines denote the  end
of  a paragraph.

Since any number of consecutive spaces are treated
like a single one, the formatting of the input
file makes no difference to
      \LaTeX,                % The \LaTeX command generates the LaTeX logo.
but it makes a difference to you.  When you use
\LaTeX, making your input file as easy to read
as possible will be a great help as you write
your document and when you change it.  This sample
file shows how you can add comments to your own input
file.

Because printing is different from typewriting,
there are a number of things that you have to do
differently when preparing an input file than if
you were just typing the document directly.
Quotation marks like
       ``this''
have to be handled specially, as do quotes within
quotes:
       ``\,`this'            % \, separates the double and single quote.
        is what I just
        wrote, not  `that'\,''.

Dashes come in three sizes: an
       intra-word
dash, a medium dash for number ranges like
       1--2,
and a punctuation
       dash---like
this.

A sentence-ending space should be larger than the
space between words within a sentence.  You
sometimes have to type special commands in
conjunction with punctuation characters to get
this right, as in the following sentence.
       Gnats, gnus, etc.\ all  % `\ ' makes an inter-word space.
       begin with G\@.         % \@ marks end-of-sentence punctuation.
You should check the spaces after periods when
reading your output to make sure you haven't
forgotten any special cases.  Generating an
ellipsis
       \ldots\               % `\ ' is needed after `\ldots' because TeX
                             % ignores spaces after command names like \ldots
                             % made from \ + letters.
                             %
                             % Note how a `%' character causes TeX to ignore
                             % the end of the input line, so these blank lines
                             % do not start a new paragraph.
                             %
with the right spacing around the periods requires
a special command.

\LaTeX\ interprets some common characters as
commands, so you must type special commands to
generate them.  These characters include the
following:
       \$ \& \% \# \{ and \}.

In printing, text is usually emphasized with an
       \emph{italic}
type style.

\begin{em}
   A long segment of text can also be emphasized
   in this way.  Text within such a segment can be
   given \emph{additional} emphasis.
\end{em}

It is sometimes necessary to prevent \LaTeX\ from
breaking a line where it might otherwise do so.
This may be at a space, as between the ``Mr.''\ and
``Jones'' in
       ``Mr.~Jones'',        % ~ produces an unbreakable interword space.
or within a word---especially when the word is a
symbol like
       \mbox{\emph{itemnum}}
that makes little sense when hyphenated across
lines.

Footnotes\footnote{This is an example of a footnote.}
pose no problem.

\LaTeX\ is good at typesetting mathematical formulas
like
       \( x-3y + z = 7 \)
or
       \( a_{1} > x^{2n} + y^{2n} > x' \)
or
       \( \ip{A}{B} = \sum_{i} a_{i} b_{i} \).
The spaces you type in a formula are
ignored.  Remember that a letter like
       $x$                   % $ ... $  and  \( ... \)  are equivalent
is a formula when it denotes a mathematical
symbol, and it should be typed as one.

\section{Displayed Text}

Text is displayed by indenting it from the left
margin.  Quotations are commonly displayed.  There
are short quotations
\begin{quote}
   This is a short quotation.  It consists of a
   single paragraph of text.  See how it is formatted.
\end{quote}
and longer ones.
\begin{quotation}
   This is a longer quotation.  It consists of two
   paragraphs of text, neither of which are
   particularly interesting.

   This is the second paragraph of the quotation.  It
   is just as dull as the first paragraph.
\end{quotation}
Another frequently-displayed structure is a list.
The following is an example of an \emph{itemized}
list.
\begin{itemize}
   \item This is the first item of an itemized list.
         Each item in the list is marked with a ``tick''.
         You don't have to worry about what kind of tick
         mark is used.

   \item This is the second item of the list.  It
         contains another list nested inside it.  The inner
         list is an \emph{enumerated} list.
         \begin{enumerate}
            \item This is the first item of an enumerated
                  list that is nested within the itemized list.

            \item This is the second item of the inner list.
                  \LaTeX\ allows you to nest lists deeper than
                  you really should.
         \end{enumerate}
         This is the rest of the second item of the outer
         list.  It is no more interesting than any other
         part of the item.
   \item This is the third item of the list.
\end{itemize}
You can even display poetry.
\begin{verse}
   There is an environment
    for verse \\             % The \\ command separates lines
   Whose features some poets % within a stanza.
   will curse.

                             % One or more blank lines separate stanzas.

   For instead of making\\
   Them do \emph{all} line breaking, \\
   It allows them to put too many words on a line when they'd rather be
   forced to be terse.
\end{verse}

Mathematical formulas may also be displayed.  A
displayed formula
is
one-line long; multiline
formulas require special formatting instructions.
   \[  \ip{\Gamma}{\psi'} = x'' + y^{2} + z_{i}^{n}\]
Don't start a paragraph with a displayed equation,
nor make one a paragraph by itself.

\end{document}               % End of document.

\section{Bench Huge Section 697}
% synthetic bench section 697
% This is a sample LaTeX input file.  (Version of 12 August 2004.)
%
% A '%' character causes TeX to ignore all remaining text on the line,
% and is used for comments like this one.

\documentclass{article}      % Specifies the document class

                             % The preamble begins here.
\title{An Example Document}  % Declares the document's title.
\author{Leslie Lamport}      % Declares the author's name.
\date{January 21, 1994}      % Deleting this command produces today's date.

\newcommand{\ip}[2]{(#1, #2)}
                             % Defines \ip{arg1}{arg2} to mean
                             % (arg1, arg2).

%\newcommand{\ip}[2]{\langle #1 | #2\rangle}
                             % This is an alternative definition of
                             % \ip that is commented out.

\begin{document}             % End of preamble and beginning of text.

\maketitle                   % Produces the title.

This is an example input file.  Comparing it with
the output it generates can show you how to
produce a simple document of your own.

\section{Ordinary Text}      % Produces section heading.  Lower-level
                             % sections are begun with similar
                             % \subsection and \subsubsection commands.

The ends  of words and sentences are marked
  by   spaces. It  doesn't matter how many
spaces    you type; one is as good as 100.  The
end of   a line counts as a space.

One   or more   blank lines denote the  end
of  a paragraph.

Since any number of consecutive spaces are treated
like a single one, the formatting of the input
file makes no difference to
      \LaTeX,                % The \LaTeX command generates the LaTeX logo.
but it makes a difference to you.  When you use
\LaTeX, making your input file as easy to read
as possible will be a great help as you write
your document and when you change it.  This sample
file shows how you can add comments to your own input
file.

Because printing is different from typewriting,
there are a number of things that you have to do
differently when preparing an input file than if
you were just typing the document directly.
Quotation marks like
       ``this''
have to be handled specially, as do quotes within
quotes:
       ``\,`this'            % \, separates the double and single quote.
        is what I just
        wrote, not  `that'\,''.

Dashes come in three sizes: an
       intra-word
dash, a medium dash for number ranges like
       1--2,
and a punctuation
       dash---like
this.

A sentence-ending space should be larger than the
space between words within a sentence.  You
sometimes have to type special commands in
conjunction with punctuation characters to get
this right, as in the following sentence.
       Gnats, gnus, etc.\ all  % `\ ' makes an inter-word space.
       begin with G\@.         % \@ marks end-of-sentence punctuation.
You should check the spaces after periods when
reading your output to make sure you haven't
forgotten any special cases.  Generating an
ellipsis
       \ldots\               % `\ ' is needed after `\ldots' because TeX
                             % ignores spaces after command names like \ldots
                             % made from \ + letters.
                             %
                             % Note how a `%' character causes TeX to ignore
                             % the end of the input line, so these blank lines
                             % do not start a new paragraph.
                             %
with the right spacing around the periods requires
a special command.

\LaTeX\ interprets some common characters as
commands, so you must type special commands to
generate them.  These characters include the
following:
       \$ \& \% \# \{ and \}.

In printing, text is usually emphasized with an
       \emph{italic}
type style.

\begin{em}
   A long segment of text can also be emphasized
   in this way.  Text within such a segment can be
   given \emph{additional} emphasis.
\end{em}

It is sometimes necessary to prevent \LaTeX\ from
breaking a line where it might otherwise do so.
This may be at a space, as between the ``Mr.''\ and
``Jones'' in
       ``Mr.~Jones'',        % ~ produces an unbreakable interword space.
or within a word---especially when the word is a
symbol like
       \mbox{\emph{itemnum}}
that makes little sense when hyphenated across
lines.

Footnotes\footnote{This is an example of a footnote.}
pose no problem.

\LaTeX\ is good at typesetting mathematical formulas
like
       \( x-3y + z = 7 \)
or
       \( a_{1} > x^{2n} + y^{2n} > x' \)
or
       \( \ip{A}{B} = \sum_{i} a_{i} b_{i} \).
The spaces you type in a formula are
ignored.  Remember that a letter like
       $x$                   % $ ... $  and  \( ... \)  are equivalent
is a formula when it denotes a mathematical
symbol, and it should be typed as one.

\section{Displayed Text}

Text is displayed by indenting it from the left
margin.  Quotations are commonly displayed.  There
are short quotations
\begin{quote}
   This is a short quotation.  It consists of a
   single paragraph of text.  See how it is formatted.
\end{quote}
and longer ones.
\begin{quotation}
   This is a longer quotation.  It consists of two
   paragraphs of text, neither of which are
   particularly interesting.

   This is the second paragraph of the quotation.  It
   is just as dull as the first paragraph.
\end{quotation}
Another frequently-displayed structure is a list.
The following is an example of an \emph{itemized}
list.
\begin{itemize}
   \item This is the first item of an itemized list.
         Each item in the list is marked with a ``tick''.
         You don't have to worry about what kind of tick
         mark is used.

   \item This is the second item of the list.  It
         contains another list nested inside it.  The inner
         list is an \emph{enumerated} list.
         \begin{enumerate}
            \item This is the first item of an enumerated
                  list that is nested within the itemized list.

            \item This is the second item of the inner list.
                  \LaTeX\ allows you to nest lists deeper than
                  you really should.
         \end{enumerate}
         This is the rest of the second item of the outer
         list.  It is no more interesting than any other
         part of the item.
   \item This is the third item of the list.
\end{itemize}
You can even display poetry.
\begin{verse}
   There is an environment
    for verse \\             % The \\ command separates lines
   Whose features some poets % within a stanza.
   will curse.

                             % One or more blank lines separate stanzas.

   For instead of making\\
   Them do \emph{all} line breaking, \\
   It allows them to put too many words on a line when they'd rather be
   forced to be terse.
\end{verse}

Mathematical formulas may also be displayed.  A
displayed formula
is
one-line long; multiline
formulas require special formatting instructions.
   \[  \ip{\Gamma}{\psi'} = x'' + y^{2} + z_{i}^{n}\]
Don't start a paragraph with a displayed equation,
nor make one a paragraph by itself.

\end{document}               % End of document.

\section{Bench Huge Section 698}
% synthetic bench section 698
% This is a sample LaTeX input file.  (Version of 12 August 2004.)
%
% A '%' character causes TeX to ignore all remaining text on the line,
% and is used for comments like this one.

\documentclass{article}      % Specifies the document class

                             % The preamble begins here.
\title{An Example Document}  % Declares the document's title.
\author{Leslie Lamport}      % Declares the author's name.
\date{January 21, 1994}      % Deleting this command produces today's date.

\newcommand{\ip}[2]{(#1, #2)}
                             % Defines \ip{arg1}{arg2} to mean
                             % (arg1, arg2).

%\newcommand{\ip}[2]{\langle #1 | #2\rangle}
                             % This is an alternative definition of
                             % \ip that is commented out.

\begin{document}             % End of preamble and beginning of text.

\maketitle                   % Produces the title.

This is an example input file.  Comparing it with
the output it generates can show you how to
produce a simple document of your own.

\section{Ordinary Text}      % Produces section heading.  Lower-level
                             % sections are begun with similar
                             % \subsection and \subsubsection commands.

The ends  of words and sentences are marked
  by   spaces. It  doesn't matter how many
spaces    you type; one is as good as 100.  The
end of   a line counts as a space.

One   or more   blank lines denote the  end
of  a paragraph.

Since any number of consecutive spaces are treated
like a single one, the formatting of the input
file makes no difference to
      \LaTeX,                % The \LaTeX command generates the LaTeX logo.
but it makes a difference to you.  When you use
\LaTeX, making your input file as easy to read
as possible will be a great help as you write
your document and when you change it.  This sample
file shows how you can add comments to your own input
file.

Because printing is different from typewriting,
there are a number of things that you have to do
differently when preparing an input file than if
you were just typing the document directly.
Quotation marks like
       ``this''
have to be handled specially, as do quotes within
quotes:
       ``\,`this'            % \, separates the double and single quote.
        is what I just
        wrote, not  `that'\,''.

Dashes come in three sizes: an
       intra-word
dash, a medium dash for number ranges like
       1--2,
and a punctuation
       dash---like
this.

A sentence-ending space should be larger than the
space between words within a sentence.  You
sometimes have to type special commands in
conjunction with punctuation characters to get
this right, as in the following sentence.
       Gnats, gnus, etc.\ all  % `\ ' makes an inter-word space.
       begin with G\@.         % \@ marks end-of-sentence punctuation.
You should check the spaces after periods when
reading your output to make sure you haven't
forgotten any special cases.  Generating an
ellipsis
       \ldots\               % `\ ' is needed after `\ldots' because TeX
                             % ignores spaces after command names like \ldots
                             % made from \ + letters.
                             %
                             % Note how a `%' character causes TeX to ignore
                             % the end of the input line, so these blank lines
                             % do not start a new paragraph.
                             %
with the right spacing around the periods requires
a special command.

\LaTeX\ interprets some common characters as
commands, so you must type special commands to
generate them.  These characters include the
following:
       \$ \& \% \# \{ and \}.

In printing, text is usually emphasized with an
       \emph{italic}
type style.

\begin{em}
   A long segment of text can also be emphasized
   in this way.  Text within such a segment can be
   given \emph{additional} emphasis.
\end{em}

It is sometimes necessary to prevent \LaTeX\ from
breaking a line where it might otherwise do so.
This may be at a space, as between the ``Mr.''\ and
``Jones'' in
       ``Mr.~Jones'',        % ~ produces an unbreakable interword space.
or within a word---especially when the word is a
symbol like
       \mbox{\emph{itemnum}}
that makes little sense when hyphenated across
lines.

Footnotes\footnote{This is an example of a footnote.}
pose no problem.

\LaTeX\ is good at typesetting mathematical formulas
like
       \( x-3y + z = 7 \)
or
       \( a_{1} > x^{2n} + y^{2n} > x' \)
or
       \( \ip{A}{B} = \sum_{i} a_{i} b_{i} \).
The spaces you type in a formula are
ignored.  Remember that a letter like
       $x$                   % $ ... $  and  \( ... \)  are equivalent
is a formula when it denotes a mathematical
symbol, and it should be typed as one.

\section{Displayed Text}

Text is displayed by indenting it from the left
margin.  Quotations are commonly displayed.  There
are short quotations
\begin{quote}
   This is a short quotation.  It consists of a
   single paragraph of text.  See how it is formatted.
\end{quote}
and longer ones.
\begin{quotation}
   This is a longer quotation.  It consists of two
   paragraphs of text, neither of which are
   particularly interesting.

   This is the second paragraph of the quotation.  It
   is just as dull as the first paragraph.
\end{quotation}
Another frequently-displayed structure is a list.
The following is an example of an \emph{itemized}
list.
\begin{itemize}
   \item This is the first item of an itemized list.
         Each item in the list is marked with a ``tick''.
         You don't have to worry about what kind of tick
         mark is used.

   \item This is the second item of the list.  It
         contains another list nested inside it.  The inner
         list is an \emph{enumerated} list.
         \begin{enumerate}
            \item This is the first item of an enumerated
                  list that is nested within the itemized list.

            \item This is the second item of the inner list.
                  \LaTeX\ allows you to nest lists deeper than
                  you really should.
         \end{enumerate}
         This is the rest of the second item of the outer
         list.  It is no more interesting than any other
         part of the item.
   \item This is the third item of the list.
\end{itemize}
You can even display poetry.
\begin{verse}
   There is an environment
    for verse \\             % The \\ command separates lines
   Whose features some poets % within a stanza.
   will curse.

                             % One or more blank lines separate stanzas.

   For instead of making\\
   Them do \emph{all} line breaking, \\
   It allows them to put too many words on a line when they'd rather be
   forced to be terse.
\end{verse}

Mathematical formulas may also be displayed.  A
displayed formula
is
one-line long; multiline
formulas require special formatting instructions.
   \[  \ip{\Gamma}{\psi'} = x'' + y^{2} + z_{i}^{n}\]
Don't start a paragraph with a displayed equation,
nor make one a paragraph by itself.

\end{document}               % End of document.

\section{Bench Huge Section 699}
% synthetic bench section 699
% This is a sample LaTeX input file.  (Version of 12 August 2004.)
%
% A '%' character causes TeX to ignore all remaining text on the line,
% and is used for comments like this one.

\documentclass{article}      % Specifies the document class

                             % The preamble begins here.
\title{An Example Document}  % Declares the document's title.
\author{Leslie Lamport}      % Declares the author's name.
\date{January 21, 1994}      % Deleting this command produces today's date.

\newcommand{\ip}[2]{(#1, #2)}
                             % Defines \ip{arg1}{arg2} to mean
                             % (arg1, arg2).

%\newcommand{\ip}[2]{\langle #1 | #2\rangle}
                             % This is an alternative definition of
                             % \ip that is commented out.

\begin{document}             % End of preamble and beginning of text.

\maketitle                   % Produces the title.

This is an example input file.  Comparing it with
the output it generates can show you how to
produce a simple document of your own.

\section{Ordinary Text}      % Produces section heading.  Lower-level
                             % sections are begun with similar
                             % \subsection and \subsubsection commands.

The ends  of words and sentences are marked
  by   spaces. It  doesn't matter how many
spaces    you type; one is as good as 100.  The
end of   a line counts as a space.

One   or more   blank lines denote the  end
of  a paragraph.

Since any number of consecutive spaces are treated
like a single one, the formatting of the input
file makes no difference to
      \LaTeX,                % The \LaTeX command generates the LaTeX logo.
but it makes a difference to you.  When you use
\LaTeX, making your input file as easy to read
as possible will be a great help as you write
your document and when you change it.  This sample
file shows how you can add comments to your own input
file.

Because printing is different from typewriting,
there are a number of things that you have to do
differently when preparing an input file than if
you were just typing the document directly.
Quotation marks like
       ``this''
have to be handled specially, as do quotes within
quotes:
       ``\,`this'            % \, separates the double and single quote.
        is what I just
        wrote, not  `that'\,''.

Dashes come in three sizes: an
       intra-word
dash, a medium dash for number ranges like
       1--2,
and a punctuation
       dash---like
this.

A sentence-ending space should be larger than the
space between words within a sentence.  You
sometimes have to type special commands in
conjunction with punctuation characters to get
this right, as in the following sentence.
       Gnats, gnus, etc.\ all  % `\ ' makes an inter-word space.
       begin with G\@.         % \@ marks end-of-sentence punctuation.
You should check the spaces after periods when
reading your output to make sure you haven't
forgotten any special cases.  Generating an
ellipsis
       \ldots\               % `\ ' is needed after `\ldots' because TeX
                             % ignores spaces after command names like \ldots
                             % made from \ + letters.
                             %
                             % Note how a `%' character causes TeX to ignore
                             % the end of the input line, so these blank lines
                             % do not start a new paragraph.
                             %
with the right spacing around the periods requires
a special command.

\LaTeX\ interprets some common characters as
commands, so you must type special commands to
generate them.  These characters include the
following:
       \$ \& \% \# \{ and \}.

In printing, text is usually emphasized with an
       \emph{italic}
type style.

\begin{em}
   A long segment of text can also be emphasized
   in this way.  Text within such a segment can be
   given \emph{additional} emphasis.
\end{em}

It is sometimes necessary to prevent \LaTeX\ from
breaking a line where it might otherwise do so.
This may be at a space, as between the ``Mr.''\ and
``Jones'' in
       ``Mr.~Jones'',        % ~ produces an unbreakable interword space.
or within a word---especially when the word is a
symbol like
       \mbox{\emph{itemnum}}
that makes little sense when hyphenated across
lines.

Footnotes\footnote{This is an example of a footnote.}
pose no problem.

\LaTeX\ is good at typesetting mathematical formulas
like
       \( x-3y + z = 7 \)
or
       \( a_{1} > x^{2n} + y^{2n} > x' \)
or
       \( \ip{A}{B} = \sum_{i} a_{i} b_{i} \).
The spaces you type in a formula are
ignored.  Remember that a letter like
       $x$                   % $ ... $  and  \( ... \)  are equivalent
is a formula when it denotes a mathematical
symbol, and it should be typed as one.

\section{Displayed Text}

Text is displayed by indenting it from the left
margin.  Quotations are commonly displayed.  There
are short quotations
\begin{quote}
   This is a short quotation.  It consists of a
   single paragraph of text.  See how it is formatted.
\end{quote}
and longer ones.
\begin{quotation}
   This is a longer quotation.  It consists of two
   paragraphs of text, neither of which are
   particularly interesting.

   This is the second paragraph of the quotation.  It
   is just as dull as the first paragraph.
\end{quotation}
Another frequently-displayed structure is a list.
The following is an example of an \emph{itemized}
list.
\begin{itemize}
   \item This is the first item of an itemized list.
         Each item in the list is marked with a ``tick''.
         You don't have to worry about what kind of tick
         mark is used.

   \item This is the second item of the list.  It
         contains another list nested inside it.  The inner
         list is an \emph{enumerated} list.
         \begin{enumerate}
            \item This is the first item of an enumerated
                  list that is nested within the itemized list.

            \item This is the second item of the inner list.
                  \LaTeX\ allows you to nest lists deeper than
                  you really should.
         \end{enumerate}
         This is the rest of the second item of the outer
         list.  It is no more interesting than any other
         part of the item.
   \item This is the third item of the list.
\end{itemize}
You can even display poetry.
\begin{verse}
   There is an environment
    for verse \\             % The \\ command separates lines
   Whose features some poets % within a stanza.
   will curse.

                             % One or more blank lines separate stanzas.

   For instead of making\\
   Them do \emph{all} line breaking, \\
   It allows them to put too many words on a line when they'd rather be
   forced to be terse.
\end{verse}

Mathematical formulas may also be displayed.  A
displayed formula
is
one-line long; multiline
formulas require special formatting instructions.
   \[  \ip{\Gamma}{\psi'} = x'' + y^{2} + z_{i}^{n}\]
Don't start a paragraph with a displayed equation,
nor make one a paragraph by itself.

\end{document}               % End of document.

\section{Bench Huge Section 700}
% synthetic bench section 700
% This is a sample LaTeX input file.  (Version of 12 August 2004.)
%
% A '%' character causes TeX to ignore all remaining text on the line,
% and is used for comments like this one.

\documentclass{article}      % Specifies the document class

                             % The preamble begins here.
\title{An Example Document}  % Declares the document's title.
\author{Leslie Lamport}      % Declares the author's name.
\date{January 21, 1994}      % Deleting this command produces today's date.

\newcommand{\ip}[2]{(#1, #2)}
                             % Defines \ip{arg1}{arg2} to mean
                             % (arg1, arg2).

%\newcommand{\ip}[2]{\langle #1 | #2\rangle}
                             % This is an alternative definition of
                             % \ip that is commented out.

\begin{document}             % End of preamble and beginning of text.

\maketitle                   % Produces the title.

This is an example input file.  Comparing it with
the output it generates can show you how to
produce a simple document of your own.

\section{Ordinary Text}      % Produces section heading.  Lower-level
                             % sections are begun with similar
                             % \subsection and \subsubsection commands.

The ends  of words and sentences are marked
  by   spaces. It  doesn't matter how many
spaces    you type; one is as good as 100.  The
end of   a line counts as a space.

One   or more   blank lines denote the  end
of  a paragraph.

Since any number of consecutive spaces are treated
like a single one, the formatting of the input
file makes no difference to
      \LaTeX,                % The \LaTeX command generates the LaTeX logo.
but it makes a difference to you.  When you use
\LaTeX, making your input file as easy to read
as possible will be a great help as you write
your document and when you change it.  This sample
file shows how you can add comments to your own input
file.

Because printing is different from typewriting,
there are a number of things that you have to do
differently when preparing an input file than if
you were just typing the document directly.
Quotation marks like
       ``this''
have to be handled specially, as do quotes within
quotes:
       ``\,`this'            % \, separates the double and single quote.
        is what I just
        wrote, not  `that'\,''.

Dashes come in three sizes: an
       intra-word
dash, a medium dash for number ranges like
       1--2,
and a punctuation
       dash---like
this.

A sentence-ending space should be larger than the
space between words within a sentence.  You
sometimes have to type special commands in
conjunction with punctuation characters to get
this right, as in the following sentence.
       Gnats, gnus, etc.\ all  % `\ ' makes an inter-word space.
       begin with G\@.         % \@ marks end-of-sentence punctuation.
You should check the spaces after periods when
reading your output to make sure you haven't
forgotten any special cases.  Generating an
ellipsis
       \ldots\               % `\ ' is needed after `\ldots' because TeX
                             % ignores spaces after command names like \ldots
                             % made from \ + letters.
                             %
                             % Note how a `%' character causes TeX to ignore
                             % the end of the input line, so these blank lines
                             % do not start a new paragraph.
                             %
with the right spacing around the periods requires
a special command.

\LaTeX\ interprets some common characters as
commands, so you must type special commands to
generate them.  These characters include the
following:
       \$ \& \% \# \{ and \}.

In printing, text is usually emphasized with an
       \emph{italic}
type style.

\begin{em}
   A long segment of text can also be emphasized
   in this way.  Text within such a segment can be
   given \emph{additional} emphasis.
\end{em}

It is sometimes necessary to prevent \LaTeX\ from
breaking a line where it might otherwise do so.
This may be at a space, as between the ``Mr.''\ and
``Jones'' in
       ``Mr.~Jones'',        % ~ produces an unbreakable interword space.
or within a word---especially when the word is a
symbol like
       \mbox{\emph{itemnum}}
that makes little sense when hyphenated across
lines.

Footnotes\footnote{This is an example of a footnote.}
pose no problem.

\LaTeX\ is good at typesetting mathematical formulas
like
       \( x-3y + z = 7 \)
or
       \( a_{1} > x^{2n} + y^{2n} > x' \)
or
       \( \ip{A}{B} = \sum_{i} a_{i} b_{i} \).
The spaces you type in a formula are
ignored.  Remember that a letter like
       $x$                   % $ ... $  and  \( ... \)  are equivalent
is a formula when it denotes a mathematical
symbol, and it should be typed as one.

\section{Displayed Text}

Text is displayed by indenting it from the left
margin.  Quotations are commonly displayed.  There
are short quotations
\begin{quote}
   This is a short quotation.  It consists of a
   single paragraph of text.  See how it is formatted.
\end{quote}
and longer ones.
\begin{quotation}
   This is a longer quotation.  It consists of two
   paragraphs of text, neither of which are
   particularly interesting.

   This is the second paragraph of the quotation.  It
   is just as dull as the first paragraph.
\end{quotation}
Another frequently-displayed structure is a list.
The following is an example of an \emph{itemized}
list.
\begin{itemize}
   \item This is the first item of an itemized list.
         Each item in the list is marked with a ``tick''.
         You don't have to worry about what kind of tick
         mark is used.

   \item This is the second item of the list.  It
         contains another list nested inside it.  The inner
         list is an \emph{enumerated} list.
         \begin{enumerate}
            \item This is the first item of an enumerated
                  list that is nested within the itemized list.

            \item This is the second item of the inner list.
                  \LaTeX\ allows you to nest lists deeper than
                  you really should.
         \end{enumerate}
         This is the rest of the second item of the outer
         list.  It is no more interesting than any other
         part of the item.
   \item This is the third item of the list.
\end{itemize}
You can even display poetry.
\begin{verse}
   There is an environment
    for verse \\             % The \\ command separates lines
   Whose features some poets % within a stanza.
   will curse.

                             % One or more blank lines separate stanzas.

   For instead of making\\
   Them do \emph{all} line breaking, \\
   It allows them to put too many words on a line when they'd rather be
   forced to be terse.
\end{verse}

Mathematical formulas may also be displayed.  A
displayed formula
is
one-line long; multiline
formulas require special formatting instructions.
   \[  \ip{\Gamma}{\psi'} = x'' + y^{2} + z_{i}^{n}\]
Don't start a paragraph with a displayed equation,
nor make one a paragraph by itself.

\end{document}               % End of document.

\section{Bench Huge Section 701}
% synthetic bench section 701
% This is a sample LaTeX input file.  (Version of 12 August 2004.)
%
% A '%' character causes TeX to ignore all remaining text on the line,
% and is used for comments like this one.

\documentclass{article}      % Specifies the document class

                             % The preamble begins here.
\title{An Example Document}  % Declares the document's title.
\author{Leslie Lamport}      % Declares the author's name.
\date{January 21, 1994}      % Deleting this command produces today's date.

\newcommand{\ip}[2]{(#1, #2)}
                             % Defines \ip{arg1}{arg2} to mean
                             % (arg1, arg2).

%\newcommand{\ip}[2]{\langle #1 | #2\rangle}
                             % This is an alternative definition of
                             % \ip that is commented out.

\begin{document}             % End of preamble and beginning of text.

\maketitle                   % Produces the title.

This is an example input file.  Comparing it with
the output it generates can show you how to
produce a simple document of your own.

\section{Ordinary Text}      % Produces section heading.  Lower-level
                             % sections are begun with similar
                             % \subsection and \subsubsection commands.

The ends  of words and sentences are marked
  by   spaces. It  doesn't matter how many
spaces    you type; one is as good as 100.  The
end of   a line counts as a space.

One   or more   blank lines denote the  end
of  a paragraph.

Since any number of consecutive spaces are treated
like a single one, the formatting of the input
file makes no difference to
      \LaTeX,                % The \LaTeX command generates the LaTeX logo.
but it makes a difference to you.  When you use
\LaTeX, making your input file as easy to read
as possible will be a great help as you write
your document and when you change it.  This sample
file shows how you can add comments to your own input
file.

Because printing is different from typewriting,
there are a number of things that you have to do
differently when preparing an input file than if
you were just typing the document directly.
Quotation marks like
       ``this''
have to be handled specially, as do quotes within
quotes:
       ``\,`this'            % \, separates the double and single quote.
        is what I just
        wrote, not  `that'\,''.

Dashes come in three sizes: an
       intra-word
dash, a medium dash for number ranges like
       1--2,
and a punctuation
       dash---like
this.

A sentence-ending space should be larger than the
space between words within a sentence.  You
sometimes have to type special commands in
conjunction with punctuation characters to get
this right, as in the following sentence.
       Gnats, gnus, etc.\ all  % `\ ' makes an inter-word space.
       begin with G\@.         % \@ marks end-of-sentence punctuation.
You should check the spaces after periods when
reading your output to make sure you haven't
forgotten any special cases.  Generating an
ellipsis
       \ldots\               % `\ ' is needed after `\ldots' because TeX
                             % ignores spaces after command names like \ldots
                             % made from \ + letters.
                             %
                             % Note how a `%' character causes TeX to ignore
                             % the end of the input line, so these blank lines
                             % do not start a new paragraph.
                             %
with the right spacing around the periods requires
a special command.

\LaTeX\ interprets some common characters as
commands, so you must type special commands to
generate them.  These characters include the
following:
       \$ \& \% \# \{ and \}.

In printing, text is usually emphasized with an
       \emph{italic}
type style.

\begin{em}
   A long segment of text can also be emphasized
   in this way.  Text within such a segment can be
   given \emph{additional} emphasis.
\end{em}

It is sometimes necessary to prevent \LaTeX\ from
breaking a line where it might otherwise do so.
This may be at a space, as between the ``Mr.''\ and
``Jones'' in
       ``Mr.~Jones'',        % ~ produces an unbreakable interword space.
or within a word---especially when the word is a
symbol like
       \mbox{\emph{itemnum}}
that makes little sense when hyphenated across
lines.

Footnotes\footnote{This is an example of a footnote.}
pose no problem.

\LaTeX\ is good at typesetting mathematical formulas
like
       \( x-3y + z = 7 \)
or
       \( a_{1} > x^{2n} + y^{2n} > x' \)
or
       \( \ip{A}{B} = \sum_{i} a_{i} b_{i} \).
The spaces you type in a formula are
ignored.  Remember that a letter like
       $x$                   % $ ... $  and  \( ... \)  are equivalent
is a formula when it denotes a mathematical
symbol, and it should be typed as one.

\section{Displayed Text}

Text is displayed by indenting it from the left
margin.  Quotations are commonly displayed.  There
are short quotations
\begin{quote}
   This is a short quotation.  It consists of a
   single paragraph of text.  See how it is formatted.
\end{quote}
and longer ones.
\begin{quotation}
   This is a longer quotation.  It consists of two
   paragraphs of text, neither of which are
   particularly interesting.

   This is the second paragraph of the quotation.  It
   is just as dull as the first paragraph.
\end{quotation}
Another frequently-displayed structure is a list.
The following is an example of an \emph{itemized}
list.
\begin{itemize}
   \item This is the first item of an itemized list.
         Each item in the list is marked with a ``tick''.
         You don't have to worry about what kind of tick
         mark is used.

   \item This is the second item of the list.  It
         contains another list nested inside it.  The inner
         list is an \emph{enumerated} list.
         \begin{enumerate}
            \item This is the first item of an enumerated
                  list that is nested within the itemized list.

            \item This is the second item of the inner list.
                  \LaTeX\ allows you to nest lists deeper than
                  you really should.
         \end{enumerate}
         This is the rest of the second item of the outer
         list.  It is no more interesting than any other
         part of the item.
   \item This is the third item of the list.
\end{itemize}
You can even display poetry.
\begin{verse}
   There is an environment
    for verse \\             % The \\ command separates lines
   Whose features some poets % within a stanza.
   will curse.

                             % One or more blank lines separate stanzas.

   For instead of making\\
   Them do \emph{all} line breaking, \\
   It allows them to put too many words on a line when they'd rather be
   forced to be terse.
\end{verse}

Mathematical formulas may also be displayed.  A
displayed formula
is
one-line long; multiline
formulas require special formatting instructions.
   \[  \ip{\Gamma}{\psi'} = x'' + y^{2} + z_{i}^{n}\]
Don't start a paragraph with a displayed equation,
nor make one a paragraph by itself.

\end{document}               % End of document.

\section{Bench Huge Section 702}
% synthetic bench section 702
% This is a sample LaTeX input file.  (Version of 12 August 2004.)
%
% A '%' character causes TeX to ignore all remaining text on the line,
% and is used for comments like this one.

\documentclass{article}      % Specifies the document class

                             % The preamble begins here.
\title{An Example Document}  % Declares the document's title.
\author{Leslie Lamport}      % Declares the author's name.
\date{January 21, 1994}      % Deleting this command produces today's date.

\newcommand{\ip}[2]{(#1, #2)}
                             % Defines \ip{arg1}{arg2} to mean
                             % (arg1, arg2).

%\newcommand{\ip}[2]{\langle #1 | #2\rangle}
                             % This is an alternative definition of
                             % \ip that is commented out.

\begin{document}             % End of preamble and beginning of text.

\maketitle                   % Produces the title.

This is an example input file.  Comparing it with
the output it generates can show you how to
produce a simple document of your own.

\section{Ordinary Text}      % Produces section heading.  Lower-level
                             % sections are begun with similar
                             % \subsection and \subsubsection commands.

The ends  of words and sentences are marked
  by   spaces. It  doesn't matter how many
spaces    you type; one is as good as 100.  The
end of   a line counts as a space.

One   or more   blank lines denote the  end
of  a paragraph.

Since any number of consecutive spaces are treated
like a single one, the formatting of the input
file makes no difference to
      \LaTeX,                % The \LaTeX command generates the LaTeX logo.
but it makes a difference to you.  When you use
\LaTeX, making your input file as easy to read
as possible will be a great help as you write
your document and when you change it.  This sample
file shows how you can add comments to your own input
file.

Because printing is different from typewriting,
there are a number of things that you have to do
differently when preparing an input file than if
you were just typing the document directly.
Quotation marks like
       ``this''
have to be handled specially, as do quotes within
quotes:
       ``\,`this'            % \, separates the double and single quote.
        is what I just
        wrote, not  `that'\,''.

Dashes come in three sizes: an
       intra-word
dash, a medium dash for number ranges like
       1--2,
and a punctuation
       dash---like
this.

A sentence-ending space should be larger than the
space between words within a sentence.  You
sometimes have to type special commands in
conjunction with punctuation characters to get
this right, as in the following sentence.
       Gnats, gnus, etc.\ all  % `\ ' makes an inter-word space.
       begin with G\@.         % \@ marks end-of-sentence punctuation.
You should check the spaces after periods when
reading your output to make sure you haven't
forgotten any special cases.  Generating an
ellipsis
       \ldots\               % `\ ' is needed after `\ldots' because TeX
                             % ignores spaces after command names like \ldots
                             % made from \ + letters.
                             %
                             % Note how a `%' character causes TeX to ignore
                             % the end of the input line, so these blank lines
                             % do not start a new paragraph.
                             %
with the right spacing around the periods requires
a special command.

\LaTeX\ interprets some common characters as
commands, so you must type special commands to
generate them.  These characters include the
following:
       \$ \& \% \# \{ and \}.

In printing, text is usually emphasized with an
       \emph{italic}
type style.

\begin{em}
   A long segment of text can also be emphasized
   in this way.  Text within such a segment can be
   given \emph{additional} emphasis.
\end{em}

It is sometimes necessary to prevent \LaTeX\ from
breaking a line where it might otherwise do so.
This may be at a space, as between the ``Mr.''\ and
``Jones'' in
       ``Mr.~Jones'',        % ~ produces an unbreakable interword space.
or within a word---especially when the word is a
symbol like
       \mbox{\emph{itemnum}}
that makes little sense when hyphenated across
lines.

Footnotes\footnote{This is an example of a footnote.}
pose no problem.

\LaTeX\ is good at typesetting mathematical formulas
like
       \( x-3y + z = 7 \)
or
       \( a_{1} > x^{2n} + y^{2n} > x' \)
or
       \( \ip{A}{B} = \sum_{i} a_{i} b_{i} \).
The spaces you type in a formula are
ignored.  Remember that a letter like
       $x$                   % $ ... $  and  \( ... \)  are equivalent
is a formula when it denotes a mathematical
symbol, and it should be typed as one.

\section{Displayed Text}

Text is displayed by indenting it from the left
margin.  Quotations are commonly displayed.  There
are short quotations
\begin{quote}
   This is a short quotation.  It consists of a
   single paragraph of text.  See how it is formatted.
\end{quote}
and longer ones.
\begin{quotation}
   This is a longer quotation.  It consists of two
   paragraphs of text, neither of which are
   particularly interesting.

   This is the second paragraph of the quotation.  It
   is just as dull as the first paragraph.
\end{quotation}
Another frequently-displayed structure is a list.
The following is an example of an \emph{itemized}
list.
\begin{itemize}
   \item This is the first item of an itemized list.
         Each item in the list is marked with a ``tick''.
         You don't have to worry about what kind of tick
         mark is used.

   \item This is the second item of the list.  It
         contains another list nested inside it.  The inner
         list is an \emph{enumerated} list.
         \begin{enumerate}
            \item This is the first item of an enumerated
                  list that is nested within the itemized list.

            \item This is the second item of the inner list.
                  \LaTeX\ allows you to nest lists deeper than
                  you really should.
         \end{enumerate}
         This is the rest of the second item of the outer
         list.  It is no more interesting than any other
         part of the item.
   \item This is the third item of the list.
\end{itemize}
You can even display poetry.
\begin{verse}
   There is an environment
    for verse \\             % The \\ command separates lines
   Whose features some poets % within a stanza.
   will curse.

                             % One or more blank lines separate stanzas.

   For instead of making\\
   Them do \emph{all} line breaking, \\
   It allows them to put too many words on a line when they'd rather be
   forced to be terse.
\end{verse}

Mathematical formulas may also be displayed.  A
displayed formula
is
one-line long; multiline
formulas require special formatting instructions.
   \[  \ip{\Gamma}{\psi'} = x'' + y^{2} + z_{i}^{n}\]
Don't start a paragraph with a displayed equation,
nor make one a paragraph by itself.

\end{document}               % End of document.

\section{Bench Huge Section 703}
% synthetic bench section 703
% This is a sample LaTeX input file.  (Version of 12 August 2004.)
%
% A '%' character causes TeX to ignore all remaining text on the line,
% and is used for comments like this one.

\documentclass{article}      % Specifies the document class

                             % The preamble begins here.
\title{An Example Document}  % Declares the document's title.
\author{Leslie Lamport}      % Declares the author's name.
\date{January 21, 1994}      % Deleting this command produces today's date.

\newcommand{\ip}[2]{(#1, #2)}
                             % Defines \ip{arg1}{arg2} to mean
                             % (arg1, arg2).

%\newcommand{\ip}[2]{\langle #1 | #2\rangle}
                             % This is an alternative definition of
                             % \ip that is commented out.

\begin{document}             % End of preamble and beginning of text.

\maketitle                   % Produces the title.

This is an example input file.  Comparing it with
the output it generates can show you how to
produce a simple document of your own.

\section{Ordinary Text}      % Produces section heading.  Lower-level
                             % sections are begun with similar
                             % \subsection and \subsubsection commands.

The ends  of words and sentences are marked
  by   spaces. It  doesn't matter how many
spaces    you type; one is as good as 100.  The
end of   a line counts as a space.

One   or more   blank lines denote the  end
of  a paragraph.

Since any number of consecutive spaces are treated
like a single one, the formatting of the input
file makes no difference to
      \LaTeX,                % The \LaTeX command generates the LaTeX logo.
but it makes a difference to you.  When you use
\LaTeX, making your input file as easy to read
as possible will be a great help as you write
your document and when you change it.  This sample
file shows how you can add comments to your own input
file.

Because printing is different from typewriting,
there are a number of things that you have to do
differently when preparing an input file than if
you were just typing the document directly.
Quotation marks like
       ``this''
have to be handled specially, as do quotes within
quotes:
       ``\,`this'            % \, separates the double and single quote.
        is what I just
        wrote, not  `that'\,''.

Dashes come in three sizes: an
       intra-word
dash, a medium dash for number ranges like
       1--2,
and a punctuation
       dash---like
this.

A sentence-ending space should be larger than the
space between words within a sentence.  You
sometimes have to type special commands in
conjunction with punctuation characters to get
this right, as in the following sentence.
       Gnats, gnus, etc.\ all  % `\ ' makes an inter-word space.
       begin with G\@.         % \@ marks end-of-sentence punctuation.
You should check the spaces after periods when
reading your output to make sure you haven't
forgotten any special cases.  Generating an
ellipsis
       \ldots\               % `\ ' is needed after `\ldots' because TeX
                             % ignores spaces after command names like \ldots
                             % made from \ + letters.
                             %
                             % Note how a `%' character causes TeX to ignore
                             % the end of the input line, so these blank lines
                             % do not start a new paragraph.
                             %
with the right spacing around the periods requires
a special command.

\LaTeX\ interprets some common characters as
commands, so you must type special commands to
generate them.  These characters include the
following:
       \$ \& \% \# \{ and \}.

In printing, text is usually emphasized with an
       \emph{italic}
type style.

\begin{em}
   A long segment of text can also be emphasized
   in this way.  Text within such a segment can be
   given \emph{additional} emphasis.
\end{em}

It is sometimes necessary to prevent \LaTeX\ from
breaking a line where it might otherwise do so.
This may be at a space, as between the ``Mr.''\ and
``Jones'' in
       ``Mr.~Jones'',        % ~ produces an unbreakable interword space.
or within a word---especially when the word is a
symbol like
       \mbox{\emph{itemnum}}
that makes little sense when hyphenated across
lines.

Footnotes\footnote{This is an example of a footnote.}
pose no problem.

\LaTeX\ is good at typesetting mathematical formulas
like
       \( x-3y + z = 7 \)
or
       \( a_{1} > x^{2n} + y^{2n} > x' \)
or
       \( \ip{A}{B} = \sum_{i} a_{i} b_{i} \).
The spaces you type in a formula are
ignored.  Remember that a letter like
       $x$                   % $ ... $  and  \( ... \)  are equivalent
is a formula when it denotes a mathematical
symbol, and it should be typed as one.

\section{Displayed Text}

Text is displayed by indenting it from the left
margin.  Quotations are commonly displayed.  There
are short quotations
\begin{quote}
   This is a short quotation.  It consists of a
   single paragraph of text.  See how it is formatted.
\end{quote}
and longer ones.
\begin{quotation}
   This is a longer quotation.  It consists of two
   paragraphs of text, neither of which are
   particularly interesting.

   This is the second paragraph of the quotation.  It
   is just as dull as the first paragraph.
\end{quotation}
Another frequently-displayed structure is a list.
The following is an example of an \emph{itemized}
list.
\begin{itemize}
   \item This is the first item of an itemized list.
         Each item in the list is marked with a ``tick''.
         You don't have to worry about what kind of tick
         mark is used.

   \item This is the second item of the list.  It
         contains another list nested inside it.  The inner
         list is an \emph{enumerated} list.
         \begin{enumerate}
            \item This is the first item of an enumerated
                  list that is nested within the itemized list.

            \item This is the second item of the inner list.
                  \LaTeX\ allows you to nest lists deeper than
                  you really should.
         \end{enumerate}
         This is the rest of the second item of the outer
         list.  It is no more interesting than any other
         part of the item.
   \item This is the third item of the list.
\end{itemize}
You can even display poetry.
\begin{verse}
   There is an environment
    for verse \\             % The \\ command separates lines
   Whose features some poets % within a stanza.
   will curse.

                             % One or more blank lines separate stanzas.

   For instead of making\\
   Them do \emph{all} line breaking, \\
   It allows them to put too many words on a line when they'd rather be
   forced to be terse.
\end{verse}

Mathematical formulas may also be displayed.  A
displayed formula
is
one-line long; multiline
formulas require special formatting instructions.
   \[  \ip{\Gamma}{\psi'} = x'' + y^{2} + z_{i}^{n}\]
Don't start a paragraph with a displayed equation,
nor make one a paragraph by itself.

\end{document}               % End of document.

\section{Bench Huge Section 704}
% synthetic bench section 704
% This is a sample LaTeX input file.  (Version of 12 August 2004.)
%
% A '%' character causes TeX to ignore all remaining text on the line,
% and is used for comments like this one.

\documentclass{article}      % Specifies the document class

                             % The preamble begins here.
\title{An Example Document}  % Declares the document's title.
\author{Leslie Lamport}      % Declares the author's name.
\date{January 21, 1994}      % Deleting this command produces today's date.

\newcommand{\ip}[2]{(#1, #2)}
                             % Defines \ip{arg1}{arg2} to mean
                             % (arg1, arg2).

%\newcommand{\ip}[2]{\langle #1 | #2\rangle}
                             % This is an alternative definition of
                             % \ip that is commented out.

\begin{document}             % End of preamble and beginning of text.

\maketitle                   % Produces the title.

This is an example input file.  Comparing it with
the output it generates can show you how to
produce a simple document of your own.

\section{Ordinary Text}      % Produces section heading.  Lower-level
                             % sections are begun with similar
                             % \subsection and \subsubsection commands.

The ends  of words and sentences are marked
  by   spaces. It  doesn't matter how many
spaces    you type; one is as good as 100.  The
end of   a line counts as a space.

One   or more   blank lines denote the  end
of  a paragraph.

Since any number of consecutive spaces are treated
like a single one, the formatting of the input
file makes no difference to
      \LaTeX,                % The \LaTeX command generates the LaTeX logo.
but it makes a difference to you.  When you use
\LaTeX, making your input file as easy to read
as possible will be a great help as you write
your document and when you change it.  This sample
file shows how you can add comments to your own input
file.

Because printing is different from typewriting,
there are a number of things that you have to do
differently when preparing an input file than if
you were just typing the document directly.
Quotation marks like
       ``this''
have to be handled specially, as do quotes within
quotes:
       ``\,`this'            % \, separates the double and single quote.
        is what I just
        wrote, not  `that'\,''.

Dashes come in three sizes: an
       intra-word
dash, a medium dash for number ranges like
       1--2,
and a punctuation
       dash---like
this.

A sentence-ending space should be larger than the
space between words within a sentence.  You
sometimes have to type special commands in
conjunction with punctuation characters to get
this right, as in the following sentence.
       Gnats, gnus, etc.\ all  % `\ ' makes an inter-word space.
       begin with G\@.         % \@ marks end-of-sentence punctuation.
You should check the spaces after periods when
reading your output to make sure you haven't
forgotten any special cases.  Generating an
ellipsis
       \ldots\               % `\ ' is needed after `\ldots' because TeX
                             % ignores spaces after command names like \ldots
                             % made from \ + letters.
                             %
                             % Note how a `%' character causes TeX to ignore
                             % the end of the input line, so these blank lines
                             % do not start a new paragraph.
                             %
with the right spacing around the periods requires
a special command.

\LaTeX\ interprets some common characters as
commands, so you must type special commands to
generate them.  These characters include the
following:
       \$ \& \% \# \{ and \}.

In printing, text is usually emphasized with an
       \emph{italic}
type style.

\begin{em}
   A long segment of text can also be emphasized
   in this way.  Text within such a segment can be
   given \emph{additional} emphasis.
\end{em}

It is sometimes necessary to prevent \LaTeX\ from
breaking a line where it might otherwise do so.
This may be at a space, as between the ``Mr.''\ and
``Jones'' in
       ``Mr.~Jones'',        % ~ produces an unbreakable interword space.
or within a word---especially when the word is a
symbol like
       \mbox{\emph{itemnum}}
that makes little sense when hyphenated across
lines.

Footnotes\footnote{This is an example of a footnote.}
pose no problem.

\LaTeX\ is good at typesetting mathematical formulas
like
       \( x-3y + z = 7 \)
or
       \( a_{1} > x^{2n} + y^{2n} > x' \)
or
       \( \ip{A}{B} = \sum_{i} a_{i} b_{i} \).
The spaces you type in a formula are
ignored.  Remember that a letter like
       $x$                   % $ ... $  and  \( ... \)  are equivalent
is a formula when it denotes a mathematical
symbol, and it should be typed as one.

\section{Displayed Text}

Text is displayed by indenting it from the left
margin.  Quotations are commonly displayed.  There
are short quotations
\begin{quote}
   This is a short quotation.  It consists of a
   single paragraph of text.  See how it is formatted.
\end{quote}
and longer ones.
\begin{quotation}
   This is a longer quotation.  It consists of two
   paragraphs of text, neither of which are
   particularly interesting.

   This is the second paragraph of the quotation.  It
   is just as dull as the first paragraph.
\end{quotation}
Another frequently-displayed structure is a list.
The following is an example of an \emph{itemized}
list.
\begin{itemize}
   \item This is the first item of an itemized list.
         Each item in the list is marked with a ``tick''.
         You don't have to worry about what kind of tick
         mark is used.

   \item This is the second item of the list.  It
         contains another list nested inside it.  The inner
         list is an \emph{enumerated} list.
         \begin{enumerate}
            \item This is the first item of an enumerated
                  list that is nested within the itemized list.

            \item This is the second item of the inner list.
                  \LaTeX\ allows you to nest lists deeper than
                  you really should.
         \end{enumerate}
         This is the rest of the second item of the outer
         list.  It is no more interesting than any other
         part of the item.
   \item This is the third item of the list.
\end{itemize}
You can even display poetry.
\begin{verse}
   There is an environment
    for verse \\             % The \\ command separates lines
   Whose features some poets % within a stanza.
   will curse.

                             % One or more blank lines separate stanzas.

   For instead of making\\
   Them do \emph{all} line breaking, \\
   It allows them to put too many words on a line when they'd rather be
   forced to be terse.
\end{verse}

Mathematical formulas may also be displayed.  A
displayed formula
is
one-line long; multiline
formulas require special formatting instructions.
   \[  \ip{\Gamma}{\psi'} = x'' + y^{2} + z_{i}^{n}\]
Don't start a paragraph with a displayed equation,
nor make one a paragraph by itself.

\end{document}               % End of document.

\section{Bench Huge Section 705}
% synthetic bench section 705
% This is a sample LaTeX input file.  (Version of 12 August 2004.)
%
% A '%' character causes TeX to ignore all remaining text on the line,
% and is used for comments like this one.

\documentclass{article}      % Specifies the document class

                             % The preamble begins here.
\title{An Example Document}  % Declares the document's title.
\author{Leslie Lamport}      % Declares the author's name.
\date{January 21, 1994}      % Deleting this command produces today's date.

\newcommand{\ip}[2]{(#1, #2)}
                             % Defines \ip{arg1}{arg2} to mean
                             % (arg1, arg2).

%\newcommand{\ip}[2]{\langle #1 | #2\rangle}
                             % This is an alternative definition of
                             % \ip that is commented out.

\begin{document}             % End of preamble and beginning of text.

\maketitle                   % Produces the title.

This is an example input file.  Comparing it with
the output it generates can show you how to
produce a simple document of your own.

\section{Ordinary Text}      % Produces section heading.  Lower-level
                             % sections are begun with similar
                             % \subsection and \subsubsection commands.

The ends  of words and sentences are marked
  by   spaces. It  doesn't matter how many
spaces    you type; one is as good as 100.  The
end of   a line counts as a space.

One   or more   blank lines denote the  end
of  a paragraph.

Since any number of consecutive spaces are treated
like a single one, the formatting of the input
file makes no difference to
      \LaTeX,                % The \LaTeX command generates the LaTeX logo.
but it makes a difference to you.  When you use
\LaTeX, making your input file as easy to read
as possible will be a great help as you write
your document and when you change it.  This sample
file shows how you can add comments to your own input
file.

Because printing is different from typewriting,
there are a number of things that you have to do
differently when preparing an input file than if
you were just typing the document directly.
Quotation marks like
       ``this''
have to be handled specially, as do quotes within
quotes:
       ``\,`this'            % \, separates the double and single quote.
        is what I just
        wrote, not  `that'\,''.

Dashes come in three sizes: an
       intra-word
dash, a medium dash for number ranges like
       1--2,
and a punctuation
       dash---like
this.

A sentence-ending space should be larger than the
space between words within a sentence.  You
sometimes have to type special commands in
conjunction with punctuation characters to get
this right, as in the following sentence.
       Gnats, gnus, etc.\ all  % `\ ' makes an inter-word space.
       begin with G\@.         % \@ marks end-of-sentence punctuation.
You should check the spaces after periods when
reading your output to make sure you haven't
forgotten any special cases.  Generating an
ellipsis
       \ldots\               % `\ ' is needed after `\ldots' because TeX
                             % ignores spaces after command names like \ldots
                             % made from \ + letters.
                             %
                             % Note how a `%' character causes TeX to ignore
                             % the end of the input line, so these blank lines
                             % do not start a new paragraph.
                             %
with the right spacing around the periods requires
a special command.

\LaTeX\ interprets some common characters as
commands, so you must type special commands to
generate them.  These characters include the
following:
       \$ \& \% \# \{ and \}.

In printing, text is usually emphasized with an
       \emph{italic}
type style.

\begin{em}
   A long segment of text can also be emphasized
   in this way.  Text within such a segment can be
   given \emph{additional} emphasis.
\end{em}

It is sometimes necessary to prevent \LaTeX\ from
breaking a line where it might otherwise do so.
This may be at a space, as between the ``Mr.''\ and
``Jones'' in
       ``Mr.~Jones'',        % ~ produces an unbreakable interword space.
or within a word---especially when the word is a
symbol like
       \mbox{\emph{itemnum}}
that makes little sense when hyphenated across
lines.

Footnotes\footnote{This is an example of a footnote.}
pose no problem.

\LaTeX\ is good at typesetting mathematical formulas
like
       \( x-3y + z = 7 \)
or
       \( a_{1} > x^{2n} + y^{2n} > x' \)
or
       \( \ip{A}{B} = \sum_{i} a_{i} b_{i} \).
The spaces you type in a formula are
ignored.  Remember that a letter like
       $x$                   % $ ... $  and  \( ... \)  are equivalent
is a formula when it denotes a mathematical
symbol, and it should be typed as one.

\section{Displayed Text}

Text is displayed by indenting it from the left
margin.  Quotations are commonly displayed.  There
are short quotations
\begin{quote}
   This is a short quotation.  It consists of a
   single paragraph of text.  See how it is formatted.
\end{quote}
and longer ones.
\begin{quotation}
   This is a longer quotation.  It consists of two
   paragraphs of text, neither of which are
   particularly interesting.

   This is the second paragraph of the quotation.  It
   is just as dull as the first paragraph.
\end{quotation}
Another frequently-displayed structure is a list.
The following is an example of an \emph{itemized}
list.
\begin{itemize}
   \item This is the first item of an itemized list.
         Each item in the list is marked with a ``tick''.
         You don't have to worry about what kind of tick
         mark is used.

   \item This is the second item of the list.  It
         contains another list nested inside it.  The inner
         list is an \emph{enumerated} list.
         \begin{enumerate}
            \item This is the first item of an enumerated
                  list that is nested within the itemized list.

            \item This is the second item of the inner list.
                  \LaTeX\ allows you to nest lists deeper than
                  you really should.
         \end{enumerate}
         This is the rest of the second item of the outer
         list.  It is no more interesting than any other
         part of the item.
   \item This is the third item of the list.
\end{itemize}
You can even display poetry.
\begin{verse}
   There is an environment
    for verse \\             % The \\ command separates lines
   Whose features some poets % within a stanza.
   will curse.

                             % One or more blank lines separate stanzas.

   For instead of making\\
   Them do \emph{all} line breaking, \\
   It allows them to put too many words on a line when they'd rather be
   forced to be terse.
\end{verse}

Mathematical formulas may also be displayed.  A
displayed formula
is
one-line long; multiline
formulas require special formatting instructions.
   \[  \ip{\Gamma}{\psi'} = x'' + y^{2} + z_{i}^{n}\]
Don't start a paragraph with a displayed equation,
nor make one a paragraph by itself.

\end{document}               % End of document.

\section{Bench Huge Section 706}
% synthetic bench section 706
% This is a sample LaTeX input file.  (Version of 12 August 2004.)
%
% A '%' character causes TeX to ignore all remaining text on the line,
% and is used for comments like this one.

\documentclass{article}      % Specifies the document class

                             % The preamble begins here.
\title{An Example Document}  % Declares the document's title.
\author{Leslie Lamport}      % Declares the author's name.
\date{January 21, 1994}      % Deleting this command produces today's date.

\newcommand{\ip}[2]{(#1, #2)}
                             % Defines \ip{arg1}{arg2} to mean
                             % (arg1, arg2).

%\newcommand{\ip}[2]{\langle #1 | #2\rangle}
                             % This is an alternative definition of
                             % \ip that is commented out.

\begin{document}             % End of preamble and beginning of text.

\maketitle                   % Produces the title.

This is an example input file.  Comparing it with
the output it generates can show you how to
produce a simple document of your own.

\section{Ordinary Text}      % Produces section heading.  Lower-level
                             % sections are begun with similar
                             % \subsection and \subsubsection commands.

The ends  of words and sentences are marked
  by   spaces. It  doesn't matter how many
spaces    you type; one is as good as 100.  The
end of   a line counts as a space.

One   or more   blank lines denote the  end
of  a paragraph.

Since any number of consecutive spaces are treated
like a single one, the formatting of the input
file makes no difference to
      \LaTeX,                % The \LaTeX command generates the LaTeX logo.
but it makes a difference to you.  When you use
\LaTeX, making your input file as easy to read
as possible will be a great help as you write
your document and when you change it.  This sample
file shows how you can add comments to your own input
file.

Because printing is different from typewriting,
there are a number of things that you have to do
differently when preparing an input file than if
you were just typing the document directly.
Quotation marks like
       ``this''
have to be handled specially, as do quotes within
quotes:
       ``\,`this'            % \, separates the double and single quote.
        is what I just
        wrote, not  `that'\,''.

Dashes come in three sizes: an
       intra-word
dash, a medium dash for number ranges like
       1--2,
and a punctuation
       dash---like
this.

A sentence-ending space should be larger than the
space between words within a sentence.  You
sometimes have to type special commands in
conjunction with punctuation characters to get
this right, as in the following sentence.
       Gnats, gnus, etc.\ all  % `\ ' makes an inter-word space.
       begin with G\@.         % \@ marks end-of-sentence punctuation.
You should check the spaces after periods when
reading your output to make sure you haven't
forgotten any special cases.  Generating an
ellipsis
       \ldots\               % `\ ' is needed after `\ldots' because TeX
                             % ignores spaces after command names like \ldots
                             % made from \ + letters.
                             %
                             % Note how a `%' character causes TeX to ignore
                             % the end of the input line, so these blank lines
                             % do not start a new paragraph.
                             %
with the right spacing around the periods requires
a special command.

\LaTeX\ interprets some common characters as
commands, so you must type special commands to
generate them.  These characters include the
following:
       \$ \& \% \# \{ and \}.

In printing, text is usually emphasized with an
       \emph{italic}
type style.

\begin{em}
   A long segment of text can also be emphasized
   in this way.  Text within such a segment can be
   given \emph{additional} emphasis.
\end{em}

It is sometimes necessary to prevent \LaTeX\ from
breaking a line where it might otherwise do so.
This may be at a space, as between the ``Mr.''\ and
``Jones'' in
       ``Mr.~Jones'',        % ~ produces an unbreakable interword space.
or within a word---especially when the word is a
symbol like
       \mbox{\emph{itemnum}}
that makes little sense when hyphenated across
lines.

Footnotes\footnote{This is an example of a footnote.}
pose no problem.

\LaTeX\ is good at typesetting mathematical formulas
like
       \( x-3y + z = 7 \)
or
       \( a_{1} > x^{2n} + y^{2n} > x' \)
or
       \( \ip{A}{B} = \sum_{i} a_{i} b_{i} \).
The spaces you type in a formula are
ignored.  Remember that a letter like
       $x$                   % $ ... $  and  \( ... \)  are equivalent
is a formula when it denotes a mathematical
symbol, and it should be typed as one.

\section{Displayed Text}

Text is displayed by indenting it from the left
margin.  Quotations are commonly displayed.  There
are short quotations
\begin{quote}
   This is a short quotation.  It consists of a
   single paragraph of text.  See how it is formatted.
\end{quote}
and longer ones.
\begin{quotation}
   This is a longer quotation.  It consists of two
   paragraphs of text, neither of which are
   particularly interesting.

   This is the second paragraph of the quotation.  It
   is just as dull as the first paragraph.
\end{quotation}
Another frequently-displayed structure is a list.
The following is an example of an \emph{itemized}
list.
\begin{itemize}
   \item This is the first item of an itemized list.
         Each item in the list is marked with a ``tick''.
         You don't have to worry about what kind of tick
         mark is used.

   \item This is the second item of the list.  It
         contains another list nested inside it.  The inner
         list is an \emph{enumerated} list.
         \begin{enumerate}
            \item This is the first item of an enumerated
                  list that is nested within the itemized list.

            \item This is the second item of the inner list.
                  \LaTeX\ allows you to nest lists deeper than
                  you really should.
         \end{enumerate}
         This is the rest of the second item of the outer
         list.  It is no more interesting than any other
         part of the item.
   \item This is the third item of the list.
\end{itemize}
You can even display poetry.
\begin{verse}
   There is an environment
    for verse \\             % The \\ command separates lines
   Whose features some poets % within a stanza.
   will curse.

                             % One or more blank lines separate stanzas.

   For instead of making\\
   Them do \emph{all} line breaking, \\
   It allows them to put too many words on a line when they'd rather be
   forced to be terse.
\end{verse}

Mathematical formulas may also be displayed.  A
displayed formula
is
one-line long; multiline
formulas require special formatting instructions.
   \[  \ip{\Gamma}{\psi'} = x'' + y^{2} + z_{i}^{n}\]
Don't start a paragraph with a displayed equation,
nor make one a paragraph by itself.

\end{document}               % End of document.

\section{Bench Huge Section 707}
% synthetic bench section 707
% This is a sample LaTeX input file.  (Version of 12 August 2004.)
%
% A '%' character causes TeX to ignore all remaining text on the line,
% and is used for comments like this one.

\documentclass{article}      % Specifies the document class

                             % The preamble begins here.
\title{An Example Document}  % Declares the document's title.
\author{Leslie Lamport}      % Declares the author's name.
\date{January 21, 1994}      % Deleting this command produces today's date.

\newcommand{\ip}[2]{(#1, #2)}
                             % Defines \ip{arg1}{arg2} to mean
                             % (arg1, arg2).

%\newcommand{\ip}[2]{\langle #1 | #2\rangle}
                             % This is an alternative definition of
                             % \ip that is commented out.

\begin{document}             % End of preamble and beginning of text.

\maketitle                   % Produces the title.

This is an example input file.  Comparing it with
the output it generates can show you how to
produce a simple document of your own.

\section{Ordinary Text}      % Produces section heading.  Lower-level
                             % sections are begun with similar
                             % \subsection and \subsubsection commands.

The ends  of words and sentences are marked
  by   spaces. It  doesn't matter how many
spaces    you type; one is as good as 100.  The
end of   a line counts as a space.

One   or more   blank lines denote the  end
of  a paragraph.

Since any number of consecutive spaces are treated
like a single one, the formatting of the input
file makes no difference to
      \LaTeX,                % The \LaTeX command generates the LaTeX logo.
but it makes a difference to you.  When you use
\LaTeX, making your input file as easy to read
as possible will be a great help as you write
your document and when you change it.  This sample
file shows how you can add comments to your own input
file.

Because printing is different from typewriting,
there are a number of things that you have to do
differently when preparing an input file than if
you were just typing the document directly.
Quotation marks like
       ``this''
have to be handled specially, as do quotes within
quotes:
       ``\,`this'            % \, separates the double and single quote.
        is what I just
        wrote, not  `that'\,''.

Dashes come in three sizes: an
       intra-word
dash, a medium dash for number ranges like
       1--2,
and a punctuation
       dash---like
this.

A sentence-ending space should be larger than the
space between words within a sentence.  You
sometimes have to type special commands in
conjunction with punctuation characters to get
this right, as in the following sentence.
       Gnats, gnus, etc.\ all  % `\ ' makes an inter-word space.
       begin with G\@.         % \@ marks end-of-sentence punctuation.
You should check the spaces after periods when
reading your output to make sure you haven't
forgotten any special cases.  Generating an
ellipsis
       \ldots\               % `\ ' is needed after `\ldots' because TeX
                             % ignores spaces after command names like \ldots
                             % made from \ + letters.
                             %
                             % Note how a `%' character causes TeX to ignore
                             % the end of the input line, so these blank lines
                             % do not start a new paragraph.
                             %
with the right spacing around the periods requires
a special command.

\LaTeX\ interprets some common characters as
commands, so you must type special commands to
generate them.  These characters include the
following:
       \$ \& \% \# \{ and \}.

In printing, text is usually emphasized with an
       \emph{italic}
type style.

\begin{em}
   A long segment of text can also be emphasized
   in this way.  Text within such a segment can be
   given \emph{additional} emphasis.
\end{em}

It is sometimes necessary to prevent \LaTeX\ from
breaking a line where it might otherwise do so.
This may be at a space, as between the ``Mr.''\ and
``Jones'' in
       ``Mr.~Jones'',        % ~ produces an unbreakable interword space.
or within a word---especially when the word is a
symbol like
       \mbox{\emph{itemnum}}
that makes little sense when hyphenated across
lines.

Footnotes\footnote{This is an example of a footnote.}
pose no problem.

\LaTeX\ is good at typesetting mathematical formulas
like
       \( x-3y + z = 7 \)
or
       \( a_{1} > x^{2n} + y^{2n} > x' \)
or
       \( \ip{A}{B} = \sum_{i} a_{i} b_{i} \).
The spaces you type in a formula are
ignored.  Remember that a letter like
       $x$                   % $ ... $  and  \( ... \)  are equivalent
is a formula when it denotes a mathematical
symbol, and it should be typed as one.

\section{Displayed Text}

Text is displayed by indenting it from the left
margin.  Quotations are commonly displayed.  There
are short quotations
\begin{quote}
   This is a short quotation.  It consists of a
   single paragraph of text.  See how it is formatted.
\end{quote}
and longer ones.
\begin{quotation}
   This is a longer quotation.  It consists of two
   paragraphs of text, neither of which are
   particularly interesting.

   This is the second paragraph of the quotation.  It
   is just as dull as the first paragraph.
\end{quotation}
Another frequently-displayed structure is a list.
The following is an example of an \emph{itemized}
list.
\begin{itemize}
   \item This is the first item of an itemized list.
         Each item in the list is marked with a ``tick''.
         You don't have to worry about what kind of tick
         mark is used.

   \item This is the second item of the list.  It
         contains another list nested inside it.  The inner
         list is an \emph{enumerated} list.
         \begin{enumerate}
            \item This is the first item of an enumerated
                  list that is nested within the itemized list.

            \item This is the second item of the inner list.
                  \LaTeX\ allows you to nest lists deeper than
                  you really should.
         \end{enumerate}
         This is the rest of the second item of the outer
         list.  It is no more interesting than any other
         part of the item.
   \item This is the third item of the list.
\end{itemize}
You can even display poetry.
\begin{verse}
   There is an environment
    for verse \\             % The \\ command separates lines
   Whose features some poets % within a stanza.
   will curse.

                             % One or more blank lines separate stanzas.

   For instead of making\\
   Them do \emph{all} line breaking, \\
   It allows them to put too many words on a line when they'd rather be
   forced to be terse.
\end{verse}

Mathematical formulas may also be displayed.  A
displayed formula
is
one-line long; multiline
formulas require special formatting instructions.
   \[  \ip{\Gamma}{\psi'} = x'' + y^{2} + z_{i}^{n}\]
Don't start a paragraph with a displayed equation,
nor make one a paragraph by itself.

\end{document}               % End of document.

\section{Bench Huge Section 708}
% synthetic bench section 708
% This is a sample LaTeX input file.  (Version of 12 August 2004.)
%
% A '%' character causes TeX to ignore all remaining text on the line,
% and is used for comments like this one.

\documentclass{article}      % Specifies the document class

                             % The preamble begins here.
\title{An Example Document}  % Declares the document's title.
\author{Leslie Lamport}      % Declares the author's name.
\date{January 21, 1994}      % Deleting this command produces today's date.

\newcommand{\ip}[2]{(#1, #2)}
                             % Defines \ip{arg1}{arg2} to mean
                             % (arg1, arg2).

%\newcommand{\ip}[2]{\langle #1 | #2\rangle}
                             % This is an alternative definition of
                             % \ip that is commented out.

\begin{document}             % End of preamble and beginning of text.

\maketitle                   % Produces the title.

This is an example input file.  Comparing it with
the output it generates can show you how to
produce a simple document of your own.

\section{Ordinary Text}      % Produces section heading.  Lower-level
                             % sections are begun with similar
                             % \subsection and \subsubsection commands.

The ends  of words and sentences are marked
  by   spaces. It  doesn't matter how many
spaces    you type; one is as good as 100.  The
end of   a line counts as a space.

One   or more   blank lines denote the  end
of  a paragraph.

Since any number of consecutive spaces are treated
like a single one, the formatting of the input
file makes no difference to
      \LaTeX,                % The \LaTeX command generates the LaTeX logo.
but it makes a difference to you.  When you use
\LaTeX, making your input file as easy to read
as possible will be a great help as you write
your document and when you change it.  This sample
file shows how you can add comments to your own input
file.

Because printing is different from typewriting,
there are a number of things that you have to do
differently when preparing an input file than if
you were just typing the document directly.
Quotation marks like
       ``this''
have to be handled specially, as do quotes within
quotes:
       ``\,`this'            % \, separates the double and single quote.
        is what I just
        wrote, not  `that'\,''.

Dashes come in three sizes: an
       intra-word
dash, a medium dash for number ranges like
       1--2,
and a punctuation
       dash---like
this.

A sentence-ending space should be larger than the
space between words within a sentence.  You
sometimes have to type special commands in
conjunction with punctuation characters to get
this right, as in the following sentence.
       Gnats, gnus, etc.\ all  % `\ ' makes an inter-word space.
       begin with G\@.         % \@ marks end-of-sentence punctuation.
You should check the spaces after periods when
reading your output to make sure you haven't
forgotten any special cases.  Generating an
ellipsis
       \ldots\               % `\ ' is needed after `\ldots' because TeX
                             % ignores spaces after command names like \ldots
                             % made from \ + letters.
                             %
                             % Note how a `%' character causes TeX to ignore
                             % the end of the input line, so these blank lines
                             % do not start a new paragraph.
                             %
with the right spacing around the periods requires
a special command.

\LaTeX\ interprets some common characters as
commands, so you must type special commands to
generate them.  These characters include the
following:
       \$ \& \% \# \{ and \}.

In printing, text is usually emphasized with an
       \emph{italic}
type style.

\begin{em}
   A long segment of text can also be emphasized
   in this way.  Text within such a segment can be
   given \emph{additional} emphasis.
\end{em}

It is sometimes necessary to prevent \LaTeX\ from
breaking a line where it might otherwise do so.
This may be at a space, as between the ``Mr.''\ and
``Jones'' in
       ``Mr.~Jones'',        % ~ produces an unbreakable interword space.
or within a word---especially when the word is a
symbol like
       \mbox{\emph{itemnum}}
that makes little sense when hyphenated across
lines.

Footnotes\footnote{This is an example of a footnote.}
pose no problem.

\LaTeX\ is good at typesetting mathematical formulas
like
       \( x-3y + z = 7 \)
or
       \( a_{1} > x^{2n} + y^{2n} > x' \)
or
       \( \ip{A}{B} = \sum_{i} a_{i} b_{i} \).
The spaces you type in a formula are
ignored.  Remember that a letter like
       $x$                   % $ ... $  and  \( ... \)  are equivalent
is a formula when it denotes a mathematical
symbol, and it should be typed as one.

\section{Displayed Text}

Text is displayed by indenting it from the left
margin.  Quotations are commonly displayed.  There
are short quotations
\begin{quote}
   This is a short quotation.  It consists of a
   single paragraph of text.  See how it is formatted.
\end{quote}
and longer ones.
\begin{quotation}
   This is a longer quotation.  It consists of two
   paragraphs of text, neither of which are
   particularly interesting.

   This is the second paragraph of the quotation.  It
   is just as dull as the first paragraph.
\end{quotation}
Another frequently-displayed structure is a list.
The following is an example of an \emph{itemized}
list.
\begin{itemize}
   \item This is the first item of an itemized list.
         Each item in the list is marked with a ``tick''.
         You don't have to worry about what kind of tick
         mark is used.

   \item This is the second item of the list.  It
         contains another list nested inside it.  The inner
         list is an \emph{enumerated} list.
         \begin{enumerate}
            \item This is the first item of an enumerated
                  list that is nested within the itemized list.

            \item This is the second item of the inner list.
                  \LaTeX\ allows you to nest lists deeper than
                  you really should.
         \end{enumerate}
         This is the rest of the second item of the outer
         list.  It is no more interesting than any other
         part of the item.
   \item This is the third item of the list.
\end{itemize}
You can even display poetry.
\begin{verse}
   There is an environment
    for verse \\             % The \\ command separates lines
   Whose features some poets % within a stanza.
   will curse.

                             % One or more blank lines separate stanzas.

   For instead of making\\
   Them do \emph{all} line breaking, \\
   It allows them to put too many words on a line when they'd rather be
   forced to be terse.
\end{verse}

Mathematical formulas may also be displayed.  A
displayed formula
is
one-line long; multiline
formulas require special formatting instructions.
   \[  \ip{\Gamma}{\psi'} = x'' + y^{2} + z_{i}^{n}\]
Don't start a paragraph with a displayed equation,
nor make one a paragraph by itself.

\end{document}               % End of document.

\section{Bench Huge Section 709}
% synthetic bench section 709
% This is a sample LaTeX input file.  (Version of 12 August 2004.)
%
% A '%' character causes TeX to ignore all remaining text on the line,
% and is used for comments like this one.

\documentclass{article}      % Specifies the document class

                             % The preamble begins here.
\title{An Example Document}  % Declares the document's title.
\author{Leslie Lamport}      % Declares the author's name.
\date{January 21, 1994}      % Deleting this command produces today's date.

\newcommand{\ip}[2]{(#1, #2)}
                             % Defines \ip{arg1}{arg2} to mean
                             % (arg1, arg2).

%\newcommand{\ip}[2]{\langle #1 | #2\rangle}
                             % This is an alternative definition of
                             % \ip that is commented out.

\begin{document}             % End of preamble and beginning of text.

\maketitle                   % Produces the title.

This is an example input file.  Comparing it with
the output it generates can show you how to
produce a simple document of your own.

\section{Ordinary Text}      % Produces section heading.  Lower-level
                             % sections are begun with similar
                             % \subsection and \subsubsection commands.

The ends  of words and sentences are marked
  by   spaces. It  doesn't matter how many
spaces    you type; one is as good as 100.  The
end of   a line counts as a space.

One   or more   blank lines denote the  end
of  a paragraph.

Since any number of consecutive spaces are treated
like a single one, the formatting of the input
file makes no difference to
      \LaTeX,                % The \LaTeX command generates the LaTeX logo.
but it makes a difference to you.  When you use
\LaTeX, making your input file as easy to read
as possible will be a great help as you write
your document and when you change it.  This sample
file shows how you can add comments to your own input
file.

Because printing is different from typewriting,
there are a number of things that you have to do
differently when preparing an input file than if
you were just typing the document directly.
Quotation marks like
       ``this''
have to be handled specially, as do quotes within
quotes:
       ``\,`this'            % \, separates the double and single quote.
        is what I just
        wrote, not  `that'\,''.

Dashes come in three sizes: an
       intra-word
dash, a medium dash for number ranges like
       1--2,
and a punctuation
       dash---like
this.

A sentence-ending space should be larger than the
space between words within a sentence.  You
sometimes have to type special commands in
conjunction with punctuation characters to get
this right, as in the following sentence.
       Gnats, gnus, etc.\ all  % `\ ' makes an inter-word space.
       begin with G\@.         % \@ marks end-of-sentence punctuation.
You should check the spaces after periods when
reading your output to make sure you haven't
forgotten any special cases.  Generating an
ellipsis
       \ldots\               % `\ ' is needed after `\ldots' because TeX
                             % ignores spaces after command names like \ldots
                             % made from \ + letters.
                             %
                             % Note how a `%' character causes TeX to ignore
                             % the end of the input line, so these blank lines
                             % do not start a new paragraph.
                             %
with the right spacing around the periods requires
a special command.

\LaTeX\ interprets some common characters as
commands, so you must type special commands to
generate them.  These characters include the
following:
       \$ \& \% \# \{ and \}.

In printing, text is usually emphasized with an
       \emph{italic}
type style.

\begin{em}
   A long segment of text can also be emphasized
   in this way.  Text within such a segment can be
   given \emph{additional} emphasis.
\end{em}

It is sometimes necessary to prevent \LaTeX\ from
breaking a line where it might otherwise do so.
This may be at a space, as between the ``Mr.''\ and
``Jones'' in
       ``Mr.~Jones'',        % ~ produces an unbreakable interword space.
or within a word---especially when the word is a
symbol like
       \mbox{\emph{itemnum}}
that makes little sense when hyphenated across
lines.

Footnotes\footnote{This is an example of a footnote.}
pose no problem.

\LaTeX\ is good at typesetting mathematical formulas
like
       \( x-3y + z = 7 \)
or
       \( a_{1} > x^{2n} + y^{2n} > x' \)
or
       \( \ip{A}{B} = \sum_{i} a_{i} b_{i} \).
The spaces you type in a formula are
ignored.  Remember that a letter like
       $x$                   % $ ... $  and  \( ... \)  are equivalent
is a formula when it denotes a mathematical
symbol, and it should be typed as one.

\section{Displayed Text}

Text is displayed by indenting it from the left
margin.  Quotations are commonly displayed.  There
are short quotations
\begin{quote}
   This is a short quotation.  It consists of a
   single paragraph of text.  See how it is formatted.
\end{quote}
and longer ones.
\begin{quotation}
   This is a longer quotation.  It consists of two
   paragraphs of text, neither of which are
   particularly interesting.

   This is the second paragraph of the quotation.  It
   is just as dull as the first paragraph.
\end{quotation}
Another frequently-displayed structure is a list.
The following is an example of an \emph{itemized}
list.
\begin{itemize}
   \item This is the first item of an itemized list.
         Each item in the list is marked with a ``tick''.
         You don't have to worry about what kind of tick
         mark is used.

   \item This is the second item of the list.  It
         contains another list nested inside it.  The inner
         list is an \emph{enumerated} list.
         \begin{enumerate}
            \item This is the first item of an enumerated
                  list that is nested within the itemized list.

            \item This is the second item of the inner list.
                  \LaTeX\ allows you to nest lists deeper than
                  you really should.
         \end{enumerate}
         This is the rest of the second item of the outer
         list.  It is no more interesting than any other
         part of the item.
   \item This is the third item of the list.
\end{itemize}
You can even display poetry.
\begin{verse}
   There is an environment
    for verse \\             % The \\ command separates lines
   Whose features some poets % within a stanza.
   will curse.

                             % One or more blank lines separate stanzas.

   For instead of making\\
   Them do \emph{all} line breaking, \\
   It allows them to put too many words on a line when they'd rather be
   forced to be terse.
\end{verse}

Mathematical formulas may also be displayed.  A
displayed formula
is
one-line long; multiline
formulas require special formatting instructions.
   \[  \ip{\Gamma}{\psi'} = x'' + y^{2} + z_{i}^{n}\]
Don't start a paragraph with a displayed equation,
nor make one a paragraph by itself.

\end{document}               % End of document.

\section{Bench Huge Section 710}
% synthetic bench section 710
% This is a sample LaTeX input file.  (Version of 12 August 2004.)
%
% A '%' character causes TeX to ignore all remaining text on the line,
% and is used for comments like this one.

\documentclass{article}      % Specifies the document class

                             % The preamble begins here.
\title{An Example Document}  % Declares the document's title.
\author{Leslie Lamport}      % Declares the author's name.
\date{January 21, 1994}      % Deleting this command produces today's date.

\newcommand{\ip}[2]{(#1, #2)}
                             % Defines \ip{arg1}{arg2} to mean
                             % (arg1, arg2).

%\newcommand{\ip}[2]{\langle #1 | #2\rangle}
                             % This is an alternative definition of
                             % \ip that is commented out.

\begin{document}             % End of preamble and beginning of text.

\maketitle                   % Produces the title.

This is an example input file.  Comparing it with
the output it generates can show you how to
produce a simple document of your own.

\section{Ordinary Text}      % Produces section heading.  Lower-level
                             % sections are begun with similar
                             % \subsection and \subsubsection commands.

The ends  of words and sentences are marked
  by   spaces. It  doesn't matter how many
spaces    you type; one is as good as 100.  The
end of   a line counts as a space.

One   or more   blank lines denote the  end
of  a paragraph.

Since any number of consecutive spaces are treated
like a single one, the formatting of the input
file makes no difference to
      \LaTeX,                % The \LaTeX command generates the LaTeX logo.
but it makes a difference to you.  When you use
\LaTeX, making your input file as easy to read
as possible will be a great help as you write
your document and when you change it.  This sample
file shows how you can add comments to your own input
file.

Because printing is different from typewriting,
there are a number of things that you have to do
differently when preparing an input file than if
you were just typing the document directly.
Quotation marks like
       ``this''
have to be handled specially, as do quotes within
quotes:
       ``\,`this'            % \, separates the double and single quote.
        is what I just
        wrote, not  `that'\,''.

Dashes come in three sizes: an
       intra-word
dash, a medium dash for number ranges like
       1--2,
and a punctuation
       dash---like
this.

A sentence-ending space should be larger than the
space between words within a sentence.  You
sometimes have to type special commands in
conjunction with punctuation characters to get
this right, as in the following sentence.
       Gnats, gnus, etc.\ all  % `\ ' makes an inter-word space.
       begin with G\@.         % \@ marks end-of-sentence punctuation.
You should check the spaces after periods when
reading your output to make sure you haven't
forgotten any special cases.  Generating an
ellipsis
       \ldots\               % `\ ' is needed after `\ldots' because TeX
                             % ignores spaces after command names like \ldots
                             % made from \ + letters.
                             %
                             % Note how a `%' character causes TeX to ignore
                             % the end of the input line, so these blank lines
                             % do not start a new paragraph.
                             %
with the right spacing around the periods requires
a special command.

\LaTeX\ interprets some common characters as
commands, so you must type special commands to
generate them.  These characters include the
following:
       \$ \& \% \# \{ and \}.

In printing, text is usually emphasized with an
       \emph{italic}
type style.

\begin{em}
   A long segment of text can also be emphasized
   in this way.  Text within such a segment can be
   given \emph{additional} emphasis.
\end{em}

It is sometimes necessary to prevent \LaTeX\ from
breaking a line where it might otherwise do so.
This may be at a space, as between the ``Mr.''\ and
``Jones'' in
       ``Mr.~Jones'',        % ~ produces an unbreakable interword space.
or within a word---especially when the word is a
symbol like
       \mbox{\emph{itemnum}}
that makes little sense when hyphenated across
lines.

Footnotes\footnote{This is an example of a footnote.}
pose no problem.

\LaTeX\ is good at typesetting mathematical formulas
like
       \( x-3y + z = 7 \)
or
       \( a_{1} > x^{2n} + y^{2n} > x' \)
or
       \( \ip{A}{B} = \sum_{i} a_{i} b_{i} \).
The spaces you type in a formula are
ignored.  Remember that a letter like
       $x$                   % $ ... $  and  \( ... \)  are equivalent
is a formula when it denotes a mathematical
symbol, and it should be typed as one.

\section{Displayed Text}

Text is displayed by indenting it from the left
margin.  Quotations are commonly displayed.  There
are short quotations
\begin{quote}
   This is a short quotation.  It consists of a
   single paragraph of text.  See how it is formatted.
\end{quote}
and longer ones.
\begin{quotation}
   This is a longer quotation.  It consists of two
   paragraphs of text, neither of which are
   particularly interesting.

   This is the second paragraph of the quotation.  It
   is just as dull as the first paragraph.
\end{quotation}
Another frequently-displayed structure is a list.
The following is an example of an \emph{itemized}
list.
\begin{itemize}
   \item This is the first item of an itemized list.
         Each item in the list is marked with a ``tick''.
         You don't have to worry about what kind of tick
         mark is used.

   \item This is the second item of the list.  It
         contains another list nested inside it.  The inner
         list is an \emph{enumerated} list.
         \begin{enumerate}
            \item This is the first item of an enumerated
                  list that is nested within the itemized list.

            \item This is the second item of the inner list.
                  \LaTeX\ allows you to nest lists deeper than
                  you really should.
         \end{enumerate}
         This is the rest of the second item of the outer
         list.  It is no more interesting than any other
         part of the item.
   \item This is the third item of the list.
\end{itemize}
You can even display poetry.
\begin{verse}
   There is an environment
    for verse \\             % The \\ command separates lines
   Whose features some poets % within a stanza.
   will curse.

                             % One or more blank lines separate stanzas.

   For instead of making\\
   Them do \emph{all} line breaking, \\
   It allows them to put too many words on a line when they'd rather be
   forced to be terse.
\end{verse}

Mathematical formulas may also be displayed.  A
displayed formula
is
one-line long; multiline
formulas require special formatting instructions.
   \[  \ip{\Gamma}{\psi'} = x'' + y^{2} + z_{i}^{n}\]
Don't start a paragraph with a displayed equation,
nor make one a paragraph by itself.

\end{document}               % End of document.

\section{Bench Huge Section 711}
% synthetic bench section 711
% This is a sample LaTeX input file.  (Version of 12 August 2004.)
%
% A '%' character causes TeX to ignore all remaining text on the line,
% and is used for comments like this one.

\documentclass{article}      % Specifies the document class

                             % The preamble begins here.
\title{An Example Document}  % Declares the document's title.
\author{Leslie Lamport}      % Declares the author's name.
\date{January 21, 1994}      % Deleting this command produces today's date.

\newcommand{\ip}[2]{(#1, #2)}
                             % Defines \ip{arg1}{arg2} to mean
                             % (arg1, arg2).

%\newcommand{\ip}[2]{\langle #1 | #2\rangle}
                             % This is an alternative definition of
                             % \ip that is commented out.

\begin{document}             % End of preamble and beginning of text.

\maketitle                   % Produces the title.

This is an example input file.  Comparing it with
the output it generates can show you how to
produce a simple document of your own.

\section{Ordinary Text}      % Produces section heading.  Lower-level
                             % sections are begun with similar
                             % \subsection and \subsubsection commands.

The ends  of words and sentences are marked
  by   spaces. It  doesn't matter how many
spaces    you type; one is as good as 100.  The
end of   a line counts as a space.

One   or more   blank lines denote the  end
of  a paragraph.

Since any number of consecutive spaces are treated
like a single one, the formatting of the input
file makes no difference to
      \LaTeX,                % The \LaTeX command generates the LaTeX logo.
but it makes a difference to you.  When you use
\LaTeX, making your input file as easy to read
as possible will be a great help as you write
your document and when you change it.  This sample
file shows how you can add comments to your own input
file.

Because printing is different from typewriting,
there are a number of things that you have to do
differently when preparing an input file than if
you were just typing the document directly.
Quotation marks like
       ``this''
have to be handled specially, as do quotes within
quotes:
       ``\,`this'            % \, separates the double and single quote.
        is what I just
        wrote, not  `that'\,''.

Dashes come in three sizes: an
       intra-word
dash, a medium dash for number ranges like
       1--2,
and a punctuation
       dash---like
this.

A sentence-ending space should be larger than the
space between words within a sentence.  You
sometimes have to type special commands in
conjunction with punctuation characters to get
this right, as in the following sentence.
       Gnats, gnus, etc.\ all  % `\ ' makes an inter-word space.
       begin with G\@.         % \@ marks end-of-sentence punctuation.
You should check the spaces after periods when
reading your output to make sure you haven't
forgotten any special cases.  Generating an
ellipsis
       \ldots\               % `\ ' is needed after `\ldots' because TeX
                             % ignores spaces after command names like \ldots
                             % made from \ + letters.
                             %
                             % Note how a `%' character causes TeX to ignore
                             % the end of the input line, so these blank lines
                             % do not start a new paragraph.
                             %
with the right spacing around the periods requires
a special command.

\LaTeX\ interprets some common characters as
commands, so you must type special commands to
generate them.  These characters include the
following:
       \$ \& \% \# \{ and \}.

In printing, text is usually emphasized with an
       \emph{italic}
type style.

\begin{em}
   A long segment of text can also be emphasized
   in this way.  Text within such a segment can be
   given \emph{additional} emphasis.
\end{em}

It is sometimes necessary to prevent \LaTeX\ from
breaking a line where it might otherwise do so.
This may be at a space, as between the ``Mr.''\ and
``Jones'' in
       ``Mr.~Jones'',        % ~ produces an unbreakable interword space.
or within a word---especially when the word is a
symbol like
       \mbox{\emph{itemnum}}
that makes little sense when hyphenated across
lines.

Footnotes\footnote{This is an example of a footnote.}
pose no problem.

\LaTeX\ is good at typesetting mathematical formulas
like
       \( x-3y + z = 7 \)
or
       \( a_{1} > x^{2n} + y^{2n} > x' \)
or
       \( \ip{A}{B} = \sum_{i} a_{i} b_{i} \).
The spaces you type in a formula are
ignored.  Remember that a letter like
       $x$                   % $ ... $  and  \( ... \)  are equivalent
is a formula when it denotes a mathematical
symbol, and it should be typed as one.

\section{Displayed Text}

Text is displayed by indenting it from the left
margin.  Quotations are commonly displayed.  There
are short quotations
\begin{quote}
   This is a short quotation.  It consists of a
   single paragraph of text.  See how it is formatted.
\end{quote}
and longer ones.
\begin{quotation}
   This is a longer quotation.  It consists of two
   paragraphs of text, neither of which are
   particularly interesting.

   This is the second paragraph of the quotation.  It
   is just as dull as the first paragraph.
\end{quotation}
Another frequently-displayed structure is a list.
The following is an example of an \emph{itemized}
list.
\begin{itemize}
   \item This is the first item of an itemized list.
         Each item in the list is marked with a ``tick''.
         You don't have to worry about what kind of tick
         mark is used.

   \item This is the second item of the list.  It
         contains another list nested inside it.  The inner
         list is an \emph{enumerated} list.
         \begin{enumerate}
            \item This is the first item of an enumerated
                  list that is nested within the itemized list.

            \item This is the second item of the inner list.
                  \LaTeX\ allows you to nest lists deeper than
                  you really should.
         \end{enumerate}
         This is the rest of the second item of the outer
         list.  It is no more interesting than any other
         part of the item.
   \item This is the third item of the list.
\end{itemize}
You can even display poetry.
\begin{verse}
   There is an environment
    for verse \\             % The \\ command separates lines
   Whose features some poets % within a stanza.
   will curse.

                             % One or more blank lines separate stanzas.

   For instead of making\\
   Them do \emph{all} line breaking, \\
   It allows them to put too many words on a line when they'd rather be
   forced to be terse.
\end{verse}

Mathematical formulas may also be displayed.  A
displayed formula
is
one-line long; multiline
formulas require special formatting instructions.
   \[  \ip{\Gamma}{\psi'} = x'' + y^{2} + z_{i}^{n}\]
Don't start a paragraph with a displayed equation,
nor make one a paragraph by itself.

\end{document}               % End of document.

\section{Bench Huge Section 712}
% synthetic bench section 712
% This is a sample LaTeX input file.  (Version of 12 August 2004.)
%
% A '%' character causes TeX to ignore all remaining text on the line,
% and is used for comments like this one.

\documentclass{article}      % Specifies the document class

                             % The preamble begins here.
\title{An Example Document}  % Declares the document's title.
\author{Leslie Lamport}      % Declares the author's name.
\date{January 21, 1994}      % Deleting this command produces today's date.

\newcommand{\ip}[2]{(#1, #2)}
                             % Defines \ip{arg1}{arg2} to mean
                             % (arg1, arg2).

%\newcommand{\ip}[2]{\langle #1 | #2\rangle}
                             % This is an alternative definition of
                             % \ip that is commented out.

\begin{document}             % End of preamble and beginning of text.

\maketitle                   % Produces the title.

This is an example input file.  Comparing it with
the output it generates can show you how to
produce a simple document of your own.

\section{Ordinary Text}      % Produces section heading.  Lower-level
                             % sections are begun with similar
                             % \subsection and \subsubsection commands.

The ends  of words and sentences are marked
  by   spaces. It  doesn't matter how many
spaces    you type; one is as good as 100.  The
end of   a line counts as a space.

One   or more   blank lines denote the  end
of  a paragraph.

Since any number of consecutive spaces are treated
like a single one, the formatting of the input
file makes no difference to
      \LaTeX,                % The \LaTeX command generates the LaTeX logo.
but it makes a difference to you.  When you use
\LaTeX, making your input file as easy to read
as possible will be a great help as you write
your document and when you change it.  This sample
file shows how you can add comments to your own input
file.

Because printing is different from typewriting,
there are a number of things that you have to do
differently when preparing an input file than if
you were just typing the document directly.
Quotation marks like
       ``this''
have to be handled specially, as do quotes within
quotes:
       ``\,`this'            % \, separates the double and single quote.
        is what I just
        wrote, not  `that'\,''.

Dashes come in three sizes: an
       intra-word
dash, a medium dash for number ranges like
       1--2,
and a punctuation
       dash---like
this.

A sentence-ending space should be larger than the
space between words within a sentence.  You
sometimes have to type special commands in
conjunction with punctuation characters to get
this right, as in the following sentence.
       Gnats, gnus, etc.\ all  % `\ ' makes an inter-word space.
       begin with G\@.         % \@ marks end-of-sentence punctuation.
You should check the spaces after periods when
reading your output to make sure you haven't
forgotten any special cases.  Generating an
ellipsis
       \ldots\               % `\ ' is needed after `\ldots' because TeX
                             % ignores spaces after command names like \ldots
                             % made from \ + letters.
                             %
                             % Note how a `%' character causes TeX to ignore
                             % the end of the input line, so these blank lines
                             % do not start a new paragraph.
                             %
with the right spacing around the periods requires
a special command.

\LaTeX\ interprets some common characters as
commands, so you must type special commands to
generate them.  These characters include the
following:
       \$ \& \% \# \{ and \}.

In printing, text is usually emphasized with an
       \emph{italic}
type style.

\begin{em}
   A long segment of text can also be emphasized
   in this way.  Text within such a segment can be
   given \emph{additional} emphasis.
\end{em}

It is sometimes necessary to prevent \LaTeX\ from
breaking a line where it might otherwise do so.
This may be at a space, as between the ``Mr.''\ and
``Jones'' in
       ``Mr.~Jones'',        % ~ produces an unbreakable interword space.
or within a word---especially when the word is a
symbol like
       \mbox{\emph{itemnum}}
that makes little sense when hyphenated across
lines.

Footnotes\footnote{This is an example of a footnote.}
pose no problem.

\LaTeX\ is good at typesetting mathematical formulas
like
       \( x-3y + z = 7 \)
or
       \( a_{1} > x^{2n} + y^{2n} > x' \)
or
       \( \ip{A}{B} = \sum_{i} a_{i} b_{i} \).
The spaces you type in a formula are
ignored.  Remember that a letter like
       $x$                   % $ ... $  and  \( ... \)  are equivalent
is a formula when it denotes a mathematical
symbol, and it should be typed as one.

\section{Displayed Text}

Text is displayed by indenting it from the left
margin.  Quotations are commonly displayed.  There
are short quotations
\begin{quote}
   This is a short quotation.  It consists of a
   single paragraph of text.  See how it is formatted.
\end{quote}
and longer ones.
\begin{quotation}
   This is a longer quotation.  It consists of two
   paragraphs of text, neither of which are
   particularly interesting.

   This is the second paragraph of the quotation.  It
   is just as dull as the first paragraph.
\end{quotation}
Another frequently-displayed structure is a list.
The following is an example of an \emph{itemized}
list.
\begin{itemize}
   \item This is the first item of an itemized list.
         Each item in the list is marked with a ``tick''.
         You don't have to worry about what kind of tick
         mark is used.

   \item This is the second item of the list.  It
         contains another list nested inside it.  The inner
         list is an \emph{enumerated} list.
         \begin{enumerate}
            \item This is the first item of an enumerated
                  list that is nested within the itemized list.

            \item This is the second item of the inner list.
                  \LaTeX\ allows you to nest lists deeper than
                  you really should.
         \end{enumerate}
         This is the rest of the second item of the outer
         list.  It is no more interesting than any other
         part of the item.
   \item This is the third item of the list.
\end{itemize}
You can even display poetry.
\begin{verse}
   There is an environment
    for verse \\             % The \\ command separates lines
   Whose features some poets % within a stanza.
   will curse.

                             % One or more blank lines separate stanzas.

   For instead of making\\
   Them do \emph{all} line breaking, \\
   It allows them to put too many words on a line when they'd rather be
   forced to be terse.
\end{verse}

Mathematical formulas may also be displayed.  A
displayed formula
is
one-line long; multiline
formulas require special formatting instructions.
   \[  \ip{\Gamma}{\psi'} = x'' + y^{2} + z_{i}^{n}\]
Don't start a paragraph with a displayed equation,
nor make one a paragraph by itself.

\end{document}               % End of document.

\section{Bench Huge Section 713}
% synthetic bench section 713
% This is a sample LaTeX input file.  (Version of 12 August 2004.)
%
% A '%' character causes TeX to ignore all remaining text on the line,
% and is used for comments like this one.

\documentclass{article}      % Specifies the document class

                             % The preamble begins here.
\title{An Example Document}  % Declares the document's title.
\author{Leslie Lamport}      % Declares the author's name.
\date{January 21, 1994}      % Deleting this command produces today's date.

\newcommand{\ip}[2]{(#1, #2)}
                             % Defines \ip{arg1}{arg2} to mean
                             % (arg1, arg2).

%\newcommand{\ip}[2]{\langle #1 | #2\rangle}
                             % This is an alternative definition of
                             % \ip that is commented out.

\begin{document}             % End of preamble and beginning of text.

\maketitle                   % Produces the title.

This is an example input file.  Comparing it with
the output it generates can show you how to
produce a simple document of your own.

\section{Ordinary Text}      % Produces section heading.  Lower-level
                             % sections are begun with similar
                             % \subsection and \subsubsection commands.

The ends  of words and sentences are marked
  by   spaces. It  doesn't matter how many
spaces    you type; one is as good as 100.  The
end of   a line counts as a space.

One   or more   blank lines denote the  end
of  a paragraph.

Since any number of consecutive spaces are treated
like a single one, the formatting of the input
file makes no difference to
      \LaTeX,                % The \LaTeX command generates the LaTeX logo.
but it makes a difference to you.  When you use
\LaTeX, making your input file as easy to read
as possible will be a great help as you write
your document and when you change it.  This sample
file shows how you can add comments to your own input
file.

Because printing is different from typewriting,
there are a number of things that you have to do
differently when preparing an input file than if
you were just typing the document directly.
Quotation marks like
       ``this''
have to be handled specially, as do quotes within
quotes:
       ``\,`this'            % \, separates the double and single quote.
        is what I just
        wrote, not  `that'\,''.

Dashes come in three sizes: an
       intra-word
dash, a medium dash for number ranges like
       1--2,
and a punctuation
       dash---like
this.

A sentence-ending space should be larger than the
space between words within a sentence.  You
sometimes have to type special commands in
conjunction with punctuation characters to get
this right, as in the following sentence.
       Gnats, gnus, etc.\ all  % `\ ' makes an inter-word space.
       begin with G\@.         % \@ marks end-of-sentence punctuation.
You should check the spaces after periods when
reading your output to make sure you haven't
forgotten any special cases.  Generating an
ellipsis
       \ldots\               % `\ ' is needed after `\ldots' because TeX
                             % ignores spaces after command names like \ldots
                             % made from \ + letters.
                             %
                             % Note how a `%' character causes TeX to ignore
                             % the end of the input line, so these blank lines
                             % do not start a new paragraph.
                             %
with the right spacing around the periods requires
a special command.

\LaTeX\ interprets some common characters as
commands, so you must type special commands to
generate them.  These characters include the
following:
       \$ \& \% \# \{ and \}.

In printing, text is usually emphasized with an
       \emph{italic}
type style.

\begin{em}
   A long segment of text can also be emphasized
   in this way.  Text within such a segment can be
   given \emph{additional} emphasis.
\end{em}

It is sometimes necessary to prevent \LaTeX\ from
breaking a line where it might otherwise do so.
This may be at a space, as between the ``Mr.''\ and
``Jones'' in
       ``Mr.~Jones'',        % ~ produces an unbreakable interword space.
or within a word---especially when the word is a
symbol like
       \mbox{\emph{itemnum}}
that makes little sense when hyphenated across
lines.

Footnotes\footnote{This is an example of a footnote.}
pose no problem.

\LaTeX\ is good at typesetting mathematical formulas
like
       \( x-3y + z = 7 \)
or
       \( a_{1} > x^{2n} + y^{2n} > x' \)
or
       \( \ip{A}{B} = \sum_{i} a_{i} b_{i} \).
The spaces you type in a formula are
ignored.  Remember that a letter like
       $x$                   % $ ... $  and  \( ... \)  are equivalent
is a formula when it denotes a mathematical
symbol, and it should be typed as one.

\section{Displayed Text}

Text is displayed by indenting it from the left
margin.  Quotations are commonly displayed.  There
are short quotations
\begin{quote}
   This is a short quotation.  It consists of a
   single paragraph of text.  See how it is formatted.
\end{quote}
and longer ones.
\begin{quotation}
   This is a longer quotation.  It consists of two
   paragraphs of text, neither of which are
   particularly interesting.

   This is the second paragraph of the quotation.  It
   is just as dull as the first paragraph.
\end{quotation}
Another frequently-displayed structure is a list.
The following is an example of an \emph{itemized}
list.
\begin{itemize}
   \item This is the first item of an itemized list.
         Each item in the list is marked with a ``tick''.
         You don't have to worry about what kind of tick
         mark is used.

   \item This is the second item of the list.  It
         contains another list nested inside it.  The inner
         list is an \emph{enumerated} list.
         \begin{enumerate}
            \item This is the first item of an enumerated
                  list that is nested within the itemized list.

            \item This is the second item of the inner list.
                  \LaTeX\ allows you to nest lists deeper than
                  you really should.
         \end{enumerate}
         This is the rest of the second item of the outer
         list.  It is no more interesting than any other
         part of the item.
   \item This is the third item of the list.
\end{itemize}
You can even display poetry.
\begin{verse}
   There is an environment
    for verse \\             % The \\ command separates lines
   Whose features some poets % within a stanza.
   will curse.

                             % One or more blank lines separate stanzas.

   For instead of making\\
   Them do \emph{all} line breaking, \\
   It allows them to put too many words on a line when they'd rather be
   forced to be terse.
\end{verse}

Mathematical formulas may also be displayed.  A
displayed formula
is
one-line long; multiline
formulas require special formatting instructions.
   \[  \ip{\Gamma}{\psi'} = x'' + y^{2} + z_{i}^{n}\]
Don't start a paragraph with a displayed equation,
nor make one a paragraph by itself.

\end{document}               % End of document.

\section{Bench Huge Section 714}
% synthetic bench section 714
% This is a sample LaTeX input file.  (Version of 12 August 2004.)
%
% A '%' character causes TeX to ignore all remaining text on the line,
% and is used for comments like this one.

\documentclass{article}      % Specifies the document class

                             % The preamble begins here.
\title{An Example Document}  % Declares the document's title.
\author{Leslie Lamport}      % Declares the author's name.
\date{January 21, 1994}      % Deleting this command produces today's date.

\newcommand{\ip}[2]{(#1, #2)}
                             % Defines \ip{arg1}{arg2} to mean
                             % (arg1, arg2).

%\newcommand{\ip}[2]{\langle #1 | #2\rangle}
                             % This is an alternative definition of
                             % \ip that is commented out.

\begin{document}             % End of preamble and beginning of text.

\maketitle                   % Produces the title.

This is an example input file.  Comparing it with
the output it generates can show you how to
produce a simple document of your own.

\section{Ordinary Text}      % Produces section heading.  Lower-level
                             % sections are begun with similar
                             % \subsection and \subsubsection commands.

The ends  of words and sentences are marked
  by   spaces. It  doesn't matter how many
spaces    you type; one is as good as 100.  The
end of   a line counts as a space.

One   or more   blank lines denote the  end
of  a paragraph.

Since any number of consecutive spaces are treated
like a single one, the formatting of the input
file makes no difference to
      \LaTeX,                % The \LaTeX command generates the LaTeX logo.
but it makes a difference to you.  When you use
\LaTeX, making your input file as easy to read
as possible will be a great help as you write
your document and when you change it.  This sample
file shows how you can add comments to your own input
file.

Because printing is different from typewriting,
there are a number of things that you have to do
differently when preparing an input file than if
you were just typing the document directly.
Quotation marks like
       ``this''
have to be handled specially, as do quotes within
quotes:
       ``\,`this'            % \, separates the double and single quote.
        is what I just
        wrote, not  `that'\,''.

Dashes come in three sizes: an
       intra-word
dash, a medium dash for number ranges like
       1--2,
and a punctuation
       dash---like
this.

A sentence-ending space should be larger than the
space between words within a sentence.  You
sometimes have to type special commands in
conjunction with punctuation characters to get
this right, as in the following sentence.
       Gnats, gnus, etc.\ all  % `\ ' makes an inter-word space.
       begin with G\@.         % \@ marks end-of-sentence punctuation.
You should check the spaces after periods when
reading your output to make sure you haven't
forgotten any special cases.  Generating an
ellipsis
       \ldots\               % `\ ' is needed after `\ldots' because TeX
                             % ignores spaces after command names like \ldots
                             % made from \ + letters.
                             %
                             % Note how a `%' character causes TeX to ignore
                             % the end of the input line, so these blank lines
                             % do not start a new paragraph.
                             %
with the right spacing around the periods requires
a special command.

\LaTeX\ interprets some common characters as
commands, so you must type special commands to
generate them.  These characters include the
following:
       \$ \& \% \# \{ and \}.

In printing, text is usually emphasized with an
       \emph{italic}
type style.

\begin{em}
   A long segment of text can also be emphasized
   in this way.  Text within such a segment can be
   given \emph{additional} emphasis.
\end{em}

It is sometimes necessary to prevent \LaTeX\ from
breaking a line where it might otherwise do so.
This may be at a space, as between the ``Mr.''\ and
``Jones'' in
       ``Mr.~Jones'',        % ~ produces an unbreakable interword space.
or within a word---especially when the word is a
symbol like
       \mbox{\emph{itemnum}}
that makes little sense when hyphenated across
lines.

Footnotes\footnote{This is an example of a footnote.}
pose no problem.

\LaTeX\ is good at typesetting mathematical formulas
like
       \( x-3y + z = 7 \)
or
       \( a_{1} > x^{2n} + y^{2n} > x' \)
or
       \( \ip{A}{B} = \sum_{i} a_{i} b_{i} \).
The spaces you type in a formula are
ignored.  Remember that a letter like
       $x$                   % $ ... $  and  \( ... \)  are equivalent
is a formula when it denotes a mathematical
symbol, and it should be typed as one.

\section{Displayed Text}

Text is displayed by indenting it from the left
margin.  Quotations are commonly displayed.  There
are short quotations
\begin{quote}
   This is a short quotation.  It consists of a
   single paragraph of text.  See how it is formatted.
\end{quote}
and longer ones.
\begin{quotation}
   This is a longer quotation.  It consists of two
   paragraphs of text, neither of which are
   particularly interesting.

   This is the second paragraph of the quotation.  It
   is just as dull as the first paragraph.
\end{quotation}
Another frequently-displayed structure is a list.
The following is an example of an \emph{itemized}
list.
\begin{itemize}
   \item This is the first item of an itemized list.
         Each item in the list is marked with a ``tick''.
         You don't have to worry about what kind of tick
         mark is used.

   \item This is the second item of the list.  It
         contains another list nested inside it.  The inner
         list is an \emph{enumerated} list.
         \begin{enumerate}
            \item This is the first item of an enumerated
                  list that is nested within the itemized list.

            \item This is the second item of the inner list.
                  \LaTeX\ allows you to nest lists deeper than
                  you really should.
         \end{enumerate}
         This is the rest of the second item of the outer
         list.  It is no more interesting than any other
         part of the item.
   \item This is the third item of the list.
\end{itemize}
You can even display poetry.
\begin{verse}
   There is an environment
    for verse \\             % The \\ command separates lines
   Whose features some poets % within a stanza.
   will curse.

                             % One or more blank lines separate stanzas.

   For instead of making\\
   Them do \emph{all} line breaking, \\
   It allows them to put too many words on a line when they'd rather be
   forced to be terse.
\end{verse}

Mathematical formulas may also be displayed.  A
displayed formula
is
one-line long; multiline
formulas require special formatting instructions.
   \[  \ip{\Gamma}{\psi'} = x'' + y^{2} + z_{i}^{n}\]
Don't start a paragraph with a displayed equation,
nor make one a paragraph by itself.

\end{document}               % End of document.

\section{Bench Huge Section 715}
% synthetic bench section 715
% This is a sample LaTeX input file.  (Version of 12 August 2004.)
%
% A '%' character causes TeX to ignore all remaining text on the line,
% and is used for comments like this one.

\documentclass{article}      % Specifies the document class

                             % The preamble begins here.
\title{An Example Document}  % Declares the document's title.
\author{Leslie Lamport}      % Declares the author's name.
\date{January 21, 1994}      % Deleting this command produces today's date.

\newcommand{\ip}[2]{(#1, #2)}
                             % Defines \ip{arg1}{arg2} to mean
                             % (arg1, arg2).

%\newcommand{\ip}[2]{\langle #1 | #2\rangle}
                             % This is an alternative definition of
                             % \ip that is commented out.

\begin{document}             % End of preamble and beginning of text.

\maketitle                   % Produces the title.

This is an example input file.  Comparing it with
the output it generates can show you how to
produce a simple document of your own.

\section{Ordinary Text}      % Produces section heading.  Lower-level
                             % sections are begun with similar
                             % \subsection and \subsubsection commands.

The ends  of words and sentences are marked
  by   spaces. It  doesn't matter how many
spaces    you type; one is as good as 100.  The
end of   a line counts as a space.

One   or more   blank lines denote the  end
of  a paragraph.

Since any number of consecutive spaces are treated
like a single one, the formatting of the input
file makes no difference to
      \LaTeX,                % The \LaTeX command generates the LaTeX logo.
but it makes a difference to you.  When you use
\LaTeX, making your input file as easy to read
as possible will be a great help as you write
your document and when you change it.  This sample
file shows how you can add comments to your own input
file.

Because printing is different from typewriting,
there are a number of things that you have to do
differently when preparing an input file than if
you were just typing the document directly.
Quotation marks like
       ``this''
have to be handled specially, as do quotes within
quotes:
       ``\,`this'            % \, separates the double and single quote.
        is what I just
        wrote, not  `that'\,''.

Dashes come in three sizes: an
       intra-word
dash, a medium dash for number ranges like
       1--2,
and a punctuation
       dash---like
this.

A sentence-ending space should be larger than the
space between words within a sentence.  You
sometimes have to type special commands in
conjunction with punctuation characters to get
this right, as in the following sentence.
       Gnats, gnus, etc.\ all  % `\ ' makes an inter-word space.
       begin with G\@.         % \@ marks end-of-sentence punctuation.
You should check the spaces after periods when
reading your output to make sure you haven't
forgotten any special cases.  Generating an
ellipsis
       \ldots\               % `\ ' is needed after `\ldots' because TeX
                             % ignores spaces after command names like \ldots
                             % made from \ + letters.
                             %
                             % Note how a `%' character causes TeX to ignore
                             % the end of the input line, so these blank lines
                             % do not start a new paragraph.
                             %
with the right spacing around the periods requires
a special command.

\LaTeX\ interprets some common characters as
commands, so you must type special commands to
generate them.  These characters include the
following:
       \$ \& \% \# \{ and \}.

In printing, text is usually emphasized with an
       \emph{italic}
type style.

\begin{em}
   A long segment of text can also be emphasized
   in this way.  Text within such a segment can be
   given \emph{additional} emphasis.
\end{em}

It is sometimes necessary to prevent \LaTeX\ from
breaking a line where it might otherwise do so.
This may be at a space, as between the ``Mr.''\ and
``Jones'' in
       ``Mr.~Jones'',        % ~ produces an unbreakable interword space.
or within a word---especially when the word is a
symbol like
       \mbox{\emph{itemnum}}
that makes little sense when hyphenated across
lines.

Footnotes\footnote{This is an example of a footnote.}
pose no problem.

\LaTeX\ is good at typesetting mathematical formulas
like
       \( x-3y + z = 7 \)
or
       \( a_{1} > x^{2n} + y^{2n} > x' \)
or
       \( \ip{A}{B} = \sum_{i} a_{i} b_{i} \).
The spaces you type in a formula are
ignored.  Remember that a letter like
       $x$                   % $ ... $  and  \( ... \)  are equivalent
is a formula when it denotes a mathematical
symbol, and it should be typed as one.

\section{Displayed Text}

Text is displayed by indenting it from the left
margin.  Quotations are commonly displayed.  There
are short quotations
\begin{quote}
   This is a short quotation.  It consists of a
   single paragraph of text.  See how it is formatted.
\end{quote}
and longer ones.
\begin{quotation}
   This is a longer quotation.  It consists of two
   paragraphs of text, neither of which are
   particularly interesting.

   This is the second paragraph of the quotation.  It
   is just as dull as the first paragraph.
\end{quotation}
Another frequently-displayed structure is a list.
The following is an example of an \emph{itemized}
list.
\begin{itemize}
   \item This is the first item of an itemized list.
         Each item in the list is marked with a ``tick''.
         You don't have to worry about what kind of tick
         mark is used.

   \item This is the second item of the list.  It
         contains another list nested inside it.  The inner
         list is an \emph{enumerated} list.
         \begin{enumerate}
            \item This is the first item of an enumerated
                  list that is nested within the itemized list.

            \item This is the second item of the inner list.
                  \LaTeX\ allows you to nest lists deeper than
                  you really should.
         \end{enumerate}
         This is the rest of the second item of the outer
         list.  It is no more interesting than any other
         part of the item.
   \item This is the third item of the list.
\end{itemize}
You can even display poetry.
\begin{verse}
   There is an environment
    for verse \\             % The \\ command separates lines
   Whose features some poets % within a stanza.
   will curse.

                             % One or more blank lines separate stanzas.

   For instead of making\\
   Them do \emph{all} line breaking, \\
   It allows them to put too many words on a line when they'd rather be
   forced to be terse.
\end{verse}

Mathematical formulas may also be displayed.  A
displayed formula
is
one-line long; multiline
formulas require special formatting instructions.
   \[  \ip{\Gamma}{\psi'} = x'' + y^{2} + z_{i}^{n}\]
Don't start a paragraph with a displayed equation,
nor make one a paragraph by itself.

\end{document}               % End of document.

\section{Bench Huge Section 716}
% synthetic bench section 716
% This is a sample LaTeX input file.  (Version of 12 August 2004.)
%
% A '%' character causes TeX to ignore all remaining text on the line,
% and is used for comments like this one.

\documentclass{article}      % Specifies the document class

                             % The preamble begins here.
\title{An Example Document}  % Declares the document's title.
\author{Leslie Lamport}      % Declares the author's name.
\date{January 21, 1994}      % Deleting this command produces today's date.

\newcommand{\ip}[2]{(#1, #2)}
                             % Defines \ip{arg1}{arg2} to mean
                             % (arg1, arg2).

%\newcommand{\ip}[2]{\langle #1 | #2\rangle}
                             % This is an alternative definition of
                             % \ip that is commented out.

\begin{document}             % End of preamble and beginning of text.

\maketitle                   % Produces the title.

This is an example input file.  Comparing it with
the output it generates can show you how to
produce a simple document of your own.

\section{Ordinary Text}      % Produces section heading.  Lower-level
                             % sections are begun with similar
                             % \subsection and \subsubsection commands.

The ends  of words and sentences are marked
  by   spaces. It  doesn't matter how many
spaces    you type; one is as good as 100.  The
end of   a line counts as a space.

One   or more   blank lines denote the  end
of  a paragraph.

Since any number of consecutive spaces are treated
like a single one, the formatting of the input
file makes no difference to
      \LaTeX,                % The \LaTeX command generates the LaTeX logo.
but it makes a difference to you.  When you use
\LaTeX, making your input file as easy to read
as possible will be a great help as you write
your document and when you change it.  This sample
file shows how you can add comments to your own input
file.

Because printing is different from typewriting,
there are a number of things that you have to do
differently when preparing an input file than if
you were just typing the document directly.
Quotation marks like
       ``this''
have to be handled specially, as do quotes within
quotes:
       ``\,`this'            % \, separates the double and single quote.
        is what I just
        wrote, not  `that'\,''.

Dashes come in three sizes: an
       intra-word
dash, a medium dash for number ranges like
       1--2,
and a punctuation
       dash---like
this.

A sentence-ending space should be larger than the
space between words within a sentence.  You
sometimes have to type special commands in
conjunction with punctuation characters to get
this right, as in the following sentence.
       Gnats, gnus, etc.\ all  % `\ ' makes an inter-word space.
       begin with G\@.         % \@ marks end-of-sentence punctuation.
You should check the spaces after periods when
reading your output to make sure you haven't
forgotten any special cases.  Generating an
ellipsis
       \ldots\               % `\ ' is needed after `\ldots' because TeX
                             % ignores spaces after command names like \ldots
                             % made from \ + letters.
                             %
                             % Note how a `%' character causes TeX to ignore
                             % the end of the input line, so these blank lines
                             % do not start a new paragraph.
                             %
with the right spacing around the periods requires
a special command.

\LaTeX\ interprets some common characters as
commands, so you must type special commands to
generate them.  These characters include the
following:
       \$ \& \% \# \{ and \}.

In printing, text is usually emphasized with an
       \emph{italic}
type style.

\begin{em}
   A long segment of text can also be emphasized
   in this way.  Text within such a segment can be
   given \emph{additional} emphasis.
\end{em}

It is sometimes necessary to prevent \LaTeX\ from
breaking a line where it might otherwise do so.
This may be at a space, as between the ``Mr.''\ and
``Jones'' in
       ``Mr.~Jones'',        % ~ produces an unbreakable interword space.
or within a word---especially when the word is a
symbol like
       \mbox{\emph{itemnum}}
that makes little sense when hyphenated across
lines.

Footnotes\footnote{This is an example of a footnote.}
pose no problem.

\LaTeX\ is good at typesetting mathematical formulas
like
       \( x-3y + z = 7 \)
or
       \( a_{1} > x^{2n} + y^{2n} > x' \)
or
       \( \ip{A}{B} = \sum_{i} a_{i} b_{i} \).
The spaces you type in a formula are
ignored.  Remember that a letter like
       $x$                   % $ ... $  and  \( ... \)  are equivalent
is a formula when it denotes a mathematical
symbol, and it should be typed as one.

\section{Displayed Text}

Text is displayed by indenting it from the left
margin.  Quotations are commonly displayed.  There
are short quotations
\begin{quote}
   This is a short quotation.  It consists of a
   single paragraph of text.  See how it is formatted.
\end{quote}
and longer ones.
\begin{quotation}
   This is a longer quotation.  It consists of two
   paragraphs of text, neither of which are
   particularly interesting.

   This is the second paragraph of the quotation.  It
   is just as dull as the first paragraph.
\end{quotation}
Another frequently-displayed structure is a list.
The following is an example of an \emph{itemized}
list.
\begin{itemize}
   \item This is the first item of an itemized list.
         Each item in the list is marked with a ``tick''.
         You don't have to worry about what kind of tick
         mark is used.

   \item This is the second item of the list.  It
         contains another list nested inside it.  The inner
         list is an \emph{enumerated} list.
         \begin{enumerate}
            \item This is the first item of an enumerated
                  list that is nested within the itemized list.

            \item This is the second item of the inner list.
                  \LaTeX\ allows you to nest lists deeper than
                  you really should.
         \end{enumerate}
         This is the rest of the second item of the outer
         list.  It is no more interesting than any other
         part of the item.
   \item This is the third item of the list.
\end{itemize}
You can even display poetry.
\begin{verse}
   There is an environment
    for verse \\             % The \\ command separates lines
   Whose features some poets % within a stanza.
   will curse.

                             % One or more blank lines separate stanzas.

   For instead of making\\
   Them do \emph{all} line breaking, \\
   It allows them to put too many words on a line when they'd rather be
   forced to be terse.
\end{verse}

Mathematical formulas may also be displayed.  A
displayed formula
is
one-line long; multiline
formulas require special formatting instructions.
   \[  \ip{\Gamma}{\psi'} = x'' + y^{2} + z_{i}^{n}\]
Don't start a paragraph with a displayed equation,
nor make one a paragraph by itself.

\end{document}               % End of document.

\section{Bench Huge Section 717}
% synthetic bench section 717
% This is a sample LaTeX input file.  (Version of 12 August 2004.)
%
% A '%' character causes TeX to ignore all remaining text on the line,
% and is used for comments like this one.

\documentclass{article}      % Specifies the document class

                             % The preamble begins here.
\title{An Example Document}  % Declares the document's title.
\author{Leslie Lamport}      % Declares the author's name.
\date{January 21, 1994}      % Deleting this command produces today's date.

\newcommand{\ip}[2]{(#1, #2)}
                             % Defines \ip{arg1}{arg2} to mean
                             % (arg1, arg2).

%\newcommand{\ip}[2]{\langle #1 | #2\rangle}
                             % This is an alternative definition of
                             % \ip that is commented out.

\begin{document}             % End of preamble and beginning of text.

\maketitle                   % Produces the title.

This is an example input file.  Comparing it with
the output it generates can show you how to
produce a simple document of your own.

\section{Ordinary Text}      % Produces section heading.  Lower-level
                             % sections are begun with similar
                             % \subsection and \subsubsection commands.

The ends  of words and sentences are marked
  by   spaces. It  doesn't matter how many
spaces    you type; one is as good as 100.  The
end of   a line counts as a space.

One   or more   blank lines denote the  end
of  a paragraph.

Since any number of consecutive spaces are treated
like a single one, the formatting of the input
file makes no difference to
      \LaTeX,                % The \LaTeX command generates the LaTeX logo.
but it makes a difference to you.  When you use
\LaTeX, making your input file as easy to read
as possible will be a great help as you write
your document and when you change it.  This sample
file shows how you can add comments to your own input
file.

Because printing is different from typewriting,
there are a number of things that you have to do
differently when preparing an input file than if
you were just typing the document directly.
Quotation marks like
       ``this''
have to be handled specially, as do quotes within
quotes:
       ``\,`this'            % \, separates the double and single quote.
        is what I just
        wrote, not  `that'\,''.

Dashes come in three sizes: an
       intra-word
dash, a medium dash for number ranges like
       1--2,
and a punctuation
       dash---like
this.

A sentence-ending space should be larger than the
space between words within a sentence.  You
sometimes have to type special commands in
conjunction with punctuation characters to get
this right, as in the following sentence.
       Gnats, gnus, etc.\ all  % `\ ' makes an inter-word space.
       begin with G\@.         % \@ marks end-of-sentence punctuation.
You should check the spaces after periods when
reading your output to make sure you haven't
forgotten any special cases.  Generating an
ellipsis
       \ldots\               % `\ ' is needed after `\ldots' because TeX
                             % ignores spaces after command names like \ldots
                             % made from \ + letters.
                             %
                             % Note how a `%' character causes TeX to ignore
                             % the end of the input line, so these blank lines
                             % do not start a new paragraph.
                             %
with the right spacing around the periods requires
a special command.

\LaTeX\ interprets some common characters as
commands, so you must type special commands to
generate them.  These characters include the
following:
       \$ \& \% \# \{ and \}.

In printing, text is usually emphasized with an
       \emph{italic}
type style.

\begin{em}
   A long segment of text can also be emphasized
   in this way.  Text within such a segment can be
   given \emph{additional} emphasis.
\end{em}

It is sometimes necessary to prevent \LaTeX\ from
breaking a line where it might otherwise do so.
This may be at a space, as between the ``Mr.''\ and
``Jones'' in
       ``Mr.~Jones'',        % ~ produces an unbreakable interword space.
or within a word---especially when the word is a
symbol like
       \mbox{\emph{itemnum}}
that makes little sense when hyphenated across
lines.

Footnotes\footnote{This is an example of a footnote.}
pose no problem.

\LaTeX\ is good at typesetting mathematical formulas
like
       \( x-3y + z = 7 \)
or
       \( a_{1} > x^{2n} + y^{2n} > x' \)
or
       \( \ip{A}{B} = \sum_{i} a_{i} b_{i} \).
The spaces you type in a formula are
ignored.  Remember that a letter like
       $x$                   % $ ... $  and  \( ... \)  are equivalent
is a formula when it denotes a mathematical
symbol, and it should be typed as one.

\section{Displayed Text}

Text is displayed by indenting it from the left
margin.  Quotations are commonly displayed.  There
are short quotations
\begin{quote}
   This is a short quotation.  It consists of a
   single paragraph of text.  See how it is formatted.
\end{quote}
and longer ones.
\begin{quotation}
   This is a longer quotation.  It consists of two
   paragraphs of text, neither of which are
   particularly interesting.

   This is the second paragraph of the quotation.  It
   is just as dull as the first paragraph.
\end{quotation}
Another frequently-displayed structure is a list.
The following is an example of an \emph{itemized}
list.
\begin{itemize}
   \item This is the first item of an itemized list.
         Each item in the list is marked with a ``tick''.
         You don't have to worry about what kind of tick
         mark is used.

   \item This is the second item of the list.  It
         contains another list nested inside it.  The inner
         list is an \emph{enumerated} list.
         \begin{enumerate}
            \item This is the first item of an enumerated
                  list that is nested within the itemized list.

            \item This is the second item of the inner list.
                  \LaTeX\ allows you to nest lists deeper than
                  you really should.
         \end{enumerate}
         This is the rest of the second item of the outer
         list.  It is no more interesting than any other
         part of the item.
   \item This is the third item of the list.
\end{itemize}
You can even display poetry.
\begin{verse}
   There is an environment
    for verse \\             % The \\ command separates lines
   Whose features some poets % within a stanza.
   will curse.

                             % One or more blank lines separate stanzas.

   For instead of making\\
   Them do \emph{all} line breaking, \\
   It allows them to put too many words on a line when they'd rather be
   forced to be terse.
\end{verse}

Mathematical formulas may also be displayed.  A
displayed formula
is
one-line long; multiline
formulas require special formatting instructions.
   \[  \ip{\Gamma}{\psi'} = x'' + y^{2} + z_{i}^{n}\]
Don't start a paragraph with a displayed equation,
nor make one a paragraph by itself.

\end{document}               % End of document.

\section{Bench Huge Section 718}
% synthetic bench section 718
% This is a sample LaTeX input file.  (Version of 12 August 2004.)
%
% A '%' character causes TeX to ignore all remaining text on the line,
% and is used for comments like this one.

\documentclass{article}      % Specifies the document class

                             % The preamble begins here.
\title{An Example Document}  % Declares the document's title.
\author{Leslie Lamport}      % Declares the author's name.
\date{January 21, 1994}      % Deleting this command produces today's date.

\newcommand{\ip}[2]{(#1, #2)}
                             % Defines \ip{arg1}{arg2} to mean
                             % (arg1, arg2).

%\newcommand{\ip}[2]{\langle #1 | #2\rangle}
                             % This is an alternative definition of
                             % \ip that is commented out.

\begin{document}             % End of preamble and beginning of text.

\maketitle                   % Produces the title.

This is an example input file.  Comparing it with
the output it generates can show you how to
produce a simple document of your own.

\section{Ordinary Text}      % Produces section heading.  Lower-level
                             % sections are begun with similar
                             % \subsection and \subsubsection commands.

The ends  of words and sentences are marked
  by   spaces. It  doesn't matter how many
spaces    you type; one is as good as 100.  The
end of   a line counts as a space.

One   or more   blank lines denote the  end
of  a paragraph.

Since any number of consecutive spaces are treated
like a single one, the formatting of the input
file makes no difference to
      \LaTeX,                % The \LaTeX command generates the LaTeX logo.
but it makes a difference to you.  When you use
\LaTeX, making your input file as easy to read
as possible will be a great help as you write
your document and when you change it.  This sample
file shows how you can add comments to your own input
file.

Because printing is different from typewriting,
there are a number of things that you have to do
differently when preparing an input file than if
you were just typing the document directly.
Quotation marks like
       ``this''
have to be handled specially, as do quotes within
quotes:
       ``\,`this'            % \, separates the double and single quote.
        is what I just
        wrote, not  `that'\,''.

Dashes come in three sizes: an
       intra-word
dash, a medium dash for number ranges like
       1--2,
and a punctuation
       dash---like
this.

A sentence-ending space should be larger than the
space between words within a sentence.  You
sometimes have to type special commands in
conjunction with punctuation characters to get
this right, as in the following sentence.
       Gnats, gnus, etc.\ all  % `\ ' makes an inter-word space.
       begin with G\@.         % \@ marks end-of-sentence punctuation.
You should check the spaces after periods when
reading your output to make sure you haven't
forgotten any special cases.  Generating an
ellipsis
       \ldots\               % `\ ' is needed after `\ldots' because TeX
                             % ignores spaces after command names like \ldots
                             % made from \ + letters.
                             %
                             % Note how a `%' character causes TeX to ignore
                             % the end of the input line, so these blank lines
                             % do not start a new paragraph.
                             %
with the right spacing around the periods requires
a special command.

\LaTeX\ interprets some common characters as
commands, so you must type special commands to
generate them.  These characters include the
following:
       \$ \& \% \# \{ and \}.

In printing, text is usually emphasized with an
       \emph{italic}
type style.

\begin{em}
   A long segment of text can also be emphasized
   in this way.  Text within such a segment can be
   given \emph{additional} emphasis.
\end{em}

It is sometimes necessary to prevent \LaTeX\ from
breaking a line where it might otherwise do so.
This may be at a space, as between the ``Mr.''\ and
``Jones'' in
       ``Mr.~Jones'',        % ~ produces an unbreakable interword space.
or within a word---especially when the word is a
symbol like
       \mbox{\emph{itemnum}}
that makes little sense when hyphenated across
lines.

Footnotes\footnote{This is an example of a footnote.}
pose no problem.

\LaTeX\ is good at typesetting mathematical formulas
like
       \( x-3y + z = 7 \)
or
       \( a_{1} > x^{2n} + y^{2n} > x' \)
or
       \( \ip{A}{B} = \sum_{i} a_{i} b_{i} \).
The spaces you type in a formula are
ignored.  Remember that a letter like
       $x$                   % $ ... $  and  \( ... \)  are equivalent
is a formula when it denotes a mathematical
symbol, and it should be typed as one.

\section{Displayed Text}

Text is displayed by indenting it from the left
margin.  Quotations are commonly displayed.  There
are short quotations
\begin{quote}
   This is a short quotation.  It consists of a
   single paragraph of text.  See how it is formatted.
\end{quote}
and longer ones.
\begin{quotation}
   This is a longer quotation.  It consists of two
   paragraphs of text, neither of which are
   particularly interesting.

   This is the second paragraph of the quotation.  It
   is just as dull as the first paragraph.
\end{quotation}
Another frequently-displayed structure is a list.
The following is an example of an \emph{itemized}
list.
\begin{itemize}
   \item This is the first item of an itemized list.
         Each item in the list is marked with a ``tick''.
         You don't have to worry about what kind of tick
         mark is used.

   \item This is the second item of the list.  It
         contains another list nested inside it.  The inner
         list is an \emph{enumerated} list.
         \begin{enumerate}
            \item This is the first item of an enumerated
                  list that is nested within the itemized list.

            \item This is the second item of the inner list.
                  \LaTeX\ allows you to nest lists deeper than
                  you really should.
         \end{enumerate}
         This is the rest of the second item of the outer
         list.  It is no more interesting than any other
         part of the item.
   \item This is the third item of the list.
\end{itemize}
You can even display poetry.
\begin{verse}
   There is an environment
    for verse \\             % The \\ command separates lines
   Whose features some poets % within a stanza.
   will curse.

                             % One or more blank lines separate stanzas.

   For instead of making\\
   Them do \emph{all} line breaking, \\
   It allows them to put too many words on a line when they'd rather be
   forced to be terse.
\end{verse}

Mathematical formulas may also be displayed.  A
displayed formula
is
one-line long; multiline
formulas require special formatting instructions.
   \[  \ip{\Gamma}{\psi'} = x'' + y^{2} + z_{i}^{n}\]
Don't start a paragraph with a displayed equation,
nor make one a paragraph by itself.

\end{document}               % End of document.

\section{Bench Huge Section 719}
% synthetic bench section 719
% This is a sample LaTeX input file.  (Version of 12 August 2004.)
%
% A '%' character causes TeX to ignore all remaining text on the line,
% and is used for comments like this one.

\documentclass{article}      % Specifies the document class

                             % The preamble begins here.
\title{An Example Document}  % Declares the document's title.
\author{Leslie Lamport}      % Declares the author's name.
\date{January 21, 1994}      % Deleting this command produces today's date.

\newcommand{\ip}[2]{(#1, #2)}
                             % Defines \ip{arg1}{arg2} to mean
                             % (arg1, arg2).

%\newcommand{\ip}[2]{\langle #1 | #2\rangle}
                             % This is an alternative definition of
                             % \ip that is commented out.

\begin{document}             % End of preamble and beginning of text.

\maketitle                   % Produces the title.

This is an example input file.  Comparing it with
the output it generates can show you how to
produce a simple document of your own.

\section{Ordinary Text}      % Produces section heading.  Lower-level
                             % sections are begun with similar
                             % \subsection and \subsubsection commands.

The ends  of words and sentences are marked
  by   spaces. It  doesn't matter how many
spaces    you type; one is as good as 100.  The
end of   a line counts as a space.

One   or more   blank lines denote the  end
of  a paragraph.

Since any number of consecutive spaces are treated
like a single one, the formatting of the input
file makes no difference to
      \LaTeX,                % The \LaTeX command generates the LaTeX logo.
but it makes a difference to you.  When you use
\LaTeX, making your input file as easy to read
as possible will be a great help as you write
your document and when you change it.  This sample
file shows how you can add comments to your own input
file.

Because printing is different from typewriting,
there are a number of things that you have to do
differently when preparing an input file than if
you were just typing the document directly.
Quotation marks like
       ``this''
have to be handled specially, as do quotes within
quotes:
       ``\,`this'            % \, separates the double and single quote.
        is what I just
        wrote, not  `that'\,''.

Dashes come in three sizes: an
       intra-word
dash, a medium dash for number ranges like
       1--2,
and a punctuation
       dash---like
this.

A sentence-ending space should be larger than the
space between words within a sentence.  You
sometimes have to type special commands in
conjunction with punctuation characters to get
this right, as in the following sentence.
       Gnats, gnus, etc.\ all  % `\ ' makes an inter-word space.
       begin with G\@.         % \@ marks end-of-sentence punctuation.
You should check the spaces after periods when
reading your output to make sure you haven't
forgotten any special cases.  Generating an
ellipsis
       \ldots\               % `\ ' is needed after `\ldots' because TeX
                             % ignores spaces after command names like \ldots
                             % made from \ + letters.
                             %
                             % Note how a `%' character causes TeX to ignore
                             % the end of the input line, so these blank lines
                             % do not start a new paragraph.
                             %
with the right spacing around the periods requires
a special command.

\LaTeX\ interprets some common characters as
commands, so you must type special commands to
generate them.  These characters include the
following:
       \$ \& \% \# \{ and \}.

In printing, text is usually emphasized with an
       \emph{italic}
type style.

\begin{em}
   A long segment of text can also be emphasized
   in this way.  Text within such a segment can be
   given \emph{additional} emphasis.
\end{em}

It is sometimes necessary to prevent \LaTeX\ from
breaking a line where it might otherwise do so.
This may be at a space, as between the ``Mr.''\ and
``Jones'' in
       ``Mr.~Jones'',        % ~ produces an unbreakable interword space.
or within a word---especially when the word is a
symbol like
       \mbox{\emph{itemnum}}
that makes little sense when hyphenated across
lines.

Footnotes\footnote{This is an example of a footnote.}
pose no problem.

\LaTeX\ is good at typesetting mathematical formulas
like
       \( x-3y + z = 7 \)
or
       \( a_{1} > x^{2n} + y^{2n} > x' \)
or
       \( \ip{A}{B} = \sum_{i} a_{i} b_{i} \).
The spaces you type in a formula are
ignored.  Remember that a letter like
       $x$                   % $ ... $  and  \( ... \)  are equivalent
is a formula when it denotes a mathematical
symbol, and it should be typed as one.

\section{Displayed Text}

Text is displayed by indenting it from the left
margin.  Quotations are commonly displayed.  There
are short quotations
\begin{quote}
   This is a short quotation.  It consists of a
   single paragraph of text.  See how it is formatted.
\end{quote}
and longer ones.
\begin{quotation}
   This is a longer quotation.  It consists of two
   paragraphs of text, neither of which are
   particularly interesting.

   This is the second paragraph of the quotation.  It
   is just as dull as the first paragraph.
\end{quotation}
Another frequently-displayed structure is a list.
The following is an example of an \emph{itemized}
list.
\begin{itemize}
   \item This is the first item of an itemized list.
         Each item in the list is marked with a ``tick''.
         You don't have to worry about what kind of tick
         mark is used.

   \item This is the second item of the list.  It
         contains another list nested inside it.  The inner
         list is an \emph{enumerated} list.
         \begin{enumerate}
            \item This is the first item of an enumerated
                  list that is nested within the itemized list.

            \item This is the second item of the inner list.
                  \LaTeX\ allows you to nest lists deeper than
                  you really should.
         \end{enumerate}
         This is the rest of the second item of the outer
         list.  It is no more interesting than any other
         part of the item.
   \item This is the third item of the list.
\end{itemize}
You can even display poetry.
\begin{verse}
   There is an environment
    for verse \\             % The \\ command separates lines
   Whose features some poets % within a stanza.
   will curse.

                             % One or more blank lines separate stanzas.

   For instead of making\\
   Them do \emph{all} line breaking, \\
   It allows them to put too many words on a line when they'd rather be
   forced to be terse.
\end{verse}

Mathematical formulas may also be displayed.  A
displayed formula
is
one-line long; multiline
formulas require special formatting instructions.
   \[  \ip{\Gamma}{\psi'} = x'' + y^{2} + z_{i}^{n}\]
Don't start a paragraph with a displayed equation,
nor make one a paragraph by itself.

\end{document}               % End of document.

\section{Bench Huge Section 720}
% synthetic bench section 720
% This is a sample LaTeX input file.  (Version of 12 August 2004.)
%
% A '%' character causes TeX to ignore all remaining text on the line,
% and is used for comments like this one.

\documentclass{article}      % Specifies the document class

                             % The preamble begins here.
\title{An Example Document}  % Declares the document's title.
\author{Leslie Lamport}      % Declares the author's name.
\date{January 21, 1994}      % Deleting this command produces today's date.

\newcommand{\ip}[2]{(#1, #2)}
                             % Defines \ip{arg1}{arg2} to mean
                             % (arg1, arg2).

%\newcommand{\ip}[2]{\langle #1 | #2\rangle}
                             % This is an alternative definition of
                             % \ip that is commented out.

\begin{document}             % End of preamble and beginning of text.

\maketitle                   % Produces the title.

This is an example input file.  Comparing it with
the output it generates can show you how to
produce a simple document of your own.

\section{Ordinary Text}      % Produces section heading.  Lower-level
                             % sections are begun with similar
                             % \subsection and \subsubsection commands.

The ends  of words and sentences are marked
  by   spaces. It  doesn't matter how many
spaces    you type; one is as good as 100.  The
end of   a line counts as a space.

One   or more   blank lines denote the  end
of  a paragraph.

Since any number of consecutive spaces are treated
like a single one, the formatting of the input
file makes no difference to
      \LaTeX,                % The \LaTeX command generates the LaTeX logo.
but it makes a difference to you.  When you use
\LaTeX, making your input file as easy to read
as possible will be a great help as you write
your document and when you change it.  This sample
file shows how you can add comments to your own input
file.

Because printing is different from typewriting,
there are a number of things that you have to do
differently when preparing an input file than if
you were just typing the document directly.
Quotation marks like
       ``this''
have to be handled specially, as do quotes within
quotes:
       ``\,`this'            % \, separates the double and single quote.
        is what I just
        wrote, not  `that'\,''.

Dashes come in three sizes: an
       intra-word
dash, a medium dash for number ranges like
       1--2,
and a punctuation
       dash---like
this.

A sentence-ending space should be larger than the
space between words within a sentence.  You
sometimes have to type special commands in
conjunction with punctuation characters to get
this right, as in the following sentence.
       Gnats, gnus, etc.\ all  % `\ ' makes an inter-word space.
       begin with G\@.         % \@ marks end-of-sentence punctuation.
You should check the spaces after periods when
reading your output to make sure you haven't
forgotten any special cases.  Generating an
ellipsis
       \ldots\               % `\ ' is needed after `\ldots' because TeX
                             % ignores spaces after command names like \ldots
                             % made from \ + letters.
                             %
                             % Note how a `%' character causes TeX to ignore
                             % the end of the input line, so these blank lines
                             % do not start a new paragraph.
                             %
with the right spacing around the periods requires
a special command.

\LaTeX\ interprets some common characters as
commands, so you must type special commands to
generate them.  These characters include the
following:
       \$ \& \% \# \{ and \}.

In printing, text is usually emphasized with an
       \emph{italic}
type style.

\begin{em}
   A long segment of text can also be emphasized
   in this way.  Text within such a segment can be
   given \emph{additional} emphasis.
\end{em}

It is sometimes necessary to prevent \LaTeX\ from
breaking a line where it might otherwise do so.
This may be at a space, as between the ``Mr.''\ and
``Jones'' in
       ``Mr.~Jones'',        % ~ produces an unbreakable interword space.
or within a word---especially when the word is a
symbol like
       \mbox{\emph{itemnum}}
that makes little sense when hyphenated across
lines.

Footnotes\footnote{This is an example of a footnote.}
pose no problem.

\LaTeX\ is good at typesetting mathematical formulas
like
       \( x-3y + z = 7 \)
or
       \( a_{1} > x^{2n} + y^{2n} > x' \)
or
       \( \ip{A}{B} = \sum_{i} a_{i} b_{i} \).
The spaces you type in a formula are
ignored.  Remember that a letter like
       $x$                   % $ ... $  and  \( ... \)  are equivalent
is a formula when it denotes a mathematical
symbol, and it should be typed as one.

\section{Displayed Text}

Text is displayed by indenting it from the left
margin.  Quotations are commonly displayed.  There
are short quotations
\begin{quote}
   This is a short quotation.  It consists of a
   single paragraph of text.  See how it is formatted.
\end{quote}
and longer ones.
\begin{quotation}
   This is a longer quotation.  It consists of two
   paragraphs of text, neither of which are
   particularly interesting.

   This is the second paragraph of the quotation.  It
   is just as dull as the first paragraph.
\end{quotation}
Another frequently-displayed structure is a list.
The following is an example of an \emph{itemized}
list.
\begin{itemize}
   \item This is the first item of an itemized list.
         Each item in the list is marked with a ``tick''.
         You don't have to worry about what kind of tick
         mark is used.

   \item This is the second item of the list.  It
         contains another list nested inside it.  The inner
         list is an \emph{enumerated} list.
         \begin{enumerate}
            \item This is the first item of an enumerated
                  list that is nested within the itemized list.

            \item This is the second item of the inner list.
                  \LaTeX\ allows you to nest lists deeper than
                  you really should.
         \end{enumerate}
         This is the rest of the second item of the outer
         list.  It is no more interesting than any other
         part of the item.
   \item This is the third item of the list.
\end{itemize}
You can even display poetry.
\begin{verse}
   There is an environment
    for verse \\             % The \\ command separates lines
   Whose features some poets % within a stanza.
   will curse.

                             % One or more blank lines separate stanzas.

   For instead of making\\
   Them do \emph{all} line breaking, \\
   It allows them to put too many words on a line when they'd rather be
   forced to be terse.
\end{verse}

Mathematical formulas may also be displayed.  A
displayed formula
is
one-line long; multiline
formulas require special formatting instructions.
   \[  \ip{\Gamma}{\psi'} = x'' + y^{2} + z_{i}^{n}\]
Don't start a paragraph with a displayed equation,
nor make one a paragraph by itself.

\end{document}               % End of document.

\section{Bench Huge Section 721}
% synthetic bench section 721
% This is a sample LaTeX input file.  (Version of 12 August 2004.)
%
% A '%' character causes TeX to ignore all remaining text on the line,
% and is used for comments like this one.

\documentclass{article}      % Specifies the document class

                             % The preamble begins here.
\title{An Example Document}  % Declares the document's title.
\author{Leslie Lamport}      % Declares the author's name.
\date{January 21, 1994}      % Deleting this command produces today's date.

\newcommand{\ip}[2]{(#1, #2)}
                             % Defines \ip{arg1}{arg2} to mean
                             % (arg1, arg2).

%\newcommand{\ip}[2]{\langle #1 | #2\rangle}
                             % This is an alternative definition of
                             % \ip that is commented out.

\begin{document}             % End of preamble and beginning of text.

\maketitle                   % Produces the title.

This is an example input file.  Comparing it with
the output it generates can show you how to
produce a simple document of your own.

\section{Ordinary Text}      % Produces section heading.  Lower-level
                             % sections are begun with similar
                             % \subsection and \subsubsection commands.

The ends  of words and sentences are marked
  by   spaces. It  doesn't matter how many
spaces    you type; one is as good as 100.  The
end of   a line counts as a space.

One   or more   blank lines denote the  end
of  a paragraph.

Since any number of consecutive spaces are treated
like a single one, the formatting of the input
file makes no difference to
      \LaTeX,                % The \LaTeX command generates the LaTeX logo.
but it makes a difference to you.  When you use
\LaTeX, making your input file as easy to read
as possible will be a great help as you write
your document and when you change it.  This sample
file shows how you can add comments to your own input
file.

Because printing is different from typewriting,
there are a number of things that you have to do
differently when preparing an input file than if
you were just typing the document directly.
Quotation marks like
       ``this''
have to be handled specially, as do quotes within
quotes:
       ``\,`this'            % \, separates the double and single quote.
        is what I just
        wrote, not  `that'\,''.

Dashes come in three sizes: an
       intra-word
dash, a medium dash for number ranges like
       1--2,
and a punctuation
       dash---like
this.

A sentence-ending space should be larger than the
space between words within a sentence.  You
sometimes have to type special commands in
conjunction with punctuation characters to get
this right, as in the following sentence.
       Gnats, gnus, etc.\ all  % `\ ' makes an inter-word space.
       begin with G\@.         % \@ marks end-of-sentence punctuation.
You should check the spaces after periods when
reading your output to make sure you haven't
forgotten any special cases.  Generating an
ellipsis
       \ldots\               % `\ ' is needed after `\ldots' because TeX
                             % ignores spaces after command names like \ldots
                             % made from \ + letters.
                             %
                             % Note how a `%' character causes TeX to ignore
                             % the end of the input line, so these blank lines
                             % do not start a new paragraph.
                             %
with the right spacing around the periods requires
a special command.

\LaTeX\ interprets some common characters as
commands, so you must type special commands to
generate them.  These characters include the
following:
       \$ \& \% \# \{ and \}.

In printing, text is usually emphasized with an
       \emph{italic}
type style.

\begin{em}
   A long segment of text can also be emphasized
   in this way.  Text within such a segment can be
   given \emph{additional} emphasis.
\end{em}

It is sometimes necessary to prevent \LaTeX\ from
breaking a line where it might otherwise do so.
This may be at a space, as between the ``Mr.''\ and
``Jones'' in
       ``Mr.~Jones'',        % ~ produces an unbreakable interword space.
or within a word---especially when the word is a
symbol like
       \mbox{\emph{itemnum}}
that makes little sense when hyphenated across
lines.

Footnotes\footnote{This is an example of a footnote.}
pose no problem.

\LaTeX\ is good at typesetting mathematical formulas
like
       \( x-3y + z = 7 \)
or
       \( a_{1} > x^{2n} + y^{2n} > x' \)
or
       \( \ip{A}{B} = \sum_{i} a_{i} b_{i} \).
The spaces you type in a formula are
ignored.  Remember that a letter like
       $x$                   % $ ... $  and  \( ... \)  are equivalent
is a formula when it denotes a mathematical
symbol, and it should be typed as one.

\section{Displayed Text}

Text is displayed by indenting it from the left
margin.  Quotations are commonly displayed.  There
are short quotations
\begin{quote}
   This is a short quotation.  It consists of a
   single paragraph of text.  See how it is formatted.
\end{quote}
and longer ones.
\begin{quotation}
   This is a longer quotation.  It consists of two
   paragraphs of text, neither of which are
   particularly interesting.

   This is the second paragraph of the quotation.  It
   is just as dull as the first paragraph.
\end{quotation}
Another frequently-displayed structure is a list.
The following is an example of an \emph{itemized}
list.
\begin{itemize}
   \item This is the first item of an itemized list.
         Each item in the list is marked with a ``tick''.
         You don't have to worry about what kind of tick
         mark is used.

   \item This is the second item of the list.  It
         contains another list nested inside it.  The inner
         list is an \emph{enumerated} list.
         \begin{enumerate}
            \item This is the first item of an enumerated
                  list that is nested within the itemized list.

            \item This is the second item of the inner list.
                  \LaTeX\ allows you to nest lists deeper than
                  you really should.
         \end{enumerate}
         This is the rest of the second item of the outer
         list.  It is no more interesting than any other
         part of the item.
   \item This is the third item of the list.
\end{itemize}
You can even display poetry.
\begin{verse}
   There is an environment
    for verse \\             % The \\ command separates lines
   Whose features some poets % within a stanza.
   will curse.

                             % One or more blank lines separate stanzas.

   For instead of making\\
   Them do \emph{all} line breaking, \\
   It allows them to put too many words on a line when they'd rather be
   forced to be terse.
\end{verse}

Mathematical formulas may also be displayed.  A
displayed formula
is
one-line long; multiline
formulas require special formatting instructions.
   \[  \ip{\Gamma}{\psi'} = x'' + y^{2} + z_{i}^{n}\]
Don't start a paragraph with a displayed equation,
nor make one a paragraph by itself.

\end{document}               % End of document.

\section{Bench Huge Section 722}
% synthetic bench section 722
% This is a sample LaTeX input file.  (Version of 12 August 2004.)
%
% A '%' character causes TeX to ignore all remaining text on the line,
% and is used for comments like this one.

\documentclass{article}      % Specifies the document class

                             % The preamble begins here.
\title{An Example Document}  % Declares the document's title.
\author{Leslie Lamport}      % Declares the author's name.
\date{January 21, 1994}      % Deleting this command produces today's date.

\newcommand{\ip}[2]{(#1, #2)}
                             % Defines \ip{arg1}{arg2} to mean
                             % (arg1, arg2).

%\newcommand{\ip}[2]{\langle #1 | #2\rangle}
                             % This is an alternative definition of
                             % \ip that is commented out.

\begin{document}             % End of preamble and beginning of text.

\maketitle                   % Produces the title.

This is an example input file.  Comparing it with
the output it generates can show you how to
produce a simple document of your own.

\section{Ordinary Text}      % Produces section heading.  Lower-level
                             % sections are begun with similar
                             % \subsection and \subsubsection commands.

The ends  of words and sentences are marked
  by   spaces. It  doesn't matter how many
spaces    you type; one is as good as 100.  The
end of   a line counts as a space.

One   or more   blank lines denote the  end
of  a paragraph.

Since any number of consecutive spaces are treated
like a single one, the formatting of the input
file makes no difference to
      \LaTeX,                % The \LaTeX command generates the LaTeX logo.
but it makes a difference to you.  When you use
\LaTeX, making your input file as easy to read
as possible will be a great help as you write
your document and when you change it.  This sample
file shows how you can add comments to your own input
file.

Because printing is different from typewriting,
there are a number of things that you have to do
differently when preparing an input file than if
you were just typing the document directly.
Quotation marks like
       ``this''
have to be handled specially, as do quotes within
quotes:
       ``\,`this'            % \, separates the double and single quote.
        is what I just
        wrote, not  `that'\,''.

Dashes come in three sizes: an
       intra-word
dash, a medium dash for number ranges like
       1--2,
and a punctuation
       dash---like
this.

A sentence-ending space should be larger than the
space between words within a sentence.  You
sometimes have to type special commands in
conjunction with punctuation characters to get
this right, as in the following sentence.
       Gnats, gnus, etc.\ all  % `\ ' makes an inter-word space.
       begin with G\@.         % \@ marks end-of-sentence punctuation.
You should check the spaces after periods when
reading your output to make sure you haven't
forgotten any special cases.  Generating an
ellipsis
       \ldots\               % `\ ' is needed after `\ldots' because TeX
                             % ignores spaces after command names like \ldots
                             % made from \ + letters.
                             %
                             % Note how a `%' character causes TeX to ignore
                             % the end of the input line, so these blank lines
                             % do not start a new paragraph.
                             %
with the right spacing around the periods requires
a special command.

\LaTeX\ interprets some common characters as
commands, so you must type special commands to
generate them.  These characters include the
following:
       \$ \& \% \# \{ and \}.

In printing, text is usually emphasized with an
       \emph{italic}
type style.

\begin{em}
   A long segment of text can also be emphasized
   in this way.  Text within such a segment can be
   given \emph{additional} emphasis.
\end{em}

It is sometimes necessary to prevent \LaTeX\ from
breaking a line where it might otherwise do so.
This may be at a space, as between the ``Mr.''\ and
``Jones'' in
       ``Mr.~Jones'',        % ~ produces an unbreakable interword space.
or within a word---especially when the word is a
symbol like
       \mbox{\emph{itemnum}}
that makes little sense when hyphenated across
lines.

Footnotes\footnote{This is an example of a footnote.}
pose no problem.

\LaTeX\ is good at typesetting mathematical formulas
like
       \( x-3y + z = 7 \)
or
       \( a_{1} > x^{2n} + y^{2n} > x' \)
or
       \( \ip{A}{B} = \sum_{i} a_{i} b_{i} \).
The spaces you type in a formula are
ignored.  Remember that a letter like
       $x$                   % $ ... $  and  \( ... \)  are equivalent
is a formula when it denotes a mathematical
symbol, and it should be typed as one.

\section{Displayed Text}

Text is displayed by indenting it from the left
margin.  Quotations are commonly displayed.  There
are short quotations
\begin{quote}
   This is a short quotation.  It consists of a
   single paragraph of text.  See how it is formatted.
\end{quote}
and longer ones.
\begin{quotation}
   This is a longer quotation.  It consists of two
   paragraphs of text, neither of which are
   particularly interesting.

   This is the second paragraph of the quotation.  It
   is just as dull as the first paragraph.
\end{quotation}
Another frequently-displayed structure is a list.
The following is an example of an \emph{itemized}
list.
\begin{itemize}
   \item This is the first item of an itemized list.
         Each item in the list is marked with a ``tick''.
         You don't have to worry about what kind of tick
         mark is used.

   \item This is the second item of the list.  It
         contains another list nested inside it.  The inner
         list is an \emph{enumerated} list.
         \begin{enumerate}
            \item This is the first item of an enumerated
                  list that is nested within the itemized list.

            \item This is the second item of the inner list.
                  \LaTeX\ allows you to nest lists deeper than
                  you really should.
         \end{enumerate}
         This is the rest of the second item of the outer
         list.  It is no more interesting than any other
         part of the item.
   \item This is the third item of the list.
\end{itemize}
You can even display poetry.
\begin{verse}
   There is an environment
    for verse \\             % The \\ command separates lines
   Whose features some poets % within a stanza.
   will curse.

                             % One or more blank lines separate stanzas.

   For instead of making\\
   Them do \emph{all} line breaking, \\
   It allows them to put too many words on a line when they'd rather be
   forced to be terse.
\end{verse}

Mathematical formulas may also be displayed.  A
displayed formula
is
one-line long; multiline
formulas require special formatting instructions.
   \[  \ip{\Gamma}{\psi'} = x'' + y^{2} + z_{i}^{n}\]
Don't start a paragraph with a displayed equation,
nor make one a paragraph by itself.

\end{document}               % End of document.

\section{Bench Huge Section 723}
% synthetic bench section 723
% This is a sample LaTeX input file.  (Version of 12 August 2004.)
%
% A '%' character causes TeX to ignore all remaining text on the line,
% and is used for comments like this one.

\documentclass{article}      % Specifies the document class

                             % The preamble begins here.
\title{An Example Document}  % Declares the document's title.
\author{Leslie Lamport}      % Declares the author's name.
\date{January 21, 1994}      % Deleting this command produces today's date.

\newcommand{\ip}[2]{(#1, #2)}
                             % Defines \ip{arg1}{arg2} to mean
                             % (arg1, arg2).

%\newcommand{\ip}[2]{\langle #1 | #2\rangle}
                             % This is an alternative definition of
                             % \ip that is commented out.

\begin{document}             % End of preamble and beginning of text.

\maketitle                   % Produces the title.

This is an example input file.  Comparing it with
the output it generates can show you how to
produce a simple document of your own.

\section{Ordinary Text}      % Produces section heading.  Lower-level
                             % sections are begun with similar
                             % \subsection and \subsubsection commands.

The ends  of words and sentences are marked
  by   spaces. It  doesn't matter how many
spaces    you type; one is as good as 100.  The
end of   a line counts as a space.

One   or more   blank lines denote the  end
of  a paragraph.

Since any number of consecutive spaces are treated
like a single one, the formatting of the input
file makes no difference to
      \LaTeX,                % The \LaTeX command generates the LaTeX logo.
but it makes a difference to you.  When you use
\LaTeX, making your input file as easy to read
as possible will be a great help as you write
your document and when you change it.  This sample
file shows how you can add comments to your own input
file.

Because printing is different from typewriting,
there are a number of things that you have to do
differently when preparing an input file than if
you were just typing the document directly.
Quotation marks like
       ``this''
have to be handled specially, as do quotes within
quotes:
       ``\,`this'            % \, separates the double and single quote.
        is what I just
        wrote, not  `that'\,''.

Dashes come in three sizes: an
       intra-word
dash, a medium dash for number ranges like
       1--2,
and a punctuation
       dash---like
this.

A sentence-ending space should be larger than the
space between words within a sentence.  You
sometimes have to type special commands in
conjunction with punctuation characters to get
this right, as in the following sentence.
       Gnats, gnus, etc.\ all  % `\ ' makes an inter-word space.
       begin with G\@.         % \@ marks end-of-sentence punctuation.
You should check the spaces after periods when
reading your output to make sure you haven't
forgotten any special cases.  Generating an
ellipsis
       \ldots\               % `\ ' is needed after `\ldots' because TeX
                             % ignores spaces after command names like \ldots
                             % made from \ + letters.
                             %
                             % Note how a `%' character causes TeX to ignore
                             % the end of the input line, so these blank lines
                             % do not start a new paragraph.
                             %
with the right spacing around the periods requires
a special command.

\LaTeX\ interprets some common characters as
commands, so you must type special commands to
generate them.  These characters include the
following:
       \$ \& \% \# \{ and \}.

In printing, text is usually emphasized with an
       \emph{italic}
type style.

\begin{em}
   A long segment of text can also be emphasized
   in this way.  Text within such a segment can be
   given \emph{additional} emphasis.
\end{em}

It is sometimes necessary to prevent \LaTeX\ from
breaking a line where it might otherwise do so.
This may be at a space, as between the ``Mr.''\ and
``Jones'' in
       ``Mr.~Jones'',        % ~ produces an unbreakable interword space.
or within a word---especially when the word is a
symbol like
       \mbox{\emph{itemnum}}
that makes little sense when hyphenated across
lines.

Footnotes\footnote{This is an example of a footnote.}
pose no problem.

\LaTeX\ is good at typesetting mathematical formulas
like
       \( x-3y + z = 7 \)
or
       \( a_{1} > x^{2n} + y^{2n} > x' \)
or
       \( \ip{A}{B} = \sum_{i} a_{i} b_{i} \).
The spaces you type in a formula are
ignored.  Remember that a letter like
       $x$                   % $ ... $  and  \( ... \)  are equivalent
is a formula when it denotes a mathematical
symbol, and it should be typed as one.

\section{Displayed Text}

Text is displayed by indenting it from the left
margin.  Quotations are commonly displayed.  There
are short quotations
\begin{quote}
   This is a short quotation.  It consists of a
   single paragraph of text.  See how it is formatted.
\end{quote}
and longer ones.
\begin{quotation}
   This is a longer quotation.  It consists of two
   paragraphs of text, neither of which are
   particularly interesting.

   This is the second paragraph of the quotation.  It
   is just as dull as the first paragraph.
\end{quotation}
Another frequently-displayed structure is a list.
The following is an example of an \emph{itemized}
list.
\begin{itemize}
   \item This is the first item of an itemized list.
         Each item in the list is marked with a ``tick''.
         You don't have to worry about what kind of tick
         mark is used.

   \item This is the second item of the list.  It
         contains another list nested inside it.  The inner
         list is an \emph{enumerated} list.
         \begin{enumerate}
            \item This is the first item of an enumerated
                  list that is nested within the itemized list.

            \item This is the second item of the inner list.
                  \LaTeX\ allows you to nest lists deeper than
                  you really should.
         \end{enumerate}
         This is the rest of the second item of the outer
         list.  It is no more interesting than any other
         part of the item.
   \item This is the third item of the list.
\end{itemize}
You can even display poetry.
\begin{verse}
   There is an environment
    for verse \\             % The \\ command separates lines
   Whose features some poets % within a stanza.
   will curse.

                             % One or more blank lines separate stanzas.

   For instead of making\\
   Them do \emph{all} line breaking, \\
   It allows them to put too many words on a line when they'd rather be
   forced to be terse.
\end{verse}

Mathematical formulas may also be displayed.  A
displayed formula
is
one-line long; multiline
formulas require special formatting instructions.
   \[  \ip{\Gamma}{\psi'} = x'' + y^{2} + z_{i}^{n}\]
Don't start a paragraph with a displayed equation,
nor make one a paragraph by itself.

\end{document}               % End of document.

\section{Bench Huge Section 724}
% synthetic bench section 724
% This is a sample LaTeX input file.  (Version of 12 August 2004.)
%
% A '%' character causes TeX to ignore all remaining text on the line,
% and is used for comments like this one.

\documentclass{article}      % Specifies the document class

                             % The preamble begins here.
\title{An Example Document}  % Declares the document's title.
\author{Leslie Lamport}      % Declares the author's name.
\date{January 21, 1994}      % Deleting this command produces today's date.

\newcommand{\ip}[2]{(#1, #2)}
                             % Defines \ip{arg1}{arg2} to mean
                             % (arg1, arg2).

%\newcommand{\ip}[2]{\langle #1 | #2\rangle}
                             % This is an alternative definition of
                             % \ip that is commented out.

\begin{document}             % End of preamble and beginning of text.

\maketitle                   % Produces the title.

This is an example input file.  Comparing it with
the output it generates can show you how to
produce a simple document of your own.

\section{Ordinary Text}      % Produces section heading.  Lower-level
                             % sections are begun with similar
                             % \subsection and \subsubsection commands.

The ends  of words and sentences are marked
  by   spaces. It  doesn't matter how many
spaces    you type; one is as good as 100.  The
end of   a line counts as a space.

One   or more   blank lines denote the  end
of  a paragraph.

Since any number of consecutive spaces are treated
like a single one, the formatting of the input
file makes no difference to
      \LaTeX,                % The \LaTeX command generates the LaTeX logo.
but it makes a difference to you.  When you use
\LaTeX, making your input file as easy to read
as possible will be a great help as you write
your document and when you change it.  This sample
file shows how you can add comments to your own input
file.

Because printing is different from typewriting,
there are a number of things that you have to do
differently when preparing an input file than if
you were just typing the document directly.
Quotation marks like
       ``this''
have to be handled specially, as do quotes within
quotes:
       ``\,`this'            % \, separates the double and single quote.
        is what I just
        wrote, not  `that'\,''.

Dashes come in three sizes: an
       intra-word
dash, a medium dash for number ranges like
       1--2,
and a punctuation
       dash---like
this.

A sentence-ending space should be larger than the
space between words within a sentence.  You
sometimes have to type special commands in
conjunction with punctuation characters to get
this right, as in the following sentence.
       Gnats, gnus, etc.\ all  % `\ ' makes an inter-word space.
       begin with G\@.         % \@ marks end-of-sentence punctuation.
You should check the spaces after periods when
reading your output to make sure you haven't
forgotten any special cases.  Generating an
ellipsis
       \ldots\               % `\ ' is needed after `\ldots' because TeX
                             % ignores spaces after command names like \ldots
                             % made from \ + letters.
                             %
                             % Note how a `%' character causes TeX to ignore
                             % the end of the input line, so these blank lines
                             % do not start a new paragraph.
                             %
with the right spacing around the periods requires
a special command.

\LaTeX\ interprets some common characters as
commands, so you must type special commands to
generate them.  These characters include the
following:
       \$ \& \% \# \{ and \}.

In printing, text is usually emphasized with an
       \emph{italic}
type style.

\begin{em}
   A long segment of text can also be emphasized
   in this way.  Text within such a segment can be
   given \emph{additional} emphasis.
\end{em}

It is sometimes necessary to prevent \LaTeX\ from
breaking a line where it might otherwise do so.
This may be at a space, as between the ``Mr.''\ and
``Jones'' in
       ``Mr.~Jones'',        % ~ produces an unbreakable interword space.
or within a word---especially when the word is a
symbol like
       \mbox{\emph{itemnum}}
that makes little sense when hyphenated across
lines.

Footnotes\footnote{This is an example of a footnote.}
pose no problem.

\LaTeX\ is good at typesetting mathematical formulas
like
       \( x-3y + z = 7 \)
or
       \( a_{1} > x^{2n} + y^{2n} > x' \)
or
       \( \ip{A}{B} = \sum_{i} a_{i} b_{i} \).
The spaces you type in a formula are
ignored.  Remember that a letter like
       $x$                   % $ ... $  and  \( ... \)  are equivalent
is a formula when it denotes a mathematical
symbol, and it should be typed as one.

\section{Displayed Text}

Text is displayed by indenting it from the left
margin.  Quotations are commonly displayed.  There
are short quotations
\begin{quote}
   This is a short quotation.  It consists of a
   single paragraph of text.  See how it is formatted.
\end{quote}
and longer ones.
\begin{quotation}
   This is a longer quotation.  It consists of two
   paragraphs of text, neither of which are
   particularly interesting.

   This is the second paragraph of the quotation.  It
   is just as dull as the first paragraph.
\end{quotation}
Another frequently-displayed structure is a list.
The following is an example of an \emph{itemized}
list.
\begin{itemize}
   \item This is the first item of an itemized list.
         Each item in the list is marked with a ``tick''.
         You don't have to worry about what kind of tick
         mark is used.

   \item This is the second item of the list.  It
         contains another list nested inside it.  The inner
         list is an \emph{enumerated} list.
         \begin{enumerate}
            \item This is the first item of an enumerated
                  list that is nested within the itemized list.

            \item This is the second item of the inner list.
                  \LaTeX\ allows you to nest lists deeper than
                  you really should.
         \end{enumerate}
         This is the rest of the second item of the outer
         list.  It is no more interesting than any other
         part of the item.
   \item This is the third item of the list.
\end{itemize}
You can even display poetry.
\begin{verse}
   There is an environment
    for verse \\             % The \\ command separates lines
   Whose features some poets % within a stanza.
   will curse.

                             % One or more blank lines separate stanzas.

   For instead of making\\
   Them do \emph{all} line breaking, \\
   It allows them to put too many words on a line when they'd rather be
   forced to be terse.
\end{verse}

Mathematical formulas may also be displayed.  A
displayed formula
is
one-line long; multiline
formulas require special formatting instructions.
   \[  \ip{\Gamma}{\psi'} = x'' + y^{2} + z_{i}^{n}\]
Don't start a paragraph with a displayed equation,
nor make one a paragraph by itself.

\end{document}               % End of document.

\section{Bench Huge Section 725}
% synthetic bench section 725
% This is a sample LaTeX input file.  (Version of 12 August 2004.)
%
% A '%' character causes TeX to ignore all remaining text on the line,
% and is used for comments like this one.

\documentclass{article}      % Specifies the document class

                             % The preamble begins here.
\title{An Example Document}  % Declares the document's title.
\author{Leslie Lamport}      % Declares the author's name.
\date{January 21, 1994}      % Deleting this command produces today's date.

\newcommand{\ip}[2]{(#1, #2)}
                             % Defines \ip{arg1}{arg2} to mean
                             % (arg1, arg2).

%\newcommand{\ip}[2]{\langle #1 | #2\rangle}
                             % This is an alternative definition of
                             % \ip that is commented out.

\begin{document}             % End of preamble and beginning of text.

\maketitle                   % Produces the title.

This is an example input file.  Comparing it with
the output it generates can show you how to
produce a simple document of your own.

\section{Ordinary Text}      % Produces section heading.  Lower-level
                             % sections are begun with similar
                             % \subsection and \subsubsection commands.

The ends  of words and sentences are marked
  by   spaces. It  doesn't matter how many
spaces    you type; one is as good as 100.  The
end of   a line counts as a space.

One   or more   blank lines denote the  end
of  a paragraph.

Since any number of consecutive spaces are treated
like a single one, the formatting of the input
file makes no difference to
      \LaTeX,                % The \LaTeX command generates the LaTeX logo.
but it makes a difference to you.  When you use
\LaTeX, making your input file as easy to read
as possible will be a great help as you write
your document and when you change it.  This sample
file shows how you can add comments to your own input
file.

Because printing is different from typewriting,
there are a number of things that you have to do
differently when preparing an input file than if
you were just typing the document directly.
Quotation marks like
       ``this''
have to be handled specially, as do quotes within
quotes:
       ``\,`this'            % \, separates the double and single quote.
        is what I just
        wrote, not  `that'\,''.

Dashes come in three sizes: an
       intra-word
dash, a medium dash for number ranges like
       1--2,
and a punctuation
       dash---like
this.

A sentence-ending space should be larger than the
space between words within a sentence.  You
sometimes have to type special commands in
conjunction with punctuation characters to get
this right, as in the following sentence.
       Gnats, gnus, etc.\ all  % `\ ' makes an inter-word space.
       begin with G\@.         % \@ marks end-of-sentence punctuation.
You should check the spaces after periods when
reading your output to make sure you haven't
forgotten any special cases.  Generating an
ellipsis
       \ldots\               % `\ ' is needed after `\ldots' because TeX
                             % ignores spaces after command names like \ldots
                             % made from \ + letters.
                             %
                             % Note how a `%' character causes TeX to ignore
                             % the end of the input line, so these blank lines
                             % do not start a new paragraph.
                             %
with the right spacing around the periods requires
a special command.

\LaTeX\ interprets some common characters as
commands, so you must type special commands to
generate them.  These characters include the
following:
       \$ \& \% \# \{ and \}.

In printing, text is usually emphasized with an
       \emph{italic}
type style.

\begin{em}
   A long segment of text can also be emphasized
   in this way.  Text within such a segment can be
   given \emph{additional} emphasis.
\end{em}

It is sometimes necessary to prevent \LaTeX\ from
breaking a line where it might otherwise do so.
This may be at a space, as between the ``Mr.''\ and
``Jones'' in
       ``Mr.~Jones'',        % ~ produces an unbreakable interword space.
or within a word---especially when the word is a
symbol like
       \mbox{\emph{itemnum}}
that makes little sense when hyphenated across
lines.

Footnotes\footnote{This is an example of a footnote.}
pose no problem.

\LaTeX\ is good at typesetting mathematical formulas
like
       \( x-3y + z = 7 \)
or
       \( a_{1} > x^{2n} + y^{2n} > x' \)
or
       \( \ip{A}{B} = \sum_{i} a_{i} b_{i} \).
The spaces you type in a formula are
ignored.  Remember that a letter like
       $x$                   % $ ... $  and  \( ... \)  are equivalent
is a formula when it denotes a mathematical
symbol, and it should be typed as one.

\section{Displayed Text}

Text is displayed by indenting it from the left
margin.  Quotations are commonly displayed.  There
are short quotations
\begin{quote}
   This is a short quotation.  It consists of a
   single paragraph of text.  See how it is formatted.
\end{quote}
and longer ones.
\begin{quotation}
   This is a longer quotation.  It consists of two
   paragraphs of text, neither of which are
   particularly interesting.

   This is the second paragraph of the quotation.  It
   is just as dull as the first paragraph.
\end{quotation}
Another frequently-displayed structure is a list.
The following is an example of an \emph{itemized}
list.
\begin{itemize}
   \item This is the first item of an itemized list.
         Each item in the list is marked with a ``tick''.
         You don't have to worry about what kind of tick
         mark is used.

   \item This is the second item of the list.  It
         contains another list nested inside it.  The inner
         list is an \emph{enumerated} list.
         \begin{enumerate}
            \item This is the first item of an enumerated
                  list that is nested within the itemized list.

            \item This is the second item of the inner list.
                  \LaTeX\ allows you to nest lists deeper than
                  you really should.
         \end{enumerate}
         This is the rest of the second item of the outer
         list.  It is no more interesting than any other
         part of the item.
   \item This is the third item of the list.
\end{itemize}
You can even display poetry.
\begin{verse}
   There is an environment
    for verse \\             % The \\ command separates lines
   Whose features some poets % within a stanza.
   will curse.

                             % One or more blank lines separate stanzas.

   For instead of making\\
   Them do \emph{all} line breaking, \\
   It allows them to put too many words on a line when they'd rather be
   forced to be terse.
\end{verse}

Mathematical formulas may also be displayed.  A
displayed formula
is
one-line long; multiline
formulas require special formatting instructions.
   \[  \ip{\Gamma}{\psi'} = x'' + y^{2} + z_{i}^{n}\]
Don't start a paragraph with a displayed equation,
nor make one a paragraph by itself.

\end{document}               % End of document.

\section{Bench Huge Section 726}
% synthetic bench section 726
% This is a sample LaTeX input file.  (Version of 12 August 2004.)
%
% A '%' character causes TeX to ignore all remaining text on the line,
% and is used for comments like this one.

\documentclass{article}      % Specifies the document class

                             % The preamble begins here.
\title{An Example Document}  % Declares the document's title.
\author{Leslie Lamport}      % Declares the author's name.
\date{January 21, 1994}      % Deleting this command produces today's date.

\newcommand{\ip}[2]{(#1, #2)}
                             % Defines \ip{arg1}{arg2} to mean
                             % (arg1, arg2).

%\newcommand{\ip}[2]{\langle #1 | #2\rangle}
                             % This is an alternative definition of
                             % \ip that is commented out.

\begin{document}             % End of preamble and beginning of text.

\maketitle                   % Produces the title.

This is an example input file.  Comparing it with
the output it generates can show you how to
produce a simple document of your own.

\section{Ordinary Text}      % Produces section heading.  Lower-level
                             % sections are begun with similar
                             % \subsection and \subsubsection commands.

The ends  of words and sentences are marked
  by   spaces. It  doesn't matter how many
spaces    you type; one is as good as 100.  The
end of   a line counts as a space.

One   or more   blank lines denote the  end
of  a paragraph.

Since any number of consecutive spaces are treated
like a single one, the formatting of the input
file makes no difference to
      \LaTeX,                % The \LaTeX command generates the LaTeX logo.
but it makes a difference to you.  When you use
\LaTeX, making your input file as easy to read
as possible will be a great help as you write
your document and when you change it.  This sample
file shows how you can add comments to your own input
file.

Because printing is different from typewriting,
there are a number of things that you have to do
differently when preparing an input file than if
you were just typing the document directly.
Quotation marks like
       ``this''
have to be handled specially, as do quotes within
quotes:
       ``\,`this'            % \, separates the double and single quote.
        is what I just
        wrote, not  `that'\,''.

Dashes come in three sizes: an
       intra-word
dash, a medium dash for number ranges like
       1--2,
and a punctuation
       dash---like
this.

A sentence-ending space should be larger than the
space between words within a sentence.  You
sometimes have to type special commands in
conjunction with punctuation characters to get
this right, as in the following sentence.
       Gnats, gnus, etc.\ all  % `\ ' makes an inter-word space.
       begin with G\@.         % \@ marks end-of-sentence punctuation.
You should check the spaces after periods when
reading your output to make sure you haven't
forgotten any special cases.  Generating an
ellipsis
       \ldots\               % `\ ' is needed after `\ldots' because TeX
                             % ignores spaces after command names like \ldots
                             % made from \ + letters.
                             %
                             % Note how a `%' character causes TeX to ignore
                             % the end of the input line, so these blank lines
                             % do not start a new paragraph.
                             %
with the right spacing around the periods requires
a special command.

\LaTeX\ interprets some common characters as
commands, so you must type special commands to
generate them.  These characters include the
following:
       \$ \& \% \# \{ and \}.

In printing, text is usually emphasized with an
       \emph{italic}
type style.

\begin{em}
   A long segment of text can also be emphasized
   in this way.  Text within such a segment can be
   given \emph{additional} emphasis.
\end{em}

It is sometimes necessary to prevent \LaTeX\ from
breaking a line where it might otherwise do so.
This may be at a space, as between the ``Mr.''\ and
``Jones'' in
       ``Mr.~Jones'',        % ~ produces an unbreakable interword space.
or within a word---especially when the word is a
symbol like
       \mbox{\emph{itemnum}}
that makes little sense when hyphenated across
lines.

Footnotes\footnote{This is an example of a footnote.}
pose no problem.

\LaTeX\ is good at typesetting mathematical formulas
like
       \( x-3y + z = 7 \)
or
       \( a_{1} > x^{2n} + y^{2n} > x' \)
or
       \( \ip{A}{B} = \sum_{i} a_{i} b_{i} \).
The spaces you type in a formula are
ignored.  Remember that a letter like
       $x$                   % $ ... $  and  \( ... \)  are equivalent
is a formula when it denotes a mathematical
symbol, and it should be typed as one.

\section{Displayed Text}

Text is displayed by indenting it from the left
margin.  Quotations are commonly displayed.  There
are short quotations
\begin{quote}
   This is a short quotation.  It consists of a
   single paragraph of text.  See how it is formatted.
\end{quote}
and longer ones.
\begin{quotation}
   This is a longer quotation.  It consists of two
   paragraphs of text, neither of which are
   particularly interesting.

   This is the second paragraph of the quotation.  It
   is just as dull as the first paragraph.
\end{quotation}
Another frequently-displayed structure is a list.
The following is an example of an \emph{itemized}
list.
\begin{itemize}
   \item This is the first item of an itemized list.
         Each item in the list is marked with a ``tick''.
         You don't have to worry about what kind of tick
         mark is used.

   \item This is the second item of the list.  It
         contains another list nested inside it.  The inner
         list is an \emph{enumerated} list.
         \begin{enumerate}
            \item This is the first item of an enumerated
                  list that is nested within the itemized list.

            \item This is the second item of the inner list.
                  \LaTeX\ allows you to nest lists deeper than
                  you really should.
         \end{enumerate}
         This is the rest of the second item of the outer
         list.  It is no more interesting than any other
         part of the item.
   \item This is the third item of the list.
\end{itemize}
You can even display poetry.
\begin{verse}
   There is an environment
    for verse \\             % The \\ command separates lines
   Whose features some poets % within a stanza.
   will curse.

                             % One or more blank lines separate stanzas.

   For instead of making\\
   Them do \emph{all} line breaking, \\
   It allows them to put too many words on a line when they'd rather be
   forced to be terse.
\end{verse}

Mathematical formulas may also be displayed.  A
displayed formula
is
one-line long; multiline
formulas require special formatting instructions.
   \[  \ip{\Gamma}{\psi'} = x'' + y^{2} + z_{i}^{n}\]
Don't start a paragraph with a displayed equation,
nor make one a paragraph by itself.

\end{document}               % End of document.

\section{Bench Huge Section 727}
% synthetic bench section 727
% This is a sample LaTeX input file.  (Version of 12 August 2004.)
%
% A '%' character causes TeX to ignore all remaining text on the line,
% and is used for comments like this one.

\documentclass{article}      % Specifies the document class

                             % The preamble begins here.
\title{An Example Document}  % Declares the document's title.
\author{Leslie Lamport}      % Declares the author's name.
\date{January 21, 1994}      % Deleting this command produces today's date.

\newcommand{\ip}[2]{(#1, #2)}
                             % Defines \ip{arg1}{arg2} to mean
                             % (arg1, arg2).

%\newcommand{\ip}[2]{\langle #1 | #2\rangle}
                             % This is an alternative definition of
                             % \ip that is commented out.

\begin{document}             % End of preamble and beginning of text.

\maketitle                   % Produces the title.

This is an example input file.  Comparing it with
the output it generates can show you how to
produce a simple document of your own.

\section{Ordinary Text}      % Produces section heading.  Lower-level
                             % sections are begun with similar
                             % \subsection and \subsubsection commands.

The ends  of words and sentences are marked
  by   spaces. It  doesn't matter how many
spaces    you type; one is as good as 100.  The
end of   a line counts as a space.

One   or more   blank lines denote the  end
of  a paragraph.

Since any number of consecutive spaces are treated
like a single one, the formatting of the input
file makes no difference to
      \LaTeX,                % The \LaTeX command generates the LaTeX logo.
but it makes a difference to you.  When you use
\LaTeX, making your input file as easy to read
as possible will be a great help as you write
your document and when you change it.  This sample
file shows how you can add comments to your own input
file.

Because printing is different from typewriting,
there are a number of things that you have to do
differently when preparing an input file than if
you were just typing the document directly.
Quotation marks like
       ``this''
have to be handled specially, as do quotes within
quotes:
       ``\,`this'            % \, separates the double and single quote.
        is what I just
        wrote, not  `that'\,''.

Dashes come in three sizes: an
       intra-word
dash, a medium dash for number ranges like
       1--2,
and a punctuation
       dash---like
this.

A sentence-ending space should be larger than the
space between words within a sentence.  You
sometimes have to type special commands in
conjunction with punctuation characters to get
this right, as in the following sentence.
       Gnats, gnus, etc.\ all  % `\ ' makes an inter-word space.
       begin with G\@.         % \@ marks end-of-sentence punctuation.
You should check the spaces after periods when
reading your output to make sure you haven't
forgotten any special cases.  Generating an
ellipsis
       \ldots\               % `\ ' is needed after `\ldots' because TeX
                             % ignores spaces after command names like \ldots
                             % made from \ + letters.
                             %
                             % Note how a `%' character causes TeX to ignore
                             % the end of the input line, so these blank lines
                             % do not start a new paragraph.
                             %
with the right spacing around the periods requires
a special command.

\LaTeX\ interprets some common characters as
commands, so you must type special commands to
generate them.  These characters include the
following:
       \$ \& \% \# \{ and \}.

In printing, text is usually emphasized with an
       \emph{italic}
type style.

\begin{em}
   A long segment of text can also be emphasized
   in this way.  Text within such a segment can be
   given \emph{additional} emphasis.
\end{em}

It is sometimes necessary to prevent \LaTeX\ from
breaking a line where it might otherwise do so.
This may be at a space, as between the ``Mr.''\ and
``Jones'' in
       ``Mr.~Jones'',        % ~ produces an unbreakable interword space.
or within a word---especially when the word is a
symbol like
       \mbox{\emph{itemnum}}
that makes little sense when hyphenated across
lines.

Footnotes\footnote{This is an example of a footnote.}
pose no problem.

\LaTeX\ is good at typesetting mathematical formulas
like
       \( x-3y + z = 7 \)
or
       \( a_{1} > x^{2n} + y^{2n} > x' \)
or
       \( \ip{A}{B} = \sum_{i} a_{i} b_{i} \).
The spaces you type in a formula are
ignored.  Remember that a letter like
       $x$                   % $ ... $  and  \( ... \)  are equivalent
is a formula when it denotes a mathematical
symbol, and it should be typed as one.

\section{Displayed Text}

Text is displayed by indenting it from the left
margin.  Quotations are commonly displayed.  There
are short quotations
\begin{quote}
   This is a short quotation.  It consists of a
   single paragraph of text.  See how it is formatted.
\end{quote}
and longer ones.
\begin{quotation}
   This is a longer quotation.  It consists of two
   paragraphs of text, neither of which are
   particularly interesting.

   This is the second paragraph of the quotation.  It
   is just as dull as the first paragraph.
\end{quotation}
Another frequently-displayed structure is a list.
The following is an example of an \emph{itemized}
list.
\begin{itemize}
   \item This is the first item of an itemized list.
         Each item in the list is marked with a ``tick''.
         You don't have to worry about what kind of tick
         mark is used.

   \item This is the second item of the list.  It
         contains another list nested inside it.  The inner
         list is an \emph{enumerated} list.
         \begin{enumerate}
            \item This is the first item of an enumerated
                  list that is nested within the itemized list.

            \item This is the second item of the inner list.
                  \LaTeX\ allows you to nest lists deeper than
                  you really should.
         \end{enumerate}
         This is the rest of the second item of the outer
         list.  It is no more interesting than any other
         part of the item.
   \item This is the third item of the list.
\end{itemize}
You can even display poetry.
\begin{verse}
   There is an environment
    for verse \\             % The \\ command separates lines
   Whose features some poets % within a stanza.
   will curse.

                             % One or more blank lines separate stanzas.

   For instead of making\\
   Them do \emph{all} line breaking, \\
   It allows them to put too many words on a line when they'd rather be
   forced to be terse.
\end{verse}

Mathematical formulas may also be displayed.  A
displayed formula
is
one-line long; multiline
formulas require special formatting instructions.
   \[  \ip{\Gamma}{\psi'} = x'' + y^{2} + z_{i}^{n}\]
Don't start a paragraph with a displayed equation,
nor make one a paragraph by itself.

\end{document}               % End of document.

\section{Bench Huge Section 728}
% synthetic bench section 728
% This is a sample LaTeX input file.  (Version of 12 August 2004.)
%
% A '%' character causes TeX to ignore all remaining text on the line,
% and is used for comments like this one.

\documentclass{article}      % Specifies the document class

                             % The preamble begins here.
\title{An Example Document}  % Declares the document's title.
\author{Leslie Lamport}      % Declares the author's name.
\date{January 21, 1994}      % Deleting this command produces today's date.

\newcommand{\ip}[2]{(#1, #2)}
                             % Defines \ip{arg1}{arg2} to mean
                             % (arg1, arg2).

%\newcommand{\ip}[2]{\langle #1 | #2\rangle}
                             % This is an alternative definition of
                             % \ip that is commented out.

\begin{document}             % End of preamble and beginning of text.

\maketitle                   % Produces the title.

This is an example input file.  Comparing it with
the output it generates can show you how to
produce a simple document of your own.

\section{Ordinary Text}      % Produces section heading.  Lower-level
                             % sections are begun with similar
                             % \subsection and \subsubsection commands.

The ends  of words and sentences are marked
  by   spaces. It  doesn't matter how many
spaces    you type; one is as good as 100.  The
end of   a line counts as a space.

One   or more   blank lines denote the  end
of  a paragraph.

Since any number of consecutive spaces are treated
like a single one, the formatting of the input
file makes no difference to
      \LaTeX,                % The \LaTeX command generates the LaTeX logo.
but it makes a difference to you.  When you use
\LaTeX, making your input file as easy to read
as possible will be a great help as you write
your document and when you change it.  This sample
file shows how you can add comments to your own input
file.

Because printing is different from typewriting,
there are a number of things that you have to do
differently when preparing an input file than if
you were just typing the document directly.
Quotation marks like
       ``this''
have to be handled specially, as do quotes within
quotes:
       ``\,`this'            % \, separates the double and single quote.
        is what I just
        wrote, not  `that'\,''.

Dashes come in three sizes: an
       intra-word
dash, a medium dash for number ranges like
       1--2,
and a punctuation
       dash---like
this.

A sentence-ending space should be larger than the
space between words within a sentence.  You
sometimes have to type special commands in
conjunction with punctuation characters to get
this right, as in the following sentence.
       Gnats, gnus, etc.\ all  % `\ ' makes an inter-word space.
       begin with G\@.         % \@ marks end-of-sentence punctuation.
You should check the spaces after periods when
reading your output to make sure you haven't
forgotten any special cases.  Generating an
ellipsis
       \ldots\               % `\ ' is needed after `\ldots' because TeX
                             % ignores spaces after command names like \ldots
                             % made from \ + letters.
                             %
                             % Note how a `%' character causes TeX to ignore
                             % the end of the input line, so these blank lines
                             % do not start a new paragraph.
                             %
with the right spacing around the periods requires
a special command.

\LaTeX\ interprets some common characters as
commands, so you must type special commands to
generate them.  These characters include the
following:
       \$ \& \% \# \{ and \}.

In printing, text is usually emphasized with an
       \emph{italic}
type style.

\begin{em}
   A long segment of text can also be emphasized
   in this way.  Text within such a segment can be
   given \emph{additional} emphasis.
\end{em}

It is sometimes necessary to prevent \LaTeX\ from
breaking a line where it might otherwise do so.
This may be at a space, as between the ``Mr.''\ and
``Jones'' in
       ``Mr.~Jones'',        % ~ produces an unbreakable interword space.
or within a word---especially when the word is a
symbol like
       \mbox{\emph{itemnum}}
that makes little sense when hyphenated across
lines.

Footnotes\footnote{This is an example of a footnote.}
pose no problem.

\LaTeX\ is good at typesetting mathematical formulas
like
       \( x-3y + z = 7 \)
or
       \( a_{1} > x^{2n} + y^{2n} > x' \)
or
       \( \ip{A}{B} = \sum_{i} a_{i} b_{i} \).
The spaces you type in a formula are
ignored.  Remember that a letter like
       $x$                   % $ ... $  and  \( ... \)  are equivalent
is a formula when it denotes a mathematical
symbol, and it should be typed as one.

\section{Displayed Text}

Text is displayed by indenting it from the left
margin.  Quotations are commonly displayed.  There
are short quotations
\begin{quote}
   This is a short quotation.  It consists of a
   single paragraph of text.  See how it is formatted.
\end{quote}
and longer ones.
\begin{quotation}
   This is a longer quotation.  It consists of two
   paragraphs of text, neither of which are
   particularly interesting.

   This is the second paragraph of the quotation.  It
   is just as dull as the first paragraph.
\end{quotation}
Another frequently-displayed structure is a list.
The following is an example of an \emph{itemized}
list.
\begin{itemize}
   \item This is the first item of an itemized list.
         Each item in the list is marked with a ``tick''.
         You don't have to worry about what kind of tick
         mark is used.

   \item This is the second item of the list.  It
         contains another list nested inside it.  The inner
         list is an \emph{enumerated} list.
         \begin{enumerate}
            \item This is the first item of an enumerated
                  list that is nested within the itemized list.

            \item This is the second item of the inner list.
                  \LaTeX\ allows you to nest lists deeper than
                  you really should.
         \end{enumerate}
         This is the rest of the second item of the outer
         list.  It is no more interesting than any other
         part of the item.
   \item This is the third item of the list.
\end{itemize}
You can even display poetry.
\begin{verse}
   There is an environment
    for verse \\             % The \\ command separates lines
   Whose features some poets % within a stanza.
   will curse.

                             % One or more blank lines separate stanzas.

   For instead of making\\
   Them do \emph{all} line breaking, \\
   It allows them to put too many words on a line when they'd rather be
   forced to be terse.
\end{verse}

Mathematical formulas may also be displayed.  A
displayed formula
is
one-line long; multiline
formulas require special formatting instructions.
   \[  \ip{\Gamma}{\psi'} = x'' + y^{2} + z_{i}^{n}\]
Don't start a paragraph with a displayed equation,
nor make one a paragraph by itself.

\end{document}               % End of document.

\section{Bench Huge Section 729}
% synthetic bench section 729
% This is a sample LaTeX input file.  (Version of 12 August 2004.)
%
% A '%' character causes TeX to ignore all remaining text on the line,
% and is used for comments like this one.

\documentclass{article}      % Specifies the document class

                             % The preamble begins here.
\title{An Example Document}  % Declares the document's title.
\author{Leslie Lamport}      % Declares the author's name.
\date{January 21, 1994}      % Deleting this command produces today's date.

\newcommand{\ip}[2]{(#1, #2)}
                             % Defines \ip{arg1}{arg2} to mean
                             % (arg1, arg2).

%\newcommand{\ip}[2]{\langle #1 | #2\rangle}
                             % This is an alternative definition of
                             % \ip that is commented out.

\begin{document}             % End of preamble and beginning of text.

\maketitle                   % Produces the title.

This is an example input file.  Comparing it with
the output it generates can show you how to
produce a simple document of your own.

\section{Ordinary Text}      % Produces section heading.  Lower-level
                             % sections are begun with similar
                             % \subsection and \subsubsection commands.

The ends  of words and sentences are marked
  by   spaces. It  doesn't matter how many
spaces    you type; one is as good as 100.  The
end of   a line counts as a space.

One   or more   blank lines denote the  end
of  a paragraph.

Since any number of consecutive spaces are treated
like a single one, the formatting of the input
file makes no difference to
      \LaTeX,                % The \LaTeX command generates the LaTeX logo.
but it makes a difference to you.  When you use
\LaTeX, making your input file as easy to read
as possible will be a great help as you write
your document and when you change it.  This sample
file shows how you can add comments to your own input
file.

Because printing is different from typewriting,
there are a number of things that you have to do
differently when preparing an input file than if
you were just typing the document directly.
Quotation marks like
       ``this''
have to be handled specially, as do quotes within
quotes:
       ``\,`this'            % \, separates the double and single quote.
        is what I just
        wrote, not  `that'\,''.

Dashes come in three sizes: an
       intra-word
dash, a medium dash for number ranges like
       1--2,
and a punctuation
       dash---like
this.

A sentence-ending space should be larger than the
space between words within a sentence.  You
sometimes have to type special commands in
conjunction with punctuation characters to get
this right, as in the following sentence.
       Gnats, gnus, etc.\ all  % `\ ' makes an inter-word space.
       begin with G\@.         % \@ marks end-of-sentence punctuation.
You should check the spaces after periods when
reading your output to make sure you haven't
forgotten any special cases.  Generating an
ellipsis
       \ldots\               % `\ ' is needed after `\ldots' because TeX
                             % ignores spaces after command names like \ldots
                             % made from \ + letters.
                             %
                             % Note how a `%' character causes TeX to ignore
                             % the end of the input line, so these blank lines
                             % do not start a new paragraph.
                             %
with the right spacing around the periods requires
a special command.

\LaTeX\ interprets some common characters as
commands, so you must type special commands to
generate them.  These characters include the
following:
       \$ \& \% \# \{ and \}.

In printing, text is usually emphasized with an
       \emph{italic}
type style.

\begin{em}
   A long segment of text can also be emphasized
   in this way.  Text within such a segment can be
   given \emph{additional} emphasis.
\end{em}

It is sometimes necessary to prevent \LaTeX\ from
breaking a line where it might otherwise do so.
This may be at a space, as between the ``Mr.''\ and
``Jones'' in
       ``Mr.~Jones'',        % ~ produces an unbreakable interword space.
or within a word---especially when the word is a
symbol like
       \mbox{\emph{itemnum}}
that makes little sense when hyphenated across
lines.

Footnotes\footnote{This is an example of a footnote.}
pose no problem.

\LaTeX\ is good at typesetting mathematical formulas
like
       \( x-3y + z = 7 \)
or
       \( a_{1} > x^{2n} + y^{2n} > x' \)
or
       \( \ip{A}{B} = \sum_{i} a_{i} b_{i} \).
The spaces you type in a formula are
ignored.  Remember that a letter like
       $x$                   % $ ... $  and  \( ... \)  are equivalent
is a formula when it denotes a mathematical
symbol, and it should be typed as one.

\section{Displayed Text}

Text is displayed by indenting it from the left
margin.  Quotations are commonly displayed.  There
are short quotations
\begin{quote}
   This is a short quotation.  It consists of a
   single paragraph of text.  See how it is formatted.
\end{quote}
and longer ones.
\begin{quotation}
   This is a longer quotation.  It consists of two
   paragraphs of text, neither of which are
   particularly interesting.

   This is the second paragraph of the quotation.  It
   is just as dull as the first paragraph.
\end{quotation}
Another frequently-displayed structure is a list.
The following is an example of an \emph{itemized}
list.
\begin{itemize}
   \item This is the first item of an itemized list.
         Each item in the list is marked with a ``tick''.
         You don't have to worry about what kind of tick
         mark is used.

   \item This is the second item of the list.  It
         contains another list nested inside it.  The inner
         list is an \emph{enumerated} list.
         \begin{enumerate}
            \item This is the first item of an enumerated
                  list that is nested within the itemized list.

            \item This is the second item of the inner list.
                  \LaTeX\ allows you to nest lists deeper than
                  you really should.
         \end{enumerate}
         This is the rest of the second item of the outer
         list.  It is no more interesting than any other
         part of the item.
   \item This is the third item of the list.
\end{itemize}
You can even display poetry.
\begin{verse}
   There is an environment
    for verse \\             % The \\ command separates lines
   Whose features some poets % within a stanza.
   will curse.

                             % One or more blank lines separate stanzas.

   For instead of making\\
   Them do \emph{all} line breaking, \\
   It allows them to put too many words on a line when they'd rather be
   forced to be terse.
\end{verse}

Mathematical formulas may also be displayed.  A
displayed formula
is
one-line long; multiline
formulas require special formatting instructions.
   \[  \ip{\Gamma}{\psi'} = x'' + y^{2} + z_{i}^{n}\]
Don't start a paragraph with a displayed equation,
nor make one a paragraph by itself.

\end{document}               % End of document.

\section{Bench Huge Section 730}
% synthetic bench section 730
% This is a sample LaTeX input file.  (Version of 12 August 2004.)
%
% A '%' character causes TeX to ignore all remaining text on the line,
% and is used for comments like this one.

\documentclass{article}      % Specifies the document class

                             % The preamble begins here.
\title{An Example Document}  % Declares the document's title.
\author{Leslie Lamport}      % Declares the author's name.
\date{January 21, 1994}      % Deleting this command produces today's date.

\newcommand{\ip}[2]{(#1, #2)}
                             % Defines \ip{arg1}{arg2} to mean
                             % (arg1, arg2).

%\newcommand{\ip}[2]{\langle #1 | #2\rangle}
                             % This is an alternative definition of
                             % \ip that is commented out.

\begin{document}             % End of preamble and beginning of text.

\maketitle                   % Produces the title.

This is an example input file.  Comparing it with
the output it generates can show you how to
produce a simple document of your own.

\section{Ordinary Text}      % Produces section heading.  Lower-level
                             % sections are begun with similar
                             % \subsection and \subsubsection commands.

The ends  of words and sentences are marked
  by   spaces. It  doesn't matter how many
spaces    you type; one is as good as 100.  The
end of   a line counts as a space.

One   or more   blank lines denote the  end
of  a paragraph.

Since any number of consecutive spaces are treated
like a single one, the formatting of the input
file makes no difference to
      \LaTeX,                % The \LaTeX command generates the LaTeX logo.
but it makes a difference to you.  When you use
\LaTeX, making your input file as easy to read
as possible will be a great help as you write
your document and when you change it.  This sample
file shows how you can add comments to your own input
file.

Because printing is different from typewriting,
there are a number of things that you have to do
differently when preparing an input file than if
you were just typing the document directly.
Quotation marks like
       ``this''
have to be handled specially, as do quotes within
quotes:
       ``\,`this'            % \, separates the double and single quote.
        is what I just
        wrote, not  `that'\,''.

Dashes come in three sizes: an
       intra-word
dash, a medium dash for number ranges like
       1--2,
and a punctuation
       dash---like
this.

A sentence-ending space should be larger than the
space between words within a sentence.  You
sometimes have to type special commands in
conjunction with punctuation characters to get
this right, as in the following sentence.
       Gnats, gnus, etc.\ all  % `\ ' makes an inter-word space.
       begin with G\@.         % \@ marks end-of-sentence punctuation.
You should check the spaces after periods when
reading your output to make sure you haven't
forgotten any special cases.  Generating an
ellipsis
       \ldots\               % `\ ' is needed after `\ldots' because TeX
                             % ignores spaces after command names like \ldots
                             % made from \ + letters.
                             %
                             % Note how a `%' character causes TeX to ignore
                             % the end of the input line, so these blank lines
                             % do not start a new paragraph.
                             %
with the right spacing around the periods requires
a special command.

\LaTeX\ interprets some common characters as
commands, so you must type special commands to
generate them.  These characters include the
following:
       \$ \& \% \# \{ and \}.

In printing, text is usually emphasized with an
       \emph{italic}
type style.

\begin{em}
   A long segment of text can also be emphasized
   in this way.  Text within such a segment can be
   given \emph{additional} emphasis.
\end{em}

It is sometimes necessary to prevent \LaTeX\ from
breaking a line where it might otherwise do so.
This may be at a space, as between the ``Mr.''\ and
``Jones'' in
       ``Mr.~Jones'',        % ~ produces an unbreakable interword space.
or within a word---especially when the word is a
symbol like
       \mbox{\emph{itemnum}}
that makes little sense when hyphenated across
lines.

Footnotes\footnote{This is an example of a footnote.}
pose no problem.

\LaTeX\ is good at typesetting mathematical formulas
like
       \( x-3y + z = 7 \)
or
       \( a_{1} > x^{2n} + y^{2n} > x' \)
or
       \( \ip{A}{B} = \sum_{i} a_{i} b_{i} \).
The spaces you type in a formula are
ignored.  Remember that a letter like
       $x$                   % $ ... $  and  \( ... \)  are equivalent
is a formula when it denotes a mathematical
symbol, and it should be typed as one.

\section{Displayed Text}

Text is displayed by indenting it from the left
margin.  Quotations are commonly displayed.  There
are short quotations
\begin{quote}
   This is a short quotation.  It consists of a
   single paragraph of text.  See how it is formatted.
\end{quote}
and longer ones.
\begin{quotation}
   This is a longer quotation.  It consists of two
   paragraphs of text, neither of which are
   particularly interesting.

   This is the second paragraph of the quotation.  It
   is just as dull as the first paragraph.
\end{quotation}
Another frequently-displayed structure is a list.
The following is an example of an \emph{itemized}
list.
\begin{itemize}
   \item This is the first item of an itemized list.
         Each item in the list is marked with a ``tick''.
         You don't have to worry about what kind of tick
         mark is used.

   \item This is the second item of the list.  It
         contains another list nested inside it.  The inner
         list is an \emph{enumerated} list.
         \begin{enumerate}
            \item This is the first item of an enumerated
                  list that is nested within the itemized list.

            \item This is the second item of the inner list.
                  \LaTeX\ allows you to nest lists deeper than
                  you really should.
         \end{enumerate}
         This is the rest of the second item of the outer
         list.  It is no more interesting than any other
         part of the item.
   \item This is the third item of the list.
\end{itemize}
You can even display poetry.
\begin{verse}
   There is an environment
    for verse \\             % The \\ command separates lines
   Whose features some poets % within a stanza.
   will curse.

                             % One or more blank lines separate stanzas.

   For instead of making\\
   Them do \emph{all} line breaking, \\
   It allows them to put too many words on a line when they'd rather be
   forced to be terse.
\end{verse}

Mathematical formulas may also be displayed.  A
displayed formula
is
one-line long; multiline
formulas require special formatting instructions.
   \[  \ip{\Gamma}{\psi'} = x'' + y^{2} + z_{i}^{n}\]
Don't start a paragraph with a displayed equation,
nor make one a paragraph by itself.

\end{document}               % End of document.

\section{Bench Huge Section 731}
% synthetic bench section 731
% This is a sample LaTeX input file.  (Version of 12 August 2004.)
%
% A '%' character causes TeX to ignore all remaining text on the line,
% and is used for comments like this one.

\documentclass{article}      % Specifies the document class

                             % The preamble begins here.
\title{An Example Document}  % Declares the document's title.
\author{Leslie Lamport}      % Declares the author's name.
\date{January 21, 1994}      % Deleting this command produces today's date.

\newcommand{\ip}[2]{(#1, #2)}
                             % Defines \ip{arg1}{arg2} to mean
                             % (arg1, arg2).

%\newcommand{\ip}[2]{\langle #1 | #2\rangle}
                             % This is an alternative definition of
                             % \ip that is commented out.

\begin{document}             % End of preamble and beginning of text.

\maketitle                   % Produces the title.

This is an example input file.  Comparing it with
the output it generates can show you how to
produce a simple document of your own.

\section{Ordinary Text}      % Produces section heading.  Lower-level
                             % sections are begun with similar
                             % \subsection and \subsubsection commands.

The ends  of words and sentences are marked
  by   spaces. It  doesn't matter how many
spaces    you type; one is as good as 100.  The
end of   a line counts as a space.

One   or more   blank lines denote the  end
of  a paragraph.

Since any number of consecutive spaces are treated
like a single one, the formatting of the input
file makes no difference to
      \LaTeX,                % The \LaTeX command generates the LaTeX logo.
but it makes a difference to you.  When you use
\LaTeX, making your input file as easy to read
as possible will be a great help as you write
your document and when you change it.  This sample
file shows how you can add comments to your own input
file.

Because printing is different from typewriting,
there are a number of things that you have to do
differently when preparing an input file than if
you were just typing the document directly.
Quotation marks like
       ``this''
have to be handled specially, as do quotes within
quotes:
       ``\,`this'            % \, separates the double and single quote.
        is what I just
        wrote, not  `that'\,''.

Dashes come in three sizes: an
       intra-word
dash, a medium dash for number ranges like
       1--2,
and a punctuation
       dash---like
this.

A sentence-ending space should be larger than the
space between words within a sentence.  You
sometimes have to type special commands in
conjunction with punctuation characters to get
this right, as in the following sentence.
       Gnats, gnus, etc.\ all  % `\ ' makes an inter-word space.
       begin with G\@.         % \@ marks end-of-sentence punctuation.
You should check the spaces after periods when
reading your output to make sure you haven't
forgotten any special cases.  Generating an
ellipsis
       \ldots\               % `\ ' is needed after `\ldots' because TeX
                             % ignores spaces after command names like \ldots
                             % made from \ + letters.
                             %
                             % Note how a `%' character causes TeX to ignore
                             % the end of the input line, so these blank lines
                             % do not start a new paragraph.
                             %
with the right spacing around the periods requires
a special command.

\LaTeX\ interprets some common characters as
commands, so you must type special commands to
generate them.  These characters include the
following:
       \$ \& \% \# \{ and \}.

In printing, text is usually emphasized with an
       \emph{italic}
type style.

\begin{em}
   A long segment of text can also be emphasized
   in this way.  Text within such a segment can be
   given \emph{additional} emphasis.
\end{em}

It is sometimes necessary to prevent \LaTeX\ from
breaking a line where it might otherwise do so.
This may be at a space, as between the ``Mr.''\ and
``Jones'' in
       ``Mr.~Jones'',        % ~ produces an unbreakable interword space.
or within a word---especially when the word is a
symbol like
       \mbox{\emph{itemnum}}
that makes little sense when hyphenated across
lines.

Footnotes\footnote{This is an example of a footnote.}
pose no problem.

\LaTeX\ is good at typesetting mathematical formulas
like
       \( x-3y + z = 7 \)
or
       \( a_{1} > x^{2n} + y^{2n} > x' \)
or
       \( \ip{A}{B} = \sum_{i} a_{i} b_{i} \).
The spaces you type in a formula are
ignored.  Remember that a letter like
       $x$                   % $ ... $  and  \( ... \)  are equivalent
is a formula when it denotes a mathematical
symbol, and it should be typed as one.

\section{Displayed Text}

Text is displayed by indenting it from the left
margin.  Quotations are commonly displayed.  There
are short quotations
\begin{quote}
   This is a short quotation.  It consists of a
   single paragraph of text.  See how it is formatted.
\end{quote}
and longer ones.
\begin{quotation}
   This is a longer quotation.  It consists of two
   paragraphs of text, neither of which are
   particularly interesting.

   This is the second paragraph of the quotation.  It
   is just as dull as the first paragraph.
\end{quotation}
Another frequently-displayed structure is a list.
The following is an example of an \emph{itemized}
list.
\begin{itemize}
   \item This is the first item of an itemized list.
         Each item in the list is marked with a ``tick''.
         You don't have to worry about what kind of tick
         mark is used.

   \item This is the second item of the list.  It
         contains another list nested inside it.  The inner
         list is an \emph{enumerated} list.
         \begin{enumerate}
            \item This is the first item of an enumerated
                  list that is nested within the itemized list.

            \item This is the second item of the inner list.
                  \LaTeX\ allows you to nest lists deeper than
                  you really should.
         \end{enumerate}
         This is the rest of the second item of the outer
         list.  It is no more interesting than any other
         part of the item.
   \item This is the third item of the list.
\end{itemize}
You can even display poetry.
\begin{verse}
   There is an environment
    for verse \\             % The \\ command separates lines
   Whose features some poets % within a stanza.
   will curse.

                             % One or more blank lines separate stanzas.

   For instead of making\\
   Them do \emph{all} line breaking, \\
   It allows them to put too many words on a line when they'd rather be
   forced to be terse.
\end{verse}

Mathematical formulas may also be displayed.  A
displayed formula
is
one-line long; multiline
formulas require special formatting instructions.
   \[  \ip{\Gamma}{\psi'} = x'' + y^{2} + z_{i}^{n}\]
Don't start a paragraph with a displayed equation,
nor make one a paragraph by itself.

\end{document}               % End of document.

\section{Bench Huge Section 732}
% synthetic bench section 732
% This is a sample LaTeX input file.  (Version of 12 August 2004.)
%
% A '%' character causes TeX to ignore all remaining text on the line,
% and is used for comments like this one.

\documentclass{article}      % Specifies the document class

                             % The preamble begins here.
\title{An Example Document}  % Declares the document's title.
\author{Leslie Lamport}      % Declares the author's name.
\date{January 21, 1994}      % Deleting this command produces today's date.

\newcommand{\ip}[2]{(#1, #2)}
                             % Defines \ip{arg1}{arg2} to mean
                             % (arg1, arg2).

%\newcommand{\ip}[2]{\langle #1 | #2\rangle}
                             % This is an alternative definition of
                             % \ip that is commented out.

\begin{document}             % End of preamble and beginning of text.

\maketitle                   % Produces the title.

This is an example input file.  Comparing it with
the output it generates can show you how to
produce a simple document of your own.

\section{Ordinary Text}      % Produces section heading.  Lower-level
                             % sections are begun with similar
                             % \subsection and \subsubsection commands.

The ends  of words and sentences are marked
  by   spaces. It  doesn't matter how many
spaces    you type; one is as good as 100.  The
end of   a line counts as a space.

One   or more   blank lines denote the  end
of  a paragraph.

Since any number of consecutive spaces are treated
like a single one, the formatting of the input
file makes no difference to
      \LaTeX,                % The \LaTeX command generates the LaTeX logo.
but it makes a difference to you.  When you use
\LaTeX, making your input file as easy to read
as possible will be a great help as you write
your document and when you change it.  This sample
file shows how you can add comments to your own input
file.

Because printing is different from typewriting,
there are a number of things that you have to do
differently when preparing an input file than if
you were just typing the document directly.
Quotation marks like
       ``this''
have to be handled specially, as do quotes within
quotes:
       ``\,`this'            % \, separates the double and single quote.
        is what I just
        wrote, not  `that'\,''.

Dashes come in three sizes: an
       intra-word
dash, a medium dash for number ranges like
       1--2,
and a punctuation
       dash---like
this.

A sentence-ending space should be larger than the
space between words within a sentence.  You
sometimes have to type special commands in
conjunction with punctuation characters to get
this right, as in the following sentence.
       Gnats, gnus, etc.\ all  % `\ ' makes an inter-word space.
       begin with G\@.         % \@ marks end-of-sentence punctuation.
You should check the spaces after periods when
reading your output to make sure you haven't
forgotten any special cases.  Generating an
ellipsis
       \ldots\               % `\ ' is needed after `\ldots' because TeX
                             % ignores spaces after command names like \ldots
                             % made from \ + letters.
                             %
                             % Note how a `%' character causes TeX to ignore
                             % the end of the input line, so these blank lines
                             % do not start a new paragraph.
                             %
with the right spacing around the periods requires
a special command.

\LaTeX\ interprets some common characters as
commands, so you must type special commands to
generate them.  These characters include the
following:
       \$ \& \% \# \{ and \}.

In printing, text is usually emphasized with an
       \emph{italic}
type style.

\begin{em}
   A long segment of text can also be emphasized
   in this way.  Text within such a segment can be
   given \emph{additional} emphasis.
\end{em}

It is sometimes necessary to prevent \LaTeX\ from
breaking a line where it might otherwise do so.
This may be at a space, as between the ``Mr.''\ and
``Jones'' in
       ``Mr.~Jones'',        % ~ produces an unbreakable interword space.
or within a word---especially when the word is a
symbol like
       \mbox{\emph{itemnum}}
that makes little sense when hyphenated across
lines.

Footnotes\footnote{This is an example of a footnote.}
pose no problem.

\LaTeX\ is good at typesetting mathematical formulas
like
       \( x-3y + z = 7 \)
or
       \( a_{1} > x^{2n} + y^{2n} > x' \)
or
       \( \ip{A}{B} = \sum_{i} a_{i} b_{i} \).
The spaces you type in a formula are
ignored.  Remember that a letter like
       $x$                   % $ ... $  and  \( ... \)  are equivalent
is a formula when it denotes a mathematical
symbol, and it should be typed as one.

\section{Displayed Text}

Text is displayed by indenting it from the left
margin.  Quotations are commonly displayed.  There
are short quotations
\begin{quote}
   This is a short quotation.  It consists of a
   single paragraph of text.  See how it is formatted.
\end{quote}
and longer ones.
\begin{quotation}
   This is a longer quotation.  It consists of two
   paragraphs of text, neither of which are
   particularly interesting.

   This is the second paragraph of the quotation.  It
   is just as dull as the first paragraph.
\end{quotation}
Another frequently-displayed structure is a list.
The following is an example of an \emph{itemized}
list.
\begin{itemize}
   \item This is the first item of an itemized list.
         Each item in the list is marked with a ``tick''.
         You don't have to worry about what kind of tick
         mark is used.

   \item This is the second item of the list.  It
         contains another list nested inside it.  The inner
         list is an \emph{enumerated} list.
         \begin{enumerate}
            \item This is the first item of an enumerated
                  list that is nested within the itemized list.

            \item This is the second item of the inner list.
                  \LaTeX\ allows you to nest lists deeper than
                  you really should.
         \end{enumerate}
         This is the rest of the second item of the outer
         list.  It is no more interesting than any other
         part of the item.
   \item This is the third item of the list.
\end{itemize}
You can even display poetry.
\begin{verse}
   There is an environment
    for verse \\             % The \\ command separates lines
   Whose features some poets % within a stanza.
   will curse.

                             % One or more blank lines separate stanzas.

   For instead of making\\
   Them do \emph{all} line breaking, \\
   It allows them to put too many words on a line when they'd rather be
   forced to be terse.
\end{verse}

Mathematical formulas may also be displayed.  A
displayed formula
is
one-line long; multiline
formulas require special formatting instructions.
   \[  \ip{\Gamma}{\psi'} = x'' + y^{2} + z_{i}^{n}\]
Don't start a paragraph with a displayed equation,
nor make one a paragraph by itself.

\end{document}               % End of document.

\section{Bench Huge Section 733}
% synthetic bench section 733
% This is a sample LaTeX input file.  (Version of 12 August 2004.)
%
% A '%' character causes TeX to ignore all remaining text on the line,
% and is used for comments like this one.

\documentclass{article}      % Specifies the document class

                             % The preamble begins here.
\title{An Example Document}  % Declares the document's title.
\author{Leslie Lamport}      % Declares the author's name.
\date{January 21, 1994}      % Deleting this command produces today's date.

\newcommand{\ip}[2]{(#1, #2)}
                             % Defines \ip{arg1}{arg2} to mean
                             % (arg1, arg2).

%\newcommand{\ip}[2]{\langle #1 | #2\rangle}
                             % This is an alternative definition of
                             % \ip that is commented out.

\begin{document}             % End of preamble and beginning of text.

\maketitle                   % Produces the title.

This is an example input file.  Comparing it with
the output it generates can show you how to
produce a simple document of your own.

\section{Ordinary Text}      % Produces section heading.  Lower-level
                             % sections are begun with similar
                             % \subsection and \subsubsection commands.

The ends  of words and sentences are marked
  by   spaces. It  doesn't matter how many
spaces    you type; one is as good as 100.  The
end of   a line counts as a space.

One   or more   blank lines denote the  end
of  a paragraph.

Since any number of consecutive spaces are treated
like a single one, the formatting of the input
file makes no difference to
      \LaTeX,                % The \LaTeX command generates the LaTeX logo.
but it makes a difference to you.  When you use
\LaTeX, making your input file as easy to read
as possible will be a great help as you write
your document and when you change it.  This sample
file shows how you can add comments to your own input
file.

Because printing is different from typewriting,
there are a number of things that you have to do
differently when preparing an input file than if
you were just typing the document directly.
Quotation marks like
       ``this''
have to be handled specially, as do quotes within
quotes:
       ``\,`this'            % \, separates the double and single quote.
        is what I just
        wrote, not  `that'\,''.

Dashes come in three sizes: an
       intra-word
dash, a medium dash for number ranges like
       1--2,
and a punctuation
       dash---like
this.

A sentence-ending space should be larger than the
space between words within a sentence.  You
sometimes have to type special commands in
conjunction with punctuation characters to get
this right, as in the following sentence.
       Gnats, gnus, etc.\ all  % `\ ' makes an inter-word space.
       begin with G\@.         % \@ marks end-of-sentence punctuation.
You should check the spaces after periods when
reading your output to make sure you haven't
forgotten any special cases.  Generating an
ellipsis
       \ldots\               % `\ ' is needed after `\ldots' because TeX
                             % ignores spaces after command names like \ldots
                             % made from \ + letters.
                             %
                             % Note how a `%' character causes TeX to ignore
                             % the end of the input line, so these blank lines
                             % do not start a new paragraph.
                             %
with the right spacing around the periods requires
a special command.

\LaTeX\ interprets some common characters as
commands, so you must type special commands to
generate them.  These characters include the
following:
       \$ \& \% \# \{ and \}.

In printing, text is usually emphasized with an
       \emph{italic}
type style.

\begin{em}
   A long segment of text can also be emphasized
   in this way.  Text within such a segment can be
   given \emph{additional} emphasis.
\end{em}

It is sometimes necessary to prevent \LaTeX\ from
breaking a line where it might otherwise do so.
This may be at a space, as between the ``Mr.''\ and
``Jones'' in
       ``Mr.~Jones'',        % ~ produces an unbreakable interword space.
or within a word---especially when the word is a
symbol like
       \mbox{\emph{itemnum}}
that makes little sense when hyphenated across
lines.

Footnotes\footnote{This is an example of a footnote.}
pose no problem.

\LaTeX\ is good at typesetting mathematical formulas
like
       \( x-3y + z = 7 \)
or
       \( a_{1} > x^{2n} + y^{2n} > x' \)
or
       \( \ip{A}{B} = \sum_{i} a_{i} b_{i} \).
The spaces you type in a formula are
ignored.  Remember that a letter like
       $x$                   % $ ... $  and  \( ... \)  are equivalent
is a formula when it denotes a mathematical
symbol, and it should be typed as one.

\section{Displayed Text}

Text is displayed by indenting it from the left
margin.  Quotations are commonly displayed.  There
are short quotations
\begin{quote}
   This is a short quotation.  It consists of a
   single paragraph of text.  See how it is formatted.
\end{quote}
and longer ones.
\begin{quotation}
   This is a longer quotation.  It consists of two
   paragraphs of text, neither of which are
   particularly interesting.

   This is the second paragraph of the quotation.  It
   is just as dull as the first paragraph.
\end{quotation}
Another frequently-displayed structure is a list.
The following is an example of an \emph{itemized}
list.
\begin{itemize}
   \item This is the first item of an itemized list.
         Each item in the list is marked with a ``tick''.
         You don't have to worry about what kind of tick
         mark is used.

   \item This is the second item of the list.  It
         contains another list nested inside it.  The inner
         list is an \emph{enumerated} list.
         \begin{enumerate}
            \item This is the first item of an enumerated
                  list that is nested within the itemized list.

            \item This is the second item of the inner list.
                  \LaTeX\ allows you to nest lists deeper than
                  you really should.
         \end{enumerate}
         This is the rest of the second item of the outer
         list.  It is no more interesting than any other
         part of the item.
   \item This is the third item of the list.
\end{itemize}
You can even display poetry.
\begin{verse}
   There is an environment
    for verse \\             % The \\ command separates lines
   Whose features some poets % within a stanza.
   will curse.

                             % One or more blank lines separate stanzas.

   For instead of making\\
   Them do \emph{all} line breaking, \\
   It allows them to put too many words on a line when they'd rather be
   forced to be terse.
\end{verse}

Mathematical formulas may also be displayed.  A
displayed formula
is
one-line long; multiline
formulas require special formatting instructions.
   \[  \ip{\Gamma}{\psi'} = x'' + y^{2} + z_{i}^{n}\]
Don't start a paragraph with a displayed equation,
nor make one a paragraph by itself.

\end{document}               % End of document.

\section{Bench Huge Section 734}
% synthetic bench section 734
% This is a sample LaTeX input file.  (Version of 12 August 2004.)
%
% A '%' character causes TeX to ignore all remaining text on the line,
% and is used for comments like this one.

\documentclass{article}      % Specifies the document class

                             % The preamble begins here.
\title{An Example Document}  % Declares the document's title.
\author{Leslie Lamport}      % Declares the author's name.
\date{January 21, 1994}      % Deleting this command produces today's date.

\newcommand{\ip}[2]{(#1, #2)}
                             % Defines \ip{arg1}{arg2} to mean
                             % (arg1, arg2).

%\newcommand{\ip}[2]{\langle #1 | #2\rangle}
                             % This is an alternative definition of
                             % \ip that is commented out.

\begin{document}             % End of preamble and beginning of text.

\maketitle                   % Produces the title.

This is an example input file.  Comparing it with
the output it generates can show you how to
produce a simple document of your own.

\section{Ordinary Text}      % Produces section heading.  Lower-level
                             % sections are begun with similar
                             % \subsection and \subsubsection commands.

The ends  of words and sentences are marked
  by   spaces. It  doesn't matter how many
spaces    you type; one is as good as 100.  The
end of   a line counts as a space.

One   or more   blank lines denote the  end
of  a paragraph.

Since any number of consecutive spaces are treated
like a single one, the formatting of the input
file makes no difference to
      \LaTeX,                % The \LaTeX command generates the LaTeX logo.
but it makes a difference to you.  When you use
\LaTeX, making your input file as easy to read
as possible will be a great help as you write
your document and when you change it.  This sample
file shows how you can add comments to your own input
file.

Because printing is different from typewriting,
there are a number of things that you have to do
differently when preparing an input file than if
you were just typing the document directly.
Quotation marks like
       ``this''
have to be handled specially, as do quotes within
quotes:
       ``\,`this'            % \, separates the double and single quote.
        is what I just
        wrote, not  `that'\,''.

Dashes come in three sizes: an
       intra-word
dash, a medium dash for number ranges like
       1--2,
and a punctuation
       dash---like
this.

A sentence-ending space should be larger than the
space between words within a sentence.  You
sometimes have to type special commands in
conjunction with punctuation characters to get
this right, as in the following sentence.
       Gnats, gnus, etc.\ all  % `\ ' makes an inter-word space.
       begin with G\@.         % \@ marks end-of-sentence punctuation.
You should check the spaces after periods when
reading your output to make sure you haven't
forgotten any special cases.  Generating an
ellipsis
       \ldots\               % `\ ' is needed after `\ldots' because TeX
                             % ignores spaces after command names like \ldots
                             % made from \ + letters.
                             %
                             % Note how a `%' character causes TeX to ignore
                             % the end of the input line, so these blank lines
                             % do not start a new paragraph.
                             %
with the right spacing around the periods requires
a special command.

\LaTeX\ interprets some common characters as
commands, so you must type special commands to
generate them.  These characters include the
following:
       \$ \& \% \# \{ and \}.

In printing, text is usually emphasized with an
       \emph{italic}
type style.

\begin{em}
   A long segment of text can also be emphasized
   in this way.  Text within such a segment can be
   given \emph{additional} emphasis.
\end{em}

It is sometimes necessary to prevent \LaTeX\ from
breaking a line where it might otherwise do so.
This may be at a space, as between the ``Mr.''\ and
``Jones'' in
       ``Mr.~Jones'',        % ~ produces an unbreakable interword space.
or within a word---especially when the word is a
symbol like
       \mbox{\emph{itemnum}}
that makes little sense when hyphenated across
lines.

Footnotes\footnote{This is an example of a footnote.}
pose no problem.

\LaTeX\ is good at typesetting mathematical formulas
like
       \( x-3y + z = 7 \)
or
       \( a_{1} > x^{2n} + y^{2n} > x' \)
or
       \( \ip{A}{B} = \sum_{i} a_{i} b_{i} \).
The spaces you type in a formula are
ignored.  Remember that a letter like
       $x$                   % $ ... $  and  \( ... \)  are equivalent
is a formula when it denotes a mathematical
symbol, and it should be typed as one.

\section{Displayed Text}

Text is displayed by indenting it from the left
margin.  Quotations are commonly displayed.  There
are short quotations
\begin{quote}
   This is a short quotation.  It consists of a
   single paragraph of text.  See how it is formatted.
\end{quote}
and longer ones.
\begin{quotation}
   This is a longer quotation.  It consists of two
   paragraphs of text, neither of which are
   particularly interesting.

   This is the second paragraph of the quotation.  It
   is just as dull as the first paragraph.
\end{quotation}
Another frequently-displayed structure is a list.
The following is an example of an \emph{itemized}
list.
\begin{itemize}
   \item This is the first item of an itemized list.
         Each item in the list is marked with a ``tick''.
         You don't have to worry about what kind of tick
         mark is used.

   \item This is the second item of the list.  It
         contains another list nested inside it.  The inner
         list is an \emph{enumerated} list.
         \begin{enumerate}
            \item This is the first item of an enumerated
                  list that is nested within the itemized list.

            \item This is the second item of the inner list.
                  \LaTeX\ allows you to nest lists deeper than
                  you really should.
         \end{enumerate}
         This is the rest of the second item of the outer
         list.  It is no more interesting than any other
         part of the item.
   \item This is the third item of the list.
\end{itemize}
You can even display poetry.
\begin{verse}
   There is an environment
    for verse \\             % The \\ command separates lines
   Whose features some poets % within a stanza.
   will curse.

                             % One or more blank lines separate stanzas.

   For instead of making\\
   Them do \emph{all} line breaking, \\
   It allows them to put too many words on a line when they'd rather be
   forced to be terse.
\end{verse}

Mathematical formulas may also be displayed.  A
displayed formula
is
one-line long; multiline
formulas require special formatting instructions.
   \[  \ip{\Gamma}{\psi'} = x'' + y^{2} + z_{i}^{n}\]
Don't start a paragraph with a displayed equation,
nor make one a paragraph by itself.

\end{document}               % End of document.

\section{Bench Huge Section 735}
% synthetic bench section 735
% This is a sample LaTeX input file.  (Version of 12 August 2004.)
%
% A '%' character causes TeX to ignore all remaining text on the line,
% and is used for comments like this one.

\documentclass{article}      % Specifies the document class

                             % The preamble begins here.
\title{An Example Document}  % Declares the document's title.
\author{Leslie Lamport}      % Declares the author's name.
\date{January 21, 1994}      % Deleting this command produces today's date.

\newcommand{\ip}[2]{(#1, #2)}
                             % Defines \ip{arg1}{arg2} to mean
                             % (arg1, arg2).

%\newcommand{\ip}[2]{\langle #1 | #2\rangle}
                             % This is an alternative definition of
                             % \ip that is commented out.

\begin{document}             % End of preamble and beginning of text.

\maketitle                   % Produces the title.

This is an example input file.  Comparing it with
the output it generates can show you how to
produce a simple document of your own.

\section{Ordinary Text}      % Produces section heading.  Lower-level
                             % sections are begun with similar
                             % \subsection and \subsubsection commands.

The ends  of words and sentences are marked
  by   spaces. It  doesn't matter how many
spaces    you type; one is as good as 100.  The
end of   a line counts as a space.

One   or more   blank lines denote the  end
of  a paragraph.

Since any number of consecutive spaces are treated
like a single one, the formatting of the input
file makes no difference to
      \LaTeX,                % The \LaTeX command generates the LaTeX logo.
but it makes a difference to you.  When you use
\LaTeX, making your input file as easy to read
as possible will be a great help as you write
your document and when you change it.  This sample
file shows how you can add comments to your own input
file.

Because printing is different from typewriting,
there are a number of things that you have to do
differently when preparing an input file than if
you were just typing the document directly.
Quotation marks like
       ``this''
have to be handled specially, as do quotes within
quotes:
       ``\,`this'            % \, separates the double and single quote.
        is what I just
        wrote, not  `that'\,''.

Dashes come in three sizes: an
       intra-word
dash, a medium dash for number ranges like
       1--2,
and a punctuation
       dash---like
this.

A sentence-ending space should be larger than the
space between words within a sentence.  You
sometimes have to type special commands in
conjunction with punctuation characters to get
this right, as in the following sentence.
       Gnats, gnus, etc.\ all  % `\ ' makes an inter-word space.
       begin with G\@.         % \@ marks end-of-sentence punctuation.
You should check the spaces after periods when
reading your output to make sure you haven't
forgotten any special cases.  Generating an
ellipsis
       \ldots\               % `\ ' is needed after `\ldots' because TeX
                             % ignores spaces after command names like \ldots
                             % made from \ + letters.
                             %
                             % Note how a `%' character causes TeX to ignore
                             % the end of the input line, so these blank lines
                             % do not start a new paragraph.
                             %
with the right spacing around the periods requires
a special command.

\LaTeX\ interprets some common characters as
commands, so you must type special commands to
generate them.  These characters include the
following:
       \$ \& \% \# \{ and \}.

In printing, text is usually emphasized with an
       \emph{italic}
type style.

\begin{em}
   A long segment of text can also be emphasized
   in this way.  Text within such a segment can be
   given \emph{additional} emphasis.
\end{em}

It is sometimes necessary to prevent \LaTeX\ from
breaking a line where it might otherwise do so.
This may be at a space, as between the ``Mr.''\ and
``Jones'' in
       ``Mr.~Jones'',        % ~ produces an unbreakable interword space.
or within a word---especially when the word is a
symbol like
       \mbox{\emph{itemnum}}
that makes little sense when hyphenated across
lines.

Footnotes\footnote{This is an example of a footnote.}
pose no problem.

\LaTeX\ is good at typesetting mathematical formulas
like
       \( x-3y + z = 7 \)
or
       \( a_{1} > x^{2n} + y^{2n} > x' \)
or
       \( \ip{A}{B} = \sum_{i} a_{i} b_{i} \).
The spaces you type in a formula are
ignored.  Remember that a letter like
       $x$                   % $ ... $  and  \( ... \)  are equivalent
is a formula when it denotes a mathematical
symbol, and it should be typed as one.

\section{Displayed Text}

Text is displayed by indenting it from the left
margin.  Quotations are commonly displayed.  There
are short quotations
\begin{quote}
   This is a short quotation.  It consists of a
   single paragraph of text.  See how it is formatted.
\end{quote}
and longer ones.
\begin{quotation}
   This is a longer quotation.  It consists of two
   paragraphs of text, neither of which are
   particularly interesting.

   This is the second paragraph of the quotation.  It
   is just as dull as the first paragraph.
\end{quotation}
Another frequently-displayed structure is a list.
The following is an example of an \emph{itemized}
list.
\begin{itemize}
   \item This is the first item of an itemized list.
         Each item in the list is marked with a ``tick''.
         You don't have to worry about what kind of tick
         mark is used.

   \item This is the second item of the list.  It
         contains another list nested inside it.  The inner
         list is an \emph{enumerated} list.
         \begin{enumerate}
            \item This is the first item of an enumerated
                  list that is nested within the itemized list.

            \item This is the second item of the inner list.
                  \LaTeX\ allows you to nest lists deeper than
                  you really should.
         \end{enumerate}
         This is the rest of the second item of the outer
         list.  It is no more interesting than any other
         part of the item.
   \item This is the third item of the list.
\end{itemize}
You can even display poetry.
\begin{verse}
   There is an environment
    for verse \\             % The \\ command separates lines
   Whose features some poets % within a stanza.
   will curse.

                             % One or more blank lines separate stanzas.

   For instead of making\\
   Them do \emph{all} line breaking, \\
   It allows them to put too many words on a line when they'd rather be
   forced to be terse.
\end{verse}

Mathematical formulas may also be displayed.  A
displayed formula
is
one-line long; multiline
formulas require special formatting instructions.
   \[  \ip{\Gamma}{\psi'} = x'' + y^{2} + z_{i}^{n}\]
Don't start a paragraph with a displayed equation,
nor make one a paragraph by itself.

\end{document}               % End of document.

\section{Bench Huge Section 736}
% synthetic bench section 736
% This is a sample LaTeX input file.  (Version of 12 August 2004.)
%
% A '%' character causes TeX to ignore all remaining text on the line,
% and is used for comments like this one.

\documentclass{article}      % Specifies the document class

                             % The preamble begins here.
\title{An Example Document}  % Declares the document's title.
\author{Leslie Lamport}      % Declares the author's name.
\date{January 21, 1994}      % Deleting this command produces today's date.

\newcommand{\ip}[2]{(#1, #2)}
                             % Defines \ip{arg1}{arg2} to mean
                             % (arg1, arg2).

%\newcommand{\ip}[2]{\langle #1 | #2\rangle}
                             % This is an alternative definition of
                             % \ip that is commented out.

\begin{document}             % End of preamble and beginning of text.

\maketitle                   % Produces the title.

This is an example input file.  Comparing it with
the output it generates can show you how to
produce a simple document of your own.

\section{Ordinary Text}      % Produces section heading.  Lower-level
                             % sections are begun with similar
                             % \subsection and \subsubsection commands.

The ends  of words and sentences are marked
  by   spaces. It  doesn't matter how many
spaces    you type; one is as good as 100.  The
end of   a line counts as a space.

One   or more   blank lines denote the  end
of  a paragraph.

Since any number of consecutive spaces are treated
like a single one, the formatting of the input
file makes no difference to
      \LaTeX,                % The \LaTeX command generates the LaTeX logo.
but it makes a difference to you.  When you use
\LaTeX, making your input file as easy to read
as possible will be a great help as you write
your document and when you change it.  This sample
file shows how you can add comments to your own input
file.

Because printing is different from typewriting,
there are a number of things that you have to do
differently when preparing an input file than if
you were just typing the document directly.
Quotation marks like
       ``this''
have to be handled specially, as do quotes within
quotes:
       ``\,`this'            % \, separates the double and single quote.
        is what I just
        wrote, not  `that'\,''.

Dashes come in three sizes: an
       intra-word
dash, a medium dash for number ranges like
       1--2,
and a punctuation
       dash---like
this.

A sentence-ending space should be larger than the
space between words within a sentence.  You
sometimes have to type special commands in
conjunction with punctuation characters to get
this right, as in the following sentence.
       Gnats, gnus, etc.\ all  % `\ ' makes an inter-word space.
       begin with G\@.         % \@ marks end-of-sentence punctuation.
You should check the spaces after periods when
reading your output to make sure you haven't
forgotten any special cases.  Generating an
ellipsis
       \ldots\               % `\ ' is needed after `\ldots' because TeX
                             % ignores spaces after command names like \ldots
                             % made from \ + letters.
                             %
                             % Note how a `%' character causes TeX to ignore
                             % the end of the input line, so these blank lines
                             % do not start a new paragraph.
                             %
with the right spacing around the periods requires
a special command.

\LaTeX\ interprets some common characters as
commands, so you must type special commands to
generate them.  These characters include the
following:
       \$ \& \% \# \{ and \}.

In printing, text is usually emphasized with an
       \emph{italic}
type style.

\begin{em}
   A long segment of text can also be emphasized
   in this way.  Text within such a segment can be
   given \emph{additional} emphasis.
\end{em}

It is sometimes necessary to prevent \LaTeX\ from
breaking a line where it might otherwise do so.
This may be at a space, as between the ``Mr.''\ and
``Jones'' in
       ``Mr.~Jones'',        % ~ produces an unbreakable interword space.
or within a word---especially when the word is a
symbol like
       \mbox{\emph{itemnum}}
that makes little sense when hyphenated across
lines.

Footnotes\footnote{This is an example of a footnote.}
pose no problem.

\LaTeX\ is good at typesetting mathematical formulas
like
       \( x-3y + z = 7 \)
or
       \( a_{1} > x^{2n} + y^{2n} > x' \)
or
       \( \ip{A}{B} = \sum_{i} a_{i} b_{i} \).
The spaces you type in a formula are
ignored.  Remember that a letter like
       $x$                   % $ ... $  and  \( ... \)  are equivalent
is a formula when it denotes a mathematical
symbol, and it should be typed as one.

\section{Displayed Text}

Text is displayed by indenting it from the left
margin.  Quotations are commonly displayed.  There
are short quotations
\begin{quote}
   This is a short quotation.  It consists of a
   single paragraph of text.  See how it is formatted.
\end{quote}
and longer ones.
\begin{quotation}
   This is a longer quotation.  It consists of two
   paragraphs of text, neither of which are
   particularly interesting.

   This is the second paragraph of the quotation.  It
   is just as dull as the first paragraph.
\end{quotation}
Another frequently-displayed structure is a list.
The following is an example of an \emph{itemized}
list.
\begin{itemize}
   \item This is the first item of an itemized list.
         Each item in the list is marked with a ``tick''.
         You don't have to worry about what kind of tick
         mark is used.

   \item This is the second item of the list.  It
         contains another list nested inside it.  The inner
         list is an \emph{enumerated} list.
         \begin{enumerate}
            \item This is the first item of an enumerated
                  list that is nested within the itemized list.

            \item This is the second item of the inner list.
                  \LaTeX\ allows you to nest lists deeper than
                  you really should.
         \end{enumerate}
         This is the rest of the second item of the outer
         list.  It is no more interesting than any other
         part of the item.
   \item This is the third item of the list.
\end{itemize}
You can even display poetry.
\begin{verse}
   There is an environment
    for verse \\             % The \\ command separates lines
   Whose features some poets % within a stanza.
   will curse.

                             % One or more blank lines separate stanzas.

   For instead of making\\
   Them do \emph{all} line breaking, \\
   It allows them to put too many words on a line when they'd rather be
   forced to be terse.
\end{verse}

Mathematical formulas may also be displayed.  A
displayed formula
is
one-line long; multiline
formulas require special formatting instructions.
   \[  \ip{\Gamma}{\psi'} = x'' + y^{2} + z_{i}^{n}\]
Don't start a paragraph with a displayed equation,
nor make one a paragraph by itself.

\end{document}               % End of document.

\section{Bench Huge Section 737}
% synthetic bench section 737
% This is a sample LaTeX input file.  (Version of 12 August 2004.)
%
% A '%' character causes TeX to ignore all remaining text on the line,
% and is used for comments like this one.

\documentclass{article}      % Specifies the document class

                             % The preamble begins here.
\title{An Example Document}  % Declares the document's title.
\author{Leslie Lamport}      % Declares the author's name.
\date{January 21, 1994}      % Deleting this command produces today's date.

\newcommand{\ip}[2]{(#1, #2)}
                             % Defines \ip{arg1}{arg2} to mean
                             % (arg1, arg2).

%\newcommand{\ip}[2]{\langle #1 | #2\rangle}
                             % This is an alternative definition of
                             % \ip that is commented out.

\begin{document}             % End of preamble and beginning of text.

\maketitle                   % Produces the title.

This is an example input file.  Comparing it with
the output it generates can show you how to
produce a simple document of your own.

\section{Ordinary Text}      % Produces section heading.  Lower-level
                             % sections are begun with similar
                             % \subsection and \subsubsection commands.

The ends  of words and sentences are marked
  by   spaces. It  doesn't matter how many
spaces    you type; one is as good as 100.  The
end of   a line counts as a space.

One   or more   blank lines denote the  end
of  a paragraph.

Since any number of consecutive spaces are treated
like a single one, the formatting of the input
file makes no difference to
      \LaTeX,                % The \LaTeX command generates the LaTeX logo.
but it makes a difference to you.  When you use
\LaTeX, making your input file as easy to read
as possible will be a great help as you write
your document and when you change it.  This sample
file shows how you can add comments to your own input
file.

Because printing is different from typewriting,
there are a number of things that you have to do
differently when preparing an input file than if
you were just typing the document directly.
Quotation marks like
       ``this''
have to be handled specially, as do quotes within
quotes:
       ``\,`this'            % \, separates the double and single quote.
        is what I just
        wrote, not  `that'\,''.

Dashes come in three sizes: an
       intra-word
dash, a medium dash for number ranges like
       1--2,
and a punctuation
       dash---like
this.

A sentence-ending space should be larger than the
space between words within a sentence.  You
sometimes have to type special commands in
conjunction with punctuation characters to get
this right, as in the following sentence.
       Gnats, gnus, etc.\ all  % `\ ' makes an inter-word space.
       begin with G\@.         % \@ marks end-of-sentence punctuation.
You should check the spaces after periods when
reading your output to make sure you haven't
forgotten any special cases.  Generating an
ellipsis
       \ldots\               % `\ ' is needed after `\ldots' because TeX
                             % ignores spaces after command names like \ldots
                             % made from \ + letters.
                             %
                             % Note how a `%' character causes TeX to ignore
                             % the end of the input line, so these blank lines
                             % do not start a new paragraph.
                             %
with the right spacing around the periods requires
a special command.

\LaTeX\ interprets some common characters as
commands, so you must type special commands to
generate them.  These characters include the
following:
       \$ \& \% \# \{ and \}.

In printing, text is usually emphasized with an
       \emph{italic}
type style.

\begin{em}
   A long segment of text can also be emphasized
   in this way.  Text within such a segment can be
   given \emph{additional} emphasis.
\end{em}

It is sometimes necessary to prevent \LaTeX\ from
breaking a line where it might otherwise do so.
This may be at a space, as between the ``Mr.''\ and
``Jones'' in
       ``Mr.~Jones'',        % ~ produces an unbreakable interword space.
or within a word---especially when the word is a
symbol like
       \mbox{\emph{itemnum}}
that makes little sense when hyphenated across
lines.

Footnotes\footnote{This is an example of a footnote.}
pose no problem.

\LaTeX\ is good at typesetting mathematical formulas
like
       \( x-3y + z = 7 \)
or
       \( a_{1} > x^{2n} + y^{2n} > x' \)
or
       \( \ip{A}{B} = \sum_{i} a_{i} b_{i} \).
The spaces you type in a formula are
ignored.  Remember that a letter like
       $x$                   % $ ... $  and  \( ... \)  are equivalent
is a formula when it denotes a mathematical
symbol, and it should be typed as one.

\section{Displayed Text}

Text is displayed by indenting it from the left
margin.  Quotations are commonly displayed.  There
are short quotations
\begin{quote}
   This is a short quotation.  It consists of a
   single paragraph of text.  See how it is formatted.
\end{quote}
and longer ones.
\begin{quotation}
   This is a longer quotation.  It consists of two
   paragraphs of text, neither of which are
   particularly interesting.

   This is the second paragraph of the quotation.  It
   is just as dull as the first paragraph.
\end{quotation}
Another frequently-displayed structure is a list.
The following is an example of an \emph{itemized}
list.
\begin{itemize}
   \item This is the first item of an itemized list.
         Each item in the list is marked with a ``tick''.
         You don't have to worry about what kind of tick
         mark is used.

   \item This is the second item of the list.  It
         contains another list nested inside it.  The inner
         list is an \emph{enumerated} list.
         \begin{enumerate}
            \item This is the first item of an enumerated
                  list that is nested within the itemized list.

            \item This is the second item of the inner list.
                  \LaTeX\ allows you to nest lists deeper than
                  you really should.
         \end{enumerate}
         This is the rest of the second item of the outer
         list.  It is no more interesting than any other
         part of the item.
   \item This is the third item of the list.
\end{itemize}
You can even display poetry.
\begin{verse}
   There is an environment
    for verse \\             % The \\ command separates lines
   Whose features some poets % within a stanza.
   will curse.

                             % One or more blank lines separate stanzas.

   For instead of making\\
   Them do \emph{all} line breaking, \\
   It allows them to put too many words on a line when they'd rather be
   forced to be terse.
\end{verse}

Mathematical formulas may also be displayed.  A
displayed formula
is
one-line long; multiline
formulas require special formatting instructions.
   \[  \ip{\Gamma}{\psi'} = x'' + y^{2} + z_{i}^{n}\]
Don't start a paragraph with a displayed equation,
nor make one a paragraph by itself.

\end{document}               % End of document.

\section{Bench Huge Section 738}
% synthetic bench section 738
% This is a sample LaTeX input file.  (Version of 12 August 2004.)
%
% A '%' character causes TeX to ignore all remaining text on the line,
% and is used for comments like this one.

\documentclass{article}      % Specifies the document class

                             % The preamble begins here.
\title{An Example Document}  % Declares the document's title.
\author{Leslie Lamport}      % Declares the author's name.
\date{January 21, 1994}      % Deleting this command produces today's date.

\newcommand{\ip}[2]{(#1, #2)}
                             % Defines \ip{arg1}{arg2} to mean
                             % (arg1, arg2).

%\newcommand{\ip}[2]{\langle #1 | #2\rangle}
                             % This is an alternative definition of
                             % \ip that is commented out.

\begin{document}             % End of preamble and beginning of text.

\maketitle                   % Produces the title.

This is an example input file.  Comparing it with
the output it generates can show you how to
produce a simple document of your own.

\section{Ordinary Text}      % Produces section heading.  Lower-level
                             % sections are begun with similar
                             % \subsection and \subsubsection commands.

The ends  of words and sentences are marked
  by   spaces. It  doesn't matter how many
spaces    you type; one is as good as 100.  The
end of   a line counts as a space.

One   or more   blank lines denote the  end
of  a paragraph.

Since any number of consecutive spaces are treated
like a single one, the formatting of the input
file makes no difference to
      \LaTeX,                % The \LaTeX command generates the LaTeX logo.
but it makes a difference to you.  When you use
\LaTeX, making your input file as easy to read
as possible will be a great help as you write
your document and when you change it.  This sample
file shows how you can add comments to your own input
file.

Because printing is different from typewriting,
there are a number of things that you have to do
differently when preparing an input file than if
you were just typing the document directly.
Quotation marks like
       ``this''
have to be handled specially, as do quotes within
quotes:
       ``\,`this'            % \, separates the double and single quote.
        is what I just
        wrote, not  `that'\,''.

Dashes come in three sizes: an
       intra-word
dash, a medium dash for number ranges like
       1--2,
and a punctuation
       dash---like
this.

A sentence-ending space should be larger than the
space between words within a sentence.  You
sometimes have to type special commands in
conjunction with punctuation characters to get
this right, as in the following sentence.
       Gnats, gnus, etc.\ all  % `\ ' makes an inter-word space.
       begin with G\@.         % \@ marks end-of-sentence punctuation.
You should check the spaces after periods when
reading your output to make sure you haven't
forgotten any special cases.  Generating an
ellipsis
       \ldots\               % `\ ' is needed after `\ldots' because TeX
                             % ignores spaces after command names like \ldots
                             % made from \ + letters.
                             %
                             % Note how a `%' character causes TeX to ignore
                             % the end of the input line, so these blank lines
                             % do not start a new paragraph.
                             %
with the right spacing around the periods requires
a special command.

\LaTeX\ interprets some common characters as
commands, so you must type special commands to
generate them.  These characters include the
following:
       \$ \& \% \# \{ and \}.

In printing, text is usually emphasized with an
       \emph{italic}
type style.

\begin{em}
   A long segment of text can also be emphasized
   in this way.  Text within such a segment can be
   given \emph{additional} emphasis.
\end{em}

It is sometimes necessary to prevent \LaTeX\ from
breaking a line where it might otherwise do so.
This may be at a space, as between the ``Mr.''\ and
``Jones'' in
       ``Mr.~Jones'',        % ~ produces an unbreakable interword space.
or within a word---especially when the word is a
symbol like
       \mbox{\emph{itemnum}}
that makes little sense when hyphenated across
lines.

Footnotes\footnote{This is an example of a footnote.}
pose no problem.

\LaTeX\ is good at typesetting mathematical formulas
like
       \( x-3y + z = 7 \)
or
       \( a_{1} > x^{2n} + y^{2n} > x' \)
or
       \( \ip{A}{B} = \sum_{i} a_{i} b_{i} \).
The spaces you type in a formula are
ignored.  Remember that a letter like
       $x$                   % $ ... $  and  \( ... \)  are equivalent
is a formula when it denotes a mathematical
symbol, and it should be typed as one.

\section{Displayed Text}

Text is displayed by indenting it from the left
margin.  Quotations are commonly displayed.  There
are short quotations
\begin{quote}
   This is a short quotation.  It consists of a
   single paragraph of text.  See how it is formatted.
\end{quote}
and longer ones.
\begin{quotation}
   This is a longer quotation.  It consists of two
   paragraphs of text, neither of which are
   particularly interesting.

   This is the second paragraph of the quotation.  It
   is just as dull as the first paragraph.
\end{quotation}
Another frequently-displayed structure is a list.
The following is an example of an \emph{itemized}
list.
\begin{itemize}
   \item This is the first item of an itemized list.
         Each item in the list is marked with a ``tick''.
         You don't have to worry about what kind of tick
         mark is used.

   \item This is the second item of the list.  It
         contains another list nested inside it.  The inner
         list is an \emph{enumerated} list.
         \begin{enumerate}
            \item This is the first item of an enumerated
                  list that is nested within the itemized list.

            \item This is the second item of the inner list.
                  \LaTeX\ allows you to nest lists deeper than
                  you really should.
         \end{enumerate}
         This is the rest of the second item of the outer
         list.  It is no more interesting than any other
         part of the item.
   \item This is the third item of the list.
\end{itemize}
You can even display poetry.
\begin{verse}
   There is an environment
    for verse \\             % The \\ command separates lines
   Whose features some poets % within a stanza.
   will curse.

                             % One or more blank lines separate stanzas.

   For instead of making\\
   Them do \emph{all} line breaking, \\
   It allows them to put too many words on a line when they'd rather be
   forced to be terse.
\end{verse}

Mathematical formulas may also be displayed.  A
displayed formula
is
one-line long; multiline
formulas require special formatting instructions.
   \[  \ip{\Gamma}{\psi'} = x'' + y^{2} + z_{i}^{n}\]
Don't start a paragraph with a displayed equation,
nor make one a paragraph by itself.

\end{document}               % End of document.

\section{Bench Huge Section 739}
% synthetic bench section 739
% This is a sample LaTeX input file.  (Version of 12 August 2004.)
%
% A '%' character causes TeX to ignore all remaining text on the line,
% and is used for comments like this one.

\documentclass{article}      % Specifies the document class

                             % The preamble begins here.
\title{An Example Document}  % Declares the document's title.
\author{Leslie Lamport}      % Declares the author's name.
\date{January 21, 1994}      % Deleting this command produces today's date.

\newcommand{\ip}[2]{(#1, #2)}
                             % Defines \ip{arg1}{arg2} to mean
                             % (arg1, arg2).

%\newcommand{\ip}[2]{\langle #1 | #2\rangle}
                             % This is an alternative definition of
                             % \ip that is commented out.

\begin{document}             % End of preamble and beginning of text.

\maketitle                   % Produces the title.

This is an example input file.  Comparing it with
the output it generates can show you how to
produce a simple document of your own.

\section{Ordinary Text}      % Produces section heading.  Lower-level
                             % sections are begun with similar
                             % \subsection and \subsubsection commands.

The ends  of words and sentences are marked
  by   spaces. It  doesn't matter how many
spaces    you type; one is as good as 100.  The
end of   a line counts as a space.

One   or more   blank lines denote the  end
of  a paragraph.

Since any number of consecutive spaces are treated
like a single one, the formatting of the input
file makes no difference to
      \LaTeX,                % The \LaTeX command generates the LaTeX logo.
but it makes a difference to you.  When you use
\LaTeX, making your input file as easy to read
as possible will be a great help as you write
your document and when you change it.  This sample
file shows how you can add comments to your own input
file.

Because printing is different from typewriting,
there are a number of things that you have to do
differently when preparing an input file than if
you were just typing the document directly.
Quotation marks like
       ``this''
have to be handled specially, as do quotes within
quotes:
       ``\,`this'            % \, separates the double and single quote.
        is what I just
        wrote, not  `that'\,''.

Dashes come in three sizes: an
       intra-word
dash, a medium dash for number ranges like
       1--2,
and a punctuation
       dash---like
this.

A sentence-ending space should be larger than the
space between words within a sentence.  You
sometimes have to type special commands in
conjunction with punctuation characters to get
this right, as in the following sentence.
       Gnats, gnus, etc.\ all  % `\ ' makes an inter-word space.
       begin with G\@.         % \@ marks end-of-sentence punctuation.
You should check the spaces after periods when
reading your output to make sure you haven't
forgotten any special cases.  Generating an
ellipsis
       \ldots\               % `\ ' is needed after `\ldots' because TeX
                             % ignores spaces after command names like \ldots
                             % made from \ + letters.
                             %
                             % Note how a `%' character causes TeX to ignore
                             % the end of the input line, so these blank lines
                             % do not start a new paragraph.
                             %
with the right spacing around the periods requires
a special command.

\LaTeX\ interprets some common characters as
commands, so you must type special commands to
generate them.  These characters include the
following:
       \$ \& \% \# \{ and \}.

In printing, text is usually emphasized with an
       \emph{italic}
type style.

\begin{em}
   A long segment of text can also be emphasized
   in this way.  Text within such a segment can be
   given \emph{additional} emphasis.
\end{em}

It is sometimes necessary to prevent \LaTeX\ from
breaking a line where it might otherwise do so.
This may be at a space, as between the ``Mr.''\ and
``Jones'' in
       ``Mr.~Jones'',        % ~ produces an unbreakable interword space.
or within a word---especially when the word is a
symbol like
       \mbox{\emph{itemnum}}
that makes little sense when hyphenated across
lines.

Footnotes\footnote{This is an example of a footnote.}
pose no problem.

\LaTeX\ is good at typesetting mathematical formulas
like
       \( x-3y + z = 7 \)
or
       \( a_{1} > x^{2n} + y^{2n} > x' \)
or
       \( \ip{A}{B} = \sum_{i} a_{i} b_{i} \).
The spaces you type in a formula are
ignored.  Remember that a letter like
       $x$                   % $ ... $  and  \( ... \)  are equivalent
is a formula when it denotes a mathematical
symbol, and it should be typed as one.

\section{Displayed Text}

Text is displayed by indenting it from the left
margin.  Quotations are commonly displayed.  There
are short quotations
\begin{quote}
   This is a short quotation.  It consists of a
   single paragraph of text.  See how it is formatted.
\end{quote}
and longer ones.
\begin{quotation}
   This is a longer quotation.  It consists of two
   paragraphs of text, neither of which are
   particularly interesting.

   This is the second paragraph of the quotation.  It
   is just as dull as the first paragraph.
\end{quotation}
Another frequently-displayed structure is a list.
The following is an example of an \emph{itemized}
list.
\begin{itemize}
   \item This is the first item of an itemized list.
         Each item in the list is marked with a ``tick''.
         You don't have to worry about what kind of tick
         mark is used.

   \item This is the second item of the list.  It
         contains another list nested inside it.  The inner
         list is an \emph{enumerated} list.
         \begin{enumerate}
            \item This is the first item of an enumerated
                  list that is nested within the itemized list.

            \item This is the second item of the inner list.
                  \LaTeX\ allows you to nest lists deeper than
                  you really should.
         \end{enumerate}
         This is the rest of the second item of the outer
         list.  It is no more interesting than any other
         part of the item.
   \item This is the third item of the list.
\end{itemize}
You can even display poetry.
\begin{verse}
   There is an environment
    for verse \\             % The \\ command separates lines
   Whose features some poets % within a stanza.
   will curse.

                             % One or more blank lines separate stanzas.

   For instead of making\\
   Them do \emph{all} line breaking, \\
   It allows them to put too many words on a line when they'd rather be
   forced to be terse.
\end{verse}

Mathematical formulas may also be displayed.  A
displayed formula
is
one-line long; multiline
formulas require special formatting instructions.
   \[  \ip{\Gamma}{\psi'} = x'' + y^{2} + z_{i}^{n}\]
Don't start a paragraph with a displayed equation,
nor make one a paragraph by itself.

\end{document}               % End of document.

\section{Bench Huge Section 740}
% synthetic bench section 740
% This is a sample LaTeX input file.  (Version of 12 August 2004.)
%
% A '%' character causes TeX to ignore all remaining text on the line,
% and is used for comments like this one.

\documentclass{article}      % Specifies the document class

                             % The preamble begins here.
\title{An Example Document}  % Declares the document's title.
\author{Leslie Lamport}      % Declares the author's name.
\date{January 21, 1994}      % Deleting this command produces today's date.

\newcommand{\ip}[2]{(#1, #2)}
                             % Defines \ip{arg1}{arg2} to mean
                             % (arg1, arg2).

%\newcommand{\ip}[2]{\langle #1 | #2\rangle}
                             % This is an alternative definition of
                             % \ip that is commented out.

\begin{document}             % End of preamble and beginning of text.

\maketitle                   % Produces the title.

This is an example input file.  Comparing it with
the output it generates can show you how to
produce a simple document of your own.

\section{Ordinary Text}      % Produces section heading.  Lower-level
                             % sections are begun with similar
                             % \subsection and \subsubsection commands.

The ends  of words and sentences are marked
  by   spaces. It  doesn't matter how many
spaces    you type; one is as good as 100.  The
end of   a line counts as a space.

One   or more   blank lines denote the  end
of  a paragraph.

Since any number of consecutive spaces are treated
like a single one, the formatting of the input
file makes no difference to
      \LaTeX,                % The \LaTeX command generates the LaTeX logo.
but it makes a difference to you.  When you use
\LaTeX, making your input file as easy to read
as possible will be a great help as you write
your document and when you change it.  This sample
file shows how you can add comments to your own input
file.

Because printing is different from typewriting,
there are a number of things that you have to do
differently when preparing an input file than if
you were just typing the document directly.
Quotation marks like
       ``this''
have to be handled specially, as do quotes within
quotes:
       ``\,`this'            % \, separates the double and single quote.
        is what I just
        wrote, not  `that'\,''.

Dashes come in three sizes: an
       intra-word
dash, a medium dash for number ranges like
       1--2,
and a punctuation
       dash---like
this.

A sentence-ending space should be larger than the
space between words within a sentence.  You
sometimes have to type special commands in
conjunction with punctuation characters to get
this right, as in the following sentence.
       Gnats, gnus, etc.\ all  % `\ ' makes an inter-word space.
       begin with G\@.         % \@ marks end-of-sentence punctuation.
You should check the spaces after periods when
reading your output to make sure you haven't
forgotten any special cases.  Generating an
ellipsis
       \ldots\               % `\ ' is needed after `\ldots' because TeX
                             % ignores spaces after command names like \ldots
                             % made from \ + letters.
                             %
                             % Note how a `%' character causes TeX to ignore
                             % the end of the input line, so these blank lines
                             % do not start a new paragraph.
                             %
with the right spacing around the periods requires
a special command.

\LaTeX\ interprets some common characters as
commands, so you must type special commands to
generate them.  These characters include the
following:
       \$ \& \% \# \{ and \}.

In printing, text is usually emphasized with an
       \emph{italic}
type style.

\begin{em}
   A long segment of text can also be emphasized
   in this way.  Text within such a segment can be
   given \emph{additional} emphasis.
\end{em}

It is sometimes necessary to prevent \LaTeX\ from
breaking a line where it might otherwise do so.
This may be at a space, as between the ``Mr.''\ and
``Jones'' in
       ``Mr.~Jones'',        % ~ produces an unbreakable interword space.
or within a word---especially when the word is a
symbol like
       \mbox{\emph{itemnum}}
that makes little sense when hyphenated across
lines.

Footnotes\footnote{This is an example of a footnote.}
pose no problem.

\LaTeX\ is good at typesetting mathematical formulas
like
       \( x-3y + z = 7 \)
or
       \( a_{1} > x^{2n} + y^{2n} > x' \)
or
       \( \ip{A}{B} = \sum_{i} a_{i} b_{i} \).
The spaces you type in a formula are
ignored.  Remember that a letter like
       $x$                   % $ ... $  and  \( ... \)  are equivalent
is a formula when it denotes a mathematical
symbol, and it should be typed as one.

\section{Displayed Text}

Text is displayed by indenting it from the left
margin.  Quotations are commonly displayed.  There
are short quotations
\begin{quote}
   This is a short quotation.  It consists of a
   single paragraph of text.  See how it is formatted.
\end{quote}
and longer ones.
\begin{quotation}
   This is a longer quotation.  It consists of two
   paragraphs of text, neither of which are
   particularly interesting.

   This is the second paragraph of the quotation.  It
   is just as dull as the first paragraph.
\end{quotation}
Another frequently-displayed structure is a list.
The following is an example of an \emph{itemized}
list.
\begin{itemize}
   \item This is the first item of an itemized list.
         Each item in the list is marked with a ``tick''.
         You don't have to worry about what kind of tick
         mark is used.

   \item This is the second item of the list.  It
         contains another list nested inside it.  The inner
         list is an \emph{enumerated} list.
         \begin{enumerate}
            \item This is the first item of an enumerated
                  list that is nested within the itemized list.

            \item This is the second item of the inner list.
                  \LaTeX\ allows you to nest lists deeper than
                  you really should.
         \end{enumerate}
         This is the rest of the second item of the outer
         list.  It is no more interesting than any other
         part of the item.
   \item This is the third item of the list.
\end{itemize}
You can even display poetry.
\begin{verse}
   There is an environment
    for verse \\             % The \\ command separates lines
   Whose features some poets % within a stanza.
   will curse.

                             % One or more blank lines separate stanzas.

   For instead of making\\
   Them do \emph{all} line breaking, \\
   It allows them to put too many words on a line when they'd rather be
   forced to be terse.
\end{verse}

Mathematical formulas may also be displayed.  A
displayed formula
is
one-line long; multiline
formulas require special formatting instructions.
   \[  \ip{\Gamma}{\psi'} = x'' + y^{2} + z_{i}^{n}\]
Don't start a paragraph with a displayed equation,
nor make one a paragraph by itself.

\end{document}               % End of document.

\section{Bench Huge Section 741}
% synthetic bench section 741
% This is a sample LaTeX input file.  (Version of 12 August 2004.)
%
% A '%' character causes TeX to ignore all remaining text on the line,
% and is used for comments like this one.

\documentclass{article}      % Specifies the document class

                             % The preamble begins here.
\title{An Example Document}  % Declares the document's title.
\author{Leslie Lamport}      % Declares the author's name.
\date{January 21, 1994}      % Deleting this command produces today's date.

\newcommand{\ip}[2]{(#1, #2)}
                             % Defines \ip{arg1}{arg2} to mean
                             % (arg1, arg2).

%\newcommand{\ip}[2]{\langle #1 | #2\rangle}
                             % This is an alternative definition of
                             % \ip that is commented out.

\begin{document}             % End of preamble and beginning of text.

\maketitle                   % Produces the title.

This is an example input file.  Comparing it with
the output it generates can show you how to
produce a simple document of your own.

\section{Ordinary Text}      % Produces section heading.  Lower-level
                             % sections are begun with similar
                             % \subsection and \subsubsection commands.

The ends  of words and sentences are marked
  by   spaces. It  doesn't matter how many
spaces    you type; one is as good as 100.  The
end of   a line counts as a space.

One   or more   blank lines denote the  end
of  a paragraph.

Since any number of consecutive spaces are treated
like a single one, the formatting of the input
file makes no difference to
      \LaTeX,                % The \LaTeX command generates the LaTeX logo.
but it makes a difference to you.  When you use
\LaTeX, making your input file as easy to read
as possible will be a great help as you write
your document and when you change it.  This sample
file shows how you can add comments to your own input
file.

Because printing is different from typewriting,
there are a number of things that you have to do
differently when preparing an input file than if
you were just typing the document directly.
Quotation marks like
       ``this''
have to be handled specially, as do quotes within
quotes:
       ``\,`this'            % \, separates the double and single quote.
        is what I just
        wrote, not  `that'\,''.

Dashes come in three sizes: an
       intra-word
dash, a medium dash for number ranges like
       1--2,
and a punctuation
       dash---like
this.

A sentence-ending space should be larger than the
space between words within a sentence.  You
sometimes have to type special commands in
conjunction with punctuation characters to get
this right, as in the following sentence.
       Gnats, gnus, etc.\ all  % `\ ' makes an inter-word space.
       begin with G\@.         % \@ marks end-of-sentence punctuation.
You should check the spaces after periods when
reading your output to make sure you haven't
forgotten any special cases.  Generating an
ellipsis
       \ldots\               % `\ ' is needed after `\ldots' because TeX
                             % ignores spaces after command names like \ldots
                             % made from \ + letters.
                             %
                             % Note how a `%' character causes TeX to ignore
                             % the end of the input line, so these blank lines
                             % do not start a new paragraph.
                             %
with the right spacing around the periods requires
a special command.

\LaTeX\ interprets some common characters as
commands, so you must type special commands to
generate them.  These characters include the
following:
       \$ \& \% \# \{ and \}.

In printing, text is usually emphasized with an
       \emph{italic}
type style.

\begin{em}
   A long segment of text can also be emphasized
   in this way.  Text within such a segment can be
   given \emph{additional} emphasis.
\end{em}

It is sometimes necessary to prevent \LaTeX\ from
breaking a line where it might otherwise do so.
This may be at a space, as between the ``Mr.''\ and
``Jones'' in
       ``Mr.~Jones'',        % ~ produces an unbreakable interword space.
or within a word---especially when the word is a
symbol like
       \mbox{\emph{itemnum}}
that makes little sense when hyphenated across
lines.

Footnotes\footnote{This is an example of a footnote.}
pose no problem.

\LaTeX\ is good at typesetting mathematical formulas
like
       \( x-3y + z = 7 \)
or
       \( a_{1} > x^{2n} + y^{2n} > x' \)
or
       \( \ip{A}{B} = \sum_{i} a_{i} b_{i} \).
The spaces you type in a formula are
ignored.  Remember that a letter like
       $x$                   % $ ... $  and  \( ... \)  are equivalent
is a formula when it denotes a mathematical
symbol, and it should be typed as one.

\section{Displayed Text}

Text is displayed by indenting it from the left
margin.  Quotations are commonly displayed.  There
are short quotations
\begin{quote}
   This is a short quotation.  It consists of a
   single paragraph of text.  See how it is formatted.
\end{quote}
and longer ones.
\begin{quotation}
   This is a longer quotation.  It consists of two
   paragraphs of text, neither of which are
   particularly interesting.

   This is the second paragraph of the quotation.  It
   is just as dull as the first paragraph.
\end{quotation}
Another frequently-displayed structure is a list.
The following is an example of an \emph{itemized}
list.
\begin{itemize}
   \item This is the first item of an itemized list.
         Each item in the list is marked with a ``tick''.
         You don't have to worry about what kind of tick
         mark is used.

   \item This is the second item of the list.  It
         contains another list nested inside it.  The inner
         list is an \emph{enumerated} list.
         \begin{enumerate}
            \item This is the first item of an enumerated
                  list that is nested within the itemized list.

            \item This is the second item of the inner list.
                  \LaTeX\ allows you to nest lists deeper than
                  you really should.
         \end{enumerate}
         This is the rest of the second item of the outer
         list.  It is no more interesting than any other
         part of the item.
   \item This is the third item of the list.
\end{itemize}
You can even display poetry.
\begin{verse}
   There is an environment
    for verse \\             % The \\ command separates lines
   Whose features some poets % within a stanza.
   will curse.

                             % One or more blank lines separate stanzas.

   For instead of making\\
   Them do \emph{all} line breaking, \\
   It allows them to put too many words on a line when they'd rather be
   forced to be terse.
\end{verse}

Mathematical formulas may also be displayed.  A
displayed formula
is
one-line long; multiline
formulas require special formatting instructions.
   \[  \ip{\Gamma}{\psi'} = x'' + y^{2} + z_{i}^{n}\]
Don't start a paragraph with a displayed equation,
nor make one a paragraph by itself.

\end{document}               % End of document.

\section{Bench Huge Section 742}
% synthetic bench section 742
% This is a sample LaTeX input file.  (Version of 12 August 2004.)
%
% A '%' character causes TeX to ignore all remaining text on the line,
% and is used for comments like this one.

\documentclass{article}      % Specifies the document class

                             % The preamble begins here.
\title{An Example Document}  % Declares the document's title.
\author{Leslie Lamport}      % Declares the author's name.
\date{January 21, 1994}      % Deleting this command produces today's date.

\newcommand{\ip}[2]{(#1, #2)}
                             % Defines \ip{arg1}{arg2} to mean
                             % (arg1, arg2).

%\newcommand{\ip}[2]{\langle #1 | #2\rangle}
                             % This is an alternative definition of
                             % \ip that is commented out.

\begin{document}             % End of preamble and beginning of text.

\maketitle                   % Produces the title.

This is an example input file.  Comparing it with
the output it generates can show you how to
produce a simple document of your own.

\section{Ordinary Text}      % Produces section heading.  Lower-level
                             % sections are begun with similar
                             % \subsection and \subsubsection commands.

The ends  of words and sentences are marked
  by   spaces. It  doesn't matter how many
spaces    you type; one is as good as 100.  The
end of   a line counts as a space.

One   or more   blank lines denote the  end
of  a paragraph.

Since any number of consecutive spaces are treated
like a single one, the formatting of the input
file makes no difference to
      \LaTeX,                % The \LaTeX command generates the LaTeX logo.
but it makes a difference to you.  When you use
\LaTeX, making your input file as easy to read
as possible will be a great help as you write
your document and when you change it.  This sample
file shows how you can add comments to your own input
file.

Because printing is different from typewriting,
there are a number of things that you have to do
differently when preparing an input file than if
you were just typing the document directly.
Quotation marks like
       ``this''
have to be handled specially, as do quotes within
quotes:
       ``\,`this'            % \, separates the double and single quote.
        is what I just
        wrote, not  `that'\,''.

Dashes come in three sizes: an
       intra-word
dash, a medium dash for number ranges like
       1--2,
and a punctuation
       dash---like
this.

A sentence-ending space should be larger than the
space between words within a sentence.  You
sometimes have to type special commands in
conjunction with punctuation characters to get
this right, as in the following sentence.
       Gnats, gnus, etc.\ all  % `\ ' makes an inter-word space.
       begin with G\@.         % \@ marks end-of-sentence punctuation.
You should check the spaces after periods when
reading your output to make sure you haven't
forgotten any special cases.  Generating an
ellipsis
       \ldots\               % `\ ' is needed after `\ldots' because TeX
                             % ignores spaces after command names like \ldots
                             % made from \ + letters.
                             %
                             % Note how a `%' character causes TeX to ignore
                             % the end of the input line, so these blank lines
                             % do not start a new paragraph.
                             %
with the right spacing around the periods requires
a special command.

\LaTeX\ interprets some common characters as
commands, so you must type special commands to
generate them.  These characters include the
following:
       \$ \& \% \# \{ and \}.

In printing, text is usually emphasized with an
       \emph{italic}
type style.

\begin{em}
   A long segment of text can also be emphasized
   in this way.  Text within such a segment can be
   given \emph{additional} emphasis.
\end{em}

It is sometimes necessary to prevent \LaTeX\ from
breaking a line where it might otherwise do so.
This may be at a space, as between the ``Mr.''\ and
``Jones'' in
       ``Mr.~Jones'',        % ~ produces an unbreakable interword space.
or within a word---especially when the word is a
symbol like
       \mbox{\emph{itemnum}}
that makes little sense when hyphenated across
lines.

Footnotes\footnote{This is an example of a footnote.}
pose no problem.

\LaTeX\ is good at typesetting mathematical formulas
like
       \( x-3y + z = 7 \)
or
       \( a_{1} > x^{2n} + y^{2n} > x' \)
or
       \( \ip{A}{B} = \sum_{i} a_{i} b_{i} \).
The spaces you type in a formula are
ignored.  Remember that a letter like
       $x$                   % $ ... $  and  \( ... \)  are equivalent
is a formula when it denotes a mathematical
symbol, and it should be typed as one.

\section{Displayed Text}

Text is displayed by indenting it from the left
margin.  Quotations are commonly displayed.  There
are short quotations
\begin{quote}
   This is a short quotation.  It consists of a
   single paragraph of text.  See how it is formatted.
\end{quote}
and longer ones.
\begin{quotation}
   This is a longer quotation.  It consists of two
   paragraphs of text, neither of which are
   particularly interesting.

   This is the second paragraph of the quotation.  It
   is just as dull as the first paragraph.
\end{quotation}
Another frequently-displayed structure is a list.
The following is an example of an \emph{itemized}
list.
\begin{itemize}
   \item This is the first item of an itemized list.
         Each item in the list is marked with a ``tick''.
         You don't have to worry about what kind of tick
         mark is used.

   \item This is the second item of the list.  It
         contains another list nested inside it.  The inner
         list is an \emph{enumerated} list.
         \begin{enumerate}
            \item This is the first item of an enumerated
                  list that is nested within the itemized list.

            \item This is the second item of the inner list.
                  \LaTeX\ allows you to nest lists deeper than
                  you really should.
         \end{enumerate}
         This is the rest of the second item of the outer
         list.  It is no more interesting than any other
         part of the item.
   \item This is the third item of the list.
\end{itemize}
You can even display poetry.
\begin{verse}
   There is an environment
    for verse \\             % The \\ command separates lines
   Whose features some poets % within a stanza.
   will curse.

                             % One or more blank lines separate stanzas.

   For instead of making\\
   Them do \emph{all} line breaking, \\
   It allows them to put too many words on a line when they'd rather be
   forced to be terse.
\end{verse}

Mathematical formulas may also be displayed.  A
displayed formula
is
one-line long; multiline
formulas require special formatting instructions.
   \[  \ip{\Gamma}{\psi'} = x'' + y^{2} + z_{i}^{n}\]
Don't start a paragraph with a displayed equation,
nor make one a paragraph by itself.

\end{document}               % End of document.

\section{Bench Huge Section 743}
% synthetic bench section 743
% This is a sample LaTeX input file.  (Version of 12 August 2004.)
%
% A '%' character causes TeX to ignore all remaining text on the line,
% and is used for comments like this one.

\documentclass{article}      % Specifies the document class

                             % The preamble begins here.
\title{An Example Document}  % Declares the document's title.
\author{Leslie Lamport}      % Declares the author's name.
\date{January 21, 1994}      % Deleting this command produces today's date.

\newcommand{\ip}[2]{(#1, #2)}
                             % Defines \ip{arg1}{arg2} to mean
                             % (arg1, arg2).

%\newcommand{\ip}[2]{\langle #1 | #2\rangle}
                             % This is an alternative definition of
                             % \ip that is commented out.

\begin{document}             % End of preamble and beginning of text.

\maketitle                   % Produces the title.

This is an example input file.  Comparing it with
the output it generates can show you how to
produce a simple document of your own.

\section{Ordinary Text}      % Produces section heading.  Lower-level
                             % sections are begun with similar
                             % \subsection and \subsubsection commands.

The ends  of words and sentences are marked
  by   spaces. It  doesn't matter how many
spaces    you type; one is as good as 100.  The
end of   a line counts as a space.

One   or more   blank lines denote the  end
of  a paragraph.

Since any number of consecutive spaces are treated
like a single one, the formatting of the input
file makes no difference to
      \LaTeX,                % The \LaTeX command generates the LaTeX logo.
but it makes a difference to you.  When you use
\LaTeX, making your input file as easy to read
as possible will be a great help as you write
your document and when you change it.  This sample
file shows how you can add comments to your own input
file.

Because printing is different from typewriting,
there are a number of things that you have to do
differently when preparing an input file than if
you were just typing the document directly.
Quotation marks like
       ``this''
have to be handled specially, as do quotes within
quotes:
       ``\,`this'            % \, separates the double and single quote.
        is what I just
        wrote, not  `that'\,''.

Dashes come in three sizes: an
       intra-word
dash, a medium dash for number ranges like
       1--2,
and a punctuation
       dash---like
this.

A sentence-ending space should be larger than the
space between words within a sentence.  You
sometimes have to type special commands in
conjunction with punctuation characters to get
this right, as in the following sentence.
       Gnats, gnus, etc.\ all  % `\ ' makes an inter-word space.
       begin with G\@.         % \@ marks end-of-sentence punctuation.
You should check the spaces after periods when
reading your output to make sure you haven't
forgotten any special cases.  Generating an
ellipsis
       \ldots\               % `\ ' is needed after `\ldots' because TeX
                             % ignores spaces after command names like \ldots
                             % made from \ + letters.
                             %
                             % Note how a `%' character causes TeX to ignore
                             % the end of the input line, so these blank lines
                             % do not start a new paragraph.
                             %
with the right spacing around the periods requires
a special command.

\LaTeX\ interprets some common characters as
commands, so you must type special commands to
generate them.  These characters include the
following:
       \$ \& \% \# \{ and \}.

In printing, text is usually emphasized with an
       \emph{italic}
type style.

\begin{em}
   A long segment of text can also be emphasized
   in this way.  Text within such a segment can be
   given \emph{additional} emphasis.
\end{em}

It is sometimes necessary to prevent \LaTeX\ from
breaking a line where it might otherwise do so.
This may be at a space, as between the ``Mr.''\ and
``Jones'' in
       ``Mr.~Jones'',        % ~ produces an unbreakable interword space.
or within a word---especially when the word is a
symbol like
       \mbox{\emph{itemnum}}
that makes little sense when hyphenated across
lines.

Footnotes\footnote{This is an example of a footnote.}
pose no problem.

\LaTeX\ is good at typesetting mathematical formulas
like
       \( x-3y + z = 7 \)
or
       \( a_{1} > x^{2n} + y^{2n} > x' \)
or
       \( \ip{A}{B} = \sum_{i} a_{i} b_{i} \).
The spaces you type in a formula are
ignored.  Remember that a letter like
       $x$                   % $ ... $  and  \( ... \)  are equivalent
is a formula when it denotes a mathematical
symbol, and it should be typed as one.

\section{Displayed Text}

Text is displayed by indenting it from the left
margin.  Quotations are commonly displayed.  There
are short quotations
\begin{quote}
   This is a short quotation.  It consists of a
   single paragraph of text.  See how it is formatted.
\end{quote}
and longer ones.
\begin{quotation}
   This is a longer quotation.  It consists of two
   paragraphs of text, neither of which are
   particularly interesting.

   This is the second paragraph of the quotation.  It
   is just as dull as the first paragraph.
\end{quotation}
Another frequently-displayed structure is a list.
The following is an example of an \emph{itemized}
list.
\begin{itemize}
   \item This is the first item of an itemized list.
         Each item in the list is marked with a ``tick''.
         You don't have to worry about what kind of tick
         mark is used.

   \item This is the second item of the list.  It
         contains another list nested inside it.  The inner
         list is an \emph{enumerated} list.
         \begin{enumerate}
            \item This is the first item of an enumerated
                  list that is nested within the itemized list.

            \item This is the second item of the inner list.
                  \LaTeX\ allows you to nest lists deeper than
                  you really should.
         \end{enumerate}
         This is the rest of the second item of the outer
         list.  It is no more interesting than any other
         part of the item.
   \item This is the third item of the list.
\end{itemize}
You can even display poetry.
\begin{verse}
   There is an environment
    for verse \\             % The \\ command separates lines
   Whose features some poets % within a stanza.
   will curse.

                             % One or more blank lines separate stanzas.

   For instead of making\\
   Them do \emph{all} line breaking, \\
   It allows them to put too many words on a line when they'd rather be
   forced to be terse.
\end{verse}

Mathematical formulas may also be displayed.  A
displayed formula
is
one-line long; multiline
formulas require special formatting instructions.
   \[  \ip{\Gamma}{\psi'} = x'' + y^{2} + z_{i}^{n}\]
Don't start a paragraph with a displayed equation,
nor make one a paragraph by itself.

\end{document}               % End of document.

\section{Bench Huge Section 744}
% synthetic bench section 744
% This is a sample LaTeX input file.  (Version of 12 August 2004.)
%
% A '%' character causes TeX to ignore all remaining text on the line,
% and is used for comments like this one.

\documentclass{article}      % Specifies the document class

                             % The preamble begins here.
\title{An Example Document}  % Declares the document's title.
\author{Leslie Lamport}      % Declares the author's name.
\date{January 21, 1994}      % Deleting this command produces today's date.

\newcommand{\ip}[2]{(#1, #2)}
                             % Defines \ip{arg1}{arg2} to mean
                             % (arg1, arg2).

%\newcommand{\ip}[2]{\langle #1 | #2\rangle}
                             % This is an alternative definition of
                             % \ip that is commented out.

\begin{document}             % End of preamble and beginning of text.

\maketitle                   % Produces the title.

This is an example input file.  Comparing it with
the output it generates can show you how to
produce a simple document of your own.

\section{Ordinary Text}      % Produces section heading.  Lower-level
                             % sections are begun with similar
                             % \subsection and \subsubsection commands.

The ends  of words and sentences are marked
  by   spaces. It  doesn't matter how many
spaces    you type; one is as good as 100.  The
end of   a line counts as a space.

One   or more   blank lines denote the  end
of  a paragraph.

Since any number of consecutive spaces are treated
like a single one, the formatting of the input
file makes no difference to
      \LaTeX,                % The \LaTeX command generates the LaTeX logo.
but it makes a difference to you.  When you use
\LaTeX, making your input file as easy to read
as possible will be a great help as you write
your document and when you change it.  This sample
file shows how you can add comments to your own input
file.

Because printing is different from typewriting,
there are a number of things that you have to do
differently when preparing an input file than if
you were just typing the document directly.
Quotation marks like
       ``this''
have to be handled specially, as do quotes within
quotes:
       ``\,`this'            % \, separates the double and single quote.
        is what I just
        wrote, not  `that'\,''.

Dashes come in three sizes: an
       intra-word
dash, a medium dash for number ranges like
       1--2,
and a punctuation
       dash---like
this.

A sentence-ending space should be larger than the
space between words within a sentence.  You
sometimes have to type special commands in
conjunction with punctuation characters to get
this right, as in the following sentence.
       Gnats, gnus, etc.\ all  % `\ ' makes an inter-word space.
       begin with G\@.         % \@ marks end-of-sentence punctuation.
You should check the spaces after periods when
reading your output to make sure you haven't
forgotten any special cases.  Generating an
ellipsis
       \ldots\               % `\ ' is needed after `\ldots' because TeX
                             % ignores spaces after command names like \ldots
                             % made from \ + letters.
                             %
                             % Note how a `%' character causes TeX to ignore
                             % the end of the input line, so these blank lines
                             % do not start a new paragraph.
                             %
with the right spacing around the periods requires
a special command.

\LaTeX\ interprets some common characters as
commands, so you must type special commands to
generate them.  These characters include the
following:
       \$ \& \% \# \{ and \}.

In printing, text is usually emphasized with an
       \emph{italic}
type style.

\begin{em}
   A long segment of text can also be emphasized
   in this way.  Text within such a segment can be
   given \emph{additional} emphasis.
\end{em}

It is sometimes necessary to prevent \LaTeX\ from
breaking a line where it might otherwise do so.
This may be at a space, as between the ``Mr.''\ and
``Jones'' in
       ``Mr.~Jones'',        % ~ produces an unbreakable interword space.
or within a word---especially when the word is a
symbol like
       \mbox{\emph{itemnum}}
that makes little sense when hyphenated across
lines.

Footnotes\footnote{This is an example of a footnote.}
pose no problem.

\LaTeX\ is good at typesetting mathematical formulas
like
       \( x-3y + z = 7 \)
or
       \( a_{1} > x^{2n} + y^{2n} > x' \)
or
       \( \ip{A}{B} = \sum_{i} a_{i} b_{i} \).
The spaces you type in a formula are
ignored.  Remember that a letter like
       $x$                   % $ ... $  and  \( ... \)  are equivalent
is a formula when it denotes a mathematical
symbol, and it should be typed as one.

\section{Displayed Text}

Text is displayed by indenting it from the left
margin.  Quotations are commonly displayed.  There
are short quotations
\begin{quote}
   This is a short quotation.  It consists of a
   single paragraph of text.  See how it is formatted.
\end{quote}
and longer ones.
\begin{quotation}
   This is a longer quotation.  It consists of two
   paragraphs of text, neither of which are
   particularly interesting.

   This is the second paragraph of the quotation.  It
   is just as dull as the first paragraph.
\end{quotation}
Another frequently-displayed structure is a list.
The following is an example of an \emph{itemized}
list.
\begin{itemize}
   \item This is the first item of an itemized list.
         Each item in the list is marked with a ``tick''.
         You don't have to worry about what kind of tick
         mark is used.

   \item This is the second item of the list.  It
         contains another list nested inside it.  The inner
         list is an \emph{enumerated} list.
         \begin{enumerate}
            \item This is the first item of an enumerated
                  list that is nested within the itemized list.

            \item This is the second item of the inner list.
                  \LaTeX\ allows you to nest lists deeper than
                  you really should.
         \end{enumerate}
         This is the rest of the second item of the outer
         list.  It is no more interesting than any other
         part of the item.
   \item This is the third item of the list.
\end{itemize}
You can even display poetry.
\begin{verse}
   There is an environment
    for verse \\             % The \\ command separates lines
   Whose features some poets % within a stanza.
   will curse.

                             % One or more blank lines separate stanzas.

   For instead of making\\
   Them do \emph{all} line breaking, \\
   It allows them to put too many words on a line when they'd rather be
   forced to be terse.
\end{verse}

Mathematical formulas may also be displayed.  A
displayed formula
is
one-line long; multiline
formulas require special formatting instructions.
   \[  \ip{\Gamma}{\psi'} = x'' + y^{2} + z_{i}^{n}\]
Don't start a paragraph with a displayed equation,
nor make one a paragraph by itself.

\end{document}               % End of document.

\section{Bench Huge Section 745}
% synthetic bench section 745
% This is a sample LaTeX input file.  (Version of 12 August 2004.)
%
% A '%' character causes TeX to ignore all remaining text on the line,
% and is used for comments like this one.

\documentclass{article}      % Specifies the document class

                             % The preamble begins here.
\title{An Example Document}  % Declares the document's title.
\author{Leslie Lamport}      % Declares the author's name.
\date{January 21, 1994}      % Deleting this command produces today's date.

\newcommand{\ip}[2]{(#1, #2)}
                             % Defines \ip{arg1}{arg2} to mean
                             % (arg1, arg2).

%\newcommand{\ip}[2]{\langle #1 | #2\rangle}
                             % This is an alternative definition of
                             % \ip that is commented out.

\begin{document}             % End of preamble and beginning of text.

\maketitle                   % Produces the title.

This is an example input file.  Comparing it with
the output it generates can show you how to
produce a simple document of your own.

\section{Ordinary Text}      % Produces section heading.  Lower-level
                             % sections are begun with similar
                             % \subsection and \subsubsection commands.

The ends  of words and sentences are marked
  by   spaces. It  doesn't matter how many
spaces    you type; one is as good as 100.  The
end of   a line counts as a space.

One   or more   blank lines denote the  end
of  a paragraph.

Since any number of consecutive spaces are treated
like a single one, the formatting of the input
file makes no difference to
      \LaTeX,                % The \LaTeX command generates the LaTeX logo.
but it makes a difference to you.  When you use
\LaTeX, making your input file as easy to read
as possible will be a great help as you write
your document and when you change it.  This sample
file shows how you can add comments to your own input
file.

Because printing is different from typewriting,
there are a number of things that you have to do
differently when preparing an input file than if
you were just typing the document directly.
Quotation marks like
       ``this''
have to be handled specially, as do quotes within
quotes:
       ``\,`this'            % \, separates the double and single quote.
        is what I just
        wrote, not  `that'\,''.

Dashes come in three sizes: an
       intra-word
dash, a medium dash for number ranges like
       1--2,
and a punctuation
       dash---like
this.

A sentence-ending space should be larger than the
space between words within a sentence.  You
sometimes have to type special commands in
conjunction with punctuation characters to get
this right, as in the following sentence.
       Gnats, gnus, etc.\ all  % `\ ' makes an inter-word space.
       begin with G\@.         % \@ marks end-of-sentence punctuation.
You should check the spaces after periods when
reading your output to make sure you haven't
forgotten any special cases.  Generating an
ellipsis
       \ldots\               % `\ ' is needed after `\ldots' because TeX
                             % ignores spaces after command names like \ldots
                             % made from \ + letters.
                             %
                             % Note how a `%' character causes TeX to ignore
                             % the end of the input line, so these blank lines
                             % do not start a new paragraph.
                             %
with the right spacing around the periods requires
a special command.

\LaTeX\ interprets some common characters as
commands, so you must type special commands to
generate them.  These characters include the
following:
       \$ \& \% \# \{ and \}.

In printing, text is usually emphasized with an
       \emph{italic}
type style.

\begin{em}
   A long segment of text can also be emphasized
   in this way.  Text within such a segment can be
   given \emph{additional} emphasis.
\end{em}

It is sometimes necessary to prevent \LaTeX\ from
breaking a line where it might otherwise do so.
This may be at a space, as between the ``Mr.''\ and
``Jones'' in
       ``Mr.~Jones'',        % ~ produces an unbreakable interword space.
or within a word---especially when the word is a
symbol like
       \mbox{\emph{itemnum}}
that makes little sense when hyphenated across
lines.

Footnotes\footnote{This is an example of a footnote.}
pose no problem.

\LaTeX\ is good at typesetting mathematical formulas
like
       \( x-3y + z = 7 \)
or
       \( a_{1} > x^{2n} + y^{2n} > x' \)
or
       \( \ip{A}{B} = \sum_{i} a_{i} b_{i} \).
The spaces you type in a formula are
ignored.  Remember that a letter like
       $x$                   % $ ... $  and  \( ... \)  are equivalent
is a formula when it denotes a mathematical
symbol, and it should be typed as one.

\section{Displayed Text}

Text is displayed by indenting it from the left
margin.  Quotations are commonly displayed.  There
are short quotations
\begin{quote}
   This is a short quotation.  It consists of a
   single paragraph of text.  See how it is formatted.
\end{quote}
and longer ones.
\begin{quotation}
   This is a longer quotation.  It consists of two
   paragraphs of text, neither of which are
   particularly interesting.

   This is the second paragraph of the quotation.  It
   is just as dull as the first paragraph.
\end{quotation}
Another frequently-displayed structure is a list.
The following is an example of an \emph{itemized}
list.
\begin{itemize}
   \item This is the first item of an itemized list.
         Each item in the list is marked with a ``tick''.
         You don't have to worry about what kind of tick
         mark is used.

   \item This is the second item of the list.  It
         contains another list nested inside it.  The inner
         list is an \emph{enumerated} list.
         \begin{enumerate}
            \item This is the first item of an enumerated
                  list that is nested within the itemized list.

            \item This is the second item of the inner list.
                  \LaTeX\ allows you to nest lists deeper than
                  you really should.
         \end{enumerate}
         This is the rest of the second item of the outer
         list.  It is no more interesting than any other
         part of the item.
   \item This is the third item of the list.
\end{itemize}
You can even display poetry.
\begin{verse}
   There is an environment
    for verse \\             % The \\ command separates lines
   Whose features some poets % within a stanza.
   will curse.

                             % One or more blank lines separate stanzas.

   For instead of making\\
   Them do \emph{all} line breaking, \\
   It allows them to put too many words on a line when they'd rather be
   forced to be terse.
\end{verse}

Mathematical formulas may also be displayed.  A
displayed formula
is
one-line long; multiline
formulas require special formatting instructions.
   \[  \ip{\Gamma}{\psi'} = x'' + y^{2} + z_{i}^{n}\]
Don't start a paragraph with a displayed equation,
nor make one a paragraph by itself.

\end{document}               % End of document.

\section{Bench Huge Section 746}
% synthetic bench section 746
% This is a sample LaTeX input file.  (Version of 12 August 2004.)
%
% A '%' character causes TeX to ignore all remaining text on the line,
% and is used for comments like this one.

\documentclass{article}      % Specifies the document class

                             % The preamble begins here.
\title{An Example Document}  % Declares the document's title.
\author{Leslie Lamport}      % Declares the author's name.
\date{January 21, 1994}      % Deleting this command produces today's date.

\newcommand{\ip}[2]{(#1, #2)}
                             % Defines \ip{arg1}{arg2} to mean
                             % (arg1, arg2).

%\newcommand{\ip}[2]{\langle #1 | #2\rangle}
                             % This is an alternative definition of
                             % \ip that is commented out.

\begin{document}             % End of preamble and beginning of text.

\maketitle                   % Produces the title.

This is an example input file.  Comparing it with
the output it generates can show you how to
produce a simple document of your own.

\section{Ordinary Text}      % Produces section heading.  Lower-level
                             % sections are begun with similar
                             % \subsection and \subsubsection commands.

The ends  of words and sentences are marked
  by   spaces. It  doesn't matter how many
spaces    you type; one is as good as 100.  The
end of   a line counts as a space.

One   or more   blank lines denote the  end
of  a paragraph.

Since any number of consecutive spaces are treated
like a single one, the formatting of the input
file makes no difference to
      \LaTeX,                % The \LaTeX command generates the LaTeX logo.
but it makes a difference to you.  When you use
\LaTeX, making your input file as easy to read
as possible will be a great help as you write
your document and when you change it.  This sample
file shows how you can add comments to your own input
file.

Because printing is different from typewriting,
there are a number of things that you have to do
differently when preparing an input file than if
you were just typing the document directly.
Quotation marks like
       ``this''
have to be handled specially, as do quotes within
quotes:
       ``\,`this'            % \, separates the double and single quote.
        is what I just
        wrote, not  `that'\,''.

Dashes come in three sizes: an
       intra-word
dash, a medium dash for number ranges like
       1--2,
and a punctuation
       dash---like
this.

A sentence-ending space should be larger than the
space between words within a sentence.  You
sometimes have to type special commands in
conjunction with punctuation characters to get
this right, as in the following sentence.
       Gnats, gnus, etc.\ all  % `\ ' makes an inter-word space.
       begin with G\@.         % \@ marks end-of-sentence punctuation.
You should check the spaces after periods when
reading your output to make sure you haven't
forgotten any special cases.  Generating an
ellipsis
       \ldots\               % `\ ' is needed after `\ldots' because TeX
                             % ignores spaces after command names like \ldots
                             % made from \ + letters.
                             %
                             % Note how a `%' character causes TeX to ignore
                             % the end of the input line, so these blank lines
                             % do not start a new paragraph.
                             %
with the right spacing around the periods requires
a special command.

\LaTeX\ interprets some common characters as
commands, so you must type special commands to
generate them.  These characters include the
following:
       \$ \& \% \# \{ and \}.

In printing, text is usually emphasized with an
       \emph{italic}
type style.

\begin{em}
   A long segment of text can also be emphasized
   in this way.  Text within such a segment can be
   given \emph{additional} emphasis.
\end{em}

It is sometimes necessary to prevent \LaTeX\ from
breaking a line where it might otherwise do so.
This may be at a space, as between the ``Mr.''\ and
``Jones'' in
       ``Mr.~Jones'',        % ~ produces an unbreakable interword space.
or within a word---especially when the word is a
symbol like
       \mbox{\emph{itemnum}}
that makes little sense when hyphenated across
lines.

Footnotes\footnote{This is an example of a footnote.}
pose no problem.

\LaTeX\ is good at typesetting mathematical formulas
like
       \( x-3y + z = 7 \)
or
       \( a_{1} > x^{2n} + y^{2n} > x' \)
or
       \( \ip{A}{B} = \sum_{i} a_{i} b_{i} \).
The spaces you type in a formula are
ignored.  Remember that a letter like
       $x$                   % $ ... $  and  \( ... \)  are equivalent
is a formula when it denotes a mathematical
symbol, and it should be typed as one.

\section{Displayed Text}

Text is displayed by indenting it from the left
margin.  Quotations are commonly displayed.  There
are short quotations
\begin{quote}
   This is a short quotation.  It consists of a
   single paragraph of text.  See how it is formatted.
\end{quote}
and longer ones.
\begin{quotation}
   This is a longer quotation.  It consists of two
   paragraphs of text, neither of which are
   particularly interesting.

   This is the second paragraph of the quotation.  It
   is just as dull as the first paragraph.
\end{quotation}
Another frequently-displayed structure is a list.
The following is an example of an \emph{itemized}
list.
\begin{itemize}
   \item This is the first item of an itemized list.
         Each item in the list is marked with a ``tick''.
         You don't have to worry about what kind of tick
         mark is used.

   \item This is the second item of the list.  It
         contains another list nested inside it.  The inner
         list is an \emph{enumerated} list.
         \begin{enumerate}
            \item This is the first item of an enumerated
                  list that is nested within the itemized list.

            \item This is the second item of the inner list.
                  \LaTeX\ allows you to nest lists deeper than
                  you really should.
         \end{enumerate}
         This is the rest of the second item of the outer
         list.  It is no more interesting than any other
         part of the item.
   \item This is the third item of the list.
\end{itemize}
You can even display poetry.
\begin{verse}
   There is an environment
    for verse \\             % The \\ command separates lines
   Whose features some poets % within a stanza.
   will curse.

                             % One or more blank lines separate stanzas.

   For instead of making\\
   Them do \emph{all} line breaking, \\
   It allows them to put too many words on a line when they'd rather be
   forced to be terse.
\end{verse}

Mathematical formulas may also be displayed.  A
displayed formula
is
one-line long; multiline
formulas require special formatting instructions.
   \[  \ip{\Gamma}{\psi'} = x'' + y^{2} + z_{i}^{n}\]
Don't start a paragraph with a displayed equation,
nor make one a paragraph by itself.

\end{document}               % End of document.

\section{Bench Huge Section 747}
% synthetic bench section 747
% This is a sample LaTeX input file.  (Version of 12 August 2004.)
%
% A '%' character causes TeX to ignore all remaining text on the line,
% and is used for comments like this one.

\documentclass{article}      % Specifies the document class

                             % The preamble begins here.
\title{An Example Document}  % Declares the document's title.
\author{Leslie Lamport}      % Declares the author's name.
\date{January 21, 1994}      % Deleting this command produces today's date.

\newcommand{\ip}[2]{(#1, #2)}
                             % Defines \ip{arg1}{arg2} to mean
                             % (arg1, arg2).

%\newcommand{\ip}[2]{\langle #1 | #2\rangle}
                             % This is an alternative definition of
                             % \ip that is commented out.

\begin{document}             % End of preamble and beginning of text.

\maketitle                   % Produces the title.

This is an example input file.  Comparing it with
the output it generates can show you how to
produce a simple document of your own.

\section{Ordinary Text}      % Produces section heading.  Lower-level
                             % sections are begun with similar
                             % \subsection and \subsubsection commands.

The ends  of words and sentences are marked
  by   spaces. It  doesn't matter how many
spaces    you type; one is as good as 100.  The
end of   a line counts as a space.

One   or more   blank lines denote the  end
of  a paragraph.

Since any number of consecutive spaces are treated
like a single one, the formatting of the input
file makes no difference to
      \LaTeX,                % The \LaTeX command generates the LaTeX logo.
but it makes a difference to you.  When you use
\LaTeX, making your input file as easy to read
as possible will be a great help as you write
your document and when you change it.  This sample
file shows how you can add comments to your own input
file.

Because printing is different from typewriting,
there are a number of things that you have to do
differently when preparing an input file than if
you were just typing the document directly.
Quotation marks like
       ``this''
have to be handled specially, as do quotes within
quotes:
       ``\,`this'            % \, separates the double and single quote.
        is what I just
        wrote, not  `that'\,''.

Dashes come in three sizes: an
       intra-word
dash, a medium dash for number ranges like
       1--2,
and a punctuation
       dash---like
this.

A sentence-ending space should be larger than the
space between words within a sentence.  You
sometimes have to type special commands in
conjunction with punctuation characters to get
this right, as in the following sentence.
       Gnats, gnus, etc.\ all  % `\ ' makes an inter-word space.
       begin with G\@.         % \@ marks end-of-sentence punctuation.
You should check the spaces after periods when
reading your output to make sure you haven't
forgotten any special cases.  Generating an
ellipsis
       \ldots\               % `\ ' is needed after `\ldots' because TeX
                             % ignores spaces after command names like \ldots
                             % made from \ + letters.
                             %
                             % Note how a `%' character causes TeX to ignore
                             % the end of the input line, so these blank lines
                             % do not start a new paragraph.
                             %
with the right spacing around the periods requires
a special command.

\LaTeX\ interprets some common characters as
commands, so you must type special commands to
generate them.  These characters include the
following:
       \$ \& \% \# \{ and \}.

In printing, text is usually emphasized with an
       \emph{italic}
type style.

\begin{em}
   A long segment of text can also be emphasized
   in this way.  Text within such a segment can be
   given \emph{additional} emphasis.
\end{em}

It is sometimes necessary to prevent \LaTeX\ from
breaking a line where it might otherwise do so.
This may be at a space, as between the ``Mr.''\ and
``Jones'' in
       ``Mr.~Jones'',        % ~ produces an unbreakable interword space.
or within a word---especially when the word is a
symbol like
       \mbox{\emph{itemnum}}
that makes little sense when hyphenated across
lines.

Footnotes\footnote{This is an example of a footnote.}
pose no problem.

\LaTeX\ is good at typesetting mathematical formulas
like
       \( x-3y + z = 7 \)
or
       \( a_{1} > x^{2n} + y^{2n} > x' \)
or
       \( \ip{A}{B} = \sum_{i} a_{i} b_{i} \).
The spaces you type in a formula are
ignored.  Remember that a letter like
       $x$                   % $ ... $  and  \( ... \)  are equivalent
is a formula when it denotes a mathematical
symbol, and it should be typed as one.

\section{Displayed Text}

Text is displayed by indenting it from the left
margin.  Quotations are commonly displayed.  There
are short quotations
\begin{quote}
   This is a short quotation.  It consists of a
   single paragraph of text.  See how it is formatted.
\end{quote}
and longer ones.
\begin{quotation}
   This is a longer quotation.  It consists of two
   paragraphs of text, neither of which are
   particularly interesting.

   This is the second paragraph of the quotation.  It
   is just as dull as the first paragraph.
\end{quotation}
Another frequently-displayed structure is a list.
The following is an example of an \emph{itemized}
list.
\begin{itemize}
   \item This is the first item of an itemized list.
         Each item in the list is marked with a ``tick''.
         You don't have to worry about what kind of tick
         mark is used.

   \item This is the second item of the list.  It
         contains another list nested inside it.  The inner
         list is an \emph{enumerated} list.
         \begin{enumerate}
            \item This is the first item of an enumerated
                  list that is nested within the itemized list.

            \item This is the second item of the inner list.
                  \LaTeX\ allows you to nest lists deeper than
                  you really should.
         \end{enumerate}
         This is the rest of the second item of the outer
         list.  It is no more interesting than any other
         part of the item.
   \item This is the third item of the list.
\end{itemize}
You can even display poetry.
\begin{verse}
   There is an environment
    for verse \\             % The \\ command separates lines
   Whose features some poets % within a stanza.
   will curse.

                             % One or more blank lines separate stanzas.

   For instead of making\\
   Them do \emph{all} line breaking, \\
   It allows them to put too many words on a line when they'd rather be
   forced to be terse.
\end{verse}

Mathematical formulas may also be displayed.  A
displayed formula
is
one-line long; multiline
formulas require special formatting instructions.
   \[  \ip{\Gamma}{\psi'} = x'' + y^{2} + z_{i}^{n}\]
Don't start a paragraph with a displayed equation,
nor make one a paragraph by itself.

\end{document}               % End of document.

\section{Bench Huge Section 748}
% synthetic bench section 748
% This is a sample LaTeX input file.  (Version of 12 August 2004.)
%
% A '%' character causes TeX to ignore all remaining text on the line,
% and is used for comments like this one.

\documentclass{article}      % Specifies the document class

                             % The preamble begins here.
\title{An Example Document}  % Declares the document's title.
\author{Leslie Lamport}      % Declares the author's name.
\date{January 21, 1994}      % Deleting this command produces today's date.

\newcommand{\ip}[2]{(#1, #2)}
                             % Defines \ip{arg1}{arg2} to mean
                             % (arg1, arg2).

%\newcommand{\ip}[2]{\langle #1 | #2\rangle}
                             % This is an alternative definition of
                             % \ip that is commented out.

\begin{document}             % End of preamble and beginning of text.

\maketitle                   % Produces the title.

This is an example input file.  Comparing it with
the output it generates can show you how to
produce a simple document of your own.

\section{Ordinary Text}      % Produces section heading.  Lower-level
                             % sections are begun with similar
                             % \subsection and \subsubsection commands.

The ends  of words and sentences are marked
  by   spaces. It  doesn't matter how many
spaces    you type; one is as good as 100.  The
end of   a line counts as a space.

One   or more   blank lines denote the  end
of  a paragraph.

Since any number of consecutive spaces are treated
like a single one, the formatting of the input
file makes no difference to
      \LaTeX,                % The \LaTeX command generates the LaTeX logo.
but it makes a difference to you.  When you use
\LaTeX, making your input file as easy to read
as possible will be a great help as you write
your document and when you change it.  This sample
file shows how you can add comments to your own input
file.

Because printing is different from typewriting,
there are a number of things that you have to do
differently when preparing an input file than if
you were just typing the document directly.
Quotation marks like
       ``this''
have to be handled specially, as do quotes within
quotes:
       ``\,`this'            % \, separates the double and single quote.
        is what I just
        wrote, not  `that'\,''.

Dashes come in three sizes: an
       intra-word
dash, a medium dash for number ranges like
       1--2,
and a punctuation
       dash---like
this.

A sentence-ending space should be larger than the
space between words within a sentence.  You
sometimes have to type special commands in
conjunction with punctuation characters to get
this right, as in the following sentence.
       Gnats, gnus, etc.\ all  % `\ ' makes an inter-word space.
       begin with G\@.         % \@ marks end-of-sentence punctuation.
You should check the spaces after periods when
reading your output to make sure you haven't
forgotten any special cases.  Generating an
ellipsis
       \ldots\               % `\ ' is needed after `\ldots' because TeX
                             % ignores spaces after command names like \ldots
                             % made from \ + letters.
                             %
                             % Note how a `%' character causes TeX to ignore
                             % the end of the input line, so these blank lines
                             % do not start a new paragraph.
                             %
with the right spacing around the periods requires
a special command.

\LaTeX\ interprets some common characters as
commands, so you must type special commands to
generate them.  These characters include the
following:
       \$ \& \% \# \{ and \}.

In printing, text is usually emphasized with an
       \emph{italic}
type style.

\begin{em}
   A long segment of text can also be emphasized
   in this way.  Text within such a segment can be
   given \emph{additional} emphasis.
\end{em}

It is sometimes necessary to prevent \LaTeX\ from
breaking a line where it might otherwise do so.
This may be at a space, as between the ``Mr.''\ and
``Jones'' in
       ``Mr.~Jones'',        % ~ produces an unbreakable interword space.
or within a word---especially when the word is a
symbol like
       \mbox{\emph{itemnum}}
that makes little sense when hyphenated across
lines.

Footnotes\footnote{This is an example of a footnote.}
pose no problem.

\LaTeX\ is good at typesetting mathematical formulas
like
       \( x-3y + z = 7 \)
or
       \( a_{1} > x^{2n} + y^{2n} > x' \)
or
       \( \ip{A}{B} = \sum_{i} a_{i} b_{i} \).
The spaces you type in a formula are
ignored.  Remember that a letter like
       $x$                   % $ ... $  and  \( ... \)  are equivalent
is a formula when it denotes a mathematical
symbol, and it should be typed as one.

\section{Displayed Text}

Text is displayed by indenting it from the left
margin.  Quotations are commonly displayed.  There
are short quotations
\begin{quote}
   This is a short quotation.  It consists of a
   single paragraph of text.  See how it is formatted.
\end{quote}
and longer ones.
\begin{quotation}
   This is a longer quotation.  It consists of two
   paragraphs of text, neither of which are
   particularly interesting.

   This is the second paragraph of the quotation.  It
   is just as dull as the first paragraph.
\end{quotation}
Another frequently-displayed structure is a list.
The following is an example of an \emph{itemized}
list.
\begin{itemize}
   \item This is the first item of an itemized list.
         Each item in the list is marked with a ``tick''.
         You don't have to worry about what kind of tick
         mark is used.

   \item This is the second item of the list.  It
         contains another list nested inside it.  The inner
         list is an \emph{enumerated} list.
         \begin{enumerate}
            \item This is the first item of an enumerated
                  list that is nested within the itemized list.

            \item This is the second item of the inner list.
                  \LaTeX\ allows you to nest lists deeper than
                  you really should.
         \end{enumerate}
         This is the rest of the second item of the outer
         list.  It is no more interesting than any other
         part of the item.
   \item This is the third item of the list.
\end{itemize}
You can even display poetry.
\begin{verse}
   There is an environment
    for verse \\             % The \\ command separates lines
   Whose features some poets % within a stanza.
   will curse.

                             % One or more blank lines separate stanzas.

   For instead of making\\
   Them do \emph{all} line breaking, \\
   It allows them to put too many words on a line when they'd rather be
   forced to be terse.
\end{verse}

Mathematical formulas may also be displayed.  A
displayed formula
is
one-line long; multiline
formulas require special formatting instructions.
   \[  \ip{\Gamma}{\psi'} = x'' + y^{2} + z_{i}^{n}\]
Don't start a paragraph with a displayed equation,
nor make one a paragraph by itself.

\end{document}               % End of document.

\section{Bench Huge Section 749}
% synthetic bench section 749
% This is a sample LaTeX input file.  (Version of 12 August 2004.)
%
% A '%' character causes TeX to ignore all remaining text on the line,
% and is used for comments like this one.

\documentclass{article}      % Specifies the document class

                             % The preamble begins here.
\title{An Example Document}  % Declares the document's title.
\author{Leslie Lamport}      % Declares the author's name.
\date{January 21, 1994}      % Deleting this command produces today's date.

\newcommand{\ip}[2]{(#1, #2)}
                             % Defines \ip{arg1}{arg2} to mean
                             % (arg1, arg2).

%\newcommand{\ip}[2]{\langle #1 | #2\rangle}
                             % This is an alternative definition of
                             % \ip that is commented out.

\begin{document}             % End of preamble and beginning of text.

\maketitle                   % Produces the title.

This is an example input file.  Comparing it with
the output it generates can show you how to
produce a simple document of your own.

\section{Ordinary Text}      % Produces section heading.  Lower-level
                             % sections are begun with similar
                             % \subsection and \subsubsection commands.

The ends  of words and sentences are marked
  by   spaces. It  doesn't matter how many
spaces    you type; one is as good as 100.  The
end of   a line counts as a space.

One   or more   blank lines denote the  end
of  a paragraph.

Since any number of consecutive spaces are treated
like a single one, the formatting of the input
file makes no difference to
      \LaTeX,                % The \LaTeX command generates the LaTeX logo.
but it makes a difference to you.  When you use
\LaTeX, making your input file as easy to read
as possible will be a great help as you write
your document and when you change it.  This sample
file shows how you can add comments to your own input
file.

Because printing is different from typewriting,
there are a number of things that you have to do
differently when preparing an input file than if
you were just typing the document directly.
Quotation marks like
       ``this''
have to be handled specially, as do quotes within
quotes:
       ``\,`this'            % \, separates the double and single quote.
        is what I just
        wrote, not  `that'\,''.

Dashes come in three sizes: an
       intra-word
dash, a medium dash for number ranges like
       1--2,
and a punctuation
       dash---like
this.

A sentence-ending space should be larger than the
space between words within a sentence.  You
sometimes have to type special commands in
conjunction with punctuation characters to get
this right, as in the following sentence.
       Gnats, gnus, etc.\ all  % `\ ' makes an inter-word space.
       begin with G\@.         % \@ marks end-of-sentence punctuation.
You should check the spaces after periods when
reading your output to make sure you haven't
forgotten any special cases.  Generating an
ellipsis
       \ldots\               % `\ ' is needed after `\ldots' because TeX
                             % ignores spaces after command names like \ldots
                             % made from \ + letters.
                             %
                             % Note how a `%' character causes TeX to ignore
                             % the end of the input line, so these blank lines
                             % do not start a new paragraph.
                             %
with the right spacing around the periods requires
a special command.

\LaTeX\ interprets some common characters as
commands, so you must type special commands to
generate them.  These characters include the
following:
       \$ \& \% \# \{ and \}.

In printing, text is usually emphasized with an
       \emph{italic}
type style.

\begin{em}
   A long segment of text can also be emphasized
   in this way.  Text within such a segment can be
   given \emph{additional} emphasis.
\end{em}

It is sometimes necessary to prevent \LaTeX\ from
breaking a line where it might otherwise do so.
This may be at a space, as between the ``Mr.''\ and
``Jones'' in
       ``Mr.~Jones'',        % ~ produces an unbreakable interword space.
or within a word---especially when the word is a
symbol like
       \mbox{\emph{itemnum}}
that makes little sense when hyphenated across
lines.

Footnotes\footnote{This is an example of a footnote.}
pose no problem.

\LaTeX\ is good at typesetting mathematical formulas
like
       \( x-3y + z = 7 \)
or
       \( a_{1} > x^{2n} + y^{2n} > x' \)
or
       \( \ip{A}{B} = \sum_{i} a_{i} b_{i} \).
The spaces you type in a formula are
ignored.  Remember that a letter like
       $x$                   % $ ... $  and  \( ... \)  are equivalent
is a formula when it denotes a mathematical
symbol, and it should be typed as one.

\section{Displayed Text}

Text is displayed by indenting it from the left
margin.  Quotations are commonly displayed.  There
are short quotations
\begin{quote}
   This is a short quotation.  It consists of a
   single paragraph of text.  See how it is formatted.
\end{quote}
and longer ones.
\begin{quotation}
   This is a longer quotation.  It consists of two
   paragraphs of text, neither of which are
   particularly interesting.

   This is the second paragraph of the quotation.  It
   is just as dull as the first paragraph.
\end{quotation}
Another frequently-displayed structure is a list.
The following is an example of an \emph{itemized}
list.
\begin{itemize}
   \item This is the first item of an itemized list.
         Each item in the list is marked with a ``tick''.
         You don't have to worry about what kind of tick
         mark is used.

   \item This is the second item of the list.  It
         contains another list nested inside it.  The inner
         list is an \emph{enumerated} list.
         \begin{enumerate}
            \item This is the first item of an enumerated
                  list that is nested within the itemized list.

            \item This is the second item of the inner list.
                  \LaTeX\ allows you to nest lists deeper than
                  you really should.
         \end{enumerate}
         This is the rest of the second item of the outer
         list.  It is no more interesting than any other
         part of the item.
   \item This is the third item of the list.
\end{itemize}
You can even display poetry.
\begin{verse}
   There is an environment
    for verse \\             % The \\ command separates lines
   Whose features some poets % within a stanza.
   will curse.

                             % One or more blank lines separate stanzas.

   For instead of making\\
   Them do \emph{all} line breaking, \\
   It allows them to put too many words on a line when they'd rather be
   forced to be terse.
\end{verse}

Mathematical formulas may also be displayed.  A
displayed formula
is
one-line long; multiline
formulas require special formatting instructions.
   \[  \ip{\Gamma}{\psi'} = x'' + y^{2} + z_{i}^{n}\]
Don't start a paragraph with a displayed equation,
nor make one a paragraph by itself.

\end{document}               % End of document.

\section{Bench Huge Section 750}
% synthetic bench section 750
% This is a sample LaTeX input file.  (Version of 12 August 2004.)
%
% A '%' character causes TeX to ignore all remaining text on the line,
% and is used for comments like this one.

\documentclass{article}      % Specifies the document class

                             % The preamble begins here.
\title{An Example Document}  % Declares the document's title.
\author{Leslie Lamport}      % Declares the author's name.
\date{January 21, 1994}      % Deleting this command produces today's date.

\newcommand{\ip}[2]{(#1, #2)}
                             % Defines \ip{arg1}{arg2} to mean
                             % (arg1, arg2).

%\newcommand{\ip}[2]{\langle #1 | #2\rangle}
                             % This is an alternative definition of
                             % \ip that is commented out.

\begin{document}             % End of preamble and beginning of text.

\maketitle                   % Produces the title.

This is an example input file.  Comparing it with
the output it generates can show you how to
produce a simple document of your own.

\section{Ordinary Text}      % Produces section heading.  Lower-level
                             % sections are begun with similar
                             % \subsection and \subsubsection commands.

The ends  of words and sentences are marked
  by   spaces. It  doesn't matter how many
spaces    you type; one is as good as 100.  The
end of   a line counts as a space.

One   or more   blank lines denote the  end
of  a paragraph.

Since any number of consecutive spaces are treated
like a single one, the formatting of the input
file makes no difference to
      \LaTeX,                % The \LaTeX command generates the LaTeX logo.
but it makes a difference to you.  When you use
\LaTeX, making your input file as easy to read
as possible will be a great help as you write
your document and when you change it.  This sample
file shows how you can add comments to your own input
file.

Because printing is different from typewriting,
there are a number of things that you have to do
differently when preparing an input file than if
you were just typing the document directly.
Quotation marks like
       ``this''
have to be handled specially, as do quotes within
quotes:
       ``\,`this'            % \, separates the double and single quote.
        is what I just
        wrote, not  `that'\,''.

Dashes come in three sizes: an
       intra-word
dash, a medium dash for number ranges like
       1--2,
and a punctuation
       dash---like
this.

A sentence-ending space should be larger than the
space between words within a sentence.  You
sometimes have to type special commands in
conjunction with punctuation characters to get
this right, as in the following sentence.
       Gnats, gnus, etc.\ all  % `\ ' makes an inter-word space.
       begin with G\@.         % \@ marks end-of-sentence punctuation.
You should check the spaces after periods when
reading your output to make sure you haven't
forgotten any special cases.  Generating an
ellipsis
       \ldots\               % `\ ' is needed after `\ldots' because TeX
                             % ignores spaces after command names like \ldots
                             % made from \ + letters.
                             %
                             % Note how a `%' character causes TeX to ignore
                             % the end of the input line, so these blank lines
                             % do not start a new paragraph.
                             %
with the right spacing around the periods requires
a special command.

\LaTeX\ interprets some common characters as
commands, so you must type special commands to
generate them.  These characters include the
following:
       \$ \& \% \# \{ and \}.

In printing, text is usually emphasized with an
       \emph{italic}
type style.

\begin{em}
   A long segment of text can also be emphasized
   in this way.  Text within such a segment can be
   given \emph{additional} emphasis.
\end{em}

It is sometimes necessary to prevent \LaTeX\ from
breaking a line where it might otherwise do so.
This may be at a space, as between the ``Mr.''\ and
``Jones'' in
       ``Mr.~Jones'',        % ~ produces an unbreakable interword space.
or within a word---especially when the word is a
symbol like
       \mbox{\emph{itemnum}}
that makes little sense when hyphenated across
lines.

Footnotes\footnote{This is an example of a footnote.}
pose no problem.

\LaTeX\ is good at typesetting mathematical formulas
like
       \( x-3y + z = 7 \)
or
       \( a_{1} > x^{2n} + y^{2n} > x' \)
or
       \( \ip{A}{B} = \sum_{i} a_{i} b_{i} \).
The spaces you type in a formula are
ignored.  Remember that a letter like
       $x$                   % $ ... $  and  \( ... \)  are equivalent
is a formula when it denotes a mathematical
symbol, and it should be typed as one.

\section{Displayed Text}

Text is displayed by indenting it from the left
margin.  Quotations are commonly displayed.  There
are short quotations
\begin{quote}
   This is a short quotation.  It consists of a
   single paragraph of text.  See how it is formatted.
\end{quote}
and longer ones.
\begin{quotation}
   This is a longer quotation.  It consists of two
   paragraphs of text, neither of which are
   particularly interesting.

   This is the second paragraph of the quotation.  It
   is just as dull as the first paragraph.
\end{quotation}
Another frequently-displayed structure is a list.
The following is an example of an \emph{itemized}
list.
\begin{itemize}
   \item This is the first item of an itemized list.
         Each item in the list is marked with a ``tick''.
         You don't have to worry about what kind of tick
         mark is used.

   \item This is the second item of the list.  It
         contains another list nested inside it.  The inner
         list is an \emph{enumerated} list.
         \begin{enumerate}
            \item This is the first item of an enumerated
                  list that is nested within the itemized list.

            \item This is the second item of the inner list.
                  \LaTeX\ allows you to nest lists deeper than
                  you really should.
         \end{enumerate}
         This is the rest of the second item of the outer
         list.  It is no more interesting than any other
         part of the item.
   \item This is the third item of the list.
\end{itemize}
You can even display poetry.
\begin{verse}
   There is an environment
    for verse \\             % The \\ command separates lines
   Whose features some poets % within a stanza.
   will curse.

                             % One or more blank lines separate stanzas.

   For instead of making\\
   Them do \emph{all} line breaking, \\
   It allows them to put too many words on a line when they'd rather be
   forced to be terse.
\end{verse}

Mathematical formulas may also be displayed.  A
displayed formula
is
one-line long; multiline
formulas require special formatting instructions.
   \[  \ip{\Gamma}{\psi'} = x'' + y^{2} + z_{i}^{n}\]
Don't start a paragraph with a displayed equation,
nor make one a paragraph by itself.

\end{document}               % End of document.

\section{Bench Huge Section 751}
% synthetic bench section 751
% This is a sample LaTeX input file.  (Version of 12 August 2004.)
%
% A '%' character causes TeX to ignore all remaining text on the line,
% and is used for comments like this one.

\documentclass{article}      % Specifies the document class

                             % The preamble begins here.
\title{An Example Document}  % Declares the document's title.
\author{Leslie Lamport}      % Declares the author's name.
\date{January 21, 1994}      % Deleting this command produces today's date.

\newcommand{\ip}[2]{(#1, #2)}
                             % Defines \ip{arg1}{arg2} to mean
                             % (arg1, arg2).

%\newcommand{\ip}[2]{\langle #1 | #2\rangle}
                             % This is an alternative definition of
                             % \ip that is commented out.

\begin{document}             % End of preamble and beginning of text.

\maketitle                   % Produces the title.

This is an example input file.  Comparing it with
the output it generates can show you how to
produce a simple document of your own.

\section{Ordinary Text}      % Produces section heading.  Lower-level
                             % sections are begun with similar
                             % \subsection and \subsubsection commands.

The ends  of words and sentences are marked
  by   spaces. It  doesn't matter how many
spaces    you type; one is as good as 100.  The
end of   a line counts as a space.

One   or more   blank lines denote the  end
of  a paragraph.

Since any number of consecutive spaces are treated
like a single one, the formatting of the input
file makes no difference to
      \LaTeX,                % The \LaTeX command generates the LaTeX logo.
but it makes a difference to you.  When you use
\LaTeX, making your input file as easy to read
as possible will be a great help as you write
your document and when you change it.  This sample
file shows how you can add comments to your own input
file.

Because printing is different from typewriting,
there are a number of things that you have to do
differently when preparing an input file than if
you were just typing the document directly.
Quotation marks like
       ``this''
have to be handled specially, as do quotes within
quotes:
       ``\,`this'            % \, separates the double and single quote.
        is what I just
        wrote, not  `that'\,''.

Dashes come in three sizes: an
       intra-word
dash, a medium dash for number ranges like
       1--2,
and a punctuation
       dash---like
this.

A sentence-ending space should be larger than the
space between words within a sentence.  You
sometimes have to type special commands in
conjunction with punctuation characters to get
this right, as in the following sentence.
       Gnats, gnus, etc.\ all  % `\ ' makes an inter-word space.
       begin with G\@.         % \@ marks end-of-sentence punctuation.
You should check the spaces after periods when
reading your output to make sure you haven't
forgotten any special cases.  Generating an
ellipsis
       \ldots\               % `\ ' is needed after `\ldots' because TeX
                             % ignores spaces after command names like \ldots
                             % made from \ + letters.
                             %
                             % Note how a `%' character causes TeX to ignore
                             % the end of the input line, so these blank lines
                             % do not start a new paragraph.
                             %
with the right spacing around the periods requires
a special command.

\LaTeX\ interprets some common characters as
commands, so you must type special commands to
generate them.  These characters include the
following:
       \$ \& \% \# \{ and \}.

In printing, text is usually emphasized with an
       \emph{italic}
type style.

\begin{em}
   A long segment of text can also be emphasized
   in this way.  Text within such a segment can be
   given \emph{additional} emphasis.
\end{em}

It is sometimes necessary to prevent \LaTeX\ from
breaking a line where it might otherwise do so.
This may be at a space, as between the ``Mr.''\ and
``Jones'' in
       ``Mr.~Jones'',        % ~ produces an unbreakable interword space.
or within a word---especially when the word is a
symbol like
       \mbox{\emph{itemnum}}
that makes little sense when hyphenated across
lines.

Footnotes\footnote{This is an example of a footnote.}
pose no problem.

\LaTeX\ is good at typesetting mathematical formulas
like
       \( x-3y + z = 7 \)
or
       \( a_{1} > x^{2n} + y^{2n} > x' \)
or
       \( \ip{A}{B} = \sum_{i} a_{i} b_{i} \).
The spaces you type in a formula are
ignored.  Remember that a letter like
       $x$                   % $ ... $  and  \( ... \)  are equivalent
is a formula when it denotes a mathematical
symbol, and it should be typed as one.

\section{Displayed Text}

Text is displayed by indenting it from the left
margin.  Quotations are commonly displayed.  There
are short quotations
\begin{quote}
   This is a short quotation.  It consists of a
   single paragraph of text.  See how it is formatted.
\end{quote}
and longer ones.
\begin{quotation}
   This is a longer quotation.  It consists of two
   paragraphs of text, neither of which are
   particularly interesting.

   This is the second paragraph of the quotation.  It
   is just as dull as the first paragraph.
\end{quotation}
Another frequently-displayed structure is a list.
The following is an example of an \emph{itemized}
list.
\begin{itemize}
   \item This is the first item of an itemized list.
         Each item in the list is marked with a ``tick''.
         You don't have to worry about what kind of tick
         mark is used.

   \item This is the second item of the list.  It
         contains another list nested inside it.  The inner
         list is an \emph{enumerated} list.
         \begin{enumerate}
            \item This is the first item of an enumerated
                  list that is nested within the itemized list.

            \item This is the second item of the inner list.
                  \LaTeX\ allows you to nest lists deeper than
                  you really should.
         \end{enumerate}
         This is the rest of the second item of the outer
         list.  It is no more interesting than any other
         part of the item.
   \item This is the third item of the list.
\end{itemize}
You can even display poetry.
\begin{verse}
   There is an environment
    for verse \\             % The \\ command separates lines
   Whose features some poets % within a stanza.
   will curse.

                             % One or more blank lines separate stanzas.

   For instead of making\\
   Them do \emph{all} line breaking, \\
   It allows them to put too many words on a line when they'd rather be
   forced to be terse.
\end{verse}

Mathematical formulas may also be displayed.  A
displayed formula
is
one-line long; multiline
formulas require special formatting instructions.
   \[  \ip{\Gamma}{\psi'} = x'' + y^{2} + z_{i}^{n}\]
Don't start a paragraph with a displayed equation,
nor make one a paragraph by itself.

\end{document}               % End of document.

\section{Bench Huge Section 752}
% synthetic bench section 752
% This is a sample LaTeX input file.  (Version of 12 August 2004.)
%
% A '%' character causes TeX to ignore all remaining text on the line,
% and is used for comments like this one.

\documentclass{article}      % Specifies the document class

                             % The preamble begins here.
\title{An Example Document}  % Declares the document's title.
\author{Leslie Lamport}      % Declares the author's name.
\date{January 21, 1994}      % Deleting this command produces today's date.

\newcommand{\ip}[2]{(#1, #2)}
                             % Defines \ip{arg1}{arg2} to mean
                             % (arg1, arg2).

%\newcommand{\ip}[2]{\langle #1 | #2\rangle}
                             % This is an alternative definition of
                             % \ip that is commented out.

\begin{document}             % End of preamble and beginning of text.

\maketitle                   % Produces the title.

This is an example input file.  Comparing it with
the output it generates can show you how to
produce a simple document of your own.

\section{Ordinary Text}      % Produces section heading.  Lower-level
                             % sections are begun with similar
                             % \subsection and \subsubsection commands.

The ends  of words and sentences are marked
  by   spaces. It  doesn't matter how many
spaces    you type; one is as good as 100.  The
end of   a line counts as a space.

One   or more   blank lines denote the  end
of  a paragraph.

Since any number of consecutive spaces are treated
like a single one, the formatting of the input
file makes no difference to
      \LaTeX,                % The \LaTeX command generates the LaTeX logo.
but it makes a difference to you.  When you use
\LaTeX, making your input file as easy to read
as possible will be a great help as you write
your document and when you change it.  This sample
file shows how you can add comments to your own input
file.

Because printing is different from typewriting,
there are a number of things that you have to do
differently when preparing an input file than if
you were just typing the document directly.
Quotation marks like
       ``this''
have to be handled specially, as do quotes within
quotes:
       ``\,`this'            % \, separates the double and single quote.
        is what I just
        wrote, not  `that'\,''.

Dashes come in three sizes: an
       intra-word
dash, a medium dash for number ranges like
       1--2,
and a punctuation
       dash---like
this.

A sentence-ending space should be larger than the
space between words within a sentence.  You
sometimes have to type special commands in
conjunction with punctuation characters to get
this right, as in the following sentence.
       Gnats, gnus, etc.\ all  % `\ ' makes an inter-word space.
       begin with G\@.         % \@ marks end-of-sentence punctuation.
You should check the spaces after periods when
reading your output to make sure you haven't
forgotten any special cases.  Generating an
ellipsis
       \ldots\               % `\ ' is needed after `\ldots' because TeX
                             % ignores spaces after command names like \ldots
                             % made from \ + letters.
                             %
                             % Note how a `%' character causes TeX to ignore
                             % the end of the input line, so these blank lines
                             % do not start a new paragraph.
                             %
with the right spacing around the periods requires
a special command.

\LaTeX\ interprets some common characters as
commands, so you must type special commands to
generate them.  These characters include the
following:
       \$ \& \% \# \{ and \}.

In printing, text is usually emphasized with an
       \emph{italic}
type style.

\begin{em}
   A long segment of text can also be emphasized
   in this way.  Text within such a segment can be
   given \emph{additional} emphasis.
\end{em}

It is sometimes necessary to prevent \LaTeX\ from
breaking a line where it might otherwise do so.
This may be at a space, as between the ``Mr.''\ and
``Jones'' in
       ``Mr.~Jones'',        % ~ produces an unbreakable interword space.
or within a word---especially when the word is a
symbol like
       \mbox{\emph{itemnum}}
that makes little sense when hyphenated across
lines.

Footnotes\footnote{This is an example of a footnote.}
pose no problem.

\LaTeX\ is good at typesetting mathematical formulas
like
       \( x-3y + z = 7 \)
or
       \( a_{1} > x^{2n} + y^{2n} > x' \)
or
       \( \ip{A}{B} = \sum_{i} a_{i} b_{i} \).
The spaces you type in a formula are
ignored.  Remember that a letter like
       $x$                   % $ ... $  and  \( ... \)  are equivalent
is a formula when it denotes a mathematical
symbol, and it should be typed as one.

\section{Displayed Text}

Text is displayed by indenting it from the left
margin.  Quotations are commonly displayed.  There
are short quotations
\begin{quote}
   This is a short quotation.  It consists of a
   single paragraph of text.  See how it is formatted.
\end{quote}
and longer ones.
\begin{quotation}
   This is a longer quotation.  It consists of two
   paragraphs of text, neither of which are
   particularly interesting.

   This is the second paragraph of the quotation.  It
   is just as dull as the first paragraph.
\end{quotation}
Another frequently-displayed structure is a list.
The following is an example of an \emph{itemized}
list.
\begin{itemize}
   \item This is the first item of an itemized list.
         Each item in the list is marked with a ``tick''.
         You don't have to worry about what kind of tick
         mark is used.

   \item This is the second item of the list.  It
         contains another list nested inside it.  The inner
         list is an \emph{enumerated} list.
         \begin{enumerate}
            \item This is the first item of an enumerated
                  list that is nested within the itemized list.

            \item This is the second item of the inner list.
                  \LaTeX\ allows you to nest lists deeper than
                  you really should.
         \end{enumerate}
         This is the rest of the second item of the outer
         list.  It is no more interesting than any other
         part of the item.
   \item This is the third item of the list.
\end{itemize}
You can even display poetry.
\begin{verse}
   There is an environment
    for verse \\             % The \\ command separates lines
   Whose features some poets % within a stanza.
   will curse.

                             % One or more blank lines separate stanzas.

   For instead of making\\
   Them do \emph{all} line breaking, \\
   It allows them to put too many words on a line when they'd rather be
   forced to be terse.
\end{verse}

Mathematical formulas may also be displayed.  A
displayed formula
is
one-line long; multiline
formulas require special formatting instructions.
   \[  \ip{\Gamma}{\psi'} = x'' + y^{2} + z_{i}^{n}\]
Don't start a paragraph with a displayed equation,
nor make one a paragraph by itself.

\end{document}               % End of document.

\section{Bench Huge Section 753}
% synthetic bench section 753
% This is a sample LaTeX input file.  (Version of 12 August 2004.)
%
% A '%' character causes TeX to ignore all remaining text on the line,
% and is used for comments like this one.

\documentclass{article}      % Specifies the document class

                             % The preamble begins here.
\title{An Example Document}  % Declares the document's title.
\author{Leslie Lamport}      % Declares the author's name.
\date{January 21, 1994}      % Deleting this command produces today's date.

\newcommand{\ip}[2]{(#1, #2)}
                             % Defines \ip{arg1}{arg2} to mean
                             % (arg1, arg2).

%\newcommand{\ip}[2]{\langle #1 | #2\rangle}
                             % This is an alternative definition of
                             % \ip that is commented out.

\begin{document}             % End of preamble and beginning of text.

\maketitle                   % Produces the title.

This is an example input file.  Comparing it with
the output it generates can show you how to
produce a simple document of your own.

\section{Ordinary Text}      % Produces section heading.  Lower-level
                             % sections are begun with similar
                             % \subsection and \subsubsection commands.

The ends  of words and sentences are marked
  by   spaces. It  doesn't matter how many
spaces    you type; one is as good as 100.  The
end of   a line counts as a space.

One   or more   blank lines denote the  end
of  a paragraph.

Since any number of consecutive spaces are treated
like a single one, the formatting of the input
file makes no difference to
      \LaTeX,                % The \LaTeX command generates the LaTeX logo.
but it makes a difference to you.  When you use
\LaTeX, making your input file as easy to read
as possible will be a great help as you write
your document and when you change it.  This sample
file shows how you can add comments to your own input
file.

Because printing is different from typewriting,
there are a number of things that you have to do
differently when preparing an input file than if
you were just typing the document directly.
Quotation marks like
       ``this''
have to be handled specially, as do quotes within
quotes:
       ``\,`this'            % \, separates the double and single quote.
        is what I just
        wrote, not  `that'\,''.

Dashes come in three sizes: an
       intra-word
dash, a medium dash for number ranges like
       1--2,
and a punctuation
       dash---like
this.

A sentence-ending space should be larger than the
space between words within a sentence.  You
sometimes have to type special commands in
conjunction with punctuation characters to get
this right, as in the following sentence.
       Gnats, gnus, etc.\ all  % `\ ' makes an inter-word space.
       begin with G\@.         % \@ marks end-of-sentence punctuation.
You should check the spaces after periods when
reading your output to make sure you haven't
forgotten any special cases.  Generating an
ellipsis
       \ldots\               % `\ ' is needed after `\ldots' because TeX
                             % ignores spaces after command names like \ldots
                             % made from \ + letters.
                             %
                             % Note how a `%' character causes TeX to ignore
                             % the end of the input line, so these blank lines
                             % do not start a new paragraph.
                             %
with the right spacing around the periods requires
a special command.

\LaTeX\ interprets some common characters as
commands, so you must type special commands to
generate them.  These characters include the
following:
       \$ \& \% \# \{ and \}.

In printing, text is usually emphasized with an
       \emph{italic}
type style.

\begin{em}
   A long segment of text can also be emphasized
   in this way.  Text within such a segment can be
   given \emph{additional} emphasis.
\end{em}

It is sometimes necessary to prevent \LaTeX\ from
breaking a line where it might otherwise do so.
This may be at a space, as between the ``Mr.''\ and
``Jones'' in
       ``Mr.~Jones'',        % ~ produces an unbreakable interword space.
or within a word---especially when the word is a
symbol like
       \mbox{\emph{itemnum}}
that makes little sense when hyphenated across
lines.

Footnotes\footnote{This is an example of a footnote.}
pose no problem.

\LaTeX\ is good at typesetting mathematical formulas
like
       \( x-3y + z = 7 \)
or
       \( a_{1} > x^{2n} + y^{2n} > x' \)
or
       \( \ip{A}{B} = \sum_{i} a_{i} b_{i} \).
The spaces you type in a formula are
ignored.  Remember that a letter like
       $x$                   % $ ... $  and  \( ... \)  are equivalent
is a formula when it denotes a mathematical
symbol, and it should be typed as one.

\section{Displayed Text}

Text is displayed by indenting it from the left
margin.  Quotations are commonly displayed.  There
are short quotations
\begin{quote}
   This is a short quotation.  It consists of a
   single paragraph of text.  See how it is formatted.
\end{quote}
and longer ones.
\begin{quotation}
   This is a longer quotation.  It consists of two
   paragraphs of text, neither of which are
   particularly interesting.

   This is the second paragraph of the quotation.  It
   is just as dull as the first paragraph.
\end{quotation}
Another frequently-displayed structure is a list.
The following is an example of an \emph{itemized}
list.
\begin{itemize}
   \item This is the first item of an itemized list.
         Each item in the list is marked with a ``tick''.
         You don't have to worry about what kind of tick
         mark is used.

   \item This is the second item of the list.  It
         contains another list nested inside it.  The inner
         list is an \emph{enumerated} list.
         \begin{enumerate}
            \item This is the first item of an enumerated
                  list that is nested within the itemized list.

            \item This is the second item of the inner list.
                  \LaTeX\ allows you to nest lists deeper than
                  you really should.
         \end{enumerate}
         This is the rest of the second item of the outer
         list.  It is no more interesting than any other
         part of the item.
   \item This is the third item of the list.
\end{itemize}
You can even display poetry.
\begin{verse}
   There is an environment
    for verse \\             % The \\ command separates lines
   Whose features some poets % within a stanza.
   will curse.

                             % One or more blank lines separate stanzas.

   For instead of making\\
   Them do \emph{all} line breaking, \\
   It allows them to put too many words on a line when they'd rather be
   forced to be terse.
\end{verse}

Mathematical formulas may also be displayed.  A
displayed formula
is
one-line long; multiline
formulas require special formatting instructions.
   \[  \ip{\Gamma}{\psi'} = x'' + y^{2} + z_{i}^{n}\]
Don't start a paragraph with a displayed equation,
nor make one a paragraph by itself.

\end{document}               % End of document.

\section{Bench Huge Section 754}
% synthetic bench section 754
% This is a sample LaTeX input file.  (Version of 12 August 2004.)
%
% A '%' character causes TeX to ignore all remaining text on the line,
% and is used for comments like this one.

\documentclass{article}      % Specifies the document class

                             % The preamble begins here.
\title{An Example Document}  % Declares the document's title.
\author{Leslie Lamport}      % Declares the author's name.
\date{January 21, 1994}      % Deleting this command produces today's date.

\newcommand{\ip}[2]{(#1, #2)}
                             % Defines \ip{arg1}{arg2} to mean
                             % (arg1, arg2).

%\newcommand{\ip}[2]{\langle #1 | #2\rangle}
                             % This is an alternative definition of
                             % \ip that is commented out.

\begin{document}             % End of preamble and beginning of text.

\maketitle                   % Produces the title.

This is an example input file.  Comparing it with
the output it generates can show you how to
produce a simple document of your own.

\section{Ordinary Text}      % Produces section heading.  Lower-level
                             % sections are begun with similar
                             % \subsection and \subsubsection commands.

The ends  of words and sentences are marked
  by   spaces. It  doesn't matter how many
spaces    you type; one is as good as 100.  The
end of   a line counts as a space.

One   or more   blank lines denote the  end
of  a paragraph.

Since any number of consecutive spaces are treated
like a single one, the formatting of the input
file makes no difference to
      \LaTeX,                % The \LaTeX command generates the LaTeX logo.
but it makes a difference to you.  When you use
\LaTeX, making your input file as easy to read
as possible will be a great help as you write
your document and when you change it.  This sample
file shows how you can add comments to your own input
file.

Because printing is different from typewriting,
there are a number of things that you have to do
differently when preparing an input file than if
you were just typing the document directly.
Quotation marks like
       ``this''
have to be handled specially, as do quotes within
quotes:
       ``\,`this'            % \, separates the double and single quote.
        is what I just
        wrote, not  `that'\,''.

Dashes come in three sizes: an
       intra-word
dash, a medium dash for number ranges like
       1--2,
and a punctuation
       dash---like
this.

A sentence-ending space should be larger than the
space between words within a sentence.  You
sometimes have to type special commands in
conjunction with punctuation characters to get
this right, as in the following sentence.
       Gnats, gnus, etc.\ all  % `\ ' makes an inter-word space.
       begin with G\@.         % \@ marks end-of-sentence punctuation.
You should check the spaces after periods when
reading your output to make sure you haven't
forgotten any special cases.  Generating an
ellipsis
       \ldots\               % `\ ' is needed after `\ldots' because TeX
                             % ignores spaces after command names like \ldots
                             % made from \ + letters.
                             %
                             % Note how a `%' character causes TeX to ignore
                             % the end of the input line, so these blank lines
                             % do not start a new paragraph.
                             %
with the right spacing around the periods requires
a special command.

\LaTeX\ interprets some common characters as
commands, so you must type special commands to
generate them.  These characters include the
following:
       \$ \& \% \# \{ and \}.

In printing, text is usually emphasized with an
       \emph{italic}
type style.

\begin{em}
   A long segment of text can also be emphasized
   in this way.  Text within such a segment can be
   given \emph{additional} emphasis.
\end{em}

It is sometimes necessary to prevent \LaTeX\ from
breaking a line where it might otherwise do so.
This may be at a space, as between the ``Mr.''\ and
``Jones'' in
       ``Mr.~Jones'',        % ~ produces an unbreakable interword space.
or within a word---especially when the word is a
symbol like
       \mbox{\emph{itemnum}}
that makes little sense when hyphenated across
lines.

Footnotes\footnote{This is an example of a footnote.}
pose no problem.

\LaTeX\ is good at typesetting mathematical formulas
like
       \( x-3y + z = 7 \)
or
       \( a_{1} > x^{2n} + y^{2n} > x' \)
or
       \( \ip{A}{B} = \sum_{i} a_{i} b_{i} \).
The spaces you type in a formula are
ignored.  Remember that a letter like
       $x$                   % $ ... $  and  \( ... \)  are equivalent
is a formula when it denotes a mathematical
symbol, and it should be typed as one.

\section{Displayed Text}

Text is displayed by indenting it from the left
margin.  Quotations are commonly displayed.  There
are short quotations
\begin{quote}
   This is a short quotation.  It consists of a
   single paragraph of text.  See how it is formatted.
\end{quote}
and longer ones.
\begin{quotation}
   This is a longer quotation.  It consists of two
   paragraphs of text, neither of which are
   particularly interesting.

   This is the second paragraph of the quotation.  It
   is just as dull as the first paragraph.
\end{quotation}
Another frequently-displayed structure is a list.
The following is an example of an \emph{itemized}
list.
\begin{itemize}
   \item This is the first item of an itemized list.
         Each item in the list is marked with a ``tick''.
         You don't have to worry about what kind of tick
         mark is used.

   \item This is the second item of the list.  It
         contains another list nested inside it.  The inner
         list is an \emph{enumerated} list.
         \begin{enumerate}
            \item This is the first item of an enumerated
                  list that is nested within the itemized list.

            \item This is the second item of the inner list.
                  \LaTeX\ allows you to nest lists deeper than
                  you really should.
         \end{enumerate}
         This is the rest of the second item of the outer
         list.  It is no more interesting than any other
         part of the item.
   \item This is the third item of the list.
\end{itemize}
You can even display poetry.
\begin{verse}
   There is an environment
    for verse \\             % The \\ command separates lines
   Whose features some poets % within a stanza.
   will curse.

                             % One or more blank lines separate stanzas.

   For instead of making\\
   Them do \emph{all} line breaking, \\
   It allows them to put too many words on a line when they'd rather be
   forced to be terse.
\end{verse}

Mathematical formulas may also be displayed.  A
displayed formula
is
one-line long; multiline
formulas require special formatting instructions.
   \[  \ip{\Gamma}{\psi'} = x'' + y^{2} + z_{i}^{n}\]
Don't start a paragraph with a displayed equation,
nor make one a paragraph by itself.

\end{document}               % End of document.

\section{Bench Huge Section 755}
% synthetic bench section 755
% This is a sample LaTeX input file.  (Version of 12 August 2004.)
%
% A '%' character causes TeX to ignore all remaining text on the line,
% and is used for comments like this one.

\documentclass{article}      % Specifies the document class

                             % The preamble begins here.
\title{An Example Document}  % Declares the document's title.
\author{Leslie Lamport}      % Declares the author's name.
\date{January 21, 1994}      % Deleting this command produces today's date.

\newcommand{\ip}[2]{(#1, #2)}
                             % Defines \ip{arg1}{arg2} to mean
                             % (arg1, arg2).

%\newcommand{\ip}[2]{\langle #1 | #2\rangle}
                             % This is an alternative definition of
                             % \ip that is commented out.

\begin{document}             % End of preamble and beginning of text.

\maketitle                   % Produces the title.

This is an example input file.  Comparing it with
the output it generates can show you how to
produce a simple document of your own.

\section{Ordinary Text}      % Produces section heading.  Lower-level
                             % sections are begun with similar
                             % \subsection and \subsubsection commands.

The ends  of words and sentences are marked
  by   spaces. It  doesn't matter how many
spaces    you type; one is as good as 100.  The
end of   a line counts as a space.

One   or more   blank lines denote the  end
of  a paragraph.

Since any number of consecutive spaces are treated
like a single one, the formatting of the input
file makes no difference to
      \LaTeX,                % The \LaTeX command generates the LaTeX logo.
but it makes a difference to you.  When you use
\LaTeX, making your input file as easy to read
as possible will be a great help as you write
your document and when you change it.  This sample
file shows how you can add comments to your own input
file.

Because printing is different from typewriting,
there are a number of things that you have to do
differently when preparing an input file than if
you were just typing the document directly.
Quotation marks like
       ``this''
have to be handled specially, as do quotes within
quotes:
       ``\,`this'            % \, separates the double and single quote.
        is what I just
        wrote, not  `that'\,''.

Dashes come in three sizes: an
       intra-word
dash, a medium dash for number ranges like
       1--2,
and a punctuation
       dash---like
this.

A sentence-ending space should be larger than the
space between words within a sentence.  You
sometimes have to type special commands in
conjunction with punctuation characters to get
this right, as in the following sentence.
       Gnats, gnus, etc.\ all  % `\ ' makes an inter-word space.
       begin with G\@.         % \@ marks end-of-sentence punctuation.
You should check the spaces after periods when
reading your output to make sure you haven't
forgotten any special cases.  Generating an
ellipsis
       \ldots\               % `\ ' is needed after `\ldots' because TeX
                             % ignores spaces after command names like \ldots
                             % made from \ + letters.
                             %
                             % Note how a `%' character causes TeX to ignore
                             % the end of the input line, so these blank lines
                             % do not start a new paragraph.
                             %
with the right spacing around the periods requires
a special command.

\LaTeX\ interprets some common characters as
commands, so you must type special commands to
generate them.  These characters include the
following:
       \$ \& \% \# \{ and \}.

In printing, text is usually emphasized with an
       \emph{italic}
type style.

\begin{em}
   A long segment of text can also be emphasized
   in this way.  Text within such a segment can be
   given \emph{additional} emphasis.
\end{em}

It is sometimes necessary to prevent \LaTeX\ from
breaking a line where it might otherwise do so.
This may be at a space, as between the ``Mr.''\ and
``Jones'' in
       ``Mr.~Jones'',        % ~ produces an unbreakable interword space.
or within a word---especially when the word is a
symbol like
       \mbox{\emph{itemnum}}
that makes little sense when hyphenated across
lines.

Footnotes\footnote{This is an example of a footnote.}
pose no problem.

\LaTeX\ is good at typesetting mathematical formulas
like
       \( x-3y + z = 7 \)
or
       \( a_{1} > x^{2n} + y^{2n} > x' \)
or
       \( \ip{A}{B} = \sum_{i} a_{i} b_{i} \).
The spaces you type in a formula are
ignored.  Remember that a letter like
       $x$                   % $ ... $  and  \( ... \)  are equivalent
is a formula when it denotes a mathematical
symbol, and it should be typed as one.

\section{Displayed Text}

Text is displayed by indenting it from the left
margin.  Quotations are commonly displayed.  There
are short quotations
\begin{quote}
   This is a short quotation.  It consists of a
   single paragraph of text.  See how it is formatted.
\end{quote}
and longer ones.
\begin{quotation}
   This is a longer quotation.  It consists of two
   paragraphs of text, neither of which are
   particularly interesting.

   This is the second paragraph of the quotation.  It
   is just as dull as the first paragraph.
\end{quotation}
Another frequently-displayed structure is a list.
The following is an example of an \emph{itemized}
list.
\begin{itemize}
   \item This is the first item of an itemized list.
         Each item in the list is marked with a ``tick''.
         You don't have to worry about what kind of tick
         mark is used.

   \item This is the second item of the list.  It
         contains another list nested inside it.  The inner
         list is an \emph{enumerated} list.
         \begin{enumerate}
            \item This is the first item of an enumerated
                  list that is nested within the itemized list.

            \item This is the second item of the inner list.
                  \LaTeX\ allows you to nest lists deeper than
                  you really should.
         \end{enumerate}
         This is the rest of the second item of the outer
         list.  It is no more interesting than any other
         part of the item.
   \item This is the third item of the list.
\end{itemize}
You can even display poetry.
\begin{verse}
   There is an environment
    for verse \\             % The \\ command separates lines
   Whose features some poets % within a stanza.
   will curse.

                             % One or more blank lines separate stanzas.

   For instead of making\\
   Them do \emph{all} line breaking, \\
   It allows them to put too many words on a line when they'd rather be
   forced to be terse.
\end{verse}

Mathematical formulas may also be displayed.  A
displayed formula
is
one-line long; multiline
formulas require special formatting instructions.
   \[  \ip{\Gamma}{\psi'} = x'' + y^{2} + z_{i}^{n}\]
Don't start a paragraph with a displayed equation,
nor make one a paragraph by itself.

\end{document}               % End of document.

\section{Bench Huge Section 756}
% synthetic bench section 756
% This is a sample LaTeX input file.  (Version of 12 August 2004.)
%
% A '%' character causes TeX to ignore all remaining text on the line,
% and is used for comments like this one.

\documentclass{article}      % Specifies the document class

                             % The preamble begins here.
\title{An Example Document}  % Declares the document's title.
\author{Leslie Lamport}      % Declares the author's name.
\date{January 21, 1994}      % Deleting this command produces today's date.

\newcommand{\ip}[2]{(#1, #2)}
                             % Defines \ip{arg1}{arg2} to mean
                             % (arg1, arg2).

%\newcommand{\ip}[2]{\langle #1 | #2\rangle}
                             % This is an alternative definition of
                             % \ip that is commented out.

\begin{document}             % End of preamble and beginning of text.

\maketitle                   % Produces the title.

This is an example input file.  Comparing it with
the output it generates can show you how to
produce a simple document of your own.

\section{Ordinary Text}      % Produces section heading.  Lower-level
                             % sections are begun with similar
                             % \subsection and \subsubsection commands.

The ends  of words and sentences are marked
  by   spaces. It  doesn't matter how many
spaces    you type; one is as good as 100.  The
end of   a line counts as a space.

One   or more   blank lines denote the  end
of  a paragraph.

Since any number of consecutive spaces are treated
like a single one, the formatting of the input
file makes no difference to
      \LaTeX,                % The \LaTeX command generates the LaTeX logo.
but it makes a difference to you.  When you use
\LaTeX, making your input file as easy to read
as possible will be a great help as you write
your document and when you change it.  This sample
file shows how you can add comments to your own input
file.

Because printing is different from typewriting,
there are a number of things that you have to do
differently when preparing an input file than if
you were just typing the document directly.
Quotation marks like
       ``this''
have to be handled specially, as do quotes within
quotes:
       ``\,`this'            % \, separates the double and single quote.
        is what I just
        wrote, not  `that'\,''.

Dashes come in three sizes: an
       intra-word
dash, a medium dash for number ranges like
       1--2,
and a punctuation
       dash---like
this.

A sentence-ending space should be larger than the
space between words within a sentence.  You
sometimes have to type special commands in
conjunction with punctuation characters to get
this right, as in the following sentence.
       Gnats, gnus, etc.\ all  % `\ ' makes an inter-word space.
       begin with G\@.         % \@ marks end-of-sentence punctuation.
You should check the spaces after periods when
reading your output to make sure you haven't
forgotten any special cases.  Generating an
ellipsis
       \ldots\               % `\ ' is needed after `\ldots' because TeX
                             % ignores spaces after command names like \ldots
                             % made from \ + letters.
                             %
                             % Note how a `%' character causes TeX to ignore
                             % the end of the input line, so these blank lines
                             % do not start a new paragraph.
                             %
with the right spacing around the periods requires
a special command.

\LaTeX\ interprets some common characters as
commands, so you must type special commands to
generate them.  These characters include the
following:
       \$ \& \% \# \{ and \}.

In printing, text is usually emphasized with an
       \emph{italic}
type style.

\begin{em}
   A long segment of text can also be emphasized
   in this way.  Text within such a segment can be
   given \emph{additional} emphasis.
\end{em}

It is sometimes necessary to prevent \LaTeX\ from
breaking a line where it might otherwise do so.
This may be at a space, as between the ``Mr.''\ and
``Jones'' in
       ``Mr.~Jones'',        % ~ produces an unbreakable interword space.
or within a word---especially when the word is a
symbol like
       \mbox{\emph{itemnum}}
that makes little sense when hyphenated across
lines.

Footnotes\footnote{This is an example of a footnote.}
pose no problem.

\LaTeX\ is good at typesetting mathematical formulas
like
       \( x-3y + z = 7 \)
or
       \( a_{1} > x^{2n} + y^{2n} > x' \)
or
       \( \ip{A}{B} = \sum_{i} a_{i} b_{i} \).
The spaces you type in a formula are
ignored.  Remember that a letter like
       $x$                   % $ ... $  and  \( ... \)  are equivalent
is a formula when it denotes a mathematical
symbol, and it should be typed as one.

\section{Displayed Text}

Text is displayed by indenting it from the left
margin.  Quotations are commonly displayed.  There
are short quotations
\begin{quote}
   This is a short quotation.  It consists of a
   single paragraph of text.  See how it is formatted.
\end{quote}
and longer ones.
\begin{quotation}
   This is a longer quotation.  It consists of two
   paragraphs of text, neither of which are
   particularly interesting.

   This is the second paragraph of the quotation.  It
   is just as dull as the first paragraph.
\end{quotation}
Another frequently-displayed structure is a list.
The following is an example of an \emph{itemized}
list.
\begin{itemize}
   \item This is the first item of an itemized list.
         Each item in the list is marked with a ``tick''.
         You don't have to worry about what kind of tick
         mark is used.

   \item This is the second item of the list.  It
         contains another list nested inside it.  The inner
         list is an \emph{enumerated} list.
         \begin{enumerate}
            \item This is the first item of an enumerated
                  list that is nested within the itemized list.

            \item This is the second item of the inner list.
                  \LaTeX\ allows you to nest lists deeper than
                  you really should.
         \end{enumerate}
         This is the rest of the second item of the outer
         list.  It is no more interesting than any other
         part of the item.
   \item This is the third item of the list.
\end{itemize}
You can even display poetry.
\begin{verse}
   There is an environment
    for verse \\             % The \\ command separates lines
   Whose features some poets % within a stanza.
   will curse.

                             % One or more blank lines separate stanzas.

   For instead of making\\
   Them do \emph{all} line breaking, \\
   It allows them to put too many words on a line when they'd rather be
   forced to be terse.
\end{verse}

Mathematical formulas may also be displayed.  A
displayed formula
is
one-line long; multiline
formulas require special formatting instructions.
   \[  \ip{\Gamma}{\psi'} = x'' + y^{2} + z_{i}^{n}\]
Don't start a paragraph with a displayed equation,
nor make one a paragraph by itself.

\end{document}               % End of document.

\section{Bench Huge Section 757}
% synthetic bench section 757
% This is a sample LaTeX input file.  (Version of 12 August 2004.)
%
% A '%' character causes TeX to ignore all remaining text on the line,
% and is used for comments like this one.

\documentclass{article}      % Specifies the document class

                             % The preamble begins here.
\title{An Example Document}  % Declares the document's title.
\author{Leslie Lamport}      % Declares the author's name.
\date{January 21, 1994}      % Deleting this command produces today's date.

\newcommand{\ip}[2]{(#1, #2)}
                             % Defines \ip{arg1}{arg2} to mean
                             % (arg1, arg2).

%\newcommand{\ip}[2]{\langle #1 | #2\rangle}
                             % This is an alternative definition of
                             % \ip that is commented out.

\begin{document}             % End of preamble and beginning of text.

\maketitle                   % Produces the title.

This is an example input file.  Comparing it with
the output it generates can show you how to
produce a simple document of your own.

\section{Ordinary Text}      % Produces section heading.  Lower-level
                             % sections are begun with similar
                             % \subsection and \subsubsection commands.

The ends  of words and sentences are marked
  by   spaces. It  doesn't matter how many
spaces    you type; one is as good as 100.  The
end of   a line counts as a space.

One   or more   blank lines denote the  end
of  a paragraph.

Since any number of consecutive spaces are treated
like a single one, the formatting of the input
file makes no difference to
      \LaTeX,                % The \LaTeX command generates the LaTeX logo.
but it makes a difference to you.  When you use
\LaTeX, making your input file as easy to read
as possible will be a great help as you write
your document and when you change it.  This sample
file shows how you can add comments to your own input
file.

Because printing is different from typewriting,
there are a number of things that you have to do
differently when preparing an input file than if
you were just typing the document directly.
Quotation marks like
       ``this''
have to be handled specially, as do quotes within
quotes:
       ``\,`this'            % \, separates the double and single quote.
        is what I just
        wrote, not  `that'\,''.

Dashes come in three sizes: an
       intra-word
dash, a medium dash for number ranges like
       1--2,
and a punctuation
       dash---like
this.

A sentence-ending space should be larger than the
space between words within a sentence.  You
sometimes have to type special commands in
conjunction with punctuation characters to get
this right, as in the following sentence.
       Gnats, gnus, etc.\ all  % `\ ' makes an inter-word space.
       begin with G\@.         % \@ marks end-of-sentence punctuation.
You should check the spaces after periods when
reading your output to make sure you haven't
forgotten any special cases.  Generating an
ellipsis
       \ldots\               % `\ ' is needed after `\ldots' because TeX
                             % ignores spaces after command names like \ldots
                             % made from \ + letters.
                             %
                             % Note how a `%' character causes TeX to ignore
                             % the end of the input line, so these blank lines
                             % do not start a new paragraph.
                             %
with the right spacing around the periods requires
a special command.

\LaTeX\ interprets some common characters as
commands, so you must type special commands to
generate them.  These characters include the
following:
       \$ \& \% \# \{ and \}.

In printing, text is usually emphasized with an
       \emph{italic}
type style.

\begin{em}
   A long segment of text can also be emphasized
   in this way.  Text within such a segment can be
   given \emph{additional} emphasis.
\end{em}

It is sometimes necessary to prevent \LaTeX\ from
breaking a line where it might otherwise do so.
This may be at a space, as between the ``Mr.''\ and
``Jones'' in
       ``Mr.~Jones'',        % ~ produces an unbreakable interword space.
or within a word---especially when the word is a
symbol like
       \mbox{\emph{itemnum}}
that makes little sense when hyphenated across
lines.

Footnotes\footnote{This is an example of a footnote.}
pose no problem.

\LaTeX\ is good at typesetting mathematical formulas
like
       \( x-3y + z = 7 \)
or
       \( a_{1} > x^{2n} + y^{2n} > x' \)
or
       \( \ip{A}{B} = \sum_{i} a_{i} b_{i} \).
The spaces you type in a formula are
ignored.  Remember that a letter like
       $x$                   % $ ... $  and  \( ... \)  are equivalent
is a formula when it denotes a mathematical
symbol, and it should be typed as one.

\section{Displayed Text}

Text is displayed by indenting it from the left
margin.  Quotations are commonly displayed.  There
are short quotations
\begin{quote}
   This is a short quotation.  It consists of a
   single paragraph of text.  See how it is formatted.
\end{quote}
and longer ones.
\begin{quotation}
   This is a longer quotation.  It consists of two
   paragraphs of text, neither of which are
   particularly interesting.

   This is the second paragraph of the quotation.  It
   is just as dull as the first paragraph.
\end{quotation}
Another frequently-displayed structure is a list.
The following is an example of an \emph{itemized}
list.
\begin{itemize}
   \item This is the first item of an itemized list.
         Each item in the list is marked with a ``tick''.
         You don't have to worry about what kind of tick
         mark is used.

   \item This is the second item of the list.  It
         contains another list nested inside it.  The inner
         list is an \emph{enumerated} list.
         \begin{enumerate}
            \item This is the first item of an enumerated
                  list that is nested within the itemized list.

            \item This is the second item of the inner list.
                  \LaTeX\ allows you to nest lists deeper than
                  you really should.
         \end{enumerate}
         This is the rest of the second item of the outer
         list.  It is no more interesting than any other
         part of the item.
   \item This is the third item of the list.
\end{itemize}
You can even display poetry.
\begin{verse}
   There is an environment
    for verse \\             % The \\ command separates lines
   Whose features some poets % within a stanza.
   will curse.

                             % One or more blank lines separate stanzas.

   For instead of making\\
   Them do \emph{all} line breaking, \\
   It allows them to put too many words on a line when they'd rather be
   forced to be terse.
\end{verse}

Mathematical formulas may also be displayed.  A
displayed formula
is
one-line long; multiline
formulas require special formatting instructions.
   \[  \ip{\Gamma}{\psi'} = x'' + y^{2} + z_{i}^{n}\]
Don't start a paragraph with a displayed equation,
nor make one a paragraph by itself.

\end{document}               % End of document.

\section{Bench Huge Section 758}
% synthetic bench section 758
% This is a sample LaTeX input file.  (Version of 12 August 2004.)
%
% A '%' character causes TeX to ignore all remaining text on the line,
% and is used for comments like this one.

\documentclass{article}      % Specifies the document class

                             % The preamble begins here.
\title{An Example Document}  % Declares the document's title.
\author{Leslie Lamport}      % Declares the author's name.
\date{January 21, 1994}      % Deleting this command produces today's date.

\newcommand{\ip}[2]{(#1, #2)}
                             % Defines \ip{arg1}{arg2} to mean
                             % (arg1, arg2).

%\newcommand{\ip}[2]{\langle #1 | #2\rangle}
                             % This is an alternative definition of
                             % \ip that is commented out.

\begin{document}             % End of preamble and beginning of text.

\maketitle                   % Produces the title.

This is an example input file.  Comparing it with
the output it generates can show you how to
produce a simple document of your own.

\section{Ordinary Text}      % Produces section heading.  Lower-level
                             % sections are begun with similar
                             % \subsection and \subsubsection commands.

The ends  of words and sentences are marked
  by   spaces. It  doesn't matter how many
spaces    you type; one is as good as 100.  The
end of   a line counts as a space.

One   or more   blank lines denote the  end
of  a paragraph.

Since any number of consecutive spaces are treated
like a single one, the formatting of the input
file makes no difference to
      \LaTeX,                % The \LaTeX command generates the LaTeX logo.
but it makes a difference to you.  When you use
\LaTeX, making your input file as easy to read
as possible will be a great help as you write
your document and when you change it.  This sample
file shows how you can add comments to your own input
file.

Because printing is different from typewriting,
there are a number of things that you have to do
differently when preparing an input file than if
you were just typing the document directly.
Quotation marks like
       ``this''
have to be handled specially, as do quotes within
quotes:
       ``\,`this'            % \, separates the double and single quote.
        is what I just
        wrote, not  `that'\,''.

Dashes come in three sizes: an
       intra-word
dash, a medium dash for number ranges like
       1--2,
and a punctuation
       dash---like
this.

A sentence-ending space should be larger than the
space between words within a sentence.  You
sometimes have to type special commands in
conjunction with punctuation characters to get
this right, as in the following sentence.
       Gnats, gnus, etc.\ all  % `\ ' makes an inter-word space.
       begin with G\@.         % \@ marks end-of-sentence punctuation.
You should check the spaces after periods when
reading your output to make sure you haven't
forgotten any special cases.  Generating an
ellipsis
       \ldots\               % `\ ' is needed after `\ldots' because TeX
                             % ignores spaces after command names like \ldots
                             % made from \ + letters.
                             %
                             % Note how a `%' character causes TeX to ignore
                             % the end of the input line, so these blank lines
                             % do not start a new paragraph.
                             %
with the right spacing around the periods requires
a special command.

\LaTeX\ interprets some common characters as
commands, so you must type special commands to
generate them.  These characters include the
following:
       \$ \& \% \# \{ and \}.

In printing, text is usually emphasized with an
       \emph{italic}
type style.

\begin{em}
   A long segment of text can also be emphasized
   in this way.  Text within such a segment can be
   given \emph{additional} emphasis.
\end{em}

It is sometimes necessary to prevent \LaTeX\ from
breaking a line where it might otherwise do so.
This may be at a space, as between the ``Mr.''\ and
``Jones'' in
       ``Mr.~Jones'',        % ~ produces an unbreakable interword space.
or within a word---especially when the word is a
symbol like
       \mbox{\emph{itemnum}}
that makes little sense when hyphenated across
lines.

Footnotes\footnote{This is an example of a footnote.}
pose no problem.

\LaTeX\ is good at typesetting mathematical formulas
like
       \( x-3y + z = 7 \)
or
       \( a_{1} > x^{2n} + y^{2n} > x' \)
or
       \( \ip{A}{B} = \sum_{i} a_{i} b_{i} \).
The spaces you type in a formula are
ignored.  Remember that a letter like
       $x$                   % $ ... $  and  \( ... \)  are equivalent
is a formula when it denotes a mathematical
symbol, and it should be typed as one.

\section{Displayed Text}

Text is displayed by indenting it from the left
margin.  Quotations are commonly displayed.  There
are short quotations
\begin{quote}
   This is a short quotation.  It consists of a
   single paragraph of text.  See how it is formatted.
\end{quote}
and longer ones.
\begin{quotation}
   This is a longer quotation.  It consists of two
   paragraphs of text, neither of which are
   particularly interesting.

   This is the second paragraph of the quotation.  It
   is just as dull as the first paragraph.
\end{quotation}
Another frequently-displayed structure is a list.
The following is an example of an \emph{itemized}
list.
\begin{itemize}
   \item This is the first item of an itemized list.
         Each item in the list is marked with a ``tick''.
         You don't have to worry about what kind of tick
         mark is used.

   \item This is the second item of the list.  It
         contains another list nested inside it.  The inner
         list is an \emph{enumerated} list.
         \begin{enumerate}
            \item This is the first item of an enumerated
                  list that is nested within the itemized list.

            \item This is the second item of the inner list.
                  \LaTeX\ allows you to nest lists deeper than
                  you really should.
         \end{enumerate}
         This is the rest of the second item of the outer
         list.  It is no more interesting than any other
         part of the item.
   \item This is the third item of the list.
\end{itemize}
You can even display poetry.
\begin{verse}
   There is an environment
    for verse \\             % The \\ command separates lines
   Whose features some poets % within a stanza.
   will curse.

                             % One or more blank lines separate stanzas.

   For instead of making\\
   Them do \emph{all} line breaking, \\
   It allows them to put too many words on a line when they'd rather be
   forced to be terse.
\end{verse}

Mathematical formulas may also be displayed.  A
displayed formula
is
one-line long; multiline
formulas require special formatting instructions.
   \[  \ip{\Gamma}{\psi'} = x'' + y^{2} + z_{i}^{n}\]
Don't start a paragraph with a displayed equation,
nor make one a paragraph by itself.

\end{document}               % End of document.

\section{Bench Huge Section 759}
% synthetic bench section 759
% This is a sample LaTeX input file.  (Version of 12 August 2004.)
%
% A '%' character causes TeX to ignore all remaining text on the line,
% and is used for comments like this one.

\documentclass{article}      % Specifies the document class

                             % The preamble begins here.
\title{An Example Document}  % Declares the document's title.
\author{Leslie Lamport}      % Declares the author's name.
\date{January 21, 1994}      % Deleting this command produces today's date.

\newcommand{\ip}[2]{(#1, #2)}
                             % Defines \ip{arg1}{arg2} to mean
                             % (arg1, arg2).

%\newcommand{\ip}[2]{\langle #1 | #2\rangle}
                             % This is an alternative definition of
                             % \ip that is commented out.

\begin{document}             % End of preamble and beginning of text.

\maketitle                   % Produces the title.

This is an example input file.  Comparing it with
the output it generates can show you how to
produce a simple document of your own.

\section{Ordinary Text}      % Produces section heading.  Lower-level
                             % sections are begun with similar
                             % \subsection and \subsubsection commands.

The ends  of words and sentences are marked
  by   spaces. It  doesn't matter how many
spaces    you type; one is as good as 100.  The
end of   a line counts as a space.

One   or more   blank lines denote the  end
of  a paragraph.

Since any number of consecutive spaces are treated
like a single one, the formatting of the input
file makes no difference to
      \LaTeX,                % The \LaTeX command generates the LaTeX logo.
but it makes a difference to you.  When you use
\LaTeX, making your input file as easy to read
as possible will be a great help as you write
your document and when you change it.  This sample
file shows how you can add comments to your own input
file.

Because printing is different from typewriting,
there are a number of things that you have to do
differently when preparing an input file than if
you were just typing the document directly.
Quotation marks like
       ``this''
have to be handled specially, as do quotes within
quotes:
       ``\,`this'            % \, separates the double and single quote.
        is what I just
        wrote, not  `that'\,''.

Dashes come in three sizes: an
       intra-word
dash, a medium dash for number ranges like
       1--2,
and a punctuation
       dash---like
this.

A sentence-ending space should be larger than the
space between words within a sentence.  You
sometimes have to type special commands in
conjunction with punctuation characters to get
this right, as in the following sentence.
       Gnats, gnus, etc.\ all  % `\ ' makes an inter-word space.
       begin with G\@.         % \@ marks end-of-sentence punctuation.
You should check the spaces after periods when
reading your output to make sure you haven't
forgotten any special cases.  Generating an
ellipsis
       \ldots\               % `\ ' is needed after `\ldots' because TeX
                             % ignores spaces after command names like \ldots
                             % made from \ + letters.
                             %
                             % Note how a `%' character causes TeX to ignore
                             % the end of the input line, so these blank lines
                             % do not start a new paragraph.
                             %
with the right spacing around the periods requires
a special command.

\LaTeX\ interprets some common characters as
commands, so you must type special commands to
generate them.  These characters include the
following:
       \$ \& \% \# \{ and \}.

In printing, text is usually emphasized with an
       \emph{italic}
type style.

\begin{em}
   A long segment of text can also be emphasized
   in this way.  Text within such a segment can be
   given \emph{additional} emphasis.
\end{em}

It is sometimes necessary to prevent \LaTeX\ from
breaking a line where it might otherwise do so.
This may be at a space, as between the ``Mr.''\ and
``Jones'' in
       ``Mr.~Jones'',        % ~ produces an unbreakable interword space.
or within a word---especially when the word is a
symbol like
       \mbox{\emph{itemnum}}
that makes little sense when hyphenated across
lines.

Footnotes\footnote{This is an example of a footnote.}
pose no problem.

\LaTeX\ is good at typesetting mathematical formulas
like
       \( x-3y + z = 7 \)
or
       \( a_{1} > x^{2n} + y^{2n} > x' \)
or
       \( \ip{A}{B} = \sum_{i} a_{i} b_{i} \).
The spaces you type in a formula are
ignored.  Remember that a letter like
       $x$                   % $ ... $  and  \( ... \)  are equivalent
is a formula when it denotes a mathematical
symbol, and it should be typed as one.

\section{Displayed Text}

Text is displayed by indenting it from the left
margin.  Quotations are commonly displayed.  There
are short quotations
\begin{quote}
   This is a short quotation.  It consists of a
   single paragraph of text.  See how it is formatted.
\end{quote}
and longer ones.
\begin{quotation}
   This is a longer quotation.  It consists of two
   paragraphs of text, neither of which are
   particularly interesting.

   This is the second paragraph of the quotation.  It
   is just as dull as the first paragraph.
\end{quotation}
Another frequently-displayed structure is a list.
The following is an example of an \emph{itemized}
list.
\begin{itemize}
   \item This is the first item of an itemized list.
         Each item in the list is marked with a ``tick''.
         You don't have to worry about what kind of tick
         mark is used.

   \item This is the second item of the list.  It
         contains another list nested inside it.  The inner
         list is an \emph{enumerated} list.
         \begin{enumerate}
            \item This is the first item of an enumerated
                  list that is nested within the itemized list.

            \item This is the second item of the inner list.
                  \LaTeX\ allows you to nest lists deeper than
                  you really should.
         \end{enumerate}
         This is the rest of the second item of the outer
         list.  It is no more interesting than any other
         part of the item.
   \item This is the third item of the list.
\end{itemize}
You can even display poetry.
\begin{verse}
   There is an environment
    for verse \\             % The \\ command separates lines
   Whose features some poets % within a stanza.
   will curse.

                             % One or more blank lines separate stanzas.

   For instead of making\\
   Them do \emph{all} line breaking, \\
   It allows them to put too many words on a line when they'd rather be
   forced to be terse.
\end{verse}

Mathematical formulas may also be displayed.  A
displayed formula
is
one-line long; multiline
formulas require special formatting instructions.
   \[  \ip{\Gamma}{\psi'} = x'' + y^{2} + z_{i}^{n}\]
Don't start a paragraph with a displayed equation,
nor make one a paragraph by itself.

\end{document}               % End of document.

\section{Bench Huge Section 760}
% synthetic bench section 760
% This is a sample LaTeX input file.  (Version of 12 August 2004.)
%
% A '%' character causes TeX to ignore all remaining text on the line,
% and is used for comments like this one.

\documentclass{article}      % Specifies the document class

                             % The preamble begins here.
\title{An Example Document}  % Declares the document's title.
\author{Leslie Lamport}      % Declares the author's name.
\date{January 21, 1994}      % Deleting this command produces today's date.

\newcommand{\ip}[2]{(#1, #2)}
                             % Defines \ip{arg1}{arg2} to mean
                             % (arg1, arg2).

%\newcommand{\ip}[2]{\langle #1 | #2\rangle}
                             % This is an alternative definition of
                             % \ip that is commented out.

\begin{document}             % End of preamble and beginning of text.

\maketitle                   % Produces the title.

This is an example input file.  Comparing it with
the output it generates can show you how to
produce a simple document of your own.

\section{Ordinary Text}      % Produces section heading.  Lower-level
                             % sections are begun with similar
                             % \subsection and \subsubsection commands.

The ends  of words and sentences are marked
  by   spaces. It  doesn't matter how many
spaces    you type; one is as good as 100.  The
end of   a line counts as a space.

One   or more   blank lines denote the  end
of  a paragraph.

Since any number of consecutive spaces are treated
like a single one, the formatting of the input
file makes no difference to
      \LaTeX,                % The \LaTeX command generates the LaTeX logo.
but it makes a difference to you.  When you use
\LaTeX, making your input file as easy to read
as possible will be a great help as you write
your document and when you change it.  This sample
file shows how you can add comments to your own input
file.

Because printing is different from typewriting,
there are a number of things that you have to do
differently when preparing an input file than if
you were just typing the document directly.
Quotation marks like
       ``this''
have to be handled specially, as do quotes within
quotes:
       ``\,`this'            % \, separates the double and single quote.
        is what I just
        wrote, not  `that'\,''.

Dashes come in three sizes: an
       intra-word
dash, a medium dash for number ranges like
       1--2,
and a punctuation
       dash---like
this.

A sentence-ending space should be larger than the
space between words within a sentence.  You
sometimes have to type special commands in
conjunction with punctuation characters to get
this right, as in the following sentence.
       Gnats, gnus, etc.\ all  % `\ ' makes an inter-word space.
       begin with G\@.         % \@ marks end-of-sentence punctuation.
You should check the spaces after periods when
reading your output to make sure you haven't
forgotten any special cases.  Generating an
ellipsis
       \ldots\               % `\ ' is needed after `\ldots' because TeX
                             % ignores spaces after command names like \ldots
                             % made from \ + letters.
                             %
                             % Note how a `%' character causes TeX to ignore
                             % the end of the input line, so these blank lines
                             % do not start a new paragraph.
                             %
with the right spacing around the periods requires
a special command.

\LaTeX\ interprets some common characters as
commands, so you must type special commands to
generate them.  These characters include the
following:
       \$ \& \% \# \{ and \}.

In printing, text is usually emphasized with an
       \emph{italic}
type style.

\begin{em}
   A long segment of text can also be emphasized
   in this way.  Text within such a segment can be
   given \emph{additional} emphasis.
\end{em}

It is sometimes necessary to prevent \LaTeX\ from
breaking a line where it might otherwise do so.
This may be at a space, as between the ``Mr.''\ and
``Jones'' in
       ``Mr.~Jones'',        % ~ produces an unbreakable interword space.
or within a word---especially when the word is a
symbol like
       \mbox{\emph{itemnum}}
that makes little sense when hyphenated across
lines.

Footnotes\footnote{This is an example of a footnote.}
pose no problem.

\LaTeX\ is good at typesetting mathematical formulas
like
       \( x-3y + z = 7 \)
or
       \( a_{1} > x^{2n} + y^{2n} > x' \)
or
       \( \ip{A}{B} = \sum_{i} a_{i} b_{i} \).
The spaces you type in a formula are
ignored.  Remember that a letter like
       $x$                   % $ ... $  and  \( ... \)  are equivalent
is a formula when it denotes a mathematical
symbol, and it should be typed as one.

\section{Displayed Text}

Text is displayed by indenting it from the left
margin.  Quotations are commonly displayed.  There
are short quotations
\begin{quote}
   This is a short quotation.  It consists of a
   single paragraph of text.  See how it is formatted.
\end{quote}
and longer ones.
\begin{quotation}
   This is a longer quotation.  It consists of two
   paragraphs of text, neither of which are
   particularly interesting.

   This is the second paragraph of the quotation.  It
   is just as dull as the first paragraph.
\end{quotation}
Another frequently-displayed structure is a list.
The following is an example of an \emph{itemized}
list.
\begin{itemize}
   \item This is the first item of an itemized list.
         Each item in the list is marked with a ``tick''.
         You don't have to worry about what kind of tick
         mark is used.

   \item This is the second item of the list.  It
         contains another list nested inside it.  The inner
         list is an \emph{enumerated} list.
         \begin{enumerate}
            \item This is the first item of an enumerated
                  list that is nested within the itemized list.

            \item This is the second item of the inner list.
                  \LaTeX\ allows you to nest lists deeper than
                  you really should.
         \end{enumerate}
         This is the rest of the second item of the outer
         list.  It is no more interesting than any other
         part of the item.
   \item This is the third item of the list.
\end{itemize}
You can even display poetry.
\begin{verse}
   There is an environment
    for verse \\             % The \\ command separates lines
   Whose features some poets % within a stanza.
   will curse.

                             % One or more blank lines separate stanzas.

   For instead of making\\
   Them do \emph{all} line breaking, \\
   It allows them to put too many words on a line when they'd rather be
   forced to be terse.
\end{verse}

Mathematical formulas may also be displayed.  A
displayed formula
is
one-line long; multiline
formulas require special formatting instructions.
   \[  \ip{\Gamma}{\psi'} = x'' + y^{2} + z_{i}^{n}\]
Don't start a paragraph with a displayed equation,
nor make one a paragraph by itself.

\end{document}               % End of document.

\section{Bench Huge Section 761}
% synthetic bench section 761
% This is a sample LaTeX input file.  (Version of 12 August 2004.)
%
% A '%' character causes TeX to ignore all remaining text on the line,
% and is used for comments like this one.

\documentclass{article}      % Specifies the document class

                             % The preamble begins here.
\title{An Example Document}  % Declares the document's title.
\author{Leslie Lamport}      % Declares the author's name.
\date{January 21, 1994}      % Deleting this command produces today's date.

\newcommand{\ip}[2]{(#1, #2)}
                             % Defines \ip{arg1}{arg2} to mean
                             % (arg1, arg2).

%\newcommand{\ip}[2]{\langle #1 | #2\rangle}
                             % This is an alternative definition of
                             % \ip that is commented out.

\begin{document}             % End of preamble and beginning of text.

\maketitle                   % Produces the title.

This is an example input file.  Comparing it with
the output it generates can show you how to
produce a simple document of your own.

\section{Ordinary Text}      % Produces section heading.  Lower-level
                             % sections are begun with similar
                             % \subsection and \subsubsection commands.

The ends  of words and sentences are marked
  by   spaces. It  doesn't matter how many
spaces    you type; one is as good as 100.  The
end of   a line counts as a space.

One   or more   blank lines denote the  end
of  a paragraph.

Since any number of consecutive spaces are treated
like a single one, the formatting of the input
file makes no difference to
      \LaTeX,                % The \LaTeX command generates the LaTeX logo.
but it makes a difference to you.  When you use
\LaTeX, making your input file as easy to read
as possible will be a great help as you write
your document and when you change it.  This sample
file shows how you can add comments to your own input
file.

Because printing is different from typewriting,
there are a number of things that you have to do
differently when preparing an input file than if
you were just typing the document directly.
Quotation marks like
       ``this''
have to be handled specially, as do quotes within
quotes:
       ``\,`this'            % \, separates the double and single quote.
        is what I just
        wrote, not  `that'\,''.

Dashes come in three sizes: an
       intra-word
dash, a medium dash for number ranges like
       1--2,
and a punctuation
       dash---like
this.

A sentence-ending space should be larger than the
space between words within a sentence.  You
sometimes have to type special commands in
conjunction with punctuation characters to get
this right, as in the following sentence.
       Gnats, gnus, etc.\ all  % `\ ' makes an inter-word space.
       begin with G\@.         % \@ marks end-of-sentence punctuation.
You should check the spaces after periods when
reading your output to make sure you haven't
forgotten any special cases.  Generating an
ellipsis
       \ldots\               % `\ ' is needed after `\ldots' because TeX
                             % ignores spaces after command names like \ldots
                             % made from \ + letters.
                             %
                             % Note how a `%' character causes TeX to ignore
                             % the end of the input line, so these blank lines
                             % do not start a new paragraph.
                             %
with the right spacing around the periods requires
a special command.

\LaTeX\ interprets some common characters as
commands, so you must type special commands to
generate them.  These characters include the
following:
       \$ \& \% \# \{ and \}.

In printing, text is usually emphasized with an
       \emph{italic}
type style.

\begin{em}
   A long segment of text can also be emphasized
   in this way.  Text within such a segment can be
   given \emph{additional} emphasis.
\end{em}

It is sometimes necessary to prevent \LaTeX\ from
breaking a line where it might otherwise do so.
This may be at a space, as between the ``Mr.''\ and
``Jones'' in
       ``Mr.~Jones'',        % ~ produces an unbreakable interword space.
or within a word---especially when the word is a
symbol like
       \mbox{\emph{itemnum}}
that makes little sense when hyphenated across
lines.

Footnotes\footnote{This is an example of a footnote.}
pose no problem.

\LaTeX\ is good at typesetting mathematical formulas
like
       \( x-3y + z = 7 \)
or
       \( a_{1} > x^{2n} + y^{2n} > x' \)
or
       \( \ip{A}{B} = \sum_{i} a_{i} b_{i} \).
The spaces you type in a formula are
ignored.  Remember that a letter like
       $x$                   % $ ... $  and  \( ... \)  are equivalent
is a formula when it denotes a mathematical
symbol, and it should be typed as one.

\section{Displayed Text}

Text is displayed by indenting it from the left
margin.  Quotations are commonly displayed.  There
are short quotations
\begin{quote}
   This is a short quotation.  It consists of a
   single paragraph of text.  See how it is formatted.
\end{quote}
and longer ones.
\begin{quotation}
   This is a longer quotation.  It consists of two
   paragraphs of text, neither of which are
   particularly interesting.

   This is the second paragraph of the quotation.  It
   is just as dull as the first paragraph.
\end{quotation}
Another frequently-displayed structure is a list.
The following is an example of an \emph{itemized}
list.
\begin{itemize}
   \item This is the first item of an itemized list.
         Each item in the list is marked with a ``tick''.
         You don't have to worry about what kind of tick
         mark is used.

   \item This is the second item of the list.  It
         contains another list nested inside it.  The inner
         list is an \emph{enumerated} list.
         \begin{enumerate}
            \item This is the first item of an enumerated
                  list that is nested within the itemized list.

            \item This is the second item of the inner list.
                  \LaTeX\ allows you to nest lists deeper than
                  you really should.
         \end{enumerate}
         This is the rest of the second item of the outer
         list.  It is no more interesting than any other
         part of the item.
   \item This is the third item of the list.
\end{itemize}
You can even display poetry.
\begin{verse}
   There is an environment
    for verse \\             % The \\ command separates lines
   Whose features some poets % within a stanza.
   will curse.

                             % One or more blank lines separate stanzas.

   For instead of making\\
   Them do \emph{all} line breaking, \\
   It allows them to put too many words on a line when they'd rather be
   forced to be terse.
\end{verse}

Mathematical formulas may also be displayed.  A
displayed formula
is
one-line long; multiline
formulas require special formatting instructions.
   \[  \ip{\Gamma}{\psi'} = x'' + y^{2} + z_{i}^{n}\]
Don't start a paragraph with a displayed equation,
nor make one a paragraph by itself.

\end{document}               % End of document.

\section{Bench Huge Section 762}
% synthetic bench section 762
% This is a sample LaTeX input file.  (Version of 12 August 2004.)
%
% A '%' character causes TeX to ignore all remaining text on the line,
% and is used for comments like this one.

\documentclass{article}      % Specifies the document class

                             % The preamble begins here.
\title{An Example Document}  % Declares the document's title.
\author{Leslie Lamport}      % Declares the author's name.
\date{January 21, 1994}      % Deleting this command produces today's date.

\newcommand{\ip}[2]{(#1, #2)}
                             % Defines \ip{arg1}{arg2} to mean
                             % (arg1, arg2).

%\newcommand{\ip}[2]{\langle #1 | #2\rangle}
                             % This is an alternative definition of
                             % \ip that is commented out.

\begin{document}             % End of preamble and beginning of text.

\maketitle                   % Produces the title.

This is an example input file.  Comparing it with
the output it generates can show you how to
produce a simple document of your own.

\section{Ordinary Text}      % Produces section heading.  Lower-level
                             % sections are begun with similar
                             % \subsection and \subsubsection commands.

The ends  of words and sentences are marked
  by   spaces. It  doesn't matter how many
spaces    you type; one is as good as 100.  The
end of   a line counts as a space.

One   or more   blank lines denote the  end
of  a paragraph.

Since any number of consecutive spaces are treated
like a single one, the formatting of the input
file makes no difference to
      \LaTeX,                % The \LaTeX command generates the LaTeX logo.
but it makes a difference to you.  When you use
\LaTeX, making your input file as easy to read
as possible will be a great help as you write
your document and when you change it.  This sample
file shows how you can add comments to your own input
file.

Because printing is different from typewriting,
there are a number of things that you have to do
differently when preparing an input file than if
you were just typing the document directly.
Quotation marks like
       ``this''
have to be handled specially, as do quotes within
quotes:
       ``\,`this'            % \, separates the double and single quote.
        is what I just
        wrote, not  `that'\,''.

Dashes come in three sizes: an
       intra-word
dash, a medium dash for number ranges like
       1--2,
and a punctuation
       dash---like
this.

A sentence-ending space should be larger than the
space between words within a sentence.  You
sometimes have to type special commands in
conjunction with punctuation characters to get
this right, as in the following sentence.
       Gnats, gnus, etc.\ all  % `\ ' makes an inter-word space.
       begin with G\@.         % \@ marks end-of-sentence punctuation.
You should check the spaces after periods when
reading your output to make sure you haven't
forgotten any special cases.  Generating an
ellipsis
       \ldots\               % `\ ' is needed after `\ldots' because TeX
                             % ignores spaces after command names like \ldots
                             % made from \ + letters.
                             %
                             % Note how a `%' character causes TeX to ignore
                             % the end of the input line, so these blank lines
                             % do not start a new paragraph.
                             %
with the right spacing around the periods requires
a special command.

\LaTeX\ interprets some common characters as
commands, so you must type special commands to
generate them.  These characters include the
following:
       \$ \& \% \# \{ and \}.

In printing, text is usually emphasized with an
       \emph{italic}
type style.

\begin{em}
   A long segment of text can also be emphasized
   in this way.  Text within such a segment can be
   given \emph{additional} emphasis.
\end{em}

It is sometimes necessary to prevent \LaTeX\ from
breaking a line where it might otherwise do so.
This may be at a space, as between the ``Mr.''\ and
``Jones'' in
       ``Mr.~Jones'',        % ~ produces an unbreakable interword space.
or within a word---especially when the word is a
symbol like
       \mbox{\emph{itemnum}}
that makes little sense when hyphenated across
lines.

Footnotes\footnote{This is an example of a footnote.}
pose no problem.

\LaTeX\ is good at typesetting mathematical formulas
like
       \( x-3y + z = 7 \)
or
       \( a_{1} > x^{2n} + y^{2n} > x' \)
or
       \( \ip{A}{B} = \sum_{i} a_{i} b_{i} \).
The spaces you type in a formula are
ignored.  Remember that a letter like
       $x$                   % $ ... $  and  \( ... \)  are equivalent
is a formula when it denotes a mathematical
symbol, and it should be typed as one.

\section{Displayed Text}

Text is displayed by indenting it from the left
margin.  Quotations are commonly displayed.  There
are short quotations
\begin{quote}
   This is a short quotation.  It consists of a
   single paragraph of text.  See how it is formatted.
\end{quote}
and longer ones.
\begin{quotation}
   This is a longer quotation.  It consists of two
   paragraphs of text, neither of which are
   particularly interesting.

   This is the second paragraph of the quotation.  It
   is just as dull as the first paragraph.
\end{quotation}
Another frequently-displayed structure is a list.
The following is an example of an \emph{itemized}
list.
\begin{itemize}
   \item This is the first item of an itemized list.
         Each item in the list is marked with a ``tick''.
         You don't have to worry about what kind of tick
         mark is used.

   \item This is the second item of the list.  It
         contains another list nested inside it.  The inner
         list is an \emph{enumerated} list.
         \begin{enumerate}
            \item This is the first item of an enumerated
                  list that is nested within the itemized list.

            \item This is the second item of the inner list.
                  \LaTeX\ allows you to nest lists deeper than
                  you really should.
         \end{enumerate}
         This is the rest of the second item of the outer
         list.  It is no more interesting than any other
         part of the item.
   \item This is the third item of the list.
\end{itemize}
You can even display poetry.
\begin{verse}
   There is an environment
    for verse \\             % The \\ command separates lines
   Whose features some poets % within a stanza.
   will curse.

                             % One or more blank lines separate stanzas.

   For instead of making\\
   Them do \emph{all} line breaking, \\
   It allows them to put too many words on a line when they'd rather be
   forced to be terse.
\end{verse}

Mathematical formulas may also be displayed.  A
displayed formula
is
one-line long; multiline
formulas require special formatting instructions.
   \[  \ip{\Gamma}{\psi'} = x'' + y^{2} + z_{i}^{n}\]
Don't start a paragraph with a displayed equation,
nor make one a paragraph by itself.

\end{document}               % End of document.

\section{Bench Huge Section 763}
% synthetic bench section 763
% This is a sample LaTeX input file.  (Version of 12 August 2004.)
%
% A '%' character causes TeX to ignore all remaining text on the line,
% and is used for comments like this one.

\documentclass{article}      % Specifies the document class

                             % The preamble begins here.
\title{An Example Document}  % Declares the document's title.
\author{Leslie Lamport}      % Declares the author's name.
\date{January 21, 1994}      % Deleting this command produces today's date.

\newcommand{\ip}[2]{(#1, #2)}
                             % Defines \ip{arg1}{arg2} to mean
                             % (arg1, arg2).

%\newcommand{\ip}[2]{\langle #1 | #2\rangle}
                             % This is an alternative definition of
                             % \ip that is commented out.

\begin{document}             % End of preamble and beginning of text.

\maketitle                   % Produces the title.

This is an example input file.  Comparing it with
the output it generates can show you how to
produce a simple document of your own.

\section{Ordinary Text}      % Produces section heading.  Lower-level
                             % sections are begun with similar
                             % \subsection and \subsubsection commands.

The ends  of words and sentences are marked
  by   spaces. It  doesn't matter how many
spaces    you type; one is as good as 100.  The
end of   a line counts as a space.

One   or more   blank lines denote the  end
of  a paragraph.

Since any number of consecutive spaces are treated
like a single one, the formatting of the input
file makes no difference to
      \LaTeX,                % The \LaTeX command generates the LaTeX logo.
but it makes a difference to you.  When you use
\LaTeX, making your input file as easy to read
as possible will be a great help as you write
your document and when you change it.  This sample
file shows how you can add comments to your own input
file.

Because printing is different from typewriting,
there are a number of things that you have to do
differently when preparing an input file than if
you were just typing the document directly.
Quotation marks like
       ``this''
have to be handled specially, as do quotes within
quotes:
       ``\,`this'            % \, separates the double and single quote.
        is what I just
        wrote, not  `that'\,''.

Dashes come in three sizes: an
       intra-word
dash, a medium dash for number ranges like
       1--2,
and a punctuation
       dash---like
this.

A sentence-ending space should be larger than the
space between words within a sentence.  You
sometimes have to type special commands in
conjunction with punctuation characters to get
this right, as in the following sentence.
       Gnats, gnus, etc.\ all  % `\ ' makes an inter-word space.
       begin with G\@.         % \@ marks end-of-sentence punctuation.
You should check the spaces after periods when
reading your output to make sure you haven't
forgotten any special cases.  Generating an
ellipsis
       \ldots\               % `\ ' is needed after `\ldots' because TeX
                             % ignores spaces after command names like \ldots
                             % made from \ + letters.
                             %
                             % Note how a `%' character causes TeX to ignore
                             % the end of the input line, so these blank lines
                             % do not start a new paragraph.
                             %
with the right spacing around the periods requires
a special command.

\LaTeX\ interprets some common characters as
commands, so you must type special commands to
generate them.  These characters include the
following:
       \$ \& \% \# \{ and \}.

In printing, text is usually emphasized with an
       \emph{italic}
type style.

\begin{em}
   A long segment of text can also be emphasized
   in this way.  Text within such a segment can be
   given \emph{additional} emphasis.
\end{em}

It is sometimes necessary to prevent \LaTeX\ from
breaking a line where it might otherwise do so.
This may be at a space, as between the ``Mr.''\ and
``Jones'' in
       ``Mr.~Jones'',        % ~ produces an unbreakable interword space.
or within a word---especially when the word is a
symbol like
       \mbox{\emph{itemnum}}
that makes little sense when hyphenated across
lines.

Footnotes\footnote{This is an example of a footnote.}
pose no problem.

\LaTeX\ is good at typesetting mathematical formulas
like
       \( x-3y + z = 7 \)
or
       \( a_{1} > x^{2n} + y^{2n} > x' \)
or
       \( \ip{A}{B} = \sum_{i} a_{i} b_{i} \).
The spaces you type in a formula are
ignored.  Remember that a letter like
       $x$                   % $ ... $  and  \( ... \)  are equivalent
is a formula when it denotes a mathematical
symbol, and it should be typed as one.

\section{Displayed Text}

Text is displayed by indenting it from the left
margin.  Quotations are commonly displayed.  There
are short quotations
\begin{quote}
   This is a short quotation.  It consists of a
   single paragraph of text.  See how it is formatted.
\end{quote}
and longer ones.
\begin{quotation}
   This is a longer quotation.  It consists of two
   paragraphs of text, neither of which are
   particularly interesting.

   This is the second paragraph of the quotation.  It
   is just as dull as the first paragraph.
\end{quotation}
Another frequently-displayed structure is a list.
The following is an example of an \emph{itemized}
list.
\begin{itemize}
   \item This is the first item of an itemized list.
         Each item in the list is marked with a ``tick''.
         You don't have to worry about what kind of tick
         mark is used.

   \item This is the second item of the list.  It
         contains another list nested inside it.  The inner
         list is an \emph{enumerated} list.
         \begin{enumerate}
            \item This is the first item of an enumerated
                  list that is nested within the itemized list.

            \item This is the second item of the inner list.
                  \LaTeX\ allows you to nest lists deeper than
                  you really should.
         \end{enumerate}
         This is the rest of the second item of the outer
         list.  It is no more interesting than any other
         part of the item.
   \item This is the third item of the list.
\end{itemize}
You can even display poetry.
\begin{verse}
   There is an environment
    for verse \\             % The \\ command separates lines
   Whose features some poets % within a stanza.
   will curse.

                             % One or more blank lines separate stanzas.

   For instead of making\\
   Them do \emph{all} line breaking, \\
   It allows them to put too many words on a line when they'd rather be
   forced to be terse.
\end{verse}

Mathematical formulas may also be displayed.  A
displayed formula
is
one-line long; multiline
formulas require special formatting instructions.
   \[  \ip{\Gamma}{\psi'} = x'' + y^{2} + z_{i}^{n}\]
Don't start a paragraph with a displayed equation,
nor make one a paragraph by itself.

\end{document}               % End of document.

\section{Bench Huge Section 764}
% synthetic bench section 764
% This is a sample LaTeX input file.  (Version of 12 August 2004.)
%
% A '%' character causes TeX to ignore all remaining text on the line,
% and is used for comments like this one.

\documentclass{article}      % Specifies the document class

                             % The preamble begins here.
\title{An Example Document}  % Declares the document's title.
\author{Leslie Lamport}      % Declares the author's name.
\date{January 21, 1994}      % Deleting this command produces today's date.

\newcommand{\ip}[2]{(#1, #2)}
                             % Defines \ip{arg1}{arg2} to mean
                             % (arg1, arg2).

%\newcommand{\ip}[2]{\langle #1 | #2\rangle}
                             % This is an alternative definition of
                             % \ip that is commented out.

\begin{document}             % End of preamble and beginning of text.

\maketitle                   % Produces the title.

This is an example input file.  Comparing it with
the output it generates can show you how to
produce a simple document of your own.

\section{Ordinary Text}      % Produces section heading.  Lower-level
                             % sections are begun with similar
                             % \subsection and \subsubsection commands.

The ends  of words and sentences are marked
  by   spaces. It  doesn't matter how many
spaces    you type; one is as good as 100.  The
end of   a line counts as a space.

One   or more   blank lines denote the  end
of  a paragraph.

Since any number of consecutive spaces are treated
like a single one, the formatting of the input
file makes no difference to
      \LaTeX,                % The \LaTeX command generates the LaTeX logo.
but it makes a difference to you.  When you use
\LaTeX, making your input file as easy to read
as possible will be a great help as you write
your document and when you change it.  This sample
file shows how you can add comments to your own input
file.

Because printing is different from typewriting,
there are a number of things that you have to do
differently when preparing an input file than if
you were just typing the document directly.
Quotation marks like
       ``this''
have to be handled specially, as do quotes within
quotes:
       ``\,`this'            % \, separates the double and single quote.
        is what I just
        wrote, not  `that'\,''.

Dashes come in three sizes: an
       intra-word
dash, a medium dash for number ranges like
       1--2,
and a punctuation
       dash---like
this.

A sentence-ending space should be larger than the
space between words within a sentence.  You
sometimes have to type special commands in
conjunction with punctuation characters to get
this right, as in the following sentence.
       Gnats, gnus, etc.\ all  % `\ ' makes an inter-word space.
       begin with G\@.         % \@ marks end-of-sentence punctuation.
You should check the spaces after periods when
reading your output to make sure you haven't
forgotten any special cases.  Generating an
ellipsis
       \ldots\               % `\ ' is needed after `\ldots' because TeX
                             % ignores spaces after command names like \ldots
                             % made from \ + letters.
                             %
                             % Note how a `%' character causes TeX to ignore
                             % the end of the input line, so these blank lines
                             % do not start a new paragraph.
                             %
with the right spacing around the periods requires
a special command.

\LaTeX\ interprets some common characters as
commands, so you must type special commands to
generate them.  These characters include the
following:
       \$ \& \% \# \{ and \}.

In printing, text is usually emphasized with an
       \emph{italic}
type style.

\begin{em}
   A long segment of text can also be emphasized
   in this way.  Text within such a segment can be
   given \emph{additional} emphasis.
\end{em}

It is sometimes necessary to prevent \LaTeX\ from
breaking a line where it might otherwise do so.
This may be at a space, as between the ``Mr.''\ and
``Jones'' in
       ``Mr.~Jones'',        % ~ produces an unbreakable interword space.
or within a word---especially when the word is a
symbol like
       \mbox{\emph{itemnum}}
that makes little sense when hyphenated across
lines.

Footnotes\footnote{This is an example of a footnote.}
pose no problem.

\LaTeX\ is good at typesetting mathematical formulas
like
       \( x-3y + z = 7 \)
or
       \( a_{1} > x^{2n} + y^{2n} > x' \)
or
       \( \ip{A}{B} = \sum_{i} a_{i} b_{i} \).
The spaces you type in a formula are
ignored.  Remember that a letter like
       $x$                   % $ ... $  and  \( ... \)  are equivalent
is a formula when it denotes a mathematical
symbol, and it should be typed as one.

\section{Displayed Text}

Text is displayed by indenting it from the left
margin.  Quotations are commonly displayed.  There
are short quotations
\begin{quote}
   This is a short quotation.  It consists of a
   single paragraph of text.  See how it is formatted.
\end{quote}
and longer ones.
\begin{quotation}
   This is a longer quotation.  It consists of two
   paragraphs of text, neither of which are
   particularly interesting.

   This is the second paragraph of the quotation.  It
   is just as dull as the first paragraph.
\end{quotation}
Another frequently-displayed structure is a list.
The following is an example of an \emph{itemized}
list.
\begin{itemize}
   \item This is the first item of an itemized list.
         Each item in the list is marked with a ``tick''.
         You don't have to worry about what kind of tick
         mark is used.

   \item This is the second item of the list.  It
         contains another list nested inside it.  The inner
         list is an \emph{enumerated} list.
         \begin{enumerate}
            \item This is the first item of an enumerated
                  list that is nested within the itemized list.

            \item This is the second item of the inner list.
                  \LaTeX\ allows you to nest lists deeper than
                  you really should.
         \end{enumerate}
         This is the rest of the second item of the outer
         list.  It is no more interesting than any other
         part of the item.
   \item This is the third item of the list.
\end{itemize}
You can even display poetry.
\begin{verse}
   There is an environment
    for verse \\             % The \\ command separates lines
   Whose features some poets % within a stanza.
   will curse.

                             % One or more blank lines separate stanzas.

   For instead of making\\
   Them do \emph{all} line breaking, \\
   It allows them to put too many words on a line when they'd rather be
   forced to be terse.
\end{verse}

Mathematical formulas may also be displayed.  A
displayed formula
is
one-line long; multiline
formulas require special formatting instructions.
   \[  \ip{\Gamma}{\psi'} = x'' + y^{2} + z_{i}^{n}\]
Don't start a paragraph with a displayed equation,
nor make one a paragraph by itself.

\end{document}               % End of document.

\section{Bench Huge Section 765}
% synthetic bench section 765
% This is a sample LaTeX input file.  (Version of 12 August 2004.)
%
% A '%' character causes TeX to ignore all remaining text on the line,
% and is used for comments like this one.

\documentclass{article}      % Specifies the document class

                             % The preamble begins here.
\title{An Example Document}  % Declares the document's title.
\author{Leslie Lamport}      % Declares the author's name.
\date{January 21, 1994}      % Deleting this command produces today's date.

\newcommand{\ip}[2]{(#1, #2)}
                             % Defines \ip{arg1}{arg2} to mean
                             % (arg1, arg2).

%\newcommand{\ip}[2]{\langle #1 | #2\rangle}
                             % This is an alternative definition of
                             % \ip that is commented out.

\begin{document}             % End of preamble and beginning of text.

\maketitle                   % Produces the title.

This is an example input file.  Comparing it with
the output it generates can show you how to
produce a simple document of your own.

\section{Ordinary Text}      % Produces section heading.  Lower-level
                             % sections are begun with similar
                             % \subsection and \subsubsection commands.

The ends  of words and sentences are marked
  by   spaces. It  doesn't matter how many
spaces    you type; one is as good as 100.  The
end of   a line counts as a space.

One   or more   blank lines denote the  end
of  a paragraph.

Since any number of consecutive spaces are treated
like a single one, the formatting of the input
file makes no difference to
      \LaTeX,                % The \LaTeX command generates the LaTeX logo.
but it makes a difference to you.  When you use
\LaTeX, making your input file as easy to read
as possible will be a great help as you write
your document and when you change it.  This sample
file shows how you can add comments to your own input
file.

Because printing is different from typewriting,
there are a number of things that you have to do
differently when preparing an input file than if
you were just typing the document directly.
Quotation marks like
       ``this''
have to be handled specially, as do quotes within
quotes:
       ``\,`this'            % \, separates the double and single quote.
        is what I just
        wrote, not  `that'\,''.

Dashes come in three sizes: an
       intra-word
dash, a medium dash for number ranges like
       1--2,
and a punctuation
       dash---like
this.

A sentence-ending space should be larger than the
space between words within a sentence.  You
sometimes have to type special commands in
conjunction with punctuation characters to get
this right, as in the following sentence.
       Gnats, gnus, etc.\ all  % `\ ' makes an inter-word space.
       begin with G\@.         % \@ marks end-of-sentence punctuation.
You should check the spaces after periods when
reading your output to make sure you haven't
forgotten any special cases.  Generating an
ellipsis
       \ldots\               % `\ ' is needed after `\ldots' because TeX
                             % ignores spaces after command names like \ldots
                             % made from \ + letters.
                             %
                             % Note how a `%' character causes TeX to ignore
                             % the end of the input line, so these blank lines
                             % do not start a new paragraph.
                             %
with the right spacing around the periods requires
a special command.

\LaTeX\ interprets some common characters as
commands, so you must type special commands to
generate them.  These characters include the
following:
       \$ \& \% \# \{ and \}.

In printing, text is usually emphasized with an
       \emph{italic}
type style.

\begin{em}
   A long segment of text can also be emphasized
   in this way.  Text within such a segment can be
   given \emph{additional} emphasis.
\end{em}

It is sometimes necessary to prevent \LaTeX\ from
breaking a line where it might otherwise do so.
This may be at a space, as between the ``Mr.''\ and
``Jones'' in
       ``Mr.~Jones'',        % ~ produces an unbreakable interword space.
or within a word---especially when the word is a
symbol like
       \mbox{\emph{itemnum}}
that makes little sense when hyphenated across
lines.

Footnotes\footnote{This is an example of a footnote.}
pose no problem.

\LaTeX\ is good at typesetting mathematical formulas
like
       \( x-3y + z = 7 \)
or
       \( a_{1} > x^{2n} + y^{2n} > x' \)
or
       \( \ip{A}{B} = \sum_{i} a_{i} b_{i} \).
The spaces you type in a formula are
ignored.  Remember that a letter like
       $x$                   % $ ... $  and  \( ... \)  are equivalent
is a formula when it denotes a mathematical
symbol, and it should be typed as one.

\section{Displayed Text}

Text is displayed by indenting it from the left
margin.  Quotations are commonly displayed.  There
are short quotations
\begin{quote}
   This is a short quotation.  It consists of a
   single paragraph of text.  See how it is formatted.
\end{quote}
and longer ones.
\begin{quotation}
   This is a longer quotation.  It consists of two
   paragraphs of text, neither of which are
   particularly interesting.

   This is the second paragraph of the quotation.  It
   is just as dull as the first paragraph.
\end{quotation}
Another frequently-displayed structure is a list.
The following is an example of an \emph{itemized}
list.
\begin{itemize}
   \item This is the first item of an itemized list.
         Each item in the list is marked with a ``tick''.
         You don't have to worry about what kind of tick
         mark is used.

   \item This is the second item of the list.  It
         contains another list nested inside it.  The inner
         list is an \emph{enumerated} list.
         \begin{enumerate}
            \item This is the first item of an enumerated
                  list that is nested within the itemized list.

            \item This is the second item of the inner list.
                  \LaTeX\ allows you to nest lists deeper than
                  you really should.
         \end{enumerate}
         This is the rest of the second item of the outer
         list.  It is no more interesting than any other
         part of the item.
   \item This is the third item of the list.
\end{itemize}
You can even display poetry.
\begin{verse}
   There is an environment
    for verse \\             % The \\ command separates lines
   Whose features some poets % within a stanza.
   will curse.

                             % One or more blank lines separate stanzas.

   For instead of making\\
   Them do \emph{all} line breaking, \\
   It allows them to put too many words on a line when they'd rather be
   forced to be terse.
\end{verse}

Mathematical formulas may also be displayed.  A
displayed formula
is
one-line long; multiline
formulas require special formatting instructions.
   \[  \ip{\Gamma}{\psi'} = x'' + y^{2} + z_{i}^{n}\]
Don't start a paragraph with a displayed equation,
nor make one a paragraph by itself.

\end{document}               % End of document.

\section{Bench Huge Section 766}
% synthetic bench section 766
% This is a sample LaTeX input file.  (Version of 12 August 2004.)
%
% A '%' character causes TeX to ignore all remaining text on the line,
% and is used for comments like this one.

\documentclass{article}      % Specifies the document class

                             % The preamble begins here.
\title{An Example Document}  % Declares the document's title.
\author{Leslie Lamport}      % Declares the author's name.
\date{January 21, 1994}      % Deleting this command produces today's date.

\newcommand{\ip}[2]{(#1, #2)}
                             % Defines \ip{arg1}{arg2} to mean
                             % (arg1, arg2).

%\newcommand{\ip}[2]{\langle #1 | #2\rangle}
                             % This is an alternative definition of
                             % \ip that is commented out.

\begin{document}             % End of preamble and beginning of text.

\maketitle                   % Produces the title.

This is an example input file.  Comparing it with
the output it generates can show you how to
produce a simple document of your own.

\section{Ordinary Text}      % Produces section heading.  Lower-level
                             % sections are begun with similar
                             % \subsection and \subsubsection commands.

The ends  of words and sentences are marked
  by   spaces. It  doesn't matter how many
spaces    you type; one is as good as 100.  The
end of   a line counts as a space.

One   or more   blank lines denote the  end
of  a paragraph.

Since any number of consecutive spaces are treated
like a single one, the formatting of the input
file makes no difference to
      \LaTeX,                % The \LaTeX command generates the LaTeX logo.
but it makes a difference to you.  When you use
\LaTeX, making your input file as easy to read
as possible will be a great help as you write
your document and when you change it.  This sample
file shows how you can add comments to your own input
file.

Because printing is different from typewriting,
there are a number of things that you have to do
differently when preparing an input file than if
you were just typing the document directly.
Quotation marks like
       ``this''
have to be handled specially, as do quotes within
quotes:
       ``\,`this'            % \, separates the double and single quote.
        is what I just
        wrote, not  `that'\,''.

Dashes come in three sizes: an
       intra-word
dash, a medium dash for number ranges like
       1--2,
and a punctuation
       dash---like
this.

A sentence-ending space should be larger than the
space between words within a sentence.  You
sometimes have to type special commands in
conjunction with punctuation characters to get
this right, as in the following sentence.
       Gnats, gnus, etc.\ all  % `\ ' makes an inter-word space.
       begin with G\@.         % \@ marks end-of-sentence punctuation.
You should check the spaces after periods when
reading your output to make sure you haven't
forgotten any special cases.  Generating an
ellipsis
       \ldots\               % `\ ' is needed after `\ldots' because TeX
                             % ignores spaces after command names like \ldots
                             % made from \ + letters.
                             %
                             % Note how a `%' character causes TeX to ignore
                             % the end of the input line, so these blank lines
                             % do not start a new paragraph.
                             %
with the right spacing around the periods requires
a special command.

\LaTeX\ interprets some common characters as
commands, so you must type special commands to
generate them.  These characters include the
following:
       \$ \& \% \# \{ and \}.

In printing, text is usually emphasized with an
       \emph{italic}
type style.

\begin{em}
   A long segment of text can also be emphasized
   in this way.  Text within such a segment can be
   given \emph{additional} emphasis.
\end{em}

It is sometimes necessary to prevent \LaTeX\ from
breaking a line where it might otherwise do so.
This may be at a space, as between the ``Mr.''\ and
``Jones'' in
       ``Mr.~Jones'',        % ~ produces an unbreakable interword space.
or within a word---especially when the word is a
symbol like
       \mbox{\emph{itemnum}}
that makes little sense when hyphenated across
lines.

Footnotes\footnote{This is an example of a footnote.}
pose no problem.

\LaTeX\ is good at typesetting mathematical formulas
like
       \( x-3y + z = 7 \)
or
       \( a_{1} > x^{2n} + y^{2n} > x' \)
or
       \( \ip{A}{B} = \sum_{i} a_{i} b_{i} \).
The spaces you type in a formula are
ignored.  Remember that a letter like
       $x$                   % $ ... $  and  \( ... \)  are equivalent
is a formula when it denotes a mathematical
symbol, and it should be typed as one.

\section{Displayed Text}

Text is displayed by indenting it from the left
margin.  Quotations are commonly displayed.  There
are short quotations
\begin{quote}
   This is a short quotation.  It consists of a
   single paragraph of text.  See how it is formatted.
\end{quote}
and longer ones.
\begin{quotation}
   This is a longer quotation.  It consists of two
   paragraphs of text, neither of which are
   particularly interesting.

   This is the second paragraph of the quotation.  It
   is just as dull as the first paragraph.
\end{quotation}
Another frequently-displayed structure is a list.
The following is an example of an \emph{itemized}
list.
\begin{itemize}
   \item This is the first item of an itemized list.
         Each item in the list is marked with a ``tick''.
         You don't have to worry about what kind of tick
         mark is used.

   \item This is the second item of the list.  It
         contains another list nested inside it.  The inner
         list is an \emph{enumerated} list.
         \begin{enumerate}
            \item This is the first item of an enumerated
                  list that is nested within the itemized list.

            \item This is the second item of the inner list.
                  \LaTeX\ allows you to nest lists deeper than
                  you really should.
         \end{enumerate}
         This is the rest of the second item of the outer
         list.  It is no more interesting than any other
         part of the item.
   \item This is the third item of the list.
\end{itemize}
You can even display poetry.
\begin{verse}
   There is an environment
    for verse \\             % The \\ command separates lines
   Whose features some poets % within a stanza.
   will curse.

                             % One or more blank lines separate stanzas.

   For instead of making\\
   Them do \emph{all} line breaking, \\
   It allows them to put too many words on a line when they'd rather be
   forced to be terse.
\end{verse}

Mathematical formulas may also be displayed.  A
displayed formula
is
one-line long; multiline
formulas require special formatting instructions.
   \[  \ip{\Gamma}{\psi'} = x'' + y^{2} + z_{i}^{n}\]
Don't start a paragraph with a displayed equation,
nor make one a paragraph by itself.

\end{document}               % End of document.

\section{Bench Huge Section 767}
% synthetic bench section 767
% This is a sample LaTeX input file.  (Version of 12 August 2004.)
%
% A '%' character causes TeX to ignore all remaining text on the line,
% and is used for comments like this one.

\documentclass{article}      % Specifies the document class

                             % The preamble begins here.
\title{An Example Document}  % Declares the document's title.
\author{Leslie Lamport}      % Declares the author's name.
\date{January 21, 1994}      % Deleting this command produces today's date.

\newcommand{\ip}[2]{(#1, #2)}
                             % Defines \ip{arg1}{arg2} to mean
                             % (arg1, arg2).

%\newcommand{\ip}[2]{\langle #1 | #2\rangle}
                             % This is an alternative definition of
                             % \ip that is commented out.

\begin{document}             % End of preamble and beginning of text.

\maketitle                   % Produces the title.

This is an example input file.  Comparing it with
the output it generates can show you how to
produce a simple document of your own.

\section{Ordinary Text}      % Produces section heading.  Lower-level
                             % sections are begun with similar
                             % \subsection and \subsubsection commands.

The ends  of words and sentences are marked
  by   spaces. It  doesn't matter how many
spaces    you type; one is as good as 100.  The
end of   a line counts as a space.

One   or more   blank lines denote the  end
of  a paragraph.

Since any number of consecutive spaces are treated
like a single one, the formatting of the input
file makes no difference to
      \LaTeX,                % The \LaTeX command generates the LaTeX logo.
but it makes a difference to you.  When you use
\LaTeX, making your input file as easy to read
as possible will be a great help as you write
your document and when you change it.  This sample
file shows how you can add comments to your own input
file.

Because printing is different from typewriting,
there are a number of things that you have to do
differently when preparing an input file than if
you were just typing the document directly.
Quotation marks like
       ``this''
have to be handled specially, as do quotes within
quotes:
       ``\,`this'            % \, separates the double and single quote.
        is what I just
        wrote, not  `that'\,''.

Dashes come in three sizes: an
       intra-word
dash, a medium dash for number ranges like
       1--2,
and a punctuation
       dash---like
this.

A sentence-ending space should be larger than the
space between words within a sentence.  You
sometimes have to type special commands in
conjunction with punctuation characters to get
this right, as in the following sentence.
       Gnats, gnus, etc.\ all  % `\ ' makes an inter-word space.
       begin with G\@.         % \@ marks end-of-sentence punctuation.
You should check the spaces after periods when
reading your output to make sure you haven't
forgotten any special cases.  Generating an
ellipsis
       \ldots\               % `\ ' is needed after `\ldots' because TeX
                             % ignores spaces after command names like \ldots
                             % made from \ + letters.
                             %
                             % Note how a `%' character causes TeX to ignore
                             % the end of the input line, so these blank lines
                             % do not start a new paragraph.
                             %
with the right spacing around the periods requires
a special command.

\LaTeX\ interprets some common characters as
commands, so you must type special commands to
generate them.  These characters include the
following:
       \$ \& \% \# \{ and \}.

In printing, text is usually emphasized with an
       \emph{italic}
type style.

\begin{em}
   A long segment of text can also be emphasized
   in this way.  Text within such a segment can be
   given \emph{additional} emphasis.
\end{em}

It is sometimes necessary to prevent \LaTeX\ from
breaking a line where it might otherwise do so.
This may be at a space, as between the ``Mr.''\ and
``Jones'' in
       ``Mr.~Jones'',        % ~ produces an unbreakable interword space.
or within a word---especially when the word is a
symbol like
       \mbox{\emph{itemnum}}
that makes little sense when hyphenated across
lines.

Footnotes\footnote{This is an example of a footnote.}
pose no problem.

\LaTeX\ is good at typesetting mathematical formulas
like
       \( x-3y + z = 7 \)
or
       \( a_{1} > x^{2n} + y^{2n} > x' \)
or
       \( \ip{A}{B} = \sum_{i} a_{i} b_{i} \).
The spaces you type in a formula are
ignored.  Remember that a letter like
       $x$                   % $ ... $  and  \( ... \)  are equivalent
is a formula when it denotes a mathematical
symbol, and it should be typed as one.

\section{Displayed Text}

Text is displayed by indenting it from the left
margin.  Quotations are commonly displayed.  There
are short quotations
\begin{quote}
   This is a short quotation.  It consists of a
   single paragraph of text.  See how it is formatted.
\end{quote}
and longer ones.
\begin{quotation}
   This is a longer quotation.  It consists of two
   paragraphs of text, neither of which are
   particularly interesting.

   This is the second paragraph of the quotation.  It
   is just as dull as the first paragraph.
\end{quotation}
Another frequently-displayed structure is a list.
The following is an example of an \emph{itemized}
list.
\begin{itemize}
   \item This is the first item of an itemized list.
         Each item in the list is marked with a ``tick''.
         You don't have to worry about what kind of tick
         mark is used.

   \item This is the second item of the list.  It
         contains another list nested inside it.  The inner
         list is an \emph{enumerated} list.
         \begin{enumerate}
            \item This is the first item of an enumerated
                  list that is nested within the itemized list.

            \item This is the second item of the inner list.
                  \LaTeX\ allows you to nest lists deeper than
                  you really should.
         \end{enumerate}
         This is the rest of the second item of the outer
         list.  It is no more interesting than any other
         part of the item.
   \item This is the third item of the list.
\end{itemize}
You can even display poetry.
\begin{verse}
   There is an environment
    for verse \\             % The \\ command separates lines
   Whose features some poets % within a stanza.
   will curse.

                             % One or more blank lines separate stanzas.

   For instead of making\\
   Them do \emph{all} line breaking, \\
   It allows them to put too many words on a line when they'd rather be
   forced to be terse.
\end{verse}

Mathematical formulas may also be displayed.  A
displayed formula
is
one-line long; multiline
formulas require special formatting instructions.
   \[  \ip{\Gamma}{\psi'} = x'' + y^{2} + z_{i}^{n}\]
Don't start a paragraph with a displayed equation,
nor make one a paragraph by itself.

\end{document}               % End of document.

\section{Bench Huge Section 768}
% synthetic bench section 768
% This is a sample LaTeX input file.  (Version of 12 August 2004.)
%
% A '%' character causes TeX to ignore all remaining text on the line,
% and is used for comments like this one.

\documentclass{article}      % Specifies the document class

                             % The preamble begins here.
\title{An Example Document}  % Declares the document's title.
\author{Leslie Lamport}      % Declares the author's name.
\date{January 21, 1994}      % Deleting this command produces today's date.

\newcommand{\ip}[2]{(#1, #2)}
                             % Defines \ip{arg1}{arg2} to mean
                             % (arg1, arg2).

%\newcommand{\ip}[2]{\langle #1 | #2\rangle}
                             % This is an alternative definition of
                             % \ip that is commented out.

\begin{document}             % End of preamble and beginning of text.

\maketitle                   % Produces the title.

This is an example input file.  Comparing it with
the output it generates can show you how to
produce a simple document of your own.

\section{Ordinary Text}      % Produces section heading.  Lower-level
                             % sections are begun with similar
                             % \subsection and \subsubsection commands.

The ends  of words and sentences are marked
  by   spaces. It  doesn't matter how many
spaces    you type; one is as good as 100.  The
end of   a line counts as a space.

One   or more   blank lines denote the  end
of  a paragraph.

Since any number of consecutive spaces are treated
like a single one, the formatting of the input
file makes no difference to
      \LaTeX,                % The \LaTeX command generates the LaTeX logo.
but it makes a difference to you.  When you use
\LaTeX, making your input file as easy to read
as possible will be a great help as you write
your document and when you change it.  This sample
file shows how you can add comments to your own input
file.

Because printing is different from typewriting,
there are a number of things that you have to do
differently when preparing an input file than if
you were just typing the document directly.
Quotation marks like
       ``this''
have to be handled specially, as do quotes within
quotes:
       ``\,`this'            % \, separates the double and single quote.
        is what I just
        wrote, not  `that'\,''.

Dashes come in three sizes: an
       intra-word
dash, a medium dash for number ranges like
       1--2,
and a punctuation
       dash---like
this.

A sentence-ending space should be larger than the
space between words within a sentence.  You
sometimes have to type special commands in
conjunction with punctuation characters to get
this right, as in the following sentence.
       Gnats, gnus, etc.\ all  % `\ ' makes an inter-word space.
       begin with G\@.         % \@ marks end-of-sentence punctuation.
You should check the spaces after periods when
reading your output to make sure you haven't
forgotten any special cases.  Generating an
ellipsis
       \ldots\               % `\ ' is needed after `\ldots' because TeX
                             % ignores spaces after command names like \ldots
                             % made from \ + letters.
                             %
                             % Note how a `%' character causes TeX to ignore
                             % the end of the input line, so these blank lines
                             % do not start a new paragraph.
                             %
with the right spacing around the periods requires
a special command.

\LaTeX\ interprets some common characters as
commands, so you must type special commands to
generate them.  These characters include the
following:
       \$ \& \% \# \{ and \}.

In printing, text is usually emphasized with an
       \emph{italic}
type style.

\begin{em}
   A long segment of text can also be emphasized
   in this way.  Text within such a segment can be
   given \emph{additional} emphasis.
\end{em}

It is sometimes necessary to prevent \LaTeX\ from
breaking a line where it might otherwise do so.
This may be at a space, as between the ``Mr.''\ and
``Jones'' in
       ``Mr.~Jones'',        % ~ produces an unbreakable interword space.
or within a word---especially when the word is a
symbol like
       \mbox{\emph{itemnum}}
that makes little sense when hyphenated across
lines.

Footnotes\footnote{This is an example of a footnote.}
pose no problem.

\LaTeX\ is good at typesetting mathematical formulas
like
       \( x-3y + z = 7 \)
or
       \( a_{1} > x^{2n} + y^{2n} > x' \)
or
       \( \ip{A}{B} = \sum_{i} a_{i} b_{i} \).
The spaces you type in a formula are
ignored.  Remember that a letter like
       $x$                   % $ ... $  and  \( ... \)  are equivalent
is a formula when it denotes a mathematical
symbol, and it should be typed as one.

\section{Displayed Text}

Text is displayed by indenting it from the left
margin.  Quotations are commonly displayed.  There
are short quotations
\begin{quote}
   This is a short quotation.  It consists of a
   single paragraph of text.  See how it is formatted.
\end{quote}
and longer ones.
\begin{quotation}
   This is a longer quotation.  It consists of two
   paragraphs of text, neither of which are
   particularly interesting.

   This is the second paragraph of the quotation.  It
   is just as dull as the first paragraph.
\end{quotation}
Another frequently-displayed structure is a list.
The following is an example of an \emph{itemized}
list.
\begin{itemize}
   \item This is the first item of an itemized list.
         Each item in the list is marked with a ``tick''.
         You don't have to worry about what kind of tick
         mark is used.

   \item This is the second item of the list.  It
         contains another list nested inside it.  The inner
         list is an \emph{enumerated} list.
         \begin{enumerate}
            \item This is the first item of an enumerated
                  list that is nested within the itemized list.

            \item This is the second item of the inner list.
                  \LaTeX\ allows you to nest lists deeper than
                  you really should.
         \end{enumerate}
         This is the rest of the second item of the outer
         list.  It is no more interesting than any other
         part of the item.
   \item This is the third item of the list.
\end{itemize}
You can even display poetry.
\begin{verse}
   There is an environment
    for verse \\             % The \\ command separates lines
   Whose features some poets % within a stanza.
   will curse.

                             % One or more blank lines separate stanzas.

   For instead of making\\
   Them do \emph{all} line breaking, \\
   It allows them to put too many words on a line when they'd rather be
   forced to be terse.
\end{verse}

Mathematical formulas may also be displayed.  A
displayed formula
is
one-line long; multiline
formulas require special formatting instructions.
   \[  \ip{\Gamma}{\psi'} = x'' + y^{2} + z_{i}^{n}\]
Don't start a paragraph with a displayed equation,
nor make one a paragraph by itself.

\end{document}               % End of document.

\section{Bench Huge Section 769}
% synthetic bench section 769
% This is a sample LaTeX input file.  (Version of 12 August 2004.)
%
% A '%' character causes TeX to ignore all remaining text on the line,
% and is used for comments like this one.

\documentclass{article}      % Specifies the document class

                             % The preamble begins here.
\title{An Example Document}  % Declares the document's title.
\author{Leslie Lamport}      % Declares the author's name.
\date{January 21, 1994}      % Deleting this command produces today's date.

\newcommand{\ip}[2]{(#1, #2)}
                             % Defines \ip{arg1}{arg2} to mean
                             % (arg1, arg2).

%\newcommand{\ip}[2]{\langle #1 | #2\rangle}
                             % This is an alternative definition of
                             % \ip that is commented out.

\begin{document}             % End of preamble and beginning of text.

\maketitle                   % Produces the title.

This is an example input file.  Comparing it with
the output it generates can show you how to
produce a simple document of your own.

\section{Ordinary Text}      % Produces section heading.  Lower-level
                             % sections are begun with similar
                             % \subsection and \subsubsection commands.

The ends  of words and sentences are marked
  by   spaces. It  doesn't matter how many
spaces    you type; one is as good as 100.  The
end of   a line counts as a space.

One   or more   blank lines denote the  end
of  a paragraph.

Since any number of consecutive spaces are treated
like a single one, the formatting of the input
file makes no difference to
      \LaTeX,                % The \LaTeX command generates the LaTeX logo.
but it makes a difference to you.  When you use
\LaTeX, making your input file as easy to read
as possible will be a great help as you write
your document and when you change it.  This sample
file shows how you can add comments to your own input
file.

Because printing is different from typewriting,
there are a number of things that you have to do
differently when preparing an input file than if
you were just typing the document directly.
Quotation marks like
       ``this''
have to be handled specially, as do quotes within
quotes:
       ``\,`this'            % \, separates the double and single quote.
        is what I just
        wrote, not  `that'\,''.

Dashes come in three sizes: an
       intra-word
dash, a medium dash for number ranges like
       1--2,
and a punctuation
       dash---like
this.

A sentence-ending space should be larger than the
space between words within a sentence.  You
sometimes have to type special commands in
conjunction with punctuation characters to get
this right, as in the following sentence.
       Gnats, gnus, etc.\ all  % `\ ' makes an inter-word space.
       begin with G\@.         % \@ marks end-of-sentence punctuation.
You should check the spaces after periods when
reading your output to make sure you haven't
forgotten any special cases.  Generating an
ellipsis
       \ldots\               % `\ ' is needed after `\ldots' because TeX
                             % ignores spaces after command names like \ldots
                             % made from \ + letters.
                             %
                             % Note how a `%' character causes TeX to ignore
                             % the end of the input line, so these blank lines
                             % do not start a new paragraph.
                             %
with the right spacing around the periods requires
a special command.

\LaTeX\ interprets some common characters as
commands, so you must type special commands to
generate them.  These characters include the
following:
       \$ \& \% \# \{ and \}.

In printing, text is usually emphasized with an
       \emph{italic}
type style.

\begin{em}
   A long segment of text can also be emphasized
   in this way.  Text within such a segment can be
   given \emph{additional} emphasis.
\end{em}

It is sometimes necessary to prevent \LaTeX\ from
breaking a line where it might otherwise do so.
This may be at a space, as between the ``Mr.''\ and
``Jones'' in
       ``Mr.~Jones'',        % ~ produces an unbreakable interword space.
or within a word---especially when the word is a
symbol like
       \mbox{\emph{itemnum}}
that makes little sense when hyphenated across
lines.

Footnotes\footnote{This is an example of a footnote.}
pose no problem.

\LaTeX\ is good at typesetting mathematical formulas
like
       \( x-3y + z = 7 \)
or
       \( a_{1} > x^{2n} + y^{2n} > x' \)
or
       \( \ip{A}{B} = \sum_{i} a_{i} b_{i} \).
The spaces you type in a formula are
ignored.  Remember that a letter like
       $x$                   % $ ... $  and  \( ... \)  are equivalent
is a formula when it denotes a mathematical
symbol, and it should be typed as one.

\section{Displayed Text}

Text is displayed by indenting it from the left
margin.  Quotations are commonly displayed.  There
are short quotations
\begin{quote}
   This is a short quotation.  It consists of a
   single paragraph of text.  See how it is formatted.
\end{quote}
and longer ones.
\begin{quotation}
   This is a longer quotation.  It consists of two
   paragraphs of text, neither of which are
   particularly interesting.

   This is the second paragraph of the quotation.  It
   is just as dull as the first paragraph.
\end{quotation}
Another frequently-displayed structure is a list.
The following is an example of an \emph{itemized}
list.
\begin{itemize}
   \item This is the first item of an itemized list.
         Each item in the list is marked with a ``tick''.
         You don't have to worry about what kind of tick
         mark is used.

   \item This is the second item of the list.  It
         contains another list nested inside it.  The inner
         list is an \emph{enumerated} list.
         \begin{enumerate}
            \item This is the first item of an enumerated
                  list that is nested within the itemized list.

            \item This is the second item of the inner list.
                  \LaTeX\ allows you to nest lists deeper than
                  you really should.
         \end{enumerate}
         This is the rest of the second item of the outer
         list.  It is no more interesting than any other
         part of the item.
   \item This is the third item of the list.
\end{itemize}
You can even display poetry.
\begin{verse}
   There is an environment
    for verse \\             % The \\ command separates lines
   Whose features some poets % within a stanza.
   will curse.

                             % One or more blank lines separate stanzas.

   For instead of making\\
   Them do \emph{all} line breaking, \\
   It allows them to put too many words on a line when they'd rather be
   forced to be terse.
\end{verse}

Mathematical formulas may also be displayed.  A
displayed formula
is
one-line long; multiline
formulas require special formatting instructions.
   \[  \ip{\Gamma}{\psi'} = x'' + y^{2} + z_{i}^{n}\]
Don't start a paragraph with a displayed equation,
nor make one a paragraph by itself.

\end{document}               % End of document.

\section{Bench Huge Section 770}
% synthetic bench section 770
% This is a sample LaTeX input file.  (Version of 12 August 2004.)
%
% A '%' character causes TeX to ignore all remaining text on the line,
% and is used for comments like this one.

\documentclass{article}      % Specifies the document class

                             % The preamble begins here.
\title{An Example Document}  % Declares the document's title.
\author{Leslie Lamport}      % Declares the author's name.
\date{January 21, 1994}      % Deleting this command produces today's date.

\newcommand{\ip}[2]{(#1, #2)}
                             % Defines \ip{arg1}{arg2} to mean
                             % (arg1, arg2).

%\newcommand{\ip}[2]{\langle #1 | #2\rangle}
                             % This is an alternative definition of
                             % \ip that is commented out.

\begin{document}             % End of preamble and beginning of text.

\maketitle                   % Produces the title.

This is an example input file.  Comparing it with
the output it generates can show you how to
produce a simple document of your own.

\section{Ordinary Text}      % Produces section heading.  Lower-level
                             % sections are begun with similar
                             % \subsection and \subsubsection commands.

The ends  of words and sentences are marked
  by   spaces. It  doesn't matter how many
spaces    you type; one is as good as 100.  The
end of   a line counts as a space.

One   or more   blank lines denote the  end
of  a paragraph.

Since any number of consecutive spaces are treated
like a single one, the formatting of the input
file makes no difference to
      \LaTeX,                % The \LaTeX command generates the LaTeX logo.
but it makes a difference to you.  When you use
\LaTeX, making your input file as easy to read
as possible will be a great help as you write
your document and when you change it.  This sample
file shows how you can add comments to your own input
file.

Because printing is different from typewriting,
there are a number of things that you have to do
differently when preparing an input file than if
you were just typing the document directly.
Quotation marks like
       ``this''
have to be handled specially, as do quotes within
quotes:
       ``\,`this'            % \, separates the double and single quote.
        is what I just
        wrote, not  `that'\,''.

Dashes come in three sizes: an
       intra-word
dash, a medium dash for number ranges like
       1--2,
and a punctuation
       dash---like
this.

A sentence-ending space should be larger than the
space between words within a sentence.  You
sometimes have to type special commands in
conjunction with punctuation characters to get
this right, as in the following sentence.
       Gnats, gnus, etc.\ all  % `\ ' makes an inter-word space.
       begin with G\@.         % \@ marks end-of-sentence punctuation.
You should check the spaces after periods when
reading your output to make sure you haven't
forgotten any special cases.  Generating an
ellipsis
       \ldots\               % `\ ' is needed after `\ldots' because TeX
                             % ignores spaces after command names like \ldots
                             % made from \ + letters.
                             %
                             % Note how a `%' character causes TeX to ignore
                             % the end of the input line, so these blank lines
                             % do not start a new paragraph.
                             %
with the right spacing around the periods requires
a special command.

\LaTeX\ interprets some common characters as
commands, so you must type special commands to
generate them.  These characters include the
following:
       \$ \& \% \# \{ and \}.

In printing, text is usually emphasized with an
       \emph{italic}
type style.

\begin{em}
   A long segment of text can also be emphasized
   in this way.  Text within such a segment can be
   given \emph{additional} emphasis.
\end{em}

It is sometimes necessary to prevent \LaTeX\ from
breaking a line where it might otherwise do so.
This may be at a space, as between the ``Mr.''\ and
``Jones'' in
       ``Mr.~Jones'',        % ~ produces an unbreakable interword space.
or within a word---especially when the word is a
symbol like
       \mbox{\emph{itemnum}}
that makes little sense when hyphenated across
lines.

Footnotes\footnote{This is an example of a footnote.}
pose no problem.

\LaTeX\ is good at typesetting mathematical formulas
like
       \( x-3y + z = 7 \)
or
       \( a_{1} > x^{2n} + y^{2n} > x' \)
or
       \( \ip{A}{B} = \sum_{i} a_{i} b_{i} \).
The spaces you type in a formula are
ignored.  Remember that a letter like
       $x$                   % $ ... $  and  \( ... \)  are equivalent
is a formula when it denotes a mathematical
symbol, and it should be typed as one.

\section{Displayed Text}

Text is displayed by indenting it from the left
margin.  Quotations are commonly displayed.  There
are short quotations
\begin{quote}
   This is a short quotation.  It consists of a
   single paragraph of text.  See how it is formatted.
\end{quote}
and longer ones.
\begin{quotation}
   This is a longer quotation.  It consists of two
   paragraphs of text, neither of which are
   particularly interesting.

   This is the second paragraph of the quotation.  It
   is just as dull as the first paragraph.
\end{quotation}
Another frequently-displayed structure is a list.
The following is an example of an \emph{itemized}
list.
\begin{itemize}
   \item This is the first item of an itemized list.
         Each item in the list is marked with a ``tick''.
         You don't have to worry about what kind of tick
         mark is used.

   \item This is the second item of the list.  It
         contains another list nested inside it.  The inner
         list is an \emph{enumerated} list.
         \begin{enumerate}
            \item This is the first item of an enumerated
                  list that is nested within the itemized list.

            \item This is the second item of the inner list.
                  \LaTeX\ allows you to nest lists deeper than
                  you really should.
         \end{enumerate}
         This is the rest of the second item of the outer
         list.  It is no more interesting than any other
         part of the item.
   \item This is the third item of the list.
\end{itemize}
You can even display poetry.
\begin{verse}
   There is an environment
    for verse \\             % The \\ command separates lines
   Whose features some poets % within a stanza.
   will curse.

                             % One or more blank lines separate stanzas.

   For instead of making\\
   Them do \emph{all} line breaking, \\
   It allows them to put too many words on a line when they'd rather be
   forced to be terse.
\end{verse}

Mathematical formulas may also be displayed.  A
displayed formula
is
one-line long; multiline
formulas require special formatting instructions.
   \[  \ip{\Gamma}{\psi'} = x'' + y^{2} + z_{i}^{n}\]
Don't start a paragraph with a displayed equation,
nor make one a paragraph by itself.

\end{document}               % End of document.

\section{Bench Huge Section 771}
% synthetic bench section 771
% This is a sample LaTeX input file.  (Version of 12 August 2004.)
%
% A '%' character causes TeX to ignore all remaining text on the line,
% and is used for comments like this one.

\documentclass{article}      % Specifies the document class

                             % The preamble begins here.
\title{An Example Document}  % Declares the document's title.
\author{Leslie Lamport}      % Declares the author's name.
\date{January 21, 1994}      % Deleting this command produces today's date.

\newcommand{\ip}[2]{(#1, #2)}
                             % Defines \ip{arg1}{arg2} to mean
                             % (arg1, arg2).

%\newcommand{\ip}[2]{\langle #1 | #2\rangle}
                             % This is an alternative definition of
                             % \ip that is commented out.

\begin{document}             % End of preamble and beginning of text.

\maketitle                   % Produces the title.

This is an example input file.  Comparing it with
the output it generates can show you how to
produce a simple document of your own.

\section{Ordinary Text}      % Produces section heading.  Lower-level
                             % sections are begun with similar
                             % \subsection and \subsubsection commands.

The ends  of words and sentences are marked
  by   spaces. It  doesn't matter how many
spaces    you type; one is as good as 100.  The
end of   a line counts as a space.

One   or more   blank lines denote the  end
of  a paragraph.

Since any number of consecutive spaces are treated
like a single one, the formatting of the input
file makes no difference to
      \LaTeX,                % The \LaTeX command generates the LaTeX logo.
but it makes a difference to you.  When you use
\LaTeX, making your input file as easy to read
as possible will be a great help as you write
your document and when you change it.  This sample
file shows how you can add comments to your own input
file.

Because printing is different from typewriting,
there are a number of things that you have to do
differently when preparing an input file than if
you were just typing the document directly.
Quotation marks like
       ``this''
have to be handled specially, as do quotes within
quotes:
       ``\,`this'            % \, separates the double and single quote.
        is what I just
        wrote, not  `that'\,''.

Dashes come in three sizes: an
       intra-word
dash, a medium dash for number ranges like
       1--2,
and a punctuation
       dash---like
this.

A sentence-ending space should be larger than the
space between words within a sentence.  You
sometimes have to type special commands in
conjunction with punctuation characters to get
this right, as in the following sentence.
       Gnats, gnus, etc.\ all  % `\ ' makes an inter-word space.
       begin with G\@.         % \@ marks end-of-sentence punctuation.
You should check the spaces after periods when
reading your output to make sure you haven't
forgotten any special cases.  Generating an
ellipsis
       \ldots\               % `\ ' is needed after `\ldots' because TeX
                             % ignores spaces after command names like \ldots
                             % made from \ + letters.
                             %
                             % Note how a `%' character causes TeX to ignore
                             % the end of the input line, so these blank lines
                             % do not start a new paragraph.
                             %
with the right spacing around the periods requires
a special command.

\LaTeX\ interprets some common characters as
commands, so you must type special commands to
generate them.  These characters include the
following:
       \$ \& \% \# \{ and \}.

In printing, text is usually emphasized with an
       \emph{italic}
type style.

\begin{em}
   A long segment of text can also be emphasized
   in this way.  Text within such a segment can be
   given \emph{additional} emphasis.
\end{em}

It is sometimes necessary to prevent \LaTeX\ from
breaking a line where it might otherwise do so.
This may be at a space, as between the ``Mr.''\ and
``Jones'' in
       ``Mr.~Jones'',        % ~ produces an unbreakable interword space.
or within a word---especially when the word is a
symbol like
       \mbox{\emph{itemnum}}
that makes little sense when hyphenated across
lines.

Footnotes\footnote{This is an example of a footnote.}
pose no problem.

\LaTeX\ is good at typesetting mathematical formulas
like
       \( x-3y + z = 7 \)
or
       \( a_{1} > x^{2n} + y^{2n} > x' \)
or
       \( \ip{A}{B} = \sum_{i} a_{i} b_{i} \).
The spaces you type in a formula are
ignored.  Remember that a letter like
       $x$                   % $ ... $  and  \( ... \)  are equivalent
is a formula when it denotes a mathematical
symbol, and it should be typed as one.

\section{Displayed Text}

Text is displayed by indenting it from the left
margin.  Quotations are commonly displayed.  There
are short quotations
\begin{quote}
   This is a short quotation.  It consists of a
   single paragraph of text.  See how it is formatted.
\end{quote}
and longer ones.
\begin{quotation}
   This is a longer quotation.  It consists of two
   paragraphs of text, neither of which are
   particularly interesting.

   This is the second paragraph of the quotation.  It
   is just as dull as the first paragraph.
\end{quotation}
Another frequently-displayed structure is a list.
The following is an example of an \emph{itemized}
list.
\begin{itemize}
   \item This is the first item of an itemized list.
         Each item in the list is marked with a ``tick''.
         You don't have to worry about what kind of tick
         mark is used.

   \item This is the second item of the list.  It
         contains another list nested inside it.  The inner
         list is an \emph{enumerated} list.
         \begin{enumerate}
            \item This is the first item of an enumerated
                  list that is nested within the itemized list.

            \item This is the second item of the inner list.
                  \LaTeX\ allows you to nest lists deeper than
                  you really should.
         \end{enumerate}
         This is the rest of the second item of the outer
         list.  It is no more interesting than any other
         part of the item.
   \item This is the third item of the list.
\end{itemize}
You can even display poetry.
\begin{verse}
   There is an environment
    for verse \\             % The \\ command separates lines
   Whose features some poets % within a stanza.
   will curse.

                             % One or more blank lines separate stanzas.

   For instead of making\\
   Them do \emph{all} line breaking, \\
   It allows them to put too many words on a line when they'd rather be
   forced to be terse.
\end{verse}

Mathematical formulas may also be displayed.  A
displayed formula
is
one-line long; multiline
formulas require special formatting instructions.
   \[  \ip{\Gamma}{\psi'} = x'' + y^{2} + z_{i}^{n}\]
Don't start a paragraph with a displayed equation,
nor make one a paragraph by itself.

\end{document}               % End of document.

\section{Bench Huge Section 772}
% synthetic bench section 772
% This is a sample LaTeX input file.  (Version of 12 August 2004.)
%
% A '%' character causes TeX to ignore all remaining text on the line,
% and is used for comments like this one.

\documentclass{article}      % Specifies the document class

                             % The preamble begins here.
\title{An Example Document}  % Declares the document's title.
\author{Leslie Lamport}      % Declares the author's name.
\date{January 21, 1994}      % Deleting this command produces today's date.

\newcommand{\ip}[2]{(#1, #2)}
                             % Defines \ip{arg1}{arg2} to mean
                             % (arg1, arg2).

%\newcommand{\ip}[2]{\langle #1 | #2\rangle}
                             % This is an alternative definition of
                             % \ip that is commented out.

\begin{document}             % End of preamble and beginning of text.

\maketitle                   % Produces the title.

This is an example input file.  Comparing it with
the output it generates can show you how to
produce a simple document of your own.

\section{Ordinary Text}      % Produces section heading.  Lower-level
                             % sections are begun with similar
                             % \subsection and \subsubsection commands.

The ends  of words and sentences are marked
  by   spaces. It  doesn't matter how many
spaces    you type; one is as good as 100.  The
end of   a line counts as a space.

One   or more   blank lines denote the  end
of  a paragraph.

Since any number of consecutive spaces are treated
like a single one, the formatting of the input
file makes no difference to
      \LaTeX,                % The \LaTeX command generates the LaTeX logo.
but it makes a difference to you.  When you use
\LaTeX, making your input file as easy to read
as possible will be a great help as you write
your document and when you change it.  This sample
file shows how you can add comments to your own input
file.

Because printing is different from typewriting,
there are a number of things that you have to do
differently when preparing an input file than if
you were just typing the document directly.
Quotation marks like
       ``this''
have to be handled specially, as do quotes within
quotes:
       ``\,`this'            % \, separates the double and single quote.
        is what I just
        wrote, not  `that'\,''.

Dashes come in three sizes: an
       intra-word
dash, a medium dash for number ranges like
       1--2,
and a punctuation
       dash---like
this.

A sentence-ending space should be larger than the
space between words within a sentence.  You
sometimes have to type special commands in
conjunction with punctuation characters to get
this right, as in the following sentence.
       Gnats, gnus, etc.\ all  % `\ ' makes an inter-word space.
       begin with G\@.         % \@ marks end-of-sentence punctuation.
You should check the spaces after periods when
reading your output to make sure you haven't
forgotten any special cases.  Generating an
ellipsis
       \ldots\               % `\ ' is needed after `\ldots' because TeX
                             % ignores spaces after command names like \ldots
                             % made from \ + letters.
                             %
                             % Note how a `%' character causes TeX to ignore
                             % the end of the input line, so these blank lines
                             % do not start a new paragraph.
                             %
with the right spacing around the periods requires
a special command.

\LaTeX\ interprets some common characters as
commands, so you must type special commands to
generate them.  These characters include the
following:
       \$ \& \% \# \{ and \}.

In printing, text is usually emphasized with an
       \emph{italic}
type style.

\begin{em}
   A long segment of text can also be emphasized
   in this way.  Text within such a segment can be
   given \emph{additional} emphasis.
\end{em}

It is sometimes necessary to prevent \LaTeX\ from
breaking a line where it might otherwise do so.
This may be at a space, as between the ``Mr.''\ and
``Jones'' in
       ``Mr.~Jones'',        % ~ produces an unbreakable interword space.
or within a word---especially when the word is a
symbol like
       \mbox{\emph{itemnum}}
that makes little sense when hyphenated across
lines.

Footnotes\footnote{This is an example of a footnote.}
pose no problem.

\LaTeX\ is good at typesetting mathematical formulas
like
       \( x-3y + z = 7 \)
or
       \( a_{1} > x^{2n} + y^{2n} > x' \)
or
       \( \ip{A}{B} = \sum_{i} a_{i} b_{i} \).
The spaces you type in a formula are
ignored.  Remember that a letter like
       $x$                   % $ ... $  and  \( ... \)  are equivalent
is a formula when it denotes a mathematical
symbol, and it should be typed as one.

\section{Displayed Text}

Text is displayed by indenting it from the left
margin.  Quotations are commonly displayed.  There
are short quotations
\begin{quote}
   This is a short quotation.  It consists of a
   single paragraph of text.  See how it is formatted.
\end{quote}
and longer ones.
\begin{quotation}
   This is a longer quotation.  It consists of two
   paragraphs of text, neither of which are
   particularly interesting.

   This is the second paragraph of the quotation.  It
   is just as dull as the first paragraph.
\end{quotation}
Another frequently-displayed structure is a list.
The following is an example of an \emph{itemized}
list.
\begin{itemize}
   \item This is the first item of an itemized list.
         Each item in the list is marked with a ``tick''.
         You don't have to worry about what kind of tick
         mark is used.

   \item This is the second item of the list.  It
         contains another list nested inside it.  The inner
         list is an \emph{enumerated} list.
         \begin{enumerate}
            \item This is the first item of an enumerated
                  list that is nested within the itemized list.

            \item This is the second item of the inner list.
                  \LaTeX\ allows you to nest lists deeper than
                  you really should.
         \end{enumerate}
         This is the rest of the second item of the outer
         list.  It is no more interesting than any other
         part of the item.
   \item This is the third item of the list.
\end{itemize}
You can even display poetry.
\begin{verse}
   There is an environment
    for verse \\             % The \\ command separates lines
   Whose features some poets % within a stanza.
   will curse.

                             % One or more blank lines separate stanzas.

   For instead of making\\
   Them do \emph{all} line breaking, \\
   It allows them to put too many words on a line when they'd rather be
   forced to be terse.
\end{verse}

Mathematical formulas may also be displayed.  A
displayed formula
is
one-line long; multiline
formulas require special formatting instructions.
   \[  \ip{\Gamma}{\psi'} = x'' + y^{2} + z_{i}^{n}\]
Don't start a paragraph with a displayed equation,
nor make one a paragraph by itself.

\end{document}               % End of document.

\section{Bench Huge Section 773}
% synthetic bench section 773
% This is a sample LaTeX input file.  (Version of 12 August 2004.)
%
% A '%' character causes TeX to ignore all remaining text on the line,
% and is used for comments like this one.

\documentclass{article}      % Specifies the document class

                             % The preamble begins here.
\title{An Example Document}  % Declares the document's title.
\author{Leslie Lamport}      % Declares the author's name.
\date{January 21, 1994}      % Deleting this command produces today's date.

\newcommand{\ip}[2]{(#1, #2)}
                             % Defines \ip{arg1}{arg2} to mean
                             % (arg1, arg2).

%\newcommand{\ip}[2]{\langle #1 | #2\rangle}
                             % This is an alternative definition of
                             % \ip that is commented out.

\begin{document}             % End of preamble and beginning of text.

\maketitle                   % Produces the title.

This is an example input file.  Comparing it with
the output it generates can show you how to
produce a simple document of your own.

\section{Ordinary Text}      % Produces section heading.  Lower-level
                             % sections are begun with similar
                             % \subsection and \subsubsection commands.

The ends  of words and sentences are marked
  by   spaces. It  doesn't matter how many
spaces    you type; one is as good as 100.  The
end of   a line counts as a space.

One   or more   blank lines denote the  end
of  a paragraph.

Since any number of consecutive spaces are treated
like a single one, the formatting of the input
file makes no difference to
      \LaTeX,                % The \LaTeX command generates the LaTeX logo.
but it makes a difference to you.  When you use
\LaTeX, making your input file as easy to read
as possible will be a great help as you write
your document and when you change it.  This sample
file shows how you can add comments to your own input
file.

Because printing is different from typewriting,
there are a number of things that you have to do
differently when preparing an input file than if
you were just typing the document directly.
Quotation marks like
       ``this''
have to be handled specially, as do quotes within
quotes:
       ``\,`this'            % \, separates the double and single quote.
        is what I just
        wrote, not  `that'\,''.

Dashes come in three sizes: an
       intra-word
dash, a medium dash for number ranges like
       1--2,
and a punctuation
       dash---like
this.

A sentence-ending space should be larger than the
space between words within a sentence.  You
sometimes have to type special commands in
conjunction with punctuation characters to get
this right, as in the following sentence.
       Gnats, gnus, etc.\ all  % `\ ' makes an inter-word space.
       begin with G\@.         % \@ marks end-of-sentence punctuation.
You should check the spaces after periods when
reading your output to make sure you haven't
forgotten any special cases.  Generating an
ellipsis
       \ldots\               % `\ ' is needed after `\ldots' because TeX
                             % ignores spaces after command names like \ldots
                             % made from \ + letters.
                             %
                             % Note how a `%' character causes TeX to ignore
                             % the end of the input line, so these blank lines
                             % do not start a new paragraph.
                             %
with the right spacing around the periods requires
a special command.

\LaTeX\ interprets some common characters as
commands, so you must type special commands to
generate them.  These characters include the
following:
       \$ \& \% \# \{ and \}.

In printing, text is usually emphasized with an
       \emph{italic}
type style.

\begin{em}
   A long segment of text can also be emphasized
   in this way.  Text within such a segment can be
   given \emph{additional} emphasis.
\end{em}

It is sometimes necessary to prevent \LaTeX\ from
breaking a line where it might otherwise do so.
This may be at a space, as between the ``Mr.''\ and
``Jones'' in
       ``Mr.~Jones'',        % ~ produces an unbreakable interword space.
or within a word---especially when the word is a
symbol like
       \mbox{\emph{itemnum}}
that makes little sense when hyphenated across
lines.

Footnotes\footnote{This is an example of a footnote.}
pose no problem.

\LaTeX\ is good at typesetting mathematical formulas
like
       \( x-3y + z = 7 \)
or
       \( a_{1} > x^{2n} + y^{2n} > x' \)
or
       \( \ip{A}{B} = \sum_{i} a_{i} b_{i} \).
The spaces you type in a formula are
ignored.  Remember that a letter like
       $x$                   % $ ... $  and  \( ... \)  are equivalent
is a formula when it denotes a mathematical
symbol, and it should be typed as one.

\section{Displayed Text}

Text is displayed by indenting it from the left
margin.  Quotations are commonly displayed.  There
are short quotations
\begin{quote}
   This is a short quotation.  It consists of a
   single paragraph of text.  See how it is formatted.
\end{quote}
and longer ones.
\begin{quotation}
   This is a longer quotation.  It consists of two
   paragraphs of text, neither of which are
   particularly interesting.

   This is the second paragraph of the quotation.  It
   is just as dull as the first paragraph.
\end{quotation}
Another frequently-displayed structure is a list.
The following is an example of an \emph{itemized}
list.
\begin{itemize}
   \item This is the first item of an itemized list.
         Each item in the list is marked with a ``tick''.
         You don't have to worry about what kind of tick
         mark is used.

   \item This is the second item of the list.  It
         contains another list nested inside it.  The inner
         list is an \emph{enumerated} list.
         \begin{enumerate}
            \item This is the first item of an enumerated
                  list that is nested within the itemized list.

            \item This is the second item of the inner list.
                  \LaTeX\ allows you to nest lists deeper than
                  you really should.
         \end{enumerate}
         This is the rest of the second item of the outer
         list.  It is no more interesting than any other
         part of the item.
   \item This is the third item of the list.
\end{itemize}
You can even display poetry.
\begin{verse}
   There is an environment
    for verse \\             % The \\ command separates lines
   Whose features some poets % within a stanza.
   will curse.

                             % One or more blank lines separate stanzas.

   For instead of making\\
   Them do \emph{all} line breaking, \\
   It allows them to put too many words on a line when they'd rather be
   forced to be terse.
\end{verse}

Mathematical formulas may also be displayed.  A
displayed formula
is
one-line long; multiline
formulas require special formatting instructions.
   \[  \ip{\Gamma}{\psi'} = x'' + y^{2} + z_{i}^{n}\]
Don't start a paragraph with a displayed equation,
nor make one a paragraph by itself.

\end{document}               % End of document.

\section{Bench Huge Section 774}
% synthetic bench section 774
% This is a sample LaTeX input file.  (Version of 12 August 2004.)
%
% A '%' character causes TeX to ignore all remaining text on the line,
% and is used for comments like this one.

\documentclass{article}      % Specifies the document class

                             % The preamble begins here.
\title{An Example Document}  % Declares the document's title.
\author{Leslie Lamport}      % Declares the author's name.
\date{January 21, 1994}      % Deleting this command produces today's date.

\newcommand{\ip}[2]{(#1, #2)}
                             % Defines \ip{arg1}{arg2} to mean
                             % (arg1, arg2).

%\newcommand{\ip}[2]{\langle #1 | #2\rangle}
                             % This is an alternative definition of
                             % \ip that is commented out.

\begin{document}             % End of preamble and beginning of text.

\maketitle                   % Produces the title.

This is an example input file.  Comparing it with
the output it generates can show you how to
produce a simple document of your own.

\section{Ordinary Text}      % Produces section heading.  Lower-level
                             % sections are begun with similar
                             % \subsection and \subsubsection commands.

The ends  of words and sentences are marked
  by   spaces. It  doesn't matter how many
spaces    you type; one is as good as 100.  The
end of   a line counts as a space.

One   or more   blank lines denote the  end
of  a paragraph.

Since any number of consecutive spaces are treated
like a single one, the formatting of the input
file makes no difference to
      \LaTeX,                % The \LaTeX command generates the LaTeX logo.
but it makes a difference to you.  When you use
\LaTeX, making your input file as easy to read
as possible will be a great help as you write
your document and when you change it.  This sample
file shows how you can add comments to your own input
file.

Because printing is different from typewriting,
there are a number of things that you have to do
differently when preparing an input file than if
you were just typing the document directly.
Quotation marks like
       ``this''
have to be handled specially, as do quotes within
quotes:
       ``\,`this'            % \, separates the double and single quote.
        is what I just
        wrote, not  `that'\,''.

Dashes come in three sizes: an
       intra-word
dash, a medium dash for number ranges like
       1--2,
and a punctuation
       dash---like
this.

A sentence-ending space should be larger than the
space between words within a sentence.  You
sometimes have to type special commands in
conjunction with punctuation characters to get
this right, as in the following sentence.
       Gnats, gnus, etc.\ all  % `\ ' makes an inter-word space.
       begin with G\@.         % \@ marks end-of-sentence punctuation.
You should check the spaces after periods when
reading your output to make sure you haven't
forgotten any special cases.  Generating an
ellipsis
       \ldots\               % `\ ' is needed after `\ldots' because TeX
                             % ignores spaces after command names like \ldots
                             % made from \ + letters.
                             %
                             % Note how a `%' character causes TeX to ignore
                             % the end of the input line, so these blank lines
                             % do not start a new paragraph.
                             %
with the right spacing around the periods requires
a special command.

\LaTeX\ interprets some common characters as
commands, so you must type special commands to
generate them.  These characters include the
following:
       \$ \& \% \# \{ and \}.

In printing, text is usually emphasized with an
       \emph{italic}
type style.

\begin{em}
   A long segment of text can also be emphasized
   in this way.  Text within such a segment can be
   given \emph{additional} emphasis.
\end{em}

It is sometimes necessary to prevent \LaTeX\ from
breaking a line where it might otherwise do so.
This may be at a space, as between the ``Mr.''\ and
``Jones'' in
       ``Mr.~Jones'',        % ~ produces an unbreakable interword space.
or within a word---especially when the word is a
symbol like
       \mbox{\emph{itemnum}}
that makes little sense when hyphenated across
lines.

Footnotes\footnote{This is an example of a footnote.}
pose no problem.

\LaTeX\ is good at typesetting mathematical formulas
like
       \( x-3y + z = 7 \)
or
       \( a_{1} > x^{2n} + y^{2n} > x' \)
or
       \( \ip{A}{B} = \sum_{i} a_{i} b_{i} \).
The spaces you type in a formula are
ignored.  Remember that a letter like
       $x$                   % $ ... $  and  \( ... \)  are equivalent
is a formula when it denotes a mathematical
symbol, and it should be typed as one.

\section{Displayed Text}

Text is displayed by indenting it from the left
margin.  Quotations are commonly displayed.  There
are short quotations
\begin{quote}
   This is a short quotation.  It consists of a
   single paragraph of text.  See how it is formatted.
\end{quote}
and longer ones.
\begin{quotation}
   This is a longer quotation.  It consists of two
   paragraphs of text, neither of which are
   particularly interesting.

   This is the second paragraph of the quotation.  It
   is just as dull as the first paragraph.
\end{quotation}
Another frequently-displayed structure is a list.
The following is an example of an \emph{itemized}
list.
\begin{itemize}
   \item This is the first item of an itemized list.
         Each item in the list is marked with a ``tick''.
         You don't have to worry about what kind of tick
         mark is used.

   \item This is the second item of the list.  It
         contains another list nested inside it.  The inner
         list is an \emph{enumerated} list.
         \begin{enumerate}
            \item This is the first item of an enumerated
                  list that is nested within the itemized list.

            \item This is the second item of the inner list.
                  \LaTeX\ allows you to nest lists deeper than
                  you really should.
         \end{enumerate}
         This is the rest of the second item of the outer
         list.  It is no more interesting than any other
         part of the item.
   \item This is the third item of the list.
\end{itemize}
You can even display poetry.
\begin{verse}
   There is an environment
    for verse \\             % The \\ command separates lines
   Whose features some poets % within a stanza.
   will curse.

                             % One or more blank lines separate stanzas.

   For instead of making\\
   Them do \emph{all} line breaking, \\
   It allows them to put too many words on a line when they'd rather be
   forced to be terse.
\end{verse}

Mathematical formulas may also be displayed.  A
displayed formula
is
one-line long; multiline
formulas require special formatting instructions.
   \[  \ip{\Gamma}{\psi'} = x'' + y^{2} + z_{i}^{n}\]
Don't start a paragraph with a displayed equation,
nor make one a paragraph by itself.

\end{document}               % End of document.

\section{Bench Huge Section 775}
% synthetic bench section 775
% This is a sample LaTeX input file.  (Version of 12 August 2004.)
%
% A '%' character causes TeX to ignore all remaining text on the line,
% and is used for comments like this one.

\documentclass{article}      % Specifies the document class

                             % The preamble begins here.
\title{An Example Document}  % Declares the document's title.
\author{Leslie Lamport}      % Declares the author's name.
\date{January 21, 1994}      % Deleting this command produces today's date.

\newcommand{\ip}[2]{(#1, #2)}
                             % Defines \ip{arg1}{arg2} to mean
                             % (arg1, arg2).

%\newcommand{\ip}[2]{\langle #1 | #2\rangle}
                             % This is an alternative definition of
                             % \ip that is commented out.

\begin{document}             % End of preamble and beginning of text.

\maketitle                   % Produces the title.

This is an example input file.  Comparing it with
the output it generates can show you how to
produce a simple document of your own.

\section{Ordinary Text}      % Produces section heading.  Lower-level
                             % sections are begun with similar
                             % \subsection and \subsubsection commands.

The ends  of words and sentences are marked
  by   spaces. It  doesn't matter how many
spaces    you type; one is as good as 100.  The
end of   a line counts as a space.

One   or more   blank lines denote the  end
of  a paragraph.

Since any number of consecutive spaces are treated
like a single one, the formatting of the input
file makes no difference to
      \LaTeX,                % The \LaTeX command generates the LaTeX logo.
but it makes a difference to you.  When you use
\LaTeX, making your input file as easy to read
as possible will be a great help as you write
your document and when you change it.  This sample
file shows how you can add comments to your own input
file.

Because printing is different from typewriting,
there are a number of things that you have to do
differently when preparing an input file than if
you were just typing the document directly.
Quotation marks like
       ``this''
have to be handled specially, as do quotes within
quotes:
       ``\,`this'            % \, separates the double and single quote.
        is what I just
        wrote, not  `that'\,''.

Dashes come in three sizes: an
       intra-word
dash, a medium dash for number ranges like
       1--2,
and a punctuation
       dash---like
this.

A sentence-ending space should be larger than the
space between words within a sentence.  You
sometimes have to type special commands in
conjunction with punctuation characters to get
this right, as in the following sentence.
       Gnats, gnus, etc.\ all  % `\ ' makes an inter-word space.
       begin with G\@.         % \@ marks end-of-sentence punctuation.
You should check the spaces after periods when
reading your output to make sure you haven't
forgotten any special cases.  Generating an
ellipsis
       \ldots\               % `\ ' is needed after `\ldots' because TeX
                             % ignores spaces after command names like \ldots
                             % made from \ + letters.
                             %
                             % Note how a `%' character causes TeX to ignore
                             % the end of the input line, so these blank lines
                             % do not start a new paragraph.
                             %
with the right spacing around the periods requires
a special command.

\LaTeX\ interprets some common characters as
commands, so you must type special commands to
generate them.  These characters include the
following:
       \$ \& \% \# \{ and \}.

In printing, text is usually emphasized with an
       \emph{italic}
type style.

\begin{em}
   A long segment of text can also be emphasized
   in this way.  Text within such a segment can be
   given \emph{additional} emphasis.
\end{em}

It is sometimes necessary to prevent \LaTeX\ from
breaking a line where it might otherwise do so.
This may be at a space, as between the ``Mr.''\ and
``Jones'' in
       ``Mr.~Jones'',        % ~ produces an unbreakable interword space.
or within a word---especially when the word is a
symbol like
       \mbox{\emph{itemnum}}
that makes little sense when hyphenated across
lines.

Footnotes\footnote{This is an example of a footnote.}
pose no problem.

\LaTeX\ is good at typesetting mathematical formulas
like
       \( x-3y + z = 7 \)
or
       \( a_{1} > x^{2n} + y^{2n} > x' \)
or
       \( \ip{A}{B} = \sum_{i} a_{i} b_{i} \).
The spaces you type in a formula are
ignored.  Remember that a letter like
       $x$                   % $ ... $  and  \( ... \)  are equivalent
is a formula when it denotes a mathematical
symbol, and it should be typed as one.

\section{Displayed Text}

Text is displayed by indenting it from the left
margin.  Quotations are commonly displayed.  There
are short quotations
\begin{quote}
   This is a short quotation.  It consists of a
   single paragraph of text.  See how it is formatted.
\end{quote}
and longer ones.
\begin{quotation}
   This is a longer quotation.  It consists of two
   paragraphs of text, neither of which are
   particularly interesting.

   This is the second paragraph of the quotation.  It
   is just as dull as the first paragraph.
\end{quotation}
Another frequently-displayed structure is a list.
The following is an example of an \emph{itemized}
list.
\begin{itemize}
   \item This is the first item of an itemized list.
         Each item in the list is marked with a ``tick''.
         You don't have to worry about what kind of tick
         mark is used.

   \item This is the second item of the list.  It
         contains another list nested inside it.  The inner
         list is an \emph{enumerated} list.
         \begin{enumerate}
            \item This is the first item of an enumerated
                  list that is nested within the itemized list.

            \item This is the second item of the inner list.
                  \LaTeX\ allows you to nest lists deeper than
                  you really should.
         \end{enumerate}
         This is the rest of the second item of the outer
         list.  It is no more interesting than any other
         part of the item.
   \item This is the third item of the list.
\end{itemize}
You can even display poetry.
\begin{verse}
   There is an environment
    for verse \\             % The \\ command separates lines
   Whose features some poets % within a stanza.
   will curse.

                             % One or more blank lines separate stanzas.

   For instead of making\\
   Them do \emph{all} line breaking, \\
   It allows them to put too many words on a line when they'd rather be
   forced to be terse.
\end{verse}

Mathematical formulas may also be displayed.  A
displayed formula
is
one-line long; multiline
formulas require special formatting instructions.
   \[  \ip{\Gamma}{\psi'} = x'' + y^{2} + z_{i}^{n}\]
Don't start a paragraph with a displayed equation,
nor make one a paragraph by itself.

\end{document}               % End of document.

\section{Bench Huge Section 776}
% synthetic bench section 776
% This is a sample LaTeX input file.  (Version of 12 August 2004.)
%
% A '%' character causes TeX to ignore all remaining text on the line,
% and is used for comments like this one.

\documentclass{article}      % Specifies the document class

                             % The preamble begins here.
\title{An Example Document}  % Declares the document's title.
\author{Leslie Lamport}      % Declares the author's name.
\date{January 21, 1994}      % Deleting this command produces today's date.

\newcommand{\ip}[2]{(#1, #2)}
                             % Defines \ip{arg1}{arg2} to mean
                             % (arg1, arg2).

%\newcommand{\ip}[2]{\langle #1 | #2\rangle}
                             % This is an alternative definition of
                             % \ip that is commented out.

\begin{document}             % End of preamble and beginning of text.

\maketitle                   % Produces the title.

This is an example input file.  Comparing it with
the output it generates can show you how to
produce a simple document of your own.

\section{Ordinary Text}      % Produces section heading.  Lower-level
                             % sections are begun with similar
                             % \subsection and \subsubsection commands.

The ends  of words and sentences are marked
  by   spaces. It  doesn't matter how many
spaces    you type; one is as good as 100.  The
end of   a line counts as a space.

One   or more   blank lines denote the  end
of  a paragraph.

Since any number of consecutive spaces are treated
like a single one, the formatting of the input
file makes no difference to
      \LaTeX,                % The \LaTeX command generates the LaTeX logo.
but it makes a difference to you.  When you use
\LaTeX, making your input file as easy to read
as possible will be a great help as you write
your document and when you change it.  This sample
file shows how you can add comments to your own input
file.

Because printing is different from typewriting,
there are a number of things that you have to do
differently when preparing an input file than if
you were just typing the document directly.
Quotation marks like
       ``this''
have to be handled specially, as do quotes within
quotes:
       ``\,`this'            % \, separates the double and single quote.
        is what I just
        wrote, not  `that'\,''.

Dashes come in three sizes: an
       intra-word
dash, a medium dash for number ranges like
       1--2,
and a punctuation
       dash---like
this.

A sentence-ending space should be larger than the
space between words within a sentence.  You
sometimes have to type special commands in
conjunction with punctuation characters to get
this right, as in the following sentence.
       Gnats, gnus, etc.\ all  % `\ ' makes an inter-word space.
       begin with G\@.         % \@ marks end-of-sentence punctuation.
You should check the spaces after periods when
reading your output to make sure you haven't
forgotten any special cases.  Generating an
ellipsis
       \ldots\               % `\ ' is needed after `\ldots' because TeX
                             % ignores spaces after command names like \ldots
                             % made from \ + letters.
                             %
                             % Note how a `%' character causes TeX to ignore
                             % the end of the input line, so these blank lines
                             % do not start a new paragraph.
                             %
with the right spacing around the periods requires
a special command.

\LaTeX\ interprets some common characters as
commands, so you must type special commands to
generate them.  These characters include the
following:
       \$ \& \% \# \{ and \}.

In printing, text is usually emphasized with an
       \emph{italic}
type style.

\begin{em}
   A long segment of text can also be emphasized
   in this way.  Text within such a segment can be
   given \emph{additional} emphasis.
\end{em}

It is sometimes necessary to prevent \LaTeX\ from
breaking a line where it might otherwise do so.
This may be at a space, as between the ``Mr.''\ and
``Jones'' in
       ``Mr.~Jones'',        % ~ produces an unbreakable interword space.
or within a word---especially when the word is a
symbol like
       \mbox{\emph{itemnum}}
that makes little sense when hyphenated across
lines.

Footnotes\footnote{This is an example of a footnote.}
pose no problem.

\LaTeX\ is good at typesetting mathematical formulas
like
       \( x-3y + z = 7 \)
or
       \( a_{1} > x^{2n} + y^{2n} > x' \)
or
       \( \ip{A}{B} = \sum_{i} a_{i} b_{i} \).
The spaces you type in a formula are
ignored.  Remember that a letter like
       $x$                   % $ ... $  and  \( ... \)  are equivalent
is a formula when it denotes a mathematical
symbol, and it should be typed as one.

\section{Displayed Text}

Text is displayed by indenting it from the left
margin.  Quotations are commonly displayed.  There
are short quotations
\begin{quote}
   This is a short quotation.  It consists of a
   single paragraph of text.  See how it is formatted.
\end{quote}
and longer ones.
\begin{quotation}
   This is a longer quotation.  It consists of two
   paragraphs of text, neither of which are
   particularly interesting.

   This is the second paragraph of the quotation.  It
   is just as dull as the first paragraph.
\end{quotation}
Another frequently-displayed structure is a list.
The following is an example of an \emph{itemized}
list.
\begin{itemize}
   \item This is the first item of an itemized list.
         Each item in the list is marked with a ``tick''.
         You don't have to worry about what kind of tick
         mark is used.

   \item This is the second item of the list.  It
         contains another list nested inside it.  The inner
         list is an \emph{enumerated} list.
         \begin{enumerate}
            \item This is the first item of an enumerated
                  list that is nested within the itemized list.

            \item This is the second item of the inner list.
                  \LaTeX\ allows you to nest lists deeper than
                  you really should.
         \end{enumerate}
         This is the rest of the second item of the outer
         list.  It is no more interesting than any other
         part of the item.
   \item This is the third item of the list.
\end{itemize}
You can even display poetry.
\begin{verse}
   There is an environment
    for verse \\             % The \\ command separates lines
   Whose features some poets % within a stanza.
   will curse.

                             % One or more blank lines separate stanzas.

   For instead of making\\
   Them do \emph{all} line breaking, \\
   It allows them to put too many words on a line when they'd rather be
   forced to be terse.
\end{verse}

Mathematical formulas may also be displayed.  A
displayed formula
is
one-line long; multiline
formulas require special formatting instructions.
   \[  \ip{\Gamma}{\psi'} = x'' + y^{2} + z_{i}^{n}\]
Don't start a paragraph with a displayed equation,
nor make one a paragraph by itself.

\end{document}               % End of document.

\section{Bench Huge Section 777}
% synthetic bench section 777
% This is a sample LaTeX input file.  (Version of 12 August 2004.)
%
% A '%' character causes TeX to ignore all remaining text on the line,
% and is used for comments like this one.

\documentclass{article}      % Specifies the document class

                             % The preamble begins here.
\title{An Example Document}  % Declares the document's title.
\author{Leslie Lamport}      % Declares the author's name.
\date{January 21, 1994}      % Deleting this command produces today's date.

\newcommand{\ip}[2]{(#1, #2)}
                             % Defines \ip{arg1}{arg2} to mean
                             % (arg1, arg2).

%\newcommand{\ip}[2]{\langle #1 | #2\rangle}
                             % This is an alternative definition of
                             % \ip that is commented out.

\begin{document}             % End of preamble and beginning of text.

\maketitle                   % Produces the title.

This is an example input file.  Comparing it with
the output it generates can show you how to
produce a simple document of your own.

\section{Ordinary Text}      % Produces section heading.  Lower-level
                             % sections are begun with similar
                             % \subsection and \subsubsection commands.

The ends  of words and sentences are marked
  by   spaces. It  doesn't matter how many
spaces    you type; one is as good as 100.  The
end of   a line counts as a space.

One   or more   blank lines denote the  end
of  a paragraph.

Since any number of consecutive spaces are treated
like a single one, the formatting of the input
file makes no difference to
      \LaTeX,                % The \LaTeX command generates the LaTeX logo.
but it makes a difference to you.  When you use
\LaTeX, making your input file as easy to read
as possible will be a great help as you write
your document and when you change it.  This sample
file shows how you can add comments to your own input
file.

Because printing is different from typewriting,
there are a number of things that you have to do
differently when preparing an input file than if
you were just typing the document directly.
Quotation marks like
       ``this''
have to be handled specially, as do quotes within
quotes:
       ``\,`this'            % \, separates the double and single quote.
        is what I just
        wrote, not  `that'\,''.

Dashes come in three sizes: an
       intra-word
dash, a medium dash for number ranges like
       1--2,
and a punctuation
       dash---like
this.

A sentence-ending space should be larger than the
space between words within a sentence.  You
sometimes have to type special commands in
conjunction with punctuation characters to get
this right, as in the following sentence.
       Gnats, gnus, etc.\ all  % `\ ' makes an inter-word space.
       begin with G\@.         % \@ marks end-of-sentence punctuation.
You should check the spaces after periods when
reading your output to make sure you haven't
forgotten any special cases.  Generating an
ellipsis
       \ldots\               % `\ ' is needed after `\ldots' because TeX
                             % ignores spaces after command names like \ldots
                             % made from \ + letters.
                             %
                             % Note how a `%' character causes TeX to ignore
                             % the end of the input line, so these blank lines
                             % do not start a new paragraph.
                             %
with the right spacing around the periods requires
a special command.

\LaTeX\ interprets some common characters as
commands, so you must type special commands to
generate them.  These characters include the
following:
       \$ \& \% \# \{ and \}.

In printing, text is usually emphasized with an
       \emph{italic}
type style.

\begin{em}
   A long segment of text can also be emphasized
   in this way.  Text within such a segment can be
   given \emph{additional} emphasis.
\end{em}

It is sometimes necessary to prevent \LaTeX\ from
breaking a line where it might otherwise do so.
This may be at a space, as between the ``Mr.''\ and
``Jones'' in
       ``Mr.~Jones'',        % ~ produces an unbreakable interword space.
or within a word---especially when the word is a
symbol like
       \mbox{\emph{itemnum}}
that makes little sense when hyphenated across
lines.

Footnotes\footnote{This is an example of a footnote.}
pose no problem.

\LaTeX\ is good at typesetting mathematical formulas
like
       \( x-3y + z = 7 \)
or
       \( a_{1} > x^{2n} + y^{2n} > x' \)
or
       \( \ip{A}{B} = \sum_{i} a_{i} b_{i} \).
The spaces you type in a formula are
ignored.  Remember that a letter like
       $x$                   % $ ... $  and  \( ... \)  are equivalent
is a formula when it denotes a mathematical
symbol, and it should be typed as one.

\section{Displayed Text}

Text is displayed by indenting it from the left
margin.  Quotations are commonly displayed.  There
are short quotations
\begin{quote}
   This is a short quotation.  It consists of a
   single paragraph of text.  See how it is formatted.
\end{quote}
and longer ones.
\begin{quotation}
   This is a longer quotation.  It consists of two
   paragraphs of text, neither of which are
   particularly interesting.

   This is the second paragraph of the quotation.  It
   is just as dull as the first paragraph.
\end{quotation}
Another frequently-displayed structure is a list.
The following is an example of an \emph{itemized}
list.
\begin{itemize}
   \item This is the first item of an itemized list.
         Each item in the list is marked with a ``tick''.
         You don't have to worry about what kind of tick
         mark is used.

   \item This is the second item of the list.  It
         contains another list nested inside it.  The inner
         list is an \emph{enumerated} list.
         \begin{enumerate}
            \item This is the first item of an enumerated
                  list that is nested within the itemized list.

            \item This is the second item of the inner list.
                  \LaTeX\ allows you to nest lists deeper than
                  you really should.
         \end{enumerate}
         This is the rest of the second item of the outer
         list.  It is no more interesting than any other
         part of the item.
   \item This is the third item of the list.
\end{itemize}
You can even display poetry.
\begin{verse}
   There is an environment
    for verse \\             % The \\ command separates lines
   Whose features some poets % within a stanza.
   will curse.

                             % One or more blank lines separate stanzas.

   For instead of making\\
   Them do \emph{all} line breaking, \\
   It allows them to put too many words on a line when they'd rather be
   forced to be terse.
\end{verse}

Mathematical formulas may also be displayed.  A
displayed formula
is
one-line long; multiline
formulas require special formatting instructions.
   \[  \ip{\Gamma}{\psi'} = x'' + y^{2} + z_{i}^{n}\]
Don't start a paragraph with a displayed equation,
nor make one a paragraph by itself.

\end{document}               % End of document.

\section{Bench Huge Section 778}
% synthetic bench section 778
% This is a sample LaTeX input file.  (Version of 12 August 2004.)
%
% A '%' character causes TeX to ignore all remaining text on the line,
% and is used for comments like this one.

\documentclass{article}      % Specifies the document class

                             % The preamble begins here.
\title{An Example Document}  % Declares the document's title.
\author{Leslie Lamport}      % Declares the author's name.
\date{January 21, 1994}      % Deleting this command produces today's date.

\newcommand{\ip}[2]{(#1, #2)}
                             % Defines \ip{arg1}{arg2} to mean
                             % (arg1, arg2).

%\newcommand{\ip}[2]{\langle #1 | #2\rangle}
                             % This is an alternative definition of
                             % \ip that is commented out.

\begin{document}             % End of preamble and beginning of text.

\maketitle                   % Produces the title.

This is an example input file.  Comparing it with
the output it generates can show you how to
produce a simple document of your own.

\section{Ordinary Text}      % Produces section heading.  Lower-level
                             % sections are begun with similar
                             % \subsection and \subsubsection commands.

The ends  of words and sentences are marked
  by   spaces. It  doesn't matter how many
spaces    you type; one is as good as 100.  The
end of   a line counts as a space.

One   or more   blank lines denote the  end
of  a paragraph.

Since any number of consecutive spaces are treated
like a single one, the formatting of the input
file makes no difference to
      \LaTeX,                % The \LaTeX command generates the LaTeX logo.
but it makes a difference to you.  When you use
\LaTeX, making your input file as easy to read
as possible will be a great help as you write
your document and when you change it.  This sample
file shows how you can add comments to your own input
file.

Because printing is different from typewriting,
there are a number of things that you have to do
differently when preparing an input file than if
you were just typing the document directly.
Quotation marks like
       ``this''
have to be handled specially, as do quotes within
quotes:
       ``\,`this'            % \, separates the double and single quote.
        is what I just
        wrote, not  `that'\,''.

Dashes come in three sizes: an
       intra-word
dash, a medium dash for number ranges like
       1--2,
and a punctuation
       dash---like
this.

A sentence-ending space should be larger than the
space between words within a sentence.  You
sometimes have to type special commands in
conjunction with punctuation characters to get
this right, as in the following sentence.
       Gnats, gnus, etc.\ all  % `\ ' makes an inter-word space.
       begin with G\@.         % \@ marks end-of-sentence punctuation.
You should check the spaces after periods when
reading your output to make sure you haven't
forgotten any special cases.  Generating an
ellipsis
       \ldots\               % `\ ' is needed after `\ldots' because TeX
                             % ignores spaces after command names like \ldots
                             % made from \ + letters.
                             %
                             % Note how a `%' character causes TeX to ignore
                             % the end of the input line, so these blank lines
                             % do not start a new paragraph.
                             %
with the right spacing around the periods requires
a special command.

\LaTeX\ interprets some common characters as
commands, so you must type special commands to
generate them.  These characters include the
following:
       \$ \& \% \# \{ and \}.

In printing, text is usually emphasized with an
       \emph{italic}
type style.

\begin{em}
   A long segment of text can also be emphasized
   in this way.  Text within such a segment can be
   given \emph{additional} emphasis.
\end{em}

It is sometimes necessary to prevent \LaTeX\ from
breaking a line where it might otherwise do so.
This may be at a space, as between the ``Mr.''\ and
``Jones'' in
       ``Mr.~Jones'',        % ~ produces an unbreakable interword space.
or within a word---especially when the word is a
symbol like
       \mbox{\emph{itemnum}}
that makes little sense when hyphenated across
lines.

Footnotes\footnote{This is an example of a footnote.}
pose no problem.

\LaTeX\ is good at typesetting mathematical formulas
like
       \( x-3y + z = 7 \)
or
       \( a_{1} > x^{2n} + y^{2n} > x' \)
or
       \( \ip{A}{B} = \sum_{i} a_{i} b_{i} \).
The spaces you type in a formula are
ignored.  Remember that a letter like
       $x$                   % $ ... $  and  \( ... \)  are equivalent
is a formula when it denotes a mathematical
symbol, and it should be typed as one.

\section{Displayed Text}

Text is displayed by indenting it from the left
margin.  Quotations are commonly displayed.  There
are short quotations
\begin{quote}
   This is a short quotation.  It consists of a
   single paragraph of text.  See how it is formatted.
\end{quote}
and longer ones.
\begin{quotation}
   This is a longer quotation.  It consists of two
   paragraphs of text, neither of which are
   particularly interesting.

   This is the second paragraph of the quotation.  It
   is just as dull as the first paragraph.
\end{quotation}
Another frequently-displayed structure is a list.
The following is an example of an \emph{itemized}
list.
\begin{itemize}
   \item This is the first item of an itemized list.
         Each item in the list is marked with a ``tick''.
         You don't have to worry about what kind of tick
         mark is used.

   \item This is the second item of the list.  It
         contains another list nested inside it.  The inner
         list is an \emph{enumerated} list.
         \begin{enumerate}
            \item This is the first item of an enumerated
                  list that is nested within the itemized list.

            \item This is the second item of the inner list.
                  \LaTeX\ allows you to nest lists deeper than
                  you really should.
         \end{enumerate}
         This is the rest of the second item of the outer
         list.  It is no more interesting than any other
         part of the item.
   \item This is the third item of the list.
\end{itemize}
You can even display poetry.
\begin{verse}
   There is an environment
    for verse \\             % The \\ command separates lines
   Whose features some poets % within a stanza.
   will curse.

                             % One or more blank lines separate stanzas.

   For instead of making\\
   Them do \emph{all} line breaking, \\
   It allows them to put too many words on a line when they'd rather be
   forced to be terse.
\end{verse}

Mathematical formulas may also be displayed.  A
displayed formula
is
one-line long; multiline
formulas require special formatting instructions.
   \[  \ip{\Gamma}{\psi'} = x'' + y^{2} + z_{i}^{n}\]
Don't start a paragraph with a displayed equation,
nor make one a paragraph by itself.

\end{document}               % End of document.

\section{Bench Huge Section 779}
% synthetic bench section 779
% This is a sample LaTeX input file.  (Version of 12 August 2004.)
%
% A '%' character causes TeX to ignore all remaining text on the line,
% and is used for comments like this one.

\documentclass{article}      % Specifies the document class

                             % The preamble begins here.
\title{An Example Document}  % Declares the document's title.
\author{Leslie Lamport}      % Declares the author's name.
\date{January 21, 1994}      % Deleting this command produces today's date.

\newcommand{\ip}[2]{(#1, #2)}
                             % Defines \ip{arg1}{arg2} to mean
                             % (arg1, arg2).

%\newcommand{\ip}[2]{\langle #1 | #2\rangle}
                             % This is an alternative definition of
                             % \ip that is commented out.

\begin{document}             % End of preamble and beginning of text.

\maketitle                   % Produces the title.

This is an example input file.  Comparing it with
the output it generates can show you how to
produce a simple document of your own.

\section{Ordinary Text}      % Produces section heading.  Lower-level
                             % sections are begun with similar
                             % \subsection and \subsubsection commands.

The ends  of words and sentences are marked
  by   spaces. It  doesn't matter how many
spaces    you type; one is as good as 100.  The
end of   a line counts as a space.

One   or more   blank lines denote the  end
of  a paragraph.

Since any number of consecutive spaces are treated
like a single one, the formatting of the input
file makes no difference to
      \LaTeX,                % The \LaTeX command generates the LaTeX logo.
but it makes a difference to you.  When you use
\LaTeX, making your input file as easy to read
as possible will be a great help as you write
your document and when you change it.  This sample
file shows how you can add comments to your own input
file.

Because printing is different from typewriting,
there are a number of things that you have to do
differently when preparing an input file than if
you were just typing the document directly.
Quotation marks like
       ``this''
have to be handled specially, as do quotes within
quotes:
       ``\,`this'            % \, separates the double and single quote.
        is what I just
        wrote, not  `that'\,''.

Dashes come in three sizes: an
       intra-word
dash, a medium dash for number ranges like
       1--2,
and a punctuation
       dash---like
this.

A sentence-ending space should be larger than the
space between words within a sentence.  You
sometimes have to type special commands in
conjunction with punctuation characters to get
this right, as in the following sentence.
       Gnats, gnus, etc.\ all  % `\ ' makes an inter-word space.
       begin with G\@.         % \@ marks end-of-sentence punctuation.
You should check the spaces after periods when
reading your output to make sure you haven't
forgotten any special cases.  Generating an
ellipsis
       \ldots\               % `\ ' is needed after `\ldots' because TeX
                             % ignores spaces after command names like \ldots
                             % made from \ + letters.
                             %
                             % Note how a `%' character causes TeX to ignore
                             % the end of the input line, so these blank lines
                             % do not start a new paragraph.
                             %
with the right spacing around the periods requires
a special command.

\LaTeX\ interprets some common characters as
commands, so you must type special commands to
generate them.  These characters include the
following:
       \$ \& \% \# \{ and \}.

In printing, text is usually emphasized with an
       \emph{italic}
type style.

\begin{em}
   A long segment of text can also be emphasized
   in this way.  Text within such a segment can be
   given \emph{additional} emphasis.
\end{em}

It is sometimes necessary to prevent \LaTeX\ from
breaking a line where it might otherwise do so.
This may be at a space, as between the ``Mr.''\ and
``Jones'' in
       ``Mr.~Jones'',        % ~ produces an unbreakable interword space.
or within a word---especially when the word is a
symbol like
       \mbox{\emph{itemnum}}
that makes little sense when hyphenated across
lines.

Footnotes\footnote{This is an example of a footnote.}
pose no problem.

\LaTeX\ is good at typesetting mathematical formulas
like
       \( x-3y + z = 7 \)
or
       \( a_{1} > x^{2n} + y^{2n} > x' \)
or
       \( \ip{A}{B} = \sum_{i} a_{i} b_{i} \).
The spaces you type in a formula are
ignored.  Remember that a letter like
       $x$                   % $ ... $  and  \( ... \)  are equivalent
is a formula when it denotes a mathematical
symbol, and it should be typed as one.

\section{Displayed Text}

Text is displayed by indenting it from the left
margin.  Quotations are commonly displayed.  There
are short quotations
\begin{quote}
   This is a short quotation.  It consists of a
   single paragraph of text.  See how it is formatted.
\end{quote}
and longer ones.
\begin{quotation}
   This is a longer quotation.  It consists of two
   paragraphs of text, neither of which are
   particularly interesting.

   This is the second paragraph of the quotation.  It
   is just as dull as the first paragraph.
\end{quotation}
Another frequently-displayed structure is a list.
The following is an example of an \emph{itemized}
list.
\begin{itemize}
   \item This is the first item of an itemized list.
         Each item in the list is marked with a ``tick''.
         You don't have to worry about what kind of tick
         mark is used.

   \item This is the second item of the list.  It
         contains another list nested inside it.  The inner
         list is an \emph{enumerated} list.
         \begin{enumerate}
            \item This is the first item of an enumerated
                  list that is nested within the itemized list.

            \item This is the second item of the inner list.
                  \LaTeX\ allows you to nest lists deeper than
                  you really should.
         \end{enumerate}
         This is the rest of the second item of the outer
         list.  It is no more interesting than any other
         part of the item.
   \item This is the third item of the list.
\end{itemize}
You can even display poetry.
\begin{verse}
   There is an environment
    for verse \\             % The \\ command separates lines
   Whose features some poets % within a stanza.
   will curse.

                             % One or more blank lines separate stanzas.

   For instead of making\\
   Them do \emph{all} line breaking, \\
   It allows them to put too many words on a line when they'd rather be
   forced to be terse.
\end{verse}

Mathematical formulas may also be displayed.  A
displayed formula
is
one-line long; multiline
formulas require special formatting instructions.
   \[  \ip{\Gamma}{\psi'} = x'' + y^{2} + z_{i}^{n}\]
Don't start a paragraph with a displayed equation,
nor make one a paragraph by itself.

\end{document}               % End of document.

\section{Bench Huge Section 780}
% synthetic bench section 780
% This is a sample LaTeX input file.  (Version of 12 August 2004.)
%
% A '%' character causes TeX to ignore all remaining text on the line,
% and is used for comments like this one.

\documentclass{article}      % Specifies the document class

                             % The preamble begins here.
\title{An Example Document}  % Declares the document's title.
\author{Leslie Lamport}      % Declares the author's name.
\date{January 21, 1994}      % Deleting this command produces today's date.

\newcommand{\ip}[2]{(#1, #2)}
                             % Defines \ip{arg1}{arg2} to mean
                             % (arg1, arg2).

%\newcommand{\ip}[2]{\langle #1 | #2\rangle}
                             % This is an alternative definition of
                             % \ip that is commented out.

\begin{document}             % End of preamble and beginning of text.

\maketitle                   % Produces the title.

This is an example input file.  Comparing it with
the output it generates can show you how to
produce a simple document of your own.

\section{Ordinary Text}      % Produces section heading.  Lower-level
                             % sections are begun with similar
                             % \subsection and \subsubsection commands.

The ends  of words and sentences are marked
  by   spaces. It  doesn't matter how many
spaces    you type; one is as good as 100.  The
end of   a line counts as a space.

One   or more   blank lines denote the  end
of  a paragraph.

Since any number of consecutive spaces are treated
like a single one, the formatting of the input
file makes no difference to
      \LaTeX,                % The \LaTeX command generates the LaTeX logo.
but it makes a difference to you.  When you use
\LaTeX, making your input file as easy to read
as possible will be a great help as you write
your document and when you change it.  This sample
file shows how you can add comments to your own input
file.

Because printing is different from typewriting,
there are a number of things that you have to do
differently when preparing an input file than if
you were just typing the document directly.
Quotation marks like
       ``this''
have to be handled specially, as do quotes within
quotes:
       ``\,`this'            % \, separates the double and single quote.
        is what I just
        wrote, not  `that'\,''.

Dashes come in three sizes: an
       intra-word
dash, a medium dash for number ranges like
       1--2,
and a punctuation
       dash---like
this.

A sentence-ending space should be larger than the
space between words within a sentence.  You
sometimes have to type special commands in
conjunction with punctuation characters to get
this right, as in the following sentence.
       Gnats, gnus, etc.\ all  % `\ ' makes an inter-word space.
       begin with G\@.         % \@ marks end-of-sentence punctuation.
You should check the spaces after periods when
reading your output to make sure you haven't
forgotten any special cases.  Generating an
ellipsis
       \ldots\               % `\ ' is needed after `\ldots' because TeX
                             % ignores spaces after command names like \ldots
                             % made from \ + letters.
                             %
                             % Note how a `%' character causes TeX to ignore
                             % the end of the input line, so these blank lines
                             % do not start a new paragraph.
                             %
with the right spacing around the periods requires
a special command.

\LaTeX\ interprets some common characters as
commands, so you must type special commands to
generate them.  These characters include the
following:
       \$ \& \% \# \{ and \}.

In printing, text is usually emphasized with an
       \emph{italic}
type style.

\begin{em}
   A long segment of text can also be emphasized
   in this way.  Text within such a segment can be
   given \emph{additional} emphasis.
\end{em}

It is sometimes necessary to prevent \LaTeX\ from
breaking a line where it might otherwise do so.
This may be at a space, as between the ``Mr.''\ and
``Jones'' in
       ``Mr.~Jones'',        % ~ produces an unbreakable interword space.
or within a word---especially when the word is a
symbol like
       \mbox{\emph{itemnum}}
that makes little sense when hyphenated across
lines.

Footnotes\footnote{This is an example of a footnote.}
pose no problem.

\LaTeX\ is good at typesetting mathematical formulas
like
       \( x-3y + z = 7 \)
or
       \( a_{1} > x^{2n} + y^{2n} > x' \)
or
       \( \ip{A}{B} = \sum_{i} a_{i} b_{i} \).
The spaces you type in a formula are
ignored.  Remember that a letter like
       $x$                   % $ ... $  and  \( ... \)  are equivalent
is a formula when it denotes a mathematical
symbol, and it should be typed as one.

\section{Displayed Text}

Text is displayed by indenting it from the left
margin.  Quotations are commonly displayed.  There
are short quotations
\begin{quote}
   This is a short quotation.  It consists of a
   single paragraph of text.  See how it is formatted.
\end{quote}
and longer ones.
\begin{quotation}
   This is a longer quotation.  It consists of two
   paragraphs of text, neither of which are
   particularly interesting.

   This is the second paragraph of the quotation.  It
   is just as dull as the first paragraph.
\end{quotation}
Another frequently-displayed structure is a list.
The following is an example of an \emph{itemized}
list.
\begin{itemize}
   \item This is the first item of an itemized list.
         Each item in the list is marked with a ``tick''.
         You don't have to worry about what kind of tick
         mark is used.

   \item This is the second item of the list.  It
         contains another list nested inside it.  The inner
         list is an \emph{enumerated} list.
         \begin{enumerate}
            \item This is the first item of an enumerated
                  list that is nested within the itemized list.

            \item This is the second item of the inner list.
                  \LaTeX\ allows you to nest lists deeper than
                  you really should.
         \end{enumerate}
         This is the rest of the second item of the outer
         list.  It is no more interesting than any other
         part of the item.
   \item This is the third item of the list.
\end{itemize}
You can even display poetry.
\begin{verse}
   There is an environment
    for verse \\             % The \\ command separates lines
   Whose features some poets % within a stanza.
   will curse.

                             % One or more blank lines separate stanzas.

   For instead of making\\
   Them do \emph{all} line breaking, \\
   It allows them to put too many words on a line when they'd rather be
   forced to be terse.
\end{verse}

Mathematical formulas may also be displayed.  A
displayed formula
is
one-line long; multiline
formulas require special formatting instructions.
   \[  \ip{\Gamma}{\psi'} = x'' + y^{2} + z_{i}^{n}\]
Don't start a paragraph with a displayed equation,
nor make one a paragraph by itself.

\end{document}               % End of document.

\section{Bench Huge Section 781}
% synthetic bench section 781
% This is a sample LaTeX input file.  (Version of 12 August 2004.)
%
% A '%' character causes TeX to ignore all remaining text on the line,
% and is used for comments like this one.

\documentclass{article}      % Specifies the document class

                             % The preamble begins here.
\title{An Example Document}  % Declares the document's title.
\author{Leslie Lamport}      % Declares the author's name.
\date{January 21, 1994}      % Deleting this command produces today's date.

\newcommand{\ip}[2]{(#1, #2)}
                             % Defines \ip{arg1}{arg2} to mean
                             % (arg1, arg2).

%\newcommand{\ip}[2]{\langle #1 | #2\rangle}
                             % This is an alternative definition of
                             % \ip that is commented out.

\begin{document}             % End of preamble and beginning of text.

\maketitle                   % Produces the title.

This is an example input file.  Comparing it with
the output it generates can show you how to
produce a simple document of your own.

\section{Ordinary Text}      % Produces section heading.  Lower-level
                             % sections are begun with similar
                             % \subsection and \subsubsection commands.

The ends  of words and sentences are marked
  by   spaces. It  doesn't matter how many
spaces    you type; one is as good as 100.  The
end of   a line counts as a space.

One   or more   blank lines denote the  end
of  a paragraph.

Since any number of consecutive spaces are treated
like a single one, the formatting of the input
file makes no difference to
      \LaTeX,                % The \LaTeX command generates the LaTeX logo.
but it makes a difference to you.  When you use
\LaTeX, making your input file as easy to read
as possible will be a great help as you write
your document and when you change it.  This sample
file shows how you can add comments to your own input
file.

Because printing is different from typewriting,
there are a number of things that you have to do
differently when preparing an input file than if
you were just typing the document directly.
Quotation marks like
       ``this''
have to be handled specially, as do quotes within
quotes:
       ``\,`this'            % \, separates the double and single quote.
        is what I just
        wrote, not  `that'\,''.

Dashes come in three sizes: an
       intra-word
dash, a medium dash for number ranges like
       1--2,
and a punctuation
       dash---like
this.

A sentence-ending space should be larger than the
space between words within a sentence.  You
sometimes have to type special commands in
conjunction with punctuation characters to get
this right, as in the following sentence.
       Gnats, gnus, etc.\ all  % `\ ' makes an inter-word space.
       begin with G\@.         % \@ marks end-of-sentence punctuation.
You should check the spaces after periods when
reading your output to make sure you haven't
forgotten any special cases.  Generating an
ellipsis
       \ldots\               % `\ ' is needed after `\ldots' because TeX
                             % ignores spaces after command names like \ldots
                             % made from \ + letters.
                             %
                             % Note how a `%' character causes TeX to ignore
                             % the end of the input line, so these blank lines
                             % do not start a new paragraph.
                             %
with the right spacing around the periods requires
a special command.

\LaTeX\ interprets some common characters as
commands, so you must type special commands to
generate them.  These characters include the
following:
       \$ \& \% \# \{ and \}.

In printing, text is usually emphasized with an
       \emph{italic}
type style.

\begin{em}
   A long segment of text can also be emphasized
   in this way.  Text within such a segment can be
   given \emph{additional} emphasis.
\end{em}

It is sometimes necessary to prevent \LaTeX\ from
breaking a line where it might otherwise do so.
This may be at a space, as between the ``Mr.''\ and
``Jones'' in
       ``Mr.~Jones'',        % ~ produces an unbreakable interword space.
or within a word---especially when the word is a
symbol like
       \mbox{\emph{itemnum}}
that makes little sense when hyphenated across
lines.

Footnotes\footnote{This is an example of a footnote.}
pose no problem.

\LaTeX\ is good at typesetting mathematical formulas
like
       \( x-3y + z = 7 \)
or
       \( a_{1} > x^{2n} + y^{2n} > x' \)
or
       \( \ip{A}{B} = \sum_{i} a_{i} b_{i} \).
The spaces you type in a formula are
ignored.  Remember that a letter like
       $x$                   % $ ... $  and  \( ... \)  are equivalent
is a formula when it denotes a mathematical
symbol, and it should be typed as one.

\section{Displayed Text}

Text is displayed by indenting it from the left
margin.  Quotations are commonly displayed.  There
are short quotations
\begin{quote}
   This is a short quotation.  It consists of a
   single paragraph of text.  See how it is formatted.
\end{quote}
and longer ones.
\begin{quotation}
   This is a longer quotation.  It consists of two
   paragraphs of text, neither of which are
   particularly interesting.

   This is the second paragraph of the quotation.  It
   is just as dull as the first paragraph.
\end{quotation}
Another frequently-displayed structure is a list.
The following is an example of an \emph{itemized}
list.
\begin{itemize}
   \item This is the first item of an itemized list.
         Each item in the list is marked with a ``tick''.
         You don't have to worry about what kind of tick
         mark is used.

   \item This is the second item of the list.  It
         contains another list nested inside it.  The inner
         list is an \emph{enumerated} list.
         \begin{enumerate}
            \item This is the first item of an enumerated
                  list that is nested within the itemized list.

            \item This is the second item of the inner list.
                  \LaTeX\ allows you to nest lists deeper than
                  you really should.
         \end{enumerate}
         This is the rest of the second item of the outer
         list.  It is no more interesting than any other
         part of the item.
   \item This is the third item of the list.
\end{itemize}
You can even display poetry.
\begin{verse}
   There is an environment
    for verse \\             % The \\ command separates lines
   Whose features some poets % within a stanza.
   will curse.

                             % One or more blank lines separate stanzas.

   For instead of making\\
   Them do \emph{all} line breaking, \\
   It allows them to put too many words on a line when they'd rather be
   forced to be terse.
\end{verse}

Mathematical formulas may also be displayed.  A
displayed formula
is
one-line long; multiline
formulas require special formatting instructions.
   \[  \ip{\Gamma}{\psi'} = x'' + y^{2} + z_{i}^{n}\]
Don't start a paragraph with a displayed equation,
nor make one a paragraph by itself.

\end{document}               % End of document.

\section{Bench Huge Section 782}
% synthetic bench section 782
% This is a sample LaTeX input file.  (Version of 12 August 2004.)
%
% A '%' character causes TeX to ignore all remaining text on the line,
% and is used for comments like this one.

\documentclass{article}      % Specifies the document class

                             % The preamble begins here.
\title{An Example Document}  % Declares the document's title.
\author{Leslie Lamport}      % Declares the author's name.
\date{January 21, 1994}      % Deleting this command produces today's date.

\newcommand{\ip}[2]{(#1, #2)}
                             % Defines \ip{arg1}{arg2} to mean
                             % (arg1, arg2).

%\newcommand{\ip}[2]{\langle #1 | #2\rangle}
                             % This is an alternative definition of
                             % \ip that is commented out.

\begin{document}             % End of preamble and beginning of text.

\maketitle                   % Produces the title.

This is an example input file.  Comparing it with
the output it generates can show you how to
produce a simple document of your own.

\section{Ordinary Text}      % Produces section heading.  Lower-level
                             % sections are begun with similar
                             % \subsection and \subsubsection commands.

The ends  of words and sentences are marked
  by   spaces. It  doesn't matter how many
spaces    you type; one is as good as 100.  The
end of   a line counts as a space.

One   or more   blank lines denote the  end
of  a paragraph.

Since any number of consecutive spaces are treated
like a single one, the formatting of the input
file makes no difference to
      \LaTeX,                % The \LaTeX command generates the LaTeX logo.
but it makes a difference to you.  When you use
\LaTeX, making your input file as easy to read
as possible will be a great help as you write
your document and when you change it.  This sample
file shows how you can add comments to your own input
file.

Because printing is different from typewriting,
there are a number of things that you have to do
differently when preparing an input file than if
you were just typing the document directly.
Quotation marks like
       ``this''
have to be handled specially, as do quotes within
quotes:
       ``\,`this'            % \, separates the double and single quote.
        is what I just
        wrote, not  `that'\,''.

Dashes come in three sizes: an
       intra-word
dash, a medium dash for number ranges like
       1--2,
and a punctuation
       dash---like
this.

A sentence-ending space should be larger than the
space between words within a sentence.  You
sometimes have to type special commands in
conjunction with punctuation characters to get
this right, as in the following sentence.
       Gnats, gnus, etc.\ all  % `\ ' makes an inter-word space.
       begin with G\@.         % \@ marks end-of-sentence punctuation.
You should check the spaces after periods when
reading your output to make sure you haven't
forgotten any special cases.  Generating an
ellipsis
       \ldots\               % `\ ' is needed after `\ldots' because TeX
                             % ignores spaces after command names like \ldots
                             % made from \ + letters.
                             %
                             % Note how a `%' character causes TeX to ignore
                             % the end of the input line, so these blank lines
                             % do not start a new paragraph.
                             %
with the right spacing around the periods requires
a special command.

\LaTeX\ interprets some common characters as
commands, so you must type special commands to
generate them.  These characters include the
following:
       \$ \& \% \# \{ and \}.

In printing, text is usually emphasized with an
       \emph{italic}
type style.

\begin{em}
   A long segment of text can also be emphasized
   in this way.  Text within such a segment can be
   given \emph{additional} emphasis.
\end{em}

It is sometimes necessary to prevent \LaTeX\ from
breaking a line where it might otherwise do so.
This may be at a space, as between the ``Mr.''\ and
``Jones'' in
       ``Mr.~Jones'',        % ~ produces an unbreakable interword space.
or within a word---especially when the word is a
symbol like
       \mbox{\emph{itemnum}}
that makes little sense when hyphenated across
lines.

Footnotes\footnote{This is an example of a footnote.}
pose no problem.

\LaTeX\ is good at typesetting mathematical formulas
like
       \( x-3y + z = 7 \)
or
       \( a_{1} > x^{2n} + y^{2n} > x' \)
or
       \( \ip{A}{B} = \sum_{i} a_{i} b_{i} \).
The spaces you type in a formula are
ignored.  Remember that a letter like
       $x$                   % $ ... $  and  \( ... \)  are equivalent
is a formula when it denotes a mathematical
symbol, and it should be typed as one.

\section{Displayed Text}

Text is displayed by indenting it from the left
margin.  Quotations are commonly displayed.  There
are short quotations
\begin{quote}
   This is a short quotation.  It consists of a
   single paragraph of text.  See how it is formatted.
\end{quote}
and longer ones.
\begin{quotation}
   This is a longer quotation.  It consists of two
   paragraphs of text, neither of which are
   particularly interesting.

   This is the second paragraph of the quotation.  It
   is just as dull as the first paragraph.
\end{quotation}
Another frequently-displayed structure is a list.
The following is an example of an \emph{itemized}
list.
\begin{itemize}
   \item This is the first item of an itemized list.
         Each item in the list is marked with a ``tick''.
         You don't have to worry about what kind of tick
         mark is used.

   \item This is the second item of the list.  It
         contains another list nested inside it.  The inner
         list is an \emph{enumerated} list.
         \begin{enumerate}
            \item This is the first item of an enumerated
                  list that is nested within the itemized list.

            \item This is the second item of the inner list.
                  \LaTeX\ allows you to nest lists deeper than
                  you really should.
         \end{enumerate}
         This is the rest of the second item of the outer
         list.  It is no more interesting than any other
         part of the item.
   \item This is the third item of the list.
\end{itemize}
You can even display poetry.
\begin{verse}
   There is an environment
    for verse \\             % The \\ command separates lines
   Whose features some poets % within a stanza.
   will curse.

                             % One or more blank lines separate stanzas.

   For instead of making\\
   Them do \emph{all} line breaking, \\
   It allows them to put too many words on a line when they'd rather be
   forced to be terse.
\end{verse}

Mathematical formulas may also be displayed.  A
displayed formula
is
one-line long; multiline
formulas require special formatting instructions.
   \[  \ip{\Gamma}{\psi'} = x'' + y^{2} + z_{i}^{n}\]
Don't start a paragraph with a displayed equation,
nor make one a paragraph by itself.

\end{document}               % End of document.

\section{Bench Huge Section 783}
% synthetic bench section 783
% This is a sample LaTeX input file.  (Version of 12 August 2004.)
%
% A '%' character causes TeX to ignore all remaining text on the line,
% and is used for comments like this one.

\documentclass{article}      % Specifies the document class

                             % The preamble begins here.
\title{An Example Document}  % Declares the document's title.
\author{Leslie Lamport}      % Declares the author's name.
\date{January 21, 1994}      % Deleting this command produces today's date.

\newcommand{\ip}[2]{(#1, #2)}
                             % Defines \ip{arg1}{arg2} to mean
                             % (arg1, arg2).

%\newcommand{\ip}[2]{\langle #1 | #2\rangle}
                             % This is an alternative definition of
                             % \ip that is commented out.

\begin{document}             % End of preamble and beginning of text.

\maketitle                   % Produces the title.

This is an example input file.  Comparing it with
the output it generates can show you how to
produce a simple document of your own.

\section{Ordinary Text}      % Produces section heading.  Lower-level
                             % sections are begun with similar
                             % \subsection and \subsubsection commands.

The ends  of words and sentences are marked
  by   spaces. It  doesn't matter how many
spaces    you type; one is as good as 100.  The
end of   a line counts as a space.

One   or more   blank lines denote the  end
of  a paragraph.

Since any number of consecutive spaces are treated
like a single one, the formatting of the input
file makes no difference to
      \LaTeX,                % The \LaTeX command generates the LaTeX logo.
but it makes a difference to you.  When you use
\LaTeX, making your input file as easy to read
as possible will be a great help as you write
your document and when you change it.  This sample
file shows how you can add comments to your own input
file.

Because printing is different from typewriting,
there are a number of things that you have to do
differently when preparing an input file than if
you were just typing the document directly.
Quotation marks like
       ``this''
have to be handled specially, as do quotes within
quotes:
       ``\,`this'            % \, separates the double and single quote.
        is what I just
        wrote, not  `that'\,''.

Dashes come in three sizes: an
       intra-word
dash, a medium dash for number ranges like
       1--2,
and a punctuation
       dash---like
this.

A sentence-ending space should be larger than the
space between words within a sentence.  You
sometimes have to type special commands in
conjunction with punctuation characters to get
this right, as in the following sentence.
       Gnats, gnus, etc.\ all  % `\ ' makes an inter-word space.
       begin with G\@.         % \@ marks end-of-sentence punctuation.
You should check the spaces after periods when
reading your output to make sure you haven't
forgotten any special cases.  Generating an
ellipsis
       \ldots\               % `\ ' is needed after `\ldots' because TeX
                             % ignores spaces after command names like \ldots
                             % made from \ + letters.
                             %
                             % Note how a `%' character causes TeX to ignore
                             % the end of the input line, so these blank lines
                             % do not start a new paragraph.
                             %
with the right spacing around the periods requires
a special command.

\LaTeX\ interprets some common characters as
commands, so you must type special commands to
generate them.  These characters include the
following:
       \$ \& \% \# \{ and \}.

In printing, text is usually emphasized with an
       \emph{italic}
type style.

\begin{em}
   A long segment of text can also be emphasized
   in this way.  Text within such a segment can be
   given \emph{additional} emphasis.
\end{em}

It is sometimes necessary to prevent \LaTeX\ from
breaking a line where it might otherwise do so.
This may be at a space, as between the ``Mr.''\ and
``Jones'' in
       ``Mr.~Jones'',        % ~ produces an unbreakable interword space.
or within a word---especially when the word is a
symbol like
       \mbox{\emph{itemnum}}
that makes little sense when hyphenated across
lines.

Footnotes\footnote{This is an example of a footnote.}
pose no problem.

\LaTeX\ is good at typesetting mathematical formulas
like
       \( x-3y + z = 7 \)
or
       \( a_{1} > x^{2n} + y^{2n} > x' \)
or
       \( \ip{A}{B} = \sum_{i} a_{i} b_{i} \).
The spaces you type in a formula are
ignored.  Remember that a letter like
       $x$                   % $ ... $  and  \( ... \)  are equivalent
is a formula when it denotes a mathematical
symbol, and it should be typed as one.

\section{Displayed Text}

Text is displayed by indenting it from the left
margin.  Quotations are commonly displayed.  There
are short quotations
\begin{quote}
   This is a short quotation.  It consists of a
   single paragraph of text.  See how it is formatted.
\end{quote}
and longer ones.
\begin{quotation}
   This is a longer quotation.  It consists of two
   paragraphs of text, neither of which are
   particularly interesting.

   This is the second paragraph of the quotation.  It
   is just as dull as the first paragraph.
\end{quotation}
Another frequently-displayed structure is a list.
The following is an example of an \emph{itemized}
list.
\begin{itemize}
   \item This is the first item of an itemized list.
         Each item in the list is marked with a ``tick''.
         You don't have to worry about what kind of tick
         mark is used.

   \item This is the second item of the list.  It
         contains another list nested inside it.  The inner
         list is an \emph{enumerated} list.
         \begin{enumerate}
            \item This is the first item of an enumerated
                  list that is nested within the itemized list.

            \item This is the second item of the inner list.
                  \LaTeX\ allows you to nest lists deeper than
                  you really should.
         \end{enumerate}
         This is the rest of the second item of the outer
         list.  It is no more interesting than any other
         part of the item.
   \item This is the third item of the list.
\end{itemize}
You can even display poetry.
\begin{verse}
   There is an environment
    for verse \\             % The \\ command separates lines
   Whose features some poets % within a stanza.
   will curse.

                             % One or more blank lines separate stanzas.

   For instead of making\\
   Them do \emph{all} line breaking, \\
   It allows them to put too many words on a line when they'd rather be
   forced to be terse.
\end{verse}

Mathematical formulas may also be displayed.  A
displayed formula
is
one-line long; multiline
formulas require special formatting instructions.
   \[  \ip{\Gamma}{\psi'} = x'' + y^{2} + z_{i}^{n}\]
Don't start a paragraph with a displayed equation,
nor make one a paragraph by itself.

\end{document}               % End of document.

\section{Bench Huge Section 784}
% synthetic bench section 784
% This is a sample LaTeX input file.  (Version of 12 August 2004.)
%
% A '%' character causes TeX to ignore all remaining text on the line,
% and is used for comments like this one.

\documentclass{article}      % Specifies the document class

                             % The preamble begins here.
\title{An Example Document}  % Declares the document's title.
\author{Leslie Lamport}      % Declares the author's name.
\date{January 21, 1994}      % Deleting this command produces today's date.

\newcommand{\ip}[2]{(#1, #2)}
                             % Defines \ip{arg1}{arg2} to mean
                             % (arg1, arg2).

%\newcommand{\ip}[2]{\langle #1 | #2\rangle}
                             % This is an alternative definition of
                             % \ip that is commented out.

\begin{document}             % End of preamble and beginning of text.

\maketitle                   % Produces the title.

This is an example input file.  Comparing it with
the output it generates can show you how to
produce a simple document of your own.

\section{Ordinary Text}      % Produces section heading.  Lower-level
                             % sections are begun with similar
                             % \subsection and \subsubsection commands.

The ends  of words and sentences are marked
  by   spaces. It  doesn't matter how many
spaces    you type; one is as good as 100.  The
end of   a line counts as a space.

One   or more   blank lines denote the  end
of  a paragraph.

Since any number of consecutive spaces are treated
like a single one, the formatting of the input
file makes no difference to
      \LaTeX,                % The \LaTeX command generates the LaTeX logo.
but it makes a difference to you.  When you use
\LaTeX, making your input file as easy to read
as possible will be a great help as you write
your document and when you change it.  This sample
file shows how you can add comments to your own input
file.

Because printing is different from typewriting,
there are a number of things that you have to do
differently when preparing an input file than if
you were just typing the document directly.
Quotation marks like
       ``this''
have to be handled specially, as do quotes within
quotes:
       ``\,`this'            % \, separates the double and single quote.
        is what I just
        wrote, not  `that'\,''.

Dashes come in three sizes: an
       intra-word
dash, a medium dash for number ranges like
       1--2,
and a punctuation
       dash---like
this.

A sentence-ending space should be larger than the
space between words within a sentence.  You
sometimes have to type special commands in
conjunction with punctuation characters to get
this right, as in the following sentence.
       Gnats, gnus, etc.\ all  % `\ ' makes an inter-word space.
       begin with G\@.         % \@ marks end-of-sentence punctuation.
You should check the spaces after periods when
reading your output to make sure you haven't
forgotten any special cases.  Generating an
ellipsis
       \ldots\               % `\ ' is needed after `\ldots' because TeX
                             % ignores spaces after command names like \ldots
                             % made from \ + letters.
                             %
                             % Note how a `%' character causes TeX to ignore
                             % the end of the input line, so these blank lines
                             % do not start a new paragraph.
                             %
with the right spacing around the periods requires
a special command.

\LaTeX\ interprets some common characters as
commands, so you must type special commands to
generate them.  These characters include the
following:
       \$ \& \% \# \{ and \}.

In printing, text is usually emphasized with an
       \emph{italic}
type style.

\begin{em}
   A long segment of text can also be emphasized
   in this way.  Text within such a segment can be
   given \emph{additional} emphasis.
\end{em}

It is sometimes necessary to prevent \LaTeX\ from
breaking a line where it might otherwise do so.
This may be at a space, as between the ``Mr.''\ and
``Jones'' in
       ``Mr.~Jones'',        % ~ produces an unbreakable interword space.
or within a word---especially when the word is a
symbol like
       \mbox{\emph{itemnum}}
that makes little sense when hyphenated across
lines.

Footnotes\footnote{This is an example of a footnote.}
pose no problem.

\LaTeX\ is good at typesetting mathematical formulas
like
       \( x-3y + z = 7 \)
or
       \( a_{1} > x^{2n} + y^{2n} > x' \)
or
       \( \ip{A}{B} = \sum_{i} a_{i} b_{i} \).
The spaces you type in a formula are
ignored.  Remember that a letter like
       $x$                   % $ ... $  and  \( ... \)  are equivalent
is a formula when it denotes a mathematical
symbol, and it should be typed as one.

\section{Displayed Text}

Text is displayed by indenting it from the left
margin.  Quotations are commonly displayed.  There
are short quotations
\begin{quote}
   This is a short quotation.  It consists of a
   single paragraph of text.  See how it is formatted.
\end{quote}
and longer ones.
\begin{quotation}
   This is a longer quotation.  It consists of two
   paragraphs of text, neither of which are
   particularly interesting.

   This is the second paragraph of the quotation.  It
   is just as dull as the first paragraph.
\end{quotation}
Another frequently-displayed structure is a list.
The following is an example of an \emph{itemized}
list.
\begin{itemize}
   \item This is the first item of an itemized list.
         Each item in the list is marked with a ``tick''.
         You don't have to worry about what kind of tick
         mark is used.

   \item This is the second item of the list.  It
         contains another list nested inside it.  The inner
         list is an \emph{enumerated} list.
         \begin{enumerate}
            \item This is the first item of an enumerated
                  list that is nested within the itemized list.

            \item This is the second item of the inner list.
                  \LaTeX\ allows you to nest lists deeper than
                  you really should.
         \end{enumerate}
         This is the rest of the second item of the outer
         list.  It is no more interesting than any other
         part of the item.
   \item This is the third item of the list.
\end{itemize}
You can even display poetry.
\begin{verse}
   There is an environment
    for verse \\             % The \\ command separates lines
   Whose features some poets % within a stanza.
   will curse.

                             % One or more blank lines separate stanzas.

   For instead of making\\
   Them do \emph{all} line breaking, \\
   It allows them to put too many words on a line when they'd rather be
   forced to be terse.
\end{verse}

Mathematical formulas may also be displayed.  A
displayed formula
is
one-line long; multiline
formulas require special formatting instructions.
   \[  \ip{\Gamma}{\psi'} = x'' + y^{2} + z_{i}^{n}\]
Don't start a paragraph with a displayed equation,
nor make one a paragraph by itself.

\end{document}               % End of document.

\section{Bench Huge Section 785}
% synthetic bench section 785
% This is a sample LaTeX input file.  (Version of 12 August 2004.)
%
% A '%' character causes TeX to ignore all remaining text on the line,
% and is used for comments like this one.

\documentclass{article}      % Specifies the document class

                             % The preamble begins here.
\title{An Example Document}  % Declares the document's title.
\author{Leslie Lamport}      % Declares the author's name.
\date{January 21, 1994}      % Deleting this command produces today's date.

\newcommand{\ip}[2]{(#1, #2)}
                             % Defines \ip{arg1}{arg2} to mean
                             % (arg1, arg2).

%\newcommand{\ip}[2]{\langle #1 | #2\rangle}
                             % This is an alternative definition of
                             % \ip that is commented out.

\begin{document}             % End of preamble and beginning of text.

\maketitle                   % Produces the title.

This is an example input file.  Comparing it with
the output it generates can show you how to
produce a simple document of your own.

\section{Ordinary Text}      % Produces section heading.  Lower-level
                             % sections are begun with similar
                             % \subsection and \subsubsection commands.

The ends  of words and sentences are marked
  by   spaces. It  doesn't matter how many
spaces    you type; one is as good as 100.  The
end of   a line counts as a space.

One   or more   blank lines denote the  end
of  a paragraph.

Since any number of consecutive spaces are treated
like a single one, the formatting of the input
file makes no difference to
      \LaTeX,                % The \LaTeX command generates the LaTeX logo.
but it makes a difference to you.  When you use
\LaTeX, making your input file as easy to read
as possible will be a great help as you write
your document and when you change it.  This sample
file shows how you can add comments to your own input
file.

Because printing is different from typewriting,
there are a number of things that you have to do
differently when preparing an input file than if
you were just typing the document directly.
Quotation marks like
       ``this''
have to be handled specially, as do quotes within
quotes:
       ``\,`this'            % \, separates the double and single quote.
        is what I just
        wrote, not  `that'\,''.

Dashes come in three sizes: an
       intra-word
dash, a medium dash for number ranges like
       1--2,
and a punctuation
       dash---like
this.

A sentence-ending space should be larger than the
space between words within a sentence.  You
sometimes have to type special commands in
conjunction with punctuation characters to get
this right, as in the following sentence.
       Gnats, gnus, etc.\ all  % `\ ' makes an inter-word space.
       begin with G\@.         % \@ marks end-of-sentence punctuation.
You should check the spaces after periods when
reading your output to make sure you haven't
forgotten any special cases.  Generating an
ellipsis
       \ldots\               % `\ ' is needed after `\ldots' because TeX
                             % ignores spaces after command names like \ldots
                             % made from \ + letters.
                             %
                             % Note how a `%' character causes TeX to ignore
                             % the end of the input line, so these blank lines
                             % do not start a new paragraph.
                             %
with the right spacing around the periods requires
a special command.

\LaTeX\ interprets some common characters as
commands, so you must type special commands to
generate them.  These characters include the
following:
       \$ \& \% \# \{ and \}.

In printing, text is usually emphasized with an
       \emph{italic}
type style.

\begin{em}
   A long segment of text can also be emphasized
   in this way.  Text within such a segment can be
   given \emph{additional} emphasis.
\end{em}

It is sometimes necessary to prevent \LaTeX\ from
breaking a line where it might otherwise do so.
This may be at a space, as between the ``Mr.''\ and
``Jones'' in
       ``Mr.~Jones'',        % ~ produces an unbreakable interword space.
or within a word---especially when the word is a
symbol like
       \mbox{\emph{itemnum}}
that makes little sense when hyphenated across
lines.

Footnotes\footnote{This is an example of a footnote.}
pose no problem.

\LaTeX\ is good at typesetting mathematical formulas
like
       \( x-3y + z = 7 \)
or
       \( a_{1} > x^{2n} + y^{2n} > x' \)
or
       \( \ip{A}{B} = \sum_{i} a_{i} b_{i} \).
The spaces you type in a formula are
ignored.  Remember that a letter like
       $x$                   % $ ... $  and  \( ... \)  are equivalent
is a formula when it denotes a mathematical
symbol, and it should be typed as one.

\section{Displayed Text}

Text is displayed by indenting it from the left
margin.  Quotations are commonly displayed.  There
are short quotations
\begin{quote}
   This is a short quotation.  It consists of a
   single paragraph of text.  See how it is formatted.
\end{quote}
and longer ones.
\begin{quotation}
   This is a longer quotation.  It consists of two
   paragraphs of text, neither of which are
   particularly interesting.

   This is the second paragraph of the quotation.  It
   is just as dull as the first paragraph.
\end{quotation}
Another frequently-displayed structure is a list.
The following is an example of an \emph{itemized}
list.
\begin{itemize}
   \item This is the first item of an itemized list.
         Each item in the list is marked with a ``tick''.
         You don't have to worry about what kind of tick
         mark is used.

   \item This is the second item of the list.  It
         contains another list nested inside it.  The inner
         list is an \emph{enumerated} list.
         \begin{enumerate}
            \item This is the first item of an enumerated
                  list that is nested within the itemized list.

            \item This is the second item of the inner list.
                  \LaTeX\ allows you to nest lists deeper than
                  you really should.
         \end{enumerate}
         This is the rest of the second item of the outer
         list.  It is no more interesting than any other
         part of the item.
   \item This is the third item of the list.
\end{itemize}
You can even display poetry.
\begin{verse}
   There is an environment
    for verse \\             % The \\ command separates lines
   Whose features some poets % within a stanza.
   will curse.

                             % One or more blank lines separate stanzas.

   For instead of making\\
   Them do \emph{all} line breaking, \\
   It allows them to put too many words on a line when they'd rather be
   forced to be terse.
\end{verse}

Mathematical formulas may also be displayed.  A
displayed formula
is
one-line long; multiline
formulas require special formatting instructions.
   \[  \ip{\Gamma}{\psi'} = x'' + y^{2} + z_{i}^{n}\]
Don't start a paragraph with a displayed equation,
nor make one a paragraph by itself.

\end{document}               % End of document.

\section{Bench Huge Section 786}
% synthetic bench section 786
% This is a sample LaTeX input file.  (Version of 12 August 2004.)
%
% A '%' character causes TeX to ignore all remaining text on the line,
% and is used for comments like this one.

\documentclass{article}      % Specifies the document class

                             % The preamble begins here.
\title{An Example Document}  % Declares the document's title.
\author{Leslie Lamport}      % Declares the author's name.
\date{January 21, 1994}      % Deleting this command produces today's date.

\newcommand{\ip}[2]{(#1, #2)}
                             % Defines \ip{arg1}{arg2} to mean
                             % (arg1, arg2).

%\newcommand{\ip}[2]{\langle #1 | #2\rangle}
                             % This is an alternative definition of
                             % \ip that is commented out.

\begin{document}             % End of preamble and beginning of text.

\maketitle                   % Produces the title.

This is an example input file.  Comparing it with
the output it generates can show you how to
produce a simple document of your own.

\section{Ordinary Text}      % Produces section heading.  Lower-level
                             % sections are begun with similar
                             % \subsection and \subsubsection commands.

The ends  of words and sentences are marked
  by   spaces. It  doesn't matter how many
spaces    you type; one is as good as 100.  The
end of   a line counts as a space.

One   or more   blank lines denote the  end
of  a paragraph.

Since any number of consecutive spaces are treated
like a single one, the formatting of the input
file makes no difference to
      \LaTeX,                % The \LaTeX command generates the LaTeX logo.
but it makes a difference to you.  When you use
\LaTeX, making your input file as easy to read
as possible will be a great help as you write
your document and when you change it.  This sample
file shows how you can add comments to your own input
file.

Because printing is different from typewriting,
there are a number of things that you have to do
differently when preparing an input file than if
you were just typing the document directly.
Quotation marks like
       ``this''
have to be handled specially, as do quotes within
quotes:
       ``\,`this'            % \, separates the double and single quote.
        is what I just
        wrote, not  `that'\,''.

Dashes come in three sizes: an
       intra-word
dash, a medium dash for number ranges like
       1--2,
and a punctuation
       dash---like
this.

A sentence-ending space should be larger than the
space between words within a sentence.  You
sometimes have to type special commands in
conjunction with punctuation characters to get
this right, as in the following sentence.
       Gnats, gnus, etc.\ all  % `\ ' makes an inter-word space.
       begin with G\@.         % \@ marks end-of-sentence punctuation.
You should check the spaces after periods when
reading your output to make sure you haven't
forgotten any special cases.  Generating an
ellipsis
       \ldots\               % `\ ' is needed after `\ldots' because TeX
                             % ignores spaces after command names like \ldots
                             % made from \ + letters.
                             %
                             % Note how a `%' character causes TeX to ignore
                             % the end of the input line, so these blank lines
                             % do not start a new paragraph.
                             %
with the right spacing around the periods requires
a special command.

\LaTeX\ interprets some common characters as
commands, so you must type special commands to
generate them.  These characters include the
following:
       \$ \& \% \# \{ and \}.

In printing, text is usually emphasized with an
       \emph{italic}
type style.

\begin{em}
   A long segment of text can also be emphasized
   in this way.  Text within such a segment can be
   given \emph{additional} emphasis.
\end{em}

It is sometimes necessary to prevent \LaTeX\ from
breaking a line where it might otherwise do so.
This may be at a space, as between the ``Mr.''\ and
``Jones'' in
       ``Mr.~Jones'',        % ~ produces an unbreakable interword space.
or within a word---especially when the word is a
symbol like
       \mbox{\emph{itemnum}}
that makes little sense when hyphenated across
lines.

Footnotes\footnote{This is an example of a footnote.}
pose no problem.

\LaTeX\ is good at typesetting mathematical formulas
like
       \( x-3y + z = 7 \)
or
       \( a_{1} > x^{2n} + y^{2n} > x' \)
or
       \( \ip{A}{B} = \sum_{i} a_{i} b_{i} \).
The spaces you type in a formula are
ignored.  Remember that a letter like
       $x$                   % $ ... $  and  \( ... \)  are equivalent
is a formula when it denotes a mathematical
symbol, and it should be typed as one.

\section{Displayed Text}

Text is displayed by indenting it from the left
margin.  Quotations are commonly displayed.  There
are short quotations
\begin{quote}
   This is a short quotation.  It consists of a
   single paragraph of text.  See how it is formatted.
\end{quote}
and longer ones.
\begin{quotation}
   This is a longer quotation.  It consists of two
   paragraphs of text, neither of which are
   particularly interesting.

   This is the second paragraph of the quotation.  It
   is just as dull as the first paragraph.
\end{quotation}
Another frequently-displayed structure is a list.
The following is an example of an \emph{itemized}
list.
\begin{itemize}
   \item This is the first item of an itemized list.
         Each item in the list is marked with a ``tick''.
         You don't have to worry about what kind of tick
         mark is used.

   \item This is the second item of the list.  It
         contains another list nested inside it.  The inner
         list is an \emph{enumerated} list.
         \begin{enumerate}
            \item This is the first item of an enumerated
                  list that is nested within the itemized list.

            \item This is the second item of the inner list.
                  \LaTeX\ allows you to nest lists deeper than
                  you really should.
         \end{enumerate}
         This is the rest of the second item of the outer
         list.  It is no more interesting than any other
         part of the item.
   \item This is the third item of the list.
\end{itemize}
You can even display poetry.
\begin{verse}
   There is an environment
    for verse \\             % The \\ command separates lines
   Whose features some poets % within a stanza.
   will curse.

                             % One or more blank lines separate stanzas.

   For instead of making\\
   Them do \emph{all} line breaking, \\
   It allows them to put too many words on a line when they'd rather be
   forced to be terse.
\end{verse}

Mathematical formulas may also be displayed.  A
displayed formula
is
one-line long; multiline
formulas require special formatting instructions.
   \[  \ip{\Gamma}{\psi'} = x'' + y^{2} + z_{i}^{n}\]
Don't start a paragraph with a displayed equation,
nor make one a paragraph by itself.

\end{document}               % End of document.

\section{Bench Huge Section 787}
% synthetic bench section 787
% This is a sample LaTeX input file.  (Version of 12 August 2004.)
%
% A '%' character causes TeX to ignore all remaining text on the line,
% and is used for comments like this one.

\documentclass{article}      % Specifies the document class

                             % The preamble begins here.
\title{An Example Document}  % Declares the document's title.
\author{Leslie Lamport}      % Declares the author's name.
\date{January 21, 1994}      % Deleting this command produces today's date.

\newcommand{\ip}[2]{(#1, #2)}
                             % Defines \ip{arg1}{arg2} to mean
                             % (arg1, arg2).

%\newcommand{\ip}[2]{\langle #1 | #2\rangle}
                             % This is an alternative definition of
                             % \ip that is commented out.

\begin{document}             % End of preamble and beginning of text.

\maketitle                   % Produces the title.

This is an example input file.  Comparing it with
the output it generates can show you how to
produce a simple document of your own.

\section{Ordinary Text}      % Produces section heading.  Lower-level
                             % sections are begun with similar
                             % \subsection and \subsubsection commands.

The ends  of words and sentences are marked
  by   spaces. It  doesn't matter how many
spaces    you type; one is as good as 100.  The
end of   a line counts as a space.

One   or more   blank lines denote the  end
of  a paragraph.

Since any number of consecutive spaces are treated
like a single one, the formatting of the input
file makes no difference to
      \LaTeX,                % The \LaTeX command generates the LaTeX logo.
but it makes a difference to you.  When you use
\LaTeX, making your input file as easy to read
as possible will be a great help as you write
your document and when you change it.  This sample
file shows how you can add comments to your own input
file.

Because printing is different from typewriting,
there are a number of things that you have to do
differently when preparing an input file than if
you were just typing the document directly.
Quotation marks like
       ``this''
have to be handled specially, as do quotes within
quotes:
       ``\,`this'            % \, separates the double and single quote.
        is what I just
        wrote, not  `that'\,''.

Dashes come in three sizes: an
       intra-word
dash, a medium dash for number ranges like
       1--2,
and a punctuation
       dash---like
this.

A sentence-ending space should be larger than the
space between words within a sentence.  You
sometimes have to type special commands in
conjunction with punctuation characters to get
this right, as in the following sentence.
       Gnats, gnus, etc.\ all  % `\ ' makes an inter-word space.
       begin with G\@.         % \@ marks end-of-sentence punctuation.
You should check the spaces after periods when
reading your output to make sure you haven't
forgotten any special cases.  Generating an
ellipsis
       \ldots\               % `\ ' is needed after `\ldots' because TeX
                             % ignores spaces after command names like \ldots
                             % made from \ + letters.
                             %
                             % Note how a `%' character causes TeX to ignore
                             % the end of the input line, so these blank lines
                             % do not start a new paragraph.
                             %
with the right spacing around the periods requires
a special command.

\LaTeX\ interprets some common characters as
commands, so you must type special commands to
generate them.  These characters include the
following:
       \$ \& \% \# \{ and \}.

In printing, text is usually emphasized with an
       \emph{italic}
type style.

\begin{em}
   A long segment of text can also be emphasized
   in this way.  Text within such a segment can be
   given \emph{additional} emphasis.
\end{em}

It is sometimes necessary to prevent \LaTeX\ from
breaking a line where it might otherwise do so.
This may be at a space, as between the ``Mr.''\ and
``Jones'' in
       ``Mr.~Jones'',        % ~ produces an unbreakable interword space.
or within a word---especially when the word is a
symbol like
       \mbox{\emph{itemnum}}
that makes little sense when hyphenated across
lines.

Footnotes\footnote{This is an example of a footnote.}
pose no problem.

\LaTeX\ is good at typesetting mathematical formulas
like
       \( x-3y + z = 7 \)
or
       \( a_{1} > x^{2n} + y^{2n} > x' \)
or
       \( \ip{A}{B} = \sum_{i} a_{i} b_{i} \).
The spaces you type in a formula are
ignored.  Remember that a letter like
       $x$                   % $ ... $  and  \( ... \)  are equivalent
is a formula when it denotes a mathematical
symbol, and it should be typed as one.

\section{Displayed Text}

Text is displayed by indenting it from the left
margin.  Quotations are commonly displayed.  There
are short quotations
\begin{quote}
   This is a short quotation.  It consists of a
   single paragraph of text.  See how it is formatted.
\end{quote}
and longer ones.
\begin{quotation}
   This is a longer quotation.  It consists of two
   paragraphs of text, neither of which are
   particularly interesting.

   This is the second paragraph of the quotation.  It
   is just as dull as the first paragraph.
\end{quotation}
Another frequently-displayed structure is a list.
The following is an example of an \emph{itemized}
list.
\begin{itemize}
   \item This is the first item of an itemized list.
         Each item in the list is marked with a ``tick''.
         You don't have to worry about what kind of tick
         mark is used.

   \item This is the second item of the list.  It
         contains another list nested inside it.  The inner
         list is an \emph{enumerated} list.
         \begin{enumerate}
            \item This is the first item of an enumerated
                  list that is nested within the itemized list.

            \item This is the second item of the inner list.
                  \LaTeX\ allows you to nest lists deeper than
                  you really should.
         \end{enumerate}
         This is the rest of the second item of the outer
         list.  It is no more interesting than any other
         part of the item.
   \item This is the third item of the list.
\end{itemize}
You can even display poetry.
\begin{verse}
   There is an environment
    for verse \\             % The \\ command separates lines
   Whose features some poets % within a stanza.
   will curse.

                             % One or more blank lines separate stanzas.

   For instead of making\\
   Them do \emph{all} line breaking, \\
   It allows them to put too many words on a line when they'd rather be
   forced to be terse.
\end{verse}

Mathematical formulas may also be displayed.  A
displayed formula
is
one-line long; multiline
formulas require special formatting instructions.
   \[  \ip{\Gamma}{\psi'} = x'' + y^{2} + z_{i}^{n}\]
Don't start a paragraph with a displayed equation,
nor make one a paragraph by itself.

\end{document}               % End of document.

\section{Bench Huge Section 788}
% synthetic bench section 788
% This is a sample LaTeX input file.  (Version of 12 August 2004.)
%
% A '%' character causes TeX to ignore all remaining text on the line,
% and is used for comments like this one.

\documentclass{article}      % Specifies the document class

                             % The preamble begins here.
\title{An Example Document}  % Declares the document's title.
\author{Leslie Lamport}      % Declares the author's name.
\date{January 21, 1994}      % Deleting this command produces today's date.

\newcommand{\ip}[2]{(#1, #2)}
                             % Defines \ip{arg1}{arg2} to mean
                             % (arg1, arg2).

%\newcommand{\ip}[2]{\langle #1 | #2\rangle}
                             % This is an alternative definition of
                             % \ip that is commented out.

\begin{document}             % End of preamble and beginning of text.

\maketitle                   % Produces the title.

This is an example input file.  Comparing it with
the output it generates can show you how to
produce a simple document of your own.

\section{Ordinary Text}      % Produces section heading.  Lower-level
                             % sections are begun with similar
                             % \subsection and \subsubsection commands.

The ends  of words and sentences are marked
  by   spaces. It  doesn't matter how many
spaces    you type; one is as good as 100.  The
end of   a line counts as a space.

One   or more   blank lines denote the  end
of  a paragraph.

Since any number of consecutive spaces are treated
like a single one, the formatting of the input
file makes no difference to
      \LaTeX,                % The \LaTeX command generates the LaTeX logo.
but it makes a difference to you.  When you use
\LaTeX, making your input file as easy to read
as possible will be a great help as you write
your document and when you change it.  This sample
file shows how you can add comments to your own input
file.

Because printing is different from typewriting,
there are a number of things that you have to do
differently when preparing an input file than if
you were just typing the document directly.
Quotation marks like
       ``this''
have to be handled specially, as do quotes within
quotes:
       ``\,`this'            % \, separates the double and single quote.
        is what I just
        wrote, not  `that'\,''.

Dashes come in three sizes: an
       intra-word
dash, a medium dash for number ranges like
       1--2,
and a punctuation
       dash---like
this.

A sentence-ending space should be larger than the
space between words within a sentence.  You
sometimes have to type special commands in
conjunction with punctuation characters to get
this right, as in the following sentence.
       Gnats, gnus, etc.\ all  % `\ ' makes an inter-word space.
       begin with G\@.         % \@ marks end-of-sentence punctuation.
You should check the spaces after periods when
reading your output to make sure you haven't
forgotten any special cases.  Generating an
ellipsis
       \ldots\               % `\ ' is needed after `\ldots' because TeX
                             % ignores spaces after command names like \ldots
                             % made from \ + letters.
                             %
                             % Note how a `%' character causes TeX to ignore
                             % the end of the input line, so these blank lines
                             % do not start a new paragraph.
                             %
with the right spacing around the periods requires
a special command.

\LaTeX\ interprets some common characters as
commands, so you must type special commands to
generate them.  These characters include the
following:
       \$ \& \% \# \{ and \}.

In printing, text is usually emphasized with an
       \emph{italic}
type style.

\begin{em}
   A long segment of text can also be emphasized
   in this way.  Text within such a segment can be
   given \emph{additional} emphasis.
\end{em}

It is sometimes necessary to prevent \LaTeX\ from
breaking a line where it might otherwise do so.
This may be at a space, as between the ``Mr.''\ and
``Jones'' in
       ``Mr.~Jones'',        % ~ produces an unbreakable interword space.
or within a word---especially when the word is a
symbol like
       \mbox{\emph{itemnum}}
that makes little sense when hyphenated across
lines.

Footnotes\footnote{This is an example of a footnote.}
pose no problem.

\LaTeX\ is good at typesetting mathematical formulas
like
       \( x-3y + z = 7 \)
or
       \( a_{1} > x^{2n} + y^{2n} > x' \)
or
       \( \ip{A}{B} = \sum_{i} a_{i} b_{i} \).
The spaces you type in a formula are
ignored.  Remember that a letter like
       $x$                   % $ ... $  and  \( ... \)  are equivalent
is a formula when it denotes a mathematical
symbol, and it should be typed as one.

\section{Displayed Text}

Text is displayed by indenting it from the left
margin.  Quotations are commonly displayed.  There
are short quotations
\begin{quote}
   This is a short quotation.  It consists of a
   single paragraph of text.  See how it is formatted.
\end{quote}
and longer ones.
\begin{quotation}
   This is a longer quotation.  It consists of two
   paragraphs of text, neither of which are
   particularly interesting.

   This is the second paragraph of the quotation.  It
   is just as dull as the first paragraph.
\end{quotation}
Another frequently-displayed structure is a list.
The following is an example of an \emph{itemized}
list.
\begin{itemize}
   \item This is the first item of an itemized list.
         Each item in the list is marked with a ``tick''.
         You don't have to worry about what kind of tick
         mark is used.

   \item This is the second item of the list.  It
         contains another list nested inside it.  The inner
         list is an \emph{enumerated} list.
         \begin{enumerate}
            \item This is the first item of an enumerated
                  list that is nested within the itemized list.

            \item This is the second item of the inner list.
                  \LaTeX\ allows you to nest lists deeper than
                  you really should.
         \end{enumerate}
         This is the rest of the second item of the outer
         list.  It is no more interesting than any other
         part of the item.
   \item This is the third item of the list.
\end{itemize}
You can even display poetry.
\begin{verse}
   There is an environment
    for verse \\             % The \\ command separates lines
   Whose features some poets % within a stanza.
   will curse.

                             % One or more blank lines separate stanzas.

   For instead of making\\
   Them do \emph{all} line breaking, \\
   It allows them to put too many words on a line when they'd rather be
   forced to be terse.
\end{verse}

Mathematical formulas may also be displayed.  A
displayed formula
is
one-line long; multiline
formulas require special formatting instructions.
   \[  \ip{\Gamma}{\psi'} = x'' + y^{2} + z_{i}^{n}\]
Don't start a paragraph with a displayed equation,
nor make one a paragraph by itself.

\end{document}               % End of document.

\section{Bench Huge Section 789}
% synthetic bench section 789
% This is a sample LaTeX input file.  (Version of 12 August 2004.)
%
% A '%' character causes TeX to ignore all remaining text on the line,
% and is used for comments like this one.

\documentclass{article}      % Specifies the document class

                             % The preamble begins here.
\title{An Example Document}  % Declares the document's title.
\author{Leslie Lamport}      % Declares the author's name.
\date{January 21, 1994}      % Deleting this command produces today's date.

\newcommand{\ip}[2]{(#1, #2)}
                             % Defines \ip{arg1}{arg2} to mean
                             % (arg1, arg2).

%\newcommand{\ip}[2]{\langle #1 | #2\rangle}
                             % This is an alternative definition of
                             % \ip that is commented out.

\begin{document}             % End of preamble and beginning of text.

\maketitle                   % Produces the title.

This is an example input file.  Comparing it with
the output it generates can show you how to
produce a simple document of your own.

\section{Ordinary Text}      % Produces section heading.  Lower-level
                             % sections are begun with similar
                             % \subsection and \subsubsection commands.

The ends  of words and sentences are marked
  by   spaces. It  doesn't matter how many
spaces    you type; one is as good as 100.  The
end of   a line counts as a space.

One   or more   blank lines denote the  end
of  a paragraph.

Since any number of consecutive spaces are treated
like a single one, the formatting of the input
file makes no difference to
      \LaTeX,                % The \LaTeX command generates the LaTeX logo.
but it makes a difference to you.  When you use
\LaTeX, making your input file as easy to read
as possible will be a great help as you write
your document and when you change it.  This sample
file shows how you can add comments to your own input
file.

Because printing is different from typewriting,
there are a number of things that you have to do
differently when preparing an input file than if
you were just typing the document directly.
Quotation marks like
       ``this''
have to be handled specially, as do quotes within
quotes:
       ``\,`this'            % \, separates the double and single quote.
        is what I just
        wrote, not  `that'\,''.

Dashes come in three sizes: an
       intra-word
dash, a medium dash for number ranges like
       1--2,
and a punctuation
       dash---like
this.

A sentence-ending space should be larger than the
space between words within a sentence.  You
sometimes have to type special commands in
conjunction with punctuation characters to get
this right, as in the following sentence.
       Gnats, gnus, etc.\ all  % `\ ' makes an inter-word space.
       begin with G\@.         % \@ marks end-of-sentence punctuation.
You should check the spaces after periods when
reading your output to make sure you haven't
forgotten any special cases.  Generating an
ellipsis
       \ldots\               % `\ ' is needed after `\ldots' because TeX
                             % ignores spaces after command names like \ldots
                             % made from \ + letters.
                             %
                             % Note how a `%' character causes TeX to ignore
                             % the end of the input line, so these blank lines
                             % do not start a new paragraph.
                             %
with the right spacing around the periods requires
a special command.

\LaTeX\ interprets some common characters as
commands, so you must type special commands to
generate them.  These characters include the
following:
       \$ \& \% \# \{ and \}.

In printing, text is usually emphasized with an
       \emph{italic}
type style.

\begin{em}
   A long segment of text can also be emphasized
   in this way.  Text within such a segment can be
   given \emph{additional} emphasis.
\end{em}

It is sometimes necessary to prevent \LaTeX\ from
breaking a line where it might otherwise do so.
This may be at a space, as between the ``Mr.''\ and
``Jones'' in
       ``Mr.~Jones'',        % ~ produces an unbreakable interword space.
or within a word---especially when the word is a
symbol like
       \mbox{\emph{itemnum}}
that makes little sense when hyphenated across
lines.

Footnotes\footnote{This is an example of a footnote.}
pose no problem.

\LaTeX\ is good at typesetting mathematical formulas
like
       \( x-3y + z = 7 \)
or
       \( a_{1} > x^{2n} + y^{2n} > x' \)
or
       \( \ip{A}{B} = \sum_{i} a_{i} b_{i} \).
The spaces you type in a formula are
ignored.  Remember that a letter like
       $x$                   % $ ... $  and  \( ... \)  are equivalent
is a formula when it denotes a mathematical
symbol, and it should be typed as one.

\section{Displayed Text}

Text is displayed by indenting it from the left
margin.  Quotations are commonly displayed.  There
are short quotations
\begin{quote}
   This is a short quotation.  It consists of a
   single paragraph of text.  See how it is formatted.
\end{quote}
and longer ones.
\begin{quotation}
   This is a longer quotation.  It consists of two
   paragraphs of text, neither of which are
   particularly interesting.

   This is the second paragraph of the quotation.  It
   is just as dull as the first paragraph.
\end{quotation}
Another frequently-displayed structure is a list.
The following is an example of an \emph{itemized}
list.
\begin{itemize}
   \item This is the first item of an itemized list.
         Each item in the list is marked with a ``tick''.
         You don't have to worry about what kind of tick
         mark is used.

   \item This is the second item of the list.  It
         contains another list nested inside it.  The inner
         list is an \emph{enumerated} list.
         \begin{enumerate}
            \item This is the first item of an enumerated
                  list that is nested within the itemized list.

            \item This is the second item of the inner list.
                  \LaTeX\ allows you to nest lists deeper than
                  you really should.
         \end{enumerate}
         This is the rest of the second item of the outer
         list.  It is no more interesting than any other
         part of the item.
   \item This is the third item of the list.
\end{itemize}
You can even display poetry.
\begin{verse}
   There is an environment
    for verse \\             % The \\ command separates lines
   Whose features some poets % within a stanza.
   will curse.

                             % One or more blank lines separate stanzas.

   For instead of making\\
   Them do \emph{all} line breaking, \\
   It allows them to put too many words on a line when they'd rather be
   forced to be terse.
\end{verse}

Mathematical formulas may also be displayed.  A
displayed formula
is
one-line long; multiline
formulas require special formatting instructions.
   \[  \ip{\Gamma}{\psi'} = x'' + y^{2} + z_{i}^{n}\]
Don't start a paragraph with a displayed equation,
nor make one a paragraph by itself.

\end{document}               % End of document.

\section{Bench Huge Section 790}
% synthetic bench section 790
% This is a sample LaTeX input file.  (Version of 12 August 2004.)
%
% A '%' character causes TeX to ignore all remaining text on the line,
% and is used for comments like this one.

\documentclass{article}      % Specifies the document class

                             % The preamble begins here.
\title{An Example Document}  % Declares the document's title.
\author{Leslie Lamport}      % Declares the author's name.
\date{January 21, 1994}      % Deleting this command produces today's date.

\newcommand{\ip}[2]{(#1, #2)}
                             % Defines \ip{arg1}{arg2} to mean
                             % (arg1, arg2).

%\newcommand{\ip}[2]{\langle #1 | #2\rangle}
                             % This is an alternative definition of
                             % \ip that is commented out.

\begin{document}             % End of preamble and beginning of text.

\maketitle                   % Produces the title.

This is an example input file.  Comparing it with
the output it generates can show you how to
produce a simple document of your own.

\section{Ordinary Text}      % Produces section heading.  Lower-level
                             % sections are begun with similar
                             % \subsection and \subsubsection commands.

The ends  of words and sentences are marked
  by   spaces. It  doesn't matter how many
spaces    you type; one is as good as 100.  The
end of   a line counts as a space.

One   or more   blank lines denote the  end
of  a paragraph.

Since any number of consecutive spaces are treated
like a single one, the formatting of the input
file makes no difference to
      \LaTeX,                % The \LaTeX command generates the LaTeX logo.
but it makes a difference to you.  When you use
\LaTeX, making your input file as easy to read
as possible will be a great help as you write
your document and when you change it.  This sample
file shows how you can add comments to your own input
file.

Because printing is different from typewriting,
there are a number of things that you have to do
differently when preparing an input file than if
you were just typing the document directly.
Quotation marks like
       ``this''
have to be handled specially, as do quotes within
quotes:
       ``\,`this'            % \, separates the double and single quote.
        is what I just
        wrote, not  `that'\,''.

Dashes come in three sizes: an
       intra-word
dash, a medium dash for number ranges like
       1--2,
and a punctuation
       dash---like
this.

A sentence-ending space should be larger than the
space between words within a sentence.  You
sometimes have to type special commands in
conjunction with punctuation characters to get
this right, as in the following sentence.
       Gnats, gnus, etc.\ all  % `\ ' makes an inter-word space.
       begin with G\@.         % \@ marks end-of-sentence punctuation.
You should check the spaces after periods when
reading your output to make sure you haven't
forgotten any special cases.  Generating an
ellipsis
       \ldots\               % `\ ' is needed after `\ldots' because TeX
                             % ignores spaces after command names like \ldots
                             % made from \ + letters.
                             %
                             % Note how a `%' character causes TeX to ignore
                             % the end of the input line, so these blank lines
                             % do not start a new paragraph.
                             %
with the right spacing around the periods requires
a special command.

\LaTeX\ interprets some common characters as
commands, so you must type special commands to
generate them.  These characters include the
following:
       \$ \& \% \# \{ and \}.

In printing, text is usually emphasized with an
       \emph{italic}
type style.

\begin{em}
   A long segment of text can also be emphasized
   in this way.  Text within such a segment can be
   given \emph{additional} emphasis.
\end{em}

It is sometimes necessary to prevent \LaTeX\ from
breaking a line where it might otherwise do so.
This may be at a space, as between the ``Mr.''\ and
``Jones'' in
       ``Mr.~Jones'',        % ~ produces an unbreakable interword space.
or within a word---especially when the word is a
symbol like
       \mbox{\emph{itemnum}}
that makes little sense when hyphenated across
lines.

Footnotes\footnote{This is an example of a footnote.}
pose no problem.

\LaTeX\ is good at typesetting mathematical formulas
like
       \( x-3y + z = 7 \)
or
       \( a_{1} > x^{2n} + y^{2n} > x' \)
or
       \( \ip{A}{B} = \sum_{i} a_{i} b_{i} \).
The spaces you type in a formula are
ignored.  Remember that a letter like
       $x$                   % $ ... $  and  \( ... \)  are equivalent
is a formula when it denotes a mathematical
symbol, and it should be typed as one.

\section{Displayed Text}

Text is displayed by indenting it from the left
margin.  Quotations are commonly displayed.  There
are short quotations
\begin{quote}
   This is a short quotation.  It consists of a
   single paragraph of text.  See how it is formatted.
\end{quote}
and longer ones.
\begin{quotation}
   This is a longer quotation.  It consists of two
   paragraphs of text, neither of which are
   particularly interesting.

   This is the second paragraph of the quotation.  It
   is just as dull as the first paragraph.
\end{quotation}
Another frequently-displayed structure is a list.
The following is an example of an \emph{itemized}
list.
\begin{itemize}
   \item This is the first item of an itemized list.
         Each item in the list is marked with a ``tick''.
         You don't have to worry about what kind of tick
         mark is used.

   \item This is the second item of the list.  It
         contains another list nested inside it.  The inner
         list is an \emph{enumerated} list.
         \begin{enumerate}
            \item This is the first item of an enumerated
                  list that is nested within the itemized list.

            \item This is the second item of the inner list.
                  \LaTeX\ allows you to nest lists deeper than
                  you really should.
         \end{enumerate}
         This is the rest of the second item of the outer
         list.  It is no more interesting than any other
         part of the item.
   \item This is the third item of the list.
\end{itemize}
You can even display poetry.
\begin{verse}
   There is an environment
    for verse \\             % The \\ command separates lines
   Whose features some poets % within a stanza.
   will curse.

                             % One or more blank lines separate stanzas.

   For instead of making\\
   Them do \emph{all} line breaking, \\
   It allows them to put too many words on a line when they'd rather be
   forced to be terse.
\end{verse}

Mathematical formulas may also be displayed.  A
displayed formula
is
one-line long; multiline
formulas require special formatting instructions.
   \[  \ip{\Gamma}{\psi'} = x'' + y^{2} + z_{i}^{n}\]
Don't start a paragraph with a displayed equation,
nor make one a paragraph by itself.

\end{document}               % End of document.

\section{Bench Huge Section 791}
% synthetic bench section 791
% This is a sample LaTeX input file.  (Version of 12 August 2004.)
%
% A '%' character causes TeX to ignore all remaining text on the line,
% and is used for comments like this one.

\documentclass{article}      % Specifies the document class

                             % The preamble begins here.
\title{An Example Document}  % Declares the document's title.
\author{Leslie Lamport}      % Declares the author's name.
\date{January 21, 1994}      % Deleting this command produces today's date.

\newcommand{\ip}[2]{(#1, #2)}
                             % Defines \ip{arg1}{arg2} to mean
                             % (arg1, arg2).

%\newcommand{\ip}[2]{\langle #1 | #2\rangle}
                             % This is an alternative definition of
                             % \ip that is commented out.

\begin{document}             % End of preamble and beginning of text.

\maketitle                   % Produces the title.

This is an example input file.  Comparing it with
the output it generates can show you how to
produce a simple document of your own.

\section{Ordinary Text}      % Produces section heading.  Lower-level
                             % sections are begun with similar
                             % \subsection and \subsubsection commands.

The ends  of words and sentences are marked
  by   spaces. It  doesn't matter how many
spaces    you type; one is as good as 100.  The
end of   a line counts as a space.

One   or more   blank lines denote the  end
of  a paragraph.

Since any number of consecutive spaces are treated
like a single one, the formatting of the input
file makes no difference to
      \LaTeX,                % The \LaTeX command generates the LaTeX logo.
but it makes a difference to you.  When you use
\LaTeX, making your input file as easy to read
as possible will be a great help as you write
your document and when you change it.  This sample
file shows how you can add comments to your own input
file.

Because printing is different from typewriting,
there are a number of things that you have to do
differently when preparing an input file than if
you were just typing the document directly.
Quotation marks like
       ``this''
have to be handled specially, as do quotes within
quotes:
       ``\,`this'            % \, separates the double and single quote.
        is what I just
        wrote, not  `that'\,''.

Dashes come in three sizes: an
       intra-word
dash, a medium dash for number ranges like
       1--2,
and a punctuation
       dash---like
this.

A sentence-ending space should be larger than the
space between words within a sentence.  You
sometimes have to type special commands in
conjunction with punctuation characters to get
this right, as in the following sentence.
       Gnats, gnus, etc.\ all  % `\ ' makes an inter-word space.
       begin with G\@.         % \@ marks end-of-sentence punctuation.
You should check the spaces after periods when
reading your output to make sure you haven't
forgotten any special cases.  Generating an
ellipsis
       \ldots\               % `\ ' is needed after `\ldots' because TeX
                             % ignores spaces after command names like \ldots
                             % made from \ + letters.
                             %
                             % Note how a `%' character causes TeX to ignore
                             % the end of the input line, so these blank lines
                             % do not start a new paragraph.
                             %
with the right spacing around the periods requires
a special command.

\LaTeX\ interprets some common characters as
commands, so you must type special commands to
generate them.  These characters include the
following:
       \$ \& \% \# \{ and \}.

In printing, text is usually emphasized with an
       \emph{italic}
type style.

\begin{em}
   A long segment of text can also be emphasized
   in this way.  Text within such a segment can be
   given \emph{additional} emphasis.
\end{em}

It is sometimes necessary to prevent \LaTeX\ from
breaking a line where it might otherwise do so.
This may be at a space, as between the ``Mr.''\ and
``Jones'' in
       ``Mr.~Jones'',        % ~ produces an unbreakable interword space.
or within a word---especially when the word is a
symbol like
       \mbox{\emph{itemnum}}
that makes little sense when hyphenated across
lines.

Footnotes\footnote{This is an example of a footnote.}
pose no problem.

\LaTeX\ is good at typesetting mathematical formulas
like
       \( x-3y + z = 7 \)
or
       \( a_{1} > x^{2n} + y^{2n} > x' \)
or
       \( \ip{A}{B} = \sum_{i} a_{i} b_{i} \).
The spaces you type in a formula are
ignored.  Remember that a letter like
       $x$                   % $ ... $  and  \( ... \)  are equivalent
is a formula when it denotes a mathematical
symbol, and it should be typed as one.

\section{Displayed Text}

Text is displayed by indenting it from the left
margin.  Quotations are commonly displayed.  There
are short quotations
\begin{quote}
   This is a short quotation.  It consists of a
   single paragraph of text.  See how it is formatted.
\end{quote}
and longer ones.
\begin{quotation}
   This is a longer quotation.  It consists of two
   paragraphs of text, neither of which are
   particularly interesting.

   This is the second paragraph of the quotation.  It
   is just as dull as the first paragraph.
\end{quotation}
Another frequently-displayed structure is a list.
The following is an example of an \emph{itemized}
list.
\begin{itemize}
   \item This is the first item of an itemized list.
         Each item in the list is marked with a ``tick''.
         You don't have to worry about what kind of tick
         mark is used.

   \item This is the second item of the list.  It
         contains another list nested inside it.  The inner
         list is an \emph{enumerated} list.
         \begin{enumerate}
            \item This is the first item of an enumerated
                  list that is nested within the itemized list.

            \item This is the second item of the inner list.
                  \LaTeX\ allows you to nest lists deeper than
                  you really should.
         \end{enumerate}
         This is the rest of the second item of the outer
         list.  It is no more interesting than any other
         part of the item.
   \item This is the third item of the list.
\end{itemize}
You can even display poetry.
\begin{verse}
   There is an environment
    for verse \\             % The \\ command separates lines
   Whose features some poets % within a stanza.
   will curse.

                             % One or more blank lines separate stanzas.

   For instead of making\\
   Them do \emph{all} line breaking, \\
   It allows them to put too many words on a line when they'd rather be
   forced to be terse.
\end{verse}

Mathematical formulas may also be displayed.  A
displayed formula
is
one-line long; multiline
formulas require special formatting instructions.
   \[  \ip{\Gamma}{\psi'} = x'' + y^{2} + z_{i}^{n}\]
Don't start a paragraph with a displayed equation,
nor make one a paragraph by itself.

\end{document}               % End of document.

\section{Bench Huge Section 792}
% synthetic bench section 792
% This is a sample LaTeX input file.  (Version of 12 August 2004.)
%
% A '%' character causes TeX to ignore all remaining text on the line,
% and is used for comments like this one.

\documentclass{article}      % Specifies the document class

                             % The preamble begins here.
\title{An Example Document}  % Declares the document's title.
\author{Leslie Lamport}      % Declares the author's name.
\date{January 21, 1994}      % Deleting this command produces today's date.

\newcommand{\ip}[2]{(#1, #2)}
                             % Defines \ip{arg1}{arg2} to mean
                             % (arg1, arg2).

%\newcommand{\ip}[2]{\langle #1 | #2\rangle}
                             % This is an alternative definition of
                             % \ip that is commented out.

\begin{document}             % End of preamble and beginning of text.

\maketitle                   % Produces the title.

This is an example input file.  Comparing it with
the output it generates can show you how to
produce a simple document of your own.

\section{Ordinary Text}      % Produces section heading.  Lower-level
                             % sections are begun with similar
                             % \subsection and \subsubsection commands.

The ends  of words and sentences are marked
  by   spaces. It  doesn't matter how many
spaces    you type; one is as good as 100.  The
end of   a line counts as a space.

One   or more   blank lines denote the  end
of  a paragraph.

Since any number of consecutive spaces are treated
like a single one, the formatting of the input
file makes no difference to
      \LaTeX,                % The \LaTeX command generates the LaTeX logo.
but it makes a difference to you.  When you use
\LaTeX, making your input file as easy to read
as possible will be a great help as you write
your document and when you change it.  This sample
file shows how you can add comments to your own input
file.

Because printing is different from typewriting,
there are a number of things that you have to do
differently when preparing an input file than if
you were just typing the document directly.
Quotation marks like
       ``this''
have to be handled specially, as do quotes within
quotes:
       ``\,`this'            % \, separates the double and single quote.
        is what I just
        wrote, not  `that'\,''.

Dashes come in three sizes: an
       intra-word
dash, a medium dash for number ranges like
       1--2,
and a punctuation
       dash---like
this.

A sentence-ending space should be larger than the
space between words within a sentence.  You
sometimes have to type special commands in
conjunction with punctuation characters to get
this right, as in the following sentence.
       Gnats, gnus, etc.\ all  % `\ ' makes an inter-word space.
       begin with G\@.         % \@ marks end-of-sentence punctuation.
You should check the spaces after periods when
reading your output to make sure you haven't
forgotten any special cases.  Generating an
ellipsis
       \ldots\               % `\ ' is needed after `\ldots' because TeX
                             % ignores spaces after command names like \ldots
                             % made from \ + letters.
                             %
                             % Note how a `%' character causes TeX to ignore
                             % the end of the input line, so these blank lines
                             % do not start a new paragraph.
                             %
with the right spacing around the periods requires
a special command.

\LaTeX\ interprets some common characters as
commands, so you must type special commands to
generate them.  These characters include the
following:
       \$ \& \% \# \{ and \}.

In printing, text is usually emphasized with an
       \emph{italic}
type style.

\begin{em}
   A long segment of text can also be emphasized
   in this way.  Text within such a segment can be
   given \emph{additional} emphasis.
\end{em}

It is sometimes necessary to prevent \LaTeX\ from
breaking a line where it might otherwise do so.
This may be at a space, as between the ``Mr.''\ and
``Jones'' in
       ``Mr.~Jones'',        % ~ produces an unbreakable interword space.
or within a word---especially when the word is a
symbol like
       \mbox{\emph{itemnum}}
that makes little sense when hyphenated across
lines.

Footnotes\footnote{This is an example of a footnote.}
pose no problem.

\LaTeX\ is good at typesetting mathematical formulas
like
       \( x-3y + z = 7 \)
or
       \( a_{1} > x^{2n} + y^{2n} > x' \)
or
       \( \ip{A}{B} = \sum_{i} a_{i} b_{i} \).
The spaces you type in a formula are
ignored.  Remember that a letter like
       $x$                   % $ ... $  and  \( ... \)  are equivalent
is a formula when it denotes a mathematical
symbol, and it should be typed as one.

\section{Displayed Text}

Text is displayed by indenting it from the left
margin.  Quotations are commonly displayed.  There
are short quotations
\begin{quote}
   This is a short quotation.  It consists of a
   single paragraph of text.  See how it is formatted.
\end{quote}
and longer ones.
\begin{quotation}
   This is a longer quotation.  It consists of two
   paragraphs of text, neither of which are
   particularly interesting.

   This is the second paragraph of the quotation.  It
   is just as dull as the first paragraph.
\end{quotation}
Another frequently-displayed structure is a list.
The following is an example of an \emph{itemized}
list.
\begin{itemize}
   \item This is the first item of an itemized list.
         Each item in the list is marked with a ``tick''.
         You don't have to worry about what kind of tick
         mark is used.

   \item This is the second item of the list.  It
         contains another list nested inside it.  The inner
         list is an \emph{enumerated} list.
         \begin{enumerate}
            \item This is the first item of an enumerated
                  list that is nested within the itemized list.

            \item This is the second item of the inner list.
                  \LaTeX\ allows you to nest lists deeper than
                  you really should.
         \end{enumerate}
         This is the rest of the second item of the outer
         list.  It is no more interesting than any other
         part of the item.
   \item This is the third item of the list.
\end{itemize}
You can even display poetry.
\begin{verse}
   There is an environment
    for verse \\             % The \\ command separates lines
   Whose features some poets % within a stanza.
   will curse.

                             % One or more blank lines separate stanzas.

   For instead of making\\
   Them do \emph{all} line breaking, \\
   It allows them to put too many words on a line when they'd rather be
   forced to be terse.
\end{verse}

Mathematical formulas may also be displayed.  A
displayed formula
is
one-line long; multiline
formulas require special formatting instructions.
   \[  \ip{\Gamma}{\psi'} = x'' + y^{2} + z_{i}^{n}\]
Don't start a paragraph with a displayed equation,
nor make one a paragraph by itself.

\end{document}               % End of document.

\section{Bench Huge Section 793}
% synthetic bench section 793
% This is a sample LaTeX input file.  (Version of 12 August 2004.)
%
% A '%' character causes TeX to ignore all remaining text on the line,
% and is used for comments like this one.

\documentclass{article}      % Specifies the document class

                             % The preamble begins here.
\title{An Example Document}  % Declares the document's title.
\author{Leslie Lamport}      % Declares the author's name.
\date{January 21, 1994}      % Deleting this command produces today's date.

\newcommand{\ip}[2]{(#1, #2)}
                             % Defines \ip{arg1}{arg2} to mean
                             % (arg1, arg2).

%\newcommand{\ip}[2]{\langle #1 | #2\rangle}
                             % This is an alternative definition of
                             % \ip that is commented out.

\begin{document}             % End of preamble and beginning of text.

\maketitle                   % Produces the title.

This is an example input file.  Comparing it with
the output it generates can show you how to
produce a simple document of your own.

\section{Ordinary Text}      % Produces section heading.  Lower-level
                             % sections are begun with similar
                             % \subsection and \subsubsection commands.

The ends  of words and sentences are marked
  by   spaces. It  doesn't matter how many
spaces    you type; one is as good as 100.  The
end of   a line counts as a space.

One   or more   blank lines denote the  end
of  a paragraph.

Since any number of consecutive spaces are treated
like a single one, the formatting of the input
file makes no difference to
      \LaTeX,                % The \LaTeX command generates the LaTeX logo.
but it makes a difference to you.  When you use
\LaTeX, making your input file as easy to read
as possible will be a great help as you write
your document and when you change it.  This sample
file shows how you can add comments to your own input
file.

Because printing is different from typewriting,
there are a number of things that you have to do
differently when preparing an input file than if
you were just typing the document directly.
Quotation marks like
       ``this''
have to be handled specially, as do quotes within
quotes:
       ``\,`this'            % \, separates the double and single quote.
        is what I just
        wrote, not  `that'\,''.

Dashes come in three sizes: an
       intra-word
dash, a medium dash for number ranges like
       1--2,
and a punctuation
       dash---like
this.

A sentence-ending space should be larger than the
space between words within a sentence.  You
sometimes have to type special commands in
conjunction with punctuation characters to get
this right, as in the following sentence.
       Gnats, gnus, etc.\ all  % `\ ' makes an inter-word space.
       begin with G\@.         % \@ marks end-of-sentence punctuation.
You should check the spaces after periods when
reading your output to make sure you haven't
forgotten any special cases.  Generating an
ellipsis
       \ldots\               % `\ ' is needed after `\ldots' because TeX
                             % ignores spaces after command names like \ldots
                             % made from \ + letters.
                             %
                             % Note how a `%' character causes TeX to ignore
                             % the end of the input line, so these blank lines
                             % do not start a new paragraph.
                             %
with the right spacing around the periods requires
a special command.

\LaTeX\ interprets some common characters as
commands, so you must type special commands to
generate them.  These characters include the
following:
       \$ \& \% \# \{ and \}.

In printing, text is usually emphasized with an
       \emph{italic}
type style.

\begin{em}
   A long segment of text can also be emphasized
   in this way.  Text within such a segment can be
   given \emph{additional} emphasis.
\end{em}

It is sometimes necessary to prevent \LaTeX\ from
breaking a line where it might otherwise do so.
This may be at a space, as between the ``Mr.''\ and
``Jones'' in
       ``Mr.~Jones'',        % ~ produces an unbreakable interword space.
or within a word---especially when the word is a
symbol like
       \mbox{\emph{itemnum}}
that makes little sense when hyphenated across
lines.

Footnotes\footnote{This is an example of a footnote.}
pose no problem.

\LaTeX\ is good at typesetting mathematical formulas
like
       \( x-3y + z = 7 \)
or
       \( a_{1} > x^{2n} + y^{2n} > x' \)
or
       \( \ip{A}{B} = \sum_{i} a_{i} b_{i} \).
The spaces you type in a formula are
ignored.  Remember that a letter like
       $x$                   % $ ... $  and  \( ... \)  are equivalent
is a formula when it denotes a mathematical
symbol, and it should be typed as one.

\section{Displayed Text}

Text is displayed by indenting it from the left
margin.  Quotations are commonly displayed.  There
are short quotations
\begin{quote}
   This is a short quotation.  It consists of a
   single paragraph of text.  See how it is formatted.
\end{quote}
and longer ones.
\begin{quotation}
   This is a longer quotation.  It consists of two
   paragraphs of text, neither of which are
   particularly interesting.

   This is the second paragraph of the quotation.  It
   is just as dull as the first paragraph.
\end{quotation}
Another frequently-displayed structure is a list.
The following is an example of an \emph{itemized}
list.
\begin{itemize}
   \item This is the first item of an itemized list.
         Each item in the list is marked with a ``tick''.
         You don't have to worry about what kind of tick
         mark is used.

   \item This is the second item of the list.  It
         contains another list nested inside it.  The inner
         list is an \emph{enumerated} list.
         \begin{enumerate}
            \item This is the first item of an enumerated
                  list that is nested within the itemized list.

            \item This is the second item of the inner list.
                  \LaTeX\ allows you to nest lists deeper than
                  you really should.
         \end{enumerate}
         This is the rest of the second item of the outer
         list.  It is no more interesting than any other
         part of the item.
   \item This is the third item of the list.
\end{itemize}
You can even display poetry.
\begin{verse}
   There is an environment
    for verse \\             % The \\ command separates lines
   Whose features some poets % within a stanza.
   will curse.

                             % One or more blank lines separate stanzas.

   For instead of making\\
   Them do \emph{all} line breaking, \\
   It allows them to put too many words on a line when they'd rather be
   forced to be terse.
\end{verse}

Mathematical formulas may also be displayed.  A
displayed formula
is
one-line long; multiline
formulas require special formatting instructions.
   \[  \ip{\Gamma}{\psi'} = x'' + y^{2} + z_{i}^{n}\]
Don't start a paragraph with a displayed equation,
nor make one a paragraph by itself.

\end{document}               % End of document.

\section{Bench Huge Section 794}
% synthetic bench section 794
% This is a sample LaTeX input file.  (Version of 12 August 2004.)
%
% A '%' character causes TeX to ignore all remaining text on the line,
% and is used for comments like this one.

\documentclass{article}      % Specifies the document class

                             % The preamble begins here.
\title{An Example Document}  % Declares the document's title.
\author{Leslie Lamport}      % Declares the author's name.
\date{January 21, 1994}      % Deleting this command produces today's date.

\newcommand{\ip}[2]{(#1, #2)}
                             % Defines \ip{arg1}{arg2} to mean
                             % (arg1, arg2).

%\newcommand{\ip}[2]{\langle #1 | #2\rangle}
                             % This is an alternative definition of
                             % \ip that is commented out.

\begin{document}             % End of preamble and beginning of text.

\maketitle                   % Produces the title.

This is an example input file.  Comparing it with
the output it generates can show you how to
produce a simple document of your own.

\section{Ordinary Text}      % Produces section heading.  Lower-level
                             % sections are begun with similar
                             % \subsection and \subsubsection commands.

The ends  of words and sentences are marked
  by   spaces. It  doesn't matter how many
spaces    you type; one is as good as 100.  The
end of   a line counts as a space.

One   or more   blank lines denote the  end
of  a paragraph.

Since any number of consecutive spaces are treated
like a single one, the formatting of the input
file makes no difference to
      \LaTeX,                % The \LaTeX command generates the LaTeX logo.
but it makes a difference to you.  When you use
\LaTeX, making your input file as easy to read
as possible will be a great help as you write
your document and when you change it.  This sample
file shows how you can add comments to your own input
file.

Because printing is different from typewriting,
there are a number of things that you have to do
differently when preparing an input file than if
you were just typing the document directly.
Quotation marks like
       ``this''
have to be handled specially, as do quotes within
quotes:
       ``\,`this'            % \, separates the double and single quote.
        is what I just
        wrote, not  `that'\,''.

Dashes come in three sizes: an
       intra-word
dash, a medium dash for number ranges like
       1--2,
and a punctuation
       dash---like
this.

A sentence-ending space should be larger than the
space between words within a sentence.  You
sometimes have to type special commands in
conjunction with punctuation characters to get
this right, as in the following sentence.
       Gnats, gnus, etc.\ all  % `\ ' makes an inter-word space.
       begin with G\@.         % \@ marks end-of-sentence punctuation.
You should check the spaces after periods when
reading your output to make sure you haven't
forgotten any special cases.  Generating an
ellipsis
       \ldots\               % `\ ' is needed after `\ldots' because TeX
                             % ignores spaces after command names like \ldots
                             % made from \ + letters.
                             %
                             % Note how a `%' character causes TeX to ignore
                             % the end of the input line, so these blank lines
                             % do not start a new paragraph.
                             %
with the right spacing around the periods requires
a special command.

\LaTeX\ interprets some common characters as
commands, so you must type special commands to
generate them.  These characters include the
following:
       \$ \& \% \# \{ and \}.

In printing, text is usually emphasized with an
       \emph{italic}
type style.

\begin{em}
   A long segment of text can also be emphasized
   in this way.  Text within such a segment can be
   given \emph{additional} emphasis.
\end{em}

It is sometimes necessary to prevent \LaTeX\ from
breaking a line where it might otherwise do so.
This may be at a space, as between the ``Mr.''\ and
``Jones'' in
       ``Mr.~Jones'',        % ~ produces an unbreakable interword space.
or within a word---especially when the word is a
symbol like
       \mbox{\emph{itemnum}}
that makes little sense when hyphenated across
lines.

Footnotes\footnote{This is an example of a footnote.}
pose no problem.

\LaTeX\ is good at typesetting mathematical formulas
like
       \( x-3y + z = 7 \)
or
       \( a_{1} > x^{2n} + y^{2n} > x' \)
or
       \( \ip{A}{B} = \sum_{i} a_{i} b_{i} \).
The spaces you type in a formula are
ignored.  Remember that a letter like
       $x$                   % $ ... $  and  \( ... \)  are equivalent
is a formula when it denotes a mathematical
symbol, and it should be typed as one.

\section{Displayed Text}

Text is displayed by indenting it from the left
margin.  Quotations are commonly displayed.  There
are short quotations
\begin{quote}
   This is a short quotation.  It consists of a
   single paragraph of text.  See how it is formatted.
\end{quote}
and longer ones.
\begin{quotation}
   This is a longer quotation.  It consists of two
   paragraphs of text, neither of which are
   particularly interesting.

   This is the second paragraph of the quotation.  It
   is just as dull as the first paragraph.
\end{quotation}
Another frequently-displayed structure is a list.
The following is an example of an \emph{itemized}
list.
\begin{itemize}
   \item This is the first item of an itemized list.
         Each item in the list is marked with a ``tick''.
         You don't have to worry about what kind of tick
         mark is used.

   \item This is the second item of the list.  It
         contains another list nested inside it.  The inner
         list is an \emph{enumerated} list.
         \begin{enumerate}
            \item This is the first item of an enumerated
                  list that is nested within the itemized list.

            \item This is the second item of the inner list.
                  \LaTeX\ allows you to nest lists deeper than
                  you really should.
         \end{enumerate}
         This is the rest of the second item of the outer
         list.  It is no more interesting than any other
         part of the item.
   \item This is the third item of the list.
\end{itemize}
You can even display poetry.
\begin{verse}
   There is an environment
    for verse \\             % The \\ command separates lines
   Whose features some poets % within a stanza.
   will curse.

                             % One or more blank lines separate stanzas.

   For instead of making\\
   Them do \emph{all} line breaking, \\
   It allows them to put too many words on a line when they'd rather be
   forced to be terse.
\end{verse}

Mathematical formulas may also be displayed.  A
displayed formula
is
one-line long; multiline
formulas require special formatting instructions.
   \[  \ip{\Gamma}{\psi'} = x'' + y^{2} + z_{i}^{n}\]
Don't start a paragraph with a displayed equation,
nor make one a paragraph by itself.

\end{document}               % End of document.

\section{Bench Huge Section 795}
% synthetic bench section 795
% This is a sample LaTeX input file.  (Version of 12 August 2004.)
%
% A '%' character causes TeX to ignore all remaining text on the line,
% and is used for comments like this one.

\documentclass{article}      % Specifies the document class

                             % The preamble begins here.
\title{An Example Document}  % Declares the document's title.
\author{Leslie Lamport}      % Declares the author's name.
\date{January 21, 1994}      % Deleting this command produces today's date.

\newcommand{\ip}[2]{(#1, #2)}
                             % Defines \ip{arg1}{arg2} to mean
                             % (arg1, arg2).

%\newcommand{\ip}[2]{\langle #1 | #2\rangle}
                             % This is an alternative definition of
                             % \ip that is commented out.

\begin{document}             % End of preamble and beginning of text.

\maketitle                   % Produces the title.

This is an example input file.  Comparing it with
the output it generates can show you how to
produce a simple document of your own.

\section{Ordinary Text}      % Produces section heading.  Lower-level
                             % sections are begun with similar
                             % \subsection and \subsubsection commands.

The ends  of words and sentences are marked
  by   spaces. It  doesn't matter how many
spaces    you type; one is as good as 100.  The
end of   a line counts as a space.

One   or more   blank lines denote the  end
of  a paragraph.

Since any number of consecutive spaces are treated
like a single one, the formatting of the input
file makes no difference to
      \LaTeX,                % The \LaTeX command generates the LaTeX logo.
but it makes a difference to you.  When you use
\LaTeX, making your input file as easy to read
as possible will be a great help as you write
your document and when you change it.  This sample
file shows how you can add comments to your own input
file.

Because printing is different from typewriting,
there are a number of things that you have to do
differently when preparing an input file than if
you were just typing the document directly.
Quotation marks like
       ``this''
have to be handled specially, as do quotes within
quotes:
       ``\,`this'            % \, separates the double and single quote.
        is what I just
        wrote, not  `that'\,''.

Dashes come in three sizes: an
       intra-word
dash, a medium dash for number ranges like
       1--2,
and a punctuation
       dash---like
this.

A sentence-ending space should be larger than the
space between words within a sentence.  You
sometimes have to type special commands in
conjunction with punctuation characters to get
this right, as in the following sentence.
       Gnats, gnus, etc.\ all  % `\ ' makes an inter-word space.
       begin with G\@.         % \@ marks end-of-sentence punctuation.
You should check the spaces after periods when
reading your output to make sure you haven't
forgotten any special cases.  Generating an
ellipsis
       \ldots\               % `\ ' is needed after `\ldots' because TeX
                             % ignores spaces after command names like \ldots
                             % made from \ + letters.
                             %
                             % Note how a `%' character causes TeX to ignore
                             % the end of the input line, so these blank lines
                             % do not start a new paragraph.
                             %
with the right spacing around the periods requires
a special command.

\LaTeX\ interprets some common characters as
commands, so you must type special commands to
generate them.  These characters include the
following:
       \$ \& \% \# \{ and \}.

In printing, text is usually emphasized with an
       \emph{italic}
type style.

\begin{em}
   A long segment of text can also be emphasized
   in this way.  Text within such a segment can be
   given \emph{additional} emphasis.
\end{em}

It is sometimes necessary to prevent \LaTeX\ from
breaking a line where it might otherwise do so.
This may be at a space, as between the ``Mr.''\ and
``Jones'' in
       ``Mr.~Jones'',        % ~ produces an unbreakable interword space.
or within a word---especially when the word is a
symbol like
       \mbox{\emph{itemnum}}
that makes little sense when hyphenated across
lines.

Footnotes\footnote{This is an example of a footnote.}
pose no problem.

\LaTeX\ is good at typesetting mathematical formulas
like
       \( x-3y + z = 7 \)
or
       \( a_{1} > x^{2n} + y^{2n} > x' \)
or
       \( \ip{A}{B} = \sum_{i} a_{i} b_{i} \).
The spaces you type in a formula are
ignored.  Remember that a letter like
       $x$                   % $ ... $  and  \( ... \)  are equivalent
is a formula when it denotes a mathematical
symbol, and it should be typed as one.

\section{Displayed Text}

Text is displayed by indenting it from the left
margin.  Quotations are commonly displayed.  There
are short quotations
\begin{quote}
   This is a short quotation.  It consists of a
   single paragraph of text.  See how it is formatted.
\end{quote}
and longer ones.
\begin{quotation}
   This is a longer quotation.  It consists of two
   paragraphs of text, neither of which are
   particularly interesting.

   This is the second paragraph of the quotation.  It
   is just as dull as the first paragraph.
\end{quotation}
Another frequently-displayed structure is a list.
The following is an example of an \emph{itemized}
list.
\begin{itemize}
   \item This is the first item of an itemized list.
         Each item in the list is marked with a ``tick''.
         You don't have to worry about what kind of tick
         mark is used.

   \item This is the second item of the list.  It
         contains another list nested inside it.  The inner
         list is an \emph{enumerated} list.
         \begin{enumerate}
            \item This is the first item of an enumerated
                  list that is nested within the itemized list.

            \item This is the second item of the inner list.
                  \LaTeX\ allows you to nest lists deeper than
                  you really should.
         \end{enumerate}
         This is the rest of the second item of the outer
         list.  It is no more interesting than any other
         part of the item.
   \item This is the third item of the list.
\end{itemize}
You can even display poetry.
\begin{verse}
   There is an environment
    for verse \\             % The \\ command separates lines
   Whose features some poets % within a stanza.
   will curse.

                             % One or more blank lines separate stanzas.

   For instead of making\\
   Them do \emph{all} line breaking, \\
   It allows them to put too many words on a line when they'd rather be
   forced to be terse.
\end{verse}

Mathematical formulas may also be displayed.  A
displayed formula
is
one-line long; multiline
formulas require special formatting instructions.
   \[  \ip{\Gamma}{\psi'} = x'' + y^{2} + z_{i}^{n}\]
Don't start a paragraph with a displayed equation,
nor make one a paragraph by itself.

\end{document}               % End of document.

\section{Bench Huge Section 796}
% synthetic bench section 796
% This is a sample LaTeX input file.  (Version of 12 August 2004.)
%
% A '%' character causes TeX to ignore all remaining text on the line,
% and is used for comments like this one.

\documentclass{article}      % Specifies the document class

                             % The preamble begins here.
\title{An Example Document}  % Declares the document's title.
\author{Leslie Lamport}      % Declares the author's name.
\date{January 21, 1994}      % Deleting this command produces today's date.

\newcommand{\ip}[2]{(#1, #2)}
                             % Defines \ip{arg1}{arg2} to mean
                             % (arg1, arg2).

%\newcommand{\ip}[2]{\langle #1 | #2\rangle}
                             % This is an alternative definition of
                             % \ip that is commented out.

\begin{document}             % End of preamble and beginning of text.

\maketitle                   % Produces the title.

This is an example input file.  Comparing it with
the output it generates can show you how to
produce a simple document of your own.

\section{Ordinary Text}      % Produces section heading.  Lower-level
                             % sections are begun with similar
                             % \subsection and \subsubsection commands.

The ends  of words and sentences are marked
  by   spaces. It  doesn't matter how many
spaces    you type; one is as good as 100.  The
end of   a line counts as a space.

One   or more   blank lines denote the  end
of  a paragraph.

Since any number of consecutive spaces are treated
like a single one, the formatting of the input
file makes no difference to
      \LaTeX,                % The \LaTeX command generates the LaTeX logo.
but it makes a difference to you.  When you use
\LaTeX, making your input file as easy to read
as possible will be a great help as you write
your document and when you change it.  This sample
file shows how you can add comments to your own input
file.

Because printing is different from typewriting,
there are a number of things that you have to do
differently when preparing an input file than if
you were just typing the document directly.
Quotation marks like
       ``this''
have to be handled specially, as do quotes within
quotes:
       ``\,`this'            % \, separates the double and single quote.
        is what I just
        wrote, not  `that'\,''.

Dashes come in three sizes: an
       intra-word
dash, a medium dash for number ranges like
       1--2,
and a punctuation
       dash---like
this.

A sentence-ending space should be larger than the
space between words within a sentence.  You
sometimes have to type special commands in
conjunction with punctuation characters to get
this right, as in the following sentence.
       Gnats, gnus, etc.\ all  % `\ ' makes an inter-word space.
       begin with G\@.         % \@ marks end-of-sentence punctuation.
You should check the spaces after periods when
reading your output to make sure you haven't
forgotten any special cases.  Generating an
ellipsis
       \ldots\               % `\ ' is needed after `\ldots' because TeX
                             % ignores spaces after command names like \ldots
                             % made from \ + letters.
                             %
                             % Note how a `%' character causes TeX to ignore
                             % the end of the input line, so these blank lines
                             % do not start a new paragraph.
                             %
with the right spacing around the periods requires
a special command.

\LaTeX\ interprets some common characters as
commands, so you must type special commands to
generate them.  These characters include the
following:
       \$ \& \% \# \{ and \}.

In printing, text is usually emphasized with an
       \emph{italic}
type style.

\begin{em}
   A long segment of text can also be emphasized
   in this way.  Text within such a segment can be
   given \emph{additional} emphasis.
\end{em}

It is sometimes necessary to prevent \LaTeX\ from
breaking a line where it might otherwise do so.
This may be at a space, as between the ``Mr.''\ and
``Jones'' in
       ``Mr.~Jones'',        % ~ produces an unbreakable interword space.
or within a word---especially when the word is a
symbol like
       \mbox{\emph{itemnum}}
that makes little sense when hyphenated across
lines.

Footnotes\footnote{This is an example of a footnote.}
pose no problem.

\LaTeX\ is good at typesetting mathematical formulas
like
       \( x-3y + z = 7 \)
or
       \( a_{1} > x^{2n} + y^{2n} > x' \)
or
       \( \ip{A}{B} = \sum_{i} a_{i} b_{i} \).
The spaces you type in a formula are
ignored.  Remember that a letter like
       $x$                   % $ ... $  and  \( ... \)  are equivalent
is a formula when it denotes a mathematical
symbol, and it should be typed as one.

\section{Displayed Text}

Text is displayed by indenting it from the left
margin.  Quotations are commonly displayed.  There
are short quotations
\begin{quote}
   This is a short quotation.  It consists of a
   single paragraph of text.  See how it is formatted.
\end{quote}
and longer ones.
\begin{quotation}
   This is a longer quotation.  It consists of two
   paragraphs of text, neither of which are
   particularly interesting.

   This is the second paragraph of the quotation.  It
   is just as dull as the first paragraph.
\end{quotation}
Another frequently-displayed structure is a list.
The following is an example of an \emph{itemized}
list.
\begin{itemize}
   \item This is the first item of an itemized list.
         Each item in the list is marked with a ``tick''.
         You don't have to worry about what kind of tick
         mark is used.

   \item This is the second item of the list.  It
         contains another list nested inside it.  The inner
         list is an \emph{enumerated} list.
         \begin{enumerate}
            \item This is the first item of an enumerated
                  list that is nested within the itemized list.

            \item This is the second item of the inner list.
                  \LaTeX\ allows you to nest lists deeper than
                  you really should.
         \end{enumerate}
         This is the rest of the second item of the outer
         list.  It is no more interesting than any other
         part of the item.
   \item This is the third item of the list.
\end{itemize}
You can even display poetry.
\begin{verse}
   There is an environment
    for verse \\             % The \\ command separates lines
   Whose features some poets % within a stanza.
   will curse.

                             % One or more blank lines separate stanzas.

   For instead of making\\
   Them do \emph{all} line breaking, \\
   It allows them to put too many words on a line when they'd rather be
   forced to be terse.
\end{verse}

Mathematical formulas may also be displayed.  A
displayed formula
is
one-line long; multiline
formulas require special formatting instructions.
   \[  \ip{\Gamma}{\psi'} = x'' + y^{2} + z_{i}^{n}\]
Don't start a paragraph with a displayed equation,
nor make one a paragraph by itself.

\end{document}               % End of document.

\section{Bench Huge Section 797}
% synthetic bench section 797
% This is a sample LaTeX input file.  (Version of 12 August 2004.)
%
% A '%' character causes TeX to ignore all remaining text on the line,
% and is used for comments like this one.

\documentclass{article}      % Specifies the document class

                             % The preamble begins here.
\title{An Example Document}  % Declares the document's title.
\author{Leslie Lamport}      % Declares the author's name.
\date{January 21, 1994}      % Deleting this command produces today's date.

\newcommand{\ip}[2]{(#1, #2)}
                             % Defines \ip{arg1}{arg2} to mean
                             % (arg1, arg2).

%\newcommand{\ip}[2]{\langle #1 | #2\rangle}
                             % This is an alternative definition of
                             % \ip that is commented out.

\begin{document}             % End of preamble and beginning of text.

\maketitle                   % Produces the title.

This is an example input file.  Comparing it with
the output it generates can show you how to
produce a simple document of your own.

\section{Ordinary Text}      % Produces section heading.  Lower-level
                             % sections are begun with similar
                             % \subsection and \subsubsection commands.

The ends  of words and sentences are marked
  by   spaces. It  doesn't matter how many
spaces    you type; one is as good as 100.  The
end of   a line counts as a space.

One   or more   blank lines denote the  end
of  a paragraph.

Since any number of consecutive spaces are treated
like a single one, the formatting of the input
file makes no difference to
      \LaTeX,                % The \LaTeX command generates the LaTeX logo.
but it makes a difference to you.  When you use
\LaTeX, making your input file as easy to read
as possible will be a great help as you write
your document and when you change it.  This sample
file shows how you can add comments to your own input
file.

Because printing is different from typewriting,
there are a number of things that you have to do
differently when preparing an input file than if
you were just typing the document directly.
Quotation marks like
       ``this''
have to be handled specially, as do quotes within
quotes:
       ``\,`this'            % \, separates the double and single quote.
        is what I just
        wrote, not  `that'\,''.

Dashes come in three sizes: an
       intra-word
dash, a medium dash for number ranges like
       1--2,
and a punctuation
       dash---like
this.

A sentence-ending space should be larger than the
space between words within a sentence.  You
sometimes have to type special commands in
conjunction with punctuation characters to get
this right, as in the following sentence.
       Gnats, gnus, etc.\ all  % `\ ' makes an inter-word space.
       begin with G\@.         % \@ marks end-of-sentence punctuation.
You should check the spaces after periods when
reading your output to make sure you haven't
forgotten any special cases.  Generating an
ellipsis
       \ldots\               % `\ ' is needed after `\ldots' because TeX
                             % ignores spaces after command names like \ldots
                             % made from \ + letters.
                             %
                             % Note how a `%' character causes TeX to ignore
                             % the end of the input line, so these blank lines
                             % do not start a new paragraph.
                             %
with the right spacing around the periods requires
a special command.

\LaTeX\ interprets some common characters as
commands, so you must type special commands to
generate them.  These characters include the
following:
       \$ \& \% \# \{ and \}.

In printing, text is usually emphasized with an
       \emph{italic}
type style.

\begin{em}
   A long segment of text can also be emphasized
   in this way.  Text within such a segment can be
   given \emph{additional} emphasis.
\end{em}

It is sometimes necessary to prevent \LaTeX\ from
breaking a line where it might otherwise do so.
This may be at a space, as between the ``Mr.''\ and
``Jones'' in
       ``Mr.~Jones'',        % ~ produces an unbreakable interword space.
or within a word---especially when the word is a
symbol like
       \mbox{\emph{itemnum}}
that makes little sense when hyphenated across
lines.

Footnotes\footnote{This is an example of a footnote.}
pose no problem.

\LaTeX\ is good at typesetting mathematical formulas
like
       \( x-3y + z = 7 \)
or
       \( a_{1} > x^{2n} + y^{2n} > x' \)
or
       \( \ip{A}{B} = \sum_{i} a_{i} b_{i} \).
The spaces you type in a formula are
ignored.  Remember that a letter like
       $x$                   % $ ... $  and  \( ... \)  are equivalent
is a formula when it denotes a mathematical
symbol, and it should be typed as one.

\section{Displayed Text}

Text is displayed by indenting it from the left
margin.  Quotations are commonly displayed.  There
are short quotations
\begin{quote}
   This is a short quotation.  It consists of a
   single paragraph of text.  See how it is formatted.
\end{quote}
and longer ones.
\begin{quotation}
   This is a longer quotation.  It consists of two
   paragraphs of text, neither of which are
   particularly interesting.

   This is the second paragraph of the quotation.  It
   is just as dull as the first paragraph.
\end{quotation}
Another frequently-displayed structure is a list.
The following is an example of an \emph{itemized}
list.
\begin{itemize}
   \item This is the first item of an itemized list.
         Each item in the list is marked with a ``tick''.
         You don't have to worry about what kind of tick
         mark is used.

   \item This is the second item of the list.  It
         contains another list nested inside it.  The inner
         list is an \emph{enumerated} list.
         \begin{enumerate}
            \item This is the first item of an enumerated
                  list that is nested within the itemized list.

            \item This is the second item of the inner list.
                  \LaTeX\ allows you to nest lists deeper than
                  you really should.
         \end{enumerate}
         This is the rest of the second item of the outer
         list.  It is no more interesting than any other
         part of the item.
   \item This is the third item of the list.
\end{itemize}
You can even display poetry.
\begin{verse}
   There is an environment
    for verse \\             % The \\ command separates lines
   Whose features some poets % within a stanza.
   will curse.

                             % One or more blank lines separate stanzas.

   For instead of making\\
   Them do \emph{all} line breaking, \\
   It allows them to put too many words on a line when they'd rather be
   forced to be terse.
\end{verse}

Mathematical formulas may also be displayed.  A
displayed formula
is
one-line long; multiline
formulas require special formatting instructions.
   \[  \ip{\Gamma}{\psi'} = x'' + y^{2} + z_{i}^{n}\]
Don't start a paragraph with a displayed equation,
nor make one a paragraph by itself.

\end{document}               % End of document.

\section{Bench Huge Section 798}
% synthetic bench section 798
% This is a sample LaTeX input file.  (Version of 12 August 2004.)
%
% A '%' character causes TeX to ignore all remaining text on the line,
% and is used for comments like this one.

\documentclass{article}      % Specifies the document class

                             % The preamble begins here.
\title{An Example Document}  % Declares the document's title.
\author{Leslie Lamport}      % Declares the author's name.
\date{January 21, 1994}      % Deleting this command produces today's date.

\newcommand{\ip}[2]{(#1, #2)}
                             % Defines \ip{arg1}{arg2} to mean
                             % (arg1, arg2).

%\newcommand{\ip}[2]{\langle #1 | #2\rangle}
                             % This is an alternative definition of
                             % \ip that is commented out.

\begin{document}             % End of preamble and beginning of text.

\maketitle                   % Produces the title.

This is an example input file.  Comparing it with
the output it generates can show you how to
produce a simple document of your own.

\section{Ordinary Text}      % Produces section heading.  Lower-level
                             % sections are begun with similar
                             % \subsection and \subsubsection commands.

The ends  of words and sentences are marked
  by   spaces. It  doesn't matter how many
spaces    you type; one is as good as 100.  The
end of   a line counts as a space.

One   or more   blank lines denote the  end
of  a paragraph.

Since any number of consecutive spaces are treated
like a single one, the formatting of the input
file makes no difference to
      \LaTeX,                % The \LaTeX command generates the LaTeX logo.
but it makes a difference to you.  When you use
\LaTeX, making your input file as easy to read
as possible will be a great help as you write
your document and when you change it.  This sample
file shows how you can add comments to your own input
file.

Because printing is different from typewriting,
there are a number of things that you have to do
differently when preparing an input file than if
you were just typing the document directly.
Quotation marks like
       ``this''
have to be handled specially, as do quotes within
quotes:
       ``\,`this'            % \, separates the double and single quote.
        is what I just
        wrote, not  `that'\,''.

Dashes come in three sizes: an
       intra-word
dash, a medium dash for number ranges like
       1--2,
and a punctuation
       dash---like
this.

A sentence-ending space should be larger than the
space between words within a sentence.  You
sometimes have to type special commands in
conjunction with punctuation characters to get
this right, as in the following sentence.
       Gnats, gnus, etc.\ all  % `\ ' makes an inter-word space.
       begin with G\@.         % \@ marks end-of-sentence punctuation.
You should check the spaces after periods when
reading your output to make sure you haven't
forgotten any special cases.  Generating an
ellipsis
       \ldots\               % `\ ' is needed after `\ldots' because TeX
                             % ignores spaces after command names like \ldots
                             % made from \ + letters.
                             %
                             % Note how a `%' character causes TeX to ignore
                             % the end of the input line, so these blank lines
                             % do not start a new paragraph.
                             %
with the right spacing around the periods requires
a special command.

\LaTeX\ interprets some common characters as
commands, so you must type special commands to
generate them.  These characters include the
following:
       \$ \& \% \# \{ and \}.

In printing, text is usually emphasized with an
       \emph{italic}
type style.

\begin{em}
   A long segment of text can also be emphasized
   in this way.  Text within such a segment can be
   given \emph{additional} emphasis.
\end{em}

It is sometimes necessary to prevent \LaTeX\ from
breaking a line where it might otherwise do so.
This may be at a space, as between the ``Mr.''\ and
``Jones'' in
       ``Mr.~Jones'',        % ~ produces an unbreakable interword space.
or within a word---especially when the word is a
symbol like
       \mbox{\emph{itemnum}}
that makes little sense when hyphenated across
lines.

Footnotes\footnote{This is an example of a footnote.}
pose no problem.

\LaTeX\ is good at typesetting mathematical formulas
like
       \( x-3y + z = 7 \)
or
       \( a_{1} > x^{2n} + y^{2n} > x' \)
or
       \( \ip{A}{B} = \sum_{i} a_{i} b_{i} \).
The spaces you type in a formula are
ignored.  Remember that a letter like
       $x$                   % $ ... $  and  \( ... \)  are equivalent
is a formula when it denotes a mathematical
symbol, and it should be typed as one.

\section{Displayed Text}

Text is displayed by indenting it from the left
margin.  Quotations are commonly displayed.  There
are short quotations
\begin{quote}
   This is a short quotation.  It consists of a
   single paragraph of text.  See how it is formatted.
\end{quote}
and longer ones.
\begin{quotation}
   This is a longer quotation.  It consists of two
   paragraphs of text, neither of which are
   particularly interesting.

   This is the second paragraph of the quotation.  It
   is just as dull as the first paragraph.
\end{quotation}
Another frequently-displayed structure is a list.
The following is an example of an \emph{itemized}
list.
\begin{itemize}
   \item This is the first item of an itemized list.
         Each item in the list is marked with a ``tick''.
         You don't have to worry about what kind of tick
         mark is used.

   \item This is the second item of the list.  It
         contains another list nested inside it.  The inner
         list is an \emph{enumerated} list.
         \begin{enumerate}
            \item This is the first item of an enumerated
                  list that is nested within the itemized list.

            \item This is the second item of the inner list.
                  \LaTeX\ allows you to nest lists deeper than
                  you really should.
         \end{enumerate}
         This is the rest of the second item of the outer
         list.  It is no more interesting than any other
         part of the item.
   \item This is the third item of the list.
\end{itemize}
You can even display poetry.
\begin{verse}
   There is an environment
    for verse \\             % The \\ command separates lines
   Whose features some poets % within a stanza.
   will curse.

                             % One or more blank lines separate stanzas.

   For instead of making\\
   Them do \emph{all} line breaking, \\
   It allows them to put too many words on a line when they'd rather be
   forced to be terse.
\end{verse}

Mathematical formulas may also be displayed.  A
displayed formula
is
one-line long; multiline
formulas require special formatting instructions.
   \[  \ip{\Gamma}{\psi'} = x'' + y^{2} + z_{i}^{n}\]
Don't start a paragraph with a displayed equation,
nor make one a paragraph by itself.

\end{document}               % End of document.

\section{Bench Huge Section 799}
% synthetic bench section 799
% This is a sample LaTeX input file.  (Version of 12 August 2004.)
%
% A '%' character causes TeX to ignore all remaining text on the line,
% and is used for comments like this one.

\documentclass{article}      % Specifies the document class

                             % The preamble begins here.
\title{An Example Document}  % Declares the document's title.
\author{Leslie Lamport}      % Declares the author's name.
\date{January 21, 1994}      % Deleting this command produces today's date.

\newcommand{\ip}[2]{(#1, #2)}
                             % Defines \ip{arg1}{arg2} to mean
                             % (arg1, arg2).

%\newcommand{\ip}[2]{\langle #1 | #2\rangle}
                             % This is an alternative definition of
                             % \ip that is commented out.

\begin{document}             % End of preamble and beginning of text.

\maketitle                   % Produces the title.

This is an example input file.  Comparing it with
the output it generates can show you how to
produce a simple document of your own.

\section{Ordinary Text}      % Produces section heading.  Lower-level
                             % sections are begun with similar
                             % \subsection and \subsubsection commands.

The ends  of words and sentences are marked
  by   spaces. It  doesn't matter how many
spaces    you type; one is as good as 100.  The
end of   a line counts as a space.

One   or more   blank lines denote the  end
of  a paragraph.

Since any number of consecutive spaces are treated
like a single one, the formatting of the input
file makes no difference to
      \LaTeX,                % The \LaTeX command generates the LaTeX logo.
but it makes a difference to you.  When you use
\LaTeX, making your input file as easy to read
as possible will be a great help as you write
your document and when you change it.  This sample
file shows how you can add comments to your own input
file.

Because printing is different from typewriting,
there are a number of things that you have to do
differently when preparing an input file than if
you were just typing the document directly.
Quotation marks like
       ``this''
have to be handled specially, as do quotes within
quotes:
       ``\,`this'            % \, separates the double and single quote.
        is what I just
        wrote, not  `that'\,''.

Dashes come in three sizes: an
       intra-word
dash, a medium dash for number ranges like
       1--2,
and a punctuation
       dash---like
this.

A sentence-ending space should be larger than the
space between words within a sentence.  You
sometimes have to type special commands in
conjunction with punctuation characters to get
this right, as in the following sentence.
       Gnats, gnus, etc.\ all  % `\ ' makes an inter-word space.
       begin with G\@.         % \@ marks end-of-sentence punctuation.
You should check the spaces after periods when
reading your output to make sure you haven't
forgotten any special cases.  Generating an
ellipsis
       \ldots\               % `\ ' is needed after `\ldots' because TeX
                             % ignores spaces after command names like \ldots
                             % made from \ + letters.
                             %
                             % Note how a `%' character causes TeX to ignore
                             % the end of the input line, so these blank lines
                             % do not start a new paragraph.
                             %
with the right spacing around the periods requires
a special command.

\LaTeX\ interprets some common characters as
commands, so you must type special commands to
generate them.  These characters include the
following:
       \$ \& \% \# \{ and \}.

In printing, text is usually emphasized with an
       \emph{italic}
type style.

\begin{em}
   A long segment of text can also be emphasized
   in this way.  Text within such a segment can be
   given \emph{additional} emphasis.
\end{em}

It is sometimes necessary to prevent \LaTeX\ from
breaking a line where it might otherwise do so.
This may be at a space, as between the ``Mr.''\ and
``Jones'' in
       ``Mr.~Jones'',        % ~ produces an unbreakable interword space.
or within a word---especially when the word is a
symbol like
       \mbox{\emph{itemnum}}
that makes little sense when hyphenated across
lines.

Footnotes\footnote{This is an example of a footnote.}
pose no problem.

\LaTeX\ is good at typesetting mathematical formulas
like
       \( x-3y + z = 7 \)
or
       \( a_{1} > x^{2n} + y^{2n} > x' \)
or
       \( \ip{A}{B} = \sum_{i} a_{i} b_{i} \).
The spaces you type in a formula are
ignored.  Remember that a letter like
       $x$                   % $ ... $  and  \( ... \)  are equivalent
is a formula when it denotes a mathematical
symbol, and it should be typed as one.

\section{Displayed Text}

Text is displayed by indenting it from the left
margin.  Quotations are commonly displayed.  There
are short quotations
\begin{quote}
   This is a short quotation.  It consists of a
   single paragraph of text.  See how it is formatted.
\end{quote}
and longer ones.
\begin{quotation}
   This is a longer quotation.  It consists of two
   paragraphs of text, neither of which are
   particularly interesting.

   This is the second paragraph of the quotation.  It
   is just as dull as the first paragraph.
\end{quotation}
Another frequently-displayed structure is a list.
The following is an example of an \emph{itemized}
list.
\begin{itemize}
   \item This is the first item of an itemized list.
         Each item in the list is marked with a ``tick''.
         You don't have to worry about what kind of tick
         mark is used.

   \item This is the second item of the list.  It
         contains another list nested inside it.  The inner
         list is an \emph{enumerated} list.
         \begin{enumerate}
            \item This is the first item of an enumerated
                  list that is nested within the itemized list.

            \item This is the second item of the inner list.
                  \LaTeX\ allows you to nest lists deeper than
                  you really should.
         \end{enumerate}
         This is the rest of the second item of the outer
         list.  It is no more interesting than any other
         part of the item.
   \item This is the third item of the list.
\end{itemize}
You can even display poetry.
\begin{verse}
   There is an environment
    for verse \\             % The \\ command separates lines
   Whose features some poets % within a stanza.
   will curse.

                             % One or more blank lines separate stanzas.

   For instead of making\\
   Them do \emph{all} line breaking, \\
   It allows them to put too many words on a line when they'd rather be
   forced to be terse.
\end{verse}

Mathematical formulas may also be displayed.  A
displayed formula
is
one-line long; multiline
formulas require special formatting instructions.
   \[  \ip{\Gamma}{\psi'} = x'' + y^{2} + z_{i}^{n}\]
Don't start a paragraph with a displayed equation,
nor make one a paragraph by itself.

\end{document}               % End of document.

\section{Bench Huge Section 800}
% synthetic bench section 800
% This is a sample LaTeX input file.  (Version of 12 August 2004.)
%
% A '%' character causes TeX to ignore all remaining text on the line,
% and is used for comments like this one.

\documentclass{article}      % Specifies the document class

                             % The preamble begins here.
\title{An Example Document}  % Declares the document's title.
\author{Leslie Lamport}      % Declares the author's name.
\date{January 21, 1994}      % Deleting this command produces today's date.

\newcommand{\ip}[2]{(#1, #2)}
                             % Defines \ip{arg1}{arg2} to mean
                             % (arg1, arg2).

%\newcommand{\ip}[2]{\langle #1 | #2\rangle}
                             % This is an alternative definition of
                             % \ip that is commented out.

\begin{document}             % End of preamble and beginning of text.

\maketitle                   % Produces the title.

This is an example input file.  Comparing it with
the output it generates can show you how to
produce a simple document of your own.

\section{Ordinary Text}      % Produces section heading.  Lower-level
                             % sections are begun with similar
                             % \subsection and \subsubsection commands.

The ends  of words and sentences are marked
  by   spaces. It  doesn't matter how many
spaces    you type; one is as good as 100.  The
end of   a line counts as a space.

One   or more   blank lines denote the  end
of  a paragraph.

Since any number of consecutive spaces are treated
like a single one, the formatting of the input
file makes no difference to
      \LaTeX,                % The \LaTeX command generates the LaTeX logo.
but it makes a difference to you.  When you use
\LaTeX, making your input file as easy to read
as possible will be a great help as you write
your document and when you change it.  This sample
file shows how you can add comments to your own input
file.

Because printing is different from typewriting,
there are a number of things that you have to do
differently when preparing an input file than if
you were just typing the document directly.
Quotation marks like
       ``this''
have to be handled specially, as do quotes within
quotes:
       ``\,`this'            % \, separates the double and single quote.
        is what I just
        wrote, not  `that'\,''.

Dashes come in three sizes: an
       intra-word
dash, a medium dash for number ranges like
       1--2,
and a punctuation
       dash---like
this.

A sentence-ending space should be larger than the
space between words within a sentence.  You
sometimes have to type special commands in
conjunction with punctuation characters to get
this right, as in the following sentence.
       Gnats, gnus, etc.\ all  % `\ ' makes an inter-word space.
       begin with G\@.         % \@ marks end-of-sentence punctuation.
You should check the spaces after periods when
reading your output to make sure you haven't
forgotten any special cases.  Generating an
ellipsis
       \ldots\               % `\ ' is needed after `\ldots' because TeX
                             % ignores spaces after command names like \ldots
                             % made from \ + letters.
                             %
                             % Note how a `%' character causes TeX to ignore
                             % the end of the input line, so these blank lines
                             % do not start a new paragraph.
                             %
with the right spacing around the periods requires
a special command.

\LaTeX\ interprets some common characters as
commands, so you must type special commands to
generate them.  These characters include the
following:
       \$ \& \% \# \{ and \}.

In printing, text is usually emphasized with an
       \emph{italic}
type style.

\begin{em}
   A long segment of text can also be emphasized
   in this way.  Text within such a segment can be
   given \emph{additional} emphasis.
\end{em}

It is sometimes necessary to prevent \LaTeX\ from
breaking a line where it might otherwise do so.
This may be at a space, as between the ``Mr.''\ and
``Jones'' in
       ``Mr.~Jones'',        % ~ produces an unbreakable interword space.
or within a word---especially when the word is a
symbol like
       \mbox{\emph{itemnum}}
that makes little sense when hyphenated across
lines.

Footnotes\footnote{This is an example of a footnote.}
pose no problem.

\LaTeX\ is good at typesetting mathematical formulas
like
       \( x-3y + z = 7 \)
or
       \( a_{1} > x^{2n} + y^{2n} > x' \)
or
       \( \ip{A}{B} = \sum_{i} a_{i} b_{i} \).
The spaces you type in a formula are
ignored.  Remember that a letter like
       $x$                   % $ ... $  and  \( ... \)  are equivalent
is a formula when it denotes a mathematical
symbol, and it should be typed as one.

\section{Displayed Text}

Text is displayed by indenting it from the left
margin.  Quotations are commonly displayed.  There
are short quotations
\begin{quote}
   This is a short quotation.  It consists of a
   single paragraph of text.  See how it is formatted.
\end{quote}
and longer ones.
\begin{quotation}
   This is a longer quotation.  It consists of two
   paragraphs of text, neither of which are
   particularly interesting.

   This is the second paragraph of the quotation.  It
   is just as dull as the first paragraph.
\end{quotation}
Another frequently-displayed structure is a list.
The following is an example of an \emph{itemized}
list.
\begin{itemize}
   \item This is the first item of an itemized list.
         Each item in the list is marked with a ``tick''.
         You don't have to worry about what kind of tick
         mark is used.

   \item This is the second item of the list.  It
         contains another list nested inside it.  The inner
         list is an \emph{enumerated} list.
         \begin{enumerate}
            \item This is the first item of an enumerated
                  list that is nested within the itemized list.

            \item This is the second item of the inner list.
                  \LaTeX\ allows you to nest lists deeper than
                  you really should.
         \end{enumerate}
         This is the rest of the second item of the outer
         list.  It is no more interesting than any other
         part of the item.
   \item This is the third item of the list.
\end{itemize}
You can even display poetry.
\begin{verse}
   There is an environment
    for verse \\             % The \\ command separates lines
   Whose features some poets % within a stanza.
   will curse.

                             % One or more blank lines separate stanzas.

   For instead of making\\
   Them do \emph{all} line breaking, \\
   It allows them to put too many words on a line when they'd rather be
   forced to be terse.
\end{verse}

Mathematical formulas may also be displayed.  A
displayed formula
is
one-line long; multiline
formulas require special formatting instructions.
   \[  \ip{\Gamma}{\psi'} = x'' + y^{2} + z_{i}^{n}\]
Don't start a paragraph with a displayed equation,
nor make one a paragraph by itself.

\end{document}               % End of document.

\section{Bench Huge Section 801}
% synthetic bench section 801
% This is a sample LaTeX input file.  (Version of 12 August 2004.)
%
% A '%' character causes TeX to ignore all remaining text on the line,
% and is used for comments like this one.

\documentclass{article}      % Specifies the document class

                             % The preamble begins here.
\title{An Example Document}  % Declares the document's title.
\author{Leslie Lamport}      % Declares the author's name.
\date{January 21, 1994}      % Deleting this command produces today's date.

\newcommand{\ip}[2]{(#1, #2)}
                             % Defines \ip{arg1}{arg2} to mean
                             % (arg1, arg2).

%\newcommand{\ip}[2]{\langle #1 | #2\rangle}
                             % This is an alternative definition of
                             % \ip that is commented out.

\begin{document}             % End of preamble and beginning of text.

\maketitle                   % Produces the title.

This is an example input file.  Comparing it with
the output it generates can show you how to
produce a simple document of your own.

\section{Ordinary Text}      % Produces section heading.  Lower-level
                             % sections are begun with similar
                             % \subsection and \subsubsection commands.

The ends  of words and sentences are marked
  by   spaces. It  doesn't matter how many
spaces    you type; one is as good as 100.  The
end of   a line counts as a space.

One   or more   blank lines denote the  end
of  a paragraph.

Since any number of consecutive spaces are treated
like a single one, the formatting of the input
file makes no difference to
      \LaTeX,                % The \LaTeX command generates the LaTeX logo.
but it makes a difference to you.  When you use
\LaTeX, making your input file as easy to read
as possible will be a great help as you write
your document and when you change it.  This sample
file shows how you can add comments to your own input
file.

Because printing is different from typewriting,
there are a number of things that you have to do
differently when preparing an input file than if
you were just typing the document directly.
Quotation marks like
       ``this''
have to be handled specially, as do quotes within
quotes:
       ``\,`this'            % \, separates the double and single quote.
        is what I just
        wrote, not  `that'\,''.

Dashes come in three sizes: an
       intra-word
dash, a medium dash for number ranges like
       1--2,
and a punctuation
       dash---like
this.

A sentence-ending space should be larger than the
space between words within a sentence.  You
sometimes have to type special commands in
conjunction with punctuation characters to get
this right, as in the following sentence.
       Gnats, gnus, etc.\ all  % `\ ' makes an inter-word space.
       begin with G\@.         % \@ marks end-of-sentence punctuation.
You should check the spaces after periods when
reading your output to make sure you haven't
forgotten any special cases.  Generating an
ellipsis
       \ldots\               % `\ ' is needed after `\ldots' because TeX
                             % ignores spaces after command names like \ldots
                             % made from \ + letters.
                             %
                             % Note how a `%' character causes TeX to ignore
                             % the end of the input line, so these blank lines
                             % do not start a new paragraph.
                             %
with the right spacing around the periods requires
a special command.

\LaTeX\ interprets some common characters as
commands, so you must type special commands to
generate them.  These characters include the
following:
       \$ \& \% \# \{ and \}.

In printing, text is usually emphasized with an
       \emph{italic}
type style.

\begin{em}
   A long segment of text can also be emphasized
   in this way.  Text within such a segment can be
   given \emph{additional} emphasis.
\end{em}

It is sometimes necessary to prevent \LaTeX\ from
breaking a line where it might otherwise do so.
This may be at a space, as between the ``Mr.''\ and
``Jones'' in
       ``Mr.~Jones'',        % ~ produces an unbreakable interword space.
or within a word---especially when the word is a
symbol like
       \mbox{\emph{itemnum}}
that makes little sense when hyphenated across
lines.

Footnotes\footnote{This is an example of a footnote.}
pose no problem.

\LaTeX\ is good at typesetting mathematical formulas
like
       \( x-3y + z = 7 \)
or
       \( a_{1} > x^{2n} + y^{2n} > x' \)
or
       \( \ip{A}{B} = \sum_{i} a_{i} b_{i} \).
The spaces you type in a formula are
ignored.  Remember that a letter like
       $x$                   % $ ... $  and  \( ... \)  are equivalent
is a formula when it denotes a mathematical
symbol, and it should be typed as one.

\section{Displayed Text}

Text is displayed by indenting it from the left
margin.  Quotations are commonly displayed.  There
are short quotations
\begin{quote}
   This is a short quotation.  It consists of a
   single paragraph of text.  See how it is formatted.
\end{quote}
and longer ones.
\begin{quotation}
   This is a longer quotation.  It consists of two
   paragraphs of text, neither of which are
   particularly interesting.

   This is the second paragraph of the quotation.  It
   is just as dull as the first paragraph.
\end{quotation}
Another frequently-displayed structure is a list.
The following is an example of an \emph{itemized}
list.
\begin{itemize}
   \item This is the first item of an itemized list.
         Each item in the list is marked with a ``tick''.
         You don't have to worry about what kind of tick
         mark is used.

   \item This is the second item of the list.  It
         contains another list nested inside it.  The inner
         list is an \emph{enumerated} list.
         \begin{enumerate}
            \item This is the first item of an enumerated
                  list that is nested within the itemized list.

            \item This is the second item of the inner list.
                  \LaTeX\ allows you to nest lists deeper than
                  you really should.
         \end{enumerate}
         This is the rest of the second item of the outer
         list.  It is no more interesting than any other
         part of the item.
   \item This is the third item of the list.
\end{itemize}
You can even display poetry.
\begin{verse}
   There is an environment
    for verse \\             % The \\ command separates lines
   Whose features some poets % within a stanza.
   will curse.

                             % One or more blank lines separate stanzas.

   For instead of making\\
   Them do \emph{all} line breaking, \\
   It allows them to put too many words on a line when they'd rather be
   forced to be terse.
\end{verse}

Mathematical formulas may also be displayed.  A
displayed formula
is
one-line long; multiline
formulas require special formatting instructions.
   \[  \ip{\Gamma}{\psi'} = x'' + y^{2} + z_{i}^{n}\]
Don't start a paragraph with a displayed equation,
nor make one a paragraph by itself.

\end{document}               % End of document.

\section{Bench Huge Section 802}
% synthetic bench section 802
% This is a sample LaTeX input file.  (Version of 12 August 2004.)
%
% A '%' character causes TeX to ignore all remaining text on the line,
% and is used for comments like this one.

\documentclass{article}      % Specifies the document class

                             % The preamble begins here.
\title{An Example Document}  % Declares the document's title.
\author{Leslie Lamport}      % Declares the author's name.
\date{January 21, 1994}      % Deleting this command produces today's date.

\newcommand{\ip}[2]{(#1, #2)}
                             % Defines \ip{arg1}{arg2} to mean
                             % (arg1, arg2).

%\newcommand{\ip}[2]{\langle #1 | #2\rangle}
                             % This is an alternative definition of
                             % \ip that is commented out.

\begin{document}             % End of preamble and beginning of text.

\maketitle                   % Produces the title.

This is an example input file.  Comparing it with
the output it generates can show you how to
produce a simple document of your own.

\section{Ordinary Text}      % Produces section heading.  Lower-level
                             % sections are begun with similar
                             % \subsection and \subsubsection commands.

The ends  of words and sentences are marked
  by   spaces. It  doesn't matter how many
spaces    you type; one is as good as 100.  The
end of   a line counts as a space.

One   or more   blank lines denote the  end
of  a paragraph.

Since any number of consecutive spaces are treated
like a single one, the formatting of the input
file makes no difference to
      \LaTeX,                % The \LaTeX command generates the LaTeX logo.
but it makes a difference to you.  When you use
\LaTeX, making your input file as easy to read
as possible will be a great help as you write
your document and when you change it.  This sample
file shows how you can add comments to your own input
file.

Because printing is different from typewriting,
there are a number of things that you have to do
differently when preparing an input file than if
you were just typing the document directly.
Quotation marks like
       ``this''
have to be handled specially, as do quotes within
quotes:
       ``\,`this'            % \, separates the double and single quote.
        is what I just
        wrote, not  `that'\,''.

Dashes come in three sizes: an
       intra-word
dash, a medium dash for number ranges like
       1--2,
and a punctuation
       dash---like
this.

A sentence-ending space should be larger than the
space between words within a sentence.  You
sometimes have to type special commands in
conjunction with punctuation characters to get
this right, as in the following sentence.
       Gnats, gnus, etc.\ all  % `\ ' makes an inter-word space.
       begin with G\@.         % \@ marks end-of-sentence punctuation.
You should check the spaces after periods when
reading your output to make sure you haven't
forgotten any special cases.  Generating an
ellipsis
       \ldots\               % `\ ' is needed after `\ldots' because TeX
                             % ignores spaces after command names like \ldots
                             % made from \ + letters.
                             %
                             % Note how a `%' character causes TeX to ignore
                             % the end of the input line, so these blank lines
                             % do not start a new paragraph.
                             %
with the right spacing around the periods requires
a special command.

\LaTeX\ interprets some common characters as
commands, so you must type special commands to
generate them.  These characters include the
following:
       \$ \& \% \# \{ and \}.

In printing, text is usually emphasized with an
       \emph{italic}
type style.

\begin{em}
   A long segment of text can also be emphasized
   in this way.  Text within such a segment can be
   given \emph{additional} emphasis.
\end{em}

It is sometimes necessary to prevent \LaTeX\ from
breaking a line where it might otherwise do so.
This may be at a space, as between the ``Mr.''\ and
``Jones'' in
       ``Mr.~Jones'',        % ~ produces an unbreakable interword space.
or within a word---especially when the word is a
symbol like
       \mbox{\emph{itemnum}}
that makes little sense when hyphenated across
lines.

Footnotes\footnote{This is an example of a footnote.}
pose no problem.

\LaTeX\ is good at typesetting mathematical formulas
like
       \( x-3y + z = 7 \)
or
       \( a_{1} > x^{2n} + y^{2n} > x' \)
or
       \( \ip{A}{B} = \sum_{i} a_{i} b_{i} \).
The spaces you type in a formula are
ignored.  Remember that a letter like
       $x$                   % $ ... $  and  \( ... \)  are equivalent
is a formula when it denotes a mathematical
symbol, and it should be typed as one.

\section{Displayed Text}

Text is displayed by indenting it from the left
margin.  Quotations are commonly displayed.  There
are short quotations
\begin{quote}
   This is a short quotation.  It consists of a
   single paragraph of text.  See how it is formatted.
\end{quote}
and longer ones.
\begin{quotation}
   This is a longer quotation.  It consists of two
   paragraphs of text, neither of which are
   particularly interesting.

   This is the second paragraph of the quotation.  It
   is just as dull as the first paragraph.
\end{quotation}
Another frequently-displayed structure is a list.
The following is an example of an \emph{itemized}
list.
\begin{itemize}
   \item This is the first item of an itemized list.
         Each item in the list is marked with a ``tick''.
         You don't have to worry about what kind of tick
         mark is used.

   \item This is the second item of the list.  It
         contains another list nested inside it.  The inner
         list is an \emph{enumerated} list.
         \begin{enumerate}
            \item This is the first item of an enumerated
                  list that is nested within the itemized list.

            \item This is the second item of the inner list.
                  \LaTeX\ allows you to nest lists deeper than
                  you really should.
         \end{enumerate}
         This is the rest of the second item of the outer
         list.  It is no more interesting than any other
         part of the item.
   \item This is the third item of the list.
\end{itemize}
You can even display poetry.
\begin{verse}
   There is an environment
    for verse \\             % The \\ command separates lines
   Whose features some poets % within a stanza.
   will curse.

                             % One or more blank lines separate stanzas.

   For instead of making\\
   Them do \emph{all} line breaking, \\
   It allows them to put too many words on a line when they'd rather be
   forced to be terse.
\end{verse}

Mathematical formulas may also be displayed.  A
displayed formula
is
one-line long; multiline
formulas require special formatting instructions.
   \[  \ip{\Gamma}{\psi'} = x'' + y^{2} + z_{i}^{n}\]
Don't start a paragraph with a displayed equation,
nor make one a paragraph by itself.

\end{document}               % End of document.

\section{Bench Huge Section 803}
% synthetic bench section 803
% This is a sample LaTeX input file.  (Version of 12 August 2004.)
%
% A '%' character causes TeX to ignore all remaining text on the line,
% and is used for comments like this one.

\documentclass{article}      % Specifies the document class

                             % The preamble begins here.
\title{An Example Document}  % Declares the document's title.
\author{Leslie Lamport}      % Declares the author's name.
\date{January 21, 1994}      % Deleting this command produces today's date.

\newcommand{\ip}[2]{(#1, #2)}
                             % Defines \ip{arg1}{arg2} to mean
                             % (arg1, arg2).

%\newcommand{\ip}[2]{\langle #1 | #2\rangle}
                             % This is an alternative definition of
                             % \ip that is commented out.

\begin{document}             % End of preamble and beginning of text.

\maketitle                   % Produces the title.

This is an example input file.  Comparing it with
the output it generates can show you how to
produce a simple document of your own.

\section{Ordinary Text}      % Produces section heading.  Lower-level
                             % sections are begun with similar
                             % \subsection and \subsubsection commands.

The ends  of words and sentences are marked
  by   spaces. It  doesn't matter how many
spaces    you type; one is as good as 100.  The
end of   a line counts as a space.

One   or more   blank lines denote the  end
of  a paragraph.

Since any number of consecutive spaces are treated
like a single one, the formatting of the input
file makes no difference to
      \LaTeX,                % The \LaTeX command generates the LaTeX logo.
but it makes a difference to you.  When you use
\LaTeX, making your input file as easy to read
as possible will be a great help as you write
your document and when you change it.  This sample
file shows how you can add comments to your own input
file.

Because printing is different from typewriting,
there are a number of things that you have to do
differently when preparing an input file than if
you were just typing the document directly.
Quotation marks like
       ``this''
have to be handled specially, as do quotes within
quotes:
       ``\,`this'            % \, separates the double and single quote.
        is what I just
        wrote, not  `that'\,''.

Dashes come in three sizes: an
       intra-word
dash, a medium dash for number ranges like
       1--2,
and a punctuation
       dash---like
this.

A sentence-ending space should be larger than the
space between words within a sentence.  You
sometimes have to type special commands in
conjunction with punctuation characters to get
this right, as in the following sentence.
       Gnats, gnus, etc.\ all  % `\ ' makes an inter-word space.
       begin with G\@.         % \@ marks end-of-sentence punctuation.
You should check the spaces after periods when
reading your output to make sure you haven't
forgotten any special cases.  Generating an
ellipsis
       \ldots\               % `\ ' is needed after `\ldots' because TeX
                             % ignores spaces after command names like \ldots
                             % made from \ + letters.
                             %
                             % Note how a `%' character causes TeX to ignore
                             % the end of the input line, so these blank lines
                             % do not start a new paragraph.
                             %
with the right spacing around the periods requires
a special command.

\LaTeX\ interprets some common characters as
commands, so you must type special commands to
generate them.  These characters include the
following:
       \$ \& \% \# \{ and \}.

In printing, text is usually emphasized with an
       \emph{italic}
type style.

\begin{em}
   A long segment of text can also be emphasized
   in this way.  Text within such a segment can be
   given \emph{additional} emphasis.
\end{em}

It is sometimes necessary to prevent \LaTeX\ from
breaking a line where it might otherwise do so.
This may be at a space, as between the ``Mr.''\ and
``Jones'' in
       ``Mr.~Jones'',        % ~ produces an unbreakable interword space.
or within a word---especially when the word is a
symbol like
       \mbox{\emph{itemnum}}
that makes little sense when hyphenated across
lines.

Footnotes\footnote{This is an example of a footnote.}
pose no problem.

\LaTeX\ is good at typesetting mathematical formulas
like
       \( x-3y + z = 7 \)
or
       \( a_{1} > x^{2n} + y^{2n} > x' \)
or
       \( \ip{A}{B} = \sum_{i} a_{i} b_{i} \).
The spaces you type in a formula are
ignored.  Remember that a letter like
       $x$                   % $ ... $  and  \( ... \)  are equivalent
is a formula when it denotes a mathematical
symbol, and it should be typed as one.

\section{Displayed Text}

Text is displayed by indenting it from the left
margin.  Quotations are commonly displayed.  There
are short quotations
\begin{quote}
   This is a short quotation.  It consists of a
   single paragraph of text.  See how it is formatted.
\end{quote}
and longer ones.
\begin{quotation}
   This is a longer quotation.  It consists of two
   paragraphs of text, neither of which are
   particularly interesting.

   This is the second paragraph of the quotation.  It
   is just as dull as the first paragraph.
\end{quotation}
Another frequently-displayed structure is a list.
The following is an example of an \emph{itemized}
list.
\begin{itemize}
   \item This is the first item of an itemized list.
         Each item in the list is marked with a ``tick''.
         You don't have to worry about what kind of tick
         mark is used.

   \item This is the second item of the list.  It
         contains another list nested inside it.  The inner
         list is an \emph{enumerated} list.
         \begin{enumerate}
            \item This is the first item of an enumerated
                  list that is nested within the itemized list.

            \item This is the second item of the inner list.
                  \LaTeX\ allows you to nest lists deeper than
                  you really should.
         \end{enumerate}
         This is the rest of the second item of the outer
         list.  It is no more interesting than any other
         part of the item.
   \item This is the third item of the list.
\end{itemize}
You can even display poetry.
\begin{verse}
   There is an environment
    for verse \\             % The \\ command separates lines
   Whose features some poets % within a stanza.
   will curse.

                             % One or more blank lines separate stanzas.

   For instead of making\\
   Them do \emph{all} line breaking, \\
   It allows them to put too many words on a line when they'd rather be
   forced to be terse.
\end{verse}

Mathematical formulas may also be displayed.  A
displayed formula
is
one-line long; multiline
formulas require special formatting instructions.
   \[  \ip{\Gamma}{\psi'} = x'' + y^{2} + z_{i}^{n}\]
Don't start a paragraph with a displayed equation,
nor make one a paragraph by itself.

\end{document}               % End of document.

\section{Bench Huge Section 804}
% synthetic bench section 804
% This is a sample LaTeX input file.  (Version of 12 August 2004.)
%
% A '%' character causes TeX to ignore all remaining text on the line,
% and is used for comments like this one.

\documentclass{article}      % Specifies the document class

                             % The preamble begins here.
\title{An Example Document}  % Declares the document's title.
\author{Leslie Lamport}      % Declares the author's name.
\date{January 21, 1994}      % Deleting this command produces today's date.

\newcommand{\ip}[2]{(#1, #2)}
                             % Defines \ip{arg1}{arg2} to mean
                             % (arg1, arg2).

%\newcommand{\ip}[2]{\langle #1 | #2\rangle}
                             % This is an alternative definition of
                             % \ip that is commented out.

\begin{document}             % End of preamble and beginning of text.

\maketitle                   % Produces the title.

This is an example input file.  Comparing it with
the output it generates can show you how to
produce a simple document of your own.

\section{Ordinary Text}      % Produces section heading.  Lower-level
                             % sections are begun with similar
                             % \subsection and \subsubsection commands.

The ends  of words and sentences are marked
  by   spaces. It  doesn't matter how many
spaces    you type; one is as good as 100.  The
end of   a line counts as a space.

One   or more   blank lines denote the  end
of  a paragraph.

Since any number of consecutive spaces are treated
like a single one, the formatting of the input
file makes no difference to
      \LaTeX,                % The \LaTeX command generates the LaTeX logo.
but it makes a difference to you.  When you use
\LaTeX, making your input file as easy to read
as possible will be a great help as you write
your document and when you change it.  This sample
file shows how you can add comments to your own input
file.

Because printing is different from typewriting,
there are a number of things that you have to do
differently when preparing an input file than if
you were just typing the document directly.
Quotation marks like
       ``this''
have to be handled specially, as do quotes within
quotes:
       ``\,`this'            % \, separates the double and single quote.
        is what I just
        wrote, not  `that'\,''.

Dashes come in three sizes: an
       intra-word
dash, a medium dash for number ranges like
       1--2,
and a punctuation
       dash---like
this.

A sentence-ending space should be larger than the
space between words within a sentence.  You
sometimes have to type special commands in
conjunction with punctuation characters to get
this right, as in the following sentence.
       Gnats, gnus, etc.\ all  % `\ ' makes an inter-word space.
       begin with G\@.         % \@ marks end-of-sentence punctuation.
You should check the spaces after periods when
reading your output to make sure you haven't
forgotten any special cases.  Generating an
ellipsis
       \ldots\               % `\ ' is needed after `\ldots' because TeX
                             % ignores spaces after command names like \ldots
                             % made from \ + letters.
                             %
                             % Note how a `%' character causes TeX to ignore
                             % the end of the input line, so these blank lines
                             % do not start a new paragraph.
                             %
with the right spacing around the periods requires
a special command.

\LaTeX\ interprets some common characters as
commands, so you must type special commands to
generate them.  These characters include the
following:
       \$ \& \% \# \{ and \}.

In printing, text is usually emphasized with an
       \emph{italic}
type style.

\begin{em}
   A long segment of text can also be emphasized
   in this way.  Text within such a segment can be
   given \emph{additional} emphasis.
\end{em}

It is sometimes necessary to prevent \LaTeX\ from
breaking a line where it might otherwise do so.
This may be at a space, as between the ``Mr.''\ and
``Jones'' in
       ``Mr.~Jones'',        % ~ produces an unbreakable interword space.
or within a word---especially when the word is a
symbol like
       \mbox{\emph{itemnum}}
that makes little sense when hyphenated across
lines.

Footnotes\footnote{This is an example of a footnote.}
pose no problem.

\LaTeX\ is good at typesetting mathematical formulas
like
       \( x-3y + z = 7 \)
or
       \( a_{1} > x^{2n} + y^{2n} > x' \)
or
       \( \ip{A}{B} = \sum_{i} a_{i} b_{i} \).
The spaces you type in a formula are
ignored.  Remember that a letter like
       $x$                   % $ ... $  and  \( ... \)  are equivalent
is a formula when it denotes a mathematical
symbol, and it should be typed as one.

\section{Displayed Text}

Text is displayed by indenting it from the left
margin.  Quotations are commonly displayed.  There
are short quotations
\begin{quote}
   This is a short quotation.  It consists of a
   single paragraph of text.  See how it is formatted.
\end{quote}
and longer ones.
\begin{quotation}
   This is a longer quotation.  It consists of two
   paragraphs of text, neither of which are
   particularly interesting.

   This is the second paragraph of the quotation.  It
   is just as dull as the first paragraph.
\end{quotation}
Another frequently-displayed structure is a list.
The following is an example of an \emph{itemized}
list.
\begin{itemize}
   \item This is the first item of an itemized list.
         Each item in the list is marked with a ``tick''.
         You don't have to worry about what kind of tick
         mark is used.

   \item This is the second item of the list.  It
         contains another list nested inside it.  The inner
         list is an \emph{enumerated} list.
         \begin{enumerate}
            \item This is the first item of an enumerated
                  list that is nested within the itemized list.

            \item This is the second item of the inner list.
                  \LaTeX\ allows you to nest lists deeper than
                  you really should.
         \end{enumerate}
         This is the rest of the second item of the outer
         list.  It is no more interesting than any other
         part of the item.
   \item This is the third item of the list.
\end{itemize}
You can even display poetry.
\begin{verse}
   There is an environment
    for verse \\             % The \\ command separates lines
   Whose features some poets % within a stanza.
   will curse.

                             % One or more blank lines separate stanzas.

   For instead of making\\
   Them do \emph{all} line breaking, \\
   It allows them to put too many words on a line when they'd rather be
   forced to be terse.
\end{verse}

Mathematical formulas may also be displayed.  A
displayed formula
is
one-line long; multiline
formulas require special formatting instructions.
   \[  \ip{\Gamma}{\psi'} = x'' + y^{2} + z_{i}^{n}\]
Don't start a paragraph with a displayed equation,
nor make one a paragraph by itself.

\end{document}               % End of document.

\section{Bench Huge Section 805}
% synthetic bench section 805
% This is a sample LaTeX input file.  (Version of 12 August 2004.)
%
% A '%' character causes TeX to ignore all remaining text on the line,
% and is used for comments like this one.

\documentclass{article}      % Specifies the document class

                             % The preamble begins here.
\title{An Example Document}  % Declares the document's title.
\author{Leslie Lamport}      % Declares the author's name.
\date{January 21, 1994}      % Deleting this command produces today's date.

\newcommand{\ip}[2]{(#1, #2)}
                             % Defines \ip{arg1}{arg2} to mean
                             % (arg1, arg2).

%\newcommand{\ip}[2]{\langle #1 | #2\rangle}
                             % This is an alternative definition of
                             % \ip that is commented out.

\begin{document}             % End of preamble and beginning of text.

\maketitle                   % Produces the title.

This is an example input file.  Comparing it with
the output it generates can show you how to
produce a simple document of your own.

\section{Ordinary Text}      % Produces section heading.  Lower-level
                             % sections are begun with similar
                             % \subsection and \subsubsection commands.

The ends  of words and sentences are marked
  by   spaces. It  doesn't matter how many
spaces    you type; one is as good as 100.  The
end of   a line counts as a space.

One   or more   blank lines denote the  end
of  a paragraph.

Since any number of consecutive spaces are treated
like a single one, the formatting of the input
file makes no difference to
      \LaTeX,                % The \LaTeX command generates the LaTeX logo.
but it makes a difference to you.  When you use
\LaTeX, making your input file as easy to read
as possible will be a great help as you write
your document and when you change it.  This sample
file shows how you can add comments to your own input
file.

Because printing is different from typewriting,
there are a number of things that you have to do
differently when preparing an input file than if
you were just typing the document directly.
Quotation marks like
       ``this''
have to be handled specially, as do quotes within
quotes:
       ``\,`this'            % \, separates the double and single quote.
        is what I just
        wrote, not  `that'\,''.

Dashes come in three sizes: an
       intra-word
dash, a medium dash for number ranges like
       1--2,
and a punctuation
       dash---like
this.

A sentence-ending space should be larger than the
space between words within a sentence.  You
sometimes have to type special commands in
conjunction with punctuation characters to get
this right, as in the following sentence.
       Gnats, gnus, etc.\ all  % `\ ' makes an inter-word space.
       begin with G\@.         % \@ marks end-of-sentence punctuation.
You should check the spaces after periods when
reading your output to make sure you haven't
forgotten any special cases.  Generating an
ellipsis
       \ldots\               % `\ ' is needed after `\ldots' because TeX
                             % ignores spaces after command names like \ldots
                             % made from \ + letters.
                             %
                             % Note how a `%' character causes TeX to ignore
                             % the end of the input line, so these blank lines
                             % do not start a new paragraph.
                             %
with the right spacing around the periods requires
a special command.

\LaTeX\ interprets some common characters as
commands, so you must type special commands to
generate them.  These characters include the
following:
       \$ \& \% \# \{ and \}.

In printing, text is usually emphasized with an
       \emph{italic}
type style.

\begin{em}
   A long segment of text can also be emphasized
   in this way.  Text within such a segment can be
   given \emph{additional} emphasis.
\end{em}

It is sometimes necessary to prevent \LaTeX\ from
breaking a line where it might otherwise do so.
This may be at a space, as between the ``Mr.''\ and
``Jones'' in
       ``Mr.~Jones'',        % ~ produces an unbreakable interword space.
or within a word---especially when the word is a
symbol like
       \mbox{\emph{itemnum}}
that makes little sense when hyphenated across
lines.

Footnotes\footnote{This is an example of a footnote.}
pose no problem.

\LaTeX\ is good at typesetting mathematical formulas
like
       \( x-3y + z = 7 \)
or
       \( a_{1} > x^{2n} + y^{2n} > x' \)
or
       \( \ip{A}{B} = \sum_{i} a_{i} b_{i} \).
The spaces you type in a formula are
ignored.  Remember that a letter like
       $x$                   % $ ... $  and  \( ... \)  are equivalent
is a formula when it denotes a mathematical
symbol, and it should be typed as one.

\section{Displayed Text}

Text is displayed by indenting it from the left
margin.  Quotations are commonly displayed.  There
are short quotations
\begin{quote}
   This is a short quotation.  It consists of a
   single paragraph of text.  See how it is formatted.
\end{quote}
and longer ones.
\begin{quotation}
   This is a longer quotation.  It consists of two
   paragraphs of text, neither of which are
   particularly interesting.

   This is the second paragraph of the quotation.  It
   is just as dull as the first paragraph.
\end{quotation}
Another frequently-displayed structure is a list.
The following is an example of an \emph{itemized}
list.
\begin{itemize}
   \item This is the first item of an itemized list.
         Each item in the list is marked with a ``tick''.
         You don't have to worry about what kind of tick
         mark is used.

   \item This is the second item of the list.  It
         contains another list nested inside it.  The inner
         list is an \emph{enumerated} list.
         \begin{enumerate}
            \item This is the first item of an enumerated
                  list that is nested within the itemized list.

            \item This is the second item of the inner list.
                  \LaTeX\ allows you to nest lists deeper than
                  you really should.
         \end{enumerate}
         This is the rest of the second item of the outer
         list.  It is no more interesting than any other
         part of the item.
   \item This is the third item of the list.
\end{itemize}
You can even display poetry.
\begin{verse}
   There is an environment
    for verse \\             % The \\ command separates lines
   Whose features some poets % within a stanza.
   will curse.

                             % One or more blank lines separate stanzas.

   For instead of making\\
   Them do \emph{all} line breaking, \\
   It allows them to put too many words on a line when they'd rather be
   forced to be terse.
\end{verse}

Mathematical formulas may also be displayed.  A
displayed formula
is
one-line long; multiline
formulas require special formatting instructions.
   \[  \ip{\Gamma}{\psi'} = x'' + y^{2} + z_{i}^{n}\]
Don't start a paragraph with a displayed equation,
nor make one a paragraph by itself.

\end{document}               % End of document.

\section{Bench Huge Section 806}
% synthetic bench section 806
% This is a sample LaTeX input file.  (Version of 12 August 2004.)
%
% A '%' character causes TeX to ignore all remaining text on the line,
% and is used for comments like this one.

\documentclass{article}      % Specifies the document class

                             % The preamble begins here.
\title{An Example Document}  % Declares the document's title.
\author{Leslie Lamport}      % Declares the author's name.
\date{January 21, 1994}      % Deleting this command produces today's date.

\newcommand{\ip}[2]{(#1, #2)}
                             % Defines \ip{arg1}{arg2} to mean
                             % (arg1, arg2).

%\newcommand{\ip}[2]{\langle #1 | #2\rangle}
                             % This is an alternative definition of
                             % \ip that is commented out.

\begin{document}             % End of preamble and beginning of text.

\maketitle                   % Produces the title.

This is an example input file.  Comparing it with
the output it generates can show you how to
produce a simple document of your own.

\section{Ordinary Text}      % Produces section heading.  Lower-level
                             % sections are begun with similar
                             % \subsection and \subsubsection commands.

The ends  of words and sentences are marked
  by   spaces. It  doesn't matter how many
spaces    you type; one is as good as 100.  The
end of   a line counts as a space.

One   or more   blank lines denote the  end
of  a paragraph.

Since any number of consecutive spaces are treated
like a single one, the formatting of the input
file makes no difference to
      \LaTeX,                % The \LaTeX command generates the LaTeX logo.
but it makes a difference to you.  When you use
\LaTeX, making your input file as easy to read
as possible will be a great help as you write
your document and when you change it.  This sample
file shows how you can add comments to your own input
file.

Because printing is different from typewriting,
there are a number of things that you have to do
differently when preparing an input file than if
you were just typing the document directly.
Quotation marks like
       ``this''
have to be handled specially, as do quotes within
quotes:
       ``\,`this'            % \, separates the double and single quote.
        is what I just
        wrote, not  `that'\,''.

Dashes come in three sizes: an
       intra-word
dash, a medium dash for number ranges like
       1--2,
and a punctuation
       dash---like
this.

A sentence-ending space should be larger than the
space between words within a sentence.  You
sometimes have to type special commands in
conjunction with punctuation characters to get
this right, as in the following sentence.
       Gnats, gnus, etc.\ all  % `\ ' makes an inter-word space.
       begin with G\@.         % \@ marks end-of-sentence punctuation.
You should check the spaces after periods when
reading your output to make sure you haven't
forgotten any special cases.  Generating an
ellipsis
       \ldots\               % `\ ' is needed after `\ldots' because TeX
                             % ignores spaces after command names like \ldots
                             % made from \ + letters.
                             %
                             % Note how a `%' character causes TeX to ignore
                             % the end of the input line, so these blank lines
                             % do not start a new paragraph.
                             %
with the right spacing around the periods requires
a special command.

\LaTeX\ interprets some common characters as
commands, so you must type special commands to
generate them.  These characters include the
following:
       \$ \& \% \# \{ and \}.

In printing, text is usually emphasized with an
       \emph{italic}
type style.

\begin{em}
   A long segment of text can also be emphasized
   in this way.  Text within such a segment can be
   given \emph{additional} emphasis.
\end{em}

It is sometimes necessary to prevent \LaTeX\ from
breaking a line where it might otherwise do so.
This may be at a space, as between the ``Mr.''\ and
``Jones'' in
       ``Mr.~Jones'',        % ~ produces an unbreakable interword space.
or within a word---especially when the word is a
symbol like
       \mbox{\emph{itemnum}}
that makes little sense when hyphenated across
lines.

Footnotes\footnote{This is an example of a footnote.}
pose no problem.

\LaTeX\ is good at typesetting mathematical formulas
like
       \( x-3y + z = 7 \)
or
       \( a_{1} > x^{2n} + y^{2n} > x' \)
or
       \( \ip{A}{B} = \sum_{i} a_{i} b_{i} \).
The spaces you type in a formula are
ignored.  Remember that a letter like
       $x$                   % $ ... $  and  \( ... \)  are equivalent
is a formula when it denotes a mathematical
symbol, and it should be typed as one.

\section{Displayed Text}

Text is displayed by indenting it from the left
margin.  Quotations are commonly displayed.  There
are short quotations
\begin{quote}
   This is a short quotation.  It consists of a
   single paragraph of text.  See how it is formatted.
\end{quote}
and longer ones.
\begin{quotation}
   This is a longer quotation.  It consists of two
   paragraphs of text, neither of which are
   particularly interesting.

   This is the second paragraph of the quotation.  It
   is just as dull as the first paragraph.
\end{quotation}
Another frequently-displayed structure is a list.
The following is an example of an \emph{itemized}
list.
\begin{itemize}
   \item This is the first item of an itemized list.
         Each item in the list is marked with a ``tick''.
         You don't have to worry about what kind of tick
         mark is used.

   \item This is the second item of the list.  It
         contains another list nested inside it.  The inner
         list is an \emph{enumerated} list.
         \begin{enumerate}
            \item This is the first item of an enumerated
                  list that is nested within the itemized list.

            \item This is the second item of the inner list.
                  \LaTeX\ allows you to nest lists deeper than
                  you really should.
         \end{enumerate}
         This is the rest of the second item of the outer
         list.  It is no more interesting than any other
         part of the item.
   \item This is the third item of the list.
\end{itemize}
You can even display poetry.
\begin{verse}
   There is an environment
    for verse \\             % The \\ command separates lines
   Whose features some poets % within a stanza.
   will curse.

                             % One or more blank lines separate stanzas.

   For instead of making\\
   Them do \emph{all} line breaking, \\
   It allows them to put too many words on a line when they'd rather be
   forced to be terse.
\end{verse}

Mathematical formulas may also be displayed.  A
displayed formula
is
one-line long; multiline
formulas require special formatting instructions.
   \[  \ip{\Gamma}{\psi'} = x'' + y^{2} + z_{i}^{n}\]
Don't start a paragraph with a displayed equation,
nor make one a paragraph by itself.

\end{document}               % End of document.

\section{Bench Huge Section 807}
% synthetic bench section 807
% This is a sample LaTeX input file.  (Version of 12 August 2004.)
%
% A '%' character causes TeX to ignore all remaining text on the line,
% and is used for comments like this one.

\documentclass{article}      % Specifies the document class

                             % The preamble begins here.
\title{An Example Document}  % Declares the document's title.
\author{Leslie Lamport}      % Declares the author's name.
\date{January 21, 1994}      % Deleting this command produces today's date.

\newcommand{\ip}[2]{(#1, #2)}
                             % Defines \ip{arg1}{arg2} to mean
                             % (arg1, arg2).

%\newcommand{\ip}[2]{\langle #1 | #2\rangle}
                             % This is an alternative definition of
                             % \ip that is commented out.

\begin{document}             % End of preamble and beginning of text.

\maketitle                   % Produces the title.

This is an example input file.  Comparing it with
the output it generates can show you how to
produce a simple document of your own.

\section{Ordinary Text}      % Produces section heading.  Lower-level
                             % sections are begun with similar
                             % \subsection and \subsubsection commands.

The ends  of words and sentences are marked
  by   spaces. It  doesn't matter how many
spaces    you type; one is as good as 100.  The
end of   a line counts as a space.

One   or more   blank lines denote the  end
of  a paragraph.

Since any number of consecutive spaces are treated
like a single one, the formatting of the input
file makes no difference to
      \LaTeX,                % The \LaTeX command generates the LaTeX logo.
but it makes a difference to you.  When you use
\LaTeX, making your input file as easy to read
as possible will be a great help as you write
your document and when you change it.  This sample
file shows how you can add comments to your own input
file.

Because printing is different from typewriting,
there are a number of things that you have to do
differently when preparing an input file than if
you were just typing the document directly.
Quotation marks like
       ``this''
have to be handled specially, as do quotes within
quotes:
       ``\,`this'            % \, separates the double and single quote.
        is what I just
        wrote, not  `that'\,''.

Dashes come in three sizes: an
       intra-word
dash, a medium dash for number ranges like
       1--2,
and a punctuation
       dash---like
this.

A sentence-ending space should be larger than the
space between words within a sentence.  You
sometimes have to type special commands in
conjunction with punctuation characters to get
this right, as in the following sentence.
       Gnats, gnus, etc.\ all  % `\ ' makes an inter-word space.
       begin with G\@.         % \@ marks end-of-sentence punctuation.
You should check the spaces after periods when
reading your output to make sure you haven't
forgotten any special cases.  Generating an
ellipsis
       \ldots\               % `\ ' is needed after `\ldots' because TeX
                             % ignores spaces after command names like \ldots
                             % made from \ + letters.
                             %
                             % Note how a `%' character causes TeX to ignore
                             % the end of the input line, so these blank lines
                             % do not start a new paragraph.
                             %
with the right spacing around the periods requires
a special command.

\LaTeX\ interprets some common characters as
commands, so you must type special commands to
generate them.  These characters include the
following:
       \$ \& \% \# \{ and \}.

In printing, text is usually emphasized with an
       \emph{italic}
type style.

\begin{em}
   A long segment of text can also be emphasized
   in this way.  Text within such a segment can be
   given \emph{additional} emphasis.
\end{em}

It is sometimes necessary to prevent \LaTeX\ from
breaking a line where it might otherwise do so.
This may be at a space, as between the ``Mr.''\ and
``Jones'' in
       ``Mr.~Jones'',        % ~ produces an unbreakable interword space.
or within a word---especially when the word is a
symbol like
       \mbox{\emph{itemnum}}
that makes little sense when hyphenated across
lines.

Footnotes\footnote{This is an example of a footnote.}
pose no problem.

\LaTeX\ is good at typesetting mathematical formulas
like
       \( x-3y + z = 7 \)
or
       \( a_{1} > x^{2n} + y^{2n} > x' \)
or
       \( \ip{A}{B} = \sum_{i} a_{i} b_{i} \).
The spaces you type in a formula are
ignored.  Remember that a letter like
       $x$                   % $ ... $  and  \( ... \)  are equivalent
is a formula when it denotes a mathematical
symbol, and it should be typed as one.

\section{Displayed Text}

Text is displayed by indenting it from the left
margin.  Quotations are commonly displayed.  There
are short quotations
\begin{quote}
   This is a short quotation.  It consists of a
   single paragraph of text.  See how it is formatted.
\end{quote}
and longer ones.
\begin{quotation}
   This is a longer quotation.  It consists of two
   paragraphs of text, neither of which are
   particularly interesting.

   This is the second paragraph of the quotation.  It
   is just as dull as the first paragraph.
\end{quotation}
Another frequently-displayed structure is a list.
The following is an example of an \emph{itemized}
list.
\begin{itemize}
   \item This is the first item of an itemized list.
         Each item in the list is marked with a ``tick''.
         You don't have to worry about what kind of tick
         mark is used.

   \item This is the second item of the list.  It
         contains another list nested inside it.  The inner
         list is an \emph{enumerated} list.
         \begin{enumerate}
            \item This is the first item of an enumerated
                  list that is nested within the itemized list.

            \item This is the second item of the inner list.
                  \LaTeX\ allows you to nest lists deeper than
                  you really should.
         \end{enumerate}
         This is the rest of the second item of the outer
         list.  It is no more interesting than any other
         part of the item.
   \item This is the third item of the list.
\end{itemize}
You can even display poetry.
\begin{verse}
   There is an environment
    for verse \\             % The \\ command separates lines
   Whose features some poets % within a stanza.
   will curse.

                             % One or more blank lines separate stanzas.

   For instead of making\\
   Them do \emph{all} line breaking, \\
   It allows them to put too many words on a line when they'd rather be
   forced to be terse.
\end{verse}

Mathematical formulas may also be displayed.  A
displayed formula
is
one-line long; multiline
formulas require special formatting instructions.
   \[  \ip{\Gamma}{\psi'} = x'' + y^{2} + z_{i}^{n}\]
Don't start a paragraph with a displayed equation,
nor make one a paragraph by itself.

\end{document}               % End of document.

\section{Bench Huge Section 808}
% synthetic bench section 808
% This is a sample LaTeX input file.  (Version of 12 August 2004.)
%
% A '%' character causes TeX to ignore all remaining text on the line,
% and is used for comments like this one.

\documentclass{article}      % Specifies the document class

                             % The preamble begins here.
\title{An Example Document}  % Declares the document's title.
\author{Leslie Lamport}      % Declares the author's name.
\date{January 21, 1994}      % Deleting this command produces today's date.

\newcommand{\ip}[2]{(#1, #2)}
                             % Defines \ip{arg1}{arg2} to mean
                             % (arg1, arg2).

%\newcommand{\ip}[2]{\langle #1 | #2\rangle}
                             % This is an alternative definition of
                             % \ip that is commented out.

\begin{document}             % End of preamble and beginning of text.

\maketitle                   % Produces the title.

This is an example input file.  Comparing it with
the output it generates can show you how to
produce a simple document of your own.

\section{Ordinary Text}      % Produces section heading.  Lower-level
                             % sections are begun with similar
                             % \subsection and \subsubsection commands.

The ends  of words and sentences are marked
  by   spaces. It  doesn't matter how many
spaces    you type; one is as good as 100.  The
end of   a line counts as a space.

One   or more   blank lines denote the  end
of  a paragraph.

Since any number of consecutive spaces are treated
like a single one, the formatting of the input
file makes no difference to
      \LaTeX,                % The \LaTeX command generates the LaTeX logo.
but it makes a difference to you.  When you use
\LaTeX, making your input file as easy to read
as possible will be a great help as you write
your document and when you change it.  This sample
file shows how you can add comments to your own input
file.

Because printing is different from typewriting,
there are a number of things that you have to do
differently when preparing an input file than if
you were just typing the document directly.
Quotation marks like
       ``this''
have to be handled specially, as do quotes within
quotes:
       ``\,`this'            % \, separates the double and single quote.
        is what I just
        wrote, not  `that'\,''.

Dashes come in three sizes: an
       intra-word
dash, a medium dash for number ranges like
       1--2,
and a punctuation
       dash---like
this.

A sentence-ending space should be larger than the
space between words within a sentence.  You
sometimes have to type special commands in
conjunction with punctuation characters to get
this right, as in the following sentence.
       Gnats, gnus, etc.\ all  % `\ ' makes an inter-word space.
       begin with G\@.         % \@ marks end-of-sentence punctuation.
You should check the spaces after periods when
reading your output to make sure you haven't
forgotten any special cases.  Generating an
ellipsis
       \ldots\               % `\ ' is needed after `\ldots' because TeX
                             % ignores spaces after command names like \ldots
                             % made from \ + letters.
                             %
                             % Note how a `%' character causes TeX to ignore
                             % the end of the input line, so these blank lines
                             % do not start a new paragraph.
                             %
with the right spacing around the periods requires
a special command.

\LaTeX\ interprets some common characters as
commands, so you must type special commands to
generate them.  These characters include the
following:
       \$ \& \% \# \{ and \}.

In printing, text is usually emphasized with an
       \emph{italic}
type style.

\begin{em}
   A long segment of text can also be emphasized
   in this way.  Text within such a segment can be
   given \emph{additional} emphasis.
\end{em}

It is sometimes necessary to prevent \LaTeX\ from
breaking a line where it might otherwise do so.
This may be at a space, as between the ``Mr.''\ and
``Jones'' in
       ``Mr.~Jones'',        % ~ produces an unbreakable interword space.
or within a word---especially when the word is a
symbol like
       \mbox{\emph{itemnum}}
that makes little sense when hyphenated across
lines.

Footnotes\footnote{This is an example of a footnote.}
pose no problem.

\LaTeX\ is good at typesetting mathematical formulas
like
       \( x-3y + z = 7 \)
or
       \( a_{1} > x^{2n} + y^{2n} > x' \)
or
       \( \ip{A}{B} = \sum_{i} a_{i} b_{i} \).
The spaces you type in a formula are
ignored.  Remember that a letter like
       $x$                   % $ ... $  and  \( ... \)  are equivalent
is a formula when it denotes a mathematical
symbol, and it should be typed as one.

\section{Displayed Text}

Text is displayed by indenting it from the left
margin.  Quotations are commonly displayed.  There
are short quotations
\begin{quote}
   This is a short quotation.  It consists of a
   single paragraph of text.  See how it is formatted.
\end{quote}
and longer ones.
\begin{quotation}
   This is a longer quotation.  It consists of two
   paragraphs of text, neither of which are
   particularly interesting.

   This is the second paragraph of the quotation.  It
   is just as dull as the first paragraph.
\end{quotation}
Another frequently-displayed structure is a list.
The following is an example of an \emph{itemized}
list.
\begin{itemize}
   \item This is the first item of an itemized list.
         Each item in the list is marked with a ``tick''.
         You don't have to worry about what kind of tick
         mark is used.

   \item This is the second item of the list.  It
         contains another list nested inside it.  The inner
         list is an \emph{enumerated} list.
         \begin{enumerate}
            \item This is the first item of an enumerated
                  list that is nested within the itemized list.

            \item This is the second item of the inner list.
                  \LaTeX\ allows you to nest lists deeper than
                  you really should.
         \end{enumerate}
         This is the rest of the second item of the outer
         list.  It is no more interesting than any other
         part of the item.
   \item This is the third item of the list.
\end{itemize}
You can even display poetry.
\begin{verse}
   There is an environment
    for verse \\             % The \\ command separates lines
   Whose features some poets % within a stanza.
   will curse.

                             % One or more blank lines separate stanzas.

   For instead of making\\
   Them do \emph{all} line breaking, \\
   It allows them to put too many words on a line when they'd rather be
   forced to be terse.
\end{verse}

Mathematical formulas may also be displayed.  A
displayed formula
is
one-line long; multiline
formulas require special formatting instructions.
   \[  \ip{\Gamma}{\psi'} = x'' + y^{2} + z_{i}^{n}\]
Don't start a paragraph with a displayed equation,
nor make one a paragraph by itself.

\end{document}               % End of document.

\section{Bench Huge Section 809}
% synthetic bench section 809
% This is a sample LaTeX input file.  (Version of 12 August 2004.)
%
% A '%' character causes TeX to ignore all remaining text on the line,
% and is used for comments like this one.

\documentclass{article}      % Specifies the document class

                             % The preamble begins here.
\title{An Example Document}  % Declares the document's title.
\author{Leslie Lamport}      % Declares the author's name.
\date{January 21, 1994}      % Deleting this command produces today's date.

\newcommand{\ip}[2]{(#1, #2)}
                             % Defines \ip{arg1}{arg2} to mean
                             % (arg1, arg2).

%\newcommand{\ip}[2]{\langle #1 | #2\rangle}
                             % This is an alternative definition of
                             % \ip that is commented out.

\begin{document}             % End of preamble and beginning of text.

\maketitle                   % Produces the title.

This is an example input file.  Comparing it with
the output it generates can show you how to
produce a simple document of your own.

\section{Ordinary Text}      % Produces section heading.  Lower-level
                             % sections are begun with similar
                             % \subsection and \subsubsection commands.

The ends  of words and sentences are marked
  by   spaces. It  doesn't matter how many
spaces    you type; one is as good as 100.  The
end of   a line counts as a space.

One   or more   blank lines denote the  end
of  a paragraph.

Since any number of consecutive spaces are treated
like a single one, the formatting of the input
file makes no difference to
      \LaTeX,                % The \LaTeX command generates the LaTeX logo.
but it makes a difference to you.  When you use
\LaTeX, making your input file as easy to read
as possible will be a great help as you write
your document and when you change it.  This sample
file shows how you can add comments to your own input
file.

Because printing is different from typewriting,
there are a number of things that you have to do
differently when preparing an input file than if
you were just typing the document directly.
Quotation marks like
       ``this''
have to be handled specially, as do quotes within
quotes:
       ``\,`this'            % \, separates the double and single quote.
        is what I just
        wrote, not  `that'\,''.

Dashes come in three sizes: an
       intra-word
dash, a medium dash for number ranges like
       1--2,
and a punctuation
       dash---like
this.

A sentence-ending space should be larger than the
space between words within a sentence.  You
sometimes have to type special commands in
conjunction with punctuation characters to get
this right, as in the following sentence.
       Gnats, gnus, etc.\ all  % `\ ' makes an inter-word space.
       begin with G\@.         % \@ marks end-of-sentence punctuation.
You should check the spaces after periods when
reading your output to make sure you haven't
forgotten any special cases.  Generating an
ellipsis
       \ldots\               % `\ ' is needed after `\ldots' because TeX
                             % ignores spaces after command names like \ldots
                             % made from \ + letters.
                             %
                             % Note how a `%' character causes TeX to ignore
                             % the end of the input line, so these blank lines
                             % do not start a new paragraph.
                             %
with the right spacing around the periods requires
a special command.

\LaTeX\ interprets some common characters as
commands, so you must type special commands to
generate them.  These characters include the
following:
       \$ \& \% \# \{ and \}.

In printing, text is usually emphasized with an
       \emph{italic}
type style.

\begin{em}
   A long segment of text can also be emphasized
   in this way.  Text within such a segment can be
   given \emph{additional} emphasis.
\end{em}

It is sometimes necessary to prevent \LaTeX\ from
breaking a line where it might otherwise do so.
This may be at a space, as between the ``Mr.''\ and
``Jones'' in
       ``Mr.~Jones'',        % ~ produces an unbreakable interword space.
or within a word---especially when the word is a
symbol like
       \mbox{\emph{itemnum}}
that makes little sense when hyphenated across
lines.

Footnotes\footnote{This is an example of a footnote.}
pose no problem.

\LaTeX\ is good at typesetting mathematical formulas
like
       \( x-3y + z = 7 \)
or
       \( a_{1} > x^{2n} + y^{2n} > x' \)
or
       \( \ip{A}{B} = \sum_{i} a_{i} b_{i} \).
The spaces you type in a formula are
ignored.  Remember that a letter like
       $x$                   % $ ... $  and  \( ... \)  are equivalent
is a formula when it denotes a mathematical
symbol, and it should be typed as one.

\section{Displayed Text}

Text is displayed by indenting it from the left
margin.  Quotations are commonly displayed.  There
are short quotations
\begin{quote}
   This is a short quotation.  It consists of a
   single paragraph of text.  See how it is formatted.
\end{quote}
and longer ones.
\begin{quotation}
   This is a longer quotation.  It consists of two
   paragraphs of text, neither of which are
   particularly interesting.

   This is the second paragraph of the quotation.  It
   is just as dull as the first paragraph.
\end{quotation}
Another frequently-displayed structure is a list.
The following is an example of an \emph{itemized}
list.
\begin{itemize}
   \item This is the first item of an itemized list.
         Each item in the list is marked with a ``tick''.
         You don't have to worry about what kind of tick
         mark is used.

   \item This is the second item of the list.  It
         contains another list nested inside it.  The inner
         list is an \emph{enumerated} list.
         \begin{enumerate}
            \item This is the first item of an enumerated
                  list that is nested within the itemized list.

            \item This is the second item of the inner list.
                  \LaTeX\ allows you to nest lists deeper than
                  you really should.
         \end{enumerate}
         This is the rest of the second item of the outer
         list.  It is no more interesting than any other
         part of the item.
   \item This is the third item of the list.
\end{itemize}
You can even display poetry.
\begin{verse}
   There is an environment
    for verse \\             % The \\ command separates lines
   Whose features some poets % within a stanza.
   will curse.

                             % One or more blank lines separate stanzas.

   For instead of making\\
   Them do \emph{all} line breaking, \\
   It allows them to put too many words on a line when they'd rather be
   forced to be terse.
\end{verse}

Mathematical formulas may also be displayed.  A
displayed formula
is
one-line long; multiline
formulas require special formatting instructions.
   \[  \ip{\Gamma}{\psi'} = x'' + y^{2} + z_{i}^{n}\]
Don't start a paragraph with a displayed equation,
nor make one a paragraph by itself.

\end{document}               % End of document.

\section{Bench Huge Section 810}
% synthetic bench section 810
% This is a sample LaTeX input file.  (Version of 12 August 2004.)
%
% A '%' character causes TeX to ignore all remaining text on the line,
% and is used for comments like this one.

\documentclass{article}      % Specifies the document class

                             % The preamble begins here.
\title{An Example Document}  % Declares the document's title.
\author{Leslie Lamport}      % Declares the author's name.
\date{January 21, 1994}      % Deleting this command produces today's date.

\newcommand{\ip}[2]{(#1, #2)}
                             % Defines \ip{arg1}{arg2} to mean
                             % (arg1, arg2).

%\newcommand{\ip}[2]{\langle #1 | #2\rangle}
                             % This is an alternative definition of
                             % \ip that is commented out.

\begin{document}             % End of preamble and beginning of text.

\maketitle                   % Produces the title.

This is an example input file.  Comparing it with
the output it generates can show you how to
produce a simple document of your own.

\section{Ordinary Text}      % Produces section heading.  Lower-level
                             % sections are begun with similar
                             % \subsection and \subsubsection commands.

The ends  of words and sentences are marked
  by   spaces. It  doesn't matter how many
spaces    you type; one is as good as 100.  The
end of   a line counts as a space.

One   or more   blank lines denote the  end
of  a paragraph.

Since any number of consecutive spaces are treated
like a single one, the formatting of the input
file makes no difference to
      \LaTeX,                % The \LaTeX command generates the LaTeX logo.
but it makes a difference to you.  When you use
\LaTeX, making your input file as easy to read
as possible will be a great help as you write
your document and when you change it.  This sample
file shows how you can add comments to your own input
file.

Because printing is different from typewriting,
there are a number of things that you have to do
differently when preparing an input file than if
you were just typing the document directly.
Quotation marks like
       ``this''
have to be handled specially, as do quotes within
quotes:
       ``\,`this'            % \, separates the double and single quote.
        is what I just
        wrote, not  `that'\,''.

Dashes come in three sizes: an
       intra-word
dash, a medium dash for number ranges like
       1--2,
and a punctuation
       dash---like
this.

A sentence-ending space should be larger than the
space between words within a sentence.  You
sometimes have to type special commands in
conjunction with punctuation characters to get
this right, as in the following sentence.
       Gnats, gnus, etc.\ all  % `\ ' makes an inter-word space.
       begin with G\@.         % \@ marks end-of-sentence punctuation.
You should check the spaces after periods when
reading your output to make sure you haven't
forgotten any special cases.  Generating an
ellipsis
       \ldots\               % `\ ' is needed after `\ldots' because TeX
                             % ignores spaces after command names like \ldots
                             % made from \ + letters.
                             %
                             % Note how a `%' character causes TeX to ignore
                             % the end of the input line, so these blank lines
                             % do not start a new paragraph.
                             %
with the right spacing around the periods requires
a special command.

\LaTeX\ interprets some common characters as
commands, so you must type special commands to
generate them.  These characters include the
following:
       \$ \& \% \# \{ and \}.

In printing, text is usually emphasized with an
       \emph{italic}
type style.

\begin{em}
   A long segment of text can also be emphasized
   in this way.  Text within such a segment can be
   given \emph{additional} emphasis.
\end{em}

It is sometimes necessary to prevent \LaTeX\ from
breaking a line where it might otherwise do so.
This may be at a space, as between the ``Mr.''\ and
``Jones'' in
       ``Mr.~Jones'',        % ~ produces an unbreakable interword space.
or within a word---especially when the word is a
symbol like
       \mbox{\emph{itemnum}}
that makes little sense when hyphenated across
lines.

Footnotes\footnote{This is an example of a footnote.}
pose no problem.

\LaTeX\ is good at typesetting mathematical formulas
like
       \( x-3y + z = 7 \)
or
       \( a_{1} > x^{2n} + y^{2n} > x' \)
or
       \( \ip{A}{B} = \sum_{i} a_{i} b_{i} \).
The spaces you type in a formula are
ignored.  Remember that a letter like
       $x$                   % $ ... $  and  \( ... \)  are equivalent
is a formula when it denotes a mathematical
symbol, and it should be typed as one.

\section{Displayed Text}

Text is displayed by indenting it from the left
margin.  Quotations are commonly displayed.  There
are short quotations
\begin{quote}
   This is a short quotation.  It consists of a
   single paragraph of text.  See how it is formatted.
\end{quote}
and longer ones.
\begin{quotation}
   This is a longer quotation.  It consists of two
   paragraphs of text, neither of which are
   particularly interesting.

   This is the second paragraph of the quotation.  It
   is just as dull as the first paragraph.
\end{quotation}
Another frequently-displayed structure is a list.
The following is an example of an \emph{itemized}
list.
\begin{itemize}
   \item This is the first item of an itemized list.
         Each item in the list is marked with a ``tick''.
         You don't have to worry about what kind of tick
         mark is used.

   \item This is the second item of the list.  It
         contains another list nested inside it.  The inner
         list is an \emph{enumerated} list.
         \begin{enumerate}
            \item This is the first item of an enumerated
                  list that is nested within the itemized list.

            \item This is the second item of the inner list.
                  \LaTeX\ allows you to nest lists deeper than
                  you really should.
         \end{enumerate}
         This is the rest of the second item of the outer
         list.  It is no more interesting than any other
         part of the item.
   \item This is the third item of the list.
\end{itemize}
You can even display poetry.
\begin{verse}
   There is an environment
    for verse \\             % The \\ command separates lines
   Whose features some poets % within a stanza.
   will curse.

                             % One or more blank lines separate stanzas.

   For instead of making\\
   Them do \emph{all} line breaking, \\
   It allows them to put too many words on a line when they'd rather be
   forced to be terse.
\end{verse}

Mathematical formulas may also be displayed.  A
displayed formula
is
one-line long; multiline
formulas require special formatting instructions.
   \[  \ip{\Gamma}{\psi'} = x'' + y^{2} + z_{i}^{n}\]
Don't start a paragraph with a displayed equation,
nor make one a paragraph by itself.

\end{document}               % End of document.

\section{Bench Huge Section 811}
% synthetic bench section 811
% This is a sample LaTeX input file.  (Version of 12 August 2004.)
%
% A '%' character causes TeX to ignore all remaining text on the line,
% and is used for comments like this one.

\documentclass{article}      % Specifies the document class

                             % The preamble begins here.
\title{An Example Document}  % Declares the document's title.
\author{Leslie Lamport}      % Declares the author's name.
\date{January 21, 1994}      % Deleting this command produces today's date.

\newcommand{\ip}[2]{(#1, #2)}
                             % Defines \ip{arg1}{arg2} to mean
                             % (arg1, arg2).

%\newcommand{\ip}[2]{\langle #1 | #2\rangle}
                             % This is an alternative definition of
                             % \ip that is commented out.

\begin{document}             % End of preamble and beginning of text.

\maketitle                   % Produces the title.

This is an example input file.  Comparing it with
the output it generates can show you how to
produce a simple document of your own.

\section{Ordinary Text}      % Produces section heading.  Lower-level
                             % sections are begun with similar
                             % \subsection and \subsubsection commands.

The ends  of words and sentences are marked
  by   spaces. It  doesn't matter how many
spaces    you type; one is as good as 100.  The
end of   a line counts as a space.

One   or more   blank lines denote the  end
of  a paragraph.

Since any number of consecutive spaces are treated
like a single one, the formatting of the input
file makes no difference to
      \LaTeX,                % The \LaTeX command generates the LaTeX logo.
but it makes a difference to you.  When you use
\LaTeX, making your input file as easy to read
as possible will be a great help as you write
your document and when you change it.  This sample
file shows how you can add comments to your own input
file.

Because printing is different from typewriting,
there are a number of things that you have to do
differently when preparing an input file than if
you were just typing the document directly.
Quotation marks like
       ``this''
have to be handled specially, as do quotes within
quotes:
       ``\,`this'            % \, separates the double and single quote.
        is what I just
        wrote, not  `that'\,''.

Dashes come in three sizes: an
       intra-word
dash, a medium dash for number ranges like
       1--2,
and a punctuation
       dash---like
this.

A sentence-ending space should be larger than the
space between words within a sentence.  You
sometimes have to type special commands in
conjunction with punctuation characters to get
this right, as in the following sentence.
       Gnats, gnus, etc.\ all  % `\ ' makes an inter-word space.
       begin with G\@.         % \@ marks end-of-sentence punctuation.
You should check the spaces after periods when
reading your output to make sure you haven't
forgotten any special cases.  Generating an
ellipsis
       \ldots\               % `\ ' is needed after `\ldots' because TeX
                             % ignores spaces after command names like \ldots
                             % made from \ + letters.
                             %
                             % Note how a `%' character causes TeX to ignore
                             % the end of the input line, so these blank lines
                             % do not start a new paragraph.
                             %
with the right spacing around the periods requires
a special command.

\LaTeX\ interprets some common characters as
commands, so you must type special commands to
generate them.  These characters include the
following:
       \$ \& \% \# \{ and \}.

In printing, text is usually emphasized with an
       \emph{italic}
type style.

\begin{em}
   A long segment of text can also be emphasized
   in this way.  Text within such a segment can be
   given \emph{additional} emphasis.
\end{em}

It is sometimes necessary to prevent \LaTeX\ from
breaking a line where it might otherwise do so.
This may be at a space, as between the ``Mr.''\ and
``Jones'' in
       ``Mr.~Jones'',        % ~ produces an unbreakable interword space.
or within a word---especially when the word is a
symbol like
       \mbox{\emph{itemnum}}
that makes little sense when hyphenated across
lines.

Footnotes\footnote{This is an example of a footnote.}
pose no problem.

\LaTeX\ is good at typesetting mathematical formulas
like
       \( x-3y + z = 7 \)
or
       \( a_{1} > x^{2n} + y^{2n} > x' \)
or
       \( \ip{A}{B} = \sum_{i} a_{i} b_{i} \).
The spaces you type in a formula are
ignored.  Remember that a letter like
       $x$                   % $ ... $  and  \( ... \)  are equivalent
is a formula when it denotes a mathematical
symbol, and it should be typed as one.

\section{Displayed Text}

Text is displayed by indenting it from the left
margin.  Quotations are commonly displayed.  There
are short quotations
\begin{quote}
   This is a short quotation.  It consists of a
   single paragraph of text.  See how it is formatted.
\end{quote}
and longer ones.
\begin{quotation}
   This is a longer quotation.  It consists of two
   paragraphs of text, neither of which are
   particularly interesting.

   This is the second paragraph of the quotation.  It
   is just as dull as the first paragraph.
\end{quotation}
Another frequently-displayed structure is a list.
The following is an example of an \emph{itemized}
list.
\begin{itemize}
   \item This is the first item of an itemized list.
         Each item in the list is marked with a ``tick''.
         You don't have to worry about what kind of tick
         mark is used.

   \item This is the second item of the list.  It
         contains another list nested inside it.  The inner
         list is an \emph{enumerated} list.
         \begin{enumerate}
            \item This is the first item of an enumerated
                  list that is nested within the itemized list.

            \item This is the second item of the inner list.
                  \LaTeX\ allows you to nest lists deeper than
                  you really should.
         \end{enumerate}
         This is the rest of the second item of the outer
         list.  It is no more interesting than any other
         part of the item.
   \item This is the third item of the list.
\end{itemize}
You can even display poetry.
\begin{verse}
   There is an environment
    for verse \\             % The \\ command separates lines
   Whose features some poets % within a stanza.
   will curse.

                             % One or more blank lines separate stanzas.

   For instead of making\\
   Them do \emph{all} line breaking, \\
   It allows them to put too many words on a line when they'd rather be
   forced to be terse.
\end{verse}

Mathematical formulas may also be displayed.  A
displayed formula
is
one-line long; multiline
formulas require special formatting instructions.
   \[  \ip{\Gamma}{\psi'} = x'' + y^{2} + z_{i}^{n}\]
Don't start a paragraph with a displayed equation,
nor make one a paragraph by itself.

\end{document}               % End of document.

\section{Bench Huge Section 812}
% synthetic bench section 812
% This is a sample LaTeX input file.  (Version of 12 August 2004.)
%
% A '%' character causes TeX to ignore all remaining text on the line,
% and is used for comments like this one.

\documentclass{article}      % Specifies the document class

                             % The preamble begins here.
\title{An Example Document}  % Declares the document's title.
\author{Leslie Lamport}      % Declares the author's name.
\date{January 21, 1994}      % Deleting this command produces today's date.

\newcommand{\ip}[2]{(#1, #2)}
                             % Defines \ip{arg1}{arg2} to mean
                             % (arg1, arg2).

%\newcommand{\ip}[2]{\langle #1 | #2\rangle}
                             % This is an alternative definition of
                             % \ip that is commented out.

\begin{document}             % End of preamble and beginning of text.

\maketitle                   % Produces the title.

This is an example input file.  Comparing it with
the output it generates can show you how to
produce a simple document of your own.

\section{Ordinary Text}      % Produces section heading.  Lower-level
                             % sections are begun with similar
                             % \subsection and \subsubsection commands.

The ends  of words and sentences are marked
  by   spaces. It  doesn't matter how many
spaces    you type; one is as good as 100.  The
end of   a line counts as a space.

One   or more   blank lines denote the  end
of  a paragraph.

Since any number of consecutive spaces are treated
like a single one, the formatting of the input
file makes no difference to
      \LaTeX,                % The \LaTeX command generates the LaTeX logo.
but it makes a difference to you.  When you use
\LaTeX, making your input file as easy to read
as possible will be a great help as you write
your document and when you change it.  This sample
file shows how you can add comments to your own input
file.

Because printing is different from typewriting,
there are a number of things that you have to do
differently when preparing an input file than if
you were just typing the document directly.
Quotation marks like
       ``this''
have to be handled specially, as do quotes within
quotes:
       ``\,`this'            % \, separates the double and single quote.
        is what I just
        wrote, not  `that'\,''.

Dashes come in three sizes: an
       intra-word
dash, a medium dash for number ranges like
       1--2,
and a punctuation
       dash---like
this.

A sentence-ending space should be larger than the
space between words within a sentence.  You
sometimes have to type special commands in
conjunction with punctuation characters to get
this right, as in the following sentence.
       Gnats, gnus, etc.\ all  % `\ ' makes an inter-word space.
       begin with G\@.         % \@ marks end-of-sentence punctuation.
You should check the spaces after periods when
reading your output to make sure you haven't
forgotten any special cases.  Generating an
ellipsis
       \ldots\               % `\ ' is needed after `\ldots' because TeX
                             % ignores spaces after command names like \ldots
                             % made from \ + letters.
                             %
                             % Note how a `%' character causes TeX to ignore
                             % the end of the input line, so these blank lines
                             % do not start a new paragraph.
                             %
with the right spacing around the periods requires
a special command.

\LaTeX\ interprets some common characters as
commands, so you must type special commands to
generate them.  These characters include the
following:
       \$ \& \% \# \{ and \}.

In printing, text is usually emphasized with an
       \emph{italic}
type style.

\begin{em}
   A long segment of text can also be emphasized
   in this way.  Text within such a segment can be
   given \emph{additional} emphasis.
\end{em}

It is sometimes necessary to prevent \LaTeX\ from
breaking a line where it might otherwise do so.
This may be at a space, as between the ``Mr.''\ and
``Jones'' in
       ``Mr.~Jones'',        % ~ produces an unbreakable interword space.
or within a word---especially when the word is a
symbol like
       \mbox{\emph{itemnum}}
that makes little sense when hyphenated across
lines.

Footnotes\footnote{This is an example of a footnote.}
pose no problem.

\LaTeX\ is good at typesetting mathematical formulas
like
       \( x-3y + z = 7 \)
or
       \( a_{1} > x^{2n} + y^{2n} > x' \)
or
       \( \ip{A}{B} = \sum_{i} a_{i} b_{i} \).
The spaces you type in a formula are
ignored.  Remember that a letter like
       $x$                   % $ ... $  and  \( ... \)  are equivalent
is a formula when it denotes a mathematical
symbol, and it should be typed as one.

\section{Displayed Text}

Text is displayed by indenting it from the left
margin.  Quotations are commonly displayed.  There
are short quotations
\begin{quote}
   This is a short quotation.  It consists of a
   single paragraph of text.  See how it is formatted.
\end{quote}
and longer ones.
\begin{quotation}
   This is a longer quotation.  It consists of two
   paragraphs of text, neither of which are
   particularly interesting.

   This is the second paragraph of the quotation.  It
   is just as dull as the first paragraph.
\end{quotation}
Another frequently-displayed structure is a list.
The following is an example of an \emph{itemized}
list.
\begin{itemize}
   \item This is the first item of an itemized list.
         Each item in the list is marked with a ``tick''.
         You don't have to worry about what kind of tick
         mark is used.

   \item This is the second item of the list.  It
         contains another list nested inside it.  The inner
         list is an \emph{enumerated} list.
         \begin{enumerate}
            \item This is the first item of an enumerated
                  list that is nested within the itemized list.

            \item This is the second item of the inner list.
                  \LaTeX\ allows you to nest lists deeper than
                  you really should.
         \end{enumerate}
         This is the rest of the second item of the outer
         list.  It is no more interesting than any other
         part of the item.
   \item This is the third item of the list.
\end{itemize}
You can even display poetry.
\begin{verse}
   There is an environment
    for verse \\             % The \\ command separates lines
   Whose features some poets % within a stanza.
   will curse.

                             % One or more blank lines separate stanzas.

   For instead of making\\
   Them do \emph{all} line breaking, \\
   It allows them to put too many words on a line when they'd rather be
   forced to be terse.
\end{verse}

Mathematical formulas may also be displayed.  A
displayed formula
is
one-line long; multiline
formulas require special formatting instructions.
   \[  \ip{\Gamma}{\psi'} = x'' + y^{2} + z_{i}^{n}\]
Don't start a paragraph with a displayed equation,
nor make one a paragraph by itself.

\end{document}               % End of document.

\section{Bench Huge Section 813}
% synthetic bench section 813
% This is a sample LaTeX input file.  (Version of 12 August 2004.)
%
% A '%' character causes TeX to ignore all remaining text on the line,
% and is used for comments like this one.

\documentclass{article}      % Specifies the document class

                             % The preamble begins here.
\title{An Example Document}  % Declares the document's title.
\author{Leslie Lamport}      % Declares the author's name.
\date{January 21, 1994}      % Deleting this command produces today's date.

\newcommand{\ip}[2]{(#1, #2)}
                             % Defines \ip{arg1}{arg2} to mean
                             % (arg1, arg2).

%\newcommand{\ip}[2]{\langle #1 | #2\rangle}
                             % This is an alternative definition of
                             % \ip that is commented out.

\begin{document}             % End of preamble and beginning of text.

\maketitle                   % Produces the title.

This is an example input file.  Comparing it with
the output it generates can show you how to
produce a simple document of your own.

\section{Ordinary Text}      % Produces section heading.  Lower-level
                             % sections are begun with similar
                             % \subsection and \subsubsection commands.

The ends  of words and sentences are marked
  by   spaces. It  doesn't matter how many
spaces    you type; one is as good as 100.  The
end of   a line counts as a space.

One   or more   blank lines denote the  end
of  a paragraph.

Since any number of consecutive spaces are treated
like a single one, the formatting of the input
file makes no difference to
      \LaTeX,                % The \LaTeX command generates the LaTeX logo.
but it makes a difference to you.  When you use
\LaTeX, making your input file as easy to read
as possible will be a great help as you write
your document and when you change it.  This sample
file shows how you can add comments to your own input
file.

Because printing is different from typewriting,
there are a number of things that you have to do
differently when preparing an input file than if
you were just typing the document directly.
Quotation marks like
       ``this''
have to be handled specially, as do quotes within
quotes:
       ``\,`this'            % \, separates the double and single quote.
        is what I just
        wrote, not  `that'\,''.

Dashes come in three sizes: an
       intra-word
dash, a medium dash for number ranges like
       1--2,
and a punctuation
       dash---like
this.

A sentence-ending space should be larger than the
space between words within a sentence.  You
sometimes have to type special commands in
conjunction with punctuation characters to get
this right, as in the following sentence.
       Gnats, gnus, etc.\ all  % `\ ' makes an inter-word space.
       begin with G\@.         % \@ marks end-of-sentence punctuation.
You should check the spaces after periods when
reading your output to make sure you haven't
forgotten any special cases.  Generating an
ellipsis
       \ldots\               % `\ ' is needed after `\ldots' because TeX
                             % ignores spaces after command names like \ldots
                             % made from \ + letters.
                             %
                             % Note how a `%' character causes TeX to ignore
                             % the end of the input line, so these blank lines
                             % do not start a new paragraph.
                             %
with the right spacing around the periods requires
a special command.

\LaTeX\ interprets some common characters as
commands, so you must type special commands to
generate them.  These characters include the
following:
       \$ \& \% \# \{ and \}.

In printing, text is usually emphasized with an
       \emph{italic}
type style.

\begin{em}
   A long segment of text can also be emphasized
   in this way.  Text within such a segment can be
   given \emph{additional} emphasis.
\end{em}

It is sometimes necessary to prevent \LaTeX\ from
breaking a line where it might otherwise do so.
This may be at a space, as between the ``Mr.''\ and
``Jones'' in
       ``Mr.~Jones'',        % ~ produces an unbreakable interword space.
or within a word---especially when the word is a
symbol like
       \mbox{\emph{itemnum}}
that makes little sense when hyphenated across
lines.

Footnotes\footnote{This is an example of a footnote.}
pose no problem.

\LaTeX\ is good at typesetting mathematical formulas
like
       \( x-3y + z = 7 \)
or
       \( a_{1} > x^{2n} + y^{2n} > x' \)
or
       \( \ip{A}{B} = \sum_{i} a_{i} b_{i} \).
The spaces you type in a formula are
ignored.  Remember that a letter like
       $x$                   % $ ... $  and  \( ... \)  are equivalent
is a formula when it denotes a mathematical
symbol, and it should be typed as one.

\section{Displayed Text}

Text is displayed by indenting it from the left
margin.  Quotations are commonly displayed.  There
are short quotations
\begin{quote}
   This is a short quotation.  It consists of a
   single paragraph of text.  See how it is formatted.
\end{quote}
and longer ones.
\begin{quotation}
   This is a longer quotation.  It consists of two
   paragraphs of text, neither of which are
   particularly interesting.

   This is the second paragraph of the quotation.  It
   is just as dull as the first paragraph.
\end{quotation}
Another frequently-displayed structure is a list.
The following is an example of an \emph{itemized}
list.
\begin{itemize}
   \item This is the first item of an itemized list.
         Each item in the list is marked with a ``tick''.
         You don't have to worry about what kind of tick
         mark is used.

   \item This is the second item of the list.  It
         contains another list nested inside it.  The inner
         list is an \emph{enumerated} list.
         \begin{enumerate}
            \item This is the first item of an enumerated
                  list that is nested within the itemized list.

            \item This is the second item of the inner list.
                  \LaTeX\ allows you to nest lists deeper than
                  you really should.
         \end{enumerate}
         This is the rest of the second item of the outer
         list.  It is no more interesting than any other
         part of the item.
   \item This is the third item of the list.
\end{itemize}
You can even display poetry.
\begin{verse}
   There is an environment
    for verse \\             % The \\ command separates lines
   Whose features some poets % within a stanza.
   will curse.

                             % One or more blank lines separate stanzas.

   For instead of making\\
   Them do \emph{all} line breaking, \\
   It allows them to put too many words on a line when they'd rather be
   forced to be terse.
\end{verse}

Mathematical formulas may also be displayed.  A
displayed formula
is
one-line long; multiline
formulas require special formatting instructions.
   \[  \ip{\Gamma}{\psi'} = x'' + y^{2} + z_{i}^{n}\]
Don't start a paragraph with a displayed equation,
nor make one a paragraph by itself.

\end{document}               % End of document.

\section{Bench Huge Section 814}
% synthetic bench section 814
% This is a sample LaTeX input file.  (Version of 12 August 2004.)
%
% A '%' character causes TeX to ignore all remaining text on the line,
% and is used for comments like this one.

\documentclass{article}      % Specifies the document class

                             % The preamble begins here.
\title{An Example Document}  % Declares the document's title.
\author{Leslie Lamport}      % Declares the author's name.
\date{January 21, 1994}      % Deleting this command produces today's date.

\newcommand{\ip}[2]{(#1, #2)}
                             % Defines \ip{arg1}{arg2} to mean
                             % (arg1, arg2).

%\newcommand{\ip}[2]{\langle #1 | #2\rangle}
                             % This is an alternative definition of
                             % \ip that is commented out.

\begin{document}             % End of preamble and beginning of text.

\maketitle                   % Produces the title.

This is an example input file.  Comparing it with
the output it generates can show you how to
produce a simple document of your own.

\section{Ordinary Text}      % Produces section heading.  Lower-level
                             % sections are begun with similar
                             % \subsection and \subsubsection commands.

The ends  of words and sentences are marked
  by   spaces. It  doesn't matter how many
spaces    you type; one is as good as 100.  The
end of   a line counts as a space.

One   or more   blank lines denote the  end
of  a paragraph.

Since any number of consecutive spaces are treated
like a single one, the formatting of the input
file makes no difference to
      \LaTeX,                % The \LaTeX command generates the LaTeX logo.
but it makes a difference to you.  When you use
\LaTeX, making your input file as easy to read
as possible will be a great help as you write
your document and when you change it.  This sample
file shows how you can add comments to your own input
file.

Because printing is different from typewriting,
there are a number of things that you have to do
differently when preparing an input file than if
you were just typing the document directly.
Quotation marks like
       ``this''
have to be handled specially, as do quotes within
quotes:
       ``\,`this'            % \, separates the double and single quote.
        is what I just
        wrote, not  `that'\,''.

Dashes come in three sizes: an
       intra-word
dash, a medium dash for number ranges like
       1--2,
and a punctuation
       dash---like
this.

A sentence-ending space should be larger than the
space between words within a sentence.  You
sometimes have to type special commands in
conjunction with punctuation characters to get
this right, as in the following sentence.
       Gnats, gnus, etc.\ all  % `\ ' makes an inter-word space.
       begin with G\@.         % \@ marks end-of-sentence punctuation.
You should check the spaces after periods when
reading your output to make sure you haven't
forgotten any special cases.  Generating an
ellipsis
       \ldots\               % `\ ' is needed after `\ldots' because TeX
                             % ignores spaces after command names like \ldots
                             % made from \ + letters.
                             %
                             % Note how a `%' character causes TeX to ignore
                             % the end of the input line, so these blank lines
                             % do not start a new paragraph.
                             %
with the right spacing around the periods requires
a special command.

\LaTeX\ interprets some common characters as
commands, so you must type special commands to
generate them.  These characters include the
following:
       \$ \& \% \# \{ and \}.

In printing, text is usually emphasized with an
       \emph{italic}
type style.

\begin{em}
   A long segment of text can also be emphasized
   in this way.  Text within such a segment can be
   given \emph{additional} emphasis.
\end{em}

It is sometimes necessary to prevent \LaTeX\ from
breaking a line where it might otherwise do so.
This may be at a space, as between the ``Mr.''\ and
``Jones'' in
       ``Mr.~Jones'',        % ~ produces an unbreakable interword space.
or within a word---especially when the word is a
symbol like
       \mbox{\emph{itemnum}}
that makes little sense when hyphenated across
lines.

Footnotes\footnote{This is an example of a footnote.}
pose no problem.

\LaTeX\ is good at typesetting mathematical formulas
like
       \( x-3y + z = 7 \)
or
       \( a_{1} > x^{2n} + y^{2n} > x' \)
or
       \( \ip{A}{B} = \sum_{i} a_{i} b_{i} \).
The spaces you type in a formula are
ignored.  Remember that a letter like
       $x$                   % $ ... $  and  \( ... \)  are equivalent
is a formula when it denotes a mathematical
symbol, and it should be typed as one.

\section{Displayed Text}

Text is displayed by indenting it from the left
margin.  Quotations are commonly displayed.  There
are short quotations
\begin{quote}
   This is a short quotation.  It consists of a
   single paragraph of text.  See how it is formatted.
\end{quote}
and longer ones.
\begin{quotation}
   This is a longer quotation.  It consists of two
   paragraphs of text, neither of which are
   particularly interesting.

   This is the second paragraph of the quotation.  It
   is just as dull as the first paragraph.
\end{quotation}
Another frequently-displayed structure is a list.
The following is an example of an \emph{itemized}
list.
\begin{itemize}
   \item This is the first item of an itemized list.
         Each item in the list is marked with a ``tick''.
         You don't have to worry about what kind of tick
         mark is used.

   \item This is the second item of the list.  It
         contains another list nested inside it.  The inner
         list is an \emph{enumerated} list.
         \begin{enumerate}
            \item This is the first item of an enumerated
                  list that is nested within the itemized list.

            \item This is the second item of the inner list.
                  \LaTeX\ allows you to nest lists deeper than
                  you really should.
         \end{enumerate}
         This is the rest of the second item of the outer
         list.  It is no more interesting than any other
         part of the item.
   \item This is the third item of the list.
\end{itemize}
You can even display poetry.
\begin{verse}
   There is an environment
    for verse \\             % The \\ command separates lines
   Whose features some poets % within a stanza.
   will curse.

                             % One or more blank lines separate stanzas.

   For instead of making\\
   Them do \emph{all} line breaking, \\
   It allows them to put too many words on a line when they'd rather be
   forced to be terse.
\end{verse}

Mathematical formulas may also be displayed.  A
displayed formula
is
one-line long; multiline
formulas require special formatting instructions.
   \[  \ip{\Gamma}{\psi'} = x'' + y^{2} + z_{i}^{n}\]
Don't start a paragraph with a displayed equation,
nor make one a paragraph by itself.

\end{document}               % End of document.

\section{Bench Huge Section 815}
% synthetic bench section 815
% This is a sample LaTeX input file.  (Version of 12 August 2004.)
%
% A '%' character causes TeX to ignore all remaining text on the line,
% and is used for comments like this one.

\documentclass{article}      % Specifies the document class

                             % The preamble begins here.
\title{An Example Document}  % Declares the document's title.
\author{Leslie Lamport}      % Declares the author's name.
\date{January 21, 1994}      % Deleting this command produces today's date.

\newcommand{\ip}[2]{(#1, #2)}
                             % Defines \ip{arg1}{arg2} to mean
                             % (arg1, arg2).

%\newcommand{\ip}[2]{\langle #1 | #2\rangle}
                             % This is an alternative definition of
                             % \ip that is commented out.

\begin{document}             % End of preamble and beginning of text.

\maketitle                   % Produces the title.

This is an example input file.  Comparing it with
the output it generates can show you how to
produce a simple document of your own.

\section{Ordinary Text}      % Produces section heading.  Lower-level
                             % sections are begun with similar
                             % \subsection and \subsubsection commands.

The ends  of words and sentences are marked
  by   spaces. It  doesn't matter how many
spaces    you type; one is as good as 100.  The
end of   a line counts as a space.

One   or more   blank lines denote the  end
of  a paragraph.

Since any number of consecutive spaces are treated
like a single one, the formatting of the input
file makes no difference to
      \LaTeX,                % The \LaTeX command generates the LaTeX logo.
but it makes a difference to you.  When you use
\LaTeX, making your input file as easy to read
as possible will be a great help as you write
your document and when you change it.  This sample
file shows how you can add comments to your own input
file.

Because printing is different from typewriting,
there are a number of things that you have to do
differently when preparing an input file than if
you were just typing the document directly.
Quotation marks like
       ``this''
have to be handled specially, as do quotes within
quotes:
       ``\,`this'            % \, separates the double and single quote.
        is what I just
        wrote, not  `that'\,''.

Dashes come in three sizes: an
       intra-word
dash, a medium dash for number ranges like
       1--2,
and a punctuation
       dash---like
this.

A sentence-ending space should be larger than the
space between words within a sentence.  You
sometimes have to type special commands in
conjunction with punctuation characters to get
this right, as in the following sentence.
       Gnats, gnus, etc.\ all  % `\ ' makes an inter-word space.
       begin with G\@.         % \@ marks end-of-sentence punctuation.
You should check the spaces after periods when
reading your output to make sure you haven't
forgotten any special cases.  Generating an
ellipsis
       \ldots\               % `\ ' is needed after `\ldots' because TeX
                             % ignores spaces after command names like \ldots
                             % made from \ + letters.
                             %
                             % Note how a `%' character causes TeX to ignore
                             % the end of the input line, so these blank lines
                             % do not start a new paragraph.
                             %
with the right spacing around the periods requires
a special command.

\LaTeX\ interprets some common characters as
commands, so you must type special commands to
generate them.  These characters include the
following:
       \$ \& \% \# \{ and \}.

In printing, text is usually emphasized with an
       \emph{italic}
type style.

\begin{em}
   A long segment of text can also be emphasized
   in this way.  Text within such a segment can be
   given \emph{additional} emphasis.
\end{em}

It is sometimes necessary to prevent \LaTeX\ from
breaking a line where it might otherwise do so.
This may be at a space, as between the ``Mr.''\ and
``Jones'' in
       ``Mr.~Jones'',        % ~ produces an unbreakable interword space.
or within a word---especially when the word is a
symbol like
       \mbox{\emph{itemnum}}
that makes little sense when hyphenated across
lines.

Footnotes\footnote{This is an example of a footnote.}
pose no problem.

\LaTeX\ is good at typesetting mathematical formulas
like
       \( x-3y + z = 7 \)
or
       \( a_{1} > x^{2n} + y^{2n} > x' \)
or
       \( \ip{A}{B} = \sum_{i} a_{i} b_{i} \).
The spaces you type in a formula are
ignored.  Remember that a letter like
       $x$                   % $ ... $  and  \( ... \)  are equivalent
is a formula when it denotes a mathematical
symbol, and it should be typed as one.

\section{Displayed Text}

Text is displayed by indenting it from the left
margin.  Quotations are commonly displayed.  There
are short quotations
\begin{quote}
   This is a short quotation.  It consists of a
   single paragraph of text.  See how it is formatted.
\end{quote}
and longer ones.
\begin{quotation}
   This is a longer quotation.  It consists of two
   paragraphs of text, neither of which are
   particularly interesting.

   This is the second paragraph of the quotation.  It
   is just as dull as the first paragraph.
\end{quotation}
Another frequently-displayed structure is a list.
The following is an example of an \emph{itemized}
list.
\begin{itemize}
   \item This is the first item of an itemized list.
         Each item in the list is marked with a ``tick''.
         You don't have to worry about what kind of tick
         mark is used.

   \item This is the second item of the list.  It
         contains another list nested inside it.  The inner
         list is an \emph{enumerated} list.
         \begin{enumerate}
            \item This is the first item of an enumerated
                  list that is nested within the itemized list.

            \item This is the second item of the inner list.
                  \LaTeX\ allows you to nest lists deeper than
                  you really should.
         \end{enumerate}
         This is the rest of the second item of the outer
         list.  It is no more interesting than any other
         part of the item.
   \item This is the third item of the list.
\end{itemize}
You can even display poetry.
\begin{verse}
   There is an environment
    for verse \\             % The \\ command separates lines
   Whose features some poets % within a stanza.
   will curse.

                             % One or more blank lines separate stanzas.

   For instead of making\\
   Them do \emph{all} line breaking, \\
   It allows them to put too many words on a line when they'd rather be
   forced to be terse.
\end{verse}

Mathematical formulas may also be displayed.  A
displayed formula
is
one-line long; multiline
formulas require special formatting instructions.
   \[  \ip{\Gamma}{\psi'} = x'' + y^{2} + z_{i}^{n}\]
Don't start a paragraph with a displayed equation,
nor make one a paragraph by itself.

\end{document}               % End of document.

\section{Bench Huge Section 816}
% synthetic bench section 816
% This is a sample LaTeX input file.  (Version of 12 August 2004.)
%
% A '%' character causes TeX to ignore all remaining text on the line,
% and is used for comments like this one.

\documentclass{article}      % Specifies the document class

                             % The preamble begins here.
\title{An Example Document}  % Declares the document's title.
\author{Leslie Lamport}      % Declares the author's name.
\date{January 21, 1994}      % Deleting this command produces today's date.

\newcommand{\ip}[2]{(#1, #2)}
                             % Defines \ip{arg1}{arg2} to mean
                             % (arg1, arg2).

%\newcommand{\ip}[2]{\langle #1 | #2\rangle}
                             % This is an alternative definition of
                             % \ip that is commented out.

\begin{document}             % End of preamble and beginning of text.

\maketitle                   % Produces the title.

This is an example input file.  Comparing it with
the output it generates can show you how to
produce a simple document of your own.

\section{Ordinary Text}      % Produces section heading.  Lower-level
                             % sections are begun with similar
                             % \subsection and \subsubsection commands.

The ends  of words and sentences are marked
  by   spaces. It  doesn't matter how many
spaces    you type; one is as good as 100.  The
end of   a line counts as a space.

One   or more   blank lines denote the  end
of  a paragraph.

Since any number of consecutive spaces are treated
like a single one, the formatting of the input
file makes no difference to
      \LaTeX,                % The \LaTeX command generates the LaTeX logo.
but it makes a difference to you.  When you use
\LaTeX, making your input file as easy to read
as possible will be a great help as you write
your document and when you change it.  This sample
file shows how you can add comments to your own input
file.

Because printing is different from typewriting,
there are a number of things that you have to do
differently when preparing an input file than if
you were just typing the document directly.
Quotation marks like
       ``this''
have to be handled specially, as do quotes within
quotes:
       ``\,`this'            % \, separates the double and single quote.
        is what I just
        wrote, not  `that'\,''.

Dashes come in three sizes: an
       intra-word
dash, a medium dash for number ranges like
       1--2,
and a punctuation
       dash---like
this.

A sentence-ending space should be larger than the
space between words within a sentence.  You
sometimes have to type special commands in
conjunction with punctuation characters to get
this right, as in the following sentence.
       Gnats, gnus, etc.\ all  % `\ ' makes an inter-word space.
       begin with G\@.         % \@ marks end-of-sentence punctuation.
You should check the spaces after periods when
reading your output to make sure you haven't
forgotten any special cases.  Generating an
ellipsis
       \ldots\               % `\ ' is needed after `\ldots' because TeX
                             % ignores spaces after command names like \ldots
                             % made from \ + letters.
                             %
                             % Note how a `%' character causes TeX to ignore
                             % the end of the input line, so these blank lines
                             % do not start a new paragraph.
                             %
with the right spacing around the periods requires
a special command.

\LaTeX\ interprets some common characters as
commands, so you must type special commands to
generate them.  These characters include the
following:
       \$ \& \% \# \{ and \}.

In printing, text is usually emphasized with an
       \emph{italic}
type style.

\begin{em}
   A long segment of text can also be emphasized
   in this way.  Text within such a segment can be
   given \emph{additional} emphasis.
\end{em}

It is sometimes necessary to prevent \LaTeX\ from
breaking a line where it might otherwise do so.
This may be at a space, as between the ``Mr.''\ and
``Jones'' in
       ``Mr.~Jones'',        % ~ produces an unbreakable interword space.
or within a word---especially when the word is a
symbol like
       \mbox{\emph{itemnum}}
that makes little sense when hyphenated across
lines.

Footnotes\footnote{This is an example of a footnote.}
pose no problem.

\LaTeX\ is good at typesetting mathematical formulas
like
       \( x-3y + z = 7 \)
or
       \( a_{1} > x^{2n} + y^{2n} > x' \)
or
       \( \ip{A}{B} = \sum_{i} a_{i} b_{i} \).
The spaces you type in a formula are
ignored.  Remember that a letter like
       $x$                   % $ ... $  and  \( ... \)  are equivalent
is a formula when it denotes a mathematical
symbol, and it should be typed as one.

\section{Displayed Text}

Text is displayed by indenting it from the left
margin.  Quotations are commonly displayed.  There
are short quotations
\begin{quote}
   This is a short quotation.  It consists of a
   single paragraph of text.  See how it is formatted.
\end{quote}
and longer ones.
\begin{quotation}
   This is a longer quotation.  It consists of two
   paragraphs of text, neither of which are
   particularly interesting.

   This is the second paragraph of the quotation.  It
   is just as dull as the first paragraph.
\end{quotation}
Another frequently-displayed structure is a list.
The following is an example of an \emph{itemized}
list.
\begin{itemize}
   \item This is the first item of an itemized list.
         Each item in the list is marked with a ``tick''.
         You don't have to worry about what kind of tick
         mark is used.

   \item This is the second item of the list.  It
         contains another list nested inside it.  The inner
         list is an \emph{enumerated} list.
         \begin{enumerate}
            \item This is the first item of an enumerated
                  list that is nested within the itemized list.

            \item This is the second item of the inner list.
                  \LaTeX\ allows you to nest lists deeper than
                  you really should.
         \end{enumerate}
         This is the rest of the second item of the outer
         list.  It is no more interesting than any other
         part of the item.
   \item This is the third item of the list.
\end{itemize}
You can even display poetry.
\begin{verse}
   There is an environment
    for verse \\             % The \\ command separates lines
   Whose features some poets % within a stanza.
   will curse.

                             % One or more blank lines separate stanzas.

   For instead of making\\
   Them do \emph{all} line breaking, \\
   It allows them to put too many words on a line when they'd rather be
   forced to be terse.
\end{verse}

Mathematical formulas may also be displayed.  A
displayed formula
is
one-line long; multiline
formulas require special formatting instructions.
   \[  \ip{\Gamma}{\psi'} = x'' + y^{2} + z_{i}^{n}\]
Don't start a paragraph with a displayed equation,
nor make one a paragraph by itself.

\end{document}               % End of document.

\section{Bench Huge Section 817}
% synthetic bench section 817
% This is a sample LaTeX input file.  (Version of 12 August 2004.)
%
% A '%' character causes TeX to ignore all remaining text on the line,
% and is used for comments like this one.

\documentclass{article}      % Specifies the document class

                             % The preamble begins here.
\title{An Example Document}  % Declares the document's title.
\author{Leslie Lamport}      % Declares the author's name.
\date{January 21, 1994}      % Deleting this command produces today's date.

\newcommand{\ip}[2]{(#1, #2)}
                             % Defines \ip{arg1}{arg2} to mean
                             % (arg1, arg2).

%\newcommand{\ip}[2]{\langle #1 | #2\rangle}
                             % This is an alternative definition of
                             % \ip that is commented out.

\begin{document}             % End of preamble and beginning of text.

\maketitle                   % Produces the title.

This is an example input file.  Comparing it with
the output it generates can show you how to
produce a simple document of your own.

\section{Ordinary Text}      % Produces section heading.  Lower-level
                             % sections are begun with similar
                             % \subsection and \subsubsection commands.

The ends  of words and sentences are marked
  by   spaces. It  doesn't matter how many
spaces    you type; one is as good as 100.  The
end of   a line counts as a space.

One   or more   blank lines denote the  end
of  a paragraph.

Since any number of consecutive spaces are treated
like a single one, the formatting of the input
file makes no difference to
      \LaTeX,                % The \LaTeX command generates the LaTeX logo.
but it makes a difference to you.  When you use
\LaTeX, making your input file as easy to read
as possible will be a great help as you write
your document and when you change it.  This sample
file shows how you can add comments to your own input
file.

Because printing is different from typewriting,
there are a number of things that you have to do
differently when preparing an input file than if
you were just typing the document directly.
Quotation marks like
       ``this''
have to be handled specially, as do quotes within
quotes:
       ``\,`this'            % \, separates the double and single quote.
        is what I just
        wrote, not  `that'\,''.

Dashes come in three sizes: an
       intra-word
dash, a medium dash for number ranges like
       1--2,
and a punctuation
       dash---like
this.

A sentence-ending space should be larger than the
space between words within a sentence.  You
sometimes have to type special commands in
conjunction with punctuation characters to get
this right, as in the following sentence.
       Gnats, gnus, etc.\ all  % `\ ' makes an inter-word space.
       begin with G\@.         % \@ marks end-of-sentence punctuation.
You should check the spaces after periods when
reading your output to make sure you haven't
forgotten any special cases.  Generating an
ellipsis
       \ldots\               % `\ ' is needed after `\ldots' because TeX
                             % ignores spaces after command names like \ldots
                             % made from \ + letters.
                             %
                             % Note how a `%' character causes TeX to ignore
                             % the end of the input line, so these blank lines
                             % do not start a new paragraph.
                             %
with the right spacing around the periods requires
a special command.

\LaTeX\ interprets some common characters as
commands, so you must type special commands to
generate them.  These characters include the
following:
       \$ \& \% \# \{ and \}.

In printing, text is usually emphasized with an
       \emph{italic}
type style.

\begin{em}
   A long segment of text can also be emphasized
   in this way.  Text within such a segment can be
   given \emph{additional} emphasis.
\end{em}

It is sometimes necessary to prevent \LaTeX\ from
breaking a line where it might otherwise do so.
This may be at a space, as between the ``Mr.''\ and
``Jones'' in
       ``Mr.~Jones'',        % ~ produces an unbreakable interword space.
or within a word---especially when the word is a
symbol like
       \mbox{\emph{itemnum}}
that makes little sense when hyphenated across
lines.

Footnotes\footnote{This is an example of a footnote.}
pose no problem.

\LaTeX\ is good at typesetting mathematical formulas
like
       \( x-3y + z = 7 \)
or
       \( a_{1} > x^{2n} + y^{2n} > x' \)
or
       \( \ip{A}{B} = \sum_{i} a_{i} b_{i} \).
The spaces you type in a formula are
ignored.  Remember that a letter like
       $x$                   % $ ... $  and  \( ... \)  are equivalent
is a formula when it denotes a mathematical
symbol, and it should be typed as one.

\section{Displayed Text}

Text is displayed by indenting it from the left
margin.  Quotations are commonly displayed.  There
are short quotations
\begin{quote}
   This is a short quotation.  It consists of a
   single paragraph of text.  See how it is formatted.
\end{quote}
and longer ones.
\begin{quotation}
   This is a longer quotation.  It consists of two
   paragraphs of text, neither of which are
   particularly interesting.

   This is the second paragraph of the quotation.  It
   is just as dull as the first paragraph.
\end{quotation}
Another frequently-displayed structure is a list.
The following is an example of an \emph{itemized}
list.
\begin{itemize}
   \item This is the first item of an itemized list.
         Each item in the list is marked with a ``tick''.
         You don't have to worry about what kind of tick
         mark is used.

   \item This is the second item of the list.  It
         contains another list nested inside it.  The inner
         list is an \emph{enumerated} list.
         \begin{enumerate}
            \item This is the first item of an enumerated
                  list that is nested within the itemized list.

            \item This is the second item of the inner list.
                  \LaTeX\ allows you to nest lists deeper than
                  you really should.
         \end{enumerate}
         This is the rest of the second item of the outer
         list.  It is no more interesting than any other
         part of the item.
   \item This is the third item of the list.
\end{itemize}
You can even display poetry.
\begin{verse}
   There is an environment
    for verse \\             % The \\ command separates lines
   Whose features some poets % within a stanza.
   will curse.

                             % One or more blank lines separate stanzas.

   For instead of making\\
   Them do \emph{all} line breaking, \\
   It allows them to put too many words on a line when they'd rather be
   forced to be terse.
\end{verse}

Mathematical formulas may also be displayed.  A
displayed formula
is
one-line long; multiline
formulas require special formatting instructions.
   \[  \ip{\Gamma}{\psi'} = x'' + y^{2} + z_{i}^{n}\]
Don't start a paragraph with a displayed equation,
nor make one a paragraph by itself.

\end{document}               % End of document.

\section{Bench Huge Section 818}
% synthetic bench section 818
% This is a sample LaTeX input file.  (Version of 12 August 2004.)
%
% A '%' character causes TeX to ignore all remaining text on the line,
% and is used for comments like this one.

\documentclass{article}      % Specifies the document class

                             % The preamble begins here.
\title{An Example Document}  % Declares the document's title.
\author{Leslie Lamport}      % Declares the author's name.
\date{January 21, 1994}      % Deleting this command produces today's date.

\newcommand{\ip}[2]{(#1, #2)}
                             % Defines \ip{arg1}{arg2} to mean
                             % (arg1, arg2).

%\newcommand{\ip}[2]{\langle #1 | #2\rangle}
                             % This is an alternative definition of
                             % \ip that is commented out.

\begin{document}             % End of preamble and beginning of text.

\maketitle                   % Produces the title.

This is an example input file.  Comparing it with
the output it generates can show you how to
produce a simple document of your own.

\section{Ordinary Text}      % Produces section heading.  Lower-level
                             % sections are begun with similar
                             % \subsection and \subsubsection commands.

The ends  of words and sentences are marked
  by   spaces. It  doesn't matter how many
spaces    you type; one is as good as 100.  The
end of   a line counts as a space.

One   or more   blank lines denote the  end
of  a paragraph.

Since any number of consecutive spaces are treated
like a single one, the formatting of the input
file makes no difference to
      \LaTeX,                % The \LaTeX command generates the LaTeX logo.
but it makes a difference to you.  When you use
\LaTeX, making your input file as easy to read
as possible will be a great help as you write
your document and when you change it.  This sample
file shows how you can add comments to your own input
file.

Because printing is different from typewriting,
there are a number of things that you have to do
differently when preparing an input file than if
you were just typing the document directly.
Quotation marks like
       ``this''
have to be handled specially, as do quotes within
quotes:
       ``\,`this'            % \, separates the double and single quote.
        is what I just
        wrote, not  `that'\,''.

Dashes come in three sizes: an
       intra-word
dash, a medium dash for number ranges like
       1--2,
and a punctuation
       dash---like
this.

A sentence-ending space should be larger than the
space between words within a sentence.  You
sometimes have to type special commands in
conjunction with punctuation characters to get
this right, as in the following sentence.
       Gnats, gnus, etc.\ all  % `\ ' makes an inter-word space.
       begin with G\@.         % \@ marks end-of-sentence punctuation.
You should check the spaces after periods when
reading your output to make sure you haven't
forgotten any special cases.  Generating an
ellipsis
       \ldots\               % `\ ' is needed after `\ldots' because TeX
                             % ignores spaces after command names like \ldots
                             % made from \ + letters.
                             %
                             % Note how a `%' character causes TeX to ignore
                             % the end of the input line, so these blank lines
                             % do not start a new paragraph.
                             %
with the right spacing around the periods requires
a special command.

\LaTeX\ interprets some common characters as
commands, so you must type special commands to
generate them.  These characters include the
following:
       \$ \& \% \# \{ and \}.

In printing, text is usually emphasized with an
       \emph{italic}
type style.

\begin{em}
   A long segment of text can also be emphasized
   in this way.  Text within such a segment can be
   given \emph{additional} emphasis.
\end{em}

It is sometimes necessary to prevent \LaTeX\ from
breaking a line where it might otherwise do so.
This may be at a space, as between the ``Mr.''\ and
``Jones'' in
       ``Mr.~Jones'',        % ~ produces an unbreakable interword space.
or within a word---especially when the word is a
symbol like
       \mbox{\emph{itemnum}}
that makes little sense when hyphenated across
lines.

Footnotes\footnote{This is an example of a footnote.}
pose no problem.

\LaTeX\ is good at typesetting mathematical formulas
like
       \( x-3y + z = 7 \)
or
       \( a_{1} > x^{2n} + y^{2n} > x' \)
or
       \( \ip{A}{B} = \sum_{i} a_{i} b_{i} \).
The spaces you type in a formula are
ignored.  Remember that a letter like
       $x$                   % $ ... $  and  \( ... \)  are equivalent
is a formula when it denotes a mathematical
symbol, and it should be typed as one.

\section{Displayed Text}

Text is displayed by indenting it from the left
margin.  Quotations are commonly displayed.  There
are short quotations
\begin{quote}
   This is a short quotation.  It consists of a
   single paragraph of text.  See how it is formatted.
\end{quote}
and longer ones.
\begin{quotation}
   This is a longer quotation.  It consists of two
   paragraphs of text, neither of which are
   particularly interesting.

   This is the second paragraph of the quotation.  It
   is just as dull as the first paragraph.
\end{quotation}
Another frequently-displayed structure is a list.
The following is an example of an \emph{itemized}
list.
\begin{itemize}
   \item This is the first item of an itemized list.
         Each item in the list is marked with a ``tick''.
         You don't have to worry about what kind of tick
         mark is used.

   \item This is the second item of the list.  It
         contains another list nested inside it.  The inner
         list is an \emph{enumerated} list.
         \begin{enumerate}
            \item This is the first item of an enumerated
                  list that is nested within the itemized list.

            \item This is the second item of the inner list.
                  \LaTeX\ allows you to nest lists deeper than
                  you really should.
         \end{enumerate}
         This is the rest of the second item of the outer
         list.  It is no more interesting than any other
         part of the item.
   \item This is the third item of the list.
\end{itemize}
You can even display poetry.
\begin{verse}
   There is an environment
    for verse \\             % The \\ command separates lines
   Whose features some poets % within a stanza.
   will curse.

                             % One or more blank lines separate stanzas.

   For instead of making\\
   Them do \emph{all} line breaking, \\
   It allows them to put too many words on a line when they'd rather be
   forced to be terse.
\end{verse}

Mathematical formulas may also be displayed.  A
displayed formula
is
one-line long; multiline
formulas require special formatting instructions.
   \[  \ip{\Gamma}{\psi'} = x'' + y^{2} + z_{i}^{n}\]
Don't start a paragraph with a displayed equation,
nor make one a paragraph by itself.

\end{document}               % End of document.

\section{Bench Huge Section 819}
% synthetic bench section 819
% This is a sample LaTeX input file.  (Version of 12 August 2004.)
%
% A '%' character causes TeX to ignore all remaining text on the line,
% and is used for comments like this one.

\documentclass{article}      % Specifies the document class

                             % The preamble begins here.
\title{An Example Document}  % Declares the document's title.
\author{Leslie Lamport}      % Declares the author's name.
\date{January 21, 1994}      % Deleting this command produces today's date.

\newcommand{\ip}[2]{(#1, #2)}
                             % Defines \ip{arg1}{arg2} to mean
                             % (arg1, arg2).

%\newcommand{\ip}[2]{\langle #1 | #2\rangle}
                             % This is an alternative definition of
                             % \ip that is commented out.

\begin{document}             % End of preamble and beginning of text.

\maketitle                   % Produces the title.

This is an example input file.  Comparing it with
the output it generates can show you how to
produce a simple document of your own.

\section{Ordinary Text}      % Produces section heading.  Lower-level
                             % sections are begun with similar
                             % \subsection and \subsubsection commands.

The ends  of words and sentences are marked
  by   spaces. It  doesn't matter how many
spaces    you type; one is as good as 100.  The
end of   a line counts as a space.

One   or more   blank lines denote the  end
of  a paragraph.

Since any number of consecutive spaces are treated
like a single one, the formatting of the input
file makes no difference to
      \LaTeX,                % The \LaTeX command generates the LaTeX logo.
but it makes a difference to you.  When you use
\LaTeX, making your input file as easy to read
as possible will be a great help as you write
your document and when you change it.  This sample
file shows how you can add comments to your own input
file.

Because printing is different from typewriting,
there are a number of things that you have to do
differently when preparing an input file than if
you were just typing the document directly.
Quotation marks like
       ``this''
have to be handled specially, as do quotes within
quotes:
       ``\,`this'            % \, separates the double and single quote.
        is what I just
        wrote, not  `that'\,''.

Dashes come in three sizes: an
       intra-word
dash, a medium dash for number ranges like
       1--2,
and a punctuation
       dash---like
this.

A sentence-ending space should be larger than the
space between words within a sentence.  You
sometimes have to type special commands in
conjunction with punctuation characters to get
this right, as in the following sentence.
       Gnats, gnus, etc.\ all  % `\ ' makes an inter-word space.
       begin with G\@.         % \@ marks end-of-sentence punctuation.
You should check the spaces after periods when
reading your output to make sure you haven't
forgotten any special cases.  Generating an
ellipsis
       \ldots\               % `\ ' is needed after `\ldots' because TeX
                             % ignores spaces after command names like \ldots
                             % made from \ + letters.
                             %
                             % Note how a `%' character causes TeX to ignore
                             % the end of the input line, so these blank lines
                             % do not start a new paragraph.
                             %
with the right spacing around the periods requires
a special command.

\LaTeX\ interprets some common characters as
commands, so you must type special commands to
generate them.  These characters include the
following:
       \$ \& \% \# \{ and \}.

In printing, text is usually emphasized with an
       \emph{italic}
type style.

\begin{em}
   A long segment of text can also be emphasized
   in this way.  Text within such a segment can be
   given \emph{additional} emphasis.
\end{em}

It is sometimes necessary to prevent \LaTeX\ from
breaking a line where it might otherwise do so.
This may be at a space, as between the ``Mr.''\ and
``Jones'' in
       ``Mr.~Jones'',        % ~ produces an unbreakable interword space.
or within a word---especially when the word is a
symbol like
       \mbox{\emph{itemnum}}
that makes little sense when hyphenated across
lines.

Footnotes\footnote{This is an example of a footnote.}
pose no problem.

\LaTeX\ is good at typesetting mathematical formulas
like
       \( x-3y + z = 7 \)
or
       \( a_{1} > x^{2n} + y^{2n} > x' \)
or
       \( \ip{A}{B} = \sum_{i} a_{i} b_{i} \).
The spaces you type in a formula are
ignored.  Remember that a letter like
       $x$                   % $ ... $  and  \( ... \)  are equivalent
is a formula when it denotes a mathematical
symbol, and it should be typed as one.

\section{Displayed Text}

Text is displayed by indenting it from the left
margin.  Quotations are commonly displayed.  There
are short quotations
\begin{quote}
   This is a short quotation.  It consists of a
   single paragraph of text.  See how it is formatted.
\end{quote}
and longer ones.
\begin{quotation}
   This is a longer quotation.  It consists of two
   paragraphs of text, neither of which are
   particularly interesting.

   This is the second paragraph of the quotation.  It
   is just as dull as the first paragraph.
\end{quotation}
Another frequently-displayed structure is a list.
The following is an example of an \emph{itemized}
list.
\begin{itemize}
   \item This is the first item of an itemized list.
         Each item in the list is marked with a ``tick''.
         You don't have to worry about what kind of tick
         mark is used.

   \item This is the second item of the list.  It
         contains another list nested inside it.  The inner
         list is an \emph{enumerated} list.
         \begin{enumerate}
            \item This is the first item of an enumerated
                  list that is nested within the itemized list.

            \item This is the second item of the inner list.
                  \LaTeX\ allows you to nest lists deeper than
                  you really should.
         \end{enumerate}
         This is the rest of the second item of the outer
         list.  It is no more interesting than any other
         part of the item.
   \item This is the third item of the list.
\end{itemize}
You can even display poetry.
\begin{verse}
   There is an environment
    for verse \\             % The \\ command separates lines
   Whose features some poets % within a stanza.
   will curse.

                             % One or more blank lines separate stanzas.

   For instead of making\\
   Them do \emph{all} line breaking, \\
   It allows them to put too many words on a line when they'd rather be
   forced to be terse.
\end{verse}

Mathematical formulas may also be displayed.  A
displayed formula
is
one-line long; multiline
formulas require special formatting instructions.
   \[  \ip{\Gamma}{\psi'} = x'' + y^{2} + z_{i}^{n}\]
Don't start a paragraph with a displayed equation,
nor make one a paragraph by itself.

\end{document}               % End of document.

\section{Bench Huge Section 820}
% synthetic bench section 820
% This is a sample LaTeX input file.  (Version of 12 August 2004.)
%
% A '%' character causes TeX to ignore all remaining text on the line,
% and is used for comments like this one.

\documentclass{article}      % Specifies the document class

                             % The preamble begins here.
\title{An Example Document}  % Declares the document's title.
\author{Leslie Lamport}      % Declares the author's name.
\date{January 21, 1994}      % Deleting this command produces today's date.

\newcommand{\ip}[2]{(#1, #2)}
                             % Defines \ip{arg1}{arg2} to mean
                             % (arg1, arg2).

%\newcommand{\ip}[2]{\langle #1 | #2\rangle}
                             % This is an alternative definition of
                             % \ip that is commented out.

\begin{document}             % End of preamble and beginning of text.

\maketitle                   % Produces the title.

This is an example input file.  Comparing it with
the output it generates can show you how to
produce a simple document of your own.

\section{Ordinary Text}      % Produces section heading.  Lower-level
                             % sections are begun with similar
                             % \subsection and \subsubsection commands.

The ends  of words and sentences are marked
  by   spaces. It  doesn't matter how many
spaces    you type; one is as good as 100.  The
end of   a line counts as a space.

One   or more   blank lines denote the  end
of  a paragraph.

Since any number of consecutive spaces are treated
like a single one, the formatting of the input
file makes no difference to
      \LaTeX,                % The \LaTeX command generates the LaTeX logo.
but it makes a difference to you.  When you use
\LaTeX, making your input file as easy to read
as possible will be a great help as you write
your document and when you change it.  This sample
file shows how you can add comments to your own input
file.

Because printing is different from typewriting,
there are a number of things that you have to do
differently when preparing an input file than if
you were just typing the document directly.
Quotation marks like
       ``this''
have to be handled specially, as do quotes within
quotes:
       ``\,`this'            % \, separates the double and single quote.
        is what I just
        wrote, not  `that'\,''.

Dashes come in three sizes: an
       intra-word
dash, a medium dash for number ranges like
       1--2,
and a punctuation
       dash---like
this.

A sentence-ending space should be larger than the
space between words within a sentence.  You
sometimes have to type special commands in
conjunction with punctuation characters to get
this right, as in the following sentence.
       Gnats, gnus, etc.\ all  % `\ ' makes an inter-word space.
       begin with G\@.         % \@ marks end-of-sentence punctuation.
You should check the spaces after periods when
reading your output to make sure you haven't
forgotten any special cases.  Generating an
ellipsis
       \ldots\               % `\ ' is needed after `\ldots' because TeX
                             % ignores spaces after command names like \ldots
                             % made from \ + letters.
                             %
                             % Note how a `%' character causes TeX to ignore
                             % the end of the input line, so these blank lines
                             % do not start a new paragraph.
                             %
with the right spacing around the periods requires
a special command.

\LaTeX\ interprets some common characters as
commands, so you must type special commands to
generate them.  These characters include the
following:
       \$ \& \% \# \{ and \}.

In printing, text is usually emphasized with an
       \emph{italic}
type style.

\begin{em}
   A long segment of text can also be emphasized
   in this way.  Text within such a segment can be
   given \emph{additional} emphasis.
\end{em}

It is sometimes necessary to prevent \LaTeX\ from
breaking a line where it might otherwise do so.
This may be at a space, as between the ``Mr.''\ and
``Jones'' in
       ``Mr.~Jones'',        % ~ produces an unbreakable interword space.
or within a word---especially when the word is a
symbol like
       \mbox{\emph{itemnum}}
that makes little sense when hyphenated across
lines.

Footnotes\footnote{This is an example of a footnote.}
pose no problem.

\LaTeX\ is good at typesetting mathematical formulas
like
       \( x-3y + z = 7 \)
or
       \( a_{1} > x^{2n} + y^{2n} > x' \)
or
       \( \ip{A}{B} = \sum_{i} a_{i} b_{i} \).
The spaces you type in a formula are
ignored.  Remember that a letter like
       $x$                   % $ ... $  and  \( ... \)  are equivalent
is a formula when it denotes a mathematical
symbol, and it should be typed as one.

\section{Displayed Text}

Text is displayed by indenting it from the left
margin.  Quotations are commonly displayed.  There
are short quotations
\begin{quote}
   This is a short quotation.  It consists of a
   single paragraph of text.  See how it is formatted.
\end{quote}
and longer ones.
\begin{quotation}
   This is a longer quotation.  It consists of two
   paragraphs of text, neither of which are
   particularly interesting.

   This is the second paragraph of the quotation.  It
   is just as dull as the first paragraph.
\end{quotation}
Another frequently-displayed structure is a list.
The following is an example of an \emph{itemized}
list.
\begin{itemize}
   \item This is the first item of an itemized list.
         Each item in the list is marked with a ``tick''.
         You don't have to worry about what kind of tick
         mark is used.

   \item This is the second item of the list.  It
         contains another list nested inside it.  The inner
         list is an \emph{enumerated} list.
         \begin{enumerate}
            \item This is the first item of an enumerated
                  list that is nested within the itemized list.

            \item This is the second item of the inner list.
                  \LaTeX\ allows you to nest lists deeper than
                  you really should.
         \end{enumerate}
         This is the rest of the second item of the outer
         list.  It is no more interesting than any other
         part of the item.
   \item This is the third item of the list.
\end{itemize}
You can even display poetry.
\begin{verse}
   There is an environment
    for verse \\             % The \\ command separates lines
   Whose features some poets % within a stanza.
   will curse.

                             % One or more blank lines separate stanzas.

   For instead of making\\
   Them do \emph{all} line breaking, \\
   It allows them to put too many words on a line when they'd rather be
   forced to be terse.
\end{verse}

Mathematical formulas may also be displayed.  A
displayed formula
is
one-line long; multiline
formulas require special formatting instructions.
   \[  \ip{\Gamma}{\psi'} = x'' + y^{2} + z_{i}^{n}\]
Don't start a paragraph with a displayed equation,
nor make one a paragraph by itself.

\end{document}               % End of document.

\section{Bench Huge Section 821}
% synthetic bench section 821
% This is a sample LaTeX input file.  (Version of 12 August 2004.)
%
% A '%' character causes TeX to ignore all remaining text on the line,
% and is used for comments like this one.

\documentclass{article}      % Specifies the document class

                             % The preamble begins here.
\title{An Example Document}  % Declares the document's title.
\author{Leslie Lamport}      % Declares the author's name.
\date{January 21, 1994}      % Deleting this command produces today's date.

\newcommand{\ip}[2]{(#1, #2)}
                             % Defines \ip{arg1}{arg2} to mean
                             % (arg1, arg2).

%\newcommand{\ip}[2]{\langle #1 | #2\rangle}
                             % This is an alternative definition of
                             % \ip that is commented out.

\begin{document}             % End of preamble and beginning of text.

\maketitle                   % Produces the title.

This is an example input file.  Comparing it with
the output it generates can show you how to
produce a simple document of your own.

\section{Ordinary Text}      % Produces section heading.  Lower-level
                             % sections are begun with similar
                             % \subsection and \subsubsection commands.

The ends  of words and sentences are marked
  by   spaces. It  doesn't matter how many
spaces    you type; one is as good as 100.  The
end of   a line counts as a space.

One   or more   blank lines denote the  end
of  a paragraph.

Since any number of consecutive spaces are treated
like a single one, the formatting of the input
file makes no difference to
      \LaTeX,                % The \LaTeX command generates the LaTeX logo.
but it makes a difference to you.  When you use
\LaTeX, making your input file as easy to read
as possible will be a great help as you write
your document and when you change it.  This sample
file shows how you can add comments to your own input
file.

Because printing is different from typewriting,
there are a number of things that you have to do
differently when preparing an input file than if
you were just typing the document directly.
Quotation marks like
       ``this''
have to be handled specially, as do quotes within
quotes:
       ``\,`this'            % \, separates the double and single quote.
        is what I just
        wrote, not  `that'\,''.

Dashes come in three sizes: an
       intra-word
dash, a medium dash for number ranges like
       1--2,
and a punctuation
       dash---like
this.

A sentence-ending space should be larger than the
space between words within a sentence.  You
sometimes have to type special commands in
conjunction with punctuation characters to get
this right, as in the following sentence.
       Gnats, gnus, etc.\ all  % `\ ' makes an inter-word space.
       begin with G\@.         % \@ marks end-of-sentence punctuation.
You should check the spaces after periods when
reading your output to make sure you haven't
forgotten any special cases.  Generating an
ellipsis
       \ldots\               % `\ ' is needed after `\ldots' because TeX
                             % ignores spaces after command names like \ldots
                             % made from \ + letters.
                             %
                             % Note how a `%' character causes TeX to ignore
                             % the end of the input line, so these blank lines
                             % do not start a new paragraph.
                             %
with the right spacing around the periods requires
a special command.

\LaTeX\ interprets some common characters as
commands, so you must type special commands to
generate them.  These characters include the
following:
       \$ \& \% \# \{ and \}.

In printing, text is usually emphasized with an
       \emph{italic}
type style.

\begin{em}
   A long segment of text can also be emphasized
   in this way.  Text within such a segment can be
   given \emph{additional} emphasis.
\end{em}

It is sometimes necessary to prevent \LaTeX\ from
breaking a line where it might otherwise do so.
This may be at a space, as between the ``Mr.''\ and
``Jones'' in
       ``Mr.~Jones'',        % ~ produces an unbreakable interword space.
or within a word---especially when the word is a
symbol like
       \mbox{\emph{itemnum}}
that makes little sense when hyphenated across
lines.

Footnotes\footnote{This is an example of a footnote.}
pose no problem.

\LaTeX\ is good at typesetting mathematical formulas
like
       \( x-3y + z = 7 \)
or
       \( a_{1} > x^{2n} + y^{2n} > x' \)
or
       \( \ip{A}{B} = \sum_{i} a_{i} b_{i} \).
The spaces you type in a formula are
ignored.  Remember that a letter like
       $x$                   % $ ... $  and  \( ... \)  are equivalent
is a formula when it denotes a mathematical
symbol, and it should be typed as one.

\section{Displayed Text}

Text is displayed by indenting it from the left
margin.  Quotations are commonly displayed.  There
are short quotations
\begin{quote}
   This is a short quotation.  It consists of a
   single paragraph of text.  See how it is formatted.
\end{quote}
and longer ones.
\begin{quotation}
   This is a longer quotation.  It consists of two
   paragraphs of text, neither of which are
   particularly interesting.

   This is the second paragraph of the quotation.  It
   is just as dull as the first paragraph.
\end{quotation}
Another frequently-displayed structure is a list.
The following is an example of an \emph{itemized}
list.
\begin{itemize}
   \item This is the first item of an itemized list.
         Each item in the list is marked with a ``tick''.
         You don't have to worry about what kind of tick
         mark is used.

   \item This is the second item of the list.  It
         contains another list nested inside it.  The inner
         list is an \emph{enumerated} list.
         \begin{enumerate}
            \item This is the first item of an enumerated
                  list that is nested within the itemized list.

            \item This is the second item of the inner list.
                  \LaTeX\ allows you to nest lists deeper than
                  you really should.
         \end{enumerate}
         This is the rest of the second item of the outer
         list.  It is no more interesting than any other
         part of the item.
   \item This is the third item of the list.
\end{itemize}
You can even display poetry.
\begin{verse}
   There is an environment
    for verse \\             % The \\ command separates lines
   Whose features some poets % within a stanza.
   will curse.

                             % One or more blank lines separate stanzas.

   For instead of making\\
   Them do \emph{all} line breaking, \\
   It allows them to put too many words on a line when they'd rather be
   forced to be terse.
\end{verse}

Mathematical formulas may also be displayed.  A
displayed formula
is
one-line long; multiline
formulas require special formatting instructions.
   \[  \ip{\Gamma}{\psi'} = x'' + y^{2} + z_{i}^{n}\]
Don't start a paragraph with a displayed equation,
nor make one a paragraph by itself.

\end{document}               % End of document.

\section{Bench Huge Section 822}
% synthetic bench section 822
% This is a sample LaTeX input file.  (Version of 12 August 2004.)
%
% A '%' character causes TeX to ignore all remaining text on the line,
% and is used for comments like this one.

\documentclass{article}      % Specifies the document class

                             % The preamble begins here.
\title{An Example Document}  % Declares the document's title.
\author{Leslie Lamport}      % Declares the author's name.
\date{January 21, 1994}      % Deleting this command produces today's date.

\newcommand{\ip}[2]{(#1, #2)}
                             % Defines \ip{arg1}{arg2} to mean
                             % (arg1, arg2).

%\newcommand{\ip}[2]{\langle #1 | #2\rangle}
                             % This is an alternative definition of
                             % \ip that is commented out.

\begin{document}             % End of preamble and beginning of text.

\maketitle                   % Produces the title.

This is an example input file.  Comparing it with
the output it generates can show you how to
produce a simple document of your own.

\section{Ordinary Text}      % Produces section heading.  Lower-level
                             % sections are begun with similar
                             % \subsection and \subsubsection commands.

The ends  of words and sentences are marked
  by   spaces. It  doesn't matter how many
spaces    you type; one is as good as 100.  The
end of   a line counts as a space.

One   or more   blank lines denote the  end
of  a paragraph.

Since any number of consecutive spaces are treated
like a single one, the formatting of the input
file makes no difference to
      \LaTeX,                % The \LaTeX command generates the LaTeX logo.
but it makes a difference to you.  When you use
\LaTeX, making your input file as easy to read
as possible will be a great help as you write
your document and when you change it.  This sample
file shows how you can add comments to your own input
file.

Because printing is different from typewriting,
there are a number of things that you have to do
differently when preparing an input file than if
you were just typing the document directly.
Quotation marks like
       ``this''
have to be handled specially, as do quotes within
quotes:
       ``\,`this'            % \, separates the double and single quote.
        is what I just
        wrote, not  `that'\,''.

Dashes come in three sizes: an
       intra-word
dash, a medium dash for number ranges like
       1--2,
and a punctuation
       dash---like
this.

A sentence-ending space should be larger than the
space between words within a sentence.  You
sometimes have to type special commands in
conjunction with punctuation characters to get
this right, as in the following sentence.
       Gnats, gnus, etc.\ all  % `\ ' makes an inter-word space.
       begin with G\@.         % \@ marks end-of-sentence punctuation.
You should check the spaces after periods when
reading your output to make sure you haven't
forgotten any special cases.  Generating an
ellipsis
       \ldots\               % `\ ' is needed after `\ldots' because TeX
                             % ignores spaces after command names like \ldots
                             % made from \ + letters.
                             %
                             % Note how a `%' character causes TeX to ignore
                             % the end of the input line, so these blank lines
                             % do not start a new paragraph.
                             %
with the right spacing around the periods requires
a special command.

\LaTeX\ interprets some common characters as
commands, so you must type special commands to
generate them.  These characters include the
following:
       \$ \& \% \# \{ and \}.

In printing, text is usually emphasized with an
       \emph{italic}
type style.

\begin{em}
   A long segment of text can also be emphasized
   in this way.  Text within such a segment can be
   given \emph{additional} emphasis.
\end{em}

It is sometimes necessary to prevent \LaTeX\ from
breaking a line where it might otherwise do so.
This may be at a space, as between the ``Mr.''\ and
``Jones'' in
       ``Mr.~Jones'',        % ~ produces an unbreakable interword space.
or within a word---especially when the word is a
symbol like
       \mbox{\emph{itemnum}}
that makes little sense when hyphenated across
lines.

Footnotes\footnote{This is an example of a footnote.}
pose no problem.

\LaTeX\ is good at typesetting mathematical formulas
like
       \( x-3y + z = 7 \)
or
       \( a_{1} > x^{2n} + y^{2n} > x' \)
or
       \( \ip{A}{B} = \sum_{i} a_{i} b_{i} \).
The spaces you type in a formula are
ignored.  Remember that a letter like
       $x$                   % $ ... $  and  \( ... \)  are equivalent
is a formula when it denotes a mathematical
symbol, and it should be typed as one.

\section{Displayed Text}

Text is displayed by indenting it from the left
margin.  Quotations are commonly displayed.  There
are short quotations
\begin{quote}
   This is a short quotation.  It consists of a
   single paragraph of text.  See how it is formatted.
\end{quote}
and longer ones.
\begin{quotation}
   This is a longer quotation.  It consists of two
   paragraphs of text, neither of which are
   particularly interesting.

   This is the second paragraph of the quotation.  It
   is just as dull as the first paragraph.
\end{quotation}
Another frequently-displayed structure is a list.
The following is an example of an \emph{itemized}
list.
\begin{itemize}
   \item This is the first item of an itemized list.
         Each item in the list is marked with a ``tick''.
         You don't have to worry about what kind of tick
         mark is used.

   \item This is the second item of the list.  It
         contains another list nested inside it.  The inner
         list is an \emph{enumerated} list.
         \begin{enumerate}
            \item This is the first item of an enumerated
                  list that is nested within the itemized list.

            \item This is the second item of the inner list.
                  \LaTeX\ allows you to nest lists deeper than
                  you really should.
         \end{enumerate}
         This is the rest of the second item of the outer
         list.  It is no more interesting than any other
         part of the item.
   \item This is the third item of the list.
\end{itemize}
You can even display poetry.
\begin{verse}
   There is an environment
    for verse \\             % The \\ command separates lines
   Whose features some poets % within a stanza.
   will curse.

                             % One or more blank lines separate stanzas.

   For instead of making\\
   Them do \emph{all} line breaking, \\
   It allows them to put too many words on a line when they'd rather be
   forced to be terse.
\end{verse}

Mathematical formulas may also be displayed.  A
displayed formula
is
one-line long; multiline
formulas require special formatting instructions.
   \[  \ip{\Gamma}{\psi'} = x'' + y^{2} + z_{i}^{n}\]
Don't start a paragraph with a displayed equation,
nor make one a paragraph by itself.

\end{document}               % End of document.

\section{Bench Huge Section 823}
% synthetic bench section 823
% This is a sample LaTeX input file.  (Version of 12 August 2004.)
%
% A '%' character causes TeX to ignore all remaining text on the line,
% and is used for comments like this one.

\documentclass{article}      % Specifies the document class

                             % The preamble begins here.
\title{An Example Document}  % Declares the document's title.
\author{Leslie Lamport}      % Declares the author's name.
\date{January 21, 1994}      % Deleting this command produces today's date.

\newcommand{\ip}[2]{(#1, #2)}
                             % Defines \ip{arg1}{arg2} to mean
                             % (arg1, arg2).

%\newcommand{\ip}[2]{\langle #1 | #2\rangle}
                             % This is an alternative definition of
                             % \ip that is commented out.

\begin{document}             % End of preamble and beginning of text.

\maketitle                   % Produces the title.

This is an example input file.  Comparing it with
the output it generates can show you how to
produce a simple document of your own.

\section{Ordinary Text}      % Produces section heading.  Lower-level
                             % sections are begun with similar
                             % \subsection and \subsubsection commands.

The ends  of words and sentences are marked
  by   spaces. It  doesn't matter how many
spaces    you type; one is as good as 100.  The
end of   a line counts as a space.

One   or more   blank lines denote the  end
of  a paragraph.

Since any number of consecutive spaces are treated
like a single one, the formatting of the input
file makes no difference to
      \LaTeX,                % The \LaTeX command generates the LaTeX logo.
but it makes a difference to you.  When you use
\LaTeX, making your input file as easy to read
as possible will be a great help as you write
your document and when you change it.  This sample
file shows how you can add comments to your own input
file.

Because printing is different from typewriting,
there are a number of things that you have to do
differently when preparing an input file than if
you were just typing the document directly.
Quotation marks like
       ``this''
have to be handled specially, as do quotes within
quotes:
       ``\,`this'            % \, separates the double and single quote.
        is what I just
        wrote, not  `that'\,''.

Dashes come in three sizes: an
       intra-word
dash, a medium dash for number ranges like
       1--2,
and a punctuation
       dash---like
this.

A sentence-ending space should be larger than the
space between words within a sentence.  You
sometimes have to type special commands in
conjunction with punctuation characters to get
this right, as in the following sentence.
       Gnats, gnus, etc.\ all  % `\ ' makes an inter-word space.
       begin with G\@.         % \@ marks end-of-sentence punctuation.
You should check the spaces after periods when
reading your output to make sure you haven't
forgotten any special cases.  Generating an
ellipsis
       \ldots\               % `\ ' is needed after `\ldots' because TeX
                             % ignores spaces after command names like \ldots
                             % made from \ + letters.
                             %
                             % Note how a `%' character causes TeX to ignore
                             % the end of the input line, so these blank lines
                             % do not start a new paragraph.
                             %
with the right spacing around the periods requires
a special command.

\LaTeX\ interprets some common characters as
commands, so you must type special commands to
generate them.  These characters include the
following:
       \$ \& \% \# \{ and \}.

In printing, text is usually emphasized with an
       \emph{italic}
type style.

\begin{em}
   A long segment of text can also be emphasized
   in this way.  Text within such a segment can be
   given \emph{additional} emphasis.
\end{em}

It is sometimes necessary to prevent \LaTeX\ from
breaking a line where it might otherwise do so.
This may be at a space, as between the ``Mr.''\ and
``Jones'' in
       ``Mr.~Jones'',        % ~ produces an unbreakable interword space.
or within a word---especially when the word is a
symbol like
       \mbox{\emph{itemnum}}
that makes little sense when hyphenated across
lines.

Footnotes\footnote{This is an example of a footnote.}
pose no problem.

\LaTeX\ is good at typesetting mathematical formulas
like
       \( x-3y + z = 7 \)
or
       \( a_{1} > x^{2n} + y^{2n} > x' \)
or
       \( \ip{A}{B} = \sum_{i} a_{i} b_{i} \).
The spaces you type in a formula are
ignored.  Remember that a letter like
       $x$                   % $ ... $  and  \( ... \)  are equivalent
is a formula when it denotes a mathematical
symbol, and it should be typed as one.

\section{Displayed Text}

Text is displayed by indenting it from the left
margin.  Quotations are commonly displayed.  There
are short quotations
\begin{quote}
   This is a short quotation.  It consists of a
   single paragraph of text.  See how it is formatted.
\end{quote}
and longer ones.
\begin{quotation}
   This is a longer quotation.  It consists of two
   paragraphs of text, neither of which are
   particularly interesting.

   This is the second paragraph of the quotation.  It
   is just as dull as the first paragraph.
\end{quotation}
Another frequently-displayed structure is a list.
The following is an example of an \emph{itemized}
list.
\begin{itemize}
   \item This is the first item of an itemized list.
         Each item in the list is marked with a ``tick''.
         You don't have to worry about what kind of tick
         mark is used.

   \item This is the second item of the list.  It
         contains another list nested inside it.  The inner
         list is an \emph{enumerated} list.
         \begin{enumerate}
            \item This is the first item of an enumerated
                  list that is nested within the itemized list.

            \item This is the second item of the inner list.
                  \LaTeX\ allows you to nest lists deeper than
                  you really should.
         \end{enumerate}
         This is the rest of the second item of the outer
         list.  It is no more interesting than any other
         part of the item.
   \item This is the third item of the list.
\end{itemize}
You can even display poetry.
\begin{verse}
   There is an environment
    for verse \\             % The \\ command separates lines
   Whose features some poets % within a stanza.
   will curse.

                             % One or more blank lines separate stanzas.

   For instead of making\\
   Them do \emph{all} line breaking, \\
   It allows them to put too many words on a line when they'd rather be
   forced to be terse.
\end{verse}

Mathematical formulas may also be displayed.  A
displayed formula
is
one-line long; multiline
formulas require special formatting instructions.
   \[  \ip{\Gamma}{\psi'} = x'' + y^{2} + z_{i}^{n}\]
Don't start a paragraph with a displayed equation,
nor make one a paragraph by itself.

\end{document}               % End of document.

\section{Bench Huge Section 824}
% synthetic bench section 824
% This is a sample LaTeX input file.  (Version of 12 August 2004.)
%
% A '%' character causes TeX to ignore all remaining text on the line,
% and is used for comments like this one.

\documentclass{article}      % Specifies the document class

                             % The preamble begins here.
\title{An Example Document}  % Declares the document's title.
\author{Leslie Lamport}      % Declares the author's name.
\date{January 21, 1994}      % Deleting this command produces today's date.

\newcommand{\ip}[2]{(#1, #2)}
                             % Defines \ip{arg1}{arg2} to mean
                             % (arg1, arg2).

%\newcommand{\ip}[2]{\langle #1 | #2\rangle}
                             % This is an alternative definition of
                             % \ip that is commented out.

\begin{document}             % End of preamble and beginning of text.

\maketitle                   % Produces the title.

This is an example input file.  Comparing it with
the output it generates can show you how to
produce a simple document of your own.

\section{Ordinary Text}      % Produces section heading.  Lower-level
                             % sections are begun with similar
                             % \subsection and \subsubsection commands.

The ends  of words and sentences are marked
  by   spaces. It  doesn't matter how many
spaces    you type; one is as good as 100.  The
end of   a line counts as a space.

One   or more   blank lines denote the  end
of  a paragraph.

Since any number of consecutive spaces are treated
like a single one, the formatting of the input
file makes no difference to
      \LaTeX,                % The \LaTeX command generates the LaTeX logo.
but it makes a difference to you.  When you use
\LaTeX, making your input file as easy to read
as possible will be a great help as you write
your document and when you change it.  This sample
file shows how you can add comments to your own input
file.

Because printing is different from typewriting,
there are a number of things that you have to do
differently when preparing an input file than if
you were just typing the document directly.
Quotation marks like
       ``this''
have to be handled specially, as do quotes within
quotes:
       ``\,`this'            % \, separates the double and single quote.
        is what I just
        wrote, not  `that'\,''.

Dashes come in three sizes: an
       intra-word
dash, a medium dash for number ranges like
       1--2,
and a punctuation
       dash---like
this.

A sentence-ending space should be larger than the
space between words within a sentence.  You
sometimes have to type special commands in
conjunction with punctuation characters to get
this right, as in the following sentence.
       Gnats, gnus, etc.\ all  % `\ ' makes an inter-word space.
       begin with G\@.         % \@ marks end-of-sentence punctuation.
You should check the spaces after periods when
reading your output to make sure you haven't
forgotten any special cases.  Generating an
ellipsis
       \ldots\               % `\ ' is needed after `\ldots' because TeX
                             % ignores spaces after command names like \ldots
                             % made from \ + letters.
                             %
                             % Note how a `%' character causes TeX to ignore
                             % the end of the input line, so these blank lines
                             % do not start a new paragraph.
                             %
with the right spacing around the periods requires
a special command.

\LaTeX\ interprets some common characters as
commands, so you must type special commands to
generate them.  These characters include the
following:
       \$ \& \% \# \{ and \}.

In printing, text is usually emphasized with an
       \emph{italic}
type style.

\begin{em}
   A long segment of text can also be emphasized
   in this way.  Text within such a segment can be
   given \emph{additional} emphasis.
\end{em}

It is sometimes necessary to prevent \LaTeX\ from
breaking a line where it might otherwise do so.
This may be at a space, as between the ``Mr.''\ and
``Jones'' in
       ``Mr.~Jones'',        % ~ produces an unbreakable interword space.
or within a word---especially when the word is a
symbol like
       \mbox{\emph{itemnum}}
that makes little sense when hyphenated across
lines.

Footnotes\footnote{This is an example of a footnote.}
pose no problem.

\LaTeX\ is good at typesetting mathematical formulas
like
       \( x-3y + z = 7 \)
or
       \( a_{1} > x^{2n} + y^{2n} > x' \)
or
       \( \ip{A}{B} = \sum_{i} a_{i} b_{i} \).
The spaces you type in a formula are
ignored.  Remember that a letter like
       $x$                   % $ ... $  and  \( ... \)  are equivalent
is a formula when it denotes a mathematical
symbol, and it should be typed as one.

\section{Displayed Text}

Text is displayed by indenting it from the left
margin.  Quotations are commonly displayed.  There
are short quotations
\begin{quote}
   This is a short quotation.  It consists of a
   single paragraph of text.  See how it is formatted.
\end{quote}
and longer ones.
\begin{quotation}
   This is a longer quotation.  It consists of two
   paragraphs of text, neither of which are
   particularly interesting.

   This is the second paragraph of the quotation.  It
   is just as dull as the first paragraph.
\end{quotation}
Another frequently-displayed structure is a list.
The following is an example of an \emph{itemized}
list.
\begin{itemize}
   \item This is the first item of an itemized list.
         Each item in the list is marked with a ``tick''.
         You don't have to worry about what kind of tick
         mark is used.

   \item This is the second item of the list.  It
         contains another list nested inside it.  The inner
         list is an \emph{enumerated} list.
         \begin{enumerate}
            \item This is the first item of an enumerated
                  list that is nested within the itemized list.

            \item This is the second item of the inner list.
                  \LaTeX\ allows you to nest lists deeper than
                  you really should.
         \end{enumerate}
         This is the rest of the second item of the outer
         list.  It is no more interesting than any other
         part of the item.
   \item This is the third item of the list.
\end{itemize}
You can even display poetry.
\begin{verse}
   There is an environment
    for verse \\             % The \\ command separates lines
   Whose features some poets % within a stanza.
   will curse.

                             % One or more blank lines separate stanzas.

   For instead of making\\
   Them do \emph{all} line breaking, \\
   It allows them to put too many words on a line when they'd rather be
   forced to be terse.
\end{verse}

Mathematical formulas may also be displayed.  A
displayed formula
is
one-line long; multiline
formulas require special formatting instructions.
   \[  \ip{\Gamma}{\psi'} = x'' + y^{2} + z_{i}^{n}\]
Don't start a paragraph with a displayed equation,
nor make one a paragraph by itself.

\end{document}               % End of document.

\section{Bench Huge Section 825}
% synthetic bench section 825
% This is a sample LaTeX input file.  (Version of 12 August 2004.)
%
% A '%' character causes TeX to ignore all remaining text on the line,
% and is used for comments like this one.

\documentclass{article}      % Specifies the document class

                             % The preamble begins here.
\title{An Example Document}  % Declares the document's title.
\author{Leslie Lamport}      % Declares the author's name.
\date{January 21, 1994}      % Deleting this command produces today's date.

\newcommand{\ip}[2]{(#1, #2)}
                             % Defines \ip{arg1}{arg2} to mean
                             % (arg1, arg2).

%\newcommand{\ip}[2]{\langle #1 | #2\rangle}
                             % This is an alternative definition of
                             % \ip that is commented out.

\begin{document}             % End of preamble and beginning of text.

\maketitle                   % Produces the title.

This is an example input file.  Comparing it with
the output it generates can show you how to
produce a simple document of your own.

\section{Ordinary Text}      % Produces section heading.  Lower-level
                             % sections are begun with similar
                             % \subsection and \subsubsection commands.

The ends  of words and sentences are marked
  by   spaces. It  doesn't matter how many
spaces    you type; one is as good as 100.  The
end of   a line counts as a space.

One   or more   blank lines denote the  end
of  a paragraph.

Since any number of consecutive spaces are treated
like a single one, the formatting of the input
file makes no difference to
      \LaTeX,                % The \LaTeX command generates the LaTeX logo.
but it makes a difference to you.  When you use
\LaTeX, making your input file as easy to read
as possible will be a great help as you write
your document and when you change it.  This sample
file shows how you can add comments to your own input
file.

Because printing is different from typewriting,
there are a number of things that you have to do
differently when preparing an input file than if
you were just typing the document directly.
Quotation marks like
       ``this''
have to be handled specially, as do quotes within
quotes:
       ``\,`this'            % \, separates the double and single quote.
        is what I just
        wrote, not  `that'\,''.

Dashes come in three sizes: an
       intra-word
dash, a medium dash for number ranges like
       1--2,
and a punctuation
       dash---like
this.

A sentence-ending space should be larger than the
space between words within a sentence.  You
sometimes have to type special commands in
conjunction with punctuation characters to get
this right, as in the following sentence.
       Gnats, gnus, etc.\ all  % `\ ' makes an inter-word space.
       begin with G\@.         % \@ marks end-of-sentence punctuation.
You should check the spaces after periods when
reading your output to make sure you haven't
forgotten any special cases.  Generating an
ellipsis
       \ldots\               % `\ ' is needed after `\ldots' because TeX
                             % ignores spaces after command names like \ldots
                             % made from \ + letters.
                             %
                             % Note how a `%' character causes TeX to ignore
                             % the end of the input line, so these blank lines
                             % do not start a new paragraph.
                             %
with the right spacing around the periods requires
a special command.

\LaTeX\ interprets some common characters as
commands, so you must type special commands to
generate them.  These characters include the
following:
       \$ \& \% \# \{ and \}.

In printing, text is usually emphasized with an
       \emph{italic}
type style.

\begin{em}
   A long segment of text can also be emphasized
   in this way.  Text within such a segment can be
   given \emph{additional} emphasis.
\end{em}

It is sometimes necessary to prevent \LaTeX\ from
breaking a line where it might otherwise do so.
This may be at a space, as between the ``Mr.''\ and
``Jones'' in
       ``Mr.~Jones'',        % ~ produces an unbreakable interword space.
or within a word---especially when the word is a
symbol like
       \mbox{\emph{itemnum}}
that makes little sense when hyphenated across
lines.

Footnotes\footnote{This is an example of a footnote.}
pose no problem.

\LaTeX\ is good at typesetting mathematical formulas
like
       \( x-3y + z = 7 \)
or
       \( a_{1} > x^{2n} + y^{2n} > x' \)
or
       \( \ip{A}{B} = \sum_{i} a_{i} b_{i} \).
The spaces you type in a formula are
ignored.  Remember that a letter like
       $x$                   % $ ... $  and  \( ... \)  are equivalent
is a formula when it denotes a mathematical
symbol, and it should be typed as one.

\section{Displayed Text}

Text is displayed by indenting it from the left
margin.  Quotations are commonly displayed.  There
are short quotations
\begin{quote}
   This is a short quotation.  It consists of a
   single paragraph of text.  See how it is formatted.
\end{quote}
and longer ones.
\begin{quotation}
   This is a longer quotation.  It consists of two
   paragraphs of text, neither of which are
   particularly interesting.

   This is the second paragraph of the quotation.  It
   is just as dull as the first paragraph.
\end{quotation}
Another frequently-displayed structure is a list.
The following is an example of an \emph{itemized}
list.
\begin{itemize}
   \item This is the first item of an itemized list.
         Each item in the list is marked with a ``tick''.
         You don't have to worry about what kind of tick
         mark is used.

   \item This is the second item of the list.  It
         contains another list nested inside it.  The inner
         list is an \emph{enumerated} list.
         \begin{enumerate}
            \item This is the first item of an enumerated
                  list that is nested within the itemized list.

            \item This is the second item of the inner list.
                  \LaTeX\ allows you to nest lists deeper than
                  you really should.
         \end{enumerate}
         This is the rest of the second item of the outer
         list.  It is no more interesting than any other
         part of the item.
   \item This is the third item of the list.
\end{itemize}
You can even display poetry.
\begin{verse}
   There is an environment
    for verse \\             % The \\ command separates lines
   Whose features some poets % within a stanza.
   will curse.

                             % One or more blank lines separate stanzas.

   For instead of making\\
   Them do \emph{all} line breaking, \\
   It allows them to put too many words on a line when they'd rather be
   forced to be terse.
\end{verse}

Mathematical formulas may also be displayed.  A
displayed formula
is
one-line long; multiline
formulas require special formatting instructions.
   \[  \ip{\Gamma}{\psi'} = x'' + y^{2} + z_{i}^{n}\]
Don't start a paragraph with a displayed equation,
nor make one a paragraph by itself.

\end{document}               % End of document.

\section{Bench Huge Section 826}
% synthetic bench section 826
% This is a sample LaTeX input file.  (Version of 12 August 2004.)
%
% A '%' character causes TeX to ignore all remaining text on the line,
% and is used for comments like this one.

\documentclass{article}      % Specifies the document class

                             % The preamble begins here.
\title{An Example Document}  % Declares the document's title.
\author{Leslie Lamport}      % Declares the author's name.
\date{January 21, 1994}      % Deleting this command produces today's date.

\newcommand{\ip}[2]{(#1, #2)}
                             % Defines \ip{arg1}{arg2} to mean
                             % (arg1, arg2).

%\newcommand{\ip}[2]{\langle #1 | #2\rangle}
                             % This is an alternative definition of
                             % \ip that is commented out.

\begin{document}             % End of preamble and beginning of text.

\maketitle                   % Produces the title.

This is an example input file.  Comparing it with
the output it generates can show you how to
produce a simple document of your own.

\section{Ordinary Text}      % Produces section heading.  Lower-level
                             % sections are begun with similar
                             % \subsection and \subsubsection commands.

The ends  of words and sentences are marked
  by   spaces. It  doesn't matter how many
spaces    you type; one is as good as 100.  The
end of   a line counts as a space.

One   or more   blank lines denote the  end
of  a paragraph.

Since any number of consecutive spaces are treated
like a single one, the formatting of the input
file makes no difference to
      \LaTeX,                % The \LaTeX command generates the LaTeX logo.
but it makes a difference to you.  When you use
\LaTeX, making your input file as easy to read
as possible will be a great help as you write
your document and when you change it.  This sample
file shows how you can add comments to your own input
file.

Because printing is different from typewriting,
there are a number of things that you have to do
differently when preparing an input file than if
you were just typing the document directly.
Quotation marks like
       ``this''
have to be handled specially, as do quotes within
quotes:
       ``\,`this'            % \, separates the double and single quote.
        is what I just
        wrote, not  `that'\,''.

Dashes come in three sizes: an
       intra-word
dash, a medium dash for number ranges like
       1--2,
and a punctuation
       dash---like
this.

A sentence-ending space should be larger than the
space between words within a sentence.  You
sometimes have to type special commands in
conjunction with punctuation characters to get
this right, as in the following sentence.
       Gnats, gnus, etc.\ all  % `\ ' makes an inter-word space.
       begin with G\@.         % \@ marks end-of-sentence punctuation.
You should check the spaces after periods when
reading your output to make sure you haven't
forgotten any special cases.  Generating an
ellipsis
       \ldots\               % `\ ' is needed after `\ldots' because TeX
                             % ignores spaces after command names like \ldots
                             % made from \ + letters.
                             %
                             % Note how a `%' character causes TeX to ignore
                             % the end of the input line, so these blank lines
                             % do not start a new paragraph.
                             %
with the right spacing around the periods requires
a special command.

\LaTeX\ interprets some common characters as
commands, so you must type special commands to
generate them.  These characters include the
following:
       \$ \& \% \# \{ and \}.

In printing, text is usually emphasized with an
       \emph{italic}
type style.

\begin{em}
   A long segment of text can also be emphasized
   in this way.  Text within such a segment can be
   given \emph{additional} emphasis.
\end{em}

It is sometimes necessary to prevent \LaTeX\ from
breaking a line where it might otherwise do so.
This may be at a space, as between the ``Mr.''\ and
``Jones'' in
       ``Mr.~Jones'',        % ~ produces an unbreakable interword space.
or within a word---especially when the word is a
symbol like
       \mbox{\emph{itemnum}}
that makes little sense when hyphenated across
lines.

Footnotes\footnote{This is an example of a footnote.}
pose no problem.

\LaTeX\ is good at typesetting mathematical formulas
like
       \( x-3y + z = 7 \)
or
       \( a_{1} > x^{2n} + y^{2n} > x' \)
or
       \( \ip{A}{B} = \sum_{i} a_{i} b_{i} \).
The spaces you type in a formula are
ignored.  Remember that a letter like
       $x$                   % $ ... $  and  \( ... \)  are equivalent
is a formula when it denotes a mathematical
symbol, and it should be typed as one.

\section{Displayed Text}

Text is displayed by indenting it from the left
margin.  Quotations are commonly displayed.  There
are short quotations
\begin{quote}
   This is a short quotation.  It consists of a
   single paragraph of text.  See how it is formatted.
\end{quote}
and longer ones.
\begin{quotation}
   This is a longer quotation.  It consists of two
   paragraphs of text, neither of which are
   particularly interesting.

   This is the second paragraph of the quotation.  It
   is just as dull as the first paragraph.
\end{quotation}
Another frequently-displayed structure is a list.
The following is an example of an \emph{itemized}
list.
\begin{itemize}
   \item This is the first item of an itemized list.
         Each item in the list is marked with a ``tick''.
         You don't have to worry about what kind of tick
         mark is used.

   \item This is the second item of the list.  It
         contains another list nested inside it.  The inner
         list is an \emph{enumerated} list.
         \begin{enumerate}
            \item This is the first item of an enumerated
                  list that is nested within the itemized list.

            \item This is the second item of the inner list.
                  \LaTeX\ allows you to nest lists deeper than
                  you really should.
         \end{enumerate}
         This is the rest of the second item of the outer
         list.  It is no more interesting than any other
         part of the item.
   \item This is the third item of the list.
\end{itemize}
You can even display poetry.
\begin{verse}
   There is an environment
    for verse \\             % The \\ command separates lines
   Whose features some poets % within a stanza.
   will curse.

                             % One or more blank lines separate stanzas.

   For instead of making\\
   Them do \emph{all} line breaking, \\
   It allows them to put too many words on a line when they'd rather be
   forced to be terse.
\end{verse}

Mathematical formulas may also be displayed.  A
displayed formula
is
one-line long; multiline
formulas require special formatting instructions.
   \[  \ip{\Gamma}{\psi'} = x'' + y^{2} + z_{i}^{n}\]
Don't start a paragraph with a displayed equation,
nor make one a paragraph by itself.

\end{document}               % End of document.

\section{Bench Huge Section 827}
% synthetic bench section 827
% This is a sample LaTeX input file.  (Version of 12 August 2004.)
%
% A '%' character causes TeX to ignore all remaining text on the line,
% and is used for comments like this one.

\documentclass{article}      % Specifies the document class

                             % The preamble begins here.
\title{An Example Document}  % Declares the document's title.
\author{Leslie Lamport}      % Declares the author's name.
\date{January 21, 1994}      % Deleting this command produces today's date.

\newcommand{\ip}[2]{(#1, #2)}
                             % Defines \ip{arg1}{arg2} to mean
                             % (arg1, arg2).

%\newcommand{\ip}[2]{\langle #1 | #2\rangle}
                             % This is an alternative definition of
                             % \ip that is commented out.

\begin{document}             % End of preamble and beginning of text.

\maketitle                   % Produces the title.

This is an example input file.  Comparing it with
the output it generates can show you how to
produce a simple document of your own.

\section{Ordinary Text}      % Produces section heading.  Lower-level
                             % sections are begun with similar
                             % \subsection and \subsubsection commands.

The ends  of words and sentences are marked
  by   spaces. It  doesn't matter how many
spaces    you type; one is as good as 100.  The
end of   a line counts as a space.

One   or more   blank lines denote the  end
of  a paragraph.

Since any number of consecutive spaces are treated
like a single one, the formatting of the input
file makes no difference to
      \LaTeX,                % The \LaTeX command generates the LaTeX logo.
but it makes a difference to you.  When you use
\LaTeX, making your input file as easy to read
as possible will be a great help as you write
your document and when you change it.  This sample
file shows how you can add comments to your own input
file.

Because printing is different from typewriting,
there are a number of things that you have to do
differently when preparing an input file than if
you were just typing the document directly.
Quotation marks like
       ``this''
have to be handled specially, as do quotes within
quotes:
       ``\,`this'            % \, separates the double and single quote.
        is what I just
        wrote, not  `that'\,''.

Dashes come in three sizes: an
       intra-word
dash, a medium dash for number ranges like
       1--2,
and a punctuation
       dash---like
this.

A sentence-ending space should be larger than the
space between words within a sentence.  You
sometimes have to type special commands in
conjunction with punctuation characters to get
this right, as in the following sentence.
       Gnats, gnus, etc.\ all  % `\ ' makes an inter-word space.
       begin with G\@.         % \@ marks end-of-sentence punctuation.
You should check the spaces after periods when
reading your output to make sure you haven't
forgotten any special cases.  Generating an
ellipsis
       \ldots\               % `\ ' is needed after `\ldots' because TeX
                             % ignores spaces after command names like \ldots
                             % made from \ + letters.
                             %
                             % Note how a `%' character causes TeX to ignore
                             % the end of the input line, so these blank lines
                             % do not start a new paragraph.
                             %
with the right spacing around the periods requires
a special command.

\LaTeX\ interprets some common characters as
commands, so you must type special commands to
generate them.  These characters include the
following:
       \$ \& \% \# \{ and \}.

In printing, text is usually emphasized with an
       \emph{italic}
type style.

\begin{em}
   A long segment of text can also be emphasized
   in this way.  Text within such a segment can be
   given \emph{additional} emphasis.
\end{em}

It is sometimes necessary to prevent \LaTeX\ from
breaking a line where it might otherwise do so.
This may be at a space, as between the ``Mr.''\ and
``Jones'' in
       ``Mr.~Jones'',        % ~ produces an unbreakable interword space.
or within a word---especially when the word is a
symbol like
       \mbox{\emph{itemnum}}
that makes little sense when hyphenated across
lines.

Footnotes\footnote{This is an example of a footnote.}
pose no problem.

\LaTeX\ is good at typesetting mathematical formulas
like
       \( x-3y + z = 7 \)
or
       \( a_{1} > x^{2n} + y^{2n} > x' \)
or
       \( \ip{A}{B} = \sum_{i} a_{i} b_{i} \).
The spaces you type in a formula are
ignored.  Remember that a letter like
       $x$                   % $ ... $  and  \( ... \)  are equivalent
is a formula when it denotes a mathematical
symbol, and it should be typed as one.

\section{Displayed Text}

Text is displayed by indenting it from the left
margin.  Quotations are commonly displayed.  There
are short quotations
\begin{quote}
   This is a short quotation.  It consists of a
   single paragraph of text.  See how it is formatted.
\end{quote}
and longer ones.
\begin{quotation}
   This is a longer quotation.  It consists of two
   paragraphs of text, neither of which are
   particularly interesting.

   This is the second paragraph of the quotation.  It
   is just as dull as the first paragraph.
\end{quotation}
Another frequently-displayed structure is a list.
The following is an example of an \emph{itemized}
list.
\begin{itemize}
   \item This is the first item of an itemized list.
         Each item in the list is marked with a ``tick''.
         You don't have to worry about what kind of tick
         mark is used.

   \item This is the second item of the list.  It
         contains another list nested inside it.  The inner
         list is an \emph{enumerated} list.
         \begin{enumerate}
            \item This is the first item of an enumerated
                  list that is nested within the itemized list.

            \item This is the second item of the inner list.
                  \LaTeX\ allows you to nest lists deeper than
                  you really should.
         \end{enumerate}
         This is the rest of the second item of the outer
         list.  It is no more interesting than any other
         part of the item.
   \item This is the third item of the list.
\end{itemize}
You can even display poetry.
\begin{verse}
   There is an environment
    for verse \\             % The \\ command separates lines
   Whose features some poets % within a stanza.
   will curse.

                             % One or more blank lines separate stanzas.

   For instead of making\\
   Them do \emph{all} line breaking, \\
   It allows them to put too many words on a line when they'd rather be
   forced to be terse.
\end{verse}

Mathematical formulas may also be displayed.  A
displayed formula
is
one-line long; multiline
formulas require special formatting instructions.
   \[  \ip{\Gamma}{\psi'} = x'' + y^{2} + z_{i}^{n}\]
Don't start a paragraph with a displayed equation,
nor make one a paragraph by itself.

\end{document}               % End of document.

\section{Bench Huge Section 828}
% synthetic bench section 828
% This is a sample LaTeX input file.  (Version of 12 August 2004.)
%
% A '%' character causes TeX to ignore all remaining text on the line,
% and is used for comments like this one.

\documentclass{article}      % Specifies the document class

                             % The preamble begins here.
\title{An Example Document}  % Declares the document's title.
\author{Leslie Lamport}      % Declares the author's name.
\date{January 21, 1994}      % Deleting this command produces today's date.

\newcommand{\ip}[2]{(#1, #2)}
                             % Defines \ip{arg1}{arg2} to mean
                             % (arg1, arg2).

%\newcommand{\ip}[2]{\langle #1 | #2\rangle}
                             % This is an alternative definition of
                             % \ip that is commented out.

\begin{document}             % End of preamble and beginning of text.

\maketitle                   % Produces the title.

This is an example input file.  Comparing it with
the output it generates can show you how to
produce a simple document of your own.

\section{Ordinary Text}      % Produces section heading.  Lower-level
                             % sections are begun with similar
                             % \subsection and \subsubsection commands.

The ends  of words and sentences are marked
  by   spaces. It  doesn't matter how many
spaces    you type; one is as good as 100.  The
end of   a line counts as a space.

One   or more   blank lines denote the  end
of  a paragraph.

Since any number of consecutive spaces are treated
like a single one, the formatting of the input
file makes no difference to
      \LaTeX,                % The \LaTeX command generates the LaTeX logo.
but it makes a difference to you.  When you use
\LaTeX, making your input file as easy to read
as possible will be a great help as you write
your document and when you change it.  This sample
file shows how you can add comments to your own input
file.

Because printing is different from typewriting,
there are a number of things that you have to do
differently when preparing an input file than if
you were just typing the document directly.
Quotation marks like
       ``this''
have to be handled specially, as do quotes within
quotes:
       ``\,`this'            % \, separates the double and single quote.
        is what I just
        wrote, not  `that'\,''.

Dashes come in three sizes: an
       intra-word
dash, a medium dash for number ranges like
       1--2,
and a punctuation
       dash---like
this.

A sentence-ending space should be larger than the
space between words within a sentence.  You
sometimes have to type special commands in
conjunction with punctuation characters to get
this right, as in the following sentence.
       Gnats, gnus, etc.\ all  % `\ ' makes an inter-word space.
       begin with G\@.         % \@ marks end-of-sentence punctuation.
You should check the spaces after periods when
reading your output to make sure you haven't
forgotten any special cases.  Generating an
ellipsis
       \ldots\               % `\ ' is needed after `\ldots' because TeX
                             % ignores spaces after command names like \ldots
                             % made from \ + letters.
                             %
                             % Note how a `%' character causes TeX to ignore
                             % the end of the input line, so these blank lines
                             % do not start a new paragraph.
                             %
with the right spacing around the periods requires
a special command.

\LaTeX\ interprets some common characters as
commands, so you must type special commands to
generate them.  These characters include the
following:
       \$ \& \% \# \{ and \}.

In printing, text is usually emphasized with an
       \emph{italic}
type style.

\begin{em}
   A long segment of text can also be emphasized
   in this way.  Text within such a segment can be
   given \emph{additional} emphasis.
\end{em}

It is sometimes necessary to prevent \LaTeX\ from
breaking a line where it might otherwise do so.
This may be at a space, as between the ``Mr.''\ and
``Jones'' in
       ``Mr.~Jones'',        % ~ produces an unbreakable interword space.
or within a word---especially when the word is a
symbol like
       \mbox{\emph{itemnum}}
that makes little sense when hyphenated across
lines.

Footnotes\footnote{This is an example of a footnote.}
pose no problem.

\LaTeX\ is good at typesetting mathematical formulas
like
       \( x-3y + z = 7 \)
or
       \( a_{1} > x^{2n} + y^{2n} > x' \)
or
       \( \ip{A}{B} = \sum_{i} a_{i} b_{i} \).
The spaces you type in a formula are
ignored.  Remember that a letter like
       $x$                   % $ ... $  and  \( ... \)  are equivalent
is a formula when it denotes a mathematical
symbol, and it should be typed as one.

\section{Displayed Text}

Text is displayed by indenting it from the left
margin.  Quotations are commonly displayed.  There
are short quotations
\begin{quote}
   This is a short quotation.  It consists of a
   single paragraph of text.  See how it is formatted.
\end{quote}
and longer ones.
\begin{quotation}
   This is a longer quotation.  It consists of two
   paragraphs of text, neither of which are
   particularly interesting.

   This is the second paragraph of the quotation.  It
   is just as dull as the first paragraph.
\end{quotation}
Another frequently-displayed structure is a list.
The following is an example of an \emph{itemized}
list.
\begin{itemize}
   \item This is the first item of an itemized list.
         Each item in the list is marked with a ``tick''.
         You don't have to worry about what kind of tick
         mark is used.

   \item This is the second item of the list.  It
         contains another list nested inside it.  The inner
         list is an \emph{enumerated} list.
         \begin{enumerate}
            \item This is the first item of an enumerated
                  list that is nested within the itemized list.

            \item This is the second item of the inner list.
                  \LaTeX\ allows you to nest lists deeper than
                  you really should.
         \end{enumerate}
         This is the rest of the second item of the outer
         list.  It is no more interesting than any other
         part of the item.
   \item This is the third item of the list.
\end{itemize}
You can even display poetry.
\begin{verse}
   There is an environment
    for verse \\             % The \\ command separates lines
   Whose features some poets % within a stanza.
   will curse.

                             % One or more blank lines separate stanzas.

   For instead of making\\
   Them do \emph{all} line breaking, \\
   It allows them to put too many words on a line when they'd rather be
   forced to be terse.
\end{verse}

Mathematical formulas may also be displayed.  A
displayed formula
is
one-line long; multiline
formulas require special formatting instructions.
   \[  \ip{\Gamma}{\psi'} = x'' + y^{2} + z_{i}^{n}\]
Don't start a paragraph with a displayed equation,
nor make one a paragraph by itself.

\end{document}               % End of document.

\section{Bench Huge Section 829}
% synthetic bench section 829
% This is a sample LaTeX input file.  (Version of 12 August 2004.)
%
% A '%' character causes TeX to ignore all remaining text on the line,
% and is used for comments like this one.

\documentclass{article}      % Specifies the document class

                             % The preamble begins here.
\title{An Example Document}  % Declares the document's title.
\author{Leslie Lamport}      % Declares the author's name.
\date{January 21, 1994}      % Deleting this command produces today's date.

\newcommand{\ip}[2]{(#1, #2)}
                             % Defines \ip{arg1}{arg2} to mean
                             % (arg1, arg2).

%\newcommand{\ip}[2]{\langle #1 | #2\rangle}
                             % This is an alternative definition of
                             % \ip that is commented out.

\begin{document}             % End of preamble and beginning of text.

\maketitle                   % Produces the title.

This is an example input file.  Comparing it with
the output it generates can show you how to
produce a simple document of your own.

\section{Ordinary Text}      % Produces section heading.  Lower-level
                             % sections are begun with similar
                             % \subsection and \subsubsection commands.

The ends  of words and sentences are marked
  by   spaces. It  doesn't matter how many
spaces    you type; one is as good as 100.  The
end of   a line counts as a space.

One   or more   blank lines denote the  end
of  a paragraph.

Since any number of consecutive spaces are treated
like a single one, the formatting of the input
file makes no difference to
      \LaTeX,                % The \LaTeX command generates the LaTeX logo.
but it makes a difference to you.  When you use
\LaTeX, making your input file as easy to read
as possible will be a great help as you write
your document and when you change it.  This sample
file shows how you can add comments to your own input
file.

Because printing is different from typewriting,
there are a number of things that you have to do
differently when preparing an input file than if
you were just typing the document directly.
Quotation marks like
       ``this''
have to be handled specially, as do quotes within
quotes:
       ``\,`this'            % \, separates the double and single quote.
        is what I just
        wrote, not  `that'\,''.

Dashes come in three sizes: an
       intra-word
dash, a medium dash for number ranges like
       1--2,
and a punctuation
       dash---like
this.

A sentence-ending space should be larger than the
space between words within a sentence.  You
sometimes have to type special commands in
conjunction with punctuation characters to get
this right, as in the following sentence.
       Gnats, gnus, etc.\ all  % `\ ' makes an inter-word space.
       begin with G\@.         % \@ marks end-of-sentence punctuation.
You should check the spaces after periods when
reading your output to make sure you haven't
forgotten any special cases.  Generating an
ellipsis
       \ldots\               % `\ ' is needed after `\ldots' because TeX
                             % ignores spaces after command names like \ldots
                             % made from \ + letters.
                             %
                             % Note how a `%' character causes TeX to ignore
                             % the end of the input line, so these blank lines
                             % do not start a new paragraph.
                             %
with the right spacing around the periods requires
a special command.

\LaTeX\ interprets some common characters as
commands, so you must type special commands to
generate them.  These characters include the
following:
       \$ \& \% \# \{ and \}.

In printing, text is usually emphasized with an
       \emph{italic}
type style.

\begin{em}
   A long segment of text can also be emphasized
   in this way.  Text within such a segment can be
   given \emph{additional} emphasis.
\end{em}

It is sometimes necessary to prevent \LaTeX\ from
breaking a line where it might otherwise do so.
This may be at a space, as between the ``Mr.''\ and
``Jones'' in
       ``Mr.~Jones'',        % ~ produces an unbreakable interword space.
or within a word---especially when the word is a
symbol like
       \mbox{\emph{itemnum}}
that makes little sense when hyphenated across
lines.

Footnotes\footnote{This is an example of a footnote.}
pose no problem.

\LaTeX\ is good at typesetting mathematical formulas
like
       \( x-3y + z = 7 \)
or
       \( a_{1} > x^{2n} + y^{2n} > x' \)
or
       \( \ip{A}{B} = \sum_{i} a_{i} b_{i} \).
The spaces you type in a formula are
ignored.  Remember that a letter like
       $x$                   % $ ... $  and  \( ... \)  are equivalent
is a formula when it denotes a mathematical
symbol, and it should be typed as one.

\section{Displayed Text}

Text is displayed by indenting it from the left
margin.  Quotations are commonly displayed.  There
are short quotations
\begin{quote}
   This is a short quotation.  It consists of a
   single paragraph of text.  See how it is formatted.
\end{quote}
and longer ones.
\begin{quotation}
   This is a longer quotation.  It consists of two
   paragraphs of text, neither of which are
   particularly interesting.

   This is the second paragraph of the quotation.  It
   is just as dull as the first paragraph.
\end{quotation}
Another frequently-displayed structure is a list.
The following is an example of an \emph{itemized}
list.
\begin{itemize}
   \item This is the first item of an itemized list.
         Each item in the list is marked with a ``tick''.
         You don't have to worry about what kind of tick
         mark is used.

   \item This is the second item of the list.  It
         contains another list nested inside it.  The inner
         list is an \emph{enumerated} list.
         \begin{enumerate}
            \item This is the first item of an enumerated
                  list that is nested within the itemized list.

            \item This is the second item of the inner list.
                  \LaTeX\ allows you to nest lists deeper than
                  you really should.
         \end{enumerate}
         This is the rest of the second item of the outer
         list.  It is no more interesting than any other
         part of the item.
   \item This is the third item of the list.
\end{itemize}
You can even display poetry.
\begin{verse}
   There is an environment
    for verse \\             % The \\ command separates lines
   Whose features some poets % within a stanza.
   will curse.

                             % One or more blank lines separate stanzas.

   For instead of making\\
   Them do \emph{all} line breaking, \\
   It allows them to put too many words on a line when they'd rather be
   forced to be terse.
\end{verse}

Mathematical formulas may also be displayed.  A
displayed formula
is
one-line long; multiline
formulas require special formatting instructions.
   \[  \ip{\Gamma}{\psi'} = x'' + y^{2} + z_{i}^{n}\]
Don't start a paragraph with a displayed equation,
nor make one a paragraph by itself.

\end{document}               % End of document.

\section{Bench Huge Section 830}
% synthetic bench section 830
% This is a sample LaTeX input file.  (Version of 12 August 2004.)
%
% A '%' character causes TeX to ignore all remaining text on the line,
% and is used for comments like this one.

\documentclass{article}      % Specifies the document class

                             % The preamble begins here.
\title{An Example Document}  % Declares the document's title.
\author{Leslie Lamport}      % Declares the author's name.
\date{January 21, 1994}      % Deleting this command produces today's date.

\newcommand{\ip}[2]{(#1, #2)}
                             % Defines \ip{arg1}{arg2} to mean
                             % (arg1, arg2).

%\newcommand{\ip}[2]{\langle #1 | #2\rangle}
                             % This is an alternative definition of
                             % \ip that is commented out.

\begin{document}             % End of preamble and beginning of text.

\maketitle                   % Produces the title.

This is an example input file.  Comparing it with
the output it generates can show you how to
produce a simple document of your own.

\section{Ordinary Text}      % Produces section heading.  Lower-level
                             % sections are begun with similar
                             % \subsection and \subsubsection commands.

The ends  of words and sentences are marked
  by   spaces. It  doesn't matter how many
spaces    you type; one is as good as 100.  The
end of   a line counts as a space.

One   or more   blank lines denote the  end
of  a paragraph.

Since any number of consecutive spaces are treated
like a single one, the formatting of the input
file makes no difference to
      \LaTeX,                % The \LaTeX command generates the LaTeX logo.
but it makes a difference to you.  When you use
\LaTeX, making your input file as easy to read
as possible will be a great help as you write
your document and when you change it.  This sample
file shows how you can add comments to your own input
file.

Because printing is different from typewriting,
there are a number of things that you have to do
differently when preparing an input file than if
you were just typing the document directly.
Quotation marks like
       ``this''
have to be handled specially, as do quotes within
quotes:
       ``\,`this'            % \, separates the double and single quote.
        is what I just
        wrote, not  `that'\,''.

Dashes come in three sizes: an
       intra-word
dash, a medium dash for number ranges like
       1--2,
and a punctuation
       dash---like
this.

A sentence-ending space should be larger than the
space between words within a sentence.  You
sometimes have to type special commands in
conjunction with punctuation characters to get
this right, as in the following sentence.
       Gnats, gnus, etc.\ all  % `\ ' makes an inter-word space.
       begin with G\@.         % \@ marks end-of-sentence punctuation.
You should check the spaces after periods when
reading your output to make sure you haven't
forgotten any special cases.  Generating an
ellipsis
       \ldots\               % `\ ' is needed after `\ldots' because TeX
                             % ignores spaces after command names like \ldots
                             % made from \ + letters.
                             %
                             % Note how a `%' character causes TeX to ignore
                             % the end of the input line, so these blank lines
                             % do not start a new paragraph.
                             %
with the right spacing around the periods requires
a special command.

\LaTeX\ interprets some common characters as
commands, so you must type special commands to
generate them.  These characters include the
following:
       \$ \& \% \# \{ and \}.

In printing, text is usually emphasized with an
       \emph{italic}
type style.

\begin{em}
   A long segment of text can also be emphasized
   in this way.  Text within such a segment can be
   given \emph{additional} emphasis.
\end{em}

It is sometimes necessary to prevent \LaTeX\ from
breaking a line where it might otherwise do so.
This may be at a space, as between the ``Mr.''\ and
``Jones'' in
       ``Mr.~Jones'',        % ~ produces an unbreakable interword space.
or within a word---especially when the word is a
symbol like
       \mbox{\emph{itemnum}}
that makes little sense when hyphenated across
lines.

Footnotes\footnote{This is an example of a footnote.}
pose no problem.

\LaTeX\ is good at typesetting mathematical formulas
like
       \( x-3y + z = 7 \)
or
       \( a_{1} > x^{2n} + y^{2n} > x' \)
or
       \( \ip{A}{B} = \sum_{i} a_{i} b_{i} \).
The spaces you type in a formula are
ignored.  Remember that a letter like
       $x$                   % $ ... $  and  \( ... \)  are equivalent
is a formula when it denotes a mathematical
symbol, and it should be typed as one.

\section{Displayed Text}

Text is displayed by indenting it from the left
margin.  Quotations are commonly displayed.  There
are short quotations
\begin{quote}
   This is a short quotation.  It consists of a
   single paragraph of text.  See how it is formatted.
\end{quote}
and longer ones.
\begin{quotation}
   This is a longer quotation.  It consists of two
   paragraphs of text, neither of which are
   particularly interesting.

   This is the second paragraph of the quotation.  It
   is just as dull as the first paragraph.
\end{quotation}
Another frequently-displayed structure is a list.
The following is an example of an \emph{itemized}
list.
\begin{itemize}
   \item This is the first item of an itemized list.
         Each item in the list is marked with a ``tick''.
         You don't have to worry about what kind of tick
         mark is used.

   \item This is the second item of the list.  It
         contains another list nested inside it.  The inner
         list is an \emph{enumerated} list.
         \begin{enumerate}
            \item This is the first item of an enumerated
                  list that is nested within the itemized list.

            \item This is the second item of the inner list.
                  \LaTeX\ allows you to nest lists deeper than
                  you really should.
         \end{enumerate}
         This is the rest of the second item of the outer
         list.  It is no more interesting than any other
         part of the item.
   \item This is the third item of the list.
\end{itemize}
You can even display poetry.
\begin{verse}
   There is an environment
    for verse \\             % The \\ command separates lines
   Whose features some poets % within a stanza.
   will curse.

                             % One or more blank lines separate stanzas.

   For instead of making\\
   Them do \emph{all} line breaking, \\
   It allows them to put too many words on a line when they'd rather be
   forced to be terse.
\end{verse}

Mathematical formulas may also be displayed.  A
displayed formula
is
one-line long; multiline
formulas require special formatting instructions.
   \[  \ip{\Gamma}{\psi'} = x'' + y^{2} + z_{i}^{n}\]
Don't start a paragraph with a displayed equation,
nor make one a paragraph by itself.

\end{document}               % End of document.

\section{Bench Huge Section 831}
% synthetic bench section 831
% This is a sample LaTeX input file.  (Version of 12 August 2004.)
%
% A '%' character causes TeX to ignore all remaining text on the line,
% and is used for comments like this one.

\documentclass{article}      % Specifies the document class

                             % The preamble begins here.
\title{An Example Document}  % Declares the document's title.
\author{Leslie Lamport}      % Declares the author's name.
\date{January 21, 1994}      % Deleting this command produces today's date.

\newcommand{\ip}[2]{(#1, #2)}
                             % Defines \ip{arg1}{arg2} to mean
                             % (arg1, arg2).

%\newcommand{\ip}[2]{\langle #1 | #2\rangle}
                             % This is an alternative definition of
                             % \ip that is commented out.

\begin{document}             % End of preamble and beginning of text.

\maketitle                   % Produces the title.

This is an example input file.  Comparing it with
the output it generates can show you how to
produce a simple document of your own.

\section{Ordinary Text}      % Produces section heading.  Lower-level
                             % sections are begun with similar
                             % \subsection and \subsubsection commands.

The ends  of words and sentences are marked
  by   spaces. It  doesn't matter how many
spaces    you type; one is as good as 100.  The
end of   a line counts as a space.

One   or more   blank lines denote the  end
of  a paragraph.

Since any number of consecutive spaces are treated
like a single one, the formatting of the input
file makes no difference to
      \LaTeX,                % The \LaTeX command generates the LaTeX logo.
but it makes a difference to you.  When you use
\LaTeX, making your input file as easy to read
as possible will be a great help as you write
your document and when you change it.  This sample
file shows how you can add comments to your own input
file.

Because printing is different from typewriting,
there are a number of things that you have to do
differently when preparing an input file than if
you were just typing the document directly.
Quotation marks like
       ``this''
have to be handled specially, as do quotes within
quotes:
       ``\,`this'            % \, separates the double and single quote.
        is what I just
        wrote, not  `that'\,''.

Dashes come in three sizes: an
       intra-word
dash, a medium dash for number ranges like
       1--2,
and a punctuation
       dash---like
this.

A sentence-ending space should be larger than the
space between words within a sentence.  You
sometimes have to type special commands in
conjunction with punctuation characters to get
this right, as in the following sentence.
       Gnats, gnus, etc.\ all  % `\ ' makes an inter-word space.
       begin with G\@.         % \@ marks end-of-sentence punctuation.
You should check the spaces after periods when
reading your output to make sure you haven't
forgotten any special cases.  Generating an
ellipsis
       \ldots\               % `\ ' is needed after `\ldots' because TeX
                             % ignores spaces after command names like \ldots
                             % made from \ + letters.
                             %
                             % Note how a `%' character causes TeX to ignore
                             % the end of the input line, so these blank lines
                             % do not start a new paragraph.
                             %
with the right spacing around the periods requires
a special command.

\LaTeX\ interprets some common characters as
commands, so you must type special commands to
generate them.  These characters include the
following:
       \$ \& \% \# \{ and \}.

In printing, text is usually emphasized with an
       \emph{italic}
type style.

\begin{em}
   A long segment of text can also be emphasized
   in this way.  Text within such a segment can be
   given \emph{additional} emphasis.
\end{em}

It is sometimes necessary to prevent \LaTeX\ from
breaking a line where it might otherwise do so.
This may be at a space, as between the ``Mr.''\ and
``Jones'' in
       ``Mr.~Jones'',        % ~ produces an unbreakable interword space.
or within a word---especially when the word is a
symbol like
       \mbox{\emph{itemnum}}
that makes little sense when hyphenated across
lines.

Footnotes\footnote{This is an example of a footnote.}
pose no problem.

\LaTeX\ is good at typesetting mathematical formulas
like
       \( x-3y + z = 7 \)
or
       \( a_{1} > x^{2n} + y^{2n} > x' \)
or
       \( \ip{A}{B} = \sum_{i} a_{i} b_{i} \).
The spaces you type in a formula are
ignored.  Remember that a letter like
       $x$                   % $ ... $  and  \( ... \)  are equivalent
is a formula when it denotes a mathematical
symbol, and it should be typed as one.

\section{Displayed Text}

Text is displayed by indenting it from the left
margin.  Quotations are commonly displayed.  There
are short quotations
\begin{quote}
   This is a short quotation.  It consists of a
   single paragraph of text.  See how it is formatted.
\end{quote}
and longer ones.
\begin{quotation}
   This is a longer quotation.  It consists of two
   paragraphs of text, neither of which are
   particularly interesting.

   This is the second paragraph of the quotation.  It
   is just as dull as the first paragraph.
\end{quotation}
Another frequently-displayed structure is a list.
The following is an example of an \emph{itemized}
list.
\begin{itemize}
   \item This is the first item of an itemized list.
         Each item in the list is marked with a ``tick''.
         You don't have to worry about what kind of tick
         mark is used.

   \item This is the second item of the list.  It
         contains another list nested inside it.  The inner
         list is an \emph{enumerated} list.
         \begin{enumerate}
            \item This is the first item of an enumerated
                  list that is nested within the itemized list.

            \item This is the second item of the inner list.
                  \LaTeX\ allows you to nest lists deeper than
                  you really should.
         \end{enumerate}
         This is the rest of the second item of the outer
         list.  It is no more interesting than any other
         part of the item.
   \item This is the third item of the list.
\end{itemize}
You can even display poetry.
\begin{verse}
   There is an environment
    for verse \\             % The \\ command separates lines
   Whose features some poets % within a stanza.
   will curse.

                             % One or more blank lines separate stanzas.

   For instead of making\\
   Them do \emph{all} line breaking, \\
   It allows them to put too many words on a line when they'd rather be
   forced to be terse.
\end{verse}

Mathematical formulas may also be displayed.  A
displayed formula
is
one-line long; multiline
formulas require special formatting instructions.
   \[  \ip{\Gamma}{\psi'} = x'' + y^{2} + z_{i}^{n}\]
Don't start a paragraph with a displayed equation,
nor make one a paragraph by itself.

\end{document}               % End of document.

\section{Bench Huge Section 832}
% synthetic bench section 832
% This is a sample LaTeX input file.  (Version of 12 August 2004.)
%
% A '%' character causes TeX to ignore all remaining text on the line,
% and is used for comments like this one.

\documentclass{article}      % Specifies the document class

                             % The preamble begins here.
\title{An Example Document}  % Declares the document's title.
\author{Leslie Lamport}      % Declares the author's name.
\date{January 21, 1994}      % Deleting this command produces today's date.

\newcommand{\ip}[2]{(#1, #2)}
                             % Defines \ip{arg1}{arg2} to mean
                             % (arg1, arg2).

%\newcommand{\ip}[2]{\langle #1 | #2\rangle}
                             % This is an alternative definition of
                             % \ip that is commented out.

\begin{document}             % End of preamble and beginning of text.

\maketitle                   % Produces the title.

This is an example input file.  Comparing it with
the output it generates can show you how to
produce a simple document of your own.

\section{Ordinary Text}      % Produces section heading.  Lower-level
                             % sections are begun with similar
                             % \subsection and \subsubsection commands.

The ends  of words and sentences are marked
  by   spaces. It  doesn't matter how many
spaces    you type; one is as good as 100.  The
end of   a line counts as a space.

One   or more   blank lines denote the  end
of  a paragraph.

Since any number of consecutive spaces are treated
like a single one, the formatting of the input
file makes no difference to
      \LaTeX,                % The \LaTeX command generates the LaTeX logo.
but it makes a difference to you.  When you use
\LaTeX, making your input file as easy to read
as possible will be a great help as you write
your document and when you change it.  This sample
file shows how you can add comments to your own input
file.

Because printing is different from typewriting,
there are a number of things that you have to do
differently when preparing an input file than if
you were just typing the document directly.
Quotation marks like
       ``this''
have to be handled specially, as do quotes within
quotes:
       ``\,`this'            % \, separates the double and single quote.
        is what I just
        wrote, not  `that'\,''.

Dashes come in three sizes: an
       intra-word
dash, a medium dash for number ranges like
       1--2,
and a punctuation
       dash---like
this.

A sentence-ending space should be larger than the
space between words within a sentence.  You
sometimes have to type special commands in
conjunction with punctuation characters to get
this right, as in the following sentence.
       Gnats, gnus, etc.\ all  % `\ ' makes an inter-word space.
       begin with G\@.         % \@ marks end-of-sentence punctuation.
You should check the spaces after periods when
reading your output to make sure you haven't
forgotten any special cases.  Generating an
ellipsis
       \ldots\               % `\ ' is needed after `\ldots' because TeX
                             % ignores spaces after command names like \ldots
                             % made from \ + letters.
                             %
                             % Note how a `%' character causes TeX to ignore
                             % the end of the input line, so these blank lines
                             % do not start a new paragraph.
                             %
with the right spacing around the periods requires
a special command.

\LaTeX\ interprets some common characters as
commands, so you must type special commands to
generate them.  These characters include the
following:
       \$ \& \% \# \{ and \}.

In printing, text is usually emphasized with an
       \emph{italic}
type style.

\begin{em}
   A long segment of text can also be emphasized
   in this way.  Text within such a segment can be
   given \emph{additional} emphasis.
\end{em}

It is sometimes necessary to prevent \LaTeX\ from
breaking a line where it might otherwise do so.
This may be at a space, as between the ``Mr.''\ and
``Jones'' in
       ``Mr.~Jones'',        % ~ produces an unbreakable interword space.
or within a word---especially when the word is a
symbol like
       \mbox{\emph{itemnum}}
that makes little sense when hyphenated across
lines.

Footnotes\footnote{This is an example of a footnote.}
pose no problem.

\LaTeX\ is good at typesetting mathematical formulas
like
       \( x-3y + z = 7 \)
or
       \( a_{1} > x^{2n} + y^{2n} > x' \)
or
       \( \ip{A}{B} = \sum_{i} a_{i} b_{i} \).
The spaces you type in a formula are
ignored.  Remember that a letter like
       $x$                   % $ ... $  and  \( ... \)  are equivalent
is a formula when it denotes a mathematical
symbol, and it should be typed as one.

\section{Displayed Text}

Text is displayed by indenting it from the left
margin.  Quotations are commonly displayed.  There
are short quotations
\begin{quote}
   This is a short quotation.  It consists of a
   single paragraph of text.  See how it is formatted.
\end{quote}
and longer ones.
\begin{quotation}
   This is a longer quotation.  It consists of two
   paragraphs of text, neither of which are
   particularly interesting.

   This is the second paragraph of the quotation.  It
   is just as dull as the first paragraph.
\end{quotation}
Another frequently-displayed structure is a list.
The following is an example of an \emph{itemized}
list.
\begin{itemize}
   \item This is the first item of an itemized list.
         Each item in the list is marked with a ``tick''.
         You don't have to worry about what kind of tick
         mark is used.

   \item This is the second item of the list.  It
         contains another list nested inside it.  The inner
         list is an \emph{enumerated} list.
         \begin{enumerate}
            \item This is the first item of an enumerated
                  list that is nested within the itemized list.

            \item This is the second item of the inner list.
                  \LaTeX\ allows you to nest lists deeper than
                  you really should.
         \end{enumerate}
         This is the rest of the second item of the outer
         list.  It is no more interesting than any other
         part of the item.
   \item This is the third item of the list.
\end{itemize}
You can even display poetry.
\begin{verse}
   There is an environment
    for verse \\             % The \\ command separates lines
   Whose features some poets % within a stanza.
   will curse.

                             % One or more blank lines separate stanzas.

   For instead of making\\
   Them do \emph{all} line breaking, \\
   It allows them to put too many words on a line when they'd rather be
   forced to be terse.
\end{verse}

Mathematical formulas may also be displayed.  A
displayed formula
is
one-line long; multiline
formulas require special formatting instructions.
   \[  \ip{\Gamma}{\psi'} = x'' + y^{2} + z_{i}^{n}\]
Don't start a paragraph with a displayed equation,
nor make one a paragraph by itself.

\end{document}               % End of document.

\section{Bench Huge Section 833}
% synthetic bench section 833
% This is a sample LaTeX input file.  (Version of 12 August 2004.)
%
% A '%' character causes TeX to ignore all remaining text on the line,
% and is used for comments like this one.

\documentclass{article}      % Specifies the document class

                             % The preamble begins here.
\title{An Example Document}  % Declares the document's title.
\author{Leslie Lamport}      % Declares the author's name.
\date{January 21, 1994}      % Deleting this command produces today's date.

\newcommand{\ip}[2]{(#1, #2)}
                             % Defines \ip{arg1}{arg2} to mean
                             % (arg1, arg2).

%\newcommand{\ip}[2]{\langle #1 | #2\rangle}
                             % This is an alternative definition of
                             % \ip that is commented out.

\begin{document}             % End of preamble and beginning of text.

\maketitle                   % Produces the title.

This is an example input file.  Comparing it with
the output it generates can show you how to
produce a simple document of your own.

\section{Ordinary Text}      % Produces section heading.  Lower-level
                             % sections are begun with similar
                             % \subsection and \subsubsection commands.

The ends  of words and sentences are marked
  by   spaces. It  doesn't matter how many
spaces    you type; one is as good as 100.  The
end of   a line counts as a space.

One   or more   blank lines denote the  end
of  a paragraph.

Since any number of consecutive spaces are treated
like a single one, the formatting of the input
file makes no difference to
      \LaTeX,                % The \LaTeX command generates the LaTeX logo.
but it makes a difference to you.  When you use
\LaTeX, making your input file as easy to read
as possible will be a great help as you write
your document and when you change it.  This sample
file shows how you can add comments to your own input
file.

Because printing is different from typewriting,
there are a number of things that you have to do
differently when preparing an input file than if
you were just typing the document directly.
Quotation marks like
       ``this''
have to be handled specially, as do quotes within
quotes:
       ``\,`this'            % \, separates the double and single quote.
        is what I just
        wrote, not  `that'\,''.

Dashes come in three sizes: an
       intra-word
dash, a medium dash for number ranges like
       1--2,
and a punctuation
       dash---like
this.

A sentence-ending space should be larger than the
space between words within a sentence.  You
sometimes have to type special commands in
conjunction with punctuation characters to get
this right, as in the following sentence.
       Gnats, gnus, etc.\ all  % `\ ' makes an inter-word space.
       begin with G\@.         % \@ marks end-of-sentence punctuation.
You should check the spaces after periods when
reading your output to make sure you haven't
forgotten any special cases.  Generating an
ellipsis
       \ldots\               % `\ ' is needed after `\ldots' because TeX
                             % ignores spaces after command names like \ldots
                             % made from \ + letters.
                             %
                             % Note how a `%' character causes TeX to ignore
                             % the end of the input line, so these blank lines
                             % do not start a new paragraph.
                             %
with the right spacing around the periods requires
a special command.

\LaTeX\ interprets some common characters as
commands, so you must type special commands to
generate them.  These characters include the
following:
       \$ \& \% \# \{ and \}.

In printing, text is usually emphasized with an
       \emph{italic}
type style.

\begin{em}
   A long segment of text can also be emphasized
   in this way.  Text within such a segment can be
   given \emph{additional} emphasis.
\end{em}

It is sometimes necessary to prevent \LaTeX\ from
breaking a line where it might otherwise do so.
This may be at a space, as between the ``Mr.''\ and
``Jones'' in
       ``Mr.~Jones'',        % ~ produces an unbreakable interword space.
or within a word---especially when the word is a
symbol like
       \mbox{\emph{itemnum}}
that makes little sense when hyphenated across
lines.

Footnotes\footnote{This is an example of a footnote.}
pose no problem.

\LaTeX\ is good at typesetting mathematical formulas
like
       \( x-3y + z = 7 \)
or
       \( a_{1} > x^{2n} + y^{2n} > x' \)
or
       \( \ip{A}{B} = \sum_{i} a_{i} b_{i} \).
The spaces you type in a formula are
ignored.  Remember that a letter like
       $x$                   % $ ... $  and  \( ... \)  are equivalent
is a formula when it denotes a mathematical
symbol, and it should be typed as one.

\section{Displayed Text}

Text is displayed by indenting it from the left
margin.  Quotations are commonly displayed.  There
are short quotations
\begin{quote}
   This is a short quotation.  It consists of a
   single paragraph of text.  See how it is formatted.
\end{quote}
and longer ones.
\begin{quotation}
   This is a longer quotation.  It consists of two
   paragraphs of text, neither of which are
   particularly interesting.

   This is the second paragraph of the quotation.  It
   is just as dull as the first paragraph.
\end{quotation}
Another frequently-displayed structure is a list.
The following is an example of an \emph{itemized}
list.
\begin{itemize}
   \item This is the first item of an itemized list.
         Each item in the list is marked with a ``tick''.
         You don't have to worry about what kind of tick
         mark is used.

   \item This is the second item of the list.  It
         contains another list nested inside it.  The inner
         list is an \emph{enumerated} list.
         \begin{enumerate}
            \item This is the first item of an enumerated
                  list that is nested within the itemized list.

            \item This is the second item of the inner list.
                  \LaTeX\ allows you to nest lists deeper than
                  you really should.
         \end{enumerate}
         This is the rest of the second item of the outer
         list.  It is no more interesting than any other
         part of the item.
   \item This is the third item of the list.
\end{itemize}
You can even display poetry.
\begin{verse}
   There is an environment
    for verse \\             % The \\ command separates lines
   Whose features some poets % within a stanza.
   will curse.

                             % One or more blank lines separate stanzas.

   For instead of making\\
   Them do \emph{all} line breaking, \\
   It allows them to put too many words on a line when they'd rather be
   forced to be terse.
\end{verse}

Mathematical formulas may also be displayed.  A
displayed formula
is
one-line long; multiline
formulas require special formatting instructions.
   \[  \ip{\Gamma}{\psi'} = x'' + y^{2} + z_{i}^{n}\]
Don't start a paragraph with a displayed equation,
nor make one a paragraph by itself.

\end{document}               % End of document.

\section{Bench Huge Section 834}
% synthetic bench section 834
% This is a sample LaTeX input file.  (Version of 12 August 2004.)
%
% A '%' character causes TeX to ignore all remaining text on the line,
% and is used for comments like this one.

\documentclass{article}      % Specifies the document class

                             % The preamble begins here.
\title{An Example Document}  % Declares the document's title.
\author{Leslie Lamport}      % Declares the author's name.
\date{January 21, 1994}      % Deleting this command produces today's date.

\newcommand{\ip}[2]{(#1, #2)}
                             % Defines \ip{arg1}{arg2} to mean
                             % (arg1, arg2).

%\newcommand{\ip}[2]{\langle #1 | #2\rangle}
                             % This is an alternative definition of
                             % \ip that is commented out.

\begin{document}             % End of preamble and beginning of text.

\maketitle                   % Produces the title.

This is an example input file.  Comparing it with
the output it generates can show you how to
produce a simple document of your own.

\section{Ordinary Text}      % Produces section heading.  Lower-level
                             % sections are begun with similar
                             % \subsection and \subsubsection commands.

The ends  of words and sentences are marked
  by   spaces. It  doesn't matter how many
spaces    you type; one is as good as 100.  The
end of   a line counts as a space.

One   or more   blank lines denote the  end
of  a paragraph.

Since any number of consecutive spaces are treated
like a single one, the formatting of the input
file makes no difference to
      \LaTeX,                % The \LaTeX command generates the LaTeX logo.
but it makes a difference to you.  When you use
\LaTeX, making your input file as easy to read
as possible will be a great help as you write
your document and when you change it.  This sample
file shows how you can add comments to your own input
file.

Because printing is different from typewriting,
there are a number of things that you have to do
differently when preparing an input file than if
you were just typing the document directly.
Quotation marks like
       ``this''
have to be handled specially, as do quotes within
quotes:
       ``\,`this'            % \, separates the double and single quote.
        is what I just
        wrote, not  `that'\,''.

Dashes come in three sizes: an
       intra-word
dash, a medium dash for number ranges like
       1--2,
and a punctuation
       dash---like
this.

A sentence-ending space should be larger than the
space between words within a sentence.  You
sometimes have to type special commands in
conjunction with punctuation characters to get
this right, as in the following sentence.
       Gnats, gnus, etc.\ all  % `\ ' makes an inter-word space.
       begin with G\@.         % \@ marks end-of-sentence punctuation.
You should check the spaces after periods when
reading your output to make sure you haven't
forgotten any special cases.  Generating an
ellipsis
       \ldots\               % `\ ' is needed after `\ldots' because TeX
                             % ignores spaces after command names like \ldots
                             % made from \ + letters.
                             %
                             % Note how a `%' character causes TeX to ignore
                             % the end of the input line, so these blank lines
                             % do not start a new paragraph.
                             %
with the right spacing around the periods requires
a special command.

\LaTeX\ interprets some common characters as
commands, so you must type special commands to
generate them.  These characters include the
following:
       \$ \& \% \# \{ and \}.

In printing, text is usually emphasized with an
       \emph{italic}
type style.

\begin{em}
   A long segment of text can also be emphasized
   in this way.  Text within such a segment can be
   given \emph{additional} emphasis.
\end{em}

It is sometimes necessary to prevent \LaTeX\ from
breaking a line where it might otherwise do so.
This may be at a space, as between the ``Mr.''\ and
``Jones'' in
       ``Mr.~Jones'',        % ~ produces an unbreakable interword space.
or within a word---especially when the word is a
symbol like
       \mbox{\emph{itemnum}}
that makes little sense when hyphenated across
lines.

Footnotes\footnote{This is an example of a footnote.}
pose no problem.

\LaTeX\ is good at typesetting mathematical formulas
like
       \( x-3y + z = 7 \)
or
       \( a_{1} > x^{2n} + y^{2n} > x' \)
or
       \( \ip{A}{B} = \sum_{i} a_{i} b_{i} \).
The spaces you type in a formula are
ignored.  Remember that a letter like
       $x$                   % $ ... $  and  \( ... \)  are equivalent
is a formula when it denotes a mathematical
symbol, and it should be typed as one.

\section{Displayed Text}

Text is displayed by indenting it from the left
margin.  Quotations are commonly displayed.  There
are short quotations
\begin{quote}
   This is a short quotation.  It consists of a
   single paragraph of text.  See how it is formatted.
\end{quote}
and longer ones.
\begin{quotation}
   This is a longer quotation.  It consists of two
   paragraphs of text, neither of which are
   particularly interesting.

   This is the second paragraph of the quotation.  It
   is just as dull as the first paragraph.
\end{quotation}
Another frequently-displayed structure is a list.
The following is an example of an \emph{itemized}
list.
\begin{itemize}
   \item This is the first item of an itemized list.
         Each item in the list is marked with a ``tick''.
         You don't have to worry about what kind of tick
         mark is used.

   \item This is the second item of the list.  It
         contains another list nested inside it.  The inner
         list is an \emph{enumerated} list.
         \begin{enumerate}
            \item This is the first item of an enumerated
                  list that is nested within the itemized list.

            \item This is the second item of the inner list.
                  \LaTeX\ allows you to nest lists deeper than
                  you really should.
         \end{enumerate}
         This is the rest of the second item of the outer
         list.  It is no more interesting than any other
         part of the item.
   \item This is the third item of the list.
\end{itemize}
You can even display poetry.
\begin{verse}
   There is an environment
    for verse \\             % The \\ command separates lines
   Whose features some poets % within a stanza.
   will curse.

                             % One or more blank lines separate stanzas.

   For instead of making\\
   Them do \emph{all} line breaking, \\
   It allows them to put too many words on a line when they'd rather be
   forced to be terse.
\end{verse}

Mathematical formulas may also be displayed.  A
displayed formula
is
one-line long; multiline
formulas require special formatting instructions.
   \[  \ip{\Gamma}{\psi'} = x'' + y^{2} + z_{i}^{n}\]
Don't start a paragraph with a displayed equation,
nor make one a paragraph by itself.

\end{document}               % End of document.

\section{Bench Huge Section 835}
% synthetic bench section 835
% This is a sample LaTeX input file.  (Version of 12 August 2004.)
%
% A '%' character causes TeX to ignore all remaining text on the line,
% and is used for comments like this one.

\documentclass{article}      % Specifies the document class

                             % The preamble begins here.
\title{An Example Document}  % Declares the document's title.
\author{Leslie Lamport}      % Declares the author's name.
\date{January 21, 1994}      % Deleting this command produces today's date.

\newcommand{\ip}[2]{(#1, #2)}
                             % Defines \ip{arg1}{arg2} to mean
                             % (arg1, arg2).

%\newcommand{\ip}[2]{\langle #1 | #2\rangle}
                             % This is an alternative definition of
                             % \ip that is commented out.

\begin{document}             % End of preamble and beginning of text.

\maketitle                   % Produces the title.

This is an example input file.  Comparing it with
the output it generates can show you how to
produce a simple document of your own.

\section{Ordinary Text}      % Produces section heading.  Lower-level
                             % sections are begun with similar
                             % \subsection and \subsubsection commands.

The ends  of words and sentences are marked
  by   spaces. It  doesn't matter how many
spaces    you type; one is as good as 100.  The
end of   a line counts as a space.

One   or more   blank lines denote the  end
of  a paragraph.

Since any number of consecutive spaces are treated
like a single one, the formatting of the input
file makes no difference to
      \LaTeX,                % The \LaTeX command generates the LaTeX logo.
but it makes a difference to you.  When you use
\LaTeX, making your input file as easy to read
as possible will be a great help as you write
your document and when you change it.  This sample
file shows how you can add comments to your own input
file.

Because printing is different from typewriting,
there are a number of things that you have to do
differently when preparing an input file than if
you were just typing the document directly.
Quotation marks like
       ``this''
have to be handled specially, as do quotes within
quotes:
       ``\,`this'            % \, separates the double and single quote.
        is what I just
        wrote, not  `that'\,''.

Dashes come in three sizes: an
       intra-word
dash, a medium dash for number ranges like
       1--2,
and a punctuation
       dash---like
this.

A sentence-ending space should be larger than the
space between words within a sentence.  You
sometimes have to type special commands in
conjunction with punctuation characters to get
this right, as in the following sentence.
       Gnats, gnus, etc.\ all  % `\ ' makes an inter-word space.
       begin with G\@.         % \@ marks end-of-sentence punctuation.
You should check the spaces after periods when
reading your output to make sure you haven't
forgotten any special cases.  Generating an
ellipsis
       \ldots\               % `\ ' is needed after `\ldots' because TeX
                             % ignores spaces after command names like \ldots
                             % made from \ + letters.
                             %
                             % Note how a `%' character causes TeX to ignore
                             % the end of the input line, so these blank lines
                             % do not start a new paragraph.
                             %
with the right spacing around the periods requires
a special command.

\LaTeX\ interprets some common characters as
commands, so you must type special commands to
generate them.  These characters include the
following:
       \$ \& \% \# \{ and \}.

In printing, text is usually emphasized with an
       \emph{italic}
type style.

\begin{em}
   A long segment of text can also be emphasized
   in this way.  Text within such a segment can be
   given \emph{additional} emphasis.
\end{em}

It is sometimes necessary to prevent \LaTeX\ from
breaking a line where it might otherwise do so.
This may be at a space, as between the ``Mr.''\ and
``Jones'' in
       ``Mr.~Jones'',        % ~ produces an unbreakable interword space.
or within a word---especially when the word is a
symbol like
       \mbox{\emph{itemnum}}
that makes little sense when hyphenated across
lines.

Footnotes\footnote{This is an example of a footnote.}
pose no problem.

\LaTeX\ is good at typesetting mathematical formulas
like
       \( x-3y + z = 7 \)
or
       \( a_{1} > x^{2n} + y^{2n} > x' \)
or
       \( \ip{A}{B} = \sum_{i} a_{i} b_{i} \).
The spaces you type in a formula are
ignored.  Remember that a letter like
       $x$                   % $ ... $  and  \( ... \)  are equivalent
is a formula when it denotes a mathematical
symbol, and it should be typed as one.

\section{Displayed Text}

Text is displayed by indenting it from the left
margin.  Quotations are commonly displayed.  There
are short quotations
\begin{quote}
   This is a short quotation.  It consists of a
   single paragraph of text.  See how it is formatted.
\end{quote}
and longer ones.
\begin{quotation}
   This is a longer quotation.  It consists of two
   paragraphs of text, neither of which are
   particularly interesting.

   This is the second paragraph of the quotation.  It
   is just as dull as the first paragraph.
\end{quotation}
Another frequently-displayed structure is a list.
The following is an example of an \emph{itemized}
list.
\begin{itemize}
   \item This is the first item of an itemized list.
         Each item in the list is marked with a ``tick''.
         You don't have to worry about what kind of tick
         mark is used.

   \item This is the second item of the list.  It
         contains another list nested inside it.  The inner
         list is an \emph{enumerated} list.
         \begin{enumerate}
            \item This is the first item of an enumerated
                  list that is nested within the itemized list.

            \item This is the second item of the inner list.
                  \LaTeX\ allows you to nest lists deeper than
                  you really should.
         \end{enumerate}
         This is the rest of the second item of the outer
         list.  It is no more interesting than any other
         part of the item.
   \item This is the third item of the list.
\end{itemize}
You can even display poetry.
\begin{verse}
   There is an environment
    for verse \\             % The \\ command separates lines
   Whose features some poets % within a stanza.
   will curse.

                             % One or more blank lines separate stanzas.

   For instead of making\\
   Them do \emph{all} line breaking, \\
   It allows them to put too many words on a line when they'd rather be
   forced to be terse.
\end{verse}

Mathematical formulas may also be displayed.  A
displayed formula
is
one-line long; multiline
formulas require special formatting instructions.
   \[  \ip{\Gamma}{\psi'} = x'' + y^{2} + z_{i}^{n}\]
Don't start a paragraph with a displayed equation,
nor make one a paragraph by itself.

\end{document}               % End of document.

\section{Bench Huge Section 836}
% synthetic bench section 836
% This is a sample LaTeX input file.  (Version of 12 August 2004.)
%
% A '%' character causes TeX to ignore all remaining text on the line,
% and is used for comments like this one.

\documentclass{article}      % Specifies the document class

                             % The preamble begins here.
\title{An Example Document}  % Declares the document's title.
\author{Leslie Lamport}      % Declares the author's name.
\date{January 21, 1994}      % Deleting this command produces today's date.

\newcommand{\ip}[2]{(#1, #2)}
                             % Defines \ip{arg1}{arg2} to mean
                             % (arg1, arg2).

%\newcommand{\ip}[2]{\langle #1 | #2\rangle}
                             % This is an alternative definition of
                             % \ip that is commented out.

\begin{document}             % End of preamble and beginning of text.

\maketitle                   % Produces the title.

This is an example input file.  Comparing it with
the output it generates can show you how to
produce a simple document of your own.

\section{Ordinary Text}      % Produces section heading.  Lower-level
                             % sections are begun with similar
                             % \subsection and \subsubsection commands.

The ends  of words and sentences are marked
  by   spaces. It  doesn't matter how many
spaces    you type; one is as good as 100.  The
end of   a line counts as a space.

One   or more   blank lines denote the  end
of  a paragraph.

Since any number of consecutive spaces are treated
like a single one, the formatting of the input
file makes no difference to
      \LaTeX,                % The \LaTeX command generates the LaTeX logo.
but it makes a difference to you.  When you use
\LaTeX, making your input file as easy to read
as possible will be a great help as you write
your document and when you change it.  This sample
file shows how you can add comments to your own input
file.

Because printing is different from typewriting,
there are a number of things that you have to do
differently when preparing an input file than if
you were just typing the document directly.
Quotation marks like
       ``this''
have to be handled specially, as do quotes within
quotes:
       ``\,`this'            % \, separates the double and single quote.
        is what I just
        wrote, not  `that'\,''.

Dashes come in three sizes: an
       intra-word
dash, a medium dash for number ranges like
       1--2,
and a punctuation
       dash---like
this.

A sentence-ending space should be larger than the
space between words within a sentence.  You
sometimes have to type special commands in
conjunction with punctuation characters to get
this right, as in the following sentence.
       Gnats, gnus, etc.\ all  % `\ ' makes an inter-word space.
       begin with G\@.         % \@ marks end-of-sentence punctuation.
You should check the spaces after periods when
reading your output to make sure you haven't
forgotten any special cases.  Generating an
ellipsis
       \ldots\               % `\ ' is needed after `\ldots' because TeX
                             % ignores spaces after command names like \ldots
                             % made from \ + letters.
                             %
                             % Note how a `%' character causes TeX to ignore
                             % the end of the input line, so these blank lines
                             % do not start a new paragraph.
                             %
with the right spacing around the periods requires
a special command.

\LaTeX\ interprets some common characters as
commands, so you must type special commands to
generate them.  These characters include the
following:
       \$ \& \% \# \{ and \}.

In printing, text is usually emphasized with an
       \emph{italic}
type style.

\begin{em}
   A long segment of text can also be emphasized
   in this way.  Text within such a segment can be
   given \emph{additional} emphasis.
\end{em}

It is sometimes necessary to prevent \LaTeX\ from
breaking a line where it might otherwise do so.
This may be at a space, as between the ``Mr.''\ and
``Jones'' in
       ``Mr.~Jones'',        % ~ produces an unbreakable interword space.
or within a word---especially when the word is a
symbol like
       \mbox{\emph{itemnum}}
that makes little sense when hyphenated across
lines.

Footnotes\footnote{This is an example of a footnote.}
pose no problem.

\LaTeX\ is good at typesetting mathematical formulas
like
       \( x-3y + z = 7 \)
or
       \( a_{1} > x^{2n} + y^{2n} > x' \)
or
       \( \ip{A}{B} = \sum_{i} a_{i} b_{i} \).
The spaces you type in a formula are
ignored.  Remember that a letter like
       $x$                   % $ ... $  and  \( ... \)  are equivalent
is a formula when it denotes a mathematical
symbol, and it should be typed as one.

\section{Displayed Text}

Text is displayed by indenting it from the left
margin.  Quotations are commonly displayed.  There
are short quotations
\begin{quote}
   This is a short quotation.  It consists of a
   single paragraph of text.  See how it is formatted.
\end{quote}
and longer ones.
\begin{quotation}
   This is a longer quotation.  It consists of two
   paragraphs of text, neither of which are
   particularly interesting.

   This is the second paragraph of the quotation.  It
   is just as dull as the first paragraph.
\end{quotation}
Another frequently-displayed structure is a list.
The following is an example of an \emph{itemized}
list.
\begin{itemize}
   \item This is the first item of an itemized list.
         Each item in the list is marked with a ``tick''.
         You don't have to worry about what kind of tick
         mark is used.

   \item This is the second item of the list.  It
         contains another list nested inside it.  The inner
         list is an \emph{enumerated} list.
         \begin{enumerate}
            \item This is the first item of an enumerated
                  list that is nested within the itemized list.

            \item This is the second item of the inner list.
                  \LaTeX\ allows you to nest lists deeper than
                  you really should.
         \end{enumerate}
         This is the rest of the second item of the outer
         list.  It is no more interesting than any other
         part of the item.
   \item This is the third item of the list.
\end{itemize}
You can even display poetry.
\begin{verse}
   There is an environment
    for verse \\             % The \\ command separates lines
   Whose features some poets % within a stanza.
   will curse.

                             % One or more blank lines separate stanzas.

   For instead of making\\
   Them do \emph{all} line breaking, \\
   It allows them to put too many words on a line when they'd rather be
   forced to be terse.
\end{verse}

Mathematical formulas may also be displayed.  A
displayed formula
is
one-line long; multiline
formulas require special formatting instructions.
   \[  \ip{\Gamma}{\psi'} = x'' + y^{2} + z_{i}^{n}\]
Don't start a paragraph with a displayed equation,
nor make one a paragraph by itself.

\end{document}               % End of document.

\section{Bench Huge Section 837}
% synthetic bench section 837
% This is a sample LaTeX input file.  (Version of 12 August 2004.)
%
% A '%' character causes TeX to ignore all remaining text on the line,
% and is used for comments like this one.

\documentclass{article}      % Specifies the document class

                             % The preamble begins here.
\title{An Example Document}  % Declares the document's title.
\author{Leslie Lamport}      % Declares the author's name.
\date{January 21, 1994}      % Deleting this command produces today's date.

\newcommand{\ip}[2]{(#1, #2)}
                             % Defines \ip{arg1}{arg2} to mean
                             % (arg1, arg2).

%\newcommand{\ip}[2]{\langle #1 | #2\rangle}
                             % This is an alternative definition of
                             % \ip that is commented out.

\begin{document}             % End of preamble and beginning of text.

\maketitle                   % Produces the title.

This is an example input file.  Comparing it with
the output it generates can show you how to
produce a simple document of your own.

\section{Ordinary Text}      % Produces section heading.  Lower-level
                             % sections are begun with similar
                             % \subsection and \subsubsection commands.

The ends  of words and sentences are marked
  by   spaces. It  doesn't matter how many
spaces    you type; one is as good as 100.  The
end of   a line counts as a space.

One   or more   blank lines denote the  end
of  a paragraph.

Since any number of consecutive spaces are treated
like a single one, the formatting of the input
file makes no difference to
      \LaTeX,                % The \LaTeX command generates the LaTeX logo.
but it makes a difference to you.  When you use
\LaTeX, making your input file as easy to read
as possible will be a great help as you write
your document and when you change it.  This sample
file shows how you can add comments to your own input
file.

Because printing is different from typewriting,
there are a number of things that you have to do
differently when preparing an input file than if
you were just typing the document directly.
Quotation marks like
       ``this''
have to be handled specially, as do quotes within
quotes:
       ``\,`this'            % \, separates the double and single quote.
        is what I just
        wrote, not  `that'\,''.

Dashes come in three sizes: an
       intra-word
dash, a medium dash for number ranges like
       1--2,
and a punctuation
       dash---like
this.

A sentence-ending space should be larger than the
space between words within a sentence.  You
sometimes have to type special commands in
conjunction with punctuation characters to get
this right, as in the following sentence.
       Gnats, gnus, etc.\ all  % `\ ' makes an inter-word space.
       begin with G\@.         % \@ marks end-of-sentence punctuation.
You should check the spaces after periods when
reading your output to make sure you haven't
forgotten any special cases.  Generating an
ellipsis
       \ldots\               % `\ ' is needed after `\ldots' because TeX
                             % ignores spaces after command names like \ldots
                             % made from \ + letters.
                             %
                             % Note how a `%' character causes TeX to ignore
                             % the end of the input line, so these blank lines
                             % do not start a new paragraph.
                             %
with the right spacing around the periods requires
a special command.

\LaTeX\ interprets some common characters as
commands, so you must type special commands to
generate them.  These characters include the
following:
       \$ \& \% \# \{ and \}.

In printing, text is usually emphasized with an
       \emph{italic}
type style.

\begin{em}
   A long segment of text can also be emphasized
   in this way.  Text within such a segment can be
   given \emph{additional} emphasis.
\end{em}

It is sometimes necessary to prevent \LaTeX\ from
breaking a line where it might otherwise do so.
This may be at a space, as between the ``Mr.''\ and
``Jones'' in
       ``Mr.~Jones'',        % ~ produces an unbreakable interword space.
or within a word---especially when the word is a
symbol like
       \mbox{\emph{itemnum}}
that makes little sense when hyphenated across
lines.

Footnotes\footnote{This is an example of a footnote.}
pose no problem.

\LaTeX\ is good at typesetting mathematical formulas
like
       \( x-3y + z = 7 \)
or
       \( a_{1} > x^{2n} + y^{2n} > x' \)
or
       \( \ip{A}{B} = \sum_{i} a_{i} b_{i} \).
The spaces you type in a formula are
ignored.  Remember that a letter like
       $x$                   % $ ... $  and  \( ... \)  are equivalent
is a formula when it denotes a mathematical
symbol, and it should be typed as one.

\section{Displayed Text}

Text is displayed by indenting it from the left
margin.  Quotations are commonly displayed.  There
are short quotations
\begin{quote}
   This is a short quotation.  It consists of a
   single paragraph of text.  See how it is formatted.
\end{quote}
and longer ones.
\begin{quotation}
   This is a longer quotation.  It consists of two
   paragraphs of text, neither of which are
   particularly interesting.

   This is the second paragraph of the quotation.  It
   is just as dull as the first paragraph.
\end{quotation}
Another frequently-displayed structure is a list.
The following is an example of an \emph{itemized}
list.
\begin{itemize}
   \item This is the first item of an itemized list.
         Each item in the list is marked with a ``tick''.
         You don't have to worry about what kind of tick
         mark is used.

   \item This is the second item of the list.  It
         contains another list nested inside it.  The inner
         list is an \emph{enumerated} list.
         \begin{enumerate}
            \item This is the first item of an enumerated
                  list that is nested within the itemized list.

            \item This is the second item of the inner list.
                  \LaTeX\ allows you to nest lists deeper than
                  you really should.
         \end{enumerate}
         This is the rest of the second item of the outer
         list.  It is no more interesting than any other
         part of the item.
   \item This is the third item of the list.
\end{itemize}
You can even display poetry.
\begin{verse}
   There is an environment
    for verse \\             % The \\ command separates lines
   Whose features some poets % within a stanza.
   will curse.

                             % One or more blank lines separate stanzas.

   For instead of making\\
   Them do \emph{all} line breaking, \\
   It allows them to put too many words on a line when they'd rather be
   forced to be terse.
\end{verse}

Mathematical formulas may also be displayed.  A
displayed formula
is
one-line long; multiline
formulas require special formatting instructions.
   \[  \ip{\Gamma}{\psi'} = x'' + y^{2} + z_{i}^{n}\]
Don't start a paragraph with a displayed equation,
nor make one a paragraph by itself.

\end{document}               % End of document.

\section{Bench Huge Section 838}
% synthetic bench section 838
% This is a sample LaTeX input file.  (Version of 12 August 2004.)
%
% A '%' character causes TeX to ignore all remaining text on the line,
% and is used for comments like this one.

\documentclass{article}      % Specifies the document class

                             % The preamble begins here.
\title{An Example Document}  % Declares the document's title.
\author{Leslie Lamport}      % Declares the author's name.
\date{January 21, 1994}      % Deleting this command produces today's date.

\newcommand{\ip}[2]{(#1, #2)}
                             % Defines \ip{arg1}{arg2} to mean
                             % (arg1, arg2).

%\newcommand{\ip}[2]{\langle #1 | #2\rangle}
                             % This is an alternative definition of
                             % \ip that is commented out.

\begin{document}             % End of preamble and beginning of text.

\maketitle                   % Produces the title.

This is an example input file.  Comparing it with
the output it generates can show you how to
produce a simple document of your own.

\section{Ordinary Text}      % Produces section heading.  Lower-level
                             % sections are begun with similar
                             % \subsection and \subsubsection commands.

The ends  of words and sentences are marked
  by   spaces. It  doesn't matter how many
spaces    you type; one is as good as 100.  The
end of   a line counts as a space.

One   or more   blank lines denote the  end
of  a paragraph.

Since any number of consecutive spaces are treated
like a single one, the formatting of the input
file makes no difference to
      \LaTeX,                % The \LaTeX command generates the LaTeX logo.
but it makes a difference to you.  When you use
\LaTeX, making your input file as easy to read
as possible will be a great help as you write
your document and when you change it.  This sample
file shows how you can add comments to your own input
file.

Because printing is different from typewriting,
there are a number of things that you have to do
differently when preparing an input file than if
you were just typing the document directly.
Quotation marks like
       ``this''
have to be handled specially, as do quotes within
quotes:
       ``\,`this'            % \, separates the double and single quote.
        is what I just
        wrote, not  `that'\,''.

Dashes come in three sizes: an
       intra-word
dash, a medium dash for number ranges like
       1--2,
and a punctuation
       dash---like
this.

A sentence-ending space should be larger than the
space between words within a sentence.  You
sometimes have to type special commands in
conjunction with punctuation characters to get
this right, as in the following sentence.
       Gnats, gnus, etc.\ all  % `\ ' makes an inter-word space.
       begin with G\@.         % \@ marks end-of-sentence punctuation.
You should check the spaces after periods when
reading your output to make sure you haven't
forgotten any special cases.  Generating an
ellipsis
       \ldots\               % `\ ' is needed after `\ldots' because TeX
                             % ignores spaces after command names like \ldots
                             % made from \ + letters.
                             %
                             % Note how a `%' character causes TeX to ignore
                             % the end of the input line, so these blank lines
                             % do not start a new paragraph.
                             %
with the right spacing around the periods requires
a special command.

\LaTeX\ interprets some common characters as
commands, so you must type special commands to
generate them.  These characters include the
following:
       \$ \& \% \# \{ and \}.

In printing, text is usually emphasized with an
       \emph{italic}
type style.

\begin{em}
   A long segment of text can also be emphasized
   in this way.  Text within such a segment can be
   given \emph{additional} emphasis.
\end{em}

It is sometimes necessary to prevent \LaTeX\ from
breaking a line where it might otherwise do so.
This may be at a space, as between the ``Mr.''\ and
``Jones'' in
       ``Mr.~Jones'',        % ~ produces an unbreakable interword space.
or within a word---especially when the word is a
symbol like
       \mbox{\emph{itemnum}}
that makes little sense when hyphenated across
lines.

Footnotes\footnote{This is an example of a footnote.}
pose no problem.

\LaTeX\ is good at typesetting mathematical formulas
like
       \( x-3y + z = 7 \)
or
       \( a_{1} > x^{2n} + y^{2n} > x' \)
or
       \( \ip{A}{B} = \sum_{i} a_{i} b_{i} \).
The spaces you type in a formula are
ignored.  Remember that a letter like
       $x$                   % $ ... $  and  \( ... \)  are equivalent
is a formula when it denotes a mathematical
symbol, and it should be typed as one.

\section{Displayed Text}

Text is displayed by indenting it from the left
margin.  Quotations are commonly displayed.  There
are short quotations
\begin{quote}
   This is a short quotation.  It consists of a
   single paragraph of text.  See how it is formatted.
\end{quote}
and longer ones.
\begin{quotation}
   This is a longer quotation.  It consists of two
   paragraphs of text, neither of which are
   particularly interesting.

   This is the second paragraph of the quotation.  It
   is just as dull as the first paragraph.
\end{quotation}
Another frequently-displayed structure is a list.
The following is an example of an \emph{itemized}
list.
\begin{itemize}
   \item This is the first item of an itemized list.
         Each item in the list is marked with a ``tick''.
         You don't have to worry about what kind of tick
         mark is used.

   \item This is the second item of the list.  It
         contains another list nested inside it.  The inner
         list is an \emph{enumerated} list.
         \begin{enumerate}
            \item This is the first item of an enumerated
                  list that is nested within the itemized list.

            \item This is the second item of the inner list.
                  \LaTeX\ allows you to nest lists deeper than
                  you really should.
         \end{enumerate}
         This is the rest of the second item of the outer
         list.  It is no more interesting than any other
         part of the item.
   \item This is the third item of the list.
\end{itemize}
You can even display poetry.
\begin{verse}
   There is an environment
    for verse \\             % The \\ command separates lines
   Whose features some poets % within a stanza.
   will curse.

                             % One or more blank lines separate stanzas.

   For instead of making\\
   Them do \emph{all} line breaking, \\
   It allows them to put too many words on a line when they'd rather be
   forced to be terse.
\end{verse}

Mathematical formulas may also be displayed.  A
displayed formula
is
one-line long; multiline
formulas require special formatting instructions.
   \[  \ip{\Gamma}{\psi'} = x'' + y^{2} + z_{i}^{n}\]
Don't start a paragraph with a displayed equation,
nor make one a paragraph by itself.

\end{document}               % End of document.

\section{Bench Huge Section 839}
% synthetic bench section 839
% This is a sample LaTeX input file.  (Version of 12 August 2004.)
%
% A '%' character causes TeX to ignore all remaining text on the line,
% and is used for comments like this one.

\documentclass{article}      % Specifies the document class

                             % The preamble begins here.
\title{An Example Document}  % Declares the document's title.
\author{Leslie Lamport}      % Declares the author's name.
\date{January 21, 1994}      % Deleting this command produces today's date.

\newcommand{\ip}[2]{(#1, #2)}
                             % Defines \ip{arg1}{arg2} to mean
                             % (arg1, arg2).

%\newcommand{\ip}[2]{\langle #1 | #2\rangle}
                             % This is an alternative definition of
                             % \ip that is commented out.

\begin{document}             % End of preamble and beginning of text.

\maketitle                   % Produces the title.

This is an example input file.  Comparing it with
the output it generates can show you how to
produce a simple document of your own.

\section{Ordinary Text}      % Produces section heading.  Lower-level
                             % sections are begun with similar
                             % \subsection and \subsubsection commands.

The ends  of words and sentences are marked
  by   spaces. It  doesn't matter how many
spaces    you type; one is as good as 100.  The
end of   a line counts as a space.

One   or more   blank lines denote the  end
of  a paragraph.

Since any number of consecutive spaces are treated
like a single one, the formatting of the input
file makes no difference to
      \LaTeX,                % The \LaTeX command generates the LaTeX logo.
but it makes a difference to you.  When you use
\LaTeX, making your input file as easy to read
as possible will be a great help as you write
your document and when you change it.  This sample
file shows how you can add comments to your own input
file.

Because printing is different from typewriting,
there are a number of things that you have to do
differently when preparing an input file than if
you were just typing the document directly.
Quotation marks like
       ``this''
have to be handled specially, as do quotes within
quotes:
       ``\,`this'            % \, separates the double and single quote.
        is what I just
        wrote, not  `that'\,''.

Dashes come in three sizes: an
       intra-word
dash, a medium dash for number ranges like
       1--2,
and a punctuation
       dash---like
this.

A sentence-ending space should be larger than the
space between words within a sentence.  You
sometimes have to type special commands in
conjunction with punctuation characters to get
this right, as in the following sentence.
       Gnats, gnus, etc.\ all  % `\ ' makes an inter-word space.
       begin with G\@.         % \@ marks end-of-sentence punctuation.
You should check the spaces after periods when
reading your output to make sure you haven't
forgotten any special cases.  Generating an
ellipsis
       \ldots\               % `\ ' is needed after `\ldots' because TeX
                             % ignores spaces after command names like \ldots
                             % made from \ + letters.
                             %
                             % Note how a `%' character causes TeX to ignore
                             % the end of the input line, so these blank lines
                             % do not start a new paragraph.
                             %
with the right spacing around the periods requires
a special command.

\LaTeX\ interprets some common characters as
commands, so you must type special commands to
generate them.  These characters include the
following:
       \$ \& \% \# \{ and \}.

In printing, text is usually emphasized with an
       \emph{italic}
type style.

\begin{em}
   A long segment of text can also be emphasized
   in this way.  Text within such a segment can be
   given \emph{additional} emphasis.
\end{em}

It is sometimes necessary to prevent \LaTeX\ from
breaking a line where it might otherwise do so.
This may be at a space, as between the ``Mr.''\ and
``Jones'' in
       ``Mr.~Jones'',        % ~ produces an unbreakable interword space.
or within a word---especially when the word is a
symbol like
       \mbox{\emph{itemnum}}
that makes little sense when hyphenated across
lines.

Footnotes\footnote{This is an example of a footnote.}
pose no problem.

\LaTeX\ is good at typesetting mathematical formulas
like
       \( x-3y + z = 7 \)
or
       \( a_{1} > x^{2n} + y^{2n} > x' \)
or
       \( \ip{A}{B} = \sum_{i} a_{i} b_{i} \).
The spaces you type in a formula are
ignored.  Remember that a letter like
       $x$                   % $ ... $  and  \( ... \)  are equivalent
is a formula when it denotes a mathematical
symbol, and it should be typed as one.

\section{Displayed Text}

Text is displayed by indenting it from the left
margin.  Quotations are commonly displayed.  There
are short quotations
\begin{quote}
   This is a short quotation.  It consists of a
   single paragraph of text.  See how it is formatted.
\end{quote}
and longer ones.
\begin{quotation}
   This is a longer quotation.  It consists of two
   paragraphs of text, neither of which are
   particularly interesting.

   This is the second paragraph of the quotation.  It
   is just as dull as the first paragraph.
\end{quotation}
Another frequently-displayed structure is a list.
The following is an example of an \emph{itemized}
list.
\begin{itemize}
   \item This is the first item of an itemized list.
         Each item in the list is marked with a ``tick''.
         You don't have to worry about what kind of tick
         mark is used.

   \item This is the second item of the list.  It
         contains another list nested inside it.  The inner
         list is an \emph{enumerated} list.
         \begin{enumerate}
            \item This is the first item of an enumerated
                  list that is nested within the itemized list.

            \item This is the second item of the inner list.
                  \LaTeX\ allows you to nest lists deeper than
                  you really should.
         \end{enumerate}
         This is the rest of the second item of the outer
         list.  It is no more interesting than any other
         part of the item.
   \item This is the third item of the list.
\end{itemize}
You can even display poetry.
\begin{verse}
   There is an environment
    for verse \\             % The \\ command separates lines
   Whose features some poets % within a stanza.
   will curse.

                             % One or more blank lines separate stanzas.

   For instead of making\\
   Them do \emph{all} line breaking, \\
   It allows them to put too many words on a line when they'd rather be
   forced to be terse.
\end{verse}

Mathematical formulas may also be displayed.  A
displayed formula
is
one-line long; multiline
formulas require special formatting instructions.
   \[  \ip{\Gamma}{\psi'} = x'' + y^{2} + z_{i}^{n}\]
Don't start a paragraph with a displayed equation,
nor make one a paragraph by itself.

\end{document}               % End of document.

\section{Bench Huge Section 840}
% synthetic bench section 840
% This is a sample LaTeX input file.  (Version of 12 August 2004.)
%
% A '%' character causes TeX to ignore all remaining text on the line,
% and is used for comments like this one.

\documentclass{article}      % Specifies the document class

                             % The preamble begins here.
\title{An Example Document}  % Declares the document's title.
\author{Leslie Lamport}      % Declares the author's name.
\date{January 21, 1994}      % Deleting this command produces today's date.

\newcommand{\ip}[2]{(#1, #2)}
                             % Defines \ip{arg1}{arg2} to mean
                             % (arg1, arg2).

%\newcommand{\ip}[2]{\langle #1 | #2\rangle}
                             % This is an alternative definition of
                             % \ip that is commented out.

\begin{document}             % End of preamble and beginning of text.

\maketitle                   % Produces the title.

This is an example input file.  Comparing it with
the output it generates can show you how to
produce a simple document of your own.

\section{Ordinary Text}      % Produces section heading.  Lower-level
                             % sections are begun with similar
                             % \subsection and \subsubsection commands.

The ends  of words and sentences are marked
  by   spaces. It  doesn't matter how many
spaces    you type; one is as good as 100.  The
end of   a line counts as a space.

One   or more   blank lines denote the  end
of  a paragraph.

Since any number of consecutive spaces are treated
like a single one, the formatting of the input
file makes no difference to
      \LaTeX,                % The \LaTeX command generates the LaTeX logo.
but it makes a difference to you.  When you use
\LaTeX, making your input file as easy to read
as possible will be a great help as you write
your document and when you change it.  This sample
file shows how you can add comments to your own input
file.

Because printing is different from typewriting,
there are a number of things that you have to do
differently when preparing an input file than if
you were just typing the document directly.
Quotation marks like
       ``this''
have to be handled specially, as do quotes within
quotes:
       ``\,`this'            % \, separates the double and single quote.
        is what I just
        wrote, not  `that'\,''.

Dashes come in three sizes: an
       intra-word
dash, a medium dash for number ranges like
       1--2,
and a punctuation
       dash---like
this.

A sentence-ending space should be larger than the
space between words within a sentence.  You
sometimes have to type special commands in
conjunction with punctuation characters to get
this right, as in the following sentence.
       Gnats, gnus, etc.\ all  % `\ ' makes an inter-word space.
       begin with G\@.         % \@ marks end-of-sentence punctuation.
You should check the spaces after periods when
reading your output to make sure you haven't
forgotten any special cases.  Generating an
ellipsis
       \ldots\               % `\ ' is needed after `\ldots' because TeX
                             % ignores spaces after command names like \ldots
                             % made from \ + letters.
                             %
                             % Note how a `%' character causes TeX to ignore
                             % the end of the input line, so these blank lines
                             % do not start a new paragraph.
                             %
with the right spacing around the periods requires
a special command.

\LaTeX\ interprets some common characters as
commands, so you must type special commands to
generate them.  These characters include the
following:
       \$ \& \% \# \{ and \}.

In printing, text is usually emphasized with an
       \emph{italic}
type style.

\begin{em}
   A long segment of text can also be emphasized
   in this way.  Text within such a segment can be
   given \emph{additional} emphasis.
\end{em}

It is sometimes necessary to prevent \LaTeX\ from
breaking a line where it might otherwise do so.
This may be at a space, as between the ``Mr.''\ and
``Jones'' in
       ``Mr.~Jones'',        % ~ produces an unbreakable interword space.
or within a word---especially when the word is a
symbol like
       \mbox{\emph{itemnum}}
that makes little sense when hyphenated across
lines.

Footnotes\footnote{This is an example of a footnote.}
pose no problem.

\LaTeX\ is good at typesetting mathematical formulas
like
       \( x-3y + z = 7 \)
or
       \( a_{1} > x^{2n} + y^{2n} > x' \)
or
       \( \ip{A}{B} = \sum_{i} a_{i} b_{i} \).
The spaces you type in a formula are
ignored.  Remember that a letter like
       $x$                   % $ ... $  and  \( ... \)  are equivalent
is a formula when it denotes a mathematical
symbol, and it should be typed as one.

\section{Displayed Text}

Text is displayed by indenting it from the left
margin.  Quotations are commonly displayed.  There
are short quotations
\begin{quote}
   This is a short quotation.  It consists of a
   single paragraph of text.  See how it is formatted.
\end{quote}
and longer ones.
\begin{quotation}
   This is a longer quotation.  It consists of two
   paragraphs of text, neither of which are
   particularly interesting.

   This is the second paragraph of the quotation.  It
   is just as dull as the first paragraph.
\end{quotation}
Another frequently-displayed structure is a list.
The following is an example of an \emph{itemized}
list.
\begin{itemize}
   \item This is the first item of an itemized list.
         Each item in the list is marked with a ``tick''.
         You don't have to worry about what kind of tick
         mark is used.

   \item This is the second item of the list.  It
         contains another list nested inside it.  The inner
         list is an \emph{enumerated} list.
         \begin{enumerate}
            \item This is the first item of an enumerated
                  list that is nested within the itemized list.

            \item This is the second item of the inner list.
                  \LaTeX\ allows you to nest lists deeper than
                  you really should.
         \end{enumerate}
         This is the rest of the second item of the outer
         list.  It is no more interesting than any other
         part of the item.
   \item This is the third item of the list.
\end{itemize}
You can even display poetry.
\begin{verse}
   There is an environment
    for verse \\             % The \\ command separates lines
   Whose features some poets % within a stanza.
   will curse.

                             % One or more blank lines separate stanzas.

   For instead of making\\
   Them do \emph{all} line breaking, \\
   It allows them to put too many words on a line when they'd rather be
   forced to be terse.
\end{verse}

Mathematical formulas may also be displayed.  A
displayed formula
is
one-line long; multiline
formulas require special formatting instructions.
   \[  \ip{\Gamma}{\psi'} = x'' + y^{2} + z_{i}^{n}\]
Don't start a paragraph with a displayed equation,
nor make one a paragraph by itself.

\end{document}               % End of document.

\section{Bench Huge Section 841}
% synthetic bench section 841
% This is a sample LaTeX input file.  (Version of 12 August 2004.)
%
% A '%' character causes TeX to ignore all remaining text on the line,
% and is used for comments like this one.

\documentclass{article}      % Specifies the document class

                             % The preamble begins here.
\title{An Example Document}  % Declares the document's title.
\author{Leslie Lamport}      % Declares the author's name.
\date{January 21, 1994}      % Deleting this command produces today's date.

\newcommand{\ip}[2]{(#1, #2)}
                             % Defines \ip{arg1}{arg2} to mean
                             % (arg1, arg2).

%\newcommand{\ip}[2]{\langle #1 | #2\rangle}
                             % This is an alternative definition of
                             % \ip that is commented out.

\begin{document}             % End of preamble and beginning of text.

\maketitle                   % Produces the title.

This is an example input file.  Comparing it with
the output it generates can show you how to
produce a simple document of your own.

\section{Ordinary Text}      % Produces section heading.  Lower-level
                             % sections are begun with similar
                             % \subsection and \subsubsection commands.

The ends  of words and sentences are marked
  by   spaces. It  doesn't matter how many
spaces    you type; one is as good as 100.  The
end of   a line counts as a space.

One   or more   blank lines denote the  end
of  a paragraph.

Since any number of consecutive spaces are treated
like a single one, the formatting of the input
file makes no difference to
      \LaTeX,                % The \LaTeX command generates the LaTeX logo.
but it makes a difference to you.  When you use
\LaTeX, making your input file as easy to read
as possible will be a great help as you write
your document and when you change it.  This sample
file shows how you can add comments to your own input
file.

Because printing is different from typewriting,
there are a number of things that you have to do
differently when preparing an input file than if
you were just typing the document directly.
Quotation marks like
       ``this''
have to be handled specially, as do quotes within
quotes:
       ``\,`this'            % \, separates the double and single quote.
        is what I just
        wrote, not  `that'\,''.

Dashes come in three sizes: an
       intra-word
dash, a medium dash for number ranges like
       1--2,
and a punctuation
       dash---like
this.

A sentence-ending space should be larger than the
space between words within a sentence.  You
sometimes have to type special commands in
conjunction with punctuation characters to get
this right, as in the following sentence.
       Gnats, gnus, etc.\ all  % `\ ' makes an inter-word space.
       begin with G\@.         % \@ marks end-of-sentence punctuation.
You should check the spaces after periods when
reading your output to make sure you haven't
forgotten any special cases.  Generating an
ellipsis
       \ldots\               % `\ ' is needed after `\ldots' because TeX
                             % ignores spaces after command names like \ldots
                             % made from \ + letters.
                             %
                             % Note how a `%' character causes TeX to ignore
                             % the end of the input line, so these blank lines
                             % do not start a new paragraph.
                             %
with the right spacing around the periods requires
a special command.

\LaTeX\ interprets some common characters as
commands, so you must type special commands to
generate them.  These characters include the
following:
       \$ \& \% \# \{ and \}.

In printing, text is usually emphasized with an
       \emph{italic}
type style.

\begin{em}
   A long segment of text can also be emphasized
   in this way.  Text within such a segment can be
   given \emph{additional} emphasis.
\end{em}

It is sometimes necessary to prevent \LaTeX\ from
breaking a line where it might otherwise do so.
This may be at a space, as between the ``Mr.''\ and
``Jones'' in
       ``Mr.~Jones'',        % ~ produces an unbreakable interword space.
or within a word---especially when the word is a
symbol like
       \mbox{\emph{itemnum}}
that makes little sense when hyphenated across
lines.

Footnotes\footnote{This is an example of a footnote.}
pose no problem.

\LaTeX\ is good at typesetting mathematical formulas
like
       \( x-3y + z = 7 \)
or
       \( a_{1} > x^{2n} + y^{2n} > x' \)
or
       \( \ip{A}{B} = \sum_{i} a_{i} b_{i} \).
The spaces you type in a formula are
ignored.  Remember that a letter like
       $x$                   % $ ... $  and  \( ... \)  are equivalent
is a formula when it denotes a mathematical
symbol, and it should be typed as one.

\section{Displayed Text}

Text is displayed by indenting it from the left
margin.  Quotations are commonly displayed.  There
are short quotations
\begin{quote}
   This is a short quotation.  It consists of a
   single paragraph of text.  See how it is formatted.
\end{quote}
and longer ones.
\begin{quotation}
   This is a longer quotation.  It consists of two
   paragraphs of text, neither of which are
   particularly interesting.

   This is the second paragraph of the quotation.  It
   is just as dull as the first paragraph.
\end{quotation}
Another frequently-displayed structure is a list.
The following is an example of an \emph{itemized}
list.
\begin{itemize}
   \item This is the first item of an itemized list.
         Each item in the list is marked with a ``tick''.
         You don't have to worry about what kind of tick
         mark is used.

   \item This is the second item of the list.  It
         contains another list nested inside it.  The inner
         list is an \emph{enumerated} list.
         \begin{enumerate}
            \item This is the first item of an enumerated
                  list that is nested within the itemized list.

            \item This is the second item of the inner list.
                  \LaTeX\ allows you to nest lists deeper than
                  you really should.
         \end{enumerate}
         This is the rest of the second item of the outer
         list.  It is no more interesting than any other
         part of the item.
   \item This is the third item of the list.
\end{itemize}
You can even display poetry.
\begin{verse}
   There is an environment
    for verse \\             % The \\ command separates lines
   Whose features some poets % within a stanza.
   will curse.

                             % One or more blank lines separate stanzas.

   For instead of making\\
   Them do \emph{all} line breaking, \\
   It allows them to put too many words on a line when they'd rather be
   forced to be terse.
\end{verse}

Mathematical formulas may also be displayed.  A
displayed formula
is
one-line long; multiline
formulas require special formatting instructions.
   \[  \ip{\Gamma}{\psi'} = x'' + y^{2} + z_{i}^{n}\]
Don't start a paragraph with a displayed equation,
nor make one a paragraph by itself.

\end{document}               % End of document.

\section{Bench Huge Section 842}
% synthetic bench section 842
% This is a sample LaTeX input file.  (Version of 12 August 2004.)
%
% A '%' character causes TeX to ignore all remaining text on the line,
% and is used for comments like this one.

\documentclass{article}      % Specifies the document class

                             % The preamble begins here.
\title{An Example Document}  % Declares the document's title.
\author{Leslie Lamport}      % Declares the author's name.
\date{January 21, 1994}      % Deleting this command produces today's date.

\newcommand{\ip}[2]{(#1, #2)}
                             % Defines \ip{arg1}{arg2} to mean
                             % (arg1, arg2).

%\newcommand{\ip}[2]{\langle #1 | #2\rangle}
                             % This is an alternative definition of
                             % \ip that is commented out.

\begin{document}             % End of preamble and beginning of text.

\maketitle                   % Produces the title.

This is an example input file.  Comparing it with
the output it generates can show you how to
produce a simple document of your own.

\section{Ordinary Text}      % Produces section heading.  Lower-level
                             % sections are begun with similar
                             % \subsection and \subsubsection commands.

The ends  of words and sentences are marked
  by   spaces. It  doesn't matter how many
spaces    you type; one is as good as 100.  The
end of   a line counts as a space.

One   or more   blank lines denote the  end
of  a paragraph.

Since any number of consecutive spaces are treated
like a single one, the formatting of the input
file makes no difference to
      \LaTeX,                % The \LaTeX command generates the LaTeX logo.
but it makes a difference to you.  When you use
\LaTeX, making your input file as easy to read
as possible will be a great help as you write
your document and when you change it.  This sample
file shows how you can add comments to your own input
file.

Because printing is different from typewriting,
there are a number of things that you have to do
differently when preparing an input file than if
you were just typing the document directly.
Quotation marks like
       ``this''
have to be handled specially, as do quotes within
quotes:
       ``\,`this'            % \, separates the double and single quote.
        is what I just
        wrote, not  `that'\,''.

Dashes come in three sizes: an
       intra-word
dash, a medium dash for number ranges like
       1--2,
and a punctuation
       dash---like
this.

A sentence-ending space should be larger than the
space between words within a sentence.  You
sometimes have to type special commands in
conjunction with punctuation characters to get
this right, as in the following sentence.
       Gnats, gnus, etc.\ all  % `\ ' makes an inter-word space.
       begin with G\@.         % \@ marks end-of-sentence punctuation.
You should check the spaces after periods when
reading your output to make sure you haven't
forgotten any special cases.  Generating an
ellipsis
       \ldots\               % `\ ' is needed after `\ldots' because TeX
                             % ignores spaces after command names like \ldots
                             % made from \ + letters.
                             %
                             % Note how a `%' character causes TeX to ignore
                             % the end of the input line, so these blank lines
                             % do not start a new paragraph.
                             %
with the right spacing around the periods requires
a special command.

\LaTeX\ interprets some common characters as
commands, so you must type special commands to
generate them.  These characters include the
following:
       \$ \& \% \# \{ and \}.

In printing, text is usually emphasized with an
       \emph{italic}
type style.

\begin{em}
   A long segment of text can also be emphasized
   in this way.  Text within such a segment can be
   given \emph{additional} emphasis.
\end{em}

It is sometimes necessary to prevent \LaTeX\ from
breaking a line where it might otherwise do so.
This may be at a space, as between the ``Mr.''\ and
``Jones'' in
       ``Mr.~Jones'',        % ~ produces an unbreakable interword space.
or within a word---especially when the word is a
symbol like
       \mbox{\emph{itemnum}}
that makes little sense when hyphenated across
lines.

Footnotes\footnote{This is an example of a footnote.}
pose no problem.

\LaTeX\ is good at typesetting mathematical formulas
like
       \( x-3y + z = 7 \)
or
       \( a_{1} > x^{2n} + y^{2n} > x' \)
or
       \( \ip{A}{B} = \sum_{i} a_{i} b_{i} \).
The spaces you type in a formula are
ignored.  Remember that a letter like
       $x$                   % $ ... $  and  \( ... \)  are equivalent
is a formula when it denotes a mathematical
symbol, and it should be typed as one.

\section{Displayed Text}

Text is displayed by indenting it from the left
margin.  Quotations are commonly displayed.  There
are short quotations
\begin{quote}
   This is a short quotation.  It consists of a
   single paragraph of text.  See how it is formatted.
\end{quote}
and longer ones.
\begin{quotation}
   This is a longer quotation.  It consists of two
   paragraphs of text, neither of which are
   particularly interesting.

   This is the second paragraph of the quotation.  It
   is just as dull as the first paragraph.
\end{quotation}
Another frequently-displayed structure is a list.
The following is an example of an \emph{itemized}
list.
\begin{itemize}
   \item This is the first item of an itemized list.
         Each item in the list is marked with a ``tick''.
         You don't have to worry about what kind of tick
         mark is used.

   \item This is the second item of the list.  It
         contains another list nested inside it.  The inner
         list is an \emph{enumerated} list.
         \begin{enumerate}
            \item This is the first item of an enumerated
                  list that is nested within the itemized list.

            \item This is the second item of the inner list.
                  \LaTeX\ allows you to nest lists deeper than
                  you really should.
         \end{enumerate}
         This is the rest of the second item of the outer
         list.  It is no more interesting than any other
         part of the item.
   \item This is the third item of the list.
\end{itemize}
You can even display poetry.
\begin{verse}
   There is an environment
    for verse \\             % The \\ command separates lines
   Whose features some poets % within a stanza.
   will curse.

                             % One or more blank lines separate stanzas.

   For instead of making\\
   Them do \emph{all} line breaking, \\
   It allows them to put too many words on a line when they'd rather be
   forced to be terse.
\end{verse}

Mathematical formulas may also be displayed.  A
displayed formula
is
one-line long; multiline
formulas require special formatting instructions.
   \[  \ip{\Gamma}{\psi'} = x'' + y^{2} + z_{i}^{n}\]
Don't start a paragraph with a displayed equation,
nor make one a paragraph by itself.

\end{document}               % End of document.

\section{Bench Huge Section 843}
% synthetic bench section 843
% This is a sample LaTeX input file.  (Version of 12 August 2004.)
%
% A '%' character causes TeX to ignore all remaining text on the line,
% and is used for comments like this one.

\documentclass{article}      % Specifies the document class

                             % The preamble begins here.
\title{An Example Document}  % Declares the document's title.
\author{Leslie Lamport}      % Declares the author's name.
\date{January 21, 1994}      % Deleting this command produces today's date.

\newcommand{\ip}[2]{(#1, #2)}
                             % Defines \ip{arg1}{arg2} to mean
                             % (arg1, arg2).

%\newcommand{\ip}[2]{\langle #1 | #2\rangle}
                             % This is an alternative definition of
                             % \ip that is commented out.

\begin{document}             % End of preamble and beginning of text.

\maketitle                   % Produces the title.

This is an example input file.  Comparing it with
the output it generates can show you how to
produce a simple document of your own.

\section{Ordinary Text}      % Produces section heading.  Lower-level
                             % sections are begun with similar
                             % \subsection and \subsubsection commands.

The ends  of words and sentences are marked
  by   spaces. It  doesn't matter how many
spaces    you type; one is as good as 100.  The
end of   a line counts as a space.

One   or more   blank lines denote the  end
of  a paragraph.

Since any number of consecutive spaces are treated
like a single one, the formatting of the input
file makes no difference to
      \LaTeX,                % The \LaTeX command generates the LaTeX logo.
but it makes a difference to you.  When you use
\LaTeX, making your input file as easy to read
as possible will be a great help as you write
your document and when you change it.  This sample
file shows how you can add comments to your own input
file.

Because printing is different from typewriting,
there are a number of things that you have to do
differently when preparing an input file than if
you were just typing the document directly.
Quotation marks like
       ``this''
have to be handled specially, as do quotes within
quotes:
       ``\,`this'            % \, separates the double and single quote.
        is what I just
        wrote, not  `that'\,''.

Dashes come in three sizes: an
       intra-word
dash, a medium dash for number ranges like
       1--2,
and a punctuation
       dash---like
this.

A sentence-ending space should be larger than the
space between words within a sentence.  You
sometimes have to type special commands in
conjunction with punctuation characters to get
this right, as in the following sentence.
       Gnats, gnus, etc.\ all  % `\ ' makes an inter-word space.
       begin with G\@.         % \@ marks end-of-sentence punctuation.
You should check the spaces after periods when
reading your output to make sure you haven't
forgotten any special cases.  Generating an
ellipsis
       \ldots\               % `\ ' is needed after `\ldots' because TeX
                             % ignores spaces after command names like \ldots
                             % made from \ + letters.
                             %
                             % Note how a `%' character causes TeX to ignore
                             % the end of the input line, so these blank lines
                             % do not start a new paragraph.
                             %
with the right spacing around the periods requires
a special command.

\LaTeX\ interprets some common characters as
commands, so you must type special commands to
generate them.  These characters include the
following:
       \$ \& \% \# \{ and \}.

In printing, text is usually emphasized with an
       \emph{italic}
type style.

\begin{em}
   A long segment of text can also be emphasized
   in this way.  Text within such a segment can be
   given \emph{additional} emphasis.
\end{em}

It is sometimes necessary to prevent \LaTeX\ from
breaking a line where it might otherwise do so.
This may be at a space, as between the ``Mr.''\ and
``Jones'' in
       ``Mr.~Jones'',        % ~ produces an unbreakable interword space.
or within a word---especially when the word is a
symbol like
       \mbox{\emph{itemnum}}
that makes little sense when hyphenated across
lines.

Footnotes\footnote{This is an example of a footnote.}
pose no problem.

\LaTeX\ is good at typesetting mathematical formulas
like
       \( x-3y + z = 7 \)
or
       \( a_{1} > x^{2n} + y^{2n} > x' \)
or
       \( \ip{A}{B} = \sum_{i} a_{i} b_{i} \).
The spaces you type in a formula are
ignored.  Remember that a letter like
       $x$                   % $ ... $  and  \( ... \)  are equivalent
is a formula when it denotes a mathematical
symbol, and it should be typed as one.

\section{Displayed Text}

Text is displayed by indenting it from the left
margin.  Quotations are commonly displayed.  There
are short quotations
\begin{quote}
   This is a short quotation.  It consists of a
   single paragraph of text.  See how it is formatted.
\end{quote}
and longer ones.
\begin{quotation}
   This is a longer quotation.  It consists of two
   paragraphs of text, neither of which are
   particularly interesting.

   This is the second paragraph of the quotation.  It
   is just as dull as the first paragraph.
\end{quotation}
Another frequently-displayed structure is a list.
The following is an example of an \emph{itemized}
list.
\begin{itemize}
   \item This is the first item of an itemized list.
         Each item in the list is marked with a ``tick''.
         You don't have to worry about what kind of tick
         mark is used.

   \item This is the second item of the list.  It
         contains another list nested inside it.  The inner
         list is an \emph{enumerated} list.
         \begin{enumerate}
            \item This is the first item of an enumerated
                  list that is nested within the itemized list.

            \item This is the second item of the inner list.
                  \LaTeX\ allows you to nest lists deeper than
                  you really should.
         \end{enumerate}
         This is the rest of the second item of the outer
         list.  It is no more interesting than any other
         part of the item.
   \item This is the third item of the list.
\end{itemize}
You can even display poetry.
\begin{verse}
   There is an environment
    for verse \\             % The \\ command separates lines
   Whose features some poets % within a stanza.
   will curse.

                             % One or more blank lines separate stanzas.

   For instead of making\\
   Them do \emph{all} line breaking, \\
   It allows them to put too many words on a line when they'd rather be
   forced to be terse.
\end{verse}

Mathematical formulas may also be displayed.  A
displayed formula
is
one-line long; multiline
formulas require special formatting instructions.
   \[  \ip{\Gamma}{\psi'} = x'' + y^{2} + z_{i}^{n}\]
Don't start a paragraph with a displayed equation,
nor make one a paragraph by itself.

\end{document}               % End of document.

\section{Bench Huge Section 844}
% synthetic bench section 844
% This is a sample LaTeX input file.  (Version of 12 August 2004.)
%
% A '%' character causes TeX to ignore all remaining text on the line,
% and is used for comments like this one.

\documentclass{article}      % Specifies the document class

                             % The preamble begins here.
\title{An Example Document}  % Declares the document's title.
\author{Leslie Lamport}      % Declares the author's name.
\date{January 21, 1994}      % Deleting this command produces today's date.

\newcommand{\ip}[2]{(#1, #2)}
                             % Defines \ip{arg1}{arg2} to mean
                             % (arg1, arg2).

%\newcommand{\ip}[2]{\langle #1 | #2\rangle}
                             % This is an alternative definition of
                             % \ip that is commented out.

\begin{document}             % End of preamble and beginning of text.

\maketitle                   % Produces the title.

This is an example input file.  Comparing it with
the output it generates can show you how to
produce a simple document of your own.

\section{Ordinary Text}      % Produces section heading.  Lower-level
                             % sections are begun with similar
                             % \subsection and \subsubsection commands.

The ends  of words and sentences are marked
  by   spaces. It  doesn't matter how many
spaces    you type; one is as good as 100.  The
end of   a line counts as a space.

One   or more   blank lines denote the  end
of  a paragraph.

Since any number of consecutive spaces are treated
like a single one, the formatting of the input
file makes no difference to
      \LaTeX,                % The \LaTeX command generates the LaTeX logo.
but it makes a difference to you.  When you use
\LaTeX, making your input file as easy to read
as possible will be a great help as you write
your document and when you change it.  This sample
file shows how you can add comments to your own input
file.

Because printing is different from typewriting,
there are a number of things that you have to do
differently when preparing an input file than if
you were just typing the document directly.
Quotation marks like
       ``this''
have to be handled specially, as do quotes within
quotes:
       ``\,`this'            % \, separates the double and single quote.
        is what I just
        wrote, not  `that'\,''.

Dashes come in three sizes: an
       intra-word
dash, a medium dash for number ranges like
       1--2,
and a punctuation
       dash---like
this.

A sentence-ending space should be larger than the
space between words within a sentence.  You
sometimes have to type special commands in
conjunction with punctuation characters to get
this right, as in the following sentence.
       Gnats, gnus, etc.\ all  % `\ ' makes an inter-word space.
       begin with G\@.         % \@ marks end-of-sentence punctuation.
You should check the spaces after periods when
reading your output to make sure you haven't
forgotten any special cases.  Generating an
ellipsis
       \ldots\               % `\ ' is needed after `\ldots' because TeX
                             % ignores spaces after command names like \ldots
                             % made from \ + letters.
                             %
                             % Note how a `%' character causes TeX to ignore
                             % the end of the input line, so these blank lines
                             % do not start a new paragraph.
                             %
with the right spacing around the periods requires
a special command.

\LaTeX\ interprets some common characters as
commands, so you must type special commands to
generate them.  These characters include the
following:
       \$ \& \% \# \{ and \}.

In printing, text is usually emphasized with an
       \emph{italic}
type style.

\begin{em}
   A long segment of text can also be emphasized
   in this way.  Text within such a segment can be
   given \emph{additional} emphasis.
\end{em}

It is sometimes necessary to prevent \LaTeX\ from
breaking a line where it might otherwise do so.
This may be at a space, as between the ``Mr.''\ and
``Jones'' in
       ``Mr.~Jones'',        % ~ produces an unbreakable interword space.
or within a word---especially when the word is a
symbol like
       \mbox{\emph{itemnum}}
that makes little sense when hyphenated across
lines.

Footnotes\footnote{This is an example of a footnote.}
pose no problem.

\LaTeX\ is good at typesetting mathematical formulas
like
       \( x-3y + z = 7 \)
or
       \( a_{1} > x^{2n} + y^{2n} > x' \)
or
       \( \ip{A}{B} = \sum_{i} a_{i} b_{i} \).
The spaces you type in a formula are
ignored.  Remember that a letter like
       $x$                   % $ ... $  and  \( ... \)  are equivalent
is a formula when it denotes a mathematical
symbol, and it should be typed as one.

\section{Displayed Text}

Text is displayed by indenting it from the left
margin.  Quotations are commonly displayed.  There
are short quotations
\begin{quote}
   This is a short quotation.  It consists of a
   single paragraph of text.  See how it is formatted.
\end{quote}
and longer ones.
\begin{quotation}
   This is a longer quotation.  It consists of two
   paragraphs of text, neither of which are
   particularly interesting.

   This is the second paragraph of the quotation.  It
   is just as dull as the first paragraph.
\end{quotation}
Another frequently-displayed structure is a list.
The following is an example of an \emph{itemized}
list.
\begin{itemize}
   \item This is the first item of an itemized list.
         Each item in the list is marked with a ``tick''.
         You don't have to worry about what kind of tick
         mark is used.

   \item This is the second item of the list.  It
         contains another list nested inside it.  The inner
         list is an \emph{enumerated} list.
         \begin{enumerate}
            \item This is the first item of an enumerated
                  list that is nested within the itemized list.

            \item This is the second item of the inner list.
                  \LaTeX\ allows you to nest lists deeper than
                  you really should.
         \end{enumerate}
         This is the rest of the second item of the outer
         list.  It is no more interesting than any other
         part of the item.
   \item This is the third item of the list.
\end{itemize}
You can even display poetry.
\begin{verse}
   There is an environment
    for verse \\             % The \\ command separates lines
   Whose features some poets % within a stanza.
   will curse.

                             % One or more blank lines separate stanzas.

   For instead of making\\
   Them do \emph{all} line breaking, \\
   It allows them to put too many words on a line when they'd rather be
   forced to be terse.
\end{verse}

Mathematical formulas may also be displayed.  A
displayed formula
is
one-line long; multiline
formulas require special formatting instructions.
   \[  \ip{\Gamma}{\psi'} = x'' + y^{2} + z_{i}^{n}\]
Don't start a paragraph with a displayed equation,
nor make one a paragraph by itself.

\end{document}               % End of document.

\section{Bench Huge Section 845}
% synthetic bench section 845
% This is a sample LaTeX input file.  (Version of 12 August 2004.)
%
% A '%' character causes TeX to ignore all remaining text on the line,
% and is used for comments like this one.

\documentclass{article}      % Specifies the document class

                             % The preamble begins here.
\title{An Example Document}  % Declares the document's title.
\author{Leslie Lamport}      % Declares the author's name.
\date{January 21, 1994}      % Deleting this command produces today's date.

\newcommand{\ip}[2]{(#1, #2)}
                             % Defines \ip{arg1}{arg2} to mean
                             % (arg1, arg2).

%\newcommand{\ip}[2]{\langle #1 | #2\rangle}
                             % This is an alternative definition of
                             % \ip that is commented out.

\begin{document}             % End of preamble and beginning of text.

\maketitle                   % Produces the title.

This is an example input file.  Comparing it with
the output it generates can show you how to
produce a simple document of your own.

\section{Ordinary Text}      % Produces section heading.  Lower-level
                             % sections are begun with similar
                             % \subsection and \subsubsection commands.

The ends  of words and sentences are marked
  by   spaces. It  doesn't matter how many
spaces    you type; one is as good as 100.  The
end of   a line counts as a space.

One   or more   blank lines denote the  end
of  a paragraph.

Since any number of consecutive spaces are treated
like a single one, the formatting of the input
file makes no difference to
      \LaTeX,                % The \LaTeX command generates the LaTeX logo.
but it makes a difference to you.  When you use
\LaTeX, making your input file as easy to read
as possible will be a great help as you write
your document and when you change it.  This sample
file shows how you can add comments to your own input
file.

Because printing is different from typewriting,
there are a number of things that you have to do
differently when preparing an input file than if
you were just typing the document directly.
Quotation marks like
       ``this''
have to be handled specially, as do quotes within
quotes:
       ``\,`this'            % \, separates the double and single quote.
        is what I just
        wrote, not  `that'\,''.

Dashes come in three sizes: an
       intra-word
dash, a medium dash for number ranges like
       1--2,
and a punctuation
       dash---like
this.

A sentence-ending space should be larger than the
space between words within a sentence.  You
sometimes have to type special commands in
conjunction with punctuation characters to get
this right, as in the following sentence.
       Gnats, gnus, etc.\ all  % `\ ' makes an inter-word space.
       begin with G\@.         % \@ marks end-of-sentence punctuation.
You should check the spaces after periods when
reading your output to make sure you haven't
forgotten any special cases.  Generating an
ellipsis
       \ldots\               % `\ ' is needed after `\ldots' because TeX
                             % ignores spaces after command names like \ldots
                             % made from \ + letters.
                             %
                             % Note how a `%' character causes TeX to ignore
                             % the end of the input line, so these blank lines
                             % do not start a new paragraph.
                             %
with the right spacing around the periods requires
a special command.

\LaTeX\ interprets some common characters as
commands, so you must type special commands to
generate them.  These characters include the
following:
       \$ \& \% \# \{ and \}.

In printing, text is usually emphasized with an
       \emph{italic}
type style.

\begin{em}
   A long segment of text can also be emphasized
   in this way.  Text within such a segment can be
   given \emph{additional} emphasis.
\end{em}

It is sometimes necessary to prevent \LaTeX\ from
breaking a line where it might otherwise do so.
This may be at a space, as between the ``Mr.''\ and
``Jones'' in
       ``Mr.~Jones'',        % ~ produces an unbreakable interword space.
or within a word---especially when the word is a
symbol like
       \mbox{\emph{itemnum}}
that makes little sense when hyphenated across
lines.

Footnotes\footnote{This is an example of a footnote.}
pose no problem.

\LaTeX\ is good at typesetting mathematical formulas
like
       \( x-3y + z = 7 \)
or
       \( a_{1} > x^{2n} + y^{2n} > x' \)
or
       \( \ip{A}{B} = \sum_{i} a_{i} b_{i} \).
The spaces you type in a formula are
ignored.  Remember that a letter like
       $x$                   % $ ... $  and  \( ... \)  are equivalent
is a formula when it denotes a mathematical
symbol, and it should be typed as one.

\section{Displayed Text}

Text is displayed by indenting it from the left
margin.  Quotations are commonly displayed.  There
are short quotations
\begin{quote}
   This is a short quotation.  It consists of a
   single paragraph of text.  See how it is formatted.
\end{quote}
and longer ones.
\begin{quotation}
   This is a longer quotation.  It consists of two
   paragraphs of text, neither of which are
   particularly interesting.

   This is the second paragraph of the quotation.  It
   is just as dull as the first paragraph.
\end{quotation}
Another frequently-displayed structure is a list.
The following is an example of an \emph{itemized}
list.
\begin{itemize}
   \item This is the first item of an itemized list.
         Each item in the list is marked with a ``tick''.
         You don't have to worry about what kind of tick
         mark is used.

   \item This is the second item of the list.  It
         contains another list nested inside it.  The inner
         list is an \emph{enumerated} list.
         \begin{enumerate}
            \item This is the first item of an enumerated
                  list that is nested within the itemized list.

            \item This is the second item of the inner list.
                  \LaTeX\ allows you to nest lists deeper than
                  you really should.
         \end{enumerate}
         This is the rest of the second item of the outer
         list.  It is no more interesting than any other
         part of the item.
   \item This is the third item of the list.
\end{itemize}
You can even display poetry.
\begin{verse}
   There is an environment
    for verse \\             % The \\ command separates lines
   Whose features some poets % within a stanza.
   will curse.

                             % One or more blank lines separate stanzas.

   For instead of making\\
   Them do \emph{all} line breaking, \\
   It allows them to put too many words on a line when they'd rather be
   forced to be terse.
\end{verse}

Mathematical formulas may also be displayed.  A
displayed formula
is
one-line long; multiline
formulas require special formatting instructions.
   \[  \ip{\Gamma}{\psi'} = x'' + y^{2} + z_{i}^{n}\]
Don't start a paragraph with a displayed equation,
nor make one a paragraph by itself.

\end{document}               % End of document.

\section{Bench Huge Section 846}
% synthetic bench section 846
% This is a sample LaTeX input file.  (Version of 12 August 2004.)
%
% A '%' character causes TeX to ignore all remaining text on the line,
% and is used for comments like this one.

\documentclass{article}      % Specifies the document class

                             % The preamble begins here.
\title{An Example Document}  % Declares the document's title.
\author{Leslie Lamport}      % Declares the author's name.
\date{January 21, 1994}      % Deleting this command produces today's date.

\newcommand{\ip}[2]{(#1, #2)}
                             % Defines \ip{arg1}{arg2} to mean
                             % (arg1, arg2).

%\newcommand{\ip}[2]{\langle #1 | #2\rangle}
                             % This is an alternative definition of
                             % \ip that is commented out.

\begin{document}             % End of preamble and beginning of text.

\maketitle                   % Produces the title.

This is an example input file.  Comparing it with
the output it generates can show you how to
produce a simple document of your own.

\section{Ordinary Text}      % Produces section heading.  Lower-level
                             % sections are begun with similar
                             % \subsection and \subsubsection commands.

The ends  of words and sentences are marked
  by   spaces. It  doesn't matter how many
spaces    you type; one is as good as 100.  The
end of   a line counts as a space.

One   or more   blank lines denote the  end
of  a paragraph.

Since any number of consecutive spaces are treated
like a single one, the formatting of the input
file makes no difference to
      \LaTeX,                % The \LaTeX command generates the LaTeX logo.
but it makes a difference to you.  When you use
\LaTeX, making your input file as easy to read
as possible will be a great help as you write
your document and when you change it.  This sample
file shows how you can add comments to your own input
file.

Because printing is different from typewriting,
there are a number of things that you have to do
differently when preparing an input file than if
you were just typing the document directly.
Quotation marks like
       ``this''
have to be handled specially, as do quotes within
quotes:
       ``\,`this'            % \, separates the double and single quote.
        is what I just
        wrote, not  `that'\,''.

Dashes come in three sizes: an
       intra-word
dash, a medium dash for number ranges like
       1--2,
and a punctuation
       dash---like
this.

A sentence-ending space should be larger than the
space between words within a sentence.  You
sometimes have to type special commands in
conjunction with punctuation characters to get
this right, as in the following sentence.
       Gnats, gnus, etc.\ all  % `\ ' makes an inter-word space.
       begin with G\@.         % \@ marks end-of-sentence punctuation.
You should check the spaces after periods when
reading your output to make sure you haven't
forgotten any special cases.  Generating an
ellipsis
       \ldots\               % `\ ' is needed after `\ldots' because TeX
                             % ignores spaces after command names like \ldots
                             % made from \ + letters.
                             %
                             % Note how a `%' character causes TeX to ignore
                             % the end of the input line, so these blank lines
                             % do not start a new paragraph.
                             %
with the right spacing around the periods requires
a special command.

\LaTeX\ interprets some common characters as
commands, so you must type special commands to
generate them.  These characters include the
following:
       \$ \& \% \# \{ and \}.

In printing, text is usually emphasized with an
       \emph{italic}
type style.

\begin{em}
   A long segment of text can also be emphasized
   in this way.  Text within such a segment can be
   given \emph{additional} emphasis.
\end{em}

It is sometimes necessary to prevent \LaTeX\ from
breaking a line where it might otherwise do so.
This may be at a space, as between the ``Mr.''\ and
``Jones'' in
       ``Mr.~Jones'',        % ~ produces an unbreakable interword space.
or within a word---especially when the word is a
symbol like
       \mbox{\emph{itemnum}}
that makes little sense when hyphenated across
lines.

Footnotes\footnote{This is an example of a footnote.}
pose no problem.

\LaTeX\ is good at typesetting mathematical formulas
like
       \( x-3y + z = 7 \)
or
       \( a_{1} > x^{2n} + y^{2n} > x' \)
or
       \( \ip{A}{B} = \sum_{i} a_{i} b_{i} \).
The spaces you type in a formula are
ignored.  Remember that a letter like
       $x$                   % $ ... $  and  \( ... \)  are equivalent
is a formula when it denotes a mathematical
symbol, and it should be typed as one.

\section{Displayed Text}

Text is displayed by indenting it from the left
margin.  Quotations are commonly displayed.  There
are short quotations
\begin{quote}
   This is a short quotation.  It consists of a
   single paragraph of text.  See how it is formatted.
\end{quote}
and longer ones.
\begin{quotation}
   This is a longer quotation.  It consists of two
   paragraphs of text, neither of which are
   particularly interesting.

   This is the second paragraph of the quotation.  It
   is just as dull as the first paragraph.
\end{quotation}
Another frequently-displayed structure is a list.
The following is an example of an \emph{itemized}
list.
\begin{itemize}
   \item This is the first item of an itemized list.
         Each item in the list is marked with a ``tick''.
         You don't have to worry about what kind of tick
         mark is used.

   \item This is the second item of the list.  It
         contains another list nested inside it.  The inner
         list is an \emph{enumerated} list.
         \begin{enumerate}
            \item This is the first item of an enumerated
                  list that is nested within the itemized list.

            \item This is the second item of the inner list.
                  \LaTeX\ allows you to nest lists deeper than
                  you really should.
         \end{enumerate}
         This is the rest of the second item of the outer
         list.  It is no more interesting than any other
         part of the item.
   \item This is the third item of the list.
\end{itemize}
You can even display poetry.
\begin{verse}
   There is an environment
    for verse \\             % The \\ command separates lines
   Whose features some poets % within a stanza.
   will curse.

                             % One or more blank lines separate stanzas.

   For instead of making\\
   Them do \emph{all} line breaking, \\
   It allows them to put too many words on a line when they'd rather be
   forced to be terse.
\end{verse}

Mathematical formulas may also be displayed.  A
displayed formula
is
one-line long; multiline
formulas require special formatting instructions.
   \[  \ip{\Gamma}{\psi'} = x'' + y^{2} + z_{i}^{n}\]
Don't start a paragraph with a displayed equation,
nor make one a paragraph by itself.

\end{document}               % End of document.

\section{Bench Huge Section 847}
% synthetic bench section 847
% This is a sample LaTeX input file.  (Version of 12 August 2004.)
%
% A '%' character causes TeX to ignore all remaining text on the line,
% and is used for comments like this one.

\documentclass{article}      % Specifies the document class

                             % The preamble begins here.
\title{An Example Document}  % Declares the document's title.
\author{Leslie Lamport}      % Declares the author's name.
\date{January 21, 1994}      % Deleting this command produces today's date.

\newcommand{\ip}[2]{(#1, #2)}
                             % Defines \ip{arg1}{arg2} to mean
                             % (arg1, arg2).

%\newcommand{\ip}[2]{\langle #1 | #2\rangle}
                             % This is an alternative definition of
                             % \ip that is commented out.

\begin{document}             % End of preamble and beginning of text.

\maketitle                   % Produces the title.

This is an example input file.  Comparing it with
the output it generates can show you how to
produce a simple document of your own.

\section{Ordinary Text}      % Produces section heading.  Lower-level
                             % sections are begun with similar
                             % \subsection and \subsubsection commands.

The ends  of words and sentences are marked
  by   spaces. It  doesn't matter how many
spaces    you type; one is as good as 100.  The
end of   a line counts as a space.

One   or more   blank lines denote the  end
of  a paragraph.

Since any number of consecutive spaces are treated
like a single one, the formatting of the input
file makes no difference to
      \LaTeX,                % The \LaTeX command generates the LaTeX logo.
but it makes a difference to you.  When you use
\LaTeX, making your input file as easy to read
as possible will be a great help as you write
your document and when you change it.  This sample
file shows how you can add comments to your own input
file.

Because printing is different from typewriting,
there are a number of things that you have to do
differently when preparing an input file than if
you were just typing the document directly.
Quotation marks like
       ``this''
have to be handled specially, as do quotes within
quotes:
       ``\,`this'            % \, separates the double and single quote.
        is what I just
        wrote, not  `that'\,''.

Dashes come in three sizes: an
       intra-word
dash, a medium dash for number ranges like
       1--2,
and a punctuation
       dash---like
this.

A sentence-ending space should be larger than the
space between words within a sentence.  You
sometimes have to type special commands in
conjunction with punctuation characters to get
this right, as in the following sentence.
       Gnats, gnus, etc.\ all  % `\ ' makes an inter-word space.
       begin with G\@.         % \@ marks end-of-sentence punctuation.
You should check the spaces after periods when
reading your output to make sure you haven't
forgotten any special cases.  Generating an
ellipsis
       \ldots\               % `\ ' is needed after `\ldots' because TeX
                             % ignores spaces after command names like \ldots
                             % made from \ + letters.
                             %
                             % Note how a `%' character causes TeX to ignore
                             % the end of the input line, so these blank lines
                             % do not start a new paragraph.
                             %
with the right spacing around the periods requires
a special command.

\LaTeX\ interprets some common characters as
commands, so you must type special commands to
generate them.  These characters include the
following:
       \$ \& \% \# \{ and \}.

In printing, text is usually emphasized with an
       \emph{italic}
type style.

\begin{em}
   A long segment of text can also be emphasized
   in this way.  Text within such a segment can be
   given \emph{additional} emphasis.
\end{em}

It is sometimes necessary to prevent \LaTeX\ from
breaking a line where it might otherwise do so.
This may be at a space, as between the ``Mr.''\ and
``Jones'' in
       ``Mr.~Jones'',        % ~ produces an unbreakable interword space.
or within a word---especially when the word is a
symbol like
       \mbox{\emph{itemnum}}
that makes little sense when hyphenated across
lines.

Footnotes\footnote{This is an example of a footnote.}
pose no problem.

\LaTeX\ is good at typesetting mathematical formulas
like
       \( x-3y + z = 7 \)
or
       \( a_{1} > x^{2n} + y^{2n} > x' \)
or
       \( \ip{A}{B} = \sum_{i} a_{i} b_{i} \).
The spaces you type in a formula are
ignored.  Remember that a letter like
       $x$                   % $ ... $  and  \( ... \)  are equivalent
is a formula when it denotes a mathematical
symbol, and it should be typed as one.

\section{Displayed Text}

Text is displayed by indenting it from the left
margin.  Quotations are commonly displayed.  There
are short quotations
\begin{quote}
   This is a short quotation.  It consists of a
   single paragraph of text.  See how it is formatted.
\end{quote}
and longer ones.
\begin{quotation}
   This is a longer quotation.  It consists of two
   paragraphs of text, neither of which are
   particularly interesting.

   This is the second paragraph of the quotation.  It
   is just as dull as the first paragraph.
\end{quotation}
Another frequently-displayed structure is a list.
The following is an example of an \emph{itemized}
list.
\begin{itemize}
   \item This is the first item of an itemized list.
         Each item in the list is marked with a ``tick''.
         You don't have to worry about what kind of tick
         mark is used.

   \item This is the second item of the list.  It
         contains another list nested inside it.  The inner
         list is an \emph{enumerated} list.
         \begin{enumerate}
            \item This is the first item of an enumerated
                  list that is nested within the itemized list.

            \item This is the second item of the inner list.
                  \LaTeX\ allows you to nest lists deeper than
                  you really should.
         \end{enumerate}
         This is the rest of the second item of the outer
         list.  It is no more interesting than any other
         part of the item.
   \item This is the third item of the list.
\end{itemize}
You can even display poetry.
\begin{verse}
   There is an environment
    for verse \\             % The \\ command separates lines
   Whose features some poets % within a stanza.
   will curse.

                             % One or more blank lines separate stanzas.

   For instead of making\\
   Them do \emph{all} line breaking, \\
   It allows them to put too many words on a line when they'd rather be
   forced to be terse.
\end{verse}

Mathematical formulas may also be displayed.  A
displayed formula
is
one-line long; multiline
formulas require special formatting instructions.
   \[  \ip{\Gamma}{\psi'} = x'' + y^{2} + z_{i}^{n}\]
Don't start a paragraph with a displayed equation,
nor make one a paragraph by itself.

\end{document}               % End of document.

\section{Bench Huge Section 848}
% synthetic bench section 848
% This is a sample LaTeX input file.  (Version of 12 August 2004.)
%
% A '%' character causes TeX to ignore all remaining text on the line,
% and is used for comments like this one.

\documentclass{article}      % Specifies the document class

                             % The preamble begins here.
\title{An Example Document}  % Declares the document's title.
\author{Leslie Lamport}      % Declares the author's name.
\date{January 21, 1994}      % Deleting this command produces today's date.

\newcommand{\ip}[2]{(#1, #2)}
                             % Defines \ip{arg1}{arg2} to mean
                             % (arg1, arg2).

%\newcommand{\ip}[2]{\langle #1 | #2\rangle}
                             % This is an alternative definition of
                             % \ip that is commented out.

\begin{document}             % End of preamble and beginning of text.

\maketitle                   % Produces the title.

This is an example input file.  Comparing it with
the output it generates can show you how to
produce a simple document of your own.

\section{Ordinary Text}      % Produces section heading.  Lower-level
                             % sections are begun with similar
                             % \subsection and \subsubsection commands.

The ends  of words and sentences are marked
  by   spaces. It  doesn't matter how many
spaces    you type; one is as good as 100.  The
end of   a line counts as a space.

One   or more   blank lines denote the  end
of  a paragraph.

Since any number of consecutive spaces are treated
like a single one, the formatting of the input
file makes no difference to
      \LaTeX,                % The \LaTeX command generates the LaTeX logo.
but it makes a difference to you.  When you use
\LaTeX, making your input file as easy to read
as possible will be a great help as you write
your document and when you change it.  This sample
file shows how you can add comments to your own input
file.

Because printing is different from typewriting,
there are a number of things that you have to do
differently when preparing an input file than if
you were just typing the document directly.
Quotation marks like
       ``this''
have to be handled specially, as do quotes within
quotes:
       ``\,`this'            % \, separates the double and single quote.
        is what I just
        wrote, not  `that'\,''.

Dashes come in three sizes: an
       intra-word
dash, a medium dash for number ranges like
       1--2,
and a punctuation
       dash---like
this.

A sentence-ending space should be larger than the
space between words within a sentence.  You
sometimes have to type special commands in
conjunction with punctuation characters to get
this right, as in the following sentence.
       Gnats, gnus, etc.\ all  % `\ ' makes an inter-word space.
       begin with G\@.         % \@ marks end-of-sentence punctuation.
You should check the spaces after periods when
reading your output to make sure you haven't
forgotten any special cases.  Generating an
ellipsis
       \ldots\               % `\ ' is needed after `\ldots' because TeX
                             % ignores spaces after command names like \ldots
                             % made from \ + letters.
                             %
                             % Note how a `%' character causes TeX to ignore
                             % the end of the input line, so these blank lines
                             % do not start a new paragraph.
                             %
with the right spacing around the periods requires
a special command.

\LaTeX\ interprets some common characters as
commands, so you must type special commands to
generate them.  These characters include the
following:
       \$ \& \% \# \{ and \}.

In printing, text is usually emphasized with an
       \emph{italic}
type style.

\begin{em}
   A long segment of text can also be emphasized
   in this way.  Text within such a segment can be
   given \emph{additional} emphasis.
\end{em}

It is sometimes necessary to prevent \LaTeX\ from
breaking a line where it might otherwise do so.
This may be at a space, as between the ``Mr.''\ and
``Jones'' in
       ``Mr.~Jones'',        % ~ produces an unbreakable interword space.
or within a word---especially when the word is a
symbol like
       \mbox{\emph{itemnum}}
that makes little sense when hyphenated across
lines.

Footnotes\footnote{This is an example of a footnote.}
pose no problem.

\LaTeX\ is good at typesetting mathematical formulas
like
       \( x-3y + z = 7 \)
or
       \( a_{1} > x^{2n} + y^{2n} > x' \)
or
       \( \ip{A}{B} = \sum_{i} a_{i} b_{i} \).
The spaces you type in a formula are
ignored.  Remember that a letter like
       $x$                   % $ ... $  and  \( ... \)  are equivalent
is a formula when it denotes a mathematical
symbol, and it should be typed as one.

\section{Displayed Text}

Text is displayed by indenting it from the left
margin.  Quotations are commonly displayed.  There
are short quotations
\begin{quote}
   This is a short quotation.  It consists of a
   single paragraph of text.  See how it is formatted.
\end{quote}
and longer ones.
\begin{quotation}
   This is a longer quotation.  It consists of two
   paragraphs of text, neither of which are
   particularly interesting.

   This is the second paragraph of the quotation.  It
   is just as dull as the first paragraph.
\end{quotation}
Another frequently-displayed structure is a list.
The following is an example of an \emph{itemized}
list.
\begin{itemize}
   \item This is the first item of an itemized list.
         Each item in the list is marked with a ``tick''.
         You don't have to worry about what kind of tick
         mark is used.

   \item This is the second item of the list.  It
         contains another list nested inside it.  The inner
         list is an \emph{enumerated} list.
         \begin{enumerate}
            \item This is the first item of an enumerated
                  list that is nested within the itemized list.

            \item This is the second item of the inner list.
                  \LaTeX\ allows you to nest lists deeper than
                  you really should.
         \end{enumerate}
         This is the rest of the second item of the outer
         list.  It is no more interesting than any other
         part of the item.
   \item This is the third item of the list.
\end{itemize}
You can even display poetry.
\begin{verse}
   There is an environment
    for verse \\             % The \\ command separates lines
   Whose features some poets % within a stanza.
   will curse.

                             % One or more blank lines separate stanzas.

   For instead of making\\
   Them do \emph{all} line breaking, \\
   It allows them to put too many words on a line when they'd rather be
   forced to be terse.
\end{verse}

Mathematical formulas may also be displayed.  A
displayed formula
is
one-line long; multiline
formulas require special formatting instructions.
   \[  \ip{\Gamma}{\psi'} = x'' + y^{2} + z_{i}^{n}\]
Don't start a paragraph with a displayed equation,
nor make one a paragraph by itself.

\end{document}               % End of document.

\section{Bench Huge Section 849}
% synthetic bench section 849
% This is a sample LaTeX input file.  (Version of 12 August 2004.)
%
% A '%' character causes TeX to ignore all remaining text on the line,
% and is used for comments like this one.

\documentclass{article}      % Specifies the document class

                             % The preamble begins here.
\title{An Example Document}  % Declares the document's title.
\author{Leslie Lamport}      % Declares the author's name.
\date{January 21, 1994}      % Deleting this command produces today's date.

\newcommand{\ip}[2]{(#1, #2)}
                             % Defines \ip{arg1}{arg2} to mean
                             % (arg1, arg2).

%\newcommand{\ip}[2]{\langle #1 | #2\rangle}
                             % This is an alternative definition of
                             % \ip that is commented out.

\begin{document}             % End of preamble and beginning of text.

\maketitle                   % Produces the title.

This is an example input file.  Comparing it with
the output it generates can show you how to
produce a simple document of your own.

\section{Ordinary Text}      % Produces section heading.  Lower-level
                             % sections are begun with similar
                             % \subsection and \subsubsection commands.

The ends  of words and sentences are marked
  by   spaces. It  doesn't matter how many
spaces    you type; one is as good as 100.  The
end of   a line counts as a space.

One   or more   blank lines denote the  end
of  a paragraph.

Since any number of consecutive spaces are treated
like a single one, the formatting of the input
file makes no difference to
      \LaTeX,                % The \LaTeX command generates the LaTeX logo.
but it makes a difference to you.  When you use
\LaTeX, making your input file as easy to read
as possible will be a great help as you write
your document and when you change it.  This sample
file shows how you can add comments to your own input
file.

Because printing is different from typewriting,
there are a number of things that you have to do
differently when preparing an input file than if
you were just typing the document directly.
Quotation marks like
       ``this''
have to be handled specially, as do quotes within
quotes:
       ``\,`this'            % \, separates the double and single quote.
        is what I just
        wrote, not  `that'\,''.

Dashes come in three sizes: an
       intra-word
dash, a medium dash for number ranges like
       1--2,
and a punctuation
       dash---like
this.

A sentence-ending space should be larger than the
space between words within a sentence.  You
sometimes have to type special commands in
conjunction with punctuation characters to get
this right, as in the following sentence.
       Gnats, gnus, etc.\ all  % `\ ' makes an inter-word space.
       begin with G\@.         % \@ marks end-of-sentence punctuation.
You should check the spaces after periods when
reading your output to make sure you haven't
forgotten any special cases.  Generating an
ellipsis
       \ldots\               % `\ ' is needed after `\ldots' because TeX
                             % ignores spaces after command names like \ldots
                             % made from \ + letters.
                             %
                             % Note how a `%' character causes TeX to ignore
                             % the end of the input line, so these blank lines
                             % do not start a new paragraph.
                             %
with the right spacing around the periods requires
a special command.

\LaTeX\ interprets some common characters as
commands, so you must type special commands to
generate them.  These characters include the
following:
       \$ \& \% \# \{ and \}.

In printing, text is usually emphasized with an
       \emph{italic}
type style.

\begin{em}
   A long segment of text can also be emphasized
   in this way.  Text within such a segment can be
   given \emph{additional} emphasis.
\end{em}

It is sometimes necessary to prevent \LaTeX\ from
breaking a line where it might otherwise do so.
This may be at a space, as between the ``Mr.''\ and
``Jones'' in
       ``Mr.~Jones'',        % ~ produces an unbreakable interword space.
or within a word---especially when the word is a
symbol like
       \mbox{\emph{itemnum}}
that makes little sense when hyphenated across
lines.

Footnotes\footnote{This is an example of a footnote.}
pose no problem.

\LaTeX\ is good at typesetting mathematical formulas
like
       \( x-3y + z = 7 \)
or
       \( a_{1} > x^{2n} + y^{2n} > x' \)
or
       \( \ip{A}{B} = \sum_{i} a_{i} b_{i} \).
The spaces you type in a formula are
ignored.  Remember that a letter like
       $x$                   % $ ... $  and  \( ... \)  are equivalent
is a formula when it denotes a mathematical
symbol, and it should be typed as one.

\section{Displayed Text}

Text is displayed by indenting it from the left
margin.  Quotations are commonly displayed.  There
are short quotations
\begin{quote}
   This is a short quotation.  It consists of a
   single paragraph of text.  See how it is formatted.
\end{quote}
and longer ones.
\begin{quotation}
   This is a longer quotation.  It consists of two
   paragraphs of text, neither of which are
   particularly interesting.

   This is the second paragraph of the quotation.  It
   is just as dull as the first paragraph.
\end{quotation}
Another frequently-displayed structure is a list.
The following is an example of an \emph{itemized}
list.
\begin{itemize}
   \item This is the first item of an itemized list.
         Each item in the list is marked with a ``tick''.
         You don't have to worry about what kind of tick
         mark is used.

   \item This is the second item of the list.  It
         contains another list nested inside it.  The inner
         list is an \emph{enumerated} list.
         \begin{enumerate}
            \item This is the first item of an enumerated
                  list that is nested within the itemized list.

            \item This is the second item of the inner list.
                  \LaTeX\ allows you to nest lists deeper than
                  you really should.
         \end{enumerate}
         This is the rest of the second item of the outer
         list.  It is no more interesting than any other
         part of the item.
   \item This is the third item of the list.
\end{itemize}
You can even display poetry.
\begin{verse}
   There is an environment
    for verse \\             % The \\ command separates lines
   Whose features some poets % within a stanza.
   will curse.

                             % One or more blank lines separate stanzas.

   For instead of making\\
   Them do \emph{all} line breaking, \\
   It allows them to put too many words on a line when they'd rather be
   forced to be terse.
\end{verse}

Mathematical formulas may also be displayed.  A
displayed formula
is
one-line long; multiline
formulas require special formatting instructions.
   \[  \ip{\Gamma}{\psi'} = x'' + y^{2} + z_{i}^{n}\]
Don't start a paragraph with a displayed equation,
nor make one a paragraph by itself.

\end{document}               % End of document.

\section{Bench Huge Section 850}
% synthetic bench section 850
% This is a sample LaTeX input file.  (Version of 12 August 2004.)
%
% A '%' character causes TeX to ignore all remaining text on the line,
% and is used for comments like this one.

\documentclass{article}      % Specifies the document class

                             % The preamble begins here.
\title{An Example Document}  % Declares the document's title.
\author{Leslie Lamport}      % Declares the author's name.
\date{January 21, 1994}      % Deleting this command produces today's date.

\newcommand{\ip}[2]{(#1, #2)}
                             % Defines \ip{arg1}{arg2} to mean
                             % (arg1, arg2).

%\newcommand{\ip}[2]{\langle #1 | #2\rangle}
                             % This is an alternative definition of
                             % \ip that is commented out.

\begin{document}             % End of preamble and beginning of text.

\maketitle                   % Produces the title.

This is an example input file.  Comparing it with
the output it generates can show you how to
produce a simple document of your own.

\section{Ordinary Text}      % Produces section heading.  Lower-level
                             % sections are begun with similar
                             % \subsection and \subsubsection commands.

The ends  of words and sentences are marked
  by   spaces. It  doesn't matter how many
spaces    you type; one is as good as 100.  The
end of   a line counts as a space.

One   or more   blank lines denote the  end
of  a paragraph.

Since any number of consecutive spaces are treated
like a single one, the formatting of the input
file makes no difference to
      \LaTeX,                % The \LaTeX command generates the LaTeX logo.
but it makes a difference to you.  When you use
\LaTeX, making your input file as easy to read
as possible will be a great help as you write
your document and when you change it.  This sample
file shows how you can add comments to your own input
file.

Because printing is different from typewriting,
there are a number of things that you have to do
differently when preparing an input file than if
you were just typing the document directly.
Quotation marks like
       ``this''
have to be handled specially, as do quotes within
quotes:
       ``\,`this'            % \, separates the double and single quote.
        is what I just
        wrote, not  `that'\,''.

Dashes come in three sizes: an
       intra-word
dash, a medium dash for number ranges like
       1--2,
and a punctuation
       dash---like
this.

A sentence-ending space should be larger than the
space between words within a sentence.  You
sometimes have to type special commands in
conjunction with punctuation characters to get
this right, as in the following sentence.
       Gnats, gnus, etc.\ all  % `\ ' makes an inter-word space.
       begin with G\@.         % \@ marks end-of-sentence punctuation.
You should check the spaces after periods when
reading your output to make sure you haven't
forgotten any special cases.  Generating an
ellipsis
       \ldots\               % `\ ' is needed after `\ldots' because TeX
                             % ignores spaces after command names like \ldots
                             % made from \ + letters.
                             %
                             % Note how a `%' character causes TeX to ignore
                             % the end of the input line, so these blank lines
                             % do not start a new paragraph.
                             %
with the right spacing around the periods requires
a special command.

\LaTeX\ interprets some common characters as
commands, so you must type special commands to
generate them.  These characters include the
following:
       \$ \& \% \# \{ and \}.

In printing, text is usually emphasized with an
       \emph{italic}
type style.

\begin{em}
   A long segment of text can also be emphasized
   in this way.  Text within such a segment can be
   given \emph{additional} emphasis.
\end{em}

It is sometimes necessary to prevent \LaTeX\ from
breaking a line where it might otherwise do so.
This may be at a space, as between the ``Mr.''\ and
``Jones'' in
       ``Mr.~Jones'',        % ~ produces an unbreakable interword space.
or within a word---especially when the word is a
symbol like
       \mbox{\emph{itemnum}}
that makes little sense when hyphenated across
lines.

Footnotes\footnote{This is an example of a footnote.}
pose no problem.

\LaTeX\ is good at typesetting mathematical formulas
like
       \( x-3y + z = 7 \)
or
       \( a_{1} > x^{2n} + y^{2n} > x' \)
or
       \( \ip{A}{B} = \sum_{i} a_{i} b_{i} \).
The spaces you type in a formula are
ignored.  Remember that a letter like
       $x$                   % $ ... $  and  \( ... \)  are equivalent
is a formula when it denotes a mathematical
symbol, and it should be typed as one.

\section{Displayed Text}

Text is displayed by indenting it from the left
margin.  Quotations are commonly displayed.  There
are short quotations
\begin{quote}
   This is a short quotation.  It consists of a
   single paragraph of text.  See how it is formatted.
\end{quote}
and longer ones.
\begin{quotation}
   This is a longer quotation.  It consists of two
   paragraphs of text, neither of which are
   particularly interesting.

   This is the second paragraph of the quotation.  It
   is just as dull as the first paragraph.
\end{quotation}
Another frequently-displayed structure is a list.
The following is an example of an \emph{itemized}
list.
\begin{itemize}
   \item This is the first item of an itemized list.
         Each item in the list is marked with a ``tick''.
         You don't have to worry about what kind of tick
         mark is used.

   \item This is the second item of the list.  It
         contains another list nested inside it.  The inner
         list is an \emph{enumerated} list.
         \begin{enumerate}
            \item This is the first item of an enumerated
                  list that is nested within the itemized list.

            \item This is the second item of the inner list.
                  \LaTeX\ allows you to nest lists deeper than
                  you really should.
         \end{enumerate}
         This is the rest of the second item of the outer
         list.  It is no more interesting than any other
         part of the item.
   \item This is the third item of the list.
\end{itemize}
You can even display poetry.
\begin{verse}
   There is an environment
    for verse \\             % The \\ command separates lines
   Whose features some poets % within a stanza.
   will curse.

                             % One or more blank lines separate stanzas.

   For instead of making\\
   Them do \emph{all} line breaking, \\
   It allows them to put too many words on a line when they'd rather be
   forced to be terse.
\end{verse}

Mathematical formulas may also be displayed.  A
displayed formula
is
one-line long; multiline
formulas require special formatting instructions.
   \[  \ip{\Gamma}{\psi'} = x'' + y^{2} + z_{i}^{n}\]
Don't start a paragraph with a displayed equation,
nor make one a paragraph by itself.

\end{document}               % End of document.

\section{Bench Huge Section 851}
% synthetic bench section 851
% This is a sample LaTeX input file.  (Version of 12 August 2004.)
%
% A '%' character causes TeX to ignore all remaining text on the line,
% and is used for comments like this one.

\documentclass{article}      % Specifies the document class

                             % The preamble begins here.
\title{An Example Document}  % Declares the document's title.
\author{Leslie Lamport}      % Declares the author's name.
\date{January 21, 1994}      % Deleting this command produces today's date.

\newcommand{\ip}[2]{(#1, #2)}
                             % Defines \ip{arg1}{arg2} to mean
                             % (arg1, arg2).

%\newcommand{\ip}[2]{\langle #1 | #2\rangle}
                             % This is an alternative definition of
                             % \ip that is commented out.

\begin{document}             % End of preamble and beginning of text.

\maketitle                   % Produces the title.

This is an example input file.  Comparing it with
the output it generates can show you how to
produce a simple document of your own.

\section{Ordinary Text}      % Produces section heading.  Lower-level
                             % sections are begun with similar
                             % \subsection and \subsubsection commands.

The ends  of words and sentences are marked
  by   spaces. It  doesn't matter how many
spaces    you type; one is as good as 100.  The
end of   a line counts as a space.

One   or more   blank lines denote the  end
of  a paragraph.

Since any number of consecutive spaces are treated
like a single one, the formatting of the input
file makes no difference to
      \LaTeX,                % The \LaTeX command generates the LaTeX logo.
but it makes a difference to you.  When you use
\LaTeX, making your input file as easy to read
as possible will be a great help as you write
your document and when you change it.  This sample
file shows how you can add comments to your own input
file.

Because printing is different from typewriting,
there are a number of things that you have to do
differently when preparing an input file than if
you were just typing the document directly.
Quotation marks like
       ``this''
have to be handled specially, as do quotes within
quotes:
       ``\,`this'            % \, separates the double and single quote.
        is what I just
        wrote, not  `that'\,''.

Dashes come in three sizes: an
       intra-word
dash, a medium dash for number ranges like
       1--2,
and a punctuation
       dash---like
this.

A sentence-ending space should be larger than the
space between words within a sentence.  You
sometimes have to type special commands in
conjunction with punctuation characters to get
this right, as in the following sentence.
       Gnats, gnus, etc.\ all  % `\ ' makes an inter-word space.
       begin with G\@.         % \@ marks end-of-sentence punctuation.
You should check the spaces after periods when
reading your output to make sure you haven't
forgotten any special cases.  Generating an
ellipsis
       \ldots\               % `\ ' is needed after `\ldots' because TeX
                             % ignores spaces after command names like \ldots
                             % made from \ + letters.
                             %
                             % Note how a `%' character causes TeX to ignore
                             % the end of the input line, so these blank lines
                             % do not start a new paragraph.
                             %
with the right spacing around the periods requires
a special command.

\LaTeX\ interprets some common characters as
commands, so you must type special commands to
generate them.  These characters include the
following:
       \$ \& \% \# \{ and \}.

In printing, text is usually emphasized with an
       \emph{italic}
type style.

\begin{em}
   A long segment of text can also be emphasized
   in this way.  Text within such a segment can be
   given \emph{additional} emphasis.
\end{em}

It is sometimes necessary to prevent \LaTeX\ from
breaking a line where it might otherwise do so.
This may be at a space, as between the ``Mr.''\ and
``Jones'' in
       ``Mr.~Jones'',        % ~ produces an unbreakable interword space.
or within a word---especially when the word is a
symbol like
       \mbox{\emph{itemnum}}
that makes little sense when hyphenated across
lines.

Footnotes\footnote{This is an example of a footnote.}
pose no problem.

\LaTeX\ is good at typesetting mathematical formulas
like
       \( x-3y + z = 7 \)
or
       \( a_{1} > x^{2n} + y^{2n} > x' \)
or
       \( \ip{A}{B} = \sum_{i} a_{i} b_{i} \).
The spaces you type in a formula are
ignored.  Remember that a letter like
       $x$                   % $ ... $  and  \( ... \)  are equivalent
is a formula when it denotes a mathematical
symbol, and it should be typed as one.

\section{Displayed Text}

Text is displayed by indenting it from the left
margin.  Quotations are commonly displayed.  There
are short quotations
\begin{quote}
   This is a short quotation.  It consists of a
   single paragraph of text.  See how it is formatted.
\end{quote}
and longer ones.
\begin{quotation}
   This is a longer quotation.  It consists of two
   paragraphs of text, neither of which are
   particularly interesting.

   This is the second paragraph of the quotation.  It
   is just as dull as the first paragraph.
\end{quotation}
Another frequently-displayed structure is a list.
The following is an example of an \emph{itemized}
list.
\begin{itemize}
   \item This is the first item of an itemized list.
         Each item in the list is marked with a ``tick''.
         You don't have to worry about what kind of tick
         mark is used.

   \item This is the second item of the list.  It
         contains another list nested inside it.  The inner
         list is an \emph{enumerated} list.
         \begin{enumerate}
            \item This is the first item of an enumerated
                  list that is nested within the itemized list.

            \item This is the second item of the inner list.
                  \LaTeX\ allows you to nest lists deeper than
                  you really should.
         \end{enumerate}
         This is the rest of the second item of the outer
         list.  It is no more interesting than any other
         part of the item.
   \item This is the third item of the list.
\end{itemize}
You can even display poetry.
\begin{verse}
   There is an environment
    for verse \\             % The \\ command separates lines
   Whose features some poets % within a stanza.
   will curse.

                             % One or more blank lines separate stanzas.

   For instead of making\\
   Them do \emph{all} line breaking, \\
   It allows them to put too many words on a line when they'd rather be
   forced to be terse.
\end{verse}

Mathematical formulas may also be displayed.  A
displayed formula
is
one-line long; multiline
formulas require special formatting instructions.
   \[  \ip{\Gamma}{\psi'} = x'' + y^{2} + z_{i}^{n}\]
Don't start a paragraph with a displayed equation,
nor make one a paragraph by itself.

\end{document}               % End of document.

\section{Bench Huge Section 852}
% synthetic bench section 852
% This is a sample LaTeX input file.  (Version of 12 August 2004.)
%
% A '%' character causes TeX to ignore all remaining text on the line,
% and is used for comments like this one.

\documentclass{article}      % Specifies the document class

                             % The preamble begins here.
\title{An Example Document}  % Declares the document's title.
\author{Leslie Lamport}      % Declares the author's name.
\date{January 21, 1994}      % Deleting this command produces today's date.

\newcommand{\ip}[2]{(#1, #2)}
                             % Defines \ip{arg1}{arg2} to mean
                             % (arg1, arg2).

%\newcommand{\ip}[2]{\langle #1 | #2\rangle}
                             % This is an alternative definition of
                             % \ip that is commented out.

\begin{document}             % End of preamble and beginning of text.

\maketitle                   % Produces the title.

This is an example input file.  Comparing it with
the output it generates can show you how to
produce a simple document of your own.

\section{Ordinary Text}      % Produces section heading.  Lower-level
                             % sections are begun with similar
                             % \subsection and \subsubsection commands.

The ends  of words and sentences are marked
  by   spaces. It  doesn't matter how many
spaces    you type; one is as good as 100.  The
end of   a line counts as a space.

One   or more   blank lines denote the  end
of  a paragraph.

Since any number of consecutive spaces are treated
like a single one, the formatting of the input
file makes no difference to
      \LaTeX,                % The \LaTeX command generates the LaTeX logo.
but it makes a difference to you.  When you use
\LaTeX, making your input file as easy to read
as possible will be a great help as you write
your document and when you change it.  This sample
file shows how you can add comments to your own input
file.

Because printing is different from typewriting,
there are a number of things that you have to do
differently when preparing an input file than if
you were just typing the document directly.
Quotation marks like
       ``this''
have to be handled specially, as do quotes within
quotes:
       ``\,`this'            % \, separates the double and single quote.
        is what I just
        wrote, not  `that'\,''.

Dashes come in three sizes: an
       intra-word
dash, a medium dash for number ranges like
       1--2,
and a punctuation
       dash---like
this.

A sentence-ending space should be larger than the
space between words within a sentence.  You
sometimes have to type special commands in
conjunction with punctuation characters to get
this right, as in the following sentence.
       Gnats, gnus, etc.\ all  % `\ ' makes an inter-word space.
       begin with G\@.         % \@ marks end-of-sentence punctuation.
You should check the spaces after periods when
reading your output to make sure you haven't
forgotten any special cases.  Generating an
ellipsis
       \ldots\               % `\ ' is needed after `\ldots' because TeX
                             % ignores spaces after command names like \ldots
                             % made from \ + letters.
                             %
                             % Note how a `%' character causes TeX to ignore
                             % the end of the input line, so these blank lines
                             % do not start a new paragraph.
                             %
with the right spacing around the periods requires
a special command.

\LaTeX\ interprets some common characters as
commands, so you must type special commands to
generate them.  These characters include the
following:
       \$ \& \% \# \{ and \}.

In printing, text is usually emphasized with an
       \emph{italic}
type style.

\begin{em}
   A long segment of text can also be emphasized
   in this way.  Text within such a segment can be
   given \emph{additional} emphasis.
\end{em}

It is sometimes necessary to prevent \LaTeX\ from
breaking a line where it might otherwise do so.
This may be at a space, as between the ``Mr.''\ and
``Jones'' in
       ``Mr.~Jones'',        % ~ produces an unbreakable interword space.
or within a word---especially when the word is a
symbol like
       \mbox{\emph{itemnum}}
that makes little sense when hyphenated across
lines.

Footnotes\footnote{This is an example of a footnote.}
pose no problem.

\LaTeX\ is good at typesetting mathematical formulas
like
       \( x-3y + z = 7 \)
or
       \( a_{1} > x^{2n} + y^{2n} > x' \)
or
       \( \ip{A}{B} = \sum_{i} a_{i} b_{i} \).
The spaces you type in a formula are
ignored.  Remember that a letter like
       $x$                   % $ ... $  and  \( ... \)  are equivalent
is a formula when it denotes a mathematical
symbol, and it should be typed as one.

\section{Displayed Text}

Text is displayed by indenting it from the left
margin.  Quotations are commonly displayed.  There
are short quotations
\begin{quote}
   This is a short quotation.  It consists of a
   single paragraph of text.  See how it is formatted.
\end{quote}
and longer ones.
\begin{quotation}
   This is a longer quotation.  It consists of two
   paragraphs of text, neither of which are
   particularly interesting.

   This is the second paragraph of the quotation.  It
   is just as dull as the first paragraph.
\end{quotation}
Another frequently-displayed structure is a list.
The following is an example of an \emph{itemized}
list.
\begin{itemize}
   \item This is the first item of an itemized list.
         Each item in the list is marked with a ``tick''.
         You don't have to worry about what kind of tick
         mark is used.

   \item This is the second item of the list.  It
         contains another list nested inside it.  The inner
         list is an \emph{enumerated} list.
         \begin{enumerate}
            \item This is the first item of an enumerated
                  list that is nested within the itemized list.

            \item This is the second item of the inner list.
                  \LaTeX\ allows you to nest lists deeper than
                  you really should.
         \end{enumerate}
         This is the rest of the second item of the outer
         list.  It is no more interesting than any other
         part of the item.
   \item This is the third item of the list.
\end{itemize}
You can even display poetry.
\begin{verse}
   There is an environment
    for verse \\             % The \\ command separates lines
   Whose features some poets % within a stanza.
   will curse.

                             % One or more blank lines separate stanzas.

   For instead of making\\
   Them do \emph{all} line breaking, \\
   It allows them to put too many words on a line when they'd rather be
   forced to be terse.
\end{verse}

Mathematical formulas may also be displayed.  A
displayed formula
is
one-line long; multiline
formulas require special formatting instructions.
   \[  \ip{\Gamma}{\psi'} = x'' + y^{2} + z_{i}^{n}\]
Don't start a paragraph with a displayed equation,
nor make one a paragraph by itself.

\end{document}               % End of document.

\section{Bench Huge Section 853}
% synthetic bench section 853
% This is a sample LaTeX input file.  (Version of 12 August 2004.)
%
% A '%' character causes TeX to ignore all remaining text on the line,
% and is used for comments like this one.

\documentclass{article}      % Specifies the document class

                             % The preamble begins here.
\title{An Example Document}  % Declares the document's title.
\author{Leslie Lamport}      % Declares the author's name.
\date{January 21, 1994}      % Deleting this command produces today's date.

\newcommand{\ip}[2]{(#1, #2)}
                             % Defines \ip{arg1}{arg2} to mean
                             % (arg1, arg2).

%\newcommand{\ip}[2]{\langle #1 | #2\rangle}
                             % This is an alternative definition of
                             % \ip that is commented out.

\begin{document}             % End of preamble and beginning of text.

\maketitle                   % Produces the title.

This is an example input file.  Comparing it with
the output it generates can show you how to
produce a simple document of your own.

\section{Ordinary Text}      % Produces section heading.  Lower-level
                             % sections are begun with similar
                             % \subsection and \subsubsection commands.

The ends  of words and sentences are marked
  by   spaces. It  doesn't matter how many
spaces    you type; one is as good as 100.  The
end of   a line counts as a space.

One   or more   blank lines denote the  end
of  a paragraph.

Since any number of consecutive spaces are treated
like a single one, the formatting of the input
file makes no difference to
      \LaTeX,                % The \LaTeX command generates the LaTeX logo.
but it makes a difference to you.  When you use
\LaTeX, making your input file as easy to read
as possible will be a great help as you write
your document and when you change it.  This sample
file shows how you can add comments to your own input
file.

Because printing is different from typewriting,
there are a number of things that you have to do
differently when preparing an input file than if
you were just typing the document directly.
Quotation marks like
       ``this''
have to be handled specially, as do quotes within
quotes:
       ``\,`this'            % \, separates the double and single quote.
        is what I just
        wrote, not  `that'\,''.

Dashes come in three sizes: an
       intra-word
dash, a medium dash for number ranges like
       1--2,
and a punctuation
       dash---like
this.

A sentence-ending space should be larger than the
space between words within a sentence.  You
sometimes have to type special commands in
conjunction with punctuation characters to get
this right, as in the following sentence.
       Gnats, gnus, etc.\ all  % `\ ' makes an inter-word space.
       begin with G\@.         % \@ marks end-of-sentence punctuation.
You should check the spaces after periods when
reading your output to make sure you haven't
forgotten any special cases.  Generating an
ellipsis
       \ldots\               % `\ ' is needed after `\ldots' because TeX
                             % ignores spaces after command names like \ldots
                             % made from \ + letters.
                             %
                             % Note how a `%' character causes TeX to ignore
                             % the end of the input line, so these blank lines
                             % do not start a new paragraph.
                             %
with the right spacing around the periods requires
a special command.

\LaTeX\ interprets some common characters as
commands, so you must type special commands to
generate them.  These characters include the
following:
       \$ \& \% \# \{ and \}.

In printing, text is usually emphasized with an
       \emph{italic}
type style.

\begin{em}
   A long segment of text can also be emphasized
   in this way.  Text within such a segment can be
   given \emph{additional} emphasis.
\end{em}

It is sometimes necessary to prevent \LaTeX\ from
breaking a line where it might otherwise do so.
This may be at a space, as between the ``Mr.''\ and
``Jones'' in
       ``Mr.~Jones'',        % ~ produces an unbreakable interword space.
or within a word---especially when the word is a
symbol like
       \mbox{\emph{itemnum}}
that makes little sense when hyphenated across
lines.

Footnotes\footnote{This is an example of a footnote.}
pose no problem.

\LaTeX\ is good at typesetting mathematical formulas
like
       \( x-3y + z = 7 \)
or
       \( a_{1} > x^{2n} + y^{2n} > x' \)
or
       \( \ip{A}{B} = \sum_{i} a_{i} b_{i} \).
The spaces you type in a formula are
ignored.  Remember that a letter like
       $x$                   % $ ... $  and  \( ... \)  are equivalent
is a formula when it denotes a mathematical
symbol, and it should be typed as one.

\section{Displayed Text}

Text is displayed by indenting it from the left
margin.  Quotations are commonly displayed.  There
are short quotations
\begin{quote}
   This is a short quotation.  It consists of a
   single paragraph of text.  See how it is formatted.
\end{quote}
and longer ones.
\begin{quotation}
   This is a longer quotation.  It consists of two
   paragraphs of text, neither of which are
   particularly interesting.

   This is the second paragraph of the quotation.  It
   is just as dull as the first paragraph.
\end{quotation}
Another frequently-displayed structure is a list.
The following is an example of an \emph{itemized}
list.
\begin{itemize}
   \item This is the first item of an itemized list.
         Each item in the list is marked with a ``tick''.
         You don't have to worry about what kind of tick
         mark is used.

   \item This is the second item of the list.  It
         contains another list nested inside it.  The inner
         list is an \emph{enumerated} list.
         \begin{enumerate}
            \item This is the first item of an enumerated
                  list that is nested within the itemized list.

            \item This is the second item of the inner list.
                  \LaTeX\ allows you to nest lists deeper than
                  you really should.
         \end{enumerate}
         This is the rest of the second item of the outer
         list.  It is no more interesting than any other
         part of the item.
   \item This is the third item of the list.
\end{itemize}
You can even display poetry.
\begin{verse}
   There is an environment
    for verse \\             % The \\ command separates lines
   Whose features some poets % within a stanza.
   will curse.

                             % One or more blank lines separate stanzas.

   For instead of making\\
   Them do \emph{all} line breaking, \\
   It allows them to put too many words on a line when they'd rather be
   forced to be terse.
\end{verse}

Mathematical formulas may also be displayed.  A
displayed formula
is
one-line long; multiline
formulas require special formatting instructions.
   \[  \ip{\Gamma}{\psi'} = x'' + y^{2} + z_{i}^{n}\]
Don't start a paragraph with a displayed equation,
nor make one a paragraph by itself.

\end{document}               % End of document.

\section{Bench Huge Section 854}
% synthetic bench section 854
% This is a sample LaTeX input file.  (Version of 12 August 2004.)
%
% A '%' character causes TeX to ignore all remaining text on the line,
% and is used for comments like this one.

\documentclass{article}      % Specifies the document class

                             % The preamble begins here.
\title{An Example Document}  % Declares the document's title.
\author{Leslie Lamport}      % Declares the author's name.
\date{January 21, 1994}      % Deleting this command produces today's date.

\newcommand{\ip}[2]{(#1, #2)}
                             % Defines \ip{arg1}{arg2} to mean
                             % (arg1, arg2).

%\newcommand{\ip}[2]{\langle #1 | #2\rangle}
                             % This is an alternative definition of
                             % \ip that is commented out.

\begin{document}             % End of preamble and beginning of text.

\maketitle                   % Produces the title.

This is an example input file.  Comparing it with
the output it generates can show you how to
produce a simple document of your own.

\section{Ordinary Text}      % Produces section heading.  Lower-level
                             % sections are begun with similar
                             % \subsection and \subsubsection commands.

The ends  of words and sentences are marked
  by   spaces. It  doesn't matter how many
spaces    you type; one is as good as 100.  The
end of   a line counts as a space.

One   or more   blank lines denote the  end
of  a paragraph.

Since any number of consecutive spaces are treated
like a single one, the formatting of the input
file makes no difference to
      \LaTeX,                % The \LaTeX command generates the LaTeX logo.
but it makes a difference to you.  When you use
\LaTeX, making your input file as easy to read
as possible will be a great help as you write
your document and when you change it.  This sample
file shows how you can add comments to your own input
file.

Because printing is different from typewriting,
there are a number of things that you have to do
differently when preparing an input file than if
you were just typing the document directly.
Quotation marks like
       ``this''
have to be handled specially, as do quotes within
quotes:
       ``\,`this'            % \, separates the double and single quote.
        is what I just
        wrote, not  `that'\,''.

Dashes come in three sizes: an
       intra-word
dash, a medium dash for number ranges like
       1--2,
and a punctuation
       dash---like
this.

A sentence-ending space should be larger than the
space between words within a sentence.  You
sometimes have to type special commands in
conjunction with punctuation characters to get
this right, as in the following sentence.
       Gnats, gnus, etc.\ all  % `\ ' makes an inter-word space.
       begin with G\@.         % \@ marks end-of-sentence punctuation.
You should check the spaces after periods when
reading your output to make sure you haven't
forgotten any special cases.  Generating an
ellipsis
       \ldots\               % `\ ' is needed after `\ldots' because TeX
                             % ignores spaces after command names like \ldots
                             % made from \ + letters.
                             %
                             % Note how a `%' character causes TeX to ignore
                             % the end of the input line, so these blank lines
                             % do not start a new paragraph.
                             %
with the right spacing around the periods requires
a special command.

\LaTeX\ interprets some common characters as
commands, so you must type special commands to
generate them.  These characters include the
following:
       \$ \& \% \# \{ and \}.

In printing, text is usually emphasized with an
       \emph{italic}
type style.

\begin{em}
   A long segment of text can also be emphasized
   in this way.  Text within such a segment can be
   given \emph{additional} emphasis.
\end{em}

It is sometimes necessary to prevent \LaTeX\ from
breaking a line where it might otherwise do so.
This may be at a space, as between the ``Mr.''\ and
``Jones'' in
       ``Mr.~Jones'',        % ~ produces an unbreakable interword space.
or within a word---especially when the word is a
symbol like
       \mbox{\emph{itemnum}}
that makes little sense when hyphenated across
lines.

Footnotes\footnote{This is an example of a footnote.}
pose no problem.

\LaTeX\ is good at typesetting mathematical formulas
like
       \( x-3y + z = 7 \)
or
       \( a_{1} > x^{2n} + y^{2n} > x' \)
or
       \( \ip{A}{B} = \sum_{i} a_{i} b_{i} \).
The spaces you type in a formula are
ignored.  Remember that a letter like
       $x$                   % $ ... $  and  \( ... \)  are equivalent
is a formula when it denotes a mathematical
symbol, and it should be typed as one.

\section{Displayed Text}

Text is displayed by indenting it from the left
margin.  Quotations are commonly displayed.  There
are short quotations
\begin{quote}
   This is a short quotation.  It consists of a
   single paragraph of text.  See how it is formatted.
\end{quote}
and longer ones.
\begin{quotation}
   This is a longer quotation.  It consists of two
   paragraphs of text, neither of which are
   particularly interesting.

   This is the second paragraph of the quotation.  It
   is just as dull as the first paragraph.
\end{quotation}
Another frequently-displayed structure is a list.
The following is an example of an \emph{itemized}
list.
\begin{itemize}
   \item This is the first item of an itemized list.
         Each item in the list is marked with a ``tick''.
         You don't have to worry about what kind of tick
         mark is used.

   \item This is the second item of the list.  It
         contains another list nested inside it.  The inner
         list is an \emph{enumerated} list.
         \begin{enumerate}
            \item This is the first item of an enumerated
                  list that is nested within the itemized list.

            \item This is the second item of the inner list.
                  \LaTeX\ allows you to nest lists deeper than
                  you really should.
         \end{enumerate}
         This is the rest of the second item of the outer
         list.  It is no more interesting than any other
         part of the item.
   \item This is the third item of the list.
\end{itemize}
You can even display poetry.
\begin{verse}
   There is an environment
    for verse \\             % The \\ command separates lines
   Whose features some poets % within a stanza.
   will curse.

                             % One or more blank lines separate stanzas.

   For instead of making\\
   Them do \emph{all} line breaking, \\
   It allows them to put too many words on a line when they'd rather be
   forced to be terse.
\end{verse}

Mathematical formulas may also be displayed.  A
displayed formula
is
one-line long; multiline
formulas require special formatting instructions.
   \[  \ip{\Gamma}{\psi'} = x'' + y^{2} + z_{i}^{n}\]
Don't start a paragraph with a displayed equation,
nor make one a paragraph by itself.

\end{document}               % End of document.

\section{Bench Huge Section 855}
% synthetic bench section 855
% This is a sample LaTeX input file.  (Version of 12 August 2004.)
%
% A '%' character causes TeX to ignore all remaining text on the line,
% and is used for comments like this one.

\documentclass{article}      % Specifies the document class

                             % The preamble begins here.
\title{An Example Document}  % Declares the document's title.
\author{Leslie Lamport}      % Declares the author's name.
\date{January 21, 1994}      % Deleting this command produces today's date.

\newcommand{\ip}[2]{(#1, #2)}
                             % Defines \ip{arg1}{arg2} to mean
                             % (arg1, arg2).

%\newcommand{\ip}[2]{\langle #1 | #2\rangle}
                             % This is an alternative definition of
                             % \ip that is commented out.

\begin{document}             % End of preamble and beginning of text.

\maketitle                   % Produces the title.

This is an example input file.  Comparing it with
the output it generates can show you how to
produce a simple document of your own.

\section{Ordinary Text}      % Produces section heading.  Lower-level
                             % sections are begun with similar
                             % \subsection and \subsubsection commands.

The ends  of words and sentences are marked
  by   spaces. It  doesn't matter how many
spaces    you type; one is as good as 100.  The
end of   a line counts as a space.

One   or more   blank lines denote the  end
of  a paragraph.

Since any number of consecutive spaces are treated
like a single one, the formatting of the input
file makes no difference to
      \LaTeX,                % The \LaTeX command generates the LaTeX logo.
but it makes a difference to you.  When you use
\LaTeX, making your input file as easy to read
as possible will be a great help as you write
your document and when you change it.  This sample
file shows how you can add comments to your own input
file.

Because printing is different from typewriting,
there are a number of things that you have to do
differently when preparing an input file than if
you were just typing the document directly.
Quotation marks like
       ``this''
have to be handled specially, as do quotes within
quotes:
       ``\,`this'            % \, separates the double and single quote.
        is what I just
        wrote, not  `that'\,''.

Dashes come in three sizes: an
       intra-word
dash, a medium dash for number ranges like
       1--2,
and a punctuation
       dash---like
this.

A sentence-ending space should be larger than the
space between words within a sentence.  You
sometimes have to type special commands in
conjunction with punctuation characters to get
this right, as in the following sentence.
       Gnats, gnus, etc.\ all  % `\ ' makes an inter-word space.
       begin with G\@.         % \@ marks end-of-sentence punctuation.
You should check the spaces after periods when
reading your output to make sure you haven't
forgotten any special cases.  Generating an
ellipsis
       \ldots\               % `\ ' is needed after `\ldots' because TeX
                             % ignores spaces after command names like \ldots
                             % made from \ + letters.
                             %
                             % Note how a `%' character causes TeX to ignore
                             % the end of the input line, so these blank lines
                             % do not start a new paragraph.
                             %
with the right spacing around the periods requires
a special command.

\LaTeX\ interprets some common characters as
commands, so you must type special commands to
generate them.  These characters include the
following:
       \$ \& \% \# \{ and \}.

In printing, text is usually emphasized with an
       \emph{italic}
type style.

\begin{em}
   A long segment of text can also be emphasized
   in this way.  Text within such a segment can be
   given \emph{additional} emphasis.
\end{em}

It is sometimes necessary to prevent \LaTeX\ from
breaking a line where it might otherwise do so.
This may be at a space, as between the ``Mr.''\ and
``Jones'' in
       ``Mr.~Jones'',        % ~ produces an unbreakable interword space.
or within a word---especially when the word is a
symbol like
       \mbox{\emph{itemnum}}
that makes little sense when hyphenated across
lines.

Footnotes\footnote{This is an example of a footnote.}
pose no problem.

\LaTeX\ is good at typesetting mathematical formulas
like
       \( x-3y + z = 7 \)
or
       \( a_{1} > x^{2n} + y^{2n} > x' \)
or
       \( \ip{A}{B} = \sum_{i} a_{i} b_{i} \).
The spaces you type in a formula are
ignored.  Remember that a letter like
       $x$                   % $ ... $  and  \( ... \)  are equivalent
is a formula when it denotes a mathematical
symbol, and it should be typed as one.

\section{Displayed Text}

Text is displayed by indenting it from the left
margin.  Quotations are commonly displayed.  There
are short quotations
\begin{quote}
   This is a short quotation.  It consists of a
   single paragraph of text.  See how it is formatted.
\end{quote}
and longer ones.
\begin{quotation}
   This is a longer quotation.  It consists of two
   paragraphs of text, neither of which are
   particularly interesting.

   This is the second paragraph of the quotation.  It
   is just as dull as the first paragraph.
\end{quotation}
Another frequently-displayed structure is a list.
The following is an example of an \emph{itemized}
list.
\begin{itemize}
   \item This is the first item of an itemized list.
         Each item in the list is marked with a ``tick''.
         You don't have to worry about what kind of tick
         mark is used.

   \item This is the second item of the list.  It
         contains another list nested inside it.  The inner
         list is an \emph{enumerated} list.
         \begin{enumerate}
            \item This is the first item of an enumerated
                  list that is nested within the itemized list.

            \item This is the second item of the inner list.
                  \LaTeX\ allows you to nest lists deeper than
                  you really should.
         \end{enumerate}
         This is the rest of the second item of the outer
         list.  It is no more interesting than any other
         part of the item.
   \item This is the third item of the list.
\end{itemize}
You can even display poetry.
\begin{verse}
   There is an environment
    for verse \\             % The \\ command separates lines
   Whose features some poets % within a stanza.
   will curse.

                             % One or more blank lines separate stanzas.

   For instead of making\\
   Them do \emph{all} line breaking, \\
   It allows them to put too many words on a line when they'd rather be
   forced to be terse.
\end{verse}

Mathematical formulas may also be displayed.  A
displayed formula
is
one-line long; multiline
formulas require special formatting instructions.
   \[  \ip{\Gamma}{\psi'} = x'' + y^{2} + z_{i}^{n}\]
Don't start a paragraph with a displayed equation,
nor make one a paragraph by itself.

\end{document}               % End of document.

\section{Bench Huge Section 856}
% synthetic bench section 856
% This is a sample LaTeX input file.  (Version of 12 August 2004.)
%
% A '%' character causes TeX to ignore all remaining text on the line,
% and is used for comments like this one.

\documentclass{article}      % Specifies the document class

                             % The preamble begins here.
\title{An Example Document}  % Declares the document's title.
\author{Leslie Lamport}      % Declares the author's name.
\date{January 21, 1994}      % Deleting this command produces today's date.

\newcommand{\ip}[2]{(#1, #2)}
                             % Defines \ip{arg1}{arg2} to mean
                             % (arg1, arg2).

%\newcommand{\ip}[2]{\langle #1 | #2\rangle}
                             % This is an alternative definition of
                             % \ip that is commented out.

\begin{document}             % End of preamble and beginning of text.

\maketitle                   % Produces the title.

This is an example input file.  Comparing it with
the output it generates can show you how to
produce a simple document of your own.

\section{Ordinary Text}      % Produces section heading.  Lower-level
                             % sections are begun with similar
                             % \subsection and \subsubsection commands.

The ends  of words and sentences are marked
  by   spaces. It  doesn't matter how many
spaces    you type; one is as good as 100.  The
end of   a line counts as a space.

One   or more   blank lines denote the  end
of  a paragraph.

Since any number of consecutive spaces are treated
like a single one, the formatting of the input
file makes no difference to
      \LaTeX,                % The \LaTeX command generates the LaTeX logo.
but it makes a difference to you.  When you use
\LaTeX, making your input file as easy to read
as possible will be a great help as you write
your document and when you change it.  This sample
file shows how you can add comments to your own input
file.

Because printing is different from typewriting,
there are a number of things that you have to do
differently when preparing an input file than if
you were just typing the document directly.
Quotation marks like
       ``this''
have to be handled specially, as do quotes within
quotes:
       ``\,`this'            % \, separates the double and single quote.
        is what I just
        wrote, not  `that'\,''.

Dashes come in three sizes: an
       intra-word
dash, a medium dash for number ranges like
       1--2,
and a punctuation
       dash---like
this.

A sentence-ending space should be larger than the
space between words within a sentence.  You
sometimes have to type special commands in
conjunction with punctuation characters to get
this right, as in the following sentence.
       Gnats, gnus, etc.\ all  % `\ ' makes an inter-word space.
       begin with G\@.         % \@ marks end-of-sentence punctuation.
You should check the spaces after periods when
reading your output to make sure you haven't
forgotten any special cases.  Generating an
ellipsis
       \ldots\               % `\ ' is needed after `\ldots' because TeX
                             % ignores spaces after command names like \ldots
                             % made from \ + letters.
                             %
                             % Note how a `%' character causes TeX to ignore
                             % the end of the input line, so these blank lines
                             % do not start a new paragraph.
                             %
with the right spacing around the periods requires
a special command.

\LaTeX\ interprets some common characters as
commands, so you must type special commands to
generate them.  These characters include the
following:
       \$ \& \% \# \{ and \}.

In printing, text is usually emphasized with an
       \emph{italic}
type style.

\begin{em}
   A long segment of text can also be emphasized
   in this way.  Text within such a segment can be
   given \emph{additional} emphasis.
\end{em}

It is sometimes necessary to prevent \LaTeX\ from
breaking a line where it might otherwise do so.
This may be at a space, as between the ``Mr.''\ and
``Jones'' in
       ``Mr.~Jones'',        % ~ produces an unbreakable interword space.
or within a word---especially when the word is a
symbol like
       \mbox{\emph{itemnum}}
that makes little sense when hyphenated across
lines.

Footnotes\footnote{This is an example of a footnote.}
pose no problem.

\LaTeX\ is good at typesetting mathematical formulas
like
       \( x-3y + z = 7 \)
or
       \( a_{1} > x^{2n} + y^{2n} > x' \)
or
       \( \ip{A}{B} = \sum_{i} a_{i} b_{i} \).
The spaces you type in a formula are
ignored.  Remember that a letter like
       $x$                   % $ ... $  and  \( ... \)  are equivalent
is a formula when it denotes a mathematical
symbol, and it should be typed as one.

\section{Displayed Text}

Text is displayed by indenting it from the left
margin.  Quotations are commonly displayed.  There
are short quotations
\begin{quote}
   This is a short quotation.  It consists of a
   single paragraph of text.  See how it is formatted.
\end{quote}
and longer ones.
\begin{quotation}
   This is a longer quotation.  It consists of two
   paragraphs of text, neither of which are
   particularly interesting.

   This is the second paragraph of the quotation.  It
   is just as dull as the first paragraph.
\end{quotation}
Another frequently-displayed structure is a list.
The following is an example of an \emph{itemized}
list.
\begin{itemize}
   \item This is the first item of an itemized list.
         Each item in the list is marked with a ``tick''.
         You don't have to worry about what kind of tick
         mark is used.

   \item This is the second item of the list.  It
         contains another list nested inside it.  The inner
         list is an \emph{enumerated} list.
         \begin{enumerate}
            \item This is the first item of an enumerated
                  list that is nested within the itemized list.

            \item This is the second item of the inner list.
                  \LaTeX\ allows you to nest lists deeper than
                  you really should.
         \end{enumerate}
         This is the rest of the second item of the outer
         list.  It is no more interesting than any other
         part of the item.
   \item This is the third item of the list.
\end{itemize}
You can even display poetry.
\begin{verse}
   There is an environment
    for verse \\             % The \\ command separates lines
   Whose features some poets % within a stanza.
   will curse.

                             % One or more blank lines separate stanzas.

   For instead of making\\
   Them do \emph{all} line breaking, \\
   It allows them to put too many words on a line when they'd rather be
   forced to be terse.
\end{verse}

Mathematical formulas may also be displayed.  A
displayed formula
is
one-line long; multiline
formulas require special formatting instructions.
   \[  \ip{\Gamma}{\psi'} = x'' + y^{2} + z_{i}^{n}\]
Don't start a paragraph with a displayed equation,
nor make one a paragraph by itself.

\end{document}               % End of document.

\section{Bench Huge Section 857}
% synthetic bench section 857
% This is a sample LaTeX input file.  (Version of 12 August 2004.)
%
% A '%' character causes TeX to ignore all remaining text on the line,
% and is used for comments like this one.

\documentclass{article}      % Specifies the document class

                             % The preamble begins here.
\title{An Example Document}  % Declares the document's title.
\author{Leslie Lamport}      % Declares the author's name.
\date{January 21, 1994}      % Deleting this command produces today's date.

\newcommand{\ip}[2]{(#1, #2)}
                             % Defines \ip{arg1}{arg2} to mean
                             % (arg1, arg2).

%\newcommand{\ip}[2]{\langle #1 | #2\rangle}
                             % This is an alternative definition of
                             % \ip that is commented out.

\begin{document}             % End of preamble and beginning of text.

\maketitle                   % Produces the title.

This is an example input file.  Comparing it with
the output it generates can show you how to
produce a simple document of your own.

\section{Ordinary Text}      % Produces section heading.  Lower-level
                             % sections are begun with similar
                             % \subsection and \subsubsection commands.

The ends  of words and sentences are marked
  by   spaces. It  doesn't matter how many
spaces    you type; one is as good as 100.  The
end of   a line counts as a space.

One   or more   blank lines denote the  end
of  a paragraph.

Since any number of consecutive spaces are treated
like a single one, the formatting of the input
file makes no difference to
      \LaTeX,                % The \LaTeX command generates the LaTeX logo.
but it makes a difference to you.  When you use
\LaTeX, making your input file as easy to read
as possible will be a great help as you write
your document and when you change it.  This sample
file shows how you can add comments to your own input
file.

Because printing is different from typewriting,
there are a number of things that you have to do
differently when preparing an input file than if
you were just typing the document directly.
Quotation marks like
       ``this''
have to be handled specially, as do quotes within
quotes:
       ``\,`this'            % \, separates the double and single quote.
        is what I just
        wrote, not  `that'\,''.

Dashes come in three sizes: an
       intra-word
dash, a medium dash for number ranges like
       1--2,
and a punctuation
       dash---like
this.

A sentence-ending space should be larger than the
space between words within a sentence.  You
sometimes have to type special commands in
conjunction with punctuation characters to get
this right, as in the following sentence.
       Gnats, gnus, etc.\ all  % `\ ' makes an inter-word space.
       begin with G\@.         % \@ marks end-of-sentence punctuation.
You should check the spaces after periods when
reading your output to make sure you haven't
forgotten any special cases.  Generating an
ellipsis
       \ldots\               % `\ ' is needed after `\ldots' because TeX
                             % ignores spaces after command names like \ldots
                             % made from \ + letters.
                             %
                             % Note how a `%' character causes TeX to ignore
                             % the end of the input line, so these blank lines
                             % do not start a new paragraph.
                             %
with the right spacing around the periods requires
a special command.

\LaTeX\ interprets some common characters as
commands, so you must type special commands to
generate them.  These characters include the
following:
       \$ \& \% \# \{ and \}.

In printing, text is usually emphasized with an
       \emph{italic}
type style.

\begin{em}
   A long segment of text can also be emphasized
   in this way.  Text within such a segment can be
   given \emph{additional} emphasis.
\end{em}

It is sometimes necessary to prevent \LaTeX\ from
breaking a line where it might otherwise do so.
This may be at a space, as between the ``Mr.''\ and
``Jones'' in
       ``Mr.~Jones'',        % ~ produces an unbreakable interword space.
or within a word---especially when the word is a
symbol like
       \mbox{\emph{itemnum}}
that makes little sense when hyphenated across
lines.

Footnotes\footnote{This is an example of a footnote.}
pose no problem.

\LaTeX\ is good at typesetting mathematical formulas
like
       \( x-3y + z = 7 \)
or
       \( a_{1} > x^{2n} + y^{2n} > x' \)
or
       \( \ip{A}{B} = \sum_{i} a_{i} b_{i} \).
The spaces you type in a formula are
ignored.  Remember that a letter like
       $x$                   % $ ... $  and  \( ... \)  are equivalent
is a formula when it denotes a mathematical
symbol, and it should be typed as one.

\section{Displayed Text}

Text is displayed by indenting it from the left
margin.  Quotations are commonly displayed.  There
are short quotations
\begin{quote}
   This is a short quotation.  It consists of a
   single paragraph of text.  See how it is formatted.
\end{quote}
and longer ones.
\begin{quotation}
   This is a longer quotation.  It consists of two
   paragraphs of text, neither of which are
   particularly interesting.

   This is the second paragraph of the quotation.  It
   is just as dull as the first paragraph.
\end{quotation}
Another frequently-displayed structure is a list.
The following is an example of an \emph{itemized}
list.
\begin{itemize}
   \item This is the first item of an itemized list.
         Each item in the list is marked with a ``tick''.
         You don't have to worry about what kind of tick
         mark is used.

   \item This is the second item of the list.  It
         contains another list nested inside it.  The inner
         list is an \emph{enumerated} list.
         \begin{enumerate}
            \item This is the first item of an enumerated
                  list that is nested within the itemized list.

            \item This is the second item of the inner list.
                  \LaTeX\ allows you to nest lists deeper than
                  you really should.
         \end{enumerate}
         This is the rest of the second item of the outer
         list.  It is no more interesting than any other
         part of the item.
   \item This is the third item of the list.
\end{itemize}
You can even display poetry.
\begin{verse}
   There is an environment
    for verse \\             % The \\ command separates lines
   Whose features some poets % within a stanza.
   will curse.

                             % One or more blank lines separate stanzas.

   For instead of making\\
   Them do \emph{all} line breaking, \\
   It allows them to put too many words on a line when they'd rather be
   forced to be terse.
\end{verse}

Mathematical formulas may also be displayed.  A
displayed formula
is
one-line long; multiline
formulas require special formatting instructions.
   \[  \ip{\Gamma}{\psi'} = x'' + y^{2} + z_{i}^{n}\]
Don't start a paragraph with a displayed equation,
nor make one a paragraph by itself.

\end{document}               % End of document.

\section{Bench Huge Section 858}
% synthetic bench section 858
% This is a sample LaTeX input file.  (Version of 12 August 2004.)
%
% A '%' character causes TeX to ignore all remaining text on the line,
% and is used for comments like this one.

\documentclass{article}      % Specifies the document class

                             % The preamble begins here.
\title{An Example Document}  % Declares the document's title.
\author{Leslie Lamport}      % Declares the author's name.
\date{January 21, 1994}      % Deleting this command produces today's date.

\newcommand{\ip}[2]{(#1, #2)}
                             % Defines \ip{arg1}{arg2} to mean
                             % (arg1, arg2).

%\newcommand{\ip}[2]{\langle #1 | #2\rangle}
                             % This is an alternative definition of
                             % \ip that is commented out.

\begin{document}             % End of preamble and beginning of text.

\maketitle                   % Produces the title.

This is an example input file.  Comparing it with
the output it generates can show you how to
produce a simple document of your own.

\section{Ordinary Text}      % Produces section heading.  Lower-level
                             % sections are begun with similar
                             % \subsection and \subsubsection commands.

The ends  of words and sentences are marked
  by   spaces. It  doesn't matter how many
spaces    you type; one is as good as 100.  The
end of   a line counts as a space.

One   or more   blank lines denote the  end
of  a paragraph.

Since any number of consecutive spaces are treated
like a single one, the formatting of the input
file makes no difference to
      \LaTeX,                % The \LaTeX command generates the LaTeX logo.
but it makes a difference to you.  When you use
\LaTeX, making your input file as easy to read
as possible will be a great help as you write
your document and when you change it.  This sample
file shows how you can add comments to your own input
file.

Because printing is different from typewriting,
there are a number of things that you have to do
differently when preparing an input file than if
you were just typing the document directly.
Quotation marks like
       ``this''
have to be handled specially, as do quotes within
quotes:
       ``\,`this'            % \, separates the double and single quote.
        is what I just
        wrote, not  `that'\,''.

Dashes come in three sizes: an
       intra-word
dash, a medium dash for number ranges like
       1--2,
and a punctuation
       dash---like
this.

A sentence-ending space should be larger than the
space between words within a sentence.  You
sometimes have to type special commands in
conjunction with punctuation characters to get
this right, as in the following sentence.
       Gnats, gnus, etc.\ all  % `\ ' makes an inter-word space.
       begin with G\@.         % \@ marks end-of-sentence punctuation.
You should check the spaces after periods when
reading your output to make sure you haven't
forgotten any special cases.  Generating an
ellipsis
       \ldots\               % `\ ' is needed after `\ldots' because TeX
                             % ignores spaces after command names like \ldots
                             % made from \ + letters.
                             %
                             % Note how a `%' character causes TeX to ignore
                             % the end of the input line, so these blank lines
                             % do not start a new paragraph.
                             %
with the right spacing around the periods requires
a special command.

\LaTeX\ interprets some common characters as
commands, so you must type special commands to
generate them.  These characters include the
following:
       \$ \& \% \# \{ and \}.

In printing, text is usually emphasized with an
       \emph{italic}
type style.

\begin{em}
   A long segment of text can also be emphasized
   in this way.  Text within such a segment can be
   given \emph{additional} emphasis.
\end{em}

It is sometimes necessary to prevent \LaTeX\ from
breaking a line where it might otherwise do so.
This may be at a space, as between the ``Mr.''\ and
``Jones'' in
       ``Mr.~Jones'',        % ~ produces an unbreakable interword space.
or within a word---especially when the word is a
symbol like
       \mbox{\emph{itemnum}}
that makes little sense when hyphenated across
lines.

Footnotes\footnote{This is an example of a footnote.}
pose no problem.

\LaTeX\ is good at typesetting mathematical formulas
like
       \( x-3y + z = 7 \)
or
       \( a_{1} > x^{2n} + y^{2n} > x' \)
or
       \( \ip{A}{B} = \sum_{i} a_{i} b_{i} \).
The spaces you type in a formula are
ignored.  Remember that a letter like
       $x$                   % $ ... $  and  \( ... \)  are equivalent
is a formula when it denotes a mathematical
symbol, and it should be typed as one.

\section{Displayed Text}

Text is displayed by indenting it from the left
margin.  Quotations are commonly displayed.  There
are short quotations
\begin{quote}
   This is a short quotation.  It consists of a
   single paragraph of text.  See how it is formatted.
\end{quote}
and longer ones.
\begin{quotation}
   This is a longer quotation.  It consists of two
   paragraphs of text, neither of which are
   particularly interesting.

   This is the second paragraph of the quotation.  It
   is just as dull as the first paragraph.
\end{quotation}
Another frequently-displayed structure is a list.
The following is an example of an \emph{itemized}
list.
\begin{itemize}
   \item This is the first item of an itemized list.
         Each item in the list is marked with a ``tick''.
         You don't have to worry about what kind of tick
         mark is used.

   \item This is the second item of the list.  It
         contains another list nested inside it.  The inner
         list is an \emph{enumerated} list.
         \begin{enumerate}
            \item This is the first item of an enumerated
                  list that is nested within the itemized list.

            \item This is the second item of the inner list.
                  \LaTeX\ allows you to nest lists deeper than
                  you really should.
         \end{enumerate}
         This is the rest of the second item of the outer
         list.  It is no more interesting than any other
         part of the item.
   \item This is the third item of the list.
\end{itemize}
You can even display poetry.
\begin{verse}
   There is an environment
    for verse \\             % The \\ command separates lines
   Whose features some poets % within a stanza.
   will curse.

                             % One or more blank lines separate stanzas.

   For instead of making\\
   Them do \emph{all} line breaking, \\
   It allows them to put too many words on a line when they'd rather be
   forced to be terse.
\end{verse}

Mathematical formulas may also be displayed.  A
displayed formula
is
one-line long; multiline
formulas require special formatting instructions.
   \[  \ip{\Gamma}{\psi'} = x'' + y^{2} + z_{i}^{n}\]
Don't start a paragraph with a displayed equation,
nor make one a paragraph by itself.

\end{document}               % End of document.

\section{Bench Huge Section 859}
% synthetic bench section 859
% This is a sample LaTeX input file.  (Version of 12 August 2004.)
%
% A '%' character causes TeX to ignore all remaining text on the line,
% and is used for comments like this one.

\documentclass{article}      % Specifies the document class

                             % The preamble begins here.
\title{An Example Document}  % Declares the document's title.
\author{Leslie Lamport}      % Declares the author's name.
\date{January 21, 1994}      % Deleting this command produces today's date.

\newcommand{\ip}[2]{(#1, #2)}
                             % Defines \ip{arg1}{arg2} to mean
                             % (arg1, arg2).

%\newcommand{\ip}[2]{\langle #1 | #2\rangle}
                             % This is an alternative definition of
                             % \ip that is commented out.

\begin{document}             % End of preamble and beginning of text.

\maketitle                   % Produces the title.

This is an example input file.  Comparing it with
the output it generates can show you how to
produce a simple document of your own.

\section{Ordinary Text}      % Produces section heading.  Lower-level
                             % sections are begun with similar
                             % \subsection and \subsubsection commands.

The ends  of words and sentences are marked
  by   spaces. It  doesn't matter how many
spaces    you type; one is as good as 100.  The
end of   a line counts as a space.

One   or more   blank lines denote the  end
of  a paragraph.

Since any number of consecutive spaces are treated
like a single one, the formatting of the input
file makes no difference to
      \LaTeX,                % The \LaTeX command generates the LaTeX logo.
but it makes a difference to you.  When you use
\LaTeX, making your input file as easy to read
as possible will be a great help as you write
your document and when you change it.  This sample
file shows how you can add comments to your own input
file.

Because printing is different from typewriting,
there are a number of things that you have to do
differently when preparing an input file than if
you were just typing the document directly.
Quotation marks like
       ``this''
have to be handled specially, as do quotes within
quotes:
       ``\,`this'            % \, separates the double and single quote.
        is what I just
        wrote, not  `that'\,''.

Dashes come in three sizes: an
       intra-word
dash, a medium dash for number ranges like
       1--2,
and a punctuation
       dash---like
this.

A sentence-ending space should be larger than the
space between words within a sentence.  You
sometimes have to type special commands in
conjunction with punctuation characters to get
this right, as in the following sentence.
       Gnats, gnus, etc.\ all  % `\ ' makes an inter-word space.
       begin with G\@.         % \@ marks end-of-sentence punctuation.
You should check the spaces after periods when
reading your output to make sure you haven't
forgotten any special cases.  Generating an
ellipsis
       \ldots\               % `\ ' is needed after `\ldots' because TeX
                             % ignores spaces after command names like \ldots
                             % made from \ + letters.
                             %
                             % Note how a `%' character causes TeX to ignore
                             % the end of the input line, so these blank lines
                             % do not start a new paragraph.
                             %
with the right spacing around the periods requires
a special command.

\LaTeX\ interprets some common characters as
commands, so you must type special commands to
generate them.  These characters include the
following:
       \$ \& \% \# \{ and \}.

In printing, text is usually emphasized with an
       \emph{italic}
type style.

\begin{em}
   A long segment of text can also be emphasized
   in this way.  Text within such a segment can be
   given \emph{additional} emphasis.
\end{em}

It is sometimes necessary to prevent \LaTeX\ from
breaking a line where it might otherwise do so.
This may be at a space, as between the ``Mr.''\ and
``Jones'' in
       ``Mr.~Jones'',        % ~ produces an unbreakable interword space.
or within a word---especially when the word is a
symbol like
       \mbox{\emph{itemnum}}
that makes little sense when hyphenated across
lines.

Footnotes\footnote{This is an example of a footnote.}
pose no problem.

\LaTeX\ is good at typesetting mathematical formulas
like
       \( x-3y + z = 7 \)
or
       \( a_{1} > x^{2n} + y^{2n} > x' \)
or
       \( \ip{A}{B} = \sum_{i} a_{i} b_{i} \).
The spaces you type in a formula are
ignored.  Remember that a letter like
       $x$                   % $ ... $  and  \( ... \)  are equivalent
is a formula when it denotes a mathematical
symbol, and it should be typed as one.

\section{Displayed Text}

Text is displayed by indenting it from the left
margin.  Quotations are commonly displayed.  There
are short quotations
\begin{quote}
   This is a short quotation.  It consists of a
   single paragraph of text.  See how it is formatted.
\end{quote}
and longer ones.
\begin{quotation}
   This is a longer quotation.  It consists of two
   paragraphs of text, neither of which are
   particularly interesting.

   This is the second paragraph of the quotation.  It
   is just as dull as the first paragraph.
\end{quotation}
Another frequently-displayed structure is a list.
The following is an example of an \emph{itemized}
list.
\begin{itemize}
   \item This is the first item of an itemized list.
         Each item in the list is marked with a ``tick''.
         You don't have to worry about what kind of tick
         mark is used.

   \item This is the second item of the list.  It
         contains another list nested inside it.  The inner
         list is an \emph{enumerated} list.
         \begin{enumerate}
            \item This is the first item of an enumerated
                  list that is nested within the itemized list.

            \item This is the second item of the inner list.
                  \LaTeX\ allows you to nest lists deeper than
                  you really should.
         \end{enumerate}
         This is the rest of the second item of the outer
         list.  It is no more interesting than any other
         part of the item.
   \item This is the third item of the list.
\end{itemize}
You can even display poetry.
\begin{verse}
   There is an environment
    for verse \\             % The \\ command separates lines
   Whose features some poets % within a stanza.
   will curse.

                             % One or more blank lines separate stanzas.

   For instead of making\\
   Them do \emph{all} line breaking, \\
   It allows them to put too many words on a line when they'd rather be
   forced to be terse.
\end{verse}

Mathematical formulas may also be displayed.  A
displayed formula
is
one-line long; multiline
formulas require special formatting instructions.
   \[  \ip{\Gamma}{\psi'} = x'' + y^{2} + z_{i}^{n}\]
Don't start a paragraph with a displayed equation,
nor make one a paragraph by itself.

\end{document}               % End of document.

\section{Bench Huge Section 860}
% synthetic bench section 860
% This is a sample LaTeX input file.  (Version of 12 August 2004.)
%
% A '%' character causes TeX to ignore all remaining text on the line,
% and is used for comments like this one.

\documentclass{article}      % Specifies the document class

                             % The preamble begins here.
\title{An Example Document}  % Declares the document's title.
\author{Leslie Lamport}      % Declares the author's name.
\date{January 21, 1994}      % Deleting this command produces today's date.

\newcommand{\ip}[2]{(#1, #2)}
                             % Defines \ip{arg1}{arg2} to mean
                             % (arg1, arg2).

%\newcommand{\ip}[2]{\langle #1 | #2\rangle}
                             % This is an alternative definition of
                             % \ip that is commented out.

\begin{document}             % End of preamble and beginning of text.

\maketitle                   % Produces the title.

This is an example input file.  Comparing it with
the output it generates can show you how to
produce a simple document of your own.

\section{Ordinary Text}      % Produces section heading.  Lower-level
                             % sections are begun with similar
                             % \subsection and \subsubsection commands.

The ends  of words and sentences are marked
  by   spaces. It  doesn't matter how many
spaces    you type; one is as good as 100.  The
end of   a line counts as a space.

One   or more   blank lines denote the  end
of  a paragraph.

Since any number of consecutive spaces are treated
like a single one, the formatting of the input
file makes no difference to
      \LaTeX,                % The \LaTeX command generates the LaTeX logo.
but it makes a difference to you.  When you use
\LaTeX, making your input file as easy to read
as possible will be a great help as you write
your document and when you change it.  This sample
file shows how you can add comments to your own input
file.

Because printing is different from typewriting,
there are a number of things that you have to do
differently when preparing an input file than if
you were just typing the document directly.
Quotation marks like
       ``this''
have to be handled specially, as do quotes within
quotes:
       ``\,`this'            % \, separates the double and single quote.
        is what I just
        wrote, not  `that'\,''.

Dashes come in three sizes: an
       intra-word
dash, a medium dash for number ranges like
       1--2,
and a punctuation
       dash---like
this.

A sentence-ending space should be larger than the
space between words within a sentence.  You
sometimes have to type special commands in
conjunction with punctuation characters to get
this right, as in the following sentence.
       Gnats, gnus, etc.\ all  % `\ ' makes an inter-word space.
       begin with G\@.         % \@ marks end-of-sentence punctuation.
You should check the spaces after periods when
reading your output to make sure you haven't
forgotten any special cases.  Generating an
ellipsis
       \ldots\               % `\ ' is needed after `\ldots' because TeX
                             % ignores spaces after command names like \ldots
                             % made from \ + letters.
                             %
                             % Note how a `%' character causes TeX to ignore
                             % the end of the input line, so these blank lines
                             % do not start a new paragraph.
                             %
with the right spacing around the periods requires
a special command.

\LaTeX\ interprets some common characters as
commands, so you must type special commands to
generate them.  These characters include the
following:
       \$ \& \% \# \{ and \}.

In printing, text is usually emphasized with an
       \emph{italic}
type style.

\begin{em}
   A long segment of text can also be emphasized
   in this way.  Text within such a segment can be
   given \emph{additional} emphasis.
\end{em}

It is sometimes necessary to prevent \LaTeX\ from
breaking a line where it might otherwise do so.
This may be at a space, as between the ``Mr.''\ and
``Jones'' in
       ``Mr.~Jones'',        % ~ produces an unbreakable interword space.
or within a word---especially when the word is a
symbol like
       \mbox{\emph{itemnum}}
that makes little sense when hyphenated across
lines.

Footnotes\footnote{This is an example of a footnote.}
pose no problem.

\LaTeX\ is good at typesetting mathematical formulas
like
       \( x-3y + z = 7 \)
or
       \( a_{1} > x^{2n} + y^{2n} > x' \)
or
       \( \ip{A}{B} = \sum_{i} a_{i} b_{i} \).
The spaces you type in a formula are
ignored.  Remember that a letter like
       $x$                   % $ ... $  and  \( ... \)  are equivalent
is a formula when it denotes a mathematical
symbol, and it should be typed as one.

\section{Displayed Text}

Text is displayed by indenting it from the left
margin.  Quotations are commonly displayed.  There
are short quotations
\begin{quote}
   This is a short quotation.  It consists of a
   single paragraph of text.  See how it is formatted.
\end{quote}
and longer ones.
\begin{quotation}
   This is a longer quotation.  It consists of two
   paragraphs of text, neither of which are
   particularly interesting.

   This is the second paragraph of the quotation.  It
   is just as dull as the first paragraph.
\end{quotation}
Another frequently-displayed structure is a list.
The following is an example of an \emph{itemized}
list.
\begin{itemize}
   \item This is the first item of an itemized list.
         Each item in the list is marked with a ``tick''.
         You don't have to worry about what kind of tick
         mark is used.

   \item This is the second item of the list.  It
         contains another list nested inside it.  The inner
         list is an \emph{enumerated} list.
         \begin{enumerate}
            \item This is the first item of an enumerated
                  list that is nested within the itemized list.

            \item This is the second item of the inner list.
                  \LaTeX\ allows you to nest lists deeper than
                  you really should.
         \end{enumerate}
         This is the rest of the second item of the outer
         list.  It is no more interesting than any other
         part of the item.
   \item This is the third item of the list.
\end{itemize}
You can even display poetry.
\begin{verse}
   There is an environment
    for verse \\             % The \\ command separates lines
   Whose features some poets % within a stanza.
   will curse.

                             % One or more blank lines separate stanzas.

   For instead of making\\
   Them do \emph{all} line breaking, \\
   It allows them to put too many words on a line when they'd rather be
   forced to be terse.
\end{verse}

Mathematical formulas may also be displayed.  A
displayed formula
is
one-line long; multiline
formulas require special formatting instructions.
   \[  \ip{\Gamma}{\psi'} = x'' + y^{2} + z_{i}^{n}\]
Don't start a paragraph with a displayed equation,
nor make one a paragraph by itself.

\end{document}               % End of document.

\section{Bench Huge Section 861}
% synthetic bench section 861
% This is a sample LaTeX input file.  (Version of 12 August 2004.)
%
% A '%' character causes TeX to ignore all remaining text on the line,
% and is used for comments like this one.

\documentclass{article}      % Specifies the document class

                             % The preamble begins here.
\title{An Example Document}  % Declares the document's title.
\author{Leslie Lamport}      % Declares the author's name.
\date{January 21, 1994}      % Deleting this command produces today's date.

\newcommand{\ip}[2]{(#1, #2)}
                             % Defines \ip{arg1}{arg2} to mean
                             % (arg1, arg2).

%\newcommand{\ip}[2]{\langle #1 | #2\rangle}
                             % This is an alternative definition of
                             % \ip that is commented out.

\begin{document}             % End of preamble and beginning of text.

\maketitle                   % Produces the title.

This is an example input file.  Comparing it with
the output it generates can show you how to
produce a simple document of your own.

\section{Ordinary Text}      % Produces section heading.  Lower-level
                             % sections are begun with similar
                             % \subsection and \subsubsection commands.

The ends  of words and sentences are marked
  by   spaces. It  doesn't matter how many
spaces    you type; one is as good as 100.  The
end of   a line counts as a space.

One   or more   blank lines denote the  end
of  a paragraph.

Since any number of consecutive spaces are treated
like a single one, the formatting of the input
file makes no difference to
      \LaTeX,                % The \LaTeX command generates the LaTeX logo.
but it makes a difference to you.  When you use
\LaTeX, making your input file as easy to read
as possible will be a great help as you write
your document and when you change it.  This sample
file shows how you can add comments to your own input
file.

Because printing is different from typewriting,
there are a number of things that you have to do
differently when preparing an input file than if
you were just typing the document directly.
Quotation marks like
       ``this''
have to be handled specially, as do quotes within
quotes:
       ``\,`this'            % \, separates the double and single quote.
        is what I just
        wrote, not  `that'\,''.

Dashes come in three sizes: an
       intra-word
dash, a medium dash for number ranges like
       1--2,
and a punctuation
       dash---like
this.

A sentence-ending space should be larger than the
space between words within a sentence.  You
sometimes have to type special commands in
conjunction with punctuation characters to get
this right, as in the following sentence.
       Gnats, gnus, etc.\ all  % `\ ' makes an inter-word space.
       begin with G\@.         % \@ marks end-of-sentence punctuation.
You should check the spaces after periods when
reading your output to make sure you haven't
forgotten any special cases.  Generating an
ellipsis
       \ldots\               % `\ ' is needed after `\ldots' because TeX
                             % ignores spaces after command names like \ldots
                             % made from \ + letters.
                             %
                             % Note how a `%' character causes TeX to ignore
                             % the end of the input line, so these blank lines
                             % do not start a new paragraph.
                             %
with the right spacing around the periods requires
a special command.

\LaTeX\ interprets some common characters as
commands, so you must type special commands to
generate them.  These characters include the
following:
       \$ \& \% \# \{ and \}.

In printing, text is usually emphasized with an
       \emph{italic}
type style.

\begin{em}
   A long segment of text can also be emphasized
   in this way.  Text within such a segment can be
   given \emph{additional} emphasis.
\end{em}

It is sometimes necessary to prevent \LaTeX\ from
breaking a line where it might otherwise do so.
This may be at a space, as between the ``Mr.''\ and
``Jones'' in
       ``Mr.~Jones'',        % ~ produces an unbreakable interword space.
or within a word---especially when the word is a
symbol like
       \mbox{\emph{itemnum}}
that makes little sense when hyphenated across
lines.

Footnotes\footnote{This is an example of a footnote.}
pose no problem.

\LaTeX\ is good at typesetting mathematical formulas
like
       \( x-3y + z = 7 \)
or
       \( a_{1} > x^{2n} + y^{2n} > x' \)
or
       \( \ip{A}{B} = \sum_{i} a_{i} b_{i} \).
The spaces you type in a formula are
ignored.  Remember that a letter like
       $x$                   % $ ... $  and  \( ... \)  are equivalent
is a formula when it denotes a mathematical
symbol, and it should be typed as one.

\section{Displayed Text}

Text is displayed by indenting it from the left
margin.  Quotations are commonly displayed.  There
are short quotations
\begin{quote}
   This is a short quotation.  It consists of a
   single paragraph of text.  See how it is formatted.
\end{quote}
and longer ones.
\begin{quotation}
   This is a longer quotation.  It consists of two
   paragraphs of text, neither of which are
   particularly interesting.

   This is the second paragraph of the quotation.  It
   is just as dull as the first paragraph.
\end{quotation}
Another frequently-displayed structure is a list.
The following is an example of an \emph{itemized}
list.
\begin{itemize}
   \item This is the first item of an itemized list.
         Each item in the list is marked with a ``tick''.
         You don't have to worry about what kind of tick
         mark is used.

   \item This is the second item of the list.  It
         contains another list nested inside it.  The inner
         list is an \emph{enumerated} list.
         \begin{enumerate}
            \item This is the first item of an enumerated
                  list that is nested within the itemized list.

            \item This is the second item of the inner list.
                  \LaTeX\ allows you to nest lists deeper than
                  you really should.
         \end{enumerate}
         This is the rest of the second item of the outer
         list.  It is no more interesting than any other
         part of the item.
   \item This is the third item of the list.
\end{itemize}
You can even display poetry.
\begin{verse}
   There is an environment
    for verse \\             % The \\ command separates lines
   Whose features some poets % within a stanza.
   will curse.

                             % One or more blank lines separate stanzas.

   For instead of making\\
   Them do \emph{all} line breaking, \\
   It allows them to put too many words on a line when they'd rather be
   forced to be terse.
\end{verse}

Mathematical formulas may also be displayed.  A
displayed formula
is
one-line long; multiline
formulas require special formatting instructions.
   \[  \ip{\Gamma}{\psi'} = x'' + y^{2} + z_{i}^{n}\]
Don't start a paragraph with a displayed equation,
nor make one a paragraph by itself.

\end{document}               % End of document.

\section{Bench Huge Section 862}
% synthetic bench section 862
% This is a sample LaTeX input file.  (Version of 12 August 2004.)
%
% A '%' character causes TeX to ignore all remaining text on the line,
% and is used for comments like this one.

\documentclass{article}      % Specifies the document class

                             % The preamble begins here.
\title{An Example Document}  % Declares the document's title.
\author{Leslie Lamport}      % Declares the author's name.
\date{January 21, 1994}      % Deleting this command produces today's date.

\newcommand{\ip}[2]{(#1, #2)}
                             % Defines \ip{arg1}{arg2} to mean
                             % (arg1, arg2).

%\newcommand{\ip}[2]{\langle #1 | #2\rangle}
                             % This is an alternative definition of
                             % \ip that is commented out.

\begin{document}             % End of preamble and beginning of text.

\maketitle                   % Produces the title.

This is an example input file.  Comparing it with
the output it generates can show you how to
produce a simple document of your own.

\section{Ordinary Text}      % Produces section heading.  Lower-level
                             % sections are begun with similar
                             % \subsection and \subsubsection commands.

The ends  of words and sentences are marked
  by   spaces. It  doesn't matter how many
spaces    you type; one is as good as 100.  The
end of   a line counts as a space.

One   or more   blank lines denote the  end
of  a paragraph.

Since any number of consecutive spaces are treated
like a single one, the formatting of the input
file makes no difference to
      \LaTeX,                % The \LaTeX command generates the LaTeX logo.
but it makes a difference to you.  When you use
\LaTeX, making your input file as easy to read
as possible will be a great help as you write
your document and when you change it.  This sample
file shows how you can add comments to your own input
file.

Because printing is different from typewriting,
there are a number of things that you have to do
differently when preparing an input file than if
you were just typing the document directly.
Quotation marks like
       ``this''
have to be handled specially, as do quotes within
quotes:
       ``\,`this'            % \, separates the double and single quote.
        is what I just
        wrote, not  `that'\,''.

Dashes come in three sizes: an
       intra-word
dash, a medium dash for number ranges like
       1--2,
and a punctuation
       dash---like
this.

A sentence-ending space should be larger than the
space between words within a sentence.  You
sometimes have to type special commands in
conjunction with punctuation characters to get
this right, as in the following sentence.
       Gnats, gnus, etc.\ all  % `\ ' makes an inter-word space.
       begin with G\@.         % \@ marks end-of-sentence punctuation.
You should check the spaces after periods when
reading your output to make sure you haven't
forgotten any special cases.  Generating an
ellipsis
       \ldots\               % `\ ' is needed after `\ldots' because TeX
                             % ignores spaces after command names like \ldots
                             % made from \ + letters.
                             %
                             % Note how a `%' character causes TeX to ignore
                             % the end of the input line, so these blank lines
                             % do not start a new paragraph.
                             %
with the right spacing around the periods requires
a special command.

\LaTeX\ interprets some common characters as
commands, so you must type special commands to
generate them.  These characters include the
following:
       \$ \& \% \# \{ and \}.

In printing, text is usually emphasized with an
       \emph{italic}
type style.

\begin{em}
   A long segment of text can also be emphasized
   in this way.  Text within such a segment can be
   given \emph{additional} emphasis.
\end{em}

It is sometimes necessary to prevent \LaTeX\ from
breaking a line where it might otherwise do so.
This may be at a space, as between the ``Mr.''\ and
``Jones'' in
       ``Mr.~Jones'',        % ~ produces an unbreakable interword space.
or within a word---especially when the word is a
symbol like
       \mbox{\emph{itemnum}}
that makes little sense when hyphenated across
lines.

Footnotes\footnote{This is an example of a footnote.}
pose no problem.

\LaTeX\ is good at typesetting mathematical formulas
like
       \( x-3y + z = 7 \)
or
       \( a_{1} > x^{2n} + y^{2n} > x' \)
or
       \( \ip{A}{B} = \sum_{i} a_{i} b_{i} \).
The spaces you type in a formula are
ignored.  Remember that a letter like
       $x$                   % $ ... $  and  \( ... \)  are equivalent
is a formula when it denotes a mathematical
symbol, and it should be typed as one.

\section{Displayed Text}

Text is displayed by indenting it from the left
margin.  Quotations are commonly displayed.  There
are short quotations
\begin{quote}
   This is a short quotation.  It consists of a
   single paragraph of text.  See how it is formatted.
\end{quote}
and longer ones.
\begin{quotation}
   This is a longer quotation.  It consists of two
   paragraphs of text, neither of which are
   particularly interesting.

   This is the second paragraph of the quotation.  It
   is just as dull as the first paragraph.
\end{quotation}
Another frequently-displayed structure is a list.
The following is an example of an \emph{itemized}
list.
\begin{itemize}
   \item This is the first item of an itemized list.
         Each item in the list is marked with a ``tick''.
         You don't have to worry about what kind of tick
         mark is used.

   \item This is the second item of the list.  It
         contains another list nested inside it.  The inner
         list is an \emph{enumerated} list.
         \begin{enumerate}
            \item This is the first item of an enumerated
                  list that is nested within the itemized list.

            \item This is the second item of the inner list.
                  \LaTeX\ allows you to nest lists deeper than
                  you really should.
         \end{enumerate}
         This is the rest of the second item of the outer
         list.  It is no more interesting than any other
         part of the item.
   \item This is the third item of the list.
\end{itemize}
You can even display poetry.
\begin{verse}
   There is an environment
    for verse \\             % The \\ command separates lines
   Whose features some poets % within a stanza.
   will curse.

                             % One or more blank lines separate stanzas.

   For instead of making\\
   Them do \emph{all} line breaking, \\
   It allows them to put too many words on a line when they'd rather be
   forced to be terse.
\end{verse}

Mathematical formulas may also be displayed.  A
displayed formula
is
one-line long; multiline
formulas require special formatting instructions.
   \[  \ip{\Gamma}{\psi'} = x'' + y^{2} + z_{i}^{n}\]
Don't start a paragraph with a displayed equation,
nor make one a paragraph by itself.

\end{document}               % End of document.

\section{Bench Huge Section 863}
% synthetic bench section 863
% This is a sample LaTeX input file.  (Version of 12 August 2004.)
%
% A '%' character causes TeX to ignore all remaining text on the line,
% and is used for comments like this one.

\documentclass{article}      % Specifies the document class

                             % The preamble begins here.
\title{An Example Document}  % Declares the document's title.
\author{Leslie Lamport}      % Declares the author's name.
\date{January 21, 1994}      % Deleting this command produces today's date.

\newcommand{\ip}[2]{(#1, #2)}
                             % Defines \ip{arg1}{arg2} to mean
                             % (arg1, arg2).

%\newcommand{\ip}[2]{\langle #1 | #2\rangle}
                             % This is an alternative definition of
                             % \ip that is commented out.

\begin{document}             % End of preamble and beginning of text.

\maketitle                   % Produces the title.

This is an example input file.  Comparing it with
the output it generates can show you how to
produce a simple document of your own.

\section{Ordinary Text}      % Produces section heading.  Lower-level
                             % sections are begun with similar
                             % \subsection and \subsubsection commands.

The ends  of words and sentences are marked
  by   spaces. It  doesn't matter how many
spaces    you type; one is as good as 100.  The
end of   a line counts as a space.

One   or more   blank lines denote the  end
of  a paragraph.

Since any number of consecutive spaces are treated
like a single one, the formatting of the input
file makes no difference to
      \LaTeX,                % The \LaTeX command generates the LaTeX logo.
but it makes a difference to you.  When you use
\LaTeX, making your input file as easy to read
as possible will be a great help as you write
your document and when you change it.  This sample
file shows how you can add comments to your own input
file.

Because printing is different from typewriting,
there are a number of things that you have to do
differently when preparing an input file than if
you were just typing the document directly.
Quotation marks like
       ``this''
have to be handled specially, as do quotes within
quotes:
       ``\,`this'            % \, separates the double and single quote.
        is what I just
        wrote, not  `that'\,''.

Dashes come in three sizes: an
       intra-word
dash, a medium dash for number ranges like
       1--2,
and a punctuation
       dash---like
this.

A sentence-ending space should be larger than the
space between words within a sentence.  You
sometimes have to type special commands in
conjunction with punctuation characters to get
this right, as in the following sentence.
       Gnats, gnus, etc.\ all  % `\ ' makes an inter-word space.
       begin with G\@.         % \@ marks end-of-sentence punctuation.
You should check the spaces after periods when
reading your output to make sure you haven't
forgotten any special cases.  Generating an
ellipsis
       \ldots\               % `\ ' is needed after `\ldots' because TeX
                             % ignores spaces after command names like \ldots
                             % made from \ + letters.
                             %
                             % Note how a `%' character causes TeX to ignore
                             % the end of the input line, so these blank lines
                             % do not start a new paragraph.
                             %
with the right spacing around the periods requires
a special command.

\LaTeX\ interprets some common characters as
commands, so you must type special commands to
generate them.  These characters include the
following:
       \$ \& \% \# \{ and \}.

In printing, text is usually emphasized with an
       \emph{italic}
type style.

\begin{em}
   A long segment of text can also be emphasized
   in this way.  Text within such a segment can be
   given \emph{additional} emphasis.
\end{em}

It is sometimes necessary to prevent \LaTeX\ from
breaking a line where it might otherwise do so.
This may be at a space, as between the ``Mr.''\ and
``Jones'' in
       ``Mr.~Jones'',        % ~ produces an unbreakable interword space.
or within a word---especially when the word is a
symbol like
       \mbox{\emph{itemnum}}
that makes little sense when hyphenated across
lines.

Footnotes\footnote{This is an example of a footnote.}
pose no problem.

\LaTeX\ is good at typesetting mathematical formulas
like
       \( x-3y + z = 7 \)
or
       \( a_{1} > x^{2n} + y^{2n} > x' \)
or
       \( \ip{A}{B} = \sum_{i} a_{i} b_{i} \).
The spaces you type in a formula are
ignored.  Remember that a letter like
       $x$                   % $ ... $  and  \( ... \)  are equivalent
is a formula when it denotes a mathematical
symbol, and it should be typed as one.

\section{Displayed Text}

Text is displayed by indenting it from the left
margin.  Quotations are commonly displayed.  There
are short quotations
\begin{quote}
   This is a short quotation.  It consists of a
   single paragraph of text.  See how it is formatted.
\end{quote}
and longer ones.
\begin{quotation}
   This is a longer quotation.  It consists of two
   paragraphs of text, neither of which are
   particularly interesting.

   This is the second paragraph of the quotation.  It
   is just as dull as the first paragraph.
\end{quotation}
Another frequently-displayed structure is a list.
The following is an example of an \emph{itemized}
list.
\begin{itemize}
   \item This is the first item of an itemized list.
         Each item in the list is marked with a ``tick''.
         You don't have to worry about what kind of tick
         mark is used.

   \item This is the second item of the list.  It
         contains another list nested inside it.  The inner
         list is an \emph{enumerated} list.
         \begin{enumerate}
            \item This is the first item of an enumerated
                  list that is nested within the itemized list.

            \item This is the second item of the inner list.
                  \LaTeX\ allows you to nest lists deeper than
                  you really should.
         \end{enumerate}
         This is the rest of the second item of the outer
         list.  It is no more interesting than any other
         part of the item.
   \item This is the third item of the list.
\end{itemize}
You can even display poetry.
\begin{verse}
   There is an environment
    for verse \\             % The \\ command separates lines
   Whose features some poets % within a stanza.
   will curse.

                             % One or more blank lines separate stanzas.

   For instead of making\\
   Them do \emph{all} line breaking, \\
   It allows them to put too many words on a line when they'd rather be
   forced to be terse.
\end{verse}

Mathematical formulas may also be displayed.  A
displayed formula
is
one-line long; multiline
formulas require special formatting instructions.
   \[  \ip{\Gamma}{\psi'} = x'' + y^{2} + z_{i}^{n}\]
Don't start a paragraph with a displayed equation,
nor make one a paragraph by itself.

\end{document}               % End of document.

\section{Bench Huge Section 864}
% synthetic bench section 864
% This is a sample LaTeX input file.  (Version of 12 August 2004.)
%
% A '%' character causes TeX to ignore all remaining text on the line,
% and is used for comments like this one.

\documentclass{article}      % Specifies the document class

                             % The preamble begins here.
\title{An Example Document}  % Declares the document's title.
\author{Leslie Lamport}      % Declares the author's name.
\date{January 21, 1994}      % Deleting this command produces today's date.

\newcommand{\ip}[2]{(#1, #2)}
                             % Defines \ip{arg1}{arg2} to mean
                             % (arg1, arg2).

%\newcommand{\ip}[2]{\langle #1 | #2\rangle}
                             % This is an alternative definition of
                             % \ip that is commented out.

\begin{document}             % End of preamble and beginning of text.

\maketitle                   % Produces the title.

This is an example input file.  Comparing it with
the output it generates can show you how to
produce a simple document of your own.

\section{Ordinary Text}      % Produces section heading.  Lower-level
                             % sections are begun with similar
                             % \subsection and \subsubsection commands.

The ends  of words and sentences are marked
  by   spaces. It  doesn't matter how many
spaces    you type; one is as good as 100.  The
end of   a line counts as a space.

One   or more   blank lines denote the  end
of  a paragraph.

Since any number of consecutive spaces are treated
like a single one, the formatting of the input
file makes no difference to
      \LaTeX,                % The \LaTeX command generates the LaTeX logo.
but it makes a difference to you.  When you use
\LaTeX, making your input file as easy to read
as possible will be a great help as you write
your document and when you change it.  This sample
file shows how you can add comments to your own input
file.

Because printing is different from typewriting,
there are a number of things that you have to do
differently when preparing an input file than if
you were just typing the document directly.
Quotation marks like
       ``this''
have to be handled specially, as do quotes within
quotes:
       ``\,`this'            % \, separates the double and single quote.
        is what I just
        wrote, not  `that'\,''.

Dashes come in three sizes: an
       intra-word
dash, a medium dash for number ranges like
       1--2,
and a punctuation
       dash---like
this.

A sentence-ending space should be larger than the
space between words within a sentence.  You
sometimes have to type special commands in
conjunction with punctuation characters to get
this right, as in the following sentence.
       Gnats, gnus, etc.\ all  % `\ ' makes an inter-word space.
       begin with G\@.         % \@ marks end-of-sentence punctuation.
You should check the spaces after periods when
reading your output to make sure you haven't
forgotten any special cases.  Generating an
ellipsis
       \ldots\               % `\ ' is needed after `\ldots' because TeX
                             % ignores spaces after command names like \ldots
                             % made from \ + letters.
                             %
                             % Note how a `%' character causes TeX to ignore
                             % the end of the input line, so these blank lines
                             % do not start a new paragraph.
                             %
with the right spacing around the periods requires
a special command.

\LaTeX\ interprets some common characters as
commands, so you must type special commands to
generate them.  These characters include the
following:
       \$ \& \% \# \{ and \}.

In printing, text is usually emphasized with an
       \emph{italic}
type style.

\begin{em}
   A long segment of text can also be emphasized
   in this way.  Text within such a segment can be
   given \emph{additional} emphasis.
\end{em}

It is sometimes necessary to prevent \LaTeX\ from
breaking a line where it might otherwise do so.
This may be at a space, as between the ``Mr.''\ and
``Jones'' in
       ``Mr.~Jones'',        % ~ produces an unbreakable interword space.
or within a word---especially when the word is a
symbol like
       \mbox{\emph{itemnum}}
that makes little sense when hyphenated across
lines.

Footnotes\footnote{This is an example of a footnote.}
pose no problem.

\LaTeX\ is good at typesetting mathematical formulas
like
       \( x-3y + z = 7 \)
or
       \( a_{1} > x^{2n} + y^{2n} > x' \)
or
       \( \ip{A}{B} = \sum_{i} a_{i} b_{i} \).
The spaces you type in a formula are
ignored.  Remember that a letter like
       $x$                   % $ ... $  and  \( ... \)  are equivalent
is a formula when it denotes a mathematical
symbol, and it should be typed as one.

\section{Displayed Text}

Text is displayed by indenting it from the left
margin.  Quotations are commonly displayed.  There
are short quotations
\begin{quote}
   This is a short quotation.  It consists of a
   single paragraph of text.  See how it is formatted.
\end{quote}
and longer ones.
\begin{quotation}
   This is a longer quotation.  It consists of two
   paragraphs of text, neither of which are
   particularly interesting.

   This is the second paragraph of the quotation.  It
   is just as dull as the first paragraph.
\end{quotation}
Another frequently-displayed structure is a list.
The following is an example of an \emph{itemized}
list.
\begin{itemize}
   \item This is the first item of an itemized list.
         Each item in the list is marked with a ``tick''.
         You don't have to worry about what kind of tick
         mark is used.

   \item This is the second item of the list.  It
         contains another list nested inside it.  The inner
         list is an \emph{enumerated} list.
         \begin{enumerate}
            \item This is the first item of an enumerated
                  list that is nested within the itemized list.

            \item This is the second item of the inner list.
                  \LaTeX\ allows you to nest lists deeper than
                  you really should.
         \end{enumerate}
         This is the rest of the second item of the outer
         list.  It is no more interesting than any other
         part of the item.
   \item This is the third item of the list.
\end{itemize}
You can even display poetry.
\begin{verse}
   There is an environment
    for verse \\             % The \\ command separates lines
   Whose features some poets % within a stanza.
   will curse.

                             % One or more blank lines separate stanzas.

   For instead of making\\
   Them do \emph{all} line breaking, \\
   It allows them to put too many words on a line when they'd rather be
   forced to be terse.
\end{verse}

Mathematical formulas may also be displayed.  A
displayed formula
is
one-line long; multiline
formulas require special formatting instructions.
   \[  \ip{\Gamma}{\psi'} = x'' + y^{2} + z_{i}^{n}\]
Don't start a paragraph with a displayed equation,
nor make one a paragraph by itself.

\end{document}               % End of document.

\section{Bench Huge Section 865}
% synthetic bench section 865
% This is a sample LaTeX input file.  (Version of 12 August 2004.)
%
% A '%' character causes TeX to ignore all remaining text on the line,
% and is used for comments like this one.

\documentclass{article}      % Specifies the document class

                             % The preamble begins here.
\title{An Example Document}  % Declares the document's title.
\author{Leslie Lamport}      % Declares the author's name.
\date{January 21, 1994}      % Deleting this command produces today's date.

\newcommand{\ip}[2]{(#1, #2)}
                             % Defines \ip{arg1}{arg2} to mean
                             % (arg1, arg2).

%\newcommand{\ip}[2]{\langle #1 | #2\rangle}
                             % This is an alternative definition of
                             % \ip that is commented out.

\begin{document}             % End of preamble and beginning of text.

\maketitle                   % Produces the title.

This is an example input file.  Comparing it with
the output it generates can show you how to
produce a simple document of your own.

\section{Ordinary Text}      % Produces section heading.  Lower-level
                             % sections are begun with similar
                             % \subsection and \subsubsection commands.

The ends  of words and sentences are marked
  by   spaces. It  doesn't matter how many
spaces    you type; one is as good as 100.  The
end of   a line counts as a space.

One   or more   blank lines denote the  end
of  a paragraph.

Since any number of consecutive spaces are treated
like a single one, the formatting of the input
file makes no difference to
      \LaTeX,                % The \LaTeX command generates the LaTeX logo.
but it makes a difference to you.  When you use
\LaTeX, making your input file as easy to read
as possible will be a great help as you write
your document and when you change it.  This sample
file shows how you can add comments to your own input
file.

Because printing is different from typewriting,
there are a number of things that you have to do
differently when preparing an input file than if
you were just typing the document directly.
Quotation marks like
       ``this''
have to be handled specially, as do quotes within
quotes:
       ``\,`this'            % \, separates the double and single quote.
        is what I just
        wrote, not  `that'\,''.

Dashes come in three sizes: an
       intra-word
dash, a medium dash for number ranges like
       1--2,
and a punctuation
       dash---like
this.

A sentence-ending space should be larger than the
space between words within a sentence.  You
sometimes have to type special commands in
conjunction with punctuation characters to get
this right, as in the following sentence.
       Gnats, gnus, etc.\ all  % `\ ' makes an inter-word space.
       begin with G\@.         % \@ marks end-of-sentence punctuation.
You should check the spaces after periods when
reading your output to make sure you haven't
forgotten any special cases.  Generating an
ellipsis
       \ldots\               % `\ ' is needed after `\ldots' because TeX
                             % ignores spaces after command names like \ldots
                             % made from \ + letters.
                             %
                             % Note how a `%' character causes TeX to ignore
                             % the end of the input line, so these blank lines
                             % do not start a new paragraph.
                             %
with the right spacing around the periods requires
a special command.

\LaTeX\ interprets some common characters as
commands, so you must type special commands to
generate them.  These characters include the
following:
       \$ \& \% \# \{ and \}.

In printing, text is usually emphasized with an
       \emph{italic}
type style.

\begin{em}
   A long segment of text can also be emphasized
   in this way.  Text within such a segment can be
   given \emph{additional} emphasis.
\end{em}

It is sometimes necessary to prevent \LaTeX\ from
breaking a line where it might otherwise do so.
This may be at a space, as between the ``Mr.''\ and
``Jones'' in
       ``Mr.~Jones'',        % ~ produces an unbreakable interword space.
or within a word---especially when the word is a
symbol like
       \mbox{\emph{itemnum}}
that makes little sense when hyphenated across
lines.

Footnotes\footnote{This is an example of a footnote.}
pose no problem.

\LaTeX\ is good at typesetting mathematical formulas
like
       \( x-3y + z = 7 \)
or
       \( a_{1} > x^{2n} + y^{2n} > x' \)
or
       \( \ip{A}{B} = \sum_{i} a_{i} b_{i} \).
The spaces you type in a formula are
ignored.  Remember that a letter like
       $x$                   % $ ... $  and  \( ... \)  are equivalent
is a formula when it denotes a mathematical
symbol, and it should be typed as one.

\section{Displayed Text}

Text is displayed by indenting it from the left
margin.  Quotations are commonly displayed.  There
are short quotations
\begin{quote}
   This is a short quotation.  It consists of a
   single paragraph of text.  See how it is formatted.
\end{quote}
and longer ones.
\begin{quotation}
   This is a longer quotation.  It consists of two
   paragraphs of text, neither of which are
   particularly interesting.

   This is the second paragraph of the quotation.  It
   is just as dull as the first paragraph.
\end{quotation}
Another frequently-displayed structure is a list.
The following is an example of an \emph{itemized}
list.
\begin{itemize}
   \item This is the first item of an itemized list.
         Each item in the list is marked with a ``tick''.
         You don't have to worry about what kind of tick
         mark is used.

   \item This is the second item of the list.  It
         contains another list nested inside it.  The inner
         list is an \emph{enumerated} list.
         \begin{enumerate}
            \item This is the first item of an enumerated
                  list that is nested within the itemized list.

            \item This is the second item of the inner list.
                  \LaTeX\ allows you to nest lists deeper than
                  you really should.
         \end{enumerate}
         This is the rest of the second item of the outer
         list.  It is no more interesting than any other
         part of the item.
   \item This is the third item of the list.
\end{itemize}
You can even display poetry.
\begin{verse}
   There is an environment
    for verse \\             % The \\ command separates lines
   Whose features some poets % within a stanza.
   will curse.

                             % One or more blank lines separate stanzas.

   For instead of making\\
   Them do \emph{all} line breaking, \\
   It allows them to put too many words on a line when they'd rather be
   forced to be terse.
\end{verse}

Mathematical formulas may also be displayed.  A
displayed formula
is
one-line long; multiline
formulas require special formatting instructions.
   \[  \ip{\Gamma}{\psi'} = x'' + y^{2} + z_{i}^{n}\]
Don't start a paragraph with a displayed equation,
nor make one a paragraph by itself.

\end{document}               % End of document.

\section{Bench Huge Section 866}
% synthetic bench section 866
% This is a sample LaTeX input file.  (Version of 12 August 2004.)
%
% A '%' character causes TeX to ignore all remaining text on the line,
% and is used for comments like this one.

\documentclass{article}      % Specifies the document class

                             % The preamble begins here.
\title{An Example Document}  % Declares the document's title.
\author{Leslie Lamport}      % Declares the author's name.
\date{January 21, 1994}      % Deleting this command produces today's date.

\newcommand{\ip}[2]{(#1, #2)}
                             % Defines \ip{arg1}{arg2} to mean
                             % (arg1, arg2).

%\newcommand{\ip}[2]{\langle #1 | #2\rangle}
                             % This is an alternative definition of
                             % \ip that is commented out.

\begin{document}             % End of preamble and beginning of text.

\maketitle                   % Produces the title.

This is an example input file.  Comparing it with
the output it generates can show you how to
produce a simple document of your own.

\section{Ordinary Text}      % Produces section heading.  Lower-level
                             % sections are begun with similar
                             % \subsection and \subsubsection commands.

The ends  of words and sentences are marked
  by   spaces. It  doesn't matter how many
spaces    you type; one is as good as 100.  The
end of   a line counts as a space.

One   or more   blank lines denote the  end
of  a paragraph.

Since any number of consecutive spaces are treated
like a single one, the formatting of the input
file makes no difference to
      \LaTeX,                % The \LaTeX command generates the LaTeX logo.
but it makes a difference to you.  When you use
\LaTeX, making your input file as easy to read
as possible will be a great help as you write
your document and when you change it.  This sample
file shows how you can add comments to your own input
file.

Because printing is different from typewriting,
there are a number of things that you have to do
differently when preparing an input file than if
you were just typing the document directly.
Quotation marks like
       ``this''
have to be handled specially, as do quotes within
quotes:
       ``\,`this'            % \, separates the double and single quote.
        is what I just
        wrote, not  `that'\,''.

Dashes come in three sizes: an
       intra-word
dash, a medium dash for number ranges like
       1--2,
and a punctuation
       dash---like
this.

A sentence-ending space should be larger than the
space between words within a sentence.  You
sometimes have to type special commands in
conjunction with punctuation characters to get
this right, as in the following sentence.
       Gnats, gnus, etc.\ all  % `\ ' makes an inter-word space.
       begin with G\@.         % \@ marks end-of-sentence punctuation.
You should check the spaces after periods when
reading your output to make sure you haven't
forgotten any special cases.  Generating an
ellipsis
       \ldots\               % `\ ' is needed after `\ldots' because TeX
                             % ignores spaces after command names like \ldots
                             % made from \ + letters.
                             %
                             % Note how a `%' character causes TeX to ignore
                             % the end of the input line, so these blank lines
                             % do not start a new paragraph.
                             %
with the right spacing around the periods requires
a special command.

\LaTeX\ interprets some common characters as
commands, so you must type special commands to
generate them.  These characters include the
following:
       \$ \& \% \# \{ and \}.

In printing, text is usually emphasized with an
       \emph{italic}
type style.

\begin{em}
   A long segment of text can also be emphasized
   in this way.  Text within such a segment can be
   given \emph{additional} emphasis.
\end{em}

It is sometimes necessary to prevent \LaTeX\ from
breaking a line where it might otherwise do so.
This may be at a space, as between the ``Mr.''\ and
``Jones'' in
       ``Mr.~Jones'',        % ~ produces an unbreakable interword space.
or within a word---especially when the word is a
symbol like
       \mbox{\emph{itemnum}}
that makes little sense when hyphenated across
lines.

Footnotes\footnote{This is an example of a footnote.}
pose no problem.

\LaTeX\ is good at typesetting mathematical formulas
like
       \( x-3y + z = 7 \)
or
       \( a_{1} > x^{2n} + y^{2n} > x' \)
or
       \( \ip{A}{B} = \sum_{i} a_{i} b_{i} \).
The spaces you type in a formula are
ignored.  Remember that a letter like
       $x$                   % $ ... $  and  \( ... \)  are equivalent
is a formula when it denotes a mathematical
symbol, and it should be typed as one.

\section{Displayed Text}

Text is displayed by indenting it from the left
margin.  Quotations are commonly displayed.  There
are short quotations
\begin{quote}
   This is a short quotation.  It consists of a
   single paragraph of text.  See how it is formatted.
\end{quote}
and longer ones.
\begin{quotation}
   This is a longer quotation.  It consists of two
   paragraphs of text, neither of which are
   particularly interesting.

   This is the second paragraph of the quotation.  It
   is just as dull as the first paragraph.
\end{quotation}
Another frequently-displayed structure is a list.
The following is an example of an \emph{itemized}
list.
\begin{itemize}
   \item This is the first item of an itemized list.
         Each item in the list is marked with a ``tick''.
         You don't have to worry about what kind of tick
         mark is used.

   \item This is the second item of the list.  It
         contains another list nested inside it.  The inner
         list is an \emph{enumerated} list.
         \begin{enumerate}
            \item This is the first item of an enumerated
                  list that is nested within the itemized list.

            \item This is the second item of the inner list.
                  \LaTeX\ allows you to nest lists deeper than
                  you really should.
         \end{enumerate}
         This is the rest of the second item of the outer
         list.  It is no more interesting than any other
         part of the item.
   \item This is the third item of the list.
\end{itemize}
You can even display poetry.
\begin{verse}
   There is an environment
    for verse \\             % The \\ command separates lines
   Whose features some poets % within a stanza.
   will curse.

                             % One or more blank lines separate stanzas.

   For instead of making\\
   Them do \emph{all} line breaking, \\
   It allows them to put too many words on a line when they'd rather be
   forced to be terse.
\end{verse}

Mathematical formulas may also be displayed.  A
displayed formula
is
one-line long; multiline
formulas require special formatting instructions.
   \[  \ip{\Gamma}{\psi'} = x'' + y^{2} + z_{i}^{n}\]
Don't start a paragraph with a displayed equation,
nor make one a paragraph by itself.

\end{document}               % End of document.

\section{Bench Huge Section 867}
% synthetic bench section 867
% This is a sample LaTeX input file.  (Version of 12 August 2004.)
%
% A '%' character causes TeX to ignore all remaining text on the line,
% and is used for comments like this one.

\documentclass{article}      % Specifies the document class

                             % The preamble begins here.
\title{An Example Document}  % Declares the document's title.
\author{Leslie Lamport}      % Declares the author's name.
\date{January 21, 1994}      % Deleting this command produces today's date.

\newcommand{\ip}[2]{(#1, #2)}
                             % Defines \ip{arg1}{arg2} to mean
                             % (arg1, arg2).

%\newcommand{\ip}[2]{\langle #1 | #2\rangle}
                             % This is an alternative definition of
                             % \ip that is commented out.

\begin{document}             % End of preamble and beginning of text.

\maketitle                   % Produces the title.

This is an example input file.  Comparing it with
the output it generates can show you how to
produce a simple document of your own.

\section{Ordinary Text}      % Produces section heading.  Lower-level
                             % sections are begun with similar
                             % \subsection and \subsubsection commands.

The ends  of words and sentences are marked
  by   spaces. It  doesn't matter how many
spaces    you type; one is as good as 100.  The
end of   a line counts as a space.

One   or more   blank lines denote the  end
of  a paragraph.

Since any number of consecutive spaces are treated
like a single one, the formatting of the input
file makes no difference to
      \LaTeX,                % The \LaTeX command generates the LaTeX logo.
but it makes a difference to you.  When you use
\LaTeX, making your input file as easy to read
as possible will be a great help as you write
your document and when you change it.  This sample
file shows how you can add comments to your own input
file.

Because printing is different from typewriting,
there are a number of things that you have to do
differently when preparing an input file than if
you were just typing the document directly.
Quotation marks like
       ``this''
have to be handled specially, as do quotes within
quotes:
       ``\,`this'            % \, separates the double and single quote.
        is what I just
        wrote, not  `that'\,''.

Dashes come in three sizes: an
       intra-word
dash, a medium dash for number ranges like
       1--2,
and a punctuation
       dash---like
this.

A sentence-ending space should be larger than the
space between words within a sentence.  You
sometimes have to type special commands in
conjunction with punctuation characters to get
this right, as in the following sentence.
       Gnats, gnus, etc.\ all  % `\ ' makes an inter-word space.
       begin with G\@.         % \@ marks end-of-sentence punctuation.
You should check the spaces after periods when
reading your output to make sure you haven't
forgotten any special cases.  Generating an
ellipsis
       \ldots\               % `\ ' is needed after `\ldots' because TeX
                             % ignores spaces after command names like \ldots
                             % made from \ + letters.
                             %
                             % Note how a `%' character causes TeX to ignore
                             % the end of the input line, so these blank lines
                             % do not start a new paragraph.
                             %
with the right spacing around the periods requires
a special command.

\LaTeX\ interprets some common characters as
commands, so you must type special commands to
generate them.  These characters include the
following:
       \$ \& \% \# \{ and \}.

In printing, text is usually emphasized with an
       \emph{italic}
type style.

\begin{em}
   A long segment of text can also be emphasized
   in this way.  Text within such a segment can be
   given \emph{additional} emphasis.
\end{em}

It is sometimes necessary to prevent \LaTeX\ from
breaking a line where it might otherwise do so.
This may be at a space, as between the ``Mr.''\ and
``Jones'' in
       ``Mr.~Jones'',        % ~ produces an unbreakable interword space.
or within a word---especially when the word is a
symbol like
       \mbox{\emph{itemnum}}
that makes little sense when hyphenated across
lines.

Footnotes\footnote{This is an example of a footnote.}
pose no problem.

\LaTeX\ is good at typesetting mathematical formulas
like
       \( x-3y + z = 7 \)
or
       \( a_{1} > x^{2n} + y^{2n} > x' \)
or
       \( \ip{A}{B} = \sum_{i} a_{i} b_{i} \).
The spaces you type in a formula are
ignored.  Remember that a letter like
       $x$                   % $ ... $  and  \( ... \)  are equivalent
is a formula when it denotes a mathematical
symbol, and it should be typed as one.

\section{Displayed Text}

Text is displayed by indenting it from the left
margin.  Quotations are commonly displayed.  There
are short quotations
\begin{quote}
   This is a short quotation.  It consists of a
   single paragraph of text.  See how it is formatted.
\end{quote}
and longer ones.
\begin{quotation}
   This is a longer quotation.  It consists of two
   paragraphs of text, neither of which are
   particularly interesting.

   This is the second paragraph of the quotation.  It
   is just as dull as the first paragraph.
\end{quotation}
Another frequently-displayed structure is a list.
The following is an example of an \emph{itemized}
list.
\begin{itemize}
   \item This is the first item of an itemized list.
         Each item in the list is marked with a ``tick''.
         You don't have to worry about what kind of tick
         mark is used.

   \item This is the second item of the list.  It
         contains another list nested inside it.  The inner
         list is an \emph{enumerated} list.
         \begin{enumerate}
            \item This is the first item of an enumerated
                  list that is nested within the itemized list.

            \item This is the second item of the inner list.
                  \LaTeX\ allows you to nest lists deeper than
                  you really should.
         \end{enumerate}
         This is the rest of the second item of the outer
         list.  It is no more interesting than any other
         part of the item.
   \item This is the third item of the list.
\end{itemize}
You can even display poetry.
\begin{verse}
   There is an environment
    for verse \\             % The \\ command separates lines
   Whose features some poets % within a stanza.
   will curse.

                             % One or more blank lines separate stanzas.

   For instead of making\\
   Them do \emph{all} line breaking, \\
   It allows them to put too many words on a line when they'd rather be
   forced to be terse.
\end{verse}

Mathematical formulas may also be displayed.  A
displayed formula
is
one-line long; multiline
formulas require special formatting instructions.
   \[  \ip{\Gamma}{\psi'} = x'' + y^{2} + z_{i}^{n}\]
Don't start a paragraph with a displayed equation,
nor make one a paragraph by itself.

\end{document}               % End of document.

\section{Bench Huge Section 868}
% synthetic bench section 868
% This is a sample LaTeX input file.  (Version of 12 August 2004.)
%
% A '%' character causes TeX to ignore all remaining text on the line,
% and is used for comments like this one.

\documentclass{article}      % Specifies the document class

                             % The preamble begins here.
\title{An Example Document}  % Declares the document's title.
\author{Leslie Lamport}      % Declares the author's name.
\date{January 21, 1994}      % Deleting this command produces today's date.

\newcommand{\ip}[2]{(#1, #2)}
                             % Defines \ip{arg1}{arg2} to mean
                             % (arg1, arg2).

%\newcommand{\ip}[2]{\langle #1 | #2\rangle}
                             % This is an alternative definition of
                             % \ip that is commented out.

\begin{document}             % End of preamble and beginning of text.

\maketitle                   % Produces the title.

This is an example input file.  Comparing it with
the output it generates can show you how to
produce a simple document of your own.

\section{Ordinary Text}      % Produces section heading.  Lower-level
                             % sections are begun with similar
                             % \subsection and \subsubsection commands.

The ends  of words and sentences are marked
  by   spaces. It  doesn't matter how many
spaces    you type; one is as good as 100.  The
end of   a line counts as a space.

One   or more   blank lines denote the  end
of  a paragraph.

Since any number of consecutive spaces are treated
like a single one, the formatting of the input
file makes no difference to
      \LaTeX,                % The \LaTeX command generates the LaTeX logo.
but it makes a difference to you.  When you use
\LaTeX, making your input file as easy to read
as possible will be a great help as you write
your document and when you change it.  This sample
file shows how you can add comments to your own input
file.

Because printing is different from typewriting,
there are a number of things that you have to do
differently when preparing an input file than if
you were just typing the document directly.
Quotation marks like
       ``this''
have to be handled specially, as do quotes within
quotes:
       ``\,`this'            % \, separates the double and single quote.
        is what I just
        wrote, not  `that'\,''.

Dashes come in three sizes: an
       intra-word
dash, a medium dash for number ranges like
       1--2,
and a punctuation
       dash---like
this.

A sentence-ending space should be larger than the
space between words within a sentence.  You
sometimes have to type special commands in
conjunction with punctuation characters to get
this right, as in the following sentence.
       Gnats, gnus, etc.\ all  % `\ ' makes an inter-word space.
       begin with G\@.         % \@ marks end-of-sentence punctuation.
You should check the spaces after periods when
reading your output to make sure you haven't
forgotten any special cases.  Generating an
ellipsis
       \ldots\               % `\ ' is needed after `\ldots' because TeX
                             % ignores spaces after command names like \ldots
                             % made from \ + letters.
                             %
                             % Note how a `%' character causes TeX to ignore
                             % the end of the input line, so these blank lines
                             % do not start a new paragraph.
                             %
with the right spacing around the periods requires
a special command.

\LaTeX\ interprets some common characters as
commands, so you must type special commands to
generate them.  These characters include the
following:
       \$ \& \% \# \{ and \}.

In printing, text is usually emphasized with an
       \emph{italic}
type style.

\begin{em}
   A long segment of text can also be emphasized
   in this way.  Text within such a segment can be
   given \emph{additional} emphasis.
\end{em}

It is sometimes necessary to prevent \LaTeX\ from
breaking a line where it might otherwise do so.
This may be at a space, as between the ``Mr.''\ and
``Jones'' in
       ``Mr.~Jones'',        % ~ produces an unbreakable interword space.
or within a word---especially when the word is a
symbol like
       \mbox{\emph{itemnum}}
that makes little sense when hyphenated across
lines.

Footnotes\footnote{This is an example of a footnote.}
pose no problem.

\LaTeX\ is good at typesetting mathematical formulas
like
       \( x-3y + z = 7 \)
or
       \( a_{1} > x^{2n} + y^{2n} > x' \)
or
       \( \ip{A}{B} = \sum_{i} a_{i} b_{i} \).
The spaces you type in a formula are
ignored.  Remember that a letter like
       $x$                   % $ ... $  and  \( ... \)  are equivalent
is a formula when it denotes a mathematical
symbol, and it should be typed as one.

\section{Displayed Text}

Text is displayed by indenting it from the left
margin.  Quotations are commonly displayed.  There
are short quotations
\begin{quote}
   This is a short quotation.  It consists of a
   single paragraph of text.  See how it is formatted.
\end{quote}
and longer ones.
\begin{quotation}
   This is a longer quotation.  It consists of two
   paragraphs of text, neither of which are
   particularly interesting.

   This is the second paragraph of the quotation.  It
   is just as dull as the first paragraph.
\end{quotation}
Another frequently-displayed structure is a list.
The following is an example of an \emph{itemized}
list.
\begin{itemize}
   \item This is the first item of an itemized list.
         Each item in the list is marked with a ``tick''.
         You don't have to worry about what kind of tick
         mark is used.

   \item This is the second item of the list.  It
         contains another list nested inside it.  The inner
         list is an \emph{enumerated} list.
         \begin{enumerate}
            \item This is the first item of an enumerated
                  list that is nested within the itemized list.

            \item This is the second item of the inner list.
                  \LaTeX\ allows you to nest lists deeper than
                  you really should.
         \end{enumerate}
         This is the rest of the second item of the outer
         list.  It is no more interesting than any other
         part of the item.
   \item This is the third item of the list.
\end{itemize}
You can even display poetry.
\begin{verse}
   There is an environment
    for verse \\             % The \\ command separates lines
   Whose features some poets % within a stanza.
   will curse.

                             % One or more blank lines separate stanzas.

   For instead of making\\
   Them do \emph{all} line breaking, \\
   It allows them to put too many words on a line when they'd rather be
   forced to be terse.
\end{verse}

Mathematical formulas may also be displayed.  A
displayed formula
is
one-line long; multiline
formulas require special formatting instructions.
   \[  \ip{\Gamma}{\psi'} = x'' + y^{2} + z_{i}^{n}\]
Don't start a paragraph with a displayed equation,
nor make one a paragraph by itself.

\end{document}               % End of document.

\section{Bench Huge Section 869}
% synthetic bench section 869
% This is a sample LaTeX input file.  (Version of 12 August 2004.)
%
% A '%' character causes TeX to ignore all remaining text on the line,
% and is used for comments like this one.

\documentclass{article}      % Specifies the document class

                             % The preamble begins here.
\title{An Example Document}  % Declares the document's title.
\author{Leslie Lamport}      % Declares the author's name.
\date{January 21, 1994}      % Deleting this command produces today's date.

\newcommand{\ip}[2]{(#1, #2)}
                             % Defines \ip{arg1}{arg2} to mean
                             % (arg1, arg2).

%\newcommand{\ip}[2]{\langle #1 | #2\rangle}
                             % This is an alternative definition of
                             % \ip that is commented out.

\begin{document}             % End of preamble and beginning of text.

\maketitle                   % Produces the title.

This is an example input file.  Comparing it with
the output it generates can show you how to
produce a simple document of your own.

\section{Ordinary Text}      % Produces section heading.  Lower-level
                             % sections are begun with similar
                             % \subsection and \subsubsection commands.

The ends  of words and sentences are marked
  by   spaces. It  doesn't matter how many
spaces    you type; one is as good as 100.  The
end of   a line counts as a space.

One   or more   blank lines denote the  end
of  a paragraph.

Since any number of consecutive spaces are treated
like a single one, the formatting of the input
file makes no difference to
      \LaTeX,                % The \LaTeX command generates the LaTeX logo.
but it makes a difference to you.  When you use
\LaTeX, making your input file as easy to read
as possible will be a great help as you write
your document and when you change it.  This sample
file shows how you can add comments to your own input
file.

Because printing is different from typewriting,
there are a number of things that you have to do
differently when preparing an input file than if
you were just typing the document directly.
Quotation marks like
       ``this''
have to be handled specially, as do quotes within
quotes:
       ``\,`this'            % \, separates the double and single quote.
        is what I just
        wrote, not  `that'\,''.

Dashes come in three sizes: an
       intra-word
dash, a medium dash for number ranges like
       1--2,
and a punctuation
       dash---like
this.

A sentence-ending space should be larger than the
space between words within a sentence.  You
sometimes have to type special commands in
conjunction with punctuation characters to get
this right, as in the following sentence.
       Gnats, gnus, etc.\ all  % `\ ' makes an inter-word space.
       begin with G\@.         % \@ marks end-of-sentence punctuation.
You should check the spaces after periods when
reading your output to make sure you haven't
forgotten any special cases.  Generating an
ellipsis
       \ldots\               % `\ ' is needed after `\ldots' because TeX
                             % ignores spaces after command names like \ldots
                             % made from \ + letters.
                             %
                             % Note how a `%' character causes TeX to ignore
                             % the end of the input line, so these blank lines
                             % do not start a new paragraph.
                             %
with the right spacing around the periods requires
a special command.

\LaTeX\ interprets some common characters as
commands, so you must type special commands to
generate them.  These characters include the
following:
       \$ \& \% \# \{ and \}.

In printing, text is usually emphasized with an
       \emph{italic}
type style.

\begin{em}
   A long segment of text can also be emphasized
   in this way.  Text within such a segment can be
   given \emph{additional} emphasis.
\end{em}

It is sometimes necessary to prevent \LaTeX\ from
breaking a line where it might otherwise do so.
This may be at a space, as between the ``Mr.''\ and
``Jones'' in
       ``Mr.~Jones'',        % ~ produces an unbreakable interword space.
or within a word---especially when the word is a
symbol like
       \mbox{\emph{itemnum}}
that makes little sense when hyphenated across
lines.

Footnotes\footnote{This is an example of a footnote.}
pose no problem.

\LaTeX\ is good at typesetting mathematical formulas
like
       \( x-3y + z = 7 \)
or
       \( a_{1} > x^{2n} + y^{2n} > x' \)
or
       \( \ip{A}{B} = \sum_{i} a_{i} b_{i} \).
The spaces you type in a formula are
ignored.  Remember that a letter like
       $x$                   % $ ... $  and  \( ... \)  are equivalent
is a formula when it denotes a mathematical
symbol, and it should be typed as one.

\section{Displayed Text}

Text is displayed by indenting it from the left
margin.  Quotations are commonly displayed.  There
are short quotations
\begin{quote}
   This is a short quotation.  It consists of a
   single paragraph of text.  See how it is formatted.
\end{quote}
and longer ones.
\begin{quotation}
   This is a longer quotation.  It consists of two
   paragraphs of text, neither of which are
   particularly interesting.

   This is the second paragraph of the quotation.  It
   is just as dull as the first paragraph.
\end{quotation}
Another frequently-displayed structure is a list.
The following is an example of an \emph{itemized}
list.
\begin{itemize}
   \item This is the first item of an itemized list.
         Each item in the list is marked with a ``tick''.
         You don't have to worry about what kind of tick
         mark is used.

   \item This is the second item of the list.  It
         contains another list nested inside it.  The inner
         list is an \emph{enumerated} list.
         \begin{enumerate}
            \item This is the first item of an enumerated
                  list that is nested within the itemized list.

            \item This is the second item of the inner list.
                  \LaTeX\ allows you to nest lists deeper than
                  you really should.
         \end{enumerate}
         This is the rest of the second item of the outer
         list.  It is no more interesting than any other
         part of the item.
   \item This is the third item of the list.
\end{itemize}
You can even display poetry.
\begin{verse}
   There is an environment
    for verse \\             % The \\ command separates lines
   Whose features some poets % within a stanza.
   will curse.

                             % One or more blank lines separate stanzas.

   For instead of making\\
   Them do \emph{all} line breaking, \\
   It allows them to put too many words on a line when they'd rather be
   forced to be terse.
\end{verse}

Mathematical formulas may also be displayed.  A
displayed formula
is
one-line long; multiline
formulas require special formatting instructions.
   \[  \ip{\Gamma}{\psi'} = x'' + y^{2} + z_{i}^{n}\]
Don't start a paragraph with a displayed equation,
nor make one a paragraph by itself.

\end{document}               % End of document.

\section{Bench Huge Section 870}
% synthetic bench section 870
% This is a sample LaTeX input file.  (Version of 12 August 2004.)
%
% A '%' character causes TeX to ignore all remaining text on the line,
% and is used for comments like this one.

\documentclass{article}      % Specifies the document class

                             % The preamble begins here.
\title{An Example Document}  % Declares the document's title.
\author{Leslie Lamport}      % Declares the author's name.
\date{January 21, 1994}      % Deleting this command produces today's date.

\newcommand{\ip}[2]{(#1, #2)}
                             % Defines \ip{arg1}{arg2} to mean
                             % (arg1, arg2).

%\newcommand{\ip}[2]{\langle #1 | #2\rangle}
                             % This is an alternative definition of
                             % \ip that is commented out.

\begin{document}             % End of preamble and beginning of text.

\maketitle                   % Produces the title.

This is an example input file.  Comparing it with
the output it generates can show you how to
produce a simple document of your own.

\section{Ordinary Text}      % Produces section heading.  Lower-level
                             % sections are begun with similar
                             % \subsection and \subsubsection commands.

The ends  of words and sentences are marked
  by   spaces. It  doesn't matter how many
spaces    you type; one is as good as 100.  The
end of   a line counts as a space.

One   or more   blank lines denote the  end
of  a paragraph.

Since any number of consecutive spaces are treated
like a single one, the formatting of the input
file makes no difference to
      \LaTeX,                % The \LaTeX command generates the LaTeX logo.
but it makes a difference to you.  When you use
\LaTeX, making your input file as easy to read
as possible will be a great help as you write
your document and when you change it.  This sample
file shows how you can add comments to your own input
file.

Because printing is different from typewriting,
there are a number of things that you have to do
differently when preparing an input file than if
you were just typing the document directly.
Quotation marks like
       ``this''
have to be handled specially, as do quotes within
quotes:
       ``\,`this'            % \, separates the double and single quote.
        is what I just
        wrote, not  `that'\,''.

Dashes come in three sizes: an
       intra-word
dash, a medium dash for number ranges like
       1--2,
and a punctuation
       dash---like
this.

A sentence-ending space should be larger than the
space between words within a sentence.  You
sometimes have to type special commands in
conjunction with punctuation characters to get
this right, as in the following sentence.
       Gnats, gnus, etc.\ all  % `\ ' makes an inter-word space.
       begin with G\@.         % \@ marks end-of-sentence punctuation.
You should check the spaces after periods when
reading your output to make sure you haven't
forgotten any special cases.  Generating an
ellipsis
       \ldots\               % `\ ' is needed after `\ldots' because TeX
                             % ignores spaces after command names like \ldots
                             % made from \ + letters.
                             %
                             % Note how a `%' character causes TeX to ignore
                             % the end of the input line, so these blank lines
                             % do not start a new paragraph.
                             %
with the right spacing around the periods requires
a special command.

\LaTeX\ interprets some common characters as
commands, so you must type special commands to
generate them.  These characters include the
following:
       \$ \& \% \# \{ and \}.

In printing, text is usually emphasized with an
       \emph{italic}
type style.

\begin{em}
   A long segment of text can also be emphasized
   in this way.  Text within such a segment can be
   given \emph{additional} emphasis.
\end{em}

It is sometimes necessary to prevent \LaTeX\ from
breaking a line where it might otherwise do so.
This may be at a space, as between the ``Mr.''\ and
``Jones'' in
       ``Mr.~Jones'',        % ~ produces an unbreakable interword space.
or within a word---especially when the word is a
symbol like
       \mbox{\emph{itemnum}}
that makes little sense when hyphenated across
lines.

Footnotes\footnote{This is an example of a footnote.}
pose no problem.

\LaTeX\ is good at typesetting mathematical formulas
like
       \( x-3y + z = 7 \)
or
       \( a_{1} > x^{2n} + y^{2n} > x' \)
or
       \( \ip{A}{B} = \sum_{i} a_{i} b_{i} \).
The spaces you type in a formula are
ignored.  Remember that a letter like
       $x$                   % $ ... $  and  \( ... \)  are equivalent
is a formula when it denotes a mathematical
symbol, and it should be typed as one.

\section{Displayed Text}

Text is displayed by indenting it from the left
margin.  Quotations are commonly displayed.  There
are short quotations
\begin{quote}
   This is a short quotation.  It consists of a
   single paragraph of text.  See how it is formatted.
\end{quote}
and longer ones.
\begin{quotation}
   This is a longer quotation.  It consists of two
   paragraphs of text, neither of which are
   particularly interesting.

   This is the second paragraph of the quotation.  It
   is just as dull as the first paragraph.
\end{quotation}
Another frequently-displayed structure is a list.
The following is an example of an \emph{itemized}
list.
\begin{itemize}
   \item This is the first item of an itemized list.
         Each item in the list is marked with a ``tick''.
         You don't have to worry about what kind of tick
         mark is used.

   \item This is the second item of the list.  It
         contains another list nested inside it.  The inner
         list is an \emph{enumerated} list.
         \begin{enumerate}
            \item This is the first item of an enumerated
                  list that is nested within the itemized list.

            \item This is the second item of the inner list.
                  \LaTeX\ allows you to nest lists deeper than
                  you really should.
         \end{enumerate}
         This is the rest of the second item of the outer
         list.  It is no more interesting than any other
         part of the item.
   \item This is the third item of the list.
\end{itemize}
You can even display poetry.
\begin{verse}
   There is an environment
    for verse \\             % The \\ command separates lines
   Whose features some poets % within a stanza.
   will curse.

                             % One or more blank lines separate stanzas.

   For instead of making\\
   Them do \emph{all} line breaking, \\
   It allows them to put too many words on a line when they'd rather be
   forced to be terse.
\end{verse}

Mathematical formulas may also be displayed.  A
displayed formula
is
one-line long; multiline
formulas require special formatting instructions.
   \[  \ip{\Gamma}{\psi'} = x'' + y^{2} + z_{i}^{n}\]
Don't start a paragraph with a displayed equation,
nor make one a paragraph by itself.

\end{document}               % End of document.

\section{Bench Huge Section 871}
% synthetic bench section 871
% This is a sample LaTeX input file.  (Version of 12 August 2004.)
%
% A '%' character causes TeX to ignore all remaining text on the line,
% and is used for comments like this one.

\documentclass{article}      % Specifies the document class

                             % The preamble begins here.
\title{An Example Document}  % Declares the document's title.
\author{Leslie Lamport}      % Declares the author's name.
\date{January 21, 1994}      % Deleting this command produces today's date.

\newcommand{\ip}[2]{(#1, #2)}
                             % Defines \ip{arg1}{arg2} to mean
                             % (arg1, arg2).

%\newcommand{\ip}[2]{\langle #1 | #2\rangle}
                             % This is an alternative definition of
                             % \ip that is commented out.

\begin{document}             % End of preamble and beginning of text.

\maketitle                   % Produces the title.

This is an example input file.  Comparing it with
the output it generates can show you how to
produce a simple document of your own.

\section{Ordinary Text}      % Produces section heading.  Lower-level
                             % sections are begun with similar
                             % \subsection and \subsubsection commands.

The ends  of words and sentences are marked
  by   spaces. It  doesn't matter how many
spaces    you type; one is as good as 100.  The
end of   a line counts as a space.

One   or more   blank lines denote the  end
of  a paragraph.

Since any number of consecutive spaces are treated
like a single one, the formatting of the input
file makes no difference to
      \LaTeX,                % The \LaTeX command generates the LaTeX logo.
but it makes a difference to you.  When you use
\LaTeX, making your input file as easy to read
as possible will be a great help as you write
your document and when you change it.  This sample
file shows how you can add comments to your own input
file.

Because printing is different from typewriting,
there are a number of things that you have to do
differently when preparing an input file than if
you were just typing the document directly.
Quotation marks like
       ``this''
have to be handled specially, as do quotes within
quotes:
       ``\,`this'            % \, separates the double and single quote.
        is what I just
        wrote, not  `that'\,''.

Dashes come in three sizes: an
       intra-word
dash, a medium dash for number ranges like
       1--2,
and a punctuation
       dash---like
this.

A sentence-ending space should be larger than the
space between words within a sentence.  You
sometimes have to type special commands in
conjunction with punctuation characters to get
this right, as in the following sentence.
       Gnats, gnus, etc.\ all  % `\ ' makes an inter-word space.
       begin with G\@.         % \@ marks end-of-sentence punctuation.
You should check the spaces after periods when
reading your output to make sure you haven't
forgotten any special cases.  Generating an
ellipsis
       \ldots\               % `\ ' is needed after `\ldots' because TeX
                             % ignores spaces after command names like \ldots
                             % made from \ + letters.
                             %
                             % Note how a `%' character causes TeX to ignore
                             % the end of the input line, so these blank lines
                             % do not start a new paragraph.
                             %
with the right spacing around the periods requires
a special command.

\LaTeX\ interprets some common characters as
commands, so you must type special commands to
generate them.  These characters include the
following:
       \$ \& \% \# \{ and \}.

In printing, text is usually emphasized with an
       \emph{italic}
type style.

\begin{em}
   A long segment of text can also be emphasized
   in this way.  Text within such a segment can be
   given \emph{additional} emphasis.
\end{em}

It is sometimes necessary to prevent \LaTeX\ from
breaking a line where it might otherwise do so.
This may be at a space, as between the ``Mr.''\ and
``Jones'' in
       ``Mr.~Jones'',        % ~ produces an unbreakable interword space.
or within a word---especially when the word is a
symbol like
       \mbox{\emph{itemnum}}
that makes little sense when hyphenated across
lines.

Footnotes\footnote{This is an example of a footnote.}
pose no problem.

\LaTeX\ is good at typesetting mathematical formulas
like
       \( x-3y + z = 7 \)
or
       \( a_{1} > x^{2n} + y^{2n} > x' \)
or
       \( \ip{A}{B} = \sum_{i} a_{i} b_{i} \).
The spaces you type in a formula are
ignored.  Remember that a letter like
       $x$                   % $ ... $  and  \( ... \)  are equivalent
is a formula when it denotes a mathematical
symbol, and it should be typed as one.

\section{Displayed Text}

Text is displayed by indenting it from the left
margin.  Quotations are commonly displayed.  There
are short quotations
\begin{quote}
   This is a short quotation.  It consists of a
   single paragraph of text.  See how it is formatted.
\end{quote}
and longer ones.
\begin{quotation}
   This is a longer quotation.  It consists of two
   paragraphs of text, neither of which are
   particularly interesting.

   This is the second paragraph of the quotation.  It
   is just as dull as the first paragraph.
\end{quotation}
Another frequently-displayed structure is a list.
The following is an example of an \emph{itemized}
list.
\begin{itemize}
   \item This is the first item of an itemized list.
         Each item in the list is marked with a ``tick''.
         You don't have to worry about what kind of tick
         mark is used.

   \item This is the second item of the list.  It
         contains another list nested inside it.  The inner
         list is an \emph{enumerated} list.
         \begin{enumerate}
            \item This is the first item of an enumerated
                  list that is nested within the itemized list.

            \item This is the second item of the inner list.
                  \LaTeX\ allows you to nest lists deeper than
                  you really should.
         \end{enumerate}
         This is the rest of the second item of the outer
         list.  It is no more interesting than any other
         part of the item.
   \item This is the third item of the list.
\end{itemize}
You can even display poetry.
\begin{verse}
   There is an environment
    for verse \\             % The \\ command separates lines
   Whose features some poets % within a stanza.
   will curse.

                             % One or more blank lines separate stanzas.

   For instead of making\\
   Them do \emph{all} line breaking, \\
   It allows them to put too many words on a line when they'd rather be
   forced to be terse.
\end{verse}

Mathematical formulas may also be displayed.  A
displayed formula
is
one-line long; multiline
formulas require special formatting instructions.
   \[  \ip{\Gamma}{\psi'} = x'' + y^{2} + z_{i}^{n}\]
Don't start a paragraph with a displayed equation,
nor make one a paragraph by itself.

\end{document}               % End of document.

\section{Bench Huge Section 872}
% synthetic bench section 872
% This is a sample LaTeX input file.  (Version of 12 August 2004.)
%
% A '%' character causes TeX to ignore all remaining text on the line,
% and is used for comments like this one.

\documentclass{article}      % Specifies the document class

                             % The preamble begins here.
\title{An Example Document}  % Declares the document's title.
\author{Leslie Lamport}      % Declares the author's name.
\date{January 21, 1994}      % Deleting this command produces today's date.

\newcommand{\ip}[2]{(#1, #2)}
                             % Defines \ip{arg1}{arg2} to mean
                             % (arg1, arg2).

%\newcommand{\ip}[2]{\langle #1 | #2\rangle}
                             % This is an alternative definition of
                             % \ip that is commented out.

\begin{document}             % End of preamble and beginning of text.

\maketitle                   % Produces the title.

This is an example input file.  Comparing it with
the output it generates can show you how to
produce a simple document of your own.

\section{Ordinary Text}      % Produces section heading.  Lower-level
                             % sections are begun with similar
                             % \subsection and \subsubsection commands.

The ends  of words and sentences are marked
  by   spaces. It  doesn't matter how many
spaces    you type; one is as good as 100.  The
end of   a line counts as a space.

One   or more   blank lines denote the  end
of  a paragraph.

Since any number of consecutive spaces are treated
like a single one, the formatting of the input
file makes no difference to
      \LaTeX,                % The \LaTeX command generates the LaTeX logo.
but it makes a difference to you.  When you use
\LaTeX, making your input file as easy to read
as possible will be a great help as you write
your document and when you change it.  This sample
file shows how you can add comments to your own input
file.

Because printing is different from typewriting,
there are a number of things that you have to do
differently when preparing an input file than if
you were just typing the document directly.
Quotation marks like
       ``this''
have to be handled specially, as do quotes within
quotes:
       ``\,`this'            % \, separates the double and single quote.
        is what I just
        wrote, not  `that'\,''.

Dashes come in three sizes: an
       intra-word
dash, a medium dash for number ranges like
       1--2,
and a punctuation
       dash---like
this.

A sentence-ending space should be larger than the
space between words within a sentence.  You
sometimes have to type special commands in
conjunction with punctuation characters to get
this right, as in the following sentence.
       Gnats, gnus, etc.\ all  % `\ ' makes an inter-word space.
       begin with G\@.         % \@ marks end-of-sentence punctuation.
You should check the spaces after periods when
reading your output to make sure you haven't
forgotten any special cases.  Generating an
ellipsis
       \ldots\               % `\ ' is needed after `\ldots' because TeX
                             % ignores spaces after command names like \ldots
                             % made from \ + letters.
                             %
                             % Note how a `%' character causes TeX to ignore
                             % the end of the input line, so these blank lines
                             % do not start a new paragraph.
                             %
with the right spacing around the periods requires
a special command.

\LaTeX\ interprets some common characters as
commands, so you must type special commands to
generate them.  These characters include the
following:
       \$ \& \% \# \{ and \}.

In printing, text is usually emphasized with an
       \emph{italic}
type style.

\begin{em}
   A long segment of text can also be emphasized
   in this way.  Text within such a segment can be
   given \emph{additional} emphasis.
\end{em}

It is sometimes necessary to prevent \LaTeX\ from
breaking a line where it might otherwise do so.
This may be at a space, as between the ``Mr.''\ and
``Jones'' in
       ``Mr.~Jones'',        % ~ produces an unbreakable interword space.
or within a word---especially when the word is a
symbol like
       \mbox{\emph{itemnum}}
that makes little sense when hyphenated across
lines.

Footnotes\footnote{This is an example of a footnote.}
pose no problem.

\LaTeX\ is good at typesetting mathematical formulas
like
       \( x-3y + z = 7 \)
or
       \( a_{1} > x^{2n} + y^{2n} > x' \)
or
       \( \ip{A}{B} = \sum_{i} a_{i} b_{i} \).
The spaces you type in a formula are
ignored.  Remember that a letter like
       $x$                   % $ ... $  and  \( ... \)  are equivalent
is a formula when it denotes a mathematical
symbol, and it should be typed as one.

\section{Displayed Text}

Text is displayed by indenting it from the left
margin.  Quotations are commonly displayed.  There
are short quotations
\begin{quote}
   This is a short quotation.  It consists of a
   single paragraph of text.  See how it is formatted.
\end{quote}
and longer ones.
\begin{quotation}
   This is a longer quotation.  It consists of two
   paragraphs of text, neither of which are
   particularly interesting.

   This is the second paragraph of the quotation.  It
   is just as dull as the first paragraph.
\end{quotation}
Another frequently-displayed structure is a list.
The following is an example of an \emph{itemized}
list.
\begin{itemize}
   \item This is the first item of an itemized list.
         Each item in the list is marked with a ``tick''.
         You don't have to worry about what kind of tick
         mark is used.

   \item This is the second item of the list.  It
         contains another list nested inside it.  The inner
         list is an \emph{enumerated} list.
         \begin{enumerate}
            \item This is the first item of an enumerated
                  list that is nested within the itemized list.

            \item This is the second item of the inner list.
                  \LaTeX\ allows you to nest lists deeper than
                  you really should.
         \end{enumerate}
         This is the rest of the second item of the outer
         list.  It is no more interesting than any other
         part of the item.
   \item This is the third item of the list.
\end{itemize}
You can even display poetry.
\begin{verse}
   There is an environment
    for verse \\             % The \\ command separates lines
   Whose features some poets % within a stanza.
   will curse.

                             % One or more blank lines separate stanzas.

   For instead of making\\
   Them do \emph{all} line breaking, \\
   It allows them to put too many words on a line when they'd rather be
   forced to be terse.
\end{verse}

Mathematical formulas may also be displayed.  A
displayed formula
is
one-line long; multiline
formulas require special formatting instructions.
   \[  \ip{\Gamma}{\psi'} = x'' + y^{2} + z_{i}^{n}\]
Don't start a paragraph with a displayed equation,
nor make one a paragraph by itself.

\end{document}               % End of document.

\section{Bench Huge Section 873}
% synthetic bench section 873
% This is a sample LaTeX input file.  (Version of 12 August 2004.)
%
% A '%' character causes TeX to ignore all remaining text on the line,
% and is used for comments like this one.

\documentclass{article}      % Specifies the document class

                             % The preamble begins here.
\title{An Example Document}  % Declares the document's title.
\author{Leslie Lamport}      % Declares the author's name.
\date{January 21, 1994}      % Deleting this command produces today's date.

\newcommand{\ip}[2]{(#1, #2)}
                             % Defines \ip{arg1}{arg2} to mean
                             % (arg1, arg2).

%\newcommand{\ip}[2]{\langle #1 | #2\rangle}
                             % This is an alternative definition of
                             % \ip that is commented out.

\begin{document}             % End of preamble and beginning of text.

\maketitle                   % Produces the title.

This is an example input file.  Comparing it with
the output it generates can show you how to
produce a simple document of your own.

\section{Ordinary Text}      % Produces section heading.  Lower-level
                             % sections are begun with similar
                             % \subsection and \subsubsection commands.

The ends  of words and sentences are marked
  by   spaces. It  doesn't matter how many
spaces    you type; one is as good as 100.  The
end of   a line counts as a space.

One   or more   blank lines denote the  end
of  a paragraph.

Since any number of consecutive spaces are treated
like a single one, the formatting of the input
file makes no difference to
      \LaTeX,                % The \LaTeX command generates the LaTeX logo.
but it makes a difference to you.  When you use
\LaTeX, making your input file as easy to read
as possible will be a great help as you write
your document and when you change it.  This sample
file shows how you can add comments to your own input
file.

Because printing is different from typewriting,
there are a number of things that you have to do
differently when preparing an input file than if
you were just typing the document directly.
Quotation marks like
       ``this''
have to be handled specially, as do quotes within
quotes:
       ``\,`this'            % \, separates the double and single quote.
        is what I just
        wrote, not  `that'\,''.

Dashes come in three sizes: an
       intra-word
dash, a medium dash for number ranges like
       1--2,
and a punctuation
       dash---like
this.

A sentence-ending space should be larger than the
space between words within a sentence.  You
sometimes have to type special commands in
conjunction with punctuation characters to get
this right, as in the following sentence.
       Gnats, gnus, etc.\ all  % `\ ' makes an inter-word space.
       begin with G\@.         % \@ marks end-of-sentence punctuation.
You should check the spaces after periods when
reading your output to make sure you haven't
forgotten any special cases.  Generating an
ellipsis
       \ldots\               % `\ ' is needed after `\ldots' because TeX
                             % ignores spaces after command names like \ldots
                             % made from \ + letters.
                             %
                             % Note how a `%' character causes TeX to ignore
                             % the end of the input line, so these blank lines
                             % do not start a new paragraph.
                             %
with the right spacing around the periods requires
a special command.

\LaTeX\ interprets some common characters as
commands, so you must type special commands to
generate them.  These characters include the
following:
       \$ \& \% \# \{ and \}.

In printing, text is usually emphasized with an
       \emph{italic}
type style.

\begin{em}
   A long segment of text can also be emphasized
   in this way.  Text within such a segment can be
   given \emph{additional} emphasis.
\end{em}

It is sometimes necessary to prevent \LaTeX\ from
breaking a line where it might otherwise do so.
This may be at a space, as between the ``Mr.''\ and
``Jones'' in
       ``Mr.~Jones'',        % ~ produces an unbreakable interword space.
or within a word---especially when the word is a
symbol like
       \mbox{\emph{itemnum}}
that makes little sense when hyphenated across
lines.

Footnotes\footnote{This is an example of a footnote.}
pose no problem.

\LaTeX\ is good at typesetting mathematical formulas
like
       \( x-3y + z = 7 \)
or
       \( a_{1} > x^{2n} + y^{2n} > x' \)
or
       \( \ip{A}{B} = \sum_{i} a_{i} b_{i} \).
The spaces you type in a formula are
ignored.  Remember that a letter like
       $x$                   % $ ... $  and  \( ... \)  are equivalent
is a formula when it denotes a mathematical
symbol, and it should be typed as one.

\section{Displayed Text}

Text is displayed by indenting it from the left
margin.  Quotations are commonly displayed.  There
are short quotations
\begin{quote}
   This is a short quotation.  It consists of a
   single paragraph of text.  See how it is formatted.
\end{quote}
and longer ones.
\begin{quotation}
   This is a longer quotation.  It consists of two
   paragraphs of text, neither of which are
   particularly interesting.

   This is the second paragraph of the quotation.  It
   is just as dull as the first paragraph.
\end{quotation}
Another frequently-displayed structure is a list.
The following is an example of an \emph{itemized}
list.
\begin{itemize}
   \item This is the first item of an itemized list.
         Each item in the list is marked with a ``tick''.
         You don't have to worry about what kind of tick
         mark is used.

   \item This is the second item of the list.  It
         contains another list nested inside it.  The inner
         list is an \emph{enumerated} list.
         \begin{enumerate}
            \item This is the first item of an enumerated
                  list that is nested within the itemized list.

            \item This is the second item of the inner list.
                  \LaTeX\ allows you to nest lists deeper than
                  you really should.
         \end{enumerate}
         This is the rest of the second item of the outer
         list.  It is no more interesting than any other
         part of the item.
   \item This is the third item of the list.
\end{itemize}
You can even display poetry.
\begin{verse}
   There is an environment
    for verse \\             % The \\ command separates lines
   Whose features some poets % within a stanza.
   will curse.

                             % One or more blank lines separate stanzas.

   For instead of making\\
   Them do \emph{all} line breaking, \\
   It allows them to put too many words on a line when they'd rather be
   forced to be terse.
\end{verse}

Mathematical formulas may also be displayed.  A
displayed formula
is
one-line long; multiline
formulas require special formatting instructions.
   \[  \ip{\Gamma}{\psi'} = x'' + y^{2} + z_{i}^{n}\]
Don't start a paragraph with a displayed equation,
nor make one a paragraph by itself.

\end{document}               % End of document.

\section{Bench Huge Section 874}
% synthetic bench section 874
% This is a sample LaTeX input file.  (Version of 12 August 2004.)
%
% A '%' character causes TeX to ignore all remaining text on the line,
% and is used for comments like this one.

\documentclass{article}      % Specifies the document class

                             % The preamble begins here.
\title{An Example Document}  % Declares the document's title.
\author{Leslie Lamport}      % Declares the author's name.
\date{January 21, 1994}      % Deleting this command produces today's date.

\newcommand{\ip}[2]{(#1, #2)}
                             % Defines \ip{arg1}{arg2} to mean
                             % (arg1, arg2).

%\newcommand{\ip}[2]{\langle #1 | #2\rangle}
                             % This is an alternative definition of
                             % \ip that is commented out.

\begin{document}             % End of preamble and beginning of text.

\maketitle                   % Produces the title.

This is an example input file.  Comparing it with
the output it generates can show you how to
produce a simple document of your own.

\section{Ordinary Text}      % Produces section heading.  Lower-level
                             % sections are begun with similar
                             % \subsection and \subsubsection commands.

The ends  of words and sentences are marked
  by   spaces. It  doesn't matter how many
spaces    you type; one is as good as 100.  The
end of   a line counts as a space.

One   or more   blank lines denote the  end
of  a paragraph.

Since any number of consecutive spaces are treated
like a single one, the formatting of the input
file makes no difference to
      \LaTeX,                % The \LaTeX command generates the LaTeX logo.
but it makes a difference to you.  When you use
\LaTeX, making your input file as easy to read
as possible will be a great help as you write
your document and when you change it.  This sample
file shows how you can add comments to your own input
file.

Because printing is different from typewriting,
there are a number of things that you have to do
differently when preparing an input file than if
you were just typing the document directly.
Quotation marks like
       ``this''
have to be handled specially, as do quotes within
quotes:
       ``\,`this'            % \, separates the double and single quote.
        is what I just
        wrote, not  `that'\,''.

Dashes come in three sizes: an
       intra-word
dash, a medium dash for number ranges like
       1--2,
and a punctuation
       dash---like
this.

A sentence-ending space should be larger than the
space between words within a sentence.  You
sometimes have to type special commands in
conjunction with punctuation characters to get
this right, as in the following sentence.
       Gnats, gnus, etc.\ all  % `\ ' makes an inter-word space.
       begin with G\@.         % \@ marks end-of-sentence punctuation.
You should check the spaces after periods when
reading your output to make sure you haven't
forgotten any special cases.  Generating an
ellipsis
       \ldots\               % `\ ' is needed after `\ldots' because TeX
                             % ignores spaces after command names like \ldots
                             % made from \ + letters.
                             %
                             % Note how a `%' character causes TeX to ignore
                             % the end of the input line, so these blank lines
                             % do not start a new paragraph.
                             %
with the right spacing around the periods requires
a special command.

\LaTeX\ interprets some common characters as
commands, so you must type special commands to
generate them.  These characters include the
following:
       \$ \& \% \# \{ and \}.

In printing, text is usually emphasized with an
       \emph{italic}
type style.

\begin{em}
   A long segment of text can also be emphasized
   in this way.  Text within such a segment can be
   given \emph{additional} emphasis.
\end{em}

It is sometimes necessary to prevent \LaTeX\ from
breaking a line where it might otherwise do so.
This may be at a space, as between the ``Mr.''\ and
``Jones'' in
       ``Mr.~Jones'',        % ~ produces an unbreakable interword space.
or within a word---especially when the word is a
symbol like
       \mbox{\emph{itemnum}}
that makes little sense when hyphenated across
lines.

Footnotes\footnote{This is an example of a footnote.}
pose no problem.

\LaTeX\ is good at typesetting mathematical formulas
like
       \( x-3y + z = 7 \)
or
       \( a_{1} > x^{2n} + y^{2n} > x' \)
or
       \( \ip{A}{B} = \sum_{i} a_{i} b_{i} \).
The spaces you type in a formula are
ignored.  Remember that a letter like
       $x$                   % $ ... $  and  \( ... \)  are equivalent
is a formula when it denotes a mathematical
symbol, and it should be typed as one.

\section{Displayed Text}

Text is displayed by indenting it from the left
margin.  Quotations are commonly displayed.  There
are short quotations
\begin{quote}
   This is a short quotation.  It consists of a
   single paragraph of text.  See how it is formatted.
\end{quote}
and longer ones.
\begin{quotation}
   This is a longer quotation.  It consists of two
   paragraphs of text, neither of which are
   particularly interesting.

   This is the second paragraph of the quotation.  It
   is just as dull as the first paragraph.
\end{quotation}
Another frequently-displayed structure is a list.
The following is an example of an \emph{itemized}
list.
\begin{itemize}
   \item This is the first item of an itemized list.
         Each item in the list is marked with a ``tick''.
         You don't have to worry about what kind of tick
         mark is used.

   \item This is the second item of the list.  It
         contains another list nested inside it.  The inner
         list is an \emph{enumerated} list.
         \begin{enumerate}
            \item This is the first item of an enumerated
                  list that is nested within the itemized list.

            \item This is the second item of the inner list.
                  \LaTeX\ allows you to nest lists deeper than
                  you really should.
         \end{enumerate}
         This is the rest of the second item of the outer
         list.  It is no more interesting than any other
         part of the item.
   \item This is the third item of the list.
\end{itemize}
You can even display poetry.
\begin{verse}
   There is an environment
    for verse \\             % The \\ command separates lines
   Whose features some poets % within a stanza.
   will curse.

                             % One or more blank lines separate stanzas.

   For instead of making\\
   Them do \emph{all} line breaking, \\
   It allows them to put too many words on a line when they'd rather be
   forced to be terse.
\end{verse}

Mathematical formulas may also be displayed.  A
displayed formula
is
one-line long; multiline
formulas require special formatting instructions.
   \[  \ip{\Gamma}{\psi'} = x'' + y^{2} + z_{i}^{n}\]
Don't start a paragraph with a displayed equation,
nor make one a paragraph by itself.

\end{document}               % End of document.

\section{Bench Huge Section 875}
% synthetic bench section 875
% This is a sample LaTeX input file.  (Version of 12 August 2004.)
%
% A '%' character causes TeX to ignore all remaining text on the line,
% and is used for comments like this one.

\documentclass{article}      % Specifies the document class

                             % The preamble begins here.
\title{An Example Document}  % Declares the document's title.
\author{Leslie Lamport}      % Declares the author's name.
\date{January 21, 1994}      % Deleting this command produces today's date.

\newcommand{\ip}[2]{(#1, #2)}
                             % Defines \ip{arg1}{arg2} to mean
                             % (arg1, arg2).

%\newcommand{\ip}[2]{\langle #1 | #2\rangle}
                             % This is an alternative definition of
                             % \ip that is commented out.

\begin{document}             % End of preamble and beginning of text.

\maketitle                   % Produces the title.

This is an example input file.  Comparing it with
the output it generates can show you how to
produce a simple document of your own.

\section{Ordinary Text}      % Produces section heading.  Lower-level
                             % sections are begun with similar
                             % \subsection and \subsubsection commands.

The ends  of words and sentences are marked
  by   spaces. It  doesn't matter how many
spaces    you type; one is as good as 100.  The
end of   a line counts as a space.

One   or more   blank lines denote the  end
of  a paragraph.

Since any number of consecutive spaces are treated
like a single one, the formatting of the input
file makes no difference to
      \LaTeX,                % The \LaTeX command generates the LaTeX logo.
but it makes a difference to you.  When you use
\LaTeX, making your input file as easy to read
as possible will be a great help as you write
your document and when you change it.  This sample
file shows how you can add comments to your own input
file.

Because printing is different from typewriting,
there are a number of things that you have to do
differently when preparing an input file than if
you were just typing the document directly.
Quotation marks like
       ``this''
have to be handled specially, as do quotes within
quotes:
       ``\,`this'            % \, separates the double and single quote.
        is what I just
        wrote, not  `that'\,''.

Dashes come in three sizes: an
       intra-word
dash, a medium dash for number ranges like
       1--2,
and a punctuation
       dash---like
this.

A sentence-ending space should be larger than the
space between words within a sentence.  You
sometimes have to type special commands in
conjunction with punctuation characters to get
this right, as in the following sentence.
       Gnats, gnus, etc.\ all  % `\ ' makes an inter-word space.
       begin with G\@.         % \@ marks end-of-sentence punctuation.
You should check the spaces after periods when
reading your output to make sure you haven't
forgotten any special cases.  Generating an
ellipsis
       \ldots\               % `\ ' is needed after `\ldots' because TeX
                             % ignores spaces after command names like \ldots
                             % made from \ + letters.
                             %
                             % Note how a `%' character causes TeX to ignore
                             % the end of the input line, so these blank lines
                             % do not start a new paragraph.
                             %
with the right spacing around the periods requires
a special command.

\LaTeX\ interprets some common characters as
commands, so you must type special commands to
generate them.  These characters include the
following:
       \$ \& \% \# \{ and \}.

In printing, text is usually emphasized with an
       \emph{italic}
type style.

\begin{em}
   A long segment of text can also be emphasized
   in this way.  Text within such a segment can be
   given \emph{additional} emphasis.
\end{em}

It is sometimes necessary to prevent \LaTeX\ from
breaking a line where it might otherwise do so.
This may be at a space, as between the ``Mr.''\ and
``Jones'' in
       ``Mr.~Jones'',        % ~ produces an unbreakable interword space.
or within a word---especially when the word is a
symbol like
       \mbox{\emph{itemnum}}
that makes little sense when hyphenated across
lines.

Footnotes\footnote{This is an example of a footnote.}
pose no problem.

\LaTeX\ is good at typesetting mathematical formulas
like
       \( x-3y + z = 7 \)
or
       \( a_{1} > x^{2n} + y^{2n} > x' \)
or
       \( \ip{A}{B} = \sum_{i} a_{i} b_{i} \).
The spaces you type in a formula are
ignored.  Remember that a letter like
       $x$                   % $ ... $  and  \( ... \)  are equivalent
is a formula when it denotes a mathematical
symbol, and it should be typed as one.

\section{Displayed Text}

Text is displayed by indenting it from the left
margin.  Quotations are commonly displayed.  There
are short quotations
\begin{quote}
   This is a short quotation.  It consists of a
   single paragraph of text.  See how it is formatted.
\end{quote}
and longer ones.
\begin{quotation}
   This is a longer quotation.  It consists of two
   paragraphs of text, neither of which are
   particularly interesting.

   This is the second paragraph of the quotation.  It
   is just as dull as the first paragraph.
\end{quotation}
Another frequently-displayed structure is a list.
The following is an example of an \emph{itemized}
list.
\begin{itemize}
   \item This is the first item of an itemized list.
         Each item in the list is marked with a ``tick''.
         You don't have to worry about what kind of tick
         mark is used.

   \item This is the second item of the list.  It
         contains another list nested inside it.  The inner
         list is an \emph{enumerated} list.
         \begin{enumerate}
            \item This is the first item of an enumerated
                  list that is nested within the itemized list.

            \item This is the second item of the inner list.
                  \LaTeX\ allows you to nest lists deeper than
                  you really should.
         \end{enumerate}
         This is the rest of the second item of the outer
         list.  It is no more interesting than any other
         part of the item.
   \item This is the third item of the list.
\end{itemize}
You can even display poetry.
\begin{verse}
   There is an environment
    for verse \\             % The \\ command separates lines
   Whose features some poets % within a stanza.
   will curse.

                             % One or more blank lines separate stanzas.

   For instead of making\\
   Them do \emph{all} line breaking, \\
   It allows them to put too many words on a line when they'd rather be
   forced to be terse.
\end{verse}

Mathematical formulas may also be displayed.  A
displayed formula
is
one-line long; multiline
formulas require special formatting instructions.
   \[  \ip{\Gamma}{\psi'} = x'' + y^{2} + z_{i}^{n}\]
Don't start a paragraph with a displayed equation,
nor make one a paragraph by itself.

\end{document}               % End of document.

\section{Bench Huge Section 876}
% synthetic bench section 876
% This is a sample LaTeX input file.  (Version of 12 August 2004.)
%
% A '%' character causes TeX to ignore all remaining text on the line,
% and is used for comments like this one.

\documentclass{article}      % Specifies the document class

                             % The preamble begins here.
\title{An Example Document}  % Declares the document's title.
\author{Leslie Lamport}      % Declares the author's name.
\date{January 21, 1994}      % Deleting this command produces today's date.

\newcommand{\ip}[2]{(#1, #2)}
                             % Defines \ip{arg1}{arg2} to mean
                             % (arg1, arg2).

%\newcommand{\ip}[2]{\langle #1 | #2\rangle}
                             % This is an alternative definition of
                             % \ip that is commented out.

\begin{document}             % End of preamble and beginning of text.

\maketitle                   % Produces the title.

This is an example input file.  Comparing it with
the output it generates can show you how to
produce a simple document of your own.

\section{Ordinary Text}      % Produces section heading.  Lower-level
                             % sections are begun with similar
                             % \subsection and \subsubsection commands.

The ends  of words and sentences are marked
  by   spaces. It  doesn't matter how many
spaces    you type; one is as good as 100.  The
end of   a line counts as a space.

One   or more   blank lines denote the  end
of  a paragraph.

Since any number of consecutive spaces are treated
like a single one, the formatting of the input
file makes no difference to
      \LaTeX,                % The \LaTeX command generates the LaTeX logo.
but it makes a difference to you.  When you use
\LaTeX, making your input file as easy to read
as possible will be a great help as you write
your document and when you change it.  This sample
file shows how you can add comments to your own input
file.

Because printing is different from typewriting,
there are a number of things that you have to do
differently when preparing an input file than if
you were just typing the document directly.
Quotation marks like
       ``this''
have to be handled specially, as do quotes within
quotes:
       ``\,`this'            % \, separates the double and single quote.
        is what I just
        wrote, not  `that'\,''.

Dashes come in three sizes: an
       intra-word
dash, a medium dash for number ranges like
       1--2,
and a punctuation
       dash---like
this.

A sentence-ending space should be larger than the
space between words within a sentence.  You
sometimes have to type special commands in
conjunction with punctuation characters to get
this right, as in the following sentence.
       Gnats, gnus, etc.\ all  % `\ ' makes an inter-word space.
       begin with G\@.         % \@ marks end-of-sentence punctuation.
You should check the spaces after periods when
reading your output to make sure you haven't
forgotten any special cases.  Generating an
ellipsis
       \ldots\               % `\ ' is needed after `\ldots' because TeX
                             % ignores spaces after command names like \ldots
                             % made from \ + letters.
                             %
                             % Note how a `%' character causes TeX to ignore
                             % the end of the input line, so these blank lines
                             % do not start a new paragraph.
                             %
with the right spacing around the periods requires
a special command.

\LaTeX\ interprets some common characters as
commands, so you must type special commands to
generate them.  These characters include the
following:
       \$ \& \% \# \{ and \}.

In printing, text is usually emphasized with an
       \emph{italic}
type style.

\begin{em}
   A long segment of text can also be emphasized
   in this way.  Text within such a segment can be
   given \emph{additional} emphasis.
\end{em}

It is sometimes necessary to prevent \LaTeX\ from
breaking a line where it might otherwise do so.
This may be at a space, as between the ``Mr.''\ and
``Jones'' in
       ``Mr.~Jones'',        % ~ produces an unbreakable interword space.
or within a word---especially when the word is a
symbol like
       \mbox{\emph{itemnum}}
that makes little sense when hyphenated across
lines.

Footnotes\footnote{This is an example of a footnote.}
pose no problem.

\LaTeX\ is good at typesetting mathematical formulas
like
       \( x-3y + z = 7 \)
or
       \( a_{1} > x^{2n} + y^{2n} > x' \)
or
       \( \ip{A}{B} = \sum_{i} a_{i} b_{i} \).
The spaces you type in a formula are
ignored.  Remember that a letter like
       $x$                   % $ ... $  and  \( ... \)  are equivalent
is a formula when it denotes a mathematical
symbol, and it should be typed as one.

\section{Displayed Text}

Text is displayed by indenting it from the left
margin.  Quotations are commonly displayed.  There
are short quotations
\begin{quote}
   This is a short quotation.  It consists of a
   single paragraph of text.  See how it is formatted.
\end{quote}
and longer ones.
\begin{quotation}
   This is a longer quotation.  It consists of two
   paragraphs of text, neither of which are
   particularly interesting.

   This is the second paragraph of the quotation.  It
   is just as dull as the first paragraph.
\end{quotation}
Another frequently-displayed structure is a list.
The following is an example of an \emph{itemized}
list.
\begin{itemize}
   \item This is the first item of an itemized list.
         Each item in the list is marked with a ``tick''.
         You don't have to worry about what kind of tick
         mark is used.

   \item This is the second item of the list.  It
         contains another list nested inside it.  The inner
         list is an \emph{enumerated} list.
         \begin{enumerate}
            \item This is the first item of an enumerated
                  list that is nested within the itemized list.

            \item This is the second item of the inner list.
                  \LaTeX\ allows you to nest lists deeper than
                  you really should.
         \end{enumerate}
         This is the rest of the second item of the outer
         list.  It is no more interesting than any other
         part of the item.
   \item This is the third item of the list.
\end{itemize}
You can even display poetry.
\begin{verse}
   There is an environment
    for verse \\             % The \\ command separates lines
   Whose features some poets % within a stanza.
   will curse.

                             % One or more blank lines separate stanzas.

   For instead of making\\
   Them do \emph{all} line breaking, \\
   It allows them to put too many words on a line when they'd rather be
   forced to be terse.
\end{verse}

Mathematical formulas may also be displayed.  A
displayed formula
is
one-line long; multiline
formulas require special formatting instructions.
   \[  \ip{\Gamma}{\psi'} = x'' + y^{2} + z_{i}^{n}\]
Don't start a paragraph with a displayed equation,
nor make one a paragraph by itself.

\end{document}               % End of document.

\section{Bench Huge Section 877}
% synthetic bench section 877
% This is a sample LaTeX input file.  (Version of 12 August 2004.)
%
% A '%' character causes TeX to ignore all remaining text on the line,
% and is used for comments like this one.

\documentclass{article}      % Specifies the document class

                             % The preamble begins here.
\title{An Example Document}  % Declares the document's title.
\author{Leslie Lamport}      % Declares the author's name.
\date{January 21, 1994}      % Deleting this command produces today's date.

\newcommand{\ip}[2]{(#1, #2)}
                             % Defines \ip{arg1}{arg2} to mean
                             % (arg1, arg2).

%\newcommand{\ip}[2]{\langle #1 | #2\rangle}
                             % This is an alternative definition of
                             % \ip that is commented out.

\begin{document}             % End of preamble and beginning of text.

\maketitle                   % Produces the title.

This is an example input file.  Comparing it with
the output it generates can show you how to
produce a simple document of your own.

\section{Ordinary Text}      % Produces section heading.  Lower-level
                             % sections are begun with similar
                             % \subsection and \subsubsection commands.

The ends  of words and sentences are marked
  by   spaces. It  doesn't matter how many
spaces    you type; one is as good as 100.  The
end of   a line counts as a space.

One   or more   blank lines denote the  end
of  a paragraph.

Since any number of consecutive spaces are treated
like a single one, the formatting of the input
file makes no difference to
      \LaTeX,                % The \LaTeX command generates the LaTeX logo.
but it makes a difference to you.  When you use
\LaTeX, making your input file as easy to read
as possible will be a great help as you write
your document and when you change it.  This sample
file shows how you can add comments to your own input
file.

Because printing is different from typewriting,
there are a number of things that you have to do
differently when preparing an input file than if
you were just typing the document directly.
Quotation marks like
       ``this''
have to be handled specially, as do quotes within
quotes:
       ``\,`this'            % \, separates the double and single quote.
        is what I just
        wrote, not  `that'\,''.

Dashes come in three sizes: an
       intra-word
dash, a medium dash for number ranges like
       1--2,
and a punctuation
       dash---like
this.

A sentence-ending space should be larger than the
space between words within a sentence.  You
sometimes have to type special commands in
conjunction with punctuation characters to get
this right, as in the following sentence.
       Gnats, gnus, etc.\ all  % `\ ' makes an inter-word space.
       begin with G\@.         % \@ marks end-of-sentence punctuation.
You should check the spaces after periods when
reading your output to make sure you haven't
forgotten any special cases.  Generating an
ellipsis
       \ldots\               % `\ ' is needed after `\ldots' because TeX
                             % ignores spaces after command names like \ldots
                             % made from \ + letters.
                             %
                             % Note how a `%' character causes TeX to ignore
                             % the end of the input line, so these blank lines
                             % do not start a new paragraph.
                             %
with the right spacing around the periods requires
a special command.

\LaTeX\ interprets some common characters as
commands, so you must type special commands to
generate them.  These characters include the
following:
       \$ \& \% \# \{ and \}.

In printing, text is usually emphasized with an
       \emph{italic}
type style.

\begin{em}
   A long segment of text can also be emphasized
   in this way.  Text within such a segment can be
   given \emph{additional} emphasis.
\end{em}

It is sometimes necessary to prevent \LaTeX\ from
breaking a line where it might otherwise do so.
This may be at a space, as between the ``Mr.''\ and
``Jones'' in
       ``Mr.~Jones'',        % ~ produces an unbreakable interword space.
or within a word---especially when the word is a
symbol like
       \mbox{\emph{itemnum}}
that makes little sense when hyphenated across
lines.

Footnotes\footnote{This is an example of a footnote.}
pose no problem.

\LaTeX\ is good at typesetting mathematical formulas
like
       \( x-3y + z = 7 \)
or
       \( a_{1} > x^{2n} + y^{2n} > x' \)
or
       \( \ip{A}{B} = \sum_{i} a_{i} b_{i} \).
The spaces you type in a formula are
ignored.  Remember that a letter like
       $x$                   % $ ... $  and  \( ... \)  are equivalent
is a formula when it denotes a mathematical
symbol, and it should be typed as one.

\section{Displayed Text}

Text is displayed by indenting it from the left
margin.  Quotations are commonly displayed.  There
are short quotations
\begin{quote}
   This is a short quotation.  It consists of a
   single paragraph of text.  See how it is formatted.
\end{quote}
and longer ones.
\begin{quotation}
   This is a longer quotation.  It consists of two
   paragraphs of text, neither of which are
   particularly interesting.

   This is the second paragraph of the quotation.  It
   is just as dull as the first paragraph.
\end{quotation}
Another frequently-displayed structure is a list.
The following is an example of an \emph{itemized}
list.
\begin{itemize}
   \item This is the first item of an itemized list.
         Each item in the list is marked with a ``tick''.
         You don't have to worry about what kind of tick
         mark is used.

   \item This is the second item of the list.  It
         contains another list nested inside it.  The inner
         list is an \emph{enumerated} list.
         \begin{enumerate}
            \item This is the first item of an enumerated
                  list that is nested within the itemized list.

            \item This is the second item of the inner list.
                  \LaTeX\ allows you to nest lists deeper than
                  you really should.
         \end{enumerate}
         This is the rest of the second item of the outer
         list.  It is no more interesting than any other
         part of the item.
   \item This is the third item of the list.
\end{itemize}
You can even display poetry.
\begin{verse}
   There is an environment
    for verse \\             % The \\ command separates lines
   Whose features some poets % within a stanza.
   will curse.

                             % One or more blank lines separate stanzas.

   For instead of making\\
   Them do \emph{all} line breaking, \\
   It allows them to put too many words on a line when they'd rather be
   forced to be terse.
\end{verse}

Mathematical formulas may also be displayed.  A
displayed formula
is
one-line long; multiline
formulas require special formatting instructions.
   \[  \ip{\Gamma}{\psi'} = x'' + y^{2} + z_{i}^{n}\]
Don't start a paragraph with a displayed equation,
nor make one a paragraph by itself.

\end{document}               % End of document.

\section{Bench Huge Section 878}
% synthetic bench section 878
% This is a sample LaTeX input file.  (Version of 12 August 2004.)
%
% A '%' character causes TeX to ignore all remaining text on the line,
% and is used for comments like this one.

\documentclass{article}      % Specifies the document class

                             % The preamble begins here.
\title{An Example Document}  % Declares the document's title.
\author{Leslie Lamport}      % Declares the author's name.
\date{January 21, 1994}      % Deleting this command produces today's date.

\newcommand{\ip}[2]{(#1, #2)}
                             % Defines \ip{arg1}{arg2} to mean
                             % (arg1, arg2).

%\newcommand{\ip}[2]{\langle #1 | #2\rangle}
                             % This is an alternative definition of
                             % \ip that is commented out.

\begin{document}             % End of preamble and beginning of text.

\maketitle                   % Produces the title.

This is an example input file.  Comparing it with
the output it generates can show you how to
produce a simple document of your own.

\section{Ordinary Text}      % Produces section heading.  Lower-level
                             % sections are begun with similar
                             % \subsection and \subsubsection commands.

The ends  of words and sentences are marked
  by   spaces. It  doesn't matter how many
spaces    you type; one is as good as 100.  The
end of   a line counts as a space.

One   or more   blank lines denote the  end
of  a paragraph.

Since any number of consecutive spaces are treated
like a single one, the formatting of the input
file makes no difference to
      \LaTeX,                % The \LaTeX command generates the LaTeX logo.
but it makes a difference to you.  When you use
\LaTeX, making your input file as easy to read
as possible will be a great help as you write
your document and when you change it.  This sample
file shows how you can add comments to your own input
file.

Because printing is different from typewriting,
there are a number of things that you have to do
differently when preparing an input file than if
you were just typing the document directly.
Quotation marks like
       ``this''
have to be handled specially, as do quotes within
quotes:
       ``\,`this'            % \, separates the double and single quote.
        is what I just
        wrote, not  `that'\,''.

Dashes come in three sizes: an
       intra-word
dash, a medium dash for number ranges like
       1--2,
and a punctuation
       dash---like
this.

A sentence-ending space should be larger than the
space between words within a sentence.  You
sometimes have to type special commands in
conjunction with punctuation characters to get
this right, as in the following sentence.
       Gnats, gnus, etc.\ all  % `\ ' makes an inter-word space.
       begin with G\@.         % \@ marks end-of-sentence punctuation.
You should check the spaces after periods when
reading your output to make sure you haven't
forgotten any special cases.  Generating an
ellipsis
       \ldots\               % `\ ' is needed after `\ldots' because TeX
                             % ignores spaces after command names like \ldots
                             % made from \ + letters.
                             %
                             % Note how a `%' character causes TeX to ignore
                             % the end of the input line, so these blank lines
                             % do not start a new paragraph.
                             %
with the right spacing around the periods requires
a special command.

\LaTeX\ interprets some common characters as
commands, so you must type special commands to
generate them.  These characters include the
following:
       \$ \& \% \# \{ and \}.

In printing, text is usually emphasized with an
       \emph{italic}
type style.

\begin{em}
   A long segment of text can also be emphasized
   in this way.  Text within such a segment can be
   given \emph{additional} emphasis.
\end{em}

It is sometimes necessary to prevent \LaTeX\ from
breaking a line where it might otherwise do so.
This may be at a space, as between the ``Mr.''\ and
``Jones'' in
       ``Mr.~Jones'',        % ~ produces an unbreakable interword space.
or within a word---especially when the word is a
symbol like
       \mbox{\emph{itemnum}}
that makes little sense when hyphenated across
lines.

Footnotes\footnote{This is an example of a footnote.}
pose no problem.

\LaTeX\ is good at typesetting mathematical formulas
like
       \( x-3y + z = 7 \)
or
       \( a_{1} > x^{2n} + y^{2n} > x' \)
or
       \( \ip{A}{B} = \sum_{i} a_{i} b_{i} \).
The spaces you type in a formula are
ignored.  Remember that a letter like
       $x$                   % $ ... $  and  \( ... \)  are equivalent
is a formula when it denotes a mathematical
symbol, and it should be typed as one.

\section{Displayed Text}

Text is displayed by indenting it from the left
margin.  Quotations are commonly displayed.  There
are short quotations
\begin{quote}
   This is a short quotation.  It consists of a
   single paragraph of text.  See how it is formatted.
\end{quote}
and longer ones.
\begin{quotation}
   This is a longer quotation.  It consists of two
   paragraphs of text, neither of which are
   particularly interesting.

   This is the second paragraph of the quotation.  It
   is just as dull as the first paragraph.
\end{quotation}
Another frequently-displayed structure is a list.
The following is an example of an \emph{itemized}
list.
\begin{itemize}
   \item This is the first item of an itemized list.
         Each item in the list is marked with a ``tick''.
         You don't have to worry about what kind of tick
         mark is used.

   \item This is the second item of the list.  It
         contains another list nested inside it.  The inner
         list is an \emph{enumerated} list.
         \begin{enumerate}
            \item This is the first item of an enumerated
                  list that is nested within the itemized list.

            \item This is the second item of the inner list.
                  \LaTeX\ allows you to nest lists deeper than
                  you really should.
         \end{enumerate}
         This is the rest of the second item of the outer
         list.  It is no more interesting than any other
         part of the item.
   \item This is the third item of the list.
\end{itemize}
You can even display poetry.
\begin{verse}
   There is an environment
    for verse \\             % The \\ command separates lines
   Whose features some poets % within a stanza.
   will curse.

                             % One or more blank lines separate stanzas.

   For instead of making\\
   Them do \emph{all} line breaking, \\
   It allows them to put too many words on a line when they'd rather be
   forced to be terse.
\end{verse}

Mathematical formulas may also be displayed.  A
displayed formula
is
one-line long; multiline
formulas require special formatting instructions.
   \[  \ip{\Gamma}{\psi'} = x'' + y^{2} + z_{i}^{n}\]
Don't start a paragraph with a displayed equation,
nor make one a paragraph by itself.

\end{document}               % End of document.

\section{Bench Huge Section 879}
% synthetic bench section 879
% This is a sample LaTeX input file.  (Version of 12 August 2004.)
%
% A '%' character causes TeX to ignore all remaining text on the line,
% and is used for comments like this one.

\documentclass{article}      % Specifies the document class

                             % The preamble begins here.
\title{An Example Document}  % Declares the document's title.
\author{Leslie Lamport}      % Declares the author's name.
\date{January 21, 1994}      % Deleting this command produces today's date.

\newcommand{\ip}[2]{(#1, #2)}
                             % Defines \ip{arg1}{arg2} to mean
                             % (arg1, arg2).

%\newcommand{\ip}[2]{\langle #1 | #2\rangle}
                             % This is an alternative definition of
                             % \ip that is commented out.

\begin{document}             % End of preamble and beginning of text.

\maketitle                   % Produces the title.

This is an example input file.  Comparing it with
the output it generates can show you how to
produce a simple document of your own.

\section{Ordinary Text}      % Produces section heading.  Lower-level
                             % sections are begun with similar
                             % \subsection and \subsubsection commands.

The ends  of words and sentences are marked
  by   spaces. It  doesn't matter how many
spaces    you type; one is as good as 100.  The
end of   a line counts as a space.

One   or more   blank lines denote the  end
of  a paragraph.

Since any number of consecutive spaces are treated
like a single one, the formatting of the input
file makes no difference to
      \LaTeX,                % The \LaTeX command generates the LaTeX logo.
but it makes a difference to you.  When you use
\LaTeX, making your input file as easy to read
as possible will be a great help as you write
your document and when you change it.  This sample
file shows how you can add comments to your own input
file.

Because printing is different from typewriting,
there are a number of things that you have to do
differently when preparing an input file than if
you were just typing the document directly.
Quotation marks like
       ``this''
have to be handled specially, as do quotes within
quotes:
       ``\,`this'            % \, separates the double and single quote.
        is what I just
        wrote, not  `that'\,''.

Dashes come in three sizes: an
       intra-word
dash, a medium dash for number ranges like
       1--2,
and a punctuation
       dash---like
this.

A sentence-ending space should be larger than the
space between words within a sentence.  You
sometimes have to type special commands in
conjunction with punctuation characters to get
this right, as in the following sentence.
       Gnats, gnus, etc.\ all  % `\ ' makes an inter-word space.
       begin with G\@.         % \@ marks end-of-sentence punctuation.
You should check the spaces after periods when
reading your output to make sure you haven't
forgotten any special cases.  Generating an
ellipsis
       \ldots\               % `\ ' is needed after `\ldots' because TeX
                             % ignores spaces after command names like \ldots
                             % made from \ + letters.
                             %
                             % Note how a `%' character causes TeX to ignore
                             % the end of the input line, so these blank lines
                             % do not start a new paragraph.
                             %
with the right spacing around the periods requires
a special command.

\LaTeX\ interprets some common characters as
commands, so you must type special commands to
generate them.  These characters include the
following:
       \$ \& \% \# \{ and \}.

In printing, text is usually emphasized with an
       \emph{italic}
type style.

\begin{em}
   A long segment of text can also be emphasized
   in this way.  Text within such a segment can be
   given \emph{additional} emphasis.
\end{em}

It is sometimes necessary to prevent \LaTeX\ from
breaking a line where it might otherwise do so.
This may be at a space, as between the ``Mr.''\ and
``Jones'' in
       ``Mr.~Jones'',        % ~ produces an unbreakable interword space.
or within a word---especially when the word is a
symbol like
       \mbox{\emph{itemnum}}
that makes little sense when hyphenated across
lines.

Footnotes\footnote{This is an example of a footnote.}
pose no problem.

\LaTeX\ is good at typesetting mathematical formulas
like
       \( x-3y + z = 7 \)
or
       \( a_{1} > x^{2n} + y^{2n} > x' \)
or
       \( \ip{A}{B} = \sum_{i} a_{i} b_{i} \).
The spaces you type in a formula are
ignored.  Remember that a letter like
       $x$                   % $ ... $  and  \( ... \)  are equivalent
is a formula when it denotes a mathematical
symbol, and it should be typed as one.

\section{Displayed Text}

Text is displayed by indenting it from the left
margin.  Quotations are commonly displayed.  There
are short quotations
\begin{quote}
   This is a short quotation.  It consists of a
   single paragraph of text.  See how it is formatted.
\end{quote}
and longer ones.
\begin{quotation}
   This is a longer quotation.  It consists of two
   paragraphs of text, neither of which are
   particularly interesting.

   This is the second paragraph of the quotation.  It
   is just as dull as the first paragraph.
\end{quotation}
Another frequently-displayed structure is a list.
The following is an example of an \emph{itemized}
list.
\begin{itemize}
   \item This is the first item of an itemized list.
         Each item in the list is marked with a ``tick''.
         You don't have to worry about what kind of tick
         mark is used.

   \item This is the second item of the list.  It
         contains another list nested inside it.  The inner
         list is an \emph{enumerated} list.
         \begin{enumerate}
            \item This is the first item of an enumerated
                  list that is nested within the itemized list.

            \item This is the second item of the inner list.
                  \LaTeX\ allows you to nest lists deeper than
                  you really should.
         \end{enumerate}
         This is the rest of the second item of the outer
         list.  It is no more interesting than any other
         part of the item.
   \item This is the third item of the list.
\end{itemize}
You can even display poetry.
\begin{verse}
   There is an environment
    for verse \\             % The \\ command separates lines
   Whose features some poets % within a stanza.
   will curse.

                             % One or more blank lines separate stanzas.

   For instead of making\\
   Them do \emph{all} line breaking, \\
   It allows them to put too many words on a line when they'd rather be
   forced to be terse.
\end{verse}

Mathematical formulas may also be displayed.  A
displayed formula
is
one-line long; multiline
formulas require special formatting instructions.
   \[  \ip{\Gamma}{\psi'} = x'' + y^{2} + z_{i}^{n}\]
Don't start a paragraph with a displayed equation,
nor make one a paragraph by itself.

\end{document}               % End of document.

\section{Bench Huge Section 880}
% synthetic bench section 880
% This is a sample LaTeX input file.  (Version of 12 August 2004.)
%
% A '%' character causes TeX to ignore all remaining text on the line,
% and is used for comments like this one.

\documentclass{article}      % Specifies the document class

                             % The preamble begins here.
\title{An Example Document}  % Declares the document's title.
\author{Leslie Lamport}      % Declares the author's name.
\date{January 21, 1994}      % Deleting this command produces today's date.

\newcommand{\ip}[2]{(#1, #2)}
                             % Defines \ip{arg1}{arg2} to mean
                             % (arg1, arg2).

%\newcommand{\ip}[2]{\langle #1 | #2\rangle}
                             % This is an alternative definition of
                             % \ip that is commented out.

\begin{document}             % End of preamble and beginning of text.

\maketitle                   % Produces the title.

This is an example input file.  Comparing it with
the output it generates can show you how to
produce a simple document of your own.

\section{Ordinary Text}      % Produces section heading.  Lower-level
                             % sections are begun with similar
                             % \subsection and \subsubsection commands.

The ends  of words and sentences are marked
  by   spaces. It  doesn't matter how many
spaces    you type; one is as good as 100.  The
end of   a line counts as a space.

One   or more   blank lines denote the  end
of  a paragraph.

Since any number of consecutive spaces are treated
like a single one, the formatting of the input
file makes no difference to
      \LaTeX,                % The \LaTeX command generates the LaTeX logo.
but it makes a difference to you.  When you use
\LaTeX, making your input file as easy to read
as possible will be a great help as you write
your document and when you change it.  This sample
file shows how you can add comments to your own input
file.

Because printing is different from typewriting,
there are a number of things that you have to do
differently when preparing an input file than if
you were just typing the document directly.
Quotation marks like
       ``this''
have to be handled specially, as do quotes within
quotes:
       ``\,`this'            % \, separates the double and single quote.
        is what I just
        wrote, not  `that'\,''.

Dashes come in three sizes: an
       intra-word
dash, a medium dash for number ranges like
       1--2,
and a punctuation
       dash---like
this.

A sentence-ending space should be larger than the
space between words within a sentence.  You
sometimes have to type special commands in
conjunction with punctuation characters to get
this right, as in the following sentence.
       Gnats, gnus, etc.\ all  % `\ ' makes an inter-word space.
       begin with G\@.         % \@ marks end-of-sentence punctuation.
You should check the spaces after periods when
reading your output to make sure you haven't
forgotten any special cases.  Generating an
ellipsis
       \ldots\               % `\ ' is needed after `\ldots' because TeX
                             % ignores spaces after command names like \ldots
                             % made from \ + letters.
                             %
                             % Note how a `%' character causes TeX to ignore
                             % the end of the input line, so these blank lines
                             % do not start a new paragraph.
                             %
with the right spacing around the periods requires
a special command.

\LaTeX\ interprets some common characters as
commands, so you must type special commands to
generate them.  These characters include the
following:
       \$ \& \% \# \{ and \}.

In printing, text is usually emphasized with an
       \emph{italic}
type style.

\begin{em}
   A long segment of text can also be emphasized
   in this way.  Text within such a segment can be
   given \emph{additional} emphasis.
\end{em}

It is sometimes necessary to prevent \LaTeX\ from
breaking a line where it might otherwise do so.
This may be at a space, as between the ``Mr.''\ and
``Jones'' in
       ``Mr.~Jones'',        % ~ produces an unbreakable interword space.
or within a word---especially when the word is a
symbol like
       \mbox{\emph{itemnum}}
that makes little sense when hyphenated across
lines.

Footnotes\footnote{This is an example of a footnote.}
pose no problem.

\LaTeX\ is good at typesetting mathematical formulas
like
       \( x-3y + z = 7 \)
or
       \( a_{1} > x^{2n} + y^{2n} > x' \)
or
       \( \ip{A}{B} = \sum_{i} a_{i} b_{i} \).
The spaces you type in a formula are
ignored.  Remember that a letter like
       $x$                   % $ ... $  and  \( ... \)  are equivalent
is a formula when it denotes a mathematical
symbol, and it should be typed as one.

\section{Displayed Text}

Text is displayed by indenting it from the left
margin.  Quotations are commonly displayed.  There
are short quotations
\begin{quote}
   This is a short quotation.  It consists of a
   single paragraph of text.  See how it is formatted.
\end{quote}
and longer ones.
\begin{quotation}
   This is a longer quotation.  It consists of two
   paragraphs of text, neither of which are
   particularly interesting.

   This is the second paragraph of the quotation.  It
   is just as dull as the first paragraph.
\end{quotation}
Another frequently-displayed structure is a list.
The following is an example of an \emph{itemized}
list.
\begin{itemize}
   \item This is the first item of an itemized list.
         Each item in the list is marked with a ``tick''.
         You don't have to worry about what kind of tick
         mark is used.

   \item This is the second item of the list.  It
         contains another list nested inside it.  The inner
         list is an \emph{enumerated} list.
         \begin{enumerate}
            \item This is the first item of an enumerated
                  list that is nested within the itemized list.

            \item This is the second item of the inner list.
                  \LaTeX\ allows you to nest lists deeper than
                  you really should.
         \end{enumerate}
         This is the rest of the second item of the outer
         list.  It is no more interesting than any other
         part of the item.
   \item This is the third item of the list.
\end{itemize}
You can even display poetry.
\begin{verse}
   There is an environment
    for verse \\             % The \\ command separates lines
   Whose features some poets % within a stanza.
   will curse.

                             % One or more blank lines separate stanzas.

   For instead of making\\
   Them do \emph{all} line breaking, \\
   It allows them to put too many words on a line when they'd rather be
   forced to be terse.
\end{verse}

Mathematical formulas may also be displayed.  A
displayed formula
is
one-line long; multiline
formulas require special formatting instructions.
   \[  \ip{\Gamma}{\psi'} = x'' + y^{2} + z_{i}^{n}\]
Don't start a paragraph with a displayed equation,
nor make one a paragraph by itself.

\end{document}               % End of document.

\section{Bench Huge Section 881}
% synthetic bench section 881
% This is a sample LaTeX input file.  (Version of 12 August 2004.)
%
% A '%' character causes TeX to ignore all remaining text on the line,
% and is used for comments like this one.

\documentclass{article}      % Specifies the document class

                             % The preamble begins here.
\title{An Example Document}  % Declares the document's title.
\author{Leslie Lamport}      % Declares the author's name.
\date{January 21, 1994}      % Deleting this command produces today's date.

\newcommand{\ip}[2]{(#1, #2)}
                             % Defines \ip{arg1}{arg2} to mean
                             % (arg1, arg2).

%\newcommand{\ip}[2]{\langle #1 | #2\rangle}
                             % This is an alternative definition of
                             % \ip that is commented out.

\begin{document}             % End of preamble and beginning of text.

\maketitle                   % Produces the title.

This is an example input file.  Comparing it with
the output it generates can show you how to
produce a simple document of your own.

\section{Ordinary Text}      % Produces section heading.  Lower-level
                             % sections are begun with similar
                             % \subsection and \subsubsection commands.

The ends  of words and sentences are marked
  by   spaces. It  doesn't matter how many
spaces    you type; one is as good as 100.  The
end of   a line counts as a space.

One   or more   blank lines denote the  end
of  a paragraph.

Since any number of consecutive spaces are treated
like a single one, the formatting of the input
file makes no difference to
      \LaTeX,                % The \LaTeX command generates the LaTeX logo.
but it makes a difference to you.  When you use
\LaTeX, making your input file as easy to read
as possible will be a great help as you write
your document and when you change it.  This sample
file shows how you can add comments to your own input
file.

Because printing is different from typewriting,
there are a number of things that you have to do
differently when preparing an input file than if
you were just typing the document directly.
Quotation marks like
       ``this''
have to be handled specially, as do quotes within
quotes:
       ``\,`this'            % \, separates the double and single quote.
        is what I just
        wrote, not  `that'\,''.

Dashes come in three sizes: an
       intra-word
dash, a medium dash for number ranges like
       1--2,
and a punctuation
       dash---like
this.

A sentence-ending space should be larger than the
space between words within a sentence.  You
sometimes have to type special commands in
conjunction with punctuation characters to get
this right, as in the following sentence.
       Gnats, gnus, etc.\ all  % `\ ' makes an inter-word space.
       begin with G\@.         % \@ marks end-of-sentence punctuation.
You should check the spaces after periods when
reading your output to make sure you haven't
forgotten any special cases.  Generating an
ellipsis
       \ldots\               % `\ ' is needed after `\ldots' because TeX
                             % ignores spaces after command names like \ldots
                             % made from \ + letters.
                             %
                             % Note how a `%' character causes TeX to ignore
                             % the end of the input line, so these blank lines
                             % do not start a new paragraph.
                             %
with the right spacing around the periods requires
a special command.

\LaTeX\ interprets some common characters as
commands, so you must type special commands to
generate them.  These characters include the
following:
       \$ \& \% \# \{ and \}.

In printing, text is usually emphasized with an
       \emph{italic}
type style.

\begin{em}
   A long segment of text can also be emphasized
   in this way.  Text within such a segment can be
   given \emph{additional} emphasis.
\end{em}

It is sometimes necessary to prevent \LaTeX\ from
breaking a line where it might otherwise do so.
This may be at a space, as between the ``Mr.''\ and
``Jones'' in
       ``Mr.~Jones'',        % ~ produces an unbreakable interword space.
or within a word---especially when the word is a
symbol like
       \mbox{\emph{itemnum}}
that makes little sense when hyphenated across
lines.

Footnotes\footnote{This is an example of a footnote.}
pose no problem.

\LaTeX\ is good at typesetting mathematical formulas
like
       \( x-3y + z = 7 \)
or
       \( a_{1} > x^{2n} + y^{2n} > x' \)
or
       \( \ip{A}{B} = \sum_{i} a_{i} b_{i} \).
The spaces you type in a formula are
ignored.  Remember that a letter like
       $x$                   % $ ... $  and  \( ... \)  are equivalent
is a formula when it denotes a mathematical
symbol, and it should be typed as one.

\section{Displayed Text}

Text is displayed by indenting it from the left
margin.  Quotations are commonly displayed.  There
are short quotations
\begin{quote}
   This is a short quotation.  It consists of a
   single paragraph of text.  See how it is formatted.
\end{quote}
and longer ones.
\begin{quotation}
   This is a longer quotation.  It consists of two
   paragraphs of text, neither of which are
   particularly interesting.

   This is the second paragraph of the quotation.  It
   is just as dull as the first paragraph.
\end{quotation}
Another frequently-displayed structure is a list.
The following is an example of an \emph{itemized}
list.
\begin{itemize}
   \item This is the first item of an itemized list.
         Each item in the list is marked with a ``tick''.
         You don't have to worry about what kind of tick
         mark is used.

   \item This is the second item of the list.  It
         contains another list nested inside it.  The inner
         list is an \emph{enumerated} list.
         \begin{enumerate}
            \item This is the first item of an enumerated
                  list that is nested within the itemized list.

            \item This is the second item of the inner list.
                  \LaTeX\ allows you to nest lists deeper than
                  you really should.
         \end{enumerate}
         This is the rest of the second item of the outer
         list.  It is no more interesting than any other
         part of the item.
   \item This is the third item of the list.
\end{itemize}
You can even display poetry.
\begin{verse}
   There is an environment
    for verse \\             % The \\ command separates lines
   Whose features some poets % within a stanza.
   will curse.

                             % One or more blank lines separate stanzas.

   For instead of making\\
   Them do \emph{all} line breaking, \\
   It allows them to put too many words on a line when they'd rather be
   forced to be terse.
\end{verse}

Mathematical formulas may also be displayed.  A
displayed formula
is
one-line long; multiline
formulas require special formatting instructions.
   \[  \ip{\Gamma}{\psi'} = x'' + y^{2} + z_{i}^{n}\]
Don't start a paragraph with a displayed equation,
nor make one a paragraph by itself.

\end{document}               % End of document.

\section{Bench Huge Section 882}
% synthetic bench section 882
% This is a sample LaTeX input file.  (Version of 12 August 2004.)
%
% A '%' character causes TeX to ignore all remaining text on the line,
% and is used for comments like this one.

\documentclass{article}      % Specifies the document class

                             % The preamble begins here.
\title{An Example Document}  % Declares the document's title.
\author{Leslie Lamport}      % Declares the author's name.
\date{January 21, 1994}      % Deleting this command produces today's date.

\newcommand{\ip}[2]{(#1, #2)}
                             % Defines \ip{arg1}{arg2} to mean
                             % (arg1, arg2).

%\newcommand{\ip}[2]{\langle #1 | #2\rangle}
                             % This is an alternative definition of
                             % \ip that is commented out.

\begin{document}             % End of preamble and beginning of text.

\maketitle                   % Produces the title.

This is an example input file.  Comparing it with
the output it generates can show you how to
produce a simple document of your own.

\section{Ordinary Text}      % Produces section heading.  Lower-level
                             % sections are begun with similar
                             % \subsection and \subsubsection commands.

The ends  of words and sentences are marked
  by   spaces. It  doesn't matter how many
spaces    you type; one is as good as 100.  The
end of   a line counts as a space.

One   or more   blank lines denote the  end
of  a paragraph.

Since any number of consecutive spaces are treated
like a single one, the formatting of the input
file makes no difference to
      \LaTeX,                % The \LaTeX command generates the LaTeX logo.
but it makes a difference to you.  When you use
\LaTeX, making your input file as easy to read
as possible will be a great help as you write
your document and when you change it.  This sample
file shows how you can add comments to your own input
file.

Because printing is different from typewriting,
there are a number of things that you have to do
differently when preparing an input file than if
you were just typing the document directly.
Quotation marks like
       ``this''
have to be handled specially, as do quotes within
quotes:
       ``\,`this'            % \, separates the double and single quote.
        is what I just
        wrote, not  `that'\,''.

Dashes come in three sizes: an
       intra-word
dash, a medium dash for number ranges like
       1--2,
and a punctuation
       dash---like
this.

A sentence-ending space should be larger than the
space between words within a sentence.  You
sometimes have to type special commands in
conjunction with punctuation characters to get
this right, as in the following sentence.
       Gnats, gnus, etc.\ all  % `\ ' makes an inter-word space.
       begin with G\@.         % \@ marks end-of-sentence punctuation.
You should check the spaces after periods when
reading your output to make sure you haven't
forgotten any special cases.  Generating an
ellipsis
       \ldots\               % `\ ' is needed after `\ldots' because TeX
                             % ignores spaces after command names like \ldots
                             % made from \ + letters.
                             %
                             % Note how a `%' character causes TeX to ignore
                             % the end of the input line, so these blank lines
                             % do not start a new paragraph.
                             %
with the right spacing around the periods requires
a special command.

\LaTeX\ interprets some common characters as
commands, so you must type special commands to
generate them.  These characters include the
following:
       \$ \& \% \# \{ and \}.

In printing, text is usually emphasized with an
       \emph{italic}
type style.

\begin{em}
   A long segment of text can also be emphasized
   in this way.  Text within such a segment can be
   given \emph{additional} emphasis.
\end{em}

It is sometimes necessary to prevent \LaTeX\ from
breaking a line where it might otherwise do so.
This may be at a space, as between the ``Mr.''\ and
``Jones'' in
       ``Mr.~Jones'',        % ~ produces an unbreakable interword space.
or within a word---especially when the word is a
symbol like
       \mbox{\emph{itemnum}}
that makes little sense when hyphenated across
lines.

Footnotes\footnote{This is an example of a footnote.}
pose no problem.

\LaTeX\ is good at typesetting mathematical formulas
like
       \( x-3y + z = 7 \)
or
       \( a_{1} > x^{2n} + y^{2n} > x' \)
or
       \( \ip{A}{B} = \sum_{i} a_{i} b_{i} \).
The spaces you type in a formula are
ignored.  Remember that a letter like
       $x$                   % $ ... $  and  \( ... \)  are equivalent
is a formula when it denotes a mathematical
symbol, and it should be typed as one.

\section{Displayed Text}

Text is displayed by indenting it from the left
margin.  Quotations are commonly displayed.  There
are short quotations
\begin{quote}
   This is a short quotation.  It consists of a
   single paragraph of text.  See how it is formatted.
\end{quote}
and longer ones.
\begin{quotation}
   This is a longer quotation.  It consists of two
   paragraphs of text, neither of which are
   particularly interesting.

   This is the second paragraph of the quotation.  It
   is just as dull as the first paragraph.
\end{quotation}
Another frequently-displayed structure is a list.
The following is an example of an \emph{itemized}
list.
\begin{itemize}
   \item This is the first item of an itemized list.
         Each item in the list is marked with a ``tick''.
         You don't have to worry about what kind of tick
         mark is used.

   \item This is the second item of the list.  It
         contains another list nested inside it.  The inner
         list is an \emph{enumerated} list.
         \begin{enumerate}
            \item This is the first item of an enumerated
                  list that is nested within the itemized list.

            \item This is the second item of the inner list.
                  \LaTeX\ allows you to nest lists deeper than
                  you really should.
         \end{enumerate}
         This is the rest of the second item of the outer
         list.  It is no more interesting than any other
         part of the item.
   \item This is the third item of the list.
\end{itemize}
You can even display poetry.
\begin{verse}
   There is an environment
    for verse \\             % The \\ command separates lines
   Whose features some poets % within a stanza.
   will curse.

                             % One or more blank lines separate stanzas.

   For instead of making\\
   Them do \emph{all} line breaking, \\
   It allows them to put too many words on a line when they'd rather be
   forced to be terse.
\end{verse}

Mathematical formulas may also be displayed.  A
displayed formula
is
one-line long; multiline
formulas require special formatting instructions.
   \[  \ip{\Gamma}{\psi'} = x'' + y^{2} + z_{i}^{n}\]
Don't start a paragraph with a displayed equation,
nor make one a paragraph by itself.

\end{document}               % End of document.

\section{Bench Huge Section 883}
% synthetic bench section 883
% This is a sample LaTeX input file.  (Version of 12 August 2004.)
%
% A '%' character causes TeX to ignore all remaining text on the line,
% and is used for comments like this one.

\documentclass{article}      % Specifies the document class

                             % The preamble begins here.
\title{An Example Document}  % Declares the document's title.
\author{Leslie Lamport}      % Declares the author's name.
\date{January 21, 1994}      % Deleting this command produces today's date.

\newcommand{\ip}[2]{(#1, #2)}
                             % Defines \ip{arg1}{arg2} to mean
                             % (arg1, arg2).

%\newcommand{\ip}[2]{\langle #1 | #2\rangle}
                             % This is an alternative definition of
                             % \ip that is commented out.

\begin{document}             % End of preamble and beginning of text.

\maketitle                   % Produces the title.

This is an example input file.  Comparing it with
the output it generates can show you how to
produce a simple document of your own.

\section{Ordinary Text}      % Produces section heading.  Lower-level
                             % sections are begun with similar
                             % \subsection and \subsubsection commands.

The ends  of words and sentences are marked
  by   spaces. It  doesn't matter how many
spaces    you type; one is as good as 100.  The
end of   a line counts as a space.

One   or more   blank lines denote the  end
of  a paragraph.

Since any number of consecutive spaces are treated
like a single one, the formatting of the input
file makes no difference to
      \LaTeX,                % The \LaTeX command generates the LaTeX logo.
but it makes a difference to you.  When you use
\LaTeX, making your input file as easy to read
as possible will be a great help as you write
your document and when you change it.  This sample
file shows how you can add comments to your own input
file.

Because printing is different from typewriting,
there are a number of things that you have to do
differently when preparing an input file than if
you were just typing the document directly.
Quotation marks like
       ``this''
have to be handled specially, as do quotes within
quotes:
       ``\,`this'            % \, separates the double and single quote.
        is what I just
        wrote, not  `that'\,''.

Dashes come in three sizes: an
       intra-word
dash, a medium dash for number ranges like
       1--2,
and a punctuation
       dash---like
this.

A sentence-ending space should be larger than the
space between words within a sentence.  You
sometimes have to type special commands in
conjunction with punctuation characters to get
this right, as in the following sentence.
       Gnats, gnus, etc.\ all  % `\ ' makes an inter-word space.
       begin with G\@.         % \@ marks end-of-sentence punctuation.
You should check the spaces after periods when
reading your output to make sure you haven't
forgotten any special cases.  Generating an
ellipsis
       \ldots\               % `\ ' is needed after `\ldots' because TeX
                             % ignores spaces after command names like \ldots
                             % made from \ + letters.
                             %
                             % Note how a `%' character causes TeX to ignore
                             % the end of the input line, so these blank lines
                             % do not start a new paragraph.
                             %
with the right spacing around the periods requires
a special command.

\LaTeX\ interprets some common characters as
commands, so you must type special commands to
generate them.  These characters include the
following:
       \$ \& \% \# \{ and \}.

In printing, text is usually emphasized with an
       \emph{italic}
type style.

\begin{em}
   A long segment of text can also be emphasized
   in this way.  Text within such a segment can be
   given \emph{additional} emphasis.
\end{em}

It is sometimes necessary to prevent \LaTeX\ from
breaking a line where it might otherwise do so.
This may be at a space, as between the ``Mr.''\ and
``Jones'' in
       ``Mr.~Jones'',        % ~ produces an unbreakable interword space.
or within a word---especially when the word is a
symbol like
       \mbox{\emph{itemnum}}
that makes little sense when hyphenated across
lines.

Footnotes\footnote{This is an example of a footnote.}
pose no problem.

\LaTeX\ is good at typesetting mathematical formulas
like
       \( x-3y + z = 7 \)
or
       \( a_{1} > x^{2n} + y^{2n} > x' \)
or
       \( \ip{A}{B} = \sum_{i} a_{i} b_{i} \).
The spaces you type in a formula are
ignored.  Remember that a letter like
       $x$                   % $ ... $  and  \( ... \)  are equivalent
is a formula when it denotes a mathematical
symbol, and it should be typed as one.

\section{Displayed Text}

Text is displayed by indenting it from the left
margin.  Quotations are commonly displayed.  There
are short quotations
\begin{quote}
   This is a short quotation.  It consists of a
   single paragraph of text.  See how it is formatted.
\end{quote}
and longer ones.
\begin{quotation}
   This is a longer quotation.  It consists of two
   paragraphs of text, neither of which are
   particularly interesting.

   This is the second paragraph of the quotation.  It
   is just as dull as the first paragraph.
\end{quotation}
Another frequently-displayed structure is a list.
The following is an example of an \emph{itemized}
list.
\begin{itemize}
   \item This is the first item of an itemized list.
         Each item in the list is marked with a ``tick''.
         You don't have to worry about what kind of tick
         mark is used.

   \item This is the second item of the list.  It
         contains another list nested inside it.  The inner
         list is an \emph{enumerated} list.
         \begin{enumerate}
            \item This is the first item of an enumerated
                  list that is nested within the itemized list.

            \item This is the second item of the inner list.
                  \LaTeX\ allows you to nest lists deeper than
                  you really should.
         \end{enumerate}
         This is the rest of the second item of the outer
         list.  It is no more interesting than any other
         part of the item.
   \item This is the third item of the list.
\end{itemize}
You can even display poetry.
\begin{verse}
   There is an environment
    for verse \\             % The \\ command separates lines
   Whose features some poets % within a stanza.
   will curse.

                             % One or more blank lines separate stanzas.

   For instead of making\\
   Them do \emph{all} line breaking, \\
   It allows them to put too many words on a line when they'd rather be
   forced to be terse.
\end{verse}

Mathematical formulas may also be displayed.  A
displayed formula
is
one-line long; multiline
formulas require special formatting instructions.
   \[  \ip{\Gamma}{\psi'} = x'' + y^{2} + z_{i}^{n}\]
Don't start a paragraph with a displayed equation,
nor make one a paragraph by itself.

\end{document}               % End of document.

\section{Bench Huge Section 884}
% synthetic bench section 884
% This is a sample LaTeX input file.  (Version of 12 August 2004.)
%
% A '%' character causes TeX to ignore all remaining text on the line,
% and is used for comments like this one.

\documentclass{article}      % Specifies the document class

                             % The preamble begins here.
\title{An Example Document}  % Declares the document's title.
\author{Leslie Lamport}      % Declares the author's name.
\date{January 21, 1994}      % Deleting this command produces today's date.

\newcommand{\ip}[2]{(#1, #2)}
                             % Defines \ip{arg1}{arg2} to mean
                             % (arg1, arg2).

%\newcommand{\ip}[2]{\langle #1 | #2\rangle}
                             % This is an alternative definition of
                             % \ip that is commented out.

\begin{document}             % End of preamble and beginning of text.

\maketitle                   % Produces the title.

This is an example input file.  Comparing it with
the output it generates can show you how to
produce a simple document of your own.

\section{Ordinary Text}      % Produces section heading.  Lower-level
                             % sections are begun with similar
                             % \subsection and \subsubsection commands.

The ends  of words and sentences are marked
  by   spaces. It  doesn't matter how many
spaces    you type; one is as good as 100.  The
end of   a line counts as a space.

One   or more   blank lines denote the  end
of  a paragraph.

Since any number of consecutive spaces are treated
like a single one, the formatting of the input
file makes no difference to
      \LaTeX,                % The \LaTeX command generates the LaTeX logo.
but it makes a difference to you.  When you use
\LaTeX, making your input file as easy to read
as possible will be a great help as you write
your document and when you change it.  This sample
file shows how you can add comments to your own input
file.

Because printing is different from typewriting,
there are a number of things that you have to do
differently when preparing an input file than if
you were just typing the document directly.
Quotation marks like
       ``this''
have to be handled specially, as do quotes within
quotes:
       ``\,`this'            % \, separates the double and single quote.
        is what I just
        wrote, not  `that'\,''.

Dashes come in three sizes: an
       intra-word
dash, a medium dash for number ranges like
       1--2,
and a punctuation
       dash---like
this.

A sentence-ending space should be larger than the
space between words within a sentence.  You
sometimes have to type special commands in
conjunction with punctuation characters to get
this right, as in the following sentence.
       Gnats, gnus, etc.\ all  % `\ ' makes an inter-word space.
       begin with G\@.         % \@ marks end-of-sentence punctuation.
You should check the spaces after periods when
reading your output to make sure you haven't
forgotten any special cases.  Generating an
ellipsis
       \ldots\               % `\ ' is needed after `\ldots' because TeX
                             % ignores spaces after command names like \ldots
                             % made from \ + letters.
                             %
                             % Note how a `%' character causes TeX to ignore
                             % the end of the input line, so these blank lines
                             % do not start a new paragraph.
                             %
with the right spacing around the periods requires
a special command.

\LaTeX\ interprets some common characters as
commands, so you must type special commands to
generate them.  These characters include the
following:
       \$ \& \% \# \{ and \}.

In printing, text is usually emphasized with an
       \emph{italic}
type style.

\begin{em}
   A long segment of text can also be emphasized
   in this way.  Text within such a segment can be
   given \emph{additional} emphasis.
\end{em}

It is sometimes necessary to prevent \LaTeX\ from
breaking a line where it might otherwise do so.
This may be at a space, as between the ``Mr.''\ and
``Jones'' in
       ``Mr.~Jones'',        % ~ produces an unbreakable interword space.
or within a word---especially when the word is a
symbol like
       \mbox{\emph{itemnum}}
that makes little sense when hyphenated across
lines.

Footnotes\footnote{This is an example of a footnote.}
pose no problem.

\LaTeX\ is good at typesetting mathematical formulas
like
       \( x-3y + z = 7 \)
or
       \( a_{1} > x^{2n} + y^{2n} > x' \)
or
       \( \ip{A}{B} = \sum_{i} a_{i} b_{i} \).
The spaces you type in a formula are
ignored.  Remember that a letter like
       $x$                   % $ ... $  and  \( ... \)  are equivalent
is a formula when it denotes a mathematical
symbol, and it should be typed as one.

\section{Displayed Text}

Text is displayed by indenting it from the left
margin.  Quotations are commonly displayed.  There
are short quotations
\begin{quote}
   This is a short quotation.  It consists of a
   single paragraph of text.  See how it is formatted.
\end{quote}
and longer ones.
\begin{quotation}
   This is a longer quotation.  It consists of two
   paragraphs of text, neither of which are
   particularly interesting.

   This is the second paragraph of the quotation.  It
   is just as dull as the first paragraph.
\end{quotation}
Another frequently-displayed structure is a list.
The following is an example of an \emph{itemized}
list.
\begin{itemize}
   \item This is the first item of an itemized list.
         Each item in the list is marked with a ``tick''.
         You don't have to worry about what kind of tick
         mark is used.

   \item This is the second item of the list.  It
         contains another list nested inside it.  The inner
         list is an \emph{enumerated} list.
         \begin{enumerate}
            \item This is the first item of an enumerated
                  list that is nested within the itemized list.

            \item This is the second item of the inner list.
                  \LaTeX\ allows you to nest lists deeper than
                  you really should.
         \end{enumerate}
         This is the rest of the second item of the outer
         list.  It is no more interesting than any other
         part of the item.
   \item This is the third item of the list.
\end{itemize}
You can even display poetry.
\begin{verse}
   There is an environment
    for verse \\             % The \\ command separates lines
   Whose features some poets % within a stanza.
   will curse.

                             % One or more blank lines separate stanzas.

   For instead of making\\
   Them do \emph{all} line breaking, \\
   It allows them to put too many words on a line when they'd rather be
   forced to be terse.
\end{verse}

Mathematical formulas may also be displayed.  A
displayed formula
is
one-line long; multiline
formulas require special formatting instructions.
   \[  \ip{\Gamma}{\psi'} = x'' + y^{2} + z_{i}^{n}\]
Don't start a paragraph with a displayed equation,
nor make one a paragraph by itself.

\end{document}               % End of document.

\section{Bench Huge Section 885}
% synthetic bench section 885
% This is a sample LaTeX input file.  (Version of 12 August 2004.)
%
% A '%' character causes TeX to ignore all remaining text on the line,
% and is used for comments like this one.

\documentclass{article}      % Specifies the document class

                             % The preamble begins here.
\title{An Example Document}  % Declares the document's title.
\author{Leslie Lamport}      % Declares the author's name.
\date{January 21, 1994}      % Deleting this command produces today's date.

\newcommand{\ip}[2]{(#1, #2)}
                             % Defines \ip{arg1}{arg2} to mean
                             % (arg1, arg2).

%\newcommand{\ip}[2]{\langle #1 | #2\rangle}
                             % This is an alternative definition of
                             % \ip that is commented out.

\begin{document}             % End of preamble and beginning of text.

\maketitle                   % Produces the title.

This is an example input file.  Comparing it with
the output it generates can show you how to
produce a simple document of your own.

\section{Ordinary Text}      % Produces section heading.  Lower-level
                             % sections are begun with similar
                             % \subsection and \subsubsection commands.

The ends  of words and sentences are marked
  by   spaces. It  doesn't matter how many
spaces    you type; one is as good as 100.  The
end of   a line counts as a space.

One   or more   blank lines denote the  end
of  a paragraph.

Since any number of consecutive spaces are treated
like a single one, the formatting of the input
file makes no difference to
      \LaTeX,                % The \LaTeX command generates the LaTeX logo.
but it makes a difference to you.  When you use
\LaTeX, making your input file as easy to read
as possible will be a great help as you write
your document and when you change it.  This sample
file shows how you can add comments to your own input
file.

Because printing is different from typewriting,
there are a number of things that you have to do
differently when preparing an input file than if
you were just typing the document directly.
Quotation marks like
       ``this''
have to be handled specially, as do quotes within
quotes:
       ``\,`this'            % \, separates the double and single quote.
        is what I just
        wrote, not  `that'\,''.

Dashes come in three sizes: an
       intra-word
dash, a medium dash for number ranges like
       1--2,
and a punctuation
       dash---like
this.

A sentence-ending space should be larger than the
space between words within a sentence.  You
sometimes have to type special commands in
conjunction with punctuation characters to get
this right, as in the following sentence.
       Gnats, gnus, etc.\ all  % `\ ' makes an inter-word space.
       begin with G\@.         % \@ marks end-of-sentence punctuation.
You should check the spaces after periods when
reading your output to make sure you haven't
forgotten any special cases.  Generating an
ellipsis
       \ldots\               % `\ ' is needed after `\ldots' because TeX
                             % ignores spaces after command names like \ldots
                             % made from \ + letters.
                             %
                             % Note how a `%' character causes TeX to ignore
                             % the end of the input line, so these blank lines
                             % do not start a new paragraph.
                             %
with the right spacing around the periods requires
a special command.

\LaTeX\ interprets some common characters as
commands, so you must type special commands to
generate them.  These characters include the
following:
       \$ \& \% \# \{ and \}.

In printing, text is usually emphasized with an
       \emph{italic}
type style.

\begin{em}
   A long segment of text can also be emphasized
   in this way.  Text within such a segment can be
   given \emph{additional} emphasis.
\end{em}

It is sometimes necessary to prevent \LaTeX\ from
breaking a line where it might otherwise do so.
This may be at a space, as between the ``Mr.''\ and
``Jones'' in
       ``Mr.~Jones'',        % ~ produces an unbreakable interword space.
or within a word---especially when the word is a
symbol like
       \mbox{\emph{itemnum}}
that makes little sense when hyphenated across
lines.

Footnotes\footnote{This is an example of a footnote.}
pose no problem.

\LaTeX\ is good at typesetting mathematical formulas
like
       \( x-3y + z = 7 \)
or
       \( a_{1} > x^{2n} + y^{2n} > x' \)
or
       \( \ip{A}{B} = \sum_{i} a_{i} b_{i} \).
The spaces you type in a formula are
ignored.  Remember that a letter like
       $x$                   % $ ... $  and  \( ... \)  are equivalent
is a formula when it denotes a mathematical
symbol, and it should be typed as one.

\section{Displayed Text}

Text is displayed by indenting it from the left
margin.  Quotations are commonly displayed.  There
are short quotations
\begin{quote}
   This is a short quotation.  It consists of a
   single paragraph of text.  See how it is formatted.
\end{quote}
and longer ones.
\begin{quotation}
   This is a longer quotation.  It consists of two
   paragraphs of text, neither of which are
   particularly interesting.

   This is the second paragraph of the quotation.  It
   is just as dull as the first paragraph.
\end{quotation}
Another frequently-displayed structure is a list.
The following is an example of an \emph{itemized}
list.
\begin{itemize}
   \item This is the first item of an itemized list.
         Each item in the list is marked with a ``tick''.
         You don't have to worry about what kind of tick
         mark is used.

   \item This is the second item of the list.  It
         contains another list nested inside it.  The inner
         list is an \emph{enumerated} list.
         \begin{enumerate}
            \item This is the first item of an enumerated
                  list that is nested within the itemized list.

            \item This is the second item of the inner list.
                  \LaTeX\ allows you to nest lists deeper than
                  you really should.
         \end{enumerate}
         This is the rest of the second item of the outer
         list.  It is no more interesting than any other
         part of the item.
   \item This is the third item of the list.
\end{itemize}
You can even display poetry.
\begin{verse}
   There is an environment
    for verse \\             % The \\ command separates lines
   Whose features some poets % within a stanza.
   will curse.

                             % One or more blank lines separate stanzas.

   For instead of making\\
   Them do \emph{all} line breaking, \\
   It allows them to put too many words on a line when they'd rather be
   forced to be terse.
\end{verse}

Mathematical formulas may also be displayed.  A
displayed formula
is
one-line long; multiline
formulas require special formatting instructions.
   \[  \ip{\Gamma}{\psi'} = x'' + y^{2} + z_{i}^{n}\]
Don't start a paragraph with a displayed equation,
nor make one a paragraph by itself.

\end{document}               % End of document.

\section{Bench Huge Section 886}
% synthetic bench section 886
% This is a sample LaTeX input file.  (Version of 12 August 2004.)
%
% A '%' character causes TeX to ignore all remaining text on the line,
% and is used for comments like this one.

\documentclass{article}      % Specifies the document class

                             % The preamble begins here.
\title{An Example Document}  % Declares the document's title.
\author{Leslie Lamport}      % Declares the author's name.
\date{January 21, 1994}      % Deleting this command produces today's date.

\newcommand{\ip}[2]{(#1, #2)}
                             % Defines \ip{arg1}{arg2} to mean
                             % (arg1, arg2).

%\newcommand{\ip}[2]{\langle #1 | #2\rangle}
                             % This is an alternative definition of
                             % \ip that is commented out.

\begin{document}             % End of preamble and beginning of text.

\maketitle                   % Produces the title.

This is an example input file.  Comparing it with
the output it generates can show you how to
produce a simple document of your own.

\section{Ordinary Text}      % Produces section heading.  Lower-level
                             % sections are begun with similar
                             % \subsection and \subsubsection commands.

The ends  of words and sentences are marked
  by   spaces. It  doesn't matter how many
spaces    you type; one is as good as 100.  The
end of   a line counts as a space.

One   or more   blank lines denote the  end
of  a paragraph.

Since any number of consecutive spaces are treated
like a single one, the formatting of the input
file makes no difference to
      \LaTeX,                % The \LaTeX command generates the LaTeX logo.
but it makes a difference to you.  When you use
\LaTeX, making your input file as easy to read
as possible will be a great help as you write
your document and when you change it.  This sample
file shows how you can add comments to your own input
file.

Because printing is different from typewriting,
there are a number of things that you have to do
differently when preparing an input file than if
you were just typing the document directly.
Quotation marks like
       ``this''
have to be handled specially, as do quotes within
quotes:
       ``\,`this'            % \, separates the double and single quote.
        is what I just
        wrote, not  `that'\,''.

Dashes come in three sizes: an
       intra-word
dash, a medium dash for number ranges like
       1--2,
and a punctuation
       dash---like
this.

A sentence-ending space should be larger than the
space between words within a sentence.  You
sometimes have to type special commands in
conjunction with punctuation characters to get
this right, as in the following sentence.
       Gnats, gnus, etc.\ all  % `\ ' makes an inter-word space.
       begin with G\@.         % \@ marks end-of-sentence punctuation.
You should check the spaces after periods when
reading your output to make sure you haven't
forgotten any special cases.  Generating an
ellipsis
       \ldots\               % `\ ' is needed after `\ldots' because TeX
                             % ignores spaces after command names like \ldots
                             % made from \ + letters.
                             %
                             % Note how a `%' character causes TeX to ignore
                             % the end of the input line, so these blank lines
                             % do not start a new paragraph.
                             %
with the right spacing around the periods requires
a special command.

\LaTeX\ interprets some common characters as
commands, so you must type special commands to
generate them.  These characters include the
following:
       \$ \& \% \# \{ and \}.

In printing, text is usually emphasized with an
       \emph{italic}
type style.

\begin{em}
   A long segment of text can also be emphasized
   in this way.  Text within such a segment can be
   given \emph{additional} emphasis.
\end{em}

It is sometimes necessary to prevent \LaTeX\ from
breaking a line where it might otherwise do so.
This may be at a space, as between the ``Mr.''\ and
``Jones'' in
       ``Mr.~Jones'',        % ~ produces an unbreakable interword space.
or within a word---especially when the word is a
symbol like
       \mbox{\emph{itemnum}}
that makes little sense when hyphenated across
lines.

Footnotes\footnote{This is an example of a footnote.}
pose no problem.

\LaTeX\ is good at typesetting mathematical formulas
like
       \( x-3y + z = 7 \)
or
       \( a_{1} > x^{2n} + y^{2n} > x' \)
or
       \( \ip{A}{B} = \sum_{i} a_{i} b_{i} \).
The spaces you type in a formula are
ignored.  Remember that a letter like
       $x$                   % $ ... $  and  \( ... \)  are equivalent
is a formula when it denotes a mathematical
symbol, and it should be typed as one.

\section{Displayed Text}

Text is displayed by indenting it from the left
margin.  Quotations are commonly displayed.  There
are short quotations
\begin{quote}
   This is a short quotation.  It consists of a
   single paragraph of text.  See how it is formatted.
\end{quote}
and longer ones.
\begin{quotation}
   This is a longer quotation.  It consists of two
   paragraphs of text, neither of which are
   particularly interesting.

   This is the second paragraph of the quotation.  It
   is just as dull as the first paragraph.
\end{quotation}
Another frequently-displayed structure is a list.
The following is an example of an \emph{itemized}
list.
\begin{itemize}
   \item This is the first item of an itemized list.
         Each item in the list is marked with a ``tick''.
         You don't have to worry about what kind of tick
         mark is used.

   \item This is the second item of the list.  It
         contains another list nested inside it.  The inner
         list is an \emph{enumerated} list.
         \begin{enumerate}
            \item This is the first item of an enumerated
                  list that is nested within the itemized list.

            \item This is the second item of the inner list.
                  \LaTeX\ allows you to nest lists deeper than
                  you really should.
         \end{enumerate}
         This is the rest of the second item of the outer
         list.  It is no more interesting than any other
         part of the item.
   \item This is the third item of the list.
\end{itemize}
You can even display poetry.
\begin{verse}
   There is an environment
    for verse \\             % The \\ command separates lines
   Whose features some poets % within a stanza.
   will curse.

                             % One or more blank lines separate stanzas.

   For instead of making\\
   Them do \emph{all} line breaking, \\
   It allows them to put too many words on a line when they'd rather be
   forced to be terse.
\end{verse}

Mathematical formulas may also be displayed.  A
displayed formula
is
one-line long; multiline
formulas require special formatting instructions.
   \[  \ip{\Gamma}{\psi'} = x'' + y^{2} + z_{i}^{n}\]
Don't start a paragraph with a displayed equation,
nor make one a paragraph by itself.

\end{document}               % End of document.

\section{Bench Huge Section 887}
% synthetic bench section 887
% This is a sample LaTeX input file.  (Version of 12 August 2004.)
%
% A '%' character causes TeX to ignore all remaining text on the line,
% and is used for comments like this one.

\documentclass{article}      % Specifies the document class

                             % The preamble begins here.
\title{An Example Document}  % Declares the document's title.
\author{Leslie Lamport}      % Declares the author's name.
\date{January 21, 1994}      % Deleting this command produces today's date.

\newcommand{\ip}[2]{(#1, #2)}
                             % Defines \ip{arg1}{arg2} to mean
                             % (arg1, arg2).

%\newcommand{\ip}[2]{\langle #1 | #2\rangle}
                             % This is an alternative definition of
                             % \ip that is commented out.

\begin{document}             % End of preamble and beginning of text.

\maketitle                   % Produces the title.

This is an example input file.  Comparing it with
the output it generates can show you how to
produce a simple document of your own.

\section{Ordinary Text}      % Produces section heading.  Lower-level
                             % sections are begun with similar
                             % \subsection and \subsubsection commands.

The ends  of words and sentences are marked
  by   spaces. It  doesn't matter how many
spaces    you type; one is as good as 100.  The
end of   a line counts as a space.

One   or more   blank lines denote the  end
of  a paragraph.

Since any number of consecutive spaces are treated
like a single one, the formatting of the input
file makes no difference to
      \LaTeX,                % The \LaTeX command generates the LaTeX logo.
but it makes a difference to you.  When you use
\LaTeX, making your input file as easy to read
as possible will be a great help as you write
your document and when you change it.  This sample
file shows how you can add comments to your own input
file.

Because printing is different from typewriting,
there are a number of things that you have to do
differently when preparing an input file than if
you were just typing the document directly.
Quotation marks like
       ``this''
have to be handled specially, as do quotes within
quotes:
       ``\,`this'            % \, separates the double and single quote.
        is what I just
        wrote, not  `that'\,''.

Dashes come in three sizes: an
       intra-word
dash, a medium dash for number ranges like
       1--2,
and a punctuation
       dash---like
this.

A sentence-ending space should be larger than the
space between words within a sentence.  You
sometimes have to type special commands in
conjunction with punctuation characters to get
this right, as in the following sentence.
       Gnats, gnus, etc.\ all  % `\ ' makes an inter-word space.
       begin with G\@.         % \@ marks end-of-sentence punctuation.
You should check the spaces after periods when
reading your output to make sure you haven't
forgotten any special cases.  Generating an
ellipsis
       \ldots\               % `\ ' is needed after `\ldots' because TeX
                             % ignores spaces after command names like \ldots
                             % made from \ + letters.
                             %
                             % Note how a `%' character causes TeX to ignore
                             % the end of the input line, so these blank lines
                             % do not start a new paragraph.
                             %
with the right spacing around the periods requires
a special command.

\LaTeX\ interprets some common characters as
commands, so you must type special commands to
generate them.  These characters include the
following:
       \$ \& \% \# \{ and \}.

In printing, text is usually emphasized with an
       \emph{italic}
type style.

\begin{em}
   A long segment of text can also be emphasized
   in this way.  Text within such a segment can be
   given \emph{additional} emphasis.
\end{em}

It is sometimes necessary to prevent \LaTeX\ from
breaking a line where it might otherwise do so.
This may be at a space, as between the ``Mr.''\ and
``Jones'' in
       ``Mr.~Jones'',        % ~ produces an unbreakable interword space.
or within a word---especially when the word is a
symbol like
       \mbox{\emph{itemnum}}
that makes little sense when hyphenated across
lines.

Footnotes\footnote{This is an example of a footnote.}
pose no problem.

\LaTeX\ is good at typesetting mathematical formulas
like
       \( x-3y + z = 7 \)
or
       \( a_{1} > x^{2n} + y^{2n} > x' \)
or
       \( \ip{A}{B} = \sum_{i} a_{i} b_{i} \).
The spaces you type in a formula are
ignored.  Remember that a letter like
       $x$                   % $ ... $  and  \( ... \)  are equivalent
is a formula when it denotes a mathematical
symbol, and it should be typed as one.

\section{Displayed Text}

Text is displayed by indenting it from the left
margin.  Quotations are commonly displayed.  There
are short quotations
\begin{quote}
   This is a short quotation.  It consists of a
   single paragraph of text.  See how it is formatted.
\end{quote}
and longer ones.
\begin{quotation}
   This is a longer quotation.  It consists of two
   paragraphs of text, neither of which are
   particularly interesting.

   This is the second paragraph of the quotation.  It
   is just as dull as the first paragraph.
\end{quotation}
Another frequently-displayed structure is a list.
The following is an example of an \emph{itemized}
list.
\begin{itemize}
   \item This is the first item of an itemized list.
         Each item in the list is marked with a ``tick''.
         You don't have to worry about what kind of tick
         mark is used.

   \item This is the second item of the list.  It
         contains another list nested inside it.  The inner
         list is an \emph{enumerated} list.
         \begin{enumerate}
            \item This is the first item of an enumerated
                  list that is nested within the itemized list.

            \item This is the second item of the inner list.
                  \LaTeX\ allows you to nest lists deeper than
                  you really should.
         \end{enumerate}
         This is the rest of the second item of the outer
         list.  It is no more interesting than any other
         part of the item.
   \item This is the third item of the list.
\end{itemize}
You can even display poetry.
\begin{verse}
   There is an environment
    for verse \\             % The \\ command separates lines
   Whose features some poets % within a stanza.
   will curse.

                             % One or more blank lines separate stanzas.

   For instead of making\\
   Them do \emph{all} line breaking, \\
   It allows them to put too many words on a line when they'd rather be
   forced to be terse.
\end{verse}

Mathematical formulas may also be displayed.  A
displayed formula
is
one-line long; multiline
formulas require special formatting instructions.
   \[  \ip{\Gamma}{\psi'} = x'' + y^{2} + z_{i}^{n}\]
Don't start a paragraph with a displayed equation,
nor make one a paragraph by itself.

\end{document}               % End of document.

\section{Bench Huge Section 888}
% synthetic bench section 888
% This is a sample LaTeX input file.  (Version of 12 August 2004.)
%
% A '%' character causes TeX to ignore all remaining text on the line,
% and is used for comments like this one.

\documentclass{article}      % Specifies the document class

                             % The preamble begins here.
\title{An Example Document}  % Declares the document's title.
\author{Leslie Lamport}      % Declares the author's name.
\date{January 21, 1994}      % Deleting this command produces today's date.

\newcommand{\ip}[2]{(#1, #2)}
                             % Defines \ip{arg1}{arg2} to mean
                             % (arg1, arg2).

%\newcommand{\ip}[2]{\langle #1 | #2\rangle}
                             % This is an alternative definition of
                             % \ip that is commented out.

\begin{document}             % End of preamble and beginning of text.

\maketitle                   % Produces the title.

This is an example input file.  Comparing it with
the output it generates can show you how to
produce a simple document of your own.

\section{Ordinary Text}      % Produces section heading.  Lower-level
                             % sections are begun with similar
                             % \subsection and \subsubsection commands.

The ends  of words and sentences are marked
  by   spaces. It  doesn't matter how many
spaces    you type; one is as good as 100.  The
end of   a line counts as a space.

One   or more   blank lines denote the  end
of  a paragraph.

Since any number of consecutive spaces are treated
like a single one, the formatting of the input
file makes no difference to
      \LaTeX,                % The \LaTeX command generates the LaTeX logo.
but it makes a difference to you.  When you use
\LaTeX, making your input file as easy to read
as possible will be a great help as you write
your document and when you change it.  This sample
file shows how you can add comments to your own input
file.

Because printing is different from typewriting,
there are a number of things that you have to do
differently when preparing an input file than if
you were just typing the document directly.
Quotation marks like
       ``this''
have to be handled specially, as do quotes within
quotes:
       ``\,`this'            % \, separates the double and single quote.
        is what I just
        wrote, not  `that'\,''.

Dashes come in three sizes: an
       intra-word
dash, a medium dash for number ranges like
       1--2,
and a punctuation
       dash---like
this.

A sentence-ending space should be larger than the
space between words within a sentence.  You
sometimes have to type special commands in
conjunction with punctuation characters to get
this right, as in the following sentence.
       Gnats, gnus, etc.\ all  % `\ ' makes an inter-word space.
       begin with G\@.         % \@ marks end-of-sentence punctuation.
You should check the spaces after periods when
reading your output to make sure you haven't
forgotten any special cases.  Generating an
ellipsis
       \ldots\               % `\ ' is needed after `\ldots' because TeX
                             % ignores spaces after command names like \ldots
                             % made from \ + letters.
                             %
                             % Note how a `%' character causes TeX to ignore
                             % the end of the input line, so these blank lines
                             % do not start a new paragraph.
                             %
with the right spacing around the periods requires
a special command.

\LaTeX\ interprets some common characters as
commands, so you must type special commands to
generate them.  These characters include the
following:
       \$ \& \% \# \{ and \}.

In printing, text is usually emphasized with an
       \emph{italic}
type style.

\begin{em}
   A long segment of text can also be emphasized
   in this way.  Text within such a segment can be
   given \emph{additional} emphasis.
\end{em}

It is sometimes necessary to prevent \LaTeX\ from
breaking a line where it might otherwise do so.
This may be at a space, as between the ``Mr.''\ and
``Jones'' in
       ``Mr.~Jones'',        % ~ produces an unbreakable interword space.
or within a word---especially when the word is a
symbol like
       \mbox{\emph{itemnum}}
that makes little sense when hyphenated across
lines.

Footnotes\footnote{This is an example of a footnote.}
pose no problem.

\LaTeX\ is good at typesetting mathematical formulas
like
       \( x-3y + z = 7 \)
or
       \( a_{1} > x^{2n} + y^{2n} > x' \)
or
       \( \ip{A}{B} = \sum_{i} a_{i} b_{i} \).
The spaces you type in a formula are
ignored.  Remember that a letter like
       $x$                   % $ ... $  and  \( ... \)  are equivalent
is a formula when it denotes a mathematical
symbol, and it should be typed as one.

\section{Displayed Text}

Text is displayed by indenting it from the left
margin.  Quotations are commonly displayed.  There
are short quotations
\begin{quote}
   This is a short quotation.  It consists of a
   single paragraph of text.  See how it is formatted.
\end{quote}
and longer ones.
\begin{quotation}
   This is a longer quotation.  It consists of two
   paragraphs of text, neither of which are
   particularly interesting.

   This is the second paragraph of the quotation.  It
   is just as dull as the first paragraph.
\end{quotation}
Another frequently-displayed structure is a list.
The following is an example of an \emph{itemized}
list.
\begin{itemize}
   \item This is the first item of an itemized list.
         Each item in the list is marked with a ``tick''.
         You don't have to worry about what kind of tick
         mark is used.

   \item This is the second item of the list.  It
         contains another list nested inside it.  The inner
         list is an \emph{enumerated} list.
         \begin{enumerate}
            \item This is the first item of an enumerated
                  list that is nested within the itemized list.

            \item This is the second item of the inner list.
                  \LaTeX\ allows you to nest lists deeper than
                  you really should.
         \end{enumerate}
         This is the rest of the second item of the outer
         list.  It is no more interesting than any other
         part of the item.
   \item This is the third item of the list.
\end{itemize}
You can even display poetry.
\begin{verse}
   There is an environment
    for verse \\             % The \\ command separates lines
   Whose features some poets % within a stanza.
   will curse.

                             % One or more blank lines separate stanzas.

   For instead of making\\
   Them do \emph{all} line breaking, \\
   It allows them to put too many words on a line when they'd rather be
   forced to be terse.
\end{verse}

Mathematical formulas may also be displayed.  A
displayed formula
is
one-line long; multiline
formulas require special formatting instructions.
   \[  \ip{\Gamma}{\psi'} = x'' + y^{2} + z_{i}^{n}\]
Don't start a paragraph with a displayed equation,
nor make one a paragraph by itself.

\end{document}               % End of document.

\section{Bench Huge Section 889}
% synthetic bench section 889
% This is a sample LaTeX input file.  (Version of 12 August 2004.)
%
% A '%' character causes TeX to ignore all remaining text on the line,
% and is used for comments like this one.

\documentclass{article}      % Specifies the document class

                             % The preamble begins here.
\title{An Example Document}  % Declares the document's title.
\author{Leslie Lamport}      % Declares the author's name.
\date{January 21, 1994}      % Deleting this command produces today's date.

\newcommand{\ip}[2]{(#1, #2)}
                             % Defines \ip{arg1}{arg2} to mean
                             % (arg1, arg2).

%\newcommand{\ip}[2]{\langle #1 | #2\rangle}
                             % This is an alternative definition of
                             % \ip that is commented out.

\begin{document}             % End of preamble and beginning of text.

\maketitle                   % Produces the title.

This is an example input file.  Comparing it with
the output it generates can show you how to
produce a simple document of your own.

\section{Ordinary Text}      % Produces section heading.  Lower-level
                             % sections are begun with similar
                             % \subsection and \subsubsection commands.

The ends  of words and sentences are marked
  by   spaces. It  doesn't matter how many
spaces    you type; one is as good as 100.  The
end of   a line counts as a space.

One   or more   blank lines denote the  end
of  a paragraph.

Since any number of consecutive spaces are treated
like a single one, the formatting of the input
file makes no difference to
      \LaTeX,                % The \LaTeX command generates the LaTeX logo.
but it makes a difference to you.  When you use
\LaTeX, making your input file as easy to read
as possible will be a great help as you write
your document and when you change it.  This sample
file shows how you can add comments to your own input
file.

Because printing is different from typewriting,
there are a number of things that you have to do
differently when preparing an input file than if
you were just typing the document directly.
Quotation marks like
       ``this''
have to be handled specially, as do quotes within
quotes:
       ``\,`this'            % \, separates the double and single quote.
        is what I just
        wrote, not  `that'\,''.

Dashes come in three sizes: an
       intra-word
dash, a medium dash for number ranges like
       1--2,
and a punctuation
       dash---like
this.

A sentence-ending space should be larger than the
space between words within a sentence.  You
sometimes have to type special commands in
conjunction with punctuation characters to get
this right, as in the following sentence.
       Gnats, gnus, etc.\ all  % `\ ' makes an inter-word space.
       begin with G\@.         % \@ marks end-of-sentence punctuation.
You should check the spaces after periods when
reading your output to make sure you haven't
forgotten any special cases.  Generating an
ellipsis
       \ldots\               % `\ ' is needed after `\ldots' because TeX
                             % ignores spaces after command names like \ldots
                             % made from \ + letters.
                             %
                             % Note how a `%' character causes TeX to ignore
                             % the end of the input line, so these blank lines
                             % do not start a new paragraph.
                             %
with the right spacing around the periods requires
a special command.

\LaTeX\ interprets some common characters as
commands, so you must type special commands to
generate them.  These characters include the
following:
       \$ \& \% \# \{ and \}.

In printing, text is usually emphasized with an
       \emph{italic}
type style.

\begin{em}
   A long segment of text can also be emphasized
   in this way.  Text within such a segment can be
   given \emph{additional} emphasis.
\end{em}

It is sometimes necessary to prevent \LaTeX\ from
breaking a line where it might otherwise do so.
This may be at a space, as between the ``Mr.''\ and
``Jones'' in
       ``Mr.~Jones'',        % ~ produces an unbreakable interword space.
or within a word---especially when the word is a
symbol like
       \mbox{\emph{itemnum}}
that makes little sense when hyphenated across
lines.

Footnotes\footnote{This is an example of a footnote.}
pose no problem.

\LaTeX\ is good at typesetting mathematical formulas
like
       \( x-3y + z = 7 \)
or
       \( a_{1} > x^{2n} + y^{2n} > x' \)
or
       \( \ip{A}{B} = \sum_{i} a_{i} b_{i} \).
The spaces you type in a formula are
ignored.  Remember that a letter like
       $x$                   % $ ... $  and  \( ... \)  are equivalent
is a formula when it denotes a mathematical
symbol, and it should be typed as one.

\section{Displayed Text}

Text is displayed by indenting it from the left
margin.  Quotations are commonly displayed.  There
are short quotations
\begin{quote}
   This is a short quotation.  It consists of a
   single paragraph of text.  See how it is formatted.
\end{quote}
and longer ones.
\begin{quotation}
   This is a longer quotation.  It consists of two
   paragraphs of text, neither of which are
   particularly interesting.

   This is the second paragraph of the quotation.  It
   is just as dull as the first paragraph.
\end{quotation}
Another frequently-displayed structure is a list.
The following is an example of an \emph{itemized}
list.
\begin{itemize}
   \item This is the first item of an itemized list.
         Each item in the list is marked with a ``tick''.
         You don't have to worry about what kind of tick
         mark is used.

   \item This is the second item of the list.  It
         contains another list nested inside it.  The inner
         list is an \emph{enumerated} list.
         \begin{enumerate}
            \item This is the first item of an enumerated
                  list that is nested within the itemized list.

            \item This is the second item of the inner list.
                  \LaTeX\ allows you to nest lists deeper than
                  you really should.
         \end{enumerate}
         This is the rest of the second item of the outer
         list.  It is no more interesting than any other
         part of the item.
   \item This is the third item of the list.
\end{itemize}
You can even display poetry.
\begin{verse}
   There is an environment
    for verse \\             % The \\ command separates lines
   Whose features some poets % within a stanza.
   will curse.

                             % One or more blank lines separate stanzas.

   For instead of making\\
   Them do \emph{all} line breaking, \\
   It allows them to put too many words on a line when they'd rather be
   forced to be terse.
\end{verse}

Mathematical formulas may also be displayed.  A
displayed formula
is
one-line long; multiline
formulas require special formatting instructions.
   \[  \ip{\Gamma}{\psi'} = x'' + y^{2} + z_{i}^{n}\]
Don't start a paragraph with a displayed equation,
nor make one a paragraph by itself.

\end{document}               % End of document.

\section{Bench Huge Section 890}
% synthetic bench section 890
% This is a sample LaTeX input file.  (Version of 12 August 2004.)
%
% A '%' character causes TeX to ignore all remaining text on the line,
% and is used for comments like this one.

\documentclass{article}      % Specifies the document class

                             % The preamble begins here.
\title{An Example Document}  % Declares the document's title.
\author{Leslie Lamport}      % Declares the author's name.
\date{January 21, 1994}      % Deleting this command produces today's date.

\newcommand{\ip}[2]{(#1, #2)}
                             % Defines \ip{arg1}{arg2} to mean
                             % (arg1, arg2).

%\newcommand{\ip}[2]{\langle #1 | #2\rangle}
                             % This is an alternative definition of
                             % \ip that is commented out.

\begin{document}             % End of preamble and beginning of text.

\maketitle                   % Produces the title.

This is an example input file.  Comparing it with
the output it generates can show you how to
produce a simple document of your own.

\section{Ordinary Text}      % Produces section heading.  Lower-level
                             % sections are begun with similar
                             % \subsection and \subsubsection commands.

The ends  of words and sentences are marked
  by   spaces. It  doesn't matter how many
spaces    you type; one is as good as 100.  The
end of   a line counts as a space.

One   or more   blank lines denote the  end
of  a paragraph.

Since any number of consecutive spaces are treated
like a single one, the formatting of the input
file makes no difference to
      \LaTeX,                % The \LaTeX command generates the LaTeX logo.
but it makes a difference to you.  When you use
\LaTeX, making your input file as easy to read
as possible will be a great help as you write
your document and when you change it.  This sample
file shows how you can add comments to your own input
file.

Because printing is different from typewriting,
there are a number of things that you have to do
differently when preparing an input file than if
you were just typing the document directly.
Quotation marks like
       ``this''
have to be handled specially, as do quotes within
quotes:
       ``\,`this'            % \, separates the double and single quote.
        is what I just
        wrote, not  `that'\,''.

Dashes come in three sizes: an
       intra-word
dash, a medium dash for number ranges like
       1--2,
and a punctuation
       dash---like
this.

A sentence-ending space should be larger than the
space between words within a sentence.  You
sometimes have to type special commands in
conjunction with punctuation characters to get
this right, as in the following sentence.
       Gnats, gnus, etc.\ all  % `\ ' makes an inter-word space.
       begin with G\@.         % \@ marks end-of-sentence punctuation.
You should check the spaces after periods when
reading your output to make sure you haven't
forgotten any special cases.  Generating an
ellipsis
       \ldots\               % `\ ' is needed after `\ldots' because TeX
                             % ignores spaces after command names like \ldots
                             % made from \ + letters.
                             %
                             % Note how a `%' character causes TeX to ignore
                             % the end of the input line, so these blank lines
                             % do not start a new paragraph.
                             %
with the right spacing around the periods requires
a special command.

\LaTeX\ interprets some common characters as
commands, so you must type special commands to
generate them.  These characters include the
following:
       \$ \& \% \# \{ and \}.

In printing, text is usually emphasized with an
       \emph{italic}
type style.

\begin{em}
   A long segment of text can also be emphasized
   in this way.  Text within such a segment can be
   given \emph{additional} emphasis.
\end{em}

It is sometimes necessary to prevent \LaTeX\ from
breaking a line where it might otherwise do so.
This may be at a space, as between the ``Mr.''\ and
``Jones'' in
       ``Mr.~Jones'',        % ~ produces an unbreakable interword space.
or within a word---especially when the word is a
symbol like
       \mbox{\emph{itemnum}}
that makes little sense when hyphenated across
lines.

Footnotes\footnote{This is an example of a footnote.}
pose no problem.

\LaTeX\ is good at typesetting mathematical formulas
like
       \( x-3y + z = 7 \)
or
       \( a_{1} > x^{2n} + y^{2n} > x' \)
or
       \( \ip{A}{B} = \sum_{i} a_{i} b_{i} \).
The spaces you type in a formula are
ignored.  Remember that a letter like
       $x$                   % $ ... $  and  \( ... \)  are equivalent
is a formula when it denotes a mathematical
symbol, and it should be typed as one.

\section{Displayed Text}

Text is displayed by indenting it from the left
margin.  Quotations are commonly displayed.  There
are short quotations
\begin{quote}
   This is a short quotation.  It consists of a
   single paragraph of text.  See how it is formatted.
\end{quote}
and longer ones.
\begin{quotation}
   This is a longer quotation.  It consists of two
   paragraphs of text, neither of which are
   particularly interesting.

   This is the second paragraph of the quotation.  It
   is just as dull as the first paragraph.
\end{quotation}
Another frequently-displayed structure is a list.
The following is an example of an \emph{itemized}
list.
\begin{itemize}
   \item This is the first item of an itemized list.
         Each item in the list is marked with a ``tick''.
         You don't have to worry about what kind of tick
         mark is used.

   \item This is the second item of the list.  It
         contains another list nested inside it.  The inner
         list is an \emph{enumerated} list.
         \begin{enumerate}
            \item This is the first item of an enumerated
                  list that is nested within the itemized list.

            \item This is the second item of the inner list.
                  \LaTeX\ allows you to nest lists deeper than
                  you really should.
         \end{enumerate}
         This is the rest of the second item of the outer
         list.  It is no more interesting than any other
         part of the item.
   \item This is the third item of the list.
\end{itemize}
You can even display poetry.
\begin{verse}
   There is an environment
    for verse \\             % The \\ command separates lines
   Whose features some poets % within a stanza.
   will curse.

                             % One or more blank lines separate stanzas.

   For instead of making\\
   Them do \emph{all} line breaking, \\
   It allows them to put too many words on a line when they'd rather be
   forced to be terse.
\end{verse}

Mathematical formulas may also be displayed.  A
displayed formula
is
one-line long; multiline
formulas require special formatting instructions.
   \[  \ip{\Gamma}{\psi'} = x'' + y^{2} + z_{i}^{n}\]
Don't start a paragraph with a displayed equation,
nor make one a paragraph by itself.

\end{document}               % End of document.

\section{Bench Huge Section 891}
% synthetic bench section 891
% This is a sample LaTeX input file.  (Version of 12 August 2004.)
%
% A '%' character causes TeX to ignore all remaining text on the line,
% and is used for comments like this one.

\documentclass{article}      % Specifies the document class

                             % The preamble begins here.
\title{An Example Document}  % Declares the document's title.
\author{Leslie Lamport}      % Declares the author's name.
\date{January 21, 1994}      % Deleting this command produces today's date.

\newcommand{\ip}[2]{(#1, #2)}
                             % Defines \ip{arg1}{arg2} to mean
                             % (arg1, arg2).

%\newcommand{\ip}[2]{\langle #1 | #2\rangle}
                             % This is an alternative definition of
                             % \ip that is commented out.

\begin{document}             % End of preamble and beginning of text.

\maketitle                   % Produces the title.

This is an example input file.  Comparing it with
the output it generates can show you how to
produce a simple document of your own.

\section{Ordinary Text}      % Produces section heading.  Lower-level
                             % sections are begun with similar
                             % \subsection and \subsubsection commands.

The ends  of words and sentences are marked
  by   spaces. It  doesn't matter how many
spaces    you type; one is as good as 100.  The
end of   a line counts as a space.

One   or more   blank lines denote the  end
of  a paragraph.

Since any number of consecutive spaces are treated
like a single one, the formatting of the input
file makes no difference to
      \LaTeX,                % The \LaTeX command generates the LaTeX logo.
but it makes a difference to you.  When you use
\LaTeX, making your input file as easy to read
as possible will be a great help as you write
your document and when you change it.  This sample
file shows how you can add comments to your own input
file.

Because printing is different from typewriting,
there are a number of things that you have to do
differently when preparing an input file than if
you were just typing the document directly.
Quotation marks like
       ``this''
have to be handled specially, as do quotes within
quotes:
       ``\,`this'            % \, separates the double and single quote.
        is what I just
        wrote, not  `that'\,''.

Dashes come in three sizes: an
       intra-word
dash, a medium dash for number ranges like
       1--2,
and a punctuation
       dash---like
this.

A sentence-ending space should be larger than the
space between words within a sentence.  You
sometimes have to type special commands in
conjunction with punctuation characters to get
this right, as in the following sentence.
       Gnats, gnus, etc.\ all  % `\ ' makes an inter-word space.
       begin with G\@.         % \@ marks end-of-sentence punctuation.
You should check the spaces after periods when
reading your output to make sure you haven't
forgotten any special cases.  Generating an
ellipsis
       \ldots\               % `\ ' is needed after `\ldots' because TeX
                             % ignores spaces after command names like \ldots
                             % made from \ + letters.
                             %
                             % Note how a `%' character causes TeX to ignore
                             % the end of the input line, so these blank lines
                             % do not start a new paragraph.
                             %
with the right spacing around the periods requires
a special command.

\LaTeX\ interprets some common characters as
commands, so you must type special commands to
generate them.  These characters include the
following:
       \$ \& \% \# \{ and \}.

In printing, text is usually emphasized with an
       \emph{italic}
type style.

\begin{em}
   A long segment of text can also be emphasized
   in this way.  Text within such a segment can be
   given \emph{additional} emphasis.
\end{em}

It is sometimes necessary to prevent \LaTeX\ from
breaking a line where it might otherwise do so.
This may be at a space, as between the ``Mr.''\ and
``Jones'' in
       ``Mr.~Jones'',        % ~ produces an unbreakable interword space.
or within a word---especially when the word is a
symbol like
       \mbox{\emph{itemnum}}
that makes little sense when hyphenated across
lines.

Footnotes\footnote{This is an example of a footnote.}
pose no problem.

\LaTeX\ is good at typesetting mathematical formulas
like
       \( x-3y + z = 7 \)
or
       \( a_{1} > x^{2n} + y^{2n} > x' \)
or
       \( \ip{A}{B} = \sum_{i} a_{i} b_{i} \).
The spaces you type in a formula are
ignored.  Remember that a letter like
       $x$                   % $ ... $  and  \( ... \)  are equivalent
is a formula when it denotes a mathematical
symbol, and it should be typed as one.

\section{Displayed Text}

Text is displayed by indenting it from the left
margin.  Quotations are commonly displayed.  There
are short quotations
\begin{quote}
   This is a short quotation.  It consists of a
   single paragraph of text.  See how it is formatted.
\end{quote}
and longer ones.
\begin{quotation}
   This is a longer quotation.  It consists of two
   paragraphs of text, neither of which are
   particularly interesting.

   This is the second paragraph of the quotation.  It
   is just as dull as the first paragraph.
\end{quotation}
Another frequently-displayed structure is a list.
The following is an example of an \emph{itemized}
list.
\begin{itemize}
   \item This is the first item of an itemized list.
         Each item in the list is marked with a ``tick''.
         You don't have to worry about what kind of tick
         mark is used.

   \item This is the second item of the list.  It
         contains another list nested inside it.  The inner
         list is an \emph{enumerated} list.
         \begin{enumerate}
            \item This is the first item of an enumerated
                  list that is nested within the itemized list.

            \item This is the second item of the inner list.
                  \LaTeX\ allows you to nest lists deeper than
                  you really should.
         \end{enumerate}
         This is the rest of the second item of the outer
         list.  It is no more interesting than any other
         part of the item.
   \item This is the third item of the list.
\end{itemize}
You can even display poetry.
\begin{verse}
   There is an environment
    for verse \\             % The \\ command separates lines
   Whose features some poets % within a stanza.
   will curse.

                             % One or more blank lines separate stanzas.

   For instead of making\\
   Them do \emph{all} line breaking, \\
   It allows them to put too many words on a line when they'd rather be
   forced to be terse.
\end{verse}

Mathematical formulas may also be displayed.  A
displayed formula
is
one-line long; multiline
formulas require special formatting instructions.
   \[  \ip{\Gamma}{\psi'} = x'' + y^{2} + z_{i}^{n}\]
Don't start a paragraph with a displayed equation,
nor make one a paragraph by itself.

\end{document}               % End of document.

\section{Bench Huge Section 892}
% synthetic bench section 892
% This is a sample LaTeX input file.  (Version of 12 August 2004.)
%
% A '%' character causes TeX to ignore all remaining text on the line,
% and is used for comments like this one.

\documentclass{article}      % Specifies the document class

                             % The preamble begins here.
\title{An Example Document}  % Declares the document's title.
\author{Leslie Lamport}      % Declares the author's name.
\date{January 21, 1994}      % Deleting this command produces today's date.

\newcommand{\ip}[2]{(#1, #2)}
                             % Defines \ip{arg1}{arg2} to mean
                             % (arg1, arg2).

%\newcommand{\ip}[2]{\langle #1 | #2\rangle}
                             % This is an alternative definition of
                             % \ip that is commented out.

\begin{document}             % End of preamble and beginning of text.

\maketitle                   % Produces the title.

This is an example input file.  Comparing it with
the output it generates can show you how to
produce a simple document of your own.

\section{Ordinary Text}      % Produces section heading.  Lower-level
                             % sections are begun with similar
                             % \subsection and \subsubsection commands.

The ends  of words and sentences are marked
  by   spaces. It  doesn't matter how many
spaces    you type; one is as good as 100.  The
end of   a line counts as a space.

One   or more   blank lines denote the  end
of  a paragraph.

Since any number of consecutive spaces are treated
like a single one, the formatting of the input
file makes no difference to
      \LaTeX,                % The \LaTeX command generates the LaTeX logo.
but it makes a difference to you.  When you use
\LaTeX, making your input file as easy to read
as possible will be a great help as you write
your document and when you change it.  This sample
file shows how you can add comments to your own input
file.

Because printing is different from typewriting,
there are a number of things that you have to do
differently when preparing an input file than if
you were just typing the document directly.
Quotation marks like
       ``this''
have to be handled specially, as do quotes within
quotes:
       ``\,`this'            % \, separates the double and single quote.
        is what I just
        wrote, not  `that'\,''.

Dashes come in three sizes: an
       intra-word
dash, a medium dash for number ranges like
       1--2,
and a punctuation
       dash---like
this.

A sentence-ending space should be larger than the
space between words within a sentence.  You
sometimes have to type special commands in
conjunction with punctuation characters to get
this right, as in the following sentence.
       Gnats, gnus, etc.\ all  % `\ ' makes an inter-word space.
       begin with G\@.         % \@ marks end-of-sentence punctuation.
You should check the spaces after periods when
reading your output to make sure you haven't
forgotten any special cases.  Generating an
ellipsis
       \ldots\               % `\ ' is needed after `\ldots' because TeX
                             % ignores spaces after command names like \ldots
                             % made from \ + letters.
                             %
                             % Note how a `%' character causes TeX to ignore
                             % the end of the input line, so these blank lines
                             % do not start a new paragraph.
                             %
with the right spacing around the periods requires
a special command.

\LaTeX\ interprets some common characters as
commands, so you must type special commands to
generate them.  These characters include the
following:
       \$ \& \% \# \{ and \}.

In printing, text is usually emphasized with an
       \emph{italic}
type style.

\begin{em}
   A long segment of text can also be emphasized
   in this way.  Text within such a segment can be
   given \emph{additional} emphasis.
\end{em}

It is sometimes necessary to prevent \LaTeX\ from
breaking a line where it might otherwise do so.
This may be at a space, as between the ``Mr.''\ and
``Jones'' in
       ``Mr.~Jones'',        % ~ produces an unbreakable interword space.
or within a word---especially when the word is a
symbol like
       \mbox{\emph{itemnum}}
that makes little sense when hyphenated across
lines.

Footnotes\footnote{This is an example of a footnote.}
pose no problem.

\LaTeX\ is good at typesetting mathematical formulas
like
       \( x-3y + z = 7 \)
or
       \( a_{1} > x^{2n} + y^{2n} > x' \)
or
       \( \ip{A}{B} = \sum_{i} a_{i} b_{i} \).
The spaces you type in a formula are
ignored.  Remember that a letter like
       $x$                   % $ ... $  and  \( ... \)  are equivalent
is a formula when it denotes a mathematical
symbol, and it should be typed as one.

\section{Displayed Text}

Text is displayed by indenting it from the left
margin.  Quotations are commonly displayed.  There
are short quotations
\begin{quote}
   This is a short quotation.  It consists of a
   single paragraph of text.  See how it is formatted.
\end{quote}
and longer ones.
\begin{quotation}
   This is a longer quotation.  It consists of two
   paragraphs of text, neither of which are
   particularly interesting.

   This is the second paragraph of the quotation.  It
   is just as dull as the first paragraph.
\end{quotation}
Another frequently-displayed structure is a list.
The following is an example of an \emph{itemized}
list.
\begin{itemize}
   \item This is the first item of an itemized list.
         Each item in the list is marked with a ``tick''.
         You don't have to worry about what kind of tick
         mark is used.

   \item This is the second item of the list.  It
         contains another list nested inside it.  The inner
         list is an \emph{enumerated} list.
         \begin{enumerate}
            \item This is the first item of an enumerated
                  list that is nested within the itemized list.

            \item This is the second item of the inner list.
                  \LaTeX\ allows you to nest lists deeper than
                  you really should.
         \end{enumerate}
         This is the rest of the second item of the outer
         list.  It is no more interesting than any other
         part of the item.
   \item This is the third item of the list.
\end{itemize}
You can even display poetry.
\begin{verse}
   There is an environment
    for verse \\             % The \\ command separates lines
   Whose features some poets % within a stanza.
   will curse.

                             % One or more blank lines separate stanzas.

   For instead of making\\
   Them do \emph{all} line breaking, \\
   It allows them to put too many words on a line when they'd rather be
   forced to be terse.
\end{verse}

Mathematical formulas may also be displayed.  A
displayed formula
is
one-line long; multiline
formulas require special formatting instructions.
   \[  \ip{\Gamma}{\psi'} = x'' + y^{2} + z_{i}^{n}\]
Don't start a paragraph with a displayed equation,
nor make one a paragraph by itself.

\end{document}               % End of document.

\section{Bench Huge Section 893}
% synthetic bench section 893
% This is a sample LaTeX input file.  (Version of 12 August 2004.)
%
% A '%' character causes TeX to ignore all remaining text on the line,
% and is used for comments like this one.

\documentclass{article}      % Specifies the document class

                             % The preamble begins here.
\title{An Example Document}  % Declares the document's title.
\author{Leslie Lamport}      % Declares the author's name.
\date{January 21, 1994}      % Deleting this command produces today's date.

\newcommand{\ip}[2]{(#1, #2)}
                             % Defines \ip{arg1}{arg2} to mean
                             % (arg1, arg2).

%\newcommand{\ip}[2]{\langle #1 | #2\rangle}
                             % This is an alternative definition of
                             % \ip that is commented out.

\begin{document}             % End of preamble and beginning of text.

\maketitle                   % Produces the title.

This is an example input file.  Comparing it with
the output it generates can show you how to
produce a simple document of your own.

\section{Ordinary Text}      % Produces section heading.  Lower-level
                             % sections are begun with similar
                             % \subsection and \subsubsection commands.

The ends  of words and sentences are marked
  by   spaces. It  doesn't matter how many
spaces    you type; one is as good as 100.  The
end of   a line counts as a space.

One   or more   blank lines denote the  end
of  a paragraph.

Since any number of consecutive spaces are treated
like a single one, the formatting of the input
file makes no difference to
      \LaTeX,                % The \LaTeX command generates the LaTeX logo.
but it makes a difference to you.  When you use
\LaTeX, making your input file as easy to read
as possible will be a great help as you write
your document and when you change it.  This sample
file shows how you can add comments to your own input
file.

Because printing is different from typewriting,
there are a number of things that you have to do
differently when preparing an input file than if
you were just typing the document directly.
Quotation marks like
       ``this''
have to be handled specially, as do quotes within
quotes:
       ``\,`this'            % \, separates the double and single quote.
        is what I just
        wrote, not  `that'\,''.

Dashes come in three sizes: an
       intra-word
dash, a medium dash for number ranges like
       1--2,
and a punctuation
       dash---like
this.

A sentence-ending space should be larger than the
space between words within a sentence.  You
sometimes have to type special commands in
conjunction with punctuation characters to get
this right, as in the following sentence.
       Gnats, gnus, etc.\ all  % `\ ' makes an inter-word space.
       begin with G\@.         % \@ marks end-of-sentence punctuation.
You should check the spaces after periods when
reading your output to make sure you haven't
forgotten any special cases.  Generating an
ellipsis
       \ldots\               % `\ ' is needed after `\ldots' because TeX
                             % ignores spaces after command names like \ldots
                             % made from \ + letters.
                             %
                             % Note how a `%' character causes TeX to ignore
                             % the end of the input line, so these blank lines
                             % do not start a new paragraph.
                             %
with the right spacing around the periods requires
a special command.

\LaTeX\ interprets some common characters as
commands, so you must type special commands to
generate them.  These characters include the
following:
       \$ \& \% \# \{ and \}.

In printing, text is usually emphasized with an
       \emph{italic}
type style.

\begin{em}
   A long segment of text can also be emphasized
   in this way.  Text within such a segment can be
   given \emph{additional} emphasis.
\end{em}

It is sometimes necessary to prevent \LaTeX\ from
breaking a line where it might otherwise do so.
This may be at a space, as between the ``Mr.''\ and
``Jones'' in
       ``Mr.~Jones'',        % ~ produces an unbreakable interword space.
or within a word---especially when the word is a
symbol like
       \mbox{\emph{itemnum}}
that makes little sense when hyphenated across
lines.

Footnotes\footnote{This is an example of a footnote.}
pose no problem.

\LaTeX\ is good at typesetting mathematical formulas
like
       \( x-3y + z = 7 \)
or
       \( a_{1} > x^{2n} + y^{2n} > x' \)
or
       \( \ip{A}{B} = \sum_{i} a_{i} b_{i} \).
The spaces you type in a formula are
ignored.  Remember that a letter like
       $x$                   % $ ... $  and  \( ... \)  are equivalent
is a formula when it denotes a mathematical
symbol, and it should be typed as one.

\section{Displayed Text}

Text is displayed by indenting it from the left
margin.  Quotations are commonly displayed.  There
are short quotations
\begin{quote}
   This is a short quotation.  It consists of a
   single paragraph of text.  See how it is formatted.
\end{quote}
and longer ones.
\begin{quotation}
   This is a longer quotation.  It consists of two
   paragraphs of text, neither of which are
   particularly interesting.

   This is the second paragraph of the quotation.  It
   is just as dull as the first paragraph.
\end{quotation}
Another frequently-displayed structure is a list.
The following is an example of an \emph{itemized}
list.
\begin{itemize}
   \item This is the first item of an itemized list.
         Each item in the list is marked with a ``tick''.
         You don't have to worry about what kind of tick
         mark is used.

   \item This is the second item of the list.  It
         contains another list nested inside it.  The inner
         list is an \emph{enumerated} list.
         \begin{enumerate}
            \item This is the first item of an enumerated
                  list that is nested within the itemized list.

            \item This is the second item of the inner list.
                  \LaTeX\ allows you to nest lists deeper than
                  you really should.
         \end{enumerate}
         This is the rest of the second item of the outer
         list.  It is no more interesting than any other
         part of the item.
   \item This is the third item of the list.
\end{itemize}
You can even display poetry.
\begin{verse}
   There is an environment
    for verse \\             % The \\ command separates lines
   Whose features some poets % within a stanza.
   will curse.

                             % One or more blank lines separate stanzas.

   For instead of making\\
   Them do \emph{all} line breaking, \\
   It allows them to put too many words on a line when they'd rather be
   forced to be terse.
\end{verse}

Mathematical formulas may also be displayed.  A
displayed formula
is
one-line long; multiline
formulas require special formatting instructions.
   \[  \ip{\Gamma}{\psi'} = x'' + y^{2} + z_{i}^{n}\]
Don't start a paragraph with a displayed equation,
nor make one a paragraph by itself.

\end{document}               % End of document.

\section{Bench Huge Section 894}
% synthetic bench section 894
% This is a sample LaTeX input file.  (Version of 12 August 2004.)
%
% A '%' character causes TeX to ignore all remaining text on the line,
% and is used for comments like this one.

\documentclass{article}      % Specifies the document class

                             % The preamble begins here.
\title{An Example Document}  % Declares the document's title.
\author{Leslie Lamport}      % Declares the author's name.
\date{January 21, 1994}      % Deleting this command produces today's date.

\newcommand{\ip}[2]{(#1, #2)}
                             % Defines \ip{arg1}{arg2} to mean
                             % (arg1, arg2).

%\newcommand{\ip}[2]{\langle #1 | #2\rangle}
                             % This is an alternative definition of
                             % \ip that is commented out.

\begin{document}             % End of preamble and beginning of text.

\maketitle                   % Produces the title.

This is an example input file.  Comparing it with
the output it generates can show you how to
produce a simple document of your own.

\section{Ordinary Text}      % Produces section heading.  Lower-level
                             % sections are begun with similar
                             % \subsection and \subsubsection commands.

The ends  of words and sentences are marked
  by   spaces. It  doesn't matter how many
spaces    you type; one is as good as 100.  The
end of   a line counts as a space.

One   or more   blank lines denote the  end
of  a paragraph.

Since any number of consecutive spaces are treated
like a single one, the formatting of the input
file makes no difference to
      \LaTeX,                % The \LaTeX command generates the LaTeX logo.
but it makes a difference to you.  When you use
\LaTeX, making your input file as easy to read
as possible will be a great help as you write
your document and when you change it.  This sample
file shows how you can add comments to your own input
file.

Because printing is different from typewriting,
there are a number of things that you have to do
differently when preparing an input file than if
you were just typing the document directly.
Quotation marks like
       ``this''
have to be handled specially, as do quotes within
quotes:
       ``\,`this'            % \, separates the double and single quote.
        is what I just
        wrote, not  `that'\,''.

Dashes come in three sizes: an
       intra-word
dash, a medium dash for number ranges like
       1--2,
and a punctuation
       dash---like
this.

A sentence-ending space should be larger than the
space between words within a sentence.  You
sometimes have to type special commands in
conjunction with punctuation characters to get
this right, as in the following sentence.
       Gnats, gnus, etc.\ all  % `\ ' makes an inter-word space.
       begin with G\@.         % \@ marks end-of-sentence punctuation.
You should check the spaces after periods when
reading your output to make sure you haven't
forgotten any special cases.  Generating an
ellipsis
       \ldots\               % `\ ' is needed after `\ldots' because TeX
                             % ignores spaces after command names like \ldots
                             % made from \ + letters.
                             %
                             % Note how a `%' character causes TeX to ignore
                             % the end of the input line, so these blank lines
                             % do not start a new paragraph.
                             %
with the right spacing around the periods requires
a special command.

\LaTeX\ interprets some common characters as
commands, so you must type special commands to
generate them.  These characters include the
following:
       \$ \& \% \# \{ and \}.

In printing, text is usually emphasized with an
       \emph{italic}
type style.

\begin{em}
   A long segment of text can also be emphasized
   in this way.  Text within such a segment can be
   given \emph{additional} emphasis.
\end{em}

It is sometimes necessary to prevent \LaTeX\ from
breaking a line where it might otherwise do so.
This may be at a space, as between the ``Mr.''\ and
``Jones'' in
       ``Mr.~Jones'',        % ~ produces an unbreakable interword space.
or within a word---especially when the word is a
symbol like
       \mbox{\emph{itemnum}}
that makes little sense when hyphenated across
lines.

Footnotes\footnote{This is an example of a footnote.}
pose no problem.

\LaTeX\ is good at typesetting mathematical formulas
like
       \( x-3y + z = 7 \)
or
       \( a_{1} > x^{2n} + y^{2n} > x' \)
or
       \( \ip{A}{B} = \sum_{i} a_{i} b_{i} \).
The spaces you type in a formula are
ignored.  Remember that a letter like
       $x$                   % $ ... $  and  \( ... \)  are equivalent
is a formula when it denotes a mathematical
symbol, and it should be typed as one.

\section{Displayed Text}

Text is displayed by indenting it from the left
margin.  Quotations are commonly displayed.  There
are short quotations
\begin{quote}
   This is a short quotation.  It consists of a
   single paragraph of text.  See how it is formatted.
\end{quote}
and longer ones.
\begin{quotation}
   This is a longer quotation.  It consists of two
   paragraphs of text, neither of which are
   particularly interesting.

   This is the second paragraph of the quotation.  It
   is just as dull as the first paragraph.
\end{quotation}
Another frequently-displayed structure is a list.
The following is an example of an \emph{itemized}
list.
\begin{itemize}
   \item This is the first item of an itemized list.
         Each item in the list is marked with a ``tick''.
         You don't have to worry about what kind of tick
         mark is used.

   \item This is the second item of the list.  It
         contains another list nested inside it.  The inner
         list is an \emph{enumerated} list.
         \begin{enumerate}
            \item This is the first item of an enumerated
                  list that is nested within the itemized list.

            \item This is the second item of the inner list.
                  \LaTeX\ allows you to nest lists deeper than
                  you really should.
         \end{enumerate}
         This is the rest of the second item of the outer
         list.  It is no more interesting than any other
         part of the item.
   \item This is the third item of the list.
\end{itemize}
You can even display poetry.
\begin{verse}
   There is an environment
    for verse \\             % The \\ command separates lines
   Whose features some poets % within a stanza.
   will curse.

                             % One or more blank lines separate stanzas.

   For instead of making\\
   Them do \emph{all} line breaking, \\
   It allows them to put too many words on a line when they'd rather be
   forced to be terse.
\end{verse}

Mathematical formulas may also be displayed.  A
displayed formula
is
one-line long; multiline
formulas require special formatting instructions.
   \[  \ip{\Gamma}{\psi'} = x'' + y^{2} + z_{i}^{n}\]
Don't start a paragraph with a displayed equation,
nor make one a paragraph by itself.

\end{document}               % End of document.

\section{Bench Huge Section 895}
% synthetic bench section 895
% This is a sample LaTeX input file.  (Version of 12 August 2004.)
%
% A '%' character causes TeX to ignore all remaining text on the line,
% and is used for comments like this one.

\documentclass{article}      % Specifies the document class

                             % The preamble begins here.
\title{An Example Document}  % Declares the document's title.
\author{Leslie Lamport}      % Declares the author's name.
\date{January 21, 1994}      % Deleting this command produces today's date.

\newcommand{\ip}[2]{(#1, #2)}
                             % Defines \ip{arg1}{arg2} to mean
                             % (arg1, arg2).

%\newcommand{\ip}[2]{\langle #1 | #2\rangle}
                             % This is an alternative definition of
                             % \ip that is commented out.

\begin{document}             % End of preamble and beginning of text.

\maketitle                   % Produces the title.

This is an example input file.  Comparing it with
the output it generates can show you how to
produce a simple document of your own.

\section{Ordinary Text}      % Produces section heading.  Lower-level
                             % sections are begun with similar
                             % \subsection and \subsubsection commands.

The ends  of words and sentences are marked
  by   spaces. It  doesn't matter how many
spaces    you type; one is as good as 100.  The
end of   a line counts as a space.

One   or more   blank lines denote the  end
of  a paragraph.

Since any number of consecutive spaces are treated
like a single one, the formatting of the input
file makes no difference to
      \LaTeX,                % The \LaTeX command generates the LaTeX logo.
but it makes a difference to you.  When you use
\LaTeX, making your input file as easy to read
as possible will be a great help as you write
your document and when you change it.  This sample
file shows how you can add comments to your own input
file.

Because printing is different from typewriting,
there are a number of things that you have to do
differently when preparing an input file than if
you were just typing the document directly.
Quotation marks like
       ``this''
have to be handled specially, as do quotes within
quotes:
       ``\,`this'            % \, separates the double and single quote.
        is what I just
        wrote, not  `that'\,''.

Dashes come in three sizes: an
       intra-word
dash, a medium dash for number ranges like
       1--2,
and a punctuation
       dash---like
this.

A sentence-ending space should be larger than the
space between words within a sentence.  You
sometimes have to type special commands in
conjunction with punctuation characters to get
this right, as in the following sentence.
       Gnats, gnus, etc.\ all  % `\ ' makes an inter-word space.
       begin with G\@.         % \@ marks end-of-sentence punctuation.
You should check the spaces after periods when
reading your output to make sure you haven't
forgotten any special cases.  Generating an
ellipsis
       \ldots\               % `\ ' is needed after `\ldots' because TeX
                             % ignores spaces after command names like \ldots
                             % made from \ + letters.
                             %
                             % Note how a `%' character causes TeX to ignore
                             % the end of the input line, so these blank lines
                             % do not start a new paragraph.
                             %
with the right spacing around the periods requires
a special command.

\LaTeX\ interprets some common characters as
commands, so you must type special commands to
generate them.  These characters include the
following:
       \$ \& \% \# \{ and \}.

In printing, text is usually emphasized with an
       \emph{italic}
type style.

\begin{em}
   A long segment of text can also be emphasized
   in this way.  Text within such a segment can be
   given \emph{additional} emphasis.
\end{em}

It is sometimes necessary to prevent \LaTeX\ from
breaking a line where it might otherwise do so.
This may be at a space, as between the ``Mr.''\ and
``Jones'' in
       ``Mr.~Jones'',        % ~ produces an unbreakable interword space.
or within a word---especially when the word is a
symbol like
       \mbox{\emph{itemnum}}
that makes little sense when hyphenated across
lines.

Footnotes\footnote{This is an example of a footnote.}
pose no problem.

\LaTeX\ is good at typesetting mathematical formulas
like
       \( x-3y + z = 7 \)
or
       \( a_{1} > x^{2n} + y^{2n} > x' \)
or
       \( \ip{A}{B} = \sum_{i} a_{i} b_{i} \).
The spaces you type in a formula are
ignored.  Remember that a letter like
       $x$                   % $ ... $  and  \( ... \)  are equivalent
is a formula when it denotes a mathematical
symbol, and it should be typed as one.

\section{Displayed Text}

Text is displayed by indenting it from the left
margin.  Quotations are commonly displayed.  There
are short quotations
\begin{quote}
   This is a short quotation.  It consists of a
   single paragraph of text.  See how it is formatted.
\end{quote}
and longer ones.
\begin{quotation}
   This is a longer quotation.  It consists of two
   paragraphs of text, neither of which are
   particularly interesting.

   This is the second paragraph of the quotation.  It
   is just as dull as the first paragraph.
\end{quotation}
Another frequently-displayed structure is a list.
The following is an example of an \emph{itemized}
list.
\begin{itemize}
   \item This is the first item of an itemized list.
         Each item in the list is marked with a ``tick''.
         You don't have to worry about what kind of tick
         mark is used.

   \item This is the second item of the list.  It
         contains another list nested inside it.  The inner
         list is an \emph{enumerated} list.
         \begin{enumerate}
            \item This is the first item of an enumerated
                  list that is nested within the itemized list.

            \item This is the second item of the inner list.
                  \LaTeX\ allows you to nest lists deeper than
                  you really should.
         \end{enumerate}
         This is the rest of the second item of the outer
         list.  It is no more interesting than any other
         part of the item.
   \item This is the third item of the list.
\end{itemize}
You can even display poetry.
\begin{verse}
   There is an environment
    for verse \\             % The \\ command separates lines
   Whose features some poets % within a stanza.
   will curse.

                             % One or more blank lines separate stanzas.

   For instead of making\\
   Them do \emph{all} line breaking, \\
   It allows them to put too many words on a line when they'd rather be
   forced to be terse.
\end{verse}

Mathematical formulas may also be displayed.  A
displayed formula
is
one-line long; multiline
formulas require special formatting instructions.
   \[  \ip{\Gamma}{\psi'} = x'' + y^{2} + z_{i}^{n}\]
Don't start a paragraph with a displayed equation,
nor make one a paragraph by itself.

\end{document}               % End of document.

\section{Bench Huge Section 896}
% synthetic bench section 896
% This is a sample LaTeX input file.  (Version of 12 August 2004.)
%
% A '%' character causes TeX to ignore all remaining text on the line,
% and is used for comments like this one.

\documentclass{article}      % Specifies the document class

                             % The preamble begins here.
\title{An Example Document}  % Declares the document's title.
\author{Leslie Lamport}      % Declares the author's name.
\date{January 21, 1994}      % Deleting this command produces today's date.

\newcommand{\ip}[2]{(#1, #2)}
                             % Defines \ip{arg1}{arg2} to mean
                             % (arg1, arg2).

%\newcommand{\ip}[2]{\langle #1 | #2\rangle}
                             % This is an alternative definition of
                             % \ip that is commented out.

\begin{document}             % End of preamble and beginning of text.

\maketitle                   % Produces the title.

This is an example input file.  Comparing it with
the output it generates can show you how to
produce a simple document of your own.

\section{Ordinary Text}      % Produces section heading.  Lower-level
                             % sections are begun with similar
                             % \subsection and \subsubsection commands.

The ends  of words and sentences are marked
  by   spaces. It  doesn't matter how many
spaces    you type; one is as good as 100.  The
end of   a line counts as a space.

One   or more   blank lines denote the  end
of  a paragraph.

Since any number of consecutive spaces are treated
like a single one, the formatting of the input
file makes no difference to
      \LaTeX,                % The \LaTeX command generates the LaTeX logo.
but it makes a difference to you.  When you use
\LaTeX, making your input file as easy to read
as possible will be a great help as you write
your document and when you change it.  This sample
file shows how you can add comments to your own input
file.

Because printing is different from typewriting,
there are a number of things that you have to do
differently when preparing an input file than if
you were just typing the document directly.
Quotation marks like
       ``this''
have to be handled specially, as do quotes within
quotes:
       ``\,`this'            % \, separates the double and single quote.
        is what I just
        wrote, not  `that'\,''.

Dashes come in three sizes: an
       intra-word
dash, a medium dash for number ranges like
       1--2,
and a punctuation
       dash---like
this.

A sentence-ending space should be larger than the
space between words within a sentence.  You
sometimes have to type special commands in
conjunction with punctuation characters to get
this right, as in the following sentence.
       Gnats, gnus, etc.\ all  % `\ ' makes an inter-word space.
       begin with G\@.         % \@ marks end-of-sentence punctuation.
You should check the spaces after periods when
reading your output to make sure you haven't
forgotten any special cases.  Generating an
ellipsis
       \ldots\               % `\ ' is needed after `\ldots' because TeX
                             % ignores spaces after command names like \ldots
                             % made from \ + letters.
                             %
                             % Note how a `%' character causes TeX to ignore
                             % the end of the input line, so these blank lines
                             % do not start a new paragraph.
                             %
with the right spacing around the periods requires
a special command.

\LaTeX\ interprets some common characters as
commands, so you must type special commands to
generate them.  These characters include the
following:
       \$ \& \% \# \{ and \}.

In printing, text is usually emphasized with an
       \emph{italic}
type style.

\begin{em}
   A long segment of text can also be emphasized
   in this way.  Text within such a segment can be
   given \emph{additional} emphasis.
\end{em}

It is sometimes necessary to prevent \LaTeX\ from
breaking a line where it might otherwise do so.
This may be at a space, as between the ``Mr.''\ and
``Jones'' in
       ``Mr.~Jones'',        % ~ produces an unbreakable interword space.
or within a word---especially when the word is a
symbol like
       \mbox{\emph{itemnum}}
that makes little sense when hyphenated across
lines.

Footnotes\footnote{This is an example of a footnote.}
pose no problem.

\LaTeX\ is good at typesetting mathematical formulas
like
       \( x-3y + z = 7 \)
or
       \( a_{1} > x^{2n} + y^{2n} > x' \)
or
       \( \ip{A}{B} = \sum_{i} a_{i} b_{i} \).
The spaces you type in a formula are
ignored.  Remember that a letter like
       $x$                   % $ ... $  and  \( ... \)  are equivalent
is a formula when it denotes a mathematical
symbol, and it should be typed as one.

\section{Displayed Text}

Text is displayed by indenting it from the left
margin.  Quotations are commonly displayed.  There
are short quotations
\begin{quote}
   This is a short quotation.  It consists of a
   single paragraph of text.  See how it is formatted.
\end{quote}
and longer ones.
\begin{quotation}
   This is a longer quotation.  It consists of two
   paragraphs of text, neither of which are
   particularly interesting.

   This is the second paragraph of the quotation.  It
   is just as dull as the first paragraph.
\end{quotation}
Another frequently-displayed structure is a list.
The following is an example of an \emph{itemized}
list.
\begin{itemize}
   \item This is the first item of an itemized list.
         Each item in the list is marked with a ``tick''.
         You don't have to worry about what kind of tick
         mark is used.

   \item This is the second item of the list.  It
         contains another list nested inside it.  The inner
         list is an \emph{enumerated} list.
         \begin{enumerate}
            \item This is the first item of an enumerated
                  list that is nested within the itemized list.

            \item This is the second item of the inner list.
                  \LaTeX\ allows you to nest lists deeper than
                  you really should.
         \end{enumerate}
         This is the rest of the second item of the outer
         list.  It is no more interesting than any other
         part of the item.
   \item This is the third item of the list.
\end{itemize}
You can even display poetry.
\begin{verse}
   There is an environment
    for verse \\             % The \\ command separates lines
   Whose features some poets % within a stanza.
   will curse.

                             % One or more blank lines separate stanzas.

   For instead of making\\
   Them do \emph{all} line breaking, \\
   It allows them to put too many words on a line when they'd rather be
   forced to be terse.
\end{verse}

Mathematical formulas may also be displayed.  A
displayed formula
is
one-line long; multiline
formulas require special formatting instructions.
   \[  \ip{\Gamma}{\psi'} = x'' + y^{2} + z_{i}^{n}\]
Don't start a paragraph with a displayed equation,
nor make one a paragraph by itself.

\end{document}               % End of document.

\section{Bench Huge Section 897}
% synthetic bench section 897
% This is a sample LaTeX input file.  (Version of 12 August 2004.)
%
% A '%' character causes TeX to ignore all remaining text on the line,
% and is used for comments like this one.

\documentclass{article}      % Specifies the document class

                             % The preamble begins here.
\title{An Example Document}  % Declares the document's title.
\author{Leslie Lamport}      % Declares the author's name.
\date{January 21, 1994}      % Deleting this command produces today's date.

\newcommand{\ip}[2]{(#1, #2)}
                             % Defines \ip{arg1}{arg2} to mean
                             % (arg1, arg2).

%\newcommand{\ip}[2]{\langle #1 | #2\rangle}
                             % This is an alternative definition of
                             % \ip that is commented out.

\begin{document}             % End of preamble and beginning of text.

\maketitle                   % Produces the title.

This is an example input file.  Comparing it with
the output it generates can show you how to
produce a simple document of your own.

\section{Ordinary Text}      % Produces section heading.  Lower-level
                             % sections are begun with similar
                             % \subsection and \subsubsection commands.

The ends  of words and sentences are marked
  by   spaces. It  doesn't matter how many
spaces    you type; one is as good as 100.  The
end of   a line counts as a space.

One   or more   blank lines denote the  end
of  a paragraph.

Since any number of consecutive spaces are treated
like a single one, the formatting of the input
file makes no difference to
      \LaTeX,                % The \LaTeX command generates the LaTeX logo.
but it makes a difference to you.  When you use
\LaTeX, making your input file as easy to read
as possible will be a great help as you write
your document and when you change it.  This sample
file shows how you can add comments to your own input
file.

Because printing is different from typewriting,
there are a number of things that you have to do
differently when preparing an input file than if
you were just typing the document directly.
Quotation marks like
       ``this''
have to be handled specially, as do quotes within
quotes:
       ``\,`this'            % \, separates the double and single quote.
        is what I just
        wrote, not  `that'\,''.

Dashes come in three sizes: an
       intra-word
dash, a medium dash for number ranges like
       1--2,
and a punctuation
       dash---like
this.

A sentence-ending space should be larger than the
space between words within a sentence.  You
sometimes have to type special commands in
conjunction with punctuation characters to get
this right, as in the following sentence.
       Gnats, gnus, etc.\ all  % `\ ' makes an inter-word space.
       begin with G\@.         % \@ marks end-of-sentence punctuation.
You should check the spaces after periods when
reading your output to make sure you haven't
forgotten any special cases.  Generating an
ellipsis
       \ldots\               % `\ ' is needed after `\ldots' because TeX
                             % ignores spaces after command names like \ldots
                             % made from \ + letters.
                             %
                             % Note how a `%' character causes TeX to ignore
                             % the end of the input line, so these blank lines
                             % do not start a new paragraph.
                             %
with the right spacing around the periods requires
a special command.

\LaTeX\ interprets some common characters as
commands, so you must type special commands to
generate them.  These characters include the
following:
       \$ \& \% \# \{ and \}.

In printing, text is usually emphasized with an
       \emph{italic}
type style.

\begin{em}
   A long segment of text can also be emphasized
   in this way.  Text within such a segment can be
   given \emph{additional} emphasis.
\end{em}

It is sometimes necessary to prevent \LaTeX\ from
breaking a line where it might otherwise do so.
This may be at a space, as between the ``Mr.''\ and
``Jones'' in
       ``Mr.~Jones'',        % ~ produces an unbreakable interword space.
or within a word---especially when the word is a
symbol like
       \mbox{\emph{itemnum}}
that makes little sense when hyphenated across
lines.

Footnotes\footnote{This is an example of a footnote.}
pose no problem.

\LaTeX\ is good at typesetting mathematical formulas
like
       \( x-3y + z = 7 \)
or
       \( a_{1} > x^{2n} + y^{2n} > x' \)
or
       \( \ip{A}{B} = \sum_{i} a_{i} b_{i} \).
The spaces you type in a formula are
ignored.  Remember that a letter like
       $x$                   % $ ... $  and  \( ... \)  are equivalent
is a formula when it denotes a mathematical
symbol, and it should be typed as one.

\section{Displayed Text}

Text is displayed by indenting it from the left
margin.  Quotations are commonly displayed.  There
are short quotations
\begin{quote}
   This is a short quotation.  It consists of a
   single paragraph of text.  See how it is formatted.
\end{quote}
and longer ones.
\begin{quotation}
   This is a longer quotation.  It consists of two
   paragraphs of text, neither of which are
   particularly interesting.

   This is the second paragraph of the quotation.  It
   is just as dull as the first paragraph.
\end{quotation}
Another frequently-displayed structure is a list.
The following is an example of an \emph{itemized}
list.
\begin{itemize}
   \item This is the first item of an itemized list.
         Each item in the list is marked with a ``tick''.
         You don't have to worry about what kind of tick
         mark is used.

   \item This is the second item of the list.  It
         contains another list nested inside it.  The inner
         list is an \emph{enumerated} list.
         \begin{enumerate}
            \item This is the first item of an enumerated
                  list that is nested within the itemized list.

            \item This is the second item of the inner list.
                  \LaTeX\ allows you to nest lists deeper than
                  you really should.
         \end{enumerate}
         This is the rest of the second item of the outer
         list.  It is no more interesting than any other
         part of the item.
   \item This is the third item of the list.
\end{itemize}
You can even display poetry.
\begin{verse}
   There is an environment
    for verse \\             % The \\ command separates lines
   Whose features some poets % within a stanza.
   will curse.

                             % One or more blank lines separate stanzas.

   For instead of making\\
   Them do \emph{all} line breaking, \\
   It allows them to put too many words on a line when they'd rather be
   forced to be terse.
\end{verse}

Mathematical formulas may also be displayed.  A
displayed formula
is
one-line long; multiline
formulas require special formatting instructions.
   \[  \ip{\Gamma}{\psi'} = x'' + y^{2} + z_{i}^{n}\]
Don't start a paragraph with a displayed equation,
nor make one a paragraph by itself.

\end{document}               % End of document.

\section{Bench Huge Section 898}
% synthetic bench section 898
% This is a sample LaTeX input file.  (Version of 12 August 2004.)
%
% A '%' character causes TeX to ignore all remaining text on the line,
% and is used for comments like this one.

\documentclass{article}      % Specifies the document class

                             % The preamble begins here.
\title{An Example Document}  % Declares the document's title.
\author{Leslie Lamport}      % Declares the author's name.
\date{January 21, 1994}      % Deleting this command produces today's date.

\newcommand{\ip}[2]{(#1, #2)}
                             % Defines \ip{arg1}{arg2} to mean
                             % (arg1, arg2).

%\newcommand{\ip}[2]{\langle #1 | #2\rangle}
                             % This is an alternative definition of
                             % \ip that is commented out.

\begin{document}             % End of preamble and beginning of text.

\maketitle                   % Produces the title.

This is an example input file.  Comparing it with
the output it generates can show you how to
produce a simple document of your own.

\section{Ordinary Text}      % Produces section heading.  Lower-level
                             % sections are begun with similar
                             % \subsection and \subsubsection commands.

The ends  of words and sentences are marked
  by   spaces. It  doesn't matter how many
spaces    you type; one is as good as 100.  The
end of   a line counts as a space.

One   or more   blank lines denote the  end
of  a paragraph.

Since any number of consecutive spaces are treated
like a single one, the formatting of the input
file makes no difference to
      \LaTeX,                % The \LaTeX command generates the LaTeX logo.
but it makes a difference to you.  When you use
\LaTeX, making your input file as easy to read
as possible will be a great help as you write
your document and when you change it.  This sample
file shows how you can add comments to your own input
file.

Because printing is different from typewriting,
there are a number of things that you have to do
differently when preparing an input file than if
you were just typing the document directly.
Quotation marks like
       ``this''
have to be handled specially, as do quotes within
quotes:
       ``\,`this'            % \, separates the double and single quote.
        is what I just
        wrote, not  `that'\,''.

Dashes come in three sizes: an
       intra-word
dash, a medium dash for number ranges like
       1--2,
and a punctuation
       dash---like
this.

A sentence-ending space should be larger than the
space between words within a sentence.  You
sometimes have to type special commands in
conjunction with punctuation characters to get
this right, as in the following sentence.
       Gnats, gnus, etc.\ all  % `\ ' makes an inter-word space.
       begin with G\@.         % \@ marks end-of-sentence punctuation.
You should check the spaces after periods when
reading your output to make sure you haven't
forgotten any special cases.  Generating an
ellipsis
       \ldots\               % `\ ' is needed after `\ldots' because TeX
                             % ignores spaces after command names like \ldots
                             % made from \ + letters.
                             %
                             % Note how a `%' character causes TeX to ignore
                             % the end of the input line, so these blank lines
                             % do not start a new paragraph.
                             %
with the right spacing around the periods requires
a special command.

\LaTeX\ interprets some common characters as
commands, so you must type special commands to
generate them.  These characters include the
following:
       \$ \& \% \# \{ and \}.

In printing, text is usually emphasized with an
       \emph{italic}
type style.

\begin{em}
   A long segment of text can also be emphasized
   in this way.  Text within such a segment can be
   given \emph{additional} emphasis.
\end{em}

It is sometimes necessary to prevent \LaTeX\ from
breaking a line where it might otherwise do so.
This may be at a space, as between the ``Mr.''\ and
``Jones'' in
       ``Mr.~Jones'',        % ~ produces an unbreakable interword space.
or within a word---especially when the word is a
symbol like
       \mbox{\emph{itemnum}}
that makes little sense when hyphenated across
lines.

Footnotes\footnote{This is an example of a footnote.}
pose no problem.

\LaTeX\ is good at typesetting mathematical formulas
like
       \( x-3y + z = 7 \)
or
       \( a_{1} > x^{2n} + y^{2n} > x' \)
or
       \( \ip{A}{B} = \sum_{i} a_{i} b_{i} \).
The spaces you type in a formula are
ignored.  Remember that a letter like
       $x$                   % $ ... $  and  \( ... \)  are equivalent
is a formula when it denotes a mathematical
symbol, and it should be typed as one.

\section{Displayed Text}

Text is displayed by indenting it from the left
margin.  Quotations are commonly displayed.  There
are short quotations
\begin{quote}
   This is a short quotation.  It consists of a
   single paragraph of text.  See how it is formatted.
\end{quote}
and longer ones.
\begin{quotation}
   This is a longer quotation.  It consists of two
   paragraphs of text, neither of which are
   particularly interesting.

   This is the second paragraph of the quotation.  It
   is just as dull as the first paragraph.
\end{quotation}
Another frequently-displayed structure is a list.
The following is an example of an \emph{itemized}
list.
\begin{itemize}
   \item This is the first item of an itemized list.
         Each item in the list is marked with a ``tick''.
         You don't have to worry about what kind of tick
         mark is used.

   \item This is the second item of the list.  It
         contains another list nested inside it.  The inner
         list is an \emph{enumerated} list.
         \begin{enumerate}
            \item This is the first item of an enumerated
                  list that is nested within the itemized list.

            \item This is the second item of the inner list.
                  \LaTeX\ allows you to nest lists deeper than
                  you really should.
         \end{enumerate}
         This is the rest of the second item of the outer
         list.  It is no more interesting than any other
         part of the item.
   \item This is the third item of the list.
\end{itemize}
You can even display poetry.
\begin{verse}
   There is an environment
    for verse \\             % The \\ command separates lines
   Whose features some poets % within a stanza.
   will curse.

                             % One or more blank lines separate stanzas.

   For instead of making\\
   Them do \emph{all} line breaking, \\
   It allows them to put too many words on a line when they'd rather be
   forced to be terse.
\end{verse}

Mathematical formulas may also be displayed.  A
displayed formula
is
one-line long; multiline
formulas require special formatting instructions.
   \[  \ip{\Gamma}{\psi'} = x'' + y^{2} + z_{i}^{n}\]
Don't start a paragraph with a displayed equation,
nor make one a paragraph by itself.

\end{document}               % End of document.

\section{Bench Huge Section 899}
% synthetic bench section 899
% This is a sample LaTeX input file.  (Version of 12 August 2004.)
%
% A '%' character causes TeX to ignore all remaining text on the line,
% and is used for comments like this one.

\documentclass{article}      % Specifies the document class

                             % The preamble begins here.
\title{An Example Document}  % Declares the document's title.
\author{Leslie Lamport}      % Declares the author's name.
\date{January 21, 1994}      % Deleting this command produces today's date.

\newcommand{\ip}[2]{(#1, #2)}
                             % Defines \ip{arg1}{arg2} to mean
                             % (arg1, arg2).

%\newcommand{\ip}[2]{\langle #1 | #2\rangle}
                             % This is an alternative definition of
                             % \ip that is commented out.

\begin{document}             % End of preamble and beginning of text.

\maketitle                   % Produces the title.

This is an example input file.  Comparing it with
the output it generates can show you how to
produce a simple document of your own.

\section{Ordinary Text}      % Produces section heading.  Lower-level
                             % sections are begun with similar
                             % \subsection and \subsubsection commands.

The ends  of words and sentences are marked
  by   spaces. It  doesn't matter how many
spaces    you type; one is as good as 100.  The
end of   a line counts as a space.

One   or more   blank lines denote the  end
of  a paragraph.

Since any number of consecutive spaces are treated
like a single one, the formatting of the input
file makes no difference to
      \LaTeX,                % The \LaTeX command generates the LaTeX logo.
but it makes a difference to you.  When you use
\LaTeX, making your input file as easy to read
as possible will be a great help as you write
your document and when you change it.  This sample
file shows how you can add comments to your own input
file.

Because printing is different from typewriting,
there are a number of things that you have to do
differently when preparing an input file than if
you were just typing the document directly.
Quotation marks like
       ``this''
have to be handled specially, as do quotes within
quotes:
       ``\,`this'            % \, separates the double and single quote.
        is what I just
        wrote, not  `that'\,''.

Dashes come in three sizes: an
       intra-word
dash, a medium dash for number ranges like
       1--2,
and a punctuation
       dash---like
this.

A sentence-ending space should be larger than the
space between words within a sentence.  You
sometimes have to type special commands in
conjunction with punctuation characters to get
this right, as in the following sentence.
       Gnats, gnus, etc.\ all  % `\ ' makes an inter-word space.
       begin with G\@.         % \@ marks end-of-sentence punctuation.
You should check the spaces after periods when
reading your output to make sure you haven't
forgotten any special cases.  Generating an
ellipsis
       \ldots\               % `\ ' is needed after `\ldots' because TeX
                             % ignores spaces after command names like \ldots
                             % made from \ + letters.
                             %
                             % Note how a `%' character causes TeX to ignore
                             % the end of the input line, so these blank lines
                             % do not start a new paragraph.
                             %
with the right spacing around the periods requires
a special command.

\LaTeX\ interprets some common characters as
commands, so you must type special commands to
generate them.  These characters include the
following:
       \$ \& \% \# \{ and \}.

In printing, text is usually emphasized with an
       \emph{italic}
type style.

\begin{em}
   A long segment of text can also be emphasized
   in this way.  Text within such a segment can be
   given \emph{additional} emphasis.
\end{em}

It is sometimes necessary to prevent \LaTeX\ from
breaking a line where it might otherwise do so.
This may be at a space, as between the ``Mr.''\ and
``Jones'' in
       ``Mr.~Jones'',        % ~ produces an unbreakable interword space.
or within a word---especially when the word is a
symbol like
       \mbox{\emph{itemnum}}
that makes little sense when hyphenated across
lines.

Footnotes\footnote{This is an example of a footnote.}
pose no problem.

\LaTeX\ is good at typesetting mathematical formulas
like
       \( x-3y + z = 7 \)
or
       \( a_{1} > x^{2n} + y^{2n} > x' \)
or
       \( \ip{A}{B} = \sum_{i} a_{i} b_{i} \).
The spaces you type in a formula are
ignored.  Remember that a letter like
       $x$                   % $ ... $  and  \( ... \)  are equivalent
is a formula when it denotes a mathematical
symbol, and it should be typed as one.

\section{Displayed Text}

Text is displayed by indenting it from the left
margin.  Quotations are commonly displayed.  There
are short quotations
\begin{quote}
   This is a short quotation.  It consists of a
   single paragraph of text.  See how it is formatted.
\end{quote}
and longer ones.
\begin{quotation}
   This is a longer quotation.  It consists of two
   paragraphs of text, neither of which are
   particularly interesting.

   This is the second paragraph of the quotation.  It
   is just as dull as the first paragraph.
\end{quotation}
Another frequently-displayed structure is a list.
The following is an example of an \emph{itemized}
list.
\begin{itemize}
   \item This is the first item of an itemized list.
         Each item in the list is marked with a ``tick''.
         You don't have to worry about what kind of tick
         mark is used.

   \item This is the second item of the list.  It
         contains another list nested inside it.  The inner
         list is an \emph{enumerated} list.
         \begin{enumerate}
            \item This is the first item of an enumerated
                  list that is nested within the itemized list.

            \item This is the second item of the inner list.
                  \LaTeX\ allows you to nest lists deeper than
                  you really should.
         \end{enumerate}
         This is the rest of the second item of the outer
         list.  It is no more interesting than any other
         part of the item.
   \item This is the third item of the list.
\end{itemize}
You can even display poetry.
\begin{verse}
   There is an environment
    for verse \\             % The \\ command separates lines
   Whose features some poets % within a stanza.
   will curse.

                             % One or more blank lines separate stanzas.

   For instead of making\\
   Them do \emph{all} line breaking, \\
   It allows them to put too many words on a line when they'd rather be
   forced to be terse.
\end{verse}

Mathematical formulas may also be displayed.  A
displayed formula
is
one-line long; multiline
formulas require special formatting instructions.
   \[  \ip{\Gamma}{\psi'} = x'' + y^{2} + z_{i}^{n}\]
Don't start a paragraph with a displayed equation,
nor make one a paragraph by itself.

\end{document}               % End of document.

\section{Bench Huge Section 900}
% synthetic bench section 900
% This is a sample LaTeX input file.  (Version of 12 August 2004.)
%
% A '%' character causes TeX to ignore all remaining text on the line,
% and is used for comments like this one.

\documentclass{article}      % Specifies the document class

                             % The preamble begins here.
\title{An Example Document}  % Declares the document's title.
\author{Leslie Lamport}      % Declares the author's name.
\date{January 21, 1994}      % Deleting this command produces today's date.

\newcommand{\ip}[2]{(#1, #2)}
                             % Defines \ip{arg1}{arg2} to mean
                             % (arg1, arg2).

%\newcommand{\ip}[2]{\langle #1 | #2\rangle}
                             % This is an alternative definition of
                             % \ip that is commented out.

\begin{document}             % End of preamble and beginning of text.

\maketitle                   % Produces the title.

This is an example input file.  Comparing it with
the output it generates can show you how to
produce a simple document of your own.

\section{Ordinary Text}      % Produces section heading.  Lower-level
                             % sections are begun with similar
                             % \subsection and \subsubsection commands.

The ends  of words and sentences are marked
  by   spaces. It  doesn't matter how many
spaces    you type; one is as good as 100.  The
end of   a line counts as a space.

One   or more   blank lines denote the  end
of  a paragraph.

Since any number of consecutive spaces are treated
like a single one, the formatting of the input
file makes no difference to
      \LaTeX,                % The \LaTeX command generates the LaTeX logo.
but it makes a difference to you.  When you use
\LaTeX, making your input file as easy to read
as possible will be a great help as you write
your document and when you change it.  This sample
file shows how you can add comments to your own input
file.

Because printing is different from typewriting,
there are a number of things that you have to do
differently when preparing an input file than if
you were just typing the document directly.
Quotation marks like
       ``this''
have to be handled specially, as do quotes within
quotes:
       ``\,`this'            % \, separates the double and single quote.
        is what I just
        wrote, not  `that'\,''.

Dashes come in three sizes: an
       intra-word
dash, a medium dash for number ranges like
       1--2,
and a punctuation
       dash---like
this.

A sentence-ending space should be larger than the
space between words within a sentence.  You
sometimes have to type special commands in
conjunction with punctuation characters to get
this right, as in the following sentence.
       Gnats, gnus, etc.\ all  % `\ ' makes an inter-word space.
       begin with G\@.         % \@ marks end-of-sentence punctuation.
You should check the spaces after periods when
reading your output to make sure you haven't
forgotten any special cases.  Generating an
ellipsis
       \ldots\               % `\ ' is needed after `\ldots' because TeX
                             % ignores spaces after command names like \ldots
                             % made from \ + letters.
                             %
                             % Note how a `%' character causes TeX to ignore
                             % the end of the input line, so these blank lines
                             % do not start a new paragraph.
                             %
with the right spacing around the periods requires
a special command.

\LaTeX\ interprets some common characters as
commands, so you must type special commands to
generate them.  These characters include the
following:
       \$ \& \% \# \{ and \}.

In printing, text is usually emphasized with an
       \emph{italic}
type style.

\begin{em}
   A long segment of text can also be emphasized
   in this way.  Text within such a segment can be
   given \emph{additional} emphasis.
\end{em}

It is sometimes necessary to prevent \LaTeX\ from
breaking a line where it might otherwise do so.
This may be at a space, as between the ``Mr.''\ and
``Jones'' in
       ``Mr.~Jones'',        % ~ produces an unbreakable interword space.
or within a word---especially when the word is a
symbol like
       \mbox{\emph{itemnum}}
that makes little sense when hyphenated across
lines.

Footnotes\footnote{This is an example of a footnote.}
pose no problem.

\LaTeX\ is good at typesetting mathematical formulas
like
       \( x-3y + z = 7 \)
or
       \( a_{1} > x^{2n} + y^{2n} > x' \)
or
       \( \ip{A}{B} = \sum_{i} a_{i} b_{i} \).
The spaces you type in a formula are
ignored.  Remember that a letter like
       $x$                   % $ ... $  and  \( ... \)  are equivalent
is a formula when it denotes a mathematical
symbol, and it should be typed as one.

\section{Displayed Text}

Text is displayed by indenting it from the left
margin.  Quotations are commonly displayed.  There
are short quotations
\begin{quote}
   This is a short quotation.  It consists of a
   single paragraph of text.  See how it is formatted.
\end{quote}
and longer ones.
\begin{quotation}
   This is a longer quotation.  It consists of two
   paragraphs of text, neither of which are
   particularly interesting.

   This is the second paragraph of the quotation.  It
   is just as dull as the first paragraph.
\end{quotation}
Another frequently-displayed structure is a list.
The following is an example of an \emph{itemized}
list.
\begin{itemize}
   \item This is the first item of an itemized list.
         Each item in the list is marked with a ``tick''.
         You don't have to worry about what kind of tick
         mark is used.

   \item This is the second item of the list.  It
         contains another list nested inside it.  The inner
         list is an \emph{enumerated} list.
         \begin{enumerate}
            \item This is the first item of an enumerated
                  list that is nested within the itemized list.

            \item This is the second item of the inner list.
                  \LaTeX\ allows you to nest lists deeper than
                  you really should.
         \end{enumerate}
         This is the rest of the second item of the outer
         list.  It is no more interesting than any other
         part of the item.
   \item This is the third item of the list.
\end{itemize}
You can even display poetry.
\begin{verse}
   There is an environment
    for verse \\             % The \\ command separates lines
   Whose features some poets % within a stanza.
   will curse.

                             % One or more blank lines separate stanzas.

   For instead of making\\
   Them do \emph{all} line breaking, \\
   It allows them to put too many words on a line when they'd rather be
   forced to be terse.
\end{verse}

Mathematical formulas may also be displayed.  A
displayed formula
is
one-line long; multiline
formulas require special formatting instructions.
   \[  \ip{\Gamma}{\psi'} = x'' + y^{2} + z_{i}^{n}\]
Don't start a paragraph with a displayed equation,
nor make one a paragraph by itself.

\end{document}               % End of document.

\section{Bench Huge Section 901}
% synthetic bench section 901
% This is a sample LaTeX input file.  (Version of 12 August 2004.)
%
% A '%' character causes TeX to ignore all remaining text on the line,
% and is used for comments like this one.

\documentclass{article}      % Specifies the document class

                             % The preamble begins here.
\title{An Example Document}  % Declares the document's title.
\author{Leslie Lamport}      % Declares the author's name.
\date{January 21, 1994}      % Deleting this command produces today's date.

\newcommand{\ip}[2]{(#1, #2)}
                             % Defines \ip{arg1}{arg2} to mean
                             % (arg1, arg2).

%\newcommand{\ip}[2]{\langle #1 | #2\rangle}
                             % This is an alternative definition of
                             % \ip that is commented out.

\begin{document}             % End of preamble and beginning of text.

\maketitle                   % Produces the title.

This is an example input file.  Comparing it with
the output it generates can show you how to
produce a simple document of your own.

\section{Ordinary Text}      % Produces section heading.  Lower-level
                             % sections are begun with similar
                             % \subsection and \subsubsection commands.

The ends  of words and sentences are marked
  by   spaces. It  doesn't matter how many
spaces    you type; one is as good as 100.  The
end of   a line counts as a space.

One   or more   blank lines denote the  end
of  a paragraph.

Since any number of consecutive spaces are treated
like a single one, the formatting of the input
file makes no difference to
      \LaTeX,                % The \LaTeX command generates the LaTeX logo.
but it makes a difference to you.  When you use
\LaTeX, making your input file as easy to read
as possible will be a great help as you write
your document and when you change it.  This sample
file shows how you can add comments to your own input
file.

Because printing is different from typewriting,
there are a number of things that you have to do
differently when preparing an input file than if
you were just typing the document directly.
Quotation marks like
       ``this''
have to be handled specially, as do quotes within
quotes:
       ``\,`this'            % \, separates the double and single quote.
        is what I just
        wrote, not  `that'\,''.

Dashes come in three sizes: an
       intra-word
dash, a medium dash for number ranges like
       1--2,
and a punctuation
       dash---like
this.

A sentence-ending space should be larger than the
space between words within a sentence.  You
sometimes have to type special commands in
conjunction with punctuation characters to get
this right, as in the following sentence.
       Gnats, gnus, etc.\ all  % `\ ' makes an inter-word space.
       begin with G\@.         % \@ marks end-of-sentence punctuation.
You should check the spaces after periods when
reading your output to make sure you haven't
forgotten any special cases.  Generating an
ellipsis
       \ldots\               % `\ ' is needed after `\ldots' because TeX
                             % ignores spaces after command names like \ldots
                             % made from \ + letters.
                             %
                             % Note how a `%' character causes TeX to ignore
                             % the end of the input line, so these blank lines
                             % do not start a new paragraph.
                             %
with the right spacing around the periods requires
a special command.

\LaTeX\ interprets some common characters as
commands, so you must type special commands to
generate them.  These characters include the
following:
       \$ \& \% \# \{ and \}.

In printing, text is usually emphasized with an
       \emph{italic}
type style.

\begin{em}
   A long segment of text can also be emphasized
   in this way.  Text within such a segment can be
   given \emph{additional} emphasis.
\end{em}

It is sometimes necessary to prevent \LaTeX\ from
breaking a line where it might otherwise do so.
This may be at a space, as between the ``Mr.''\ and
``Jones'' in
       ``Mr.~Jones'',        % ~ produces an unbreakable interword space.
or within a word---especially when the word is a
symbol like
       \mbox{\emph{itemnum}}
that makes little sense when hyphenated across
lines.

Footnotes\footnote{This is an example of a footnote.}
pose no problem.

\LaTeX\ is good at typesetting mathematical formulas
like
       \( x-3y + z = 7 \)
or
       \( a_{1} > x^{2n} + y^{2n} > x' \)
or
       \( \ip{A}{B} = \sum_{i} a_{i} b_{i} \).
The spaces you type in a formula are
ignored.  Remember that a letter like
       $x$                   % $ ... $  and  \( ... \)  are equivalent
is a formula when it denotes a mathematical
symbol, and it should be typed as one.

\section{Displayed Text}

Text is displayed by indenting it from the left
margin.  Quotations are commonly displayed.  There
are short quotations
\begin{quote}
   This is a short quotation.  It consists of a
   single paragraph of text.  See how it is formatted.
\end{quote}
and longer ones.
\begin{quotation}
   This is a longer quotation.  It consists of two
   paragraphs of text, neither of which are
   particularly interesting.

   This is the second paragraph of the quotation.  It
   is just as dull as the first paragraph.
\end{quotation}
Another frequently-displayed structure is a list.
The following is an example of an \emph{itemized}
list.
\begin{itemize}
   \item This is the first item of an itemized list.
         Each item in the list is marked with a ``tick''.
         You don't have to worry about what kind of tick
         mark is used.

   \item This is the second item of the list.  It
         contains another list nested inside it.  The inner
         list is an \emph{enumerated} list.
         \begin{enumerate}
            \item This is the first item of an enumerated
                  list that is nested within the itemized list.

            \item This is the second item of the inner list.
                  \LaTeX\ allows you to nest lists deeper than
                  you really should.
         \end{enumerate}
         This is the rest of the second item of the outer
         list.  It is no more interesting than any other
         part of the item.
   \item This is the third item of the list.
\end{itemize}
You can even display poetry.
\begin{verse}
   There is an environment
    for verse \\             % The \\ command separates lines
   Whose features some poets % within a stanza.
   will curse.

                             % One or more blank lines separate stanzas.

   For instead of making\\
   Them do \emph{all} line breaking, \\
   It allows them to put too many words on a line when they'd rather be
   forced to be terse.
\end{verse}

Mathematical formulas may also be displayed.  A
displayed formula
is
one-line long; multiline
formulas require special formatting instructions.
   \[  \ip{\Gamma}{\psi'} = x'' + y^{2} + z_{i}^{n}\]
Don't start a paragraph with a displayed equation,
nor make one a paragraph by itself.

\end{document}               % End of document.

\section{Bench Huge Section 902}
% synthetic bench section 902
% This is a sample LaTeX input file.  (Version of 12 August 2004.)
%
% A '%' character causes TeX to ignore all remaining text on the line,
% and is used for comments like this one.

\documentclass{article}      % Specifies the document class

                             % The preamble begins here.
\title{An Example Document}  % Declares the document's title.
\author{Leslie Lamport}      % Declares the author's name.
\date{January 21, 1994}      % Deleting this command produces today's date.

\newcommand{\ip}[2]{(#1, #2)}
                             % Defines \ip{arg1}{arg2} to mean
                             % (arg1, arg2).

%\newcommand{\ip}[2]{\langle #1 | #2\rangle}
                             % This is an alternative definition of
                             % \ip that is commented out.

\begin{document}             % End of preamble and beginning of text.

\maketitle                   % Produces the title.

This is an example input file.  Comparing it with
the output it generates can show you how to
produce a simple document of your own.

\section{Ordinary Text}      % Produces section heading.  Lower-level
                             % sections are begun with similar
                             % \subsection and \subsubsection commands.

The ends  of words and sentences are marked
  by   spaces. It  doesn't matter how many
spaces    you type; one is as good as 100.  The
end of   a line counts as a space.

One   or more   blank lines denote the  end
of  a paragraph.

Since any number of consecutive spaces are treated
like a single one, the formatting of the input
file makes no difference to
      \LaTeX,                % The \LaTeX command generates the LaTeX logo.
but it makes a difference to you.  When you use
\LaTeX, making your input file as easy to read
as possible will be a great help as you write
your document and when you change it.  This sample
file shows how you can add comments to your own input
file.

Because printing is different from typewriting,
there are a number of things that you have to do
differently when preparing an input file than if
you were just typing the document directly.
Quotation marks like
       ``this''
have to be handled specially, as do quotes within
quotes:
       ``\,`this'            % \, separates the double and single quote.
        is what I just
        wrote, not  `that'\,''.

Dashes come in three sizes: an
       intra-word
dash, a medium dash for number ranges like
       1--2,
and a punctuation
       dash---like
this.

A sentence-ending space should be larger than the
space between words within a sentence.  You
sometimes have to type special commands in
conjunction with punctuation characters to get
this right, as in the following sentence.
       Gnats, gnus, etc.\ all  % `\ ' makes an inter-word space.
       begin with G\@.         % \@ marks end-of-sentence punctuation.
You should check the spaces after periods when
reading your output to make sure you haven't
forgotten any special cases.  Generating an
ellipsis
       \ldots\               % `\ ' is needed after `\ldots' because TeX
                             % ignores spaces after command names like \ldots
                             % made from \ + letters.
                             %
                             % Note how a `%' character causes TeX to ignore
                             % the end of the input line, so these blank lines
                             % do not start a new paragraph.
                             %
with the right spacing around the periods requires
a special command.

\LaTeX\ interprets some common characters as
commands, so you must type special commands to
generate them.  These characters include the
following:
       \$ \& \% \# \{ and \}.

In printing, text is usually emphasized with an
       \emph{italic}
type style.

\begin{em}
   A long segment of text can also be emphasized
   in this way.  Text within such a segment can be
   given \emph{additional} emphasis.
\end{em}

It is sometimes necessary to prevent \LaTeX\ from
breaking a line where it might otherwise do so.
This may be at a space, as between the ``Mr.''\ and
``Jones'' in
       ``Mr.~Jones'',        % ~ produces an unbreakable interword space.
or within a word---especially when the word is a
symbol like
       \mbox{\emph{itemnum}}
that makes little sense when hyphenated across
lines.

Footnotes\footnote{This is an example of a footnote.}
pose no problem.

\LaTeX\ is good at typesetting mathematical formulas
like
       \( x-3y + z = 7 \)
or
       \( a_{1} > x^{2n} + y^{2n} > x' \)
or
       \( \ip{A}{B} = \sum_{i} a_{i} b_{i} \).
The spaces you type in a formula are
ignored.  Remember that a letter like
       $x$                   % $ ... $  and  \( ... \)  are equivalent
is a formula when it denotes a mathematical
symbol, and it should be typed as one.

\section{Displayed Text}

Text is displayed by indenting it from the left
margin.  Quotations are commonly displayed.  There
are short quotations
\begin{quote}
   This is a short quotation.  It consists of a
   single paragraph of text.  See how it is formatted.
\end{quote}
and longer ones.
\begin{quotation}
   This is a longer quotation.  It consists of two
   paragraphs of text, neither of which are
   particularly interesting.

   This is the second paragraph of the quotation.  It
   is just as dull as the first paragraph.
\end{quotation}
Another frequently-displayed structure is a list.
The following is an example of an \emph{itemized}
list.
\begin{itemize}
   \item This is the first item of an itemized list.
         Each item in the list is marked with a ``tick''.
         You don't have to worry about what kind of tick
         mark is used.

   \item This is the second item of the list.  It
         contains another list nested inside it.  The inner
         list is an \emph{enumerated} list.
         \begin{enumerate}
            \item This is the first item of an enumerated
                  list that is nested within the itemized list.

            \item This is the second item of the inner list.
                  \LaTeX\ allows you to nest lists deeper than
                  you really should.
         \end{enumerate}
         This is the rest of the second item of the outer
         list.  It is no more interesting than any other
         part of the item.
   \item This is the third item of the list.
\end{itemize}
You can even display poetry.
\begin{verse}
   There is an environment
    for verse \\             % The \\ command separates lines
   Whose features some poets % within a stanza.
   will curse.

                             % One or more blank lines separate stanzas.

   For instead of making\\
   Them do \emph{all} line breaking, \\
   It allows them to put too many words on a line when they'd rather be
   forced to be terse.
\end{verse}

Mathematical formulas may also be displayed.  A
displayed formula
is
one-line long; multiline
formulas require special formatting instructions.
   \[  \ip{\Gamma}{\psi'} = x'' + y^{2} + z_{i}^{n}\]
Don't start a paragraph with a displayed equation,
nor make one a paragraph by itself.

\end{document}               % End of document.

\section{Bench Huge Section 903}
% synthetic bench section 903
% This is a sample LaTeX input file.  (Version of 12 August 2004.)
%
% A '%' character causes TeX to ignore all remaining text on the line,
% and is used for comments like this one.

\documentclass{article}      % Specifies the document class

                             % The preamble begins here.
\title{An Example Document}  % Declares the document's title.
\author{Leslie Lamport}      % Declares the author's name.
\date{January 21, 1994}      % Deleting this command produces today's date.

\newcommand{\ip}[2]{(#1, #2)}
                             % Defines \ip{arg1}{arg2} to mean
                             % (arg1, arg2).

%\newcommand{\ip}[2]{\langle #1 | #2\rangle}
                             % This is an alternative definition of
                             % \ip that is commented out.

\begin{document}             % End of preamble and beginning of text.

\maketitle                   % Produces the title.

This is an example input file.  Comparing it with
the output it generates can show you how to
produce a simple document of your own.

\section{Ordinary Text}      % Produces section heading.  Lower-level
                             % sections are begun with similar
                             % \subsection and \subsubsection commands.

The ends  of words and sentences are marked
  by   spaces. It  doesn't matter how many
spaces    you type; one is as good as 100.  The
end of   a line counts as a space.

One   or more   blank lines denote the  end
of  a paragraph.

Since any number of consecutive spaces are treated
like a single one, the formatting of the input
file makes no difference to
      \LaTeX,                % The \LaTeX command generates the LaTeX logo.
but it makes a difference to you.  When you use
\LaTeX, making your input file as easy to read
as possible will be a great help as you write
your document and when you change it.  This sample
file shows how you can add comments to your own input
file.

Because printing is different from typewriting,
there are a number of things that you have to do
differently when preparing an input file than if
you were just typing the document directly.
Quotation marks like
       ``this''
have to be handled specially, as do quotes within
quotes:
       ``\,`this'            % \, separates the double and single quote.
        is what I just
        wrote, not  `that'\,''.

Dashes come in three sizes: an
       intra-word
dash, a medium dash for number ranges like
       1--2,
and a punctuation
       dash---like
this.

A sentence-ending space should be larger than the
space between words within a sentence.  You
sometimes have to type special commands in
conjunction with punctuation characters to get
this right, as in the following sentence.
       Gnats, gnus, etc.\ all  % `\ ' makes an inter-word space.
       begin with G\@.         % \@ marks end-of-sentence punctuation.
You should check the spaces after periods when
reading your output to make sure you haven't
forgotten any special cases.  Generating an
ellipsis
       \ldots\               % `\ ' is needed after `\ldots' because TeX
                             % ignores spaces after command names like \ldots
                             % made from \ + letters.
                             %
                             % Note how a `%' character causes TeX to ignore
                             % the end of the input line, so these blank lines
                             % do not start a new paragraph.
                             %
with the right spacing around the periods requires
a special command.

\LaTeX\ interprets some common characters as
commands, so you must type special commands to
generate them.  These characters include the
following:
       \$ \& \% \# \{ and \}.

In printing, text is usually emphasized with an
       \emph{italic}
type style.

\begin{em}
   A long segment of text can also be emphasized
   in this way.  Text within such a segment can be
   given \emph{additional} emphasis.
\end{em}

It is sometimes necessary to prevent \LaTeX\ from
breaking a line where it might otherwise do so.
This may be at a space, as between the ``Mr.''\ and
``Jones'' in
       ``Mr.~Jones'',        % ~ produces an unbreakable interword space.
or within a word---especially when the word is a
symbol like
       \mbox{\emph{itemnum}}
that makes little sense when hyphenated across
lines.

Footnotes\footnote{This is an example of a footnote.}
pose no problem.

\LaTeX\ is good at typesetting mathematical formulas
like
       \( x-3y + z = 7 \)
or
       \( a_{1} > x^{2n} + y^{2n} > x' \)
or
       \( \ip{A}{B} = \sum_{i} a_{i} b_{i} \).
The spaces you type in a formula are
ignored.  Remember that a letter like
       $x$                   % $ ... $  and  \( ... \)  are equivalent
is a formula when it denotes a mathematical
symbol, and it should be typed as one.

\section{Displayed Text}

Text is displayed by indenting it from the left
margin.  Quotations are commonly displayed.  There
are short quotations
\begin{quote}
   This is a short quotation.  It consists of a
   single paragraph of text.  See how it is formatted.
\end{quote}
and longer ones.
\begin{quotation}
   This is a longer quotation.  It consists of two
   paragraphs of text, neither of which are
   particularly interesting.

   This is the second paragraph of the quotation.  It
   is just as dull as the first paragraph.
\end{quotation}
Another frequently-displayed structure is a list.
The following is an example of an \emph{itemized}
list.
\begin{itemize}
   \item This is the first item of an itemized list.
         Each item in the list is marked with a ``tick''.
         You don't have to worry about what kind of tick
         mark is used.

   \item This is the second item of the list.  It
         contains another list nested inside it.  The inner
         list is an \emph{enumerated} list.
         \begin{enumerate}
            \item This is the first item of an enumerated
                  list that is nested within the itemized list.

            \item This is the second item of the inner list.
                  \LaTeX\ allows you to nest lists deeper than
                  you really should.
         \end{enumerate}
         This is the rest of the second item of the outer
         list.  It is no more interesting than any other
         part of the item.
   \item This is the third item of the list.
\end{itemize}
You can even display poetry.
\begin{verse}
   There is an environment
    for verse \\             % The \\ command separates lines
   Whose features some poets % within a stanza.
   will curse.

                             % One or more blank lines separate stanzas.

   For instead of making\\
   Them do \emph{all} line breaking, \\
   It allows them to put too many words on a line when they'd rather be
   forced to be terse.
\end{verse}

Mathematical formulas may also be displayed.  A
displayed formula
is
one-line long; multiline
formulas require special formatting instructions.
   \[  \ip{\Gamma}{\psi'} = x'' + y^{2} + z_{i}^{n}\]
Don't start a paragraph with a displayed equation,
nor make one a paragraph by itself.

\end{document}               % End of document.

\section{Bench Huge Section 904}
% synthetic bench section 904
% This is a sample LaTeX input file.  (Version of 12 August 2004.)
%
% A '%' character causes TeX to ignore all remaining text on the line,
% and is used for comments like this one.

\documentclass{article}      % Specifies the document class

                             % The preamble begins here.
\title{An Example Document}  % Declares the document's title.
\author{Leslie Lamport}      % Declares the author's name.
\date{January 21, 1994}      % Deleting this command produces today's date.

\newcommand{\ip}[2]{(#1, #2)}
                             % Defines \ip{arg1}{arg2} to mean
                             % (arg1, arg2).

%\newcommand{\ip}[2]{\langle #1 | #2\rangle}
                             % This is an alternative definition of
                             % \ip that is commented out.

\begin{document}             % End of preamble and beginning of text.

\maketitle                   % Produces the title.

This is an example input file.  Comparing it with
the output it generates can show you how to
produce a simple document of your own.

\section{Ordinary Text}      % Produces section heading.  Lower-level
                             % sections are begun with similar
                             % \subsection and \subsubsection commands.

The ends  of words and sentences are marked
  by   spaces. It  doesn't matter how many
spaces    you type; one is as good as 100.  The
end of   a line counts as a space.

One   or more   blank lines denote the  end
of  a paragraph.

Since any number of consecutive spaces are treated
like a single one, the formatting of the input
file makes no difference to
      \LaTeX,                % The \LaTeX command generates the LaTeX logo.
but it makes a difference to you.  When you use
\LaTeX, making your input file as easy to read
as possible will be a great help as you write
your document and when you change it.  This sample
file shows how you can add comments to your own input
file.

Because printing is different from typewriting,
there are a number of things that you have to do
differently when preparing an input file than if
you were just typing the document directly.
Quotation marks like
       ``this''
have to be handled specially, as do quotes within
quotes:
       ``\,`this'            % \, separates the double and single quote.
        is what I just
        wrote, not  `that'\,''.

Dashes come in three sizes: an
       intra-word
dash, a medium dash for number ranges like
       1--2,
and a punctuation
       dash---like
this.

A sentence-ending space should be larger than the
space between words within a sentence.  You
sometimes have to type special commands in
conjunction with punctuation characters to get
this right, as in the following sentence.
       Gnats, gnus, etc.\ all  % `\ ' makes an inter-word space.
       begin with G\@.         % \@ marks end-of-sentence punctuation.
You should check the spaces after periods when
reading your output to make sure you haven't
forgotten any special cases.  Generating an
ellipsis
       \ldots\               % `\ ' is needed after `\ldots' because TeX
                             % ignores spaces after command names like \ldots
                             % made from \ + letters.
                             %
                             % Note how a `%' character causes TeX to ignore
                             % the end of the input line, so these blank lines
                             % do not start a new paragraph.
                             %
with the right spacing around the periods requires
a special command.

\LaTeX\ interprets some common characters as
commands, so you must type special commands to
generate them.  These characters include the
following:
       \$ \& \% \# \{ and \}.

In printing, text is usually emphasized with an
       \emph{italic}
type style.

\begin{em}
   A long segment of text can also be emphasized
   in this way.  Text within such a segment can be
   given \emph{additional} emphasis.
\end{em}

It is sometimes necessary to prevent \LaTeX\ from
breaking a line where it might otherwise do so.
This may be at a space, as between the ``Mr.''\ and
``Jones'' in
       ``Mr.~Jones'',        % ~ produces an unbreakable interword space.
or within a word---especially when the word is a
symbol like
       \mbox{\emph{itemnum}}
that makes little sense when hyphenated across
lines.

Footnotes\footnote{This is an example of a footnote.}
pose no problem.

\LaTeX\ is good at typesetting mathematical formulas
like
       \( x-3y + z = 7 \)
or
       \( a_{1} > x^{2n} + y^{2n} > x' \)
or
       \( \ip{A}{B} = \sum_{i} a_{i} b_{i} \).
The spaces you type in a formula are
ignored.  Remember that a letter like
       $x$                   % $ ... $  and  \( ... \)  are equivalent
is a formula when it denotes a mathematical
symbol, and it should be typed as one.

\section{Displayed Text}

Text is displayed by indenting it from the left
margin.  Quotations are commonly displayed.  There
are short quotations
\begin{quote}
   This is a short quotation.  It consists of a
   single paragraph of text.  See how it is formatted.
\end{quote}
and longer ones.
\begin{quotation}
   This is a longer quotation.  It consists of two
   paragraphs of text, neither of which are
   particularly interesting.

   This is the second paragraph of the quotation.  It
   is just as dull as the first paragraph.
\end{quotation}
Another frequently-displayed structure is a list.
The following is an example of an \emph{itemized}
list.
\begin{itemize}
   \item This is the first item of an itemized list.
         Each item in the list is marked with a ``tick''.
         You don't have to worry about what kind of tick
         mark is used.

   \item This is the second item of the list.  It
         contains another list nested inside it.  The inner
         list is an \emph{enumerated} list.
         \begin{enumerate}
            \item This is the first item of an enumerated
                  list that is nested within the itemized list.

            \item This is the second item of the inner list.
                  \LaTeX\ allows you to nest lists deeper than
                  you really should.
         \end{enumerate}
         This is the rest of the second item of the outer
         list.  It is no more interesting than any other
         part of the item.
   \item This is the third item of the list.
\end{itemize}
You can even display poetry.
\begin{verse}
   There is an environment
    for verse \\             % The \\ command separates lines
   Whose features some poets % within a stanza.
   will curse.

                             % One or more blank lines separate stanzas.

   For instead of making\\
   Them do \emph{all} line breaking, \\
   It allows them to put too many words on a line when they'd rather be
   forced to be terse.
\end{verse}

Mathematical formulas may also be displayed.  A
displayed formula
is
one-line long; multiline
formulas require special formatting instructions.
   \[  \ip{\Gamma}{\psi'} = x'' + y^{2} + z_{i}^{n}\]
Don't start a paragraph with a displayed equation,
nor make one a paragraph by itself.

\end{document}               % End of document.

\section{Bench Huge Section 905}
% synthetic bench section 905
% This is a sample LaTeX input file.  (Version of 12 August 2004.)
%
% A '%' character causes TeX to ignore all remaining text on the line,
% and is used for comments like this one.

\documentclass{article}      % Specifies the document class

                             % The preamble begins here.
\title{An Example Document}  % Declares the document's title.
\author{Leslie Lamport}      % Declares the author's name.
\date{January 21, 1994}      % Deleting this command produces today's date.

\newcommand{\ip}[2]{(#1, #2)}
                             % Defines \ip{arg1}{arg2} to mean
                             % (arg1, arg2).

%\newcommand{\ip}[2]{\langle #1 | #2\rangle}
                             % This is an alternative definition of
                             % \ip that is commented out.

\begin{document}             % End of preamble and beginning of text.

\maketitle                   % Produces the title.

This is an example input file.  Comparing it with
the output it generates can show you how to
produce a simple document of your own.

\section{Ordinary Text}      % Produces section heading.  Lower-level
                             % sections are begun with similar
                             % \subsection and \subsubsection commands.

The ends  of words and sentences are marked
  by   spaces. It  doesn't matter how many
spaces    you type; one is as good as 100.  The
end of   a line counts as a space.

One   or more   blank lines denote the  end
of  a paragraph.

Since any number of consecutive spaces are treated
like a single one, the formatting of the input
file makes no difference to
      \LaTeX,                % The \LaTeX command generates the LaTeX logo.
but it makes a difference to you.  When you use
\LaTeX, making your input file as easy to read
as possible will be a great help as you write
your document and when you change it.  This sample
file shows how you can add comments to your own input
file.

Because printing is different from typewriting,
there are a number of things that you have to do
differently when preparing an input file than if
you were just typing the document directly.
Quotation marks like
       ``this''
have to be handled specially, as do quotes within
quotes:
       ``\,`this'            % \, separates the double and single quote.
        is what I just
        wrote, not  `that'\,''.

Dashes come in three sizes: an
       intra-word
dash, a medium dash for number ranges like
       1--2,
and a punctuation
       dash---like
this.

A sentence-ending space should be larger than the
space between words within a sentence.  You
sometimes have to type special commands in
conjunction with punctuation characters to get
this right, as in the following sentence.
       Gnats, gnus, etc.\ all  % `\ ' makes an inter-word space.
       begin with G\@.         % \@ marks end-of-sentence punctuation.
You should check the spaces after periods when
reading your output to make sure you haven't
forgotten any special cases.  Generating an
ellipsis
       \ldots\               % `\ ' is needed after `\ldots' because TeX
                             % ignores spaces after command names like \ldots
                             % made from \ + letters.
                             %
                             % Note how a `%' character causes TeX to ignore
                             % the end of the input line, so these blank lines
                             % do not start a new paragraph.
                             %
with the right spacing around the periods requires
a special command.

\LaTeX\ interprets some common characters as
commands, so you must type special commands to
generate them.  These characters include the
following:
       \$ \& \% \# \{ and \}.

In printing, text is usually emphasized with an
       \emph{italic}
type style.

\begin{em}
   A long segment of text can also be emphasized
   in this way.  Text within such a segment can be
   given \emph{additional} emphasis.
\end{em}

It is sometimes necessary to prevent \LaTeX\ from
breaking a line where it might otherwise do so.
This may be at a space, as between the ``Mr.''\ and
``Jones'' in
       ``Mr.~Jones'',        % ~ produces an unbreakable interword space.
or within a word---especially when the word is a
symbol like
       \mbox{\emph{itemnum}}
that makes little sense when hyphenated across
lines.

Footnotes\footnote{This is an example of a footnote.}
pose no problem.

\LaTeX\ is good at typesetting mathematical formulas
like
       \( x-3y + z = 7 \)
or
       \( a_{1} > x^{2n} + y^{2n} > x' \)
or
       \( \ip{A}{B} = \sum_{i} a_{i} b_{i} \).
The spaces you type in a formula are
ignored.  Remember that a letter like
       $x$                   % $ ... $  and  \( ... \)  are equivalent
is a formula when it denotes a mathematical
symbol, and it should be typed as one.

\section{Displayed Text}

Text is displayed by indenting it from the left
margin.  Quotations are commonly displayed.  There
are short quotations
\begin{quote}
   This is a short quotation.  It consists of a
   single paragraph of text.  See how it is formatted.
\end{quote}
and longer ones.
\begin{quotation}
   This is a longer quotation.  It consists of two
   paragraphs of text, neither of which are
   particularly interesting.

   This is the second paragraph of the quotation.  It
   is just as dull as the first paragraph.
\end{quotation}
Another frequently-displayed structure is a list.
The following is an example of an \emph{itemized}
list.
\begin{itemize}
   \item This is the first item of an itemized list.
         Each item in the list is marked with a ``tick''.
         You don't have to worry about what kind of tick
         mark is used.

   \item This is the second item of the list.  It
         contains another list nested inside it.  The inner
         list is an \emph{enumerated} list.
         \begin{enumerate}
            \item This is the first item of an enumerated
                  list that is nested within the itemized list.

            \item This is the second item of the inner list.
                  \LaTeX\ allows you to nest lists deeper than
                  you really should.
         \end{enumerate}
         This is the rest of the second item of the outer
         list.  It is no more interesting than any other
         part of the item.
   \item This is the third item of the list.
\end{itemize}
You can even display poetry.
\begin{verse}
   There is an environment
    for verse \\             % The \\ command separates lines
   Whose features some poets % within a stanza.
   will curse.

                             % One or more blank lines separate stanzas.

   For instead of making\\
   Them do \emph{all} line breaking, \\
   It allows them to put too many words on a line when they'd rather be
   forced to be terse.
\end{verse}

Mathematical formulas may also be displayed.  A
displayed formula
is
one-line long; multiline
formulas require special formatting instructions.
   \[  \ip{\Gamma}{\psi'} = x'' + y^{2} + z_{i}^{n}\]
Don't start a paragraph with a displayed equation,
nor make one a paragraph by itself.

\end{document}               % End of document.

\section{Bench Huge Section 906}
% synthetic bench section 906
% This is a sample LaTeX input file.  (Version of 12 August 2004.)
%
% A '%' character causes TeX to ignore all remaining text on the line,
% and is used for comments like this one.

\documentclass{article}      % Specifies the document class

                             % The preamble begins here.
\title{An Example Document}  % Declares the document's title.
\author{Leslie Lamport}      % Declares the author's name.
\date{January 21, 1994}      % Deleting this command produces today's date.

\newcommand{\ip}[2]{(#1, #2)}
                             % Defines \ip{arg1}{arg2} to mean
                             % (arg1, arg2).

%\newcommand{\ip}[2]{\langle #1 | #2\rangle}
                             % This is an alternative definition of
                             % \ip that is commented out.

\begin{document}             % End of preamble and beginning of text.

\maketitle                   % Produces the title.

This is an example input file.  Comparing it with
the output it generates can show you how to
produce a simple document of your own.

\section{Ordinary Text}      % Produces section heading.  Lower-level
                             % sections are begun with similar
                             % \subsection and \subsubsection commands.

The ends  of words and sentences are marked
  by   spaces. It  doesn't matter how many
spaces    you type; one is as good as 100.  The
end of   a line counts as a space.

One   or more   blank lines denote the  end
of  a paragraph.

Since any number of consecutive spaces are treated
like a single one, the formatting of the input
file makes no difference to
      \LaTeX,                % The \LaTeX command generates the LaTeX logo.
but it makes a difference to you.  When you use
\LaTeX, making your input file as easy to read
as possible will be a great help as you write
your document and when you change it.  This sample
file shows how you can add comments to your own input
file.

Because printing is different from typewriting,
there are a number of things that you have to do
differently when preparing an input file than if
you were just typing the document directly.
Quotation marks like
       ``this''
have to be handled specially, as do quotes within
quotes:
       ``\,`this'            % \, separates the double and single quote.
        is what I just
        wrote, not  `that'\,''.

Dashes come in three sizes: an
       intra-word
dash, a medium dash for number ranges like
       1--2,
and a punctuation
       dash---like
this.

A sentence-ending space should be larger than the
space between words within a sentence.  You
sometimes have to type special commands in
conjunction with punctuation characters to get
this right, as in the following sentence.
       Gnats, gnus, etc.\ all  % `\ ' makes an inter-word space.
       begin with G\@.         % \@ marks end-of-sentence punctuation.
You should check the spaces after periods when
reading your output to make sure you haven't
forgotten any special cases.  Generating an
ellipsis
       \ldots\               % `\ ' is needed after `\ldots' because TeX
                             % ignores spaces after command names like \ldots
                             % made from \ + letters.
                             %
                             % Note how a `%' character causes TeX to ignore
                             % the end of the input line, so these blank lines
                             % do not start a new paragraph.
                             %
with the right spacing around the periods requires
a special command.

\LaTeX\ interprets some common characters as
commands, so you must type special commands to
generate them.  These characters include the
following:
       \$ \& \% \# \{ and \}.

In printing, text is usually emphasized with an
       \emph{italic}
type style.

\begin{em}
   A long segment of text can also be emphasized
   in this way.  Text within such a segment can be
   given \emph{additional} emphasis.
\end{em}

It is sometimes necessary to prevent \LaTeX\ from
breaking a line where it might otherwise do so.
This may be at a space, as between the ``Mr.''\ and
``Jones'' in
       ``Mr.~Jones'',        % ~ produces an unbreakable interword space.
or within a word---especially when the word is a
symbol like
       \mbox{\emph{itemnum}}
that makes little sense when hyphenated across
lines.

Footnotes\footnote{This is an example of a footnote.}
pose no problem.

\LaTeX\ is good at typesetting mathematical formulas
like
       \( x-3y + z = 7 \)
or
       \( a_{1} > x^{2n} + y^{2n} > x' \)
or
       \( \ip{A}{B} = \sum_{i} a_{i} b_{i} \).
The spaces you type in a formula are
ignored.  Remember that a letter like
       $x$                   % $ ... $  and  \( ... \)  are equivalent
is a formula when it denotes a mathematical
symbol, and it should be typed as one.

\section{Displayed Text}

Text is displayed by indenting it from the left
margin.  Quotations are commonly displayed.  There
are short quotations
\begin{quote}
   This is a short quotation.  It consists of a
   single paragraph of text.  See how it is formatted.
\end{quote}
and longer ones.
\begin{quotation}
   This is a longer quotation.  It consists of two
   paragraphs of text, neither of which are
   particularly interesting.

   This is the second paragraph of the quotation.  It
   is just as dull as the first paragraph.
\end{quotation}
Another frequently-displayed structure is a list.
The following is an example of an \emph{itemized}
list.
\begin{itemize}
   \item This is the first item of an itemized list.
         Each item in the list is marked with a ``tick''.
         You don't have to worry about what kind of tick
         mark is used.

   \item This is the second item of the list.  It
         contains another list nested inside it.  The inner
         list is an \emph{enumerated} list.
         \begin{enumerate}
            \item This is the first item of an enumerated
                  list that is nested within the itemized list.

            \item This is the second item of the inner list.
                  \LaTeX\ allows you to nest lists deeper than
                  you really should.
         \end{enumerate}
         This is the rest of the second item of the outer
         list.  It is no more interesting than any other
         part of the item.
   \item This is the third item of the list.
\end{itemize}
You can even display poetry.
\begin{verse}
   There is an environment
    for verse \\             % The \\ command separates lines
   Whose features some poets % within a stanza.
   will curse.

                             % One or more blank lines separate stanzas.

   For instead of making\\
   Them do \emph{all} line breaking, \\
   It allows them to put too many words on a line when they'd rather be
   forced to be terse.
\end{verse}

Mathematical formulas may also be displayed.  A
displayed formula
is
one-line long; multiline
formulas require special formatting instructions.
   \[  \ip{\Gamma}{\psi'} = x'' + y^{2} + z_{i}^{n}\]
Don't start a paragraph with a displayed equation,
nor make one a paragraph by itself.

\end{document}               % End of document.

\section{Bench Huge Section 907}
% synthetic bench section 907
% This is a sample LaTeX input file.  (Version of 12 August 2004.)
%
% A '%' character causes TeX to ignore all remaining text on the line,
% and is used for comments like this one.

\documentclass{article}      % Specifies the document class

                             % The preamble begins here.
\title{An Example Document}  % Declares the document's title.
\author{Leslie Lamport}      % Declares the author's name.
\date{January 21, 1994}      % Deleting this command produces today's date.

\newcommand{\ip}[2]{(#1, #2)}
                             % Defines \ip{arg1}{arg2} to mean
                             % (arg1, arg2).

%\newcommand{\ip}[2]{\langle #1 | #2\rangle}
                             % This is an alternative definition of
                             % \ip that is commented out.

\begin{document}             % End of preamble and beginning of text.

\maketitle                   % Produces the title.

This is an example input file.  Comparing it with
the output it generates can show you how to
produce a simple document of your own.

\section{Ordinary Text}      % Produces section heading.  Lower-level
                             % sections are begun with similar
                             % \subsection and \subsubsection commands.

The ends  of words and sentences are marked
  by   spaces. It  doesn't matter how many
spaces    you type; one is as good as 100.  The
end of   a line counts as a space.

One   or more   blank lines denote the  end
of  a paragraph.

Since any number of consecutive spaces are treated
like a single one, the formatting of the input
file makes no difference to
      \LaTeX,                % The \LaTeX command generates the LaTeX logo.
but it makes a difference to you.  When you use
\LaTeX, making your input file as easy to read
as possible will be a great help as you write
your document and when you change it.  This sample
file shows how you can add comments to your own input
file.

Because printing is different from typewriting,
there are a number of things that you have to do
differently when preparing an input file than if
you were just typing the document directly.
Quotation marks like
       ``this''
have to be handled specially, as do quotes within
quotes:
       ``\,`this'            % \, separates the double and single quote.
        is what I just
        wrote, not  `that'\,''.

Dashes come in three sizes: an
       intra-word
dash, a medium dash for number ranges like
       1--2,
and a punctuation
       dash---like
this.

A sentence-ending space should be larger than the
space between words within a sentence.  You
sometimes have to type special commands in
conjunction with punctuation characters to get
this right, as in the following sentence.
       Gnats, gnus, etc.\ all  % `\ ' makes an inter-word space.
       begin with G\@.         % \@ marks end-of-sentence punctuation.
You should check the spaces after periods when
reading your output to make sure you haven't
forgotten any special cases.  Generating an
ellipsis
       \ldots\               % `\ ' is needed after `\ldots' because TeX
                             % ignores spaces after command names like \ldots
                             % made from \ + letters.
                             %
                             % Note how a `%' character causes TeX to ignore
                             % the end of the input line, so these blank lines
                             % do not start a new paragraph.
                             %
with the right spacing around the periods requires
a special command.

\LaTeX\ interprets some common characters as
commands, so you must type special commands to
generate them.  These characters include the
following:
       \$ \& \% \# \{ and \}.

In printing, text is usually emphasized with an
       \emph{italic}
type style.

\begin{em}
   A long segment of text can also be emphasized
   in this way.  Text within such a segment can be
   given \emph{additional} emphasis.
\end{em}

It is sometimes necessary to prevent \LaTeX\ from
breaking a line where it might otherwise do so.
This may be at a space, as between the ``Mr.''\ and
``Jones'' in
       ``Mr.~Jones'',        % ~ produces an unbreakable interword space.
or within a word---especially when the word is a
symbol like
       \mbox{\emph{itemnum}}
that makes little sense when hyphenated across
lines.

Footnotes\footnote{This is an example of a footnote.}
pose no problem.

\LaTeX\ is good at typesetting mathematical formulas
like
       \( x-3y + z = 7 \)
or
       \( a_{1} > x^{2n} + y^{2n} > x' \)
or
       \( \ip{A}{B} = \sum_{i} a_{i} b_{i} \).
The spaces you type in a formula are
ignored.  Remember that a letter like
       $x$                   % $ ... $  and  \( ... \)  are equivalent
is a formula when it denotes a mathematical
symbol, and it should be typed as one.

\section{Displayed Text}

Text is displayed by indenting it from the left
margin.  Quotations are commonly displayed.  There
are short quotations
\begin{quote}
   This is a short quotation.  It consists of a
   single paragraph of text.  See how it is formatted.
\end{quote}
and longer ones.
\begin{quotation}
   This is a longer quotation.  It consists of two
   paragraphs of text, neither of which are
   particularly interesting.

   This is the second paragraph of the quotation.  It
   is just as dull as the first paragraph.
\end{quotation}
Another frequently-displayed structure is a list.
The following is an example of an \emph{itemized}
list.
\begin{itemize}
   \item This is the first item of an itemized list.
         Each item in the list is marked with a ``tick''.
         You don't have to worry about what kind of tick
         mark is used.

   \item This is the second item of the list.  It
         contains another list nested inside it.  The inner
         list is an \emph{enumerated} list.
         \begin{enumerate}
            \item This is the first item of an enumerated
                  list that is nested within the itemized list.

            \item This is the second item of the inner list.
                  \LaTeX\ allows you to nest lists deeper than
                  you really should.
         \end{enumerate}
         This is the rest of the second item of the outer
         list.  It is no more interesting than any other
         part of the item.
   \item This is the third item of the list.
\end{itemize}
You can even display poetry.
\begin{verse}
   There is an environment
    for verse \\             % The \\ command separates lines
   Whose features some poets % within a stanza.
   will curse.

                             % One or more blank lines separate stanzas.

   For instead of making\\
   Them do \emph{all} line breaking, \\
   It allows them to put too many words on a line when they'd rather be
   forced to be terse.
\end{verse}

Mathematical formulas may also be displayed.  A
displayed formula
is
one-line long; multiline
formulas require special formatting instructions.
   \[  \ip{\Gamma}{\psi'} = x'' + y^{2} + z_{i}^{n}\]
Don't start a paragraph with a displayed equation,
nor make one a paragraph by itself.

\end{document}               % End of document.

\section{Bench Huge Section 908}
% synthetic bench section 908
% This is a sample LaTeX input file.  (Version of 12 August 2004.)
%
% A '%' character causes TeX to ignore all remaining text on the line,
% and is used for comments like this one.

\documentclass{article}      % Specifies the document class

                             % The preamble begins here.
\title{An Example Document}  % Declares the document's title.
\author{Leslie Lamport}      % Declares the author's name.
\date{January 21, 1994}      % Deleting this command produces today's date.

\newcommand{\ip}[2]{(#1, #2)}
                             % Defines \ip{arg1}{arg2} to mean
                             % (arg1, arg2).

%\newcommand{\ip}[2]{\langle #1 | #2\rangle}
                             % This is an alternative definition of
                             % \ip that is commented out.

\begin{document}             % End of preamble and beginning of text.

\maketitle                   % Produces the title.

This is an example input file.  Comparing it with
the output it generates can show you how to
produce a simple document of your own.

\section{Ordinary Text}      % Produces section heading.  Lower-level
                             % sections are begun with similar
                             % \subsection and \subsubsection commands.

The ends  of words and sentences are marked
  by   spaces. It  doesn't matter how many
spaces    you type; one is as good as 100.  The
end of   a line counts as a space.

One   or more   blank lines denote the  end
of  a paragraph.

Since any number of consecutive spaces are treated
like a single one, the formatting of the input
file makes no difference to
      \LaTeX,                % The \LaTeX command generates the LaTeX logo.
but it makes a difference to you.  When you use
\LaTeX, making your input file as easy to read
as possible will be a great help as you write
your document and when you change it.  This sample
file shows how you can add comments to your own input
file.

Because printing is different from typewriting,
there are a number of things that you have to do
differently when preparing an input file than if
you were just typing the document directly.
Quotation marks like
       ``this''
have to be handled specially, as do quotes within
quotes:
       ``\,`this'            % \, separates the double and single quote.
        is what I just
        wrote, not  `that'\,''.

Dashes come in three sizes: an
       intra-word
dash, a medium dash for number ranges like
       1--2,
and a punctuation
       dash---like
this.

A sentence-ending space should be larger than the
space between words within a sentence.  You
sometimes have to type special commands in
conjunction with punctuation characters to get
this right, as in the following sentence.
       Gnats, gnus, etc.\ all  % `\ ' makes an inter-word space.
       begin with G\@.         % \@ marks end-of-sentence punctuation.
You should check the spaces after periods when
reading your output to make sure you haven't
forgotten any special cases.  Generating an
ellipsis
       \ldots\               % `\ ' is needed after `\ldots' because TeX
                             % ignores spaces after command names like \ldots
                             % made from \ + letters.
                             %
                             % Note how a `%' character causes TeX to ignore
                             % the end of the input line, so these blank lines
                             % do not start a new paragraph.
                             %
with the right spacing around the periods requires
a special command.

\LaTeX\ interprets some common characters as
commands, so you must type special commands to
generate them.  These characters include the
following:
       \$ \& \% \# \{ and \}.

In printing, text is usually emphasized with an
       \emph{italic}
type style.

\begin{em}
   A long segment of text can also be emphasized
   in this way.  Text within such a segment can be
   given \emph{additional} emphasis.
\end{em}

It is sometimes necessary to prevent \LaTeX\ from
breaking a line where it might otherwise do so.
This may be at a space, as between the ``Mr.''\ and
``Jones'' in
       ``Mr.~Jones'',        % ~ produces an unbreakable interword space.
or within a word---especially when the word is a
symbol like
       \mbox{\emph{itemnum}}
that makes little sense when hyphenated across
lines.

Footnotes\footnote{This is an example of a footnote.}
pose no problem.

\LaTeX\ is good at typesetting mathematical formulas
like
       \( x-3y + z = 7 \)
or
       \( a_{1} > x^{2n} + y^{2n} > x' \)
or
       \( \ip{A}{B} = \sum_{i} a_{i} b_{i} \).
The spaces you type in a formula are
ignored.  Remember that a letter like
       $x$                   % $ ... $  and  \( ... \)  are equivalent
is a formula when it denotes a mathematical
symbol, and it should be typed as one.

\section{Displayed Text}

Text is displayed by indenting it from the left
margin.  Quotations are commonly displayed.  There
are short quotations
\begin{quote}
   This is a short quotation.  It consists of a
   single paragraph of text.  See how it is formatted.
\end{quote}
and longer ones.
\begin{quotation}
   This is a longer quotation.  It consists of two
   paragraphs of text, neither of which are
   particularly interesting.

   This is the second paragraph of the quotation.  It
   is just as dull as the first paragraph.
\end{quotation}
Another frequently-displayed structure is a list.
The following is an example of an \emph{itemized}
list.
\begin{itemize}
   \item This is the first item of an itemized list.
         Each item in the list is marked with a ``tick''.
         You don't have to worry about what kind of tick
         mark is used.

   \item This is the second item of the list.  It
         contains another list nested inside it.  The inner
         list is an \emph{enumerated} list.
         \begin{enumerate}
            \item This is the first item of an enumerated
                  list that is nested within the itemized list.

            \item This is the second item of the inner list.
                  \LaTeX\ allows you to nest lists deeper than
                  you really should.
         \end{enumerate}
         This is the rest of the second item of the outer
         list.  It is no more interesting than any other
         part of the item.
   \item This is the third item of the list.
\end{itemize}
You can even display poetry.
\begin{verse}
   There is an environment
    for verse \\             % The \\ command separates lines
   Whose features some poets % within a stanza.
   will curse.

                             % One or more blank lines separate stanzas.

   For instead of making\\
   Them do \emph{all} line breaking, \\
   It allows them to put too many words on a line when they'd rather be
   forced to be terse.
\end{verse}

Mathematical formulas may also be displayed.  A
displayed formula
is
one-line long; multiline
formulas require special formatting instructions.
   \[  \ip{\Gamma}{\psi'} = x'' + y^{2} + z_{i}^{n}\]
Don't start a paragraph with a displayed equation,
nor make one a paragraph by itself.

\end{document}               % End of document.

\section{Bench Huge Section 909}
% synthetic bench section 909
% This is a sample LaTeX input file.  (Version of 12 August 2004.)
%
% A '%' character causes TeX to ignore all remaining text on the line,
% and is used for comments like this one.

\documentclass{article}      % Specifies the document class

                             % The preamble begins here.
\title{An Example Document}  % Declares the document's title.
\author{Leslie Lamport}      % Declares the author's name.
\date{January 21, 1994}      % Deleting this command produces today's date.

\newcommand{\ip}[2]{(#1, #2)}
                             % Defines \ip{arg1}{arg2} to mean
                             % (arg1, arg2).

%\newcommand{\ip}[2]{\langle #1 | #2\rangle}
                             % This is an alternative definition of
                             % \ip that is commented out.

\begin{document}             % End of preamble and beginning of text.

\maketitle                   % Produces the title.

This is an example input file.  Comparing it with
the output it generates can show you how to
produce a simple document of your own.

\section{Ordinary Text}      % Produces section heading.  Lower-level
                             % sections are begun with similar
                             % \subsection and \subsubsection commands.

The ends  of words and sentences are marked
  by   spaces. It  doesn't matter how many
spaces    you type; one is as good as 100.  The
end of   a line counts as a space.

One   or more   blank lines denote the  end
of  a paragraph.

Since any number of consecutive spaces are treated
like a single one, the formatting of the input
file makes no difference to
      \LaTeX,                % The \LaTeX command generates the LaTeX logo.
but it makes a difference to you.  When you use
\LaTeX, making your input file as easy to read
as possible will be a great help as you write
your document and when you change it.  This sample
file shows how you can add comments to your own input
file.

Because printing is different from typewriting,
there are a number of things that you have to do
differently when preparing an input file than if
you were just typing the document directly.
Quotation marks like
       ``this''
have to be handled specially, as do quotes within
quotes:
       ``\,`this'            % \, separates the double and single quote.
        is what I just
        wrote, not  `that'\,''.

Dashes come in three sizes: an
       intra-word
dash, a medium dash for number ranges like
       1--2,
and a punctuation
       dash---like
this.

A sentence-ending space should be larger than the
space between words within a sentence.  You
sometimes have to type special commands in
conjunction with punctuation characters to get
this right, as in the following sentence.
       Gnats, gnus, etc.\ all  % `\ ' makes an inter-word space.
       begin with G\@.         % \@ marks end-of-sentence punctuation.
You should check the spaces after periods when
reading your output to make sure you haven't
forgotten any special cases.  Generating an
ellipsis
       \ldots\               % `\ ' is needed after `\ldots' because TeX
                             % ignores spaces after command names like \ldots
                             % made from \ + letters.
                             %
                             % Note how a `%' character causes TeX to ignore
                             % the end of the input line, so these blank lines
                             % do not start a new paragraph.
                             %
with the right spacing around the periods requires
a special command.

\LaTeX\ interprets some common characters as
commands, so you must type special commands to
generate them.  These characters include the
following:
       \$ \& \% \# \{ and \}.

In printing, text is usually emphasized with an
       \emph{italic}
type style.

\begin{em}
   A long segment of text can also be emphasized
   in this way.  Text within such a segment can be
   given \emph{additional} emphasis.
\end{em}

It is sometimes necessary to prevent \LaTeX\ from
breaking a line where it might otherwise do so.
This may be at a space, as between the ``Mr.''\ and
``Jones'' in
       ``Mr.~Jones'',        % ~ produces an unbreakable interword space.
or within a word---especially when the word is a
symbol like
       \mbox{\emph{itemnum}}
that makes little sense when hyphenated across
lines.

Footnotes\footnote{This is an example of a footnote.}
pose no problem.

\LaTeX\ is good at typesetting mathematical formulas
like
       \( x-3y + z = 7 \)
or
       \( a_{1} > x^{2n} + y^{2n} > x' \)
or
       \( \ip{A}{B} = \sum_{i} a_{i} b_{i} \).
The spaces you type in a formula are
ignored.  Remember that a letter like
       $x$                   % $ ... $  and  \( ... \)  are equivalent
is a formula when it denotes a mathematical
symbol, and it should be typed as one.

\section{Displayed Text}

Text is displayed by indenting it from the left
margin.  Quotations are commonly displayed.  There
are short quotations
\begin{quote}
   This is a short quotation.  It consists of a
   single paragraph of text.  See how it is formatted.
\end{quote}
and longer ones.
\begin{quotation}
   This is a longer quotation.  It consists of two
   paragraphs of text, neither of which are
   particularly interesting.

   This is the second paragraph of the quotation.  It
   is just as dull as the first paragraph.
\end{quotation}
Another frequently-displayed structure is a list.
The following is an example of an \emph{itemized}
list.
\begin{itemize}
   \item This is the first item of an itemized list.
         Each item in the list is marked with a ``tick''.
         You don't have to worry about what kind of tick
         mark is used.

   \item This is the second item of the list.  It
         contains another list nested inside it.  The inner
         list is an \emph{enumerated} list.
         \begin{enumerate}
            \item This is the first item of an enumerated
                  list that is nested within the itemized list.

            \item This is the second item of the inner list.
                  \LaTeX\ allows you to nest lists deeper than
                  you really should.
         \end{enumerate}
         This is the rest of the second item of the outer
         list.  It is no more interesting than any other
         part of the item.
   \item This is the third item of the list.
\end{itemize}
You can even display poetry.
\begin{verse}
   There is an environment
    for verse \\             % The \\ command separates lines
   Whose features some poets % within a stanza.
   will curse.

                             % One or more blank lines separate stanzas.

   For instead of making\\
   Them do \emph{all} line breaking, \\
   It allows them to put too many words on a line when they'd rather be
   forced to be terse.
\end{verse}

Mathematical formulas may also be displayed.  A
displayed formula
is
one-line long; multiline
formulas require special formatting instructions.
   \[  \ip{\Gamma}{\psi'} = x'' + y^{2} + z_{i}^{n}\]
Don't start a paragraph with a displayed equation,
nor make one a paragraph by itself.

\end{document}               % End of document.

\section{Bench Huge Section 910}
% synthetic bench section 910
% This is a sample LaTeX input file.  (Version of 12 August 2004.)
%
% A '%' character causes TeX to ignore all remaining text on the line,
% and is used for comments like this one.

\documentclass{article}      % Specifies the document class

                             % The preamble begins here.
\title{An Example Document}  % Declares the document's title.
\author{Leslie Lamport}      % Declares the author's name.
\date{January 21, 1994}      % Deleting this command produces today's date.

\newcommand{\ip}[2]{(#1, #2)}
                             % Defines \ip{arg1}{arg2} to mean
                             % (arg1, arg2).

%\newcommand{\ip}[2]{\langle #1 | #2\rangle}
                             % This is an alternative definition of
                             % \ip that is commented out.

\begin{document}             % End of preamble and beginning of text.

\maketitle                   % Produces the title.

This is an example input file.  Comparing it with
the output it generates can show you how to
produce a simple document of your own.

\section{Ordinary Text}      % Produces section heading.  Lower-level
                             % sections are begun with similar
                             % \subsection and \subsubsection commands.

The ends  of words and sentences are marked
  by   spaces. It  doesn't matter how many
spaces    you type; one is as good as 100.  The
end of   a line counts as a space.

One   or more   blank lines denote the  end
of  a paragraph.

Since any number of consecutive spaces are treated
like a single one, the formatting of the input
file makes no difference to
      \LaTeX,                % The \LaTeX command generates the LaTeX logo.
but it makes a difference to you.  When you use
\LaTeX, making your input file as easy to read
as possible will be a great help as you write
your document and when you change it.  This sample
file shows how you can add comments to your own input
file.

Because printing is different from typewriting,
there are a number of things that you have to do
differently when preparing an input file than if
you were just typing the document directly.
Quotation marks like
       ``this''
have to be handled specially, as do quotes within
quotes:
       ``\,`this'            % \, separates the double and single quote.
        is what I just
        wrote, not  `that'\,''.

Dashes come in three sizes: an
       intra-word
dash, a medium dash for number ranges like
       1--2,
and a punctuation
       dash---like
this.

A sentence-ending space should be larger than the
space between words within a sentence.  You
sometimes have to type special commands in
conjunction with punctuation characters to get
this right, as in the following sentence.
       Gnats, gnus, etc.\ all  % `\ ' makes an inter-word space.
       begin with G\@.         % \@ marks end-of-sentence punctuation.
You should check the spaces after periods when
reading your output to make sure you haven't
forgotten any special cases.  Generating an
ellipsis
       \ldots\               % `\ ' is needed after `\ldots' because TeX
                             % ignores spaces after command names like \ldots
                             % made from \ + letters.
                             %
                             % Note how a `%' character causes TeX to ignore
                             % the end of the input line, so these blank lines
                             % do not start a new paragraph.
                             %
with the right spacing around the periods requires
a special command.

\LaTeX\ interprets some common characters as
commands, so you must type special commands to
generate them.  These characters include the
following:
       \$ \& \% \# \{ and \}.

In printing, text is usually emphasized with an
       \emph{italic}
type style.

\begin{em}
   A long segment of text can also be emphasized
   in this way.  Text within such a segment can be
   given \emph{additional} emphasis.
\end{em}

It is sometimes necessary to prevent \LaTeX\ from
breaking a line where it might otherwise do so.
This may be at a space, as between the ``Mr.''\ and
``Jones'' in
       ``Mr.~Jones'',        % ~ produces an unbreakable interword space.
or within a word---especially when the word is a
symbol like
       \mbox{\emph{itemnum}}
that makes little sense when hyphenated across
lines.

Footnotes\footnote{This is an example of a footnote.}
pose no problem.

\LaTeX\ is good at typesetting mathematical formulas
like
       \( x-3y + z = 7 \)
or
       \( a_{1} > x^{2n} + y^{2n} > x' \)
or
       \( \ip{A}{B} = \sum_{i} a_{i} b_{i} \).
The spaces you type in a formula are
ignored.  Remember that a letter like
       $x$                   % $ ... $  and  \( ... \)  are equivalent
is a formula when it denotes a mathematical
symbol, and it should be typed as one.

\section{Displayed Text}

Text is displayed by indenting it from the left
margin.  Quotations are commonly displayed.  There
are short quotations
\begin{quote}
   This is a short quotation.  It consists of a
   single paragraph of text.  See how it is formatted.
\end{quote}
and longer ones.
\begin{quotation}
   This is a longer quotation.  It consists of two
   paragraphs of text, neither of which are
   particularly interesting.

   This is the second paragraph of the quotation.  It
   is just as dull as the first paragraph.
\end{quotation}
Another frequently-displayed structure is a list.
The following is an example of an \emph{itemized}
list.
\begin{itemize}
   \item This is the first item of an itemized list.
         Each item in the list is marked with a ``tick''.
         You don't have to worry about what kind of tick
         mark is used.

   \item This is the second item of the list.  It
         contains another list nested inside it.  The inner
         list is an \emph{enumerated} list.
         \begin{enumerate}
            \item This is the first item of an enumerated
                  list that is nested within the itemized list.

            \item This is the second item of the inner list.
                  \LaTeX\ allows you to nest lists deeper than
                  you really should.
         \end{enumerate}
         This is the rest of the second item of the outer
         list.  It is no more interesting than any other
         part of the item.
   \item This is the third item of the list.
\end{itemize}
You can even display poetry.
\begin{verse}
   There is an environment
    for verse \\             % The \\ command separates lines
   Whose features some poets % within a stanza.
   will curse.

                             % One or more blank lines separate stanzas.

   For instead of making\\
   Them do \emph{all} line breaking, \\
   It allows them to put too many words on a line when they'd rather be
   forced to be terse.
\end{verse}

Mathematical formulas may also be displayed.  A
displayed formula
is
one-line long; multiline
formulas require special formatting instructions.
   \[  \ip{\Gamma}{\psi'} = x'' + y^{2} + z_{i}^{n}\]
Don't start a paragraph with a displayed equation,
nor make one a paragraph by itself.

\end{document}               % End of document.

\section{Bench Huge Section 911}
% synthetic bench section 911
% This is a sample LaTeX input file.  (Version of 12 August 2004.)
%
% A '%' character causes TeX to ignore all remaining text on the line,
% and is used for comments like this one.

\documentclass{article}      % Specifies the document class

                             % The preamble begins here.
\title{An Example Document}  % Declares the document's title.
\author{Leslie Lamport}      % Declares the author's name.
\date{January 21, 1994}      % Deleting this command produces today's date.

\newcommand{\ip}[2]{(#1, #2)}
                             % Defines \ip{arg1}{arg2} to mean
                             % (arg1, arg2).

%\newcommand{\ip}[2]{\langle #1 | #2\rangle}
                             % This is an alternative definition of
                             % \ip that is commented out.

\begin{document}             % End of preamble and beginning of text.

\maketitle                   % Produces the title.

This is an example input file.  Comparing it with
the output it generates can show you how to
produce a simple document of your own.

\section{Ordinary Text}      % Produces section heading.  Lower-level
                             % sections are begun with similar
                             % \subsection and \subsubsection commands.

The ends  of words and sentences are marked
  by   spaces. It  doesn't matter how many
spaces    you type; one is as good as 100.  The
end of   a line counts as a space.

One   or more   blank lines denote the  end
of  a paragraph.

Since any number of consecutive spaces are treated
like a single one, the formatting of the input
file makes no difference to
      \LaTeX,                % The \LaTeX command generates the LaTeX logo.
but it makes a difference to you.  When you use
\LaTeX, making your input file as easy to read
as possible will be a great help as you write
your document and when you change it.  This sample
file shows how you can add comments to your own input
file.

Because printing is different from typewriting,
there are a number of things that you have to do
differently when preparing an input file than if
you were just typing the document directly.
Quotation marks like
       ``this''
have to be handled specially, as do quotes within
quotes:
       ``\,`this'            % \, separates the double and single quote.
        is what I just
        wrote, not  `that'\,''.

Dashes come in three sizes: an
       intra-word
dash, a medium dash for number ranges like
       1--2,
and a punctuation
       dash---like
this.

A sentence-ending space should be larger than the
space between words within a sentence.  You
sometimes have to type special commands in
conjunction with punctuation characters to get
this right, as in the following sentence.
       Gnats, gnus, etc.\ all  % `\ ' makes an inter-word space.
       begin with G\@.         % \@ marks end-of-sentence punctuation.
You should check the spaces after periods when
reading your output to make sure you haven't
forgotten any special cases.  Generating an
ellipsis
       \ldots\               % `\ ' is needed after `\ldots' because TeX
                             % ignores spaces after command names like \ldots
                             % made from \ + letters.
                             %
                             % Note how a `%' character causes TeX to ignore
                             % the end of the input line, so these blank lines
                             % do not start a new paragraph.
                             %
with the right spacing around the periods requires
a special command.

\LaTeX\ interprets some common characters as
commands, so you must type special commands to
generate them.  These characters include the
following:
       \$ \& \% \# \{ and \}.

In printing, text is usually emphasized with an
       \emph{italic}
type style.

\begin{em}
   A long segment of text can also be emphasized
   in this way.  Text within such a segment can be
   given \emph{additional} emphasis.
\end{em}

It is sometimes necessary to prevent \LaTeX\ from
breaking a line where it might otherwise do so.
This may be at a space, as between the ``Mr.''\ and
``Jones'' in
       ``Mr.~Jones'',        % ~ produces an unbreakable interword space.
or within a word---especially when the word is a
symbol like
       \mbox{\emph{itemnum}}
that makes little sense when hyphenated across
lines.

Footnotes\footnote{This is an example of a footnote.}
pose no problem.

\LaTeX\ is good at typesetting mathematical formulas
like
       \( x-3y + z = 7 \)
or
       \( a_{1} > x^{2n} + y^{2n} > x' \)
or
       \( \ip{A}{B} = \sum_{i} a_{i} b_{i} \).
The spaces you type in a formula are
ignored.  Remember that a letter like
       $x$                   % $ ... $  and  \( ... \)  are equivalent
is a formula when it denotes a mathematical
symbol, and it should be typed as one.

\section{Displayed Text}

Text is displayed by indenting it from the left
margin.  Quotations are commonly displayed.  There
are short quotations
\begin{quote}
   This is a short quotation.  It consists of a
   single paragraph of text.  See how it is formatted.
\end{quote}
and longer ones.
\begin{quotation}
   This is a longer quotation.  It consists of two
   paragraphs of text, neither of which are
   particularly interesting.

   This is the second paragraph of the quotation.  It
   is just as dull as the first paragraph.
\end{quotation}
Another frequently-displayed structure is a list.
The following is an example of an \emph{itemized}
list.
\begin{itemize}
   \item This is the first item of an itemized list.
         Each item in the list is marked with a ``tick''.
         You don't have to worry about what kind of tick
         mark is used.

   \item This is the second item of the list.  It
         contains another list nested inside it.  The inner
         list is an \emph{enumerated} list.
         \begin{enumerate}
            \item This is the first item of an enumerated
                  list that is nested within the itemized list.

            \item This is the second item of the inner list.
                  \LaTeX\ allows you to nest lists deeper than
                  you really should.
         \end{enumerate}
         This is the rest of the second item of the outer
         list.  It is no more interesting than any other
         part of the item.
   \item This is the third item of the list.
\end{itemize}
You can even display poetry.
\begin{verse}
   There is an environment
    for verse \\             % The \\ command separates lines
   Whose features some poets % within a stanza.
   will curse.

                             % One or more blank lines separate stanzas.

   For instead of making\\
   Them do \emph{all} line breaking, \\
   It allows them to put too many words on a line when they'd rather be
   forced to be terse.
\end{verse}

Mathematical formulas may also be displayed.  A
displayed formula
is
one-line long; multiline
formulas require special formatting instructions.
   \[  \ip{\Gamma}{\psi'} = x'' + y^{2} + z_{i}^{n}\]
Don't start a paragraph with a displayed equation,
nor make one a paragraph by itself.

\end{document}               % End of document.

\section{Bench Huge Section 912}
% synthetic bench section 912
% This is a sample LaTeX input file.  (Version of 12 August 2004.)
%
% A '%' character causes TeX to ignore all remaining text on the line,
% and is used for comments like this one.

\documentclass{article}      % Specifies the document class

                             % The preamble begins here.
\title{An Example Document}  % Declares the document's title.
\author{Leslie Lamport}      % Declares the author's name.
\date{January 21, 1994}      % Deleting this command produces today's date.

\newcommand{\ip}[2]{(#1, #2)}
                             % Defines \ip{arg1}{arg2} to mean
                             % (arg1, arg2).

%\newcommand{\ip}[2]{\langle #1 | #2\rangle}
                             % This is an alternative definition of
                             % \ip that is commented out.

\begin{document}             % End of preamble and beginning of text.

\maketitle                   % Produces the title.

This is an example input file.  Comparing it with
the output it generates can show you how to
produce a simple document of your own.

\section{Ordinary Text}      % Produces section heading.  Lower-level
                             % sections are begun with similar
                             % \subsection and \subsubsection commands.

The ends  of words and sentences are marked
  by   spaces. It  doesn't matter how many
spaces    you type; one is as good as 100.  The
end of   a line counts as a space.

One   or more   blank lines denote the  end
of  a paragraph.

Since any number of consecutive spaces are treated
like a single one, the formatting of the input
file makes no difference to
      \LaTeX,                % The \LaTeX command generates the LaTeX logo.
but it makes a difference to you.  When you use
\LaTeX, making your input file as easy to read
as possible will be a great help as you write
your document and when you change it.  This sample
file shows how you can add comments to your own input
file.

Because printing is different from typewriting,
there are a number of things that you have to do
differently when preparing an input file than if
you were just typing the document directly.
Quotation marks like
       ``this''
have to be handled specially, as do quotes within
quotes:
       ``\,`this'            % \, separates the double and single quote.
        is what I just
        wrote, not  `that'\,''.

Dashes come in three sizes: an
       intra-word
dash, a medium dash for number ranges like
       1--2,
and a punctuation
       dash---like
this.

A sentence-ending space should be larger than the
space between words within a sentence.  You
sometimes have to type special commands in
conjunction with punctuation characters to get
this right, as in the following sentence.
       Gnats, gnus, etc.\ all  % `\ ' makes an inter-word space.
       begin with G\@.         % \@ marks end-of-sentence punctuation.
You should check the spaces after periods when
reading your output to make sure you haven't
forgotten any special cases.  Generating an
ellipsis
       \ldots\               % `\ ' is needed after `\ldots' because TeX
                             % ignores spaces after command names like \ldots
                             % made from \ + letters.
                             %
                             % Note how a `%' character causes TeX to ignore
                             % the end of the input line, so these blank lines
                             % do not start a new paragraph.
                             %
with the right spacing around the periods requires
a special command.

\LaTeX\ interprets some common characters as
commands, so you must type special commands to
generate them.  These characters include the
following:
       \$ \& \% \# \{ and \}.

In printing, text is usually emphasized with an
       \emph{italic}
type style.

\begin{em}
   A long segment of text can also be emphasized
   in this way.  Text within such a segment can be
   given \emph{additional} emphasis.
\end{em}

It is sometimes necessary to prevent \LaTeX\ from
breaking a line where it might otherwise do so.
This may be at a space, as between the ``Mr.''\ and
``Jones'' in
       ``Mr.~Jones'',        % ~ produces an unbreakable interword space.
or within a word---especially when the word is a
symbol like
       \mbox{\emph{itemnum}}
that makes little sense when hyphenated across
lines.

Footnotes\footnote{This is an example of a footnote.}
pose no problem.

\LaTeX\ is good at typesetting mathematical formulas
like
       \( x-3y + z = 7 \)
or
       \( a_{1} > x^{2n} + y^{2n} > x' \)
or
       \( \ip{A}{B} = \sum_{i} a_{i} b_{i} \).
The spaces you type in a formula are
ignored.  Remember that a letter like
       $x$                   % $ ... $  and  \( ... \)  are equivalent
is a formula when it denotes a mathematical
symbol, and it should be typed as one.

\section{Displayed Text}

Text is displayed by indenting it from the left
margin.  Quotations are commonly displayed.  There
are short quotations
\begin{quote}
   This is a short quotation.  It consists of a
   single paragraph of text.  See how it is formatted.
\end{quote}
and longer ones.
\begin{quotation}
   This is a longer quotation.  It consists of two
   paragraphs of text, neither of which are
   particularly interesting.

   This is the second paragraph of the quotation.  It
   is just as dull as the first paragraph.
\end{quotation}
Another frequently-displayed structure is a list.
The following is an example of an \emph{itemized}
list.
\begin{itemize}
   \item This is the first item of an itemized list.
         Each item in the list is marked with a ``tick''.
         You don't have to worry about what kind of tick
         mark is used.

   \item This is the second item of the list.  It
         contains another list nested inside it.  The inner
         list is an \emph{enumerated} list.
         \begin{enumerate}
            \item This is the first item of an enumerated
                  list that is nested within the itemized list.

            \item This is the second item of the inner list.
                  \LaTeX\ allows you to nest lists deeper than
                  you really should.
         \end{enumerate}
         This is the rest of the second item of the outer
         list.  It is no more interesting than any other
         part of the item.
   \item This is the third item of the list.
\end{itemize}
You can even display poetry.
\begin{verse}
   There is an environment
    for verse \\             % The \\ command separates lines
   Whose features some poets % within a stanza.
   will curse.

                             % One or more blank lines separate stanzas.

   For instead of making\\
   Them do \emph{all} line breaking, \\
   It allows them to put too many words on a line when they'd rather be
   forced to be terse.
\end{verse}

Mathematical formulas may also be displayed.  A
displayed formula
is
one-line long; multiline
formulas require special formatting instructions.
   \[  \ip{\Gamma}{\psi'} = x'' + y^{2} + z_{i}^{n}\]
Don't start a paragraph with a displayed equation,
nor make one a paragraph by itself.

\end{document}               % End of document.

\section{Bench Huge Section 913}
% synthetic bench section 913
% This is a sample LaTeX input file.  (Version of 12 August 2004.)
%
% A '%' character causes TeX to ignore all remaining text on the line,
% and is used for comments like this one.

\documentclass{article}      % Specifies the document class

                             % The preamble begins here.
\title{An Example Document}  % Declares the document's title.
\author{Leslie Lamport}      % Declares the author's name.
\date{January 21, 1994}      % Deleting this command produces today's date.

\newcommand{\ip}[2]{(#1, #2)}
                             % Defines \ip{arg1}{arg2} to mean
                             % (arg1, arg2).

%\newcommand{\ip}[2]{\langle #1 | #2\rangle}
                             % This is an alternative definition of
                             % \ip that is commented out.

\begin{document}             % End of preamble and beginning of text.

\maketitle                   % Produces the title.

This is an example input file.  Comparing it with
the output it generates can show you how to
produce a simple document of your own.

\section{Ordinary Text}      % Produces section heading.  Lower-level
                             % sections are begun with similar
                             % \subsection and \subsubsection commands.

The ends  of words and sentences are marked
  by   spaces. It  doesn't matter how many
spaces    you type; one is as good as 100.  The
end of   a line counts as a space.

One   or more   blank lines denote the  end
of  a paragraph.

Since any number of consecutive spaces are treated
like a single one, the formatting of the input
file makes no difference to
      \LaTeX,                % The \LaTeX command generates the LaTeX logo.
but it makes a difference to you.  When you use
\LaTeX, making your input file as easy to read
as possible will be a great help as you write
your document and when you change it.  This sample
file shows how you can add comments to your own input
file.

Because printing is different from typewriting,
there are a number of things that you have to do
differently when preparing an input file than if
you were just typing the document directly.
Quotation marks like
       ``this''
have to be handled specially, as do quotes within
quotes:
       ``\,`this'            % \, separates the double and single quote.
        is what I just
        wrote, not  `that'\,''.

Dashes come in three sizes: an
       intra-word
dash, a medium dash for number ranges like
       1--2,
and a punctuation
       dash---like
this.

A sentence-ending space should be larger than the
space between words within a sentence.  You
sometimes have to type special commands in
conjunction with punctuation characters to get
this right, as in the following sentence.
       Gnats, gnus, etc.\ all  % `\ ' makes an inter-word space.
       begin with G\@.         % \@ marks end-of-sentence punctuation.
You should check the spaces after periods when
reading your output to make sure you haven't
forgotten any special cases.  Generating an
ellipsis
       \ldots\               % `\ ' is needed after `\ldots' because TeX
                             % ignores spaces after command names like \ldots
                             % made from \ + letters.
                             %
                             % Note how a `%' character causes TeX to ignore
                             % the end of the input line, so these blank lines
                             % do not start a new paragraph.
                             %
with the right spacing around the periods requires
a special command.

\LaTeX\ interprets some common characters as
commands, so you must type special commands to
generate them.  These characters include the
following:
       \$ \& \% \# \{ and \}.

In printing, text is usually emphasized with an
       \emph{italic}
type style.

\begin{em}
   A long segment of text can also be emphasized
   in this way.  Text within such a segment can be
   given \emph{additional} emphasis.
\end{em}

It is sometimes necessary to prevent \LaTeX\ from
breaking a line where it might otherwise do so.
This may be at a space, as between the ``Mr.''\ and
``Jones'' in
       ``Mr.~Jones'',        % ~ produces an unbreakable interword space.
or within a word---especially when the word is a
symbol like
       \mbox{\emph{itemnum}}
that makes little sense when hyphenated across
lines.

Footnotes\footnote{This is an example of a footnote.}
pose no problem.

\LaTeX\ is good at typesetting mathematical formulas
like
       \( x-3y + z = 7 \)
or
       \( a_{1} > x^{2n} + y^{2n} > x' \)
or
       \( \ip{A}{B} = \sum_{i} a_{i} b_{i} \).
The spaces you type in a formula are
ignored.  Remember that a letter like
       $x$                   % $ ... $  and  \( ... \)  are equivalent
is a formula when it denotes a mathematical
symbol, and it should be typed as one.

\section{Displayed Text}

Text is displayed by indenting it from the left
margin.  Quotations are commonly displayed.  There
are short quotations
\begin{quote}
   This is a short quotation.  It consists of a
   single paragraph of text.  See how it is formatted.
\end{quote}
and longer ones.
\begin{quotation}
   This is a longer quotation.  It consists of two
   paragraphs of text, neither of which are
   particularly interesting.

   This is the second paragraph of the quotation.  It
   is just as dull as the first paragraph.
\end{quotation}
Another frequently-displayed structure is a list.
The following is an example of an \emph{itemized}
list.
\begin{itemize}
   \item This is the first item of an itemized list.
         Each item in the list is marked with a ``tick''.
         You don't have to worry about what kind of tick
         mark is used.

   \item This is the second item of the list.  It
         contains another list nested inside it.  The inner
         list is an \emph{enumerated} list.
         \begin{enumerate}
            \item This is the first item of an enumerated
                  list that is nested within the itemized list.

            \item This is the second item of the inner list.
                  \LaTeX\ allows you to nest lists deeper than
                  you really should.
         \end{enumerate}
         This is the rest of the second item of the outer
         list.  It is no more interesting than any other
         part of the item.
   \item This is the third item of the list.
\end{itemize}
You can even display poetry.
\begin{verse}
   There is an environment
    for verse \\             % The \\ command separates lines
   Whose features some poets % within a stanza.
   will curse.

                             % One or more blank lines separate stanzas.

   For instead of making\\
   Them do \emph{all} line breaking, \\
   It allows them to put too many words on a line when they'd rather be
   forced to be terse.
\end{verse}

Mathematical formulas may also be displayed.  A
displayed formula
is
one-line long; multiline
formulas require special formatting instructions.
   \[  \ip{\Gamma}{\psi'} = x'' + y^{2} + z_{i}^{n}\]
Don't start a paragraph with a displayed equation,
nor make one a paragraph by itself.

\end{document}               % End of document.

\section{Bench Huge Section 914}
% synthetic bench section 914
% This is a sample LaTeX input file.  (Version of 12 August 2004.)
%
% A '%' character causes TeX to ignore all remaining text on the line,
% and is used for comments like this one.

\documentclass{article}      % Specifies the document class

                             % The preamble begins here.
\title{An Example Document}  % Declares the document's title.
\author{Leslie Lamport}      % Declares the author's name.
\date{January 21, 1994}      % Deleting this command produces today's date.

\newcommand{\ip}[2]{(#1, #2)}
                             % Defines \ip{arg1}{arg2} to mean
                             % (arg1, arg2).

%\newcommand{\ip}[2]{\langle #1 | #2\rangle}
                             % This is an alternative definition of
                             % \ip that is commented out.

\begin{document}             % End of preamble and beginning of text.

\maketitle                   % Produces the title.

This is an example input file.  Comparing it with
the output it generates can show you how to
produce a simple document of your own.

\section{Ordinary Text}      % Produces section heading.  Lower-level
                             % sections are begun with similar
                             % \subsection and \subsubsection commands.

The ends  of words and sentences are marked
  by   spaces. It  doesn't matter how many
spaces    you type; one is as good as 100.  The
end of   a line counts as a space.

One   or more   blank lines denote the  end
of  a paragraph.

Since any number of consecutive spaces are treated
like a single one, the formatting of the input
file makes no difference to
      \LaTeX,                % The \LaTeX command generates the LaTeX logo.
but it makes a difference to you.  When you use
\LaTeX, making your input file as easy to read
as possible will be a great help as you write
your document and when you change it.  This sample
file shows how you can add comments to your own input
file.

Because printing is different from typewriting,
there are a number of things that you have to do
differently when preparing an input file than if
you were just typing the document directly.
Quotation marks like
       ``this''
have to be handled specially, as do quotes within
quotes:
       ``\,`this'            % \, separates the double and single quote.
        is what I just
        wrote, not  `that'\,''.

Dashes come in three sizes: an
       intra-word
dash, a medium dash for number ranges like
       1--2,
and a punctuation
       dash---like
this.

A sentence-ending space should be larger than the
space between words within a sentence.  You
sometimes have to type special commands in
conjunction with punctuation characters to get
this right, as in the following sentence.
       Gnats, gnus, etc.\ all  % `\ ' makes an inter-word space.
       begin with G\@.         % \@ marks end-of-sentence punctuation.
You should check the spaces after periods when
reading your output to make sure you haven't
forgotten any special cases.  Generating an
ellipsis
       \ldots\               % `\ ' is needed after `\ldots' because TeX
                             % ignores spaces after command names like \ldots
                             % made from \ + letters.
                             %
                             % Note how a `%' character causes TeX to ignore
                             % the end of the input line, so these blank lines
                             % do not start a new paragraph.
                             %
with the right spacing around the periods requires
a special command.

\LaTeX\ interprets some common characters as
commands, so you must type special commands to
generate them.  These characters include the
following:
       \$ \& \% \# \{ and \}.

In printing, text is usually emphasized with an
       \emph{italic}
type style.

\begin{em}
   A long segment of text can also be emphasized
   in this way.  Text within such a segment can be
   given \emph{additional} emphasis.
\end{em}

It is sometimes necessary to prevent \LaTeX\ from
breaking a line where it might otherwise do so.
This may be at a space, as between the ``Mr.''\ and
``Jones'' in
       ``Mr.~Jones'',        % ~ produces an unbreakable interword space.
or within a word---especially when the word is a
symbol like
       \mbox{\emph{itemnum}}
that makes little sense when hyphenated across
lines.

Footnotes\footnote{This is an example of a footnote.}
pose no problem.

\LaTeX\ is good at typesetting mathematical formulas
like
       \( x-3y + z = 7 \)
or
       \( a_{1} > x^{2n} + y^{2n} > x' \)
or
       \( \ip{A}{B} = \sum_{i} a_{i} b_{i} \).
The spaces you type in a formula are
ignored.  Remember that a letter like
       $x$                   % $ ... $  and  \( ... \)  are equivalent
is a formula when it denotes a mathematical
symbol, and it should be typed as one.

\section{Displayed Text}

Text is displayed by indenting it from the left
margin.  Quotations are commonly displayed.  There
are short quotations
\begin{quote}
   This is a short quotation.  It consists of a
   single paragraph of text.  See how it is formatted.
\end{quote}
and longer ones.
\begin{quotation}
   This is a longer quotation.  It consists of two
   paragraphs of text, neither of which are
   particularly interesting.

   This is the second paragraph of the quotation.  It
   is just as dull as the first paragraph.
\end{quotation}
Another frequently-displayed structure is a list.
The following is an example of an \emph{itemized}
list.
\begin{itemize}
   \item This is the first item of an itemized list.
         Each item in the list is marked with a ``tick''.
         You don't have to worry about what kind of tick
         mark is used.

   \item This is the second item of the list.  It
         contains another list nested inside it.  The inner
         list is an \emph{enumerated} list.
         \begin{enumerate}
            \item This is the first item of an enumerated
                  list that is nested within the itemized list.

            \item This is the second item of the inner list.
                  \LaTeX\ allows you to nest lists deeper than
                  you really should.
         \end{enumerate}
         This is the rest of the second item of the outer
         list.  It is no more interesting than any other
         part of the item.
   \item This is the third item of the list.
\end{itemize}
You can even display poetry.
\begin{verse}
   There is an environment
    for verse \\             % The \\ command separates lines
   Whose features some poets % within a stanza.
   will curse.

                             % One or more blank lines separate stanzas.

   For instead of making\\
   Them do \emph{all} line breaking, \\
   It allows them to put too many words on a line when they'd rather be
   forced to be terse.
\end{verse}

Mathematical formulas may also be displayed.  A
displayed formula
is
one-line long; multiline
formulas require special formatting instructions.
   \[  \ip{\Gamma}{\psi'} = x'' + y^{2} + z_{i}^{n}\]
Don't start a paragraph with a displayed equation,
nor make one a paragraph by itself.

\end{document}               % End of document.

\section{Bench Huge Section 915}
% synthetic bench section 915
% This is a sample LaTeX input file.  (Version of 12 August 2004.)
%
% A '%' character causes TeX to ignore all remaining text on the line,
% and is used for comments like this one.

\documentclass{article}      % Specifies the document class

                             % The preamble begins here.
\title{An Example Document}  % Declares the document's title.
\author{Leslie Lamport}      % Declares the author's name.
\date{January 21, 1994}      % Deleting this command produces today's date.

\newcommand{\ip}[2]{(#1, #2)}
                             % Defines \ip{arg1}{arg2} to mean
                             % (arg1, arg2).

%\newcommand{\ip}[2]{\langle #1 | #2\rangle}
                             % This is an alternative definition of
                             % \ip that is commented out.

\begin{document}             % End of preamble and beginning of text.

\maketitle                   % Produces the title.

This is an example input file.  Comparing it with
the output it generates can show you how to
produce a simple document of your own.

\section{Ordinary Text}      % Produces section heading.  Lower-level
                             % sections are begun with similar
                             % \subsection and \subsubsection commands.

The ends  of words and sentences are marked
  by   spaces. It  doesn't matter how many
spaces    you type; one is as good as 100.  The
end of   a line counts as a space.

One   or more   blank lines denote the  end
of  a paragraph.

Since any number of consecutive spaces are treated
like a single one, the formatting of the input
file makes no difference to
      \LaTeX,                % The \LaTeX command generates the LaTeX logo.
but it makes a difference to you.  When you use
\LaTeX, making your input file as easy to read
as possible will be a great help as you write
your document and when you change it.  This sample
file shows how you can add comments to your own input
file.

Because printing is different from typewriting,
there are a number of things that you have to do
differently when preparing an input file than if
you were just typing the document directly.
Quotation marks like
       ``this''
have to be handled specially, as do quotes within
quotes:
       ``\,`this'            % \, separates the double and single quote.
        is what I just
        wrote, not  `that'\,''.

Dashes come in three sizes: an
       intra-word
dash, a medium dash for number ranges like
       1--2,
and a punctuation
       dash---like
this.

A sentence-ending space should be larger than the
space between words within a sentence.  You
sometimes have to type special commands in
conjunction with punctuation characters to get
this right, as in the following sentence.
       Gnats, gnus, etc.\ all  % `\ ' makes an inter-word space.
       begin with G\@.         % \@ marks end-of-sentence punctuation.
You should check the spaces after periods when
reading your output to make sure you haven't
forgotten any special cases.  Generating an
ellipsis
       \ldots\               % `\ ' is needed after `\ldots' because TeX
                             % ignores spaces after command names like \ldots
                             % made from \ + letters.
                             %
                             % Note how a `%' character causes TeX to ignore
                             % the end of the input line, so these blank lines
                             % do not start a new paragraph.
                             %
with the right spacing around the periods requires
a special command.

\LaTeX\ interprets some common characters as
commands, so you must type special commands to
generate them.  These characters include the
following:
       \$ \& \% \# \{ and \}.

In printing, text is usually emphasized with an
       \emph{italic}
type style.

\begin{em}
   A long segment of text can also be emphasized
   in this way.  Text within such a segment can be
   given \emph{additional} emphasis.
\end{em}

It is sometimes necessary to prevent \LaTeX\ from
breaking a line where it might otherwise do so.
This may be at a space, as between the ``Mr.''\ and
``Jones'' in
       ``Mr.~Jones'',        % ~ produces an unbreakable interword space.
or within a word---especially when the word is a
symbol like
       \mbox{\emph{itemnum}}
that makes little sense when hyphenated across
lines.

Footnotes\footnote{This is an example of a footnote.}
pose no problem.

\LaTeX\ is good at typesetting mathematical formulas
like
       \( x-3y + z = 7 \)
or
       \( a_{1} > x^{2n} + y^{2n} > x' \)
or
       \( \ip{A}{B} = \sum_{i} a_{i} b_{i} \).
The spaces you type in a formula are
ignored.  Remember that a letter like
       $x$                   % $ ... $  and  \( ... \)  are equivalent
is a formula when it denotes a mathematical
symbol, and it should be typed as one.

\section{Displayed Text}

Text is displayed by indenting it from the left
margin.  Quotations are commonly displayed.  There
are short quotations
\begin{quote}
   This is a short quotation.  It consists of a
   single paragraph of text.  See how it is formatted.
\end{quote}
and longer ones.
\begin{quotation}
   This is a longer quotation.  It consists of two
   paragraphs of text, neither of which are
   particularly interesting.

   This is the second paragraph of the quotation.  It
   is just as dull as the first paragraph.
\end{quotation}
Another frequently-displayed structure is a list.
The following is an example of an \emph{itemized}
list.
\begin{itemize}
   \item This is the first item of an itemized list.
         Each item in the list is marked with a ``tick''.
         You don't have to worry about what kind of tick
         mark is used.

   \item This is the second item of the list.  It
         contains another list nested inside it.  The inner
         list is an \emph{enumerated} list.
         \begin{enumerate}
            \item This is the first item of an enumerated
                  list that is nested within the itemized list.

            \item This is the second item of the inner list.
                  \LaTeX\ allows you to nest lists deeper than
                  you really should.
         \end{enumerate}
         This is the rest of the second item of the outer
         list.  It is no more interesting than any other
         part of the item.
   \item This is the third item of the list.
\end{itemize}
You can even display poetry.
\begin{verse}
   There is an environment
    for verse \\             % The \\ command separates lines
   Whose features some poets % within a stanza.
   will curse.

                             % One or more blank lines separate stanzas.

   For instead of making\\
   Them do \emph{all} line breaking, \\
   It allows them to put too many words on a line when they'd rather be
   forced to be terse.
\end{verse}

Mathematical formulas may also be displayed.  A
displayed formula
is
one-line long; multiline
formulas require special formatting instructions.
   \[  \ip{\Gamma}{\psi'} = x'' + y^{2} + z_{i}^{n}\]
Don't start a paragraph with a displayed equation,
nor make one a paragraph by itself.

\end{document}               % End of document.

\section{Bench Huge Section 916}
% synthetic bench section 916
% This is a sample LaTeX input file.  (Version of 12 August 2004.)
%
% A '%' character causes TeX to ignore all remaining text on the line,
% and is used for comments like this one.

\documentclass{article}      % Specifies the document class

                             % The preamble begins here.
\title{An Example Document}  % Declares the document's title.
\author{Leslie Lamport}      % Declares the author's name.
\date{January 21, 1994}      % Deleting this command produces today's date.

\newcommand{\ip}[2]{(#1, #2)}
                             % Defines \ip{arg1}{arg2} to mean
                             % (arg1, arg2).

%\newcommand{\ip}[2]{\langle #1 | #2\rangle}
                             % This is an alternative definition of
                             % \ip that is commented out.

\begin{document}             % End of preamble and beginning of text.

\maketitle                   % Produces the title.

This is an example input file.  Comparing it with
the output it generates can show you how to
produce a simple document of your own.

\section{Ordinary Text}      % Produces section heading.  Lower-level
                             % sections are begun with similar
                             % \subsection and \subsubsection commands.

The ends  of words and sentences are marked
  by   spaces. It  doesn't matter how many
spaces    you type; one is as good as 100.  The
end of   a line counts as a space.

One   or more   blank lines denote the  end
of  a paragraph.

Since any number of consecutive spaces are treated
like a single one, the formatting of the input
file makes no difference to
      \LaTeX,                % The \LaTeX command generates the LaTeX logo.
but it makes a difference to you.  When you use
\LaTeX, making your input file as easy to read
as possible will be a great help as you write
your document and when you change it.  This sample
file shows how you can add comments to your own input
file.

Because printing is different from typewriting,
there are a number of things that you have to do
differently when preparing an input file than if
you were just typing the document directly.
Quotation marks like
       ``this''
have to be handled specially, as do quotes within
quotes:
       ``\,`this'            % \, separates the double and single quote.
        is what I just
        wrote, not  `that'\,''.

Dashes come in three sizes: an
       intra-word
dash, a medium dash for number ranges like
       1--2,
and a punctuation
       dash---like
this.

A sentence-ending space should be larger than the
space between words within a sentence.  You
sometimes have to type special commands in
conjunction with punctuation characters to get
this right, as in the following sentence.
       Gnats, gnus, etc.\ all  % `\ ' makes an inter-word space.
       begin with G\@.         % \@ marks end-of-sentence punctuation.
You should check the spaces after periods when
reading your output to make sure you haven't
forgotten any special cases.  Generating an
ellipsis
       \ldots\               % `\ ' is needed after `\ldots' because TeX
                             % ignores spaces after command names like \ldots
                             % made from \ + letters.
                             %
                             % Note how a `%' character causes TeX to ignore
                             % the end of the input line, so these blank lines
                             % do not start a new paragraph.
                             %
with the right spacing around the periods requires
a special command.

\LaTeX\ interprets some common characters as
commands, so you must type special commands to
generate them.  These characters include the
following:
       \$ \& \% \# \{ and \}.

In printing, text is usually emphasized with an
       \emph{italic}
type style.

\begin{em}
   A long segment of text can also be emphasized
   in this way.  Text within such a segment can be
   given \emph{additional} emphasis.
\end{em}

It is sometimes necessary to prevent \LaTeX\ from
breaking a line where it might otherwise do so.
This may be at a space, as between the ``Mr.''\ and
``Jones'' in
       ``Mr.~Jones'',        % ~ produces an unbreakable interword space.
or within a word---especially when the word is a
symbol like
       \mbox{\emph{itemnum}}
that makes little sense when hyphenated across
lines.

Footnotes\footnote{This is an example of a footnote.}
pose no problem.

\LaTeX\ is good at typesetting mathematical formulas
like
       \( x-3y + z = 7 \)
or
       \( a_{1} > x^{2n} + y^{2n} > x' \)
or
       \( \ip{A}{B} = \sum_{i} a_{i} b_{i} \).
The spaces you type in a formula are
ignored.  Remember that a letter like
       $x$                   % $ ... $  and  \( ... \)  are equivalent
is a formula when it denotes a mathematical
symbol, and it should be typed as one.

\section{Displayed Text}

Text is displayed by indenting it from the left
margin.  Quotations are commonly displayed.  There
are short quotations
\begin{quote}
   This is a short quotation.  It consists of a
   single paragraph of text.  See how it is formatted.
\end{quote}
and longer ones.
\begin{quotation}
   This is a longer quotation.  It consists of two
   paragraphs of text, neither of which are
   particularly interesting.

   This is the second paragraph of the quotation.  It
   is just as dull as the first paragraph.
\end{quotation}
Another frequently-displayed structure is a list.
The following is an example of an \emph{itemized}
list.
\begin{itemize}
   \item This is the first item of an itemized list.
         Each item in the list is marked with a ``tick''.
         You don't have to worry about what kind of tick
         mark is used.

   \item This is the second item of the list.  It
         contains another list nested inside it.  The inner
         list is an \emph{enumerated} list.
         \begin{enumerate}
            \item This is the first item of an enumerated
                  list that is nested within the itemized list.

            \item This is the second item of the inner list.
                  \LaTeX\ allows you to nest lists deeper than
                  you really should.
         \end{enumerate}
         This is the rest of the second item of the outer
         list.  It is no more interesting than any other
         part of the item.
   \item This is the third item of the list.
\end{itemize}
You can even display poetry.
\begin{verse}
   There is an environment
    for verse \\             % The \\ command separates lines
   Whose features some poets % within a stanza.
   will curse.

                             % One or more blank lines separate stanzas.

   For instead of making\\
   Them do \emph{all} line breaking, \\
   It allows them to put too many words on a line when they'd rather be
   forced to be terse.
\end{verse}

Mathematical formulas may also be displayed.  A
displayed formula
is
one-line long; multiline
formulas require special formatting instructions.
   \[  \ip{\Gamma}{\psi'} = x'' + y^{2} + z_{i}^{n}\]
Don't start a paragraph with a displayed equation,
nor make one a paragraph by itself.

\end{document}               % End of document.

\section{Bench Huge Section 917}
% synthetic bench section 917
% This is a sample LaTeX input file.  (Version of 12 August 2004.)
%
% A '%' character causes TeX to ignore all remaining text on the line,
% and is used for comments like this one.

\documentclass{article}      % Specifies the document class

                             % The preamble begins here.
\title{An Example Document}  % Declares the document's title.
\author{Leslie Lamport}      % Declares the author's name.
\date{January 21, 1994}      % Deleting this command produces today's date.

\newcommand{\ip}[2]{(#1, #2)}
                             % Defines \ip{arg1}{arg2} to mean
                             % (arg1, arg2).

%\newcommand{\ip}[2]{\langle #1 | #2\rangle}
                             % This is an alternative definition of
                             % \ip that is commented out.

\begin{document}             % End of preamble and beginning of text.

\maketitle                   % Produces the title.

This is an example input file.  Comparing it with
the output it generates can show you how to
produce a simple document of your own.

\section{Ordinary Text}      % Produces section heading.  Lower-level
                             % sections are begun with similar
                             % \subsection and \subsubsection commands.

The ends  of words and sentences are marked
  by   spaces. It  doesn't matter how many
spaces    you type; one is as good as 100.  The
end of   a line counts as a space.

One   or more   blank lines denote the  end
of  a paragraph.

Since any number of consecutive spaces are treated
like a single one, the formatting of the input
file makes no difference to
      \LaTeX,                % The \LaTeX command generates the LaTeX logo.
but it makes a difference to you.  When you use
\LaTeX, making your input file as easy to read
as possible will be a great help as you write
your document and when you change it.  This sample
file shows how you can add comments to your own input
file.

Because printing is different from typewriting,
there are a number of things that you have to do
differently when preparing an input file than if
you were just typing the document directly.
Quotation marks like
       ``this''
have to be handled specially, as do quotes within
quotes:
       ``\,`this'            % \, separates the double and single quote.
        is what I just
        wrote, not  `that'\,''.

Dashes come in three sizes: an
       intra-word
dash, a medium dash for number ranges like
       1--2,
and a punctuation
       dash---like
this.

A sentence-ending space should be larger than the
space between words within a sentence.  You
sometimes have to type special commands in
conjunction with punctuation characters to get
this right, as in the following sentence.
       Gnats, gnus, etc.\ all  % `\ ' makes an inter-word space.
       begin with G\@.         % \@ marks end-of-sentence punctuation.
You should check the spaces after periods when
reading your output to make sure you haven't
forgotten any special cases.  Generating an
ellipsis
       \ldots\               % `\ ' is needed after `\ldots' because TeX
                             % ignores spaces after command names like \ldots
                             % made from \ + letters.
                             %
                             % Note how a `%' character causes TeX to ignore
                             % the end of the input line, so these blank lines
                             % do not start a new paragraph.
                             %
with the right spacing around the periods requires
a special command.

\LaTeX\ interprets some common characters as
commands, so you must type special commands to
generate them.  These characters include the
following:
       \$ \& \% \# \{ and \}.

In printing, text is usually emphasized with an
       \emph{italic}
type style.

\begin{em}
   A long segment of text can also be emphasized
   in this way.  Text within such a segment can be
   given \emph{additional} emphasis.
\end{em}

It is sometimes necessary to prevent \LaTeX\ from
breaking a line where it might otherwise do so.
This may be at a space, as between the ``Mr.''\ and
``Jones'' in
       ``Mr.~Jones'',        % ~ produces an unbreakable interword space.
or within a word---especially when the word is a
symbol like
       \mbox{\emph{itemnum}}
that makes little sense when hyphenated across
lines.

Footnotes\footnote{This is an example of a footnote.}
pose no problem.

\LaTeX\ is good at typesetting mathematical formulas
like
       \( x-3y + z = 7 \)
or
       \( a_{1} > x^{2n} + y^{2n} > x' \)
or
       \( \ip{A}{B} = \sum_{i} a_{i} b_{i} \).
The spaces you type in a formula are
ignored.  Remember that a letter like
       $x$                   % $ ... $  and  \( ... \)  are equivalent
is a formula when it denotes a mathematical
symbol, and it should be typed as one.

\section{Displayed Text}

Text is displayed by indenting it from the left
margin.  Quotations are commonly displayed.  There
are short quotations
\begin{quote}
   This is a short quotation.  It consists of a
   single paragraph of text.  See how it is formatted.
\end{quote}
and longer ones.
\begin{quotation}
   This is a longer quotation.  It consists of two
   paragraphs of text, neither of which are
   particularly interesting.

   This is the second paragraph of the quotation.  It
   is just as dull as the first paragraph.
\end{quotation}
Another frequently-displayed structure is a list.
The following is an example of an \emph{itemized}
list.
\begin{itemize}
   \item This is the first item of an itemized list.
         Each item in the list is marked with a ``tick''.
         You don't have to worry about what kind of tick
         mark is used.

   \item This is the second item of the list.  It
         contains another list nested inside it.  The inner
         list is an \emph{enumerated} list.
         \begin{enumerate}
            \item This is the first item of an enumerated
                  list that is nested within the itemized list.

            \item This is the second item of the inner list.
                  \LaTeX\ allows you to nest lists deeper than
                  you really should.
         \end{enumerate}
         This is the rest of the second item of the outer
         list.  It is no more interesting than any other
         part of the item.
   \item This is the third item of the list.
\end{itemize}
You can even display poetry.
\begin{verse}
   There is an environment
    for verse \\             % The \\ command separates lines
   Whose features some poets % within a stanza.
   will curse.

                             % One or more blank lines separate stanzas.

   For instead of making\\
   Them do \emph{all} line breaking, \\
   It allows them to put too many words on a line when they'd rather be
   forced to be terse.
\end{verse}

Mathematical formulas may also be displayed.  A
displayed formula
is
one-line long; multiline
formulas require special formatting instructions.
   \[  \ip{\Gamma}{\psi'} = x'' + y^{2} + z_{i}^{n}\]
Don't start a paragraph with a displayed equation,
nor make one a paragraph by itself.

\end{document}               % End of document.

\section{Bench Huge Section 918}
% synthetic bench section 918
% This is a sample LaTeX input file.  (Version of 12 August 2004.)
%
% A '%' character causes TeX to ignore all remaining text on the line,
% and is used for comments like this one.

\documentclass{article}      % Specifies the document class

                             % The preamble begins here.
\title{An Example Document}  % Declares the document's title.
\author{Leslie Lamport}      % Declares the author's name.
\date{January 21, 1994}      % Deleting this command produces today's date.

\newcommand{\ip}[2]{(#1, #2)}
                             % Defines \ip{arg1}{arg2} to mean
                             % (arg1, arg2).

%\newcommand{\ip}[2]{\langle #1 | #2\rangle}
                             % This is an alternative definition of
                             % \ip that is commented out.

\begin{document}             % End of preamble and beginning of text.

\maketitle                   % Produces the title.

This is an example input file.  Comparing it with
the output it generates can show you how to
produce a simple document of your own.

\section{Ordinary Text}      % Produces section heading.  Lower-level
                             % sections are begun with similar
                             % \subsection and \subsubsection commands.

The ends  of words and sentences are marked
  by   spaces. It  doesn't matter how many
spaces    you type; one is as good as 100.  The
end of   a line counts as a space.

One   or more   blank lines denote the  end
of  a paragraph.

Since any number of consecutive spaces are treated
like a single one, the formatting of the input
file makes no difference to
      \LaTeX,                % The \LaTeX command generates the LaTeX logo.
but it makes a difference to you.  When you use
\LaTeX, making your input file as easy to read
as possible will be a great help as you write
your document and when you change it.  This sample
file shows how you can add comments to your own input
file.

Because printing is different from typewriting,
there are a number of things that you have to do
differently when preparing an input file than if
you were just typing the document directly.
Quotation marks like
       ``this''
have to be handled specially, as do quotes within
quotes:
       ``\,`this'            % \, separates the double and single quote.
        is what I just
        wrote, not  `that'\,''.

Dashes come in three sizes: an
       intra-word
dash, a medium dash for number ranges like
       1--2,
and a punctuation
       dash---like
this.

A sentence-ending space should be larger than the
space between words within a sentence.  You
sometimes have to type special commands in
conjunction with punctuation characters to get
this right, as in the following sentence.
       Gnats, gnus, etc.\ all  % `\ ' makes an inter-word space.
       begin with G\@.         % \@ marks end-of-sentence punctuation.
You should check the spaces after periods when
reading your output to make sure you haven't
forgotten any special cases.  Generating an
ellipsis
       \ldots\               % `\ ' is needed after `\ldots' because TeX
                             % ignores spaces after command names like \ldots
                             % made from \ + letters.
                             %
                             % Note how a `%' character causes TeX to ignore
                             % the end of the input line, so these blank lines
                             % do not start a new paragraph.
                             %
with the right spacing around the periods requires
a special command.

\LaTeX\ interprets some common characters as
commands, so you must type special commands to
generate them.  These characters include the
following:
       \$ \& \% \# \{ and \}.

In printing, text is usually emphasized with an
       \emph{italic}
type style.

\begin{em}
   A long segment of text can also be emphasized
   in this way.  Text within such a segment can be
   given \emph{additional} emphasis.
\end{em}

It is sometimes necessary to prevent \LaTeX\ from
breaking a line where it might otherwise do so.
This may be at a space, as between the ``Mr.''\ and
``Jones'' in
       ``Mr.~Jones'',        % ~ produces an unbreakable interword space.
or within a word---especially when the word is a
symbol like
       \mbox{\emph{itemnum}}
that makes little sense when hyphenated across
lines.

Footnotes\footnote{This is an example of a footnote.}
pose no problem.

\LaTeX\ is good at typesetting mathematical formulas
like
       \( x-3y + z = 7 \)
or
       \( a_{1} > x^{2n} + y^{2n} > x' \)
or
       \( \ip{A}{B} = \sum_{i} a_{i} b_{i} \).
The spaces you type in a formula are
ignored.  Remember that a letter like
       $x$                   % $ ... $  and  \( ... \)  are equivalent
is a formula when it denotes a mathematical
symbol, and it should be typed as one.

\section{Displayed Text}

Text is displayed by indenting it from the left
margin.  Quotations are commonly displayed.  There
are short quotations
\begin{quote}
   This is a short quotation.  It consists of a
   single paragraph of text.  See how it is formatted.
\end{quote}
and longer ones.
\begin{quotation}
   This is a longer quotation.  It consists of two
   paragraphs of text, neither of which are
   particularly interesting.

   This is the second paragraph of the quotation.  It
   is just as dull as the first paragraph.
\end{quotation}
Another frequently-displayed structure is a list.
The following is an example of an \emph{itemized}
list.
\begin{itemize}
   \item This is the first item of an itemized list.
         Each item in the list is marked with a ``tick''.
         You don't have to worry about what kind of tick
         mark is used.

   \item This is the second item of the list.  It
         contains another list nested inside it.  The inner
         list is an \emph{enumerated} list.
         \begin{enumerate}
            \item This is the first item of an enumerated
                  list that is nested within the itemized list.

            \item This is the second item of the inner list.
                  \LaTeX\ allows you to nest lists deeper than
                  you really should.
         \end{enumerate}
         This is the rest of the second item of the outer
         list.  It is no more interesting than any other
         part of the item.
   \item This is the third item of the list.
\end{itemize}
You can even display poetry.
\begin{verse}
   There is an environment
    for verse \\             % The \\ command separates lines
   Whose features some poets % within a stanza.
   will curse.

                             % One or more blank lines separate stanzas.

   For instead of making\\
   Them do \emph{all} line breaking, \\
   It allows them to put too many words on a line when they'd rather be
   forced to be terse.
\end{verse}

Mathematical formulas may also be displayed.  A
displayed formula
is
one-line long; multiline
formulas require special formatting instructions.
   \[  \ip{\Gamma}{\psi'} = x'' + y^{2} + z_{i}^{n}\]
Don't start a paragraph with a displayed equation,
nor make one a paragraph by itself.

\end{document}               % End of document.

\section{Bench Huge Section 919}
% synthetic bench section 919
% This is a sample LaTeX input file.  (Version of 12 August 2004.)
%
% A '%' character causes TeX to ignore all remaining text on the line,
% and is used for comments like this one.

\documentclass{article}      % Specifies the document class

                             % The preamble begins here.
\title{An Example Document}  % Declares the document's title.
\author{Leslie Lamport}      % Declares the author's name.
\date{January 21, 1994}      % Deleting this command produces today's date.

\newcommand{\ip}[2]{(#1, #2)}
                             % Defines \ip{arg1}{arg2} to mean
                             % (arg1, arg2).

%\newcommand{\ip}[2]{\langle #1 | #2\rangle}
                             % This is an alternative definition of
                             % \ip that is commented out.

\begin{document}             % End of preamble and beginning of text.

\maketitle                   % Produces the title.

This is an example input file.  Comparing it with
the output it generates can show you how to
produce a simple document of your own.

\section{Ordinary Text}      % Produces section heading.  Lower-level
                             % sections are begun with similar
                             % \subsection and \subsubsection commands.

The ends  of words and sentences are marked
  by   spaces. It  doesn't matter how many
spaces    you type; one is as good as 100.  The
end of   a line counts as a space.

One   or more   blank lines denote the  end
of  a paragraph.

Since any number of consecutive spaces are treated
like a single one, the formatting of the input
file makes no difference to
      \LaTeX,                % The \LaTeX command generates the LaTeX logo.
but it makes a difference to you.  When you use
\LaTeX, making your input file as easy to read
as possible will be a great help as you write
your document and when you change it.  This sample
file shows how you can add comments to your own input
file.

Because printing is different from typewriting,
there are a number of things that you have to do
differently when preparing an input file than if
you were just typing the document directly.
Quotation marks like
       ``this''
have to be handled specially, as do quotes within
quotes:
       ``\,`this'            % \, separates the double and single quote.
        is what I just
        wrote, not  `that'\,''.

Dashes come in three sizes: an
       intra-word
dash, a medium dash for number ranges like
       1--2,
and a punctuation
       dash---like
this.

A sentence-ending space should be larger than the
space between words within a sentence.  You
sometimes have to type special commands in
conjunction with punctuation characters to get
this right, as in the following sentence.
       Gnats, gnus, etc.\ all  % `\ ' makes an inter-word space.
       begin with G\@.         % \@ marks end-of-sentence punctuation.
You should check the spaces after periods when
reading your output to make sure you haven't
forgotten any special cases.  Generating an
ellipsis
       \ldots\               % `\ ' is needed after `\ldots' because TeX
                             % ignores spaces after command names like \ldots
                             % made from \ + letters.
                             %
                             % Note how a `%' character causes TeX to ignore
                             % the end of the input line, so these blank lines
                             % do not start a new paragraph.
                             %
with the right spacing around the periods requires
a special command.

\LaTeX\ interprets some common characters as
commands, so you must type special commands to
generate them.  These characters include the
following:
       \$ \& \% \# \{ and \}.

In printing, text is usually emphasized with an
       \emph{italic}
type style.

\begin{em}
   A long segment of text can also be emphasized
   in this way.  Text within such a segment can be
   given \emph{additional} emphasis.
\end{em}

It is sometimes necessary to prevent \LaTeX\ from
breaking a line where it might otherwise do so.
This may be at a space, as between the ``Mr.''\ and
``Jones'' in
       ``Mr.~Jones'',        % ~ produces an unbreakable interword space.
or within a word---especially when the word is a
symbol like
       \mbox{\emph{itemnum}}
that makes little sense when hyphenated across
lines.

Footnotes\footnote{This is an example of a footnote.}
pose no problem.

\LaTeX\ is good at typesetting mathematical formulas
like
       \( x-3y + z = 7 \)
or
       \( a_{1} > x^{2n} + y^{2n} > x' \)
or
       \( \ip{A}{B} = \sum_{i} a_{i} b_{i} \).
The spaces you type in a formula are
ignored.  Remember that a letter like
       $x$                   % $ ... $  and  \( ... \)  are equivalent
is a formula when it denotes a mathematical
symbol, and it should be typed as one.

\section{Displayed Text}

Text is displayed by indenting it from the left
margin.  Quotations are commonly displayed.  There
are short quotations
\begin{quote}
   This is a short quotation.  It consists of a
   single paragraph of text.  See how it is formatted.
\end{quote}
and longer ones.
\begin{quotation}
   This is a longer quotation.  It consists of two
   paragraphs of text, neither of which are
   particularly interesting.

   This is the second paragraph of the quotation.  It
   is just as dull as the first paragraph.
\end{quotation}
Another frequently-displayed structure is a list.
The following is an example of an \emph{itemized}
list.
\begin{itemize}
   \item This is the first item of an itemized list.
         Each item in the list is marked with a ``tick''.
         You don't have to worry about what kind of tick
         mark is used.

   \item This is the second item of the list.  It
         contains another list nested inside it.  The inner
         list is an \emph{enumerated} list.
         \begin{enumerate}
            \item This is the first item of an enumerated
                  list that is nested within the itemized list.

            \item This is the second item of the inner list.
                  \LaTeX\ allows you to nest lists deeper than
                  you really should.
         \end{enumerate}
         This is the rest of the second item of the outer
         list.  It is no more interesting than any other
         part of the item.
   \item This is the third item of the list.
\end{itemize}
You can even display poetry.
\begin{verse}
   There is an environment
    for verse \\             % The \\ command separates lines
   Whose features some poets % within a stanza.
   will curse.

                             % One or more blank lines separate stanzas.

   For instead of making\\
   Them do \emph{all} line breaking, \\
   It allows them to put too many words on a line when they'd rather be
   forced to be terse.
\end{verse}

Mathematical formulas may also be displayed.  A
displayed formula
is
one-line long; multiline
formulas require special formatting instructions.
   \[  \ip{\Gamma}{\psi'} = x'' + y^{2} + z_{i}^{n}\]
Don't start a paragraph with a displayed equation,
nor make one a paragraph by itself.

\end{document}               % End of document.

\section{Bench Huge Section 920}
% synthetic bench section 920
% This is a sample LaTeX input file.  (Version of 12 August 2004.)
%
% A '%' character causes TeX to ignore all remaining text on the line,
% and is used for comments like this one.

\documentclass{article}      % Specifies the document class

                             % The preamble begins here.
\title{An Example Document}  % Declares the document's title.
\author{Leslie Lamport}      % Declares the author's name.
\date{January 21, 1994}      % Deleting this command produces today's date.

\newcommand{\ip}[2]{(#1, #2)}
                             % Defines \ip{arg1}{arg2} to mean
                             % (arg1, arg2).

%\newcommand{\ip}[2]{\langle #1 | #2\rangle}
                             % This is an alternative definition of
                             % \ip that is commented out.

\begin{document}             % End of preamble and beginning of text.

\maketitle                   % Produces the title.

This is an example input file.  Comparing it with
the output it generates can show you how to
produce a simple document of your own.

\section{Ordinary Text}      % Produces section heading.  Lower-level
                             % sections are begun with similar
                             % \subsection and \subsubsection commands.

The ends  of words and sentences are marked
  by   spaces. It  doesn't matter how many
spaces    you type; one is as good as 100.  The
end of   a line counts as a space.

One   or more   blank lines denote the  end
of  a paragraph.

Since any number of consecutive spaces are treated
like a single one, the formatting of the input
file makes no difference to
      \LaTeX,                % The \LaTeX command generates the LaTeX logo.
but it makes a difference to you.  When you use
\LaTeX, making your input file as easy to read
as possible will be a great help as you write
your document and when you change it.  This sample
file shows how you can add comments to your own input
file.

Because printing is different from typewriting,
there are a number of things that you have to do
differently when preparing an input file than if
you were just typing the document directly.
Quotation marks like
       ``this''
have to be handled specially, as do quotes within
quotes:
       ``\,`this'            % \, separates the double and single quote.
        is what I just
        wrote, not  `that'\,''.

Dashes come in three sizes: an
       intra-word
dash, a medium dash for number ranges like
       1--2,
and a punctuation
       dash---like
this.

A sentence-ending space should be larger than the
space between words within a sentence.  You
sometimes have to type special commands in
conjunction with punctuation characters to get
this right, as in the following sentence.
       Gnats, gnus, etc.\ all  % `\ ' makes an inter-word space.
       begin with G\@.         % \@ marks end-of-sentence punctuation.
You should check the spaces after periods when
reading your output to make sure you haven't
forgotten any special cases.  Generating an
ellipsis
       \ldots\               % `\ ' is needed after `\ldots' because TeX
                             % ignores spaces after command names like \ldots
                             % made from \ + letters.
                             %
                             % Note how a `%' character causes TeX to ignore
                             % the end of the input line, so these blank lines
                             % do not start a new paragraph.
                             %
with the right spacing around the periods requires
a special command.

\LaTeX\ interprets some common characters as
commands, so you must type special commands to
generate them.  These characters include the
following:
       \$ \& \% \# \{ and \}.

In printing, text is usually emphasized with an
       \emph{italic}
type style.

\begin{em}
   A long segment of text can also be emphasized
   in this way.  Text within such a segment can be
   given \emph{additional} emphasis.
\end{em}

It is sometimes necessary to prevent \LaTeX\ from
breaking a line where it might otherwise do so.
This may be at a space, as between the ``Mr.''\ and
``Jones'' in
       ``Mr.~Jones'',        % ~ produces an unbreakable interword space.
or within a word---especially when the word is a
symbol like
       \mbox{\emph{itemnum}}
that makes little sense when hyphenated across
lines.

Footnotes\footnote{This is an example of a footnote.}
pose no problem.

\LaTeX\ is good at typesetting mathematical formulas
like
       \( x-3y + z = 7 \)
or
       \( a_{1} > x^{2n} + y^{2n} > x' \)
or
       \( \ip{A}{B} = \sum_{i} a_{i} b_{i} \).
The spaces you type in a formula are
ignored.  Remember that a letter like
       $x$                   % $ ... $  and  \( ... \)  are equivalent
is a formula when it denotes a mathematical
symbol, and it should be typed as one.

\section{Displayed Text}

Text is displayed by indenting it from the left
margin.  Quotations are commonly displayed.  There
are short quotations
\begin{quote}
   This is a short quotation.  It consists of a
   single paragraph of text.  See how it is formatted.
\end{quote}
and longer ones.
\begin{quotation}
   This is a longer quotation.  It consists of two
   paragraphs of text, neither of which are
   particularly interesting.

   This is the second paragraph of the quotation.  It
   is just as dull as the first paragraph.
\end{quotation}
Another frequently-displayed structure is a list.
The following is an example of an \emph{itemized}
list.
\begin{itemize}
   \item This is the first item of an itemized list.
         Each item in the list is marked with a ``tick''.
         You don't have to worry about what kind of tick
         mark is used.

   \item This is the second item of the list.  It
         contains another list nested inside it.  The inner
         list is an \emph{enumerated} list.
         \begin{enumerate}
            \item This is the first item of an enumerated
                  list that is nested within the itemized list.

            \item This is the second item of the inner list.
                  \LaTeX\ allows you to nest lists deeper than
                  you really should.
         \end{enumerate}
         This is the rest of the second item of the outer
         list.  It is no more interesting than any other
         part of the item.
   \item This is the third item of the list.
\end{itemize}
You can even display poetry.
\begin{verse}
   There is an environment
    for verse \\             % The \\ command separates lines
   Whose features some poets % within a stanza.
   will curse.

                             % One or more blank lines separate stanzas.

   For instead of making\\
   Them do \emph{all} line breaking, \\
   It allows them to put too many words on a line when they'd rather be
   forced to be terse.
\end{verse}

Mathematical formulas may also be displayed.  A
displayed formula
is
one-line long; multiline
formulas require special formatting instructions.
   \[  \ip{\Gamma}{\psi'} = x'' + y^{2} + z_{i}^{n}\]
Don't start a paragraph with a displayed equation,
nor make one a paragraph by itself.

\end{document}               % End of document.

\section{Bench Huge Section 921}
% synthetic bench section 921
% This is a sample LaTeX input file.  (Version of 12 August 2004.)
%
% A '%' character causes TeX to ignore all remaining text on the line,
% and is used for comments like this one.

\documentclass{article}      % Specifies the document class

                             % The preamble begins here.
\title{An Example Document}  % Declares the document's title.
\author{Leslie Lamport}      % Declares the author's name.
\date{January 21, 1994}      % Deleting this command produces today's date.

\newcommand{\ip}[2]{(#1, #2)}
                             % Defines \ip{arg1}{arg2} to mean
                             % (arg1, arg2).

%\newcommand{\ip}[2]{\langle #1 | #2\rangle}
                             % This is an alternative definition of
                             % \ip that is commented out.

\begin{document}             % End of preamble and beginning of text.

\maketitle                   % Produces the title.

This is an example input file.  Comparing it with
the output it generates can show you how to
produce a simple document of your own.

\section{Ordinary Text}      % Produces section heading.  Lower-level
                             % sections are begun with similar
                             % \subsection and \subsubsection commands.

The ends  of words and sentences are marked
  by   spaces. It  doesn't matter how many
spaces    you type; one is as good as 100.  The
end of   a line counts as a space.

One   or more   blank lines denote the  end
of  a paragraph.

Since any number of consecutive spaces are treated
like a single one, the formatting of the input
file makes no difference to
      \LaTeX,                % The \LaTeX command generates the LaTeX logo.
but it makes a difference to you.  When you use
\LaTeX, making your input file as easy to read
as possible will be a great help as you write
your document and when you change it.  This sample
file shows how you can add comments to your own input
file.

Because printing is different from typewriting,
there are a number of things that you have to do
differently when preparing an input file than if
you were just typing the document directly.
Quotation marks like
       ``this''
have to be handled specially, as do quotes within
quotes:
       ``\,`this'            % \, separates the double and single quote.
        is what I just
        wrote, not  `that'\,''.

Dashes come in three sizes: an
       intra-word
dash, a medium dash for number ranges like
       1--2,
and a punctuation
       dash---like
this.

A sentence-ending space should be larger than the
space between words within a sentence.  You
sometimes have to type special commands in
conjunction with punctuation characters to get
this right, as in the following sentence.
       Gnats, gnus, etc.\ all  % `\ ' makes an inter-word space.
       begin with G\@.         % \@ marks end-of-sentence punctuation.
You should check the spaces after periods when
reading your output to make sure you haven't
forgotten any special cases.  Generating an
ellipsis
       \ldots\               % `\ ' is needed after `\ldots' because TeX
                             % ignores spaces after command names like \ldots
                             % made from \ + letters.
                             %
                             % Note how a `%' character causes TeX to ignore
                             % the end of the input line, so these blank lines
                             % do not start a new paragraph.
                             %
with the right spacing around the periods requires
a special command.

\LaTeX\ interprets some common characters as
commands, so you must type special commands to
generate them.  These characters include the
following:
       \$ \& \% \# \{ and \}.

In printing, text is usually emphasized with an
       \emph{italic}
type style.

\begin{em}
   A long segment of text can also be emphasized
   in this way.  Text within such a segment can be
   given \emph{additional} emphasis.
\end{em}

It is sometimes necessary to prevent \LaTeX\ from
breaking a line where it might otherwise do so.
This may be at a space, as between the ``Mr.''\ and
``Jones'' in
       ``Mr.~Jones'',        % ~ produces an unbreakable interword space.
or within a word---especially when the word is a
symbol like
       \mbox{\emph{itemnum}}
that makes little sense when hyphenated across
lines.

Footnotes\footnote{This is an example of a footnote.}
pose no problem.

\LaTeX\ is good at typesetting mathematical formulas
like
       \( x-3y + z = 7 \)
or
       \( a_{1} > x^{2n} + y^{2n} > x' \)
or
       \( \ip{A}{B} = \sum_{i} a_{i} b_{i} \).
The spaces you type in a formula are
ignored.  Remember that a letter like
       $x$                   % $ ... $  and  \( ... \)  are equivalent
is a formula when it denotes a mathematical
symbol, and it should be typed as one.

\section{Displayed Text}

Text is displayed by indenting it from the left
margin.  Quotations are commonly displayed.  There
are short quotations
\begin{quote}
   This is a short quotation.  It consists of a
   single paragraph of text.  See how it is formatted.
\end{quote}
and longer ones.
\begin{quotation}
   This is a longer quotation.  It consists of two
   paragraphs of text, neither of which are
   particularly interesting.

   This is the second paragraph of the quotation.  It
   is just as dull as the first paragraph.
\end{quotation}
Another frequently-displayed structure is a list.
The following is an example of an \emph{itemized}
list.
\begin{itemize}
   \item This is the first item of an itemized list.
         Each item in the list is marked with a ``tick''.
         You don't have to worry about what kind of tick
         mark is used.

   \item This is the second item of the list.  It
         contains another list nested inside it.  The inner
         list is an \emph{enumerated} list.
         \begin{enumerate}
            \item This is the first item of an enumerated
                  list that is nested within the itemized list.

            \item This is the second item of the inner list.
                  \LaTeX\ allows you to nest lists deeper than
                  you really should.
         \end{enumerate}
         This is the rest of the second item of the outer
         list.  It is no more interesting than any other
         part of the item.
   \item This is the third item of the list.
\end{itemize}
You can even display poetry.
\begin{verse}
   There is an environment
    for verse \\             % The \\ command separates lines
   Whose features some poets % within a stanza.
   will curse.

                             % One or more blank lines separate stanzas.

   For instead of making\\
   Them do \emph{all} line breaking, \\
   It allows them to put too many words on a line when they'd rather be
   forced to be terse.
\end{verse}

Mathematical formulas may also be displayed.  A
displayed formula
is
one-line long; multiline
formulas require special formatting instructions.
   \[  \ip{\Gamma}{\psi'} = x'' + y^{2} + z_{i}^{n}\]
Don't start a paragraph with a displayed equation,
nor make one a paragraph by itself.

\end{document}               % End of document.

\section{Bench Huge Section 922}
% synthetic bench section 922
% This is a sample LaTeX input file.  (Version of 12 August 2004.)
%
% A '%' character causes TeX to ignore all remaining text on the line,
% and is used for comments like this one.

\documentclass{article}      % Specifies the document class

                             % The preamble begins here.
\title{An Example Document}  % Declares the document's title.
\author{Leslie Lamport}      % Declares the author's name.
\date{January 21, 1994}      % Deleting this command produces today's date.

\newcommand{\ip}[2]{(#1, #2)}
                             % Defines \ip{arg1}{arg2} to mean
                             % (arg1, arg2).

%\newcommand{\ip}[2]{\langle #1 | #2\rangle}
                             % This is an alternative definition of
                             % \ip that is commented out.

\begin{document}             % End of preamble and beginning of text.

\maketitle                   % Produces the title.

This is an example input file.  Comparing it with
the output it generates can show you how to
produce a simple document of your own.

\section{Ordinary Text}      % Produces section heading.  Lower-level
                             % sections are begun with similar
                             % \subsection and \subsubsection commands.

The ends  of words and sentences are marked
  by   spaces. It  doesn't matter how many
spaces    you type; one is as good as 100.  The
end of   a line counts as a space.

One   or more   blank lines denote the  end
of  a paragraph.

Since any number of consecutive spaces are treated
like a single one, the formatting of the input
file makes no difference to
      \LaTeX,                % The \LaTeX command generates the LaTeX logo.
but it makes a difference to you.  When you use
\LaTeX, making your input file as easy to read
as possible will be a great help as you write
your document and when you change it.  This sample
file shows how you can add comments to your own input
file.

Because printing is different from typewriting,
there are a number of things that you have to do
differently when preparing an input file than if
you were just typing the document directly.
Quotation marks like
       ``this''
have to be handled specially, as do quotes within
quotes:
       ``\,`this'            % \, separates the double and single quote.
        is what I just
        wrote, not  `that'\,''.

Dashes come in three sizes: an
       intra-word
dash, a medium dash for number ranges like
       1--2,
and a punctuation
       dash---like
this.

A sentence-ending space should be larger than the
space between words within a sentence.  You
sometimes have to type special commands in
conjunction with punctuation characters to get
this right, as in the following sentence.
       Gnats, gnus, etc.\ all  % `\ ' makes an inter-word space.
       begin with G\@.         % \@ marks end-of-sentence punctuation.
You should check the spaces after periods when
reading your output to make sure you haven't
forgotten any special cases.  Generating an
ellipsis
       \ldots\               % `\ ' is needed after `\ldots' because TeX
                             % ignores spaces after command names like \ldots
                             % made from \ + letters.
                             %
                             % Note how a `%' character causes TeX to ignore
                             % the end of the input line, so these blank lines
                             % do not start a new paragraph.
                             %
with the right spacing around the periods requires
a special command.

\LaTeX\ interprets some common characters as
commands, so you must type special commands to
generate them.  These characters include the
following:
       \$ \& \% \# \{ and \}.

In printing, text is usually emphasized with an
       \emph{italic}
type style.

\begin{em}
   A long segment of text can also be emphasized
   in this way.  Text within such a segment can be
   given \emph{additional} emphasis.
\end{em}

It is sometimes necessary to prevent \LaTeX\ from
breaking a line where it might otherwise do so.
This may be at a space, as between the ``Mr.''\ and
``Jones'' in
       ``Mr.~Jones'',        % ~ produces an unbreakable interword space.
or within a word---especially when the word is a
symbol like
       \mbox{\emph{itemnum}}
that makes little sense when hyphenated across
lines.

Footnotes\footnote{This is an example of a footnote.}
pose no problem.

\LaTeX\ is good at typesetting mathematical formulas
like
       \( x-3y + z = 7 \)
or
       \( a_{1} > x^{2n} + y^{2n} > x' \)
or
       \( \ip{A}{B} = \sum_{i} a_{i} b_{i} \).
The spaces you type in a formula are
ignored.  Remember that a letter like
       $x$                   % $ ... $  and  \( ... \)  are equivalent
is a formula when it denotes a mathematical
symbol, and it should be typed as one.

\section{Displayed Text}

Text is displayed by indenting it from the left
margin.  Quotations are commonly displayed.  There
are short quotations
\begin{quote}
   This is a short quotation.  It consists of a
   single paragraph of text.  See how it is formatted.
\end{quote}
and longer ones.
\begin{quotation}
   This is a longer quotation.  It consists of two
   paragraphs of text, neither of which are
   particularly interesting.

   This is the second paragraph of the quotation.  It
   is just as dull as the first paragraph.
\end{quotation}
Another frequently-displayed structure is a list.
The following is an example of an \emph{itemized}
list.
\begin{itemize}
   \item This is the first item of an itemized list.
         Each item in the list is marked with a ``tick''.
         You don't have to worry about what kind of tick
         mark is used.

   \item This is the second item of the list.  It
         contains another list nested inside it.  The inner
         list is an \emph{enumerated} list.
         \begin{enumerate}
            \item This is the first item of an enumerated
                  list that is nested within the itemized list.

            \item This is the second item of the inner list.
                  \LaTeX\ allows you to nest lists deeper than
                  you really should.
         \end{enumerate}
         This is the rest of the second item of the outer
         list.  It is no more interesting than any other
         part of the item.
   \item This is the third item of the list.
\end{itemize}
You can even display poetry.
\begin{verse}
   There is an environment
    for verse \\             % The \\ command separates lines
   Whose features some poets % within a stanza.
   will curse.

                             % One or more blank lines separate stanzas.

   For instead of making\\
   Them do \emph{all} line breaking, \\
   It allows them to put too many words on a line when they'd rather be
   forced to be terse.
\end{verse}

Mathematical formulas may also be displayed.  A
displayed formula
is
one-line long; multiline
formulas require special formatting instructions.
   \[  \ip{\Gamma}{\psi'} = x'' + y^{2} + z_{i}^{n}\]
Don't start a paragraph with a displayed equation,
nor make one a paragraph by itself.

\end{document}               % End of document.

\section{Bench Huge Section 923}
% synthetic bench section 923
% This is a sample LaTeX input file.  (Version of 12 August 2004.)
%
% A '%' character causes TeX to ignore all remaining text on the line,
% and is used for comments like this one.

\documentclass{article}      % Specifies the document class

                             % The preamble begins here.
\title{An Example Document}  % Declares the document's title.
\author{Leslie Lamport}      % Declares the author's name.
\date{January 21, 1994}      % Deleting this command produces today's date.

\newcommand{\ip}[2]{(#1, #2)}
                             % Defines \ip{arg1}{arg2} to mean
                             % (arg1, arg2).

%\newcommand{\ip}[2]{\langle #1 | #2\rangle}
                             % This is an alternative definition of
                             % \ip that is commented out.

\begin{document}             % End of preamble and beginning of text.

\maketitle                   % Produces the title.

This is an example input file.  Comparing it with
the output it generates can show you how to
produce a simple document of your own.

\section{Ordinary Text}      % Produces section heading.  Lower-level
                             % sections are begun with similar
                             % \subsection and \subsubsection commands.

The ends  of words and sentences are marked
  by   spaces. It  doesn't matter how many
spaces    you type; one is as good as 100.  The
end of   a line counts as a space.

One   or more   blank lines denote the  end
of  a paragraph.

Since any number of consecutive spaces are treated
like a single one, the formatting of the input
file makes no difference to
      \LaTeX,                % The \LaTeX command generates the LaTeX logo.
but it makes a difference to you.  When you use
\LaTeX, making your input file as easy to read
as possible will be a great help as you write
your document and when you change it.  This sample
file shows how you can add comments to your own input
file.

Because printing is different from typewriting,
there are a number of things that you have to do
differently when preparing an input file than if
you were just typing the document directly.
Quotation marks like
       ``this''
have to be handled specially, as do quotes within
quotes:
       ``\,`this'            % \, separates the double and single quote.
        is what I just
        wrote, not  `that'\,''.

Dashes come in three sizes: an
       intra-word
dash, a medium dash for number ranges like
       1--2,
and a punctuation
       dash---like
this.

A sentence-ending space should be larger than the
space between words within a sentence.  You
sometimes have to type special commands in
conjunction with punctuation characters to get
this right, as in the following sentence.
       Gnats, gnus, etc.\ all  % `\ ' makes an inter-word space.
       begin with G\@.         % \@ marks end-of-sentence punctuation.
You should check the spaces after periods when
reading your output to make sure you haven't
forgotten any special cases.  Generating an
ellipsis
       \ldots\               % `\ ' is needed after `\ldots' because TeX
                             % ignores spaces after command names like \ldots
                             % made from \ + letters.
                             %
                             % Note how a `%' character causes TeX to ignore
                             % the end of the input line, so these blank lines
                             % do not start a new paragraph.
                             %
with the right spacing around the periods requires
a special command.

\LaTeX\ interprets some common characters as
commands, so you must type special commands to
generate them.  These characters include the
following:
       \$ \& \% \# \{ and \}.

In printing, text is usually emphasized with an
       \emph{italic}
type style.

\begin{em}
   A long segment of text can also be emphasized
   in this way.  Text within such a segment can be
   given \emph{additional} emphasis.
\end{em}

It is sometimes necessary to prevent \LaTeX\ from
breaking a line where it might otherwise do so.
This may be at a space, as between the ``Mr.''\ and
``Jones'' in
       ``Mr.~Jones'',        % ~ produces an unbreakable interword space.
or within a word---especially when the word is a
symbol like
       \mbox{\emph{itemnum}}
that makes little sense when hyphenated across
lines.

Footnotes\footnote{This is an example of a footnote.}
pose no problem.

\LaTeX\ is good at typesetting mathematical formulas
like
       \( x-3y + z = 7 \)
or
       \( a_{1} > x^{2n} + y^{2n} > x' \)
or
       \( \ip{A}{B} = \sum_{i} a_{i} b_{i} \).
The spaces you type in a formula are
ignored.  Remember that a letter like
       $x$                   % $ ... $  and  \( ... \)  are equivalent
is a formula when it denotes a mathematical
symbol, and it should be typed as one.

\section{Displayed Text}

Text is displayed by indenting it from the left
margin.  Quotations are commonly displayed.  There
are short quotations
\begin{quote}
   This is a short quotation.  It consists of a
   single paragraph of text.  See how it is formatted.
\end{quote}
and longer ones.
\begin{quotation}
   This is a longer quotation.  It consists of two
   paragraphs of text, neither of which are
   particularly interesting.

   This is the second paragraph of the quotation.  It
   is just as dull as the first paragraph.
\end{quotation}
Another frequently-displayed structure is a list.
The following is an example of an \emph{itemized}
list.
\begin{itemize}
   \item This is the first item of an itemized list.
         Each item in the list is marked with a ``tick''.
         You don't have to worry about what kind of tick
         mark is used.

   \item This is the second item of the list.  It
         contains another list nested inside it.  The inner
         list is an \emph{enumerated} list.
         \begin{enumerate}
            \item This is the first item of an enumerated
                  list that is nested within the itemized list.

            \item This is the second item of the inner list.
                  \LaTeX\ allows you to nest lists deeper than
                  you really should.
         \end{enumerate}
         This is the rest of the second item of the outer
         list.  It is no more interesting than any other
         part of the item.
   \item This is the third item of the list.
\end{itemize}
You can even display poetry.
\begin{verse}
   There is an environment
    for verse \\             % The \\ command separates lines
   Whose features some poets % within a stanza.
   will curse.

                             % One or more blank lines separate stanzas.

   For instead of making\\
   Them do \emph{all} line breaking, \\
   It allows them to put too many words on a line when they'd rather be
   forced to be terse.
\end{verse}

Mathematical formulas may also be displayed.  A
displayed formula
is
one-line long; multiline
formulas require special formatting instructions.
   \[  \ip{\Gamma}{\psi'} = x'' + y^{2} + z_{i}^{n}\]
Don't start a paragraph with a displayed equation,
nor make one a paragraph by itself.

\end{document}               % End of document.

\section{Bench Huge Section 924}
% synthetic bench section 924
% This is a sample LaTeX input file.  (Version of 12 August 2004.)
%
% A '%' character causes TeX to ignore all remaining text on the line,
% and is used for comments like this one.

\documentclass{article}      % Specifies the document class

                             % The preamble begins here.
\title{An Example Document}  % Declares the document's title.
\author{Leslie Lamport}      % Declares the author's name.
\date{January 21, 1994}      % Deleting this command produces today's date.

\newcommand{\ip}[2]{(#1, #2)}
                             % Defines \ip{arg1}{arg2} to mean
                             % (arg1, arg2).

%\newcommand{\ip}[2]{\langle #1 | #2\rangle}
                             % This is an alternative definition of
                             % \ip that is commented out.

\begin{document}             % End of preamble and beginning of text.

\maketitle                   % Produces the title.

This is an example input file.  Comparing it with
the output it generates can show you how to
produce a simple document of your own.

\section{Ordinary Text}      % Produces section heading.  Lower-level
                             % sections are begun with similar
                             % \subsection and \subsubsection commands.

The ends  of words and sentences are marked
  by   spaces. It  doesn't matter how many
spaces    you type; one is as good as 100.  The
end of   a line counts as a space.

One   or more   blank lines denote the  end
of  a paragraph.

Since any number of consecutive spaces are treated
like a single one, the formatting of the input
file makes no difference to
      \LaTeX,                % The \LaTeX command generates the LaTeX logo.
but it makes a difference to you.  When you use
\LaTeX, making your input file as easy to read
as possible will be a great help as you write
your document and when you change it.  This sample
file shows how you can add comments to your own input
file.

Because printing is different from typewriting,
there are a number of things that you have to do
differently when preparing an input file than if
you were just typing the document directly.
Quotation marks like
       ``this''
have to be handled specially, as do quotes within
quotes:
       ``\,`this'            % \, separates the double and single quote.
        is what I just
        wrote, not  `that'\,''.

Dashes come in three sizes: an
       intra-word
dash, a medium dash for number ranges like
       1--2,
and a punctuation
       dash---like
this.

A sentence-ending space should be larger than the
space between words within a sentence.  You
sometimes have to type special commands in
conjunction with punctuation characters to get
this right, as in the following sentence.
       Gnats, gnus, etc.\ all  % `\ ' makes an inter-word space.
       begin with G\@.         % \@ marks end-of-sentence punctuation.
You should check the spaces after periods when
reading your output to make sure you haven't
forgotten any special cases.  Generating an
ellipsis
       \ldots\               % `\ ' is needed after `\ldots' because TeX
                             % ignores spaces after command names like \ldots
                             % made from \ + letters.
                             %
                             % Note how a `%' character causes TeX to ignore
                             % the end of the input line, so these blank lines
                             % do not start a new paragraph.
                             %
with the right spacing around the periods requires
a special command.

\LaTeX\ interprets some common characters as
commands, so you must type special commands to
generate them.  These characters include the
following:
       \$ \& \% \# \{ and \}.

In printing, text is usually emphasized with an
       \emph{italic}
type style.

\begin{em}
   A long segment of text can also be emphasized
   in this way.  Text within such a segment can be
   given \emph{additional} emphasis.
\end{em}

It is sometimes necessary to prevent \LaTeX\ from
breaking a line where it might otherwise do so.
This may be at a space, as between the ``Mr.''\ and
``Jones'' in
       ``Mr.~Jones'',        % ~ produces an unbreakable interword space.
or within a word---especially when the word is a
symbol like
       \mbox{\emph{itemnum}}
that makes little sense when hyphenated across
lines.

Footnotes\footnote{This is an example of a footnote.}
pose no problem.

\LaTeX\ is good at typesetting mathematical formulas
like
       \( x-3y + z = 7 \)
or
       \( a_{1} > x^{2n} + y^{2n} > x' \)
or
       \( \ip{A}{B} = \sum_{i} a_{i} b_{i} \).
The spaces you type in a formula are
ignored.  Remember that a letter like
       $x$                   % $ ... $  and  \( ... \)  are equivalent
is a formula when it denotes a mathematical
symbol, and it should be typed as one.

\section{Displayed Text}

Text is displayed by indenting it from the left
margin.  Quotations are commonly displayed.  There
are short quotations
\begin{quote}
   This is a short quotation.  It consists of a
   single paragraph of text.  See how it is formatted.
\end{quote}
and longer ones.
\begin{quotation}
   This is a longer quotation.  It consists of two
   paragraphs of text, neither of which are
   particularly interesting.

   This is the second paragraph of the quotation.  It
   is just as dull as the first paragraph.
\end{quotation}
Another frequently-displayed structure is a list.
The following is an example of an \emph{itemized}
list.
\begin{itemize}
   \item This is the first item of an itemized list.
         Each item in the list is marked with a ``tick''.
         You don't have to worry about what kind of tick
         mark is used.

   \item This is the second item of the list.  It
         contains another list nested inside it.  The inner
         list is an \emph{enumerated} list.
         \begin{enumerate}
            \item This is the first item of an enumerated
                  list that is nested within the itemized list.

            \item This is the second item of the inner list.
                  \LaTeX\ allows you to nest lists deeper than
                  you really should.
         \end{enumerate}
         This is the rest of the second item of the outer
         list.  It is no more interesting than any other
         part of the item.
   \item This is the third item of the list.
\end{itemize}
You can even display poetry.
\begin{verse}
   There is an environment
    for verse \\             % The \\ command separates lines
   Whose features some poets % within a stanza.
   will curse.

                             % One or more blank lines separate stanzas.

   For instead of making\\
   Them do \emph{all} line breaking, \\
   It allows them to put too many words on a line when they'd rather be
   forced to be terse.
\end{verse}

Mathematical formulas may also be displayed.  A
displayed formula
is
one-line long; multiline
formulas require special formatting instructions.
   \[  \ip{\Gamma}{\psi'} = x'' + y^{2} + z_{i}^{n}\]
Don't start a paragraph with a displayed equation,
nor make one a paragraph by itself.

\end{document}               % End of document.

\section{Bench Huge Section 925}
% synthetic bench section 925
% This is a sample LaTeX input file.  (Version of 12 August 2004.)
%
% A '%' character causes TeX to ignore all remaining text on the line,
% and is used for comments like this one.

\documentclass{article}      % Specifies the document class

                             % The preamble begins here.
\title{An Example Document}  % Declares the document's title.
\author{Leslie Lamport}      % Declares the author's name.
\date{January 21, 1994}      % Deleting this command produces today's date.

\newcommand{\ip}[2]{(#1, #2)}
                             % Defines \ip{arg1}{arg2} to mean
                             % (arg1, arg2).

%\newcommand{\ip}[2]{\langle #1 | #2\rangle}
                             % This is an alternative definition of
                             % \ip that is commented out.

\begin{document}             % End of preamble and beginning of text.

\maketitle                   % Produces the title.

This is an example input file.  Comparing it with
the output it generates can show you how to
produce a simple document of your own.

\section{Ordinary Text}      % Produces section heading.  Lower-level
                             % sections are begun with similar
                             % \subsection and \subsubsection commands.

The ends  of words and sentences are marked
  by   spaces. It  doesn't matter how many
spaces    you type; one is as good as 100.  The
end of   a line counts as a space.

One   or more   blank lines denote the  end
of  a paragraph.

Since any number of consecutive spaces are treated
like a single one, the formatting of the input
file makes no difference to
      \LaTeX,                % The \LaTeX command generates the LaTeX logo.
but it makes a difference to you.  When you use
\LaTeX, making your input file as easy to read
as possible will be a great help as you write
your document and when you change it.  This sample
file shows how you can add comments to your own input
file.

Because printing is different from typewriting,
there are a number of things that you have to do
differently when preparing an input file than if
you were just typing the document directly.
Quotation marks like
       ``this''
have to be handled specially, as do quotes within
quotes:
       ``\,`this'            % \, separates the double and single quote.
        is what I just
        wrote, not  `that'\,''.

Dashes come in three sizes: an
       intra-word
dash, a medium dash for number ranges like
       1--2,
and a punctuation
       dash---like
this.

A sentence-ending space should be larger than the
space between words within a sentence.  You
sometimes have to type special commands in
conjunction with punctuation characters to get
this right, as in the following sentence.
       Gnats, gnus, etc.\ all  % `\ ' makes an inter-word space.
       begin with G\@.         % \@ marks end-of-sentence punctuation.
You should check the spaces after periods when
reading your output to make sure you haven't
forgotten any special cases.  Generating an
ellipsis
       \ldots\               % `\ ' is needed after `\ldots' because TeX
                             % ignores spaces after command names like \ldots
                             % made from \ + letters.
                             %
                             % Note how a `%' character causes TeX to ignore
                             % the end of the input line, so these blank lines
                             % do not start a new paragraph.
                             %
with the right spacing around the periods requires
a special command.

\LaTeX\ interprets some common characters as
commands, so you must type special commands to
generate them.  These characters include the
following:
       \$ \& \% \# \{ and \}.

In printing, text is usually emphasized with an
       \emph{italic}
type style.

\begin{em}
   A long segment of text can also be emphasized
   in this way.  Text within such a segment can be
   given \emph{additional} emphasis.
\end{em}

It is sometimes necessary to prevent \LaTeX\ from
breaking a line where it might otherwise do so.
This may be at a space, as between the ``Mr.''\ and
``Jones'' in
       ``Mr.~Jones'',        % ~ produces an unbreakable interword space.
or within a word---especially when the word is a
symbol like
       \mbox{\emph{itemnum}}
that makes little sense when hyphenated across
lines.

Footnotes\footnote{This is an example of a footnote.}
pose no problem.

\LaTeX\ is good at typesetting mathematical formulas
like
       \( x-3y + z = 7 \)
or
       \( a_{1} > x^{2n} + y^{2n} > x' \)
or
       \( \ip{A}{B} = \sum_{i} a_{i} b_{i} \).
The spaces you type in a formula are
ignored.  Remember that a letter like
       $x$                   % $ ... $  and  \( ... \)  are equivalent
is a formula when it denotes a mathematical
symbol, and it should be typed as one.

\section{Displayed Text}

Text is displayed by indenting it from the left
margin.  Quotations are commonly displayed.  There
are short quotations
\begin{quote}
   This is a short quotation.  It consists of a
   single paragraph of text.  See how it is formatted.
\end{quote}
and longer ones.
\begin{quotation}
   This is a longer quotation.  It consists of two
   paragraphs of text, neither of which are
   particularly interesting.

   This is the second paragraph of the quotation.  It
   is just as dull as the first paragraph.
\end{quotation}
Another frequently-displayed structure is a list.
The following is an example of an \emph{itemized}
list.
\begin{itemize}
   \item This is the first item of an itemized list.
         Each item in the list is marked with a ``tick''.
         You don't have to worry about what kind of tick
         mark is used.

   \item This is the second item of the list.  It
         contains another list nested inside it.  The inner
         list is an \emph{enumerated} list.
         \begin{enumerate}
            \item This is the first item of an enumerated
                  list that is nested within the itemized list.

            \item This is the second item of the inner list.
                  \LaTeX\ allows you to nest lists deeper than
                  you really should.
         \end{enumerate}
         This is the rest of the second item of the outer
         list.  It is no more interesting than any other
         part of the item.
   \item This is the third item of the list.
\end{itemize}
You can even display poetry.
\begin{verse}
   There is an environment
    for verse \\             % The \\ command separates lines
   Whose features some poets % within a stanza.
   will curse.

                             % One or more blank lines separate stanzas.

   For instead of making\\
   Them do \emph{all} line breaking, \\
   It allows them to put too many words on a line when they'd rather be
   forced to be terse.
\end{verse}

Mathematical formulas may also be displayed.  A
displayed formula
is
one-line long; multiline
formulas require special formatting instructions.
   \[  \ip{\Gamma}{\psi'} = x'' + y^{2} + z_{i}^{n}\]
Don't start a paragraph with a displayed equation,
nor make one a paragraph by itself.

\end{document}               % End of document.

\section{Bench Huge Section 926}
% synthetic bench section 926
% This is a sample LaTeX input file.  (Version of 12 August 2004.)
%
% A '%' character causes TeX to ignore all remaining text on the line,
% and is used for comments like this one.

\documentclass{article}      % Specifies the document class

                             % The preamble begins here.
\title{An Example Document}  % Declares the document's title.
\author{Leslie Lamport}      % Declares the author's name.
\date{January 21, 1994}      % Deleting this command produces today's date.

\newcommand{\ip}[2]{(#1, #2)}
                             % Defines \ip{arg1}{arg2} to mean
                             % (arg1, arg2).

%\newcommand{\ip}[2]{\langle #1 | #2\rangle}
                             % This is an alternative definition of
                             % \ip that is commented out.

\begin{document}             % End of preamble and beginning of text.

\maketitle                   % Produces the title.

This is an example input file.  Comparing it with
the output it generates can show you how to
produce a simple document of your own.

\section{Ordinary Text}      % Produces section heading.  Lower-level
                             % sections are begun with similar
                             % \subsection and \subsubsection commands.

The ends  of words and sentences are marked
  by   spaces. It  doesn't matter how many
spaces    you type; one is as good as 100.  The
end of   a line counts as a space.

One   or more   blank lines denote the  end
of  a paragraph.

Since any number of consecutive spaces are treated
like a single one, the formatting of the input
file makes no difference to
      \LaTeX,                % The \LaTeX command generates the LaTeX logo.
but it makes a difference to you.  When you use
\LaTeX, making your input file as easy to read
as possible will be a great help as you write
your document and when you change it.  This sample
file shows how you can add comments to your own input
file.

Because printing is different from typewriting,
there are a number of things that you have to do
differently when preparing an input file than if
you were just typing the document directly.
Quotation marks like
       ``this''
have to be handled specially, as do quotes within
quotes:
       ``\,`this'            % \, separates the double and single quote.
        is what I just
        wrote, not  `that'\,''.

Dashes come in three sizes: an
       intra-word
dash, a medium dash for number ranges like
       1--2,
and a punctuation
       dash---like
this.

A sentence-ending space should be larger than the
space between words within a sentence.  You
sometimes have to type special commands in
conjunction with punctuation characters to get
this right, as in the following sentence.
       Gnats, gnus, etc.\ all  % `\ ' makes an inter-word space.
       begin with G\@.         % \@ marks end-of-sentence punctuation.
You should check the spaces after periods when
reading your output to make sure you haven't
forgotten any special cases.  Generating an
ellipsis
       \ldots\               % `\ ' is needed after `\ldots' because TeX
                             % ignores spaces after command names like \ldots
                             % made from \ + letters.
                             %
                             % Note how a `%' character causes TeX to ignore
                             % the end of the input line, so these blank lines
                             % do not start a new paragraph.
                             %
with the right spacing around the periods requires
a special command.

\LaTeX\ interprets some common characters as
commands, so you must type special commands to
generate them.  These characters include the
following:
       \$ \& \% \# \{ and \}.

In printing, text is usually emphasized with an
       \emph{italic}
type style.

\begin{em}
   A long segment of text can also be emphasized
   in this way.  Text within such a segment can be
   given \emph{additional} emphasis.
\end{em}

It is sometimes necessary to prevent \LaTeX\ from
breaking a line where it might otherwise do so.
This may be at a space, as between the ``Mr.''\ and
``Jones'' in
       ``Mr.~Jones'',        % ~ produces an unbreakable interword space.
or within a word---especially when the word is a
symbol like
       \mbox{\emph{itemnum}}
that makes little sense when hyphenated across
lines.

Footnotes\footnote{This is an example of a footnote.}
pose no problem.

\LaTeX\ is good at typesetting mathematical formulas
like
       \( x-3y + z = 7 \)
or
       \( a_{1} > x^{2n} + y^{2n} > x' \)
or
       \( \ip{A}{B} = \sum_{i} a_{i} b_{i} \).
The spaces you type in a formula are
ignored.  Remember that a letter like
       $x$                   % $ ... $  and  \( ... \)  are equivalent
is a formula when it denotes a mathematical
symbol, and it should be typed as one.

\section{Displayed Text}

Text is displayed by indenting it from the left
margin.  Quotations are commonly displayed.  There
are short quotations
\begin{quote}
   This is a short quotation.  It consists of a
   single paragraph of text.  See how it is formatted.
\end{quote}
and longer ones.
\begin{quotation}
   This is a longer quotation.  It consists of two
   paragraphs of text, neither of which are
   particularly interesting.

   This is the second paragraph of the quotation.  It
   is just as dull as the first paragraph.
\end{quotation}
Another frequently-displayed structure is a list.
The following is an example of an \emph{itemized}
list.
\begin{itemize}
   \item This is the first item of an itemized list.
         Each item in the list is marked with a ``tick''.
         You don't have to worry about what kind of tick
         mark is used.

   \item This is the second item of the list.  It
         contains another list nested inside it.  The inner
         list is an \emph{enumerated} list.
         \begin{enumerate}
            \item This is the first item of an enumerated
                  list that is nested within the itemized list.

            \item This is the second item of the inner list.
                  \LaTeX\ allows you to nest lists deeper than
                  you really should.
         \end{enumerate}
         This is the rest of the second item of the outer
         list.  It is no more interesting than any other
         part of the item.
   \item This is the third item of the list.
\end{itemize}
You can even display poetry.
\begin{verse}
   There is an environment
    for verse \\             % The \\ command separates lines
   Whose features some poets % within a stanza.
   will curse.

                             % One or more blank lines separate stanzas.

   For instead of making\\
   Them do \emph{all} line breaking, \\
   It allows them to put too many words on a line when they'd rather be
   forced to be terse.
\end{verse}

Mathematical formulas may also be displayed.  A
displayed formula
is
one-line long; multiline
formulas require special formatting instructions.
   \[  \ip{\Gamma}{\psi'} = x'' + y^{2} + z_{i}^{n}\]
Don't start a paragraph with a displayed equation,
nor make one a paragraph by itself.

\end{document}               % End of document.

\section{Bench Huge Section 927}
% synthetic bench section 927
% This is a sample LaTeX input file.  (Version of 12 August 2004.)
%
% A '%' character causes TeX to ignore all remaining text on the line,
% and is used for comments like this one.

\documentclass{article}      % Specifies the document class

                             % The preamble begins here.
\title{An Example Document}  % Declares the document's title.
\author{Leslie Lamport}      % Declares the author's name.
\date{January 21, 1994}      % Deleting this command produces today's date.

\newcommand{\ip}[2]{(#1, #2)}
                             % Defines \ip{arg1}{arg2} to mean
                             % (arg1, arg2).

%\newcommand{\ip}[2]{\langle #1 | #2\rangle}
                             % This is an alternative definition of
                             % \ip that is commented out.

\begin{document}             % End of preamble and beginning of text.

\maketitle                   % Produces the title.

This is an example input file.  Comparing it with
the output it generates can show you how to
produce a simple document of your own.

\section{Ordinary Text}      % Produces section heading.  Lower-level
                             % sections are begun with similar
                             % \subsection and \subsubsection commands.

The ends  of words and sentences are marked
  by   spaces. It  doesn't matter how many
spaces    you type; one is as good as 100.  The
end of   a line counts as a space.

One   or more   blank lines denote the  end
of  a paragraph.

Since any number of consecutive spaces are treated
like a single one, the formatting of the input
file makes no difference to
      \LaTeX,                % The \LaTeX command generates the LaTeX logo.
but it makes a difference to you.  When you use
\LaTeX, making your input file as easy to read
as possible will be a great help as you write
your document and when you change it.  This sample
file shows how you can add comments to your own input
file.

Because printing is different from typewriting,
there are a number of things that you have to do
differently when preparing an input file than if
you were just typing the document directly.
Quotation marks like
       ``this''
have to be handled specially, as do quotes within
quotes:
       ``\,`this'            % \, separates the double and single quote.
        is what I just
        wrote, not  `that'\,''.

Dashes come in three sizes: an
       intra-word
dash, a medium dash for number ranges like
       1--2,
and a punctuation
       dash---like
this.

A sentence-ending space should be larger than the
space between words within a sentence.  You
sometimes have to type special commands in
conjunction with punctuation characters to get
this right, as in the following sentence.
       Gnats, gnus, etc.\ all  % `\ ' makes an inter-word space.
       begin with G\@.         % \@ marks end-of-sentence punctuation.
You should check the spaces after periods when
reading your output to make sure you haven't
forgotten any special cases.  Generating an
ellipsis
       \ldots\               % `\ ' is needed after `\ldots' because TeX
                             % ignores spaces after command names like \ldots
                             % made from \ + letters.
                             %
                             % Note how a `%' character causes TeX to ignore
                             % the end of the input line, so these blank lines
                             % do not start a new paragraph.
                             %
with the right spacing around the periods requires
a special command.

\LaTeX\ interprets some common characters as
commands, so you must type special commands to
generate them.  These characters include the
following:
       \$ \& \% \# \{ and \}.

In printing, text is usually emphasized with an
       \emph{italic}
type style.

\begin{em}
   A long segment of text can also be emphasized
   in this way.  Text within such a segment can be
   given \emph{additional} emphasis.
\end{em}

It is sometimes necessary to prevent \LaTeX\ from
breaking a line where it might otherwise do so.
This may be at a space, as between the ``Mr.''\ and
``Jones'' in
       ``Mr.~Jones'',        % ~ produces an unbreakable interword space.
or within a word---especially when the word is a
symbol like
       \mbox{\emph{itemnum}}
that makes little sense when hyphenated across
lines.

Footnotes\footnote{This is an example of a footnote.}
pose no problem.

\LaTeX\ is good at typesetting mathematical formulas
like
       \( x-3y + z = 7 \)
or
       \( a_{1} > x^{2n} + y^{2n} > x' \)
or
       \( \ip{A}{B} = \sum_{i} a_{i} b_{i} \).
The spaces you type in a formula are
ignored.  Remember that a letter like
       $x$                   % $ ... $  and  \( ... \)  are equivalent
is a formula when it denotes a mathematical
symbol, and it should be typed as one.

\section{Displayed Text}

Text is displayed by indenting it from the left
margin.  Quotations are commonly displayed.  There
are short quotations
\begin{quote}
   This is a short quotation.  It consists of a
   single paragraph of text.  See how it is formatted.
\end{quote}
and longer ones.
\begin{quotation}
   This is a longer quotation.  It consists of two
   paragraphs of text, neither of which are
   particularly interesting.

   This is the second paragraph of the quotation.  It
   is just as dull as the first paragraph.
\end{quotation}
Another frequently-displayed structure is a list.
The following is an example of an \emph{itemized}
list.
\begin{itemize}
   \item This is the first item of an itemized list.
         Each item in the list is marked with a ``tick''.
         You don't have to worry about what kind of tick
         mark is used.

   \item This is the second item of the list.  It
         contains another list nested inside it.  The inner
         list is an \emph{enumerated} list.
         \begin{enumerate}
            \item This is the first item of an enumerated
                  list that is nested within the itemized list.

            \item This is the second item of the inner list.
                  \LaTeX\ allows you to nest lists deeper than
                  you really should.
         \end{enumerate}
         This is the rest of the second item of the outer
         list.  It is no more interesting than any other
         part of the item.
   \item This is the third item of the list.
\end{itemize}
You can even display poetry.
\begin{verse}
   There is an environment
    for verse \\             % The \\ command separates lines
   Whose features some poets % within a stanza.
   will curse.

                             % One or more blank lines separate stanzas.

   For instead of making\\
   Them do \emph{all} line breaking, \\
   It allows them to put too many words on a line when they'd rather be
   forced to be terse.
\end{verse}

Mathematical formulas may also be displayed.  A
displayed formula
is
one-line long; multiline
formulas require special formatting instructions.
   \[  \ip{\Gamma}{\psi'} = x'' + y^{2} + z_{i}^{n}\]
Don't start a paragraph with a displayed equation,
nor make one a paragraph by itself.

\end{document}               % End of document.

\section{Bench Huge Section 928}
% synthetic bench section 928
% This is a sample LaTeX input file.  (Version of 12 August 2004.)
%
% A '%' character causes TeX to ignore all remaining text on the line,
% and is used for comments like this one.

\documentclass{article}      % Specifies the document class

                             % The preamble begins here.
\title{An Example Document}  % Declares the document's title.
\author{Leslie Lamport}      % Declares the author's name.
\date{January 21, 1994}      % Deleting this command produces today's date.

\newcommand{\ip}[2]{(#1, #2)}
                             % Defines \ip{arg1}{arg2} to mean
                             % (arg1, arg2).

%\newcommand{\ip}[2]{\langle #1 | #2\rangle}
                             % This is an alternative definition of
                             % \ip that is commented out.

\begin{document}             % End of preamble and beginning of text.

\maketitle                   % Produces the title.

This is an example input file.  Comparing it with
the output it generates can show you how to
produce a simple document of your own.

\section{Ordinary Text}      % Produces section heading.  Lower-level
                             % sections are begun with similar
                             % \subsection and \subsubsection commands.

The ends  of words and sentences are marked
  by   spaces. It  doesn't matter how many
spaces    you type; one is as good as 100.  The
end of   a line counts as a space.

One   or more   blank lines denote the  end
of  a paragraph.

Since any number of consecutive spaces are treated
like a single one, the formatting of the input
file makes no difference to
      \LaTeX,                % The \LaTeX command generates the LaTeX logo.
but it makes a difference to you.  When you use
\LaTeX, making your input file as easy to read
as possible will be a great help as you write
your document and when you change it.  This sample
file shows how you can add comments to your own input
file.

Because printing is different from typewriting,
there are a number of things that you have to do
differently when preparing an input file than if
you were just typing the document directly.
Quotation marks like
       ``this''
have to be handled specially, as do quotes within
quotes:
       ``\,`this'            % \, separates the double and single quote.
        is what I just
        wrote, not  `that'\,''.

Dashes come in three sizes: an
       intra-word
dash, a medium dash for number ranges like
       1--2,
and a punctuation
       dash---like
this.

A sentence-ending space should be larger than the
space between words within a sentence.  You
sometimes have to type special commands in
conjunction with punctuation characters to get
this right, as in the following sentence.
       Gnats, gnus, etc.\ all  % `\ ' makes an inter-word space.
       begin with G\@.         % \@ marks end-of-sentence punctuation.
You should check the spaces after periods when
reading your output to make sure you haven't
forgotten any special cases.  Generating an
ellipsis
       \ldots\               % `\ ' is needed after `\ldots' because TeX
                             % ignores spaces after command names like \ldots
                             % made from \ + letters.
                             %
                             % Note how a `%' character causes TeX to ignore
                             % the end of the input line, so these blank lines
                             % do not start a new paragraph.
                             %
with the right spacing around the periods requires
a special command.

\LaTeX\ interprets some common characters as
commands, so you must type special commands to
generate them.  These characters include the
following:
       \$ \& \% \# \{ and \}.

In printing, text is usually emphasized with an
       \emph{italic}
type style.

\begin{em}
   A long segment of text can also be emphasized
   in this way.  Text within such a segment can be
   given \emph{additional} emphasis.
\end{em}

It is sometimes necessary to prevent \LaTeX\ from
breaking a line where it might otherwise do so.
This may be at a space, as between the ``Mr.''\ and
``Jones'' in
       ``Mr.~Jones'',        % ~ produces an unbreakable interword space.
or within a word---especially when the word is a
symbol like
       \mbox{\emph{itemnum}}
that makes little sense when hyphenated across
lines.

Footnotes\footnote{This is an example of a footnote.}
pose no problem.

\LaTeX\ is good at typesetting mathematical formulas
like
       \( x-3y + z = 7 \)
or
       \( a_{1} > x^{2n} + y^{2n} > x' \)
or
       \( \ip{A}{B} = \sum_{i} a_{i} b_{i} \).
The spaces you type in a formula are
ignored.  Remember that a letter like
       $x$                   % $ ... $  and  \( ... \)  are equivalent
is a formula when it denotes a mathematical
symbol, and it should be typed as one.

\section{Displayed Text}

Text is displayed by indenting it from the left
margin.  Quotations are commonly displayed.  There
are short quotations
\begin{quote}
   This is a short quotation.  It consists of a
   single paragraph of text.  See how it is formatted.
\end{quote}
and longer ones.
\begin{quotation}
   This is a longer quotation.  It consists of two
   paragraphs of text, neither of which are
   particularly interesting.

   This is the second paragraph of the quotation.  It
   is just as dull as the first paragraph.
\end{quotation}
Another frequently-displayed structure is a list.
The following is an example of an \emph{itemized}
list.
\begin{itemize}
   \item This is the first item of an itemized list.
         Each item in the list is marked with a ``tick''.
         You don't have to worry about what kind of tick
         mark is used.

   \item This is the second item of the list.  It
         contains another list nested inside it.  The inner
         list is an \emph{enumerated} list.
         \begin{enumerate}
            \item This is the first item of an enumerated
                  list that is nested within the itemized list.

            \item This is the second item of the inner list.
                  \LaTeX\ allows you to nest lists deeper than
                  you really should.
         \end{enumerate}
         This is the rest of the second item of the outer
         list.  It is no more interesting than any other
         part of the item.
   \item This is the third item of the list.
\end{itemize}
You can even display poetry.
\begin{verse}
   There is an environment
    for verse \\             % The \\ command separates lines
   Whose features some poets % within a stanza.
   will curse.

                             % One or more blank lines separate stanzas.

   For instead of making\\
   Them do \emph{all} line breaking, \\
   It allows them to put too many words on a line when they'd rather be
   forced to be terse.
\end{verse}

Mathematical formulas may also be displayed.  A
displayed formula
is
one-line long; multiline
formulas require special formatting instructions.
   \[  \ip{\Gamma}{\psi'} = x'' + y^{2} + z_{i}^{n}\]
Don't start a paragraph with a displayed equation,
nor make one a paragraph by itself.

\end{document}               % End of document.

\section{Bench Huge Section 929}
% synthetic bench section 929
% This is a sample LaTeX input file.  (Version of 12 August 2004.)
%
% A '%' character causes TeX to ignore all remaining text on the line,
% and is used for comments like this one.

\documentclass{article}      % Specifies the document class

                             % The preamble begins here.
\title{An Example Document}  % Declares the document's title.
\author{Leslie Lamport}      % Declares the author's name.
\date{January 21, 1994}      % Deleting this command produces today's date.

\newcommand{\ip}[2]{(#1, #2)}
                             % Defines \ip{arg1}{arg2} to mean
                             % (arg1, arg2).

%\newcommand{\ip}[2]{\langle #1 | #2\rangle}
                             % This is an alternative definition of
                             % \ip that is commented out.

\begin{document}             % End of preamble and beginning of text.

\maketitle                   % Produces the title.

This is an example input file.  Comparing it with
the output it generates can show you how to
produce a simple document of your own.

\section{Ordinary Text}      % Produces section heading.  Lower-level
                             % sections are begun with similar
                             % \subsection and \subsubsection commands.

The ends  of words and sentences are marked
  by   spaces. It  doesn't matter how many
spaces    you type; one is as good as 100.  The
end of   a line counts as a space.

One   or more   blank lines denote the  end
of  a paragraph.

Since any number of consecutive spaces are treated
like a single one, the formatting of the input
file makes no difference to
      \LaTeX,                % The \LaTeX command generates the LaTeX logo.
but it makes a difference to you.  When you use
\LaTeX, making your input file as easy to read
as possible will be a great help as you write
your document and when you change it.  This sample
file shows how you can add comments to your own input
file.

Because printing is different from typewriting,
there are a number of things that you have to do
differently when preparing an input file than if
you were just typing the document directly.
Quotation marks like
       ``this''
have to be handled specially, as do quotes within
quotes:
       ``\,`this'            % \, separates the double and single quote.
        is what I just
        wrote, not  `that'\,''.

Dashes come in three sizes: an
       intra-word
dash, a medium dash for number ranges like
       1--2,
and a punctuation
       dash---like
this.

A sentence-ending space should be larger than the
space between words within a sentence.  You
sometimes have to type special commands in
conjunction with punctuation characters to get
this right, as in the following sentence.
       Gnats, gnus, etc.\ all  % `\ ' makes an inter-word space.
       begin with G\@.         % \@ marks end-of-sentence punctuation.
You should check the spaces after periods when
reading your output to make sure you haven't
forgotten any special cases.  Generating an
ellipsis
       \ldots\               % `\ ' is needed after `\ldots' because TeX
                             % ignores spaces after command names like \ldots
                             % made from \ + letters.
                             %
                             % Note how a `%' character causes TeX to ignore
                             % the end of the input line, so these blank lines
                             % do not start a new paragraph.
                             %
with the right spacing around the periods requires
a special command.

\LaTeX\ interprets some common characters as
commands, so you must type special commands to
generate them.  These characters include the
following:
       \$ \& \% \# \{ and \}.

In printing, text is usually emphasized with an
       \emph{italic}
type style.

\begin{em}
   A long segment of text can also be emphasized
   in this way.  Text within such a segment can be
   given \emph{additional} emphasis.
\end{em}

It is sometimes necessary to prevent \LaTeX\ from
breaking a line where it might otherwise do so.
This may be at a space, as between the ``Mr.''\ and
``Jones'' in
       ``Mr.~Jones'',        % ~ produces an unbreakable interword space.
or within a word---especially when the word is a
symbol like
       \mbox{\emph{itemnum}}
that makes little sense when hyphenated across
lines.

Footnotes\footnote{This is an example of a footnote.}
pose no problem.

\LaTeX\ is good at typesetting mathematical formulas
like
       \( x-3y + z = 7 \)
or
       \( a_{1} > x^{2n} + y^{2n} > x' \)
or
       \( \ip{A}{B} = \sum_{i} a_{i} b_{i} \).
The spaces you type in a formula are
ignored.  Remember that a letter like
       $x$                   % $ ... $  and  \( ... \)  are equivalent
is a formula when it denotes a mathematical
symbol, and it should be typed as one.

\section{Displayed Text}

Text is displayed by indenting it from the left
margin.  Quotations are commonly displayed.  There
are short quotations
\begin{quote}
   This is a short quotation.  It consists of a
   single paragraph of text.  See how it is formatted.
\end{quote}
and longer ones.
\begin{quotation}
   This is a longer quotation.  It consists of two
   paragraphs of text, neither of which are
   particularly interesting.

   This is the second paragraph of the quotation.  It
   is just as dull as the first paragraph.
\end{quotation}
Another frequently-displayed structure is a list.
The following is an example of an \emph{itemized}
list.
\begin{itemize}
   \item This is the first item of an itemized list.
         Each item in the list is marked with a ``tick''.
         You don't have to worry about what kind of tick
         mark is used.

   \item This is the second item of the list.  It
         contains another list nested inside it.  The inner
         list is an \emph{enumerated} list.
         \begin{enumerate}
            \item This is the first item of an enumerated
                  list that is nested within the itemized list.

            \item This is the second item of the inner list.
                  \LaTeX\ allows you to nest lists deeper than
                  you really should.
         \end{enumerate}
         This is the rest of the second item of the outer
         list.  It is no more interesting than any other
         part of the item.
   \item This is the third item of the list.
\end{itemize}
You can even display poetry.
\begin{verse}
   There is an environment
    for verse \\             % The \\ command separates lines
   Whose features some poets % within a stanza.
   will curse.

                             % One or more blank lines separate stanzas.

   For instead of making\\
   Them do \emph{all} line breaking, \\
   It allows them to put too many words on a line when they'd rather be
   forced to be terse.
\end{verse}

Mathematical formulas may also be displayed.  A
displayed formula
is
one-line long; multiline
formulas require special formatting instructions.
   \[  \ip{\Gamma}{\psi'} = x'' + y^{2} + z_{i}^{n}\]
Don't start a paragraph with a displayed equation,
nor make one a paragraph by itself.

\end{document}               % End of document.

\section{Bench Huge Section 930}
% synthetic bench section 930
% This is a sample LaTeX input file.  (Version of 12 August 2004.)
%
% A '%' character causes TeX to ignore all remaining text on the line,
% and is used for comments like this one.

\documentclass{article}      % Specifies the document class

                             % The preamble begins here.
\title{An Example Document}  % Declares the document's title.
\author{Leslie Lamport}      % Declares the author's name.
\date{January 21, 1994}      % Deleting this command produces today's date.

\newcommand{\ip}[2]{(#1, #2)}
                             % Defines \ip{arg1}{arg2} to mean
                             % (arg1, arg2).

%\newcommand{\ip}[2]{\langle #1 | #2\rangle}
                             % This is an alternative definition of
                             % \ip that is commented out.

\begin{document}             % End of preamble and beginning of text.

\maketitle                   % Produces the title.

This is an example input file.  Comparing it with
the output it generates can show you how to
produce a simple document of your own.

\section{Ordinary Text}      % Produces section heading.  Lower-level
                             % sections are begun with similar
                             % \subsection and \subsubsection commands.

The ends  of words and sentences are marked
  by   spaces. It  doesn't matter how many
spaces    you type; one is as good as 100.  The
end of   a line counts as a space.

One   or more   blank lines denote the  end
of  a paragraph.

Since any number of consecutive spaces are treated
like a single one, the formatting of the input
file makes no difference to
      \LaTeX,                % The \LaTeX command generates the LaTeX logo.
but it makes a difference to you.  When you use
\LaTeX, making your input file as easy to read
as possible will be a great help as you write
your document and when you change it.  This sample
file shows how you can add comments to your own input
file.

Because printing is different from typewriting,
there are a number of things that you have to do
differently when preparing an input file than if
you were just typing the document directly.
Quotation marks like
       ``this''
have to be handled specially, as do quotes within
quotes:
       ``\,`this'            % \, separates the double and single quote.
        is what I just
        wrote, not  `that'\,''.

Dashes come in three sizes: an
       intra-word
dash, a medium dash for number ranges like
       1--2,
and a punctuation
       dash---like
this.

A sentence-ending space should be larger than the
space between words within a sentence.  You
sometimes have to type special commands in
conjunction with punctuation characters to get
this right, as in the following sentence.
       Gnats, gnus, etc.\ all  % `\ ' makes an inter-word space.
       begin with G\@.         % \@ marks end-of-sentence punctuation.
You should check the spaces after periods when
reading your output to make sure you haven't
forgotten any special cases.  Generating an
ellipsis
       \ldots\               % `\ ' is needed after `\ldots' because TeX
                             % ignores spaces after command names like \ldots
                             % made from \ + letters.
                             %
                             % Note how a `%' character causes TeX to ignore
                             % the end of the input line, so these blank lines
                             % do not start a new paragraph.
                             %
with the right spacing around the periods requires
a special command.

\LaTeX\ interprets some common characters as
commands, so you must type special commands to
generate them.  These characters include the
following:
       \$ \& \% \# \{ and \}.

In printing, text is usually emphasized with an
       \emph{italic}
type style.

\begin{em}
   A long segment of text can also be emphasized
   in this way.  Text within such a segment can be
   given \emph{additional} emphasis.
\end{em}

It is sometimes necessary to prevent \LaTeX\ from
breaking a line where it might otherwise do so.
This may be at a space, as between the ``Mr.''\ and
``Jones'' in
       ``Mr.~Jones'',        % ~ produces an unbreakable interword space.
or within a word---especially when the word is a
symbol like
       \mbox{\emph{itemnum}}
that makes little sense when hyphenated across
lines.

Footnotes\footnote{This is an example of a footnote.}
pose no problem.

\LaTeX\ is good at typesetting mathematical formulas
like
       \( x-3y + z = 7 \)
or
       \( a_{1} > x^{2n} + y^{2n} > x' \)
or
       \( \ip{A}{B} = \sum_{i} a_{i} b_{i} \).
The spaces you type in a formula are
ignored.  Remember that a letter like
       $x$                   % $ ... $  and  \( ... \)  are equivalent
is a formula when it denotes a mathematical
symbol, and it should be typed as one.

\section{Displayed Text}

Text is displayed by indenting it from the left
margin.  Quotations are commonly displayed.  There
are short quotations
\begin{quote}
   This is a short quotation.  It consists of a
   single paragraph of text.  See how it is formatted.
\end{quote}
and longer ones.
\begin{quotation}
   This is a longer quotation.  It consists of two
   paragraphs of text, neither of which are
   particularly interesting.

   This is the second paragraph of the quotation.  It
   is just as dull as the first paragraph.
\end{quotation}
Another frequently-displayed structure is a list.
The following is an example of an \emph{itemized}
list.
\begin{itemize}
   \item This is the first item of an itemized list.
         Each item in the list is marked with a ``tick''.
         You don't have to worry about what kind of tick
         mark is used.

   \item This is the second item of the list.  It
         contains another list nested inside it.  The inner
         list is an \emph{enumerated} list.
         \begin{enumerate}
            \item This is the first item of an enumerated
                  list that is nested within the itemized list.

            \item This is the second item of the inner list.
                  \LaTeX\ allows you to nest lists deeper than
                  you really should.
         \end{enumerate}
         This is the rest of the second item of the outer
         list.  It is no more interesting than any other
         part of the item.
   \item This is the third item of the list.
\end{itemize}
You can even display poetry.
\begin{verse}
   There is an environment
    for verse \\             % The \\ command separates lines
   Whose features some poets % within a stanza.
   will curse.

                             % One or more blank lines separate stanzas.

   For instead of making\\
   Them do \emph{all} line breaking, \\
   It allows them to put too many words on a line when they'd rather be
   forced to be terse.
\end{verse}

Mathematical formulas may also be displayed.  A
displayed formula
is
one-line long; multiline
formulas require special formatting instructions.
   \[  \ip{\Gamma}{\psi'} = x'' + y^{2} + z_{i}^{n}\]
Don't start a paragraph with a displayed equation,
nor make one a paragraph by itself.

\end{document}               % End of document.

\section{Bench Huge Section 931}
% synthetic bench section 931
% This is a sample LaTeX input file.  (Version of 12 August 2004.)
%
% A '%' character causes TeX to ignore all remaining text on the line,
% and is used for comments like this one.

\documentclass{article}      % Specifies the document class

                             % The preamble begins here.
\title{An Example Document}  % Declares the document's title.
\author{Leslie Lamport}      % Declares the author's name.
\date{January 21, 1994}      % Deleting this command produces today's date.

\newcommand{\ip}[2]{(#1, #2)}
                             % Defines \ip{arg1}{arg2} to mean
                             % (arg1, arg2).

%\newcommand{\ip}[2]{\langle #1 | #2\rangle}
                             % This is an alternative definition of
                             % \ip that is commented out.

\begin{document}             % End of preamble and beginning of text.

\maketitle                   % Produces the title.

This is an example input file.  Comparing it with
the output it generates can show you how to
produce a simple document of your own.

\section{Ordinary Text}      % Produces section heading.  Lower-level
                             % sections are begun with similar
                             % \subsection and \subsubsection commands.

The ends  of words and sentences are marked
  by   spaces. It  doesn't matter how many
spaces    you type; one is as good as 100.  The
end of   a line counts as a space.

One   or more   blank lines denote the  end
of  a paragraph.

Since any number of consecutive spaces are treated
like a single one, the formatting of the input
file makes no difference to
      \LaTeX,                % The \LaTeX command generates the LaTeX logo.
but it makes a difference to you.  When you use
\LaTeX, making your input file as easy to read
as possible will be a great help as you write
your document and when you change it.  This sample
file shows how you can add comments to your own input
file.

Because printing is different from typewriting,
there are a number of things that you have to do
differently when preparing an input file than if
you were just typing the document directly.
Quotation marks like
       ``this''
have to be handled specially, as do quotes within
quotes:
       ``\,`this'            % \, separates the double and single quote.
        is what I just
        wrote, not  `that'\,''.

Dashes come in three sizes: an
       intra-word
dash, a medium dash for number ranges like
       1--2,
and a punctuation
       dash---like
this.

A sentence-ending space should be larger than the
space between words within a sentence.  You
sometimes have to type special commands in
conjunction with punctuation characters to get
this right, as in the following sentence.
       Gnats, gnus, etc.\ all  % `\ ' makes an inter-word space.
       begin with G\@.         % \@ marks end-of-sentence punctuation.
You should check the spaces after periods when
reading your output to make sure you haven't
forgotten any special cases.  Generating an
ellipsis
       \ldots\               % `\ ' is needed after `\ldots' because TeX
                             % ignores spaces after command names like \ldots
                             % made from \ + letters.
                             %
                             % Note how a `%' character causes TeX to ignore
                             % the end of the input line, so these blank lines
                             % do not start a new paragraph.
                             %
with the right spacing around the periods requires
a special command.

\LaTeX\ interprets some common characters as
commands, so you must type special commands to
generate them.  These characters include the
following:
       \$ \& \% \# \{ and \}.

In printing, text is usually emphasized with an
       \emph{italic}
type style.

\begin{em}
   A long segment of text can also be emphasized
   in this way.  Text within such a segment can be
   given \emph{additional} emphasis.
\end{em}

It is sometimes necessary to prevent \LaTeX\ from
breaking a line where it might otherwise do so.
This may be at a space, as between the ``Mr.''\ and
``Jones'' in
       ``Mr.~Jones'',        % ~ produces an unbreakable interword space.
or within a word---especially when the word is a
symbol like
       \mbox{\emph{itemnum}}
that makes little sense when hyphenated across
lines.

Footnotes\footnote{This is an example of a footnote.}
pose no problem.

\LaTeX\ is good at typesetting mathematical formulas
like
       \( x-3y + z = 7 \)
or
       \( a_{1} > x^{2n} + y^{2n} > x' \)
or
       \( \ip{A}{B} = \sum_{i} a_{i} b_{i} \).
The spaces you type in a formula are
ignored.  Remember that a letter like
       $x$                   % $ ... $  and  \( ... \)  are equivalent
is a formula when it denotes a mathematical
symbol, and it should be typed as one.

\section{Displayed Text}

Text is displayed by indenting it from the left
margin.  Quotations are commonly displayed.  There
are short quotations
\begin{quote}
   This is a short quotation.  It consists of a
   single paragraph of text.  See how it is formatted.
\end{quote}
and longer ones.
\begin{quotation}
   This is a longer quotation.  It consists of two
   paragraphs of text, neither of which are
   particularly interesting.

   This is the second paragraph of the quotation.  It
   is just as dull as the first paragraph.
\end{quotation}
Another frequently-displayed structure is a list.
The following is an example of an \emph{itemized}
list.
\begin{itemize}
   \item This is the first item of an itemized list.
         Each item in the list is marked with a ``tick''.
         You don't have to worry about what kind of tick
         mark is used.

   \item This is the second item of the list.  It
         contains another list nested inside it.  The inner
         list is an \emph{enumerated} list.
         \begin{enumerate}
            \item This is the first item of an enumerated
                  list that is nested within the itemized list.

            \item This is the second item of the inner list.
                  \LaTeX\ allows you to nest lists deeper than
                  you really should.
         \end{enumerate}
         This is the rest of the second item of the outer
         list.  It is no more interesting than any other
         part of the item.
   \item This is the third item of the list.
\end{itemize}
You can even display poetry.
\begin{verse}
   There is an environment
    for verse \\             % The \\ command separates lines
   Whose features some poets % within a stanza.
   will curse.

                             % One or more blank lines separate stanzas.

   For instead of making\\
   Them do \emph{all} line breaking, \\
   It allows them to put too many words on a line when they'd rather be
   forced to be terse.
\end{verse}

Mathematical formulas may also be displayed.  A
displayed formula
is
one-line long; multiline
formulas require special formatting instructions.
   \[  \ip{\Gamma}{\psi'} = x'' + y^{2} + z_{i}^{n}\]
Don't start a paragraph with a displayed equation,
nor make one a paragraph by itself.

\end{document}               % End of document.

\section{Bench Huge Section 932}
% synthetic bench section 932
% This is a sample LaTeX input file.  (Version of 12 August 2004.)
%
% A '%' character causes TeX to ignore all remaining text on the line,
% and is used for comments like this one.

\documentclass{article}      % Specifies the document class

                             % The preamble begins here.
\title{An Example Document}  % Declares the document's title.
\author{Leslie Lamport}      % Declares the author's name.
\date{January 21, 1994}      % Deleting this command produces today's date.

\newcommand{\ip}[2]{(#1, #2)}
                             % Defines \ip{arg1}{arg2} to mean
                             % (arg1, arg2).

%\newcommand{\ip}[2]{\langle #1 | #2\rangle}
                             % This is an alternative definition of
                             % \ip that is commented out.

\begin{document}             % End of preamble and beginning of text.

\maketitle                   % Produces the title.

This is an example input file.  Comparing it with
the output it generates can show you how to
produce a simple document of your own.

\section{Ordinary Text}      % Produces section heading.  Lower-level
                             % sections are begun with similar
                             % \subsection and \subsubsection commands.

The ends  of words and sentences are marked
  by   spaces. It  doesn't matter how many
spaces    you type; one is as good as 100.  The
end of   a line counts as a space.

One   or more   blank lines denote the  end
of  a paragraph.

Since any number of consecutive spaces are treated
like a single one, the formatting of the input
file makes no difference to
      \LaTeX,                % The \LaTeX command generates the LaTeX logo.
but it makes a difference to you.  When you use
\LaTeX, making your input file as easy to read
as possible will be a great help as you write
your document and when you change it.  This sample
file shows how you can add comments to your own input
file.

Because printing is different from typewriting,
there are a number of things that you have to do
differently when preparing an input file than if
you were just typing the document directly.
Quotation marks like
       ``this''
have to be handled specially, as do quotes within
quotes:
       ``\,`this'            % \, separates the double and single quote.
        is what I just
        wrote, not  `that'\,''.

Dashes come in three sizes: an
       intra-word
dash, a medium dash for number ranges like
       1--2,
and a punctuation
       dash---like
this.

A sentence-ending space should be larger than the
space between words within a sentence.  You
sometimes have to type special commands in
conjunction with punctuation characters to get
this right, as in the following sentence.
       Gnats, gnus, etc.\ all  % `\ ' makes an inter-word space.
       begin with G\@.         % \@ marks end-of-sentence punctuation.
You should check the spaces after periods when
reading your output to make sure you haven't
forgotten any special cases.  Generating an
ellipsis
       \ldots\               % `\ ' is needed after `\ldots' because TeX
                             % ignores spaces after command names like \ldots
                             % made from \ + letters.
                             %
                             % Note how a `%' character causes TeX to ignore
                             % the end of the input line, so these blank lines
                             % do not start a new paragraph.
                             %
with the right spacing around the periods requires
a special command.

\LaTeX\ interprets some common characters as
commands, so you must type special commands to
generate them.  These characters include the
following:
       \$ \& \% \# \{ and \}.

In printing, text is usually emphasized with an
       \emph{italic}
type style.

\begin{em}
   A long segment of text can also be emphasized
   in this way.  Text within such a segment can be
   given \emph{additional} emphasis.
\end{em}

It is sometimes necessary to prevent \LaTeX\ from
breaking a line where it might otherwise do so.
This may be at a space, as between the ``Mr.''\ and
``Jones'' in
       ``Mr.~Jones'',        % ~ produces an unbreakable interword space.
or within a word---especially when the word is a
symbol like
       \mbox{\emph{itemnum}}
that makes little sense when hyphenated across
lines.

Footnotes\footnote{This is an example of a footnote.}
pose no problem.

\LaTeX\ is good at typesetting mathematical formulas
like
       \( x-3y + z = 7 \)
or
       \( a_{1} > x^{2n} + y^{2n} > x' \)
or
       \( \ip{A}{B} = \sum_{i} a_{i} b_{i} \).
The spaces you type in a formula are
ignored.  Remember that a letter like
       $x$                   % $ ... $  and  \( ... \)  are equivalent
is a formula when it denotes a mathematical
symbol, and it should be typed as one.

\section{Displayed Text}

Text is displayed by indenting it from the left
margin.  Quotations are commonly displayed.  There
are short quotations
\begin{quote}
   This is a short quotation.  It consists of a
   single paragraph of text.  See how it is formatted.
\end{quote}
and longer ones.
\begin{quotation}
   This is a longer quotation.  It consists of two
   paragraphs of text, neither of which are
   particularly interesting.

   This is the second paragraph of the quotation.  It
   is just as dull as the first paragraph.
\end{quotation}
Another frequently-displayed structure is a list.
The following is an example of an \emph{itemized}
list.
\begin{itemize}
   \item This is the first item of an itemized list.
         Each item in the list is marked with a ``tick''.
         You don't have to worry about what kind of tick
         mark is used.

   \item This is the second item of the list.  It
         contains another list nested inside it.  The inner
         list is an \emph{enumerated} list.
         \begin{enumerate}
            \item This is the first item of an enumerated
                  list that is nested within the itemized list.

            \item This is the second item of the inner list.
                  \LaTeX\ allows you to nest lists deeper than
                  you really should.
         \end{enumerate}
         This is the rest of the second item of the outer
         list.  It is no more interesting than any other
         part of the item.
   \item This is the third item of the list.
\end{itemize}
You can even display poetry.
\begin{verse}
   There is an environment
    for verse \\             % The \\ command separates lines
   Whose features some poets % within a stanza.
   will curse.

                             % One or more blank lines separate stanzas.

   For instead of making\\
   Them do \emph{all} line breaking, \\
   It allows them to put too many words on a line when they'd rather be
   forced to be terse.
\end{verse}

Mathematical formulas may also be displayed.  A
displayed formula
is
one-line long; multiline
formulas require special formatting instructions.
   \[  \ip{\Gamma}{\psi'} = x'' + y^{2} + z_{i}^{n}\]
Don't start a paragraph with a displayed equation,
nor make one a paragraph by itself.

\end{document}               % End of document.

\section{Bench Huge Section 933}
% synthetic bench section 933
% This is a sample LaTeX input file.  (Version of 12 August 2004.)
%
% A '%' character causes TeX to ignore all remaining text on the line,
% and is used for comments like this one.

\documentclass{article}      % Specifies the document class

                             % The preamble begins here.
\title{An Example Document}  % Declares the document's title.
\author{Leslie Lamport}      % Declares the author's name.
\date{January 21, 1994}      % Deleting this command produces today's date.

\newcommand{\ip}[2]{(#1, #2)}
                             % Defines \ip{arg1}{arg2} to mean
                             % (arg1, arg2).

%\newcommand{\ip}[2]{\langle #1 | #2\rangle}
                             % This is an alternative definition of
                             % \ip that is commented out.

\begin{document}             % End of preamble and beginning of text.

\maketitle                   % Produces the title.

This is an example input file.  Comparing it with
the output it generates can show you how to
produce a simple document of your own.

\section{Ordinary Text}      % Produces section heading.  Lower-level
                             % sections are begun with similar
                             % \subsection and \subsubsection commands.

The ends  of words and sentences are marked
  by   spaces. It  doesn't matter how many
spaces    you type; one is as good as 100.  The
end of   a line counts as a space.

One   or more   blank lines denote the  end
of  a paragraph.

Since any number of consecutive spaces are treated
like a single one, the formatting of the input
file makes no difference to
      \LaTeX,                % The \LaTeX command generates the LaTeX logo.
but it makes a difference to you.  When you use
\LaTeX, making your input file as easy to read
as possible will be a great help as you write
your document and when you change it.  This sample
file shows how you can add comments to your own input
file.

Because printing is different from typewriting,
there are a number of things that you have to do
differently when preparing an input file than if
you were just typing the document directly.
Quotation marks like
       ``this''
have to be handled specially, as do quotes within
quotes:
       ``\,`this'            % \, separates the double and single quote.
        is what I just
        wrote, not  `that'\,''.

Dashes come in three sizes: an
       intra-word
dash, a medium dash for number ranges like
       1--2,
and a punctuation
       dash---like
this.

A sentence-ending space should be larger than the
space between words within a sentence.  You
sometimes have to type special commands in
conjunction with punctuation characters to get
this right, as in the following sentence.
       Gnats, gnus, etc.\ all  % `\ ' makes an inter-word space.
       begin with G\@.         % \@ marks end-of-sentence punctuation.
You should check the spaces after periods when
reading your output to make sure you haven't
forgotten any special cases.  Generating an
ellipsis
       \ldots\               % `\ ' is needed after `\ldots' because TeX
                             % ignores spaces after command names like \ldots
                             % made from \ + letters.
                             %
                             % Note how a `%' character causes TeX to ignore
                             % the end of the input line, so these blank lines
                             % do not start a new paragraph.
                             %
with the right spacing around the periods requires
a special command.

\LaTeX\ interprets some common characters as
commands, so you must type special commands to
generate them.  These characters include the
following:
       \$ \& \% \# \{ and \}.

In printing, text is usually emphasized with an
       \emph{italic}
type style.

\begin{em}
   A long segment of text can also be emphasized
   in this way.  Text within such a segment can be
   given \emph{additional} emphasis.
\end{em}

It is sometimes necessary to prevent \LaTeX\ from
breaking a line where it might otherwise do so.
This may be at a space, as between the ``Mr.''\ and
``Jones'' in
       ``Mr.~Jones'',        % ~ produces an unbreakable interword space.
or within a word---especially when the word is a
symbol like
       \mbox{\emph{itemnum}}
that makes little sense when hyphenated across
lines.

Footnotes\footnote{This is an example of a footnote.}
pose no problem.

\LaTeX\ is good at typesetting mathematical formulas
like
       \( x-3y + z = 7 \)
or
       \( a_{1} > x^{2n} + y^{2n} > x' \)
or
       \( \ip{A}{B} = \sum_{i} a_{i} b_{i} \).
The spaces you type in a formula are
ignored.  Remember that a letter like
       $x$                   % $ ... $  and  \( ... \)  are equivalent
is a formula when it denotes a mathematical
symbol, and it should be typed as one.

\section{Displayed Text}

Text is displayed by indenting it from the left
margin.  Quotations are commonly displayed.  There
are short quotations
\begin{quote}
   This is a short quotation.  It consists of a
   single paragraph of text.  See how it is formatted.
\end{quote}
and longer ones.
\begin{quotation}
   This is a longer quotation.  It consists of two
   paragraphs of text, neither of which are
   particularly interesting.

   This is the second paragraph of the quotation.  It
   is just as dull as the first paragraph.
\end{quotation}
Another frequently-displayed structure is a list.
The following is an example of an \emph{itemized}
list.
\begin{itemize}
   \item This is the first item of an itemized list.
         Each item in the list is marked with a ``tick''.
         You don't have to worry about what kind of tick
         mark is used.

   \item This is the second item of the list.  It
         contains another list nested inside it.  The inner
         list is an \emph{enumerated} list.
         \begin{enumerate}
            \item This is the first item of an enumerated
                  list that is nested within the itemized list.

            \item This is the second item of the inner list.
                  \LaTeX\ allows you to nest lists deeper than
                  you really should.
         \end{enumerate}
         This is the rest of the second item of the outer
         list.  It is no more interesting than any other
         part of the item.
   \item This is the third item of the list.
\end{itemize}
You can even display poetry.
\begin{verse}
   There is an environment
    for verse \\             % The \\ command separates lines
   Whose features some poets % within a stanza.
   will curse.

                             % One or more blank lines separate stanzas.

   For instead of making\\
   Them do \emph{all} line breaking, \\
   It allows them to put too many words on a line when they'd rather be
   forced to be terse.
\end{verse}

Mathematical formulas may also be displayed.  A
displayed formula
is
one-line long; multiline
formulas require special formatting instructions.
   \[  \ip{\Gamma}{\psi'} = x'' + y^{2} + z_{i}^{n}\]
Don't start a paragraph with a displayed equation,
nor make one a paragraph by itself.

\end{document}               % End of document.

\section{Bench Huge Section 934}
% synthetic bench section 934
% This is a sample LaTeX input file.  (Version of 12 August 2004.)
%
% A '%' character causes TeX to ignore all remaining text on the line,
% and is used for comments like this one.

\documentclass{article}      % Specifies the document class

                             % The preamble begins here.
\title{An Example Document}  % Declares the document's title.
\author{Leslie Lamport}      % Declares the author's name.
\date{January 21, 1994}      % Deleting this command produces today's date.

\newcommand{\ip}[2]{(#1, #2)}
                             % Defines \ip{arg1}{arg2} to mean
                             % (arg1, arg2).

%\newcommand{\ip}[2]{\langle #1 | #2\rangle}
                             % This is an alternative definition of
                             % \ip that is commented out.

\begin{document}             % End of preamble and beginning of text.

\maketitle                   % Produces the title.

This is an example input file.  Comparing it with
the output it generates can show you how to
produce a simple document of your own.

\section{Ordinary Text}      % Produces section heading.  Lower-level
                             % sections are begun with similar
                             % \subsection and \subsubsection commands.

The ends  of words and sentences are marked
  by   spaces. It  doesn't matter how many
spaces    you type; one is as good as 100.  The
end of   a line counts as a space.

One   or more   blank lines denote the  end
of  a paragraph.

Since any number of consecutive spaces are treated
like a single one, the formatting of the input
file makes no difference to
      \LaTeX,                % The \LaTeX command generates the LaTeX logo.
but it makes a difference to you.  When you use
\LaTeX, making your input file as easy to read
as possible will be a great help as you write
your document and when you change it.  This sample
file shows how you can add comments to your own input
file.

Because printing is different from typewriting,
there are a number of things that you have to do
differently when preparing an input file than if
you were just typing the document directly.
Quotation marks like
       ``this''
have to be handled specially, as do quotes within
quotes:
       ``\,`this'            % \, separates the double and single quote.
        is what I just
        wrote, not  `that'\,''.

Dashes come in three sizes: an
       intra-word
dash, a medium dash for number ranges like
       1--2,
and a punctuation
       dash---like
this.

A sentence-ending space should be larger than the
space between words within a sentence.  You
sometimes have to type special commands in
conjunction with punctuation characters to get
this right, as in the following sentence.
       Gnats, gnus, etc.\ all  % `\ ' makes an inter-word space.
       begin with G\@.         % \@ marks end-of-sentence punctuation.
You should check the spaces after periods when
reading your output to make sure you haven't
forgotten any special cases.  Generating an
ellipsis
       \ldots\               % `\ ' is needed after `\ldots' because TeX
                             % ignores spaces after command names like \ldots
                             % made from \ + letters.
                             %
                             % Note how a `%' character causes TeX to ignore
                             % the end of the input line, so these blank lines
                             % do not start a new paragraph.
                             %
with the right spacing around the periods requires
a special command.

\LaTeX\ interprets some common characters as
commands, so you must type special commands to
generate them.  These characters include the
following:
       \$ \& \% \# \{ and \}.

In printing, text is usually emphasized with an
       \emph{italic}
type style.

\begin{em}
   A long segment of text can also be emphasized
   in this way.  Text within such a segment can be
   given \emph{additional} emphasis.
\end{em}

It is sometimes necessary to prevent \LaTeX\ from
breaking a line where it might otherwise do so.
This may be at a space, as between the ``Mr.''\ and
``Jones'' in
       ``Mr.~Jones'',        % ~ produces an unbreakable interword space.
or within a word---especially when the word is a
symbol like
       \mbox{\emph{itemnum}}
that makes little sense when hyphenated across
lines.

Footnotes\footnote{This is an example of a footnote.}
pose no problem.

\LaTeX\ is good at typesetting mathematical formulas
like
       \( x-3y + z = 7 \)
or
       \( a_{1} > x^{2n} + y^{2n} > x' \)
or
       \( \ip{A}{B} = \sum_{i} a_{i} b_{i} \).
The spaces you type in a formula are
ignored.  Remember that a letter like
       $x$                   % $ ... $  and  \( ... \)  are equivalent
is a formula when it denotes a mathematical
symbol, and it should be typed as one.

\section{Displayed Text}

Text is displayed by indenting it from the left
margin.  Quotations are commonly displayed.  There
are short quotations
\begin{quote}
   This is a short quotation.  It consists of a
   single paragraph of text.  See how it is formatted.
\end{quote}
and longer ones.
\begin{quotation}
   This is a longer quotation.  It consists of two
   paragraphs of text, neither of which are
   particularly interesting.

   This is the second paragraph of the quotation.  It
   is just as dull as the first paragraph.
\end{quotation}
Another frequently-displayed structure is a list.
The following is an example of an \emph{itemized}
list.
\begin{itemize}
   \item This is the first item of an itemized list.
         Each item in the list is marked with a ``tick''.
         You don't have to worry about what kind of tick
         mark is used.

   \item This is the second item of the list.  It
         contains another list nested inside it.  The inner
         list is an \emph{enumerated} list.
         \begin{enumerate}
            \item This is the first item of an enumerated
                  list that is nested within the itemized list.

            \item This is the second item of the inner list.
                  \LaTeX\ allows you to nest lists deeper than
                  you really should.
         \end{enumerate}
         This is the rest of the second item of the outer
         list.  It is no more interesting than any other
         part of the item.
   \item This is the third item of the list.
\end{itemize}
You can even display poetry.
\begin{verse}
   There is an environment
    for verse \\             % The \\ command separates lines
   Whose features some poets % within a stanza.
   will curse.

                             % One or more blank lines separate stanzas.

   For instead of making\\
   Them do \emph{all} line breaking, \\
   It allows them to put too many words on a line when they'd rather be
   forced to be terse.
\end{verse}

Mathematical formulas may also be displayed.  A
displayed formula
is
one-line long; multiline
formulas require special formatting instructions.
   \[  \ip{\Gamma}{\psi'} = x'' + y^{2} + z_{i}^{n}\]
Don't start a paragraph with a displayed equation,
nor make one a paragraph by itself.

\end{document}               % End of document.

\section{Bench Huge Section 935}
% synthetic bench section 935
% This is a sample LaTeX input file.  (Version of 12 August 2004.)
%
% A '%' character causes TeX to ignore all remaining text on the line,
% and is used for comments like this one.

\documentclass{article}      % Specifies the document class

                             % The preamble begins here.
\title{An Example Document}  % Declares the document's title.
\author{Leslie Lamport}      % Declares the author's name.
\date{January 21, 1994}      % Deleting this command produces today's date.

\newcommand{\ip}[2]{(#1, #2)}
                             % Defines \ip{arg1}{arg2} to mean
                             % (arg1, arg2).

%\newcommand{\ip}[2]{\langle #1 | #2\rangle}
                             % This is an alternative definition of
                             % \ip that is commented out.

\begin{document}             % End of preamble and beginning of text.

\maketitle                   % Produces the title.

This is an example input file.  Comparing it with
the output it generates can show you how to
produce a simple document of your own.

\section{Ordinary Text}      % Produces section heading.  Lower-level
                             % sections are begun with similar
                             % \subsection and \subsubsection commands.

The ends  of words and sentences are marked
  by   spaces. It  doesn't matter how many
spaces    you type; one is as good as 100.  The
end of   a line counts as a space.

One   or more   blank lines denote the  end
of  a paragraph.

Since any number of consecutive spaces are treated
like a single one, the formatting of the input
file makes no difference to
      \LaTeX,                % The \LaTeX command generates the LaTeX logo.
but it makes a difference to you.  When you use
\LaTeX, making your input file as easy to read
as possible will be a great help as you write
your document and when you change it.  This sample
file shows how you can add comments to your own input
file.

Because printing is different from typewriting,
there are a number of things that you have to do
differently when preparing an input file than if
you were just typing the document directly.
Quotation marks like
       ``this''
have to be handled specially, as do quotes within
quotes:
       ``\,`this'            % \, separates the double and single quote.
        is what I just
        wrote, not  `that'\,''.

Dashes come in three sizes: an
       intra-word
dash, a medium dash for number ranges like
       1--2,
and a punctuation
       dash---like
this.

A sentence-ending space should be larger than the
space between words within a sentence.  You
sometimes have to type special commands in
conjunction with punctuation characters to get
this right, as in the following sentence.
       Gnats, gnus, etc.\ all  % `\ ' makes an inter-word space.
       begin with G\@.         % \@ marks end-of-sentence punctuation.
You should check the spaces after periods when
reading your output to make sure you haven't
forgotten any special cases.  Generating an
ellipsis
       \ldots\               % `\ ' is needed after `\ldots' because TeX
                             % ignores spaces after command names like \ldots
                             % made from \ + letters.
                             %
                             % Note how a `%' character causes TeX to ignore
                             % the end of the input line, so these blank lines
                             % do not start a new paragraph.
                             %
with the right spacing around the periods requires
a special command.

\LaTeX\ interprets some common characters as
commands, so you must type special commands to
generate them.  These characters include the
following:
       \$ \& \% \# \{ and \}.

In printing, text is usually emphasized with an
       \emph{italic}
type style.

\begin{em}
   A long segment of text can also be emphasized
   in this way.  Text within such a segment can be
   given \emph{additional} emphasis.
\end{em}

It is sometimes necessary to prevent \LaTeX\ from
breaking a line where it might otherwise do so.
This may be at a space, as between the ``Mr.''\ and
``Jones'' in
       ``Mr.~Jones'',        % ~ produces an unbreakable interword space.
or within a word---especially when the word is a
symbol like
       \mbox{\emph{itemnum}}
that makes little sense when hyphenated across
lines.

Footnotes\footnote{This is an example of a footnote.}
pose no problem.

\LaTeX\ is good at typesetting mathematical formulas
like
       \( x-3y + z = 7 \)
or
       \( a_{1} > x^{2n} + y^{2n} > x' \)
or
       \( \ip{A}{B} = \sum_{i} a_{i} b_{i} \).
The spaces you type in a formula are
ignored.  Remember that a letter like
       $x$                   % $ ... $  and  \( ... \)  are equivalent
is a formula when it denotes a mathematical
symbol, and it should be typed as one.

\section{Displayed Text}

Text is displayed by indenting it from the left
margin.  Quotations are commonly displayed.  There
are short quotations
\begin{quote}
   This is a short quotation.  It consists of a
   single paragraph of text.  See how it is formatted.
\end{quote}
and longer ones.
\begin{quotation}
   This is a longer quotation.  It consists of two
   paragraphs of text, neither of which are
   particularly interesting.

   This is the second paragraph of the quotation.  It
   is just as dull as the first paragraph.
\end{quotation}
Another frequently-displayed structure is a list.
The following is an example of an \emph{itemized}
list.
\begin{itemize}
   \item This is the first item of an itemized list.
         Each item in the list is marked with a ``tick''.
         You don't have to worry about what kind of tick
         mark is used.

   \item This is the second item of the list.  It
         contains another list nested inside it.  The inner
         list is an \emph{enumerated} list.
         \begin{enumerate}
            \item This is the first item of an enumerated
                  list that is nested within the itemized list.

            \item This is the second item of the inner list.
                  \LaTeX\ allows you to nest lists deeper than
                  you really should.
         \end{enumerate}
         This is the rest of the second item of the outer
         list.  It is no more interesting than any other
         part of the item.
   \item This is the third item of the list.
\end{itemize}
You can even display poetry.
\begin{verse}
   There is an environment
    for verse \\             % The \\ command separates lines
   Whose features some poets % within a stanza.
   will curse.

                             % One or more blank lines separate stanzas.

   For instead of making\\
   Them do \emph{all} line breaking, \\
   It allows them to put too many words on a line when they'd rather be
   forced to be terse.
\end{verse}

Mathematical formulas may also be displayed.  A
displayed formula
is
one-line long; multiline
formulas require special formatting instructions.
   \[  \ip{\Gamma}{\psi'} = x'' + y^{2} + z_{i}^{n}\]
Don't start a paragraph with a displayed equation,
nor make one a paragraph by itself.

\end{document}               % End of document.

\section{Bench Huge Section 936}
% synthetic bench section 936
% This is a sample LaTeX input file.  (Version of 12 August 2004.)
%
% A '%' character causes TeX to ignore all remaining text on the line,
% and is used for comments like this one.

\documentclass{article}      % Specifies the document class

                             % The preamble begins here.
\title{An Example Document}  % Declares the document's title.
\author{Leslie Lamport}      % Declares the author's name.
\date{January 21, 1994}      % Deleting this command produces today's date.

\newcommand{\ip}[2]{(#1, #2)}
                             % Defines \ip{arg1}{arg2} to mean
                             % (arg1, arg2).

%\newcommand{\ip}[2]{\langle #1 | #2\rangle}
                             % This is an alternative definition of
                             % \ip that is commented out.

\begin{document}             % End of preamble and beginning of text.

\maketitle                   % Produces the title.

This is an example input file.  Comparing it with
the output it generates can show you how to
produce a simple document of your own.

\section{Ordinary Text}      % Produces section heading.  Lower-level
                             % sections are begun with similar
                             % \subsection and \subsubsection commands.

The ends  of words and sentences are marked
  by   spaces. It  doesn't matter how many
spaces    you type; one is as good as 100.  The
end of   a line counts as a space.

One   or more   blank lines denote the  end
of  a paragraph.

Since any number of consecutive spaces are treated
like a single one, the formatting of the input
file makes no difference to
      \LaTeX,                % The \LaTeX command generates the LaTeX logo.
but it makes a difference to you.  When you use
\LaTeX, making your input file as easy to read
as possible will be a great help as you write
your document and when you change it.  This sample
file shows how you can add comments to your own input
file.

Because printing is different from typewriting,
there are a number of things that you have to do
differently when preparing an input file than if
you were just typing the document directly.
Quotation marks like
       ``this''
have to be handled specially, as do quotes within
quotes:
       ``\,`this'            % \, separates the double and single quote.
        is what I just
        wrote, not  `that'\,''.

Dashes come in three sizes: an
       intra-word
dash, a medium dash for number ranges like
       1--2,
and a punctuation
       dash---like
this.

A sentence-ending space should be larger than the
space between words within a sentence.  You
sometimes have to type special commands in
conjunction with punctuation characters to get
this right, as in the following sentence.
       Gnats, gnus, etc.\ all  % `\ ' makes an inter-word space.
       begin with G\@.         % \@ marks end-of-sentence punctuation.
You should check the spaces after periods when
reading your output to make sure you haven't
forgotten any special cases.  Generating an
ellipsis
       \ldots\               % `\ ' is needed after `\ldots' because TeX
                             % ignores spaces after command names like \ldots
                             % made from \ + letters.
                             %
                             % Note how a `%' character causes TeX to ignore
                             % the end of the input line, so these blank lines
                             % do not start a new paragraph.
                             %
with the right spacing around the periods requires
a special command.

\LaTeX\ interprets some common characters as
commands, so you must type special commands to
generate them.  These characters include the
following:
       \$ \& \% \# \{ and \}.

In printing, text is usually emphasized with an
       \emph{italic}
type style.

\begin{em}
   A long segment of text can also be emphasized
   in this way.  Text within such a segment can be
   given \emph{additional} emphasis.
\end{em}

It is sometimes necessary to prevent \LaTeX\ from
breaking a line where it might otherwise do so.
This may be at a space, as between the ``Mr.''\ and
``Jones'' in
       ``Mr.~Jones'',        % ~ produces an unbreakable interword space.
or within a word---especially when the word is a
symbol like
       \mbox{\emph{itemnum}}
that makes little sense when hyphenated across
lines.

Footnotes\footnote{This is an example of a footnote.}
pose no problem.

\LaTeX\ is good at typesetting mathematical formulas
like
       \( x-3y + z = 7 \)
or
       \( a_{1} > x^{2n} + y^{2n} > x' \)
or
       \( \ip{A}{B} = \sum_{i} a_{i} b_{i} \).
The spaces you type in a formula are
ignored.  Remember that a letter like
       $x$                   % $ ... $  and  \( ... \)  are equivalent
is a formula when it denotes a mathematical
symbol, and it should be typed as one.

\section{Displayed Text}

Text is displayed by indenting it from the left
margin.  Quotations are commonly displayed.  There
are short quotations
\begin{quote}
   This is a short quotation.  It consists of a
   single paragraph of text.  See how it is formatted.
\end{quote}
and longer ones.
\begin{quotation}
   This is a longer quotation.  It consists of two
   paragraphs of text, neither of which are
   particularly interesting.

   This is the second paragraph of the quotation.  It
   is just as dull as the first paragraph.
\end{quotation}
Another frequently-displayed structure is a list.
The following is an example of an \emph{itemized}
list.
\begin{itemize}
   \item This is the first item of an itemized list.
         Each item in the list is marked with a ``tick''.
         You don't have to worry about what kind of tick
         mark is used.

   \item This is the second item of the list.  It
         contains another list nested inside it.  The inner
         list is an \emph{enumerated} list.
         \begin{enumerate}
            \item This is the first item of an enumerated
                  list that is nested within the itemized list.

            \item This is the second item of the inner list.
                  \LaTeX\ allows you to nest lists deeper than
                  you really should.
         \end{enumerate}
         This is the rest of the second item of the outer
         list.  It is no more interesting than any other
         part of the item.
   \item This is the third item of the list.
\end{itemize}
You can even display poetry.
\begin{verse}
   There is an environment
    for verse \\             % The \\ command separates lines
   Whose features some poets % within a stanza.
   will curse.

                             % One or more blank lines separate stanzas.

   For instead of making\\
   Them do \emph{all} line breaking, \\
   It allows them to put too many words on a line when they'd rather be
   forced to be terse.
\end{verse}

Mathematical formulas may also be displayed.  A
displayed formula
is
one-line long; multiline
formulas require special formatting instructions.
   \[  \ip{\Gamma}{\psi'} = x'' + y^{2} + z_{i}^{n}\]
Don't start a paragraph with a displayed equation,
nor make one a paragraph by itself.

\end{document}               % End of document.

\section{Bench Huge Section 937}
% synthetic bench section 937
% This is a sample LaTeX input file.  (Version of 12 August 2004.)
%
% A '%' character causes TeX to ignore all remaining text on the line,
% and is used for comments like this one.

\documentclass{article}      % Specifies the document class

                             % The preamble begins here.
\title{An Example Document}  % Declares the document's title.
\author{Leslie Lamport}      % Declares the author's name.
\date{January 21, 1994}      % Deleting this command produces today's date.

\newcommand{\ip}[2]{(#1, #2)}
                             % Defines \ip{arg1}{arg2} to mean
                             % (arg1, arg2).

%\newcommand{\ip}[2]{\langle #1 | #2\rangle}
                             % This is an alternative definition of
                             % \ip that is commented out.

\begin{document}             % End of preamble and beginning of text.

\maketitle                   % Produces the title.

This is an example input file.  Comparing it with
the output it generates can show you how to
produce a simple document of your own.

\section{Ordinary Text}      % Produces section heading.  Lower-level
                             % sections are begun with similar
                             % \subsection and \subsubsection commands.

The ends  of words and sentences are marked
  by   spaces. It  doesn't matter how many
spaces    you type; one is as good as 100.  The
end of   a line counts as a space.

One   or more   blank lines denote the  end
of  a paragraph.

Since any number of consecutive spaces are treated
like a single one, the formatting of the input
file makes no difference to
      \LaTeX,                % The \LaTeX command generates the LaTeX logo.
but it makes a difference to you.  When you use
\LaTeX, making your input file as easy to read
as possible will be a great help as you write
your document and when you change it.  This sample
file shows how you can add comments to your own input
file.

Because printing is different from typewriting,
there are a number of things that you have to do
differently when preparing an input file than if
you were just typing the document directly.
Quotation marks like
       ``this''
have to be handled specially, as do quotes within
quotes:
       ``\,`this'            % \, separates the double and single quote.
        is what I just
        wrote, not  `that'\,''.

Dashes come in three sizes: an
       intra-word
dash, a medium dash for number ranges like
       1--2,
and a punctuation
       dash---like
this.

A sentence-ending space should be larger than the
space between words within a sentence.  You
sometimes have to type special commands in
conjunction with punctuation characters to get
this right, as in the following sentence.
       Gnats, gnus, etc.\ all  % `\ ' makes an inter-word space.
       begin with G\@.         % \@ marks end-of-sentence punctuation.
You should check the spaces after periods when
reading your output to make sure you haven't
forgotten any special cases.  Generating an
ellipsis
       \ldots\               % `\ ' is needed after `\ldots' because TeX
                             % ignores spaces after command names like \ldots
                             % made from \ + letters.
                             %
                             % Note how a `%' character causes TeX to ignore
                             % the end of the input line, so these blank lines
                             % do not start a new paragraph.
                             %
with the right spacing around the periods requires
a special command.

\LaTeX\ interprets some common characters as
commands, so you must type special commands to
generate them.  These characters include the
following:
       \$ \& \% \# \{ and \}.

In printing, text is usually emphasized with an
       \emph{italic}
type style.

\begin{em}
   A long segment of text can also be emphasized
   in this way.  Text within such a segment can be
   given \emph{additional} emphasis.
\end{em}

It is sometimes necessary to prevent \LaTeX\ from
breaking a line where it might otherwise do so.
This may be at a space, as between the ``Mr.''\ and
``Jones'' in
       ``Mr.~Jones'',        % ~ produces an unbreakable interword space.
or within a word---especially when the word is a
symbol like
       \mbox{\emph{itemnum}}
that makes little sense when hyphenated across
lines.

Footnotes\footnote{This is an example of a footnote.}
pose no problem.

\LaTeX\ is good at typesetting mathematical formulas
like
       \( x-3y + z = 7 \)
or
       \( a_{1} > x^{2n} + y^{2n} > x' \)
or
       \( \ip{A}{B} = \sum_{i} a_{i} b_{i} \).
The spaces you type in a formula are
ignored.  Remember that a letter like
       $x$                   % $ ... $  and  \( ... \)  are equivalent
is a formula when it denotes a mathematical
symbol, and it should be typed as one.

\section{Displayed Text}

Text is displayed by indenting it from the left
margin.  Quotations are commonly displayed.  There
are short quotations
\begin{quote}
   This is a short quotation.  It consists of a
   single paragraph of text.  See how it is formatted.
\end{quote}
and longer ones.
\begin{quotation}
   This is a longer quotation.  It consists of two
   paragraphs of text, neither of which are
   particularly interesting.

   This is the second paragraph of the quotation.  It
   is just as dull as the first paragraph.
\end{quotation}
Another frequently-displayed structure is a list.
The following is an example of an \emph{itemized}
list.
\begin{itemize}
   \item This is the first item of an itemized list.
         Each item in the list is marked with a ``tick''.
         You don't have to worry about what kind of tick
         mark is used.

   \item This is the second item of the list.  It
         contains another list nested inside it.  The inner
         list is an \emph{enumerated} list.
         \begin{enumerate}
            \item This is the first item of an enumerated
                  list that is nested within the itemized list.

            \item This is the second item of the inner list.
                  \LaTeX\ allows you to nest lists deeper than
                  you really should.
         \end{enumerate}
         This is the rest of the second item of the outer
         list.  It is no more interesting than any other
         part of the item.
   \item This is the third item of the list.
\end{itemize}
You can even display poetry.
\begin{verse}
   There is an environment
    for verse \\             % The \\ command separates lines
   Whose features some poets % within a stanza.
   will curse.

                             % One or more blank lines separate stanzas.

   For instead of making\\
   Them do \emph{all} line breaking, \\
   It allows them to put too many words on a line when they'd rather be
   forced to be terse.
\end{verse}

Mathematical formulas may also be displayed.  A
displayed formula
is
one-line long; multiline
formulas require special formatting instructions.
   \[  \ip{\Gamma}{\psi'} = x'' + y^{2} + z_{i}^{n}\]
Don't start a paragraph with a displayed equation,
nor make one a paragraph by itself.

\end{document}               % End of document.

\section{Bench Huge Section 938}
% synthetic bench section 938
% This is a sample LaTeX input file.  (Version of 12 August 2004.)
%
% A '%' character causes TeX to ignore all remaining text on the line,
% and is used for comments like this one.

\documentclass{article}      % Specifies the document class

                             % The preamble begins here.
\title{An Example Document}  % Declares the document's title.
\author{Leslie Lamport}      % Declares the author's name.
\date{January 21, 1994}      % Deleting this command produces today's date.

\newcommand{\ip}[2]{(#1, #2)}
                             % Defines \ip{arg1}{arg2} to mean
                             % (arg1, arg2).

%\newcommand{\ip}[2]{\langle #1 | #2\rangle}
                             % This is an alternative definition of
                             % \ip that is commented out.

\begin{document}             % End of preamble and beginning of text.

\maketitle                   % Produces the title.

This is an example input file.  Comparing it with
the output it generates can show you how to
produce a simple document of your own.

\section{Ordinary Text}      % Produces section heading.  Lower-level
                             % sections are begun with similar
                             % \subsection and \subsubsection commands.

The ends  of words and sentences are marked
  by   spaces. It  doesn't matter how many
spaces    you type; one is as good as 100.  The
end of   a line counts as a space.

One   or more   blank lines denote the  end
of  a paragraph.

Since any number of consecutive spaces are treated
like a single one, the formatting of the input
file makes no difference to
      \LaTeX,                % The \LaTeX command generates the LaTeX logo.
but it makes a difference to you.  When you use
\LaTeX, making your input file as easy to read
as possible will be a great help as you write
your document and when you change it.  This sample
file shows how you can add comments to your own input
file.

Because printing is different from typewriting,
there are a number of things that you have to do
differently when preparing an input file than if
you were just typing the document directly.
Quotation marks like
       ``this''
have to be handled specially, as do quotes within
quotes:
       ``\,`this'            % \, separates the double and single quote.
        is what I just
        wrote, not  `that'\,''.

Dashes come in three sizes: an
       intra-word
dash, a medium dash for number ranges like
       1--2,
and a punctuation
       dash---like
this.

A sentence-ending space should be larger than the
space between words within a sentence.  You
sometimes have to type special commands in
conjunction with punctuation characters to get
this right, as in the following sentence.
       Gnats, gnus, etc.\ all  % `\ ' makes an inter-word space.
       begin with G\@.         % \@ marks end-of-sentence punctuation.
You should check the spaces after periods when
reading your output to make sure you haven't
forgotten any special cases.  Generating an
ellipsis
       \ldots\               % `\ ' is needed after `\ldots' because TeX
                             % ignores spaces after command names like \ldots
                             % made from \ + letters.
                             %
                             % Note how a `%' character causes TeX to ignore
                             % the end of the input line, so these blank lines
                             % do not start a new paragraph.
                             %
with the right spacing around the periods requires
a special command.

\LaTeX\ interprets some common characters as
commands, so you must type special commands to
generate them.  These characters include the
following:
       \$ \& \% \# \{ and \}.

In printing, text is usually emphasized with an
       \emph{italic}
type style.

\begin{em}
   A long segment of text can also be emphasized
   in this way.  Text within such a segment can be
   given \emph{additional} emphasis.
\end{em}

It is sometimes necessary to prevent \LaTeX\ from
breaking a line where it might otherwise do so.
This may be at a space, as between the ``Mr.''\ and
``Jones'' in
       ``Mr.~Jones'',        % ~ produces an unbreakable interword space.
or within a word---especially when the word is a
symbol like
       \mbox{\emph{itemnum}}
that makes little sense when hyphenated across
lines.

Footnotes\footnote{This is an example of a footnote.}
pose no problem.

\LaTeX\ is good at typesetting mathematical formulas
like
       \( x-3y + z = 7 \)
or
       \( a_{1} > x^{2n} + y^{2n} > x' \)
or
       \( \ip{A}{B} = \sum_{i} a_{i} b_{i} \).
The spaces you type in a formula are
ignored.  Remember that a letter like
       $x$                   % $ ... $  and  \( ... \)  are equivalent
is a formula when it denotes a mathematical
symbol, and it should be typed as one.

\section{Displayed Text}

Text is displayed by indenting it from the left
margin.  Quotations are commonly displayed.  There
are short quotations
\begin{quote}
   This is a short quotation.  It consists of a
   single paragraph of text.  See how it is formatted.
\end{quote}
and longer ones.
\begin{quotation}
   This is a longer quotation.  It consists of two
   paragraphs of text, neither of which are
   particularly interesting.

   This is the second paragraph of the quotation.  It
   is just as dull as the first paragraph.
\end{quotation}
Another frequently-displayed structure is a list.
The following is an example of an \emph{itemized}
list.
\begin{itemize}
   \item This is the first item of an itemized list.
         Each item in the list is marked with a ``tick''.
         You don't have to worry about what kind of tick
         mark is used.

   \item This is the second item of the list.  It
         contains another list nested inside it.  The inner
         list is an \emph{enumerated} list.
         \begin{enumerate}
            \item This is the first item of an enumerated
                  list that is nested within the itemized list.

            \item This is the second item of the inner list.
                  \LaTeX\ allows you to nest lists deeper than
                  you really should.
         \end{enumerate}
         This is the rest of the second item of the outer
         list.  It is no more interesting than any other
         part of the item.
   \item This is the third item of the list.
\end{itemize}
You can even display poetry.
\begin{verse}
   There is an environment
    for verse \\             % The \\ command separates lines
   Whose features some poets % within a stanza.
   will curse.

                             % One or more blank lines separate stanzas.

   For instead of making\\
   Them do \emph{all} line breaking, \\
   It allows them to put too many words on a line when they'd rather be
   forced to be terse.
\end{verse}

Mathematical formulas may also be displayed.  A
displayed formula
is
one-line long; multiline
formulas require special formatting instructions.
   \[  \ip{\Gamma}{\psi'} = x'' + y^{2} + z_{i}^{n}\]
Don't start a paragraph with a displayed equation,
nor make one a paragraph by itself.

\end{document}               % End of document.

\section{Bench Huge Section 939}
% synthetic bench section 939
% This is a sample LaTeX input file.  (Version of 12 August 2004.)
%
% A '%' character causes TeX to ignore all remaining text on the line,
% and is used for comments like this one.

\documentclass{article}      % Specifies the document class

                             % The preamble begins here.
\title{An Example Document}  % Declares the document's title.
\author{Leslie Lamport}      % Declares the author's name.
\date{January 21, 1994}      % Deleting this command produces today's date.

\newcommand{\ip}[2]{(#1, #2)}
                             % Defines \ip{arg1}{arg2} to mean
                             % (arg1, arg2).

%\newcommand{\ip}[2]{\langle #1 | #2\rangle}
                             % This is an alternative definition of
                             % \ip that is commented out.

\begin{document}             % End of preamble and beginning of text.

\maketitle                   % Produces the title.

This is an example input file.  Comparing it with
the output it generates can show you how to
produce a simple document of your own.

\section{Ordinary Text}      % Produces section heading.  Lower-level
                             % sections are begun with similar
                             % \subsection and \subsubsection commands.

The ends  of words and sentences are marked
  by   spaces. It  doesn't matter how many
spaces    you type; one is as good as 100.  The
end of   a line counts as a space.

One   or more   blank lines denote the  end
of  a paragraph.

Since any number of consecutive spaces are treated
like a single one, the formatting of the input
file makes no difference to
      \LaTeX,                % The \LaTeX command generates the LaTeX logo.
but it makes a difference to you.  When you use
\LaTeX, making your input file as easy to read
as possible will be a great help as you write
your document and when you change it.  This sample
file shows how you can add comments to your own input
file.

Because printing is different from typewriting,
there are a number of things that you have to do
differently when preparing an input file than if
you were just typing the document directly.
Quotation marks like
       ``this''
have to be handled specially, as do quotes within
quotes:
       ``\,`this'            % \, separates the double and single quote.
        is what I just
        wrote, not  `that'\,''.

Dashes come in three sizes: an
       intra-word
dash, a medium dash for number ranges like
       1--2,
and a punctuation
       dash---like
this.

A sentence-ending space should be larger than the
space between words within a sentence.  You
sometimes have to type special commands in
conjunction with punctuation characters to get
this right, as in the following sentence.
       Gnats, gnus, etc.\ all  % `\ ' makes an inter-word space.
       begin with G\@.         % \@ marks end-of-sentence punctuation.
You should check the spaces after periods when
reading your output to make sure you haven't
forgotten any special cases.  Generating an
ellipsis
       \ldots\               % `\ ' is needed after `\ldots' because TeX
                             % ignores spaces after command names like \ldots
                             % made from \ + letters.
                             %
                             % Note how a `%' character causes TeX to ignore
                             % the end of the input line, so these blank lines
                             % do not start a new paragraph.
                             %
with the right spacing around the periods requires
a special command.

\LaTeX\ interprets some common characters as
commands, so you must type special commands to
generate them.  These characters include the
following:
       \$ \& \% \# \{ and \}.

In printing, text is usually emphasized with an
       \emph{italic}
type style.

\begin{em}
   A long segment of text can also be emphasized
   in this way.  Text within such a segment can be
   given \emph{additional} emphasis.
\end{em}

It is sometimes necessary to prevent \LaTeX\ from
breaking a line where it might otherwise do so.
This may be at a space, as between the ``Mr.''\ and
``Jones'' in
       ``Mr.~Jones'',        % ~ produces an unbreakable interword space.
or within a word---especially when the word is a
symbol like
       \mbox{\emph{itemnum}}
that makes little sense when hyphenated across
lines.

Footnotes\footnote{This is an example of a footnote.}
pose no problem.

\LaTeX\ is good at typesetting mathematical formulas
like
       \( x-3y + z = 7 \)
or
       \( a_{1} > x^{2n} + y^{2n} > x' \)
or
       \( \ip{A}{B} = \sum_{i} a_{i} b_{i} \).
The spaces you type in a formula are
ignored.  Remember that a letter like
       $x$                   % $ ... $  and  \( ... \)  are equivalent
is a formula when it denotes a mathematical
symbol, and it should be typed as one.

\section{Displayed Text}

Text is displayed by indenting it from the left
margin.  Quotations are commonly displayed.  There
are short quotations
\begin{quote}
   This is a short quotation.  It consists of a
   single paragraph of text.  See how it is formatted.
\end{quote}
and longer ones.
\begin{quotation}
   This is a longer quotation.  It consists of two
   paragraphs of text, neither of which are
   particularly interesting.

   This is the second paragraph of the quotation.  It
   is just as dull as the first paragraph.
\end{quotation}
Another frequently-displayed structure is a list.
The following is an example of an \emph{itemized}
list.
\begin{itemize}
   \item This is the first item of an itemized list.
         Each item in the list is marked with a ``tick''.
         You don't have to worry about what kind of tick
         mark is used.

   \item This is the second item of the list.  It
         contains another list nested inside it.  The inner
         list is an \emph{enumerated} list.
         \begin{enumerate}
            \item This is the first item of an enumerated
                  list that is nested within the itemized list.

            \item This is the second item of the inner list.
                  \LaTeX\ allows you to nest lists deeper than
                  you really should.
         \end{enumerate}
         This is the rest of the second item of the outer
         list.  It is no more interesting than any other
         part of the item.
   \item This is the third item of the list.
\end{itemize}
You can even display poetry.
\begin{verse}
   There is an environment
    for verse \\             % The \\ command separates lines
   Whose features some poets % within a stanza.
   will curse.

                             % One or more blank lines separate stanzas.

   For instead of making\\
   Them do \emph{all} line breaking, \\
   It allows them to put too many words on a line when they'd rather be
   forced to be terse.
\end{verse}

Mathematical formulas may also be displayed.  A
displayed formula
is
one-line long; multiline
formulas require special formatting instructions.
   \[  \ip{\Gamma}{\psi'} = x'' + y^{2} + z_{i}^{n}\]
Don't start a paragraph with a displayed equation,
nor make one a paragraph by itself.

\end{document}               % End of document.

\section{Bench Huge Section 940}
% synthetic bench section 940
% This is a sample LaTeX input file.  (Version of 12 August 2004.)
%
% A '%' character causes TeX to ignore all remaining text on the line,
% and is used for comments like this one.

\documentclass{article}      % Specifies the document class

                             % The preamble begins here.
\title{An Example Document}  % Declares the document's title.
\author{Leslie Lamport}      % Declares the author's name.
\date{January 21, 1994}      % Deleting this command produces today's date.

\newcommand{\ip}[2]{(#1, #2)}
                             % Defines \ip{arg1}{arg2} to mean
                             % (arg1, arg2).

%\newcommand{\ip}[2]{\langle #1 | #2\rangle}
                             % This is an alternative definition of
                             % \ip that is commented out.

\begin{document}             % End of preamble and beginning of text.

\maketitle                   % Produces the title.

This is an example input file.  Comparing it with
the output it generates can show you how to
produce a simple document of your own.

\section{Ordinary Text}      % Produces section heading.  Lower-level
                             % sections are begun with similar
                             % \subsection and \subsubsection commands.

The ends  of words and sentences are marked
  by   spaces. It  doesn't matter how many
spaces    you type; one is as good as 100.  The
end of   a line counts as a space.

One   or more   blank lines denote the  end
of  a paragraph.

Since any number of consecutive spaces are treated
like a single one, the formatting of the input
file makes no difference to
      \LaTeX,                % The \LaTeX command generates the LaTeX logo.
but it makes a difference to you.  When you use
\LaTeX, making your input file as easy to read
as possible will be a great help as you write
your document and when you change it.  This sample
file shows how you can add comments to your own input
file.

Because printing is different from typewriting,
there are a number of things that you have to do
differently when preparing an input file than if
you were just typing the document directly.
Quotation marks like
       ``this''
have to be handled specially, as do quotes within
quotes:
       ``\,`this'            % \, separates the double and single quote.
        is what I just
        wrote, not  `that'\,''.

Dashes come in three sizes: an
       intra-word
dash, a medium dash for number ranges like
       1--2,
and a punctuation
       dash---like
this.

A sentence-ending space should be larger than the
space between words within a sentence.  You
sometimes have to type special commands in
conjunction with punctuation characters to get
this right, as in the following sentence.
       Gnats, gnus, etc.\ all  % `\ ' makes an inter-word space.
       begin with G\@.         % \@ marks end-of-sentence punctuation.
You should check the spaces after periods when
reading your output to make sure you haven't
forgotten any special cases.  Generating an
ellipsis
       \ldots\               % `\ ' is needed after `\ldots' because TeX
                             % ignores spaces after command names like \ldots
                             % made from \ + letters.
                             %
                             % Note how a `%' character causes TeX to ignore
                             % the end of the input line, so these blank lines
                             % do not start a new paragraph.
                             %
with the right spacing around the periods requires
a special command.

\LaTeX\ interprets some common characters as
commands, so you must type special commands to
generate them.  These characters include the
following:
       \$ \& \% \# \{ and \}.

In printing, text is usually emphasized with an
       \emph{italic}
type style.

\begin{em}
   A long segment of text can also be emphasized
   in this way.  Text within such a segment can be
   given \emph{additional} emphasis.
\end{em}

It is sometimes necessary to prevent \LaTeX\ from
breaking a line where it might otherwise do so.
This may be at a space, as between the ``Mr.''\ and
``Jones'' in
       ``Mr.~Jones'',        % ~ produces an unbreakable interword space.
or within a word---especially when the word is a
symbol like
       \mbox{\emph{itemnum}}
that makes little sense when hyphenated across
lines.

Footnotes\footnote{This is an example of a footnote.}
pose no problem.

\LaTeX\ is good at typesetting mathematical formulas
like
       \( x-3y + z = 7 \)
or
       \( a_{1} > x^{2n} + y^{2n} > x' \)
or
       \( \ip{A}{B} = \sum_{i} a_{i} b_{i} \).
The spaces you type in a formula are
ignored.  Remember that a letter like
       $x$                   % $ ... $  and  \( ... \)  are equivalent
is a formula when it denotes a mathematical
symbol, and it should be typed as one.

\section{Displayed Text}

Text is displayed by indenting it from the left
margin.  Quotations are commonly displayed.  There
are short quotations
\begin{quote}
   This is a short quotation.  It consists of a
   single paragraph of text.  See how it is formatted.
\end{quote}
and longer ones.
\begin{quotation}
   This is a longer quotation.  It consists of two
   paragraphs of text, neither of which are
   particularly interesting.

   This is the second paragraph of the quotation.  It
   is just as dull as the first paragraph.
\end{quotation}
Another frequently-displayed structure is a list.
The following is an example of an \emph{itemized}
list.
\begin{itemize}
   \item This is the first item of an itemized list.
         Each item in the list is marked with a ``tick''.
         You don't have to worry about what kind of tick
         mark is used.

   \item This is the second item of the list.  It
         contains another list nested inside it.  The inner
         list is an \emph{enumerated} list.
         \begin{enumerate}
            \item This is the first item of an enumerated
                  list that is nested within the itemized list.

            \item This is the second item of the inner list.
                  \LaTeX\ allows you to nest lists deeper than
                  you really should.
         \end{enumerate}
         This is the rest of the second item of the outer
         list.  It is no more interesting than any other
         part of the item.
   \item This is the third item of the list.
\end{itemize}
You can even display poetry.
\begin{verse}
   There is an environment
    for verse \\             % The \\ command separates lines
   Whose features some poets % within a stanza.
   will curse.

                             % One or more blank lines separate stanzas.

   For instead of making\\
   Them do \emph{all} line breaking, \\
   It allows them to put too many words on a line when they'd rather be
   forced to be terse.
\end{verse}

Mathematical formulas may also be displayed.  A
displayed formula
is
one-line long; multiline
formulas require special formatting instructions.
   \[  \ip{\Gamma}{\psi'} = x'' + y^{2} + z_{i}^{n}\]
Don't start a paragraph with a displayed equation,
nor make one a paragraph by itself.

\end{document}               % End of document.

\section{Bench Huge Section 941}
% synthetic bench section 941
% This is a sample LaTeX input file.  (Version of 12 August 2004.)
%
% A '%' character causes TeX to ignore all remaining text on the line,
% and is used for comments like this one.

\documentclass{article}      % Specifies the document class

                             % The preamble begins here.
\title{An Example Document}  % Declares the document's title.
\author{Leslie Lamport}      % Declares the author's name.
\date{January 21, 1994}      % Deleting this command produces today's date.

\newcommand{\ip}[2]{(#1, #2)}
                             % Defines \ip{arg1}{arg2} to mean
                             % (arg1, arg2).

%\newcommand{\ip}[2]{\langle #1 | #2\rangle}
                             % This is an alternative definition of
                             % \ip that is commented out.

\begin{document}             % End of preamble and beginning of text.

\maketitle                   % Produces the title.

This is an example input file.  Comparing it with
the output it generates can show you how to
produce a simple document of your own.

\section{Ordinary Text}      % Produces section heading.  Lower-level
                             % sections are begun with similar
                             % \subsection and \subsubsection commands.

The ends  of words and sentences are marked
  by   spaces. It  doesn't matter how many
spaces    you type; one is as good as 100.  The
end of   a line counts as a space.

One   or more   blank lines denote the  end
of  a paragraph.

Since any number of consecutive spaces are treated
like a single one, the formatting of the input
file makes no difference to
      \LaTeX,                % The \LaTeX command generates the LaTeX logo.
but it makes a difference to you.  When you use
\LaTeX, making your input file as easy to read
as possible will be a great help as you write
your document and when you change it.  This sample
file shows how you can add comments to your own input
file.

Because printing is different from typewriting,
there are a number of things that you have to do
differently when preparing an input file than if
you were just typing the document directly.
Quotation marks like
       ``this''
have to be handled specially, as do quotes within
quotes:
       ``\,`this'            % \, separates the double and single quote.
        is what I just
        wrote, not  `that'\,''.

Dashes come in three sizes: an
       intra-word
dash, a medium dash for number ranges like
       1--2,
and a punctuation
       dash---like
this.

A sentence-ending space should be larger than the
space between words within a sentence.  You
sometimes have to type special commands in
conjunction with punctuation characters to get
this right, as in the following sentence.
       Gnats, gnus, etc.\ all  % `\ ' makes an inter-word space.
       begin with G\@.         % \@ marks end-of-sentence punctuation.
You should check the spaces after periods when
reading your output to make sure you haven't
forgotten any special cases.  Generating an
ellipsis
       \ldots\               % `\ ' is needed after `\ldots' because TeX
                             % ignores spaces after command names like \ldots
                             % made from \ + letters.
                             %
                             % Note how a `%' character causes TeX to ignore
                             % the end of the input line, so these blank lines
                             % do not start a new paragraph.
                             %
with the right spacing around the periods requires
a special command.

\LaTeX\ interprets some common characters as
commands, so you must type special commands to
generate them.  These characters include the
following:
       \$ \& \% \# \{ and \}.

In printing, text is usually emphasized with an
       \emph{italic}
type style.

\begin{em}
   A long segment of text can also be emphasized
   in this way.  Text within such a segment can be
   given \emph{additional} emphasis.
\end{em}

It is sometimes necessary to prevent \LaTeX\ from
breaking a line where it might otherwise do so.
This may be at a space, as between the ``Mr.''\ and
``Jones'' in
       ``Mr.~Jones'',        % ~ produces an unbreakable interword space.
or within a word---especially when the word is a
symbol like
       \mbox{\emph{itemnum}}
that makes little sense when hyphenated across
lines.

Footnotes\footnote{This is an example of a footnote.}
pose no problem.

\LaTeX\ is good at typesetting mathematical formulas
like
       \( x-3y + z = 7 \)
or
       \( a_{1} > x^{2n} + y^{2n} > x' \)
or
       \( \ip{A}{B} = \sum_{i} a_{i} b_{i} \).
The spaces you type in a formula are
ignored.  Remember that a letter like
       $x$                   % $ ... $  and  \( ... \)  are equivalent
is a formula when it denotes a mathematical
symbol, and it should be typed as one.

\section{Displayed Text}

Text is displayed by indenting it from the left
margin.  Quotations are commonly displayed.  There
are short quotations
\begin{quote}
   This is a short quotation.  It consists of a
   single paragraph of text.  See how it is formatted.
\end{quote}
and longer ones.
\begin{quotation}
   This is a longer quotation.  It consists of two
   paragraphs of text, neither of which are
   particularly interesting.

   This is the second paragraph of the quotation.  It
   is just as dull as the first paragraph.
\end{quotation}
Another frequently-displayed structure is a list.
The following is an example of an \emph{itemized}
list.
\begin{itemize}
   \item This is the first item of an itemized list.
         Each item in the list is marked with a ``tick''.
         You don't have to worry about what kind of tick
         mark is used.

   \item This is the second item of the list.  It
         contains another list nested inside it.  The inner
         list is an \emph{enumerated} list.
         \begin{enumerate}
            \item This is the first item of an enumerated
                  list that is nested within the itemized list.

            \item This is the second item of the inner list.
                  \LaTeX\ allows you to nest lists deeper than
                  you really should.
         \end{enumerate}
         This is the rest of the second item of the outer
         list.  It is no more interesting than any other
         part of the item.
   \item This is the third item of the list.
\end{itemize}
You can even display poetry.
\begin{verse}
   There is an environment
    for verse \\             % The \\ command separates lines
   Whose features some poets % within a stanza.
   will curse.

                             % One or more blank lines separate stanzas.

   For instead of making\\
   Them do \emph{all} line breaking, \\
   It allows them to put too many words on a line when they'd rather be
   forced to be terse.
\end{verse}

Mathematical formulas may also be displayed.  A
displayed formula
is
one-line long; multiline
formulas require special formatting instructions.
   \[  \ip{\Gamma}{\psi'} = x'' + y^{2} + z_{i}^{n}\]
Don't start a paragraph with a displayed equation,
nor make one a paragraph by itself.

\end{document}               % End of document.

\section{Bench Huge Section 942}
% synthetic bench section 942
% This is a sample LaTeX input file.  (Version of 12 August 2004.)
%
% A '%' character causes TeX to ignore all remaining text on the line,
% and is used for comments like this one.

\documentclass{article}      % Specifies the document class

                             % The preamble begins here.
\title{An Example Document}  % Declares the document's title.
\author{Leslie Lamport}      % Declares the author's name.
\date{January 21, 1994}      % Deleting this command produces today's date.

\newcommand{\ip}[2]{(#1, #2)}
                             % Defines \ip{arg1}{arg2} to mean
                             % (arg1, arg2).

%\newcommand{\ip}[2]{\langle #1 | #2\rangle}
                             % This is an alternative definition of
                             % \ip that is commented out.

\begin{document}             % End of preamble and beginning of text.

\maketitle                   % Produces the title.

This is an example input file.  Comparing it with
the output it generates can show you how to
produce a simple document of your own.

\section{Ordinary Text}      % Produces section heading.  Lower-level
                             % sections are begun with similar
                             % \subsection and \subsubsection commands.

The ends  of words and sentences are marked
  by   spaces. It  doesn't matter how many
spaces    you type; one is as good as 100.  The
end of   a line counts as a space.

One   or more   blank lines denote the  end
of  a paragraph.

Since any number of consecutive spaces are treated
like a single one, the formatting of the input
file makes no difference to
      \LaTeX,                % The \LaTeX command generates the LaTeX logo.
but it makes a difference to you.  When you use
\LaTeX, making your input file as easy to read
as possible will be a great help as you write
your document and when you change it.  This sample
file shows how you can add comments to your own input
file.

Because printing is different from typewriting,
there are a number of things that you have to do
differently when preparing an input file than if
you were just typing the document directly.
Quotation marks like
       ``this''
have to be handled specially, as do quotes within
quotes:
       ``\,`this'            % \, separates the double and single quote.
        is what I just
        wrote, not  `that'\,''.

Dashes come in three sizes: an
       intra-word
dash, a medium dash for number ranges like
       1--2,
and a punctuation
       dash---like
this.

A sentence-ending space should be larger than the
space between words within a sentence.  You
sometimes have to type special commands in
conjunction with punctuation characters to get
this right, as in the following sentence.
       Gnats, gnus, etc.\ all  % `\ ' makes an inter-word space.
       begin with G\@.         % \@ marks end-of-sentence punctuation.
You should check the spaces after periods when
reading your output to make sure you haven't
forgotten any special cases.  Generating an
ellipsis
       \ldots\               % `\ ' is needed after `\ldots' because TeX
                             % ignores spaces after command names like \ldots
                             % made from \ + letters.
                             %
                             % Note how a `%' character causes TeX to ignore
                             % the end of the input line, so these blank lines
                             % do not start a new paragraph.
                             %
with the right spacing around the periods requires
a special command.

\LaTeX\ interprets some common characters as
commands, so you must type special commands to
generate them.  These characters include the
following:
       \$ \& \% \# \{ and \}.

In printing, text is usually emphasized with an
       \emph{italic}
type style.

\begin{em}
   A long segment of text can also be emphasized
   in this way.  Text within such a segment can be
   given \emph{additional} emphasis.
\end{em}

It is sometimes necessary to prevent \LaTeX\ from
breaking a line where it might otherwise do so.
This may be at a space, as between the ``Mr.''\ and
``Jones'' in
       ``Mr.~Jones'',        % ~ produces an unbreakable interword space.
or within a word---especially when the word is a
symbol like
       \mbox{\emph{itemnum}}
that makes little sense when hyphenated across
lines.

Footnotes\footnote{This is an example of a footnote.}
pose no problem.

\LaTeX\ is good at typesetting mathematical formulas
like
       \( x-3y + z = 7 \)
or
       \( a_{1} > x^{2n} + y^{2n} > x' \)
or
       \( \ip{A}{B} = \sum_{i} a_{i} b_{i} \).
The spaces you type in a formula are
ignored.  Remember that a letter like
       $x$                   % $ ... $  and  \( ... \)  are equivalent
is a formula when it denotes a mathematical
symbol, and it should be typed as one.

\section{Displayed Text}

Text is displayed by indenting it from the left
margin.  Quotations are commonly displayed.  There
are short quotations
\begin{quote}
   This is a short quotation.  It consists of a
   single paragraph of text.  See how it is formatted.
\end{quote}
and longer ones.
\begin{quotation}
   This is a longer quotation.  It consists of two
   paragraphs of text, neither of which are
   particularly interesting.

   This is the second paragraph of the quotation.  It
   is just as dull as the first paragraph.
\end{quotation}
Another frequently-displayed structure is a list.
The following is an example of an \emph{itemized}
list.
\begin{itemize}
   \item This is the first item of an itemized list.
         Each item in the list is marked with a ``tick''.
         You don't have to worry about what kind of tick
         mark is used.

   \item This is the second item of the list.  It
         contains another list nested inside it.  The inner
         list is an \emph{enumerated} list.
         \begin{enumerate}
            \item This is the first item of an enumerated
                  list that is nested within the itemized list.

            \item This is the second item of the inner list.
                  \LaTeX\ allows you to nest lists deeper than
                  you really should.
         \end{enumerate}
         This is the rest of the second item of the outer
         list.  It is no more interesting than any other
         part of the item.
   \item This is the third item of the list.
\end{itemize}
You can even display poetry.
\begin{verse}
   There is an environment
    for verse \\             % The \\ command separates lines
   Whose features some poets % within a stanza.
   will curse.

                             % One or more blank lines separate stanzas.

   For instead of making\\
   Them do \emph{all} line breaking, \\
   It allows them to put too many words on a line when they'd rather be
   forced to be terse.
\end{verse}

Mathematical formulas may also be displayed.  A
displayed formula
is
one-line long; multiline
formulas require special formatting instructions.
   \[  \ip{\Gamma}{\psi'} = x'' + y^{2} + z_{i}^{n}\]
Don't start a paragraph with a displayed equation,
nor make one a paragraph by itself.

\end{document}               % End of document.

\section{Bench Huge Section 943}
% synthetic bench section 943
% This is a sample LaTeX input file.  (Version of 12 August 2004.)
%
% A '%' character causes TeX to ignore all remaining text on the line,
% and is used for comments like this one.

\documentclass{article}      % Specifies the document class

                             % The preamble begins here.
\title{An Example Document}  % Declares the document's title.
\author{Leslie Lamport}      % Declares the author's name.
\date{January 21, 1994}      % Deleting this command produces today's date.

\newcommand{\ip}[2]{(#1, #2)}
                             % Defines \ip{arg1}{arg2} to mean
                             % (arg1, arg2).

%\newcommand{\ip}[2]{\langle #1 | #2\rangle}
                             % This is an alternative definition of
                             % \ip that is commented out.

\begin{document}             % End of preamble and beginning of text.

\maketitle                   % Produces the title.

This is an example input file.  Comparing it with
the output it generates can show you how to
produce a simple document of your own.

\section{Ordinary Text}      % Produces section heading.  Lower-level
                             % sections are begun with similar
                             % \subsection and \subsubsection commands.

The ends  of words and sentences are marked
  by   spaces. It  doesn't matter how many
spaces    you type; one is as good as 100.  The
end of   a line counts as a space.

One   or more   blank lines denote the  end
of  a paragraph.

Since any number of consecutive spaces are treated
like a single one, the formatting of the input
file makes no difference to
      \LaTeX,                % The \LaTeX command generates the LaTeX logo.
but it makes a difference to you.  When you use
\LaTeX, making your input file as easy to read
as possible will be a great help as you write
your document and when you change it.  This sample
file shows how you can add comments to your own input
file.

Because printing is different from typewriting,
there are a number of things that you have to do
differently when preparing an input file than if
you were just typing the document directly.
Quotation marks like
       ``this''
have to be handled specially, as do quotes within
quotes:
       ``\,`this'            % \, separates the double and single quote.
        is what I just
        wrote, not  `that'\,''.

Dashes come in three sizes: an
       intra-word
dash, a medium dash for number ranges like
       1--2,
and a punctuation
       dash---like
this.

A sentence-ending space should be larger than the
space between words within a sentence.  You
sometimes have to type special commands in
conjunction with punctuation characters to get
this right, as in the following sentence.
       Gnats, gnus, etc.\ all  % `\ ' makes an inter-word space.
       begin with G\@.         % \@ marks end-of-sentence punctuation.
You should check the spaces after periods when
reading your output to make sure you haven't
forgotten any special cases.  Generating an
ellipsis
       \ldots\               % `\ ' is needed after `\ldots' because TeX
                             % ignores spaces after command names like \ldots
                             % made from \ + letters.
                             %
                             % Note how a `%' character causes TeX to ignore
                             % the end of the input line, so these blank lines
                             % do not start a new paragraph.
                             %
with the right spacing around the periods requires
a special command.

\LaTeX\ interprets some common characters as
commands, so you must type special commands to
generate them.  These characters include the
following:
       \$ \& \% \# \{ and \}.

In printing, text is usually emphasized with an
       \emph{italic}
type style.

\begin{em}
   A long segment of text can also be emphasized
   in this way.  Text within such a segment can be
   given \emph{additional} emphasis.
\end{em}

It is sometimes necessary to prevent \LaTeX\ from
breaking a line where it might otherwise do so.
This may be at a space, as between the ``Mr.''\ and
``Jones'' in
       ``Mr.~Jones'',        % ~ produces an unbreakable interword space.
or within a word---especially when the word is a
symbol like
       \mbox{\emph{itemnum}}
that makes little sense when hyphenated across
lines.

Footnotes\footnote{This is an example of a footnote.}
pose no problem.

\LaTeX\ is good at typesetting mathematical formulas
like
       \( x-3y + z = 7 \)
or
       \( a_{1} > x^{2n} + y^{2n} > x' \)
or
       \( \ip{A}{B} = \sum_{i} a_{i} b_{i} \).
The spaces you type in a formula are
ignored.  Remember that a letter like
       $x$                   % $ ... $  and  \( ... \)  are equivalent
is a formula when it denotes a mathematical
symbol, and it should be typed as one.

\section{Displayed Text}

Text is displayed by indenting it from the left
margin.  Quotations are commonly displayed.  There
are short quotations
\begin{quote}
   This is a short quotation.  It consists of a
   single paragraph of text.  See how it is formatted.
\end{quote}
and longer ones.
\begin{quotation}
   This is a longer quotation.  It consists of two
   paragraphs of text, neither of which are
   particularly interesting.

   This is the second paragraph of the quotation.  It
   is just as dull as the first paragraph.
\end{quotation}
Another frequently-displayed structure is a list.
The following is an example of an \emph{itemized}
list.
\begin{itemize}
   \item This is the first item of an itemized list.
         Each item in the list is marked with a ``tick''.
         You don't have to worry about what kind of tick
         mark is used.

   \item This is the second item of the list.  It
         contains another list nested inside it.  The inner
         list is an \emph{enumerated} list.
         \begin{enumerate}
            \item This is the first item of an enumerated
                  list that is nested within the itemized list.

            \item This is the second item of the inner list.
                  \LaTeX\ allows you to nest lists deeper than
                  you really should.
         \end{enumerate}
         This is the rest of the second item of the outer
         list.  It is no more interesting than any other
         part of the item.
   \item This is the third item of the list.
\end{itemize}
You can even display poetry.
\begin{verse}
   There is an environment
    for verse \\             % The \\ command separates lines
   Whose features some poets % within a stanza.
   will curse.

                             % One or more blank lines separate stanzas.

   For instead of making\\
   Them do \emph{all} line breaking, \\
   It allows them to put too many words on a line when they'd rather be
   forced to be terse.
\end{verse}

Mathematical formulas may also be displayed.  A
displayed formula
is
one-line long; multiline
formulas require special formatting instructions.
   \[  \ip{\Gamma}{\psi'} = x'' + y^{2} + z_{i}^{n}\]
Don't start a paragraph with a displayed equation,
nor make one a paragraph by itself.

\end{document}               % End of document.

\section{Bench Huge Section 944}
% synthetic bench section 944
% This is a sample LaTeX input file.  (Version of 12 August 2004.)
%
% A '%' character causes TeX to ignore all remaining text on the line,
% and is used for comments like this one.

\documentclass{article}      % Specifies the document class

                             % The preamble begins here.
\title{An Example Document}  % Declares the document's title.
\author{Leslie Lamport}      % Declares the author's name.
\date{January 21, 1994}      % Deleting this command produces today's date.

\newcommand{\ip}[2]{(#1, #2)}
                             % Defines \ip{arg1}{arg2} to mean
                             % (arg1, arg2).

%\newcommand{\ip}[2]{\langle #1 | #2\rangle}
                             % This is an alternative definition of
                             % \ip that is commented out.

\begin{document}             % End of preamble and beginning of text.

\maketitle                   % Produces the title.

This is an example input file.  Comparing it with
the output it generates can show you how to
produce a simple document of your own.

\section{Ordinary Text}      % Produces section heading.  Lower-level
                             % sections are begun with similar
                             % \subsection and \subsubsection commands.

The ends  of words and sentences are marked
  by   spaces. It  doesn't matter how many
spaces    you type; one is as good as 100.  The
end of   a line counts as a space.

One   or more   blank lines denote the  end
of  a paragraph.

Since any number of consecutive spaces are treated
like a single one, the formatting of the input
file makes no difference to
      \LaTeX,                % The \LaTeX command generates the LaTeX logo.
but it makes a difference to you.  When you use
\LaTeX, making your input file as easy to read
as possible will be a great help as you write
your document and when you change it.  This sample
file shows how you can add comments to your own input
file.

Because printing is different from typewriting,
there are a number of things that you have to do
differently when preparing an input file than if
you were just typing the document directly.
Quotation marks like
       ``this''
have to be handled specially, as do quotes within
quotes:
       ``\,`this'            % \, separates the double and single quote.
        is what I just
        wrote, not  `that'\,''.

Dashes come in three sizes: an
       intra-word
dash, a medium dash for number ranges like
       1--2,
and a punctuation
       dash---like
this.

A sentence-ending space should be larger than the
space between words within a sentence.  You
sometimes have to type special commands in
conjunction with punctuation characters to get
this right, as in the following sentence.
       Gnats, gnus, etc.\ all  % `\ ' makes an inter-word space.
       begin with G\@.         % \@ marks end-of-sentence punctuation.
You should check the spaces after periods when
reading your output to make sure you haven't
forgotten any special cases.  Generating an
ellipsis
       \ldots\               % `\ ' is needed after `\ldots' because TeX
                             % ignores spaces after command names like \ldots
                             % made from \ + letters.
                             %
                             % Note how a `%' character causes TeX to ignore
                             % the end of the input line, so these blank lines
                             % do not start a new paragraph.
                             %
with the right spacing around the periods requires
a special command.

\LaTeX\ interprets some common characters as
commands, so you must type special commands to
generate them.  These characters include the
following:
       \$ \& \% \# \{ and \}.

In printing, text is usually emphasized with an
       \emph{italic}
type style.

\begin{em}
   A long segment of text can also be emphasized
   in this way.  Text within such a segment can be
   given \emph{additional} emphasis.
\end{em}

It is sometimes necessary to prevent \LaTeX\ from
breaking a line where it might otherwise do so.
This may be at a space, as between the ``Mr.''\ and
``Jones'' in
       ``Mr.~Jones'',        % ~ produces an unbreakable interword space.
or within a word---especially when the word is a
symbol like
       \mbox{\emph{itemnum}}
that makes little sense when hyphenated across
lines.

Footnotes\footnote{This is an example of a footnote.}
pose no problem.

\LaTeX\ is good at typesetting mathematical formulas
like
       \( x-3y + z = 7 \)
or
       \( a_{1} > x^{2n} + y^{2n} > x' \)
or
       \( \ip{A}{B} = \sum_{i} a_{i} b_{i} \).
The spaces you type in a formula are
ignored.  Remember that a letter like
       $x$                   % $ ... $  and  \( ... \)  are equivalent
is a formula when it denotes a mathematical
symbol, and it should be typed as one.

\section{Displayed Text}

Text is displayed by indenting it from the left
margin.  Quotations are commonly displayed.  There
are short quotations
\begin{quote}
   This is a short quotation.  It consists of a
   single paragraph of text.  See how it is formatted.
\end{quote}
and longer ones.
\begin{quotation}
   This is a longer quotation.  It consists of two
   paragraphs of text, neither of which are
   particularly interesting.

   This is the second paragraph of the quotation.  It
   is just as dull as the first paragraph.
\end{quotation}
Another frequently-displayed structure is a list.
The following is an example of an \emph{itemized}
list.
\begin{itemize}
   \item This is the first item of an itemized list.
         Each item in the list is marked with a ``tick''.
         You don't have to worry about what kind of tick
         mark is used.

   \item This is the second item of the list.  It
         contains another list nested inside it.  The inner
         list is an \emph{enumerated} list.
         \begin{enumerate}
            \item This is the first item of an enumerated
                  list that is nested within the itemized list.

            \item This is the second item of the inner list.
                  \LaTeX\ allows you to nest lists deeper than
                  you really should.
         \end{enumerate}
         This is the rest of the second item of the outer
         list.  It is no more interesting than any other
         part of the item.
   \item This is the third item of the list.
\end{itemize}
You can even display poetry.
\begin{verse}
   There is an environment
    for verse \\             % The \\ command separates lines
   Whose features some poets % within a stanza.
   will curse.

                             % One or more blank lines separate stanzas.

   For instead of making\\
   Them do \emph{all} line breaking, \\
   It allows them to put too many words on a line when they'd rather be
   forced to be terse.
\end{verse}

Mathematical formulas may also be displayed.  A
displayed formula
is
one-line long; multiline
formulas require special formatting instructions.
   \[  \ip{\Gamma}{\psi'} = x'' + y^{2} + z_{i}^{n}\]
Don't start a paragraph with a displayed equation,
nor make one a paragraph by itself.

\end{document}               % End of document.

\section{Bench Huge Section 945}
% synthetic bench section 945
% This is a sample LaTeX input file.  (Version of 12 August 2004.)
%
% A '%' character causes TeX to ignore all remaining text on the line,
% and is used for comments like this one.

\documentclass{article}      % Specifies the document class

                             % The preamble begins here.
\title{An Example Document}  % Declares the document's title.
\author{Leslie Lamport}      % Declares the author's name.
\date{January 21, 1994}      % Deleting this command produces today's date.

\newcommand{\ip}[2]{(#1, #2)}
                             % Defines \ip{arg1}{arg2} to mean
                             % (arg1, arg2).

%\newcommand{\ip}[2]{\langle #1 | #2\rangle}
                             % This is an alternative definition of
                             % \ip that is commented out.

\begin{document}             % End of preamble and beginning of text.

\maketitle                   % Produces the title.

This is an example input file.  Comparing it with
the output it generates can show you how to
produce a simple document of your own.

\section{Ordinary Text}      % Produces section heading.  Lower-level
                             % sections are begun with similar
                             % \subsection and \subsubsection commands.

The ends  of words and sentences are marked
  by   spaces. It  doesn't matter how many
spaces    you type; one is as good as 100.  The
end of   a line counts as a space.

One   or more   blank lines denote the  end
of  a paragraph.

Since any number of consecutive spaces are treated
like a single one, the formatting of the input
file makes no difference to
      \LaTeX,                % The \LaTeX command generates the LaTeX logo.
but it makes a difference to you.  When you use
\LaTeX, making your input file as easy to read
as possible will be a great help as you write
your document and when you change it.  This sample
file shows how you can add comments to your own input
file.

Because printing is different from typewriting,
there are a number of things that you have to do
differently when preparing an input file than if
you were just typing the document directly.
Quotation marks like
       ``this''
have to be handled specially, as do quotes within
quotes:
       ``\,`this'            % \, separates the double and single quote.
        is what I just
        wrote, not  `that'\,''.

Dashes come in three sizes: an
       intra-word
dash, a medium dash for number ranges like
       1--2,
and a punctuation
       dash---like
this.

A sentence-ending space should be larger than the
space between words within a sentence.  You
sometimes have to type special commands in
conjunction with punctuation characters to get
this right, as in the following sentence.
       Gnats, gnus, etc.\ all  % `\ ' makes an inter-word space.
       begin with G\@.         % \@ marks end-of-sentence punctuation.
You should check the spaces after periods when
reading your output to make sure you haven't
forgotten any special cases.  Generating an
ellipsis
       \ldots\               % `\ ' is needed after `\ldots' because TeX
                             % ignores spaces after command names like \ldots
                             % made from \ + letters.
                             %
                             % Note how a `%' character causes TeX to ignore
                             % the end of the input line, so these blank lines
                             % do not start a new paragraph.
                             %
with the right spacing around the periods requires
a special command.

\LaTeX\ interprets some common characters as
commands, so you must type special commands to
generate them.  These characters include the
following:
       \$ \& \% \# \{ and \}.

In printing, text is usually emphasized with an
       \emph{italic}
type style.

\begin{em}
   A long segment of text can also be emphasized
   in this way.  Text within such a segment can be
   given \emph{additional} emphasis.
\end{em}

It is sometimes necessary to prevent \LaTeX\ from
breaking a line where it might otherwise do so.
This may be at a space, as between the ``Mr.''\ and
``Jones'' in
       ``Mr.~Jones'',        % ~ produces an unbreakable interword space.
or within a word---especially when the word is a
symbol like
       \mbox{\emph{itemnum}}
that makes little sense when hyphenated across
lines.

Footnotes\footnote{This is an example of a footnote.}
pose no problem.

\LaTeX\ is good at typesetting mathematical formulas
like
       \( x-3y + z = 7 \)
or
       \( a_{1} > x^{2n} + y^{2n} > x' \)
or
       \( \ip{A}{B} = \sum_{i} a_{i} b_{i} \).
The spaces you type in a formula are
ignored.  Remember that a letter like
       $x$                   % $ ... $  and  \( ... \)  are equivalent
is a formula when it denotes a mathematical
symbol, and it should be typed as one.

\section{Displayed Text}

Text is displayed by indenting it from the left
margin.  Quotations are commonly displayed.  There
are short quotations
\begin{quote}
   This is a short quotation.  It consists of a
   single paragraph of text.  See how it is formatted.
\end{quote}
and longer ones.
\begin{quotation}
   This is a longer quotation.  It consists of two
   paragraphs of text, neither of which are
   particularly interesting.

   This is the second paragraph of the quotation.  It
   is just as dull as the first paragraph.
\end{quotation}
Another frequently-displayed structure is a list.
The following is an example of an \emph{itemized}
list.
\begin{itemize}
   \item This is the first item of an itemized list.
         Each item in the list is marked with a ``tick''.
         You don't have to worry about what kind of tick
         mark is used.

   \item This is the second item of the list.  It
         contains another list nested inside it.  The inner
         list is an \emph{enumerated} list.
         \begin{enumerate}
            \item This is the first item of an enumerated
                  list that is nested within the itemized list.

            \item This is the second item of the inner list.
                  \LaTeX\ allows you to nest lists deeper than
                  you really should.
         \end{enumerate}
         This is the rest of the second item of the outer
         list.  It is no more interesting than any other
         part of the item.
   \item This is the third item of the list.
\end{itemize}
You can even display poetry.
\begin{verse}
   There is an environment
    for verse \\             % The \\ command separates lines
   Whose features some poets % within a stanza.
   will curse.

                             % One or more blank lines separate stanzas.

   For instead of making\\
   Them do \emph{all} line breaking, \\
   It allows them to put too many words on a line when they'd rather be
   forced to be terse.
\end{verse}

Mathematical formulas may also be displayed.  A
displayed formula
is
one-line long; multiline
formulas require special formatting instructions.
   \[  \ip{\Gamma}{\psi'} = x'' + y^{2} + z_{i}^{n}\]
Don't start a paragraph with a displayed equation,
nor make one a paragraph by itself.

\end{document}               % End of document.

\section{Bench Huge Section 946}
% synthetic bench section 946
% This is a sample LaTeX input file.  (Version of 12 August 2004.)
%
% A '%' character causes TeX to ignore all remaining text on the line,
% and is used for comments like this one.

\documentclass{article}      % Specifies the document class

                             % The preamble begins here.
\title{An Example Document}  % Declares the document's title.
\author{Leslie Lamport}      % Declares the author's name.
\date{January 21, 1994}      % Deleting this command produces today's date.

\newcommand{\ip}[2]{(#1, #2)}
                             % Defines \ip{arg1}{arg2} to mean
                             % (arg1, arg2).

%\newcommand{\ip}[2]{\langle #1 | #2\rangle}
                             % This is an alternative definition of
                             % \ip that is commented out.

\begin{document}             % End of preamble and beginning of text.

\maketitle                   % Produces the title.

This is an example input file.  Comparing it with
the output it generates can show you how to
produce a simple document of your own.

\section{Ordinary Text}      % Produces section heading.  Lower-level
                             % sections are begun with similar
                             % \subsection and \subsubsection commands.

The ends  of words and sentences are marked
  by   spaces. It  doesn't matter how many
spaces    you type; one is as good as 100.  The
end of   a line counts as a space.

One   or more   blank lines denote the  end
of  a paragraph.

Since any number of consecutive spaces are treated
like a single one, the formatting of the input
file makes no difference to
      \LaTeX,                % The \LaTeX command generates the LaTeX logo.
but it makes a difference to you.  When you use
\LaTeX, making your input file as easy to read
as possible will be a great help as you write
your document and when you change it.  This sample
file shows how you can add comments to your own input
file.

Because printing is different from typewriting,
there are a number of things that you have to do
differently when preparing an input file than if
you were just typing the document directly.
Quotation marks like
       ``this''
have to be handled specially, as do quotes within
quotes:
       ``\,`this'            % \, separates the double and single quote.
        is what I just
        wrote, not  `that'\,''.

Dashes come in three sizes: an
       intra-word
dash, a medium dash for number ranges like
       1--2,
and a punctuation
       dash---like
this.

A sentence-ending space should be larger than the
space between words within a sentence.  You
sometimes have to type special commands in
conjunction with punctuation characters to get
this right, as in the following sentence.
       Gnats, gnus, etc.\ all  % `\ ' makes an inter-word space.
       begin with G\@.         % \@ marks end-of-sentence punctuation.
You should check the spaces after periods when
reading your output to make sure you haven't
forgotten any special cases.  Generating an
ellipsis
       \ldots\               % `\ ' is needed after `\ldots' because TeX
                             % ignores spaces after command names like \ldots
                             % made from \ + letters.
                             %
                             % Note how a `%' character causes TeX to ignore
                             % the end of the input line, so these blank lines
                             % do not start a new paragraph.
                             %
with the right spacing around the periods requires
a special command.

\LaTeX\ interprets some common characters as
commands, so you must type special commands to
generate them.  These characters include the
following:
       \$ \& \% \# \{ and \}.

In printing, text is usually emphasized with an
       \emph{italic}
type style.

\begin{em}
   A long segment of text can also be emphasized
   in this way.  Text within such a segment can be
   given \emph{additional} emphasis.
\end{em}

It is sometimes necessary to prevent \LaTeX\ from
breaking a line where it might otherwise do so.
This may be at a space, as between the ``Mr.''\ and
``Jones'' in
       ``Mr.~Jones'',        % ~ produces an unbreakable interword space.
or within a word---especially when the word is a
symbol like
       \mbox{\emph{itemnum}}
that makes little sense when hyphenated across
lines.

Footnotes\footnote{This is an example of a footnote.}
pose no problem.

\LaTeX\ is good at typesetting mathematical formulas
like
       \( x-3y + z = 7 \)
or
       \( a_{1} > x^{2n} + y^{2n} > x' \)
or
       \( \ip{A}{B} = \sum_{i} a_{i} b_{i} \).
The spaces you type in a formula are
ignored.  Remember that a letter like
       $x$                   % $ ... $  and  \( ... \)  are equivalent
is a formula when it denotes a mathematical
symbol, and it should be typed as one.

\section{Displayed Text}

Text is displayed by indenting it from the left
margin.  Quotations are commonly displayed.  There
are short quotations
\begin{quote}
   This is a short quotation.  It consists of a
   single paragraph of text.  See how it is formatted.
\end{quote}
and longer ones.
\begin{quotation}
   This is a longer quotation.  It consists of two
   paragraphs of text, neither of which are
   particularly interesting.

   This is the second paragraph of the quotation.  It
   is just as dull as the first paragraph.
\end{quotation}
Another frequently-displayed structure is a list.
The following is an example of an \emph{itemized}
list.
\begin{itemize}
   \item This is the first item of an itemized list.
         Each item in the list is marked with a ``tick''.
         You don't have to worry about what kind of tick
         mark is used.

   \item This is the second item of the list.  It
         contains another list nested inside it.  The inner
         list is an \emph{enumerated} list.
         \begin{enumerate}
            \item This is the first item of an enumerated
                  list that is nested within the itemized list.

            \item This is the second item of the inner list.
                  \LaTeX\ allows you to nest lists deeper than
                  you really should.
         \end{enumerate}
         This is the rest of the second item of the outer
         list.  It is no more interesting than any other
         part of the item.
   \item This is the third item of the list.
\end{itemize}
You can even display poetry.
\begin{verse}
   There is an environment
    for verse \\             % The \\ command separates lines
   Whose features some poets % within a stanza.
   will curse.

                             % One or more blank lines separate stanzas.

   For instead of making\\
   Them do \emph{all} line breaking, \\
   It allows them to put too many words on a line when they'd rather be
   forced to be terse.
\end{verse}

Mathematical formulas may also be displayed.  A
displayed formula
is
one-line long; multiline
formulas require special formatting instructions.
   \[  \ip{\Gamma}{\psi'} = x'' + y^{2} + z_{i}^{n}\]
Don't start a paragraph with a displayed equation,
nor make one a paragraph by itself.

\end{document}               % End of document.

\section{Bench Huge Section 947}
% synthetic bench section 947
% This is a sample LaTeX input file.  (Version of 12 August 2004.)
%
% A '%' character causes TeX to ignore all remaining text on the line,
% and is used for comments like this one.

\documentclass{article}      % Specifies the document class

                             % The preamble begins here.
\title{An Example Document}  % Declares the document's title.
\author{Leslie Lamport}      % Declares the author's name.
\date{January 21, 1994}      % Deleting this command produces today's date.

\newcommand{\ip}[2]{(#1, #2)}
                             % Defines \ip{arg1}{arg2} to mean
                             % (arg1, arg2).

%\newcommand{\ip}[2]{\langle #1 | #2\rangle}
                             % This is an alternative definition of
                             % \ip that is commented out.

\begin{document}             % End of preamble and beginning of text.

\maketitle                   % Produces the title.

This is an example input file.  Comparing it with
the output it generates can show you how to
produce a simple document of your own.

\section{Ordinary Text}      % Produces section heading.  Lower-level
                             % sections are begun with similar
                             % \subsection and \subsubsection commands.

The ends  of words and sentences are marked
  by   spaces. It  doesn't matter how many
spaces    you type; one is as good as 100.  The
end of   a line counts as a space.

One   or more   blank lines denote the  end
of  a paragraph.

Since any number of consecutive spaces are treated
like a single one, the formatting of the input
file makes no difference to
      \LaTeX,                % The \LaTeX command generates the LaTeX logo.
but it makes a difference to you.  When you use
\LaTeX, making your input file as easy to read
as possible will be a great help as you write
your document and when you change it.  This sample
file shows how you can add comments to your own input
file.

Because printing is different from typewriting,
there are a number of things that you have to do
differently when preparing an input file than if
you were just typing the document directly.
Quotation marks like
       ``this''
have to be handled specially, as do quotes within
quotes:
       ``\,`this'            % \, separates the double and single quote.
        is what I just
        wrote, not  `that'\,''.

Dashes come in three sizes: an
       intra-word
dash, a medium dash for number ranges like
       1--2,
and a punctuation
       dash---like
this.

A sentence-ending space should be larger than the
space between words within a sentence.  You
sometimes have to type special commands in
conjunction with punctuation characters to get
this right, as in the following sentence.
       Gnats, gnus, etc.\ all  % `\ ' makes an inter-word space.
       begin with G\@.         % \@ marks end-of-sentence punctuation.
You should check the spaces after periods when
reading your output to make sure you haven't
forgotten any special cases.  Generating an
ellipsis
       \ldots\               % `\ ' is needed after `\ldots' because TeX
                             % ignores spaces after command names like \ldots
                             % made from \ + letters.
                             %
                             % Note how a `%' character causes TeX to ignore
                             % the end of the input line, so these blank lines
                             % do not start a new paragraph.
                             %
with the right spacing around the periods requires
a special command.

\LaTeX\ interprets some common characters as
commands, so you must type special commands to
generate them.  These characters include the
following:
       \$ \& \% \# \{ and \}.

In printing, text is usually emphasized with an
       \emph{italic}
type style.

\begin{em}
   A long segment of text can also be emphasized
   in this way.  Text within such a segment can be
   given \emph{additional} emphasis.
\end{em}

It is sometimes necessary to prevent \LaTeX\ from
breaking a line where it might otherwise do so.
This may be at a space, as between the ``Mr.''\ and
``Jones'' in
       ``Mr.~Jones'',        % ~ produces an unbreakable interword space.
or within a word---especially when the word is a
symbol like
       \mbox{\emph{itemnum}}
that makes little sense when hyphenated across
lines.

Footnotes\footnote{This is an example of a footnote.}
pose no problem.

\LaTeX\ is good at typesetting mathematical formulas
like
       \( x-3y + z = 7 \)
or
       \( a_{1} > x^{2n} + y^{2n} > x' \)
or
       \( \ip{A}{B} = \sum_{i} a_{i} b_{i} \).
The spaces you type in a formula are
ignored.  Remember that a letter like
       $x$                   % $ ... $  and  \( ... \)  are equivalent
is a formula when it denotes a mathematical
symbol, and it should be typed as one.

\section{Displayed Text}

Text is displayed by indenting it from the left
margin.  Quotations are commonly displayed.  There
are short quotations
\begin{quote}
   This is a short quotation.  It consists of a
   single paragraph of text.  See how it is formatted.
\end{quote}
and longer ones.
\begin{quotation}
   This is a longer quotation.  It consists of two
   paragraphs of text, neither of which are
   particularly interesting.

   This is the second paragraph of the quotation.  It
   is just as dull as the first paragraph.
\end{quotation}
Another frequently-displayed structure is a list.
The following is an example of an \emph{itemized}
list.
\begin{itemize}
   \item This is the first item of an itemized list.
         Each item in the list is marked with a ``tick''.
         You don't have to worry about what kind of tick
         mark is used.

   \item This is the second item of the list.  It
         contains another list nested inside it.  The inner
         list is an \emph{enumerated} list.
         \begin{enumerate}
            \item This is the first item of an enumerated
                  list that is nested within the itemized list.

            \item This is the second item of the inner list.
                  \LaTeX\ allows you to nest lists deeper than
                  you really should.
         \end{enumerate}
         This is the rest of the second item of the outer
         list.  It is no more interesting than any other
         part of the item.
   \item This is the third item of the list.
\end{itemize}
You can even display poetry.
\begin{verse}
   There is an environment
    for verse \\             % The \\ command separates lines
   Whose features some poets % within a stanza.
   will curse.

                             % One or more blank lines separate stanzas.

   For instead of making\\
   Them do \emph{all} line breaking, \\
   It allows them to put too many words on a line when they'd rather be
   forced to be terse.
\end{verse}

Mathematical formulas may also be displayed.  A
displayed formula
is
one-line long; multiline
formulas require special formatting instructions.
   \[  \ip{\Gamma}{\psi'} = x'' + y^{2} + z_{i}^{n}\]
Don't start a paragraph with a displayed equation,
nor make one a paragraph by itself.

\end{document}               % End of document.

\section{Bench Huge Section 948}
% synthetic bench section 948
% This is a sample LaTeX input file.  (Version of 12 August 2004.)
%
% A '%' character causes TeX to ignore all remaining text on the line,
% and is used for comments like this one.

\documentclass{article}      % Specifies the document class

                             % The preamble begins here.
\title{An Example Document}  % Declares the document's title.
\author{Leslie Lamport}      % Declares the author's name.
\date{January 21, 1994}      % Deleting this command produces today's date.

\newcommand{\ip}[2]{(#1, #2)}
                             % Defines \ip{arg1}{arg2} to mean
                             % (arg1, arg2).

%\newcommand{\ip}[2]{\langle #1 | #2\rangle}
                             % This is an alternative definition of
                             % \ip that is commented out.

\begin{document}             % End of preamble and beginning of text.

\maketitle                   % Produces the title.

This is an example input file.  Comparing it with
the output it generates can show you how to
produce a simple document of your own.

\section{Ordinary Text}      % Produces section heading.  Lower-level
                             % sections are begun with similar
                             % \subsection and \subsubsection commands.

The ends  of words and sentences are marked
  by   spaces. It  doesn't matter how many
spaces    you type; one is as good as 100.  The
end of   a line counts as a space.

One   or more   blank lines denote the  end
of  a paragraph.

Since any number of consecutive spaces are treated
like a single one, the formatting of the input
file makes no difference to
      \LaTeX,                % The \LaTeX command generates the LaTeX logo.
but it makes a difference to you.  When you use
\LaTeX, making your input file as easy to read
as possible will be a great help as you write
your document and when you change it.  This sample
file shows how you can add comments to your own input
file.

Because printing is different from typewriting,
there are a number of things that you have to do
differently when preparing an input file than if
you were just typing the document directly.
Quotation marks like
       ``this''
have to be handled specially, as do quotes within
quotes:
       ``\,`this'            % \, separates the double and single quote.
        is what I just
        wrote, not  `that'\,''.

Dashes come in three sizes: an
       intra-word
dash, a medium dash for number ranges like
       1--2,
and a punctuation
       dash---like
this.

A sentence-ending space should be larger than the
space between words within a sentence.  You
sometimes have to type special commands in
conjunction with punctuation characters to get
this right, as in the following sentence.
       Gnats, gnus, etc.\ all  % `\ ' makes an inter-word space.
       begin with G\@.         % \@ marks end-of-sentence punctuation.
You should check the spaces after periods when
reading your output to make sure you haven't
forgotten any special cases.  Generating an
ellipsis
       \ldots\               % `\ ' is needed after `\ldots' because TeX
                             % ignores spaces after command names like \ldots
                             % made from \ + letters.
                             %
                             % Note how a `%' character causes TeX to ignore
                             % the end of the input line, so these blank lines
                             % do not start a new paragraph.
                             %
with the right spacing around the periods requires
a special command.

\LaTeX\ interprets some common characters as
commands, so you must type special commands to
generate them.  These characters include the
following:
       \$ \& \% \# \{ and \}.

In printing, text is usually emphasized with an
       \emph{italic}
type style.

\begin{em}
   A long segment of text can also be emphasized
   in this way.  Text within such a segment can be
   given \emph{additional} emphasis.
\end{em}

It is sometimes necessary to prevent \LaTeX\ from
breaking a line where it might otherwise do so.
This may be at a space, as between the ``Mr.''\ and
``Jones'' in
       ``Mr.~Jones'',        % ~ produces an unbreakable interword space.
or within a word---especially when the word is a
symbol like
       \mbox{\emph{itemnum}}
that makes little sense when hyphenated across
lines.

Footnotes\footnote{This is an example of a footnote.}
pose no problem.

\LaTeX\ is good at typesetting mathematical formulas
like
       \( x-3y + z = 7 \)
or
       \( a_{1} > x^{2n} + y^{2n} > x' \)
or
       \( \ip{A}{B} = \sum_{i} a_{i} b_{i} \).
The spaces you type in a formula are
ignored.  Remember that a letter like
       $x$                   % $ ... $  and  \( ... \)  are equivalent
is a formula when it denotes a mathematical
symbol, and it should be typed as one.

\section{Displayed Text}

Text is displayed by indenting it from the left
margin.  Quotations are commonly displayed.  There
are short quotations
\begin{quote}
   This is a short quotation.  It consists of a
   single paragraph of text.  See how it is formatted.
\end{quote}
and longer ones.
\begin{quotation}
   This is a longer quotation.  It consists of two
   paragraphs of text, neither of which are
   particularly interesting.

   This is the second paragraph of the quotation.  It
   is just as dull as the first paragraph.
\end{quotation}
Another frequently-displayed structure is a list.
The following is an example of an \emph{itemized}
list.
\begin{itemize}
   \item This is the first item of an itemized list.
         Each item in the list is marked with a ``tick''.
         You don't have to worry about what kind of tick
         mark is used.

   \item This is the second item of the list.  It
         contains another list nested inside it.  The inner
         list is an \emph{enumerated} list.
         \begin{enumerate}
            \item This is the first item of an enumerated
                  list that is nested within the itemized list.

            \item This is the second item of the inner list.
                  \LaTeX\ allows you to nest lists deeper than
                  you really should.
         \end{enumerate}
         This is the rest of the second item of the outer
         list.  It is no more interesting than any other
         part of the item.
   \item This is the third item of the list.
\end{itemize}
You can even display poetry.
\begin{verse}
   There is an environment
    for verse \\             % The \\ command separates lines
   Whose features some poets % within a stanza.
   will curse.

                             % One or more blank lines separate stanzas.

   For instead of making\\
   Them do \emph{all} line breaking, \\
   It allows them to put too many words on a line when they'd rather be
   forced to be terse.
\end{verse}

Mathematical formulas may also be displayed.  A
displayed formula
is
one-line long; multiline
formulas require special formatting instructions.
   \[  \ip{\Gamma}{\psi'} = x'' + y^{2} + z_{i}^{n}\]
Don't start a paragraph with a displayed equation,
nor make one a paragraph by itself.

\end{document}               % End of document.

\section{Bench Huge Section 949}
% synthetic bench section 949
% This is a sample LaTeX input file.  (Version of 12 August 2004.)
%
% A '%' character causes TeX to ignore all remaining text on the line,
% and is used for comments like this one.

\documentclass{article}      % Specifies the document class

                             % The preamble begins here.
\title{An Example Document}  % Declares the document's title.
\author{Leslie Lamport}      % Declares the author's name.
\date{January 21, 1994}      % Deleting this command produces today's date.

\newcommand{\ip}[2]{(#1, #2)}
                             % Defines \ip{arg1}{arg2} to mean
                             % (arg1, arg2).

%\newcommand{\ip}[2]{\langle #1 | #2\rangle}
                             % This is an alternative definition of
                             % \ip that is commented out.

\begin{document}             % End of preamble and beginning of text.

\maketitle                   % Produces the title.

This is an example input file.  Comparing it with
the output it generates can show you how to
produce a simple document of your own.

\section{Ordinary Text}      % Produces section heading.  Lower-level
                             % sections are begun with similar
                             % \subsection and \subsubsection commands.

The ends  of words and sentences are marked
  by   spaces. It  doesn't matter how many
spaces    you type; one is as good as 100.  The
end of   a line counts as a space.

One   or more   blank lines denote the  end
of  a paragraph.

Since any number of consecutive spaces are treated
like a single one, the formatting of the input
file makes no difference to
      \LaTeX,                % The \LaTeX command generates the LaTeX logo.
but it makes a difference to you.  When you use
\LaTeX, making your input file as easy to read
as possible will be a great help as you write
your document and when you change it.  This sample
file shows how you can add comments to your own input
file.

Because printing is different from typewriting,
there are a number of things that you have to do
differently when preparing an input file than if
you were just typing the document directly.
Quotation marks like
       ``this''
have to be handled specially, as do quotes within
quotes:
       ``\,`this'            % \, separates the double and single quote.
        is what I just
        wrote, not  `that'\,''.

Dashes come in three sizes: an
       intra-word
dash, a medium dash for number ranges like
       1--2,
and a punctuation
       dash---like
this.

A sentence-ending space should be larger than the
space between words within a sentence.  You
sometimes have to type special commands in
conjunction with punctuation characters to get
this right, as in the following sentence.
       Gnats, gnus, etc.\ all  % `\ ' makes an inter-word space.
       begin with G\@.         % \@ marks end-of-sentence punctuation.
You should check the spaces after periods when
reading your output to make sure you haven't
forgotten any special cases.  Generating an
ellipsis
       \ldots\               % `\ ' is needed after `\ldots' because TeX
                             % ignores spaces after command names like \ldots
                             % made from \ + letters.
                             %
                             % Note how a `%' character causes TeX to ignore
                             % the end of the input line, so these blank lines
                             % do not start a new paragraph.
                             %
with the right spacing around the periods requires
a special command.

\LaTeX\ interprets some common characters as
commands, so you must type special commands to
generate them.  These characters include the
following:
       \$ \& \% \# \{ and \}.

In printing, text is usually emphasized with an
       \emph{italic}
type style.

\begin{em}
   A long segment of text can also be emphasized
   in this way.  Text within such a segment can be
   given \emph{additional} emphasis.
\end{em}

It is sometimes necessary to prevent \LaTeX\ from
breaking a line where it might otherwise do so.
This may be at a space, as between the ``Mr.''\ and
``Jones'' in
       ``Mr.~Jones'',        % ~ produces an unbreakable interword space.
or within a word---especially when the word is a
symbol like
       \mbox{\emph{itemnum}}
that makes little sense when hyphenated across
lines.

Footnotes\footnote{This is an example of a footnote.}
pose no problem.

\LaTeX\ is good at typesetting mathematical formulas
like
       \( x-3y + z = 7 \)
or
       \( a_{1} > x^{2n} + y^{2n} > x' \)
or
       \( \ip{A}{B} = \sum_{i} a_{i} b_{i} \).
The spaces you type in a formula are
ignored.  Remember that a letter like
       $x$                   % $ ... $  and  \( ... \)  are equivalent
is a formula when it denotes a mathematical
symbol, and it should be typed as one.

\section{Displayed Text}

Text is displayed by indenting it from the left
margin.  Quotations are commonly displayed.  There
are short quotations
\begin{quote}
   This is a short quotation.  It consists of a
   single paragraph of text.  See how it is formatted.
\end{quote}
and longer ones.
\begin{quotation}
   This is a longer quotation.  It consists of two
   paragraphs of text, neither of which are
   particularly interesting.

   This is the second paragraph of the quotation.  It
   is just as dull as the first paragraph.
\end{quotation}
Another frequently-displayed structure is a list.
The following is an example of an \emph{itemized}
list.
\begin{itemize}
   \item This is the first item of an itemized list.
         Each item in the list is marked with a ``tick''.
         You don't have to worry about what kind of tick
         mark is used.

   \item This is the second item of the list.  It
         contains another list nested inside it.  The inner
         list is an \emph{enumerated} list.
         \begin{enumerate}
            \item This is the first item of an enumerated
                  list that is nested within the itemized list.

            \item This is the second item of the inner list.
                  \LaTeX\ allows you to nest lists deeper than
                  you really should.
         \end{enumerate}
         This is the rest of the second item of the outer
         list.  It is no more interesting than any other
         part of the item.
   \item This is the third item of the list.
\end{itemize}
You can even display poetry.
\begin{verse}
   There is an environment
    for verse \\             % The \\ command separates lines
   Whose features some poets % within a stanza.
   will curse.

                             % One or more blank lines separate stanzas.

   For instead of making\\
   Them do \emph{all} line breaking, \\
   It allows them to put too many words on a line when they'd rather be
   forced to be terse.
\end{verse}

Mathematical formulas may also be displayed.  A
displayed formula
is
one-line long; multiline
formulas require special formatting instructions.
   \[  \ip{\Gamma}{\psi'} = x'' + y^{2} + z_{i}^{n}\]
Don't start a paragraph with a displayed equation,
nor make one a paragraph by itself.

\end{document}               % End of document.

\section{Bench Huge Section 950}
% synthetic bench section 950
% This is a sample LaTeX input file.  (Version of 12 August 2004.)
%
% A '%' character causes TeX to ignore all remaining text on the line,
% and is used for comments like this one.

\documentclass{article}      % Specifies the document class

                             % The preamble begins here.
\title{An Example Document}  % Declares the document's title.
\author{Leslie Lamport}      % Declares the author's name.
\date{January 21, 1994}      % Deleting this command produces today's date.

\newcommand{\ip}[2]{(#1, #2)}
                             % Defines \ip{arg1}{arg2} to mean
                             % (arg1, arg2).

%\newcommand{\ip}[2]{\langle #1 | #2\rangle}
                             % This is an alternative definition of
                             % \ip that is commented out.

\begin{document}             % End of preamble and beginning of text.

\maketitle                   % Produces the title.

This is an example input file.  Comparing it with
the output it generates can show you how to
produce a simple document of your own.

\section{Ordinary Text}      % Produces section heading.  Lower-level
                             % sections are begun with similar
                             % \subsection and \subsubsection commands.

The ends  of words and sentences are marked
  by   spaces. It  doesn't matter how many
spaces    you type; one is as good as 100.  The
end of   a line counts as a space.

One   or more   blank lines denote the  end
of  a paragraph.

Since any number of consecutive spaces are treated
like a single one, the formatting of the input
file makes no difference to
      \LaTeX,                % The \LaTeX command generates the LaTeX logo.
but it makes a difference to you.  When you use
\LaTeX, making your input file as easy to read
as possible will be a great help as you write
your document and when you change it.  This sample
file shows how you can add comments to your own input
file.

Because printing is different from typewriting,
there are a number of things that you have to do
differently when preparing an input file than if
you were just typing the document directly.
Quotation marks like
       ``this''
have to be handled specially, as do quotes within
quotes:
       ``\,`this'            % \, separates the double and single quote.
        is what I just
        wrote, not  `that'\,''.

Dashes come in three sizes: an
       intra-word
dash, a medium dash for number ranges like
       1--2,
and a punctuation
       dash---like
this.

A sentence-ending space should be larger than the
space between words within a sentence.  You
sometimes have to type special commands in
conjunction with punctuation characters to get
this right, as in the following sentence.
       Gnats, gnus, etc.\ all  % `\ ' makes an inter-word space.
       begin with G\@.         % \@ marks end-of-sentence punctuation.
You should check the spaces after periods when
reading your output to make sure you haven't
forgotten any special cases.  Generating an
ellipsis
       \ldots\               % `\ ' is needed after `\ldots' because TeX
                             % ignores spaces after command names like \ldots
                             % made from \ + letters.
                             %
                             % Note how a `%' character causes TeX to ignore
                             % the end of the input line, so these blank lines
                             % do not start a new paragraph.
                             %
with the right spacing around the periods requires
a special command.

\LaTeX\ interprets some common characters as
commands, so you must type special commands to
generate them.  These characters include the
following:
       \$ \& \% \# \{ and \}.

In printing, text is usually emphasized with an
       \emph{italic}
type style.

\begin{em}
   A long segment of text can also be emphasized
   in this way.  Text within such a segment can be
   given \emph{additional} emphasis.
\end{em}

It is sometimes necessary to prevent \LaTeX\ from
breaking a line where it might otherwise do so.
This may be at a space, as between the ``Mr.''\ and
``Jones'' in
       ``Mr.~Jones'',        % ~ produces an unbreakable interword space.
or within a word---especially when the word is a
symbol like
       \mbox{\emph{itemnum}}
that makes little sense when hyphenated across
lines.

Footnotes\footnote{This is an example of a footnote.}
pose no problem.

\LaTeX\ is good at typesetting mathematical formulas
like
       \( x-3y + z = 7 \)
or
       \( a_{1} > x^{2n} + y^{2n} > x' \)
or
       \( \ip{A}{B} = \sum_{i} a_{i} b_{i} \).
The spaces you type in a formula are
ignored.  Remember that a letter like
       $x$                   % $ ... $  and  \( ... \)  are equivalent
is a formula when it denotes a mathematical
symbol, and it should be typed as one.

\section{Displayed Text}

Text is displayed by indenting it from the left
margin.  Quotations are commonly displayed.  There
are short quotations
\begin{quote}
   This is a short quotation.  It consists of a
   single paragraph of text.  See how it is formatted.
\end{quote}
and longer ones.
\begin{quotation}
   This is a longer quotation.  It consists of two
   paragraphs of text, neither of which are
   particularly interesting.

   This is the second paragraph of the quotation.  It
   is just as dull as the first paragraph.
\end{quotation}
Another frequently-displayed structure is a list.
The following is an example of an \emph{itemized}
list.
\begin{itemize}
   \item This is the first item of an itemized list.
         Each item in the list is marked with a ``tick''.
         You don't have to worry about what kind of tick
         mark is used.

   \item This is the second item of the list.  It
         contains another list nested inside it.  The inner
         list is an \emph{enumerated} list.
         \begin{enumerate}
            \item This is the first item of an enumerated
                  list that is nested within the itemized list.

            \item This is the second item of the inner list.
                  \LaTeX\ allows you to nest lists deeper than
                  you really should.
         \end{enumerate}
         This is the rest of the second item of the outer
         list.  It is no more interesting than any other
         part of the item.
   \item This is the third item of the list.
\end{itemize}
You can even display poetry.
\begin{verse}
   There is an environment
    for verse \\             % The \\ command separates lines
   Whose features some poets % within a stanza.
   will curse.

                             % One or more blank lines separate stanzas.

   For instead of making\\
   Them do \emph{all} line breaking, \\
   It allows them to put too many words on a line when they'd rather be
   forced to be terse.
\end{verse}

Mathematical formulas may also be displayed.  A
displayed formula
is
one-line long; multiline
formulas require special formatting instructions.
   \[  \ip{\Gamma}{\psi'} = x'' + y^{2} + z_{i}^{n}\]
Don't start a paragraph with a displayed equation,
nor make one a paragraph by itself.

\end{document}               % End of document.

\section{Bench Huge Section 951}
% synthetic bench section 951
% This is a sample LaTeX input file.  (Version of 12 August 2004.)
%
% A '%' character causes TeX to ignore all remaining text on the line,
% and is used for comments like this one.

\documentclass{article}      % Specifies the document class

                             % The preamble begins here.
\title{An Example Document}  % Declares the document's title.
\author{Leslie Lamport}      % Declares the author's name.
\date{January 21, 1994}      % Deleting this command produces today's date.

\newcommand{\ip}[2]{(#1, #2)}
                             % Defines \ip{arg1}{arg2} to mean
                             % (arg1, arg2).

%\newcommand{\ip}[2]{\langle #1 | #2\rangle}
                             % This is an alternative definition of
                             % \ip that is commented out.

\begin{document}             % End of preamble and beginning of text.

\maketitle                   % Produces the title.

This is an example input file.  Comparing it with
the output it generates can show you how to
produce a simple document of your own.

\section{Ordinary Text}      % Produces section heading.  Lower-level
                             % sections are begun with similar
                             % \subsection and \subsubsection commands.

The ends  of words and sentences are marked
  by   spaces. It  doesn't matter how many
spaces    you type; one is as good as 100.  The
end of   a line counts as a space.

One   or more   blank lines denote the  end
of  a paragraph.

Since any number of consecutive spaces are treated
like a single one, the formatting of the input
file makes no difference to
      \LaTeX,                % The \LaTeX command generates the LaTeX logo.
but it makes a difference to you.  When you use
\LaTeX, making your input file as easy to read
as possible will be a great help as you write
your document and when you change it.  This sample
file shows how you can add comments to your own input
file.

Because printing is different from typewriting,
there are a number of things that you have to do
differently when preparing an input file than if
you were just typing the document directly.
Quotation marks like
       ``this''
have to be handled specially, as do quotes within
quotes:
       ``\,`this'            % \, separates the double and single quote.
        is what I just
        wrote, not  `that'\,''.

Dashes come in three sizes: an
       intra-word
dash, a medium dash for number ranges like
       1--2,
and a punctuation
       dash---like
this.

A sentence-ending space should be larger than the
space between words within a sentence.  You
sometimes have to type special commands in
conjunction with punctuation characters to get
this right, as in the following sentence.
       Gnats, gnus, etc.\ all  % `\ ' makes an inter-word space.
       begin with G\@.         % \@ marks end-of-sentence punctuation.
You should check the spaces after periods when
reading your output to make sure you haven't
forgotten any special cases.  Generating an
ellipsis
       \ldots\               % `\ ' is needed after `\ldots' because TeX
                             % ignores spaces after command names like \ldots
                             % made from \ + letters.
                             %
                             % Note how a `%' character causes TeX to ignore
                             % the end of the input line, so these blank lines
                             % do not start a new paragraph.
                             %
with the right spacing around the periods requires
a special command.

\LaTeX\ interprets some common characters as
commands, so you must type special commands to
generate them.  These characters include the
following:
       \$ \& \% \# \{ and \}.

In printing, text is usually emphasized with an
       \emph{italic}
type style.

\begin{em}
   A long segment of text can also be emphasized
   in this way.  Text within such a segment can be
   given \emph{additional} emphasis.
\end{em}

It is sometimes necessary to prevent \LaTeX\ from
breaking a line where it might otherwise do so.
This may be at a space, as between the ``Mr.''\ and
``Jones'' in
       ``Mr.~Jones'',        % ~ produces an unbreakable interword space.
or within a word---especially when the word is a
symbol like
       \mbox{\emph{itemnum}}
that makes little sense when hyphenated across
lines.

Footnotes\footnote{This is an example of a footnote.}
pose no problem.

\LaTeX\ is good at typesetting mathematical formulas
like
       \( x-3y + z = 7 \)
or
       \( a_{1} > x^{2n} + y^{2n} > x' \)
or
       \( \ip{A}{B} = \sum_{i} a_{i} b_{i} \).
The spaces you type in a formula are
ignored.  Remember that a letter like
       $x$                   % $ ... $  and  \( ... \)  are equivalent
is a formula when it denotes a mathematical
symbol, and it should be typed as one.

\section{Displayed Text}

Text is displayed by indenting it from the left
margin.  Quotations are commonly displayed.  There
are short quotations
\begin{quote}
   This is a short quotation.  It consists of a
   single paragraph of text.  See how it is formatted.
\end{quote}
and longer ones.
\begin{quotation}
   This is a longer quotation.  It consists of two
   paragraphs of text, neither of which are
   particularly interesting.

   This is the second paragraph of the quotation.  It
   is just as dull as the first paragraph.
\end{quotation}
Another frequently-displayed structure is a list.
The following is an example of an \emph{itemized}
list.
\begin{itemize}
   \item This is the first item of an itemized list.
         Each item in the list is marked with a ``tick''.
         You don't have to worry about what kind of tick
         mark is used.

   \item This is the second item of the list.  It
         contains another list nested inside it.  The inner
         list is an \emph{enumerated} list.
         \begin{enumerate}
            \item This is the first item of an enumerated
                  list that is nested within the itemized list.

            \item This is the second item of the inner list.
                  \LaTeX\ allows you to nest lists deeper than
                  you really should.
         \end{enumerate}
         This is the rest of the second item of the outer
         list.  It is no more interesting than any other
         part of the item.
   \item This is the third item of the list.
\end{itemize}
You can even display poetry.
\begin{verse}
   There is an environment
    for verse \\             % The \\ command separates lines
   Whose features some poets % within a stanza.
   will curse.

                             % One or more blank lines separate stanzas.

   For instead of making\\
   Them do \emph{all} line breaking, \\
   It allows them to put too many words on a line when they'd rather be
   forced to be terse.
\end{verse}

Mathematical formulas may also be displayed.  A
displayed formula
is
one-line long; multiline
formulas require special formatting instructions.
   \[  \ip{\Gamma}{\psi'} = x'' + y^{2} + z_{i}^{n}\]
Don't start a paragraph with a displayed equation,
nor make one a paragraph by itself.

\end{document}               % End of document.

\section{Bench Huge Section 952}
% synthetic bench section 952
% This is a sample LaTeX input file.  (Version of 12 August 2004.)
%
% A '%' character causes TeX to ignore all remaining text on the line,
% and is used for comments like this one.

\documentclass{article}      % Specifies the document class

                             % The preamble begins here.
\title{An Example Document}  % Declares the document's title.
\author{Leslie Lamport}      % Declares the author's name.
\date{January 21, 1994}      % Deleting this command produces today's date.

\newcommand{\ip}[2]{(#1, #2)}
                             % Defines \ip{arg1}{arg2} to mean
                             % (arg1, arg2).

%\newcommand{\ip}[2]{\langle #1 | #2\rangle}
                             % This is an alternative definition of
                             % \ip that is commented out.

\begin{document}             % End of preamble and beginning of text.

\maketitle                   % Produces the title.

This is an example input file.  Comparing it with
the output it generates can show you how to
produce a simple document of your own.

\section{Ordinary Text}      % Produces section heading.  Lower-level
                             % sections are begun with similar
                             % \subsection and \subsubsection commands.

The ends  of words and sentences are marked
  by   spaces. It  doesn't matter how many
spaces    you type; one is as good as 100.  The
end of   a line counts as a space.

One   or more   blank lines denote the  end
of  a paragraph.

Since any number of consecutive spaces are treated
like a single one, the formatting of the input
file makes no difference to
      \LaTeX,                % The \LaTeX command generates the LaTeX logo.
but it makes a difference to you.  When you use
\LaTeX, making your input file as easy to read
as possible will be a great help as you write
your document and when you change it.  This sample
file shows how you can add comments to your own input
file.

Because printing is different from typewriting,
there are a number of things that you have to do
differently when preparing an input file than if
you were just typing the document directly.
Quotation marks like
       ``this''
have to be handled specially, as do quotes within
quotes:
       ``\,`this'            % \, separates the double and single quote.
        is what I just
        wrote, not  `that'\,''.

Dashes come in three sizes: an
       intra-word
dash, a medium dash for number ranges like
       1--2,
and a punctuation
       dash---like
this.

A sentence-ending space should be larger than the
space between words within a sentence.  You
sometimes have to type special commands in
conjunction with punctuation characters to get
this right, as in the following sentence.
       Gnats, gnus, etc.\ all  % `\ ' makes an inter-word space.
       begin with G\@.         % \@ marks end-of-sentence punctuation.
You should check the spaces after periods when
reading your output to make sure you haven't
forgotten any special cases.  Generating an
ellipsis
       \ldots\               % `\ ' is needed after `\ldots' because TeX
                             % ignores spaces after command names like \ldots
                             % made from \ + letters.
                             %
                             % Note how a `%' character causes TeX to ignore
                             % the end of the input line, so these blank lines
                             % do not start a new paragraph.
                             %
with the right spacing around the periods requires
a special command.

\LaTeX\ interprets some common characters as
commands, so you must type special commands to
generate them.  These characters include the
following:
       \$ \& \% \# \{ and \}.

In printing, text is usually emphasized with an
       \emph{italic}
type style.

\begin{em}
   A long segment of text can also be emphasized
   in this way.  Text within such a segment can be
   given \emph{additional} emphasis.
\end{em}

It is sometimes necessary to prevent \LaTeX\ from
breaking a line where it might otherwise do so.
This may be at a space, as between the ``Mr.''\ and
``Jones'' in
       ``Mr.~Jones'',        % ~ produces an unbreakable interword space.
or within a word---especially when the word is a
symbol like
       \mbox{\emph{itemnum}}
that makes little sense when hyphenated across
lines.

Footnotes\footnote{This is an example of a footnote.}
pose no problem.

\LaTeX\ is good at typesetting mathematical formulas
like
       \( x-3y + z = 7 \)
or
       \( a_{1} > x^{2n} + y^{2n} > x' \)
or
       \( \ip{A}{B} = \sum_{i} a_{i} b_{i} \).
The spaces you type in a formula are
ignored.  Remember that a letter like
       $x$                   % $ ... $  and  \( ... \)  are equivalent
is a formula when it denotes a mathematical
symbol, and it should be typed as one.

\section{Displayed Text}

Text is displayed by indenting it from the left
margin.  Quotations are commonly displayed.  There
are short quotations
\begin{quote}
   This is a short quotation.  It consists of a
   single paragraph of text.  See how it is formatted.
\end{quote}
and longer ones.
\begin{quotation}
   This is a longer quotation.  It consists of two
   paragraphs of text, neither of which are
   particularly interesting.

   This is the second paragraph of the quotation.  It
   is just as dull as the first paragraph.
\end{quotation}
Another frequently-displayed structure is a list.
The following is an example of an \emph{itemized}
list.
\begin{itemize}
   \item This is the first item of an itemized list.
         Each item in the list is marked with a ``tick''.
         You don't have to worry about what kind of tick
         mark is used.

   \item This is the second item of the list.  It
         contains another list nested inside it.  The inner
         list is an \emph{enumerated} list.
         \begin{enumerate}
            \item This is the first item of an enumerated
                  list that is nested within the itemized list.

            \item This is the second item of the inner list.
                  \LaTeX\ allows you to nest lists deeper than
                  you really should.
         \end{enumerate}
         This is the rest of the second item of the outer
         list.  It is no more interesting than any other
         part of the item.
   \item This is the third item of the list.
\end{itemize}
You can even display poetry.
\begin{verse}
   There is an environment
    for verse \\             % The \\ command separates lines
   Whose features some poets % within a stanza.
   will curse.

                             % One or more blank lines separate stanzas.

   For instead of making\\
   Them do \emph{all} line breaking, \\
   It allows them to put too many words on a line when they'd rather be
   forced to be terse.
\end{verse}

Mathematical formulas may also be displayed.  A
displayed formula
is
one-line long; multiline
formulas require special formatting instructions.
   \[  \ip{\Gamma}{\psi'} = x'' + y^{2} + z_{i}^{n}\]
Don't start a paragraph with a displayed equation,
nor make one a paragraph by itself.

\end{document}               % End of document.

\section{Bench Huge Section 953}
% synthetic bench section 953
% This is a sample LaTeX input file.  (Version of 12 August 2004.)
%
% A '%' character causes TeX to ignore all remaining text on the line,
% and is used for comments like this one.

\documentclass{article}      % Specifies the document class

                             % The preamble begins here.
\title{An Example Document}  % Declares the document's title.
\author{Leslie Lamport}      % Declares the author's name.
\date{January 21, 1994}      % Deleting this command produces today's date.

\newcommand{\ip}[2]{(#1, #2)}
                             % Defines \ip{arg1}{arg2} to mean
                             % (arg1, arg2).

%\newcommand{\ip}[2]{\langle #1 | #2\rangle}
                             % This is an alternative definition of
                             % \ip that is commented out.

\begin{document}             % End of preamble and beginning of text.

\maketitle                   % Produces the title.

This is an example input file.  Comparing it with
the output it generates can show you how to
produce a simple document of your own.

\section{Ordinary Text}      % Produces section heading.  Lower-level
                             % sections are begun with similar
                             % \subsection and \subsubsection commands.

The ends  of words and sentences are marked
  by   spaces. It  doesn't matter how many
spaces    you type; one is as good as 100.  The
end of   a line counts as a space.

One   or more   blank lines denote the  end
of  a paragraph.

Since any number of consecutive spaces are treated
like a single one, the formatting of the input
file makes no difference to
      \LaTeX,                % The \LaTeX command generates the LaTeX logo.
but it makes a difference to you.  When you use
\LaTeX, making your input file as easy to read
as possible will be a great help as you write
your document and when you change it.  This sample
file shows how you can add comments to your own input
file.

Because printing is different from typewriting,
there are a number of things that you have to do
differently when preparing an input file than if
you were just typing the document directly.
Quotation marks like
       ``this''
have to be handled specially, as do quotes within
quotes:
       ``\,`this'            % \, separates the double and single quote.
        is what I just
        wrote, not  `that'\,''.

Dashes come in three sizes: an
       intra-word
dash, a medium dash for number ranges like
       1--2,
and a punctuation
       dash---like
this.

A sentence-ending space should be larger than the
space between words within a sentence.  You
sometimes have to type special commands in
conjunction with punctuation characters to get
this right, as in the following sentence.
       Gnats, gnus, etc.\ all  % `\ ' makes an inter-word space.
       begin with G\@.         % \@ marks end-of-sentence punctuation.
You should check the spaces after periods when
reading your output to make sure you haven't
forgotten any special cases.  Generating an
ellipsis
       \ldots\               % `\ ' is needed after `\ldots' because TeX
                             % ignores spaces after command names like \ldots
                             % made from \ + letters.
                             %
                             % Note how a `%' character causes TeX to ignore
                             % the end of the input line, so these blank lines
                             % do not start a new paragraph.
                             %
with the right spacing around the periods requires
a special command.

\LaTeX\ interprets some common characters as
commands, so you must type special commands to
generate them.  These characters include the
following:
       \$ \& \% \# \{ and \}.

In printing, text is usually emphasized with an
       \emph{italic}
type style.

\begin{em}
   A long segment of text can also be emphasized
   in this way.  Text within such a segment can be
   given \emph{additional} emphasis.
\end{em}

It is sometimes necessary to prevent \LaTeX\ from
breaking a line where it might otherwise do so.
This may be at a space, as between the ``Mr.''\ and
``Jones'' in
       ``Mr.~Jones'',        % ~ produces an unbreakable interword space.
or within a word---especially when the word is a
symbol like
       \mbox{\emph{itemnum}}
that makes little sense when hyphenated across
lines.

Footnotes\footnote{This is an example of a footnote.}
pose no problem.

\LaTeX\ is good at typesetting mathematical formulas
like
       \( x-3y + z = 7 \)
or
       \( a_{1} > x^{2n} + y^{2n} > x' \)
or
       \( \ip{A}{B} = \sum_{i} a_{i} b_{i} \).
The spaces you type in a formula are
ignored.  Remember that a letter like
       $x$                   % $ ... $  and  \( ... \)  are equivalent
is a formula when it denotes a mathematical
symbol, and it should be typed as one.

\section{Displayed Text}

Text is displayed by indenting it from the left
margin.  Quotations are commonly displayed.  There
are short quotations
\begin{quote}
   This is a short quotation.  It consists of a
   single paragraph of text.  See how it is formatted.
\end{quote}
and longer ones.
\begin{quotation}
   This is a longer quotation.  It consists of two
   paragraphs of text, neither of which are
   particularly interesting.

   This is the second paragraph of the quotation.  It
   is just as dull as the first paragraph.
\end{quotation}
Another frequently-displayed structure is a list.
The following is an example of an \emph{itemized}
list.
\begin{itemize}
   \item This is the first item of an itemized list.
         Each item in the list is marked with a ``tick''.
         You don't have to worry about what kind of tick
         mark is used.

   \item This is the second item of the list.  It
         contains another list nested inside it.  The inner
         list is an \emph{enumerated} list.
         \begin{enumerate}
            \item This is the first item of an enumerated
                  list that is nested within the itemized list.

            \item This is the second item of the inner list.
                  \LaTeX\ allows you to nest lists deeper than
                  you really should.
         \end{enumerate}
         This is the rest of the second item of the outer
         list.  It is no more interesting than any other
         part of the item.
   \item This is the third item of the list.
\end{itemize}
You can even display poetry.
\begin{verse}
   There is an environment
    for verse \\             % The \\ command separates lines
   Whose features some poets % within a stanza.
   will curse.

                             % One or more blank lines separate stanzas.

   For instead of making\\
   Them do \emph{all} line breaking, \\
   It allows them to put too many words on a line when they'd rather be
   forced to be terse.
\end{verse}

Mathematical formulas may also be displayed.  A
displayed formula
is
one-line long; multiline
formulas require special formatting instructions.
   \[  \ip{\Gamma}{\psi'} = x'' + y^{2} + z_{i}^{n}\]
Don't start a paragraph with a displayed equation,
nor make one a paragraph by itself.

\end{document}               % End of document.

\section{Bench Huge Section 954}
% synthetic bench section 954
% This is a sample LaTeX input file.  (Version of 12 August 2004.)
%
% A '%' character causes TeX to ignore all remaining text on the line,
% and is used for comments like this one.

\documentclass{article}      % Specifies the document class

                             % The preamble begins here.
\title{An Example Document}  % Declares the document's title.
\author{Leslie Lamport}      % Declares the author's name.
\date{January 21, 1994}      % Deleting this command produces today's date.

\newcommand{\ip}[2]{(#1, #2)}
                             % Defines \ip{arg1}{arg2} to mean
                             % (arg1, arg2).

%\newcommand{\ip}[2]{\langle #1 | #2\rangle}
                             % This is an alternative definition of
                             % \ip that is commented out.

\begin{document}             % End of preamble and beginning of text.

\maketitle                   % Produces the title.

This is an example input file.  Comparing it with
the output it generates can show you how to
produce a simple document of your own.

\section{Ordinary Text}      % Produces section heading.  Lower-level
                             % sections are begun with similar
                             % \subsection and \subsubsection commands.

The ends  of words and sentences are marked
  by   spaces. It  doesn't matter how many
spaces    you type; one is as good as 100.  The
end of   a line counts as a space.

One   or more   blank lines denote the  end
of  a paragraph.

Since any number of consecutive spaces are treated
like a single one, the formatting of the input
file makes no difference to
      \LaTeX,                % The \LaTeX command generates the LaTeX logo.
but it makes a difference to you.  When you use
\LaTeX, making your input file as easy to read
as possible will be a great help as you write
your document and when you change it.  This sample
file shows how you can add comments to your own input
file.

Because printing is different from typewriting,
there are a number of things that you have to do
differently when preparing an input file than if
you were just typing the document directly.
Quotation marks like
       ``this''
have to be handled specially, as do quotes within
quotes:
       ``\,`this'            % \, separates the double and single quote.
        is what I just
        wrote, not  `that'\,''.

Dashes come in three sizes: an
       intra-word
dash, a medium dash for number ranges like
       1--2,
and a punctuation
       dash---like
this.

A sentence-ending space should be larger than the
space between words within a sentence.  You
sometimes have to type special commands in
conjunction with punctuation characters to get
this right, as in the following sentence.
       Gnats, gnus, etc.\ all  % `\ ' makes an inter-word space.
       begin with G\@.         % \@ marks end-of-sentence punctuation.
You should check the spaces after periods when
reading your output to make sure you haven't
forgotten any special cases.  Generating an
ellipsis
       \ldots\               % `\ ' is needed after `\ldots' because TeX
                             % ignores spaces after command names like \ldots
                             % made from \ + letters.
                             %
                             % Note how a `%' character causes TeX to ignore
                             % the end of the input line, so these blank lines
                             % do not start a new paragraph.
                             %
with the right spacing around the periods requires
a special command.

\LaTeX\ interprets some common characters as
commands, so you must type special commands to
generate them.  These characters include the
following:
       \$ \& \% \# \{ and \}.

In printing, text is usually emphasized with an
       \emph{italic}
type style.

\begin{em}
   A long segment of text can also be emphasized
   in this way.  Text within such a segment can be
   given \emph{additional} emphasis.
\end{em}

It is sometimes necessary to prevent \LaTeX\ from
breaking a line where it might otherwise do so.
This may be at a space, as between the ``Mr.''\ and
``Jones'' in
       ``Mr.~Jones'',        % ~ produces an unbreakable interword space.
or within a word---especially when the word is a
symbol like
       \mbox{\emph{itemnum}}
that makes little sense when hyphenated across
lines.

Footnotes\footnote{This is an example of a footnote.}
pose no problem.

\LaTeX\ is good at typesetting mathematical formulas
like
       \( x-3y + z = 7 \)
or
       \( a_{1} > x^{2n} + y^{2n} > x' \)
or
       \( \ip{A}{B} = \sum_{i} a_{i} b_{i} \).
The spaces you type in a formula are
ignored.  Remember that a letter like
       $x$                   % $ ... $  and  \( ... \)  are equivalent
is a formula when it denotes a mathematical
symbol, and it should be typed as one.

\section{Displayed Text}

Text is displayed by indenting it from the left
margin.  Quotations are commonly displayed.  There
are short quotations
\begin{quote}
   This is a short quotation.  It consists of a
   single paragraph of text.  See how it is formatted.
\end{quote}
and longer ones.
\begin{quotation}
   This is a longer quotation.  It consists of two
   paragraphs of text, neither of which are
   particularly interesting.

   This is the second paragraph of the quotation.  It
   is just as dull as the first paragraph.
\end{quotation}
Another frequently-displayed structure is a list.
The following is an example of an \emph{itemized}
list.
\begin{itemize}
   \item This is the first item of an itemized list.
         Each item in the list is marked with a ``tick''.
         You don't have to worry about what kind of tick
         mark is used.

   \item This is the second item of the list.  It
         contains another list nested inside it.  The inner
         list is an \emph{enumerated} list.
         \begin{enumerate}
            \item This is the first item of an enumerated
                  list that is nested within the itemized list.

            \item This is the second item of the inner list.
                  \LaTeX\ allows you to nest lists deeper than
                  you really should.
         \end{enumerate}
         This is the rest of the second item of the outer
         list.  It is no more interesting than any other
         part of the item.
   \item This is the third item of the list.
\end{itemize}
You can even display poetry.
\begin{verse}
   There is an environment
    for verse \\             % The \\ command separates lines
   Whose features some poets % within a stanza.
   will curse.

                             % One or more blank lines separate stanzas.

   For instead of making\\
   Them do \emph{all} line breaking, \\
   It allows them to put too many words on a line when they'd rather be
   forced to be terse.
\end{verse}

Mathematical formulas may also be displayed.  A
displayed formula
is
one-line long; multiline
formulas require special formatting instructions.
   \[  \ip{\Gamma}{\psi'} = x'' + y^{2} + z_{i}^{n}\]
Don't start a paragraph with a displayed equation,
nor make one a paragraph by itself.

\end{document}               % End of document.

\section{Bench Huge Section 955}
% synthetic bench section 955
% This is a sample LaTeX input file.  (Version of 12 August 2004.)
%
% A '%' character causes TeX to ignore all remaining text on the line,
% and is used for comments like this one.

\documentclass{article}      % Specifies the document class

                             % The preamble begins here.
\title{An Example Document}  % Declares the document's title.
\author{Leslie Lamport}      % Declares the author's name.
\date{January 21, 1994}      % Deleting this command produces today's date.

\newcommand{\ip}[2]{(#1, #2)}
                             % Defines \ip{arg1}{arg2} to mean
                             % (arg1, arg2).

%\newcommand{\ip}[2]{\langle #1 | #2\rangle}
                             % This is an alternative definition of
                             % \ip that is commented out.

\begin{document}             % End of preamble and beginning of text.

\maketitle                   % Produces the title.

This is an example input file.  Comparing it with
the output it generates can show you how to
produce a simple document of your own.

\section{Ordinary Text}      % Produces section heading.  Lower-level
                             % sections are begun with similar
                             % \subsection and \subsubsection commands.

The ends  of words and sentences are marked
  by   spaces. It  doesn't matter how many
spaces    you type; one is as good as 100.  The
end of   a line counts as a space.

One   or more   blank lines denote the  end
of  a paragraph.

Since any number of consecutive spaces are treated
like a single one, the formatting of the input
file makes no difference to
      \LaTeX,                % The \LaTeX command generates the LaTeX logo.
but it makes a difference to you.  When you use
\LaTeX, making your input file as easy to read
as possible will be a great help as you write
your document and when you change it.  This sample
file shows how you can add comments to your own input
file.

Because printing is different from typewriting,
there are a number of things that you have to do
differently when preparing an input file than if
you were just typing the document directly.
Quotation marks like
       ``this''
have to be handled specially, as do quotes within
quotes:
       ``\,`this'            % \, separates the double and single quote.
        is what I just
        wrote, not  `that'\,''.

Dashes come in three sizes: an
       intra-word
dash, a medium dash for number ranges like
       1--2,
and a punctuation
       dash---like
this.

A sentence-ending space should be larger than the
space between words within a sentence.  You
sometimes have to type special commands in
conjunction with punctuation characters to get
this right, as in the following sentence.
       Gnats, gnus, etc.\ all  % `\ ' makes an inter-word space.
       begin with G\@.         % \@ marks end-of-sentence punctuation.
You should check the spaces after periods when
reading your output to make sure you haven't
forgotten any special cases.  Generating an
ellipsis
       \ldots\               % `\ ' is needed after `\ldots' because TeX
                             % ignores spaces after command names like \ldots
                             % made from \ + letters.
                             %
                             % Note how a `%' character causes TeX to ignore
                             % the end of the input line, so these blank lines
                             % do not start a new paragraph.
                             %
with the right spacing around the periods requires
a special command.

\LaTeX\ interprets some common characters as
commands, so you must type special commands to
generate them.  These characters include the
following:
       \$ \& \% \# \{ and \}.

In printing, text is usually emphasized with an
       \emph{italic}
type style.

\begin{em}
   A long segment of text can also be emphasized
   in this way.  Text within such a segment can be
   given \emph{additional} emphasis.
\end{em}

It is sometimes necessary to prevent \LaTeX\ from
breaking a line where it might otherwise do so.
This may be at a space, as between the ``Mr.''\ and
``Jones'' in
       ``Mr.~Jones'',        % ~ produces an unbreakable interword space.
or within a word---especially when the word is a
symbol like
       \mbox{\emph{itemnum}}
that makes little sense when hyphenated across
lines.

Footnotes\footnote{This is an example of a footnote.}
pose no problem.

\LaTeX\ is good at typesetting mathematical formulas
like
       \( x-3y + z = 7 \)
or
       \( a_{1} > x^{2n} + y^{2n} > x' \)
or
       \( \ip{A}{B} = \sum_{i} a_{i} b_{i} \).
The spaces you type in a formula are
ignored.  Remember that a letter like
       $x$                   % $ ... $  and  \( ... \)  are equivalent
is a formula when it denotes a mathematical
symbol, and it should be typed as one.

\section{Displayed Text}

Text is displayed by indenting it from the left
margin.  Quotations are commonly displayed.  There
are short quotations
\begin{quote}
   This is a short quotation.  It consists of a
   single paragraph of text.  See how it is formatted.
\end{quote}
and longer ones.
\begin{quotation}
   This is a longer quotation.  It consists of two
   paragraphs of text, neither of which are
   particularly interesting.

   This is the second paragraph of the quotation.  It
   is just as dull as the first paragraph.
\end{quotation}
Another frequently-displayed structure is a list.
The following is an example of an \emph{itemized}
list.
\begin{itemize}
   \item This is the first item of an itemized list.
         Each item in the list is marked with a ``tick''.
         You don't have to worry about what kind of tick
         mark is used.

   \item This is the second item of the list.  It
         contains another list nested inside it.  The inner
         list is an \emph{enumerated} list.
         \begin{enumerate}
            \item This is the first item of an enumerated
                  list that is nested within the itemized list.

            \item This is the second item of the inner list.
                  \LaTeX\ allows you to nest lists deeper than
                  you really should.
         \end{enumerate}
         This is the rest of the second item of the outer
         list.  It is no more interesting than any other
         part of the item.
   \item This is the third item of the list.
\end{itemize}
You can even display poetry.
\begin{verse}
   There is an environment
    for verse \\             % The \\ command separates lines
   Whose features some poets % within a stanza.
   will curse.

                             % One or more blank lines separate stanzas.

   For instead of making\\
   Them do \emph{all} line breaking, \\
   It allows them to put too many words on a line when they'd rather be
   forced to be terse.
\end{verse}

Mathematical formulas may also be displayed.  A
displayed formula
is
one-line long; multiline
formulas require special formatting instructions.
   \[  \ip{\Gamma}{\psi'} = x'' + y^{2} + z_{i}^{n}\]
Don't start a paragraph with a displayed equation,
nor make one a paragraph by itself.

\end{document}               % End of document.

\section{Bench Huge Section 956}
% synthetic bench section 956
% This is a sample LaTeX input file.  (Version of 12 August 2004.)
%
% A '%' character causes TeX to ignore all remaining text on the line,
% and is used for comments like this one.

\documentclass{article}      % Specifies the document class

                             % The preamble begins here.
\title{An Example Document}  % Declares the document's title.
\author{Leslie Lamport}      % Declares the author's name.
\date{January 21, 1994}      % Deleting this command produces today's date.

\newcommand{\ip}[2]{(#1, #2)}
                             % Defines \ip{arg1}{arg2} to mean
                             % (arg1, arg2).

%\newcommand{\ip}[2]{\langle #1 | #2\rangle}
                             % This is an alternative definition of
                             % \ip that is commented out.

\begin{document}             % End of preamble and beginning of text.

\maketitle                   % Produces the title.

This is an example input file.  Comparing it with
the output it generates can show you how to
produce a simple document of your own.

\section{Ordinary Text}      % Produces section heading.  Lower-level
                             % sections are begun with similar
                             % \subsection and \subsubsection commands.

The ends  of words and sentences are marked
  by   spaces. It  doesn't matter how many
spaces    you type; one is as good as 100.  The
end of   a line counts as a space.

One   or more   blank lines denote the  end
of  a paragraph.

Since any number of consecutive spaces are treated
like a single one, the formatting of the input
file makes no difference to
      \LaTeX,                % The \LaTeX command generates the LaTeX logo.
but it makes a difference to you.  When you use
\LaTeX, making your input file as easy to read
as possible will be a great help as you write
your document and when you change it.  This sample
file shows how you can add comments to your own input
file.

Because printing is different from typewriting,
there are a number of things that you have to do
differently when preparing an input file than if
you were just typing the document directly.
Quotation marks like
       ``this''
have to be handled specially, as do quotes within
quotes:
       ``\,`this'            % \, separates the double and single quote.
        is what I just
        wrote, not  `that'\,''.

Dashes come in three sizes: an
       intra-word
dash, a medium dash for number ranges like
       1--2,
and a punctuation
       dash---like
this.

A sentence-ending space should be larger than the
space between words within a sentence.  You
sometimes have to type special commands in
conjunction with punctuation characters to get
this right, as in the following sentence.
       Gnats, gnus, etc.\ all  % `\ ' makes an inter-word space.
       begin with G\@.         % \@ marks end-of-sentence punctuation.
You should check the spaces after periods when
reading your output to make sure you haven't
forgotten any special cases.  Generating an
ellipsis
       \ldots\               % `\ ' is needed after `\ldots' because TeX
                             % ignores spaces after command names like \ldots
                             % made from \ + letters.
                             %
                             % Note how a `%' character causes TeX to ignore
                             % the end of the input line, so these blank lines
                             % do not start a new paragraph.
                             %
with the right spacing around the periods requires
a special command.

\LaTeX\ interprets some common characters as
commands, so you must type special commands to
generate them.  These characters include the
following:
       \$ \& \% \# \{ and \}.

In printing, text is usually emphasized with an
       \emph{italic}
type style.

\begin{em}
   A long segment of text can also be emphasized
   in this way.  Text within such a segment can be
   given \emph{additional} emphasis.
\end{em}

It is sometimes necessary to prevent \LaTeX\ from
breaking a line where it might otherwise do so.
This may be at a space, as between the ``Mr.''\ and
``Jones'' in
       ``Mr.~Jones'',        % ~ produces an unbreakable interword space.
or within a word---especially when the word is a
symbol like
       \mbox{\emph{itemnum}}
that makes little sense when hyphenated across
lines.

Footnotes\footnote{This is an example of a footnote.}
pose no problem.

\LaTeX\ is good at typesetting mathematical formulas
like
       \( x-3y + z = 7 \)
or
       \( a_{1} > x^{2n} + y^{2n} > x' \)
or
       \( \ip{A}{B} = \sum_{i} a_{i} b_{i} \).
The spaces you type in a formula are
ignored.  Remember that a letter like
       $x$                   % $ ... $  and  \( ... \)  are equivalent
is a formula when it denotes a mathematical
symbol, and it should be typed as one.

\section{Displayed Text}

Text is displayed by indenting it from the left
margin.  Quotations are commonly displayed.  There
are short quotations
\begin{quote}
   This is a short quotation.  It consists of a
   single paragraph of text.  See how it is formatted.
\end{quote}
and longer ones.
\begin{quotation}
   This is a longer quotation.  It consists of two
   paragraphs of text, neither of which are
   particularly interesting.

   This is the second paragraph of the quotation.  It
   is just as dull as the first paragraph.
\end{quotation}
Another frequently-displayed structure is a list.
The following is an example of an \emph{itemized}
list.
\begin{itemize}
   \item This is the first item of an itemized list.
         Each item in the list is marked with a ``tick''.
         You don't have to worry about what kind of tick
         mark is used.

   \item This is the second item of the list.  It
         contains another list nested inside it.  The inner
         list is an \emph{enumerated} list.
         \begin{enumerate}
            \item This is the first item of an enumerated
                  list that is nested within the itemized list.

            \item This is the second item of the inner list.
                  \LaTeX\ allows you to nest lists deeper than
                  you really should.
         \end{enumerate}
         This is the rest of the second item of the outer
         list.  It is no more interesting than any other
         part of the item.
   \item This is the third item of the list.
\end{itemize}
You can even display poetry.
\begin{verse}
   There is an environment
    for verse \\             % The \\ command separates lines
   Whose features some poets % within a stanza.
   will curse.

                             % One or more blank lines separate stanzas.

   For instead of making\\
   Them do \emph{all} line breaking, \\
   It allows them to put too many words on a line when they'd rather be
   forced to be terse.
\end{verse}

Mathematical formulas may also be displayed.  A
displayed formula
is
one-line long; multiline
formulas require special formatting instructions.
   \[  \ip{\Gamma}{\psi'} = x'' + y^{2} + z_{i}^{n}\]
Don't start a paragraph with a displayed equation,
nor make one a paragraph by itself.

\end{document}               % End of document.

\section{Bench Huge Section 957}
% synthetic bench section 957
% This is a sample LaTeX input file.  (Version of 12 August 2004.)
%
% A '%' character causes TeX to ignore all remaining text on the line,
% and is used for comments like this one.

\documentclass{article}      % Specifies the document class

                             % The preamble begins here.
\title{An Example Document}  % Declares the document's title.
\author{Leslie Lamport}      % Declares the author's name.
\date{January 21, 1994}      % Deleting this command produces today's date.

\newcommand{\ip}[2]{(#1, #2)}
                             % Defines \ip{arg1}{arg2} to mean
                             % (arg1, arg2).

%\newcommand{\ip}[2]{\langle #1 | #2\rangle}
                             % This is an alternative definition of
                             % \ip that is commented out.

\begin{document}             % End of preamble and beginning of text.

\maketitle                   % Produces the title.

This is an example input file.  Comparing it with
the output it generates can show you how to
produce a simple document of your own.

\section{Ordinary Text}      % Produces section heading.  Lower-level
                             % sections are begun with similar
                             % \subsection and \subsubsection commands.

The ends  of words and sentences are marked
  by   spaces. It  doesn't matter how many
spaces    you type; one is as good as 100.  The
end of   a line counts as a space.

One   or more   blank lines denote the  end
of  a paragraph.

Since any number of consecutive spaces are treated
like a single one, the formatting of the input
file makes no difference to
      \LaTeX,                % The \LaTeX command generates the LaTeX logo.
but it makes a difference to you.  When you use
\LaTeX, making your input file as easy to read
as possible will be a great help as you write
your document and when you change it.  This sample
file shows how you can add comments to your own input
file.

Because printing is different from typewriting,
there are a number of things that you have to do
differently when preparing an input file than if
you were just typing the document directly.
Quotation marks like
       ``this''
have to be handled specially, as do quotes within
quotes:
       ``\,`this'            % \, separates the double and single quote.
        is what I just
        wrote, not  `that'\,''.

Dashes come in three sizes: an
       intra-word
dash, a medium dash for number ranges like
       1--2,
and a punctuation
       dash---like
this.

A sentence-ending space should be larger than the
space between words within a sentence.  You
sometimes have to type special commands in
conjunction with punctuation characters to get
this right, as in the following sentence.
       Gnats, gnus, etc.\ all  % `\ ' makes an inter-word space.
       begin with G\@.         % \@ marks end-of-sentence punctuation.
You should check the spaces after periods when
reading your output to make sure you haven't
forgotten any special cases.  Generating an
ellipsis
       \ldots\               % `\ ' is needed after `\ldots' because TeX
                             % ignores spaces after command names like \ldots
                             % made from \ + letters.
                             %
                             % Note how a `%' character causes TeX to ignore
                             % the end of the input line, so these blank lines
                             % do not start a new paragraph.
                             %
with the right spacing around the periods requires
a special command.

\LaTeX\ interprets some common characters as
commands, so you must type special commands to
generate them.  These characters include the
following:
       \$ \& \% \# \{ and \}.

In printing, text is usually emphasized with an
       \emph{italic}
type style.

\begin{em}
   A long segment of text can also be emphasized
   in this way.  Text within such a segment can be
   given \emph{additional} emphasis.
\end{em}

It is sometimes necessary to prevent \LaTeX\ from
breaking a line where it might otherwise do so.
This may be at a space, as between the ``Mr.''\ and
``Jones'' in
       ``Mr.~Jones'',        % ~ produces an unbreakable interword space.
or within a word---especially when the word is a
symbol like
       \mbox{\emph{itemnum}}
that makes little sense when hyphenated across
lines.

Footnotes\footnote{This is an example of a footnote.}
pose no problem.

\LaTeX\ is good at typesetting mathematical formulas
like
       \( x-3y + z = 7 \)
or
       \( a_{1} > x^{2n} + y^{2n} > x' \)
or
       \( \ip{A}{B} = \sum_{i} a_{i} b_{i} \).
The spaces you type in a formula are
ignored.  Remember that a letter like
       $x$                   % $ ... $  and  \( ... \)  are equivalent
is a formula when it denotes a mathematical
symbol, and it should be typed as one.

\section{Displayed Text}

Text is displayed by indenting it from the left
margin.  Quotations are commonly displayed.  There
are short quotations
\begin{quote}
   This is a short quotation.  It consists of a
   single paragraph of text.  See how it is formatted.
\end{quote}
and longer ones.
\begin{quotation}
   This is a longer quotation.  It consists of two
   paragraphs of text, neither of which are
   particularly interesting.

   This is the second paragraph of the quotation.  It
   is just as dull as the first paragraph.
\end{quotation}
Another frequently-displayed structure is a list.
The following is an example of an \emph{itemized}
list.
\begin{itemize}
   \item This is the first item of an itemized list.
         Each item in the list is marked with a ``tick''.
         You don't have to worry about what kind of tick
         mark is used.

   \item This is the second item of the list.  It
         contains another list nested inside it.  The inner
         list is an \emph{enumerated} list.
         \begin{enumerate}
            \item This is the first item of an enumerated
                  list that is nested within the itemized list.

            \item This is the second item of the inner list.
                  \LaTeX\ allows you to nest lists deeper than
                  you really should.
         \end{enumerate}
         This is the rest of the second item of the outer
         list.  It is no more interesting than any other
         part of the item.
   \item This is the third item of the list.
\end{itemize}
You can even display poetry.
\begin{verse}
   There is an environment
    for verse \\             % The \\ command separates lines
   Whose features some poets % within a stanza.
   will curse.

                             % One or more blank lines separate stanzas.

   For instead of making\\
   Them do \emph{all} line breaking, \\
   It allows them to put too many words on a line when they'd rather be
   forced to be terse.
\end{verse}

Mathematical formulas may also be displayed.  A
displayed formula
is
one-line long; multiline
formulas require special formatting instructions.
   \[  \ip{\Gamma}{\psi'} = x'' + y^{2} + z_{i}^{n}\]
Don't start a paragraph with a displayed equation,
nor make one a paragraph by itself.

\end{document}               % End of document.

\section{Bench Huge Section 958}
% synthetic bench section 958
% This is a sample LaTeX input file.  (Version of 12 August 2004.)
%
% A '%' character causes TeX to ignore all remaining text on the line,
% and is used for comments like this one.

\documentclass{article}      % Specifies the document class

                             % The preamble begins here.
\title{An Example Document}  % Declares the document's title.
\author{Leslie Lamport}      % Declares the author's name.
\date{January 21, 1994}      % Deleting this command produces today's date.

\newcommand{\ip}[2]{(#1, #2)}
                             % Defines \ip{arg1}{arg2} to mean
                             % (arg1, arg2).

%\newcommand{\ip}[2]{\langle #1 | #2\rangle}
                             % This is an alternative definition of
                             % \ip that is commented out.

\begin{document}             % End of preamble and beginning of text.

\maketitle                   % Produces the title.

This is an example input file.  Comparing it with
the output it generates can show you how to
produce a simple document of your own.

\section{Ordinary Text}      % Produces section heading.  Lower-level
                             % sections are begun with similar
                             % \subsection and \subsubsection commands.

The ends  of words and sentences are marked
  by   spaces. It  doesn't matter how many
spaces    you type; one is as good as 100.  The
end of   a line counts as a space.

One   or more   blank lines denote the  end
of  a paragraph.

Since any number of consecutive spaces are treated
like a single one, the formatting of the input
file makes no difference to
      \LaTeX,                % The \LaTeX command generates the LaTeX logo.
but it makes a difference to you.  When you use
\LaTeX, making your input file as easy to read
as possible will be a great help as you write
your document and when you change it.  This sample
file shows how you can add comments to your own input
file.

Because printing is different from typewriting,
there are a number of things that you have to do
differently when preparing an input file than if
you were just typing the document directly.
Quotation marks like
       ``this''
have to be handled specially, as do quotes within
quotes:
       ``\,`this'            % \, separates the double and single quote.
        is what I just
        wrote, not  `that'\,''.

Dashes come in three sizes: an
       intra-word
dash, a medium dash for number ranges like
       1--2,
and a punctuation
       dash---like
this.

A sentence-ending space should be larger than the
space between words within a sentence.  You
sometimes have to type special commands in
conjunction with punctuation characters to get
this right, as in the following sentence.
       Gnats, gnus, etc.\ all  % `\ ' makes an inter-word space.
       begin with G\@.         % \@ marks end-of-sentence punctuation.
You should check the spaces after periods when
reading your output to make sure you haven't
forgotten any special cases.  Generating an
ellipsis
       \ldots\               % `\ ' is needed after `\ldots' because TeX
                             % ignores spaces after command names like \ldots
                             % made from \ + letters.
                             %
                             % Note how a `%' character causes TeX to ignore
                             % the end of the input line, so these blank lines
                             % do not start a new paragraph.
                             %
with the right spacing around the periods requires
a special command.

\LaTeX\ interprets some common characters as
commands, so you must type special commands to
generate them.  These characters include the
following:
       \$ \& \% \# \{ and \}.

In printing, text is usually emphasized with an
       \emph{italic}
type style.

\begin{em}
   A long segment of text can also be emphasized
   in this way.  Text within such a segment can be
   given \emph{additional} emphasis.
\end{em}

It is sometimes necessary to prevent \LaTeX\ from
breaking a line where it might otherwise do so.
This may be at a space, as between the ``Mr.''\ and
``Jones'' in
       ``Mr.~Jones'',        % ~ produces an unbreakable interword space.
or within a word---especially when the word is a
symbol like
       \mbox{\emph{itemnum}}
that makes little sense when hyphenated across
lines.

Footnotes\footnote{This is an example of a footnote.}
pose no problem.

\LaTeX\ is good at typesetting mathematical formulas
like
       \( x-3y + z = 7 \)
or
       \( a_{1} > x^{2n} + y^{2n} > x' \)
or
       \( \ip{A}{B} = \sum_{i} a_{i} b_{i} \).
The spaces you type in a formula are
ignored.  Remember that a letter like
       $x$                   % $ ... $  and  \( ... \)  are equivalent
is a formula when it denotes a mathematical
symbol, and it should be typed as one.

\section{Displayed Text}

Text is displayed by indenting it from the left
margin.  Quotations are commonly displayed.  There
are short quotations
\begin{quote}
   This is a short quotation.  It consists of a
   single paragraph of text.  See how it is formatted.
\end{quote}
and longer ones.
\begin{quotation}
   This is a longer quotation.  It consists of two
   paragraphs of text, neither of which are
   particularly interesting.

   This is the second paragraph of the quotation.  It
   is just as dull as the first paragraph.
\end{quotation}
Another frequently-displayed structure is a list.
The following is an example of an \emph{itemized}
list.
\begin{itemize}
   \item This is the first item of an itemized list.
         Each item in the list is marked with a ``tick''.
         You don't have to worry about what kind of tick
         mark is used.

   \item This is the second item of the list.  It
         contains another list nested inside it.  The inner
         list is an \emph{enumerated} list.
         \begin{enumerate}
            \item This is the first item of an enumerated
                  list that is nested within the itemized list.

            \item This is the second item of the inner list.
                  \LaTeX\ allows you to nest lists deeper than
                  you really should.
         \end{enumerate}
         This is the rest of the second item of the outer
         list.  It is no more interesting than any other
         part of the item.
   \item This is the third item of the list.
\end{itemize}
You can even display poetry.
\begin{verse}
   There is an environment
    for verse \\             % The \\ command separates lines
   Whose features some poets % within a stanza.
   will curse.

                             % One or more blank lines separate stanzas.

   For instead of making\\
   Them do \emph{all} line breaking, \\
   It allows them to put too many words on a line when they'd rather be
   forced to be terse.
\end{verse}

Mathematical formulas may also be displayed.  A
displayed formula
is
one-line long; multiline
formulas require special formatting instructions.
   \[  \ip{\Gamma}{\psi'} = x'' + y^{2} + z_{i}^{n}\]
Don't start a paragraph with a displayed equation,
nor make one a paragraph by itself.

\end{document}               % End of document.

\section{Bench Huge Section 959}
% synthetic bench section 959
% This is a sample LaTeX input file.  (Version of 12 August 2004.)
%
% A '%' character causes TeX to ignore all remaining text on the line,
% and is used for comments like this one.

\documentclass{article}      % Specifies the document class

                             % The preamble begins here.
\title{An Example Document}  % Declares the document's title.
\author{Leslie Lamport}      % Declares the author's name.
\date{January 21, 1994}      % Deleting this command produces today's date.

\newcommand{\ip}[2]{(#1, #2)}
                             % Defines \ip{arg1}{arg2} to mean
                             % (arg1, arg2).

%\newcommand{\ip}[2]{\langle #1 | #2\rangle}
                             % This is an alternative definition of
                             % \ip that is commented out.

\begin{document}             % End of preamble and beginning of text.

\maketitle                   % Produces the title.

This is an example input file.  Comparing it with
the output it generates can show you how to
produce a simple document of your own.

\section{Ordinary Text}      % Produces section heading.  Lower-level
                             % sections are begun with similar
                             % \subsection and \subsubsection commands.

The ends  of words and sentences are marked
  by   spaces. It  doesn't matter how many
spaces    you type; one is as good as 100.  The
end of   a line counts as a space.

One   or more   blank lines denote the  end
of  a paragraph.

Since any number of consecutive spaces are treated
like a single one, the formatting of the input
file makes no difference to
      \LaTeX,                % The \LaTeX command generates the LaTeX logo.
but it makes a difference to you.  When you use
\LaTeX, making your input file as easy to read
as possible will be a great help as you write
your document and when you change it.  This sample
file shows how you can add comments to your own input
file.

Because printing is different from typewriting,
there are a number of things that you have to do
differently when preparing an input file than if
you were just typing the document directly.
Quotation marks like
       ``this''
have to be handled specially, as do quotes within
quotes:
       ``\,`this'            % \, separates the double and single quote.
        is what I just
        wrote, not  `that'\,''.

Dashes come in three sizes: an
       intra-word
dash, a medium dash for number ranges like
       1--2,
and a punctuation
       dash---like
this.

A sentence-ending space should be larger than the
space between words within a sentence.  You
sometimes have to type special commands in
conjunction with punctuation characters to get
this right, as in the following sentence.
       Gnats, gnus, etc.\ all  % `\ ' makes an inter-word space.
       begin with G\@.         % \@ marks end-of-sentence punctuation.
You should check the spaces after periods when
reading your output to make sure you haven't
forgotten any special cases.  Generating an
ellipsis
       \ldots\               % `\ ' is needed after `\ldots' because TeX
                             % ignores spaces after command names like \ldots
                             % made from \ + letters.
                             %
                             % Note how a `%' character causes TeX to ignore
                             % the end of the input line, so these blank lines
                             % do not start a new paragraph.
                             %
with the right spacing around the periods requires
a special command.

\LaTeX\ interprets some common characters as
commands, so you must type special commands to
generate them.  These characters include the
following:
       \$ \& \% \# \{ and \}.

In printing, text is usually emphasized with an
       \emph{italic}
type style.

\begin{em}
   A long segment of text can also be emphasized
   in this way.  Text within such a segment can be
   given \emph{additional} emphasis.
\end{em}

It is sometimes necessary to prevent \LaTeX\ from
breaking a line where it might otherwise do so.
This may be at a space, as between the ``Mr.''\ and
``Jones'' in
       ``Mr.~Jones'',        % ~ produces an unbreakable interword space.
or within a word---especially when the word is a
symbol like
       \mbox{\emph{itemnum}}
that makes little sense when hyphenated across
lines.

Footnotes\footnote{This is an example of a footnote.}
pose no problem.

\LaTeX\ is good at typesetting mathematical formulas
like
       \( x-3y + z = 7 \)
or
       \( a_{1} > x^{2n} + y^{2n} > x' \)
or
       \( \ip{A}{B} = \sum_{i} a_{i} b_{i} \).
The spaces you type in a formula are
ignored.  Remember that a letter like
       $x$                   % $ ... $  and  \( ... \)  are equivalent
is a formula when it denotes a mathematical
symbol, and it should be typed as one.

\section{Displayed Text}

Text is displayed by indenting it from the left
margin.  Quotations are commonly displayed.  There
are short quotations
\begin{quote}
   This is a short quotation.  It consists of a
   single paragraph of text.  See how it is formatted.
\end{quote}
and longer ones.
\begin{quotation}
   This is a longer quotation.  It consists of two
   paragraphs of text, neither of which are
   particularly interesting.

   This is the second paragraph of the quotation.  It
   is just as dull as the first paragraph.
\end{quotation}
Another frequently-displayed structure is a list.
The following is an example of an \emph{itemized}
list.
\begin{itemize}
   \item This is the first item of an itemized list.
         Each item in the list is marked with a ``tick''.
         You don't have to worry about what kind of tick
         mark is used.

   \item This is the second item of the list.  It
         contains another list nested inside it.  The inner
         list is an \emph{enumerated} list.
         \begin{enumerate}
            \item This is the first item of an enumerated
                  list that is nested within the itemized list.

            \item This is the second item of the inner list.
                  \LaTeX\ allows you to nest lists deeper than
                  you really should.
         \end{enumerate}
         This is the rest of the second item of the outer
         list.  It is no more interesting than any other
         part of the item.
   \item This is the third item of the list.
\end{itemize}
You can even display poetry.
\begin{verse}
   There is an environment
    for verse \\             % The \\ command separates lines
   Whose features some poets % within a stanza.
   will curse.

                             % One or more blank lines separate stanzas.

   For instead of making\\
   Them do \emph{all} line breaking, \\
   It allows them to put too many words on a line when they'd rather be
   forced to be terse.
\end{verse}

Mathematical formulas may also be displayed.  A
displayed formula
is
one-line long; multiline
formulas require special formatting instructions.
   \[  \ip{\Gamma}{\psi'} = x'' + y^{2} + z_{i}^{n}\]
Don't start a paragraph with a displayed equation,
nor make one a paragraph by itself.

\end{document}               % End of document.

\section{Bench Huge Section 960}
% synthetic bench section 960
% This is a sample LaTeX input file.  (Version of 12 August 2004.)
%
% A '%' character causes TeX to ignore all remaining text on the line,
% and is used for comments like this one.

\documentclass{article}      % Specifies the document class

                             % The preamble begins here.
\title{An Example Document}  % Declares the document's title.
\author{Leslie Lamport}      % Declares the author's name.
\date{January 21, 1994}      % Deleting this command produces today's date.

\newcommand{\ip}[2]{(#1, #2)}
                             % Defines \ip{arg1}{arg2} to mean
                             % (arg1, arg2).

%\newcommand{\ip}[2]{\langle #1 | #2\rangle}
                             % This is an alternative definition of
                             % \ip that is commented out.

\begin{document}             % End of preamble and beginning of text.

\maketitle                   % Produces the title.

This is an example input file.  Comparing it with
the output it generates can show you how to
produce a simple document of your own.

\section{Ordinary Text}      % Produces section heading.  Lower-level
                             % sections are begun with similar
                             % \subsection and \subsubsection commands.

The ends  of words and sentences are marked
  by   spaces. It  doesn't matter how many
spaces    you type; one is as good as 100.  The
end of   a line counts as a space.

One   or more   blank lines denote the  end
of  a paragraph.

Since any number of consecutive spaces are treated
like a single one, the formatting of the input
file makes no difference to
      \LaTeX,                % The \LaTeX command generates the LaTeX logo.
but it makes a difference to you.  When you use
\LaTeX, making your input file as easy to read
as possible will be a great help as you write
your document and when you change it.  This sample
file shows how you can add comments to your own input
file.

Because printing is different from typewriting,
there are a number of things that you have to do
differently when preparing an input file than if
you were just typing the document directly.
Quotation marks like
       ``this''
have to be handled specially, as do quotes within
quotes:
       ``\,`this'            % \, separates the double and single quote.
        is what I just
        wrote, not  `that'\,''.

Dashes come in three sizes: an
       intra-word
dash, a medium dash for number ranges like
       1--2,
and a punctuation
       dash---like
this.

A sentence-ending space should be larger than the
space between words within a sentence.  You
sometimes have to type special commands in
conjunction with punctuation characters to get
this right, as in the following sentence.
       Gnats, gnus, etc.\ all  % `\ ' makes an inter-word space.
       begin with G\@.         % \@ marks end-of-sentence punctuation.
You should check the spaces after periods when
reading your output to make sure you haven't
forgotten any special cases.  Generating an
ellipsis
       \ldots\               % `\ ' is needed after `\ldots' because TeX
                             % ignores spaces after command names like \ldots
                             % made from \ + letters.
                             %
                             % Note how a `%' character causes TeX to ignore
                             % the end of the input line, so these blank lines
                             % do not start a new paragraph.
                             %
with the right spacing around the periods requires
a special command.

\LaTeX\ interprets some common characters as
commands, so you must type special commands to
generate them.  These characters include the
following:
       \$ \& \% \# \{ and \}.

In printing, text is usually emphasized with an
       \emph{italic}
type style.

\begin{em}
   A long segment of text can also be emphasized
   in this way.  Text within such a segment can be
   given \emph{additional} emphasis.
\end{em}

It is sometimes necessary to prevent \LaTeX\ from
breaking a line where it might otherwise do so.
This may be at a space, as between the ``Mr.''\ and
``Jones'' in
       ``Mr.~Jones'',        % ~ produces an unbreakable interword space.
or within a word---especially when the word is a
symbol like
       \mbox{\emph{itemnum}}
that makes little sense when hyphenated across
lines.

Footnotes\footnote{This is an example of a footnote.}
pose no problem.

\LaTeX\ is good at typesetting mathematical formulas
like
       \( x-3y + z = 7 \)
or
       \( a_{1} > x^{2n} + y^{2n} > x' \)
or
       \( \ip{A}{B} = \sum_{i} a_{i} b_{i} \).
The spaces you type in a formula are
ignored.  Remember that a letter like
       $x$                   % $ ... $  and  \( ... \)  are equivalent
is a formula when it denotes a mathematical
symbol, and it should be typed as one.

\section{Displayed Text}

Text is displayed by indenting it from the left
margin.  Quotations are commonly displayed.  There
are short quotations
\begin{quote}
   This is a short quotation.  It consists of a
   single paragraph of text.  See how it is formatted.
\end{quote}
and longer ones.
\begin{quotation}
   This is a longer quotation.  It consists of two
   paragraphs of text, neither of which are
   particularly interesting.

   This is the second paragraph of the quotation.  It
   is just as dull as the first paragraph.
\end{quotation}
Another frequently-displayed structure is a list.
The following is an example of an \emph{itemized}
list.
\begin{itemize}
   \item This is the first item of an itemized list.
         Each item in the list is marked with a ``tick''.
         You don't have to worry about what kind of tick
         mark is used.

   \item This is the second item of the list.  It
         contains another list nested inside it.  The inner
         list is an \emph{enumerated} list.
         \begin{enumerate}
            \item This is the first item of an enumerated
                  list that is nested within the itemized list.

            \item This is the second item of the inner list.
                  \LaTeX\ allows you to nest lists deeper than
                  you really should.
         \end{enumerate}
         This is the rest of the second item of the outer
         list.  It is no more interesting than any other
         part of the item.
   \item This is the third item of the list.
\end{itemize}
You can even display poetry.
\begin{verse}
   There is an environment
    for verse \\             % The \\ command separates lines
   Whose features some poets % within a stanza.
   will curse.

                             % One or more blank lines separate stanzas.

   For instead of making\\
   Them do \emph{all} line breaking, \\
   It allows them to put too many words on a line when they'd rather be
   forced to be terse.
\end{verse}

Mathematical formulas may also be displayed.  A
displayed formula
is
one-line long; multiline
formulas require special formatting instructions.
   \[  \ip{\Gamma}{\psi'} = x'' + y^{2} + z_{i}^{n}\]
Don't start a paragraph with a displayed equation,
nor make one a paragraph by itself.

\end{document}               % End of document.

\section{Bench Huge Section 961}
% synthetic bench section 961
% This is a sample LaTeX input file.  (Version of 12 August 2004.)
%
% A '%' character causes TeX to ignore all remaining text on the line,
% and is used for comments like this one.

\documentclass{article}      % Specifies the document class

                             % The preamble begins here.
\title{An Example Document}  % Declares the document's title.
\author{Leslie Lamport}      % Declares the author's name.
\date{January 21, 1994}      % Deleting this command produces today's date.

\newcommand{\ip}[2]{(#1, #2)}
                             % Defines \ip{arg1}{arg2} to mean
                             % (arg1, arg2).

%\newcommand{\ip}[2]{\langle #1 | #2\rangle}
                             % This is an alternative definition of
                             % \ip that is commented out.

\begin{document}             % End of preamble and beginning of text.

\maketitle                   % Produces the title.

This is an example input file.  Comparing it with
the output it generates can show you how to
produce a simple document of your own.

\section{Ordinary Text}      % Produces section heading.  Lower-level
                             % sections are begun with similar
                             % \subsection and \subsubsection commands.

The ends  of words and sentences are marked
  by   spaces. It  doesn't matter how many
spaces    you type; one is as good as 100.  The
end of   a line counts as a space.

One   or more   blank lines denote the  end
of  a paragraph.

Since any number of consecutive spaces are treated
like a single one, the formatting of the input
file makes no difference to
      \LaTeX,                % The \LaTeX command generates the LaTeX logo.
but it makes a difference to you.  When you use
\LaTeX, making your input file as easy to read
as possible will be a great help as you write
your document and when you change it.  This sample
file shows how you can add comments to your own input
file.

Because printing is different from typewriting,
there are a number of things that you have to do
differently when preparing an input file than if
you were just typing the document directly.
Quotation marks like
       ``this''
have to be handled specially, as do quotes within
quotes:
       ``\,`this'            % \, separates the double and single quote.
        is what I just
        wrote, not  `that'\,''.

Dashes come in three sizes: an
       intra-word
dash, a medium dash for number ranges like
       1--2,
and a punctuation
       dash---like
this.

A sentence-ending space should be larger than the
space between words within a sentence.  You
sometimes have to type special commands in
conjunction with punctuation characters to get
this right, as in the following sentence.
       Gnats, gnus, etc.\ all  % `\ ' makes an inter-word space.
       begin with G\@.         % \@ marks end-of-sentence punctuation.
You should check the spaces after periods when
reading your output to make sure you haven't
forgotten any special cases.  Generating an
ellipsis
       \ldots\               % `\ ' is needed after `\ldots' because TeX
                             % ignores spaces after command names like \ldots
                             % made from \ + letters.
                             %
                             % Note how a `%' character causes TeX to ignore
                             % the end of the input line, so these blank lines
                             % do not start a new paragraph.
                             %
with the right spacing around the periods requires
a special command.

\LaTeX\ interprets some common characters as
commands, so you must type special commands to
generate them.  These characters include the
following:
       \$ \& \% \# \{ and \}.

In printing, text is usually emphasized with an
       \emph{italic}
type style.

\begin{em}
   A long segment of text can also be emphasized
   in this way.  Text within such a segment can be
   given \emph{additional} emphasis.
\end{em}

It is sometimes necessary to prevent \LaTeX\ from
breaking a line where it might otherwise do so.
This may be at a space, as between the ``Mr.''\ and
``Jones'' in
       ``Mr.~Jones'',        % ~ produces an unbreakable interword space.
or within a word---especially when the word is a
symbol like
       \mbox{\emph{itemnum}}
that makes little sense when hyphenated across
lines.

Footnotes\footnote{This is an example of a footnote.}
pose no problem.

\LaTeX\ is good at typesetting mathematical formulas
like
       \( x-3y + z = 7 \)
or
       \( a_{1} > x^{2n} + y^{2n} > x' \)
or
       \( \ip{A}{B} = \sum_{i} a_{i} b_{i} \).
The spaces you type in a formula are
ignored.  Remember that a letter like
       $x$                   % $ ... $  and  \( ... \)  are equivalent
is a formula when it denotes a mathematical
symbol, and it should be typed as one.

\section{Displayed Text}

Text is displayed by indenting it from the left
margin.  Quotations are commonly displayed.  There
are short quotations
\begin{quote}
   This is a short quotation.  It consists of a
   single paragraph of text.  See how it is formatted.
\end{quote}
and longer ones.
\begin{quotation}
   This is a longer quotation.  It consists of two
   paragraphs of text, neither of which are
   particularly interesting.

   This is the second paragraph of the quotation.  It
   is just as dull as the first paragraph.
\end{quotation}
Another frequently-displayed structure is a list.
The following is an example of an \emph{itemized}
list.
\begin{itemize}
   \item This is the first item of an itemized list.
         Each item in the list is marked with a ``tick''.
         You don't have to worry about what kind of tick
         mark is used.

   \item This is the second item of the list.  It
         contains another list nested inside it.  The inner
         list is an \emph{enumerated} list.
         \begin{enumerate}
            \item This is the first item of an enumerated
                  list that is nested within the itemized list.

            \item This is the second item of the inner list.
                  \LaTeX\ allows you to nest lists deeper than
                  you really should.
         \end{enumerate}
         This is the rest of the second item of the outer
         list.  It is no more interesting than any other
         part of the item.
   \item This is the third item of the list.
\end{itemize}
You can even display poetry.
\begin{verse}
   There is an environment
    for verse \\             % The \\ command separates lines
   Whose features some poets % within a stanza.
   will curse.

                             % One or more blank lines separate stanzas.

   For instead of making\\
   Them do \emph{all} line breaking, \\
   It allows them to put too many words on a line when they'd rather be
   forced to be terse.
\end{verse}

Mathematical formulas may also be displayed.  A
displayed formula
is
one-line long; multiline
formulas require special formatting instructions.
   \[  \ip{\Gamma}{\psi'} = x'' + y^{2} + z_{i}^{n}\]
Don't start a paragraph with a displayed equation,
nor make one a paragraph by itself.

\end{document}               % End of document.

\section{Bench Huge Section 962}
% synthetic bench section 962
% This is a sample LaTeX input file.  (Version of 12 August 2004.)
%
% A '%' character causes TeX to ignore all remaining text on the line,
% and is used for comments like this one.

\documentclass{article}      % Specifies the document class

                             % The preamble begins here.
\title{An Example Document}  % Declares the document's title.
\author{Leslie Lamport}      % Declares the author's name.
\date{January 21, 1994}      % Deleting this command produces today's date.

\newcommand{\ip}[2]{(#1, #2)}
                             % Defines \ip{arg1}{arg2} to mean
                             % (arg1, arg2).

%\newcommand{\ip}[2]{\langle #1 | #2\rangle}
                             % This is an alternative definition of
                             % \ip that is commented out.

\begin{document}             % End of preamble and beginning of text.

\maketitle                   % Produces the title.

This is an example input file.  Comparing it with
the output it generates can show you how to
produce a simple document of your own.

\section{Ordinary Text}      % Produces section heading.  Lower-level
                             % sections are begun with similar
                             % \subsection and \subsubsection commands.

The ends  of words and sentences are marked
  by   spaces. It  doesn't matter how many
spaces    you type; one is as good as 100.  The
end of   a line counts as a space.

One   or more   blank lines denote the  end
of  a paragraph.

Since any number of consecutive spaces are treated
like a single one, the formatting of the input
file makes no difference to
      \LaTeX,                % The \LaTeX command generates the LaTeX logo.
but it makes a difference to you.  When you use
\LaTeX, making your input file as easy to read
as possible will be a great help as you write
your document and when you change it.  This sample
file shows how you can add comments to your own input
file.

Because printing is different from typewriting,
there are a number of things that you have to do
differently when preparing an input file than if
you were just typing the document directly.
Quotation marks like
       ``this''
have to be handled specially, as do quotes within
quotes:
       ``\,`this'            % \, separates the double and single quote.
        is what I just
        wrote, not  `that'\,''.

Dashes come in three sizes: an
       intra-word
dash, a medium dash for number ranges like
       1--2,
and a punctuation
       dash---like
this.

A sentence-ending space should be larger than the
space between words within a sentence.  You
sometimes have to type special commands in
conjunction with punctuation characters to get
this right, as in the following sentence.
       Gnats, gnus, etc.\ all  % `\ ' makes an inter-word space.
       begin with G\@.         % \@ marks end-of-sentence punctuation.
You should check the spaces after periods when
reading your output to make sure you haven't
forgotten any special cases.  Generating an
ellipsis
       \ldots\               % `\ ' is needed after `\ldots' because TeX
                             % ignores spaces after command names like \ldots
                             % made from \ + letters.
                             %
                             % Note how a `%' character causes TeX to ignore
                             % the end of the input line, so these blank lines
                             % do not start a new paragraph.
                             %
with the right spacing around the periods requires
a special command.

\LaTeX\ interprets some common characters as
commands, so you must type special commands to
generate them.  These characters include the
following:
       \$ \& \% \# \{ and \}.

In printing, text is usually emphasized with an
       \emph{italic}
type style.

\begin{em}
   A long segment of text can also be emphasized
   in this way.  Text within such a segment can be
   given \emph{additional} emphasis.
\end{em}

It is sometimes necessary to prevent \LaTeX\ from
breaking a line where it might otherwise do so.
This may be at a space, as between the ``Mr.''\ and
``Jones'' in
       ``Mr.~Jones'',        % ~ produces an unbreakable interword space.
or within a word---especially when the word is a
symbol like
       \mbox{\emph{itemnum}}
that makes little sense when hyphenated across
lines.

Footnotes\footnote{This is an example of a footnote.}
pose no problem.

\LaTeX\ is good at typesetting mathematical formulas
like
       \( x-3y + z = 7 \)
or
       \( a_{1} > x^{2n} + y^{2n} > x' \)
or
       \( \ip{A}{B} = \sum_{i} a_{i} b_{i} \).
The spaces you type in a formula are
ignored.  Remember that a letter like
       $x$                   % $ ... $  and  \( ... \)  are equivalent
is a formula when it denotes a mathematical
symbol, and it should be typed as one.

\section{Displayed Text}

Text is displayed by indenting it from the left
margin.  Quotations are commonly displayed.  There
are short quotations
\begin{quote}
   This is a short quotation.  It consists of a
   single paragraph of text.  See how it is formatted.
\end{quote}
and longer ones.
\begin{quotation}
   This is a longer quotation.  It consists of two
   paragraphs of text, neither of which are
   particularly interesting.

   This is the second paragraph of the quotation.  It
   is just as dull as the first paragraph.
\end{quotation}
Another frequently-displayed structure is a list.
The following is an example of an \emph{itemized}
list.
\begin{itemize}
   \item This is the first item of an itemized list.
         Each item in the list is marked with a ``tick''.
         You don't have to worry about what kind of tick
         mark is used.

   \item This is the second item of the list.  It
         contains another list nested inside it.  The inner
         list is an \emph{enumerated} list.
         \begin{enumerate}
            \item This is the first item of an enumerated
                  list that is nested within the itemized list.

            \item This is the second item of the inner list.
                  \LaTeX\ allows you to nest lists deeper than
                  you really should.
         \end{enumerate}
         This is the rest of the second item of the outer
         list.  It is no more interesting than any other
         part of the item.
   \item This is the third item of the list.
\end{itemize}
You can even display poetry.
\begin{verse}
   There is an environment
    for verse \\             % The \\ command separates lines
   Whose features some poets % within a stanza.
   will curse.

                             % One or more blank lines separate stanzas.

   For instead of making\\
   Them do \emph{all} line breaking, \\
   It allows them to put too many words on a line when they'd rather be
   forced to be terse.
\end{verse}

Mathematical formulas may also be displayed.  A
displayed formula
is
one-line long; multiline
formulas require special formatting instructions.
   \[  \ip{\Gamma}{\psi'} = x'' + y^{2} + z_{i}^{n}\]
Don't start a paragraph with a displayed equation,
nor make one a paragraph by itself.

\end{document}               % End of document.

\section{Bench Huge Section 963}
% synthetic bench section 963
% This is a sample LaTeX input file.  (Version of 12 August 2004.)
%
% A '%' character causes TeX to ignore all remaining text on the line,
% and is used for comments like this one.

\documentclass{article}      % Specifies the document class

                             % The preamble begins here.
\title{An Example Document}  % Declares the document's title.
\author{Leslie Lamport}      % Declares the author's name.
\date{January 21, 1994}      % Deleting this command produces today's date.

\newcommand{\ip}[2]{(#1, #2)}
                             % Defines \ip{arg1}{arg2} to mean
                             % (arg1, arg2).

%\newcommand{\ip}[2]{\langle #1 | #2\rangle}
                             % This is an alternative definition of
                             % \ip that is commented out.

\begin{document}             % End of preamble and beginning of text.

\maketitle                   % Produces the title.

This is an example input file.  Comparing it with
the output it generates can show you how to
produce a simple document of your own.

\section{Ordinary Text}      % Produces section heading.  Lower-level
                             % sections are begun with similar
                             % \subsection and \subsubsection commands.

The ends  of words and sentences are marked
  by   spaces. It  doesn't matter how many
spaces    you type; one is as good as 100.  The
end of   a line counts as a space.

One   or more   blank lines denote the  end
of  a paragraph.

Since any number of consecutive spaces are treated
like a single one, the formatting of the input
file makes no difference to
      \LaTeX,                % The \LaTeX command generates the LaTeX logo.
but it makes a difference to you.  When you use
\LaTeX, making your input file as easy to read
as possible will be a great help as you write
your document and when you change it.  This sample
file shows how you can add comments to your own input
file.

Because printing is different from typewriting,
there are a number of things that you have to do
differently when preparing an input file than if
you were just typing the document directly.
Quotation marks like
       ``this''
have to be handled specially, as do quotes within
quotes:
       ``\,`this'            % \, separates the double and single quote.
        is what I just
        wrote, not  `that'\,''.

Dashes come in three sizes: an
       intra-word
dash, a medium dash for number ranges like
       1--2,
and a punctuation
       dash---like
this.

A sentence-ending space should be larger than the
space between words within a sentence.  You
sometimes have to type special commands in
conjunction with punctuation characters to get
this right, as in the following sentence.
       Gnats, gnus, etc.\ all  % `\ ' makes an inter-word space.
       begin with G\@.         % \@ marks end-of-sentence punctuation.
You should check the spaces after periods when
reading your output to make sure you haven't
forgotten any special cases.  Generating an
ellipsis
       \ldots\               % `\ ' is needed after `\ldots' because TeX
                             % ignores spaces after command names like \ldots
                             % made from \ + letters.
                             %
                             % Note how a `%' character causes TeX to ignore
                             % the end of the input line, so these blank lines
                             % do not start a new paragraph.
                             %
with the right spacing around the periods requires
a special command.

\LaTeX\ interprets some common characters as
commands, so you must type special commands to
generate them.  These characters include the
following:
       \$ \& \% \# \{ and \}.

In printing, text is usually emphasized with an
       \emph{italic}
type style.

\begin{em}
   A long segment of text can also be emphasized
   in this way.  Text within such a segment can be
   given \emph{additional} emphasis.
\end{em}

It is sometimes necessary to prevent \LaTeX\ from
breaking a line where it might otherwise do so.
This may be at a space, as between the ``Mr.''\ and
``Jones'' in
       ``Mr.~Jones'',        % ~ produces an unbreakable interword space.
or within a word---especially when the word is a
symbol like
       \mbox{\emph{itemnum}}
that makes little sense when hyphenated across
lines.

Footnotes\footnote{This is an example of a footnote.}
pose no problem.

\LaTeX\ is good at typesetting mathematical formulas
like
       \( x-3y + z = 7 \)
or
       \( a_{1} > x^{2n} + y^{2n} > x' \)
or
       \( \ip{A}{B} = \sum_{i} a_{i} b_{i} \).
The spaces you type in a formula are
ignored.  Remember that a letter like
       $x$                   % $ ... $  and  \( ... \)  are equivalent
is a formula when it denotes a mathematical
symbol, and it should be typed as one.

\section{Displayed Text}

Text is displayed by indenting it from the left
margin.  Quotations are commonly displayed.  There
are short quotations
\begin{quote}
   This is a short quotation.  It consists of a
   single paragraph of text.  See how it is formatted.
\end{quote}
and longer ones.
\begin{quotation}
   This is a longer quotation.  It consists of two
   paragraphs of text, neither of which are
   particularly interesting.

   This is the second paragraph of the quotation.  It
   is just as dull as the first paragraph.
\end{quotation}
Another frequently-displayed structure is a list.
The following is an example of an \emph{itemized}
list.
\begin{itemize}
   \item This is the first item of an itemized list.
         Each item in the list is marked with a ``tick''.
         You don't have to worry about what kind of tick
         mark is used.

   \item This is the second item of the list.  It
         contains another list nested inside it.  The inner
         list is an \emph{enumerated} list.
         \begin{enumerate}
            \item This is the first item of an enumerated
                  list that is nested within the itemized list.

            \item This is the second item of the inner list.
                  \LaTeX\ allows you to nest lists deeper than
                  you really should.
         \end{enumerate}
         This is the rest of the second item of the outer
         list.  It is no more interesting than any other
         part of the item.
   \item This is the third item of the list.
\end{itemize}
You can even display poetry.
\begin{verse}
   There is an environment
    for verse \\             % The \\ command separates lines
   Whose features some poets % within a stanza.
   will curse.

                             % One or more blank lines separate stanzas.

   For instead of making\\
   Them do \emph{all} line breaking, \\
   It allows them to put too many words on a line when they'd rather be
   forced to be terse.
\end{verse}

Mathematical formulas may also be displayed.  A
displayed formula
is
one-line long; multiline
formulas require special formatting instructions.
   \[  \ip{\Gamma}{\psi'} = x'' + y^{2} + z_{i}^{n}\]
Don't start a paragraph with a displayed equation,
nor make one a paragraph by itself.

\end{document}               % End of document.

\section{Bench Huge Section 964}
% synthetic bench section 964
% This is a sample LaTeX input file.  (Version of 12 August 2004.)
%
% A '%' character causes TeX to ignore all remaining text on the line,
% and is used for comments like this one.

\documentclass{article}      % Specifies the document class

                             % The preamble begins here.
\title{An Example Document}  % Declares the document's title.
\author{Leslie Lamport}      % Declares the author's name.
\date{January 21, 1994}      % Deleting this command produces today's date.

\newcommand{\ip}[2]{(#1, #2)}
                             % Defines \ip{arg1}{arg2} to mean
                             % (arg1, arg2).

%\newcommand{\ip}[2]{\langle #1 | #2\rangle}
                             % This is an alternative definition of
                             % \ip that is commented out.

\begin{document}             % End of preamble and beginning of text.

\maketitle                   % Produces the title.

This is an example input file.  Comparing it with
the output it generates can show you how to
produce a simple document of your own.

\section{Ordinary Text}      % Produces section heading.  Lower-level
                             % sections are begun with similar
                             % \subsection and \subsubsection commands.

The ends  of words and sentences are marked
  by   spaces. It  doesn't matter how many
spaces    you type; one is as good as 100.  The
end of   a line counts as a space.

One   or more   blank lines denote the  end
of  a paragraph.

Since any number of consecutive spaces are treated
like a single one, the formatting of the input
file makes no difference to
      \LaTeX,                % The \LaTeX command generates the LaTeX logo.
but it makes a difference to you.  When you use
\LaTeX, making your input file as easy to read
as possible will be a great help as you write
your document and when you change it.  This sample
file shows how you can add comments to your own input
file.

Because printing is different from typewriting,
there are a number of things that you have to do
differently when preparing an input file than if
you were just typing the document directly.
Quotation marks like
       ``this''
have to be handled specially, as do quotes within
quotes:
       ``\,`this'            % \, separates the double and single quote.
        is what I just
        wrote, not  `that'\,''.

Dashes come in three sizes: an
       intra-word
dash, a medium dash for number ranges like
       1--2,
and a punctuation
       dash---like
this.

A sentence-ending space should be larger than the
space between words within a sentence.  You
sometimes have to type special commands in
conjunction with punctuation characters to get
this right, as in the following sentence.
       Gnats, gnus, etc.\ all  % `\ ' makes an inter-word space.
       begin with G\@.         % \@ marks end-of-sentence punctuation.
You should check the spaces after periods when
reading your output to make sure you haven't
forgotten any special cases.  Generating an
ellipsis
       \ldots\               % `\ ' is needed after `\ldots' because TeX
                             % ignores spaces after command names like \ldots
                             % made from \ + letters.
                             %
                             % Note how a `%' character causes TeX to ignore
                             % the end of the input line, so these blank lines
                             % do not start a new paragraph.
                             %
with the right spacing around the periods requires
a special command.

\LaTeX\ interprets some common characters as
commands, so you must type special commands to
generate them.  These characters include the
following:
       \$ \& \% \# \{ and \}.

In printing, text is usually emphasized with an
       \emph{italic}
type style.

\begin{em}
   A long segment of text can also be emphasized
   in this way.  Text within such a segment can be
   given \emph{additional} emphasis.
\end{em}

It is sometimes necessary to prevent \LaTeX\ from
breaking a line where it might otherwise do so.
This may be at a space, as between the ``Mr.''\ and
``Jones'' in
       ``Mr.~Jones'',        % ~ produces an unbreakable interword space.
or within a word---especially when the word is a
symbol like
       \mbox{\emph{itemnum}}
that makes little sense when hyphenated across
lines.

Footnotes\footnote{This is an example of a footnote.}
pose no problem.

\LaTeX\ is good at typesetting mathematical formulas
like
       \( x-3y + z = 7 \)
or
       \( a_{1} > x^{2n} + y^{2n} > x' \)
or
       \( \ip{A}{B} = \sum_{i} a_{i} b_{i} \).
The spaces you type in a formula are
ignored.  Remember that a letter like
       $x$                   % $ ... $  and  \( ... \)  are equivalent
is a formula when it denotes a mathematical
symbol, and it should be typed as one.

\section{Displayed Text}

Text is displayed by indenting it from the left
margin.  Quotations are commonly displayed.  There
are short quotations
\begin{quote}
   This is a short quotation.  It consists of a
   single paragraph of text.  See how it is formatted.
\end{quote}
and longer ones.
\begin{quotation}
   This is a longer quotation.  It consists of two
   paragraphs of text, neither of which are
   particularly interesting.

   This is the second paragraph of the quotation.  It
   is just as dull as the first paragraph.
\end{quotation}
Another frequently-displayed structure is a list.
The following is an example of an \emph{itemized}
list.
\begin{itemize}
   \item This is the first item of an itemized list.
         Each item in the list is marked with a ``tick''.
         You don't have to worry about what kind of tick
         mark is used.

   \item This is the second item of the list.  It
         contains another list nested inside it.  The inner
         list is an \emph{enumerated} list.
         \begin{enumerate}
            \item This is the first item of an enumerated
                  list that is nested within the itemized list.

            \item This is the second item of the inner list.
                  \LaTeX\ allows you to nest lists deeper than
                  you really should.
         \end{enumerate}
         This is the rest of the second item of the outer
         list.  It is no more interesting than any other
         part of the item.
   \item This is the third item of the list.
\end{itemize}
You can even display poetry.
\begin{verse}
   There is an environment
    for verse \\             % The \\ command separates lines
   Whose features some poets % within a stanza.
   will curse.

                             % One or more blank lines separate stanzas.

   For instead of making\\
   Them do \emph{all} line breaking, \\
   It allows them to put too many words on a line when they'd rather be
   forced to be terse.
\end{verse}

Mathematical formulas may also be displayed.  A
displayed formula
is
one-line long; multiline
formulas require special formatting instructions.
   \[  \ip{\Gamma}{\psi'} = x'' + y^{2} + z_{i}^{n}\]
Don't start a paragraph with a displayed equation,
nor make one a paragraph by itself.

\end{document}               % End of document.

\section{Bench Huge Section 965}
% synthetic bench section 965
% This is a sample LaTeX input file.  (Version of 12 August 2004.)
%
% A '%' character causes TeX to ignore all remaining text on the line,
% and is used for comments like this one.

\documentclass{article}      % Specifies the document class

                             % The preamble begins here.
\title{An Example Document}  % Declares the document's title.
\author{Leslie Lamport}      % Declares the author's name.
\date{January 21, 1994}      % Deleting this command produces today's date.

\newcommand{\ip}[2]{(#1, #2)}
                             % Defines \ip{arg1}{arg2} to mean
                             % (arg1, arg2).

%\newcommand{\ip}[2]{\langle #1 | #2\rangle}
                             % This is an alternative definition of
                             % \ip that is commented out.

\begin{document}             % End of preamble and beginning of text.

\maketitle                   % Produces the title.

This is an example input file.  Comparing it with
the output it generates can show you how to
produce a simple document of your own.

\section{Ordinary Text}      % Produces section heading.  Lower-level
                             % sections are begun with similar
                             % \subsection and \subsubsection commands.

The ends  of words and sentences are marked
  by   spaces. It  doesn't matter how many
spaces    you type; one is as good as 100.  The
end of   a line counts as a space.

One   or more   blank lines denote the  end
of  a paragraph.

Since any number of consecutive spaces are treated
like a single one, the formatting of the input
file makes no difference to
      \LaTeX,                % The \LaTeX command generates the LaTeX logo.
but it makes a difference to you.  When you use
\LaTeX, making your input file as easy to read
as possible will be a great help as you write
your document and when you change it.  This sample
file shows how you can add comments to your own input
file.

Because printing is different from typewriting,
there are a number of things that you have to do
differently when preparing an input file than if
you were just typing the document directly.
Quotation marks like
       ``this''
have to be handled specially, as do quotes within
quotes:
       ``\,`this'            % \, separates the double and single quote.
        is what I just
        wrote, not  `that'\,''.

Dashes come in three sizes: an
       intra-word
dash, a medium dash for number ranges like
       1--2,
and a punctuation
       dash---like
this.

A sentence-ending space should be larger than the
space between words within a sentence.  You
sometimes have to type special commands in
conjunction with punctuation characters to get
this right, as in the following sentence.
       Gnats, gnus, etc.\ all  % `\ ' makes an inter-word space.
       begin with G\@.         % \@ marks end-of-sentence punctuation.
You should check the spaces after periods when
reading your output to make sure you haven't
forgotten any special cases.  Generating an
ellipsis
       \ldots\               % `\ ' is needed after `\ldots' because TeX
                             % ignores spaces after command names like \ldots
                             % made from \ + letters.
                             %
                             % Note how a `%' character causes TeX to ignore
                             % the end of the input line, so these blank lines
                             % do not start a new paragraph.
                             %
with the right spacing around the periods requires
a special command.

\LaTeX\ interprets some common characters as
commands, so you must type special commands to
generate them.  These characters include the
following:
       \$ \& \% \# \{ and \}.

In printing, text is usually emphasized with an
       \emph{italic}
type style.

\begin{em}
   A long segment of text can also be emphasized
   in this way.  Text within such a segment can be
   given \emph{additional} emphasis.
\end{em}

It is sometimes necessary to prevent \LaTeX\ from
breaking a line where it might otherwise do so.
This may be at a space, as between the ``Mr.''\ and
``Jones'' in
       ``Mr.~Jones'',        % ~ produces an unbreakable interword space.
or within a word---especially when the word is a
symbol like
       \mbox{\emph{itemnum}}
that makes little sense when hyphenated across
lines.

Footnotes\footnote{This is an example of a footnote.}
pose no problem.

\LaTeX\ is good at typesetting mathematical formulas
like
       \( x-3y + z = 7 \)
or
       \( a_{1} > x^{2n} + y^{2n} > x' \)
or
       \( \ip{A}{B} = \sum_{i} a_{i} b_{i} \).
The spaces you type in a formula are
ignored.  Remember that a letter like
       $x$                   % $ ... $  and  \( ... \)  are equivalent
is a formula when it denotes a mathematical
symbol, and it should be typed as one.

\section{Displayed Text}

Text is displayed by indenting it from the left
margin.  Quotations are commonly displayed.  There
are short quotations
\begin{quote}
   This is a short quotation.  It consists of a
   single paragraph of text.  See how it is formatted.
\end{quote}
and longer ones.
\begin{quotation}
   This is a longer quotation.  It consists of two
   paragraphs of text, neither of which are
   particularly interesting.

   This is the second paragraph of the quotation.  It
   is just as dull as the first paragraph.
\end{quotation}
Another frequently-displayed structure is a list.
The following is an example of an \emph{itemized}
list.
\begin{itemize}
   \item This is the first item of an itemized list.
         Each item in the list is marked with a ``tick''.
         You don't have to worry about what kind of tick
         mark is used.

   \item This is the second item of the list.  It
         contains another list nested inside it.  The inner
         list is an \emph{enumerated} list.
         \begin{enumerate}
            \item This is the first item of an enumerated
                  list that is nested within the itemized list.

            \item This is the second item of the inner list.
                  \LaTeX\ allows you to nest lists deeper than
                  you really should.
         \end{enumerate}
         This is the rest of the second item of the outer
         list.  It is no more interesting than any other
         part of the item.
   \item This is the third item of the list.
\end{itemize}
You can even display poetry.
\begin{verse}
   There is an environment
    for verse \\             % The \\ command separates lines
   Whose features some poets % within a stanza.
   will curse.

                             % One or more blank lines separate stanzas.

   For instead of making\\
   Them do \emph{all} line breaking, \\
   It allows them to put too many words on a line when they'd rather be
   forced to be terse.
\end{verse}

Mathematical formulas may also be displayed.  A
displayed formula
is
one-line long; multiline
formulas require special formatting instructions.
   \[  \ip{\Gamma}{\psi'} = x'' + y^{2} + z_{i}^{n}\]
Don't start a paragraph with a displayed equation,
nor make one a paragraph by itself.

\end{document}               % End of document.

\section{Bench Huge Section 966}
% synthetic bench section 966
% This is a sample LaTeX input file.  (Version of 12 August 2004.)
%
% A '%' character causes TeX to ignore all remaining text on the line,
% and is used for comments like this one.

\documentclass{article}      % Specifies the document class

                             % The preamble begins here.
\title{An Example Document}  % Declares the document's title.
\author{Leslie Lamport}      % Declares the author's name.
\date{January 21, 1994}      % Deleting this command produces today's date.

\newcommand{\ip}[2]{(#1, #2)}
                             % Defines \ip{arg1}{arg2} to mean
                             % (arg1, arg2).

%\newcommand{\ip}[2]{\langle #1 | #2\rangle}
                             % This is an alternative definition of
                             % \ip that is commented out.

\begin{document}             % End of preamble and beginning of text.

\maketitle                   % Produces the title.

This is an example input file.  Comparing it with
the output it generates can show you how to
produce a simple document of your own.

\section{Ordinary Text}      % Produces section heading.  Lower-level
                             % sections are begun with similar
                             % \subsection and \subsubsection commands.

The ends  of words and sentences are marked
  by   spaces. It  doesn't matter how many
spaces    you type; one is as good as 100.  The
end of   a line counts as a space.

One   or more   blank lines denote the  end
of  a paragraph.

Since any number of consecutive spaces are treated
like a single one, the formatting of the input
file makes no difference to
      \LaTeX,                % The \LaTeX command generates the LaTeX logo.
but it makes a difference to you.  When you use
\LaTeX, making your input file as easy to read
as possible will be a great help as you write
your document and when you change it.  This sample
file shows how you can add comments to your own input
file.

Because printing is different from typewriting,
there are a number of things that you have to do
differently when preparing an input file than if
you were just typing the document directly.
Quotation marks like
       ``this''
have to be handled specially, as do quotes within
quotes:
       ``\,`this'            % \, separates the double and single quote.
        is what I just
        wrote, not  `that'\,''.

Dashes come in three sizes: an
       intra-word
dash, a medium dash for number ranges like
       1--2,
and a punctuation
       dash---like
this.

A sentence-ending space should be larger than the
space between words within a sentence.  You
sometimes have to type special commands in
conjunction with punctuation characters to get
this right, as in the following sentence.
       Gnats, gnus, etc.\ all  % `\ ' makes an inter-word space.
       begin with G\@.         % \@ marks end-of-sentence punctuation.
You should check the spaces after periods when
reading your output to make sure you haven't
forgotten any special cases.  Generating an
ellipsis
       \ldots\               % `\ ' is needed after `\ldots' because TeX
                             % ignores spaces after command names like \ldots
                             % made from \ + letters.
                             %
                             % Note how a `%' character causes TeX to ignore
                             % the end of the input line, so these blank lines
                             % do not start a new paragraph.
                             %
with the right spacing around the periods requires
a special command.

\LaTeX\ interprets some common characters as
commands, so you must type special commands to
generate them.  These characters include the
following:
       \$ \& \% \# \{ and \}.

In printing, text is usually emphasized with an
       \emph{italic}
type style.

\begin{em}
   A long segment of text can also be emphasized
   in this way.  Text within such a segment can be
   given \emph{additional} emphasis.
\end{em}

It is sometimes necessary to prevent \LaTeX\ from
breaking a line where it might otherwise do so.
This may be at a space, as between the ``Mr.''\ and
``Jones'' in
       ``Mr.~Jones'',        % ~ produces an unbreakable interword space.
or within a word---especially when the word is a
symbol like
       \mbox{\emph{itemnum}}
that makes little sense when hyphenated across
lines.

Footnotes\footnote{This is an example of a footnote.}
pose no problem.

\LaTeX\ is good at typesetting mathematical formulas
like
       \( x-3y + z = 7 \)
or
       \( a_{1} > x^{2n} + y^{2n} > x' \)
or
       \( \ip{A}{B} = \sum_{i} a_{i} b_{i} \).
The spaces you type in a formula are
ignored.  Remember that a letter like
       $x$                   % $ ... $  and  \( ... \)  are equivalent
is a formula when it denotes a mathematical
symbol, and it should be typed as one.

\section{Displayed Text}

Text is displayed by indenting it from the left
margin.  Quotations are commonly displayed.  There
are short quotations
\begin{quote}
   This is a short quotation.  It consists of a
   single paragraph of text.  See how it is formatted.
\end{quote}
and longer ones.
\begin{quotation}
   This is a longer quotation.  It consists of two
   paragraphs of text, neither of which are
   particularly interesting.

   This is the second paragraph of the quotation.  It
   is just as dull as the first paragraph.
\end{quotation}
Another frequently-displayed structure is a list.
The following is an example of an \emph{itemized}
list.
\begin{itemize}
   \item This is the first item of an itemized list.
         Each item in the list is marked with a ``tick''.
         You don't have to worry about what kind of tick
         mark is used.

   \item This is the second item of the list.  It
         contains another list nested inside it.  The inner
         list is an \emph{enumerated} list.
         \begin{enumerate}
            \item This is the first item of an enumerated
                  list that is nested within the itemized list.

            \item This is the second item of the inner list.
                  \LaTeX\ allows you to nest lists deeper than
                  you really should.
         \end{enumerate}
         This is the rest of the second item of the outer
         list.  It is no more interesting than any other
         part of the item.
   \item This is the third item of the list.
\end{itemize}
You can even display poetry.
\begin{verse}
   There is an environment
    for verse \\             % The \\ command separates lines
   Whose features some poets % within a stanza.
   will curse.

                             % One or more blank lines separate stanzas.

   For instead of making\\
   Them do \emph{all} line breaking, \\
   It allows them to put too many words on a line when they'd rather be
   forced to be terse.
\end{verse}

Mathematical formulas may also be displayed.  A
displayed formula
is
one-line long; multiline
formulas require special formatting instructions.
   \[  \ip{\Gamma}{\psi'} = x'' + y^{2} + z_{i}^{n}\]
Don't start a paragraph with a displayed equation,
nor make one a paragraph by itself.

\end{document}               % End of document.

\section{Bench Huge Section 967}
% synthetic bench section 967
% This is a sample LaTeX input file.  (Version of 12 August 2004.)
%
% A '%' character causes TeX to ignore all remaining text on the line,
% and is used for comments like this one.

\documentclass{article}      % Specifies the document class

                             % The preamble begins here.
\title{An Example Document}  % Declares the document's title.
\author{Leslie Lamport}      % Declares the author's name.
\date{January 21, 1994}      % Deleting this command produces today's date.

\newcommand{\ip}[2]{(#1, #2)}
                             % Defines \ip{arg1}{arg2} to mean
                             % (arg1, arg2).

%\newcommand{\ip}[2]{\langle #1 | #2\rangle}
                             % This is an alternative definition of
                             % \ip that is commented out.

\begin{document}             % End of preamble and beginning of text.

\maketitle                   % Produces the title.

This is an example input file.  Comparing it with
the output it generates can show you how to
produce a simple document of your own.

\section{Ordinary Text}      % Produces section heading.  Lower-level
                             % sections are begun with similar
                             % \subsection and \subsubsection commands.

The ends  of words and sentences are marked
  by   spaces. It  doesn't matter how many
spaces    you type; one is as good as 100.  The
end of   a line counts as a space.

One   or more   blank lines denote the  end
of  a paragraph.

Since any number of consecutive spaces are treated
like a single one, the formatting of the input
file makes no difference to
      \LaTeX,                % The \LaTeX command generates the LaTeX logo.
but it makes a difference to you.  When you use
\LaTeX, making your input file as easy to read
as possible will be a great help as you write
your document and when you change it.  This sample
file shows how you can add comments to your own input
file.

Because printing is different from typewriting,
there are a number of things that you have to do
differently when preparing an input file than if
you were just typing the document directly.
Quotation marks like
       ``this''
have to be handled specially, as do quotes within
quotes:
       ``\,`this'            % \, separates the double and single quote.
        is what I just
        wrote, not  `that'\,''.

Dashes come in three sizes: an
       intra-word
dash, a medium dash for number ranges like
       1--2,
and a punctuation
       dash---like
this.

A sentence-ending space should be larger than the
space between words within a sentence.  You
sometimes have to type special commands in
conjunction with punctuation characters to get
this right, as in the following sentence.
       Gnats, gnus, etc.\ all  % `\ ' makes an inter-word space.
       begin with G\@.         % \@ marks end-of-sentence punctuation.
You should check the spaces after periods when
reading your output to make sure you haven't
forgotten any special cases.  Generating an
ellipsis
       \ldots\               % `\ ' is needed after `\ldots' because TeX
                             % ignores spaces after command names like \ldots
                             % made from \ + letters.
                             %
                             % Note how a `%' character causes TeX to ignore
                             % the end of the input line, so these blank lines
                             % do not start a new paragraph.
                             %
with the right spacing around the periods requires
a special command.

\LaTeX\ interprets some common characters as
commands, so you must type special commands to
generate them.  These characters include the
following:
       \$ \& \% \# \{ and \}.

In printing, text is usually emphasized with an
       \emph{italic}
type style.

\begin{em}
   A long segment of text can also be emphasized
   in this way.  Text within such a segment can be
   given \emph{additional} emphasis.
\end{em}

It is sometimes necessary to prevent \LaTeX\ from
breaking a line where it might otherwise do so.
This may be at a space, as between the ``Mr.''\ and
``Jones'' in
       ``Mr.~Jones'',        % ~ produces an unbreakable interword space.
or within a word---especially when the word is a
symbol like
       \mbox{\emph{itemnum}}
that makes little sense when hyphenated across
lines.

Footnotes\footnote{This is an example of a footnote.}
pose no problem.

\LaTeX\ is good at typesetting mathematical formulas
like
       \( x-3y + z = 7 \)
or
       \( a_{1} > x^{2n} + y^{2n} > x' \)
or
       \( \ip{A}{B} = \sum_{i} a_{i} b_{i} \).
The spaces you type in a formula are
ignored.  Remember that a letter like
       $x$                   % $ ... $  and  \( ... \)  are equivalent
is a formula when it denotes a mathematical
symbol, and it should be typed as one.

\section{Displayed Text}

Text is displayed by indenting it from the left
margin.  Quotations are commonly displayed.  There
are short quotations
\begin{quote}
   This is a short quotation.  It consists of a
   single paragraph of text.  See how it is formatted.
\end{quote}
and longer ones.
\begin{quotation}
   This is a longer quotation.  It consists of two
   paragraphs of text, neither of which are
   particularly interesting.

   This is the second paragraph of the quotation.  It
   is just as dull as the first paragraph.
\end{quotation}
Another frequently-displayed structure is a list.
The following is an example of an \emph{itemized}
list.
\begin{itemize}
   \item This is the first item of an itemized list.
         Each item in the list is marked with a ``tick''.
         You don't have to worry about what kind of tick
         mark is used.

   \item This is the second item of the list.  It
         contains another list nested inside it.  The inner
         list is an \emph{enumerated} list.
         \begin{enumerate}
            \item This is the first item of an enumerated
                  list that is nested within the itemized list.

            \item This is the second item of the inner list.
                  \LaTeX\ allows you to nest lists deeper than
                  you really should.
         \end{enumerate}
         This is the rest of the second item of the outer
         list.  It is no more interesting than any other
         part of the item.
   \item This is the third item of the list.
\end{itemize}
You can even display poetry.
\begin{verse}
   There is an environment
    for verse \\             % The \\ command separates lines
   Whose features some poets % within a stanza.
   will curse.

                             % One or more blank lines separate stanzas.

   For instead of making\\
   Them do \emph{all} line breaking, \\
   It allows them to put too many words on a line when they'd rather be
   forced to be terse.
\end{verse}

Mathematical formulas may also be displayed.  A
displayed formula
is
one-line long; multiline
formulas require special formatting instructions.
   \[  \ip{\Gamma}{\psi'} = x'' + y^{2} + z_{i}^{n}\]
Don't start a paragraph with a displayed equation,
nor make one a paragraph by itself.

\end{document}               % End of document.

\section{Bench Huge Section 968}
% synthetic bench section 968
% This is a sample LaTeX input file.  (Version of 12 August 2004.)
%
% A '%' character causes TeX to ignore all remaining text on the line,
% and is used for comments like this one.

\documentclass{article}      % Specifies the document class

                             % The preamble begins here.
\title{An Example Document}  % Declares the document's title.
\author{Leslie Lamport}      % Declares the author's name.
\date{January 21, 1994}      % Deleting this command produces today's date.

\newcommand{\ip}[2]{(#1, #2)}
                             % Defines \ip{arg1}{arg2} to mean
                             % (arg1, arg2).

%\newcommand{\ip}[2]{\langle #1 | #2\rangle}
                             % This is an alternative definition of
                             % \ip that is commented out.

\begin{document}             % End of preamble and beginning of text.

\maketitle                   % Produces the title.

This is an example input file.  Comparing it with
the output it generates can show you how to
produce a simple document of your own.

\section{Ordinary Text}      % Produces section heading.  Lower-level
                             % sections are begun with similar
                             % \subsection and \subsubsection commands.

The ends  of words and sentences are marked
  by   spaces. It  doesn't matter how many
spaces    you type; one is as good as 100.  The
end of   a line counts as a space.

One   or more   blank lines denote the  end
of  a paragraph.

Since any number of consecutive spaces are treated
like a single one, the formatting of the input
file makes no difference to
      \LaTeX,                % The \LaTeX command generates the LaTeX logo.
but it makes a difference to you.  When you use
\LaTeX, making your input file as easy to read
as possible will be a great help as you write
your document and when you change it.  This sample
file shows how you can add comments to your own input
file.

Because printing is different from typewriting,
there are a number of things that you have to do
differently when preparing an input file than if
you were just typing the document directly.
Quotation marks like
       ``this''
have to be handled specially, as do quotes within
quotes:
       ``\,`this'            % \, separates the double and single quote.
        is what I just
        wrote, not  `that'\,''.

Dashes come in three sizes: an
       intra-word
dash, a medium dash for number ranges like
       1--2,
and a punctuation
       dash---like
this.

A sentence-ending space should be larger than the
space between words within a sentence.  You
sometimes have to type special commands in
conjunction with punctuation characters to get
this right, as in the following sentence.
       Gnats, gnus, etc.\ all  % `\ ' makes an inter-word space.
       begin with G\@.         % \@ marks end-of-sentence punctuation.
You should check the spaces after periods when
reading your output to make sure you haven't
forgotten any special cases.  Generating an
ellipsis
       \ldots\               % `\ ' is needed after `\ldots' because TeX
                             % ignores spaces after command names like \ldots
                             % made from \ + letters.
                             %
                             % Note how a `%' character causes TeX to ignore
                             % the end of the input line, so these blank lines
                             % do not start a new paragraph.
                             %
with the right spacing around the periods requires
a special command.

\LaTeX\ interprets some common characters as
commands, so you must type special commands to
generate them.  These characters include the
following:
       \$ \& \% \# \{ and \}.

In printing, text is usually emphasized with an
       \emph{italic}
type style.

\begin{em}
   A long segment of text can also be emphasized
   in this way.  Text within such a segment can be
   given \emph{additional} emphasis.
\end{em}

It is sometimes necessary to prevent \LaTeX\ from
breaking a line where it might otherwise do so.
This may be at a space, as between the ``Mr.''\ and
``Jones'' in
       ``Mr.~Jones'',        % ~ produces an unbreakable interword space.
or within a word---especially when the word is a
symbol like
       \mbox{\emph{itemnum}}
that makes little sense when hyphenated across
lines.

Footnotes\footnote{This is an example of a footnote.}
pose no problem.

\LaTeX\ is good at typesetting mathematical formulas
like
       \( x-3y + z = 7 \)
or
       \( a_{1} > x^{2n} + y^{2n} > x' \)
or
       \( \ip{A}{B} = \sum_{i} a_{i} b_{i} \).
The spaces you type in a formula are
ignored.  Remember that a letter like
       $x$                   % $ ... $  and  \( ... \)  are equivalent
is a formula when it denotes a mathematical
symbol, and it should be typed as one.

\section{Displayed Text}

Text is displayed by indenting it from the left
margin.  Quotations are commonly displayed.  There
are short quotations
\begin{quote}
   This is a short quotation.  It consists of a
   single paragraph of text.  See how it is formatted.
\end{quote}
and longer ones.
\begin{quotation}
   This is a longer quotation.  It consists of two
   paragraphs of text, neither of which are
   particularly interesting.

   This is the second paragraph of the quotation.  It
   is just as dull as the first paragraph.
\end{quotation}
Another frequently-displayed structure is a list.
The following is an example of an \emph{itemized}
list.
\begin{itemize}
   \item This is the first item of an itemized list.
         Each item in the list is marked with a ``tick''.
         You don't have to worry about what kind of tick
         mark is used.

   \item This is the second item of the list.  It
         contains another list nested inside it.  The inner
         list is an \emph{enumerated} list.
         \begin{enumerate}
            \item This is the first item of an enumerated
                  list that is nested within the itemized list.

            \item This is the second item of the inner list.
                  \LaTeX\ allows you to nest lists deeper than
                  you really should.
         \end{enumerate}
         This is the rest of the second item of the outer
         list.  It is no more interesting than any other
         part of the item.
   \item This is the third item of the list.
\end{itemize}
You can even display poetry.
\begin{verse}
   There is an environment
    for verse \\             % The \\ command separates lines
   Whose features some poets % within a stanza.
   will curse.

                             % One or more blank lines separate stanzas.

   For instead of making\\
   Them do \emph{all} line breaking, \\
   It allows them to put too many words on a line when they'd rather be
   forced to be terse.
\end{verse}

Mathematical formulas may also be displayed.  A
displayed formula
is
one-line long; multiline
formulas require special formatting instructions.
   \[  \ip{\Gamma}{\psi'} = x'' + y^{2} + z_{i}^{n}\]
Don't start a paragraph with a displayed equation,
nor make one a paragraph by itself.

\end{document}               % End of document.

\section{Bench Huge Section 969}
% synthetic bench section 969
% This is a sample LaTeX input file.  (Version of 12 August 2004.)
%
% A '%' character causes TeX to ignore all remaining text on the line,
% and is used for comments like this one.

\documentclass{article}      % Specifies the document class

                             % The preamble begins here.
\title{An Example Document}  % Declares the document's title.
\author{Leslie Lamport}      % Declares the author's name.
\date{January 21, 1994}      % Deleting this command produces today's date.

\newcommand{\ip}[2]{(#1, #2)}
                             % Defines \ip{arg1}{arg2} to mean
                             % (arg1, arg2).

%\newcommand{\ip}[2]{\langle #1 | #2\rangle}
                             % This is an alternative definition of
                             % \ip that is commented out.

\begin{document}             % End of preamble and beginning of text.

\maketitle                   % Produces the title.

This is an example input file.  Comparing it with
the output it generates can show you how to
produce a simple document of your own.

\section{Ordinary Text}      % Produces section heading.  Lower-level
                             % sections are begun with similar
                             % \subsection and \subsubsection commands.

The ends  of words and sentences are marked
  by   spaces. It  doesn't matter how many
spaces    you type; one is as good as 100.  The
end of   a line counts as a space.

One   or more   blank lines denote the  end
of  a paragraph.

Since any number of consecutive spaces are treated
like a single one, the formatting of the input
file makes no difference to
      \LaTeX,                % The \LaTeX command generates the LaTeX logo.
but it makes a difference to you.  When you use
\LaTeX, making your input file as easy to read
as possible will be a great help as you write
your document and when you change it.  This sample
file shows how you can add comments to your own input
file.

Because printing is different from typewriting,
there are a number of things that you have to do
differently when preparing an input file than if
you were just typing the document directly.
Quotation marks like
       ``this''
have to be handled specially, as do quotes within
quotes:
       ``\,`this'            % \, separates the double and single quote.
        is what I just
        wrote, not  `that'\,''.

Dashes come in three sizes: an
       intra-word
dash, a medium dash for number ranges like
       1--2,
and a punctuation
       dash---like
this.

A sentence-ending space should be larger than the
space between words within a sentence.  You
sometimes have to type special commands in
conjunction with punctuation characters to get
this right, as in the following sentence.
       Gnats, gnus, etc.\ all  % `\ ' makes an inter-word space.
       begin with G\@.         % \@ marks end-of-sentence punctuation.
You should check the spaces after periods when
reading your output to make sure you haven't
forgotten any special cases.  Generating an
ellipsis
       \ldots\               % `\ ' is needed after `\ldots' because TeX
                             % ignores spaces after command names like \ldots
                             % made from \ + letters.
                             %
                             % Note how a `%' character causes TeX to ignore
                             % the end of the input line, so these blank lines
                             % do not start a new paragraph.
                             %
with the right spacing around the periods requires
a special command.

\LaTeX\ interprets some common characters as
commands, so you must type special commands to
generate them.  These characters include the
following:
       \$ \& \% \# \{ and \}.

In printing, text is usually emphasized with an
       \emph{italic}
type style.

\begin{em}
   A long segment of text can also be emphasized
   in this way.  Text within such a segment can be
   given \emph{additional} emphasis.
\end{em}

It is sometimes necessary to prevent \LaTeX\ from
breaking a line where it might otherwise do so.
This may be at a space, as between the ``Mr.''\ and
``Jones'' in
       ``Mr.~Jones'',        % ~ produces an unbreakable interword space.
or within a word---especially when the word is a
symbol like
       \mbox{\emph{itemnum}}
that makes little sense when hyphenated across
lines.

Footnotes\footnote{This is an example of a footnote.}
pose no problem.

\LaTeX\ is good at typesetting mathematical formulas
like
       \( x-3y + z = 7 \)
or
       \( a_{1} > x^{2n} + y^{2n} > x' \)
or
       \( \ip{A}{B} = \sum_{i} a_{i} b_{i} \).
The spaces you type in a formula are
ignored.  Remember that a letter like
       $x$                   % $ ... $  and  \( ... \)  are equivalent
is a formula when it denotes a mathematical
symbol, and it should be typed as one.

\section{Displayed Text}

Text is displayed by indenting it from the left
margin.  Quotations are commonly displayed.  There
are short quotations
\begin{quote}
   This is a short quotation.  It consists of a
   single paragraph of text.  See how it is formatted.
\end{quote}
and longer ones.
\begin{quotation}
   This is a longer quotation.  It consists of two
   paragraphs of text, neither of which are
   particularly interesting.

   This is the second paragraph of the quotation.  It
   is just as dull as the first paragraph.
\end{quotation}
Another frequently-displayed structure is a list.
The following is an example of an \emph{itemized}
list.
\begin{itemize}
   \item This is the first item of an itemized list.
         Each item in the list is marked with a ``tick''.
         You don't have to worry about what kind of tick
         mark is used.

   \item This is the second item of the list.  It
         contains another list nested inside it.  The inner
         list is an \emph{enumerated} list.
         \begin{enumerate}
            \item This is the first item of an enumerated
                  list that is nested within the itemized list.

            \item This is the second item of the inner list.
                  \LaTeX\ allows you to nest lists deeper than
                  you really should.
         \end{enumerate}
         This is the rest of the second item of the outer
         list.  It is no more interesting than any other
         part of the item.
   \item This is the third item of the list.
\end{itemize}
You can even display poetry.
\begin{verse}
   There is an environment
    for verse \\             % The \\ command separates lines
   Whose features some poets % within a stanza.
   will curse.

                             % One or more blank lines separate stanzas.

   For instead of making\\
   Them do \emph{all} line breaking, \\
   It allows them to put too many words on a line when they'd rather be
   forced to be terse.
\end{verse}

Mathematical formulas may also be displayed.  A
displayed formula
is
one-line long; multiline
formulas require special formatting instructions.
   \[  \ip{\Gamma}{\psi'} = x'' + y^{2} + z_{i}^{n}\]
Don't start a paragraph with a displayed equation,
nor make one a paragraph by itself.

\end{document}               % End of document.

\section{Bench Huge Section 970}
% synthetic bench section 970
% This is a sample LaTeX input file.  (Version of 12 August 2004.)
%
% A '%' character causes TeX to ignore all remaining text on the line,
% and is used for comments like this one.

\documentclass{article}      % Specifies the document class

                             % The preamble begins here.
\title{An Example Document}  % Declares the document's title.
\author{Leslie Lamport}      % Declares the author's name.
\date{January 21, 1994}      % Deleting this command produces today's date.

\newcommand{\ip}[2]{(#1, #2)}
                             % Defines \ip{arg1}{arg2} to mean
                             % (arg1, arg2).

%\newcommand{\ip}[2]{\langle #1 | #2\rangle}
                             % This is an alternative definition of
                             % \ip that is commented out.

\begin{document}             % End of preamble and beginning of text.

\maketitle                   % Produces the title.

This is an example input file.  Comparing it with
the output it generates can show you how to
produce a simple document of your own.

\section{Ordinary Text}      % Produces section heading.  Lower-level
                             % sections are begun with similar
                             % \subsection and \subsubsection commands.

The ends  of words and sentences are marked
  by   spaces. It  doesn't matter how many
spaces    you type; one is as good as 100.  The
end of   a line counts as a space.

One   or more   blank lines denote the  end
of  a paragraph.

Since any number of consecutive spaces are treated
like a single one, the formatting of the input
file makes no difference to
      \LaTeX,                % The \LaTeX command generates the LaTeX logo.
but it makes a difference to you.  When you use
\LaTeX, making your input file as easy to read
as possible will be a great help as you write
your document and when you change it.  This sample
file shows how you can add comments to your own input
file.

Because printing is different from typewriting,
there are a number of things that you have to do
differently when preparing an input file than if
you were just typing the document directly.
Quotation marks like
       ``this''
have to be handled specially, as do quotes within
quotes:
       ``\,`this'            % \, separates the double and single quote.
        is what I just
        wrote, not  `that'\,''.

Dashes come in three sizes: an
       intra-word
dash, a medium dash for number ranges like
       1--2,
and a punctuation
       dash---like
this.

A sentence-ending space should be larger than the
space between words within a sentence.  You
sometimes have to type special commands in
conjunction with punctuation characters to get
this right, as in the following sentence.
       Gnats, gnus, etc.\ all  % `\ ' makes an inter-word space.
       begin with G\@.         % \@ marks end-of-sentence punctuation.
You should check the spaces after periods when
reading your output to make sure you haven't
forgotten any special cases.  Generating an
ellipsis
       \ldots\               % `\ ' is needed after `\ldots' because TeX
                             % ignores spaces after command names like \ldots
                             % made from \ + letters.
                             %
                             % Note how a `%' character causes TeX to ignore
                             % the end of the input line, so these blank lines
                             % do not start a new paragraph.
                             %
with the right spacing around the periods requires
a special command.

\LaTeX\ interprets some common characters as
commands, so you must type special commands to
generate them.  These characters include the
following:
       \$ \& \% \# \{ and \}.

In printing, text is usually emphasized with an
       \emph{italic}
type style.

\begin{em}
   A long segment of text can also be emphasized
   in this way.  Text within such a segment can be
   given \emph{additional} emphasis.
\end{em}

It is sometimes necessary to prevent \LaTeX\ from
breaking a line where it might otherwise do so.
This may be at a space, as between the ``Mr.''\ and
``Jones'' in
       ``Mr.~Jones'',        % ~ produces an unbreakable interword space.
or within a word---especially when the word is a
symbol like
       \mbox{\emph{itemnum}}
that makes little sense when hyphenated across
lines.

Footnotes\footnote{This is an example of a footnote.}
pose no problem.

\LaTeX\ is good at typesetting mathematical formulas
like
       \( x-3y + z = 7 \)
or
       \( a_{1} > x^{2n} + y^{2n} > x' \)
or
       \( \ip{A}{B} = \sum_{i} a_{i} b_{i} \).
The spaces you type in a formula are
ignored.  Remember that a letter like
       $x$                   % $ ... $  and  \( ... \)  are equivalent
is a formula when it denotes a mathematical
symbol, and it should be typed as one.

\section{Displayed Text}

Text is displayed by indenting it from the left
margin.  Quotations are commonly displayed.  There
are short quotations
\begin{quote}
   This is a short quotation.  It consists of a
   single paragraph of text.  See how it is formatted.
\end{quote}
and longer ones.
\begin{quotation}
   This is a longer quotation.  It consists of two
   paragraphs of text, neither of which are
   particularly interesting.

   This is the second paragraph of the quotation.  It
   is just as dull as the first paragraph.
\end{quotation}
Another frequently-displayed structure is a list.
The following is an example of an \emph{itemized}
list.
\begin{itemize}
   \item This is the first item of an itemized list.
         Each item in the list is marked with a ``tick''.
         You don't have to worry about what kind of tick
         mark is used.

   \item This is the second item of the list.  It
         contains another list nested inside it.  The inner
         list is an \emph{enumerated} list.
         \begin{enumerate}
            \item This is the first item of an enumerated
                  list that is nested within the itemized list.

            \item This is the second item of the inner list.
                  \LaTeX\ allows you to nest lists deeper than
                  you really should.
         \end{enumerate}
         This is the rest of the second item of the outer
         list.  It is no more interesting than any other
         part of the item.
   \item This is the third item of the list.
\end{itemize}
You can even display poetry.
\begin{verse}
   There is an environment
    for verse \\             % The \\ command separates lines
   Whose features some poets % within a stanza.
   will curse.

                             % One or more blank lines separate stanzas.

   For instead of making\\
   Them do \emph{all} line breaking, \\
   It allows them to put too many words on a line when they'd rather be
   forced to be terse.
\end{verse}

Mathematical formulas may also be displayed.  A
displayed formula
is
one-line long; multiline
formulas require special formatting instructions.
   \[  \ip{\Gamma}{\psi'} = x'' + y^{2} + z_{i}^{n}\]
Don't start a paragraph with a displayed equation,
nor make one a paragraph by itself.

\end{document}               % End of document.

\section{Bench Huge Section 971}
% synthetic bench section 971
% This is a sample LaTeX input file.  (Version of 12 August 2004.)
%
% A '%' character causes TeX to ignore all remaining text on the line,
% and is used for comments like this one.

\documentclass{article}      % Specifies the document class

                             % The preamble begins here.
\title{An Example Document}  % Declares the document's title.
\author{Leslie Lamport}      % Declares the author's name.
\date{January 21, 1994}      % Deleting this command produces today's date.

\newcommand{\ip}[2]{(#1, #2)}
                             % Defines \ip{arg1}{arg2} to mean
                             % (arg1, arg2).

%\newcommand{\ip}[2]{\langle #1 | #2\rangle}
                             % This is an alternative definition of
                             % \ip that is commented out.

\begin{document}             % End of preamble and beginning of text.

\maketitle                   % Produces the title.

This is an example input file.  Comparing it with
the output it generates can show you how to
produce a simple document of your own.

\section{Ordinary Text}      % Produces section heading.  Lower-level
                             % sections are begun with similar
                             % \subsection and \subsubsection commands.

The ends  of words and sentences are marked
  by   spaces. It  doesn't matter how many
spaces    you type; one is as good as 100.  The
end of   a line counts as a space.

One   or more   blank lines denote the  end
of  a paragraph.

Since any number of consecutive spaces are treated
like a single one, the formatting of the input
file makes no difference to
      \LaTeX,                % The \LaTeX command generates the LaTeX logo.
but it makes a difference to you.  When you use
\LaTeX, making your input file as easy to read
as possible will be a great help as you write
your document and when you change it.  This sample
file shows how you can add comments to your own input
file.

Because printing is different from typewriting,
there are a number of things that you have to do
differently when preparing an input file than if
you were just typing the document directly.
Quotation marks like
       ``this''
have to be handled specially, as do quotes within
quotes:
       ``\,`this'            % \, separates the double and single quote.
        is what I just
        wrote, not  `that'\,''.

Dashes come in three sizes: an
       intra-word
dash, a medium dash for number ranges like
       1--2,
and a punctuation
       dash---like
this.

A sentence-ending space should be larger than the
space between words within a sentence.  You
sometimes have to type special commands in
conjunction with punctuation characters to get
this right, as in the following sentence.
       Gnats, gnus, etc.\ all  % `\ ' makes an inter-word space.
       begin with G\@.         % \@ marks end-of-sentence punctuation.
You should check the spaces after periods when
reading your output to make sure you haven't
forgotten any special cases.  Generating an
ellipsis
       \ldots\               % `\ ' is needed after `\ldots' because TeX
                             % ignores spaces after command names like \ldots
                             % made from \ + letters.
                             %
                             % Note how a `%' character causes TeX to ignore
                             % the end of the input line, so these blank lines
                             % do not start a new paragraph.
                             %
with the right spacing around the periods requires
a special command.

\LaTeX\ interprets some common characters as
commands, so you must type special commands to
generate them.  These characters include the
following:
       \$ \& \% \# \{ and \}.

In printing, text is usually emphasized with an
       \emph{italic}
type style.

\begin{em}
   A long segment of text can also be emphasized
   in this way.  Text within such a segment can be
   given \emph{additional} emphasis.
\end{em}

It is sometimes necessary to prevent \LaTeX\ from
breaking a line where it might otherwise do so.
This may be at a space, as between the ``Mr.''\ and
``Jones'' in
       ``Mr.~Jones'',        % ~ produces an unbreakable interword space.
or within a word---especially when the word is a
symbol like
       \mbox{\emph{itemnum}}
that makes little sense when hyphenated across
lines.

Footnotes\footnote{This is an example of a footnote.}
pose no problem.

\LaTeX\ is good at typesetting mathematical formulas
like
       \( x-3y + z = 7 \)
or
       \( a_{1} > x^{2n} + y^{2n} > x' \)
or
       \( \ip{A}{B} = \sum_{i} a_{i} b_{i} \).
The spaces you type in a formula are
ignored.  Remember that a letter like
       $x$                   % $ ... $  and  \( ... \)  are equivalent
is a formula when it denotes a mathematical
symbol, and it should be typed as one.

\section{Displayed Text}

Text is displayed by indenting it from the left
margin.  Quotations are commonly displayed.  There
are short quotations
\begin{quote}
   This is a short quotation.  It consists of a
   single paragraph of text.  See how it is formatted.
\end{quote}
and longer ones.
\begin{quotation}
   This is a longer quotation.  It consists of two
   paragraphs of text, neither of which are
   particularly interesting.

   This is the second paragraph of the quotation.  It
   is just as dull as the first paragraph.
\end{quotation}
Another frequently-displayed structure is a list.
The following is an example of an \emph{itemized}
list.
\begin{itemize}
   \item This is the first item of an itemized list.
         Each item in the list is marked with a ``tick''.
         You don't have to worry about what kind of tick
         mark is used.

   \item This is the second item of the list.  It
         contains another list nested inside it.  The inner
         list is an \emph{enumerated} list.
         \begin{enumerate}
            \item This is the first item of an enumerated
                  list that is nested within the itemized list.

            \item This is the second item of the inner list.
                  \LaTeX\ allows you to nest lists deeper than
                  you really should.
         \end{enumerate}
         This is the rest of the second item of the outer
         list.  It is no more interesting than any other
         part of the item.
   \item This is the third item of the list.
\end{itemize}
You can even display poetry.
\begin{verse}
   There is an environment
    for verse \\             % The \\ command separates lines
   Whose features some poets % within a stanza.
   will curse.

                             % One or more blank lines separate stanzas.

   For instead of making\\
   Them do \emph{all} line breaking, \\
   It allows them to put too many words on a line when they'd rather be
   forced to be terse.
\end{verse}

Mathematical formulas may also be displayed.  A
displayed formula
is
one-line long; multiline
formulas require special formatting instructions.
   \[  \ip{\Gamma}{\psi'} = x'' + y^{2} + z_{i}^{n}\]
Don't start a paragraph with a displayed equation,
nor make one a paragraph by itself.

\end{document}               % End of document.

\section{Bench Huge Section 972}
% synthetic bench section 972
% This is a sample LaTeX input file.  (Version of 12 August 2004.)
%
% A '%' character causes TeX to ignore all remaining text on the line,
% and is used for comments like this one.

\documentclass{article}      % Specifies the document class

                             % The preamble begins here.
\title{An Example Document}  % Declares the document's title.
\author{Leslie Lamport}      % Declares the author's name.
\date{January 21, 1994}      % Deleting this command produces today's date.

\newcommand{\ip}[2]{(#1, #2)}
                             % Defines \ip{arg1}{arg2} to mean
                             % (arg1, arg2).

%\newcommand{\ip}[2]{\langle #1 | #2\rangle}
                             % This is an alternative definition of
                             % \ip that is commented out.

\begin{document}             % End of preamble and beginning of text.

\maketitle                   % Produces the title.

This is an example input file.  Comparing it with
the output it generates can show you how to
produce a simple document of your own.

\section{Ordinary Text}      % Produces section heading.  Lower-level
                             % sections are begun with similar
                             % \subsection and \subsubsection commands.

The ends  of words and sentences are marked
  by   spaces. It  doesn't matter how many
spaces    you type; one is as good as 100.  The
end of   a line counts as a space.

One   or more   blank lines denote the  end
of  a paragraph.

Since any number of consecutive spaces are treated
like a single one, the formatting of the input
file makes no difference to
      \LaTeX,                % The \LaTeX command generates the LaTeX logo.
but it makes a difference to you.  When you use
\LaTeX, making your input file as easy to read
as possible will be a great help as you write
your document and when you change it.  This sample
file shows how you can add comments to your own input
file.

Because printing is different from typewriting,
there are a number of things that you have to do
differently when preparing an input file than if
you were just typing the document directly.
Quotation marks like
       ``this''
have to be handled specially, as do quotes within
quotes:
       ``\,`this'            % \, separates the double and single quote.
        is what I just
        wrote, not  `that'\,''.

Dashes come in three sizes: an
       intra-word
dash, a medium dash for number ranges like
       1--2,
and a punctuation
       dash---like
this.

A sentence-ending space should be larger than the
space between words within a sentence.  You
sometimes have to type special commands in
conjunction with punctuation characters to get
this right, as in the following sentence.
       Gnats, gnus, etc.\ all  % `\ ' makes an inter-word space.
       begin with G\@.         % \@ marks end-of-sentence punctuation.
You should check the spaces after periods when
reading your output to make sure you haven't
forgotten any special cases.  Generating an
ellipsis
       \ldots\               % `\ ' is needed after `\ldots' because TeX
                             % ignores spaces after command names like \ldots
                             % made from \ + letters.
                             %
                             % Note how a `%' character causes TeX to ignore
                             % the end of the input line, so these blank lines
                             % do not start a new paragraph.
                             %
with the right spacing around the periods requires
a special command.

\LaTeX\ interprets some common characters as
commands, so you must type special commands to
generate them.  These characters include the
following:
       \$ \& \% \# \{ and \}.

In printing, text is usually emphasized with an
       \emph{italic}
type style.

\begin{em}
   A long segment of text can also be emphasized
   in this way.  Text within such a segment can be
   given \emph{additional} emphasis.
\end{em}

It is sometimes necessary to prevent \LaTeX\ from
breaking a line where it might otherwise do so.
This may be at a space, as between the ``Mr.''\ and
``Jones'' in
       ``Mr.~Jones'',        % ~ produces an unbreakable interword space.
or within a word---especially when the word is a
symbol like
       \mbox{\emph{itemnum}}
that makes little sense when hyphenated across
lines.

Footnotes\footnote{This is an example of a footnote.}
pose no problem.

\LaTeX\ is good at typesetting mathematical formulas
like
       \( x-3y + z = 7 \)
or
       \( a_{1} > x^{2n} + y^{2n} > x' \)
or
       \( \ip{A}{B} = \sum_{i} a_{i} b_{i} \).
The spaces you type in a formula are
ignored.  Remember that a letter like
       $x$                   % $ ... $  and  \( ... \)  are equivalent
is a formula when it denotes a mathematical
symbol, and it should be typed as one.

\section{Displayed Text}

Text is displayed by indenting it from the left
margin.  Quotations are commonly displayed.  There
are short quotations
\begin{quote}
   This is a short quotation.  It consists of a
   single paragraph of text.  See how it is formatted.
\end{quote}
and longer ones.
\begin{quotation}
   This is a longer quotation.  It consists of two
   paragraphs of text, neither of which are
   particularly interesting.

   This is the second paragraph of the quotation.  It
   is just as dull as the first paragraph.
\end{quotation}
Another frequently-displayed structure is a list.
The following is an example of an \emph{itemized}
list.
\begin{itemize}
   \item This is the first item of an itemized list.
         Each item in the list is marked with a ``tick''.
         You don't have to worry about what kind of tick
         mark is used.

   \item This is the second item of the list.  It
         contains another list nested inside it.  The inner
         list is an \emph{enumerated} list.
         \begin{enumerate}
            \item This is the first item of an enumerated
                  list that is nested within the itemized list.

            \item This is the second item of the inner list.
                  \LaTeX\ allows you to nest lists deeper than
                  you really should.
         \end{enumerate}
         This is the rest of the second item of the outer
         list.  It is no more interesting than any other
         part of the item.
   \item This is the third item of the list.
\end{itemize}
You can even display poetry.
\begin{verse}
   There is an environment
    for verse \\             % The \\ command separates lines
   Whose features some poets % within a stanza.
   will curse.

                             % One or more blank lines separate stanzas.

   For instead of making\\
   Them do \emph{all} line breaking, \\
   It allows them to put too many words on a line when they'd rather be
   forced to be terse.
\end{verse}

Mathematical formulas may also be displayed.  A
displayed formula
is
one-line long; multiline
formulas require special formatting instructions.
   \[  \ip{\Gamma}{\psi'} = x'' + y^{2} + z_{i}^{n}\]
Don't start a paragraph with a displayed equation,
nor make one a paragraph by itself.

\end{document}               % End of document.

\section{Bench Huge Section 973}
% synthetic bench section 973
% This is a sample LaTeX input file.  (Version of 12 August 2004.)
%
% A '%' character causes TeX to ignore all remaining text on the line,
% and is used for comments like this one.

\documentclass{article}      % Specifies the document class

                             % The preamble begins here.
\title{An Example Document}  % Declares the document's title.
\author{Leslie Lamport}      % Declares the author's name.
\date{January 21, 1994}      % Deleting this command produces today's date.

\newcommand{\ip}[2]{(#1, #2)}
                             % Defines \ip{arg1}{arg2} to mean
                             % (arg1, arg2).

%\newcommand{\ip}[2]{\langle #1 | #2\rangle}
                             % This is an alternative definition of
                             % \ip that is commented out.

\begin{document}             % End of preamble and beginning of text.

\maketitle                   % Produces the title.

This is an example input file.  Comparing it with
the output it generates can show you how to
produce a simple document of your own.

\section{Ordinary Text}      % Produces section heading.  Lower-level
                             % sections are begun with similar
                             % \subsection and \subsubsection commands.

The ends  of words and sentences are marked
  by   spaces. It  doesn't matter how many
spaces    you type; one is as good as 100.  The
end of   a line counts as a space.

One   or more   blank lines denote the  end
of  a paragraph.

Since any number of consecutive spaces are treated
like a single one, the formatting of the input
file makes no difference to
      \LaTeX,                % The \LaTeX command generates the LaTeX logo.
but it makes a difference to you.  When you use
\LaTeX, making your input file as easy to read
as possible will be a great help as you write
your document and when you change it.  This sample
file shows how you can add comments to your own input
file.

Because printing is different from typewriting,
there are a number of things that you have to do
differently when preparing an input file than if
you were just typing the document directly.
Quotation marks like
       ``this''
have to be handled specially, as do quotes within
quotes:
       ``\,`this'            % \, separates the double and single quote.
        is what I just
        wrote, not  `that'\,''.

Dashes come in three sizes: an
       intra-word
dash, a medium dash for number ranges like
       1--2,
and a punctuation
       dash---like
this.

A sentence-ending space should be larger than the
space between words within a sentence.  You
sometimes have to type special commands in
conjunction with punctuation characters to get
this right, as in the following sentence.
       Gnats, gnus, etc.\ all  % `\ ' makes an inter-word space.
       begin with G\@.         % \@ marks end-of-sentence punctuation.
You should check the spaces after periods when
reading your output to make sure you haven't
forgotten any special cases.  Generating an
ellipsis
       \ldots\               % `\ ' is needed after `\ldots' because TeX
                             % ignores spaces after command names like \ldots
                             % made from \ + letters.
                             %
                             % Note how a `%' character causes TeX to ignore
                             % the end of the input line, so these blank lines
                             % do not start a new paragraph.
                             %
with the right spacing around the periods requires
a special command.

\LaTeX\ interprets some common characters as
commands, so you must type special commands to
generate them.  These characters include the
following:
       \$ \& \% \# \{ and \}.

In printing, text is usually emphasized with an
       \emph{italic}
type style.

\begin{em}
   A long segment of text can also be emphasized
   in this way.  Text within such a segment can be
   given \emph{additional} emphasis.
\end{em}

It is sometimes necessary to prevent \LaTeX\ from
breaking a line where it might otherwise do so.
This may be at a space, as between the ``Mr.''\ and
``Jones'' in
       ``Mr.~Jones'',        % ~ produces an unbreakable interword space.
or within a word---especially when the word is a
symbol like
       \mbox{\emph{itemnum}}
that makes little sense when hyphenated across
lines.

Footnotes\footnote{This is an example of a footnote.}
pose no problem.

\LaTeX\ is good at typesetting mathematical formulas
like
       \( x-3y + z = 7 \)
or
       \( a_{1} > x^{2n} + y^{2n} > x' \)
or
       \( \ip{A}{B} = \sum_{i} a_{i} b_{i} \).
The spaces you type in a formula are
ignored.  Remember that a letter like
       $x$                   % $ ... $  and  \( ... \)  are equivalent
is a formula when it denotes a mathematical
symbol, and it should be typed as one.

\section{Displayed Text}

Text is displayed by indenting it from the left
margin.  Quotations are commonly displayed.  There
are short quotations
\begin{quote}
   This is a short quotation.  It consists of a
   single paragraph of text.  See how it is formatted.
\end{quote}
and longer ones.
\begin{quotation}
   This is a longer quotation.  It consists of two
   paragraphs of text, neither of which are
   particularly interesting.

   This is the second paragraph of the quotation.  It
   is just as dull as the first paragraph.
\end{quotation}
Another frequently-displayed structure is a list.
The following is an example of an \emph{itemized}
list.
\begin{itemize}
   \item This is the first item of an itemized list.
         Each item in the list is marked with a ``tick''.
         You don't have to worry about what kind of tick
         mark is used.

   \item This is the second item of the list.  It
         contains another list nested inside it.  The inner
         list is an \emph{enumerated} list.
         \begin{enumerate}
            \item This is the first item of an enumerated
                  list that is nested within the itemized list.

            \item This is the second item of the inner list.
                  \LaTeX\ allows you to nest lists deeper than
                  you really should.
         \end{enumerate}
         This is the rest of the second item of the outer
         list.  It is no more interesting than any other
         part of the item.
   \item This is the third item of the list.
\end{itemize}
You can even display poetry.
\begin{verse}
   There is an environment
    for verse \\             % The \\ command separates lines
   Whose features some poets % within a stanza.
   will curse.

                             % One or more blank lines separate stanzas.

   For instead of making\\
   Them do \emph{all} line breaking, \\
   It allows them to put too many words on a line when they'd rather be
   forced to be terse.
\end{verse}

Mathematical formulas may also be displayed.  A
displayed formula
is
one-line long; multiline
formulas require special formatting instructions.
   \[  \ip{\Gamma}{\psi'} = x'' + y^{2} + z_{i}^{n}\]
Don't start a paragraph with a displayed equation,
nor make one a paragraph by itself.

\end{document}               % End of document.

\section{Bench Huge Section 974}
% synthetic bench section 974
% This is a sample LaTeX input file.  (Version of 12 August 2004.)
%
% A '%' character causes TeX to ignore all remaining text on the line,
% and is used for comments like this one.

\documentclass{article}      % Specifies the document class

                             % The preamble begins here.
\title{An Example Document}  % Declares the document's title.
\author{Leslie Lamport}      % Declares the author's name.
\date{January 21, 1994}      % Deleting this command produces today's date.

\newcommand{\ip}[2]{(#1, #2)}
                             % Defines \ip{arg1}{arg2} to mean
                             % (arg1, arg2).

%\newcommand{\ip}[2]{\langle #1 | #2\rangle}
                             % This is an alternative definition of
                             % \ip that is commented out.

\begin{document}             % End of preamble and beginning of text.

\maketitle                   % Produces the title.

This is an example input file.  Comparing it with
the output it generates can show you how to
produce a simple document of your own.

\section{Ordinary Text}      % Produces section heading.  Lower-level
                             % sections are begun with similar
                             % \subsection and \subsubsection commands.

The ends  of words and sentences are marked
  by   spaces. It  doesn't matter how many
spaces    you type; one is as good as 100.  The
end of   a line counts as a space.

One   or more   blank lines denote the  end
of  a paragraph.

Since any number of consecutive spaces are treated
like a single one, the formatting of the input
file makes no difference to
      \LaTeX,                % The \LaTeX command generates the LaTeX logo.
but it makes a difference to you.  When you use
\LaTeX, making your input file as easy to read
as possible will be a great help as you write
your document and when you change it.  This sample
file shows how you can add comments to your own input
file.

Because printing is different from typewriting,
there are a number of things that you have to do
differently when preparing an input file than if
you were just typing the document directly.
Quotation marks like
       ``this''
have to be handled specially, as do quotes within
quotes:
       ``\,`this'            % \, separates the double and single quote.
        is what I just
        wrote, not  `that'\,''.

Dashes come in three sizes: an
       intra-word
dash, a medium dash for number ranges like
       1--2,
and a punctuation
       dash---like
this.

A sentence-ending space should be larger than the
space between words within a sentence.  You
sometimes have to type special commands in
conjunction with punctuation characters to get
this right, as in the following sentence.
       Gnats, gnus, etc.\ all  % `\ ' makes an inter-word space.
       begin with G\@.         % \@ marks end-of-sentence punctuation.
You should check the spaces after periods when
reading your output to make sure you haven't
forgotten any special cases.  Generating an
ellipsis
       \ldots\               % `\ ' is needed after `\ldots' because TeX
                             % ignores spaces after command names like \ldots
                             % made from \ + letters.
                             %
                             % Note how a `%' character causes TeX to ignore
                             % the end of the input line, so these blank lines
                             % do not start a new paragraph.
                             %
with the right spacing around the periods requires
a special command.

\LaTeX\ interprets some common characters as
commands, so you must type special commands to
generate them.  These characters include the
following:
       \$ \& \% \# \{ and \}.

In printing, text is usually emphasized with an
       \emph{italic}
type style.

\begin{em}
   A long segment of text can also be emphasized
   in this way.  Text within such a segment can be
   given \emph{additional} emphasis.
\end{em}

It is sometimes necessary to prevent \LaTeX\ from
breaking a line where it might otherwise do so.
This may be at a space, as between the ``Mr.''\ and
``Jones'' in
       ``Mr.~Jones'',        % ~ produces an unbreakable interword space.
or within a word---especially when the word is a
symbol like
       \mbox{\emph{itemnum}}
that makes little sense when hyphenated across
lines.

Footnotes\footnote{This is an example of a footnote.}
pose no problem.

\LaTeX\ is good at typesetting mathematical formulas
like
       \( x-3y + z = 7 \)
or
       \( a_{1} > x^{2n} + y^{2n} > x' \)
or
       \( \ip{A}{B} = \sum_{i} a_{i} b_{i} \).
The spaces you type in a formula are
ignored.  Remember that a letter like
       $x$                   % $ ... $  and  \( ... \)  are equivalent
is a formula when it denotes a mathematical
symbol, and it should be typed as one.

\section{Displayed Text}

Text is displayed by indenting it from the left
margin.  Quotations are commonly displayed.  There
are short quotations
\begin{quote}
   This is a short quotation.  It consists of a
   single paragraph of text.  See how it is formatted.
\end{quote}
and longer ones.
\begin{quotation}
   This is a longer quotation.  It consists of two
   paragraphs of text, neither of which are
   particularly interesting.

   This is the second paragraph of the quotation.  It
   is just as dull as the first paragraph.
\end{quotation}
Another frequently-displayed structure is a list.
The following is an example of an \emph{itemized}
list.
\begin{itemize}
   \item This is the first item of an itemized list.
         Each item in the list is marked with a ``tick''.
         You don't have to worry about what kind of tick
         mark is used.

   \item This is the second item of the list.  It
         contains another list nested inside it.  The inner
         list is an \emph{enumerated} list.
         \begin{enumerate}
            \item This is the first item of an enumerated
                  list that is nested within the itemized list.

            \item This is the second item of the inner list.
                  \LaTeX\ allows you to nest lists deeper than
                  you really should.
         \end{enumerate}
         This is the rest of the second item of the outer
         list.  It is no more interesting than any other
         part of the item.
   \item This is the third item of the list.
\end{itemize}
You can even display poetry.
\begin{verse}
   There is an environment
    for verse \\             % The \\ command separates lines
   Whose features some poets % within a stanza.
   will curse.

                             % One or more blank lines separate stanzas.

   For instead of making\\
   Them do \emph{all} line breaking, \\
   It allows them to put too many words on a line when they'd rather be
   forced to be terse.
\end{verse}

Mathematical formulas may also be displayed.  A
displayed formula
is
one-line long; multiline
formulas require special formatting instructions.
   \[  \ip{\Gamma}{\psi'} = x'' + y^{2} + z_{i}^{n}\]
Don't start a paragraph with a displayed equation,
nor make one a paragraph by itself.

\end{document}               % End of document.

\section{Bench Huge Section 975}
% synthetic bench section 975
% This is a sample LaTeX input file.  (Version of 12 August 2004.)
%
% A '%' character causes TeX to ignore all remaining text on the line,
% and is used for comments like this one.

\documentclass{article}      % Specifies the document class

                             % The preamble begins here.
\title{An Example Document}  % Declares the document's title.
\author{Leslie Lamport}      % Declares the author's name.
\date{January 21, 1994}      % Deleting this command produces today's date.

\newcommand{\ip}[2]{(#1, #2)}
                             % Defines \ip{arg1}{arg2} to mean
                             % (arg1, arg2).

%\newcommand{\ip}[2]{\langle #1 | #2\rangle}
                             % This is an alternative definition of
                             % \ip that is commented out.

\begin{document}             % End of preamble and beginning of text.

\maketitle                   % Produces the title.

This is an example input file.  Comparing it with
the output it generates can show you how to
produce a simple document of your own.

\section{Ordinary Text}      % Produces section heading.  Lower-level
                             % sections are begun with similar
                             % \subsection and \subsubsection commands.

The ends  of words and sentences are marked
  by   spaces. It  doesn't matter how many
spaces    you type; one is as good as 100.  The
end of   a line counts as a space.

One   or more   blank lines denote the  end
of  a paragraph.

Since any number of consecutive spaces are treated
like a single one, the formatting of the input
file makes no difference to
      \LaTeX,                % The \LaTeX command generates the LaTeX logo.
but it makes a difference to you.  When you use
\LaTeX, making your input file as easy to read
as possible will be a great help as you write
your document and when you change it.  This sample
file shows how you can add comments to your own input
file.

Because printing is different from typewriting,
there are a number of things that you have to do
differently when preparing an input file than if
you were just typing the document directly.
Quotation marks like
       ``this''
have to be handled specially, as do quotes within
quotes:
       ``\,`this'            % \, separates the double and single quote.
        is what I just
        wrote, not  `that'\,''.

Dashes come in three sizes: an
       intra-word
dash, a medium dash for number ranges like
       1--2,
and a punctuation
       dash---like
this.

A sentence-ending space should be larger than the
space between words within a sentence.  You
sometimes have to type special commands in
conjunction with punctuation characters to get
this right, as in the following sentence.
       Gnats, gnus, etc.\ all  % `\ ' makes an inter-word space.
       begin with G\@.         % \@ marks end-of-sentence punctuation.
You should check the spaces after periods when
reading your output to make sure you haven't
forgotten any special cases.  Generating an
ellipsis
       \ldots\               % `\ ' is needed after `\ldots' because TeX
                             % ignores spaces after command names like \ldots
                             % made from \ + letters.
                             %
                             % Note how a `%' character causes TeX to ignore
                             % the end of the input line, so these blank lines
                             % do not start a new paragraph.
                             %
with the right spacing around the periods requires
a special command.

\LaTeX\ interprets some common characters as
commands, so you must type special commands to
generate them.  These characters include the
following:
       \$ \& \% \# \{ and \}.

In printing, text is usually emphasized with an
       \emph{italic}
type style.

\begin{em}
   A long segment of text can also be emphasized
   in this way.  Text within such a segment can be
   given \emph{additional} emphasis.
\end{em}

It is sometimes necessary to prevent \LaTeX\ from
breaking a line where it might otherwise do so.
This may be at a space, as between the ``Mr.''\ and
``Jones'' in
       ``Mr.~Jones'',        % ~ produces an unbreakable interword space.
or within a word---especially when the word is a
symbol like
       \mbox{\emph{itemnum}}
that makes little sense when hyphenated across
lines.

Footnotes\footnote{This is an example of a footnote.}
pose no problem.

\LaTeX\ is good at typesetting mathematical formulas
like
       \( x-3y + z = 7 \)
or
       \( a_{1} > x^{2n} + y^{2n} > x' \)
or
       \( \ip{A}{B} = \sum_{i} a_{i} b_{i} \).
The spaces you type in a formula are
ignored.  Remember that a letter like
       $x$                   % $ ... $  and  \( ... \)  are equivalent
is a formula when it denotes a mathematical
symbol, and it should be typed as one.

\section{Displayed Text}

Text is displayed by indenting it from the left
margin.  Quotations are commonly displayed.  There
are short quotations
\begin{quote}
   This is a short quotation.  It consists of a
   single paragraph of text.  See how it is formatted.
\end{quote}
and longer ones.
\begin{quotation}
   This is a longer quotation.  It consists of two
   paragraphs of text, neither of which are
   particularly interesting.

   This is the second paragraph of the quotation.  It
   is just as dull as the first paragraph.
\end{quotation}
Another frequently-displayed structure is a list.
The following is an example of an \emph{itemized}
list.
\begin{itemize}
   \item This is the first item of an itemized list.
         Each item in the list is marked with a ``tick''.
         You don't have to worry about what kind of tick
         mark is used.

   \item This is the second item of the list.  It
         contains another list nested inside it.  The inner
         list is an \emph{enumerated} list.
         \begin{enumerate}
            \item This is the first item of an enumerated
                  list that is nested within the itemized list.

            \item This is the second item of the inner list.
                  \LaTeX\ allows you to nest lists deeper than
                  you really should.
         \end{enumerate}
         This is the rest of the second item of the outer
         list.  It is no more interesting than any other
         part of the item.
   \item This is the third item of the list.
\end{itemize}
You can even display poetry.
\begin{verse}
   There is an environment
    for verse \\             % The \\ command separates lines
   Whose features some poets % within a stanza.
   will curse.

                             % One or more blank lines separate stanzas.

   For instead of making\\
   Them do \emph{all} line breaking, \\
   It allows them to put too many words on a line when they'd rather be
   forced to be terse.
\end{verse}

Mathematical formulas may also be displayed.  A
displayed formula
is
one-line long; multiline
formulas require special formatting instructions.
   \[  \ip{\Gamma}{\psi'} = x'' + y^{2} + z_{i}^{n}\]
Don't start a paragraph with a displayed equation,
nor make one a paragraph by itself.

\end{document}               % End of document.

\section{Bench Huge Section 976}
% synthetic bench section 976
% This is a sample LaTeX input file.  (Version of 12 August 2004.)
%
% A '%' character causes TeX to ignore all remaining text on the line,
% and is used for comments like this one.

\documentclass{article}      % Specifies the document class

                             % The preamble begins here.
\title{An Example Document}  % Declares the document's title.
\author{Leslie Lamport}      % Declares the author's name.
\date{January 21, 1994}      % Deleting this command produces today's date.

\newcommand{\ip}[2]{(#1, #2)}
                             % Defines \ip{arg1}{arg2} to mean
                             % (arg1, arg2).

%\newcommand{\ip}[2]{\langle #1 | #2\rangle}
                             % This is an alternative definition of
                             % \ip that is commented out.

\begin{document}             % End of preamble and beginning of text.

\maketitle                   % Produces the title.

This is an example input file.  Comparing it with
the output it generates can show you how to
produce a simple document of your own.

\section{Ordinary Text}      % Produces section heading.  Lower-level
                             % sections are begun with similar
                             % \subsection and \subsubsection commands.

The ends  of words and sentences are marked
  by   spaces. It  doesn't matter how many
spaces    you type; one is as good as 100.  The
end of   a line counts as a space.

One   or more   blank lines denote the  end
of  a paragraph.

Since any number of consecutive spaces are treated
like a single one, the formatting of the input
file makes no difference to
      \LaTeX,                % The \LaTeX command generates the LaTeX logo.
but it makes a difference to you.  When you use
\LaTeX, making your input file as easy to read
as possible will be a great help as you write
your document and when you change it.  This sample
file shows how you can add comments to your own input
file.

Because printing is different from typewriting,
there are a number of things that you have to do
differently when preparing an input file than if
you were just typing the document directly.
Quotation marks like
       ``this''
have to be handled specially, as do quotes within
quotes:
       ``\,`this'            % \, separates the double and single quote.
        is what I just
        wrote, not  `that'\,''.

Dashes come in three sizes: an
       intra-word
dash, a medium dash for number ranges like
       1--2,
and a punctuation
       dash---like
this.

A sentence-ending space should be larger than the
space between words within a sentence.  You
sometimes have to type special commands in
conjunction with punctuation characters to get
this right, as in the following sentence.
       Gnats, gnus, etc.\ all  % `\ ' makes an inter-word space.
       begin with G\@.         % \@ marks end-of-sentence punctuation.
You should check the spaces after periods when
reading your output to make sure you haven't
forgotten any special cases.  Generating an
ellipsis
       \ldots\               % `\ ' is needed after `\ldots' because TeX
                             % ignores spaces after command names like \ldots
                             % made from \ + letters.
                             %
                             % Note how a `%' character causes TeX to ignore
                             % the end of the input line, so these blank lines
                             % do not start a new paragraph.
                             %
with the right spacing around the periods requires
a special command.

\LaTeX\ interprets some common characters as
commands, so you must type special commands to
generate them.  These characters include the
following:
       \$ \& \% \# \{ and \}.

In printing, text is usually emphasized with an
       \emph{italic}
type style.

\begin{em}
   A long segment of text can also be emphasized
   in this way.  Text within such a segment can be
   given \emph{additional} emphasis.
\end{em}

It is sometimes necessary to prevent \LaTeX\ from
breaking a line where it might otherwise do so.
This may be at a space, as between the ``Mr.''\ and
``Jones'' in
       ``Mr.~Jones'',        % ~ produces an unbreakable interword space.
or within a word---especially when the word is a
symbol like
       \mbox{\emph{itemnum}}
that makes little sense when hyphenated across
lines.

Footnotes\footnote{This is an example of a footnote.}
pose no problem.

\LaTeX\ is good at typesetting mathematical formulas
like
       \( x-3y + z = 7 \)
or
       \( a_{1} > x^{2n} + y^{2n} > x' \)
or
       \( \ip{A}{B} = \sum_{i} a_{i} b_{i} \).
The spaces you type in a formula are
ignored.  Remember that a letter like
       $x$                   % $ ... $  and  \( ... \)  are equivalent
is a formula when it denotes a mathematical
symbol, and it should be typed as one.

\section{Displayed Text}

Text is displayed by indenting it from the left
margin.  Quotations are commonly displayed.  There
are short quotations
\begin{quote}
   This is a short quotation.  It consists of a
   single paragraph of text.  See how it is formatted.
\end{quote}
and longer ones.
\begin{quotation}
   This is a longer quotation.  It consists of two
   paragraphs of text, neither of which are
   particularly interesting.

   This is the second paragraph of the quotation.  It
   is just as dull as the first paragraph.
\end{quotation}
Another frequently-displayed structure is a list.
The following is an example of an \emph{itemized}
list.
\begin{itemize}
   \item This is the first item of an itemized list.
         Each item in the list is marked with a ``tick''.
         You don't have to worry about what kind of tick
         mark is used.

   \item This is the second item of the list.  It
         contains another list nested inside it.  The inner
         list is an \emph{enumerated} list.
         \begin{enumerate}
            \item This is the first item of an enumerated
                  list that is nested within the itemized list.

            \item This is the second item of the inner list.
                  \LaTeX\ allows you to nest lists deeper than
                  you really should.
         \end{enumerate}
         This is the rest of the second item of the outer
         list.  It is no more interesting than any other
         part of the item.
   \item This is the third item of the list.
\end{itemize}
You can even display poetry.
\begin{verse}
   There is an environment
    for verse \\             % The \\ command separates lines
   Whose features some poets % within a stanza.
   will curse.

                             % One or more blank lines separate stanzas.

   For instead of making\\
   Them do \emph{all} line breaking, \\
   It allows them to put too many words on a line when they'd rather be
   forced to be terse.
\end{verse}

Mathematical formulas may also be displayed.  A
displayed formula
is
one-line long; multiline
formulas require special formatting instructions.
   \[  \ip{\Gamma}{\psi'} = x'' + y^{2} + z_{i}^{n}\]
Don't start a paragraph with a displayed equation,
nor make one a paragraph by itself.

\end{document}               % End of document.

\section{Bench Huge Section 977}
% synthetic bench section 977
% This is a sample LaTeX input file.  (Version of 12 August 2004.)
%
% A '%' character causes TeX to ignore all remaining text on the line,
% and is used for comments like this one.

\documentclass{article}      % Specifies the document class

                             % The preamble begins here.
\title{An Example Document}  % Declares the document's title.
\author{Leslie Lamport}      % Declares the author's name.
\date{January 21, 1994}      % Deleting this command produces today's date.

\newcommand{\ip}[2]{(#1, #2)}
                             % Defines \ip{arg1}{arg2} to mean
                             % (arg1, arg2).

%\newcommand{\ip}[2]{\langle #1 | #2\rangle}
                             % This is an alternative definition of
                             % \ip that is commented out.

\begin{document}             % End of preamble and beginning of text.

\maketitle                   % Produces the title.

This is an example input file.  Comparing it with
the output it generates can show you how to
produce a simple document of your own.

\section{Ordinary Text}      % Produces section heading.  Lower-level
                             % sections are begun with similar
                             % \subsection and \subsubsection commands.

The ends  of words and sentences are marked
  by   spaces. It  doesn't matter how many
spaces    you type; one is as good as 100.  The
end of   a line counts as a space.

One   or more   blank lines denote the  end
of  a paragraph.

Since any number of consecutive spaces are treated
like a single one, the formatting of the input
file makes no difference to
      \LaTeX,                % The \LaTeX command generates the LaTeX logo.
but it makes a difference to you.  When you use
\LaTeX, making your input file as easy to read
as possible will be a great help as you write
your document and when you change it.  This sample
file shows how you can add comments to your own input
file.

Because printing is different from typewriting,
there are a number of things that you have to do
differently when preparing an input file than if
you were just typing the document directly.
Quotation marks like
       ``this''
have to be handled specially, as do quotes within
quotes:
       ``\,`this'            % \, separates the double and single quote.
        is what I just
        wrote, not  `that'\,''.

Dashes come in three sizes: an
       intra-word
dash, a medium dash for number ranges like
       1--2,
and a punctuation
       dash---like
this.

A sentence-ending space should be larger than the
space between words within a sentence.  You
sometimes have to type special commands in
conjunction with punctuation characters to get
this right, as in the following sentence.
       Gnats, gnus, etc.\ all  % `\ ' makes an inter-word space.
       begin with G\@.         % \@ marks end-of-sentence punctuation.
You should check the spaces after periods when
reading your output to make sure you haven't
forgotten any special cases.  Generating an
ellipsis
       \ldots\               % `\ ' is needed after `\ldots' because TeX
                             % ignores spaces after command names like \ldots
                             % made from \ + letters.
                             %
                             % Note how a `%' character causes TeX to ignore
                             % the end of the input line, so these blank lines
                             % do not start a new paragraph.
                             %
with the right spacing around the periods requires
a special command.

\LaTeX\ interprets some common characters as
commands, so you must type special commands to
generate them.  These characters include the
following:
       \$ \& \% \# \{ and \}.

In printing, text is usually emphasized with an
       \emph{italic}
type style.

\begin{em}
   A long segment of text can also be emphasized
   in this way.  Text within such a segment can be
   given \emph{additional} emphasis.
\end{em}

It is sometimes necessary to prevent \LaTeX\ from
breaking a line where it might otherwise do so.
This may be at a space, as between the ``Mr.''\ and
``Jones'' in
       ``Mr.~Jones'',        % ~ produces an unbreakable interword space.
or within a word---especially when the word is a
symbol like
       \mbox{\emph{itemnum}}
that makes little sense when hyphenated across
lines.

Footnotes\footnote{This is an example of a footnote.}
pose no problem.

\LaTeX\ is good at typesetting mathematical formulas
like
       \( x-3y + z = 7 \)
or
       \( a_{1} > x^{2n} + y^{2n} > x' \)
or
       \( \ip{A}{B} = \sum_{i} a_{i} b_{i} \).
The spaces you type in a formula are
ignored.  Remember that a letter like
       $x$                   % $ ... $  and  \( ... \)  are equivalent
is a formula when it denotes a mathematical
symbol, and it should be typed as one.

\section{Displayed Text}

Text is displayed by indenting it from the left
margin.  Quotations are commonly displayed.  There
are short quotations
\begin{quote}
   This is a short quotation.  It consists of a
   single paragraph of text.  See how it is formatted.
\end{quote}
and longer ones.
\begin{quotation}
   This is a longer quotation.  It consists of two
   paragraphs of text, neither of which are
   particularly interesting.

   This is the second paragraph of the quotation.  It
   is just as dull as the first paragraph.
\end{quotation}
Another frequently-displayed structure is a list.
The following is an example of an \emph{itemized}
list.
\begin{itemize}
   \item This is the first item of an itemized list.
         Each item in the list is marked with a ``tick''.
         You don't have to worry about what kind of tick
         mark is used.

   \item This is the second item of the list.  It
         contains another list nested inside it.  The inner
         list is an \emph{enumerated} list.
         \begin{enumerate}
            \item This is the first item of an enumerated
                  list that is nested within the itemized list.

            \item This is the second item of the inner list.
                  \LaTeX\ allows you to nest lists deeper than
                  you really should.
         \end{enumerate}
         This is the rest of the second item of the outer
         list.  It is no more interesting than any other
         part of the item.
   \item This is the third item of the list.
\end{itemize}
You can even display poetry.
\begin{verse}
   There is an environment
    for verse \\             % The \\ command separates lines
   Whose features some poets % within a stanza.
   will curse.

                             % One or more blank lines separate stanzas.

   For instead of making\\
   Them do \emph{all} line breaking, \\
   It allows them to put too many words on a line when they'd rather be
   forced to be terse.
\end{verse}

Mathematical formulas may also be displayed.  A
displayed formula
is
one-line long; multiline
formulas require special formatting instructions.
   \[  \ip{\Gamma}{\psi'} = x'' + y^{2} + z_{i}^{n}\]
Don't start a paragraph with a displayed equation,
nor make one a paragraph by itself.

\end{document}               % End of document.

\section{Bench Huge Section 978}
% synthetic bench section 978
% This is a sample LaTeX input file.  (Version of 12 August 2004.)
%
% A '%' character causes TeX to ignore all remaining text on the line,
% and is used for comments like this one.

\documentclass{article}      % Specifies the document class

                             % The preamble begins here.
\title{An Example Document}  % Declares the document's title.
\author{Leslie Lamport}      % Declares the author's name.
\date{January 21, 1994}      % Deleting this command produces today's date.

\newcommand{\ip}[2]{(#1, #2)}
                             % Defines \ip{arg1}{arg2} to mean
                             % (arg1, arg2).

%\newcommand{\ip}[2]{\langle #1 | #2\rangle}
                             % This is an alternative definition of
                             % \ip that is commented out.

\begin{document}             % End of preamble and beginning of text.

\maketitle                   % Produces the title.

This is an example input file.  Comparing it with
the output it generates can show you how to
produce a simple document of your own.

\section{Ordinary Text}      % Produces section heading.  Lower-level
                             % sections are begun with similar
                             % \subsection and \subsubsection commands.

The ends  of words and sentences are marked
  by   spaces. It  doesn't matter how many
spaces    you type; one is as good as 100.  The
end of   a line counts as a space.

One   or more   blank lines denote the  end
of  a paragraph.

Since any number of consecutive spaces are treated
like a single one, the formatting of the input
file makes no difference to
      \LaTeX,                % The \LaTeX command generates the LaTeX logo.
but it makes a difference to you.  When you use
\LaTeX, making your input file as easy to read
as possible will be a great help as you write
your document and when you change it.  This sample
file shows how you can add comments to your own input
file.

Because printing is different from typewriting,
there are a number of things that you have to do
differently when preparing an input file than if
you were just typing the document directly.
Quotation marks like
       ``this''
have to be handled specially, as do quotes within
quotes:
       ``\,`this'            % \, separates the double and single quote.
        is what I just
        wrote, not  `that'\,''.

Dashes come in three sizes: an
       intra-word
dash, a medium dash for number ranges like
       1--2,
and a punctuation
       dash---like
this.

A sentence-ending space should be larger than the
space between words within a sentence.  You
sometimes have to type special commands in
conjunction with punctuation characters to get
this right, as in the following sentence.
       Gnats, gnus, etc.\ all  % `\ ' makes an inter-word space.
       begin with G\@.         % \@ marks end-of-sentence punctuation.
You should check the spaces after periods when
reading your output to make sure you haven't
forgotten any special cases.  Generating an
ellipsis
       \ldots\               % `\ ' is needed after `\ldots' because TeX
                             % ignores spaces after command names like \ldots
                             % made from \ + letters.
                             %
                             % Note how a `%' character causes TeX to ignore
                             % the end of the input line, so these blank lines
                             % do not start a new paragraph.
                             %
with the right spacing around the periods requires
a special command.

\LaTeX\ interprets some common characters as
commands, so you must type special commands to
generate them.  These characters include the
following:
       \$ \& \% \# \{ and \}.

In printing, text is usually emphasized with an
       \emph{italic}
type style.

\begin{em}
   A long segment of text can also be emphasized
   in this way.  Text within such a segment can be
   given \emph{additional} emphasis.
\end{em}

It is sometimes necessary to prevent \LaTeX\ from
breaking a line where it might otherwise do so.
This may be at a space, as between the ``Mr.''\ and
``Jones'' in
       ``Mr.~Jones'',        % ~ produces an unbreakable interword space.
or within a word---especially when the word is a
symbol like
       \mbox{\emph{itemnum}}
that makes little sense when hyphenated across
lines.

Footnotes\footnote{This is an example of a footnote.}
pose no problem.

\LaTeX\ is good at typesetting mathematical formulas
like
       \( x-3y + z = 7 \)
or
       \( a_{1} > x^{2n} + y^{2n} > x' \)
or
       \( \ip{A}{B} = \sum_{i} a_{i} b_{i} \).
The spaces you type in a formula are
ignored.  Remember that a letter like
       $x$                   % $ ... $  and  \( ... \)  are equivalent
is a formula when it denotes a mathematical
symbol, and it should be typed as one.

\section{Displayed Text}

Text is displayed by indenting it from the left
margin.  Quotations are commonly displayed.  There
are short quotations
\begin{quote}
   This is a short quotation.  It consists of a
   single paragraph of text.  See how it is formatted.
\end{quote}
and longer ones.
\begin{quotation}
   This is a longer quotation.  It consists of two
   paragraphs of text, neither of which are
   particularly interesting.

   This is the second paragraph of the quotation.  It
   is just as dull as the first paragraph.
\end{quotation}
Another frequently-displayed structure is a list.
The following is an example of an \emph{itemized}
list.
\begin{itemize}
   \item This is the first item of an itemized list.
         Each item in the list is marked with a ``tick''.
         You don't have to worry about what kind of tick
         mark is used.

   \item This is the second item of the list.  It
         contains another list nested inside it.  The inner
         list is an \emph{enumerated} list.
         \begin{enumerate}
            \item This is the first item of an enumerated
                  list that is nested within the itemized list.

            \item This is the second item of the inner list.
                  \LaTeX\ allows you to nest lists deeper than
                  you really should.
         \end{enumerate}
         This is the rest of the second item of the outer
         list.  It is no more interesting than any other
         part of the item.
   \item This is the third item of the list.
\end{itemize}
You can even display poetry.
\begin{verse}
   There is an environment
    for verse \\             % The \\ command separates lines
   Whose features some poets % within a stanza.
   will curse.

                             % One or more blank lines separate stanzas.

   For instead of making\\
   Them do \emph{all} line breaking, \\
   It allows them to put too many words on a line when they'd rather be
   forced to be terse.
\end{verse}

Mathematical formulas may also be displayed.  A
displayed formula
is
one-line long; multiline
formulas require special formatting instructions.
   \[  \ip{\Gamma}{\psi'} = x'' + y^{2} + z_{i}^{n}\]
Don't start a paragraph with a displayed equation,
nor make one a paragraph by itself.

\end{document}               % End of document.

\section{Bench Huge Section 979}
% synthetic bench section 979
% This is a sample LaTeX input file.  (Version of 12 August 2004.)
%
% A '%' character causes TeX to ignore all remaining text on the line,
% and is used for comments like this one.

\documentclass{article}      % Specifies the document class

                             % The preamble begins here.
\title{An Example Document}  % Declares the document's title.
\author{Leslie Lamport}      % Declares the author's name.
\date{January 21, 1994}      % Deleting this command produces today's date.

\newcommand{\ip}[2]{(#1, #2)}
                             % Defines \ip{arg1}{arg2} to mean
                             % (arg1, arg2).

%\newcommand{\ip}[2]{\langle #1 | #2\rangle}
                             % This is an alternative definition of
                             % \ip that is commented out.

\begin{document}             % End of preamble and beginning of text.

\maketitle                   % Produces the title.

This is an example input file.  Comparing it with
the output it generates can show you how to
produce a simple document of your own.

\section{Ordinary Text}      % Produces section heading.  Lower-level
                             % sections are begun with similar
                             % \subsection and \subsubsection commands.

The ends  of words and sentences are marked
  by   spaces. It  doesn't matter how many
spaces    you type; one is as good as 100.  The
end of   a line counts as a space.

One   or more   blank lines denote the  end
of  a paragraph.

Since any number of consecutive spaces are treated
like a single one, the formatting of the input
file makes no difference to
      \LaTeX,                % The \LaTeX command generates the LaTeX logo.
but it makes a difference to you.  When you use
\LaTeX, making your input file as easy to read
as possible will be a great help as you write
your document and when you change it.  This sample
file shows how you can add comments to your own input
file.

Because printing is different from typewriting,
there are a number of things that you have to do
differently when preparing an input file than if
you were just typing the document directly.
Quotation marks like
       ``this''
have to be handled specially, as do quotes within
quotes:
       ``\,`this'            % \, separates the double and single quote.
        is what I just
        wrote, not  `that'\,''.

Dashes come in three sizes: an
       intra-word
dash, a medium dash for number ranges like
       1--2,
and a punctuation
       dash---like
this.

A sentence-ending space should be larger than the
space between words within a sentence.  You
sometimes have to type special commands in
conjunction with punctuation characters to get
this right, as in the following sentence.
       Gnats, gnus, etc.\ all  % `\ ' makes an inter-word space.
       begin with G\@.         % \@ marks end-of-sentence punctuation.
You should check the spaces after periods when
reading your output to make sure you haven't
forgotten any special cases.  Generating an
ellipsis
       \ldots\               % `\ ' is needed after `\ldots' because TeX
                             % ignores spaces after command names like \ldots
                             % made from \ + letters.
                             %
                             % Note how a `%' character causes TeX to ignore
                             % the end of the input line, so these blank lines
                             % do not start a new paragraph.
                             %
with the right spacing around the periods requires
a special command.

\LaTeX\ interprets some common characters as
commands, so you must type special commands to
generate them.  These characters include the
following:
       \$ \& \% \# \{ and \}.

In printing, text is usually emphasized with an
       \emph{italic}
type style.

\begin{em}
   A long segment of text can also be emphasized
   in this way.  Text within such a segment can be
   given \emph{additional} emphasis.
\end{em}

It is sometimes necessary to prevent \LaTeX\ from
breaking a line where it might otherwise do so.
This may be at a space, as between the ``Mr.''\ and
``Jones'' in
       ``Mr.~Jones'',        % ~ produces an unbreakable interword space.
or within a word---especially when the word is a
symbol like
       \mbox{\emph{itemnum}}
that makes little sense when hyphenated across
lines.

Footnotes\footnote{This is an example of a footnote.}
pose no problem.

\LaTeX\ is good at typesetting mathematical formulas
like
       \( x-3y + z = 7 \)
or
       \( a_{1} > x^{2n} + y^{2n} > x' \)
or
       \( \ip{A}{B} = \sum_{i} a_{i} b_{i} \).
The spaces you type in a formula are
ignored.  Remember that a letter like
       $x$                   % $ ... $  and  \( ... \)  are equivalent
is a formula when it denotes a mathematical
symbol, and it should be typed as one.

\section{Displayed Text}

Text is displayed by indenting it from the left
margin.  Quotations are commonly displayed.  There
are short quotations
\begin{quote}
   This is a short quotation.  It consists of a
   single paragraph of text.  See how it is formatted.
\end{quote}
and longer ones.
\begin{quotation}
   This is a longer quotation.  It consists of two
   paragraphs of text, neither of which are
   particularly interesting.

   This is the second paragraph of the quotation.  It
   is just as dull as the first paragraph.
\end{quotation}
Another frequently-displayed structure is a list.
The following is an example of an \emph{itemized}
list.
\begin{itemize}
   \item This is the first item of an itemized list.
         Each item in the list is marked with a ``tick''.
         You don't have to worry about what kind of tick
         mark is used.

   \item This is the second item of the list.  It
         contains another list nested inside it.  The inner
         list is an \emph{enumerated} list.
         \begin{enumerate}
            \item This is the first item of an enumerated
                  list that is nested within the itemized list.

            \item This is the second item of the inner list.
                  \LaTeX\ allows you to nest lists deeper than
                  you really should.
         \end{enumerate}
         This is the rest of the second item of the outer
         list.  It is no more interesting than any other
         part of the item.
   \item This is the third item of the list.
\end{itemize}
You can even display poetry.
\begin{verse}
   There is an environment
    for verse \\             % The \\ command separates lines
   Whose features some poets % within a stanza.
   will curse.

                             % One or more blank lines separate stanzas.

   For instead of making\\
   Them do \emph{all} line breaking, \\
   It allows them to put too many words on a line when they'd rather be
   forced to be terse.
\end{verse}

Mathematical formulas may also be displayed.  A
displayed formula
is
one-line long; multiline
formulas require special formatting instructions.
   \[  \ip{\Gamma}{\psi'} = x'' + y^{2} + z_{i}^{n}\]
Don't start a paragraph with a displayed equation,
nor make one a paragraph by itself.

\end{document}               % End of document.

\section{Bench Huge Section 980}
% synthetic bench section 980
% This is a sample LaTeX input file.  (Version of 12 August 2004.)
%
% A '%' character causes TeX to ignore all remaining text on the line,
% and is used for comments like this one.

\documentclass{article}      % Specifies the document class

                             % The preamble begins here.
\title{An Example Document}  % Declares the document's title.
\author{Leslie Lamport}      % Declares the author's name.
\date{January 21, 1994}      % Deleting this command produces today's date.

\newcommand{\ip}[2]{(#1, #2)}
                             % Defines \ip{arg1}{arg2} to mean
                             % (arg1, arg2).

%\newcommand{\ip}[2]{\langle #1 | #2\rangle}
                             % This is an alternative definition of
                             % \ip that is commented out.

\begin{document}             % End of preamble and beginning of text.

\maketitle                   % Produces the title.

This is an example input file.  Comparing it with
the output it generates can show you how to
produce a simple document of your own.

\section{Ordinary Text}      % Produces section heading.  Lower-level
                             % sections are begun with similar
                             % \subsection and \subsubsection commands.

The ends  of words and sentences are marked
  by   spaces. It  doesn't matter how many
spaces    you type; one is as good as 100.  The
end of   a line counts as a space.

One   or more   blank lines denote the  end
of  a paragraph.

Since any number of consecutive spaces are treated
like a single one, the formatting of the input
file makes no difference to
      \LaTeX,                % The \LaTeX command generates the LaTeX logo.
but it makes a difference to you.  When you use
\LaTeX, making your input file as easy to read
as possible will be a great help as you write
your document and when you change it.  This sample
file shows how you can add comments to your own input
file.

Because printing is different from typewriting,
there are a number of things that you have to do
differently when preparing an input file than if
you were just typing the document directly.
Quotation marks like
       ``this''
have to be handled specially, as do quotes within
quotes:
       ``\,`this'            % \, separates the double and single quote.
        is what I just
        wrote, not  `that'\,''.

Dashes come in three sizes: an
       intra-word
dash, a medium dash for number ranges like
       1--2,
and a punctuation
       dash---like
this.

A sentence-ending space should be larger than the
space between words within a sentence.  You
sometimes have to type special commands in
conjunction with punctuation characters to get
this right, as in the following sentence.
       Gnats, gnus, etc.\ all  % `\ ' makes an inter-word space.
       begin with G\@.         % \@ marks end-of-sentence punctuation.
You should check the spaces after periods when
reading your output to make sure you haven't
forgotten any special cases.  Generating an
ellipsis
       \ldots\               % `\ ' is needed after `\ldots' because TeX
                             % ignores spaces after command names like \ldots
                             % made from \ + letters.
                             %
                             % Note how a `%' character causes TeX to ignore
                             % the end of the input line, so these blank lines
                             % do not start a new paragraph.
                             %
with the right spacing around the periods requires
a special command.

\LaTeX\ interprets some common characters as
commands, so you must type special commands to
generate them.  These characters include the
following:
       \$ \& \% \# \{ and \}.

In printing, text is usually emphasized with an
       \emph{italic}
type style.

\begin{em}
   A long segment of text can also be emphasized
   in this way.  Text within such a segment can be
   given \emph{additional} emphasis.
\end{em}

It is sometimes necessary to prevent \LaTeX\ from
breaking a line where it might otherwise do so.
This may be at a space, as between the ``Mr.''\ and
``Jones'' in
       ``Mr.~Jones'',        % ~ produces an unbreakable interword space.
or within a word---especially when the word is a
symbol like
       \mbox{\emph{itemnum}}
that makes little sense when hyphenated across
lines.

Footnotes\footnote{This is an example of a footnote.}
pose no problem.

\LaTeX\ is good at typesetting mathematical formulas
like
       \( x-3y + z = 7 \)
or
       \( a_{1} > x^{2n} + y^{2n} > x' \)
or
       \( \ip{A}{B} = \sum_{i} a_{i} b_{i} \).
The spaces you type in a formula are
ignored.  Remember that a letter like
       $x$                   % $ ... $  and  \( ... \)  are equivalent
is a formula when it denotes a mathematical
symbol, and it should be typed as one.

\section{Displayed Text}

Text is displayed by indenting it from the left
margin.  Quotations are commonly displayed.  There
are short quotations
\begin{quote}
   This is a short quotation.  It consists of a
   single paragraph of text.  See how it is formatted.
\end{quote}
and longer ones.
\begin{quotation}
   This is a longer quotation.  It consists of two
   paragraphs of text, neither of which are
   particularly interesting.

   This is the second paragraph of the quotation.  It
   is just as dull as the first paragraph.
\end{quotation}
Another frequently-displayed structure is a list.
The following is an example of an \emph{itemized}
list.
\begin{itemize}
   \item This is the first item of an itemized list.
         Each item in the list is marked with a ``tick''.
         You don't have to worry about what kind of tick
         mark is used.

   \item This is the second item of the list.  It
         contains another list nested inside it.  The inner
         list is an \emph{enumerated} list.
         \begin{enumerate}
            \item This is the first item of an enumerated
                  list that is nested within the itemized list.

            \item This is the second item of the inner list.
                  \LaTeX\ allows you to nest lists deeper than
                  you really should.
         \end{enumerate}
         This is the rest of the second item of the outer
         list.  It is no more interesting than any other
         part of the item.
   \item This is the third item of the list.
\end{itemize}
You can even display poetry.
\begin{verse}
   There is an environment
    for verse \\             % The \\ command separates lines
   Whose features some poets % within a stanza.
   will curse.

                             % One or more blank lines separate stanzas.

   For instead of making\\
   Them do \emph{all} line breaking, \\
   It allows them to put too many words on a line when they'd rather be
   forced to be terse.
\end{verse}

Mathematical formulas may also be displayed.  A
displayed formula
is
one-line long; multiline
formulas require special formatting instructions.
   \[  \ip{\Gamma}{\psi'} = x'' + y^{2} + z_{i}^{n}\]
Don't start a paragraph with a displayed equation,
nor make one a paragraph by itself.

\end{document}               % End of document.

\section{Bench Huge Section 981}
% synthetic bench section 981
% This is a sample LaTeX input file.  (Version of 12 August 2004.)
%
% A '%' character causes TeX to ignore all remaining text on the line,
% and is used for comments like this one.

\documentclass{article}      % Specifies the document class

                             % The preamble begins here.
\title{An Example Document}  % Declares the document's title.
\author{Leslie Lamport}      % Declares the author's name.
\date{January 21, 1994}      % Deleting this command produces today's date.

\newcommand{\ip}[2]{(#1, #2)}
                             % Defines \ip{arg1}{arg2} to mean
                             % (arg1, arg2).

%\newcommand{\ip}[2]{\langle #1 | #2\rangle}
                             % This is an alternative definition of
                             % \ip that is commented out.

\begin{document}             % End of preamble and beginning of text.

\maketitle                   % Produces the title.

This is an example input file.  Comparing it with
the output it generates can show you how to
produce a simple document of your own.

\section{Ordinary Text}      % Produces section heading.  Lower-level
                             % sections are begun with similar
                             % \subsection and \subsubsection commands.

The ends  of words and sentences are marked
  by   spaces. It  doesn't matter how many
spaces    you type; one is as good as 100.  The
end of   a line counts as a space.

One   or more   blank lines denote the  end
of  a paragraph.

Since any number of consecutive spaces are treated
like a single one, the formatting of the input
file makes no difference to
      \LaTeX,                % The \LaTeX command generates the LaTeX logo.
but it makes a difference to you.  When you use
\LaTeX, making your input file as easy to read
as possible will be a great help as you write
your document and when you change it.  This sample
file shows how you can add comments to your own input
file.

Because printing is different from typewriting,
there are a number of things that you have to do
differently when preparing an input file than if
you were just typing the document directly.
Quotation marks like
       ``this''
have to be handled specially, as do quotes within
quotes:
       ``\,`this'            % \, separates the double and single quote.
        is what I just
        wrote, not  `that'\,''.

Dashes come in three sizes: an
       intra-word
dash, a medium dash for number ranges like
       1--2,
and a punctuation
       dash---like
this.

A sentence-ending space should be larger than the
space between words within a sentence.  You
sometimes have to type special commands in
conjunction with punctuation characters to get
this right, as in the following sentence.
       Gnats, gnus, etc.\ all  % `\ ' makes an inter-word space.
       begin with G\@.         % \@ marks end-of-sentence punctuation.
You should check the spaces after periods when
reading your output to make sure you haven't
forgotten any special cases.  Generating an
ellipsis
       \ldots\               % `\ ' is needed after `\ldots' because TeX
                             % ignores spaces after command names like \ldots
                             % made from \ + letters.
                             %
                             % Note how a `%' character causes TeX to ignore
                             % the end of the input line, so these blank lines
                             % do not start a new paragraph.
                             %
with the right spacing around the periods requires
a special command.

\LaTeX\ interprets some common characters as
commands, so you must type special commands to
generate them.  These characters include the
following:
       \$ \& \% \# \{ and \}.

In printing, text is usually emphasized with an
       \emph{italic}
type style.

\begin{em}
   A long segment of text can also be emphasized
   in this way.  Text within such a segment can be
   given \emph{additional} emphasis.
\end{em}

It is sometimes necessary to prevent \LaTeX\ from
breaking a line where it might otherwise do so.
This may be at a space, as between the ``Mr.''\ and
``Jones'' in
       ``Mr.~Jones'',        % ~ produces an unbreakable interword space.
or within a word---especially when the word is a
symbol like
       \mbox{\emph{itemnum}}
that makes little sense when hyphenated across
lines.

Footnotes\footnote{This is an example of a footnote.}
pose no problem.

\LaTeX\ is good at typesetting mathematical formulas
like
       \( x-3y + z = 7 \)
or
       \( a_{1} > x^{2n} + y^{2n} > x' \)
or
       \( \ip{A}{B} = \sum_{i} a_{i} b_{i} \).
The spaces you type in a formula are
ignored.  Remember that a letter like
       $x$                   % $ ... $  and  \( ... \)  are equivalent
is a formula when it denotes a mathematical
symbol, and it should be typed as one.

\section{Displayed Text}

Text is displayed by indenting it from the left
margin.  Quotations are commonly displayed.  There
are short quotations
\begin{quote}
   This is a short quotation.  It consists of a
   single paragraph of text.  See how it is formatted.
\end{quote}
and longer ones.
\begin{quotation}
   This is a longer quotation.  It consists of two
   paragraphs of text, neither of which are
   particularly interesting.

   This is the second paragraph of the quotation.  It
   is just as dull as the first paragraph.
\end{quotation}
Another frequently-displayed structure is a list.
The following is an example of an \emph{itemized}
list.
\begin{itemize}
   \item This is the first item of an itemized list.
         Each item in the list is marked with a ``tick''.
         You don't have to worry about what kind of tick
         mark is used.

   \item This is the second item of the list.  It
         contains another list nested inside it.  The inner
         list is an \emph{enumerated} list.
         \begin{enumerate}
            \item This is the first item of an enumerated
                  list that is nested within the itemized list.

            \item This is the second item of the inner list.
                  \LaTeX\ allows you to nest lists deeper than
                  you really should.
         \end{enumerate}
         This is the rest of the second item of the outer
         list.  It is no more interesting than any other
         part of the item.
   \item This is the third item of the list.
\end{itemize}
You can even display poetry.
\begin{verse}
   There is an environment
    for verse \\             % The \\ command separates lines
   Whose features some poets % within a stanza.
   will curse.

                             % One or more blank lines separate stanzas.

   For instead of making\\
   Them do \emph{all} line breaking, \\
   It allows them to put too many words on a line when they'd rather be
   forced to be terse.
\end{verse}

Mathematical formulas may also be displayed.  A
displayed formula
is
one-line long; multiline
formulas require special formatting instructions.
   \[  \ip{\Gamma}{\psi'} = x'' + y^{2} + z_{i}^{n}\]
Don't start a paragraph with a displayed equation,
nor make one a paragraph by itself.

\end{document}               % End of document.

\section{Bench Huge Section 982}
% synthetic bench section 982
% This is a sample LaTeX input file.  (Version of 12 August 2004.)
%
% A '%' character causes TeX to ignore all remaining text on the line,
% and is used for comments like this one.

\documentclass{article}      % Specifies the document class

                             % The preamble begins here.
\title{An Example Document}  % Declares the document's title.
\author{Leslie Lamport}      % Declares the author's name.
\date{January 21, 1994}      % Deleting this command produces today's date.

\newcommand{\ip}[2]{(#1, #2)}
                             % Defines \ip{arg1}{arg2} to mean
                             % (arg1, arg2).

%\newcommand{\ip}[2]{\langle #1 | #2\rangle}
                             % This is an alternative definition of
                             % \ip that is commented out.

\begin{document}             % End of preamble and beginning of text.

\maketitle                   % Produces the title.

This is an example input file.  Comparing it with
the output it generates can show you how to
produce a simple document of your own.

\section{Ordinary Text}      % Produces section heading.  Lower-level
                             % sections are begun with similar
                             % \subsection and \subsubsection commands.

The ends  of words and sentences are marked
  by   spaces. It  doesn't matter how many
spaces    you type; one is as good as 100.  The
end of   a line counts as a space.

One   or more   blank lines denote the  end
of  a paragraph.

Since any number of consecutive spaces are treated
like a single one, the formatting of the input
file makes no difference to
      \LaTeX,                % The \LaTeX command generates the LaTeX logo.
but it makes a difference to you.  When you use
\LaTeX, making your input file as easy to read
as possible will be a great help as you write
your document and when you change it.  This sample
file shows how you can add comments to your own input
file.

Because printing is different from typewriting,
there are a number of things that you have to do
differently when preparing an input file than if
you were just typing the document directly.
Quotation marks like
       ``this''
have to be handled specially, as do quotes within
quotes:
       ``\,`this'            % \, separates the double and single quote.
        is what I just
        wrote, not  `that'\,''.

Dashes come in three sizes: an
       intra-word
dash, a medium dash for number ranges like
       1--2,
and a punctuation
       dash---like
this.

A sentence-ending space should be larger than the
space between words within a sentence.  You
sometimes have to type special commands in
conjunction with punctuation characters to get
this right, as in the following sentence.
       Gnats, gnus, etc.\ all  % `\ ' makes an inter-word space.
       begin with G\@.         % \@ marks end-of-sentence punctuation.
You should check the spaces after periods when
reading your output to make sure you haven't
forgotten any special cases.  Generating an
ellipsis
       \ldots\               % `\ ' is needed after `\ldots' because TeX
                             % ignores spaces after command names like \ldots
                             % made from \ + letters.
                             %
                             % Note how a `%' character causes TeX to ignore
                             % the end of the input line, so these blank lines
                             % do not start a new paragraph.
                             %
with the right spacing around the periods requires
a special command.

\LaTeX\ interprets some common characters as
commands, so you must type special commands to
generate them.  These characters include the
following:
       \$ \& \% \# \{ and \}.

In printing, text is usually emphasized with an
       \emph{italic}
type style.

\begin{em}
   A long segment of text can also be emphasized
   in this way.  Text within such a segment can be
   given \emph{additional} emphasis.
\end{em}

It is sometimes necessary to prevent \LaTeX\ from
breaking a line where it might otherwise do so.
This may be at a space, as between the ``Mr.''\ and
``Jones'' in
       ``Mr.~Jones'',        % ~ produces an unbreakable interword space.
or within a word---especially when the word is a
symbol like
       \mbox{\emph{itemnum}}
that makes little sense when hyphenated across
lines.

Footnotes\footnote{This is an example of a footnote.}
pose no problem.

\LaTeX\ is good at typesetting mathematical formulas
like
       \( x-3y + z = 7 \)
or
       \( a_{1} > x^{2n} + y^{2n} > x' \)
or
       \( \ip{A}{B} = \sum_{i} a_{i} b_{i} \).
The spaces you type in a formula are
ignored.  Remember that a letter like
       $x$                   % $ ... $  and  \( ... \)  are equivalent
is a formula when it denotes a mathematical
symbol, and it should be typed as one.

\section{Displayed Text}

Text is displayed by indenting it from the left
margin.  Quotations are commonly displayed.  There
are short quotations
\begin{quote}
   This is a short quotation.  It consists of a
   single paragraph of text.  See how it is formatted.
\end{quote}
and longer ones.
\begin{quotation}
   This is a longer quotation.  It consists of two
   paragraphs of text, neither of which are
   particularly interesting.

   This is the second paragraph of the quotation.  It
   is just as dull as the first paragraph.
\end{quotation}
Another frequently-displayed structure is a list.
The following is an example of an \emph{itemized}
list.
\begin{itemize}
   \item This is the first item of an itemized list.
         Each item in the list is marked with a ``tick''.
         You don't have to worry about what kind of tick
         mark is used.

   \item This is the second item of the list.  It
         contains another list nested inside it.  The inner
         list is an \emph{enumerated} list.
         \begin{enumerate}
            \item This is the first item of an enumerated
                  list that is nested within the itemized list.

            \item This is the second item of the inner list.
                  \LaTeX\ allows you to nest lists deeper than
                  you really should.
         \end{enumerate}
         This is the rest of the second item of the outer
         list.  It is no more interesting than any other
         part of the item.
   \item This is the third item of the list.
\end{itemize}
You can even display poetry.
\begin{verse}
   There is an environment
    for verse \\             % The \\ command separates lines
   Whose features some poets % within a stanza.
   will curse.

                             % One or more blank lines separate stanzas.

   For instead of making\\
   Them do \emph{all} line breaking, \\
   It allows them to put too many words on a line when they'd rather be
   forced to be terse.
\end{verse}

Mathematical formulas may also be displayed.  A
displayed formula
is
one-line long; multiline
formulas require special formatting instructions.
   \[  \ip{\Gamma}{\psi'} = x'' + y^{2} + z_{i}^{n}\]
Don't start a paragraph with a displayed equation,
nor make one a paragraph by itself.

\end{document}               % End of document.

\section{Bench Huge Section 983}
% synthetic bench section 983
% This is a sample LaTeX input file.  (Version of 12 August 2004.)
%
% A '%' character causes TeX to ignore all remaining text on the line,
% and is used for comments like this one.

\documentclass{article}      % Specifies the document class

                             % The preamble begins here.
\title{An Example Document}  % Declares the document's title.
\author{Leslie Lamport}      % Declares the author's name.
\date{January 21, 1994}      % Deleting this command produces today's date.

\newcommand{\ip}[2]{(#1, #2)}
                             % Defines \ip{arg1}{arg2} to mean
                             % (arg1, arg2).

%\newcommand{\ip}[2]{\langle #1 | #2\rangle}
                             % This is an alternative definition of
                             % \ip that is commented out.

\begin{document}             % End of preamble and beginning of text.

\maketitle                   % Produces the title.

This is an example input file.  Comparing it with
the output it generates can show you how to
produce a simple document of your own.

\section{Ordinary Text}      % Produces section heading.  Lower-level
                             % sections are begun with similar
                             % \subsection and \subsubsection commands.

The ends  of words and sentences are marked
  by   spaces. It  doesn't matter how many
spaces    you type; one is as good as 100.  The
end of   a line counts as a space.

One   or more   blank lines denote the  end
of  a paragraph.

Since any number of consecutive spaces are treated
like a single one, the formatting of the input
file makes no difference to
      \LaTeX,                % The \LaTeX command generates the LaTeX logo.
but it makes a difference to you.  When you use
\LaTeX, making your input file as easy to read
as possible will be a great help as you write
your document and when you change it.  This sample
file shows how you can add comments to your own input
file.

Because printing is different from typewriting,
there are a number of things that you have to do
differently when preparing an input file than if
you were just typing the document directly.
Quotation marks like
       ``this''
have to be handled specially, as do quotes within
quotes:
       ``\,`this'            % \, separates the double and single quote.
        is what I just
        wrote, not  `that'\,''.

Dashes come in three sizes: an
       intra-word
dash, a medium dash for number ranges like
       1--2,
and a punctuation
       dash---like
this.

A sentence-ending space should be larger than the
space between words within a sentence.  You
sometimes have to type special commands in
conjunction with punctuation characters to get
this right, as in the following sentence.
       Gnats, gnus, etc.\ all  % `\ ' makes an inter-word space.
       begin with G\@.         % \@ marks end-of-sentence punctuation.
You should check the spaces after periods when
reading your output to make sure you haven't
forgotten any special cases.  Generating an
ellipsis
       \ldots\               % `\ ' is needed after `\ldots' because TeX
                             % ignores spaces after command names like \ldots
                             % made from \ + letters.
                             %
                             % Note how a `%' character causes TeX to ignore
                             % the end of the input line, so these blank lines
                             % do not start a new paragraph.
                             %
with the right spacing around the periods requires
a special command.

\LaTeX\ interprets some common characters as
commands, so you must type special commands to
generate them.  These characters include the
following:
       \$ \& \% \# \{ and \}.

In printing, text is usually emphasized with an
       \emph{italic}
type style.

\begin{em}
   A long segment of text can also be emphasized
   in this way.  Text within such a segment can be
   given \emph{additional} emphasis.
\end{em}

It is sometimes necessary to prevent \LaTeX\ from
breaking a line where it might otherwise do so.
This may be at a space, as between the ``Mr.''\ and
``Jones'' in
       ``Mr.~Jones'',        % ~ produces an unbreakable interword space.
or within a word---especially when the word is a
symbol like
       \mbox{\emph{itemnum}}
that makes little sense when hyphenated across
lines.

Footnotes\footnote{This is an example of a footnote.}
pose no problem.

\LaTeX\ is good at typesetting mathematical formulas
like
       \( x-3y + z = 7 \)
or
       \( a_{1} > x^{2n} + y^{2n} > x' \)
or
       \( \ip{A}{B} = \sum_{i} a_{i} b_{i} \).
The spaces you type in a formula are
ignored.  Remember that a letter like
       $x$                   % $ ... $  and  \( ... \)  are equivalent
is a formula when it denotes a mathematical
symbol, and it should be typed as one.

\section{Displayed Text}

Text is displayed by indenting it from the left
margin.  Quotations are commonly displayed.  There
are short quotations
\begin{quote}
   This is a short quotation.  It consists of a
   single paragraph of text.  See how it is formatted.
\end{quote}
and longer ones.
\begin{quotation}
   This is a longer quotation.  It consists of two
   paragraphs of text, neither of which are
   particularly interesting.

   This is the second paragraph of the quotation.  It
   is just as dull as the first paragraph.
\end{quotation}
Another frequently-displayed structure is a list.
The following is an example of an \emph{itemized}
list.
\begin{itemize}
   \item This is the first item of an itemized list.
         Each item in the list is marked with a ``tick''.
         You don't have to worry about what kind of tick
         mark is used.

   \item This is the second item of the list.  It
         contains another list nested inside it.  The inner
         list is an \emph{enumerated} list.
         \begin{enumerate}
            \item This is the first item of an enumerated
                  list that is nested within the itemized list.

            \item This is the second item of the inner list.
                  \LaTeX\ allows you to nest lists deeper than
                  you really should.
         \end{enumerate}
         This is the rest of the second item of the outer
         list.  It is no more interesting than any other
         part of the item.
   \item This is the third item of the list.
\end{itemize}
You can even display poetry.
\begin{verse}
   There is an environment
    for verse \\             % The \\ command separates lines
   Whose features some poets % within a stanza.
   will curse.

                             % One or more blank lines separate stanzas.

   For instead of making\\
   Them do \emph{all} line breaking, \\
   It allows them to put too many words on a line when they'd rather be
   forced to be terse.
\end{verse}

Mathematical formulas may also be displayed.  A
displayed formula
is
one-line long; multiline
formulas require special formatting instructions.
   \[  \ip{\Gamma}{\psi'} = x'' + y^{2} + z_{i}^{n}\]
Don't start a paragraph with a displayed equation,
nor make one a paragraph by itself.

\end{document}               % End of document.

\section{Bench Huge Section 984}
% synthetic bench section 984
% This is a sample LaTeX input file.  (Version of 12 August 2004.)
%
% A '%' character causes TeX to ignore all remaining text on the line,
% and is used for comments like this one.

\documentclass{article}      % Specifies the document class

                             % The preamble begins here.
\title{An Example Document}  % Declares the document's title.
\author{Leslie Lamport}      % Declares the author's name.
\date{January 21, 1994}      % Deleting this command produces today's date.

\newcommand{\ip}[2]{(#1, #2)}
                             % Defines \ip{arg1}{arg2} to mean
                             % (arg1, arg2).

%\newcommand{\ip}[2]{\langle #1 | #2\rangle}
                             % This is an alternative definition of
                             % \ip that is commented out.

\begin{document}             % End of preamble and beginning of text.

\maketitle                   % Produces the title.

This is an example input file.  Comparing it with
the output it generates can show you how to
produce a simple document of your own.

\section{Ordinary Text}      % Produces section heading.  Lower-level
                             % sections are begun with similar
                             % \subsection and \subsubsection commands.

The ends  of words and sentences are marked
  by   spaces. It  doesn't matter how many
spaces    you type; one is as good as 100.  The
end of   a line counts as a space.

One   or more   blank lines denote the  end
of  a paragraph.

Since any number of consecutive spaces are treated
like a single one, the formatting of the input
file makes no difference to
      \LaTeX,                % The \LaTeX command generates the LaTeX logo.
but it makes a difference to you.  When you use
\LaTeX, making your input file as easy to read
as possible will be a great help as you write
your document and when you change it.  This sample
file shows how you can add comments to your own input
file.

Because printing is different from typewriting,
there are a number of things that you have to do
differently when preparing an input file than if
you were just typing the document directly.
Quotation marks like
       ``this''
have to be handled specially, as do quotes within
quotes:
       ``\,`this'            % \, separates the double and single quote.
        is what I just
        wrote, not  `that'\,''.

Dashes come in three sizes: an
       intra-word
dash, a medium dash for number ranges like
       1--2,
and a punctuation
       dash---like
this.

A sentence-ending space should be larger than the
space between words within a sentence.  You
sometimes have to type special commands in
conjunction with punctuation characters to get
this right, as in the following sentence.
       Gnats, gnus, etc.\ all  % `\ ' makes an inter-word space.
       begin with G\@.         % \@ marks end-of-sentence punctuation.
You should check the spaces after periods when
reading your output to make sure you haven't
forgotten any special cases.  Generating an
ellipsis
       \ldots\               % `\ ' is needed after `\ldots' because TeX
                             % ignores spaces after command names like \ldots
                             % made from \ + letters.
                             %
                             % Note how a `%' character causes TeX to ignore
                             % the end of the input line, so these blank lines
                             % do not start a new paragraph.
                             %
with the right spacing around the periods requires
a special command.

\LaTeX\ interprets some common characters as
commands, so you must type special commands to
generate them.  These characters include the
following:
       \$ \& \% \# \{ and \}.

In printing, text is usually emphasized with an
       \emph{italic}
type style.

\begin{em}
   A long segment of text can also be emphasized
   in this way.  Text within such a segment can be
   given \emph{additional} emphasis.
\end{em}

It is sometimes necessary to prevent \LaTeX\ from
breaking a line where it might otherwise do so.
This may be at a space, as between the ``Mr.''\ and
``Jones'' in
       ``Mr.~Jones'',        % ~ produces an unbreakable interword space.
or within a word---especially when the word is a
symbol like
       \mbox{\emph{itemnum}}
that makes little sense when hyphenated across
lines.

Footnotes\footnote{This is an example of a footnote.}
pose no problem.

\LaTeX\ is good at typesetting mathematical formulas
like
       \( x-3y + z = 7 \)
or
       \( a_{1} > x^{2n} + y^{2n} > x' \)
or
       \( \ip{A}{B} = \sum_{i} a_{i} b_{i} \).
The spaces you type in a formula are
ignored.  Remember that a letter like
       $x$                   % $ ... $  and  \( ... \)  are equivalent
is a formula when it denotes a mathematical
symbol, and it should be typed as one.

\section{Displayed Text}

Text is displayed by indenting it from the left
margin.  Quotations are commonly displayed.  There
are short quotations
\begin{quote}
   This is a short quotation.  It consists of a
   single paragraph of text.  See how it is formatted.
\end{quote}
and longer ones.
\begin{quotation}
   This is a longer quotation.  It consists of two
   paragraphs of text, neither of which are
   particularly interesting.

   This is the second paragraph of the quotation.  It
   is just as dull as the first paragraph.
\end{quotation}
Another frequently-displayed structure is a list.
The following is an example of an \emph{itemized}
list.
\begin{itemize}
   \item This is the first item of an itemized list.
         Each item in the list is marked with a ``tick''.
         You don't have to worry about what kind of tick
         mark is used.

   \item This is the second item of the list.  It
         contains another list nested inside it.  The inner
         list is an \emph{enumerated} list.
         \begin{enumerate}
            \item This is the first item of an enumerated
                  list that is nested within the itemized list.

            \item This is the second item of the inner list.
                  \LaTeX\ allows you to nest lists deeper than
                  you really should.
         \end{enumerate}
         This is the rest of the second item of the outer
         list.  It is no more interesting than any other
         part of the item.
   \item This is the third item of the list.
\end{itemize}
You can even display poetry.
\begin{verse}
   There is an environment
    for verse \\             % The \\ command separates lines
   Whose features some poets % within a stanza.
   will curse.

                             % One or more blank lines separate stanzas.

   For instead of making\\
   Them do \emph{all} line breaking, \\
   It allows them to put too many words on a line when they'd rather be
   forced to be terse.
\end{verse}

Mathematical formulas may also be displayed.  A
displayed formula
is
one-line long; multiline
formulas require special formatting instructions.
   \[  \ip{\Gamma}{\psi'} = x'' + y^{2} + z_{i}^{n}\]
Don't start a paragraph with a displayed equation,
nor make one a paragraph by itself.

\end{document}               % End of document.

\section{Bench Huge Section 985}
% synthetic bench section 985
% This is a sample LaTeX input file.  (Version of 12 August 2004.)
%
% A '%' character causes TeX to ignore all remaining text on the line,
% and is used for comments like this one.

\documentclass{article}      % Specifies the document class

                             % The preamble begins here.
\title{An Example Document}  % Declares the document's title.
\author{Leslie Lamport}      % Declares the author's name.
\date{January 21, 1994}      % Deleting this command produces today's date.

\newcommand{\ip}[2]{(#1, #2)}
                             % Defines \ip{arg1}{arg2} to mean
                             % (arg1, arg2).

%\newcommand{\ip}[2]{\langle #1 | #2\rangle}
                             % This is an alternative definition of
                             % \ip that is commented out.

\begin{document}             % End of preamble and beginning of text.

\maketitle                   % Produces the title.

This is an example input file.  Comparing it with
the output it generates can show you how to
produce a simple document of your own.

\section{Ordinary Text}      % Produces section heading.  Lower-level
                             % sections are begun with similar
                             % \subsection and \subsubsection commands.

The ends  of words and sentences are marked
  by   spaces. It  doesn't matter how many
spaces    you type; one is as good as 100.  The
end of   a line counts as a space.

One   or more   blank lines denote the  end
of  a paragraph.

Since any number of consecutive spaces are treated
like a single one, the formatting of the input
file makes no difference to
      \LaTeX,                % The \LaTeX command generates the LaTeX logo.
but it makes a difference to you.  When you use
\LaTeX, making your input file as easy to read
as possible will be a great help as you write
your document and when you change it.  This sample
file shows how you can add comments to your own input
file.

Because printing is different from typewriting,
there are a number of things that you have to do
differently when preparing an input file than if
you were just typing the document directly.
Quotation marks like
       ``this''
have to be handled specially, as do quotes within
quotes:
       ``\,`this'            % \, separates the double and single quote.
        is what I just
        wrote, not  `that'\,''.

Dashes come in three sizes: an
       intra-word
dash, a medium dash for number ranges like
       1--2,
and a punctuation
       dash---like
this.

A sentence-ending space should be larger than the
space between words within a sentence.  You
sometimes have to type special commands in
conjunction with punctuation characters to get
this right, as in the following sentence.
       Gnats, gnus, etc.\ all  % `\ ' makes an inter-word space.
       begin with G\@.         % \@ marks end-of-sentence punctuation.
You should check the spaces after periods when
reading your output to make sure you haven't
forgotten any special cases.  Generating an
ellipsis
       \ldots\               % `\ ' is needed after `\ldots' because TeX
                             % ignores spaces after command names like \ldots
                             % made from \ + letters.
                             %
                             % Note how a `%' character causes TeX to ignore
                             % the end of the input line, so these blank lines
                             % do not start a new paragraph.
                             %
with the right spacing around the periods requires
a special command.

\LaTeX\ interprets some common characters as
commands, so you must type special commands to
generate them.  These characters include the
following:
       \$ \& \% \# \{ and \}.

In printing, text is usually emphasized with an
       \emph{italic}
type style.

\begin{em}
   A long segment of text can also be emphasized
   in this way.  Text within such a segment can be
   given \emph{additional} emphasis.
\end{em}

It is sometimes necessary to prevent \LaTeX\ from
breaking a line where it might otherwise do so.
This may be at a space, as between the ``Mr.''\ and
``Jones'' in
       ``Mr.~Jones'',        % ~ produces an unbreakable interword space.
or within a word---especially when the word is a
symbol like
       \mbox{\emph{itemnum}}
that makes little sense when hyphenated across
lines.

Footnotes\footnote{This is an example of a footnote.}
pose no problem.

\LaTeX\ is good at typesetting mathematical formulas
like
       \( x-3y + z = 7 \)
or
       \( a_{1} > x^{2n} + y^{2n} > x' \)
or
       \( \ip{A}{B} = \sum_{i} a_{i} b_{i} \).
The spaces you type in a formula are
ignored.  Remember that a letter like
       $x$                   % $ ... $  and  \( ... \)  are equivalent
is a formula when it denotes a mathematical
symbol, and it should be typed as one.

\section{Displayed Text}

Text is displayed by indenting it from the left
margin.  Quotations are commonly displayed.  There
are short quotations
\begin{quote}
   This is a short quotation.  It consists of a
   single paragraph of text.  See how it is formatted.
\end{quote}
and longer ones.
\begin{quotation}
   This is a longer quotation.  It consists of two
   paragraphs of text, neither of which are
   particularly interesting.

   This is the second paragraph of the quotation.  It
   is just as dull as the first paragraph.
\end{quotation}
Another frequently-displayed structure is a list.
The following is an example of an \emph{itemized}
list.
\begin{itemize}
   \item This is the first item of an itemized list.
         Each item in the list is marked with a ``tick''.
         You don't have to worry about what kind of tick
         mark is used.

   \item This is the second item of the list.  It
         contains another list nested inside it.  The inner
         list is an \emph{enumerated} list.
         \begin{enumerate}
            \item This is the first item of an enumerated
                  list that is nested within the itemized list.

            \item This is the second item of the inner list.
                  \LaTeX\ allows you to nest lists deeper than
                  you really should.
         \end{enumerate}
         This is the rest of the second item of the outer
         list.  It is no more interesting than any other
         part of the item.
   \item This is the third item of the list.
\end{itemize}
You can even display poetry.
\begin{verse}
   There is an environment
    for verse \\             % The \\ command separates lines
   Whose features some poets % within a stanza.
   will curse.

                             % One or more blank lines separate stanzas.

   For instead of making\\
   Them do \emph{all} line breaking, \\
   It allows them to put too many words on a line when they'd rather be
   forced to be terse.
\end{verse}

Mathematical formulas may also be displayed.  A
displayed formula
is
one-line long; multiline
formulas require special formatting instructions.
   \[  \ip{\Gamma}{\psi'} = x'' + y^{2} + z_{i}^{n}\]
Don't start a paragraph with a displayed equation,
nor make one a paragraph by itself.

\end{document}               % End of document.

\section{Bench Huge Section 986}
% synthetic bench section 986
% This is a sample LaTeX input file.  (Version of 12 August 2004.)
%
% A '%' character causes TeX to ignore all remaining text on the line,
% and is used for comments like this one.

\documentclass{article}      % Specifies the document class

                             % The preamble begins here.
\title{An Example Document}  % Declares the document's title.
\author{Leslie Lamport}      % Declares the author's name.
\date{January 21, 1994}      % Deleting this command produces today's date.

\newcommand{\ip}[2]{(#1, #2)}
                             % Defines \ip{arg1}{arg2} to mean
                             % (arg1, arg2).

%\newcommand{\ip}[2]{\langle #1 | #2\rangle}
                             % This is an alternative definition of
                             % \ip that is commented out.

\begin{document}             % End of preamble and beginning of text.

\maketitle                   % Produces the title.

This is an example input file.  Comparing it with
the output it generates can show you how to
produce a simple document of your own.

\section{Ordinary Text}      % Produces section heading.  Lower-level
                             % sections are begun with similar
                             % \subsection and \subsubsection commands.

The ends  of words and sentences are marked
  by   spaces. It  doesn't matter how many
spaces    you type; one is as good as 100.  The
end of   a line counts as a space.

One   or more   blank lines denote the  end
of  a paragraph.

Since any number of consecutive spaces are treated
like a single one, the formatting of the input
file makes no difference to
      \LaTeX,                % The \LaTeX command generates the LaTeX logo.
but it makes a difference to you.  When you use
\LaTeX, making your input file as easy to read
as possible will be a great help as you write
your document and when you change it.  This sample
file shows how you can add comments to your own input
file.

Because printing is different from typewriting,
there are a number of things that you have to do
differently when preparing an input file than if
you were just typing the document directly.
Quotation marks like
       ``this''
have to be handled specially, as do quotes within
quotes:
       ``\,`this'            % \, separates the double and single quote.
        is what I just
        wrote, not  `that'\,''.

Dashes come in three sizes: an
       intra-word
dash, a medium dash for number ranges like
       1--2,
and a punctuation
       dash---like
this.

A sentence-ending space should be larger than the
space between words within a sentence.  You
sometimes have to type special commands in
conjunction with punctuation characters to get
this right, as in the following sentence.
       Gnats, gnus, etc.\ all  % `\ ' makes an inter-word space.
       begin with G\@.         % \@ marks end-of-sentence punctuation.
You should check the spaces after periods when
reading your output to make sure you haven't
forgotten any special cases.  Generating an
ellipsis
       \ldots\               % `\ ' is needed after `\ldots' because TeX
                             % ignores spaces after command names like \ldots
                             % made from \ + letters.
                             %
                             % Note how a `%' character causes TeX to ignore
                             % the end of the input line, so these blank lines
                             % do not start a new paragraph.
                             %
with the right spacing around the periods requires
a special command.

\LaTeX\ interprets some common characters as
commands, so you must type special commands to
generate them.  These characters include the
following:
       \$ \& \% \# \{ and \}.

In printing, text is usually emphasized with an
       \emph{italic}
type style.

\begin{em}
   A long segment of text can also be emphasized
   in this way.  Text within such a segment can be
   given \emph{additional} emphasis.
\end{em}

It is sometimes necessary to prevent \LaTeX\ from
breaking a line where it might otherwise do so.
This may be at a space, as between the ``Mr.''\ and
``Jones'' in
       ``Mr.~Jones'',        % ~ produces an unbreakable interword space.
or within a word---especially when the word is a
symbol like
       \mbox{\emph{itemnum}}
that makes little sense when hyphenated across
lines.

Footnotes\footnote{This is an example of a footnote.}
pose no problem.

\LaTeX\ is good at typesetting mathematical formulas
like
       \( x-3y + z = 7 \)
or
       \( a_{1} > x^{2n} + y^{2n} > x' \)
or
       \( \ip{A}{B} = \sum_{i} a_{i} b_{i} \).
The spaces you type in a formula are
ignored.  Remember that a letter like
       $x$                   % $ ... $  and  \( ... \)  are equivalent
is a formula when it denotes a mathematical
symbol, and it should be typed as one.

\section{Displayed Text}

Text is displayed by indenting it from the left
margin.  Quotations are commonly displayed.  There
are short quotations
\begin{quote}
   This is a short quotation.  It consists of a
   single paragraph of text.  See how it is formatted.
\end{quote}
and longer ones.
\begin{quotation}
   This is a longer quotation.  It consists of two
   paragraphs of text, neither of which are
   particularly interesting.

   This is the second paragraph of the quotation.  It
   is just as dull as the first paragraph.
\end{quotation}
Another frequently-displayed structure is a list.
The following is an example of an \emph{itemized}
list.
\begin{itemize}
   \item This is the first item of an itemized list.
         Each item in the list is marked with a ``tick''.
         You don't have to worry about what kind of tick
         mark is used.

   \item This is the second item of the list.  It
         contains another list nested inside it.  The inner
         list is an \emph{enumerated} list.
         \begin{enumerate}
            \item This is the first item of an enumerated
                  list that is nested within the itemized list.

            \item This is the second item of the inner list.
                  \LaTeX\ allows you to nest lists deeper than
                  you really should.
         \end{enumerate}
         This is the rest of the second item of the outer
         list.  It is no more interesting than any other
         part of the item.
   \item This is the third item of the list.
\end{itemize}
You can even display poetry.
\begin{verse}
   There is an environment
    for verse \\             % The \\ command separates lines
   Whose features some poets % within a stanza.
   will curse.

                             % One or more blank lines separate stanzas.

   For instead of making\\
   Them do \emph{all} line breaking, \\
   It allows them to put too many words on a line when they'd rather be
   forced to be terse.
\end{verse}

Mathematical formulas may also be displayed.  A
displayed formula
is
one-line long; multiline
formulas require special formatting instructions.
   \[  \ip{\Gamma}{\psi'} = x'' + y^{2} + z_{i}^{n}\]
Don't start a paragraph with a displayed equation,
nor make one a paragraph by itself.

\end{document}               % End of document.

\section{Bench Huge Section 987}
% synthetic bench section 987
% This is a sample LaTeX input file.  (Version of 12 August 2004.)
%
% A '%' character causes TeX to ignore all remaining text on the line,
% and is used for comments like this one.

\documentclass{article}      % Specifies the document class

                             % The preamble begins here.
\title{An Example Document}  % Declares the document's title.
\author{Leslie Lamport}      % Declares the author's name.
\date{January 21, 1994}      % Deleting this command produces today's date.

\newcommand{\ip}[2]{(#1, #2)}
                             % Defines \ip{arg1}{arg2} to mean
                             % (arg1, arg2).

%\newcommand{\ip}[2]{\langle #1 | #2\rangle}
                             % This is an alternative definition of
                             % \ip that is commented out.

\begin{document}             % End of preamble and beginning of text.

\maketitle                   % Produces the title.

This is an example input file.  Comparing it with
the output it generates can show you how to
produce a simple document of your own.

\section{Ordinary Text}      % Produces section heading.  Lower-level
                             % sections are begun with similar
                             % \subsection and \subsubsection commands.

The ends  of words and sentences are marked
  by   spaces. It  doesn't matter how many
spaces    you type; one is as good as 100.  The
end of   a line counts as a space.

One   or more   blank lines denote the  end
of  a paragraph.

Since any number of consecutive spaces are treated
like a single one, the formatting of the input
file makes no difference to
      \LaTeX,                % The \LaTeX command generates the LaTeX logo.
but it makes a difference to you.  When you use
\LaTeX, making your input file as easy to read
as possible will be a great help as you write
your document and when you change it.  This sample
file shows how you can add comments to your own input
file.

Because printing is different from typewriting,
there are a number of things that you have to do
differently when preparing an input file than if
you were just typing the document directly.
Quotation marks like
       ``this''
have to be handled specially, as do quotes within
quotes:
       ``\,`this'            % \, separates the double and single quote.
        is what I just
        wrote, not  `that'\,''.

Dashes come in three sizes: an
       intra-word
dash, a medium dash for number ranges like
       1--2,
and a punctuation
       dash---like
this.

A sentence-ending space should be larger than the
space between words within a sentence.  You
sometimes have to type special commands in
conjunction with punctuation characters to get
this right, as in the following sentence.
       Gnats, gnus, etc.\ all  % `\ ' makes an inter-word space.
       begin with G\@.         % \@ marks end-of-sentence punctuation.
You should check the spaces after periods when
reading your output to make sure you haven't
forgotten any special cases.  Generating an
ellipsis
       \ldots\               % `\ ' is needed after `\ldots' because TeX
                             % ignores spaces after command names like \ldots
                             % made from \ + letters.
                             %
                             % Note how a `%' character causes TeX to ignore
                             % the end of the input line, so these blank lines
                             % do not start a new paragraph.
                             %
with the right spacing around the periods requires
a special command.

\LaTeX\ interprets some common characters as
commands, so you must type special commands to
generate them.  These characters include the
following:
       \$ \& \% \# \{ and \}.

In printing, text is usually emphasized with an
       \emph{italic}
type style.

\begin{em}
   A long segment of text can also be emphasized
   in this way.  Text within such a segment can be
   given \emph{additional} emphasis.
\end{em}

It is sometimes necessary to prevent \LaTeX\ from
breaking a line where it might otherwise do so.
This may be at a space, as between the ``Mr.''\ and
``Jones'' in
       ``Mr.~Jones'',        % ~ produces an unbreakable interword space.
or within a word---especially when the word is a
symbol like
       \mbox{\emph{itemnum}}
that makes little sense when hyphenated across
lines.

Footnotes\footnote{This is an example of a footnote.}
pose no problem.

\LaTeX\ is good at typesetting mathematical formulas
like
       \( x-3y + z = 7 \)
or
       \( a_{1} > x^{2n} + y^{2n} > x' \)
or
       \( \ip{A}{B} = \sum_{i} a_{i} b_{i} \).
The spaces you type in a formula are
ignored.  Remember that a letter like
       $x$                   % $ ... $  and  \( ... \)  are equivalent
is a formula when it denotes a mathematical
symbol, and it should be typed as one.

\section{Displayed Text}

Text is displayed by indenting it from the left
margin.  Quotations are commonly displayed.  There
are short quotations
\begin{quote}
   This is a short quotation.  It consists of a
   single paragraph of text.  See how it is formatted.
\end{quote}
and longer ones.
\begin{quotation}
   This is a longer quotation.  It consists of two
   paragraphs of text, neither of which are
   particularly interesting.

   This is the second paragraph of the quotation.  It
   is just as dull as the first paragraph.
\end{quotation}
Another frequently-displayed structure is a list.
The following is an example of an \emph{itemized}
list.
\begin{itemize}
   \item This is the first item of an itemized list.
         Each item in the list is marked with a ``tick''.
         You don't have to worry about what kind of tick
         mark is used.

   \item This is the second item of the list.  It
         contains another list nested inside it.  The inner
         list is an \emph{enumerated} list.
         \begin{enumerate}
            \item This is the first item of an enumerated
                  list that is nested within the itemized list.

            \item This is the second item of the inner list.
                  \LaTeX\ allows you to nest lists deeper than
                  you really should.
         \end{enumerate}
         This is the rest of the second item of the outer
         list.  It is no more interesting than any other
         part of the item.
   \item This is the third item of the list.
\end{itemize}
You can even display poetry.
\begin{verse}
   There is an environment
    for verse \\             % The \\ command separates lines
   Whose features some poets % within a stanza.
   will curse.

                             % One or more blank lines separate stanzas.

   For instead of making\\
   Them do \emph{all} line breaking, \\
   It allows them to put too many words on a line when they'd rather be
   forced to be terse.
\end{verse}

Mathematical formulas may also be displayed.  A
displayed formula
is
one-line long; multiline
formulas require special formatting instructions.
   \[  \ip{\Gamma}{\psi'} = x'' + y^{2} + z_{i}^{n}\]
Don't start a paragraph with a displayed equation,
nor make one a paragraph by itself.

\end{document}               % End of document.

\section{Bench Huge Section 988}
% synthetic bench section 988
% This is a sample LaTeX input file.  (Version of 12 August 2004.)
%
% A '%' character causes TeX to ignore all remaining text on the line,
% and is used for comments like this one.

\documentclass{article}      % Specifies the document class

                             % The preamble begins here.
\title{An Example Document}  % Declares the document's title.
\author{Leslie Lamport}      % Declares the author's name.
\date{January 21, 1994}      % Deleting this command produces today's date.

\newcommand{\ip}[2]{(#1, #2)}
                             % Defines \ip{arg1}{arg2} to mean
                             % (arg1, arg2).

%\newcommand{\ip}[2]{\langle #1 | #2\rangle}
                             % This is an alternative definition of
                             % \ip that is commented out.

\begin{document}             % End of preamble and beginning of text.

\maketitle                   % Produces the title.

This is an example input file.  Comparing it with
the output it generates can show you how to
produce a simple document of your own.

\section{Ordinary Text}      % Produces section heading.  Lower-level
                             % sections are begun with similar
                             % \subsection and \subsubsection commands.

The ends  of words and sentences are marked
  by   spaces. It  doesn't matter how many
spaces    you type; one is as good as 100.  The
end of   a line counts as a space.

One   or more   blank lines denote the  end
of  a paragraph.

Since any number of consecutive spaces are treated
like a single one, the formatting of the input
file makes no difference to
      \LaTeX,                % The \LaTeX command generates the LaTeX logo.
but it makes a difference to you.  When you use
\LaTeX, making your input file as easy to read
as possible will be a great help as you write
your document and when you change it.  This sample
file shows how you can add comments to your own input
file.

Because printing is different from typewriting,
there are a number of things that you have to do
differently when preparing an input file than if
you were just typing the document directly.
Quotation marks like
       ``this''
have to be handled specially, as do quotes within
quotes:
       ``\,`this'            % \, separates the double and single quote.
        is what I just
        wrote, not  `that'\,''.

Dashes come in three sizes: an
       intra-word
dash, a medium dash for number ranges like
       1--2,
and a punctuation
       dash---like
this.

A sentence-ending space should be larger than the
space between words within a sentence.  You
sometimes have to type special commands in
conjunction with punctuation characters to get
this right, as in the following sentence.
       Gnats, gnus, etc.\ all  % `\ ' makes an inter-word space.
       begin with G\@.         % \@ marks end-of-sentence punctuation.
You should check the spaces after periods when
reading your output to make sure you haven't
forgotten any special cases.  Generating an
ellipsis
       \ldots\               % `\ ' is needed after `\ldots' because TeX
                             % ignores spaces after command names like \ldots
                             % made from \ + letters.
                             %
                             % Note how a `%' character causes TeX to ignore
                             % the end of the input line, so these blank lines
                             % do not start a new paragraph.
                             %
with the right spacing around the periods requires
a special command.

\LaTeX\ interprets some common characters as
commands, so you must type special commands to
generate them.  These characters include the
following:
       \$ \& \% \# \{ and \}.

In printing, text is usually emphasized with an
       \emph{italic}
type style.

\begin{em}
   A long segment of text can also be emphasized
   in this way.  Text within such a segment can be
   given \emph{additional} emphasis.
\end{em}

It is sometimes necessary to prevent \LaTeX\ from
breaking a line where it might otherwise do so.
This may be at a space, as between the ``Mr.''\ and
``Jones'' in
       ``Mr.~Jones'',        % ~ produces an unbreakable interword space.
or within a word---especially when the word is a
symbol like
       \mbox{\emph{itemnum}}
that makes little sense when hyphenated across
lines.

Footnotes\footnote{This is an example of a footnote.}
pose no problem.

\LaTeX\ is good at typesetting mathematical formulas
like
       \( x-3y + z = 7 \)
or
       \( a_{1} > x^{2n} + y^{2n} > x' \)
or
       \( \ip{A}{B} = \sum_{i} a_{i} b_{i} \).
The spaces you type in a formula are
ignored.  Remember that a letter like
       $x$                   % $ ... $  and  \( ... \)  are equivalent
is a formula when it denotes a mathematical
symbol, and it should be typed as one.

\section{Displayed Text}

Text is displayed by indenting it from the left
margin.  Quotations are commonly displayed.  There
are short quotations
\begin{quote}
   This is a short quotation.  It consists of a
   single paragraph of text.  See how it is formatted.
\end{quote}
and longer ones.
\begin{quotation}
   This is a longer quotation.  It consists of two
   paragraphs of text, neither of which are
   particularly interesting.

   This is the second paragraph of the quotation.  It
   is just as dull as the first paragraph.
\end{quotation}
Another frequently-displayed structure is a list.
The following is an example of an \emph{itemized}
list.
\begin{itemize}
   \item This is the first item of an itemized list.
         Each item in the list is marked with a ``tick''.
         You don't have to worry about what kind of tick
         mark is used.

   \item This is the second item of the list.  It
         contains another list nested inside it.  The inner
         list is an \emph{enumerated} list.
         \begin{enumerate}
            \item This is the first item of an enumerated
                  list that is nested within the itemized list.

            \item This is the second item of the inner list.
                  \LaTeX\ allows you to nest lists deeper than
                  you really should.
         \end{enumerate}
         This is the rest of the second item of the outer
         list.  It is no more interesting than any other
         part of the item.
   \item This is the third item of the list.
\end{itemize}
You can even display poetry.
\begin{verse}
   There is an environment
    for verse \\             % The \\ command separates lines
   Whose features some poets % within a stanza.
   will curse.

                             % One or more blank lines separate stanzas.

   For instead of making\\
   Them do \emph{all} line breaking, \\
   It allows them to put too many words on a line when they'd rather be
   forced to be terse.
\end{verse}

Mathematical formulas may also be displayed.  A
displayed formula
is
one-line long; multiline
formulas require special formatting instructions.
   \[  \ip{\Gamma}{\psi'} = x'' + y^{2} + z_{i}^{n}\]
Don't start a paragraph with a displayed equation,
nor make one a paragraph by itself.

\end{document}               % End of document.

\section{Bench Huge Section 989}
% synthetic bench section 989
% This is a sample LaTeX input file.  (Version of 12 August 2004.)
%
% A '%' character causes TeX to ignore all remaining text on the line,
% and is used for comments like this one.

\documentclass{article}      % Specifies the document class

                             % The preamble begins here.
\title{An Example Document}  % Declares the document's title.
\author{Leslie Lamport}      % Declares the author's name.
\date{January 21, 1994}      % Deleting this command produces today's date.

\newcommand{\ip}[2]{(#1, #2)}
                             % Defines \ip{arg1}{arg2} to mean
                             % (arg1, arg2).

%\newcommand{\ip}[2]{\langle #1 | #2\rangle}
                             % This is an alternative definition of
                             % \ip that is commented out.

\begin{document}             % End of preamble and beginning of text.

\maketitle                   % Produces the title.

This is an example input file.  Comparing it with
the output it generates can show you how to
produce a simple document of your own.

\section{Ordinary Text}      % Produces section heading.  Lower-level
                             % sections are begun with similar
                             % \subsection and \subsubsection commands.

The ends  of words and sentences are marked
  by   spaces. It  doesn't matter how many
spaces    you type; one is as good as 100.  The
end of   a line counts as a space.

One   or more   blank lines denote the  end
of  a paragraph.

Since any number of consecutive spaces are treated
like a single one, the formatting of the input
file makes no difference to
      \LaTeX,                % The \LaTeX command generates the LaTeX logo.
but it makes a difference to you.  When you use
\LaTeX, making your input file as easy to read
as possible will be a great help as you write
your document and when you change it.  This sample
file shows how you can add comments to your own input
file.

Because printing is different from typewriting,
there are a number of things that you have to do
differently when preparing an input file than if
you were just typing the document directly.
Quotation marks like
       ``this''
have to be handled specially, as do quotes within
quotes:
       ``\,`this'            % \, separates the double and single quote.
        is what I just
        wrote, not  `that'\,''.

Dashes come in three sizes: an
       intra-word
dash, a medium dash for number ranges like
       1--2,
and a punctuation
       dash---like
this.

A sentence-ending space should be larger than the
space between words within a sentence.  You
sometimes have to type special commands in
conjunction with punctuation characters to get
this right, as in the following sentence.
       Gnats, gnus, etc.\ all  % `\ ' makes an inter-word space.
       begin with G\@.         % \@ marks end-of-sentence punctuation.
You should check the spaces after periods when
reading your output to make sure you haven't
forgotten any special cases.  Generating an
ellipsis
       \ldots\               % `\ ' is needed after `\ldots' because TeX
                             % ignores spaces after command names like \ldots
                             % made from \ + letters.
                             %
                             % Note how a `%' character causes TeX to ignore
                             % the end of the input line, so these blank lines
                             % do not start a new paragraph.
                             %
with the right spacing around the periods requires
a special command.

\LaTeX\ interprets some common characters as
commands, so you must type special commands to
generate them.  These characters include the
following:
       \$ \& \% \# \{ and \}.

In printing, text is usually emphasized with an
       \emph{italic}
type style.

\begin{em}
   A long segment of text can also be emphasized
   in this way.  Text within such a segment can be
   given \emph{additional} emphasis.
\end{em}

It is sometimes necessary to prevent \LaTeX\ from
breaking a line where it might otherwise do so.
This may be at a space, as between the ``Mr.''\ and
``Jones'' in
       ``Mr.~Jones'',        % ~ produces an unbreakable interword space.
or within a word---especially when the word is a
symbol like
       \mbox{\emph{itemnum}}
that makes little sense when hyphenated across
lines.

Footnotes\footnote{This is an example of a footnote.}
pose no problem.

\LaTeX\ is good at typesetting mathematical formulas
like
       \( x-3y + z = 7 \)
or
       \( a_{1} > x^{2n} + y^{2n} > x' \)
or
       \( \ip{A}{B} = \sum_{i} a_{i} b_{i} \).
The spaces you type in a formula are
ignored.  Remember that a letter like
       $x$                   % $ ... $  and  \( ... \)  are equivalent
is a formula when it denotes a mathematical
symbol, and it should be typed as one.

\section{Displayed Text}

Text is displayed by indenting it from the left
margin.  Quotations are commonly displayed.  There
are short quotations
\begin{quote}
   This is a short quotation.  It consists of a
   single paragraph of text.  See how it is formatted.
\end{quote}
and longer ones.
\begin{quotation}
   This is a longer quotation.  It consists of two
   paragraphs of text, neither of which are
   particularly interesting.

   This is the second paragraph of the quotation.  It
   is just as dull as the first paragraph.
\end{quotation}
Another frequently-displayed structure is a list.
The following is an example of an \emph{itemized}
list.
\begin{itemize}
   \item This is the first item of an itemized list.
         Each item in the list is marked with a ``tick''.
         You don't have to worry about what kind of tick
         mark is used.

   \item This is the second item of the list.  It
         contains another list nested inside it.  The inner
         list is an \emph{enumerated} list.
         \begin{enumerate}
            \item This is the first item of an enumerated
                  list that is nested within the itemized list.

            \item This is the second item of the inner list.
                  \LaTeX\ allows you to nest lists deeper than
                  you really should.
         \end{enumerate}
         This is the rest of the second item of the outer
         list.  It is no more interesting than any other
         part of the item.
   \item This is the third item of the list.
\end{itemize}
You can even display poetry.
\begin{verse}
   There is an environment
    for verse \\             % The \\ command separates lines
   Whose features some poets % within a stanza.
   will curse.

                             % One or more blank lines separate stanzas.

   For instead of making\\
   Them do \emph{all} line breaking, \\
   It allows them to put too many words on a line when they'd rather be
   forced to be terse.
\end{verse}

Mathematical formulas may also be displayed.  A
displayed formula
is
one-line long; multiline
formulas require special formatting instructions.
   \[  \ip{\Gamma}{\psi'} = x'' + y^{2} + z_{i}^{n}\]
Don't start a paragraph with a displayed equation,
nor make one a paragraph by itself.

\end{document}               % End of document.

\section{Bench Huge Section 990}
% synthetic bench section 990
% This is a sample LaTeX input file.  (Version of 12 August 2004.)
%
% A '%' character causes TeX to ignore all remaining text on the line,
% and is used for comments like this one.

\documentclass{article}      % Specifies the document class

                             % The preamble begins here.
\title{An Example Document}  % Declares the document's title.
\author{Leslie Lamport}      % Declares the author's name.
\date{January 21, 1994}      % Deleting this command produces today's date.

\newcommand{\ip}[2]{(#1, #2)}
                             % Defines \ip{arg1}{arg2} to mean
                             % (arg1, arg2).

%\newcommand{\ip}[2]{\langle #1 | #2\rangle}
                             % This is an alternative definition of
                             % \ip that is commented out.

\begin{document}             % End of preamble and beginning of text.

\maketitle                   % Produces the title.

This is an example input file.  Comparing it with
the output it generates can show you how to
produce a simple document of your own.

\section{Ordinary Text}      % Produces section heading.  Lower-level
                             % sections are begun with similar
                             % \subsection and \subsubsection commands.

The ends  of words and sentences are marked
  by   spaces. It  doesn't matter how many
spaces    you type; one is as good as 100.  The
end of   a line counts as a space.

One   or more   blank lines denote the  end
of  a paragraph.

Since any number of consecutive spaces are treated
like a single one, the formatting of the input
file makes no difference to
      \LaTeX,                % The \LaTeX command generates the LaTeX logo.
but it makes a difference to you.  When you use
\LaTeX, making your input file as easy to read
as possible will be a great help as you write
your document and when you change it.  This sample
file shows how you can add comments to your own input
file.

Because printing is different from typewriting,
there are a number of things that you have to do
differently when preparing an input file than if
you were just typing the document directly.
Quotation marks like
       ``this''
have to be handled specially, as do quotes within
quotes:
       ``\,`this'            % \, separates the double and single quote.
        is what I just
        wrote, not  `that'\,''.

Dashes come in three sizes: an
       intra-word
dash, a medium dash for number ranges like
       1--2,
and a punctuation
       dash---like
this.

A sentence-ending space should be larger than the
space between words within a sentence.  You
sometimes have to type special commands in
conjunction with punctuation characters to get
this right, as in the following sentence.
       Gnats, gnus, etc.\ all  % `\ ' makes an inter-word space.
       begin with G\@.         % \@ marks end-of-sentence punctuation.
You should check the spaces after periods when
reading your output to make sure you haven't
forgotten any special cases.  Generating an
ellipsis
       \ldots\               % `\ ' is needed after `\ldots' because TeX
                             % ignores spaces after command names like \ldots
                             % made from \ + letters.
                             %
                             % Note how a `%' character causes TeX to ignore
                             % the end of the input line, so these blank lines
                             % do not start a new paragraph.
                             %
with the right spacing around the periods requires
a special command.

\LaTeX\ interprets some common characters as
commands, so you must type special commands to
generate them.  These characters include the
following:
       \$ \& \% \# \{ and \}.

In printing, text is usually emphasized with an
       \emph{italic}
type style.

\begin{em}
   A long segment of text can also be emphasized
   in this way.  Text within such a segment can be
   given \emph{additional} emphasis.
\end{em}

It is sometimes necessary to prevent \LaTeX\ from
breaking a line where it might otherwise do so.
This may be at a space, as between the ``Mr.''\ and
``Jones'' in
       ``Mr.~Jones'',        % ~ produces an unbreakable interword space.
or within a word---especially when the word is a
symbol like
       \mbox{\emph{itemnum}}
that makes little sense when hyphenated across
lines.

Footnotes\footnote{This is an example of a footnote.}
pose no problem.

\LaTeX\ is good at typesetting mathematical formulas
like
       \( x-3y + z = 7 \)
or
       \( a_{1} > x^{2n} + y^{2n} > x' \)
or
       \( \ip{A}{B} = \sum_{i} a_{i} b_{i} \).
The spaces you type in a formula are
ignored.  Remember that a letter like
       $x$                   % $ ... $  and  \( ... \)  are equivalent
is a formula when it denotes a mathematical
symbol, and it should be typed as one.

\section{Displayed Text}

Text is displayed by indenting it from the left
margin.  Quotations are commonly displayed.  There
are short quotations
\begin{quote}
   This is a short quotation.  It consists of a
   single paragraph of text.  See how it is formatted.
\end{quote}
and longer ones.
\begin{quotation}
   This is a longer quotation.  It consists of two
   paragraphs of text, neither of which are
   particularly interesting.

   This is the second paragraph of the quotation.  It
   is just as dull as the first paragraph.
\end{quotation}
Another frequently-displayed structure is a list.
The following is an example of an \emph{itemized}
list.
\begin{itemize}
   \item This is the first item of an itemized list.
         Each item in the list is marked with a ``tick''.
         You don't have to worry about what kind of tick
         mark is used.

   \item This is the second item of the list.  It
         contains another list nested inside it.  The inner
         list is an \emph{enumerated} list.
         \begin{enumerate}
            \item This is the first item of an enumerated
                  list that is nested within the itemized list.

            \item This is the second item of the inner list.
                  \LaTeX\ allows you to nest lists deeper than
                  you really should.
         \end{enumerate}
         This is the rest of the second item of the outer
         list.  It is no more interesting than any other
         part of the item.
   \item This is the third item of the list.
\end{itemize}
You can even display poetry.
\begin{verse}
   There is an environment
    for verse \\             % The \\ command separates lines
   Whose features some poets % within a stanza.
   will curse.

                             % One or more blank lines separate stanzas.

   For instead of making\\
   Them do \emph{all} line breaking, \\
   It allows them to put too many words on a line when they'd rather be
   forced to be terse.
\end{verse}

Mathematical formulas may also be displayed.  A
displayed formula
is
one-line long; multiline
formulas require special formatting instructions.
   \[  \ip{\Gamma}{\psi'} = x'' + y^{2} + z_{i}^{n}\]
Don't start a paragraph with a displayed equation,
nor make one a paragraph by itself.

\end{document}               % End of document.

\section{Bench Huge Section 991}
% synthetic bench section 991
% This is a sample LaTeX input file.  (Version of 12 August 2004.)
%
% A '%' character causes TeX to ignore all remaining text on the line,
% and is used for comments like this one.

\documentclass{article}      % Specifies the document class

                             % The preamble begins here.
\title{An Example Document}  % Declares the document's title.
\author{Leslie Lamport}      % Declares the author's name.
\date{January 21, 1994}      % Deleting this command produces today's date.

\newcommand{\ip}[2]{(#1, #2)}
                             % Defines \ip{arg1}{arg2} to mean
                             % (arg1, arg2).

%\newcommand{\ip}[2]{\langle #1 | #2\rangle}
                             % This is an alternative definition of
                             % \ip that is commented out.

\begin{document}             % End of preamble and beginning of text.

\maketitle                   % Produces the title.

This is an example input file.  Comparing it with
the output it generates can show you how to
produce a simple document of your own.

\section{Ordinary Text}      % Produces section heading.  Lower-level
                             % sections are begun with similar
                             % \subsection and \subsubsection commands.

The ends  of words and sentences are marked
  by   spaces. It  doesn't matter how many
spaces    you type; one is as good as 100.  The
end of   a line counts as a space.

One   or more   blank lines denote the  end
of  a paragraph.

Since any number of consecutive spaces are treated
like a single one, the formatting of the input
file makes no difference to
      \LaTeX,                % The \LaTeX command generates the LaTeX logo.
but it makes a difference to you.  When you use
\LaTeX, making your input file as easy to read
as possible will be a great help as you write
your document and when you change it.  This sample
file shows how you can add comments to your own input
file.

Because printing is different from typewriting,
there are a number of things that you have to do
differently when preparing an input file than if
you were just typing the document directly.
Quotation marks like
       ``this''
have to be handled specially, as do quotes within
quotes:
       ``\,`this'            % \, separates the double and single quote.
        is what I just
        wrote, not  `that'\,''.

Dashes come in three sizes: an
       intra-word
dash, a medium dash for number ranges like
       1--2,
and a punctuation
       dash---like
this.

A sentence-ending space should be larger than the
space between words within a sentence.  You
sometimes have to type special commands in
conjunction with punctuation characters to get
this right, as in the following sentence.
       Gnats, gnus, etc.\ all  % `\ ' makes an inter-word space.
       begin with G\@.         % \@ marks end-of-sentence punctuation.
You should check the spaces after periods when
reading your output to make sure you haven't
forgotten any special cases.  Generating an
ellipsis
       \ldots\               % `\ ' is needed after `\ldots' because TeX
                             % ignores spaces after command names like \ldots
                             % made from \ + letters.
                             %
                             % Note how a `%' character causes TeX to ignore
                             % the end of the input line, so these blank lines
                             % do not start a new paragraph.
                             %
with the right spacing around the periods requires
a special command.

\LaTeX\ interprets some common characters as
commands, so you must type special commands to
generate them.  These characters include the
following:
       \$ \& \% \# \{ and \}.

In printing, text is usually emphasized with an
       \emph{italic}
type style.

\begin{em}
   A long segment of text can also be emphasized
   in this way.  Text within such a segment can be
   given \emph{additional} emphasis.
\end{em}

It is sometimes necessary to prevent \LaTeX\ from
breaking a line where it might otherwise do so.
This may be at a space, as between the ``Mr.''\ and
``Jones'' in
       ``Mr.~Jones'',        % ~ produces an unbreakable interword space.
or within a word---especially when the word is a
symbol like
       \mbox{\emph{itemnum}}
that makes little sense when hyphenated across
lines.

Footnotes\footnote{This is an example of a footnote.}
pose no problem.

\LaTeX\ is good at typesetting mathematical formulas
like
       \( x-3y + z = 7 \)
or
       \( a_{1} > x^{2n} + y^{2n} > x' \)
or
       \( \ip{A}{B} = \sum_{i} a_{i} b_{i} \).
The spaces you type in a formula are
ignored.  Remember that a letter like
       $x$                   % $ ... $  and  \( ... \)  are equivalent
is a formula when it denotes a mathematical
symbol, and it should be typed as one.

\section{Displayed Text}

Text is displayed by indenting it from the left
margin.  Quotations are commonly displayed.  There
are short quotations
\begin{quote}
   This is a short quotation.  It consists of a
   single paragraph of text.  See how it is formatted.
\end{quote}
and longer ones.
\begin{quotation}
   This is a longer quotation.  It consists of two
   paragraphs of text, neither of which are
   particularly interesting.

   This is the second paragraph of the quotation.  It
   is just as dull as the first paragraph.
\end{quotation}
Another frequently-displayed structure is a list.
The following is an example of an \emph{itemized}
list.
\begin{itemize}
   \item This is the first item of an itemized list.
         Each item in the list is marked with a ``tick''.
         You don't have to worry about what kind of tick
         mark is used.

   \item This is the second item of the list.  It
         contains another list nested inside it.  The inner
         list is an \emph{enumerated} list.
         \begin{enumerate}
            \item This is the first item of an enumerated
                  list that is nested within the itemized list.

            \item This is the second item of the inner list.
                  \LaTeX\ allows you to nest lists deeper than
                  you really should.
         \end{enumerate}
         This is the rest of the second item of the outer
         list.  It is no more interesting than any other
         part of the item.
   \item This is the third item of the list.
\end{itemize}
You can even display poetry.
\begin{verse}
   There is an environment
    for verse \\             % The \\ command separates lines
   Whose features some poets % within a stanza.
   will curse.

                             % One or more blank lines separate stanzas.

   For instead of making\\
   Them do \emph{all} line breaking, \\
   It allows them to put too many words on a line when they'd rather be
   forced to be terse.
\end{verse}

Mathematical formulas may also be displayed.  A
displayed formula
is
one-line long; multiline
formulas require special formatting instructions.
   \[  \ip{\Gamma}{\psi'} = x'' + y^{2} + z_{i}^{n}\]
Don't start a paragraph with a displayed equation,
nor make one a paragraph by itself.

\end{document}               % End of document.

\section{Bench Huge Section 992}
% synthetic bench section 992
% This is a sample LaTeX input file.  (Version of 12 August 2004.)
%
% A '%' character causes TeX to ignore all remaining text on the line,
% and is used for comments like this one.

\documentclass{article}      % Specifies the document class

                             % The preamble begins here.
\title{An Example Document}  % Declares the document's title.
\author{Leslie Lamport}      % Declares the author's name.
\date{January 21, 1994}      % Deleting this command produces today's date.

\newcommand{\ip}[2]{(#1, #2)}
                             % Defines \ip{arg1}{arg2} to mean
                             % (arg1, arg2).

%\newcommand{\ip}[2]{\langle #1 | #2\rangle}
                             % This is an alternative definition of
                             % \ip that is commented out.

\begin{document}             % End of preamble and beginning of text.

\maketitle                   % Produces the title.

This is an example input file.  Comparing it with
the output it generates can show you how to
produce a simple document of your own.

\section{Ordinary Text}      % Produces section heading.  Lower-level
                             % sections are begun with similar
                             % \subsection and \subsubsection commands.

The ends  of words and sentences are marked
  by   spaces. It  doesn't matter how many
spaces    you type; one is as good as 100.  The
end of   a line counts as a space.

One   or more   blank lines denote the  end
of  a paragraph.

Since any number of consecutive spaces are treated
like a single one, the formatting of the input
file makes no difference to
      \LaTeX,                % The \LaTeX command generates the LaTeX logo.
but it makes a difference to you.  When you use
\LaTeX, making your input file as easy to read
as possible will be a great help as you write
your document and when you change it.  This sample
file shows how you can add comments to your own input
file.

Because printing is different from typewriting,
there are a number of things that you have to do
differently when preparing an input file than if
you were just typing the document directly.
Quotation marks like
       ``this''
have to be handled specially, as do quotes within
quotes:
       ``\,`this'            % \, separates the double and single quote.
        is what I just
        wrote, not  `that'\,''.

Dashes come in three sizes: an
       intra-word
dash, a medium dash for number ranges like
       1--2,
and a punctuation
       dash---like
this.

A sentence-ending space should be larger than the
space between words within a sentence.  You
sometimes have to type special commands in
conjunction with punctuation characters to get
this right, as in the following sentence.
       Gnats, gnus, etc.\ all  % `\ ' makes an inter-word space.
       begin with G\@.         % \@ marks end-of-sentence punctuation.
You should check the spaces after periods when
reading your output to make sure you haven't
forgotten any special cases.  Generating an
ellipsis
       \ldots\               % `\ ' is needed after `\ldots' because TeX
                             % ignores spaces after command names like \ldots
                             % made from \ + letters.
                             %
                             % Note how a `%' character causes TeX to ignore
                             % the end of the input line, so these blank lines
                             % do not start a new paragraph.
                             %
with the right spacing around the periods requires
a special command.

\LaTeX\ interprets some common characters as
commands, so you must type special commands to
generate them.  These characters include the
following:
       \$ \& \% \# \{ and \}.

In printing, text is usually emphasized with an
       \emph{italic}
type style.

\begin{em}
   A long segment of text can also be emphasized
   in this way.  Text within such a segment can be
   given \emph{additional} emphasis.
\end{em}

It is sometimes necessary to prevent \LaTeX\ from
breaking a line where it might otherwise do so.
This may be at a space, as between the ``Mr.''\ and
``Jones'' in
       ``Mr.~Jones'',        % ~ produces an unbreakable interword space.
or within a word---especially when the word is a
symbol like
       \mbox{\emph{itemnum}}
that makes little sense when hyphenated across
lines.

Footnotes\footnote{This is an example of a footnote.}
pose no problem.

\LaTeX\ is good at typesetting mathematical formulas
like
       \( x-3y + z = 7 \)
or
       \( a_{1} > x^{2n} + y^{2n} > x' \)
or
       \( \ip{A}{B} = \sum_{i} a_{i} b_{i} \).
The spaces you type in a formula are
ignored.  Remember that a letter like
       $x$                   % $ ... $  and  \( ... \)  are equivalent
is a formula when it denotes a mathematical
symbol, and it should be typed as one.

\section{Displayed Text}

Text is displayed by indenting it from the left
margin.  Quotations are commonly displayed.  There
are short quotations
\begin{quote}
   This is a short quotation.  It consists of a
   single paragraph of text.  See how it is formatted.
\end{quote}
and longer ones.
\begin{quotation}
   This is a longer quotation.  It consists of two
   paragraphs of text, neither of which are
   particularly interesting.

   This is the second paragraph of the quotation.  It
   is just as dull as the first paragraph.
\end{quotation}
Another frequently-displayed structure is a list.
The following is an example of an \emph{itemized}
list.
\begin{itemize}
   \item This is the first item of an itemized list.
         Each item in the list is marked with a ``tick''.
         You don't have to worry about what kind of tick
         mark is used.

   \item This is the second item of the list.  It
         contains another list nested inside it.  The inner
         list is an \emph{enumerated} list.
         \begin{enumerate}
            \item This is the first item of an enumerated
                  list that is nested within the itemized list.

            \item This is the second item of the inner list.
                  \LaTeX\ allows you to nest lists deeper than
                  you really should.
         \end{enumerate}
         This is the rest of the second item of the outer
         list.  It is no more interesting than any other
         part of the item.
   \item This is the third item of the list.
\end{itemize}
You can even display poetry.
\begin{verse}
   There is an environment
    for verse \\             % The \\ command separates lines
   Whose features some poets % within a stanza.
   will curse.

                             % One or more blank lines separate stanzas.

   For instead of making\\
   Them do \emph{all} line breaking, \\
   It allows them to put too many words on a line when they'd rather be
   forced to be terse.
\end{verse}

Mathematical formulas may also be displayed.  A
displayed formula
is
one-line long; multiline
formulas require special formatting instructions.
   \[  \ip{\Gamma}{\psi'} = x'' + y^{2} + z_{i}^{n}\]
Don't start a paragraph with a displayed equation,
nor make one a paragraph by itself.

\end{document}               % End of document.

\section{Bench Huge Section 993}
% synthetic bench section 993
% This is a sample LaTeX input file.  (Version of 12 August 2004.)
%
% A '%' character causes TeX to ignore all remaining text on the line,
% and is used for comments like this one.

\documentclass{article}      % Specifies the document class

                             % The preamble begins here.
\title{An Example Document}  % Declares the document's title.
\author{Leslie Lamport}      % Declares the author's name.
\date{January 21, 1994}      % Deleting this command produces today's date.

\newcommand{\ip}[2]{(#1, #2)}
                             % Defines \ip{arg1}{arg2} to mean
                             % (arg1, arg2).

%\newcommand{\ip}[2]{\langle #1 | #2\rangle}
                             % This is an alternative definition of
                             % \ip that is commented out.

\begin{document}             % End of preamble and beginning of text.

\maketitle                   % Produces the title.

This is an example input file.  Comparing it with
the output it generates can show you how to
produce a simple document of your own.

\section{Ordinary Text}      % Produces section heading.  Lower-level
                             % sections are begun with similar
                             % \subsection and \subsubsection commands.

The ends  of words and sentences are marked
  by   spaces. It  doesn't matter how many
spaces    you type; one is as good as 100.  The
end of   a line counts as a space.

One   or more   blank lines denote the  end
of  a paragraph.

Since any number of consecutive spaces are treated
like a single one, the formatting of the input
file makes no difference to
      \LaTeX,                % The \LaTeX command generates the LaTeX logo.
but it makes a difference to you.  When you use
\LaTeX, making your input file as easy to read
as possible will be a great help as you write
your document and when you change it.  This sample
file shows how you can add comments to your own input
file.

Because printing is different from typewriting,
there are a number of things that you have to do
differently when preparing an input file than if
you were just typing the document directly.
Quotation marks like
       ``this''
have to be handled specially, as do quotes within
quotes:
       ``\,`this'            % \, separates the double and single quote.
        is what I just
        wrote, not  `that'\,''.

Dashes come in three sizes: an
       intra-word
dash, a medium dash for number ranges like
       1--2,
and a punctuation
       dash---like
this.

A sentence-ending space should be larger than the
space between words within a sentence.  You
sometimes have to type special commands in
conjunction with punctuation characters to get
this right, as in the following sentence.
       Gnats, gnus, etc.\ all  % `\ ' makes an inter-word space.
       begin with G\@.         % \@ marks end-of-sentence punctuation.
You should check the spaces after periods when
reading your output to make sure you haven't
forgotten any special cases.  Generating an
ellipsis
       \ldots\               % `\ ' is needed after `\ldots' because TeX
                             % ignores spaces after command names like \ldots
                             % made from \ + letters.
                             %
                             % Note how a `%' character causes TeX to ignore
                             % the end of the input line, so these blank lines
                             % do not start a new paragraph.
                             %
with the right spacing around the periods requires
a special command.

\LaTeX\ interprets some common characters as
commands, so you must type special commands to
generate them.  These characters include the
following:
       \$ \& \% \# \{ and \}.

In printing, text is usually emphasized with an
       \emph{italic}
type style.

\begin{em}
   A long segment of text can also be emphasized
   in this way.  Text within such a segment can be
   given \emph{additional} emphasis.
\end{em}

It is sometimes necessary to prevent \LaTeX\ from
breaking a line where it might otherwise do so.
This may be at a space, as between the ``Mr.''\ and
``Jones'' in
       ``Mr.~Jones'',        % ~ produces an unbreakable interword space.
or within a word---especially when the word is a
symbol like
       \mbox{\emph{itemnum}}
that makes little sense when hyphenated across
lines.

Footnotes\footnote{This is an example of a footnote.}
pose no problem.

\LaTeX\ is good at typesetting mathematical formulas
like
       \( x-3y + z = 7 \)
or
       \( a_{1} > x^{2n} + y^{2n} > x' \)
or
       \( \ip{A}{B} = \sum_{i} a_{i} b_{i} \).
The spaces you type in a formula are
ignored.  Remember that a letter like
       $x$                   % $ ... $  and  \( ... \)  are equivalent
is a formula when it denotes a mathematical
symbol, and it should be typed as one.

\section{Displayed Text}

Text is displayed by indenting it from the left
margin.  Quotations are commonly displayed.  There
are short quotations
\begin{quote}
   This is a short quotation.  It consists of a
   single paragraph of text.  See how it is formatted.
\end{quote}
and longer ones.
\begin{quotation}
   This is a longer quotation.  It consists of two
   paragraphs of text, neither of which are
   particularly interesting.

   This is the second paragraph of the quotation.  It
   is just as dull as the first paragraph.
\end{quotation}
Another frequently-displayed structure is a list.
The following is an example of an \emph{itemized}
list.
\begin{itemize}
   \item This is the first item of an itemized list.
         Each item in the list is marked with a ``tick''.
         You don't have to worry about what kind of tick
         mark is used.

   \item This is the second item of the list.  It
         contains another list nested inside it.  The inner
         list is an \emph{enumerated} list.
         \begin{enumerate}
            \item This is the first item of an enumerated
                  list that is nested within the itemized list.

            \item This is the second item of the inner list.
                  \LaTeX\ allows you to nest lists deeper than
                  you really should.
         \end{enumerate}
         This is the rest of the second item of the outer
         list.  It is no more interesting than any other
         part of the item.
   \item This is the third item of the list.
\end{itemize}
You can even display poetry.
\begin{verse}
   There is an environment
    for verse \\             % The \\ command separates lines
   Whose features some poets % within a stanza.
   will curse.

                             % One or more blank lines separate stanzas.

   For instead of making\\
   Them do \emph{all} line breaking, \\
   It allows them to put too many words on a line when they'd rather be
   forced to be terse.
\end{verse}

Mathematical formulas may also be displayed.  A
displayed formula
is
one-line long; multiline
formulas require special formatting instructions.
   \[  \ip{\Gamma}{\psi'} = x'' + y^{2} + z_{i}^{n}\]
Don't start a paragraph with a displayed equation,
nor make one a paragraph by itself.

\end{document}               % End of document.

\section{Bench Huge Section 994}
% synthetic bench section 994
% This is a sample LaTeX input file.  (Version of 12 August 2004.)
%
% A '%' character causes TeX to ignore all remaining text on the line,
% and is used for comments like this one.

\documentclass{article}      % Specifies the document class

                             % The preamble begins here.
\title{An Example Document}  % Declares the document's title.
\author{Leslie Lamport}      % Declares the author's name.
\date{January 21, 1994}      % Deleting this command produces today's date.

\newcommand{\ip}[2]{(#1, #2)}
                             % Defines \ip{arg1}{arg2} to mean
                             % (arg1, arg2).

%\newcommand{\ip}[2]{\langle #1 | #2\rangle}
                             % This is an alternative definition of
                             % \ip that is commented out.

\begin{document}             % End of preamble and beginning of text.

\maketitle                   % Produces the title.

This is an example input file.  Comparing it with
the output it generates can show you how to
produce a simple document of your own.

\section{Ordinary Text}      % Produces section heading.  Lower-level
                             % sections are begun with similar
                             % \subsection and \subsubsection commands.

The ends  of words and sentences are marked
  by   spaces. It  doesn't matter how many
spaces    you type; one is as good as 100.  The
end of   a line counts as a space.

One   or more   blank lines denote the  end
of  a paragraph.

Since any number of consecutive spaces are treated
like a single one, the formatting of the input
file makes no difference to
      \LaTeX,                % The \LaTeX command generates the LaTeX logo.
but it makes a difference to you.  When you use
\LaTeX, making your input file as easy to read
as possible will be a great help as you write
your document and when you change it.  This sample
file shows how you can add comments to your own input
file.

Because printing is different from typewriting,
there are a number of things that you have to do
differently when preparing an input file than if
you were just typing the document directly.
Quotation marks like
       ``this''
have to be handled specially, as do quotes within
quotes:
       ``\,`this'            % \, separates the double and single quote.
        is what I just
        wrote, not  `that'\,''.

Dashes come in three sizes: an
       intra-word
dash, a medium dash for number ranges like
       1--2,
and a punctuation
       dash---like
this.

A sentence-ending space should be larger than the
space between words within a sentence.  You
sometimes have to type special commands in
conjunction with punctuation characters to get
this right, as in the following sentence.
       Gnats, gnus, etc.\ all  % `\ ' makes an inter-word space.
       begin with G\@.         % \@ marks end-of-sentence punctuation.
You should check the spaces after periods when
reading your output to make sure you haven't
forgotten any special cases.  Generating an
ellipsis
       \ldots\               % `\ ' is needed after `\ldots' because TeX
                             % ignores spaces after command names like \ldots
                             % made from \ + letters.
                             %
                             % Note how a `%' character causes TeX to ignore
                             % the end of the input line, so these blank lines
                             % do not start a new paragraph.
                             %
with the right spacing around the periods requires
a special command.

\LaTeX\ interprets some common characters as
commands, so you must type special commands to
generate them.  These characters include the
following:
       \$ \& \% \# \{ and \}.

In printing, text is usually emphasized with an
       \emph{italic}
type style.

\begin{em}
   A long segment of text can also be emphasized
   in this way.  Text within such a segment can be
   given \emph{additional} emphasis.
\end{em}

It is sometimes necessary to prevent \LaTeX\ from
breaking a line where it might otherwise do so.
This may be at a space, as between the ``Mr.''\ and
``Jones'' in
       ``Mr.~Jones'',        % ~ produces an unbreakable interword space.
or within a word---especially when the word is a
symbol like
       \mbox{\emph{itemnum}}
that makes little sense when hyphenated across
lines.

Footnotes\footnote{This is an example of a footnote.}
pose no problem.

\LaTeX\ is good at typesetting mathematical formulas
like
       \( x-3y + z = 7 \)
or
       \( a_{1} > x^{2n} + y^{2n} > x' \)
or
       \( \ip{A}{B} = \sum_{i} a_{i} b_{i} \).
The spaces you type in a formula are
ignored.  Remember that a letter like
       $x$                   % $ ... $  and  \( ... \)  are equivalent
is a formula when it denotes a mathematical
symbol, and it should be typed as one.

\section{Displayed Text}

Text is displayed by indenting it from the left
margin.  Quotations are commonly displayed.  There
are short quotations
\begin{quote}
   This is a short quotation.  It consists of a
   single paragraph of text.  See how it is formatted.
\end{quote}
and longer ones.
\begin{quotation}
   This is a longer quotation.  It consists of two
   paragraphs of text, neither of which are
   particularly interesting.

   This is the second paragraph of the quotation.  It
   is just as dull as the first paragraph.
\end{quotation}
Another frequently-displayed structure is a list.
The following is an example of an \emph{itemized}
list.
\begin{itemize}
   \item This is the first item of an itemized list.
         Each item in the list is marked with a ``tick''.
         You don't have to worry about what kind of tick
         mark is used.

   \item This is the second item of the list.  It
         contains another list nested inside it.  The inner
         list is an \emph{enumerated} list.
         \begin{enumerate}
            \item This is the first item of an enumerated
                  list that is nested within the itemized list.

            \item This is the second item of the inner list.
                  \LaTeX\ allows you to nest lists deeper than
                  you really should.
         \end{enumerate}
         This is the rest of the second item of the outer
         list.  It is no more interesting than any other
         part of the item.
   \item This is the third item of the list.
\end{itemize}
You can even display poetry.
\begin{verse}
   There is an environment
    for verse \\             % The \\ command separates lines
   Whose features some poets % within a stanza.
   will curse.

                             % One or more blank lines separate stanzas.

   For instead of making\\
   Them do \emph{all} line breaking, \\
   It allows them to put too many words on a line when they'd rather be
   forced to be terse.
\end{verse}

Mathematical formulas may also be displayed.  A
displayed formula
is
one-line long; multiline
formulas require special formatting instructions.
   \[  \ip{\Gamma}{\psi'} = x'' + y^{2} + z_{i}^{n}\]
Don't start a paragraph with a displayed equation,
nor make one a paragraph by itself.

\end{document}               % End of document.

\section{Bench Huge Section 995}
% synthetic bench section 995
% This is a sample LaTeX input file.  (Version of 12 August 2004.)
%
% A '%' character causes TeX to ignore all remaining text on the line,
% and is used for comments like this one.

\documentclass{article}      % Specifies the document class

                             % The preamble begins here.
\title{An Example Document}  % Declares the document's title.
\author{Leslie Lamport}      % Declares the author's name.
\date{January 21, 1994}      % Deleting this command produces today's date.

\newcommand{\ip}[2]{(#1, #2)}
                             % Defines \ip{arg1}{arg2} to mean
                             % (arg1, arg2).

%\newcommand{\ip}[2]{\langle #1 | #2\rangle}
                             % This is an alternative definition of
                             % \ip that is commented out.

\begin{document}             % End of preamble and beginning of text.

\maketitle                   % Produces the title.

This is an example input file.  Comparing it with
the output it generates can show you how to
produce a simple document of your own.

\section{Ordinary Text}      % Produces section heading.  Lower-level
                             % sections are begun with similar
                             % \subsection and \subsubsection commands.

The ends  of words and sentences are marked
  by   spaces. It  doesn't matter how many
spaces    you type; one is as good as 100.  The
end of   a line counts as a space.

One   or more   blank lines denote the  end
of  a paragraph.

Since any number of consecutive spaces are treated
like a single one, the formatting of the input
file makes no difference to
      \LaTeX,                % The \LaTeX command generates the LaTeX logo.
but it makes a difference to you.  When you use
\LaTeX, making your input file as easy to read
as possible will be a great help as you write
your document and when you change it.  This sample
file shows how you can add comments to your own input
file.

Because printing is different from typewriting,
there are a number of things that you have to do
differently when preparing an input file than if
you were just typing the document directly.
Quotation marks like
       ``this''
have to be handled specially, as do quotes within
quotes:
       ``\,`this'            % \, separates the double and single quote.
        is what I just
        wrote, not  `that'\,''.

Dashes come in three sizes: an
       intra-word
dash, a medium dash for number ranges like
       1--2,
and a punctuation
       dash---like
this.

A sentence-ending space should be larger than the
space between words within a sentence.  You
sometimes have to type special commands in
conjunction with punctuation characters to get
this right, as in the following sentence.
       Gnats, gnus, etc.\ all  % `\ ' makes an inter-word space.
       begin with G\@.         % \@ marks end-of-sentence punctuation.
You should check the spaces after periods when
reading your output to make sure you haven't
forgotten any special cases.  Generating an
ellipsis
       \ldots\               % `\ ' is needed after `\ldots' because TeX
                             % ignores spaces after command names like \ldots
                             % made from \ + letters.
                             %
                             % Note how a `%' character causes TeX to ignore
                             % the end of the input line, so these blank lines
                             % do not start a new paragraph.
                             %
with the right spacing around the periods requires
a special command.

\LaTeX\ interprets some common characters as
commands, so you must type special commands to
generate them.  These characters include the
following:
       \$ \& \% \# \{ and \}.

In printing, text is usually emphasized with an
       \emph{italic}
type style.

\begin{em}
   A long segment of text can also be emphasized
   in this way.  Text within such a segment can be
   given \emph{additional} emphasis.
\end{em}

It is sometimes necessary to prevent \LaTeX\ from
breaking a line where it might otherwise do so.
This may be at a space, as between the ``Mr.''\ and
``Jones'' in
       ``Mr.~Jones'',        % ~ produces an unbreakable interword space.
or within a word---especially when the word is a
symbol like
       \mbox{\emph{itemnum}}
that makes little sense when hyphenated across
lines.

Footnotes\footnote{This is an example of a footnote.}
pose no problem.

\LaTeX\ is good at typesetting mathematical formulas
like
       \( x-3y + z = 7 \)
or
       \( a_{1} > x^{2n} + y^{2n} > x' \)
or
       \( \ip{A}{B} = \sum_{i} a_{i} b_{i} \).
The spaces you type in a formula are
ignored.  Remember that a letter like
       $x$                   % $ ... $  and  \( ... \)  are equivalent
is a formula when it denotes a mathematical
symbol, and it should be typed as one.

\section{Displayed Text}

Text is displayed by indenting it from the left
margin.  Quotations are commonly displayed.  There
are short quotations
\begin{quote}
   This is a short quotation.  It consists of a
   single paragraph of text.  See how it is formatted.
\end{quote}
and longer ones.
\begin{quotation}
   This is a longer quotation.  It consists of two
   paragraphs of text, neither of which are
   particularly interesting.

   This is the second paragraph of the quotation.  It
   is just as dull as the first paragraph.
\end{quotation}
Another frequently-displayed structure is a list.
The following is an example of an \emph{itemized}
list.
\begin{itemize}
   \item This is the first item of an itemized list.
         Each item in the list is marked with a ``tick''.
         You don't have to worry about what kind of tick
         mark is used.

   \item This is the second item of the list.  It
         contains another list nested inside it.  The inner
         list is an \emph{enumerated} list.
         \begin{enumerate}
            \item This is the first item of an enumerated
                  list that is nested within the itemized list.

            \item This is the second item of the inner list.
                  \LaTeX\ allows you to nest lists deeper than
                  you really should.
         \end{enumerate}
         This is the rest of the second item of the outer
         list.  It is no more interesting than any other
         part of the item.
   \item This is the third item of the list.
\end{itemize}
You can even display poetry.
\begin{verse}
   There is an environment
    for verse \\             % The \\ command separates lines
   Whose features some poets % within a stanza.
   will curse.

                             % One or more blank lines separate stanzas.

   For instead of making\\
   Them do \emph{all} line breaking, \\
   It allows them to put too many words on a line when they'd rather be
   forced to be terse.
\end{verse}

Mathematical formulas may also be displayed.  A
displayed formula
is
one-line long; multiline
formulas require special formatting instructions.
   \[  \ip{\Gamma}{\psi'} = x'' + y^{2} + z_{i}^{n}\]
Don't start a paragraph with a displayed equation,
nor make one a paragraph by itself.

\end{document}               % End of document.

\section{Bench Huge Section 996}
% synthetic bench section 996
% This is a sample LaTeX input file.  (Version of 12 August 2004.)
%
% A '%' character causes TeX to ignore all remaining text on the line,
% and is used for comments like this one.

\documentclass{article}      % Specifies the document class

                             % The preamble begins here.
\title{An Example Document}  % Declares the document's title.
\author{Leslie Lamport}      % Declares the author's name.
\date{January 21, 1994}      % Deleting this command produces today's date.

\newcommand{\ip}[2]{(#1, #2)}
                             % Defines \ip{arg1}{arg2} to mean
                             % (arg1, arg2).

%\newcommand{\ip}[2]{\langle #1 | #2\rangle}
                             % This is an alternative definition of
                             % \ip that is commented out.

\begin{document}             % End of preamble and beginning of text.

\maketitle                   % Produces the title.

This is an example input file.  Comparing it with
the output it generates can show you how to
produce a simple document of your own.

\section{Ordinary Text}      % Produces section heading.  Lower-level
                             % sections are begun with similar
                             % \subsection and \subsubsection commands.

The ends  of words and sentences are marked
  by   spaces. It  doesn't matter how many
spaces    you type; one is as good as 100.  The
end of   a line counts as a space.

One   or more   blank lines denote the  end
of  a paragraph.

Since any number of consecutive spaces are treated
like a single one, the formatting of the input
file makes no difference to
      \LaTeX,                % The \LaTeX command generates the LaTeX logo.
but it makes a difference to you.  When you use
\LaTeX, making your input file as easy to read
as possible will be a great help as you write
your document and when you change it.  This sample
file shows how you can add comments to your own input
file.

Because printing is different from typewriting,
there are a number of things that you have to do
differently when preparing an input file than if
you were just typing the document directly.
Quotation marks like
       ``this''
have to be handled specially, as do quotes within
quotes:
       ``\,`this'            % \, separates the double and single quote.
        is what I just
        wrote, not  `that'\,''.

Dashes come in three sizes: an
       intra-word
dash, a medium dash for number ranges like
       1--2,
and a punctuation
       dash---like
this.

A sentence-ending space should be larger than the
space between words within a sentence.  You
sometimes have to type special commands in
conjunction with punctuation characters to get
this right, as in the following sentence.
       Gnats, gnus, etc.\ all  % `\ ' makes an inter-word space.
       begin with G\@.         % \@ marks end-of-sentence punctuation.
You should check the spaces after periods when
reading your output to make sure you haven't
forgotten any special cases.  Generating an
ellipsis
       \ldots\               % `\ ' is needed after `\ldots' because TeX
                             % ignores spaces after command names like \ldots
                             % made from \ + letters.
                             %
                             % Note how a `%' character causes TeX to ignore
                             % the end of the input line, so these blank lines
                             % do not start a new paragraph.
                             %
with the right spacing around the periods requires
a special command.

\LaTeX\ interprets some common characters as
commands, so you must type special commands to
generate them.  These characters include the
following:
       \$ \& \% \# \{ and \}.

In printing, text is usually emphasized with an
       \emph{italic}
type style.

\begin{em}
   A long segment of text can also be emphasized
   in this way.  Text within such a segment can be
   given \emph{additional} emphasis.
\end{em}

It is sometimes necessary to prevent \LaTeX\ from
breaking a line where it might otherwise do so.
This may be at a space, as between the ``Mr.''\ and
``Jones'' in
       ``Mr.~Jones'',        % ~ produces an unbreakable interword space.
or within a word---especially when the word is a
symbol like
       \mbox{\emph{itemnum}}
that makes little sense when hyphenated across
lines.

Footnotes\footnote{This is an example of a footnote.}
pose no problem.

\LaTeX\ is good at typesetting mathematical formulas
like
       \( x-3y + z = 7 \)
or
       \( a_{1} > x^{2n} + y^{2n} > x' \)
or
       \( \ip{A}{B} = \sum_{i} a_{i} b_{i} \).
The spaces you type in a formula are
ignored.  Remember that a letter like
       $x$                   % $ ... $  and  \( ... \)  are equivalent
is a formula when it denotes a mathematical
symbol, and it should be typed as one.

\section{Displayed Text}

Text is displayed by indenting it from the left
margin.  Quotations are commonly displayed.  There
are short quotations
\begin{quote}
   This is a short quotation.  It consists of a
   single paragraph of text.  See how it is formatted.
\end{quote}
and longer ones.
\begin{quotation}
   This is a longer quotation.  It consists of two
   paragraphs of text, neither of which are
   particularly interesting.

   This is the second paragraph of the quotation.  It
   is just as dull as the first paragraph.
\end{quotation}
Another frequently-displayed structure is a list.
The following is an example of an \emph{itemized}
list.
\begin{itemize}
   \item This is the first item of an itemized list.
         Each item in the list is marked with a ``tick''.
         You don't have to worry about what kind of tick
         mark is used.

   \item This is the second item of the list.  It
         contains another list nested inside it.  The inner
         list is an \emph{enumerated} list.
         \begin{enumerate}
            \item This is the first item of an enumerated
                  list that is nested within the itemized list.

            \item This is the second item of the inner list.
                  \LaTeX\ allows you to nest lists deeper than
                  you really should.
         \end{enumerate}
         This is the rest of the second item of the outer
         list.  It is no more interesting than any other
         part of the item.
   \item This is the third item of the list.
\end{itemize}
You can even display poetry.
\begin{verse}
   There is an environment
    for verse \\             % The \\ command separates lines
   Whose features some poets % within a stanza.
   will curse.

                             % One or more blank lines separate stanzas.

   For instead of making\\
   Them do \emph{all} line breaking, \\
   It allows them to put too many words on a line when they'd rather be
   forced to be terse.
\end{verse}

Mathematical formulas may also be displayed.  A
displayed formula
is
one-line long; multiline
formulas require special formatting instructions.
   \[  \ip{\Gamma}{\psi'} = x'' + y^{2} + z_{i}^{n}\]
Don't start a paragraph with a displayed equation,
nor make one a paragraph by itself.

\end{document}               % End of document.

\section{Bench Huge Section 997}
% synthetic bench section 997
% This is a sample LaTeX input file.  (Version of 12 August 2004.)
%
% A '%' character causes TeX to ignore all remaining text on the line,
% and is used for comments like this one.

\documentclass{article}      % Specifies the document class

                             % The preamble begins here.
\title{An Example Document}  % Declares the document's title.
\author{Leslie Lamport}      % Declares the author's name.
\date{January 21, 1994}      % Deleting this command produces today's date.

\newcommand{\ip}[2]{(#1, #2)}
                             % Defines \ip{arg1}{arg2} to mean
                             % (arg1, arg2).

%\newcommand{\ip}[2]{\langle #1 | #2\rangle}
                             % This is an alternative definition of
                             % \ip that is commented out.

\begin{document}             % End of preamble and beginning of text.

\maketitle                   % Produces the title.

This is an example input file.  Comparing it with
the output it generates can show you how to
produce a simple document of your own.

\section{Ordinary Text}      % Produces section heading.  Lower-level
                             % sections are begun with similar
                             % \subsection and \subsubsection commands.

The ends  of words and sentences are marked
  by   spaces. It  doesn't matter how many
spaces    you type; one is as good as 100.  The
end of   a line counts as a space.

One   or more   blank lines denote the  end
of  a paragraph.

Since any number of consecutive spaces are treated
like a single one, the formatting of the input
file makes no difference to
      \LaTeX,                % The \LaTeX command generates the LaTeX logo.
but it makes a difference to you.  When you use
\LaTeX, making your input file as easy to read
as possible will be a great help as you write
your document and when you change it.  This sample
file shows how you can add comments to your own input
file.

Because printing is different from typewriting,
there are a number of things that you have to do
differently when preparing an input file than if
you were just typing the document directly.
Quotation marks like
       ``this''
have to be handled specially, as do quotes within
quotes:
       ``\,`this'            % \, separates the double and single quote.
        is what I just
        wrote, not  `that'\,''.

Dashes come in three sizes: an
       intra-word
dash, a medium dash for number ranges like
       1--2,
and a punctuation
       dash---like
this.

A sentence-ending space should be larger than the
space between words within a sentence.  You
sometimes have to type special commands in
conjunction with punctuation characters to get
this right, as in the following sentence.
       Gnats, gnus, etc.\ all  % `\ ' makes an inter-word space.
       begin with G\@.         % \@ marks end-of-sentence punctuation.
You should check the spaces after periods when
reading your output to make sure you haven't
forgotten any special cases.  Generating an
ellipsis
       \ldots\               % `\ ' is needed after `\ldots' because TeX
                             % ignores spaces after command names like \ldots
                             % made from \ + letters.
                             %
                             % Note how a `%' character causes TeX to ignore
                             % the end of the input line, so these blank lines
                             % do not start a new paragraph.
                             %
with the right spacing around the periods requires
a special command.

\LaTeX\ interprets some common characters as
commands, so you must type special commands to
generate them.  These characters include the
following:
       \$ \& \% \# \{ and \}.

In printing, text is usually emphasized with an
       \emph{italic}
type style.

\begin{em}
   A long segment of text can also be emphasized
   in this way.  Text within such a segment can be
   given \emph{additional} emphasis.
\end{em}

It is sometimes necessary to prevent \LaTeX\ from
breaking a line where it might otherwise do so.
This may be at a space, as between the ``Mr.''\ and
``Jones'' in
       ``Mr.~Jones'',        % ~ produces an unbreakable interword space.
or within a word---especially when the word is a
symbol like
       \mbox{\emph{itemnum}}
that makes little sense when hyphenated across
lines.

Footnotes\footnote{This is an example of a footnote.}
pose no problem.

\LaTeX\ is good at typesetting mathematical formulas
like
       \( x-3y + z = 7 \)
or
       \( a_{1} > x^{2n} + y^{2n} > x' \)
or
       \( \ip{A}{B} = \sum_{i} a_{i} b_{i} \).
The spaces you type in a formula are
ignored.  Remember that a letter like
       $x$                   % $ ... $  and  \( ... \)  are equivalent
is a formula when it denotes a mathematical
symbol, and it should be typed as one.

\section{Displayed Text}

Text is displayed by indenting it from the left
margin.  Quotations are commonly displayed.  There
are short quotations
\begin{quote}
   This is a short quotation.  It consists of a
   single paragraph of text.  See how it is formatted.
\end{quote}
and longer ones.
\begin{quotation}
   This is a longer quotation.  It consists of two
   paragraphs of text, neither of which are
   particularly interesting.

   This is the second paragraph of the quotation.  It
   is just as dull as the first paragraph.
\end{quotation}
Another frequently-displayed structure is a list.
The following is an example of an \emph{itemized}
list.
\begin{itemize}
   \item This is the first item of an itemized list.
         Each item in the list is marked with a ``tick''.
         You don't have to worry about what kind of tick
         mark is used.

   \item This is the second item of the list.  It
         contains another list nested inside it.  The inner
         list is an \emph{enumerated} list.
         \begin{enumerate}
            \item This is the first item of an enumerated
                  list that is nested within the itemized list.

            \item This is the second item of the inner list.
                  \LaTeX\ allows you to nest lists deeper than
                  you really should.
         \end{enumerate}
         This is the rest of the second item of the outer
         list.  It is no more interesting than any other
         part of the item.
   \item This is the third item of the list.
\end{itemize}
You can even display poetry.
\begin{verse}
   There is an environment
    for verse \\             % The \\ command separates lines
   Whose features some poets % within a stanza.
   will curse.

                             % One or more blank lines separate stanzas.

   For instead of making\\
   Them do \emph{all} line breaking, \\
   It allows them to put too many words on a line when they'd rather be
   forced to be terse.
\end{verse}

Mathematical formulas may also be displayed.  A
displayed formula
is
one-line long; multiline
formulas require special formatting instructions.
   \[  \ip{\Gamma}{\psi'} = x'' + y^{2} + z_{i}^{n}\]
Don't start a paragraph with a displayed equation,
nor make one a paragraph by itself.

\end{document}               % End of document.

\section{Bench Huge Section 998}
% synthetic bench section 998
% This is a sample LaTeX input file.  (Version of 12 August 2004.)
%
% A '%' character causes TeX to ignore all remaining text on the line,
% and is used for comments like this one.

\documentclass{article}      % Specifies the document class

                             % The preamble begins here.
\title{An Example Document}  % Declares the document's title.
\author{Leslie Lamport}      % Declares the author's name.
\date{January 21, 1994}      % Deleting this command produces today's date.

\newcommand{\ip}[2]{(#1, #2)}
                             % Defines \ip{arg1}{arg2} to mean
                             % (arg1, arg2).

%\newcommand{\ip}[2]{\langle #1 | #2\rangle}
                             % This is an alternative definition of
                             % \ip that is commented out.

\begin{document}             % End of preamble and beginning of text.

\maketitle                   % Produces the title.

This is an example input file.  Comparing it with
the output it generates can show you how to
produce a simple document of your own.

\section{Ordinary Text}      % Produces section heading.  Lower-level
                             % sections are begun with similar
                             % \subsection and \subsubsection commands.

The ends  of words and sentences are marked
  by   spaces. It  doesn't matter how many
spaces    you type; one is as good as 100.  The
end of   a line counts as a space.

One   or more   blank lines denote the  end
of  a paragraph.

Since any number of consecutive spaces are treated
like a single one, the formatting of the input
file makes no difference to
      \LaTeX,                % The \LaTeX command generates the LaTeX logo.
but it makes a difference to you.  When you use
\LaTeX, making your input file as easy to read
as possible will be a great help as you write
your document and when you change it.  This sample
file shows how you can add comments to your own input
file.

Because printing is different from typewriting,
there are a number of things that you have to do
differently when preparing an input file than if
you were just typing the document directly.
Quotation marks like
       ``this''
have to be handled specially, as do quotes within
quotes:
       ``\,`this'            % \, separates the double and single quote.
        is what I just
        wrote, not  `that'\,''.

Dashes come in three sizes: an
       intra-word
dash, a medium dash for number ranges like
       1--2,
and a punctuation
       dash---like
this.

A sentence-ending space should be larger than the
space between words within a sentence.  You
sometimes have to type special commands in
conjunction with punctuation characters to get
this right, as in the following sentence.
       Gnats, gnus, etc.\ all  % `\ ' makes an inter-word space.
       begin with G\@.         % \@ marks end-of-sentence punctuation.
You should check the spaces after periods when
reading your output to make sure you haven't
forgotten any special cases.  Generating an
ellipsis
       \ldots\               % `\ ' is needed after `\ldots' because TeX
                             % ignores spaces after command names like \ldots
                             % made from \ + letters.
                             %
                             % Note how a `%' character causes TeX to ignore
                             % the end of the input line, so these blank lines
                             % do not start a new paragraph.
                             %
with the right spacing around the periods requires
a special command.

\LaTeX\ interprets some common characters as
commands, so you must type special commands to
generate them.  These characters include the
following:
       \$ \& \% \# \{ and \}.

In printing, text is usually emphasized with an
       \emph{italic}
type style.

\begin{em}
   A long segment of text can also be emphasized
   in this way.  Text within such a segment can be
   given \emph{additional} emphasis.
\end{em}

It is sometimes necessary to prevent \LaTeX\ from
breaking a line where it might otherwise do so.
This may be at a space, as between the ``Mr.''\ and
``Jones'' in
       ``Mr.~Jones'',        % ~ produces an unbreakable interword space.
or within a word---especially when the word is a
symbol like
       \mbox{\emph{itemnum}}
that makes little sense when hyphenated across
lines.

Footnotes\footnote{This is an example of a footnote.}
pose no problem.

\LaTeX\ is good at typesetting mathematical formulas
like
       \( x-3y + z = 7 \)
or
       \( a_{1} > x^{2n} + y^{2n} > x' \)
or
       \( \ip{A}{B} = \sum_{i} a_{i} b_{i} \).
The spaces you type in a formula are
ignored.  Remember that a letter like
       $x$                   % $ ... $  and  \( ... \)  are equivalent
is a formula when it denotes a mathematical
symbol, and it should be typed as one.

\section{Displayed Text}

Text is displayed by indenting it from the left
margin.  Quotations are commonly displayed.  There
are short quotations
\begin{quote}
   This is a short quotation.  It consists of a
   single paragraph of text.  See how it is formatted.
\end{quote}
and longer ones.
\begin{quotation}
   This is a longer quotation.  It consists of two
   paragraphs of text, neither of which are
   particularly interesting.

   This is the second paragraph of the quotation.  It
   is just as dull as the first paragraph.
\end{quotation}
Another frequently-displayed structure is a list.
The following is an example of an \emph{itemized}
list.
\begin{itemize}
   \item This is the first item of an itemized list.
         Each item in the list is marked with a ``tick''.
         You don't have to worry about what kind of tick
         mark is used.

   \item This is the second item of the list.  It
         contains another list nested inside it.  The inner
         list is an \emph{enumerated} list.
         \begin{enumerate}
            \item This is the first item of an enumerated
                  list that is nested within the itemized list.

            \item This is the second item of the inner list.
                  \LaTeX\ allows you to nest lists deeper than
                  you really should.
         \end{enumerate}
         This is the rest of the second item of the outer
         list.  It is no more interesting than any other
         part of the item.
   \item This is the third item of the list.
\end{itemize}
You can even display poetry.
\begin{verse}
   There is an environment
    for verse \\             % The \\ command separates lines
   Whose features some poets % within a stanza.
   will curse.

                             % One or more blank lines separate stanzas.

   For instead of making\\
   Them do \emph{all} line breaking, \\
   It allows them to put too many words on a line when they'd rather be
   forced to be terse.
\end{verse}

Mathematical formulas may also be displayed.  A
displayed formula
is
one-line long; multiline
formulas require special formatting instructions.
   \[  \ip{\Gamma}{\psi'} = x'' + y^{2} + z_{i}^{n}\]
Don't start a paragraph with a displayed equation,
nor make one a paragraph by itself.

\end{document}               % End of document.

\section{Bench Huge Section 999}
% synthetic bench section 999
% This is a sample LaTeX input file.  (Version of 12 August 2004.)
%
% A '%' character causes TeX to ignore all remaining text on the line,
% and is used for comments like this one.

\documentclass{article}      % Specifies the document class

                             % The preamble begins here.
\title{An Example Document}  % Declares the document's title.
\author{Leslie Lamport}      % Declares the author's name.
\date{January 21, 1994}      % Deleting this command produces today's date.

\newcommand{\ip}[2]{(#1, #2)}
                             % Defines \ip{arg1}{arg2} to mean
                             % (arg1, arg2).

%\newcommand{\ip}[2]{\langle #1 | #2\rangle}
                             % This is an alternative definition of
                             % \ip that is commented out.

\begin{document}             % End of preamble and beginning of text.

\maketitle                   % Produces the title.

This is an example input file.  Comparing it with
the output it generates can show you how to
produce a simple document of your own.

\section{Ordinary Text}      % Produces section heading.  Lower-level
                             % sections are begun with similar
                             % \subsection and \subsubsection commands.

The ends  of words and sentences are marked
  by   spaces. It  doesn't matter how many
spaces    you type; one is as good as 100.  The
end of   a line counts as a space.

One   or more   blank lines denote the  end
of  a paragraph.

Since any number of consecutive spaces are treated
like a single one, the formatting of the input
file makes no difference to
      \LaTeX,                % The \LaTeX command generates the LaTeX logo.
but it makes a difference to you.  When you use
\LaTeX, making your input file as easy to read
as possible will be a great help as you write
your document and when you change it.  This sample
file shows how you can add comments to your own input
file.

Because printing is different from typewriting,
there are a number of things that you have to do
differently when preparing an input file than if
you were just typing the document directly.
Quotation marks like
       ``this''
have to be handled specially, as do quotes within
quotes:
       ``\,`this'            % \, separates the double and single quote.
        is what I just
        wrote, not  `that'\,''.

Dashes come in three sizes: an
       intra-word
dash, a medium dash for number ranges like
       1--2,
and a punctuation
       dash---like
this.

A sentence-ending space should be larger than the
space between words within a sentence.  You
sometimes have to type special commands in
conjunction with punctuation characters to get
this right, as in the following sentence.
       Gnats, gnus, etc.\ all  % `\ ' makes an inter-word space.
       begin with G\@.         % \@ marks end-of-sentence punctuation.
You should check the spaces after periods when
reading your output to make sure you haven't
forgotten any special cases.  Generating an
ellipsis
       \ldots\               % `\ ' is needed after `\ldots' because TeX
                             % ignores spaces after command names like \ldots
                             % made from \ + letters.
                             %
                             % Note how a `%' character causes TeX to ignore
                             % the end of the input line, so these blank lines
                             % do not start a new paragraph.
                             %
with the right spacing around the periods requires
a special command.

\LaTeX\ interprets some common characters as
commands, so you must type special commands to
generate them.  These characters include the
following:
       \$ \& \% \# \{ and \}.

In printing, text is usually emphasized with an
       \emph{italic}
type style.

\begin{em}
   A long segment of text can also be emphasized
   in this way.  Text within such a segment can be
   given \emph{additional} emphasis.
\end{em}

It is sometimes necessary to prevent \LaTeX\ from
breaking a line where it might otherwise do so.
This may be at a space, as between the ``Mr.''\ and
``Jones'' in
       ``Mr.~Jones'',        % ~ produces an unbreakable interword space.
or within a word---especially when the word is a
symbol like
       \mbox{\emph{itemnum}}
that makes little sense when hyphenated across
lines.

Footnotes\footnote{This is an example of a footnote.}
pose no problem.

\LaTeX\ is good at typesetting mathematical formulas
like
       \( x-3y + z = 7 \)
or
       \( a_{1} > x^{2n} + y^{2n} > x' \)
or
       \( \ip{A}{B} = \sum_{i} a_{i} b_{i} \).
The spaces you type in a formula are
ignored.  Remember that a letter like
       $x$                   % $ ... $  and  \( ... \)  are equivalent
is a formula when it denotes a mathematical
symbol, and it should be typed as one.

\section{Displayed Text}

Text is displayed by indenting it from the left
margin.  Quotations are commonly displayed.  There
are short quotations
\begin{quote}
   This is a short quotation.  It consists of a
   single paragraph of text.  See how it is formatted.
\end{quote}
and longer ones.
\begin{quotation}
   This is a longer quotation.  It consists of two
   paragraphs of text, neither of which are
   particularly interesting.

   This is the second paragraph of the quotation.  It
   is just as dull as the first paragraph.
\end{quotation}
Another frequently-displayed structure is a list.
The following is an example of an \emph{itemized}
list.
\begin{itemize}
   \item This is the first item of an itemized list.
         Each item in the list is marked with a ``tick''.
         You don't have to worry about what kind of tick
         mark is used.

   \item This is the second item of the list.  It
         contains another list nested inside it.  The inner
         list is an \emph{enumerated} list.
         \begin{enumerate}
            \item This is the first item of an enumerated
                  list that is nested within the itemized list.

            \item This is the second item of the inner list.
                  \LaTeX\ allows you to nest lists deeper than
                  you really should.
         \end{enumerate}
         This is the rest of the second item of the outer
         list.  It is no more interesting than any other
         part of the item.
   \item This is the third item of the list.
\end{itemize}
You can even display poetry.
\begin{verse}
   There is an environment
    for verse \\             % The \\ command separates lines
   Whose features some poets % within a stanza.
   will curse.

                             % One or more blank lines separate stanzas.

   For instead of making\\
   Them do \emph{all} line breaking, \\
   It allows them to put too many words on a line when they'd rather be
   forced to be terse.
\end{verse}

Mathematical formulas may also be displayed.  A
displayed formula
is
one-line long; multiline
formulas require special formatting instructions.
   \[  \ip{\Gamma}{\psi'} = x'' + y^{2} + z_{i}^{n}\]
Don't start a paragraph with a displayed equation,
nor make one a paragraph by itself.

\end{document}               % End of document.

\section{Bench Huge Section 1000}
% synthetic bench section 1000
% This is a sample LaTeX input file.  (Version of 12 August 2004.)
%
% A '%' character causes TeX to ignore all remaining text on the line,
% and is used for comments like this one.

\documentclass{article}      % Specifies the document class

                             % The preamble begins here.
\title{An Example Document}  % Declares the document's title.
\author{Leslie Lamport}      % Declares the author's name.
\date{January 21, 1994}      % Deleting this command produces today's date.

\newcommand{\ip}[2]{(#1, #2)}
                             % Defines \ip{arg1}{arg2} to mean
                             % (arg1, arg2).

%\newcommand{\ip}[2]{\langle #1 | #2\rangle}
                             % This is an alternative definition of
                             % \ip that is commented out.

\begin{document}             % End of preamble and beginning of text.

\maketitle                   % Produces the title.

This is an example input file.  Comparing it with
the output it generates can show you how to
produce a simple document of your own.

\section{Ordinary Text}      % Produces section heading.  Lower-level
                             % sections are begun with similar
                             % \subsection and \subsubsection commands.

The ends  of words and sentences are marked
  by   spaces. It  doesn't matter how many
spaces    you type; one is as good as 100.  The
end of   a line counts as a space.

One   or more   blank lines denote the  end
of  a paragraph.

Since any number of consecutive spaces are treated
like a single one, the formatting of the input
file makes no difference to
      \LaTeX,                % The \LaTeX command generates the LaTeX logo.
but it makes a difference to you.  When you use
\LaTeX, making your input file as easy to read
as possible will be a great help as you write
your document and when you change it.  This sample
file shows how you can add comments to your own input
file.

Because printing is different from typewriting,
there are a number of things that you have to do
differently when preparing an input file than if
you were just typing the document directly.
Quotation marks like
       ``this''
have to be handled specially, as do quotes within
quotes:
       ``\,`this'            % \, separates the double and single quote.
        is what I just
        wrote, not  `that'\,''.

Dashes come in three sizes: an
       intra-word
dash, a medium dash for number ranges like
       1--2,
and a punctuation
       dash---like
this.

A sentence-ending space should be larger than the
space between words within a sentence.  You
sometimes have to type special commands in
conjunction with punctuation characters to get
this right, as in the following sentence.
       Gnats, gnus, etc.\ all  % `\ ' makes an inter-word space.
       begin with G\@.         % \@ marks end-of-sentence punctuation.
You should check the spaces after periods when
reading your output to make sure you haven't
forgotten any special cases.  Generating an
ellipsis
       \ldots\               % `\ ' is needed after `\ldots' because TeX
                             % ignores spaces after command names like \ldots
                             % made from \ + letters.
                             %
                             % Note how a `%' character causes TeX to ignore
                             % the end of the input line, so these blank lines
                             % do not start a new paragraph.
                             %
with the right spacing around the periods requires
a special command.

\LaTeX\ interprets some common characters as
commands, so you must type special commands to
generate them.  These characters include the
following:
       \$ \& \% \# \{ and \}.

In printing, text is usually emphasized with an
       \emph{italic}
type style.

\begin{em}
   A long segment of text can also be emphasized
   in this way.  Text within such a segment can be
   given \emph{additional} emphasis.
\end{em}

It is sometimes necessary to prevent \LaTeX\ from
breaking a line where it might otherwise do so.
This may be at a space, as between the ``Mr.''\ and
``Jones'' in
       ``Mr.~Jones'',        % ~ produces an unbreakable interword space.
or within a word---especially when the word is a
symbol like
       \mbox{\emph{itemnum}}
that makes little sense when hyphenated across
lines.

Footnotes\footnote{This is an example of a footnote.}
pose no problem.

\LaTeX\ is good at typesetting mathematical formulas
like
       \( x-3y + z = 7 \)
or
       \( a_{1} > x^{2n} + y^{2n} > x' \)
or
       \( \ip{A}{B} = \sum_{i} a_{i} b_{i} \).
The spaces you type in a formula are
ignored.  Remember that a letter like
       $x$                   % $ ... $  and  \( ... \)  are equivalent
is a formula when it denotes a mathematical
symbol, and it should be typed as one.

\section{Displayed Text}

Text is displayed by indenting it from the left
margin.  Quotations are commonly displayed.  There
are short quotations
\begin{quote}
   This is a short quotation.  It consists of a
   single paragraph of text.  See how it is formatted.
\end{quote}
and longer ones.
\begin{quotation}
   This is a longer quotation.  It consists of two
   paragraphs of text, neither of which are
   particularly interesting.

   This is the second paragraph of the quotation.  It
   is just as dull as the first paragraph.
\end{quotation}
Another frequently-displayed structure is a list.
The following is an example of an \emph{itemized}
list.
\begin{itemize}
   \item This is the first item of an itemized list.
         Each item in the list is marked with a ``tick''.
         You don't have to worry about what kind of tick
         mark is used.

   \item This is the second item of the list.  It
         contains another list nested inside it.  The inner
         list is an \emph{enumerated} list.
         \begin{enumerate}
            \item This is the first item of an enumerated
                  list that is nested within the itemized list.

            \item This is the second item of the inner list.
                  \LaTeX\ allows you to nest lists deeper than
                  you really should.
         \end{enumerate}
         This is the rest of the second item of the outer
         list.  It is no more interesting than any other
         part of the item.
   \item This is the third item of the list.
\end{itemize}
You can even display poetry.
\begin{verse}
   There is an environment
    for verse \\             % The \\ command separates lines
   Whose features some poets % within a stanza.
   will curse.

                             % One or more blank lines separate stanzas.

   For instead of making\\
   Them do \emph{all} line breaking, \\
   It allows them to put too many words on a line when they'd rather be
   forced to be terse.
\end{verse}

Mathematical formulas may also be displayed.  A
displayed formula
is
one-line long; multiline
formulas require special formatting instructions.
   \[  \ip{\Gamma}{\psi'} = x'' + y^{2} + z_{i}^{n}\]
Don't start a paragraph with a displayed equation,
nor make one a paragraph by itself.

\end{document}               % End of document.

\section{Bench Huge Section 1001}
% synthetic bench section 1001
% This is a sample LaTeX input file.  (Version of 12 August 2004.)
%
% A '%' character causes TeX to ignore all remaining text on the line,
% and is used for comments like this one.

\documentclass{article}      % Specifies the document class

                             % The preamble begins here.
\title{An Example Document}  % Declares the document's title.
\author{Leslie Lamport}      % Declares the author's name.
\date{January 21, 1994}      % Deleting this command produces today's date.

\newcommand{\ip}[2]{(#1, #2)}
                             % Defines \ip{arg1}{arg2} to mean
                             % (arg1, arg2).

%\newcommand{\ip}[2]{\langle #1 | #2\rangle}
                             % This is an alternative definition of
                             % \ip that is commented out.

\begin{document}             % End of preamble and beginning of text.

\maketitle                   % Produces the title.

This is an example input file.  Comparing it with
the output it generates can show you how to
produce a simple document of your own.

\section{Ordinary Text}      % Produces section heading.  Lower-level
                             % sections are begun with similar
                             % \subsection and \subsubsection commands.

The ends  of words and sentences are marked
  by   spaces. It  doesn't matter how many
spaces    you type; one is as good as 100.  The
end of   a line counts as a space.

One   or more   blank lines denote the  end
of  a paragraph.

Since any number of consecutive spaces are treated
like a single one, the formatting of the input
file makes no difference to
      \LaTeX,                % The \LaTeX command generates the LaTeX logo.
but it makes a difference to you.  When you use
\LaTeX, making your input file as easy to read
as possible will be a great help as you write
your document and when you change it.  This sample
file shows how you can add comments to your own input
file.

Because printing is different from typewriting,
there are a number of things that you have to do
differently when preparing an input file than if
you were just typing the document directly.
Quotation marks like
       ``this''
have to be handled specially, as do quotes within
quotes:
       ``\,`this'            % \, separates the double and single quote.
        is what I just
        wrote, not  `that'\,''.

Dashes come in three sizes: an
       intra-word
dash, a medium dash for number ranges like
       1--2,
and a punctuation
       dash---like
this.

A sentence-ending space should be larger than the
space between words within a sentence.  You
sometimes have to type special commands in
conjunction with punctuation characters to get
this right, as in the following sentence.
       Gnats, gnus, etc.\ all  % `\ ' makes an inter-word space.
       begin with G\@.         % \@ marks end-of-sentence punctuation.
You should check the spaces after periods when
reading your output to make sure you haven't
forgotten any special cases.  Generating an
ellipsis
       \ldots\               % `\ ' is needed after `\ldots' because TeX
                             % ignores spaces after command names like \ldots
                             % made from \ + letters.
                             %
                             % Note how a `%' character causes TeX to ignore
                             % the end of the input line, so these blank lines
                             % do not start a new paragraph.
                             %
with the right spacing around the periods requires
a special command.

\LaTeX\ interprets some common characters as
commands, so you must type special commands to
generate them.  These characters include the
following:
       \$ \& \% \# \{ and \}.

In printing, text is usually emphasized with an
       \emph{italic}
type style.

\begin{em}
   A long segment of text can also be emphasized
   in this way.  Text within such a segment can be
   given \emph{additional} emphasis.
\end{em}

It is sometimes necessary to prevent \LaTeX\ from
breaking a line where it might otherwise do so.
This may be at a space, as between the ``Mr.''\ and
``Jones'' in
       ``Mr.~Jones'',        % ~ produces an unbreakable interword space.
or within a word---especially when the word is a
symbol like
       \mbox{\emph{itemnum}}
that makes little sense when hyphenated across
lines.

Footnotes\footnote{This is an example of a footnote.}
pose no problem.

\LaTeX\ is good at typesetting mathematical formulas
like
       \( x-3y + z = 7 \)
or
       \( a_{1} > x^{2n} + y^{2n} > x' \)
or
       \( \ip{A}{B} = \sum_{i} a_{i} b_{i} \).
The spaces you type in a formula are
ignored.  Remember that a letter like
       $x$                   % $ ... $  and  \( ... \)  are equivalent
is a formula when it denotes a mathematical
symbol, and it should be typed as one.

\section{Displayed Text}

Text is displayed by indenting it from the left
margin.  Quotations are commonly displayed.  There
are short quotations
\begin{quote}
   This is a short quotation.  It consists of a
   single paragraph of text.  See how it is formatted.
\end{quote}
and longer ones.
\begin{quotation}
   This is a longer quotation.  It consists of two
   paragraphs of text, neither of which are
   particularly interesting.

   This is the second paragraph of the quotation.  It
   is just as dull as the first paragraph.
\end{quotation}
Another frequently-displayed structure is a list.
The following is an example of an \emph{itemized}
list.
\begin{itemize}
   \item This is the first item of an itemized list.
         Each item in the list is marked with a ``tick''.
         You don't have to worry about what kind of tick
         mark is used.

   \item This is the second item of the list.  It
         contains another list nested inside it.  The inner
         list is an \emph{enumerated} list.
         \begin{enumerate}
            \item This is the first item of an enumerated
                  list that is nested within the itemized list.

            \item This is the second item of the inner list.
                  \LaTeX\ allows you to nest lists deeper than
                  you really should.
         \end{enumerate}
         This is the rest of the second item of the outer
         list.  It is no more interesting than any other
         part of the item.
   \item This is the third item of the list.
\end{itemize}
You can even display poetry.
\begin{verse}
   There is an environment
    for verse \\             % The \\ command separates lines
   Whose features some poets % within a stanza.
   will curse.

                             % One or more blank lines separate stanzas.

   For instead of making\\
   Them do \emph{all} line breaking, \\
   It allows them to put too many words on a line when they'd rather be
   forced to be terse.
\end{verse}

Mathematical formulas may also be displayed.  A
displayed formula
is
one-line long; multiline
formulas require special formatting instructions.
   \[  \ip{\Gamma}{\psi'} = x'' + y^{2} + z_{i}^{n}\]
Don't start a paragraph with a displayed equation,
nor make one a paragraph by itself.

\end{document}               % End of document.

\section{Bench Huge Section 1002}
% synthetic bench section 1002
% This is a sample LaTeX input file.  (Version of 12 August 2004.)
%
% A '%' character causes TeX to ignore all remaining text on the line,
% and is used for comments like this one.

\documentclass{article}      % Specifies the document class

                             % The preamble begins here.
\title{An Example Document}  % Declares the document's title.
\author{Leslie Lamport}      % Declares the author's name.
\date{January 21, 1994}      % Deleting this command produces today's date.

\newcommand{\ip}[2]{(#1, #2)}
                             % Defines \ip{arg1}{arg2} to mean
                             % (arg1, arg2).

%\newcommand{\ip}[2]{\langle #1 | #2\rangle}
                             % This is an alternative definition of
                             % \ip that is commented out.

\begin{document}             % End of preamble and beginning of text.

\maketitle                   % Produces the title.

This is an example input file.  Comparing it with
the output it generates can show you how to
produce a simple document of your own.

\section{Ordinary Text}      % Produces section heading.  Lower-level
                             % sections are begun with similar
                             % \subsection and \subsubsection commands.

The ends  of words and sentences are marked
  by   spaces. It  doesn't matter how many
spaces    you type; one is as good as 100.  The
end of   a line counts as a space.

One   or more   blank lines denote the  end
of  a paragraph.

Since any number of consecutive spaces are treated
like a single one, the formatting of the input
file makes no difference to
      \LaTeX,                % The \LaTeX command generates the LaTeX logo.
but it makes a difference to you.  When you use
\LaTeX, making your input file as easy to read
as possible will be a great help as you write
your document and when you change it.  This sample
file shows how you can add comments to your own input
file.

Because printing is different from typewriting,
there are a number of things that you have to do
differently when preparing an input file than if
you were just typing the document directly.
Quotation marks like
       ``this''
have to be handled specially, as do quotes within
quotes:
       ``\,`this'            % \, separates the double and single quote.
        is what I just
        wrote, not  `that'\,''.

Dashes come in three sizes: an
       intra-word
dash, a medium dash for number ranges like
       1--2,
and a punctuation
       dash---like
this.

A sentence-ending space should be larger than the
space between words within a sentence.  You
sometimes have to type special commands in
conjunction with punctuation characters to get
this right, as in the following sentence.
       Gnats, gnus, etc.\ all  % `\ ' makes an inter-word space.
       begin with G\@.         % \@ marks end-of-sentence punctuation.
You should check the spaces after periods when
reading your output to make sure you haven't
forgotten any special cases.  Generating an
ellipsis
       \ldots\               % `\ ' is needed after `\ldots' because TeX
                             % ignores spaces after command names like \ldots
                             % made from \ + letters.
                             %
                             % Note how a `%' character causes TeX to ignore
                             % the end of the input line, so these blank lines
                             % do not start a new paragraph.
                             %
with the right spacing around the periods requires
a special command.

\LaTeX\ interprets some common characters as
commands, so you must type special commands to
generate them.  These characters include the
following:
       \$ \& \% \# \{ and \}.

In printing, text is usually emphasized with an
       \emph{italic}
type style.

\begin{em}
   A long segment of text can also be emphasized
   in this way.  Text within such a segment can be
   given \emph{additional} emphasis.
\end{em}

It is sometimes necessary to prevent \LaTeX\ from
breaking a line where it might otherwise do so.
This may be at a space, as between the ``Mr.''\ and
``Jones'' in
       ``Mr.~Jones'',        % ~ produces an unbreakable interword space.
or within a word---especially when the word is a
symbol like
       \mbox{\emph{itemnum}}
that makes little sense when hyphenated across
lines.

Footnotes\footnote{This is an example of a footnote.}
pose no problem.

\LaTeX\ is good at typesetting mathematical formulas
like
       \( x-3y + z = 7 \)
or
       \( a_{1} > x^{2n} + y^{2n} > x' \)
or
       \( \ip{A}{B} = \sum_{i} a_{i} b_{i} \).
The spaces you type in a formula are
ignored.  Remember that a letter like
       $x$                   % $ ... $  and  \( ... \)  are equivalent
is a formula when it denotes a mathematical
symbol, and it should be typed as one.

\section{Displayed Text}

Text is displayed by indenting it from the left
margin.  Quotations are commonly displayed.  There
are short quotations
\begin{quote}
   This is a short quotation.  It consists of a
   single paragraph of text.  See how it is formatted.
\end{quote}
and longer ones.
\begin{quotation}
   This is a longer quotation.  It consists of two
   paragraphs of text, neither of which are
   particularly interesting.

   This is the second paragraph of the quotation.  It
   is just as dull as the first paragraph.
\end{quotation}
Another frequently-displayed structure is a list.
The following is an example of an \emph{itemized}
list.
\begin{itemize}
   \item This is the first item of an itemized list.
         Each item in the list is marked with a ``tick''.
         You don't have to worry about what kind of tick
         mark is used.

   \item This is the second item of the list.  It
         contains another list nested inside it.  The inner
         list is an \emph{enumerated} list.
         \begin{enumerate}
            \item This is the first item of an enumerated
                  list that is nested within the itemized list.

            \item This is the second item of the inner list.
                  \LaTeX\ allows you to nest lists deeper than
                  you really should.
         \end{enumerate}
         This is the rest of the second item of the outer
         list.  It is no more interesting than any other
         part of the item.
   \item This is the third item of the list.
\end{itemize}
You can even display poetry.
\begin{verse}
   There is an environment
    for verse \\             % The \\ command separates lines
   Whose features some poets % within a stanza.
   will curse.

                             % One or more blank lines separate stanzas.

   For instead of making\\
   Them do \emph{all} line breaking, \\
   It allows them to put too many words on a line when they'd rather be
   forced to be terse.
\end{verse}

Mathematical formulas may also be displayed.  A
displayed formula
is
one-line long; multiline
formulas require special formatting instructions.
   \[  \ip{\Gamma}{\psi'} = x'' + y^{2} + z_{i}^{n}\]
Don't start a paragraph with a displayed equation,
nor make one a paragraph by itself.

\end{document}               % End of document.

\section{Bench Huge Section 1003}
% synthetic bench section 1003
% This is a sample LaTeX input file.  (Version of 12 August 2004.)
%
% A '%' character causes TeX to ignore all remaining text on the line,
% and is used for comments like this one.

\documentclass{article}      % Specifies the document class

                             % The preamble begins here.
\title{An Example Document}  % Declares the document's title.
\author{Leslie Lamport}      % Declares the author's name.
\date{January 21, 1994}      % Deleting this command produces today's date.

\newcommand{\ip}[2]{(#1, #2)}
                             % Defines \ip{arg1}{arg2} to mean
                             % (arg1, arg2).

%\newcommand{\ip}[2]{\langle #1 | #2\rangle}
                             % This is an alternative definition of
                             % \ip that is commented out.

\begin{document}             % End of preamble and beginning of text.

\maketitle                   % Produces the title.

This is an example input file.  Comparing it with
the output it generates can show you how to
produce a simple document of your own.

\section{Ordinary Text}      % Produces section heading.  Lower-level
                             % sections are begun with similar
                             % \subsection and \subsubsection commands.

The ends  of words and sentences are marked
  by   spaces. It  doesn't matter how many
spaces    you type; one is as good as 100.  The
end of   a line counts as a space.

One   or more   blank lines denote the  end
of  a paragraph.

Since any number of consecutive spaces are treated
like a single one, the formatting of the input
file makes no difference to
      \LaTeX,                % The \LaTeX command generates the LaTeX logo.
but it makes a difference to you.  When you use
\LaTeX, making your input file as easy to read
as possible will be a great help as you write
your document and when you change it.  This sample
file shows how you can add comments to your own input
file.

Because printing is different from typewriting,
there are a number of things that you have to do
differently when preparing an input file than if
you were just typing the document directly.
Quotation marks like
       ``this''
have to be handled specially, as do quotes within
quotes:
       ``\,`this'            % \, separates the double and single quote.
        is what I just
        wrote, not  `that'\,''.

Dashes come in three sizes: an
       intra-word
dash, a medium dash for number ranges like
       1--2,
and a punctuation
       dash---like
this.

A sentence-ending space should be larger than the
space between words within a sentence.  You
sometimes have to type special commands in
conjunction with punctuation characters to get
this right, as in the following sentence.
       Gnats, gnus, etc.\ all  % `\ ' makes an inter-word space.
       begin with G\@.         % \@ marks end-of-sentence punctuation.
You should check the spaces after periods when
reading your output to make sure you haven't
forgotten any special cases.  Generating an
ellipsis
       \ldots\               % `\ ' is needed after `\ldots' because TeX
                             % ignores spaces after command names like \ldots
                             % made from \ + letters.
                             %
                             % Note how a `%' character causes TeX to ignore
                             % the end of the input line, so these blank lines
                             % do not start a new paragraph.
                             %
with the right spacing around the periods requires
a special command.

\LaTeX\ interprets some common characters as
commands, so you must type special commands to
generate them.  These characters include the
following:
       \$ \& \% \# \{ and \}.

In printing, text is usually emphasized with an
       \emph{italic}
type style.

\begin{em}
   A long segment of text can also be emphasized
   in this way.  Text within such a segment can be
   given \emph{additional} emphasis.
\end{em}

It is sometimes necessary to prevent \LaTeX\ from
breaking a line where it might otherwise do so.
This may be at a space, as between the ``Mr.''\ and
``Jones'' in
       ``Mr.~Jones'',        % ~ produces an unbreakable interword space.
or within a word---especially when the word is a
symbol like
       \mbox{\emph{itemnum}}
that makes little sense when hyphenated across
lines.

Footnotes\footnote{This is an example of a footnote.}
pose no problem.

\LaTeX\ is good at typesetting mathematical formulas
like
       \( x-3y + z = 7 \)
or
       \( a_{1} > x^{2n} + y^{2n} > x' \)
or
       \( \ip{A}{B} = \sum_{i} a_{i} b_{i} \).
The spaces you type in a formula are
ignored.  Remember that a letter like
       $x$                   % $ ... $  and  \( ... \)  are equivalent
is a formula when it denotes a mathematical
symbol, and it should be typed as one.

\section{Displayed Text}

Text is displayed by indenting it from the left
margin.  Quotations are commonly displayed.  There
are short quotations
\begin{quote}
   This is a short quotation.  It consists of a
   single paragraph of text.  See how it is formatted.
\end{quote}
and longer ones.
\begin{quotation}
   This is a longer quotation.  It consists of two
   paragraphs of text, neither of which are
   particularly interesting.

   This is the second paragraph of the quotation.  It
   is just as dull as the first paragraph.
\end{quotation}
Another frequently-displayed structure is a list.
The following is an example of an \emph{itemized}
list.
\begin{itemize}
   \item This is the first item of an itemized list.
         Each item in the list is marked with a ``tick''.
         You don't have to worry about what kind of tick
         mark is used.

   \item This is the second item of the list.  It
         contains another list nested inside it.  The inner
         list is an \emph{enumerated} list.
         \begin{enumerate}
            \item This is the first item of an enumerated
                  list that is nested within the itemized list.

            \item This is the second item of the inner list.
                  \LaTeX\ allows you to nest lists deeper than
                  you really should.
         \end{enumerate}
         This is the rest of the second item of the outer
         list.  It is no more interesting than any other
         part of the item.
   \item This is the third item of the list.
\end{itemize}
You can even display poetry.
\begin{verse}
   There is an environment
    for verse \\             % The \\ command separates lines
   Whose features some poets % within a stanza.
   will curse.

                             % One or more blank lines separate stanzas.

   For instead of making\\
   Them do \emph{all} line breaking, \\
   It allows them to put too many words on a line when they'd rather be
   forced to be terse.
\end{verse}

Mathematical formulas may also be displayed.  A
displayed formula
is
one-line long; multiline
formulas require special formatting instructions.
   \[  \ip{\Gamma}{\psi'} = x'' + y^{2} + z_{i}^{n}\]
Don't start a paragraph with a displayed equation,
nor make one a paragraph by itself.

\end{document}               % End of document.

\section{Bench Huge Section 1004}
% synthetic bench section 1004
% This is a sample LaTeX input file.  (Version of 12 August 2004.)
%
% A '%' character causes TeX to ignore all remaining text on the line,
% and is used for comments like this one.

\documentclass{article}      % Specifies the document class

                             % The preamble begins here.
\title{An Example Document}  % Declares the document's title.
\author{Leslie Lamport}      % Declares the author's name.
\date{January 21, 1994}      % Deleting this command produces today's date.

\newcommand{\ip}[2]{(#1, #2)}
                             % Defines \ip{arg1}{arg2} to mean
                             % (arg1, arg2).

%\newcommand{\ip}[2]{\langle #1 | #2\rangle}
                             % This is an alternative definition of
                             % \ip that is commented out.

\begin{document}             % End of preamble and beginning of text.

\maketitle                   % Produces the title.

This is an example input file.  Comparing it with
the output it generates can show you how to
produce a simple document of your own.

\section{Ordinary Text}      % Produces section heading.  Lower-level
                             % sections are begun with similar
                             % \subsection and \subsubsection commands.

The ends  of words and sentences are marked
  by   spaces. It  doesn't matter how many
spaces    you type; one is as good as 100.  The
end of   a line counts as a space.

One   or more   blank lines denote the  end
of  a paragraph.

Since any number of consecutive spaces are treated
like a single one, the formatting of the input
file makes no difference to
      \LaTeX,                % The \LaTeX command generates the LaTeX logo.
but it makes a difference to you.  When you use
\LaTeX, making your input file as easy to read
as possible will be a great help as you write
your document and when you change it.  This sample
file shows how you can add comments to your own input
file.

Because printing is different from typewriting,
there are a number of things that you have to do
differently when preparing an input file than if
you were just typing the document directly.
Quotation marks like
       ``this''
have to be handled specially, as do quotes within
quotes:
       ``\,`this'            % \, separates the double and single quote.
        is what I just
        wrote, not  `that'\,''.

Dashes come in three sizes: an
       intra-word
dash, a medium dash for number ranges like
       1--2,
and a punctuation
       dash---like
this.

A sentence-ending space should be larger than the
space between words within a sentence.  You
sometimes have to type special commands in
conjunction with punctuation characters to get
this right, as in the following sentence.
       Gnats, gnus, etc.\ all  % `\ ' makes an inter-word space.
       begin with G\@.         % \@ marks end-of-sentence punctuation.
You should check the spaces after periods when
reading your output to make sure you haven't
forgotten any special cases.  Generating an
ellipsis
       \ldots\               % `\ ' is needed after `\ldots' because TeX
                             % ignores spaces after command names like \ldots
                             % made from \ + letters.
                             %
                             % Note how a `%' character causes TeX to ignore
                             % the end of the input line, so these blank lines
                             % do not start a new paragraph.
                             %
with the right spacing around the periods requires
a special command.

\LaTeX\ interprets some common characters as
commands, so you must type special commands to
generate them.  These characters include the
following:
       \$ \& \% \# \{ and \}.

In printing, text is usually emphasized with an
       \emph{italic}
type style.

\begin{em}
   A long segment of text can also be emphasized
   in this way.  Text within such a segment can be
   given \emph{additional} emphasis.
\end{em}

It is sometimes necessary to prevent \LaTeX\ from
breaking a line where it might otherwise do so.
This may be at a space, as between the ``Mr.''\ and
``Jones'' in
       ``Mr.~Jones'',        % ~ produces an unbreakable interword space.
or within a word---especially when the word is a
symbol like
       \mbox{\emph{itemnum}}
that makes little sense when hyphenated across
lines.

Footnotes\footnote{This is an example of a footnote.}
pose no problem.

\LaTeX\ is good at typesetting mathematical formulas
like
       \( x-3y + z = 7 \)
or
       \( a_{1} > x^{2n} + y^{2n} > x' \)
or
       \( \ip{A}{B} = \sum_{i} a_{i} b_{i} \).
The spaces you type in a formula are
ignored.  Remember that a letter like
       $x$                   % $ ... $  and  \( ... \)  are equivalent
is a formula when it denotes a mathematical
symbol, and it should be typed as one.

\section{Displayed Text}

Text is displayed by indenting it from the left
margin.  Quotations are commonly displayed.  There
are short quotations
\begin{quote}
   This is a short quotation.  It consists of a
   single paragraph of text.  See how it is formatted.
\end{quote}
and longer ones.
\begin{quotation}
   This is a longer quotation.  It consists of two
   paragraphs of text, neither of which are
   particularly interesting.

   This is the second paragraph of the quotation.  It
   is just as dull as the first paragraph.
\end{quotation}
Another frequently-displayed structure is a list.
The following is an example of an \emph{itemized}
list.
\begin{itemize}
   \item This is the first item of an itemized list.
         Each item in the list is marked with a ``tick''.
         You don't have to worry about what kind of tick
         mark is used.

   \item This is the second item of the list.  It
         contains another list nested inside it.  The inner
         list is an \emph{enumerated} list.
         \begin{enumerate}
            \item This is the first item of an enumerated
                  list that is nested within the itemized list.

            \item This is the second item of the inner list.
                  \LaTeX\ allows you to nest lists deeper than
                  you really should.
         \end{enumerate}
         This is the rest of the second item of the outer
         list.  It is no more interesting than any other
         part of the item.
   \item This is the third item of the list.
\end{itemize}
You can even display poetry.
\begin{verse}
   There is an environment
    for verse \\             % The \\ command separates lines
   Whose features some poets % within a stanza.
   will curse.

                             % One or more blank lines separate stanzas.

   For instead of making\\
   Them do \emph{all} line breaking, \\
   It allows them to put too many words on a line when they'd rather be
   forced to be terse.
\end{verse}

Mathematical formulas may also be displayed.  A
displayed formula
is
one-line long; multiline
formulas require special formatting instructions.
   \[  \ip{\Gamma}{\psi'} = x'' + y^{2} + z_{i}^{n}\]
Don't start a paragraph with a displayed equation,
nor make one a paragraph by itself.

\end{document}               % End of document.

\section{Bench Huge Section 1005}
% synthetic bench section 1005
% This is a sample LaTeX input file.  (Version of 12 August 2004.)
%
% A '%' character causes TeX to ignore all remaining text on the line,
% and is used for comments like this one.

\documentclass{article}      % Specifies the document class

                             % The preamble begins here.
\title{An Example Document}  % Declares the document's title.
\author{Leslie Lamport}      % Declares the author's name.
\date{January 21, 1994}      % Deleting this command produces today's date.

\newcommand{\ip}[2]{(#1, #2)}
                             % Defines \ip{arg1}{arg2} to mean
                             % (arg1, arg2).

%\newcommand{\ip}[2]{\langle #1 | #2\rangle}
                             % This is an alternative definition of
                             % \ip that is commented out.

\begin{document}             % End of preamble and beginning of text.

\maketitle                   % Produces the title.

This is an example input file.  Comparing it with
the output it generates can show you how to
produce a simple document of your own.

\section{Ordinary Text}      % Produces section heading.  Lower-level
                             % sections are begun with similar
                             % \subsection and \subsubsection commands.

The ends  of words and sentences are marked
  by   spaces. It  doesn't matter how many
spaces    you type; one is as good as 100.  The
end of   a line counts as a space.

One   or more   blank lines denote the  end
of  a paragraph.

Since any number of consecutive spaces are treated
like a single one, the formatting of the input
file makes no difference to
      \LaTeX,                % The \LaTeX command generates the LaTeX logo.
but it makes a difference to you.  When you use
\LaTeX, making your input file as easy to read
as possible will be a great help as you write
your document and when you change it.  This sample
file shows how you can add comments to your own input
file.

Because printing is different from typewriting,
there are a number of things that you have to do
differently when preparing an input file than if
you were just typing the document directly.
Quotation marks like
       ``this''
have to be handled specially, as do quotes within
quotes:
       ``\,`this'            % \, separates the double and single quote.
        is what I just
        wrote, not  `that'\,''.

Dashes come in three sizes: an
       intra-word
dash, a medium dash for number ranges like
       1--2,
and a punctuation
       dash---like
this.

A sentence-ending space should be larger than the
space between words within a sentence.  You
sometimes have to type special commands in
conjunction with punctuation characters to get
this right, as in the following sentence.
       Gnats, gnus, etc.\ all  % `\ ' makes an inter-word space.
       begin with G\@.         % \@ marks end-of-sentence punctuation.
You should check the spaces after periods when
reading your output to make sure you haven't
forgotten any special cases.  Generating an
ellipsis
       \ldots\               % `\ ' is needed after `\ldots' because TeX
                             % ignores spaces after command names like \ldots
                             % made from \ + letters.
                             %
                             % Note how a `%' character causes TeX to ignore
                             % the end of the input line, so these blank lines
                             % do not start a new paragraph.
                             %
with the right spacing around the periods requires
a special command.

\LaTeX\ interprets some common characters as
commands, so you must type special commands to
generate them.  These characters include the
following:
       \$ \& \% \# \{ and \}.

In printing, text is usually emphasized with an
       \emph{italic}
type style.

\begin{em}
   A long segment of text can also be emphasized
   in this way.  Text within such a segment can be
   given \emph{additional} emphasis.
\end{em}

It is sometimes necessary to prevent \LaTeX\ from
breaking a line where it might otherwise do so.
This may be at a space, as between the ``Mr.''\ and
``Jones'' in
       ``Mr.~Jones'',        % ~ produces an unbreakable interword space.
or within a word---especially when the word is a
symbol like
       \mbox{\emph{itemnum}}
that makes little sense when hyphenated across
lines.

Footnotes\footnote{This is an example of a footnote.}
pose no problem.

\LaTeX\ is good at typesetting mathematical formulas
like
       \( x-3y + z = 7 \)
or
       \( a_{1} > x^{2n} + y^{2n} > x' \)
or
       \( \ip{A}{B} = \sum_{i} a_{i} b_{i} \).
The spaces you type in a formula are
ignored.  Remember that a letter like
       $x$                   % $ ... $  and  \( ... \)  are equivalent
is a formula when it denotes a mathematical
symbol, and it should be typed as one.

\section{Displayed Text}

Text is displayed by indenting it from the left
margin.  Quotations are commonly displayed.  There
are short quotations
\begin{quote}
   This is a short quotation.  It consists of a
   single paragraph of text.  See how it is formatted.
\end{quote}
and longer ones.
\begin{quotation}
   This is a longer quotation.  It consists of two
   paragraphs of text, neither of which are
   particularly interesting.

   This is the second paragraph of the quotation.  It
   is just as dull as the first paragraph.
\end{quotation}
Another frequently-displayed structure is a list.
The following is an example of an \emph{itemized}
list.
\begin{itemize}
   \item This is the first item of an itemized list.
         Each item in the list is marked with a ``tick''.
         You don't have to worry about what kind of tick
         mark is used.

   \item This is the second item of the list.  It
         contains another list nested inside it.  The inner
         list is an \emph{enumerated} list.
         \begin{enumerate}
            \item This is the first item of an enumerated
                  list that is nested within the itemized list.

            \item This is the second item of the inner list.
                  \LaTeX\ allows you to nest lists deeper than
                  you really should.
         \end{enumerate}
         This is the rest of the second item of the outer
         list.  It is no more interesting than any other
         part of the item.
   \item This is the third item of the list.
\end{itemize}
You can even display poetry.
\begin{verse}
   There is an environment
    for verse \\             % The \\ command separates lines
   Whose features some poets % within a stanza.
   will curse.

                             % One or more blank lines separate stanzas.

   For instead of making\\
   Them do \emph{all} line breaking, \\
   It allows them to put too many words on a line when they'd rather be
   forced to be terse.
\end{verse}

Mathematical formulas may also be displayed.  A
displayed formula
is
one-line long; multiline
formulas require special formatting instructions.
   \[  \ip{\Gamma}{\psi'} = x'' + y^{2} + z_{i}^{n}\]
Don't start a paragraph with a displayed equation,
nor make one a paragraph by itself.

\end{document}               % End of document.

\section{Bench Huge Section 1006}
% synthetic bench section 1006
% This is a sample LaTeX input file.  (Version of 12 August 2004.)
%
% A '%' character causes TeX to ignore all remaining text on the line,
% and is used for comments like this one.

\documentclass{article}      % Specifies the document class

                             % The preamble begins here.
\title{An Example Document}  % Declares the document's title.
\author{Leslie Lamport}      % Declares the author's name.
\date{January 21, 1994}      % Deleting this command produces today's date.

\newcommand{\ip}[2]{(#1, #2)}
                             % Defines \ip{arg1}{arg2} to mean
                             % (arg1, arg2).

%\newcommand{\ip}[2]{\langle #1 | #2\rangle}
                             % This is an alternative definition of
                             % \ip that is commented out.

\begin{document}             % End of preamble and beginning of text.

\maketitle                   % Produces the title.

This is an example input file.  Comparing it with
the output it generates can show you how to
produce a simple document of your own.

\section{Ordinary Text}      % Produces section heading.  Lower-level
                             % sections are begun with similar
                             % \subsection and \subsubsection commands.

The ends  of words and sentences are marked
  by   spaces. It  doesn't matter how many
spaces    you type; one is as good as 100.  The
end of   a line counts as a space.

One   or more   blank lines denote the  end
of  a paragraph.

Since any number of consecutive spaces are treated
like a single one, the formatting of the input
file makes no difference to
      \LaTeX,                % The \LaTeX command generates the LaTeX logo.
but it makes a difference to you.  When you use
\LaTeX, making your input file as easy to read
as possible will be a great help as you write
your document and when you change it.  This sample
file shows how you can add comments to your own input
file.

Because printing is different from typewriting,
there are a number of things that you have to do
differently when preparing an input file than if
you were just typing the document directly.
Quotation marks like
       ``this''
have to be handled specially, as do quotes within
quotes:
       ``\,`this'            % \, separates the double and single quote.
        is what I just
        wrote, not  `that'\,''.

Dashes come in three sizes: an
       intra-word
dash, a medium dash for number ranges like
       1--2,
and a punctuation
       dash---like
this.

A sentence-ending space should be larger than the
space between words within a sentence.  You
sometimes have to type special commands in
conjunction with punctuation characters to get
this right, as in the following sentence.
       Gnats, gnus, etc.\ all  % `\ ' makes an inter-word space.
       begin with G\@.         % \@ marks end-of-sentence punctuation.
You should check the spaces after periods when
reading your output to make sure you haven't
forgotten any special cases.  Generating an
ellipsis
       \ldots\               % `\ ' is needed after `\ldots' because TeX
                             % ignores spaces after command names like \ldots
                             % made from \ + letters.
                             %
                             % Note how a `%' character causes TeX to ignore
                             % the end of the input line, so these blank lines
                             % do not start a new paragraph.
                             %
with the right spacing around the periods requires
a special command.

\LaTeX\ interprets some common characters as
commands, so you must type special commands to
generate them.  These characters include the
following:
       \$ \& \% \# \{ and \}.

In printing, text is usually emphasized with an
       \emph{italic}
type style.

\begin{em}
   A long segment of text can also be emphasized
   in this way.  Text within such a segment can be
   given \emph{additional} emphasis.
\end{em}

It is sometimes necessary to prevent \LaTeX\ from
breaking a line where it might otherwise do so.
This may be at a space, as between the ``Mr.''\ and
``Jones'' in
       ``Mr.~Jones'',        % ~ produces an unbreakable interword space.
or within a word---especially when the word is a
symbol like
       \mbox{\emph{itemnum}}
that makes little sense when hyphenated across
lines.

Footnotes\footnote{This is an example of a footnote.}
pose no problem.

\LaTeX\ is good at typesetting mathematical formulas
like
       \( x-3y + z = 7 \)
or
       \( a_{1} > x^{2n} + y^{2n} > x' \)
or
       \( \ip{A}{B} = \sum_{i} a_{i} b_{i} \).
The spaces you type in a formula are
ignored.  Remember that a letter like
       $x$                   % $ ... $  and  \( ... \)  are equivalent
is a formula when it denotes a mathematical
symbol, and it should be typed as one.

\section{Displayed Text}

Text is displayed by indenting it from the left
margin.  Quotations are commonly displayed.  There
are short quotations
\begin{quote}
   This is a short quotation.  It consists of a
   single paragraph of text.  See how it is formatted.
\end{quote}
and longer ones.
\begin{quotation}
   This is a longer quotation.  It consists of two
   paragraphs of text, neither of which are
   particularly interesting.

   This is the second paragraph of the quotation.  It
   is just as dull as the first paragraph.
\end{quotation}
Another frequently-displayed structure is a list.
The following is an example of an \emph{itemized}
list.
\begin{itemize}
   \item This is the first item of an itemized list.
         Each item in the list is marked with a ``tick''.
         You don't have to worry about what kind of tick
         mark is used.

   \item This is the second item of the list.  It
         contains another list nested inside it.  The inner
         list is an \emph{enumerated} list.
         \begin{enumerate}
            \item This is the first item of an enumerated
                  list that is nested within the itemized list.

            \item This is the second item of the inner list.
                  \LaTeX\ allows you to nest lists deeper than
                  you really should.
         \end{enumerate}
         This is the rest of the second item of the outer
         list.  It is no more interesting than any other
         part of the item.
   \item This is the third item of the list.
\end{itemize}
You can even display poetry.
\begin{verse}
   There is an environment
    for verse \\             % The \\ command separates lines
   Whose features some poets % within a stanza.
   will curse.

                             % One or more blank lines separate stanzas.

   For instead of making\\
   Them do \emph{all} line breaking, \\
   It allows them to put too many words on a line when they'd rather be
   forced to be terse.
\end{verse}

Mathematical formulas may also be displayed.  A
displayed formula
is
one-line long; multiline
formulas require special formatting instructions.
   \[  \ip{\Gamma}{\psi'} = x'' + y^{2} + z_{i}^{n}\]
Don't start a paragraph with a displayed equation,
nor make one a paragraph by itself.

\end{document}               % End of document.

\section{Bench Huge Section 1007}
% synthetic bench section 1007
% This is a sample LaTeX input file.  (Version of 12 August 2004.)
%
% A '%' character causes TeX to ignore all remaining text on the line,
% and is used for comments like this one.

\documentclass{article}      % Specifies the document class

                             % The preamble begins here.
\title{An Example Document}  % Declares the document's title.
\author{Leslie Lamport}      % Declares the author's name.
\date{January 21, 1994}      % Deleting this command produces today's date.

\newcommand{\ip}[2]{(#1, #2)}
                             % Defines \ip{arg1}{arg2} to mean
                             % (arg1, arg2).

%\newcommand{\ip}[2]{\langle #1 | #2\rangle}
                             % This is an alternative definition of
                             % \ip that is commented out.

\begin{document}             % End of preamble and beginning of text.

\maketitle                   % Produces the title.

This is an example input file.  Comparing it with
the output it generates can show you how to
produce a simple document of your own.

\section{Ordinary Text}      % Produces section heading.  Lower-level
                             % sections are begun with similar
                             % \subsection and \subsubsection commands.

The ends  of words and sentences are marked
  by   spaces. It  doesn't matter how many
spaces    you type; one is as good as 100.  The
end of   a line counts as a space.

One   or more   blank lines denote the  end
of  a paragraph.

Since any number of consecutive spaces are treated
like a single one, the formatting of the input
file makes no difference to
      \LaTeX,                % The \LaTeX command generates the LaTeX logo.
but it makes a difference to you.  When you use
\LaTeX, making your input file as easy to read
as possible will be a great help as you write
your document and when you change it.  This sample
file shows how you can add comments to your own input
file.

Because printing is different from typewriting,
there are a number of things that you have to do
differently when preparing an input file than if
you were just typing the document directly.
Quotation marks like
       ``this''
have to be handled specially, as do quotes within
quotes:
       ``\,`this'            % \, separates the double and single quote.
        is what I just
        wrote, not  `that'\,''.

Dashes come in three sizes: an
       intra-word
dash, a medium dash for number ranges like
       1--2,
and a punctuation
       dash---like
this.

A sentence-ending space should be larger than the
space between words within a sentence.  You
sometimes have to type special commands in
conjunction with punctuation characters to get
this right, as in the following sentence.
       Gnats, gnus, etc.\ all  % `\ ' makes an inter-word space.
       begin with G\@.         % \@ marks end-of-sentence punctuation.
You should check the spaces after periods when
reading your output to make sure you haven't
forgotten any special cases.  Generating an
ellipsis
       \ldots\               % `\ ' is needed after `\ldots' because TeX
                             % ignores spaces after command names like \ldots
                             % made from \ + letters.
                             %
                             % Note how a `%' character causes TeX to ignore
                             % the end of the input line, so these blank lines
                             % do not start a new paragraph.
                             %
with the right spacing around the periods requires
a special command.

\LaTeX\ interprets some common characters as
commands, so you must type special commands to
generate them.  These characters include the
following:
       \$ \& \% \# \{ and \}.

In printing, text is usually emphasized with an
       \emph{italic}
type style.

\begin{em}
   A long segment of text can also be emphasized
   in this way.  Text within such a segment can be
   given \emph{additional} emphasis.
\end{em}

It is sometimes necessary to prevent \LaTeX\ from
breaking a line where it might otherwise do so.
This may be at a space, as between the ``Mr.''\ and
``Jones'' in
       ``Mr.~Jones'',        % ~ produces an unbreakable interword space.
or within a word---especially when the word is a
symbol like
       \mbox{\emph{itemnum}}
that makes little sense when hyphenated across
lines.

Footnotes\footnote{This is an example of a footnote.}
pose no problem.

\LaTeX\ is good at typesetting mathematical formulas
like
       \( x-3y + z = 7 \)
or
       \( a_{1} > x^{2n} + y^{2n} > x' \)
or
       \( \ip{A}{B} = \sum_{i} a_{i} b_{i} \).
The spaces you type in a formula are
ignored.  Remember that a letter like
       $x$                   % $ ... $  and  \( ... \)  are equivalent
is a formula when it denotes a mathematical
symbol, and it should be typed as one.

\section{Displayed Text}

Text is displayed by indenting it from the left
margin.  Quotations are commonly displayed.  There
are short quotations
\begin{quote}
   This is a short quotation.  It consists of a
   single paragraph of text.  See how it is formatted.
\end{quote}
and longer ones.
\begin{quotation}
   This is a longer quotation.  It consists of two
   paragraphs of text, neither of which are
   particularly interesting.

   This is the second paragraph of the quotation.  It
   is just as dull as the first paragraph.
\end{quotation}
Another frequently-displayed structure is a list.
The following is an example of an \emph{itemized}
list.
\begin{itemize}
   \item This is the first item of an itemized list.
         Each item in the list is marked with a ``tick''.
         You don't have to worry about what kind of tick
         mark is used.

   \item This is the second item of the list.  It
         contains another list nested inside it.  The inner
         list is an \emph{enumerated} list.
         \begin{enumerate}
            \item This is the first item of an enumerated
                  list that is nested within the itemized list.

            \item This is the second item of the inner list.
                  \LaTeX\ allows you to nest lists deeper than
                  you really should.
         \end{enumerate}
         This is the rest of the second item of the outer
         list.  It is no more interesting than any other
         part of the item.
   \item This is the third item of the list.
\end{itemize}
You can even display poetry.
\begin{verse}
   There is an environment
    for verse \\             % The \\ command separates lines
   Whose features some poets % within a stanza.
   will curse.

                             % One or more blank lines separate stanzas.

   For instead of making\\
   Them do \emph{all} line breaking, \\
   It allows them to put too many words on a line when they'd rather be
   forced to be terse.
\end{verse}

Mathematical formulas may also be displayed.  A
displayed formula
is
one-line long; multiline
formulas require special formatting instructions.
   \[  \ip{\Gamma}{\psi'} = x'' + y^{2} + z_{i}^{n}\]
Don't start a paragraph with a displayed equation,
nor make one a paragraph by itself.

\end{document}               % End of document.

\section{Bench Huge Section 1008}
% synthetic bench section 1008
% This is a sample LaTeX input file.  (Version of 12 August 2004.)
%
% A '%' character causes TeX to ignore all remaining text on the line,
% and is used for comments like this one.

\documentclass{article}      % Specifies the document class

                             % The preamble begins here.
\title{An Example Document}  % Declares the document's title.
\author{Leslie Lamport}      % Declares the author's name.
\date{January 21, 1994}      % Deleting this command produces today's date.

\newcommand{\ip}[2]{(#1, #2)}
                             % Defines \ip{arg1}{arg2} to mean
                             % (arg1, arg2).

%\newcommand{\ip}[2]{\langle #1 | #2\rangle}
                             % This is an alternative definition of
                             % \ip that is commented out.

\begin{document}             % End of preamble and beginning of text.

\maketitle                   % Produces the title.

This is an example input file.  Comparing it with
the output it generates can show you how to
produce a simple document of your own.

\section{Ordinary Text}      % Produces section heading.  Lower-level
                             % sections are begun with similar
                             % \subsection and \subsubsection commands.

The ends  of words and sentences are marked
  by   spaces. It  doesn't matter how many
spaces    you type; one is as good as 100.  The
end of   a line counts as a space.

One   or more   blank lines denote the  end
of  a paragraph.

Since any number of consecutive spaces are treated
like a single one, the formatting of the input
file makes no difference to
      \LaTeX,                % The \LaTeX command generates the LaTeX logo.
but it makes a difference to you.  When you use
\LaTeX, making your input file as easy to read
as possible will be a great help as you write
your document and when you change it.  This sample
file shows how you can add comments to your own input
file.

Because printing is different from typewriting,
there are a number of things that you have to do
differently when preparing an input file than if
you were just typing the document directly.
Quotation marks like
       ``this''
have to be handled specially, as do quotes within
quotes:
       ``\,`this'            % \, separates the double and single quote.
        is what I just
        wrote, not  `that'\,''.

Dashes come in three sizes: an
       intra-word
dash, a medium dash for number ranges like
       1--2,
and a punctuation
       dash---like
this.

A sentence-ending space should be larger than the
space between words within a sentence.  You
sometimes have to type special commands in
conjunction with punctuation characters to get
this right, as in the following sentence.
       Gnats, gnus, etc.\ all  % `\ ' makes an inter-word space.
       begin with G\@.         % \@ marks end-of-sentence punctuation.
You should check the spaces after periods when
reading your output to make sure you haven't
forgotten any special cases.  Generating an
ellipsis
       \ldots\               % `\ ' is needed after `\ldots' because TeX
                             % ignores spaces after command names like \ldots
                             % made from \ + letters.
                             %
                             % Note how a `%' character causes TeX to ignore
                             % the end of the input line, so these blank lines
                             % do not start a new paragraph.
                             %
with the right spacing around the periods requires
a special command.

\LaTeX\ interprets some common characters as
commands, so you must type special commands to
generate them.  These characters include the
following:
       \$ \& \% \# \{ and \}.

In printing, text is usually emphasized with an
       \emph{italic}
type style.

\begin{em}
   A long segment of text can also be emphasized
   in this way.  Text within such a segment can be
   given \emph{additional} emphasis.
\end{em}

It is sometimes necessary to prevent \LaTeX\ from
breaking a line where it might otherwise do so.
This may be at a space, as between the ``Mr.''\ and
``Jones'' in
       ``Mr.~Jones'',        % ~ produces an unbreakable interword space.
or within a word---especially when the word is a
symbol like
       \mbox{\emph{itemnum}}
that makes little sense when hyphenated across
lines.

Footnotes\footnote{This is an example of a footnote.}
pose no problem.

\LaTeX\ is good at typesetting mathematical formulas
like
       \( x-3y + z = 7 \)
or
       \( a_{1} > x^{2n} + y^{2n} > x' \)
or
       \( \ip{A}{B} = \sum_{i} a_{i} b_{i} \).
The spaces you type in a formula are
ignored.  Remember that a letter like
       $x$                   % $ ... $  and  \( ... \)  are equivalent
is a formula when it denotes a mathematical
symbol, and it should be typed as one.

\section{Displayed Text}

Text is displayed by indenting it from the left
margin.  Quotations are commonly displayed.  There
are short quotations
\begin{quote}
   This is a short quotation.  It consists of a
   single paragraph of text.  See how it is formatted.
\end{quote}
and longer ones.
\begin{quotation}
   This is a longer quotation.  It consists of two
   paragraphs of text, neither of which are
   particularly interesting.

   This is the second paragraph of the quotation.  It
   is just as dull as the first paragraph.
\end{quotation}
Another frequently-displayed structure is a list.
The following is an example of an \emph{itemized}
list.
\begin{itemize}
   \item This is the first item of an itemized list.
         Each item in the list is marked with a ``tick''.
         You don't have to worry about what kind of tick
         mark is used.

   \item This is the second item of the list.  It
         contains another list nested inside it.  The inner
         list is an \emph{enumerated} list.
         \begin{enumerate}
            \item This is the first item of an enumerated
                  list that is nested within the itemized list.

            \item This is the second item of the inner list.
                  \LaTeX\ allows you to nest lists deeper than
                  you really should.
         \end{enumerate}
         This is the rest of the second item of the outer
         list.  It is no more interesting than any other
         part of the item.
   \item This is the third item of the list.
\end{itemize}
You can even display poetry.
\begin{verse}
   There is an environment
    for verse \\             % The \\ command separates lines
   Whose features some poets % within a stanza.
   will curse.

                             % One or more blank lines separate stanzas.

   For instead of making\\
   Them do \emph{all} line breaking, \\
   It allows them to put too many words on a line when they'd rather be
   forced to be terse.
\end{verse}

Mathematical formulas may also be displayed.  A
displayed formula
is
one-line long; multiline
formulas require special formatting instructions.
   \[  \ip{\Gamma}{\psi'} = x'' + y^{2} + z_{i}^{n}\]
Don't start a paragraph with a displayed equation,
nor make one a paragraph by itself.

\end{document}               % End of document.

\section{Bench Huge Section 1009}
% synthetic bench section 1009
% This is a sample LaTeX input file.  (Version of 12 August 2004.)
%
% A '%' character causes TeX to ignore all remaining text on the line,
% and is used for comments like this one.

\documentclass{article}      % Specifies the document class

                             % The preamble begins here.
\title{An Example Document}  % Declares the document's title.
\author{Leslie Lamport}      % Declares the author's name.
\date{January 21, 1994}      % Deleting this command produces today's date.

\newcommand{\ip}[2]{(#1, #2)}
                             % Defines \ip{arg1}{arg2} to mean
                             % (arg1, arg2).

%\newcommand{\ip}[2]{\langle #1 | #2\rangle}
                             % This is an alternative definition of
                             % \ip that is commented out.

\begin{document}             % End of preamble and beginning of text.

\maketitle                   % Produces the title.

This is an example input file.  Comparing it with
the output it generates can show you how to
produce a simple document of your own.

\section{Ordinary Text}      % Produces section heading.  Lower-level
                             % sections are begun with similar
                             % \subsection and \subsubsection commands.

The ends  of words and sentences are marked
  by   spaces. It  doesn't matter how many
spaces    you type; one is as good as 100.  The
end of   a line counts as a space.

One   or more   blank lines denote the  end
of  a paragraph.

Since any number of consecutive spaces are treated
like a single one, the formatting of the input
file makes no difference to
      \LaTeX,                % The \LaTeX command generates the LaTeX logo.
but it makes a difference to you.  When you use
\LaTeX, making your input file as easy to read
as possible will be a great help as you write
your document and when you change it.  This sample
file shows how you can add comments to your own input
file.

Because printing is different from typewriting,
there are a number of things that you have to do
differently when preparing an input file than if
you were just typing the document directly.
Quotation marks like
       ``this''
have to be handled specially, as do quotes within
quotes:
       ``\,`this'            % \, separates the double and single quote.
        is what I just
        wrote, not  `that'\,''.

Dashes come in three sizes: an
       intra-word
dash, a medium dash for number ranges like
       1--2,
and a punctuation
       dash---like
this.

A sentence-ending space should be larger than the
space between words within a sentence.  You
sometimes have to type special commands in
conjunction with punctuation characters to get
this right, as in the following sentence.
       Gnats, gnus, etc.\ all  % `\ ' makes an inter-word space.
       begin with G\@.         % \@ marks end-of-sentence punctuation.
You should check the spaces after periods when
reading your output to make sure you haven't
forgotten any special cases.  Generating an
ellipsis
       \ldots\               % `\ ' is needed after `\ldots' because TeX
                             % ignores spaces after command names like \ldots
                             % made from \ + letters.
                             %
                             % Note how a `%' character causes TeX to ignore
                             % the end of the input line, so these blank lines
                             % do not start a new paragraph.
                             %
with the right spacing around the periods requires
a special command.

\LaTeX\ interprets some common characters as
commands, so you must type special commands to
generate them.  These characters include the
following:
       \$ \& \% \# \{ and \}.

In printing, text is usually emphasized with an
       \emph{italic}
type style.

\begin{em}
   A long segment of text can also be emphasized
   in this way.  Text within such a segment can be
   given \emph{additional} emphasis.
\end{em}

It is sometimes necessary to prevent \LaTeX\ from
breaking a line where it might otherwise do so.
This may be at a space, as between the ``Mr.''\ and
``Jones'' in
       ``Mr.~Jones'',        % ~ produces an unbreakable interword space.
or within a word---especially when the word is a
symbol like
       \mbox{\emph{itemnum}}
that makes little sense when hyphenated across
lines.

Footnotes\footnote{This is an example of a footnote.}
pose no problem.

\LaTeX\ is good at typesetting mathematical formulas
like
       \( x-3y + z = 7 \)
or
       \( a_{1} > x^{2n} + y^{2n} > x' \)
or
       \( \ip{A}{B} = \sum_{i} a_{i} b_{i} \).
The spaces you type in a formula are
ignored.  Remember that a letter like
       $x$                   % $ ... $  and  \( ... \)  are equivalent
is a formula when it denotes a mathematical
symbol, and it should be typed as one.

\section{Displayed Text}

Text is displayed by indenting it from the left
margin.  Quotations are commonly displayed.  There
are short quotations
\begin{quote}
   This is a short quotation.  It consists of a
   single paragraph of text.  See how it is formatted.
\end{quote}
and longer ones.
\begin{quotation}
   This is a longer quotation.  It consists of two
   paragraphs of text, neither of which are
   particularly interesting.

   This is the second paragraph of the quotation.  It
   is just as dull as the first paragraph.
\end{quotation}
Another frequently-displayed structure is a list.
The following is an example of an \emph{itemized}
list.
\begin{itemize}
   \item This is the first item of an itemized list.
         Each item in the list is marked with a ``tick''.
         You don't have to worry about what kind of tick
         mark is used.

   \item This is the second item of the list.  It
         contains another list nested inside it.  The inner
         list is an \emph{enumerated} list.
         \begin{enumerate}
            \item This is the first item of an enumerated
                  list that is nested within the itemized list.

            \item This is the second item of the inner list.
                  \LaTeX\ allows you to nest lists deeper than
                  you really should.
         \end{enumerate}
         This is the rest of the second item of the outer
         list.  It is no more interesting than any other
         part of the item.
   \item This is the third item of the list.
\end{itemize}
You can even display poetry.
\begin{verse}
   There is an environment
    for verse \\             % The \\ command separates lines
   Whose features some poets % within a stanza.
   will curse.

                             % One or more blank lines separate stanzas.

   For instead of making\\
   Them do \emph{all} line breaking, \\
   It allows them to put too many words on a line when they'd rather be
   forced to be terse.
\end{verse}

Mathematical formulas may also be displayed.  A
displayed formula
is
one-line long; multiline
formulas require special formatting instructions.
   \[  \ip{\Gamma}{\psi'} = x'' + y^{2} + z_{i}^{n}\]
Don't start a paragraph with a displayed equation,
nor make one a paragraph by itself.

\end{document}               % End of document.

\section{Bench Huge Section 1010}
% synthetic bench section 1010
% This is a sample LaTeX input file.  (Version of 12 August 2004.)
%
% A '%' character causes TeX to ignore all remaining text on the line,
% and is used for comments like this one.

\documentclass{article}      % Specifies the document class

                             % The preamble begins here.
\title{An Example Document}  % Declares the document's title.
\author{Leslie Lamport}      % Declares the author's name.
\date{January 21, 1994}      % Deleting this command produces today's date.

\newcommand{\ip}[2]{(#1, #2)}
                             % Defines \ip{arg1}{arg2} to mean
                             % (arg1, arg2).

%\newcommand{\ip}[2]{\langle #1 | #2\rangle}
                             % This is an alternative definition of
                             % \ip that is commented out.

\begin{document}             % End of preamble and beginning of text.

\maketitle                   % Produces the title.

This is an example input file.  Comparing it with
the output it generates can show you how to
produce a simple document of your own.

\section{Ordinary Text}      % Produces section heading.  Lower-level
                             % sections are begun with similar
                             % \subsection and \subsubsection commands.

The ends  of words and sentences are marked
  by   spaces. It  doesn't matter how many
spaces    you type; one is as good as 100.  The
end of   a line counts as a space.

One   or more   blank lines denote the  end
of  a paragraph.

Since any number of consecutive spaces are treated
like a single one, the formatting of the input
file makes no difference to
      \LaTeX,                % The \LaTeX command generates the LaTeX logo.
but it makes a difference to you.  When you use
\LaTeX, making your input file as easy to read
as possible will be a great help as you write
your document and when you change it.  This sample
file shows how you can add comments to your own input
file.

Because printing is different from typewriting,
there are a number of things that you have to do
differently when preparing an input file than if
you were just typing the document directly.
Quotation marks like
       ``this''
have to be handled specially, as do quotes within
quotes:
       ``\,`this'            % \, separates the double and single quote.
        is what I just
        wrote, not  `that'\,''.

Dashes come in three sizes: an
       intra-word
dash, a medium dash for number ranges like
       1--2,
and a punctuation
       dash---like
this.

A sentence-ending space should be larger than the
space between words within a sentence.  You
sometimes have to type special commands in
conjunction with punctuation characters to get
this right, as in the following sentence.
       Gnats, gnus, etc.\ all  % `\ ' makes an inter-word space.
       begin with G\@.         % \@ marks end-of-sentence punctuation.
You should check the spaces after periods when
reading your output to make sure you haven't
forgotten any special cases.  Generating an
ellipsis
       \ldots\               % `\ ' is needed after `\ldots' because TeX
                             % ignores spaces after command names like \ldots
                             % made from \ + letters.
                             %
                             % Note how a `%' character causes TeX to ignore
                             % the end of the input line, so these blank lines
                             % do not start a new paragraph.
                             %
with the right spacing around the periods requires
a special command.

\LaTeX\ interprets some common characters as
commands, so you must type special commands to
generate them.  These characters include the
following:
       \$ \& \% \# \{ and \}.

In printing, text is usually emphasized with an
       \emph{italic}
type style.

\begin{em}
   A long segment of text can also be emphasized
   in this way.  Text within such a segment can be
   given \emph{additional} emphasis.
\end{em}

It is sometimes necessary to prevent \LaTeX\ from
breaking a line where it might otherwise do so.
This may be at a space, as between the ``Mr.''\ and
``Jones'' in
       ``Mr.~Jones'',        % ~ produces an unbreakable interword space.
or within a word---especially when the word is a
symbol like
       \mbox{\emph{itemnum}}
that makes little sense when hyphenated across
lines.

Footnotes\footnote{This is an example of a footnote.}
pose no problem.

\LaTeX\ is good at typesetting mathematical formulas
like
       \( x-3y + z = 7 \)
or
       \( a_{1} > x^{2n} + y^{2n} > x' \)
or
       \( \ip{A}{B} = \sum_{i} a_{i} b_{i} \).
The spaces you type in a formula are
ignored.  Remember that a letter like
       $x$                   % $ ... $  and  \( ... \)  are equivalent
is a formula when it denotes a mathematical
symbol, and it should be typed as one.

\section{Displayed Text}

Text is displayed by indenting it from the left
margin.  Quotations are commonly displayed.  There
are short quotations
\begin{quote}
   This is a short quotation.  It consists of a
   single paragraph of text.  See how it is formatted.
\end{quote}
and longer ones.
\begin{quotation}
   This is a longer quotation.  It consists of two
   paragraphs of text, neither of which are
   particularly interesting.

   This is the second paragraph of the quotation.  It
   is just as dull as the first paragraph.
\end{quotation}
Another frequently-displayed structure is a list.
The following is an example of an \emph{itemized}
list.
\begin{itemize}
   \item This is the first item of an itemized list.
         Each item in the list is marked with a ``tick''.
         You don't have to worry about what kind of tick
         mark is used.

   \item This is the second item of the list.  It
         contains another list nested inside it.  The inner
         list is an \emph{enumerated} list.
         \begin{enumerate}
            \item This is the first item of an enumerated
                  list that is nested within the itemized list.

            \item This is the second item of the inner list.
                  \LaTeX\ allows you to nest lists deeper than
                  you really should.
         \end{enumerate}
         This is the rest of the second item of the outer
         list.  It is no more interesting than any other
         part of the item.
   \item This is the third item of the list.
\end{itemize}
You can even display poetry.
\begin{verse}
   There is an environment
    for verse \\             % The \\ command separates lines
   Whose features some poets % within a stanza.
   will curse.

                             % One or more blank lines separate stanzas.

   For instead of making\\
   Them do \emph{all} line breaking, \\
   It allows them to put too many words on a line when they'd rather be
   forced to be terse.
\end{verse}

Mathematical formulas may also be displayed.  A
displayed formula
is
one-line long; multiline
formulas require special formatting instructions.
   \[  \ip{\Gamma}{\psi'} = x'' + y^{2} + z_{i}^{n}\]
Don't start a paragraph with a displayed equation,
nor make one a paragraph by itself.

\end{document}               % End of document.

\section{Bench Huge Section 1011}
% synthetic bench section 1011
% This is a sample LaTeX input file.  (Version of 12 August 2004.)
%
% A '%' character causes TeX to ignore all remaining text on the line,
% and is used for comments like this one.

\documentclass{article}      % Specifies the document class

                             % The preamble begins here.
\title{An Example Document}  % Declares the document's title.
\author{Leslie Lamport}      % Declares the author's name.
\date{January 21, 1994}      % Deleting this command produces today's date.

\newcommand{\ip}[2]{(#1, #2)}
                             % Defines \ip{arg1}{arg2} to mean
                             % (arg1, arg2).

%\newcommand{\ip}[2]{\langle #1 | #2\rangle}
                             % This is an alternative definition of
                             % \ip that is commented out.

\begin{document}             % End of preamble and beginning of text.

\maketitle                   % Produces the title.

This is an example input file.  Comparing it with
the output it generates can show you how to
produce a simple document of your own.

\section{Ordinary Text}      % Produces section heading.  Lower-level
                             % sections are begun with similar
                             % \subsection and \subsubsection commands.

The ends  of words and sentences are marked
  by   spaces. It  doesn't matter how many
spaces    you type; one is as good as 100.  The
end of   a line counts as a space.

One   or more   blank lines denote the  end
of  a paragraph.

Since any number of consecutive spaces are treated
like a single one, the formatting of the input
file makes no difference to
      \LaTeX,                % The \LaTeX command generates the LaTeX logo.
but it makes a difference to you.  When you use
\LaTeX, making your input file as easy to read
as possible will be a great help as you write
your document and when you change it.  This sample
file shows how you can add comments to your own input
file.

Because printing is different from typewriting,
there are a number of things that you have to do
differently when preparing an input file than if
you were just typing the document directly.
Quotation marks like
       ``this''
have to be handled specially, as do quotes within
quotes:
       ``\,`this'            % \, separates the double and single quote.
        is what I just
        wrote, not  `that'\,''.

Dashes come in three sizes: an
       intra-word
dash, a medium dash for number ranges like
       1--2,
and a punctuation
       dash---like
this.

A sentence-ending space should be larger than the
space between words within a sentence.  You
sometimes have to type special commands in
conjunction with punctuation characters to get
this right, as in the following sentence.
       Gnats, gnus, etc.\ all  % `\ ' makes an inter-word space.
       begin with G\@.         % \@ marks end-of-sentence punctuation.
You should check the spaces after periods when
reading your output to make sure you haven't
forgotten any special cases.  Generating an
ellipsis
       \ldots\               % `\ ' is needed after `\ldots' because TeX
                             % ignores spaces after command names like \ldots
                             % made from \ + letters.
                             %
                             % Note how a `%' character causes TeX to ignore
                             % the end of the input line, so these blank lines
                             % do not start a new paragraph.
                             %
with the right spacing around the periods requires
a special command.

\LaTeX\ interprets some common characters as
commands, so you must type special commands to
generate them.  These characters include the
following:
       \$ \& \% \# \{ and \}.

In printing, text is usually emphasized with an
       \emph{italic}
type style.

\begin{em}
   A long segment of text can also be emphasized
   in this way.  Text within such a segment can be
   given \emph{additional} emphasis.
\end{em}

It is sometimes necessary to prevent \LaTeX\ from
breaking a line where it might otherwise do so.
This may be at a space, as between the ``Mr.''\ and
``Jones'' in
       ``Mr.~Jones'',        % ~ produces an unbreakable interword space.
or within a word---especially when the word is a
symbol like
       \mbox{\emph{itemnum}}
that makes little sense when hyphenated across
lines.

Footnotes\footnote{This is an example of a footnote.}
pose no problem.

\LaTeX\ is good at typesetting mathematical formulas
like
       \( x-3y + z = 7 \)
or
       \( a_{1} > x^{2n} + y^{2n} > x' \)
or
       \( \ip{A}{B} = \sum_{i} a_{i} b_{i} \).
The spaces you type in a formula are
ignored.  Remember that a letter like
       $x$                   % $ ... $  and  \( ... \)  are equivalent
is a formula when it denotes a mathematical
symbol, and it should be typed as one.

\section{Displayed Text}

Text is displayed by indenting it from the left
margin.  Quotations are commonly displayed.  There
are short quotations
\begin{quote}
   This is a short quotation.  It consists of a
   single paragraph of text.  See how it is formatted.
\end{quote}
and longer ones.
\begin{quotation}
   This is a longer quotation.  It consists of two
   paragraphs of text, neither of which are
   particularly interesting.

   This is the second paragraph of the quotation.  It
   is just as dull as the first paragraph.
\end{quotation}
Another frequently-displayed structure is a list.
The following is an example of an \emph{itemized}
list.
\begin{itemize}
   \item This is the first item of an itemized list.
         Each item in the list is marked with a ``tick''.
         You don't have to worry about what kind of tick
         mark is used.

   \item This is the second item of the list.  It
         contains another list nested inside it.  The inner
         list is an \emph{enumerated} list.
         \begin{enumerate}
            \item This is the first item of an enumerated
                  list that is nested within the itemized list.

            \item This is the second item of the inner list.
                  \LaTeX\ allows you to nest lists deeper than
                  you really should.
         \end{enumerate}
         This is the rest of the second item of the outer
         list.  It is no more interesting than any other
         part of the item.
   \item This is the third item of the list.
\end{itemize}
You can even display poetry.
\begin{verse}
   There is an environment
    for verse \\             % The \\ command separates lines
   Whose features some poets % within a stanza.
   will curse.

                             % One or more blank lines separate stanzas.

   For instead of making\\
   Them do \emph{all} line breaking, \\
   It allows them to put too many words on a line when they'd rather be
   forced to be terse.
\end{verse}

Mathematical formulas may also be displayed.  A
displayed formula
is
one-line long; multiline
formulas require special formatting instructions.
   \[  \ip{\Gamma}{\psi'} = x'' + y^{2} + z_{i}^{n}\]
Don't start a paragraph with a displayed equation,
nor make one a paragraph by itself.

\end{document}               % End of document.

\section{Bench Huge Section 1012}
% synthetic bench section 1012
% This is a sample LaTeX input file.  (Version of 12 August 2004.)
%
% A '%' character causes TeX to ignore all remaining text on the line,
% and is used for comments like this one.

\documentclass{article}      % Specifies the document class

                             % The preamble begins here.
\title{An Example Document}  % Declares the document's title.
\author{Leslie Lamport}      % Declares the author's name.
\date{January 21, 1994}      % Deleting this command produces today's date.

\newcommand{\ip}[2]{(#1, #2)}
                             % Defines \ip{arg1}{arg2} to mean
                             % (arg1, arg2).

%\newcommand{\ip}[2]{\langle #1 | #2\rangle}
                             % This is an alternative definition of
                             % \ip that is commented out.

\begin{document}             % End of preamble and beginning of text.

\maketitle                   % Produces the title.

This is an example input file.  Comparing it with
the output it generates can show you how to
produce a simple document of your own.

\section{Ordinary Text}      % Produces section heading.  Lower-level
                             % sections are begun with similar
                             % \subsection and \subsubsection commands.

The ends  of words and sentences are marked
  by   spaces. It  doesn't matter how many
spaces    you type; one is as good as 100.  The
end of   a line counts as a space.

One   or more   blank lines denote the  end
of  a paragraph.

Since any number of consecutive spaces are treated
like a single one, the formatting of the input
file makes no difference to
      \LaTeX,                % The \LaTeX command generates the LaTeX logo.
but it makes a difference to you.  When you use
\LaTeX, making your input file as easy to read
as possible will be a great help as you write
your document and when you change it.  This sample
file shows how you can add comments to your own input
file.

Because printing is different from typewriting,
there are a number of things that you have to do
differently when preparing an input file than if
you were just typing the document directly.
Quotation marks like
       ``this''
have to be handled specially, as do quotes within
quotes:
       ``\,`this'            % \, separates the double and single quote.
        is what I just
        wrote, not  `that'\,''.

Dashes come in three sizes: an
       intra-word
dash, a medium dash for number ranges like
       1--2,
and a punctuation
       dash---like
this.

A sentence-ending space should be larger than the
space between words within a sentence.  You
sometimes have to type special commands in
conjunction with punctuation characters to get
this right, as in the following sentence.
       Gnats, gnus, etc.\ all  % `\ ' makes an inter-word space.
       begin with G\@.         % \@ marks end-of-sentence punctuation.
You should check the spaces after periods when
reading your output to make sure you haven't
forgotten any special cases.  Generating an
ellipsis
       \ldots\               % `\ ' is needed after `\ldots' because TeX
                             % ignores spaces after command names like \ldots
                             % made from \ + letters.
                             %
                             % Note how a `%' character causes TeX to ignore
                             % the end of the input line, so these blank lines
                             % do not start a new paragraph.
                             %
with the right spacing around the periods requires
a special command.

\LaTeX\ interprets some common characters as
commands, so you must type special commands to
generate them.  These characters include the
following:
       \$ \& \% \# \{ and \}.

In printing, text is usually emphasized with an
       \emph{italic}
type style.

\begin{em}
   A long segment of text can also be emphasized
   in this way.  Text within such a segment can be
   given \emph{additional} emphasis.
\end{em}

It is sometimes necessary to prevent \LaTeX\ from
breaking a line where it might otherwise do so.
This may be at a space, as between the ``Mr.''\ and
``Jones'' in
       ``Mr.~Jones'',        % ~ produces an unbreakable interword space.
or within a word---especially when the word is a
symbol like
       \mbox{\emph{itemnum}}
that makes little sense when hyphenated across
lines.

Footnotes\footnote{This is an example of a footnote.}
pose no problem.

\LaTeX\ is good at typesetting mathematical formulas
like
       \( x-3y + z = 7 \)
or
       \( a_{1} > x^{2n} + y^{2n} > x' \)
or
       \( \ip{A}{B} = \sum_{i} a_{i} b_{i} \).
The spaces you type in a formula are
ignored.  Remember that a letter like
       $x$                   % $ ... $  and  \( ... \)  are equivalent
is a formula when it denotes a mathematical
symbol, and it should be typed as one.

\section{Displayed Text}

Text is displayed by indenting it from the left
margin.  Quotations are commonly displayed.  There
are short quotations
\begin{quote}
   This is a short quotation.  It consists of a
   single paragraph of text.  See how it is formatted.
\end{quote}
and longer ones.
\begin{quotation}
   This is a longer quotation.  It consists of two
   paragraphs of text, neither of which are
   particularly interesting.

   This is the second paragraph of the quotation.  It
   is just as dull as the first paragraph.
\end{quotation}
Another frequently-displayed structure is a list.
The following is an example of an \emph{itemized}
list.
\begin{itemize}
   \item This is the first item of an itemized list.
         Each item in the list is marked with a ``tick''.
         You don't have to worry about what kind of tick
         mark is used.

   \item This is the second item of the list.  It
         contains another list nested inside it.  The inner
         list is an \emph{enumerated} list.
         \begin{enumerate}
            \item This is the first item of an enumerated
                  list that is nested within the itemized list.

            \item This is the second item of the inner list.
                  \LaTeX\ allows you to nest lists deeper than
                  you really should.
         \end{enumerate}
         This is the rest of the second item of the outer
         list.  It is no more interesting than any other
         part of the item.
   \item This is the third item of the list.
\end{itemize}
You can even display poetry.
\begin{verse}
   There is an environment
    for verse \\             % The \\ command separates lines
   Whose features some poets % within a stanza.
   will curse.

                             % One or more blank lines separate stanzas.

   For instead of making\\
   Them do \emph{all} line breaking, \\
   It allows them to put too many words on a line when they'd rather be
   forced to be terse.
\end{verse}

Mathematical formulas may also be displayed.  A
displayed formula
is
one-line long; multiline
formulas require special formatting instructions.
   \[  \ip{\Gamma}{\psi'} = x'' + y^{2} + z_{i}^{n}\]
Don't start a paragraph with a displayed equation,
nor make one a paragraph by itself.

\end{document}               % End of document.

\section{Bench Huge Section 1013}
% synthetic bench section 1013
% This is a sample LaTeX input file.  (Version of 12 August 2004.)
%
% A '%' character causes TeX to ignore all remaining text on the line,
% and is used for comments like this one.

\documentclass{article}      % Specifies the document class

                             % The preamble begins here.
\title{An Example Document}  % Declares the document's title.
\author{Leslie Lamport}      % Declares the author's name.
\date{January 21, 1994}      % Deleting this command produces today's date.

\newcommand{\ip}[2]{(#1, #2)}
                             % Defines \ip{arg1}{arg2} to mean
                             % (arg1, arg2).

%\newcommand{\ip}[2]{\langle #1 | #2\rangle}
                             % This is an alternative definition of
                             % \ip that is commented out.

\begin{document}             % End of preamble and beginning of text.

\maketitle                   % Produces the title.

This is an example input file.  Comparing it with
the output it generates can show you how to
produce a simple document of your own.

\section{Ordinary Text}      % Produces section heading.  Lower-level
                             % sections are begun with similar
                             % \subsection and \subsubsection commands.

The ends  of words and sentences are marked
  by   spaces. It  doesn't matter how many
spaces    you type; one is as good as 100.  The
end of   a line counts as a space.

One   or more   blank lines denote the  end
of  a paragraph.

Since any number of consecutive spaces are treated
like a single one, the formatting of the input
file makes no difference to
      \LaTeX,                % The \LaTeX command generates the LaTeX logo.
but it makes a difference to you.  When you use
\LaTeX, making your input file as easy to read
as possible will be a great help as you write
your document and when you change it.  This sample
file shows how you can add comments to your own input
file.

Because printing is different from typewriting,
there are a number of things that you have to do
differently when preparing an input file than if
you were just typing the document directly.
Quotation marks like
       ``this''
have to be handled specially, as do quotes within
quotes:
       ``\,`this'            % \, separates the double and single quote.
        is what I just
        wrote, not  `that'\,''.

Dashes come in three sizes: an
       intra-word
dash, a medium dash for number ranges like
       1--2,
and a punctuation
       dash---like
this.

A sentence-ending space should be larger than the
space between words within a sentence.  You
sometimes have to type special commands in
conjunction with punctuation characters to get
this right, as in the following sentence.
       Gnats, gnus, etc.\ all  % `\ ' makes an inter-word space.
       begin with G\@.         % \@ marks end-of-sentence punctuation.
You should check the spaces after periods when
reading your output to make sure you haven't
forgotten any special cases.  Generating an
ellipsis
       \ldots\               % `\ ' is needed after `\ldots' because TeX
                             % ignores spaces after command names like \ldots
                             % made from \ + letters.
                             %
                             % Note how a `%' character causes TeX to ignore
                             % the end of the input line, so these blank lines
                             % do not start a new paragraph.
                             %
with the right spacing around the periods requires
a special command.

\LaTeX\ interprets some common characters as
commands, so you must type special commands to
generate them.  These characters include the
following:
       \$ \& \% \# \{ and \}.

In printing, text is usually emphasized with an
       \emph{italic}
type style.

\begin{em}
   A long segment of text can also be emphasized
   in this way.  Text within such a segment can be
   given \emph{additional} emphasis.
\end{em}

It is sometimes necessary to prevent \LaTeX\ from
breaking a line where it might otherwise do so.
This may be at a space, as between the ``Mr.''\ and
``Jones'' in
       ``Mr.~Jones'',        % ~ produces an unbreakable interword space.
or within a word---especially when the word is a
symbol like
       \mbox{\emph{itemnum}}
that makes little sense when hyphenated across
lines.

Footnotes\footnote{This is an example of a footnote.}
pose no problem.

\LaTeX\ is good at typesetting mathematical formulas
like
       \( x-3y + z = 7 \)
or
       \( a_{1} > x^{2n} + y^{2n} > x' \)
or
       \( \ip{A}{B} = \sum_{i} a_{i} b_{i} \).
The spaces you type in a formula are
ignored.  Remember that a letter like
       $x$                   % $ ... $  and  \( ... \)  are equivalent
is a formula when it denotes a mathematical
symbol, and it should be typed as one.

\section{Displayed Text}

Text is displayed by indenting it from the left
margin.  Quotations are commonly displayed.  There
are short quotations
\begin{quote}
   This is a short quotation.  It consists of a
   single paragraph of text.  See how it is formatted.
\end{quote}
and longer ones.
\begin{quotation}
   This is a longer quotation.  It consists of two
   paragraphs of text, neither of which are
   particularly interesting.

   This is the second paragraph of the quotation.  It
   is just as dull as the first paragraph.
\end{quotation}
Another frequently-displayed structure is a list.
The following is an example of an \emph{itemized}
list.
\begin{itemize}
   \item This is the first item of an itemized list.
         Each item in the list is marked with a ``tick''.
         You don't have to worry about what kind of tick
         mark is used.

   \item This is the second item of the list.  It
         contains another list nested inside it.  The inner
         list is an \emph{enumerated} list.
         \begin{enumerate}
            \item This is the first item of an enumerated
                  list that is nested within the itemized list.

            \item This is the second item of the inner list.
                  \LaTeX\ allows you to nest lists deeper than
                  you really should.
         \end{enumerate}
         This is the rest of the second item of the outer
         list.  It is no more interesting than any other
         part of the item.
   \item This is the third item of the list.
\end{itemize}
You can even display poetry.
\begin{verse}
   There is an environment
    for verse \\             % The \\ command separates lines
   Whose features some poets % within a stanza.
   will curse.

                             % One or more blank lines separate stanzas.

   For instead of making\\
   Them do \emph{all} line breaking, \\
   It allows them to put too many words on a line when they'd rather be
   forced to be terse.
\end{verse}

Mathematical formulas may also be displayed.  A
displayed formula
is
one-line long; multiline
formulas require special formatting instructions.
   \[  \ip{\Gamma}{\psi'} = x'' + y^{2} + z_{i}^{n}\]
Don't start a paragraph with a displayed equation,
nor make one a paragraph by itself.

\end{document}               % End of document.

\section{Bench Huge Section 1014}
% synthetic bench section 1014
% This is a sample LaTeX input file.  (Version of 12 August 2004.)
%
% A '%' character causes TeX to ignore all remaining text on the line,
% and is used for comments like this one.

\documentclass{article}      % Specifies the document class

                             % The preamble begins here.
\title{An Example Document}  % Declares the document's title.
\author{Leslie Lamport}      % Declares the author's name.
\date{January 21, 1994}      % Deleting this command produces today's date.

\newcommand{\ip}[2]{(#1, #2)}
                             % Defines \ip{arg1}{arg2} to mean
                             % (arg1, arg2).

%\newcommand{\ip}[2]{\langle #1 | #2\rangle}
                             % This is an alternative definition of
                             % \ip that is commented out.

\begin{document}             % End of preamble and beginning of text.

\maketitle                   % Produces the title.

This is an example input file.  Comparing it with
the output it generates can show you how to
produce a simple document of your own.

\section{Ordinary Text}      % Produces section heading.  Lower-level
                             % sections are begun with similar
                             % \subsection and \subsubsection commands.

The ends  of words and sentences are marked
  by   spaces. It  doesn't matter how many
spaces    you type; one is as good as 100.  The
end of   a line counts as a space.

One   or more   blank lines denote the  end
of  a paragraph.

Since any number of consecutive spaces are treated
like a single one, the formatting of the input
file makes no difference to
      \LaTeX,                % The \LaTeX command generates the LaTeX logo.
but it makes a difference to you.  When you use
\LaTeX, making your input file as easy to read
as possible will be a great help as you write
your document and when you change it.  This sample
file shows how you can add comments to your own input
file.

Because printing is different from typewriting,
there are a number of things that you have to do
differently when preparing an input file than if
you were just typing the document directly.
Quotation marks like
       ``this''
have to be handled specially, as do quotes within
quotes:
       ``\,`this'            % \, separates the double and single quote.
        is what I just
        wrote, not  `that'\,''.

Dashes come in three sizes: an
       intra-word
dash, a medium dash for number ranges like
       1--2,
and a punctuation
       dash---like
this.

A sentence-ending space should be larger than the
space between words within a sentence.  You
sometimes have to type special commands in
conjunction with punctuation characters to get
this right, as in the following sentence.
       Gnats, gnus, etc.\ all  % `\ ' makes an inter-word space.
       begin with G\@.         % \@ marks end-of-sentence punctuation.
You should check the spaces after periods when
reading your output to make sure you haven't
forgotten any special cases.  Generating an
ellipsis
       \ldots\               % `\ ' is needed after `\ldots' because TeX
                             % ignores spaces after command names like \ldots
                             % made from \ + letters.
                             %
                             % Note how a `%' character causes TeX to ignore
                             % the end of the input line, so these blank lines
                             % do not start a new paragraph.
                             %
with the right spacing around the periods requires
a special command.

\LaTeX\ interprets some common characters as
commands, so you must type special commands to
generate them.  These characters include the
following:
       \$ \& \% \# \{ and \}.

In printing, text is usually emphasized with an
       \emph{italic}
type style.

\begin{em}
   A long segment of text can also be emphasized
   in this way.  Text within such a segment can be
   given \emph{additional} emphasis.
\end{em}

It is sometimes necessary to prevent \LaTeX\ from
breaking a line where it might otherwise do so.
This may be at a space, as between the ``Mr.''\ and
``Jones'' in
       ``Mr.~Jones'',        % ~ produces an unbreakable interword space.
or within a word---especially when the word is a
symbol like
       \mbox{\emph{itemnum}}
that makes little sense when hyphenated across
lines.

Footnotes\footnote{This is an example of a footnote.}
pose no problem.

\LaTeX\ is good at typesetting mathematical formulas
like
       \( x-3y + z = 7 \)
or
       \( a_{1} > x^{2n} + y^{2n} > x' \)
or
       \( \ip{A}{B} = \sum_{i} a_{i} b_{i} \).
The spaces you type in a formula are
ignored.  Remember that a letter like
       $x$                   % $ ... $  and  \( ... \)  are equivalent
is a formula when it denotes a mathematical
symbol, and it should be typed as one.

\section{Displayed Text}

Text is displayed by indenting it from the left
margin.  Quotations are commonly displayed.  There
are short quotations
\begin{quote}
   This is a short quotation.  It consists of a
   single paragraph of text.  See how it is formatted.
\end{quote}
and longer ones.
\begin{quotation}
   This is a longer quotation.  It consists of two
   paragraphs of text, neither of which are
   particularly interesting.

   This is the second paragraph of the quotation.  It
   is just as dull as the first paragraph.
\end{quotation}
Another frequently-displayed structure is a list.
The following is an example of an \emph{itemized}
list.
\begin{itemize}
   \item This is the first item of an itemized list.
         Each item in the list is marked with a ``tick''.
         You don't have to worry about what kind of tick
         mark is used.

   \item This is the second item of the list.  It
         contains another list nested inside it.  The inner
         list is an \emph{enumerated} list.
         \begin{enumerate}
            \item This is the first item of an enumerated
                  list that is nested within the itemized list.

            \item This is the second item of the inner list.
                  \LaTeX\ allows you to nest lists deeper than
                  you really should.
         \end{enumerate}
         This is the rest of the second item of the outer
         list.  It is no more interesting than any other
         part of the item.
   \item This is the third item of the list.
\end{itemize}
You can even display poetry.
\begin{verse}
   There is an environment
    for verse \\             % The \\ command separates lines
   Whose features some poets % within a stanza.
   will curse.

                             % One or more blank lines separate stanzas.

   For instead of making\\
   Them do \emph{all} line breaking, \\
   It allows them to put too many words on a line when they'd rather be
   forced to be terse.
\end{verse}

Mathematical formulas may also be displayed.  A
displayed formula
is
one-line long; multiline
formulas require special formatting instructions.
   \[  \ip{\Gamma}{\psi'} = x'' + y^{2} + z_{i}^{n}\]
Don't start a paragraph with a displayed equation,
nor make one a paragraph by itself.

\end{document}               % End of document.

\section{Bench Huge Section 1015}
% synthetic bench section 1015
% This is a sample LaTeX input file.  (Version of 12 August 2004.)
%
% A '%' character causes TeX to ignore all remaining text on the line,
% and is used for comments like this one.

\documentclass{article}      % Specifies the document class

                             % The preamble begins here.
\title{An Example Document}  % Declares the document's title.
\author{Leslie Lamport}      % Declares the author's name.
\date{January 21, 1994}      % Deleting this command produces today's date.

\newcommand{\ip}[2]{(#1, #2)}
                             % Defines \ip{arg1}{arg2} to mean
                             % (arg1, arg2).

%\newcommand{\ip}[2]{\langle #1 | #2\rangle}
                             % This is an alternative definition of
                             % \ip that is commented out.

\begin{document}             % End of preamble and beginning of text.

\maketitle                   % Produces the title.

This is an example input file.  Comparing it with
the output it generates can show you how to
produce a simple document of your own.

\section{Ordinary Text}      % Produces section heading.  Lower-level
                             % sections are begun with similar
                             % \subsection and \subsubsection commands.

The ends  of words and sentences are marked
  by   spaces. It  doesn't matter how many
spaces    you type; one is as good as 100.  The
end of   a line counts as a space.

One   or more   blank lines denote the  end
of  a paragraph.

Since any number of consecutive spaces are treated
like a single one, the formatting of the input
file makes no difference to
      \LaTeX,                % The \LaTeX command generates the LaTeX logo.
but it makes a difference to you.  When you use
\LaTeX, making your input file as easy to read
as possible will be a great help as you write
your document and when you change it.  This sample
file shows how you can add comments to your own input
file.

Because printing is different from typewriting,
there are a number of things that you have to do
differently when preparing an input file than if
you were just typing the document directly.
Quotation marks like
       ``this''
have to be handled specially, as do quotes within
quotes:
       ``\,`this'            % \, separates the double and single quote.
        is what I just
        wrote, not  `that'\,''.

Dashes come in three sizes: an
       intra-word
dash, a medium dash for number ranges like
       1--2,
and a punctuation
       dash---like
this.

A sentence-ending space should be larger than the
space between words within a sentence.  You
sometimes have to type special commands in
conjunction with punctuation characters to get
this right, as in the following sentence.
       Gnats, gnus, etc.\ all  % `\ ' makes an inter-word space.
       begin with G\@.         % \@ marks end-of-sentence punctuation.
You should check the spaces after periods when
reading your output to make sure you haven't
forgotten any special cases.  Generating an
ellipsis
       \ldots\               % `\ ' is needed after `\ldots' because TeX
                             % ignores spaces after command names like \ldots
                             % made from \ + letters.
                             %
                             % Note how a `%' character causes TeX to ignore
                             % the end of the input line, so these blank lines
                             % do not start a new paragraph.
                             %
with the right spacing around the periods requires
a special command.

\LaTeX\ interprets some common characters as
commands, so you must type special commands to
generate them.  These characters include the
following:
       \$ \& \% \# \{ and \}.

In printing, text is usually emphasized with an
       \emph{italic}
type style.

\begin{em}
   A long segment of text can also be emphasized
   in this way.  Text within such a segment can be
   given \emph{additional} emphasis.
\end{em}

It is sometimes necessary to prevent \LaTeX\ from
breaking a line where it might otherwise do so.
This may be at a space, as between the ``Mr.''\ and
``Jones'' in
       ``Mr.~Jones'',        % ~ produces an unbreakable interword space.
or within a word---especially when the word is a
symbol like
       \mbox{\emph{itemnum}}
that makes little sense when hyphenated across
lines.

Footnotes\footnote{This is an example of a footnote.}
pose no problem.

\LaTeX\ is good at typesetting mathematical formulas
like
       \( x-3y + z = 7 \)
or
       \( a_{1} > x^{2n} + y^{2n} > x' \)
or
       \( \ip{A}{B} = \sum_{i} a_{i} b_{i} \).
The spaces you type in a formula are
ignored.  Remember that a letter like
       $x$                   % $ ... $  and  \( ... \)  are equivalent
is a formula when it denotes a mathematical
symbol, and it should be typed as one.

\section{Displayed Text}

Text is displayed by indenting it from the left
margin.  Quotations are commonly displayed.  There
are short quotations
\begin{quote}
   This is a short quotation.  It consists of a
   single paragraph of text.  See how it is formatted.
\end{quote}
and longer ones.
\begin{quotation}
   This is a longer quotation.  It consists of two
   paragraphs of text, neither of which are
   particularly interesting.

   This is the second paragraph of the quotation.  It
   is just as dull as the first paragraph.
\end{quotation}
Another frequently-displayed structure is a list.
The following is an example of an \emph{itemized}
list.
\begin{itemize}
   \item This is the first item of an itemized list.
         Each item in the list is marked with a ``tick''.
         You don't have to worry about what kind of tick
         mark is used.

   \item This is the second item of the list.  It
         contains another list nested inside it.  The inner
         list is an \emph{enumerated} list.
         \begin{enumerate}
            \item This is the first item of an enumerated
                  list that is nested within the itemized list.

            \item This is the second item of the inner list.
                  \LaTeX\ allows you to nest lists deeper than
                  you really should.
         \end{enumerate}
         This is the rest of the second item of the outer
         list.  It is no more interesting than any other
         part of the item.
   \item This is the third item of the list.
\end{itemize}
You can even display poetry.
\begin{verse}
   There is an environment
    for verse \\             % The \\ command separates lines
   Whose features some poets % within a stanza.
   will curse.

                             % One or more blank lines separate stanzas.

   For instead of making\\
   Them do \emph{all} line breaking, \\
   It allows them to put too many words on a line when they'd rather be
   forced to be terse.
\end{verse}

Mathematical formulas may also be displayed.  A
displayed formula
is
one-line long; multiline
formulas require special formatting instructions.
   \[  \ip{\Gamma}{\psi'} = x'' + y^{2} + z_{i}^{n}\]
Don't start a paragraph with a displayed equation,
nor make one a paragraph by itself.

\end{document}               % End of document.

\section{Bench Huge Section 1016}
% synthetic bench section 1016
% This is a sample LaTeX input file.  (Version of 12 August 2004.)
%
% A '%' character causes TeX to ignore all remaining text on the line,
% and is used for comments like this one.

\documentclass{article}      % Specifies the document class

                             % The preamble begins here.
\title{An Example Document}  % Declares the document's title.
\author{Leslie Lamport}      % Declares the author's name.
\date{January 21, 1994}      % Deleting this command produces today's date.

\newcommand{\ip}[2]{(#1, #2)}
                             % Defines \ip{arg1}{arg2} to mean
                             % (arg1, arg2).

%\newcommand{\ip}[2]{\langle #1 | #2\rangle}
                             % This is an alternative definition of
                             % \ip that is commented out.

\begin{document}             % End of preamble and beginning of text.

\maketitle                   % Produces the title.

This is an example input file.  Comparing it with
the output it generates can show you how to
produce a simple document of your own.

\section{Ordinary Text}      % Produces section heading.  Lower-level
                             % sections are begun with similar
                             % \subsection and \subsubsection commands.

The ends  of words and sentences are marked
  by   spaces. It  doesn't matter how many
spaces    you type; one is as good as 100.  The
end of   a line counts as a space.

One   or more   blank lines denote the  end
of  a paragraph.

Since any number of consecutive spaces are treated
like a single one, the formatting of the input
file makes no difference to
      \LaTeX,                % The \LaTeX command generates the LaTeX logo.
but it makes a difference to you.  When you use
\LaTeX, making your input file as easy to read
as possible will be a great help as you write
your document and when you change it.  This sample
file shows how you can add comments to your own input
file.

Because printing is different from typewriting,
there are a number of things that you have to do
differently when preparing an input file than if
you were just typing the document directly.
Quotation marks like
       ``this''
have to be handled specially, as do quotes within
quotes:
       ``\,`this'            % \, separates the double and single quote.
        is what I just
        wrote, not  `that'\,''.

Dashes come in three sizes: an
       intra-word
dash, a medium dash for number ranges like
       1--2,
and a punctuation
       dash---like
this.

A sentence-ending space should be larger than the
space between words within a sentence.  You
sometimes have to type special commands in
conjunction with punctuation characters to get
this right, as in the following sentence.
       Gnats, gnus, etc.\ all  % `\ ' makes an inter-word space.
       begin with G\@.         % \@ marks end-of-sentence punctuation.
You should check the spaces after periods when
reading your output to make sure you haven't
forgotten any special cases.  Generating an
ellipsis
       \ldots\               % `\ ' is needed after `\ldots' because TeX
                             % ignores spaces after command names like \ldots
                             % made from \ + letters.
                             %
                             % Note how a `%' character causes TeX to ignore
                             % the end of the input line, so these blank lines
                             % do not start a new paragraph.
                             %
with the right spacing around the periods requires
a special command.

\LaTeX\ interprets some common characters as
commands, so you must type special commands to
generate them.  These characters include the
following:
       \$ \& \% \# \{ and \}.

In printing, text is usually emphasized with an
       \emph{italic}
type style.

\begin{em}
   A long segment of text can also be emphasized
   in this way.  Text within such a segment can be
   given \emph{additional} emphasis.
\end{em}

It is sometimes necessary to prevent \LaTeX\ from
breaking a line where it might otherwise do so.
This may be at a space, as between the ``Mr.''\ and
``Jones'' in
       ``Mr.~Jones'',        % ~ produces an unbreakable interword space.
or within a word---especially when the word is a
symbol like
       \mbox{\emph{itemnum}}
that makes little sense when hyphenated across
lines.

Footnotes\footnote{This is an example of a footnote.}
pose no problem.

\LaTeX\ is good at typesetting mathematical formulas
like
       \( x-3y + z = 7 \)
or
       \( a_{1} > x^{2n} + y^{2n} > x' \)
or
       \( \ip{A}{B} = \sum_{i} a_{i} b_{i} \).
The spaces you type in a formula are
ignored.  Remember that a letter like
       $x$                   % $ ... $  and  \( ... \)  are equivalent
is a formula when it denotes a mathematical
symbol, and it should be typed as one.

\section{Displayed Text}

Text is displayed by indenting it from the left
margin.  Quotations are commonly displayed.  There
are short quotations
\begin{quote}
   This is a short quotation.  It consists of a
   single paragraph of text.  See how it is formatted.
\end{quote}
and longer ones.
\begin{quotation}
   This is a longer quotation.  It consists of two
   paragraphs of text, neither of which are
   particularly interesting.

   This is the second paragraph of the quotation.  It
   is just as dull as the first paragraph.
\end{quotation}
Another frequently-displayed structure is a list.
The following is an example of an \emph{itemized}
list.
\begin{itemize}
   \item This is the first item of an itemized list.
         Each item in the list is marked with a ``tick''.
         You don't have to worry about what kind of tick
         mark is used.

   \item This is the second item of the list.  It
         contains another list nested inside it.  The inner
         list is an \emph{enumerated} list.
         \begin{enumerate}
            \item This is the first item of an enumerated
                  list that is nested within the itemized list.

            \item This is the second item of the inner list.
                  \LaTeX\ allows you to nest lists deeper than
                  you really should.
         \end{enumerate}
         This is the rest of the second item of the outer
         list.  It is no more interesting than any other
         part of the item.
   \item This is the third item of the list.
\end{itemize}
You can even display poetry.
\begin{verse}
   There is an environment
    for verse \\             % The \\ command separates lines
   Whose features some poets % within a stanza.
   will curse.

                             % One or more blank lines separate stanzas.

   For instead of making\\
   Them do \emph{all} line breaking, \\
   It allows them to put too many words on a line when they'd rather be
   forced to be terse.
\end{verse}

Mathematical formulas may also be displayed.  A
displayed formula
is
one-line long; multiline
formulas require special formatting instructions.
   \[  \ip{\Gamma}{\psi'} = x'' + y^{2} + z_{i}^{n}\]
Don't start a paragraph with a displayed equation,
nor make one a paragraph by itself.

\end{document}               % End of document.

\section{Bench Huge Section 1017}
% synthetic bench section 1017
% This is a sample LaTeX input file.  (Version of 12 August 2004.)
%
% A '%' character causes TeX to ignore all remaining text on the line,
% and is used for comments like this one.

\documentclass{article}      % Specifies the document class

                             % The preamble begins here.
\title{An Example Document}  % Declares the document's title.
\author{Leslie Lamport}      % Declares the author's name.
\date{January 21, 1994}      % Deleting this command produces today's date.

\newcommand{\ip}[2]{(#1, #2)}
                             % Defines \ip{arg1}{arg2} to mean
                             % (arg1, arg2).

%\newcommand{\ip}[2]{\langle #1 | #2\rangle}
                             % This is an alternative definition of
                             % \ip that is commented out.

\begin{document}             % End of preamble and beginning of text.

\maketitle                   % Produces the title.

This is an example input file.  Comparing it with
the output it generates can show you how to
produce a simple document of your own.

\section{Ordinary Text}      % Produces section heading.  Lower-level
                             % sections are begun with similar
                             % \subsection and \subsubsection commands.

The ends  of words and sentences are marked
  by   spaces. It  doesn't matter how many
spaces    you type; one is as good as 100.  The
end of   a line counts as a space.

One   or more   blank lines denote the  end
of  a paragraph.

Since any number of consecutive spaces are treated
like a single one, the formatting of the input
file makes no difference to
      \LaTeX,                % The \LaTeX command generates the LaTeX logo.
but it makes a difference to you.  When you use
\LaTeX, making your input file as easy to read
as possible will be a great help as you write
your document and when you change it.  This sample
file shows how you can add comments to your own input
file.

Because printing is different from typewriting,
there are a number of things that you have to do
differently when preparing an input file than if
you were just typing the document directly.
Quotation marks like
       ``this''
have to be handled specially, as do quotes within
quotes:
       ``\,`this'            % \, separates the double and single quote.
        is what I just
        wrote, not  `that'\,''.

Dashes come in three sizes: an
       intra-word
dash, a medium dash for number ranges like
       1--2,
and a punctuation
       dash---like
this.

A sentence-ending space should be larger than the
space between words within a sentence.  You
sometimes have to type special commands in
conjunction with punctuation characters to get
this right, as in the following sentence.
       Gnats, gnus, etc.\ all  % `\ ' makes an inter-word space.
       begin with G\@.         % \@ marks end-of-sentence punctuation.
You should check the spaces after periods when
reading your output to make sure you haven't
forgotten any special cases.  Generating an
ellipsis
       \ldots\               % `\ ' is needed after `\ldots' because TeX
                             % ignores spaces after command names like \ldots
                             % made from \ + letters.
                             %
                             % Note how a `%' character causes TeX to ignore
                             % the end of the input line, so these blank lines
                             % do not start a new paragraph.
                             %
with the right spacing around the periods requires
a special command.

\LaTeX\ interprets some common characters as
commands, so you must type special commands to
generate them.  These characters include the
following:
       \$ \& \% \# \{ and \}.

In printing, text is usually emphasized with an
       \emph{italic}
type style.

\begin{em}
   A long segment of text can also be emphasized
   in this way.  Text within such a segment can be
   given \emph{additional} emphasis.
\end{em}

It is sometimes necessary to prevent \LaTeX\ from
breaking a line where it might otherwise do so.
This may be at a space, as between the ``Mr.''\ and
``Jones'' in
       ``Mr.~Jones'',        % ~ produces an unbreakable interword space.
or within a word---especially when the word is a
symbol like
       \mbox{\emph{itemnum}}
that makes little sense when hyphenated across
lines.

Footnotes\footnote{This is an example of a footnote.}
pose no problem.

\LaTeX\ is good at typesetting mathematical formulas
like
       \( x-3y + z = 7 \)
or
       \( a_{1} > x^{2n} + y^{2n} > x' \)
or
       \( \ip{A}{B} = \sum_{i} a_{i} b_{i} \).
The spaces you type in a formula are
ignored.  Remember that a letter like
       $x$                   % $ ... $  and  \( ... \)  are equivalent
is a formula when it denotes a mathematical
symbol, and it should be typed as one.

\section{Displayed Text}

Text is displayed by indenting it from the left
margin.  Quotations are commonly displayed.  There
are short quotations
\begin{quote}
   This is a short quotation.  It consists of a
   single paragraph of text.  See how it is formatted.
\end{quote}
and longer ones.
\begin{quotation}
   This is a longer quotation.  It consists of two
   paragraphs of text, neither of which are
   particularly interesting.

   This is the second paragraph of the quotation.  It
   is just as dull as the first paragraph.
\end{quotation}
Another frequently-displayed structure is a list.
The following is an example of an \emph{itemized}
list.
\begin{itemize}
   \item This is the first item of an itemized list.
         Each item in the list is marked with a ``tick''.
         You don't have to worry about what kind of tick
         mark is used.

   \item This is the second item of the list.  It
         contains another list nested inside it.  The inner
         list is an \emph{enumerated} list.
         \begin{enumerate}
            \item This is the first item of an enumerated
                  list that is nested within the itemized list.

            \item This is the second item of the inner list.
                  \LaTeX\ allows you to nest lists deeper than
                  you really should.
         \end{enumerate}
         This is the rest of the second item of the outer
         list.  It is no more interesting than any other
         part of the item.
   \item This is the third item of the list.
\end{itemize}
You can even display poetry.
\begin{verse}
   There is an environment
    for verse \\             % The \\ command separates lines
   Whose features some poets % within a stanza.
   will curse.

                             % One or more blank lines separate stanzas.

   For instead of making\\
   Them do \emph{all} line breaking, \\
   It allows them to put too many words on a line when they'd rather be
   forced to be terse.
\end{verse}

Mathematical formulas may also be displayed.  A
displayed formula
is
one-line long; multiline
formulas require special formatting instructions.
   \[  \ip{\Gamma}{\psi'} = x'' + y^{2} + z_{i}^{n}\]
Don't start a paragraph with a displayed equation,
nor make one a paragraph by itself.

\end{document}               % End of document.

\section{Bench Huge Section 1018}
% synthetic bench section 1018
% This is a sample LaTeX input file.  (Version of 12 August 2004.)
%
% A '%' character causes TeX to ignore all remaining text on the line,
% and is used for comments like this one.

\documentclass{article}      % Specifies the document class

                             % The preamble begins here.
\title{An Example Document}  % Declares the document's title.
\author{Leslie Lamport}      % Declares the author's name.
\date{January 21, 1994}      % Deleting this command produces today's date.

\newcommand{\ip}[2]{(#1, #2)}
                             % Defines \ip{arg1}{arg2} to mean
                             % (arg1, arg2).

%\newcommand{\ip}[2]{\langle #1 | #2\rangle}
                             % This is an alternative definition of
                             % \ip that is commented out.

\begin{document}             % End of preamble and beginning of text.

\maketitle                   % Produces the title.

This is an example input file.  Comparing it with
the output it generates can show you how to
produce a simple document of your own.

\section{Ordinary Text}      % Produces section heading.  Lower-level
                             % sections are begun with similar
                             % \subsection and \subsubsection commands.

The ends  of words and sentences are marked
  by   spaces. It  doesn't matter how many
spaces    you type; one is as good as 100.  The
end of   a line counts as a space.

One   or more   blank lines denote the  end
of  a paragraph.

Since any number of consecutive spaces are treated
like a single one, the formatting of the input
file makes no difference to
      \LaTeX,                % The \LaTeX command generates the LaTeX logo.
but it makes a difference to you.  When you use
\LaTeX, making your input file as easy to read
as possible will be a great help as you write
your document and when you change it.  This sample
file shows how you can add comments to your own input
file.

Because printing is different from typewriting,
there are a number of things that you have to do
differently when preparing an input file than if
you were just typing the document directly.
Quotation marks like
       ``this''
have to be handled specially, as do quotes within
quotes:
       ``\,`this'            % \, separates the double and single quote.
        is what I just
        wrote, not  `that'\,''.

Dashes come in three sizes: an
       intra-word
dash, a medium dash for number ranges like
       1--2,
and a punctuation
       dash---like
this.

A sentence-ending space should be larger than the
space between words within a sentence.  You
sometimes have to type special commands in
conjunction with punctuation characters to get
this right, as in the following sentence.
       Gnats, gnus, etc.\ all  % `\ ' makes an inter-word space.
       begin with G\@.         % \@ marks end-of-sentence punctuation.
You should check the spaces after periods when
reading your output to make sure you haven't
forgotten any special cases.  Generating an
ellipsis
       \ldots\               % `\ ' is needed after `\ldots' because TeX
                             % ignores spaces after command names like \ldots
                             % made from \ + letters.
                             %
                             % Note how a `%' character causes TeX to ignore
                             % the end of the input line, so these blank lines
                             % do not start a new paragraph.
                             %
with the right spacing around the periods requires
a special command.

\LaTeX\ interprets some common characters as
commands, so you must type special commands to
generate them.  These characters include the
following:
       \$ \& \% \# \{ and \}.

In printing, text is usually emphasized with an
       \emph{italic}
type style.

\begin{em}
   A long segment of text can also be emphasized
   in this way.  Text within such a segment can be
   given \emph{additional} emphasis.
\end{em}

It is sometimes necessary to prevent \LaTeX\ from
breaking a line where it might otherwise do so.
This may be at a space, as between the ``Mr.''\ and
``Jones'' in
       ``Mr.~Jones'',        % ~ produces an unbreakable interword space.
or within a word---especially when the word is a
symbol like
       \mbox{\emph{itemnum}}
that makes little sense when hyphenated across
lines.

Footnotes\footnote{This is an example of a footnote.}
pose no problem.

\LaTeX\ is good at typesetting mathematical formulas
like
       \( x-3y + z = 7 \)
or
       \( a_{1} > x^{2n} + y^{2n} > x' \)
or
       \( \ip{A}{B} = \sum_{i} a_{i} b_{i} \).
The spaces you type in a formula are
ignored.  Remember that a letter like
       $x$                   % $ ... $  and  \( ... \)  are equivalent
is a formula when it denotes a mathematical
symbol, and it should be typed as one.

\section{Displayed Text}

Text is displayed by indenting it from the left
margin.  Quotations are commonly displayed.  There
are short quotations
\begin{quote}
   This is a short quotation.  It consists of a
   single paragraph of text.  See how it is formatted.
\end{quote}
and longer ones.
\begin{quotation}
   This is a longer quotation.  It consists of two
   paragraphs of text, neither of which are
   particularly interesting.

   This is the second paragraph of the quotation.  It
   is just as dull as the first paragraph.
\end{quotation}
Another frequently-displayed structure is a list.
The following is an example of an \emph{itemized}
list.
\begin{itemize}
   \item This is the first item of an itemized list.
         Each item in the list is marked with a ``tick''.
         You don't have to worry about what kind of tick
         mark is used.

   \item This is the second item of the list.  It
         contains another list nested inside it.  The inner
         list is an \emph{enumerated} list.
         \begin{enumerate}
            \item This is the first item of an enumerated
                  list that is nested within the itemized list.

            \item This is the second item of the inner list.
                  \LaTeX\ allows you to nest lists deeper than
                  you really should.
         \end{enumerate}
         This is the rest of the second item of the outer
         list.  It is no more interesting than any other
         part of the item.
   \item This is the third item of the list.
\end{itemize}
You can even display poetry.
\begin{verse}
   There is an environment
    for verse \\             % The \\ command separates lines
   Whose features some poets % within a stanza.
   will curse.

                             % One or more blank lines separate stanzas.

   For instead of making\\
   Them do \emph{all} line breaking, \\
   It allows them to put too many words on a line when they'd rather be
   forced to be terse.
\end{verse}

Mathematical formulas may also be displayed.  A
displayed formula
is
one-line long; multiline
formulas require special formatting instructions.
   \[  \ip{\Gamma}{\psi'} = x'' + y^{2} + z_{i}^{n}\]
Don't start a paragraph with a displayed equation,
nor make one a paragraph by itself.

\end{document}               % End of document.

\section{Bench Huge Section 1019}
% synthetic bench section 1019
% This is a sample LaTeX input file.  (Version of 12 August 2004.)
%
% A '%' character causes TeX to ignore all remaining text on the line,
% and is used for comments like this one.

\documentclass{article}      % Specifies the document class

                             % The preamble begins here.
\title{An Example Document}  % Declares the document's title.
\author{Leslie Lamport}      % Declares the author's name.
\date{January 21, 1994}      % Deleting this command produces today's date.

\newcommand{\ip}[2]{(#1, #2)}
                             % Defines \ip{arg1}{arg2} to mean
                             % (arg1, arg2).

%\newcommand{\ip}[2]{\langle #1 | #2\rangle}
                             % This is an alternative definition of
                             % \ip that is commented out.

\begin{document}             % End of preamble and beginning of text.

\maketitle                   % Produces the title.

This is an example input file.  Comparing it with
the output it generates can show you how to
produce a simple document of your own.

\section{Ordinary Text}      % Produces section heading.  Lower-level
                             % sections are begun with similar
                             % \subsection and \subsubsection commands.

The ends  of words and sentences are marked
  by   spaces. It  doesn't matter how many
spaces    you type; one is as good as 100.  The
end of   a line counts as a space.

One   or more   blank lines denote the  end
of  a paragraph.

Since any number of consecutive spaces are treated
like a single one, the formatting of the input
file makes no difference to
      \LaTeX,                % The \LaTeX command generates the LaTeX logo.
but it makes a difference to you.  When you use
\LaTeX, making your input file as easy to read
as possible will be a great help as you write
your document and when you change it.  This sample
file shows how you can add comments to your own input
file.

Because printing is different from typewriting,
there are a number of things that you have to do
differently when preparing an input file than if
you were just typing the document directly.
Quotation marks like
       ``this''
have to be handled specially, as do quotes within
quotes:
       ``\,`this'            % \, separates the double and single quote.
        is what I just
        wrote, not  `that'\,''.

Dashes come in three sizes: an
       intra-word
dash, a medium dash for number ranges like
       1--2,
and a punctuation
       dash---like
this.

A sentence-ending space should be larger than the
space between words within a sentence.  You
sometimes have to type special commands in
conjunction with punctuation characters to get
this right, as in the following sentence.
       Gnats, gnus, etc.\ all  % `\ ' makes an inter-word space.
       begin with G\@.         % \@ marks end-of-sentence punctuation.
You should check the spaces after periods when
reading your output to make sure you haven't
forgotten any special cases.  Generating an
ellipsis
       \ldots\               % `\ ' is needed after `\ldots' because TeX
                             % ignores spaces after command names like \ldots
                             % made from \ + letters.
                             %
                             % Note how a `%' character causes TeX to ignore
                             % the end of the input line, so these blank lines
                             % do not start a new paragraph.
                             %
with the right spacing around the periods requires
a special command.

\LaTeX\ interprets some common characters as
commands, so you must type special commands to
generate them.  These characters include the
following:
       \$ \& \% \# \{ and \}.

In printing, text is usually emphasized with an
       \emph{italic}
type style.

\begin{em}
   A long segment of text can also be emphasized
   in this way.  Text within such a segment can be
   given \emph{additional} emphasis.
\end{em}

It is sometimes necessary to prevent \LaTeX\ from
breaking a line where it might otherwise do so.
This may be at a space, as between the ``Mr.''\ and
``Jones'' in
       ``Mr.~Jones'',        % ~ produces an unbreakable interword space.
or within a word---especially when the word is a
symbol like
       \mbox{\emph{itemnum}}
that makes little sense when hyphenated across
lines.

Footnotes\footnote{This is an example of a footnote.}
pose no problem.

\LaTeX\ is good at typesetting mathematical formulas
like
       \( x-3y + z = 7 \)
or
       \( a_{1} > x^{2n} + y^{2n} > x' \)
or
       \( \ip{A}{B} = \sum_{i} a_{i} b_{i} \).
The spaces you type in a formula are
ignored.  Remember that a letter like
       $x$                   % $ ... $  and  \( ... \)  are equivalent
is a formula when it denotes a mathematical
symbol, and it should be typed as one.

\section{Displayed Text}

Text is displayed by indenting it from the left
margin.  Quotations are commonly displayed.  There
are short quotations
\begin{quote}
   This is a short quotation.  It consists of a
   single paragraph of text.  See how it is formatted.
\end{quote}
and longer ones.
\begin{quotation}
   This is a longer quotation.  It consists of two
   paragraphs of text, neither of which are
   particularly interesting.

   This is the second paragraph of the quotation.  It
   is just as dull as the first paragraph.
\end{quotation}
Another frequently-displayed structure is a list.
The following is an example of an \emph{itemized}
list.
\begin{itemize}
   \item This is the first item of an itemized list.
         Each item in the list is marked with a ``tick''.
         You don't have to worry about what kind of tick
         mark is used.

   \item This is the second item of the list.  It
         contains another list nested inside it.  The inner
         list is an \emph{enumerated} list.
         \begin{enumerate}
            \item This is the first item of an enumerated
                  list that is nested within the itemized list.

            \item This is the second item of the inner list.
                  \LaTeX\ allows you to nest lists deeper than
                  you really should.
         \end{enumerate}
         This is the rest of the second item of the outer
         list.  It is no more interesting than any other
         part of the item.
   \item This is the third item of the list.
\end{itemize}
You can even display poetry.
\begin{verse}
   There is an environment
    for verse \\             % The \\ command separates lines
   Whose features some poets % within a stanza.
   will curse.

                             % One or more blank lines separate stanzas.

   For instead of making\\
   Them do \emph{all} line breaking, \\
   It allows them to put too many words on a line when they'd rather be
   forced to be terse.
\end{verse}

Mathematical formulas may also be displayed.  A
displayed formula
is
one-line long; multiline
formulas require special formatting instructions.
   \[  \ip{\Gamma}{\psi'} = x'' + y^{2} + z_{i}^{n}\]
Don't start a paragraph with a displayed equation,
nor make one a paragraph by itself.

\end{document}               % End of document.

\section{Bench Huge Section 1020}
% synthetic bench section 1020
% This is a sample LaTeX input file.  (Version of 12 August 2004.)
%
% A '%' character causes TeX to ignore all remaining text on the line,
% and is used for comments like this one.

\documentclass{article}      % Specifies the document class

                             % The preamble begins here.
\title{An Example Document}  % Declares the document's title.
\author{Leslie Lamport}      % Declares the author's name.
\date{January 21, 1994}      % Deleting this command produces today's date.

\newcommand{\ip}[2]{(#1, #2)}
                             % Defines \ip{arg1}{arg2} to mean
                             % (arg1, arg2).

%\newcommand{\ip}[2]{\langle #1 | #2\rangle}
                             % This is an alternative definition of
                             % \ip that is commented out.

\begin{document}             % End of preamble and beginning of text.

\maketitle                   % Produces the title.

This is an example input file.  Comparing it with
the output it generates can show you how to
produce a simple document of your own.

\section{Ordinary Text}      % Produces section heading.  Lower-level
                             % sections are begun with similar
                             % \subsection and \subsubsection commands.

The ends  of words and sentences are marked
  by   spaces. It  doesn't matter how many
spaces    you type; one is as good as 100.  The
end of   a line counts as a space.

One   or more   blank lines denote the  end
of  a paragraph.

Since any number of consecutive spaces are treated
like a single one, the formatting of the input
file makes no difference to
      \LaTeX,                % The \LaTeX command generates the LaTeX logo.
but it makes a difference to you.  When you use
\LaTeX, making your input file as easy to read
as possible will be a great help as you write
your document and when you change it.  This sample
file shows how you can add comments to your own input
file.

Because printing is different from typewriting,
there are a number of things that you have to do
differently when preparing an input file than if
you were just typing the document directly.
Quotation marks like
       ``this''
have to be handled specially, as do quotes within
quotes:
       ``\,`this'            % \, separates the double and single quote.
        is what I just
        wrote, not  `that'\,''.

Dashes come in three sizes: an
       intra-word
dash, a medium dash for number ranges like
       1--2,
and a punctuation
       dash---like
this.

A sentence-ending space should be larger than the
space between words within a sentence.  You
sometimes have to type special commands in
conjunction with punctuation characters to get
this right, as in the following sentence.
       Gnats, gnus, etc.\ all  % `\ ' makes an inter-word space.
       begin with G\@.         % \@ marks end-of-sentence punctuation.
You should check the spaces after periods when
reading your output to make sure you haven't
forgotten any special cases.  Generating an
ellipsis
       \ldots\               % `\ ' is needed after `\ldots' because TeX
                             % ignores spaces after command names like \ldots
                             % made from \ + letters.
                             %
                             % Note how a `%' character causes TeX to ignore
                             % the end of the input line, so these blank lines
                             % do not start a new paragraph.
                             %
with the right spacing around the periods requires
a special command.

\LaTeX\ interprets some common characters as
commands, so you must type special commands to
generate them.  These characters include the
following:
       \$ \& \% \# \{ and \}.

In printing, text is usually emphasized with an
       \emph{italic}
type style.

\begin{em}
   A long segment of text can also be emphasized
   in this way.  Text within such a segment can be
   given \emph{additional} emphasis.
\end{em}

It is sometimes necessary to prevent \LaTeX\ from
breaking a line where it might otherwise do so.
This may be at a space, as between the ``Mr.''\ and
``Jones'' in
       ``Mr.~Jones'',        % ~ produces an unbreakable interword space.
or within a word---especially when the word is a
symbol like
       \mbox{\emph{itemnum}}
that makes little sense when hyphenated across
lines.

Footnotes\footnote{This is an example of a footnote.}
pose no problem.

\LaTeX\ is good at typesetting mathematical formulas
like
       \( x-3y + z = 7 \)
or
       \( a_{1} > x^{2n} + y^{2n} > x' \)
or
       \( \ip{A}{B} = \sum_{i} a_{i} b_{i} \).
The spaces you type in a formula are
ignored.  Remember that a letter like
       $x$                   % $ ... $  and  \( ... \)  are equivalent
is a formula when it denotes a mathematical
symbol, and it should be typed as one.

\section{Displayed Text}

Text is displayed by indenting it from the left
margin.  Quotations are commonly displayed.  There
are short quotations
\begin{quote}
   This is a short quotation.  It consists of a
   single paragraph of text.  See how it is formatted.
\end{quote}
and longer ones.
\begin{quotation}
   This is a longer quotation.  It consists of two
   paragraphs of text, neither of which are
   particularly interesting.

   This is the second paragraph of the quotation.  It
   is just as dull as the first paragraph.
\end{quotation}
Another frequently-displayed structure is a list.
The following is an example of an \emph{itemized}
list.
\begin{itemize}
   \item This is the first item of an itemized list.
         Each item in the list is marked with a ``tick''.
         You don't have to worry about what kind of tick
         mark is used.

   \item This is the second item of the list.  It
         contains another list nested inside it.  The inner
         list is an \emph{enumerated} list.
         \begin{enumerate}
            \item This is the first item of an enumerated
                  list that is nested within the itemized list.

            \item This is the second item of the inner list.
                  \LaTeX\ allows you to nest lists deeper than
                  you really should.
         \end{enumerate}
         This is the rest of the second item of the outer
         list.  It is no more interesting than any other
         part of the item.
   \item This is the third item of the list.
\end{itemize}
You can even display poetry.
\begin{verse}
   There is an environment
    for verse \\             % The \\ command separates lines
   Whose features some poets % within a stanza.
   will curse.

                             % One or more blank lines separate stanzas.

   For instead of making\\
   Them do \emph{all} line breaking, \\
   It allows them to put too many words on a line when they'd rather be
   forced to be terse.
\end{verse}

Mathematical formulas may also be displayed.  A
displayed formula
is
one-line long; multiline
formulas require special formatting instructions.
   \[  \ip{\Gamma}{\psi'} = x'' + y^{2} + z_{i}^{n}\]
Don't start a paragraph with a displayed equation,
nor make one a paragraph by itself.

\end{document}               % End of document.

\section{Bench Huge Section 1021}
% synthetic bench section 1021
% This is a sample LaTeX input file.  (Version of 12 August 2004.)
%
% A '%' character causes TeX to ignore all remaining text on the line,
% and is used for comments like this one.

\documentclass{article}      % Specifies the document class

                             % The preamble begins here.
\title{An Example Document}  % Declares the document's title.
\author{Leslie Lamport}      % Declares the author's name.
\date{January 21, 1994}      % Deleting this command produces today's date.

\newcommand{\ip}[2]{(#1, #2)}
                             % Defines \ip{arg1}{arg2} to mean
                             % (arg1, arg2).

%\newcommand{\ip}[2]{\langle #1 | #2\rangle}
                             % This is an alternative definition of
                             % \ip that is commented out.

\begin{document}             % End of preamble and beginning of text.

\maketitle                   % Produces the title.

This is an example input file.  Comparing it with
the output it generates can show you how to
produce a simple document of your own.

\section{Ordinary Text}      % Produces section heading.  Lower-level
                             % sections are begun with similar
                             % \subsection and \subsubsection commands.

The ends  of words and sentences are marked
  by   spaces. It  doesn't matter how many
spaces    you type; one is as good as 100.  The
end of   a line counts as a space.

One   or more   blank lines denote the  end
of  a paragraph.

Since any number of consecutive spaces are treated
like a single one, the formatting of the input
file makes no difference to
      \LaTeX,                % The \LaTeX command generates the LaTeX logo.
but it makes a difference to you.  When you use
\LaTeX, making your input file as easy to read
as possible will be a great help as you write
your document and when you change it.  This sample
file shows how you can add comments to your own input
file.

Because printing is different from typewriting,
there are a number of things that you have to do
differently when preparing an input file than if
you were just typing the document directly.
Quotation marks like
       ``this''
have to be handled specially, as do quotes within
quotes:
       ``\,`this'            % \, separates the double and single quote.
        is what I just
        wrote, not  `that'\,''.

Dashes come in three sizes: an
       intra-word
dash, a medium dash for number ranges like
       1--2,
and a punctuation
       dash---like
this.

A sentence-ending space should be larger than the
space between words within a sentence.  You
sometimes have to type special commands in
conjunction with punctuation characters to get
this right, as in the following sentence.
       Gnats, gnus, etc.\ all  % `\ ' makes an inter-word space.
       begin with G\@.         % \@ marks end-of-sentence punctuation.
You should check the spaces after periods when
reading your output to make sure you haven't
forgotten any special cases.  Generating an
ellipsis
       \ldots\               % `\ ' is needed after `\ldots' because TeX
                             % ignores spaces after command names like \ldots
                             % made from \ + letters.
                             %
                             % Note how a `%' character causes TeX to ignore
                             % the end of the input line, so these blank lines
                             % do not start a new paragraph.
                             %
with the right spacing around the periods requires
a special command.

\LaTeX\ interprets some common characters as
commands, so you must type special commands to
generate them.  These characters include the
following:
       \$ \& \% \# \{ and \}.

In printing, text is usually emphasized with an
       \emph{italic}
type style.

\begin{em}
   A long segment of text can also be emphasized
   in this way.  Text within such a segment can be
   given \emph{additional} emphasis.
\end{em}

It is sometimes necessary to prevent \LaTeX\ from
breaking a line where it might otherwise do so.
This may be at a space, as between the ``Mr.''\ and
``Jones'' in
       ``Mr.~Jones'',        % ~ produces an unbreakable interword space.
or within a word---especially when the word is a
symbol like
       \mbox{\emph{itemnum}}
that makes little sense when hyphenated across
lines.

Footnotes\footnote{This is an example of a footnote.}
pose no problem.

\LaTeX\ is good at typesetting mathematical formulas
like
       \( x-3y + z = 7 \)
or
       \( a_{1} > x^{2n} + y^{2n} > x' \)
or
       \( \ip{A}{B} = \sum_{i} a_{i} b_{i} \).
The spaces you type in a formula are
ignored.  Remember that a letter like
       $x$                   % $ ... $  and  \( ... \)  are equivalent
is a formula when it denotes a mathematical
symbol, and it should be typed as one.

\section{Displayed Text}

Text is displayed by indenting it from the left
margin.  Quotations are commonly displayed.  There
are short quotations
\begin{quote}
   This is a short quotation.  It consists of a
   single paragraph of text.  See how it is formatted.
\end{quote}
and longer ones.
\begin{quotation}
   This is a longer quotation.  It consists of two
   paragraphs of text, neither of which are
   particularly interesting.

   This is the second paragraph of the quotation.  It
   is just as dull as the first paragraph.
\end{quotation}
Another frequently-displayed structure is a list.
The following is an example of an \emph{itemized}
list.
\begin{itemize}
   \item This is the first item of an itemized list.
         Each item in the list is marked with a ``tick''.
         You don't have to worry about what kind of tick
         mark is used.

   \item This is the second item of the list.  It
         contains another list nested inside it.  The inner
         list is an \emph{enumerated} list.
         \begin{enumerate}
            \item This is the first item of an enumerated
                  list that is nested within the itemized list.

            \item This is the second item of the inner list.
                  \LaTeX\ allows you to nest lists deeper than
                  you really should.
         \end{enumerate}
         This is the rest of the second item of the outer
         list.  It is no more interesting than any other
         part of the item.
   \item This is the third item of the list.
\end{itemize}
You can even display poetry.
\begin{verse}
   There is an environment
    for verse \\             % The \\ command separates lines
   Whose features some poets % within a stanza.
   will curse.

                             % One or more blank lines separate stanzas.

   For instead of making\\
   Them do \emph{all} line breaking, \\
   It allows them to put too many words on a line when they'd rather be
   forced to be terse.
\end{verse}

Mathematical formulas may also be displayed.  A
displayed formula
is
one-line long; multiline
formulas require special formatting instructions.
   \[  \ip{\Gamma}{\psi'} = x'' + y^{2} + z_{i}^{n}\]
Don't start a paragraph with a displayed equation,
nor make one a paragraph by itself.

\end{document}               % End of document.

\section{Bench Huge Section 1022}
% synthetic bench section 1022
% This is a sample LaTeX input file.  (Version of 12 August 2004.)
%
% A '%' character causes TeX to ignore all remaining text on the line,
% and is used for comments like this one.

\documentclass{article}      % Specifies the document class

                             % The preamble begins here.
\title{An Example Document}  % Declares the document's title.
\author{Leslie Lamport}      % Declares the author's name.
\date{January 21, 1994}      % Deleting this command produces today's date.

\newcommand{\ip}[2]{(#1, #2)}
                             % Defines \ip{arg1}{arg2} to mean
                             % (arg1, arg2).

%\newcommand{\ip}[2]{\langle #1 | #2\rangle}
                             % This is an alternative definition of
                             % \ip that is commented out.

\begin{document}             % End of preamble and beginning of text.

\maketitle                   % Produces the title.

This is an example input file.  Comparing it with
the output it generates can show you how to
produce a simple document of your own.

\section{Ordinary Text}      % Produces section heading.  Lower-level
                             % sections are begun with similar
                             % \subsection and \subsubsection commands.

The ends  of words and sentences are marked
  by   spaces. It  doesn't matter how many
spaces    you type; one is as good as 100.  The
end of   a line counts as a space.

One   or more   blank lines denote the  end
of  a paragraph.

Since any number of consecutive spaces are treated
like a single one, the formatting of the input
file makes no difference to
      \LaTeX,                % The \LaTeX command generates the LaTeX logo.
but it makes a difference to you.  When you use
\LaTeX, making your input file as easy to read
as possible will be a great help as you write
your document and when you change it.  This sample
file shows how you can add comments to your own input
file.

Because printing is different from typewriting,
there are a number of things that you have to do
differently when preparing an input file than if
you were just typing the document directly.
Quotation marks like
       ``this''
have to be handled specially, as do quotes within
quotes:
       ``\,`this'            % \, separates the double and single quote.
        is what I just
        wrote, not  `that'\,''.

Dashes come in three sizes: an
       intra-word
dash, a medium dash for number ranges like
       1--2,
and a punctuation
       dash---like
this.

A sentence-ending space should be larger than the
space between words within a sentence.  You
sometimes have to type special commands in
conjunction with punctuation characters to get
this right, as in the following sentence.
       Gnats, gnus, etc.\ all  % `\ ' makes an inter-word space.
       begin with G\@.         % \@ marks end-of-sentence punctuation.
You should check the spaces after periods when
reading your output to make sure you haven't
forgotten any special cases.  Generating an
ellipsis
       \ldots\               % `\ ' is needed after `\ldots' because TeX
                             % ignores spaces after command names like \ldots
                             % made from \ + letters.
                             %
                             % Note how a `%' character causes TeX to ignore
                             % the end of the input line, so these blank lines
                             % do not start a new paragraph.
                             %
with the right spacing around the periods requires
a special command.

\LaTeX\ interprets some common characters as
commands, so you must type special commands to
generate them.  These characters include the
following:
       \$ \& \% \# \{ and \}.

In printing, text is usually emphasized with an
       \emph{italic}
type style.

\begin{em}
   A long segment of text can also be emphasized
   in this way.  Text within such a segment can be
   given \emph{additional} emphasis.
\end{em}

It is sometimes necessary to prevent \LaTeX\ from
breaking a line where it might otherwise do so.
This may be at a space, as between the ``Mr.''\ and
``Jones'' in
       ``Mr.~Jones'',        % ~ produces an unbreakable interword space.
or within a word---especially when the word is a
symbol like
       \mbox{\emph{itemnum}}
that makes little sense when hyphenated across
lines.

Footnotes\footnote{This is an example of a footnote.}
pose no problem.

\LaTeX\ is good at typesetting mathematical formulas
like
       \( x-3y + z = 7 \)
or
       \( a_{1} > x^{2n} + y^{2n} > x' \)
or
       \( \ip{A}{B} = \sum_{i} a_{i} b_{i} \).
The spaces you type in a formula are
ignored.  Remember that a letter like
       $x$                   % $ ... $  and  \( ... \)  are equivalent
is a formula when it denotes a mathematical
symbol, and it should be typed as one.

\section{Displayed Text}

Text is displayed by indenting it from the left
margin.  Quotations are commonly displayed.  There
are short quotations
\begin{quote}
   This is a short quotation.  It consists of a
   single paragraph of text.  See how it is formatted.
\end{quote}
and longer ones.
\begin{quotation}
   This is a longer quotation.  It consists of two
   paragraphs of text, neither of which are
   particularly interesting.

   This is the second paragraph of the quotation.  It
   is just as dull as the first paragraph.
\end{quotation}
Another frequently-displayed structure is a list.
The following is an example of an \emph{itemized}
list.
\begin{itemize}
   \item This is the first item of an itemized list.
         Each item in the list is marked with a ``tick''.
         You don't have to worry about what kind of tick
         mark is used.

   \item This is the second item of the list.  It
         contains another list nested inside it.  The inner
         list is an \emph{enumerated} list.
         \begin{enumerate}
            \item This is the first item of an enumerated
                  list that is nested within the itemized list.

            \item This is the second item of the inner list.
                  \LaTeX\ allows you to nest lists deeper than
                  you really should.
         \end{enumerate}
         This is the rest of the second item of the outer
         list.  It is no more interesting than any other
         part of the item.
   \item This is the third item of the list.
\end{itemize}
You can even display poetry.
\begin{verse}
   There is an environment
    for verse \\             % The \\ command separates lines
   Whose features some poets % within a stanza.
   will curse.

                             % One or more blank lines separate stanzas.

   For instead of making\\
   Them do \emph{all} line breaking, \\
   It allows them to put too many words on a line when they'd rather be
   forced to be terse.
\end{verse}

Mathematical formulas may also be displayed.  A
displayed formula
is
one-line long; multiline
formulas require special formatting instructions.
   \[  \ip{\Gamma}{\psi'} = x'' + y^{2} + z_{i}^{n}\]
Don't start a paragraph with a displayed equation,
nor make one a paragraph by itself.

\end{document}               % End of document.

\section{Bench Huge Section 1023}
% synthetic bench section 1023
% This is a sample LaTeX input file.  (Version of 12 August 2004.)
%
% A '%' character causes TeX to ignore all remaining text on the line,
% and is used for comments like this one.

\documentclass{article}      % Specifies the document class

                             % The preamble begins here.
\title{An Example Document}  % Declares the document's title.
\author{Leslie Lamport}      % Declares the author's name.
\date{January 21, 1994}      % Deleting this command produces today's date.

\newcommand{\ip}[2]{(#1, #2)}
                             % Defines \ip{arg1}{arg2} to mean
                             % (arg1, arg2).

%\newcommand{\ip}[2]{\langle #1 | #2\rangle}
                             % This is an alternative definition of
                             % \ip that is commented out.

\begin{document}             % End of preamble and beginning of text.

\maketitle                   % Produces the title.

This is an example input file.  Comparing it with
the output it generates can show you how to
produce a simple document of your own.

\section{Ordinary Text}      % Produces section heading.  Lower-level
                             % sections are begun with similar
                             % \subsection and \subsubsection commands.

The ends  of words and sentences are marked
  by   spaces. It  doesn't matter how many
spaces    you type; one is as good as 100.  The
end of   a line counts as a space.

One   or more   blank lines denote the  end
of  a paragraph.

Since any number of consecutive spaces are treated
like a single one, the formatting of the input
file makes no difference to
      \LaTeX,                % The \LaTeX command generates the LaTeX logo.
but it makes a difference to you.  When you use
\LaTeX, making your input file as easy to read
as possible will be a great help as you write
your document and when you change it.  This sample
file shows how you can add comments to your own input
file.

Because printing is different from typewriting,
there are a number of things that you have to do
differently when preparing an input file than if
you were just typing the document directly.
Quotation marks like
       ``this''
have to be handled specially, as do quotes within
quotes:
       ``\,`this'            % \, separates the double and single quote.
        is what I just
        wrote, not  `that'\,''.

Dashes come in three sizes: an
       intra-word
dash, a medium dash for number ranges like
       1--2,
and a punctuation
       dash---like
this.

A sentence-ending space should be larger than the
space between words within a sentence.  You
sometimes have to type special commands in
conjunction with punctuation characters to get
this right, as in the following sentence.
       Gnats, gnus, etc.\ all  % `\ ' makes an inter-word space.
       begin with G\@.         % \@ marks end-of-sentence punctuation.
You should check the spaces after periods when
reading your output to make sure you haven't
forgotten any special cases.  Generating an
ellipsis
       \ldots\               % `\ ' is needed after `\ldots' because TeX
                             % ignores spaces after command names like \ldots
                             % made from \ + letters.
                             %
                             % Note how a `%' character causes TeX to ignore
                             % the end of the input line, so these blank lines
                             % do not start a new paragraph.
                             %
with the right spacing around the periods requires
a special command.

\LaTeX\ interprets some common characters as
commands, so you must type special commands to
generate them.  These characters include the
following:
       \$ \& \% \# \{ and \}.

In printing, text is usually emphasized with an
       \emph{italic}
type style.

\begin{em}
   A long segment of text can also be emphasized
   in this way.  Text within such a segment can be
   given \emph{additional} emphasis.
\end{em}

It is sometimes necessary to prevent \LaTeX\ from
breaking a line where it might otherwise do so.
This may be at a space, as between the ``Mr.''\ and
``Jones'' in
       ``Mr.~Jones'',        % ~ produces an unbreakable interword space.
or within a word---especially when the word is a
symbol like
       \mbox{\emph{itemnum}}
that makes little sense when hyphenated across
lines.

Footnotes\footnote{This is an example of a footnote.}
pose no problem.

\LaTeX\ is good at typesetting mathematical formulas
like
       \( x-3y + z = 7 \)
or
       \( a_{1} > x^{2n} + y^{2n} > x' \)
or
       \( \ip{A}{B} = \sum_{i} a_{i} b_{i} \).
The spaces you type in a formula are
ignored.  Remember that a letter like
       $x$                   % $ ... $  and  \( ... \)  are equivalent
is a formula when it denotes a mathematical
symbol, and it should be typed as one.

\section{Displayed Text}

Text is displayed by indenting it from the left
margin.  Quotations are commonly displayed.  There
are short quotations
\begin{quote}
   This is a short quotation.  It consists of a
   single paragraph of text.  See how it is formatted.
\end{quote}
and longer ones.
\begin{quotation}
   This is a longer quotation.  It consists of two
   paragraphs of text, neither of which are
   particularly interesting.

   This is the second paragraph of the quotation.  It
   is just as dull as the first paragraph.
\end{quotation}
Another frequently-displayed structure is a list.
The following is an example of an \emph{itemized}
list.
\begin{itemize}
   \item This is the first item of an itemized list.
         Each item in the list is marked with a ``tick''.
         You don't have to worry about what kind of tick
         mark is used.

   \item This is the second item of the list.  It
         contains another list nested inside it.  The inner
         list is an \emph{enumerated} list.
         \begin{enumerate}
            \item This is the first item of an enumerated
                  list that is nested within the itemized list.

            \item This is the second item of the inner list.
                  \LaTeX\ allows you to nest lists deeper than
                  you really should.
         \end{enumerate}
         This is the rest of the second item of the outer
         list.  It is no more interesting than any other
         part of the item.
   \item This is the third item of the list.
\end{itemize}
You can even display poetry.
\begin{verse}
   There is an environment
    for verse \\             % The \\ command separates lines
   Whose features some poets % within a stanza.
   will curse.

                             % One or more blank lines separate stanzas.

   For instead of making\\
   Them do \emph{all} line breaking, \\
   It allows them to put too many words on a line when they'd rather be
   forced to be terse.
\end{verse}

Mathematical formulas may also be displayed.  A
displayed formula
is
one-line long; multiline
formulas require special formatting instructions.
   \[  \ip{\Gamma}{\psi'} = x'' + y^{2} + z_{i}^{n}\]
Don't start a paragraph with a displayed equation,
nor make one a paragraph by itself.

\end{document}               % End of document.

\section{Bench Huge Section 1024}
% synthetic bench section 1024
% This is a sample LaTeX input file.  (Version of 12 August 2004.)
%
% A '%' character causes TeX to ignore all remaining text on the line,
% and is used for comments like this one.

\documentclass{article}      % Specifies the document class

                             % The preamble begins here.
\title{An Example Document}  % Declares the document's title.
\author{Leslie Lamport}      % Declares the author's name.
\date{January 21, 1994}      % Deleting this command produces today's date.

\newcommand{\ip}[2]{(#1, #2)}
                             % Defines \ip{arg1}{arg2} to mean
                             % (arg1, arg2).

%\newcommand{\ip}[2]{\langle #1 | #2\rangle}
                             % This is an alternative definition of
                             % \ip that is commented out.

\begin{document}             % End of preamble and beginning of text.

\maketitle                   % Produces the title.

This is an example input file.  Comparing it with
the output it generates can show you how to
produce a simple document of your own.

\section{Ordinary Text}      % Produces section heading.  Lower-level
                             % sections are begun with similar
                             % \subsection and \subsubsection commands.

The ends  of words and sentences are marked
  by   spaces. It  doesn't matter how many
spaces    you type; one is as good as 100.  The
end of   a line counts as a space.

One   or more   blank lines denote the  end
of  a paragraph.

Since any number of consecutive spaces are treated
like a single one, the formatting of the input
file makes no difference to
      \LaTeX,                % The \LaTeX command generates the LaTeX logo.
but it makes a difference to you.  When you use
\LaTeX, making your input file as easy to read
as possible will be a great help as you write
your document and when you change it.  This sample
file shows how you can add comments to your own input
file.

Because printing is different from typewriting,
there are a number of things that you have to do
differently when preparing an input file than if
you were just typing the document directly.
Quotation marks like
       ``this''
have to be handled specially, as do quotes within
quotes:
       ``\,`this'            % \, separates the double and single quote.
        is what I just
        wrote, not  `that'\,''.

Dashes come in three sizes: an
       intra-word
dash, a medium dash for number ranges like
       1--2,
and a punctuation
       dash---like
this.

A sentence-ending space should be larger than the
space between words within a sentence.  You
sometimes have to type special commands in
conjunction with punctuation characters to get
this right, as in the following sentence.
       Gnats, gnus, etc.\ all  % `\ ' makes an inter-word space.
       begin with G\@.         % \@ marks end-of-sentence punctuation.
You should check the spaces after periods when
reading your output to make sure you haven't
forgotten any special cases.  Generating an
ellipsis
       \ldots\               % `\ ' is needed after `\ldots' because TeX
                             % ignores spaces after command names like \ldots
                             % made from \ + letters.
                             %
                             % Note how a `%' character causes TeX to ignore
                             % the end of the input line, so these blank lines
                             % do not start a new paragraph.
                             %
with the right spacing around the periods requires
a special command.

\LaTeX\ interprets some common characters as
commands, so you must type special commands to
generate them.  These characters include the
following:
       \$ \& \% \# \{ and \}.

In printing, text is usually emphasized with an
       \emph{italic}
type style.

\begin{em}
   A long segment of text can also be emphasized
   in this way.  Text within such a segment can be
   given \emph{additional} emphasis.
\end{em}

It is sometimes necessary to prevent \LaTeX\ from
breaking a line where it might otherwise do so.
This may be at a space, as between the ``Mr.''\ and
``Jones'' in
       ``Mr.~Jones'',        % ~ produces an unbreakable interword space.
or within a word---especially when the word is a
symbol like
       \mbox{\emph{itemnum}}
that makes little sense when hyphenated across
lines.

Footnotes\footnote{This is an example of a footnote.}
pose no problem.

\LaTeX\ is good at typesetting mathematical formulas
like
       \( x-3y + z = 7 \)
or
       \( a_{1} > x^{2n} + y^{2n} > x' \)
or
       \( \ip{A}{B} = \sum_{i} a_{i} b_{i} \).
The spaces you type in a formula are
ignored.  Remember that a letter like
       $x$                   % $ ... $  and  \( ... \)  are equivalent
is a formula when it denotes a mathematical
symbol, and it should be typed as one.

\section{Displayed Text}

Text is displayed by indenting it from the left
margin.  Quotations are commonly displayed.  There
are short quotations
\begin{quote}
   This is a short quotation.  It consists of a
   single paragraph of text.  See how it is formatted.
\end{quote}
and longer ones.
\begin{quotation}
   This is a longer quotation.  It consists of two
   paragraphs of text, neither of which are
   particularly interesting.

   This is the second paragraph of the quotation.  It
   is just as dull as the first paragraph.
\end{quotation}
Another frequently-displayed structure is a list.
The following is an example of an \emph{itemized}
list.
\begin{itemize}
   \item This is the first item of an itemized list.
         Each item in the list is marked with a ``tick''.
         You don't have to worry about what kind of tick
         mark is used.

   \item This is the second item of the list.  It
         contains another list nested inside it.  The inner
         list is an \emph{enumerated} list.
         \begin{enumerate}
            \item This is the first item of an enumerated
                  list that is nested within the itemized list.

            \item This is the second item of the inner list.
                  \LaTeX\ allows you to nest lists deeper than
                  you really should.
         \end{enumerate}
         This is the rest of the second item of the outer
         list.  It is no more interesting than any other
         part of the item.
   \item This is the third item of the list.
\end{itemize}
You can even display poetry.
\begin{verse}
   There is an environment
    for verse \\             % The \\ command separates lines
   Whose features some poets % within a stanza.
   will curse.

                             % One or more blank lines separate stanzas.

   For instead of making\\
   Them do \emph{all} line breaking, \\
   It allows them to put too many words on a line when they'd rather be
   forced to be terse.
\end{verse}

Mathematical formulas may also be displayed.  A
displayed formula
is
one-line long; multiline
formulas require special formatting instructions.
   \[  \ip{\Gamma}{\psi'} = x'' + y^{2} + z_{i}^{n}\]
Don't start a paragraph with a displayed equation,
nor make one a paragraph by itself.

\end{document}               % End of document.

\section{Bench Huge Section 1025}
% synthetic bench section 1025
% This is a sample LaTeX input file.  (Version of 12 August 2004.)
%
% A '%' character causes TeX to ignore all remaining text on the line,
% and is used for comments like this one.

\documentclass{article}      % Specifies the document class

                             % The preamble begins here.
\title{An Example Document}  % Declares the document's title.
\author{Leslie Lamport}      % Declares the author's name.
\date{January 21, 1994}      % Deleting this command produces today's date.

\newcommand{\ip}[2]{(#1, #2)}
                             % Defines \ip{arg1}{arg2} to mean
                             % (arg1, arg2).

%\newcommand{\ip}[2]{\langle #1 | #2\rangle}
                             % This is an alternative definition of
                             % \ip that is commented out.

\begin{document}             % End of preamble and beginning of text.

\maketitle                   % Produces the title.

This is an example input file.  Comparing it with
the output it generates can show you how to
produce a simple document of your own.

\section{Ordinary Text}      % Produces section heading.  Lower-level
                             % sections are begun with similar
                             % \subsection and \subsubsection commands.

The ends  of words and sentences are marked
  by   spaces. It  doesn't matter how many
spaces    you type; one is as good as 100.  The
end of   a line counts as a space.

One   or more   blank lines denote the  end
of  a paragraph.

Since any number of consecutive spaces are treated
like a single one, the formatting of the input
file makes no difference to
      \LaTeX,                % The \LaTeX command generates the LaTeX logo.
but it makes a difference to you.  When you use
\LaTeX, making your input file as easy to read
as possible will be a great help as you write
your document and when you change it.  This sample
file shows how you can add comments to your own input
file.

Because printing is different from typewriting,
there are a number of things that you have to do
differently when preparing an input file than if
you were just typing the document directly.
Quotation marks like
       ``this''
have to be handled specially, as do quotes within
quotes:
       ``\,`this'            % \, separates the double and single quote.
        is what I just
        wrote, not  `that'\,''.

Dashes come in three sizes: an
       intra-word
dash, a medium dash for number ranges like
       1--2,
and a punctuation
       dash---like
this.

A sentence-ending space should be larger than the
space between words within a sentence.  You
sometimes have to type special commands in
conjunction with punctuation characters to get
this right, as in the following sentence.
       Gnats, gnus, etc.\ all  % `\ ' makes an inter-word space.
       begin with G\@.         % \@ marks end-of-sentence punctuation.
You should check the spaces after periods when
reading your output to make sure you haven't
forgotten any special cases.  Generating an
ellipsis
       \ldots\               % `\ ' is needed after `\ldots' because TeX
                             % ignores spaces after command names like \ldots
                             % made from \ + letters.
                             %
                             % Note how a `%' character causes TeX to ignore
                             % the end of the input line, so these blank lines
                             % do not start a new paragraph.
                             %
with the right spacing around the periods requires
a special command.

\LaTeX\ interprets some common characters as
commands, so you must type special commands to
generate them.  These characters include the
following:
       \$ \& \% \# \{ and \}.

In printing, text is usually emphasized with an
       \emph{italic}
type style.

\begin{em}
   A long segment of text can also be emphasized
   in this way.  Text within such a segment can be
   given \emph{additional} emphasis.
\end{em}

It is sometimes necessary to prevent \LaTeX\ from
breaking a line where it might otherwise do so.
This may be at a space, as between the ``Mr.''\ and
``Jones'' in
       ``Mr.~Jones'',        % ~ produces an unbreakable interword space.
or within a word---especially when the word is a
symbol like
       \mbox{\emph{itemnum}}
that makes little sense when hyphenated across
lines.

Footnotes\footnote{This is an example of a footnote.}
pose no problem.

\LaTeX\ is good at typesetting mathematical formulas
like
       \( x-3y + z = 7 \)
or
       \( a_{1} > x^{2n} + y^{2n} > x' \)
or
       \( \ip{A}{B} = \sum_{i} a_{i} b_{i} \).
The spaces you type in a formula are
ignored.  Remember that a letter like
       $x$                   % $ ... $  and  \( ... \)  are equivalent
is a formula when it denotes a mathematical
symbol, and it should be typed as one.

\section{Displayed Text}

Text is displayed by indenting it from the left
margin.  Quotations are commonly displayed.  There
are short quotations
\begin{quote}
   This is a short quotation.  It consists of a
   single paragraph of text.  See how it is formatted.
\end{quote}
and longer ones.
\begin{quotation}
   This is a longer quotation.  It consists of two
   paragraphs of text, neither of which are
   particularly interesting.

   This is the second paragraph of the quotation.  It
   is just as dull as the first paragraph.
\end{quotation}
Another frequently-displayed structure is a list.
The following is an example of an \emph{itemized}
list.
\begin{itemize}
   \item This is the first item of an itemized list.
         Each item in the list is marked with a ``tick''.
         You don't have to worry about what kind of tick
         mark is used.

   \item This is the second item of the list.  It
         contains another list nested inside it.  The inner
         list is an \emph{enumerated} list.
         \begin{enumerate}
            \item This is the first item of an enumerated
                  list that is nested within the itemized list.

            \item This is the second item of the inner list.
                  \LaTeX\ allows you to nest lists deeper than
                  you really should.
         \end{enumerate}
         This is the rest of the second item of the outer
         list.  It is no more interesting than any other
         part of the item.
   \item This is the third item of the list.
\end{itemize}
You can even display poetry.
\begin{verse}
   There is an environment
    for verse \\             % The \\ command separates lines
   Whose features some poets % within a stanza.
   will curse.

                             % One or more blank lines separate stanzas.

   For instead of making\\
   Them do \emph{all} line breaking, \\
   It allows them to put too many words on a line when they'd rather be
   forced to be terse.
\end{verse}

Mathematical formulas may also be displayed.  A
displayed formula
is
one-line long; multiline
formulas require special formatting instructions.
   \[  \ip{\Gamma}{\psi'} = x'' + y^{2} + z_{i}^{n}\]
Don't start a paragraph with a displayed equation,
nor make one a paragraph by itself.

\end{document}               % End of document.

\section{Bench Huge Section 1026}
% synthetic bench section 1026
% This is a sample LaTeX input file.  (Version of 12 August 2004.)
%
% A '%' character causes TeX to ignore all remaining text on the line,
% and is used for comments like this one.

\documentclass{article}      % Specifies the document class

                             % The preamble begins here.
\title{An Example Document}  % Declares the document's title.
\author{Leslie Lamport}      % Declares the author's name.
\date{January 21, 1994}      % Deleting this command produces today's date.

\newcommand{\ip}[2]{(#1, #2)}
                             % Defines \ip{arg1}{arg2} to mean
                             % (arg1, arg2).

%\newcommand{\ip}[2]{\langle #1 | #2\rangle}
                             % This is an alternative definition of
                             % \ip that is commented out.

\begin{document}             % End of preamble and beginning of text.

\maketitle                   % Produces the title.

This is an example input file.  Comparing it with
the output it generates can show you how to
produce a simple document of your own.

\section{Ordinary Text}      % Produces section heading.  Lower-level
                             % sections are begun with similar
                             % \subsection and \subsubsection commands.

The ends  of words and sentences are marked
  by   spaces. It  doesn't matter how many
spaces    you type; one is as good as 100.  The
end of   a line counts as a space.

One   or more   blank lines denote the  end
of  a paragraph.

Since any number of consecutive spaces are treated
like a single one, the formatting of the input
file makes no difference to
      \LaTeX,                % The \LaTeX command generates the LaTeX logo.
but it makes a difference to you.  When you use
\LaTeX, making your input file as easy to read
as possible will be a great help as you write
your document and when you change it.  This sample
file shows how you can add comments to your own input
file.

Because printing is different from typewriting,
there are a number of things that you have to do
differently when preparing an input file than if
you were just typing the document directly.
Quotation marks like
       ``this''
have to be handled specially, as do quotes within
quotes:
       ``\,`this'            % \, separates the double and single quote.
        is what I just
        wrote, not  `that'\,''.

Dashes come in three sizes: an
       intra-word
dash, a medium dash for number ranges like
       1--2,
and a punctuation
       dash---like
this.

A sentence-ending space should be larger than the
space between words within a sentence.  You
sometimes have to type special commands in
conjunction with punctuation characters to get
this right, as in the following sentence.
       Gnats, gnus, etc.\ all  % `\ ' makes an inter-word space.
       begin with G\@.         % \@ marks end-of-sentence punctuation.
You should check the spaces after periods when
reading your output to make sure you haven't
forgotten any special cases.  Generating an
ellipsis
       \ldots\               % `\ ' is needed after `\ldots' because TeX
                             % ignores spaces after command names like \ldots
                             % made from \ + letters.
                             %
                             % Note how a `%' character causes TeX to ignore
                             % the end of the input line, so these blank lines
                             % do not start a new paragraph.
                             %
with the right spacing around the periods requires
a special command.

\LaTeX\ interprets some common characters as
commands, so you must type special commands to
generate them.  These characters include the
following:
       \$ \& \% \# \{ and \}.

In printing, text is usually emphasized with an
       \emph{italic}
type style.

\begin{em}
   A long segment of text can also be emphasized
   in this way.  Text within such a segment can be
   given \emph{additional} emphasis.
\end{em}

It is sometimes necessary to prevent \LaTeX\ from
breaking a line where it might otherwise do so.
This may be at a space, as between the ``Mr.''\ and
``Jones'' in
       ``Mr.~Jones'',        % ~ produces an unbreakable interword space.
or within a word---especially when the word is a
symbol like
       \mbox{\emph{itemnum}}
that makes little sense when hyphenated across
lines.

Footnotes\footnote{This is an example of a footnote.}
pose no problem.

\LaTeX\ is good at typesetting mathematical formulas
like
       \( x-3y + z = 7 \)
or
       \( a_{1} > x^{2n} + y^{2n} > x' \)
or
       \( \ip{A}{B} = \sum_{i} a_{i} b_{i} \).
The spaces you type in a formula are
ignored.  Remember that a letter like
       $x$                   % $ ... $  and  \( ... \)  are equivalent
is a formula when it denotes a mathematical
symbol, and it should be typed as one.

\section{Displayed Text}

Text is displayed by indenting it from the left
margin.  Quotations are commonly displayed.  There
are short quotations
\begin{quote}
   This is a short quotation.  It consists of a
   single paragraph of text.  See how it is formatted.
\end{quote}
and longer ones.
\begin{quotation}
   This is a longer quotation.  It consists of two
   paragraphs of text, neither of which are
   particularly interesting.

   This is the second paragraph of the quotation.  It
   is just as dull as the first paragraph.
\end{quotation}
Another frequently-displayed structure is a list.
The following is an example of an \emph{itemized}
list.
\begin{itemize}
   \item This is the first item of an itemized list.
         Each item in the list is marked with a ``tick''.
         You don't have to worry about what kind of tick
         mark is used.

   \item This is the second item of the list.  It
         contains another list nested inside it.  The inner
         list is an \emph{enumerated} list.
         \begin{enumerate}
            \item This is the first item of an enumerated
                  list that is nested within the itemized list.

            \item This is the second item of the inner list.
                  \LaTeX\ allows you to nest lists deeper than
                  you really should.
         \end{enumerate}
         This is the rest of the second item of the outer
         list.  It is no more interesting than any other
         part of the item.
   \item This is the third item of the list.
\end{itemize}
You can even display poetry.
\begin{verse}
   There is an environment
    for verse \\             % The \\ command separates lines
   Whose features some poets % within a stanza.
   will curse.

                             % One or more blank lines separate stanzas.

   For instead of making\\
   Them do \emph{all} line breaking, \\
   It allows them to put too many words on a line when they'd rather be
   forced to be terse.
\end{verse}

Mathematical formulas may also be displayed.  A
displayed formula
is
one-line long; multiline
formulas require special formatting instructions.
   \[  \ip{\Gamma}{\psi'} = x'' + y^{2} + z_{i}^{n}\]
Don't start a paragraph with a displayed equation,
nor make one a paragraph by itself.

\end{document}               % End of document.

\section{Bench Huge Section 1027}
% synthetic bench section 1027
% This is a sample LaTeX input file.  (Version of 12 August 2004.)
%
% A '%' character causes TeX to ignore all remaining text on the line,
% and is used for comments like this one.

\documentclass{article}      % Specifies the document class

                             % The preamble begins here.
\title{An Example Document}  % Declares the document's title.
\author{Leslie Lamport}      % Declares the author's name.
\date{January 21, 1994}      % Deleting this command produces today's date.

\newcommand{\ip}[2]{(#1, #2)}
                             % Defines \ip{arg1}{arg2} to mean
                             % (arg1, arg2).

%\newcommand{\ip}[2]{\langle #1 | #2\rangle}
                             % This is an alternative definition of
                             % \ip that is commented out.

\begin{document}             % End of preamble and beginning of text.

\maketitle                   % Produces the title.

This is an example input file.  Comparing it with
the output it generates can show you how to
produce a simple document of your own.

\section{Ordinary Text}      % Produces section heading.  Lower-level
                             % sections are begun with similar
                             % \subsection and \subsubsection commands.

The ends  of words and sentences are marked
  by   spaces. It  doesn't matter how many
spaces    you type; one is as good as 100.  The
end of   a line counts as a space.

One   or more   blank lines denote the  end
of  a paragraph.

Since any number of consecutive spaces are treated
like a single one, the formatting of the input
file makes no difference to
      \LaTeX,                % The \LaTeX command generates the LaTeX logo.
but it makes a difference to you.  When you use
\LaTeX, making your input file as easy to read
as possible will be a great help as you write
your document and when you change it.  This sample
file shows how you can add comments to your own input
file.

Because printing is different from typewriting,
there are a number of things that you have to do
differently when preparing an input file than if
you were just typing the document directly.
Quotation marks like
       ``this''
have to be handled specially, as do quotes within
quotes:
       ``\,`this'            % \, separates the double and single quote.
        is what I just
        wrote, not  `that'\,''.

Dashes come in three sizes: an
       intra-word
dash, a medium dash for number ranges like
       1--2,
and a punctuation
       dash---like
this.

A sentence-ending space should be larger than the
space between words within a sentence.  You
sometimes have to type special commands in
conjunction with punctuation characters to get
this right, as in the following sentence.
       Gnats, gnus, etc.\ all  % `\ ' makes an inter-word space.
       begin with G\@.         % \@ marks end-of-sentence punctuation.
You should check the spaces after periods when
reading your output to make sure you haven't
forgotten any special cases.  Generating an
ellipsis
       \ldots\               % `\ ' is needed after `\ldots' because TeX
                             % ignores spaces after command names like \ldots
                             % made from \ + letters.
                             %
                             % Note how a `%' character causes TeX to ignore
                             % the end of the input line, so these blank lines
                             % do not start a new paragraph.
                             %
with the right spacing around the periods requires
a special command.

\LaTeX\ interprets some common characters as
commands, so you must type special commands to
generate them.  These characters include the
following:
       \$ \& \% \# \{ and \}.

In printing, text is usually emphasized with an
       \emph{italic}
type style.

\begin{em}
   A long segment of text can also be emphasized
   in this way.  Text within such a segment can be
   given \emph{additional} emphasis.
\end{em}

It is sometimes necessary to prevent \LaTeX\ from
breaking a line where it might otherwise do so.
This may be at a space, as between the ``Mr.''\ and
``Jones'' in
       ``Mr.~Jones'',        % ~ produces an unbreakable interword space.
or within a word---especially when the word is a
symbol like
       \mbox{\emph{itemnum}}
that makes little sense when hyphenated across
lines.

Footnotes\footnote{This is an example of a footnote.}
pose no problem.

\LaTeX\ is good at typesetting mathematical formulas
like
       \( x-3y + z = 7 \)
or
       \( a_{1} > x^{2n} + y^{2n} > x' \)
or
       \( \ip{A}{B} = \sum_{i} a_{i} b_{i} \).
The spaces you type in a formula are
ignored.  Remember that a letter like
       $x$                   % $ ... $  and  \( ... \)  are equivalent
is a formula when it denotes a mathematical
symbol, and it should be typed as one.

\section{Displayed Text}

Text is displayed by indenting it from the left
margin.  Quotations are commonly displayed.  There
are short quotations
\begin{quote}
   This is a short quotation.  It consists of a
   single paragraph of text.  See how it is formatted.
\end{quote}
and longer ones.
\begin{quotation}
   This is a longer quotation.  It consists of two
   paragraphs of text, neither of which are
   particularly interesting.

   This is the second paragraph of the quotation.  It
   is just as dull as the first paragraph.
\end{quotation}
Another frequently-displayed structure is a list.
The following is an example of an \emph{itemized}
list.
\begin{itemize}
   \item This is the first item of an itemized list.
         Each item in the list is marked with a ``tick''.
         You don't have to worry about what kind of tick
         mark is used.

   \item This is the second item of the list.  It
         contains another list nested inside it.  The inner
         list is an \emph{enumerated} list.
         \begin{enumerate}
            \item This is the first item of an enumerated
                  list that is nested within the itemized list.

            \item This is the second item of the inner list.
                  \LaTeX\ allows you to nest lists deeper than
                  you really should.
         \end{enumerate}
         This is the rest of the second item of the outer
         list.  It is no more interesting than any other
         part of the item.
   \item This is the third item of the list.
\end{itemize}
You can even display poetry.
\begin{verse}
   There is an environment
    for verse \\             % The \\ command separates lines
   Whose features some poets % within a stanza.
   will curse.

                             % One or more blank lines separate stanzas.

   For instead of making\\
   Them do \emph{all} line breaking, \\
   It allows them to put too many words on a line when they'd rather be
   forced to be terse.
\end{verse}

Mathematical formulas may also be displayed.  A
displayed formula
is
one-line long; multiline
formulas require special formatting instructions.
   \[  \ip{\Gamma}{\psi'} = x'' + y^{2} + z_{i}^{n}\]
Don't start a paragraph with a displayed equation,
nor make one a paragraph by itself.

\end{document}               % End of document.

\section{Bench Huge Section 1028}
% synthetic bench section 1028
% This is a sample LaTeX input file.  (Version of 12 August 2004.)
%
% A '%' character causes TeX to ignore all remaining text on the line,
% and is used for comments like this one.

\documentclass{article}      % Specifies the document class

                             % The preamble begins here.
\title{An Example Document}  % Declares the document's title.
\author{Leslie Lamport}      % Declares the author's name.
\date{January 21, 1994}      % Deleting this command produces today's date.

\newcommand{\ip}[2]{(#1, #2)}
                             % Defines \ip{arg1}{arg2} to mean
                             % (arg1, arg2).

%\newcommand{\ip}[2]{\langle #1 | #2\rangle}
                             % This is an alternative definition of
                             % \ip that is commented out.

\begin{document}             % End of preamble and beginning of text.

\maketitle                   % Produces the title.

This is an example input file.  Comparing it with
the output it generates can show you how to
produce a simple document of your own.

\section{Ordinary Text}      % Produces section heading.  Lower-level
                             % sections are begun with similar
                             % \subsection and \subsubsection commands.

The ends  of words and sentences are marked
  by   spaces. It  doesn't matter how many
spaces    you type; one is as good as 100.  The
end of   a line counts as a space.

One   or more   blank lines denote the  end
of  a paragraph.

Since any number of consecutive spaces are treated
like a single one, the formatting of the input
file makes no difference to
      \LaTeX,                % The \LaTeX command generates the LaTeX logo.
but it makes a difference to you.  When you use
\LaTeX, making your input file as easy to read
as possible will be a great help as you write
your document and when you change it.  This sample
file shows how you can add comments to your own input
file.

Because printing is different from typewriting,
there are a number of things that you have to do
differently when preparing an input file than if
you were just typing the document directly.
Quotation marks like
       ``this''
have to be handled specially, as do quotes within
quotes:
       ``\,`this'            % \, separates the double and single quote.
        is what I just
        wrote, not  `that'\,''.

Dashes come in three sizes: an
       intra-word
dash, a medium dash for number ranges like
       1--2,
and a punctuation
       dash---like
this.

A sentence-ending space should be larger than the
space between words within a sentence.  You
sometimes have to type special commands in
conjunction with punctuation characters to get
this right, as in the following sentence.
       Gnats, gnus, etc.\ all  % `\ ' makes an inter-word space.
       begin with G\@.         % \@ marks end-of-sentence punctuation.
You should check the spaces after periods when
reading your output to make sure you haven't
forgotten any special cases.  Generating an
ellipsis
       \ldots\               % `\ ' is needed after `\ldots' because TeX
                             % ignores spaces after command names like \ldots
                             % made from \ + letters.
                             %
                             % Note how a `%' character causes TeX to ignore
                             % the end of the input line, so these blank lines
                             % do not start a new paragraph.
                             %
with the right spacing around the periods requires
a special command.

\LaTeX\ interprets some common characters as
commands, so you must type special commands to
generate them.  These characters include the
following:
       \$ \& \% \# \{ and \}.

In printing, text is usually emphasized with an
       \emph{italic}
type style.

\begin{em}
   A long segment of text can also be emphasized
   in this way.  Text within such a segment can be
   given \emph{additional} emphasis.
\end{em}

It is sometimes necessary to prevent \LaTeX\ from
breaking a line where it might otherwise do so.
This may be at a space, as between the ``Mr.''\ and
``Jones'' in
       ``Mr.~Jones'',        % ~ produces an unbreakable interword space.
or within a word---especially when the word is a
symbol like
       \mbox{\emph{itemnum}}
that makes little sense when hyphenated across
lines.

Footnotes\footnote{This is an example of a footnote.}
pose no problem.

\LaTeX\ is good at typesetting mathematical formulas
like
       \( x-3y + z = 7 \)
or
       \( a_{1} > x^{2n} + y^{2n} > x' \)
or
       \( \ip{A}{B} = \sum_{i} a_{i} b_{i} \).
The spaces you type in a formula are
ignored.  Remember that a letter like
       $x$                   % $ ... $  and  \( ... \)  are equivalent
is a formula when it denotes a mathematical
symbol, and it should be typed as one.

\section{Displayed Text}

Text is displayed by indenting it from the left
margin.  Quotations are commonly displayed.  There
are short quotations
\begin{quote}
   This is a short quotation.  It consists of a
   single paragraph of text.  See how it is formatted.
\end{quote}
and longer ones.
\begin{quotation}
   This is a longer quotation.  It consists of two
   paragraphs of text, neither of which are
   particularly interesting.

   This is the second paragraph of the quotation.  It
   is just as dull as the first paragraph.
\end{quotation}
Another frequently-displayed structure is a list.
The following is an example of an \emph{itemized}
list.
\begin{itemize}
   \item This is the first item of an itemized list.
         Each item in the list is marked with a ``tick''.
         You don't have to worry about what kind of tick
         mark is used.

   \item This is the second item of the list.  It
         contains another list nested inside it.  The inner
         list is an \emph{enumerated} list.
         \begin{enumerate}
            \item This is the first item of an enumerated
                  list that is nested within the itemized list.

            \item This is the second item of the inner list.
                  \LaTeX\ allows you to nest lists deeper than
                  you really should.
         \end{enumerate}
         This is the rest of the second item of the outer
         list.  It is no more interesting than any other
         part of the item.
   \item This is the third item of the list.
\end{itemize}
You can even display poetry.
\begin{verse}
   There is an environment
    for verse \\             % The \\ command separates lines
   Whose features some poets % within a stanza.
   will curse.

                             % One or more blank lines separate stanzas.

   For instead of making\\
   Them do \emph{all} line breaking, \\
   It allows them to put too many words on a line when they'd rather be
   forced to be terse.
\end{verse}

Mathematical formulas may also be displayed.  A
displayed formula
is
one-line long; multiline
formulas require special formatting instructions.
   \[  \ip{\Gamma}{\psi'} = x'' + y^{2} + z_{i}^{n}\]
Don't start a paragraph with a displayed equation,
nor make one a paragraph by itself.

\end{document}               % End of document.

\section{Bench Huge Section 1029}
% synthetic bench section 1029
% This is a sample LaTeX input file.  (Version of 12 August 2004.)
%
% A '%' character causes TeX to ignore all remaining text on the line,
% and is used for comments like this one.

\documentclass{article}      % Specifies the document class

                             % The preamble begins here.
\title{An Example Document}  % Declares the document's title.
\author{Leslie Lamport}      % Declares the author's name.
\date{January 21, 1994}      % Deleting this command produces today's date.

\newcommand{\ip}[2]{(#1, #2)}
                             % Defines \ip{arg1}{arg2} to mean
                             % (arg1, arg2).

%\newcommand{\ip}[2]{\langle #1 | #2\rangle}
                             % This is an alternative definition of
                             % \ip that is commented out.

\begin{document}             % End of preamble and beginning of text.

\maketitle                   % Produces the title.

This is an example input file.  Comparing it with
the output it generates can show you how to
produce a simple document of your own.

\section{Ordinary Text}      % Produces section heading.  Lower-level
                             % sections are begun with similar
                             % \subsection and \subsubsection commands.

The ends  of words and sentences are marked
  by   spaces. It  doesn't matter how many
spaces    you type; one is as good as 100.  The
end of   a line counts as a space.

One   or more   blank lines denote the  end
of  a paragraph.

Since any number of consecutive spaces are treated
like a single one, the formatting of the input
file makes no difference to
      \LaTeX,                % The \LaTeX command generates the LaTeX logo.
but it makes a difference to you.  When you use
\LaTeX, making your input file as easy to read
as possible will be a great help as you write
your document and when you change it.  This sample
file shows how you can add comments to your own input
file.

Because printing is different from typewriting,
there are a number of things that you have to do
differently when preparing an input file than if
you were just typing the document directly.
Quotation marks like
       ``this''
have to be handled specially, as do quotes within
quotes:
       ``\,`this'            % \, separates the double and single quote.
        is what I just
        wrote, not  `that'\,''.

Dashes come in three sizes: an
       intra-word
dash, a medium dash for number ranges like
       1--2,
and a punctuation
       dash---like
this.

A sentence-ending space should be larger than the
space between words within a sentence.  You
sometimes have to type special commands in
conjunction with punctuation characters to get
this right, as in the following sentence.
       Gnats, gnus, etc.\ all  % `\ ' makes an inter-word space.
       begin with G\@.         % \@ marks end-of-sentence punctuation.
You should check the spaces after periods when
reading your output to make sure you haven't
forgotten any special cases.  Generating an
ellipsis
       \ldots\               % `\ ' is needed after `\ldots' because TeX
                             % ignores spaces after command names like \ldots
                             % made from \ + letters.
                             %
                             % Note how a `%' character causes TeX to ignore
                             % the end of the input line, so these blank lines
                             % do not start a new paragraph.
                             %
with the right spacing around the periods requires
a special command.

\LaTeX\ interprets some common characters as
commands, so you must type special commands to
generate them.  These characters include the
following:
       \$ \& \% \# \{ and \}.

In printing, text is usually emphasized with an
       \emph{italic}
type style.

\begin{em}
   A long segment of text can also be emphasized
   in this way.  Text within such a segment can be
   given \emph{additional} emphasis.
\end{em}

It is sometimes necessary to prevent \LaTeX\ from
breaking a line where it might otherwise do so.
This may be at a space, as between the ``Mr.''\ and
``Jones'' in
       ``Mr.~Jones'',        % ~ produces an unbreakable interword space.
or within a word---especially when the word is a
symbol like
       \mbox{\emph{itemnum}}
that makes little sense when hyphenated across
lines.

Footnotes\footnote{This is an example of a footnote.}
pose no problem.

\LaTeX\ is good at typesetting mathematical formulas
like
       \( x-3y + z = 7 \)
or
       \( a_{1} > x^{2n} + y^{2n} > x' \)
or
       \( \ip{A}{B} = \sum_{i} a_{i} b_{i} \).
The spaces you type in a formula are
ignored.  Remember that a letter like
       $x$                   % $ ... $  and  \( ... \)  are equivalent
is a formula when it denotes a mathematical
symbol, and it should be typed as one.

\section{Displayed Text}

Text is displayed by indenting it from the left
margin.  Quotations are commonly displayed.  There
are short quotations
\begin{quote}
   This is a short quotation.  It consists of a
   single paragraph of text.  See how it is formatted.
\end{quote}
and longer ones.
\begin{quotation}
   This is a longer quotation.  It consists of two
   paragraphs of text, neither of which are
   particularly interesting.

   This is the second paragraph of the quotation.  It
   is just as dull as the first paragraph.
\end{quotation}
Another frequently-displayed structure is a list.
The following is an example of an \emph{itemized}
list.
\begin{itemize}
   \item This is the first item of an itemized list.
         Each item in the list is marked with a ``tick''.
         You don't have to worry about what kind of tick
         mark is used.

   \item This is the second item of the list.  It
         contains another list nested inside it.  The inner
         list is an \emph{enumerated} list.
         \begin{enumerate}
            \item This is the first item of an enumerated
                  list that is nested within the itemized list.

            \item This is the second item of the inner list.
                  \LaTeX\ allows you to nest lists deeper than
                  you really should.
         \end{enumerate}
         This is the rest of the second item of the outer
         list.  It is no more interesting than any other
         part of the item.
   \item This is the third item of the list.
\end{itemize}
You can even display poetry.
\begin{verse}
   There is an environment
    for verse \\             % The \\ command separates lines
   Whose features some poets % within a stanza.
   will curse.

                             % One or more blank lines separate stanzas.

   For instead of making\\
   Them do \emph{all} line breaking, \\
   It allows them to put too many words on a line when they'd rather be
   forced to be terse.
\end{verse}

Mathematical formulas may also be displayed.  A
displayed formula
is
one-line long; multiline
formulas require special formatting instructions.
   \[  \ip{\Gamma}{\psi'} = x'' + y^{2} + z_{i}^{n}\]
Don't start a paragraph with a displayed equation,
nor make one a paragraph by itself.

\end{document}               % End of document.

\section{Bench Huge Section 1030}
% synthetic bench section 1030
% This is a sample LaTeX input file.  (Version of 12 August 2004.)
%
% A '%' character causes TeX to ignore all remaining text on the line,
% and is used for comments like this one.

\documentclass{article}      % Specifies the document class

                             % The preamble begins here.
\title{An Example Document}  % Declares the document's title.
\author{Leslie Lamport}      % Declares the author's name.
\date{January 21, 1994}      % Deleting this command produces today's date.

\newcommand{\ip}[2]{(#1, #2)}
                             % Defines \ip{arg1}{arg2} to mean
                             % (arg1, arg2).

%\newcommand{\ip}[2]{\langle #1 | #2\rangle}
                             % This is an alternative definition of
                             % \ip that is commented out.

\begin{document}             % End of preamble and beginning of text.

\maketitle                   % Produces the title.

This is an example input file.  Comparing it with
the output it generates can show you how to
produce a simple document of your own.

\section{Ordinary Text}      % Produces section heading.  Lower-level
                             % sections are begun with similar
                             % \subsection and \subsubsection commands.

The ends  of words and sentences are marked
  by   spaces. It  doesn't matter how many
spaces    you type; one is as good as 100.  The
end of   a line counts as a space.

One   or more   blank lines denote the  end
of  a paragraph.

Since any number of consecutive spaces are treated
like a single one, the formatting of the input
file makes no difference to
      \LaTeX,                % The \LaTeX command generates the LaTeX logo.
but it makes a difference to you.  When you use
\LaTeX, making your input file as easy to read
as possible will be a great help as you write
your document and when you change it.  This sample
file shows how you can add comments to your own input
file.

Because printing is different from typewriting,
there are a number of things that you have to do
differently when preparing an input file than if
you were just typing the document directly.
Quotation marks like
       ``this''
have to be handled specially, as do quotes within
quotes:
       ``\,`this'            % \, separates the double and single quote.
        is what I just
        wrote, not  `that'\,''.

Dashes come in three sizes: an
       intra-word
dash, a medium dash for number ranges like
       1--2,
and a punctuation
       dash---like
this.

A sentence-ending space should be larger than the
space between words within a sentence.  You
sometimes have to type special commands in
conjunction with punctuation characters to get
this right, as in the following sentence.
       Gnats, gnus, etc.\ all  % `\ ' makes an inter-word space.
       begin with G\@.         % \@ marks end-of-sentence punctuation.
You should check the spaces after periods when
reading your output to make sure you haven't
forgotten any special cases.  Generating an
ellipsis
       \ldots\               % `\ ' is needed after `\ldots' because TeX
                             % ignores spaces after command names like \ldots
                             % made from \ + letters.
                             %
                             % Note how a `%' character causes TeX to ignore
                             % the end of the input line, so these blank lines
                             % do not start a new paragraph.
                             %
with the right spacing around the periods requires
a special command.

\LaTeX\ interprets some common characters as
commands, so you must type special commands to
generate them.  These characters include the
following:
       \$ \& \% \# \{ and \}.

In printing, text is usually emphasized with an
       \emph{italic}
type style.

\begin{em}
   A long segment of text can also be emphasized
   in this way.  Text within such a segment can be
   given \emph{additional} emphasis.
\end{em}

It is sometimes necessary to prevent \LaTeX\ from
breaking a line where it might otherwise do so.
This may be at a space, as between the ``Mr.''\ and
``Jones'' in
       ``Mr.~Jones'',        % ~ produces an unbreakable interword space.
or within a word---especially when the word is a
symbol like
       \mbox{\emph{itemnum}}
that makes little sense when hyphenated across
lines.

Footnotes\footnote{This is an example of a footnote.}
pose no problem.

\LaTeX\ is good at typesetting mathematical formulas
like
       \( x-3y + z = 7 \)
or
       \( a_{1} > x^{2n} + y^{2n} > x' \)
or
       \( \ip{A}{B} = \sum_{i} a_{i} b_{i} \).
The spaces you type in a formula are
ignored.  Remember that a letter like
       $x$                   % $ ... $  and  \( ... \)  are equivalent
is a formula when it denotes a mathematical
symbol, and it should be typed as one.

\section{Displayed Text}

Text is displayed by indenting it from the left
margin.  Quotations are commonly displayed.  There
are short quotations
\begin{quote}
   This is a short quotation.  It consists of a
   single paragraph of text.  See how it is formatted.
\end{quote}
and longer ones.
\begin{quotation}
   This is a longer quotation.  It consists of two
   paragraphs of text, neither of which are
   particularly interesting.

   This is the second paragraph of the quotation.  It
   is just as dull as the first paragraph.
\end{quotation}
Another frequently-displayed structure is a list.
The following is an example of an \emph{itemized}
list.
\begin{itemize}
   \item This is the first item of an itemized list.
         Each item in the list is marked with a ``tick''.
         You don't have to worry about what kind of tick
         mark is used.

   \item This is the second item of the list.  It
         contains another list nested inside it.  The inner
         list is an \emph{enumerated} list.
         \begin{enumerate}
            \item This is the first item of an enumerated
                  list that is nested within the itemized list.

            \item This is the second item of the inner list.
                  \LaTeX\ allows you to nest lists deeper than
                  you really should.
         \end{enumerate}
         This is the rest of the second item of the outer
         list.  It is no more interesting than any other
         part of the item.
   \item This is the third item of the list.
\end{itemize}
You can even display poetry.
\begin{verse}
   There is an environment
    for verse \\             % The \\ command separates lines
   Whose features some poets % within a stanza.
   will curse.

                             % One or more blank lines separate stanzas.

   For instead of making\\
   Them do \emph{all} line breaking, \\
   It allows them to put too many words on a line when they'd rather be
   forced to be terse.
\end{verse}

Mathematical formulas may also be displayed.  A
displayed formula
is
one-line long; multiline
formulas require special formatting instructions.
   \[  \ip{\Gamma}{\psi'} = x'' + y^{2} + z_{i}^{n}\]
Don't start a paragraph with a displayed equation,
nor make one a paragraph by itself.

\end{document}               % End of document.

\section{Bench Huge Section 1031}
% synthetic bench section 1031
% This is a sample LaTeX input file.  (Version of 12 August 2004.)
%
% A '%' character causes TeX to ignore all remaining text on the line,
% and is used for comments like this one.

\documentclass{article}      % Specifies the document class

                             % The preamble begins here.
\title{An Example Document}  % Declares the document's title.
\author{Leslie Lamport}      % Declares the author's name.
\date{January 21, 1994}      % Deleting this command produces today's date.

\newcommand{\ip}[2]{(#1, #2)}
                             % Defines \ip{arg1}{arg2} to mean
                             % (arg1, arg2).

%\newcommand{\ip}[2]{\langle #1 | #2\rangle}
                             % This is an alternative definition of
                             % \ip that is commented out.

\begin{document}             % End of preamble and beginning of text.

\maketitle                   % Produces the title.

This is an example input file.  Comparing it with
the output it generates can show you how to
produce a simple document of your own.

\section{Ordinary Text}      % Produces section heading.  Lower-level
                             % sections are begun with similar
                             % \subsection and \subsubsection commands.

The ends  of words and sentences are marked
  by   spaces. It  doesn't matter how many
spaces    you type; one is as good as 100.  The
end of   a line counts as a space.

One   or more   blank lines denote the  end
of  a paragraph.

Since any number of consecutive spaces are treated
like a single one, the formatting of the input
file makes no difference to
      \LaTeX,                % The \LaTeX command generates the LaTeX logo.
but it makes a difference to you.  When you use
\LaTeX, making your input file as easy to read
as possible will be a great help as you write
your document and when you change it.  This sample
file shows how you can add comments to your own input
file.

Because printing is different from typewriting,
there are a number of things that you have to do
differently when preparing an input file than if
you were just typing the document directly.
Quotation marks like
       ``this''
have to be handled specially, as do quotes within
quotes:
       ``\,`this'            % \, separates the double and single quote.
        is what I just
        wrote, not  `that'\,''.

Dashes come in three sizes: an
       intra-word
dash, a medium dash for number ranges like
       1--2,
and a punctuation
       dash---like
this.

A sentence-ending space should be larger than the
space between words within a sentence.  You
sometimes have to type special commands in
conjunction with punctuation characters to get
this right, as in the following sentence.
       Gnats, gnus, etc.\ all  % `\ ' makes an inter-word space.
       begin with G\@.         % \@ marks end-of-sentence punctuation.
You should check the spaces after periods when
reading your output to make sure you haven't
forgotten any special cases.  Generating an
ellipsis
       \ldots\               % `\ ' is needed after `\ldots' because TeX
                             % ignores spaces after command names like \ldots
                             % made from \ + letters.
                             %
                             % Note how a `%' character causes TeX to ignore
                             % the end of the input line, so these blank lines
                             % do not start a new paragraph.
                             %
with the right spacing around the periods requires
a special command.

\LaTeX\ interprets some common characters as
commands, so you must type special commands to
generate them.  These characters include the
following:
       \$ \& \% \# \{ and \}.

In printing, text is usually emphasized with an
       \emph{italic}
type style.

\begin{em}
   A long segment of text can also be emphasized
   in this way.  Text within such a segment can be
   given \emph{additional} emphasis.
\end{em}

It is sometimes necessary to prevent \LaTeX\ from
breaking a line where it might otherwise do so.
This may be at a space, as between the ``Mr.''\ and
``Jones'' in
       ``Mr.~Jones'',        % ~ produces an unbreakable interword space.
or within a word---especially when the word is a
symbol like
       \mbox{\emph{itemnum}}
that makes little sense when hyphenated across
lines.

Footnotes\footnote{This is an example of a footnote.}
pose no problem.

\LaTeX\ is good at typesetting mathematical formulas
like
       \( x-3y + z = 7 \)
or
       \( a_{1} > x^{2n} + y^{2n} > x' \)
or
       \( \ip{A}{B} = \sum_{i} a_{i} b_{i} \).
The spaces you type in a formula are
ignored.  Remember that a letter like
       $x$                   % $ ... $  and  \( ... \)  are equivalent
is a formula when it denotes a mathematical
symbol, and it should be typed as one.

\section{Displayed Text}

Text is displayed by indenting it from the left
margin.  Quotations are commonly displayed.  There
are short quotations
\begin{quote}
   This is a short quotation.  It consists of a
   single paragraph of text.  See how it is formatted.
\end{quote}
and longer ones.
\begin{quotation}
   This is a longer quotation.  It consists of two
   paragraphs of text, neither of which are
   particularly interesting.

   This is the second paragraph of the quotation.  It
   is just as dull as the first paragraph.
\end{quotation}
Another frequently-displayed structure is a list.
The following is an example of an \emph{itemized}
list.
\begin{itemize}
   \item This is the first item of an itemized list.
         Each item in the list is marked with a ``tick''.
         You don't have to worry about what kind of tick
         mark is used.

   \item This is the second item of the list.  It
         contains another list nested inside it.  The inner
         list is an \emph{enumerated} list.
         \begin{enumerate}
            \item This is the first item of an enumerated
                  list that is nested within the itemized list.

            \item This is the second item of the inner list.
                  \LaTeX\ allows you to nest lists deeper than
                  you really should.
         \end{enumerate}
         This is the rest of the second item of the outer
         list.  It is no more interesting than any other
         part of the item.
   \item This is the third item of the list.
\end{itemize}
You can even display poetry.
\begin{verse}
   There is an environment
    for verse \\             % The \\ command separates lines
   Whose features some poets % within a stanza.
   will curse.

                             % One or more blank lines separate stanzas.

   For instead of making\\
   Them do \emph{all} line breaking, \\
   It allows them to put too many words on a line when they'd rather be
   forced to be terse.
\end{verse}

Mathematical formulas may also be displayed.  A
displayed formula
is
one-line long; multiline
formulas require special formatting instructions.
   \[  \ip{\Gamma}{\psi'} = x'' + y^{2} + z_{i}^{n}\]
Don't start a paragraph with a displayed equation,
nor make one a paragraph by itself.

\end{document}               % End of document.

\section{Bench Huge Section 1032}
% synthetic bench section 1032
% This is a sample LaTeX input file.  (Version of 12 August 2004.)
%
% A '%' character causes TeX to ignore all remaining text on the line,
% and is used for comments like this one.

\documentclass{article}      % Specifies the document class

                             % The preamble begins here.
\title{An Example Document}  % Declares the document's title.
\author{Leslie Lamport}      % Declares the author's name.
\date{January 21, 1994}      % Deleting this command produces today's date.

\newcommand{\ip}[2]{(#1, #2)}
                             % Defines \ip{arg1}{arg2} to mean
                             % (arg1, arg2).

%\newcommand{\ip}[2]{\langle #1 | #2\rangle}
                             % This is an alternative definition of
                             % \ip that is commented out.

\begin{document}             % End of preamble and beginning of text.

\maketitle                   % Produces the title.

This is an example input file.  Comparing it with
the output it generates can show you how to
produce a simple document of your own.

\section{Ordinary Text}      % Produces section heading.  Lower-level
                             % sections are begun with similar
                             % \subsection and \subsubsection commands.

The ends  of words and sentences are marked
  by   spaces. It  doesn't matter how many
spaces    you type; one is as good as 100.  The
end of   a line counts as a space.

One   or more   blank lines denote the  end
of  a paragraph.

Since any number of consecutive spaces are treated
like a single one, the formatting of the input
file makes no difference to
      \LaTeX,                % The \LaTeX command generates the LaTeX logo.
but it makes a difference to you.  When you use
\LaTeX, making your input file as easy to read
as possible will be a great help as you write
your document and when you change it.  This sample
file shows how you can add comments to your own input
file.

Because printing is different from typewriting,
there are a number of things that you have to do
differently when preparing an input file than if
you were just typing the document directly.
Quotation marks like
       ``this''
have to be handled specially, as do quotes within
quotes:
       ``\,`this'            % \, separates the double and single quote.
        is what I just
        wrote, not  `that'\,''.

Dashes come in three sizes: an
       intra-word
dash, a medium dash for number ranges like
       1--2,
and a punctuation
       dash---like
this.

A sentence-ending space should be larger than the
space between words within a sentence.  You
sometimes have to type special commands in
conjunction with punctuation characters to get
this right, as in the following sentence.
       Gnats, gnus, etc.\ all  % `\ ' makes an inter-word space.
       begin with G\@.         % \@ marks end-of-sentence punctuation.
You should check the spaces after periods when
reading your output to make sure you haven't
forgotten any special cases.  Generating an
ellipsis
       \ldots\               % `\ ' is needed after `\ldots' because TeX
                             % ignores spaces after command names like \ldots
                             % made from \ + letters.
                             %
                             % Note how a `%' character causes TeX to ignore
                             % the end of the input line, so these blank lines
                             % do not start a new paragraph.
                             %
with the right spacing around the periods requires
a special command.

\LaTeX\ interprets some common characters as
commands, so you must type special commands to
generate them.  These characters include the
following:
       \$ \& \% \# \{ and \}.

In printing, text is usually emphasized with an
       \emph{italic}
type style.

\begin{em}
   A long segment of text can also be emphasized
   in this way.  Text within such a segment can be
   given \emph{additional} emphasis.
\end{em}

It is sometimes necessary to prevent \LaTeX\ from
breaking a line where it might otherwise do so.
This may be at a space, as between the ``Mr.''\ and
``Jones'' in
       ``Mr.~Jones'',        % ~ produces an unbreakable interword space.
or within a word---especially when the word is a
symbol like
       \mbox{\emph{itemnum}}
that makes little sense when hyphenated across
lines.

Footnotes\footnote{This is an example of a footnote.}
pose no problem.

\LaTeX\ is good at typesetting mathematical formulas
like
       \( x-3y + z = 7 \)
or
       \( a_{1} > x^{2n} + y^{2n} > x' \)
or
       \( \ip{A}{B} = \sum_{i} a_{i} b_{i} \).
The spaces you type in a formula are
ignored.  Remember that a letter like
       $x$                   % $ ... $  and  \( ... \)  are equivalent
is a formula when it denotes a mathematical
symbol, and it should be typed as one.

\section{Displayed Text}

Text is displayed by indenting it from the left
margin.  Quotations are commonly displayed.  There
are short quotations
\begin{quote}
   This is a short quotation.  It consists of a
   single paragraph of text.  See how it is formatted.
\end{quote}
and longer ones.
\begin{quotation}
   This is a longer quotation.  It consists of two
   paragraphs of text, neither of which are
   particularly interesting.

   This is the second paragraph of the quotation.  It
   is just as dull as the first paragraph.
\end{quotation}
Another frequently-displayed structure is a list.
The following is an example of an \emph{itemized}
list.
\begin{itemize}
   \item This is the first item of an itemized list.
         Each item in the list is marked with a ``tick''.
         You don't have to worry about what kind of tick
         mark is used.

   \item This is the second item of the list.  It
         contains another list nested inside it.  The inner
         list is an \emph{enumerated} list.
         \begin{enumerate}
            \item This is the first item of an enumerated
                  list that is nested within the itemized list.

            \item This is the second item of the inner list.
                  \LaTeX\ allows you to nest lists deeper than
                  you really should.
         \end{enumerate}
         This is the rest of the second item of the outer
         list.  It is no more interesting than any other
         part of the item.
   \item This is the third item of the list.
\end{itemize}
You can even display poetry.
\begin{verse}
   There is an environment
    for verse \\             % The \\ command separates lines
   Whose features some poets % within a stanza.
   will curse.

                             % One or more blank lines separate stanzas.

   For instead of making\\
   Them do \emph{all} line breaking, \\
   It allows them to put too many words on a line when they'd rather be
   forced to be terse.
\end{verse}

Mathematical formulas may also be displayed.  A
displayed formula
is
one-line long; multiline
formulas require special formatting instructions.
   \[  \ip{\Gamma}{\psi'} = x'' + y^{2} + z_{i}^{n}\]
Don't start a paragraph with a displayed equation,
nor make one a paragraph by itself.

\end{document}               % End of document.

\section{Bench Huge Section 1033}
% synthetic bench section 1033
% This is a sample LaTeX input file.  (Version of 12 August 2004.)
%
% A '%' character causes TeX to ignore all remaining text on the line,
% and is used for comments like this one.

\documentclass{article}      % Specifies the document class

                             % The preamble begins here.
\title{An Example Document}  % Declares the document's title.
\author{Leslie Lamport}      % Declares the author's name.
\date{January 21, 1994}      % Deleting this command produces today's date.

\newcommand{\ip}[2]{(#1, #2)}
                             % Defines \ip{arg1}{arg2} to mean
                             % (arg1, arg2).

%\newcommand{\ip}[2]{\langle #1 | #2\rangle}
                             % This is an alternative definition of
                             % \ip that is commented out.

\begin{document}             % End of preamble and beginning of text.

\maketitle                   % Produces the title.

This is an example input file.  Comparing it with
the output it generates can show you how to
produce a simple document of your own.

\section{Ordinary Text}      % Produces section heading.  Lower-level
                             % sections are begun with similar
                             % \subsection and \subsubsection commands.

The ends  of words and sentences are marked
  by   spaces. It  doesn't matter how many
spaces    you type; one is as good as 100.  The
end of   a line counts as a space.

One   or more   blank lines denote the  end
of  a paragraph.

Since any number of consecutive spaces are treated
like a single one, the formatting of the input
file makes no difference to
      \LaTeX,                % The \LaTeX command generates the LaTeX logo.
but it makes a difference to you.  When you use
\LaTeX, making your input file as easy to read
as possible will be a great help as you write
your document and when you change it.  This sample
file shows how you can add comments to your own input
file.

Because printing is different from typewriting,
there are a number of things that you have to do
differently when preparing an input file than if
you were just typing the document directly.
Quotation marks like
       ``this''
have to be handled specially, as do quotes within
quotes:
       ``\,`this'            % \, separates the double and single quote.
        is what I just
        wrote, not  `that'\,''.

Dashes come in three sizes: an
       intra-word
dash, a medium dash for number ranges like
       1--2,
and a punctuation
       dash---like
this.

A sentence-ending space should be larger than the
space between words within a sentence.  You
sometimes have to type special commands in
conjunction with punctuation characters to get
this right, as in the following sentence.
       Gnats, gnus, etc.\ all  % `\ ' makes an inter-word space.
       begin with G\@.         % \@ marks end-of-sentence punctuation.
You should check the spaces after periods when
reading your output to make sure you haven't
forgotten any special cases.  Generating an
ellipsis
       \ldots\               % `\ ' is needed after `\ldots' because TeX
                             % ignores spaces after command names like \ldots
                             % made from \ + letters.
                             %
                             % Note how a `%' character causes TeX to ignore
                             % the end of the input line, so these blank lines
                             % do not start a new paragraph.
                             %
with the right spacing around the periods requires
a special command.

\LaTeX\ interprets some common characters as
commands, so you must type special commands to
generate them.  These characters include the
following:
       \$ \& \% \# \{ and \}.

In printing, text is usually emphasized with an
       \emph{italic}
type style.

\begin{em}
   A long segment of text can also be emphasized
   in this way.  Text within such a segment can be
   given \emph{additional} emphasis.
\end{em}

It is sometimes necessary to prevent \LaTeX\ from
breaking a line where it might otherwise do so.
This may be at a space, as between the ``Mr.''\ and
``Jones'' in
       ``Mr.~Jones'',        % ~ produces an unbreakable interword space.
or within a word---especially when the word is a
symbol like
       \mbox{\emph{itemnum}}
that makes little sense when hyphenated across
lines.

Footnotes\footnote{This is an example of a footnote.}
pose no problem.

\LaTeX\ is good at typesetting mathematical formulas
like
       \( x-3y + z = 7 \)
or
       \( a_{1} > x^{2n} + y^{2n} > x' \)
or
       \( \ip{A}{B} = \sum_{i} a_{i} b_{i} \).
The spaces you type in a formula are
ignored.  Remember that a letter like
       $x$                   % $ ... $  and  \( ... \)  are equivalent
is a formula when it denotes a mathematical
symbol, and it should be typed as one.

\section{Displayed Text}

Text is displayed by indenting it from the left
margin.  Quotations are commonly displayed.  There
are short quotations
\begin{quote}
   This is a short quotation.  It consists of a
   single paragraph of text.  See how it is formatted.
\end{quote}
and longer ones.
\begin{quotation}
   This is a longer quotation.  It consists of two
   paragraphs of text, neither of which are
   particularly interesting.

   This is the second paragraph of the quotation.  It
   is just as dull as the first paragraph.
\end{quotation}
Another frequently-displayed structure is a list.
The following is an example of an \emph{itemized}
list.
\begin{itemize}
   \item This is the first item of an itemized list.
         Each item in the list is marked with a ``tick''.
         You don't have to worry about what kind of tick
         mark is used.

   \item This is the second item of the list.  It
         contains another list nested inside it.  The inner
         list is an \emph{enumerated} list.
         \begin{enumerate}
            \item This is the first item of an enumerated
                  list that is nested within the itemized list.

            \item This is the second item of the inner list.
                  \LaTeX\ allows you to nest lists deeper than
                  you really should.
         \end{enumerate}
         This is the rest of the second item of the outer
         list.  It is no more interesting than any other
         part of the item.
   \item This is the third item of the list.
\end{itemize}
You can even display poetry.
\begin{verse}
   There is an environment
    for verse \\             % The \\ command separates lines
   Whose features some poets % within a stanza.
   will curse.

                             % One or more blank lines separate stanzas.

   For instead of making\\
   Them do \emph{all} line breaking, \\
   It allows them to put too many words on a line when they'd rather be
   forced to be terse.
\end{verse}

Mathematical formulas may also be displayed.  A
displayed formula
is
one-line long; multiline
formulas require special formatting instructions.
   \[  \ip{\Gamma}{\psi'} = x'' + y^{2} + z_{i}^{n}\]
Don't start a paragraph with a displayed equation,
nor make one a paragraph by itself.

\end{document}               % End of document.

\section{Bench Huge Section 1034}
% synthetic bench section 1034
% This is a sample LaTeX input file.  (Version of 12 August 2004.)
%
% A '%' character causes TeX to ignore all remaining text on the line,
% and is used for comments like this one.

\documentclass{article}      % Specifies the document class

                             % The preamble begins here.
\title{An Example Document}  % Declares the document's title.
\author{Leslie Lamport}      % Declares the author's name.
\date{January 21, 1994}      % Deleting this command produces today's date.

\newcommand{\ip}[2]{(#1, #2)}
                             % Defines \ip{arg1}{arg2} to mean
                             % (arg1, arg2).

%\newcommand{\ip}[2]{\langle #1 | #2\rangle}
                             % This is an alternative definition of
                             % \ip that is commented out.

\begin{document}             % End of preamble and beginning of text.

\maketitle                   % Produces the title.

This is an example input file.  Comparing it with
the output it generates can show you how to
produce a simple document of your own.

\section{Ordinary Text}      % Produces section heading.  Lower-level
                             % sections are begun with similar
                             % \subsection and \subsubsection commands.

The ends  of words and sentences are marked
  by   spaces. It  doesn't matter how many
spaces    you type; one is as good as 100.  The
end of   a line counts as a space.

One   or more   blank lines denote the  end
of  a paragraph.

Since any number of consecutive spaces are treated
like a single one, the formatting of the input
file makes no difference to
      \LaTeX,                % The \LaTeX command generates the LaTeX logo.
but it makes a difference to you.  When you use
\LaTeX, making your input file as easy to read
as possible will be a great help as you write
your document and when you change it.  This sample
file shows how you can add comments to your own input
file.

Because printing is different from typewriting,
there are a number of things that you have to do
differently when preparing an input file than if
you were just typing the document directly.
Quotation marks like
       ``this''
have to be handled specially, as do quotes within
quotes:
       ``\,`this'            % \, separates the double and single quote.
        is what I just
        wrote, not  `that'\,''.

Dashes come in three sizes: an
       intra-word
dash, a medium dash for number ranges like
       1--2,
and a punctuation
       dash---like
this.

A sentence-ending space should be larger than the
space between words within a sentence.  You
sometimes have to type special commands in
conjunction with punctuation characters to get
this right, as in the following sentence.
       Gnats, gnus, etc.\ all  % `\ ' makes an inter-word space.
       begin with G\@.         % \@ marks end-of-sentence punctuation.
You should check the spaces after periods when
reading your output to make sure you haven't
forgotten any special cases.  Generating an
ellipsis
       \ldots\               % `\ ' is needed after `\ldots' because TeX
                             % ignores spaces after command names like \ldots
                             % made from \ + letters.
                             %
                             % Note how a `%' character causes TeX to ignore
                             % the end of the input line, so these blank lines
                             % do not start a new paragraph.
                             %
with the right spacing around the periods requires
a special command.

\LaTeX\ interprets some common characters as
commands, so you must type special commands to
generate them.  These characters include the
following:
       \$ \& \% \# \{ and \}.

In printing, text is usually emphasized with an
       \emph{italic}
type style.

\begin{em}
   A long segment of text can also be emphasized
   in this way.  Text within such a segment can be
   given \emph{additional} emphasis.
\end{em}

It is sometimes necessary to prevent \LaTeX\ from
breaking a line where it might otherwise do so.
This may be at a space, as between the ``Mr.''\ and
``Jones'' in
       ``Mr.~Jones'',        % ~ produces an unbreakable interword space.
or within a word---especially when the word is a
symbol like
       \mbox{\emph{itemnum}}
that makes little sense when hyphenated across
lines.

Footnotes\footnote{This is an example of a footnote.}
pose no problem.

\LaTeX\ is good at typesetting mathematical formulas
like
       \( x-3y + z = 7 \)
or
       \( a_{1} > x^{2n} + y^{2n} > x' \)
or
       \( \ip{A}{B} = \sum_{i} a_{i} b_{i} \).
The spaces you type in a formula are
ignored.  Remember that a letter like
       $x$                   % $ ... $  and  \( ... \)  are equivalent
is a formula when it denotes a mathematical
symbol, and it should be typed as one.

\section{Displayed Text}

Text is displayed by indenting it from the left
margin.  Quotations are commonly displayed.  There
are short quotations
\begin{quote}
   This is a short quotation.  It consists of a
   single paragraph of text.  See how it is formatted.
\end{quote}
and longer ones.
\begin{quotation}
   This is a longer quotation.  It consists of two
   paragraphs of text, neither of which are
   particularly interesting.

   This is the second paragraph of the quotation.  It
   is just as dull as the first paragraph.
\end{quotation}
Another frequently-displayed structure is a list.
The following is an example of an \emph{itemized}
list.
\begin{itemize}
   \item This is the first item of an itemized list.
         Each item in the list is marked with a ``tick''.
         You don't have to worry about what kind of tick
         mark is used.

   \item This is the second item of the list.  It
         contains another list nested inside it.  The inner
         list is an \emph{enumerated} list.
         \begin{enumerate}
            \item This is the first item of an enumerated
                  list that is nested within the itemized list.

            \item This is the second item of the inner list.
                  \LaTeX\ allows you to nest lists deeper than
                  you really should.
         \end{enumerate}
         This is the rest of the second item of the outer
         list.  It is no more interesting than any other
         part of the item.
   \item This is the third item of the list.
\end{itemize}
You can even display poetry.
\begin{verse}
   There is an environment
    for verse \\             % The \\ command separates lines
   Whose features some poets % within a stanza.
   will curse.

                             % One or more blank lines separate stanzas.

   For instead of making\\
   Them do \emph{all} line breaking, \\
   It allows them to put too many words on a line when they'd rather be
   forced to be terse.
\end{verse}

Mathematical formulas may also be displayed.  A
displayed formula
is
one-line long; multiline
formulas require special formatting instructions.
   \[  \ip{\Gamma}{\psi'} = x'' + y^{2} + z_{i}^{n}\]
Don't start a paragraph with a displayed equation,
nor make one a paragraph by itself.

\end{document}               % End of document.

\section{Bench Huge Section 1035}
% synthetic bench section 1035
% This is a sample LaTeX input file.  (Version of 12 August 2004.)
%
% A '%' character causes TeX to ignore all remaining text on the line,
% and is used for comments like this one.

\documentclass{article}      % Specifies the document class

                             % The preamble begins here.
\title{An Example Document}  % Declares the document's title.
\author{Leslie Lamport}      % Declares the author's name.
\date{January 21, 1994}      % Deleting this command produces today's date.

\newcommand{\ip}[2]{(#1, #2)}
                             % Defines \ip{arg1}{arg2} to mean
                             % (arg1, arg2).

%\newcommand{\ip}[2]{\langle #1 | #2\rangle}
                             % This is an alternative definition of
                             % \ip that is commented out.

\begin{document}             % End of preamble and beginning of text.

\maketitle                   % Produces the title.

This is an example input file.  Comparing it with
the output it generates can show you how to
produce a simple document of your own.

\section{Ordinary Text}      % Produces section heading.  Lower-level
                             % sections are begun with similar
                             % \subsection and \subsubsection commands.

The ends  of words and sentences are marked
  by   spaces. It  doesn't matter how many
spaces    you type; one is as good as 100.  The
end of   a line counts as a space.

One   or more   blank lines denote the  end
of  a paragraph.

Since any number of consecutive spaces are treated
like a single one, the formatting of the input
file makes no difference to
      \LaTeX,                % The \LaTeX command generates the LaTeX logo.
but it makes a difference to you.  When you use
\LaTeX, making your input file as easy to read
as possible will be a great help as you write
your document and when you change it.  This sample
file shows how you can add comments to your own input
file.

Because printing is different from typewriting,
there are a number of things that you have to do
differently when preparing an input file than if
you were just typing the document directly.
Quotation marks like
       ``this''
have to be handled specially, as do quotes within
quotes:
       ``\,`this'            % \, separates the double and single quote.
        is what I just
        wrote, not  `that'\,''.

Dashes come in three sizes: an
       intra-word
dash, a medium dash for number ranges like
       1--2,
and a punctuation
       dash---like
this.

A sentence-ending space should be larger than the
space between words within a sentence.  You
sometimes have to type special commands in
conjunction with punctuation characters to get
this right, as in the following sentence.
       Gnats, gnus, etc.\ all  % `\ ' makes an inter-word space.
       begin with G\@.         % \@ marks end-of-sentence punctuation.
You should check the spaces after periods when
reading your output to make sure you haven't
forgotten any special cases.  Generating an
ellipsis
       \ldots\               % `\ ' is needed after `\ldots' because TeX
                             % ignores spaces after command names like \ldots
                             % made from \ + letters.
                             %
                             % Note how a `%' character causes TeX to ignore
                             % the end of the input line, so these blank lines
                             % do not start a new paragraph.
                             %
with the right spacing around the periods requires
a special command.

\LaTeX\ interprets some common characters as
commands, so you must type special commands to
generate them.  These characters include the
following:
       \$ \& \% \# \{ and \}.

In printing, text is usually emphasized with an
       \emph{italic}
type style.

\begin{em}
   A long segment of text can also be emphasized
   in this way.  Text within such a segment can be
   given \emph{additional} emphasis.
\end{em}

It is sometimes necessary to prevent \LaTeX\ from
breaking a line where it might otherwise do so.
This may be at a space, as between the ``Mr.''\ and
``Jones'' in
       ``Mr.~Jones'',        % ~ produces an unbreakable interword space.
or within a word---especially when the word is a
symbol like
       \mbox{\emph{itemnum}}
that makes little sense when hyphenated across
lines.

Footnotes\footnote{This is an example of a footnote.}
pose no problem.

\LaTeX\ is good at typesetting mathematical formulas
like
       \( x-3y + z = 7 \)
or
       \( a_{1} > x^{2n} + y^{2n} > x' \)
or
       \( \ip{A}{B} = \sum_{i} a_{i} b_{i} \).
The spaces you type in a formula are
ignored.  Remember that a letter like
       $x$                   % $ ... $  and  \( ... \)  are equivalent
is a formula when it denotes a mathematical
symbol, and it should be typed as one.

\section{Displayed Text}

Text is displayed by indenting it from the left
margin.  Quotations are commonly displayed.  There
are short quotations
\begin{quote}
   This is a short quotation.  It consists of a
   single paragraph of text.  See how it is formatted.
\end{quote}
and longer ones.
\begin{quotation}
   This is a longer quotation.  It consists of two
   paragraphs of text, neither of which are
   particularly interesting.

   This is the second paragraph of the quotation.  It
   is just as dull as the first paragraph.
\end{quotation}
Another frequently-displayed structure is a list.
The following is an example of an \emph{itemized}
list.
\begin{itemize}
   \item This is the first item of an itemized list.
         Each item in the list is marked with a ``tick''.
         You don't have to worry about what kind of tick
         mark is used.

   \item This is the second item of the list.  It
         contains another list nested inside it.  The inner
         list is an \emph{enumerated} list.
         \begin{enumerate}
            \item This is the first item of an enumerated
                  list that is nested within the itemized list.

            \item This is the second item of the inner list.
                  \LaTeX\ allows you to nest lists deeper than
                  you really should.
         \end{enumerate}
         This is the rest of the second item of the outer
         list.  It is no more interesting than any other
         part of the item.
   \item This is the third item of the list.
\end{itemize}
You can even display poetry.
\begin{verse}
   There is an environment
    for verse \\             % The \\ command separates lines
   Whose features some poets % within a stanza.
   will curse.

                             % One or more blank lines separate stanzas.

   For instead of making\\
   Them do \emph{all} line breaking, \\
   It allows them to put too many words on a line when they'd rather be
   forced to be terse.
\end{verse}

Mathematical formulas may also be displayed.  A
displayed formula
is
one-line long; multiline
formulas require special formatting instructions.
   \[  \ip{\Gamma}{\psi'} = x'' + y^{2} + z_{i}^{n}\]
Don't start a paragraph with a displayed equation,
nor make one a paragraph by itself.

\end{document}               % End of document.

\section{Bench Huge Section 1036}
% synthetic bench section 1036
% This is a sample LaTeX input file.  (Version of 12 August 2004.)
%
% A '%' character causes TeX to ignore all remaining text on the line,
% and is used for comments like this one.

\documentclass{article}      % Specifies the document class

                             % The preamble begins here.
\title{An Example Document}  % Declares the document's title.
\author{Leslie Lamport}      % Declares the author's name.
\date{January 21, 1994}      % Deleting this command produces today's date.

\newcommand{\ip}[2]{(#1, #2)}
                             % Defines \ip{arg1}{arg2} to mean
                             % (arg1, arg2).

%\newcommand{\ip}[2]{\langle #1 | #2\rangle}
                             % This is an alternative definition of
                             % \ip that is commented out.

\begin{document}             % End of preamble and beginning of text.

\maketitle                   % Produces the title.

This is an example input file.  Comparing it with
the output it generates can show you how to
produce a simple document of your own.

\section{Ordinary Text}      % Produces section heading.  Lower-level
                             % sections are begun with similar
                             % \subsection and \subsubsection commands.

The ends  of words and sentences are marked
  by   spaces. It  doesn't matter how many
spaces    you type; one is as good as 100.  The
end of   a line counts as a space.

One   or more   blank lines denote the  end
of  a paragraph.

Since any number of consecutive spaces are treated
like a single one, the formatting of the input
file makes no difference to
      \LaTeX,                % The \LaTeX command generates the LaTeX logo.
but it makes a difference to you.  When you use
\LaTeX, making your input file as easy to read
as possible will be a great help as you write
your document and when you change it.  This sample
file shows how you can add comments to your own input
file.

Because printing is different from typewriting,
there are a number of things that you have to do
differently when preparing an input file than if
you were just typing the document directly.
Quotation marks like
       ``this''
have to be handled specially, as do quotes within
quotes:
       ``\,`this'            % \, separates the double and single quote.
        is what I just
        wrote, not  `that'\,''.

Dashes come in three sizes: an
       intra-word
dash, a medium dash for number ranges like
       1--2,
and a punctuation
       dash---like
this.

A sentence-ending space should be larger than the
space between words within a sentence.  You
sometimes have to type special commands in
conjunction with punctuation characters to get
this right, as in the following sentence.
       Gnats, gnus, etc.\ all  % `\ ' makes an inter-word space.
       begin with G\@.         % \@ marks end-of-sentence punctuation.
You should check the spaces after periods when
reading your output to make sure you haven't
forgotten any special cases.  Generating an
ellipsis
       \ldots\               % `\ ' is needed after `\ldots' because TeX
                             % ignores spaces after command names like \ldots
                             % made from \ + letters.
                             %
                             % Note how a `%' character causes TeX to ignore
                             % the end of the input line, so these blank lines
                             % do not start a new paragraph.
                             %
with the right spacing around the periods requires
a special command.

\LaTeX\ interprets some common characters as
commands, so you must type special commands to
generate them.  These characters include the
following:
       \$ \& \% \# \{ and \}.

In printing, text is usually emphasized with an
       \emph{italic}
type style.

\begin{em}
   A long segment of text can also be emphasized
   in this way.  Text within such a segment can be
   given \emph{additional} emphasis.
\end{em}

It is sometimes necessary to prevent \LaTeX\ from
breaking a line where it might otherwise do so.
This may be at a space, as between the ``Mr.''\ and
``Jones'' in
       ``Mr.~Jones'',        % ~ produces an unbreakable interword space.
or within a word---especially when the word is a
symbol like
       \mbox{\emph{itemnum}}
that makes little sense when hyphenated across
lines.

Footnotes\footnote{This is an example of a footnote.}
pose no problem.

\LaTeX\ is good at typesetting mathematical formulas
like
       \( x-3y + z = 7 \)
or
       \( a_{1} > x^{2n} + y^{2n} > x' \)
or
       \( \ip{A}{B} = \sum_{i} a_{i} b_{i} \).
The spaces you type in a formula are
ignored.  Remember that a letter like
       $x$                   % $ ... $  and  \( ... \)  are equivalent
is a formula when it denotes a mathematical
symbol, and it should be typed as one.

\section{Displayed Text}

Text is displayed by indenting it from the left
margin.  Quotations are commonly displayed.  There
are short quotations
\begin{quote}
   This is a short quotation.  It consists of a
   single paragraph of text.  See how it is formatted.
\end{quote}
and longer ones.
\begin{quotation}
   This is a longer quotation.  It consists of two
   paragraphs of text, neither of which are
   particularly interesting.

   This is the second paragraph of the quotation.  It
   is just as dull as the first paragraph.
\end{quotation}
Another frequently-displayed structure is a list.
The following is an example of an \emph{itemized}
list.
\begin{itemize}
   \item This is the first item of an itemized list.
         Each item in the list is marked with a ``tick''.
         You don't have to worry about what kind of tick
         mark is used.

   \item This is the second item of the list.  It
         contains another list nested inside it.  The inner
         list is an \emph{enumerated} list.
         \begin{enumerate}
            \item This is the first item of an enumerated
                  list that is nested within the itemized list.

            \item This is the second item of the inner list.
                  \LaTeX\ allows you to nest lists deeper than
                  you really should.
         \end{enumerate}
         This is the rest of the second item of the outer
         list.  It is no more interesting than any other
         part of the item.
   \item This is the third item of the list.
\end{itemize}
You can even display poetry.
\begin{verse}
   There is an environment
    for verse \\             % The \\ command separates lines
   Whose features some poets % within a stanza.
   will curse.

                             % One or more blank lines separate stanzas.

   For instead of making\\
   Them do \emph{all} line breaking, \\
   It allows them to put too many words on a line when they'd rather be
   forced to be terse.
\end{verse}

Mathematical formulas may also be displayed.  A
displayed formula
is
one-line long; multiline
formulas require special formatting instructions.
   \[  \ip{\Gamma}{\psi'} = x'' + y^{2} + z_{i}^{n}\]
Don't start a paragraph with a displayed equation,
nor make one a paragraph by itself.

\end{document}               % End of document.

\section{Bench Huge Section 1037}
% synthetic bench section 1037
% This is a sample LaTeX input file.  (Version of 12 August 2004.)
%
% A '%' character causes TeX to ignore all remaining text on the line,
% and is used for comments like this one.

\documentclass{article}      % Specifies the document class

                             % The preamble begins here.
\title{An Example Document}  % Declares the document's title.
\author{Leslie Lamport}      % Declares the author's name.
\date{January 21, 1994}      % Deleting this command produces today's date.

\newcommand{\ip}[2]{(#1, #2)}
                             % Defines \ip{arg1}{arg2} to mean
                             % (arg1, arg2).

%\newcommand{\ip}[2]{\langle #1 | #2\rangle}
                             % This is an alternative definition of
                             % \ip that is commented out.

\begin{document}             % End of preamble and beginning of text.

\maketitle                   % Produces the title.

This is an example input file.  Comparing it with
the output it generates can show you how to
produce a simple document of your own.

\section{Ordinary Text}      % Produces section heading.  Lower-level
                             % sections are begun with similar
                             % \subsection and \subsubsection commands.

The ends  of words and sentences are marked
  by   spaces. It  doesn't matter how many
spaces    you type; one is as good as 100.  The
end of   a line counts as a space.

One   or more   blank lines denote the  end
of  a paragraph.

Since any number of consecutive spaces are treated
like a single one, the formatting of the input
file makes no difference to
      \LaTeX,                % The \LaTeX command generates the LaTeX logo.
but it makes a difference to you.  When you use
\LaTeX, making your input file as easy to read
as possible will be a great help as you write
your document and when you change it.  This sample
file shows how you can add comments to your own input
file.

Because printing is different from typewriting,
there are a number of things that you have to do
differently when preparing an input file than if
you were just typing the document directly.
Quotation marks like
       ``this''
have to be handled specially, as do quotes within
quotes:
       ``\,`this'            % \, separates the double and single quote.
        is what I just
        wrote, not  `that'\,''.

Dashes come in three sizes: an
       intra-word
dash, a medium dash for number ranges like
       1--2,
and a punctuation
       dash---like
this.

A sentence-ending space should be larger than the
space between words within a sentence.  You
sometimes have to type special commands in
conjunction with punctuation characters to get
this right, as in the following sentence.
       Gnats, gnus, etc.\ all  % `\ ' makes an inter-word space.
       begin with G\@.         % \@ marks end-of-sentence punctuation.
You should check the spaces after periods when
reading your output to make sure you haven't
forgotten any special cases.  Generating an
ellipsis
       \ldots\               % `\ ' is needed after `\ldots' because TeX
                             % ignores spaces after command names like \ldots
                             % made from \ + letters.
                             %
                             % Note how a `%' character causes TeX to ignore
                             % the end of the input line, so these blank lines
                             % do not start a new paragraph.
                             %
with the right spacing around the periods requires
a special command.

\LaTeX\ interprets some common characters as
commands, so you must type special commands to
generate them.  These characters include the
following:
       \$ \& \% \# \{ and \}.

In printing, text is usually emphasized with an
       \emph{italic}
type style.

\begin{em}
   A long segment of text can also be emphasized
   in this way.  Text within such a segment can be
   given \emph{additional} emphasis.
\end{em}

It is sometimes necessary to prevent \LaTeX\ from
breaking a line where it might otherwise do so.
This may be at a space, as between the ``Mr.''\ and
``Jones'' in
       ``Mr.~Jones'',        % ~ produces an unbreakable interword space.
or within a word---especially when the word is a
symbol like
       \mbox{\emph{itemnum}}
that makes little sense when hyphenated across
lines.

Footnotes\footnote{This is an example of a footnote.}
pose no problem.

\LaTeX\ is good at typesetting mathematical formulas
like
       \( x-3y + z = 7 \)
or
       \( a_{1} > x^{2n} + y^{2n} > x' \)
or
       \( \ip{A}{B} = \sum_{i} a_{i} b_{i} \).
The spaces you type in a formula are
ignored.  Remember that a letter like
       $x$                   % $ ... $  and  \( ... \)  are equivalent
is a formula when it denotes a mathematical
symbol, and it should be typed as one.

\section{Displayed Text}

Text is displayed by indenting it from the left
margin.  Quotations are commonly displayed.  There
are short quotations
\begin{quote}
   This is a short quotation.  It consists of a
   single paragraph of text.  See how it is formatted.
\end{quote}
and longer ones.
\begin{quotation}
   This is a longer quotation.  It consists of two
   paragraphs of text, neither of which are
   particularly interesting.

   This is the second paragraph of the quotation.  It
   is just as dull as the first paragraph.
\end{quotation}
Another frequently-displayed structure is a list.
The following is an example of an \emph{itemized}
list.
\begin{itemize}
   \item This is the first item of an itemized list.
         Each item in the list is marked with a ``tick''.
         You don't have to worry about what kind of tick
         mark is used.

   \item This is the second item of the list.  It
         contains another list nested inside it.  The inner
         list is an \emph{enumerated} list.
         \begin{enumerate}
            \item This is the first item of an enumerated
                  list that is nested within the itemized list.

            \item This is the second item of the inner list.
                  \LaTeX\ allows you to nest lists deeper than
                  you really should.
         \end{enumerate}
         This is the rest of the second item of the outer
         list.  It is no more interesting than any other
         part of the item.
   \item This is the third item of the list.
\end{itemize}
You can even display poetry.
\begin{verse}
   There is an environment
    for verse \\             % The \\ command separates lines
   Whose features some poets % within a stanza.
   will curse.

                             % One or more blank lines separate stanzas.

   For instead of making\\
   Them do \emph{all} line breaking, \\
   It allows them to put too many words on a line when they'd rather be
   forced to be terse.
\end{verse}

Mathematical formulas may also be displayed.  A
displayed formula
is
one-line long; multiline
formulas require special formatting instructions.
   \[  \ip{\Gamma}{\psi'} = x'' + y^{2} + z_{i}^{n}\]
Don't start a paragraph with a displayed equation,
nor make one a paragraph by itself.

\end{document}               % End of document.

\section{Bench Huge Section 1038}
% synthetic bench section 1038
% This is a sample LaTeX input file.  (Version of 12 August 2004.)
%
% A '%' character causes TeX to ignore all remaining text on the line,
% and is used for comments like this one.

\documentclass{article}      % Specifies the document class

                             % The preamble begins here.
\title{An Example Document}  % Declares the document's title.
\author{Leslie Lamport}      % Declares the author's name.
\date{January 21, 1994}      % Deleting this command produces today's date.

\newcommand{\ip}[2]{(#1, #2)}
                             % Defines \ip{arg1}{arg2} to mean
                             % (arg1, arg2).

%\newcommand{\ip}[2]{\langle #1 | #2\rangle}
                             % This is an alternative definition of
                             % \ip that is commented out.

\begin{document}             % End of preamble and beginning of text.

\maketitle                   % Produces the title.

This is an example input file.  Comparing it with
the output it generates can show you how to
produce a simple document of your own.

\section{Ordinary Text}      % Produces section heading.  Lower-level
                             % sections are begun with similar
                             % \subsection and \subsubsection commands.

The ends  of words and sentences are marked
  by   spaces. It  doesn't matter how many
spaces    you type; one is as good as 100.  The
end of   a line counts as a space.

One   or more   blank lines denote the  end
of  a paragraph.

Since any number of consecutive spaces are treated
like a single one, the formatting of the input
file makes no difference to
      \LaTeX,                % The \LaTeX command generates the LaTeX logo.
but it makes a difference to you.  When you use
\LaTeX, making your input file as easy to read
as possible will be a great help as you write
your document and when you change it.  This sample
file shows how you can add comments to your own input
file.

Because printing is different from typewriting,
there are a number of things that you have to do
differently when preparing an input file than if
you were just typing the document directly.
Quotation marks like
       ``this''
have to be handled specially, as do quotes within
quotes:
       ``\,`this'            % \, separates the double and single quote.
        is what I just
        wrote, not  `that'\,''.

Dashes come in three sizes: an
       intra-word
dash, a medium dash for number ranges like
       1--2,
and a punctuation
       dash---like
this.

A sentence-ending space should be larger than the
space between words within a sentence.  You
sometimes have to type special commands in
conjunction with punctuation characters to get
this right, as in the following sentence.
       Gnats, gnus, etc.\ all  % `\ ' makes an inter-word space.
       begin with G\@.         % \@ marks end-of-sentence punctuation.
You should check the spaces after periods when
reading your output to make sure you haven't
forgotten any special cases.  Generating an
ellipsis
       \ldots\               % `\ ' is needed after `\ldots' because TeX
                             % ignores spaces after command names like \ldots
                             % made from \ + letters.
                             %
                             % Note how a `%' character causes TeX to ignore
                             % the end of the input line, so these blank lines
                             % do not start a new paragraph.
                             %
with the right spacing around the periods requires
a special command.

\LaTeX\ interprets some common characters as
commands, so you must type special commands to
generate them.  These characters include the
following:
       \$ \& \% \# \{ and \}.

In printing, text is usually emphasized with an
       \emph{italic}
type style.

\begin{em}
   A long segment of text can also be emphasized
   in this way.  Text within such a segment can be
   given \emph{additional} emphasis.
\end{em}

It is sometimes necessary to prevent \LaTeX\ from
breaking a line where it might otherwise do so.
This may be at a space, as between the ``Mr.''\ and
``Jones'' in
       ``Mr.~Jones'',        % ~ produces an unbreakable interword space.
or within a word---especially when the word is a
symbol like
       \mbox{\emph{itemnum}}
that makes little sense when hyphenated across
lines.

Footnotes\footnote{This is an example of a footnote.}
pose no problem.

\LaTeX\ is good at typesetting mathematical formulas
like
       \( x-3y + z = 7 \)
or
       \( a_{1} > x^{2n} + y^{2n} > x' \)
or
       \( \ip{A}{B} = \sum_{i} a_{i} b_{i} \).
The spaces you type in a formula are
ignored.  Remember that a letter like
       $x$                   % $ ... $  and  \( ... \)  are equivalent
is a formula when it denotes a mathematical
symbol, and it should be typed as one.

\section{Displayed Text}

Text is displayed by indenting it from the left
margin.  Quotations are commonly displayed.  There
are short quotations
\begin{quote}
   This is a short quotation.  It consists of a
   single paragraph of text.  See how it is formatted.
\end{quote}
and longer ones.
\begin{quotation}
   This is a longer quotation.  It consists of two
   paragraphs of text, neither of which are
   particularly interesting.

   This is the second paragraph of the quotation.  It
   is just as dull as the first paragraph.
\end{quotation}
Another frequently-displayed structure is a list.
The following is an example of an \emph{itemized}
list.
\begin{itemize}
   \item This is the first item of an itemized list.
         Each item in the list is marked with a ``tick''.
         You don't have to worry about what kind of tick
         mark is used.

   \item This is the second item of the list.  It
         contains another list nested inside it.  The inner
         list is an \emph{enumerated} list.
         \begin{enumerate}
            \item This is the first item of an enumerated
                  list that is nested within the itemized list.

            \item This is the second item of the inner list.
                  \LaTeX\ allows you to nest lists deeper than
                  you really should.
         \end{enumerate}
         This is the rest of the second item of the outer
         list.  It is no more interesting than any other
         part of the item.
   \item This is the third item of the list.
\end{itemize}
You can even display poetry.
\begin{verse}
   There is an environment
    for verse \\             % The \\ command separates lines
   Whose features some poets % within a stanza.
   will curse.

                             % One or more blank lines separate stanzas.

   For instead of making\\
   Them do \emph{all} line breaking, \\
   It allows them to put too many words on a line when they'd rather be
   forced to be terse.
\end{verse}

Mathematical formulas may also be displayed.  A
displayed formula
is
one-line long; multiline
formulas require special formatting instructions.
   \[  \ip{\Gamma}{\psi'} = x'' + y^{2} + z_{i}^{n}\]
Don't start a paragraph with a displayed equation,
nor make one a paragraph by itself.

\end{document}               % End of document.

\section{Bench Huge Section 1039}
% synthetic bench section 1039
% This is a sample LaTeX input file.  (Version of 12 August 2004.)
%
% A '%' character causes TeX to ignore all remaining text on the line,
% and is used for comments like this one.

\documentclass{article}      % Specifies the document class

                             % The preamble begins here.
\title{An Example Document}  % Declares the document's title.
\author{Leslie Lamport}      % Declares the author's name.
\date{January 21, 1994}      % Deleting this command produces today's date.

\newcommand{\ip}[2]{(#1, #2)}
                             % Defines \ip{arg1}{arg2} to mean
                             % (arg1, arg2).

%\newcommand{\ip}[2]{\langle #1 | #2\rangle}
                             % This is an alternative definition of
                             % \ip that is commented out.

\begin{document}             % End of preamble and beginning of text.

\maketitle                   % Produces the title.

This is an example input file.  Comparing it with
the output it generates can show you how to
produce a simple document of your own.

\section{Ordinary Text}      % Produces section heading.  Lower-level
                             % sections are begun with similar
                             % \subsection and \subsubsection commands.

The ends  of words and sentences are marked
  by   spaces. It  doesn't matter how many
spaces    you type; one is as good as 100.  The
end of   a line counts as a space.

One   or more   blank lines denote the  end
of  a paragraph.

Since any number of consecutive spaces are treated
like a single one, the formatting of the input
file makes no difference to
      \LaTeX,                % The \LaTeX command generates the LaTeX logo.
but it makes a difference to you.  When you use
\LaTeX, making your input file as easy to read
as possible will be a great help as you write
your document and when you change it.  This sample
file shows how you can add comments to your own input
file.

Because printing is different from typewriting,
there are a number of things that you have to do
differently when preparing an input file than if
you were just typing the document directly.
Quotation marks like
       ``this''
have to be handled specially, as do quotes within
quotes:
       ``\,`this'            % \, separates the double and single quote.
        is what I just
        wrote, not  `that'\,''.

Dashes come in three sizes: an
       intra-word
dash, a medium dash for number ranges like
       1--2,
and a punctuation
       dash---like
this.

A sentence-ending space should be larger than the
space between words within a sentence.  You
sometimes have to type special commands in
conjunction with punctuation characters to get
this right, as in the following sentence.
       Gnats, gnus, etc.\ all  % `\ ' makes an inter-word space.
       begin with G\@.         % \@ marks end-of-sentence punctuation.
You should check the spaces after periods when
reading your output to make sure you haven't
forgotten any special cases.  Generating an
ellipsis
       \ldots\               % `\ ' is needed after `\ldots' because TeX
                             % ignores spaces after command names like \ldots
                             % made from \ + letters.
                             %
                             % Note how a `%' character causes TeX to ignore
                             % the end of the input line, so these blank lines
                             % do not start a new paragraph.
                             %
with the right spacing around the periods requires
a special command.

\LaTeX\ interprets some common characters as
commands, so you must type special commands to
generate them.  These characters include the
following:
       \$ \& \% \# \{ and \}.

In printing, text is usually emphasized with an
       \emph{italic}
type style.

\begin{em}
   A long segment of text can also be emphasized
   in this way.  Text within such a segment can be
   given \emph{additional} emphasis.
\end{em}

It is sometimes necessary to prevent \LaTeX\ from
breaking a line where it might otherwise do so.
This may be at a space, as between the ``Mr.''\ and
``Jones'' in
       ``Mr.~Jones'',        % ~ produces an unbreakable interword space.
or within a word---especially when the word is a
symbol like
       \mbox{\emph{itemnum}}
that makes little sense when hyphenated across
lines.

Footnotes\footnote{This is an example of a footnote.}
pose no problem.

\LaTeX\ is good at typesetting mathematical formulas
like
       \( x-3y + z = 7 \)
or
       \( a_{1} > x^{2n} + y^{2n} > x' \)
or
       \( \ip{A}{B} = \sum_{i} a_{i} b_{i} \).
The spaces you type in a formula are
ignored.  Remember that a letter like
       $x$                   % $ ... $  and  \( ... \)  are equivalent
is a formula when it denotes a mathematical
symbol, and it should be typed as one.

\section{Displayed Text}

Text is displayed by indenting it from the left
margin.  Quotations are commonly displayed.  There
are short quotations
\begin{quote}
   This is a short quotation.  It consists of a
   single paragraph of text.  See how it is formatted.
\end{quote}
and longer ones.
\begin{quotation}
   This is a longer quotation.  It consists of two
   paragraphs of text, neither of which are
   particularly interesting.

   This is the second paragraph of the quotation.  It
   is just as dull as the first paragraph.
\end{quotation}
Another frequently-displayed structure is a list.
The following is an example of an \emph{itemized}
list.
\begin{itemize}
   \item This is the first item of an itemized list.
         Each item in the list is marked with a ``tick''.
         You don't have to worry about what kind of tick
         mark is used.

   \item This is the second item of the list.  It
         contains another list nested inside it.  The inner
         list is an \emph{enumerated} list.
         \begin{enumerate}
            \item This is the first item of an enumerated
                  list that is nested within the itemized list.

            \item This is the second item of the inner list.
                  \LaTeX\ allows you to nest lists deeper than
                  you really should.
         \end{enumerate}
         This is the rest of the second item of the outer
         list.  It is no more interesting than any other
         part of the item.
   \item This is the third item of the list.
\end{itemize}
You can even display poetry.
\begin{verse}
   There is an environment
    for verse \\             % The \\ command separates lines
   Whose features some poets % within a stanza.
   will curse.

                             % One or more blank lines separate stanzas.

   For instead of making\\
   Them do \emph{all} line breaking, \\
   It allows them to put too many words on a line when they'd rather be
   forced to be terse.
\end{verse}

Mathematical formulas may also be displayed.  A
displayed formula
is
one-line long; multiline
formulas require special formatting instructions.
   \[  \ip{\Gamma}{\psi'} = x'' + y^{2} + z_{i}^{n}\]
Don't start a paragraph with a displayed equation,
nor make one a paragraph by itself.

\end{document}               % End of document.

\section{Bench Huge Section 1040}
% synthetic bench section 1040
% This is a sample LaTeX input file.  (Version of 12 August 2004.)
%
% A '%' character causes TeX to ignore all remaining text on the line,
% and is used for comments like this one.

\documentclass{article}      % Specifies the document class

                             % The preamble begins here.
\title{An Example Document}  % Declares the document's title.
\author{Leslie Lamport}      % Declares the author's name.
\date{January 21, 1994}      % Deleting this command produces today's date.

\newcommand{\ip}[2]{(#1, #2)}
                             % Defines \ip{arg1}{arg2} to mean
                             % (arg1, arg2).

%\newcommand{\ip}[2]{\langle #1 | #2\rangle}
                             % This is an alternative definition of
                             % \ip that is commented out.

\begin{document}             % End of preamble and beginning of text.

\maketitle                   % Produces the title.

This is an example input file.  Comparing it with
the output it generates can show you how to
produce a simple document of your own.

\section{Ordinary Text}      % Produces section heading.  Lower-level
                             % sections are begun with similar
                             % \subsection and \subsubsection commands.

The ends  of words and sentences are marked
  by   spaces. It  doesn't matter how many
spaces    you type; one is as good as 100.  The
end of   a line counts as a space.

One   or more   blank lines denote the  end
of  a paragraph.

Since any number of consecutive spaces are treated
like a single one, the formatting of the input
file makes no difference to
      \LaTeX,                % The \LaTeX command generates the LaTeX logo.
but it makes a difference to you.  When you use
\LaTeX, making your input file as easy to read
as possible will be a great help as you write
your document and when you change it.  This sample
file shows how you can add comments to your own input
file.

Because printing is different from typewriting,
there are a number of things that you have to do
differently when preparing an input file than if
you were just typing the document directly.
Quotation marks like
       ``this''
have to be handled specially, as do quotes within
quotes:
       ``\,`this'            % \, separates the double and single quote.
        is what I just
        wrote, not  `that'\,''.

Dashes come in three sizes: an
       intra-word
dash, a medium dash for number ranges like
       1--2,
and a punctuation
       dash---like
this.

A sentence-ending space should be larger than the
space between words within a sentence.  You
sometimes have to type special commands in
conjunction with punctuation characters to get
this right, as in the following sentence.
       Gnats, gnus, etc.\ all  % `\ ' makes an inter-word space.
       begin with G\@.         % \@ marks end-of-sentence punctuation.
You should check the spaces after periods when
reading your output to make sure you haven't
forgotten any special cases.  Generating an
ellipsis
       \ldots\               % `\ ' is needed after `\ldots' because TeX
                             % ignores spaces after command names like \ldots
                             % made from \ + letters.
                             %
                             % Note how a `%' character causes TeX to ignore
                             % the end of the input line, so these blank lines
                             % do not start a new paragraph.
                             %
with the right spacing around the periods requires
a special command.

\LaTeX\ interprets some common characters as
commands, so you must type special commands to
generate them.  These characters include the
following:
       \$ \& \% \# \{ and \}.

In printing, text is usually emphasized with an
       \emph{italic}
type style.

\begin{em}
   A long segment of text can also be emphasized
   in this way.  Text within such a segment can be
   given \emph{additional} emphasis.
\end{em}

It is sometimes necessary to prevent \LaTeX\ from
breaking a line where it might otherwise do so.
This may be at a space, as between the ``Mr.''\ and
``Jones'' in
       ``Mr.~Jones'',        % ~ produces an unbreakable interword space.
or within a word---especially when the word is a
symbol like
       \mbox{\emph{itemnum}}
that makes little sense when hyphenated across
lines.

Footnotes\footnote{This is an example of a footnote.}
pose no problem.

\LaTeX\ is good at typesetting mathematical formulas
like
       \( x-3y + z = 7 \)
or
       \( a_{1} > x^{2n} + y^{2n} > x' \)
or
       \( \ip{A}{B} = \sum_{i} a_{i} b_{i} \).
The spaces you type in a formula are
ignored.  Remember that a letter like
       $x$                   % $ ... $  and  \( ... \)  are equivalent
is a formula when it denotes a mathematical
symbol, and it should be typed as one.

\section{Displayed Text}

Text is displayed by indenting it from the left
margin.  Quotations are commonly displayed.  There
are short quotations
\begin{quote}
   This is a short quotation.  It consists of a
   single paragraph of text.  See how it is formatted.
\end{quote}
and longer ones.
\begin{quotation}
   This is a longer quotation.  It consists of two
   paragraphs of text, neither of which are
   particularly interesting.

   This is the second paragraph of the quotation.  It
   is just as dull as the first paragraph.
\end{quotation}
Another frequently-displayed structure is a list.
The following is an example of an \emph{itemized}
list.
\begin{itemize}
   \item This is the first item of an itemized list.
         Each item in the list is marked with a ``tick''.
         You don't have to worry about what kind of tick
         mark is used.

   \item This is the second item of the list.  It
         contains another list nested inside it.  The inner
         list is an \emph{enumerated} list.
         \begin{enumerate}
            \item This is the first item of an enumerated
                  list that is nested within the itemized list.

            \item This is the second item of the inner list.
                  \LaTeX\ allows you to nest lists deeper than
                  you really should.
         \end{enumerate}
         This is the rest of the second item of the outer
         list.  It is no more interesting than any other
         part of the item.
   \item This is the third item of the list.
\end{itemize}
You can even display poetry.
\begin{verse}
   There is an environment
    for verse \\             % The \\ command separates lines
   Whose features some poets % within a stanza.
   will curse.

                             % One or more blank lines separate stanzas.

   For instead of making\\
   Them do \emph{all} line breaking, \\
   It allows them to put too many words on a line when they'd rather be
   forced to be terse.
\end{verse}

Mathematical formulas may also be displayed.  A
displayed formula
is
one-line long; multiline
formulas require special formatting instructions.
   \[  \ip{\Gamma}{\psi'} = x'' + y^{2} + z_{i}^{n}\]
Don't start a paragraph with a displayed equation,
nor make one a paragraph by itself.

\end{document}               % End of document.

\section{Bench Huge Section 1041}
% synthetic bench section 1041
% This is a sample LaTeX input file.  (Version of 12 August 2004.)
%
% A '%' character causes TeX to ignore all remaining text on the line,
% and is used for comments like this one.

\documentclass{article}      % Specifies the document class

                             % The preamble begins here.
\title{An Example Document}  % Declares the document's title.
\author{Leslie Lamport}      % Declares the author's name.
\date{January 21, 1994}      % Deleting this command produces today's date.

\newcommand{\ip}[2]{(#1, #2)}
                             % Defines \ip{arg1}{arg2} to mean
                             % (arg1, arg2).

%\newcommand{\ip}[2]{\langle #1 | #2\rangle}
                             % This is an alternative definition of
                             % \ip that is commented out.

\begin{document}             % End of preamble and beginning of text.

\maketitle                   % Produces the title.

This is an example input file.  Comparing it with
the output it generates can show you how to
produce a simple document of your own.

\section{Ordinary Text}      % Produces section heading.  Lower-level
                             % sections are begun with similar
                             % \subsection and \subsubsection commands.

The ends  of words and sentences are marked
  by   spaces. It  doesn't matter how many
spaces    you type; one is as good as 100.  The
end of   a line counts as a space.

One   or more   blank lines denote the  end
of  a paragraph.

Since any number of consecutive spaces are treated
like a single one, the formatting of the input
file makes no difference to
      \LaTeX,                % The \LaTeX command generates the LaTeX logo.
but it makes a difference to you.  When you use
\LaTeX, making your input file as easy to read
as possible will be a great help as you write
your document and when you change it.  This sample
file shows how you can add comments to your own input
file.

Because printing is different from typewriting,
there are a number of things that you have to do
differently when preparing an input file than if
you were just typing the document directly.
Quotation marks like
       ``this''
have to be handled specially, as do quotes within
quotes:
       ``\,`this'            % \, separates the double and single quote.
        is what I just
        wrote, not  `that'\,''.

Dashes come in three sizes: an
       intra-word
dash, a medium dash for number ranges like
       1--2,
and a punctuation
       dash---like
this.

A sentence-ending space should be larger than the
space between words within a sentence.  You
sometimes have to type special commands in
conjunction with punctuation characters to get
this right, as in the following sentence.
       Gnats, gnus, etc.\ all  % `\ ' makes an inter-word space.
       begin with G\@.         % \@ marks end-of-sentence punctuation.
You should check the spaces after periods when
reading your output to make sure you haven't
forgotten any special cases.  Generating an
ellipsis
       \ldots\               % `\ ' is needed after `\ldots' because TeX
                             % ignores spaces after command names like \ldots
                             % made from \ + letters.
                             %
                             % Note how a `%' character causes TeX to ignore
                             % the end of the input line, so these blank lines
                             % do not start a new paragraph.
                             %
with the right spacing around the periods requires
a special command.

\LaTeX\ interprets some common characters as
commands, so you must type special commands to
generate them.  These characters include the
following:
       \$ \& \% \# \{ and \}.

In printing, text is usually emphasized with an
       \emph{italic}
type style.

\begin{em}
   A long segment of text can also be emphasized
   in this way.  Text within such a segment can be
   given \emph{additional} emphasis.
\end{em}

It is sometimes necessary to prevent \LaTeX\ from
breaking a line where it might otherwise do so.
This may be at a space, as between the ``Mr.''\ and
``Jones'' in
       ``Mr.~Jones'',        % ~ produces an unbreakable interword space.
or within a word---especially when the word is a
symbol like
       \mbox{\emph{itemnum}}
that makes little sense when hyphenated across
lines.

Footnotes\footnote{This is an example of a footnote.}
pose no problem.

\LaTeX\ is good at typesetting mathematical formulas
like
       \( x-3y + z = 7 \)
or
       \( a_{1} > x^{2n} + y^{2n} > x' \)
or
       \( \ip{A}{B} = \sum_{i} a_{i} b_{i} \).
The spaces you type in a formula are
ignored.  Remember that a letter like
       $x$                   % $ ... $  and  \( ... \)  are equivalent
is a formula when it denotes a mathematical
symbol, and it should be typed as one.

\section{Displayed Text}

Text is displayed by indenting it from the left
margin.  Quotations are commonly displayed.  There
are short quotations
\begin{quote}
   This is a short quotation.  It consists of a
   single paragraph of text.  See how it is formatted.
\end{quote}
and longer ones.
\begin{quotation}
   This is a longer quotation.  It consists of two
   paragraphs of text, neither of which are
   particularly interesting.

   This is the second paragraph of the quotation.  It
   is just as dull as the first paragraph.
\end{quotation}
Another frequently-displayed structure is a list.
The following is an example of an \emph{itemized}
list.
\begin{itemize}
   \item This is the first item of an itemized list.
         Each item in the list is marked with a ``tick''.
         You don't have to worry about what kind of tick
         mark is used.

   \item This is the second item of the list.  It
         contains another list nested inside it.  The inner
         list is an \emph{enumerated} list.
         \begin{enumerate}
            \item This is the first item of an enumerated
                  list that is nested within the itemized list.

            \item This is the second item of the inner list.
                  \LaTeX\ allows you to nest lists deeper than
                  you really should.
         \end{enumerate}
         This is the rest of the second item of the outer
         list.  It is no more interesting than any other
         part of the item.
   \item This is the third item of the list.
\end{itemize}
You can even display poetry.
\begin{verse}
   There is an environment
    for verse \\             % The \\ command separates lines
   Whose features some poets % within a stanza.
   will curse.

                             % One or more blank lines separate stanzas.

   For instead of making\\
   Them do \emph{all} line breaking, \\
   It allows them to put too many words on a line when they'd rather be
   forced to be terse.
\end{verse}

Mathematical formulas may also be displayed.  A
displayed formula
is
one-line long; multiline
formulas require special formatting instructions.
   \[  \ip{\Gamma}{\psi'} = x'' + y^{2} + z_{i}^{n}\]
Don't start a paragraph with a displayed equation,
nor make one a paragraph by itself.

\end{document}               % End of document.

\section{Bench Huge Section 1042}
% synthetic bench section 1042
% This is a sample LaTeX input file.  (Version of 12 August 2004.)
%
% A '%' character causes TeX to ignore all remaining text on the line,
% and is used for comments like this one.

\documentclass{article}      % Specifies the document class

                             % The preamble begins here.
\title{An Example Document}  % Declares the document's title.
\author{Leslie Lamport}      % Declares the author's name.
\date{January 21, 1994}      % Deleting this command produces today's date.

\newcommand{\ip}[2]{(#1, #2)}
                             % Defines \ip{arg1}{arg2} to mean
                             % (arg1, arg2).

%\newcommand{\ip}[2]{\langle #1 | #2\rangle}
                             % This is an alternative definition of
                             % \ip that is commented out.

\begin{document}             % End of preamble and beginning of text.

\maketitle                   % Produces the title.

This is an example input file.  Comparing it with
the output it generates can show you how to
produce a simple document of your own.

\section{Ordinary Text}      % Produces section heading.  Lower-level
                             % sections are begun with similar
                             % \subsection and \subsubsection commands.

The ends  of words and sentences are marked
  by   spaces. It  doesn't matter how many
spaces    you type; one is as good as 100.  The
end of   a line counts as a space.

One   or more   blank lines denote the  end
of  a paragraph.

Since any number of consecutive spaces are treated
like a single one, the formatting of the input
file makes no difference to
      \LaTeX,                % The \LaTeX command generates the LaTeX logo.
but it makes a difference to you.  When you use
\LaTeX, making your input file as easy to read
as possible will be a great help as you write
your document and when you change it.  This sample
file shows how you can add comments to your own input
file.

Because printing is different from typewriting,
there are a number of things that you have to do
differently when preparing an input file than if
you were just typing the document directly.
Quotation marks like
       ``this''
have to be handled specially, as do quotes within
quotes:
       ``\,`this'            % \, separates the double and single quote.
        is what I just
        wrote, not  `that'\,''.

Dashes come in three sizes: an
       intra-word
dash, a medium dash for number ranges like
       1--2,
and a punctuation
       dash---like
this.

A sentence-ending space should be larger than the
space between words within a sentence.  You
sometimes have to type special commands in
conjunction with punctuation characters to get
this right, as in the following sentence.
       Gnats, gnus, etc.\ all  % `\ ' makes an inter-word space.
       begin with G\@.         % \@ marks end-of-sentence punctuation.
You should check the spaces after periods when
reading your output to make sure you haven't
forgotten any special cases.  Generating an
ellipsis
       \ldots\               % `\ ' is needed after `\ldots' because TeX
                             % ignores spaces after command names like \ldots
                             % made from \ + letters.
                             %
                             % Note how a `%' character causes TeX to ignore
                             % the end of the input line, so these blank lines
                             % do not start a new paragraph.
                             %
with the right spacing around the periods requires
a special command.

\LaTeX\ interprets some common characters as
commands, so you must type special commands to
generate them.  These characters include the
following:
       \$ \& \% \# \{ and \}.

In printing, text is usually emphasized with an
       \emph{italic}
type style.

\begin{em}
   A long segment of text can also be emphasized
   in this way.  Text within such a segment can be
   given \emph{additional} emphasis.
\end{em}

It is sometimes necessary to prevent \LaTeX\ from
breaking a line where it might otherwise do so.
This may be at a space, as between the ``Mr.''\ and
``Jones'' in
       ``Mr.~Jones'',        % ~ produces an unbreakable interword space.
or within a word---especially when the word is a
symbol like
       \mbox{\emph{itemnum}}
that makes little sense when hyphenated across
lines.

Footnotes\footnote{This is an example of a footnote.}
pose no problem.

\LaTeX\ is good at typesetting mathematical formulas
like
       \( x-3y + z = 7 \)
or
       \( a_{1} > x^{2n} + y^{2n} > x' \)
or
       \( \ip{A}{B} = \sum_{i} a_{i} b_{i} \).
The spaces you type in a formula are
ignored.  Remember that a letter like
       $x$                   % $ ... $  and  \( ... \)  are equivalent
is a formula when it denotes a mathematical
symbol, and it should be typed as one.

\section{Displayed Text}

Text is displayed by indenting it from the left
margin.  Quotations are commonly displayed.  There
are short quotations
\begin{quote}
   This is a short quotation.  It consists of a
   single paragraph of text.  See how it is formatted.
\end{quote}
and longer ones.
\begin{quotation}
   This is a longer quotation.  It consists of two
   paragraphs of text, neither of which are
   particularly interesting.

   This is the second paragraph of the quotation.  It
   is just as dull as the first paragraph.
\end{quotation}
Another frequently-displayed structure is a list.
The following is an example of an \emph{itemized}
list.
\begin{itemize}
   \item This is the first item of an itemized list.
         Each item in the list is marked with a ``tick''.
         You don't have to worry about what kind of tick
         mark is used.

   \item This is the second item of the list.  It
         contains another list nested inside it.  The inner
         list is an \emph{enumerated} list.
         \begin{enumerate}
            \item This is the first item of an enumerated
                  list that is nested within the itemized list.

            \item This is the second item of the inner list.
                  \LaTeX\ allows you to nest lists deeper than
                  you really should.
         \end{enumerate}
         This is the rest of the second item of the outer
         list.  It is no more interesting than any other
         part of the item.
   \item This is the third item of the list.
\end{itemize}
You can even display poetry.
\begin{verse}
   There is an environment
    for verse \\             % The \\ command separates lines
   Whose features some poets % within a stanza.
   will curse.

                             % One or more blank lines separate stanzas.

   For instead of making\\
   Them do \emph{all} line breaking, \\
   It allows them to put too many words on a line when they'd rather be
   forced to be terse.
\end{verse}

Mathematical formulas may also be displayed.  A
displayed formula
is
one-line long; multiline
formulas require special formatting instructions.
   \[  \ip{\Gamma}{\psi'} = x'' + y^{2} + z_{i}^{n}\]
Don't start a paragraph with a displayed equation,
nor make one a paragraph by itself.

\end{document}               % End of document.

\section{Bench Huge Section 1043}
% synthetic bench section 1043
% This is a sample LaTeX input file.  (Version of 12 August 2004.)
%
% A '%' character causes TeX to ignore all remaining text on the line,
% and is used for comments like this one.

\documentclass{article}      % Specifies the document class

                             % The preamble begins here.
\title{An Example Document}  % Declares the document's title.
\author{Leslie Lamport}      % Declares the author's name.
\date{January 21, 1994}      % Deleting this command produces today's date.

\newcommand{\ip}[2]{(#1, #2)}
                             % Defines \ip{arg1}{arg2} to mean
                             % (arg1, arg2).

%\newcommand{\ip}[2]{\langle #1 | #2\rangle}
                             % This is an alternative definition of
                             % \ip that is commented out.

\begin{document}             % End of preamble and beginning of text.

\maketitle                   % Produces the title.

This is an example input file.  Comparing it with
the output it generates can show you how to
produce a simple document of your own.

\section{Ordinary Text}      % Produces section heading.  Lower-level
                             % sections are begun with similar
                             % \subsection and \subsubsection commands.

The ends  of words and sentences are marked
  by   spaces. It  doesn't matter how many
spaces    you type; one is as good as 100.  The
end of   a line counts as a space.

One   or more   blank lines denote the  end
of  a paragraph.

Since any number of consecutive spaces are treated
like a single one, the formatting of the input
file makes no difference to
      \LaTeX,                % The \LaTeX command generates the LaTeX logo.
but it makes a difference to you.  When you use
\LaTeX, making your input file as easy to read
as possible will be a great help as you write
your document and when you change it.  This sample
file shows how you can add comments to your own input
file.

Because printing is different from typewriting,
there are a number of things that you have to do
differently when preparing an input file than if
you were just typing the document directly.
Quotation marks like
       ``this''
have to be handled specially, as do quotes within
quotes:
       ``\,`this'            % \, separates the double and single quote.
        is what I just
        wrote, not  `that'\,''.

Dashes come in three sizes: an
       intra-word
dash, a medium dash for number ranges like
       1--2,
and a punctuation
       dash---like
this.

A sentence-ending space should be larger than the
space between words within a sentence.  You
sometimes have to type special commands in
conjunction with punctuation characters to get
this right, as in the following sentence.
       Gnats, gnus, etc.\ all  % `\ ' makes an inter-word space.
       begin with G\@.         % \@ marks end-of-sentence punctuation.
You should check the spaces after periods when
reading your output to make sure you haven't
forgotten any special cases.  Generating an
ellipsis
       \ldots\               % `\ ' is needed after `\ldots' because TeX
                             % ignores spaces after command names like \ldots
                             % made from \ + letters.
                             %
                             % Note how a `%' character causes TeX to ignore
                             % the end of the input line, so these blank lines
                             % do not start a new paragraph.
                             %
with the right spacing around the periods requires
a special command.

\LaTeX\ interprets some common characters as
commands, so you must type special commands to
generate them.  These characters include the
following:
       \$ \& \% \# \{ and \}.

In printing, text is usually emphasized with an
       \emph{italic}
type style.

\begin{em}
   A long segment of text can also be emphasized
   in this way.  Text within such a segment can be
   given \emph{additional} emphasis.
\end{em}

It is sometimes necessary to prevent \LaTeX\ from
breaking a line where it might otherwise do so.
This may be at a space, as between the ``Mr.''\ and
``Jones'' in
       ``Mr.~Jones'',        % ~ produces an unbreakable interword space.
or within a word---especially when the word is a
symbol like
       \mbox{\emph{itemnum}}
that makes little sense when hyphenated across
lines.

Footnotes\footnote{This is an example of a footnote.}
pose no problem.

\LaTeX\ is good at typesetting mathematical formulas
like
       \( x-3y + z = 7 \)
or
       \( a_{1} > x^{2n} + y^{2n} > x' \)
or
       \( \ip{A}{B} = \sum_{i} a_{i} b_{i} \).
The spaces you type in a formula are
ignored.  Remember that a letter like
       $x$                   % $ ... $  and  \( ... \)  are equivalent
is a formula when it denotes a mathematical
symbol, and it should be typed as one.

\section{Displayed Text}

Text is displayed by indenting it from the left
margin.  Quotations are commonly displayed.  There
are short quotations
\begin{quote}
   This is a short quotation.  It consists of a
   single paragraph of text.  See how it is formatted.
\end{quote}
and longer ones.
\begin{quotation}
   This is a longer quotation.  It consists of two
   paragraphs of text, neither of which are
   particularly interesting.

   This is the second paragraph of the quotation.  It
   is just as dull as the first paragraph.
\end{quotation}
Another frequently-displayed structure is a list.
The following is an example of an \emph{itemized}
list.
\begin{itemize}
   \item This is the first item of an itemized list.
         Each item in the list is marked with a ``tick''.
         You don't have to worry about what kind of tick
         mark is used.

   \item This is the second item of the list.  It
         contains another list nested inside it.  The inner
         list is an \emph{enumerated} list.
         \begin{enumerate}
            \item This is the first item of an enumerated
                  list that is nested within the itemized list.

            \item This is the second item of the inner list.
                  \LaTeX\ allows you to nest lists deeper than
                  you really should.
         \end{enumerate}
         This is the rest of the second item of the outer
         list.  It is no more interesting than any other
         part of the item.
   \item This is the third item of the list.
\end{itemize}
You can even display poetry.
\begin{verse}
   There is an environment
    for verse \\             % The \\ command separates lines
   Whose features some poets % within a stanza.
   will curse.

                             % One or more blank lines separate stanzas.

   For instead of making\\
   Them do \emph{all} line breaking, \\
   It allows them to put too many words on a line when they'd rather be
   forced to be terse.
\end{verse}

Mathematical formulas may also be displayed.  A
displayed formula
is
one-line long; multiline
formulas require special formatting instructions.
   \[  \ip{\Gamma}{\psi'} = x'' + y^{2} + z_{i}^{n}\]
Don't start a paragraph with a displayed equation,
nor make one a paragraph by itself.

\end{document}               % End of document.

\section{Bench Huge Section 1044}
% synthetic bench section 1044
% This is a sample LaTeX input file.  (Version of 12 August 2004.)
%
% A '%' character causes TeX to ignore all remaining text on the line,
% and is used for comments like this one.

\documentclass{article}      % Specifies the document class

                             % The preamble begins here.
\title{An Example Document}  % Declares the document's title.
\author{Leslie Lamport}      % Declares the author's name.
\date{January 21, 1994}      % Deleting this command produces today's date.

\newcommand{\ip}[2]{(#1, #2)}
                             % Defines \ip{arg1}{arg2} to mean
                             % (arg1, arg2).

%\newcommand{\ip}[2]{\langle #1 | #2\rangle}
                             % This is an alternative definition of
                             % \ip that is commented out.

\begin{document}             % End of preamble and beginning of text.

\maketitle                   % Produces the title.

This is an example input file.  Comparing it with
the output it generates can show you how to
produce a simple document of your own.

\section{Ordinary Text}      % Produces section heading.  Lower-level
                             % sections are begun with similar
                             % \subsection and \subsubsection commands.

The ends  of words and sentences are marked
  by   spaces. It  doesn't matter how many
spaces    you type; one is as good as 100.  The
end of   a line counts as a space.

One   or more   blank lines denote the  end
of  a paragraph.

Since any number of consecutive spaces are treated
like a single one, the formatting of the input
file makes no difference to
      \LaTeX,                % The \LaTeX command generates the LaTeX logo.
but it makes a difference to you.  When you use
\LaTeX, making your input file as easy to read
as possible will be a great help as you write
your document and when you change it.  This sample
file shows how you can add comments to your own input
file.

Because printing is different from typewriting,
there are a number of things that you have to do
differently when preparing an input file than if
you were just typing the document directly.
Quotation marks like
       ``this''
have to be handled specially, as do quotes within
quotes:
       ``\,`this'            % \, separates the double and single quote.
        is what I just
        wrote, not  `that'\,''.

Dashes come in three sizes: an
       intra-word
dash, a medium dash for number ranges like
       1--2,
and a punctuation
       dash---like
this.

A sentence-ending space should be larger than the
space between words within a sentence.  You
sometimes have to type special commands in
conjunction with punctuation characters to get
this right, as in the following sentence.
       Gnats, gnus, etc.\ all  % `\ ' makes an inter-word space.
       begin with G\@.         % \@ marks end-of-sentence punctuation.
You should check the spaces after periods when
reading your output to make sure you haven't
forgotten any special cases.  Generating an
ellipsis
       \ldots\               % `\ ' is needed after `\ldots' because TeX
                             % ignores spaces after command names like \ldots
                             % made from \ + letters.
                             %
                             % Note how a `%' character causes TeX to ignore
                             % the end of the input line, so these blank lines
                             % do not start a new paragraph.
                             %
with the right spacing around the periods requires
a special command.

\LaTeX\ interprets some common characters as
commands, so you must type special commands to
generate them.  These characters include the
following:
       \$ \& \% \# \{ and \}.

In printing, text is usually emphasized with an
       \emph{italic}
type style.

\begin{em}
   A long segment of text can also be emphasized
   in this way.  Text within such a segment can be
   given \emph{additional} emphasis.
\end{em}

It is sometimes necessary to prevent \LaTeX\ from
breaking a line where it might otherwise do so.
This may be at a space, as between the ``Mr.''\ and
``Jones'' in
       ``Mr.~Jones'',        % ~ produces an unbreakable interword space.
or within a word---especially when the word is a
symbol like
       \mbox{\emph{itemnum}}
that makes little sense when hyphenated across
lines.

Footnotes\footnote{This is an example of a footnote.}
pose no problem.

\LaTeX\ is good at typesetting mathematical formulas
like
       \( x-3y + z = 7 \)
or
       \( a_{1} > x^{2n} + y^{2n} > x' \)
or
       \( \ip{A}{B} = \sum_{i} a_{i} b_{i} \).
The spaces you type in a formula are
ignored.  Remember that a letter like
       $x$                   % $ ... $  and  \( ... \)  are equivalent
is a formula when it denotes a mathematical
symbol, and it should be typed as one.

\section{Displayed Text}

Text is displayed by indenting it from the left
margin.  Quotations are commonly displayed.  There
are short quotations
\begin{quote}
   This is a short quotation.  It consists of a
   single paragraph of text.  See how it is formatted.
\end{quote}
and longer ones.
\begin{quotation}
   This is a longer quotation.  It consists of two
   paragraphs of text, neither of which are
   particularly interesting.

   This is the second paragraph of the quotation.  It
   is just as dull as the first paragraph.
\end{quotation}
Another frequently-displayed structure is a list.
The following is an example of an \emph{itemized}
list.
\begin{itemize}
   \item This is the first item of an itemized list.
         Each item in the list is marked with a ``tick''.
         You don't have to worry about what kind of tick
         mark is used.

   \item This is the second item of the list.  It
         contains another list nested inside it.  The inner
         list is an \emph{enumerated} list.
         \begin{enumerate}
            \item This is the first item of an enumerated
                  list that is nested within the itemized list.

            \item This is the second item of the inner list.
                  \LaTeX\ allows you to nest lists deeper than
                  you really should.
         \end{enumerate}
         This is the rest of the second item of the outer
         list.  It is no more interesting than any other
         part of the item.
   \item This is the third item of the list.
\end{itemize}
You can even display poetry.
\begin{verse}
   There is an environment
    for verse \\             % The \\ command separates lines
   Whose features some poets % within a stanza.
   will curse.

                             % One or more blank lines separate stanzas.

   For instead of making\\
   Them do \emph{all} line breaking, \\
   It allows them to put too many words on a line when they'd rather be
   forced to be terse.
\end{verse}

Mathematical formulas may also be displayed.  A
displayed formula
is
one-line long; multiline
formulas require special formatting instructions.
   \[  \ip{\Gamma}{\psi'} = x'' + y^{2} + z_{i}^{n}\]
Don't start a paragraph with a displayed equation,
nor make one a paragraph by itself.

\end{document}               % End of document.

\section{Bench Huge Section 1045}
% synthetic bench section 1045
% This is a sample LaTeX input file.  (Version of 12 August 2004.)
%
% A '%' character causes TeX to ignore all remaining text on the line,
% and is used for comments like this one.

\documentclass{article}      % Specifies the document class

                             % The preamble begins here.
\title{An Example Document}  % Declares the document's title.
\author{Leslie Lamport}      % Declares the author's name.
\date{January 21, 1994}      % Deleting this command produces today's date.

\newcommand{\ip}[2]{(#1, #2)}
                             % Defines \ip{arg1}{arg2} to mean
                             % (arg1, arg2).

%\newcommand{\ip}[2]{\langle #1 | #2\rangle}
                             % This is an alternative definition of
                             % \ip that is commented out.

\begin{document}             % End of preamble and beginning of text.

\maketitle                   % Produces the title.

This is an example input file.  Comparing it with
the output it generates can show you how to
produce a simple document of your own.

\section{Ordinary Text}      % Produces section heading.  Lower-level
                             % sections are begun with similar
                             % \subsection and \subsubsection commands.

The ends  of words and sentences are marked
  by   spaces. It  doesn't matter how many
spaces    you type; one is as good as 100.  The
end of   a line counts as a space.

One   or more   blank lines denote the  end
of  a paragraph.

Since any number of consecutive spaces are treated
like a single one, the formatting of the input
file makes no difference to
      \LaTeX,                % The \LaTeX command generates the LaTeX logo.
but it makes a difference to you.  When you use
\LaTeX, making your input file as easy to read
as possible will be a great help as you write
your document and when you change it.  This sample
file shows how you can add comments to your own input
file.

Because printing is different from typewriting,
there are a number of things that you have to do
differently when preparing an input file than if
you were just typing the document directly.
Quotation marks like
       ``this''
have to be handled specially, as do quotes within
quotes:
       ``\,`this'            % \, separates the double and single quote.
        is what I just
        wrote, not  `that'\,''.

Dashes come in three sizes: an
       intra-word
dash, a medium dash for number ranges like
       1--2,
and a punctuation
       dash---like
this.

A sentence-ending space should be larger than the
space between words within a sentence.  You
sometimes have to type special commands in
conjunction with punctuation characters to get
this right, as in the following sentence.
       Gnats, gnus, etc.\ all  % `\ ' makes an inter-word space.
       begin with G\@.         % \@ marks end-of-sentence punctuation.
You should check the spaces after periods when
reading your output to make sure you haven't
forgotten any special cases.  Generating an
ellipsis
       \ldots\               % `\ ' is needed after `\ldots' because TeX
                             % ignores spaces after command names like \ldots
                             % made from \ + letters.
                             %
                             % Note how a `%' character causes TeX to ignore
                             % the end of the input line, so these blank lines
                             % do not start a new paragraph.
                             %
with the right spacing around the periods requires
a special command.

\LaTeX\ interprets some common characters as
commands, so you must type special commands to
generate them.  These characters include the
following:
       \$ \& \% \# \{ and \}.

In printing, text is usually emphasized with an
       \emph{italic}
type style.

\begin{em}
   A long segment of text can also be emphasized
   in this way.  Text within such a segment can be
   given \emph{additional} emphasis.
\end{em}

It is sometimes necessary to prevent \LaTeX\ from
breaking a line where it might otherwise do so.
This may be at a space, as between the ``Mr.''\ and
``Jones'' in
       ``Mr.~Jones'',        % ~ produces an unbreakable interword space.
or within a word---especially when the word is a
symbol like
       \mbox{\emph{itemnum}}
that makes little sense when hyphenated across
lines.

Footnotes\footnote{This is an example of a footnote.}
pose no problem.

\LaTeX\ is good at typesetting mathematical formulas
like
       \( x-3y + z = 7 \)
or
       \( a_{1} > x^{2n} + y^{2n} > x' \)
or
       \( \ip{A}{B} = \sum_{i} a_{i} b_{i} \).
The spaces you type in a formula are
ignored.  Remember that a letter like
       $x$                   % $ ... $  and  \( ... \)  are equivalent
is a formula when it denotes a mathematical
symbol, and it should be typed as one.

\section{Displayed Text}

Text is displayed by indenting it from the left
margin.  Quotations are commonly displayed.  There
are short quotations
\begin{quote}
   This is a short quotation.  It consists of a
   single paragraph of text.  See how it is formatted.
\end{quote}
and longer ones.
\begin{quotation}
   This is a longer quotation.  It consists of two
   paragraphs of text, neither of which are
   particularly interesting.

   This is the second paragraph of the quotation.  It
   is just as dull as the first paragraph.
\end{quotation}
Another frequently-displayed structure is a list.
The following is an example of an \emph{itemized}
list.
\begin{itemize}
   \item This is the first item of an itemized list.
         Each item in the list is marked with a ``tick''.
         You don't have to worry about what kind of tick
         mark is used.

   \item This is the second item of the list.  It
         contains another list nested inside it.  The inner
         list is an \emph{enumerated} list.
         \begin{enumerate}
            \item This is the first item of an enumerated
                  list that is nested within the itemized list.

            \item This is the second item of the inner list.
                  \LaTeX\ allows you to nest lists deeper than
                  you really should.
         \end{enumerate}
         This is the rest of the second item of the outer
         list.  It is no more interesting than any other
         part of the item.
   \item This is the third item of the list.
\end{itemize}
You can even display poetry.
\begin{verse}
   There is an environment
    for verse \\             % The \\ command separates lines
   Whose features some poets % within a stanza.
   will curse.

                             % One or more blank lines separate stanzas.

   For instead of making\\
   Them do \emph{all} line breaking, \\
   It allows them to put too many words on a line when they'd rather be
   forced to be terse.
\end{verse}

Mathematical formulas may also be displayed.  A
displayed formula
is
one-line long; multiline
formulas require special formatting instructions.
   \[  \ip{\Gamma}{\psi'} = x'' + y^{2} + z_{i}^{n}\]
Don't start a paragraph with a displayed equation,
nor make one a paragraph by itself.

\end{document}               % End of document.

\section{Bench Huge Section 1046}
% synthetic bench section 1046
% This is a sample LaTeX input file.  (Version of 12 August 2004.)
%
% A '%' character causes TeX to ignore all remaining text on the line,
% and is used for comments like this one.

\documentclass{article}      % Specifies the document class

                             % The preamble begins here.
\title{An Example Document}  % Declares the document's title.
\author{Leslie Lamport}      % Declares the author's name.
\date{January 21, 1994}      % Deleting this command produces today's date.

\newcommand{\ip}[2]{(#1, #2)}
                             % Defines \ip{arg1}{arg2} to mean
                             % (arg1, arg2).

%\newcommand{\ip}[2]{\langle #1 | #2\rangle}
                             % This is an alternative definition of
                             % \ip that is commented out.

\begin{document}             % End of preamble and beginning of text.

\maketitle                   % Produces the title.

This is an example input file.  Comparing it with
the output it generates can show you how to
produce a simple document of your own.

\section{Ordinary Text}      % Produces section heading.  Lower-level
                             % sections are begun with similar
                             % \subsection and \subsubsection commands.

The ends  of words and sentences are marked
  by   spaces. It  doesn't matter how many
spaces    you type; one is as good as 100.  The
end of   a line counts as a space.

One   or more   blank lines denote the  end
of  a paragraph.

Since any number of consecutive spaces are treated
like a single one, the formatting of the input
file makes no difference to
      \LaTeX,                % The \LaTeX command generates the LaTeX logo.
but it makes a difference to you.  When you use
\LaTeX, making your input file as easy to read
as possible will be a great help as you write
your document and when you change it.  This sample
file shows how you can add comments to your own input
file.

Because printing is different from typewriting,
there are a number of things that you have to do
differently when preparing an input file than if
you were just typing the document directly.
Quotation marks like
       ``this''
have to be handled specially, as do quotes within
quotes:
       ``\,`this'            % \, separates the double and single quote.
        is what I just
        wrote, not  `that'\,''.

Dashes come in three sizes: an
       intra-word
dash, a medium dash for number ranges like
       1--2,
and a punctuation
       dash---like
this.

A sentence-ending space should be larger than the
space between words within a sentence.  You
sometimes have to type special commands in
conjunction with punctuation characters to get
this right, as in the following sentence.
       Gnats, gnus, etc.\ all  % `\ ' makes an inter-word space.
       begin with G\@.         % \@ marks end-of-sentence punctuation.
You should check the spaces after periods when
reading your output to make sure you haven't
forgotten any special cases.  Generating an
ellipsis
       \ldots\               % `\ ' is needed after `\ldots' because TeX
                             % ignores spaces after command names like \ldots
                             % made from \ + letters.
                             %
                             % Note how a `%' character causes TeX to ignore
                             % the end of the input line, so these blank lines
                             % do not start a new paragraph.
                             %
with the right spacing around the periods requires
a special command.

\LaTeX\ interprets some common characters as
commands, so you must type special commands to
generate them.  These characters include the
following:
       \$ \& \% \# \{ and \}.

In printing, text is usually emphasized with an
       \emph{italic}
type style.

\begin{em}
   A long segment of text can also be emphasized
   in this way.  Text within such a segment can be
   given \emph{additional} emphasis.
\end{em}

It is sometimes necessary to prevent \LaTeX\ from
breaking a line where it might otherwise do so.
This may be at a space, as between the ``Mr.''\ and
``Jones'' in
       ``Mr.~Jones'',        % ~ produces an unbreakable interword space.
or within a word---especially when the word is a
symbol like
       \mbox{\emph{itemnum}}
that makes little sense when hyphenated across
lines.

Footnotes\footnote{This is an example of a footnote.}
pose no problem.

\LaTeX\ is good at typesetting mathematical formulas
like
       \( x-3y + z = 7 \)
or
       \( a_{1} > x^{2n} + y^{2n} > x' \)
or
       \( \ip{A}{B} = \sum_{i} a_{i} b_{i} \).
The spaces you type in a formula are
ignored.  Remember that a letter like
       $x$                   % $ ... $  and  \( ... \)  are equivalent
is a formula when it denotes a mathematical
symbol, and it should be typed as one.

\section{Displayed Text}

Text is displayed by indenting it from the left
margin.  Quotations are commonly displayed.  There
are short quotations
\begin{quote}
   This is a short quotation.  It consists of a
   single paragraph of text.  See how it is formatted.
\end{quote}
and longer ones.
\begin{quotation}
   This is a longer quotation.  It consists of two
   paragraphs of text, neither of which are
   particularly interesting.

   This is the second paragraph of the quotation.  It
   is just as dull as the first paragraph.
\end{quotation}
Another frequently-displayed structure is a list.
The following is an example of an \emph{itemized}
list.
\begin{itemize}
   \item This is the first item of an itemized list.
         Each item in the list is marked with a ``tick''.
         You don't have to worry about what kind of tick
         mark is used.

   \item This is the second item of the list.  It
         contains another list nested inside it.  The inner
         list is an \emph{enumerated} list.
         \begin{enumerate}
            \item This is the first item of an enumerated
                  list that is nested within the itemized list.

            \item This is the second item of the inner list.
                  \LaTeX\ allows you to nest lists deeper than
                  you really should.
         \end{enumerate}
         This is the rest of the second item of the outer
         list.  It is no more interesting than any other
         part of the item.
   \item This is the third item of the list.
\end{itemize}
You can even display poetry.
\begin{verse}
   There is an environment
    for verse \\             % The \\ command separates lines
   Whose features some poets % within a stanza.
   will curse.

                             % One or more blank lines separate stanzas.

   For instead of making\\
   Them do \emph{all} line breaking, \\
   It allows them to put too many words on a line when they'd rather be
   forced to be terse.
\end{verse}

Mathematical formulas may also be displayed.  A
displayed formula
is
one-line long; multiline
formulas require special formatting instructions.
   \[  \ip{\Gamma}{\psi'} = x'' + y^{2} + z_{i}^{n}\]
Don't start a paragraph with a displayed equation,
nor make one a paragraph by itself.

\end{document}               % End of document.

\section{Bench Huge Section 1047}
% synthetic bench section 1047
% This is a sample LaTeX input file.  (Version of 12 August 2004.)
%
% A '%' character causes TeX to ignore all remaining text on the line,
% and is used for comments like this one.

\documentclass{article}      % Specifies the document class

                             % The preamble begins here.
\title{An Example Document}  % Declares the document's title.
\author{Leslie Lamport}      % Declares the author's name.
\date{January 21, 1994}      % Deleting this command produces today's date.

\newcommand{\ip}[2]{(#1, #2)}
                             % Defines \ip{arg1}{arg2} to mean
                             % (arg1, arg2).

%\newcommand{\ip}[2]{\langle #1 | #2\rangle}
                             % This is an alternative definition of
                             % \ip that is commented out.

\begin{document}             % End of preamble and beginning of text.

\maketitle                   % Produces the title.

This is an example input file.  Comparing it with
the output it generates can show you how to
produce a simple document of your own.

\section{Ordinary Text}      % Produces section heading.  Lower-level
                             % sections are begun with similar
                             % \subsection and \subsubsection commands.

The ends  of words and sentences are marked
  by   spaces. It  doesn't matter how many
spaces    you type; one is as good as 100.  The
end of   a line counts as a space.

One   or more   blank lines denote the  end
of  a paragraph.

Since any number of consecutive spaces are treated
like a single one, the formatting of the input
file makes no difference to
      \LaTeX,                % The \LaTeX command generates the LaTeX logo.
but it makes a difference to you.  When you use
\LaTeX, making your input file as easy to read
as possible will be a great help as you write
your document and when you change it.  This sample
file shows how you can add comments to your own input
file.

Because printing is different from typewriting,
there are a number of things that you have to do
differently when preparing an input file than if
you were just typing the document directly.
Quotation marks like
       ``this''
have to be handled specially, as do quotes within
quotes:
       ``\,`this'            % \, separates the double and single quote.
        is what I just
        wrote, not  `that'\,''.

Dashes come in three sizes: an
       intra-word
dash, a medium dash for number ranges like
       1--2,
and a punctuation
       dash---like
this.

A sentence-ending space should be larger than the
space between words within a sentence.  You
sometimes have to type special commands in
conjunction with punctuation characters to get
this right, as in the following sentence.
       Gnats, gnus, etc.\ all  % `\ ' makes an inter-word space.
       begin with G\@.         % \@ marks end-of-sentence punctuation.
You should check the spaces after periods when
reading your output to make sure you haven't
forgotten any special cases.  Generating an
ellipsis
       \ldots\               % `\ ' is needed after `\ldots' because TeX
                             % ignores spaces after command names like \ldots
                             % made from \ + letters.
                             %
                             % Note how a `%' character causes TeX to ignore
                             % the end of the input line, so these blank lines
                             % do not start a new paragraph.
                             %
with the right spacing around the periods requires
a special command.

\LaTeX\ interprets some common characters as
commands, so you must type special commands to
generate them.  These characters include the
following:
       \$ \& \% \# \{ and \}.

In printing, text is usually emphasized with an
       \emph{italic}
type style.

\begin{em}
   A long segment of text can also be emphasized
   in this way.  Text within such a segment can be
   given \emph{additional} emphasis.
\end{em}

It is sometimes necessary to prevent \LaTeX\ from
breaking a line where it might otherwise do so.
This may be at a space, as between the ``Mr.''\ and
``Jones'' in
       ``Mr.~Jones'',        % ~ produces an unbreakable interword space.
or within a word---especially when the word is a
symbol like
       \mbox{\emph{itemnum}}
that makes little sense when hyphenated across
lines.

Footnotes\footnote{This is an example of a footnote.}
pose no problem.

\LaTeX\ is good at typesetting mathematical formulas
like
       \( x-3y + z = 7 \)
or
       \( a_{1} > x^{2n} + y^{2n} > x' \)
or
       \( \ip{A}{B} = \sum_{i} a_{i} b_{i} \).
The spaces you type in a formula are
ignored.  Remember that a letter like
       $x$                   % $ ... $  and  \( ... \)  are equivalent
is a formula when it denotes a mathematical
symbol, and it should be typed as one.

\section{Displayed Text}

Text is displayed by indenting it from the left
margin.  Quotations are commonly displayed.  There
are short quotations
\begin{quote}
   This is a short quotation.  It consists of a
   single paragraph of text.  See how it is formatted.
\end{quote}
and longer ones.
\begin{quotation}
   This is a longer quotation.  It consists of two
   paragraphs of text, neither of which are
   particularly interesting.

   This is the second paragraph of the quotation.  It
   is just as dull as the first paragraph.
\end{quotation}
Another frequently-displayed structure is a list.
The following is an example of an \emph{itemized}
list.
\begin{itemize}
   \item This is the first item of an itemized list.
         Each item in the list is marked with a ``tick''.
         You don't have to worry about what kind of tick
         mark is used.

   \item This is the second item of the list.  It
         contains another list nested inside it.  The inner
         list is an \emph{enumerated} list.
         \begin{enumerate}
            \item This is the first item of an enumerated
                  list that is nested within the itemized list.

            \item This is the second item of the inner list.
                  \LaTeX\ allows you to nest lists deeper than
                  you really should.
         \end{enumerate}
         This is the rest of the second item of the outer
         list.  It is no more interesting than any other
         part of the item.
   \item This is the third item of the list.
\end{itemize}
You can even display poetry.
\begin{verse}
   There is an environment
    for verse \\             % The \\ command separates lines
   Whose features some poets % within a stanza.
   will curse.

                             % One or more blank lines separate stanzas.

   For instead of making\\
   Them do \emph{all} line breaking, \\
   It allows them to put too many words on a line when they'd rather be
   forced to be terse.
\end{verse}

Mathematical formulas may also be displayed.  A
displayed formula
is
one-line long; multiline
formulas require special formatting instructions.
   \[  \ip{\Gamma}{\psi'} = x'' + y^{2} + z_{i}^{n}\]
Don't start a paragraph with a displayed equation,
nor make one a paragraph by itself.

\end{document}               % End of document.

\section{Bench Huge Section 1048}
% synthetic bench section 1048
% This is a sample LaTeX input file.  (Version of 12 August 2004.)
%
% A '%' character causes TeX to ignore all remaining text on the line,
% and is used for comments like this one.

\documentclass{article}      % Specifies the document class

                             % The preamble begins here.
\title{An Example Document}  % Declares the document's title.
\author{Leslie Lamport}      % Declares the author's name.
\date{January 21, 1994}      % Deleting this command produces today's date.

\newcommand{\ip}[2]{(#1, #2)}
                             % Defines \ip{arg1}{arg2} to mean
                             % (arg1, arg2).

%\newcommand{\ip}[2]{\langle #1 | #2\rangle}
                             % This is an alternative definition of
                             % \ip that is commented out.

\begin{document}             % End of preamble and beginning of text.

\maketitle                   % Produces the title.

This is an example input file.  Comparing it with
the output it generates can show you how to
produce a simple document of your own.

\section{Ordinary Text}      % Produces section heading.  Lower-level
                             % sections are begun with similar
                             % \subsection and \subsubsection commands.

The ends  of words and sentences are marked
  by   spaces. It  doesn't matter how many
spaces    you type; one is as good as 100.  The
end of   a line counts as a space.

One   or more   blank lines denote the  end
of  a paragraph.

Since any number of consecutive spaces are treated
like a single one, the formatting of the input
file makes no difference to
      \LaTeX,                % The \LaTeX command generates the LaTeX logo.
but it makes a difference to you.  When you use
\LaTeX, making your input file as easy to read
as possible will be a great help as you write
your document and when you change it.  This sample
file shows how you can add comments to your own input
file.

Because printing is different from typewriting,
there are a number of things that you have to do
differently when preparing an input file than if
you were just typing the document directly.
Quotation marks like
       ``this''
have to be handled specially, as do quotes within
quotes:
       ``\,`this'            % \, separates the double and single quote.
        is what I just
        wrote, not  `that'\,''.

Dashes come in three sizes: an
       intra-word
dash, a medium dash for number ranges like
       1--2,
and a punctuation
       dash---like
this.

A sentence-ending space should be larger than the
space between words within a sentence.  You
sometimes have to type special commands in
conjunction with punctuation characters to get
this right, as in the following sentence.
       Gnats, gnus, etc.\ all  % `\ ' makes an inter-word space.
       begin with G\@.         % \@ marks end-of-sentence punctuation.
You should check the spaces after periods when
reading your output to make sure you haven't
forgotten any special cases.  Generating an
ellipsis
       \ldots\               % `\ ' is needed after `\ldots' because TeX
                             % ignores spaces after command names like \ldots
                             % made from \ + letters.
                             %
                             % Note how a `%' character causes TeX to ignore
                             % the end of the input line, so these blank lines
                             % do not start a new paragraph.
                             %
with the right spacing around the periods requires
a special command.

\LaTeX\ interprets some common characters as
commands, so you must type special commands to
generate them.  These characters include the
following:
       \$ \& \% \# \{ and \}.

In printing, text is usually emphasized with an
       \emph{italic}
type style.

\begin{em}
   A long segment of text can also be emphasized
   in this way.  Text within such a segment can be
   given \emph{additional} emphasis.
\end{em}

It is sometimes necessary to prevent \LaTeX\ from
breaking a line where it might otherwise do so.
This may be at a space, as between the ``Mr.''\ and
``Jones'' in
       ``Mr.~Jones'',        % ~ produces an unbreakable interword space.
or within a word---especially when the word is a
symbol like
       \mbox{\emph{itemnum}}
that makes little sense when hyphenated across
lines.

Footnotes\footnote{This is an example of a footnote.}
pose no problem.

\LaTeX\ is good at typesetting mathematical formulas
like
       \( x-3y + z = 7 \)
or
       \( a_{1} > x^{2n} + y^{2n} > x' \)
or
       \( \ip{A}{B} = \sum_{i} a_{i} b_{i} \).
The spaces you type in a formula are
ignored.  Remember that a letter like
       $x$                   % $ ... $  and  \( ... \)  are equivalent
is a formula when it denotes a mathematical
symbol, and it should be typed as one.

\section{Displayed Text}

Text is displayed by indenting it from the left
margin.  Quotations are commonly displayed.  There
are short quotations
\begin{quote}
   This is a short quotation.  It consists of a
   single paragraph of text.  See how it is formatted.
\end{quote}
and longer ones.
\begin{quotation}
   This is a longer quotation.  It consists of two
   paragraphs of text, neither of which are
   particularly interesting.

   This is the second paragraph of the quotation.  It
   is just as dull as the first paragraph.
\end{quotation}
Another frequently-displayed structure is a list.
The following is an example of an \emph{itemized}
list.
\begin{itemize}
   \item This is the first item of an itemized list.
         Each item in the list is marked with a ``tick''.
         You don't have to worry about what kind of tick
         mark is used.

   \item This is the second item of the list.  It
         contains another list nested inside it.  The inner
         list is an \emph{enumerated} list.
         \begin{enumerate}
            \item This is the first item of an enumerated
                  list that is nested within the itemized list.

            \item This is the second item of the inner list.
                  \LaTeX\ allows you to nest lists deeper than
                  you really should.
         \end{enumerate}
         This is the rest of the second item of the outer
         list.  It is no more interesting than any other
         part of the item.
   \item This is the third item of the list.
\end{itemize}
You can even display poetry.
\begin{verse}
   There is an environment
    for verse \\             % The \\ command separates lines
   Whose features some poets % within a stanza.
   will curse.

                             % One or more blank lines separate stanzas.

   For instead of making\\
   Them do \emph{all} line breaking, \\
   It allows them to put too many words on a line when they'd rather be
   forced to be terse.
\end{verse}

Mathematical formulas may also be displayed.  A
displayed formula
is
one-line long; multiline
formulas require special formatting instructions.
   \[  \ip{\Gamma}{\psi'} = x'' + y^{2} + z_{i}^{n}\]
Don't start a paragraph with a displayed equation,
nor make one a paragraph by itself.

\end{document}               % End of document.

\section{Bench Huge Section 1049}
% synthetic bench section 1049
% This is a sample LaTeX input file.  (Version of 12 August 2004.)
%
% A '%' character causes TeX to ignore all remaining text on the line,
% and is used for comments like this one.

\documentclass{article}      % Specifies the document class

                             % The preamble begins here.
\title{An Example Document}  % Declares the document's title.
\author{Leslie Lamport}      % Declares the author's name.
\date{January 21, 1994}      % Deleting this command produces today's date.

\newcommand{\ip}[2]{(#1, #2)}
                             % Defines \ip{arg1}{arg2} to mean
                             % (arg1, arg2).

%\newcommand{\ip}[2]{\langle #1 | #2\rangle}
                             % This is an alternative definition of
                             % \ip that is commented out.

\begin{document}             % End of preamble and beginning of text.

\maketitle                   % Produces the title.

This is an example input file.  Comparing it with
the output it generates can show you how to
produce a simple document of your own.

\section{Ordinary Text}      % Produces section heading.  Lower-level
                             % sections are begun with similar
                             % \subsection and \subsubsection commands.

The ends  of words and sentences are marked
  by   spaces. It  doesn't matter how many
spaces    you type; one is as good as 100.  The
end of   a line counts as a space.

One   or more   blank lines denote the  end
of  a paragraph.

Since any number of consecutive spaces are treated
like a single one, the formatting of the input
file makes no difference to
      \LaTeX,                % The \LaTeX command generates the LaTeX logo.
but it makes a difference to you.  When you use
\LaTeX, making your input file as easy to read
as possible will be a great help as you write
your document and when you change it.  This sample
file shows how you can add comments to your own input
file.

Because printing is different from typewriting,
there are a number of things that you have to do
differently when preparing an input file than if
you were just typing the document directly.
Quotation marks like
       ``this''
have to be handled specially, as do quotes within
quotes:
       ``\,`this'            % \, separates the double and single quote.
        is what I just
        wrote, not  `that'\,''.

Dashes come in three sizes: an
       intra-word
dash, a medium dash for number ranges like
       1--2,
and a punctuation
       dash---like
this.

A sentence-ending space should be larger than the
space between words within a sentence.  You
sometimes have to type special commands in
conjunction with punctuation characters to get
this right, as in the following sentence.
       Gnats, gnus, etc.\ all  % `\ ' makes an inter-word space.
       begin with G\@.         % \@ marks end-of-sentence punctuation.
You should check the spaces after periods when
reading your output to make sure you haven't
forgotten any special cases.  Generating an
ellipsis
       \ldots\               % `\ ' is needed after `\ldots' because TeX
                             % ignores spaces after command names like \ldots
                             % made from \ + letters.
                             %
                             % Note how a `%' character causes TeX to ignore
                             % the end of the input line, so these blank lines
                             % do not start a new paragraph.
                             %
with the right spacing around the periods requires
a special command.

\LaTeX\ interprets some common characters as
commands, so you must type special commands to
generate them.  These characters include the
following:
       \$ \& \% \# \{ and \}.

In printing, text is usually emphasized with an
       \emph{italic}
type style.

\begin{em}
   A long segment of text can also be emphasized
   in this way.  Text within such a segment can be
   given \emph{additional} emphasis.
\end{em}

It is sometimes necessary to prevent \LaTeX\ from
breaking a line where it might otherwise do so.
This may be at a space, as between the ``Mr.''\ and
``Jones'' in
       ``Mr.~Jones'',        % ~ produces an unbreakable interword space.
or within a word---especially when the word is a
symbol like
       \mbox{\emph{itemnum}}
that makes little sense when hyphenated across
lines.

Footnotes\footnote{This is an example of a footnote.}
pose no problem.

\LaTeX\ is good at typesetting mathematical formulas
like
       \( x-3y + z = 7 \)
or
       \( a_{1} > x^{2n} + y^{2n} > x' \)
or
       \( \ip{A}{B} = \sum_{i} a_{i} b_{i} \).
The spaces you type in a formula are
ignored.  Remember that a letter like
       $x$                   % $ ... $  and  \( ... \)  are equivalent
is a formula when it denotes a mathematical
symbol, and it should be typed as one.

\section{Displayed Text}

Text is displayed by indenting it from the left
margin.  Quotations are commonly displayed.  There
are short quotations
\begin{quote}
   This is a short quotation.  It consists of a
   single paragraph of text.  See how it is formatted.
\end{quote}
and longer ones.
\begin{quotation}
   This is a longer quotation.  It consists of two
   paragraphs of text, neither of which are
   particularly interesting.

   This is the second paragraph of the quotation.  It
   is just as dull as the first paragraph.
\end{quotation}
Another frequently-displayed structure is a list.
The following is an example of an \emph{itemized}
list.
\begin{itemize}
   \item This is the first item of an itemized list.
         Each item in the list is marked with a ``tick''.
         You don't have to worry about what kind of tick
         mark is used.

   \item This is the second item of the list.  It
         contains another list nested inside it.  The inner
         list is an \emph{enumerated} list.
         \begin{enumerate}
            \item This is the first item of an enumerated
                  list that is nested within the itemized list.

            \item This is the second item of the inner list.
                  \LaTeX\ allows you to nest lists deeper than
                  you really should.
         \end{enumerate}
         This is the rest of the second item of the outer
         list.  It is no more interesting than any other
         part of the item.
   \item This is the third item of the list.
\end{itemize}
You can even display poetry.
\begin{verse}
   There is an environment
    for verse \\             % The \\ command separates lines
   Whose features some poets % within a stanza.
   will curse.

                             % One or more blank lines separate stanzas.

   For instead of making\\
   Them do \emph{all} line breaking, \\
   It allows them to put too many words on a line when they'd rather be
   forced to be terse.
\end{verse}

Mathematical formulas may also be displayed.  A
displayed formula
is
one-line long; multiline
formulas require special formatting instructions.
   \[  \ip{\Gamma}{\psi'} = x'' + y^{2} + z_{i}^{n}\]
Don't start a paragraph with a displayed equation,
nor make one a paragraph by itself.

\end{document}               % End of document.

\section{Bench Huge Section 1050}
% synthetic bench section 1050
% This is a sample LaTeX input file.  (Version of 12 August 2004.)
%
% A '%' character causes TeX to ignore all remaining text on the line,
% and is used for comments like this one.

\documentclass{article}      % Specifies the document class

                             % The preamble begins here.
\title{An Example Document}  % Declares the document's title.
\author{Leslie Lamport}      % Declares the author's name.
\date{January 21, 1994}      % Deleting this command produces today's date.

\newcommand{\ip}[2]{(#1, #2)}
                             % Defines \ip{arg1}{arg2} to mean
                             % (arg1, arg2).

%\newcommand{\ip}[2]{\langle #1 | #2\rangle}
                             % This is an alternative definition of
                             % \ip that is commented out.

\begin{document}             % End of preamble and beginning of text.

\maketitle                   % Produces the title.

This is an example input file.  Comparing it with
the output it generates can show you how to
produce a simple document of your own.

\section{Ordinary Text}      % Produces section heading.  Lower-level
                             % sections are begun with similar
                             % \subsection and \subsubsection commands.

The ends  of words and sentences are marked
  by   spaces. It  doesn't matter how many
spaces    you type; one is as good as 100.  The
end of   a line counts as a space.

One   or more   blank lines denote the  end
of  a paragraph.

Since any number of consecutive spaces are treated
like a single one, the formatting of the input
file makes no difference to
      \LaTeX,                % The \LaTeX command generates the LaTeX logo.
but it makes a difference to you.  When you use
\LaTeX, making your input file as easy to read
as possible will be a great help as you write
your document and when you change it.  This sample
file shows how you can add comments to your own input
file.

Because printing is different from typewriting,
there are a number of things that you have to do
differently when preparing an input file than if
you were just typing the document directly.
Quotation marks like
       ``this''
have to be handled specially, as do quotes within
quotes:
       ``\,`this'            % \, separates the double and single quote.
        is what I just
        wrote, not  `that'\,''.

Dashes come in three sizes: an
       intra-word
dash, a medium dash for number ranges like
       1--2,
and a punctuation
       dash---like
this.

A sentence-ending space should be larger than the
space between words within a sentence.  You
sometimes have to type special commands in
conjunction with punctuation characters to get
this right, as in the following sentence.
       Gnats, gnus, etc.\ all  % `\ ' makes an inter-word space.
       begin with G\@.         % \@ marks end-of-sentence punctuation.
You should check the spaces after periods when
reading your output to make sure you haven't
forgotten any special cases.  Generating an
ellipsis
       \ldots\               % `\ ' is needed after `\ldots' because TeX
                             % ignores spaces after command names like \ldots
                             % made from \ + letters.
                             %
                             % Note how a `%' character causes TeX to ignore
                             % the end of the input line, so these blank lines
                             % do not start a new paragraph.
                             %
with the right spacing around the periods requires
a special command.

\LaTeX\ interprets some common characters as
commands, so you must type special commands to
generate them.  These characters include the
following:
       \$ \& \% \# \{ and \}.

In printing, text is usually emphasized with an
       \emph{italic}
type style.

\begin{em}
   A long segment of text can also be emphasized
   in this way.  Text within such a segment can be
   given \emph{additional} emphasis.
\end{em}

It is sometimes necessary to prevent \LaTeX\ from
breaking a line where it might otherwise do so.
This may be at a space, as between the ``Mr.''\ and
``Jones'' in
       ``Mr.~Jones'',        % ~ produces an unbreakable interword space.
or within a word---especially when the word is a
symbol like
       \mbox{\emph{itemnum}}
that makes little sense when hyphenated across
lines.

Footnotes\footnote{This is an example of a footnote.}
pose no problem.

\LaTeX\ is good at typesetting mathematical formulas
like
       \( x-3y + z = 7 \)
or
       \( a_{1} > x^{2n} + y^{2n} > x' \)
or
       \( \ip{A}{B} = \sum_{i} a_{i} b_{i} \).
The spaces you type in a formula are
ignored.  Remember that a letter like
       $x$                   % $ ... $  and  \( ... \)  are equivalent
is a formula when it denotes a mathematical
symbol, and it should be typed as one.

\section{Displayed Text}

Text is displayed by indenting it from the left
margin.  Quotations are commonly displayed.  There
are short quotations
\begin{quote}
   This is a short quotation.  It consists of a
   single paragraph of text.  See how it is formatted.
\end{quote}
and longer ones.
\begin{quotation}
   This is a longer quotation.  It consists of two
   paragraphs of text, neither of which are
   particularly interesting.

   This is the second paragraph of the quotation.  It
   is just as dull as the first paragraph.
\end{quotation}
Another frequently-displayed structure is a list.
The following is an example of an \emph{itemized}
list.
\begin{itemize}
   \item This is the first item of an itemized list.
         Each item in the list is marked with a ``tick''.
         You don't have to worry about what kind of tick
         mark is used.

   \item This is the second item of the list.  It
         contains another list nested inside it.  The inner
         list is an \emph{enumerated} list.
         \begin{enumerate}
            \item This is the first item of an enumerated
                  list that is nested within the itemized list.

            \item This is the second item of the inner list.
                  \LaTeX\ allows you to nest lists deeper than
                  you really should.
         \end{enumerate}
         This is the rest of the second item of the outer
         list.  It is no more interesting than any other
         part of the item.
   \item This is the third item of the list.
\end{itemize}
You can even display poetry.
\begin{verse}
   There is an environment
    for verse \\             % The \\ command separates lines
   Whose features some poets % within a stanza.
   will curse.

                             % One or more blank lines separate stanzas.

   For instead of making\\
   Them do \emph{all} line breaking, \\
   It allows them to put too many words on a line when they'd rather be
   forced to be terse.
\end{verse}

Mathematical formulas may also be displayed.  A
displayed formula
is
one-line long; multiline
formulas require special formatting instructions.
   \[  \ip{\Gamma}{\psi'} = x'' + y^{2} + z_{i}^{n}\]
Don't start a paragraph with a displayed equation,
nor make one a paragraph by itself.

\end{document}               % End of document.

\section{Bench Huge Section 1051}
% synthetic bench section 1051
% This is a sample LaTeX input file.  (Version of 12 August 2004.)
%
% A '%' character causes TeX to ignore all remaining text on the line,
% and is used for comments like this one.

\documentclass{article}      % Specifies the document class

                             % The preamble begins here.
\title{An Example Document}  % Declares the document's title.
\author{Leslie Lamport}      % Declares the author's name.
\date{January 21, 1994}      % Deleting this command produces today's date.

\newcommand{\ip}[2]{(#1, #2)}
                             % Defines \ip{arg1}{arg2} to mean
                             % (arg1, arg2).

%\newcommand{\ip}[2]{\langle #1 | #2\rangle}
                             % This is an alternative definition of
                             % \ip that is commented out.

\begin{document}             % End of preamble and beginning of text.

\maketitle                   % Produces the title.

This is an example input file.  Comparing it with
the output it generates can show you how to
produce a simple document of your own.

\section{Ordinary Text}      % Produces section heading.  Lower-level
                             % sections are begun with similar
                             % \subsection and \subsubsection commands.

The ends  of words and sentences are marked
  by   spaces. It  doesn't matter how many
spaces    you type; one is as good as 100.  The
end of   a line counts as a space.

One   or more   blank lines denote the  end
of  a paragraph.

Since any number of consecutive spaces are treated
like a single one, the formatting of the input
file makes no difference to
      \LaTeX,                % The \LaTeX command generates the LaTeX logo.
but it makes a difference to you.  When you use
\LaTeX, making your input file as easy to read
as possible will be a great help as you write
your document and when you change it.  This sample
file shows how you can add comments to your own input
file.

Because printing is different from typewriting,
there are a number of things that you have to do
differently when preparing an input file than if
you were just typing the document directly.
Quotation marks like
       ``this''
have to be handled specially, as do quotes within
quotes:
       ``\,`this'            % \, separates the double and single quote.
        is what I just
        wrote, not  `that'\,''.

Dashes come in three sizes: an
       intra-word
dash, a medium dash for number ranges like
       1--2,
and a punctuation
       dash---like
this.

A sentence-ending space should be larger than the
space between words within a sentence.  You
sometimes have to type special commands in
conjunction with punctuation characters to get
this right, as in the following sentence.
       Gnats, gnus, etc.\ all  % `\ ' makes an inter-word space.
       begin with G\@.         % \@ marks end-of-sentence punctuation.
You should check the spaces after periods when
reading your output to make sure you haven't
forgotten any special cases.  Generating an
ellipsis
       \ldots\               % `\ ' is needed after `\ldots' because TeX
                             % ignores spaces after command names like \ldots
                             % made from \ + letters.
                             %
                             % Note how a `%' character causes TeX to ignore
                             % the end of the input line, so these blank lines
                             % do not start a new paragraph.
                             %
with the right spacing around the periods requires
a special command.

\LaTeX\ interprets some common characters as
commands, so you must type special commands to
generate them.  These characters include the
following:
       \$ \& \% \# \{ and \}.

In printing, text is usually emphasized with an
       \emph{italic}
type style.

\begin{em}
   A long segment of text can also be emphasized
   in this way.  Text within such a segment can be
   given \emph{additional} emphasis.
\end{em}

It is sometimes necessary to prevent \LaTeX\ from
breaking a line where it might otherwise do so.
This may be at a space, as between the ``Mr.''\ and
``Jones'' in
       ``Mr.~Jones'',        % ~ produces an unbreakable interword space.
or within a word---especially when the word is a
symbol like
       \mbox{\emph{itemnum}}
that makes little sense when hyphenated across
lines.

Footnotes\footnote{This is an example of a footnote.}
pose no problem.

\LaTeX\ is good at typesetting mathematical formulas
like
       \( x-3y + z = 7 \)
or
       \( a_{1} > x^{2n} + y^{2n} > x' \)
or
       \( \ip{A}{B} = \sum_{i} a_{i} b_{i} \).
The spaces you type in a formula are
ignored.  Remember that a letter like
       $x$                   % $ ... $  and  \( ... \)  are equivalent
is a formula when it denotes a mathematical
symbol, and it should be typed as one.

\section{Displayed Text}

Text is displayed by indenting it from the left
margin.  Quotations are commonly displayed.  There
are short quotations
\begin{quote}
   This is a short quotation.  It consists of a
   single paragraph of text.  See how it is formatted.
\end{quote}
and longer ones.
\begin{quotation}
   This is a longer quotation.  It consists of two
   paragraphs of text, neither of which are
   particularly interesting.

   This is the second paragraph of the quotation.  It
   is just as dull as the first paragraph.
\end{quotation}
Another frequently-displayed structure is a list.
The following is an example of an \emph{itemized}
list.
\begin{itemize}
   \item This is the first item of an itemized list.
         Each item in the list is marked with a ``tick''.
         You don't have to worry about what kind of tick
         mark is used.

   \item This is the second item of the list.  It
         contains another list nested inside it.  The inner
         list is an \emph{enumerated} list.
         \begin{enumerate}
            \item This is the first item of an enumerated
                  list that is nested within the itemized list.

            \item This is the second item of the inner list.
                  \LaTeX\ allows you to nest lists deeper than
                  you really should.
         \end{enumerate}
         This is the rest of the second item of the outer
         list.  It is no more interesting than any other
         part of the item.
   \item This is the third item of the list.
\end{itemize}
You can even display poetry.
\begin{verse}
   There is an environment
    for verse \\             % The \\ command separates lines
   Whose features some poets % within a stanza.
   will curse.

                             % One or more blank lines separate stanzas.

   For instead of making\\
   Them do \emph{all} line breaking, \\
   It allows them to put too many words on a line when they'd rather be
   forced to be terse.
\end{verse}

Mathematical formulas may also be displayed.  A
displayed formula
is
one-line long; multiline
formulas require special formatting instructions.
   \[  \ip{\Gamma}{\psi'} = x'' + y^{2} + z_{i}^{n}\]
Don't start a paragraph with a displayed equation,
nor make one a paragraph by itself.

\end{document}               % End of document.

\section{Bench Huge Section 1052}
% synthetic bench section 1052
% This is a sample LaTeX input file.  (Version of 12 August 2004.)
%
% A '%' character causes TeX to ignore all remaining text on the line,
% and is used for comments like this one.

\documentclass{article}      % Specifies the document class

                             % The preamble begins here.
\title{An Example Document}  % Declares the document's title.
\author{Leslie Lamport}      % Declares the author's name.
\date{January 21, 1994}      % Deleting this command produces today's date.

\newcommand{\ip}[2]{(#1, #2)}
                             % Defines \ip{arg1}{arg2} to mean
                             % (arg1, arg2).

%\newcommand{\ip}[2]{\langle #1 | #2\rangle}
                             % This is an alternative definition of
                             % \ip that is commented out.

\begin{document}             % End of preamble and beginning of text.

\maketitle                   % Produces the title.

This is an example input file.  Comparing it with
the output it generates can show you how to
produce a simple document of your own.

\section{Ordinary Text}      % Produces section heading.  Lower-level
                             % sections are begun with similar
                             % \subsection and \subsubsection commands.

The ends  of words and sentences are marked
  by   spaces. It  doesn't matter how many
spaces    you type; one is as good as 100.  The
end of   a line counts as a space.

One   or more   blank lines denote the  end
of  a paragraph.

Since any number of consecutive spaces are treated
like a single one, the formatting of the input
file makes no difference to
      \LaTeX,                % The \LaTeX command generates the LaTeX logo.
but it makes a difference to you.  When you use
\LaTeX, making your input file as easy to read
as possible will be a great help as you write
your document and when you change it.  This sample
file shows how you can add comments to your own input
file.

Because printing is different from typewriting,
there are a number of things that you have to do
differently when preparing an input file than if
you were just typing the document directly.
Quotation marks like
       ``this''
have to be handled specially, as do quotes within
quotes:
       ``\,`this'            % \, separates the double and single quote.
        is what I just
        wrote, not  `that'\,''.

Dashes come in three sizes: an
       intra-word
dash, a medium dash for number ranges like
       1--2,
and a punctuation
       dash---like
this.

A sentence-ending space should be larger than the
space between words within a sentence.  You
sometimes have to type special commands in
conjunction with punctuation characters to get
this right, as in the following sentence.
       Gnats, gnus, etc.\ all  % `\ ' makes an inter-word space.
       begin with G\@.         % \@ marks end-of-sentence punctuation.
You should check the spaces after periods when
reading your output to make sure you haven't
forgotten any special cases.  Generating an
ellipsis
       \ldots\               % `\ ' is needed after `\ldots' because TeX
                             % ignores spaces after command names like \ldots
                             % made from \ + letters.
                             %
                             % Note how a `%' character causes TeX to ignore
                             % the end of the input line, so these blank lines
                             % do not start a new paragraph.
                             %
with the right spacing around the periods requires
a special command.

\LaTeX\ interprets some common characters as
commands, so you must type special commands to
generate them.  These characters include the
following:
       \$ \& \% \# \{ and \}.

In printing, text is usually emphasized with an
       \emph{italic}
type style.

\begin{em}
   A long segment of text can also be emphasized
   in this way.  Text within such a segment can be
   given \emph{additional} emphasis.
\end{em}

It is sometimes necessary to prevent \LaTeX\ from
breaking a line where it might otherwise do so.
This may be at a space, as between the ``Mr.''\ and
``Jones'' in
       ``Mr.~Jones'',        % ~ produces an unbreakable interword space.
or within a word---especially when the word is a
symbol like
       \mbox{\emph{itemnum}}
that makes little sense when hyphenated across
lines.

Footnotes\footnote{This is an example of a footnote.}
pose no problem.

\LaTeX\ is good at typesetting mathematical formulas
like
       \( x-3y + z = 7 \)
or
       \( a_{1} > x^{2n} + y^{2n} > x' \)
or
       \( \ip{A}{B} = \sum_{i} a_{i} b_{i} \).
The spaces you type in a formula are
ignored.  Remember that a letter like
       $x$                   % $ ... $  and  \( ... \)  are equivalent
is a formula when it denotes a mathematical
symbol, and it should be typed as one.

\section{Displayed Text}

Text is displayed by indenting it from the left
margin.  Quotations are commonly displayed.  There
are short quotations
\begin{quote}
   This is a short quotation.  It consists of a
   single paragraph of text.  See how it is formatted.
\end{quote}
and longer ones.
\begin{quotation}
   This is a longer quotation.  It consists of two
   paragraphs of text, neither of which are
   particularly interesting.

   This is the second paragraph of the quotation.  It
   is just as dull as the first paragraph.
\end{quotation}
Another frequently-displayed structure is a list.
The following is an example of an \emph{itemized}
list.
\begin{itemize}
   \item This is the first item of an itemized list.
         Each item in the list is marked with a ``tick''.
         You don't have to worry about what kind of tick
         mark is used.

   \item This is the second item of the list.  It
         contains another list nested inside it.  The inner
         list is an \emph{enumerated} list.
         \begin{enumerate}
            \item This is the first item of an enumerated
                  list that is nested within the itemized list.

            \item This is the second item of the inner list.
                  \LaTeX\ allows you to nest lists deeper than
                  you really should.
         \end{enumerate}
         This is the rest of the second item of the outer
         list.  It is no more interesting than any other
         part of the item.
   \item This is the third item of the list.
\end{itemize}
You can even display poetry.
\begin{verse}
   There is an environment
    for verse \\             % The \\ command separates lines
   Whose features some poets % within a stanza.
   will curse.

                             % One or more blank lines separate stanzas.

   For instead of making\\
   Them do \emph{all} line breaking, \\
   It allows them to put too many words on a line when they'd rather be
   forced to be terse.
\end{verse}

Mathematical formulas may also be displayed.  A
displayed formula
is
one-line long; multiline
formulas require special formatting instructions.
   \[  \ip{\Gamma}{\psi'} = x'' + y^{2} + z_{i}^{n}\]
Don't start a paragraph with a displayed equation,
nor make one a paragraph by itself.

\end{document}               % End of document.

\section{Bench Huge Section 1053}
% synthetic bench section 1053
% This is a sample LaTeX input file.  (Version of 12 August 2004.)
%
% A '%' character causes TeX to ignore all remaining text on the line,
% and is used for comments like this one.

\documentclass{article}      % Specifies the document class

                             % The preamble begins here.
\title{An Example Document}  % Declares the document's title.
\author{Leslie Lamport}      % Declares the author's name.
\date{January 21, 1994}      % Deleting this command produces today's date.

\newcommand{\ip}[2]{(#1, #2)}
                             % Defines \ip{arg1}{arg2} to mean
                             % (arg1, arg2).

%\newcommand{\ip}[2]{\langle #1 | #2\rangle}
                             % This is an alternative definition of
                             % \ip that is commented out.

\begin{document}             % End of preamble and beginning of text.

\maketitle                   % Produces the title.

This is an example input file.  Comparing it with
the output it generates can show you how to
produce a simple document of your own.

\section{Ordinary Text}      % Produces section heading.  Lower-level
                             % sections are begun with similar
                             % \subsection and \subsubsection commands.

The ends  of words and sentences are marked
  by   spaces. It  doesn't matter how many
spaces    you type; one is as good as 100.  The
end of   a line counts as a space.

One   or more   blank lines denote the  end
of  a paragraph.

Since any number of consecutive spaces are treated
like a single one, the formatting of the input
file makes no difference to
      \LaTeX,                % The \LaTeX command generates the LaTeX logo.
but it makes a difference to you.  When you use
\LaTeX, making your input file as easy to read
as possible will be a great help as you write
your document and when you change it.  This sample
file shows how you can add comments to your own input
file.

Because printing is different from typewriting,
there are a number of things that you have to do
differently when preparing an input file than if
you were just typing the document directly.
Quotation marks like
       ``this''
have to be handled specially, as do quotes within
quotes:
       ``\,`this'            % \, separates the double and single quote.
        is what I just
        wrote, not  `that'\,''.

Dashes come in three sizes: an
       intra-word
dash, a medium dash for number ranges like
       1--2,
and a punctuation
       dash---like
this.

A sentence-ending space should be larger than the
space between words within a sentence.  You
sometimes have to type special commands in
conjunction with punctuation characters to get
this right, as in the following sentence.
       Gnats, gnus, etc.\ all  % `\ ' makes an inter-word space.
       begin with G\@.         % \@ marks end-of-sentence punctuation.
You should check the spaces after periods when
reading your output to make sure you haven't
forgotten any special cases.  Generating an
ellipsis
       \ldots\               % `\ ' is needed after `\ldots' because TeX
                             % ignores spaces after command names like \ldots
                             % made from \ + letters.
                             %
                             % Note how a `%' character causes TeX to ignore
                             % the end of the input line, so these blank lines
                             % do not start a new paragraph.
                             %
with the right spacing around the periods requires
a special command.

\LaTeX\ interprets some common characters as
commands, so you must type special commands to
generate them.  These characters include the
following:
       \$ \& \% \# \{ and \}.

In printing, text is usually emphasized with an
       \emph{italic}
type style.

\begin{em}
   A long segment of text can also be emphasized
   in this way.  Text within such a segment can be
   given \emph{additional} emphasis.
\end{em}

It is sometimes necessary to prevent \LaTeX\ from
breaking a line where it might otherwise do so.
This may be at a space, as between the ``Mr.''\ and
``Jones'' in
       ``Mr.~Jones'',        % ~ produces an unbreakable interword space.
or within a word---especially when the word is a
symbol like
       \mbox{\emph{itemnum}}
that makes little sense when hyphenated across
lines.

Footnotes\footnote{This is an example of a footnote.}
pose no problem.

\LaTeX\ is good at typesetting mathematical formulas
like
       \( x-3y + z = 7 \)
or
       \( a_{1} > x^{2n} + y^{2n} > x' \)
or
       \( \ip{A}{B} = \sum_{i} a_{i} b_{i} \).
The spaces you type in a formula are
ignored.  Remember that a letter like
       $x$                   % $ ... $  and  \( ... \)  are equivalent
is a formula when it denotes a mathematical
symbol, and it should be typed as one.

\section{Displayed Text}

Text is displayed by indenting it from the left
margin.  Quotations are commonly displayed.  There
are short quotations
\begin{quote}
   This is a short quotation.  It consists of a
   single paragraph of text.  See how it is formatted.
\end{quote}
and longer ones.
\begin{quotation}
   This is a longer quotation.  It consists of two
   paragraphs of text, neither of which are
   particularly interesting.

   This is the second paragraph of the quotation.  It
   is just as dull as the first paragraph.
\end{quotation}
Another frequently-displayed structure is a list.
The following is an example of an \emph{itemized}
list.
\begin{itemize}
   \item This is the first item of an itemized list.
         Each item in the list is marked with a ``tick''.
         You don't have to worry about what kind of tick
         mark is used.

   \item This is the second item of the list.  It
         contains another list nested inside it.  The inner
         list is an \emph{enumerated} list.
         \begin{enumerate}
            \item This is the first item of an enumerated
                  list that is nested within the itemized list.

            \item This is the second item of the inner list.
                  \LaTeX\ allows you to nest lists deeper than
                  you really should.
         \end{enumerate}
         This is the rest of the second item of the outer
         list.  It is no more interesting than any other
         part of the item.
   \item This is the third item of the list.
\end{itemize}
You can even display poetry.
\begin{verse}
   There is an environment
    for verse \\             % The \\ command separates lines
   Whose features some poets % within a stanza.
   will curse.

                             % One or more blank lines separate stanzas.

   For instead of making\\
   Them do \emph{all} line breaking, \\
   It allows them to put too many words on a line when they'd rather be
   forced to be terse.
\end{verse}

Mathematical formulas may also be displayed.  A
displayed formula
is
one-line long; multiline
formulas require special formatting instructions.
   \[  \ip{\Gamma}{\psi'} = x'' + y^{2} + z_{i}^{n}\]
Don't start a paragraph with a displayed equation,
nor make one a paragraph by itself.

\end{document}               % End of document.

\section{Bench Huge Section 1054}
% synthetic bench section 1054
% This is a sample LaTeX input file.  (Version of 12 August 2004.)
%
% A '%' character causes TeX to ignore all remaining text on the line,
% and is used for comments like this one.

\documentclass{article}      % Specifies the document class

                             % The preamble begins here.
\title{An Example Document}  % Declares the document's title.
\author{Leslie Lamport}      % Declares the author's name.
\date{January 21, 1994}      % Deleting this command produces today's date.

\newcommand{\ip}[2]{(#1, #2)}
                             % Defines \ip{arg1}{arg2} to mean
                             % (arg1, arg2).

%\newcommand{\ip}[2]{\langle #1 | #2\rangle}
                             % This is an alternative definition of
                             % \ip that is commented out.

\begin{document}             % End of preamble and beginning of text.

\maketitle                   % Produces the title.

This is an example input file.  Comparing it with
the output it generates can show you how to
produce a simple document of your own.

\section{Ordinary Text}      % Produces section heading.  Lower-level
                             % sections are begun with similar
                             % \subsection and \subsubsection commands.

The ends  of words and sentences are marked
  by   spaces. It  doesn't matter how many
spaces    you type; one is as good as 100.  The
end of   a line counts as a space.

One   or more   blank lines denote the  end
of  a paragraph.

Since any number of consecutive spaces are treated
like a single one, the formatting of the input
file makes no difference to
      \LaTeX,                % The \LaTeX command generates the LaTeX logo.
but it makes a difference to you.  When you use
\LaTeX, making your input file as easy to read
as possible will be a great help as you write
your document and when you change it.  This sample
file shows how you can add comments to your own input
file.

Because printing is different from typewriting,
there are a number of things that you have to do
differently when preparing an input file than if
you were just typing the document directly.
Quotation marks like
       ``this''
have to be handled specially, as do quotes within
quotes:
       ``\,`this'            % \, separates the double and single quote.
        is what I just
        wrote, not  `that'\,''.

Dashes come in three sizes: an
       intra-word
dash, a medium dash for number ranges like
       1--2,
and a punctuation
       dash---like
this.

A sentence-ending space should be larger than the
space between words within a sentence.  You
sometimes have to type special commands in
conjunction with punctuation characters to get
this right, as in the following sentence.
       Gnats, gnus, etc.\ all  % `\ ' makes an inter-word space.
       begin with G\@.         % \@ marks end-of-sentence punctuation.
You should check the spaces after periods when
reading your output to make sure you haven't
forgotten any special cases.  Generating an
ellipsis
       \ldots\               % `\ ' is needed after `\ldots' because TeX
                             % ignores spaces after command names like \ldots
                             % made from \ + letters.
                             %
                             % Note how a `%' character causes TeX to ignore
                             % the end of the input line, so these blank lines
                             % do not start a new paragraph.
                             %
with the right spacing around the periods requires
a special command.

\LaTeX\ interprets some common characters as
commands, so you must type special commands to
generate them.  These characters include the
following:
       \$ \& \% \# \{ and \}.

In printing, text is usually emphasized with an
       \emph{italic}
type style.

\begin{em}
   A long segment of text can also be emphasized
   in this way.  Text within such a segment can be
   given \emph{additional} emphasis.
\end{em}

It is sometimes necessary to prevent \LaTeX\ from
breaking a line where it might otherwise do so.
This may be at a space, as between the ``Mr.''\ and
``Jones'' in
       ``Mr.~Jones'',        % ~ produces an unbreakable interword space.
or within a word---especially when the word is a
symbol like
       \mbox{\emph{itemnum}}
that makes little sense when hyphenated across
lines.

Footnotes\footnote{This is an example of a footnote.}
pose no problem.

\LaTeX\ is good at typesetting mathematical formulas
like
       \( x-3y + z = 7 \)
or
       \( a_{1} > x^{2n} + y^{2n} > x' \)
or
       \( \ip{A}{B} = \sum_{i} a_{i} b_{i} \).
The spaces you type in a formula are
ignored.  Remember that a letter like
       $x$                   % $ ... $  and  \( ... \)  are equivalent
is a formula when it denotes a mathematical
symbol, and it should be typed as one.

\section{Displayed Text}

Text is displayed by indenting it from the left
margin.  Quotations are commonly displayed.  There
are short quotations
\begin{quote}
   This is a short quotation.  It consists of a
   single paragraph of text.  See how it is formatted.
\end{quote}
and longer ones.
\begin{quotation}
   This is a longer quotation.  It consists of two
   paragraphs of text, neither of which are
   particularly interesting.

   This is the second paragraph of the quotation.  It
   is just as dull as the first paragraph.
\end{quotation}
Another frequently-displayed structure is a list.
The following is an example of an \emph{itemized}
list.
\begin{itemize}
   \item This is the first item of an itemized list.
         Each item in the list is marked with a ``tick''.
         You don't have to worry about what kind of tick
         mark is used.

   \item This is the second item of the list.  It
         contains another list nested inside it.  The inner
         list is an \emph{enumerated} list.
         \begin{enumerate}
            \item This is the first item of an enumerated
                  list that is nested within the itemized list.

            \item This is the second item of the inner list.
                  \LaTeX\ allows you to nest lists deeper than
                  you really should.
         \end{enumerate}
         This is the rest of the second item of the outer
         list.  It is no more interesting than any other
         part of the item.
   \item This is the third item of the list.
\end{itemize}
You can even display poetry.
\begin{verse}
   There is an environment
    for verse \\             % The \\ command separates lines
   Whose features some poets % within a stanza.
   will curse.

                             % One or more blank lines separate stanzas.

   For instead of making\\
   Them do \emph{all} line breaking, \\
   It allows them to put too many words on a line when they'd rather be
   forced to be terse.
\end{verse}

Mathematical formulas may also be displayed.  A
displayed formula
is
one-line long; multiline
formulas require special formatting instructions.
   \[  \ip{\Gamma}{\psi'} = x'' + y^{2} + z_{i}^{n}\]
Don't start a paragraph with a displayed equation,
nor make one a paragraph by itself.

\end{document}               % End of document.

\section{Bench Huge Section 1055}
% synthetic bench section 1055
% This is a sample LaTeX input file.  (Version of 12 August 2004.)
%
% A '%' character causes TeX to ignore all remaining text on the line,
% and is used for comments like this one.

\documentclass{article}      % Specifies the document class

                             % The preamble begins here.
\title{An Example Document}  % Declares the document's title.
\author{Leslie Lamport}      % Declares the author's name.
\date{January 21, 1994}      % Deleting this command produces today's date.

\newcommand{\ip}[2]{(#1, #2)}
                             % Defines \ip{arg1}{arg2} to mean
                             % (arg1, arg2).

%\newcommand{\ip}[2]{\langle #1 | #2\rangle}
                             % This is an alternative definition of
                             % \ip that is commented out.

\begin{document}             % End of preamble and beginning of text.

\maketitle                   % Produces the title.

This is an example input file.  Comparing it with
the output it generates can show you how to
produce a simple document of your own.

\section{Ordinary Text}      % Produces section heading.  Lower-level
                             % sections are begun with similar
                             % \subsection and \subsubsection commands.

The ends  of words and sentences are marked
  by   spaces. It  doesn't matter how many
spaces    you type; one is as good as 100.  The
end of   a line counts as a space.

One   or more   blank lines denote the  end
of  a paragraph.

Since any number of consecutive spaces are treated
like a single one, the formatting of the input
file makes no difference to
      \LaTeX,                % The \LaTeX command generates the LaTeX logo.
but it makes a difference to you.  When you use
\LaTeX, making your input file as easy to read
as possible will be a great help as you write
your document and when you change it.  This sample
file shows how you can add comments to your own input
file.

Because printing is different from typewriting,
there are a number of things that you have to do
differently when preparing an input file than if
you were just typing the document directly.
Quotation marks like
       ``this''
have to be handled specially, as do quotes within
quotes:
       ``\,`this'            % \, separates the double and single quote.
        is what I just
        wrote, not  `that'\,''.

Dashes come in three sizes: an
       intra-word
dash, a medium dash for number ranges like
       1--2,
and a punctuation
       dash---like
this.

A sentence-ending space should be larger than the
space between words within a sentence.  You
sometimes have to type special commands in
conjunction with punctuation characters to get
this right, as in the following sentence.
       Gnats, gnus, etc.\ all  % `\ ' makes an inter-word space.
       begin with G\@.         % \@ marks end-of-sentence punctuation.
You should check the spaces after periods when
reading your output to make sure you haven't
forgotten any special cases.  Generating an
ellipsis
       \ldots\               % `\ ' is needed after `\ldots' because TeX
                             % ignores spaces after command names like \ldots
                             % made from \ + letters.
                             %
                             % Note how a `%' character causes TeX to ignore
                             % the end of the input line, so these blank lines
                             % do not start a new paragraph.
                             %
with the right spacing around the periods requires
a special command.

\LaTeX\ interprets some common characters as
commands, so you must type special commands to
generate them.  These characters include the
following:
       \$ \& \% \# \{ and \}.

In printing, text is usually emphasized with an
       \emph{italic}
type style.

\begin{em}
   A long segment of text can also be emphasized
   in this way.  Text within such a segment can be
   given \emph{additional} emphasis.
\end{em}

It is sometimes necessary to prevent \LaTeX\ from
breaking a line where it might otherwise do so.
This may be at a space, as between the ``Mr.''\ and
``Jones'' in
       ``Mr.~Jones'',        % ~ produces an unbreakable interword space.
or within a word---especially when the word is a
symbol like
       \mbox{\emph{itemnum}}
that makes little sense when hyphenated across
lines.

Footnotes\footnote{This is an example of a footnote.}
pose no problem.

\LaTeX\ is good at typesetting mathematical formulas
like
       \( x-3y + z = 7 \)
or
       \( a_{1} > x^{2n} + y^{2n} > x' \)
or
       \( \ip{A}{B} = \sum_{i} a_{i} b_{i} \).
The spaces you type in a formula are
ignored.  Remember that a letter like
       $x$                   % $ ... $  and  \( ... \)  are equivalent
is a formula when it denotes a mathematical
symbol, and it should be typed as one.

\section{Displayed Text}

Text is displayed by indenting it from the left
margin.  Quotations are commonly displayed.  There
are short quotations
\begin{quote}
   This is a short quotation.  It consists of a
   single paragraph of text.  See how it is formatted.
\end{quote}
and longer ones.
\begin{quotation}
   This is a longer quotation.  It consists of two
   paragraphs of text, neither of which are
   particularly interesting.

   This is the second paragraph of the quotation.  It
   is just as dull as the first paragraph.
\end{quotation}
Another frequently-displayed structure is a list.
The following is an example of an \emph{itemized}
list.
\begin{itemize}
   \item This is the first item of an itemized list.
         Each item in the list is marked with a ``tick''.
         You don't have to worry about what kind of tick
         mark is used.

   \item This is the second item of the list.  It
         contains another list nested inside it.  The inner
         list is an \emph{enumerated} list.
         \begin{enumerate}
            \item This is the first item of an enumerated
                  list that is nested within the itemized list.

            \item This is the second item of the inner list.
                  \LaTeX\ allows you to nest lists deeper than
                  you really should.
         \end{enumerate}
         This is the rest of the second item of the outer
         list.  It is no more interesting than any other
         part of the item.
   \item This is the third item of the list.
\end{itemize}
You can even display poetry.
\begin{verse}
   There is an environment
    for verse \\             % The \\ command separates lines
   Whose features some poets % within a stanza.
   will curse.

                             % One or more blank lines separate stanzas.

   For instead of making\\
   Them do \emph{all} line breaking, \\
   It allows them to put too many words on a line when they'd rather be
   forced to be terse.
\end{verse}

Mathematical formulas may also be displayed.  A
displayed formula
is
one-line long; multiline
formulas require special formatting instructions.
   \[  \ip{\Gamma}{\psi'} = x'' + y^{2} + z_{i}^{n}\]
Don't start a paragraph with a displayed equation,
nor make one a paragraph by itself.

\end{document}               % End of document.

\section{Bench Huge Section 1056}
% synthetic bench section 1056
% This is a sample LaTeX input file.  (Version of 12 August 2004.)
%
% A '%' character causes TeX to ignore all remaining text on the line,
% and is used for comments like this one.

\documentclass{article}      % Specifies the document class

                             % The preamble begins here.
\title{An Example Document}  % Declares the document's title.
\author{Leslie Lamport}      % Declares the author's name.
\date{January 21, 1994}      % Deleting this command produces today's date.

\newcommand{\ip}[2]{(#1, #2)}
                             % Defines \ip{arg1}{arg2} to mean
                             % (arg1, arg2).

%\newcommand{\ip}[2]{\langle #1 | #2\rangle}
                             % This is an alternative definition of
                             % \ip that is commented out.

\begin{document}             % End of preamble and beginning of text.

\maketitle                   % Produces the title.

This is an example input file.  Comparing it with
the output it generates can show you how to
produce a simple document of your own.

\section{Ordinary Text}      % Produces section heading.  Lower-level
                             % sections are begun with similar
                             % \subsection and \subsubsection commands.

The ends  of words and sentences are marked
  by   spaces. It  doesn't matter how many
spaces    you type; one is as good as 100.  The
end of   a line counts as a space.

One   or more   blank lines denote the  end
of  a paragraph.

Since any number of consecutive spaces are treated
like a single one, the formatting of the input
file makes no difference to
      \LaTeX,                % The \LaTeX command generates the LaTeX logo.
but it makes a difference to you.  When you use
\LaTeX, making your input file as easy to read
as possible will be a great help as you write
your document and when you change it.  This sample
file shows how you can add comments to your own input
file.

Because printing is different from typewriting,
there are a number of things that you have to do
differently when preparing an input file than if
you were just typing the document directly.
Quotation marks like
       ``this''
have to be handled specially, as do quotes within
quotes:
       ``\,`this'            % \, separates the double and single quote.
        is what I just
        wrote, not  `that'\,''.

Dashes come in three sizes: an
       intra-word
dash, a medium dash for number ranges like
       1--2,
and a punctuation
       dash---like
this.

A sentence-ending space should be larger than the
space between words within a sentence.  You
sometimes have to type special commands in
conjunction with punctuation characters to get
this right, as in the following sentence.
       Gnats, gnus, etc.\ all  % `\ ' makes an inter-word space.
       begin with G\@.         % \@ marks end-of-sentence punctuation.
You should check the spaces after periods when
reading your output to make sure you haven't
forgotten any special cases.  Generating an
ellipsis
       \ldots\               % `\ ' is needed after `\ldots' because TeX
                             % ignores spaces after command names like \ldots
                             % made from \ + letters.
                             %
                             % Note how a `%' character causes TeX to ignore
                             % the end of the input line, so these blank lines
                             % do not start a new paragraph.
                             %
with the right spacing around the periods requires
a special command.

\LaTeX\ interprets some common characters as
commands, so you must type special commands to
generate them.  These characters include the
following:
       \$ \& \% \# \{ and \}.

In printing, text is usually emphasized with an
       \emph{italic}
type style.

\begin{em}
   A long segment of text can also be emphasized
   in this way.  Text within such a segment can be
   given \emph{additional} emphasis.
\end{em}

It is sometimes necessary to prevent \LaTeX\ from
breaking a line where it might otherwise do so.
This may be at a space, as between the ``Mr.''\ and
``Jones'' in
       ``Mr.~Jones'',        % ~ produces an unbreakable interword space.
or within a word---especially when the word is a
symbol like
       \mbox{\emph{itemnum}}
that makes little sense when hyphenated across
lines.

Footnotes\footnote{This is an example of a footnote.}
pose no problem.

\LaTeX\ is good at typesetting mathematical formulas
like
       \( x-3y + z = 7 \)
or
       \( a_{1} > x^{2n} + y^{2n} > x' \)
or
       \( \ip{A}{B} = \sum_{i} a_{i} b_{i} \).
The spaces you type in a formula are
ignored.  Remember that a letter like
       $x$                   % $ ... $  and  \( ... \)  are equivalent
is a formula when it denotes a mathematical
symbol, and it should be typed as one.

\section{Displayed Text}

Text is displayed by indenting it from the left
margin.  Quotations are commonly displayed.  There
are short quotations
\begin{quote}
   This is a short quotation.  It consists of a
   single paragraph of text.  See how it is formatted.
\end{quote}
and longer ones.
\begin{quotation}
   This is a longer quotation.  It consists of two
   paragraphs of text, neither of which are
   particularly interesting.

   This is the second paragraph of the quotation.  It
   is just as dull as the first paragraph.
\end{quotation}
Another frequently-displayed structure is a list.
The following is an example of an \emph{itemized}
list.
\begin{itemize}
   \item This is the first item of an itemized list.
         Each item in the list is marked with a ``tick''.
         You don't have to worry about what kind of tick
         mark is used.

   \item This is the second item of the list.  It
         contains another list nested inside it.  The inner
         list is an \emph{enumerated} list.
         \begin{enumerate}
            \item This is the first item of an enumerated
                  list that is nested within the itemized list.

            \item This is the second item of the inner list.
                  \LaTeX\ allows you to nest lists deeper than
                  you really should.
         \end{enumerate}
         This is the rest of the second item of the outer
         list.  It is no more interesting than any other
         part of the item.
   \item This is the third item of the list.
\end{itemize}
You can even display poetry.
\begin{verse}
   There is an environment
    for verse \\             % The \\ command separates lines
   Whose features some poets % within a stanza.
   will curse.

                             % One or more blank lines separate stanzas.

   For instead of making\\
   Them do \emph{all} line breaking, \\
   It allows them to put too many words on a line when they'd rather be
   forced to be terse.
\end{verse}

Mathematical formulas may also be displayed.  A
displayed formula
is
one-line long; multiline
formulas require special formatting instructions.
   \[  \ip{\Gamma}{\psi'} = x'' + y^{2} + z_{i}^{n}\]
Don't start a paragraph with a displayed equation,
nor make one a paragraph by itself.

\end{document}               % End of document.

\section{Bench Huge Section 1057}
% synthetic bench section 1057
% This is a sample LaTeX input file.  (Version of 12 August 2004.)
%
% A '%' character causes TeX to ignore all remaining text on the line,
% and is used for comments like this one.

\documentclass{article}      % Specifies the document class

                             % The preamble begins here.
\title{An Example Document}  % Declares the document's title.
\author{Leslie Lamport}      % Declares the author's name.
\date{January 21, 1994}      % Deleting this command produces today's date.

\newcommand{\ip}[2]{(#1, #2)}
                             % Defines \ip{arg1}{arg2} to mean
                             % (arg1, arg2).

%\newcommand{\ip}[2]{\langle #1 | #2\rangle}
                             % This is an alternative definition of
                             % \ip that is commented out.

\begin{document}             % End of preamble and beginning of text.

\maketitle                   % Produces the title.

This is an example input file.  Comparing it with
the output it generates can show you how to
produce a simple document of your own.

\section{Ordinary Text}      % Produces section heading.  Lower-level
                             % sections are begun with similar
                             % \subsection and \subsubsection commands.

The ends  of words and sentences are marked
  by   spaces. It  doesn't matter how many
spaces    you type; one is as good as 100.  The
end of   a line counts as a space.

One   or more   blank lines denote the  end
of  a paragraph.

Since any number of consecutive spaces are treated
like a single one, the formatting of the input
file makes no difference to
      \LaTeX,                % The \LaTeX command generates the LaTeX logo.
but it makes a difference to you.  When you use
\LaTeX, making your input file as easy to read
as possible will be a great help as you write
your document and when you change it.  This sample
file shows how you can add comments to your own input
file.

Because printing is different from typewriting,
there are a number of things that you have to do
differently when preparing an input file than if
you were just typing the document directly.
Quotation marks like
       ``this''
have to be handled specially, as do quotes within
quotes:
       ``\,`this'            % \, separates the double and single quote.
        is what I just
        wrote, not  `that'\,''.

Dashes come in three sizes: an
       intra-word
dash, a medium dash for number ranges like
       1--2,
and a punctuation
       dash---like
this.

A sentence-ending space should be larger than the
space between words within a sentence.  You
sometimes have to type special commands in
conjunction with punctuation characters to get
this right, as in the following sentence.
       Gnats, gnus, etc.\ all  % `\ ' makes an inter-word space.
       begin with G\@.         % \@ marks end-of-sentence punctuation.
You should check the spaces after periods when
reading your output to make sure you haven't
forgotten any special cases.  Generating an
ellipsis
       \ldots\               % `\ ' is needed after `\ldots' because TeX
                             % ignores spaces after command names like \ldots
                             % made from \ + letters.
                             %
                             % Note how a `%' character causes TeX to ignore
                             % the end of the input line, so these blank lines
                             % do not start a new paragraph.
                             %
with the right spacing around the periods requires
a special command.

\LaTeX\ interprets some common characters as
commands, so you must type special commands to
generate them.  These characters include the
following:
       \$ \& \% \# \{ and \}.

In printing, text is usually emphasized with an
       \emph{italic}
type style.

\begin{em}
   A long segment of text can also be emphasized
   in this way.  Text within such a segment can be
   given \emph{additional} emphasis.
\end{em}

It is sometimes necessary to prevent \LaTeX\ from
breaking a line where it might otherwise do so.
This may be at a space, as between the ``Mr.''\ and
``Jones'' in
       ``Mr.~Jones'',        % ~ produces an unbreakable interword space.
or within a word---especially when the word is a
symbol like
       \mbox{\emph{itemnum}}
that makes little sense when hyphenated across
lines.

Footnotes\footnote{This is an example of a footnote.}
pose no problem.

\LaTeX\ is good at typesetting mathematical formulas
like
       \( x-3y + z = 7 \)
or
       \( a_{1} > x^{2n} + y^{2n} > x' \)
or
       \( \ip{A}{B} = \sum_{i} a_{i} b_{i} \).
The spaces you type in a formula are
ignored.  Remember that a letter like
       $x$                   % $ ... $  and  \( ... \)  are equivalent
is a formula when it denotes a mathematical
symbol, and it should be typed as one.

\section{Displayed Text}

Text is displayed by indenting it from the left
margin.  Quotations are commonly displayed.  There
are short quotations
\begin{quote}
   This is a short quotation.  It consists of a
   single paragraph of text.  See how it is formatted.
\end{quote}
and longer ones.
\begin{quotation}
   This is a longer quotation.  It consists of two
   paragraphs of text, neither of which are
   particularly interesting.

   This is the second paragraph of the quotation.  It
   is just as dull as the first paragraph.
\end{quotation}
Another frequently-displayed structure is a list.
The following is an example of an \emph{itemized}
list.
\begin{itemize}
   \item This is the first item of an itemized list.
         Each item in the list is marked with a ``tick''.
         You don't have to worry about what kind of tick
         mark is used.

   \item This is the second item of the list.  It
         contains another list nested inside it.  The inner
         list is an \emph{enumerated} list.
         \begin{enumerate}
            \item This is the first item of an enumerated
                  list that is nested within the itemized list.

            \item This is the second item of the inner list.
                  \LaTeX\ allows you to nest lists deeper than
                  you really should.
         \end{enumerate}
         This is the rest of the second item of the outer
         list.  It is no more interesting than any other
         part of the item.
   \item This is the third item of the list.
\end{itemize}
You can even display poetry.
\begin{verse}
   There is an environment
    for verse \\             % The \\ command separates lines
   Whose features some poets % within a stanza.
   will curse.

                             % One or more blank lines separate stanzas.

   For instead of making\\
   Them do \emph{all} line breaking, \\
   It allows them to put too many words on a line when they'd rather be
   forced to be terse.
\end{verse}

Mathematical formulas may also be displayed.  A
displayed formula
is
one-line long; multiline
formulas require special formatting instructions.
   \[  \ip{\Gamma}{\psi'} = x'' + y^{2} + z_{i}^{n}\]
Don't start a paragraph with a displayed equation,
nor make one a paragraph by itself.

\end{document}               % End of document.

\section{Bench Huge Section 1058}
% synthetic bench section 1058
% This is a sample LaTeX input file.  (Version of 12 August 2004.)
%
% A '%' character causes TeX to ignore all remaining text on the line,
% and is used for comments like this one.

\documentclass{article}      % Specifies the document class

                             % The preamble begins here.
\title{An Example Document}  % Declares the document's title.
\author{Leslie Lamport}      % Declares the author's name.
\date{January 21, 1994}      % Deleting this command produces today's date.

\newcommand{\ip}[2]{(#1, #2)}
                             % Defines \ip{arg1}{arg2} to mean
                             % (arg1, arg2).

%\newcommand{\ip}[2]{\langle #1 | #2\rangle}
                             % This is an alternative definition of
                             % \ip that is commented out.

\begin{document}             % End of preamble and beginning of text.

\maketitle                   % Produces the title.

This is an example input file.  Comparing it with
the output it generates can show you how to
produce a simple document of your own.

\section{Ordinary Text}      % Produces section heading.  Lower-level
                             % sections are begun with similar
                             % \subsection and \subsubsection commands.

The ends  of words and sentences are marked
  by   spaces. It  doesn't matter how many
spaces    you type; one is as good as 100.  The
end of   a line counts as a space.

One   or more   blank lines denote the  end
of  a paragraph.

Since any number of consecutive spaces are treated
like a single one, the formatting of the input
file makes no difference to
      \LaTeX,                % The \LaTeX command generates the LaTeX logo.
but it makes a difference to you.  When you use
\LaTeX, making your input file as easy to read
as possible will be a great help as you write
your document and when you change it.  This sample
file shows how you can add comments to your own input
file.

Because printing is different from typewriting,
there are a number of things that you have to do
differently when preparing an input file than if
you were just typing the document directly.
Quotation marks like
       ``this''
have to be handled specially, as do quotes within
quotes:
       ``\,`this'            % \, separates the double and single quote.
        is what I just
        wrote, not  `that'\,''.

Dashes come in three sizes: an
       intra-word
dash, a medium dash for number ranges like
       1--2,
and a punctuation
       dash---like
this.

A sentence-ending space should be larger than the
space between words within a sentence.  You
sometimes have to type special commands in
conjunction with punctuation characters to get
this right, as in the following sentence.
       Gnats, gnus, etc.\ all  % `\ ' makes an inter-word space.
       begin with G\@.         % \@ marks end-of-sentence punctuation.
You should check the spaces after periods when
reading your output to make sure you haven't
forgotten any special cases.  Generating an
ellipsis
       \ldots\               % `\ ' is needed after `\ldots' because TeX
                             % ignores spaces after command names like \ldots
                             % made from \ + letters.
                             %
                             % Note how a `%' character causes TeX to ignore
                             % the end of the input line, so these blank lines
                             % do not start a new paragraph.
                             %
with the right spacing around the periods requires
a special command.

\LaTeX\ interprets some common characters as
commands, so you must type special commands to
generate them.  These characters include the
following:
       \$ \& \% \# \{ and \}.

In printing, text is usually emphasized with an
       \emph{italic}
type style.

\begin{em}
   A long segment of text can also be emphasized
   in this way.  Text within such a segment can be
   given \emph{additional} emphasis.
\end{em}

It is sometimes necessary to prevent \LaTeX\ from
breaking a line where it might otherwise do so.
This may be at a space, as between the ``Mr.''\ and
``Jones'' in
       ``Mr.~Jones'',        % ~ produces an unbreakable interword space.
or within a word---especially when the word is a
symbol like
       \mbox{\emph{itemnum}}
that makes little sense when hyphenated across
lines.

Footnotes\footnote{This is an example of a footnote.}
pose no problem.

\LaTeX\ is good at typesetting mathematical formulas
like
       \( x-3y + z = 7 \)
or
       \( a_{1} > x^{2n} + y^{2n} > x' \)
or
       \( \ip{A}{B} = \sum_{i} a_{i} b_{i} \).
The spaces you type in a formula are
ignored.  Remember that a letter like
       $x$                   % $ ... $  and  \( ... \)  are equivalent
is a formula when it denotes a mathematical
symbol, and it should be typed as one.

\section{Displayed Text}

Text is displayed by indenting it from the left
margin.  Quotations are commonly displayed.  There
are short quotations
\begin{quote}
   This is a short quotation.  It consists of a
   single paragraph of text.  See how it is formatted.
\end{quote}
and longer ones.
\begin{quotation}
   This is a longer quotation.  It consists of two
   paragraphs of text, neither of which are
   particularly interesting.

   This is the second paragraph of the quotation.  It
   is just as dull as the first paragraph.
\end{quotation}
Another frequently-displayed structure is a list.
The following is an example of an \emph{itemized}
list.
\begin{itemize}
   \item This is the first item of an itemized list.
         Each item in the list is marked with a ``tick''.
         You don't have to worry about what kind of tick
         mark is used.

   \item This is the second item of the list.  It
         contains another list nested inside it.  The inner
         list is an \emph{enumerated} list.
         \begin{enumerate}
            \item This is the first item of an enumerated
                  list that is nested within the itemized list.

            \item This is the second item of the inner list.
                  \LaTeX\ allows you to nest lists deeper than
                  you really should.
         \end{enumerate}
         This is the rest of the second item of the outer
         list.  It is no more interesting than any other
         part of the item.
   \item This is the third item of the list.
\end{itemize}
You can even display poetry.
\begin{verse}
   There is an environment
    for verse \\             % The \\ command separates lines
   Whose features some poets % within a stanza.
   will curse.

                             % One or more blank lines separate stanzas.

   For instead of making\\
   Them do \emph{all} line breaking, \\
   It allows them to put too many words on a line when they'd rather be
   forced to be terse.
\end{verse}

Mathematical formulas may also be displayed.  A
displayed formula
is
one-line long; multiline
formulas require special formatting instructions.
   \[  \ip{\Gamma}{\psi'} = x'' + y^{2} + z_{i}^{n}\]
Don't start a paragraph with a displayed equation,
nor make one a paragraph by itself.

\end{document}               % End of document.

\section{Bench Huge Section 1059}
% synthetic bench section 1059
% This is a sample LaTeX input file.  (Version of 12 August 2004.)
%
% A '%' character causes TeX to ignore all remaining text on the line,
% and is used for comments like this one.

\documentclass{article}      % Specifies the document class

                             % The preamble begins here.
\title{An Example Document}  % Declares the document's title.
\author{Leslie Lamport}      % Declares the author's name.
\date{January 21, 1994}      % Deleting this command produces today's date.

\newcommand{\ip}[2]{(#1, #2)}
                             % Defines \ip{arg1}{arg2} to mean
                             % (arg1, arg2).

%\newcommand{\ip}[2]{\langle #1 | #2\rangle}
                             % This is an alternative definition of
                             % \ip that is commented out.

\begin{document}             % End of preamble and beginning of text.

\maketitle                   % Produces the title.

This is an example input file.  Comparing it with
the output it generates can show you how to
produce a simple document of your own.

\section{Ordinary Text}      % Produces section heading.  Lower-level
                             % sections are begun with similar
                             % \subsection and \subsubsection commands.

The ends  of words and sentences are marked
  by   spaces. It  doesn't matter how many
spaces    you type; one is as good as 100.  The
end of   a line counts as a space.

One   or more   blank lines denote the  end
of  a paragraph.

Since any number of consecutive spaces are treated
like a single one, the formatting of the input
file makes no difference to
      \LaTeX,                % The \LaTeX command generates the LaTeX logo.
but it makes a difference to you.  When you use
\LaTeX, making your input file as easy to read
as possible will be a great help as you write
your document and when you change it.  This sample
file shows how you can add comments to your own input
file.

Because printing is different from typewriting,
there are a number of things that you have to do
differently when preparing an input file than if
you were just typing the document directly.
Quotation marks like
       ``this''
have to be handled specially, as do quotes within
quotes:
       ``\,`this'            % \, separates the double and single quote.
        is what I just
        wrote, not  `that'\,''.

Dashes come in three sizes: an
       intra-word
dash, a medium dash for number ranges like
       1--2,
and a punctuation
       dash---like
this.

A sentence-ending space should be larger than the
space between words within a sentence.  You
sometimes have to type special commands in
conjunction with punctuation characters to get
this right, as in the following sentence.
       Gnats, gnus, etc.\ all  % `\ ' makes an inter-word space.
       begin with G\@.         % \@ marks end-of-sentence punctuation.
You should check the spaces after periods when
reading your output to make sure you haven't
forgotten any special cases.  Generating an
ellipsis
       \ldots\               % `\ ' is needed after `\ldots' because TeX
                             % ignores spaces after command names like \ldots
                             % made from \ + letters.
                             %
                             % Note how a `%' character causes TeX to ignore
                             % the end of the input line, so these blank lines
                             % do not start a new paragraph.
                             %
with the right spacing around the periods requires
a special command.

\LaTeX\ interprets some common characters as
commands, so you must type special commands to
generate them.  These characters include the
following:
       \$ \& \% \# \{ and \}.

In printing, text is usually emphasized with an
       \emph{italic}
type style.

\begin{em}
   A long segment of text can also be emphasized
   in this way.  Text within such a segment can be
   given \emph{additional} emphasis.
\end{em}

It is sometimes necessary to prevent \LaTeX\ from
breaking a line where it might otherwise do so.
This may be at a space, as between the ``Mr.''\ and
``Jones'' in
       ``Mr.~Jones'',        % ~ produces an unbreakable interword space.
or within a word---especially when the word is a
symbol like
       \mbox{\emph{itemnum}}
that makes little sense when hyphenated across
lines.

Footnotes\footnote{This is an example of a footnote.}
pose no problem.

\LaTeX\ is good at typesetting mathematical formulas
like
       \( x-3y + z = 7 \)
or
       \( a_{1} > x^{2n} + y^{2n} > x' \)
or
       \( \ip{A}{B} = \sum_{i} a_{i} b_{i} \).
The spaces you type in a formula are
ignored.  Remember that a letter like
       $x$                   % $ ... $  and  \( ... \)  are equivalent
is a formula when it denotes a mathematical
symbol, and it should be typed as one.

\section{Displayed Text}

Text is displayed by indenting it from the left
margin.  Quotations are commonly displayed.  There
are short quotations
\begin{quote}
   This is a short quotation.  It consists of a
   single paragraph of text.  See how it is formatted.
\end{quote}
and longer ones.
\begin{quotation}
   This is a longer quotation.  It consists of two
   paragraphs of text, neither of which are
   particularly interesting.

   This is the second paragraph of the quotation.  It
   is just as dull as the first paragraph.
\end{quotation}
Another frequently-displayed structure is a list.
The following is an example of an \emph{itemized}
list.
\begin{itemize}
   \item This is the first item of an itemized list.
         Each item in the list is marked with a ``tick''.
         You don't have to worry about what kind of tick
         mark is used.

   \item This is the second item of the list.  It
         contains another list nested inside it.  The inner
         list is an \emph{enumerated} list.
         \begin{enumerate}
            \item This is the first item of an enumerated
                  list that is nested within the itemized list.

            \item This is the second item of the inner list.
                  \LaTeX\ allows you to nest lists deeper than
                  you really should.
         \end{enumerate}
         This is the rest of the second item of the outer
         list.  It is no more interesting than any other
         part of the item.
   \item This is the third item of the list.
\end{itemize}
You can even display poetry.
\begin{verse}
   There is an environment
    for verse \\             % The \\ command separates lines
   Whose features some poets % within a stanza.
   will curse.

                             % One or more blank lines separate stanzas.

   For instead of making\\
   Them do \emph{all} line breaking, \\
   It allows them to put too many words on a line when they'd rather be
   forced to be terse.
\end{verse}

Mathematical formulas may also be displayed.  A
displayed formula
is
one-line long; multiline
formulas require special formatting instructions.
   \[  \ip{\Gamma}{\psi'} = x'' + y^{2} + z_{i}^{n}\]
Don't start a paragraph with a displayed equation,
nor make one a paragraph by itself.

\end{document}               % End of document.

\section{Bench Huge Section 1060}
% synthetic bench section 1060
% This is a sample LaTeX input file.  (Version of 12 August 2004.)
%
% A '%' character causes TeX to ignore all remaining text on the line,
% and is used for comments like this one.

\documentclass{article}      % Specifies the document class

                             % The preamble begins here.
\title{An Example Document}  % Declares the document's title.
\author{Leslie Lamport}      % Declares the author's name.
\date{January 21, 1994}      % Deleting this command produces today's date.

\newcommand{\ip}[2]{(#1, #2)}
                             % Defines \ip{arg1}{arg2} to mean
                             % (arg1, arg2).

%\newcommand{\ip}[2]{\langle #1 | #2\rangle}
                             % This is an alternative definition of
                             % \ip that is commented out.

\begin{document}             % End of preamble and beginning of text.

\maketitle                   % Produces the title.

This is an example input file.  Comparing it with
the output it generates can show you how to
produce a simple document of your own.

\section{Ordinary Text}      % Produces section heading.  Lower-level
                             % sections are begun with similar
                             % \subsection and \subsubsection commands.

The ends  of words and sentences are marked
  by   spaces. It  doesn't matter how many
spaces    you type; one is as good as 100.  The
end of   a line counts as a space.

One   or more   blank lines denote the  end
of  a paragraph.

Since any number of consecutive spaces are treated
like a single one, the formatting of the input
file makes no difference to
      \LaTeX,                % The \LaTeX command generates the LaTeX logo.
but it makes a difference to you.  When you use
\LaTeX, making your input file as easy to read
as possible will be a great help as you write
your document and when you change it.  This sample
file shows how you can add comments to your own input
file.

Because printing is different from typewriting,
there are a number of things that you have to do
differently when preparing an input file than if
you were just typing the document directly.
Quotation marks like
       ``this''
have to be handled specially, as do quotes within
quotes:
       ``\,`this'            % \, separates the double and single quote.
        is what I just
        wrote, not  `that'\,''.

Dashes come in three sizes: an
       intra-word
dash, a medium dash for number ranges like
       1--2,
and a punctuation
       dash---like
this.

A sentence-ending space should be larger than the
space between words within a sentence.  You
sometimes have to type special commands in
conjunction with punctuation characters to get
this right, as in the following sentence.
       Gnats, gnus, etc.\ all  % `\ ' makes an inter-word space.
       begin with G\@.         % \@ marks end-of-sentence punctuation.
You should check the spaces after periods when
reading your output to make sure you haven't
forgotten any special cases.  Generating an
ellipsis
       \ldots\               % `\ ' is needed after `\ldots' because TeX
                             % ignores spaces after command names like \ldots
                             % made from \ + letters.
                             %
                             % Note how a `%' character causes TeX to ignore
                             % the end of the input line, so these blank lines
                             % do not start a new paragraph.
                             %
with the right spacing around the periods requires
a special command.

\LaTeX\ interprets some common characters as
commands, so you must type special commands to
generate them.  These characters include the
following:
       \$ \& \% \# \{ and \}.

In printing, text is usually emphasized with an
       \emph{italic}
type style.

\begin{em}
   A long segment of text can also be emphasized
   in this way.  Text within such a segment can be
   given \emph{additional} emphasis.
\end{em}

It is sometimes necessary to prevent \LaTeX\ from
breaking a line where it might otherwise do so.
This may be at a space, as between the ``Mr.''\ and
``Jones'' in
       ``Mr.~Jones'',        % ~ produces an unbreakable interword space.
or within a word---especially when the word is a
symbol like
       \mbox{\emph{itemnum}}
that makes little sense when hyphenated across
lines.

Footnotes\footnote{This is an example of a footnote.}
pose no problem.

\LaTeX\ is good at typesetting mathematical formulas
like
       \( x-3y + z = 7 \)
or
       \( a_{1} > x^{2n} + y^{2n} > x' \)
or
       \( \ip{A}{B} = \sum_{i} a_{i} b_{i} \).
The spaces you type in a formula are
ignored.  Remember that a letter like
       $x$                   % $ ... $  and  \( ... \)  are equivalent
is a formula when it denotes a mathematical
symbol, and it should be typed as one.

\section{Displayed Text}

Text is displayed by indenting it from the left
margin.  Quotations are commonly displayed.  There
are short quotations
\begin{quote}
   This is a short quotation.  It consists of a
   single paragraph of text.  See how it is formatted.
\end{quote}
and longer ones.
\begin{quotation}
   This is a longer quotation.  It consists of two
   paragraphs of text, neither of which are
   particularly interesting.

   This is the second paragraph of the quotation.  It
   is just as dull as the first paragraph.
\end{quotation}
Another frequently-displayed structure is a list.
The following is an example of an \emph{itemized}
list.
\begin{itemize}
   \item This is the first item of an itemized list.
         Each item in the list is marked with a ``tick''.
         You don't have to worry about what kind of tick
         mark is used.

   \item This is the second item of the list.  It
         contains another list nested inside it.  The inner
         list is an \emph{enumerated} list.
         \begin{enumerate}
            \item This is the first item of an enumerated
                  list that is nested within the itemized list.

            \item This is the second item of the inner list.
                  \LaTeX\ allows you to nest lists deeper than
                  you really should.
         \end{enumerate}
         This is the rest of the second item of the outer
         list.  It is no more interesting than any other
         part of the item.
   \item This is the third item of the list.
\end{itemize}
You can even display poetry.
\begin{verse}
   There is an environment
    for verse \\             % The \\ command separates lines
   Whose features some poets % within a stanza.
   will curse.

                             % One or more blank lines separate stanzas.

   For instead of making\\
   Them do \emph{all} line breaking, \\
   It allows them to put too many words on a line when they'd rather be
   forced to be terse.
\end{verse}

Mathematical formulas may also be displayed.  A
displayed formula
is
one-line long; multiline
formulas require special formatting instructions.
   \[  \ip{\Gamma}{\psi'} = x'' + y^{2} + z_{i}^{n}\]
Don't start a paragraph with a displayed equation,
nor make one a paragraph by itself.

\end{document}               % End of document.

\section{Bench Huge Section 1061}
% synthetic bench section 1061
% This is a sample LaTeX input file.  (Version of 12 August 2004.)
%
% A '%' character causes TeX to ignore all remaining text on the line,
% and is used for comments like this one.

\documentclass{article}      % Specifies the document class

                             % The preamble begins here.
\title{An Example Document}  % Declares the document's title.
\author{Leslie Lamport}      % Declares the author's name.
\date{January 21, 1994}      % Deleting this command produces today's date.

\newcommand{\ip}[2]{(#1, #2)}
                             % Defines \ip{arg1}{arg2} to mean
                             % (arg1, arg2).

%\newcommand{\ip}[2]{\langle #1 | #2\rangle}
                             % This is an alternative definition of
                             % \ip that is commented out.

\begin{document}             % End of preamble and beginning of text.

\maketitle                   % Produces the title.

This is an example input file.  Comparing it with
the output it generates can show you how to
produce a simple document of your own.

\section{Ordinary Text}      % Produces section heading.  Lower-level
                             % sections are begun with similar
                             % \subsection and \subsubsection commands.

The ends  of words and sentences are marked
  by   spaces. It  doesn't matter how many
spaces    you type; one is as good as 100.  The
end of   a line counts as a space.

One   or more   blank lines denote the  end
of  a paragraph.

Since any number of consecutive spaces are treated
like a single one, the formatting of the input
file makes no difference to
      \LaTeX,                % The \LaTeX command generates the LaTeX logo.
but it makes a difference to you.  When you use
\LaTeX, making your input file as easy to read
as possible will be a great help as you write
your document and when you change it.  This sample
file shows how you can add comments to your own input
file.

Because printing is different from typewriting,
there are a number of things that you have to do
differently when preparing an input file than if
you were just typing the document directly.
Quotation marks like
       ``this''
have to be handled specially, as do quotes within
quotes:
       ``\,`this'            % \, separates the double and single quote.
        is what I just
        wrote, not  `that'\,''.

Dashes come in three sizes: an
       intra-word
dash, a medium dash for number ranges like
       1--2,
and a punctuation
       dash---like
this.

A sentence-ending space should be larger than the
space between words within a sentence.  You
sometimes have to type special commands in
conjunction with punctuation characters to get
this right, as in the following sentence.
       Gnats, gnus, etc.\ all  % `\ ' makes an inter-word space.
       begin with G\@.         % \@ marks end-of-sentence punctuation.
You should check the spaces after periods when
reading your output to make sure you haven't
forgotten any special cases.  Generating an
ellipsis
       \ldots\               % `\ ' is needed after `\ldots' because TeX
                             % ignores spaces after command names like \ldots
                             % made from \ + letters.
                             %
                             % Note how a `%' character causes TeX to ignore
                             % the end of the input line, so these blank lines
                             % do not start a new paragraph.
                             %
with the right spacing around the periods requires
a special command.

\LaTeX\ interprets some common characters as
commands, so you must type special commands to
generate them.  These characters include the
following:
       \$ \& \% \# \{ and \}.

In printing, text is usually emphasized with an
       \emph{italic}
type style.

\begin{em}
   A long segment of text can also be emphasized
   in this way.  Text within such a segment can be
   given \emph{additional} emphasis.
\end{em}

It is sometimes necessary to prevent \LaTeX\ from
breaking a line where it might otherwise do so.
This may be at a space, as between the ``Mr.''\ and
``Jones'' in
       ``Mr.~Jones'',        % ~ produces an unbreakable interword space.
or within a word---especially when the word is a
symbol like
       \mbox{\emph{itemnum}}
that makes little sense when hyphenated across
lines.

Footnotes\footnote{This is an example of a footnote.}
pose no problem.

\LaTeX\ is good at typesetting mathematical formulas
like
       \( x-3y + z = 7 \)
or
       \( a_{1} > x^{2n} + y^{2n} > x' \)
or
       \( \ip{A}{B} = \sum_{i} a_{i} b_{i} \).
The spaces you type in a formula are
ignored.  Remember that a letter like
       $x$                   % $ ... $  and  \( ... \)  are equivalent
is a formula when it denotes a mathematical
symbol, and it should be typed as one.

\section{Displayed Text}

Text is displayed by indenting it from the left
margin.  Quotations are commonly displayed.  There
are short quotations
\begin{quote}
   This is a short quotation.  It consists of a
   single paragraph of text.  See how it is formatted.
\end{quote}
and longer ones.
\begin{quotation}
   This is a longer quotation.  It consists of two
   paragraphs of text, neither of which are
   particularly interesting.

   This is the second paragraph of the quotation.  It
   is just as dull as the first paragraph.
\end{quotation}
Another frequently-displayed structure is a list.
The following is an example of an \emph{itemized}
list.
\begin{itemize}
   \item This is the first item of an itemized list.
         Each item in the list is marked with a ``tick''.
         You don't have to worry about what kind of tick
         mark is used.

   \item This is the second item of the list.  It
         contains another list nested inside it.  The inner
         list is an \emph{enumerated} list.
         \begin{enumerate}
            \item This is the first item of an enumerated
                  list that is nested within the itemized list.

            \item This is the second item of the inner list.
                  \LaTeX\ allows you to nest lists deeper than
                  you really should.
         \end{enumerate}
         This is the rest of the second item of the outer
         list.  It is no more interesting than any other
         part of the item.
   \item This is the third item of the list.
\end{itemize}
You can even display poetry.
\begin{verse}
   There is an environment
    for verse \\             % The \\ command separates lines
   Whose features some poets % within a stanza.
   will curse.

                             % One or more blank lines separate stanzas.

   For instead of making\\
   Them do \emph{all} line breaking, \\
   It allows them to put too many words on a line when they'd rather be
   forced to be terse.
\end{verse}

Mathematical formulas may also be displayed.  A
displayed formula
is
one-line long; multiline
formulas require special formatting instructions.
   \[  \ip{\Gamma}{\psi'} = x'' + y^{2} + z_{i}^{n}\]
Don't start a paragraph with a displayed equation,
nor make one a paragraph by itself.

\end{document}               % End of document.

\section{Bench Huge Section 1062}
% synthetic bench section 1062
% This is a sample LaTeX input file.  (Version of 12 August 2004.)
%
% A '%' character causes TeX to ignore all remaining text on the line,
% and is used for comments like this one.

\documentclass{article}      % Specifies the document class

                             % The preamble begins here.
\title{An Example Document}  % Declares the document's title.
\author{Leslie Lamport}      % Declares the author's name.
\date{January 21, 1994}      % Deleting this command produces today's date.

\newcommand{\ip}[2]{(#1, #2)}
                             % Defines \ip{arg1}{arg2} to mean
                             % (arg1, arg2).

%\newcommand{\ip}[2]{\langle #1 | #2\rangle}
                             % This is an alternative definition of
                             % \ip that is commented out.

\begin{document}             % End of preamble and beginning of text.

\maketitle                   % Produces the title.

This is an example input file.  Comparing it with
the output it generates can show you how to
produce a simple document of your own.

\section{Ordinary Text}      % Produces section heading.  Lower-level
                             % sections are begun with similar
                             % \subsection and \subsubsection commands.

The ends  of words and sentences are marked
  by   spaces. It  doesn't matter how many
spaces    you type; one is as good as 100.  The
end of   a line counts as a space.

One   or more   blank lines denote the  end
of  a paragraph.

Since any number of consecutive spaces are treated
like a single one, the formatting of the input
file makes no difference to
      \LaTeX,                % The \LaTeX command generates the LaTeX logo.
but it makes a difference to you.  When you use
\LaTeX, making your input file as easy to read
as possible will be a great help as you write
your document and when you change it.  This sample
file shows how you can add comments to your own input
file.

Because printing is different from typewriting,
there are a number of things that you have to do
differently when preparing an input file than if
you were just typing the document directly.
Quotation marks like
       ``this''
have to be handled specially, as do quotes within
quotes:
       ``\,`this'            % \, separates the double and single quote.
        is what I just
        wrote, not  `that'\,''.

Dashes come in three sizes: an
       intra-word
dash, a medium dash for number ranges like
       1--2,
and a punctuation
       dash---like
this.

A sentence-ending space should be larger than the
space between words within a sentence.  You
sometimes have to type special commands in
conjunction with punctuation characters to get
this right, as in the following sentence.
       Gnats, gnus, etc.\ all  % `\ ' makes an inter-word space.
       begin with G\@.         % \@ marks end-of-sentence punctuation.
You should check the spaces after periods when
reading your output to make sure you haven't
forgotten any special cases.  Generating an
ellipsis
       \ldots\               % `\ ' is needed after `\ldots' because TeX
                             % ignores spaces after command names like \ldots
                             % made from \ + letters.
                             %
                             % Note how a `%' character causes TeX to ignore
                             % the end of the input line, so these blank lines
                             % do not start a new paragraph.
                             %
with the right spacing around the periods requires
a special command.

\LaTeX\ interprets some common characters as
commands, so you must type special commands to
generate them.  These characters include the
following:
       \$ \& \% \# \{ and \}.

In printing, text is usually emphasized with an
       \emph{italic}
type style.

\begin{em}
   A long segment of text can also be emphasized
   in this way.  Text within such a segment can be
   given \emph{additional} emphasis.
\end{em}

It is sometimes necessary to prevent \LaTeX\ from
breaking a line where it might otherwise do so.
This may be at a space, as between the ``Mr.''\ and
``Jones'' in
       ``Mr.~Jones'',        % ~ produces an unbreakable interword space.
or within a word---especially when the word is a
symbol like
       \mbox{\emph{itemnum}}
that makes little sense when hyphenated across
lines.

Footnotes\footnote{This is an example of a footnote.}
pose no problem.

\LaTeX\ is good at typesetting mathematical formulas
like
       \( x-3y + z = 7 \)
or
       \( a_{1} > x^{2n} + y^{2n} > x' \)
or
       \( \ip{A}{B} = \sum_{i} a_{i} b_{i} \).
The spaces you type in a formula are
ignored.  Remember that a letter like
       $x$                   % $ ... $  and  \( ... \)  are equivalent
is a formula when it denotes a mathematical
symbol, and it should be typed as one.

\section{Displayed Text}

Text is displayed by indenting it from the left
margin.  Quotations are commonly displayed.  There
are short quotations
\begin{quote}
   This is a short quotation.  It consists of a
   single paragraph of text.  See how it is formatted.
\end{quote}
and longer ones.
\begin{quotation}
   This is a longer quotation.  It consists of two
   paragraphs of text, neither of which are
   particularly interesting.

   This is the second paragraph of the quotation.  It
   is just as dull as the first paragraph.
\end{quotation}
Another frequently-displayed structure is a list.
The following is an example of an \emph{itemized}
list.
\begin{itemize}
   \item This is the first item of an itemized list.
         Each item in the list is marked with a ``tick''.
         You don't have to worry about what kind of tick
         mark is used.

   \item This is the second item of the list.  It
         contains another list nested inside it.  The inner
         list is an \emph{enumerated} list.
         \begin{enumerate}
            \item This is the first item of an enumerated
                  list that is nested within the itemized list.

            \item This is the second item of the inner list.
                  \LaTeX\ allows you to nest lists deeper than
                  you really should.
         \end{enumerate}
         This is the rest of the second item of the outer
         list.  It is no more interesting than any other
         part of the item.
   \item This is the third item of the list.
\end{itemize}
You can even display poetry.
\begin{verse}
   There is an environment
    for verse \\             % The \\ command separates lines
   Whose features some poets % within a stanza.
   will curse.

                             % One or more blank lines separate stanzas.

   For instead of making\\
   Them do \emph{all} line breaking, \\
   It allows them to put too many words on a line when they'd rather be
   forced to be terse.
\end{verse}

Mathematical formulas may also be displayed.  A
displayed formula
is
one-line long; multiline
formulas require special formatting instructions.
   \[  \ip{\Gamma}{\psi'} = x'' + y^{2} + z_{i}^{n}\]
Don't start a paragraph with a displayed equation,
nor make one a paragraph by itself.

\end{document}               % End of document.

\section{Bench Huge Section 1063}
% synthetic bench section 1063
% This is a sample LaTeX input file.  (Version of 12 August 2004.)
%
% A '%' character causes TeX to ignore all remaining text on the line,
% and is used for comments like this one.

\documentclass{article}      % Specifies the document class

                             % The preamble begins here.
\title{An Example Document}  % Declares the document's title.
\author{Leslie Lamport}      % Declares the author's name.
\date{January 21, 1994}      % Deleting this command produces today's date.

\newcommand{\ip}[2]{(#1, #2)}
                             % Defines \ip{arg1}{arg2} to mean
                             % (arg1, arg2).

%\newcommand{\ip}[2]{\langle #1 | #2\rangle}
                             % This is an alternative definition of
                             % \ip that is commented out.

\begin{document}             % End of preamble and beginning of text.

\maketitle                   % Produces the title.

This is an example input file.  Comparing it with
the output it generates can show you how to
produce a simple document of your own.

\section{Ordinary Text}      % Produces section heading.  Lower-level
                             % sections are begun with similar
                             % \subsection and \subsubsection commands.

The ends  of words and sentences are marked
  by   spaces. It  doesn't matter how many
spaces    you type; one is as good as 100.  The
end of   a line counts as a space.

One   or more   blank lines denote the  end
of  a paragraph.

Since any number of consecutive spaces are treated
like a single one, the formatting of the input
file makes no difference to
      \LaTeX,                % The \LaTeX command generates the LaTeX logo.
but it makes a difference to you.  When you use
\LaTeX, making your input file as easy to read
as possible will be a great help as you write
your document and when you change it.  This sample
file shows how you can add comments to your own input
file.

Because printing is different from typewriting,
there are a number of things that you have to do
differently when preparing an input file than if
you were just typing the document directly.
Quotation marks like
       ``this''
have to be handled specially, as do quotes within
quotes:
       ``\,`this'            % \, separates the double and single quote.
        is what I just
        wrote, not  `that'\,''.

Dashes come in three sizes: an
       intra-word
dash, a medium dash for number ranges like
       1--2,
and a punctuation
       dash---like
this.

A sentence-ending space should be larger than the
space between words within a sentence.  You
sometimes have to type special commands in
conjunction with punctuation characters to get
this right, as in the following sentence.
       Gnats, gnus, etc.\ all  % `\ ' makes an inter-word space.
       begin with G\@.         % \@ marks end-of-sentence punctuation.
You should check the spaces after periods when
reading your output to make sure you haven't
forgotten any special cases.  Generating an
ellipsis
       \ldots\               % `\ ' is needed after `\ldots' because TeX
                             % ignores spaces after command names like \ldots
                             % made from \ + letters.
                             %
                             % Note how a `%' character causes TeX to ignore
                             % the end of the input line, so these blank lines
                             % do not start a new paragraph.
                             %
with the right spacing around the periods requires
a special command.

\LaTeX\ interprets some common characters as
commands, so you must type special commands to
generate them.  These characters include the
following:
       \$ \& \% \# \{ and \}.

In printing, text is usually emphasized with an
       \emph{italic}
type style.

\begin{em}
   A long segment of text can also be emphasized
   in this way.  Text within such a segment can be
   given \emph{additional} emphasis.
\end{em}

It is sometimes necessary to prevent \LaTeX\ from
breaking a line where it might otherwise do so.
This may be at a space, as between the ``Mr.''\ and
``Jones'' in
       ``Mr.~Jones'',        % ~ produces an unbreakable interword space.
or within a word---especially when the word is a
symbol like
       \mbox{\emph{itemnum}}
that makes little sense when hyphenated across
lines.

Footnotes\footnote{This is an example of a footnote.}
pose no problem.

\LaTeX\ is good at typesetting mathematical formulas
like
       \( x-3y + z = 7 \)
or
       \( a_{1} > x^{2n} + y^{2n} > x' \)
or
       \( \ip{A}{B} = \sum_{i} a_{i} b_{i} \).
The spaces you type in a formula are
ignored.  Remember that a letter like
       $x$                   % $ ... $  and  \( ... \)  are equivalent
is a formula when it denotes a mathematical
symbol, and it should be typed as one.

\section{Displayed Text}

Text is displayed by indenting it from the left
margin.  Quotations are commonly displayed.  There
are short quotations
\begin{quote}
   This is a short quotation.  It consists of a
   single paragraph of text.  See how it is formatted.
\end{quote}
and longer ones.
\begin{quotation}
   This is a longer quotation.  It consists of two
   paragraphs of text, neither of which are
   particularly interesting.

   This is the second paragraph of the quotation.  It
   is just as dull as the first paragraph.
\end{quotation}
Another frequently-displayed structure is a list.
The following is an example of an \emph{itemized}
list.
\begin{itemize}
   \item This is the first item of an itemized list.
         Each item in the list is marked with a ``tick''.
         You don't have to worry about what kind of tick
         mark is used.

   \item This is the second item of the list.  It
         contains another list nested inside it.  The inner
         list is an \emph{enumerated} list.
         \begin{enumerate}
            \item This is the first item of an enumerated
                  list that is nested within the itemized list.

            \item This is the second item of the inner list.
                  \LaTeX\ allows you to nest lists deeper than
                  you really should.
         \end{enumerate}
         This is the rest of the second item of the outer
         list.  It is no more interesting than any other
         part of the item.
   \item This is the third item of the list.
\end{itemize}
You can even display poetry.
\begin{verse}
   There is an environment
    for verse \\             % The \\ command separates lines
   Whose features some poets % within a stanza.
   will curse.

                             % One or more blank lines separate stanzas.

   For instead of making\\
   Them do \emph{all} line breaking, \\
   It allows them to put too many words on a line when they'd rather be
   forced to be terse.
\end{verse}

Mathematical formulas may also be displayed.  A
displayed formula
is
one-line long; multiline
formulas require special formatting instructions.
   \[  \ip{\Gamma}{\psi'} = x'' + y^{2} + z_{i}^{n}\]
Don't start a paragraph with a displayed equation,
nor make one a paragraph by itself.

\end{document}               % End of document.

\section{Bench Huge Section 1064}
% synthetic bench section 1064
% This is a sample LaTeX input file.  (Version of 12 August 2004.)
%
% A '%' character causes TeX to ignore all remaining text on the line,
% and is used for comments like this one.

\documentclass{article}      % Specifies the document class

                             % The preamble begins here.
\title{An Example Document}  % Declares the document's title.
\author{Leslie Lamport}      % Declares the author's name.
\date{January 21, 1994}      % Deleting this command produces today's date.

\newcommand{\ip}[2]{(#1, #2)}
                             % Defines \ip{arg1}{arg2} to mean
                             % (arg1, arg2).

%\newcommand{\ip}[2]{\langle #1 | #2\rangle}
                             % This is an alternative definition of
                             % \ip that is commented out.

\begin{document}             % End of preamble and beginning of text.

\maketitle                   % Produces the title.

This is an example input file.  Comparing it with
the output it generates can show you how to
produce a simple document of your own.

\section{Ordinary Text}      % Produces section heading.  Lower-level
                             % sections are begun with similar
                             % \subsection and \subsubsection commands.

The ends  of words and sentences are marked
  by   spaces. It  doesn't matter how many
spaces    you type; one is as good as 100.  The
end of   a line counts as a space.

One   or more   blank lines denote the  end
of  a paragraph.

Since any number of consecutive spaces are treated
like a single one, the formatting of the input
file makes no difference to
      \LaTeX,                % The \LaTeX command generates the LaTeX logo.
but it makes a difference to you.  When you use
\LaTeX, making your input file as easy to read
as possible will be a great help as you write
your document and when you change it.  This sample
file shows how you can add comments to your own input
file.

Because printing is different from typewriting,
there are a number of things that you have to do
differently when preparing an input file than if
you were just typing the document directly.
Quotation marks like
       ``this''
have to be handled specially, as do quotes within
quotes:
       ``\,`this'            % \, separates the double and single quote.
        is what I just
        wrote, not  `that'\,''.

Dashes come in three sizes: an
       intra-word
dash, a medium dash for number ranges like
       1--2,
and a punctuation
       dash---like
this.

A sentence-ending space should be larger than the
space between words within a sentence.  You
sometimes have to type special commands in
conjunction with punctuation characters to get
this right, as in the following sentence.
       Gnats, gnus, etc.\ all  % `\ ' makes an inter-word space.
       begin with G\@.         % \@ marks end-of-sentence punctuation.
You should check the spaces after periods when
reading your output to make sure you haven't
forgotten any special cases.  Generating an
ellipsis
       \ldots\               % `\ ' is needed after `\ldots' because TeX
                             % ignores spaces after command names like \ldots
                             % made from \ + letters.
                             %
                             % Note how a `%' character causes TeX to ignore
                             % the end of the input line, so these blank lines
                             % do not start a new paragraph.
                             %
with the right spacing around the periods requires
a special command.

\LaTeX\ interprets some common characters as
commands, so you must type special commands to
generate them.  These characters include the
following:
       \$ \& \% \# \{ and \}.

In printing, text is usually emphasized with an
       \emph{italic}
type style.

\begin{em}
   A long segment of text can also be emphasized
   in this way.  Text within such a segment can be
   given \emph{additional} emphasis.
\end{em}

It is sometimes necessary to prevent \LaTeX\ from
breaking a line where it might otherwise do so.
This may be at a space, as between the ``Mr.''\ and
``Jones'' in
       ``Mr.~Jones'',        % ~ produces an unbreakable interword space.
or within a word---especially when the word is a
symbol like
       \mbox{\emph{itemnum}}
that makes little sense when hyphenated across
lines.

Footnotes\footnote{This is an example of a footnote.}
pose no problem.

\LaTeX\ is good at typesetting mathematical formulas
like
       \( x-3y + z = 7 \)
or
       \( a_{1} > x^{2n} + y^{2n} > x' \)
or
       \( \ip{A}{B} = \sum_{i} a_{i} b_{i} \).
The spaces you type in a formula are
ignored.  Remember that a letter like
       $x$                   % $ ... $  and  \( ... \)  are equivalent
is a formula when it denotes a mathematical
symbol, and it should be typed as one.

\section{Displayed Text}

Text is displayed by indenting it from the left
margin.  Quotations are commonly displayed.  There
are short quotations
\begin{quote}
   This is a short quotation.  It consists of a
   single paragraph of text.  See how it is formatted.
\end{quote}
and longer ones.
\begin{quotation}
   This is a longer quotation.  It consists of two
   paragraphs of text, neither of which are
   particularly interesting.

   This is the second paragraph of the quotation.  It
   is just as dull as the first paragraph.
\end{quotation}
Another frequently-displayed structure is a list.
The following is an example of an \emph{itemized}
list.
\begin{itemize}
   \item This is the first item of an itemized list.
         Each item in the list is marked with a ``tick''.
         You don't have to worry about what kind of tick
         mark is used.

   \item This is the second item of the list.  It
         contains another list nested inside it.  The inner
         list is an \emph{enumerated} list.
         \begin{enumerate}
            \item This is the first item of an enumerated
                  list that is nested within the itemized list.

            \item This is the second item of the inner list.
                  \LaTeX\ allows you to nest lists deeper than
                  you really should.
         \end{enumerate}
         This is the rest of the second item of the outer
         list.  It is no more interesting than any other
         part of the item.
   \item This is the third item of the list.
\end{itemize}
You can even display poetry.
\begin{verse}
   There is an environment
    for verse \\             % The \\ command separates lines
   Whose features some poets % within a stanza.
   will curse.

                             % One or more blank lines separate stanzas.

   For instead of making\\
   Them do \emph{all} line breaking, \\
   It allows them to put too many words on a line when they'd rather be
   forced to be terse.
\end{verse}

Mathematical formulas may also be displayed.  A
displayed formula
is
one-line long; multiline
formulas require special formatting instructions.
   \[  \ip{\Gamma}{\psi'} = x'' + y^{2} + z_{i}^{n}\]
Don't start a paragraph with a displayed equation,
nor make one a paragraph by itself.

\end{document}               % End of document.

\section{Bench Huge Section 1065}
% synthetic bench section 1065
% This is a sample LaTeX input file.  (Version of 12 August 2004.)
%
% A '%' character causes TeX to ignore all remaining text on the line,
% and is used for comments like this one.

\documentclass{article}      % Specifies the document class

                             % The preamble begins here.
\title{An Example Document}  % Declares the document's title.
\author{Leslie Lamport}      % Declares the author's name.
\date{January 21, 1994}      % Deleting this command produces today's date.

\newcommand{\ip}[2]{(#1, #2)}
                             % Defines \ip{arg1}{arg2} to mean
                             % (arg1, arg2).

%\newcommand{\ip}[2]{\langle #1 | #2\rangle}
                             % This is an alternative definition of
                             % \ip that is commented out.

\begin{document}             % End of preamble and beginning of text.

\maketitle                   % Produces the title.

This is an example input file.  Comparing it with
the output it generates can show you how to
produce a simple document of your own.

\section{Ordinary Text}      % Produces section heading.  Lower-level
                             % sections are begun with similar
                             % \subsection and \subsubsection commands.

The ends  of words and sentences are marked
  by   spaces. It  doesn't matter how many
spaces    you type; one is as good as 100.  The
end of   a line counts as a space.

One   or more   blank lines denote the  end
of  a paragraph.

Since any number of consecutive spaces are treated
like a single one, the formatting of the input
file makes no difference to
      \LaTeX,                % The \LaTeX command generates the LaTeX logo.
but it makes a difference to you.  When you use
\LaTeX, making your input file as easy to read
as possible will be a great help as you write
your document and when you change it.  This sample
file shows how you can add comments to your own input
file.

Because printing is different from typewriting,
there are a number of things that you have to do
differently when preparing an input file than if
you were just typing the document directly.
Quotation marks like
       ``this''
have to be handled specially, as do quotes within
quotes:
       ``\,`this'            % \, separates the double and single quote.
        is what I just
        wrote, not  `that'\,''.

Dashes come in three sizes: an
       intra-word
dash, a medium dash for number ranges like
       1--2,
and a punctuation
       dash---like
this.

A sentence-ending space should be larger than the
space between words within a sentence.  You
sometimes have to type special commands in
conjunction with punctuation characters to get
this right, as in the following sentence.
       Gnats, gnus, etc.\ all  % `\ ' makes an inter-word space.
       begin with G\@.         % \@ marks end-of-sentence punctuation.
You should check the spaces after periods when
reading your output to make sure you haven't
forgotten any special cases.  Generating an
ellipsis
       \ldots\               % `\ ' is needed after `\ldots' because TeX
                             % ignores spaces after command names like \ldots
                             % made from \ + letters.
                             %
                             % Note how a `%' character causes TeX to ignore
                             % the end of the input line, so these blank lines
                             % do not start a new paragraph.
                             %
with the right spacing around the periods requires
a special command.

\LaTeX\ interprets some common characters as
commands, so you must type special commands to
generate them.  These characters include the
following:
       \$ \& \% \# \{ and \}.

In printing, text is usually emphasized with an
       \emph{italic}
type style.

\begin{em}
   A long segment of text can also be emphasized
   in this way.  Text within such a segment can be
   given \emph{additional} emphasis.
\end{em}

It is sometimes necessary to prevent \LaTeX\ from
breaking a line where it might otherwise do so.
This may be at a space, as between the ``Mr.''\ and
``Jones'' in
       ``Mr.~Jones'',        % ~ produces an unbreakable interword space.
or within a word---especially when the word is a
symbol like
       \mbox{\emph{itemnum}}
that makes little sense when hyphenated across
lines.

Footnotes\footnote{This is an example of a footnote.}
pose no problem.

\LaTeX\ is good at typesetting mathematical formulas
like
       \( x-3y + z = 7 \)
or
       \( a_{1} > x^{2n} + y^{2n} > x' \)
or
       \( \ip{A}{B} = \sum_{i} a_{i} b_{i} \).
The spaces you type in a formula are
ignored.  Remember that a letter like
       $x$                   % $ ... $  and  \( ... \)  are equivalent
is a formula when it denotes a mathematical
symbol, and it should be typed as one.

\section{Displayed Text}

Text is displayed by indenting it from the left
margin.  Quotations are commonly displayed.  There
are short quotations
\begin{quote}
   This is a short quotation.  It consists of a
   single paragraph of text.  See how it is formatted.
\end{quote}
and longer ones.
\begin{quotation}
   This is a longer quotation.  It consists of two
   paragraphs of text, neither of which are
   particularly interesting.

   This is the second paragraph of the quotation.  It
   is just as dull as the first paragraph.
\end{quotation}
Another frequently-displayed structure is a list.
The following is an example of an \emph{itemized}
list.
\begin{itemize}
   \item This is the first item of an itemized list.
         Each item in the list is marked with a ``tick''.
         You don't have to worry about what kind of tick
         mark is used.

   \item This is the second item of the list.  It
         contains another list nested inside it.  The inner
         list is an \emph{enumerated} list.
         \begin{enumerate}
            \item This is the first item of an enumerated
                  list that is nested within the itemized list.

            \item This is the second item of the inner list.
                  \LaTeX\ allows you to nest lists deeper than
                  you really should.
         \end{enumerate}
         This is the rest of the second item of the outer
         list.  It is no more interesting than any other
         part of the item.
   \item This is the third item of the list.
\end{itemize}
You can even display poetry.
\begin{verse}
   There is an environment
    for verse \\             % The \\ command separates lines
   Whose features some poets % within a stanza.
   will curse.

                             % One or more blank lines separate stanzas.

   For instead of making\\
   Them do \emph{all} line breaking, \\
   It allows them to put too many words on a line when they'd rather be
   forced to be terse.
\end{verse}

Mathematical formulas may also be displayed.  A
displayed formula
is
one-line long; multiline
formulas require special formatting instructions.
   \[  \ip{\Gamma}{\psi'} = x'' + y^{2} + z_{i}^{n}\]
Don't start a paragraph with a displayed equation,
nor make one a paragraph by itself.

\end{document}               % End of document.

\section{Bench Huge Section 1066}
% synthetic bench section 1066
% This is a sample LaTeX input file.  (Version of 12 August 2004.)
%
% A '%' character causes TeX to ignore all remaining text on the line,
% and is used for comments like this one.

\documentclass{article}      % Specifies the document class

                             % The preamble begins here.
\title{An Example Document}  % Declares the document's title.
\author{Leslie Lamport}      % Declares the author's name.
\date{January 21, 1994}      % Deleting this command produces today's date.

\newcommand{\ip}[2]{(#1, #2)}
                             % Defines \ip{arg1}{arg2} to mean
                             % (arg1, arg2).

%\newcommand{\ip}[2]{\langle #1 | #2\rangle}
                             % This is an alternative definition of
                             % \ip that is commented out.

\begin{document}             % End of preamble and beginning of text.

\maketitle                   % Produces the title.

This is an example input file.  Comparing it with
the output it generates can show you how to
produce a simple document of your own.

\section{Ordinary Text}      % Produces section heading.  Lower-level
                             % sections are begun with similar
                             % \subsection and \subsubsection commands.

The ends  of words and sentences are marked
  by   spaces. It  doesn't matter how many
spaces    you type; one is as good as 100.  The
end of   a line counts as a space.

One   or more   blank lines denote the  end
of  a paragraph.

Since any number of consecutive spaces are treated
like a single one, the formatting of the input
file makes no difference to
      \LaTeX,                % The \LaTeX command generates the LaTeX logo.
but it makes a difference to you.  When you use
\LaTeX, making your input file as easy to read
as possible will be a great help as you write
your document and when you change it.  This sample
file shows how you can add comments to your own input
file.

Because printing is different from typewriting,
there are a number of things that you have to do
differently when preparing an input file than if
you were just typing the document directly.
Quotation marks like
       ``this''
have to be handled specially, as do quotes within
quotes:
       ``\,`this'            % \, separates the double and single quote.
        is what I just
        wrote, not  `that'\,''.

Dashes come in three sizes: an
       intra-word
dash, a medium dash for number ranges like
       1--2,
and a punctuation
       dash---like
this.

A sentence-ending space should be larger than the
space between words within a sentence.  You
sometimes have to type special commands in
conjunction with punctuation characters to get
this right, as in the following sentence.
       Gnats, gnus, etc.\ all  % `\ ' makes an inter-word space.
       begin with G\@.         % \@ marks end-of-sentence punctuation.
You should check the spaces after periods when
reading your output to make sure you haven't
forgotten any special cases.  Generating an
ellipsis
       \ldots\               % `\ ' is needed after `\ldots' because TeX
                             % ignores spaces after command names like \ldots
                             % made from \ + letters.
                             %
                             % Note how a `%' character causes TeX to ignore
                             % the end of the input line, so these blank lines
                             % do not start a new paragraph.
                             %
with the right spacing around the periods requires
a special command.

\LaTeX\ interprets some common characters as
commands, so you must type special commands to
generate them.  These characters include the
following:
       \$ \& \% \# \{ and \}.

In printing, text is usually emphasized with an
       \emph{italic}
type style.

\begin{em}
   A long segment of text can also be emphasized
   in this way.  Text within such a segment can be
   given \emph{additional} emphasis.
\end{em}

It is sometimes necessary to prevent \LaTeX\ from
breaking a line where it might otherwise do so.
This may be at a space, as between the ``Mr.''\ and
``Jones'' in
       ``Mr.~Jones'',        % ~ produces an unbreakable interword space.
or within a word---especially when the word is a
symbol like
       \mbox{\emph{itemnum}}
that makes little sense when hyphenated across
lines.

Footnotes\footnote{This is an example of a footnote.}
pose no problem.

\LaTeX\ is good at typesetting mathematical formulas
like
       \( x-3y + z = 7 \)
or
       \( a_{1} > x^{2n} + y^{2n} > x' \)
or
       \( \ip{A}{B} = \sum_{i} a_{i} b_{i} \).
The spaces you type in a formula are
ignored.  Remember that a letter like
       $x$                   % $ ... $  and  \( ... \)  are equivalent
is a formula when it denotes a mathematical
symbol, and it should be typed as one.

\section{Displayed Text}

Text is displayed by indenting it from the left
margin.  Quotations are commonly displayed.  There
are short quotations
\begin{quote}
   This is a short quotation.  It consists of a
   single paragraph of text.  See how it is formatted.
\end{quote}
and longer ones.
\begin{quotation}
   This is a longer quotation.  It consists of two
   paragraphs of text, neither of which are
   particularly interesting.

   This is the second paragraph of the quotation.  It
   is just as dull as the first paragraph.
\end{quotation}
Another frequently-displayed structure is a list.
The following is an example of an \emph{itemized}
list.
\begin{itemize}
   \item This is the first item of an itemized list.
         Each item in the list is marked with a ``tick''.
         You don't have to worry about what kind of tick
         mark is used.

   \item This is the second item of the list.  It
         contains another list nested inside it.  The inner
         list is an \emph{enumerated} list.
         \begin{enumerate}
            \item This is the first item of an enumerated
                  list that is nested within the itemized list.

            \item This is the second item of the inner list.
                  \LaTeX\ allows you to nest lists deeper than
                  you really should.
         \end{enumerate}
         This is the rest of the second item of the outer
         list.  It is no more interesting than any other
         part of the item.
   \item This is the third item of the list.
\end{itemize}
You can even display poetry.
\begin{verse}
   There is an environment
    for verse \\             % The \\ command separates lines
   Whose features some poets % within a stanza.
   will curse.

                             % One or more blank lines separate stanzas.

   For instead of making\\
   Them do \emph{all} line breaking, \\
   It allows them to put too many words on a line when they'd rather be
   forced to be terse.
\end{verse}

Mathematical formulas may also be displayed.  A
displayed formula
is
one-line long; multiline
formulas require special formatting instructions.
   \[  \ip{\Gamma}{\psi'} = x'' + y^{2} + z_{i}^{n}\]
Don't start a paragraph with a displayed equation,
nor make one a paragraph by itself.

\end{document}               % End of document.

\section{Bench Huge Section 1067}
% synthetic bench section 1067
% This is a sample LaTeX input file.  (Version of 12 August 2004.)
%
% A '%' character causes TeX to ignore all remaining text on the line,
% and is used for comments like this one.

\documentclass{article}      % Specifies the document class

                             % The preamble begins here.
\title{An Example Document}  % Declares the document's title.
\author{Leslie Lamport}      % Declares the author's name.
\date{January 21, 1994}      % Deleting this command produces today's date.

\newcommand{\ip}[2]{(#1, #2)}
                             % Defines \ip{arg1}{arg2} to mean
                             % (arg1, arg2).

%\newcommand{\ip}[2]{\langle #1 | #2\rangle}
                             % This is an alternative definition of
                             % \ip that is commented out.

\begin{document}             % End of preamble and beginning of text.

\maketitle                   % Produces the title.

This is an example input file.  Comparing it with
the output it generates can show you how to
produce a simple document of your own.

\section{Ordinary Text}      % Produces section heading.  Lower-level
                             % sections are begun with similar
                             % \subsection and \subsubsection commands.

The ends  of words and sentences are marked
  by   spaces. It  doesn't matter how many
spaces    you type; one is as good as 100.  The
end of   a line counts as a space.

One   or more   blank lines denote the  end
of  a paragraph.

Since any number of consecutive spaces are treated
like a single one, the formatting of the input
file makes no difference to
      \LaTeX,                % The \LaTeX command generates the LaTeX logo.
but it makes a difference to you.  When you use
\LaTeX, making your input file as easy to read
as possible will be a great help as you write
your document and when you change it.  This sample
file shows how you can add comments to your own input
file.

Because printing is different from typewriting,
there are a number of things that you have to do
differently when preparing an input file than if
you were just typing the document directly.
Quotation marks like
       ``this''
have to be handled specially, as do quotes within
quotes:
       ``\,`this'            % \, separates the double and single quote.
        is what I just
        wrote, not  `that'\,''.

Dashes come in three sizes: an
       intra-word
dash, a medium dash for number ranges like
       1--2,
and a punctuation
       dash---like
this.

A sentence-ending space should be larger than the
space between words within a sentence.  You
sometimes have to type special commands in
conjunction with punctuation characters to get
this right, as in the following sentence.
       Gnats, gnus, etc.\ all  % `\ ' makes an inter-word space.
       begin with G\@.         % \@ marks end-of-sentence punctuation.
You should check the spaces after periods when
reading your output to make sure you haven't
forgotten any special cases.  Generating an
ellipsis
       \ldots\               % `\ ' is needed after `\ldots' because TeX
                             % ignores spaces after command names like \ldots
                             % made from \ + letters.
                             %
                             % Note how a `%' character causes TeX to ignore
                             % the end of the input line, so these blank lines
                             % do not start a new paragraph.
                             %
with the right spacing around the periods requires
a special command.

\LaTeX\ interprets some common characters as
commands, so you must type special commands to
generate them.  These characters include the
following:
       \$ \& \% \# \{ and \}.

In printing, text is usually emphasized with an
       \emph{italic}
type style.

\begin{em}
   A long segment of text can also be emphasized
   in this way.  Text within such a segment can be
   given \emph{additional} emphasis.
\end{em}

It is sometimes necessary to prevent \LaTeX\ from
breaking a line where it might otherwise do so.
This may be at a space, as between the ``Mr.''\ and
``Jones'' in
       ``Mr.~Jones'',        % ~ produces an unbreakable interword space.
or within a word---especially when the word is a
symbol like
       \mbox{\emph{itemnum}}
that makes little sense when hyphenated across
lines.

Footnotes\footnote{This is an example of a footnote.}
pose no problem.

\LaTeX\ is good at typesetting mathematical formulas
like
       \( x-3y + z = 7 \)
or
       \( a_{1} > x^{2n} + y^{2n} > x' \)
or
       \( \ip{A}{B} = \sum_{i} a_{i} b_{i} \).
The spaces you type in a formula are
ignored.  Remember that a letter like
       $x$                   % $ ... $  and  \( ... \)  are equivalent
is a formula when it denotes a mathematical
symbol, and it should be typed as one.

\section{Displayed Text}

Text is displayed by indenting it from the left
margin.  Quotations are commonly displayed.  There
are short quotations
\begin{quote}
   This is a short quotation.  It consists of a
   single paragraph of text.  See how it is formatted.
\end{quote}
and longer ones.
\begin{quotation}
   This is a longer quotation.  It consists of two
   paragraphs of text, neither of which are
   particularly interesting.

   This is the second paragraph of the quotation.  It
   is just as dull as the first paragraph.
\end{quotation}
Another frequently-displayed structure is a list.
The following is an example of an \emph{itemized}
list.
\begin{itemize}
   \item This is the first item of an itemized list.
         Each item in the list is marked with a ``tick''.
         You don't have to worry about what kind of tick
         mark is used.

   \item This is the second item of the list.  It
         contains another list nested inside it.  The inner
         list is an \emph{enumerated} list.
         \begin{enumerate}
            \item This is the first item of an enumerated
                  list that is nested within the itemized list.

            \item This is the second item of the inner list.
                  \LaTeX\ allows you to nest lists deeper than
                  you really should.
         \end{enumerate}
         This is the rest of the second item of the outer
         list.  It is no more interesting than any other
         part of the item.
   \item This is the third item of the list.
\end{itemize}
You can even display poetry.
\begin{verse}
   There is an environment
    for verse \\             % The \\ command separates lines
   Whose features some poets % within a stanza.
   will curse.

                             % One or more blank lines separate stanzas.

   For instead of making\\
   Them do \emph{all} line breaking, \\
   It allows them to put too many words on a line when they'd rather be
   forced to be terse.
\end{verse}

Mathematical formulas may also be displayed.  A
displayed formula
is
one-line long; multiline
formulas require special formatting instructions.
   \[  \ip{\Gamma}{\psi'} = x'' + y^{2} + z_{i}^{n}\]
Don't start a paragraph with a displayed equation,
nor make one a paragraph by itself.

\end{document}               % End of document.

\section{Bench Huge Section 1068}
% synthetic bench section 1068
% This is a sample LaTeX input file.  (Version of 12 August 2004.)
%
% A '%' character causes TeX to ignore all remaining text on the line,
% and is used for comments like this one.

\documentclass{article}      % Specifies the document class

                             % The preamble begins here.
\title{An Example Document}  % Declares the document's title.
\author{Leslie Lamport}      % Declares the author's name.
\date{January 21, 1994}      % Deleting this command produces today's date.

\newcommand{\ip}[2]{(#1, #2)}
                             % Defines \ip{arg1}{arg2} to mean
                             % (arg1, arg2).

%\newcommand{\ip}[2]{\langle #1 | #2\rangle}
                             % This is an alternative definition of
                             % \ip that is commented out.

\begin{document}             % End of preamble and beginning of text.

\maketitle                   % Produces the title.

This is an example input file.  Comparing it with
the output it generates can show you how to
produce a simple document of your own.

\section{Ordinary Text}      % Produces section heading.  Lower-level
                             % sections are begun with similar
                             % \subsection and \subsubsection commands.

The ends  of words and sentences are marked
  by   spaces. It  doesn't matter how many
spaces    you type; one is as good as 100.  The
end of   a line counts as a space.

One   or more   blank lines denote the  end
of  a paragraph.

Since any number of consecutive spaces are treated
like a single one, the formatting of the input
file makes no difference to
      \LaTeX,                % The \LaTeX command generates the LaTeX logo.
but it makes a difference to you.  When you use
\LaTeX, making your input file as easy to read
as possible will be a great help as you write
your document and when you change it.  This sample
file shows how you can add comments to your own input
file.

Because printing is different from typewriting,
there are a number of things that you have to do
differently when preparing an input file than if
you were just typing the document directly.
Quotation marks like
       ``this''
have to be handled specially, as do quotes within
quotes:
       ``\,`this'            % \, separates the double and single quote.
        is what I just
        wrote, not  `that'\,''.

Dashes come in three sizes: an
       intra-word
dash, a medium dash for number ranges like
       1--2,
and a punctuation
       dash---like
this.

A sentence-ending space should be larger than the
space between words within a sentence.  You
sometimes have to type special commands in
conjunction with punctuation characters to get
this right, as in the following sentence.
       Gnats, gnus, etc.\ all  % `\ ' makes an inter-word space.
       begin with G\@.         % \@ marks end-of-sentence punctuation.
You should check the spaces after periods when
reading your output to make sure you haven't
forgotten any special cases.  Generating an
ellipsis
       \ldots\               % `\ ' is needed after `\ldots' because TeX
                             % ignores spaces after command names like \ldots
                             % made from \ + letters.
                             %
                             % Note how a `%' character causes TeX to ignore
                             % the end of the input line, so these blank lines
                             % do not start a new paragraph.
                             %
with the right spacing around the periods requires
a special command.

\LaTeX\ interprets some common characters as
commands, so you must type special commands to
generate them.  These characters include the
following:
       \$ \& \% \# \{ and \}.

In printing, text is usually emphasized with an
       \emph{italic}
type style.

\begin{em}
   A long segment of text can also be emphasized
   in this way.  Text within such a segment can be
   given \emph{additional} emphasis.
\end{em}

It is sometimes necessary to prevent \LaTeX\ from
breaking a line where it might otherwise do so.
This may be at a space, as between the ``Mr.''\ and
``Jones'' in
       ``Mr.~Jones'',        % ~ produces an unbreakable interword space.
or within a word---especially when the word is a
symbol like
       \mbox{\emph{itemnum}}
that makes little sense when hyphenated across
lines.

Footnotes\footnote{This is an example of a footnote.}
pose no problem.

\LaTeX\ is good at typesetting mathematical formulas
like
       \( x-3y + z = 7 \)
or
       \( a_{1} > x^{2n} + y^{2n} > x' \)
or
       \( \ip{A}{B} = \sum_{i} a_{i} b_{i} \).
The spaces you type in a formula are
ignored.  Remember that a letter like
       $x$                   % $ ... $  and  \( ... \)  are equivalent
is a formula when it denotes a mathematical
symbol, and it should be typed as one.

\section{Displayed Text}

Text is displayed by indenting it from the left
margin.  Quotations are commonly displayed.  There
are short quotations
\begin{quote}
   This is a short quotation.  It consists of a
   single paragraph of text.  See how it is formatted.
\end{quote}
and longer ones.
\begin{quotation}
   This is a longer quotation.  It consists of two
   paragraphs of text, neither of which are
   particularly interesting.

   This is the second paragraph of the quotation.  It
   is just as dull as the first paragraph.
\end{quotation}
Another frequently-displayed structure is a list.
The following is an example of an \emph{itemized}
list.
\begin{itemize}
   \item This is the first item of an itemized list.
         Each item in the list is marked with a ``tick''.
         You don't have to worry about what kind of tick
         mark is used.

   \item This is the second item of the list.  It
         contains another list nested inside it.  The inner
         list is an \emph{enumerated} list.
         \begin{enumerate}
            \item This is the first item of an enumerated
                  list that is nested within the itemized list.

            \item This is the second item of the inner list.
                  \LaTeX\ allows you to nest lists deeper than
                  you really should.
         \end{enumerate}
         This is the rest of the second item of the outer
         list.  It is no more interesting than any other
         part of the item.
   \item This is the third item of the list.
\end{itemize}
You can even display poetry.
\begin{verse}
   There is an environment
    for verse \\             % The \\ command separates lines
   Whose features some poets % within a stanza.
   will curse.

                             % One or more blank lines separate stanzas.

   For instead of making\\
   Them do \emph{all} line breaking, \\
   It allows them to put too many words on a line when they'd rather be
   forced to be terse.
\end{verse}

Mathematical formulas may also be displayed.  A
displayed formula
is
one-line long; multiline
formulas require special formatting instructions.
   \[  \ip{\Gamma}{\psi'} = x'' + y^{2} + z_{i}^{n}\]
Don't start a paragraph with a displayed equation,
nor make one a paragraph by itself.

\end{document}               % End of document.

\section{Bench Huge Section 1069}
% synthetic bench section 1069
% This is a sample LaTeX input file.  (Version of 12 August 2004.)
%
% A '%' character causes TeX to ignore all remaining text on the line,
% and is used for comments like this one.

\documentclass{article}      % Specifies the document class

                             % The preamble begins here.
\title{An Example Document}  % Declares the document's title.
\author{Leslie Lamport}      % Declares the author's name.
\date{January 21, 1994}      % Deleting this command produces today's date.

\newcommand{\ip}[2]{(#1, #2)}
                             % Defines \ip{arg1}{arg2} to mean
                             % (arg1, arg2).

%\newcommand{\ip}[2]{\langle #1 | #2\rangle}
                             % This is an alternative definition of
                             % \ip that is commented out.

\begin{document}             % End of preamble and beginning of text.

\maketitle                   % Produces the title.

This is an example input file.  Comparing it with
the output it generates can show you how to
produce a simple document of your own.

\section{Ordinary Text}      % Produces section heading.  Lower-level
                             % sections are begun with similar
                             % \subsection and \subsubsection commands.

The ends  of words and sentences are marked
  by   spaces. It  doesn't matter how many
spaces    you type; one is as good as 100.  The
end of   a line counts as a space.

One   or more   blank lines denote the  end
of  a paragraph.

Since any number of consecutive spaces are treated
like a single one, the formatting of the input
file makes no difference to
      \LaTeX,                % The \LaTeX command generates the LaTeX logo.
but it makes a difference to you.  When you use
\LaTeX, making your input file as easy to read
as possible will be a great help as you write
your document and when you change it.  This sample
file shows how you can add comments to your own input
file.

Because printing is different from typewriting,
there are a number of things that you have to do
differently when preparing an input file than if
you were just typing the document directly.
Quotation marks like
       ``this''
have to be handled specially, as do quotes within
quotes:
       ``\,`this'            % \, separates the double and single quote.
        is what I just
        wrote, not  `that'\,''.

Dashes come in three sizes: an
       intra-word
dash, a medium dash for number ranges like
       1--2,
and a punctuation
       dash---like
this.

A sentence-ending space should be larger than the
space between words within a sentence.  You
sometimes have to type special commands in
conjunction with punctuation characters to get
this right, as in the following sentence.
       Gnats, gnus, etc.\ all  % `\ ' makes an inter-word space.
       begin with G\@.         % \@ marks end-of-sentence punctuation.
You should check the spaces after periods when
reading your output to make sure you haven't
forgotten any special cases.  Generating an
ellipsis
       \ldots\               % `\ ' is needed after `\ldots' because TeX
                             % ignores spaces after command names like \ldots
                             % made from \ + letters.
                             %
                             % Note how a `%' character causes TeX to ignore
                             % the end of the input line, so these blank lines
                             % do not start a new paragraph.
                             %
with the right spacing around the periods requires
a special command.

\LaTeX\ interprets some common characters as
commands, so you must type special commands to
generate them.  These characters include the
following:
       \$ \& \% \# \{ and \}.

In printing, text is usually emphasized with an
       \emph{italic}
type style.

\begin{em}
   A long segment of text can also be emphasized
   in this way.  Text within such a segment can be
   given \emph{additional} emphasis.
\end{em}

It is sometimes necessary to prevent \LaTeX\ from
breaking a line where it might otherwise do so.
This may be at a space, as between the ``Mr.''\ and
``Jones'' in
       ``Mr.~Jones'',        % ~ produces an unbreakable interword space.
or within a word---especially when the word is a
symbol like
       \mbox{\emph{itemnum}}
that makes little sense when hyphenated across
lines.

Footnotes\footnote{This is an example of a footnote.}
pose no problem.

\LaTeX\ is good at typesetting mathematical formulas
like
       \( x-3y + z = 7 \)
or
       \( a_{1} > x^{2n} + y^{2n} > x' \)
or
       \( \ip{A}{B} = \sum_{i} a_{i} b_{i} \).
The spaces you type in a formula are
ignored.  Remember that a letter like
       $x$                   % $ ... $  and  \( ... \)  are equivalent
is a formula when it denotes a mathematical
symbol, and it should be typed as one.

\section{Displayed Text}

Text is displayed by indenting it from the left
margin.  Quotations are commonly displayed.  There
are short quotations
\begin{quote}
   This is a short quotation.  It consists of a
   single paragraph of text.  See how it is formatted.
\end{quote}
and longer ones.
\begin{quotation}
   This is a longer quotation.  It consists of two
   paragraphs of text, neither of which are
   particularly interesting.

   This is the second paragraph of the quotation.  It
   is just as dull as the first paragraph.
\end{quotation}
Another frequently-displayed structure is a list.
The following is an example of an \emph{itemized}
list.
\begin{itemize}
   \item This is the first item of an itemized list.
         Each item in the list is marked with a ``tick''.
         You don't have to worry about what kind of tick
         mark is used.

   \item This is the second item of the list.  It
         contains another list nested inside it.  The inner
         list is an \emph{enumerated} list.
         \begin{enumerate}
            \item This is the first item of an enumerated
                  list that is nested within the itemized list.

            \item This is the second item of the inner list.
                  \LaTeX\ allows you to nest lists deeper than
                  you really should.
         \end{enumerate}
         This is the rest of the second item of the outer
         list.  It is no more interesting than any other
         part of the item.
   \item This is the third item of the list.
\end{itemize}
You can even display poetry.
\begin{verse}
   There is an environment
    for verse \\             % The \\ command separates lines
   Whose features some poets % within a stanza.
   will curse.

                             % One or more blank lines separate stanzas.

   For instead of making\\
   Them do \emph{all} line breaking, \\
   It allows them to put too many words on a line when they'd rather be
   forced to be terse.
\end{verse}

Mathematical formulas may also be displayed.  A
displayed formula
is
one-line long; multiline
formulas require special formatting instructions.
   \[  \ip{\Gamma}{\psi'} = x'' + y^{2} + z_{i}^{n}\]
Don't start a paragraph with a displayed equation,
nor make one a paragraph by itself.

\end{document}               % End of document.

\section{Bench Huge Section 1070}
% synthetic bench section 1070
% This is a sample LaTeX input file.  (Version of 12 August 2004.)
%
% A '%' character causes TeX to ignore all remaining text on the line,
% and is used for comments like this one.

\documentclass{article}      % Specifies the document class

                             % The preamble begins here.
\title{An Example Document}  % Declares the document's title.
\author{Leslie Lamport}      % Declares the author's name.
\date{January 21, 1994}      % Deleting this command produces today's date.

\newcommand{\ip}[2]{(#1, #2)}
                             % Defines \ip{arg1}{arg2} to mean
                             % (arg1, arg2).

%\newcommand{\ip}[2]{\langle #1 | #2\rangle}
                             % This is an alternative definition of
                             % \ip that is commented out.

\begin{document}             % End of preamble and beginning of text.

\maketitle                   % Produces the title.

This is an example input file.  Comparing it with
the output it generates can show you how to
produce a simple document of your own.

\section{Ordinary Text}      % Produces section heading.  Lower-level
                             % sections are begun with similar
                             % \subsection and \subsubsection commands.

The ends  of words and sentences are marked
  by   spaces. It  doesn't matter how many
spaces    you type; one is as good as 100.  The
end of   a line counts as a space.

One   or more   blank lines denote the  end
of  a paragraph.

Since any number of consecutive spaces are treated
like a single one, the formatting of the input
file makes no difference to
      \LaTeX,                % The \LaTeX command generates the LaTeX logo.
but it makes a difference to you.  When you use
\LaTeX, making your input file as easy to read
as possible will be a great help as you write
your document and when you change it.  This sample
file shows how you can add comments to your own input
file.

Because printing is different from typewriting,
there are a number of things that you have to do
differently when preparing an input file than if
you were just typing the document directly.
Quotation marks like
       ``this''
have to be handled specially, as do quotes within
quotes:
       ``\,`this'            % \, separates the double and single quote.
        is what I just
        wrote, not  `that'\,''.

Dashes come in three sizes: an
       intra-word
dash, a medium dash for number ranges like
       1--2,
and a punctuation
       dash---like
this.

A sentence-ending space should be larger than the
space between words within a sentence.  You
sometimes have to type special commands in
conjunction with punctuation characters to get
this right, as in the following sentence.
       Gnats, gnus, etc.\ all  % `\ ' makes an inter-word space.
       begin with G\@.         % \@ marks end-of-sentence punctuation.
You should check the spaces after periods when
reading your output to make sure you haven't
forgotten any special cases.  Generating an
ellipsis
       \ldots\               % `\ ' is needed after `\ldots' because TeX
                             % ignores spaces after command names like \ldots
                             % made from \ + letters.
                             %
                             % Note how a `%' character causes TeX to ignore
                             % the end of the input line, so these blank lines
                             % do not start a new paragraph.
                             %
with the right spacing around the periods requires
a special command.

\LaTeX\ interprets some common characters as
commands, so you must type special commands to
generate them.  These characters include the
following:
       \$ \& \% \# \{ and \}.

In printing, text is usually emphasized with an
       \emph{italic}
type style.

\begin{em}
   A long segment of text can also be emphasized
   in this way.  Text within such a segment can be
   given \emph{additional} emphasis.
\end{em}

It is sometimes necessary to prevent \LaTeX\ from
breaking a line where it might otherwise do so.
This may be at a space, as between the ``Mr.''\ and
``Jones'' in
       ``Mr.~Jones'',        % ~ produces an unbreakable interword space.
or within a word---especially when the word is a
symbol like
       \mbox{\emph{itemnum}}
that makes little sense when hyphenated across
lines.

Footnotes\footnote{This is an example of a footnote.}
pose no problem.

\LaTeX\ is good at typesetting mathematical formulas
like
       \( x-3y + z = 7 \)
or
       \( a_{1} > x^{2n} + y^{2n} > x' \)
or
       \( \ip{A}{B} = \sum_{i} a_{i} b_{i} \).
The spaces you type in a formula are
ignored.  Remember that a letter like
       $x$                   % $ ... $  and  \( ... \)  are equivalent
is a formula when it denotes a mathematical
symbol, and it should be typed as one.

\section{Displayed Text}

Text is displayed by indenting it from the left
margin.  Quotations are commonly displayed.  There
are short quotations
\begin{quote}
   This is a short quotation.  It consists of a
   single paragraph of text.  See how it is formatted.
\end{quote}
and longer ones.
\begin{quotation}
   This is a longer quotation.  It consists of two
   paragraphs of text, neither of which are
   particularly interesting.

   This is the second paragraph of the quotation.  It
   is just as dull as the first paragraph.
\end{quotation}
Another frequently-displayed structure is a list.
The following is an example of an \emph{itemized}
list.
\begin{itemize}
   \item This is the first item of an itemized list.
         Each item in the list is marked with a ``tick''.
         You don't have to worry about what kind of tick
         mark is used.

   \item This is the second item of the list.  It
         contains another list nested inside it.  The inner
         list is an \emph{enumerated} list.
         \begin{enumerate}
            \item This is the first item of an enumerated
                  list that is nested within the itemized list.

            \item This is the second item of the inner list.
                  \LaTeX\ allows you to nest lists deeper than
                  you really should.
         \end{enumerate}
         This is the rest of the second item of the outer
         list.  It is no more interesting than any other
         part of the item.
   \item This is the third item of the list.
\end{itemize}
You can even display poetry.
\begin{verse}
   There is an environment
    for verse \\             % The \\ command separates lines
   Whose features some poets % within a stanza.
   will curse.

                             % One or more blank lines separate stanzas.

   For instead of making\\
   Them do \emph{all} line breaking, \\
   It allows them to put too many words on a line when they'd rather be
   forced to be terse.
\end{verse}

Mathematical formulas may also be displayed.  A
displayed formula
is
one-line long; multiline
formulas require special formatting instructions.
   \[  \ip{\Gamma}{\psi'} = x'' + y^{2} + z_{i}^{n}\]
Don't start a paragraph with a displayed equation,
nor make one a paragraph by itself.

\end{document}               % End of document.

\section{Bench Huge Section 1071}
% synthetic bench section 1071
% This is a sample LaTeX input file.  (Version of 12 August 2004.)
%
% A '%' character causes TeX to ignore all remaining text on the line,
% and is used for comments like this one.

\documentclass{article}      % Specifies the document class

                             % The preamble begins here.
\title{An Example Document}  % Declares the document's title.
\author{Leslie Lamport}      % Declares the author's name.
\date{January 21, 1994}      % Deleting this command produces today's date.

\newcommand{\ip}[2]{(#1, #2)}
                             % Defines \ip{arg1}{arg2} to mean
                             % (arg1, arg2).

%\newcommand{\ip}[2]{\langle #1 | #2\rangle}
                             % This is an alternative definition of
                             % \ip that is commented out.

\begin{document}             % End of preamble and beginning of text.

\maketitle                   % Produces the title.

This is an example input file.  Comparing it with
the output it generates can show you how to
produce a simple document of your own.

\section{Ordinary Text}      % Produces section heading.  Lower-level
                             % sections are begun with similar
                             % \subsection and \subsubsection commands.

The ends  of words and sentences are marked
  by   spaces. It  doesn't matter how many
spaces    you type; one is as good as 100.  The
end of   a line counts as a space.

One   or more   blank lines denote the  end
of  a paragraph.

Since any number of consecutive spaces are treated
like a single one, the formatting of the input
file makes no difference to
      \LaTeX,                % The \LaTeX command generates the LaTeX logo.
but it makes a difference to you.  When you use
\LaTeX, making your input file as easy to read
as possible will be a great help as you write
your document and when you change it.  This sample
file shows how you can add comments to your own input
file.

Because printing is different from typewriting,
there are a number of things that you have to do
differently when preparing an input file than if
you were just typing the document directly.
Quotation marks like
       ``this''
have to be handled specially, as do quotes within
quotes:
       ``\,`this'            % \, separates the double and single quote.
        is what I just
        wrote, not  `that'\,''.

Dashes come in three sizes: an
       intra-word
dash, a medium dash for number ranges like
       1--2,
and a punctuation
       dash---like
this.

A sentence-ending space should be larger than the
space between words within a sentence.  You
sometimes have to type special commands in
conjunction with punctuation characters to get
this right, as in the following sentence.
       Gnats, gnus, etc.\ all  % `\ ' makes an inter-word space.
       begin with G\@.         % \@ marks end-of-sentence punctuation.
You should check the spaces after periods when
reading your output to make sure you haven't
forgotten any special cases.  Generating an
ellipsis
       \ldots\               % `\ ' is needed after `\ldots' because TeX
                             % ignores spaces after command names like \ldots
                             % made from \ + letters.
                             %
                             % Note how a `%' character causes TeX to ignore
                             % the end of the input line, so these blank lines
                             % do not start a new paragraph.
                             %
with the right spacing around the periods requires
a special command.

\LaTeX\ interprets some common characters as
commands, so you must type special commands to
generate them.  These characters include the
following:
       \$ \& \% \# \{ and \}.

In printing, text is usually emphasized with an
       \emph{italic}
type style.

\begin{em}
   A long segment of text can also be emphasized
   in this way.  Text within such a segment can be
   given \emph{additional} emphasis.
\end{em}

It is sometimes necessary to prevent \LaTeX\ from
breaking a line where it might otherwise do so.
This may be at a space, as between the ``Mr.''\ and
``Jones'' in
       ``Mr.~Jones'',        % ~ produces an unbreakable interword space.
or within a word---especially when the word is a
symbol like
       \mbox{\emph{itemnum}}
that makes little sense when hyphenated across
lines.

Footnotes\footnote{This is an example of a footnote.}
pose no problem.

\LaTeX\ is good at typesetting mathematical formulas
like
       \( x-3y + z = 7 \)
or
       \( a_{1} > x^{2n} + y^{2n} > x' \)
or
       \( \ip{A}{B} = \sum_{i} a_{i} b_{i} \).
The spaces you type in a formula are
ignored.  Remember that a letter like
       $x$                   % $ ... $  and  \( ... \)  are equivalent
is a formula when it denotes a mathematical
symbol, and it should be typed as one.

\section{Displayed Text}

Text is displayed by indenting it from the left
margin.  Quotations are commonly displayed.  There
are short quotations
\begin{quote}
   This is a short quotation.  It consists of a
   single paragraph of text.  See how it is formatted.
\end{quote}
and longer ones.
\begin{quotation}
   This is a longer quotation.  It consists of two
   paragraphs of text, neither of which are
   particularly interesting.

   This is the second paragraph of the quotation.  It
   is just as dull as the first paragraph.
\end{quotation}
Another frequently-displayed structure is a list.
The following is an example of an \emph{itemized}
list.
\begin{itemize}
   \item This is the first item of an itemized list.
         Each item in the list is marked with a ``tick''.
         You don't have to worry about what kind of tick
         mark is used.

   \item This is the second item of the list.  It
         contains another list nested inside it.  The inner
         list is an \emph{enumerated} list.
         \begin{enumerate}
            \item This is the first item of an enumerated
                  list that is nested within the itemized list.

            \item This is the second item of the inner list.
                  \LaTeX\ allows you to nest lists deeper than
                  you really should.
         \end{enumerate}
         This is the rest of the second item of the outer
         list.  It is no more interesting than any other
         part of the item.
   \item This is the third item of the list.
\end{itemize}
You can even display poetry.
\begin{verse}
   There is an environment
    for verse \\             % The \\ command separates lines
   Whose features some poets % within a stanza.
   will curse.

                             % One or more blank lines separate stanzas.

   For instead of making\\
   Them do \emph{all} line breaking, \\
   It allows them to put too many words on a line when they'd rather be
   forced to be terse.
\end{verse}

Mathematical formulas may also be displayed.  A
displayed formula
is
one-line long; multiline
formulas require special formatting instructions.
   \[  \ip{\Gamma}{\psi'} = x'' + y^{2} + z_{i}^{n}\]
Don't start a paragraph with a displayed equation,
nor make one a paragraph by itself.

\end{document}               % End of document.

\section{Bench Huge Section 1072}
% synthetic bench section 1072
% This is a sample LaTeX input file.  (Version of 12 August 2004.)
%
% A '%' character causes TeX to ignore all remaining text on the line,
% and is used for comments like this one.

\documentclass{article}      % Specifies the document class

                             % The preamble begins here.
\title{An Example Document}  % Declares the document's title.
\author{Leslie Lamport}      % Declares the author's name.
\date{January 21, 1994}      % Deleting this command produces today's date.

\newcommand{\ip}[2]{(#1, #2)}
                             % Defines \ip{arg1}{arg2} to mean
                             % (arg1, arg2).

%\newcommand{\ip}[2]{\langle #1 | #2\rangle}
                             % This is an alternative definition of
                             % \ip that is commented out.

\begin{document}             % End of preamble and beginning of text.

\maketitle                   % Produces the title.

This is an example input file.  Comparing it with
the output it generates can show you how to
produce a simple document of your own.

\section{Ordinary Text}      % Produces section heading.  Lower-level
                             % sections are begun with similar
                             % \subsection and \subsubsection commands.

The ends  of words and sentences are marked
  by   spaces. It  doesn't matter how many
spaces    you type; one is as good as 100.  The
end of   a line counts as a space.

One   or more   blank lines denote the  end
of  a paragraph.

Since any number of consecutive spaces are treated
like a single one, the formatting of the input
file makes no difference to
      \LaTeX,                % The \LaTeX command generates the LaTeX logo.
but it makes a difference to you.  When you use
\LaTeX, making your input file as easy to read
as possible will be a great help as you write
your document and when you change it.  This sample
file shows how you can add comments to your own input
file.

Because printing is different from typewriting,
there are a number of things that you have to do
differently when preparing an input file than if
you were just typing the document directly.
Quotation marks like
       ``this''
have to be handled specially, as do quotes within
quotes:
       ``\,`this'            % \, separates the double and single quote.
        is what I just
        wrote, not  `that'\,''.

Dashes come in three sizes: an
       intra-word
dash, a medium dash for number ranges like
       1--2,
and a punctuation
       dash---like
this.

A sentence-ending space should be larger than the
space between words within a sentence.  You
sometimes have to type special commands in
conjunction with punctuation characters to get
this right, as in the following sentence.
       Gnats, gnus, etc.\ all  % `\ ' makes an inter-word space.
       begin with G\@.         % \@ marks end-of-sentence punctuation.
You should check the spaces after periods when
reading your output to make sure you haven't
forgotten any special cases.  Generating an
ellipsis
       \ldots\               % `\ ' is needed after `\ldots' because TeX
                             % ignores spaces after command names like \ldots
                             % made from \ + letters.
                             %
                             % Note how a `%' character causes TeX to ignore
                             % the end of the input line, so these blank lines
                             % do not start a new paragraph.
                             %
with the right spacing around the periods requires
a special command.

\LaTeX\ interprets some common characters as
commands, so you must type special commands to
generate them.  These characters include the
following:
       \$ \& \% \# \{ and \}.

In printing, text is usually emphasized with an
       \emph{italic}
type style.

\begin{em}
   A long segment of text can also be emphasized
   in this way.  Text within such a segment can be
   given \emph{additional} emphasis.
\end{em}

It is sometimes necessary to prevent \LaTeX\ from
breaking a line where it might otherwise do so.
This may be at a space, as between the ``Mr.''\ and
``Jones'' in
       ``Mr.~Jones'',        % ~ produces an unbreakable interword space.
or within a word---especially when the word is a
symbol like
       \mbox{\emph{itemnum}}
that makes little sense when hyphenated across
lines.

Footnotes\footnote{This is an example of a footnote.}
pose no problem.

\LaTeX\ is good at typesetting mathematical formulas
like
       \( x-3y + z = 7 \)
or
       \( a_{1} > x^{2n} + y^{2n} > x' \)
or
       \( \ip{A}{B} = \sum_{i} a_{i} b_{i} \).
The spaces you type in a formula are
ignored.  Remember that a letter like
       $x$                   % $ ... $  and  \( ... \)  are equivalent
is a formula when it denotes a mathematical
symbol, and it should be typed as one.

\section{Displayed Text}

Text is displayed by indenting it from the left
margin.  Quotations are commonly displayed.  There
are short quotations
\begin{quote}
   This is a short quotation.  It consists of a
   single paragraph of text.  See how it is formatted.
\end{quote}
and longer ones.
\begin{quotation}
   This is a longer quotation.  It consists of two
   paragraphs of text, neither of which are
   particularly interesting.

   This is the second paragraph of the quotation.  It
   is just as dull as the first paragraph.
\end{quotation}
Another frequently-displayed structure is a list.
The following is an example of an \emph{itemized}
list.
\begin{itemize}
   \item This is the first item of an itemized list.
         Each item in the list is marked with a ``tick''.
         You don't have to worry about what kind of tick
         mark is used.

   \item This is the second item of the list.  It
         contains another list nested inside it.  The inner
         list is an \emph{enumerated} list.
         \begin{enumerate}
            \item This is the first item of an enumerated
                  list that is nested within the itemized list.

            \item This is the second item of the inner list.
                  \LaTeX\ allows you to nest lists deeper than
                  you really should.
         \end{enumerate}
         This is the rest of the second item of the outer
         list.  It is no more interesting than any other
         part of the item.
   \item This is the third item of the list.
\end{itemize}
You can even display poetry.
\begin{verse}
   There is an environment
    for verse \\             % The \\ command separates lines
   Whose features some poets % within a stanza.
   will curse.

                             % One or more blank lines separate stanzas.

   For instead of making\\
   Them do \emph{all} line breaking, \\
   It allows them to put too many words on a line when they'd rather be
   forced to be terse.
\end{verse}

Mathematical formulas may also be displayed.  A
displayed formula
is
one-line long; multiline
formulas require special formatting instructions.
   \[  \ip{\Gamma}{\psi'} = x'' + y^{2} + z_{i}^{n}\]
Don't start a paragraph with a displayed equation,
nor make one a paragraph by itself.

\end{document}               % End of document.

\section{Bench Huge Section 1073}
% synthetic bench section 1073
% This is a sample LaTeX input file.  (Version of 12 August 2004.)
%
% A '%' character causes TeX to ignore all remaining text on the line,
% and is used for comments like this one.

\documentclass{article}      % Specifies the document class

                             % The preamble begins here.
\title{An Example Document}  % Declares the document's title.
\author{Leslie Lamport}      % Declares the author's name.
\date{January 21, 1994}      % Deleting this command produces today's date.

\newcommand{\ip}[2]{(#1, #2)}
                             % Defines \ip{arg1}{arg2} to mean
                             % (arg1, arg2).

%\newcommand{\ip}[2]{\langle #1 | #2\rangle}
                             % This is an alternative definition of
                             % \ip that is commented out.

\begin{document}             % End of preamble and beginning of text.

\maketitle                   % Produces the title.

This is an example input file.  Comparing it with
the output it generates can show you how to
produce a simple document of your own.

\section{Ordinary Text}      % Produces section heading.  Lower-level
                             % sections are begun with similar
                             % \subsection and \subsubsection commands.

The ends  of words and sentences are marked
  by   spaces. It  doesn't matter how many
spaces    you type; one is as good as 100.  The
end of   a line counts as a space.

One   or more   blank lines denote the  end
of  a paragraph.

Since any number of consecutive spaces are treated
like a single one, the formatting of the input
file makes no difference to
      \LaTeX,                % The \LaTeX command generates the LaTeX logo.
but it makes a difference to you.  When you use
\LaTeX, making your input file as easy to read
as possible will be a great help as you write
your document and when you change it.  This sample
file shows how you can add comments to your own input
file.

Because printing is different from typewriting,
there are a number of things that you have to do
differently when preparing an input file than if
you were just typing the document directly.
Quotation marks like
       ``this''
have to be handled specially, as do quotes within
quotes:
       ``\,`this'            % \, separates the double and single quote.
        is what I just
        wrote, not  `that'\,''.

Dashes come in three sizes: an
       intra-word
dash, a medium dash for number ranges like
       1--2,
and a punctuation
       dash---like
this.

A sentence-ending space should be larger than the
space between words within a sentence.  You
sometimes have to type special commands in
conjunction with punctuation characters to get
this right, as in the following sentence.
       Gnats, gnus, etc.\ all  % `\ ' makes an inter-word space.
       begin with G\@.         % \@ marks end-of-sentence punctuation.
You should check the spaces after periods when
reading your output to make sure you haven't
forgotten any special cases.  Generating an
ellipsis
       \ldots\               % `\ ' is needed after `\ldots' because TeX
                             % ignores spaces after command names like \ldots
                             % made from \ + letters.
                             %
                             % Note how a `%' character causes TeX to ignore
                             % the end of the input line, so these blank lines
                             % do not start a new paragraph.
                             %
with the right spacing around the periods requires
a special command.

\LaTeX\ interprets some common characters as
commands, so you must type special commands to
generate them.  These characters include the
following:
       \$ \& \% \# \{ and \}.

In printing, text is usually emphasized with an
       \emph{italic}
type style.

\begin{em}
   A long segment of text can also be emphasized
   in this way.  Text within such a segment can be
   given \emph{additional} emphasis.
\end{em}

It is sometimes necessary to prevent \LaTeX\ from
breaking a line where it might otherwise do so.
This may be at a space, as between the ``Mr.''\ and
``Jones'' in
       ``Mr.~Jones'',        % ~ produces an unbreakable interword space.
or within a word---especially when the word is a
symbol like
       \mbox{\emph{itemnum}}
that makes little sense when hyphenated across
lines.

Footnotes\footnote{This is an example of a footnote.}
pose no problem.

\LaTeX\ is good at typesetting mathematical formulas
like
       \( x-3y + z = 7 \)
or
       \( a_{1} > x^{2n} + y^{2n} > x' \)
or
       \( \ip{A}{B} = \sum_{i} a_{i} b_{i} \).
The spaces you type in a formula are
ignored.  Remember that a letter like
       $x$                   % $ ... $  and  \( ... \)  are equivalent
is a formula when it denotes a mathematical
symbol, and it should be typed as one.

\section{Displayed Text}

Text is displayed by indenting it from the left
margin.  Quotations are commonly displayed.  There
are short quotations
\begin{quote}
   This is a short quotation.  It consists of a
   single paragraph of text.  See how it is formatted.
\end{quote}
and longer ones.
\begin{quotation}
   This is a longer quotation.  It consists of two
   paragraphs of text, neither of which are
   particularly interesting.

   This is the second paragraph of the quotation.  It
   is just as dull as the first paragraph.
\end{quotation}
Another frequently-displayed structure is a list.
The following is an example of an \emph{itemized}
list.
\begin{itemize}
   \item This is the first item of an itemized list.
         Each item in the list is marked with a ``tick''.
         You don't have to worry about what kind of tick
         mark is used.

   \item This is the second item of the list.  It
         contains another list nested inside it.  The inner
         list is an \emph{enumerated} list.
         \begin{enumerate}
            \item This is the first item of an enumerated
                  list that is nested within the itemized list.

            \item This is the second item of the inner list.
                  \LaTeX\ allows you to nest lists deeper than
                  you really should.
         \end{enumerate}
         This is the rest of the second item of the outer
         list.  It is no more interesting than any other
         part of the item.
   \item This is the third item of the list.
\end{itemize}
You can even display poetry.
\begin{verse}
   There is an environment
    for verse \\             % The \\ command separates lines
   Whose features some poets % within a stanza.
   will curse.

                             % One or more blank lines separate stanzas.

   For instead of making\\
   Them do \emph{all} line breaking, \\
   It allows them to put too many words on a line when they'd rather be
   forced to be terse.
\end{verse}

Mathematical formulas may also be displayed.  A
displayed formula
is
one-line long; multiline
formulas require special formatting instructions.
   \[  \ip{\Gamma}{\psi'} = x'' + y^{2} + z_{i}^{n}\]
Don't start a paragraph with a displayed equation,
nor make one a paragraph by itself.

\end{document}               % End of document.

\section{Bench Huge Section 1074}
% synthetic bench section 1074
% This is a sample LaTeX input file.  (Version of 12 August 2004.)
%
% A '%' character causes TeX to ignore all remaining text on the line,
% and is used for comments like this one.

\documentclass{article}      % Specifies the document class

                             % The preamble begins here.
\title{An Example Document}  % Declares the document's title.
\author{Leslie Lamport}      % Declares the author's name.
\date{January 21, 1994}      % Deleting this command produces today's date.

\newcommand{\ip}[2]{(#1, #2)}
                             % Defines \ip{arg1}{arg2} to mean
                             % (arg1, arg2).

%\newcommand{\ip}[2]{\langle #1 | #2\rangle}
                             % This is an alternative definition of
                             % \ip that is commented out.

\begin{document}             % End of preamble and beginning of text.

\maketitle                   % Produces the title.

This is an example input file.  Comparing it with
the output it generates can show you how to
produce a simple document of your own.

\section{Ordinary Text}      % Produces section heading.  Lower-level
                             % sections are begun with similar
                             % \subsection and \subsubsection commands.

The ends  of words and sentences are marked
  by   spaces. It  doesn't matter how many
spaces    you type; one is as good as 100.  The
end of   a line counts as a space.

One   or more   blank lines denote the  end
of  a paragraph.

Since any number of consecutive spaces are treated
like a single one, the formatting of the input
file makes no difference to
      \LaTeX,                % The \LaTeX command generates the LaTeX logo.
but it makes a difference to you.  When you use
\LaTeX, making your input file as easy to read
as possible will be a great help as you write
your document and when you change it.  This sample
file shows how you can add comments to your own input
file.

Because printing is different from typewriting,
there are a number of things that you have to do
differently when preparing an input file than if
you were just typing the document directly.
Quotation marks like
       ``this''
have to be handled specially, as do quotes within
quotes:
       ``\,`this'            % \, separates the double and single quote.
        is what I just
        wrote, not  `that'\,''.

Dashes come in three sizes: an
       intra-word
dash, a medium dash for number ranges like
       1--2,
and a punctuation
       dash---like
this.

A sentence-ending space should be larger than the
space between words within a sentence.  You
sometimes have to type special commands in
conjunction with punctuation characters to get
this right, as in the following sentence.
       Gnats, gnus, etc.\ all  % `\ ' makes an inter-word space.
       begin with G\@.         % \@ marks end-of-sentence punctuation.
You should check the spaces after periods when
reading your output to make sure you haven't
forgotten any special cases.  Generating an
ellipsis
       \ldots\               % `\ ' is needed after `\ldots' because TeX
                             % ignores spaces after command names like \ldots
                             % made from \ + letters.
                             %
                             % Note how a `%' character causes TeX to ignore
                             % the end of the input line, so these blank lines
                             % do not start a new paragraph.
                             %
with the right spacing around the periods requires
a special command.

\LaTeX\ interprets some common characters as
commands, so you must type special commands to
generate them.  These characters include the
following:
       \$ \& \% \# \{ and \}.

In printing, text is usually emphasized with an
       \emph{italic}
type style.

\begin{em}
   A long segment of text can also be emphasized
   in this way.  Text within such a segment can be
   given \emph{additional} emphasis.
\end{em}

It is sometimes necessary to prevent \LaTeX\ from
breaking a line where it might otherwise do so.
This may be at a space, as between the ``Mr.''\ and
``Jones'' in
       ``Mr.~Jones'',        % ~ produces an unbreakable interword space.
or within a word---especially when the word is a
symbol like
       \mbox{\emph{itemnum}}
that makes little sense when hyphenated across
lines.

Footnotes\footnote{This is an example of a footnote.}
pose no problem.

\LaTeX\ is good at typesetting mathematical formulas
like
       \( x-3y + z = 7 \)
or
       \( a_{1} > x^{2n} + y^{2n} > x' \)
or
       \( \ip{A}{B} = \sum_{i} a_{i} b_{i} \).
The spaces you type in a formula are
ignored.  Remember that a letter like
       $x$                   % $ ... $  and  \( ... \)  are equivalent
is a formula when it denotes a mathematical
symbol, and it should be typed as one.

\section{Displayed Text}

Text is displayed by indenting it from the left
margin.  Quotations are commonly displayed.  There
are short quotations
\begin{quote}
   This is a short quotation.  It consists of a
   single paragraph of text.  See how it is formatted.
\end{quote}
and longer ones.
\begin{quotation}
   This is a longer quotation.  It consists of two
   paragraphs of text, neither of which are
   particularly interesting.

   This is the second paragraph of the quotation.  It
   is just as dull as the first paragraph.
\end{quotation}
Another frequently-displayed structure is a list.
The following is an example of an \emph{itemized}
list.
\begin{itemize}
   \item This is the first item of an itemized list.
         Each item in the list is marked with a ``tick''.
         You don't have to worry about what kind of tick
         mark is used.

   \item This is the second item of the list.  It
         contains another list nested inside it.  The inner
         list is an \emph{enumerated} list.
         \begin{enumerate}
            \item This is the first item of an enumerated
                  list that is nested within the itemized list.

            \item This is the second item of the inner list.
                  \LaTeX\ allows you to nest lists deeper than
                  you really should.
         \end{enumerate}
         This is the rest of the second item of the outer
         list.  It is no more interesting than any other
         part of the item.
   \item This is the third item of the list.
\end{itemize}
You can even display poetry.
\begin{verse}
   There is an environment
    for verse \\             % The \\ command separates lines
   Whose features some poets % within a stanza.
   will curse.

                             % One or more blank lines separate stanzas.

   For instead of making\\
   Them do \emph{all} line breaking, \\
   It allows them to put too many words on a line when they'd rather be
   forced to be terse.
\end{verse}

Mathematical formulas may also be displayed.  A
displayed formula
is
one-line long; multiline
formulas require special formatting instructions.
   \[  \ip{\Gamma}{\psi'} = x'' + y^{2} + z_{i}^{n}\]
Don't start a paragraph with a displayed equation,
nor make one a paragraph by itself.

\end{document}               % End of document.

\section{Bench Huge Section 1075}
% synthetic bench section 1075
% This is a sample LaTeX input file.  (Version of 12 August 2004.)
%
% A '%' character causes TeX to ignore all remaining text on the line,
% and is used for comments like this one.

\documentclass{article}      % Specifies the document class

                             % The preamble begins here.
\title{An Example Document}  % Declares the document's title.
\author{Leslie Lamport}      % Declares the author's name.
\date{January 21, 1994}      % Deleting this command produces today's date.

\newcommand{\ip}[2]{(#1, #2)}
                             % Defines \ip{arg1}{arg2} to mean
                             % (arg1, arg2).

%\newcommand{\ip}[2]{\langle #1 | #2\rangle}
                             % This is an alternative definition of
                             % \ip that is commented out.

\begin{document}             % End of preamble and beginning of text.

\maketitle                   % Produces the title.

This is an example input file.  Comparing it with
the output it generates can show you how to
produce a simple document of your own.

\section{Ordinary Text}      % Produces section heading.  Lower-level
                             % sections are begun with similar
                             % \subsection and \subsubsection commands.

The ends  of words and sentences are marked
  by   spaces. It  doesn't matter how many
spaces    you type; one is as good as 100.  The
end of   a line counts as a space.

One   or more   blank lines denote the  end
of  a paragraph.

Since any number of consecutive spaces are treated
like a single one, the formatting of the input
file makes no difference to
      \LaTeX,                % The \LaTeX command generates the LaTeX logo.
but it makes a difference to you.  When you use
\LaTeX, making your input file as easy to read
as possible will be a great help as you write
your document and when you change it.  This sample
file shows how you can add comments to your own input
file.

Because printing is different from typewriting,
there are a number of things that you have to do
differently when preparing an input file than if
you were just typing the document directly.
Quotation marks like
       ``this''
have to be handled specially, as do quotes within
quotes:
       ``\,`this'            % \, separates the double and single quote.
        is what I just
        wrote, not  `that'\,''.

Dashes come in three sizes: an
       intra-word
dash, a medium dash for number ranges like
       1--2,
and a punctuation
       dash---like
this.

A sentence-ending space should be larger than the
space between words within a sentence.  You
sometimes have to type special commands in
conjunction with punctuation characters to get
this right, as in the following sentence.
       Gnats, gnus, etc.\ all  % `\ ' makes an inter-word space.
       begin with G\@.         % \@ marks end-of-sentence punctuation.
You should check the spaces after periods when
reading your output to make sure you haven't
forgotten any special cases.  Generating an
ellipsis
       \ldots\               % `\ ' is needed after `\ldots' because TeX
                             % ignores spaces after command names like \ldots
                             % made from \ + letters.
                             %
                             % Note how a `%' character causes TeX to ignore
                             % the end of the input line, so these blank lines
                             % do not start a new paragraph.
                             %
with the right spacing around the periods requires
a special command.

\LaTeX\ interprets some common characters as
commands, so you must type special commands to
generate them.  These characters include the
following:
       \$ \& \% \# \{ and \}.

In printing, text is usually emphasized with an
       \emph{italic}
type style.

\begin{em}
   A long segment of text can also be emphasized
   in this way.  Text within such a segment can be
   given \emph{additional} emphasis.
\end{em}

It is sometimes necessary to prevent \LaTeX\ from
breaking a line where it might otherwise do so.
This may be at a space, as between the ``Mr.''\ and
``Jones'' in
       ``Mr.~Jones'',        % ~ produces an unbreakable interword space.
or within a word---especially when the word is a
symbol like
       \mbox{\emph{itemnum}}
that makes little sense when hyphenated across
lines.

Footnotes\footnote{This is an example of a footnote.}
pose no problem.

\LaTeX\ is good at typesetting mathematical formulas
like
       \( x-3y + z = 7 \)
or
       \( a_{1} > x^{2n} + y^{2n} > x' \)
or
       \( \ip{A}{B} = \sum_{i} a_{i} b_{i} \).
The spaces you type in a formula are
ignored.  Remember that a letter like
       $x$                   % $ ... $  and  \( ... \)  are equivalent
is a formula when it denotes a mathematical
symbol, and it should be typed as one.

\section{Displayed Text}

Text is displayed by indenting it from the left
margin.  Quotations are commonly displayed.  There
are short quotations
\begin{quote}
   This is a short quotation.  It consists of a
   single paragraph of text.  See how it is formatted.
\end{quote}
and longer ones.
\begin{quotation}
   This is a longer quotation.  It consists of two
   paragraphs of text, neither of which are
   particularly interesting.

   This is the second paragraph of the quotation.  It
   is just as dull as the first paragraph.
\end{quotation}
Another frequently-displayed structure is a list.
The following is an example of an \emph{itemized}
list.
\begin{itemize}
   \item This is the first item of an itemized list.
         Each item in the list is marked with a ``tick''.
         You don't have to worry about what kind of tick
         mark is used.

   \item This is the second item of the list.  It
         contains another list nested inside it.  The inner
         list is an \emph{enumerated} list.
         \begin{enumerate}
            \item This is the first item of an enumerated
                  list that is nested within the itemized list.

            \item This is the second item of the inner list.
                  \LaTeX\ allows you to nest lists deeper than
                  you really should.
         \end{enumerate}
         This is the rest of the second item of the outer
         list.  It is no more interesting than any other
         part of the item.
   \item This is the third item of the list.
\end{itemize}
You can even display poetry.
\begin{verse}
   There is an environment
    for verse \\             % The \\ command separates lines
   Whose features some poets % within a stanza.
   will curse.

                             % One or more blank lines separate stanzas.

   For instead of making\\
   Them do \emph{all} line breaking, \\
   It allows them to put too many words on a line when they'd rather be
   forced to be terse.
\end{verse}

Mathematical formulas may also be displayed.  A
displayed formula
is
one-line long; multiline
formulas require special formatting instructions.
   \[  \ip{\Gamma}{\psi'} = x'' + y^{2} + z_{i}^{n}\]
Don't start a paragraph with a displayed equation,
nor make one a paragraph by itself.

\end{document}               % End of document.

\section{Bench Huge Section 1076}
% synthetic bench section 1076
% This is a sample LaTeX input file.  (Version of 12 August 2004.)
%
% A '%' character causes TeX to ignore all remaining text on the line,
% and is used for comments like this one.

\documentclass{article}      % Specifies the document class

                             % The preamble begins here.
\title{An Example Document}  % Declares the document's title.
\author{Leslie Lamport}      % Declares the author's name.
\date{January 21, 1994}      % Deleting this command produces today's date.

\newcommand{\ip}[2]{(#1, #2)}
                             % Defines \ip{arg1}{arg2} to mean
                             % (arg1, arg2).

%\newcommand{\ip}[2]{\langle #1 | #2\rangle}
                             % This is an alternative definition of
                             % \ip that is commented out.

\begin{document}             % End of preamble and beginning of text.

\maketitle                   % Produces the title.

This is an example input file.  Comparing it with
the output it generates can show you how to
produce a simple document of your own.

\section{Ordinary Text}      % Produces section heading.  Lower-level
                             % sections are begun with similar
                             % \subsection and \subsubsection commands.

The ends  of words and sentences are marked
  by   spaces. It  doesn't matter how many
spaces    you type; one is as good as 100.  The
end of   a line counts as a space.

One   or more   blank lines denote the  end
of  a paragraph.

Since any number of consecutive spaces are treated
like a single one, the formatting of the input
file makes no difference to
      \LaTeX,                % The \LaTeX command generates the LaTeX logo.
but it makes a difference to you.  When you use
\LaTeX, making your input file as easy to read
as possible will be a great help as you write
your document and when you change it.  This sample
file shows how you can add comments to your own input
file.

Because printing is different from typewriting,
there are a number of things that you have to do
differently when preparing an input file than if
you were just typing the document directly.
Quotation marks like
       ``this''
have to be handled specially, as do quotes within
quotes:
       ``\,`this'            % \, separates the double and single quote.
        is what I just
        wrote, not  `that'\,''.

Dashes come in three sizes: an
       intra-word
dash, a medium dash for number ranges like
       1--2,
and a punctuation
       dash---like
this.

A sentence-ending space should be larger than the
space between words within a sentence.  You
sometimes have to type special commands in
conjunction with punctuation characters to get
this right, as in the following sentence.
       Gnats, gnus, etc.\ all  % `\ ' makes an inter-word space.
       begin with G\@.         % \@ marks end-of-sentence punctuation.
You should check the spaces after periods when
reading your output to make sure you haven't
forgotten any special cases.  Generating an
ellipsis
       \ldots\               % `\ ' is needed after `\ldots' because TeX
                             % ignores spaces after command names like \ldots
                             % made from \ + letters.
                             %
                             % Note how a `%' character causes TeX to ignore
                             % the end of the input line, so these blank lines
                             % do not start a new paragraph.
                             %
with the right spacing around the periods requires
a special command.

\LaTeX\ interprets some common characters as
commands, so you must type special commands to
generate them.  These characters include the
following:
       \$ \& \% \# \{ and \}.

In printing, text is usually emphasized with an
       \emph{italic}
type style.

\begin{em}
   A long segment of text can also be emphasized
   in this way.  Text within such a segment can be
   given \emph{additional} emphasis.
\end{em}

It is sometimes necessary to prevent \LaTeX\ from
breaking a line where it might otherwise do so.
This may be at a space, as between the ``Mr.''\ and
``Jones'' in
       ``Mr.~Jones'',        % ~ produces an unbreakable interword space.
or within a word---especially when the word is a
symbol like
       \mbox{\emph{itemnum}}
that makes little sense when hyphenated across
lines.

Footnotes\footnote{This is an example of a footnote.}
pose no problem.

\LaTeX\ is good at typesetting mathematical formulas
like
       \( x-3y + z = 7 \)
or
       \( a_{1} > x^{2n} + y^{2n} > x' \)
or
       \( \ip{A}{B} = \sum_{i} a_{i} b_{i} \).
The spaces you type in a formula are
ignored.  Remember that a letter like
       $x$                   % $ ... $  and  \( ... \)  are equivalent
is a formula when it denotes a mathematical
symbol, and it should be typed as one.

\section{Displayed Text}

Text is displayed by indenting it from the left
margin.  Quotations are commonly displayed.  There
are short quotations
\begin{quote}
   This is a short quotation.  It consists of a
   single paragraph of text.  See how it is formatted.
\end{quote}
and longer ones.
\begin{quotation}
   This is a longer quotation.  It consists of two
   paragraphs of text, neither of which are
   particularly interesting.

   This is the second paragraph of the quotation.  It
   is just as dull as the first paragraph.
\end{quotation}
Another frequently-displayed structure is a list.
The following is an example of an \emph{itemized}
list.
\begin{itemize}
   \item This is the first item of an itemized list.
         Each item in the list is marked with a ``tick''.
         You don't have to worry about what kind of tick
         mark is used.

   \item This is the second item of the list.  It
         contains another list nested inside it.  The inner
         list is an \emph{enumerated} list.
         \begin{enumerate}
            \item This is the first item of an enumerated
                  list that is nested within the itemized list.

            \item This is the second item of the inner list.
                  \LaTeX\ allows you to nest lists deeper than
                  you really should.
         \end{enumerate}
         This is the rest of the second item of the outer
         list.  It is no more interesting than any other
         part of the item.
   \item This is the third item of the list.
\end{itemize}
You can even display poetry.
\begin{verse}
   There is an environment
    for verse \\             % The \\ command separates lines
   Whose features some poets % within a stanza.
   will curse.

                             % One or more blank lines separate stanzas.

   For instead of making\\
   Them do \emph{all} line breaking, \\
   It allows them to put too many words on a line when they'd rather be
   forced to be terse.
\end{verse}

Mathematical formulas may also be displayed.  A
displayed formula
is
one-line long; multiline
formulas require special formatting instructions.
   \[  \ip{\Gamma}{\psi'} = x'' + y^{2} + z_{i}^{n}\]
Don't start a paragraph with a displayed equation,
nor make one a paragraph by itself.

\end{document}               % End of document.

\section{Bench Huge Section 1077}
% synthetic bench section 1077
% This is a sample LaTeX input file.  (Version of 12 August 2004.)
%
% A '%' character causes TeX to ignore all remaining text on the line,
% and is used for comments like this one.

\documentclass{article}      % Specifies the document class

                             % The preamble begins here.
\title{An Example Document}  % Declares the document's title.
\author{Leslie Lamport}      % Declares the author's name.
\date{January 21, 1994}      % Deleting this command produces today's date.

\newcommand{\ip}[2]{(#1, #2)}
                             % Defines \ip{arg1}{arg2} to mean
                             % (arg1, arg2).

%\newcommand{\ip}[2]{\langle #1 | #2\rangle}
                             % This is an alternative definition of
                             % \ip that is commented out.

\begin{document}             % End of preamble and beginning of text.

\maketitle                   % Produces the title.

This is an example input file.  Comparing it with
the output it generates can show you how to
produce a simple document of your own.

\section{Ordinary Text}      % Produces section heading.  Lower-level
                             % sections are begun with similar
                             % \subsection and \subsubsection commands.

The ends  of words and sentences are marked
  by   spaces. It  doesn't matter how many
spaces    you type; one is as good as 100.  The
end of   a line counts as a space.

One   or more   blank lines denote the  end
of  a paragraph.

Since any number of consecutive spaces are treated
like a single one, the formatting of the input
file makes no difference to
      \LaTeX,                % The \LaTeX command generates the LaTeX logo.
but it makes a difference to you.  When you use
\LaTeX, making your input file as easy to read
as possible will be a great help as you write
your document and when you change it.  This sample
file shows how you can add comments to your own input
file.

Because printing is different from typewriting,
there are a number of things that you have to do
differently when preparing an input file than if
you were just typing the document directly.
Quotation marks like
       ``this''
have to be handled specially, as do quotes within
quotes:
       ``\,`this'            % \, separates the double and single quote.
        is what I just
        wrote, not  `that'\,''.

Dashes come in three sizes: an
       intra-word
dash, a medium dash for number ranges like
       1--2,
and a punctuation
       dash---like
this.

A sentence-ending space should be larger than the
space between words within a sentence.  You
sometimes have to type special commands in
conjunction with punctuation characters to get
this right, as in the following sentence.
       Gnats, gnus, etc.\ all  % `\ ' makes an inter-word space.
       begin with G\@.         % \@ marks end-of-sentence punctuation.
You should check the spaces after periods when
reading your output to make sure you haven't
forgotten any special cases.  Generating an
ellipsis
       \ldots\               % `\ ' is needed after `\ldots' because TeX
                             % ignores spaces after command names like \ldots
                             % made from \ + letters.
                             %
                             % Note how a `%' character causes TeX to ignore
                             % the end of the input line, so these blank lines
                             % do not start a new paragraph.
                             %
with the right spacing around the periods requires
a special command.

\LaTeX\ interprets some common characters as
commands, so you must type special commands to
generate them.  These characters include the
following:
       \$ \& \% \# \{ and \}.

In printing, text is usually emphasized with an
       \emph{italic}
type style.

\begin{em}
   A long segment of text can also be emphasized
   in this way.  Text within such a segment can be
   given \emph{additional} emphasis.
\end{em}

It is sometimes necessary to prevent \LaTeX\ from
breaking a line where it might otherwise do so.
This may be at a space, as between the ``Mr.''\ and
``Jones'' in
       ``Mr.~Jones'',        % ~ produces an unbreakable interword space.
or within a word---especially when the word is a
symbol like
       \mbox{\emph{itemnum}}
that makes little sense when hyphenated across
lines.

Footnotes\footnote{This is an example of a footnote.}
pose no problem.

\LaTeX\ is good at typesetting mathematical formulas
like
       \( x-3y + z = 7 \)
or
       \( a_{1} > x^{2n} + y^{2n} > x' \)
or
       \( \ip{A}{B} = \sum_{i} a_{i} b_{i} \).
The spaces you type in a formula are
ignored.  Remember that a letter like
       $x$                   % $ ... $  and  \( ... \)  are equivalent
is a formula when it denotes a mathematical
symbol, and it should be typed as one.

\section{Displayed Text}

Text is displayed by indenting it from the left
margin.  Quotations are commonly displayed.  There
are short quotations
\begin{quote}
   This is a short quotation.  It consists of a
   single paragraph of text.  See how it is formatted.
\end{quote}
and longer ones.
\begin{quotation}
   This is a longer quotation.  It consists of two
   paragraphs of text, neither of which are
   particularly interesting.

   This is the second paragraph of the quotation.  It
   is just as dull as the first paragraph.
\end{quotation}
Another frequently-displayed structure is a list.
The following is an example of an \emph{itemized}
list.
\begin{itemize}
   \item This is the first item of an itemized list.
         Each item in the list is marked with a ``tick''.
         You don't have to worry about what kind of tick
         mark is used.

   \item This is the second item of the list.  It
         contains another list nested inside it.  The inner
         list is an \emph{enumerated} list.
         \begin{enumerate}
            \item This is the first item of an enumerated
                  list that is nested within the itemized list.

            \item This is the second item of the inner list.
                  \LaTeX\ allows you to nest lists deeper than
                  you really should.
         \end{enumerate}
         This is the rest of the second item of the outer
         list.  It is no more interesting than any other
         part of the item.
   \item This is the third item of the list.
\end{itemize}
You can even display poetry.
\begin{verse}
   There is an environment
    for verse \\             % The \\ command separates lines
   Whose features some poets % within a stanza.
   will curse.

                             % One or more blank lines separate stanzas.

   For instead of making\\
   Them do \emph{all} line breaking, \\
   It allows them to put too many words on a line when they'd rather be
   forced to be terse.
\end{verse}

Mathematical formulas may also be displayed.  A
displayed formula
is
one-line long; multiline
formulas require special formatting instructions.
   \[  \ip{\Gamma}{\psi'} = x'' + y^{2} + z_{i}^{n}\]
Don't start a paragraph with a displayed equation,
nor make one a paragraph by itself.

\end{document}               % End of document.

\section{Bench Huge Section 1078}
% synthetic bench section 1078
% This is a sample LaTeX input file.  (Version of 12 August 2004.)
%
% A '%' character causes TeX to ignore all remaining text on the line,
% and is used for comments like this one.

\documentclass{article}      % Specifies the document class

                             % The preamble begins here.
\title{An Example Document}  % Declares the document's title.
\author{Leslie Lamport}      % Declares the author's name.
\date{January 21, 1994}      % Deleting this command produces today's date.

\newcommand{\ip}[2]{(#1, #2)}
                             % Defines \ip{arg1}{arg2} to mean
                             % (arg1, arg2).

%\newcommand{\ip}[2]{\langle #1 | #2\rangle}
                             % This is an alternative definition of
                             % \ip that is commented out.

\begin{document}             % End of preamble and beginning of text.

\maketitle                   % Produces the title.

This is an example input file.  Comparing it with
the output it generates can show you how to
produce a simple document of your own.

\section{Ordinary Text}      % Produces section heading.  Lower-level
                             % sections are begun with similar
                             % \subsection and \subsubsection commands.

The ends  of words and sentences are marked
  by   spaces. It  doesn't matter how many
spaces    you type; one is as good as 100.  The
end of   a line counts as a space.

One   or more   blank lines denote the  end
of  a paragraph.

Since any number of consecutive spaces are treated
like a single one, the formatting of the input
file makes no difference to
      \LaTeX,                % The \LaTeX command generates the LaTeX logo.
but it makes a difference to you.  When you use
\LaTeX, making your input file as easy to read
as possible will be a great help as you write
your document and when you change it.  This sample
file shows how you can add comments to your own input
file.

Because printing is different from typewriting,
there are a number of things that you have to do
differently when preparing an input file than if
you were just typing the document directly.
Quotation marks like
       ``this''
have to be handled specially, as do quotes within
quotes:
       ``\,`this'            % \, separates the double and single quote.
        is what I just
        wrote, not  `that'\,''.

Dashes come in three sizes: an
       intra-word
dash, a medium dash for number ranges like
       1--2,
and a punctuation
       dash---like
this.

A sentence-ending space should be larger than the
space between words within a sentence.  You
sometimes have to type special commands in
conjunction with punctuation characters to get
this right, as in the following sentence.
       Gnats, gnus, etc.\ all  % `\ ' makes an inter-word space.
       begin with G\@.         % \@ marks end-of-sentence punctuation.
You should check the spaces after periods when
reading your output to make sure you haven't
forgotten any special cases.  Generating an
ellipsis
       \ldots\               % `\ ' is needed after `\ldots' because TeX
                             % ignores spaces after command names like \ldots
                             % made from \ + letters.
                             %
                             % Note how a `%' character causes TeX to ignore
                             % the end of the input line, so these blank lines
                             % do not start a new paragraph.
                             %
with the right spacing around the periods requires
a special command.

\LaTeX\ interprets some common characters as
commands, so you must type special commands to
generate them.  These characters include the
following:
       \$ \& \% \# \{ and \}.

In printing, text is usually emphasized with an
       \emph{italic}
type style.

\begin{em}
   A long segment of text can also be emphasized
   in this way.  Text within such a segment can be
   given \emph{additional} emphasis.
\end{em}

It is sometimes necessary to prevent \LaTeX\ from
breaking a line where it might otherwise do so.
This may be at a space, as between the ``Mr.''\ and
``Jones'' in
       ``Mr.~Jones'',        % ~ produces an unbreakable interword space.
or within a word---especially when the word is a
symbol like
       \mbox{\emph{itemnum}}
that makes little sense when hyphenated across
lines.

Footnotes\footnote{This is an example of a footnote.}
pose no problem.

\LaTeX\ is good at typesetting mathematical formulas
like
       \( x-3y + z = 7 \)
or
       \( a_{1} > x^{2n} + y^{2n} > x' \)
or
       \( \ip{A}{B} = \sum_{i} a_{i} b_{i} \).
The spaces you type in a formula are
ignored.  Remember that a letter like
       $x$                   % $ ... $  and  \( ... \)  are equivalent
is a formula when it denotes a mathematical
symbol, and it should be typed as one.

\section{Displayed Text}

Text is displayed by indenting it from the left
margin.  Quotations are commonly displayed.  There
are short quotations
\begin{quote}
   This is a short quotation.  It consists of a
   single paragraph of text.  See how it is formatted.
\end{quote}
and longer ones.
\begin{quotation}
   This is a longer quotation.  It consists of two
   paragraphs of text, neither of which are
   particularly interesting.

   This is the second paragraph of the quotation.  It
   is just as dull as the first paragraph.
\end{quotation}
Another frequently-displayed structure is a list.
The following is an example of an \emph{itemized}
list.
\begin{itemize}
   \item This is the first item of an itemized list.
         Each item in the list is marked with a ``tick''.
         You don't have to worry about what kind of tick
         mark is used.

   \item This is the second item of the list.  It
         contains another list nested inside it.  The inner
         list is an \emph{enumerated} list.
         \begin{enumerate}
            \item This is the first item of an enumerated
                  list that is nested within the itemized list.

            \item This is the second item of the inner list.
                  \LaTeX\ allows you to nest lists deeper than
                  you really should.
         \end{enumerate}
         This is the rest of the second item of the outer
         list.  It is no more interesting than any other
         part of the item.
   \item This is the third item of the list.
\end{itemize}
You can even display poetry.
\begin{verse}
   There is an environment
    for verse \\             % The \\ command separates lines
   Whose features some poets % within a stanza.
   will curse.

                             % One or more blank lines separate stanzas.

   For instead of making\\
   Them do \emph{all} line breaking, \\
   It allows them to put too many words on a line when they'd rather be
   forced to be terse.
\end{verse}

Mathematical formulas may also be displayed.  A
displayed formula
is
one-line long; multiline
formulas require special formatting instructions.
   \[  \ip{\Gamma}{\psi'} = x'' + y^{2} + z_{i}^{n}\]
Don't start a paragraph with a displayed equation,
nor make one a paragraph by itself.

\end{document}               % End of document.

\section{Bench Huge Section 1079}
% synthetic bench section 1079
% This is a sample LaTeX input file.  (Version of 12 August 2004.)
%
% A '%' character causes TeX to ignore all remaining text on the line,
% and is used for comments like this one.

\documentclass{article}      % Specifies the document class

                             % The preamble begins here.
\title{An Example Document}  % Declares the document's title.
\author{Leslie Lamport}      % Declares the author's name.
\date{January 21, 1994}      % Deleting this command produces today's date.

\newcommand{\ip}[2]{(#1, #2)}
                             % Defines \ip{arg1}{arg2} to mean
                             % (arg1, arg2).

%\newcommand{\ip}[2]{\langle #1 | #2\rangle}
                             % This is an alternative definition of
                             % \ip that is commented out.

\begin{document}             % End of preamble and beginning of text.

\maketitle                   % Produces the title.

This is an example input file.  Comparing it with
the output it generates can show you how to
produce a simple document of your own.

\section{Ordinary Text}      % Produces section heading.  Lower-level
                             % sections are begun with similar
                             % \subsection and \subsubsection commands.

The ends  of words and sentences are marked
  by   spaces. It  doesn't matter how many
spaces    you type; one is as good as 100.  The
end of   a line counts as a space.

One   or more   blank lines denote the  end
of  a paragraph.

Since any number of consecutive spaces are treated
like a single one, the formatting of the input
file makes no difference to
      \LaTeX,                % The \LaTeX command generates the LaTeX logo.
but it makes a difference to you.  When you use
\LaTeX, making your input file as easy to read
as possible will be a great help as you write
your document and when you change it.  This sample
file shows how you can add comments to your own input
file.

Because printing is different from typewriting,
there are a number of things that you have to do
differently when preparing an input file than if
you were just typing the document directly.
Quotation marks like
       ``this''
have to be handled specially, as do quotes within
quotes:
       ``\,`this'            % \, separates the double and single quote.
        is what I just
        wrote, not  `that'\,''.

Dashes come in three sizes: an
       intra-word
dash, a medium dash for number ranges like
       1--2,
and a punctuation
       dash---like
this.

A sentence-ending space should be larger than the
space between words within a sentence.  You
sometimes have to type special commands in
conjunction with punctuation characters to get
this right, as in the following sentence.
       Gnats, gnus, etc.\ all  % `\ ' makes an inter-word space.
       begin with G\@.         % \@ marks end-of-sentence punctuation.
You should check the spaces after periods when
reading your output to make sure you haven't
forgotten any special cases.  Generating an
ellipsis
       \ldots\               % `\ ' is needed after `\ldots' because TeX
                             % ignores spaces after command names like \ldots
                             % made from \ + letters.
                             %
                             % Note how a `%' character causes TeX to ignore
                             % the end of the input line, so these blank lines
                             % do not start a new paragraph.
                             %
with the right spacing around the periods requires
a special command.

\LaTeX\ interprets some common characters as
commands, so you must type special commands to
generate them.  These characters include the
following:
       \$ \& \% \# \{ and \}.

In printing, text is usually emphasized with an
       \emph{italic}
type style.

\begin{em}
   A long segment of text can also be emphasized
   in this way.  Text within such a segment can be
   given \emph{additional} emphasis.
\end{em}

It is sometimes necessary to prevent \LaTeX\ from
breaking a line where it might otherwise do so.
This may be at a space, as between the ``Mr.''\ and
``Jones'' in
       ``Mr.~Jones'',        % ~ produces an unbreakable interword space.
or within a word---especially when the word is a
symbol like
       \mbox{\emph{itemnum}}
that makes little sense when hyphenated across
lines.

Footnotes\footnote{This is an example of a footnote.}
pose no problem.

\LaTeX\ is good at typesetting mathematical formulas
like
       \( x-3y + z = 7 \)
or
       \( a_{1} > x^{2n} + y^{2n} > x' \)
or
       \( \ip{A}{B} = \sum_{i} a_{i} b_{i} \).
The spaces you type in a formula are
ignored.  Remember that a letter like
       $x$                   % $ ... $  and  \( ... \)  are equivalent
is a formula when it denotes a mathematical
symbol, and it should be typed as one.

\section{Displayed Text}

Text is displayed by indenting it from the left
margin.  Quotations are commonly displayed.  There
are short quotations
\begin{quote}
   This is a short quotation.  It consists of a
   single paragraph of text.  See how it is formatted.
\end{quote}
and longer ones.
\begin{quotation}
   This is a longer quotation.  It consists of two
   paragraphs of text, neither of which are
   particularly interesting.

   This is the second paragraph of the quotation.  It
   is just as dull as the first paragraph.
\end{quotation}
Another frequently-displayed structure is a list.
The following is an example of an \emph{itemized}
list.
\begin{itemize}
   \item This is the first item of an itemized list.
         Each item in the list is marked with a ``tick''.
         You don't have to worry about what kind of tick
         mark is used.

   \item This is the second item of the list.  It
         contains another list nested inside it.  The inner
         list is an \emph{enumerated} list.
         \begin{enumerate}
            \item This is the first item of an enumerated
                  list that is nested within the itemized list.

            \item This is the second item of the inner list.
                  \LaTeX\ allows you to nest lists deeper than
                  you really should.
         \end{enumerate}
         This is the rest of the second item of the outer
         list.  It is no more interesting than any other
         part of the item.
   \item This is the third item of the list.
\end{itemize}
You can even display poetry.
\begin{verse}
   There is an environment
    for verse \\             % The \\ command separates lines
   Whose features some poets % within a stanza.
   will curse.

                             % One or more blank lines separate stanzas.

   For instead of making\\
   Them do \emph{all} line breaking, \\
   It allows them to put too many words on a line when they'd rather be
   forced to be terse.
\end{verse}

Mathematical formulas may also be displayed.  A
displayed formula
is
one-line long; multiline
formulas require special formatting instructions.
   \[  \ip{\Gamma}{\psi'} = x'' + y^{2} + z_{i}^{n}\]
Don't start a paragraph with a displayed equation,
nor make one a paragraph by itself.

\end{document}               % End of document.

\section{Bench Huge Section 1080}
% synthetic bench section 1080
% This is a sample LaTeX input file.  (Version of 12 August 2004.)
%
% A '%' character causes TeX to ignore all remaining text on the line,
% and is used for comments like this one.

\documentclass{article}      % Specifies the document class

                             % The preamble begins here.
\title{An Example Document}  % Declares the document's title.
\author{Leslie Lamport}      % Declares the author's name.
\date{January 21, 1994}      % Deleting this command produces today's date.

\newcommand{\ip}[2]{(#1, #2)}
                             % Defines \ip{arg1}{arg2} to mean
                             % (arg1, arg2).

%\newcommand{\ip}[2]{\langle #1 | #2\rangle}
                             % This is an alternative definition of
                             % \ip that is commented out.

\begin{document}             % End of preamble and beginning of text.

\maketitle                   % Produces the title.

This is an example input file.  Comparing it with
the output it generates can show you how to
produce a simple document of your own.

\section{Ordinary Text}      % Produces section heading.  Lower-level
                             % sections are begun with similar
                             % \subsection and \subsubsection commands.

The ends  of words and sentences are marked
  by   spaces. It  doesn't matter how many
spaces    you type; one is as good as 100.  The
end of   a line counts as a space.

One   or more   blank lines denote the  end
of  a paragraph.

Since any number of consecutive spaces are treated
like a single one, the formatting of the input
file makes no difference to
      \LaTeX,                % The \LaTeX command generates the LaTeX logo.
but it makes a difference to you.  When you use
\LaTeX, making your input file as easy to read
as possible will be a great help as you write
your document and when you change it.  This sample
file shows how you can add comments to your own input
file.

Because printing is different from typewriting,
there are a number of things that you have to do
differently when preparing an input file than if
you were just typing the document directly.
Quotation marks like
       ``this''
have to be handled specially, as do quotes within
quotes:
       ``\,`this'            % \, separates the double and single quote.
        is what I just
        wrote, not  `that'\,''.

Dashes come in three sizes: an
       intra-word
dash, a medium dash for number ranges like
       1--2,
and a punctuation
       dash---like
this.

A sentence-ending space should be larger than the
space between words within a sentence.  You
sometimes have to type special commands in
conjunction with punctuation characters to get
this right, as in the following sentence.
       Gnats, gnus, etc.\ all  % `\ ' makes an inter-word space.
       begin with G\@.         % \@ marks end-of-sentence punctuation.
You should check the spaces after periods when
reading your output to make sure you haven't
forgotten any special cases.  Generating an
ellipsis
       \ldots\               % `\ ' is needed after `\ldots' because TeX
                             % ignores spaces after command names like \ldots
                             % made from \ + letters.
                             %
                             % Note how a `%' character causes TeX to ignore
                             % the end of the input line, so these blank lines
                             % do not start a new paragraph.
                             %
with the right spacing around the periods requires
a special command.

\LaTeX\ interprets some common characters as
commands, so you must type special commands to
generate them.  These characters include the
following:
       \$ \& \% \# \{ and \}.

In printing, text is usually emphasized with an
       \emph{italic}
type style.

\begin{em}
   A long segment of text can also be emphasized
   in this way.  Text within such a segment can be
   given \emph{additional} emphasis.
\end{em}

It is sometimes necessary to prevent \LaTeX\ from
breaking a line where it might otherwise do so.
This may be at a space, as between the ``Mr.''\ and
``Jones'' in
       ``Mr.~Jones'',        % ~ produces an unbreakable interword space.
or within a word---especially when the word is a
symbol like
       \mbox{\emph{itemnum}}
that makes little sense when hyphenated across
lines.

Footnotes\footnote{This is an example of a footnote.}
pose no problem.

\LaTeX\ is good at typesetting mathematical formulas
like
       \( x-3y + z = 7 \)
or
       \( a_{1} > x^{2n} + y^{2n} > x' \)
or
       \( \ip{A}{B} = \sum_{i} a_{i} b_{i} \).
The spaces you type in a formula are
ignored.  Remember that a letter like
       $x$                   % $ ... $  and  \( ... \)  are equivalent
is a formula when it denotes a mathematical
symbol, and it should be typed as one.

\section{Displayed Text}

Text is displayed by indenting it from the left
margin.  Quotations are commonly displayed.  There
are short quotations
\begin{quote}
   This is a short quotation.  It consists of a
   single paragraph of text.  See how it is formatted.
\end{quote}
and longer ones.
\begin{quotation}
   This is a longer quotation.  It consists of two
   paragraphs of text, neither of which are
   particularly interesting.

   This is the second paragraph of the quotation.  It
   is just as dull as the first paragraph.
\end{quotation}
Another frequently-displayed structure is a list.
The following is an example of an \emph{itemized}
list.
\begin{itemize}
   \item This is the first item of an itemized list.
         Each item in the list is marked with a ``tick''.
         You don't have to worry about what kind of tick
         mark is used.

   \item This is the second item of the list.  It
         contains another list nested inside it.  The inner
         list is an \emph{enumerated} list.
         \begin{enumerate}
            \item This is the first item of an enumerated
                  list that is nested within the itemized list.

            \item This is the second item of the inner list.
                  \LaTeX\ allows you to nest lists deeper than
                  you really should.
         \end{enumerate}
         This is the rest of the second item of the outer
         list.  It is no more interesting than any other
         part of the item.
   \item This is the third item of the list.
\end{itemize}
You can even display poetry.
\begin{verse}
   There is an environment
    for verse \\             % The \\ command separates lines
   Whose features some poets % within a stanza.
   will curse.

                             % One or more blank lines separate stanzas.

   For instead of making\\
   Them do \emph{all} line breaking, \\
   It allows them to put too many words on a line when they'd rather be
   forced to be terse.
\end{verse}

Mathematical formulas may also be displayed.  A
displayed formula
is
one-line long; multiline
formulas require special formatting instructions.
   \[  \ip{\Gamma}{\psi'} = x'' + y^{2} + z_{i}^{n}\]
Don't start a paragraph with a displayed equation,
nor make one a paragraph by itself.

\end{document}               % End of document.

\section{Bench Huge Section 1081}
% synthetic bench section 1081
% This is a sample LaTeX input file.  (Version of 12 August 2004.)
%
% A '%' character causes TeX to ignore all remaining text on the line,
% and is used for comments like this one.

\documentclass{article}      % Specifies the document class

                             % The preamble begins here.
\title{An Example Document}  % Declares the document's title.
\author{Leslie Lamport}      % Declares the author's name.
\date{January 21, 1994}      % Deleting this command produces today's date.

\newcommand{\ip}[2]{(#1, #2)}
                             % Defines \ip{arg1}{arg2} to mean
                             % (arg1, arg2).

%\newcommand{\ip}[2]{\langle #1 | #2\rangle}
                             % This is an alternative definition of
                             % \ip that is commented out.

\begin{document}             % End of preamble and beginning of text.

\maketitle                   % Produces the title.

This is an example input file.  Comparing it with
the output it generates can show you how to
produce a simple document of your own.

\section{Ordinary Text}      % Produces section heading.  Lower-level
                             % sections are begun with similar
                             % \subsection and \subsubsection commands.

The ends  of words and sentences are marked
  by   spaces. It  doesn't matter how many
spaces    you type; one is as good as 100.  The
end of   a line counts as a space.

One   or more   blank lines denote the  end
of  a paragraph.

Since any number of consecutive spaces are treated
like a single one, the formatting of the input
file makes no difference to
      \LaTeX,                % The \LaTeX command generates the LaTeX logo.
but it makes a difference to you.  When you use
\LaTeX, making your input file as easy to read
as possible will be a great help as you write
your document and when you change it.  This sample
file shows how you can add comments to your own input
file.

Because printing is different from typewriting,
there are a number of things that you have to do
differently when preparing an input file than if
you were just typing the document directly.
Quotation marks like
       ``this''
have to be handled specially, as do quotes within
quotes:
       ``\,`this'            % \, separates the double and single quote.
        is what I just
        wrote, not  `that'\,''.

Dashes come in three sizes: an
       intra-word
dash, a medium dash for number ranges like
       1--2,
and a punctuation
       dash---like
this.

A sentence-ending space should be larger than the
space between words within a sentence.  You
sometimes have to type special commands in
conjunction with punctuation characters to get
this right, as in the following sentence.
       Gnats, gnus, etc.\ all  % `\ ' makes an inter-word space.
       begin with G\@.         % \@ marks end-of-sentence punctuation.
You should check the spaces after periods when
reading your output to make sure you haven't
forgotten any special cases.  Generating an
ellipsis
       \ldots\               % `\ ' is needed after `\ldots' because TeX
                             % ignores spaces after command names like \ldots
                             % made from \ + letters.
                             %
                             % Note how a `%' character causes TeX to ignore
                             % the end of the input line, so these blank lines
                             % do not start a new paragraph.
                             %
with the right spacing around the periods requires
a special command.

\LaTeX\ interprets some common characters as
commands, so you must type special commands to
generate them.  These characters include the
following:
       \$ \& \% \# \{ and \}.

In printing, text is usually emphasized with an
       \emph{italic}
type style.

\begin{em}
   A long segment of text can also be emphasized
   in this way.  Text within such a segment can be
   given \emph{additional} emphasis.
\end{em}

It is sometimes necessary to prevent \LaTeX\ from
breaking a line where it might otherwise do so.
This may be at a space, as between the ``Mr.''\ and
``Jones'' in
       ``Mr.~Jones'',        % ~ produces an unbreakable interword space.
or within a word---especially when the word is a
symbol like
       \mbox{\emph{itemnum}}
that makes little sense when hyphenated across
lines.

Footnotes\footnote{This is an example of a footnote.}
pose no problem.

\LaTeX\ is good at typesetting mathematical formulas
like
       \( x-3y + z = 7 \)
or
       \( a_{1} > x^{2n} + y^{2n} > x' \)
or
       \( \ip{A}{B} = \sum_{i} a_{i} b_{i} \).
The spaces you type in a formula are
ignored.  Remember that a letter like
       $x$                   % $ ... $  and  \( ... \)  are equivalent
is a formula when it denotes a mathematical
symbol, and it should be typed as one.

\section{Displayed Text}

Text is displayed by indenting it from the left
margin.  Quotations are commonly displayed.  There
are short quotations
\begin{quote}
   This is a short quotation.  It consists of a
   single paragraph of text.  See how it is formatted.
\end{quote}
and longer ones.
\begin{quotation}
   This is a longer quotation.  It consists of two
   paragraphs of text, neither of which are
   particularly interesting.

   This is the second paragraph of the quotation.  It
   is just as dull as the first paragraph.
\end{quotation}
Another frequently-displayed structure is a list.
The following is an example of an \emph{itemized}
list.
\begin{itemize}
   \item This is the first item of an itemized list.
         Each item in the list is marked with a ``tick''.
         You don't have to worry about what kind of tick
         mark is used.

   \item This is the second item of the list.  It
         contains another list nested inside it.  The inner
         list is an \emph{enumerated} list.
         \begin{enumerate}
            \item This is the first item of an enumerated
                  list that is nested within the itemized list.

            \item This is the second item of the inner list.
                  \LaTeX\ allows you to nest lists deeper than
                  you really should.
         \end{enumerate}
         This is the rest of the second item of the outer
         list.  It is no more interesting than any other
         part of the item.
   \item This is the third item of the list.
\end{itemize}
You can even display poetry.
\begin{verse}
   There is an environment
    for verse \\             % The \\ command separates lines
   Whose features some poets % within a stanza.
   will curse.

                             % One or more blank lines separate stanzas.

   For instead of making\\
   Them do \emph{all} line breaking, \\
   It allows them to put too many words on a line when they'd rather be
   forced to be terse.
\end{verse}

Mathematical formulas may also be displayed.  A
displayed formula
is
one-line long; multiline
formulas require special formatting instructions.
   \[  \ip{\Gamma}{\psi'} = x'' + y^{2} + z_{i}^{n}\]
Don't start a paragraph with a displayed equation,
nor make one a paragraph by itself.

\end{document}               % End of document.

\section{Bench Huge Section 1082}
% synthetic bench section 1082
% This is a sample LaTeX input file.  (Version of 12 August 2004.)
%
% A '%' character causes TeX to ignore all remaining text on the line,
% and is used for comments like this one.

\documentclass{article}      % Specifies the document class

                             % The preamble begins here.
\title{An Example Document}  % Declares the document's title.
\author{Leslie Lamport}      % Declares the author's name.
\date{January 21, 1994}      % Deleting this command produces today's date.

\newcommand{\ip}[2]{(#1, #2)}
                             % Defines \ip{arg1}{arg2} to mean
                             % (arg1, arg2).

%\newcommand{\ip}[2]{\langle #1 | #2\rangle}
                             % This is an alternative definition of
                             % \ip that is commented out.

\begin{document}             % End of preamble and beginning of text.

\maketitle                   % Produces the title.

This is an example input file.  Comparing it with
the output it generates can show you how to
produce a simple document of your own.

\section{Ordinary Text}      % Produces section heading.  Lower-level
                             % sections are begun with similar
                             % \subsection and \subsubsection commands.

The ends  of words and sentences are marked
  by   spaces. It  doesn't matter how many
spaces    you type; one is as good as 100.  The
end of   a line counts as a space.

One   or more   blank lines denote the  end
of  a paragraph.

Since any number of consecutive spaces are treated
like a single one, the formatting of the input
file makes no difference to
      \LaTeX,                % The \LaTeX command generates the LaTeX logo.
but it makes a difference to you.  When you use
\LaTeX, making your input file as easy to read
as possible will be a great help as you write
your document and when you change it.  This sample
file shows how you can add comments to your own input
file.

Because printing is different from typewriting,
there are a number of things that you have to do
differently when preparing an input file than if
you were just typing the document directly.
Quotation marks like
       ``this''
have to be handled specially, as do quotes within
quotes:
       ``\,`this'            % \, separates the double and single quote.
        is what I just
        wrote, not  `that'\,''.

Dashes come in three sizes: an
       intra-word
dash, a medium dash for number ranges like
       1--2,
and a punctuation
       dash---like
this.

A sentence-ending space should be larger than the
space between words within a sentence.  You
sometimes have to type special commands in
conjunction with punctuation characters to get
this right, as in the following sentence.
       Gnats, gnus, etc.\ all  % `\ ' makes an inter-word space.
       begin with G\@.         % \@ marks end-of-sentence punctuation.
You should check the spaces after periods when
reading your output to make sure you haven't
forgotten any special cases.  Generating an
ellipsis
       \ldots\               % `\ ' is needed after `\ldots' because TeX
                             % ignores spaces after command names like \ldots
                             % made from \ + letters.
                             %
                             % Note how a `%' character causes TeX to ignore
                             % the end of the input line, so these blank lines
                             % do not start a new paragraph.
                             %
with the right spacing around the periods requires
a special command.

\LaTeX\ interprets some common characters as
commands, so you must type special commands to
generate them.  These characters include the
following:
       \$ \& \% \# \{ and \}.

In printing, text is usually emphasized with an
       \emph{italic}
type style.

\begin{em}
   A long segment of text can also be emphasized
   in this way.  Text within such a segment can be
   given \emph{additional} emphasis.
\end{em}

It is sometimes necessary to prevent \LaTeX\ from
breaking a line where it might otherwise do so.
This may be at a space, as between the ``Mr.''\ and
``Jones'' in
       ``Mr.~Jones'',        % ~ produces an unbreakable interword space.
or within a word---especially when the word is a
symbol like
       \mbox{\emph{itemnum}}
that makes little sense when hyphenated across
lines.

Footnotes\footnote{This is an example of a footnote.}
pose no problem.

\LaTeX\ is good at typesetting mathematical formulas
like
       \( x-3y + z = 7 \)
or
       \( a_{1} > x^{2n} + y^{2n} > x' \)
or
       \( \ip{A}{B} = \sum_{i} a_{i} b_{i} \).
The spaces you type in a formula are
ignored.  Remember that a letter like
       $x$                   % $ ... $  and  \( ... \)  are equivalent
is a formula when it denotes a mathematical
symbol, and it should be typed as one.

\section{Displayed Text}

Text is displayed by indenting it from the left
margin.  Quotations are commonly displayed.  There
are short quotations
\begin{quote}
   This is a short quotation.  It consists of a
   single paragraph of text.  See how it is formatted.
\end{quote}
and longer ones.
\begin{quotation}
   This is a longer quotation.  It consists of two
   paragraphs of text, neither of which are
   particularly interesting.

   This is the second paragraph of the quotation.  It
   is just as dull as the first paragraph.
\end{quotation}
Another frequently-displayed structure is a list.
The following is an example of an \emph{itemized}
list.
\begin{itemize}
   \item This is the first item of an itemized list.
         Each item in the list is marked with a ``tick''.
         You don't have to worry about what kind of tick
         mark is used.

   \item This is the second item of the list.  It
         contains another list nested inside it.  The inner
         list is an \emph{enumerated} list.
         \begin{enumerate}
            \item This is the first item of an enumerated
                  list that is nested within the itemized list.

            \item This is the second item of the inner list.
                  \LaTeX\ allows you to nest lists deeper than
                  you really should.
         \end{enumerate}
         This is the rest of the second item of the outer
         list.  It is no more interesting than any other
         part of the item.
   \item This is the third item of the list.
\end{itemize}
You can even display poetry.
\begin{verse}
   There is an environment
    for verse \\             % The \\ command separates lines
   Whose features some poets % within a stanza.
   will curse.

                             % One or more blank lines separate stanzas.

   For instead of making\\
   Them do \emph{all} line breaking, \\
   It allows them to put too many words on a line when they'd rather be
   forced to be terse.
\end{verse}

Mathematical formulas may also be displayed.  A
displayed formula
is
one-line long; multiline
formulas require special formatting instructions.
   \[  \ip{\Gamma}{\psi'} = x'' + y^{2} + z_{i}^{n}\]
Don't start a paragraph with a displayed equation,
nor make one a paragraph by itself.

\end{document}               % End of document.

\section{Bench Huge Section 1083}
% synthetic bench section 1083
% This is a sample LaTeX input file.  (Version of 12 August 2004.)
%
% A '%' character causes TeX to ignore all remaining text on the line,
% and is used for comments like this one.

\documentclass{article}      % Specifies the document class

                             % The preamble begins here.
\title{An Example Document}  % Declares the document's title.
\author{Leslie Lamport}      % Declares the author's name.
\date{January 21, 1994}      % Deleting this command produces today's date.

\newcommand{\ip}[2]{(#1, #2)}
                             % Defines \ip{arg1}{arg2} to mean
                             % (arg1, arg2).

%\newcommand{\ip}[2]{\langle #1 | #2\rangle}
                             % This is an alternative definition of
                             % \ip that is commented out.

\begin{document}             % End of preamble and beginning of text.

\maketitle                   % Produces the title.

This is an example input file.  Comparing it with
the output it generates can show you how to
produce a simple document of your own.

\section{Ordinary Text}      % Produces section heading.  Lower-level
                             % sections are begun with similar
                             % \subsection and \subsubsection commands.

The ends  of words and sentences are marked
  by   spaces. It  doesn't matter how many
spaces    you type; one is as good as 100.  The
end of   a line counts as a space.

One   or more   blank lines denote the  end
of  a paragraph.

Since any number of consecutive spaces are treated
like a single one, the formatting of the input
file makes no difference to
      \LaTeX,                % The \LaTeX command generates the LaTeX logo.
but it makes a difference to you.  When you use
\LaTeX, making your input file as easy to read
as possible will be a great help as you write
your document and when you change it.  This sample
file shows how you can add comments to your own input
file.

Because printing is different from typewriting,
there are a number of things that you have to do
differently when preparing an input file than if
you were just typing the document directly.
Quotation marks like
       ``this''
have to be handled specially, as do quotes within
quotes:
       ``\,`this'            % \, separates the double and single quote.
        is what I just
        wrote, not  `that'\,''.

Dashes come in three sizes: an
       intra-word
dash, a medium dash for number ranges like
       1--2,
and a punctuation
       dash---like
this.

A sentence-ending space should be larger than the
space between words within a sentence.  You
sometimes have to type special commands in
conjunction with punctuation characters to get
this right, as in the following sentence.
       Gnats, gnus, etc.\ all  % `\ ' makes an inter-word space.
       begin with G\@.         % \@ marks end-of-sentence punctuation.
You should check the spaces after periods when
reading your output to make sure you haven't
forgotten any special cases.  Generating an
ellipsis
       \ldots\               % `\ ' is needed after `\ldots' because TeX
                             % ignores spaces after command names like \ldots
                             % made from \ + letters.
                             %
                             % Note how a `%' character causes TeX to ignore
                             % the end of the input line, so these blank lines
                             % do not start a new paragraph.
                             %
with the right spacing around the periods requires
a special command.

\LaTeX\ interprets some common characters as
commands, so you must type special commands to
generate them.  These characters include the
following:
       \$ \& \% \# \{ and \}.

In printing, text is usually emphasized with an
       \emph{italic}
type style.

\begin{em}
   A long segment of text can also be emphasized
   in this way.  Text within such a segment can be
   given \emph{additional} emphasis.
\end{em}

It is sometimes necessary to prevent \LaTeX\ from
breaking a line where it might otherwise do so.
This may be at a space, as between the ``Mr.''\ and
``Jones'' in
       ``Mr.~Jones'',        % ~ produces an unbreakable interword space.
or within a word---especially when the word is a
symbol like
       \mbox{\emph{itemnum}}
that makes little sense when hyphenated across
lines.

Footnotes\footnote{This is an example of a footnote.}
pose no problem.

\LaTeX\ is good at typesetting mathematical formulas
like
       \( x-3y + z = 7 \)
or
       \( a_{1} > x^{2n} + y^{2n} > x' \)
or
       \( \ip{A}{B} = \sum_{i} a_{i} b_{i} \).
The spaces you type in a formula are
ignored.  Remember that a letter like
       $x$                   % $ ... $  and  \( ... \)  are equivalent
is a formula when it denotes a mathematical
symbol, and it should be typed as one.

\section{Displayed Text}

Text is displayed by indenting it from the left
margin.  Quotations are commonly displayed.  There
are short quotations
\begin{quote}
   This is a short quotation.  It consists of a
   single paragraph of text.  See how it is formatted.
\end{quote}
and longer ones.
\begin{quotation}
   This is a longer quotation.  It consists of two
   paragraphs of text, neither of which are
   particularly interesting.

   This is the second paragraph of the quotation.  It
   is just as dull as the first paragraph.
\end{quotation}
Another frequently-displayed structure is a list.
The following is an example of an \emph{itemized}
list.
\begin{itemize}
   \item This is the first item of an itemized list.
         Each item in the list is marked with a ``tick''.
         You don't have to worry about what kind of tick
         mark is used.

   \item This is the second item of the list.  It
         contains another list nested inside it.  The inner
         list is an \emph{enumerated} list.
         \begin{enumerate}
            \item This is the first item of an enumerated
                  list that is nested within the itemized list.

            \item This is the second item of the inner list.
                  \LaTeX\ allows you to nest lists deeper than
                  you really should.
         \end{enumerate}
         This is the rest of the second item of the outer
         list.  It is no more interesting than any other
         part of the item.
   \item This is the third item of the list.
\end{itemize}
You can even display poetry.
\begin{verse}
   There is an environment
    for verse \\             % The \\ command separates lines
   Whose features some poets % within a stanza.
   will curse.

                             % One or more blank lines separate stanzas.

   For instead of making\\
   Them do \emph{all} line breaking, \\
   It allows them to put too many words on a line when they'd rather be
   forced to be terse.
\end{verse}

Mathematical formulas may also be displayed.  A
displayed formula
is
one-line long; multiline
formulas require special formatting instructions.
   \[  \ip{\Gamma}{\psi'} = x'' + y^{2} + z_{i}^{n}\]
Don't start a paragraph with a displayed equation,
nor make one a paragraph by itself.

\end{document}               % End of document.

\section{Bench Huge Section 1084}
% synthetic bench section 1084
% This is a sample LaTeX input file.  (Version of 12 August 2004.)
%
% A '%' character causes TeX to ignore all remaining text on the line,
% and is used for comments like this one.

\documentclass{article}      % Specifies the document class

                             % The preamble begins here.
\title{An Example Document}  % Declares the document's title.
\author{Leslie Lamport}      % Declares the author's name.
\date{January 21, 1994}      % Deleting this command produces today's date.

\newcommand{\ip}[2]{(#1, #2)}
                             % Defines \ip{arg1}{arg2} to mean
                             % (arg1, arg2).

%\newcommand{\ip}[2]{\langle #1 | #2\rangle}
                             % This is an alternative definition of
                             % \ip that is commented out.

\begin{document}             % End of preamble and beginning of text.

\maketitle                   % Produces the title.

This is an example input file.  Comparing it with
the output it generates can show you how to
produce a simple document of your own.

\section{Ordinary Text}      % Produces section heading.  Lower-level
                             % sections are begun with similar
                             % \subsection and \subsubsection commands.

The ends  of words and sentences are marked
  by   spaces. It  doesn't matter how many
spaces    you type; one is as good as 100.  The
end of   a line counts as a space.

One   or more   blank lines denote the  end
of  a paragraph.

Since any number of consecutive spaces are treated
like a single one, the formatting of the input
file makes no difference to
      \LaTeX,                % The \LaTeX command generates the LaTeX logo.
but it makes a difference to you.  When you use
\LaTeX, making your input file as easy to read
as possible will be a great help as you write
your document and when you change it.  This sample
file shows how you can add comments to your own input
file.

Because printing is different from typewriting,
there are a number of things that you have to do
differently when preparing an input file than if
you were just typing the document directly.
Quotation marks like
       ``this''
have to be handled specially, as do quotes within
quotes:
       ``\,`this'            % \, separates the double and single quote.
        is what I just
        wrote, not  `that'\,''.

Dashes come in three sizes: an
       intra-word
dash, a medium dash for number ranges like
       1--2,
and a punctuation
       dash---like
this.

A sentence-ending space should be larger than the
space between words within a sentence.  You
sometimes have to type special commands in
conjunction with punctuation characters to get
this right, as in the following sentence.
       Gnats, gnus, etc.\ all  % `\ ' makes an inter-word space.
       begin with G\@.         % \@ marks end-of-sentence punctuation.
You should check the spaces after periods when
reading your output to make sure you haven't
forgotten any special cases.  Generating an
ellipsis
       \ldots\               % `\ ' is needed after `\ldots' because TeX
                             % ignores spaces after command names like \ldots
                             % made from \ + letters.
                             %
                             % Note how a `%' character causes TeX to ignore
                             % the end of the input line, so these blank lines
                             % do not start a new paragraph.
                             %
with the right spacing around the periods requires
a special command.

\LaTeX\ interprets some common characters as
commands, so you must type special commands to
generate them.  These characters include the
following:
       \$ \& \% \# \{ and \}.

In printing, text is usually emphasized with an
       \emph{italic}
type style.

\begin{em}
   A long segment of text can also be emphasized
   in this way.  Text within such a segment can be
   given \emph{additional} emphasis.
\end{em}

It is sometimes necessary to prevent \LaTeX\ from
breaking a line where it might otherwise do so.
This may be at a space, as between the ``Mr.''\ and
``Jones'' in
       ``Mr.~Jones'',        % ~ produces an unbreakable interword space.
or within a word---especially when the word is a
symbol like
       \mbox{\emph{itemnum}}
that makes little sense when hyphenated across
lines.

Footnotes\footnote{This is an example of a footnote.}
pose no problem.

\LaTeX\ is good at typesetting mathematical formulas
like
       \( x-3y + z = 7 \)
or
       \( a_{1} > x^{2n} + y^{2n} > x' \)
or
       \( \ip{A}{B} = \sum_{i} a_{i} b_{i} \).
The spaces you type in a formula are
ignored.  Remember that a letter like
       $x$                   % $ ... $  and  \( ... \)  are equivalent
is a formula when it denotes a mathematical
symbol, and it should be typed as one.

\section{Displayed Text}

Text is displayed by indenting it from the left
margin.  Quotations are commonly displayed.  There
are short quotations
\begin{quote}
   This is a short quotation.  It consists of a
   single paragraph of text.  See how it is formatted.
\end{quote}
and longer ones.
\begin{quotation}
   This is a longer quotation.  It consists of two
   paragraphs of text, neither of which are
   particularly interesting.

   This is the second paragraph of the quotation.  It
   is just as dull as the first paragraph.
\end{quotation}
Another frequently-displayed structure is a list.
The following is an example of an \emph{itemized}
list.
\begin{itemize}
   \item This is the first item of an itemized list.
         Each item in the list is marked with a ``tick''.
         You don't have to worry about what kind of tick
         mark is used.

   \item This is the second item of the list.  It
         contains another list nested inside it.  The inner
         list is an \emph{enumerated} list.
         \begin{enumerate}
            \item This is the first item of an enumerated
                  list that is nested within the itemized list.

            \item This is the second item of the inner list.
                  \LaTeX\ allows you to nest lists deeper than
                  you really should.
         \end{enumerate}
         This is the rest of the second item of the outer
         list.  It is no more interesting than any other
         part of the item.
   \item This is the third item of the list.
\end{itemize}
You can even display poetry.
\begin{verse}
   There is an environment
    for verse \\             % The \\ command separates lines
   Whose features some poets % within a stanza.
   will curse.

                             % One or more blank lines separate stanzas.

   For instead of making\\
   Them do \emph{all} line breaking, \\
   It allows them to put too many words on a line when they'd rather be
   forced to be terse.
\end{verse}

Mathematical formulas may also be displayed.  A
displayed formula
is
one-line long; multiline
formulas require special formatting instructions.
   \[  \ip{\Gamma}{\psi'} = x'' + y^{2} + z_{i}^{n}\]
Don't start a paragraph with a displayed equation,
nor make one a paragraph by itself.

\end{document}               % End of document.

\section{Bench Huge Section 1085}
% synthetic bench section 1085
% This is a sample LaTeX input file.  (Version of 12 August 2004.)
%
% A '%' character causes TeX to ignore all remaining text on the line,
% and is used for comments like this one.

\documentclass{article}      % Specifies the document class

                             % The preamble begins here.
\title{An Example Document}  % Declares the document's title.
\author{Leslie Lamport}      % Declares the author's name.
\date{January 21, 1994}      % Deleting this command produces today's date.

\newcommand{\ip}[2]{(#1, #2)}
                             % Defines \ip{arg1}{arg2} to mean
                             % (arg1, arg2).

%\newcommand{\ip}[2]{\langle #1 | #2\rangle}
                             % This is an alternative definition of
                             % \ip that is commented out.

\begin{document}             % End of preamble and beginning of text.

\maketitle                   % Produces the title.

This is an example input file.  Comparing it with
the output it generates can show you how to
produce a simple document of your own.

\section{Ordinary Text}      % Produces section heading.  Lower-level
                             % sections are begun with similar
                             % \subsection and \subsubsection commands.

The ends  of words and sentences are marked
  by   spaces. It  doesn't matter how many
spaces    you type; one is as good as 100.  The
end of   a line counts as a space.

One   or more   blank lines denote the  end
of  a paragraph.

Since any number of consecutive spaces are treated
like a single one, the formatting of the input
file makes no difference to
      \LaTeX,                % The \LaTeX command generates the LaTeX logo.
but it makes a difference to you.  When you use
\LaTeX, making your input file as easy to read
as possible will be a great help as you write
your document and when you change it.  This sample
file shows how you can add comments to your own input
file.

Because printing is different from typewriting,
there are a number of things that you have to do
differently when preparing an input file than if
you were just typing the document directly.
Quotation marks like
       ``this''
have to be handled specially, as do quotes within
quotes:
       ``\,`this'            % \, separates the double and single quote.
        is what I just
        wrote, not  `that'\,''.

Dashes come in three sizes: an
       intra-word
dash, a medium dash for number ranges like
       1--2,
and a punctuation
       dash---like
this.

A sentence-ending space should be larger than the
space between words within a sentence.  You
sometimes have to type special commands in
conjunction with punctuation characters to get
this right, as in the following sentence.
       Gnats, gnus, etc.\ all  % `\ ' makes an inter-word space.
       begin with G\@.         % \@ marks end-of-sentence punctuation.
You should check the spaces after periods when
reading your output to make sure you haven't
forgotten any special cases.  Generating an
ellipsis
       \ldots\               % `\ ' is needed after `\ldots' because TeX
                             % ignores spaces after command names like \ldots
                             % made from \ + letters.
                             %
                             % Note how a `%' character causes TeX to ignore
                             % the end of the input line, so these blank lines
                             % do not start a new paragraph.
                             %
with the right spacing around the periods requires
a special command.

\LaTeX\ interprets some common characters as
commands, so you must type special commands to
generate them.  These characters include the
following:
       \$ \& \% \# \{ and \}.

In printing, text is usually emphasized with an
       \emph{italic}
type style.

\begin{em}
   A long segment of text can also be emphasized
   in this way.  Text within such a segment can be
   given \emph{additional} emphasis.
\end{em}

It is sometimes necessary to prevent \LaTeX\ from
breaking a line where it might otherwise do so.
This may be at a space, as between the ``Mr.''\ and
``Jones'' in
       ``Mr.~Jones'',        % ~ produces an unbreakable interword space.
or within a word---especially when the word is a
symbol like
       \mbox{\emph{itemnum}}
that makes little sense when hyphenated across
lines.

Footnotes\footnote{This is an example of a footnote.}
pose no problem.

\LaTeX\ is good at typesetting mathematical formulas
like
       \( x-3y + z = 7 \)
or
       \( a_{1} > x^{2n} + y^{2n} > x' \)
or
       \( \ip{A}{B} = \sum_{i} a_{i} b_{i} \).
The spaces you type in a formula are
ignored.  Remember that a letter like
       $x$                   % $ ... $  and  \( ... \)  are equivalent
is a formula when it denotes a mathematical
symbol, and it should be typed as one.

\section{Displayed Text}

Text is displayed by indenting it from the left
margin.  Quotations are commonly displayed.  There
are short quotations
\begin{quote}
   This is a short quotation.  It consists of a
   single paragraph of text.  See how it is formatted.
\end{quote}
and longer ones.
\begin{quotation}
   This is a longer quotation.  It consists of two
   paragraphs of text, neither of which are
   particularly interesting.

   This is the second paragraph of the quotation.  It
   is just as dull as the first paragraph.
\end{quotation}
Another frequently-displayed structure is a list.
The following is an example of an \emph{itemized}
list.
\begin{itemize}
   \item This is the first item of an itemized list.
         Each item in the list is marked with a ``tick''.
         You don't have to worry about what kind of tick
         mark is used.

   \item This is the second item of the list.  It
         contains another list nested inside it.  The inner
         list is an \emph{enumerated} list.
         \begin{enumerate}
            \item This is the first item of an enumerated
                  list that is nested within the itemized list.

            \item This is the second item of the inner list.
                  \LaTeX\ allows you to nest lists deeper than
                  you really should.
         \end{enumerate}
         This is the rest of the second item of the outer
         list.  It is no more interesting than any other
         part of the item.
   \item This is the third item of the list.
\end{itemize}
You can even display poetry.
\begin{verse}
   There is an environment
    for verse \\             % The \\ command separates lines
   Whose features some poets % within a stanza.
   will curse.

                             % One or more blank lines separate stanzas.

   For instead of making\\
   Them do \emph{all} line breaking, \\
   It allows them to put too many words on a line when they'd rather be
   forced to be terse.
\end{verse}

Mathematical formulas may also be displayed.  A
displayed formula
is
one-line long; multiline
formulas require special formatting instructions.
   \[  \ip{\Gamma}{\psi'} = x'' + y^{2} + z_{i}^{n}\]
Don't start a paragraph with a displayed equation,
nor make one a paragraph by itself.

\end{document}               % End of document.

\section{Bench Huge Section 1086}
% synthetic bench section 1086
% This is a sample LaTeX input file.  (Version of 12 August 2004.)
%
% A '%' character causes TeX to ignore all remaining text on the line,
% and is used for comments like this one.

\documentclass{article}      % Specifies the document class

                             % The preamble begins here.
\title{An Example Document}  % Declares the document's title.
\author{Leslie Lamport}      % Declares the author's name.
\date{January 21, 1994}      % Deleting this command produces today's date.

\newcommand{\ip}[2]{(#1, #2)}
                             % Defines \ip{arg1}{arg2} to mean
                             % (arg1, arg2).

%\newcommand{\ip}[2]{\langle #1 | #2\rangle}
                             % This is an alternative definition of
                             % \ip that is commented out.

\begin{document}             % End of preamble and beginning of text.

\maketitle                   % Produces the title.

This is an example input file.  Comparing it with
the output it generates can show you how to
produce a simple document of your own.

\section{Ordinary Text}      % Produces section heading.  Lower-level
                             % sections are begun with similar
                             % \subsection and \subsubsection commands.

The ends  of words and sentences are marked
  by   spaces. It  doesn't matter how many
spaces    you type; one is as good as 100.  The
end of   a line counts as a space.

One   or more   blank lines denote the  end
of  a paragraph.

Since any number of consecutive spaces are treated
like a single one, the formatting of the input
file makes no difference to
      \LaTeX,                % The \LaTeX command generates the LaTeX logo.
but it makes a difference to you.  When you use
\LaTeX, making your input file as easy to read
as possible will be a great help as you write
your document and when you change it.  This sample
file shows how you can add comments to your own input
file.

Because printing is different from typewriting,
there are a number of things that you have to do
differently when preparing an input file than if
you were just typing the document directly.
Quotation marks like
       ``this''
have to be handled specially, as do quotes within
quotes:
       ``\,`this'            % \, separates the double and single quote.
        is what I just
        wrote, not  `that'\,''.

Dashes come in three sizes: an
       intra-word
dash, a medium dash for number ranges like
       1--2,
and a punctuation
       dash---like
this.

A sentence-ending space should be larger than the
space between words within a sentence.  You
sometimes have to type special commands in
conjunction with punctuation characters to get
this right, as in the following sentence.
       Gnats, gnus, etc.\ all  % `\ ' makes an inter-word space.
       begin with G\@.         % \@ marks end-of-sentence punctuation.
You should check the spaces after periods when
reading your output to make sure you haven't
forgotten any special cases.  Generating an
ellipsis
       \ldots\               % `\ ' is needed after `\ldots' because TeX
                             % ignores spaces after command names like \ldots
                             % made from \ + letters.
                             %
                             % Note how a `%' character causes TeX to ignore
                             % the end of the input line, so these blank lines
                             % do not start a new paragraph.
                             %
with the right spacing around the periods requires
a special command.

\LaTeX\ interprets some common characters as
commands, so you must type special commands to
generate them.  These characters include the
following:
       \$ \& \% \# \{ and \}.

In printing, text is usually emphasized with an
       \emph{italic}
type style.

\begin{em}
   A long segment of text can also be emphasized
   in this way.  Text within such a segment can be
   given \emph{additional} emphasis.
\end{em}

It is sometimes necessary to prevent \LaTeX\ from
breaking a line where it might otherwise do so.
This may be at a space, as between the ``Mr.''\ and
``Jones'' in
       ``Mr.~Jones'',        % ~ produces an unbreakable interword space.
or within a word---especially when the word is a
symbol like
       \mbox{\emph{itemnum}}
that makes little sense when hyphenated across
lines.

Footnotes\footnote{This is an example of a footnote.}
pose no problem.

\LaTeX\ is good at typesetting mathematical formulas
like
       \( x-3y + z = 7 \)
or
       \( a_{1} > x^{2n} + y^{2n} > x' \)
or
       \( \ip{A}{B} = \sum_{i} a_{i} b_{i} \).
The spaces you type in a formula are
ignored.  Remember that a letter like
       $x$                   % $ ... $  and  \( ... \)  are equivalent
is a formula when it denotes a mathematical
symbol, and it should be typed as one.

\section{Displayed Text}

Text is displayed by indenting it from the left
margin.  Quotations are commonly displayed.  There
are short quotations
\begin{quote}
   This is a short quotation.  It consists of a
   single paragraph of text.  See how it is formatted.
\end{quote}
and longer ones.
\begin{quotation}
   This is a longer quotation.  It consists of two
   paragraphs of text, neither of which are
   particularly interesting.

   This is the second paragraph of the quotation.  It
   is just as dull as the first paragraph.
\end{quotation}
Another frequently-displayed structure is a list.
The following is an example of an \emph{itemized}
list.
\begin{itemize}
   \item This is the first item of an itemized list.
         Each item in the list is marked with a ``tick''.
         You don't have to worry about what kind of tick
         mark is used.

   \item This is the second item of the list.  It
         contains another list nested inside it.  The inner
         list is an \emph{enumerated} list.
         \begin{enumerate}
            \item This is the first item of an enumerated
                  list that is nested within the itemized list.

            \item This is the second item of the inner list.
                  \LaTeX\ allows you to nest lists deeper than
                  you really should.
         \end{enumerate}
         This is the rest of the second item of the outer
         list.  It is no more interesting than any other
         part of the item.
   \item This is the third item of the list.
\end{itemize}
You can even display poetry.
\begin{verse}
   There is an environment
    for verse \\             % The \\ command separates lines
   Whose features some poets % within a stanza.
   will curse.

                             % One or more blank lines separate stanzas.

   For instead of making\\
   Them do \emph{all} line breaking, \\
   It allows them to put too many words on a line when they'd rather be
   forced to be terse.
\end{verse}

Mathematical formulas may also be displayed.  A
displayed formula
is
one-line long; multiline
formulas require special formatting instructions.
   \[  \ip{\Gamma}{\psi'} = x'' + y^{2} + z_{i}^{n}\]
Don't start a paragraph with a displayed equation,
nor make one a paragraph by itself.

\end{document}               % End of document.

\section{Bench Huge Section 1087}
% synthetic bench section 1087
% This is a sample LaTeX input file.  (Version of 12 August 2004.)
%
% A '%' character causes TeX to ignore all remaining text on the line,
% and is used for comments like this one.

\documentclass{article}      % Specifies the document class

                             % The preamble begins here.
\title{An Example Document}  % Declares the document's title.
\author{Leslie Lamport}      % Declares the author's name.
\date{January 21, 1994}      % Deleting this command produces today's date.

\newcommand{\ip}[2]{(#1, #2)}
                             % Defines \ip{arg1}{arg2} to mean
                             % (arg1, arg2).

%\newcommand{\ip}[2]{\langle #1 | #2\rangle}
                             % This is an alternative definition of
                             % \ip that is commented out.

\begin{document}             % End of preamble and beginning of text.

\maketitle                   % Produces the title.

This is an example input file.  Comparing it with
the output it generates can show you how to
produce a simple document of your own.

\section{Ordinary Text}      % Produces section heading.  Lower-level
                             % sections are begun with similar
                             % \subsection and \subsubsection commands.

The ends  of words and sentences are marked
  by   spaces. It  doesn't matter how many
spaces    you type; one is as good as 100.  The
end of   a line counts as a space.

One   or more   blank lines denote the  end
of  a paragraph.

Since any number of consecutive spaces are treated
like a single one, the formatting of the input
file makes no difference to
      \LaTeX,                % The \LaTeX command generates the LaTeX logo.
but it makes a difference to you.  When you use
\LaTeX, making your input file as easy to read
as possible will be a great help as you write
your document and when you change it.  This sample
file shows how you can add comments to your own input
file.

Because printing is different from typewriting,
there are a number of things that you have to do
differently when preparing an input file than if
you were just typing the document directly.
Quotation marks like
       ``this''
have to be handled specially, as do quotes within
quotes:
       ``\,`this'            % \, separates the double and single quote.
        is what I just
        wrote, not  `that'\,''.

Dashes come in three sizes: an
       intra-word
dash, a medium dash for number ranges like
       1--2,
and a punctuation
       dash---like
this.

A sentence-ending space should be larger than the
space between words within a sentence.  You
sometimes have to type special commands in
conjunction with punctuation characters to get
this right, as in the following sentence.
       Gnats, gnus, etc.\ all  % `\ ' makes an inter-word space.
       begin with G\@.         % \@ marks end-of-sentence punctuation.
You should check the spaces after periods when
reading your output to make sure you haven't
forgotten any special cases.  Generating an
ellipsis
       \ldots\               % `\ ' is needed after `\ldots' because TeX
                             % ignores spaces after command names like \ldots
                             % made from \ + letters.
                             %
                             % Note how a `%' character causes TeX to ignore
                             % the end of the input line, so these blank lines
                             % do not start a new paragraph.
                             %
with the right spacing around the periods requires
a special command.

\LaTeX\ interprets some common characters as
commands, so you must type special commands to
generate them.  These characters include the
following:
       \$ \& \% \# \{ and \}.

In printing, text is usually emphasized with an
       \emph{italic}
type style.

\begin{em}
   A long segment of text can also be emphasized
   in this way.  Text within such a segment can be
   given \emph{additional} emphasis.
\end{em}

It is sometimes necessary to prevent \LaTeX\ from
breaking a line where it might otherwise do so.
This may be at a space, as between the ``Mr.''\ and
``Jones'' in
       ``Mr.~Jones'',        % ~ produces an unbreakable interword space.
or within a word---especially when the word is a
symbol like
       \mbox{\emph{itemnum}}
that makes little sense when hyphenated across
lines.

Footnotes\footnote{This is an example of a footnote.}
pose no problem.

\LaTeX\ is good at typesetting mathematical formulas
like
       \( x-3y + z = 7 \)
or
       \( a_{1} > x^{2n} + y^{2n} > x' \)
or
       \( \ip{A}{B} = \sum_{i} a_{i} b_{i} \).
The spaces you type in a formula are
ignored.  Remember that a letter like
       $x$                   % $ ... $  and  \( ... \)  are equivalent
is a formula when it denotes a mathematical
symbol, and it should be typed as one.

\section{Displayed Text}

Text is displayed by indenting it from the left
margin.  Quotations are commonly displayed.  There
are short quotations
\begin{quote}
   This is a short quotation.  It consists of a
   single paragraph of text.  See how it is formatted.
\end{quote}
and longer ones.
\begin{quotation}
   This is a longer quotation.  It consists of two
   paragraphs of text, neither of which are
   particularly interesting.

   This is the second paragraph of the quotation.  It
   is just as dull as the first paragraph.
\end{quotation}
Another frequently-displayed structure is a list.
The following is an example of an \emph{itemized}
list.
\begin{itemize}
   \item This is the first item of an itemized list.
         Each item in the list is marked with a ``tick''.
         You don't have to worry about what kind of tick
         mark is used.

   \item This is the second item of the list.  It
         contains another list nested inside it.  The inner
         list is an \emph{enumerated} list.
         \begin{enumerate}
            \item This is the first item of an enumerated
                  list that is nested within the itemized list.

            \item This is the second item of the inner list.
                  \LaTeX\ allows you to nest lists deeper than
                  you really should.
         \end{enumerate}
         This is the rest of the second item of the outer
         list.  It is no more interesting than any other
         part of the item.
   \item This is the third item of the list.
\end{itemize}
You can even display poetry.
\begin{verse}
   There is an environment
    for verse \\             % The \\ command separates lines
   Whose features some poets % within a stanza.
   will curse.

                             % One or more blank lines separate stanzas.

   For instead of making\\
   Them do \emph{all} line breaking, \\
   It allows them to put too many words on a line when they'd rather be
   forced to be terse.
\end{verse}

Mathematical formulas may also be displayed.  A
displayed formula
is
one-line long; multiline
formulas require special formatting instructions.
   \[  \ip{\Gamma}{\psi'} = x'' + y^{2} + z_{i}^{n}\]
Don't start a paragraph with a displayed equation,
nor make one a paragraph by itself.

\end{document}               % End of document.

\section{Bench Huge Section 1088}
% synthetic bench section 1088
% This is a sample LaTeX input file.  (Version of 12 August 2004.)
%
% A '%' character causes TeX to ignore all remaining text on the line,
% and is used for comments like this one.

\documentclass{article}      % Specifies the document class

                             % The preamble begins here.
\title{An Example Document}  % Declares the document's title.
\author{Leslie Lamport}      % Declares the author's name.
\date{January 21, 1994}      % Deleting this command produces today's date.

\newcommand{\ip}[2]{(#1, #2)}
                             % Defines \ip{arg1}{arg2} to mean
                             % (arg1, arg2).

%\newcommand{\ip}[2]{\langle #1 | #2\rangle}
                             % This is an alternative definition of
                             % \ip that is commented out.

\begin{document}             % End of preamble and beginning of text.

\maketitle                   % Produces the title.

This is an example input file.  Comparing it with
the output it generates can show you how to
produce a simple document of your own.

\section{Ordinary Text}      % Produces section heading.  Lower-level
                             % sections are begun with similar
                             % \subsection and \subsubsection commands.

The ends  of words and sentences are marked
  by   spaces. It  doesn't matter how many
spaces    you type; one is as good as 100.  The
end of   a line counts as a space.

One   or more   blank lines denote the  end
of  a paragraph.

Since any number of consecutive spaces are treated
like a single one, the formatting of the input
file makes no difference to
      \LaTeX,                % The \LaTeX command generates the LaTeX logo.
but it makes a difference to you.  When you use
\LaTeX, making your input file as easy to read
as possible will be a great help as you write
your document and when you change it.  This sample
file shows how you can add comments to your own input
file.

Because printing is different from typewriting,
there are a number of things that you have to do
differently when preparing an input file than if
you were just typing the document directly.
Quotation marks like
       ``this''
have to be handled specially, as do quotes within
quotes:
       ``\,`this'            % \, separates the double and single quote.
        is what I just
        wrote, not  `that'\,''.

Dashes come in three sizes: an
       intra-word
dash, a medium dash for number ranges like
       1--2,
and a punctuation
       dash---like
this.

A sentence-ending space should be larger than the
space between words within a sentence.  You
sometimes have to type special commands in
conjunction with punctuation characters to get
this right, as in the following sentence.
       Gnats, gnus, etc.\ all  % `\ ' makes an inter-word space.
       begin with G\@.         % \@ marks end-of-sentence punctuation.
You should check the spaces after periods when
reading your output to make sure you haven't
forgotten any special cases.  Generating an
ellipsis
       \ldots\               % `\ ' is needed after `\ldots' because TeX
                             % ignores spaces after command names like \ldots
                             % made from \ + letters.
                             %
                             % Note how a `%' character causes TeX to ignore
                             % the end of the input line, so these blank lines
                             % do not start a new paragraph.
                             %
with the right spacing around the periods requires
a special command.

\LaTeX\ interprets some common characters as
commands, so you must type special commands to
generate them.  These characters include the
following:
       \$ \& \% \# \{ and \}.

In printing, text is usually emphasized with an
       \emph{italic}
type style.

\begin{em}
   A long segment of text can also be emphasized
   in this way.  Text within such a segment can be
   given \emph{additional} emphasis.
\end{em}

It is sometimes necessary to prevent \LaTeX\ from
breaking a line where it might otherwise do so.
This may be at a space, as between the ``Mr.''\ and
``Jones'' in
       ``Mr.~Jones'',        % ~ produces an unbreakable interword space.
or within a word---especially when the word is a
symbol like
       \mbox{\emph{itemnum}}
that makes little sense when hyphenated across
lines.

Footnotes\footnote{This is an example of a footnote.}
pose no problem.

\LaTeX\ is good at typesetting mathematical formulas
like
       \( x-3y + z = 7 \)
or
       \( a_{1} > x^{2n} + y^{2n} > x' \)
or
       \( \ip{A}{B} = \sum_{i} a_{i} b_{i} \).
The spaces you type in a formula are
ignored.  Remember that a letter like
       $x$                   % $ ... $  and  \( ... \)  are equivalent
is a formula when it denotes a mathematical
symbol, and it should be typed as one.

\section{Displayed Text}

Text is displayed by indenting it from the left
margin.  Quotations are commonly displayed.  There
are short quotations
\begin{quote}
   This is a short quotation.  It consists of a
   single paragraph of text.  See how it is formatted.
\end{quote}
and longer ones.
\begin{quotation}
   This is a longer quotation.  It consists of two
   paragraphs of text, neither of which are
   particularly interesting.

   This is the second paragraph of the quotation.  It
   is just as dull as the first paragraph.
\end{quotation}
Another frequently-displayed structure is a list.
The following is an example of an \emph{itemized}
list.
\begin{itemize}
   \item This is the first item of an itemized list.
         Each item in the list is marked with a ``tick''.
         You don't have to worry about what kind of tick
         mark is used.

   \item This is the second item of the list.  It
         contains another list nested inside it.  The inner
         list is an \emph{enumerated} list.
         \begin{enumerate}
            \item This is the first item of an enumerated
                  list that is nested within the itemized list.

            \item This is the second item of the inner list.
                  \LaTeX\ allows you to nest lists deeper than
                  you really should.
         \end{enumerate}
         This is the rest of the second item of the outer
         list.  It is no more interesting than any other
         part of the item.
   \item This is the third item of the list.
\end{itemize}
You can even display poetry.
\begin{verse}
   There is an environment
    for verse \\             % The \\ command separates lines
   Whose features some poets % within a stanza.
   will curse.

                             % One or more blank lines separate stanzas.

   For instead of making\\
   Them do \emph{all} line breaking, \\
   It allows them to put too many words on a line when they'd rather be
   forced to be terse.
\end{verse}

Mathematical formulas may also be displayed.  A
displayed formula
is
one-line long; multiline
formulas require special formatting instructions.
   \[  \ip{\Gamma}{\psi'} = x'' + y^{2} + z_{i}^{n}\]
Don't start a paragraph with a displayed equation,
nor make one a paragraph by itself.

\end{document}               % End of document.

\section{Bench Huge Section 1089}
% synthetic bench section 1089
% This is a sample LaTeX input file.  (Version of 12 August 2004.)
%
% A '%' character causes TeX to ignore all remaining text on the line,
% and is used for comments like this one.

\documentclass{article}      % Specifies the document class

                             % The preamble begins here.
\title{An Example Document}  % Declares the document's title.
\author{Leslie Lamport}      % Declares the author's name.
\date{January 21, 1994}      % Deleting this command produces today's date.

\newcommand{\ip}[2]{(#1, #2)}
                             % Defines \ip{arg1}{arg2} to mean
                             % (arg1, arg2).

%\newcommand{\ip}[2]{\langle #1 | #2\rangle}
                             % This is an alternative definition of
                             % \ip that is commented out.

\begin{document}             % End of preamble and beginning of text.

\maketitle                   % Produces the title.

This is an example input file.  Comparing it with
the output it generates can show you how to
produce a simple document of your own.

\section{Ordinary Text}      % Produces section heading.  Lower-level
                             % sections are begun with similar
                             % \subsection and \subsubsection commands.

The ends  of words and sentences are marked
  by   spaces. It  doesn't matter how many
spaces    you type; one is as good as 100.  The
end of   a line counts as a space.

One   or more   blank lines denote the  end
of  a paragraph.

Since any number of consecutive spaces are treated
like a single one, the formatting of the input
file makes no difference to
      \LaTeX,                % The \LaTeX command generates the LaTeX logo.
but it makes a difference to you.  When you use
\LaTeX, making your input file as easy to read
as possible will be a great help as you write
your document and when you change it.  This sample
file shows how you can add comments to your own input
file.

Because printing is different from typewriting,
there are a number of things that you have to do
differently when preparing an input file than if
you were just typing the document directly.
Quotation marks like
       ``this''
have to be handled specially, as do quotes within
quotes:
       ``\,`this'            % \, separates the double and single quote.
        is what I just
        wrote, not  `that'\,''.

Dashes come in three sizes: an
       intra-word
dash, a medium dash for number ranges like
       1--2,
and a punctuation
       dash---like
this.

A sentence-ending space should be larger than the
space between words within a sentence.  You
sometimes have to type special commands in
conjunction with punctuation characters to get
this right, as in the following sentence.
       Gnats, gnus, etc.\ all  % `\ ' makes an inter-word space.
       begin with G\@.         % \@ marks end-of-sentence punctuation.
You should check the spaces after periods when
reading your output to make sure you haven't
forgotten any special cases.  Generating an
ellipsis
       \ldots\               % `\ ' is needed after `\ldots' because TeX
                             % ignores spaces after command names like \ldots
                             % made from \ + letters.
                             %
                             % Note how a `%' character causes TeX to ignore
                             % the end of the input line, so these blank lines
                             % do not start a new paragraph.
                             %
with the right spacing around the periods requires
a special command.

\LaTeX\ interprets some common characters as
commands, so you must type special commands to
generate them.  These characters include the
following:
       \$ \& \% \# \{ and \}.

In printing, text is usually emphasized with an
       \emph{italic}
type style.

\begin{em}
   A long segment of text can also be emphasized
   in this way.  Text within such a segment can be
   given \emph{additional} emphasis.
\end{em}

It is sometimes necessary to prevent \LaTeX\ from
breaking a line where it might otherwise do so.
This may be at a space, as between the ``Mr.''\ and
``Jones'' in
       ``Mr.~Jones'',        % ~ produces an unbreakable interword space.
or within a word---especially when the word is a
symbol like
       \mbox{\emph{itemnum}}
that makes little sense when hyphenated across
lines.

Footnotes\footnote{This is an example of a footnote.}
pose no problem.

\LaTeX\ is good at typesetting mathematical formulas
like
       \( x-3y + z = 7 \)
or
       \( a_{1} > x^{2n} + y^{2n} > x' \)
or
       \( \ip{A}{B} = \sum_{i} a_{i} b_{i} \).
The spaces you type in a formula are
ignored.  Remember that a letter like
       $x$                   % $ ... $  and  \( ... \)  are equivalent
is a formula when it denotes a mathematical
symbol, and it should be typed as one.

\section{Displayed Text}

Text is displayed by indenting it from the left
margin.  Quotations are commonly displayed.  There
are short quotations
\begin{quote}
   This is a short quotation.  It consists of a
   single paragraph of text.  See how it is formatted.
\end{quote}
and longer ones.
\begin{quotation}
   This is a longer quotation.  It consists of two
   paragraphs of text, neither of which are
   particularly interesting.

   This is the second paragraph of the quotation.  It
   is just as dull as the first paragraph.
\end{quotation}
Another frequently-displayed structure is a list.
The following is an example of an \emph{itemized}
list.
\begin{itemize}
   \item This is the first item of an itemized list.
         Each item in the list is marked with a ``tick''.
         You don't have to worry about what kind of tick
         mark is used.

   \item This is the second item of the list.  It
         contains another list nested inside it.  The inner
         list is an \emph{enumerated} list.
         \begin{enumerate}
            \item This is the first item of an enumerated
                  list that is nested within the itemized list.

            \item This is the second item of the inner list.
                  \LaTeX\ allows you to nest lists deeper than
                  you really should.
         \end{enumerate}
         This is the rest of the second item of the outer
         list.  It is no more interesting than any other
         part of the item.
   \item This is the third item of the list.
\end{itemize}
You can even display poetry.
\begin{verse}
   There is an environment
    for verse \\             % The \\ command separates lines
   Whose features some poets % within a stanza.
   will curse.

                             % One or more blank lines separate stanzas.

   For instead of making\\
   Them do \emph{all} line breaking, \\
   It allows them to put too many words on a line when they'd rather be
   forced to be terse.
\end{verse}

Mathematical formulas may also be displayed.  A
displayed formula
is
one-line long; multiline
formulas require special formatting instructions.
   \[  \ip{\Gamma}{\psi'} = x'' + y^{2} + z_{i}^{n}\]
Don't start a paragraph with a displayed equation,
nor make one a paragraph by itself.

\end{document}               % End of document.

\section{Bench Huge Section 1090}
% synthetic bench section 1090
% This is a sample LaTeX input file.  (Version of 12 August 2004.)
%
% A '%' character causes TeX to ignore all remaining text on the line,
% and is used for comments like this one.

\documentclass{article}      % Specifies the document class

                             % The preamble begins here.
\title{An Example Document}  % Declares the document's title.
\author{Leslie Lamport}      % Declares the author's name.
\date{January 21, 1994}      % Deleting this command produces today's date.

\newcommand{\ip}[2]{(#1, #2)}
                             % Defines \ip{arg1}{arg2} to mean
                             % (arg1, arg2).

%\newcommand{\ip}[2]{\langle #1 | #2\rangle}
                             % This is an alternative definition of
                             % \ip that is commented out.

\begin{document}             % End of preamble and beginning of text.

\maketitle                   % Produces the title.

This is an example input file.  Comparing it with
the output it generates can show you how to
produce a simple document of your own.

\section{Ordinary Text}      % Produces section heading.  Lower-level
                             % sections are begun with similar
                             % \subsection and \subsubsection commands.

The ends  of words and sentences are marked
  by   spaces. It  doesn't matter how many
spaces    you type; one is as good as 100.  The
end of   a line counts as a space.

One   or more   blank lines denote the  end
of  a paragraph.

Since any number of consecutive spaces are treated
like a single one, the formatting of the input
file makes no difference to
      \LaTeX,                % The \LaTeX command generates the LaTeX logo.
but it makes a difference to you.  When you use
\LaTeX, making your input file as easy to read
as possible will be a great help as you write
your document and when you change it.  This sample
file shows how you can add comments to your own input
file.

Because printing is different from typewriting,
there are a number of things that you have to do
differently when preparing an input file than if
you were just typing the document directly.
Quotation marks like
       ``this''
have to be handled specially, as do quotes within
quotes:
       ``\,`this'            % \, separates the double and single quote.
        is what I just
        wrote, not  `that'\,''.

Dashes come in three sizes: an
       intra-word
dash, a medium dash for number ranges like
       1--2,
and a punctuation
       dash---like
this.

A sentence-ending space should be larger than the
space between words within a sentence.  You
sometimes have to type special commands in
conjunction with punctuation characters to get
this right, as in the following sentence.
       Gnats, gnus, etc.\ all  % `\ ' makes an inter-word space.
       begin with G\@.         % \@ marks end-of-sentence punctuation.
You should check the spaces after periods when
reading your output to make sure you haven't
forgotten any special cases.  Generating an
ellipsis
       \ldots\               % `\ ' is needed after `\ldots' because TeX
                             % ignores spaces after command names like \ldots
                             % made from \ + letters.
                             %
                             % Note how a `%' character causes TeX to ignore
                             % the end of the input line, so these blank lines
                             % do not start a new paragraph.
                             %
with the right spacing around the periods requires
a special command.

\LaTeX\ interprets some common characters as
commands, so you must type special commands to
generate them.  These characters include the
following:
       \$ \& \% \# \{ and \}.

In printing, text is usually emphasized with an
       \emph{italic}
type style.

\begin{em}
   A long segment of text can also be emphasized
   in this way.  Text within such a segment can be
   given \emph{additional} emphasis.
\end{em}

It is sometimes necessary to prevent \LaTeX\ from
breaking a line where it might otherwise do so.
This may be at a space, as between the ``Mr.''\ and
``Jones'' in
       ``Mr.~Jones'',        % ~ produces an unbreakable interword space.
or within a word---especially when the word is a
symbol like
       \mbox{\emph{itemnum}}
that makes little sense when hyphenated across
lines.

Footnotes\footnote{This is an example of a footnote.}
pose no problem.

\LaTeX\ is good at typesetting mathematical formulas
like
       \( x-3y + z = 7 \)
or
       \( a_{1} > x^{2n} + y^{2n} > x' \)
or
       \( \ip{A}{B} = \sum_{i} a_{i} b_{i} \).
The spaces you type in a formula are
ignored.  Remember that a letter like
       $x$                   % $ ... $  and  \( ... \)  are equivalent
is a formula when it denotes a mathematical
symbol, and it should be typed as one.

\section{Displayed Text}

Text is displayed by indenting it from the left
margin.  Quotations are commonly displayed.  There
are short quotations
\begin{quote}
   This is a short quotation.  It consists of a
   single paragraph of text.  See how it is formatted.
\end{quote}
and longer ones.
\begin{quotation}
   This is a longer quotation.  It consists of two
   paragraphs of text, neither of which are
   particularly interesting.

   This is the second paragraph of the quotation.  It
   is just as dull as the first paragraph.
\end{quotation}
Another frequently-displayed structure is a list.
The following is an example of an \emph{itemized}
list.
\begin{itemize}
   \item This is the first item of an itemized list.
         Each item in the list is marked with a ``tick''.
         You don't have to worry about what kind of tick
         mark is used.

   \item This is the second item of the list.  It
         contains another list nested inside it.  The inner
         list is an \emph{enumerated} list.
         \begin{enumerate}
            \item This is the first item of an enumerated
                  list that is nested within the itemized list.

            \item This is the second item of the inner list.
                  \LaTeX\ allows you to nest lists deeper than
                  you really should.
         \end{enumerate}
         This is the rest of the second item of the outer
         list.  It is no more interesting than any other
         part of the item.
   \item This is the third item of the list.
\end{itemize}
You can even display poetry.
\begin{verse}
   There is an environment
    for verse \\             % The \\ command separates lines
   Whose features some poets % within a stanza.
   will curse.

                             % One or more blank lines separate stanzas.

   For instead of making\\
   Them do \emph{all} line breaking, \\
   It allows them to put too many words on a line when they'd rather be
   forced to be terse.
\end{verse}

Mathematical formulas may also be displayed.  A
displayed formula
is
one-line long; multiline
formulas require special formatting instructions.
   \[  \ip{\Gamma}{\psi'} = x'' + y^{2} + z_{i}^{n}\]
Don't start a paragraph with a displayed equation,
nor make one a paragraph by itself.

\end{document}               % End of document.

\section{Bench Huge Section 1091}
% synthetic bench section 1091
% This is a sample LaTeX input file.  (Version of 12 August 2004.)
%
% A '%' character causes TeX to ignore all remaining text on the line,
% and is used for comments like this one.

\documentclass{article}      % Specifies the document class

                             % The preamble begins here.
\title{An Example Document}  % Declares the document's title.
\author{Leslie Lamport}      % Declares the author's name.
\date{January 21, 1994}      % Deleting this command produces today's date.

\newcommand{\ip}[2]{(#1, #2)}
                             % Defines \ip{arg1}{arg2} to mean
                             % (arg1, arg2).

%\newcommand{\ip}[2]{\langle #1 | #2\rangle}
                             % This is an alternative definition of
                             % \ip that is commented out.

\begin{document}             % End of preamble and beginning of text.

\maketitle                   % Produces the title.

This is an example input file.  Comparing it with
the output it generates can show you how to
produce a simple document of your own.

\section{Ordinary Text}      % Produces section heading.  Lower-level
                             % sections are begun with similar
                             % \subsection and \subsubsection commands.

The ends  of words and sentences are marked
  by   spaces. It  doesn't matter how many
spaces    you type; one is as good as 100.  The
end of   a line counts as a space.

One   or more   blank lines denote the  end
of  a paragraph.

Since any number of consecutive spaces are treated
like a single one, the formatting of the input
file makes no difference to
      \LaTeX,                % The \LaTeX command generates the LaTeX logo.
but it makes a difference to you.  When you use
\LaTeX, making your input file as easy to read
as possible will be a great help as you write
your document and when you change it.  This sample
file shows how you can add comments to your own input
file.

Because printing is different from typewriting,
there are a number of things that you have to do
differently when preparing an input file than if
you were just typing the document directly.
Quotation marks like
       ``this''
have to be handled specially, as do quotes within
quotes:
       ``\,`this'            % \, separates the double and single quote.
        is what I just
        wrote, not  `that'\,''.

Dashes come in three sizes: an
       intra-word
dash, a medium dash for number ranges like
       1--2,
and a punctuation
       dash---like
this.

A sentence-ending space should be larger than the
space between words within a sentence.  You
sometimes have to type special commands in
conjunction with punctuation characters to get
this right, as in the following sentence.
       Gnats, gnus, etc.\ all  % `\ ' makes an inter-word space.
       begin with G\@.         % \@ marks end-of-sentence punctuation.
You should check the spaces after periods when
reading your output to make sure you haven't
forgotten any special cases.  Generating an
ellipsis
       \ldots\               % `\ ' is needed after `\ldots' because TeX
                             % ignores spaces after command names like \ldots
                             % made from \ + letters.
                             %
                             % Note how a `%' character causes TeX to ignore
                             % the end of the input line, so these blank lines
                             % do not start a new paragraph.
                             %
with the right spacing around the periods requires
a special command.

\LaTeX\ interprets some common characters as
commands, so you must type special commands to
generate them.  These characters include the
following:
       \$ \& \% \# \{ and \}.

In printing, text is usually emphasized with an
       \emph{italic}
type style.

\begin{em}
   A long segment of text can also be emphasized
   in this way.  Text within such a segment can be
   given \emph{additional} emphasis.
\end{em}

It is sometimes necessary to prevent \LaTeX\ from
breaking a line where it might otherwise do so.
This may be at a space, as between the ``Mr.''\ and
``Jones'' in
       ``Mr.~Jones'',        % ~ produces an unbreakable interword space.
or within a word---especially when the word is a
symbol like
       \mbox{\emph{itemnum}}
that makes little sense when hyphenated across
lines.

Footnotes\footnote{This is an example of a footnote.}
pose no problem.

\LaTeX\ is good at typesetting mathematical formulas
like
       \( x-3y + z = 7 \)
or
       \( a_{1} > x^{2n} + y^{2n} > x' \)
or
       \( \ip{A}{B} = \sum_{i} a_{i} b_{i} \).
The spaces you type in a formula are
ignored.  Remember that a letter like
       $x$                   % $ ... $  and  \( ... \)  are equivalent
is a formula when it denotes a mathematical
symbol, and it should be typed as one.

\section{Displayed Text}

Text is displayed by indenting it from the left
margin.  Quotations are commonly displayed.  There
are short quotations
\begin{quote}
   This is a short quotation.  It consists of a
   single paragraph of text.  See how it is formatted.
\end{quote}
and longer ones.
\begin{quotation}
   This is a longer quotation.  It consists of two
   paragraphs of text, neither of which are
   particularly interesting.

   This is the second paragraph of the quotation.  It
   is just as dull as the first paragraph.
\end{quotation}
Another frequently-displayed structure is a list.
The following is an example of an \emph{itemized}
list.
\begin{itemize}
   \item This is the first item of an itemized list.
         Each item in the list is marked with a ``tick''.
         You don't have to worry about what kind of tick
         mark is used.

   \item This is the second item of the list.  It
         contains another list nested inside it.  The inner
         list is an \emph{enumerated} list.
         \begin{enumerate}
            \item This is the first item of an enumerated
                  list that is nested within the itemized list.

            \item This is the second item of the inner list.
                  \LaTeX\ allows you to nest lists deeper than
                  you really should.
         \end{enumerate}
         This is the rest of the second item of the outer
         list.  It is no more interesting than any other
         part of the item.
   \item This is the third item of the list.
\end{itemize}
You can even display poetry.
\begin{verse}
   There is an environment
    for verse \\             % The \\ command separates lines
   Whose features some poets % within a stanza.
   will curse.

                             % One or more blank lines separate stanzas.

   For instead of making\\
   Them do \emph{all} line breaking, \\
   It allows them to put too many words on a line when they'd rather be
   forced to be terse.
\end{verse}

Mathematical formulas may also be displayed.  A
displayed formula
is
one-line long; multiline
formulas require special formatting instructions.
   \[  \ip{\Gamma}{\psi'} = x'' + y^{2} + z_{i}^{n}\]
Don't start a paragraph with a displayed equation,
nor make one a paragraph by itself.

\end{document}               % End of document.

\section{Bench Huge Section 1092}
% synthetic bench section 1092
% This is a sample LaTeX input file.  (Version of 12 August 2004.)
%
% A '%' character causes TeX to ignore all remaining text on the line,
% and is used for comments like this one.

\documentclass{article}      % Specifies the document class

                             % The preamble begins here.
\title{An Example Document}  % Declares the document's title.
\author{Leslie Lamport}      % Declares the author's name.
\date{January 21, 1994}      % Deleting this command produces today's date.

\newcommand{\ip}[2]{(#1, #2)}
                             % Defines \ip{arg1}{arg2} to mean
                             % (arg1, arg2).

%\newcommand{\ip}[2]{\langle #1 | #2\rangle}
                             % This is an alternative definition of
                             % \ip that is commented out.

\begin{document}             % End of preamble and beginning of text.

\maketitle                   % Produces the title.

This is an example input file.  Comparing it with
the output it generates can show you how to
produce a simple document of your own.

\section{Ordinary Text}      % Produces section heading.  Lower-level
                             % sections are begun with similar
                             % \subsection and \subsubsection commands.

The ends  of words and sentences are marked
  by   spaces. It  doesn't matter how many
spaces    you type; one is as good as 100.  The
end of   a line counts as a space.

One   or more   blank lines denote the  end
of  a paragraph.

Since any number of consecutive spaces are treated
like a single one, the formatting of the input
file makes no difference to
      \LaTeX,                % The \LaTeX command generates the LaTeX logo.
but it makes a difference to you.  When you use
\LaTeX, making your input file as easy to read
as possible will be a great help as you write
your document and when you change it.  This sample
file shows how you can add comments to your own input
file.

Because printing is different from typewriting,
there are a number of things that you have to do
differently when preparing an input file than if
you were just typing the document directly.
Quotation marks like
       ``this''
have to be handled specially, as do quotes within
quotes:
       ``\,`this'            % \, separates the double and single quote.
        is what I just
        wrote, not  `that'\,''.

Dashes come in three sizes: an
       intra-word
dash, a medium dash for number ranges like
       1--2,
and a punctuation
       dash---like
this.

A sentence-ending space should be larger than the
space between words within a sentence.  You
sometimes have to type special commands in
conjunction with punctuation characters to get
this right, as in the following sentence.
       Gnats, gnus, etc.\ all  % `\ ' makes an inter-word space.
       begin with G\@.         % \@ marks end-of-sentence punctuation.
You should check the spaces after periods when
reading your output to make sure you haven't
forgotten any special cases.  Generating an
ellipsis
       \ldots\               % `\ ' is needed after `\ldots' because TeX
                             % ignores spaces after command names like \ldots
                             % made from \ + letters.
                             %
                             % Note how a `%' character causes TeX to ignore
                             % the end of the input line, so these blank lines
                             % do not start a new paragraph.
                             %
with the right spacing around the periods requires
a special command.

\LaTeX\ interprets some common characters as
commands, so you must type special commands to
generate them.  These characters include the
following:
       \$ \& \% \# \{ and \}.

In printing, text is usually emphasized with an
       \emph{italic}
type style.

\begin{em}
   A long segment of text can also be emphasized
   in this way.  Text within such a segment can be
   given \emph{additional} emphasis.
\end{em}

It is sometimes necessary to prevent \LaTeX\ from
breaking a line where it might otherwise do so.
This may be at a space, as between the ``Mr.''\ and
``Jones'' in
       ``Mr.~Jones'',        % ~ produces an unbreakable interword space.
or within a word---especially when the word is a
symbol like
       \mbox{\emph{itemnum}}
that makes little sense when hyphenated across
lines.

Footnotes\footnote{This is an example of a footnote.}
pose no problem.

\LaTeX\ is good at typesetting mathematical formulas
like
       \( x-3y + z = 7 \)
or
       \( a_{1} > x^{2n} + y^{2n} > x' \)
or
       \( \ip{A}{B} = \sum_{i} a_{i} b_{i} \).
The spaces you type in a formula are
ignored.  Remember that a letter like
       $x$                   % $ ... $  and  \( ... \)  are equivalent
is a formula when it denotes a mathematical
symbol, and it should be typed as one.

\section{Displayed Text}

Text is displayed by indenting it from the left
margin.  Quotations are commonly displayed.  There
are short quotations
\begin{quote}
   This is a short quotation.  It consists of a
   single paragraph of text.  See how it is formatted.
\end{quote}
and longer ones.
\begin{quotation}
   This is a longer quotation.  It consists of two
   paragraphs of text, neither of which are
   particularly interesting.

   This is the second paragraph of the quotation.  It
   is just as dull as the first paragraph.
\end{quotation}
Another frequently-displayed structure is a list.
The following is an example of an \emph{itemized}
list.
\begin{itemize}
   \item This is the first item of an itemized list.
         Each item in the list is marked with a ``tick''.
         You don't have to worry about what kind of tick
         mark is used.

   \item This is the second item of the list.  It
         contains another list nested inside it.  The inner
         list is an \emph{enumerated} list.
         \begin{enumerate}
            \item This is the first item of an enumerated
                  list that is nested within the itemized list.

            \item This is the second item of the inner list.
                  \LaTeX\ allows you to nest lists deeper than
                  you really should.
         \end{enumerate}
         This is the rest of the second item of the outer
         list.  It is no more interesting than any other
         part of the item.
   \item This is the third item of the list.
\end{itemize}
You can even display poetry.
\begin{verse}
   There is an environment
    for verse \\             % The \\ command separates lines
   Whose features some poets % within a stanza.
   will curse.

                             % One or more blank lines separate stanzas.

   For instead of making\\
   Them do \emph{all} line breaking, \\
   It allows them to put too many words on a line when they'd rather be
   forced to be terse.
\end{verse}

Mathematical formulas may also be displayed.  A
displayed formula
is
one-line long; multiline
formulas require special formatting instructions.
   \[  \ip{\Gamma}{\psi'} = x'' + y^{2} + z_{i}^{n}\]
Don't start a paragraph with a displayed equation,
nor make one a paragraph by itself.

\end{document}               % End of document.

\section{Bench Huge Section 1093}
% synthetic bench section 1093
% This is a sample LaTeX input file.  (Version of 12 August 2004.)
%
% A '%' character causes TeX to ignore all remaining text on the line,
% and is used for comments like this one.

\documentclass{article}      % Specifies the document class

                             % The preamble begins here.
\title{An Example Document}  % Declares the document's title.
\author{Leslie Lamport}      % Declares the author's name.
\date{January 21, 1994}      % Deleting this command produces today's date.

\newcommand{\ip}[2]{(#1, #2)}
                             % Defines \ip{arg1}{arg2} to mean
                             % (arg1, arg2).

%\newcommand{\ip}[2]{\langle #1 | #2\rangle}
                             % This is an alternative definition of
                             % \ip that is commented out.

\begin{document}             % End of preamble and beginning of text.

\maketitle                   % Produces the title.

This is an example input file.  Comparing it with
the output it generates can show you how to
produce a simple document of your own.

\section{Ordinary Text}      % Produces section heading.  Lower-level
                             % sections are begun with similar
                             % \subsection and \subsubsection commands.

The ends  of words and sentences are marked
  by   spaces. It  doesn't matter how many
spaces    you type; one is as good as 100.  The
end of   a line counts as a space.

One   or more   blank lines denote the  end
of  a paragraph.

Since any number of consecutive spaces are treated
like a single one, the formatting of the input
file makes no difference to
      \LaTeX,                % The \LaTeX command generates the LaTeX logo.
but it makes a difference to you.  When you use
\LaTeX, making your input file as easy to read
as possible will be a great help as you write
your document and when you change it.  This sample
file shows how you can add comments to your own input
file.

Because printing is different from typewriting,
there are a number of things that you have to do
differently when preparing an input file than if
you were just typing the document directly.
Quotation marks like
       ``this''
have to be handled specially, as do quotes within
quotes:
       ``\,`this'            % \, separates the double and single quote.
        is what I just
        wrote, not  `that'\,''.

Dashes come in three sizes: an
       intra-word
dash, a medium dash for number ranges like
       1--2,
and a punctuation
       dash---like
this.

A sentence-ending space should be larger than the
space between words within a sentence.  You
sometimes have to type special commands in
conjunction with punctuation characters to get
this right, as in the following sentence.
       Gnats, gnus, etc.\ all  % `\ ' makes an inter-word space.
       begin with G\@.         % \@ marks end-of-sentence punctuation.
You should check the spaces after periods when
reading your output to make sure you haven't
forgotten any special cases.  Generating an
ellipsis
       \ldots\               % `\ ' is needed after `\ldots' because TeX
                             % ignores spaces after command names like \ldots
                             % made from \ + letters.
                             %
                             % Note how a `%' character causes TeX to ignore
                             % the end of the input line, so these blank lines
                             % do not start a new paragraph.
                             %
with the right spacing around the periods requires
a special command.

\LaTeX\ interprets some common characters as
commands, so you must type special commands to
generate them.  These characters include the
following:
       \$ \& \% \# \{ and \}.

In printing, text is usually emphasized with an
       \emph{italic}
type style.

\begin{em}
   A long segment of text can also be emphasized
   in this way.  Text within such a segment can be
   given \emph{additional} emphasis.
\end{em}

It is sometimes necessary to prevent \LaTeX\ from
breaking a line where it might otherwise do so.
This may be at a space, as between the ``Mr.''\ and
``Jones'' in
       ``Mr.~Jones'',        % ~ produces an unbreakable interword space.
or within a word---especially when the word is a
symbol like
       \mbox{\emph{itemnum}}
that makes little sense when hyphenated across
lines.

Footnotes\footnote{This is an example of a footnote.}
pose no problem.

\LaTeX\ is good at typesetting mathematical formulas
like
       \( x-3y + z = 7 \)
or
       \( a_{1} > x^{2n} + y^{2n} > x' \)
or
       \( \ip{A}{B} = \sum_{i} a_{i} b_{i} \).
The spaces you type in a formula are
ignored.  Remember that a letter like
       $x$                   % $ ... $  and  \( ... \)  are equivalent
is a formula when it denotes a mathematical
symbol, and it should be typed as one.

\section{Displayed Text}

Text is displayed by indenting it from the left
margin.  Quotations are commonly displayed.  There
are short quotations
\begin{quote}
   This is a short quotation.  It consists of a
   single paragraph of text.  See how it is formatted.
\end{quote}
and longer ones.
\begin{quotation}
   This is a longer quotation.  It consists of two
   paragraphs of text, neither of which are
   particularly interesting.

   This is the second paragraph of the quotation.  It
   is just as dull as the first paragraph.
\end{quotation}
Another frequently-displayed structure is a list.
The following is an example of an \emph{itemized}
list.
\begin{itemize}
   \item This is the first item of an itemized list.
         Each item in the list is marked with a ``tick''.
         You don't have to worry about what kind of tick
         mark is used.

   \item This is the second item of the list.  It
         contains another list nested inside it.  The inner
         list is an \emph{enumerated} list.
         \begin{enumerate}
            \item This is the first item of an enumerated
                  list that is nested within the itemized list.

            \item This is the second item of the inner list.
                  \LaTeX\ allows you to nest lists deeper than
                  you really should.
         \end{enumerate}
         This is the rest of the second item of the outer
         list.  It is no more interesting than any other
         part of the item.
   \item This is the third item of the list.
\end{itemize}
You can even display poetry.
\begin{verse}
   There is an environment
    for verse \\             % The \\ command separates lines
   Whose features some poets % within a stanza.
   will curse.

                             % One or more blank lines separate stanzas.

   For instead of making\\
   Them do \emph{all} line breaking, \\
   It allows them to put too many words on a line when they'd rather be
   forced to be terse.
\end{verse}

Mathematical formulas may also be displayed.  A
displayed formula
is
one-line long; multiline
formulas require special formatting instructions.
   \[  \ip{\Gamma}{\psi'} = x'' + y^{2} + z_{i}^{n}\]
Don't start a paragraph with a displayed equation,
nor make one a paragraph by itself.

\end{document}               % End of document.

\section{Bench Huge Section 1094}
% synthetic bench section 1094
% This is a sample LaTeX input file.  (Version of 12 August 2004.)
%
% A '%' character causes TeX to ignore all remaining text on the line,
% and is used for comments like this one.

\documentclass{article}      % Specifies the document class

                             % The preamble begins here.
\title{An Example Document}  % Declares the document's title.
\author{Leslie Lamport}      % Declares the author's name.
\date{January 21, 1994}      % Deleting this command produces today's date.

\newcommand{\ip}[2]{(#1, #2)}
                             % Defines \ip{arg1}{arg2} to mean
                             % (arg1, arg2).

%\newcommand{\ip}[2]{\langle #1 | #2\rangle}
                             % This is an alternative definition of
                             % \ip that is commented out.

\begin{document}             % End of preamble and beginning of text.

\maketitle                   % Produces the title.

This is an example input file.  Comparing it with
the output it generates can show you how to
produce a simple document of your own.

\section{Ordinary Text}      % Produces section heading.  Lower-level
                             % sections are begun with similar
                             % \subsection and \subsubsection commands.

The ends  of words and sentences are marked
  by   spaces. It  doesn't matter how many
spaces    you type; one is as good as 100.  The
end of   a line counts as a space.

One   or more   blank lines denote the  end
of  a paragraph.

Since any number of consecutive spaces are treated
like a single one, the formatting of the input
file makes no difference to
      \LaTeX,                % The \LaTeX command generates the LaTeX logo.
but it makes a difference to you.  When you use
\LaTeX, making your input file as easy to read
as possible will be a great help as you write
your document and when you change it.  This sample
file shows how you can add comments to your own input
file.

Because printing is different from typewriting,
there are a number of things that you have to do
differently when preparing an input file than if
you were just typing the document directly.
Quotation marks like
       ``this''
have to be handled specially, as do quotes within
quotes:
       ``\,`this'            % \, separates the double and single quote.
        is what I just
        wrote, not  `that'\,''.

Dashes come in three sizes: an
       intra-word
dash, a medium dash for number ranges like
       1--2,
and a punctuation
       dash---like
this.

A sentence-ending space should be larger than the
space between words within a sentence.  You
sometimes have to type special commands in
conjunction with punctuation characters to get
this right, as in the following sentence.
       Gnats, gnus, etc.\ all  % `\ ' makes an inter-word space.
       begin with G\@.         % \@ marks end-of-sentence punctuation.
You should check the spaces after periods when
reading your output to make sure you haven't
forgotten any special cases.  Generating an
ellipsis
       \ldots\               % `\ ' is needed after `\ldots' because TeX
                             % ignores spaces after command names like \ldots
                             % made from \ + letters.
                             %
                             % Note how a `%' character causes TeX to ignore
                             % the end of the input line, so these blank lines
                             % do not start a new paragraph.
                             %
with the right spacing around the periods requires
a special command.

\LaTeX\ interprets some common characters as
commands, so you must type special commands to
generate them.  These characters include the
following:
       \$ \& \% \# \{ and \}.

In printing, text is usually emphasized with an
       \emph{italic}
type style.

\begin{em}
   A long segment of text can also be emphasized
   in this way.  Text within such a segment can be
   given \emph{additional} emphasis.
\end{em}

It is sometimes necessary to prevent \LaTeX\ from
breaking a line where it might otherwise do so.
This may be at a space, as between the ``Mr.''\ and
``Jones'' in
       ``Mr.~Jones'',        % ~ produces an unbreakable interword space.
or within a word---especially when the word is a
symbol like
       \mbox{\emph{itemnum}}
that makes little sense when hyphenated across
lines.

Footnotes\footnote{This is an example of a footnote.}
pose no problem.

\LaTeX\ is good at typesetting mathematical formulas
like
       \( x-3y + z = 7 \)
or
       \( a_{1} > x^{2n} + y^{2n} > x' \)
or
       \( \ip{A}{B} = \sum_{i} a_{i} b_{i} \).
The spaces you type in a formula are
ignored.  Remember that a letter like
       $x$                   % $ ... $  and  \( ... \)  are equivalent
is a formula when it denotes a mathematical
symbol, and it should be typed as one.

\section{Displayed Text}

Text is displayed by indenting it from the left
margin.  Quotations are commonly displayed.  There
are short quotations
\begin{quote}
   This is a short quotation.  It consists of a
   single paragraph of text.  See how it is formatted.
\end{quote}
and longer ones.
\begin{quotation}
   This is a longer quotation.  It consists of two
   paragraphs of text, neither of which are
   particularly interesting.

   This is the second paragraph of the quotation.  It
   is just as dull as the first paragraph.
\end{quotation}
Another frequently-displayed structure is a list.
The following is an example of an \emph{itemized}
list.
\begin{itemize}
   \item This is the first item of an itemized list.
         Each item in the list is marked with a ``tick''.
         You don't have to worry about what kind of tick
         mark is used.

   \item This is the second item of the list.  It
         contains another list nested inside it.  The inner
         list is an \emph{enumerated} list.
         \begin{enumerate}
            \item This is the first item of an enumerated
                  list that is nested within the itemized list.

            \item This is the second item of the inner list.
                  \LaTeX\ allows you to nest lists deeper than
                  you really should.
         \end{enumerate}
         This is the rest of the second item of the outer
         list.  It is no more interesting than any other
         part of the item.
   \item This is the third item of the list.
\end{itemize}
You can even display poetry.
\begin{verse}
   There is an environment
    for verse \\             % The \\ command separates lines
   Whose features some poets % within a stanza.
   will curse.

                             % One or more blank lines separate stanzas.

   For instead of making\\
   Them do \emph{all} line breaking, \\
   It allows them to put too many words on a line when they'd rather be
   forced to be terse.
\end{verse}

Mathematical formulas may also be displayed.  A
displayed formula
is
one-line long; multiline
formulas require special formatting instructions.
   \[  \ip{\Gamma}{\psi'} = x'' + y^{2} + z_{i}^{n}\]
Don't start a paragraph with a displayed equation,
nor make one a paragraph by itself.

\end{document}               % End of document.

\section{Bench Huge Section 1095}
% synthetic bench section 1095
% This is a sample LaTeX input file.  (Version of 12 August 2004.)
%
% A '%' character causes TeX to ignore all remaining text on the line,
% and is used for comments like this one.

\documentclass{article}      % Specifies the document class

                             % The preamble begins here.
\title{An Example Document}  % Declares the document's title.
\author{Leslie Lamport}      % Declares the author's name.
\date{January 21, 1994}      % Deleting this command produces today's date.

\newcommand{\ip}[2]{(#1, #2)}
                             % Defines \ip{arg1}{arg2} to mean
                             % (arg1, arg2).

%\newcommand{\ip}[2]{\langle #1 | #2\rangle}
                             % This is an alternative definition of
                             % \ip that is commented out.

\begin{document}             % End of preamble and beginning of text.

\maketitle                   % Produces the title.

This is an example input file.  Comparing it with
the output it generates can show you how to
produce a simple document of your own.

\section{Ordinary Text}      % Produces section heading.  Lower-level
                             % sections are begun with similar
                             % \subsection and \subsubsection commands.

The ends  of words and sentences are marked
  by   spaces. It  doesn't matter how many
spaces    you type; one is as good as 100.  The
end of   a line counts as a space.

One   or more   blank lines denote the  end
of  a paragraph.

Since any number of consecutive spaces are treated
like a single one, the formatting of the input
file makes no difference to
      \LaTeX,                % The \LaTeX command generates the LaTeX logo.
but it makes a difference to you.  When you use
\LaTeX, making your input file as easy to read
as possible will be a great help as you write
your document and when you change it.  This sample
file shows how you can add comments to your own input
file.

Because printing is different from typewriting,
there are a number of things that you have to do
differently when preparing an input file than if
you were just typing the document directly.
Quotation marks like
       ``this''
have to be handled specially, as do quotes within
quotes:
       ``\,`this'            % \, separates the double and single quote.
        is what I just
        wrote, not  `that'\,''.

Dashes come in three sizes: an
       intra-word
dash, a medium dash for number ranges like
       1--2,
and a punctuation
       dash---like
this.

A sentence-ending space should be larger than the
space between words within a sentence.  You
sometimes have to type special commands in
conjunction with punctuation characters to get
this right, as in the following sentence.
       Gnats, gnus, etc.\ all  % `\ ' makes an inter-word space.
       begin with G\@.         % \@ marks end-of-sentence punctuation.
You should check the spaces after periods when
reading your output to make sure you haven't
forgotten any special cases.  Generating an
ellipsis
       \ldots\               % `\ ' is needed after `\ldots' because TeX
                             % ignores spaces after command names like \ldots
                             % made from \ + letters.
                             %
                             % Note how a `%' character causes TeX to ignore
                             % the end of the input line, so these blank lines
                             % do not start a new paragraph.
                             %
with the right spacing around the periods requires
a special command.

\LaTeX\ interprets some common characters as
commands, so you must type special commands to
generate them.  These characters include the
following:
       \$ \& \% \# \{ and \}.

In printing, text is usually emphasized with an
       \emph{italic}
type style.

\begin{em}
   A long segment of text can also be emphasized
   in this way.  Text within such a segment can be
   given \emph{additional} emphasis.
\end{em}

It is sometimes necessary to prevent \LaTeX\ from
breaking a line where it might otherwise do so.
This may be at a space, as between the ``Mr.''\ and
``Jones'' in
       ``Mr.~Jones'',        % ~ produces an unbreakable interword space.
or within a word---especially when the word is a
symbol like
       \mbox{\emph{itemnum}}
that makes little sense when hyphenated across
lines.

Footnotes\footnote{This is an example of a footnote.}
pose no problem.

\LaTeX\ is good at typesetting mathematical formulas
like
       \( x-3y + z = 7 \)
or
       \( a_{1} > x^{2n} + y^{2n} > x' \)
or
       \( \ip{A}{B} = \sum_{i} a_{i} b_{i} \).
The spaces you type in a formula are
ignored.  Remember that a letter like
       $x$                   % $ ... $  and  \( ... \)  are equivalent
is a formula when it denotes a mathematical
symbol, and it should be typed as one.

\section{Displayed Text}

Text is displayed by indenting it from the left
margin.  Quotations are commonly displayed.  There
are short quotations
\begin{quote}
   This is a short quotation.  It consists of a
   single paragraph of text.  See how it is formatted.
\end{quote}
and longer ones.
\begin{quotation}
   This is a longer quotation.  It consists of two
   paragraphs of text, neither of which are
   particularly interesting.

   This is the second paragraph of the quotation.  It
   is just as dull as the first paragraph.
\end{quotation}
Another frequently-displayed structure is a list.
The following is an example of an \emph{itemized}
list.
\begin{itemize}
   \item This is the first item of an itemized list.
         Each item in the list is marked with a ``tick''.
         You don't have to worry about what kind of tick
         mark is used.

   \item This is the second item of the list.  It
         contains another list nested inside it.  The inner
         list is an \emph{enumerated} list.
         \begin{enumerate}
            \item This is the first item of an enumerated
                  list that is nested within the itemized list.

            \item This is the second item of the inner list.
                  \LaTeX\ allows you to nest lists deeper than
                  you really should.
         \end{enumerate}
         This is the rest of the second item of the outer
         list.  It is no more interesting than any other
         part of the item.
   \item This is the third item of the list.
\end{itemize}
You can even display poetry.
\begin{verse}
   There is an environment
    for verse \\             % The \\ command separates lines
   Whose features some poets % within a stanza.
   will curse.

                             % One or more blank lines separate stanzas.

   For instead of making\\
   Them do \emph{all} line breaking, \\
   It allows them to put too many words on a line when they'd rather be
   forced to be terse.
\end{verse}

Mathematical formulas may also be displayed.  A
displayed formula
is
one-line long; multiline
formulas require special formatting instructions.
   \[  \ip{\Gamma}{\psi'} = x'' + y^{2} + z_{i}^{n}\]
Don't start a paragraph with a displayed equation,
nor make one a paragraph by itself.

\end{document}               % End of document.

\section{Bench Huge Section 1096}
% synthetic bench section 1096
% This is a sample LaTeX input file.  (Version of 12 August 2004.)
%
% A '%' character causes TeX to ignore all remaining text on the line,
% and is used for comments like this one.

\documentclass{article}      % Specifies the document class

                             % The preamble begins here.
\title{An Example Document}  % Declares the document's title.
\author{Leslie Lamport}      % Declares the author's name.
\date{January 21, 1994}      % Deleting this command produces today's date.

\newcommand{\ip}[2]{(#1, #2)}
                             % Defines \ip{arg1}{arg2} to mean
                             % (arg1, arg2).

%\newcommand{\ip}[2]{\langle #1 | #2\rangle}
                             % This is an alternative definition of
                             % \ip that is commented out.

\begin{document}             % End of preamble and beginning of text.

\maketitle                   % Produces the title.

This is an example input file.  Comparing it with
the output it generates can show you how to
produce a simple document of your own.

\section{Ordinary Text}      % Produces section heading.  Lower-level
                             % sections are begun with similar
                             % \subsection and \subsubsection commands.

The ends  of words and sentences are marked
  by   spaces. It  doesn't matter how many
spaces    you type; one is as good as 100.  The
end of   a line counts as a space.

One   or more   blank lines denote the  end
of  a paragraph.

Since any number of consecutive spaces are treated
like a single one, the formatting of the input
file makes no difference to
      \LaTeX,                % The \LaTeX command generates the LaTeX logo.
but it makes a difference to you.  When you use
\LaTeX, making your input file as easy to read
as possible will be a great help as you write
your document and when you change it.  This sample
file shows how you can add comments to your own input
file.

Because printing is different from typewriting,
there are a number of things that you have to do
differently when preparing an input file than if
you were just typing the document directly.
Quotation marks like
       ``this''
have to be handled specially, as do quotes within
quotes:
       ``\,`this'            % \, separates the double and single quote.
        is what I just
        wrote, not  `that'\,''.

Dashes come in three sizes: an
       intra-word
dash, a medium dash for number ranges like
       1--2,
and a punctuation
       dash---like
this.

A sentence-ending space should be larger than the
space between words within a sentence.  You
sometimes have to type special commands in
conjunction with punctuation characters to get
this right, as in the following sentence.
       Gnats, gnus, etc.\ all  % `\ ' makes an inter-word space.
       begin with G\@.         % \@ marks end-of-sentence punctuation.
You should check the spaces after periods when
reading your output to make sure you haven't
forgotten any special cases.  Generating an
ellipsis
       \ldots\               % `\ ' is needed after `\ldots' because TeX
                             % ignores spaces after command names like \ldots
                             % made from \ + letters.
                             %
                             % Note how a `%' character causes TeX to ignore
                             % the end of the input line, so these blank lines
                             % do not start a new paragraph.
                             %
with the right spacing around the periods requires
a special command.

\LaTeX\ interprets some common characters as
commands, so you must type special commands to
generate them.  These characters include the
following:
       \$ \& \% \# \{ and \}.

In printing, text is usually emphasized with an
       \emph{italic}
type style.

\begin{em}
   A long segment of text can also be emphasized
   in this way.  Text within such a segment can be
   given \emph{additional} emphasis.
\end{em}

It is sometimes necessary to prevent \LaTeX\ from
breaking a line where it might otherwise do so.
This may be at a space, as between the ``Mr.''\ and
``Jones'' in
       ``Mr.~Jones'',        % ~ produces an unbreakable interword space.
or within a word---especially when the word is a
symbol like
       \mbox{\emph{itemnum}}
that makes little sense when hyphenated across
lines.

Footnotes\footnote{This is an example of a footnote.}
pose no problem.

\LaTeX\ is good at typesetting mathematical formulas
like
       \( x-3y + z = 7 \)
or
       \( a_{1} > x^{2n} + y^{2n} > x' \)
or
       \( \ip{A}{B} = \sum_{i} a_{i} b_{i} \).
The spaces you type in a formula are
ignored.  Remember that a letter like
       $x$                   % $ ... $  and  \( ... \)  are equivalent
is a formula when it denotes a mathematical
symbol, and it should be typed as one.

\section{Displayed Text}

Text is displayed by indenting it from the left
margin.  Quotations are commonly displayed.  There
are short quotations
\begin{quote}
   This is a short quotation.  It consists of a
   single paragraph of text.  See how it is formatted.
\end{quote}
and longer ones.
\begin{quotation}
   This is a longer quotation.  It consists of two
   paragraphs of text, neither of which are
   particularly interesting.

   This is the second paragraph of the quotation.  It
   is just as dull as the first paragraph.
\end{quotation}
Another frequently-displayed structure is a list.
The following is an example of an \emph{itemized}
list.
\begin{itemize}
   \item This is the first item of an itemized list.
         Each item in the list is marked with a ``tick''.
         You don't have to worry about what kind of tick
         mark is used.

   \item This is the second item of the list.  It
         contains another list nested inside it.  The inner
         list is an \emph{enumerated} list.
         \begin{enumerate}
            \item This is the first item of an enumerated
                  list that is nested within the itemized list.

            \item This is the second item of the inner list.
                  \LaTeX\ allows you to nest lists deeper than
                  you really should.
         \end{enumerate}
         This is the rest of the second item of the outer
         list.  It is no more interesting than any other
         part of the item.
   \item This is the third item of the list.
\end{itemize}
You can even display poetry.
\begin{verse}
   There is an environment
    for verse \\             % The \\ command separates lines
   Whose features some poets % within a stanza.
   will curse.

                             % One or more blank lines separate stanzas.

   For instead of making\\
   Them do \emph{all} line breaking, \\
   It allows them to put too many words on a line when they'd rather be
   forced to be terse.
\end{verse}

Mathematical formulas may also be displayed.  A
displayed formula
is
one-line long; multiline
formulas require special formatting instructions.
   \[  \ip{\Gamma}{\psi'} = x'' + y^{2} + z_{i}^{n}\]
Don't start a paragraph with a displayed equation,
nor make one a paragraph by itself.

\end{document}               % End of document.

\section{Bench Huge Section 1097}
% synthetic bench section 1097
% This is a sample LaTeX input file.  (Version of 12 August 2004.)
%
% A '%' character causes TeX to ignore all remaining text on the line,
% and is used for comments like this one.

\documentclass{article}      % Specifies the document class

                             % The preamble begins here.
\title{An Example Document}  % Declares the document's title.
\author{Leslie Lamport}      % Declares the author's name.
\date{January 21, 1994}      % Deleting this command produces today's date.

\newcommand{\ip}[2]{(#1, #2)}
                             % Defines \ip{arg1}{arg2} to mean
                             % (arg1, arg2).

%\newcommand{\ip}[2]{\langle #1 | #2\rangle}
                             % This is an alternative definition of
                             % \ip that is commented out.

\begin{document}             % End of preamble and beginning of text.

\maketitle                   % Produces the title.

This is an example input file.  Comparing it with
the output it generates can show you how to
produce a simple document of your own.

\section{Ordinary Text}      % Produces section heading.  Lower-level
                             % sections are begun with similar
                             % \subsection and \subsubsection commands.

The ends  of words and sentences are marked
  by   spaces. It  doesn't matter how many
spaces    you type; one is as good as 100.  The
end of   a line counts as a space.

One   or more   blank lines denote the  end
of  a paragraph.

Since any number of consecutive spaces are treated
like a single one, the formatting of the input
file makes no difference to
      \LaTeX,                % The \LaTeX command generates the LaTeX logo.
but it makes a difference to you.  When you use
\LaTeX, making your input file as easy to read
as possible will be a great help as you write
your document and when you change it.  This sample
file shows how you can add comments to your own input
file.

Because printing is different from typewriting,
there are a number of things that you have to do
differently when preparing an input file than if
you were just typing the document directly.
Quotation marks like
       ``this''
have to be handled specially, as do quotes within
quotes:
       ``\,`this'            % \, separates the double and single quote.
        is what I just
        wrote, not  `that'\,''.

Dashes come in three sizes: an
       intra-word
dash, a medium dash for number ranges like
       1--2,
and a punctuation
       dash---like
this.

A sentence-ending space should be larger than the
space between words within a sentence.  You
sometimes have to type special commands in
conjunction with punctuation characters to get
this right, as in the following sentence.
       Gnats, gnus, etc.\ all  % `\ ' makes an inter-word space.
       begin with G\@.         % \@ marks end-of-sentence punctuation.
You should check the spaces after periods when
reading your output to make sure you haven't
forgotten any special cases.  Generating an
ellipsis
       \ldots\               % `\ ' is needed after `\ldots' because TeX
                             % ignores spaces after command names like \ldots
                             % made from \ + letters.
                             %
                             % Note how a `%' character causes TeX to ignore
                             % the end of the input line, so these blank lines
                             % do not start a new paragraph.
                             %
with the right spacing around the periods requires
a special command.

\LaTeX\ interprets some common characters as
commands, so you must type special commands to
generate them.  These characters include the
following:
       \$ \& \% \# \{ and \}.

In printing, text is usually emphasized with an
       \emph{italic}
type style.

\begin{em}
   A long segment of text can also be emphasized
   in this way.  Text within such a segment can be
   given \emph{additional} emphasis.
\end{em}

It is sometimes necessary to prevent \LaTeX\ from
breaking a line where it might otherwise do so.
This may be at a space, as between the ``Mr.''\ and
``Jones'' in
       ``Mr.~Jones'',        % ~ produces an unbreakable interword space.
or within a word---especially when the word is a
symbol like
       \mbox{\emph{itemnum}}
that makes little sense when hyphenated across
lines.

Footnotes\footnote{This is an example of a footnote.}
pose no problem.

\LaTeX\ is good at typesetting mathematical formulas
like
       \( x-3y + z = 7 \)
or
       \( a_{1} > x^{2n} + y^{2n} > x' \)
or
       \( \ip{A}{B} = \sum_{i} a_{i} b_{i} \).
The spaces you type in a formula are
ignored.  Remember that a letter like
       $x$                   % $ ... $  and  \( ... \)  are equivalent
is a formula when it denotes a mathematical
symbol, and it should be typed as one.

\section{Displayed Text}

Text is displayed by indenting it from the left
margin.  Quotations are commonly displayed.  There
are short quotations
\begin{quote}
   This is a short quotation.  It consists of a
   single paragraph of text.  See how it is formatted.
\end{quote}
and longer ones.
\begin{quotation}
   This is a longer quotation.  It consists of two
   paragraphs of text, neither of which are
   particularly interesting.

   This is the second paragraph of the quotation.  It
   is just as dull as the first paragraph.
\end{quotation}
Another frequently-displayed structure is a list.
The following is an example of an \emph{itemized}
list.
\begin{itemize}
   \item This is the first item of an itemized list.
         Each item in the list is marked with a ``tick''.
         You don't have to worry about what kind of tick
         mark is used.

   \item This is the second item of the list.  It
         contains another list nested inside it.  The inner
         list is an \emph{enumerated} list.
         \begin{enumerate}
            \item This is the first item of an enumerated
                  list that is nested within the itemized list.

            \item This is the second item of the inner list.
                  \LaTeX\ allows you to nest lists deeper than
                  you really should.
         \end{enumerate}
         This is the rest of the second item of the outer
         list.  It is no more interesting than any other
         part of the item.
   \item This is the third item of the list.
\end{itemize}
You can even display poetry.
\begin{verse}
   There is an environment
    for verse \\             % The \\ command separates lines
   Whose features some poets % within a stanza.
   will curse.

                             % One or more blank lines separate stanzas.

   For instead of making\\
   Them do \emph{all} line breaking, \\
   It allows them to put too many words on a line when they'd rather be
   forced to be terse.
\end{verse}

Mathematical formulas may also be displayed.  A
displayed formula
is
one-line long; multiline
formulas require special formatting instructions.
   \[  \ip{\Gamma}{\psi'} = x'' + y^{2} + z_{i}^{n}\]
Don't start a paragraph with a displayed equation,
nor make one a paragraph by itself.

\end{document}               % End of document.

\section{Bench Huge Section 1098}
% synthetic bench section 1098
% This is a sample LaTeX input file.  (Version of 12 August 2004.)
%
% A '%' character causes TeX to ignore all remaining text on the line,
% and is used for comments like this one.

\documentclass{article}      % Specifies the document class

                             % The preamble begins here.
\title{An Example Document}  % Declares the document's title.
\author{Leslie Lamport}      % Declares the author's name.
\date{January 21, 1994}      % Deleting this command produces today's date.

\newcommand{\ip}[2]{(#1, #2)}
                             % Defines \ip{arg1}{arg2} to mean
                             % (arg1, arg2).

%\newcommand{\ip}[2]{\langle #1 | #2\rangle}
                             % This is an alternative definition of
                             % \ip that is commented out.

\begin{document}             % End of preamble and beginning of text.

\maketitle                   % Produces the title.

This is an example input file.  Comparing it with
the output it generates can show you how to
produce a simple document of your own.

\section{Ordinary Text}      % Produces section heading.  Lower-level
                             % sections are begun with similar
                             % \subsection and \subsubsection commands.

The ends  of words and sentences are marked
  by   spaces. It  doesn't matter how many
spaces    you type; one is as good as 100.  The
end of   a line counts as a space.

One   or more   blank lines denote the  end
of  a paragraph.

Since any number of consecutive spaces are treated
like a single one, the formatting of the input
file makes no difference to
      \LaTeX,                % The \LaTeX command generates the LaTeX logo.
but it makes a difference to you.  When you use
\LaTeX, making your input file as easy to read
as possible will be a great help as you write
your document and when you change it.  This sample
file shows how you can add comments to your own input
file.

Because printing is different from typewriting,
there are a number of things that you have to do
differently when preparing an input file than if
you were just typing the document directly.
Quotation marks like
       ``this''
have to be handled specially, as do quotes within
quotes:
       ``\,`this'            % \, separates the double and single quote.
        is what I just
        wrote, not  `that'\,''.

Dashes come in three sizes: an
       intra-word
dash, a medium dash for number ranges like
       1--2,
and a punctuation
       dash---like
this.

A sentence-ending space should be larger than the
space between words within a sentence.  You
sometimes have to type special commands in
conjunction with punctuation characters to get
this right, as in the following sentence.
       Gnats, gnus, etc.\ all  % `\ ' makes an inter-word space.
       begin with G\@.         % \@ marks end-of-sentence punctuation.
You should check the spaces after periods when
reading your output to make sure you haven't
forgotten any special cases.  Generating an
ellipsis
       \ldots\               % `\ ' is needed after `\ldots' because TeX
                             % ignores spaces after command names like \ldots
                             % made from \ + letters.
                             %
                             % Note how a `%' character causes TeX to ignore
                             % the end of the input line, so these blank lines
                             % do not start a new paragraph.
                             %
with the right spacing around the periods requires
a special command.

\LaTeX\ interprets some common characters as
commands, so you must type special commands to
generate them.  These characters include the
following:
       \$ \& \% \# \{ and \}.

In printing, text is usually emphasized with an
       \emph{italic}
type style.

\begin{em}
   A long segment of text can also be emphasized
   in this way.  Text within such a segment can be
   given \emph{additional} emphasis.
\end{em}

It is sometimes necessary to prevent \LaTeX\ from
breaking a line where it might otherwise do so.
This may be at a space, as between the ``Mr.''\ and
``Jones'' in
       ``Mr.~Jones'',        % ~ produces an unbreakable interword space.
or within a word---especially when the word is a
symbol like
       \mbox{\emph{itemnum}}
that makes little sense when hyphenated across
lines.

Footnotes\footnote{This is an example of a footnote.}
pose no problem.

\LaTeX\ is good at typesetting mathematical formulas
like
       \( x-3y + z = 7 \)
or
       \( a_{1} > x^{2n} + y^{2n} > x' \)
or
       \( \ip{A}{B} = \sum_{i} a_{i} b_{i} \).
The spaces you type in a formula are
ignored.  Remember that a letter like
       $x$                   % $ ... $  and  \( ... \)  are equivalent
is a formula when it denotes a mathematical
symbol, and it should be typed as one.

\section{Displayed Text}

Text is displayed by indenting it from the left
margin.  Quotations are commonly displayed.  There
are short quotations
\begin{quote}
   This is a short quotation.  It consists of a
   single paragraph of text.  See how it is formatted.
\end{quote}
and longer ones.
\begin{quotation}
   This is a longer quotation.  It consists of two
   paragraphs of text, neither of which are
   particularly interesting.

   This is the second paragraph of the quotation.  It
   is just as dull as the first paragraph.
\end{quotation}
Another frequently-displayed structure is a list.
The following is an example of an \emph{itemized}
list.
\begin{itemize}
   \item This is the first item of an itemized list.
         Each item in the list is marked with a ``tick''.
         You don't have to worry about what kind of tick
         mark is used.

   \item This is the second item of the list.  It
         contains another list nested inside it.  The inner
         list is an \emph{enumerated} list.
         \begin{enumerate}
            \item This is the first item of an enumerated
                  list that is nested within the itemized list.

            \item This is the second item of the inner list.
                  \LaTeX\ allows you to nest lists deeper than
                  you really should.
         \end{enumerate}
         This is the rest of the second item of the outer
         list.  It is no more interesting than any other
         part of the item.
   \item This is the third item of the list.
\end{itemize}
You can even display poetry.
\begin{verse}
   There is an environment
    for verse \\             % The \\ command separates lines
   Whose features some poets % within a stanza.
   will curse.

                             % One or more blank lines separate stanzas.

   For instead of making\\
   Them do \emph{all} line breaking, \\
   It allows them to put too many words on a line when they'd rather be
   forced to be terse.
\end{verse}

Mathematical formulas may also be displayed.  A
displayed formula
is
one-line long; multiline
formulas require special formatting instructions.
   \[  \ip{\Gamma}{\psi'} = x'' + y^{2} + z_{i}^{n}\]
Don't start a paragraph with a displayed equation,
nor make one a paragraph by itself.

\end{document}               % End of document.

\section{Bench Huge Section 1099}
% synthetic bench section 1099
% This is a sample LaTeX input file.  (Version of 12 August 2004.)
%
% A '%' character causes TeX to ignore all remaining text on the line,
% and is used for comments like this one.

\documentclass{article}      % Specifies the document class

                             % The preamble begins here.
\title{An Example Document}  % Declares the document's title.
\author{Leslie Lamport}      % Declares the author's name.
\date{January 21, 1994}      % Deleting this command produces today's date.

\newcommand{\ip}[2]{(#1, #2)}
                             % Defines \ip{arg1}{arg2} to mean
                             % (arg1, arg2).

%\newcommand{\ip}[2]{\langle #1 | #2\rangle}
                             % This is an alternative definition of
                             % \ip that is commented out.

\begin{document}             % End of preamble and beginning of text.

\maketitle                   % Produces the title.

This is an example input file.  Comparing it with
the output it generates can show you how to
produce a simple document of your own.

\section{Ordinary Text}      % Produces section heading.  Lower-level
                             % sections are begun with similar
                             % \subsection and \subsubsection commands.

The ends  of words and sentences are marked
  by   spaces. It  doesn't matter how many
spaces    you type; one is as good as 100.  The
end of   a line counts as a space.

One   or more   blank lines denote the  end
of  a paragraph.

Since any number of consecutive spaces are treated
like a single one, the formatting of the input
file makes no difference to
      \LaTeX,                % The \LaTeX command generates the LaTeX logo.
but it makes a difference to you.  When you use
\LaTeX, making your input file as easy to read
as possible will be a great help as you write
your document and when you change it.  This sample
file shows how you can add comments to your own input
file.

Because printing is different from typewriting,
there are a number of things that you have to do
differently when preparing an input file than if
you were just typing the document directly.
Quotation marks like
       ``this''
have to be handled specially, as do quotes within
quotes:
       ``\,`this'            % \, separates the double and single quote.
        is what I just
        wrote, not  `that'\,''.

Dashes come in three sizes: an
       intra-word
dash, a medium dash for number ranges like
       1--2,
and a punctuation
       dash---like
this.

A sentence-ending space should be larger than the
space between words within a sentence.  You
sometimes have to type special commands in
conjunction with punctuation characters to get
this right, as in the following sentence.
       Gnats, gnus, etc.\ all  % `\ ' makes an inter-word space.
       begin with G\@.         % \@ marks end-of-sentence punctuation.
You should check the spaces after periods when
reading your output to make sure you haven't
forgotten any special cases.  Generating an
ellipsis
       \ldots\               % `\ ' is needed after `\ldots' because TeX
                             % ignores spaces after command names like \ldots
                             % made from \ + letters.
                             %
                             % Note how a `%' character causes TeX to ignore
                             % the end of the input line, so these blank lines
                             % do not start a new paragraph.
                             %
with the right spacing around the periods requires
a special command.

\LaTeX\ interprets some common characters as
commands, so you must type special commands to
generate them.  These characters include the
following:
       \$ \& \% \# \{ and \}.

In printing, text is usually emphasized with an
       \emph{italic}
type style.

\begin{em}
   A long segment of text can also be emphasized
   in this way.  Text within such a segment can be
   given \emph{additional} emphasis.
\end{em}

It is sometimes necessary to prevent \LaTeX\ from
breaking a line where it might otherwise do so.
This may be at a space, as between the ``Mr.''\ and
``Jones'' in
       ``Mr.~Jones'',        % ~ produces an unbreakable interword space.
or within a word---especially when the word is a
symbol like
       \mbox{\emph{itemnum}}
that makes little sense when hyphenated across
lines.

Footnotes\footnote{This is an example of a footnote.}
pose no problem.

\LaTeX\ is good at typesetting mathematical formulas
like
       \( x-3y + z = 7 \)
or
       \( a_{1} > x^{2n} + y^{2n} > x' \)
or
       \( \ip{A}{B} = \sum_{i} a_{i} b_{i} \).
The spaces you type in a formula are
ignored.  Remember that a letter like
       $x$                   % $ ... $  and  \( ... \)  are equivalent
is a formula when it denotes a mathematical
symbol, and it should be typed as one.

\section{Displayed Text}

Text is displayed by indenting it from the left
margin.  Quotations are commonly displayed.  There
are short quotations
\begin{quote}
   This is a short quotation.  It consists of a
   single paragraph of text.  See how it is formatted.
\end{quote}
and longer ones.
\begin{quotation}
   This is a longer quotation.  It consists of two
   paragraphs of text, neither of which are
   particularly interesting.

   This is the second paragraph of the quotation.  It
   is just as dull as the first paragraph.
\end{quotation}
Another frequently-displayed structure is a list.
The following is an example of an \emph{itemized}
list.
\begin{itemize}
   \item This is the first item of an itemized list.
         Each item in the list is marked with a ``tick''.
         You don't have to worry about what kind of tick
         mark is used.

   \item This is the second item of the list.  It
         contains another list nested inside it.  The inner
         list is an \emph{enumerated} list.
         \begin{enumerate}
            \item This is the first item of an enumerated
                  list that is nested within the itemized list.

            \item This is the second item of the inner list.
                  \LaTeX\ allows you to nest lists deeper than
                  you really should.
         \end{enumerate}
         This is the rest of the second item of the outer
         list.  It is no more interesting than any other
         part of the item.
   \item This is the third item of the list.
\end{itemize}
You can even display poetry.
\begin{verse}
   There is an environment
    for verse \\             % The \\ command separates lines
   Whose features some poets % within a stanza.
   will curse.

                             % One or more blank lines separate stanzas.

   For instead of making\\
   Them do \emph{all} line breaking, \\
   It allows them to put too many words on a line when they'd rather be
   forced to be terse.
\end{verse}

Mathematical formulas may also be displayed.  A
displayed formula
is
one-line long; multiline
formulas require special formatting instructions.
   \[  \ip{\Gamma}{\psi'} = x'' + y^{2} + z_{i}^{n}\]
Don't start a paragraph with a displayed equation,
nor make one a paragraph by itself.

\end{document}               % End of document.

\section{Bench Huge Section 1100}
% synthetic bench section 1100
% This is a sample LaTeX input file.  (Version of 12 August 2004.)
%
% A '%' character causes TeX to ignore all remaining text on the line,
% and is used for comments like this one.

\documentclass{article}      % Specifies the document class

                             % The preamble begins here.
\title{An Example Document}  % Declares the document's title.
\author{Leslie Lamport}      % Declares the author's name.
\date{January 21, 1994}      % Deleting this command produces today's date.

\newcommand{\ip}[2]{(#1, #2)}
                             % Defines \ip{arg1}{arg2} to mean
                             % (arg1, arg2).

%\newcommand{\ip}[2]{\langle #1 | #2\rangle}
                             % This is an alternative definition of
                             % \ip that is commented out.

\begin{document}             % End of preamble and beginning of text.

\maketitle                   % Produces the title.

This is an example input file.  Comparing it with
the output it generates can show you how to
produce a simple document of your own.

\section{Ordinary Text}      % Produces section heading.  Lower-level
                             % sections are begun with similar
                             % \subsection and \subsubsection commands.

The ends  of words and sentences are marked
  by   spaces. It  doesn't matter how many
spaces    you type; one is as good as 100.  The
end of   a line counts as a space.

One   or more   blank lines denote the  end
of  a paragraph.

Since any number of consecutive spaces are treated
like a single one, the formatting of the input
file makes no difference to
      \LaTeX,                % The \LaTeX command generates the LaTeX logo.
but it makes a difference to you.  When you use
\LaTeX, making your input file as easy to read
as possible will be a great help as you write
your document and when you change it.  This sample
file shows how you can add comments to your own input
file.

Because printing is different from typewriting,
there are a number of things that you have to do
differently when preparing an input file than if
you were just typing the document directly.
Quotation marks like
       ``this''
have to be handled specially, as do quotes within
quotes:
       ``\,`this'            % \, separates the double and single quote.
        is what I just
        wrote, not  `that'\,''.

Dashes come in three sizes: an
       intra-word
dash, a medium dash for number ranges like
       1--2,
and a punctuation
       dash---like
this.

A sentence-ending space should be larger than the
space between words within a sentence.  You
sometimes have to type special commands in
conjunction with punctuation characters to get
this right, as in the following sentence.
       Gnats, gnus, etc.\ all  % `\ ' makes an inter-word space.
       begin with G\@.         % \@ marks end-of-sentence punctuation.
You should check the spaces after periods when
reading your output to make sure you haven't
forgotten any special cases.  Generating an
ellipsis
       \ldots\               % `\ ' is needed after `\ldots' because TeX
                             % ignores spaces after command names like \ldots
                             % made from \ + letters.
                             %
                             % Note how a `%' character causes TeX to ignore
                             % the end of the input line, so these blank lines
                             % do not start a new paragraph.
                             %
with the right spacing around the periods requires
a special command.

\LaTeX\ interprets some common characters as
commands, so you must type special commands to
generate them.  These characters include the
following:
       \$ \& \% \# \{ and \}.

In printing, text is usually emphasized with an
       \emph{italic}
type style.

\begin{em}
   A long segment of text can also be emphasized
   in this way.  Text within such a segment can be
   given \emph{additional} emphasis.
\end{em}

It is sometimes necessary to prevent \LaTeX\ from
breaking a line where it might otherwise do so.
This may be at a space, as between the ``Mr.''\ and
``Jones'' in
       ``Mr.~Jones'',        % ~ produces an unbreakable interword space.
or within a word---especially when the word is a
symbol like
       \mbox{\emph{itemnum}}
that makes little sense when hyphenated across
lines.

Footnotes\footnote{This is an example of a footnote.}
pose no problem.

\LaTeX\ is good at typesetting mathematical formulas
like
       \( x-3y + z = 7 \)
or
       \( a_{1} > x^{2n} + y^{2n} > x' \)
or
       \( \ip{A}{B} = \sum_{i} a_{i} b_{i} \).
The spaces you type in a formula are
ignored.  Remember that a letter like
       $x$                   % $ ... $  and  \( ... \)  are equivalent
is a formula when it denotes a mathematical
symbol, and it should be typed as one.

\section{Displayed Text}

Text is displayed by indenting it from the left
margin.  Quotations are commonly displayed.  There
are short quotations
\begin{quote}
   This is a short quotation.  It consists of a
   single paragraph of text.  See how it is formatted.
\end{quote}
and longer ones.
\begin{quotation}
   This is a longer quotation.  It consists of two
   paragraphs of text, neither of which are
   particularly interesting.

   This is the second paragraph of the quotation.  It
   is just as dull as the first paragraph.
\end{quotation}
Another frequently-displayed structure is a list.
The following is an example of an \emph{itemized}
list.
\begin{itemize}
   \item This is the first item of an itemized list.
         Each item in the list is marked with a ``tick''.
         You don't have to worry about what kind of tick
         mark is used.

   \item This is the second item of the list.  It
         contains another list nested inside it.  The inner
         list is an \emph{enumerated} list.
         \begin{enumerate}
            \item This is the first item of an enumerated
                  list that is nested within the itemized list.

            \item This is the second item of the inner list.
                  \LaTeX\ allows you to nest lists deeper than
                  you really should.
         \end{enumerate}
         This is the rest of the second item of the outer
         list.  It is no more interesting than any other
         part of the item.
   \item This is the third item of the list.
\end{itemize}
You can even display poetry.
\begin{verse}
   There is an environment
    for verse \\             % The \\ command separates lines
   Whose features some poets % within a stanza.
   will curse.

                             % One or more blank lines separate stanzas.

   For instead of making\\
   Them do \emph{all} line breaking, \\
   It allows them to put too many words on a line when they'd rather be
   forced to be terse.
\end{verse}

Mathematical formulas may also be displayed.  A
displayed formula
is
one-line long; multiline
formulas require special formatting instructions.
   \[  \ip{\Gamma}{\psi'} = x'' + y^{2} + z_{i}^{n}\]
Don't start a paragraph with a displayed equation,
nor make one a paragraph by itself.

\end{document}               % End of document.

\section{Bench Huge Section 1101}
% synthetic bench section 1101
% This is a sample LaTeX input file.  (Version of 12 August 2004.)
%
% A '%' character causes TeX to ignore all remaining text on the line,
% and is used for comments like this one.

\documentclass{article}      % Specifies the document class

                             % The preamble begins here.
\title{An Example Document}  % Declares the document's title.
\author{Leslie Lamport}      % Declares the author's name.
\date{January 21, 1994}      % Deleting this command produces today's date.

\newcommand{\ip}[2]{(#1, #2)}
                             % Defines \ip{arg1}{arg2} to mean
                             % (arg1, arg2).

%\newcommand{\ip}[2]{\langle #1 | #2\rangle}
                             % This is an alternative definition of
                             % \ip that is commented out.

\begin{document}             % End of preamble and beginning of text.

\maketitle                   % Produces the title.

This is an example input file.  Comparing it with
the output it generates can show you how to
produce a simple document of your own.

\section{Ordinary Text}      % Produces section heading.  Lower-level
                             % sections are begun with similar
                             % \subsection and \subsubsection commands.

The ends  of words and sentences are marked
  by   spaces. It  doesn't matter how many
spaces    you type; one is as good as 100.  The
end of   a line counts as a space.

One   or more   blank lines denote the  end
of  a paragraph.

Since any number of consecutive spaces are treated
like a single one, the formatting of the input
file makes no difference to
      \LaTeX,                % The \LaTeX command generates the LaTeX logo.
but it makes a difference to you.  When you use
\LaTeX, making your input file as easy to read
as possible will be a great help as you write
your document and when you change it.  This sample
file shows how you can add comments to your own input
file.

Because printing is different from typewriting,
there are a number of things that you have to do
differently when preparing an input file than if
you were just typing the document directly.
Quotation marks like
       ``this''
have to be handled specially, as do quotes within
quotes:
       ``\,`this'            % \, separates the double and single quote.
        is what I just
        wrote, not  `that'\,''.

Dashes come in three sizes: an
       intra-word
dash, a medium dash for number ranges like
       1--2,
and a punctuation
       dash---like
this.

A sentence-ending space should be larger than the
space between words within a sentence.  You
sometimes have to type special commands in
conjunction with punctuation characters to get
this right, as in the following sentence.
       Gnats, gnus, etc.\ all  % `\ ' makes an inter-word space.
       begin with G\@.         % \@ marks end-of-sentence punctuation.
You should check the spaces after periods when
reading your output to make sure you haven't
forgotten any special cases.  Generating an
ellipsis
       \ldots\               % `\ ' is needed after `\ldots' because TeX
                             % ignores spaces after command names like \ldots
                             % made from \ + letters.
                             %
                             % Note how a `%' character causes TeX to ignore
                             % the end of the input line, so these blank lines
                             % do not start a new paragraph.
                             %
with the right spacing around the periods requires
a special command.

\LaTeX\ interprets some common characters as
commands, so you must type special commands to
generate them.  These characters include the
following:
       \$ \& \% \# \{ and \}.

In printing, text is usually emphasized with an
       \emph{italic}
type style.

\begin{em}
   A long segment of text can also be emphasized
   in this way.  Text within such a segment can be
   given \emph{additional} emphasis.
\end{em}

It is sometimes necessary to prevent \LaTeX\ from
breaking a line where it might otherwise do so.
This may be at a space, as between the ``Mr.''\ and
``Jones'' in
       ``Mr.~Jones'',        % ~ produces an unbreakable interword space.
or within a word---especially when the word is a
symbol like
       \mbox{\emph{itemnum}}
that makes little sense when hyphenated across
lines.

Footnotes\footnote{This is an example of a footnote.}
pose no problem.

\LaTeX\ is good at typesetting mathematical formulas
like
       \( x-3y + z = 7 \)
or
       \( a_{1} > x^{2n} + y^{2n} > x' \)
or
       \( \ip{A}{B} = \sum_{i} a_{i} b_{i} \).
The spaces you type in a formula are
ignored.  Remember that a letter like
       $x$                   % $ ... $  and  \( ... \)  are equivalent
is a formula when it denotes a mathematical
symbol, and it should be typed as one.

\section{Displayed Text}

Text is displayed by indenting it from the left
margin.  Quotations are commonly displayed.  There
are short quotations
\begin{quote}
   This is a short quotation.  It consists of a
   single paragraph of text.  See how it is formatted.
\end{quote}
and longer ones.
\begin{quotation}
   This is a longer quotation.  It consists of two
   paragraphs of text, neither of which are
   particularly interesting.

   This is the second paragraph of the quotation.  It
   is just as dull as the first paragraph.
\end{quotation}
Another frequently-displayed structure is a list.
The following is an example of an \emph{itemized}
list.
\begin{itemize}
   \item This is the first item of an itemized list.
         Each item in the list is marked with a ``tick''.
         You don't have to worry about what kind of tick
         mark is used.

   \item This is the second item of the list.  It
         contains another list nested inside it.  The inner
         list is an \emph{enumerated} list.
         \begin{enumerate}
            \item This is the first item of an enumerated
                  list that is nested within the itemized list.

            \item This is the second item of the inner list.
                  \LaTeX\ allows you to nest lists deeper than
                  you really should.
         \end{enumerate}
         This is the rest of the second item of the outer
         list.  It is no more interesting than any other
         part of the item.
   \item This is the third item of the list.
\end{itemize}
You can even display poetry.
\begin{verse}
   There is an environment
    for verse \\             % The \\ command separates lines
   Whose features some poets % within a stanza.
   will curse.

                             % One or more blank lines separate stanzas.

   For instead of making\\
   Them do \emph{all} line breaking, \\
   It allows them to put too many words on a line when they'd rather be
   forced to be terse.
\end{verse}

Mathematical formulas may also be displayed.  A
displayed formula
is
one-line long; multiline
formulas require special formatting instructions.
   \[  \ip{\Gamma}{\psi'} = x'' + y^{2} + z_{i}^{n}\]
Don't start a paragraph with a displayed equation,
nor make one a paragraph by itself.

\end{document}               % End of document.

\section{Bench Huge Section 1102}
% synthetic bench section 1102
% This is a sample LaTeX input file.  (Version of 12 August 2004.)
%
% A '%' character causes TeX to ignore all remaining text on the line,
% and is used for comments like this one.

\documentclass{article}      % Specifies the document class

                             % The preamble begins here.
\title{An Example Document}  % Declares the document's title.
\author{Leslie Lamport}      % Declares the author's name.
\date{January 21, 1994}      % Deleting this command produces today's date.

\newcommand{\ip}[2]{(#1, #2)}
                             % Defines \ip{arg1}{arg2} to mean
                             % (arg1, arg2).

%\newcommand{\ip}[2]{\langle #1 | #2\rangle}
                             % This is an alternative definition of
                             % \ip that is commented out.

\begin{document}             % End of preamble and beginning of text.

\maketitle                   % Produces the title.

This is an example input file.  Comparing it with
the output it generates can show you how to
produce a simple document of your own.

\section{Ordinary Text}      % Produces section heading.  Lower-level
                             % sections are begun with similar
                             % \subsection and \subsubsection commands.

The ends  of words and sentences are marked
  by   spaces. It  doesn't matter how many
spaces    you type; one is as good as 100.  The
end of   a line counts as a space.

One   or more   blank lines denote the  end
of  a paragraph.

Since any number of consecutive spaces are treated
like a single one, the formatting of the input
file makes no difference to
      \LaTeX,                % The \LaTeX command generates the LaTeX logo.
but it makes a difference to you.  When you use
\LaTeX, making your input file as easy to read
as possible will be a great help as you write
your document and when you change it.  This sample
file shows how you can add comments to your own input
file.

Because printing is different from typewriting,
there are a number of things that you have to do
differently when preparing an input file than if
you were just typing the document directly.
Quotation marks like
       ``this''
have to be handled specially, as do quotes within
quotes:
       ``\,`this'            % \, separates the double and single quote.
        is what I just
        wrote, not  `that'\,''.

Dashes come in three sizes: an
       intra-word
dash, a medium dash for number ranges like
       1--2,
and a punctuation
       dash---like
this.

A sentence-ending space should be larger than the
space between words within a sentence.  You
sometimes have to type special commands in
conjunction with punctuation characters to get
this right, as in the following sentence.
       Gnats, gnus, etc.\ all  % `\ ' makes an inter-word space.
       begin with G\@.         % \@ marks end-of-sentence punctuation.
You should check the spaces after periods when
reading your output to make sure you haven't
forgotten any special cases.  Generating an
ellipsis
       \ldots\               % `\ ' is needed after `\ldots' because TeX
                             % ignores spaces after command names like \ldots
                             % made from \ + letters.
                             %
                             % Note how a `%' character causes TeX to ignore
                             % the end of the input line, so these blank lines
                             % do not start a new paragraph.
                             %
with the right spacing around the periods requires
a special command.

\LaTeX\ interprets some common characters as
commands, so you must type special commands to
generate them.  These characters include the
following:
       \$ \& \% \# \{ and \}.

In printing, text is usually emphasized with an
       \emph{italic}
type style.

\begin{em}
   A long segment of text can also be emphasized
   in this way.  Text within such a segment can be
   given \emph{additional} emphasis.
\end{em}

It is sometimes necessary to prevent \LaTeX\ from
breaking a line where it might otherwise do so.
This may be at a space, as between the ``Mr.''\ and
``Jones'' in
       ``Mr.~Jones'',        % ~ produces an unbreakable interword space.
or within a word---especially when the word is a
symbol like
       \mbox{\emph{itemnum}}
that makes little sense when hyphenated across
lines.

Footnotes\footnote{This is an example of a footnote.}
pose no problem.

\LaTeX\ is good at typesetting mathematical formulas
like
       \( x-3y + z = 7 \)
or
       \( a_{1} > x^{2n} + y^{2n} > x' \)
or
       \( \ip{A}{B} = \sum_{i} a_{i} b_{i} \).
The spaces you type in a formula are
ignored.  Remember that a letter like
       $x$                   % $ ... $  and  \( ... \)  are equivalent
is a formula when it denotes a mathematical
symbol, and it should be typed as one.

\section{Displayed Text}

Text is displayed by indenting it from the left
margin.  Quotations are commonly displayed.  There
are short quotations
\begin{quote}
   This is a short quotation.  It consists of a
   single paragraph of text.  See how it is formatted.
\end{quote}
and longer ones.
\begin{quotation}
   This is a longer quotation.  It consists of two
   paragraphs of text, neither of which are
   particularly interesting.

   This is the second paragraph of the quotation.  It
   is just as dull as the first paragraph.
\end{quotation}
Another frequently-displayed structure is a list.
The following is an example of an \emph{itemized}
list.
\begin{itemize}
   \item This is the first item of an itemized list.
         Each item in the list is marked with a ``tick''.
         You don't have to worry about what kind of tick
         mark is used.

   \item This is the second item of the list.  It
         contains another list nested inside it.  The inner
         list is an \emph{enumerated} list.
         \begin{enumerate}
            \item This is the first item of an enumerated
                  list that is nested within the itemized list.

            \item This is the second item of the inner list.
                  \LaTeX\ allows you to nest lists deeper than
                  you really should.
         \end{enumerate}
         This is the rest of the second item of the outer
         list.  It is no more interesting than any other
         part of the item.
   \item This is the third item of the list.
\end{itemize}
You can even display poetry.
\begin{verse}
   There is an environment
    for verse \\             % The \\ command separates lines
   Whose features some poets % within a stanza.
   will curse.

                             % One or more blank lines separate stanzas.

   For instead of making\\
   Them do \emph{all} line breaking, \\
   It allows them to put too many words on a line when they'd rather be
   forced to be terse.
\end{verse}

Mathematical formulas may also be displayed.  A
displayed formula
is
one-line long; multiline
formulas require special formatting instructions.
   \[  \ip{\Gamma}{\psi'} = x'' + y^{2} + z_{i}^{n}\]
Don't start a paragraph with a displayed equation,
nor make one a paragraph by itself.

\end{document}               % End of document.

\section{Bench Huge Section 1103}
% synthetic bench section 1103
% This is a sample LaTeX input file.  (Version of 12 August 2004.)
%
% A '%' character causes TeX to ignore all remaining text on the line,
% and is used for comments like this one.

\documentclass{article}      % Specifies the document class

                             % The preamble begins here.
\title{An Example Document}  % Declares the document's title.
\author{Leslie Lamport}      % Declares the author's name.
\date{January 21, 1994}      % Deleting this command produces today's date.

\newcommand{\ip}[2]{(#1, #2)}
                             % Defines \ip{arg1}{arg2} to mean
                             % (arg1, arg2).

%\newcommand{\ip}[2]{\langle #1 | #2\rangle}
                             % This is an alternative definition of
                             % \ip that is commented out.

\begin{document}             % End of preamble and beginning of text.

\maketitle                   % Produces the title.

This is an example input file.  Comparing it with
the output it generates can show you how to
produce a simple document of your own.

\section{Ordinary Text}      % Produces section heading.  Lower-level
                             % sections are begun with similar
                             % \subsection and \subsubsection commands.

The ends  of words and sentences are marked
  by   spaces. It  doesn't matter how many
spaces    you type; one is as good as 100.  The
end of   a line counts as a space.

One   or more   blank lines denote the  end
of  a paragraph.

Since any number of consecutive spaces are treated
like a single one, the formatting of the input
file makes no difference to
      \LaTeX,                % The \LaTeX command generates the LaTeX logo.
but it makes a difference to you.  When you use
\LaTeX, making your input file as easy to read
as possible will be a great help as you write
your document and when you change it.  This sample
file shows how you can add comments to your own input
file.

Because printing is different from typewriting,
there are a number of things that you have to do
differently when preparing an input file than if
you were just typing the document directly.
Quotation marks like
       ``this''
have to be handled specially, as do quotes within
quotes:
       ``\,`this'            % \, separates the double and single quote.
        is what I just
        wrote, not  `that'\,''.

Dashes come in three sizes: an
       intra-word
dash, a medium dash for number ranges like
       1--2,
and a punctuation
       dash---like
this.

A sentence-ending space should be larger than the
space between words within a sentence.  You
sometimes have to type special commands in
conjunction with punctuation characters to get
this right, as in the following sentence.
       Gnats, gnus, etc.\ all  % `\ ' makes an inter-word space.
       begin with G\@.         % \@ marks end-of-sentence punctuation.
You should check the spaces after periods when
reading your output to make sure you haven't
forgotten any special cases.  Generating an
ellipsis
       \ldots\               % `\ ' is needed after `\ldots' because TeX
                             % ignores spaces after command names like \ldots
                             % made from \ + letters.
                             %
                             % Note how a `%' character causes TeX to ignore
                             % the end of the input line, so these blank lines
                             % do not start a new paragraph.
                             %
with the right spacing around the periods requires
a special command.

\LaTeX\ interprets some common characters as
commands, so you must type special commands to
generate them.  These characters include the
following:
       \$ \& \% \# \{ and \}.

In printing, text is usually emphasized with an
       \emph{italic}
type style.

\begin{em}
   A long segment of text can also be emphasized
   in this way.  Text within such a segment can be
   given \emph{additional} emphasis.
\end{em}

It is sometimes necessary to prevent \LaTeX\ from
breaking a line where it might otherwise do so.
This may be at a space, as between the ``Mr.''\ and
``Jones'' in
       ``Mr.~Jones'',        % ~ produces an unbreakable interword space.
or within a word---especially when the word is a
symbol like
       \mbox{\emph{itemnum}}
that makes little sense when hyphenated across
lines.

Footnotes\footnote{This is an example of a footnote.}
pose no problem.

\LaTeX\ is good at typesetting mathematical formulas
like
       \( x-3y + z = 7 \)
or
       \( a_{1} > x^{2n} + y^{2n} > x' \)
or
       \( \ip{A}{B} = \sum_{i} a_{i} b_{i} \).
The spaces you type in a formula are
ignored.  Remember that a letter like
       $x$                   % $ ... $  and  \( ... \)  are equivalent
is a formula when it denotes a mathematical
symbol, and it should be typed as one.

\section{Displayed Text}

Text is displayed by indenting it from the left
margin.  Quotations are commonly displayed.  There
are short quotations
\begin{quote}
   This is a short quotation.  It consists of a
   single paragraph of text.  See how it is formatted.
\end{quote}
and longer ones.
\begin{quotation}
   This is a longer quotation.  It consists of two
   paragraphs of text, neither of which are
   particularly interesting.

   This is the second paragraph of the quotation.  It
   is just as dull as the first paragraph.
\end{quotation}
Another frequently-displayed structure is a list.
The following is an example of an \emph{itemized}
list.
\begin{itemize}
   \item This is the first item of an itemized list.
         Each item in the list is marked with a ``tick''.
         You don't have to worry about what kind of tick
         mark is used.

   \item This is the second item of the list.  It
         contains another list nested inside it.  The inner
         list is an \emph{enumerated} list.
         \begin{enumerate}
            \item This is the first item of an enumerated
                  list that is nested within the itemized list.

            \item This is the second item of the inner list.
                  \LaTeX\ allows you to nest lists deeper than
                  you really should.
         \end{enumerate}
         This is the rest of the second item of the outer
         list.  It is no more interesting than any other
         part of the item.
   \item This is the third item of the list.
\end{itemize}
You can even display poetry.
\begin{verse}
   There is an environment
    for verse \\             % The \\ command separates lines
   Whose features some poets % within a stanza.
   will curse.

                             % One or more blank lines separate stanzas.

   For instead of making\\
   Them do \emph{all} line breaking, \\
   It allows them to put too many words on a line when they'd rather be
   forced to be terse.
\end{verse}

Mathematical formulas may also be displayed.  A
displayed formula
is
one-line long; multiline
formulas require special formatting instructions.
   \[  \ip{\Gamma}{\psi'} = x'' + y^{2} + z_{i}^{n}\]
Don't start a paragraph with a displayed equation,
nor make one a paragraph by itself.

\end{document}               % End of document.

\section{Bench Huge Section 1104}
% synthetic bench section 1104
% This is a sample LaTeX input file.  (Version of 12 August 2004.)
%
% A '%' character causes TeX to ignore all remaining text on the line,
% and is used for comments like this one.

\documentclass{article}      % Specifies the document class

                             % The preamble begins here.
\title{An Example Document}  % Declares the document's title.
\author{Leslie Lamport}      % Declares the author's name.
\date{January 21, 1994}      % Deleting this command produces today's date.

\newcommand{\ip}[2]{(#1, #2)}
                             % Defines \ip{arg1}{arg2} to mean
                             % (arg1, arg2).

%\newcommand{\ip}[2]{\langle #1 | #2\rangle}
                             % This is an alternative definition of
                             % \ip that is commented out.

\begin{document}             % End of preamble and beginning of text.

\maketitle                   % Produces the title.

This is an example input file.  Comparing it with
the output it generates can show you how to
produce a simple document of your own.

\section{Ordinary Text}      % Produces section heading.  Lower-level
                             % sections are begun with similar
                             % \subsection and \subsubsection commands.

The ends  of words and sentences are marked
  by   spaces. It  doesn't matter how many
spaces    you type; one is as good as 100.  The
end of   a line counts as a space.

One   or more   blank lines denote the  end
of  a paragraph.

Since any number of consecutive spaces are treated
like a single one, the formatting of the input
file makes no difference to
      \LaTeX,                % The \LaTeX command generates the LaTeX logo.
but it makes a difference to you.  When you use
\LaTeX, making your input file as easy to read
as possible will be a great help as you write
your document and when you change it.  This sample
file shows how you can add comments to your own input
file.

Because printing is different from typewriting,
there are a number of things that you have to do
differently when preparing an input file than if
you were just typing the document directly.
Quotation marks like
       ``this''
have to be handled specially, as do quotes within
quotes:
       ``\,`this'            % \, separates the double and single quote.
        is what I just
        wrote, not  `that'\,''.

Dashes come in three sizes: an
       intra-word
dash, a medium dash for number ranges like
       1--2,
and a punctuation
       dash---like
this.

A sentence-ending space should be larger than the
space between words within a sentence.  You
sometimes have to type special commands in
conjunction with punctuation characters to get
this right, as in the following sentence.
       Gnats, gnus, etc.\ all  % `\ ' makes an inter-word space.
       begin with G\@.         % \@ marks end-of-sentence punctuation.
You should check the spaces after periods when
reading your output to make sure you haven't
forgotten any special cases.  Generating an
ellipsis
       \ldots\               % `\ ' is needed after `\ldots' because TeX
                             % ignores spaces after command names like \ldots
                             % made from \ + letters.
                             %
                             % Note how a `%' character causes TeX to ignore
                             % the end of the input line, so these blank lines
                             % do not start a new paragraph.
                             %
with the right spacing around the periods requires
a special command.

\LaTeX\ interprets some common characters as
commands, so you must type special commands to
generate them.  These characters include the
following:
       \$ \& \% \# \{ and \}.

In printing, text is usually emphasized with an
       \emph{italic}
type style.

\begin{em}
   A long segment of text can also be emphasized
   in this way.  Text within such a segment can be
   given \emph{additional} emphasis.
\end{em}

It is sometimes necessary to prevent \LaTeX\ from
breaking a line where it might otherwise do so.
This may be at a space, as between the ``Mr.''\ and
``Jones'' in
       ``Mr.~Jones'',        % ~ produces an unbreakable interword space.
or within a word---especially when the word is a
symbol like
       \mbox{\emph{itemnum}}
that makes little sense when hyphenated across
lines.

Footnotes\footnote{This is an example of a footnote.}
pose no problem.

\LaTeX\ is good at typesetting mathematical formulas
like
       \( x-3y + z = 7 \)
or
       \( a_{1} > x^{2n} + y^{2n} > x' \)
or
       \( \ip{A}{B} = \sum_{i} a_{i} b_{i} \).
The spaces you type in a formula are
ignored.  Remember that a letter like
       $x$                   % $ ... $  and  \( ... \)  are equivalent
is a formula when it denotes a mathematical
symbol, and it should be typed as one.

\section{Displayed Text}

Text is displayed by indenting it from the left
margin.  Quotations are commonly displayed.  There
are short quotations
\begin{quote}
   This is a short quotation.  It consists of a
   single paragraph of text.  See how it is formatted.
\end{quote}
and longer ones.
\begin{quotation}
   This is a longer quotation.  It consists of two
   paragraphs of text, neither of which are
   particularly interesting.

   This is the second paragraph of the quotation.  It
   is just as dull as the first paragraph.
\end{quotation}
Another frequently-displayed structure is a list.
The following is an example of an \emph{itemized}
list.
\begin{itemize}
   \item This is the first item of an itemized list.
         Each item in the list is marked with a ``tick''.
         You don't have to worry about what kind of tick
         mark is used.

   \item This is the second item of the list.  It
         contains another list nested inside it.  The inner
         list is an \emph{enumerated} list.
         \begin{enumerate}
            \item This is the first item of an enumerated
                  list that is nested within the itemized list.

            \item This is the second item of the inner list.
                  \LaTeX\ allows you to nest lists deeper than
                  you really should.
         \end{enumerate}
         This is the rest of the second item of the outer
         list.  It is no more interesting than any other
         part of the item.
   \item This is the third item of the list.
\end{itemize}
You can even display poetry.
\begin{verse}
   There is an environment
    for verse \\             % The \\ command separates lines
   Whose features some poets % within a stanza.
   will curse.

                             % One or more blank lines separate stanzas.

   For instead of making\\
   Them do \emph{all} line breaking, \\
   It allows them to put too many words on a line when they'd rather be
   forced to be terse.
\end{verse}

Mathematical formulas may also be displayed.  A
displayed formula
is
one-line long; multiline
formulas require special formatting instructions.
   \[  \ip{\Gamma}{\psi'} = x'' + y^{2} + z_{i}^{n}\]
Don't start a paragraph with a displayed equation,
nor make one a paragraph by itself.

\end{document}               % End of document.

\section{Bench Huge Section 1105}
% synthetic bench section 1105
% This is a sample LaTeX input file.  (Version of 12 August 2004.)
%
% A '%' character causes TeX to ignore all remaining text on the line,
% and is used for comments like this one.

\documentclass{article}      % Specifies the document class

                             % The preamble begins here.
\title{An Example Document}  % Declares the document's title.
\author{Leslie Lamport}      % Declares the author's name.
\date{January 21, 1994}      % Deleting this command produces today's date.

\newcommand{\ip}[2]{(#1, #2)}
                             % Defines \ip{arg1}{arg2} to mean
                             % (arg1, arg2).

%\newcommand{\ip}[2]{\langle #1 | #2\rangle}
                             % This is an alternative definition of
                             % \ip that is commented out.

\begin{document}             % End of preamble and beginning of text.

\maketitle                   % Produces the title.

This is an example input file.  Comparing it with
the output it generates can show you how to
produce a simple document of your own.

\section{Ordinary Text}      % Produces section heading.  Lower-level
                             % sections are begun with similar
                             % \subsection and \subsubsection commands.

The ends  of words and sentences are marked
  by   spaces. It  doesn't matter how many
spaces    you type; one is as good as 100.  The
end of   a line counts as a space.

One   or more   blank lines denote the  end
of  a paragraph.

Since any number of consecutive spaces are treated
like a single one, the formatting of the input
file makes no difference to
      \LaTeX,                % The \LaTeX command generates the LaTeX logo.
but it makes a difference to you.  When you use
\LaTeX, making your input file as easy to read
as possible will be a great help as you write
your document and when you change it.  This sample
file shows how you can add comments to your own input
file.

Because printing is different from typewriting,
there are a number of things that you have to do
differently when preparing an input file than if
you were just typing the document directly.
Quotation marks like
       ``this''
have to be handled specially, as do quotes within
quotes:
       ``\,`this'            % \, separates the double and single quote.
        is what I just
        wrote, not  `that'\,''.

Dashes come in three sizes: an
       intra-word
dash, a medium dash for number ranges like
       1--2,
and a punctuation
       dash---like
this.

A sentence-ending space should be larger than the
space between words within a sentence.  You
sometimes have to type special commands in
conjunction with punctuation characters to get
this right, as in the following sentence.
       Gnats, gnus, etc.\ all  % `\ ' makes an inter-word space.
       begin with G\@.         % \@ marks end-of-sentence punctuation.
You should check the spaces after periods when
reading your output to make sure you haven't
forgotten any special cases.  Generating an
ellipsis
       \ldots\               % `\ ' is needed after `\ldots' because TeX
                             % ignores spaces after command names like \ldots
                             % made from \ + letters.
                             %
                             % Note how a `%' character causes TeX to ignore
                             % the end of the input line, so these blank lines
                             % do not start a new paragraph.
                             %
with the right spacing around the periods requires
a special command.

\LaTeX\ interprets some common characters as
commands, so you must type special commands to
generate them.  These characters include the
following:
       \$ \& \% \# \{ and \}.

In printing, text is usually emphasized with an
       \emph{italic}
type style.

\begin{em}
   A long segment of text can also be emphasized
   in this way.  Text within such a segment can be
   given \emph{additional} emphasis.
\end{em}

It is sometimes necessary to prevent \LaTeX\ from
breaking a line where it might otherwise do so.
This may be at a space, as between the ``Mr.''\ and
``Jones'' in
       ``Mr.~Jones'',        % ~ produces an unbreakable interword space.
or within a word---especially when the word is a
symbol like
       \mbox{\emph{itemnum}}
that makes little sense when hyphenated across
lines.

Footnotes\footnote{This is an example of a footnote.}
pose no problem.

\LaTeX\ is good at typesetting mathematical formulas
like
       \( x-3y + z = 7 \)
or
       \( a_{1} > x^{2n} + y^{2n} > x' \)
or
       \( \ip{A}{B} = \sum_{i} a_{i} b_{i} \).
The spaces you type in a formula are
ignored.  Remember that a letter like
       $x$                   % $ ... $  and  \( ... \)  are equivalent
is a formula when it denotes a mathematical
symbol, and it should be typed as one.

\section{Displayed Text}

Text is displayed by indenting it from the left
margin.  Quotations are commonly displayed.  There
are short quotations
\begin{quote}
   This is a short quotation.  It consists of a
   single paragraph of text.  See how it is formatted.
\end{quote}
and longer ones.
\begin{quotation}
   This is a longer quotation.  It consists of two
   paragraphs of text, neither of which are
   particularly interesting.

   This is the second paragraph of the quotation.  It
   is just as dull as the first paragraph.
\end{quotation}
Another frequently-displayed structure is a list.
The following is an example of an \emph{itemized}
list.
\begin{itemize}
   \item This is the first item of an itemized list.
         Each item in the list is marked with a ``tick''.
         You don't have to worry about what kind of tick
         mark is used.

   \item This is the second item of the list.  It
         contains another list nested inside it.  The inner
         list is an \emph{enumerated} list.
         \begin{enumerate}
            \item This is the first item of an enumerated
                  list that is nested within the itemized list.

            \item This is the second item of the inner list.
                  \LaTeX\ allows you to nest lists deeper than
                  you really should.
         \end{enumerate}
         This is the rest of the second item of the outer
         list.  It is no more interesting than any other
         part of the item.
   \item This is the third item of the list.
\end{itemize}
You can even display poetry.
\begin{verse}
   There is an environment
    for verse \\             % The \\ command separates lines
   Whose features some poets % within a stanza.
   will curse.

                             % One or more blank lines separate stanzas.

   For instead of making\\
   Them do \emph{all} line breaking, \\
   It allows them to put too many words on a line when they'd rather be
   forced to be terse.
\end{verse}

Mathematical formulas may also be displayed.  A
displayed formula
is
one-line long; multiline
formulas require special formatting instructions.
   \[  \ip{\Gamma}{\psi'} = x'' + y^{2} + z_{i}^{n}\]
Don't start a paragraph with a displayed equation,
nor make one a paragraph by itself.

\end{document}               % End of document.

\section{Bench Huge Section 1106}
% synthetic bench section 1106
% This is a sample LaTeX input file.  (Version of 12 August 2004.)
%
% A '%' character causes TeX to ignore all remaining text on the line,
% and is used for comments like this one.

\documentclass{article}      % Specifies the document class

                             % The preamble begins here.
\title{An Example Document}  % Declares the document's title.
\author{Leslie Lamport}      % Declares the author's name.
\date{January 21, 1994}      % Deleting this command produces today's date.

\newcommand{\ip}[2]{(#1, #2)}
                             % Defines \ip{arg1}{arg2} to mean
                             % (arg1, arg2).

%\newcommand{\ip}[2]{\langle #1 | #2\rangle}
                             % This is an alternative definition of
                             % \ip that is commented out.

\begin{document}             % End of preamble and beginning of text.

\maketitle                   % Produces the title.

This is an example input file.  Comparing it with
the output it generates can show you how to
produce a simple document of your own.

\section{Ordinary Text}      % Produces section heading.  Lower-level
                             % sections are begun with similar
                             % \subsection and \subsubsection commands.

The ends  of words and sentences are marked
  by   spaces. It  doesn't matter how many
spaces    you type; one is as good as 100.  The
end of   a line counts as a space.

One   or more   blank lines denote the  end
of  a paragraph.

Since any number of consecutive spaces are treated
like a single one, the formatting of the input
file makes no difference to
      \LaTeX,                % The \LaTeX command generates the LaTeX logo.
but it makes a difference to you.  When you use
\LaTeX, making your input file as easy to read
as possible will be a great help as you write
your document and when you change it.  This sample
file shows how you can add comments to your own input
file.

Because printing is different from typewriting,
there are a number of things that you have to do
differently when preparing an input file than if
you were just typing the document directly.
Quotation marks like
       ``this''
have to be handled specially, as do quotes within
quotes:
       ``\,`this'            % \, separates the double and single quote.
        is what I just
        wrote, not  `that'\,''.

Dashes come in three sizes: an
       intra-word
dash, a medium dash for number ranges like
       1--2,
and a punctuation
       dash---like
this.

A sentence-ending space should be larger than the
space between words within a sentence.  You
sometimes have to type special commands in
conjunction with punctuation characters to get
this right, as in the following sentence.
       Gnats, gnus, etc.\ all  % `\ ' makes an inter-word space.
       begin with G\@.         % \@ marks end-of-sentence punctuation.
You should check the spaces after periods when
reading your output to make sure you haven't
forgotten any special cases.  Generating an
ellipsis
       \ldots\               % `\ ' is needed after `\ldots' because TeX
                             % ignores spaces after command names like \ldots
                             % made from \ + letters.
                             %
                             % Note how a `%' character causes TeX to ignore
                             % the end of the input line, so these blank lines
                             % do not start a new paragraph.
                             %
with the right spacing around the periods requires
a special command.

\LaTeX\ interprets some common characters as
commands, so you must type special commands to
generate them.  These characters include the
following:
       \$ \& \% \# \{ and \}.

In printing, text is usually emphasized with an
       \emph{italic}
type style.

\begin{em}
   A long segment of text can also be emphasized
   in this way.  Text within such a segment can be
   given \emph{additional} emphasis.
\end{em}

It is sometimes necessary to prevent \LaTeX\ from
breaking a line where it might otherwise do so.
This may be at a space, as between the ``Mr.''\ and
``Jones'' in
       ``Mr.~Jones'',        % ~ produces an unbreakable interword space.
or within a word---especially when the word is a
symbol like
       \mbox{\emph{itemnum}}
that makes little sense when hyphenated across
lines.

Footnotes\footnote{This is an example of a footnote.}
pose no problem.

\LaTeX\ is good at typesetting mathematical formulas
like
       \( x-3y + z = 7 \)
or
       \( a_{1} > x^{2n} + y^{2n} > x' \)
or
       \( \ip{A}{B} = \sum_{i} a_{i} b_{i} \).
The spaces you type in a formula are
ignored.  Remember that a letter like
       $x$                   % $ ... $  and  \( ... \)  are equivalent
is a formula when it denotes a mathematical
symbol, and it should be typed as one.

\section{Displayed Text}

Text is displayed by indenting it from the left
margin.  Quotations are commonly displayed.  There
are short quotations
\begin{quote}
   This is a short quotation.  It consists of a
   single paragraph of text.  See how it is formatted.
\end{quote}
and longer ones.
\begin{quotation}
   This is a longer quotation.  It consists of two
   paragraphs of text, neither of which are
   particularly interesting.

   This is the second paragraph of the quotation.  It
   is just as dull as the first paragraph.
\end{quotation}
Another frequently-displayed structure is a list.
The following is an example of an \emph{itemized}
list.
\begin{itemize}
   \item This is the first item of an itemized list.
         Each item in the list is marked with a ``tick''.
         You don't have to worry about what kind of tick
         mark is used.

   \item This is the second item of the list.  It
         contains another list nested inside it.  The inner
         list is an \emph{enumerated} list.
         \begin{enumerate}
            \item This is the first item of an enumerated
                  list that is nested within the itemized list.

            \item This is the second item of the inner list.
                  \LaTeX\ allows you to nest lists deeper than
                  you really should.
         \end{enumerate}
         This is the rest of the second item of the outer
         list.  It is no more interesting than any other
         part of the item.
   \item This is the third item of the list.
\end{itemize}
You can even display poetry.
\begin{verse}
   There is an environment
    for verse \\             % The \\ command separates lines
   Whose features some poets % within a stanza.
   will curse.

                             % One or more blank lines separate stanzas.

   For instead of making\\
   Them do \emph{all} line breaking, \\
   It allows them to put too many words on a line when they'd rather be
   forced to be terse.
\end{verse}

Mathematical formulas may also be displayed.  A
displayed formula
is
one-line long; multiline
formulas require special formatting instructions.
   \[  \ip{\Gamma}{\psi'} = x'' + y^{2} + z_{i}^{n}\]
Don't start a paragraph with a displayed equation,
nor make one a paragraph by itself.

\end{document}               % End of document.

\section{Bench Huge Section 1107}
% synthetic bench section 1107
% This is a sample LaTeX input file.  (Version of 12 August 2004.)
%
% A '%' character causes TeX to ignore all remaining text on the line,
% and is used for comments like this one.

\documentclass{article}      % Specifies the document class

                             % The preamble begins here.
\title{An Example Document}  % Declares the document's title.
\author{Leslie Lamport}      % Declares the author's name.
\date{January 21, 1994}      % Deleting this command produces today's date.

\newcommand{\ip}[2]{(#1, #2)}
                             % Defines \ip{arg1}{arg2} to mean
                             % (arg1, arg2).

%\newcommand{\ip}[2]{\langle #1 | #2\rangle}
                             % This is an alternative definition of
                             % \ip that is commented out.

\begin{document}             % End of preamble and beginning of text.

\maketitle                   % Produces the title.

This is an example input file.  Comparing it with
the output it generates can show you how to
produce a simple document of your own.

\section{Ordinary Text}      % Produces section heading.  Lower-level
                             % sections are begun with similar
                             % \subsection and \subsubsection commands.

The ends  of words and sentences are marked
  by   spaces. It  doesn't matter how many
spaces    you type; one is as good as 100.  The
end of   a line counts as a space.

One   or more   blank lines denote the  end
of  a paragraph.

Since any number of consecutive spaces are treated
like a single one, the formatting of the input
file makes no difference to
      \LaTeX,                % The \LaTeX command generates the LaTeX logo.
but it makes a difference to you.  When you use
\LaTeX, making your input file as easy to read
as possible will be a great help as you write
your document and when you change it.  This sample
file shows how you can add comments to your own input
file.

Because printing is different from typewriting,
there are a number of things that you have to do
differently when preparing an input file than if
you were just typing the document directly.
Quotation marks like
       ``this''
have to be handled specially, as do quotes within
quotes:
       ``\,`this'            % \, separates the double and single quote.
        is what I just
        wrote, not  `that'\,''.

Dashes come in three sizes: an
       intra-word
dash, a medium dash for number ranges like
       1--2,
and a punctuation
       dash---like
this.

A sentence-ending space should be larger than the
space between words within a sentence.  You
sometimes have to type special commands in
conjunction with punctuation characters to get
this right, as in the following sentence.
       Gnats, gnus, etc.\ all  % `\ ' makes an inter-word space.
       begin with G\@.         % \@ marks end-of-sentence punctuation.
You should check the spaces after periods when
reading your output to make sure you haven't
forgotten any special cases.  Generating an
ellipsis
       \ldots\               % `\ ' is needed after `\ldots' because TeX
                             % ignores spaces after command names like \ldots
                             % made from \ + letters.
                             %
                             % Note how a `%' character causes TeX to ignore
                             % the end of the input line, so these blank lines
                             % do not start a new paragraph.
                             %
with the right spacing around the periods requires
a special command.

\LaTeX\ interprets some common characters as
commands, so you must type special commands to
generate them.  These characters include the
following:
       \$ \& \% \# \{ and \}.

In printing, text is usually emphasized with an
       \emph{italic}
type style.

\begin{em}
   A long segment of text can also be emphasized
   in this way.  Text within such a segment can be
   given \emph{additional} emphasis.
\end{em}

It is sometimes necessary to prevent \LaTeX\ from
breaking a line where it might otherwise do so.
This may be at a space, as between the ``Mr.''\ and
``Jones'' in
       ``Mr.~Jones'',        % ~ produces an unbreakable interword space.
or within a word---especially when the word is a
symbol like
       \mbox{\emph{itemnum}}
that makes little sense when hyphenated across
lines.

Footnotes\footnote{This is an example of a footnote.}
pose no problem.

\LaTeX\ is good at typesetting mathematical formulas
like
       \( x-3y + z = 7 \)
or
       \( a_{1} > x^{2n} + y^{2n} > x' \)
or
       \( \ip{A}{B} = \sum_{i} a_{i} b_{i} \).
The spaces you type in a formula are
ignored.  Remember that a letter like
       $x$                   % $ ... $  and  \( ... \)  are equivalent
is a formula when it denotes a mathematical
symbol, and it should be typed as one.

\section{Displayed Text}

Text is displayed by indenting it from the left
margin.  Quotations are commonly displayed.  There
are short quotations
\begin{quote}
   This is a short quotation.  It consists of a
   single paragraph of text.  See how it is formatted.
\end{quote}
and longer ones.
\begin{quotation}
   This is a longer quotation.  It consists of two
   paragraphs of text, neither of which are
   particularly interesting.

   This is the second paragraph of the quotation.  It
   is just as dull as the first paragraph.
\end{quotation}
Another frequently-displayed structure is a list.
The following is an example of an \emph{itemized}
list.
\begin{itemize}
   \item This is the first item of an itemized list.
         Each item in the list is marked with a ``tick''.
         You don't have to worry about what kind of tick
         mark is used.

   \item This is the second item of the list.  It
         contains another list nested inside it.  The inner
         list is an \emph{enumerated} list.
         \begin{enumerate}
            \item This is the first item of an enumerated
                  list that is nested within the itemized list.

            \item This is the second item of the inner list.
                  \LaTeX\ allows you to nest lists deeper than
                  you really should.
         \end{enumerate}
         This is the rest of the second item of the outer
         list.  It is no more interesting than any other
         part of the item.
   \item This is the third item of the list.
\end{itemize}
You can even display poetry.
\begin{verse}
   There is an environment
    for verse \\             % The \\ command separates lines
   Whose features some poets % within a stanza.
   will curse.

                             % One or more blank lines separate stanzas.

   For instead of making\\
   Them do \emph{all} line breaking, \\
   It allows them to put too many words on a line when they'd rather be
   forced to be terse.
\end{verse}

Mathematical formulas may also be displayed.  A
displayed formula
is
one-line long; multiline
formulas require special formatting instructions.
   \[  \ip{\Gamma}{\psi'} = x'' + y^{2} + z_{i}^{n}\]
Don't start a paragraph with a displayed equation,
nor make one a paragraph by itself.

\end{document}               % End of document.

\section{Bench Huge Section 1108}
% synthetic bench section 1108
% This is a sample LaTeX input file.  (Version of 12 August 2004.)
%
% A '%' character causes TeX to ignore all remaining text on the line,
% and is used for comments like this one.

\documentclass{article}      % Specifies the document class

                             % The preamble begins here.
\title{An Example Document}  % Declares the document's title.
\author{Leslie Lamport}      % Declares the author's name.
\date{January 21, 1994}      % Deleting this command produces today's date.

\newcommand{\ip}[2]{(#1, #2)}
                             % Defines \ip{arg1}{arg2} to mean
                             % (arg1, arg2).

%\newcommand{\ip}[2]{\langle #1 | #2\rangle}
                             % This is an alternative definition of
                             % \ip that is commented out.

\begin{document}             % End of preamble and beginning of text.

\maketitle                   % Produces the title.

This is an example input file.  Comparing it with
the output it generates can show you how to
produce a simple document of your own.

\section{Ordinary Text}      % Produces section heading.  Lower-level
                             % sections are begun with similar
                             % \subsection and \subsubsection commands.

The ends  of words and sentences are marked
  by   spaces. It  doesn't matter how many
spaces    you type; one is as good as 100.  The
end of   a line counts as a space.

One   or more   blank lines denote the  end
of  a paragraph.

Since any number of consecutive spaces are treated
like a single one, the formatting of the input
file makes no difference to
      \LaTeX,                % The \LaTeX command generates the LaTeX logo.
but it makes a difference to you.  When you use
\LaTeX, making your input file as easy to read
as possible will be a great help as you write
your document and when you change it.  This sample
file shows how you can add comments to your own input
file.

Because printing is different from typewriting,
there are a number of things that you have to do
differently when preparing an input file than if
you were just typing the document directly.
Quotation marks like
       ``this''
have to be handled specially, as do quotes within
quotes:
       ``\,`this'            % \, separates the double and single quote.
        is what I just
        wrote, not  `that'\,''.

Dashes come in three sizes: an
       intra-word
dash, a medium dash for number ranges like
       1--2,
and a punctuation
       dash---like
this.

A sentence-ending space should be larger than the
space between words within a sentence.  You
sometimes have to type special commands in
conjunction with punctuation characters to get
this right, as in the following sentence.
       Gnats, gnus, etc.\ all  % `\ ' makes an inter-word space.
       begin with G\@.         % \@ marks end-of-sentence punctuation.
You should check the spaces after periods when
reading your output to make sure you haven't
forgotten any special cases.  Generating an
ellipsis
       \ldots\               % `\ ' is needed after `\ldots' because TeX
                             % ignores spaces after command names like \ldots
                             % made from \ + letters.
                             %
                             % Note how a `%' character causes TeX to ignore
                             % the end of the input line, so these blank lines
                             % do not start a new paragraph.
                             %
with the right spacing around the periods requires
a special command.

\LaTeX\ interprets some common characters as
commands, so you must type special commands to
generate them.  These characters include the
following:
       \$ \& \% \# \{ and \}.

In printing, text is usually emphasized with an
       \emph{italic}
type style.

\begin{em}
   A long segment of text can also be emphasized
   in this way.  Text within such a segment can be
   given \emph{additional} emphasis.
\end{em}

It is sometimes necessary to prevent \LaTeX\ from
breaking a line where it might otherwise do so.
This may be at a space, as between the ``Mr.''\ and
``Jones'' in
       ``Mr.~Jones'',        % ~ produces an unbreakable interword space.
or within a word---especially when the word is a
symbol like
       \mbox{\emph{itemnum}}
that makes little sense when hyphenated across
lines.

Footnotes\footnote{This is an example of a footnote.}
pose no problem.

\LaTeX\ is good at typesetting mathematical formulas
like
       \( x-3y + z = 7 \)
or
       \( a_{1} > x^{2n} + y^{2n} > x' \)
or
       \( \ip{A}{B} = \sum_{i} a_{i} b_{i} \).
The spaces you type in a formula are
ignored.  Remember that a letter like
       $x$                   % $ ... $  and  \( ... \)  are equivalent
is a formula when it denotes a mathematical
symbol, and it should be typed as one.

\section{Displayed Text}

Text is displayed by indenting it from the left
margin.  Quotations are commonly displayed.  There
are short quotations
\begin{quote}
   This is a short quotation.  It consists of a
   single paragraph of text.  See how it is formatted.
\end{quote}
and longer ones.
\begin{quotation}
   This is a longer quotation.  It consists of two
   paragraphs of text, neither of which are
   particularly interesting.

   This is the second paragraph of the quotation.  It
   is just as dull as the first paragraph.
\end{quotation}
Another frequently-displayed structure is a list.
The following is an example of an \emph{itemized}
list.
\begin{itemize}
   \item This is the first item of an itemized list.
         Each item in the list is marked with a ``tick''.
         You don't have to worry about what kind of tick
         mark is used.

   \item This is the second item of the list.  It
         contains another list nested inside it.  The inner
         list is an \emph{enumerated} list.
         \begin{enumerate}
            \item This is the first item of an enumerated
                  list that is nested within the itemized list.

            \item This is the second item of the inner list.
                  \LaTeX\ allows you to nest lists deeper than
                  you really should.
         \end{enumerate}
         This is the rest of the second item of the outer
         list.  It is no more interesting than any other
         part of the item.
   \item This is the third item of the list.
\end{itemize}
You can even display poetry.
\begin{verse}
   There is an environment
    for verse \\             % The \\ command separates lines
   Whose features some poets % within a stanza.
   will curse.

                             % One or more blank lines separate stanzas.

   For instead of making\\
   Them do \emph{all} line breaking, \\
   It allows them to put too many words on a line when they'd rather be
   forced to be terse.
\end{verse}

Mathematical formulas may also be displayed.  A
displayed formula
is
one-line long; multiline
formulas require special formatting instructions.
   \[  \ip{\Gamma}{\psi'} = x'' + y^{2} + z_{i}^{n}\]
Don't start a paragraph with a displayed equation,
nor make one a paragraph by itself.

\end{document}               % End of document.

\section{Bench Huge Section 1109}
% synthetic bench section 1109
% This is a sample LaTeX input file.  (Version of 12 August 2004.)
%
% A '%' character causes TeX to ignore all remaining text on the line,
% and is used for comments like this one.

\documentclass{article}      % Specifies the document class

                             % The preamble begins here.
\title{An Example Document}  % Declares the document's title.
\author{Leslie Lamport}      % Declares the author's name.
\date{January 21, 1994}      % Deleting this command produces today's date.

\newcommand{\ip}[2]{(#1, #2)}
                             % Defines \ip{arg1}{arg2} to mean
                             % (arg1, arg2).

%\newcommand{\ip}[2]{\langle #1 | #2\rangle}
                             % This is an alternative definition of
                             % \ip that is commented out.

\begin{document}             % End of preamble and beginning of text.

\maketitle                   % Produces the title.

This is an example input file.  Comparing it with
the output it generates can show you how to
produce a simple document of your own.

\section{Ordinary Text}      % Produces section heading.  Lower-level
                             % sections are begun with similar
                             % \subsection and \subsubsection commands.

The ends  of words and sentences are marked
  by   spaces. It  doesn't matter how many
spaces    you type; one is as good as 100.  The
end of   a line counts as a space.

One   or more   blank lines denote the  end
of  a paragraph.

Since any number of consecutive spaces are treated
like a single one, the formatting of the input
file makes no difference to
      \LaTeX,                % The \LaTeX command generates the LaTeX logo.
but it makes a difference to you.  When you use
\LaTeX, making your input file as easy to read
as possible will be a great help as you write
your document and when you change it.  This sample
file shows how you can add comments to your own input
file.

Because printing is different from typewriting,
there are a number of things that you have to do
differently when preparing an input file than if
you were just typing the document directly.
Quotation marks like
       ``this''
have to be handled specially, as do quotes within
quotes:
       ``\,`this'            % \, separates the double and single quote.
        is what I just
        wrote, not  `that'\,''.

Dashes come in three sizes: an
       intra-word
dash, a medium dash for number ranges like
       1--2,
and a punctuation
       dash---like
this.

A sentence-ending space should be larger than the
space between words within a sentence.  You
sometimes have to type special commands in
conjunction with punctuation characters to get
this right, as in the following sentence.
       Gnats, gnus, etc.\ all  % `\ ' makes an inter-word space.
       begin with G\@.         % \@ marks end-of-sentence punctuation.
You should check the spaces after periods when
reading your output to make sure you haven't
forgotten any special cases.  Generating an
ellipsis
       \ldots\               % `\ ' is needed after `\ldots' because TeX
                             % ignores spaces after command names like \ldots
                             % made from \ + letters.
                             %
                             % Note how a `%' character causes TeX to ignore
                             % the end of the input line, so these blank lines
                             % do not start a new paragraph.
                             %
with the right spacing around the periods requires
a special command.

\LaTeX\ interprets some common characters as
commands, so you must type special commands to
generate them.  These characters include the
following:
       \$ \& \% \# \{ and \}.

In printing, text is usually emphasized with an
       \emph{italic}
type style.

\begin{em}
   A long segment of text can also be emphasized
   in this way.  Text within such a segment can be
   given \emph{additional} emphasis.
\end{em}

It is sometimes necessary to prevent \LaTeX\ from
breaking a line where it might otherwise do so.
This may be at a space, as between the ``Mr.''\ and
``Jones'' in
       ``Mr.~Jones'',        % ~ produces an unbreakable interword space.
or within a word---especially when the word is a
symbol like
       \mbox{\emph{itemnum}}
that makes little sense when hyphenated across
lines.

Footnotes\footnote{This is an example of a footnote.}
pose no problem.

\LaTeX\ is good at typesetting mathematical formulas
like
       \( x-3y + z = 7 \)
or
       \( a_{1} > x^{2n} + y^{2n} > x' \)
or
       \( \ip{A}{B} = \sum_{i} a_{i} b_{i} \).
The spaces you type in a formula are
ignored.  Remember that a letter like
       $x$                   % $ ... $  and  \( ... \)  are equivalent
is a formula when it denotes a mathematical
symbol, and it should be typed as one.

\section{Displayed Text}

Text is displayed by indenting it from the left
margin.  Quotations are commonly displayed.  There
are short quotations
\begin{quote}
   This is a short quotation.  It consists of a
   single paragraph of text.  See how it is formatted.
\end{quote}
and longer ones.
\begin{quotation}
   This is a longer quotation.  It consists of two
   paragraphs of text, neither of which are
   particularly interesting.

   This is the second paragraph of the quotation.  It
   is just as dull as the first paragraph.
\end{quotation}
Another frequently-displayed structure is a list.
The following is an example of an \emph{itemized}
list.
\begin{itemize}
   \item This is the first item of an itemized list.
         Each item in the list is marked with a ``tick''.
         You don't have to worry about what kind of tick
         mark is used.

   \item This is the second item of the list.  It
         contains another list nested inside it.  The inner
         list is an \emph{enumerated} list.
         \begin{enumerate}
            \item This is the first item of an enumerated
                  list that is nested within the itemized list.

            \item This is the second item of the inner list.
                  \LaTeX\ allows you to nest lists deeper than
                  you really should.
         \end{enumerate}
         This is the rest of the second item of the outer
         list.  It is no more interesting than any other
         part of the item.
   \item This is the third item of the list.
\end{itemize}
You can even display poetry.
\begin{verse}
   There is an environment
    for verse \\             % The \\ command separates lines
   Whose features some poets % within a stanza.
   will curse.

                             % One or more blank lines separate stanzas.

   For instead of making\\
   Them do \emph{all} line breaking, \\
   It allows them to put too many words on a line when they'd rather be
   forced to be terse.
\end{verse}

Mathematical formulas may also be displayed.  A
displayed formula
is
one-line long; multiline
formulas require special formatting instructions.
   \[  \ip{\Gamma}{\psi'} = x'' + y^{2} + z_{i}^{n}\]
Don't start a paragraph with a displayed equation,
nor make one a paragraph by itself.

\end{document}               % End of document.

\section{Bench Huge Section 1110}
% synthetic bench section 1110
% This is a sample LaTeX input file.  (Version of 12 August 2004.)
%
% A '%' character causes TeX to ignore all remaining text on the line,
% and is used for comments like this one.

\documentclass{article}      % Specifies the document class

                             % The preamble begins here.
\title{An Example Document}  % Declares the document's title.
\author{Leslie Lamport}      % Declares the author's name.
\date{January 21, 1994}      % Deleting this command produces today's date.

\newcommand{\ip}[2]{(#1, #2)}
                             % Defines \ip{arg1}{arg2} to mean
                             % (arg1, arg2).

%\newcommand{\ip}[2]{\langle #1 | #2\rangle}
                             % This is an alternative definition of
                             % \ip that is commented out.

\begin{document}             % End of preamble and beginning of text.

\maketitle                   % Produces the title.

This is an example input file.  Comparing it with
the output it generates can show you how to
produce a simple document of your own.

\section{Ordinary Text}      % Produces section heading.  Lower-level
                             % sections are begun with similar
                             % \subsection and \subsubsection commands.

The ends  of words and sentences are marked
  by   spaces. It  doesn't matter how many
spaces    you type; one is as good as 100.  The
end of   a line counts as a space.

One   or more   blank lines denote the  end
of  a paragraph.

Since any number of consecutive spaces are treated
like a single one, the formatting of the input
file makes no difference to
      \LaTeX,                % The \LaTeX command generates the LaTeX logo.
but it makes a difference to you.  When you use
\LaTeX, making your input file as easy to read
as possible will be a great help as you write
your document and when you change it.  This sample
file shows how you can add comments to your own input
file.

Because printing is different from typewriting,
there are a number of things that you have to do
differently when preparing an input file than if
you were just typing the document directly.
Quotation marks like
       ``this''
have to be handled specially, as do quotes within
quotes:
       ``\,`this'            % \, separates the double and single quote.
        is what I just
        wrote, not  `that'\,''.

Dashes come in three sizes: an
       intra-word
dash, a medium dash for number ranges like
       1--2,
and a punctuation
       dash---like
this.

A sentence-ending space should be larger than the
space between words within a sentence.  You
sometimes have to type special commands in
conjunction with punctuation characters to get
this right, as in the following sentence.
       Gnats, gnus, etc.\ all  % `\ ' makes an inter-word space.
       begin with G\@.         % \@ marks end-of-sentence punctuation.
You should check the spaces after periods when
reading your output to make sure you haven't
forgotten any special cases.  Generating an
ellipsis
       \ldots\               % `\ ' is needed after `\ldots' because TeX
                             % ignores spaces after command names like \ldots
                             % made from \ + letters.
                             %
                             % Note how a `%' character causes TeX to ignore
                             % the end of the input line, so these blank lines
                             % do not start a new paragraph.
                             %
with the right spacing around the periods requires
a special command.

\LaTeX\ interprets some common characters as
commands, so you must type special commands to
generate them.  These characters include the
following:
       \$ \& \% \# \{ and \}.

In printing, text is usually emphasized with an
       \emph{italic}
type style.

\begin{em}
   A long segment of text can also be emphasized
   in this way.  Text within such a segment can be
   given \emph{additional} emphasis.
\end{em}

It is sometimes necessary to prevent \LaTeX\ from
breaking a line where it might otherwise do so.
This may be at a space, as between the ``Mr.''\ and
``Jones'' in
       ``Mr.~Jones'',        % ~ produces an unbreakable interword space.
or within a word---especially when the word is a
symbol like
       \mbox{\emph{itemnum}}
that makes little sense when hyphenated across
lines.

Footnotes\footnote{This is an example of a footnote.}
pose no problem.

\LaTeX\ is good at typesetting mathematical formulas
like
       \( x-3y + z = 7 \)
or
       \( a_{1} > x^{2n} + y^{2n} > x' \)
or
       \( \ip{A}{B} = \sum_{i} a_{i} b_{i} \).
The spaces you type in a formula are
ignored.  Remember that a letter like
       $x$                   % $ ... $  and  \( ... \)  are equivalent
is a formula when it denotes a mathematical
symbol, and it should be typed as one.

\section{Displayed Text}

Text is displayed by indenting it from the left
margin.  Quotations are commonly displayed.  There
are short quotations
\begin{quote}
   This is a short quotation.  It consists of a
   single paragraph of text.  See how it is formatted.
\end{quote}
and longer ones.
\begin{quotation}
   This is a longer quotation.  It consists of two
   paragraphs of text, neither of which are
   particularly interesting.

   This is the second paragraph of the quotation.  It
   is just as dull as the first paragraph.
\end{quotation}
Another frequently-displayed structure is a list.
The following is an example of an \emph{itemized}
list.
\begin{itemize}
   \item This is the first item of an itemized list.
         Each item in the list is marked with a ``tick''.
         You don't have to worry about what kind of tick
         mark is used.

   \item This is the second item of the list.  It
         contains another list nested inside it.  The inner
         list is an \emph{enumerated} list.
         \begin{enumerate}
            \item This is the first item of an enumerated
                  list that is nested within the itemized list.

            \item This is the second item of the inner list.
                  \LaTeX\ allows you to nest lists deeper than
                  you really should.
         \end{enumerate}
         This is the rest of the second item of the outer
         list.  It is no more interesting than any other
         part of the item.
   \item This is the third item of the list.
\end{itemize}
You can even display poetry.
\begin{verse}
   There is an environment
    for verse \\             % The \\ command separates lines
   Whose features some poets % within a stanza.
   will curse.

                             % One or more blank lines separate stanzas.

   For instead of making\\
   Them do \emph{all} line breaking, \\
   It allows them to put too many words on a line when they'd rather be
   forced to be terse.
\end{verse}

Mathematical formulas may also be displayed.  A
displayed formula
is
one-line long; multiline
formulas require special formatting instructions.
   \[  \ip{\Gamma}{\psi'} = x'' + y^{2} + z_{i}^{n}\]
Don't start a paragraph with a displayed equation,
nor make one a paragraph by itself.

\end{document}               % End of document.

\section{Bench Huge Section 1111}
% synthetic bench section 1111
% This is a sample LaTeX input file.  (Version of 12 August 2004.)
%
% A '%' character causes TeX to ignore all remaining text on the line,
% and is used for comments like this one.

\documentclass{article}      % Specifies the document class

                             % The preamble begins here.
\title{An Example Document}  % Declares the document's title.
\author{Leslie Lamport}      % Declares the author's name.
\date{January 21, 1994}      % Deleting this command produces today's date.

\newcommand{\ip}[2]{(#1, #2)}
                             % Defines \ip{arg1}{arg2} to mean
                             % (arg1, arg2).

%\newcommand{\ip}[2]{\langle #1 | #2\rangle}
                             % This is an alternative definition of
                             % \ip that is commented out.

\begin{document}             % End of preamble and beginning of text.

\maketitle                   % Produces the title.

This is an example input file.  Comparing it with
the output it generates can show you how to
produce a simple document of your own.

\section{Ordinary Text}      % Produces section heading.  Lower-level
                             % sections are begun with similar
                             % \subsection and \subsubsection commands.

The ends  of words and sentences are marked
  by   spaces. It  doesn't matter how many
spaces    you type; one is as good as 100.  The
end of   a line counts as a space.

One   or more   blank lines denote the  end
of  a paragraph.

Since any number of consecutive spaces are treated
like a single one, the formatting of the input
file makes no difference to
      \LaTeX,                % The \LaTeX command generates the LaTeX logo.
but it makes a difference to you.  When you use
\LaTeX, making your input file as easy to read
as possible will be a great help as you write
your document and when you change it.  This sample
file shows how you can add comments to your own input
file.

Because printing is different from typewriting,
there are a number of things that you have to do
differently when preparing an input file than if
you were just typing the document directly.
Quotation marks like
       ``this''
have to be handled specially, as do quotes within
quotes:
       ``\,`this'            % \, separates the double and single quote.
        is what I just
        wrote, not  `that'\,''.

Dashes come in three sizes: an
       intra-word
dash, a medium dash for number ranges like
       1--2,
and a punctuation
       dash---like
this.

A sentence-ending space should be larger than the
space between words within a sentence.  You
sometimes have to type special commands in
conjunction with punctuation characters to get
this right, as in the following sentence.
       Gnats, gnus, etc.\ all  % `\ ' makes an inter-word space.
       begin with G\@.         % \@ marks end-of-sentence punctuation.
You should check the spaces after periods when
reading your output to make sure you haven't
forgotten any special cases.  Generating an
ellipsis
       \ldots\               % `\ ' is needed after `\ldots' because TeX
                             % ignores spaces after command names like \ldots
                             % made from \ + letters.
                             %
                             % Note how a `%' character causes TeX to ignore
                             % the end of the input line, so these blank lines
                             % do not start a new paragraph.
                             %
with the right spacing around the periods requires
a special command.

\LaTeX\ interprets some common characters as
commands, so you must type special commands to
generate them.  These characters include the
following:
       \$ \& \% \# \{ and \}.

In printing, text is usually emphasized with an
       \emph{italic}
type style.

\begin{em}
   A long segment of text can also be emphasized
   in this way.  Text within such a segment can be
   given \emph{additional} emphasis.
\end{em}

It is sometimes necessary to prevent \LaTeX\ from
breaking a line where it might otherwise do so.
This may be at a space, as between the ``Mr.''\ and
``Jones'' in
       ``Mr.~Jones'',        % ~ produces an unbreakable interword space.
or within a word---especially when the word is a
symbol like
       \mbox{\emph{itemnum}}
that makes little sense when hyphenated across
lines.

Footnotes\footnote{This is an example of a footnote.}
pose no problem.

\LaTeX\ is good at typesetting mathematical formulas
like
       \( x-3y + z = 7 \)
or
       \( a_{1} > x^{2n} + y^{2n} > x' \)
or
       \( \ip{A}{B} = \sum_{i} a_{i} b_{i} \).
The spaces you type in a formula are
ignored.  Remember that a letter like
       $x$                   % $ ... $  and  \( ... \)  are equivalent
is a formula when it denotes a mathematical
symbol, and it should be typed as one.

\section{Displayed Text}

Text is displayed by indenting it from the left
margin.  Quotations are commonly displayed.  There
are short quotations
\begin{quote}
   This is a short quotation.  It consists of a
   single paragraph of text.  See how it is formatted.
\end{quote}
and longer ones.
\begin{quotation}
   This is a longer quotation.  It consists of two
   paragraphs of text, neither of which are
   particularly interesting.

   This is the second paragraph of the quotation.  It
   is just as dull as the first paragraph.
\end{quotation}
Another frequently-displayed structure is a list.
The following is an example of an \emph{itemized}
list.
\begin{itemize}
   \item This is the first item of an itemized list.
         Each item in the list is marked with a ``tick''.
         You don't have to worry about what kind of tick
         mark is used.

   \item This is the second item of the list.  It
         contains another list nested inside it.  The inner
         list is an \emph{enumerated} list.
         \begin{enumerate}
            \item This is the first item of an enumerated
                  list that is nested within the itemized list.

            \item This is the second item of the inner list.
                  \LaTeX\ allows you to nest lists deeper than
                  you really should.
         \end{enumerate}
         This is the rest of the second item of the outer
         list.  It is no more interesting than any other
         part of the item.
   \item This is the third item of the list.
\end{itemize}
You can even display poetry.
\begin{verse}
   There is an environment
    for verse \\             % The \\ command separates lines
   Whose features some poets % within a stanza.
   will curse.

                             % One or more blank lines separate stanzas.

   For instead of making\\
   Them do \emph{all} line breaking, \\
   It allows them to put too many words on a line when they'd rather be
   forced to be terse.
\end{verse}

Mathematical formulas may also be displayed.  A
displayed formula
is
one-line long; multiline
formulas require special formatting instructions.
   \[  \ip{\Gamma}{\psi'} = x'' + y^{2} + z_{i}^{n}\]
Don't start a paragraph with a displayed equation,
nor make one a paragraph by itself.

\end{document}               % End of document.

\section{Bench Huge Section 1112}
% synthetic bench section 1112
% This is a sample LaTeX input file.  (Version of 12 August 2004.)
%
% A '%' character causes TeX to ignore all remaining text on the line,
% and is used for comments like this one.

\documentclass{article}      % Specifies the document class

                             % The preamble begins here.
\title{An Example Document}  % Declares the document's title.
\author{Leslie Lamport}      % Declares the author's name.
\date{January 21, 1994}      % Deleting this command produces today's date.

\newcommand{\ip}[2]{(#1, #2)}
                             % Defines \ip{arg1}{arg2} to mean
                             % (arg1, arg2).

%\newcommand{\ip}[2]{\langle #1 | #2\rangle}
                             % This is an alternative definition of
                             % \ip that is commented out.

\begin{document}             % End of preamble and beginning of text.

\maketitle                   % Produces the title.

This is an example input file.  Comparing it with
the output it generates can show you how to
produce a simple document of your own.

\section{Ordinary Text}      % Produces section heading.  Lower-level
                             % sections are begun with similar
                             % \subsection and \subsubsection commands.

The ends  of words and sentences are marked
  by   spaces. It  doesn't matter how many
spaces    you type; one is as good as 100.  The
end of   a line counts as a space.

One   or more   blank lines denote the  end
of  a paragraph.

Since any number of consecutive spaces are treated
like a single one, the formatting of the input
file makes no difference to
      \LaTeX,                % The \LaTeX command generates the LaTeX logo.
but it makes a difference to you.  When you use
\LaTeX, making your input file as easy to read
as possible will be a great help as you write
your document and when you change it.  This sample
file shows how you can add comments to your own input
file.

Because printing is different from typewriting,
there are a number of things that you have to do
differently when preparing an input file than if
you were just typing the document directly.
Quotation marks like
       ``this''
have to be handled specially, as do quotes within
quotes:
       ``\,`this'            % \, separates the double and single quote.
        is what I just
        wrote, not  `that'\,''.

Dashes come in three sizes: an
       intra-word
dash, a medium dash for number ranges like
       1--2,
and a punctuation
       dash---like
this.

A sentence-ending space should be larger than the
space between words within a sentence.  You
sometimes have to type special commands in
conjunction with punctuation characters to get
this right, as in the following sentence.
       Gnats, gnus, etc.\ all  % `\ ' makes an inter-word space.
       begin with G\@.         % \@ marks end-of-sentence punctuation.
You should check the spaces after periods when
reading your output to make sure you haven't
forgotten any special cases.  Generating an
ellipsis
       \ldots\               % `\ ' is needed after `\ldots' because TeX
                             % ignores spaces after command names like \ldots
                             % made from \ + letters.
                             %
                             % Note how a `%' character causes TeX to ignore
                             % the end of the input line, so these blank lines
                             % do not start a new paragraph.
                             %
with the right spacing around the periods requires
a special command.

\LaTeX\ interprets some common characters as
commands, so you must type special commands to
generate them.  These characters include the
following:
       \$ \& \% \# \{ and \}.

In printing, text is usually emphasized with an
       \emph{italic}
type style.

\begin{em}
   A long segment of text can also be emphasized
   in this way.  Text within such a segment can be
   given \emph{additional} emphasis.
\end{em}

It is sometimes necessary to prevent \LaTeX\ from
breaking a line where it might otherwise do so.
This may be at a space, as between the ``Mr.''\ and
``Jones'' in
       ``Mr.~Jones'',        % ~ produces an unbreakable interword space.
or within a word---especially when the word is a
symbol like
       \mbox{\emph{itemnum}}
that makes little sense when hyphenated across
lines.

Footnotes\footnote{This is an example of a footnote.}
pose no problem.

\LaTeX\ is good at typesetting mathematical formulas
like
       \( x-3y + z = 7 \)
or
       \( a_{1} > x^{2n} + y^{2n} > x' \)
or
       \( \ip{A}{B} = \sum_{i} a_{i} b_{i} \).
The spaces you type in a formula are
ignored.  Remember that a letter like
       $x$                   % $ ... $  and  \( ... \)  are equivalent
is a formula when it denotes a mathematical
symbol, and it should be typed as one.

\section{Displayed Text}

Text is displayed by indenting it from the left
margin.  Quotations are commonly displayed.  There
are short quotations
\begin{quote}
   This is a short quotation.  It consists of a
   single paragraph of text.  See how it is formatted.
\end{quote}
and longer ones.
\begin{quotation}
   This is a longer quotation.  It consists of two
   paragraphs of text, neither of which are
   particularly interesting.

   This is the second paragraph of the quotation.  It
   is just as dull as the first paragraph.
\end{quotation}
Another frequently-displayed structure is a list.
The following is an example of an \emph{itemized}
list.
\begin{itemize}
   \item This is the first item of an itemized list.
         Each item in the list is marked with a ``tick''.
         You don't have to worry about what kind of tick
         mark is used.

   \item This is the second item of the list.  It
         contains another list nested inside it.  The inner
         list is an \emph{enumerated} list.
         \begin{enumerate}
            \item This is the first item of an enumerated
                  list that is nested within the itemized list.

            \item This is the second item of the inner list.
                  \LaTeX\ allows you to nest lists deeper than
                  you really should.
         \end{enumerate}
         This is the rest of the second item of the outer
         list.  It is no more interesting than any other
         part of the item.
   \item This is the third item of the list.
\end{itemize}
You can even display poetry.
\begin{verse}
   There is an environment
    for verse \\             % The \\ command separates lines
   Whose features some poets % within a stanza.
   will curse.

                             % One or more blank lines separate stanzas.

   For instead of making\\
   Them do \emph{all} line breaking, \\
   It allows them to put too many words on a line when they'd rather be
   forced to be terse.
\end{verse}

Mathematical formulas may also be displayed.  A
displayed formula
is
one-line long; multiline
formulas require special formatting instructions.
   \[  \ip{\Gamma}{\psi'} = x'' + y^{2} + z_{i}^{n}\]
Don't start a paragraph with a displayed equation,
nor make one a paragraph by itself.

\end{document}               % End of document.

\section{Bench Huge Section 1113}
% synthetic bench section 1113
% This is a sample LaTeX input file.  (Version of 12 August 2004.)
%
% A '%' character causes TeX to ignore all remaining text on the line,
% and is used for comments like this one.

\documentclass{article}      % Specifies the document class

                             % The preamble begins here.
\title{An Example Document}  % Declares the document's title.
\author{Leslie Lamport}      % Declares the author's name.
\date{January 21, 1994}      % Deleting this command produces today's date.

\newcommand{\ip}[2]{(#1, #2)}
                             % Defines \ip{arg1}{arg2} to mean
                             % (arg1, arg2).

%\newcommand{\ip}[2]{\langle #1 | #2\rangle}
                             % This is an alternative definition of
                             % \ip that is commented out.

\begin{document}             % End of preamble and beginning of text.

\maketitle                   % Produces the title.

This is an example input file.  Comparing it with
the output it generates can show you how to
produce a simple document of your own.

\section{Ordinary Text}      % Produces section heading.  Lower-level
                             % sections are begun with similar
                             % \subsection and \subsubsection commands.

The ends  of words and sentences are marked
  by   spaces. It  doesn't matter how many
spaces    you type; one is as good as 100.  The
end of   a line counts as a space.

One   or more   blank lines denote the  end
of  a paragraph.

Since any number of consecutive spaces are treated
like a single one, the formatting of the input
file makes no difference to
      \LaTeX,                % The \LaTeX command generates the LaTeX logo.
but it makes a difference to you.  When you use
\LaTeX, making your input file as easy to read
as possible will be a great help as you write
your document and when you change it.  This sample
file shows how you can add comments to your own input
file.

Because printing is different from typewriting,
there are a number of things that you have to do
differently when preparing an input file than if
you were just typing the document directly.
Quotation marks like
       ``this''
have to be handled specially, as do quotes within
quotes:
       ``\,`this'            % \, separates the double and single quote.
        is what I just
        wrote, not  `that'\,''.

Dashes come in three sizes: an
       intra-word
dash, a medium dash for number ranges like
       1--2,
and a punctuation
       dash---like
this.

A sentence-ending space should be larger than the
space between words within a sentence.  You
sometimes have to type special commands in
conjunction with punctuation characters to get
this right, as in the following sentence.
       Gnats, gnus, etc.\ all  % `\ ' makes an inter-word space.
       begin with G\@.         % \@ marks end-of-sentence punctuation.
You should check the spaces after periods when
reading your output to make sure you haven't
forgotten any special cases.  Generating an
ellipsis
       \ldots\               % `\ ' is needed after `\ldots' because TeX
                             % ignores spaces after command names like \ldots
                             % made from \ + letters.
                             %
                             % Note how a `%' character causes TeX to ignore
                             % the end of the input line, so these blank lines
                             % do not start a new paragraph.
                             %
with the right spacing around the periods requires
a special command.

\LaTeX\ interprets some common characters as
commands, so you must type special commands to
generate them.  These characters include the
following:
       \$ \& \% \# \{ and \}.

In printing, text is usually emphasized with an
       \emph{italic}
type style.

\begin{em}
   A long segment of text can also be emphasized
   in this way.  Text within such a segment can be
   given \emph{additional} emphasis.
\end{em}

It is sometimes necessary to prevent \LaTeX\ from
breaking a line where it might otherwise do so.
This may be at a space, as between the ``Mr.''\ and
``Jones'' in
       ``Mr.~Jones'',        % ~ produces an unbreakable interword space.
or within a word---especially when the word is a
symbol like
       \mbox{\emph{itemnum}}
that makes little sense when hyphenated across
lines.

Footnotes\footnote{This is an example of a footnote.}
pose no problem.

\LaTeX\ is good at typesetting mathematical formulas
like
       \( x-3y + z = 7 \)
or
       \( a_{1} > x^{2n} + y^{2n} > x' \)
or
       \( \ip{A}{B} = \sum_{i} a_{i} b_{i} \).
The spaces you type in a formula are
ignored.  Remember that a letter like
       $x$                   % $ ... $  and  \( ... \)  are equivalent
is a formula when it denotes a mathematical
symbol, and it should be typed as one.

\section{Displayed Text}

Text is displayed by indenting it from the left
margin.  Quotations are commonly displayed.  There
are short quotations
\begin{quote}
   This is a short quotation.  It consists of a
   single paragraph of text.  See how it is formatted.
\end{quote}
and longer ones.
\begin{quotation}
   This is a longer quotation.  It consists of two
   paragraphs of text, neither of which are
   particularly interesting.

   This is the second paragraph of the quotation.  It
   is just as dull as the first paragraph.
\end{quotation}
Another frequently-displayed structure is a list.
The following is an example of an \emph{itemized}
list.
\begin{itemize}
   \item This is the first item of an itemized list.
         Each item in the list is marked with a ``tick''.
         You don't have to worry about what kind of tick
         mark is used.

   \item This is the second item of the list.  It
         contains another list nested inside it.  The inner
         list is an \emph{enumerated} list.
         \begin{enumerate}
            \item This is the first item of an enumerated
                  list that is nested within the itemized list.

            \item This is the second item of the inner list.
                  \LaTeX\ allows you to nest lists deeper than
                  you really should.
         \end{enumerate}
         This is the rest of the second item of the outer
         list.  It is no more interesting than any other
         part of the item.
   \item This is the third item of the list.
\end{itemize}
You can even display poetry.
\begin{verse}
   There is an environment
    for verse \\             % The \\ command separates lines
   Whose features some poets % within a stanza.
   will curse.

                             % One or more blank lines separate stanzas.

   For instead of making\\
   Them do \emph{all} line breaking, \\
   It allows them to put too many words on a line when they'd rather be
   forced to be terse.
\end{verse}

Mathematical formulas may also be displayed.  A
displayed formula
is
one-line long; multiline
formulas require special formatting instructions.
   \[  \ip{\Gamma}{\psi'} = x'' + y^{2} + z_{i}^{n}\]
Don't start a paragraph with a displayed equation,
nor make one a paragraph by itself.

\end{document}               % End of document.

\section{Bench Huge Section 1114}
% synthetic bench section 1114
% This is a sample LaTeX input file.  (Version of 12 August 2004.)
%
% A '%' character causes TeX to ignore all remaining text on the line,
% and is used for comments like this one.

\documentclass{article}      % Specifies the document class

                             % The preamble begins here.
\title{An Example Document}  % Declares the document's title.
\author{Leslie Lamport}      % Declares the author's name.
\date{January 21, 1994}      % Deleting this command produces today's date.

\newcommand{\ip}[2]{(#1, #2)}
                             % Defines \ip{arg1}{arg2} to mean
                             % (arg1, arg2).

%\newcommand{\ip}[2]{\langle #1 | #2\rangle}
                             % This is an alternative definition of
                             % \ip that is commented out.

\begin{document}             % End of preamble and beginning of text.

\maketitle                   % Produces the title.

This is an example input file.  Comparing it with
the output it generates can show you how to
produce a simple document of your own.

\section{Ordinary Text}      % Produces section heading.  Lower-level
                             % sections are begun with similar
                             % \subsection and \subsubsection commands.

The ends  of words and sentences are marked
  by   spaces. It  doesn't matter how many
spaces    you type; one is as good as 100.  The
end of   a line counts as a space.

One   or more   blank lines denote the  end
of  a paragraph.

Since any number of consecutive spaces are treated
like a single one, the formatting of the input
file makes no difference to
      \LaTeX,                % The \LaTeX command generates the LaTeX logo.
but it makes a difference to you.  When you use
\LaTeX, making your input file as easy to read
as possible will be a great help as you write
your document and when you change it.  This sample
file shows how you can add comments to your own input
file.

Because printing is different from typewriting,
there are a number of things that you have to do
differently when preparing an input file than if
you were just typing the document directly.
Quotation marks like
       ``this''
have to be handled specially, as do quotes within
quotes:
       ``\,`this'            % \, separates the double and single quote.
        is what I just
        wrote, not  `that'\,''.

Dashes come in three sizes: an
       intra-word
dash, a medium dash for number ranges like
       1--2,
and a punctuation
       dash---like
this.

A sentence-ending space should be larger than the
space between words within a sentence.  You
sometimes have to type special commands in
conjunction with punctuation characters to get
this right, as in the following sentence.
       Gnats, gnus, etc.\ all  % `\ ' makes an inter-word space.
       begin with G\@.         % \@ marks end-of-sentence punctuation.
You should check the spaces after periods when
reading your output to make sure you haven't
forgotten any special cases.  Generating an
ellipsis
       \ldots\               % `\ ' is needed after `\ldots' because TeX
                             % ignores spaces after command names like \ldots
                             % made from \ + letters.
                             %
                             % Note how a `%' character causes TeX to ignore
                             % the end of the input line, so these blank lines
                             % do not start a new paragraph.
                             %
with the right spacing around the periods requires
a special command.

\LaTeX\ interprets some common characters as
commands, so you must type special commands to
generate them.  These characters include the
following:
       \$ \& \% \# \{ and \}.

In printing, text is usually emphasized with an
       \emph{italic}
type style.

\begin{em}
   A long segment of text can also be emphasized
   in this way.  Text within such a segment can be
   given \emph{additional} emphasis.
\end{em}

It is sometimes necessary to prevent \LaTeX\ from
breaking a line where it might otherwise do so.
This may be at a space, as between the ``Mr.''\ and
``Jones'' in
       ``Mr.~Jones'',        % ~ produces an unbreakable interword space.
or within a word---especially when the word is a
symbol like
       \mbox{\emph{itemnum}}
that makes little sense when hyphenated across
lines.

Footnotes\footnote{This is an example of a footnote.}
pose no problem.

\LaTeX\ is good at typesetting mathematical formulas
like
       \( x-3y + z = 7 \)
or
       \( a_{1} > x^{2n} + y^{2n} > x' \)
or
       \( \ip{A}{B} = \sum_{i} a_{i} b_{i} \).
The spaces you type in a formula are
ignored.  Remember that a letter like
       $x$                   % $ ... $  and  \( ... \)  are equivalent
is a formula when it denotes a mathematical
symbol, and it should be typed as one.

\section{Displayed Text}

Text is displayed by indenting it from the left
margin.  Quotations are commonly displayed.  There
are short quotations
\begin{quote}
   This is a short quotation.  It consists of a
   single paragraph of text.  See how it is formatted.
\end{quote}
and longer ones.
\begin{quotation}
   This is a longer quotation.  It consists of two
   paragraphs of text, neither of which are
   particularly interesting.

   This is the second paragraph of the quotation.  It
   is just as dull as the first paragraph.
\end{quotation}
Another frequently-displayed structure is a list.
The following is an example of an \emph{itemized}
list.
\begin{itemize}
   \item This is the first item of an itemized list.
         Each item in the list is marked with a ``tick''.
         You don't have to worry about what kind of tick
         mark is used.

   \item This is the second item of the list.  It
         contains another list nested inside it.  The inner
         list is an \emph{enumerated} list.
         \begin{enumerate}
            \item This is the first item of an enumerated
                  list that is nested within the itemized list.

            \item This is the second item of the inner list.
                  \LaTeX\ allows you to nest lists deeper than
                  you really should.
         \end{enumerate}
         This is the rest of the second item of the outer
         list.  It is no more interesting than any other
         part of the item.
   \item This is the third item of the list.
\end{itemize}
You can even display poetry.
\begin{verse}
   There is an environment
    for verse \\             % The \\ command separates lines
   Whose features some poets % within a stanza.
   will curse.

                             % One or more blank lines separate stanzas.

   For instead of making\\
   Them do \emph{all} line breaking, \\
   It allows them to put too many words on a line when they'd rather be
   forced to be terse.
\end{verse}

Mathematical formulas may also be displayed.  A
displayed formula
is
one-line long; multiline
formulas require special formatting instructions.
   \[  \ip{\Gamma}{\psi'} = x'' + y^{2} + z_{i}^{n}\]
Don't start a paragraph with a displayed equation,
nor make one a paragraph by itself.

\end{document}               % End of document.

\section{Bench Huge Section 1115}
% synthetic bench section 1115
% This is a sample LaTeX input file.  (Version of 12 August 2004.)
%
% A '%' character causes TeX to ignore all remaining text on the line,
% and is used for comments like this one.

\documentclass{article}      % Specifies the document class

                             % The preamble begins here.
\title{An Example Document}  % Declares the document's title.
\author{Leslie Lamport}      % Declares the author's name.
\date{January 21, 1994}      % Deleting this command produces today's date.

\newcommand{\ip}[2]{(#1, #2)}
                             % Defines \ip{arg1}{arg2} to mean
                             % (arg1, arg2).

%\newcommand{\ip}[2]{\langle #1 | #2\rangle}
                             % This is an alternative definition of
                             % \ip that is commented out.

\begin{document}             % End of preamble and beginning of text.

\maketitle                   % Produces the title.

This is an example input file.  Comparing it with
the output it generates can show you how to
produce a simple document of your own.

\section{Ordinary Text}      % Produces section heading.  Lower-level
                             % sections are begun with similar
                             % \subsection and \subsubsection commands.

The ends  of words and sentences are marked
  by   spaces. It  doesn't matter how many
spaces    you type; one is as good as 100.  The
end of   a line counts as a space.

One   or more   blank lines denote the  end
of  a paragraph.

Since any number of consecutive spaces are treated
like a single one, the formatting of the input
file makes no difference to
      \LaTeX,                % The \LaTeX command generates the LaTeX logo.
but it makes a difference to you.  When you use
\LaTeX, making your input file as easy to read
as possible will be a great help as you write
your document and when you change it.  This sample
file shows how you can add comments to your own input
file.

Because printing is different from typewriting,
there are a number of things that you have to do
differently when preparing an input file than if
you were just typing the document directly.
Quotation marks like
       ``this''
have to be handled specially, as do quotes within
quotes:
       ``\,`this'            % \, separates the double and single quote.
        is what I just
        wrote, not  `that'\,''.

Dashes come in three sizes: an
       intra-word
dash, a medium dash for number ranges like
       1--2,
and a punctuation
       dash---like
this.

A sentence-ending space should be larger than the
space between words within a sentence.  You
sometimes have to type special commands in
conjunction with punctuation characters to get
this right, as in the following sentence.
       Gnats, gnus, etc.\ all  % `\ ' makes an inter-word space.
       begin with G\@.         % \@ marks end-of-sentence punctuation.
You should check the spaces after periods when
reading your output to make sure you haven't
forgotten any special cases.  Generating an
ellipsis
       \ldots\               % `\ ' is needed after `\ldots' because TeX
                             % ignores spaces after command names like \ldots
                             % made from \ + letters.
                             %
                             % Note how a `%' character causes TeX to ignore
                             % the end of the input line, so these blank lines
                             % do not start a new paragraph.
                             %
with the right spacing around the periods requires
a special command.

\LaTeX\ interprets some common characters as
commands, so you must type special commands to
generate them.  These characters include the
following:
       \$ \& \% \# \{ and \}.

In printing, text is usually emphasized with an
       \emph{italic}
type style.

\begin{em}
   A long segment of text can also be emphasized
   in this way.  Text within such a segment can be
   given \emph{additional} emphasis.
\end{em}

It is sometimes necessary to prevent \LaTeX\ from
breaking a line where it might otherwise do so.
This may be at a space, as between the ``Mr.''\ and
``Jones'' in
       ``Mr.~Jones'',        % ~ produces an unbreakable interword space.
or within a word---especially when the word is a
symbol like
       \mbox{\emph{itemnum}}
that makes little sense when hyphenated across
lines.

Footnotes\footnote{This is an example of a footnote.}
pose no problem.

\LaTeX\ is good at typesetting mathematical formulas
like
       \( x-3y + z = 7 \)
or
       \( a_{1} > x^{2n} + y^{2n} > x' \)
or
       \( \ip{A}{B} = \sum_{i} a_{i} b_{i} \).
The spaces you type in a formula are
ignored.  Remember that a letter like
       $x$                   % $ ... $  and  \( ... \)  are equivalent
is a formula when it denotes a mathematical
symbol, and it should be typed as one.

\section{Displayed Text}

Text is displayed by indenting it from the left
margin.  Quotations are commonly displayed.  There
are short quotations
\begin{quote}
   This is a short quotation.  It consists of a
   single paragraph of text.  See how it is formatted.
\end{quote}
and longer ones.
\begin{quotation}
   This is a longer quotation.  It consists of two
   paragraphs of text, neither of which are
   particularly interesting.

   This is the second paragraph of the quotation.  It
   is just as dull as the first paragraph.
\end{quotation}
Another frequently-displayed structure is a list.
The following is an example of an \emph{itemized}
list.
\begin{itemize}
   \item This is the first item of an itemized list.
         Each item in the list is marked with a ``tick''.
         You don't have to worry about what kind of tick
         mark is used.

   \item This is the second item of the list.  It
         contains another list nested inside it.  The inner
         list is an \emph{enumerated} list.
         \begin{enumerate}
            \item This is the first item of an enumerated
                  list that is nested within the itemized list.

            \item This is the second item of the inner list.
                  \LaTeX\ allows you to nest lists deeper than
                  you really should.
         \end{enumerate}
         This is the rest of the second item of the outer
         list.  It is no more interesting than any other
         part of the item.
   \item This is the third item of the list.
\end{itemize}
You can even display poetry.
\begin{verse}
   There is an environment
    for verse \\             % The \\ command separates lines
   Whose features some poets % within a stanza.
   will curse.

                             % One or more blank lines separate stanzas.

   For instead of making\\
   Them do \emph{all} line breaking, \\
   It allows them to put too many words on a line when they'd rather be
   forced to be terse.
\end{verse}

Mathematical formulas may also be displayed.  A
displayed formula
is
one-line long; multiline
formulas require special formatting instructions.
   \[  \ip{\Gamma}{\psi'} = x'' + y^{2} + z_{i}^{n}\]
Don't start a paragraph with a displayed equation,
nor make one a paragraph by itself.

\end{document}               % End of document.

\section{Bench Huge Section 1116}
% synthetic bench section 1116
% This is a sample LaTeX input file.  (Version of 12 August 2004.)
%
% A '%' character causes TeX to ignore all remaining text on the line,
% and is used for comments like this one.

\documentclass{article}      % Specifies the document class

                             % The preamble begins here.
\title{An Example Document}  % Declares the document's title.
\author{Leslie Lamport}      % Declares the author's name.
\date{January 21, 1994}      % Deleting this command produces today's date.

\newcommand{\ip}[2]{(#1, #2)}
                             % Defines \ip{arg1}{arg2} to mean
                             % (arg1, arg2).

%\newcommand{\ip}[2]{\langle #1 | #2\rangle}
                             % This is an alternative definition of
                             % \ip that is commented out.

\begin{document}             % End of preamble and beginning of text.

\maketitle                   % Produces the title.

This is an example input file.  Comparing it with
the output it generates can show you how to
produce a simple document of your own.

\section{Ordinary Text}      % Produces section heading.  Lower-level
                             % sections are begun with similar
                             % \subsection and \subsubsection commands.

The ends  of words and sentences are marked
  by   spaces. It  doesn't matter how many
spaces    you type; one is as good as 100.  The
end of   a line counts as a space.

One   or more   blank lines denote the  end
of  a paragraph.

Since any number of consecutive spaces are treated
like a single one, the formatting of the input
file makes no difference to
      \LaTeX,                % The \LaTeX command generates the LaTeX logo.
but it makes a difference to you.  When you use
\LaTeX, making your input file as easy to read
as possible will be a great help as you write
your document and when you change it.  This sample
file shows how you can add comments to your own input
file.

Because printing is different from typewriting,
there are a number of things that you have to do
differently when preparing an input file than if
you were just typing the document directly.
Quotation marks like
       ``this''
have to be handled specially, as do quotes within
quotes:
       ``\,`this'            % \, separates the double and single quote.
        is what I just
        wrote, not  `that'\,''.

Dashes come in three sizes: an
       intra-word
dash, a medium dash for number ranges like
       1--2,
and a punctuation
       dash---like
this.

A sentence-ending space should be larger than the
space between words within a sentence.  You
sometimes have to type special commands in
conjunction with punctuation characters to get
this right, as in the following sentence.
       Gnats, gnus, etc.\ all  % `\ ' makes an inter-word space.
       begin with G\@.         % \@ marks end-of-sentence punctuation.
You should check the spaces after periods when
reading your output to make sure you haven't
forgotten any special cases.  Generating an
ellipsis
       \ldots\               % `\ ' is needed after `\ldots' because TeX
                             % ignores spaces after command names like \ldots
                             % made from \ + letters.
                             %
                             % Note how a `%' character causes TeX to ignore
                             % the end of the input line, so these blank lines
                             % do not start a new paragraph.
                             %
with the right spacing around the periods requires
a special command.

\LaTeX\ interprets some common characters as
commands, so you must type special commands to
generate them.  These characters include the
following:
       \$ \& \% \# \{ and \}.

In printing, text is usually emphasized with an
       \emph{italic}
type style.

\begin{em}
   A long segment of text can also be emphasized
   in this way.  Text within such a segment can be
   given \emph{additional} emphasis.
\end{em}

It is sometimes necessary to prevent \LaTeX\ from
breaking a line where it might otherwise do so.
This may be at a space, as between the ``Mr.''\ and
``Jones'' in
       ``Mr.~Jones'',        % ~ produces an unbreakable interword space.
or within a word---especially when the word is a
symbol like
       \mbox{\emph{itemnum}}
that makes little sense when hyphenated across
lines.

Footnotes\footnote{This is an example of a footnote.}
pose no problem.

\LaTeX\ is good at typesetting mathematical formulas
like
       \( x-3y + z = 7 \)
or
       \( a_{1} > x^{2n} + y^{2n} > x' \)
or
       \( \ip{A}{B} = \sum_{i} a_{i} b_{i} \).
The spaces you type in a formula are
ignored.  Remember that a letter like
       $x$                   % $ ... $  and  \( ... \)  are equivalent
is a formula when it denotes a mathematical
symbol, and it should be typed as one.

\section{Displayed Text}

Text is displayed by indenting it from the left
margin.  Quotations are commonly displayed.  There
are short quotations
\begin{quote}
   This is a short quotation.  It consists of a
   single paragraph of text.  See how it is formatted.
\end{quote}
and longer ones.
\begin{quotation}
   This is a longer quotation.  It consists of two
   paragraphs of text, neither of which are
   particularly interesting.

   This is the second paragraph of the quotation.  It
   is just as dull as the first paragraph.
\end{quotation}
Another frequently-displayed structure is a list.
The following is an example of an \emph{itemized}
list.
\begin{itemize}
   \item This is the first item of an itemized list.
         Each item in the list is marked with a ``tick''.
         You don't have to worry about what kind of tick
         mark is used.

   \item This is the second item of the list.  It
         contains another list nested inside it.  The inner
         list is an \emph{enumerated} list.
         \begin{enumerate}
            \item This is the first item of an enumerated
                  list that is nested within the itemized list.

            \item This is the second item of the inner list.
                  \LaTeX\ allows you to nest lists deeper than
                  you really should.
         \end{enumerate}
         This is the rest of the second item of the outer
         list.  It is no more interesting than any other
         part of the item.
   \item This is the third item of the list.
\end{itemize}
You can even display poetry.
\begin{verse}
   There is an environment
    for verse \\             % The \\ command separates lines
   Whose features some poets % within a stanza.
   will curse.

                             % One or more blank lines separate stanzas.

   For instead of making\\
   Them do \emph{all} line breaking, \\
   It allows them to put too many words on a line when they'd rather be
   forced to be terse.
\end{verse}

Mathematical formulas may also be displayed.  A
displayed formula
is
one-line long; multiline
formulas require special formatting instructions.
   \[  \ip{\Gamma}{\psi'} = x'' + y^{2} + z_{i}^{n}\]
Don't start a paragraph with a displayed equation,
nor make one a paragraph by itself.

\end{document}               % End of document.

\section{Bench Huge Section 1117}
% synthetic bench section 1117
% This is a sample LaTeX input file.  (Version of 12 August 2004.)
%
% A '%' character causes TeX to ignore all remaining text on the line,
% and is used for comments like this one.

\documentclass{article}      % Specifies the document class

                             % The preamble begins here.
\title{An Example Document}  % Declares the document's title.
\author{Leslie Lamport}      % Declares the author's name.
\date{January 21, 1994}      % Deleting this command produces today's date.

\newcommand{\ip}[2]{(#1, #2)}
                             % Defines \ip{arg1}{arg2} to mean
                             % (arg1, arg2).

%\newcommand{\ip}[2]{\langle #1 | #2\rangle}
                             % This is an alternative definition of
                             % \ip that is commented out.

\begin{document}             % End of preamble and beginning of text.

\maketitle                   % Produces the title.

This is an example input file.  Comparing it with
the output it generates can show you how to
produce a simple document of your own.

\section{Ordinary Text}      % Produces section heading.  Lower-level
                             % sections are begun with similar
                             % \subsection and \subsubsection commands.

The ends  of words and sentences are marked
  by   spaces. It  doesn't matter how many
spaces    you type; one is as good as 100.  The
end of   a line counts as a space.

One   or more   blank lines denote the  end
of  a paragraph.

Since any number of consecutive spaces are treated
like a single one, the formatting of the input
file makes no difference to
      \LaTeX,                % The \LaTeX command generates the LaTeX logo.
but it makes a difference to you.  When you use
\LaTeX, making your input file as easy to read
as possible will be a great help as you write
your document and when you change it.  This sample
file shows how you can add comments to your own input
file.

Because printing is different from typewriting,
there are a number of things that you have to do
differently when preparing an input file than if
you were just typing the document directly.
Quotation marks like
       ``this''
have to be handled specially, as do quotes within
quotes:
       ``\,`this'            % \, separates the double and single quote.
        is what I just
        wrote, not  `that'\,''.

Dashes come in three sizes: an
       intra-word
dash, a medium dash for number ranges like
       1--2,
and a punctuation
       dash---like
this.

A sentence-ending space should be larger than the
space between words within a sentence.  You
sometimes have to type special commands in
conjunction with punctuation characters to get
this right, as in the following sentence.
       Gnats, gnus, etc.\ all  % `\ ' makes an inter-word space.
       begin with G\@.         % \@ marks end-of-sentence punctuation.
You should check the spaces after periods when
reading your output to make sure you haven't
forgotten any special cases.  Generating an
ellipsis
       \ldots\               % `\ ' is needed after `\ldots' because TeX
                             % ignores spaces after command names like \ldots
                             % made from \ + letters.
                             %
                             % Note how a `%' character causes TeX to ignore
                             % the end of the input line, so these blank lines
                             % do not start a new paragraph.
                             %
with the right spacing around the periods requires
a special command.

\LaTeX\ interprets some common characters as
commands, so you must type special commands to
generate them.  These characters include the
following:
       \$ \& \% \# \{ and \}.

In printing, text is usually emphasized with an
       \emph{italic}
type style.

\begin{em}
   A long segment of text can also be emphasized
   in this way.  Text within such a segment can be
   given \emph{additional} emphasis.
\end{em}

It is sometimes necessary to prevent \LaTeX\ from
breaking a line where it might otherwise do so.
This may be at a space, as between the ``Mr.''\ and
``Jones'' in
       ``Mr.~Jones'',        % ~ produces an unbreakable interword space.
or within a word---especially when the word is a
symbol like
       \mbox{\emph{itemnum}}
that makes little sense when hyphenated across
lines.

Footnotes\footnote{This is an example of a footnote.}
pose no problem.

\LaTeX\ is good at typesetting mathematical formulas
like
       \( x-3y + z = 7 \)
or
       \( a_{1} > x^{2n} + y^{2n} > x' \)
or
       \( \ip{A}{B} = \sum_{i} a_{i} b_{i} \).
The spaces you type in a formula are
ignored.  Remember that a letter like
       $x$                   % $ ... $  and  \( ... \)  are equivalent
is a formula when it denotes a mathematical
symbol, and it should be typed as one.

\section{Displayed Text}

Text is displayed by indenting it from the left
margin.  Quotations are commonly displayed.  There
are short quotations
\begin{quote}
   This is a short quotation.  It consists of a
   single paragraph of text.  See how it is formatted.
\end{quote}
and longer ones.
\begin{quotation}
   This is a longer quotation.  It consists of two
   paragraphs of text, neither of which are
   particularly interesting.

   This is the second paragraph of the quotation.  It
   is just as dull as the first paragraph.
\end{quotation}
Another frequently-displayed structure is a list.
The following is an example of an \emph{itemized}
list.
\begin{itemize}
   \item This is the first item of an itemized list.
         Each item in the list is marked with a ``tick''.
         You don't have to worry about what kind of tick
         mark is used.

   \item This is the second item of the list.  It
         contains another list nested inside it.  The inner
         list is an \emph{enumerated} list.
         \begin{enumerate}
            \item This is the first item of an enumerated
                  list that is nested within the itemized list.

            \item This is the second item of the inner list.
                  \LaTeX\ allows you to nest lists deeper than
                  you really should.
         \end{enumerate}
         This is the rest of the second item of the outer
         list.  It is no more interesting than any other
         part of the item.
   \item This is the third item of the list.
\end{itemize}
You can even display poetry.
\begin{verse}
   There is an environment
    for verse \\             % The \\ command separates lines
   Whose features some poets % within a stanza.
   will curse.

                             % One or more blank lines separate stanzas.

   For instead of making\\
   Them do \emph{all} line breaking, \\
   It allows them to put too many words on a line when they'd rather be
   forced to be terse.
\end{verse}

Mathematical formulas may also be displayed.  A
displayed formula
is
one-line long; multiline
formulas require special formatting instructions.
   \[  \ip{\Gamma}{\psi'} = x'' + y^{2} + z_{i}^{n}\]
Don't start a paragraph with a displayed equation,
nor make one a paragraph by itself.

\end{document}               % End of document.

\section{Bench Huge Section 1118}
% synthetic bench section 1118
% This is a sample LaTeX input file.  (Version of 12 August 2004.)
%
% A '%' character causes TeX to ignore all remaining text on the line,
% and is used for comments like this one.

\documentclass{article}      % Specifies the document class

                             % The preamble begins here.
\title{An Example Document}  % Declares the document's title.
\author{Leslie Lamport}      % Declares the author's name.
\date{January 21, 1994}      % Deleting this command produces today's date.

\newcommand{\ip}[2]{(#1, #2)}
                             % Defines \ip{arg1}{arg2} to mean
                             % (arg1, arg2).

%\newcommand{\ip}[2]{\langle #1 | #2\rangle}
                             % This is an alternative definition of
                             % \ip that is commented out.

\begin{document}             % End of preamble and beginning of text.

\maketitle                   % Produces the title.

This is an example input file.  Comparing it with
the output it generates can show you how to
produce a simple document of your own.

\section{Ordinary Text}      % Produces section heading.  Lower-level
                             % sections are begun with similar
                             % \subsection and \subsubsection commands.

The ends  of words and sentences are marked
  by   spaces. It  doesn't matter how many
spaces    you type; one is as good as 100.  The
end of   a line counts as a space.

One   or more   blank lines denote the  end
of  a paragraph.

Since any number of consecutive spaces are treated
like a single one, the formatting of the input
file makes no difference to
      \LaTeX,                % The \LaTeX command generates the LaTeX logo.
but it makes a difference to you.  When you use
\LaTeX, making your input file as easy to read
as possible will be a great help as you write
your document and when you change it.  This sample
file shows how you can add comments to your own input
file.

Because printing is different from typewriting,
there are a number of things that you have to do
differently when preparing an input file than if
you were just typing the document directly.
Quotation marks like
       ``this''
have to be handled specially, as do quotes within
quotes:
       ``\,`this'            % \, separates the double and single quote.
        is what I just
        wrote, not  `that'\,''.

Dashes come in three sizes: an
       intra-word
dash, a medium dash for number ranges like
       1--2,
and a punctuation
       dash---like
this.

A sentence-ending space should be larger than the
space between words within a sentence.  You
sometimes have to type special commands in
conjunction with punctuation characters to get
this right, as in the following sentence.
       Gnats, gnus, etc.\ all  % `\ ' makes an inter-word space.
       begin with G\@.         % \@ marks end-of-sentence punctuation.
You should check the spaces after periods when
reading your output to make sure you haven't
forgotten any special cases.  Generating an
ellipsis
       \ldots\               % `\ ' is needed after `\ldots' because TeX
                             % ignores spaces after command names like \ldots
                             % made from \ + letters.
                             %
                             % Note how a `%' character causes TeX to ignore
                             % the end of the input line, so these blank lines
                             % do not start a new paragraph.
                             %
with the right spacing around the periods requires
a special command.

\LaTeX\ interprets some common characters as
commands, so you must type special commands to
generate them.  These characters include the
following:
       \$ \& \% \# \{ and \}.

In printing, text is usually emphasized with an
       \emph{italic}
type style.

\begin{em}
   A long segment of text can also be emphasized
   in this way.  Text within such a segment can be
   given \emph{additional} emphasis.
\end{em}

It is sometimes necessary to prevent \LaTeX\ from
breaking a line where it might otherwise do so.
This may be at a space, as between the ``Mr.''\ and
``Jones'' in
       ``Mr.~Jones'',        % ~ produces an unbreakable interword space.
or within a word---especially when the word is a
symbol like
       \mbox{\emph{itemnum}}
that makes little sense when hyphenated across
lines.

Footnotes\footnote{This is an example of a footnote.}
pose no problem.

\LaTeX\ is good at typesetting mathematical formulas
like
       \( x-3y + z = 7 \)
or
       \( a_{1} > x^{2n} + y^{2n} > x' \)
or
       \( \ip{A}{B} = \sum_{i} a_{i} b_{i} \).
The spaces you type in a formula are
ignored.  Remember that a letter like
       $x$                   % $ ... $  and  \( ... \)  are equivalent
is a formula when it denotes a mathematical
symbol, and it should be typed as one.

\section{Displayed Text}

Text is displayed by indenting it from the left
margin.  Quotations are commonly displayed.  There
are short quotations
\begin{quote}
   This is a short quotation.  It consists of a
   single paragraph of text.  See how it is formatted.
\end{quote}
and longer ones.
\begin{quotation}
   This is a longer quotation.  It consists of two
   paragraphs of text, neither of which are
   particularly interesting.

   This is the second paragraph of the quotation.  It
   is just as dull as the first paragraph.
\end{quotation}
Another frequently-displayed structure is a list.
The following is an example of an \emph{itemized}
list.
\begin{itemize}
   \item This is the first item of an itemized list.
         Each item in the list is marked with a ``tick''.
         You don't have to worry about what kind of tick
         mark is used.

   \item This is the second item of the list.  It
         contains another list nested inside it.  The inner
         list is an \emph{enumerated} list.
         \begin{enumerate}
            \item This is the first item of an enumerated
                  list that is nested within the itemized list.

            \item This is the second item of the inner list.
                  \LaTeX\ allows you to nest lists deeper than
                  you really should.
         \end{enumerate}
         This is the rest of the second item of the outer
         list.  It is no more interesting than any other
         part of the item.
   \item This is the third item of the list.
\end{itemize}
You can even display poetry.
\begin{verse}
   There is an environment
    for verse \\             % The \\ command separates lines
   Whose features some poets % within a stanza.
   will curse.

                             % One or more blank lines separate stanzas.

   For instead of making\\
   Them do \emph{all} line breaking, \\
   It allows them to put too many words on a line when they'd rather be
   forced to be terse.
\end{verse}

Mathematical formulas may also be displayed.  A
displayed formula
is
one-line long; multiline
formulas require special formatting instructions.
   \[  \ip{\Gamma}{\psi'} = x'' + y^{2} + z_{i}^{n}\]
Don't start a paragraph with a displayed equation,
nor make one a paragraph by itself.

\end{document}               % End of document.

\section{Bench Huge Section 1119}
% synthetic bench section 1119
% This is a sample LaTeX input file.  (Version of 12 August 2004.)
%
% A '%' character causes TeX to ignore all remaining text on the line,
% and is used for comments like this one.

\documentclass{article}      % Specifies the document class

                             % The preamble begins here.
\title{An Example Document}  % Declares the document's title.
\author{Leslie Lamport}      % Declares the author's name.
\date{January 21, 1994}      % Deleting this command produces today's date.

\newcommand{\ip}[2]{(#1, #2)}
                             % Defines \ip{arg1}{arg2} to mean
                             % (arg1, arg2).

%\newcommand{\ip}[2]{\langle #1 | #2\rangle}
                             % This is an alternative definition of
                             % \ip that is commented out.

\begin{document}             % End of preamble and beginning of text.

\maketitle                   % Produces the title.

This is an example input file.  Comparing it with
the output it generates can show you how to
produce a simple document of your own.

\section{Ordinary Text}      % Produces section heading.  Lower-level
                             % sections are begun with similar
                             % \subsection and \subsubsection commands.

The ends  of words and sentences are marked
  by   spaces. It  doesn't matter how many
spaces    you type; one is as good as 100.  The
end of   a line counts as a space.

One   or more   blank lines denote the  end
of  a paragraph.

Since any number of consecutive spaces are treated
like a single one, the formatting of the input
file makes no difference to
      \LaTeX,                % The \LaTeX command generates the LaTeX logo.
but it makes a difference to you.  When you use
\LaTeX, making your input file as easy to read
as possible will be a great help as you write
your document and when you change it.  This sample
file shows how you can add comments to your own input
file.

Because printing is different from typewriting,
there are a number of things that you have to do
differently when preparing an input file than if
you were just typing the document directly.
Quotation marks like
       ``this''
have to be handled specially, as do quotes within
quotes:
       ``\,`this'            % \, separates the double and single quote.
        is what I just
        wrote, not  `that'\,''.

Dashes come in three sizes: an
       intra-word
dash, a medium dash for number ranges like
       1--2,
and a punctuation
       dash---like
this.

A sentence-ending space should be larger than the
space between words within a sentence.  You
sometimes have to type special commands in
conjunction with punctuation characters to get
this right, as in the following sentence.
       Gnats, gnus, etc.\ all  % `\ ' makes an inter-word space.
       begin with G\@.         % \@ marks end-of-sentence punctuation.
You should check the spaces after periods when
reading your output to make sure you haven't
forgotten any special cases.  Generating an
ellipsis
       \ldots\               % `\ ' is needed after `\ldots' because TeX
                             % ignores spaces after command names like \ldots
                             % made from \ + letters.
                             %
                             % Note how a `%' character causes TeX to ignore
                             % the end of the input line, so these blank lines
                             % do not start a new paragraph.
                             %
with the right spacing around the periods requires
a special command.

\LaTeX\ interprets some common characters as
commands, so you must type special commands to
generate them.  These characters include the
following:
       \$ \& \% \# \{ and \}.

In printing, text is usually emphasized with an
       \emph{italic}
type style.

\begin{em}
   A long segment of text can also be emphasized
   in this way.  Text within such a segment can be
   given \emph{additional} emphasis.
\end{em}

It is sometimes necessary to prevent \LaTeX\ from
breaking a line where it might otherwise do so.
This may be at a space, as between the ``Mr.''\ and
``Jones'' in
       ``Mr.~Jones'',        % ~ produces an unbreakable interword space.
or within a word---especially when the word is a
symbol like
       \mbox{\emph{itemnum}}
that makes little sense when hyphenated across
lines.

Footnotes\footnote{This is an example of a footnote.}
pose no problem.

\LaTeX\ is good at typesetting mathematical formulas
like
       \( x-3y + z = 7 \)
or
       \( a_{1} > x^{2n} + y^{2n} > x' \)
or
       \( \ip{A}{B} = \sum_{i} a_{i} b_{i} \).
The spaces you type in a formula are
ignored.  Remember that a letter like
       $x$                   % $ ... $  and  \( ... \)  are equivalent
is a formula when it denotes a mathematical
symbol, and it should be typed as one.

\section{Displayed Text}

Text is displayed by indenting it from the left
margin.  Quotations are commonly displayed.  There
are short quotations
\begin{quote}
   This is a short quotation.  It consists of a
   single paragraph of text.  See how it is formatted.
\end{quote}
and longer ones.
\begin{quotation}
   This is a longer quotation.  It consists of two
   paragraphs of text, neither of which are
   particularly interesting.

   This is the second paragraph of the quotation.  It
   is just as dull as the first paragraph.
\end{quotation}
Another frequently-displayed structure is a list.
The following is an example of an \emph{itemized}
list.
\begin{itemize}
   \item This is the first item of an itemized list.
         Each item in the list is marked with a ``tick''.
         You don't have to worry about what kind of tick
         mark is used.

   \item This is the second item of the list.  It
         contains another list nested inside it.  The inner
         list is an \emph{enumerated} list.
         \begin{enumerate}
            \item This is the first item of an enumerated
                  list that is nested within the itemized list.

            \item This is the second item of the inner list.
                  \LaTeX\ allows you to nest lists deeper than
                  you really should.
         \end{enumerate}
         This is the rest of the second item of the outer
         list.  It is no more interesting than any other
         part of the item.
   \item This is the third item of the list.
\end{itemize}
You can even display poetry.
\begin{verse}
   There is an environment
    for verse \\             % The \\ command separates lines
   Whose features some poets % within a stanza.
   will curse.

                             % One or more blank lines separate stanzas.

   For instead of making\\
   Them do \emph{all} line breaking, \\
   It allows them to put too many words on a line when they'd rather be
   forced to be terse.
\end{verse}

Mathematical formulas may also be displayed.  A
displayed formula
is
one-line long; multiline
formulas require special formatting instructions.
   \[  \ip{\Gamma}{\psi'} = x'' + y^{2} + z_{i}^{n}\]
Don't start a paragraph with a displayed equation,
nor make one a paragraph by itself.

\end{document}               % End of document.

\section{Bench Huge Section 1120}
% synthetic bench section 1120
% This is a sample LaTeX input file.  (Version of 12 August 2004.)
%
% A '%' character causes TeX to ignore all remaining text on the line,
% and is used for comments like this one.

\documentclass{article}      % Specifies the document class

                             % The preamble begins here.
\title{An Example Document}  % Declares the document's title.
\author{Leslie Lamport}      % Declares the author's name.
\date{January 21, 1994}      % Deleting this command produces today's date.

\newcommand{\ip}[2]{(#1, #2)}
                             % Defines \ip{arg1}{arg2} to mean
                             % (arg1, arg2).

%\newcommand{\ip}[2]{\langle #1 | #2\rangle}
                             % This is an alternative definition of
                             % \ip that is commented out.

\begin{document}             % End of preamble and beginning of text.

\maketitle                   % Produces the title.

This is an example input file.  Comparing it with
the output it generates can show you how to
produce a simple document of your own.

\section{Ordinary Text}      % Produces section heading.  Lower-level
                             % sections are begun with similar
                             % \subsection and \subsubsection commands.

The ends  of words and sentences are marked
  by   spaces. It  doesn't matter how many
spaces    you type; one is as good as 100.  The
end of   a line counts as a space.

One   or more   blank lines denote the  end
of  a paragraph.

Since any number of consecutive spaces are treated
like a single one, the formatting of the input
file makes no difference to
      \LaTeX,                % The \LaTeX command generates the LaTeX logo.
but it makes a difference to you.  When you use
\LaTeX, making your input file as easy to read
as possible will be a great help as you write
your document and when you change it.  This sample
file shows how you can add comments to your own input
file.

Because printing is different from typewriting,
there are a number of things that you have to do
differently when preparing an input file than if
you were just typing the document directly.
Quotation marks like
       ``this''
have to be handled specially, as do quotes within
quotes:
       ``\,`this'            % \, separates the double and single quote.
        is what I just
        wrote, not  `that'\,''.

Dashes come in three sizes: an
       intra-word
dash, a medium dash for number ranges like
       1--2,
and a punctuation
       dash---like
this.

A sentence-ending space should be larger than the
space between words within a sentence.  You
sometimes have to type special commands in
conjunction with punctuation characters to get
this right, as in the following sentence.
       Gnats, gnus, etc.\ all  % `\ ' makes an inter-word space.
       begin with G\@.         % \@ marks end-of-sentence punctuation.
You should check the spaces after periods when
reading your output to make sure you haven't
forgotten any special cases.  Generating an
ellipsis
       \ldots\               % `\ ' is needed after `\ldots' because TeX
                             % ignores spaces after command names like \ldots
                             % made from \ + letters.
                             %
                             % Note how a `%' character causes TeX to ignore
                             % the end of the input line, so these blank lines
                             % do not start a new paragraph.
                             %
with the right spacing around the periods requires
a special command.

\LaTeX\ interprets some common characters as
commands, so you must type special commands to
generate them.  These characters include the
following:
       \$ \& \% \# \{ and \}.

In printing, text is usually emphasized with an
       \emph{italic}
type style.

\begin{em}
   A long segment of text can also be emphasized
   in this way.  Text within such a segment can be
   given \emph{additional} emphasis.
\end{em}

It is sometimes necessary to prevent \LaTeX\ from
breaking a line where it might otherwise do so.
This may be at a space, as between the ``Mr.''\ and
``Jones'' in
       ``Mr.~Jones'',        % ~ produces an unbreakable interword space.
or within a word---especially when the word is a
symbol like
       \mbox{\emph{itemnum}}
that makes little sense when hyphenated across
lines.

Footnotes\footnote{This is an example of a footnote.}
pose no problem.

\LaTeX\ is good at typesetting mathematical formulas
like
       \( x-3y + z = 7 \)
or
       \( a_{1} > x^{2n} + y^{2n} > x' \)
or
       \( \ip{A}{B} = \sum_{i} a_{i} b_{i} \).
The spaces you type in a formula are
ignored.  Remember that a letter like
       $x$                   % $ ... $  and  \( ... \)  are equivalent
is a formula when it denotes a mathematical
symbol, and it should be typed as one.

\section{Displayed Text}

Text is displayed by indenting it from the left
margin.  Quotations are commonly displayed.  There
are short quotations
\begin{quote}
   This is a short quotation.  It consists of a
   single paragraph of text.  See how it is formatted.
\end{quote}
and longer ones.
\begin{quotation}
   This is a longer quotation.  It consists of two
   paragraphs of text, neither of which are
   particularly interesting.

   This is the second paragraph of the quotation.  It
   is just as dull as the first paragraph.
\end{quotation}
Another frequently-displayed structure is a list.
The following is an example of an \emph{itemized}
list.
\begin{itemize}
   \item This is the first item of an itemized list.
         Each item in the list is marked with a ``tick''.
         You don't have to worry about what kind of tick
         mark is used.

   \item This is the second item of the list.  It
         contains another list nested inside it.  The inner
         list is an \emph{enumerated} list.
         \begin{enumerate}
            \item This is the first item of an enumerated
                  list that is nested within the itemized list.

            \item This is the second item of the inner list.
                  \LaTeX\ allows you to nest lists deeper than
                  you really should.
         \end{enumerate}
         This is the rest of the second item of the outer
         list.  It is no more interesting than any other
         part of the item.
   \item This is the third item of the list.
\end{itemize}
You can even display poetry.
\begin{verse}
   There is an environment
    for verse \\             % The \\ command separates lines
   Whose features some poets % within a stanza.
   will curse.

                             % One or more blank lines separate stanzas.

   For instead of making\\
   Them do \emph{all} line breaking, \\
   It allows them to put too many words on a line when they'd rather be
   forced to be terse.
\end{verse}

Mathematical formulas may also be displayed.  A
displayed formula
is
one-line long; multiline
formulas require special formatting instructions.
   \[  \ip{\Gamma}{\psi'} = x'' + y^{2} + z_{i}^{n}\]
Don't start a paragraph with a displayed equation,
nor make one a paragraph by itself.

\end{document}               % End of document.

\section{Bench Huge Section 1121}
% synthetic bench section 1121
% This is a sample LaTeX input file.  (Version of 12 August 2004.)
%
% A '%' character causes TeX to ignore all remaining text on the line,
% and is used for comments like this one.

\documentclass{article}      % Specifies the document class

                             % The preamble begins here.
\title{An Example Document}  % Declares the document's title.
\author{Leslie Lamport}      % Declares the author's name.
\date{January 21, 1994}      % Deleting this command produces today's date.

\newcommand{\ip}[2]{(#1, #2)}
                             % Defines \ip{arg1}{arg2} to mean
                             % (arg1, arg2).

%\newcommand{\ip}[2]{\langle #1 | #2\rangle}
                             % This is an alternative definition of
                             % \ip that is commented out.

\begin{document}             % End of preamble and beginning of text.

\maketitle                   % Produces the title.

This is an example input file.  Comparing it with
the output it generates can show you how to
produce a simple document of your own.

\section{Ordinary Text}      % Produces section heading.  Lower-level
                             % sections are begun with similar
                             % \subsection and \subsubsection commands.

The ends  of words and sentences are marked
  by   spaces. It  doesn't matter how many
spaces    you type; one is as good as 100.  The
end of   a line counts as a space.

One   or more   blank lines denote the  end
of  a paragraph.

Since any number of consecutive spaces are treated
like a single one, the formatting of the input
file makes no difference to
      \LaTeX,                % The \LaTeX command generates the LaTeX logo.
but it makes a difference to you.  When you use
\LaTeX, making your input file as easy to read
as possible will be a great help as you write
your document and when you change it.  This sample
file shows how you can add comments to your own input
file.

Because printing is different from typewriting,
there are a number of things that you have to do
differently when preparing an input file than if
you were just typing the document directly.
Quotation marks like
       ``this''
have to be handled specially, as do quotes within
quotes:
       ``\,`this'            % \, separates the double and single quote.
        is what I just
        wrote, not  `that'\,''.

Dashes come in three sizes: an
       intra-word
dash, a medium dash for number ranges like
       1--2,
and a punctuation
       dash---like
this.

A sentence-ending space should be larger than the
space between words within a sentence.  You
sometimes have to type special commands in
conjunction with punctuation characters to get
this right, as in the following sentence.
       Gnats, gnus, etc.\ all  % `\ ' makes an inter-word space.
       begin with G\@.         % \@ marks end-of-sentence punctuation.
You should check the spaces after periods when
reading your output to make sure you haven't
forgotten any special cases.  Generating an
ellipsis
       \ldots\               % `\ ' is needed after `\ldots' because TeX
                             % ignores spaces after command names like \ldots
                             % made from \ + letters.
                             %
                             % Note how a `%' character causes TeX to ignore
                             % the end of the input line, so these blank lines
                             % do not start a new paragraph.
                             %
with the right spacing around the periods requires
a special command.

\LaTeX\ interprets some common characters as
commands, so you must type special commands to
generate them.  These characters include the
following:
       \$ \& \% \# \{ and \}.

In printing, text is usually emphasized with an
       \emph{italic}
type style.

\begin{em}
   A long segment of text can also be emphasized
   in this way.  Text within such a segment can be
   given \emph{additional} emphasis.
\end{em}

It is sometimes necessary to prevent \LaTeX\ from
breaking a line where it might otherwise do so.
This may be at a space, as between the ``Mr.''\ and
``Jones'' in
       ``Mr.~Jones'',        % ~ produces an unbreakable interword space.
or within a word---especially when the word is a
symbol like
       \mbox{\emph{itemnum}}
that makes little sense when hyphenated across
lines.

Footnotes\footnote{This is an example of a footnote.}
pose no problem.

\LaTeX\ is good at typesetting mathematical formulas
like
       \( x-3y + z = 7 \)
or
       \( a_{1} > x^{2n} + y^{2n} > x' \)
or
       \( \ip{A}{B} = \sum_{i} a_{i} b_{i} \).
The spaces you type in a formula are
ignored.  Remember that a letter like
       $x$                   % $ ... $  and  \( ... \)  are equivalent
is a formula when it denotes a mathematical
symbol, and it should be typed as one.

\section{Displayed Text}

Text is displayed by indenting it from the left
margin.  Quotations are commonly displayed.  There
are short quotations
\begin{quote}
   This is a short quotation.  It consists of a
   single paragraph of text.  See how it is formatted.
\end{quote}
and longer ones.
\begin{quotation}
   This is a longer quotation.  It consists of two
   paragraphs of text, neither of which are
   particularly interesting.

   This is the second paragraph of the quotation.  It
   is just as dull as the first paragraph.
\end{quotation}
Another frequently-displayed structure is a list.
The following is an example of an \emph{itemized}
list.
\begin{itemize}
   \item This is the first item of an itemized list.
         Each item in the list is marked with a ``tick''.
         You don't have to worry about what kind of tick
         mark is used.

   \item This is the second item of the list.  It
         contains another list nested inside it.  The inner
         list is an \emph{enumerated} list.
         \begin{enumerate}
            \item This is the first item of an enumerated
                  list that is nested within the itemized list.

            \item This is the second item of the inner list.
                  \LaTeX\ allows you to nest lists deeper than
                  you really should.
         \end{enumerate}
         This is the rest of the second item of the outer
         list.  It is no more interesting than any other
         part of the item.
   \item This is the third item of the list.
\end{itemize}
You can even display poetry.
\begin{verse}
   There is an environment
    for verse \\             % The \\ command separates lines
   Whose features some poets % within a stanza.
   will curse.

                             % One or more blank lines separate stanzas.

   For instead of making\\
   Them do \emph{all} line breaking, \\
   It allows them to put too many words on a line when they'd rather be
   forced to be terse.
\end{verse}

Mathematical formulas may also be displayed.  A
displayed formula
is
one-line long; multiline
formulas require special formatting instructions.
   \[  \ip{\Gamma}{\psi'} = x'' + y^{2} + z_{i}^{n}\]
Don't start a paragraph with a displayed equation,
nor make one a paragraph by itself.

\end{document}               % End of document.

\section{Bench Huge Section 1122}
% synthetic bench section 1122
% This is a sample LaTeX input file.  (Version of 12 August 2004.)
%
% A '%' character causes TeX to ignore all remaining text on the line,
% and is used for comments like this one.

\documentclass{article}      % Specifies the document class

                             % The preamble begins here.
\title{An Example Document}  % Declares the document's title.
\author{Leslie Lamport}      % Declares the author's name.
\date{January 21, 1994}      % Deleting this command produces today's date.

\newcommand{\ip}[2]{(#1, #2)}
                             % Defines \ip{arg1}{arg2} to mean
                             % (arg1, arg2).

%\newcommand{\ip}[2]{\langle #1 | #2\rangle}
                             % This is an alternative definition of
                             % \ip that is commented out.

\begin{document}             % End of preamble and beginning of text.

\maketitle                   % Produces the title.

This is an example input file.  Comparing it with
the output it generates can show you how to
produce a simple document of your own.

\section{Ordinary Text}      % Produces section heading.  Lower-level
                             % sections are begun with similar
                             % \subsection and \subsubsection commands.

The ends  of words and sentences are marked
  by   spaces. It  doesn't matter how many
spaces    you type; one is as good as 100.  The
end of   a line counts as a space.

One   or more   blank lines denote the  end
of  a paragraph.

Since any number of consecutive spaces are treated
like a single one, the formatting of the input
file makes no difference to
      \LaTeX,                % The \LaTeX command generates the LaTeX logo.
but it makes a difference to you.  When you use
\LaTeX, making your input file as easy to read
as possible will be a great help as you write
your document and when you change it.  This sample
file shows how you can add comments to your own input
file.

Because printing is different from typewriting,
there are a number of things that you have to do
differently when preparing an input file than if
you were just typing the document directly.
Quotation marks like
       ``this''
have to be handled specially, as do quotes within
quotes:
       ``\,`this'            % \, separates the double and single quote.
        is what I just
        wrote, not  `that'\,''.

Dashes come in three sizes: an
       intra-word
dash, a medium dash for number ranges like
       1--2,
and a punctuation
       dash---like
this.

A sentence-ending space should be larger than the
space between words within a sentence.  You
sometimes have to type special commands in
conjunction with punctuation characters to get
this right, as in the following sentence.
       Gnats, gnus, etc.\ all  % `\ ' makes an inter-word space.
       begin with G\@.         % \@ marks end-of-sentence punctuation.
You should check the spaces after periods when
reading your output to make sure you haven't
forgotten any special cases.  Generating an
ellipsis
       \ldots\               % `\ ' is needed after `\ldots' because TeX
                             % ignores spaces after command names like \ldots
                             % made from \ + letters.
                             %
                             % Note how a `%' character causes TeX to ignore
                             % the end of the input line, so these blank lines
                             % do not start a new paragraph.
                             %
with the right spacing around the periods requires
a special command.

\LaTeX\ interprets some common characters as
commands, so you must type special commands to
generate them.  These characters include the
following:
       \$ \& \% \# \{ and \}.

In printing, text is usually emphasized with an
       \emph{italic}
type style.

\begin{em}
   A long segment of text can also be emphasized
   in this way.  Text within such a segment can be
   given \emph{additional} emphasis.
\end{em}

It is sometimes necessary to prevent \LaTeX\ from
breaking a line where it might otherwise do so.
This may be at a space, as between the ``Mr.''\ and
``Jones'' in
       ``Mr.~Jones'',        % ~ produces an unbreakable interword space.
or within a word---especially when the word is a
symbol like
       \mbox{\emph{itemnum}}
that makes little sense when hyphenated across
lines.

Footnotes\footnote{This is an example of a footnote.}
pose no problem.

\LaTeX\ is good at typesetting mathematical formulas
like
       \( x-3y + z = 7 \)
or
       \( a_{1} > x^{2n} + y^{2n} > x' \)
or
       \( \ip{A}{B} = \sum_{i} a_{i} b_{i} \).
The spaces you type in a formula are
ignored.  Remember that a letter like
       $x$                   % $ ... $  and  \( ... \)  are equivalent
is a formula when it denotes a mathematical
symbol, and it should be typed as one.

\section{Displayed Text}

Text is displayed by indenting it from the left
margin.  Quotations are commonly displayed.  There
are short quotations
\begin{quote}
   This is a short quotation.  It consists of a
   single paragraph of text.  See how it is formatted.
\end{quote}
and longer ones.
\begin{quotation}
   This is a longer quotation.  It consists of two
   paragraphs of text, neither of which are
   particularly interesting.

   This is the second paragraph of the quotation.  It
   is just as dull as the first paragraph.
\end{quotation}
Another frequently-displayed structure is a list.
The following is an example of an \emph{itemized}
list.
\begin{itemize}
   \item This is the first item of an itemized list.
         Each item in the list is marked with a ``tick''.
         You don't have to worry about what kind of tick
         mark is used.

   \item This is the second item of the list.  It
         contains another list nested inside it.  The inner
         list is an \emph{enumerated} list.
         \begin{enumerate}
            \item This is the first item of an enumerated
                  list that is nested within the itemized list.

            \item This is the second item of the inner list.
                  \LaTeX\ allows you to nest lists deeper than
                  you really should.
         \end{enumerate}
         This is the rest of the second item of the outer
         list.  It is no more interesting than any other
         part of the item.
   \item This is the third item of the list.
\end{itemize}
You can even display poetry.
\begin{verse}
   There is an environment
    for verse \\             % The \\ command separates lines
   Whose features some poets % within a stanza.
   will curse.

                             % One or more blank lines separate stanzas.

   For instead of making\\
   Them do \emph{all} line breaking, \\
   It allows them to put too many words on a line when they'd rather be
   forced to be terse.
\end{verse}

Mathematical formulas may also be displayed.  A
displayed formula
is
one-line long; multiline
formulas require special formatting instructions.
   \[  \ip{\Gamma}{\psi'} = x'' + y^{2} + z_{i}^{n}\]
Don't start a paragraph with a displayed equation,
nor make one a paragraph by itself.

\end{document}               % End of document.

\section{Bench Huge Section 1123}
% synthetic bench section 1123
% This is a sample LaTeX input file.  (Version of 12 August 2004.)
%
% A '%' character causes TeX to ignore all remaining text on the line,
% and is used for comments like this one.

\documentclass{article}      % Specifies the document class

                             % The preamble begins here.
\title{An Example Document}  % Declares the document's title.
\author{Leslie Lamport}      % Declares the author's name.
\date{January 21, 1994}      % Deleting this command produces today's date.

\newcommand{\ip}[2]{(#1, #2)}
                             % Defines \ip{arg1}{arg2} to mean
                             % (arg1, arg2).

%\newcommand{\ip}[2]{\langle #1 | #2\rangle}
                             % This is an alternative definition of
                             % \ip that is commented out.

\begin{document}             % End of preamble and beginning of text.

\maketitle                   % Produces the title.

This is an example input file.  Comparing it with
the output it generates can show you how to
produce a simple document of your own.

\section{Ordinary Text}      % Produces section heading.  Lower-level
                             % sections are begun with similar
                             % \subsection and \subsubsection commands.

The ends  of words and sentences are marked
  by   spaces. It  doesn't matter how many
spaces    you type; one is as good as 100.  The
end of   a line counts as a space.

One   or more   blank lines denote the  end
of  a paragraph.

Since any number of consecutive spaces are treated
like a single one, the formatting of the input
file makes no difference to
      \LaTeX,                % The \LaTeX command generates the LaTeX logo.
but it makes a difference to you.  When you use
\LaTeX, making your input file as easy to read
as possible will be a great help as you write
your document and when you change it.  This sample
file shows how you can add comments to your own input
file.

Because printing is different from typewriting,
there are a number of things that you have to do
differently when preparing an input file than if
you were just typing the document directly.
Quotation marks like
       ``this''
have to be handled specially, as do quotes within
quotes:
       ``\,`this'            % \, separates the double and single quote.
        is what I just
        wrote, not  `that'\,''.

Dashes come in three sizes: an
       intra-word
dash, a medium dash for number ranges like
       1--2,
and a punctuation
       dash---like
this.

A sentence-ending space should be larger than the
space between words within a sentence.  You
sometimes have to type special commands in
conjunction with punctuation characters to get
this right, as in the following sentence.
       Gnats, gnus, etc.\ all  % `\ ' makes an inter-word space.
       begin with G\@.         % \@ marks end-of-sentence punctuation.
You should check the spaces after periods when
reading your output to make sure you haven't
forgotten any special cases.  Generating an
ellipsis
       \ldots\               % `\ ' is needed after `\ldots' because TeX
                             % ignores spaces after command names like \ldots
                             % made from \ + letters.
                             %
                             % Note how a `%' character causes TeX to ignore
                             % the end of the input line, so these blank lines
                             % do not start a new paragraph.
                             %
with the right spacing around the periods requires
a special command.

\LaTeX\ interprets some common characters as
commands, so you must type special commands to
generate them.  These characters include the
following:
       \$ \& \% \# \{ and \}.

In printing, text is usually emphasized with an
       \emph{italic}
type style.

\begin{em}
   A long segment of text can also be emphasized
   in this way.  Text within such a segment can be
   given \emph{additional} emphasis.
\end{em}

It is sometimes necessary to prevent \LaTeX\ from
breaking a line where it might otherwise do so.
This may be at a space, as between the ``Mr.''\ and
``Jones'' in
       ``Mr.~Jones'',        % ~ produces an unbreakable interword space.
or within a word---especially when the word is a
symbol like
       \mbox{\emph{itemnum}}
that makes little sense when hyphenated across
lines.

Footnotes\footnote{This is an example of a footnote.}
pose no problem.

\LaTeX\ is good at typesetting mathematical formulas
like
       \( x-3y + z = 7 \)
or
       \( a_{1} > x^{2n} + y^{2n} > x' \)
or
       \( \ip{A}{B} = \sum_{i} a_{i} b_{i} \).
The spaces you type in a formula are
ignored.  Remember that a letter like
       $x$                   % $ ... $  and  \( ... \)  are equivalent
is a formula when it denotes a mathematical
symbol, and it should be typed as one.

\section{Displayed Text}

Text is displayed by indenting it from the left
margin.  Quotations are commonly displayed.  There
are short quotations
\begin{quote}
   This is a short quotation.  It consists of a
   single paragraph of text.  See how it is formatted.
\end{quote}
and longer ones.
\begin{quotation}
   This is a longer quotation.  It consists of two
   paragraphs of text, neither of which are
   particularly interesting.

   This is the second paragraph of the quotation.  It
   is just as dull as the first paragraph.
\end{quotation}
Another frequently-displayed structure is a list.
The following is an example of an \emph{itemized}
list.
\begin{itemize}
   \item This is the first item of an itemized list.
         Each item in the list is marked with a ``tick''.
         You don't have to worry about what kind of tick
         mark is used.

   \item This is the second item of the list.  It
         contains another list nested inside it.  The inner
         list is an \emph{enumerated} list.
         \begin{enumerate}
            \item This is the first item of an enumerated
                  list that is nested within the itemized list.

            \item This is the second item of the inner list.
                  \LaTeX\ allows you to nest lists deeper than
                  you really should.
         \end{enumerate}
         This is the rest of the second item of the outer
         list.  It is no more interesting than any other
         part of the item.
   \item This is the third item of the list.
\end{itemize}
You can even display poetry.
\begin{verse}
   There is an environment
    for verse \\             % The \\ command separates lines
   Whose features some poets % within a stanza.
   will curse.

                             % One or more blank lines separate stanzas.

   For instead of making\\
   Them do \emph{all} line breaking, \\
   It allows them to put too many words on a line when they'd rather be
   forced to be terse.
\end{verse}

Mathematical formulas may also be displayed.  A
displayed formula
is
one-line long; multiline
formulas require special formatting instructions.
   \[  \ip{\Gamma}{\psi'} = x'' + y^{2} + z_{i}^{n}\]
Don't start a paragraph with a displayed equation,
nor make one a paragraph by itself.

\end{document}               % End of document.

\section{Bench Huge Section 1124}
% synthetic bench section 1124
% This is a sample LaTeX input file.  (Version of 12 August 2004.)
%
% A '%' character causes TeX to ignore all remaining text on the line,
% and is used for comments like this one.

\documentclass{article}      % Specifies the document class

                             % The preamble begins here.
\title{An Example Document}  % Declares the document's title.
\author{Leslie Lamport}      % Declares the author's name.
\date{January 21, 1994}      % Deleting this command produces today's date.

\newcommand{\ip}[2]{(#1, #2)}
                             % Defines \ip{arg1}{arg2} to mean
                             % (arg1, arg2).

%\newcommand{\ip}[2]{\langle #1 | #2\rangle}
                             % This is an alternative definition of
                             % \ip that is commented out.

\begin{document}             % End of preamble and beginning of text.

\maketitle                   % Produces the title.

This is an example input file.  Comparing it with
the output it generates can show you how to
produce a simple document of your own.

\section{Ordinary Text}      % Produces section heading.  Lower-level
                             % sections are begun with similar
                             % \subsection and \subsubsection commands.

The ends  of words and sentences are marked
  by   spaces. It  doesn't matter how many
spaces    you type; one is as good as 100.  The
end of   a line counts as a space.

One   or more   blank lines denote the  end
of  a paragraph.

Since any number of consecutive spaces are treated
like a single one, the formatting of the input
file makes no difference to
      \LaTeX,                % The \LaTeX command generates the LaTeX logo.
but it makes a difference to you.  When you use
\LaTeX, making your input file as easy to read
as possible will be a great help as you write
your document and when you change it.  This sample
file shows how you can add comments to your own input
file.

Because printing is different from typewriting,
there are a number of things that you have to do
differently when preparing an input file than if
you were just typing the document directly.
Quotation marks like
       ``this''
have to be handled specially, as do quotes within
quotes:
       ``\,`this'            % \, separates the double and single quote.
        is what I just
        wrote, not  `that'\,''.

Dashes come in three sizes: an
       intra-word
dash, a medium dash for number ranges like
       1--2,
and a punctuation
       dash---like
this.

A sentence-ending space should be larger than the
space between words within a sentence.  You
sometimes have to type special commands in
conjunction with punctuation characters to get
this right, as in the following sentence.
       Gnats, gnus, etc.\ all  % `\ ' makes an inter-word space.
       begin with G\@.         % \@ marks end-of-sentence punctuation.
You should check the spaces after periods when
reading your output to make sure you haven't
forgotten any special cases.  Generating an
ellipsis
       \ldots\               % `\ ' is needed after `\ldots' because TeX
                             % ignores spaces after command names like \ldots
                             % made from \ + letters.
                             %
                             % Note how a `%' character causes TeX to ignore
                             % the end of the input line, so these blank lines
                             % do not start a new paragraph.
                             %
with the right spacing around the periods requires
a special command.

\LaTeX\ interprets some common characters as
commands, so you must type special commands to
generate them.  These characters include the
following:
       \$ \& \% \# \{ and \}.

In printing, text is usually emphasized with an
       \emph{italic}
type style.

\begin{em}
   A long segment of text can also be emphasized
   in this way.  Text within such a segment can be
   given \emph{additional} emphasis.
\end{em}

It is sometimes necessary to prevent \LaTeX\ from
breaking a line where it might otherwise do so.
This may be at a space, as between the ``Mr.''\ and
``Jones'' in
       ``Mr.~Jones'',        % ~ produces an unbreakable interword space.
or within a word---especially when the word is a
symbol like
       \mbox{\emph{itemnum}}
that makes little sense when hyphenated across
lines.

Footnotes\footnote{This is an example of a footnote.}
pose no problem.

\LaTeX\ is good at typesetting mathematical formulas
like
       \( x-3y + z = 7 \)
or
       \( a_{1} > x^{2n} + y^{2n} > x' \)
or
       \( \ip{A}{B} = \sum_{i} a_{i} b_{i} \).
The spaces you type in a formula are
ignored.  Remember that a letter like
       $x$                   % $ ... $  and  \( ... \)  are equivalent
is a formula when it denotes a mathematical
symbol, and it should be typed as one.

\section{Displayed Text}

Text is displayed by indenting it from the left
margin.  Quotations are commonly displayed.  There
are short quotations
\begin{quote}
   This is a short quotation.  It consists of a
   single paragraph of text.  See how it is formatted.
\end{quote}
and longer ones.
\begin{quotation}
   This is a longer quotation.  It consists of two
   paragraphs of text, neither of which are
   particularly interesting.

   This is the second paragraph of the quotation.  It
   is just as dull as the first paragraph.
\end{quotation}
Another frequently-displayed structure is a list.
The following is an example of an \emph{itemized}
list.
\begin{itemize}
   \item This is the first item of an itemized list.
         Each item in the list is marked with a ``tick''.
         You don't have to worry about what kind of tick
         mark is used.

   \item This is the second item of the list.  It
         contains another list nested inside it.  The inner
         list is an \emph{enumerated} list.
         \begin{enumerate}
            \item This is the first item of an enumerated
                  list that is nested within the itemized list.

            \item This is the second item of the inner list.
                  \LaTeX\ allows you to nest lists deeper than
                  you really should.
         \end{enumerate}
         This is the rest of the second item of the outer
         list.  It is no more interesting than any other
         part of the item.
   \item This is the third item of the list.
\end{itemize}
You can even display poetry.
\begin{verse}
   There is an environment
    for verse \\             % The \\ command separates lines
   Whose features some poets % within a stanza.
   will curse.

                             % One or more blank lines separate stanzas.

   For instead of making\\
   Them do \emph{all} line breaking, \\
   It allows them to put too many words on a line when they'd rather be
   forced to be terse.
\end{verse}

Mathematical formulas may also be displayed.  A
displayed formula
is
one-line long; multiline
formulas require special formatting instructions.
   \[  \ip{\Gamma}{\psi'} = x'' + y^{2} + z_{i}^{n}\]
Don't start a paragraph with a displayed equation,
nor make one a paragraph by itself.

\end{document}               % End of document.

\section{Bench Huge Section 1125}
% synthetic bench section 1125
% This is a sample LaTeX input file.  (Version of 12 August 2004.)
%
% A '%' character causes TeX to ignore all remaining text on the line,
% and is used for comments like this one.

\documentclass{article}      % Specifies the document class

                             % The preamble begins here.
\title{An Example Document}  % Declares the document's title.
\author{Leslie Lamport}      % Declares the author's name.
\date{January 21, 1994}      % Deleting this command produces today's date.

\newcommand{\ip}[2]{(#1, #2)}
                             % Defines \ip{arg1}{arg2} to mean
                             % (arg1, arg2).

%\newcommand{\ip}[2]{\langle #1 | #2\rangle}
                             % This is an alternative definition of
                             % \ip that is commented out.

\begin{document}             % End of preamble and beginning of text.

\maketitle                   % Produces the title.

This is an example input file.  Comparing it with
the output it generates can show you how to
produce a simple document of your own.

\section{Ordinary Text}      % Produces section heading.  Lower-level
                             % sections are begun with similar
                             % \subsection and \subsubsection commands.

The ends  of words and sentences are marked
  by   spaces. It  doesn't matter how many
spaces    you type; one is as good as 100.  The
end of   a line counts as a space.

One   or more   blank lines denote the  end
of  a paragraph.

Since any number of consecutive spaces are treated
like a single one, the formatting of the input
file makes no difference to
      \LaTeX,                % The \LaTeX command generates the LaTeX logo.
but it makes a difference to you.  When you use
\LaTeX, making your input file as easy to read
as possible will be a great help as you write
your document and when you change it.  This sample
file shows how you can add comments to your own input
file.

Because printing is different from typewriting,
there are a number of things that you have to do
differently when preparing an input file than if
you were just typing the document directly.
Quotation marks like
       ``this''
have to be handled specially, as do quotes within
quotes:
       ``\,`this'            % \, separates the double and single quote.
        is what I just
        wrote, not  `that'\,''.

Dashes come in three sizes: an
       intra-word
dash, a medium dash for number ranges like
       1--2,
and a punctuation
       dash---like
this.

A sentence-ending space should be larger than the
space between words within a sentence.  You
sometimes have to type special commands in
conjunction with punctuation characters to get
this right, as in the following sentence.
       Gnats, gnus, etc.\ all  % `\ ' makes an inter-word space.
       begin with G\@.         % \@ marks end-of-sentence punctuation.
You should check the spaces after periods when
reading your output to make sure you haven't
forgotten any special cases.  Generating an
ellipsis
       \ldots\               % `\ ' is needed after `\ldots' because TeX
                             % ignores spaces after command names like \ldots
                             % made from \ + letters.
                             %
                             % Note how a `%' character causes TeX to ignore
                             % the end of the input line, so these blank lines
                             % do not start a new paragraph.
                             %
with the right spacing around the periods requires
a special command.

\LaTeX\ interprets some common characters as
commands, so you must type special commands to
generate them.  These characters include the
following:
       \$ \& \% \# \{ and \}.

In printing, text is usually emphasized with an
       \emph{italic}
type style.

\begin{em}
   A long segment of text can also be emphasized
   in this way.  Text within such a segment can be
   given \emph{additional} emphasis.
\end{em}

It is sometimes necessary to prevent \LaTeX\ from
breaking a line where it might otherwise do so.
This may be at a space, as between the ``Mr.''\ and
``Jones'' in
       ``Mr.~Jones'',        % ~ produces an unbreakable interword space.
or within a word---especially when the word is a
symbol like
       \mbox{\emph{itemnum}}
that makes little sense when hyphenated across
lines.

Footnotes\footnote{This is an example of a footnote.}
pose no problem.

\LaTeX\ is good at typesetting mathematical formulas
like
       \( x-3y + z = 7 \)
or
       \( a_{1} > x^{2n} + y^{2n} > x' \)
or
       \( \ip{A}{B} = \sum_{i} a_{i} b_{i} \).
The spaces you type in a formula are
ignored.  Remember that a letter like
       $x$                   % $ ... $  and  \( ... \)  are equivalent
is a formula when it denotes a mathematical
symbol, and it should be typed as one.

\section{Displayed Text}

Text is displayed by indenting it from the left
margin.  Quotations are commonly displayed.  There
are short quotations
\begin{quote}
   This is a short quotation.  It consists of a
   single paragraph of text.  See how it is formatted.
\end{quote}
and longer ones.
\begin{quotation}
   This is a longer quotation.  It consists of two
   paragraphs of text, neither of which are
   particularly interesting.

   This is the second paragraph of the quotation.  It
   is just as dull as the first paragraph.
\end{quotation}
Another frequently-displayed structure is a list.
The following is an example of an \emph{itemized}
list.
\begin{itemize}
   \item This is the first item of an itemized list.
         Each item in the list is marked with a ``tick''.
         You don't have to worry about what kind of tick
         mark is used.

   \item This is the second item of the list.  It
         contains another list nested inside it.  The inner
         list is an \emph{enumerated} list.
         \begin{enumerate}
            \item This is the first item of an enumerated
                  list that is nested within the itemized list.

            \item This is the second item of the inner list.
                  \LaTeX\ allows you to nest lists deeper than
                  you really should.
         \end{enumerate}
         This is the rest of the second item of the outer
         list.  It is no more interesting than any other
         part of the item.
   \item This is the third item of the list.
\end{itemize}
You can even display poetry.
\begin{verse}
   There is an environment
    for verse \\             % The \\ command separates lines
   Whose features some poets % within a stanza.
   will curse.

                             % One or more blank lines separate stanzas.

   For instead of making\\
   Them do \emph{all} line breaking, \\
   It allows them to put too many words on a line when they'd rather be
   forced to be terse.
\end{verse}

Mathematical formulas may also be displayed.  A
displayed formula
is
one-line long; multiline
formulas require special formatting instructions.
   \[  \ip{\Gamma}{\psi'} = x'' + y^{2} + z_{i}^{n}\]
Don't start a paragraph with a displayed equation,
nor make one a paragraph by itself.

\end{document}               % End of document.

\section{Bench Huge Section 1126}
% synthetic bench section 1126
% This is a sample LaTeX input file.  (Version of 12 August 2004.)
%
% A '%' character causes TeX to ignore all remaining text on the line,
% and is used for comments like this one.

\documentclass{article}      % Specifies the document class

                             % The preamble begins here.
\title{An Example Document}  % Declares the document's title.
\author{Leslie Lamport}      % Declares the author's name.
\date{January 21, 1994}      % Deleting this command produces today's date.

\newcommand{\ip}[2]{(#1, #2)}
                             % Defines \ip{arg1}{arg2} to mean
                             % (arg1, arg2).

%\newcommand{\ip}[2]{\langle #1 | #2\rangle}
                             % This is an alternative definition of
                             % \ip that is commented out.

\begin{document}             % End of preamble and beginning of text.

\maketitle                   % Produces the title.

This is an example input file.  Comparing it with
the output it generates can show you how to
produce a simple document of your own.

\section{Ordinary Text}      % Produces section heading.  Lower-level
                             % sections are begun with similar
                             % \subsection and \subsubsection commands.

The ends  of words and sentences are marked
  by   spaces. It  doesn't matter how many
spaces    you type; one is as good as 100.  The
end of   a line counts as a space.

One   or more   blank lines denote the  end
of  a paragraph.

Since any number of consecutive spaces are treated
like a single one, the formatting of the input
file makes no difference to
      \LaTeX,                % The \LaTeX command generates the LaTeX logo.
but it makes a difference to you.  When you use
\LaTeX, making your input file as easy to read
as possible will be a great help as you write
your document and when you change it.  This sample
file shows how you can add comments to your own input
file.

Because printing is different from typewriting,
there are a number of things that you have to do
differently when preparing an input file than if
you were just typing the document directly.
Quotation marks like
       ``this''
have to be handled specially, as do quotes within
quotes:
       ``\,`this'            % \, separates the double and single quote.
        is what I just
        wrote, not  `that'\,''.

Dashes come in three sizes: an
       intra-word
dash, a medium dash for number ranges like
       1--2,
and a punctuation
       dash---like
this.

A sentence-ending space should be larger than the
space between words within a sentence.  You
sometimes have to type special commands in
conjunction with punctuation characters to get
this right, as in the following sentence.
       Gnats, gnus, etc.\ all  % `\ ' makes an inter-word space.
       begin with G\@.         % \@ marks end-of-sentence punctuation.
You should check the spaces after periods when
reading your output to make sure you haven't
forgotten any special cases.  Generating an
ellipsis
       \ldots\               % `\ ' is needed after `\ldots' because TeX
                             % ignores spaces after command names like \ldots
                             % made from \ + letters.
                             %
                             % Note how a `%' character causes TeX to ignore
                             % the end of the input line, so these blank lines
                             % do not start a new paragraph.
                             %
with the right spacing around the periods requires
a special command.

\LaTeX\ interprets some common characters as
commands, so you must type special commands to
generate them.  These characters include the
following:
       \$ \& \% \# \{ and \}.

In printing, text is usually emphasized with an
       \emph{italic}
type style.

\begin{em}
   A long segment of text can also be emphasized
   in this way.  Text within such a segment can be
   given \emph{additional} emphasis.
\end{em}

It is sometimes necessary to prevent \LaTeX\ from
breaking a line where it might otherwise do so.
This may be at a space, as between the ``Mr.''\ and
``Jones'' in
       ``Mr.~Jones'',        % ~ produces an unbreakable interword space.
or within a word---especially when the word is a
symbol like
       \mbox{\emph{itemnum}}
that makes little sense when hyphenated across
lines.

Footnotes\footnote{This is an example of a footnote.}
pose no problem.

\LaTeX\ is good at typesetting mathematical formulas
like
       \( x-3y + z = 7 \)
or
       \( a_{1} > x^{2n} + y^{2n} > x' \)
or
       \( \ip{A}{B} = \sum_{i} a_{i} b_{i} \).
The spaces you type in a formula are
ignored.  Remember that a letter like
       $x$                   % $ ... $  and  \( ... \)  are equivalent
is a formula when it denotes a mathematical
symbol, and it should be typed as one.

\section{Displayed Text}

Text is displayed by indenting it from the left
margin.  Quotations are commonly displayed.  There
are short quotations
\begin{quote}
   This is a short quotation.  It consists of a
   single paragraph of text.  See how it is formatted.
\end{quote}
and longer ones.
\begin{quotation}
   This is a longer quotation.  It consists of two
   paragraphs of text, neither of which are
   particularly interesting.

   This is the second paragraph of the quotation.  It
   is just as dull as the first paragraph.
\end{quotation}
Another frequently-displayed structure is a list.
The following is an example of an \emph{itemized}
list.
\begin{itemize}
   \item This is the first item of an itemized list.
         Each item in the list is marked with a ``tick''.
         You don't have to worry about what kind of tick
         mark is used.

   \item This is the second item of the list.  It
         contains another list nested inside it.  The inner
         list is an \emph{enumerated} list.
         \begin{enumerate}
            \item This is the first item of an enumerated
                  list that is nested within the itemized list.

            \item This is the second item of the inner list.
                  \LaTeX\ allows you to nest lists deeper than
                  you really should.
         \end{enumerate}
         This is the rest of the second item of the outer
         list.  It is no more interesting than any other
         part of the item.
   \item This is the third item of the list.
\end{itemize}
You can even display poetry.
\begin{verse}
   There is an environment
    for verse \\             % The \\ command separates lines
   Whose features some poets % within a stanza.
   will curse.

                             % One or more blank lines separate stanzas.

   For instead of making\\
   Them do \emph{all} line breaking, \\
   It allows them to put too many words on a line when they'd rather be
   forced to be terse.
\end{verse}

Mathematical formulas may also be displayed.  A
displayed formula
is
one-line long; multiline
formulas require special formatting instructions.
   \[  \ip{\Gamma}{\psi'} = x'' + y^{2} + z_{i}^{n}\]
Don't start a paragraph with a displayed equation,
nor make one a paragraph by itself.

\end{document}               % End of document.

\section{Bench Huge Section 1127}
% synthetic bench section 1127
% This is a sample LaTeX input file.  (Version of 12 August 2004.)
%
% A '%' character causes TeX to ignore all remaining text on the line,
% and is used for comments like this one.

\documentclass{article}      % Specifies the document class

                             % The preamble begins here.
\title{An Example Document}  % Declares the document's title.
\author{Leslie Lamport}      % Declares the author's name.
\date{January 21, 1994}      % Deleting this command produces today's date.

\newcommand{\ip}[2]{(#1, #2)}
                             % Defines \ip{arg1}{arg2} to mean
                             % (arg1, arg2).

%\newcommand{\ip}[2]{\langle #1 | #2\rangle}
                             % This is an alternative definition of
                             % \ip that is commented out.

\begin{document}             % End of preamble and beginning of text.

\maketitle                   % Produces the title.

This is an example input file.  Comparing it with
the output it generates can show you how to
produce a simple document of your own.

\section{Ordinary Text}      % Produces section heading.  Lower-level
                             % sections are begun with similar
                             % \subsection and \subsubsection commands.

The ends  of words and sentences are marked
  by   spaces. It  doesn't matter how many
spaces    you type; one is as good as 100.  The
end of   a line counts as a space.

One   or more   blank lines denote the  end
of  a paragraph.

Since any number of consecutive spaces are treated
like a single one, the formatting of the input
file makes no difference to
      \LaTeX,                % The \LaTeX command generates the LaTeX logo.
but it makes a difference to you.  When you use
\LaTeX, making your input file as easy to read
as possible will be a great help as you write
your document and when you change it.  This sample
file shows how you can add comments to your own input
file.

Because printing is different from typewriting,
there are a number of things that you have to do
differently when preparing an input file than if
you were just typing the document directly.
Quotation marks like
       ``this''
have to be handled specially, as do quotes within
quotes:
       ``\,`this'            % \, separates the double and single quote.
        is what I just
        wrote, not  `that'\,''.

Dashes come in three sizes: an
       intra-word
dash, a medium dash for number ranges like
       1--2,
and a punctuation
       dash---like
this.

A sentence-ending space should be larger than the
space between words within a sentence.  You
sometimes have to type special commands in
conjunction with punctuation characters to get
this right, as in the following sentence.
       Gnats, gnus, etc.\ all  % `\ ' makes an inter-word space.
       begin with G\@.         % \@ marks end-of-sentence punctuation.
You should check the spaces after periods when
reading your output to make sure you haven't
forgotten any special cases.  Generating an
ellipsis
       \ldots\               % `\ ' is needed after `\ldots' because TeX
                             % ignores spaces after command names like \ldots
                             % made from \ + letters.
                             %
                             % Note how a `%' character causes TeX to ignore
                             % the end of the input line, so these blank lines
                             % do not start a new paragraph.
                             %
with the right spacing around the periods requires
a special command.

\LaTeX\ interprets some common characters as
commands, so you must type special commands to
generate them.  These characters include the
following:
       \$ \& \% \# \{ and \}.

In printing, text is usually emphasized with an
       \emph{italic}
type style.

\begin{em}
   A long segment of text can also be emphasized
   in this way.  Text within such a segment can be
   given \emph{additional} emphasis.
\end{em}

It is sometimes necessary to prevent \LaTeX\ from
breaking a line where it might otherwise do so.
This may be at a space, as between the ``Mr.''\ and
``Jones'' in
       ``Mr.~Jones'',        % ~ produces an unbreakable interword space.
or within a word---especially when the word is a
symbol like
       \mbox{\emph{itemnum}}
that makes little sense when hyphenated across
lines.

Footnotes\footnote{This is an example of a footnote.}
pose no problem.

\LaTeX\ is good at typesetting mathematical formulas
like
       \( x-3y + z = 7 \)
or
       \( a_{1} > x^{2n} + y^{2n} > x' \)
or
       \( \ip{A}{B} = \sum_{i} a_{i} b_{i} \).
The spaces you type in a formula are
ignored.  Remember that a letter like
       $x$                   % $ ... $  and  \( ... \)  are equivalent
is a formula when it denotes a mathematical
symbol, and it should be typed as one.

\section{Displayed Text}

Text is displayed by indenting it from the left
margin.  Quotations are commonly displayed.  There
are short quotations
\begin{quote}
   This is a short quotation.  It consists of a
   single paragraph of text.  See how it is formatted.
\end{quote}
and longer ones.
\begin{quotation}
   This is a longer quotation.  It consists of two
   paragraphs of text, neither of which are
   particularly interesting.

   This is the second paragraph of the quotation.  It
   is just as dull as the first paragraph.
\end{quotation}
Another frequently-displayed structure is a list.
The following is an example of an \emph{itemized}
list.
\begin{itemize}
   \item This is the first item of an itemized list.
         Each item in the list is marked with a ``tick''.
         You don't have to worry about what kind of tick
         mark is used.

   \item This is the second item of the list.  It
         contains another list nested inside it.  The inner
         list is an \emph{enumerated} list.
         \begin{enumerate}
            \item This is the first item of an enumerated
                  list that is nested within the itemized list.

            \item This is the second item of the inner list.
                  \LaTeX\ allows you to nest lists deeper than
                  you really should.
         \end{enumerate}
         This is the rest of the second item of the outer
         list.  It is no more interesting than any other
         part of the item.
   \item This is the third item of the list.
\end{itemize}
You can even display poetry.
\begin{verse}
   There is an environment
    for verse \\             % The \\ command separates lines
   Whose features some poets % within a stanza.
   will curse.

                             % One or more blank lines separate stanzas.

   For instead of making\\
   Them do \emph{all} line breaking, \\
   It allows them to put too many words on a line when they'd rather be
   forced to be terse.
\end{verse}

Mathematical formulas may also be displayed.  A
displayed formula
is
one-line long; multiline
formulas require special formatting instructions.
   \[  \ip{\Gamma}{\psi'} = x'' + y^{2} + z_{i}^{n}\]
Don't start a paragraph with a displayed equation,
nor make one a paragraph by itself.

\end{document}               % End of document.

\section{Bench Huge Section 1128}
% synthetic bench section 1128
% This is a sample LaTeX input file.  (Version of 12 August 2004.)
%
% A '%' character causes TeX to ignore all remaining text on the line,
% and is used for comments like this one.

\documentclass{article}      % Specifies the document class

                             % The preamble begins here.
\title{An Example Document}  % Declares the document's title.
\author{Leslie Lamport}      % Declares the author's name.
\date{January 21, 1994}      % Deleting this command produces today's date.

\newcommand{\ip}[2]{(#1, #2)}
                             % Defines \ip{arg1}{arg2} to mean
                             % (arg1, arg2).

%\newcommand{\ip}[2]{\langle #1 | #2\rangle}
                             % This is an alternative definition of
                             % \ip that is commented out.

\begin{document}             % End of preamble and beginning of text.

\maketitle                   % Produces the title.

This is an example input file.  Comparing it with
the output it generates can show you how to
produce a simple document of your own.

\section{Ordinary Text}      % Produces section heading.  Lower-level
                             % sections are begun with similar
                             % \subsection and \subsubsection commands.

The ends  of words and sentences are marked
  by   spaces. It  doesn't matter how many
spaces    you type; one is as good as 100.  The
end of   a line counts as a space.

One   or more   blank lines denote the  end
of  a paragraph.

Since any number of consecutive spaces are treated
like a single one, the formatting of the input
file makes no difference to
      \LaTeX,                % The \LaTeX command generates the LaTeX logo.
but it makes a difference to you.  When you use
\LaTeX, making your input file as easy to read
as possible will be a great help as you write
your document and when you change it.  This sample
file shows how you can add comments to your own input
file.

Because printing is different from typewriting,
there are a number of things that you have to do
differently when preparing an input file than if
you were just typing the document directly.
Quotation marks like
       ``this''
have to be handled specially, as do quotes within
quotes:
       ``\,`this'            % \, separates the double and single quote.
        is what I just
        wrote, not  `that'\,''.

Dashes come in three sizes: an
       intra-word
dash, a medium dash for number ranges like
       1--2,
and a punctuation
       dash---like
this.

A sentence-ending space should be larger than the
space between words within a sentence.  You
sometimes have to type special commands in
conjunction with punctuation characters to get
this right, as in the following sentence.
       Gnats, gnus, etc.\ all  % `\ ' makes an inter-word space.
       begin with G\@.         % \@ marks end-of-sentence punctuation.
You should check the spaces after periods when
reading your output to make sure you haven't
forgotten any special cases.  Generating an
ellipsis
       \ldots\               % `\ ' is needed after `\ldots' because TeX
                             % ignores spaces after command names like \ldots
                             % made from \ + letters.
                             %
                             % Note how a `%' character causes TeX to ignore
                             % the end of the input line, so these blank lines
                             % do not start a new paragraph.
                             %
with the right spacing around the periods requires
a special command.

\LaTeX\ interprets some common characters as
commands, so you must type special commands to
generate them.  These characters include the
following:
       \$ \& \% \# \{ and \}.

In printing, text is usually emphasized with an
       \emph{italic}
type style.

\begin{em}
   A long segment of text can also be emphasized
   in this way.  Text within such a segment can be
   given \emph{additional} emphasis.
\end{em}

It is sometimes necessary to prevent \LaTeX\ from
breaking a line where it might otherwise do so.
This may be at a space, as between the ``Mr.''\ and
``Jones'' in
       ``Mr.~Jones'',        % ~ produces an unbreakable interword space.
or within a word---especially when the word is a
symbol like
       \mbox{\emph{itemnum}}
that makes little sense when hyphenated across
lines.

Footnotes\footnote{This is an example of a footnote.}
pose no problem.

\LaTeX\ is good at typesetting mathematical formulas
like
       \( x-3y + z = 7 \)
or
       \( a_{1} > x^{2n} + y^{2n} > x' \)
or
       \( \ip{A}{B} = \sum_{i} a_{i} b_{i} \).
The spaces you type in a formula are
ignored.  Remember that a letter like
       $x$                   % $ ... $  and  \( ... \)  are equivalent
is a formula when it denotes a mathematical
symbol, and it should be typed as one.

\section{Displayed Text}

Text is displayed by indenting it from the left
margin.  Quotations are commonly displayed.  There
are short quotations
\begin{quote}
   This is a short quotation.  It consists of a
   single paragraph of text.  See how it is formatted.
\end{quote}
and longer ones.
\begin{quotation}
   This is a longer quotation.  It consists of two
   paragraphs of text, neither of which are
   particularly interesting.

   This is the second paragraph of the quotation.  It
   is just as dull as the first paragraph.
\end{quotation}
Another frequently-displayed structure is a list.
The following is an example of an \emph{itemized}
list.
\begin{itemize}
   \item This is the first item of an itemized list.
         Each item in the list is marked with a ``tick''.
         You don't have to worry about what kind of tick
         mark is used.

   \item This is the second item of the list.  It
         contains another list nested inside it.  The inner
         list is an \emph{enumerated} list.
         \begin{enumerate}
            \item This is the first item of an enumerated
                  list that is nested within the itemized list.

            \item This is the second item of the inner list.
                  \LaTeX\ allows you to nest lists deeper than
                  you really should.
         \end{enumerate}
         This is the rest of the second item of the outer
         list.  It is no more interesting than any other
         part of the item.
   \item This is the third item of the list.
\end{itemize}
You can even display poetry.
\begin{verse}
   There is an environment
    for verse \\             % The \\ command separates lines
   Whose features some poets % within a stanza.
   will curse.

                             % One or more blank lines separate stanzas.

   For instead of making\\
   Them do \emph{all} line breaking, \\
   It allows them to put too many words on a line when they'd rather be
   forced to be terse.
\end{verse}

Mathematical formulas may also be displayed.  A
displayed formula
is
one-line long; multiline
formulas require special formatting instructions.
   \[  \ip{\Gamma}{\psi'} = x'' + y^{2} + z_{i}^{n}\]
Don't start a paragraph with a displayed equation,
nor make one a paragraph by itself.

\end{document}               % End of document.

\section{Bench Huge Section 1129}
% synthetic bench section 1129
% This is a sample LaTeX input file.  (Version of 12 August 2004.)
%
% A '%' character causes TeX to ignore all remaining text on the line,
% and is used for comments like this one.

\documentclass{article}      % Specifies the document class

                             % The preamble begins here.
\title{An Example Document}  % Declares the document's title.
\author{Leslie Lamport}      % Declares the author's name.
\date{January 21, 1994}      % Deleting this command produces today's date.

\newcommand{\ip}[2]{(#1, #2)}
                             % Defines \ip{arg1}{arg2} to mean
                             % (arg1, arg2).

%\newcommand{\ip}[2]{\langle #1 | #2\rangle}
                             % This is an alternative definition of
                             % \ip that is commented out.

\begin{document}             % End of preamble and beginning of text.

\maketitle                   % Produces the title.

This is an example input file.  Comparing it with
the output it generates can show you how to
produce a simple document of your own.

\section{Ordinary Text}      % Produces section heading.  Lower-level
                             % sections are begun with similar
                             % \subsection and \subsubsection commands.

The ends  of words and sentences are marked
  by   spaces. It  doesn't matter how many
spaces    you type; one is as good as 100.  The
end of   a line counts as a space.

One   or more   blank lines denote the  end
of  a paragraph.

Since any number of consecutive spaces are treated
like a single one, the formatting of the input
file makes no difference to
      \LaTeX,                % The \LaTeX command generates the LaTeX logo.
but it makes a difference to you.  When you use
\LaTeX, making your input file as easy to read
as possible will be a great help as you write
your document and when you change it.  This sample
file shows how you can add comments to your own input
file.

Because printing is different from typewriting,
there are a number of things that you have to do
differently when preparing an input file than if
you were just typing the document directly.
Quotation marks like
       ``this''
have to be handled specially, as do quotes within
quotes:
       ``\,`this'            % \, separates the double and single quote.
        is what I just
        wrote, not  `that'\,''.

Dashes come in three sizes: an
       intra-word
dash, a medium dash for number ranges like
       1--2,
and a punctuation
       dash---like
this.

A sentence-ending space should be larger than the
space between words within a sentence.  You
sometimes have to type special commands in
conjunction with punctuation characters to get
this right, as in the following sentence.
       Gnats, gnus, etc.\ all  % `\ ' makes an inter-word space.
       begin with G\@.         % \@ marks end-of-sentence punctuation.
You should check the spaces after periods when
reading your output to make sure you haven't
forgotten any special cases.  Generating an
ellipsis
       \ldots\               % `\ ' is needed after `\ldots' because TeX
                             % ignores spaces after command names like \ldots
                             % made from \ + letters.
                             %
                             % Note how a `%' character causes TeX to ignore
                             % the end of the input line, so these blank lines
                             % do not start a new paragraph.
                             %
with the right spacing around the periods requires
a special command.

\LaTeX\ interprets some common characters as
commands, so you must type special commands to
generate them.  These characters include the
following:
       \$ \& \% \# \{ and \}.

In printing, text is usually emphasized with an
       \emph{italic}
type style.

\begin{em}
   A long segment of text can also be emphasized
   in this way.  Text within such a segment can be
   given \emph{additional} emphasis.
\end{em}

It is sometimes necessary to prevent \LaTeX\ from
breaking a line where it might otherwise do so.
This may be at a space, as between the ``Mr.''\ and
``Jones'' in
       ``Mr.~Jones'',        % ~ produces an unbreakable interword space.
or within a word---especially when the word is a
symbol like
       \mbox{\emph{itemnum}}
that makes little sense when hyphenated across
lines.

Footnotes\footnote{This is an example of a footnote.}
pose no problem.

\LaTeX\ is good at typesetting mathematical formulas
like
       \( x-3y + z = 7 \)
or
       \( a_{1} > x^{2n} + y^{2n} > x' \)
or
       \( \ip{A}{B} = \sum_{i} a_{i} b_{i} \).
The spaces you type in a formula are
ignored.  Remember that a letter like
       $x$                   % $ ... $  and  \( ... \)  are equivalent
is a formula when it denotes a mathematical
symbol, and it should be typed as one.

\section{Displayed Text}

Text is displayed by indenting it from the left
margin.  Quotations are commonly displayed.  There
are short quotations
\begin{quote}
   This is a short quotation.  It consists of a
   single paragraph of text.  See how it is formatted.
\end{quote}
and longer ones.
\begin{quotation}
   This is a longer quotation.  It consists of two
   paragraphs of text, neither of which are
   particularly interesting.

   This is the second paragraph of the quotation.  It
   is just as dull as the first paragraph.
\end{quotation}
Another frequently-displayed structure is a list.
The following is an example of an \emph{itemized}
list.
\begin{itemize}
   \item This is the first item of an itemized list.
         Each item in the list is marked with a ``tick''.
         You don't have to worry about what kind of tick
         mark is used.

   \item This is the second item of the list.  It
         contains another list nested inside it.  The inner
         list is an \emph{enumerated} list.
         \begin{enumerate}
            \item This is the first item of an enumerated
                  list that is nested within the itemized list.

            \item This is the second item of the inner list.
                  \LaTeX\ allows you to nest lists deeper than
                  you really should.
         \end{enumerate}
         This is the rest of the second item of the outer
         list.  It is no more interesting than any other
         part of the item.
   \item This is the third item of the list.
\end{itemize}
You can even display poetry.
\begin{verse}
   There is an environment
    for verse \\             % The \\ command separates lines
   Whose features some poets % within a stanza.
   will curse.

                             % One or more blank lines separate stanzas.

   For instead of making\\
   Them do \emph{all} line breaking, \\
   It allows them to put too many words on a line when they'd rather be
   forced to be terse.
\end{verse}

Mathematical formulas may also be displayed.  A
displayed formula
is
one-line long; multiline
formulas require special formatting instructions.
   \[  \ip{\Gamma}{\psi'} = x'' + y^{2} + z_{i}^{n}\]
Don't start a paragraph with a displayed equation,
nor make one a paragraph by itself.

\end{document}               % End of document.

\section{Bench Huge Section 1130}
% synthetic bench section 1130
% This is a sample LaTeX input file.  (Version of 12 August 2004.)
%
% A '%' character causes TeX to ignore all remaining text on the line,
% and is used for comments like this one.

\documentclass{article}      % Specifies the document class

                             % The preamble begins here.
\title{An Example Document}  % Declares the document's title.
\author{Leslie Lamport}      % Declares the author's name.
\date{January 21, 1994}      % Deleting this command produces today's date.

\newcommand{\ip}[2]{(#1, #2)}
                             % Defines \ip{arg1}{arg2} to mean
                             % (arg1, arg2).

%\newcommand{\ip}[2]{\langle #1 | #2\rangle}
                             % This is an alternative definition of
                             % \ip that is commented out.

\begin{document}             % End of preamble and beginning of text.

\maketitle                   % Produces the title.

This is an example input file.  Comparing it with
the output it generates can show you how to
produce a simple document of your own.

\section{Ordinary Text}      % Produces section heading.  Lower-level
                             % sections are begun with similar
                             % \subsection and \subsubsection commands.

The ends  of words and sentences are marked
  by   spaces. It  doesn't matter how many
spaces    you type; one is as good as 100.  The
end of   a line counts as a space.

One   or more   blank lines denote the  end
of  a paragraph.

Since any number of consecutive spaces are treated
like a single one, the formatting of the input
file makes no difference to
      \LaTeX,                % The \LaTeX command generates the LaTeX logo.
but it makes a difference to you.  When you use
\LaTeX, making your input file as easy to read
as possible will be a great help as you write
your document and when you change it.  This sample
file shows how you can add comments to your own input
file.

Because printing is different from typewriting,
there are a number of things that you have to do
differently when preparing an input file than if
you were just typing the document directly.
Quotation marks like
       ``this''
have to be handled specially, as do quotes within
quotes:
       ``\,`this'            % \, separates the double and single quote.
        is what I just
        wrote, not  `that'\,''.

Dashes come in three sizes: an
       intra-word
dash, a medium dash for number ranges like
       1--2,
and a punctuation
       dash---like
this.

A sentence-ending space should be larger than the
space between words within a sentence.  You
sometimes have to type special commands in
conjunction with punctuation characters to get
this right, as in the following sentence.
       Gnats, gnus, etc.\ all  % `\ ' makes an inter-word space.
       begin with G\@.         % \@ marks end-of-sentence punctuation.
You should check the spaces after periods when
reading your output to make sure you haven't
forgotten any special cases.  Generating an
ellipsis
       \ldots\               % `\ ' is needed after `\ldots' because TeX
                             % ignores spaces after command names like \ldots
                             % made from \ + letters.
                             %
                             % Note how a `%' character causes TeX to ignore
                             % the end of the input line, so these blank lines
                             % do not start a new paragraph.
                             %
with the right spacing around the periods requires
a special command.

\LaTeX\ interprets some common characters as
commands, so you must type special commands to
generate them.  These characters include the
following:
       \$ \& \% \# \{ and \}.

In printing, text is usually emphasized with an
       \emph{italic}
type style.

\begin{em}
   A long segment of text can also be emphasized
   in this way.  Text within such a segment can be
   given \emph{additional} emphasis.
\end{em}

It is sometimes necessary to prevent \LaTeX\ from
breaking a line where it might otherwise do so.
This may be at a space, as between the ``Mr.''\ and
``Jones'' in
       ``Mr.~Jones'',        % ~ produces an unbreakable interword space.
or within a word---especially when the word is a
symbol like
       \mbox{\emph{itemnum}}
that makes little sense when hyphenated across
lines.

Footnotes\footnote{This is an example of a footnote.}
pose no problem.

\LaTeX\ is good at typesetting mathematical formulas
like
       \( x-3y + z = 7 \)
or
       \( a_{1} > x^{2n} + y^{2n} > x' \)
or
       \( \ip{A}{B} = \sum_{i} a_{i} b_{i} \).
The spaces you type in a formula are
ignored.  Remember that a letter like
       $x$                   % $ ... $  and  \( ... \)  are equivalent
is a formula when it denotes a mathematical
symbol, and it should be typed as one.

\section{Displayed Text}

Text is displayed by indenting it from the left
margin.  Quotations are commonly displayed.  There
are short quotations
\begin{quote}
   This is a short quotation.  It consists of a
   single paragraph of text.  See how it is formatted.
\end{quote}
and longer ones.
\begin{quotation}
   This is a longer quotation.  It consists of two
   paragraphs of text, neither of which are
   particularly interesting.

   This is the second paragraph of the quotation.  It
   is just as dull as the first paragraph.
\end{quotation}
Another frequently-displayed structure is a list.
The following is an example of an \emph{itemized}
list.
\begin{itemize}
   \item This is the first item of an itemized list.
         Each item in the list is marked with a ``tick''.
         You don't have to worry about what kind of tick
         mark is used.

   \item This is the second item of the list.  It
         contains another list nested inside it.  The inner
         list is an \emph{enumerated} list.
         \begin{enumerate}
            \item This is the first item of an enumerated
                  list that is nested within the itemized list.

            \item This is the second item of the inner list.
                  \LaTeX\ allows you to nest lists deeper than
                  you really should.
         \end{enumerate}
         This is the rest of the second item of the outer
         list.  It is no more interesting than any other
         part of the item.
   \item This is the third item of the list.
\end{itemize}
You can even display poetry.
\begin{verse}
   There is an environment
    for verse \\             % The \\ command separates lines
   Whose features some poets % within a stanza.
   will curse.

                             % One or more blank lines separate stanzas.

   For instead of making\\
   Them do \emph{all} line breaking, \\
   It allows them to put too many words on a line when they'd rather be
   forced to be terse.
\end{verse}

Mathematical formulas may also be displayed.  A
displayed formula
is
one-line long; multiline
formulas require special formatting instructions.
   \[  \ip{\Gamma}{\psi'} = x'' + y^{2} + z_{i}^{n}\]
Don't start a paragraph with a displayed equation,
nor make one a paragraph by itself.

\end{document}               % End of document.

\section{Bench Huge Section 1131}
% synthetic bench section 1131
% This is a sample LaTeX input file.  (Version of 12 August 2004.)
%
% A '%' character causes TeX to ignore all remaining text on the line,
% and is used for comments like this one.

\documentclass{article}      % Specifies the document class

                             % The preamble begins here.
\title{An Example Document}  % Declares the document's title.
\author{Leslie Lamport}      % Declares the author's name.
\date{January 21, 1994}      % Deleting this command produces today's date.

\newcommand{\ip}[2]{(#1, #2)}
                             % Defines \ip{arg1}{arg2} to mean
                             % (arg1, arg2).

%\newcommand{\ip}[2]{\langle #1 | #2\rangle}
                             % This is an alternative definition of
                             % \ip that is commented out.

\begin{document}             % End of preamble and beginning of text.

\maketitle                   % Produces the title.

This is an example input file.  Comparing it with
the output it generates can show you how to
produce a simple document of your own.

\section{Ordinary Text}      % Produces section heading.  Lower-level
                             % sections are begun with similar
                             % \subsection and \subsubsection commands.

The ends  of words and sentences are marked
  by   spaces. It  doesn't matter how many
spaces    you type; one is as good as 100.  The
end of   a line counts as a space.

One   or more   blank lines denote the  end
of  a paragraph.

Since any number of consecutive spaces are treated
like a single one, the formatting of the input
file makes no difference to
      \LaTeX,                % The \LaTeX command generates the LaTeX logo.
but it makes a difference to you.  When you use
\LaTeX, making your input file as easy to read
as possible will be a great help as you write
your document and when you change it.  This sample
file shows how you can add comments to your own input
file.

Because printing is different from typewriting,
there are a number of things that you have to do
differently when preparing an input file than if
you were just typing the document directly.
Quotation marks like
       ``this''
have to be handled specially, as do quotes within
quotes:
       ``\,`this'            % \, separates the double and single quote.
        is what I just
        wrote, not  `that'\,''.

Dashes come in three sizes: an
       intra-word
dash, a medium dash for number ranges like
       1--2,
and a punctuation
       dash---like
this.

A sentence-ending space should be larger than the
space between words within a sentence.  You
sometimes have to type special commands in
conjunction with punctuation characters to get
this right, as in the following sentence.
       Gnats, gnus, etc.\ all  % `\ ' makes an inter-word space.
       begin with G\@.         % \@ marks end-of-sentence punctuation.
You should check the spaces after periods when
reading your output to make sure you haven't
forgotten any special cases.  Generating an
ellipsis
       \ldots\               % `\ ' is needed after `\ldots' because TeX
                             % ignores spaces after command names like \ldots
                             % made from \ + letters.
                             %
                             % Note how a `%' character causes TeX to ignore
                             % the end of the input line, so these blank lines
                             % do not start a new paragraph.
                             %
with the right spacing around the periods requires
a special command.

\LaTeX\ interprets some common characters as
commands, so you must type special commands to
generate them.  These characters include the
following:
       \$ \& \% \# \{ and \}.

In printing, text is usually emphasized with an
       \emph{italic}
type style.

\begin{em}
   A long segment of text can also be emphasized
   in this way.  Text within such a segment can be
   given \emph{additional} emphasis.
\end{em}

It is sometimes necessary to prevent \LaTeX\ from
breaking a line where it might otherwise do so.
This may be at a space, as between the ``Mr.''\ and
``Jones'' in
       ``Mr.~Jones'',        % ~ produces an unbreakable interword space.
or within a word---especially when the word is a
symbol like
       \mbox{\emph{itemnum}}
that makes little sense when hyphenated across
lines.

Footnotes\footnote{This is an example of a footnote.}
pose no problem.

\LaTeX\ is good at typesetting mathematical formulas
like
       \( x-3y + z = 7 \)
or
       \( a_{1} > x^{2n} + y^{2n} > x' \)
or
       \( \ip{A}{B} = \sum_{i} a_{i} b_{i} \).
The spaces you type in a formula are
ignored.  Remember that a letter like
       $x$                   % $ ... $  and  \( ... \)  are equivalent
is a formula when it denotes a mathematical
symbol, and it should be typed as one.

\section{Displayed Text}

Text is displayed by indenting it from the left
margin.  Quotations are commonly displayed.  There
are short quotations
\begin{quote}
   This is a short quotation.  It consists of a
   single paragraph of text.  See how it is formatted.
\end{quote}
and longer ones.
\begin{quotation}
   This is a longer quotation.  It consists of two
   paragraphs of text, neither of which are
   particularly interesting.

   This is the second paragraph of the quotation.  It
   is just as dull as the first paragraph.
\end{quotation}
Another frequently-displayed structure is a list.
The following is an example of an \emph{itemized}
list.
\begin{itemize}
   \item This is the first item of an itemized list.
         Each item in the list is marked with a ``tick''.
         You don't have to worry about what kind of tick
         mark is used.

   \item This is the second item of the list.  It
         contains another list nested inside it.  The inner
         list is an \emph{enumerated} list.
         \begin{enumerate}
            \item This is the first item of an enumerated
                  list that is nested within the itemized list.

            \item This is the second item of the inner list.
                  \LaTeX\ allows you to nest lists deeper than
                  you really should.
         \end{enumerate}
         This is the rest of the second item of the outer
         list.  It is no more interesting than any other
         part of the item.
   \item This is the third item of the list.
\end{itemize}
You can even display poetry.
\begin{verse}
   There is an environment
    for verse \\             % The \\ command separates lines
   Whose features some poets % within a stanza.
   will curse.

                             % One or more blank lines separate stanzas.

   For instead of making\\
   Them do \emph{all} line breaking, \\
   It allows them to put too many words on a line when they'd rather be
   forced to be terse.
\end{verse}

Mathematical formulas may also be displayed.  A
displayed formula
is
one-line long; multiline
formulas require special formatting instructions.
   \[  \ip{\Gamma}{\psi'} = x'' + y^{2} + z_{i}^{n}\]
Don't start a paragraph with a displayed equation,
nor make one a paragraph by itself.

\end{document}               % End of document.

\section{Bench Huge Section 1132}
% synthetic bench section 1132
% This is a sample LaTeX input file.  (Version of 12 August 2004.)
%
% A '%' character causes TeX to ignore all remaining text on the line,
% and is used for comments like this one.

\documentclass{article}      % Specifies the document class

                             % The preamble begins here.
\title{An Example Document}  % Declares the document's title.
\author{Leslie Lamport}      % Declares the author's name.
\date{January 21, 1994}      % Deleting this command produces today's date.

\newcommand{\ip}[2]{(#1, #2)}
                             % Defines \ip{arg1}{arg2} to mean
                             % (arg1, arg2).

%\newcommand{\ip}[2]{\langle #1 | #2\rangle}
                             % This is an alternative definition of
                             % \ip that is commented out.

\begin{document}             % End of preamble and beginning of text.

\maketitle                   % Produces the title.

This is an example input file.  Comparing it with
the output it generates can show you how to
produce a simple document of your own.

\section{Ordinary Text}      % Produces section heading.  Lower-level
                             % sections are begun with similar
                             % \subsection and \subsubsection commands.

The ends  of words and sentences are marked
  by   spaces. It  doesn't matter how many
spaces    you type; one is as good as 100.  The
end of   a line counts as a space.

One   or more   blank lines denote the  end
of  a paragraph.

Since any number of consecutive spaces are treated
like a single one, the formatting of the input
file makes no difference to
      \LaTeX,                % The \LaTeX command generates the LaTeX logo.
but it makes a difference to you.  When you use
\LaTeX, making your input file as easy to read
as possible will be a great help as you write
your document and when you change it.  This sample
file shows how you can add comments to your own input
file.

Because printing is different from typewriting,
there are a number of things that you have to do
differently when preparing an input file than if
you were just typing the document directly.
Quotation marks like
       ``this''
have to be handled specially, as do quotes within
quotes:
       ``\,`this'            % \, separates the double and single quote.
        is what I just
        wrote, not  `that'\,''.

Dashes come in three sizes: an
       intra-word
dash, a medium dash for number ranges like
       1--2,
and a punctuation
       dash---like
this.

A sentence-ending space should be larger than the
space between words within a sentence.  You
sometimes have to type special commands in
conjunction with punctuation characters to get
this right, as in the following sentence.
       Gnats, gnus, etc.\ all  % `\ ' makes an inter-word space.
       begin with G\@.         % \@ marks end-of-sentence punctuation.
You should check the spaces after periods when
reading your output to make sure you haven't
forgotten any special cases.  Generating an
ellipsis
       \ldots\               % `\ ' is needed after `\ldots' because TeX
                             % ignores spaces after command names like \ldots
                             % made from \ + letters.
                             %
                             % Note how a `%' character causes TeX to ignore
                             % the end of the input line, so these blank lines
                             % do not start a new paragraph.
                             %
with the right spacing around the periods requires
a special command.

\LaTeX\ interprets some common characters as
commands, so you must type special commands to
generate them.  These characters include the
following:
       \$ \& \% \# \{ and \}.

In printing, text is usually emphasized with an
       \emph{italic}
type style.

\begin{em}
   A long segment of text can also be emphasized
   in this way.  Text within such a segment can be
   given \emph{additional} emphasis.
\end{em}

It is sometimes necessary to prevent \LaTeX\ from
breaking a line where it might otherwise do so.
This may be at a space, as between the ``Mr.''\ and
``Jones'' in
       ``Mr.~Jones'',        % ~ produces an unbreakable interword space.
or within a word---especially when the word is a
symbol like
       \mbox{\emph{itemnum}}
that makes little sense when hyphenated across
lines.

Footnotes\footnote{This is an example of a footnote.}
pose no problem.

\LaTeX\ is good at typesetting mathematical formulas
like
       \( x-3y + z = 7 \)
or
       \( a_{1} > x^{2n} + y^{2n} > x' \)
or
       \( \ip{A}{B} = \sum_{i} a_{i} b_{i} \).
The spaces you type in a formula are
ignored.  Remember that a letter like
       $x$                   % $ ... $  and  \( ... \)  are equivalent
is a formula when it denotes a mathematical
symbol, and it should be typed as one.

\section{Displayed Text}

Text is displayed by indenting it from the left
margin.  Quotations are commonly displayed.  There
are short quotations
\begin{quote}
   This is a short quotation.  It consists of a
   single paragraph of text.  See how it is formatted.
\end{quote}
and longer ones.
\begin{quotation}
   This is a longer quotation.  It consists of two
   paragraphs of text, neither of which are
   particularly interesting.

   This is the second paragraph of the quotation.  It
   is just as dull as the first paragraph.
\end{quotation}
Another frequently-displayed structure is a list.
The following is an example of an \emph{itemized}
list.
\begin{itemize}
   \item This is the first item of an itemized list.
         Each item in the list is marked with a ``tick''.
         You don't have to worry about what kind of tick
         mark is used.

   \item This is the second item of the list.  It
         contains another list nested inside it.  The inner
         list is an \emph{enumerated} list.
         \begin{enumerate}
            \item This is the first item of an enumerated
                  list that is nested within the itemized list.

            \item This is the second item of the inner list.
                  \LaTeX\ allows you to nest lists deeper than
                  you really should.
         \end{enumerate}
         This is the rest of the second item of the outer
         list.  It is no more interesting than any other
         part of the item.
   \item This is the third item of the list.
\end{itemize}
You can even display poetry.
\begin{verse}
   There is an environment
    for verse \\             % The \\ command separates lines
   Whose features some poets % within a stanza.
   will curse.

                             % One or more blank lines separate stanzas.

   For instead of making\\
   Them do \emph{all} line breaking, \\
   It allows them to put too many words on a line when they'd rather be
   forced to be terse.
\end{verse}

Mathematical formulas may also be displayed.  A
displayed formula
is
one-line long; multiline
formulas require special formatting instructions.
   \[  \ip{\Gamma}{\psi'} = x'' + y^{2} + z_{i}^{n}\]
Don't start a paragraph with a displayed equation,
nor make one a paragraph by itself.

\end{document}               % End of document.

\section{Bench Huge Section 1133}
% synthetic bench section 1133
% This is a sample LaTeX input file.  (Version of 12 August 2004.)
%
% A '%' character causes TeX to ignore all remaining text on the line,
% and is used for comments like this one.

\documentclass{article}      % Specifies the document class

                             % The preamble begins here.
\title{An Example Document}  % Declares the document's title.
\author{Leslie Lamport}      % Declares the author's name.
\date{January 21, 1994}      % Deleting this command produces today's date.

\newcommand{\ip}[2]{(#1, #2)}
                             % Defines \ip{arg1}{arg2} to mean
                             % (arg1, arg2).

%\newcommand{\ip}[2]{\langle #1 | #2\rangle}
                             % This is an alternative definition of
                             % \ip that is commented out.

\begin{document}             % End of preamble and beginning of text.

\maketitle                   % Produces the title.

This is an example input file.  Comparing it with
the output it generates can show you how to
produce a simple document of your own.

\section{Ordinary Text}      % Produces section heading.  Lower-level
                             % sections are begun with similar
                             % \subsection and \subsubsection commands.

The ends  of words and sentences are marked
  by   spaces. It  doesn't matter how many
spaces    you type; one is as good as 100.  The
end of   a line counts as a space.

One   or more   blank lines denote the  end
of  a paragraph.

Since any number of consecutive spaces are treated
like a single one, the formatting of the input
file makes no difference to
      \LaTeX,                % The \LaTeX command generates the LaTeX logo.
but it makes a difference to you.  When you use
\LaTeX, making your input file as easy to read
as possible will be a great help as you write
your document and when you change it.  This sample
file shows how you can add comments to your own input
file.

Because printing is different from typewriting,
there are a number of things that you have to do
differently when preparing an input file than if
you were just typing the document directly.
Quotation marks like
       ``this''
have to be handled specially, as do quotes within
quotes:
       ``\,`this'            % \, separates the double and single quote.
        is what I just
        wrote, not  `that'\,''.

Dashes come in three sizes: an
       intra-word
dash, a medium dash for number ranges like
       1--2,
and a punctuation
       dash---like
this.

A sentence-ending space should be larger than the
space between words within a sentence.  You
sometimes have to type special commands in
conjunction with punctuation characters to get
this right, as in the following sentence.
       Gnats, gnus, etc.\ all  % `\ ' makes an inter-word space.
       begin with G\@.         % \@ marks end-of-sentence punctuation.
You should check the spaces after periods when
reading your output to make sure you haven't
forgotten any special cases.  Generating an
ellipsis
       \ldots\               % `\ ' is needed after `\ldots' because TeX
                             % ignores spaces after command names like \ldots
                             % made from \ + letters.
                             %
                             % Note how a `%' character causes TeX to ignore
                             % the end of the input line, so these blank lines
                             % do not start a new paragraph.
                             %
with the right spacing around the periods requires
a special command.

\LaTeX\ interprets some common characters as
commands, so you must type special commands to
generate them.  These characters include the
following:
       \$ \& \% \# \{ and \}.

In printing, text is usually emphasized with an
       \emph{italic}
type style.

\begin{em}
   A long segment of text can also be emphasized
   in this way.  Text within such a segment can be
   given \emph{additional} emphasis.
\end{em}

It is sometimes necessary to prevent \LaTeX\ from
breaking a line where it might otherwise do so.
This may be at a space, as between the ``Mr.''\ and
``Jones'' in
       ``Mr.~Jones'',        % ~ produces an unbreakable interword space.
or within a word---especially when the word is a
symbol like
       \mbox{\emph{itemnum}}
that makes little sense when hyphenated across
lines.

Footnotes\footnote{This is an example of a footnote.}
pose no problem.

\LaTeX\ is good at typesetting mathematical formulas
like
       \( x-3y + z = 7 \)
or
       \( a_{1} > x^{2n} + y^{2n} > x' \)
or
       \( \ip{A}{B} = \sum_{i} a_{i} b_{i} \).
The spaces you type in a formula are
ignored.  Remember that a letter like
       $x$                   % $ ... $  and  \( ... \)  are equivalent
is a formula when it denotes a mathematical
symbol, and it should be typed as one.

\section{Displayed Text}

Text is displayed by indenting it from the left
margin.  Quotations are commonly displayed.  There
are short quotations
\begin{quote}
   This is a short quotation.  It consists of a
   single paragraph of text.  See how it is formatted.
\end{quote}
and longer ones.
\begin{quotation}
   This is a longer quotation.  It consists of two
   paragraphs of text, neither of which are
   particularly interesting.

   This is the second paragraph of the quotation.  It
   is just as dull as the first paragraph.
\end{quotation}
Another frequently-displayed structure is a list.
The following is an example of an \emph{itemized}
list.
\begin{itemize}
   \item This is the first item of an itemized list.
         Each item in the list is marked with a ``tick''.
         You don't have to worry about what kind of tick
         mark is used.

   \item This is the second item of the list.  It
         contains another list nested inside it.  The inner
         list is an \emph{enumerated} list.
         \begin{enumerate}
            \item This is the first item of an enumerated
                  list that is nested within the itemized list.

            \item This is the second item of the inner list.
                  \LaTeX\ allows you to nest lists deeper than
                  you really should.
         \end{enumerate}
         This is the rest of the second item of the outer
         list.  It is no more interesting than any other
         part of the item.
   \item This is the third item of the list.
\end{itemize}
You can even display poetry.
\begin{verse}
   There is an environment
    for verse \\             % The \\ command separates lines
   Whose features some poets % within a stanza.
   will curse.

                             % One or more blank lines separate stanzas.

   For instead of making\\
   Them do \emph{all} line breaking, \\
   It allows them to put too many words on a line when they'd rather be
   forced to be terse.
\end{verse}

Mathematical formulas may also be displayed.  A
displayed formula
is
one-line long; multiline
formulas require special formatting instructions.
   \[  \ip{\Gamma}{\psi'} = x'' + y^{2} + z_{i}^{n}\]
Don't start a paragraph with a displayed equation,
nor make one a paragraph by itself.

\end{document}               % End of document.

\section{Bench Huge Section 1134}
% synthetic bench section 1134
% This is a sample LaTeX input file.  (Version of 12 August 2004.)
%
% A '%' character causes TeX to ignore all remaining text on the line,
% and is used for comments like this one.

\documentclass{article}      % Specifies the document class

                             % The preamble begins here.
\title{An Example Document}  % Declares the document's title.
\author{Leslie Lamport}      % Declares the author's name.
\date{January 21, 1994}      % Deleting this command produces today's date.

\newcommand{\ip}[2]{(#1, #2)}
                             % Defines \ip{arg1}{arg2} to mean
                             % (arg1, arg2).

%\newcommand{\ip}[2]{\langle #1 | #2\rangle}
                             % This is an alternative definition of
                             % \ip that is commented out.

\begin{document}             % End of preamble and beginning of text.

\maketitle                   % Produces the title.

This is an example input file.  Comparing it with
the output it generates can show you how to
produce a simple document of your own.

\section{Ordinary Text}      % Produces section heading.  Lower-level
                             % sections are begun with similar
                             % \subsection and \subsubsection commands.

The ends  of words and sentences are marked
  by   spaces. It  doesn't matter how many
spaces    you type; one is as good as 100.  The
end of   a line counts as a space.

One   or more   blank lines denote the  end
of  a paragraph.

Since any number of consecutive spaces are treated
like a single one, the formatting of the input
file makes no difference to
      \LaTeX,                % The \LaTeX command generates the LaTeX logo.
but it makes a difference to you.  When you use
\LaTeX, making your input file as easy to read
as possible will be a great help as you write
your document and when you change it.  This sample
file shows how you can add comments to your own input
file.

Because printing is different from typewriting,
there are a number of things that you have to do
differently when preparing an input file than if
you were just typing the document directly.
Quotation marks like
       ``this''
have to be handled specially, as do quotes within
quotes:
       ``\,`this'            % \, separates the double and single quote.
        is what I just
        wrote, not  `that'\,''.

Dashes come in three sizes: an
       intra-word
dash, a medium dash for number ranges like
       1--2,
and a punctuation
       dash---like
this.

A sentence-ending space should be larger than the
space between words within a sentence.  You
sometimes have to type special commands in
conjunction with punctuation characters to get
this right, as in the following sentence.
       Gnats, gnus, etc.\ all  % `\ ' makes an inter-word space.
       begin with G\@.         % \@ marks end-of-sentence punctuation.
You should check the spaces after periods when
reading your output to make sure you haven't
forgotten any special cases.  Generating an
ellipsis
       \ldots\               % `\ ' is needed after `\ldots' because TeX
                             % ignores spaces after command names like \ldots
                             % made from \ + letters.
                             %
                             % Note how a `%' character causes TeX to ignore
                             % the end of the input line, so these blank lines
                             % do not start a new paragraph.
                             %
with the right spacing around the periods requires
a special command.

\LaTeX\ interprets some common characters as
commands, so you must type special commands to
generate them.  These characters include the
following:
       \$ \& \% \# \{ and \}.

In printing, text is usually emphasized with an
       \emph{italic}
type style.

\begin{em}
   A long segment of text can also be emphasized
   in this way.  Text within such a segment can be
   given \emph{additional} emphasis.
\end{em}

It is sometimes necessary to prevent \LaTeX\ from
breaking a line where it might otherwise do so.
This may be at a space, as between the ``Mr.''\ and
``Jones'' in
       ``Mr.~Jones'',        % ~ produces an unbreakable interword space.
or within a word---especially when the word is a
symbol like
       \mbox{\emph{itemnum}}
that makes little sense when hyphenated across
lines.

Footnotes\footnote{This is an example of a footnote.}
pose no problem.

\LaTeX\ is good at typesetting mathematical formulas
like
       \( x-3y + z = 7 \)
or
       \( a_{1} > x^{2n} + y^{2n} > x' \)
or
       \( \ip{A}{B} = \sum_{i} a_{i} b_{i} \).
The spaces you type in a formula are
ignored.  Remember that a letter like
       $x$                   % $ ... $  and  \( ... \)  are equivalent
is a formula when it denotes a mathematical
symbol, and it should be typed as one.

\section{Displayed Text}

Text is displayed by indenting it from the left
margin.  Quotations are commonly displayed.  There
are short quotations
\begin{quote}
   This is a short quotation.  It consists of a
   single paragraph of text.  See how it is formatted.
\end{quote}
and longer ones.
\begin{quotation}
   This is a longer quotation.  It consists of two
   paragraphs of text, neither of which are
   particularly interesting.

   This is the second paragraph of the quotation.  It
   is just as dull as the first paragraph.
\end{quotation}
Another frequently-displayed structure is a list.
The following is an example of an \emph{itemized}
list.
\begin{itemize}
   \item This is the first item of an itemized list.
         Each item in the list is marked with a ``tick''.
         You don't have to worry about what kind of tick
         mark is used.

   \item This is the second item of the list.  It
         contains another list nested inside it.  The inner
         list is an \emph{enumerated} list.
         \begin{enumerate}
            \item This is the first item of an enumerated
                  list that is nested within the itemized list.

            \item This is the second item of the inner list.
                  \LaTeX\ allows you to nest lists deeper than
                  you really should.
         \end{enumerate}
         This is the rest of the second item of the outer
         list.  It is no more interesting than any other
         part of the item.
   \item This is the third item of the list.
\end{itemize}
You can even display poetry.
\begin{verse}
   There is an environment
    for verse \\             % The \\ command separates lines
   Whose features some poets % within a stanza.
   will curse.

                             % One or more blank lines separate stanzas.

   For instead of making\\
   Them do \emph{all} line breaking, \\
   It allows them to put too many words on a line when they'd rather be
   forced to be terse.
\end{verse}

Mathematical formulas may also be displayed.  A
displayed formula
is
one-line long; multiline
formulas require special formatting instructions.
   \[  \ip{\Gamma}{\psi'} = x'' + y^{2} + z_{i}^{n}\]
Don't start a paragraph with a displayed equation,
nor make one a paragraph by itself.

\end{document}               % End of document.

\section{Bench Huge Section 1135}
% synthetic bench section 1135
% This is a sample LaTeX input file.  (Version of 12 August 2004.)
%
% A '%' character causes TeX to ignore all remaining text on the line,
% and is used for comments like this one.

\documentclass{article}      % Specifies the document class

                             % The preamble begins here.
\title{An Example Document}  % Declares the document's title.
\author{Leslie Lamport}      % Declares the author's name.
\date{January 21, 1994}      % Deleting this command produces today's date.

\newcommand{\ip}[2]{(#1, #2)}
                             % Defines \ip{arg1}{arg2} to mean
                             % (arg1, arg2).

%\newcommand{\ip}[2]{\langle #1 | #2\rangle}
                             % This is an alternative definition of
                             % \ip that is commented out.

\begin{document}             % End of preamble and beginning of text.

\maketitle                   % Produces the title.

This is an example input file.  Comparing it with
the output it generates can show you how to
produce a simple document of your own.

\section{Ordinary Text}      % Produces section heading.  Lower-level
                             % sections are begun with similar
                             % \subsection and \subsubsection commands.

The ends  of words and sentences are marked
  by   spaces. It  doesn't matter how many
spaces    you type; one is as good as 100.  The
end of   a line counts as a space.

One   or more   blank lines denote the  end
of  a paragraph.

Since any number of consecutive spaces are treated
like a single one, the formatting of the input
file makes no difference to
      \LaTeX,                % The \LaTeX command generates the LaTeX logo.
but it makes a difference to you.  When you use
\LaTeX, making your input file as easy to read
as possible will be a great help as you write
your document and when you change it.  This sample
file shows how you can add comments to your own input
file.

Because printing is different from typewriting,
there are a number of things that you have to do
differently when preparing an input file than if
you were just typing the document directly.
Quotation marks like
       ``this''
have to be handled specially, as do quotes within
quotes:
       ``\,`this'            % \, separates the double and single quote.
        is what I just
        wrote, not  `that'\,''.

Dashes come in three sizes: an
       intra-word
dash, a medium dash for number ranges like
       1--2,
and a punctuation
       dash---like
this.

A sentence-ending space should be larger than the
space between words within a sentence.  You
sometimes have to type special commands in
conjunction with punctuation characters to get
this right, as in the following sentence.
       Gnats, gnus, etc.\ all  % `\ ' makes an inter-word space.
       begin with G\@.         % \@ marks end-of-sentence punctuation.
You should check the spaces after periods when
reading your output to make sure you haven't
forgotten any special cases.  Generating an
ellipsis
       \ldots\               % `\ ' is needed after `\ldots' because TeX
                             % ignores spaces after command names like \ldots
                             % made from \ + letters.
                             %
                             % Note how a `%' character causes TeX to ignore
                             % the end of the input line, so these blank lines
                             % do not start a new paragraph.
                             %
with the right spacing around the periods requires
a special command.

\LaTeX\ interprets some common characters as
commands, so you must type special commands to
generate them.  These characters include the
following:
       \$ \& \% \# \{ and \}.

In printing, text is usually emphasized with an
       \emph{italic}
type style.

\begin{em}
   A long segment of text can also be emphasized
   in this way.  Text within such a segment can be
   given \emph{additional} emphasis.
\end{em}

It is sometimes necessary to prevent \LaTeX\ from
breaking a line where it might otherwise do so.
This may be at a space, as between the ``Mr.''\ and
``Jones'' in
       ``Mr.~Jones'',        % ~ produces an unbreakable interword space.
or within a word---especially when the word is a
symbol like
       \mbox{\emph{itemnum}}
that makes little sense when hyphenated across
lines.

Footnotes\footnote{This is an example of a footnote.}
pose no problem.

\LaTeX\ is good at typesetting mathematical formulas
like
       \( x-3y + z = 7 \)
or
       \( a_{1} > x^{2n} + y^{2n} > x' \)
or
       \( \ip{A}{B} = \sum_{i} a_{i} b_{i} \).
The spaces you type in a formula are
ignored.  Remember that a letter like
       $x$                   % $ ... $  and  \( ... \)  are equivalent
is a formula when it denotes a mathematical
symbol, and it should be typed as one.

\section{Displayed Text}

Text is displayed by indenting it from the left
margin.  Quotations are commonly displayed.  There
are short quotations
\begin{quote}
   This is a short quotation.  It consists of a
   single paragraph of text.  See how it is formatted.
\end{quote}
and longer ones.
\begin{quotation}
   This is a longer quotation.  It consists of two
   paragraphs of text, neither of which are
   particularly interesting.

   This is the second paragraph of the quotation.  It
   is just as dull as the first paragraph.
\end{quotation}
Another frequently-displayed structure is a list.
The following is an example of an \emph{itemized}
list.
\begin{itemize}
   \item This is the first item of an itemized list.
         Each item in the list is marked with a ``tick''.
         You don't have to worry about what kind of tick
         mark is used.

   \item This is the second item of the list.  It
         contains another list nested inside it.  The inner
         list is an \emph{enumerated} list.
         \begin{enumerate}
            \item This is the first item of an enumerated
                  list that is nested within the itemized list.

            \item This is the second item of the inner list.
                  \LaTeX\ allows you to nest lists deeper than
                  you really should.
         \end{enumerate}
         This is the rest of the second item of the outer
         list.  It is no more interesting than any other
         part of the item.
   \item This is the third item of the list.
\end{itemize}
You can even display poetry.
\begin{verse}
   There is an environment
    for verse \\             % The \\ command separates lines
   Whose features some poets % within a stanza.
   will curse.

                             % One or more blank lines separate stanzas.

   For instead of making\\
   Them do \emph{all} line breaking, \\
   It allows them to put too many words on a line when they'd rather be
   forced to be terse.
\end{verse}

Mathematical formulas may also be displayed.  A
displayed formula
is
one-line long; multiline
formulas require special formatting instructions.
   \[  \ip{\Gamma}{\psi'} = x'' + y^{2} + z_{i}^{n}\]
Don't start a paragraph with a displayed equation,
nor make one a paragraph by itself.

\end{document}               % End of document.

\section{Bench Huge Section 1136}
% synthetic bench section 1136
% This is a sample LaTeX input file.  (Version of 12 August 2004.)
%
% A '%' character causes TeX to ignore all remaining text on the line,
% and is used for comments like this one.

\documentclass{article}      % Specifies the document class

                             % The preamble begins here.
\title{An Example Document}  % Declares the document's title.
\author{Leslie Lamport}      % Declares the author's name.
\date{January 21, 1994}      % Deleting this command produces today's date.

\newcommand{\ip}[2]{(#1, #2)}
                             % Defines \ip{arg1}{arg2} to mean
                             % (arg1, arg2).

%\newcommand{\ip}[2]{\langle #1 | #2\rangle}
                             % This is an alternative definition of
                             % \ip that is commented out.

\begin{document}             % End of preamble and beginning of text.

\maketitle                   % Produces the title.

This is an example input file.  Comparing it with
the output it generates can show you how to
produce a simple document of your own.

\section{Ordinary Text}      % Produces section heading.  Lower-level
                             % sections are begun with similar
                             % \subsection and \subsubsection commands.

The ends  of words and sentences are marked
  by   spaces. It  doesn't matter how many
spaces    you type; one is as good as 100.  The
end of   a line counts as a space.

One   or more   blank lines denote the  end
of  a paragraph.

Since any number of consecutive spaces are treated
like a single one, the formatting of the input
file makes no difference to
      \LaTeX,                % The \LaTeX command generates the LaTeX logo.
but it makes a difference to you.  When you use
\LaTeX, making your input file as easy to read
as possible will be a great help as you write
your document and when you change it.  This sample
file shows how you can add comments to your own input
file.

Because printing is different from typewriting,
there are a number of things that you have to do
differently when preparing an input file than if
you were just typing the document directly.
Quotation marks like
       ``this''
have to be handled specially, as do quotes within
quotes:
       ``\,`this'            % \, separates the double and single quote.
        is what I just
        wrote, not  `that'\,''.

Dashes come in three sizes: an
       intra-word
dash, a medium dash for number ranges like
       1--2,
and a punctuation
       dash---like
this.

A sentence-ending space should be larger than the
space between words within a sentence.  You
sometimes have to type special commands in
conjunction with punctuation characters to get
this right, as in the following sentence.
       Gnats, gnus, etc.\ all  % `\ ' makes an inter-word space.
       begin with G\@.         % \@ marks end-of-sentence punctuation.
You should check the spaces after periods when
reading your output to make sure you haven't
forgotten any special cases.  Generating an
ellipsis
       \ldots\               % `\ ' is needed after `\ldots' because TeX
                             % ignores spaces after command names like \ldots
                             % made from \ + letters.
                             %
                             % Note how a `%' character causes TeX to ignore
                             % the end of the input line, so these blank lines
                             % do not start a new paragraph.
                             %
with the right spacing around the periods requires
a special command.

\LaTeX\ interprets some common characters as
commands, so you must type special commands to
generate them.  These characters include the
following:
       \$ \& \% \# \{ and \}.

In printing, text is usually emphasized with an
       \emph{italic}
type style.

\begin{em}
   A long segment of text can also be emphasized
   in this way.  Text within such a segment can be
   given \emph{additional} emphasis.
\end{em}

It is sometimes necessary to prevent \LaTeX\ from
breaking a line where it might otherwise do so.
This may be at a space, as between the ``Mr.''\ and
``Jones'' in
       ``Mr.~Jones'',        % ~ produces an unbreakable interword space.
or within a word---especially when the word is a
symbol like
       \mbox{\emph{itemnum}}
that makes little sense when hyphenated across
lines.

Footnotes\footnote{This is an example of a footnote.}
pose no problem.

\LaTeX\ is good at typesetting mathematical formulas
like
       \( x-3y + z = 7 \)
or
       \( a_{1} > x^{2n} + y^{2n} > x' \)
or
       \( \ip{A}{B} = \sum_{i} a_{i} b_{i} \).
The spaces you type in a formula are
ignored.  Remember that a letter like
       $x$                   % $ ... $  and  \( ... \)  are equivalent
is a formula when it denotes a mathematical
symbol, and it should be typed as one.

\section{Displayed Text}

Text is displayed by indenting it from the left
margin.  Quotations are commonly displayed.  There
are short quotations
\begin{quote}
   This is a short quotation.  It consists of a
   single paragraph of text.  See how it is formatted.
\end{quote}
and longer ones.
\begin{quotation}
   This is a longer quotation.  It consists of two
   paragraphs of text, neither of which are
   particularly interesting.

   This is the second paragraph of the quotation.  It
   is just as dull as the first paragraph.
\end{quotation}
Another frequently-displayed structure is a list.
The following is an example of an \emph{itemized}
list.
\begin{itemize}
   \item This is the first item of an itemized list.
         Each item in the list is marked with a ``tick''.
         You don't have to worry about what kind of tick
         mark is used.

   \item This is the second item of the list.  It
         contains another list nested inside it.  The inner
         list is an \emph{enumerated} list.
         \begin{enumerate}
            \item This is the first item of an enumerated
                  list that is nested within the itemized list.

            \item This is the second item of the inner list.
                  \LaTeX\ allows you to nest lists deeper than
                  you really should.
         \end{enumerate}
         This is the rest of the second item of the outer
         list.  It is no more interesting than any other
         part of the item.
   \item This is the third item of the list.
\end{itemize}
You can even display poetry.
\begin{verse}
   There is an environment
    for verse \\             % The \\ command separates lines
   Whose features some poets % within a stanza.
   will curse.

                             % One or more blank lines separate stanzas.

   For instead of making\\
   Them do \emph{all} line breaking, \\
   It allows them to put too many words on a line when they'd rather be
   forced to be terse.
\end{verse}

Mathematical formulas may also be displayed.  A
displayed formula
is
one-line long; multiline
formulas require special formatting instructions.
   \[  \ip{\Gamma}{\psi'} = x'' + y^{2} + z_{i}^{n}\]
Don't start a paragraph with a displayed equation,
nor make one a paragraph by itself.

\end{document}               % End of document.

\section{Bench Huge Section 1137}
% synthetic bench section 1137
% This is a sample LaTeX input file.  (Version of 12 August 2004.)
%
% A '%' character causes TeX to ignore all remaining text on the line,
% and is used for comments like this one.

\documentclass{article}      % Specifies the document class

                             % The preamble begins here.
\title{An Example Document}  % Declares the document's title.
\author{Leslie Lamport}      % Declares the author's name.
\date{January 21, 1994}      % Deleting this command produces today's date.

\newcommand{\ip}[2]{(#1, #2)}
                             % Defines \ip{arg1}{arg2} to mean
                             % (arg1, arg2).

%\newcommand{\ip}[2]{\langle #1 | #2\rangle}
                             % This is an alternative definition of
                             % \ip that is commented out.

\begin{document}             % End of preamble and beginning of text.

\maketitle                   % Produces the title.

This is an example input file.  Comparing it with
the output it generates can show you how to
produce a simple document of your own.

\section{Ordinary Text}      % Produces section heading.  Lower-level
                             % sections are begun with similar
                             % \subsection and \subsubsection commands.

The ends  of words and sentences are marked
  by   spaces. It  doesn't matter how many
spaces    you type; one is as good as 100.  The
end of   a line counts as a space.

One   or more   blank lines denote the  end
of  a paragraph.

Since any number of consecutive spaces are treated
like a single one, the formatting of the input
file makes no difference to
      \LaTeX,                % The \LaTeX command generates the LaTeX logo.
but it makes a difference to you.  When you use
\LaTeX, making your input file as easy to read
as possible will be a great help as you write
your document and when you change it.  This sample
file shows how you can add comments to your own input
file.

Because printing is different from typewriting,
there are a number of things that you have to do
differently when preparing an input file than if
you were just typing the document directly.
Quotation marks like
       ``this''
have to be handled specially, as do quotes within
quotes:
       ``\,`this'            % \, separates the double and single quote.
        is what I just
        wrote, not  `that'\,''.

Dashes come in three sizes: an
       intra-word
dash, a medium dash for number ranges like
       1--2,
and a punctuation
       dash---like
this.

A sentence-ending space should be larger than the
space between words within a sentence.  You
sometimes have to type special commands in
conjunction with punctuation characters to get
this right, as in the following sentence.
       Gnats, gnus, etc.\ all  % `\ ' makes an inter-word space.
       begin with G\@.         % \@ marks end-of-sentence punctuation.
You should check the spaces after periods when
reading your output to make sure you haven't
forgotten any special cases.  Generating an
ellipsis
       \ldots\               % `\ ' is needed after `\ldots' because TeX
                             % ignores spaces after command names like \ldots
                             % made from \ + letters.
                             %
                             % Note how a `%' character causes TeX to ignore
                             % the end of the input line, so these blank lines
                             % do not start a new paragraph.
                             %
with the right spacing around the periods requires
a special command.

\LaTeX\ interprets some common characters as
commands, so you must type special commands to
generate them.  These characters include the
following:
       \$ \& \% \# \{ and \}.

In printing, text is usually emphasized with an
       \emph{italic}
type style.

\begin{em}
   A long segment of text can also be emphasized
   in this way.  Text within such a segment can be
   given \emph{additional} emphasis.
\end{em}

It is sometimes necessary to prevent \LaTeX\ from
breaking a line where it might otherwise do so.
This may be at a space, as between the ``Mr.''\ and
``Jones'' in
       ``Mr.~Jones'',        % ~ produces an unbreakable interword space.
or within a word---especially when the word is a
symbol like
       \mbox{\emph{itemnum}}
that makes little sense when hyphenated across
lines.

Footnotes\footnote{This is an example of a footnote.}
pose no problem.

\LaTeX\ is good at typesetting mathematical formulas
like
       \( x-3y + z = 7 \)
or
       \( a_{1} > x^{2n} + y^{2n} > x' \)
or
       \( \ip{A}{B} = \sum_{i} a_{i} b_{i} \).
The spaces you type in a formula are
ignored.  Remember that a letter like
       $x$                   % $ ... $  and  \( ... \)  are equivalent
is a formula when it denotes a mathematical
symbol, and it should be typed as one.

\section{Displayed Text}

Text is displayed by indenting it from the left
margin.  Quotations are commonly displayed.  There
are short quotations
\begin{quote}
   This is a short quotation.  It consists of a
   single paragraph of text.  See how it is formatted.
\end{quote}
and longer ones.
\begin{quotation}
   This is a longer quotation.  It consists of two
   paragraphs of text, neither of which are
   particularly interesting.

   This is the second paragraph of the quotation.  It
   is just as dull as the first paragraph.
\end{quotation}
Another frequently-displayed structure is a list.
The following is an example of an \emph{itemized}
list.
\begin{itemize}
   \item This is the first item of an itemized list.
         Each item in the list is marked with a ``tick''.
         You don't have to worry about what kind of tick
         mark is used.

   \item This is the second item of the list.  It
         contains another list nested inside it.  The inner
         list is an \emph{enumerated} list.
         \begin{enumerate}
            \item This is the first item of an enumerated
                  list that is nested within the itemized list.

            \item This is the second item of the inner list.
                  \LaTeX\ allows you to nest lists deeper than
                  you really should.
         \end{enumerate}
         This is the rest of the second item of the outer
         list.  It is no more interesting than any other
         part of the item.
   \item This is the third item of the list.
\end{itemize}
You can even display poetry.
\begin{verse}
   There is an environment
    for verse \\             % The \\ command separates lines
   Whose features some poets % within a stanza.
   will curse.

                             % One or more blank lines separate stanzas.

   For instead of making\\
   Them do \emph{all} line breaking, \\
   It allows them to put too many words on a line when they'd rather be
   forced to be terse.
\end{verse}

Mathematical formulas may also be displayed.  A
displayed formula
is
one-line long; multiline
formulas require special formatting instructions.
   \[  \ip{\Gamma}{\psi'} = x'' + y^{2} + z_{i}^{n}\]
Don't start a paragraph with a displayed equation,
nor make one a paragraph by itself.

\end{document}               % End of document.

\section{Bench Huge Section 1138}
% synthetic bench section 1138
% This is a sample LaTeX input file.  (Version of 12 August 2004.)
%
% A '%' character causes TeX to ignore all remaining text on the line,
% and is used for comments like this one.

\documentclass{article}      % Specifies the document class

                             % The preamble begins here.
\title{An Example Document}  % Declares the document's title.
\author{Leslie Lamport}      % Declares the author's name.
\date{January 21, 1994}      % Deleting this command produces today's date.

\newcommand{\ip}[2]{(#1, #2)}
                             % Defines \ip{arg1}{arg2} to mean
                             % (arg1, arg2).

%\newcommand{\ip}[2]{\langle #1 | #2\rangle}
                             % This is an alternative definition of
                             % \ip that is commented out.

\begin{document}             % End of preamble and beginning of text.

\maketitle                   % Produces the title.

This is an example input file.  Comparing it with
the output it generates can show you how to
produce a simple document of your own.

\section{Ordinary Text}      % Produces section heading.  Lower-level
                             % sections are begun with similar
                             % \subsection and \subsubsection commands.

The ends  of words and sentences are marked
  by   spaces. It  doesn't matter how many
spaces    you type; one is as good as 100.  The
end of   a line counts as a space.

One   or more   blank lines denote the  end
of  a paragraph.

Since any number of consecutive spaces are treated
like a single one, the formatting of the input
file makes no difference to
      \LaTeX,                % The \LaTeX command generates the LaTeX logo.
but it makes a difference to you.  When you use
\LaTeX, making your input file as easy to read
as possible will be a great help as you write
your document and when you change it.  This sample
file shows how you can add comments to your own input
file.

Because printing is different from typewriting,
there are a number of things that you have to do
differently when preparing an input file than if
you were just typing the document directly.
Quotation marks like
       ``this''
have to be handled specially, as do quotes within
quotes:
       ``\,`this'            % \, separates the double and single quote.
        is what I just
        wrote, not  `that'\,''.

Dashes come in three sizes: an
       intra-word
dash, a medium dash for number ranges like
       1--2,
and a punctuation
       dash---like
this.

A sentence-ending space should be larger than the
space between words within a sentence.  You
sometimes have to type special commands in
conjunction with punctuation characters to get
this right, as in the following sentence.
       Gnats, gnus, etc.\ all  % `\ ' makes an inter-word space.
       begin with G\@.         % \@ marks end-of-sentence punctuation.
You should check the spaces after periods when
reading your output to make sure you haven't
forgotten any special cases.  Generating an
ellipsis
       \ldots\               % `\ ' is needed after `\ldots' because TeX
                             % ignores spaces after command names like \ldots
                             % made from \ + letters.
                             %
                             % Note how a `%' character causes TeX to ignore
                             % the end of the input line, so these blank lines
                             % do not start a new paragraph.
                             %
with the right spacing around the periods requires
a special command.

\LaTeX\ interprets some common characters as
commands, so you must type special commands to
generate them.  These characters include the
following:
       \$ \& \% \# \{ and \}.

In printing, text is usually emphasized with an
       \emph{italic}
type style.

\begin{em}
   A long segment of text can also be emphasized
   in this way.  Text within such a segment can be
   given \emph{additional} emphasis.
\end{em}

It is sometimes necessary to prevent \LaTeX\ from
breaking a line where it might otherwise do so.
This may be at a space, as between the ``Mr.''\ and
``Jones'' in
       ``Mr.~Jones'',        % ~ produces an unbreakable interword space.
or within a word---especially when the word is a
symbol like
       \mbox{\emph{itemnum}}
that makes little sense when hyphenated across
lines.

Footnotes\footnote{This is an example of a footnote.}
pose no problem.

\LaTeX\ is good at typesetting mathematical formulas
like
       \( x-3y + z = 7 \)
or
       \( a_{1} > x^{2n} + y^{2n} > x' \)
or
       \( \ip{A}{B} = \sum_{i} a_{i} b_{i} \).
The spaces you type in a formula are
ignored.  Remember that a letter like
       $x$                   % $ ... $  and  \( ... \)  are equivalent
is a formula when it denotes a mathematical
symbol, and it should be typed as one.

\section{Displayed Text}

Text is displayed by indenting it from the left
margin.  Quotations are commonly displayed.  There
are short quotations
\begin{quote}
   This is a short quotation.  It consists of a
   single paragraph of text.  See how it is formatted.
\end{quote}
and longer ones.
\begin{quotation}
   This is a longer quotation.  It consists of two
   paragraphs of text, neither of which are
   particularly interesting.

   This is the second paragraph of the quotation.  It
   is just as dull as the first paragraph.
\end{quotation}
Another frequently-displayed structure is a list.
The following is an example of an \emph{itemized}
list.
\begin{itemize}
   \item This is the first item of an itemized list.
         Each item in the list is marked with a ``tick''.
         You don't have to worry about what kind of tick
         mark is used.

   \item This is the second item of the list.  It
         contains another list nested inside it.  The inner
         list is an \emph{enumerated} list.
         \begin{enumerate}
            \item This is the first item of an enumerated
                  list that is nested within the itemized list.

            \item This is the second item of the inner list.
                  \LaTeX\ allows you to nest lists deeper than
                  you really should.
         \end{enumerate}
         This is the rest of the second item of the outer
         list.  It is no more interesting than any other
         part of the item.
   \item This is the third item of the list.
\end{itemize}
You can even display poetry.
\begin{verse}
   There is an environment
    for verse \\             % The \\ command separates lines
   Whose features some poets % within a stanza.
   will curse.

                             % One or more blank lines separate stanzas.

   For instead of making\\
   Them do \emph{all} line breaking, \\
   It allows them to put too many words on a line when they'd rather be
   forced to be terse.
\end{verse}

Mathematical formulas may also be displayed.  A
displayed formula
is
one-line long; multiline
formulas require special formatting instructions.
   \[  \ip{\Gamma}{\psi'} = x'' + y^{2} + z_{i}^{n}\]
Don't start a paragraph with a displayed equation,
nor make one a paragraph by itself.

\end{document}               % End of document.

\section{Bench Huge Section 1139}
% synthetic bench section 1139
% This is a sample LaTeX input file.  (Version of 12 August 2004.)
%
% A '%' character causes TeX to ignore all remaining text on the line,
% and is used for comments like this one.

\documentclass{article}      % Specifies the document class

                             % The preamble begins here.
\title{An Example Document}  % Declares the document's title.
\author{Leslie Lamport}      % Declares the author's name.
\date{January 21, 1994}      % Deleting this command produces today's date.

\newcommand{\ip}[2]{(#1, #2)}
                             % Defines \ip{arg1}{arg2} to mean
                             % (arg1, arg2).

%\newcommand{\ip}[2]{\langle #1 | #2\rangle}
                             % This is an alternative definition of
                             % \ip that is commented out.

\begin{document}             % End of preamble and beginning of text.

\maketitle                   % Produces the title.

This is an example input file.  Comparing it with
the output it generates can show you how to
produce a simple document of your own.

\section{Ordinary Text}      % Produces section heading.  Lower-level
                             % sections are begun with similar
                             % \subsection and \subsubsection commands.

The ends  of words and sentences are marked
  by   spaces. It  doesn't matter how many
spaces    you type; one is as good as 100.  The
end of   a line counts as a space.

One   or more   blank lines denote the  end
of  a paragraph.

Since any number of consecutive spaces are treated
like a single one, the formatting of the input
file makes no difference to
      \LaTeX,                % The \LaTeX command generates the LaTeX logo.
but it makes a difference to you.  When you use
\LaTeX, making your input file as easy to read
as possible will be a great help as you write
your document and when you change it.  This sample
file shows how you can add comments to your own input
file.

Because printing is different from typewriting,
there are a number of things that you have to do
differently when preparing an input file than if
you were just typing the document directly.
Quotation marks like
       ``this''
have to be handled specially, as do quotes within
quotes:
       ``\,`this'            % \, separates the double and single quote.
        is what I just
        wrote, not  `that'\,''.

Dashes come in three sizes: an
       intra-word
dash, a medium dash for number ranges like
       1--2,
and a punctuation
       dash---like
this.

A sentence-ending space should be larger than the
space between words within a sentence.  You
sometimes have to type special commands in
conjunction with punctuation characters to get
this right, as in the following sentence.
       Gnats, gnus, etc.\ all  % `\ ' makes an inter-word space.
       begin with G\@.         % \@ marks end-of-sentence punctuation.
You should check the spaces after periods when
reading your output to make sure you haven't
forgotten any special cases.  Generating an
ellipsis
       \ldots\               % `\ ' is needed after `\ldots' because TeX
                             % ignores spaces after command names like \ldots
                             % made from \ + letters.
                             %
                             % Note how a `%' character causes TeX to ignore
                             % the end of the input line, so these blank lines
                             % do not start a new paragraph.
                             %
with the right spacing around the periods requires
a special command.

\LaTeX\ interprets some common characters as
commands, so you must type special commands to
generate them.  These characters include the
following:
       \$ \& \% \# \{ and \}.

In printing, text is usually emphasized with an
       \emph{italic}
type style.

\begin{em}
   A long segment of text can also be emphasized
   in this way.  Text within such a segment can be
   given \emph{additional} emphasis.
\end{em}

It is sometimes necessary to prevent \LaTeX\ from
breaking a line where it might otherwise do so.
This may be at a space, as between the ``Mr.''\ and
``Jones'' in
       ``Mr.~Jones'',        % ~ produces an unbreakable interword space.
or within a word---especially when the word is a
symbol like
       \mbox{\emph{itemnum}}
that makes little sense when hyphenated across
lines.

Footnotes\footnote{This is an example of a footnote.}
pose no problem.

\LaTeX\ is good at typesetting mathematical formulas
like
       \( x-3y + z = 7 \)
or
       \( a_{1} > x^{2n} + y^{2n} > x' \)
or
       \( \ip{A}{B} = \sum_{i} a_{i} b_{i} \).
The spaces you type in a formula are
ignored.  Remember that a letter like
       $x$                   % $ ... $  and  \( ... \)  are equivalent
is a formula when it denotes a mathematical
symbol, and it should be typed as one.

\section{Displayed Text}

Text is displayed by indenting it from the left
margin.  Quotations are commonly displayed.  There
are short quotations
\begin{quote}
   This is a short quotation.  It consists of a
   single paragraph of text.  See how it is formatted.
\end{quote}
and longer ones.
\begin{quotation}
   This is a longer quotation.  It consists of two
   paragraphs of text, neither of which are
   particularly interesting.

   This is the second paragraph of the quotation.  It
   is just as dull as the first paragraph.
\end{quotation}
Another frequently-displayed structure is a list.
The following is an example of an \emph{itemized}
list.
\begin{itemize}
   \item This is the first item of an itemized list.
         Each item in the list is marked with a ``tick''.
         You don't have to worry about what kind of tick
         mark is used.

   \item This is the second item of the list.  It
         contains another list nested inside it.  The inner
         list is an \emph{enumerated} list.
         \begin{enumerate}
            \item This is the first item of an enumerated
                  list that is nested within the itemized list.

            \item This is the second item of the inner list.
                  \LaTeX\ allows you to nest lists deeper than
                  you really should.
         \end{enumerate}
         This is the rest of the second item of the outer
         list.  It is no more interesting than any other
         part of the item.
   \item This is the third item of the list.
\end{itemize}
You can even display poetry.
\begin{verse}
   There is an environment
    for verse \\             % The \\ command separates lines
   Whose features some poets % within a stanza.
   will curse.

                             % One or more blank lines separate stanzas.

   For instead of making\\
   Them do \emph{all} line breaking, \\
   It allows them to put too many words on a line when they'd rather be
   forced to be terse.
\end{verse}

Mathematical formulas may also be displayed.  A
displayed formula
is
one-line long; multiline
formulas require special formatting instructions.
   \[  \ip{\Gamma}{\psi'} = x'' + y^{2} + z_{i}^{n}\]
Don't start a paragraph with a displayed equation,
nor make one a paragraph by itself.

\end{document}               % End of document.

\section{Bench Huge Section 1140}
% synthetic bench section 1140
% This is a sample LaTeX input file.  (Version of 12 August 2004.)
%
% A '%' character causes TeX to ignore all remaining text on the line,
% and is used for comments like this one.

\documentclass{article}      % Specifies the document class

                             % The preamble begins here.
\title{An Example Document}  % Declares the document's title.
\author{Leslie Lamport}      % Declares the author's name.
\date{January 21, 1994}      % Deleting this command produces today's date.

\newcommand{\ip}[2]{(#1, #2)}
                             % Defines \ip{arg1}{arg2} to mean
                             % (arg1, arg2).

%\newcommand{\ip}[2]{\langle #1 | #2\rangle}
                             % This is an alternative definition of
                             % \ip that is commented out.

\begin{document}             % End of preamble and beginning of text.

\maketitle                   % Produces the title.

This is an example input file.  Comparing it with
the output it generates can show you how to
produce a simple document of your own.

\section{Ordinary Text}      % Produces section heading.  Lower-level
                             % sections are begun with similar
                             % \subsection and \subsubsection commands.

The ends  of words and sentences are marked
  by   spaces. It  doesn't matter how many
spaces    you type; one is as good as 100.  The
end of   a line counts as a space.

One   or more   blank lines denote the  end
of  a paragraph.

Since any number of consecutive spaces are treated
like a single one, the formatting of the input
file makes no difference to
      \LaTeX,                % The \LaTeX command generates the LaTeX logo.
but it makes a difference to you.  When you use
\LaTeX, making your input file as easy to read
as possible will be a great help as you write
your document and when you change it.  This sample
file shows how you can add comments to your own input
file.

Because printing is different from typewriting,
there are a number of things that you have to do
differently when preparing an input file than if
you were just typing the document directly.
Quotation marks like
       ``this''
have to be handled specially, as do quotes within
quotes:
       ``\,`this'            % \, separates the double and single quote.
        is what I just
        wrote, not  `that'\,''.

Dashes come in three sizes: an
       intra-word
dash, a medium dash for number ranges like
       1--2,
and a punctuation
       dash---like
this.

A sentence-ending space should be larger than the
space between words within a sentence.  You
sometimes have to type special commands in
conjunction with punctuation characters to get
this right, as in the following sentence.
       Gnats, gnus, etc.\ all  % `\ ' makes an inter-word space.
       begin with G\@.         % \@ marks end-of-sentence punctuation.
You should check the spaces after periods when
reading your output to make sure you haven't
forgotten any special cases.  Generating an
ellipsis
       \ldots\               % `\ ' is needed after `\ldots' because TeX
                             % ignores spaces after command names like \ldots
                             % made from \ + letters.
                             %
                             % Note how a `%' character causes TeX to ignore
                             % the end of the input line, so these blank lines
                             % do not start a new paragraph.
                             %
with the right spacing around the periods requires
a special command.

\LaTeX\ interprets some common characters as
commands, so you must type special commands to
generate them.  These characters include the
following:
       \$ \& \% \# \{ and \}.

In printing, text is usually emphasized with an
       \emph{italic}
type style.

\begin{em}
   A long segment of text can also be emphasized
   in this way.  Text within such a segment can be
   given \emph{additional} emphasis.
\end{em}

It is sometimes necessary to prevent \LaTeX\ from
breaking a line where it might otherwise do so.
This may be at a space, as between the ``Mr.''\ and
``Jones'' in
       ``Mr.~Jones'',        % ~ produces an unbreakable interword space.
or within a word---especially when the word is a
symbol like
       \mbox{\emph{itemnum}}
that makes little sense when hyphenated across
lines.

Footnotes\footnote{This is an example of a footnote.}
pose no problem.

\LaTeX\ is good at typesetting mathematical formulas
like
       \( x-3y + z = 7 \)
or
       \( a_{1} > x^{2n} + y^{2n} > x' \)
or
       \( \ip{A}{B} = \sum_{i} a_{i} b_{i} \).
The spaces you type in a formula are
ignored.  Remember that a letter like
       $x$                   % $ ... $  and  \( ... \)  are equivalent
is a formula when it denotes a mathematical
symbol, and it should be typed as one.

\section{Displayed Text}

Text is displayed by indenting it from the left
margin.  Quotations are commonly displayed.  There
are short quotations
\begin{quote}
   This is a short quotation.  It consists of a
   single paragraph of text.  See how it is formatted.
\end{quote}
and longer ones.
\begin{quotation}
   This is a longer quotation.  It consists of two
   paragraphs of text, neither of which are
   particularly interesting.

   This is the second paragraph of the quotation.  It
   is just as dull as the first paragraph.
\end{quotation}
Another frequently-displayed structure is a list.
The following is an example of an \emph{itemized}
list.
\begin{itemize}
   \item This is the first item of an itemized list.
         Each item in the list is marked with a ``tick''.
         You don't have to worry about what kind of tick
         mark is used.

   \item This is the second item of the list.  It
         contains another list nested inside it.  The inner
         list is an \emph{enumerated} list.
         \begin{enumerate}
            \item This is the first item of an enumerated
                  list that is nested within the itemized list.

            \item This is the second item of the inner list.
                  \LaTeX\ allows you to nest lists deeper than
                  you really should.
         \end{enumerate}
         This is the rest of the second item of the outer
         list.  It is no more interesting than any other
         part of the item.
   \item This is the third item of the list.
\end{itemize}
You can even display poetry.
\begin{verse}
   There is an environment
    for verse \\             % The \\ command separates lines
   Whose features some poets % within a stanza.
   will curse.

                             % One or more blank lines separate stanzas.

   For instead of making\\
   Them do \emph{all} line breaking, \\
   It allows them to put too many words on a line when they'd rather be
   forced to be terse.
\end{verse}

Mathematical formulas may also be displayed.  A
displayed formula
is
one-line long; multiline
formulas require special formatting instructions.
   \[  \ip{\Gamma}{\psi'} = x'' + y^{2} + z_{i}^{n}\]
Don't start a paragraph with a displayed equation,
nor make one a paragraph by itself.

\end{document}               % End of document.

\section{Bench Huge Section 1141}
% synthetic bench section 1141
% This is a sample LaTeX input file.  (Version of 12 August 2004.)
%
% A '%' character causes TeX to ignore all remaining text on the line,
% and is used for comments like this one.

\documentclass{article}      % Specifies the document class

                             % The preamble begins here.
\title{An Example Document}  % Declares the document's title.
\author{Leslie Lamport}      % Declares the author's name.
\date{January 21, 1994}      % Deleting this command produces today's date.

\newcommand{\ip}[2]{(#1, #2)}
                             % Defines \ip{arg1}{arg2} to mean
                             % (arg1, arg2).

%\newcommand{\ip}[2]{\langle #1 | #2\rangle}
                             % This is an alternative definition of
                             % \ip that is commented out.

\begin{document}             % End of preamble and beginning of text.

\maketitle                   % Produces the title.

This is an example input file.  Comparing it with
the output it generates can show you how to
produce a simple document of your own.

\section{Ordinary Text}      % Produces section heading.  Lower-level
                             % sections are begun with similar
                             % \subsection and \subsubsection commands.

The ends  of words and sentences are marked
  by   spaces. It  doesn't matter how many
spaces    you type; one is as good as 100.  The
end of   a line counts as a space.

One   or more   blank lines denote the  end
of  a paragraph.

Since any number of consecutive spaces are treated
like a single one, the formatting of the input
file makes no difference to
      \LaTeX,                % The \LaTeX command generates the LaTeX logo.
but it makes a difference to you.  When you use
\LaTeX, making your input file as easy to read
as possible will be a great help as you write
your document and when you change it.  This sample
file shows how you can add comments to your own input
file.

Because printing is different from typewriting,
there are a number of things that you have to do
differently when preparing an input file than if
you were just typing the document directly.
Quotation marks like
       ``this''
have to be handled specially, as do quotes within
quotes:
       ``\,`this'            % \, separates the double and single quote.
        is what I just
        wrote, not  `that'\,''.

Dashes come in three sizes: an
       intra-word
dash, a medium dash for number ranges like
       1--2,
and a punctuation
       dash---like
this.

A sentence-ending space should be larger than the
space between words within a sentence.  You
sometimes have to type special commands in
conjunction with punctuation characters to get
this right, as in the following sentence.
       Gnats, gnus, etc.\ all  % `\ ' makes an inter-word space.
       begin with G\@.         % \@ marks end-of-sentence punctuation.
You should check the spaces after periods when
reading your output to make sure you haven't
forgotten any special cases.  Generating an
ellipsis
       \ldots\               % `\ ' is needed after `\ldots' because TeX
                             % ignores spaces after command names like \ldots
                             % made from \ + letters.
                             %
                             % Note how a `%' character causes TeX to ignore
                             % the end of the input line, so these blank lines
                             % do not start a new paragraph.
                             %
with the right spacing around the periods requires
a special command.

\LaTeX\ interprets some common characters as
commands, so you must type special commands to
generate them.  These characters include the
following:
       \$ \& \% \# \{ and \}.

In printing, text is usually emphasized with an
       \emph{italic}
type style.

\begin{em}
   A long segment of text can also be emphasized
   in this way.  Text within such a segment can be
   given \emph{additional} emphasis.
\end{em}

It is sometimes necessary to prevent \LaTeX\ from
breaking a line where it might otherwise do so.
This may be at a space, as between the ``Mr.''\ and
``Jones'' in
       ``Mr.~Jones'',        % ~ produces an unbreakable interword space.
or within a word---especially when the word is a
symbol like
       \mbox{\emph{itemnum}}
that makes little sense when hyphenated across
lines.

Footnotes\footnote{This is an example of a footnote.}
pose no problem.

\LaTeX\ is good at typesetting mathematical formulas
like
       \( x-3y + z = 7 \)
or
       \( a_{1} > x^{2n} + y^{2n} > x' \)
or
       \( \ip{A}{B} = \sum_{i} a_{i} b_{i} \).
The spaces you type in a formula are
ignored.  Remember that a letter like
       $x$                   % $ ... $  and  \( ... \)  are equivalent
is a formula when it denotes a mathematical
symbol, and it should be typed as one.

\section{Displayed Text}

Text is displayed by indenting it from the left
margin.  Quotations are commonly displayed.  There
are short quotations
\begin{quote}
   This is a short quotation.  It consists of a
   single paragraph of text.  See how it is formatted.
\end{quote}
and longer ones.
\begin{quotation}
   This is a longer quotation.  It consists of two
   paragraphs of text, neither of which are
   particularly interesting.

   This is the second paragraph of the quotation.  It
   is just as dull as the first paragraph.
\end{quotation}
Another frequently-displayed structure is a list.
The following is an example of an \emph{itemized}
list.
\begin{itemize}
   \item This is the first item of an itemized list.
         Each item in the list is marked with a ``tick''.
         You don't have to worry about what kind of tick
         mark is used.

   \item This is the second item of the list.  It
         contains another list nested inside it.  The inner
         list is an \emph{enumerated} list.
         \begin{enumerate}
            \item This is the first item of an enumerated
                  list that is nested within the itemized list.

            \item This is the second item of the inner list.
                  \LaTeX\ allows you to nest lists deeper than
                  you really should.
         \end{enumerate}
         This is the rest of the second item of the outer
         list.  It is no more interesting than any other
         part of the item.
   \item This is the third item of the list.
\end{itemize}
You can even display poetry.
\begin{verse}
   There is an environment
    for verse \\             % The \\ command separates lines
   Whose features some poets % within a stanza.
   will curse.

                             % One or more blank lines separate stanzas.

   For instead of making\\
   Them do \emph{all} line breaking, \\
   It allows them to put too many words on a line when they'd rather be
   forced to be terse.
\end{verse}

Mathematical formulas may also be displayed.  A
displayed formula
is
one-line long; multiline
formulas require special formatting instructions.
   \[  \ip{\Gamma}{\psi'} = x'' + y^{2} + z_{i}^{n}\]
Don't start a paragraph with a displayed equation,
nor make one a paragraph by itself.

\end{document}               % End of document.

\section{Bench Huge Section 1142}
% synthetic bench section 1142
% This is a sample LaTeX input file.  (Version of 12 August 2004.)
%
% A '%' character causes TeX to ignore all remaining text on the line,
% and is used for comments like this one.

\documentclass{article}      % Specifies the document class

                             % The preamble begins here.
\title{An Example Document}  % Declares the document's title.
\author{Leslie Lamport}      % Declares the author's name.
\date{January 21, 1994}      % Deleting this command produces today's date.

\newcommand{\ip}[2]{(#1, #2)}
                             % Defines \ip{arg1}{arg2} to mean
                             % (arg1, arg2).

%\newcommand{\ip}[2]{\langle #1 | #2\rangle}
                             % This is an alternative definition of
                             % \ip that is commented out.

\begin{document}             % End of preamble and beginning of text.

\maketitle                   % Produces the title.

This is an example input file.  Comparing it with
the output it generates can show you how to
produce a simple document of your own.

\section{Ordinary Text}      % Produces section heading.  Lower-level
                             % sections are begun with similar
                             % \subsection and \subsubsection commands.

The ends  of words and sentences are marked
  by   spaces. It  doesn't matter how many
spaces    you type; one is as good as 100.  The
end of   a line counts as a space.

One   or more   blank lines denote the  end
of  a paragraph.

Since any number of consecutive spaces are treated
like a single one, the formatting of the input
file makes no difference to
      \LaTeX,                % The \LaTeX command generates the LaTeX logo.
but it makes a difference to you.  When you use
\LaTeX, making your input file as easy to read
as possible will be a great help as you write
your document and when you change it.  This sample
file shows how you can add comments to your own input
file.

Because printing is different from typewriting,
there are a number of things that you have to do
differently when preparing an input file than if
you were just typing the document directly.
Quotation marks like
       ``this''
have to be handled specially, as do quotes within
quotes:
       ``\,`this'            % \, separates the double and single quote.
        is what I just
        wrote, not  `that'\,''.

Dashes come in three sizes: an
       intra-word
dash, a medium dash for number ranges like
       1--2,
and a punctuation
       dash---like
this.

A sentence-ending space should be larger than the
space between words within a sentence.  You
sometimes have to type special commands in
conjunction with punctuation characters to get
this right, as in the following sentence.
       Gnats, gnus, etc.\ all  % `\ ' makes an inter-word space.
       begin with G\@.         % \@ marks end-of-sentence punctuation.
You should check the spaces after periods when
reading your output to make sure you haven't
forgotten any special cases.  Generating an
ellipsis
       \ldots\               % `\ ' is needed after `\ldots' because TeX
                             % ignores spaces after command names like \ldots
                             % made from \ + letters.
                             %
                             % Note how a `%' character causes TeX to ignore
                             % the end of the input line, so these blank lines
                             % do not start a new paragraph.
                             %
with the right spacing around the periods requires
a special command.

\LaTeX\ interprets some common characters as
commands, so you must type special commands to
generate them.  These characters include the
following:
       \$ \& \% \# \{ and \}.

In printing, text is usually emphasized with an
       \emph{italic}
type style.

\begin{em}
   A long segment of text can also be emphasized
   in this way.  Text within such a segment can be
   given \emph{additional} emphasis.
\end{em}

It is sometimes necessary to prevent \LaTeX\ from
breaking a line where it might otherwise do so.
This may be at a space, as between the ``Mr.''\ and
``Jones'' in
       ``Mr.~Jones'',        % ~ produces an unbreakable interword space.
or within a word---especially when the word is a
symbol like
       \mbox{\emph{itemnum}}
that makes little sense when hyphenated across
lines.

Footnotes\footnote{This is an example of a footnote.}
pose no problem.

\LaTeX\ is good at typesetting mathematical formulas
like
       \( x-3y + z = 7 \)
or
       \( a_{1} > x^{2n} + y^{2n} > x' \)
or
       \( \ip{A}{B} = \sum_{i} a_{i} b_{i} \).
The spaces you type in a formula are
ignored.  Remember that a letter like
       $x$                   % $ ... $  and  \( ... \)  are equivalent
is a formula when it denotes a mathematical
symbol, and it should be typed as one.

\section{Displayed Text}

Text is displayed by indenting it from the left
margin.  Quotations are commonly displayed.  There
are short quotations
\begin{quote}
   This is a short quotation.  It consists of a
   single paragraph of text.  See how it is formatted.
\end{quote}
and longer ones.
\begin{quotation}
   This is a longer quotation.  It consists of two
   paragraphs of text, neither of which are
   particularly interesting.

   This is the second paragraph of the quotation.  It
   is just as dull as the first paragraph.
\end{quotation}
Another frequently-displayed structure is a list.
The following is an example of an \emph{itemized}
list.
\begin{itemize}
   \item This is the first item of an itemized list.
         Each item in the list is marked with a ``tick''.
         You don't have to worry about what kind of tick
         mark is used.

   \item This is the second item of the list.  It
         contains another list nested inside it.  The inner
         list is an \emph{enumerated} list.
         \begin{enumerate}
            \item This is the first item of an enumerated
                  list that is nested within the itemized list.

            \item This is the second item of the inner list.
                  \LaTeX\ allows you to nest lists deeper than
                  you really should.
         \end{enumerate}
         This is the rest of the second item of the outer
         list.  It is no more interesting than any other
         part of the item.
   \item This is the third item of the list.
\end{itemize}
You can even display poetry.
\begin{verse}
   There is an environment
    for verse \\             % The \\ command separates lines
   Whose features some poets % within a stanza.
   will curse.

                             % One or more blank lines separate stanzas.

   For instead of making\\
   Them do \emph{all} line breaking, \\
   It allows them to put too many words on a line when they'd rather be
   forced to be terse.
\end{verse}

Mathematical formulas may also be displayed.  A
displayed formula
is
one-line long; multiline
formulas require special formatting instructions.
   \[  \ip{\Gamma}{\psi'} = x'' + y^{2} + z_{i}^{n}\]
Don't start a paragraph with a displayed equation,
nor make one a paragraph by itself.

\end{document}               % End of document.

\section{Bench Huge Section 1143}
% synthetic bench section 1143
% This is a sample LaTeX input file.  (Version of 12 August 2004.)
%
% A '%' character causes TeX to ignore all remaining text on the line,
% and is used for comments like this one.

\documentclass{article}      % Specifies the document class

                             % The preamble begins here.
\title{An Example Document}  % Declares the document's title.
\author{Leslie Lamport}      % Declares the author's name.
\date{January 21, 1994}      % Deleting this command produces today's date.

\newcommand{\ip}[2]{(#1, #2)}
                             % Defines \ip{arg1}{arg2} to mean
                             % (arg1, arg2).

%\newcommand{\ip}[2]{\langle #1 | #2\rangle}
                             % This is an alternative definition of
                             % \ip that is commented out.

\begin{document}             % End of preamble and beginning of text.

\maketitle                   % Produces the title.

This is an example input file.  Comparing it with
the output it generates can show you how to
produce a simple document of your own.

\section{Ordinary Text}      % Produces section heading.  Lower-level
                             % sections are begun with similar
                             % \subsection and \subsubsection commands.

The ends  of words and sentences are marked
  by   spaces. It  doesn't matter how many
spaces    you type; one is as good as 100.  The
end of   a line counts as a space.

One   or more   blank lines denote the  end
of  a paragraph.

Since any number of consecutive spaces are treated
like a single one, the formatting of the input
file makes no difference to
      \LaTeX,                % The \LaTeX command generates the LaTeX logo.
but it makes a difference to you.  When you use
\LaTeX, making your input file as easy to read
as possible will be a great help as you write
your document and when you change it.  This sample
file shows how you can add comments to your own input
file.

Because printing is different from typewriting,
there are a number of things that you have to do
differently when preparing an input file than if
you were just typing the document directly.
Quotation marks like
       ``this''
have to be handled specially, as do quotes within
quotes:
       ``\,`this'            % \, separates the double and single quote.
        is what I just
        wrote, not  `that'\,''.

Dashes come in three sizes: an
       intra-word
dash, a medium dash for number ranges like
       1--2,
and a punctuation
       dash---like
this.

A sentence-ending space should be larger than the
space between words within a sentence.  You
sometimes have to type special commands in
conjunction with punctuation characters to get
this right, as in the following sentence.
       Gnats, gnus, etc.\ all  % `\ ' makes an inter-word space.
       begin with G\@.         % \@ marks end-of-sentence punctuation.
You should check the spaces after periods when
reading your output to make sure you haven't
forgotten any special cases.  Generating an
ellipsis
       \ldots\               % `\ ' is needed after `\ldots' because TeX
                             % ignores spaces after command names like \ldots
                             % made from \ + letters.
                             %
                             % Note how a `%' character causes TeX to ignore
                             % the end of the input line, so these blank lines
                             % do not start a new paragraph.
                             %
with the right spacing around the periods requires
a special command.

\LaTeX\ interprets some common characters as
commands, so you must type special commands to
generate them.  These characters include the
following:
       \$ \& \% \# \{ and \}.

In printing, text is usually emphasized with an
       \emph{italic}
type style.

\begin{em}
   A long segment of text can also be emphasized
   in this way.  Text within such a segment can be
   given \emph{additional} emphasis.
\end{em}

It is sometimes necessary to prevent \LaTeX\ from
breaking a line where it might otherwise do so.
This may be at a space, as between the ``Mr.''\ and
``Jones'' in
       ``Mr.~Jones'',        % ~ produces an unbreakable interword space.
or within a word---especially when the word is a
symbol like
       \mbox{\emph{itemnum}}
that makes little sense when hyphenated across
lines.

Footnotes\footnote{This is an example of a footnote.}
pose no problem.

\LaTeX\ is good at typesetting mathematical formulas
like
       \( x-3y + z = 7 \)
or
       \( a_{1} > x^{2n} + y^{2n} > x' \)
or
       \( \ip{A}{B} = \sum_{i} a_{i} b_{i} \).
The spaces you type in a formula are
ignored.  Remember that a letter like
       $x$                   % $ ... $  and  \( ... \)  are equivalent
is a formula when it denotes a mathematical
symbol, and it should be typed as one.

\section{Displayed Text}

Text is displayed by indenting it from the left
margin.  Quotations are commonly displayed.  There
are short quotations
\begin{quote}
   This is a short quotation.  It consists of a
   single paragraph of text.  See how it is formatted.
\end{quote}
and longer ones.
\begin{quotation}
   This is a longer quotation.  It consists of two
   paragraphs of text, neither of which are
   particularly interesting.

   This is the second paragraph of the quotation.  It
   is just as dull as the first paragraph.
\end{quotation}
Another frequently-displayed structure is a list.
The following is an example of an \emph{itemized}
list.
\begin{itemize}
   \item This is the first item of an itemized list.
         Each item in the list is marked with a ``tick''.
         You don't have to worry about what kind of tick
         mark is used.

   \item This is the second item of the list.  It
         contains another list nested inside it.  The inner
         list is an \emph{enumerated} list.
         \begin{enumerate}
            \item This is the first item of an enumerated
                  list that is nested within the itemized list.

            \item This is the second item of the inner list.
                  \LaTeX\ allows you to nest lists deeper than
                  you really should.
         \end{enumerate}
         This is the rest of the second item of the outer
         list.  It is no more interesting than any other
         part of the item.
   \item This is the third item of the list.
\end{itemize}
You can even display poetry.
\begin{verse}
   There is an environment
    for verse \\             % The \\ command separates lines
   Whose features some poets % within a stanza.
   will curse.

                             % One or more blank lines separate stanzas.

   For instead of making\\
   Them do \emph{all} line breaking, \\
   It allows them to put too many words on a line when they'd rather be
   forced to be terse.
\end{verse}

Mathematical formulas may also be displayed.  A
displayed formula
is
one-line long; multiline
formulas require special formatting instructions.
   \[  \ip{\Gamma}{\psi'} = x'' + y^{2} + z_{i}^{n}\]
Don't start a paragraph with a displayed equation,
nor make one a paragraph by itself.

\end{document}               % End of document.

\section{Bench Huge Section 1144}
% synthetic bench section 1144
% This is a sample LaTeX input file.  (Version of 12 August 2004.)
%
% A '%' character causes TeX to ignore all remaining text on the line,
% and is used for comments like this one.

\documentclass{article}      % Specifies the document class

                             % The preamble begins here.
\title{An Example Document}  % Declares the document's title.
\author{Leslie Lamport}      % Declares the author's name.
\date{January 21, 1994}      % Deleting this command produces today's date.

\newcommand{\ip}[2]{(#1, #2)}
                             % Defines \ip{arg1}{arg2} to mean
                             % (arg1, arg2).

%\newcommand{\ip}[2]{\langle #1 | #2\rangle}
                             % This is an alternative definition of
                             % \ip that is commented out.

\begin{document}             % End of preamble and beginning of text.

\maketitle                   % Produces the title.

This is an example input file.  Comparing it with
the output it generates can show you how to
produce a simple document of your own.

\section{Ordinary Text}      % Produces section heading.  Lower-level
                             % sections are begun with similar
                             % \subsection and \subsubsection commands.

The ends  of words and sentences are marked
  by   spaces. It  doesn't matter how many
spaces    you type; one is as good as 100.  The
end of   a line counts as a space.

One   or more   blank lines denote the  end
of  a paragraph.

Since any number of consecutive spaces are treated
like a single one, the formatting of the input
file makes no difference to
      \LaTeX,                % The \LaTeX command generates the LaTeX logo.
but it makes a difference to you.  When you use
\LaTeX, making your input file as easy to read
as possible will be a great help as you write
your document and when you change it.  This sample
file shows how you can add comments to your own input
file.

Because printing is different from typewriting,
there are a number of things that you have to do
differently when preparing an input file than if
you were just typing the document directly.
Quotation marks like
       ``this''
have to be handled specially, as do quotes within
quotes:
       ``\,`this'            % \, separates the double and single quote.
        is what I just
        wrote, not  `that'\,''.

Dashes come in three sizes: an
       intra-word
dash, a medium dash for number ranges like
       1--2,
and a punctuation
       dash---like
this.

A sentence-ending space should be larger than the
space between words within a sentence.  You
sometimes have to type special commands in
conjunction with punctuation characters to get
this right, as in the following sentence.
       Gnats, gnus, etc.\ all  % `\ ' makes an inter-word space.
       begin with G\@.         % \@ marks end-of-sentence punctuation.
You should check the spaces after periods when
reading your output to make sure you haven't
forgotten any special cases.  Generating an
ellipsis
       \ldots\               % `\ ' is needed after `\ldots' because TeX
                             % ignores spaces after command names like \ldots
                             % made from \ + letters.
                             %
                             % Note how a `%' character causes TeX to ignore
                             % the end of the input line, so these blank lines
                             % do not start a new paragraph.
                             %
with the right spacing around the periods requires
a special command.

\LaTeX\ interprets some common characters as
commands, so you must type special commands to
generate them.  These characters include the
following:
       \$ \& \% \# \{ and \}.

In printing, text is usually emphasized with an
       \emph{italic}
type style.

\begin{em}
   A long segment of text can also be emphasized
   in this way.  Text within such a segment can be
   given \emph{additional} emphasis.
\end{em}

It is sometimes necessary to prevent \LaTeX\ from
breaking a line where it might otherwise do so.
This may be at a space, as between the ``Mr.''\ and
``Jones'' in
       ``Mr.~Jones'',        % ~ produces an unbreakable interword space.
or within a word---especially when the word is a
symbol like
       \mbox{\emph{itemnum}}
that makes little sense when hyphenated across
lines.

Footnotes\footnote{This is an example of a footnote.}
pose no problem.

\LaTeX\ is good at typesetting mathematical formulas
like
       \( x-3y + z = 7 \)
or
       \( a_{1} > x^{2n} + y^{2n} > x' \)
or
       \( \ip{A}{B} = \sum_{i} a_{i} b_{i} \).
The spaces you type in a formula are
ignored.  Remember that a letter like
       $x$                   % $ ... $  and  \( ... \)  are equivalent
is a formula when it denotes a mathematical
symbol, and it should be typed as one.

\section{Displayed Text}

Text is displayed by indenting it from the left
margin.  Quotations are commonly displayed.  There
are short quotations
\begin{quote}
   This is a short quotation.  It consists of a
   single paragraph of text.  See how it is formatted.
\end{quote}
and longer ones.
\begin{quotation}
   This is a longer quotation.  It consists of two
   paragraphs of text, neither of which are
   particularly interesting.

   This is the second paragraph of the quotation.  It
   is just as dull as the first paragraph.
\end{quotation}
Another frequently-displayed structure is a list.
The following is an example of an \emph{itemized}
list.
\begin{itemize}
   \item This is the first item of an itemized list.
         Each item in the list is marked with a ``tick''.
         You don't have to worry about what kind of tick
         mark is used.

   \item This is the second item of the list.  It
         contains another list nested inside it.  The inner
         list is an \emph{enumerated} list.
         \begin{enumerate}
            \item This is the first item of an enumerated
                  list that is nested within the itemized list.

            \item This is the second item of the inner list.
                  \LaTeX\ allows you to nest lists deeper than
                  you really should.
         \end{enumerate}
         This is the rest of the second item of the outer
         list.  It is no more interesting than any other
         part of the item.
   \item This is the third item of the list.
\end{itemize}
You can even display poetry.
\begin{verse}
   There is an environment
    for verse \\             % The \\ command separates lines
   Whose features some poets % within a stanza.
   will curse.

                             % One or more blank lines separate stanzas.

   For instead of making\\
   Them do \emph{all} line breaking, \\
   It allows them to put too many words on a line when they'd rather be
   forced to be terse.
\end{verse}

Mathematical formulas may also be displayed.  A
displayed formula
is
one-line long; multiline
formulas require special formatting instructions.
   \[  \ip{\Gamma}{\psi'} = x'' + y^{2} + z_{i}^{n}\]
Don't start a paragraph with a displayed equation,
nor make one a paragraph by itself.

\end{document}               % End of document.

\section{Bench Huge Section 1145}
% synthetic bench section 1145
% This is a sample LaTeX input file.  (Version of 12 August 2004.)
%
% A '%' character causes TeX to ignore all remaining text on the line,
% and is used for comments like this one.

\documentclass{article}      % Specifies the document class

                             % The preamble begins here.
\title{An Example Document}  % Declares the document's title.
\author{Leslie Lamport}      % Declares the author's name.
\date{January 21, 1994}      % Deleting this command produces today's date.

\newcommand{\ip}[2]{(#1, #2)}
                             % Defines \ip{arg1}{arg2} to mean
                             % (arg1, arg2).

%\newcommand{\ip}[2]{\langle #1 | #2\rangle}
                             % This is an alternative definition of
                             % \ip that is commented out.

\begin{document}             % End of preamble and beginning of text.

\maketitle                   % Produces the title.

This is an example input file.  Comparing it with
the output it generates can show you how to
produce a simple document of your own.

\section{Ordinary Text}      % Produces section heading.  Lower-level
                             % sections are begun with similar
                             % \subsection and \subsubsection commands.

The ends  of words and sentences are marked
  by   spaces. It  doesn't matter how many
spaces    you type; one is as good as 100.  The
end of   a line counts as a space.

One   or more   blank lines denote the  end
of  a paragraph.

Since any number of consecutive spaces are treated
like a single one, the formatting of the input
file makes no difference to
      \LaTeX,                % The \LaTeX command generates the LaTeX logo.
but it makes a difference to you.  When you use
\LaTeX, making your input file as easy to read
as possible will be a great help as you write
your document and when you change it.  This sample
file shows how you can add comments to your own input
file.

Because printing is different from typewriting,
there are a number of things that you have to do
differently when preparing an input file than if
you were just typing the document directly.
Quotation marks like
       ``this''
have to be handled specially, as do quotes within
quotes:
       ``\,`this'            % \, separates the double and single quote.
        is what I just
        wrote, not  `that'\,''.

Dashes come in three sizes: an
       intra-word
dash, a medium dash for number ranges like
       1--2,
and a punctuation
       dash---like
this.

A sentence-ending space should be larger than the
space between words within a sentence.  You
sometimes have to type special commands in
conjunction with punctuation characters to get
this right, as in the following sentence.
       Gnats, gnus, etc.\ all  % `\ ' makes an inter-word space.
       begin with G\@.         % \@ marks end-of-sentence punctuation.
You should check the spaces after periods when
reading your output to make sure you haven't
forgotten any special cases.  Generating an
ellipsis
       \ldots\               % `\ ' is needed after `\ldots' because TeX
                             % ignores spaces after command names like \ldots
                             % made from \ + letters.
                             %
                             % Note how a `%' character causes TeX to ignore
                             % the end of the input line, so these blank lines
                             % do not start a new paragraph.
                             %
with the right spacing around the periods requires
a special command.

\LaTeX\ interprets some common characters as
commands, so you must type special commands to
generate them.  These characters include the
following:
       \$ \& \% \# \{ and \}.

In printing, text is usually emphasized with an
       \emph{italic}
type style.

\begin{em}
   A long segment of text can also be emphasized
   in this way.  Text within such a segment can be
   given \emph{additional} emphasis.
\end{em}

It is sometimes necessary to prevent \LaTeX\ from
breaking a line where it might otherwise do so.
This may be at a space, as between the ``Mr.''\ and
``Jones'' in
       ``Mr.~Jones'',        % ~ produces an unbreakable interword space.
or within a word---especially when the word is a
symbol like
       \mbox{\emph{itemnum}}
that makes little sense when hyphenated across
lines.

Footnotes\footnote{This is an example of a footnote.}
pose no problem.

\LaTeX\ is good at typesetting mathematical formulas
like
       \( x-3y + z = 7 \)
or
       \( a_{1} > x^{2n} + y^{2n} > x' \)
or
       \( \ip{A}{B} = \sum_{i} a_{i} b_{i} \).
The spaces you type in a formula are
ignored.  Remember that a letter like
       $x$                   % $ ... $  and  \( ... \)  are equivalent
is a formula when it denotes a mathematical
symbol, and it should be typed as one.

\section{Displayed Text}

Text is displayed by indenting it from the left
margin.  Quotations are commonly displayed.  There
are short quotations
\begin{quote}
   This is a short quotation.  It consists of a
   single paragraph of text.  See how it is formatted.
\end{quote}
and longer ones.
\begin{quotation}
   This is a longer quotation.  It consists of two
   paragraphs of text, neither of which are
   particularly interesting.

   This is the second paragraph of the quotation.  It
   is just as dull as the first paragraph.
\end{quotation}
Another frequently-displayed structure is a list.
The following is an example of an \emph{itemized}
list.
\begin{itemize}
   \item This is the first item of an itemized list.
         Each item in the list is marked with a ``tick''.
         You don't have to worry about what kind of tick
         mark is used.

   \item This is the second item of the list.  It
         contains another list nested inside it.  The inner
         list is an \emph{enumerated} list.
         \begin{enumerate}
            \item This is the first item of an enumerated
                  list that is nested within the itemized list.

            \item This is the second item of the inner list.
                  \LaTeX\ allows you to nest lists deeper than
                  you really should.
         \end{enumerate}
         This is the rest of the second item of the outer
         list.  It is no more interesting than any other
         part of the item.
   \item This is the third item of the list.
\end{itemize}
You can even display poetry.
\begin{verse}
   There is an environment
    for verse \\             % The \\ command separates lines
   Whose features some poets % within a stanza.
   will curse.

                             % One or more blank lines separate stanzas.

   For instead of making\\
   Them do \emph{all} line breaking, \\
   It allows them to put too many words on a line when they'd rather be
   forced to be terse.
\end{verse}

Mathematical formulas may also be displayed.  A
displayed formula
is
one-line long; multiline
formulas require special formatting instructions.
   \[  \ip{\Gamma}{\psi'} = x'' + y^{2} + z_{i}^{n}\]
Don't start a paragraph with a displayed equation,
nor make one a paragraph by itself.

\end{document}               % End of document.

\section{Bench Huge Section 1146}
% synthetic bench section 1146
% This is a sample LaTeX input file.  (Version of 12 August 2004.)
%
% A '%' character causes TeX to ignore all remaining text on the line,
% and is used for comments like this one.

\documentclass{article}      % Specifies the document class

                             % The preamble begins here.
\title{An Example Document}  % Declares the document's title.
\author{Leslie Lamport}      % Declares the author's name.
\date{January 21, 1994}      % Deleting this command produces today's date.

\newcommand{\ip}[2]{(#1, #2)}
                             % Defines \ip{arg1}{arg2} to mean
                             % (arg1, arg2).

%\newcommand{\ip}[2]{\langle #1 | #2\rangle}
                             % This is an alternative definition of
                             % \ip that is commented out.

\begin{document}             % End of preamble and beginning of text.

\maketitle                   % Produces the title.

This is an example input file.  Comparing it with
the output it generates can show you how to
produce a simple document of your own.

\section{Ordinary Text}      % Produces section heading.  Lower-level
                             % sections are begun with similar
                             % \subsection and \subsubsection commands.

The ends  of words and sentences are marked
  by   spaces. It  doesn't matter how many
spaces    you type; one is as good as 100.  The
end of   a line counts as a space.

One   or more   blank lines denote the  end
of  a paragraph.

Since any number of consecutive spaces are treated
like a single one, the formatting of the input
file makes no difference to
      \LaTeX,                % The \LaTeX command generates the LaTeX logo.
but it makes a difference to you.  When you use
\LaTeX, making your input file as easy to read
as possible will be a great help as you write
your document and when you change it.  This sample
file shows how you can add comments to your own input
file.

Because printing is different from typewriting,
there are a number of things that you have to do
differently when preparing an input file than if
you were just typing the document directly.
Quotation marks like
       ``this''
have to be handled specially, as do quotes within
quotes:
       ``\,`this'            % \, separates the double and single quote.
        is what I just
        wrote, not  `that'\,''.

Dashes come in three sizes: an
       intra-word
dash, a medium dash for number ranges like
       1--2,
and a punctuation
       dash---like
this.

A sentence-ending space should be larger than the
space between words within a sentence.  You
sometimes have to type special commands in
conjunction with punctuation characters to get
this right, as in the following sentence.
       Gnats, gnus, etc.\ all  % `\ ' makes an inter-word space.
       begin with G\@.         % \@ marks end-of-sentence punctuation.
You should check the spaces after periods when
reading your output to make sure you haven't
forgotten any special cases.  Generating an
ellipsis
       \ldots\               % `\ ' is needed after `\ldots' because TeX
                             % ignores spaces after command names like \ldots
                             % made from \ + letters.
                             %
                             % Note how a `%' character causes TeX to ignore
                             % the end of the input line, so these blank lines
                             % do not start a new paragraph.
                             %
with the right spacing around the periods requires
a special command.

\LaTeX\ interprets some common characters as
commands, so you must type special commands to
generate them.  These characters include the
following:
       \$ \& \% \# \{ and \}.

In printing, text is usually emphasized with an
       \emph{italic}
type style.

\begin{em}
   A long segment of text can also be emphasized
   in this way.  Text within such a segment can be
   given \emph{additional} emphasis.
\end{em}

It is sometimes necessary to prevent \LaTeX\ from
breaking a line where it might otherwise do so.
This may be at a space, as between the ``Mr.''\ and
``Jones'' in
       ``Mr.~Jones'',        % ~ produces an unbreakable interword space.
or within a word---especially when the word is a
symbol like
       \mbox{\emph{itemnum}}
that makes little sense when hyphenated across
lines.

Footnotes\footnote{This is an example of a footnote.}
pose no problem.

\LaTeX\ is good at typesetting mathematical formulas
like
       \( x-3y + z = 7 \)
or
       \( a_{1} > x^{2n} + y^{2n} > x' \)
or
       \( \ip{A}{B} = \sum_{i} a_{i} b_{i} \).
The spaces you type in a formula are
ignored.  Remember that a letter like
       $x$                   % $ ... $  and  \( ... \)  are equivalent
is a formula when it denotes a mathematical
symbol, and it should be typed as one.

\section{Displayed Text}

Text is displayed by indenting it from the left
margin.  Quotations are commonly displayed.  There
are short quotations
\begin{quote}
   This is a short quotation.  It consists of a
   single paragraph of text.  See how it is formatted.
\end{quote}
and longer ones.
\begin{quotation}
   This is a longer quotation.  It consists of two
   paragraphs of text, neither of which are
   particularly interesting.

   This is the second paragraph of the quotation.  It
   is just as dull as the first paragraph.
\end{quotation}
Another frequently-displayed structure is a list.
The following is an example of an \emph{itemized}
list.
\begin{itemize}
   \item This is the first item of an itemized list.
         Each item in the list is marked with a ``tick''.
         You don't have to worry about what kind of tick
         mark is used.

   \item This is the second item of the list.  It
         contains another list nested inside it.  The inner
         list is an \emph{enumerated} list.
         \begin{enumerate}
            \item This is the first item of an enumerated
                  list that is nested within the itemized list.

            \item This is the second item of the inner list.
                  \LaTeX\ allows you to nest lists deeper than
                  you really should.
         \end{enumerate}
         This is the rest of the second item of the outer
         list.  It is no more interesting than any other
         part of the item.
   \item This is the third item of the list.
\end{itemize}
You can even display poetry.
\begin{verse}
   There is an environment
    for verse \\             % The \\ command separates lines
   Whose features some poets % within a stanza.
   will curse.

                             % One or more blank lines separate stanzas.

   For instead of making\\
   Them do \emph{all} line breaking, \\
   It allows them to put too many words on a line when they'd rather be
   forced to be terse.
\end{verse}

Mathematical formulas may also be displayed.  A
displayed formula
is
one-line long; multiline
formulas require special formatting instructions.
   \[  \ip{\Gamma}{\psi'} = x'' + y^{2} + z_{i}^{n}\]
Don't start a paragraph with a displayed equation,
nor make one a paragraph by itself.

\end{document}               % End of document.

\section{Bench Huge Section 1147}
% synthetic bench section 1147
% This is a sample LaTeX input file.  (Version of 12 August 2004.)
%
% A '%' character causes TeX to ignore all remaining text on the line,
% and is used for comments like this one.

\documentclass{article}      % Specifies the document class

                             % The preamble begins here.
\title{An Example Document}  % Declares the document's title.
\author{Leslie Lamport}      % Declares the author's name.
\date{January 21, 1994}      % Deleting this command produces today's date.

\newcommand{\ip}[2]{(#1, #2)}
                             % Defines \ip{arg1}{arg2} to mean
                             % (arg1, arg2).

%\newcommand{\ip}[2]{\langle #1 | #2\rangle}
                             % This is an alternative definition of
                             % \ip that is commented out.

\begin{document}             % End of preamble and beginning of text.

\maketitle                   % Produces the title.

This is an example input file.  Comparing it with
the output it generates can show you how to
produce a simple document of your own.

\section{Ordinary Text}      % Produces section heading.  Lower-level
                             % sections are begun with similar
                             % \subsection and \subsubsection commands.

The ends  of words and sentences are marked
  by   spaces. It  doesn't matter how many
spaces    you type; one is as good as 100.  The
end of   a line counts as a space.

One   or more   blank lines denote the  end
of  a paragraph.

Since any number of consecutive spaces are treated
like a single one, the formatting of the input
file makes no difference to
      \LaTeX,                % The \LaTeX command generates the LaTeX logo.
but it makes a difference to you.  When you use
\LaTeX, making your input file as easy to read
as possible will be a great help as you write
your document and when you change it.  This sample
file shows how you can add comments to your own input
file.

Because printing is different from typewriting,
there are a number of things that you have to do
differently when preparing an input file than if
you were just typing the document directly.
Quotation marks like
       ``this''
have to be handled specially, as do quotes within
quotes:
       ``\,`this'            % \, separates the double and single quote.
        is what I just
        wrote, not  `that'\,''.

Dashes come in three sizes: an
       intra-word
dash, a medium dash for number ranges like
       1--2,
and a punctuation
       dash---like
this.

A sentence-ending space should be larger than the
space between words within a sentence.  You
sometimes have to type special commands in
conjunction with punctuation characters to get
this right, as in the following sentence.
       Gnats, gnus, etc.\ all  % `\ ' makes an inter-word space.
       begin with G\@.         % \@ marks end-of-sentence punctuation.
You should check the spaces after periods when
reading your output to make sure you haven't
forgotten any special cases.  Generating an
ellipsis
       \ldots\               % `\ ' is needed after `\ldots' because TeX
                             % ignores spaces after command names like \ldots
                             % made from \ + letters.
                             %
                             % Note how a `%' character causes TeX to ignore
                             % the end of the input line, so these blank lines
                             % do not start a new paragraph.
                             %
with the right spacing around the periods requires
a special command.

\LaTeX\ interprets some common characters as
commands, so you must type special commands to
generate them.  These characters include the
following:
       \$ \& \% \# \{ and \}.

In printing, text is usually emphasized with an
       \emph{italic}
type style.

\begin{em}
   A long segment of text can also be emphasized
   in this way.  Text within such a segment can be
   given \emph{additional} emphasis.
\end{em}

It is sometimes necessary to prevent \LaTeX\ from
breaking a line where it might otherwise do so.
This may be at a space, as between the ``Mr.''\ and
``Jones'' in
       ``Mr.~Jones'',        % ~ produces an unbreakable interword space.
or within a word---especially when the word is a
symbol like
       \mbox{\emph{itemnum}}
that makes little sense when hyphenated across
lines.

Footnotes\footnote{This is an example of a footnote.}
pose no problem.

\LaTeX\ is good at typesetting mathematical formulas
like
       \( x-3y + z = 7 \)
or
       \( a_{1} > x^{2n} + y^{2n} > x' \)
or
       \( \ip{A}{B} = \sum_{i} a_{i} b_{i} \).
The spaces you type in a formula are
ignored.  Remember that a letter like
       $x$                   % $ ... $  and  \( ... \)  are equivalent
is a formula when it denotes a mathematical
symbol, and it should be typed as one.

\section{Displayed Text}

Text is displayed by indenting it from the left
margin.  Quotations are commonly displayed.  There
are short quotations
\begin{quote}
   This is a short quotation.  It consists of a
   single paragraph of text.  See how it is formatted.
\end{quote}
and longer ones.
\begin{quotation}
   This is a longer quotation.  It consists of two
   paragraphs of text, neither of which are
   particularly interesting.

   This is the second paragraph of the quotation.  It
   is just as dull as the first paragraph.
\end{quotation}
Another frequently-displayed structure is a list.
The following is an example of an \emph{itemized}
list.
\begin{itemize}
   \item This is the first item of an itemized list.
         Each item in the list is marked with a ``tick''.
         You don't have to worry about what kind of tick
         mark is used.

   \item This is the second item of the list.  It
         contains another list nested inside it.  The inner
         list is an \emph{enumerated} list.
         \begin{enumerate}
            \item This is the first item of an enumerated
                  list that is nested within the itemized list.

            \item This is the second item of the inner list.
                  \LaTeX\ allows you to nest lists deeper than
                  you really should.
         \end{enumerate}
         This is the rest of the second item of the outer
         list.  It is no more interesting than any other
         part of the item.
   \item This is the third item of the list.
\end{itemize}
You can even display poetry.
\begin{verse}
   There is an environment
    for verse \\             % The \\ command separates lines
   Whose features some poets % within a stanza.
   will curse.

                             % One or more blank lines separate stanzas.

   For instead of making\\
   Them do \emph{all} line breaking, \\
   It allows them to put too many words on a line when they'd rather be
   forced to be terse.
\end{verse}

Mathematical formulas may also be displayed.  A
displayed formula
is
one-line long; multiline
formulas require special formatting instructions.
   \[  \ip{\Gamma}{\psi'} = x'' + y^{2} + z_{i}^{n}\]
Don't start a paragraph with a displayed equation,
nor make one a paragraph by itself.

\end{document}               % End of document.

\section{Bench Huge Section 1148}
% synthetic bench section 1148
% This is a sample LaTeX input file.  (Version of 12 August 2004.)
%
% A '%' character causes TeX to ignore all remaining text on the line,
% and is used for comments like this one.

\documentclass{article}      % Specifies the document class

                             % The preamble begins here.
\title{An Example Document}  % Declares the document's title.
\author{Leslie Lamport}      % Declares the author's name.
\date{January 21, 1994}      % Deleting this command produces today's date.

\newcommand{\ip}[2]{(#1, #2)}
                             % Defines \ip{arg1}{arg2} to mean
                             % (arg1, arg2).

%\newcommand{\ip}[2]{\langle #1 | #2\rangle}
                             % This is an alternative definition of
                             % \ip that is commented out.

\begin{document}             % End of preamble and beginning of text.

\maketitle                   % Produces the title.

This is an example input file.  Comparing it with
the output it generates can show you how to
produce a simple document of your own.

\section{Ordinary Text}      % Produces section heading.  Lower-level
                             % sections are begun with similar
                             % \subsection and \subsubsection commands.

The ends  of words and sentences are marked
  by   spaces. It  doesn't matter how many
spaces    you type; one is as good as 100.  The
end of   a line counts as a space.

One   or more   blank lines denote the  end
of  a paragraph.

Since any number of consecutive spaces are treated
like a single one, the formatting of the input
file makes no difference to
      \LaTeX,                % The \LaTeX command generates the LaTeX logo.
but it makes a difference to you.  When you use
\LaTeX, making your input file as easy to read
as possible will be a great help as you write
your document and when you change it.  This sample
file shows how you can add comments to your own input
file.

Because printing is different from typewriting,
there are a number of things that you have to do
differently when preparing an input file than if
you were just typing the document directly.
Quotation marks like
       ``this''
have to be handled specially, as do quotes within
quotes:
       ``\,`this'            % \, separates the double and single quote.
        is what I just
        wrote, not  `that'\,''.

Dashes come in three sizes: an
       intra-word
dash, a medium dash for number ranges like
       1--2,
and a punctuation
       dash---like
this.

A sentence-ending space should be larger than the
space between words within a sentence.  You
sometimes have to type special commands in
conjunction with punctuation characters to get
this right, as in the following sentence.
       Gnats, gnus, etc.\ all  % `\ ' makes an inter-word space.
       begin with G\@.         % \@ marks end-of-sentence punctuation.
You should check the spaces after periods when
reading your output to make sure you haven't
forgotten any special cases.  Generating an
ellipsis
       \ldots\               % `\ ' is needed after `\ldots' because TeX
                             % ignores spaces after command names like \ldots
                             % made from \ + letters.
                             %
                             % Note how a `%' character causes TeX to ignore
                             % the end of the input line, so these blank lines
                             % do not start a new paragraph.
                             %
with the right spacing around the periods requires
a special command.

\LaTeX\ interprets some common characters as
commands, so you must type special commands to
generate them.  These characters include the
following:
       \$ \& \% \# \{ and \}.

In printing, text is usually emphasized with an
       \emph{italic}
type style.

\begin{em}
   A long segment of text can also be emphasized
   in this way.  Text within such a segment can be
   given \emph{additional} emphasis.
\end{em}

It is sometimes necessary to prevent \LaTeX\ from
breaking a line where it might otherwise do so.
This may be at a space, as between the ``Mr.''\ and
``Jones'' in
       ``Mr.~Jones'',        % ~ produces an unbreakable interword space.
or within a word---especially when the word is a
symbol like
       \mbox{\emph{itemnum}}
that makes little sense when hyphenated across
lines.

Footnotes\footnote{This is an example of a footnote.}
pose no problem.

\LaTeX\ is good at typesetting mathematical formulas
like
       \( x-3y + z = 7 \)
or
       \( a_{1} > x^{2n} + y^{2n} > x' \)
or
       \( \ip{A}{B} = \sum_{i} a_{i} b_{i} \).
The spaces you type in a formula are
ignored.  Remember that a letter like
       $x$                   % $ ... $  and  \( ... \)  are equivalent
is a formula when it denotes a mathematical
symbol, and it should be typed as one.

\section{Displayed Text}

Text is displayed by indenting it from the left
margin.  Quotations are commonly displayed.  There
are short quotations
\begin{quote}
   This is a short quotation.  It consists of a
   single paragraph of text.  See how it is formatted.
\end{quote}
and longer ones.
\begin{quotation}
   This is a longer quotation.  It consists of two
   paragraphs of text, neither of which are
   particularly interesting.

   This is the second paragraph of the quotation.  It
   is just as dull as the first paragraph.
\end{quotation}
Another frequently-displayed structure is a list.
The following is an example of an \emph{itemized}
list.
\begin{itemize}
   \item This is the first item of an itemized list.
         Each item in the list is marked with a ``tick''.
         You don't have to worry about what kind of tick
         mark is used.

   \item This is the second item of the list.  It
         contains another list nested inside it.  The inner
         list is an \emph{enumerated} list.
         \begin{enumerate}
            \item This is the first item of an enumerated
                  list that is nested within the itemized list.

            \item This is the second item of the inner list.
                  \LaTeX\ allows you to nest lists deeper than
                  you really should.
         \end{enumerate}
         This is the rest of the second item of the outer
         list.  It is no more interesting than any other
         part of the item.
   \item This is the third item of the list.
\end{itemize}
You can even display poetry.
\begin{verse}
   There is an environment
    for verse \\             % The \\ command separates lines
   Whose features some poets % within a stanza.
   will curse.

                             % One or more blank lines separate stanzas.

   For instead of making\\
   Them do \emph{all} line breaking, \\
   It allows them to put too many words on a line when they'd rather be
   forced to be terse.
\end{verse}

Mathematical formulas may also be displayed.  A
displayed formula
is
one-line long; multiline
formulas require special formatting instructions.
   \[  \ip{\Gamma}{\psi'} = x'' + y^{2} + z_{i}^{n}\]
Don't start a paragraph with a displayed equation,
nor make one a paragraph by itself.

\end{document}               % End of document.

\section{Bench Huge Section 1149}
% synthetic bench section 1149
% This is a sample LaTeX input file.  (Version of 12 August 2004.)
%
% A '%' character causes TeX to ignore all remaining text on the line,
% and is used for comments like this one.

\documentclass{article}      % Specifies the document class

                             % The preamble begins here.
\title{An Example Document}  % Declares the document's title.
\author{Leslie Lamport}      % Declares the author's name.
\date{January 21, 1994}      % Deleting this command produces today's date.

\newcommand{\ip}[2]{(#1, #2)}
                             % Defines \ip{arg1}{arg2} to mean
                             % (arg1, arg2).

%\newcommand{\ip}[2]{\langle #1 | #2\rangle}
                             % This is an alternative definition of
                             % \ip that is commented out.

\begin{document}             % End of preamble and beginning of text.

\maketitle                   % Produces the title.

This is an example input file.  Comparing it with
the output it generates can show you how to
produce a simple document of your own.

\section{Ordinary Text}      % Produces section heading.  Lower-level
                             % sections are begun with similar
                             % \subsection and \subsubsection commands.

The ends  of words and sentences are marked
  by   spaces. It  doesn't matter how many
spaces    you type; one is as good as 100.  The
end of   a line counts as a space.

One   or more   blank lines denote the  end
of  a paragraph.

Since any number of consecutive spaces are treated
like a single one, the formatting of the input
file makes no difference to
      \LaTeX,                % The \LaTeX command generates the LaTeX logo.
but it makes a difference to you.  When you use
\LaTeX, making your input file as easy to read
as possible will be a great help as you write
your document and when you change it.  This sample
file shows how you can add comments to your own input
file.

Because printing is different from typewriting,
there are a number of things that you have to do
differently when preparing an input file than if
you were just typing the document directly.
Quotation marks like
       ``this''
have to be handled specially, as do quotes within
quotes:
       ``\,`this'            % \, separates the double and single quote.
        is what I just
        wrote, not  `that'\,''.

Dashes come in three sizes: an
       intra-word
dash, a medium dash for number ranges like
       1--2,
and a punctuation
       dash---like
this.

A sentence-ending space should be larger than the
space between words within a sentence.  You
sometimes have to type special commands in
conjunction with punctuation characters to get
this right, as in the following sentence.
       Gnats, gnus, etc.\ all  % `\ ' makes an inter-word space.
       begin with G\@.         % \@ marks end-of-sentence punctuation.
You should check the spaces after periods when
reading your output to make sure you haven't
forgotten any special cases.  Generating an
ellipsis
       \ldots\               % `\ ' is needed after `\ldots' because TeX
                             % ignores spaces after command names like \ldots
                             % made from \ + letters.
                             %
                             % Note how a `%' character causes TeX to ignore
                             % the end of the input line, so these blank lines
                             % do not start a new paragraph.
                             %
with the right spacing around the periods requires
a special command.

\LaTeX\ interprets some common characters as
commands, so you must type special commands to
generate them.  These characters include the
following:
       \$ \& \% \# \{ and \}.

In printing, text is usually emphasized with an
       \emph{italic}
type style.

\begin{em}
   A long segment of text can also be emphasized
   in this way.  Text within such a segment can be
   given \emph{additional} emphasis.
\end{em}

It is sometimes necessary to prevent \LaTeX\ from
breaking a line where it might otherwise do so.
This may be at a space, as between the ``Mr.''\ and
``Jones'' in
       ``Mr.~Jones'',        % ~ produces an unbreakable interword space.
or within a word---especially when the word is a
symbol like
       \mbox{\emph{itemnum}}
that makes little sense when hyphenated across
lines.

Footnotes\footnote{This is an example of a footnote.}
pose no problem.

\LaTeX\ is good at typesetting mathematical formulas
like
       \( x-3y + z = 7 \)
or
       \( a_{1} > x^{2n} + y^{2n} > x' \)
or
       \( \ip{A}{B} = \sum_{i} a_{i} b_{i} \).
The spaces you type in a formula are
ignored.  Remember that a letter like
       $x$                   % $ ... $  and  \( ... \)  are equivalent
is a formula when it denotes a mathematical
symbol, and it should be typed as one.

\section{Displayed Text}

Text is displayed by indenting it from the left
margin.  Quotations are commonly displayed.  There
are short quotations
\begin{quote}
   This is a short quotation.  It consists of a
   single paragraph of text.  See how it is formatted.
\end{quote}
and longer ones.
\begin{quotation}
   This is a longer quotation.  It consists of two
   paragraphs of text, neither of which are
   particularly interesting.

   This is the second paragraph of the quotation.  It
   is just as dull as the first paragraph.
\end{quotation}
Another frequently-displayed structure is a list.
The following is an example of an \emph{itemized}
list.
\begin{itemize}
   \item This is the first item of an itemized list.
         Each item in the list is marked with a ``tick''.
         You don't have to worry about what kind of tick
         mark is used.

   \item This is the second item of the list.  It
         contains another list nested inside it.  The inner
         list is an \emph{enumerated} list.
         \begin{enumerate}
            \item This is the first item of an enumerated
                  list that is nested within the itemized list.

            \item This is the second item of the inner list.
                  \LaTeX\ allows you to nest lists deeper than
                  you really should.
         \end{enumerate}
         This is the rest of the second item of the outer
         list.  It is no more interesting than any other
         part of the item.
   \item This is the third item of the list.
\end{itemize}
You can even display poetry.
\begin{verse}
   There is an environment
    for verse \\             % The \\ command separates lines
   Whose features some poets % within a stanza.
   will curse.

                             % One or more blank lines separate stanzas.

   For instead of making\\
   Them do \emph{all} line breaking, \\
   It allows them to put too many words on a line when they'd rather be
   forced to be terse.
\end{verse}

Mathematical formulas may also be displayed.  A
displayed formula
is
one-line long; multiline
formulas require special formatting instructions.
   \[  \ip{\Gamma}{\psi'} = x'' + y^{2} + z_{i}^{n}\]
Don't start a paragraph with a displayed equation,
nor make one a paragraph by itself.

\end{document}               % End of document.

\section{Bench Huge Section 1150}
% synthetic bench section 1150
% This is a sample LaTeX input file.  (Version of 12 August 2004.)
%
% A '%' character causes TeX to ignore all remaining text on the line,
% and is used for comments like this one.

\documentclass{article}      % Specifies the document class

                             % The preamble begins here.
\title{An Example Document}  % Declares the document's title.
\author{Leslie Lamport}      % Declares the author's name.
\date{January 21, 1994}      % Deleting this command produces today's date.

\newcommand{\ip}[2]{(#1, #2)}
                             % Defines \ip{arg1}{arg2} to mean
                             % (arg1, arg2).

%\newcommand{\ip}[2]{\langle #1 | #2\rangle}
                             % This is an alternative definition of
                             % \ip that is commented out.

\begin{document}             % End of preamble and beginning of text.

\maketitle                   % Produces the title.

This is an example input file.  Comparing it with
the output it generates can show you how to
produce a simple document of your own.

\section{Ordinary Text}      % Produces section heading.  Lower-level
                             % sections are begun with similar
                             % \subsection and \subsubsection commands.

The ends  of words and sentences are marked
  by   spaces. It  doesn't matter how many
spaces    you type; one is as good as 100.  The
end of   a line counts as a space.

One   or more   blank lines denote the  end
of  a paragraph.

Since any number of consecutive spaces are treated
like a single one, the formatting of the input
file makes no difference to
      \LaTeX,                % The \LaTeX command generates the LaTeX logo.
but it makes a difference to you.  When you use
\LaTeX, making your input file as easy to read
as possible will be a great help as you write
your document and when you change it.  This sample
file shows how you can add comments to your own input
file.

Because printing is different from typewriting,
there are a number of things that you have to do
differently when preparing an input file than if
you were just typing the document directly.
Quotation marks like
       ``this''
have to be handled specially, as do quotes within
quotes:
       ``\,`this'            % \, separates the double and single quote.
        is what I just
        wrote, not  `that'\,''.

Dashes come in three sizes: an
       intra-word
dash, a medium dash for number ranges like
       1--2,
and a punctuation
       dash---like
this.

A sentence-ending space should be larger than the
space between words within a sentence.  You
sometimes have to type special commands in
conjunction with punctuation characters to get
this right, as in the following sentence.
       Gnats, gnus, etc.\ all  % `\ ' makes an inter-word space.
       begin with G\@.         % \@ marks end-of-sentence punctuation.
You should check the spaces after periods when
reading your output to make sure you haven't
forgotten any special cases.  Generating an
ellipsis
       \ldots\               % `\ ' is needed after `\ldots' because TeX
                             % ignores spaces after command names like \ldots
                             % made from \ + letters.
                             %
                             % Note how a `%' character causes TeX to ignore
                             % the end of the input line, so these blank lines
                             % do not start a new paragraph.
                             %
with the right spacing around the periods requires
a special command.

\LaTeX\ interprets some common characters as
commands, so you must type special commands to
generate them.  These characters include the
following:
       \$ \& \% \# \{ and \}.

In printing, text is usually emphasized with an
       \emph{italic}
type style.

\begin{em}
   A long segment of text can also be emphasized
   in this way.  Text within such a segment can be
   given \emph{additional} emphasis.
\end{em}

It is sometimes necessary to prevent \LaTeX\ from
breaking a line where it might otherwise do so.
This may be at a space, as between the ``Mr.''\ and
``Jones'' in
       ``Mr.~Jones'',        % ~ produces an unbreakable interword space.
or within a word---especially when the word is a
symbol like
       \mbox{\emph{itemnum}}
that makes little sense when hyphenated across
lines.

Footnotes\footnote{This is an example of a footnote.}
pose no problem.

\LaTeX\ is good at typesetting mathematical formulas
like
       \( x-3y + z = 7 \)
or
       \( a_{1} > x^{2n} + y^{2n} > x' \)
or
       \( \ip{A}{B} = \sum_{i} a_{i} b_{i} \).
The spaces you type in a formula are
ignored.  Remember that a letter like
       $x$                   % $ ... $  and  \( ... \)  are equivalent
is a formula when it denotes a mathematical
symbol, and it should be typed as one.

\section{Displayed Text}

Text is displayed by indenting it from the left
margin.  Quotations are commonly displayed.  There
are short quotations
\begin{quote}
   This is a short quotation.  It consists of a
   single paragraph of text.  See how it is formatted.
\end{quote}
and longer ones.
\begin{quotation}
   This is a longer quotation.  It consists of two
   paragraphs of text, neither of which are
   particularly interesting.

   This is the second paragraph of the quotation.  It
   is just as dull as the first paragraph.
\end{quotation}
Another frequently-displayed structure is a list.
The following is an example of an \emph{itemized}
list.
\begin{itemize}
   \item This is the first item of an itemized list.
         Each item in the list is marked with a ``tick''.
         You don't have to worry about what kind of tick
         mark is used.

   \item This is the second item of the list.  It
         contains another list nested inside it.  The inner
         list is an \emph{enumerated} list.
         \begin{enumerate}
            \item This is the first item of an enumerated
                  list that is nested within the itemized list.

            \item This is the second item of the inner list.
                  \LaTeX\ allows you to nest lists deeper than
                  you really should.
         \end{enumerate}
         This is the rest of the second item of the outer
         list.  It is no more interesting than any other
         part of the item.
   \item This is the third item of the list.
\end{itemize}
You can even display poetry.
\begin{verse}
   There is an environment
    for verse \\             % The \\ command separates lines
   Whose features some poets % within a stanza.
   will curse.

                             % One or more blank lines separate stanzas.

   For instead of making\\
   Them do \emph{all} line breaking, \\
   It allows them to put too many words on a line when they'd rather be
   forced to be terse.
\end{verse}

Mathematical formulas may also be displayed.  A
displayed formula
is
one-line long; multiline
formulas require special formatting instructions.
   \[  \ip{\Gamma}{\psi'} = x'' + y^{2} + z_{i}^{n}\]
Don't start a paragraph with a displayed equation,
nor make one a paragraph by itself.

\end{document}               % End of document.

\section{Bench Huge Section 1151}
% synthetic bench section 1151
% This is a sample LaTeX input file.  (Version of 12 August 2004.)
%
% A '%' character causes TeX to ignore all remaining text on the line,
% and is used for comments like this one.

\documentclass{article}      % Specifies the document class

                             % The preamble begins here.
\title{An Example Document}  % Declares the document's title.
\author{Leslie Lamport}      % Declares the author's name.
\date{January 21, 1994}      % Deleting this command produces today's date.

\newcommand{\ip}[2]{(#1, #2)}
                             % Defines \ip{arg1}{arg2} to mean
                             % (arg1, arg2).

%\newcommand{\ip}[2]{\langle #1 | #2\rangle}
                             % This is an alternative definition of
                             % \ip that is commented out.

\begin{document}             % End of preamble and beginning of text.

\maketitle                   % Produces the title.

This is an example input file.  Comparing it with
the output it generates can show you how to
produce a simple document of your own.

\section{Ordinary Text}      % Produces section heading.  Lower-level
                             % sections are begun with similar
                             % \subsection and \subsubsection commands.

The ends  of words and sentences are marked
  by   spaces. It  doesn't matter how many
spaces    you type; one is as good as 100.  The
end of   a line counts as a space.

One   or more   blank lines denote the  end
of  a paragraph.

Since any number of consecutive spaces are treated
like a single one, the formatting of the input
file makes no difference to
      \LaTeX,                % The \LaTeX command generates the LaTeX logo.
but it makes a difference to you.  When you use
\LaTeX, making your input file as easy to read
as possible will be a great help as you write
your document and when you change it.  This sample
file shows how you can add comments to your own input
file.

Because printing is different from typewriting,
there are a number of things that you have to do
differently when preparing an input file than if
you were just typing the document directly.
Quotation marks like
       ``this''
have to be handled specially, as do quotes within
quotes:
       ``\,`this'            % \, separates the double and single quote.
        is what I just
        wrote, not  `that'\,''.

Dashes come in three sizes: an
       intra-word
dash, a medium dash for number ranges like
       1--2,
and a punctuation
       dash---like
this.

A sentence-ending space should be larger than the
space between words within a sentence.  You
sometimes have to type special commands in
conjunction with punctuation characters to get
this right, as in the following sentence.
       Gnats, gnus, etc.\ all  % `\ ' makes an inter-word space.
       begin with G\@.         % \@ marks end-of-sentence punctuation.
You should check the spaces after periods when
reading your output to make sure you haven't
forgotten any special cases.  Generating an
ellipsis
       \ldots\               % `\ ' is needed after `\ldots' because TeX
                             % ignores spaces after command names like \ldots
                             % made from \ + letters.
                             %
                             % Note how a `%' character causes TeX to ignore
                             % the end of the input line, so these blank lines
                             % do not start a new paragraph.
                             %
with the right spacing around the periods requires
a special command.

\LaTeX\ interprets some common characters as
commands, so you must type special commands to
generate them.  These characters include the
following:
       \$ \& \% \# \{ and \}.

In printing, text is usually emphasized with an
       \emph{italic}
type style.

\begin{em}
   A long segment of text can also be emphasized
   in this way.  Text within such a segment can be
   given \emph{additional} emphasis.
\end{em}

It is sometimes necessary to prevent \LaTeX\ from
breaking a line where it might otherwise do so.
This may be at a space, as between the ``Mr.''\ and
``Jones'' in
       ``Mr.~Jones'',        % ~ produces an unbreakable interword space.
or within a word---especially when the word is a
symbol like
       \mbox{\emph{itemnum}}
that makes little sense when hyphenated across
lines.

Footnotes\footnote{This is an example of a footnote.}
pose no problem.

\LaTeX\ is good at typesetting mathematical formulas
like
       \( x-3y + z = 7 \)
or
       \( a_{1} > x^{2n} + y^{2n} > x' \)
or
       \( \ip{A}{B} = \sum_{i} a_{i} b_{i} \).
The spaces you type in a formula are
ignored.  Remember that a letter like
       $x$                   % $ ... $  and  \( ... \)  are equivalent
is a formula when it denotes a mathematical
symbol, and it should be typed as one.

\section{Displayed Text}

Text is displayed by indenting it from the left
margin.  Quotations are commonly displayed.  There
are short quotations
\begin{quote}
   This is a short quotation.  It consists of a
   single paragraph of text.  See how it is formatted.
\end{quote}
and longer ones.
\begin{quotation}
   This is a longer quotation.  It consists of two
   paragraphs of text, neither of which are
   particularly interesting.

   This is the second paragraph of the quotation.  It
   is just as dull as the first paragraph.
\end{quotation}
Another frequently-displayed structure is a list.
The following is an example of an \emph{itemized}
list.
\begin{itemize}
   \item This is the first item of an itemized list.
         Each item in the list is marked with a ``tick''.
         You don't have to worry about what kind of tick
         mark is used.

   \item This is the second item of the list.  It
         contains another list nested inside it.  The inner
         list is an \emph{enumerated} list.
         \begin{enumerate}
            \item This is the first item of an enumerated
                  list that is nested within the itemized list.

            \item This is the second item of the inner list.
                  \LaTeX\ allows you to nest lists deeper than
                  you really should.
         \end{enumerate}
         This is the rest of the second item of the outer
         list.  It is no more interesting than any other
         part of the item.
   \item This is the third item of the list.
\end{itemize}
You can even display poetry.
\begin{verse}
   There is an environment
    for verse \\             % The \\ command separates lines
   Whose features some poets % within a stanza.
   will curse.

                             % One or more blank lines separate stanzas.

   For instead of making\\
   Them do \emph{all} line breaking, \\
   It allows them to put too many words on a line when they'd rather be
   forced to be terse.
\end{verse}

Mathematical formulas may also be displayed.  A
displayed formula
is
one-line long; multiline
formulas require special formatting instructions.
   \[  \ip{\Gamma}{\psi'} = x'' + y^{2} + z_{i}^{n}\]
Don't start a paragraph with a displayed equation,
nor make one a paragraph by itself.

\end{document}               % End of document.

\section{Bench Huge Section 1152}
% synthetic bench section 1152
% This is a sample LaTeX input file.  (Version of 12 August 2004.)
%
% A '%' character causes TeX to ignore all remaining text on the line,
% and is used for comments like this one.

\documentclass{article}      % Specifies the document class

                             % The preamble begins here.
\title{An Example Document}  % Declares the document's title.
\author{Leslie Lamport}      % Declares the author's name.
\date{January 21, 1994}      % Deleting this command produces today's date.

\newcommand{\ip}[2]{(#1, #2)}
                             % Defines \ip{arg1}{arg2} to mean
                             % (arg1, arg2).

%\newcommand{\ip}[2]{\langle #1 | #2\rangle}
                             % This is an alternative definition of
                             % \ip that is commented out.

\begin{document}             % End of preamble and beginning of text.

\maketitle                   % Produces the title.

This is an example input file.  Comparing it with
the output it generates can show you how to
produce a simple document of your own.

\section{Ordinary Text}      % Produces section heading.  Lower-level
                             % sections are begun with similar
                             % \subsection and \subsubsection commands.

The ends  of words and sentences are marked
  by   spaces. It  doesn't matter how many
spaces    you type; one is as good as 100.  The
end of   a line counts as a space.

One   or more   blank lines denote the  end
of  a paragraph.

Since any number of consecutive spaces are treated
like a single one, the formatting of the input
file makes no difference to
      \LaTeX,                % The \LaTeX command generates the LaTeX logo.
but it makes a difference to you.  When you use
\LaTeX, making your input file as easy to read
as possible will be a great help as you write
your document and when you change it.  This sample
file shows how you can add comments to your own input
file.

Because printing is different from typewriting,
there are a number of things that you have to do
differently when preparing an input file than if
you were just typing the document directly.
Quotation marks like
       ``this''
have to be handled specially, as do quotes within
quotes:
       ``\,`this'            % \, separates the double and single quote.
        is what I just
        wrote, not  `that'\,''.

Dashes come in three sizes: an
       intra-word
dash, a medium dash for number ranges like
       1--2,
and a punctuation
       dash---like
this.

A sentence-ending space should be larger than the
space between words within a sentence.  You
sometimes have to type special commands in
conjunction with punctuation characters to get
this right, as in the following sentence.
       Gnats, gnus, etc.\ all  % `\ ' makes an inter-word space.
       begin with G\@.         % \@ marks end-of-sentence punctuation.
You should check the spaces after periods when
reading your output to make sure you haven't
forgotten any special cases.  Generating an
ellipsis
       \ldots\               % `\ ' is needed after `\ldots' because TeX
                             % ignores spaces after command names like \ldots
                             % made from \ + letters.
                             %
                             % Note how a `%' character causes TeX to ignore
                             % the end of the input line, so these blank lines
                             % do not start a new paragraph.
                             %
with the right spacing around the periods requires
a special command.

\LaTeX\ interprets some common characters as
commands, so you must type special commands to
generate them.  These characters include the
following:
       \$ \& \% \# \{ and \}.

In printing, text is usually emphasized with an
       \emph{italic}
type style.

\begin{em}
   A long segment of text can also be emphasized
   in this way.  Text within such a segment can be
   given \emph{additional} emphasis.
\end{em}

It is sometimes necessary to prevent \LaTeX\ from
breaking a line where it might otherwise do so.
This may be at a space, as between the ``Mr.''\ and
``Jones'' in
       ``Mr.~Jones'',        % ~ produces an unbreakable interword space.
or within a word---especially when the word is a
symbol like
       \mbox{\emph{itemnum}}
that makes little sense when hyphenated across
lines.

Footnotes\footnote{This is an example of a footnote.}
pose no problem.

\LaTeX\ is good at typesetting mathematical formulas
like
       \( x-3y + z = 7 \)
or
       \( a_{1} > x^{2n} + y^{2n} > x' \)
or
       \( \ip{A}{B} = \sum_{i} a_{i} b_{i} \).
The spaces you type in a formula are
ignored.  Remember that a letter like
       $x$                   % $ ... $  and  \( ... \)  are equivalent
is a formula when it denotes a mathematical
symbol, and it should be typed as one.

\section{Displayed Text}

Text is displayed by indenting it from the left
margin.  Quotations are commonly displayed.  There
are short quotations
\begin{quote}
   This is a short quotation.  It consists of a
   single paragraph of text.  See how it is formatted.
\end{quote}
and longer ones.
\begin{quotation}
   This is a longer quotation.  It consists of two
   paragraphs of text, neither of which are
   particularly interesting.

   This is the second paragraph of the quotation.  It
   is just as dull as the first paragraph.
\end{quotation}
Another frequently-displayed structure is a list.
The following is an example of an \emph{itemized}
list.
\begin{itemize}
   \item This is the first item of an itemized list.
         Each item in the list is marked with a ``tick''.
         You don't have to worry about what kind of tick
         mark is used.

   \item This is the second item of the list.  It
         contains another list nested inside it.  The inner
         list is an \emph{enumerated} list.
         \begin{enumerate}
            \item This is the first item of an enumerated
                  list that is nested within the itemized list.

            \item This is the second item of the inner list.
                  \LaTeX\ allows you to nest lists deeper than
                  you really should.
         \end{enumerate}
         This is the rest of the second item of the outer
         list.  It is no more interesting than any other
         part of the item.
   \item This is the third item of the list.
\end{itemize}
You can even display poetry.
\begin{verse}
   There is an environment
    for verse \\             % The \\ command separates lines
   Whose features some poets % within a stanza.
   will curse.

                             % One or more blank lines separate stanzas.

   For instead of making\\
   Them do \emph{all} line breaking, \\
   It allows them to put too many words on a line when they'd rather be
   forced to be terse.
\end{verse}

Mathematical formulas may also be displayed.  A
displayed formula
is
one-line long; multiline
formulas require special formatting instructions.
   \[  \ip{\Gamma}{\psi'} = x'' + y^{2} + z_{i}^{n}\]
Don't start a paragraph with a displayed equation,
nor make one a paragraph by itself.

\end{document}               % End of document.

\section{Bench Huge Section 1153}
% synthetic bench section 1153
% This is a sample LaTeX input file.  (Version of 12 August 2004.)
%
% A '%' character causes TeX to ignore all remaining text on the line,
% and is used for comments like this one.

\documentclass{article}      % Specifies the document class

                             % The preamble begins here.
\title{An Example Document}  % Declares the document's title.
\author{Leslie Lamport}      % Declares the author's name.
\date{January 21, 1994}      % Deleting this command produces today's date.

\newcommand{\ip}[2]{(#1, #2)}
                             % Defines \ip{arg1}{arg2} to mean
                             % (arg1, arg2).

%\newcommand{\ip}[2]{\langle #1 | #2\rangle}
                             % This is an alternative definition of
                             % \ip that is commented out.

\begin{document}             % End of preamble and beginning of text.

\maketitle                   % Produces the title.

This is an example input file.  Comparing it with
the output it generates can show you how to
produce a simple document of your own.

\section{Ordinary Text}      % Produces section heading.  Lower-level
                             % sections are begun with similar
                             % \subsection and \subsubsection commands.

The ends  of words and sentences are marked
  by   spaces. It  doesn't matter how many
spaces    you type; one is as good as 100.  The
end of   a line counts as a space.

One   or more   blank lines denote the  end
of  a paragraph.

Since any number of consecutive spaces are treated
like a single one, the formatting of the input
file makes no difference to
      \LaTeX,                % The \LaTeX command generates the LaTeX logo.
but it makes a difference to you.  When you use
\LaTeX, making your input file as easy to read
as possible will be a great help as you write
your document and when you change it.  This sample
file shows how you can add comments to your own input
file.

Because printing is different from typewriting,
there are a number of things that you have to do
differently when preparing an input file than if
you were just typing the document directly.
Quotation marks like
       ``this''
have to be handled specially, as do quotes within
quotes:
       ``\,`this'            % \, separates the double and single quote.
        is what I just
        wrote, not  `that'\,''.

Dashes come in three sizes: an
       intra-word
dash, a medium dash for number ranges like
       1--2,
and a punctuation
       dash---like
this.

A sentence-ending space should be larger than the
space between words within a sentence.  You
sometimes have to type special commands in
conjunction with punctuation characters to get
this right, as in the following sentence.
       Gnats, gnus, etc.\ all  % `\ ' makes an inter-word space.
       begin with G\@.         % \@ marks end-of-sentence punctuation.
You should check the spaces after periods when
reading your output to make sure you haven't
forgotten any special cases.  Generating an
ellipsis
       \ldots\               % `\ ' is needed after `\ldots' because TeX
                             % ignores spaces after command names like \ldots
                             % made from \ + letters.
                             %
                             % Note how a `%' character causes TeX to ignore
                             % the end of the input line, so these blank lines
                             % do not start a new paragraph.
                             %
with the right spacing around the periods requires
a special command.

\LaTeX\ interprets some common characters as
commands, so you must type special commands to
generate them.  These characters include the
following:
       \$ \& \% \# \{ and \}.

In printing, text is usually emphasized with an
       \emph{italic}
type style.

\begin{em}
   A long segment of text can also be emphasized
   in this way.  Text within such a segment can be
   given \emph{additional} emphasis.
\end{em}

It is sometimes necessary to prevent \LaTeX\ from
breaking a line where it might otherwise do so.
This may be at a space, as between the ``Mr.''\ and
``Jones'' in
       ``Mr.~Jones'',        % ~ produces an unbreakable interword space.
or within a word---especially when the word is a
symbol like
       \mbox{\emph{itemnum}}
that makes little sense when hyphenated across
lines.

Footnotes\footnote{This is an example of a footnote.}
pose no problem.

\LaTeX\ is good at typesetting mathematical formulas
like
       \( x-3y + z = 7 \)
or
       \( a_{1} > x^{2n} + y^{2n} > x' \)
or
       \( \ip{A}{B} = \sum_{i} a_{i} b_{i} \).
The spaces you type in a formula are
ignored.  Remember that a letter like
       $x$                   % $ ... $  and  \( ... \)  are equivalent
is a formula when it denotes a mathematical
symbol, and it should be typed as one.

\section{Displayed Text}

Text is displayed by indenting it from the left
margin.  Quotations are commonly displayed.  There
are short quotations
\begin{quote}
   This is a short quotation.  It consists of a
   single paragraph of text.  See how it is formatted.
\end{quote}
and longer ones.
\begin{quotation}
   This is a longer quotation.  It consists of two
   paragraphs of text, neither of which are
   particularly interesting.

   This is the second paragraph of the quotation.  It
   is just as dull as the first paragraph.
\end{quotation}
Another frequently-displayed structure is a list.
The following is an example of an \emph{itemized}
list.
\begin{itemize}
   \item This is the first item of an itemized list.
         Each item in the list is marked with a ``tick''.
         You don't have to worry about what kind of tick
         mark is used.

   \item This is the second item of the list.  It
         contains another list nested inside it.  The inner
         list is an \emph{enumerated} list.
         \begin{enumerate}
            \item This is the first item of an enumerated
                  list that is nested within the itemized list.

            \item This is the second item of the inner list.
                  \LaTeX\ allows you to nest lists deeper than
                  you really should.
         \end{enumerate}
         This is the rest of the second item of the outer
         list.  It is no more interesting than any other
         part of the item.
   \item This is the third item of the list.
\end{itemize}
You can even display poetry.
\begin{verse}
   There is an environment
    for verse \\             % The \\ command separates lines
   Whose features some poets % within a stanza.
   will curse.

                             % One or more blank lines separate stanzas.

   For instead of making\\
   Them do \emph{all} line breaking, \\
   It allows them to put too many words on a line when they'd rather be
   forced to be terse.
\end{verse}

Mathematical formulas may also be displayed.  A
displayed formula
is
one-line long; multiline
formulas require special formatting instructions.
   \[  \ip{\Gamma}{\psi'} = x'' + y^{2} + z_{i}^{n}\]
Don't start a paragraph with a displayed equation,
nor make one a paragraph by itself.

\end{document}               % End of document.

\section{Bench Huge Section 1154}
% synthetic bench section 1154
% This is a sample LaTeX input file.  (Version of 12 August 2004.)
%
% A '%' character causes TeX to ignore all remaining text on the line,
% and is used for comments like this one.

\documentclass{article}      % Specifies the document class

                             % The preamble begins here.
\title{An Example Document}  % Declares the document's title.
\author{Leslie Lamport}      % Declares the author's name.
\date{January 21, 1994}      % Deleting this command produces today's date.

\newcommand{\ip}[2]{(#1, #2)}
                             % Defines \ip{arg1}{arg2} to mean
                             % (arg1, arg2).

%\newcommand{\ip}[2]{\langle #1 | #2\rangle}
                             % This is an alternative definition of
                             % \ip that is commented out.

\begin{document}             % End of preamble and beginning of text.

\maketitle                   % Produces the title.

This is an example input file.  Comparing it with
the output it generates can show you how to
produce a simple document of your own.

\section{Ordinary Text}      % Produces section heading.  Lower-level
                             % sections are begun with similar
                             % \subsection and \subsubsection commands.

The ends  of words and sentences are marked
  by   spaces. It  doesn't matter how many
spaces    you type; one is as good as 100.  The
end of   a line counts as a space.

One   or more   blank lines denote the  end
of  a paragraph.

Since any number of consecutive spaces are treated
like a single one, the formatting of the input
file makes no difference to
      \LaTeX,                % The \LaTeX command generates the LaTeX logo.
but it makes a difference to you.  When you use
\LaTeX, making your input file as easy to read
as possible will be a great help as you write
your document and when you change it.  This sample
file shows how you can add comments to your own input
file.

Because printing is different from typewriting,
there are a number of things that you have to do
differently when preparing an input file than if
you were just typing the document directly.
Quotation marks like
       ``this''
have to be handled specially, as do quotes within
quotes:
       ``\,`this'            % \, separates the double and single quote.
        is what I just
        wrote, not  `that'\,''.

Dashes come in three sizes: an
       intra-word
dash, a medium dash for number ranges like
       1--2,
and a punctuation
       dash---like
this.

A sentence-ending space should be larger than the
space between words within a sentence.  You
sometimes have to type special commands in
conjunction with punctuation characters to get
this right, as in the following sentence.
       Gnats, gnus, etc.\ all  % `\ ' makes an inter-word space.
       begin with G\@.         % \@ marks end-of-sentence punctuation.
You should check the spaces after periods when
reading your output to make sure you haven't
forgotten any special cases.  Generating an
ellipsis
       \ldots\               % `\ ' is needed after `\ldots' because TeX
                             % ignores spaces after command names like \ldots
                             % made from \ + letters.
                             %
                             % Note how a `%' character causes TeX to ignore
                             % the end of the input line, so these blank lines
                             % do not start a new paragraph.
                             %
with the right spacing around the periods requires
a special command.

\LaTeX\ interprets some common characters as
commands, so you must type special commands to
generate them.  These characters include the
following:
       \$ \& \% \# \{ and \}.

In printing, text is usually emphasized with an
       \emph{italic}
type style.

\begin{em}
   A long segment of text can also be emphasized
   in this way.  Text within such a segment can be
   given \emph{additional} emphasis.
\end{em}

It is sometimes necessary to prevent \LaTeX\ from
breaking a line where it might otherwise do so.
This may be at a space, as between the ``Mr.''\ and
``Jones'' in
       ``Mr.~Jones'',        % ~ produces an unbreakable interword space.
or within a word---especially when the word is a
symbol like
       \mbox{\emph{itemnum}}
that makes little sense when hyphenated across
lines.

Footnotes\footnote{This is an example of a footnote.}
pose no problem.

\LaTeX\ is good at typesetting mathematical formulas
like
       \( x-3y + z = 7 \)
or
       \( a_{1} > x^{2n} + y^{2n} > x' \)
or
       \( \ip{A}{B} = \sum_{i} a_{i} b_{i} \).
The spaces you type in a formula are
ignored.  Remember that a letter like
       $x$                   % $ ... $  and  \( ... \)  are equivalent
is a formula when it denotes a mathematical
symbol, and it should be typed as one.

\section{Displayed Text}

Text is displayed by indenting it from the left
margin.  Quotations are commonly displayed.  There
are short quotations
\begin{quote}
   This is a short quotation.  It consists of a
   single paragraph of text.  See how it is formatted.
\end{quote}
and longer ones.
\begin{quotation}
   This is a longer quotation.  It consists of two
   paragraphs of text, neither of which are
   particularly interesting.

   This is the second paragraph of the quotation.  It
   is just as dull as the first paragraph.
\end{quotation}
Another frequently-displayed structure is a list.
The following is an example of an \emph{itemized}
list.
\begin{itemize}
   \item This is the first item of an itemized list.
         Each item in the list is marked with a ``tick''.
         You don't have to worry about what kind of tick
         mark is used.

   \item This is the second item of the list.  It
         contains another list nested inside it.  The inner
         list is an \emph{enumerated} list.
         \begin{enumerate}
            \item This is the first item of an enumerated
                  list that is nested within the itemized list.

            \item This is the second item of the inner list.
                  \LaTeX\ allows you to nest lists deeper than
                  you really should.
         \end{enumerate}
         This is the rest of the second item of the outer
         list.  It is no more interesting than any other
         part of the item.
   \item This is the third item of the list.
\end{itemize}
You can even display poetry.
\begin{verse}
   There is an environment
    for verse \\             % The \\ command separates lines
   Whose features some poets % within a stanza.
   will curse.

                             % One or more blank lines separate stanzas.

   For instead of making\\
   Them do \emph{all} line breaking, \\
   It allows them to put too many words on a line when they'd rather be
   forced to be terse.
\end{verse}

Mathematical formulas may also be displayed.  A
displayed formula
is
one-line long; multiline
formulas require special formatting instructions.
   \[  \ip{\Gamma}{\psi'} = x'' + y^{2} + z_{i}^{n}\]
Don't start a paragraph with a displayed equation,
nor make one a paragraph by itself.

\end{document}               % End of document.

\section{Bench Huge Section 1155}
% synthetic bench section 1155
% This is a sample LaTeX input file.  (Version of 12 August 2004.)
%
% A '%' character causes TeX to ignore all remaining text on the line,
% and is used for comments like this one.

\documentclass{article}      % Specifies the document class

                             % The preamble begins here.
\title{An Example Document}  % Declares the document's title.
\author{Leslie Lamport}      % Declares the author's name.
\date{January 21, 1994}      % Deleting this command produces today's date.

\newcommand{\ip}[2]{(#1, #2)}
                             % Defines \ip{arg1}{arg2} to mean
                             % (arg1, arg2).

%\newcommand{\ip}[2]{\langle #1 | #2\rangle}
                             % This is an alternative definition of
                             % \ip that is commented out.

\begin{document}             % End of preamble and beginning of text.

\maketitle                   % Produces the title.

This is an example input file.  Comparing it with
the output it generates can show you how to
produce a simple document of your own.

\section{Ordinary Text}      % Produces section heading.  Lower-level
                             % sections are begun with similar
                             % \subsection and \subsubsection commands.

The ends  of words and sentences are marked
  by   spaces. It  doesn't matter how many
spaces    you type; one is as good as 100.  The
end of   a line counts as a space.

One   or more   blank lines denote the  end
of  a paragraph.

Since any number of consecutive spaces are treated
like a single one, the formatting of the input
file makes no difference to
      \LaTeX,                % The \LaTeX command generates the LaTeX logo.
but it makes a difference to you.  When you use
\LaTeX, making your input file as easy to read
as possible will be a great help as you write
your document and when you change it.  This sample
file shows how you can add comments to your own input
file.

Because printing is different from typewriting,
there are a number of things that you have to do
differently when preparing an input file than if
you were just typing the document directly.
Quotation marks like
       ``this''
have to be handled specially, as do quotes within
quotes:
       ``\,`this'            % \, separates the double and single quote.
        is what I just
        wrote, not  `that'\,''.

Dashes come in three sizes: an
       intra-word
dash, a medium dash for number ranges like
       1--2,
and a punctuation
       dash---like
this.

A sentence-ending space should be larger than the
space between words within a sentence.  You
sometimes have to type special commands in
conjunction with punctuation characters to get
this right, as in the following sentence.
       Gnats, gnus, etc.\ all  % `\ ' makes an inter-word space.
       begin with G\@.         % \@ marks end-of-sentence punctuation.
You should check the spaces after periods when
reading your output to make sure you haven't
forgotten any special cases.  Generating an
ellipsis
       \ldots\               % `\ ' is needed after `\ldots' because TeX
                             % ignores spaces after command names like \ldots
                             % made from \ + letters.
                             %
                             % Note how a `%' character causes TeX to ignore
                             % the end of the input line, so these blank lines
                             % do not start a new paragraph.
                             %
with the right spacing around the periods requires
a special command.

\LaTeX\ interprets some common characters as
commands, so you must type special commands to
generate them.  These characters include the
following:
       \$ \& \% \# \{ and \}.

In printing, text is usually emphasized with an
       \emph{italic}
type style.

\begin{em}
   A long segment of text can also be emphasized
   in this way.  Text within such a segment can be
   given \emph{additional} emphasis.
\end{em}

It is sometimes necessary to prevent \LaTeX\ from
breaking a line where it might otherwise do so.
This may be at a space, as between the ``Mr.''\ and
``Jones'' in
       ``Mr.~Jones'',        % ~ produces an unbreakable interword space.
or within a word---especially when the word is a
symbol like
       \mbox{\emph{itemnum}}
that makes little sense when hyphenated across
lines.

Footnotes\footnote{This is an example of a footnote.}
pose no problem.

\LaTeX\ is good at typesetting mathematical formulas
like
       \( x-3y + z = 7 \)
or
       \( a_{1} > x^{2n} + y^{2n} > x' \)
or
       \( \ip{A}{B} = \sum_{i} a_{i} b_{i} \).
The spaces you type in a formula are
ignored.  Remember that a letter like
       $x$                   % $ ... $  and  \( ... \)  are equivalent
is a formula when it denotes a mathematical
symbol, and it should be typed as one.

\section{Displayed Text}

Text is displayed by indenting it from the left
margin.  Quotations are commonly displayed.  There
are short quotations
\begin{quote}
   This is a short quotation.  It consists of a
   single paragraph of text.  See how it is formatted.
\end{quote}
and longer ones.
\begin{quotation}
   This is a longer quotation.  It consists of two
   paragraphs of text, neither of which are
   particularly interesting.

   This is the second paragraph of the quotation.  It
   is just as dull as the first paragraph.
\end{quotation}
Another frequently-displayed structure is a list.
The following is an example of an \emph{itemized}
list.
\begin{itemize}
   \item This is the first item of an itemized list.
         Each item in the list is marked with a ``tick''.
         You don't have to worry about what kind of tick
         mark is used.

   \item This is the second item of the list.  It
         contains another list nested inside it.  The inner
         list is an \emph{enumerated} list.
         \begin{enumerate}
            \item This is the first item of an enumerated
                  list that is nested within the itemized list.

            \item This is the second item of the inner list.
                  \LaTeX\ allows you to nest lists deeper than
                  you really should.
         \end{enumerate}
         This is the rest of the second item of the outer
         list.  It is no more interesting than any other
         part of the item.
   \item This is the third item of the list.
\end{itemize}
You can even display poetry.
\begin{verse}
   There is an environment
    for verse \\             % The \\ command separates lines
   Whose features some poets % within a stanza.
   will curse.

                             % One or more blank lines separate stanzas.

   For instead of making\\
   Them do \emph{all} line breaking, \\
   It allows them to put too many words on a line when they'd rather be
   forced to be terse.
\end{verse}

Mathematical formulas may also be displayed.  A
displayed formula
is
one-line long; multiline
formulas require special formatting instructions.
   \[  \ip{\Gamma}{\psi'} = x'' + y^{2} + z_{i}^{n}\]
Don't start a paragraph with a displayed equation,
nor make one a paragraph by itself.

\end{document}               % End of document.

\section{Bench Huge Section 1156}
% synthetic bench section 1156
% This is a sample LaTeX input file.  (Version of 12 August 2004.)
%
% A '%' character causes TeX to ignore all remaining text on the line,
% and is used for comments like this one.

\documentclass{article}      % Specifies the document class

                             % The preamble begins here.
\title{An Example Document}  % Declares the document's title.
\author{Leslie Lamport}      % Declares the author's name.
\date{January 21, 1994}      % Deleting this command produces today's date.

\newcommand{\ip}[2]{(#1, #2)}
                             % Defines \ip{arg1}{arg2} to mean
                             % (arg1, arg2).

%\newcommand{\ip}[2]{\langle #1 | #2\rangle}
                             % This is an alternative definition of
                             % \ip that is commented out.

\begin{document}             % End of preamble and beginning of text.

\maketitle                   % Produces the title.

This is an example input file.  Comparing it with
the output it generates can show you how to
produce a simple document of your own.

\section{Ordinary Text}      % Produces section heading.  Lower-level
                             % sections are begun with similar
                             % \subsection and \subsubsection commands.

The ends  of words and sentences are marked
  by   spaces. It  doesn't matter how many
spaces    you type; one is as good as 100.  The
end of   a line counts as a space.

One   or more   blank lines denote the  end
of  a paragraph.

Since any number of consecutive spaces are treated
like a single one, the formatting of the input
file makes no difference to
      \LaTeX,                % The \LaTeX command generates the LaTeX logo.
but it makes a difference to you.  When you use
\LaTeX, making your input file as easy to read
as possible will be a great help as you write
your document and when you change it.  This sample
file shows how you can add comments to your own input
file.

Because printing is different from typewriting,
there are a number of things that you have to do
differently when preparing an input file than if
you were just typing the document directly.
Quotation marks like
       ``this''
have to be handled specially, as do quotes within
quotes:
       ``\,`this'            % \, separates the double and single quote.
        is what I just
        wrote, not  `that'\,''.

Dashes come in three sizes: an
       intra-word
dash, a medium dash for number ranges like
       1--2,
and a punctuation
       dash---like
this.

A sentence-ending space should be larger than the
space between words within a sentence.  You
sometimes have to type special commands in
conjunction with punctuation characters to get
this right, as in the following sentence.
       Gnats, gnus, etc.\ all  % `\ ' makes an inter-word space.
       begin with G\@.         % \@ marks end-of-sentence punctuation.
You should check the spaces after periods when
reading your output to make sure you haven't
forgotten any special cases.  Generating an
ellipsis
       \ldots\               % `\ ' is needed after `\ldots' because TeX
                             % ignores spaces after command names like \ldots
                             % made from \ + letters.
                             %
                             % Note how a `%' character causes TeX to ignore
                             % the end of the input line, so these blank lines
                             % do not start a new paragraph.
                             %
with the right spacing around the periods requires
a special command.

\LaTeX\ interprets some common characters as
commands, so you must type special commands to
generate them.  These characters include the
following:
       \$ \& \% \# \{ and \}.

In printing, text is usually emphasized with an
       \emph{italic}
type style.

\begin{em}
   A long segment of text can also be emphasized
   in this way.  Text within such a segment can be
   given \emph{additional} emphasis.
\end{em}

It is sometimes necessary to prevent \LaTeX\ from
breaking a line where it might otherwise do so.
This may be at a space, as between the ``Mr.''\ and
``Jones'' in
       ``Mr.~Jones'',        % ~ produces an unbreakable interword space.
or within a word---especially when the word is a
symbol like
       \mbox{\emph{itemnum}}
that makes little sense when hyphenated across
lines.

Footnotes\footnote{This is an example of a footnote.}
pose no problem.

\LaTeX\ is good at typesetting mathematical formulas
like
       \( x-3y + z = 7 \)
or
       \( a_{1} > x^{2n} + y^{2n} > x' \)
or
       \( \ip{A}{B} = \sum_{i} a_{i} b_{i} \).
The spaces you type in a formula are
ignored.  Remember that a letter like
       $x$                   % $ ... $  and  \( ... \)  are equivalent
is a formula when it denotes a mathematical
symbol, and it should be typed as one.

\section{Displayed Text}

Text is displayed by indenting it from the left
margin.  Quotations are commonly displayed.  There
are short quotations
\begin{quote}
   This is a short quotation.  It consists of a
   single paragraph of text.  See how it is formatted.
\end{quote}
and longer ones.
\begin{quotation}
   This is a longer quotation.  It consists of two
   paragraphs of text, neither of which are
   particularly interesting.

   This is the second paragraph of the quotation.  It
   is just as dull as the first paragraph.
\end{quotation}
Another frequently-displayed structure is a list.
The following is an example of an \emph{itemized}
list.
\begin{itemize}
   \item This is the first item of an itemized list.
         Each item in the list is marked with a ``tick''.
         You don't have to worry about what kind of tick
         mark is used.

   \item This is the second item of the list.  It
         contains another list nested inside it.  The inner
         list is an \emph{enumerated} list.
         \begin{enumerate}
            \item This is the first item of an enumerated
                  list that is nested within the itemized list.

            \item This is the second item of the inner list.
                  \LaTeX\ allows you to nest lists deeper than
                  you really should.
         \end{enumerate}
         This is the rest of the second item of the outer
         list.  It is no more interesting than any other
         part of the item.
   \item This is the third item of the list.
\end{itemize}
You can even display poetry.
\begin{verse}
   There is an environment
    for verse \\             % The \\ command separates lines
   Whose features some poets % within a stanza.
   will curse.

                             % One or more blank lines separate stanzas.

   For instead of making\\
   Them do \emph{all} line breaking, \\
   It allows them to put too many words on a line when they'd rather be
   forced to be terse.
\end{verse}

Mathematical formulas may also be displayed.  A
displayed formula
is
one-line long; multiline
formulas require special formatting instructions.
   \[  \ip{\Gamma}{\psi'} = x'' + y^{2} + z_{i}^{n}\]
Don't start a paragraph with a displayed equation,
nor make one a paragraph by itself.

\end{document}               % End of document.

\section{Bench Huge Section 1157}
% synthetic bench section 1157
% This is a sample LaTeX input file.  (Version of 12 August 2004.)
%
% A '%' character causes TeX to ignore all remaining text on the line,
% and is used for comments like this one.

\documentclass{article}      % Specifies the document class

                             % The preamble begins here.
\title{An Example Document}  % Declares the document's title.
\author{Leslie Lamport}      % Declares the author's name.
\date{January 21, 1994}      % Deleting this command produces today's date.

\newcommand{\ip}[2]{(#1, #2)}
                             % Defines \ip{arg1}{arg2} to mean
                             % (arg1, arg2).

%\newcommand{\ip}[2]{\langle #1 | #2\rangle}
                             % This is an alternative definition of
                             % \ip that is commented out.

\begin{document}             % End of preamble and beginning of text.

\maketitle                   % Produces the title.

This is an example input file.  Comparing it with
the output it generates can show you how to
produce a simple document of your own.

\section{Ordinary Text}      % Produces section heading.  Lower-level
                             % sections are begun with similar
                             % \subsection and \subsubsection commands.

The ends  of words and sentences are marked
  by   spaces. It  doesn't matter how many
spaces    you type; one is as good as 100.  The
end of   a line counts as a space.

One   or more   blank lines denote the  end
of  a paragraph.

Since any number of consecutive spaces are treated
like a single one, the formatting of the input
file makes no difference to
      \LaTeX,                % The \LaTeX command generates the LaTeX logo.
but it makes a difference to you.  When you use
\LaTeX, making your input file as easy to read
as possible will be a great help as you write
your document and when you change it.  This sample
file shows how you can add comments to your own input
file.

Because printing is different from typewriting,
there are a number of things that you have to do
differently when preparing an input file than if
you were just typing the document directly.
Quotation marks like
       ``this''
have to be handled specially, as do quotes within
quotes:
       ``\,`this'            % \, separates the double and single quote.
        is what I just
        wrote, not  `that'\,''.

Dashes come in three sizes: an
       intra-word
dash, a medium dash for number ranges like
       1--2,
and a punctuation
       dash---like
this.

A sentence-ending space should be larger than the
space between words within a sentence.  You
sometimes have to type special commands in
conjunction with punctuation characters to get
this right, as in the following sentence.
       Gnats, gnus, etc.\ all  % `\ ' makes an inter-word space.
       begin with G\@.         % \@ marks end-of-sentence punctuation.
You should check the spaces after periods when
reading your output to make sure you haven't
forgotten any special cases.  Generating an
ellipsis
       \ldots\               % `\ ' is needed after `\ldots' because TeX
                             % ignores spaces after command names like \ldots
                             % made from \ + letters.
                             %
                             % Note how a `%' character causes TeX to ignore
                             % the end of the input line, so these blank lines
                             % do not start a new paragraph.
                             %
with the right spacing around the periods requires
a special command.

\LaTeX\ interprets some common characters as
commands, so you must type special commands to
generate them.  These characters include the
following:
       \$ \& \% \# \{ and \}.

In printing, text is usually emphasized with an
       \emph{italic}
type style.

\begin{em}
   A long segment of text can also be emphasized
   in this way.  Text within such a segment can be
   given \emph{additional} emphasis.
\end{em}

It is sometimes necessary to prevent \LaTeX\ from
breaking a line where it might otherwise do so.
This may be at a space, as between the ``Mr.''\ and
``Jones'' in
       ``Mr.~Jones'',        % ~ produces an unbreakable interword space.
or within a word---especially when the word is a
symbol like
       \mbox{\emph{itemnum}}
that makes little sense when hyphenated across
lines.

Footnotes\footnote{This is an example of a footnote.}
pose no problem.

\LaTeX\ is good at typesetting mathematical formulas
like
       \( x-3y + z = 7 \)
or
       \( a_{1} > x^{2n} + y^{2n} > x' \)
or
       \( \ip{A}{B} = \sum_{i} a_{i} b_{i} \).
The spaces you type in a formula are
ignored.  Remember that a letter like
       $x$                   % $ ... $  and  \( ... \)  are equivalent
is a formula when it denotes a mathematical
symbol, and it should be typed as one.

\section{Displayed Text}

Text is displayed by indenting it from the left
margin.  Quotations are commonly displayed.  There
are short quotations
\begin{quote}
   This is a short quotation.  It consists of a
   single paragraph of text.  See how it is formatted.
\end{quote}
and longer ones.
\begin{quotation}
   This is a longer quotation.  It consists of two
   paragraphs of text, neither of which are
   particularly interesting.

   This is the second paragraph of the quotation.  It
   is just as dull as the first paragraph.
\end{quotation}
Another frequently-displayed structure is a list.
The following is an example of an \emph{itemized}
list.
\begin{itemize}
   \item This is the first item of an itemized list.
         Each item in the list is marked with a ``tick''.
         You don't have to worry about what kind of tick
         mark is used.

   \item This is the second item of the list.  It
         contains another list nested inside it.  The inner
         list is an \emph{enumerated} list.
         \begin{enumerate}
            \item This is the first item of an enumerated
                  list that is nested within the itemized list.

            \item This is the second item of the inner list.
                  \LaTeX\ allows you to nest lists deeper than
                  you really should.
         \end{enumerate}
         This is the rest of the second item of the outer
         list.  It is no more interesting than any other
         part of the item.
   \item This is the third item of the list.
\end{itemize}
You can even display poetry.
\begin{verse}
   There is an environment
    for verse \\             % The \\ command separates lines
   Whose features some poets % within a stanza.
   will curse.

                             % One or more blank lines separate stanzas.

   For instead of making\\
   Them do \emph{all} line breaking, \\
   It allows them to put too many words on a line when they'd rather be
   forced to be terse.
\end{verse}

Mathematical formulas may also be displayed.  A
displayed formula
is
one-line long; multiline
formulas require special formatting instructions.
   \[  \ip{\Gamma}{\psi'} = x'' + y^{2} + z_{i}^{n}\]
Don't start a paragraph with a displayed equation,
nor make one a paragraph by itself.

\end{document}               % End of document.

\section{Bench Huge Section 1158}
% synthetic bench section 1158
% This is a sample LaTeX input file.  (Version of 12 August 2004.)
%
% A '%' character causes TeX to ignore all remaining text on the line,
% and is used for comments like this one.

\documentclass{article}      % Specifies the document class

                             % The preamble begins here.
\title{An Example Document}  % Declares the document's title.
\author{Leslie Lamport}      % Declares the author's name.
\date{January 21, 1994}      % Deleting this command produces today's date.

\newcommand{\ip}[2]{(#1, #2)}
                             % Defines \ip{arg1}{arg2} to mean
                             % (arg1, arg2).

%\newcommand{\ip}[2]{\langle #1 | #2\rangle}
                             % This is an alternative definition of
                             % \ip that is commented out.

\begin{document}             % End of preamble and beginning of text.

\maketitle                   % Produces the title.

This is an example input file.  Comparing it with
the output it generates can show you how to
produce a simple document of your own.

\section{Ordinary Text}      % Produces section heading.  Lower-level
                             % sections are begun with similar
                             % \subsection and \subsubsection commands.

The ends  of words and sentences are marked
  by   spaces. It  doesn't matter how many
spaces    you type; one is as good as 100.  The
end of   a line counts as a space.

One   or more   blank lines denote the  end
of  a paragraph.

Since any number of consecutive spaces are treated
like a single one, the formatting of the input
file makes no difference to
      \LaTeX,                % The \LaTeX command generates the LaTeX logo.
but it makes a difference to you.  When you use
\LaTeX, making your input file as easy to read
as possible will be a great help as you write
your document and when you change it.  This sample
file shows how you can add comments to your own input
file.

Because printing is different from typewriting,
there are a number of things that you have to do
differently when preparing an input file than if
you were just typing the document directly.
Quotation marks like
       ``this''
have to be handled specially, as do quotes within
quotes:
       ``\,`this'            % \, separates the double and single quote.
        is what I just
        wrote, not  `that'\,''.

Dashes come in three sizes: an
       intra-word
dash, a medium dash for number ranges like
       1--2,
and a punctuation
       dash---like
this.

A sentence-ending space should be larger than the
space between words within a sentence.  You
sometimes have to type special commands in
conjunction with punctuation characters to get
this right, as in the following sentence.
       Gnats, gnus, etc.\ all  % `\ ' makes an inter-word space.
       begin with G\@.         % \@ marks end-of-sentence punctuation.
You should check the spaces after periods when
reading your output to make sure you haven't
forgotten any special cases.  Generating an
ellipsis
       \ldots\               % `\ ' is needed after `\ldots' because TeX
                             % ignores spaces after command names like \ldots
                             % made from \ + letters.
                             %
                             % Note how a `%' character causes TeX to ignore
                             % the end of the input line, so these blank lines
                             % do not start a new paragraph.
                             %
with the right spacing around the periods requires
a special command.

\LaTeX\ interprets some common characters as
commands, so you must type special commands to
generate them.  These characters include the
following:
       \$ \& \% \# \{ and \}.

In printing, text is usually emphasized with an
       \emph{italic}
type style.

\begin{em}
   A long segment of text can also be emphasized
   in this way.  Text within such a segment can be
   given \emph{additional} emphasis.
\end{em}

It is sometimes necessary to prevent \LaTeX\ from
breaking a line where it might otherwise do so.
This may be at a space, as between the ``Mr.''\ and
``Jones'' in
       ``Mr.~Jones'',        % ~ produces an unbreakable interword space.
or within a word---especially when the word is a
symbol like
       \mbox{\emph{itemnum}}
that makes little sense when hyphenated across
lines.

Footnotes\footnote{This is an example of a footnote.}
pose no problem.

\LaTeX\ is good at typesetting mathematical formulas
like
       \( x-3y + z = 7 \)
or
       \( a_{1} > x^{2n} + y^{2n} > x' \)
or
       \( \ip{A}{B} = \sum_{i} a_{i} b_{i} \).
The spaces you type in a formula are
ignored.  Remember that a letter like
       $x$                   % $ ... $  and  \( ... \)  are equivalent
is a formula when it denotes a mathematical
symbol, and it should be typed as one.

\section{Displayed Text}

Text is displayed by indenting it from the left
margin.  Quotations are commonly displayed.  There
are short quotations
\begin{quote}
   This is a short quotation.  It consists of a
   single paragraph of text.  See how it is formatted.
\end{quote}
and longer ones.
\begin{quotation}
   This is a longer quotation.  It consists of two
   paragraphs of text, neither of which are
   particularly interesting.

   This is the second paragraph of the quotation.  It
   is just as dull as the first paragraph.
\end{quotation}
Another frequently-displayed structure is a list.
The following is an example of an \emph{itemized}
list.
\begin{itemize}
   \item This is the first item of an itemized list.
         Each item in the list is marked with a ``tick''.
         You don't have to worry about what kind of tick
         mark is used.

   \item This is the second item of the list.  It
         contains another list nested inside it.  The inner
         list is an \emph{enumerated} list.
         \begin{enumerate}
            \item This is the first item of an enumerated
                  list that is nested within the itemized list.

            \item This is the second item of the inner list.
                  \LaTeX\ allows you to nest lists deeper than
                  you really should.
         \end{enumerate}
         This is the rest of the second item of the outer
         list.  It is no more interesting than any other
         part of the item.
   \item This is the third item of the list.
\end{itemize}
You can even display poetry.
\begin{verse}
   There is an environment
    for verse \\             % The \\ command separates lines
   Whose features some poets % within a stanza.
   will curse.

                             % One or more blank lines separate stanzas.

   For instead of making\\
   Them do \emph{all} line breaking, \\
   It allows them to put too many words on a line when they'd rather be
   forced to be terse.
\end{verse}

Mathematical formulas may also be displayed.  A
displayed formula
is
one-line long; multiline
formulas require special formatting instructions.
   \[  \ip{\Gamma}{\psi'} = x'' + y^{2} + z_{i}^{n}\]
Don't start a paragraph with a displayed equation,
nor make one a paragraph by itself.

\end{document}               % End of document.

\section{Bench Huge Section 1159}
% synthetic bench section 1159
% This is a sample LaTeX input file.  (Version of 12 August 2004.)
%
% A '%' character causes TeX to ignore all remaining text on the line,
% and is used for comments like this one.

\documentclass{article}      % Specifies the document class

                             % The preamble begins here.
\title{An Example Document}  % Declares the document's title.
\author{Leslie Lamport}      % Declares the author's name.
\date{January 21, 1994}      % Deleting this command produces today's date.

\newcommand{\ip}[2]{(#1, #2)}
                             % Defines \ip{arg1}{arg2} to mean
                             % (arg1, arg2).

%\newcommand{\ip}[2]{\langle #1 | #2\rangle}
                             % This is an alternative definition of
                             % \ip that is commented out.

\begin{document}             % End of preamble and beginning of text.

\maketitle                   % Produces the title.

This is an example input file.  Comparing it with
the output it generates can show you how to
produce a simple document of your own.

\section{Ordinary Text}      % Produces section heading.  Lower-level
                             % sections are begun with similar
                             % \subsection and \subsubsection commands.

The ends  of words and sentences are marked
  by   spaces. It  doesn't matter how many
spaces    you type; one is as good as 100.  The
end of   a line counts as a space.

One   or more   blank lines denote the  end
of  a paragraph.

Since any number of consecutive spaces are treated
like a single one, the formatting of the input
file makes no difference to
      \LaTeX,                % The \LaTeX command generates the LaTeX logo.
but it makes a difference to you.  When you use
\LaTeX, making your input file as easy to read
as possible will be a great help as you write
your document and when you change it.  This sample
file shows how you can add comments to your own input
file.

Because printing is different from typewriting,
there are a number of things that you have to do
differently when preparing an input file than if
you were just typing the document directly.
Quotation marks like
       ``this''
have to be handled specially, as do quotes within
quotes:
       ``\,`this'            % \, separates the double and single quote.
        is what I just
        wrote, not  `that'\,''.

Dashes come in three sizes: an
       intra-word
dash, a medium dash for number ranges like
       1--2,
and a punctuation
       dash---like
this.

A sentence-ending space should be larger than the
space between words within a sentence.  You
sometimes have to type special commands in
conjunction with punctuation characters to get
this right, as in the following sentence.
       Gnats, gnus, etc.\ all  % `\ ' makes an inter-word space.
       begin with G\@.         % \@ marks end-of-sentence punctuation.
You should check the spaces after periods when
reading your output to make sure you haven't
forgotten any special cases.  Generating an
ellipsis
       \ldots\               % `\ ' is needed after `\ldots' because TeX
                             % ignores spaces after command names like \ldots
                             % made from \ + letters.
                             %
                             % Note how a `%' character causes TeX to ignore
                             % the end of the input line, so these blank lines
                             % do not start a new paragraph.
                             %
with the right spacing around the periods requires
a special command.

\LaTeX\ interprets some common characters as
commands, so you must type special commands to
generate them.  These characters include the
following:
       \$ \& \% \# \{ and \}.

In printing, text is usually emphasized with an
       \emph{italic}
type style.

\begin{em}
   A long segment of text can also be emphasized
   in this way.  Text within such a segment can be
   given \emph{additional} emphasis.
\end{em}

It is sometimes necessary to prevent \LaTeX\ from
breaking a line where it might otherwise do so.
This may be at a space, as between the ``Mr.''\ and
``Jones'' in
       ``Mr.~Jones'',        % ~ produces an unbreakable interword space.
or within a word---especially when the word is a
symbol like
       \mbox{\emph{itemnum}}
that makes little sense when hyphenated across
lines.

Footnotes\footnote{This is an example of a footnote.}
pose no problem.

\LaTeX\ is good at typesetting mathematical formulas
like
       \( x-3y + z = 7 \)
or
       \( a_{1} > x^{2n} + y^{2n} > x' \)
or
       \( \ip{A}{B} = \sum_{i} a_{i} b_{i} \).
The spaces you type in a formula are
ignored.  Remember that a letter like
       $x$                   % $ ... $  and  \( ... \)  are equivalent
is a formula when it denotes a mathematical
symbol, and it should be typed as one.

\section{Displayed Text}

Text is displayed by indenting it from the left
margin.  Quotations are commonly displayed.  There
are short quotations
\begin{quote}
   This is a short quotation.  It consists of a
   single paragraph of text.  See how it is formatted.
\end{quote}
and longer ones.
\begin{quotation}
   This is a longer quotation.  It consists of two
   paragraphs of text, neither of which are
   particularly interesting.

   This is the second paragraph of the quotation.  It
   is just as dull as the first paragraph.
\end{quotation}
Another frequently-displayed structure is a list.
The following is an example of an \emph{itemized}
list.
\begin{itemize}
   \item This is the first item of an itemized list.
         Each item in the list is marked with a ``tick''.
         You don't have to worry about what kind of tick
         mark is used.

   \item This is the second item of the list.  It
         contains another list nested inside it.  The inner
         list is an \emph{enumerated} list.
         \begin{enumerate}
            \item This is the first item of an enumerated
                  list that is nested within the itemized list.

            \item This is the second item of the inner list.
                  \LaTeX\ allows you to nest lists deeper than
                  you really should.
         \end{enumerate}
         This is the rest of the second item of the outer
         list.  It is no more interesting than any other
         part of the item.
   \item This is the third item of the list.
\end{itemize}
You can even display poetry.
\begin{verse}
   There is an environment
    for verse \\             % The \\ command separates lines
   Whose features some poets % within a stanza.
   will curse.

                             % One or more blank lines separate stanzas.

   For instead of making\\
   Them do \emph{all} line breaking, \\
   It allows them to put too many words on a line when they'd rather be
   forced to be terse.
\end{verse}

Mathematical formulas may also be displayed.  A
displayed formula
is
one-line long; multiline
formulas require special formatting instructions.
   \[  \ip{\Gamma}{\psi'} = x'' + y^{2} + z_{i}^{n}\]
Don't start a paragraph with a displayed equation,
nor make one a paragraph by itself.

\end{document}               % End of document.

\section{Bench Huge Section 1160}
% synthetic bench section 1160
% This is a sample LaTeX input file.  (Version of 12 August 2004.)
%
% A '%' character causes TeX to ignore all remaining text on the line,
% and is used for comments like this one.

\documentclass{article}      % Specifies the document class

                             % The preamble begins here.
\title{An Example Document}  % Declares the document's title.
\author{Leslie Lamport}      % Declares the author's name.
\date{January 21, 1994}      % Deleting this command produces today's date.

\newcommand{\ip}[2]{(#1, #2)}
                             % Defines \ip{arg1}{arg2} to mean
                             % (arg1, arg2).

%\newcommand{\ip}[2]{\langle #1 | #2\rangle}
                             % This is an alternative definition of
                             % \ip that is commented out.

\begin{document}             % End of preamble and beginning of text.

\maketitle                   % Produces the title.

This is an example input file.  Comparing it with
the output it generates can show you how to
produce a simple document of your own.

\section{Ordinary Text}      % Produces section heading.  Lower-level
                             % sections are begun with similar
                             % \subsection and \subsubsection commands.

The ends  of words and sentences are marked
  by   spaces. It  doesn't matter how many
spaces    you type; one is as good as 100.  The
end of   a line counts as a space.

One   or more   blank lines denote the  end
of  a paragraph.

Since any number of consecutive spaces are treated
like a single one, the formatting of the input
file makes no difference to
      \LaTeX,                % The \LaTeX command generates the LaTeX logo.
but it makes a difference to you.  When you use
\LaTeX, making your input file as easy to read
as possible will be a great help as you write
your document and when you change it.  This sample
file shows how you can add comments to your own input
file.

Because printing is different from typewriting,
there are a number of things that you have to do
differently when preparing an input file than if
you were just typing the document directly.
Quotation marks like
       ``this''
have to be handled specially, as do quotes within
quotes:
       ``\,`this'            % \, separates the double and single quote.
        is what I just
        wrote, not  `that'\,''.

Dashes come in three sizes: an
       intra-word
dash, a medium dash for number ranges like
       1--2,
and a punctuation
       dash---like
this.

A sentence-ending space should be larger than the
space between words within a sentence.  You
sometimes have to type special commands in
conjunction with punctuation characters to get
this right, as in the following sentence.
       Gnats, gnus, etc.\ all  % `\ ' makes an inter-word space.
       begin with G\@.         % \@ marks end-of-sentence punctuation.
You should check the spaces after periods when
reading your output to make sure you haven't
forgotten any special cases.  Generating an
ellipsis
       \ldots\               % `\ ' is needed after `\ldots' because TeX
                             % ignores spaces after command names like \ldots
                             % made from \ + letters.
                             %
                             % Note how a `%' character causes TeX to ignore
                             % the end of the input line, so these blank lines
                             % do not start a new paragraph.
                             %
with the right spacing around the periods requires
a special command.

\LaTeX\ interprets some common characters as
commands, so you must type special commands to
generate them.  These characters include the
following:
       \$ \& \% \# \{ and \}.

In printing, text is usually emphasized with an
       \emph{italic}
type style.

\begin{em}
   A long segment of text can also be emphasized
   in this way.  Text within such a segment can be
   given \emph{additional} emphasis.
\end{em}

It is sometimes necessary to prevent \LaTeX\ from
breaking a line where it might otherwise do so.
This may be at a space, as between the ``Mr.''\ and
``Jones'' in
       ``Mr.~Jones'',        % ~ produces an unbreakable interword space.
or within a word---especially when the word is a
symbol like
       \mbox{\emph{itemnum}}
that makes little sense when hyphenated across
lines.

Footnotes\footnote{This is an example of a footnote.}
pose no problem.

\LaTeX\ is good at typesetting mathematical formulas
like
       \( x-3y + z = 7 \)
or
       \( a_{1} > x^{2n} + y^{2n} > x' \)
or
       \( \ip{A}{B} = \sum_{i} a_{i} b_{i} \).
The spaces you type in a formula are
ignored.  Remember that a letter like
       $x$                   % $ ... $  and  \( ... \)  are equivalent
is a formula when it denotes a mathematical
symbol, and it should be typed as one.

\section{Displayed Text}

Text is displayed by indenting it from the left
margin.  Quotations are commonly displayed.  There
are short quotations
\begin{quote}
   This is a short quotation.  It consists of a
   single paragraph of text.  See how it is formatted.
\end{quote}
and longer ones.
\begin{quotation}
   This is a longer quotation.  It consists of two
   paragraphs of text, neither of which are
   particularly interesting.

   This is the second paragraph of the quotation.  It
   is just as dull as the first paragraph.
\end{quotation}
Another frequently-displayed structure is a list.
The following is an example of an \emph{itemized}
list.
\begin{itemize}
   \item This is the first item of an itemized list.
         Each item in the list is marked with a ``tick''.
         You don't have to worry about what kind of tick
         mark is used.

   \item This is the second item of the list.  It
         contains another list nested inside it.  The inner
         list is an \emph{enumerated} list.
         \begin{enumerate}
            \item This is the first item of an enumerated
                  list that is nested within the itemized list.

            \item This is the second item of the inner list.
                  \LaTeX\ allows you to nest lists deeper than
                  you really should.
         \end{enumerate}
         This is the rest of the second item of the outer
         list.  It is no more interesting than any other
         part of the item.
   \item This is the third item of the list.
\end{itemize}
You can even display poetry.
\begin{verse}
   There is an environment
    for verse \\             % The \\ command separates lines
   Whose features some poets % within a stanza.
   will curse.

                             % One or more blank lines separate stanzas.

   For instead of making\\
   Them do \emph{all} line breaking, \\
   It allows them to put too many words on a line when they'd rather be
   forced to be terse.
\end{verse}

Mathematical formulas may also be displayed.  A
displayed formula
is
one-line long; multiline
formulas require special formatting instructions.
   \[  \ip{\Gamma}{\psi'} = x'' + y^{2} + z_{i}^{n}\]
Don't start a paragraph with a displayed equation,
nor make one a paragraph by itself.

\end{document}               % End of document.

\section{Bench Huge Section 1161}
% synthetic bench section 1161
% This is a sample LaTeX input file.  (Version of 12 August 2004.)
%
% A '%' character causes TeX to ignore all remaining text on the line,
% and is used for comments like this one.

\documentclass{article}      % Specifies the document class

                             % The preamble begins here.
\title{An Example Document}  % Declares the document's title.
\author{Leslie Lamport}      % Declares the author's name.
\date{January 21, 1994}      % Deleting this command produces today's date.

\newcommand{\ip}[2]{(#1, #2)}
                             % Defines \ip{arg1}{arg2} to mean
                             % (arg1, arg2).

%\newcommand{\ip}[2]{\langle #1 | #2\rangle}
                             % This is an alternative definition of
                             % \ip that is commented out.

\begin{document}             % End of preamble and beginning of text.

\maketitle                   % Produces the title.

This is an example input file.  Comparing it with
the output it generates can show you how to
produce a simple document of your own.

\section{Ordinary Text}      % Produces section heading.  Lower-level
                             % sections are begun with similar
                             % \subsection and \subsubsection commands.

The ends  of words and sentences are marked
  by   spaces. It  doesn't matter how many
spaces    you type; one is as good as 100.  The
end of   a line counts as a space.

One   or more   blank lines denote the  end
of  a paragraph.

Since any number of consecutive spaces are treated
like a single one, the formatting of the input
file makes no difference to
      \LaTeX,                % The \LaTeX command generates the LaTeX logo.
but it makes a difference to you.  When you use
\LaTeX, making your input file as easy to read
as possible will be a great help as you write
your document and when you change it.  This sample
file shows how you can add comments to your own input
file.

Because printing is different from typewriting,
there are a number of things that you have to do
differently when preparing an input file than if
you were just typing the document directly.
Quotation marks like
       ``this''
have to be handled specially, as do quotes within
quotes:
       ``\,`this'            % \, separates the double and single quote.
        is what I just
        wrote, not  `that'\,''.

Dashes come in three sizes: an
       intra-word
dash, a medium dash for number ranges like
       1--2,
and a punctuation
       dash---like
this.

A sentence-ending space should be larger than the
space between words within a sentence.  You
sometimes have to type special commands in
conjunction with punctuation characters to get
this right, as in the following sentence.
       Gnats, gnus, etc.\ all  % `\ ' makes an inter-word space.
       begin with G\@.         % \@ marks end-of-sentence punctuation.
You should check the spaces after periods when
reading your output to make sure you haven't
forgotten any special cases.  Generating an
ellipsis
       \ldots\               % `\ ' is needed after `\ldots' because TeX
                             % ignores spaces after command names like \ldots
                             % made from \ + letters.
                             %
                             % Note how a `%' character causes TeX to ignore
                             % the end of the input line, so these blank lines
                             % do not start a new paragraph.
                             %
with the right spacing around the periods requires
a special command.

\LaTeX\ interprets some common characters as
commands, so you must type special commands to
generate them.  These characters include the
following:
       \$ \& \% \# \{ and \}.

In printing, text is usually emphasized with an
       \emph{italic}
type style.

\begin{em}
   A long segment of text can also be emphasized
   in this way.  Text within such a segment can be
   given \emph{additional} emphasis.
\end{em}

It is sometimes necessary to prevent \LaTeX\ from
breaking a line where it might otherwise do so.
This may be at a space, as between the ``Mr.''\ and
``Jones'' in
       ``Mr.~Jones'',        % ~ produces an unbreakable interword space.
or within a word---especially when the word is a
symbol like
       \mbox{\emph{itemnum}}
that makes little sense when hyphenated across
lines.

Footnotes\footnote{This is an example of a footnote.}
pose no problem.

\LaTeX\ is good at typesetting mathematical formulas
like
       \( x-3y + z = 7 \)
or
       \( a_{1} > x^{2n} + y^{2n} > x' \)
or
       \( \ip{A}{B} = \sum_{i} a_{i} b_{i} \).
The spaces you type in a formula are
ignored.  Remember that a letter like
       $x$                   % $ ... $  and  \( ... \)  are equivalent
is a formula when it denotes a mathematical
symbol, and it should be typed as one.

\section{Displayed Text}

Text is displayed by indenting it from the left
margin.  Quotations are commonly displayed.  There
are short quotations
\begin{quote}
   This is a short quotation.  It consists of a
   single paragraph of text.  See how it is formatted.
\end{quote}
and longer ones.
\begin{quotation}
   This is a longer quotation.  It consists of two
   paragraphs of text, neither of which are
   particularly interesting.

   This is the second paragraph of the quotation.  It
   is just as dull as the first paragraph.
\end{quotation}
Another frequently-displayed structure is a list.
The following is an example of an \emph{itemized}
list.
\begin{itemize}
   \item This is the first item of an itemized list.
         Each item in the list is marked with a ``tick''.
         You don't have to worry about what kind of tick
         mark is used.

   \item This is the second item of the list.  It
         contains another list nested inside it.  The inner
         list is an \emph{enumerated} list.
         \begin{enumerate}
            \item This is the first item of an enumerated
                  list that is nested within the itemized list.

            \item This is the second item of the inner list.
                  \LaTeX\ allows you to nest lists deeper than
                  you really should.
         \end{enumerate}
         This is the rest of the second item of the outer
         list.  It is no more interesting than any other
         part of the item.
   \item This is the third item of the list.
\end{itemize}
You can even display poetry.
\begin{verse}
   There is an environment
    for verse \\             % The \\ command separates lines
   Whose features some poets % within a stanza.
   will curse.

                             % One or more blank lines separate stanzas.

   For instead of making\\
   Them do \emph{all} line breaking, \\
   It allows them to put too many words on a line when they'd rather be
   forced to be terse.
\end{verse}

Mathematical formulas may also be displayed.  A
displayed formula
is
one-line long; multiline
formulas require special formatting instructions.
   \[  \ip{\Gamma}{\psi'} = x'' + y^{2} + z_{i}^{n}\]
Don't start a paragraph with a displayed equation,
nor make one a paragraph by itself.

\end{document}               % End of document.

\section{Bench Huge Section 1162}
% synthetic bench section 1162
% This is a sample LaTeX input file.  (Version of 12 August 2004.)
%
% A '%' character causes TeX to ignore all remaining text on the line,
% and is used for comments like this one.

\documentclass{article}      % Specifies the document class

                             % The preamble begins here.
\title{An Example Document}  % Declares the document's title.
\author{Leslie Lamport}      % Declares the author's name.
\date{January 21, 1994}      % Deleting this command produces today's date.

\newcommand{\ip}[2]{(#1, #2)}
                             % Defines \ip{arg1}{arg2} to mean
                             % (arg1, arg2).

%\newcommand{\ip}[2]{\langle #1 | #2\rangle}
                             % This is an alternative definition of
                             % \ip that is commented out.

\begin{document}             % End of preamble and beginning of text.

\maketitle                   % Produces the title.

This is an example input file.  Comparing it with
the output it generates can show you how to
produce a simple document of your own.

\section{Ordinary Text}      % Produces section heading.  Lower-level
                             % sections are begun with similar
                             % \subsection and \subsubsection commands.

The ends  of words and sentences are marked
  by   spaces. It  doesn't matter how many
spaces    you type; one is as good as 100.  The
end of   a line counts as a space.

One   or more   blank lines denote the  end
of  a paragraph.

Since any number of consecutive spaces are treated
like a single one, the formatting of the input
file makes no difference to
      \LaTeX,                % The \LaTeX command generates the LaTeX logo.
but it makes a difference to you.  When you use
\LaTeX, making your input file as easy to read
as possible will be a great help as you write
your document and when you change it.  This sample
file shows how you can add comments to your own input
file.

Because printing is different from typewriting,
there are a number of things that you have to do
differently when preparing an input file than if
you were just typing the document directly.
Quotation marks like
       ``this''
have to be handled specially, as do quotes within
quotes:
       ``\,`this'            % \, separates the double and single quote.
        is what I just
        wrote, not  `that'\,''.

Dashes come in three sizes: an
       intra-word
dash, a medium dash for number ranges like
       1--2,
and a punctuation
       dash---like
this.

A sentence-ending space should be larger than the
space between words within a sentence.  You
sometimes have to type special commands in
conjunction with punctuation characters to get
this right, as in the following sentence.
       Gnats, gnus, etc.\ all  % `\ ' makes an inter-word space.
       begin with G\@.         % \@ marks end-of-sentence punctuation.
You should check the spaces after periods when
reading your output to make sure you haven't
forgotten any special cases.  Generating an
ellipsis
       \ldots\               % `\ ' is needed after `\ldots' because TeX
                             % ignores spaces after command names like \ldots
                             % made from \ + letters.
                             %
                             % Note how a `%' character causes TeX to ignore
                             % the end of the input line, so these blank lines
                             % do not start a new paragraph.
                             %
with the right spacing around the periods requires
a special command.

\LaTeX\ interprets some common characters as
commands, so you must type special commands to
generate them.  These characters include the
following:
       \$ \& \% \# \{ and \}.

In printing, text is usually emphasized with an
       \emph{italic}
type style.

\begin{em}
   A long segment of text can also be emphasized
   in this way.  Text within such a segment can be
   given \emph{additional} emphasis.
\end{em}

It is sometimes necessary to prevent \LaTeX\ from
breaking a line where it might otherwise do so.
This may be at a space, as between the ``Mr.''\ and
``Jones'' in
       ``Mr.~Jones'',        % ~ produces an unbreakable interword space.
or within a word---especially when the word is a
symbol like
       \mbox{\emph{itemnum}}
that makes little sense when hyphenated across
lines.

Footnotes\footnote{This is an example of a footnote.}
pose no problem.

\LaTeX\ is good at typesetting mathematical formulas
like
       \( x-3y + z = 7 \)
or
       \( a_{1} > x^{2n} + y^{2n} > x' \)
or
       \( \ip{A}{B} = \sum_{i} a_{i} b_{i} \).
The spaces you type in a formula are
ignored.  Remember that a letter like
       $x$                   % $ ... $  and  \( ... \)  are equivalent
is a formula when it denotes a mathematical
symbol, and it should be typed as one.

\section{Displayed Text}

Text is displayed by indenting it from the left
margin.  Quotations are commonly displayed.  There
are short quotations
\begin{quote}
   This is a short quotation.  It consists of a
   single paragraph of text.  See how it is formatted.
\end{quote}
and longer ones.
\begin{quotation}
   This is a longer quotation.  It consists of two
   paragraphs of text, neither of which are
   particularly interesting.

   This is the second paragraph of the quotation.  It
   is just as dull as the first paragraph.
\end{quotation}
Another frequently-displayed structure is a list.
The following is an example of an \emph{itemized}
list.
\begin{itemize}
   \item This is the first item of an itemized list.
         Each item in the list is marked with a ``tick''.
         You don't have to worry about what kind of tick
         mark is used.

   \item This is the second item of the list.  It
         contains another list nested inside it.  The inner
         list is an \emph{enumerated} list.
         \begin{enumerate}
            \item This is the first item of an enumerated
                  list that is nested within the itemized list.

            \item This is the second item of the inner list.
                  \LaTeX\ allows you to nest lists deeper than
                  you really should.
         \end{enumerate}
         This is the rest of the second item of the outer
         list.  It is no more interesting than any other
         part of the item.
   \item This is the third item of the list.
\end{itemize}
You can even display poetry.
\begin{verse}
   There is an environment
    for verse \\             % The \\ command separates lines
   Whose features some poets % within a stanza.
   will curse.

                             % One or more blank lines separate stanzas.

   For instead of making\\
   Them do \emph{all} line breaking, \\
   It allows them to put too many words on a line when they'd rather be
   forced to be terse.
\end{verse}

Mathematical formulas may also be displayed.  A
displayed formula
is
one-line long; multiline
formulas require special formatting instructions.
   \[  \ip{\Gamma}{\psi'} = x'' + y^{2} + z_{i}^{n}\]
Don't start a paragraph with a displayed equation,
nor make one a paragraph by itself.

\end{document}               % End of document.

\section{Bench Huge Section 1163}
% synthetic bench section 1163
% This is a sample LaTeX input file.  (Version of 12 August 2004.)
%
% A '%' character causes TeX to ignore all remaining text on the line,
% and is used for comments like this one.

\documentclass{article}      % Specifies the document class

                             % The preamble begins here.
\title{An Example Document}  % Declares the document's title.
\author{Leslie Lamport}      % Declares the author's name.
\date{January 21, 1994}      % Deleting this command produces today's date.

\newcommand{\ip}[2]{(#1, #2)}
                             % Defines \ip{arg1}{arg2} to mean
                             % (arg1, arg2).

%\newcommand{\ip}[2]{\langle #1 | #2\rangle}
                             % This is an alternative definition of
                             % \ip that is commented out.

\begin{document}             % End of preamble and beginning of text.

\maketitle                   % Produces the title.

This is an example input file.  Comparing it with
the output it generates can show you how to
produce a simple document of your own.

\section{Ordinary Text}      % Produces section heading.  Lower-level
                             % sections are begun with similar
                             % \subsection and \subsubsection commands.

The ends  of words and sentences are marked
  by   spaces. It  doesn't matter how many
spaces    you type; one is as good as 100.  The
end of   a line counts as a space.

One   or more   blank lines denote the  end
of  a paragraph.

Since any number of consecutive spaces are treated
like a single one, the formatting of the input
file makes no difference to
      \LaTeX,                % The \LaTeX command generates the LaTeX logo.
but it makes a difference to you.  When you use
\LaTeX, making your input file as easy to read
as possible will be a great help as you write
your document and when you change it.  This sample
file shows how you can add comments to your own input
file.

Because printing is different from typewriting,
there are a number of things that you have to do
differently when preparing an input file than if
you were just typing the document directly.
Quotation marks like
       ``this''
have to be handled specially, as do quotes within
quotes:
       ``\,`this'            % \, separates the double and single quote.
        is what I just
        wrote, not  `that'\,''.

Dashes come in three sizes: an
       intra-word
dash, a medium dash for number ranges like
       1--2,
and a punctuation
       dash---like
this.

A sentence-ending space should be larger than the
space between words within a sentence.  You
sometimes have to type special commands in
conjunction with punctuation characters to get
this right, as in the following sentence.
       Gnats, gnus, etc.\ all  % `\ ' makes an inter-word space.
       begin with G\@.         % \@ marks end-of-sentence punctuation.
You should check the spaces after periods when
reading your output to make sure you haven't
forgotten any special cases.  Generating an
ellipsis
       \ldots\               % `\ ' is needed after `\ldots' because TeX
                             % ignores spaces after command names like \ldots
                             % made from \ + letters.
                             %
                             % Note how a `%' character causes TeX to ignore
                             % the end of the input line, so these blank lines
                             % do not start a new paragraph.
                             %
with the right spacing around the periods requires
a special command.

\LaTeX\ interprets some common characters as
commands, so you must type special commands to
generate them.  These characters include the
following:
       \$ \& \% \# \{ and \}.

In printing, text is usually emphasized with an
       \emph{italic}
type style.

\begin{em}
   A long segment of text can also be emphasized
   in this way.  Text within such a segment can be
   given \emph{additional} emphasis.
\end{em}

It is sometimes necessary to prevent \LaTeX\ from
breaking a line where it might otherwise do so.
This may be at a space, as between the ``Mr.''\ and
``Jones'' in
       ``Mr.~Jones'',        % ~ produces an unbreakable interword space.
or within a word---especially when the word is a
symbol like
       \mbox{\emph{itemnum}}
that makes little sense when hyphenated across
lines.

Footnotes\footnote{This is an example of a footnote.}
pose no problem.

\LaTeX\ is good at typesetting mathematical formulas
like
       \( x-3y + z = 7 \)
or
       \( a_{1} > x^{2n} + y^{2n} > x' \)
or
       \( \ip{A}{B} = \sum_{i} a_{i} b_{i} \).
The spaces you type in a formula are
ignored.  Remember that a letter like
       $x$                   % $ ... $  and  \( ... \)  are equivalent
is a formula when it denotes a mathematical
symbol, and it should be typed as one.

\section{Displayed Text}

Text is displayed by indenting it from the left
margin.  Quotations are commonly displayed.  There
are short quotations
\begin{quote}
   This is a short quotation.  It consists of a
   single paragraph of text.  See how it is formatted.
\end{quote}
and longer ones.
\begin{quotation}
   This is a longer quotation.  It consists of two
   paragraphs of text, neither of which are
   particularly interesting.

   This is the second paragraph of the quotation.  It
   is just as dull as the first paragraph.
\end{quotation}
Another frequently-displayed structure is a list.
The following is an example of an \emph{itemized}
list.
\begin{itemize}
   \item This is the first item of an itemized list.
         Each item in the list is marked with a ``tick''.
         You don't have to worry about what kind of tick
         mark is used.

   \item This is the second item of the list.  It
         contains another list nested inside it.  The inner
         list is an \emph{enumerated} list.
         \begin{enumerate}
            \item This is the first item of an enumerated
                  list that is nested within the itemized list.

            \item This is the second item of the inner list.
                  \LaTeX\ allows you to nest lists deeper than
                  you really should.
         \end{enumerate}
         This is the rest of the second item of the outer
         list.  It is no more interesting than any other
         part of the item.
   \item This is the third item of the list.
\end{itemize}
You can even display poetry.
\begin{verse}
   There is an environment
    for verse \\             % The \\ command separates lines
   Whose features some poets % within a stanza.
   will curse.

                             % One or more blank lines separate stanzas.

   For instead of making\\
   Them do \emph{all} line breaking, \\
   It allows them to put too many words on a line when they'd rather be
   forced to be terse.
\end{verse}

Mathematical formulas may also be displayed.  A
displayed formula
is
one-line long; multiline
formulas require special formatting instructions.
   \[  \ip{\Gamma}{\psi'} = x'' + y^{2} + z_{i}^{n}\]
Don't start a paragraph with a displayed equation,
nor make one a paragraph by itself.

\end{document}               % End of document.

\section{Bench Huge Section 1164}
% synthetic bench section 1164
% This is a sample LaTeX input file.  (Version of 12 August 2004.)
%
% A '%' character causes TeX to ignore all remaining text on the line,
% and is used for comments like this one.

\documentclass{article}      % Specifies the document class

                             % The preamble begins here.
\title{An Example Document}  % Declares the document's title.
\author{Leslie Lamport}      % Declares the author's name.
\date{January 21, 1994}      % Deleting this command produces today's date.

\newcommand{\ip}[2]{(#1, #2)}
                             % Defines \ip{arg1}{arg2} to mean
                             % (arg1, arg2).

%\newcommand{\ip}[2]{\langle #1 | #2\rangle}
                             % This is an alternative definition of
                             % \ip that is commented out.

\begin{document}             % End of preamble and beginning of text.

\maketitle                   % Produces the title.

This is an example input file.  Comparing it with
the output it generates can show you how to
produce a simple document of your own.

\section{Ordinary Text}      % Produces section heading.  Lower-level
                             % sections are begun with similar
                             % \subsection and \subsubsection commands.

The ends  of words and sentences are marked
  by   spaces. It  doesn't matter how many
spaces    you type; one is as good as 100.  The
end of   a line counts as a space.

One   or more   blank lines denote the  end
of  a paragraph.

Since any number of consecutive spaces are treated
like a single one, the formatting of the input
file makes no difference to
      \LaTeX,                % The \LaTeX command generates the LaTeX logo.
but it makes a difference to you.  When you use
\LaTeX, making your input file as easy to read
as possible will be a great help as you write
your document and when you change it.  This sample
file shows how you can add comments to your own input
file.

Because printing is different from typewriting,
there are a number of things that you have to do
differently when preparing an input file than if
you were just typing the document directly.
Quotation marks like
       ``this''
have to be handled specially, as do quotes within
quotes:
       ``\,`this'            % \, separates the double and single quote.
        is what I just
        wrote, not  `that'\,''.

Dashes come in three sizes: an
       intra-word
dash, a medium dash for number ranges like
       1--2,
and a punctuation
       dash---like
this.

A sentence-ending space should be larger than the
space between words within a sentence.  You
sometimes have to type special commands in
conjunction with punctuation characters to get
this right, as in the following sentence.
       Gnats, gnus, etc.\ all  % `\ ' makes an inter-word space.
       begin with G\@.         % \@ marks end-of-sentence punctuation.
You should check the spaces after periods when
reading your output to make sure you haven't
forgotten any special cases.  Generating an
ellipsis
       \ldots\               % `\ ' is needed after `\ldots' because TeX
                             % ignores spaces after command names like \ldots
                             % made from \ + letters.
                             %
                             % Note how a `%' character causes TeX to ignore
                             % the end of the input line, so these blank lines
                             % do not start a new paragraph.
                             %
with the right spacing around the periods requires
a special command.

\LaTeX\ interprets some common characters as
commands, so you must type special commands to
generate them.  These characters include the
following:
       \$ \& \% \# \{ and \}.

In printing, text is usually emphasized with an
       \emph{italic}
type style.

\begin{em}
   A long segment of text can also be emphasized
   in this way.  Text within such a segment can be
   given \emph{additional} emphasis.
\end{em}

It is sometimes necessary to prevent \LaTeX\ from
breaking a line where it might otherwise do so.
This may be at a space, as between the ``Mr.''\ and
``Jones'' in
       ``Mr.~Jones'',        % ~ produces an unbreakable interword space.
or within a word---especially when the word is a
symbol like
       \mbox{\emph{itemnum}}
that makes little sense when hyphenated across
lines.

Footnotes\footnote{This is an example of a footnote.}
pose no problem.

\LaTeX\ is good at typesetting mathematical formulas
like
       \( x-3y + z = 7 \)
or
       \( a_{1} > x^{2n} + y^{2n} > x' \)
or
       \( \ip{A}{B} = \sum_{i} a_{i} b_{i} \).
The spaces you type in a formula are
ignored.  Remember that a letter like
       $x$                   % $ ... $  and  \( ... \)  are equivalent
is a formula when it denotes a mathematical
symbol, and it should be typed as one.

\section{Displayed Text}

Text is displayed by indenting it from the left
margin.  Quotations are commonly displayed.  There
are short quotations
\begin{quote}
   This is a short quotation.  It consists of a
   single paragraph of text.  See how it is formatted.
\end{quote}
and longer ones.
\begin{quotation}
   This is a longer quotation.  It consists of two
   paragraphs of text, neither of which are
   particularly interesting.

   This is the second paragraph of the quotation.  It
   is just as dull as the first paragraph.
\end{quotation}
Another frequently-displayed structure is a list.
The following is an example of an \emph{itemized}
list.
\begin{itemize}
   \item This is the first item of an itemized list.
         Each item in the list is marked with a ``tick''.
         You don't have to worry about what kind of tick
         mark is used.

   \item This is the second item of the list.  It
         contains another list nested inside it.  The inner
         list is an \emph{enumerated} list.
         \begin{enumerate}
            \item This is the first item of an enumerated
                  list that is nested within the itemized list.

            \item This is the second item of the inner list.
                  \LaTeX\ allows you to nest lists deeper than
                  you really should.
         \end{enumerate}
         This is the rest of the second item of the outer
         list.  It is no more interesting than any other
         part of the item.
   \item This is the third item of the list.
\end{itemize}
You can even display poetry.
\begin{verse}
   There is an environment
    for verse \\             % The \\ command separates lines
   Whose features some poets % within a stanza.
   will curse.

                             % One or more blank lines separate stanzas.

   For instead of making\\
   Them do \emph{all} line breaking, \\
   It allows them to put too many words on a line when they'd rather be
   forced to be terse.
\end{verse}

Mathematical formulas may also be displayed.  A
displayed formula
is
one-line long; multiline
formulas require special formatting instructions.
   \[  \ip{\Gamma}{\psi'} = x'' + y^{2} + z_{i}^{n}\]
Don't start a paragraph with a displayed equation,
nor make one a paragraph by itself.

\end{document}               % End of document.

\section{Bench Huge Section 1165}
% synthetic bench section 1165
% This is a sample LaTeX input file.  (Version of 12 August 2004.)
%
% A '%' character causes TeX to ignore all remaining text on the line,
% and is used for comments like this one.

\documentclass{article}      % Specifies the document class

                             % The preamble begins here.
\title{An Example Document}  % Declares the document's title.
\author{Leslie Lamport}      % Declares the author's name.
\date{January 21, 1994}      % Deleting this command produces today's date.

\newcommand{\ip}[2]{(#1, #2)}
                             % Defines \ip{arg1}{arg2} to mean
                             % (arg1, arg2).

%\newcommand{\ip}[2]{\langle #1 | #2\rangle}
                             % This is an alternative definition of
                             % \ip that is commented out.

\begin{document}             % End of preamble and beginning of text.

\maketitle                   % Produces the title.

This is an example input file.  Comparing it with
the output it generates can show you how to
produce a simple document of your own.

\section{Ordinary Text}      % Produces section heading.  Lower-level
                             % sections are begun with similar
                             % \subsection and \subsubsection commands.

The ends  of words and sentences are marked
  by   spaces. It  doesn't matter how many
spaces    you type; one is as good as 100.  The
end of   a line counts as a space.

One   or more   blank lines denote the  end
of  a paragraph.

Since any number of consecutive spaces are treated
like a single one, the formatting of the input
file makes no difference to
      \LaTeX,                % The \LaTeX command generates the LaTeX logo.
but it makes a difference to you.  When you use
\LaTeX, making your input file as easy to read
as possible will be a great help as you write
your document and when you change it.  This sample
file shows how you can add comments to your own input
file.

Because printing is different from typewriting,
there are a number of things that you have to do
differently when preparing an input file than if
you were just typing the document directly.
Quotation marks like
       ``this''
have to be handled specially, as do quotes within
quotes:
       ``\,`this'            % \, separates the double and single quote.
        is what I just
        wrote, not  `that'\,''.

Dashes come in three sizes: an
       intra-word
dash, a medium dash for number ranges like
       1--2,
and a punctuation
       dash---like
this.

A sentence-ending space should be larger than the
space between words within a sentence.  You
sometimes have to type special commands in
conjunction with punctuation characters to get
this right, as in the following sentence.
       Gnats, gnus, etc.\ all  % `\ ' makes an inter-word space.
       begin with G\@.         % \@ marks end-of-sentence punctuation.
You should check the spaces after periods when
reading your output to make sure you haven't
forgotten any special cases.  Generating an
ellipsis
       \ldots\               % `\ ' is needed after `\ldots' because TeX
                             % ignores spaces after command names like \ldots
                             % made from \ + letters.
                             %
                             % Note how a `%' character causes TeX to ignore
                             % the end of the input line, so these blank lines
                             % do not start a new paragraph.
                             %
with the right spacing around the periods requires
a special command.

\LaTeX\ interprets some common characters as
commands, so you must type special commands to
generate them.  These characters include the
following:
       \$ \& \% \# \{ and \}.

In printing, text is usually emphasized with an
       \emph{italic}
type style.

\begin{em}
   A long segment of text can also be emphasized
   in this way.  Text within such a segment can be
   given \emph{additional} emphasis.
\end{em}

It is sometimes necessary to prevent \LaTeX\ from
breaking a line where it might otherwise do so.
This may be at a space, as between the ``Mr.''\ and
``Jones'' in
       ``Mr.~Jones'',        % ~ produces an unbreakable interword space.
or within a word---especially when the word is a
symbol like
       \mbox{\emph{itemnum}}
that makes little sense when hyphenated across
lines.

Footnotes\footnote{This is an example of a footnote.}
pose no problem.

\LaTeX\ is good at typesetting mathematical formulas
like
       \( x-3y + z = 7 \)
or
       \( a_{1} > x^{2n} + y^{2n} > x' \)
or
       \( \ip{A}{B} = \sum_{i} a_{i} b_{i} \).
The spaces you type in a formula are
ignored.  Remember that a letter like
       $x$                   % $ ... $  and  \( ... \)  are equivalent
is a formula when it denotes a mathematical
symbol, and it should be typed as one.

\section{Displayed Text}

Text is displayed by indenting it from the left
margin.  Quotations are commonly displayed.  There
are short quotations
\begin{quote}
   This is a short quotation.  It consists of a
   single paragraph of text.  See how it is formatted.
\end{quote}
and longer ones.
\begin{quotation}
   This is a longer quotation.  It consists of two
   paragraphs of text, neither of which are
   particularly interesting.

   This is the second paragraph of the quotation.  It
   is just as dull as the first paragraph.
\end{quotation}
Another frequently-displayed structure is a list.
The following is an example of an \emph{itemized}
list.
\begin{itemize}
   \item This is the first item of an itemized list.
         Each item in the list is marked with a ``tick''.
         You don't have to worry about what kind of tick
         mark is used.

   \item This is the second item of the list.  It
         contains another list nested inside it.  The inner
         list is an \emph{enumerated} list.
         \begin{enumerate}
            \item This is the first item of an enumerated
                  list that is nested within the itemized list.

            \item This is the second item of the inner list.
                  \LaTeX\ allows you to nest lists deeper than
                  you really should.
         \end{enumerate}
         This is the rest of the second item of the outer
         list.  It is no more interesting than any other
         part of the item.
   \item This is the third item of the list.
\end{itemize}
You can even display poetry.
\begin{verse}
   There is an environment
    for verse \\             % The \\ command separates lines
   Whose features some poets % within a stanza.
   will curse.

                             % One or more blank lines separate stanzas.

   For instead of making\\
   Them do \emph{all} line breaking, \\
   It allows them to put too many words on a line when they'd rather be
   forced to be terse.
\end{verse}

Mathematical formulas may also be displayed.  A
displayed formula
is
one-line long; multiline
formulas require special formatting instructions.
   \[  \ip{\Gamma}{\psi'} = x'' + y^{2} + z_{i}^{n}\]
Don't start a paragraph with a displayed equation,
nor make one a paragraph by itself.

\end{document}               % End of document.

\section{Bench Huge Section 1166}
% synthetic bench section 1166
% This is a sample LaTeX input file.  (Version of 12 August 2004.)
%
% A '%' character causes TeX to ignore all remaining text on the line,
% and is used for comments like this one.

\documentclass{article}      % Specifies the document class

                             % The preamble begins here.
\title{An Example Document}  % Declares the document's title.
\author{Leslie Lamport}      % Declares the author's name.
\date{January 21, 1994}      % Deleting this command produces today's date.

\newcommand{\ip}[2]{(#1, #2)}
                             % Defines \ip{arg1}{arg2} to mean
                             % (arg1, arg2).

%\newcommand{\ip}[2]{\langle #1 | #2\rangle}
                             % This is an alternative definition of
                             % \ip that is commented out.

\begin{document}             % End of preamble and beginning of text.

\maketitle                   % Produces the title.

This is an example input file.  Comparing it with
the output it generates can show you how to
produce a simple document of your own.

\section{Ordinary Text}      % Produces section heading.  Lower-level
                             % sections are begun with similar
                             % \subsection and \subsubsection commands.

The ends  of words and sentences are marked
  by   spaces. It  doesn't matter how many
spaces    you type; one is as good as 100.  The
end of   a line counts as a space.

One   or more   blank lines denote the  end
of  a paragraph.

Since any number of consecutive spaces are treated
like a single one, the formatting of the input
file makes no difference to
      \LaTeX,                % The \LaTeX command generates the LaTeX logo.
but it makes a difference to you.  When you use
\LaTeX, making your input file as easy to read
as possible will be a great help as you write
your document and when you change it.  This sample
file shows how you can add comments to your own input
file.

Because printing is different from typewriting,
there are a number of things that you have to do
differently when preparing an input file than if
you were just typing the document directly.
Quotation marks like
       ``this''
have to be handled specially, as do quotes within
quotes:
       ``\,`this'            % \, separates the double and single quote.
        is what I just
        wrote, not  `that'\,''.

Dashes come in three sizes: an
       intra-word
dash, a medium dash for number ranges like
       1--2,
and a punctuation
       dash---like
this.

A sentence-ending space should be larger than the
space between words within a sentence.  You
sometimes have to type special commands in
conjunction with punctuation characters to get
this right, as in the following sentence.
       Gnats, gnus, etc.\ all  % `\ ' makes an inter-word space.
       begin with G\@.         % \@ marks end-of-sentence punctuation.
You should check the spaces after periods when
reading your output to make sure you haven't
forgotten any special cases.  Generating an
ellipsis
       \ldots\               % `\ ' is needed after `\ldots' because TeX
                             % ignores spaces after command names like \ldots
                             % made from \ + letters.
                             %
                             % Note how a `%' character causes TeX to ignore
                             % the end of the input line, so these blank lines
                             % do not start a new paragraph.
                             %
with the right spacing around the periods requires
a special command.

\LaTeX\ interprets some common characters as
commands, so you must type special commands to
generate them.  These characters include the
following:
       \$ \& \% \# \{ and \}.

In printing, text is usually emphasized with an
       \emph{italic}
type style.

\begin{em}
   A long segment of text can also be emphasized
   in this way.  Text within such a segment can be
   given \emph{additional} emphasis.
\end{em}

It is sometimes necessary to prevent \LaTeX\ from
breaking a line where it might otherwise do so.
This may be at a space, as between the ``Mr.''\ and
``Jones'' in
       ``Mr.~Jones'',        % ~ produces an unbreakable interword space.
or within a word---especially when the word is a
symbol like
       \mbox{\emph{itemnum}}
that makes little sense when hyphenated across
lines.

Footnotes\footnote{This is an example of a footnote.}
pose no problem.

\LaTeX\ is good at typesetting mathematical formulas
like
       \( x-3y + z = 7 \)
or
       \( a_{1} > x^{2n} + y^{2n} > x' \)
or
       \( \ip{A}{B} = \sum_{i} a_{i} b_{i} \).
The spaces you type in a formula are
ignored.  Remember that a letter like
       $x$                   % $ ... $  and  \( ... \)  are equivalent
is a formula when it denotes a mathematical
symbol, and it should be typed as one.

\section{Displayed Text}

Text is displayed by indenting it from the left
margin.  Quotations are commonly displayed.  There
are short quotations
\begin{quote}
   This is a short quotation.  It consists of a
   single paragraph of text.  See how it is formatted.
\end{quote}
and longer ones.
\begin{quotation}
   This is a longer quotation.  It consists of two
   paragraphs of text, neither of which are
   particularly interesting.

   This is the second paragraph of the quotation.  It
   is just as dull as the first paragraph.
\end{quotation}
Another frequently-displayed structure is a list.
The following is an example of an \emph{itemized}
list.
\begin{itemize}
   \item This is the first item of an itemized list.
         Each item in the list is marked with a ``tick''.
         You don't have to worry about what kind of tick
         mark is used.

   \item This is the second item of the list.  It
         contains another list nested inside it.  The inner
         list is an \emph{enumerated} list.
         \begin{enumerate}
            \item This is the first item of an enumerated
                  list that is nested within the itemized list.

            \item This is the second item of the inner list.
                  \LaTeX\ allows you to nest lists deeper than
                  you really should.
         \end{enumerate}
         This is the rest of the second item of the outer
         list.  It is no more interesting than any other
         part of the item.
   \item This is the third item of the list.
\end{itemize}
You can even display poetry.
\begin{verse}
   There is an environment
    for verse \\             % The \\ command separates lines
   Whose features some poets % within a stanza.
   will curse.

                             % One or more blank lines separate stanzas.

   For instead of making\\
   Them do \emph{all} line breaking, \\
   It allows them to put too many words on a line when they'd rather be
   forced to be terse.
\end{verse}

Mathematical formulas may also be displayed.  A
displayed formula
is
one-line long; multiline
formulas require special formatting instructions.
   \[  \ip{\Gamma}{\psi'} = x'' + y^{2} + z_{i}^{n}\]
Don't start a paragraph with a displayed equation,
nor make one a paragraph by itself.

\end{document}               % End of document.

\section{Bench Huge Section 1167}
% synthetic bench section 1167
% This is a sample LaTeX input file.  (Version of 12 August 2004.)
%
% A '%' character causes TeX to ignore all remaining text on the line,
% and is used for comments like this one.

\documentclass{article}      % Specifies the document class

                             % The preamble begins here.
\title{An Example Document}  % Declares the document's title.
\author{Leslie Lamport}      % Declares the author's name.
\date{January 21, 1994}      % Deleting this command produces today's date.

\newcommand{\ip}[2]{(#1, #2)}
                             % Defines \ip{arg1}{arg2} to mean
                             % (arg1, arg2).

%\newcommand{\ip}[2]{\langle #1 | #2\rangle}
                             % This is an alternative definition of
                             % \ip that is commented out.

\begin{document}             % End of preamble and beginning of text.

\maketitle                   % Produces the title.

This is an example input file.  Comparing it with
the output it generates can show you how to
produce a simple document of your own.

\section{Ordinary Text}      % Produces section heading.  Lower-level
                             % sections are begun with similar
                             % \subsection and \subsubsection commands.

The ends  of words and sentences are marked
  by   spaces. It  doesn't matter how many
spaces    you type; one is as good as 100.  The
end of   a line counts as a space.

One   or more   blank lines denote the  end
of  a paragraph.

Since any number of consecutive spaces are treated
like a single one, the formatting of the input
file makes no difference to
      \LaTeX,                % The \LaTeX command generates the LaTeX logo.
but it makes a difference to you.  When you use
\LaTeX, making your input file as easy to read
as possible will be a great help as you write
your document and when you change it.  This sample
file shows how you can add comments to your own input
file.

Because printing is different from typewriting,
there are a number of things that you have to do
differently when preparing an input file than if
you were just typing the document directly.
Quotation marks like
       ``this''
have to be handled specially, as do quotes within
quotes:
       ``\,`this'            % \, separates the double and single quote.
        is what I just
        wrote, not  `that'\,''.

Dashes come in three sizes: an
       intra-word
dash, a medium dash for number ranges like
       1--2,
and a punctuation
       dash---like
this.

A sentence-ending space should be larger than the
space between words within a sentence.  You
sometimes have to type special commands in
conjunction with punctuation characters to get
this right, as in the following sentence.
       Gnats, gnus, etc.\ all  % `\ ' makes an inter-word space.
       begin with G\@.         % \@ marks end-of-sentence punctuation.
You should check the spaces after periods when
reading your output to make sure you haven't
forgotten any special cases.  Generating an
ellipsis
       \ldots\               % `\ ' is needed after `\ldots' because TeX
                             % ignores spaces after command names like \ldots
                             % made from \ + letters.
                             %
                             % Note how a `%' character causes TeX to ignore
                             % the end of the input line, so these blank lines
                             % do not start a new paragraph.
                             %
with the right spacing around the periods requires
a special command.

\LaTeX\ interprets some common characters as
commands, so you must type special commands to
generate them.  These characters include the
following:
       \$ \& \% \# \{ and \}.

In printing, text is usually emphasized with an
       \emph{italic}
type style.

\begin{em}
   A long segment of text can also be emphasized
   in this way.  Text within such a segment can be
   given \emph{additional} emphasis.
\end{em}

It is sometimes necessary to prevent \LaTeX\ from
breaking a line where it might otherwise do so.
This may be at a space, as between the ``Mr.''\ and
``Jones'' in
       ``Mr.~Jones'',        % ~ produces an unbreakable interword space.
or within a word---especially when the word is a
symbol like
       \mbox{\emph{itemnum}}
that makes little sense when hyphenated across
lines.

Footnotes\footnote{This is an example of a footnote.}
pose no problem.

\LaTeX\ is good at typesetting mathematical formulas
like
       \( x-3y + z = 7 \)
or
       \( a_{1} > x^{2n} + y^{2n} > x' \)
or
       \( \ip{A}{B} = \sum_{i} a_{i} b_{i} \).
The spaces you type in a formula are
ignored.  Remember that a letter like
       $x$                   % $ ... $  and  \( ... \)  are equivalent
is a formula when it denotes a mathematical
symbol, and it should be typed as one.

\section{Displayed Text}

Text is displayed by indenting it from the left
margin.  Quotations are commonly displayed.  There
are short quotations
\begin{quote}
   This is a short quotation.  It consists of a
   single paragraph of text.  See how it is formatted.
\end{quote}
and longer ones.
\begin{quotation}
   This is a longer quotation.  It consists of two
   paragraphs of text, neither of which are
   particularly interesting.

   This is the second paragraph of the quotation.  It
   is just as dull as the first paragraph.
\end{quotation}
Another frequently-displayed structure is a list.
The following is an example of an \emph{itemized}
list.
\begin{itemize}
   \item This is the first item of an itemized list.
         Each item in the list is marked with a ``tick''.
         You don't have to worry about what kind of tick
         mark is used.

   \item This is the second item of the list.  It
         contains another list nested inside it.  The inner
         list is an \emph{enumerated} list.
         \begin{enumerate}
            \item This is the first item of an enumerated
                  list that is nested within the itemized list.

            \item This is the second item of the inner list.
                  \LaTeX\ allows you to nest lists deeper than
                  you really should.
         \end{enumerate}
         This is the rest of the second item of the outer
         list.  It is no more interesting than any other
         part of the item.
   \item This is the third item of the list.
\end{itemize}
You can even display poetry.
\begin{verse}
   There is an environment
    for verse \\             % The \\ command separates lines
   Whose features some poets % within a stanza.
   will curse.

                             % One or more blank lines separate stanzas.

   For instead of making\\
   Them do \emph{all} line breaking, \\
   It allows them to put too many words on a line when they'd rather be
   forced to be terse.
\end{verse}

Mathematical formulas may also be displayed.  A
displayed formula
is
one-line long; multiline
formulas require special formatting instructions.
   \[  \ip{\Gamma}{\psi'} = x'' + y^{2} + z_{i}^{n}\]
Don't start a paragraph with a displayed equation,
nor make one a paragraph by itself.

\end{document}               % End of document.

\section{Bench Huge Section 1168}
% synthetic bench section 1168
% This is a sample LaTeX input file.  (Version of 12 August 2004.)
%
% A '%' character causes TeX to ignore all remaining text on the line,
% and is used for comments like this one.

\documentclass{article}      % Specifies the document class

                             % The preamble begins here.
\title{An Example Document}  % Declares the document's title.
\author{Leslie Lamport}      % Declares the author's name.
\date{January 21, 1994}      % Deleting this command produces today's date.

\newcommand{\ip}[2]{(#1, #2)}
                             % Defines \ip{arg1}{arg2} to mean
                             % (arg1, arg2).

%\newcommand{\ip}[2]{\langle #1 | #2\rangle}
                             % This is an alternative definition of
                             % \ip that is commented out.

\begin{document}             % End of preamble and beginning of text.

\maketitle                   % Produces the title.

This is an example input file.  Comparing it with
the output it generates can show you how to
produce a simple document of your own.

\section{Ordinary Text}      % Produces section heading.  Lower-level
                             % sections are begun with similar
                             % \subsection and \subsubsection commands.

The ends  of words and sentences are marked
  by   spaces. It  doesn't matter how many
spaces    you type; one is as good as 100.  The
end of   a line counts as a space.

One   or more   blank lines denote the  end
of  a paragraph.

Since any number of consecutive spaces are treated
like a single one, the formatting of the input
file makes no difference to
      \LaTeX,                % The \LaTeX command generates the LaTeX logo.
but it makes a difference to you.  When you use
\LaTeX, making your input file as easy to read
as possible will be a great help as you write
your document and when you change it.  This sample
file shows how you can add comments to your own input
file.

Because printing is different from typewriting,
there are a number of things that you have to do
differently when preparing an input file than if
you were just typing the document directly.
Quotation marks like
       ``this''
have to be handled specially, as do quotes within
quotes:
       ``\,`this'            % \, separates the double and single quote.
        is what I just
        wrote, not  `that'\,''.

Dashes come in three sizes: an
       intra-word
dash, a medium dash for number ranges like
       1--2,
and a punctuation
       dash---like
this.

A sentence-ending space should be larger than the
space between words within a sentence.  You
sometimes have to type special commands in
conjunction with punctuation characters to get
this right, as in the following sentence.
       Gnats, gnus, etc.\ all  % `\ ' makes an inter-word space.
       begin with G\@.         % \@ marks end-of-sentence punctuation.
You should check the spaces after periods when
reading your output to make sure you haven't
forgotten any special cases.  Generating an
ellipsis
       \ldots\               % `\ ' is needed after `\ldots' because TeX
                             % ignores spaces after command names like \ldots
                             % made from \ + letters.
                             %
                             % Note how a `%' character causes TeX to ignore
                             % the end of the input line, so these blank lines
                             % do not start a new paragraph.
                             %
with the right spacing around the periods requires
a special command.

\LaTeX\ interprets some common characters as
commands, so you must type special commands to
generate them.  These characters include the
following:
       \$ \& \% \# \{ and \}.

In printing, text is usually emphasized with an
       \emph{italic}
type style.

\begin{em}
   A long segment of text can also be emphasized
   in this way.  Text within such a segment can be
   given \emph{additional} emphasis.
\end{em}

It is sometimes necessary to prevent \LaTeX\ from
breaking a line where it might otherwise do so.
This may be at a space, as between the ``Mr.''\ and
``Jones'' in
       ``Mr.~Jones'',        % ~ produces an unbreakable interword space.
or within a word---especially when the word is a
symbol like
       \mbox{\emph{itemnum}}
that makes little sense when hyphenated across
lines.

Footnotes\footnote{This is an example of a footnote.}
pose no problem.

\LaTeX\ is good at typesetting mathematical formulas
like
       \( x-3y + z = 7 \)
or
       \( a_{1} > x^{2n} + y^{2n} > x' \)
or
       \( \ip{A}{B} = \sum_{i} a_{i} b_{i} \).
The spaces you type in a formula are
ignored.  Remember that a letter like
       $x$                   % $ ... $  and  \( ... \)  are equivalent
is a formula when it denotes a mathematical
symbol, and it should be typed as one.

\section{Displayed Text}

Text is displayed by indenting it from the left
margin.  Quotations are commonly displayed.  There
are short quotations
\begin{quote}
   This is a short quotation.  It consists of a
   single paragraph of text.  See how it is formatted.
\end{quote}
and longer ones.
\begin{quotation}
   This is a longer quotation.  It consists of two
   paragraphs of text, neither of which are
   particularly interesting.

   This is the second paragraph of the quotation.  It
   is just as dull as the first paragraph.
\end{quotation}
Another frequently-displayed structure is a list.
The following is an example of an \emph{itemized}
list.
\begin{itemize}
   \item This is the first item of an itemized list.
         Each item in the list is marked with a ``tick''.
         You don't have to worry about what kind of tick
         mark is used.

   \item This is the second item of the list.  It
         contains another list nested inside it.  The inner
         list is an \emph{enumerated} list.
         \begin{enumerate}
            \item This is the first item of an enumerated
                  list that is nested within the itemized list.

            \item This is the second item of the inner list.
                  \LaTeX\ allows you to nest lists deeper than
                  you really should.
         \end{enumerate}
         This is the rest of the second item of the outer
         list.  It is no more interesting than any other
         part of the item.
   \item This is the third item of the list.
\end{itemize}
You can even display poetry.
\begin{verse}
   There is an environment
    for verse \\             % The \\ command separates lines
   Whose features some poets % within a stanza.
   will curse.

                             % One or more blank lines separate stanzas.

   For instead of making\\
   Them do \emph{all} line breaking, \\
   It allows them to put too many words on a line when they'd rather be
   forced to be terse.
\end{verse}

Mathematical formulas may also be displayed.  A
displayed formula
is
one-line long; multiline
formulas require special formatting instructions.
   \[  \ip{\Gamma}{\psi'} = x'' + y^{2} + z_{i}^{n}\]
Don't start a paragraph with a displayed equation,
nor make one a paragraph by itself.

\end{document}               % End of document.

\section{Bench Huge Section 1169}
% synthetic bench section 1169
% This is a sample LaTeX input file.  (Version of 12 August 2004.)
%
% A '%' character causes TeX to ignore all remaining text on the line,
% and is used for comments like this one.

\documentclass{article}      % Specifies the document class

                             % The preamble begins here.
\title{An Example Document}  % Declares the document's title.
\author{Leslie Lamport}      % Declares the author's name.
\date{January 21, 1994}      % Deleting this command produces today's date.

\newcommand{\ip}[2]{(#1, #2)}
                             % Defines \ip{arg1}{arg2} to mean
                             % (arg1, arg2).

%\newcommand{\ip}[2]{\langle #1 | #2\rangle}
                             % This is an alternative definition of
                             % \ip that is commented out.

\begin{document}             % End of preamble and beginning of text.

\maketitle                   % Produces the title.

This is an example input file.  Comparing it with
the output it generates can show you how to
produce a simple document of your own.

\section{Ordinary Text}      % Produces section heading.  Lower-level
                             % sections are begun with similar
                             % \subsection and \subsubsection commands.

The ends  of words and sentences are marked
  by   spaces. It  doesn't matter how many
spaces    you type; one is as good as 100.  The
end of   a line counts as a space.

One   or more   blank lines denote the  end
of  a paragraph.

Since any number of consecutive spaces are treated
like a single one, the formatting of the input
file makes no difference to
      \LaTeX,                % The \LaTeX command generates the LaTeX logo.
but it makes a difference to you.  When you use
\LaTeX, making your input file as easy to read
as possible will be a great help as you write
your document and when you change it.  This sample
file shows how you can add comments to your own input
file.

Because printing is different from typewriting,
there are a number of things that you have to do
differently when preparing an input file than if
you were just typing the document directly.
Quotation marks like
       ``this''
have to be handled specially, as do quotes within
quotes:
       ``\,`this'            % \, separates the double and single quote.
        is what I just
        wrote, not  `that'\,''.

Dashes come in three sizes: an
       intra-word
dash, a medium dash for number ranges like
       1--2,
and a punctuation
       dash---like
this.

A sentence-ending space should be larger than the
space between words within a sentence.  You
sometimes have to type special commands in
conjunction with punctuation characters to get
this right, as in the following sentence.
       Gnats, gnus, etc.\ all  % `\ ' makes an inter-word space.
       begin with G\@.         % \@ marks end-of-sentence punctuation.
You should check the spaces after periods when
reading your output to make sure you haven't
forgotten any special cases.  Generating an
ellipsis
       \ldots\               % `\ ' is needed after `\ldots' because TeX
                             % ignores spaces after command names like \ldots
                             % made from \ + letters.
                             %
                             % Note how a `%' character causes TeX to ignore
                             % the end of the input line, so these blank lines
                             % do not start a new paragraph.
                             %
with the right spacing around the periods requires
a special command.

\LaTeX\ interprets some common characters as
commands, so you must type special commands to
generate them.  These characters include the
following:
       \$ \& \% \# \{ and \}.

In printing, text is usually emphasized with an
       \emph{italic}
type style.

\begin{em}
   A long segment of text can also be emphasized
   in this way.  Text within such a segment can be
   given \emph{additional} emphasis.
\end{em}

It is sometimes necessary to prevent \LaTeX\ from
breaking a line where it might otherwise do so.
This may be at a space, as between the ``Mr.''\ and
``Jones'' in
       ``Mr.~Jones'',        % ~ produces an unbreakable interword space.
or within a word---especially when the word is a
symbol like
       \mbox{\emph{itemnum}}
that makes little sense when hyphenated across
lines.

Footnotes\footnote{This is an example of a footnote.}
pose no problem.

\LaTeX\ is good at typesetting mathematical formulas
like
       \( x-3y + z = 7 \)
or
       \( a_{1} > x^{2n} + y^{2n} > x' \)
or
       \( \ip{A}{B} = \sum_{i} a_{i} b_{i} \).
The spaces you type in a formula are
ignored.  Remember that a letter like
       $x$                   % $ ... $  and  \( ... \)  are equivalent
is a formula when it denotes a mathematical
symbol, and it should be typed as one.

\section{Displayed Text}

Text is displayed by indenting it from the left
margin.  Quotations are commonly displayed.  There
are short quotations
\begin{quote}
   This is a short quotation.  It consists of a
   single paragraph of text.  See how it is formatted.
\end{quote}
and longer ones.
\begin{quotation}
   This is a longer quotation.  It consists of two
   paragraphs of text, neither of which are
   particularly interesting.

   This is the second paragraph of the quotation.  It
   is just as dull as the first paragraph.
\end{quotation}
Another frequently-displayed structure is a list.
The following is an example of an \emph{itemized}
list.
\begin{itemize}
   \item This is the first item of an itemized list.
         Each item in the list is marked with a ``tick''.
         You don't have to worry about what kind of tick
         mark is used.

   \item This is the second item of the list.  It
         contains another list nested inside it.  The inner
         list is an \emph{enumerated} list.
         \begin{enumerate}
            \item This is the first item of an enumerated
                  list that is nested within the itemized list.

            \item This is the second item of the inner list.
                  \LaTeX\ allows you to nest lists deeper than
                  you really should.
         \end{enumerate}
         This is the rest of the second item of the outer
         list.  It is no more interesting than any other
         part of the item.
   \item This is the third item of the list.
\end{itemize}
You can even display poetry.
\begin{verse}
   There is an environment
    for verse \\             % The \\ command separates lines
   Whose features some poets % within a stanza.
   will curse.

                             % One or more blank lines separate stanzas.

   For instead of making\\
   Them do \emph{all} line breaking, \\
   It allows them to put too many words on a line when they'd rather be
   forced to be terse.
\end{verse}

Mathematical formulas may also be displayed.  A
displayed formula
is
one-line long; multiline
formulas require special formatting instructions.
   \[  \ip{\Gamma}{\psi'} = x'' + y^{2} + z_{i}^{n}\]
Don't start a paragraph with a displayed equation,
nor make one a paragraph by itself.

\end{document}               % End of document.

\section{Bench Huge Section 1170}
% synthetic bench section 1170
% This is a sample LaTeX input file.  (Version of 12 August 2004.)
%
% A '%' character causes TeX to ignore all remaining text on the line,
% and is used for comments like this one.

\documentclass{article}      % Specifies the document class

                             % The preamble begins here.
\title{An Example Document}  % Declares the document's title.
\author{Leslie Lamport}      % Declares the author's name.
\date{January 21, 1994}      % Deleting this command produces today's date.

\newcommand{\ip}[2]{(#1, #2)}
                             % Defines \ip{arg1}{arg2} to mean
                             % (arg1, arg2).

%\newcommand{\ip}[2]{\langle #1 | #2\rangle}
                             % This is an alternative definition of
                             % \ip that is commented out.

\begin{document}             % End of preamble and beginning of text.

\maketitle                   % Produces the title.

This is an example input file.  Comparing it with
the output it generates can show you how to
produce a simple document of your own.

\section{Ordinary Text}      % Produces section heading.  Lower-level
                             % sections are begun with similar
                             % \subsection and \subsubsection commands.

The ends  of words and sentences are marked
  by   spaces. It  doesn't matter how many
spaces    you type; one is as good as 100.  The
end of   a line counts as a space.

One   or more   blank lines denote the  end
of  a paragraph.

Since any number of consecutive spaces are treated
like a single one, the formatting of the input
file makes no difference to
      \LaTeX,                % The \LaTeX command generates the LaTeX logo.
but it makes a difference to you.  When you use
\LaTeX, making your input file as easy to read
as possible will be a great help as you write
your document and when you change it.  This sample
file shows how you can add comments to your own input
file.

Because printing is different from typewriting,
there are a number of things that you have to do
differently when preparing an input file than if
you were just typing the document directly.
Quotation marks like
       ``this''
have to be handled specially, as do quotes within
quotes:
       ``\,`this'            % \, separates the double and single quote.
        is what I just
        wrote, not  `that'\,''.

Dashes come in three sizes: an
       intra-word
dash, a medium dash for number ranges like
       1--2,
and a punctuation
       dash---like
this.

A sentence-ending space should be larger than the
space between words within a sentence.  You
sometimes have to type special commands in
conjunction with punctuation characters to get
this right, as in the following sentence.
       Gnats, gnus, etc.\ all  % `\ ' makes an inter-word space.
       begin with G\@.         % \@ marks end-of-sentence punctuation.
You should check the spaces after periods when
reading your output to make sure you haven't
forgotten any special cases.  Generating an
ellipsis
       \ldots\               % `\ ' is needed after `\ldots' because TeX
                             % ignores spaces after command names like \ldots
                             % made from \ + letters.
                             %
                             % Note how a `%' character causes TeX to ignore
                             % the end of the input line, so these blank lines
                             % do not start a new paragraph.
                             %
with the right spacing around the periods requires
a special command.

\LaTeX\ interprets some common characters as
commands, so you must type special commands to
generate them.  These characters include the
following:
       \$ \& \% \# \{ and \}.

In printing, text is usually emphasized with an
       \emph{italic}
type style.

\begin{em}
   A long segment of text can also be emphasized
   in this way.  Text within such a segment can be
   given \emph{additional} emphasis.
\end{em}

It is sometimes necessary to prevent \LaTeX\ from
breaking a line where it might otherwise do so.
This may be at a space, as between the ``Mr.''\ and
``Jones'' in
       ``Mr.~Jones'',        % ~ produces an unbreakable interword space.
or within a word---especially when the word is a
symbol like
       \mbox{\emph{itemnum}}
that makes little sense when hyphenated across
lines.

Footnotes\footnote{This is an example of a footnote.}
pose no problem.

\LaTeX\ is good at typesetting mathematical formulas
like
       \( x-3y + z = 7 \)
or
       \( a_{1} > x^{2n} + y^{2n} > x' \)
or
       \( \ip{A}{B} = \sum_{i} a_{i} b_{i} \).
The spaces you type in a formula are
ignored.  Remember that a letter like
       $x$                   % $ ... $  and  \( ... \)  are equivalent
is a formula when it denotes a mathematical
symbol, and it should be typed as one.

\section{Displayed Text}

Text is displayed by indenting it from the left
margin.  Quotations are commonly displayed.  There
are short quotations
\begin{quote}
   This is a short quotation.  It consists of a
   single paragraph of text.  See how it is formatted.
\end{quote}
and longer ones.
\begin{quotation}
   This is a longer quotation.  It consists of two
   paragraphs of text, neither of which are
   particularly interesting.

   This is the second paragraph of the quotation.  It
   is just as dull as the first paragraph.
\end{quotation}
Another frequently-displayed structure is a list.
The following is an example of an \emph{itemized}
list.
\begin{itemize}
   \item This is the first item of an itemized list.
         Each item in the list is marked with a ``tick''.
         You don't have to worry about what kind of tick
         mark is used.

   \item This is the second item of the list.  It
         contains another list nested inside it.  The inner
         list is an \emph{enumerated} list.
         \begin{enumerate}
            \item This is the first item of an enumerated
                  list that is nested within the itemized list.

            \item This is the second item of the inner list.
                  \LaTeX\ allows you to nest lists deeper than
                  you really should.
         \end{enumerate}
         This is the rest of the second item of the outer
         list.  It is no more interesting than any other
         part of the item.
   \item This is the third item of the list.
\end{itemize}
You can even display poetry.
\begin{verse}
   There is an environment
    for verse \\             % The \\ command separates lines
   Whose features some poets % within a stanza.
   will curse.

                             % One or more blank lines separate stanzas.

   For instead of making\\
   Them do \emph{all} line breaking, \\
   It allows them to put too many words on a line when they'd rather be
   forced to be terse.
\end{verse}

Mathematical formulas may also be displayed.  A
displayed formula
is
one-line long; multiline
formulas require special formatting instructions.
   \[  \ip{\Gamma}{\psi'} = x'' + y^{2} + z_{i}^{n}\]
Don't start a paragraph with a displayed equation,
nor make one a paragraph by itself.

\end{document}               % End of document.

\section{Bench Huge Section 1171}
% synthetic bench section 1171
% This is a sample LaTeX input file.  (Version of 12 August 2004.)
%
% A '%' character causes TeX to ignore all remaining text on the line,
% and is used for comments like this one.

\documentclass{article}      % Specifies the document class

                             % The preamble begins here.
\title{An Example Document}  % Declares the document's title.
\author{Leslie Lamport}      % Declares the author's name.
\date{January 21, 1994}      % Deleting this command produces today's date.

\newcommand{\ip}[2]{(#1, #2)}
                             % Defines \ip{arg1}{arg2} to mean
                             % (arg1, arg2).

%\newcommand{\ip}[2]{\langle #1 | #2\rangle}
                             % This is an alternative definition of
                             % \ip that is commented out.

\begin{document}             % End of preamble and beginning of text.

\maketitle                   % Produces the title.

This is an example input file.  Comparing it with
the output it generates can show you how to
produce a simple document of your own.

\section{Ordinary Text}      % Produces section heading.  Lower-level
                             % sections are begun with similar
                             % \subsection and \subsubsection commands.

The ends  of words and sentences are marked
  by   spaces. It  doesn't matter how many
spaces    you type; one is as good as 100.  The
end of   a line counts as a space.

One   or more   blank lines denote the  end
of  a paragraph.

Since any number of consecutive spaces are treated
like a single one, the formatting of the input
file makes no difference to
      \LaTeX,                % The \LaTeX command generates the LaTeX logo.
but it makes a difference to you.  When you use
\LaTeX, making your input file as easy to read
as possible will be a great help as you write
your document and when you change it.  This sample
file shows how you can add comments to your own input
file.

Because printing is different from typewriting,
there are a number of things that you have to do
differently when preparing an input file than if
you were just typing the document directly.
Quotation marks like
       ``this''
have to be handled specially, as do quotes within
quotes:
       ``\,`this'            % \, separates the double and single quote.
        is what I just
        wrote, not  `that'\,''.

Dashes come in three sizes: an
       intra-word
dash, a medium dash for number ranges like
       1--2,
and a punctuation
       dash---like
this.

A sentence-ending space should be larger than the
space between words within a sentence.  You
sometimes have to type special commands in
conjunction with punctuation characters to get
this right, as in the following sentence.
       Gnats, gnus, etc.\ all  % `\ ' makes an inter-word space.
       begin with G\@.         % \@ marks end-of-sentence punctuation.
You should check the spaces after periods when
reading your output to make sure you haven't
forgotten any special cases.  Generating an
ellipsis
       \ldots\               % `\ ' is needed after `\ldots' because TeX
                             % ignores spaces after command names like \ldots
                             % made from \ + letters.
                             %
                             % Note how a `%' character causes TeX to ignore
                             % the end of the input line, so these blank lines
                             % do not start a new paragraph.
                             %
with the right spacing around the periods requires
a special command.

\LaTeX\ interprets some common characters as
commands, so you must type special commands to
generate them.  These characters include the
following:
       \$ \& \% \# \{ and \}.

In printing, text is usually emphasized with an
       \emph{italic}
type style.

\begin{em}
   A long segment of text can also be emphasized
   in this way.  Text within such a segment can be
   given \emph{additional} emphasis.
\end{em}

It is sometimes necessary to prevent \LaTeX\ from
breaking a line where it might otherwise do so.
This may be at a space, as between the ``Mr.''\ and
``Jones'' in
       ``Mr.~Jones'',        % ~ produces an unbreakable interword space.
or within a word---especially when the word is a
symbol like
       \mbox{\emph{itemnum}}
that makes little sense when hyphenated across
lines.

Footnotes\footnote{This is an example of a footnote.}
pose no problem.

\LaTeX\ is good at typesetting mathematical formulas
like
       \( x-3y + z = 7 \)
or
       \( a_{1} > x^{2n} + y^{2n} > x' \)
or
       \( \ip{A}{B} = \sum_{i} a_{i} b_{i} \).
The spaces you type in a formula are
ignored.  Remember that a letter like
       $x$                   % $ ... $  and  \( ... \)  are equivalent
is a formula when it denotes a mathematical
symbol, and it should be typed as one.

\section{Displayed Text}

Text is displayed by indenting it from the left
margin.  Quotations are commonly displayed.  There
are short quotations
\begin{quote}
   This is a short quotation.  It consists of a
   single paragraph of text.  See how it is formatted.
\end{quote}
and longer ones.
\begin{quotation}
   This is a longer quotation.  It consists of two
   paragraphs of text, neither of which are
   particularly interesting.

   This is the second paragraph of the quotation.  It
   is just as dull as the first paragraph.
\end{quotation}
Another frequently-displayed structure is a list.
The following is an example of an \emph{itemized}
list.
\begin{itemize}
   \item This is the first item of an itemized list.
         Each item in the list is marked with a ``tick''.
         You don't have to worry about what kind of tick
         mark is used.

   \item This is the second item of the list.  It
         contains another list nested inside it.  The inner
         list is an \emph{enumerated} list.
         \begin{enumerate}
            \item This is the first item of an enumerated
                  list that is nested within the itemized list.

            \item This is the second item of the inner list.
                  \LaTeX\ allows you to nest lists deeper than
                  you really should.
         \end{enumerate}
         This is the rest of the second item of the outer
         list.  It is no more interesting than any other
         part of the item.
   \item This is the third item of the list.
\end{itemize}
You can even display poetry.
\begin{verse}
   There is an environment
    for verse \\             % The \\ command separates lines
   Whose features some poets % within a stanza.
   will curse.

                             % One or more blank lines separate stanzas.

   For instead of making\\
   Them do \emph{all} line breaking, \\
   It allows them to put too many words on a line when they'd rather be
   forced to be terse.
\end{verse}

Mathematical formulas may also be displayed.  A
displayed formula
is
one-line long; multiline
formulas require special formatting instructions.
   \[  \ip{\Gamma}{\psi'} = x'' + y^{2} + z_{i}^{n}\]
Don't start a paragraph with a displayed equation,
nor make one a paragraph by itself.

\end{document}               % End of document.

\section{Bench Huge Section 1172}
% synthetic bench section 1172
% This is a sample LaTeX input file.  (Version of 12 August 2004.)
%
% A '%' character causes TeX to ignore all remaining text on the line,
% and is used for comments like this one.

\documentclass{article}      % Specifies the document class

                             % The preamble begins here.
\title{An Example Document}  % Declares the document's title.
\author{Leslie Lamport}      % Declares the author's name.
\date{January 21, 1994}      % Deleting this command produces today's date.

\newcommand{\ip}[2]{(#1, #2)}
                             % Defines \ip{arg1}{arg2} to mean
                             % (arg1, arg2).

%\newcommand{\ip}[2]{\langle #1 | #2\rangle}
                             % This is an alternative definition of
                             % \ip that is commented out.

\begin{document}             % End of preamble and beginning of text.

\maketitle                   % Produces the title.

This is an example input file.  Comparing it with
the output it generates can show you how to
produce a simple document of your own.

\section{Ordinary Text}      % Produces section heading.  Lower-level
                             % sections are begun with similar
                             % \subsection and \subsubsection commands.

The ends  of words and sentences are marked
  by   spaces. It  doesn't matter how many
spaces    you type; one is as good as 100.  The
end of   a line counts as a space.

One   or more   blank lines denote the  end
of  a paragraph.

Since any number of consecutive spaces are treated
like a single one, the formatting of the input
file makes no difference to
      \LaTeX,                % The \LaTeX command generates the LaTeX logo.
but it makes a difference to you.  When you use
\LaTeX, making your input file as easy to read
as possible will be a great help as you write
your document and when you change it.  This sample
file shows how you can add comments to your own input
file.

Because printing is different from typewriting,
there are a number of things that you have to do
differently when preparing an input file than if
you were just typing the document directly.
Quotation marks like
       ``this''
have to be handled specially, as do quotes within
quotes:
       ``\,`this'            % \, separates the double and single quote.
        is what I just
        wrote, not  `that'\,''.

Dashes come in three sizes: an
       intra-word
dash, a medium dash for number ranges like
       1--2,
and a punctuation
       dash---like
this.

A sentence-ending space should be larger than the
space between words within a sentence.  You
sometimes have to type special commands in
conjunction with punctuation characters to get
this right, as in the following sentence.
       Gnats, gnus, etc.\ all  % `\ ' makes an inter-word space.
       begin with G\@.         % \@ marks end-of-sentence punctuation.
You should check the spaces after periods when
reading your output to make sure you haven't
forgotten any special cases.  Generating an
ellipsis
       \ldots\               % `\ ' is needed after `\ldots' because TeX
                             % ignores spaces after command names like \ldots
                             % made from \ + letters.
                             %
                             % Note how a `%' character causes TeX to ignore
                             % the end of the input line, so these blank lines
                             % do not start a new paragraph.
                             %
with the right spacing around the periods requires
a special command.

\LaTeX\ interprets some common characters as
commands, so you must type special commands to
generate them.  These characters include the
following:
       \$ \& \% \# \{ and \}.

In printing, text is usually emphasized with an
       \emph{italic}
type style.

\begin{em}
   A long segment of text can also be emphasized
   in this way.  Text within such a segment can be
   given \emph{additional} emphasis.
\end{em}

It is sometimes necessary to prevent \LaTeX\ from
breaking a line where it might otherwise do so.
This may be at a space, as between the ``Mr.''\ and
``Jones'' in
       ``Mr.~Jones'',        % ~ produces an unbreakable interword space.
or within a word---especially when the word is a
symbol like
       \mbox{\emph{itemnum}}
that makes little sense when hyphenated across
lines.

Footnotes\footnote{This is an example of a footnote.}
pose no problem.

\LaTeX\ is good at typesetting mathematical formulas
like
       \( x-3y + z = 7 \)
or
       \( a_{1} > x^{2n} + y^{2n} > x' \)
or
       \( \ip{A}{B} = \sum_{i} a_{i} b_{i} \).
The spaces you type in a formula are
ignored.  Remember that a letter like
       $x$                   % $ ... $  and  \( ... \)  are equivalent
is a formula when it denotes a mathematical
symbol, and it should be typed as one.

\section{Displayed Text}

Text is displayed by indenting it from the left
margin.  Quotations are commonly displayed.  There
are short quotations
\begin{quote}
   This is a short quotation.  It consists of a
   single paragraph of text.  See how it is formatted.
\end{quote}
and longer ones.
\begin{quotation}
   This is a longer quotation.  It consists of two
   paragraphs of text, neither of which are
   particularly interesting.

   This is the second paragraph of the quotation.  It
   is just as dull as the first paragraph.
\end{quotation}
Another frequently-displayed structure is a list.
The following is an example of an \emph{itemized}
list.
\begin{itemize}
   \item This is the first item of an itemized list.
         Each item in the list is marked with a ``tick''.
         You don't have to worry about what kind of tick
         mark is used.

   \item This is the second item of the list.  It
         contains another list nested inside it.  The inner
         list is an \emph{enumerated} list.
         \begin{enumerate}
            \item This is the first item of an enumerated
                  list that is nested within the itemized list.

            \item This is the second item of the inner list.
                  \LaTeX\ allows you to nest lists deeper than
                  you really should.
         \end{enumerate}
         This is the rest of the second item of the outer
         list.  It is no more interesting than any other
         part of the item.
   \item This is the third item of the list.
\end{itemize}
You can even display poetry.
\begin{verse}
   There is an environment
    for verse \\             % The \\ command separates lines
   Whose features some poets % within a stanza.
   will curse.

                             % One or more blank lines separate stanzas.

   For instead of making\\
   Them do \emph{all} line breaking, \\
   It allows them to put too many words on a line when they'd rather be
   forced to be terse.
\end{verse}

Mathematical formulas may also be displayed.  A
displayed formula
is
one-line long; multiline
formulas require special formatting instructions.
   \[  \ip{\Gamma}{\psi'} = x'' + y^{2} + z_{i}^{n}\]
Don't start a paragraph with a displayed equation,
nor make one a paragraph by itself.

\end{document}               % End of document.

\section{Bench Huge Section 1173}
% synthetic bench section 1173
% This is a sample LaTeX input file.  (Version of 12 August 2004.)
%
% A '%' character causes TeX to ignore all remaining text on the line,
% and is used for comments like this one.

\documentclass{article}      % Specifies the document class

                             % The preamble begins here.
\title{An Example Document}  % Declares the document's title.
\author{Leslie Lamport}      % Declares the author's name.
\date{January 21, 1994}      % Deleting this command produces today's date.

\newcommand{\ip}[2]{(#1, #2)}
                             % Defines \ip{arg1}{arg2} to mean
                             % (arg1, arg2).

%\newcommand{\ip}[2]{\langle #1 | #2\rangle}
                             % This is an alternative definition of
                             % \ip that is commented out.

\begin{document}             % End of preamble and beginning of text.

\maketitle                   % Produces the title.

This is an example input file.  Comparing it with
the output it generates can show you how to
produce a simple document of your own.

\section{Ordinary Text}      % Produces section heading.  Lower-level
                             % sections are begun with similar
                             % \subsection and \subsubsection commands.

The ends  of words and sentences are marked
  by   spaces. It  doesn't matter how many
spaces    you type; one is as good as 100.  The
end of   a line counts as a space.

One   or more   blank lines denote the  end
of  a paragraph.

Since any number of consecutive spaces are treated
like a single one, the formatting of the input
file makes no difference to
      \LaTeX,                % The \LaTeX command generates the LaTeX logo.
but it makes a difference to you.  When you use
\LaTeX, making your input file as easy to read
as possible will be a great help as you write
your document and when you change it.  This sample
file shows how you can add comments to your own input
file.

Because printing is different from typewriting,
there are a number of things that you have to do
differently when preparing an input file than if
you were just typing the document directly.
Quotation marks like
       ``this''
have to be handled specially, as do quotes within
quotes:
       ``\,`this'            % \, separates the double and single quote.
        is what I just
        wrote, not  `that'\,''.

Dashes come in three sizes: an
       intra-word
dash, a medium dash for number ranges like
       1--2,
and a punctuation
       dash---like
this.

A sentence-ending space should be larger than the
space between words within a sentence.  You
sometimes have to type special commands in
conjunction with punctuation characters to get
this right, as in the following sentence.
       Gnats, gnus, etc.\ all  % `\ ' makes an inter-word space.
       begin with G\@.         % \@ marks end-of-sentence punctuation.
You should check the spaces after periods when
reading your output to make sure you haven't
forgotten any special cases.  Generating an
ellipsis
       \ldots\               % `\ ' is needed after `\ldots' because TeX
                             % ignores spaces after command names like \ldots
                             % made from \ + letters.
                             %
                             % Note how a `%' character causes TeX to ignore
                             % the end of the input line, so these blank lines
                             % do not start a new paragraph.
                             %
with the right spacing around the periods requires
a special command.

\LaTeX\ interprets some common characters as
commands, so you must type special commands to
generate them.  These characters include the
following:
       \$ \& \% \# \{ and \}.

In printing, text is usually emphasized with an
       \emph{italic}
type style.

\begin{em}
   A long segment of text can also be emphasized
   in this way.  Text within such a segment can be
   given \emph{additional} emphasis.
\end{em}

It is sometimes necessary to prevent \LaTeX\ from
breaking a line where it might otherwise do so.
This may be at a space, as between the ``Mr.''\ and
``Jones'' in
       ``Mr.~Jones'',        % ~ produces an unbreakable interword space.
or within a word---especially when the word is a
symbol like
       \mbox{\emph{itemnum}}
that makes little sense when hyphenated across
lines.

Footnotes\footnote{This is an example of a footnote.}
pose no problem.

\LaTeX\ is good at typesetting mathematical formulas
like
       \( x-3y + z = 7 \)
or
       \( a_{1} > x^{2n} + y^{2n} > x' \)
or
       \( \ip{A}{B} = \sum_{i} a_{i} b_{i} \).
The spaces you type in a formula are
ignored.  Remember that a letter like
       $x$                   % $ ... $  and  \( ... \)  are equivalent
is a formula when it denotes a mathematical
symbol, and it should be typed as one.

\section{Displayed Text}

Text is displayed by indenting it from the left
margin.  Quotations are commonly displayed.  There
are short quotations
\begin{quote}
   This is a short quotation.  It consists of a
   single paragraph of text.  See how it is formatted.
\end{quote}
and longer ones.
\begin{quotation}
   This is a longer quotation.  It consists of two
   paragraphs of text, neither of which are
   particularly interesting.

   This is the second paragraph of the quotation.  It
   is just as dull as the first paragraph.
\end{quotation}
Another frequently-displayed structure is a list.
The following is an example of an \emph{itemized}
list.
\begin{itemize}
   \item This is the first item of an itemized list.
         Each item in the list is marked with a ``tick''.
         You don't have to worry about what kind of tick
         mark is used.

   \item This is the second item of the list.  It
         contains another list nested inside it.  The inner
         list is an \emph{enumerated} list.
         \begin{enumerate}
            \item This is the first item of an enumerated
                  list that is nested within the itemized list.

            \item This is the second item of the inner list.
                  \LaTeX\ allows you to nest lists deeper than
                  you really should.
         \end{enumerate}
         This is the rest of the second item of the outer
         list.  It is no more interesting than any other
         part of the item.
   \item This is the third item of the list.
\end{itemize}
You can even display poetry.
\begin{verse}
   There is an environment
    for verse \\             % The \\ command separates lines
   Whose features some poets % within a stanza.
   will curse.

                             % One or more blank lines separate stanzas.

   For instead of making\\
   Them do \emph{all} line breaking, \\
   It allows them to put too many words on a line when they'd rather be
   forced to be terse.
\end{verse}

Mathematical formulas may also be displayed.  A
displayed formula
is
one-line long; multiline
formulas require special formatting instructions.
   \[  \ip{\Gamma}{\psi'} = x'' + y^{2} + z_{i}^{n}\]
Don't start a paragraph with a displayed equation,
nor make one a paragraph by itself.

\end{document}               % End of document.

\section{Bench Huge Section 1174}
% synthetic bench section 1174
% This is a sample LaTeX input file.  (Version of 12 August 2004.)
%
% A '%' character causes TeX to ignore all remaining text on the line,
% and is used for comments like this one.

\documentclass{article}      % Specifies the document class

                             % The preamble begins here.
\title{An Example Document}  % Declares the document's title.
\author{Leslie Lamport}      % Declares the author's name.
\date{January 21, 1994}      % Deleting this command produces today's date.

\newcommand{\ip}[2]{(#1, #2)}
                             % Defines \ip{arg1}{arg2} to mean
                             % (arg1, arg2).

%\newcommand{\ip}[2]{\langle #1 | #2\rangle}
                             % This is an alternative definition of
                             % \ip that is commented out.

\begin{document}             % End of preamble and beginning of text.

\maketitle                   % Produces the title.

This is an example input file.  Comparing it with
the output it generates can show you how to
produce a simple document of your own.

\section{Ordinary Text}      % Produces section heading.  Lower-level
                             % sections are begun with similar
                             % \subsection and \subsubsection commands.

The ends  of words and sentences are marked
  by   spaces. It  doesn't matter how many
spaces    you type; one is as good as 100.  The
end of   a line counts as a space.

One   or more   blank lines denote the  end
of  a paragraph.

Since any number of consecutive spaces are treated
like a single one, the formatting of the input
file makes no difference to
      \LaTeX,                % The \LaTeX command generates the LaTeX logo.
but it makes a difference to you.  When you use
\LaTeX, making your input file as easy to read
as possible will be a great help as you write
your document and when you change it.  This sample
file shows how you can add comments to your own input
file.

Because printing is different from typewriting,
there are a number of things that you have to do
differently when preparing an input file than if
you were just typing the document directly.
Quotation marks like
       ``this''
have to be handled specially, as do quotes within
quotes:
       ``\,`this'            % \, separates the double and single quote.
        is what I just
        wrote, not  `that'\,''.

Dashes come in three sizes: an
       intra-word
dash, a medium dash for number ranges like
       1--2,
and a punctuation
       dash---like
this.

A sentence-ending space should be larger than the
space between words within a sentence.  You
sometimes have to type special commands in
conjunction with punctuation characters to get
this right, as in the following sentence.
       Gnats, gnus, etc.\ all  % `\ ' makes an inter-word space.
       begin with G\@.         % \@ marks end-of-sentence punctuation.
You should check the spaces after periods when
reading your output to make sure you haven't
forgotten any special cases.  Generating an
ellipsis
       \ldots\               % `\ ' is needed after `\ldots' because TeX
                             % ignores spaces after command names like \ldots
                             % made from \ + letters.
                             %
                             % Note how a `%' character causes TeX to ignore
                             % the end of the input line, so these blank lines
                             % do not start a new paragraph.
                             %
with the right spacing around the periods requires
a special command.

\LaTeX\ interprets some common characters as
commands, so you must type special commands to
generate them.  These characters include the
following:
       \$ \& \% \# \{ and \}.

In printing, text is usually emphasized with an
       \emph{italic}
type style.

\begin{em}
   A long segment of text can also be emphasized
   in this way.  Text within such a segment can be
   given \emph{additional} emphasis.
\end{em}

It is sometimes necessary to prevent \LaTeX\ from
breaking a line where it might otherwise do so.
This may be at a space, as between the ``Mr.''\ and
``Jones'' in
       ``Mr.~Jones'',        % ~ produces an unbreakable interword space.
or within a word---especially when the word is a
symbol like
       \mbox{\emph{itemnum}}
that makes little sense when hyphenated across
lines.

Footnotes\footnote{This is an example of a footnote.}
pose no problem.

\LaTeX\ is good at typesetting mathematical formulas
like
       \( x-3y + z = 7 \)
or
       \( a_{1} > x^{2n} + y^{2n} > x' \)
or
       \( \ip{A}{B} = \sum_{i} a_{i} b_{i} \).
The spaces you type in a formula are
ignored.  Remember that a letter like
       $x$                   % $ ... $  and  \( ... \)  are equivalent
is a formula when it denotes a mathematical
symbol, and it should be typed as one.

\section{Displayed Text}

Text is displayed by indenting it from the left
margin.  Quotations are commonly displayed.  There
are short quotations
\begin{quote}
   This is a short quotation.  It consists of a
   single paragraph of text.  See how it is formatted.
\end{quote}
and longer ones.
\begin{quotation}
   This is a longer quotation.  It consists of two
   paragraphs of text, neither of which are
   particularly interesting.

   This is the second paragraph of the quotation.  It
   is just as dull as the first paragraph.
\end{quotation}
Another frequently-displayed structure is a list.
The following is an example of an \emph{itemized}
list.
\begin{itemize}
   \item This is the first item of an itemized list.
         Each item in the list is marked with a ``tick''.
         You don't have to worry about what kind of tick
         mark is used.

   \item This is the second item of the list.  It
         contains another list nested inside it.  The inner
         list is an \emph{enumerated} list.
         \begin{enumerate}
            \item This is the first item of an enumerated
                  list that is nested within the itemized list.

            \item This is the second item of the inner list.
                  \LaTeX\ allows you to nest lists deeper than
                  you really should.
         \end{enumerate}
         This is the rest of the second item of the outer
         list.  It is no more interesting than any other
         part of the item.
   \item This is the third item of the list.
\end{itemize}
You can even display poetry.
\begin{verse}
   There is an environment
    for verse \\             % The \\ command separates lines
   Whose features some poets % within a stanza.
   will curse.

                             % One or more blank lines separate stanzas.

   For instead of making\\
   Them do \emph{all} line breaking, \\
   It allows them to put too many words on a line when they'd rather be
   forced to be terse.
\end{verse}

Mathematical formulas may also be displayed.  A
displayed formula
is
one-line long; multiline
formulas require special formatting instructions.
   \[  \ip{\Gamma}{\psi'} = x'' + y^{2} + z_{i}^{n}\]
Don't start a paragraph with a displayed equation,
nor make one a paragraph by itself.

\end{document}               % End of document.

\section{Bench Huge Section 1175}
% synthetic bench section 1175
% This is a sample LaTeX input file.  (Version of 12 August 2004.)
%
% A '%' character causes TeX to ignore all remaining text on the line,
% and is used for comments like this one.

\documentclass{article}      % Specifies the document class

                             % The preamble begins here.
\title{An Example Document}  % Declares the document's title.
\author{Leslie Lamport}      % Declares the author's name.
\date{January 21, 1994}      % Deleting this command produces today's date.

\newcommand{\ip}[2]{(#1, #2)}
                             % Defines \ip{arg1}{arg2} to mean
                             % (arg1, arg2).

%\newcommand{\ip}[2]{\langle #1 | #2\rangle}
                             % This is an alternative definition of
                             % \ip that is commented out.

\begin{document}             % End of preamble and beginning of text.

\maketitle                   % Produces the title.

This is an example input file.  Comparing it with
the output it generates can show you how to
produce a simple document of your own.

\section{Ordinary Text}      % Produces section heading.  Lower-level
                             % sections are begun with similar
                             % \subsection and \subsubsection commands.

The ends  of words and sentences are marked
  by   spaces. It  doesn't matter how many
spaces    you type; one is as good as 100.  The
end of   a line counts as a space.

One   or more   blank lines denote the  end
of  a paragraph.

Since any number of consecutive spaces are treated
like a single one, the formatting of the input
file makes no difference to
      \LaTeX,                % The \LaTeX command generates the LaTeX logo.
but it makes a difference to you.  When you use
\LaTeX, making your input file as easy to read
as possible will be a great help as you write
your document and when you change it.  This sample
file shows how you can add comments to your own input
file.

Because printing is different from typewriting,
there are a number of things that you have to do
differently when preparing an input file than if
you were just typing the document directly.
Quotation marks like
       ``this''
have to be handled specially, as do quotes within
quotes:
       ``\,`this'            % \, separates the double and single quote.
        is what I just
        wrote, not  `that'\,''.

Dashes come in three sizes: an
       intra-word
dash, a medium dash for number ranges like
       1--2,
and a punctuation
       dash---like
this.

A sentence-ending space should be larger than the
space between words within a sentence.  You
sometimes have to type special commands in
conjunction with punctuation characters to get
this right, as in the following sentence.
       Gnats, gnus, etc.\ all  % `\ ' makes an inter-word space.
       begin with G\@.         % \@ marks end-of-sentence punctuation.
You should check the spaces after periods when
reading your output to make sure you haven't
forgotten any special cases.  Generating an
ellipsis
       \ldots\               % `\ ' is needed after `\ldots' because TeX
                             % ignores spaces after command names like \ldots
                             % made from \ + letters.
                             %
                             % Note how a `%' character causes TeX to ignore
                             % the end of the input line, so these blank lines
                             % do not start a new paragraph.
                             %
with the right spacing around the periods requires
a special command.

\LaTeX\ interprets some common characters as
commands, so you must type special commands to
generate them.  These characters include the
following:
       \$ \& \% \# \{ and \}.

In printing, text is usually emphasized with an
       \emph{italic}
type style.

\begin{em}
   A long segment of text can also be emphasized
   in this way.  Text within such a segment can be
   given \emph{additional} emphasis.
\end{em}

It is sometimes necessary to prevent \LaTeX\ from
breaking a line where it might otherwise do so.
This may be at a space, as between the ``Mr.''\ and
``Jones'' in
       ``Mr.~Jones'',        % ~ produces an unbreakable interword space.
or within a word---especially when the word is a
symbol like
       \mbox{\emph{itemnum}}
that makes little sense when hyphenated across
lines.

Footnotes\footnote{This is an example of a footnote.}
pose no problem.

\LaTeX\ is good at typesetting mathematical formulas
like
       \( x-3y + z = 7 \)
or
       \( a_{1} > x^{2n} + y^{2n} > x' \)
or
       \( \ip{A}{B} = \sum_{i} a_{i} b_{i} \).
The spaces you type in a formula are
ignored.  Remember that a letter like
       $x$                   % $ ... $  and  \( ... \)  are equivalent
is a formula when it denotes a mathematical
symbol, and it should be typed as one.

\section{Displayed Text}

Text is displayed by indenting it from the left
margin.  Quotations are commonly displayed.  There
are short quotations
\begin{quote}
   This is a short quotation.  It consists of a
   single paragraph of text.  See how it is formatted.
\end{quote}
and longer ones.
\begin{quotation}
   This is a longer quotation.  It consists of two
   paragraphs of text, neither of which are
   particularly interesting.

   This is the second paragraph of the quotation.  It
   is just as dull as the first paragraph.
\end{quotation}
Another frequently-displayed structure is a list.
The following is an example of an \emph{itemized}
list.
\begin{itemize}
   \item This is the first item of an itemized list.
         Each item in the list is marked with a ``tick''.
         You don't have to worry about what kind of tick
         mark is used.

   \item This is the second item of the list.  It
         contains another list nested inside it.  The inner
         list is an \emph{enumerated} list.
         \begin{enumerate}
            \item This is the first item of an enumerated
                  list that is nested within the itemized list.

            \item This is the second item of the inner list.
                  \LaTeX\ allows you to nest lists deeper than
                  you really should.
         \end{enumerate}
         This is the rest of the second item of the outer
         list.  It is no more interesting than any other
         part of the item.
   \item This is the third item of the list.
\end{itemize}
You can even display poetry.
\begin{verse}
   There is an environment
    for verse \\             % The \\ command separates lines
   Whose features some poets % within a stanza.
   will curse.

                             % One or more blank lines separate stanzas.

   For instead of making\\
   Them do \emph{all} line breaking, \\
   It allows them to put too many words on a line when they'd rather be
   forced to be terse.
\end{verse}

Mathematical formulas may also be displayed.  A
displayed formula
is
one-line long; multiline
formulas require special formatting instructions.
   \[  \ip{\Gamma}{\psi'} = x'' + y^{2} + z_{i}^{n}\]
Don't start a paragraph with a displayed equation,
nor make one a paragraph by itself.

\end{document}               % End of document.

\section{Bench Huge Section 1176}
% synthetic bench section 1176
% This is a sample LaTeX input file.  (Version of 12 August 2004.)
%
% A '%' character causes TeX to ignore all remaining text on the line,
% and is used for comments like this one.

\documentclass{article}      % Specifies the document class

                             % The preamble begins here.
\title{An Example Document}  % Declares the document's title.
\author{Leslie Lamport}      % Declares the author's name.
\date{January 21, 1994}      % Deleting this command produces today's date.

\newcommand{\ip}[2]{(#1, #2)}
                             % Defines \ip{arg1}{arg2} to mean
                             % (arg1, arg2).

%\newcommand{\ip}[2]{\langle #1 | #2\rangle}
                             % This is an alternative definition of
                             % \ip that is commented out.

\begin{document}             % End of preamble and beginning of text.

\maketitle                   % Produces the title.

This is an example input file.  Comparing it with
the output it generates can show you how to
produce a simple document of your own.

\section{Ordinary Text}      % Produces section heading.  Lower-level
                             % sections are begun with similar
                             % \subsection and \subsubsection commands.

The ends  of words and sentences are marked
  by   spaces. It  doesn't matter how many
spaces    you type; one is as good as 100.  The
end of   a line counts as a space.

One   or more   blank lines denote the  end
of  a paragraph.

Since any number of consecutive spaces are treated
like a single one, the formatting of the input
file makes no difference to
      \LaTeX,                % The \LaTeX command generates the LaTeX logo.
but it makes a difference to you.  When you use
\LaTeX, making your input file as easy to read
as possible will be a great help as you write
your document and when you change it.  This sample
file shows how you can add comments to your own input
file.

Because printing is different from typewriting,
there are a number of things that you have to do
differently when preparing an input file than if
you were just typing the document directly.
Quotation marks like
       ``this''
have to be handled specially, as do quotes within
quotes:
       ``\,`this'            % \, separates the double and single quote.
        is what I just
        wrote, not  `that'\,''.

Dashes come in three sizes: an
       intra-word
dash, a medium dash for number ranges like
       1--2,
and a punctuation
       dash---like
this.

A sentence-ending space should be larger than the
space between words within a sentence.  You
sometimes have to type special commands in
conjunction with punctuation characters to get
this right, as in the following sentence.
       Gnats, gnus, etc.\ all  % `\ ' makes an inter-word space.
       begin with G\@.         % \@ marks end-of-sentence punctuation.
You should check the spaces after periods when
reading your output to make sure you haven't
forgotten any special cases.  Generating an
ellipsis
       \ldots\               % `\ ' is needed after `\ldots' because TeX
                             % ignores spaces after command names like \ldots
                             % made from \ + letters.
                             %
                             % Note how a `%' character causes TeX to ignore
                             % the end of the input line, so these blank lines
                             % do not start a new paragraph.
                             %
with the right spacing around the periods requires
a special command.

\LaTeX\ interprets some common characters as
commands, so you must type special commands to
generate them.  These characters include the
following:
       \$ \& \% \# \{ and \}.

In printing, text is usually emphasized with an
       \emph{italic}
type style.

\begin{em}
   A long segment of text can also be emphasized
   in this way.  Text within such a segment can be
   given \emph{additional} emphasis.
\end{em}

It is sometimes necessary to prevent \LaTeX\ from
breaking a line where it might otherwise do so.
This may be at a space, as between the ``Mr.''\ and
``Jones'' in
       ``Mr.~Jones'',        % ~ produces an unbreakable interword space.
or within a word---especially when the word is a
symbol like
       \mbox{\emph{itemnum}}
that makes little sense when hyphenated across
lines.

Footnotes\footnote{This is an example of a footnote.}
pose no problem.

\LaTeX\ is good at typesetting mathematical formulas
like
       \( x-3y + z = 7 \)
or
       \( a_{1} > x^{2n} + y^{2n} > x' \)
or
       \( \ip{A}{B} = \sum_{i} a_{i} b_{i} \).
The spaces you type in a formula are
ignored.  Remember that a letter like
       $x$                   % $ ... $  and  \( ... \)  are equivalent
is a formula when it denotes a mathematical
symbol, and it should be typed as one.

\section{Displayed Text}

Text is displayed by indenting it from the left
margin.  Quotations are commonly displayed.  There
are short quotations
\begin{quote}
   This is a short quotation.  It consists of a
   single paragraph of text.  See how it is formatted.
\end{quote}
and longer ones.
\begin{quotation}
   This is a longer quotation.  It consists of two
   paragraphs of text, neither of which are
   particularly interesting.

   This is the second paragraph of the quotation.  It
   is just as dull as the first paragraph.
\end{quotation}
Another frequently-displayed structure is a list.
The following is an example of an \emph{itemized}
list.
\begin{itemize}
   \item This is the first item of an itemized list.
         Each item in the list is marked with a ``tick''.
         You don't have to worry about what kind of tick
         mark is used.

   \item This is the second item of the list.  It
         contains another list nested inside it.  The inner
         list is an \emph{enumerated} list.
         \begin{enumerate}
            \item This is the first item of an enumerated
                  list that is nested within the itemized list.

            \item This is the second item of the inner list.
                  \LaTeX\ allows you to nest lists deeper than
                  you really should.
         \end{enumerate}
         This is the rest of the second item of the outer
         list.  It is no more interesting than any other
         part of the item.
   \item This is the third item of the list.
\end{itemize}
You can even display poetry.
\begin{verse}
   There is an environment
    for verse \\             % The \\ command separates lines
   Whose features some poets % within a stanza.
   will curse.

                             % One or more blank lines separate stanzas.

   For instead of making\\
   Them do \emph{all} line breaking, \\
   It allows them to put too many words on a line when they'd rather be
   forced to be terse.
\end{verse}

Mathematical formulas may also be displayed.  A
displayed formula
is
one-line long; multiline
formulas require special formatting instructions.
   \[  \ip{\Gamma}{\psi'} = x'' + y^{2} + z_{i}^{n}\]
Don't start a paragraph with a displayed equation,
nor make one a paragraph by itself.

\end{document}               % End of document.

\section{Bench Huge Section 1177}
% synthetic bench section 1177
% This is a sample LaTeX input file.  (Version of 12 August 2004.)
%
% A '%' character causes TeX to ignore all remaining text on the line,
% and is used for comments like this one.

\documentclass{article}      % Specifies the document class

                             % The preamble begins here.
\title{An Example Document}  % Declares the document's title.
\author{Leslie Lamport}      % Declares the author's name.
\date{January 21, 1994}      % Deleting this command produces today's date.

\newcommand{\ip}[2]{(#1, #2)}
                             % Defines \ip{arg1}{arg2} to mean
                             % (arg1, arg2).

%\newcommand{\ip}[2]{\langle #1 | #2\rangle}
                             % This is an alternative definition of
                             % \ip that is commented out.

\begin{document}             % End of preamble and beginning of text.

\maketitle                   % Produces the title.

This is an example input file.  Comparing it with
the output it generates can show you how to
produce a simple document of your own.

\section{Ordinary Text}      % Produces section heading.  Lower-level
                             % sections are begun with similar
                             % \subsection and \subsubsection commands.

The ends  of words and sentences are marked
  by   spaces. It  doesn't matter how many
spaces    you type; one is as good as 100.  The
end of   a line counts as a space.

One   or more   blank lines denote the  end
of  a paragraph.

Since any number of consecutive spaces are treated
like a single one, the formatting of the input
file makes no difference to
      \LaTeX,                % The \LaTeX command generates the LaTeX logo.
but it makes a difference to you.  When you use
\LaTeX, making your input file as easy to read
as possible will be a great help as you write
your document and when you change it.  This sample
file shows how you can add comments to your own input
file.

Because printing is different from typewriting,
there are a number of things that you have to do
differently when preparing an input file than if
you were just typing the document directly.
Quotation marks like
       ``this''
have to be handled specially, as do quotes within
quotes:
       ``\,`this'            % \, separates the double and single quote.
        is what I just
        wrote, not  `that'\,''.

Dashes come in three sizes: an
       intra-word
dash, a medium dash for number ranges like
       1--2,
and a punctuation
       dash---like
this.

A sentence-ending space should be larger than the
space between words within a sentence.  You
sometimes have to type special commands in
conjunction with punctuation characters to get
this right, as in the following sentence.
       Gnats, gnus, etc.\ all  % `\ ' makes an inter-word space.
       begin with G\@.         % \@ marks end-of-sentence punctuation.
You should check the spaces after periods when
reading your output to make sure you haven't
forgotten any special cases.  Generating an
ellipsis
       \ldots\               % `\ ' is needed after `\ldots' because TeX
                             % ignores spaces after command names like \ldots
                             % made from \ + letters.
                             %
                             % Note how a `%' character causes TeX to ignore
                             % the end of the input line, so these blank lines
                             % do not start a new paragraph.
                             %
with the right spacing around the periods requires
a special command.

\LaTeX\ interprets some common characters as
commands, so you must type special commands to
generate them.  These characters include the
following:
       \$ \& \% \# \{ and \}.

In printing, text is usually emphasized with an
       \emph{italic}
type style.

\begin{em}
   A long segment of text can also be emphasized
   in this way.  Text within such a segment can be
   given \emph{additional} emphasis.
\end{em}

It is sometimes necessary to prevent \LaTeX\ from
breaking a line where it might otherwise do so.
This may be at a space, as between the ``Mr.''\ and
``Jones'' in
       ``Mr.~Jones'',        % ~ produces an unbreakable interword space.
or within a word---especially when the word is a
symbol like
       \mbox{\emph{itemnum}}
that makes little sense when hyphenated across
lines.

Footnotes\footnote{This is an example of a footnote.}
pose no problem.

\LaTeX\ is good at typesetting mathematical formulas
like
       \( x-3y + z = 7 \)
or
       \( a_{1} > x^{2n} + y^{2n} > x' \)
or
       \( \ip{A}{B} = \sum_{i} a_{i} b_{i} \).
The spaces you type in a formula are
ignored.  Remember that a letter like
       $x$                   % $ ... $  and  \( ... \)  are equivalent
is a formula when it denotes a mathematical
symbol, and it should be typed as one.

\section{Displayed Text}

Text is displayed by indenting it from the left
margin.  Quotations are commonly displayed.  There
are short quotations
\begin{quote}
   This is a short quotation.  It consists of a
   single paragraph of text.  See how it is formatted.
\end{quote}
and longer ones.
\begin{quotation}
   This is a longer quotation.  It consists of two
   paragraphs of text, neither of which are
   particularly interesting.

   This is the second paragraph of the quotation.  It
   is just as dull as the first paragraph.
\end{quotation}
Another frequently-displayed structure is a list.
The following is an example of an \emph{itemized}
list.
\begin{itemize}
   \item This is the first item of an itemized list.
         Each item in the list is marked with a ``tick''.
         You don't have to worry about what kind of tick
         mark is used.

   \item This is the second item of the list.  It
         contains another list nested inside it.  The inner
         list is an \emph{enumerated} list.
         \begin{enumerate}
            \item This is the first item of an enumerated
                  list that is nested within the itemized list.

            \item This is the second item of the inner list.
                  \LaTeX\ allows you to nest lists deeper than
                  you really should.
         \end{enumerate}
         This is the rest of the second item of the outer
         list.  It is no more interesting than any other
         part of the item.
   \item This is the third item of the list.
\end{itemize}
You can even display poetry.
\begin{verse}
   There is an environment
    for verse \\             % The \\ command separates lines
   Whose features some poets % within a stanza.
   will curse.

                             % One or more blank lines separate stanzas.

   For instead of making\\
   Them do \emph{all} line breaking, \\
   It allows them to put too many words on a line when they'd rather be
   forced to be terse.
\end{verse}

Mathematical formulas may also be displayed.  A
displayed formula
is
one-line long; multiline
formulas require special formatting instructions.
   \[  \ip{\Gamma}{\psi'} = x'' + y^{2} + z_{i}^{n}\]
Don't start a paragraph with a displayed equation,
nor make one a paragraph by itself.

\end{document}               % End of document.

\section{Bench Huge Section 1178}
% synthetic bench section 1178
% This is a sample LaTeX input file.  (Version of 12 August 2004.)
%
% A '%' character causes TeX to ignore all remaining text on the line,
% and is used for comments like this one.

\documentclass{article}      % Specifies the document class

                             % The preamble begins here.
\title{An Example Document}  % Declares the document's title.
\author{Leslie Lamport}      % Declares the author's name.
\date{January 21, 1994}      % Deleting this command produces today's date.

\newcommand{\ip}[2]{(#1, #2)}
                             % Defines \ip{arg1}{arg2} to mean
                             % (arg1, arg2).

%\newcommand{\ip}[2]{\langle #1 | #2\rangle}
                             % This is an alternative definition of
                             % \ip that is commented out.

\begin{document}             % End of preamble and beginning of text.

\maketitle                   % Produces the title.

This is an example input file.  Comparing it with
the output it generates can show you how to
produce a simple document of your own.

\section{Ordinary Text}      % Produces section heading.  Lower-level
                             % sections are begun with similar
                             % \subsection and \subsubsection commands.

The ends  of words and sentences are marked
  by   spaces. It  doesn't matter how many
spaces    you type; one is as good as 100.  The
end of   a line counts as a space.

One   or more   blank lines denote the  end
of  a paragraph.

Since any number of consecutive spaces are treated
like a single one, the formatting of the input
file makes no difference to
      \LaTeX,                % The \LaTeX command generates the LaTeX logo.
but it makes a difference to you.  When you use
\LaTeX, making your input file as easy to read
as possible will be a great help as you write
your document and when you change it.  This sample
file shows how you can add comments to your own input
file.

Because printing is different from typewriting,
there are a number of things that you have to do
differently when preparing an input file than if
you were just typing the document directly.
Quotation marks like
       ``this''
have to be handled specially, as do quotes within
quotes:
       ``\,`this'            % \, separates the double and single quote.
        is what I just
        wrote, not  `that'\,''.

Dashes come in three sizes: an
       intra-word
dash, a medium dash for number ranges like
       1--2,
and a punctuation
       dash---like
this.

A sentence-ending space should be larger than the
space between words within a sentence.  You
sometimes have to type special commands in
conjunction with punctuation characters to get
this right, as in the following sentence.
       Gnats, gnus, etc.\ all  % `\ ' makes an inter-word space.
       begin with G\@.         % \@ marks end-of-sentence punctuation.
You should check the spaces after periods when
reading your output to make sure you haven't
forgotten any special cases.  Generating an
ellipsis
       \ldots\               % `\ ' is needed after `\ldots' because TeX
                             % ignores spaces after command names like \ldots
                             % made from \ + letters.
                             %
                             % Note how a `%' character causes TeX to ignore
                             % the end of the input line, so these blank lines
                             % do not start a new paragraph.
                             %
with the right spacing around the periods requires
a special command.

\LaTeX\ interprets some common characters as
commands, so you must type special commands to
generate them.  These characters include the
following:
       \$ \& \% \# \{ and \}.

In printing, text is usually emphasized with an
       \emph{italic}
type style.

\begin{em}
   A long segment of text can also be emphasized
   in this way.  Text within such a segment can be
   given \emph{additional} emphasis.
\end{em}

It is sometimes necessary to prevent \LaTeX\ from
breaking a line where it might otherwise do so.
This may be at a space, as between the ``Mr.''\ and
``Jones'' in
       ``Mr.~Jones'',        % ~ produces an unbreakable interword space.
or within a word---especially when the word is a
symbol like
       \mbox{\emph{itemnum}}
that makes little sense when hyphenated across
lines.

Footnotes\footnote{This is an example of a footnote.}
pose no problem.

\LaTeX\ is good at typesetting mathematical formulas
like
       \( x-3y + z = 7 \)
or
       \( a_{1} > x^{2n} + y^{2n} > x' \)
or
       \( \ip{A}{B} = \sum_{i} a_{i} b_{i} \).
The spaces you type in a formula are
ignored.  Remember that a letter like
       $x$                   % $ ... $  and  \( ... \)  are equivalent
is a formula when it denotes a mathematical
symbol, and it should be typed as one.

\section{Displayed Text}

Text is displayed by indenting it from the left
margin.  Quotations are commonly displayed.  There
are short quotations
\begin{quote}
   This is a short quotation.  It consists of a
   single paragraph of text.  See how it is formatted.
\end{quote}
and longer ones.
\begin{quotation}
   This is a longer quotation.  It consists of two
   paragraphs of text, neither of which are
   particularly interesting.

   This is the second paragraph of the quotation.  It
   is just as dull as the first paragraph.
\end{quotation}
Another frequently-displayed structure is a list.
The following is an example of an \emph{itemized}
list.
\begin{itemize}
   \item This is the first item of an itemized list.
         Each item in the list is marked with a ``tick''.
         You don't have to worry about what kind of tick
         mark is used.

   \item This is the second item of the list.  It
         contains another list nested inside it.  The inner
         list is an \emph{enumerated} list.
         \begin{enumerate}
            \item This is the first item of an enumerated
                  list that is nested within the itemized list.

            \item This is the second item of the inner list.
                  \LaTeX\ allows you to nest lists deeper than
                  you really should.
         \end{enumerate}
         This is the rest of the second item of the outer
         list.  It is no more interesting than any other
         part of the item.
   \item This is the third item of the list.
\end{itemize}
You can even display poetry.
\begin{verse}
   There is an environment
    for verse \\             % The \\ command separates lines
   Whose features some poets % within a stanza.
   will curse.

                             % One or more blank lines separate stanzas.

   For instead of making\\
   Them do \emph{all} line breaking, \\
   It allows them to put too many words on a line when they'd rather be
   forced to be terse.
\end{verse}

Mathematical formulas may also be displayed.  A
displayed formula
is
one-line long; multiline
formulas require special formatting instructions.
   \[  \ip{\Gamma}{\psi'} = x'' + y^{2} + z_{i}^{n}\]
Don't start a paragraph with a displayed equation,
nor make one a paragraph by itself.

\end{document}               % End of document.

\section{Bench Huge Section 1179}
% synthetic bench section 1179
% This is a sample LaTeX input file.  (Version of 12 August 2004.)
%
% A '%' character causes TeX to ignore all remaining text on the line,
% and is used for comments like this one.

\documentclass{article}      % Specifies the document class

                             % The preamble begins here.
\title{An Example Document}  % Declares the document's title.
\author{Leslie Lamport}      % Declares the author's name.
\date{January 21, 1994}      % Deleting this command produces today's date.

\newcommand{\ip}[2]{(#1, #2)}
                             % Defines \ip{arg1}{arg2} to mean
                             % (arg1, arg2).

%\newcommand{\ip}[2]{\langle #1 | #2\rangle}
                             % This is an alternative definition of
                             % \ip that is commented out.

\begin{document}             % End of preamble and beginning of text.

\maketitle                   % Produces the title.

This is an example input file.  Comparing it with
the output it generates can show you how to
produce a simple document of your own.

\section{Ordinary Text}      % Produces section heading.  Lower-level
                             % sections are begun with similar
                             % \subsection and \subsubsection commands.

The ends  of words and sentences are marked
  by   spaces. It  doesn't matter how many
spaces    you type; one is as good as 100.  The
end of   a line counts as a space.

One   or more   blank lines denote the  end
of  a paragraph.

Since any number of consecutive spaces are treated
like a single one, the formatting of the input
file makes no difference to
      \LaTeX,                % The \LaTeX command generates the LaTeX logo.
but it makes a difference to you.  When you use
\LaTeX, making your input file as easy to read
as possible will be a great help as you write
your document and when you change it.  This sample
file shows how you can add comments to your own input
file.

Because printing is different from typewriting,
there are a number of things that you have to do
differently when preparing an input file than if
you were just typing the document directly.
Quotation marks like
       ``this''
have to be handled specially, as do quotes within
quotes:
       ``\,`this'            % \, separates the double and single quote.
        is what I just
        wrote, not  `that'\,''.

Dashes come in three sizes: an
       intra-word
dash, a medium dash for number ranges like
       1--2,
and a punctuation
       dash---like
this.

A sentence-ending space should be larger than the
space between words within a sentence.  You
sometimes have to type special commands in
conjunction with punctuation characters to get
this right, as in the following sentence.
       Gnats, gnus, etc.\ all  % `\ ' makes an inter-word space.
       begin with G\@.         % \@ marks end-of-sentence punctuation.
You should check the spaces after periods when
reading your output to make sure you haven't
forgotten any special cases.  Generating an
ellipsis
       \ldots\               % `\ ' is needed after `\ldots' because TeX
                             % ignores spaces after command names like \ldots
                             % made from \ + letters.
                             %
                             % Note how a `%' character causes TeX to ignore
                             % the end of the input line, so these blank lines
                             % do not start a new paragraph.
                             %
with the right spacing around the periods requires
a special command.

\LaTeX\ interprets some common characters as
commands, so you must type special commands to
generate them.  These characters include the
following:
       \$ \& \% \# \{ and \}.

In printing, text is usually emphasized with an
       \emph{italic}
type style.

\begin{em}
   A long segment of text can also be emphasized
   in this way.  Text within such a segment can be
   given \emph{additional} emphasis.
\end{em}

It is sometimes necessary to prevent \LaTeX\ from
breaking a line where it might otherwise do so.
This may be at a space, as between the ``Mr.''\ and
``Jones'' in
       ``Mr.~Jones'',        % ~ produces an unbreakable interword space.
or within a word---especially when the word is a
symbol like
       \mbox{\emph{itemnum}}
that makes little sense when hyphenated across
lines.

Footnotes\footnote{This is an example of a footnote.}
pose no problem.

\LaTeX\ is good at typesetting mathematical formulas
like
       \( x-3y + z = 7 \)
or
       \( a_{1} > x^{2n} + y^{2n} > x' \)
or
       \( \ip{A}{B} = \sum_{i} a_{i} b_{i} \).
The spaces you type in a formula are
ignored.  Remember that a letter like
       $x$                   % $ ... $  and  \( ... \)  are equivalent
is a formula when it denotes a mathematical
symbol, and it should be typed as one.

\section{Displayed Text}

Text is displayed by indenting it from the left
margin.  Quotations are commonly displayed.  There
are short quotations
\begin{quote}
   This is a short quotation.  It consists of a
   single paragraph of text.  See how it is formatted.
\end{quote}
and longer ones.
\begin{quotation}
   This is a longer quotation.  It consists of two
   paragraphs of text, neither of which are
   particularly interesting.

   This is the second paragraph of the quotation.  It
   is just as dull as the first paragraph.
\end{quotation}
Another frequently-displayed structure is a list.
The following is an example of an \emph{itemized}
list.
\begin{itemize}
   \item This is the first item of an itemized list.
         Each item in the list is marked with a ``tick''.
         You don't have to worry about what kind of tick
         mark is used.

   \item This is the second item of the list.  It
         contains another list nested inside it.  The inner
         list is an \emph{enumerated} list.
         \begin{enumerate}
            \item This is the first item of an enumerated
                  list that is nested within the itemized list.

            \item This is the second item of the inner list.
                  \LaTeX\ allows you to nest lists deeper than
                  you really should.
         \end{enumerate}
         This is the rest of the second item of the outer
         list.  It is no more interesting than any other
         part of the item.
   \item This is the third item of the list.
\end{itemize}
You can even display poetry.
\begin{verse}
   There is an environment
    for verse \\             % The \\ command separates lines
   Whose features some poets % within a stanza.
   will curse.

                             % One or more blank lines separate stanzas.

   For instead of making\\
   Them do \emph{all} line breaking, \\
   It allows them to put too many words on a line when they'd rather be
   forced to be terse.
\end{verse}

Mathematical formulas may also be displayed.  A
displayed formula
is
one-line long; multiline
formulas require special formatting instructions.
   \[  \ip{\Gamma}{\psi'} = x'' + y^{2} + z_{i}^{n}\]
Don't start a paragraph with a displayed equation,
nor make one a paragraph by itself.

\end{document}               % End of document.

\section{Bench Huge Section 1180}
% synthetic bench section 1180
% This is a sample LaTeX input file.  (Version of 12 August 2004.)
%
% A '%' character causes TeX to ignore all remaining text on the line,
% and is used for comments like this one.

\documentclass{article}      % Specifies the document class

                             % The preamble begins here.
\title{An Example Document}  % Declares the document's title.
\author{Leslie Lamport}      % Declares the author's name.
\date{January 21, 1994}      % Deleting this command produces today's date.

\newcommand{\ip}[2]{(#1, #2)}
                             % Defines \ip{arg1}{arg2} to mean
                             % (arg1, arg2).

%\newcommand{\ip}[2]{\langle #1 | #2\rangle}
                             % This is an alternative definition of
                             % \ip that is commented out.

\begin{document}             % End of preamble and beginning of text.

\maketitle                   % Produces the title.

This is an example input file.  Comparing it with
the output it generates can show you how to
produce a simple document of your own.

\section{Ordinary Text}      % Produces section heading.  Lower-level
                             % sections are begun with similar
                             % \subsection and \subsubsection commands.

The ends  of words and sentences are marked
  by   spaces. It  doesn't matter how many
spaces    you type; one is as good as 100.  The
end of   a line counts as a space.

One   or more   blank lines denote the  end
of  a paragraph.

Since any number of consecutive spaces are treated
like a single one, the formatting of the input
file makes no difference to
      \LaTeX,                % The \LaTeX command generates the LaTeX logo.
but it makes a difference to you.  When you use
\LaTeX, making your input file as easy to read
as possible will be a great help as you write
your document and when you change it.  This sample
file shows how you can add comments to your own input
file.

Because printing is different from typewriting,
there are a number of things that you have to do
differently when preparing an input file than if
you were just typing the document directly.
Quotation marks like
       ``this''
have to be handled specially, as do quotes within
quotes:
       ``\,`this'            % \, separates the double and single quote.
        is what I just
        wrote, not  `that'\,''.

Dashes come in three sizes: an
       intra-word
dash, a medium dash for number ranges like
       1--2,
and a punctuation
       dash---like
this.

A sentence-ending space should be larger than the
space between words within a sentence.  You
sometimes have to type special commands in
conjunction with punctuation characters to get
this right, as in the following sentence.
       Gnats, gnus, etc.\ all  % `\ ' makes an inter-word space.
       begin with G\@.         % \@ marks end-of-sentence punctuation.
You should check the spaces after periods when
reading your output to make sure you haven't
forgotten any special cases.  Generating an
ellipsis
       \ldots\               % `\ ' is needed after `\ldots' because TeX
                             % ignores spaces after command names like \ldots
                             % made from \ + letters.
                             %
                             % Note how a `%' character causes TeX to ignore
                             % the end of the input line, so these blank lines
                             % do not start a new paragraph.
                             %
with the right spacing around the periods requires
a special command.

\LaTeX\ interprets some common characters as
commands, so you must type special commands to
generate them.  These characters include the
following:
       \$ \& \% \# \{ and \}.

In printing, text is usually emphasized with an
       \emph{italic}
type style.

\begin{em}
   A long segment of text can also be emphasized
   in this way.  Text within such a segment can be
   given \emph{additional} emphasis.
\end{em}

It is sometimes necessary to prevent \LaTeX\ from
breaking a line where it might otherwise do so.
This may be at a space, as between the ``Mr.''\ and
``Jones'' in
       ``Mr.~Jones'',        % ~ produces an unbreakable interword space.
or within a word---especially when the word is a
symbol like
       \mbox{\emph{itemnum}}
that makes little sense when hyphenated across
lines.

Footnotes\footnote{This is an example of a footnote.}
pose no problem.

\LaTeX\ is good at typesetting mathematical formulas
like
       \( x-3y + z = 7 \)
or
       \( a_{1} > x^{2n} + y^{2n} > x' \)
or
       \( \ip{A}{B} = \sum_{i} a_{i} b_{i} \).
The spaces you type in a formula are
ignored.  Remember that a letter like
       $x$                   % $ ... $  and  \( ... \)  are equivalent
is a formula when it denotes a mathematical
symbol, and it should be typed as one.

\section{Displayed Text}

Text is displayed by indenting it from the left
margin.  Quotations are commonly displayed.  There
are short quotations
\begin{quote}
   This is a short quotation.  It consists of a
   single paragraph of text.  See how it is formatted.
\end{quote}
and longer ones.
\begin{quotation}
   This is a longer quotation.  It consists of two
   paragraphs of text, neither of which are
   particularly interesting.

   This is the second paragraph of the quotation.  It
   is just as dull as the first paragraph.
\end{quotation}
Another frequently-displayed structure is a list.
The following is an example of an \emph{itemized}
list.
\begin{itemize}
   \item This is the first item of an itemized list.
         Each item in the list is marked with a ``tick''.
         You don't have to worry about what kind of tick
         mark is used.

   \item This is the second item of the list.  It
         contains another list nested inside it.  The inner
         list is an \emph{enumerated} list.
         \begin{enumerate}
            \item This is the first item of an enumerated
                  list that is nested within the itemized list.

            \item This is the second item of the inner list.
                  \LaTeX\ allows you to nest lists deeper than
                  you really should.
         \end{enumerate}
         This is the rest of the second item of the outer
         list.  It is no more interesting than any other
         part of the item.
   \item This is the third item of the list.
\end{itemize}
You can even display poetry.
\begin{verse}
   There is an environment
    for verse \\             % The \\ command separates lines
   Whose features some poets % within a stanza.
   will curse.

                             % One or more blank lines separate stanzas.

   For instead of making\\
   Them do \emph{all} line breaking, \\
   It allows them to put too many words on a line when they'd rather be
   forced to be terse.
\end{verse}

Mathematical formulas may also be displayed.  A
displayed formula
is
one-line long; multiline
formulas require special formatting instructions.
   \[  \ip{\Gamma}{\psi'} = x'' + y^{2} + z_{i}^{n}\]
Don't start a paragraph with a displayed equation,
nor make one a paragraph by itself.

\end{document}               % End of document.

\section{Bench Huge Section 1181}
% synthetic bench section 1181
% This is a sample LaTeX input file.  (Version of 12 August 2004.)
%
% A '%' character causes TeX to ignore all remaining text on the line,
% and is used for comments like this one.

\documentclass{article}      % Specifies the document class

                             % The preamble begins here.
\title{An Example Document}  % Declares the document's title.
\author{Leslie Lamport}      % Declares the author's name.
\date{January 21, 1994}      % Deleting this command produces today's date.

\newcommand{\ip}[2]{(#1, #2)}
                             % Defines \ip{arg1}{arg2} to mean
                             % (arg1, arg2).

%\newcommand{\ip}[2]{\langle #1 | #2\rangle}
                             % This is an alternative definition of
                             % \ip that is commented out.

\begin{document}             % End of preamble and beginning of text.

\maketitle                   % Produces the title.

This is an example input file.  Comparing it with
the output it generates can show you how to
produce a simple document of your own.

\section{Ordinary Text}      % Produces section heading.  Lower-level
                             % sections are begun with similar
                             % \subsection and \subsubsection commands.

The ends  of words and sentences are marked
  by   spaces. It  doesn't matter how many
spaces    you type; one is as good as 100.  The
end of   a line counts as a space.

One   or more   blank lines denote the  end
of  a paragraph.

Since any number of consecutive spaces are treated
like a single one, the formatting of the input
file makes no difference to
      \LaTeX,                % The \LaTeX command generates the LaTeX logo.
but it makes a difference to you.  When you use
\LaTeX, making your input file as easy to read
as possible will be a great help as you write
your document and when you change it.  This sample
file shows how you can add comments to your own input
file.

Because printing is different from typewriting,
there are a number of things that you have to do
differently when preparing an input file than if
you were just typing the document directly.
Quotation marks like
       ``this''
have to be handled specially, as do quotes within
quotes:
       ``\,`this'            % \, separates the double and single quote.
        is what I just
        wrote, not  `that'\,''.

Dashes come in three sizes: an
       intra-word
dash, a medium dash for number ranges like
       1--2,
and a punctuation
       dash---like
this.

A sentence-ending space should be larger than the
space between words within a sentence.  You
sometimes have to type special commands in
conjunction with punctuation characters to get
this right, as in the following sentence.
       Gnats, gnus, etc.\ all  % `\ ' makes an inter-word space.
       begin with G\@.         % \@ marks end-of-sentence punctuation.
You should check the spaces after periods when
reading your output to make sure you haven't
forgotten any special cases.  Generating an
ellipsis
       \ldots\               % `\ ' is needed after `\ldots' because TeX
                             % ignores spaces after command names like \ldots
                             % made from \ + letters.
                             %
                             % Note how a `%' character causes TeX to ignore
                             % the end of the input line, so these blank lines
                             % do not start a new paragraph.
                             %
with the right spacing around the periods requires
a special command.

\LaTeX\ interprets some common characters as
commands, so you must type special commands to
generate them.  These characters include the
following:
       \$ \& \% \# \{ and \}.

In printing, text is usually emphasized with an
       \emph{italic}
type style.

\begin{em}
   A long segment of text can also be emphasized
   in this way.  Text within such a segment can be
   given \emph{additional} emphasis.
\end{em}

It is sometimes necessary to prevent \LaTeX\ from
breaking a line where it might otherwise do so.
This may be at a space, as between the ``Mr.''\ and
``Jones'' in
       ``Mr.~Jones'',        % ~ produces an unbreakable interword space.
or within a word---especially when the word is a
symbol like
       \mbox{\emph{itemnum}}
that makes little sense when hyphenated across
lines.

Footnotes\footnote{This is an example of a footnote.}
pose no problem.

\LaTeX\ is good at typesetting mathematical formulas
like
       \( x-3y + z = 7 \)
or
       \( a_{1} > x^{2n} + y^{2n} > x' \)
or
       \( \ip{A}{B} = \sum_{i} a_{i} b_{i} \).
The spaces you type in a formula are
ignored.  Remember that a letter like
       $x$                   % $ ... $  and  \( ... \)  are equivalent
is a formula when it denotes a mathematical
symbol, and it should be typed as one.

\section{Displayed Text}

Text is displayed by indenting it from the left
margin.  Quotations are commonly displayed.  There
are short quotations
\begin{quote}
   This is a short quotation.  It consists of a
   single paragraph of text.  See how it is formatted.
\end{quote}
and longer ones.
\begin{quotation}
   This is a longer quotation.  It consists of two
   paragraphs of text, neither of which are
   particularly interesting.

   This is the second paragraph of the quotation.  It
   is just as dull as the first paragraph.
\end{quotation}
Another frequently-displayed structure is a list.
The following is an example of an \emph{itemized}
list.
\begin{itemize}
   \item This is the first item of an itemized list.
         Each item in the list is marked with a ``tick''.
         You don't have to worry about what kind of tick
         mark is used.

   \item This is the second item of the list.  It
         contains another list nested inside it.  The inner
         list is an \emph{enumerated} list.
         \begin{enumerate}
            \item This is the first item of an enumerated
                  list that is nested within the itemized list.

            \item This is the second item of the inner list.
                  \LaTeX\ allows you to nest lists deeper than
                  you really should.
         \end{enumerate}
         This is the rest of the second item of the outer
         list.  It is no more interesting than any other
         part of the item.
   \item This is the third item of the list.
\end{itemize}
You can even display poetry.
\begin{verse}
   There is an environment
    for verse \\             % The \\ command separates lines
   Whose features some poets % within a stanza.
   will curse.

                             % One or more blank lines separate stanzas.

   For instead of making\\
   Them do \emph{all} line breaking, \\
   It allows them to put too many words on a line when they'd rather be
   forced to be terse.
\end{verse}

Mathematical formulas may also be displayed.  A
displayed formula
is
one-line long; multiline
formulas require special formatting instructions.
   \[  \ip{\Gamma}{\psi'} = x'' + y^{2} + z_{i}^{n}\]
Don't start a paragraph with a displayed equation,
nor make one a paragraph by itself.

\end{document}               % End of document.

\section{Bench Huge Section 1182}
% synthetic bench section 1182
% This is a sample LaTeX input file.  (Version of 12 August 2004.)
%
% A '%' character causes TeX to ignore all remaining text on the line,
% and is used for comments like this one.

\documentclass{article}      % Specifies the document class

                             % The preamble begins here.
\title{An Example Document}  % Declares the document's title.
\author{Leslie Lamport}      % Declares the author's name.
\date{January 21, 1994}      % Deleting this command produces today's date.

\newcommand{\ip}[2]{(#1, #2)}
                             % Defines \ip{arg1}{arg2} to mean
                             % (arg1, arg2).

%\newcommand{\ip}[2]{\langle #1 | #2\rangle}
                             % This is an alternative definition of
                             % \ip that is commented out.

\begin{document}             % End of preamble and beginning of text.

\maketitle                   % Produces the title.

This is an example input file.  Comparing it with
the output it generates can show you how to
produce a simple document of your own.

\section{Ordinary Text}      % Produces section heading.  Lower-level
                             % sections are begun with similar
                             % \subsection and \subsubsection commands.

The ends  of words and sentences are marked
  by   spaces. It  doesn't matter how many
spaces    you type; one is as good as 100.  The
end of   a line counts as a space.

One   or more   blank lines denote the  end
of  a paragraph.

Since any number of consecutive spaces are treated
like a single one, the formatting of the input
file makes no difference to
      \LaTeX,                % The \LaTeX command generates the LaTeX logo.
but it makes a difference to you.  When you use
\LaTeX, making your input file as easy to read
as possible will be a great help as you write
your document and when you change it.  This sample
file shows how you can add comments to your own input
file.

Because printing is different from typewriting,
there are a number of things that you have to do
differently when preparing an input file than if
you were just typing the document directly.
Quotation marks like
       ``this''
have to be handled specially, as do quotes within
quotes:
       ``\,`this'            % \, separates the double and single quote.
        is what I just
        wrote, not  `that'\,''.

Dashes come in three sizes: an
       intra-word
dash, a medium dash for number ranges like
       1--2,
and a punctuation
       dash---like
this.

A sentence-ending space should be larger than the
space between words within a sentence.  You
sometimes have to type special commands in
conjunction with punctuation characters to get
this right, as in the following sentence.
       Gnats, gnus, etc.\ all  % `\ ' makes an inter-word space.
       begin with G\@.         % \@ marks end-of-sentence punctuation.
You should check the spaces after periods when
reading your output to make sure you haven't
forgotten any special cases.  Generating an
ellipsis
       \ldots\               % `\ ' is needed after `\ldots' because TeX
                             % ignores spaces after command names like \ldots
                             % made from \ + letters.
                             %
                             % Note how a `%' character causes TeX to ignore
                             % the end of the input line, so these blank lines
                             % do not start a new paragraph.
                             %
with the right spacing around the periods requires
a special command.

\LaTeX\ interprets some common characters as
commands, so you must type special commands to
generate them.  These characters include the
following:
       \$ \& \% \# \{ and \}.

In printing, text is usually emphasized with an
       \emph{italic}
type style.

\begin{em}
   A long segment of text can also be emphasized
   in this way.  Text within such a segment can be
   given \emph{additional} emphasis.
\end{em}

It is sometimes necessary to prevent \LaTeX\ from
breaking a line where it might otherwise do so.
This may be at a space, as between the ``Mr.''\ and
``Jones'' in
       ``Mr.~Jones'',        % ~ produces an unbreakable interword space.
or within a word---especially when the word is a
symbol like
       \mbox{\emph{itemnum}}
that makes little sense when hyphenated across
lines.

Footnotes\footnote{This is an example of a footnote.}
pose no problem.

\LaTeX\ is good at typesetting mathematical formulas
like
       \( x-3y + z = 7 \)
or
       \( a_{1} > x^{2n} + y^{2n} > x' \)
or
       \( \ip{A}{B} = \sum_{i} a_{i} b_{i} \).
The spaces you type in a formula are
ignored.  Remember that a letter like
       $x$                   % $ ... $  and  \( ... \)  are equivalent
is a formula when it denotes a mathematical
symbol, and it should be typed as one.

\section{Displayed Text}

Text is displayed by indenting it from the left
margin.  Quotations are commonly displayed.  There
are short quotations
\begin{quote}
   This is a short quotation.  It consists of a
   single paragraph of text.  See how it is formatted.
\end{quote}
and longer ones.
\begin{quotation}
   This is a longer quotation.  It consists of two
   paragraphs of text, neither of which are
   particularly interesting.

   This is the second paragraph of the quotation.  It
   is just as dull as the first paragraph.
\end{quotation}
Another frequently-displayed structure is a list.
The following is an example of an \emph{itemized}
list.
\begin{itemize}
   \item This is the first item of an itemized list.
         Each item in the list is marked with a ``tick''.
         You don't have to worry about what kind of tick
         mark is used.

   \item This is the second item of the list.  It
         contains another list nested inside it.  The inner
         list is an \emph{enumerated} list.
         \begin{enumerate}
            \item This is the first item of an enumerated
                  list that is nested within the itemized list.

            \item This is the second item of the inner list.
                  \LaTeX\ allows you to nest lists deeper than
                  you really should.
         \end{enumerate}
         This is the rest of the second item of the outer
         list.  It is no more interesting than any other
         part of the item.
   \item This is the third item of the list.
\end{itemize}
You can even display poetry.
\begin{verse}
   There is an environment
    for verse \\             % The \\ command separates lines
   Whose features some poets % within a stanza.
   will curse.

                             % One or more blank lines separate stanzas.

   For instead of making\\
   Them do \emph{all} line breaking, \\
   It allows them to put too many words on a line when they'd rather be
   forced to be terse.
\end{verse}

Mathematical formulas may also be displayed.  A
displayed formula
is
one-line long; multiline
formulas require special formatting instructions.
   \[  \ip{\Gamma}{\psi'} = x'' + y^{2} + z_{i}^{n}\]
Don't start a paragraph with a displayed equation,
nor make one a paragraph by itself.

\end{document}               % End of document.

\section{Bench Huge Section 1183}
% synthetic bench section 1183
% This is a sample LaTeX input file.  (Version of 12 August 2004.)
%
% A '%' character causes TeX to ignore all remaining text on the line,
% and is used for comments like this one.

\documentclass{article}      % Specifies the document class

                             % The preamble begins here.
\title{An Example Document}  % Declares the document's title.
\author{Leslie Lamport}      % Declares the author's name.
\date{January 21, 1994}      % Deleting this command produces today's date.

\newcommand{\ip}[2]{(#1, #2)}
                             % Defines \ip{arg1}{arg2} to mean
                             % (arg1, arg2).

%\newcommand{\ip}[2]{\langle #1 | #2\rangle}
                             % This is an alternative definition of
                             % \ip that is commented out.

\begin{document}             % End of preamble and beginning of text.

\maketitle                   % Produces the title.

This is an example input file.  Comparing it with
the output it generates can show you how to
produce a simple document of your own.

\section{Ordinary Text}      % Produces section heading.  Lower-level
                             % sections are begun with similar
                             % \subsection and \subsubsection commands.

The ends  of words and sentences are marked
  by   spaces. It  doesn't matter how many
spaces    you type; one is as good as 100.  The
end of   a line counts as a space.

One   or more   blank lines denote the  end
of  a paragraph.

Since any number of consecutive spaces are treated
like a single one, the formatting of the input
file makes no difference to
      \LaTeX,                % The \LaTeX command generates the LaTeX logo.
but it makes a difference to you.  When you use
\LaTeX, making your input file as easy to read
as possible will be a great help as you write
your document and when you change it.  This sample
file shows how you can add comments to your own input
file.

Because printing is different from typewriting,
there are a number of things that you have to do
differently when preparing an input file than if
you were just typing the document directly.
Quotation marks like
       ``this''
have to be handled specially, as do quotes within
quotes:
       ``\,`this'            % \, separates the double and single quote.
        is what I just
        wrote, not  `that'\,''.

Dashes come in three sizes: an
       intra-word
dash, a medium dash for number ranges like
       1--2,
and a punctuation
       dash---like
this.

A sentence-ending space should be larger than the
space between words within a sentence.  You
sometimes have to type special commands in
conjunction with punctuation characters to get
this right, as in the following sentence.
       Gnats, gnus, etc.\ all  % `\ ' makes an inter-word space.
       begin with G\@.         % \@ marks end-of-sentence punctuation.
You should check the spaces after periods when
reading your output to make sure you haven't
forgotten any special cases.  Generating an
ellipsis
       \ldots\               % `\ ' is needed after `\ldots' because TeX
                             % ignores spaces after command names like \ldots
                             % made from \ + letters.
                             %
                             % Note how a `%' character causes TeX to ignore
                             % the end of the input line, so these blank lines
                             % do not start a new paragraph.
                             %
with the right spacing around the periods requires
a special command.

\LaTeX\ interprets some common characters as
commands, so you must type special commands to
generate them.  These characters include the
following:
       \$ \& \% \# \{ and \}.

In printing, text is usually emphasized with an
       \emph{italic}
type style.

\begin{em}
   A long segment of text can also be emphasized
   in this way.  Text within such a segment can be
   given \emph{additional} emphasis.
\end{em}

It is sometimes necessary to prevent \LaTeX\ from
breaking a line where it might otherwise do so.
This may be at a space, as between the ``Mr.''\ and
``Jones'' in
       ``Mr.~Jones'',        % ~ produces an unbreakable interword space.
or within a word---especially when the word is a
symbol like
       \mbox{\emph{itemnum}}
that makes little sense when hyphenated across
lines.

Footnotes\footnote{This is an example of a footnote.}
pose no problem.

\LaTeX\ is good at typesetting mathematical formulas
like
       \( x-3y + z = 7 \)
or
       \( a_{1} > x^{2n} + y^{2n} > x' \)
or
       \( \ip{A}{B} = \sum_{i} a_{i} b_{i} \).
The spaces you type in a formula are
ignored.  Remember that a letter like
       $x$                   % $ ... $  and  \( ... \)  are equivalent
is a formula when it denotes a mathematical
symbol, and it should be typed as one.

\section{Displayed Text}

Text is displayed by indenting it from the left
margin.  Quotations are commonly displayed.  There
are short quotations
\begin{quote}
   This is a short quotation.  It consists of a
   single paragraph of text.  See how it is formatted.
\end{quote}
and longer ones.
\begin{quotation}
   This is a longer quotation.  It consists of two
   paragraphs of text, neither of which are
   particularly interesting.

   This is the second paragraph of the quotation.  It
   is just as dull as the first paragraph.
\end{quotation}
Another frequently-displayed structure is a list.
The following is an example of an \emph{itemized}
list.
\begin{itemize}
   \item This is the first item of an itemized list.
         Each item in the list is marked with a ``tick''.
         You don't have to worry about what kind of tick
         mark is used.

   \item This is the second item of the list.  It
         contains another list nested inside it.  The inner
         list is an \emph{enumerated} list.
         \begin{enumerate}
            \item This is the first item of an enumerated
                  list that is nested within the itemized list.

            \item This is the second item of the inner list.
                  \LaTeX\ allows you to nest lists deeper than
                  you really should.
         \end{enumerate}
         This is the rest of the second item of the outer
         list.  It is no more interesting than any other
         part of the item.
   \item This is the third item of the list.
\end{itemize}
You can even display poetry.
\begin{verse}
   There is an environment
    for verse \\             % The \\ command separates lines
   Whose features some poets % within a stanza.
   will curse.

                             % One or more blank lines separate stanzas.

   For instead of making\\
   Them do \emph{all} line breaking, \\
   It allows them to put too many words on a line when they'd rather be
   forced to be terse.
\end{verse}

Mathematical formulas may also be displayed.  A
displayed formula
is
one-line long; multiline
formulas require special formatting instructions.
   \[  \ip{\Gamma}{\psi'} = x'' + y^{2} + z_{i}^{n}\]
Don't start a paragraph with a displayed equation,
nor make one a paragraph by itself.

\end{document}               % End of document.

\section{Bench Huge Section 1184}
% synthetic bench section 1184
% This is a sample LaTeX input file.  (Version of 12 August 2004.)
%
% A '%' character causes TeX to ignore all remaining text on the line,
% and is used for comments like this one.

\documentclass{article}      % Specifies the document class

                             % The preamble begins here.
\title{An Example Document}  % Declares the document's title.
\author{Leslie Lamport}      % Declares the author's name.
\date{January 21, 1994}      % Deleting this command produces today's date.

\newcommand{\ip}[2]{(#1, #2)}
                             % Defines \ip{arg1}{arg2} to mean
                             % (arg1, arg2).

%\newcommand{\ip}[2]{\langle #1 | #2\rangle}
                             % This is an alternative definition of
                             % \ip that is commented out.

\begin{document}             % End of preamble and beginning of text.

\maketitle                   % Produces the title.

This is an example input file.  Comparing it with
the output it generates can show you how to
produce a simple document of your own.

\section{Ordinary Text}      % Produces section heading.  Lower-level
                             % sections are begun with similar
                             % \subsection and \subsubsection commands.

The ends  of words and sentences are marked
  by   spaces. It  doesn't matter how many
spaces    you type; one is as good as 100.  The
end of   a line counts as a space.

One   or more   blank lines denote the  end
of  a paragraph.

Since any number of consecutive spaces are treated
like a single one, the formatting of the input
file makes no difference to
      \LaTeX,                % The \LaTeX command generates the LaTeX logo.
but it makes a difference to you.  When you use
\LaTeX, making your input file as easy to read
as possible will be a great help as you write
your document and when you change it.  This sample
file shows how you can add comments to your own input
file.

Because printing is different from typewriting,
there are a number of things that you have to do
differently when preparing an input file than if
you were just typing the document directly.
Quotation marks like
       ``this''
have to be handled specially, as do quotes within
quotes:
       ``\,`this'            % \, separates the double and single quote.
        is what I just
        wrote, not  `that'\,''.

Dashes come in three sizes: an
       intra-word
dash, a medium dash for number ranges like
       1--2,
and a punctuation
       dash---like
this.

A sentence-ending space should be larger than the
space between words within a sentence.  You
sometimes have to type special commands in
conjunction with punctuation characters to get
this right, as in the following sentence.
       Gnats, gnus, etc.\ all  % `\ ' makes an inter-word space.
       begin with G\@.         % \@ marks end-of-sentence punctuation.
You should check the spaces after periods when
reading your output to make sure you haven't
forgotten any special cases.  Generating an
ellipsis
       \ldots\               % `\ ' is needed after `\ldots' because TeX
                             % ignores spaces after command names like \ldots
                             % made from \ + letters.
                             %
                             % Note how a `%' character causes TeX to ignore
                             % the end of the input line, so these blank lines
                             % do not start a new paragraph.
                             %
with the right spacing around the periods requires
a special command.

\LaTeX\ interprets some common characters as
commands, so you must type special commands to
generate them.  These characters include the
following:
       \$ \& \% \# \{ and \}.

In printing, text is usually emphasized with an
       \emph{italic}
type style.

\begin{em}
   A long segment of text can also be emphasized
   in this way.  Text within such a segment can be
   given \emph{additional} emphasis.
\end{em}

It is sometimes necessary to prevent \LaTeX\ from
breaking a line where it might otherwise do so.
This may be at a space, as between the ``Mr.''\ and
``Jones'' in
       ``Mr.~Jones'',        % ~ produces an unbreakable interword space.
or within a word---especially when the word is a
symbol like
       \mbox{\emph{itemnum}}
that makes little sense when hyphenated across
lines.

Footnotes\footnote{This is an example of a footnote.}
pose no problem.

\LaTeX\ is good at typesetting mathematical formulas
like
       \( x-3y + z = 7 \)
or
       \( a_{1} > x^{2n} + y^{2n} > x' \)
or
       \( \ip{A}{B} = \sum_{i} a_{i} b_{i} \).
The spaces you type in a formula are
ignored.  Remember that a letter like
       $x$                   % $ ... $  and  \( ... \)  are equivalent
is a formula when it denotes a mathematical
symbol, and it should be typed as one.

\section{Displayed Text}

Text is displayed by indenting it from the left
margin.  Quotations are commonly displayed.  There
are short quotations
\begin{quote}
   This is a short quotation.  It consists of a
   single paragraph of text.  See how it is formatted.
\end{quote}
and longer ones.
\begin{quotation}
   This is a longer quotation.  It consists of two
   paragraphs of text, neither of which are
   particularly interesting.

   This is the second paragraph of the quotation.  It
   is just as dull as the first paragraph.
\end{quotation}
Another frequently-displayed structure is a list.
The following is an example of an \emph{itemized}
list.
\begin{itemize}
   \item This is the first item of an itemized list.
         Each item in the list is marked with a ``tick''.
         You don't have to worry about what kind of tick
         mark is used.

   \item This is the second item of the list.  It
         contains another list nested inside it.  The inner
         list is an \emph{enumerated} list.
         \begin{enumerate}
            \item This is the first item of an enumerated
                  list that is nested within the itemized list.

            \item This is the second item of the inner list.
                  \LaTeX\ allows you to nest lists deeper than
                  you really should.
         \end{enumerate}
         This is the rest of the second item of the outer
         list.  It is no more interesting than any other
         part of the item.
   \item This is the third item of the list.
\end{itemize}
You can even display poetry.
\begin{verse}
   There is an environment
    for verse \\             % The \\ command separates lines
   Whose features some poets % within a stanza.
   will curse.

                             % One or more blank lines separate stanzas.

   For instead of making\\
   Them do \emph{all} line breaking, \\
   It allows them to put too many words on a line when they'd rather be
   forced to be terse.
\end{verse}

Mathematical formulas may also be displayed.  A
displayed formula
is
one-line long; multiline
formulas require special formatting instructions.
   \[  \ip{\Gamma}{\psi'} = x'' + y^{2} + z_{i}^{n}\]
Don't start a paragraph with a displayed equation,
nor make one a paragraph by itself.

\end{document}               % End of document.

\section{Bench Huge Section 1185}
% synthetic bench section 1185
% This is a sample LaTeX input file.  (Version of 12 August 2004.)
%
% A '%' character causes TeX to ignore all remaining text on the line,
% and is used for comments like this one.

\documentclass{article}      % Specifies the document class

                             % The preamble begins here.
\title{An Example Document}  % Declares the document's title.
\author{Leslie Lamport}      % Declares the author's name.
\date{January 21, 1994}      % Deleting this command produces today's date.

\newcommand{\ip}[2]{(#1, #2)}
                             % Defines \ip{arg1}{arg2} to mean
                             % (arg1, arg2).

%\newcommand{\ip}[2]{\langle #1 | #2\rangle}
                             % This is an alternative definition of
                             % \ip that is commented out.

\begin{document}             % End of preamble and beginning of text.

\maketitle                   % Produces the title.

This is an example input file.  Comparing it with
the output it generates can show you how to
produce a simple document of your own.

\section{Ordinary Text}      % Produces section heading.  Lower-level
                             % sections are begun with similar
                             % \subsection and \subsubsection commands.

The ends  of words and sentences are marked
  by   spaces. It  doesn't matter how many
spaces    you type; one is as good as 100.  The
end of   a line counts as a space.

One   or more   blank lines denote the  end
of  a paragraph.

Since any number of consecutive spaces are treated
like a single one, the formatting of the input
file makes no difference to
      \LaTeX,                % The \LaTeX command generates the LaTeX logo.
but it makes a difference to you.  When you use
\LaTeX, making your input file as easy to read
as possible will be a great help as you write
your document and when you change it.  This sample
file shows how you can add comments to your own input
file.

Because printing is different from typewriting,
there are a number of things that you have to do
differently when preparing an input file than if
you were just typing the document directly.
Quotation marks like
       ``this''
have to be handled specially, as do quotes within
quotes:
       ``\,`this'            % \, separates the double and single quote.
        is what I just
        wrote, not  `that'\,''.

Dashes come in three sizes: an
       intra-word
dash, a medium dash for number ranges like
       1--2,
and a punctuation
       dash---like
this.

A sentence-ending space should be larger than the
space between words within a sentence.  You
sometimes have to type special commands in
conjunction with punctuation characters to get
this right, as in the following sentence.
       Gnats, gnus, etc.\ all  % `\ ' makes an inter-word space.
       begin with G\@.         % \@ marks end-of-sentence punctuation.
You should check the spaces after periods when
reading your output to make sure you haven't
forgotten any special cases.  Generating an
ellipsis
       \ldots\               % `\ ' is needed after `\ldots' because TeX
                             % ignores spaces after command names like \ldots
                             % made from \ + letters.
                             %
                             % Note how a `%' character causes TeX to ignore
                             % the end of the input line, so these blank lines
                             % do not start a new paragraph.
                             %
with the right spacing around the periods requires
a special command.

\LaTeX\ interprets some common characters as
commands, so you must type special commands to
generate them.  These characters include the
following:
       \$ \& \% \# \{ and \}.

In printing, text is usually emphasized with an
       \emph{italic}
type style.

\begin{em}
   A long segment of text can also be emphasized
   in this way.  Text within such a segment can be
   given \emph{additional} emphasis.
\end{em}

It is sometimes necessary to prevent \LaTeX\ from
breaking a line where it might otherwise do so.
This may be at a space, as between the ``Mr.''\ and
``Jones'' in
       ``Mr.~Jones'',        % ~ produces an unbreakable interword space.
or within a word---especially when the word is a
symbol like
       \mbox{\emph{itemnum}}
that makes little sense when hyphenated across
lines.

Footnotes\footnote{This is an example of a footnote.}
pose no problem.

\LaTeX\ is good at typesetting mathematical formulas
like
       \( x-3y + z = 7 \)
or
       \( a_{1} > x^{2n} + y^{2n} > x' \)
or
       \( \ip{A}{B} = \sum_{i} a_{i} b_{i} \).
The spaces you type in a formula are
ignored.  Remember that a letter like
       $x$                   % $ ... $  and  \( ... \)  are equivalent
is a formula when it denotes a mathematical
symbol, and it should be typed as one.

\section{Displayed Text}

Text is displayed by indenting it from the left
margin.  Quotations are commonly displayed.  There
are short quotations
\begin{quote}
   This is a short quotation.  It consists of a
   single paragraph of text.  See how it is formatted.
\end{quote}
and longer ones.
\begin{quotation}
   This is a longer quotation.  It consists of two
   paragraphs of text, neither of which are
   particularly interesting.

   This is the second paragraph of the quotation.  It
   is just as dull as the first paragraph.
\end{quotation}
Another frequently-displayed structure is a list.
The following is an example of an \emph{itemized}
list.
\begin{itemize}
   \item This is the first item of an itemized list.
         Each item in the list is marked with a ``tick''.
         You don't have to worry about what kind of tick
         mark is used.

   \item This is the second item of the list.  It
         contains another list nested inside it.  The inner
         list is an \emph{enumerated} list.
         \begin{enumerate}
            \item This is the first item of an enumerated
                  list that is nested within the itemized list.

            \item This is the second item of the inner list.
                  \LaTeX\ allows you to nest lists deeper than
                  you really should.
         \end{enumerate}
         This is the rest of the second item of the outer
         list.  It is no more interesting than any other
         part of the item.
   \item This is the third item of the list.
\end{itemize}
You can even display poetry.
\begin{verse}
   There is an environment
    for verse \\             % The \\ command separates lines
   Whose features some poets % within a stanza.
   will curse.

                             % One or more blank lines separate stanzas.

   For instead of making\\
   Them do \emph{all} line breaking, \\
   It allows them to put too many words on a line when they'd rather be
   forced to be terse.
\end{verse}

Mathematical formulas may also be displayed.  A
displayed formula
is
one-line long; multiline
formulas require special formatting instructions.
   \[  \ip{\Gamma}{\psi'} = x'' + y^{2} + z_{i}^{n}\]
Don't start a paragraph with a displayed equation,
nor make one a paragraph by itself.

\end{document}               % End of document.

\section{Bench Huge Section 1186}
% synthetic bench section 1186
% This is a sample LaTeX input file.  (Version of 12 August 2004.)
%
% A '%' character causes TeX to ignore all remaining text on the line,
% and is used for comments like this one.

\documentclass{article}      % Specifies the document class

                             % The preamble begins here.
\title{An Example Document}  % Declares the document's title.
\author{Leslie Lamport}      % Declares the author's name.
\date{January 21, 1994}      % Deleting this command produces today's date.

\newcommand{\ip}[2]{(#1, #2)}
                             % Defines \ip{arg1}{arg2} to mean
                             % (arg1, arg2).

%\newcommand{\ip}[2]{\langle #1 | #2\rangle}
                             % This is an alternative definition of
                             % \ip that is commented out.

\begin{document}             % End of preamble and beginning of text.

\maketitle                   % Produces the title.

This is an example input file.  Comparing it with
the output it generates can show you how to
produce a simple document of your own.

\section{Ordinary Text}      % Produces section heading.  Lower-level
                             % sections are begun with similar
                             % \subsection and \subsubsection commands.

The ends  of words and sentences are marked
  by   spaces. It  doesn't matter how many
spaces    you type; one is as good as 100.  The
end of   a line counts as a space.

One   or more   blank lines denote the  end
of  a paragraph.

Since any number of consecutive spaces are treated
like a single one, the formatting of the input
file makes no difference to
      \LaTeX,                % The \LaTeX command generates the LaTeX logo.
but it makes a difference to you.  When you use
\LaTeX, making your input file as easy to read
as possible will be a great help as you write
your document and when you change it.  This sample
file shows how you can add comments to your own input
file.

Because printing is different from typewriting,
there are a number of things that you have to do
differently when preparing an input file than if
you were just typing the document directly.
Quotation marks like
       ``this''
have to be handled specially, as do quotes within
quotes:
       ``\,`this'            % \, separates the double and single quote.
        is what I just
        wrote, not  `that'\,''.

Dashes come in three sizes: an
       intra-word
dash, a medium dash for number ranges like
       1--2,
and a punctuation
       dash---like
this.

A sentence-ending space should be larger than the
space between words within a sentence.  You
sometimes have to type special commands in
conjunction with punctuation characters to get
this right, as in the following sentence.
       Gnats, gnus, etc.\ all  % `\ ' makes an inter-word space.
       begin with G\@.         % \@ marks end-of-sentence punctuation.
You should check the spaces after periods when
reading your output to make sure you haven't
forgotten any special cases.  Generating an
ellipsis
       \ldots\               % `\ ' is needed after `\ldots' because TeX
                             % ignores spaces after command names like \ldots
                             % made from \ + letters.
                             %
                             % Note how a `%' character causes TeX to ignore
                             % the end of the input line, so these blank lines
                             % do not start a new paragraph.
                             %
with the right spacing around the periods requires
a special command.

\LaTeX\ interprets some common characters as
commands, so you must type special commands to
generate them.  These characters include the
following:
       \$ \& \% \# \{ and \}.

In printing, text is usually emphasized with an
       \emph{italic}
type style.

\begin{em}
   A long segment of text can also be emphasized
   in this way.  Text within such a segment can be
   given \emph{additional} emphasis.
\end{em}

It is sometimes necessary to prevent \LaTeX\ from
breaking a line where it might otherwise do so.
This may be at a space, as between the ``Mr.''\ and
``Jones'' in
       ``Mr.~Jones'',        % ~ produces an unbreakable interword space.
or within a word---especially when the word is a
symbol like
       \mbox{\emph{itemnum}}
that makes little sense when hyphenated across
lines.

Footnotes\footnote{This is an example of a footnote.}
pose no problem.

\LaTeX\ is good at typesetting mathematical formulas
like
       \( x-3y + z = 7 \)
or
       \( a_{1} > x^{2n} + y^{2n} > x' \)
or
       \( \ip{A}{B} = \sum_{i} a_{i} b_{i} \).
The spaces you type in a formula are
ignored.  Remember that a letter like
       $x$                   % $ ... $  and  \( ... \)  are equivalent
is a formula when it denotes a mathematical
symbol, and it should be typed as one.

\section{Displayed Text}

Text is displayed by indenting it from the left
margin.  Quotations are commonly displayed.  There
are short quotations
\begin{quote}
   This is a short quotation.  It consists of a
   single paragraph of text.  See how it is formatted.
\end{quote}
and longer ones.
\begin{quotation}
   This is a longer quotation.  It consists of two
   paragraphs of text, neither of which are
   particularly interesting.

   This is the second paragraph of the quotation.  It
   is just as dull as the first paragraph.
\end{quotation}
Another frequently-displayed structure is a list.
The following is an example of an \emph{itemized}
list.
\begin{itemize}
   \item This is the first item of an itemized list.
         Each item in the list is marked with a ``tick''.
         You don't have to worry about what kind of tick
         mark is used.

   \item This is the second item of the list.  It
         contains another list nested inside it.  The inner
         list is an \emph{enumerated} list.
         \begin{enumerate}
            \item This is the first item of an enumerated
                  list that is nested within the itemized list.

            \item This is the second item of the inner list.
                  \LaTeX\ allows you to nest lists deeper than
                  you really should.
         \end{enumerate}
         This is the rest of the second item of the outer
         list.  It is no more interesting than any other
         part of the item.
   \item This is the third item of the list.
\end{itemize}
You can even display poetry.
\begin{verse}
   There is an environment
    for verse \\             % The \\ command separates lines
   Whose features some poets % within a stanza.
   will curse.

                             % One or more blank lines separate stanzas.

   For instead of making\\
   Them do \emph{all} line breaking, \\
   It allows them to put too many words on a line when they'd rather be
   forced to be terse.
\end{verse}

Mathematical formulas may also be displayed.  A
displayed formula
is
one-line long; multiline
formulas require special formatting instructions.
   \[  \ip{\Gamma}{\psi'} = x'' + y^{2} + z_{i}^{n}\]
Don't start a paragraph with a displayed equation,
nor make one a paragraph by itself.

\end{document}               % End of document.

\section{Bench Huge Section 1187}
% synthetic bench section 1187
% This is a sample LaTeX input file.  (Version of 12 August 2004.)
%
% A '%' character causes TeX to ignore all remaining text on the line,
% and is used for comments like this one.

\documentclass{article}      % Specifies the document class

                             % The preamble begins here.
\title{An Example Document}  % Declares the document's title.
\author{Leslie Lamport}      % Declares the author's name.
\date{January 21, 1994}      % Deleting this command produces today's date.

\newcommand{\ip}[2]{(#1, #2)}
                             % Defines \ip{arg1}{arg2} to mean
                             % (arg1, arg2).

%\newcommand{\ip}[2]{\langle #1 | #2\rangle}
                             % This is an alternative definition of
                             % \ip that is commented out.

\begin{document}             % End of preamble and beginning of text.

\maketitle                   % Produces the title.

This is an example input file.  Comparing it with
the output it generates can show you how to
produce a simple document of your own.

\section{Ordinary Text}      % Produces section heading.  Lower-level
                             % sections are begun with similar
                             % \subsection and \subsubsection commands.

The ends  of words and sentences are marked
  by   spaces. It  doesn't matter how many
spaces    you type; one is as good as 100.  The
end of   a line counts as a space.

One   or more   blank lines denote the  end
of  a paragraph.

Since any number of consecutive spaces are treated
like a single one, the formatting of the input
file makes no difference to
      \LaTeX,                % The \LaTeX command generates the LaTeX logo.
but it makes a difference to you.  When you use
\LaTeX, making your input file as easy to read
as possible will be a great help as you write
your document and when you change it.  This sample
file shows how you can add comments to your own input
file.

Because printing is different from typewriting,
there are a number of things that you have to do
differently when preparing an input file than if
you were just typing the document directly.
Quotation marks like
       ``this''
have to be handled specially, as do quotes within
quotes:
       ``\,`this'            % \, separates the double and single quote.
        is what I just
        wrote, not  `that'\,''.

Dashes come in three sizes: an
       intra-word
dash, a medium dash for number ranges like
       1--2,
and a punctuation
       dash---like
this.

A sentence-ending space should be larger than the
space between words within a sentence.  You
sometimes have to type special commands in
conjunction with punctuation characters to get
this right, as in the following sentence.
       Gnats, gnus, etc.\ all  % `\ ' makes an inter-word space.
       begin with G\@.         % \@ marks end-of-sentence punctuation.
You should check the spaces after periods when
reading your output to make sure you haven't
forgotten any special cases.  Generating an
ellipsis
       \ldots\               % `\ ' is needed after `\ldots' because TeX
                             % ignores spaces after command names like \ldots
                             % made from \ + letters.
                             %
                             % Note how a `%' character causes TeX to ignore
                             % the end of the input line, so these blank lines
                             % do not start a new paragraph.
                             %
with the right spacing around the periods requires
a special command.

\LaTeX\ interprets some common characters as
commands, so you must type special commands to
generate them.  These characters include the
following:
       \$ \& \% \# \{ and \}.

In printing, text is usually emphasized with an
       \emph{italic}
type style.

\begin{em}
   A long segment of text can also be emphasized
   in this way.  Text within such a segment can be
   given \emph{additional} emphasis.
\end{em}

It is sometimes necessary to prevent \LaTeX\ from
breaking a line where it might otherwise do so.
This may be at a space, as between the ``Mr.''\ and
``Jones'' in
       ``Mr.~Jones'',        % ~ produces an unbreakable interword space.
or within a word---especially when the word is a
symbol like
       \mbox{\emph{itemnum}}
that makes little sense when hyphenated across
lines.

Footnotes\footnote{This is an example of a footnote.}
pose no problem.

\LaTeX\ is good at typesetting mathematical formulas
like
       \( x-3y + z = 7 \)
or
       \( a_{1} > x^{2n} + y^{2n} > x' \)
or
       \( \ip{A}{B} = \sum_{i} a_{i} b_{i} \).
The spaces you type in a formula are
ignored.  Remember that a letter like
       $x$                   % $ ... $  and  \( ... \)  are equivalent
is a formula when it denotes a mathematical
symbol, and it should be typed as one.

\section{Displayed Text}

Text is displayed by indenting it from the left
margin.  Quotations are commonly displayed.  There
are short quotations
\begin{quote}
   This is a short quotation.  It consists of a
   single paragraph of text.  See how it is formatted.
\end{quote}
and longer ones.
\begin{quotation}
   This is a longer quotation.  It consists of two
   paragraphs of text, neither of which are
   particularly interesting.

   This is the second paragraph of the quotation.  It
   is just as dull as the first paragraph.
\end{quotation}
Another frequently-displayed structure is a list.
The following is an example of an \emph{itemized}
list.
\begin{itemize}
   \item This is the first item of an itemized list.
         Each item in the list is marked with a ``tick''.
         You don't have to worry about what kind of tick
         mark is used.

   \item This is the second item of the list.  It
         contains another list nested inside it.  The inner
         list is an \emph{enumerated} list.
         \begin{enumerate}
            \item This is the first item of an enumerated
                  list that is nested within the itemized list.

            \item This is the second item of the inner list.
                  \LaTeX\ allows you to nest lists deeper than
                  you really should.
         \end{enumerate}
         This is the rest of the second item of the outer
         list.  It is no more interesting than any other
         part of the item.
   \item This is the third item of the list.
\end{itemize}
You can even display poetry.
\begin{verse}
   There is an environment
    for verse \\             % The \\ command separates lines
   Whose features some poets % within a stanza.
   will curse.

                             % One or more blank lines separate stanzas.

   For instead of making\\
   Them do \emph{all} line breaking, \\
   It allows them to put too many words on a line when they'd rather be
   forced to be terse.
\end{verse}

Mathematical formulas may also be displayed.  A
displayed formula
is
one-line long; multiline
formulas require special formatting instructions.
   \[  \ip{\Gamma}{\psi'} = x'' + y^{2} + z_{i}^{n}\]
Don't start a paragraph with a displayed equation,
nor make one a paragraph by itself.

\end{document}               % End of document.

\section{Bench Huge Section 1188}
% synthetic bench section 1188
% This is a sample LaTeX input file.  (Version of 12 August 2004.)
%
% A '%' character causes TeX to ignore all remaining text on the line,
% and is used for comments like this one.

\documentclass{article}      % Specifies the document class

                             % The preamble begins here.
\title{An Example Document}  % Declares the document's title.
\author{Leslie Lamport}      % Declares the author's name.
\date{January 21, 1994}      % Deleting this command produces today's date.

\newcommand{\ip}[2]{(#1, #2)}
                             % Defines \ip{arg1}{arg2} to mean
                             % (arg1, arg2).

%\newcommand{\ip}[2]{\langle #1 | #2\rangle}
                             % This is an alternative definition of
                             % \ip that is commented out.

\begin{document}             % End of preamble and beginning of text.

\maketitle                   % Produces the title.

This is an example input file.  Comparing it with
the output it generates can show you how to
produce a simple document of your own.

\section{Ordinary Text}      % Produces section heading.  Lower-level
                             % sections are begun with similar
                             % \subsection and \subsubsection commands.

The ends  of words and sentences are marked
  by   spaces. It  doesn't matter how many
spaces    you type; one is as good as 100.  The
end of   a line counts as a space.

One   or more   blank lines denote the  end
of  a paragraph.

Since any number of consecutive spaces are treated
like a single one, the formatting of the input
file makes no difference to
      \LaTeX,                % The \LaTeX command generates the LaTeX logo.
but it makes a difference to you.  When you use
\LaTeX, making your input file as easy to read
as possible will be a great help as you write
your document and when you change it.  This sample
file shows how you can add comments to your own input
file.

Because printing is different from typewriting,
there are a number of things that you have to do
differently when preparing an input file than if
you were just typing the document directly.
Quotation marks like
       ``this''
have to be handled specially, as do quotes within
quotes:
       ``\,`this'            % \, separates the double and single quote.
        is what I just
        wrote, not  `that'\,''.

Dashes come in three sizes: an
       intra-word
dash, a medium dash for number ranges like
       1--2,
and a punctuation
       dash---like
this.

A sentence-ending space should be larger than the
space between words within a sentence.  You
sometimes have to type special commands in
conjunction with punctuation characters to get
this right, as in the following sentence.
       Gnats, gnus, etc.\ all  % `\ ' makes an inter-word space.
       begin with G\@.         % \@ marks end-of-sentence punctuation.
You should check the spaces after periods when
reading your output to make sure you haven't
forgotten any special cases.  Generating an
ellipsis
       \ldots\               % `\ ' is needed after `\ldots' because TeX
                             % ignores spaces after command names like \ldots
                             % made from \ + letters.
                             %
                             % Note how a `%' character causes TeX to ignore
                             % the end of the input line, so these blank lines
                             % do not start a new paragraph.
                             %
with the right spacing around the periods requires
a special command.

\LaTeX\ interprets some common characters as
commands, so you must type special commands to
generate them.  These characters include the
following:
       \$ \& \% \# \{ and \}.

In printing, text is usually emphasized with an
       \emph{italic}
type style.

\begin{em}
   A long segment of text can also be emphasized
   in this way.  Text within such a segment can be
   given \emph{additional} emphasis.
\end{em}

It is sometimes necessary to prevent \LaTeX\ from
breaking a line where it might otherwise do so.
This may be at a space, as between the ``Mr.''\ and
``Jones'' in
       ``Mr.~Jones'',        % ~ produces an unbreakable interword space.
or within a word---especially when the word is a
symbol like
       \mbox{\emph{itemnum}}
that makes little sense when hyphenated across
lines.

Footnotes\footnote{This is an example of a footnote.}
pose no problem.

\LaTeX\ is good at typesetting mathematical formulas
like
       \( x-3y + z = 7 \)
or
       \( a_{1} > x^{2n} + y^{2n} > x' \)
or
       \( \ip{A}{B} = \sum_{i} a_{i} b_{i} \).
The spaces you type in a formula are
ignored.  Remember that a letter like
       $x$                   % $ ... $  and  \( ... \)  are equivalent
is a formula when it denotes a mathematical
symbol, and it should be typed as one.

\section{Displayed Text}

Text is displayed by indenting it from the left
margin.  Quotations are commonly displayed.  There
are short quotations
\begin{quote}
   This is a short quotation.  It consists of a
   single paragraph of text.  See how it is formatted.
\end{quote}
and longer ones.
\begin{quotation}
   This is a longer quotation.  It consists of two
   paragraphs of text, neither of which are
   particularly interesting.

   This is the second paragraph of the quotation.  It
   is just as dull as the first paragraph.
\end{quotation}
Another frequently-displayed structure is a list.
The following is an example of an \emph{itemized}
list.
\begin{itemize}
   \item This is the first item of an itemized list.
         Each item in the list is marked with a ``tick''.
         You don't have to worry about what kind of tick
         mark is used.

   \item This is the second item of the list.  It
         contains another list nested inside it.  The inner
         list is an \emph{enumerated} list.
         \begin{enumerate}
            \item This is the first item of an enumerated
                  list that is nested within the itemized list.

            \item This is the second item of the inner list.
                  \LaTeX\ allows you to nest lists deeper than
                  you really should.
         \end{enumerate}
         This is the rest of the second item of the outer
         list.  It is no more interesting than any other
         part of the item.
   \item This is the third item of the list.
\end{itemize}
You can even display poetry.
\begin{verse}
   There is an environment
    for verse \\             % The \\ command separates lines
   Whose features some poets % within a stanza.
   will curse.

                             % One or more blank lines separate stanzas.

   For instead of making\\
   Them do \emph{all} line breaking, \\
   It allows them to put too many words on a line when they'd rather be
   forced to be terse.
\end{verse}

Mathematical formulas may also be displayed.  A
displayed formula
is
one-line long; multiline
formulas require special formatting instructions.
   \[  \ip{\Gamma}{\psi'} = x'' + y^{2} + z_{i}^{n}\]
Don't start a paragraph with a displayed equation,
nor make one a paragraph by itself.

\end{document}               % End of document.

\section{Bench Huge Section 1189}
% synthetic bench section 1189
% This is a sample LaTeX input file.  (Version of 12 August 2004.)
%
% A '%' character causes TeX to ignore all remaining text on the line,
% and is used for comments like this one.

\documentclass{article}      % Specifies the document class

                             % The preamble begins here.
\title{An Example Document}  % Declares the document's title.
\author{Leslie Lamport}      % Declares the author's name.
\date{January 21, 1994}      % Deleting this command produces today's date.

\newcommand{\ip}[2]{(#1, #2)}
                             % Defines \ip{arg1}{arg2} to mean
                             % (arg1, arg2).

%\newcommand{\ip}[2]{\langle #1 | #2\rangle}
                             % This is an alternative definition of
                             % \ip that is commented out.

\begin{document}             % End of preamble and beginning of text.

\maketitle                   % Produces the title.

This is an example input file.  Comparing it with
the output it generates can show you how to
produce a simple document of your own.

\section{Ordinary Text}      % Produces section heading.  Lower-level
                             % sections are begun with similar
                             % \subsection and \subsubsection commands.

The ends  of words and sentences are marked
  by   spaces. It  doesn't matter how many
spaces    you type; one is as good as 100.  The
end of   a line counts as a space.

One   or more   blank lines denote the  end
of  a paragraph.

Since any number of consecutive spaces are treated
like a single one, the formatting of the input
file makes no difference to
      \LaTeX,                % The \LaTeX command generates the LaTeX logo.
but it makes a difference to you.  When you use
\LaTeX, making your input file as easy to read
as possible will be a great help as you write
your document and when you change it.  This sample
file shows how you can add comments to your own input
file.

Because printing is different from typewriting,
there are a number of things that you have to do
differently when preparing an input file than if
you were just typing the document directly.
Quotation marks like
       ``this''
have to be handled specially, as do quotes within
quotes:
       ``\,`this'            % \, separates the double and single quote.
        is what I just
        wrote, not  `that'\,''.

Dashes come in three sizes: an
       intra-word
dash, a medium dash for number ranges like
       1--2,
and a punctuation
       dash---like
this.

A sentence-ending space should be larger than the
space between words within a sentence.  You
sometimes have to type special commands in
conjunction with punctuation characters to get
this right, as in the following sentence.
       Gnats, gnus, etc.\ all  % `\ ' makes an inter-word space.
       begin with G\@.         % \@ marks end-of-sentence punctuation.
You should check the spaces after periods when
reading your output to make sure you haven't
forgotten any special cases.  Generating an
ellipsis
       \ldots\               % `\ ' is needed after `\ldots' because TeX
                             % ignores spaces after command names like \ldots
                             % made from \ + letters.
                             %
                             % Note how a `%' character causes TeX to ignore
                             % the end of the input line, so these blank lines
                             % do not start a new paragraph.
                             %
with the right spacing around the periods requires
a special command.

\LaTeX\ interprets some common characters as
commands, so you must type special commands to
generate them.  These characters include the
following:
       \$ \& \% \# \{ and \}.

In printing, text is usually emphasized with an
       \emph{italic}
type style.

\begin{em}
   A long segment of text can also be emphasized
   in this way.  Text within such a segment can be
   given \emph{additional} emphasis.
\end{em}

It is sometimes necessary to prevent \LaTeX\ from
breaking a line where it might otherwise do so.
This may be at a space, as between the ``Mr.''\ and
``Jones'' in
       ``Mr.~Jones'',        % ~ produces an unbreakable interword space.
or within a word---especially when the word is a
symbol like
       \mbox{\emph{itemnum}}
that makes little sense when hyphenated across
lines.

Footnotes\footnote{This is an example of a footnote.}
pose no problem.

\LaTeX\ is good at typesetting mathematical formulas
like
       \( x-3y + z = 7 \)
or
       \( a_{1} > x^{2n} + y^{2n} > x' \)
or
       \( \ip{A}{B} = \sum_{i} a_{i} b_{i} \).
The spaces you type in a formula are
ignored.  Remember that a letter like
       $x$                   % $ ... $  and  \( ... \)  are equivalent
is a formula when it denotes a mathematical
symbol, and it should be typed as one.

\section{Displayed Text}

Text is displayed by indenting it from the left
margin.  Quotations are commonly displayed.  There
are short quotations
\begin{quote}
   This is a short quotation.  It consists of a
   single paragraph of text.  See how it is formatted.
\end{quote}
and longer ones.
\begin{quotation}
   This is a longer quotation.  It consists of two
   paragraphs of text, neither of which are
   particularly interesting.

   This is the second paragraph of the quotation.  It
   is just as dull as the first paragraph.
\end{quotation}
Another frequently-displayed structure is a list.
The following is an example of an \emph{itemized}
list.
\begin{itemize}
   \item This is the first item of an itemized list.
         Each item in the list is marked with a ``tick''.
         You don't have to worry about what kind of tick
         mark is used.

   \item This is the second item of the list.  It
         contains another list nested inside it.  The inner
         list is an \emph{enumerated} list.
         \begin{enumerate}
            \item This is the first item of an enumerated
                  list that is nested within the itemized list.

            \item This is the second item of the inner list.
                  \LaTeX\ allows you to nest lists deeper than
                  you really should.
         \end{enumerate}
         This is the rest of the second item of the outer
         list.  It is no more interesting than any other
         part of the item.
   \item This is the third item of the list.
\end{itemize}
You can even display poetry.
\begin{verse}
   There is an environment
    for verse \\             % The \\ command separates lines
   Whose features some poets % within a stanza.
   will curse.

                             % One or more blank lines separate stanzas.

   For instead of making\\
   Them do \emph{all} line breaking, \\
   It allows them to put too many words on a line when they'd rather be
   forced to be terse.
\end{verse}

Mathematical formulas may also be displayed.  A
displayed formula
is
one-line long; multiline
formulas require special formatting instructions.
   \[  \ip{\Gamma}{\psi'} = x'' + y^{2} + z_{i}^{n}\]
Don't start a paragraph with a displayed equation,
nor make one a paragraph by itself.

\end{document}               % End of document.

\section{Bench Huge Section 1190}
% synthetic bench section 1190
% This is a sample LaTeX input file.  (Version of 12 August 2004.)
%
% A '%' character causes TeX to ignore all remaining text on the line,
% and is used for comments like this one.

\documentclass{article}      % Specifies the document class

                             % The preamble begins here.
\title{An Example Document}  % Declares the document's title.
\author{Leslie Lamport}      % Declares the author's name.
\date{January 21, 1994}      % Deleting this command produces today's date.

\newcommand{\ip}[2]{(#1, #2)}
                             % Defines \ip{arg1}{arg2} to mean
                             % (arg1, arg2).

%\newcommand{\ip}[2]{\langle #1 | #2\rangle}
                             % This is an alternative definition of
                             % \ip that is commented out.

\begin{document}             % End of preamble and beginning of text.

\maketitle                   % Produces the title.

This is an example input file.  Comparing it with
the output it generates can show you how to
produce a simple document of your own.

\section{Ordinary Text}      % Produces section heading.  Lower-level
                             % sections are begun with similar
                             % \subsection and \subsubsection commands.

The ends  of words and sentences are marked
  by   spaces. It  doesn't matter how many
spaces    you type; one is as good as 100.  The
end of   a line counts as a space.

One   or more   blank lines denote the  end
of  a paragraph.

Since any number of consecutive spaces are treated
like a single one, the formatting of the input
file makes no difference to
      \LaTeX,                % The \LaTeX command generates the LaTeX logo.
but it makes a difference to you.  When you use
\LaTeX, making your input file as easy to read
as possible will be a great help as you write
your document and when you change it.  This sample
file shows how you can add comments to your own input
file.

Because printing is different from typewriting,
there are a number of things that you have to do
differently when preparing an input file than if
you were just typing the document directly.
Quotation marks like
       ``this''
have to be handled specially, as do quotes within
quotes:
       ``\,`this'            % \, separates the double and single quote.
        is what I just
        wrote, not  `that'\,''.

Dashes come in three sizes: an
       intra-word
dash, a medium dash for number ranges like
       1--2,
and a punctuation
       dash---like
this.

A sentence-ending space should be larger than the
space between words within a sentence.  You
sometimes have to type special commands in
conjunction with punctuation characters to get
this right, as in the following sentence.
       Gnats, gnus, etc.\ all  % `\ ' makes an inter-word space.
       begin with G\@.         % \@ marks end-of-sentence punctuation.
You should check the spaces after periods when
reading your output to make sure you haven't
forgotten any special cases.  Generating an
ellipsis
       \ldots\               % `\ ' is needed after `\ldots' because TeX
                             % ignores spaces after command names like \ldots
                             % made from \ + letters.
                             %
                             % Note how a `%' character causes TeX to ignore
                             % the end of the input line, so these blank lines
                             % do not start a new paragraph.
                             %
with the right spacing around the periods requires
a special command.

\LaTeX\ interprets some common characters as
commands, so you must type special commands to
generate them.  These characters include the
following:
       \$ \& \% \# \{ and \}.

In printing, text is usually emphasized with an
       \emph{italic}
type style.

\begin{em}
   A long segment of text can also be emphasized
   in this way.  Text within such a segment can be
   given \emph{additional} emphasis.
\end{em}

It is sometimes necessary to prevent \LaTeX\ from
breaking a line where it might otherwise do so.
This may be at a space, as between the ``Mr.''\ and
``Jones'' in
       ``Mr.~Jones'',        % ~ produces an unbreakable interword space.
or within a word---especially when the word is a
symbol like
       \mbox{\emph{itemnum}}
that makes little sense when hyphenated across
lines.

Footnotes\footnote{This is an example of a footnote.}
pose no problem.

\LaTeX\ is good at typesetting mathematical formulas
like
       \( x-3y + z = 7 \)
or
       \( a_{1} > x^{2n} + y^{2n} > x' \)
or
       \( \ip{A}{B} = \sum_{i} a_{i} b_{i} \).
The spaces you type in a formula are
ignored.  Remember that a letter like
       $x$                   % $ ... $  and  \( ... \)  are equivalent
is a formula when it denotes a mathematical
symbol, and it should be typed as one.

\section{Displayed Text}

Text is displayed by indenting it from the left
margin.  Quotations are commonly displayed.  There
are short quotations
\begin{quote}
   This is a short quotation.  It consists of a
   single paragraph of text.  See how it is formatted.
\end{quote}
and longer ones.
\begin{quotation}
   This is a longer quotation.  It consists of two
   paragraphs of text, neither of which are
   particularly interesting.

   This is the second paragraph of the quotation.  It
   is just as dull as the first paragraph.
\end{quotation}
Another frequently-displayed structure is a list.
The following is an example of an \emph{itemized}
list.
\begin{itemize}
   \item This is the first item of an itemized list.
         Each item in the list is marked with a ``tick''.
         You don't have to worry about what kind of tick
         mark is used.

   \item This is the second item of the list.  It
         contains another list nested inside it.  The inner
         list is an \emph{enumerated} list.
         \begin{enumerate}
            \item This is the first item of an enumerated
                  list that is nested within the itemized list.

            \item This is the second item of the inner list.
                  \LaTeX\ allows you to nest lists deeper than
                  you really should.
         \end{enumerate}
         This is the rest of the second item of the outer
         list.  It is no more interesting than any other
         part of the item.
   \item This is the third item of the list.
\end{itemize}
You can even display poetry.
\begin{verse}
   There is an environment
    for verse \\             % The \\ command separates lines
   Whose features some poets % within a stanza.
   will curse.

                             % One or more blank lines separate stanzas.

   For instead of making\\
   Them do \emph{all} line breaking, \\
   It allows them to put too many words on a line when they'd rather be
   forced to be terse.
\end{verse}

Mathematical formulas may also be displayed.  A
displayed formula
is
one-line long; multiline
formulas require special formatting instructions.
   \[  \ip{\Gamma}{\psi'} = x'' + y^{2} + z_{i}^{n}\]
Don't start a paragraph with a displayed equation,
nor make one a paragraph by itself.

\end{document}               % End of document.

\section{Bench Huge Section 1191}
% synthetic bench section 1191
% This is a sample LaTeX input file.  (Version of 12 August 2004.)
%
% A '%' character causes TeX to ignore all remaining text on the line,
% and is used for comments like this one.

\documentclass{article}      % Specifies the document class

                             % The preamble begins here.
\title{An Example Document}  % Declares the document's title.
\author{Leslie Lamport}      % Declares the author's name.
\date{January 21, 1994}      % Deleting this command produces today's date.

\newcommand{\ip}[2]{(#1, #2)}
                             % Defines \ip{arg1}{arg2} to mean
                             % (arg1, arg2).

%\newcommand{\ip}[2]{\langle #1 | #2\rangle}
                             % This is an alternative definition of
                             % \ip that is commented out.

\begin{document}             % End of preamble and beginning of text.

\maketitle                   % Produces the title.

This is an example input file.  Comparing it with
the output it generates can show you how to
produce a simple document of your own.

\section{Ordinary Text}      % Produces section heading.  Lower-level
                             % sections are begun with similar
                             % \subsection and \subsubsection commands.

The ends  of words and sentences are marked
  by   spaces. It  doesn't matter how many
spaces    you type; one is as good as 100.  The
end of   a line counts as a space.

One   or more   blank lines denote the  end
of  a paragraph.

Since any number of consecutive spaces are treated
like a single one, the formatting of the input
file makes no difference to
      \LaTeX,                % The \LaTeX command generates the LaTeX logo.
but it makes a difference to you.  When you use
\LaTeX, making your input file as easy to read
as possible will be a great help as you write
your document and when you change it.  This sample
file shows how you can add comments to your own input
file.

Because printing is different from typewriting,
there are a number of things that you have to do
differently when preparing an input file than if
you were just typing the document directly.
Quotation marks like
       ``this''
have to be handled specially, as do quotes within
quotes:
       ``\,`this'            % \, separates the double and single quote.
        is what I just
        wrote, not  `that'\,''.

Dashes come in three sizes: an
       intra-word
dash, a medium dash for number ranges like
       1--2,
and a punctuation
       dash---like
this.

A sentence-ending space should be larger than the
space between words within a sentence.  You
sometimes have to type special commands in
conjunction with punctuation characters to get
this right, as in the following sentence.
       Gnats, gnus, etc.\ all  % `\ ' makes an inter-word space.
       begin with G\@.         % \@ marks end-of-sentence punctuation.
You should check the spaces after periods when
reading your output to make sure you haven't
forgotten any special cases.  Generating an
ellipsis
       \ldots\               % `\ ' is needed after `\ldots' because TeX
                             % ignores spaces after command names like \ldots
                             % made from \ + letters.
                             %
                             % Note how a `%' character causes TeX to ignore
                             % the end of the input line, so these blank lines
                             % do not start a new paragraph.
                             %
with the right spacing around the periods requires
a special command.

\LaTeX\ interprets some common characters as
commands, so you must type special commands to
generate them.  These characters include the
following:
       \$ \& \% \# \{ and \}.

In printing, text is usually emphasized with an
       \emph{italic}
type style.

\begin{em}
   A long segment of text can also be emphasized
   in this way.  Text within such a segment can be
   given \emph{additional} emphasis.
\end{em}

It is sometimes necessary to prevent \LaTeX\ from
breaking a line where it might otherwise do so.
This may be at a space, as between the ``Mr.''\ and
``Jones'' in
       ``Mr.~Jones'',        % ~ produces an unbreakable interword space.
or within a word---especially when the word is a
symbol like
       \mbox{\emph{itemnum}}
that makes little sense when hyphenated across
lines.

Footnotes\footnote{This is an example of a footnote.}
pose no problem.

\LaTeX\ is good at typesetting mathematical formulas
like
       \( x-3y + z = 7 \)
or
       \( a_{1} > x^{2n} + y^{2n} > x' \)
or
       \( \ip{A}{B} = \sum_{i} a_{i} b_{i} \).
The spaces you type in a formula are
ignored.  Remember that a letter like
       $x$                   % $ ... $  and  \( ... \)  are equivalent
is a formula when it denotes a mathematical
symbol, and it should be typed as one.

\section{Displayed Text}

Text is displayed by indenting it from the left
margin.  Quotations are commonly displayed.  There
are short quotations
\begin{quote}
   This is a short quotation.  It consists of a
   single paragraph of text.  See how it is formatted.
\end{quote}
and longer ones.
\begin{quotation}
   This is a longer quotation.  It consists of two
   paragraphs of text, neither of which are
   particularly interesting.

   This is the second paragraph of the quotation.  It
   is just as dull as the first paragraph.
\end{quotation}
Another frequently-displayed structure is a list.
The following is an example of an \emph{itemized}
list.
\begin{itemize}
   \item This is the first item of an itemized list.
         Each item in the list is marked with a ``tick''.
         You don't have to worry about what kind of tick
         mark is used.

   \item This is the second item of the list.  It
         contains another list nested inside it.  The inner
         list is an \emph{enumerated} list.
         \begin{enumerate}
            \item This is the first item of an enumerated
                  list that is nested within the itemized list.

            \item This is the second item of the inner list.
                  \LaTeX\ allows you to nest lists deeper than
                  you really should.
         \end{enumerate}
         This is the rest of the second item of the outer
         list.  It is no more interesting than any other
         part of the item.
   \item This is the third item of the list.
\end{itemize}
You can even display poetry.
\begin{verse}
   There is an environment
    for verse \\             % The \\ command separates lines
   Whose features some poets % within a stanza.
   will curse.

                             % One or more blank lines separate stanzas.

   For instead of making\\
   Them do \emph{all} line breaking, \\
   It allows them to put too many words on a line when they'd rather be
   forced to be terse.
\end{verse}

Mathematical formulas may also be displayed.  A
displayed formula
is
one-line long; multiline
formulas require special formatting instructions.
   \[  \ip{\Gamma}{\psi'} = x'' + y^{2} + z_{i}^{n}\]
Don't start a paragraph with a displayed equation,
nor make one a paragraph by itself.

\end{document}               % End of document.

\section{Bench Huge Section 1192}
% synthetic bench section 1192
% This is a sample LaTeX input file.  (Version of 12 August 2004.)
%
% A '%' character causes TeX to ignore all remaining text on the line,
% and is used for comments like this one.

\documentclass{article}      % Specifies the document class

                             % The preamble begins here.
\title{An Example Document}  % Declares the document's title.
\author{Leslie Lamport}      % Declares the author's name.
\date{January 21, 1994}      % Deleting this command produces today's date.

\newcommand{\ip}[2]{(#1, #2)}
                             % Defines \ip{arg1}{arg2} to mean
                             % (arg1, arg2).

%\newcommand{\ip}[2]{\langle #1 | #2\rangle}
                             % This is an alternative definition of
                             % \ip that is commented out.

\begin{document}             % End of preamble and beginning of text.

\maketitle                   % Produces the title.

This is an example input file.  Comparing it with
the output it generates can show you how to
produce a simple document of your own.

\section{Ordinary Text}      % Produces section heading.  Lower-level
                             % sections are begun with similar
                             % \subsection and \subsubsection commands.

The ends  of words and sentences are marked
  by   spaces. It  doesn't matter how many
spaces    you type; one is as good as 100.  The
end of   a line counts as a space.

One   or more   blank lines denote the  end
of  a paragraph.

Since any number of consecutive spaces are treated
like a single one, the formatting of the input
file makes no difference to
      \LaTeX,                % The \LaTeX command generates the LaTeX logo.
but it makes a difference to you.  When you use
\LaTeX, making your input file as easy to read
as possible will be a great help as you write
your document and when you change it.  This sample
file shows how you can add comments to your own input
file.

Because printing is different from typewriting,
there are a number of things that you have to do
differently when preparing an input file than if
you were just typing the document directly.
Quotation marks like
       ``this''
have to be handled specially, as do quotes within
quotes:
       ``\,`this'            % \, separates the double and single quote.
        is what I just
        wrote, not  `that'\,''.

Dashes come in three sizes: an
       intra-word
dash, a medium dash for number ranges like
       1--2,
and a punctuation
       dash---like
this.

A sentence-ending space should be larger than the
space between words within a sentence.  You
sometimes have to type special commands in
conjunction with punctuation characters to get
this right, as in the following sentence.
       Gnats, gnus, etc.\ all  % `\ ' makes an inter-word space.
       begin with G\@.         % \@ marks end-of-sentence punctuation.
You should check the spaces after periods when
reading your output to make sure you haven't
forgotten any special cases.  Generating an
ellipsis
       \ldots\               % `\ ' is needed after `\ldots' because TeX
                             % ignores spaces after command names like \ldots
                             % made from \ + letters.
                             %
                             % Note how a `%' character causes TeX to ignore
                             % the end of the input line, so these blank lines
                             % do not start a new paragraph.
                             %
with the right spacing around the periods requires
a special command.

\LaTeX\ interprets some common characters as
commands, so you must type special commands to
generate them.  These characters include the
following:
       \$ \& \% \# \{ and \}.

In printing, text is usually emphasized with an
       \emph{italic}
type style.

\begin{em}
   A long segment of text can also be emphasized
   in this way.  Text within such a segment can be
   given \emph{additional} emphasis.
\end{em}

It is sometimes necessary to prevent \LaTeX\ from
breaking a line where it might otherwise do so.
This may be at a space, as between the ``Mr.''\ and
``Jones'' in
       ``Mr.~Jones'',        % ~ produces an unbreakable interword space.
or within a word---especially when the word is a
symbol like
       \mbox{\emph{itemnum}}
that makes little sense when hyphenated across
lines.

Footnotes\footnote{This is an example of a footnote.}
pose no problem.

\LaTeX\ is good at typesetting mathematical formulas
like
       \( x-3y + z = 7 \)
or
       \( a_{1} > x^{2n} + y^{2n} > x' \)
or
       \( \ip{A}{B} = \sum_{i} a_{i} b_{i} \).
The spaces you type in a formula are
ignored.  Remember that a letter like
       $x$                   % $ ... $  and  \( ... \)  are equivalent
is a formula when it denotes a mathematical
symbol, and it should be typed as one.

\section{Displayed Text}

Text is displayed by indenting it from the left
margin.  Quotations are commonly displayed.  There
are short quotations
\begin{quote}
   This is a short quotation.  It consists of a
   single paragraph of text.  See how it is formatted.
\end{quote}
and longer ones.
\begin{quotation}
   This is a longer quotation.  It consists of two
   paragraphs of text, neither of which are
   particularly interesting.

   This is the second paragraph of the quotation.  It
   is just as dull as the first paragraph.
\end{quotation}
Another frequently-displayed structure is a list.
The following is an example of an \emph{itemized}
list.
\begin{itemize}
   \item This is the first item of an itemized list.
         Each item in the list is marked with a ``tick''.
         You don't have to worry about what kind of tick
         mark is used.

   \item This is the second item of the list.  It
         contains another list nested inside it.  The inner
         list is an \emph{enumerated} list.
         \begin{enumerate}
            \item This is the first item of an enumerated
                  list that is nested within the itemized list.

            \item This is the second item of the inner list.
                  \LaTeX\ allows you to nest lists deeper than
                  you really should.
         \end{enumerate}
         This is the rest of the second item of the outer
         list.  It is no more interesting than any other
         part of the item.
   \item This is the third item of the list.
\end{itemize}
You can even display poetry.
\begin{verse}
   There is an environment
    for verse \\             % The \\ command separates lines
   Whose features some poets % within a stanza.
   will curse.

                             % One or more blank lines separate stanzas.

   For instead of making\\
   Them do \emph{all} line breaking, \\
   It allows them to put too many words on a line when they'd rather be
   forced to be terse.
\end{verse}

Mathematical formulas may also be displayed.  A
displayed formula
is
one-line long; multiline
formulas require special formatting instructions.
   \[  \ip{\Gamma}{\psi'} = x'' + y^{2} + z_{i}^{n}\]
Don't start a paragraph with a displayed equation,
nor make one a paragraph by itself.

\end{document}               % End of document.

\section{Bench Huge Section 1193}
% synthetic bench section 1193
% This is a sample LaTeX input file.  (Version of 12 August 2004.)
%
% A '%' character causes TeX to ignore all remaining text on the line,
% and is used for comments like this one.

\documentclass{article}      % Specifies the document class

                             % The preamble begins here.
\title{An Example Document}  % Declares the document's title.
\author{Leslie Lamport}      % Declares the author's name.
\date{January 21, 1994}      % Deleting this command produces today's date.

\newcommand{\ip}[2]{(#1, #2)}
                             % Defines \ip{arg1}{arg2} to mean
                             % (arg1, arg2).

%\newcommand{\ip}[2]{\langle #1 | #2\rangle}
                             % This is an alternative definition of
                             % \ip that is commented out.

\begin{document}             % End of preamble and beginning of text.

\maketitle                   % Produces the title.

This is an example input file.  Comparing it with
the output it generates can show you how to
produce a simple document of your own.

\section{Ordinary Text}      % Produces section heading.  Lower-level
                             % sections are begun with similar
                             % \subsection and \subsubsection commands.

The ends  of words and sentences are marked
  by   spaces. It  doesn't matter how many
spaces    you type; one is as good as 100.  The
end of   a line counts as a space.

One   or more   blank lines denote the  end
of  a paragraph.

Since any number of consecutive spaces are treated
like a single one, the formatting of the input
file makes no difference to
      \LaTeX,                % The \LaTeX command generates the LaTeX logo.
but it makes a difference to you.  When you use
\LaTeX, making your input file as easy to read
as possible will be a great help as you write
your document and when you change it.  This sample
file shows how you can add comments to your own input
file.

Because printing is different from typewriting,
there are a number of things that you have to do
differently when preparing an input file than if
you were just typing the document directly.
Quotation marks like
       ``this''
have to be handled specially, as do quotes within
quotes:
       ``\,`this'            % \, separates the double and single quote.
        is what I just
        wrote, not  `that'\,''.

Dashes come in three sizes: an
       intra-word
dash, a medium dash for number ranges like
       1--2,
and a punctuation
       dash---like
this.

A sentence-ending space should be larger than the
space between words within a sentence.  You
sometimes have to type special commands in
conjunction with punctuation characters to get
this right, as in the following sentence.
       Gnats, gnus, etc.\ all  % `\ ' makes an inter-word space.
       begin with G\@.         % \@ marks end-of-sentence punctuation.
You should check the spaces after periods when
reading your output to make sure you haven't
forgotten any special cases.  Generating an
ellipsis
       \ldots\               % `\ ' is needed after `\ldots' because TeX
                             % ignores spaces after command names like \ldots
                             % made from \ + letters.
                             %
                             % Note how a `%' character causes TeX to ignore
                             % the end of the input line, so these blank lines
                             % do not start a new paragraph.
                             %
with the right spacing around the periods requires
a special command.

\LaTeX\ interprets some common characters as
commands, so you must type special commands to
generate them.  These characters include the
following:
       \$ \& \% \# \{ and \}.

In printing, text is usually emphasized with an
       \emph{italic}
type style.

\begin{em}
   A long segment of text can also be emphasized
   in this way.  Text within such a segment can be
   given \emph{additional} emphasis.
\end{em}

It is sometimes necessary to prevent \LaTeX\ from
breaking a line where it might otherwise do so.
This may be at a space, as between the ``Mr.''\ and
``Jones'' in
       ``Mr.~Jones'',        % ~ produces an unbreakable interword space.
or within a word---especially when the word is a
symbol like
       \mbox{\emph{itemnum}}
that makes little sense when hyphenated across
lines.

Footnotes\footnote{This is an example of a footnote.}
pose no problem.

\LaTeX\ is good at typesetting mathematical formulas
like
       \( x-3y + z = 7 \)
or
       \( a_{1} > x^{2n} + y^{2n} > x' \)
or
       \( \ip{A}{B} = \sum_{i} a_{i} b_{i} \).
The spaces you type in a formula are
ignored.  Remember that a letter like
       $x$                   % $ ... $  and  \( ... \)  are equivalent
is a formula when it denotes a mathematical
symbol, and it should be typed as one.

\section{Displayed Text}

Text is displayed by indenting it from the left
margin.  Quotations are commonly displayed.  There
are short quotations
\begin{quote}
   This is a short quotation.  It consists of a
   single paragraph of text.  See how it is formatted.
\end{quote}
and longer ones.
\begin{quotation}
   This is a longer quotation.  It consists of two
   paragraphs of text, neither of which are
   particularly interesting.

   This is the second paragraph of the quotation.  It
   is just as dull as the first paragraph.
\end{quotation}
Another frequently-displayed structure is a list.
The following is an example of an \emph{itemized}
list.
\begin{itemize}
   \item This is the first item of an itemized list.
         Each item in the list is marked with a ``tick''.
         You don't have to worry about what kind of tick
         mark is used.

   \item This is the second item of the list.  It
         contains another list nested inside it.  The inner
         list is an \emph{enumerated} list.
         \begin{enumerate}
            \item This is the first item of an enumerated
                  list that is nested within the itemized list.

            \item This is the second item of the inner list.
                  \LaTeX\ allows you to nest lists deeper than
                  you really should.
         \end{enumerate}
         This is the rest of the second item of the outer
         list.  It is no more interesting than any other
         part of the item.
   \item This is the third item of the list.
\end{itemize}
You can even display poetry.
\begin{verse}
   There is an environment
    for verse \\             % The \\ command separates lines
   Whose features some poets % within a stanza.
   will curse.

                             % One or more blank lines separate stanzas.

   For instead of making\\
   Them do \emph{all} line breaking, \\
   It allows them to put too many words on a line when they'd rather be
   forced to be terse.
\end{verse}

Mathematical formulas may also be displayed.  A
displayed formula
is
one-line long; multiline
formulas require special formatting instructions.
   \[  \ip{\Gamma}{\psi'} = x'' + y^{2} + z_{i}^{n}\]
Don't start a paragraph with a displayed equation,
nor make one a paragraph by itself.

\end{document}               % End of document.

\section{Bench Huge Section 1194}
% synthetic bench section 1194
% This is a sample LaTeX input file.  (Version of 12 August 2004.)
%
% A '%' character causes TeX to ignore all remaining text on the line,
% and is used for comments like this one.

\documentclass{article}      % Specifies the document class

                             % The preamble begins here.
\title{An Example Document}  % Declares the document's title.
\author{Leslie Lamport}      % Declares the author's name.
\date{January 21, 1994}      % Deleting this command produces today's date.

\newcommand{\ip}[2]{(#1, #2)}
                             % Defines \ip{arg1}{arg2} to mean
                             % (arg1, arg2).

%\newcommand{\ip}[2]{\langle #1 | #2\rangle}
                             % This is an alternative definition of
                             % \ip that is commented out.

\begin{document}             % End of preamble and beginning of text.

\maketitle                   % Produces the title.

This is an example input file.  Comparing it with
the output it generates can show you how to
produce a simple document of your own.

\section{Ordinary Text}      % Produces section heading.  Lower-level
                             % sections are begun with similar
                             % \subsection and \subsubsection commands.

The ends  of words and sentences are marked
  by   spaces. It  doesn't matter how many
spaces    you type; one is as good as 100.  The
end of   a line counts as a space.

One   or more   blank lines denote the  end
of  a paragraph.

Since any number of consecutive spaces are treated
like a single one, the formatting of the input
file makes no difference to
      \LaTeX,                % The \LaTeX command generates the LaTeX logo.
but it makes a difference to you.  When you use
\LaTeX, making your input file as easy to read
as possible will be a great help as you write
your document and when you change it.  This sample
file shows how you can add comments to your own input
file.

Because printing is different from typewriting,
there are a number of things that you have to do
differently when preparing an input file than if
you were just typing the document directly.
Quotation marks like
       ``this''
have to be handled specially, as do quotes within
quotes:
       ``\,`this'            % \, separates the double and single quote.
        is what I just
        wrote, not  `that'\,''.

Dashes come in three sizes: an
       intra-word
dash, a medium dash for number ranges like
       1--2,
and a punctuation
       dash---like
this.

A sentence-ending space should be larger than the
space between words within a sentence.  You
sometimes have to type special commands in
conjunction with punctuation characters to get
this right, as in the following sentence.
       Gnats, gnus, etc.\ all  % `\ ' makes an inter-word space.
       begin with G\@.         % \@ marks end-of-sentence punctuation.
You should check the spaces after periods when
reading your output to make sure you haven't
forgotten any special cases.  Generating an
ellipsis
       \ldots\               % `\ ' is needed after `\ldots' because TeX
                             % ignores spaces after command names like \ldots
                             % made from \ + letters.
                             %
                             % Note how a `%' character causes TeX to ignore
                             % the end of the input line, so these blank lines
                             % do not start a new paragraph.
                             %
with the right spacing around the periods requires
a special command.

\LaTeX\ interprets some common characters as
commands, so you must type special commands to
generate them.  These characters include the
following:
       \$ \& \% \# \{ and \}.

In printing, text is usually emphasized with an
       \emph{italic}
type style.

\begin{em}
   A long segment of text can also be emphasized
   in this way.  Text within such a segment can be
   given \emph{additional} emphasis.
\end{em}

It is sometimes necessary to prevent \LaTeX\ from
breaking a line where it might otherwise do so.
This may be at a space, as between the ``Mr.''\ and
``Jones'' in
       ``Mr.~Jones'',        % ~ produces an unbreakable interword space.
or within a word---especially when the word is a
symbol like
       \mbox{\emph{itemnum}}
that makes little sense when hyphenated across
lines.

Footnotes\footnote{This is an example of a footnote.}
pose no problem.

\LaTeX\ is good at typesetting mathematical formulas
like
       \( x-3y + z = 7 \)
or
       \( a_{1} > x^{2n} + y^{2n} > x' \)
or
       \( \ip{A}{B} = \sum_{i} a_{i} b_{i} \).
The spaces you type in a formula are
ignored.  Remember that a letter like
       $x$                   % $ ... $  and  \( ... \)  are equivalent
is a formula when it denotes a mathematical
symbol, and it should be typed as one.

\section{Displayed Text}

Text is displayed by indenting it from the left
margin.  Quotations are commonly displayed.  There
are short quotations
\begin{quote}
   This is a short quotation.  It consists of a
   single paragraph of text.  See how it is formatted.
\end{quote}
and longer ones.
\begin{quotation}
   This is a longer quotation.  It consists of two
   paragraphs of text, neither of which are
   particularly interesting.

   This is the second paragraph of the quotation.  It
   is just as dull as the first paragraph.
\end{quotation}
Another frequently-displayed structure is a list.
The following is an example of an \emph{itemized}
list.
\begin{itemize}
   \item This is the first item of an itemized list.
         Each item in the list is marked with a ``tick''.
         You don't have to worry about what kind of tick
         mark is used.

   \item This is the second item of the list.  It
         contains another list nested inside it.  The inner
         list is an \emph{enumerated} list.
         \begin{enumerate}
            \item This is the first item of an enumerated
                  list that is nested within the itemized list.

            \item This is the second item of the inner list.
                  \LaTeX\ allows you to nest lists deeper than
                  you really should.
         \end{enumerate}
         This is the rest of the second item of the outer
         list.  It is no more interesting than any other
         part of the item.
   \item This is the third item of the list.
\end{itemize}
You can even display poetry.
\begin{verse}
   There is an environment
    for verse \\             % The \\ command separates lines
   Whose features some poets % within a stanza.
   will curse.

                             % One or more blank lines separate stanzas.

   For instead of making\\
   Them do \emph{all} line breaking, \\
   It allows them to put too many words on a line when they'd rather be
   forced to be terse.
\end{verse}

Mathematical formulas may also be displayed.  A
displayed formula
is
one-line long; multiline
formulas require special formatting instructions.
   \[  \ip{\Gamma}{\psi'} = x'' + y^{2} + z_{i}^{n}\]
Don't start a paragraph with a displayed equation,
nor make one a paragraph by itself.

\end{document}               % End of document.

\section{Bench Huge Section 1195}
% synthetic bench section 1195
% This is a sample LaTeX input file.  (Version of 12 August 2004.)
%
% A '%' character causes TeX to ignore all remaining text on the line,
% and is used for comments like this one.

\documentclass{article}      % Specifies the document class

                             % The preamble begins here.
\title{An Example Document}  % Declares the document's title.
\author{Leslie Lamport}      % Declares the author's name.
\date{January 21, 1994}      % Deleting this command produces today's date.

\newcommand{\ip}[2]{(#1, #2)}
                             % Defines \ip{arg1}{arg2} to mean
                             % (arg1, arg2).

%\newcommand{\ip}[2]{\langle #1 | #2\rangle}
                             % This is an alternative definition of
                             % \ip that is commented out.

\begin{document}             % End of preamble and beginning of text.

\maketitle                   % Produces the title.

This is an example input file.  Comparing it with
the output it generates can show you how to
produce a simple document of your own.

\section{Ordinary Text}      % Produces section heading.  Lower-level
                             % sections are begun with similar
                             % \subsection and \subsubsection commands.

The ends  of words and sentences are marked
  by   spaces. It  doesn't matter how many
spaces    you type; one is as good as 100.  The
end of   a line counts as a space.

One   or more   blank lines denote the  end
of  a paragraph.

Since any number of consecutive spaces are treated
like a single one, the formatting of the input
file makes no difference to
      \LaTeX,                % The \LaTeX command generates the LaTeX logo.
but it makes a difference to you.  When you use
\LaTeX, making your input file as easy to read
as possible will be a great help as you write
your document and when you change it.  This sample
file shows how you can add comments to your own input
file.

Because printing is different from typewriting,
there are a number of things that you have to do
differently when preparing an input file than if
you were just typing the document directly.
Quotation marks like
       ``this''
have to be handled specially, as do quotes within
quotes:
       ``\,`this'            % \, separates the double and single quote.
        is what I just
        wrote, not  `that'\,''.

Dashes come in three sizes: an
       intra-word
dash, a medium dash for number ranges like
       1--2,
and a punctuation
       dash---like
this.

A sentence-ending space should be larger than the
space between words within a sentence.  You
sometimes have to type special commands in
conjunction with punctuation characters to get
this right, as in the following sentence.
       Gnats, gnus, etc.\ all  % `\ ' makes an inter-word space.
       begin with G\@.         % \@ marks end-of-sentence punctuation.
You should check the spaces after periods when
reading your output to make sure you haven't
forgotten any special cases.  Generating an
ellipsis
       \ldots\               % `\ ' is needed after `\ldots' because TeX
                             % ignores spaces after command names like \ldots
                             % made from \ + letters.
                             %
                             % Note how a `%' character causes TeX to ignore
                             % the end of the input line, so these blank lines
                             % do not start a new paragraph.
                             %
with the right spacing around the periods requires
a special command.

\LaTeX\ interprets some common characters as
commands, so you must type special commands to
generate them.  These characters include the
following:
       \$ \& \% \# \{ and \}.

In printing, text is usually emphasized with an
       \emph{italic}
type style.

\begin{em}
   A long segment of text can also be emphasized
   in this way.  Text within such a segment can be
   given \emph{additional} emphasis.
\end{em}

It is sometimes necessary to prevent \LaTeX\ from
breaking a line where it might otherwise do so.
This may be at a space, as between the ``Mr.''\ and
``Jones'' in
       ``Mr.~Jones'',        % ~ produces an unbreakable interword space.
or within a word---especially when the word is a
symbol like
       \mbox{\emph{itemnum}}
that makes little sense when hyphenated across
lines.

Footnotes\footnote{This is an example of a footnote.}
pose no problem.

\LaTeX\ is good at typesetting mathematical formulas
like
       \( x-3y + z = 7 \)
or
       \( a_{1} > x^{2n} + y^{2n} > x' \)
or
       \( \ip{A}{B} = \sum_{i} a_{i} b_{i} \).
The spaces you type in a formula are
ignored.  Remember that a letter like
       $x$                   % $ ... $  and  \( ... \)  are equivalent
is a formula when it denotes a mathematical
symbol, and it should be typed as one.

\section{Displayed Text}

Text is displayed by indenting it from the left
margin.  Quotations are commonly displayed.  There
are short quotations
\begin{quote}
   This is a short quotation.  It consists of a
   single paragraph of text.  See how it is formatted.
\end{quote}
and longer ones.
\begin{quotation}
   This is a longer quotation.  It consists of two
   paragraphs of text, neither of which are
   particularly interesting.

   This is the second paragraph of the quotation.  It
   is just as dull as the first paragraph.
\end{quotation}
Another frequently-displayed structure is a list.
The following is an example of an \emph{itemized}
list.
\begin{itemize}
   \item This is the first item of an itemized list.
         Each item in the list is marked with a ``tick''.
         You don't have to worry about what kind of tick
         mark is used.

   \item This is the second item of the list.  It
         contains another list nested inside it.  The inner
         list is an \emph{enumerated} list.
         \begin{enumerate}
            \item This is the first item of an enumerated
                  list that is nested within the itemized list.

            \item This is the second item of the inner list.
                  \LaTeX\ allows you to nest lists deeper than
                  you really should.
         \end{enumerate}
         This is the rest of the second item of the outer
         list.  It is no more interesting than any other
         part of the item.
   \item This is the third item of the list.
\end{itemize}
You can even display poetry.
\begin{verse}
   There is an environment
    for verse \\             % The \\ command separates lines
   Whose features some poets % within a stanza.
   will curse.

                             % One or more blank lines separate stanzas.

   For instead of making\\
   Them do \emph{all} line breaking, \\
   It allows them to put too many words on a line when they'd rather be
   forced to be terse.
\end{verse}

Mathematical formulas may also be displayed.  A
displayed formula
is
one-line long; multiline
formulas require special formatting instructions.
   \[  \ip{\Gamma}{\psi'} = x'' + y^{2} + z_{i}^{n}\]
Don't start a paragraph with a displayed equation,
nor make one a paragraph by itself.

\end{document}               % End of document.

\section{Bench Huge Section 1196}
% synthetic bench section 1196
% This is a sample LaTeX input file.  (Version of 12 August 2004.)
%
% A '%' character causes TeX to ignore all remaining text on the line,
% and is used for comments like this one.

\documentclass{article}      % Specifies the document class

                             % The preamble begins here.
\title{An Example Document}  % Declares the document's title.
\author{Leslie Lamport}      % Declares the author's name.
\date{January 21, 1994}      % Deleting this command produces today's date.

\newcommand{\ip}[2]{(#1, #2)}
                             % Defines \ip{arg1}{arg2} to mean
                             % (arg1, arg2).

%\newcommand{\ip}[2]{\langle #1 | #2\rangle}
                             % This is an alternative definition of
                             % \ip that is commented out.

\begin{document}             % End of preamble and beginning of text.

\maketitle                   % Produces the title.

This is an example input file.  Comparing it with
the output it generates can show you how to
produce a simple document of your own.

\section{Ordinary Text}      % Produces section heading.  Lower-level
                             % sections are begun with similar
                             % \subsection and \subsubsection commands.

The ends  of words and sentences are marked
  by   spaces. It  doesn't matter how many
spaces    you type; one is as good as 100.  The
end of   a line counts as a space.

One   or more   blank lines denote the  end
of  a paragraph.

Since any number of consecutive spaces are treated
like a single one, the formatting of the input
file makes no difference to
      \LaTeX,                % The \LaTeX command generates the LaTeX logo.
but it makes a difference to you.  When you use
\LaTeX, making your input file as easy to read
as possible will be a great help as you write
your document and when you change it.  This sample
file shows how you can add comments to your own input
file.

Because printing is different from typewriting,
there are a number of things that you have to do
differently when preparing an input file than if
you were just typing the document directly.
Quotation marks like
       ``this''
have to be handled specially, as do quotes within
quotes:
       ``\,`this'            % \, separates the double and single quote.
        is what I just
        wrote, not  `that'\,''.

Dashes come in three sizes: an
       intra-word
dash, a medium dash for number ranges like
       1--2,
and a punctuation
       dash---like
this.

A sentence-ending space should be larger than the
space between words within a sentence.  You
sometimes have to type special commands in
conjunction with punctuation characters to get
this right, as in the following sentence.
       Gnats, gnus, etc.\ all  % `\ ' makes an inter-word space.
       begin with G\@.         % \@ marks end-of-sentence punctuation.
You should check the spaces after periods when
reading your output to make sure you haven't
forgotten any special cases.  Generating an
ellipsis
       \ldots\               % `\ ' is needed after `\ldots' because TeX
                             % ignores spaces after command names like \ldots
                             % made from \ + letters.
                             %
                             % Note how a `%' character causes TeX to ignore
                             % the end of the input line, so these blank lines
                             % do not start a new paragraph.
                             %
with the right spacing around the periods requires
a special command.

\LaTeX\ interprets some common characters as
commands, so you must type special commands to
generate them.  These characters include the
following:
       \$ \& \% \# \{ and \}.

In printing, text is usually emphasized with an
       \emph{italic}
type style.

\begin{em}
   A long segment of text can also be emphasized
   in this way.  Text within such a segment can be
   given \emph{additional} emphasis.
\end{em}

It is sometimes necessary to prevent \LaTeX\ from
breaking a line where it might otherwise do so.
This may be at a space, as between the ``Mr.''\ and
``Jones'' in
       ``Mr.~Jones'',        % ~ produces an unbreakable interword space.
or within a word---especially when the word is a
symbol like
       \mbox{\emph{itemnum}}
that makes little sense when hyphenated across
lines.

Footnotes\footnote{This is an example of a footnote.}
pose no problem.

\LaTeX\ is good at typesetting mathematical formulas
like
       \( x-3y + z = 7 \)
or
       \( a_{1} > x^{2n} + y^{2n} > x' \)
or
       \( \ip{A}{B} = \sum_{i} a_{i} b_{i} \).
The spaces you type in a formula are
ignored.  Remember that a letter like
       $x$                   % $ ... $  and  \( ... \)  are equivalent
is a formula when it denotes a mathematical
symbol, and it should be typed as one.

\section{Displayed Text}

Text is displayed by indenting it from the left
margin.  Quotations are commonly displayed.  There
are short quotations
\begin{quote}
   This is a short quotation.  It consists of a
   single paragraph of text.  See how it is formatted.
\end{quote}
and longer ones.
\begin{quotation}
   This is a longer quotation.  It consists of two
   paragraphs of text, neither of which are
   particularly interesting.

   This is the second paragraph of the quotation.  It
   is just as dull as the first paragraph.
\end{quotation}
Another frequently-displayed structure is a list.
The following is an example of an \emph{itemized}
list.
\begin{itemize}
   \item This is the first item of an itemized list.
         Each item in the list is marked with a ``tick''.
         You don't have to worry about what kind of tick
         mark is used.

   \item This is the second item of the list.  It
         contains another list nested inside it.  The inner
         list is an \emph{enumerated} list.
         \begin{enumerate}
            \item This is the first item of an enumerated
                  list that is nested within the itemized list.

            \item This is the second item of the inner list.
                  \LaTeX\ allows you to nest lists deeper than
                  you really should.
         \end{enumerate}
         This is the rest of the second item of the outer
         list.  It is no more interesting than any other
         part of the item.
   \item This is the third item of the list.
\end{itemize}
You can even display poetry.
\begin{verse}
   There is an environment
    for verse \\             % The \\ command separates lines
   Whose features some poets % within a stanza.
   will curse.

                             % One or more blank lines separate stanzas.

   For instead of making\\
   Them do \emph{all} line breaking, \\
   It allows them to put too many words on a line when they'd rather be
   forced to be terse.
\end{verse}

Mathematical formulas may also be displayed.  A
displayed formula
is
one-line long; multiline
formulas require special formatting instructions.
   \[  \ip{\Gamma}{\psi'} = x'' + y^{2} + z_{i}^{n}\]
Don't start a paragraph with a displayed equation,
nor make one a paragraph by itself.

\end{document}               % End of document.

\section{Bench Huge Section 1197}
% synthetic bench section 1197
% This is a sample LaTeX input file.  (Version of 12 August 2004.)
%
% A '%' character causes TeX to ignore all remaining text on the line,
% and is used for comments like this one.

\documentclass{article}      % Specifies the document class

                             % The preamble begins here.
\title{An Example Document}  % Declares the document's title.
\author{Leslie Lamport}      % Declares the author's name.
\date{January 21, 1994}      % Deleting this command produces today's date.

\newcommand{\ip}[2]{(#1, #2)}
                             % Defines \ip{arg1}{arg2} to mean
                             % (arg1, arg2).

%\newcommand{\ip}[2]{\langle #1 | #2\rangle}
                             % This is an alternative definition of
                             % \ip that is commented out.

\begin{document}             % End of preamble and beginning of text.

\maketitle                   % Produces the title.

This is an example input file.  Comparing it with
the output it generates can show you how to
produce a simple document of your own.

\section{Ordinary Text}      % Produces section heading.  Lower-level
                             % sections are begun with similar
                             % \subsection and \subsubsection commands.

The ends  of words and sentences are marked
  by   spaces. It  doesn't matter how many
spaces    you type; one is as good as 100.  The
end of   a line counts as a space.

One   or more   blank lines denote the  end
of  a paragraph.

Since any number of consecutive spaces are treated
like a single one, the formatting of the input
file makes no difference to
      \LaTeX,                % The \LaTeX command generates the LaTeX logo.
but it makes a difference to you.  When you use
\LaTeX, making your input file as easy to read
as possible will be a great help as you write
your document and when you change it.  This sample
file shows how you can add comments to your own input
file.

Because printing is different from typewriting,
there are a number of things that you have to do
differently when preparing an input file than if
you were just typing the document directly.
Quotation marks like
       ``this''
have to be handled specially, as do quotes within
quotes:
       ``\,`this'            % \, separates the double and single quote.
        is what I just
        wrote, not  `that'\,''.

Dashes come in three sizes: an
       intra-word
dash, a medium dash for number ranges like
       1--2,
and a punctuation
       dash---like
this.

A sentence-ending space should be larger than the
space between words within a sentence.  You
sometimes have to type special commands in
conjunction with punctuation characters to get
this right, as in the following sentence.
       Gnats, gnus, etc.\ all  % `\ ' makes an inter-word space.
       begin with G\@.         % \@ marks end-of-sentence punctuation.
You should check the spaces after periods when
reading your output to make sure you haven't
forgotten any special cases.  Generating an
ellipsis
       \ldots\               % `\ ' is needed after `\ldots' because TeX
                             % ignores spaces after command names like \ldots
                             % made from \ + letters.
                             %
                             % Note how a `%' character causes TeX to ignore
                             % the end of the input line, so these blank lines
                             % do not start a new paragraph.
                             %
with the right spacing around the periods requires
a special command.

\LaTeX\ interprets some common characters as
commands, so you must type special commands to
generate them.  These characters include the
following:
       \$ \& \% \# \{ and \}.

In printing, text is usually emphasized with an
       \emph{italic}
type style.

\begin{em}
   A long segment of text can also be emphasized
   in this way.  Text within such a segment can be
   given \emph{additional} emphasis.
\end{em}

It is sometimes necessary to prevent \LaTeX\ from
breaking a line where it might otherwise do so.
This may be at a space, as between the ``Mr.''\ and
``Jones'' in
       ``Mr.~Jones'',        % ~ produces an unbreakable interword space.
or within a word---especially when the word is a
symbol like
       \mbox{\emph{itemnum}}
that makes little sense when hyphenated across
lines.

Footnotes\footnote{This is an example of a footnote.}
pose no problem.

\LaTeX\ is good at typesetting mathematical formulas
like
       \( x-3y + z = 7 \)
or
       \( a_{1} > x^{2n} + y^{2n} > x' \)
or
       \( \ip{A}{B} = \sum_{i} a_{i} b_{i} \).
The spaces you type in a formula are
ignored.  Remember that a letter like
       $x$                   % $ ... $  and  \( ... \)  are equivalent
is a formula when it denotes a mathematical
symbol, and it should be typed as one.

\section{Displayed Text}

Text is displayed by indenting it from the left
margin.  Quotations are commonly displayed.  There
are short quotations
\begin{quote}
   This is a short quotation.  It consists of a
   single paragraph of text.  See how it is formatted.
\end{quote}
and longer ones.
\begin{quotation}
   This is a longer quotation.  It consists of two
   paragraphs of text, neither of which are
   particularly interesting.

   This is the second paragraph of the quotation.  It
   is just as dull as the first paragraph.
\end{quotation}
Another frequently-displayed structure is a list.
The following is an example of an \emph{itemized}
list.
\begin{itemize}
   \item This is the first item of an itemized list.
         Each item in the list is marked with a ``tick''.
         You don't have to worry about what kind of tick
         mark is used.

   \item This is the second item of the list.  It
         contains another list nested inside it.  The inner
         list is an \emph{enumerated} list.
         \begin{enumerate}
            \item This is the first item of an enumerated
                  list that is nested within the itemized list.

            \item This is the second item of the inner list.
                  \LaTeX\ allows you to nest lists deeper than
                  you really should.
         \end{enumerate}
         This is the rest of the second item of the outer
         list.  It is no more interesting than any other
         part of the item.
   \item This is the third item of the list.
\end{itemize}
You can even display poetry.
\begin{verse}
   There is an environment
    for verse \\             % The \\ command separates lines
   Whose features some poets % within a stanza.
   will curse.

                             % One or more blank lines separate stanzas.

   For instead of making\\
   Them do \emph{all} line breaking, \\
   It allows them to put too many words on a line when they'd rather be
   forced to be terse.
\end{verse}

Mathematical formulas may also be displayed.  A
displayed formula
is
one-line long; multiline
formulas require special formatting instructions.
   \[  \ip{\Gamma}{\psi'} = x'' + y^{2} + z_{i}^{n}\]
Don't start a paragraph with a displayed equation,
nor make one a paragraph by itself.

\end{document}               % End of document.

\section{Bench Huge Section 1198}
% synthetic bench section 1198
% This is a sample LaTeX input file.  (Version of 12 August 2004.)
%
% A '%' character causes TeX to ignore all remaining text on the line,
% and is used for comments like this one.

\documentclass{article}      % Specifies the document class

                             % The preamble begins here.
\title{An Example Document}  % Declares the document's title.
\author{Leslie Lamport}      % Declares the author's name.
\date{January 21, 1994}      % Deleting this command produces today's date.

\newcommand{\ip}[2]{(#1, #2)}
                             % Defines \ip{arg1}{arg2} to mean
                             % (arg1, arg2).

%\newcommand{\ip}[2]{\langle #1 | #2\rangle}
                             % This is an alternative definition of
                             % \ip that is commented out.

\begin{document}             % End of preamble and beginning of text.

\maketitle                   % Produces the title.

This is an example input file.  Comparing it with
the output it generates can show you how to
produce a simple document of your own.

\section{Ordinary Text}      % Produces section heading.  Lower-level
                             % sections are begun with similar
                             % \subsection and \subsubsection commands.

The ends  of words and sentences are marked
  by   spaces. It  doesn't matter how many
spaces    you type; one is as good as 100.  The
end of   a line counts as a space.

One   or more   blank lines denote the  end
of  a paragraph.

Since any number of consecutive spaces are treated
like a single one, the formatting of the input
file makes no difference to
      \LaTeX,                % The \LaTeX command generates the LaTeX logo.
but it makes a difference to you.  When you use
\LaTeX, making your input file as easy to read
as possible will be a great help as you write
your document and when you change it.  This sample
file shows how you can add comments to your own input
file.

Because printing is different from typewriting,
there are a number of things that you have to do
differently when preparing an input file than if
you were just typing the document directly.
Quotation marks like
       ``this''
have to be handled specially, as do quotes within
quotes:
       ``\,`this'            % \, separates the double and single quote.
        is what I just
        wrote, not  `that'\,''.

Dashes come in three sizes: an
       intra-word
dash, a medium dash for number ranges like
       1--2,
and a punctuation
       dash---like
this.

A sentence-ending space should be larger than the
space between words within a sentence.  You
sometimes have to type special commands in
conjunction with punctuation characters to get
this right, as in the following sentence.
       Gnats, gnus, etc.\ all  % `\ ' makes an inter-word space.
       begin with G\@.         % \@ marks end-of-sentence punctuation.
You should check the spaces after periods when
reading your output to make sure you haven't
forgotten any special cases.  Generating an
ellipsis
       \ldots\               % `\ ' is needed after `\ldots' because TeX
                             % ignores spaces after command names like \ldots
                             % made from \ + letters.
                             %
                             % Note how a `%' character causes TeX to ignore
                             % the end of the input line, so these blank lines
                             % do not start a new paragraph.
                             %
with the right spacing around the periods requires
a special command.

\LaTeX\ interprets some common characters as
commands, so you must type special commands to
generate them.  These characters include the
following:
       \$ \& \% \# \{ and \}.

In printing, text is usually emphasized with an
       \emph{italic}
type style.

\begin{em}
   A long segment of text can also be emphasized
   in this way.  Text within such a segment can be
   given \emph{additional} emphasis.
\end{em}

It is sometimes necessary to prevent \LaTeX\ from
breaking a line where it might otherwise do so.
This may be at a space, as between the ``Mr.''\ and
``Jones'' in
       ``Mr.~Jones'',        % ~ produces an unbreakable interword space.
or within a word---especially when the word is a
symbol like
       \mbox{\emph{itemnum}}
that makes little sense when hyphenated across
lines.

Footnotes\footnote{This is an example of a footnote.}
pose no problem.

\LaTeX\ is good at typesetting mathematical formulas
like
       \( x-3y + z = 7 \)
or
       \( a_{1} > x^{2n} + y^{2n} > x' \)
or
       \( \ip{A}{B} = \sum_{i} a_{i} b_{i} \).
The spaces you type in a formula are
ignored.  Remember that a letter like
       $x$                   % $ ... $  and  \( ... \)  are equivalent
is a formula when it denotes a mathematical
symbol, and it should be typed as one.

\section{Displayed Text}

Text is displayed by indenting it from the left
margin.  Quotations are commonly displayed.  There
are short quotations
\begin{quote}
   This is a short quotation.  It consists of a
   single paragraph of text.  See how it is formatted.
\end{quote}
and longer ones.
\begin{quotation}
   This is a longer quotation.  It consists of two
   paragraphs of text, neither of which are
   particularly interesting.

   This is the second paragraph of the quotation.  It
   is just as dull as the first paragraph.
\end{quotation}
Another frequently-displayed structure is a list.
The following is an example of an \emph{itemized}
list.
\begin{itemize}
   \item This is the first item of an itemized list.
         Each item in the list is marked with a ``tick''.
         You don't have to worry about what kind of tick
         mark is used.

   \item This is the second item of the list.  It
         contains another list nested inside it.  The inner
         list is an \emph{enumerated} list.
         \begin{enumerate}
            \item This is the first item of an enumerated
                  list that is nested within the itemized list.

            \item This is the second item of the inner list.
                  \LaTeX\ allows you to nest lists deeper than
                  you really should.
         \end{enumerate}
         This is the rest of the second item of the outer
         list.  It is no more interesting than any other
         part of the item.
   \item This is the third item of the list.
\end{itemize}
You can even display poetry.
\begin{verse}
   There is an environment
    for verse \\             % The \\ command separates lines
   Whose features some poets % within a stanza.
   will curse.

                             % One or more blank lines separate stanzas.

   For instead of making\\
   Them do \emph{all} line breaking, \\
   It allows them to put too many words on a line when they'd rather be
   forced to be terse.
\end{verse}

Mathematical formulas may also be displayed.  A
displayed formula
is
one-line long; multiline
formulas require special formatting instructions.
   \[  \ip{\Gamma}{\psi'} = x'' + y^{2} + z_{i}^{n}\]
Don't start a paragraph with a displayed equation,
nor make one a paragraph by itself.

\end{document}               % End of document.

\section{Bench Huge Section 1199}
% synthetic bench section 1199
% This is a sample LaTeX input file.  (Version of 12 August 2004.)
%
% A '%' character causes TeX to ignore all remaining text on the line,
% and is used for comments like this one.

\documentclass{article}      % Specifies the document class

                             % The preamble begins here.
\title{An Example Document}  % Declares the document's title.
\author{Leslie Lamport}      % Declares the author's name.
\date{January 21, 1994}      % Deleting this command produces today's date.

\newcommand{\ip}[2]{(#1, #2)}
                             % Defines \ip{arg1}{arg2} to mean
                             % (arg1, arg2).

%\newcommand{\ip}[2]{\langle #1 | #2\rangle}
                             % This is an alternative definition of
                             % \ip that is commented out.

\begin{document}             % End of preamble and beginning of text.

\maketitle                   % Produces the title.

This is an example input file.  Comparing it with
the output it generates can show you how to
produce a simple document of your own.

\section{Ordinary Text}      % Produces section heading.  Lower-level
                             % sections are begun with similar
                             % \subsection and \subsubsection commands.

The ends  of words and sentences are marked
  by   spaces. It  doesn't matter how many
spaces    you type; one is as good as 100.  The
end of   a line counts as a space.

One   or more   blank lines denote the  end
of  a paragraph.

Since any number of consecutive spaces are treated
like a single one, the formatting of the input
file makes no difference to
      \LaTeX,                % The \LaTeX command generates the LaTeX logo.
but it makes a difference to you.  When you use
\LaTeX, making your input file as easy to read
as possible will be a great help as you write
your document and when you change it.  This sample
file shows how you can add comments to your own input
file.

Because printing is different from typewriting,
there are a number of things that you have to do
differently when preparing an input file than if
you were just typing the document directly.
Quotation marks like
       ``this''
have to be handled specially, as do quotes within
quotes:
       ``\,`this'            % \, separates the double and single quote.
        is what I just
        wrote, not  `that'\,''.

Dashes come in three sizes: an
       intra-word
dash, a medium dash for number ranges like
       1--2,
and a punctuation
       dash---like
this.

A sentence-ending space should be larger than the
space between words within a sentence.  You
sometimes have to type special commands in
conjunction with punctuation characters to get
this right, as in the following sentence.
       Gnats, gnus, etc.\ all  % `\ ' makes an inter-word space.
       begin with G\@.         % \@ marks end-of-sentence punctuation.
You should check the spaces after periods when
reading your output to make sure you haven't
forgotten any special cases.  Generating an
ellipsis
       \ldots\               % `\ ' is needed after `\ldots' because TeX
                             % ignores spaces after command names like \ldots
                             % made from \ + letters.
                             %
                             % Note how a `%' character causes TeX to ignore
                             % the end of the input line, so these blank lines
                             % do not start a new paragraph.
                             %
with the right spacing around the periods requires
a special command.

\LaTeX\ interprets some common characters as
commands, so you must type special commands to
generate them.  These characters include the
following:
       \$ \& \% \# \{ and \}.

In printing, text is usually emphasized with an
       \emph{italic}
type style.

\begin{em}
   A long segment of text can also be emphasized
   in this way.  Text within such a segment can be
   given \emph{additional} emphasis.
\end{em}

It is sometimes necessary to prevent \LaTeX\ from
breaking a line where it might otherwise do so.
This may be at a space, as between the ``Mr.''\ and
``Jones'' in
       ``Mr.~Jones'',        % ~ produces an unbreakable interword space.
or within a word---especially when the word is a
symbol like
       \mbox{\emph{itemnum}}
that makes little sense when hyphenated across
lines.

Footnotes\footnote{This is an example of a footnote.}
pose no problem.

\LaTeX\ is good at typesetting mathematical formulas
like
       \( x-3y + z = 7 \)
or
       \( a_{1} > x^{2n} + y^{2n} > x' \)
or
       \( \ip{A}{B} = \sum_{i} a_{i} b_{i} \).
The spaces you type in a formula are
ignored.  Remember that a letter like
       $x$                   % $ ... $  and  \( ... \)  are equivalent
is a formula when it denotes a mathematical
symbol, and it should be typed as one.

\section{Displayed Text}

Text is displayed by indenting it from the left
margin.  Quotations are commonly displayed.  There
are short quotations
\begin{quote}
   This is a short quotation.  It consists of a
   single paragraph of text.  See how it is formatted.
\end{quote}
and longer ones.
\begin{quotation}
   This is a longer quotation.  It consists of two
   paragraphs of text, neither of which are
   particularly interesting.

   This is the second paragraph of the quotation.  It
   is just as dull as the first paragraph.
\end{quotation}
Another frequently-displayed structure is a list.
The following is an example of an \emph{itemized}
list.
\begin{itemize}
   \item This is the first item of an itemized list.
         Each item in the list is marked with a ``tick''.
         You don't have to worry about what kind of tick
         mark is used.

   \item This is the second item of the list.  It
         contains another list nested inside it.  The inner
         list is an \emph{enumerated} list.
         \begin{enumerate}
            \item This is the first item of an enumerated
                  list that is nested within the itemized list.

            \item This is the second item of the inner list.
                  \LaTeX\ allows you to nest lists deeper than
                  you really should.
         \end{enumerate}
         This is the rest of the second item of the outer
         list.  It is no more interesting than any other
         part of the item.
   \item This is the third item of the list.
\end{itemize}
You can even display poetry.
\begin{verse}
   There is an environment
    for verse \\             % The \\ command separates lines
   Whose features some poets % within a stanza.
   will curse.

                             % One or more blank lines separate stanzas.

   For instead of making\\
   Them do \emph{all} line breaking, \\
   It allows them to put too many words on a line when they'd rather be
   forced to be terse.
\end{verse}

Mathematical formulas may also be displayed.  A
displayed formula
is
one-line long; multiline
formulas require special formatting instructions.
   \[  \ip{\Gamma}{\psi'} = x'' + y^{2} + z_{i}^{n}\]
Don't start a paragraph with a displayed equation,
nor make one a paragraph by itself.

\end{document}               % End of document.

\section{Bench Huge Section 1200}
% synthetic bench section 1200
% This is a sample LaTeX input file.  (Version of 12 August 2004.)
%
% A '%' character causes TeX to ignore all remaining text on the line,
% and is used for comments like this one.

\documentclass{article}      % Specifies the document class

                             % The preamble begins here.
\title{An Example Document}  % Declares the document's title.
\author{Leslie Lamport}      % Declares the author's name.
\date{January 21, 1994}      % Deleting this command produces today's date.

\newcommand{\ip}[2]{(#1, #2)}
                             % Defines \ip{arg1}{arg2} to mean
                             % (arg1, arg2).

%\newcommand{\ip}[2]{\langle #1 | #2\rangle}
                             % This is an alternative definition of
                             % \ip that is commented out.

\begin{document}             % End of preamble and beginning of text.

\maketitle                   % Produces the title.

This is an example input file.  Comparing it with
the output it generates can show you how to
produce a simple document of your own.

\section{Ordinary Text}      % Produces section heading.  Lower-level
                             % sections are begun with similar
                             % \subsection and \subsubsection commands.

The ends  of words and sentences are marked
  by   spaces. It  doesn't matter how many
spaces    you type; one is as good as 100.  The
end of   a line counts as a space.

One   or more   blank lines denote the  end
of  a paragraph.

Since any number of consecutive spaces are treated
like a single one, the formatting of the input
file makes no difference to
      \LaTeX,                % The \LaTeX command generates the LaTeX logo.
but it makes a difference to you.  When you use
\LaTeX, making your input file as easy to read
as possible will be a great help as you write
your document and when you change it.  This sample
file shows how you can add comments to your own input
file.

Because printing is different from typewriting,
there are a number of things that you have to do
differently when preparing an input file than if
you were just typing the document directly.
Quotation marks like
       ``this''
have to be handled specially, as do quotes within
quotes:
       ``\,`this'            % \, separates the double and single quote.
        is what I just
        wrote, not  `that'\,''.

Dashes come in three sizes: an
       intra-word
dash, a medium dash for number ranges like
       1--2,
and a punctuation
       dash---like
this.

A sentence-ending space should be larger than the
space between words within a sentence.  You
sometimes have to type special commands in
conjunction with punctuation characters to get
this right, as in the following sentence.
       Gnats, gnus, etc.\ all  % `\ ' makes an inter-word space.
       begin with G\@.         % \@ marks end-of-sentence punctuation.
You should check the spaces after periods when
reading your output to make sure you haven't
forgotten any special cases.  Generating an
ellipsis
       \ldots\               % `\ ' is needed after `\ldots' because TeX
                             % ignores spaces after command names like \ldots
                             % made from \ + letters.
                             %
                             % Note how a `%' character causes TeX to ignore
                             % the end of the input line, so these blank lines
                             % do not start a new paragraph.
                             %
with the right spacing around the periods requires
a special command.

\LaTeX\ interprets some common characters as
commands, so you must type special commands to
generate them.  These characters include the
following:
       \$ \& \% \# \{ and \}.

In printing, text is usually emphasized with an
       \emph{italic}
type style.

\begin{em}
   A long segment of text can also be emphasized
   in this way.  Text within such a segment can be
   given \emph{additional} emphasis.
\end{em}

It is sometimes necessary to prevent \LaTeX\ from
breaking a line where it might otherwise do so.
This may be at a space, as between the ``Mr.''\ and
``Jones'' in
       ``Mr.~Jones'',        % ~ produces an unbreakable interword space.
or within a word---especially when the word is a
symbol like
       \mbox{\emph{itemnum}}
that makes little sense when hyphenated across
lines.

Footnotes\footnote{This is an example of a footnote.}
pose no problem.

\LaTeX\ is good at typesetting mathematical formulas
like
       \( x-3y + z = 7 \)
or
       \( a_{1} > x^{2n} + y^{2n} > x' \)
or
       \( \ip{A}{B} = \sum_{i} a_{i} b_{i} \).
The spaces you type in a formula are
ignored.  Remember that a letter like
       $x$                   % $ ... $  and  \( ... \)  are equivalent
is a formula when it denotes a mathematical
symbol, and it should be typed as one.

\section{Displayed Text}

Text is displayed by indenting it from the left
margin.  Quotations are commonly displayed.  There
are short quotations
\begin{quote}
   This is a short quotation.  It consists of a
   single paragraph of text.  See how it is formatted.
\end{quote}
and longer ones.
\begin{quotation}
   This is a longer quotation.  It consists of two
   paragraphs of text, neither of which are
   particularly interesting.

   This is the second paragraph of the quotation.  It
   is just as dull as the first paragraph.
\end{quotation}
Another frequently-displayed structure is a list.
The following is an example of an \emph{itemized}
list.
\begin{itemize}
   \item This is the first item of an itemized list.
         Each item in the list is marked with a ``tick''.
         You don't have to worry about what kind of tick
         mark is used.

   \item This is the second item of the list.  It
         contains another list nested inside it.  The inner
         list is an \emph{enumerated} list.
         \begin{enumerate}
            \item This is the first item of an enumerated
                  list that is nested within the itemized list.

            \item This is the second item of the inner list.
                  \LaTeX\ allows you to nest lists deeper than
                  you really should.
         \end{enumerate}
         This is the rest of the second item of the outer
         list.  It is no more interesting than any other
         part of the item.
   \item This is the third item of the list.
\end{itemize}
You can even display poetry.
\begin{verse}
   There is an environment
    for verse \\             % The \\ command separates lines
   Whose features some poets % within a stanza.
   will curse.

                             % One or more blank lines separate stanzas.

   For instead of making\\
   Them do \emph{all} line breaking, \\
   It allows them to put too many words on a line when they'd rather be
   forced to be terse.
\end{verse}

Mathematical formulas may also be displayed.  A
displayed formula
is
one-line long; multiline
formulas require special formatting instructions.
   \[  \ip{\Gamma}{\psi'} = x'' + y^{2} + z_{i}^{n}\]
Don't start a paragraph with a displayed equation,
nor make one a paragraph by itself.

\end{document}               % End of document.

\section{Bench Huge Section 1201}
% synthetic bench section 1201
% This is a sample LaTeX input file.  (Version of 12 August 2004.)
%
% A '%' character causes TeX to ignore all remaining text on the line,
% and is used for comments like this one.

\documentclass{article}      % Specifies the document class

                             % The preamble begins here.
\title{An Example Document}  % Declares the document's title.
\author{Leslie Lamport}      % Declares the author's name.
\date{January 21, 1994}      % Deleting this command produces today's date.

\newcommand{\ip}[2]{(#1, #2)}
                             % Defines \ip{arg1}{arg2} to mean
                             % (arg1, arg2).

%\newcommand{\ip}[2]{\langle #1 | #2\rangle}
                             % This is an alternative definition of
                             % \ip that is commented out.

\begin{document}             % End of preamble and beginning of text.

\maketitle                   % Produces the title.

This is an example input file.  Comparing it with
the output it generates can show you how to
produce a simple document of your own.

\section{Ordinary Text}      % Produces section heading.  Lower-level
                             % sections are begun with similar
                             % \subsection and \subsubsection commands.

The ends  of words and sentences are marked
  by   spaces. It  doesn't matter how many
spaces    you type; one is as good as 100.  The
end of   a line counts as a space.

One   or more   blank lines denote the  end
of  a paragraph.

Since any number of consecutive spaces are treated
like a single one, the formatting of the input
file makes no difference to
      \LaTeX,                % The \LaTeX command generates the LaTeX logo.
but it makes a difference to you.  When you use
\LaTeX, making your input file as easy to read
as possible will be a great help as you write
your document and when you change it.  This sample
file shows how you can add comments to your own input
file.

Because printing is different from typewriting,
there are a number of things that you have to do
differently when preparing an input file than if
you were just typing the document directly.
Quotation marks like
       ``this''
have to be handled specially, as do quotes within
quotes:
       ``\,`this'            % \, separates the double and single quote.
        is what I just
        wrote, not  `that'\,''.

Dashes come in three sizes: an
       intra-word
dash, a medium dash for number ranges like
       1--2,
and a punctuation
       dash---like
this.

A sentence-ending space should be larger than the
space between words within a sentence.  You
sometimes have to type special commands in
conjunction with punctuation characters to get
this right, as in the following sentence.
       Gnats, gnus, etc.\ all  % `\ ' makes an inter-word space.
       begin with G\@.         % \@ marks end-of-sentence punctuation.
You should check the spaces after periods when
reading your output to make sure you haven't
forgotten any special cases.  Generating an
ellipsis
       \ldots\               % `\ ' is needed after `\ldots' because TeX
                             % ignores spaces after command names like \ldots
                             % made from \ + letters.
                             %
                             % Note how a `%' character causes TeX to ignore
                             % the end of the input line, so these blank lines
                             % do not start a new paragraph.
                             %
with the right spacing around the periods requires
a special command.

\LaTeX\ interprets some common characters as
commands, so you must type special commands to
generate them.  These characters include the
following:
       \$ \& \% \# \{ and \}.

In printing, text is usually emphasized with an
       \emph{italic}
type style.

\begin{em}
   A long segment of text can also be emphasized
   in this way.  Text within such a segment can be
   given \emph{additional} emphasis.
\end{em}

It is sometimes necessary to prevent \LaTeX\ from
breaking a line where it might otherwise do so.
This may be at a space, as between the ``Mr.''\ and
``Jones'' in
       ``Mr.~Jones'',        % ~ produces an unbreakable interword space.
or within a word---especially when the word is a
symbol like
       \mbox{\emph{itemnum}}
that makes little sense when hyphenated across
lines.

Footnotes\footnote{This is an example of a footnote.}
pose no problem.

\LaTeX\ is good at typesetting mathematical formulas
like
       \( x-3y + z = 7 \)
or
       \( a_{1} > x^{2n} + y^{2n} > x' \)
or
       \( \ip{A}{B} = \sum_{i} a_{i} b_{i} \).
The spaces you type in a formula are
ignored.  Remember that a letter like
       $x$                   % $ ... $  and  \( ... \)  are equivalent
is a formula when it denotes a mathematical
symbol, and it should be typed as one.

\section{Displayed Text}

Text is displayed by indenting it from the left
margin.  Quotations are commonly displayed.  There
are short quotations
\begin{quote}
   This is a short quotation.  It consists of a
   single paragraph of text.  See how it is formatted.
\end{quote}
and longer ones.
\begin{quotation}
   This is a longer quotation.  It consists of two
   paragraphs of text, neither of which are
   particularly interesting.

   This is the second paragraph of the quotation.  It
   is just as dull as the first paragraph.
\end{quotation}
Another frequently-displayed structure is a list.
The following is an example of an \emph{itemized}
list.
\begin{itemize}
   \item This is the first item of an itemized list.
         Each item in the list is marked with a ``tick''.
         You don't have to worry about what kind of tick
         mark is used.

   \item This is the second item of the list.  It
         contains another list nested inside it.  The inner
         list is an \emph{enumerated} list.
         \begin{enumerate}
            \item This is the first item of an enumerated
                  list that is nested within the itemized list.

            \item This is the second item of the inner list.
                  \LaTeX\ allows you to nest lists deeper than
                  you really should.
         \end{enumerate}
         This is the rest of the second item of the outer
         list.  It is no more interesting than any other
         part of the item.
   \item This is the third item of the list.
\end{itemize}
You can even display poetry.
\begin{verse}
   There is an environment
    for verse \\             % The \\ command separates lines
   Whose features some poets % within a stanza.
   will curse.

                             % One or more blank lines separate stanzas.

   For instead of making\\
   Them do \emph{all} line breaking, \\
   It allows them to put too many words on a line when they'd rather be
   forced to be terse.
\end{verse}

Mathematical formulas may also be displayed.  A
displayed formula
is
one-line long; multiline
formulas require special formatting instructions.
   \[  \ip{\Gamma}{\psi'} = x'' + y^{2} + z_{i}^{n}\]
Don't start a paragraph with a displayed equation,
nor make one a paragraph by itself.

\end{document}               % End of document.

\section{Bench Huge Section 1202}
% synthetic bench section 1202
% This is a sample LaTeX input file.  (Version of 12 August 2004.)
%
% A '%' character causes TeX to ignore all remaining text on the line,
% and is used for comments like this one.

\documentclass{article}      % Specifies the document class

                             % The preamble begins here.
\title{An Example Document}  % Declares the document's title.
\author{Leslie Lamport}      % Declares the author's name.
\date{January 21, 1994}      % Deleting this command produces today's date.

\newcommand{\ip}[2]{(#1, #2)}
                             % Defines \ip{arg1}{arg2} to mean
                             % (arg1, arg2).

%\newcommand{\ip}[2]{\langle #1 | #2\rangle}
                             % This is an alternative definition of
                             % \ip that is commented out.

\begin{document}             % End of preamble and beginning of text.

\maketitle                   % Produces the title.

This is an example input file.  Comparing it with
the output it generates can show you how to
produce a simple document of your own.

\section{Ordinary Text}      % Produces section heading.  Lower-level
                             % sections are begun with similar
                             % \subsection and \subsubsection commands.

The ends  of words and sentences are marked
  by   spaces. It  doesn't matter how many
spaces    you type; one is as good as 100.  The
end of   a line counts as a space.

One   or more   blank lines denote the  end
of  a paragraph.

Since any number of consecutive spaces are treated
like a single one, the formatting of the input
file makes no difference to
      \LaTeX,                % The \LaTeX command generates the LaTeX logo.
but it makes a difference to you.  When you use
\LaTeX, making your input file as easy to read
as possible will be a great help as you write
your document and when you change it.  This sample
file shows how you can add comments to your own input
file.

Because printing is different from typewriting,
there are a number of things that you have to do
differently when preparing an input file than if
you were just typing the document directly.
Quotation marks like
       ``this''
have to be handled specially, as do quotes within
quotes:
       ``\,`this'            % \, separates the double and single quote.
        is what I just
        wrote, not  `that'\,''.

Dashes come in three sizes: an
       intra-word
dash, a medium dash for number ranges like
       1--2,
and a punctuation
       dash---like
this.

A sentence-ending space should be larger than the
space between words within a sentence.  You
sometimes have to type special commands in
conjunction with punctuation characters to get
this right, as in the following sentence.
       Gnats, gnus, etc.\ all  % `\ ' makes an inter-word space.
       begin with G\@.         % \@ marks end-of-sentence punctuation.
You should check the spaces after periods when
reading your output to make sure you haven't
forgotten any special cases.  Generating an
ellipsis
       \ldots\               % `\ ' is needed after `\ldots' because TeX
                             % ignores spaces after command names like \ldots
                             % made from \ + letters.
                             %
                             % Note how a `%' character causes TeX to ignore
                             % the end of the input line, so these blank lines
                             % do not start a new paragraph.
                             %
with the right spacing around the periods requires
a special command.

\LaTeX\ interprets some common characters as
commands, so you must type special commands to
generate them.  These characters include the
following:
       \$ \& \% \# \{ and \}.

In printing, text is usually emphasized with an
       \emph{italic}
type style.

\begin{em}
   A long segment of text can also be emphasized
   in this way.  Text within such a segment can be
   given \emph{additional} emphasis.
\end{em}

It is sometimes necessary to prevent \LaTeX\ from
breaking a line where it might otherwise do so.
This may be at a space, as between the ``Mr.''\ and
``Jones'' in
       ``Mr.~Jones'',        % ~ produces an unbreakable interword space.
or within a word---especially when the word is a
symbol like
       \mbox{\emph{itemnum}}
that makes little sense when hyphenated across
lines.

Footnotes\footnote{This is an example of a footnote.}
pose no problem.

\LaTeX\ is good at typesetting mathematical formulas
like
       \( x-3y + z = 7 \)
or
       \( a_{1} > x^{2n} + y^{2n} > x' \)
or
       \( \ip{A}{B} = \sum_{i} a_{i} b_{i} \).
The spaces you type in a formula are
ignored.  Remember that a letter like
       $x$                   % $ ... $  and  \( ... \)  are equivalent
is a formula when it denotes a mathematical
symbol, and it should be typed as one.

\section{Displayed Text}

Text is displayed by indenting it from the left
margin.  Quotations are commonly displayed.  There
are short quotations
\begin{quote}
   This is a short quotation.  It consists of a
   single paragraph of text.  See how it is formatted.
\end{quote}
and longer ones.
\begin{quotation}
   This is a longer quotation.  It consists of two
   paragraphs of text, neither of which are
   particularly interesting.

   This is the second paragraph of the quotation.  It
   is just as dull as the first paragraph.
\end{quotation}
Another frequently-displayed structure is a list.
The following is an example of an \emph{itemized}
list.
\begin{itemize}
   \item This is the first item of an itemized list.
         Each item in the list is marked with a ``tick''.
         You don't have to worry about what kind of tick
         mark is used.

   \item This is the second item of the list.  It
         contains another list nested inside it.  The inner
         list is an \emph{enumerated} list.
         \begin{enumerate}
            \item This is the first item of an enumerated
                  list that is nested within the itemized list.

            \item This is the second item of the inner list.
                  \LaTeX\ allows you to nest lists deeper than
                  you really should.
         \end{enumerate}
         This is the rest of the second item of the outer
         list.  It is no more interesting than any other
         part of the item.
   \item This is the third item of the list.
\end{itemize}
You can even display poetry.
\begin{verse}
   There is an environment
    for verse \\             % The \\ command separates lines
   Whose features some poets % within a stanza.
   will curse.

                             % One or more blank lines separate stanzas.

   For instead of making\\
   Them do \emph{all} line breaking, \\
   It allows them to put too many words on a line when they'd rather be
   forced to be terse.
\end{verse}

Mathematical formulas may also be displayed.  A
displayed formula
is
one-line long; multiline
formulas require special formatting instructions.
   \[  \ip{\Gamma}{\psi'} = x'' + y^{2} + z_{i}^{n}\]
Don't start a paragraph with a displayed equation,
nor make one a paragraph by itself.

\end{document}               % End of document.

\section{Bench Huge Section 1203}
% synthetic bench section 1203
% This is a sample LaTeX input file.  (Version of 12 August 2004.)
%
% A '%' character causes TeX to ignore all remaining text on the line,
% and is used for comments like this one.

\documentclass{article}      % Specifies the document class

                             % The preamble begins here.
\title{An Example Document}  % Declares the document's title.
\author{Leslie Lamport}      % Declares the author's name.
\date{January 21, 1994}      % Deleting this command produces today's date.

\newcommand{\ip}[2]{(#1, #2)}
                             % Defines \ip{arg1}{arg2} to mean
                             % (arg1, arg2).

%\newcommand{\ip}[2]{\langle #1 | #2\rangle}
                             % This is an alternative definition of
                             % \ip that is commented out.

\begin{document}             % End of preamble and beginning of text.

\maketitle                   % Produces the title.

This is an example input file.  Comparing it with
the output it generates can show you how to
produce a simple document of your own.

\section{Ordinary Text}      % Produces section heading.  Lower-level
                             % sections are begun with similar
                             % \subsection and \subsubsection commands.

The ends  of words and sentences are marked
  by   spaces. It  doesn't matter how many
spaces    you type; one is as good as 100.  The
end of   a line counts as a space.

One   or more   blank lines denote the  end
of  a paragraph.

Since any number of consecutive spaces are treated
like a single one, the formatting of the input
file makes no difference to
      \LaTeX,                % The \LaTeX command generates the LaTeX logo.
but it makes a difference to you.  When you use
\LaTeX, making your input file as easy to read
as possible will be a great help as you write
your document and when you change it.  This sample
file shows how you can add comments to your own input
file.

Because printing is different from typewriting,
there are a number of things that you have to do
differently when preparing an input file than if
you were just typing the document directly.
Quotation marks like
       ``this''
have to be handled specially, as do quotes within
quotes:
       ``\,`this'            % \, separates the double and single quote.
        is what I just
        wrote, not  `that'\,''.

Dashes come in three sizes: an
       intra-word
dash, a medium dash for number ranges like
       1--2,
and a punctuation
       dash---like
this.

A sentence-ending space should be larger than the
space between words within a sentence.  You
sometimes have to type special commands in
conjunction with punctuation characters to get
this right, as in the following sentence.
       Gnats, gnus, etc.\ all  % `\ ' makes an inter-word space.
       begin with G\@.         % \@ marks end-of-sentence punctuation.
You should check the spaces after periods when
reading your output to make sure you haven't
forgotten any special cases.  Generating an
ellipsis
       \ldots\               % `\ ' is needed after `\ldots' because TeX
                             % ignores spaces after command names like \ldots
                             % made from \ + letters.
                             %
                             % Note how a `%' character causes TeX to ignore
                             % the end of the input line, so these blank lines
                             % do not start a new paragraph.
                             %
with the right spacing around the periods requires
a special command.

\LaTeX\ interprets some common characters as
commands, so you must type special commands to
generate them.  These characters include the
following:
       \$ \& \% \# \{ and \}.

In printing, text is usually emphasized with an
       \emph{italic}
type style.

\begin{em}
   A long segment of text can also be emphasized
   in this way.  Text within such a segment can be
   given \emph{additional} emphasis.
\end{em}

It is sometimes necessary to prevent \LaTeX\ from
breaking a line where it might otherwise do so.
This may be at a space, as between the ``Mr.''\ and
``Jones'' in
       ``Mr.~Jones'',        % ~ produces an unbreakable interword space.
or within a word---especially when the word is a
symbol like
       \mbox{\emph{itemnum}}
that makes little sense when hyphenated across
lines.

Footnotes\footnote{This is an example of a footnote.}
pose no problem.

\LaTeX\ is good at typesetting mathematical formulas
like
       \( x-3y + z = 7 \)
or
       \( a_{1} > x^{2n} + y^{2n} > x' \)
or
       \( \ip{A}{B} = \sum_{i} a_{i} b_{i} \).
The spaces you type in a formula are
ignored.  Remember that a letter like
       $x$                   % $ ... $  and  \( ... \)  are equivalent
is a formula when it denotes a mathematical
symbol, and it should be typed as one.

\section{Displayed Text}

Text is displayed by indenting it from the left
margin.  Quotations are commonly displayed.  There
are short quotations
\begin{quote}
   This is a short quotation.  It consists of a
   single paragraph of text.  See how it is formatted.
\end{quote}
and longer ones.
\begin{quotation}
   This is a longer quotation.  It consists of two
   paragraphs of text, neither of which are
   particularly interesting.

   This is the second paragraph of the quotation.  It
   is just as dull as the first paragraph.
\end{quotation}
Another frequently-displayed structure is a list.
The following is an example of an \emph{itemized}
list.
\begin{itemize}
   \item This is the first item of an itemized list.
         Each item in the list is marked with a ``tick''.
         You don't have to worry about what kind of tick
         mark is used.

   \item This is the second item of the list.  It
         contains another list nested inside it.  The inner
         list is an \emph{enumerated} list.
         \begin{enumerate}
            \item This is the first item of an enumerated
                  list that is nested within the itemized list.

            \item This is the second item of the inner list.
                  \LaTeX\ allows you to nest lists deeper than
                  you really should.
         \end{enumerate}
         This is the rest of the second item of the outer
         list.  It is no more interesting than any other
         part of the item.
   \item This is the third item of the list.
\end{itemize}
You can even display poetry.
\begin{verse}
   There is an environment
    for verse \\             % The \\ command separates lines
   Whose features some poets % within a stanza.
   will curse.

                             % One or more blank lines separate stanzas.

   For instead of making\\
   Them do \emph{all} line breaking, \\
   It allows them to put too many words on a line when they'd rather be
   forced to be terse.
\end{verse}

Mathematical formulas may also be displayed.  A
displayed formula
is
one-line long; multiline
formulas require special formatting instructions.
   \[  \ip{\Gamma}{\psi'} = x'' + y^{2} + z_{i}^{n}\]
Don't start a paragraph with a displayed equation,
nor make one a paragraph by itself.

\end{document}               % End of document.

\section{Bench Huge Section 1204}
% synthetic bench section 1204
% This is a sample LaTeX input file.  (Version of 12 August 2004.)
%
% A '%' character causes TeX to ignore all remaining text on the line,
% and is used for comments like this one.

\documentclass{article}      % Specifies the document class

                             % The preamble begins here.
\title{An Example Document}  % Declares the document's title.
\author{Leslie Lamport}      % Declares the author's name.
\date{January 21, 1994}      % Deleting this command produces today's date.

\newcommand{\ip}[2]{(#1, #2)}
                             % Defines \ip{arg1}{arg2} to mean
                             % (arg1, arg2).

%\newcommand{\ip}[2]{\langle #1 | #2\rangle}
                             % This is an alternative definition of
                             % \ip that is commented out.

\begin{document}             % End of preamble and beginning of text.

\maketitle                   % Produces the title.

This is an example input file.  Comparing it with
the output it generates can show you how to
produce a simple document of your own.

\section{Ordinary Text}      % Produces section heading.  Lower-level
                             % sections are begun with similar
                             % \subsection and \subsubsection commands.

The ends  of words and sentences are marked
  by   spaces. It  doesn't matter how many
spaces    you type; one is as good as 100.  The
end of   a line counts as a space.

One   or more   blank lines denote the  end
of  a paragraph.

Since any number of consecutive spaces are treated
like a single one, the formatting of the input
file makes no difference to
      \LaTeX,                % The \LaTeX command generates the LaTeX logo.
but it makes a difference to you.  When you use
\LaTeX, making your input file as easy to read
as possible will be a great help as you write
your document and when you change it.  This sample
file shows how you can add comments to your own input
file.

Because printing is different from typewriting,
there are a number of things that you have to do
differently when preparing an input file than if
you were just typing the document directly.
Quotation marks like
       ``this''
have to be handled specially, as do quotes within
quotes:
       ``\,`this'            % \, separates the double and single quote.
        is what I just
        wrote, not  `that'\,''.

Dashes come in three sizes: an
       intra-word
dash, a medium dash for number ranges like
       1--2,
and a punctuation
       dash---like
this.

A sentence-ending space should be larger than the
space between words within a sentence.  You
sometimes have to type special commands in
conjunction with punctuation characters to get
this right, as in the following sentence.
       Gnats, gnus, etc.\ all  % `\ ' makes an inter-word space.
       begin with G\@.         % \@ marks end-of-sentence punctuation.
You should check the spaces after periods when
reading your output to make sure you haven't
forgotten any special cases.  Generating an
ellipsis
       \ldots\               % `\ ' is needed after `\ldots' because TeX
                             % ignores spaces after command names like \ldots
                             % made from \ + letters.
                             %
                             % Note how a `%' character causes TeX to ignore
                             % the end of the input line, so these blank lines
                             % do not start a new paragraph.
                             %
with the right spacing around the periods requires
a special command.

\LaTeX\ interprets some common characters as
commands, so you must type special commands to
generate them.  These characters include the
following:
       \$ \& \% \# \{ and \}.

In printing, text is usually emphasized with an
       \emph{italic}
type style.

\begin{em}
   A long segment of text can also be emphasized
   in this way.  Text within such a segment can be
   given \emph{additional} emphasis.
\end{em}

It is sometimes necessary to prevent \LaTeX\ from
breaking a line where it might otherwise do so.
This may be at a space, as between the ``Mr.''\ and
``Jones'' in
       ``Mr.~Jones'',        % ~ produces an unbreakable interword space.
or within a word---especially when the word is a
symbol like
       \mbox{\emph{itemnum}}
that makes little sense when hyphenated across
lines.

Footnotes\footnote{This is an example of a footnote.}
pose no problem.

\LaTeX\ is good at typesetting mathematical formulas
like
       \( x-3y + z = 7 \)
or
       \( a_{1} > x^{2n} + y^{2n} > x' \)
or
       \( \ip{A}{B} = \sum_{i} a_{i} b_{i} \).
The spaces you type in a formula are
ignored.  Remember that a letter like
       $x$                   % $ ... $  and  \( ... \)  are equivalent
is a formula when it denotes a mathematical
symbol, and it should be typed as one.

\section{Displayed Text}

Text is displayed by indenting it from the left
margin.  Quotations are commonly displayed.  There
are short quotations
\begin{quote}
   This is a short quotation.  It consists of a
   single paragraph of text.  See how it is formatted.
\end{quote}
and longer ones.
\begin{quotation}
   This is a longer quotation.  It consists of two
   paragraphs of text, neither of which are
   particularly interesting.

   This is the second paragraph of the quotation.  It
   is just as dull as the first paragraph.
\end{quotation}
Another frequently-displayed structure is a list.
The following is an example of an \emph{itemized}
list.
\begin{itemize}
   \item This is the first item of an itemized list.
         Each item in the list is marked with a ``tick''.
         You don't have to worry about what kind of tick
         mark is used.

   \item This is the second item of the list.  It
         contains another list nested inside it.  The inner
         list is an \emph{enumerated} list.
         \begin{enumerate}
            \item This is the first item of an enumerated
                  list that is nested within the itemized list.

            \item This is the second item of the inner list.
                  \LaTeX\ allows you to nest lists deeper than
                  you really should.
         \end{enumerate}
         This is the rest of the second item of the outer
         list.  It is no more interesting than any other
         part of the item.
   \item This is the third item of the list.
\end{itemize}
You can even display poetry.
\begin{verse}
   There is an environment
    for verse \\             % The \\ command separates lines
   Whose features some poets % within a stanza.
   will curse.

                             % One or more blank lines separate stanzas.

   For instead of making\\
   Them do \emph{all} line breaking, \\
   It allows them to put too many words on a line when they'd rather be
   forced to be terse.
\end{verse}

Mathematical formulas may also be displayed.  A
displayed formula
is
one-line long; multiline
formulas require special formatting instructions.
   \[  \ip{\Gamma}{\psi'} = x'' + y^{2} + z_{i}^{n}\]
Don't start a paragraph with a displayed equation,
nor make one a paragraph by itself.

\end{document}               % End of document.

\section{Bench Huge Section 1205}
% synthetic bench section 1205
% This is a sample LaTeX input file.  (Version of 12 August 2004.)
%
% A '%' character causes TeX to ignore all remaining text on the line,
% and is used for comments like this one.

\documentclass{article}      % Specifies the document class

                             % The preamble begins here.
\title{An Example Document}  % Declares the document's title.
\author{Leslie Lamport}      % Declares the author's name.
\date{January 21, 1994}      % Deleting this command produces today's date.

\newcommand{\ip}[2]{(#1, #2)}
                             % Defines \ip{arg1}{arg2} to mean
                             % (arg1, arg2).

%\newcommand{\ip}[2]{\langle #1 | #2\rangle}
                             % This is an alternative definition of
                             % \ip that is commented out.

\begin{document}             % End of preamble and beginning of text.

\maketitle                   % Produces the title.

This is an example input file.  Comparing it with
the output it generates can show you how to
produce a simple document of your own.

\section{Ordinary Text}      % Produces section heading.  Lower-level
                             % sections are begun with similar
                             % \subsection and \subsubsection commands.

The ends  of words and sentences are marked
  by   spaces. It  doesn't matter how many
spaces    you type; one is as good as 100.  The
end of   a line counts as a space.

One   or more   blank lines denote the  end
of  a paragraph.

Since any number of consecutive spaces are treated
like a single one, the formatting of the input
file makes no difference to
      \LaTeX,                % The \LaTeX command generates the LaTeX logo.
but it makes a difference to you.  When you use
\LaTeX, making your input file as easy to read
as possible will be a great help as you write
your document and when you change it.  This sample
file shows how you can add comments to your own input
file.

Because printing is different from typewriting,
there are a number of things that you have to do
differently when preparing an input file than if
you were just typing the document directly.
Quotation marks like
       ``this''
have to be handled specially, as do quotes within
quotes:
       ``\,`this'            % \, separates the double and single quote.
        is what I just
        wrote, not  `that'\,''.

Dashes come in three sizes: an
       intra-word
dash, a medium dash for number ranges like
       1--2,
and a punctuation
       dash---like
this.

A sentence-ending space should be larger than the
space between words within a sentence.  You
sometimes have to type special commands in
conjunction with punctuation characters to get
this right, as in the following sentence.
       Gnats, gnus, etc.\ all  % `\ ' makes an inter-word space.
       begin with G\@.         % \@ marks end-of-sentence punctuation.
You should check the spaces after periods when
reading your output to make sure you haven't
forgotten any special cases.  Generating an
ellipsis
       \ldots\               % `\ ' is needed after `\ldots' because TeX
                             % ignores spaces after command names like \ldots
                             % made from \ + letters.
                             %
                             % Note how a `%' character causes TeX to ignore
                             % the end of the input line, so these blank lines
                             % do not start a new paragraph.
                             %
with the right spacing around the periods requires
a special command.

\LaTeX\ interprets some common characters as
commands, so you must type special commands to
generate them.  These characters include the
following:
       \$ \& \% \# \{ and \}.

In printing, text is usually emphasized with an
       \emph{italic}
type style.

\begin{em}
   A long segment of text can also be emphasized
   in this way.  Text within such a segment can be
   given \emph{additional} emphasis.
\end{em}

It is sometimes necessary to prevent \LaTeX\ from
breaking a line where it might otherwise do so.
This may be at a space, as between the ``Mr.''\ and
``Jones'' in
       ``Mr.~Jones'',        % ~ produces an unbreakable interword space.
or within a word---especially when the word is a
symbol like
       \mbox{\emph{itemnum}}
that makes little sense when hyphenated across
lines.

Footnotes\footnote{This is an example of a footnote.}
pose no problem.

\LaTeX\ is good at typesetting mathematical formulas
like
       \( x-3y + z = 7 \)
or
       \( a_{1} > x^{2n} + y^{2n} > x' \)
or
       \( \ip{A}{B} = \sum_{i} a_{i} b_{i} \).
The spaces you type in a formula are
ignored.  Remember that a letter like
       $x$                   % $ ... $  and  \( ... \)  are equivalent
is a formula when it denotes a mathematical
symbol, and it should be typed as one.

\section{Displayed Text}

Text is displayed by indenting it from the left
margin.  Quotations are commonly displayed.  There
are short quotations
\begin{quote}
   This is a short quotation.  It consists of a
   single paragraph of text.  See how it is formatted.
\end{quote}
and longer ones.
\begin{quotation}
   This is a longer quotation.  It consists of two
   paragraphs of text, neither of which are
   particularly interesting.

   This is the second paragraph of the quotation.  It
   is just as dull as the first paragraph.
\end{quotation}
Another frequently-displayed structure is a list.
The following is an example of an \emph{itemized}
list.
\begin{itemize}
   \item This is the first item of an itemized list.
         Each item in the list is marked with a ``tick''.
         You don't have to worry about what kind of tick
         mark is used.

   \item This is the second item of the list.  It
         contains another list nested inside it.  The inner
         list is an \emph{enumerated} list.
         \begin{enumerate}
            \item This is the first item of an enumerated
                  list that is nested within the itemized list.

            \item This is the second item of the inner list.
                  \LaTeX\ allows you to nest lists deeper than
                  you really should.
         \end{enumerate}
         This is the rest of the second item of the outer
         list.  It is no more interesting than any other
         part of the item.
   \item This is the third item of the list.
\end{itemize}
You can even display poetry.
\begin{verse}
   There is an environment
    for verse \\             % The \\ command separates lines
   Whose features some poets % within a stanza.
   will curse.

                             % One or more blank lines separate stanzas.

   For instead of making\\
   Them do \emph{all} line breaking, \\
   It allows them to put too many words on a line when they'd rather be
   forced to be terse.
\end{verse}

Mathematical formulas may also be displayed.  A
displayed formula
is
one-line long; multiline
formulas require special formatting instructions.
   \[  \ip{\Gamma}{\psi'} = x'' + y^{2} + z_{i}^{n}\]
Don't start a paragraph with a displayed equation,
nor make one a paragraph by itself.

\end{document}               % End of document.

\section{Bench Huge Section 1206}
% synthetic bench section 1206
% This is a sample LaTeX input file.  (Version of 12 August 2004.)
%
% A '%' character causes TeX to ignore all remaining text on the line,
% and is used for comments like this one.

\documentclass{article}      % Specifies the document class

                             % The preamble begins here.
\title{An Example Document}  % Declares the document's title.
\author{Leslie Lamport}      % Declares the author's name.
\date{January 21, 1994}      % Deleting this command produces today's date.

\newcommand{\ip}[2]{(#1, #2)}
                             % Defines \ip{arg1}{arg2} to mean
                             % (arg1, arg2).

%\newcommand{\ip}[2]{\langle #1 | #2\rangle}
                             % This is an alternative definition of
                             % \ip that is commented out.

\begin{document}             % End of preamble and beginning of text.

\maketitle                   % Produces the title.

This is an example input file.  Comparing it with
the output it generates can show you how to
produce a simple document of your own.

\section{Ordinary Text}      % Produces section heading.  Lower-level
                             % sections are begun with similar
                             % \subsection and \subsubsection commands.

The ends  of words and sentences are marked
  by   spaces. It  doesn't matter how many
spaces    you type; one is as good as 100.  The
end of   a line counts as a space.

One   or more   blank lines denote the  end
of  a paragraph.

Since any number of consecutive spaces are treated
like a single one, the formatting of the input
file makes no difference to
      \LaTeX,                % The \LaTeX command generates the LaTeX logo.
but it makes a difference to you.  When you use
\LaTeX, making your input file as easy to read
as possible will be a great help as you write
your document and when you change it.  This sample
file shows how you can add comments to your own input
file.

Because printing is different from typewriting,
there are a number of things that you have to do
differently when preparing an input file than if
you were just typing the document directly.
Quotation marks like
       ``this''
have to be handled specially, as do quotes within
quotes:
       ``\,`this'            % \, separates the double and single quote.
        is what I just
        wrote, not  `that'\,''.

Dashes come in three sizes: an
       intra-word
dash, a medium dash for number ranges like
       1--2,
and a punctuation
       dash---like
this.

A sentence-ending space should be larger than the
space between words within a sentence.  You
sometimes have to type special commands in
conjunction with punctuation characters to get
this right, as in the following sentence.
       Gnats, gnus, etc.\ all  % `\ ' makes an inter-word space.
       begin with G\@.         % \@ marks end-of-sentence punctuation.
You should check the spaces after periods when
reading your output to make sure you haven't
forgotten any special cases.  Generating an
ellipsis
       \ldots\               % `\ ' is needed after `\ldots' because TeX
                             % ignores spaces after command names like \ldots
                             % made from \ + letters.
                             %
                             % Note how a `%' character causes TeX to ignore
                             % the end of the input line, so these blank lines
                             % do not start a new paragraph.
                             %
with the right spacing around the periods requires
a special command.

\LaTeX\ interprets some common characters as
commands, so you must type special commands to
generate them.  These characters include the
following:
       \$ \& \% \# \{ and \}.

In printing, text is usually emphasized with an
       \emph{italic}
type style.

\begin{em}
   A long segment of text can also be emphasized
   in this way.  Text within such a segment can be
   given \emph{additional} emphasis.
\end{em}

It is sometimes necessary to prevent \LaTeX\ from
breaking a line where it might otherwise do so.
This may be at a space, as between the ``Mr.''\ and
``Jones'' in
       ``Mr.~Jones'',        % ~ produces an unbreakable interword space.
or within a word---especially when the word is a
symbol like
       \mbox{\emph{itemnum}}
that makes little sense when hyphenated across
lines.

Footnotes\footnote{This is an example of a footnote.}
pose no problem.

\LaTeX\ is good at typesetting mathematical formulas
like
       \( x-3y + z = 7 \)
or
       \( a_{1} > x^{2n} + y^{2n} > x' \)
or
       \( \ip{A}{B} = \sum_{i} a_{i} b_{i} \).
The spaces you type in a formula are
ignored.  Remember that a letter like
       $x$                   % $ ... $  and  \( ... \)  are equivalent
is a formula when it denotes a mathematical
symbol, and it should be typed as one.

\section{Displayed Text}

Text is displayed by indenting it from the left
margin.  Quotations are commonly displayed.  There
are short quotations
\begin{quote}
   This is a short quotation.  It consists of a
   single paragraph of text.  See how it is formatted.
\end{quote}
and longer ones.
\begin{quotation}
   This is a longer quotation.  It consists of two
   paragraphs of text, neither of which are
   particularly interesting.

   This is the second paragraph of the quotation.  It
   is just as dull as the first paragraph.
\end{quotation}
Another frequently-displayed structure is a list.
The following is an example of an \emph{itemized}
list.
\begin{itemize}
   \item This is the first item of an itemized list.
         Each item in the list is marked with a ``tick''.
         You don't have to worry about what kind of tick
         mark is used.

   \item This is the second item of the list.  It
         contains another list nested inside it.  The inner
         list is an \emph{enumerated} list.
         \begin{enumerate}
            \item This is the first item of an enumerated
                  list that is nested within the itemized list.

            \item This is the second item of the inner list.
                  \LaTeX\ allows you to nest lists deeper than
                  you really should.
         \end{enumerate}
         This is the rest of the second item of the outer
         list.  It is no more interesting than any other
         part of the item.
   \item This is the third item of the list.
\end{itemize}
You can even display poetry.
\begin{verse}
   There is an environment
    for verse \\             % The \\ command separates lines
   Whose features some poets % within a stanza.
   will curse.

                             % One or more blank lines separate stanzas.

   For instead of making\\
   Them do \emph{all} line breaking, \\
   It allows them to put too many words on a line when they'd rather be
   forced to be terse.
\end{verse}

Mathematical formulas may also be displayed.  A
displayed formula
is
one-line long; multiline
formulas require special formatting instructions.
   \[  \ip{\Gamma}{\psi'} = x'' + y^{2} + z_{i}^{n}\]
Don't start a paragraph with a displayed equation,
nor make one a paragraph by itself.

\end{document}               % End of document.

\section{Bench Huge Section 1207}
% synthetic bench section 1207
% This is a sample LaTeX input file.  (Version of 12 August 2004.)
%
% A '%' character causes TeX to ignore all remaining text on the line,
% and is used for comments like this one.

\documentclass{article}      % Specifies the document class

                             % The preamble begins here.
\title{An Example Document}  % Declares the document's title.
\author{Leslie Lamport}      % Declares the author's name.
\date{January 21, 1994}      % Deleting this command produces today's date.

\newcommand{\ip}[2]{(#1, #2)}
                             % Defines \ip{arg1}{arg2} to mean
                             % (arg1, arg2).

%\newcommand{\ip}[2]{\langle #1 | #2\rangle}
                             % This is an alternative definition of
                             % \ip that is commented out.

\begin{document}             % End of preamble and beginning of text.

\maketitle                   % Produces the title.

This is an example input file.  Comparing it with
the output it generates can show you how to
produce a simple document of your own.

\section{Ordinary Text}      % Produces section heading.  Lower-level
                             % sections are begun with similar
                             % \subsection and \subsubsection commands.

The ends  of words and sentences are marked
  by   spaces. It  doesn't matter how many
spaces    you type; one is as good as 100.  The
end of   a line counts as a space.

One   or more   blank lines denote the  end
of  a paragraph.

Since any number of consecutive spaces are treated
like a single one, the formatting of the input
file makes no difference to
      \LaTeX,                % The \LaTeX command generates the LaTeX logo.
but it makes a difference to you.  When you use
\LaTeX, making your input file as easy to read
as possible will be a great help as you write
your document and when you change it.  This sample
file shows how you can add comments to your own input
file.

Because printing is different from typewriting,
there are a number of things that you have to do
differently when preparing an input file than if
you were just typing the document directly.
Quotation marks like
       ``this''
have to be handled specially, as do quotes within
quotes:
       ``\,`this'            % \, separates the double and single quote.
        is what I just
        wrote, not  `that'\,''.

Dashes come in three sizes: an
       intra-word
dash, a medium dash for number ranges like
       1--2,
and a punctuation
       dash---like
this.

A sentence-ending space should be larger than the
space between words within a sentence.  You
sometimes have to type special commands in
conjunction with punctuation characters to get
this right, as in the following sentence.
       Gnats, gnus, etc.\ all  % `\ ' makes an inter-word space.
       begin with G\@.         % \@ marks end-of-sentence punctuation.
You should check the spaces after periods when
reading your output to make sure you haven't
forgotten any special cases.  Generating an
ellipsis
       \ldots\               % `\ ' is needed after `\ldots' because TeX
                             % ignores spaces after command names like \ldots
                             % made from \ + letters.
                             %
                             % Note how a `%' character causes TeX to ignore
                             % the end of the input line, so these blank lines
                             % do not start a new paragraph.
                             %
with the right spacing around the periods requires
a special command.

\LaTeX\ interprets some common characters as
commands, so you must type special commands to
generate them.  These characters include the
following:
       \$ \& \% \# \{ and \}.

In printing, text is usually emphasized with an
       \emph{italic}
type style.

\begin{em}
   A long segment of text can also be emphasized
   in this way.  Text within such a segment can be
   given \emph{additional} emphasis.
\end{em}

It is sometimes necessary to prevent \LaTeX\ from
breaking a line where it might otherwise do so.
This may be at a space, as between the ``Mr.''\ and
``Jones'' in
       ``Mr.~Jones'',        % ~ produces an unbreakable interword space.
or within a word---especially when the word is a
symbol like
       \mbox{\emph{itemnum}}
that makes little sense when hyphenated across
lines.

Footnotes\footnote{This is an example of a footnote.}
pose no problem.

\LaTeX\ is good at typesetting mathematical formulas
like
       \( x-3y + z = 7 \)
or
       \( a_{1} > x^{2n} + y^{2n} > x' \)
or
       \( \ip{A}{B} = \sum_{i} a_{i} b_{i} \).
The spaces you type in a formula are
ignored.  Remember that a letter like
       $x$                   % $ ... $  and  \( ... \)  are equivalent
is a formula when it denotes a mathematical
symbol, and it should be typed as one.

\section{Displayed Text}

Text is displayed by indenting it from the left
margin.  Quotations are commonly displayed.  There
are short quotations
\begin{quote}
   This is a short quotation.  It consists of a
   single paragraph of text.  See how it is formatted.
\end{quote}
and longer ones.
\begin{quotation}
   This is a longer quotation.  It consists of two
   paragraphs of text, neither of which are
   particularly interesting.

   This is the second paragraph of the quotation.  It
   is just as dull as the first paragraph.
\end{quotation}
Another frequently-displayed structure is a list.
The following is an example of an \emph{itemized}
list.
\begin{itemize}
   \item This is the first item of an itemized list.
         Each item in the list is marked with a ``tick''.
         You don't have to worry about what kind of tick
         mark is used.

   \item This is the second item of the list.  It
         contains another list nested inside it.  The inner
         list is an \emph{enumerated} list.
         \begin{enumerate}
            \item This is the first item of an enumerated
                  list that is nested within the itemized list.

            \item This is the second item of the inner list.
                  \LaTeX\ allows you to nest lists deeper than
                  you really should.
         \end{enumerate}
         This is the rest of the second item of the outer
         list.  It is no more interesting than any other
         part of the item.
   \item This is the third item of the list.
\end{itemize}
You can even display poetry.
\begin{verse}
   There is an environment
    for verse \\             % The \\ command separates lines
   Whose features some poets % within a stanza.
   will curse.

                             % One or more blank lines separate stanzas.

   For instead of making\\
   Them do \emph{all} line breaking, \\
   It allows them to put too many words on a line when they'd rather be
   forced to be terse.
\end{verse}

Mathematical formulas may also be displayed.  A
displayed formula
is
one-line long; multiline
formulas require special formatting instructions.
   \[  \ip{\Gamma}{\psi'} = x'' + y^{2} + z_{i}^{n}\]
Don't start a paragraph with a displayed equation,
nor make one a paragraph by itself.

\end{document}               % End of document.

\section{Bench Huge Section 1208}
% synthetic bench section 1208
% This is a sample LaTeX input file.  (Version of 12 August 2004.)
%
% A '%' character causes TeX to ignore all remaining text on the line,
% and is used for comments like this one.

\documentclass{article}      % Specifies the document class

                             % The preamble begins here.
\title{An Example Document}  % Declares the document's title.
\author{Leslie Lamport}      % Declares the author's name.
\date{January 21, 1994}      % Deleting this command produces today's date.

\newcommand{\ip}[2]{(#1, #2)}
                             % Defines \ip{arg1}{arg2} to mean
                             % (arg1, arg2).

%\newcommand{\ip}[2]{\langle #1 | #2\rangle}
                             % This is an alternative definition of
                             % \ip that is commented out.

\begin{document}             % End of preamble and beginning of text.

\maketitle                   % Produces the title.

This is an example input file.  Comparing it with
the output it generates can show you how to
produce a simple document of your own.

\section{Ordinary Text}      % Produces section heading.  Lower-level
                             % sections are begun with similar
                             % \subsection and \subsubsection commands.

The ends  of words and sentences are marked
  by   spaces. It  doesn't matter how many
spaces    you type; one is as good as 100.  The
end of   a line counts as a space.

One   or more   blank lines denote the  end
of  a paragraph.

Since any number of consecutive spaces are treated
like a single one, the formatting of the input
file makes no difference to
      \LaTeX,                % The \LaTeX command generates the LaTeX logo.
but it makes a difference to you.  When you use
\LaTeX, making your input file as easy to read
as possible will be a great help as you write
your document and when you change it.  This sample
file shows how you can add comments to your own input
file.

Because printing is different from typewriting,
there are a number of things that you have to do
differently when preparing an input file than if
you were just typing the document directly.
Quotation marks like
       ``this''
have to be handled specially, as do quotes within
quotes:
       ``\,`this'            % \, separates the double and single quote.
        is what I just
        wrote, not  `that'\,''.

Dashes come in three sizes: an
       intra-word
dash, a medium dash for number ranges like
       1--2,
and a punctuation
       dash---like
this.

A sentence-ending space should be larger than the
space between words within a sentence.  You
sometimes have to type special commands in
conjunction with punctuation characters to get
this right, as in the following sentence.
       Gnats, gnus, etc.\ all  % `\ ' makes an inter-word space.
       begin with G\@.         % \@ marks end-of-sentence punctuation.
You should check the spaces after periods when
reading your output to make sure you haven't
forgotten any special cases.  Generating an
ellipsis
       \ldots\               % `\ ' is needed after `\ldots' because TeX
                             % ignores spaces after command names like \ldots
                             % made from \ + letters.
                             %
                             % Note how a `%' character causes TeX to ignore
                             % the end of the input line, so these blank lines
                             % do not start a new paragraph.
                             %
with the right spacing around the periods requires
a special command.

\LaTeX\ interprets some common characters as
commands, so you must type special commands to
generate them.  These characters include the
following:
       \$ \& \% \# \{ and \}.

In printing, text is usually emphasized with an
       \emph{italic}
type style.

\begin{em}
   A long segment of text can also be emphasized
   in this way.  Text within such a segment can be
   given \emph{additional} emphasis.
\end{em}

It is sometimes necessary to prevent \LaTeX\ from
breaking a line where it might otherwise do so.
This may be at a space, as between the ``Mr.''\ and
``Jones'' in
       ``Mr.~Jones'',        % ~ produces an unbreakable interword space.
or within a word---especially when the word is a
symbol like
       \mbox{\emph{itemnum}}
that makes little sense when hyphenated across
lines.

Footnotes\footnote{This is an example of a footnote.}
pose no problem.

\LaTeX\ is good at typesetting mathematical formulas
like
       \( x-3y + z = 7 \)
or
       \( a_{1} > x^{2n} + y^{2n} > x' \)
or
       \( \ip{A}{B} = \sum_{i} a_{i} b_{i} \).
The spaces you type in a formula are
ignored.  Remember that a letter like
       $x$                   % $ ... $  and  \( ... \)  are equivalent
is a formula when it denotes a mathematical
symbol, and it should be typed as one.

\section{Displayed Text}

Text is displayed by indenting it from the left
margin.  Quotations are commonly displayed.  There
are short quotations
\begin{quote}
   This is a short quotation.  It consists of a
   single paragraph of text.  See how it is formatted.
\end{quote}
and longer ones.
\begin{quotation}
   This is a longer quotation.  It consists of two
   paragraphs of text, neither of which are
   particularly interesting.

   This is the second paragraph of the quotation.  It
   is just as dull as the first paragraph.
\end{quotation}
Another frequently-displayed structure is a list.
The following is an example of an \emph{itemized}
list.
\begin{itemize}
   \item This is the first item of an itemized list.
         Each item in the list is marked with a ``tick''.
         You don't have to worry about what kind of tick
         mark is used.

   \item This is the second item of the list.  It
         contains another list nested inside it.  The inner
         list is an \emph{enumerated} list.
         \begin{enumerate}
            \item This is the first item of an enumerated
                  list that is nested within the itemized list.

            \item This is the second item of the inner list.
                  \LaTeX\ allows you to nest lists deeper than
                  you really should.
         \end{enumerate}
         This is the rest of the second item of the outer
         list.  It is no more interesting than any other
         part of the item.
   \item This is the third item of the list.
\end{itemize}
You can even display poetry.
\begin{verse}
   There is an environment
    for verse \\             % The \\ command separates lines
   Whose features some poets % within a stanza.
   will curse.

                             % One or more blank lines separate stanzas.

   For instead of making\\
   Them do \emph{all} line breaking, \\
   It allows them to put too many words on a line when they'd rather be
   forced to be terse.
\end{verse}

Mathematical formulas may also be displayed.  A
displayed formula
is
one-line long; multiline
formulas require special formatting instructions.
   \[  \ip{\Gamma}{\psi'} = x'' + y^{2} + z_{i}^{n}\]
Don't start a paragraph with a displayed equation,
nor make one a paragraph by itself.

\end{document}               % End of document.

\section{Bench Huge Section 1209}
% synthetic bench section 1209
% This is a sample LaTeX input file.  (Version of 12 August 2004.)
%
% A '%' character causes TeX to ignore all remaining text on the line,
% and is used for comments like this one.

\documentclass{article}      % Specifies the document class

                             % The preamble begins here.
\title{An Example Document}  % Declares the document's title.
\author{Leslie Lamport}      % Declares the author's name.
\date{January 21, 1994}      % Deleting this command produces today's date.

\newcommand{\ip}[2]{(#1, #2)}
                             % Defines \ip{arg1}{arg2} to mean
                             % (arg1, arg2).

%\newcommand{\ip}[2]{\langle #1 | #2\rangle}
                             % This is an alternative definition of
                             % \ip that is commented out.

\begin{document}             % End of preamble and beginning of text.

\maketitle                   % Produces the title.

This is an example input file.  Comparing it with
the output it generates can show you how to
produce a simple document of your own.

\section{Ordinary Text}      % Produces section heading.  Lower-level
                             % sections are begun with similar
                             % \subsection and \subsubsection commands.

The ends  of words and sentences are marked
  by   spaces. It  doesn't matter how many
spaces    you type; one is as good as 100.  The
end of   a line counts as a space.

One   or more   blank lines denote the  end
of  a paragraph.

Since any number of consecutive spaces are treated
like a single one, the formatting of the input
file makes no difference to
      \LaTeX,                % The \LaTeX command generates the LaTeX logo.
but it makes a difference to you.  When you use
\LaTeX, making your input file as easy to read
as possible will be a great help as you write
your document and when you change it.  This sample
file shows how you can add comments to your own input
file.

Because printing is different from typewriting,
there are a number of things that you have to do
differently when preparing an input file than if
you were just typing the document directly.
Quotation marks like
       ``this''
have to be handled specially, as do quotes within
quotes:
       ``\,`this'            % \, separates the double and single quote.
        is what I just
        wrote, not  `that'\,''.

Dashes come in three sizes: an
       intra-word
dash, a medium dash for number ranges like
       1--2,
and a punctuation
       dash---like
this.

A sentence-ending space should be larger than the
space between words within a sentence.  You
sometimes have to type special commands in
conjunction with punctuation characters to get
this right, as in the following sentence.
       Gnats, gnus, etc.\ all  % `\ ' makes an inter-word space.
       begin with G\@.         % \@ marks end-of-sentence punctuation.
You should check the spaces after periods when
reading your output to make sure you haven't
forgotten any special cases.  Generating an
ellipsis
       \ldots\               % `\ ' is needed after `\ldots' because TeX
                             % ignores spaces after command names like \ldots
                             % made from \ + letters.
                             %
                             % Note how a `%' character causes TeX to ignore
                             % the end of the input line, so these blank lines
                             % do not start a new paragraph.
                             %
with the right spacing around the periods requires
a special command.

\LaTeX\ interprets some common characters as
commands, so you must type special commands to
generate them.  These characters include the
following:
       \$ \& \% \# \{ and \}.

In printing, text is usually emphasized with an
       \emph{italic}
type style.

\begin{em}
   A long segment of text can also be emphasized
   in this way.  Text within such a segment can be
   given \emph{additional} emphasis.
\end{em}

It is sometimes necessary to prevent \LaTeX\ from
breaking a line where it might otherwise do so.
This may be at a space, as between the ``Mr.''\ and
``Jones'' in
       ``Mr.~Jones'',        % ~ produces an unbreakable interword space.
or within a word---especially when the word is a
symbol like
       \mbox{\emph{itemnum}}
that makes little sense when hyphenated across
lines.

Footnotes\footnote{This is an example of a footnote.}
pose no problem.

\LaTeX\ is good at typesetting mathematical formulas
like
       \( x-3y + z = 7 \)
or
       \( a_{1} > x^{2n} + y^{2n} > x' \)
or
       \( \ip{A}{B} = \sum_{i} a_{i} b_{i} \).
The spaces you type in a formula are
ignored.  Remember that a letter like
       $x$                   % $ ... $  and  \( ... \)  are equivalent
is a formula when it denotes a mathematical
symbol, and it should be typed as one.

\section{Displayed Text}

Text is displayed by indenting it from the left
margin.  Quotations are commonly displayed.  There
are short quotations
\begin{quote}
   This is a short quotation.  It consists of a
   single paragraph of text.  See how it is formatted.
\end{quote}
and longer ones.
\begin{quotation}
   This is a longer quotation.  It consists of two
   paragraphs of text, neither of which are
   particularly interesting.

   This is the second paragraph of the quotation.  It
   is just as dull as the first paragraph.
\end{quotation}
Another frequently-displayed structure is a list.
The following is an example of an \emph{itemized}
list.
\begin{itemize}
   \item This is the first item of an itemized list.
         Each item in the list is marked with a ``tick''.
         You don't have to worry about what kind of tick
         mark is used.

   \item This is the second item of the list.  It
         contains another list nested inside it.  The inner
         list is an \emph{enumerated} list.
         \begin{enumerate}
            \item This is the first item of an enumerated
                  list that is nested within the itemized list.

            \item This is the second item of the inner list.
                  \LaTeX\ allows you to nest lists deeper than
                  you really should.
         \end{enumerate}
         This is the rest of the second item of the outer
         list.  It is no more interesting than any other
         part of the item.
   \item This is the third item of the list.
\end{itemize}
You can even display poetry.
\begin{verse}
   There is an environment
    for verse \\             % The \\ command separates lines
   Whose features some poets % within a stanza.
   will curse.

                             % One or more blank lines separate stanzas.

   For instead of making\\
   Them do \emph{all} line breaking, \\
   It allows them to put too many words on a line when they'd rather be
   forced to be terse.
\end{verse}

Mathematical formulas may also be displayed.  A
displayed formula
is
one-line long; multiline
formulas require special formatting instructions.
   \[  \ip{\Gamma}{\psi'} = x'' + y^{2} + z_{i}^{n}\]
Don't start a paragraph with a displayed equation,
nor make one a paragraph by itself.

\end{document}               % End of document.

\section{Bench Huge Section 1210}
% synthetic bench section 1210
% This is a sample LaTeX input file.  (Version of 12 August 2004.)
%
% A '%' character causes TeX to ignore all remaining text on the line,
% and is used for comments like this one.

\documentclass{article}      % Specifies the document class

                             % The preamble begins here.
\title{An Example Document}  % Declares the document's title.
\author{Leslie Lamport}      % Declares the author's name.
\date{January 21, 1994}      % Deleting this command produces today's date.

\newcommand{\ip}[2]{(#1, #2)}
                             % Defines \ip{arg1}{arg2} to mean
                             % (arg1, arg2).

%\newcommand{\ip}[2]{\langle #1 | #2\rangle}
                             % This is an alternative definition of
                             % \ip that is commented out.

\begin{document}             % End of preamble and beginning of text.

\maketitle                   % Produces the title.

This is an example input file.  Comparing it with
the output it generates can show you how to
produce a simple document of your own.

\section{Ordinary Text}      % Produces section heading.  Lower-level
                             % sections are begun with similar
                             % \subsection and \subsubsection commands.

The ends  of words and sentences are marked
  by   spaces. It  doesn't matter how many
spaces    you type; one is as good as 100.  The
end of   a line counts as a space.

One   or more   blank lines denote the  end
of  a paragraph.

Since any number of consecutive spaces are treated
like a single one, the formatting of the input
file makes no difference to
      \LaTeX,                % The \LaTeX command generates the LaTeX logo.
but it makes a difference to you.  When you use
\LaTeX, making your input file as easy to read
as possible will be a great help as you write
your document and when you change it.  This sample
file shows how you can add comments to your own input
file.

Because printing is different from typewriting,
there are a number of things that you have to do
differently when preparing an input file than if
you were just typing the document directly.
Quotation marks like
       ``this''
have to be handled specially, as do quotes within
quotes:
       ``\,`this'            % \, separates the double and single quote.
        is what I just
        wrote, not  `that'\,''.

Dashes come in three sizes: an
       intra-word
dash, a medium dash for number ranges like
       1--2,
and a punctuation
       dash---like
this.

A sentence-ending space should be larger than the
space between words within a sentence.  You
sometimes have to type special commands in
conjunction with punctuation characters to get
this right, as in the following sentence.
       Gnats, gnus, etc.\ all  % `\ ' makes an inter-word space.
       begin with G\@.         % \@ marks end-of-sentence punctuation.
You should check the spaces after periods when
reading your output to make sure you haven't
forgotten any special cases.  Generating an
ellipsis
       \ldots\               % `\ ' is needed after `\ldots' because TeX
                             % ignores spaces after command names like \ldots
                             % made from \ + letters.
                             %
                             % Note how a `%' character causes TeX to ignore
                             % the end of the input line, so these blank lines
                             % do not start a new paragraph.
                             %
with the right spacing around the periods requires
a special command.

\LaTeX\ interprets some common characters as
commands, so you must type special commands to
generate them.  These characters include the
following:
       \$ \& \% \# \{ and \}.

In printing, text is usually emphasized with an
       \emph{italic}
type style.

\begin{em}
   A long segment of text can also be emphasized
   in this way.  Text within such a segment can be
   given \emph{additional} emphasis.
\end{em}

It is sometimes necessary to prevent \LaTeX\ from
breaking a line where it might otherwise do so.
This may be at a space, as between the ``Mr.''\ and
``Jones'' in
       ``Mr.~Jones'',        % ~ produces an unbreakable interword space.
or within a word---especially when the word is a
symbol like
       \mbox{\emph{itemnum}}
that makes little sense when hyphenated across
lines.

Footnotes\footnote{This is an example of a footnote.}
pose no problem.

\LaTeX\ is good at typesetting mathematical formulas
like
       \( x-3y + z = 7 \)
or
       \( a_{1} > x^{2n} + y^{2n} > x' \)
or
       \( \ip{A}{B} = \sum_{i} a_{i} b_{i} \).
The spaces you type in a formula are
ignored.  Remember that a letter like
       $x$                   % $ ... $  and  \( ... \)  are equivalent
is a formula when it denotes a mathematical
symbol, and it should be typed as one.

\section{Displayed Text}

Text is displayed by indenting it from the left
margin.  Quotations are commonly displayed.  There
are short quotations
\begin{quote}
   This is a short quotation.  It consists of a
   single paragraph of text.  See how it is formatted.
\end{quote}
and longer ones.
\begin{quotation}
   This is a longer quotation.  It consists of two
   paragraphs of text, neither of which are
   particularly interesting.

   This is the second paragraph of the quotation.  It
   is just as dull as the first paragraph.
\end{quotation}
Another frequently-displayed structure is a list.
The following is an example of an \emph{itemized}
list.
\begin{itemize}
   \item This is the first item of an itemized list.
         Each item in the list is marked with a ``tick''.
         You don't have to worry about what kind of tick
         mark is used.

   \item This is the second item of the list.  It
         contains another list nested inside it.  The inner
         list is an \emph{enumerated} list.
         \begin{enumerate}
            \item This is the first item of an enumerated
                  list that is nested within the itemized list.

            \item This is the second item of the inner list.
                  \LaTeX\ allows you to nest lists deeper than
                  you really should.
         \end{enumerate}
         This is the rest of the second item of the outer
         list.  It is no more interesting than any other
         part of the item.
   \item This is the third item of the list.
\end{itemize}
You can even display poetry.
\begin{verse}
   There is an environment
    for verse \\             % The \\ command separates lines
   Whose features some poets % within a stanza.
   will curse.

                             % One or more blank lines separate stanzas.

   For instead of making\\
   Them do \emph{all} line breaking, \\
   It allows them to put too many words on a line when they'd rather be
   forced to be terse.
\end{verse}

Mathematical formulas may also be displayed.  A
displayed formula
is
one-line long; multiline
formulas require special formatting instructions.
   \[  \ip{\Gamma}{\psi'} = x'' + y^{2} + z_{i}^{n}\]
Don't start a paragraph with a displayed equation,
nor make one a paragraph by itself.

\end{document}               % End of document.

\section{Bench Huge Section 1211}
% synthetic bench section 1211
% This is a sample LaTeX input file.  (Version of 12 August 2004.)
%
% A '%' character causes TeX to ignore all remaining text on the line,
% and is used for comments like this one.

\documentclass{article}      % Specifies the document class

                             % The preamble begins here.
\title{An Example Document}  % Declares the document's title.
\author{Leslie Lamport}      % Declares the author's name.
\date{January 21, 1994}      % Deleting this command produces today's date.

\newcommand{\ip}[2]{(#1, #2)}
                             % Defines \ip{arg1}{arg2} to mean
                             % (arg1, arg2).

%\newcommand{\ip}[2]{\langle #1 | #2\rangle}
                             % This is an alternative definition of
                             % \ip that is commented out.

\begin{document}             % End of preamble and beginning of text.

\maketitle                   % Produces the title.

This is an example input file.  Comparing it with
the output it generates can show you how to
produce a simple document of your own.

\section{Ordinary Text}      % Produces section heading.  Lower-level
                             % sections are begun with similar
                             % \subsection and \subsubsection commands.

The ends  of words and sentences are marked
  by   spaces. It  doesn't matter how many
spaces    you type; one is as good as 100.  The
end of   a line counts as a space.

One   or more   blank lines denote the  end
of  a paragraph.

Since any number of consecutive spaces are treated
like a single one, the formatting of the input
file makes no difference to
      \LaTeX,                % The \LaTeX command generates the LaTeX logo.
but it makes a difference to you.  When you use
\LaTeX, making your input file as easy to read
as possible will be a great help as you write
your document and when you change it.  This sample
file shows how you can add comments to your own input
file.

Because printing is different from typewriting,
there are a number of things that you have to do
differently when preparing an input file than if
you were just typing the document directly.
Quotation marks like
       ``this''
have to be handled specially, as do quotes within
quotes:
       ``\,`this'            % \, separates the double and single quote.
        is what I just
        wrote, not  `that'\,''.

Dashes come in three sizes: an
       intra-word
dash, a medium dash for number ranges like
       1--2,
and a punctuation
       dash---like
this.

A sentence-ending space should be larger than the
space between words within a sentence.  You
sometimes have to type special commands in
conjunction with punctuation characters to get
this right, as in the following sentence.
       Gnats, gnus, etc.\ all  % `\ ' makes an inter-word space.
       begin with G\@.         % \@ marks end-of-sentence punctuation.
You should check the spaces after periods when
reading your output to make sure you haven't
forgotten any special cases.  Generating an
ellipsis
       \ldots\               % `\ ' is needed after `\ldots' because TeX
                             % ignores spaces after command names like \ldots
                             % made from \ + letters.
                             %
                             % Note how a `%' character causes TeX to ignore
                             % the end of the input line, so these blank lines
                             % do not start a new paragraph.
                             %
with the right spacing around the periods requires
a special command.

\LaTeX\ interprets some common characters as
commands, so you must type special commands to
generate them.  These characters include the
following:
       \$ \& \% \# \{ and \}.

In printing, text is usually emphasized with an
       \emph{italic}
type style.

\begin{em}
   A long segment of text can also be emphasized
   in this way.  Text within such a segment can be
   given \emph{additional} emphasis.
\end{em}

It is sometimes necessary to prevent \LaTeX\ from
breaking a line where it might otherwise do so.
This may be at a space, as between the ``Mr.''\ and
``Jones'' in
       ``Mr.~Jones'',        % ~ produces an unbreakable interword space.
or within a word---especially when the word is a
symbol like
       \mbox{\emph{itemnum}}
that makes little sense when hyphenated across
lines.

Footnotes\footnote{This is an example of a footnote.}
pose no problem.

\LaTeX\ is good at typesetting mathematical formulas
like
       \( x-3y + z = 7 \)
or
       \( a_{1} > x^{2n} + y^{2n} > x' \)
or
       \( \ip{A}{B} = \sum_{i} a_{i} b_{i} \).
The spaces you type in a formula are
ignored.  Remember that a letter like
       $x$                   % $ ... $  and  \( ... \)  are equivalent
is a formula when it denotes a mathematical
symbol, and it should be typed as one.

\section{Displayed Text}

Text is displayed by indenting it from the left
margin.  Quotations are commonly displayed.  There
are short quotations
\begin{quote}
   This is a short quotation.  It consists of a
   single paragraph of text.  See how it is formatted.
\end{quote}
and longer ones.
\begin{quotation}
   This is a longer quotation.  It consists of two
   paragraphs of text, neither of which are
   particularly interesting.

   This is the second paragraph of the quotation.  It
   is just as dull as the first paragraph.
\end{quotation}
Another frequently-displayed structure is a list.
The following is an example of an \emph{itemized}
list.
\begin{itemize}
   \item This is the first item of an itemized list.
         Each item in the list is marked with a ``tick''.
         You don't have to worry about what kind of tick
         mark is used.

   \item This is the second item of the list.  It
         contains another list nested inside it.  The inner
         list is an \emph{enumerated} list.
         \begin{enumerate}
            \item This is the first item of an enumerated
                  list that is nested within the itemized list.

            \item This is the second item of the inner list.
                  \LaTeX\ allows you to nest lists deeper than
                  you really should.
         \end{enumerate}
         This is the rest of the second item of the outer
         list.  It is no more interesting than any other
         part of the item.
   \item This is the third item of the list.
\end{itemize}
You can even display poetry.
\begin{verse}
   There is an environment
    for verse \\             % The \\ command separates lines
   Whose features some poets % within a stanza.
   will curse.

                             % One or more blank lines separate stanzas.

   For instead of making\\
   Them do \emph{all} line breaking, \\
   It allows them to put too many words on a line when they'd rather be
   forced to be terse.
\end{verse}

Mathematical formulas may also be displayed.  A
displayed formula
is
one-line long; multiline
formulas require special formatting instructions.
   \[  \ip{\Gamma}{\psi'} = x'' + y^{2} + z_{i}^{n}\]
Don't start a paragraph with a displayed equation,
nor make one a paragraph by itself.

\end{document}               % End of document.

\section{Bench Huge Section 1212}
% synthetic bench section 1212
% This is a sample LaTeX input file.  (Version of 12 August 2004.)
%
% A '%' character causes TeX to ignore all remaining text on the line,
% and is used for comments like this one.

\documentclass{article}      % Specifies the document class

                             % The preamble begins here.
\title{An Example Document}  % Declares the document's title.
\author{Leslie Lamport}      % Declares the author's name.
\date{January 21, 1994}      % Deleting this command produces today's date.

\newcommand{\ip}[2]{(#1, #2)}
                             % Defines \ip{arg1}{arg2} to mean
                             % (arg1, arg2).

%\newcommand{\ip}[2]{\langle #1 | #2\rangle}
                             % This is an alternative definition of
                             % \ip that is commented out.

\begin{document}             % End of preamble and beginning of text.

\maketitle                   % Produces the title.

This is an example input file.  Comparing it with
the output it generates can show you how to
produce a simple document of your own.

\section{Ordinary Text}      % Produces section heading.  Lower-level
                             % sections are begun with similar
                             % \subsection and \subsubsection commands.

The ends  of words and sentences are marked
  by   spaces. It  doesn't matter how many
spaces    you type; one is as good as 100.  The
end of   a line counts as a space.

One   or more   blank lines denote the  end
of  a paragraph.

Since any number of consecutive spaces are treated
like a single one, the formatting of the input
file makes no difference to
      \LaTeX,                % The \LaTeX command generates the LaTeX logo.
but it makes a difference to you.  When you use
\LaTeX, making your input file as easy to read
as possible will be a great help as you write
your document and when you change it.  This sample
file shows how you can add comments to your own input
file.

Because printing is different from typewriting,
there are a number of things that you have to do
differently when preparing an input file than if
you were just typing the document directly.
Quotation marks like
       ``this''
have to be handled specially, as do quotes within
quotes:
       ``\,`this'            % \, separates the double and single quote.
        is what I just
        wrote, not  `that'\,''.

Dashes come in three sizes: an
       intra-word
dash, a medium dash for number ranges like
       1--2,
and a punctuation
       dash---like
this.

A sentence-ending space should be larger than the
space between words within a sentence.  You
sometimes have to type special commands in
conjunction with punctuation characters to get
this right, as in the following sentence.
       Gnats, gnus, etc.\ all  % `\ ' makes an inter-word space.
       begin with G\@.         % \@ marks end-of-sentence punctuation.
You should check the spaces after periods when
reading your output to make sure you haven't
forgotten any special cases.  Generating an
ellipsis
       \ldots\               % `\ ' is needed after `\ldots' because TeX
                             % ignores spaces after command names like \ldots
                             % made from \ + letters.
                             %
                             % Note how a `%' character causes TeX to ignore
                             % the end of the input line, so these blank lines
                             % do not start a new paragraph.
                             %
with the right spacing around the periods requires
a special command.

\LaTeX\ interprets some common characters as
commands, so you must type special commands to
generate them.  These characters include the
following:
       \$ \& \% \# \{ and \}.

In printing, text is usually emphasized with an
       \emph{italic}
type style.

\begin{em}
   A long segment of text can also be emphasized
   in this way.  Text within such a segment can be
   given \emph{additional} emphasis.
\end{em}

It is sometimes necessary to prevent \LaTeX\ from
breaking a line where it might otherwise do so.
This may be at a space, as between the ``Mr.''\ and
``Jones'' in
       ``Mr.~Jones'',        % ~ produces an unbreakable interword space.
or within a word---especially when the word is a
symbol like
       \mbox{\emph{itemnum}}
that makes little sense when hyphenated across
lines.

Footnotes\footnote{This is an example of a footnote.}
pose no problem.

\LaTeX\ is good at typesetting mathematical formulas
like
       \( x-3y + z = 7 \)
or
       \( a_{1} > x^{2n} + y^{2n} > x' \)
or
       \( \ip{A}{B} = \sum_{i} a_{i} b_{i} \).
The spaces you type in a formula are
ignored.  Remember that a letter like
       $x$                   % $ ... $  and  \( ... \)  are equivalent
is a formula when it denotes a mathematical
symbol, and it should be typed as one.

\section{Displayed Text}

Text is displayed by indenting it from the left
margin.  Quotations are commonly displayed.  There
are short quotations
\begin{quote}
   This is a short quotation.  It consists of a
   single paragraph of text.  See how it is formatted.
\end{quote}
and longer ones.
\begin{quotation}
   This is a longer quotation.  It consists of two
   paragraphs of text, neither of which are
   particularly interesting.

   This is the second paragraph of the quotation.  It
   is just as dull as the first paragraph.
\end{quotation}
Another frequently-displayed structure is a list.
The following is an example of an \emph{itemized}
list.
\begin{itemize}
   \item This is the first item of an itemized list.
         Each item in the list is marked with a ``tick''.
         You don't have to worry about what kind of tick
         mark is used.

   \item This is the second item of the list.  It
         contains another list nested inside it.  The inner
         list is an \emph{enumerated} list.
         \begin{enumerate}
            \item This is the first item of an enumerated
                  list that is nested within the itemized list.

            \item This is the second item of the inner list.
                  \LaTeX\ allows you to nest lists deeper than
                  you really should.
         \end{enumerate}
         This is the rest of the second item of the outer
         list.  It is no more interesting than any other
         part of the item.
   \item This is the third item of the list.
\end{itemize}
You can even display poetry.
\begin{verse}
   There is an environment
    for verse \\             % The \\ command separates lines
   Whose features some poets % within a stanza.
   will curse.

                             % One or more blank lines separate stanzas.

   For instead of making\\
   Them do \emph{all} line breaking, \\
   It allows them to put too many words on a line when they'd rather be
   forced to be terse.
\end{verse}

Mathematical formulas may also be displayed.  A
displayed formula
is
one-line long; multiline
formulas require special formatting instructions.
   \[  \ip{\Gamma}{\psi'} = x'' + y^{2} + z_{i}^{n}\]
Don't start a paragraph with a displayed equation,
nor make one a paragraph by itself.

\end{document}               % End of document.

\section{Bench Huge Section 1213}
% synthetic bench section 1213
% This is a sample LaTeX input file.  (Version of 12 August 2004.)
%
% A '%' character causes TeX to ignore all remaining text on the line,
% and is used for comments like this one.

\documentclass{article}      % Specifies the document class

                             % The preamble begins here.
\title{An Example Document}  % Declares the document's title.
\author{Leslie Lamport}      % Declares the author's name.
\date{January 21, 1994}      % Deleting this command produces today's date.

\newcommand{\ip}[2]{(#1, #2)}
                             % Defines \ip{arg1}{arg2} to mean
                             % (arg1, arg2).

%\newcommand{\ip}[2]{\langle #1 | #2\rangle}
                             % This is an alternative definition of
                             % \ip that is commented out.

\begin{document}             % End of preamble and beginning of text.

\maketitle                   % Produces the title.

This is an example input file.  Comparing it with
the output it generates can show you how to
produce a simple document of your own.

\section{Ordinary Text}      % Produces section heading.  Lower-level
                             % sections are begun with similar
                             % \subsection and \subsubsection commands.

The ends  of words and sentences are marked
  by   spaces. It  doesn't matter how many
spaces    you type; one is as good as 100.  The
end of   a line counts as a space.

One   or more   blank lines denote the  end
of  a paragraph.

Since any number of consecutive spaces are treated
like a single one, the formatting of the input
file makes no difference to
      \LaTeX,                % The \LaTeX command generates the LaTeX logo.
but it makes a difference to you.  When you use
\LaTeX, making your input file as easy to read
as possible will be a great help as you write
your document and when you change it.  This sample
file shows how you can add comments to your own input
file.

Because printing is different from typewriting,
there are a number of things that you have to do
differently when preparing an input file than if
you were just typing the document directly.
Quotation marks like
       ``this''
have to be handled specially, as do quotes within
quotes:
       ``\,`this'            % \, separates the double and single quote.
        is what I just
        wrote, not  `that'\,''.

Dashes come in three sizes: an
       intra-word
dash, a medium dash for number ranges like
       1--2,
and a punctuation
       dash---like
this.

A sentence-ending space should be larger than the
space between words within a sentence.  You
sometimes have to type special commands in
conjunction with punctuation characters to get
this right, as in the following sentence.
       Gnats, gnus, etc.\ all  % `\ ' makes an inter-word space.
       begin with G\@.         % \@ marks end-of-sentence punctuation.
You should check the spaces after periods when
reading your output to make sure you haven't
forgotten any special cases.  Generating an
ellipsis
       \ldots\               % `\ ' is needed after `\ldots' because TeX
                             % ignores spaces after command names like \ldots
                             % made from \ + letters.
                             %
                             % Note how a `%' character causes TeX to ignore
                             % the end of the input line, so these blank lines
                             % do not start a new paragraph.
                             %
with the right spacing around the periods requires
a special command.

\LaTeX\ interprets some common characters as
commands, so you must type special commands to
generate them.  These characters include the
following:
       \$ \& \% \# \{ and \}.

In printing, text is usually emphasized with an
       \emph{italic}
type style.

\begin{em}
   A long segment of text can also be emphasized
   in this way.  Text within such a segment can be
   given \emph{additional} emphasis.
\end{em}

It is sometimes necessary to prevent \LaTeX\ from
breaking a line where it might otherwise do so.
This may be at a space, as between the ``Mr.''\ and
``Jones'' in
       ``Mr.~Jones'',        % ~ produces an unbreakable interword space.
or within a word---especially when the word is a
symbol like
       \mbox{\emph{itemnum}}
that makes little sense when hyphenated across
lines.

Footnotes\footnote{This is an example of a footnote.}
pose no problem.

\LaTeX\ is good at typesetting mathematical formulas
like
       \( x-3y + z = 7 \)
or
       \( a_{1} > x^{2n} + y^{2n} > x' \)
or
       \( \ip{A}{B} = \sum_{i} a_{i} b_{i} \).
The spaces you type in a formula are
ignored.  Remember that a letter like
       $x$                   % $ ... $  and  \( ... \)  are equivalent
is a formula when it denotes a mathematical
symbol, and it should be typed as one.

\section{Displayed Text}

Text is displayed by indenting it from the left
margin.  Quotations are commonly displayed.  There
are short quotations
\begin{quote}
   This is a short quotation.  It consists of a
   single paragraph of text.  See how it is formatted.
\end{quote}
and longer ones.
\begin{quotation}
   This is a longer quotation.  It consists of two
   paragraphs of text, neither of which are
   particularly interesting.

   This is the second paragraph of the quotation.  It
   is just as dull as the first paragraph.
\end{quotation}
Another frequently-displayed structure is a list.
The following is an example of an \emph{itemized}
list.
\begin{itemize}
   \item This is the first item of an itemized list.
         Each item in the list is marked with a ``tick''.
         You don't have to worry about what kind of tick
         mark is used.

   \item This is the second item of the list.  It
         contains another list nested inside it.  The inner
         list is an \emph{enumerated} list.
         \begin{enumerate}
            \item This is the first item of an enumerated
                  list that is nested within the itemized list.

            \item This is the second item of the inner list.
                  \LaTeX\ allows you to nest lists deeper than
                  you really should.
         \end{enumerate}
         This is the rest of the second item of the outer
         list.  It is no more interesting than any other
         part of the item.
   \item This is the third item of the list.
\end{itemize}
You can even display poetry.
\begin{verse}
   There is an environment
    for verse \\             % The \\ command separates lines
   Whose features some poets % within a stanza.
   will curse.

                             % One or more blank lines separate stanzas.

   For instead of making\\
   Them do \emph{all} line breaking, \\
   It allows them to put too many words on a line when they'd rather be
   forced to be terse.
\end{verse}

Mathematical formulas may also be displayed.  A
displayed formula
is
one-line long; multiline
formulas require special formatting instructions.
   \[  \ip{\Gamma}{\psi'} = x'' + y^{2} + z_{i}^{n}\]
Don't start a paragraph with a displayed equation,
nor make one a paragraph by itself.

\end{document}               % End of document.

\section{Bench Huge Section 1214}
% synthetic bench section 1214
% This is a sample LaTeX input file.  (Version of 12 August 2004.)
%
% A '%' character causes TeX to ignore all remaining text on the line,
% and is used for comments like this one.

\documentclass{article}      % Specifies the document class

                             % The preamble begins here.
\title{An Example Document}  % Declares the document's title.
\author{Leslie Lamport}      % Declares the author's name.
\date{January 21, 1994}      % Deleting this command produces today's date.

\newcommand{\ip}[2]{(#1, #2)}
                             % Defines \ip{arg1}{arg2} to mean
                             % (arg1, arg2).

%\newcommand{\ip}[2]{\langle #1 | #2\rangle}
                             % This is an alternative definition of
                             % \ip that is commented out.

\begin{document}             % End of preamble and beginning of text.

\maketitle                   % Produces the title.

This is an example input file.  Comparing it with
the output it generates can show you how to
produce a simple document of your own.

\section{Ordinary Text}      % Produces section heading.  Lower-level
                             % sections are begun with similar
                             % \subsection and \subsubsection commands.

The ends  of words and sentences are marked
  by   spaces. It  doesn't matter how many
spaces    you type; one is as good as 100.  The
end of   a line counts as a space.

One   or more   blank lines denote the  end
of  a paragraph.

Since any number of consecutive spaces are treated
like a single one, the formatting of the input
file makes no difference to
      \LaTeX,                % The \LaTeX command generates the LaTeX logo.
but it makes a difference to you.  When you use
\LaTeX, making your input file as easy to read
as possible will be a great help as you write
your document and when you change it.  This sample
file shows how you can add comments to your own input
file.

Because printing is different from typewriting,
there are a number of things that you have to do
differently when preparing an input file than if
you were just typing the document directly.
Quotation marks like
       ``this''
have to be handled specially, as do quotes within
quotes:
       ``\,`this'            % \, separates the double and single quote.
        is what I just
        wrote, not  `that'\,''.

Dashes come in three sizes: an
       intra-word
dash, a medium dash for number ranges like
       1--2,
and a punctuation
       dash---like
this.

A sentence-ending space should be larger than the
space between words within a sentence.  You
sometimes have to type special commands in
conjunction with punctuation characters to get
this right, as in the following sentence.
       Gnats, gnus, etc.\ all  % `\ ' makes an inter-word space.
       begin with G\@.         % \@ marks end-of-sentence punctuation.
You should check the spaces after periods when
reading your output to make sure you haven't
forgotten any special cases.  Generating an
ellipsis
       \ldots\               % `\ ' is needed after `\ldots' because TeX
                             % ignores spaces after command names like \ldots
                             % made from \ + letters.
                             %
                             % Note how a `%' character causes TeX to ignore
                             % the end of the input line, so these blank lines
                             % do not start a new paragraph.
                             %
with the right spacing around the periods requires
a special command.

\LaTeX\ interprets some common characters as
commands, so you must type special commands to
generate them.  These characters include the
following:
       \$ \& \% \# \{ and \}.

In printing, text is usually emphasized with an
       \emph{italic}
type style.

\begin{em}
   A long segment of text can also be emphasized
   in this way.  Text within such a segment can be
   given \emph{additional} emphasis.
\end{em}

It is sometimes necessary to prevent \LaTeX\ from
breaking a line where it might otherwise do so.
This may be at a space, as between the ``Mr.''\ and
``Jones'' in
       ``Mr.~Jones'',        % ~ produces an unbreakable interword space.
or within a word---especially when the word is a
symbol like
       \mbox{\emph{itemnum}}
that makes little sense when hyphenated across
lines.

Footnotes\footnote{This is an example of a footnote.}
pose no problem.

\LaTeX\ is good at typesetting mathematical formulas
like
       \( x-3y + z = 7 \)
or
       \( a_{1} > x^{2n} + y^{2n} > x' \)
or
       \( \ip{A}{B} = \sum_{i} a_{i} b_{i} \).
The spaces you type in a formula are
ignored.  Remember that a letter like
       $x$                   % $ ... $  and  \( ... \)  are equivalent
is a formula when it denotes a mathematical
symbol, and it should be typed as one.

\section{Displayed Text}

Text is displayed by indenting it from the left
margin.  Quotations are commonly displayed.  There
are short quotations
\begin{quote}
   This is a short quotation.  It consists of a
   single paragraph of text.  See how it is formatted.
\end{quote}
and longer ones.
\begin{quotation}
   This is a longer quotation.  It consists of two
   paragraphs of text, neither of which are
   particularly interesting.

   This is the second paragraph of the quotation.  It
   is just as dull as the first paragraph.
\end{quotation}
Another frequently-displayed structure is a list.
The following is an example of an \emph{itemized}
list.
\begin{itemize}
   \item This is the first item of an itemized list.
         Each item in the list is marked with a ``tick''.
         You don't have to worry about what kind of tick
         mark is used.

   \item This is the second item of the list.  It
         contains another list nested inside it.  The inner
         list is an \emph{enumerated} list.
         \begin{enumerate}
            \item This is the first item of an enumerated
                  list that is nested within the itemized list.

            \item This is the second item of the inner list.
                  \LaTeX\ allows you to nest lists deeper than
                  you really should.
         \end{enumerate}
         This is the rest of the second item of the outer
         list.  It is no more interesting than any other
         part of the item.
   \item This is the third item of the list.
\end{itemize}
You can even display poetry.
\begin{verse}
   There is an environment
    for verse \\             % The \\ command separates lines
   Whose features some poets % within a stanza.
   will curse.

                             % One or more blank lines separate stanzas.

   For instead of making\\
   Them do \emph{all} line breaking, \\
   It allows them to put too many words on a line when they'd rather be
   forced to be terse.
\end{verse}

Mathematical formulas may also be displayed.  A
displayed formula
is
one-line long; multiline
formulas require special formatting instructions.
   \[  \ip{\Gamma}{\psi'} = x'' + y^{2} + z_{i}^{n}\]
Don't start a paragraph with a displayed equation,
nor make one a paragraph by itself.

\end{document}               % End of document.

\section{Bench Huge Section 1215}
% synthetic bench section 1215
% This is a sample LaTeX input file.  (Version of 12 August 2004.)
%
% A '%' character causes TeX to ignore all remaining text on the line,
% and is used for comments like this one.

\documentclass{article}      % Specifies the document class

                             % The preamble begins here.
\title{An Example Document}  % Declares the document's title.
\author{Leslie Lamport}      % Declares the author's name.
\date{January 21, 1994}      % Deleting this command produces today's date.

\newcommand{\ip}[2]{(#1, #2)}
                             % Defines \ip{arg1}{arg2} to mean
                             % (arg1, arg2).

%\newcommand{\ip}[2]{\langle #1 | #2\rangle}
                             % This is an alternative definition of
                             % \ip that is commented out.

\begin{document}             % End of preamble and beginning of text.

\maketitle                   % Produces the title.

This is an example input file.  Comparing it with
the output it generates can show you how to
produce a simple document of your own.

\section{Ordinary Text}      % Produces section heading.  Lower-level
                             % sections are begun with similar
                             % \subsection and \subsubsection commands.

The ends  of words and sentences are marked
  by   spaces. It  doesn't matter how many
spaces    you type; one is as good as 100.  The
end of   a line counts as a space.

One   or more   blank lines denote the  end
of  a paragraph.

Since any number of consecutive spaces are treated
like a single one, the formatting of the input
file makes no difference to
      \LaTeX,                % The \LaTeX command generates the LaTeX logo.
but it makes a difference to you.  When you use
\LaTeX, making your input file as easy to read
as possible will be a great help as you write
your document and when you change it.  This sample
file shows how you can add comments to your own input
file.

Because printing is different from typewriting,
there are a number of things that you have to do
differently when preparing an input file than if
you were just typing the document directly.
Quotation marks like
       ``this''
have to be handled specially, as do quotes within
quotes:
       ``\,`this'            % \, separates the double and single quote.
        is what I just
        wrote, not  `that'\,''.

Dashes come in three sizes: an
       intra-word
dash, a medium dash for number ranges like
       1--2,
and a punctuation
       dash---like
this.

A sentence-ending space should be larger than the
space between words within a sentence.  You
sometimes have to type special commands in
conjunction with punctuation characters to get
this right, as in the following sentence.
       Gnats, gnus, etc.\ all  % `\ ' makes an inter-word space.
       begin with G\@.         % \@ marks end-of-sentence punctuation.
You should check the spaces after periods when
reading your output to make sure you haven't
forgotten any special cases.  Generating an
ellipsis
       \ldots\               % `\ ' is needed after `\ldots' because TeX
                             % ignores spaces after command names like \ldots
                             % made from \ + letters.
                             %
                             % Note how a `%' character causes TeX to ignore
                             % the end of the input line, so these blank lines
                             % do not start a new paragraph.
                             %
with the right spacing around the periods requires
a special command.

\LaTeX\ interprets some common characters as
commands, so you must type special commands to
generate them.  These characters include the
following:
       \$ \& \% \# \{ and \}.

In printing, text is usually emphasized with an
       \emph{italic}
type style.

\begin{em}
   A long segment of text can also be emphasized
   in this way.  Text within such a segment can be
   given \emph{additional} emphasis.
\end{em}

It is sometimes necessary to prevent \LaTeX\ from
breaking a line where it might otherwise do so.
This may be at a space, as between the ``Mr.''\ and
``Jones'' in
       ``Mr.~Jones'',        % ~ produces an unbreakable interword space.
or within a word---especially when the word is a
symbol like
       \mbox{\emph{itemnum}}
that makes little sense when hyphenated across
lines.

Footnotes\footnote{This is an example of a footnote.}
pose no problem.

\LaTeX\ is good at typesetting mathematical formulas
like
       \( x-3y + z = 7 \)
or
       \( a_{1} > x^{2n} + y^{2n} > x' \)
or
       \( \ip{A}{B} = \sum_{i} a_{i} b_{i} \).
The spaces you type in a formula are
ignored.  Remember that a letter like
       $x$                   % $ ... $  and  \( ... \)  are equivalent
is a formula when it denotes a mathematical
symbol, and it should be typed as one.

\section{Displayed Text}

Text is displayed by indenting it from the left
margin.  Quotations are commonly displayed.  There
are short quotations
\begin{quote}
   This is a short quotation.  It consists of a
   single paragraph of text.  See how it is formatted.
\end{quote}
and longer ones.
\begin{quotation}
   This is a longer quotation.  It consists of two
   paragraphs of text, neither of which are
   particularly interesting.

   This is the second paragraph of the quotation.  It
   is just as dull as the first paragraph.
\end{quotation}
Another frequently-displayed structure is a list.
The following is an example of an \emph{itemized}
list.
\begin{itemize}
   \item This is the first item of an itemized list.
         Each item in the list is marked with a ``tick''.
         You don't have to worry about what kind of tick
         mark is used.

   \item This is the second item of the list.  It
         contains another list nested inside it.  The inner
         list is an \emph{enumerated} list.
         \begin{enumerate}
            \item This is the first item of an enumerated
                  list that is nested within the itemized list.

            \item This is the second item of the inner list.
                  \LaTeX\ allows you to nest lists deeper than
                  you really should.
         \end{enumerate}
         This is the rest of the second item of the outer
         list.  It is no more interesting than any other
         part of the item.
   \item This is the third item of the list.
\end{itemize}
You can even display poetry.
\begin{verse}
   There is an environment
    for verse \\             % The \\ command separates lines
   Whose features some poets % within a stanza.
   will curse.

                             % One or more blank lines separate stanzas.

   For instead of making\\
   Them do \emph{all} line breaking, \\
   It allows them to put too many words on a line when they'd rather be
   forced to be terse.
\end{verse}

Mathematical formulas may also be displayed.  A
displayed formula
is
one-line long; multiline
formulas require special formatting instructions.
   \[  \ip{\Gamma}{\psi'} = x'' + y^{2} + z_{i}^{n}\]
Don't start a paragraph with a displayed equation,
nor make one a paragraph by itself.

\end{document}               % End of document.

\section{Bench Huge Section 1216}
% synthetic bench section 1216
% This is a sample LaTeX input file.  (Version of 12 August 2004.)
%
% A '%' character causes TeX to ignore all remaining text on the line,
% and is used for comments like this one.

\documentclass{article}      % Specifies the document class

                             % The preamble begins here.
\title{An Example Document}  % Declares the document's title.
\author{Leslie Lamport}      % Declares the author's name.
\date{January 21, 1994}      % Deleting this command produces today's date.

\newcommand{\ip}[2]{(#1, #2)}
                             % Defines \ip{arg1}{arg2} to mean
                             % (arg1, arg2).

%\newcommand{\ip}[2]{\langle #1 | #2\rangle}
                             % This is an alternative definition of
                             % \ip that is commented out.

\begin{document}             % End of preamble and beginning of text.

\maketitle                   % Produces the title.

This is an example input file.  Comparing it with
the output it generates can show you how to
produce a simple document of your own.

\section{Ordinary Text}      % Produces section heading.  Lower-level
                             % sections are begun with similar
                             % \subsection and \subsubsection commands.

The ends  of words and sentences are marked
  by   spaces. It  doesn't matter how many
spaces    you type; one is as good as 100.  The
end of   a line counts as a space.

One   or more   blank lines denote the  end
of  a paragraph.

Since any number of consecutive spaces are treated
like a single one, the formatting of the input
file makes no difference to
      \LaTeX,                % The \LaTeX command generates the LaTeX logo.
but it makes a difference to you.  When you use
\LaTeX, making your input file as easy to read
as possible will be a great help as you write
your document and when you change it.  This sample
file shows how you can add comments to your own input
file.

Because printing is different from typewriting,
there are a number of things that you have to do
differently when preparing an input file than if
you were just typing the document directly.
Quotation marks like
       ``this''
have to be handled specially, as do quotes within
quotes:
       ``\,`this'            % \, separates the double and single quote.
        is what I just
        wrote, not  `that'\,''.

Dashes come in three sizes: an
       intra-word
dash, a medium dash for number ranges like
       1--2,
and a punctuation
       dash---like
this.

A sentence-ending space should be larger than the
space between words within a sentence.  You
sometimes have to type special commands in
conjunction with punctuation characters to get
this right, as in the following sentence.
       Gnats, gnus, etc.\ all  % `\ ' makes an inter-word space.
       begin with G\@.         % \@ marks end-of-sentence punctuation.
You should check the spaces after periods when
reading your output to make sure you haven't
forgotten any special cases.  Generating an
ellipsis
       \ldots\               % `\ ' is needed after `\ldots' because TeX
                             % ignores spaces after command names like \ldots
                             % made from \ + letters.
                             %
                             % Note how a `%' character causes TeX to ignore
                             % the end of the input line, so these blank lines
                             % do not start a new paragraph.
                             %
with the right spacing around the periods requires
a special command.

\LaTeX\ interprets some common characters as
commands, so you must type special commands to
generate them.  These characters include the
following:
       \$ \& \% \# \{ and \}.

In printing, text is usually emphasized with an
       \emph{italic}
type style.

\begin{em}
   A long segment of text can also be emphasized
   in this way.  Text within such a segment can be
   given \emph{additional} emphasis.
\end{em}

It is sometimes necessary to prevent \LaTeX\ from
breaking a line where it might otherwise do so.
This may be at a space, as between the ``Mr.''\ and
``Jones'' in
       ``Mr.~Jones'',        % ~ produces an unbreakable interword space.
or within a word---especially when the word is a
symbol like
       \mbox{\emph{itemnum}}
that makes little sense when hyphenated across
lines.

Footnotes\footnote{This is an example of a footnote.}
pose no problem.

\LaTeX\ is good at typesetting mathematical formulas
like
       \( x-3y + z = 7 \)
or
       \( a_{1} > x^{2n} + y^{2n} > x' \)
or
       \( \ip{A}{B} = \sum_{i} a_{i} b_{i} \).
The spaces you type in a formula are
ignored.  Remember that a letter like
       $x$                   % $ ... $  and  \( ... \)  are equivalent
is a formula when it denotes a mathematical
symbol, and it should be typed as one.

\section{Displayed Text}

Text is displayed by indenting it from the left
margin.  Quotations are commonly displayed.  There
are short quotations
\begin{quote}
   This is a short quotation.  It consists of a
   single paragraph of text.  See how it is formatted.
\end{quote}
and longer ones.
\begin{quotation}
   This is a longer quotation.  It consists of two
   paragraphs of text, neither of which are
   particularly interesting.

   This is the second paragraph of the quotation.  It
   is just as dull as the first paragraph.
\end{quotation}
Another frequently-displayed structure is a list.
The following is an example of an \emph{itemized}
list.
\begin{itemize}
   \item This is the first item of an itemized list.
         Each item in the list is marked with a ``tick''.
         You don't have to worry about what kind of tick
         mark is used.

   \item This is the second item of the list.  It
         contains another list nested inside it.  The inner
         list is an \emph{enumerated} list.
         \begin{enumerate}
            \item This is the first item of an enumerated
                  list that is nested within the itemized list.

            \item This is the second item of the inner list.
                  \LaTeX\ allows you to nest lists deeper than
                  you really should.
         \end{enumerate}
         This is the rest of the second item of the outer
         list.  It is no more interesting than any other
         part of the item.
   \item This is the third item of the list.
\end{itemize}
You can even display poetry.
\begin{verse}
   There is an environment
    for verse \\             % The \\ command separates lines
   Whose features some poets % within a stanza.
   will curse.

                             % One or more blank lines separate stanzas.

   For instead of making\\
   Them do \emph{all} line breaking, \\
   It allows them to put too many words on a line when they'd rather be
   forced to be terse.
\end{verse}

Mathematical formulas may also be displayed.  A
displayed formula
is
one-line long; multiline
formulas require special formatting instructions.
   \[  \ip{\Gamma}{\psi'} = x'' + y^{2} + z_{i}^{n}\]
Don't start a paragraph with a displayed equation,
nor make one a paragraph by itself.

\end{document}               % End of document.

\section{Bench Huge Section 1217}
% synthetic bench section 1217
% This is a sample LaTeX input file.  (Version of 12 August 2004.)
%
% A '%' character causes TeX to ignore all remaining text on the line,
% and is used for comments like this one.

\documentclass{article}      % Specifies the document class

                             % The preamble begins here.
\title{An Example Document}  % Declares the document's title.
\author{Leslie Lamport}      % Declares the author's name.
\date{January 21, 1994}      % Deleting this command produces today's date.

\newcommand{\ip}[2]{(#1, #2)}
                             % Defines \ip{arg1}{arg2} to mean
                             % (arg1, arg2).

%\newcommand{\ip}[2]{\langle #1 | #2\rangle}
                             % This is an alternative definition of
                             % \ip that is commented out.

\begin{document}             % End of preamble and beginning of text.

\maketitle                   % Produces the title.

This is an example input file.  Comparing it with
the output it generates can show you how to
produce a simple document of your own.

\section{Ordinary Text}      % Produces section heading.  Lower-level
                             % sections are begun with similar
                             % \subsection and \subsubsection commands.

The ends  of words and sentences are marked
  by   spaces. It  doesn't matter how many
spaces    you type; one is as good as 100.  The
end of   a line counts as a space.

One   or more   blank lines denote the  end
of  a paragraph.

Since any number of consecutive spaces are treated
like a single one, the formatting of the input
file makes no difference to
      \LaTeX,                % The \LaTeX command generates the LaTeX logo.
but it makes a difference to you.  When you use
\LaTeX, making your input file as easy to read
as possible will be a great help as you write
your document and when you change it.  This sample
file shows how you can add comments to your own input
file.

Because printing is different from typewriting,
there are a number of things that you have to do
differently when preparing an input file than if
you were just typing the document directly.
Quotation marks like
       ``this''
have to be handled specially, as do quotes within
quotes:
       ``\,`this'            % \, separates the double and single quote.
        is what I just
        wrote, not  `that'\,''.

Dashes come in three sizes: an
       intra-word
dash, a medium dash for number ranges like
       1--2,
and a punctuation
       dash---like
this.

A sentence-ending space should be larger than the
space between words within a sentence.  You
sometimes have to type special commands in
conjunction with punctuation characters to get
this right, as in the following sentence.
       Gnats, gnus, etc.\ all  % `\ ' makes an inter-word space.
       begin with G\@.         % \@ marks end-of-sentence punctuation.
You should check the spaces after periods when
reading your output to make sure you haven't
forgotten any special cases.  Generating an
ellipsis
       \ldots\               % `\ ' is needed after `\ldots' because TeX
                             % ignores spaces after command names like \ldots
                             % made from \ + letters.
                             %
                             % Note how a `%' character causes TeX to ignore
                             % the end of the input line, so these blank lines
                             % do not start a new paragraph.
                             %
with the right spacing around the periods requires
a special command.

\LaTeX\ interprets some common characters as
commands, so you must type special commands to
generate them.  These characters include the
following:
       \$ \& \% \# \{ and \}.

In printing, text is usually emphasized with an
       \emph{italic}
type style.

\begin{em}
   A long segment of text can also be emphasized
   in this way.  Text within such a segment can be
   given \emph{additional} emphasis.
\end{em}

It is sometimes necessary to prevent \LaTeX\ from
breaking a line where it might otherwise do so.
This may be at a space, as between the ``Mr.''\ and
``Jones'' in
       ``Mr.~Jones'',        % ~ produces an unbreakable interword space.
or within a word---especially when the word is a
symbol like
       \mbox{\emph{itemnum}}
that makes little sense when hyphenated across
lines.

Footnotes\footnote{This is an example of a footnote.}
pose no problem.

\LaTeX\ is good at typesetting mathematical formulas
like
       \( x-3y + z = 7 \)
or
       \( a_{1} > x^{2n} + y^{2n} > x' \)
or
       \( \ip{A}{B} = \sum_{i} a_{i} b_{i} \).
The spaces you type in a formula are
ignored.  Remember that a letter like
       $x$                   % $ ... $  and  \( ... \)  are equivalent
is a formula when it denotes a mathematical
symbol, and it should be typed as one.

\section{Displayed Text}

Text is displayed by indenting it from the left
margin.  Quotations are commonly displayed.  There
are short quotations
\begin{quote}
   This is a short quotation.  It consists of a
   single paragraph of text.  See how it is formatted.
\end{quote}
and longer ones.
\begin{quotation}
   This is a longer quotation.  It consists of two
   paragraphs of text, neither of which are
   particularly interesting.

   This is the second paragraph of the quotation.  It
   is just as dull as the first paragraph.
\end{quotation}
Another frequently-displayed structure is a list.
The following is an example of an \emph{itemized}
list.
\begin{itemize}
   \item This is the first item of an itemized list.
         Each item in the list is marked with a ``tick''.
         You don't have to worry about what kind of tick
         mark is used.

   \item This is the second item of the list.  It
         contains another list nested inside it.  The inner
         list is an \emph{enumerated} list.
         \begin{enumerate}
            \item This is the first item of an enumerated
                  list that is nested within the itemized list.

            \item This is the second item of the inner list.
                  \LaTeX\ allows you to nest lists deeper than
                  you really should.
         \end{enumerate}
         This is the rest of the second item of the outer
         list.  It is no more interesting than any other
         part of the item.
   \item This is the third item of the list.
\end{itemize}
You can even display poetry.
\begin{verse}
   There is an environment
    for verse \\             % The \\ command separates lines
   Whose features some poets % within a stanza.
   will curse.

                             % One or more blank lines separate stanzas.

   For instead of making\\
   Them do \emph{all} line breaking, \\
   It allows them to put too many words on a line when they'd rather be
   forced to be terse.
\end{verse}

Mathematical formulas may also be displayed.  A
displayed formula
is
one-line long; multiline
formulas require special formatting instructions.
   \[  \ip{\Gamma}{\psi'} = x'' + y^{2} + z_{i}^{n}\]
Don't start a paragraph with a displayed equation,
nor make one a paragraph by itself.

\end{document}               % End of document.

\section{Bench Huge Section 1218}
% synthetic bench section 1218
% This is a sample LaTeX input file.  (Version of 12 August 2004.)
%
% A '%' character causes TeX to ignore all remaining text on the line,
% and is used for comments like this one.

\documentclass{article}      % Specifies the document class

                             % The preamble begins here.
\title{An Example Document}  % Declares the document's title.
\author{Leslie Lamport}      % Declares the author's name.
\date{January 21, 1994}      % Deleting this command produces today's date.

\newcommand{\ip}[2]{(#1, #2)}
                             % Defines \ip{arg1}{arg2} to mean
                             % (arg1, arg2).

%\newcommand{\ip}[2]{\langle #1 | #2\rangle}
                             % This is an alternative definition of
                             % \ip that is commented out.

\begin{document}             % End of preamble and beginning of text.

\maketitle                   % Produces the title.

This is an example input file.  Comparing it with
the output it generates can show you how to
produce a simple document of your own.

\section{Ordinary Text}      % Produces section heading.  Lower-level
                             % sections are begun with similar
                             % \subsection and \subsubsection commands.

The ends  of words and sentences are marked
  by   spaces. It  doesn't matter how many
spaces    you type; one is as good as 100.  The
end of   a line counts as a space.

One   or more   blank lines denote the  end
of  a paragraph.

Since any number of consecutive spaces are treated
like a single one, the formatting of the input
file makes no difference to
      \LaTeX,                % The \LaTeX command generates the LaTeX logo.
but it makes a difference to you.  When you use
\LaTeX, making your input file as easy to read
as possible will be a great help as you write
your document and when you change it.  This sample
file shows how you can add comments to your own input
file.

Because printing is different from typewriting,
there are a number of things that you have to do
differently when preparing an input file than if
you were just typing the document directly.
Quotation marks like
       ``this''
have to be handled specially, as do quotes within
quotes:
       ``\,`this'            % \, separates the double and single quote.
        is what I just
        wrote, not  `that'\,''.

Dashes come in three sizes: an
       intra-word
dash, a medium dash for number ranges like
       1--2,
and a punctuation
       dash---like
this.

A sentence-ending space should be larger than the
space between words within a sentence.  You
sometimes have to type special commands in
conjunction with punctuation characters to get
this right, as in the following sentence.
       Gnats, gnus, etc.\ all  % `\ ' makes an inter-word space.
       begin with G\@.         % \@ marks end-of-sentence punctuation.
You should check the spaces after periods when
reading your output to make sure you haven't
forgotten any special cases.  Generating an
ellipsis
       \ldots\               % `\ ' is needed after `\ldots' because TeX
                             % ignores spaces after command names like \ldots
                             % made from \ + letters.
                             %
                             % Note how a `%' character causes TeX to ignore
                             % the end of the input line, so these blank lines
                             % do not start a new paragraph.
                             %
with the right spacing around the periods requires
a special command.

\LaTeX\ interprets some common characters as
commands, so you must type special commands to
generate them.  These characters include the
following:
       \$ \& \% \# \{ and \}.

In printing, text is usually emphasized with an
       \emph{italic}
type style.

\begin{em}
   A long segment of text can also be emphasized
   in this way.  Text within such a segment can be
   given \emph{additional} emphasis.
\end{em}

It is sometimes necessary to prevent \LaTeX\ from
breaking a line where it might otherwise do so.
This may be at a space, as between the ``Mr.''\ and
``Jones'' in
       ``Mr.~Jones'',        % ~ produces an unbreakable interword space.
or within a word---especially when the word is a
symbol like
       \mbox{\emph{itemnum}}
that makes little sense when hyphenated across
lines.

Footnotes\footnote{This is an example of a footnote.}
pose no problem.

\LaTeX\ is good at typesetting mathematical formulas
like
       \( x-3y + z = 7 \)
or
       \( a_{1} > x^{2n} + y^{2n} > x' \)
or
       \( \ip{A}{B} = \sum_{i} a_{i} b_{i} \).
The spaces you type in a formula are
ignored.  Remember that a letter like
       $x$                   % $ ... $  and  \( ... \)  are equivalent
is a formula when it denotes a mathematical
symbol, and it should be typed as one.

\section{Displayed Text}

Text is displayed by indenting it from the left
margin.  Quotations are commonly displayed.  There
are short quotations
\begin{quote}
   This is a short quotation.  It consists of a
   single paragraph of text.  See how it is formatted.
\end{quote}
and longer ones.
\begin{quotation}
   This is a longer quotation.  It consists of two
   paragraphs of text, neither of which are
   particularly interesting.

   This is the second paragraph of the quotation.  It
   is just as dull as the first paragraph.
\end{quotation}
Another frequently-displayed structure is a list.
The following is an example of an \emph{itemized}
list.
\begin{itemize}
   \item This is the first item of an itemized list.
         Each item in the list is marked with a ``tick''.
         You don't have to worry about what kind of tick
         mark is used.

   \item This is the second item of the list.  It
         contains another list nested inside it.  The inner
         list is an \emph{enumerated} list.
         \begin{enumerate}
            \item This is the first item of an enumerated
                  list that is nested within the itemized list.

            \item This is the second item of the inner list.
                  \LaTeX\ allows you to nest lists deeper than
                  you really should.
         \end{enumerate}
         This is the rest of the second item of the outer
         list.  It is no more interesting than any other
         part of the item.
   \item This is the third item of the list.
\end{itemize}
You can even display poetry.
\begin{verse}
   There is an environment
    for verse \\             % The \\ command separates lines
   Whose features some poets % within a stanza.
   will curse.

                             % One or more blank lines separate stanzas.

   For instead of making\\
   Them do \emph{all} line breaking, \\
   It allows them to put too many words on a line when they'd rather be
   forced to be terse.
\end{verse}

Mathematical formulas may also be displayed.  A
displayed formula
is
one-line long; multiline
formulas require special formatting instructions.
   \[  \ip{\Gamma}{\psi'} = x'' + y^{2} + z_{i}^{n}\]
Don't start a paragraph with a displayed equation,
nor make one a paragraph by itself.

\end{document}               % End of document.

\section{Bench Huge Section 1219}
% synthetic bench section 1219
% This is a sample LaTeX input file.  (Version of 12 August 2004.)
%
% A '%' character causes TeX to ignore all remaining text on the line,
% and is used for comments like this one.

\documentclass{article}      % Specifies the document class

                             % The preamble begins here.
\title{An Example Document}  % Declares the document's title.
\author{Leslie Lamport}      % Declares the author's name.
\date{January 21, 1994}      % Deleting this command produces today's date.

\newcommand{\ip}[2]{(#1, #2)}
                             % Defines \ip{arg1}{arg2} to mean
                             % (arg1, arg2).

%\newcommand{\ip}[2]{\langle #1 | #2\rangle}
                             % This is an alternative definition of
                             % \ip that is commented out.

\begin{document}             % End of preamble and beginning of text.

\maketitle                   % Produces the title.

This is an example input file.  Comparing it with
the output it generates can show you how to
produce a simple document of your own.

\section{Ordinary Text}      % Produces section heading.  Lower-level
                             % sections are begun with similar
                             % \subsection and \subsubsection commands.

The ends  of words and sentences are marked
  by   spaces. It  doesn't matter how many
spaces    you type; one is as good as 100.  The
end of   a line counts as a space.

One   or more   blank lines denote the  end
of  a paragraph.

Since any number of consecutive spaces are treated
like a single one, the formatting of the input
file makes no difference to
      \LaTeX,                % The \LaTeX command generates the LaTeX logo.
but it makes a difference to you.  When you use
\LaTeX, making your input file as easy to read
as possible will be a great help as you write
your document and when you change it.  This sample
file shows how you can add comments to your own input
file.

Because printing is different from typewriting,
there are a number of things that you have to do
differently when preparing an input file than if
you were just typing the document directly.
Quotation marks like
       ``this''
have to be handled specially, as do quotes within
quotes:
       ``\,`this'            % \, separates the double and single quote.
        is what I just
        wrote, not  `that'\,''.

Dashes come in three sizes: an
       intra-word
dash, a medium dash for number ranges like
       1--2,
and a punctuation
       dash---like
this.

A sentence-ending space should be larger than the
space between words within a sentence.  You
sometimes have to type special commands in
conjunction with punctuation characters to get
this right, as in the following sentence.
       Gnats, gnus, etc.\ all  % `\ ' makes an inter-word space.
       begin with G\@.         % \@ marks end-of-sentence punctuation.
You should check the spaces after periods when
reading your output to make sure you haven't
forgotten any special cases.  Generating an
ellipsis
       \ldots\               % `\ ' is needed after `\ldots' because TeX
                             % ignores spaces after command names like \ldots
                             % made from \ + letters.
                             %
                             % Note how a `%' character causes TeX to ignore
                             % the end of the input line, so these blank lines
                             % do not start a new paragraph.
                             %
with the right spacing around the periods requires
a special command.

\LaTeX\ interprets some common characters as
commands, so you must type special commands to
generate them.  These characters include the
following:
       \$ \& \% \# \{ and \}.

In printing, text is usually emphasized with an
       \emph{italic}
type style.

\begin{em}
   A long segment of text can also be emphasized
   in this way.  Text within such a segment can be
   given \emph{additional} emphasis.
\end{em}

It is sometimes necessary to prevent \LaTeX\ from
breaking a line where it might otherwise do so.
This may be at a space, as between the ``Mr.''\ and
``Jones'' in
       ``Mr.~Jones'',        % ~ produces an unbreakable interword space.
or within a word---especially when the word is a
symbol like
       \mbox{\emph{itemnum}}
that makes little sense when hyphenated across
lines.

Footnotes\footnote{This is an example of a footnote.}
pose no problem.

\LaTeX\ is good at typesetting mathematical formulas
like
       \( x-3y + z = 7 \)
or
       \( a_{1} > x^{2n} + y^{2n} > x' \)
or
       \( \ip{A}{B} = \sum_{i} a_{i} b_{i} \).
The spaces you type in a formula are
ignored.  Remember that a letter like
       $x$                   % $ ... $  and  \( ... \)  are equivalent
is a formula when it denotes a mathematical
symbol, and it should be typed as one.

\section{Displayed Text}

Text is displayed by indenting it from the left
margin.  Quotations are commonly displayed.  There
are short quotations
\begin{quote}
   This is a short quotation.  It consists of a
   single paragraph of text.  See how it is formatted.
\end{quote}
and longer ones.
\begin{quotation}
   This is a longer quotation.  It consists of two
   paragraphs of text, neither of which are
   particularly interesting.

   This is the second paragraph of the quotation.  It
   is just as dull as the first paragraph.
\end{quotation}
Another frequently-displayed structure is a list.
The following is an example of an \emph{itemized}
list.
\begin{itemize}
   \item This is the first item of an itemized list.
         Each item in the list is marked with a ``tick''.
         You don't have to worry about what kind of tick
         mark is used.

   \item This is the second item of the list.  It
         contains another list nested inside it.  The inner
         list is an \emph{enumerated} list.
         \begin{enumerate}
            \item This is the first item of an enumerated
                  list that is nested within the itemized list.

            \item This is the second item of the inner list.
                  \LaTeX\ allows you to nest lists deeper than
                  you really should.
         \end{enumerate}
         This is the rest of the second item of the outer
         list.  It is no more interesting than any other
         part of the item.
   \item This is the third item of the list.
\end{itemize}
You can even display poetry.
\begin{verse}
   There is an environment
    for verse \\             % The \\ command separates lines
   Whose features some poets % within a stanza.
   will curse.

                             % One or more blank lines separate stanzas.

   For instead of making\\
   Them do \emph{all} line breaking, \\
   It allows them to put too many words on a line when they'd rather be
   forced to be terse.
\end{verse}

Mathematical formulas may also be displayed.  A
displayed formula
is
one-line long; multiline
formulas require special formatting instructions.
   \[  \ip{\Gamma}{\psi'} = x'' + y^{2} + z_{i}^{n}\]
Don't start a paragraph with a displayed equation,
nor make one a paragraph by itself.

\end{document}               % End of document.

\section{Bench Huge Section 1220}
% synthetic bench section 1220
% This is a sample LaTeX input file.  (Version of 12 August 2004.)
%
% A '%' character causes TeX to ignore all remaining text on the line,
% and is used for comments like this one.

\documentclass{article}      % Specifies the document class

                             % The preamble begins here.
\title{An Example Document}  % Declares the document's title.
\author{Leslie Lamport}      % Declares the author's name.
\date{January 21, 1994}      % Deleting this command produces today's date.

\newcommand{\ip}[2]{(#1, #2)}
                             % Defines \ip{arg1}{arg2} to mean
                             % (arg1, arg2).

%\newcommand{\ip}[2]{\langle #1 | #2\rangle}
                             % This is an alternative definition of
                             % \ip that is commented out.

\begin{document}             % End of preamble and beginning of text.

\maketitle                   % Produces the title.

This is an example input file.  Comparing it with
the output it generates can show you how to
produce a simple document of your own.

\section{Ordinary Text}      % Produces section heading.  Lower-level
                             % sections are begun with similar
                             % \subsection and \subsubsection commands.

The ends  of words and sentences are marked
  by   spaces. It  doesn't matter how many
spaces    you type; one is as good as 100.  The
end of   a line counts as a space.

One   or more   blank lines denote the  end
of  a paragraph.

Since any number of consecutive spaces are treated
like a single one, the formatting of the input
file makes no difference to
      \LaTeX,                % The \LaTeX command generates the LaTeX logo.
but it makes a difference to you.  When you use
\LaTeX, making your input file as easy to read
as possible will be a great help as you write
your document and when you change it.  This sample
file shows how you can add comments to your own input
file.

Because printing is different from typewriting,
there are a number of things that you have to do
differently when preparing an input file than if
you were just typing the document directly.
Quotation marks like
       ``this''
have to be handled specially, as do quotes within
quotes:
       ``\,`this'            % \, separates the double and single quote.
        is what I just
        wrote, not  `that'\,''.

Dashes come in three sizes: an
       intra-word
dash, a medium dash for number ranges like
       1--2,
and a punctuation
       dash---like
this.

A sentence-ending space should be larger than the
space between words within a sentence.  You
sometimes have to type special commands in
conjunction with punctuation characters to get
this right, as in the following sentence.
       Gnats, gnus, etc.\ all  % `\ ' makes an inter-word space.
       begin with G\@.         % \@ marks end-of-sentence punctuation.
You should check the spaces after periods when
reading your output to make sure you haven't
forgotten any special cases.  Generating an
ellipsis
       \ldots\               % `\ ' is needed after `\ldots' because TeX
                             % ignores spaces after command names like \ldots
                             % made from \ + letters.
                             %
                             % Note how a `%' character causes TeX to ignore
                             % the end of the input line, so these blank lines
                             % do not start a new paragraph.
                             %
with the right spacing around the periods requires
a special command.

\LaTeX\ interprets some common characters as
commands, so you must type special commands to
generate them.  These characters include the
following:
       \$ \& \% \# \{ and \}.

In printing, text is usually emphasized with an
       \emph{italic}
type style.

\begin{em}
   A long segment of text can also be emphasized
   in this way.  Text within such a segment can be
   given \emph{additional} emphasis.
\end{em}

It is sometimes necessary to prevent \LaTeX\ from
breaking a line where it might otherwise do so.
This may be at a space, as between the ``Mr.''\ and
``Jones'' in
       ``Mr.~Jones'',        % ~ produces an unbreakable interword space.
or within a word---especially when the word is a
symbol like
       \mbox{\emph{itemnum}}
that makes little sense when hyphenated across
lines.

Footnotes\footnote{This is an example of a footnote.}
pose no problem.

\LaTeX\ is good at typesetting mathematical formulas
like
       \( x-3y + z = 7 \)
or
       \( a_{1} > x^{2n} + y^{2n} > x' \)
or
       \( \ip{A}{B} = \sum_{i} a_{i} b_{i} \).
The spaces you type in a formula are
ignored.  Remember that a letter like
       $x$                   % $ ... $  and  \( ... \)  are equivalent
is a formula when it denotes a mathematical
symbol, and it should be typed as one.

\section{Displayed Text}

Text is displayed by indenting it from the left
margin.  Quotations are commonly displayed.  There
are short quotations
\begin{quote}
   This is a short quotation.  It consists of a
   single paragraph of text.  See how it is formatted.
\end{quote}
and longer ones.
\begin{quotation}
   This is a longer quotation.  It consists of two
   paragraphs of text, neither of which are
   particularly interesting.

   This is the second paragraph of the quotation.  It
   is just as dull as the first paragraph.
\end{quotation}
Another frequently-displayed structure is a list.
The following is an example of an \emph{itemized}
list.
\begin{itemize}
   \item This is the first item of an itemized list.
         Each item in the list is marked with a ``tick''.
         You don't have to worry about what kind of tick
         mark is used.

   \item This is the second item of the list.  It
         contains another list nested inside it.  The inner
         list is an \emph{enumerated} list.
         \begin{enumerate}
            \item This is the first item of an enumerated
                  list that is nested within the itemized list.

            \item This is the second item of the inner list.
                  \LaTeX\ allows you to nest lists deeper than
                  you really should.
         \end{enumerate}
         This is the rest of the second item of the outer
         list.  It is no more interesting than any other
         part of the item.
   \item This is the third item of the list.
\end{itemize}
You can even display poetry.
\begin{verse}
   There is an environment
    for verse \\             % The \\ command separates lines
   Whose features some poets % within a stanza.
   will curse.

                             % One or more blank lines separate stanzas.

   For instead of making\\
   Them do \emph{all} line breaking, \\
   It allows them to put too many words on a line when they'd rather be
   forced to be terse.
\end{verse}

Mathematical formulas may also be displayed.  A
displayed formula
is
one-line long; multiline
formulas require special formatting instructions.
   \[  \ip{\Gamma}{\psi'} = x'' + y^{2} + z_{i}^{n}\]
Don't start a paragraph with a displayed equation,
nor make one a paragraph by itself.

\end{document}               % End of document.

\section{Bench Huge Section 1221}
% synthetic bench section 1221
% This is a sample LaTeX input file.  (Version of 12 August 2004.)
%
% A '%' character causes TeX to ignore all remaining text on the line,
% and is used for comments like this one.

\documentclass{article}      % Specifies the document class

                             % The preamble begins here.
\title{An Example Document}  % Declares the document's title.
\author{Leslie Lamport}      % Declares the author's name.
\date{January 21, 1994}      % Deleting this command produces today's date.

\newcommand{\ip}[2]{(#1, #2)}
                             % Defines \ip{arg1}{arg2} to mean
                             % (arg1, arg2).

%\newcommand{\ip}[2]{\langle #1 | #2\rangle}
                             % This is an alternative definition of
                             % \ip that is commented out.

\begin{document}             % End of preamble and beginning of text.

\maketitle                   % Produces the title.

This is an example input file.  Comparing it with
the output it generates can show you how to
produce a simple document of your own.

\section{Ordinary Text}      % Produces section heading.  Lower-level
                             % sections are begun with similar
                             % \subsection and \subsubsection commands.

The ends  of words and sentences are marked
  by   spaces. It  doesn't matter how many
spaces    you type; one is as good as 100.  The
end of   a line counts as a space.

One   or more   blank lines denote the  end
of  a paragraph.

Since any number of consecutive spaces are treated
like a single one, the formatting of the input
file makes no difference to
      \LaTeX,                % The \LaTeX command generates the LaTeX logo.
but it makes a difference to you.  When you use
\LaTeX, making your input file as easy to read
as possible will be a great help as you write
your document and when you change it.  This sample
file shows how you can add comments to your own input
file.

Because printing is different from typewriting,
there are a number of things that you have to do
differently when preparing an input file than if
you were just typing the document directly.
Quotation marks like
       ``this''
have to be handled specially, as do quotes within
quotes:
       ``\,`this'            % \, separates the double and single quote.
        is what I just
        wrote, not  `that'\,''.

Dashes come in three sizes: an
       intra-word
dash, a medium dash for number ranges like
       1--2,
and a punctuation
       dash---like
this.

A sentence-ending space should be larger than the
space between words within a sentence.  You
sometimes have to type special commands in
conjunction with punctuation characters to get
this right, as in the following sentence.
       Gnats, gnus, etc.\ all  % `\ ' makes an inter-word space.
       begin with G\@.         % \@ marks end-of-sentence punctuation.
You should check the spaces after periods when
reading your output to make sure you haven't
forgotten any special cases.  Generating an
ellipsis
       \ldots\               % `\ ' is needed after `\ldots' because TeX
                             % ignores spaces after command names like \ldots
                             % made from \ + letters.
                             %
                             % Note how a `%' character causes TeX to ignore
                             % the end of the input line, so these blank lines
                             % do not start a new paragraph.
                             %
with the right spacing around the periods requires
a special command.

\LaTeX\ interprets some common characters as
commands, so you must type special commands to
generate them.  These characters include the
following:
       \$ \& \% \# \{ and \}.

In printing, text is usually emphasized with an
       \emph{italic}
type style.

\begin{em}
   A long segment of text can also be emphasized
   in this way.  Text within such a segment can be
   given \emph{additional} emphasis.
\end{em}

It is sometimes necessary to prevent \LaTeX\ from
breaking a line where it might otherwise do so.
This may be at a space, as between the ``Mr.''\ and
``Jones'' in
       ``Mr.~Jones'',        % ~ produces an unbreakable interword space.
or within a word---especially when the word is a
symbol like
       \mbox{\emph{itemnum}}
that makes little sense when hyphenated across
lines.

Footnotes\footnote{This is an example of a footnote.}
pose no problem.

\LaTeX\ is good at typesetting mathematical formulas
like
       \( x-3y + z = 7 \)
or
       \( a_{1} > x^{2n} + y^{2n} > x' \)
or
       \( \ip{A}{B} = \sum_{i} a_{i} b_{i} \).
The spaces you type in a formula are
ignored.  Remember that a letter like
       $x$                   % $ ... $  and  \( ... \)  are equivalent
is a formula when it denotes a mathematical
symbol, and it should be typed as one.

\section{Displayed Text}

Text is displayed by indenting it from the left
margin.  Quotations are commonly displayed.  There
are short quotations
\begin{quote}
   This is a short quotation.  It consists of a
   single paragraph of text.  See how it is formatted.
\end{quote}
and longer ones.
\begin{quotation}
   This is a longer quotation.  It consists of two
   paragraphs of text, neither of which are
   particularly interesting.

   This is the second paragraph of the quotation.  It
   is just as dull as the first paragraph.
\end{quotation}
Another frequently-displayed structure is a list.
The following is an example of an \emph{itemized}
list.
\begin{itemize}
   \item This is the first item of an itemized list.
         Each item in the list is marked with a ``tick''.
         You don't have to worry about what kind of tick
         mark is used.

   \item This is the second item of the list.  It
         contains another list nested inside it.  The inner
         list is an \emph{enumerated} list.
         \begin{enumerate}
            \item This is the first item of an enumerated
                  list that is nested within the itemized list.

            \item This is the second item of the inner list.
                  \LaTeX\ allows you to nest lists deeper than
                  you really should.
         \end{enumerate}
         This is the rest of the second item of the outer
         list.  It is no more interesting than any other
         part of the item.
   \item This is the third item of the list.
\end{itemize}
You can even display poetry.
\begin{verse}
   There is an environment
    for verse \\             % The \\ command separates lines
   Whose features some poets % within a stanza.
   will curse.

                             % One or more blank lines separate stanzas.

   For instead of making\\
   Them do \emph{all} line breaking, \\
   It allows them to put too many words on a line when they'd rather be
   forced to be terse.
\end{verse}

Mathematical formulas may also be displayed.  A
displayed formula
is
one-line long; multiline
formulas require special formatting instructions.
   \[  \ip{\Gamma}{\psi'} = x'' + y^{2} + z_{i}^{n}\]
Don't start a paragraph with a displayed equation,
nor make one a paragraph by itself.

\end{document}               % End of document.

\section{Bench Huge Section 1222}
% synthetic bench section 1222
% This is a sample LaTeX input file.  (Version of 12 August 2004.)
%
% A '%' character causes TeX to ignore all remaining text on the line,
% and is used for comments like this one.

\documentclass{article}      % Specifies the document class

                             % The preamble begins here.
\title{An Example Document}  % Declares the document's title.
\author{Leslie Lamport}      % Declares the author's name.
\date{January 21, 1994}      % Deleting this command produces today's date.

\newcommand{\ip}[2]{(#1, #2)}
                             % Defines \ip{arg1}{arg2} to mean
                             % (arg1, arg2).

%\newcommand{\ip}[2]{\langle #1 | #2\rangle}
                             % This is an alternative definition of
                             % \ip that is commented out.

\begin{document}             % End of preamble and beginning of text.

\maketitle                   % Produces the title.

This is an example input file.  Comparing it with
the output it generates can show you how to
produce a simple document of your own.

\section{Ordinary Text}      % Produces section heading.  Lower-level
                             % sections are begun with similar
                             % \subsection and \subsubsection commands.

The ends  of words and sentences are marked
  by   spaces. It  doesn't matter how many
spaces    you type; one is as good as 100.  The
end of   a line counts as a space.

One   or more   blank lines denote the  end
of  a paragraph.

Since any number of consecutive spaces are treated
like a single one, the formatting of the input
file makes no difference to
      \LaTeX,                % The \LaTeX command generates the LaTeX logo.
but it makes a difference to you.  When you use
\LaTeX, making your input file as easy to read
as possible will be a great help as you write
your document and when you change it.  This sample
file shows how you can add comments to your own input
file.

Because printing is different from typewriting,
there are a number of things that you have to do
differently when preparing an input file than if
you were just typing the document directly.
Quotation marks like
       ``this''
have to be handled specially, as do quotes within
quotes:
       ``\,`this'            % \, separates the double and single quote.
        is what I just
        wrote, not  `that'\,''.

Dashes come in three sizes: an
       intra-word
dash, a medium dash for number ranges like
       1--2,
and a punctuation
       dash---like
this.

A sentence-ending space should be larger than the
space between words within a sentence.  You
sometimes have to type special commands in
conjunction with punctuation characters to get
this right, as in the following sentence.
       Gnats, gnus, etc.\ all  % `\ ' makes an inter-word space.
       begin with G\@.         % \@ marks end-of-sentence punctuation.
You should check the spaces after periods when
reading your output to make sure you haven't
forgotten any special cases.  Generating an
ellipsis
       \ldots\               % `\ ' is needed after `\ldots' because TeX
                             % ignores spaces after command names like \ldots
                             % made from \ + letters.
                             %
                             % Note how a `%' character causes TeX to ignore
                             % the end of the input line, so these blank lines
                             % do not start a new paragraph.
                             %
with the right spacing around the periods requires
a special command.

\LaTeX\ interprets some common characters as
commands, so you must type special commands to
generate them.  These characters include the
following:
       \$ \& \% \# \{ and \}.

In printing, text is usually emphasized with an
       \emph{italic}
type style.

\begin{em}
   A long segment of text can also be emphasized
   in this way.  Text within such a segment can be
   given \emph{additional} emphasis.
\end{em}

It is sometimes necessary to prevent \LaTeX\ from
breaking a line where it might otherwise do so.
This may be at a space, as between the ``Mr.''\ and
``Jones'' in
       ``Mr.~Jones'',        % ~ produces an unbreakable interword space.
or within a word---especially when the word is a
symbol like
       \mbox{\emph{itemnum}}
that makes little sense when hyphenated across
lines.

Footnotes\footnote{This is an example of a footnote.}
pose no problem.

\LaTeX\ is good at typesetting mathematical formulas
like
       \( x-3y + z = 7 \)
or
       \( a_{1} > x^{2n} + y^{2n} > x' \)
or
       \( \ip{A}{B} = \sum_{i} a_{i} b_{i} \).
The spaces you type in a formula are
ignored.  Remember that a letter like
       $x$                   % $ ... $  and  \( ... \)  are equivalent
is a formula when it denotes a mathematical
symbol, and it should be typed as one.

\section{Displayed Text}

Text is displayed by indenting it from the left
margin.  Quotations are commonly displayed.  There
are short quotations
\begin{quote}
   This is a short quotation.  It consists of a
   single paragraph of text.  See how it is formatted.
\end{quote}
and longer ones.
\begin{quotation}
   This is a longer quotation.  It consists of two
   paragraphs of text, neither of which are
   particularly interesting.

   This is the second paragraph of the quotation.  It
   is just as dull as the first paragraph.
\end{quotation}
Another frequently-displayed structure is a list.
The following is an example of an \emph{itemized}
list.
\begin{itemize}
   \item This is the first item of an itemized list.
         Each item in the list is marked with a ``tick''.
         You don't have to worry about what kind of tick
         mark is used.

   \item This is the second item of the list.  It
         contains another list nested inside it.  The inner
         list is an \emph{enumerated} list.
         \begin{enumerate}
            \item This is the first item of an enumerated
                  list that is nested within the itemized list.

            \item This is the second item of the inner list.
                  \LaTeX\ allows you to nest lists deeper than
                  you really should.
         \end{enumerate}
         This is the rest of the second item of the outer
         list.  It is no more interesting than any other
         part of the item.
   \item This is the third item of the list.
\end{itemize}
You can even display poetry.
\begin{verse}
   There is an environment
    for verse \\             % The \\ command separates lines
   Whose features some poets % within a stanza.
   will curse.

                             % One or more blank lines separate stanzas.

   For instead of making\\
   Them do \emph{all} line breaking, \\
   It allows them to put too many words on a line when they'd rather be
   forced to be terse.
\end{verse}

Mathematical formulas may also be displayed.  A
displayed formula
is
one-line long; multiline
formulas require special formatting instructions.
   \[  \ip{\Gamma}{\psi'} = x'' + y^{2} + z_{i}^{n}\]
Don't start a paragraph with a displayed equation,
nor make one a paragraph by itself.

\end{document}               % End of document.

\section{Bench Huge Section 1223}
% synthetic bench section 1223
% This is a sample LaTeX input file.  (Version of 12 August 2004.)
%
% A '%' character causes TeX to ignore all remaining text on the line,
% and is used for comments like this one.

\documentclass{article}      % Specifies the document class

                             % The preamble begins here.
\title{An Example Document}  % Declares the document's title.
\author{Leslie Lamport}      % Declares the author's name.
\date{January 21, 1994}      % Deleting this command produces today's date.

\newcommand{\ip}[2]{(#1, #2)}
                             % Defines \ip{arg1}{arg2} to mean
                             % (arg1, arg2).

%\newcommand{\ip}[2]{\langle #1 | #2\rangle}
                             % This is an alternative definition of
                             % \ip that is commented out.

\begin{document}             % End of preamble and beginning of text.

\maketitle                   % Produces the title.

This is an example input file.  Comparing it with
the output it generates can show you how to
produce a simple document of your own.

\section{Ordinary Text}      % Produces section heading.  Lower-level
                             % sections are begun with similar
                             % \subsection and \subsubsection commands.

The ends  of words and sentences are marked
  by   spaces. It  doesn't matter how many
spaces    you type; one is as good as 100.  The
end of   a line counts as a space.

One   or more   blank lines denote the  end
of  a paragraph.

Since any number of consecutive spaces are treated
like a single one, the formatting of the input
file makes no difference to
      \LaTeX,                % The \LaTeX command generates the LaTeX logo.
but it makes a difference to you.  When you use
\LaTeX, making your input file as easy to read
as possible will be a great help as you write
your document and when you change it.  This sample
file shows how you can add comments to your own input
file.

Because printing is different from typewriting,
there are a number of things that you have to do
differently when preparing an input file than if
you were just typing the document directly.
Quotation marks like
       ``this''
have to be handled specially, as do quotes within
quotes:
       ``\,`this'            % \, separates the double and single quote.
        is what I just
        wrote, not  `that'\,''.

Dashes come in three sizes: an
       intra-word
dash, a medium dash for number ranges like
       1--2,
and a punctuation
       dash---like
this.

A sentence-ending space should be larger than the
space between words within a sentence.  You
sometimes have to type special commands in
conjunction with punctuation characters to get
this right, as in the following sentence.
       Gnats, gnus, etc.\ all  % `\ ' makes an inter-word space.
       begin with G\@.         % \@ marks end-of-sentence punctuation.
You should check the spaces after periods when
reading your output to make sure you haven't
forgotten any special cases.  Generating an
ellipsis
       \ldots\               % `\ ' is needed after `\ldots' because TeX
                             % ignores spaces after command names like \ldots
                             % made from \ + letters.
                             %
                             % Note how a `%' character causes TeX to ignore
                             % the end of the input line, so these blank lines
                             % do not start a new paragraph.
                             %
with the right spacing around the periods requires
a special command.

\LaTeX\ interprets some common characters as
commands, so you must type special commands to
generate them.  These characters include the
following:
       \$ \& \% \# \{ and \}.

In printing, text is usually emphasized with an
       \emph{italic}
type style.

\begin{em}
   A long segment of text can also be emphasized
   in this way.  Text within such a segment can be
   given \emph{additional} emphasis.
\end{em}

It is sometimes necessary to prevent \LaTeX\ from
breaking a line where it might otherwise do so.
This may be at a space, as between the ``Mr.''\ and
``Jones'' in
       ``Mr.~Jones'',        % ~ produces an unbreakable interword space.
or within a word---especially when the word is a
symbol like
       \mbox{\emph{itemnum}}
that makes little sense when hyphenated across
lines.

Footnotes\footnote{This is an example of a footnote.}
pose no problem.

\LaTeX\ is good at typesetting mathematical formulas
like
       \( x-3y + z = 7 \)
or
       \( a_{1} > x^{2n} + y^{2n} > x' \)
or
       \( \ip{A}{B} = \sum_{i} a_{i} b_{i} \).
The spaces you type in a formula are
ignored.  Remember that a letter like
       $x$                   % $ ... $  and  \( ... \)  are equivalent
is a formula when it denotes a mathematical
symbol, and it should be typed as one.

\section{Displayed Text}

Text is displayed by indenting it from the left
margin.  Quotations are commonly displayed.  There
are short quotations
\begin{quote}
   This is a short quotation.  It consists of a
   single paragraph of text.  See how it is formatted.
\end{quote}
and longer ones.
\begin{quotation}
   This is a longer quotation.  It consists of two
   paragraphs of text, neither of which are
   particularly interesting.

   This is the second paragraph of the quotation.  It
   is just as dull as the first paragraph.
\end{quotation}
Another frequently-displayed structure is a list.
The following is an example of an \emph{itemized}
list.
\begin{itemize}
   \item This is the first item of an itemized list.
         Each item in the list is marked with a ``tick''.
         You don't have to worry about what kind of tick
         mark is used.

   \item This is the second item of the list.  It
         contains another list nested inside it.  The inner
         list is an \emph{enumerated} list.
         \begin{enumerate}
            \item This is the first item of an enumerated
                  list that is nested within the itemized list.

            \item This is the second item of the inner list.
                  \LaTeX\ allows you to nest lists deeper than
                  you really should.
         \end{enumerate}
         This is the rest of the second item of the outer
         list.  It is no more interesting than any other
         part of the item.
   \item This is the third item of the list.
\end{itemize}
You can even display poetry.
\begin{verse}
   There is an environment
    for verse \\             % The \\ command separates lines
   Whose features some poets % within a stanza.
   will curse.

                             % One or more blank lines separate stanzas.

   For instead of making\\
   Them do \emph{all} line breaking, \\
   It allows them to put too many words on a line when they'd rather be
   forced to be terse.
\end{verse}

Mathematical formulas may also be displayed.  A
displayed formula
is
one-line long; multiline
formulas require special formatting instructions.
   \[  \ip{\Gamma}{\psi'} = x'' + y^{2} + z_{i}^{n}\]
Don't start a paragraph with a displayed equation,
nor make one a paragraph by itself.

\end{document}               % End of document.

\section{Bench Huge Section 1224}
% synthetic bench section 1224
% This is a sample LaTeX input file.  (Version of 12 August 2004.)
%
% A '%' character causes TeX to ignore all remaining text on the line,
% and is used for comments like this one.

\documentclass{article}      % Specifies the document class

                             % The preamble begins here.
\title{An Example Document}  % Declares the document's title.
\author{Leslie Lamport}      % Declares the author's name.
\date{January 21, 1994}      % Deleting this command produces today's date.

\newcommand{\ip}[2]{(#1, #2)}
                             % Defines \ip{arg1}{arg2} to mean
                             % (arg1, arg2).

%\newcommand{\ip}[2]{\langle #1 | #2\rangle}
                             % This is an alternative definition of
                             % \ip that is commented out.

\begin{document}             % End of preamble and beginning of text.

\maketitle                   % Produces the title.

This is an example input file.  Comparing it with
the output it generates can show you how to
produce a simple document of your own.

\section{Ordinary Text}      % Produces section heading.  Lower-level
                             % sections are begun with similar
                             % \subsection and \subsubsection commands.

The ends  of words and sentences are marked
  by   spaces. It  doesn't matter how many
spaces    you type; one is as good as 100.  The
end of   a line counts as a space.

One   or more   blank lines denote the  end
of  a paragraph.

Since any number of consecutive spaces are treated
like a single one, the formatting of the input
file makes no difference to
      \LaTeX,                % The \LaTeX command generates the LaTeX logo.
but it makes a difference to you.  When you use
\LaTeX, making your input file as easy to read
as possible will be a great help as you write
your document and when you change it.  This sample
file shows how you can add comments to your own input
file.

Because printing is different from typewriting,
there are a number of things that you have to do
differently when preparing an input file than if
you were just typing the document directly.
Quotation marks like
       ``this''
have to be handled specially, as do quotes within
quotes:
       ``\,`this'            % \, separates the double and single quote.
        is what I just
        wrote, not  `that'\,''.

Dashes come in three sizes: an
       intra-word
dash, a medium dash for number ranges like
       1--2,
and a punctuation
       dash---like
this.

A sentence-ending space should be larger than the
space between words within a sentence.  You
sometimes have to type special commands in
conjunction with punctuation characters to get
this right, as in the following sentence.
       Gnats, gnus, etc.\ all  % `\ ' makes an inter-word space.
       begin with G\@.         % \@ marks end-of-sentence punctuation.
You should check the spaces after periods when
reading your output to make sure you haven't
forgotten any special cases.  Generating an
ellipsis
       \ldots\               % `\ ' is needed after `\ldots' because TeX
                             % ignores spaces after command names like \ldots
                             % made from \ + letters.
                             %
                             % Note how a `%' character causes TeX to ignore
                             % the end of the input line, so these blank lines
                             % do not start a new paragraph.
                             %
with the right spacing around the periods requires
a special command.

\LaTeX\ interprets some common characters as
commands, so you must type special commands to
generate them.  These characters include the
following:
       \$ \& \% \# \{ and \}.

In printing, text is usually emphasized with an
       \emph{italic}
type style.

\begin{em}
   A long segment of text can also be emphasized
   in this way.  Text within such a segment can be
   given \emph{additional} emphasis.
\end{em}

It is sometimes necessary to prevent \LaTeX\ from
breaking a line where it might otherwise do so.
This may be at a space, as between the ``Mr.''\ and
``Jones'' in
       ``Mr.~Jones'',        % ~ produces an unbreakable interword space.
or within a word---especially when the word is a
symbol like
       \mbox{\emph{itemnum}}
that makes little sense when hyphenated across
lines.

Footnotes\footnote{This is an example of a footnote.}
pose no problem.

\LaTeX\ is good at typesetting mathematical formulas
like
       \( x-3y + z = 7 \)
or
       \( a_{1} > x^{2n} + y^{2n} > x' \)
or
       \( \ip{A}{B} = \sum_{i} a_{i} b_{i} \).
The spaces you type in a formula are
ignored.  Remember that a letter like
       $x$                   % $ ... $  and  \( ... \)  are equivalent
is a formula when it denotes a mathematical
symbol, and it should be typed as one.

\section{Displayed Text}

Text is displayed by indenting it from the left
margin.  Quotations are commonly displayed.  There
are short quotations
\begin{quote}
   This is a short quotation.  It consists of a
   single paragraph of text.  See how it is formatted.
\end{quote}
and longer ones.
\begin{quotation}
   This is a longer quotation.  It consists of two
   paragraphs of text, neither of which are
   particularly interesting.

   This is the second paragraph of the quotation.  It
   is just as dull as the first paragraph.
\end{quotation}
Another frequently-displayed structure is a list.
The following is an example of an \emph{itemized}
list.
\begin{itemize}
   \item This is the first item of an itemized list.
         Each item in the list is marked with a ``tick''.
         You don't have to worry about what kind of tick
         mark is used.

   \item This is the second item of the list.  It
         contains another list nested inside it.  The inner
         list is an \emph{enumerated} list.
         \begin{enumerate}
            \item This is the first item of an enumerated
                  list that is nested within the itemized list.

            \item This is the second item of the inner list.
                  \LaTeX\ allows you to nest lists deeper than
                  you really should.
         \end{enumerate}
         This is the rest of the second item of the outer
         list.  It is no more interesting than any other
         part of the item.
   \item This is the third item of the list.
\end{itemize}
You can even display poetry.
\begin{verse}
   There is an environment
    for verse \\             % The \\ command separates lines
   Whose features some poets % within a stanza.
   will curse.

                             % One or more blank lines separate stanzas.

   For instead of making\\
   Them do \emph{all} line breaking, \\
   It allows them to put too many words on a line when they'd rather be
   forced to be terse.
\end{verse}

Mathematical formulas may also be displayed.  A
displayed formula
is
one-line long; multiline
formulas require special formatting instructions.
   \[  \ip{\Gamma}{\psi'} = x'' + y^{2} + z_{i}^{n}\]
Don't start a paragraph with a displayed equation,
nor make one a paragraph by itself.

\end{document}               % End of document.

\section{Bench Huge Section 1225}
% synthetic bench section 1225
% This is a sample LaTeX input file.  (Version of 12 August 2004.)
%
% A '%' character causes TeX to ignore all remaining text on the line,
% and is used for comments like this one.

\documentclass{article}      % Specifies the document class

                             % The preamble begins here.
\title{An Example Document}  % Declares the document's title.
\author{Leslie Lamport}      % Declares the author's name.
\date{January 21, 1994}      % Deleting this command produces today's date.

\newcommand{\ip}[2]{(#1, #2)}
                             % Defines \ip{arg1}{arg2} to mean
                             % (arg1, arg2).

%\newcommand{\ip}[2]{\langle #1 | #2\rangle}
                             % This is an alternative definition of
                             % \ip that is commented out.

\begin{document}             % End of preamble and beginning of text.

\maketitle                   % Produces the title.

This is an example input file.  Comparing it with
the output it generates can show you how to
produce a simple document of your own.

\section{Ordinary Text}      % Produces section heading.  Lower-level
                             % sections are begun with similar
                             % \subsection and \subsubsection commands.

The ends  of words and sentences are marked
  by   spaces. It  doesn't matter how many
spaces    you type; one is as good as 100.  The
end of   a line counts as a space.

One   or more   blank lines denote the  end
of  a paragraph.

Since any number of consecutive spaces are treated
like a single one, the formatting of the input
file makes no difference to
      \LaTeX,                % The \LaTeX command generates the LaTeX logo.
but it makes a difference to you.  When you use
\LaTeX, making your input file as easy to read
as possible will be a great help as you write
your document and when you change it.  This sample
file shows how you can add comments to your own input
file.

Because printing is different from typewriting,
there are a number of things that you have to do
differently when preparing an input file than if
you were just typing the document directly.
Quotation marks like
       ``this''
have to be handled specially, as do quotes within
quotes:
       ``\,`this'            % \, separates the double and single quote.
        is what I just
        wrote, not  `that'\,''.

Dashes come in three sizes: an
       intra-word
dash, a medium dash for number ranges like
       1--2,
and a punctuation
       dash---like
this.

A sentence-ending space should be larger than the
space between words within a sentence.  You
sometimes have to type special commands in
conjunction with punctuation characters to get
this right, as in the following sentence.
       Gnats, gnus, etc.\ all  % `\ ' makes an inter-word space.
       begin with G\@.         % \@ marks end-of-sentence punctuation.
You should check the spaces after periods when
reading your output to make sure you haven't
forgotten any special cases.  Generating an
ellipsis
       \ldots\               % `\ ' is needed after `\ldots' because TeX
                             % ignores spaces after command names like \ldots
                             % made from \ + letters.
                             %
                             % Note how a `%' character causes TeX to ignore
                             % the end of the input line, so these blank lines
                             % do not start a new paragraph.
                             %
with the right spacing around the periods requires
a special command.

\LaTeX\ interprets some common characters as
commands, so you must type special commands to
generate them.  These characters include the
following:
       \$ \& \% \# \{ and \}.

In printing, text is usually emphasized with an
       \emph{italic}
type style.

\begin{em}
   A long segment of text can also be emphasized
   in this way.  Text within such a segment can be
   given \emph{additional} emphasis.
\end{em}

It is sometimes necessary to prevent \LaTeX\ from
breaking a line where it might otherwise do so.
This may be at a space, as between the ``Mr.''\ and
``Jones'' in
       ``Mr.~Jones'',        % ~ produces an unbreakable interword space.
or within a word---especially when the word is a
symbol like
       \mbox{\emph{itemnum}}
that makes little sense when hyphenated across
lines.

Footnotes\footnote{This is an example of a footnote.}
pose no problem.

\LaTeX\ is good at typesetting mathematical formulas
like
       \( x-3y + z = 7 \)
or
       \( a_{1} > x^{2n} + y^{2n} > x' \)
or
       \( \ip{A}{B} = \sum_{i} a_{i} b_{i} \).
The spaces you type in a formula are
ignored.  Remember that a letter like
       $x$                   % $ ... $  and  \( ... \)  are equivalent
is a formula when it denotes a mathematical
symbol, and it should be typed as one.

\section{Displayed Text}

Text is displayed by indenting it from the left
margin.  Quotations are commonly displayed.  There
are short quotations
\begin{quote}
   This is a short quotation.  It consists of a
   single paragraph of text.  See how it is formatted.
\end{quote}
and longer ones.
\begin{quotation}
   This is a longer quotation.  It consists of two
   paragraphs of text, neither of which are
   particularly interesting.

   This is the second paragraph of the quotation.  It
   is just as dull as the first paragraph.
\end{quotation}
Another frequently-displayed structure is a list.
The following is an example of an \emph{itemized}
list.
\begin{itemize}
   \item This is the first item of an itemized list.
         Each item in the list is marked with a ``tick''.
         You don't have to worry about what kind of tick
         mark is used.

   \item This is the second item of the list.  It
         contains another list nested inside it.  The inner
         list is an \emph{enumerated} list.
         \begin{enumerate}
            \item This is the first item of an enumerated
                  list that is nested within the itemized list.

            \item This is the second item of the inner list.
                  \LaTeX\ allows you to nest lists deeper than
                  you really should.
         \end{enumerate}
         This is the rest of the second item of the outer
         list.  It is no more interesting than any other
         part of the item.
   \item This is the third item of the list.
\end{itemize}
You can even display poetry.
\begin{verse}
   There is an environment
    for verse \\             % The \\ command separates lines
   Whose features some poets % within a stanza.
   will curse.

                             % One or more blank lines separate stanzas.

   For instead of making\\
   Them do \emph{all} line breaking, \\
   It allows them to put too many words on a line when they'd rather be
   forced to be terse.
\end{verse}

Mathematical formulas may also be displayed.  A
displayed formula
is
one-line long; multiline
formulas require special formatting instructions.
   \[  \ip{\Gamma}{\psi'} = x'' + y^{2} + z_{i}^{n}\]
Don't start a paragraph with a displayed equation,
nor make one a paragraph by itself.

\end{document}               % End of document.

\section{Bench Huge Section 1226}
% synthetic bench section 1226
% This is a sample LaTeX input file.  (Version of 12 August 2004.)
%
% A '%' character causes TeX to ignore all remaining text on the line,
% and is used for comments like this one.

\documentclass{article}      % Specifies the document class

                             % The preamble begins here.
\title{An Example Document}  % Declares the document's title.
\author{Leslie Lamport}      % Declares the author's name.
\date{January 21, 1994}      % Deleting this command produces today's date.

\newcommand{\ip}[2]{(#1, #2)}
                             % Defines \ip{arg1}{arg2} to mean
                             % (arg1, arg2).

%\newcommand{\ip}[2]{\langle #1 | #2\rangle}
                             % This is an alternative definition of
                             % \ip that is commented out.

\begin{document}             % End of preamble and beginning of text.

\maketitle                   % Produces the title.

This is an example input file.  Comparing it with
the output it generates can show you how to
produce a simple document of your own.

\section{Ordinary Text}      % Produces section heading.  Lower-level
                             % sections are begun with similar
                             % \subsection and \subsubsection commands.

The ends  of words and sentences are marked
  by   spaces. It  doesn't matter how many
spaces    you type; one is as good as 100.  The
end of   a line counts as a space.

One   or more   blank lines denote the  end
of  a paragraph.

Since any number of consecutive spaces are treated
like a single one, the formatting of the input
file makes no difference to
      \LaTeX,                % The \LaTeX command generates the LaTeX logo.
but it makes a difference to you.  When you use
\LaTeX, making your input file as easy to read
as possible will be a great help as you write
your document and when you change it.  This sample
file shows how you can add comments to your own input
file.

Because printing is different from typewriting,
there are a number of things that you have to do
differently when preparing an input file than if
you were just typing the document directly.
Quotation marks like
       ``this''
have to be handled specially, as do quotes within
quotes:
       ``\,`this'            % \, separates the double and single quote.
        is what I just
        wrote, not  `that'\,''.

Dashes come in three sizes: an
       intra-word
dash, a medium dash for number ranges like
       1--2,
and a punctuation
       dash---like
this.

A sentence-ending space should be larger than the
space between words within a sentence.  You
sometimes have to type special commands in
conjunction with punctuation characters to get
this right, as in the following sentence.
       Gnats, gnus, etc.\ all  % `\ ' makes an inter-word space.
       begin with G\@.         % \@ marks end-of-sentence punctuation.
You should check the spaces after periods when
reading your output to make sure you haven't
forgotten any special cases.  Generating an
ellipsis
       \ldots\               % `\ ' is needed after `\ldots' because TeX
                             % ignores spaces after command names like \ldots
                             % made from \ + letters.
                             %
                             % Note how a `%' character causes TeX to ignore
                             % the end of the input line, so these blank lines
                             % do not start a new paragraph.
                             %
with the right spacing around the periods requires
a special command.

\LaTeX\ interprets some common characters as
commands, so you must type special commands to
generate them.  These characters include the
following:
       \$ \& \% \# \{ and \}.

In printing, text is usually emphasized with an
       \emph{italic}
type style.

\begin{em}
   A long segment of text can also be emphasized
   in this way.  Text within such a segment can be
   given \emph{additional} emphasis.
\end{em}

It is sometimes necessary to prevent \LaTeX\ from
breaking a line where it might otherwise do so.
This may be at a space, as between the ``Mr.''\ and
``Jones'' in
       ``Mr.~Jones'',        % ~ produces an unbreakable interword space.
or within a word---especially when the word is a
symbol like
       \mbox{\emph{itemnum}}
that makes little sense when hyphenated across
lines.

Footnotes\footnote{This is an example of a footnote.}
pose no problem.

\LaTeX\ is good at typesetting mathematical formulas
like
       \( x-3y + z = 7 \)
or
       \( a_{1} > x^{2n} + y^{2n} > x' \)
or
       \( \ip{A}{B} = \sum_{i} a_{i} b_{i} \).
The spaces you type in a formula are
ignored.  Remember that a letter like
       $x$                   % $ ... $  and  \( ... \)  are equivalent
is a formula when it denotes a mathematical
symbol, and it should be typed as one.

\section{Displayed Text}

Text is displayed by indenting it from the left
margin.  Quotations are commonly displayed.  There
are short quotations
\begin{quote}
   This is a short quotation.  It consists of a
   single paragraph of text.  See how it is formatted.
\end{quote}
and longer ones.
\begin{quotation}
   This is a longer quotation.  It consists of two
   paragraphs of text, neither of which are
   particularly interesting.

   This is the second paragraph of the quotation.  It
   is just as dull as the first paragraph.
\end{quotation}
Another frequently-displayed structure is a list.
The following is an example of an \emph{itemized}
list.
\begin{itemize}
   \item This is the first item of an itemized list.
         Each item in the list is marked with a ``tick''.
         You don't have to worry about what kind of tick
         mark is used.

   \item This is the second item of the list.  It
         contains another list nested inside it.  The inner
         list is an \emph{enumerated} list.
         \begin{enumerate}
            \item This is the first item of an enumerated
                  list that is nested within the itemized list.

            \item This is the second item of the inner list.
                  \LaTeX\ allows you to nest lists deeper than
                  you really should.
         \end{enumerate}
         This is the rest of the second item of the outer
         list.  It is no more interesting than any other
         part of the item.
   \item This is the third item of the list.
\end{itemize}
You can even display poetry.
\begin{verse}
   There is an environment
    for verse \\             % The \\ command separates lines
   Whose features some poets % within a stanza.
   will curse.

                             % One or more blank lines separate stanzas.

   For instead of making\\
   Them do \emph{all} line breaking, \\
   It allows them to put too many words on a line when they'd rather be
   forced to be terse.
\end{verse}

Mathematical formulas may also be displayed.  A
displayed formula
is
one-line long; multiline
formulas require special formatting instructions.
   \[  \ip{\Gamma}{\psi'} = x'' + y^{2} + z_{i}^{n}\]
Don't start a paragraph with a displayed equation,
nor make one a paragraph by itself.

\end{document}               % End of document.

\section{Bench Huge Section 1227}
% synthetic bench section 1227
% This is a sample LaTeX input file.  (Version of 12 August 2004.)
%
% A '%' character causes TeX to ignore all remaining text on the line,
% and is used for comments like this one.

\documentclass{article}      % Specifies the document class

                             % The preamble begins here.
\title{An Example Document}  % Declares the document's title.
\author{Leslie Lamport}      % Declares the author's name.
\date{January 21, 1994}      % Deleting this command produces today's date.

\newcommand{\ip}[2]{(#1, #2)}
                             % Defines \ip{arg1}{arg2} to mean
                             % (arg1, arg2).

%\newcommand{\ip}[2]{\langle #1 | #2\rangle}
                             % This is an alternative definition of
                             % \ip that is commented out.

\begin{document}             % End of preamble and beginning of text.

\maketitle                   % Produces the title.

This is an example input file.  Comparing it with
the output it generates can show you how to
produce a simple document of your own.

\section{Ordinary Text}      % Produces section heading.  Lower-level
                             % sections are begun with similar
                             % \subsection and \subsubsection commands.

The ends  of words and sentences are marked
  by   spaces. It  doesn't matter how many
spaces    you type; one is as good as 100.  The
end of   a line counts as a space.

One   or more   blank lines denote the  end
of  a paragraph.

Since any number of consecutive spaces are treated
like a single one, the formatting of the input
file makes no difference to
      \LaTeX,                % The \LaTeX command generates the LaTeX logo.
but it makes a difference to you.  When you use
\LaTeX, making your input file as easy to read
as possible will be a great help as you write
your document and when you change it.  This sample
file shows how you can add comments to your own input
file.

Because printing is different from typewriting,
there are a number of things that you have to do
differently when preparing an input file than if
you were just typing the document directly.
Quotation marks like
       ``this''
have to be handled specially, as do quotes within
quotes:
       ``\,`this'            % \, separates the double and single quote.
        is what I just
        wrote, not  `that'\,''.

Dashes come in three sizes: an
       intra-word
dash, a medium dash for number ranges like
       1--2,
and a punctuation
       dash---like
this.

A sentence-ending space should be larger than the
space between words within a sentence.  You
sometimes have to type special commands in
conjunction with punctuation characters to get
this right, as in the following sentence.
       Gnats, gnus, etc.\ all  % `\ ' makes an inter-word space.
       begin with G\@.         % \@ marks end-of-sentence punctuation.
You should check the spaces after periods when
reading your output to make sure you haven't
forgotten any special cases.  Generating an
ellipsis
       \ldots\               % `\ ' is needed after `\ldots' because TeX
                             % ignores spaces after command names like \ldots
                             % made from \ + letters.
                             %
                             % Note how a `%' character causes TeX to ignore
                             % the end of the input line, so these blank lines
                             % do not start a new paragraph.
                             %
with the right spacing around the periods requires
a special command.

\LaTeX\ interprets some common characters as
commands, so you must type special commands to
generate them.  These characters include the
following:
       \$ \& \% \# \{ and \}.

In printing, text is usually emphasized with an
       \emph{italic}
type style.

\begin{em}
   A long segment of text can also be emphasized
   in this way.  Text within such a segment can be
   given \emph{additional} emphasis.
\end{em}

It is sometimes necessary to prevent \LaTeX\ from
breaking a line where it might otherwise do so.
This may be at a space, as between the ``Mr.''\ and
``Jones'' in
       ``Mr.~Jones'',        % ~ produces an unbreakable interword space.
or within a word---especially when the word is a
symbol like
       \mbox{\emph{itemnum}}
that makes little sense when hyphenated across
lines.

Footnotes\footnote{This is an example of a footnote.}
pose no problem.

\LaTeX\ is good at typesetting mathematical formulas
like
       \( x-3y + z = 7 \)
or
       \( a_{1} > x^{2n} + y^{2n} > x' \)
or
       \( \ip{A}{B} = \sum_{i} a_{i} b_{i} \).
The spaces you type in a formula are
ignored.  Remember that a letter like
       $x$                   % $ ... $  and  \( ... \)  are equivalent
is a formula when it denotes a mathematical
symbol, and it should be typed as one.

\section{Displayed Text}

Text is displayed by indenting it from the left
margin.  Quotations are commonly displayed.  There
are short quotations
\begin{quote}
   This is a short quotation.  It consists of a
   single paragraph of text.  See how it is formatted.
\end{quote}
and longer ones.
\begin{quotation}
   This is a longer quotation.  It consists of two
   paragraphs of text, neither of which are
   particularly interesting.

   This is the second paragraph of the quotation.  It
   is just as dull as the first paragraph.
\end{quotation}
Another frequently-displayed structure is a list.
The following is an example of an \emph{itemized}
list.
\begin{itemize}
   \item This is the first item of an itemized list.
         Each item in the list is marked with a ``tick''.
         You don't have to worry about what kind of tick
         mark is used.

   \item This is the second item of the list.  It
         contains another list nested inside it.  The inner
         list is an \emph{enumerated} list.
         \begin{enumerate}
            \item This is the first item of an enumerated
                  list that is nested within the itemized list.

            \item This is the second item of the inner list.
                  \LaTeX\ allows you to nest lists deeper than
                  you really should.
         \end{enumerate}
         This is the rest of the second item of the outer
         list.  It is no more interesting than any other
         part of the item.
   \item This is the third item of the list.
\end{itemize}
You can even display poetry.
\begin{verse}
   There is an environment
    for verse \\             % The \\ command separates lines
   Whose features some poets % within a stanza.
   will curse.

                             % One or more blank lines separate stanzas.

   For instead of making\\
   Them do \emph{all} line breaking, \\
   It allows them to put too many words on a line when they'd rather be
   forced to be terse.
\end{verse}

Mathematical formulas may also be displayed.  A
displayed formula
is
one-line long; multiline
formulas require special formatting instructions.
   \[  \ip{\Gamma}{\psi'} = x'' + y^{2} + z_{i}^{n}\]
Don't start a paragraph with a displayed equation,
nor make one a paragraph by itself.

\end{document}               % End of document.

\section{Bench Huge Section 1228}
% synthetic bench section 1228
% This is a sample LaTeX input file.  (Version of 12 August 2004.)
%
% A '%' character causes TeX to ignore all remaining text on the line,
% and is used for comments like this one.

\documentclass{article}      % Specifies the document class

                             % The preamble begins here.
\title{An Example Document}  % Declares the document's title.
\author{Leslie Lamport}      % Declares the author's name.
\date{January 21, 1994}      % Deleting this command produces today's date.

\newcommand{\ip}[2]{(#1, #2)}
                             % Defines \ip{arg1}{arg2} to mean
                             % (arg1, arg2).

%\newcommand{\ip}[2]{\langle #1 | #2\rangle}
                             % This is an alternative definition of
                             % \ip that is commented out.

\begin{document}             % End of preamble and beginning of text.

\maketitle                   % Produces the title.

This is an example input file.  Comparing it with
the output it generates can show you how to
produce a simple document of your own.

\section{Ordinary Text}      % Produces section heading.  Lower-level
                             % sections are begun with similar
                             % \subsection and \subsubsection commands.

The ends  of words and sentences are marked
  by   spaces. It  doesn't matter how many
spaces    you type; one is as good as 100.  The
end of   a line counts as a space.

One   or more   blank lines denote the  end
of  a paragraph.

Since any number of consecutive spaces are treated
like a single one, the formatting of the input
file makes no difference to
      \LaTeX,                % The \LaTeX command generates the LaTeX logo.
but it makes a difference to you.  When you use
\LaTeX, making your input file as easy to read
as possible will be a great help as you write
your document and when you change it.  This sample
file shows how you can add comments to your own input
file.

Because printing is different from typewriting,
there are a number of things that you have to do
differently when preparing an input file than if
you were just typing the document directly.
Quotation marks like
       ``this''
have to be handled specially, as do quotes within
quotes:
       ``\,`this'            % \, separates the double and single quote.
        is what I just
        wrote, not  `that'\,''.

Dashes come in three sizes: an
       intra-word
dash, a medium dash for number ranges like
       1--2,
and a punctuation
       dash---like
this.

A sentence-ending space should be larger than the
space between words within a sentence.  You
sometimes have to type special commands in
conjunction with punctuation characters to get
this right, as in the following sentence.
       Gnats, gnus, etc.\ all  % `\ ' makes an inter-word space.
       begin with G\@.         % \@ marks end-of-sentence punctuation.
You should check the spaces after periods when
reading your output to make sure you haven't
forgotten any special cases.  Generating an
ellipsis
       \ldots\               % `\ ' is needed after `\ldots' because TeX
                             % ignores spaces after command names like \ldots
                             % made from \ + letters.
                             %
                             % Note how a `%' character causes TeX to ignore
                             % the end of the input line, so these blank lines
                             % do not start a new paragraph.
                             %
with the right spacing around the periods requires
a special command.

\LaTeX\ interprets some common characters as
commands, so you must type special commands to
generate them.  These characters include the
following:
       \$ \& \% \# \{ and \}.

In printing, text is usually emphasized with an
       \emph{italic}
type style.

\begin{em}
   A long segment of text can also be emphasized
   in this way.  Text within such a segment can be
   given \emph{additional} emphasis.
\end{em}

It is sometimes necessary to prevent \LaTeX\ from
breaking a line where it might otherwise do so.
This may be at a space, as between the ``Mr.''\ and
``Jones'' in
       ``Mr.~Jones'',        % ~ produces an unbreakable interword space.
or within a word---especially when the word is a
symbol like
       \mbox{\emph{itemnum}}
that makes little sense when hyphenated across
lines.

Footnotes\footnote{This is an example of a footnote.}
pose no problem.

\LaTeX\ is good at typesetting mathematical formulas
like
       \( x-3y + z = 7 \)
or
       \( a_{1} > x^{2n} + y^{2n} > x' \)
or
       \( \ip{A}{B} = \sum_{i} a_{i} b_{i} \).
The spaces you type in a formula are
ignored.  Remember that a letter like
       $x$                   % $ ... $  and  \( ... \)  are equivalent
is a formula when it denotes a mathematical
symbol, and it should be typed as one.

\section{Displayed Text}

Text is displayed by indenting it from the left
margin.  Quotations are commonly displayed.  There
are short quotations
\begin{quote}
   This is a short quotation.  It consists of a
   single paragraph of text.  See how it is formatted.
\end{quote}
and longer ones.
\begin{quotation}
   This is a longer quotation.  It consists of two
   paragraphs of text, neither of which are
   particularly interesting.

   This is the second paragraph of the quotation.  It
   is just as dull as the first paragraph.
\end{quotation}
Another frequently-displayed structure is a list.
The following is an example of an \emph{itemized}
list.
\begin{itemize}
   \item This is the first item of an itemized list.
         Each item in the list is marked with a ``tick''.
         You don't have to worry about what kind of tick
         mark is used.

   \item This is the second item of the list.  It
         contains another list nested inside it.  The inner
         list is an \emph{enumerated} list.
         \begin{enumerate}
            \item This is the first item of an enumerated
                  list that is nested within the itemized list.

            \item This is the second item of the inner list.
                  \LaTeX\ allows you to nest lists deeper than
                  you really should.
         \end{enumerate}
         This is the rest of the second item of the outer
         list.  It is no more interesting than any other
         part of the item.
   \item This is the third item of the list.
\end{itemize}
You can even display poetry.
\begin{verse}
   There is an environment
    for verse \\             % The \\ command separates lines
   Whose features some poets % within a stanza.
   will curse.

                             % One or more blank lines separate stanzas.

   For instead of making\\
   Them do \emph{all} line breaking, \\
   It allows them to put too many words on a line when they'd rather be
   forced to be terse.
\end{verse}

Mathematical formulas may also be displayed.  A
displayed formula
is
one-line long; multiline
formulas require special formatting instructions.
   \[  \ip{\Gamma}{\psi'} = x'' + y^{2} + z_{i}^{n}\]
Don't start a paragraph with a displayed equation,
nor make one a paragraph by itself.

\end{document}               % End of document.

\section{Bench Huge Section 1229}
% synthetic bench section 1229
% This is a sample LaTeX input file.  (Version of 12 August 2004.)
%
% A '%' character causes TeX to ignore all remaining text on the line,
% and is used for comments like this one.

\documentclass{article}      % Specifies the document class

                             % The preamble begins here.
\title{An Example Document}  % Declares the document's title.
\author{Leslie Lamport}      % Declares the author's name.
\date{January 21, 1994}      % Deleting this command produces today's date.

\newcommand{\ip}[2]{(#1, #2)}
                             % Defines \ip{arg1}{arg2} to mean
                             % (arg1, arg2).

%\newcommand{\ip}[2]{\langle #1 | #2\rangle}
                             % This is an alternative definition of
                             % \ip that is commented out.

\begin{document}             % End of preamble and beginning of text.

\maketitle                   % Produces the title.

This is an example input file.  Comparing it with
the output it generates can show you how to
produce a simple document of your own.

\section{Ordinary Text}      % Produces section heading.  Lower-level
                             % sections are begun with similar
                             % \subsection and \subsubsection commands.

The ends  of words and sentences are marked
  by   spaces. It  doesn't matter how many
spaces    you type; one is as good as 100.  The
end of   a line counts as a space.

One   or more   blank lines denote the  end
of  a paragraph.

Since any number of consecutive spaces are treated
like a single one, the formatting of the input
file makes no difference to
      \LaTeX,                % The \LaTeX command generates the LaTeX logo.
but it makes a difference to you.  When you use
\LaTeX, making your input file as easy to read
as possible will be a great help as you write
your document and when you change it.  This sample
file shows how you can add comments to your own input
file.

Because printing is different from typewriting,
there are a number of things that you have to do
differently when preparing an input file than if
you were just typing the document directly.
Quotation marks like
       ``this''
have to be handled specially, as do quotes within
quotes:
       ``\,`this'            % \, separates the double and single quote.
        is what I just
        wrote, not  `that'\,''.

Dashes come in three sizes: an
       intra-word
dash, a medium dash for number ranges like
       1--2,
and a punctuation
       dash---like
this.

A sentence-ending space should be larger than the
space between words within a sentence.  You
sometimes have to type special commands in
conjunction with punctuation characters to get
this right, as in the following sentence.
       Gnats, gnus, etc.\ all  % `\ ' makes an inter-word space.
       begin with G\@.         % \@ marks end-of-sentence punctuation.
You should check the spaces after periods when
reading your output to make sure you haven't
forgotten any special cases.  Generating an
ellipsis
       \ldots\               % `\ ' is needed after `\ldots' because TeX
                             % ignores spaces after command names like \ldots
                             % made from \ + letters.
                             %
                             % Note how a `%' character causes TeX to ignore
                             % the end of the input line, so these blank lines
                             % do not start a new paragraph.
                             %
with the right spacing around the periods requires
a special command.

\LaTeX\ interprets some common characters as
commands, so you must type special commands to
generate them.  These characters include the
following:
       \$ \& \% \# \{ and \}.

In printing, text is usually emphasized with an
       \emph{italic}
type style.

\begin{em}
   A long segment of text can also be emphasized
   in this way.  Text within such a segment can be
   given \emph{additional} emphasis.
\end{em}

It is sometimes necessary to prevent \LaTeX\ from
breaking a line where it might otherwise do so.
This may be at a space, as between the ``Mr.''\ and
``Jones'' in
       ``Mr.~Jones'',        % ~ produces an unbreakable interword space.
or within a word---especially when the word is a
symbol like
       \mbox{\emph{itemnum}}
that makes little sense when hyphenated across
lines.

Footnotes\footnote{This is an example of a footnote.}
pose no problem.

\LaTeX\ is good at typesetting mathematical formulas
like
       \( x-3y + z = 7 \)
or
       \( a_{1} > x^{2n} + y^{2n} > x' \)
or
       \( \ip{A}{B} = \sum_{i} a_{i} b_{i} \).
The spaces you type in a formula are
ignored.  Remember that a letter like
       $x$                   % $ ... $  and  \( ... \)  are equivalent
is a formula when it denotes a mathematical
symbol, and it should be typed as one.

\section{Displayed Text}

Text is displayed by indenting it from the left
margin.  Quotations are commonly displayed.  There
are short quotations
\begin{quote}
   This is a short quotation.  It consists of a
   single paragraph of text.  See how it is formatted.
\end{quote}
and longer ones.
\begin{quotation}
   This is a longer quotation.  It consists of two
   paragraphs of text, neither of which are
   particularly interesting.

   This is the second paragraph of the quotation.  It
   is just as dull as the first paragraph.
\end{quotation}
Another frequently-displayed structure is a list.
The following is an example of an \emph{itemized}
list.
\begin{itemize}
   \item This is the first item of an itemized list.
         Each item in the list is marked with a ``tick''.
         You don't have to worry about what kind of tick
         mark is used.

   \item This is the second item of the list.  It
         contains another list nested inside it.  The inner
         list is an \emph{enumerated} list.
         \begin{enumerate}
            \item This is the first item of an enumerated
                  list that is nested within the itemized list.

            \item This is the second item of the inner list.
                  \LaTeX\ allows you to nest lists deeper than
                  you really should.
         \end{enumerate}
         This is the rest of the second item of the outer
         list.  It is no more interesting than any other
         part of the item.
   \item This is the third item of the list.
\end{itemize}
You can even display poetry.
\begin{verse}
   There is an environment
    for verse \\             % The \\ command separates lines
   Whose features some poets % within a stanza.
   will curse.

                             % One or more blank lines separate stanzas.

   For instead of making\\
   Them do \emph{all} line breaking, \\
   It allows them to put too many words on a line when they'd rather be
   forced to be terse.
\end{verse}

Mathematical formulas may also be displayed.  A
displayed formula
is
one-line long; multiline
formulas require special formatting instructions.
   \[  \ip{\Gamma}{\psi'} = x'' + y^{2} + z_{i}^{n}\]
Don't start a paragraph with a displayed equation,
nor make one a paragraph by itself.

\end{document}               % End of document.

\section{Bench Huge Section 1230}
% synthetic bench section 1230
% This is a sample LaTeX input file.  (Version of 12 August 2004.)
%
% A '%' character causes TeX to ignore all remaining text on the line,
% and is used for comments like this one.

\documentclass{article}      % Specifies the document class

                             % The preamble begins here.
\title{An Example Document}  % Declares the document's title.
\author{Leslie Lamport}      % Declares the author's name.
\date{January 21, 1994}      % Deleting this command produces today's date.

\newcommand{\ip}[2]{(#1, #2)}
                             % Defines \ip{arg1}{arg2} to mean
                             % (arg1, arg2).

%\newcommand{\ip}[2]{\langle #1 | #2\rangle}
                             % This is an alternative definition of
                             % \ip that is commented out.

\begin{document}             % End of preamble and beginning of text.

\maketitle                   % Produces the title.

This is an example input file.  Comparing it with
the output it generates can show you how to
produce a simple document of your own.

\section{Ordinary Text}      % Produces section heading.  Lower-level
                             % sections are begun with similar
                             % \subsection and \subsubsection commands.

The ends  of words and sentences are marked
  by   spaces. It  doesn't matter how many
spaces    you type; one is as good as 100.  The
end of   a line counts as a space.

One   or more   blank lines denote the  end
of  a paragraph.

Since any number of consecutive spaces are treated
like a single one, the formatting of the input
file makes no difference to
      \LaTeX,                % The \LaTeX command generates the LaTeX logo.
but it makes a difference to you.  When you use
\LaTeX, making your input file as easy to read
as possible will be a great help as you write
your document and when you change it.  This sample
file shows how you can add comments to your own input
file.

Because printing is different from typewriting,
there are a number of things that you have to do
differently when preparing an input file than if
you were just typing the document directly.
Quotation marks like
       ``this''
have to be handled specially, as do quotes within
quotes:
       ``\,`this'            % \, separates the double and single quote.
        is what I just
        wrote, not  `that'\,''.

Dashes come in three sizes: an
       intra-word
dash, a medium dash for number ranges like
       1--2,
and a punctuation
       dash---like
this.

A sentence-ending space should be larger than the
space between words within a sentence.  You
sometimes have to type special commands in
conjunction with punctuation characters to get
this right, as in the following sentence.
       Gnats, gnus, etc.\ all  % `\ ' makes an inter-word space.
       begin with G\@.         % \@ marks end-of-sentence punctuation.
You should check the spaces after periods when
reading your output to make sure you haven't
forgotten any special cases.  Generating an
ellipsis
       \ldots\               % `\ ' is needed after `\ldots' because TeX
                             % ignores spaces after command names like \ldots
                             % made from \ + letters.
                             %
                             % Note how a `%' character causes TeX to ignore
                             % the end of the input line, so these blank lines
                             % do not start a new paragraph.
                             %
with the right spacing around the periods requires
a special command.

\LaTeX\ interprets some common characters as
commands, so you must type special commands to
generate them.  These characters include the
following:
       \$ \& \% \# \{ and \}.

In printing, text is usually emphasized with an
       \emph{italic}
type style.

\begin{em}
   A long segment of text can also be emphasized
   in this way.  Text within such a segment can be
   given \emph{additional} emphasis.
\end{em}

It is sometimes necessary to prevent \LaTeX\ from
breaking a line where it might otherwise do so.
This may be at a space, as between the ``Mr.''\ and
``Jones'' in
       ``Mr.~Jones'',        % ~ produces an unbreakable interword space.
or within a word---especially when the word is a
symbol like
       \mbox{\emph{itemnum}}
that makes little sense when hyphenated across
lines.

Footnotes\footnote{This is an example of a footnote.}
pose no problem.

\LaTeX\ is good at typesetting mathematical formulas
like
       \( x-3y + z = 7 \)
or
       \( a_{1} > x^{2n} + y^{2n} > x' \)
or
       \( \ip{A}{B} = \sum_{i} a_{i} b_{i} \).
The spaces you type in a formula are
ignored.  Remember that a letter like
       $x$                   % $ ... $  and  \( ... \)  are equivalent
is a formula when it denotes a mathematical
symbol, and it should be typed as one.

\section{Displayed Text}

Text is displayed by indenting it from the left
margin.  Quotations are commonly displayed.  There
are short quotations
\begin{quote}
   This is a short quotation.  It consists of a
   single paragraph of text.  See how it is formatted.
\end{quote}
and longer ones.
\begin{quotation}
   This is a longer quotation.  It consists of two
   paragraphs of text, neither of which are
   particularly interesting.

   This is the second paragraph of the quotation.  It
   is just as dull as the first paragraph.
\end{quotation}
Another frequently-displayed structure is a list.
The following is an example of an \emph{itemized}
list.
\begin{itemize}
   \item This is the first item of an itemized list.
         Each item in the list is marked with a ``tick''.
         You don't have to worry about what kind of tick
         mark is used.

   \item This is the second item of the list.  It
         contains another list nested inside it.  The inner
         list is an \emph{enumerated} list.
         \begin{enumerate}
            \item This is the first item of an enumerated
                  list that is nested within the itemized list.

            \item This is the second item of the inner list.
                  \LaTeX\ allows you to nest lists deeper than
                  you really should.
         \end{enumerate}
         This is the rest of the second item of the outer
         list.  It is no more interesting than any other
         part of the item.
   \item This is the third item of the list.
\end{itemize}
You can even display poetry.
\begin{verse}
   There is an environment
    for verse \\             % The \\ command separates lines
   Whose features some poets % within a stanza.
   will curse.

                             % One or more blank lines separate stanzas.

   For instead of making\\
   Them do \emph{all} line breaking, \\
   It allows them to put too many words on a line when they'd rather be
   forced to be terse.
\end{verse}

Mathematical formulas may also be displayed.  A
displayed formula
is
one-line long; multiline
formulas require special formatting instructions.
   \[  \ip{\Gamma}{\psi'} = x'' + y^{2} + z_{i}^{n}\]
Don't start a paragraph with a displayed equation,
nor make one a paragraph by itself.

\end{document}               % End of document.

\section{Bench Huge Section 1231}
% synthetic bench section 1231
% This is a sample LaTeX input file.  (Version of 12 August 2004.)
%
% A '%' character causes TeX to ignore all remaining text on the line,
% and is used for comments like this one.

\documentclass{article}      % Specifies the document class

                             % The preamble begins here.
\title{An Example Document}  % Declares the document's title.
\author{Leslie Lamport}      % Declares the author's name.
\date{January 21, 1994}      % Deleting this command produces today's date.

\newcommand{\ip}[2]{(#1, #2)}
                             % Defines \ip{arg1}{arg2} to mean
                             % (arg1, arg2).

%\newcommand{\ip}[2]{\langle #1 | #2\rangle}
                             % This is an alternative definition of
                             % \ip that is commented out.

\begin{document}             % End of preamble and beginning of text.

\maketitle                   % Produces the title.

This is an example input file.  Comparing it with
the output it generates can show you how to
produce a simple document of your own.

\section{Ordinary Text}      % Produces section heading.  Lower-level
                             % sections are begun with similar
                             % \subsection and \subsubsection commands.

The ends  of words and sentences are marked
  by   spaces. It  doesn't matter how many
spaces    you type; one is as good as 100.  The
end of   a line counts as a space.

One   or more   blank lines denote the  end
of  a paragraph.

Since any number of consecutive spaces are treated
like a single one, the formatting of the input
file makes no difference to
      \LaTeX,                % The \LaTeX command generates the LaTeX logo.
but it makes a difference to you.  When you use
\LaTeX, making your input file as easy to read
as possible will be a great help as you write
your document and when you change it.  This sample
file shows how you can add comments to your own input
file.

Because printing is different from typewriting,
there are a number of things that you have to do
differently when preparing an input file than if
you were just typing the document directly.
Quotation marks like
       ``this''
have to be handled specially, as do quotes within
quotes:
       ``\,`this'            % \, separates the double and single quote.
        is what I just
        wrote, not  `that'\,''.

Dashes come in three sizes: an
       intra-word
dash, a medium dash for number ranges like
       1--2,
and a punctuation
       dash---like
this.

A sentence-ending space should be larger than the
space between words within a sentence.  You
sometimes have to type special commands in
conjunction with punctuation characters to get
this right, as in the following sentence.
       Gnats, gnus, etc.\ all  % `\ ' makes an inter-word space.
       begin with G\@.         % \@ marks end-of-sentence punctuation.
You should check the spaces after periods when
reading your output to make sure you haven't
forgotten any special cases.  Generating an
ellipsis
       \ldots\               % `\ ' is needed after `\ldots' because TeX
                             % ignores spaces after command names like \ldots
                             % made from \ + letters.
                             %
                             % Note how a `%' character causes TeX to ignore
                             % the end of the input line, so these blank lines
                             % do not start a new paragraph.
                             %
with the right spacing around the periods requires
a special command.

\LaTeX\ interprets some common characters as
commands, so you must type special commands to
generate them.  These characters include the
following:
       \$ \& \% \# \{ and \}.

In printing, text is usually emphasized with an
       \emph{italic}
type style.

\begin{em}
   A long segment of text can also be emphasized
   in this way.  Text within such a segment can be
   given \emph{additional} emphasis.
\end{em}

It is sometimes necessary to prevent \LaTeX\ from
breaking a line where it might otherwise do so.
This may be at a space, as between the ``Mr.''\ and
``Jones'' in
       ``Mr.~Jones'',        % ~ produces an unbreakable interword space.
or within a word---especially when the word is a
symbol like
       \mbox{\emph{itemnum}}
that makes little sense when hyphenated across
lines.

Footnotes\footnote{This is an example of a footnote.}
pose no problem.

\LaTeX\ is good at typesetting mathematical formulas
like
       \( x-3y + z = 7 \)
or
       \( a_{1} > x^{2n} + y^{2n} > x' \)
or
       \( \ip{A}{B} = \sum_{i} a_{i} b_{i} \).
The spaces you type in a formula are
ignored.  Remember that a letter like
       $x$                   % $ ... $  and  \( ... \)  are equivalent
is a formula when it denotes a mathematical
symbol, and it should be typed as one.

\section{Displayed Text}

Text is displayed by indenting it from the left
margin.  Quotations are commonly displayed.  There
are short quotations
\begin{quote}
   This is a short quotation.  It consists of a
   single paragraph of text.  See how it is formatted.
\end{quote}
and longer ones.
\begin{quotation}
   This is a longer quotation.  It consists of two
   paragraphs of text, neither of which are
   particularly interesting.

   This is the second paragraph of the quotation.  It
   is just as dull as the first paragraph.
\end{quotation}
Another frequently-displayed structure is a list.
The following is an example of an \emph{itemized}
list.
\begin{itemize}
   \item This is the first item of an itemized list.
         Each item in the list is marked with a ``tick''.
         You don't have to worry about what kind of tick
         mark is used.

   \item This is the second item of the list.  It
         contains another list nested inside it.  The inner
         list is an \emph{enumerated} list.
         \begin{enumerate}
            \item This is the first item of an enumerated
                  list that is nested within the itemized list.

            \item This is the second item of the inner list.
                  \LaTeX\ allows you to nest lists deeper than
                  you really should.
         \end{enumerate}
         This is the rest of the second item of the outer
         list.  It is no more interesting than any other
         part of the item.
   \item This is the third item of the list.
\end{itemize}
You can even display poetry.
\begin{verse}
   There is an environment
    for verse \\             % The \\ command separates lines
   Whose features some poets % within a stanza.
   will curse.

                             % One or more blank lines separate stanzas.

   For instead of making\\
   Them do \emph{all} line breaking, \\
   It allows them to put too many words on a line when they'd rather be
   forced to be terse.
\end{verse}

Mathematical formulas may also be displayed.  A
displayed formula
is
one-line long; multiline
formulas require special formatting instructions.
   \[  \ip{\Gamma}{\psi'} = x'' + y^{2} + z_{i}^{n}\]
Don't start a paragraph with a displayed equation,
nor make one a paragraph by itself.

\end{document}               % End of document.

\section{Bench Huge Section 1232}
% synthetic bench section 1232
% This is a sample LaTeX input file.  (Version of 12 August 2004.)
%
% A '%' character causes TeX to ignore all remaining text on the line,
% and is used for comments like this one.

\documentclass{article}      % Specifies the document class

                             % The preamble begins here.
\title{An Example Document}  % Declares the document's title.
\author{Leslie Lamport}      % Declares the author's name.
\date{January 21, 1994}      % Deleting this command produces today's date.

\newcommand{\ip}[2]{(#1, #2)}
                             % Defines \ip{arg1}{arg2} to mean
                             % (arg1, arg2).

%\newcommand{\ip}[2]{\langle #1 | #2\rangle}
                             % This is an alternative definition of
                             % \ip that is commented out.

\begin{document}             % End of preamble and beginning of text.

\maketitle                   % Produces the title.

This is an example input file.  Comparing it with
the output it generates can show you how to
produce a simple document of your own.

\section{Ordinary Text}      % Produces section heading.  Lower-level
                             % sections are begun with similar
                             % \subsection and \subsubsection commands.

The ends  of words and sentences are marked
  by   spaces. It  doesn't matter how many
spaces    you type; one is as good as 100.  The
end of   a line counts as a space.

One   or more   blank lines denote the  end
of  a paragraph.

Since any number of consecutive spaces are treated
like a single one, the formatting of the input
file makes no difference to
      \LaTeX,                % The \LaTeX command generates the LaTeX logo.
but it makes a difference to you.  When you use
\LaTeX, making your input file as easy to read
as possible will be a great help as you write
your document and when you change it.  This sample
file shows how you can add comments to your own input
file.

Because printing is different from typewriting,
there are a number of things that you have to do
differently when preparing an input file than if
you were just typing the document directly.
Quotation marks like
       ``this''
have to be handled specially, as do quotes within
quotes:
       ``\,`this'            % \, separates the double and single quote.
        is what I just
        wrote, not  `that'\,''.

Dashes come in three sizes: an
       intra-word
dash, a medium dash for number ranges like
       1--2,
and a punctuation
       dash---like
this.

A sentence-ending space should be larger than the
space between words within a sentence.  You
sometimes have to type special commands in
conjunction with punctuation characters to get
this right, as in the following sentence.
       Gnats, gnus, etc.\ all  % `\ ' makes an inter-word space.
       begin with G\@.         % \@ marks end-of-sentence punctuation.
You should check the spaces after periods when
reading your output to make sure you haven't
forgotten any special cases.  Generating an
ellipsis
       \ldots\               % `\ ' is needed after `\ldots' because TeX
                             % ignores spaces after command names like \ldots
                             % made from \ + letters.
                             %
                             % Note how a `%' character causes TeX to ignore
                             % the end of the input line, so these blank lines
                             % do not start a new paragraph.
                             %
with the right spacing around the periods requires
a special command.

\LaTeX\ interprets some common characters as
commands, so you must type special commands to
generate them.  These characters include the
following:
       \$ \& \% \# \{ and \}.

In printing, text is usually emphasized with an
       \emph{italic}
type style.

\begin{em}
   A long segment of text can also be emphasized
   in this way.  Text within such a segment can be
   given \emph{additional} emphasis.
\end{em}

It is sometimes necessary to prevent \LaTeX\ from
breaking a line where it might otherwise do so.
This may be at a space, as between the ``Mr.''\ and
``Jones'' in
       ``Mr.~Jones'',        % ~ produces an unbreakable interword space.
or within a word---especially when the word is a
symbol like
       \mbox{\emph{itemnum}}
that makes little sense when hyphenated across
lines.

Footnotes\footnote{This is an example of a footnote.}
pose no problem.

\LaTeX\ is good at typesetting mathematical formulas
like
       \( x-3y + z = 7 \)
or
       \( a_{1} > x^{2n} + y^{2n} > x' \)
or
       \( \ip{A}{B} = \sum_{i} a_{i} b_{i} \).
The spaces you type in a formula are
ignored.  Remember that a letter like
       $x$                   % $ ... $  and  \( ... \)  are equivalent
is a formula when it denotes a mathematical
symbol, and it should be typed as one.

\section{Displayed Text}

Text is displayed by indenting it from the left
margin.  Quotations are commonly displayed.  There
are short quotations
\begin{quote}
   This is a short quotation.  It consists of a
   single paragraph of text.  See how it is formatted.
\end{quote}
and longer ones.
\begin{quotation}
   This is a longer quotation.  It consists of two
   paragraphs of text, neither of which are
   particularly interesting.

   This is the second paragraph of the quotation.  It
   is just as dull as the first paragraph.
\end{quotation}
Another frequently-displayed structure is a list.
The following is an example of an \emph{itemized}
list.
\begin{itemize}
   \item This is the first item of an itemized list.
         Each item in the list is marked with a ``tick''.
         You don't have to worry about what kind of tick
         mark is used.

   \item This is the second item of the list.  It
         contains another list nested inside it.  The inner
         list is an \emph{enumerated} list.
         \begin{enumerate}
            \item This is the first item of an enumerated
                  list that is nested within the itemized list.

            \item This is the second item of the inner list.
                  \LaTeX\ allows you to nest lists deeper than
                  you really should.
         \end{enumerate}
         This is the rest of the second item of the outer
         list.  It is no more interesting than any other
         part of the item.
   \item This is the third item of the list.
\end{itemize}
You can even display poetry.
\begin{verse}
   There is an environment
    for verse \\             % The \\ command separates lines
   Whose features some poets % within a stanza.
   will curse.

                             % One or more blank lines separate stanzas.

   For instead of making\\
   Them do \emph{all} line breaking, \\
   It allows them to put too many words on a line when they'd rather be
   forced to be terse.
\end{verse}

Mathematical formulas may also be displayed.  A
displayed formula
is
one-line long; multiline
formulas require special formatting instructions.
   \[  \ip{\Gamma}{\psi'} = x'' + y^{2} + z_{i}^{n}\]
Don't start a paragraph with a displayed equation,
nor make one a paragraph by itself.

\end{document}               % End of document.

\section{Bench Huge Section 1233}
% synthetic bench section 1233
% This is a sample LaTeX input file.  (Version of 12 August 2004.)
%
% A '%' character causes TeX to ignore all remaining text on the line,
% and is used for comments like this one.

\documentclass{article}      % Specifies the document class

                             % The preamble begins here.
\title{An Example Document}  % Declares the document's title.
\author{Leslie Lamport}      % Declares the author's name.
\date{January 21, 1994}      % Deleting this command produces today's date.

\newcommand{\ip}[2]{(#1, #2)}
                             % Defines \ip{arg1}{arg2} to mean
                             % (arg1, arg2).

%\newcommand{\ip}[2]{\langle #1 | #2\rangle}
                             % This is an alternative definition of
                             % \ip that is commented out.

\begin{document}             % End of preamble and beginning of text.

\maketitle                   % Produces the title.

This is an example input file.  Comparing it with
the output it generates can show you how to
produce a simple document of your own.

\section{Ordinary Text}      % Produces section heading.  Lower-level
                             % sections are begun with similar
                             % \subsection and \subsubsection commands.

The ends  of words and sentences are marked
  by   spaces. It  doesn't matter how many
spaces    you type; one is as good as 100.  The
end of   a line counts as a space.

One   or more   blank lines denote the  end
of  a paragraph.

Since any number of consecutive spaces are treated
like a single one, the formatting of the input
file makes no difference to
      \LaTeX,                % The \LaTeX command generates the LaTeX logo.
but it makes a difference to you.  When you use
\LaTeX, making your input file as easy to read
as possible will be a great help as you write
your document and when you change it.  This sample
file shows how you can add comments to your own input
file.

Because printing is different from typewriting,
there are a number of things that you have to do
differently when preparing an input file than if
you were just typing the document directly.
Quotation marks like
       ``this''
have to be handled specially, as do quotes within
quotes:
       ``\,`this'            % \, separates the double and single quote.
        is what I just
        wrote, not  `that'\,''.

Dashes come in three sizes: an
       intra-word
dash, a medium dash for number ranges like
       1--2,
and a punctuation
       dash---like
this.

A sentence-ending space should be larger than the
space between words within a sentence.  You
sometimes have to type special commands in
conjunction with punctuation characters to get
this right, as in the following sentence.
       Gnats, gnus, etc.\ all  % `\ ' makes an inter-word space.
       begin with G\@.         % \@ marks end-of-sentence punctuation.
You should check the spaces after periods when
reading your output to make sure you haven't
forgotten any special cases.  Generating an
ellipsis
       \ldots\               % `\ ' is needed after `\ldots' because TeX
                             % ignores spaces after command names like \ldots
                             % made from \ + letters.
                             %
                             % Note how a `%' character causes TeX to ignore
                             % the end of the input line, so these blank lines
                             % do not start a new paragraph.
                             %
with the right spacing around the periods requires
a special command.

\LaTeX\ interprets some common characters as
commands, so you must type special commands to
generate them.  These characters include the
following:
       \$ \& \% \# \{ and \}.

In printing, text is usually emphasized with an
       \emph{italic}
type style.

\begin{em}
   A long segment of text can also be emphasized
   in this way.  Text within such a segment can be
   given \emph{additional} emphasis.
\end{em}

It is sometimes necessary to prevent \LaTeX\ from
breaking a line where it might otherwise do so.
This may be at a space, as between the ``Mr.''\ and
``Jones'' in
       ``Mr.~Jones'',        % ~ produces an unbreakable interword space.
or within a word---especially when the word is a
symbol like
       \mbox{\emph{itemnum}}
that makes little sense when hyphenated across
lines.

Footnotes\footnote{This is an example of a footnote.}
pose no problem.

\LaTeX\ is good at typesetting mathematical formulas
like
       \( x-3y + z = 7 \)
or
       \( a_{1} > x^{2n} + y^{2n} > x' \)
or
       \( \ip{A}{B} = \sum_{i} a_{i} b_{i} \).
The spaces you type in a formula are
ignored.  Remember that a letter like
       $x$                   % $ ... $  and  \( ... \)  are equivalent
is a formula when it denotes a mathematical
symbol, and it should be typed as one.

\section{Displayed Text}

Text is displayed by indenting it from the left
margin.  Quotations are commonly displayed.  There
are short quotations
\begin{quote}
   This is a short quotation.  It consists of a
   single paragraph of text.  See how it is formatted.
\end{quote}
and longer ones.
\begin{quotation}
   This is a longer quotation.  It consists of two
   paragraphs of text, neither of which are
   particularly interesting.

   This is the second paragraph of the quotation.  It
   is just as dull as the first paragraph.
\end{quotation}
Another frequently-displayed structure is a list.
The following is an example of an \emph{itemized}
list.
\begin{itemize}
   \item This is the first item of an itemized list.
         Each item in the list is marked with a ``tick''.
         You don't have to worry about what kind of tick
         mark is used.

   \item This is the second item of the list.  It
         contains another list nested inside it.  The inner
         list is an \emph{enumerated} list.
         \begin{enumerate}
            \item This is the first item of an enumerated
                  list that is nested within the itemized list.

            \item This is the second item of the inner list.
                  \LaTeX\ allows you to nest lists deeper than
                  you really should.
         \end{enumerate}
         This is the rest of the second item of the outer
         list.  It is no more interesting than any other
         part of the item.
   \item This is the third item of the list.
\end{itemize}
You can even display poetry.
\begin{verse}
   There is an environment
    for verse \\             % The \\ command separates lines
   Whose features some poets % within a stanza.
   will curse.

                             % One or more blank lines separate stanzas.

   For instead of making\\
   Them do \emph{all} line breaking, \\
   It allows them to put too many words on a line when they'd rather be
   forced to be terse.
\end{verse}

Mathematical formulas may also be displayed.  A
displayed formula
is
one-line long; multiline
formulas require special formatting instructions.
   \[  \ip{\Gamma}{\psi'} = x'' + y^{2} + z_{i}^{n}\]
Don't start a paragraph with a displayed equation,
nor make one a paragraph by itself.

\end{document}               % End of document.

\section{Bench Huge Section 1234}
% synthetic bench section 1234
% This is a sample LaTeX input file.  (Version of 12 August 2004.)
%
% A '%' character causes TeX to ignore all remaining text on the line,
% and is used for comments like this one.

\documentclass{article}      % Specifies the document class

                             % The preamble begins here.
\title{An Example Document}  % Declares the document's title.
\author{Leslie Lamport}      % Declares the author's name.
\date{January 21, 1994}      % Deleting this command produces today's date.

\newcommand{\ip}[2]{(#1, #2)}
                             % Defines \ip{arg1}{arg2} to mean
                             % (arg1, arg2).

%\newcommand{\ip}[2]{\langle #1 | #2\rangle}
                             % This is an alternative definition of
                             % \ip that is commented out.

\begin{document}             % End of preamble and beginning of text.

\maketitle                   % Produces the title.

This is an example input file.  Comparing it with
the output it generates can show you how to
produce a simple document of your own.

\section{Ordinary Text}      % Produces section heading.  Lower-level
                             % sections are begun with similar
                             % \subsection and \subsubsection commands.

The ends  of words and sentences are marked
  by   spaces. It  doesn't matter how many
spaces    you type; one is as good as 100.  The
end of   a line counts as a space.

One   or more   blank lines denote the  end
of  a paragraph.

Since any number of consecutive spaces are treated
like a single one, the formatting of the input
file makes no difference to
      \LaTeX,                % The \LaTeX command generates the LaTeX logo.
but it makes a difference to you.  When you use
\LaTeX, making your input file as easy to read
as possible will be a great help as you write
your document and when you change it.  This sample
file shows how you can add comments to your own input
file.

Because printing is different from typewriting,
there are a number of things that you have to do
differently when preparing an input file than if
you were just typing the document directly.
Quotation marks like
       ``this''
have to be handled specially, as do quotes within
quotes:
       ``\,`this'            % \, separates the double and single quote.
        is what I just
        wrote, not  `that'\,''.

Dashes come in three sizes: an
       intra-word
dash, a medium dash for number ranges like
       1--2,
and a punctuation
       dash---like
this.

A sentence-ending space should be larger than the
space between words within a sentence.  You
sometimes have to type special commands in
conjunction with punctuation characters to get
this right, as in the following sentence.
       Gnats, gnus, etc.\ all  % `\ ' makes an inter-word space.
       begin with G\@.         % \@ marks end-of-sentence punctuation.
You should check the spaces after periods when
reading your output to make sure you haven't
forgotten any special cases.  Generating an
ellipsis
       \ldots\               % `\ ' is needed after `\ldots' because TeX
                             % ignores spaces after command names like \ldots
                             % made from \ + letters.
                             %
                             % Note how a `%' character causes TeX to ignore
                             % the end of the input line, so these blank lines
                             % do not start a new paragraph.
                             %
with the right spacing around the periods requires
a special command.

\LaTeX\ interprets some common characters as
commands, so you must type special commands to
generate them.  These characters include the
following:
       \$ \& \% \# \{ and \}.

In printing, text is usually emphasized with an
       \emph{italic}
type style.

\begin{em}
   A long segment of text can also be emphasized
   in this way.  Text within such a segment can be
   given \emph{additional} emphasis.
\end{em}

It is sometimes necessary to prevent \LaTeX\ from
breaking a line where it might otherwise do so.
This may be at a space, as between the ``Mr.''\ and
``Jones'' in
       ``Mr.~Jones'',        % ~ produces an unbreakable interword space.
or within a word---especially when the word is a
symbol like
       \mbox{\emph{itemnum}}
that makes little sense when hyphenated across
lines.

Footnotes\footnote{This is an example of a footnote.}
pose no problem.

\LaTeX\ is good at typesetting mathematical formulas
like
       \( x-3y + z = 7 \)
or
       \( a_{1} > x^{2n} + y^{2n} > x' \)
or
       \( \ip{A}{B} = \sum_{i} a_{i} b_{i} \).
The spaces you type in a formula are
ignored.  Remember that a letter like
       $x$                   % $ ... $  and  \( ... \)  are equivalent
is a formula when it denotes a mathematical
symbol, and it should be typed as one.

\section{Displayed Text}

Text is displayed by indenting it from the left
margin.  Quotations are commonly displayed.  There
are short quotations
\begin{quote}
   This is a short quotation.  It consists of a
   single paragraph of text.  See how it is formatted.
\end{quote}
and longer ones.
\begin{quotation}
   This is a longer quotation.  It consists of two
   paragraphs of text, neither of which are
   particularly interesting.

   This is the second paragraph of the quotation.  It
   is just as dull as the first paragraph.
\end{quotation}
Another frequently-displayed structure is a list.
The following is an example of an \emph{itemized}
list.
\begin{itemize}
   \item This is the first item of an itemized list.
         Each item in the list is marked with a ``tick''.
         You don't have to worry about what kind of tick
         mark is used.

   \item This is the second item of the list.  It
         contains another list nested inside it.  The inner
         list is an \emph{enumerated} list.
         \begin{enumerate}
            \item This is the first item of an enumerated
                  list that is nested within the itemized list.

            \item This is the second item of the inner list.
                  \LaTeX\ allows you to nest lists deeper than
                  you really should.
         \end{enumerate}
         This is the rest of the second item of the outer
         list.  It is no more interesting than any other
         part of the item.
   \item This is the third item of the list.
\end{itemize}
You can even display poetry.
\begin{verse}
   There is an environment
    for verse \\             % The \\ command separates lines
   Whose features some poets % within a stanza.
   will curse.

                             % One or more blank lines separate stanzas.

   For instead of making\\
   Them do \emph{all} line breaking, \\
   It allows them to put too many words on a line when they'd rather be
   forced to be terse.
\end{verse}

Mathematical formulas may also be displayed.  A
displayed formula
is
one-line long; multiline
formulas require special formatting instructions.
   \[  \ip{\Gamma}{\psi'} = x'' + y^{2} + z_{i}^{n}\]
Don't start a paragraph with a displayed equation,
nor make one a paragraph by itself.

\end{document}               % End of document.

\section{Bench Huge Section 1235}
% synthetic bench section 1235
% This is a sample LaTeX input file.  (Version of 12 August 2004.)
%
% A '%' character causes TeX to ignore all remaining text on the line,
% and is used for comments like this one.

\documentclass{article}      % Specifies the document class

                             % The preamble begins here.
\title{An Example Document}  % Declares the document's title.
\author{Leslie Lamport}      % Declares the author's name.
\date{January 21, 1994}      % Deleting this command produces today's date.

\newcommand{\ip}[2]{(#1, #2)}
                             % Defines \ip{arg1}{arg2} to mean
                             % (arg1, arg2).

%\newcommand{\ip}[2]{\langle #1 | #2\rangle}
                             % This is an alternative definition of
                             % \ip that is commented out.

\begin{document}             % End of preamble and beginning of text.

\maketitle                   % Produces the title.

This is an example input file.  Comparing it with
the output it generates can show you how to
produce a simple document of your own.

\section{Ordinary Text}      % Produces section heading.  Lower-level
                             % sections are begun with similar
                             % \subsection and \subsubsection commands.

The ends  of words and sentences are marked
  by   spaces. It  doesn't matter how many
spaces    you type; one is as good as 100.  The
end of   a line counts as a space.

One   or more   blank lines denote the  end
of  a paragraph.

Since any number of consecutive spaces are treated
like a single one, the formatting of the input
file makes no difference to
      \LaTeX,                % The \LaTeX command generates the LaTeX logo.
but it makes a difference to you.  When you use
\LaTeX, making your input file as easy to read
as possible will be a great help as you write
your document and when you change it.  This sample
file shows how you can add comments to your own input
file.

Because printing is different from typewriting,
there are a number of things that you have to do
differently when preparing an input file than if
you were just typing the document directly.
Quotation marks like
       ``this''
have to be handled specially, as do quotes within
quotes:
       ``\,`this'            % \, separates the double and single quote.
        is what I just
        wrote, not  `that'\,''.

Dashes come in three sizes: an
       intra-word
dash, a medium dash for number ranges like
       1--2,
and a punctuation
       dash---like
this.

A sentence-ending space should be larger than the
space between words within a sentence.  You
sometimes have to type special commands in
conjunction with punctuation characters to get
this right, as in the following sentence.
       Gnats, gnus, etc.\ all  % `\ ' makes an inter-word space.
       begin with G\@.         % \@ marks end-of-sentence punctuation.
You should check the spaces after periods when
reading your output to make sure you haven't
forgotten any special cases.  Generating an
ellipsis
       \ldots\               % `\ ' is needed after `\ldots' because TeX
                             % ignores spaces after command names like \ldots
                             % made from \ + letters.
                             %
                             % Note how a `%' character causes TeX to ignore
                             % the end of the input line, so these blank lines
                             % do not start a new paragraph.
                             %
with the right spacing around the periods requires
a special command.

\LaTeX\ interprets some common characters as
commands, so you must type special commands to
generate them.  These characters include the
following:
       \$ \& \% \# \{ and \}.

In printing, text is usually emphasized with an
       \emph{italic}
type style.

\begin{em}
   A long segment of text can also be emphasized
   in this way.  Text within such a segment can be
   given \emph{additional} emphasis.
\end{em}

It is sometimes necessary to prevent \LaTeX\ from
breaking a line where it might otherwise do so.
This may be at a space, as between the ``Mr.''\ and
``Jones'' in
       ``Mr.~Jones'',        % ~ produces an unbreakable interword space.
or within a word---especially when the word is a
symbol like
       \mbox{\emph{itemnum}}
that makes little sense when hyphenated across
lines.

Footnotes\footnote{This is an example of a footnote.}
pose no problem.

\LaTeX\ is good at typesetting mathematical formulas
like
       \( x-3y + z = 7 \)
or
       \( a_{1} > x^{2n} + y^{2n} > x' \)
or
       \( \ip{A}{B} = \sum_{i} a_{i} b_{i} \).
The spaces you type in a formula are
ignored.  Remember that a letter like
       $x$                   % $ ... $  and  \( ... \)  are equivalent
is a formula when it denotes a mathematical
symbol, and it should be typed as one.

\section{Displayed Text}

Text is displayed by indenting it from the left
margin.  Quotations are commonly displayed.  There
are short quotations
\begin{quote}
   This is a short quotation.  It consists of a
   single paragraph of text.  See how it is formatted.
\end{quote}
and longer ones.
\begin{quotation}
   This is a longer quotation.  It consists of two
   paragraphs of text, neither of which are
   particularly interesting.

   This is the second paragraph of the quotation.  It
   is just as dull as the first paragraph.
\end{quotation}
Another frequently-displayed structure is a list.
The following is an example of an \emph{itemized}
list.
\begin{itemize}
   \item This is the first item of an itemized list.
         Each item in the list is marked with a ``tick''.
         You don't have to worry about what kind of tick
         mark is used.

   \item This is the second item of the list.  It
         contains another list nested inside it.  The inner
         list is an \emph{enumerated} list.
         \begin{enumerate}
            \item This is the first item of an enumerated
                  list that is nested within the itemized list.

            \item This is the second item of the inner list.
                  \LaTeX\ allows you to nest lists deeper than
                  you really should.
         \end{enumerate}
         This is the rest of the second item of the outer
         list.  It is no more interesting than any other
         part of the item.
   \item This is the third item of the list.
\end{itemize}
You can even display poetry.
\begin{verse}
   There is an environment
    for verse \\             % The \\ command separates lines
   Whose features some poets % within a stanza.
   will curse.

                             % One or more blank lines separate stanzas.

   For instead of making\\
   Them do \emph{all} line breaking, \\
   It allows them to put too many words on a line when they'd rather be
   forced to be terse.
\end{verse}

Mathematical formulas may also be displayed.  A
displayed formula
is
one-line long; multiline
formulas require special formatting instructions.
   \[  \ip{\Gamma}{\psi'} = x'' + y^{2} + z_{i}^{n}\]
Don't start a paragraph with a displayed equation,
nor make one a paragraph by itself.

\end{document}               % End of document.

\section{Bench Huge Section 1236}
% synthetic bench section 1236
% This is a sample LaTeX input file.  (Version of 12 August 2004.)
%
% A '%' character causes TeX to ignore all remaining text on the line,
% and is used for comments like this one.

\documentclass{article}      % Specifies the document class

                             % The preamble begins here.
\title{An Example Document}  % Declares the document's title.
\author{Leslie Lamport}      % Declares the author's name.
\date{January 21, 1994}      % Deleting this command produces today's date.

\newcommand{\ip}[2]{(#1, #2)}
                             % Defines \ip{arg1}{arg2} to mean
                             % (arg1, arg2).

%\newcommand{\ip}[2]{\langle #1 | #2\rangle}
                             % This is an alternative definition of
                             % \ip that is commented out.

\begin{document}             % End of preamble and beginning of text.

\maketitle                   % Produces the title.

This is an example input file.  Comparing it with
the output it generates can show you how to
produce a simple document of your own.

\section{Ordinary Text}      % Produces section heading.  Lower-level
                             % sections are begun with similar
                             % \subsection and \subsubsection commands.

The ends  of words and sentences are marked
  by   spaces. It  doesn't matter how many
spaces    you type; one is as good as 100.  The
end of   a line counts as a space.

One   or more   blank lines denote the  end
of  a paragraph.

Since any number of consecutive spaces are treated
like a single one, the formatting of the input
file makes no difference to
      \LaTeX,                % The \LaTeX command generates the LaTeX logo.
but it makes a difference to you.  When you use
\LaTeX, making your input file as easy to read
as possible will be a great help as you write
your document and when you change it.  This sample
file shows how you can add comments to your own input
file.

Because printing is different from typewriting,
there are a number of things that you have to do
differently when preparing an input file than if
you were just typing the document directly.
Quotation marks like
       ``this''
have to be handled specially, as do quotes within
quotes:
       ``\,`this'            % \, separates the double and single quote.
        is what I just
        wrote, not  `that'\,''.

Dashes come in three sizes: an
       intra-word
dash, a medium dash for number ranges like
       1--2,
and a punctuation
       dash---like
this.

A sentence-ending space should be larger than the
space between words within a sentence.  You
sometimes have to type special commands in
conjunction with punctuation characters to get
this right, as in the following sentence.
       Gnats, gnus, etc.\ all  % `\ ' makes an inter-word space.
       begin with G\@.         % \@ marks end-of-sentence punctuation.
You should check the spaces after periods when
reading your output to make sure you haven't
forgotten any special cases.  Generating an
ellipsis
       \ldots\               % `\ ' is needed after `\ldots' because TeX
                             % ignores spaces after command names like \ldots
                             % made from \ + letters.
                             %
                             % Note how a `%' character causes TeX to ignore
                             % the end of the input line, so these blank lines
                             % do not start a new paragraph.
                             %
with the right spacing around the periods requires
a special command.

\LaTeX\ interprets some common characters as
commands, so you must type special commands to
generate them.  These characters include the
following:
       \$ \& \% \# \{ and \}.

In printing, text is usually emphasized with an
       \emph{italic}
type style.

\begin{em}
   A long segment of text can also be emphasized
   in this way.  Text within such a segment can be
   given \emph{additional} emphasis.
\end{em}

It is sometimes necessary to prevent \LaTeX\ from
breaking a line where it might otherwise do so.
This may be at a space, as between the ``Mr.''\ and
``Jones'' in
       ``Mr.~Jones'',        % ~ produces an unbreakable interword space.
or within a word---especially when the word is a
symbol like
       \mbox{\emph{itemnum}}
that makes little sense when hyphenated across
lines.

Footnotes\footnote{This is an example of a footnote.}
pose no problem.

\LaTeX\ is good at typesetting mathematical formulas
like
       \( x-3y + z = 7 \)
or
       \( a_{1} > x^{2n} + y^{2n} > x' \)
or
       \( \ip{A}{B} = \sum_{i} a_{i} b_{i} \).
The spaces you type in a formula are
ignored.  Remember that a letter like
       $x$                   % $ ... $  and  \( ... \)  are equivalent
is a formula when it denotes a mathematical
symbol, and it should be typed as one.

\section{Displayed Text}

Text is displayed by indenting it from the left
margin.  Quotations are commonly displayed.  There
are short quotations
\begin{quote}
   This is a short quotation.  It consists of a
   single paragraph of text.  See how it is formatted.
\end{quote}
and longer ones.
\begin{quotation}
   This is a longer quotation.  It consists of two
   paragraphs of text, neither of which are
   particularly interesting.

   This is the second paragraph of the quotation.  It
   is just as dull as the first paragraph.
\end{quotation}
Another frequently-displayed structure is a list.
The following is an example of an \emph{itemized}
list.
\begin{itemize}
   \item This is the first item of an itemized list.
         Each item in the list is marked with a ``tick''.
         You don't have to worry about what kind of tick
         mark is used.

   \item This is the second item of the list.  It
         contains another list nested inside it.  The inner
         list is an \emph{enumerated} list.
         \begin{enumerate}
            \item This is the first item of an enumerated
                  list that is nested within the itemized list.

            \item This is the second item of the inner list.
                  \LaTeX\ allows you to nest lists deeper than
                  you really should.
         \end{enumerate}
         This is the rest of the second item of the outer
         list.  It is no more interesting than any other
         part of the item.
   \item This is the third item of the list.
\end{itemize}
You can even display poetry.
\begin{verse}
   There is an environment
    for verse \\             % The \\ command separates lines
   Whose features some poets % within a stanza.
   will curse.

                             % One or more blank lines separate stanzas.

   For instead of making\\
   Them do \emph{all} line breaking, \\
   It allows them to put too many words on a line when they'd rather be
   forced to be terse.
\end{verse}

Mathematical formulas may also be displayed.  A
displayed formula
is
one-line long; multiline
formulas require special formatting instructions.
   \[  \ip{\Gamma}{\psi'} = x'' + y^{2} + z_{i}^{n}\]
Don't start a paragraph with a displayed equation,
nor make one a paragraph by itself.

\end{document}               % End of document.

\section{Bench Huge Section 1237}
% synthetic bench section 1237
% This is a sample LaTeX input file.  (Version of 12 August 2004.)
%
% A '%' character causes TeX to ignore all remaining text on the line,
% and is used for comments like this one.

\documentclass{article}      % Specifies the document class

                             % The preamble begins here.
\title{An Example Document}  % Declares the document's title.
\author{Leslie Lamport}      % Declares the author's name.
\date{January 21, 1994}      % Deleting this command produces today's date.

\newcommand{\ip}[2]{(#1, #2)}
                             % Defines \ip{arg1}{arg2} to mean
                             % (arg1, arg2).

%\newcommand{\ip}[2]{\langle #1 | #2\rangle}
                             % This is an alternative definition of
                             % \ip that is commented out.

\begin{document}             % End of preamble and beginning of text.

\maketitle                   % Produces the title.

This is an example input file.  Comparing it with
the output it generates can show you how to
produce a simple document of your own.

\section{Ordinary Text}      % Produces section heading.  Lower-level
                             % sections are begun with similar
                             % \subsection and \subsubsection commands.

The ends  of words and sentences are marked
  by   spaces. It  doesn't matter how many
spaces    you type; one is as good as 100.  The
end of   a line counts as a space.

One   or more   blank lines denote the  end
of  a paragraph.

Since any number of consecutive spaces are treated
like a single one, the formatting of the input
file makes no difference to
      \LaTeX,                % The \LaTeX command generates the LaTeX logo.
but it makes a difference to you.  When you use
\LaTeX, making your input file as easy to read
as possible will be a great help as you write
your document and when you change it.  This sample
file shows how you can add comments to your own input
file.

Because printing is different from typewriting,
there are a number of things that you have to do
differently when preparing an input file than if
you were just typing the document directly.
Quotation marks like
       ``this''
have to be handled specially, as do quotes within
quotes:
       ``\,`this'            % \, separates the double and single quote.
        is what I just
        wrote, not  `that'\,''.

Dashes come in three sizes: an
       intra-word
dash, a medium dash for number ranges like
       1--2,
and a punctuation
       dash---like
this.

A sentence-ending space should be larger than the
space between words within a sentence.  You
sometimes have to type special commands in
conjunction with punctuation characters to get
this right, as in the following sentence.
       Gnats, gnus, etc.\ all  % `\ ' makes an inter-word space.
       begin with G\@.         % \@ marks end-of-sentence punctuation.
You should check the spaces after periods when
reading your output to make sure you haven't
forgotten any special cases.  Generating an
ellipsis
       \ldots\               % `\ ' is needed after `\ldots' because TeX
                             % ignores spaces after command names like \ldots
                             % made from \ + letters.
                             %
                             % Note how a `%' character causes TeX to ignore
                             % the end of the input line, so these blank lines
                             % do not start a new paragraph.
                             %
with the right spacing around the periods requires
a special command.

\LaTeX\ interprets some common characters as
commands, so you must type special commands to
generate them.  These characters include the
following:
       \$ \& \% \# \{ and \}.

In printing, text is usually emphasized with an
       \emph{italic}
type style.

\begin{em}
   A long segment of text can also be emphasized
   in this way.  Text within such a segment can be
   given \emph{additional} emphasis.
\end{em}

It is sometimes necessary to prevent \LaTeX\ from
breaking a line where it might otherwise do so.
This may be at a space, as between the ``Mr.''\ and
``Jones'' in
       ``Mr.~Jones'',        % ~ produces an unbreakable interword space.
or within a word---especially when the word is a
symbol like
       \mbox{\emph{itemnum}}
that makes little sense when hyphenated across
lines.

Footnotes\footnote{This is an example of a footnote.}
pose no problem.

\LaTeX\ is good at typesetting mathematical formulas
like
       \( x-3y + z = 7 \)
or
       \( a_{1} > x^{2n} + y^{2n} > x' \)
or
       \( \ip{A}{B} = \sum_{i} a_{i} b_{i} \).
The spaces you type in a formula are
ignored.  Remember that a letter like
       $x$                   % $ ... $  and  \( ... \)  are equivalent
is a formula when it denotes a mathematical
symbol, and it should be typed as one.

\section{Displayed Text}

Text is displayed by indenting it from the left
margin.  Quotations are commonly displayed.  There
are short quotations
\begin{quote}
   This is a short quotation.  It consists of a
   single paragraph of text.  See how it is formatted.
\end{quote}
and longer ones.
\begin{quotation}
   This is a longer quotation.  It consists of two
   paragraphs of text, neither of which are
   particularly interesting.

   This is the second paragraph of the quotation.  It
   is just as dull as the first paragraph.
\end{quotation}
Another frequently-displayed structure is a list.
The following is an example of an \emph{itemized}
list.
\begin{itemize}
   \item This is the first item of an itemized list.
         Each item in the list is marked with a ``tick''.
         You don't have to worry about what kind of tick
         mark is used.

   \item This is the second item of the list.  It
         contains another list nested inside it.  The inner
         list is an \emph{enumerated} list.
         \begin{enumerate}
            \item This is the first item of an enumerated
                  list that is nested within the itemized list.

            \item This is the second item of the inner list.
                  \LaTeX\ allows you to nest lists deeper than
                  you really should.
         \end{enumerate}
         This is the rest of the second item of the outer
         list.  It is no more interesting than any other
         part of the item.
   \item This is the third item of the list.
\end{itemize}
You can even display poetry.
\begin{verse}
   There is an environment
    for verse \\             % The \\ command separates lines
   Whose features some poets % within a stanza.
   will curse.

                             % One or more blank lines separate stanzas.

   For instead of making\\
   Them do \emph{all} line breaking, \\
   It allows them to put too many words on a line when they'd rather be
   forced to be terse.
\end{verse}

Mathematical formulas may also be displayed.  A
displayed formula
is
one-line long; multiline
formulas require special formatting instructions.
   \[  \ip{\Gamma}{\psi'} = x'' + y^{2} + z_{i}^{n}\]
Don't start a paragraph with a displayed equation,
nor make one a paragraph by itself.

\end{document}               % End of document.

\section{Bench Huge Section 1238}
% synthetic bench section 1238
% This is a sample LaTeX input file.  (Version of 12 August 2004.)
%
% A '%' character causes TeX to ignore all remaining text on the line,
% and is used for comments like this one.

\documentclass{article}      % Specifies the document class

                             % The preamble begins here.
\title{An Example Document}  % Declares the document's title.
\author{Leslie Lamport}      % Declares the author's name.
\date{January 21, 1994}      % Deleting this command produces today's date.

\newcommand{\ip}[2]{(#1, #2)}
                             % Defines \ip{arg1}{arg2} to mean
                             % (arg1, arg2).

%\newcommand{\ip}[2]{\langle #1 | #2\rangle}
                             % This is an alternative definition of
                             % \ip that is commented out.

\begin{document}             % End of preamble and beginning of text.

\maketitle                   % Produces the title.

This is an example input file.  Comparing it with
the output it generates can show you how to
produce a simple document of your own.

\section{Ordinary Text}      % Produces section heading.  Lower-level
                             % sections are begun with similar
                             % \subsection and \subsubsection commands.

The ends  of words and sentences are marked
  by   spaces. It  doesn't matter how many
spaces    you type; one is as good as 100.  The
end of   a line counts as a space.

One   or more   blank lines denote the  end
of  a paragraph.

Since any number of consecutive spaces are treated
like a single one, the formatting of the input
file makes no difference to
      \LaTeX,                % The \LaTeX command generates the LaTeX logo.
but it makes a difference to you.  When you use
\LaTeX, making your input file as easy to read
as possible will be a great help as you write
your document and when you change it.  This sample
file shows how you can add comments to your own input
file.

Because printing is different from typewriting,
there are a number of things that you have to do
differently when preparing an input file than if
you were just typing the document directly.
Quotation marks like
       ``this''
have to be handled specially, as do quotes within
quotes:
       ``\,`this'            % \, separates the double and single quote.
        is what I just
        wrote, not  `that'\,''.

Dashes come in three sizes: an
       intra-word
dash, a medium dash for number ranges like
       1--2,
and a punctuation
       dash---like
this.

A sentence-ending space should be larger than the
space between words within a sentence.  You
sometimes have to type special commands in
conjunction with punctuation characters to get
this right, as in the following sentence.
       Gnats, gnus, etc.\ all  % `\ ' makes an inter-word space.
       begin with G\@.         % \@ marks end-of-sentence punctuation.
You should check the spaces after periods when
reading your output to make sure you haven't
forgotten any special cases.  Generating an
ellipsis
       \ldots\               % `\ ' is needed after `\ldots' because TeX
                             % ignores spaces after command names like \ldots
                             % made from \ + letters.
                             %
                             % Note how a `%' character causes TeX to ignore
                             % the end of the input line, so these blank lines
                             % do not start a new paragraph.
                             %
with the right spacing around the periods requires
a special command.

\LaTeX\ interprets some common characters as
commands, so you must type special commands to
generate them.  These characters include the
following:
       \$ \& \% \# \{ and \}.

In printing, text is usually emphasized with an
       \emph{italic}
type style.

\begin{em}
   A long segment of text can also be emphasized
   in this way.  Text within such a segment can be
   given \emph{additional} emphasis.
\end{em}

It is sometimes necessary to prevent \LaTeX\ from
breaking a line where it might otherwise do so.
This may be at a space, as between the ``Mr.''\ and
``Jones'' in
       ``Mr.~Jones'',        % ~ produces an unbreakable interword space.
or within a word---especially when the word is a
symbol like
       \mbox{\emph{itemnum}}
that makes little sense when hyphenated across
lines.

Footnotes\footnote{This is an example of a footnote.}
pose no problem.

\LaTeX\ is good at typesetting mathematical formulas
like
       \( x-3y + z = 7 \)
or
       \( a_{1} > x^{2n} + y^{2n} > x' \)
or
       \( \ip{A}{B} = \sum_{i} a_{i} b_{i} \).
The spaces you type in a formula are
ignored.  Remember that a letter like
       $x$                   % $ ... $  and  \( ... \)  are equivalent
is a formula when it denotes a mathematical
symbol, and it should be typed as one.

\section{Displayed Text}

Text is displayed by indenting it from the left
margin.  Quotations are commonly displayed.  There
are short quotations
\begin{quote}
   This is a short quotation.  It consists of a
   single paragraph of text.  See how it is formatted.
\end{quote}
and longer ones.
\begin{quotation}
   This is a longer quotation.  It consists of two
   paragraphs of text, neither of which are
   particularly interesting.

   This is the second paragraph of the quotation.  It
   is just as dull as the first paragraph.
\end{quotation}
Another frequently-displayed structure is a list.
The following is an example of an \emph{itemized}
list.
\begin{itemize}
   \item This is the first item of an itemized list.
         Each item in the list is marked with a ``tick''.
         You don't have to worry about what kind of tick
         mark is used.

   \item This is the second item of the list.  It
         contains another list nested inside it.  The inner
         list is an \emph{enumerated} list.
         \begin{enumerate}
            \item This is the first item of an enumerated
                  list that is nested within the itemized list.

            \item This is the second item of the inner list.
                  \LaTeX\ allows you to nest lists deeper than
                  you really should.
         \end{enumerate}
         This is the rest of the second item of the outer
         list.  It is no more interesting than any other
         part of the item.
   \item This is the third item of the list.
\end{itemize}
You can even display poetry.
\begin{verse}
   There is an environment
    for verse \\             % The \\ command separates lines
   Whose features some poets % within a stanza.
   will curse.

                             % One or more blank lines separate stanzas.

   For instead of making\\
   Them do \emph{all} line breaking, \\
   It allows them to put too many words on a line when they'd rather be
   forced to be terse.
\end{verse}

Mathematical formulas may also be displayed.  A
displayed formula
is
one-line long; multiline
formulas require special formatting instructions.
   \[  \ip{\Gamma}{\psi'} = x'' + y^{2} + z_{i}^{n}\]
Don't start a paragraph with a displayed equation,
nor make one a paragraph by itself.

\end{document}               % End of document.

\section{Bench Huge Section 1239}
% synthetic bench section 1239
% This is a sample LaTeX input file.  (Version of 12 August 2004.)
%
% A '%' character causes TeX to ignore all remaining text on the line,
% and is used for comments like this one.

\documentclass{article}      % Specifies the document class

                             % The preamble begins here.
\title{An Example Document}  % Declares the document's title.
\author{Leslie Lamport}      % Declares the author's name.
\date{January 21, 1994}      % Deleting this command produces today's date.

\newcommand{\ip}[2]{(#1, #2)}
                             % Defines \ip{arg1}{arg2} to mean
                             % (arg1, arg2).

%\newcommand{\ip}[2]{\langle #1 | #2\rangle}
                             % This is an alternative definition of
                             % \ip that is commented out.

\begin{document}             % End of preamble and beginning of text.

\maketitle                   % Produces the title.

This is an example input file.  Comparing it with
the output it generates can show you how to
produce a simple document of your own.

\section{Ordinary Text}      % Produces section heading.  Lower-level
                             % sections are begun with similar
                             % \subsection and \subsubsection commands.

The ends  of words and sentences are marked
  by   spaces. It  doesn't matter how many
spaces    you type; one is as good as 100.  The
end of   a line counts as a space.

One   or more   blank lines denote the  end
of  a paragraph.

Since any number of consecutive spaces are treated
like a single one, the formatting of the input
file makes no difference to
      \LaTeX,                % The \LaTeX command generates the LaTeX logo.
but it makes a difference to you.  When you use
\LaTeX, making your input file as easy to read
as possible will be a great help as you write
your document and when you change it.  This sample
file shows how you can add comments to your own input
file.

Because printing is different from typewriting,
there are a number of things that you have to do
differently when preparing an input file than if
you were just typing the document directly.
Quotation marks like
       ``this''
have to be handled specially, as do quotes within
quotes:
       ``\,`this'            % \, separates the double and single quote.
        is what I just
        wrote, not  `that'\,''.

Dashes come in three sizes: an
       intra-word
dash, a medium dash for number ranges like
       1--2,
and a punctuation
       dash---like
this.

A sentence-ending space should be larger than the
space between words within a sentence.  You
sometimes have to type special commands in
conjunction with punctuation characters to get
this right, as in the following sentence.
       Gnats, gnus, etc.\ all  % `\ ' makes an inter-word space.
       begin with G\@.         % \@ marks end-of-sentence punctuation.
You should check the spaces after periods when
reading your output to make sure you haven't
forgotten any special cases.  Generating an
ellipsis
       \ldots\               % `\ ' is needed after `\ldots' because TeX
                             % ignores spaces after command names like \ldots
                             % made from \ + letters.
                             %
                             % Note how a `%' character causes TeX to ignore
                             % the end of the input line, so these blank lines
                             % do not start a new paragraph.
                             %
with the right spacing around the periods requires
a special command.

\LaTeX\ interprets some common characters as
commands, so you must type special commands to
generate them.  These characters include the
following:
       \$ \& \% \# \{ and \}.

In printing, text is usually emphasized with an
       \emph{italic}
type style.

\begin{em}
   A long segment of text can also be emphasized
   in this way.  Text within such a segment can be
   given \emph{additional} emphasis.
\end{em}

It is sometimes necessary to prevent \LaTeX\ from
breaking a line where it might otherwise do so.
This may be at a space, as between the ``Mr.''\ and
``Jones'' in
       ``Mr.~Jones'',        % ~ produces an unbreakable interword space.
or within a word---especially when the word is a
symbol like
       \mbox{\emph{itemnum}}
that makes little sense when hyphenated across
lines.

Footnotes\footnote{This is an example of a footnote.}
pose no problem.

\LaTeX\ is good at typesetting mathematical formulas
like
       \( x-3y + z = 7 \)
or
       \( a_{1} > x^{2n} + y^{2n} > x' \)
or
       \( \ip{A}{B} = \sum_{i} a_{i} b_{i} \).
The spaces you type in a formula are
ignored.  Remember that a letter like
       $x$                   % $ ... $  and  \( ... \)  are equivalent
is a formula when it denotes a mathematical
symbol, and it should be typed as one.

\section{Displayed Text}

Text is displayed by indenting it from the left
margin.  Quotations are commonly displayed.  There
are short quotations
\begin{quote}
   This is a short quotation.  It consists of a
   single paragraph of text.  See how it is formatted.
\end{quote}
and longer ones.
\begin{quotation}
   This is a longer quotation.  It consists of two
   paragraphs of text, neither of which are
   particularly interesting.

   This is the second paragraph of the quotation.  It
   is just as dull as the first paragraph.
\end{quotation}
Another frequently-displayed structure is a list.
The following is an example of an \emph{itemized}
list.
\begin{itemize}
   \item This is the first item of an itemized list.
         Each item in the list is marked with a ``tick''.
         You don't have to worry about what kind of tick
         mark is used.

   \item This is the second item of the list.  It
         contains another list nested inside it.  The inner
         list is an \emph{enumerated} list.
         \begin{enumerate}
            \item This is the first item of an enumerated
                  list that is nested within the itemized list.

            \item This is the second item of the inner list.
                  \LaTeX\ allows you to nest lists deeper than
                  you really should.
         \end{enumerate}
         This is the rest of the second item of the outer
         list.  It is no more interesting than any other
         part of the item.
   \item This is the third item of the list.
\end{itemize}
You can even display poetry.
\begin{verse}
   There is an environment
    for verse \\             % The \\ command separates lines
   Whose features some poets % within a stanza.
   will curse.

                             % One or more blank lines separate stanzas.

   For instead of making\\
   Them do \emph{all} line breaking, \\
   It allows them to put too many words on a line when they'd rather be
   forced to be terse.
\end{verse}

Mathematical formulas may also be displayed.  A
displayed formula
is
one-line long; multiline
formulas require special formatting instructions.
   \[  \ip{\Gamma}{\psi'} = x'' + y^{2} + z_{i}^{n}\]
Don't start a paragraph with a displayed equation,
nor make one a paragraph by itself.

\end{document}               % End of document.

\section{Bench Huge Section 1240}
% synthetic bench section 1240
% This is a sample LaTeX input file.  (Version of 12 August 2004.)
%
% A '%' character causes TeX to ignore all remaining text on the line,
% and is used for comments like this one.

\documentclass{article}      % Specifies the document class

                             % The preamble begins here.
\title{An Example Document}  % Declares the document's title.
\author{Leslie Lamport}      % Declares the author's name.
\date{January 21, 1994}      % Deleting this command produces today's date.

\newcommand{\ip}[2]{(#1, #2)}
                             % Defines \ip{arg1}{arg2} to mean
                             % (arg1, arg2).

%\newcommand{\ip}[2]{\langle #1 | #2\rangle}
                             % This is an alternative definition of
                             % \ip that is commented out.

\begin{document}             % End of preamble and beginning of text.

\maketitle                   % Produces the title.

This is an example input file.  Comparing it with
the output it generates can show you how to
produce a simple document of your own.

\section{Ordinary Text}      % Produces section heading.  Lower-level
                             % sections are begun with similar
                             % \subsection and \subsubsection commands.

The ends  of words and sentences are marked
  by   spaces. It  doesn't matter how many
spaces    you type; one is as good as 100.  The
end of   a line counts as a space.

One   or more   blank lines denote the  end
of  a paragraph.

Since any number of consecutive spaces are treated
like a single one, the formatting of the input
file makes no difference to
      \LaTeX,                % The \LaTeX command generates the LaTeX logo.
but it makes a difference to you.  When you use
\LaTeX, making your input file as easy to read
as possible will be a great help as you write
your document and when you change it.  This sample
file shows how you can add comments to your own input
file.

Because printing is different from typewriting,
there are a number of things that you have to do
differently when preparing an input file than if
you were just typing the document directly.
Quotation marks like
       ``this''
have to be handled specially, as do quotes within
quotes:
       ``\,`this'            % \, separates the double and single quote.
        is what I just
        wrote, not  `that'\,''.

Dashes come in three sizes: an
       intra-word
dash, a medium dash for number ranges like
       1--2,
and a punctuation
       dash---like
this.

A sentence-ending space should be larger than the
space between words within a sentence.  You
sometimes have to type special commands in
conjunction with punctuation characters to get
this right, as in the following sentence.
       Gnats, gnus, etc.\ all  % `\ ' makes an inter-word space.
       begin with G\@.         % \@ marks end-of-sentence punctuation.
You should check the spaces after periods when
reading your output to make sure you haven't
forgotten any special cases.  Generating an
ellipsis
       \ldots\               % `\ ' is needed after `\ldots' because TeX
                             % ignores spaces after command names like \ldots
                             % made from \ + letters.
                             %
                             % Note how a `%' character causes TeX to ignore
                             % the end of the input line, so these blank lines
                             % do not start a new paragraph.
                             %
with the right spacing around the periods requires
a special command.

\LaTeX\ interprets some common characters as
commands, so you must type special commands to
generate them.  These characters include the
following:
       \$ \& \% \# \{ and \}.

In printing, text is usually emphasized with an
       \emph{italic}
type style.

\begin{em}
   A long segment of text can also be emphasized
   in this way.  Text within such a segment can be
   given \emph{additional} emphasis.
\end{em}

It is sometimes necessary to prevent \LaTeX\ from
breaking a line where it might otherwise do so.
This may be at a space, as between the ``Mr.''\ and
``Jones'' in
       ``Mr.~Jones'',        % ~ produces an unbreakable interword space.
or within a word---especially when the word is a
symbol like
       \mbox{\emph{itemnum}}
that makes little sense when hyphenated across
lines.

Footnotes\footnote{This is an example of a footnote.}
pose no problem.

\LaTeX\ is good at typesetting mathematical formulas
like
       \( x-3y + z = 7 \)
or
       \( a_{1} > x^{2n} + y^{2n} > x' \)
or
       \( \ip{A}{B} = \sum_{i} a_{i} b_{i} \).
The spaces you type in a formula are
ignored.  Remember that a letter like
       $x$                   % $ ... $  and  \( ... \)  are equivalent
is a formula when it denotes a mathematical
symbol, and it should be typed as one.

\section{Displayed Text}

Text is displayed by indenting it from the left
margin.  Quotations are commonly displayed.  There
are short quotations
\begin{quote}
   This is a short quotation.  It consists of a
   single paragraph of text.  See how it is formatted.
\end{quote}
and longer ones.
\begin{quotation}
   This is a longer quotation.  It consists of two
   paragraphs of text, neither of which are
   particularly interesting.

   This is the second paragraph of the quotation.  It
   is just as dull as the first paragraph.
\end{quotation}
Another frequently-displayed structure is a list.
The following is an example of an \emph{itemized}
list.
\begin{itemize}
   \item This is the first item of an itemized list.
         Each item in the list is marked with a ``tick''.
         You don't have to worry about what kind of tick
         mark is used.

   \item This is the second item of the list.  It
         contains another list nested inside it.  The inner
         list is an \emph{enumerated} list.
         \begin{enumerate}
            \item This is the first item of an enumerated
                  list that is nested within the itemized list.

            \item This is the second item of the inner list.
                  \LaTeX\ allows you to nest lists deeper than
                  you really should.
         \end{enumerate}
         This is the rest of the second item of the outer
         list.  It is no more interesting than any other
         part of the item.
   \item This is the third item of the list.
\end{itemize}
You can even display poetry.
\begin{verse}
   There is an environment
    for verse \\             % The \\ command separates lines
   Whose features some poets % within a stanza.
   will curse.

                             % One or more blank lines separate stanzas.

   For instead of making\\
   Them do \emph{all} line breaking, \\
   It allows them to put too many words on a line when they'd rather be
   forced to be terse.
\end{verse}

Mathematical formulas may also be displayed.  A
displayed formula
is
one-line long; multiline
formulas require special formatting instructions.
   \[  \ip{\Gamma}{\psi'} = x'' + y^{2} + z_{i}^{n}\]
Don't start a paragraph with a displayed equation,
nor make one a paragraph by itself.

\end{document}               % End of document.

\section{Bench Huge Section 1241}
% synthetic bench section 1241
% This is a sample LaTeX input file.  (Version of 12 August 2004.)
%
% A '%' character causes TeX to ignore all remaining text on the line,
% and is used for comments like this one.

\documentclass{article}      % Specifies the document class

                             % The preamble begins here.
\title{An Example Document}  % Declares the document's title.
\author{Leslie Lamport}      % Declares the author's name.
\date{January 21, 1994}      % Deleting this command produces today's date.

\newcommand{\ip}[2]{(#1, #2)}
                             % Defines \ip{arg1}{arg2} to mean
                             % (arg1, arg2).

%\newcommand{\ip}[2]{\langle #1 | #2\rangle}
                             % This is an alternative definition of
                             % \ip that is commented out.

\begin{document}             % End of preamble and beginning of text.

\maketitle                   % Produces the title.

This is an example input file.  Comparing it with
the output it generates can show you how to
produce a simple document of your own.

\section{Ordinary Text}      % Produces section heading.  Lower-level
                             % sections are begun with similar
                             % \subsection and \subsubsection commands.

The ends  of words and sentences are marked
  by   spaces. It  doesn't matter how many
spaces    you type; one is as good as 100.  The
end of   a line counts as a space.

One   or more   blank lines denote the  end
of  a paragraph.

Since any number of consecutive spaces are treated
like a single one, the formatting of the input
file makes no difference to
      \LaTeX,                % The \LaTeX command generates the LaTeX logo.
but it makes a difference to you.  When you use
\LaTeX, making your input file as easy to read
as possible will be a great help as you write
your document and when you change it.  This sample
file shows how you can add comments to your own input
file.

Because printing is different from typewriting,
there are a number of things that you have to do
differently when preparing an input file than if
you were just typing the document directly.
Quotation marks like
       ``this''
have to be handled specially, as do quotes within
quotes:
       ``\,`this'            % \, separates the double and single quote.
        is what I just
        wrote, not  `that'\,''.

Dashes come in three sizes: an
       intra-word
dash, a medium dash for number ranges like
       1--2,
and a punctuation
       dash---like
this.

A sentence-ending space should be larger than the
space between words within a sentence.  You
sometimes have to type special commands in
conjunction with punctuation characters to get
this right, as in the following sentence.
       Gnats, gnus, etc.\ all  % `\ ' makes an inter-word space.
       begin with G\@.         % \@ marks end-of-sentence punctuation.
You should check the spaces after periods when
reading your output to make sure you haven't
forgotten any special cases.  Generating an
ellipsis
       \ldots\               % `\ ' is needed after `\ldots' because TeX
                             % ignores spaces after command names like \ldots
                             % made from \ + letters.
                             %
                             % Note how a `%' character causes TeX to ignore
                             % the end of the input line, so these blank lines
                             % do not start a new paragraph.
                             %
with the right spacing around the periods requires
a special command.

\LaTeX\ interprets some common characters as
commands, so you must type special commands to
generate them.  These characters include the
following:
       \$ \& \% \# \{ and \}.

In printing, text is usually emphasized with an
       \emph{italic}
type style.

\begin{em}
   A long segment of text can also be emphasized
   in this way.  Text within such a segment can be
   given \emph{additional} emphasis.
\end{em}

It is sometimes necessary to prevent \LaTeX\ from
breaking a line where it might otherwise do so.
This may be at a space, as between the ``Mr.''\ and
``Jones'' in
       ``Mr.~Jones'',        % ~ produces an unbreakable interword space.
or within a word---especially when the word is a
symbol like
       \mbox{\emph{itemnum}}
that makes little sense when hyphenated across
lines.

Footnotes\footnote{This is an example of a footnote.}
pose no problem.

\LaTeX\ is good at typesetting mathematical formulas
like
       \( x-3y + z = 7 \)
or
       \( a_{1} > x^{2n} + y^{2n} > x' \)
or
       \( \ip{A}{B} = \sum_{i} a_{i} b_{i} \).
The spaces you type in a formula are
ignored.  Remember that a letter like
       $x$                   % $ ... $  and  \( ... \)  are equivalent
is a formula when it denotes a mathematical
symbol, and it should be typed as one.

\section{Displayed Text}

Text is displayed by indenting it from the left
margin.  Quotations are commonly displayed.  There
are short quotations
\begin{quote}
   This is a short quotation.  It consists of a
   single paragraph of text.  See how it is formatted.
\end{quote}
and longer ones.
\begin{quotation}
   This is a longer quotation.  It consists of two
   paragraphs of text, neither of which are
   particularly interesting.

   This is the second paragraph of the quotation.  It
   is just as dull as the first paragraph.
\end{quotation}
Another frequently-displayed structure is a list.
The following is an example of an \emph{itemized}
list.
\begin{itemize}
   \item This is the first item of an itemized list.
         Each item in the list is marked with a ``tick''.
         You don't have to worry about what kind of tick
         mark is used.

   \item This is the second item of the list.  It
         contains another list nested inside it.  The inner
         list is an \emph{enumerated} list.
         \begin{enumerate}
            \item This is the first item of an enumerated
                  list that is nested within the itemized list.

            \item This is the second item of the inner list.
                  \LaTeX\ allows you to nest lists deeper than
                  you really should.
         \end{enumerate}
         This is the rest of the second item of the outer
         list.  It is no more interesting than any other
         part of the item.
   \item This is the third item of the list.
\end{itemize}
You can even display poetry.
\begin{verse}
   There is an environment
    for verse \\             % The \\ command separates lines
   Whose features some poets % within a stanza.
   will curse.

                             % One or more blank lines separate stanzas.

   For instead of making\\
   Them do \emph{all} line breaking, \\
   It allows them to put too many words on a line when they'd rather be
   forced to be terse.
\end{verse}

Mathematical formulas may also be displayed.  A
displayed formula
is
one-line long; multiline
formulas require special formatting instructions.
   \[  \ip{\Gamma}{\psi'} = x'' + y^{2} + z_{i}^{n}\]
Don't start a paragraph with a displayed equation,
nor make one a paragraph by itself.

\end{document}               % End of document.

\section{Bench Huge Section 1242}
% synthetic bench section 1242
% This is a sample LaTeX input file.  (Version of 12 August 2004.)
%
% A '%' character causes TeX to ignore all remaining text on the line,
% and is used for comments like this one.

\documentclass{article}      % Specifies the document class

                             % The preamble begins here.
\title{An Example Document}  % Declares the document's title.
\author{Leslie Lamport}      % Declares the author's name.
\date{January 21, 1994}      % Deleting this command produces today's date.

\newcommand{\ip}[2]{(#1, #2)}
                             % Defines \ip{arg1}{arg2} to mean
                             % (arg1, arg2).

%\newcommand{\ip}[2]{\langle #1 | #2\rangle}
                             % This is an alternative definition of
                             % \ip that is commented out.

\begin{document}             % End of preamble and beginning of text.

\maketitle                   % Produces the title.

This is an example input file.  Comparing it with
the output it generates can show you how to
produce a simple document of your own.

\section{Ordinary Text}      % Produces section heading.  Lower-level
                             % sections are begun with similar
                             % \subsection and \subsubsection commands.

The ends  of words and sentences are marked
  by   spaces. It  doesn't matter how many
spaces    you type; one is as good as 100.  The
end of   a line counts as a space.

One   or more   blank lines denote the  end
of  a paragraph.

Since any number of consecutive spaces are treated
like a single one, the formatting of the input
file makes no difference to
      \LaTeX,                % The \LaTeX command generates the LaTeX logo.
but it makes a difference to you.  When you use
\LaTeX, making your input file as easy to read
as possible will be a great help as you write
your document and when you change it.  This sample
file shows how you can add comments to your own input
file.

Because printing is different from typewriting,
there are a number of things that you have to do
differently when preparing an input file than if
you were just typing the document directly.
Quotation marks like
       ``this''
have to be handled specially, as do quotes within
quotes:
       ``\,`this'            % \, separates the double and single quote.
        is what I just
        wrote, not  `that'\,''.

Dashes come in three sizes: an
       intra-word
dash, a medium dash for number ranges like
       1--2,
and a punctuation
       dash---like
this.

A sentence-ending space should be larger than the
space between words within a sentence.  You
sometimes have to type special commands in
conjunction with punctuation characters to get
this right, as in the following sentence.
       Gnats, gnus, etc.\ all  % `\ ' makes an inter-word space.
       begin with G\@.         % \@ marks end-of-sentence punctuation.
You should check the spaces after periods when
reading your output to make sure you haven't
forgotten any special cases.  Generating an
ellipsis
       \ldots\               % `\ ' is needed after `\ldots' because TeX
                             % ignores spaces after command names like \ldots
                             % made from \ + letters.
                             %
                             % Note how a `%' character causes TeX to ignore
                             % the end of the input line, so these blank lines
                             % do not start a new paragraph.
                             %
with the right spacing around the periods requires
a special command.

\LaTeX\ interprets some common characters as
commands, so you must type special commands to
generate them.  These characters include the
following:
       \$ \& \% \# \{ and \}.

In printing, text is usually emphasized with an
       \emph{italic}
type style.

\begin{em}
   A long segment of text can also be emphasized
   in this way.  Text within such a segment can be
   given \emph{additional} emphasis.
\end{em}

It is sometimes necessary to prevent \LaTeX\ from
breaking a line where it might otherwise do so.
This may be at a space, as between the ``Mr.''\ and
``Jones'' in
       ``Mr.~Jones'',        % ~ produces an unbreakable interword space.
or within a word---especially when the word is a
symbol like
       \mbox{\emph{itemnum}}
that makes little sense when hyphenated across
lines.

Footnotes\footnote{This is an example of a footnote.}
pose no problem.

\LaTeX\ is good at typesetting mathematical formulas
like
       \( x-3y + z = 7 \)
or
       \( a_{1} > x^{2n} + y^{2n} > x' \)
or
       \( \ip{A}{B} = \sum_{i} a_{i} b_{i} \).
The spaces you type in a formula are
ignored.  Remember that a letter like
       $x$                   % $ ... $  and  \( ... \)  are equivalent
is a formula when it denotes a mathematical
symbol, and it should be typed as one.

\section{Displayed Text}

Text is displayed by indenting it from the left
margin.  Quotations are commonly displayed.  There
are short quotations
\begin{quote}
   This is a short quotation.  It consists of a
   single paragraph of text.  See how it is formatted.
\end{quote}
and longer ones.
\begin{quotation}
   This is a longer quotation.  It consists of two
   paragraphs of text, neither of which are
   particularly interesting.

   This is the second paragraph of the quotation.  It
   is just as dull as the first paragraph.
\end{quotation}
Another frequently-displayed structure is a list.
The following is an example of an \emph{itemized}
list.
\begin{itemize}
   \item This is the first item of an itemized list.
         Each item in the list is marked with a ``tick''.
         You don't have to worry about what kind of tick
         mark is used.

   \item This is the second item of the list.  It
         contains another list nested inside it.  The inner
         list is an \emph{enumerated} list.
         \begin{enumerate}
            \item This is the first item of an enumerated
                  list that is nested within the itemized list.

            \item This is the second item of the inner list.
                  \LaTeX\ allows you to nest lists deeper than
                  you really should.
         \end{enumerate}
         This is the rest of the second item of the outer
         list.  It is no more interesting than any other
         part of the item.
   \item This is the third item of the list.
\end{itemize}
You can even display poetry.
\begin{verse}
   There is an environment
    for verse \\             % The \\ command separates lines
   Whose features some poets % within a stanza.
   will curse.

                             % One or more blank lines separate stanzas.

   For instead of making\\
   Them do \emph{all} line breaking, \\
   It allows them to put too many words on a line when they'd rather be
   forced to be terse.
\end{verse}

Mathematical formulas may also be displayed.  A
displayed formula
is
one-line long; multiline
formulas require special formatting instructions.
   \[  \ip{\Gamma}{\psi'} = x'' + y^{2} + z_{i}^{n}\]
Don't start a paragraph with a displayed equation,
nor make one a paragraph by itself.

\end{document}               % End of document.

\section{Bench Huge Section 1243}
% synthetic bench section 1243
% This is a sample LaTeX input file.  (Version of 12 August 2004.)
%
% A '%' character causes TeX to ignore all remaining text on the line,
% and is used for comments like this one.

\documentclass{article}      % Specifies the document class

                             % The preamble begins here.
\title{An Example Document}  % Declares the document's title.
\author{Leslie Lamport}      % Declares the author's name.
\date{January 21, 1994}      % Deleting this command produces today's date.

\newcommand{\ip}[2]{(#1, #2)}
                             % Defines \ip{arg1}{arg2} to mean
                             % (arg1, arg2).

%\newcommand{\ip}[2]{\langle #1 | #2\rangle}
                             % This is an alternative definition of
                             % \ip that is commented out.

\begin{document}             % End of preamble and beginning of text.

\maketitle                   % Produces the title.

This is an example input file.  Comparing it with
the output it generates can show you how to
produce a simple document of your own.

\section{Ordinary Text}      % Produces section heading.  Lower-level
                             % sections are begun with similar
                             % \subsection and \subsubsection commands.

The ends  of words and sentences are marked
  by   spaces. It  doesn't matter how many
spaces    you type; one is as good as 100.  The
end of   a line counts as a space.

One   or more   blank lines denote the  end
of  a paragraph.

Since any number of consecutive spaces are treated
like a single one, the formatting of the input
file makes no difference to
      \LaTeX,                % The \LaTeX command generates the LaTeX logo.
but it makes a difference to you.  When you use
\LaTeX, making your input file as easy to read
as possible will be a great help as you write
your document and when you change it.  This sample
file shows how you can add comments to your own input
file.

Because printing is different from typewriting,
there are a number of things that you have to do
differently when preparing an input file than if
you were just typing the document directly.
Quotation marks like
       ``this''
have to be handled specially, as do quotes within
quotes:
       ``\,`this'            % \, separates the double and single quote.
        is what I just
        wrote, not  `that'\,''.

Dashes come in three sizes: an
       intra-word
dash, a medium dash for number ranges like
       1--2,
and a punctuation
       dash---like
this.

A sentence-ending space should be larger than the
space between words within a sentence.  You
sometimes have to type special commands in
conjunction with punctuation characters to get
this right, as in the following sentence.
       Gnats, gnus, etc.\ all  % `\ ' makes an inter-word space.
       begin with G\@.         % \@ marks end-of-sentence punctuation.
You should check the spaces after periods when
reading your output to make sure you haven't
forgotten any special cases.  Generating an
ellipsis
       \ldots\               % `\ ' is needed after `\ldots' because TeX
                             % ignores spaces after command names like \ldots
                             % made from \ + letters.
                             %
                             % Note how a `%' character causes TeX to ignore
                             % the end of the input line, so these blank lines
                             % do not start a new paragraph.
                             %
with the right spacing around the periods requires
a special command.

\LaTeX\ interprets some common characters as
commands, so you must type special commands to
generate them.  These characters include the
following:
       \$ \& \% \# \{ and \}.

In printing, text is usually emphasized with an
       \emph{italic}
type style.

\begin{em}
   A long segment of text can also be emphasized
   in this way.  Text within such a segment can be
   given \emph{additional} emphasis.
\end{em}

It is sometimes necessary to prevent \LaTeX\ from
breaking a line where it might otherwise do so.
This may be at a space, as between the ``Mr.''\ and
``Jones'' in
       ``Mr.~Jones'',        % ~ produces an unbreakable interword space.
or within a word---especially when the word is a
symbol like
       \mbox{\emph{itemnum}}
that makes little sense when hyphenated across
lines.

Footnotes\footnote{This is an example of a footnote.}
pose no problem.

\LaTeX\ is good at typesetting mathematical formulas
like
       \( x-3y + z = 7 \)
or
       \( a_{1} > x^{2n} + y^{2n} > x' \)
or
       \( \ip{A}{B} = \sum_{i} a_{i} b_{i} \).
The spaces you type in a formula are
ignored.  Remember that a letter like
       $x$                   % $ ... $  and  \( ... \)  are equivalent
is a formula when it denotes a mathematical
symbol, and it should be typed as one.

\section{Displayed Text}

Text is displayed by indenting it from the left
margin.  Quotations are commonly displayed.  There
are short quotations
\begin{quote}
   This is a short quotation.  It consists of a
   single paragraph of text.  See how it is formatted.
\end{quote}
and longer ones.
\begin{quotation}
   This is a longer quotation.  It consists of two
   paragraphs of text, neither of which are
   particularly interesting.

   This is the second paragraph of the quotation.  It
   is just as dull as the first paragraph.
\end{quotation}
Another frequently-displayed structure is a list.
The following is an example of an \emph{itemized}
list.
\begin{itemize}
   \item This is the first item of an itemized list.
         Each item in the list is marked with a ``tick''.
         You don't have to worry about what kind of tick
         mark is used.

   \item This is the second item of the list.  It
         contains another list nested inside it.  The inner
         list is an \emph{enumerated} list.
         \begin{enumerate}
            \item This is the first item of an enumerated
                  list that is nested within the itemized list.

            \item This is the second item of the inner list.
                  \LaTeX\ allows you to nest lists deeper than
                  you really should.
         \end{enumerate}
         This is the rest of the second item of the outer
         list.  It is no more interesting than any other
         part of the item.
   \item This is the third item of the list.
\end{itemize}
You can even display poetry.
\begin{verse}
   There is an environment
    for verse \\             % The \\ command separates lines
   Whose features some poets % within a stanza.
   will curse.

                             % One or more blank lines separate stanzas.

   For instead of making\\
   Them do \emph{all} line breaking, \\
   It allows them to put too many words on a line when they'd rather be
   forced to be terse.
\end{verse}

Mathematical formulas may also be displayed.  A
displayed formula
is
one-line long; multiline
formulas require special formatting instructions.
   \[  \ip{\Gamma}{\psi'} = x'' + y^{2} + z_{i}^{n}\]
Don't start a paragraph with a displayed equation,
nor make one a paragraph by itself.

\end{document}               % End of document.

\section{Bench Huge Section 1244}
% synthetic bench section 1244
% This is a sample LaTeX input file.  (Version of 12 August 2004.)
%
% A '%' character causes TeX to ignore all remaining text on the line,
% and is used for comments like this one.

\documentclass{article}      % Specifies the document class

                             % The preamble begins here.
\title{An Example Document}  % Declares the document's title.
\author{Leslie Lamport}      % Declares the author's name.
\date{January 21, 1994}      % Deleting this command produces today's date.

\newcommand{\ip}[2]{(#1, #2)}
                             % Defines \ip{arg1}{arg2} to mean
                             % (arg1, arg2).

%\newcommand{\ip}[2]{\langle #1 | #2\rangle}
                             % This is an alternative definition of
                             % \ip that is commented out.

\begin{document}             % End of preamble and beginning of text.

\maketitle                   % Produces the title.

This is an example input file.  Comparing it with
the output it generates can show you how to
produce a simple document of your own.

\section{Ordinary Text}      % Produces section heading.  Lower-level
                             % sections are begun with similar
                             % \subsection and \subsubsection commands.

The ends  of words and sentences are marked
  by   spaces. It  doesn't matter how many
spaces    you type; one is as good as 100.  The
end of   a line counts as a space.

One   or more   blank lines denote the  end
of  a paragraph.

Since any number of consecutive spaces are treated
like a single one, the formatting of the input
file makes no difference to
      \LaTeX,                % The \LaTeX command generates the LaTeX logo.
but it makes a difference to you.  When you use
\LaTeX, making your input file as easy to read
as possible will be a great help as you write
your document and when you change it.  This sample
file shows how you can add comments to your own input
file.

Because printing is different from typewriting,
there are a number of things that you have to do
differently when preparing an input file than if
you were just typing the document directly.
Quotation marks like
       ``this''
have to be handled specially, as do quotes within
quotes:
       ``\,`this'            % \, separates the double and single quote.
        is what I just
        wrote, not  `that'\,''.

Dashes come in three sizes: an
       intra-word
dash, a medium dash for number ranges like
       1--2,
and a punctuation
       dash---like
this.

A sentence-ending space should be larger than the
space between words within a sentence.  You
sometimes have to type special commands in
conjunction with punctuation characters to get
this right, as in the following sentence.
       Gnats, gnus, etc.\ all  % `\ ' makes an inter-word space.
       begin with G\@.         % \@ marks end-of-sentence punctuation.
You should check the spaces after periods when
reading your output to make sure you haven't
forgotten any special cases.  Generating an
ellipsis
       \ldots\               % `\ ' is needed after `\ldots' because TeX
                             % ignores spaces after command names like \ldots
                             % made from \ + letters.
                             %
                             % Note how a `%' character causes TeX to ignore
                             % the end of the input line, so these blank lines
                             % do not start a new paragraph.
                             %
with the right spacing around the periods requires
a special command.

\LaTeX\ interprets some common characters as
commands, so you must type special commands to
generate them.  These characters include the
following:
       \$ \& \% \# \{ and \}.

In printing, text is usually emphasized with an
       \emph{italic}
type style.

\begin{em}
   A long segment of text can also be emphasized
   in this way.  Text within such a segment can be
   given \emph{additional} emphasis.
\end{em}

It is sometimes necessary to prevent \LaTeX\ from
breaking a line where it might otherwise do so.
This may be at a space, as between the ``Mr.''\ and
``Jones'' in
       ``Mr.~Jones'',        % ~ produces an unbreakable interword space.
or within a word---especially when the word is a
symbol like
       \mbox{\emph{itemnum}}
that makes little sense when hyphenated across
lines.

Footnotes\footnote{This is an example of a footnote.}
pose no problem.

\LaTeX\ is good at typesetting mathematical formulas
like
       \( x-3y + z = 7 \)
or
       \( a_{1} > x^{2n} + y^{2n} > x' \)
or
       \( \ip{A}{B} = \sum_{i} a_{i} b_{i} \).
The spaces you type in a formula are
ignored.  Remember that a letter like
       $x$                   % $ ... $  and  \( ... \)  are equivalent
is a formula when it denotes a mathematical
symbol, and it should be typed as one.

\section{Displayed Text}

Text is displayed by indenting it from the left
margin.  Quotations are commonly displayed.  There
are short quotations
\begin{quote}
   This is a short quotation.  It consists of a
   single paragraph of text.  See how it is formatted.
\end{quote}
and longer ones.
\begin{quotation}
   This is a longer quotation.  It consists of two
   paragraphs of text, neither of which are
   particularly interesting.

   This is the second paragraph of the quotation.  It
   is just as dull as the first paragraph.
\end{quotation}
Another frequently-displayed structure is a list.
The following is an example of an \emph{itemized}
list.
\begin{itemize}
   \item This is the first item of an itemized list.
         Each item in the list is marked with a ``tick''.
         You don't have to worry about what kind of tick
         mark is used.

   \item This is the second item of the list.  It
         contains another list nested inside it.  The inner
         list is an \emph{enumerated} list.
         \begin{enumerate}
            \item This is the first item of an enumerated
                  list that is nested within the itemized list.

            \item This is the second item of the inner list.
                  \LaTeX\ allows you to nest lists deeper than
                  you really should.
         \end{enumerate}
         This is the rest of the second item of the outer
         list.  It is no more interesting than any other
         part of the item.
   \item This is the third item of the list.
\end{itemize}
You can even display poetry.
\begin{verse}
   There is an environment
    for verse \\             % The \\ command separates lines
   Whose features some poets % within a stanza.
   will curse.

                             % One or more blank lines separate stanzas.

   For instead of making\\
   Them do \emph{all} line breaking, \\
   It allows them to put too many words on a line when they'd rather be
   forced to be terse.
\end{verse}

Mathematical formulas may also be displayed.  A
displayed formula
is
one-line long; multiline
formulas require special formatting instructions.
   \[  \ip{\Gamma}{\psi'} = x'' + y^{2} + z_{i}^{n}\]
Don't start a paragraph with a displayed equation,
nor make one a paragraph by itself.

\end{document}               % End of document.

\section{Bench Huge Section 1245}
% synthetic bench section 1245
% This is a sample LaTeX input file.  (Version of 12 August 2004.)
%
% A '%' character causes TeX to ignore all remaining text on the line,
% and is used for comments like this one.

\documentclass{article}      % Specifies the document class

                             % The preamble begins here.
\title{An Example Document}  % Declares the document's title.
\author{Leslie Lamport}      % Declares the author's name.
\date{January 21, 1994}      % Deleting this command produces today's date.

\newcommand{\ip}[2]{(#1, #2)}
                             % Defines \ip{arg1}{arg2} to mean
                             % (arg1, arg2).

%\newcommand{\ip}[2]{\langle #1 | #2\rangle}
                             % This is an alternative definition of
                             % \ip that is commented out.

\begin{document}             % End of preamble and beginning of text.

\maketitle                   % Produces the title.

This is an example input file.  Comparing it with
the output it generates can show you how to
produce a simple document of your own.

\section{Ordinary Text}      % Produces section heading.  Lower-level
                             % sections are begun with similar
                             % \subsection and \subsubsection commands.

The ends  of words and sentences are marked
  by   spaces. It  doesn't matter how many
spaces    you type; one is as good as 100.  The
end of   a line counts as a space.

One   or more   blank lines denote the  end
of  a paragraph.

Since any number of consecutive spaces are treated
like a single one, the formatting of the input
file makes no difference to
      \LaTeX,                % The \LaTeX command generates the LaTeX logo.
but it makes a difference to you.  When you use
\LaTeX, making your input file as easy to read
as possible will be a great help as you write
your document and when you change it.  This sample
file shows how you can add comments to your own input
file.

Because printing is different from typewriting,
there are a number of things that you have to do
differently when preparing an input file than if
you were just typing the document directly.
Quotation marks like
       ``this''
have to be handled specially, as do quotes within
quotes:
       ``\,`this'            % \, separates the double and single quote.
        is what I just
        wrote, not  `that'\,''.

Dashes come in three sizes: an
       intra-word
dash, a medium dash for number ranges like
       1--2,
and a punctuation
       dash---like
this.

A sentence-ending space should be larger than the
space between words within a sentence.  You
sometimes have to type special commands in
conjunction with punctuation characters to get
this right, as in the following sentence.
       Gnats, gnus, etc.\ all  % `\ ' makes an inter-word space.
       begin with G\@.         % \@ marks end-of-sentence punctuation.
You should check the spaces after periods when
reading your output to make sure you haven't
forgotten any special cases.  Generating an
ellipsis
       \ldots\               % `\ ' is needed after `\ldots' because TeX
                             % ignores spaces after command names like \ldots
                             % made from \ + letters.
                             %
                             % Note how a `%' character causes TeX to ignore
                             % the end of the input line, so these blank lines
                             % do not start a new paragraph.
                             %
with the right spacing around the periods requires
a special command.

\LaTeX\ interprets some common characters as
commands, so you must type special commands to
generate them.  These characters include the
following:
       \$ \& \% \# \{ and \}.

In printing, text is usually emphasized with an
       \emph{italic}
type style.

\begin{em}
   A long segment of text can also be emphasized
   in this way.  Text within such a segment can be
   given \emph{additional} emphasis.
\end{em}

It is sometimes necessary to prevent \LaTeX\ from
breaking a line where it might otherwise do so.
This may be at a space, as between the ``Mr.''\ and
``Jones'' in
       ``Mr.~Jones'',        % ~ produces an unbreakable interword space.
or within a word---especially when the word is a
symbol like
       \mbox{\emph{itemnum}}
that makes little sense when hyphenated across
lines.

Footnotes\footnote{This is an example of a footnote.}
pose no problem.

\LaTeX\ is good at typesetting mathematical formulas
like
       \( x-3y + z = 7 \)
or
       \( a_{1} > x^{2n} + y^{2n} > x' \)
or
       \( \ip{A}{B} = \sum_{i} a_{i} b_{i} \).
The spaces you type in a formula are
ignored.  Remember that a letter like
       $x$                   % $ ... $  and  \( ... \)  are equivalent
is a formula when it denotes a mathematical
symbol, and it should be typed as one.

\section{Displayed Text}

Text is displayed by indenting it from the left
margin.  Quotations are commonly displayed.  There
are short quotations
\begin{quote}
   This is a short quotation.  It consists of a
   single paragraph of text.  See how it is formatted.
\end{quote}
and longer ones.
\begin{quotation}
   This is a longer quotation.  It consists of two
   paragraphs of text, neither of which are
   particularly interesting.

   This is the second paragraph of the quotation.  It
   is just as dull as the first paragraph.
\end{quotation}
Another frequently-displayed structure is a list.
The following is an example of an \emph{itemized}
list.
\begin{itemize}
   \item This is the first item of an itemized list.
         Each item in the list is marked with a ``tick''.
         You don't have to worry about what kind of tick
         mark is used.

   \item This is the second item of the list.  It
         contains another list nested inside it.  The inner
         list is an \emph{enumerated} list.
         \begin{enumerate}
            \item This is the first item of an enumerated
                  list that is nested within the itemized list.

            \item This is the second item of the inner list.
                  \LaTeX\ allows you to nest lists deeper than
                  you really should.
         \end{enumerate}
         This is the rest of the second item of the outer
         list.  It is no more interesting than any other
         part of the item.
   \item This is the third item of the list.
\end{itemize}
You can even display poetry.
\begin{verse}
   There is an environment
    for verse \\             % The \\ command separates lines
   Whose features some poets % within a stanza.
   will curse.

                             % One or more blank lines separate stanzas.

   For instead of making\\
   Them do \emph{all} line breaking, \\
   It allows them to put too many words on a line when they'd rather be
   forced to be terse.
\end{verse}

Mathematical formulas may also be displayed.  A
displayed formula
is
one-line long; multiline
formulas require special formatting instructions.
   \[  \ip{\Gamma}{\psi'} = x'' + y^{2} + z_{i}^{n}\]
Don't start a paragraph with a displayed equation,
nor make one a paragraph by itself.

\end{document}               % End of document.

\section{Bench Huge Section 1246}
% synthetic bench section 1246
% This is a sample LaTeX input file.  (Version of 12 August 2004.)
%
% A '%' character causes TeX to ignore all remaining text on the line,
% and is used for comments like this one.

\documentclass{article}      % Specifies the document class

                             % The preamble begins here.
\title{An Example Document}  % Declares the document's title.
\author{Leslie Lamport}      % Declares the author's name.
\date{January 21, 1994}      % Deleting this command produces today's date.

\newcommand{\ip}[2]{(#1, #2)}
                             % Defines \ip{arg1}{arg2} to mean
                             % (arg1, arg2).

%\newcommand{\ip}[2]{\langle #1 | #2\rangle}
                             % This is an alternative definition of
                             % \ip that is commented out.

\begin{document}             % End of preamble and beginning of text.

\maketitle                   % Produces the title.

This is an example input file.  Comparing it with
the output it generates can show you how to
produce a simple document of your own.

\section{Ordinary Text}      % Produces section heading.  Lower-level
                             % sections are begun with similar
                             % \subsection and \subsubsection commands.

The ends  of words and sentences are marked
  by   spaces. It  doesn't matter how many
spaces    you type; one is as good as 100.  The
end of   a line counts as a space.

One   or more   blank lines denote the  end
of  a paragraph.

Since any number of consecutive spaces are treated
like a single one, the formatting of the input
file makes no difference to
      \LaTeX,                % The \LaTeX command generates the LaTeX logo.
but it makes a difference to you.  When you use
\LaTeX, making your input file as easy to read
as possible will be a great help as you write
your document and when you change it.  This sample
file shows how you can add comments to your own input
file.

Because printing is different from typewriting,
there are a number of things that you have to do
differently when preparing an input file than if
you were just typing the document directly.
Quotation marks like
       ``this''
have to be handled specially, as do quotes within
quotes:
       ``\,`this'            % \, separates the double and single quote.
        is what I just
        wrote, not  `that'\,''.

Dashes come in three sizes: an
       intra-word
dash, a medium dash for number ranges like
       1--2,
and a punctuation
       dash---like
this.

A sentence-ending space should be larger than the
space between words within a sentence.  You
sometimes have to type special commands in
conjunction with punctuation characters to get
this right, as in the following sentence.
       Gnats, gnus, etc.\ all  % `\ ' makes an inter-word space.
       begin with G\@.         % \@ marks end-of-sentence punctuation.
You should check the spaces after periods when
reading your output to make sure you haven't
forgotten any special cases.  Generating an
ellipsis
       \ldots\               % `\ ' is needed after `\ldots' because TeX
                             % ignores spaces after command names like \ldots
                             % made from \ + letters.
                             %
                             % Note how a `%' character causes TeX to ignore
                             % the end of the input line, so these blank lines
                             % do not start a new paragraph.
                             %
with the right spacing around the periods requires
a special command.

\LaTeX\ interprets some common characters as
commands, so you must type special commands to
generate them.  These characters include the
following:
       \$ \& \% \# \{ and \}.

In printing, text is usually emphasized with an
       \emph{italic}
type style.

\begin{em}
   A long segment of text can also be emphasized
   in this way.  Text within such a segment can be
   given \emph{additional} emphasis.
\end{em}

It is sometimes necessary to prevent \LaTeX\ from
breaking a line where it might otherwise do so.
This may be at a space, as between the ``Mr.''\ and
``Jones'' in
       ``Mr.~Jones'',        % ~ produces an unbreakable interword space.
or within a word---especially when the word is a
symbol like
       \mbox{\emph{itemnum}}
that makes little sense when hyphenated across
lines.

Footnotes\footnote{This is an example of a footnote.}
pose no problem.

\LaTeX\ is good at typesetting mathematical formulas
like
       \( x-3y + z = 7 \)
or
       \( a_{1} > x^{2n} + y^{2n} > x' \)
or
       \( \ip{A}{B} = \sum_{i} a_{i} b_{i} \).
The spaces you type in a formula are
ignored.  Remember that a letter like
       $x$                   % $ ... $  and  \( ... \)  are equivalent
is a formula when it denotes a mathematical
symbol, and it should be typed as one.

\section{Displayed Text}

Text is displayed by indenting it from the left
margin.  Quotations are commonly displayed.  There
are short quotations
\begin{quote}
   This is a short quotation.  It consists of a
   single paragraph of text.  See how it is formatted.
\end{quote}
and longer ones.
\begin{quotation}
   This is a longer quotation.  It consists of two
   paragraphs of text, neither of which are
   particularly interesting.

   This is the second paragraph of the quotation.  It
   is just as dull as the first paragraph.
\end{quotation}
Another frequently-displayed structure is a list.
The following is an example of an \emph{itemized}
list.
\begin{itemize}
   \item This is the first item of an itemized list.
         Each item in the list is marked with a ``tick''.
         You don't have to worry about what kind of tick
         mark is used.

   \item This is the second item of the list.  It
         contains another list nested inside it.  The inner
         list is an \emph{enumerated} list.
         \begin{enumerate}
            \item This is the first item of an enumerated
                  list that is nested within the itemized list.

            \item This is the second item of the inner list.
                  \LaTeX\ allows you to nest lists deeper than
                  you really should.
         \end{enumerate}
         This is the rest of the second item of the outer
         list.  It is no more interesting than any other
         part of the item.
   \item This is the third item of the list.
\end{itemize}
You can even display poetry.
\begin{verse}
   There is an environment
    for verse \\             % The \\ command separates lines
   Whose features some poets % within a stanza.
   will curse.

                             % One or more blank lines separate stanzas.

   For instead of making\\
   Them do \emph{all} line breaking, \\
   It allows them to put too many words on a line when they'd rather be
   forced to be terse.
\end{verse}

Mathematical formulas may also be displayed.  A
displayed formula
is
one-line long; multiline
formulas require special formatting instructions.
   \[  \ip{\Gamma}{\psi'} = x'' + y^{2} + z_{i}^{n}\]
Don't start a paragraph with a displayed equation,
nor make one a paragraph by itself.

\end{document}               % End of document.

\section{Bench Huge Section 1247}
% synthetic bench section 1247
% This is a sample LaTeX input file.  (Version of 12 August 2004.)
%
% A '%' character causes TeX to ignore all remaining text on the line,
% and is used for comments like this one.

\documentclass{article}      % Specifies the document class

                             % The preamble begins here.
\title{An Example Document}  % Declares the document's title.
\author{Leslie Lamport}      % Declares the author's name.
\date{January 21, 1994}      % Deleting this command produces today's date.

\newcommand{\ip}[2]{(#1, #2)}
                             % Defines \ip{arg1}{arg2} to mean
                             % (arg1, arg2).

%\newcommand{\ip}[2]{\langle #1 | #2\rangle}
                             % This is an alternative definition of
                             % \ip that is commented out.

\begin{document}             % End of preamble and beginning of text.

\maketitle                   % Produces the title.

This is an example input file.  Comparing it with
the output it generates can show you how to
produce a simple document of your own.

\section{Ordinary Text}      % Produces section heading.  Lower-level
                             % sections are begun with similar
                             % \subsection and \subsubsection commands.

The ends  of words and sentences are marked
  by   spaces. It  doesn't matter how many
spaces    you type; one is as good as 100.  The
end of   a line counts as a space.

One   or more   blank lines denote the  end
of  a paragraph.

Since any number of consecutive spaces are treated
like a single one, the formatting of the input
file makes no difference to
      \LaTeX,                % The \LaTeX command generates the LaTeX logo.
but it makes a difference to you.  When you use
\LaTeX, making your input file as easy to read
as possible will be a great help as you write
your document and when you change it.  This sample
file shows how you can add comments to your own input
file.

Because printing is different from typewriting,
there are a number of things that you have to do
differently when preparing an input file than if
you were just typing the document directly.
Quotation marks like
       ``this''
have to be handled specially, as do quotes within
quotes:
       ``\,`this'            % \, separates the double and single quote.
        is what I just
        wrote, not  `that'\,''.

Dashes come in three sizes: an
       intra-word
dash, a medium dash for number ranges like
       1--2,
and a punctuation
       dash---like
this.

A sentence-ending space should be larger than the
space between words within a sentence.  You
sometimes have to type special commands in
conjunction with punctuation characters to get
this right, as in the following sentence.
       Gnats, gnus, etc.\ all  % `\ ' makes an inter-word space.
       begin with G\@.         % \@ marks end-of-sentence punctuation.
You should check the spaces after periods when
reading your output to make sure you haven't
forgotten any special cases.  Generating an
ellipsis
       \ldots\               % `\ ' is needed after `\ldots' because TeX
                             % ignores spaces after command names like \ldots
                             % made from \ + letters.
                             %
                             % Note how a `%' character causes TeX to ignore
                             % the end of the input line, so these blank lines
                             % do not start a new paragraph.
                             %
with the right spacing around the periods requires
a special command.

\LaTeX\ interprets some common characters as
commands, so you must type special commands to
generate them.  These characters include the
following:
       \$ \& \% \# \{ and \}.

In printing, text is usually emphasized with an
       \emph{italic}
type style.

\begin{em}
   A long segment of text can also be emphasized
   in this way.  Text within such a segment can be
   given \emph{additional} emphasis.
\end{em}

It is sometimes necessary to prevent \LaTeX\ from
breaking a line where it might otherwise do so.
This may be at a space, as between the ``Mr.''\ and
``Jones'' in
       ``Mr.~Jones'',        % ~ produces an unbreakable interword space.
or within a word---especially when the word is a
symbol like
       \mbox{\emph{itemnum}}
that makes little sense when hyphenated across
lines.

Footnotes\footnote{This is an example of a footnote.}
pose no problem.

\LaTeX\ is good at typesetting mathematical formulas
like
       \( x-3y + z = 7 \)
or
       \( a_{1} > x^{2n} + y^{2n} > x' \)
or
       \( \ip{A}{B} = \sum_{i} a_{i} b_{i} \).
The spaces you type in a formula are
ignored.  Remember that a letter like
       $x$                   % $ ... $  and  \( ... \)  are equivalent
is a formula when it denotes a mathematical
symbol, and it should be typed as one.

\section{Displayed Text}

Text is displayed by indenting it from the left
margin.  Quotations are commonly displayed.  There
are short quotations
\begin{quote}
   This is a short quotation.  It consists of a
   single paragraph of text.  See how it is formatted.
\end{quote}
and longer ones.
\begin{quotation}
   This is a longer quotation.  It consists of two
   paragraphs of text, neither of which are
   particularly interesting.

   This is the second paragraph of the quotation.  It
   is just as dull as the first paragraph.
\end{quotation}
Another frequently-displayed structure is a list.
The following is an example of an \emph{itemized}
list.
\begin{itemize}
   \item This is the first item of an itemized list.
         Each item in the list is marked with a ``tick''.
         You don't have to worry about what kind of tick
         mark is used.

   \item This is the second item of the list.  It
         contains another list nested inside it.  The inner
         list is an \emph{enumerated} list.
         \begin{enumerate}
            \item This is the first item of an enumerated
                  list that is nested within the itemized list.

            \item This is the second item of the inner list.
                  \LaTeX\ allows you to nest lists deeper than
                  you really should.
         \end{enumerate}
         This is the rest of the second item of the outer
         list.  It is no more interesting than any other
         part of the item.
   \item This is the third item of the list.
\end{itemize}
You can even display poetry.
\begin{verse}
   There is an environment
    for verse \\             % The \\ command separates lines
   Whose features some poets % within a stanza.
   will curse.

                             % One or more blank lines separate stanzas.

   For instead of making\\
   Them do \emph{all} line breaking, \\
   It allows them to put too many words on a line when they'd rather be
   forced to be terse.
\end{verse}

Mathematical formulas may also be displayed.  A
displayed formula
is
one-line long; multiline
formulas require special formatting instructions.
   \[  \ip{\Gamma}{\psi'} = x'' + y^{2} + z_{i}^{n}\]
Don't start a paragraph with a displayed equation,
nor make one a paragraph by itself.

\end{document}               % End of document.

\section{Bench Huge Section 1248}
% synthetic bench section 1248
% This is a sample LaTeX input file.  (Version of 12 August 2004.)
%
% A '%' character causes TeX to ignore all remaining text on the line,
% and is used for comments like this one.

\documentclass{article}      % Specifies the document class

                             % The preamble begins here.
\title{An Example Document}  % Declares the document's title.
\author{Leslie Lamport}      % Declares the author's name.
\date{January 21, 1994}      % Deleting this command produces today's date.

\newcommand{\ip}[2]{(#1, #2)}
                             % Defines \ip{arg1}{arg2} to mean
                             % (arg1, arg2).

%\newcommand{\ip}[2]{\langle #1 | #2\rangle}
                             % This is an alternative definition of
                             % \ip that is commented out.

\begin{document}             % End of preamble and beginning of text.

\maketitle                   % Produces the title.

This is an example input file.  Comparing it with
the output it generates can show you how to
produce a simple document of your own.

\section{Ordinary Text}      % Produces section heading.  Lower-level
                             % sections are begun with similar
                             % \subsection and \subsubsection commands.

The ends  of words and sentences are marked
  by   spaces. It  doesn't matter how many
spaces    you type; one is as good as 100.  The
end of   a line counts as a space.

One   or more   blank lines denote the  end
of  a paragraph.

Since any number of consecutive spaces are treated
like a single one, the formatting of the input
file makes no difference to
      \LaTeX,                % The \LaTeX command generates the LaTeX logo.
but it makes a difference to you.  When you use
\LaTeX, making your input file as easy to read
as possible will be a great help as you write
your document and when you change it.  This sample
file shows how you can add comments to your own input
file.

Because printing is different from typewriting,
there are a number of things that you have to do
differently when preparing an input file than if
you were just typing the document directly.
Quotation marks like
       ``this''
have to be handled specially, as do quotes within
quotes:
       ``\,`this'            % \, separates the double and single quote.
        is what I just
        wrote, not  `that'\,''.

Dashes come in three sizes: an
       intra-word
dash, a medium dash for number ranges like
       1--2,
and a punctuation
       dash---like
this.

A sentence-ending space should be larger than the
space between words within a sentence.  You
sometimes have to type special commands in
conjunction with punctuation characters to get
this right, as in the following sentence.
       Gnats, gnus, etc.\ all  % `\ ' makes an inter-word space.
       begin with G\@.         % \@ marks end-of-sentence punctuation.
You should check the spaces after periods when
reading your output to make sure you haven't
forgotten any special cases.  Generating an
ellipsis
       \ldots\               % `\ ' is needed after `\ldots' because TeX
                             % ignores spaces after command names like \ldots
                             % made from \ + letters.
                             %
                             % Note how a `%' character causes TeX to ignore
                             % the end of the input line, so these blank lines
                             % do not start a new paragraph.
                             %
with the right spacing around the periods requires
a special command.

\LaTeX\ interprets some common characters as
commands, so you must type special commands to
generate them.  These characters include the
following:
       \$ \& \% \# \{ and \}.

In printing, text is usually emphasized with an
       \emph{italic}
type style.

\begin{em}
   A long segment of text can also be emphasized
   in this way.  Text within such a segment can be
   given \emph{additional} emphasis.
\end{em}

It is sometimes necessary to prevent \LaTeX\ from
breaking a line where it might otherwise do so.
This may be at a space, as between the ``Mr.''\ and
``Jones'' in
       ``Mr.~Jones'',        % ~ produces an unbreakable interword space.
or within a word---especially when the word is a
symbol like
       \mbox{\emph{itemnum}}
that makes little sense when hyphenated across
lines.

Footnotes\footnote{This is an example of a footnote.}
pose no problem.

\LaTeX\ is good at typesetting mathematical formulas
like
       \( x-3y + z = 7 \)
or
       \( a_{1} > x^{2n} + y^{2n} > x' \)
or
       \( \ip{A}{B} = \sum_{i} a_{i} b_{i} \).
The spaces you type in a formula are
ignored.  Remember that a letter like
       $x$                   % $ ... $  and  \( ... \)  are equivalent
is a formula when it denotes a mathematical
symbol, and it should be typed as one.

\section{Displayed Text}

Text is displayed by indenting it from the left
margin.  Quotations are commonly displayed.  There
are short quotations
\begin{quote}
   This is a short quotation.  It consists of a
   single paragraph of text.  See how it is formatted.
\end{quote}
and longer ones.
\begin{quotation}
   This is a longer quotation.  It consists of two
   paragraphs of text, neither of which are
   particularly interesting.

   This is the second paragraph of the quotation.  It
   is just as dull as the first paragraph.
\end{quotation}
Another frequently-displayed structure is a list.
The following is an example of an \emph{itemized}
list.
\begin{itemize}
   \item This is the first item of an itemized list.
         Each item in the list is marked with a ``tick''.
         You don't have to worry about what kind of tick
         mark is used.

   \item This is the second item of the list.  It
         contains another list nested inside it.  The inner
         list is an \emph{enumerated} list.
         \begin{enumerate}
            \item This is the first item of an enumerated
                  list that is nested within the itemized list.

            \item This is the second item of the inner list.
                  \LaTeX\ allows you to nest lists deeper than
                  you really should.
         \end{enumerate}
         This is the rest of the second item of the outer
         list.  It is no more interesting than any other
         part of the item.
   \item This is the third item of the list.
\end{itemize}
You can even display poetry.
\begin{verse}
   There is an environment
    for verse \\             % The \\ command separates lines
   Whose features some poets % within a stanza.
   will curse.

                             % One or more blank lines separate stanzas.

   For instead of making\\
   Them do \emph{all} line breaking, \\
   It allows them to put too many words on a line when they'd rather be
   forced to be terse.
\end{verse}

Mathematical formulas may also be displayed.  A
displayed formula
is
one-line long; multiline
formulas require special formatting instructions.
   \[  \ip{\Gamma}{\psi'} = x'' + y^{2} + z_{i}^{n}\]
Don't start a paragraph with a displayed equation,
nor make one a paragraph by itself.

\end{document}               % End of document.

\section{Bench Huge Section 1249}
% synthetic bench section 1249
% This is a sample LaTeX input file.  (Version of 12 August 2004.)
%
% A '%' character causes TeX to ignore all remaining text on the line,
% and is used for comments like this one.

\documentclass{article}      % Specifies the document class

                             % The preamble begins here.
\title{An Example Document}  % Declares the document's title.
\author{Leslie Lamport}      % Declares the author's name.
\date{January 21, 1994}      % Deleting this command produces today's date.

\newcommand{\ip}[2]{(#1, #2)}
                             % Defines \ip{arg1}{arg2} to mean
                             % (arg1, arg2).

%\newcommand{\ip}[2]{\langle #1 | #2\rangle}
                             % This is an alternative definition of
                             % \ip that is commented out.

\begin{document}             % End of preamble and beginning of text.

\maketitle                   % Produces the title.

This is an example input file.  Comparing it with
the output it generates can show you how to
produce a simple document of your own.

\section{Ordinary Text}      % Produces section heading.  Lower-level
                             % sections are begun with similar
                             % \subsection and \subsubsection commands.

The ends  of words and sentences are marked
  by   spaces. It  doesn't matter how many
spaces    you type; one is as good as 100.  The
end of   a line counts as a space.

One   or more   blank lines denote the  end
of  a paragraph.

Since any number of consecutive spaces are treated
like a single one, the formatting of the input
file makes no difference to
      \LaTeX,                % The \LaTeX command generates the LaTeX logo.
but it makes a difference to you.  When you use
\LaTeX, making your input file as easy to read
as possible will be a great help as you write
your document and when you change it.  This sample
file shows how you can add comments to your own input
file.

Because printing is different from typewriting,
there are a number of things that you have to do
differently when preparing an input file than if
you were just typing the document directly.
Quotation marks like
       ``this''
have to be handled specially, as do quotes within
quotes:
       ``\,`this'            % \, separates the double and single quote.
        is what I just
        wrote, not  `that'\,''.

Dashes come in three sizes: an
       intra-word
dash, a medium dash for number ranges like
       1--2,
and a punctuation
       dash---like
this.

A sentence-ending space should be larger than the
space between words within a sentence.  You
sometimes have to type special commands in
conjunction with punctuation characters to get
this right, as in the following sentence.
       Gnats, gnus, etc.\ all  % `\ ' makes an inter-word space.
       begin with G\@.         % \@ marks end-of-sentence punctuation.
You should check the spaces after periods when
reading your output to make sure you haven't
forgotten any special cases.  Generating an
ellipsis
       \ldots\               % `\ ' is needed after `\ldots' because TeX
                             % ignores spaces after command names like \ldots
                             % made from \ + letters.
                             %
                             % Note how a `%' character causes TeX to ignore
                             % the end of the input line, so these blank lines
                             % do not start a new paragraph.
                             %
with the right spacing around the periods requires
a special command.

\LaTeX\ interprets some common characters as
commands, so you must type special commands to
generate them.  These characters include the
following:
       \$ \& \% \# \{ and \}.

In printing, text is usually emphasized with an
       \emph{italic}
type style.

\begin{em}
   A long segment of text can also be emphasized
   in this way.  Text within such a segment can be
   given \emph{additional} emphasis.
\end{em}

It is sometimes necessary to prevent \LaTeX\ from
breaking a line where it might otherwise do so.
This may be at a space, as between the ``Mr.''\ and
``Jones'' in
       ``Mr.~Jones'',        % ~ produces an unbreakable interword space.
or within a word---especially when the word is a
symbol like
       \mbox{\emph{itemnum}}
that makes little sense when hyphenated across
lines.

Footnotes\footnote{This is an example of a footnote.}
pose no problem.

\LaTeX\ is good at typesetting mathematical formulas
like
       \( x-3y + z = 7 \)
or
       \( a_{1} > x^{2n} + y^{2n} > x' \)
or
       \( \ip{A}{B} = \sum_{i} a_{i} b_{i} \).
The spaces you type in a formula are
ignored.  Remember that a letter like
       $x$                   % $ ... $  and  \( ... \)  are equivalent
is a formula when it denotes a mathematical
symbol, and it should be typed as one.

\section{Displayed Text}

Text is displayed by indenting it from the left
margin.  Quotations are commonly displayed.  There
are short quotations
\begin{quote}
   This is a short quotation.  It consists of a
   single paragraph of text.  See how it is formatted.
\end{quote}
and longer ones.
\begin{quotation}
   This is a longer quotation.  It consists of two
   paragraphs of text, neither of which are
   particularly interesting.

   This is the second paragraph of the quotation.  It
   is just as dull as the first paragraph.
\end{quotation}
Another frequently-displayed structure is a list.
The following is an example of an \emph{itemized}
list.
\begin{itemize}
   \item This is the first item of an itemized list.
         Each item in the list is marked with a ``tick''.
         You don't have to worry about what kind of tick
         mark is used.

   \item This is the second item of the list.  It
         contains another list nested inside it.  The inner
         list is an \emph{enumerated} list.
         \begin{enumerate}
            \item This is the first item of an enumerated
                  list that is nested within the itemized list.

            \item This is the second item of the inner list.
                  \LaTeX\ allows you to nest lists deeper than
                  you really should.
         \end{enumerate}
         This is the rest of the second item of the outer
         list.  It is no more interesting than any other
         part of the item.
   \item This is the third item of the list.
\end{itemize}
You can even display poetry.
\begin{verse}
   There is an environment
    for verse \\             % The \\ command separates lines
   Whose features some poets % within a stanza.
   will curse.

                             % One or more blank lines separate stanzas.

   For instead of making\\
   Them do \emph{all} line breaking, \\
   It allows them to put too many words on a line when they'd rather be
   forced to be terse.
\end{verse}

Mathematical formulas may also be displayed.  A
displayed formula
is
one-line long; multiline
formulas require special formatting instructions.
   \[  \ip{\Gamma}{\psi'} = x'' + y^{2} + z_{i}^{n}\]
Don't start a paragraph with a displayed equation,
nor make one a paragraph by itself.

\end{document}               % End of document.

\section{Bench Huge Section 1250}
% synthetic bench section 1250
% This is a sample LaTeX input file.  (Version of 12 August 2004.)
%
% A '%' character causes TeX to ignore all remaining text on the line,
% and is used for comments like this one.

\documentclass{article}      % Specifies the document class

                             % The preamble begins here.
\title{An Example Document}  % Declares the document's title.
\author{Leslie Lamport}      % Declares the author's name.
\date{January 21, 1994}      % Deleting this command produces today's date.

\newcommand{\ip}[2]{(#1, #2)}
                             % Defines \ip{arg1}{arg2} to mean
                             % (arg1, arg2).

%\newcommand{\ip}[2]{\langle #1 | #2\rangle}
                             % This is an alternative definition of
                             % \ip that is commented out.

\begin{document}             % End of preamble and beginning of text.

\maketitle                   % Produces the title.

This is an example input file.  Comparing it with
the output it generates can show you how to
produce a simple document of your own.

\section{Ordinary Text}      % Produces section heading.  Lower-level
                             % sections are begun with similar
                             % \subsection and \subsubsection commands.

The ends  of words and sentences are marked
  by   spaces. It  doesn't matter how many
spaces    you type; one is as good as 100.  The
end of   a line counts as a space.

One   or more   blank lines denote the  end
of  a paragraph.

Since any number of consecutive spaces are treated
like a single one, the formatting of the input
file makes no difference to
      \LaTeX,                % The \LaTeX command generates the LaTeX logo.
but it makes a difference to you.  When you use
\LaTeX, making your input file as easy to read
as possible will be a great help as you write
your document and when you change it.  This sample
file shows how you can add comments to your own input
file.

Because printing is different from typewriting,
there are a number of things that you have to do
differently when preparing an input file than if
you were just typing the document directly.
Quotation marks like
       ``this''
have to be handled specially, as do quotes within
quotes:
       ``\,`this'            % \, separates the double and single quote.
        is what I just
        wrote, not  `that'\,''.

Dashes come in three sizes: an
       intra-word
dash, a medium dash for number ranges like
       1--2,
and a punctuation
       dash---like
this.

A sentence-ending space should be larger than the
space between words within a sentence.  You
sometimes have to type special commands in
conjunction with punctuation characters to get
this right, as in the following sentence.
       Gnats, gnus, etc.\ all  % `\ ' makes an inter-word space.
       begin with G\@.         % \@ marks end-of-sentence punctuation.
You should check the spaces after periods when
reading your output to make sure you haven't
forgotten any special cases.  Generating an
ellipsis
       \ldots\               % `\ ' is needed after `\ldots' because TeX
                             % ignores spaces after command names like \ldots
                             % made from \ + letters.
                             %
                             % Note how a `%' character causes TeX to ignore
                             % the end of the input line, so these blank lines
                             % do not start a new paragraph.
                             %
with the right spacing around the periods requires
a special command.

\LaTeX\ interprets some common characters as
commands, so you must type special commands to
generate them.  These characters include the
following:
       \$ \& \% \# \{ and \}.

In printing, text is usually emphasized with an
       \emph{italic}
type style.

\begin{em}
   A long segment of text can also be emphasized
   in this way.  Text within such a segment can be
   given \emph{additional} emphasis.
\end{em}

It is sometimes necessary to prevent \LaTeX\ from
breaking a line where it might otherwise do so.
This may be at a space, as between the ``Mr.''\ and
``Jones'' in
       ``Mr.~Jones'',        % ~ produces an unbreakable interword space.
or within a word---especially when the word is a
symbol like
       \mbox{\emph{itemnum}}
that makes little sense when hyphenated across
lines.

Footnotes\footnote{This is an example of a footnote.}
pose no problem.

\LaTeX\ is good at typesetting mathematical formulas
like
       \( x-3y + z = 7 \)
or
       \( a_{1} > x^{2n} + y^{2n} > x' \)
or
       \( \ip{A}{B} = \sum_{i} a_{i} b_{i} \).
The spaces you type in a formula are
ignored.  Remember that a letter like
       $x$                   % $ ... $  and  \( ... \)  are equivalent
is a formula when it denotes a mathematical
symbol, and it should be typed as one.

\section{Displayed Text}

Text is displayed by indenting it from the left
margin.  Quotations are commonly displayed.  There
are short quotations
\begin{quote}
   This is a short quotation.  It consists of a
   single paragraph of text.  See how it is formatted.
\end{quote}
and longer ones.
\begin{quotation}
   This is a longer quotation.  It consists of two
   paragraphs of text, neither of which are
   particularly interesting.

   This is the second paragraph of the quotation.  It
   is just as dull as the first paragraph.
\end{quotation}
Another frequently-displayed structure is a list.
The following is an example of an \emph{itemized}
list.
\begin{itemize}
   \item This is the first item of an itemized list.
         Each item in the list is marked with a ``tick''.
         You don't have to worry about what kind of tick
         mark is used.

   \item This is the second item of the list.  It
         contains another list nested inside it.  The inner
         list is an \emph{enumerated} list.
         \begin{enumerate}
            \item This is the first item of an enumerated
                  list that is nested within the itemized list.

            \item This is the second item of the inner list.
                  \LaTeX\ allows you to nest lists deeper than
                  you really should.
         \end{enumerate}
         This is the rest of the second item of the outer
         list.  It is no more interesting than any other
         part of the item.
   \item This is the third item of the list.
\end{itemize}
You can even display poetry.
\begin{verse}
   There is an environment
    for verse \\             % The \\ command separates lines
   Whose features some poets % within a stanza.
   will curse.

                             % One or more blank lines separate stanzas.

   For instead of making\\
   Them do \emph{all} line breaking, \\
   It allows them to put too many words on a line when they'd rather be
   forced to be terse.
\end{verse}

Mathematical formulas may also be displayed.  A
displayed formula
is
one-line long; multiline
formulas require special formatting instructions.
   \[  \ip{\Gamma}{\psi'} = x'' + y^{2} + z_{i}^{n}\]
Don't start a paragraph with a displayed equation,
nor make one a paragraph by itself.

\end{document}               % End of document.

\section{Bench Huge Section 1251}
% synthetic bench section 1251
% This is a sample LaTeX input file.  (Version of 12 August 2004.)
%
% A '%' character causes TeX to ignore all remaining text on the line,
% and is used for comments like this one.

\documentclass{article}      % Specifies the document class

                             % The preamble begins here.
\title{An Example Document}  % Declares the document's title.
\author{Leslie Lamport}      % Declares the author's name.
\date{January 21, 1994}      % Deleting this command produces today's date.

\newcommand{\ip}[2]{(#1, #2)}
                             % Defines \ip{arg1}{arg2} to mean
                             % (arg1, arg2).

%\newcommand{\ip}[2]{\langle #1 | #2\rangle}
                             % This is an alternative definition of
                             % \ip that is commented out.

\begin{document}             % End of preamble and beginning of text.

\maketitle                   % Produces the title.

This is an example input file.  Comparing it with
the output it generates can show you how to
produce a simple document of your own.

\section{Ordinary Text}      % Produces section heading.  Lower-level
                             % sections are begun with similar
                             % \subsection and \subsubsection commands.

The ends  of words and sentences are marked
  by   spaces. It  doesn't matter how many
spaces    you type; one is as good as 100.  The
end of   a line counts as a space.

One   or more   blank lines denote the  end
of  a paragraph.

Since any number of consecutive spaces are treated
like a single one, the formatting of the input
file makes no difference to
      \LaTeX,                % The \LaTeX command generates the LaTeX logo.
but it makes a difference to you.  When you use
\LaTeX, making your input file as easy to read
as possible will be a great help as you write
your document and when you change it.  This sample
file shows how you can add comments to your own input
file.

Because printing is different from typewriting,
there are a number of things that you have to do
differently when preparing an input file than if
you were just typing the document directly.
Quotation marks like
       ``this''
have to be handled specially, as do quotes within
quotes:
       ``\,`this'            % \, separates the double and single quote.
        is what I just
        wrote, not  `that'\,''.

Dashes come in three sizes: an
       intra-word
dash, a medium dash for number ranges like
       1--2,
and a punctuation
       dash---like
this.

A sentence-ending space should be larger than the
space between words within a sentence.  You
sometimes have to type special commands in
conjunction with punctuation characters to get
this right, as in the following sentence.
       Gnats, gnus, etc.\ all  % `\ ' makes an inter-word space.
       begin with G\@.         % \@ marks end-of-sentence punctuation.
You should check the spaces after periods when
reading your output to make sure you haven't
forgotten any special cases.  Generating an
ellipsis
       \ldots\               % `\ ' is needed after `\ldots' because TeX
                             % ignores spaces after command names like \ldots
                             % made from \ + letters.
                             %
                             % Note how a `%' character causes TeX to ignore
                             % the end of the input line, so these blank lines
                             % do not start a new paragraph.
                             %
with the right spacing around the periods requires
a special command.

\LaTeX\ interprets some common characters as
commands, so you must type special commands to
generate them.  These characters include the
following:
       \$ \& \% \# \{ and \}.

In printing, text is usually emphasized with an
       \emph{italic}
type style.

\begin{em}
   A long segment of text can also be emphasized
   in this way.  Text within such a segment can be
   given \emph{additional} emphasis.
\end{em}

It is sometimes necessary to prevent \LaTeX\ from
breaking a line where it might otherwise do so.
This may be at a space, as between the ``Mr.''\ and
``Jones'' in
       ``Mr.~Jones'',        % ~ produces an unbreakable interword space.
or within a word---especially when the word is a
symbol like
       \mbox{\emph{itemnum}}
that makes little sense when hyphenated across
lines.

Footnotes\footnote{This is an example of a footnote.}
pose no problem.

\LaTeX\ is good at typesetting mathematical formulas
like
       \( x-3y + z = 7 \)
or
       \( a_{1} > x^{2n} + y^{2n} > x' \)
or
       \( \ip{A}{B} = \sum_{i} a_{i} b_{i} \).
The spaces you type in a formula are
ignored.  Remember that a letter like
       $x$                   % $ ... $  and  \( ... \)  are equivalent
is a formula when it denotes a mathematical
symbol, and it should be typed as one.

\section{Displayed Text}

Text is displayed by indenting it from the left
margin.  Quotations are commonly displayed.  There
are short quotations
\begin{quote}
   This is a short quotation.  It consists of a
   single paragraph of text.  See how it is formatted.
\end{quote}
and longer ones.
\begin{quotation}
   This is a longer quotation.  It consists of two
   paragraphs of text, neither of which are
   particularly interesting.

   This is the second paragraph of the quotation.  It
   is just as dull as the first paragraph.
\end{quotation}
Another frequently-displayed structure is a list.
The following is an example of an \emph{itemized}
list.
\begin{itemize}
   \item This is the first item of an itemized list.
         Each item in the list is marked with a ``tick''.
         You don't have to worry about what kind of tick
         mark is used.

   \item This is the second item of the list.  It
         contains another list nested inside it.  The inner
         list is an \emph{enumerated} list.
         \begin{enumerate}
            \item This is the first item of an enumerated
                  list that is nested within the itemized list.

            \item This is the second item of the inner list.
                  \LaTeX\ allows you to nest lists deeper than
                  you really should.
         \end{enumerate}
         This is the rest of the second item of the outer
         list.  It is no more interesting than any other
         part of the item.
   \item This is the third item of the list.
\end{itemize}
You can even display poetry.
\begin{verse}
   There is an environment
    for verse \\             % The \\ command separates lines
   Whose features some poets % within a stanza.
   will curse.

                             % One or more blank lines separate stanzas.

   For instead of making\\
   Them do \emph{all} line breaking, \\
   It allows them to put too many words on a line when they'd rather be
   forced to be terse.
\end{verse}

Mathematical formulas may also be displayed.  A
displayed formula
is
one-line long; multiline
formulas require special formatting instructions.
   \[  \ip{\Gamma}{\psi'} = x'' + y^{2} + z_{i}^{n}\]
Don't start a paragraph with a displayed equation,
nor make one a paragraph by itself.

\end{document}               % End of document.

\section{Bench Huge Section 1252}
% synthetic bench section 1252
% This is a sample LaTeX input file.  (Version of 12 August 2004.)
%
% A '%' character causes TeX to ignore all remaining text on the line,
% and is used for comments like this one.

\documentclass{article}      % Specifies the document class

                             % The preamble begins here.
\title{An Example Document}  % Declares the document's title.
\author{Leslie Lamport}      % Declares the author's name.
\date{January 21, 1994}      % Deleting this command produces today's date.

\newcommand{\ip}[2]{(#1, #2)}
                             % Defines \ip{arg1}{arg2} to mean
                             % (arg1, arg2).

%\newcommand{\ip}[2]{\langle #1 | #2\rangle}
                             % This is an alternative definition of
                             % \ip that is commented out.

\begin{document}             % End of preamble and beginning of text.

\maketitle                   % Produces the title.

This is an example input file.  Comparing it with
the output it generates can show you how to
produce a simple document of your own.

\section{Ordinary Text}      % Produces section heading.  Lower-level
                             % sections are begun with similar
                             % \subsection and \subsubsection commands.

The ends  of words and sentences are marked
  by   spaces. It  doesn't matter how many
spaces    you type; one is as good as 100.  The
end of   a line counts as a space.

One   or more   blank lines denote the  end
of  a paragraph.

Since any number of consecutive spaces are treated
like a single one, the formatting of the input
file makes no difference to
      \LaTeX,                % The \LaTeX command generates the LaTeX logo.
but it makes a difference to you.  When you use
\LaTeX, making your input file as easy to read
as possible will be a great help as you write
your document and when you change it.  This sample
file shows how you can add comments to your own input
file.

Because printing is different from typewriting,
there are a number of things that you have to do
differently when preparing an input file than if
you were just typing the document directly.
Quotation marks like
       ``this''
have to be handled specially, as do quotes within
quotes:
       ``\,`this'            % \, separates the double and single quote.
        is what I just
        wrote, not  `that'\,''.

Dashes come in three sizes: an
       intra-word
dash, a medium dash for number ranges like
       1--2,
and a punctuation
       dash---like
this.

A sentence-ending space should be larger than the
space between words within a sentence.  You
sometimes have to type special commands in
conjunction with punctuation characters to get
this right, as in the following sentence.
       Gnats, gnus, etc.\ all  % `\ ' makes an inter-word space.
       begin with G\@.         % \@ marks end-of-sentence punctuation.
You should check the spaces after periods when
reading your output to make sure you haven't
forgotten any special cases.  Generating an
ellipsis
       \ldots\               % `\ ' is needed after `\ldots' because TeX
                             % ignores spaces after command names like \ldots
                             % made from \ + letters.
                             %
                             % Note how a `%' character causes TeX to ignore
                             % the end of the input line, so these blank lines
                             % do not start a new paragraph.
                             %
with the right spacing around the periods requires
a special command.

\LaTeX\ interprets some common characters as
commands, so you must type special commands to
generate them.  These characters include the
following:
       \$ \& \% \# \{ and \}.

In printing, text is usually emphasized with an
       \emph{italic}
type style.

\begin{em}
   A long segment of text can also be emphasized
   in this way.  Text within such a segment can be
   given \emph{additional} emphasis.
\end{em}

It is sometimes necessary to prevent \LaTeX\ from
breaking a line where it might otherwise do so.
This may be at a space, as between the ``Mr.''\ and
``Jones'' in
       ``Mr.~Jones'',        % ~ produces an unbreakable interword space.
or within a word---especially when the word is a
symbol like
       \mbox{\emph{itemnum}}
that makes little sense when hyphenated across
lines.

Footnotes\footnote{This is an example of a footnote.}
pose no problem.

\LaTeX\ is good at typesetting mathematical formulas
like
       \( x-3y + z = 7 \)
or
       \( a_{1} > x^{2n} + y^{2n} > x' \)
or
       \( \ip{A}{B} = \sum_{i} a_{i} b_{i} \).
The spaces you type in a formula are
ignored.  Remember that a letter like
       $x$                   % $ ... $  and  \( ... \)  are equivalent
is a formula when it denotes a mathematical
symbol, and it should be typed as one.

\section{Displayed Text}

Text is displayed by indenting it from the left
margin.  Quotations are commonly displayed.  There
are short quotations
\begin{quote}
   This is a short quotation.  It consists of a
   single paragraph of text.  See how it is formatted.
\end{quote}
and longer ones.
\begin{quotation}
   This is a longer quotation.  It consists of two
   paragraphs of text, neither of which are
   particularly interesting.

   This is the second paragraph of the quotation.  It
   is just as dull as the first paragraph.
\end{quotation}
Another frequently-displayed structure is a list.
The following is an example of an \emph{itemized}
list.
\begin{itemize}
   \item This is the first item of an itemized list.
         Each item in the list is marked with a ``tick''.
         You don't have to worry about what kind of tick
         mark is used.

   \item This is the second item of the list.  It
         contains another list nested inside it.  The inner
         list is an \emph{enumerated} list.
         \begin{enumerate}
            \item This is the first item of an enumerated
                  list that is nested within the itemized list.

            \item This is the second item of the inner list.
                  \LaTeX\ allows you to nest lists deeper than
                  you really should.
         \end{enumerate}
         This is the rest of the second item of the outer
         list.  It is no more interesting than any other
         part of the item.
   \item This is the third item of the list.
\end{itemize}
You can even display poetry.
\begin{verse}
   There is an environment
    for verse \\             % The \\ command separates lines
   Whose features some poets % within a stanza.
   will curse.

                             % One or more blank lines separate stanzas.

   For instead of making\\
   Them do \emph{all} line breaking, \\
   It allows them to put too many words on a line when they'd rather be
   forced to be terse.
\end{verse}

Mathematical formulas may also be displayed.  A
displayed formula
is
one-line long; multiline
formulas require special formatting instructions.
   \[  \ip{\Gamma}{\psi'} = x'' + y^{2} + z_{i}^{n}\]
Don't start a paragraph with a displayed equation,
nor make one a paragraph by itself.

\end{document}               % End of document.

\section{Bench Huge Section 1253}
% synthetic bench section 1253
% This is a sample LaTeX input file.  (Version of 12 August 2004.)
%
% A '%' character causes TeX to ignore all remaining text on the line,
% and is used for comments like this one.

\documentclass{article}      % Specifies the document class

                             % The preamble begins here.
\title{An Example Document}  % Declares the document's title.
\author{Leslie Lamport}      % Declares the author's name.
\date{January 21, 1994}      % Deleting this command produces today's date.

\newcommand{\ip}[2]{(#1, #2)}
                             % Defines \ip{arg1}{arg2} to mean
                             % (arg1, arg2).

%\newcommand{\ip}[2]{\langle #1 | #2\rangle}
                             % This is an alternative definition of
                             % \ip that is commented out.

\begin{document}             % End of preamble and beginning of text.

\maketitle                   % Produces the title.

This is an example input file.  Comparing it with
the output it generates can show you how to
produce a simple document of your own.

\section{Ordinary Text}      % Produces section heading.  Lower-level
                             % sections are begun with similar
                             % \subsection and \subsubsection commands.

The ends  of words and sentences are marked
  by   spaces. It  doesn't matter how many
spaces    you type; one is as good as 100.  The
end of   a line counts as a space.

One   or more   blank lines denote the  end
of  a paragraph.

Since any number of consecutive spaces are treated
like a single one, the formatting of the input
file makes no difference to
      \LaTeX,                % The \LaTeX command generates the LaTeX logo.
but it makes a difference to you.  When you use
\LaTeX, making your input file as easy to read
as possible will be a great help as you write
your document and when you change it.  This sample
file shows how you can add comments to your own input
file.

Because printing is different from typewriting,
there are a number of things that you have to do
differently when preparing an input file than if
you were just typing the document directly.
Quotation marks like
       ``this''
have to be handled specially, as do quotes within
quotes:
       ``\,`this'            % \, separates the double and single quote.
        is what I just
        wrote, not  `that'\,''.

Dashes come in three sizes: an
       intra-word
dash, a medium dash for number ranges like
       1--2,
and a punctuation
       dash---like
this.

A sentence-ending space should be larger than the
space between words within a sentence.  You
sometimes have to type special commands in
conjunction with punctuation characters to get
this right, as in the following sentence.
       Gnats, gnus, etc.\ all  % `\ ' makes an inter-word space.
       begin with G\@.         % \@ marks end-of-sentence punctuation.
You should check the spaces after periods when
reading your output to make sure you haven't
forgotten any special cases.  Generating an
ellipsis
       \ldots\               % `\ ' is needed after `\ldots' because TeX
                             % ignores spaces after command names like \ldots
                             % made from \ + letters.
                             %
                             % Note how a `%' character causes TeX to ignore
                             % the end of the input line, so these blank lines
                             % do not start a new paragraph.
                             %
with the right spacing around the periods requires
a special command.

\LaTeX\ interprets some common characters as
commands, so you must type special commands to
generate them.  These characters include the
following:
       \$ \& \% \# \{ and \}.

In printing, text is usually emphasized with an
       \emph{italic}
type style.

\begin{em}
   A long segment of text can also be emphasized
   in this way.  Text within such a segment can be
   given \emph{additional} emphasis.
\end{em}

It is sometimes necessary to prevent \LaTeX\ from
breaking a line where it might otherwise do so.
This may be at a space, as between the ``Mr.''\ and
``Jones'' in
       ``Mr.~Jones'',        % ~ produces an unbreakable interword space.
or within a word---especially when the word is a
symbol like
       \mbox{\emph{itemnum}}
that makes little sense when hyphenated across
lines.

Footnotes\footnote{This is an example of a footnote.}
pose no problem.

\LaTeX\ is good at typesetting mathematical formulas
like
       \( x-3y + z = 7 \)
or
       \( a_{1} > x^{2n} + y^{2n} > x' \)
or
       \( \ip{A}{B} = \sum_{i} a_{i} b_{i} \).
The spaces you type in a formula are
ignored.  Remember that a letter like
       $x$                   % $ ... $  and  \( ... \)  are equivalent
is a formula when it denotes a mathematical
symbol, and it should be typed as one.

\section{Displayed Text}

Text is displayed by indenting it from the left
margin.  Quotations are commonly displayed.  There
are short quotations
\begin{quote}
   This is a short quotation.  It consists of a
   single paragraph of text.  See how it is formatted.
\end{quote}
and longer ones.
\begin{quotation}
   This is a longer quotation.  It consists of two
   paragraphs of text, neither of which are
   particularly interesting.

   This is the second paragraph of the quotation.  It
   is just as dull as the first paragraph.
\end{quotation}
Another frequently-displayed structure is a list.
The following is an example of an \emph{itemized}
list.
\begin{itemize}
   \item This is the first item of an itemized list.
         Each item in the list is marked with a ``tick''.
         You don't have to worry about what kind of tick
         mark is used.

   \item This is the second item of the list.  It
         contains another list nested inside it.  The inner
         list is an \emph{enumerated} list.
         \begin{enumerate}
            \item This is the first item of an enumerated
                  list that is nested within the itemized list.

            \item This is the second item of the inner list.
                  \LaTeX\ allows you to nest lists deeper than
                  you really should.
         \end{enumerate}
         This is the rest of the second item of the outer
         list.  It is no more interesting than any other
         part of the item.
   \item This is the third item of the list.
\end{itemize}
You can even display poetry.
\begin{verse}
   There is an environment
    for verse \\             % The \\ command separates lines
   Whose features some poets % within a stanza.
   will curse.

                             % One or more blank lines separate stanzas.

   For instead of making\\
   Them do \emph{all} line breaking, \\
   It allows them to put too many words on a line when they'd rather be
   forced to be terse.
\end{verse}

Mathematical formulas may also be displayed.  A
displayed formula
is
one-line long; multiline
formulas require special formatting instructions.
   \[  \ip{\Gamma}{\psi'} = x'' + y^{2} + z_{i}^{n}\]
Don't start a paragraph with a displayed equation,
nor make one a paragraph by itself.

\end{document}               % End of document.

\section{Bench Huge Section 1254}
% synthetic bench section 1254
% This is a sample LaTeX input file.  (Version of 12 August 2004.)
%
% A '%' character causes TeX to ignore all remaining text on the line,
% and is used for comments like this one.

\documentclass{article}      % Specifies the document class

                             % The preamble begins here.
\title{An Example Document}  % Declares the document's title.
\author{Leslie Lamport}      % Declares the author's name.
\date{January 21, 1994}      % Deleting this command produces today's date.

\newcommand{\ip}[2]{(#1, #2)}
                             % Defines \ip{arg1}{arg2} to mean
                             % (arg1, arg2).

%\newcommand{\ip}[2]{\langle #1 | #2\rangle}
                             % This is an alternative definition of
                             % \ip that is commented out.

\begin{document}             % End of preamble and beginning of text.

\maketitle                   % Produces the title.

This is an example input file.  Comparing it with
the output it generates can show you how to
produce a simple document of your own.

\section{Ordinary Text}      % Produces section heading.  Lower-level
                             % sections are begun with similar
                             % \subsection and \subsubsection commands.

The ends  of words and sentences are marked
  by   spaces. It  doesn't matter how many
spaces    you type; one is as good as 100.  The
end of   a line counts as a space.

One   or more   blank lines denote the  end
of  a paragraph.

Since any number of consecutive spaces are treated
like a single one, the formatting of the input
file makes no difference to
      \LaTeX,                % The \LaTeX command generates the LaTeX logo.
but it makes a difference to you.  When you use
\LaTeX, making your input file as easy to read
as possible will be a great help as you write
your document and when you change it.  This sample
file shows how you can add comments to your own input
file.

Because printing is different from typewriting,
there are a number of things that you have to do
differently when preparing an input file than if
you were just typing the document directly.
Quotation marks like
       ``this''
have to be handled specially, as do quotes within
quotes:
       ``\,`this'            % \, separates the double and single quote.
        is what I just
        wrote, not  `that'\,''.

Dashes come in three sizes: an
       intra-word
dash, a medium dash for number ranges like
       1--2,
and a punctuation
       dash---like
this.

A sentence-ending space should be larger than the
space between words within a sentence.  You
sometimes have to type special commands in
conjunction with punctuation characters to get
this right, as in the following sentence.
       Gnats, gnus, etc.\ all  % `\ ' makes an inter-word space.
       begin with G\@.         % \@ marks end-of-sentence punctuation.
You should check the spaces after periods when
reading your output to make sure you haven't
forgotten any special cases.  Generating an
ellipsis
       \ldots\               % `\ ' is needed after `\ldots' because TeX
                             % ignores spaces after command names like \ldots
                             % made from \ + letters.
                             %
                             % Note how a `%' character causes TeX to ignore
                             % the end of the input line, so these blank lines
                             % do not start a new paragraph.
                             %
with the right spacing around the periods requires
a special command.

\LaTeX\ interprets some common characters as
commands, so you must type special commands to
generate them.  These characters include the
following:
       \$ \& \% \# \{ and \}.

In printing, text is usually emphasized with an
       \emph{italic}
type style.

\begin{em}
   A long segment of text can also be emphasized
   in this way.  Text within such a segment can be
   given \emph{additional} emphasis.
\end{em}

It is sometimes necessary to prevent \LaTeX\ from
breaking a line where it might otherwise do so.
This may be at a space, as between the ``Mr.''\ and
``Jones'' in
       ``Mr.~Jones'',        % ~ produces an unbreakable interword space.
or within a word---especially when the word is a
symbol like
       \mbox{\emph{itemnum}}
that makes little sense when hyphenated across
lines.

Footnotes\footnote{This is an example of a footnote.}
pose no problem.

\LaTeX\ is good at typesetting mathematical formulas
like
       \( x-3y + z = 7 \)
or
       \( a_{1} > x^{2n} + y^{2n} > x' \)
or
       \( \ip{A}{B} = \sum_{i} a_{i} b_{i} \).
The spaces you type in a formula are
ignored.  Remember that a letter like
       $x$                   % $ ... $  and  \( ... \)  are equivalent
is a formula when it denotes a mathematical
symbol, and it should be typed as one.

\section{Displayed Text}

Text is displayed by indenting it from the left
margin.  Quotations are commonly displayed.  There
are short quotations
\begin{quote}
   This is a short quotation.  It consists of a
   single paragraph of text.  See how it is formatted.
\end{quote}
and longer ones.
\begin{quotation}
   This is a longer quotation.  It consists of two
   paragraphs of text, neither of which are
   particularly interesting.

   This is the second paragraph of the quotation.  It
   is just as dull as the first paragraph.
\end{quotation}
Another frequently-displayed structure is a list.
The following is an example of an \emph{itemized}
list.
\begin{itemize}
   \item This is the first item of an itemized list.
         Each item in the list is marked with a ``tick''.
         You don't have to worry about what kind of tick
         mark is used.

   \item This is the second item of the list.  It
         contains another list nested inside it.  The inner
         list is an \emph{enumerated} list.
         \begin{enumerate}
            \item This is the first item of an enumerated
                  list that is nested within the itemized list.

            \item This is the second item of the inner list.
                  \LaTeX\ allows you to nest lists deeper than
                  you really should.
         \end{enumerate}
         This is the rest of the second item of the outer
         list.  It is no more interesting than any other
         part of the item.
   \item This is the third item of the list.
\end{itemize}
You can even display poetry.
\begin{verse}
   There is an environment
    for verse \\             % The \\ command separates lines
   Whose features some poets % within a stanza.
   will curse.

                             % One or more blank lines separate stanzas.

   For instead of making\\
   Them do \emph{all} line breaking, \\
   It allows them to put too many words on a line when they'd rather be
   forced to be terse.
\end{verse}

Mathematical formulas may also be displayed.  A
displayed formula
is
one-line long; multiline
formulas require special formatting instructions.
   \[  \ip{\Gamma}{\psi'} = x'' + y^{2} + z_{i}^{n}\]
Don't start a paragraph with a displayed equation,
nor make one a paragraph by itself.

\end{document}               % End of document.

\section{Bench Huge Section 1255}
% synthetic bench section 1255
% This is a sample LaTeX input file.  (Version of 12 August 2004.)
%
% A '%' character causes TeX to ignore all remaining text on the line,
% and is used for comments like this one.

\documentclass{article}      % Specifies the document class

                             % The preamble begins here.
\title{An Example Document}  % Declares the document's title.
\author{Leslie Lamport}      % Declares the author's name.
\date{January 21, 1994}      % Deleting this command produces today's date.

\newcommand{\ip}[2]{(#1, #2)}
                             % Defines \ip{arg1}{arg2} to mean
                             % (arg1, arg2).

%\newcommand{\ip}[2]{\langle #1 | #2\rangle}
                             % This is an alternative definition of
                             % \ip that is commented out.

\begin{document}             % End of preamble and beginning of text.

\maketitle                   % Produces the title.

This is an example input file.  Comparing it with
the output it generates can show you how to
produce a simple document of your own.

\section{Ordinary Text}      % Produces section heading.  Lower-level
                             % sections are begun with similar
                             % \subsection and \subsubsection commands.

The ends  of words and sentences are marked
  by   spaces. It  doesn't matter how many
spaces    you type; one is as good as 100.  The
end of   a line counts as a space.

One   or more   blank lines denote the  end
of  a paragraph.

Since any number of consecutive spaces are treated
like a single one, the formatting of the input
file makes no difference to
      \LaTeX,                % The \LaTeX command generates the LaTeX logo.
but it makes a difference to you.  When you use
\LaTeX, making your input file as easy to read
as possible will be a great help as you write
your document and when you change it.  This sample
file shows how you can add comments to your own input
file.

Because printing is different from typewriting,
there are a number of things that you have to do
differently when preparing an input file than if
you were just typing the document directly.
Quotation marks like
       ``this''
have to be handled specially, as do quotes within
quotes:
       ``\,`this'            % \, separates the double and single quote.
        is what I just
        wrote, not  `that'\,''.

Dashes come in three sizes: an
       intra-word
dash, a medium dash for number ranges like
       1--2,
and a punctuation
       dash---like
this.

A sentence-ending space should be larger than the
space between words within a sentence.  You
sometimes have to type special commands in
conjunction with punctuation characters to get
this right, as in the following sentence.
       Gnats, gnus, etc.\ all  % `\ ' makes an inter-word space.
       begin with G\@.         % \@ marks end-of-sentence punctuation.
You should check the spaces after periods when
reading your output to make sure you haven't
forgotten any special cases.  Generating an
ellipsis
       \ldots\               % `\ ' is needed after `\ldots' because TeX
                             % ignores spaces after command names like \ldots
                             % made from \ + letters.
                             %
                             % Note how a `%' character causes TeX to ignore
                             % the end of the input line, so these blank lines
                             % do not start a new paragraph.
                             %
with the right spacing around the periods requires
a special command.

\LaTeX\ interprets some common characters as
commands, so you must type special commands to
generate them.  These characters include the
following:
       \$ \& \% \# \{ and \}.

In printing, text is usually emphasized with an
       \emph{italic}
type style.

\begin{em}
   A long segment of text can also be emphasized
   in this way.  Text within such a segment can be
   given \emph{additional} emphasis.
\end{em}

It is sometimes necessary to prevent \LaTeX\ from
breaking a line where it might otherwise do so.
This may be at a space, as between the ``Mr.''\ and
``Jones'' in
       ``Mr.~Jones'',        % ~ produces an unbreakable interword space.
or within a word---especially when the word is a
symbol like
       \mbox{\emph{itemnum}}
that makes little sense when hyphenated across
lines.

Footnotes\footnote{This is an example of a footnote.}
pose no problem.

\LaTeX\ is good at typesetting mathematical formulas
like
       \( x-3y + z = 7 \)
or
       \( a_{1} > x^{2n} + y^{2n} > x' \)
or
       \( \ip{A}{B} = \sum_{i} a_{i} b_{i} \).
The spaces you type in a formula are
ignored.  Remember that a letter like
       $x$                   % $ ... $  and  \( ... \)  are equivalent
is a formula when it denotes a mathematical
symbol, and it should be typed as one.

\section{Displayed Text}

Text is displayed by indenting it from the left
margin.  Quotations are commonly displayed.  There
are short quotations
\begin{quote}
   This is a short quotation.  It consists of a
   single paragraph of text.  See how it is formatted.
\end{quote}
and longer ones.
\begin{quotation}
   This is a longer quotation.  It consists of two
   paragraphs of text, neither of which are
   particularly interesting.

   This is the second paragraph of the quotation.  It
   is just as dull as the first paragraph.
\end{quotation}
Another frequently-displayed structure is a list.
The following is an example of an \emph{itemized}
list.
\begin{itemize}
   \item This is the first item of an itemized list.
         Each item in the list is marked with a ``tick''.
         You don't have to worry about what kind of tick
         mark is used.

   \item This is the second item of the list.  It
         contains another list nested inside it.  The inner
         list is an \emph{enumerated} list.
         \begin{enumerate}
            \item This is the first item of an enumerated
                  list that is nested within the itemized list.

            \item This is the second item of the inner list.
                  \LaTeX\ allows you to nest lists deeper than
                  you really should.
         \end{enumerate}
         This is the rest of the second item of the outer
         list.  It is no more interesting than any other
         part of the item.
   \item This is the third item of the list.
\end{itemize}
You can even display poetry.
\begin{verse}
   There is an environment
    for verse \\             % The \\ command separates lines
   Whose features some poets % within a stanza.
   will curse.

                             % One or more blank lines separate stanzas.

   For instead of making\\
   Them do \emph{all} line breaking, \\
   It allows them to put too many words on a line when they'd rather be
   forced to be terse.
\end{verse}

Mathematical formulas may also be displayed.  A
displayed formula
is
one-line long; multiline
formulas require special formatting instructions.
   \[  \ip{\Gamma}{\psi'} = x'' + y^{2} + z_{i}^{n}\]
Don't start a paragraph with a displayed equation,
nor make one a paragraph by itself.

\end{document}               % End of document.

\section{Bench Huge Section 1256}
% synthetic bench section 1256
% This is a sample LaTeX input file.  (Version of 12 August 2004.)
%
% A '%' character causes TeX to ignore all remaining text on the line,
% and is used for comments like this one.

\documentclass{article}      % Specifies the document class

                             % The preamble begins here.
\title{An Example Document}  % Declares the document's title.
\author{Leslie Lamport}      % Declares the author's name.
\date{January 21, 1994}      % Deleting this command produces today's date.

\newcommand{\ip}[2]{(#1, #2)}
                             % Defines \ip{arg1}{arg2} to mean
                             % (arg1, arg2).

%\newcommand{\ip}[2]{\langle #1 | #2\rangle}
                             % This is an alternative definition of
                             % \ip that is commented out.

\begin{document}             % End of preamble and beginning of text.

\maketitle                   % Produces the title.

This is an example input file.  Comparing it with
the output it generates can show you how to
produce a simple document of your own.

\section{Ordinary Text}      % Produces section heading.  Lower-level
                             % sections are begun with similar
                             % \subsection and \subsubsection commands.

The ends  of words and sentences are marked
  by   spaces. It  doesn't matter how many
spaces    you type; one is as good as 100.  The
end of   a line counts as a space.

One   or more   blank lines denote the  end
of  a paragraph.

Since any number of consecutive spaces are treated
like a single one, the formatting of the input
file makes no difference to
      \LaTeX,                % The \LaTeX command generates the LaTeX logo.
but it makes a difference to you.  When you use
\LaTeX, making your input file as easy to read
as possible will be a great help as you write
your document and when you change it.  This sample
file shows how you can add comments to your own input
file.

Because printing is different from typewriting,
there are a number of things that you have to do
differently when preparing an input file than if
you were just typing the document directly.
Quotation marks like
       ``this''
have to be handled specially, as do quotes within
quotes:
       ``\,`this'            % \, separates the double and single quote.
        is what I just
        wrote, not  `that'\,''.

Dashes come in three sizes: an
       intra-word
dash, a medium dash for number ranges like
       1--2,
and a punctuation
       dash---like
this.

A sentence-ending space should be larger than the
space between words within a sentence.  You
sometimes have to type special commands in
conjunction with punctuation characters to get
this right, as in the following sentence.
       Gnats, gnus, etc.\ all  % `\ ' makes an inter-word space.
       begin with G\@.         % \@ marks end-of-sentence punctuation.
You should check the spaces after periods when
reading your output to make sure you haven't
forgotten any special cases.  Generating an
ellipsis
       \ldots\               % `\ ' is needed after `\ldots' because TeX
                             % ignores spaces after command names like \ldots
                             % made from \ + letters.
                             %
                             % Note how a `%' character causes TeX to ignore
                             % the end of the input line, so these blank lines
                             % do not start a new paragraph.
                             %
with the right spacing around the periods requires
a special command.

\LaTeX\ interprets some common characters as
commands, so you must type special commands to
generate them.  These characters include the
following:
       \$ \& \% \# \{ and \}.

In printing, text is usually emphasized with an
       \emph{italic}
type style.

\begin{em}
   A long segment of text can also be emphasized
   in this way.  Text within such a segment can be
   given \emph{additional} emphasis.
\end{em}

It is sometimes necessary to prevent \LaTeX\ from
breaking a line where it might otherwise do so.
This may be at a space, as between the ``Mr.''\ and
``Jones'' in
       ``Mr.~Jones'',        % ~ produces an unbreakable interword space.
or within a word---especially when the word is a
symbol like
       \mbox{\emph{itemnum}}
that makes little sense when hyphenated across
lines.

Footnotes\footnote{This is an example of a footnote.}
pose no problem.

\LaTeX\ is good at typesetting mathematical formulas
like
       \( x-3y + z = 7 \)
or
       \( a_{1} > x^{2n} + y^{2n} > x' \)
or
       \( \ip{A}{B} = \sum_{i} a_{i} b_{i} \).
The spaces you type in a formula are
ignored.  Remember that a letter like
       $x$                   % $ ... $  and  \( ... \)  are equivalent
is a formula when it denotes a mathematical
symbol, and it should be typed as one.

\section{Displayed Text}

Text is displayed by indenting it from the left
margin.  Quotations are commonly displayed.  There
are short quotations
\begin{quote}
   This is a short quotation.  It consists of a
   single paragraph of text.  See how it is formatted.
\end{quote}
and longer ones.
\begin{quotation}
   This is a longer quotation.  It consists of two
   paragraphs of text, neither of which are
   particularly interesting.

   This is the second paragraph of the quotation.  It
   is just as dull as the first paragraph.
\end{quotation}
Another frequently-displayed structure is a list.
The following is an example of an \emph{itemized}
list.
\begin{itemize}
   \item This is the first item of an itemized list.
         Each item in the list is marked with a ``tick''.
         You don't have to worry about what kind of tick
         mark is used.

   \item This is the second item of the list.  It
         contains another list nested inside it.  The inner
         list is an \emph{enumerated} list.
         \begin{enumerate}
            \item This is the first item of an enumerated
                  list that is nested within the itemized list.

            \item This is the second item of the inner list.
                  \LaTeX\ allows you to nest lists deeper than
                  you really should.
         \end{enumerate}
         This is the rest of the second item of the outer
         list.  It is no more interesting than any other
         part of the item.
   \item This is the third item of the list.
\end{itemize}
You can even display poetry.
\begin{verse}
   There is an environment
    for verse \\             % The \\ command separates lines
   Whose features some poets % within a stanza.
   will curse.

                             % One or more blank lines separate stanzas.

   For instead of making\\
   Them do \emph{all} line breaking, \\
   It allows them to put too many words on a line when they'd rather be
   forced to be terse.
\end{verse}

Mathematical formulas may also be displayed.  A
displayed formula
is
one-line long; multiline
formulas require special formatting instructions.
   \[  \ip{\Gamma}{\psi'} = x'' + y^{2} + z_{i}^{n}\]
Don't start a paragraph with a displayed equation,
nor make one a paragraph by itself.

\end{document}               % End of document.

\section{Bench Huge Section 1257}
% synthetic bench section 1257
% This is a sample LaTeX input file.  (Version of 12 August 2004.)
%
% A '%' character causes TeX to ignore all remaining text on the line,
% and is used for comments like this one.

\documentclass{article}      % Specifies the document class

                             % The preamble begins here.
\title{An Example Document}  % Declares the document's title.
\author{Leslie Lamport}      % Declares the author's name.
\date{January 21, 1994}      % Deleting this command produces today's date.

\newcommand{\ip}[2]{(#1, #2)}
                             % Defines \ip{arg1}{arg2} to mean
                             % (arg1, arg2).

%\newcommand{\ip}[2]{\langle #1 | #2\rangle}
                             % This is an alternative definition of
                             % \ip that is commented out.

\begin{document}             % End of preamble and beginning of text.

\maketitle                   % Produces the title.

This is an example input file.  Comparing it with
the output it generates can show you how to
produce a simple document of your own.

\section{Ordinary Text}      % Produces section heading.  Lower-level
                             % sections are begun with similar
                             % \subsection and \subsubsection commands.

The ends  of words and sentences are marked
  by   spaces. It  doesn't matter how many
spaces    you type; one is as good as 100.  The
end of   a line counts as a space.

One   or more   blank lines denote the  end
of  a paragraph.

Since any number of consecutive spaces are treated
like a single one, the formatting of the input
file makes no difference to
      \LaTeX,                % The \LaTeX command generates the LaTeX logo.
but it makes a difference to you.  When you use
\LaTeX, making your input file as easy to read
as possible will be a great help as you write
your document and when you change it.  This sample
file shows how you can add comments to your own input
file.

Because printing is different from typewriting,
there are a number of things that you have to do
differently when preparing an input file than if
you were just typing the document directly.
Quotation marks like
       ``this''
have to be handled specially, as do quotes within
quotes:
       ``\,`this'            % \, separates the double and single quote.
        is what I just
        wrote, not  `that'\,''.

Dashes come in three sizes: an
       intra-word
dash, a medium dash for number ranges like
       1--2,
and a punctuation
       dash---like
this.

A sentence-ending space should be larger than the
space between words within a sentence.  You
sometimes have to type special commands in
conjunction with punctuation characters to get
this right, as in the following sentence.
       Gnats, gnus, etc.\ all  % `\ ' makes an inter-word space.
       begin with G\@.         % \@ marks end-of-sentence punctuation.
You should check the spaces after periods when
reading your output to make sure you haven't
forgotten any special cases.  Generating an
ellipsis
       \ldots\               % `\ ' is needed after `\ldots' because TeX
                             % ignores spaces after command names like \ldots
                             % made from \ + letters.
                             %
                             % Note how a `%' character causes TeX to ignore
                             % the end of the input line, so these blank lines
                             % do not start a new paragraph.
                             %
with the right spacing around the periods requires
a special command.

\LaTeX\ interprets some common characters as
commands, so you must type special commands to
generate them.  These characters include the
following:
       \$ \& \% \# \{ and \}.

In printing, text is usually emphasized with an
       \emph{italic}
type style.

\begin{em}
   A long segment of text can also be emphasized
   in this way.  Text within such a segment can be
   given \emph{additional} emphasis.
\end{em}

It is sometimes necessary to prevent \LaTeX\ from
breaking a line where it might otherwise do so.
This may be at a space, as between the ``Mr.''\ and
``Jones'' in
       ``Mr.~Jones'',        % ~ produces an unbreakable interword space.
or within a word---especially when the word is a
symbol like
       \mbox{\emph{itemnum}}
that makes little sense when hyphenated across
lines.

Footnotes\footnote{This is an example of a footnote.}
pose no problem.

\LaTeX\ is good at typesetting mathematical formulas
like
       \( x-3y + z = 7 \)
or
       \( a_{1} > x^{2n} + y^{2n} > x' \)
or
       \( \ip{A}{B} = \sum_{i} a_{i} b_{i} \).
The spaces you type in a formula are
ignored.  Remember that a letter like
       $x$                   % $ ... $  and  \( ... \)  are equivalent
is a formula when it denotes a mathematical
symbol, and it should be typed as one.

\section{Displayed Text}

Text is displayed by indenting it from the left
margin.  Quotations are commonly displayed.  There
are short quotations
\begin{quote}
   This is a short quotation.  It consists of a
   single paragraph of text.  See how it is formatted.
\end{quote}
and longer ones.
\begin{quotation}
   This is a longer quotation.  It consists of two
   paragraphs of text, neither of which are
   particularly interesting.

   This is the second paragraph of the quotation.  It
   is just as dull as the first paragraph.
\end{quotation}
Another frequently-displayed structure is a list.
The following is an example of an \emph{itemized}
list.
\begin{itemize}
   \item This is the first item of an itemized list.
         Each item in the list is marked with a ``tick''.
         You don't have to worry about what kind of tick
         mark is used.

   \item This is the second item of the list.  It
         contains another list nested inside it.  The inner
         list is an \emph{enumerated} list.
         \begin{enumerate}
            \item This is the first item of an enumerated
                  list that is nested within the itemized list.

            \item This is the second item of the inner list.
                  \LaTeX\ allows you to nest lists deeper than
                  you really should.
         \end{enumerate}
         This is the rest of the second item of the outer
         list.  It is no more interesting than any other
         part of the item.
   \item This is the third item of the list.
\end{itemize}
You can even display poetry.
\begin{verse}
   There is an environment
    for verse \\             % The \\ command separates lines
   Whose features some poets % within a stanza.
   will curse.

                             % One or more blank lines separate stanzas.

   For instead of making\\
   Them do \emph{all} line breaking, \\
   It allows them to put too many words on a line when they'd rather be
   forced to be terse.
\end{verse}

Mathematical formulas may also be displayed.  A
displayed formula
is
one-line long; multiline
formulas require special formatting instructions.
   \[  \ip{\Gamma}{\psi'} = x'' + y^{2} + z_{i}^{n}\]
Don't start a paragraph with a displayed equation,
nor make one a paragraph by itself.

\end{document}               % End of document.

\section{Bench Huge Section 1258}
% synthetic bench section 1258
% This is a sample LaTeX input file.  (Version of 12 August 2004.)
%
% A '%' character causes TeX to ignore all remaining text on the line,
% and is used for comments like this one.

\documentclass{article}      % Specifies the document class

                             % The preamble begins here.
\title{An Example Document}  % Declares the document's title.
\author{Leslie Lamport}      % Declares the author's name.
\date{January 21, 1994}      % Deleting this command produces today's date.

\newcommand{\ip}[2]{(#1, #2)}
                             % Defines \ip{arg1}{arg2} to mean
                             % (arg1, arg2).

%\newcommand{\ip}[2]{\langle #1 | #2\rangle}
                             % This is an alternative definition of
                             % \ip that is commented out.

\begin{document}             % End of preamble and beginning of text.

\maketitle                   % Produces the title.

This is an example input file.  Comparing it with
the output it generates can show you how to
produce a simple document of your own.

\section{Ordinary Text}      % Produces section heading.  Lower-level
                             % sections are begun with similar
                             % \subsection and \subsubsection commands.

The ends  of words and sentences are marked
  by   spaces. It  doesn't matter how many
spaces    you type; one is as good as 100.  The
end of   a line counts as a space.

One   or more   blank lines denote the  end
of  a paragraph.

Since any number of consecutive spaces are treated
like a single one, the formatting of the input
file makes no difference to
      \LaTeX,                % The \LaTeX command generates the LaTeX logo.
but it makes a difference to you.  When you use
\LaTeX, making your input file as easy to read
as possible will be a great help as you write
your document and when you change it.  This sample
file shows how you can add comments to your own input
file.

Because printing is different from typewriting,
there are a number of things that you have to do
differently when preparing an input file than if
you were just typing the document directly.
Quotation marks like
       ``this''
have to be handled specially, as do quotes within
quotes:
       ``\,`this'            % \, separates the double and single quote.
        is what I just
        wrote, not  `that'\,''.

Dashes come in three sizes: an
       intra-word
dash, a medium dash for number ranges like
       1--2,
and a punctuation
       dash---like
this.

A sentence-ending space should be larger than the
space between words within a sentence.  You
sometimes have to type special commands in
conjunction with punctuation characters to get
this right, as in the following sentence.
       Gnats, gnus, etc.\ all  % `\ ' makes an inter-word space.
       begin with G\@.         % \@ marks end-of-sentence punctuation.
You should check the spaces after periods when
reading your output to make sure you haven't
forgotten any special cases.  Generating an
ellipsis
       \ldots\               % `\ ' is needed after `\ldots' because TeX
                             % ignores spaces after command names like \ldots
                             % made from \ + letters.
                             %
                             % Note how a `%' character causes TeX to ignore
                             % the end of the input line, so these blank lines
                             % do not start a new paragraph.
                             %
with the right spacing around the periods requires
a special command.

\LaTeX\ interprets some common characters as
commands, so you must type special commands to
generate them.  These characters include the
following:
       \$ \& \% \# \{ and \}.

In printing, text is usually emphasized with an
       \emph{italic}
type style.

\begin{em}
   A long segment of text can also be emphasized
   in this way.  Text within such a segment can be
   given \emph{additional} emphasis.
\end{em}

It is sometimes necessary to prevent \LaTeX\ from
breaking a line where it might otherwise do so.
This may be at a space, as between the ``Mr.''\ and
``Jones'' in
       ``Mr.~Jones'',        % ~ produces an unbreakable interword space.
or within a word---especially when the word is a
symbol like
       \mbox{\emph{itemnum}}
that makes little sense when hyphenated across
lines.

Footnotes\footnote{This is an example of a footnote.}
pose no problem.

\LaTeX\ is good at typesetting mathematical formulas
like
       \( x-3y + z = 7 \)
or
       \( a_{1} > x^{2n} + y^{2n} > x' \)
or
       \( \ip{A}{B} = \sum_{i} a_{i} b_{i} \).
The spaces you type in a formula are
ignored.  Remember that a letter like
       $x$                   % $ ... $  and  \( ... \)  are equivalent
is a formula when it denotes a mathematical
symbol, and it should be typed as one.

\section{Displayed Text}

Text is displayed by indenting it from the left
margin.  Quotations are commonly displayed.  There
are short quotations
\begin{quote}
   This is a short quotation.  It consists of a
   single paragraph of text.  See how it is formatted.
\end{quote}
and longer ones.
\begin{quotation}
   This is a longer quotation.  It consists of two
   paragraphs of text, neither of which are
   particularly interesting.

   This is the second paragraph of the quotation.  It
   is just as dull as the first paragraph.
\end{quotation}
Another frequently-displayed structure is a list.
The following is an example of an \emph{itemized}
list.
\begin{itemize}
   \item This is the first item of an itemized list.
         Each item in the list is marked with a ``tick''.
         You don't have to worry about what kind of tick
         mark is used.

   \item This is the second item of the list.  It
         contains another list nested inside it.  The inner
         list is an \emph{enumerated} list.
         \begin{enumerate}
            \item This is the first item of an enumerated
                  list that is nested within the itemized list.

            \item This is the second item of the inner list.
                  \LaTeX\ allows you to nest lists deeper than
                  you really should.
         \end{enumerate}
         This is the rest of the second item of the outer
         list.  It is no more interesting than any other
         part of the item.
   \item This is the third item of the list.
\end{itemize}
You can even display poetry.
\begin{verse}
   There is an environment
    for verse \\             % The \\ command separates lines
   Whose features some poets % within a stanza.
   will curse.

                             % One or more blank lines separate stanzas.

   For instead of making\\
   Them do \emph{all} line breaking, \\
   It allows them to put too many words on a line when they'd rather be
   forced to be terse.
\end{verse}

Mathematical formulas may also be displayed.  A
displayed formula
is
one-line long; multiline
formulas require special formatting instructions.
   \[  \ip{\Gamma}{\psi'} = x'' + y^{2} + z_{i}^{n}\]
Don't start a paragraph with a displayed equation,
nor make one a paragraph by itself.

\end{document}               % End of document.

\section{Bench Huge Section 1259}
% synthetic bench section 1259
% This is a sample LaTeX input file.  (Version of 12 August 2004.)
%
% A '%' character causes TeX to ignore all remaining text on the line,
% and is used for comments like this one.

\documentclass{article}      % Specifies the document class

                             % The preamble begins here.
\title{An Example Document}  % Declares the document's title.
\author{Leslie Lamport}      % Declares the author's name.
\date{January 21, 1994}      % Deleting this command produces today's date.

\newcommand{\ip}[2]{(#1, #2)}
                             % Defines \ip{arg1}{arg2} to mean
                             % (arg1, arg2).

%\newcommand{\ip}[2]{\langle #1 | #2\rangle}
                             % This is an alternative definition of
                             % \ip that is commented out.

\begin{document}             % End of preamble and beginning of text.

\maketitle                   % Produces the title.

This is an example input file.  Comparing it with
the output it generates can show you how to
produce a simple document of your own.

\section{Ordinary Text}      % Produces section heading.  Lower-level
                             % sections are begun with similar
                             % \subsection and \subsubsection commands.

The ends  of words and sentences are marked
  by   spaces. It  doesn't matter how many
spaces    you type; one is as good as 100.  The
end of   a line counts as a space.

One   or more   blank lines denote the  end
of  a paragraph.

Since any number of consecutive spaces are treated
like a single one, the formatting of the input
file makes no difference to
      \LaTeX,                % The \LaTeX command generates the LaTeX logo.
but it makes a difference to you.  When you use
\LaTeX, making your input file as easy to read
as possible will be a great help as you write
your document and when you change it.  This sample
file shows how you can add comments to your own input
file.

Because printing is different from typewriting,
there are a number of things that you have to do
differently when preparing an input file than if
you were just typing the document directly.
Quotation marks like
       ``this''
have to be handled specially, as do quotes within
quotes:
       ``\,`this'            % \, separates the double and single quote.
        is what I just
        wrote, not  `that'\,''.

Dashes come in three sizes: an
       intra-word
dash, a medium dash for number ranges like
       1--2,
and a punctuation
       dash---like
this.

A sentence-ending space should be larger than the
space between words within a sentence.  You
sometimes have to type special commands in
conjunction with punctuation characters to get
this right, as in the following sentence.
       Gnats, gnus, etc.\ all  % `\ ' makes an inter-word space.
       begin with G\@.         % \@ marks end-of-sentence punctuation.
You should check the spaces after periods when
reading your output to make sure you haven't
forgotten any special cases.  Generating an
ellipsis
       \ldots\               % `\ ' is needed after `\ldots' because TeX
                             % ignores spaces after command names like \ldots
                             % made from \ + letters.
                             %
                             % Note how a `%' character causes TeX to ignore
                             % the end of the input line, so these blank lines
                             % do not start a new paragraph.
                             %
with the right spacing around the periods requires
a special command.

\LaTeX\ interprets some common characters as
commands, so you must type special commands to
generate them.  These characters include the
following:
       \$ \& \% \# \{ and \}.

In printing, text is usually emphasized with an
       \emph{italic}
type style.

\begin{em}
   A long segment of text can also be emphasized
   in this way.  Text within such a segment can be
   given \emph{additional} emphasis.
\end{em}

It is sometimes necessary to prevent \LaTeX\ from
breaking a line where it might otherwise do so.
This may be at a space, as between the ``Mr.''\ and
``Jones'' in
       ``Mr.~Jones'',        % ~ produces an unbreakable interword space.
or within a word---especially when the word is a
symbol like
       \mbox{\emph{itemnum}}
that makes little sense when hyphenated across
lines.

Footnotes\footnote{This is an example of a footnote.}
pose no problem.

\LaTeX\ is good at typesetting mathematical formulas
like
       \( x-3y + z = 7 \)
or
       \( a_{1} > x^{2n} + y^{2n} > x' \)
or
       \( \ip{A}{B} = \sum_{i} a_{i} b_{i} \).
The spaces you type in a formula are
ignored.  Remember that a letter like
       $x$                   % $ ... $  and  \( ... \)  are equivalent
is a formula when it denotes a mathematical
symbol, and it should be typed as one.

\section{Displayed Text}

Text is displayed by indenting it from the left
margin.  Quotations are commonly displayed.  There
are short quotations
\begin{quote}
   This is a short quotation.  It consists of a
   single paragraph of text.  See how it is formatted.
\end{quote}
and longer ones.
\begin{quotation}
   This is a longer quotation.  It consists of two
   paragraphs of text, neither of which are
   particularly interesting.

   This is the second paragraph of the quotation.  It
   is just as dull as the first paragraph.
\end{quotation}
Another frequently-displayed structure is a list.
The following is an example of an \emph{itemized}
list.
\begin{itemize}
   \item This is the first item of an itemized list.
         Each item in the list is marked with a ``tick''.
         You don't have to worry about what kind of tick
         mark is used.

   \item This is the second item of the list.  It
         contains another list nested inside it.  The inner
         list is an \emph{enumerated} list.
         \begin{enumerate}
            \item This is the first item of an enumerated
                  list that is nested within the itemized list.

            \item This is the second item of the inner list.
                  \LaTeX\ allows you to nest lists deeper than
                  you really should.
         \end{enumerate}
         This is the rest of the second item of the outer
         list.  It is no more interesting than any other
         part of the item.
   \item This is the third item of the list.
\end{itemize}
You can even display poetry.
\begin{verse}
   There is an environment
    for verse \\             % The \\ command separates lines
   Whose features some poets % within a stanza.
   will curse.

                             % One or more blank lines separate stanzas.

   For instead of making\\
   Them do \emph{all} line breaking, \\
   It allows them to put too many words on a line when they'd rather be
   forced to be terse.
\end{verse}

Mathematical formulas may also be displayed.  A
displayed formula
is
one-line long; multiline
formulas require special formatting instructions.
   \[  \ip{\Gamma}{\psi'} = x'' + y^{2} + z_{i}^{n}\]
Don't start a paragraph with a displayed equation,
nor make one a paragraph by itself.

\end{document}               % End of document.

\section{Bench Huge Section 1260}
% synthetic bench section 1260
% This is a sample LaTeX input file.  (Version of 12 August 2004.)
%
% A '%' character causes TeX to ignore all remaining text on the line,
% and is used for comments like this one.

\documentclass{article}      % Specifies the document class

                             % The preamble begins here.
\title{An Example Document}  % Declares the document's title.
\author{Leslie Lamport}      % Declares the author's name.
\date{January 21, 1994}      % Deleting this command produces today's date.

\newcommand{\ip}[2]{(#1, #2)}
                             % Defines \ip{arg1}{arg2} to mean
                             % (arg1, arg2).

%\newcommand{\ip}[2]{\langle #1 | #2\rangle}
                             % This is an alternative definition of
                             % \ip that is commented out.

\begin{document}             % End of preamble and beginning of text.

\maketitle                   % Produces the title.

This is an example input file.  Comparing it with
the output it generates can show you how to
produce a simple document of your own.

\section{Ordinary Text}      % Produces section heading.  Lower-level
                             % sections are begun with similar
                             % \subsection and \subsubsection commands.

The ends  of words and sentences are marked
  by   spaces. It  doesn't matter how many
spaces    you type; one is as good as 100.  The
end of   a line counts as a space.

One   or more   blank lines denote the  end
of  a paragraph.

Since any number of consecutive spaces are treated
like a single one, the formatting of the input
file makes no difference to
      \LaTeX,                % The \LaTeX command generates the LaTeX logo.
but it makes a difference to you.  When you use
\LaTeX, making your input file as easy to read
as possible will be a great help as you write
your document and when you change it.  This sample
file shows how you can add comments to your own input
file.

Because printing is different from typewriting,
there are a number of things that you have to do
differently when preparing an input file than if
you were just typing the document directly.
Quotation marks like
       ``this''
have to be handled specially, as do quotes within
quotes:
       ``\,`this'            % \, separates the double and single quote.
        is what I just
        wrote, not  `that'\,''.

Dashes come in three sizes: an
       intra-word
dash, a medium dash for number ranges like
       1--2,
and a punctuation
       dash---like
this.

A sentence-ending space should be larger than the
space between words within a sentence.  You
sometimes have to type special commands in
conjunction with punctuation characters to get
this right, as in the following sentence.
       Gnats, gnus, etc.\ all  % `\ ' makes an inter-word space.
       begin with G\@.         % \@ marks end-of-sentence punctuation.
You should check the spaces after periods when
reading your output to make sure you haven't
forgotten any special cases.  Generating an
ellipsis
       \ldots\               % `\ ' is needed after `\ldots' because TeX
                             % ignores spaces after command names like \ldots
                             % made from \ + letters.
                             %
                             % Note how a `%' character causes TeX to ignore
                             % the end of the input line, so these blank lines
                             % do not start a new paragraph.
                             %
with the right spacing around the periods requires
a special command.

\LaTeX\ interprets some common characters as
commands, so you must type special commands to
generate them.  These characters include the
following:
       \$ \& \% \# \{ and \}.

In printing, text is usually emphasized with an
       \emph{italic}
type style.

\begin{em}
   A long segment of text can also be emphasized
   in this way.  Text within such a segment can be
   given \emph{additional} emphasis.
\end{em}

It is sometimes necessary to prevent \LaTeX\ from
breaking a line where it might otherwise do so.
This may be at a space, as between the ``Mr.''\ and
``Jones'' in
       ``Mr.~Jones'',        % ~ produces an unbreakable interword space.
or within a word---especially when the word is a
symbol like
       \mbox{\emph{itemnum}}
that makes little sense when hyphenated across
lines.

Footnotes\footnote{This is an example of a footnote.}
pose no problem.

\LaTeX\ is good at typesetting mathematical formulas
like
       \( x-3y + z = 7 \)
or
       \( a_{1} > x^{2n} + y^{2n} > x' \)
or
       \( \ip{A}{B} = \sum_{i} a_{i} b_{i} \).
The spaces you type in a formula are
ignored.  Remember that a letter like
       $x$                   % $ ... $  and  \( ... \)  are equivalent
is a formula when it denotes a mathematical
symbol, and it should be typed as one.

\section{Displayed Text}

Text is displayed by indenting it from the left
margin.  Quotations are commonly displayed.  There
are short quotations
\begin{quote}
   This is a short quotation.  It consists of a
   single paragraph of text.  See how it is formatted.
\end{quote}
and longer ones.
\begin{quotation}
   This is a longer quotation.  It consists of two
   paragraphs of text, neither of which are
   particularly interesting.

   This is the second paragraph of the quotation.  It
   is just as dull as the first paragraph.
\end{quotation}
Another frequently-displayed structure is a list.
The following is an example of an \emph{itemized}
list.
\begin{itemize}
   \item This is the first item of an itemized list.
         Each item in the list is marked with a ``tick''.
         You don't have to worry about what kind of tick
         mark is used.

   \item This is the second item of the list.  It
         contains another list nested inside it.  The inner
         list is an \emph{enumerated} list.
         \begin{enumerate}
            \item This is the first item of an enumerated
                  list that is nested within the itemized list.

            \item This is the second item of the inner list.
                  \LaTeX\ allows you to nest lists deeper than
                  you really should.
         \end{enumerate}
         This is the rest of the second item of the outer
         list.  It is no more interesting than any other
         part of the item.
   \item This is the third item of the list.
\end{itemize}
You can even display poetry.
\begin{verse}
   There is an environment
    for verse \\             % The \\ command separates lines
   Whose features some poets % within a stanza.
   will curse.

                             % One or more blank lines separate stanzas.

   For instead of making\\
   Them do \emph{all} line breaking, \\
   It allows them to put too many words on a line when they'd rather be
   forced to be terse.
\end{verse}

Mathematical formulas may also be displayed.  A
displayed formula
is
one-line long; multiline
formulas require special formatting instructions.
   \[  \ip{\Gamma}{\psi'} = x'' + y^{2} + z_{i}^{n}\]
Don't start a paragraph with a displayed equation,
nor make one a paragraph by itself.

\end{document}               % End of document.

\section{Bench Huge Section 1261}
% synthetic bench section 1261
% This is a sample LaTeX input file.  (Version of 12 August 2004.)
%
% A '%' character causes TeX to ignore all remaining text on the line,
% and is used for comments like this one.

\documentclass{article}      % Specifies the document class

                             % The preamble begins here.
\title{An Example Document}  % Declares the document's title.
\author{Leslie Lamport}      % Declares the author's name.
\date{January 21, 1994}      % Deleting this command produces today's date.

\newcommand{\ip}[2]{(#1, #2)}
                             % Defines \ip{arg1}{arg2} to mean
                             % (arg1, arg2).

%\newcommand{\ip}[2]{\langle #1 | #2\rangle}
                             % This is an alternative definition of
                             % \ip that is commented out.

\begin{document}             % End of preamble and beginning of text.

\maketitle                   % Produces the title.

This is an example input file.  Comparing it with
the output it generates can show you how to
produce a simple document of your own.

\section{Ordinary Text}      % Produces section heading.  Lower-level
                             % sections are begun with similar
                             % \subsection and \subsubsection commands.

The ends  of words and sentences are marked
  by   spaces. It  doesn't matter how many
spaces    you type; one is as good as 100.  The
end of   a line counts as a space.

One   or more   blank lines denote the  end
of  a paragraph.

Since any number of consecutive spaces are treated
like a single one, the formatting of the input
file makes no difference to
      \LaTeX,                % The \LaTeX command generates the LaTeX logo.
but it makes a difference to you.  When you use
\LaTeX, making your input file as easy to read
as possible will be a great help as you write
your document and when you change it.  This sample
file shows how you can add comments to your own input
file.

Because printing is different from typewriting,
there are a number of things that you have to do
differently when preparing an input file than if
you were just typing the document directly.
Quotation marks like
       ``this''
have to be handled specially, as do quotes within
quotes:
       ``\,`this'            % \, separates the double and single quote.
        is what I just
        wrote, not  `that'\,''.

Dashes come in three sizes: an
       intra-word
dash, a medium dash for number ranges like
       1--2,
and a punctuation
       dash---like
this.

A sentence-ending space should be larger than the
space between words within a sentence.  You
sometimes have to type special commands in
conjunction with punctuation characters to get
this right, as in the following sentence.
       Gnats, gnus, etc.\ all  % `\ ' makes an inter-word space.
       begin with G\@.         % \@ marks end-of-sentence punctuation.
You should check the spaces after periods when
reading your output to make sure you haven't
forgotten any special cases.  Generating an
ellipsis
       \ldots\               % `\ ' is needed after `\ldots' because TeX
                             % ignores spaces after command names like \ldots
                             % made from \ + letters.
                             %
                             % Note how a `%' character causes TeX to ignore
                             % the end of the input line, so these blank lines
                             % do not start a new paragraph.
                             %
with the right spacing around the periods requires
a special command.

\LaTeX\ interprets some common characters as
commands, so you must type special commands to
generate them.  These characters include the
following:
       \$ \& \% \# \{ and \}.

In printing, text is usually emphasized with an
       \emph{italic}
type style.

\begin{em}
   A long segment of text can also be emphasized
   in this way.  Text within such a segment can be
   given \emph{additional} emphasis.
\end{em}

It is sometimes necessary to prevent \LaTeX\ from
breaking a line where it might otherwise do so.
This may be at a space, as between the ``Mr.''\ and
``Jones'' in
       ``Mr.~Jones'',        % ~ produces an unbreakable interword space.
or within a word---especially when the word is a
symbol like
       \mbox{\emph{itemnum}}
that makes little sense when hyphenated across
lines.

Footnotes\footnote{This is an example of a footnote.}
pose no problem.

\LaTeX\ is good at typesetting mathematical formulas
like
       \( x-3y + z = 7 \)
or
       \( a_{1} > x^{2n} + y^{2n} > x' \)
or
       \( \ip{A}{B} = \sum_{i} a_{i} b_{i} \).
The spaces you type in a formula are
ignored.  Remember that a letter like
       $x$                   % $ ... $  and  \( ... \)  are equivalent
is a formula when it denotes a mathematical
symbol, and it should be typed as one.

\section{Displayed Text}

Text is displayed by indenting it from the left
margin.  Quotations are commonly displayed.  There
are short quotations
\begin{quote}
   This is a short quotation.  It consists of a
   single paragraph of text.  See how it is formatted.
\end{quote}
and longer ones.
\begin{quotation}
   This is a longer quotation.  It consists of two
   paragraphs of text, neither of which are
   particularly interesting.

   This is the second paragraph of the quotation.  It
   is just as dull as the first paragraph.
\end{quotation}
Another frequently-displayed structure is a list.
The following is an example of an \emph{itemized}
list.
\begin{itemize}
   \item This is the first item of an itemized list.
         Each item in the list is marked with a ``tick''.
         You don't have to worry about what kind of tick
         mark is used.

   \item This is the second item of the list.  It
         contains another list nested inside it.  The inner
         list is an \emph{enumerated} list.
         \begin{enumerate}
            \item This is the first item of an enumerated
                  list that is nested within the itemized list.

            \item This is the second item of the inner list.
                  \LaTeX\ allows you to nest lists deeper than
                  you really should.
         \end{enumerate}
         This is the rest of the second item of the outer
         list.  It is no more interesting than any other
         part of the item.
   \item This is the third item of the list.
\end{itemize}
You can even display poetry.
\begin{verse}
   There is an environment
    for verse \\             % The \\ command separates lines
   Whose features some poets % within a stanza.
   will curse.

                             % One or more blank lines separate stanzas.

   For instead of making\\
   Them do \emph{all} line breaking, \\
   It allows them to put too many words on a line when they'd rather be
   forced to be terse.
\end{verse}

Mathematical formulas may also be displayed.  A
displayed formula
is
one-line long; multiline
formulas require special formatting instructions.
   \[  \ip{\Gamma}{\psi'} = x'' + y^{2} + z_{i}^{n}\]
Don't start a paragraph with a displayed equation,
nor make one a paragraph by itself.

\end{document}               % End of document.

\section{Bench Huge Section 1262}
% synthetic bench section 1262
% This is a sample LaTeX input file.  (Version of 12 August 2004.)
%
% A '%' character causes TeX to ignore all remaining text on the line,
% and is used for comments like this one.

\documentclass{article}      % Specifies the document class

                             % The preamble begins here.
\title{An Example Document}  % Declares the document's title.
\author{Leslie Lamport}      % Declares the author's name.
\date{January 21, 1994}      % Deleting this command produces today's date.

\newcommand{\ip}[2]{(#1, #2)}
                             % Defines \ip{arg1}{arg2} to mean
                             % (arg1, arg2).

%\newcommand{\ip}[2]{\langle #1 | #2\rangle}
                             % This is an alternative definition of
                             % \ip that is commented out.

\begin{document}             % End of preamble and beginning of text.

\maketitle                   % Produces the title.

This is an example input file.  Comparing it with
the output it generates can show you how to
produce a simple document of your own.

\section{Ordinary Text}      % Produces section heading.  Lower-level
                             % sections are begun with similar
                             % \subsection and \subsubsection commands.

The ends  of words and sentences are marked
  by   spaces. It  doesn't matter how many
spaces    you type; one is as good as 100.  The
end of   a line counts as a space.

One   or more   blank lines denote the  end
of  a paragraph.

Since any number of consecutive spaces are treated
like a single one, the formatting of the input
file makes no difference to
      \LaTeX,                % The \LaTeX command generates the LaTeX logo.
but it makes a difference to you.  When you use
\LaTeX, making your input file as easy to read
as possible will be a great help as you write
your document and when you change it.  This sample
file shows how you can add comments to your own input
file.

Because printing is different from typewriting,
there are a number of things that you have to do
differently when preparing an input file than if
you were just typing the document directly.
Quotation marks like
       ``this''
have to be handled specially, as do quotes within
quotes:
       ``\,`this'            % \, separates the double and single quote.
        is what I just
        wrote, not  `that'\,''.

Dashes come in three sizes: an
       intra-word
dash, a medium dash for number ranges like
       1--2,
and a punctuation
       dash---like
this.

A sentence-ending space should be larger than the
space between words within a sentence.  You
sometimes have to type special commands in
conjunction with punctuation characters to get
this right, as in the following sentence.
       Gnats, gnus, etc.\ all  % `\ ' makes an inter-word space.
       begin with G\@.         % \@ marks end-of-sentence punctuation.
You should check the spaces after periods when
reading your output to make sure you haven't
forgotten any special cases.  Generating an
ellipsis
       \ldots\               % `\ ' is needed after `\ldots' because TeX
                             % ignores spaces after command names like \ldots
                             % made from \ + letters.
                             %
                             % Note how a `%' character causes TeX to ignore
                             % the end of the input line, so these blank lines
                             % do not start a new paragraph.
                             %
with the right spacing around the periods requires
a special command.

\LaTeX\ interprets some common characters as
commands, so you must type special commands to
generate them.  These characters include the
following:
       \$ \& \% \# \{ and \}.

In printing, text is usually emphasized with an
       \emph{italic}
type style.

\begin{em}
   A long segment of text can also be emphasized
   in this way.  Text within such a segment can be
   given \emph{additional} emphasis.
\end{em}

It is sometimes necessary to prevent \LaTeX\ from
breaking a line where it might otherwise do so.
This may be at a space, as between the ``Mr.''\ and
``Jones'' in
       ``Mr.~Jones'',        % ~ produces an unbreakable interword space.
or within a word---especially when the word is a
symbol like
       \mbox{\emph{itemnum}}
that makes little sense when hyphenated across
lines.

Footnotes\footnote{This is an example of a footnote.}
pose no problem.

\LaTeX\ is good at typesetting mathematical formulas
like
       \( x-3y + z = 7 \)
or
       \( a_{1} > x^{2n} + y^{2n} > x' \)
or
       \( \ip{A}{B} = \sum_{i} a_{i} b_{i} \).
The spaces you type in a formula are
ignored.  Remember that a letter like
       $x$                   % $ ... $  and  \( ... \)  are equivalent
is a formula when it denotes a mathematical
symbol, and it should be typed as one.

\section{Displayed Text}

Text is displayed by indenting it from the left
margin.  Quotations are commonly displayed.  There
are short quotations
\begin{quote}
   This is a short quotation.  It consists of a
   single paragraph of text.  See how it is formatted.
\end{quote}
and longer ones.
\begin{quotation}
   This is a longer quotation.  It consists of two
   paragraphs of text, neither of which are
   particularly interesting.

   This is the second paragraph of the quotation.  It
   is just as dull as the first paragraph.
\end{quotation}
Another frequently-displayed structure is a list.
The following is an example of an \emph{itemized}
list.
\begin{itemize}
   \item This is the first item of an itemized list.
         Each item in the list is marked with a ``tick''.
         You don't have to worry about what kind of tick
         mark is used.

   \item This is the second item of the list.  It
         contains another list nested inside it.  The inner
         list is an \emph{enumerated} list.
         \begin{enumerate}
            \item This is the first item of an enumerated
                  list that is nested within the itemized list.

            \item This is the second item of the inner list.
                  \LaTeX\ allows you to nest lists deeper than
                  you really should.
         \end{enumerate}
         This is the rest of the second item of the outer
         list.  It is no more interesting than any other
         part of the item.
   \item This is the third item of the list.
\end{itemize}
You can even display poetry.
\begin{verse}
   There is an environment
    for verse \\             % The \\ command separates lines
   Whose features some poets % within a stanza.
   will curse.

                             % One or more blank lines separate stanzas.

   For instead of making\\
   Them do \emph{all} line breaking, \\
   It allows them to put too many words on a line when they'd rather be
   forced to be terse.
\end{verse}

Mathematical formulas may also be displayed.  A
displayed formula
is
one-line long; multiline
formulas require special formatting instructions.
   \[  \ip{\Gamma}{\psi'} = x'' + y^{2} + z_{i}^{n}\]
Don't start a paragraph with a displayed equation,
nor make one a paragraph by itself.

\end{document}               % End of document.

\section{Bench Huge Section 1263}
% synthetic bench section 1263
% This is a sample LaTeX input file.  (Version of 12 August 2004.)
%
% A '%' character causes TeX to ignore all remaining text on the line,
% and is used for comments like this one.

\documentclass{article}      % Specifies the document class

                             % The preamble begins here.
\title{An Example Document}  % Declares the document's title.
\author{Leslie Lamport}      % Declares the author's name.
\date{January 21, 1994}      % Deleting this command produces today's date.

\newcommand{\ip}[2]{(#1, #2)}
                             % Defines \ip{arg1}{arg2} to mean
                             % (arg1, arg2).

%\newcommand{\ip}[2]{\langle #1 | #2\rangle}
                             % This is an alternative definition of
                             % \ip that is commented out.

\begin{document}             % End of preamble and beginning of text.

\maketitle                   % Produces the title.

This is an example input file.  Comparing it with
the output it generates can show you how to
produce a simple document of your own.

\section{Ordinary Text}      % Produces section heading.  Lower-level
                             % sections are begun with similar
                             % \subsection and \subsubsection commands.

The ends  of words and sentences are marked
  by   spaces. It  doesn't matter how many
spaces    you type; one is as good as 100.  The
end of   a line counts as a space.

One   or more   blank lines denote the  end
of  a paragraph.

Since any number of consecutive spaces are treated
like a single one, the formatting of the input
file makes no difference to
      \LaTeX,                % The \LaTeX command generates the LaTeX logo.
but it makes a difference to you.  When you use
\LaTeX, making your input file as easy to read
as possible will be a great help as you write
your document and when you change it.  This sample
file shows how you can add comments to your own input
file.

Because printing is different from typewriting,
there are a number of things that you have to do
differently when preparing an input file than if
you were just typing the document directly.
Quotation marks like
       ``this''
have to be handled specially, as do quotes within
quotes:
       ``\,`this'            % \, separates the double and single quote.
        is what I just
        wrote, not  `that'\,''.

Dashes come in three sizes: an
       intra-word
dash, a medium dash for number ranges like
       1--2,
and a punctuation
       dash---like
this.

A sentence-ending space should be larger than the
space between words within a sentence.  You
sometimes have to type special commands in
conjunction with punctuation characters to get
this right, as in the following sentence.
       Gnats, gnus, etc.\ all  % `\ ' makes an inter-word space.
       begin with G\@.         % \@ marks end-of-sentence punctuation.
You should check the spaces after periods when
reading your output to make sure you haven't
forgotten any special cases.  Generating an
ellipsis
       \ldots\               % `\ ' is needed after `\ldots' because TeX
                             % ignores spaces after command names like \ldots
                             % made from \ + letters.
                             %
                             % Note how a `%' character causes TeX to ignore
                             % the end of the input line, so these blank lines
                             % do not start a new paragraph.
                             %
with the right spacing around the periods requires
a special command.

\LaTeX\ interprets some common characters as
commands, so you must type special commands to
generate them.  These characters include the
following:
       \$ \& \% \# \{ and \}.

In printing, text is usually emphasized with an
       \emph{italic}
type style.

\begin{em}
   A long segment of text can also be emphasized
   in this way.  Text within such a segment can be
   given \emph{additional} emphasis.
\end{em}

It is sometimes necessary to prevent \LaTeX\ from
breaking a line where it might otherwise do so.
This may be at a space, as between the ``Mr.''\ and
``Jones'' in
       ``Mr.~Jones'',        % ~ produces an unbreakable interword space.
or within a word---especially when the word is a
symbol like
       \mbox{\emph{itemnum}}
that makes little sense when hyphenated across
lines.

Footnotes\footnote{This is an example of a footnote.}
pose no problem.

\LaTeX\ is good at typesetting mathematical formulas
like
       \( x-3y + z = 7 \)
or
       \( a_{1} > x^{2n} + y^{2n} > x' \)
or
       \( \ip{A}{B} = \sum_{i} a_{i} b_{i} \).
The spaces you type in a formula are
ignored.  Remember that a letter like
       $x$                   % $ ... $  and  \( ... \)  are equivalent
is a formula when it denotes a mathematical
symbol, and it should be typed as one.

\section{Displayed Text}

Text is displayed by indenting it from the left
margin.  Quotations are commonly displayed.  There
are short quotations
\begin{quote}
   This is a short quotation.  It consists of a
   single paragraph of text.  See how it is formatted.
\end{quote}
and longer ones.
\begin{quotation}
   This is a longer quotation.  It consists of two
   paragraphs of text, neither of which are
   particularly interesting.

   This is the second paragraph of the quotation.  It
   is just as dull as the first paragraph.
\end{quotation}
Another frequently-displayed structure is a list.
The following is an example of an \emph{itemized}
list.
\begin{itemize}
   \item This is the first item of an itemized list.
         Each item in the list is marked with a ``tick''.
         You don't have to worry about what kind of tick
         mark is used.

   \item This is the second item of the list.  It
         contains another list nested inside it.  The inner
         list is an \emph{enumerated} list.
         \begin{enumerate}
            \item This is the first item of an enumerated
                  list that is nested within the itemized list.

            \item This is the second item of the inner list.
                  \LaTeX\ allows you to nest lists deeper than
                  you really should.
         \end{enumerate}
         This is the rest of the second item of the outer
         list.  It is no more interesting than any other
         part of the item.
   \item This is the third item of the list.
\end{itemize}
You can even display poetry.
\begin{verse}
   There is an environment
    for verse \\             % The \\ command separates lines
   Whose features some poets % within a stanza.
   will curse.

                             % One or more blank lines separate stanzas.

   For instead of making\\
   Them do \emph{all} line breaking, \\
   It allows them to put too many words on a line when they'd rather be
   forced to be terse.
\end{verse}

Mathematical formulas may also be displayed.  A
displayed formula
is
one-line long; multiline
formulas require special formatting instructions.
   \[  \ip{\Gamma}{\psi'} = x'' + y^{2} + z_{i}^{n}\]
Don't start a paragraph with a displayed equation,
nor make one a paragraph by itself.

\end{document}               % End of document.

\section{Bench Huge Section 1264}
% synthetic bench section 1264
% This is a sample LaTeX input file.  (Version of 12 August 2004.)
%
% A '%' character causes TeX to ignore all remaining text on the line,
% and is used for comments like this one.

\documentclass{article}      % Specifies the document class

                             % The preamble begins here.
\title{An Example Document}  % Declares the document's title.
\author{Leslie Lamport}      % Declares the author's name.
\date{January 21, 1994}      % Deleting this command produces today's date.

\newcommand{\ip}[2]{(#1, #2)}
                             % Defines \ip{arg1}{arg2} to mean
                             % (arg1, arg2).

%\newcommand{\ip}[2]{\langle #1 | #2\rangle}
                             % This is an alternative definition of
                             % \ip that is commented out.

\begin{document}             % End of preamble and beginning of text.

\maketitle                   % Produces the title.

This is an example input file.  Comparing it with
the output it generates can show you how to
produce a simple document of your own.

\section{Ordinary Text}      % Produces section heading.  Lower-level
                             % sections are begun with similar
                             % \subsection and \subsubsection commands.

The ends  of words and sentences are marked
  by   spaces. It  doesn't matter how many
spaces    you type; one is as good as 100.  The
end of   a line counts as a space.

One   or more   blank lines denote the  end
of  a paragraph.

Since any number of consecutive spaces are treated
like a single one, the formatting of the input
file makes no difference to
      \LaTeX,                % The \LaTeX command generates the LaTeX logo.
but it makes a difference to you.  When you use
\LaTeX, making your input file as easy to read
as possible will be a great help as you write
your document and when you change it.  This sample
file shows how you can add comments to your own input
file.

Because printing is different from typewriting,
there are a number of things that you have to do
differently when preparing an input file than if
you were just typing the document directly.
Quotation marks like
       ``this''
have to be handled specially, as do quotes within
quotes:
       ``\,`this'            % \, separates the double and single quote.
        is what I just
        wrote, not  `that'\,''.

Dashes come in three sizes: an
       intra-word
dash, a medium dash for number ranges like
       1--2,
and a punctuation
       dash---like
this.

A sentence-ending space should be larger than the
space between words within a sentence.  You
sometimes have to type special commands in
conjunction with punctuation characters to get
this right, as in the following sentence.
       Gnats, gnus, etc.\ all  % `\ ' makes an inter-word space.
       begin with G\@.         % \@ marks end-of-sentence punctuation.
You should check the spaces after periods when
reading your output to make sure you haven't
forgotten any special cases.  Generating an
ellipsis
       \ldots\               % `\ ' is needed after `\ldots' because TeX
                             % ignores spaces after command names like \ldots
                             % made from \ + letters.
                             %
                             % Note how a `%' character causes TeX to ignore
                             % the end of the input line, so these blank lines
                             % do not start a new paragraph.
                             %
with the right spacing around the periods requires
a special command.

\LaTeX\ interprets some common characters as
commands, so you must type special commands to
generate them.  These characters include the
following:
       \$ \& \% \# \{ and \}.

In printing, text is usually emphasized with an
       \emph{italic}
type style.

\begin{em}
   A long segment of text can also be emphasized
   in this way.  Text within such a segment can be
   given \emph{additional} emphasis.
\end{em}

It is sometimes necessary to prevent \LaTeX\ from
breaking a line where it might otherwise do so.
This may be at a space, as between the ``Mr.''\ and
``Jones'' in
       ``Mr.~Jones'',        % ~ produces an unbreakable interword space.
or within a word---especially when the word is a
symbol like
       \mbox{\emph{itemnum}}
that makes little sense when hyphenated across
lines.

Footnotes\footnote{This is an example of a footnote.}
pose no problem.

\LaTeX\ is good at typesetting mathematical formulas
like
       \( x-3y + z = 7 \)
or
       \( a_{1} > x^{2n} + y^{2n} > x' \)
or
       \( \ip{A}{B} = \sum_{i} a_{i} b_{i} \).
The spaces you type in a formula are
ignored.  Remember that a letter like
       $x$                   % $ ... $  and  \( ... \)  are equivalent
is a formula when it denotes a mathematical
symbol, and it should be typed as one.

\section{Displayed Text}

Text is displayed by indenting it from the left
margin.  Quotations are commonly displayed.  There
are short quotations
\begin{quote}
   This is a short quotation.  It consists of a
   single paragraph of text.  See how it is formatted.
\end{quote}
and longer ones.
\begin{quotation}
   This is a longer quotation.  It consists of two
   paragraphs of text, neither of which are
   particularly interesting.

   This is the second paragraph of the quotation.  It
   is just as dull as the first paragraph.
\end{quotation}
Another frequently-displayed structure is a list.
The following is an example of an \emph{itemized}
list.
\begin{itemize}
   \item This is the first item of an itemized list.
         Each item in the list is marked with a ``tick''.
         You don't have to worry about what kind of tick
         mark is used.

   \item This is the second item of the list.  It
         contains another list nested inside it.  The inner
         list is an \emph{enumerated} list.
         \begin{enumerate}
            \item This is the first item of an enumerated
                  list that is nested within the itemized list.

            \item This is the second item of the inner list.
                  \LaTeX\ allows you to nest lists deeper than
                  you really should.
         \end{enumerate}
         This is the rest of the second item of the outer
         list.  It is no more interesting than any other
         part of the item.
   \item This is the third item of the list.
\end{itemize}
You can even display poetry.
\begin{verse}
   There is an environment
    for verse \\             % The \\ command separates lines
   Whose features some poets % within a stanza.
   will curse.

                             % One or more blank lines separate stanzas.

   For instead of making\\
   Them do \emph{all} line breaking, \\
   It allows them to put too many words on a line when they'd rather be
   forced to be terse.
\end{verse}

Mathematical formulas may also be displayed.  A
displayed formula
is
one-line long; multiline
formulas require special formatting instructions.
   \[  \ip{\Gamma}{\psi'} = x'' + y^{2} + z_{i}^{n}\]
Don't start a paragraph with a displayed equation,
nor make one a paragraph by itself.

\end{document}               % End of document.

\section{Bench Huge Section 1265}
% synthetic bench section 1265
% This is a sample LaTeX input file.  (Version of 12 August 2004.)
%
% A '%' character causes TeX to ignore all remaining text on the line,
% and is used for comments like this one.

\documentclass{article}      % Specifies the document class

                             % The preamble begins here.
\title{An Example Document}  % Declares the document's title.
\author{Leslie Lamport}      % Declares the author's name.
\date{January 21, 1994}      % Deleting this command produces today's date.

\newcommand{\ip}[2]{(#1, #2)}
                             % Defines \ip{arg1}{arg2} to mean
                             % (arg1, arg2).

%\newcommand{\ip}[2]{\langle #1 | #2\rangle}
                             % This is an alternative definition of
                             % \ip that is commented out.

\begin{document}             % End of preamble and beginning of text.

\maketitle                   % Produces the title.

This is an example input file.  Comparing it with
the output it generates can show you how to
produce a simple document of your own.

\section{Ordinary Text}      % Produces section heading.  Lower-level
                             % sections are begun with similar
                             % \subsection and \subsubsection commands.

The ends  of words and sentences are marked
  by   spaces. It  doesn't matter how many
spaces    you type; one is as good as 100.  The
end of   a line counts as a space.

One   or more   blank lines denote the  end
of  a paragraph.

Since any number of consecutive spaces are treated
like a single one, the formatting of the input
file makes no difference to
      \LaTeX,                % The \LaTeX command generates the LaTeX logo.
but it makes a difference to you.  When you use
\LaTeX, making your input file as easy to read
as possible will be a great help as you write
your document and when you change it.  This sample
file shows how you can add comments to your own input
file.

Because printing is different from typewriting,
there are a number of things that you have to do
differently when preparing an input file than if
you were just typing the document directly.
Quotation marks like
       ``this''
have to be handled specially, as do quotes within
quotes:
       ``\,`this'            % \, separates the double and single quote.
        is what I just
        wrote, not  `that'\,''.

Dashes come in three sizes: an
       intra-word
dash, a medium dash for number ranges like
       1--2,
and a punctuation
       dash---like
this.

A sentence-ending space should be larger than the
space between words within a sentence.  You
sometimes have to type special commands in
conjunction with punctuation characters to get
this right, as in the following sentence.
       Gnats, gnus, etc.\ all  % `\ ' makes an inter-word space.
       begin with G\@.         % \@ marks end-of-sentence punctuation.
You should check the spaces after periods when
reading your output to make sure you haven't
forgotten any special cases.  Generating an
ellipsis
       \ldots\               % `\ ' is needed after `\ldots' because TeX
                             % ignores spaces after command names like \ldots
                             % made from \ + letters.
                             %
                             % Note how a `%' character causes TeX to ignore
                             % the end of the input line, so these blank lines
                             % do not start a new paragraph.
                             %
with the right spacing around the periods requires
a special command.

\LaTeX\ interprets some common characters as
commands, so you must type special commands to
generate them.  These characters include the
following:
       \$ \& \% \# \{ and \}.

In printing, text is usually emphasized with an
       \emph{italic}
type style.

\begin{em}
   A long segment of text can also be emphasized
   in this way.  Text within such a segment can be
   given \emph{additional} emphasis.
\end{em}

It is sometimes necessary to prevent \LaTeX\ from
breaking a line where it might otherwise do so.
This may be at a space, as between the ``Mr.''\ and
``Jones'' in
       ``Mr.~Jones'',        % ~ produces an unbreakable interword space.
or within a word---especially when the word is a
symbol like
       \mbox{\emph{itemnum}}
that makes little sense when hyphenated across
lines.

Footnotes\footnote{This is an example of a footnote.}
pose no problem.

\LaTeX\ is good at typesetting mathematical formulas
like
       \( x-3y + z = 7 \)
or
       \( a_{1} > x^{2n} + y^{2n} > x' \)
or
       \( \ip{A}{B} = \sum_{i} a_{i} b_{i} \).
The spaces you type in a formula are
ignored.  Remember that a letter like
       $x$                   % $ ... $  and  \( ... \)  are equivalent
is a formula when it denotes a mathematical
symbol, and it should be typed as one.

\section{Displayed Text}

Text is displayed by indenting it from the left
margin.  Quotations are commonly displayed.  There
are short quotations
\begin{quote}
   This is a short quotation.  It consists of a
   single paragraph of text.  See how it is formatted.
\end{quote}
and longer ones.
\begin{quotation}
   This is a longer quotation.  It consists of two
   paragraphs of text, neither of which are
   particularly interesting.

   This is the second paragraph of the quotation.  It
   is just as dull as the first paragraph.
\end{quotation}
Another frequently-displayed structure is a list.
The following is an example of an \emph{itemized}
list.
\begin{itemize}
   \item This is the first item of an itemized list.
         Each item in the list is marked with a ``tick''.
         You don't have to worry about what kind of tick
         mark is used.

   \item This is the second item of the list.  It
         contains another list nested inside it.  The inner
         list is an \emph{enumerated} list.
         \begin{enumerate}
            \item This is the first item of an enumerated
                  list that is nested within the itemized list.

            \item This is the second item of the inner list.
                  \LaTeX\ allows you to nest lists deeper than
                  you really should.
         \end{enumerate}
         This is the rest of the second item of the outer
         list.  It is no more interesting than any other
         part of the item.
   \item This is the third item of the list.
\end{itemize}
You can even display poetry.
\begin{verse}
   There is an environment
    for verse \\             % The \\ command separates lines
   Whose features some poets % within a stanza.
   will curse.

                             % One or more blank lines separate stanzas.

   For instead of making\\
   Them do \emph{all} line breaking, \\
   It allows them to put too many words on a line when they'd rather be
   forced to be terse.
\end{verse}

Mathematical formulas may also be displayed.  A
displayed formula
is
one-line long; multiline
formulas require special formatting instructions.
   \[  \ip{\Gamma}{\psi'} = x'' + y^{2} + z_{i}^{n}\]
Don't start a paragraph with a displayed equation,
nor make one a paragraph by itself.

\end{document}               % End of document.

\section{Bench Huge Section 1266}
% synthetic bench section 1266
% This is a sample LaTeX input file.  (Version of 12 August 2004.)
%
% A '%' character causes TeX to ignore all remaining text on the line,
% and is used for comments like this one.

\documentclass{article}      % Specifies the document class

                             % The preamble begins here.
\title{An Example Document}  % Declares the document's title.
\author{Leslie Lamport}      % Declares the author's name.
\date{January 21, 1994}      % Deleting this command produces today's date.

\newcommand{\ip}[2]{(#1, #2)}
                             % Defines \ip{arg1}{arg2} to mean
                             % (arg1, arg2).

%\newcommand{\ip}[2]{\langle #1 | #2\rangle}
                             % This is an alternative definition of
                             % \ip that is commented out.

\begin{document}             % End of preamble and beginning of text.

\maketitle                   % Produces the title.

This is an example input file.  Comparing it with
the output it generates can show you how to
produce a simple document of your own.

\section{Ordinary Text}      % Produces section heading.  Lower-level
                             % sections are begun with similar
                             % \subsection and \subsubsection commands.

The ends  of words and sentences are marked
  by   spaces. It  doesn't matter how many
spaces    you type; one is as good as 100.  The
end of   a line counts as a space.

One   or more   blank lines denote the  end
of  a paragraph.

Since any number of consecutive spaces are treated
like a single one, the formatting of the input
file makes no difference to
      \LaTeX,                % The \LaTeX command generates the LaTeX logo.
but it makes a difference to you.  When you use
\LaTeX, making your input file as easy to read
as possible will be a great help as you write
your document and when you change it.  This sample
file shows how you can add comments to your own input
file.

Because printing is different from typewriting,
there are a number of things that you have to do
differently when preparing an input file than if
you were just typing the document directly.
Quotation marks like
       ``this''
have to be handled specially, as do quotes within
quotes:
       ``\,`this'            % \, separates the double and single quote.
        is what I just
        wrote, not  `that'\,''.

Dashes come in three sizes: an
       intra-word
dash, a medium dash for number ranges like
       1--2,
and a punctuation
       dash---like
this.

A sentence-ending space should be larger than the
space between words within a sentence.  You
sometimes have to type special commands in
conjunction with punctuation characters to get
this right, as in the following sentence.
       Gnats, gnus, etc.\ all  % `\ ' makes an inter-word space.
       begin with G\@.         % \@ marks end-of-sentence punctuation.
You should check the spaces after periods when
reading your output to make sure you haven't
forgotten any special cases.  Generating an
ellipsis
       \ldots\               % `\ ' is needed after `\ldots' because TeX
                             % ignores spaces after command names like \ldots
                             % made from \ + letters.
                             %
                             % Note how a `%' character causes TeX to ignore
                             % the end of the input line, so these blank lines
                             % do not start a new paragraph.
                             %
with the right spacing around the periods requires
a special command.

\LaTeX\ interprets some common characters as
commands, so you must type special commands to
generate them.  These characters include the
following:
       \$ \& \% \# \{ and \}.

In printing, text is usually emphasized with an
       \emph{italic}
type style.

\begin{em}
   A long segment of text can also be emphasized
   in this way.  Text within such a segment can be
   given \emph{additional} emphasis.
\end{em}

It is sometimes necessary to prevent \LaTeX\ from
breaking a line where it might otherwise do so.
This may be at a space, as between the ``Mr.''\ and
``Jones'' in
       ``Mr.~Jones'',        % ~ produces an unbreakable interword space.
or within a word---especially when the word is a
symbol like
       \mbox{\emph{itemnum}}
that makes little sense when hyphenated across
lines.

Footnotes\footnote{This is an example of a footnote.}
pose no problem.

\LaTeX\ is good at typesetting mathematical formulas
like
       \( x-3y + z = 7 \)
or
       \( a_{1} > x^{2n} + y^{2n} > x' \)
or
       \( \ip{A}{B} = \sum_{i} a_{i} b_{i} \).
The spaces you type in a formula are
ignored.  Remember that a letter like
       $x$                   % $ ... $  and  \( ... \)  are equivalent
is a formula when it denotes a mathematical
symbol, and it should be typed as one.

\section{Displayed Text}

Text is displayed by indenting it from the left
margin.  Quotations are commonly displayed.  There
are short quotations
\begin{quote}
   This is a short quotation.  It consists of a
   single paragraph of text.  See how it is formatted.
\end{quote}
and longer ones.
\begin{quotation}
   This is a longer quotation.  It consists of two
   paragraphs of text, neither of which are
   particularly interesting.

   This is the second paragraph of the quotation.  It
   is just as dull as the first paragraph.
\end{quotation}
Another frequently-displayed structure is a list.
The following is an example of an \emph{itemized}
list.
\begin{itemize}
   \item This is the first item of an itemized list.
         Each item in the list is marked with a ``tick''.
         You don't have to worry about what kind of tick
         mark is used.

   \item This is the second item of the list.  It
         contains another list nested inside it.  The inner
         list is an \emph{enumerated} list.
         \begin{enumerate}
            \item This is the first item of an enumerated
                  list that is nested within the itemized list.

            \item This is the second item of the inner list.
                  \LaTeX\ allows you to nest lists deeper than
                  you really should.
         \end{enumerate}
         This is the rest of the second item of the outer
         list.  It is no more interesting than any other
         part of the item.
   \item This is the third item of the list.
\end{itemize}
You can even display poetry.
\begin{verse}
   There is an environment
    for verse \\             % The \\ command separates lines
   Whose features some poets % within a stanza.
   will curse.

                             % One or more blank lines separate stanzas.

   For instead of making\\
   Them do \emph{all} line breaking, \\
   It allows them to put too many words on a line when they'd rather be
   forced to be terse.
\end{verse}

Mathematical formulas may also be displayed.  A
displayed formula
is
one-line long; multiline
formulas require special formatting instructions.
   \[  \ip{\Gamma}{\psi'} = x'' + y^{2} + z_{i}^{n}\]
Don't start a paragraph with a displayed equation,
nor make one a paragraph by itself.

\end{document}               % End of document.

\section{Bench Huge Section 1267}
% synthetic bench section 1267
% This is a sample LaTeX input file.  (Version of 12 August 2004.)
%
% A '%' character causes TeX to ignore all remaining text on the line,
% and is used for comments like this one.

\documentclass{article}      % Specifies the document class

                             % The preamble begins here.
\title{An Example Document}  % Declares the document's title.
\author{Leslie Lamport}      % Declares the author's name.
\date{January 21, 1994}      % Deleting this command produces today's date.

\newcommand{\ip}[2]{(#1, #2)}
                             % Defines \ip{arg1}{arg2} to mean
                             % (arg1, arg2).

%\newcommand{\ip}[2]{\langle #1 | #2\rangle}
                             % This is an alternative definition of
                             % \ip that is commented out.

\begin{document}             % End of preamble and beginning of text.

\maketitle                   % Produces the title.

This is an example input file.  Comparing it with
the output it generates can show you how to
produce a simple document of your own.

\section{Ordinary Text}      % Produces section heading.  Lower-level
                             % sections are begun with similar
                             % \subsection and \subsubsection commands.

The ends  of words and sentences are marked
  by   spaces. It  doesn't matter how many
spaces    you type; one is as good as 100.  The
end of   a line counts as a space.

One   or more   blank lines denote the  end
of  a paragraph.

Since any number of consecutive spaces are treated
like a single one, the formatting of the input
file makes no difference to
      \LaTeX,                % The \LaTeX command generates the LaTeX logo.
but it makes a difference to you.  When you use
\LaTeX, making your input file as easy to read
as possible will be a great help as you write
your document and when you change it.  This sample
file shows how you can add comments to your own input
file.

Because printing is different from typewriting,
there are a number of things that you have to do
differently when preparing an input file than if
you were just typing the document directly.
Quotation marks like
       ``this''
have to be handled specially, as do quotes within
quotes:
       ``\,`this'            % \, separates the double and single quote.
        is what I just
        wrote, not  `that'\,''.

Dashes come in three sizes: an
       intra-word
dash, a medium dash for number ranges like
       1--2,
and a punctuation
       dash---like
this.

A sentence-ending space should be larger than the
space between words within a sentence.  You
sometimes have to type special commands in
conjunction with punctuation characters to get
this right, as in the following sentence.
       Gnats, gnus, etc.\ all  % `\ ' makes an inter-word space.
       begin with G\@.         % \@ marks end-of-sentence punctuation.
You should check the spaces after periods when
reading your output to make sure you haven't
forgotten any special cases.  Generating an
ellipsis
       \ldots\               % `\ ' is needed after `\ldots' because TeX
                             % ignores spaces after command names like \ldots
                             % made from \ + letters.
                             %
                             % Note how a `%' character causes TeX to ignore
                             % the end of the input line, so these blank lines
                             % do not start a new paragraph.
                             %
with the right spacing around the periods requires
a special command.

\LaTeX\ interprets some common characters as
commands, so you must type special commands to
generate them.  These characters include the
following:
       \$ \& \% \# \{ and \}.

In printing, text is usually emphasized with an
       \emph{italic}
type style.

\begin{em}
   A long segment of text can also be emphasized
   in this way.  Text within such a segment can be
   given \emph{additional} emphasis.
\end{em}

It is sometimes necessary to prevent \LaTeX\ from
breaking a line where it might otherwise do so.
This may be at a space, as between the ``Mr.''\ and
``Jones'' in
       ``Mr.~Jones'',        % ~ produces an unbreakable interword space.
or within a word---especially when the word is a
symbol like
       \mbox{\emph{itemnum}}
that makes little sense when hyphenated across
lines.

Footnotes\footnote{This is an example of a footnote.}
pose no problem.

\LaTeX\ is good at typesetting mathematical formulas
like
       \( x-3y + z = 7 \)
or
       \( a_{1} > x^{2n} + y^{2n} > x' \)
or
       \( \ip{A}{B} = \sum_{i} a_{i} b_{i} \).
The spaces you type in a formula are
ignored.  Remember that a letter like
       $x$                   % $ ... $  and  \( ... \)  are equivalent
is a formula when it denotes a mathematical
symbol, and it should be typed as one.

\section{Displayed Text}

Text is displayed by indenting it from the left
margin.  Quotations are commonly displayed.  There
are short quotations
\begin{quote}
   This is a short quotation.  It consists of a
   single paragraph of text.  See how it is formatted.
\end{quote}
and longer ones.
\begin{quotation}
   This is a longer quotation.  It consists of two
   paragraphs of text, neither of which are
   particularly interesting.

   This is the second paragraph of the quotation.  It
   is just as dull as the first paragraph.
\end{quotation}
Another frequently-displayed structure is a list.
The following is an example of an \emph{itemized}
list.
\begin{itemize}
   \item This is the first item of an itemized list.
         Each item in the list is marked with a ``tick''.
         You don't have to worry about what kind of tick
         mark is used.

   \item This is the second item of the list.  It
         contains another list nested inside it.  The inner
         list is an \emph{enumerated} list.
         \begin{enumerate}
            \item This is the first item of an enumerated
                  list that is nested within the itemized list.

            \item This is the second item of the inner list.
                  \LaTeX\ allows you to nest lists deeper than
                  you really should.
         \end{enumerate}
         This is the rest of the second item of the outer
         list.  It is no more interesting than any other
         part of the item.
   \item This is the third item of the list.
\end{itemize}
You can even display poetry.
\begin{verse}
   There is an environment
    for verse \\             % The \\ command separates lines
   Whose features some poets % within a stanza.
   will curse.

                             % One or more blank lines separate stanzas.

   For instead of making\\
   Them do \emph{all} line breaking, \\
   It allows them to put too many words on a line when they'd rather be
   forced to be terse.
\end{verse}

Mathematical formulas may also be displayed.  A
displayed formula
is
one-line long; multiline
formulas require special formatting instructions.
   \[  \ip{\Gamma}{\psi'} = x'' + y^{2} + z_{i}^{n}\]
Don't start a paragraph with a displayed equation,
nor make one a paragraph by itself.

\end{document}               % End of document.

\section{Bench Huge Section 1268}
% synthetic bench section 1268
% This is a sample LaTeX input file.  (Version of 12 August 2004.)
%
% A '%' character causes TeX to ignore all remaining text on the line,
% and is used for comments like this one.

\documentclass{article}      % Specifies the document class

                             % The preamble begins here.
\title{An Example Document}  % Declares the document's title.
\author{Leslie Lamport}      % Declares the author's name.
\date{January 21, 1994}      % Deleting this command produces today's date.

\newcommand{\ip}[2]{(#1, #2)}
                             % Defines \ip{arg1}{arg2} to mean
                             % (arg1, arg2).

%\newcommand{\ip}[2]{\langle #1 | #2\rangle}
                             % This is an alternative definition of
                             % \ip that is commented out.

\begin{document}             % End of preamble and beginning of text.

\maketitle                   % Produces the title.

This is an example input file.  Comparing it with
the output it generates can show you how to
produce a simple document of your own.

\section{Ordinary Text}      % Produces section heading.  Lower-level
                             % sections are begun with similar
                             % \subsection and \subsubsection commands.

The ends  of words and sentences are marked
  by   spaces. It  doesn't matter how many
spaces    you type; one is as good as 100.  The
end of   a line counts as a space.

One   or more   blank lines denote the  end
of  a paragraph.

Since any number of consecutive spaces are treated
like a single one, the formatting of the input
file makes no difference to
      \LaTeX,                % The \LaTeX command generates the LaTeX logo.
but it makes a difference to you.  When you use
\LaTeX, making your input file as easy to read
as possible will be a great help as you write
your document and when you change it.  This sample
file shows how you can add comments to your own input
file.

Because printing is different from typewriting,
there are a number of things that you have to do
differently when preparing an input file than if
you were just typing the document directly.
Quotation marks like
       ``this''
have to be handled specially, as do quotes within
quotes:
       ``\,`this'            % \, separates the double and single quote.
        is what I just
        wrote, not  `that'\,''.

Dashes come in three sizes: an
       intra-word
dash, a medium dash for number ranges like
       1--2,
and a punctuation
       dash---like
this.

A sentence-ending space should be larger than the
space between words within a sentence.  You
sometimes have to type special commands in
conjunction with punctuation characters to get
this right, as in the following sentence.
       Gnats, gnus, etc.\ all  % `\ ' makes an inter-word space.
       begin with G\@.         % \@ marks end-of-sentence punctuation.
You should check the spaces after periods when
reading your output to make sure you haven't
forgotten any special cases.  Generating an
ellipsis
       \ldots\               % `\ ' is needed after `\ldots' because TeX
                             % ignores spaces after command names like \ldots
                             % made from \ + letters.
                             %
                             % Note how a `%' character causes TeX to ignore
                             % the end of the input line, so these blank lines
                             % do not start a new paragraph.
                             %
with the right spacing around the periods requires
a special command.

\LaTeX\ interprets some common characters as
commands, so you must type special commands to
generate them.  These characters include the
following:
       \$ \& \% \# \{ and \}.

In printing, text is usually emphasized with an
       \emph{italic}
type style.

\begin{em}
   A long segment of text can also be emphasized
   in this way.  Text within such a segment can be
   given \emph{additional} emphasis.
\end{em}

It is sometimes necessary to prevent \LaTeX\ from
breaking a line where it might otherwise do so.
This may be at a space, as between the ``Mr.''\ and
``Jones'' in
       ``Mr.~Jones'',        % ~ produces an unbreakable interword space.
or within a word---especially when the word is a
symbol like
       \mbox{\emph{itemnum}}
that makes little sense when hyphenated across
lines.

Footnotes\footnote{This is an example of a footnote.}
pose no problem.

\LaTeX\ is good at typesetting mathematical formulas
like
       \( x-3y + z = 7 \)
or
       \( a_{1} > x^{2n} + y^{2n} > x' \)
or
       \( \ip{A}{B} = \sum_{i} a_{i} b_{i} \).
The spaces you type in a formula are
ignored.  Remember that a letter like
       $x$                   % $ ... $  and  \( ... \)  are equivalent
is a formula when it denotes a mathematical
symbol, and it should be typed as one.

\section{Displayed Text}

Text is displayed by indenting it from the left
margin.  Quotations are commonly displayed.  There
are short quotations
\begin{quote}
   This is a short quotation.  It consists of a
   single paragraph of text.  See how it is formatted.
\end{quote}
and longer ones.
\begin{quotation}
   This is a longer quotation.  It consists of two
   paragraphs of text, neither of which are
   particularly interesting.

   This is the second paragraph of the quotation.  It
   is just as dull as the first paragraph.
\end{quotation}
Another frequently-displayed structure is a list.
The following is an example of an \emph{itemized}
list.
\begin{itemize}
   \item This is the first item of an itemized list.
         Each item in the list is marked with a ``tick''.
         You don't have to worry about what kind of tick
         mark is used.

   \item This is the second item of the list.  It
         contains another list nested inside it.  The inner
         list is an \emph{enumerated} list.
         \begin{enumerate}
            \item This is the first item of an enumerated
                  list that is nested within the itemized list.

            \item This is the second item of the inner list.
                  \LaTeX\ allows you to nest lists deeper than
                  you really should.
         \end{enumerate}
         This is the rest of the second item of the outer
         list.  It is no more interesting than any other
         part of the item.
   \item This is the third item of the list.
\end{itemize}
You can even display poetry.
\begin{verse}
   There is an environment
    for verse \\             % The \\ command separates lines
   Whose features some poets % within a stanza.
   will curse.

                             % One or more blank lines separate stanzas.

   For instead of making\\
   Them do \emph{all} line breaking, \\
   It allows them to put too many words on a line when they'd rather be
   forced to be terse.
\end{verse}

Mathematical formulas may also be displayed.  A
displayed formula
is
one-line long; multiline
formulas require special formatting instructions.
   \[  \ip{\Gamma}{\psi'} = x'' + y^{2} + z_{i}^{n}\]
Don't start a paragraph with a displayed equation,
nor make one a paragraph by itself.

\end{document}               % End of document.

\section{Bench Huge Section 1269}
% synthetic bench section 1269
% This is a sample LaTeX input file.  (Version of 12 August 2004.)
%
% A '%' character causes TeX to ignore all remaining text on the line,
% and is used for comments like this one.

\documentclass{article}      % Specifies the document class

                             % The preamble begins here.
\title{An Example Document}  % Declares the document's title.
\author{Leslie Lamport}      % Declares the author's name.
\date{January 21, 1994}      % Deleting this command produces today's date.

\newcommand{\ip}[2]{(#1, #2)}
                             % Defines \ip{arg1}{arg2} to mean
                             % (arg1, arg2).

%\newcommand{\ip}[2]{\langle #1 | #2\rangle}
                             % This is an alternative definition of
                             % \ip that is commented out.

\begin{document}             % End of preamble and beginning of text.

\maketitle                   % Produces the title.

This is an example input file.  Comparing it with
the output it generates can show you how to
produce a simple document of your own.

\section{Ordinary Text}      % Produces section heading.  Lower-level
                             % sections are begun with similar
                             % \subsection and \subsubsection commands.

The ends  of words and sentences are marked
  by   spaces. It  doesn't matter how many
spaces    you type; one is as good as 100.  The
end of   a line counts as a space.

One   or more   blank lines denote the  end
of  a paragraph.

Since any number of consecutive spaces are treated
like a single one, the formatting of the input
file makes no difference to
      \LaTeX,                % The \LaTeX command generates the LaTeX logo.
but it makes a difference to you.  When you use
\LaTeX, making your input file as easy to read
as possible will be a great help as you write
your document and when you change it.  This sample
file shows how you can add comments to your own input
file.

Because printing is different from typewriting,
there are a number of things that you have to do
differently when preparing an input file than if
you were just typing the document directly.
Quotation marks like
       ``this''
have to be handled specially, as do quotes within
quotes:
       ``\,`this'            % \, separates the double and single quote.
        is what I just
        wrote, not  `that'\,''.

Dashes come in three sizes: an
       intra-word
dash, a medium dash for number ranges like
       1--2,
and a punctuation
       dash---like
this.

A sentence-ending space should be larger than the
space between words within a sentence.  You
sometimes have to type special commands in
conjunction with punctuation characters to get
this right, as in the following sentence.
       Gnats, gnus, etc.\ all  % `\ ' makes an inter-word space.
       begin with G\@.         % \@ marks end-of-sentence punctuation.
You should check the spaces after periods when
reading your output to make sure you haven't
forgotten any special cases.  Generating an
ellipsis
       \ldots\               % `\ ' is needed after `\ldots' because TeX
                             % ignores spaces after command names like \ldots
                             % made from \ + letters.
                             %
                             % Note how a `%' character causes TeX to ignore
                             % the end of the input line, so these blank lines
                             % do not start a new paragraph.
                             %
with the right spacing around the periods requires
a special command.

\LaTeX\ interprets some common characters as
commands, so you must type special commands to
generate them.  These characters include the
following:
       \$ \& \% \# \{ and \}.

In printing, text is usually emphasized with an
       \emph{italic}
type style.

\begin{em}
   A long segment of text can also be emphasized
   in this way.  Text within such a segment can be
   given \emph{additional} emphasis.
\end{em}

It is sometimes necessary to prevent \LaTeX\ from
breaking a line where it might otherwise do so.
This may be at a space, as between the ``Mr.''\ and
``Jones'' in
       ``Mr.~Jones'',        % ~ produces an unbreakable interword space.
or within a word---especially when the word is a
symbol like
       \mbox{\emph{itemnum}}
that makes little sense when hyphenated across
lines.

Footnotes\footnote{This is an example of a footnote.}
pose no problem.

\LaTeX\ is good at typesetting mathematical formulas
like
       \( x-3y + z = 7 \)
or
       \( a_{1} > x^{2n} + y^{2n} > x' \)
or
       \( \ip{A}{B} = \sum_{i} a_{i} b_{i} \).
The spaces you type in a formula are
ignored.  Remember that a letter like
       $x$                   % $ ... $  and  \( ... \)  are equivalent
is a formula when it denotes a mathematical
symbol, and it should be typed as one.

\section{Displayed Text}

Text is displayed by indenting it from the left
margin.  Quotations are commonly displayed.  There
are short quotations
\begin{quote}
   This is a short quotation.  It consists of a
   single paragraph of text.  See how it is formatted.
\end{quote}
and longer ones.
\begin{quotation}
   This is a longer quotation.  It consists of two
   paragraphs of text, neither of which are
   particularly interesting.

   This is the second paragraph of the quotation.  It
   is just as dull as the first paragraph.
\end{quotation}
Another frequently-displayed structure is a list.
The following is an example of an \emph{itemized}
list.
\begin{itemize}
   \item This is the first item of an itemized list.
         Each item in the list is marked with a ``tick''.
         You don't have to worry about what kind of tick
         mark is used.

   \item This is the second item of the list.  It
         contains another list nested inside it.  The inner
         list is an \emph{enumerated} list.
         \begin{enumerate}
            \item This is the first item of an enumerated
                  list that is nested within the itemized list.

            \item This is the second item of the inner list.
                  \LaTeX\ allows you to nest lists deeper than
                  you really should.
         \end{enumerate}
         This is the rest of the second item of the outer
         list.  It is no more interesting than any other
         part of the item.
   \item This is the third item of the list.
\end{itemize}
You can even display poetry.
\begin{verse}
   There is an environment
    for verse \\             % The \\ command separates lines
   Whose features some poets % within a stanza.
   will curse.

                             % One or more blank lines separate stanzas.

   For instead of making\\
   Them do \emph{all} line breaking, \\
   It allows them to put too many words on a line when they'd rather be
   forced to be terse.
\end{verse}

Mathematical formulas may also be displayed.  A
displayed formula
is
one-line long; multiline
formulas require special formatting instructions.
   \[  \ip{\Gamma}{\psi'} = x'' + y^{2} + z_{i}^{n}\]
Don't start a paragraph with a displayed equation,
nor make one a paragraph by itself.

\end{document}               % End of document.

\section{Bench Huge Section 1270}
% synthetic bench section 1270
% This is a sample LaTeX input file.  (Version of 12 August 2004.)
%
% A '%' character causes TeX to ignore all remaining text on the line,
% and is used for comments like this one.

\documentclass{article}      % Specifies the document class

                             % The preamble begins here.
\title{An Example Document}  % Declares the document's title.
\author{Leslie Lamport}      % Declares the author's name.
\date{January 21, 1994}      % Deleting this command produces today's date.

\newcommand{\ip}[2]{(#1, #2)}
                             % Defines \ip{arg1}{arg2} to mean
                             % (arg1, arg2).

%\newcommand{\ip}[2]{\langle #1 | #2\rangle}
                             % This is an alternative definition of
                             % \ip that is commented out.

\begin{document}             % End of preamble and beginning of text.

\maketitle                   % Produces the title.

This is an example input file.  Comparing it with
the output it generates can show you how to
produce a simple document of your own.

\section{Ordinary Text}      % Produces section heading.  Lower-level
                             % sections are begun with similar
                             % \subsection and \subsubsection commands.

The ends  of words and sentences are marked
  by   spaces. It  doesn't matter how many
spaces    you type; one is as good as 100.  The
end of   a line counts as a space.

One   or more   blank lines denote the  end
of  a paragraph.

Since any number of consecutive spaces are treated
like a single one, the formatting of the input
file makes no difference to
      \LaTeX,                % The \LaTeX command generates the LaTeX logo.
but it makes a difference to you.  When you use
\LaTeX, making your input file as easy to read
as possible will be a great help as you write
your document and when you change it.  This sample
file shows how you can add comments to your own input
file.

Because printing is different from typewriting,
there are a number of things that you have to do
differently when preparing an input file than if
you were just typing the document directly.
Quotation marks like
       ``this''
have to be handled specially, as do quotes within
quotes:
       ``\,`this'            % \, separates the double and single quote.
        is what I just
        wrote, not  `that'\,''.

Dashes come in three sizes: an
       intra-word
dash, a medium dash for number ranges like
       1--2,
and a punctuation
       dash---like
this.

A sentence-ending space should be larger than the
space between words within a sentence.  You
sometimes have to type special commands in
conjunction with punctuation characters to get
this right, as in the following sentence.
       Gnats, gnus, etc.\ all  % `\ ' makes an inter-word space.
       begin with G\@.         % \@ marks end-of-sentence punctuation.
You should check the spaces after periods when
reading your output to make sure you haven't
forgotten any special cases.  Generating an
ellipsis
       \ldots\               % `\ ' is needed after `\ldots' because TeX
                             % ignores spaces after command names like \ldots
                             % made from \ + letters.
                             %
                             % Note how a `%' character causes TeX to ignore
                             % the end of the input line, so these blank lines
                             % do not start a new paragraph.
                             %
with the right spacing around the periods requires
a special command.

\LaTeX\ interprets some common characters as
commands, so you must type special commands to
generate them.  These characters include the
following:
       \$ \& \% \# \{ and \}.

In printing, text is usually emphasized with an
       \emph{italic}
type style.

\begin{em}
   A long segment of text can also be emphasized
   in this way.  Text within such a segment can be
   given \emph{additional} emphasis.
\end{em}

It is sometimes necessary to prevent \LaTeX\ from
breaking a line where it might otherwise do so.
This may be at a space, as between the ``Mr.''\ and
``Jones'' in
       ``Mr.~Jones'',        % ~ produces an unbreakable interword space.
or within a word---especially when the word is a
symbol like
       \mbox{\emph{itemnum}}
that makes little sense when hyphenated across
lines.

Footnotes\footnote{This is an example of a footnote.}
pose no problem.

\LaTeX\ is good at typesetting mathematical formulas
like
       \( x-3y + z = 7 \)
or
       \( a_{1} > x^{2n} + y^{2n} > x' \)
or
       \( \ip{A}{B} = \sum_{i} a_{i} b_{i} \).
The spaces you type in a formula are
ignored.  Remember that a letter like
       $x$                   % $ ... $  and  \( ... \)  are equivalent
is a formula when it denotes a mathematical
symbol, and it should be typed as one.

\section{Displayed Text}

Text is displayed by indenting it from the left
margin.  Quotations are commonly displayed.  There
are short quotations
\begin{quote}
   This is a short quotation.  It consists of a
   single paragraph of text.  See how it is formatted.
\end{quote}
and longer ones.
\begin{quotation}
   This is a longer quotation.  It consists of two
   paragraphs of text, neither of which are
   particularly interesting.

   This is the second paragraph of the quotation.  It
   is just as dull as the first paragraph.
\end{quotation}
Another frequently-displayed structure is a list.
The following is an example of an \emph{itemized}
list.
\begin{itemize}
   \item This is the first item of an itemized list.
         Each item in the list is marked with a ``tick''.
         You don't have to worry about what kind of tick
         mark is used.

   \item This is the second item of the list.  It
         contains another list nested inside it.  The inner
         list is an \emph{enumerated} list.
         \begin{enumerate}
            \item This is the first item of an enumerated
                  list that is nested within the itemized list.

            \item This is the second item of the inner list.
                  \LaTeX\ allows you to nest lists deeper than
                  you really should.
         \end{enumerate}
         This is the rest of the second item of the outer
         list.  It is no more interesting than any other
         part of the item.
   \item This is the third item of the list.
\end{itemize}
You can even display poetry.
\begin{verse}
   There is an environment
    for verse \\             % The \\ command separates lines
   Whose features some poets % within a stanza.
   will curse.

                             % One or more blank lines separate stanzas.

   For instead of making\\
   Them do \emph{all} line breaking, \\
   It allows them to put too many words on a line when they'd rather be
   forced to be terse.
\end{verse}

Mathematical formulas may also be displayed.  A
displayed formula
is
one-line long; multiline
formulas require special formatting instructions.
   \[  \ip{\Gamma}{\psi'} = x'' + y^{2} + z_{i}^{n}\]
Don't start a paragraph with a displayed equation,
nor make one a paragraph by itself.

\end{document}               % End of document.

\section{Bench Huge Section 1271}
% synthetic bench section 1271
% This is a sample LaTeX input file.  (Version of 12 August 2004.)
%
% A '%' character causes TeX to ignore all remaining text on the line,
% and is used for comments like this one.

\documentclass{article}      % Specifies the document class

                             % The preamble begins here.
\title{An Example Document}  % Declares the document's title.
\author{Leslie Lamport}      % Declares the author's name.
\date{January 21, 1994}      % Deleting this command produces today's date.

\newcommand{\ip}[2]{(#1, #2)}
                             % Defines \ip{arg1}{arg2} to mean
                             % (arg1, arg2).

%\newcommand{\ip}[2]{\langle #1 | #2\rangle}
                             % This is an alternative definition of
                             % \ip that is commented out.

\begin{document}             % End of preamble and beginning of text.

\maketitle                   % Produces the title.

This is an example input file.  Comparing it with
the output it generates can show you how to
produce a simple document of your own.

\section{Ordinary Text}      % Produces section heading.  Lower-level
                             % sections are begun with similar
                             % \subsection and \subsubsection commands.

The ends  of words and sentences are marked
  by   spaces. It  doesn't matter how many
spaces    you type; one is as good as 100.  The
end of   a line counts as a space.

One   or more   blank lines denote the  end
of  a paragraph.

Since any number of consecutive spaces are treated
like a single one, the formatting of the input
file makes no difference to
      \LaTeX,                % The \LaTeX command generates the LaTeX logo.
but it makes a difference to you.  When you use
\LaTeX, making your input file as easy to read
as possible will be a great help as you write
your document and when you change it.  This sample
file shows how you can add comments to your own input
file.

Because printing is different from typewriting,
there are a number of things that you have to do
differently when preparing an input file than if
you were just typing the document directly.
Quotation marks like
       ``this''
have to be handled specially, as do quotes within
quotes:
       ``\,`this'            % \, separates the double and single quote.
        is what I just
        wrote, not  `that'\,''.

Dashes come in three sizes: an
       intra-word
dash, a medium dash for number ranges like
       1--2,
and a punctuation
       dash---like
this.

A sentence-ending space should be larger than the
space between words within a sentence.  You
sometimes have to type special commands in
conjunction with punctuation characters to get
this right, as in the following sentence.
       Gnats, gnus, etc.\ all  % `\ ' makes an inter-word space.
       begin with G\@.         % \@ marks end-of-sentence punctuation.
You should check the spaces after periods when
reading your output to make sure you haven't
forgotten any special cases.  Generating an
ellipsis
       \ldots\               % `\ ' is needed after `\ldots' because TeX
                             % ignores spaces after command names like \ldots
                             % made from \ + letters.
                             %
                             % Note how a `%' character causes TeX to ignore
                             % the end of the input line, so these blank lines
                             % do not start a new paragraph.
                             %
with the right spacing around the periods requires
a special command.

\LaTeX\ interprets some common characters as
commands, so you must type special commands to
generate them.  These characters include the
following:
       \$ \& \% \# \{ and \}.

In printing, text is usually emphasized with an
       \emph{italic}
type style.

\begin{em}
   A long segment of text can also be emphasized
   in this way.  Text within such a segment can be
   given \emph{additional} emphasis.
\end{em}

It is sometimes necessary to prevent \LaTeX\ from
breaking a line where it might otherwise do so.
This may be at a space, as between the ``Mr.''\ and
``Jones'' in
       ``Mr.~Jones'',        % ~ produces an unbreakable interword space.
or within a word---especially when the word is a
symbol like
       \mbox{\emph{itemnum}}
that makes little sense when hyphenated across
lines.

Footnotes\footnote{This is an example of a footnote.}
pose no problem.

\LaTeX\ is good at typesetting mathematical formulas
like
       \( x-3y + z = 7 \)
or
       \( a_{1} > x^{2n} + y^{2n} > x' \)
or
       \( \ip{A}{B} = \sum_{i} a_{i} b_{i} \).
The spaces you type in a formula are
ignored.  Remember that a letter like
       $x$                   % $ ... $  and  \( ... \)  are equivalent
is a formula when it denotes a mathematical
symbol, and it should be typed as one.

\section{Displayed Text}

Text is displayed by indenting it from the left
margin.  Quotations are commonly displayed.  There
are short quotations
\begin{quote}
   This is a short quotation.  It consists of a
   single paragraph of text.  See how it is formatted.
\end{quote}
and longer ones.
\begin{quotation}
   This is a longer quotation.  It consists of two
   paragraphs of text, neither of which are
   particularly interesting.

   This is the second paragraph of the quotation.  It
   is just as dull as the first paragraph.
\end{quotation}
Another frequently-displayed structure is a list.
The following is an example of an \emph{itemized}
list.
\begin{itemize}
   \item This is the first item of an itemized list.
         Each item in the list is marked with a ``tick''.
         You don't have to worry about what kind of tick
         mark is used.

   \item This is the second item of the list.  It
         contains another list nested inside it.  The inner
         list is an \emph{enumerated} list.
         \begin{enumerate}
            \item This is the first item of an enumerated
                  list that is nested within the itemized list.

            \item This is the second item of the inner list.
                  \LaTeX\ allows you to nest lists deeper than
                  you really should.
         \end{enumerate}
         This is the rest of the second item of the outer
         list.  It is no more interesting than any other
         part of the item.
   \item This is the third item of the list.
\end{itemize}
You can even display poetry.
\begin{verse}
   There is an environment
    for verse \\             % The \\ command separates lines
   Whose features some poets % within a stanza.
   will curse.

                             % One or more blank lines separate stanzas.

   For instead of making\\
   Them do \emph{all} line breaking, \\
   It allows them to put too many words on a line when they'd rather be
   forced to be terse.
\end{verse}

Mathematical formulas may also be displayed.  A
displayed formula
is
one-line long; multiline
formulas require special formatting instructions.
   \[  \ip{\Gamma}{\psi'} = x'' + y^{2} + z_{i}^{n}\]
Don't start a paragraph with a displayed equation,
nor make one a paragraph by itself.

\end{document}               % End of document.

\section{Bench Huge Section 1272}
% synthetic bench section 1272
% This is a sample LaTeX input file.  (Version of 12 August 2004.)
%
% A '%' character causes TeX to ignore all remaining text on the line,
% and is used for comments like this one.

\documentclass{article}      % Specifies the document class

                             % The preamble begins here.
\title{An Example Document}  % Declares the document's title.
\author{Leslie Lamport}      % Declares the author's name.
\date{January 21, 1994}      % Deleting this command produces today's date.

\newcommand{\ip}[2]{(#1, #2)}
                             % Defines \ip{arg1}{arg2} to mean
                             % (arg1, arg2).

%\newcommand{\ip}[2]{\langle #1 | #2\rangle}
                             % This is an alternative definition of
                             % \ip that is commented out.

\begin{document}             % End of preamble and beginning of text.

\maketitle                   % Produces the title.

This is an example input file.  Comparing it with
the output it generates can show you how to
produce a simple document of your own.

\section{Ordinary Text}      % Produces section heading.  Lower-level
                             % sections are begun with similar
                             % \subsection and \subsubsection commands.

The ends  of words and sentences are marked
  by   spaces. It  doesn't matter how many
spaces    you type; one is as good as 100.  The
end of   a line counts as a space.

One   or more   blank lines denote the  end
of  a paragraph.

Since any number of consecutive spaces are treated
like a single one, the formatting of the input
file makes no difference to
      \LaTeX,                % The \LaTeX command generates the LaTeX logo.
but it makes a difference to you.  When you use
\LaTeX, making your input file as easy to read
as possible will be a great help as you write
your document and when you change it.  This sample
file shows how you can add comments to your own input
file.

Because printing is different from typewriting,
there are a number of things that you have to do
differently when preparing an input file than if
you were just typing the document directly.
Quotation marks like
       ``this''
have to be handled specially, as do quotes within
quotes:
       ``\,`this'            % \, separates the double and single quote.
        is what I just
        wrote, not  `that'\,''.

Dashes come in three sizes: an
       intra-word
dash, a medium dash for number ranges like
       1--2,
and a punctuation
       dash---like
this.

A sentence-ending space should be larger than the
space between words within a sentence.  You
sometimes have to type special commands in
conjunction with punctuation characters to get
this right, as in the following sentence.
       Gnats, gnus, etc.\ all  % `\ ' makes an inter-word space.
       begin with G\@.         % \@ marks end-of-sentence punctuation.
You should check the spaces after periods when
reading your output to make sure you haven't
forgotten any special cases.  Generating an
ellipsis
       \ldots\               % `\ ' is needed after `\ldots' because TeX
                             % ignores spaces after command names like \ldots
                             % made from \ + letters.
                             %
                             % Note how a `%' character causes TeX to ignore
                             % the end of the input line, so these blank lines
                             % do not start a new paragraph.
                             %
with the right spacing around the periods requires
a special command.

\LaTeX\ interprets some common characters as
commands, so you must type special commands to
generate them.  These characters include the
following:
       \$ \& \% \# \{ and \}.

In printing, text is usually emphasized with an
       \emph{italic}
type style.

\begin{em}
   A long segment of text can also be emphasized
   in this way.  Text within such a segment can be
   given \emph{additional} emphasis.
\end{em}

It is sometimes necessary to prevent \LaTeX\ from
breaking a line where it might otherwise do so.
This may be at a space, as between the ``Mr.''\ and
``Jones'' in
       ``Mr.~Jones'',        % ~ produces an unbreakable interword space.
or within a word---especially when the word is a
symbol like
       \mbox{\emph{itemnum}}
that makes little sense when hyphenated across
lines.

Footnotes\footnote{This is an example of a footnote.}
pose no problem.

\LaTeX\ is good at typesetting mathematical formulas
like
       \( x-3y + z = 7 \)
or
       \( a_{1} > x^{2n} + y^{2n} > x' \)
or
       \( \ip{A}{B} = \sum_{i} a_{i} b_{i} \).
The spaces you type in a formula are
ignored.  Remember that a letter like
       $x$                   % $ ... $  and  \( ... \)  are equivalent
is a formula when it denotes a mathematical
symbol, and it should be typed as one.

\section{Displayed Text}

Text is displayed by indenting it from the left
margin.  Quotations are commonly displayed.  There
are short quotations
\begin{quote}
   This is a short quotation.  It consists of a
   single paragraph of text.  See how it is formatted.
\end{quote}
and longer ones.
\begin{quotation}
   This is a longer quotation.  It consists of two
   paragraphs of text, neither of which are
   particularly interesting.

   This is the second paragraph of the quotation.  It
   is just as dull as the first paragraph.
\end{quotation}
Another frequently-displayed structure is a list.
The following is an example of an \emph{itemized}
list.
\begin{itemize}
   \item This is the first item of an itemized list.
         Each item in the list is marked with a ``tick''.
         You don't have to worry about what kind of tick
         mark is used.

   \item This is the second item of the list.  It
         contains another list nested inside it.  The inner
         list is an \emph{enumerated} list.
         \begin{enumerate}
            \item This is the first item of an enumerated
                  list that is nested within the itemized list.

            \item This is the second item of the inner list.
                  \LaTeX\ allows you to nest lists deeper than
                  you really should.
         \end{enumerate}
         This is the rest of the second item of the outer
         list.  It is no more interesting than any other
         part of the item.
   \item This is the third item of the list.
\end{itemize}
You can even display poetry.
\begin{verse}
   There is an environment
    for verse \\             % The \\ command separates lines
   Whose features some poets % within a stanza.
   will curse.

                             % One or more blank lines separate stanzas.

   For instead of making\\
   Them do \emph{all} line breaking, \\
   It allows them to put too many words on a line when they'd rather be
   forced to be terse.
\end{verse}

Mathematical formulas may also be displayed.  A
displayed formula
is
one-line long; multiline
formulas require special formatting instructions.
   \[  \ip{\Gamma}{\psi'} = x'' + y^{2} + z_{i}^{n}\]
Don't start a paragraph with a displayed equation,
nor make one a paragraph by itself.

\end{document}               % End of document.

\section{Bench Huge Section 1273}
% synthetic bench section 1273
% This is a sample LaTeX input file.  (Version of 12 August 2004.)
%
% A '%' character causes TeX to ignore all remaining text on the line,
% and is used for comments like this one.

\documentclass{article}      % Specifies the document class

                             % The preamble begins here.
\title{An Example Document}  % Declares the document's title.
\author{Leslie Lamport}      % Declares the author's name.
\date{January 21, 1994}      % Deleting this command produces today's date.

\newcommand{\ip}[2]{(#1, #2)}
                             % Defines \ip{arg1}{arg2} to mean
                             % (arg1, arg2).

%\newcommand{\ip}[2]{\langle #1 | #2\rangle}
                             % This is an alternative definition of
                             % \ip that is commented out.

\begin{document}             % End of preamble and beginning of text.

\maketitle                   % Produces the title.

This is an example input file.  Comparing it with
the output it generates can show you how to
produce a simple document of your own.

\section{Ordinary Text}      % Produces section heading.  Lower-level
                             % sections are begun with similar
                             % \subsection and \subsubsection commands.

The ends  of words and sentences are marked
  by   spaces. It  doesn't matter how many
spaces    you type; one is as good as 100.  The
end of   a line counts as a space.

One   or more   blank lines denote the  end
of  a paragraph.

Since any number of consecutive spaces are treated
like a single one, the formatting of the input
file makes no difference to
      \LaTeX,                % The \LaTeX command generates the LaTeX logo.
but it makes a difference to you.  When you use
\LaTeX, making your input file as easy to read
as possible will be a great help as you write
your document and when you change it.  This sample
file shows how you can add comments to your own input
file.

Because printing is different from typewriting,
there are a number of things that you have to do
differently when preparing an input file than if
you were just typing the document directly.
Quotation marks like
       ``this''
have to be handled specially, as do quotes within
quotes:
       ``\,`this'            % \, separates the double and single quote.
        is what I just
        wrote, not  `that'\,''.

Dashes come in three sizes: an
       intra-word
dash, a medium dash for number ranges like
       1--2,
and a punctuation
       dash---like
this.

A sentence-ending space should be larger than the
space between words within a sentence.  You
sometimes have to type special commands in
conjunction with punctuation characters to get
this right, as in the following sentence.
       Gnats, gnus, etc.\ all  % `\ ' makes an inter-word space.
       begin with G\@.         % \@ marks end-of-sentence punctuation.
You should check the spaces after periods when
reading your output to make sure you haven't
forgotten any special cases.  Generating an
ellipsis
       \ldots\               % `\ ' is needed after `\ldots' because TeX
                             % ignores spaces after command names like \ldots
                             % made from \ + letters.
                             %
                             % Note how a `%' character causes TeX to ignore
                             % the end of the input line, so these blank lines
                             % do not start a new paragraph.
                             %
with the right spacing around the periods requires
a special command.

\LaTeX\ interprets some common characters as
commands, so you must type special commands to
generate them.  These characters include the
following:
       \$ \& \% \# \{ and \}.

In printing, text is usually emphasized with an
       \emph{italic}
type style.

\begin{em}
   A long segment of text can also be emphasized
   in this way.  Text within such a segment can be
   given \emph{additional} emphasis.
\end{em}

It is sometimes necessary to prevent \LaTeX\ from
breaking a line where it might otherwise do so.
This may be at a space, as between the ``Mr.''\ and
``Jones'' in
       ``Mr.~Jones'',        % ~ produces an unbreakable interword space.
or within a word---especially when the word is a
symbol like
       \mbox{\emph{itemnum}}
that makes little sense when hyphenated across
lines.

Footnotes\footnote{This is an example of a footnote.}
pose no problem.

\LaTeX\ is good at typesetting mathematical formulas
like
       \( x-3y + z = 7 \)
or
       \( a_{1} > x^{2n} + y^{2n} > x' \)
or
       \( \ip{A}{B} = \sum_{i} a_{i} b_{i} \).
The spaces you type in a formula are
ignored.  Remember that a letter like
       $x$                   % $ ... $  and  \( ... \)  are equivalent
is a formula when it denotes a mathematical
symbol, and it should be typed as one.

\section{Displayed Text}

Text is displayed by indenting it from the left
margin.  Quotations are commonly displayed.  There
are short quotations
\begin{quote}
   This is a short quotation.  It consists of a
   single paragraph of text.  See how it is formatted.
\end{quote}
and longer ones.
\begin{quotation}
   This is a longer quotation.  It consists of two
   paragraphs of text, neither of which are
   particularly interesting.

   This is the second paragraph of the quotation.  It
   is just as dull as the first paragraph.
\end{quotation}
Another frequently-displayed structure is a list.
The following is an example of an \emph{itemized}
list.
\begin{itemize}
   \item This is the first item of an itemized list.
         Each item in the list is marked with a ``tick''.
         You don't have to worry about what kind of tick
         mark is used.

   \item This is the second item of the list.  It
         contains another list nested inside it.  The inner
         list is an \emph{enumerated} list.
         \begin{enumerate}
            \item This is the first item of an enumerated
                  list that is nested within the itemized list.

            \item This is the second item of the inner list.
                  \LaTeX\ allows you to nest lists deeper than
                  you really should.
         \end{enumerate}
         This is the rest of the second item of the outer
         list.  It is no more interesting than any other
         part of the item.
   \item This is the third item of the list.
\end{itemize}
You can even display poetry.
\begin{verse}
   There is an environment
    for verse \\             % The \\ command separates lines
   Whose features some poets % within a stanza.
   will curse.

                             % One or more blank lines separate stanzas.

   For instead of making\\
   Them do \emph{all} line breaking, \\
   It allows them to put too many words on a line when they'd rather be
   forced to be terse.
\end{verse}

Mathematical formulas may also be displayed.  A
displayed formula
is
one-line long; multiline
formulas require special formatting instructions.
   \[  \ip{\Gamma}{\psi'} = x'' + y^{2} + z_{i}^{n}\]
Don't start a paragraph with a displayed equation,
nor make one a paragraph by itself.

\end{document}               % End of document.

\section{Bench Huge Section 1274}
% synthetic bench section 1274
% This is a sample LaTeX input file.  (Version of 12 August 2004.)
%
% A '%' character causes TeX to ignore all remaining text on the line,
% and is used for comments like this one.

\documentclass{article}      % Specifies the document class

                             % The preamble begins here.
\title{An Example Document}  % Declares the document's title.
\author{Leslie Lamport}      % Declares the author's name.
\date{January 21, 1994}      % Deleting this command produces today's date.

\newcommand{\ip}[2]{(#1, #2)}
                             % Defines \ip{arg1}{arg2} to mean
                             % (arg1, arg2).

%\newcommand{\ip}[2]{\langle #1 | #2\rangle}
                             % This is an alternative definition of
                             % \ip that is commented out.

\begin{document}             % End of preamble and beginning of text.

\maketitle                   % Produces the title.

This is an example input file.  Comparing it with
the output it generates can show you how to
produce a simple document of your own.

\section{Ordinary Text}      % Produces section heading.  Lower-level
                             % sections are begun with similar
                             % \subsection and \subsubsection commands.

The ends  of words and sentences are marked
  by   spaces. It  doesn't matter how many
spaces    you type; one is as good as 100.  The
end of   a line counts as a space.

One   or more   blank lines denote the  end
of  a paragraph.

Since any number of consecutive spaces are treated
like a single one, the formatting of the input
file makes no difference to
      \LaTeX,                % The \LaTeX command generates the LaTeX logo.
but it makes a difference to you.  When you use
\LaTeX, making your input file as easy to read
as possible will be a great help as you write
your document and when you change it.  This sample
file shows how you can add comments to your own input
file.

Because printing is different from typewriting,
there are a number of things that you have to do
differently when preparing an input file than if
you were just typing the document directly.
Quotation marks like
       ``this''
have to be handled specially, as do quotes within
quotes:
       ``\,`this'            % \, separates the double and single quote.
        is what I just
        wrote, not  `that'\,''.

Dashes come in three sizes: an
       intra-word
dash, a medium dash for number ranges like
       1--2,
and a punctuation
       dash---like
this.

A sentence-ending space should be larger than the
space between words within a sentence.  You
sometimes have to type special commands in
conjunction with punctuation characters to get
this right, as in the following sentence.
       Gnats, gnus, etc.\ all  % `\ ' makes an inter-word space.
       begin with G\@.         % \@ marks end-of-sentence punctuation.
You should check the spaces after periods when
reading your output to make sure you haven't
forgotten any special cases.  Generating an
ellipsis
       \ldots\               % `\ ' is needed after `\ldots' because TeX
                             % ignores spaces after command names like \ldots
                             % made from \ + letters.
                             %
                             % Note how a `%' character causes TeX to ignore
                             % the end of the input line, so these blank lines
                             % do not start a new paragraph.
                             %
with the right spacing around the periods requires
a special command.

\LaTeX\ interprets some common characters as
commands, so you must type special commands to
generate them.  These characters include the
following:
       \$ \& \% \# \{ and \}.

In printing, text is usually emphasized with an
       \emph{italic}
type style.

\begin{em}
   A long segment of text can also be emphasized
   in this way.  Text within such a segment can be
   given \emph{additional} emphasis.
\end{em}

It is sometimes necessary to prevent \LaTeX\ from
breaking a line where it might otherwise do so.
This may be at a space, as between the ``Mr.''\ and
``Jones'' in
       ``Mr.~Jones'',        % ~ produces an unbreakable interword space.
or within a word---especially when the word is a
symbol like
       \mbox{\emph{itemnum}}
that makes little sense when hyphenated across
lines.

Footnotes\footnote{This is an example of a footnote.}
pose no problem.

\LaTeX\ is good at typesetting mathematical formulas
like
       \( x-3y + z = 7 \)
or
       \( a_{1} > x^{2n} + y^{2n} > x' \)
or
       \( \ip{A}{B} = \sum_{i} a_{i} b_{i} \).
The spaces you type in a formula are
ignored.  Remember that a letter like
       $x$                   % $ ... $  and  \( ... \)  are equivalent
is a formula when it denotes a mathematical
symbol, and it should be typed as one.

\section{Displayed Text}

Text is displayed by indenting it from the left
margin.  Quotations are commonly displayed.  There
are short quotations
\begin{quote}
   This is a short quotation.  It consists of a
   single paragraph of text.  See how it is formatted.
\end{quote}
and longer ones.
\begin{quotation}
   This is a longer quotation.  It consists of two
   paragraphs of text, neither of which are
   particularly interesting.

   This is the second paragraph of the quotation.  It
   is just as dull as the first paragraph.
\end{quotation}
Another frequently-displayed structure is a list.
The following is an example of an \emph{itemized}
list.
\begin{itemize}
   \item This is the first item of an itemized list.
         Each item in the list is marked with a ``tick''.
         You don't have to worry about what kind of tick
         mark is used.

   \item This is the second item of the list.  It
         contains another list nested inside it.  The inner
         list is an \emph{enumerated} list.
         \begin{enumerate}
            \item This is the first item of an enumerated
                  list that is nested within the itemized list.

            \item This is the second item of the inner list.
                  \LaTeX\ allows you to nest lists deeper than
                  you really should.
         \end{enumerate}
         This is the rest of the second item of the outer
         list.  It is no more interesting than any other
         part of the item.
   \item This is the third item of the list.
\end{itemize}
You can even display poetry.
\begin{verse}
   There is an environment
    for verse \\             % The \\ command separates lines
   Whose features some poets % within a stanza.
   will curse.

                             % One or more blank lines separate stanzas.

   For instead of making\\
   Them do \emph{all} line breaking, \\
   It allows them to put too many words on a line when they'd rather be
   forced to be terse.
\end{verse}

Mathematical formulas may also be displayed.  A
displayed formula
is
one-line long; multiline
formulas require special formatting instructions.
   \[  \ip{\Gamma}{\psi'} = x'' + y^{2} + z_{i}^{n}\]
Don't start a paragraph with a displayed equation,
nor make one a paragraph by itself.

\end{document}               % End of document.

\section{Bench Huge Section 1275}
% synthetic bench section 1275
% This is a sample LaTeX input file.  (Version of 12 August 2004.)
%
% A '%' character causes TeX to ignore all remaining text on the line,
% and is used for comments like this one.

\documentclass{article}      % Specifies the document class

                             % The preamble begins here.
\title{An Example Document}  % Declares the document's title.
\author{Leslie Lamport}      % Declares the author's name.
\date{January 21, 1994}      % Deleting this command produces today's date.

\newcommand{\ip}[2]{(#1, #2)}
                             % Defines \ip{arg1}{arg2} to mean
                             % (arg1, arg2).

%\newcommand{\ip}[2]{\langle #1 | #2\rangle}
                             % This is an alternative definition of
                             % \ip that is commented out.

\begin{document}             % End of preamble and beginning of text.

\maketitle                   % Produces the title.

This is an example input file.  Comparing it with
the output it generates can show you how to
produce a simple document of your own.

\section{Ordinary Text}      % Produces section heading.  Lower-level
                             % sections are begun with similar
                             % \subsection and \subsubsection commands.

The ends  of words and sentences are marked
  by   spaces. It  doesn't matter how many
spaces    you type; one is as good as 100.  The
end of   a line counts as a space.

One   or more   blank lines denote the  end
of  a paragraph.

Since any number of consecutive spaces are treated
like a single one, the formatting of the input
file makes no difference to
      \LaTeX,                % The \LaTeX command generates the LaTeX logo.
but it makes a difference to you.  When you use
\LaTeX, making your input file as easy to read
as possible will be a great help as you write
your document and when you change it.  This sample
file shows how you can add comments to your own input
file.

Because printing is different from typewriting,
there are a number of things that you have to do
differently when preparing an input file than if
you were just typing the document directly.
Quotation marks like
       ``this''
have to be handled specially, as do quotes within
quotes:
       ``\,`this'            % \, separates the double and single quote.
        is what I just
        wrote, not  `that'\,''.

Dashes come in three sizes: an
       intra-word
dash, a medium dash for number ranges like
       1--2,
and a punctuation
       dash---like
this.

A sentence-ending space should be larger than the
space between words within a sentence.  You
sometimes have to type special commands in
conjunction with punctuation characters to get
this right, as in the following sentence.
       Gnats, gnus, etc.\ all  % `\ ' makes an inter-word space.
       begin with G\@.         % \@ marks end-of-sentence punctuation.
You should check the spaces after periods when
reading your output to make sure you haven't
forgotten any special cases.  Generating an
ellipsis
       \ldots\               % `\ ' is needed after `\ldots' because TeX
                             % ignores spaces after command names like \ldots
                             % made from \ + letters.
                             %
                             % Note how a `%' character causes TeX to ignore
                             % the end of the input line, so these blank lines
                             % do not start a new paragraph.
                             %
with the right spacing around the periods requires
a special command.

\LaTeX\ interprets some common characters as
commands, so you must type special commands to
generate them.  These characters include the
following:
       \$ \& \% \# \{ and \}.

In printing, text is usually emphasized with an
       \emph{italic}
type style.

\begin{em}
   A long segment of text can also be emphasized
   in this way.  Text within such a segment can be
   given \emph{additional} emphasis.
\end{em}

It is sometimes necessary to prevent \LaTeX\ from
breaking a line where it might otherwise do so.
This may be at a space, as between the ``Mr.''\ and
``Jones'' in
       ``Mr.~Jones'',        % ~ produces an unbreakable interword space.
or within a word---especially when the word is a
symbol like
       \mbox{\emph{itemnum}}
that makes little sense when hyphenated across
lines.

Footnotes\footnote{This is an example of a footnote.}
pose no problem.

\LaTeX\ is good at typesetting mathematical formulas
like
       \( x-3y + z = 7 \)
or
       \( a_{1} > x^{2n} + y^{2n} > x' \)
or
       \( \ip{A}{B} = \sum_{i} a_{i} b_{i} \).
The spaces you type in a formula are
ignored.  Remember that a letter like
       $x$                   % $ ... $  and  \( ... \)  are equivalent
is a formula when it denotes a mathematical
symbol, and it should be typed as one.

\section{Displayed Text}

Text is displayed by indenting it from the left
margin.  Quotations are commonly displayed.  There
are short quotations
\begin{quote}
   This is a short quotation.  It consists of a
   single paragraph of text.  See how it is formatted.
\end{quote}
and longer ones.
\begin{quotation}
   This is a longer quotation.  It consists of two
   paragraphs of text, neither of which are
   particularly interesting.

   This is the second paragraph of the quotation.  It
   is just as dull as the first paragraph.
\end{quotation}
Another frequently-displayed structure is a list.
The following is an example of an \emph{itemized}
list.
\begin{itemize}
   \item This is the first item of an itemized list.
         Each item in the list is marked with a ``tick''.
         You don't have to worry about what kind of tick
         mark is used.

   \item This is the second item of the list.  It
         contains another list nested inside it.  The inner
         list is an \emph{enumerated} list.
         \begin{enumerate}
            \item This is the first item of an enumerated
                  list that is nested within the itemized list.

            \item This is the second item of the inner list.
                  \LaTeX\ allows you to nest lists deeper than
                  you really should.
         \end{enumerate}
         This is the rest of the second item of the outer
         list.  It is no more interesting than any other
         part of the item.
   \item This is the third item of the list.
\end{itemize}
You can even display poetry.
\begin{verse}
   There is an environment
    for verse \\             % The \\ command separates lines
   Whose features some poets % within a stanza.
   will curse.

                             % One or more blank lines separate stanzas.

   For instead of making\\
   Them do \emph{all} line breaking, \\
   It allows them to put too many words on a line when they'd rather be
   forced to be terse.
\end{verse}

Mathematical formulas may also be displayed.  A
displayed formula
is
one-line long; multiline
formulas require special formatting instructions.
   \[  \ip{\Gamma}{\psi'} = x'' + y^{2} + z_{i}^{n}\]
Don't start a paragraph with a displayed equation,
nor make one a paragraph by itself.

\end{document}               % End of document.

\section{Bench Huge Section 1276}
% synthetic bench section 1276
% This is a sample LaTeX input file.  (Version of 12 August 2004.)
%
% A '%' character causes TeX to ignore all remaining text on the line,
% and is used for comments like this one.

\documentclass{article}      % Specifies the document class

                             % The preamble begins here.
\title{An Example Document}  % Declares the document's title.
\author{Leslie Lamport}      % Declares the author's name.
\date{January 21, 1994}      % Deleting this command produces today's date.

\newcommand{\ip}[2]{(#1, #2)}
                             % Defines \ip{arg1}{arg2} to mean
                             % (arg1, arg2).

%\newcommand{\ip}[2]{\langle #1 | #2\rangle}
                             % This is an alternative definition of
                             % \ip that is commented out.

\begin{document}             % End of preamble and beginning of text.

\maketitle                   % Produces the title.

This is an example input file.  Comparing it with
the output it generates can show you how to
produce a simple document of your own.

\section{Ordinary Text}      % Produces section heading.  Lower-level
                             % sections are begun with similar
                             % \subsection and \subsubsection commands.

The ends  of words and sentences are marked
  by   spaces. It  doesn't matter how many
spaces    you type; one is as good as 100.  The
end of   a line counts as a space.

One   or more   blank lines denote the  end
of  a paragraph.

Since any number of consecutive spaces are treated
like a single one, the formatting of the input
file makes no difference to
      \LaTeX,                % The \LaTeX command generates the LaTeX logo.
but it makes a difference to you.  When you use
\LaTeX, making your input file as easy to read
as possible will be a great help as you write
your document and when you change it.  This sample
file shows how you can add comments to your own input
file.

Because printing is different from typewriting,
there are a number of things that you have to do
differently when preparing an input file than if
you were just typing the document directly.
Quotation marks like
       ``this''
have to be handled specially, as do quotes within
quotes:
       ``\,`this'            % \, separates the double and single quote.
        is what I just
        wrote, not  `that'\,''.

Dashes come in three sizes: an
       intra-word
dash, a medium dash for number ranges like
       1--2,
and a punctuation
       dash---like
this.

A sentence-ending space should be larger than the
space between words within a sentence.  You
sometimes have to type special commands in
conjunction with punctuation characters to get
this right, as in the following sentence.
       Gnats, gnus, etc.\ all  % `\ ' makes an inter-word space.
       begin with G\@.         % \@ marks end-of-sentence punctuation.
You should check the spaces after periods when
reading your output to make sure you haven't
forgotten any special cases.  Generating an
ellipsis
       \ldots\               % `\ ' is needed after `\ldots' because TeX
                             % ignores spaces after command names like \ldots
                             % made from \ + letters.
                             %
                             % Note how a `%' character causes TeX to ignore
                             % the end of the input line, so these blank lines
                             % do not start a new paragraph.
                             %
with the right spacing around the periods requires
a special command.

\LaTeX\ interprets some common characters as
commands, so you must type special commands to
generate them.  These characters include the
following:
       \$ \& \% \# \{ and \}.

In printing, text is usually emphasized with an
       \emph{italic}
type style.

\begin{em}
   A long segment of text can also be emphasized
   in this way.  Text within such a segment can be
   given \emph{additional} emphasis.
\end{em}

It is sometimes necessary to prevent \LaTeX\ from
breaking a line where it might otherwise do so.
This may be at a space, as between the ``Mr.''\ and
``Jones'' in
       ``Mr.~Jones'',        % ~ produces an unbreakable interword space.
or within a word---especially when the word is a
symbol like
       \mbox{\emph{itemnum}}
that makes little sense when hyphenated across
lines.

Footnotes\footnote{This is an example of a footnote.}
pose no problem.

\LaTeX\ is good at typesetting mathematical formulas
like
       \( x-3y + z = 7 \)
or
       \( a_{1} > x^{2n} + y^{2n} > x' \)
or
       \( \ip{A}{B} = \sum_{i} a_{i} b_{i} \).
The spaces you type in a formula are
ignored.  Remember that a letter like
       $x$                   % $ ... $  and  \( ... \)  are equivalent
is a formula when it denotes a mathematical
symbol, and it should be typed as one.

\section{Displayed Text}

Text is displayed by indenting it from the left
margin.  Quotations are commonly displayed.  There
are short quotations
\begin{quote}
   This is a short quotation.  It consists of a
   single paragraph of text.  See how it is formatted.
\end{quote}
and longer ones.
\begin{quotation}
   This is a longer quotation.  It consists of two
   paragraphs of text, neither of which are
   particularly interesting.

   This is the second paragraph of the quotation.  It
   is just as dull as the first paragraph.
\end{quotation}
Another frequently-displayed structure is a list.
The following is an example of an \emph{itemized}
list.
\begin{itemize}
   \item This is the first item of an itemized list.
         Each item in the list is marked with a ``tick''.
         You don't have to worry about what kind of tick
         mark is used.

   \item This is the second item of the list.  It
         contains another list nested inside it.  The inner
         list is an \emph{enumerated} list.
         \begin{enumerate}
            \item This is the first item of an enumerated
                  list that is nested within the itemized list.

            \item This is the second item of the inner list.
                  \LaTeX\ allows you to nest lists deeper than
                  you really should.
         \end{enumerate}
         This is the rest of the second item of the outer
         list.  It is no more interesting than any other
         part of the item.
   \item This is the third item of the list.
\end{itemize}
You can even display poetry.
\begin{verse}
   There is an environment
    for verse \\             % The \\ command separates lines
   Whose features some poets % within a stanza.
   will curse.

                             % One or more blank lines separate stanzas.

   For instead of making\\
   Them do \emph{all} line breaking, \\
   It allows them to put too many words on a line when they'd rather be
   forced to be terse.
\end{verse}

Mathematical formulas may also be displayed.  A
displayed formula
is
one-line long; multiline
formulas require special formatting instructions.
   \[  \ip{\Gamma}{\psi'} = x'' + y^{2} + z_{i}^{n}\]
Don't start a paragraph with a displayed equation,
nor make one a paragraph by itself.

\end{document}               % End of document.

\section{Bench Huge Section 1277}
% synthetic bench section 1277
% This is a sample LaTeX input file.  (Version of 12 August 2004.)
%
% A '%' character causes TeX to ignore all remaining text on the line,
% and is used for comments like this one.

\documentclass{article}      % Specifies the document class

                             % The preamble begins here.
\title{An Example Document}  % Declares the document's title.
\author{Leslie Lamport}      % Declares the author's name.
\date{January 21, 1994}      % Deleting this command produces today's date.

\newcommand{\ip}[2]{(#1, #2)}
                             % Defines \ip{arg1}{arg2} to mean
                             % (arg1, arg2).

%\newcommand{\ip}[2]{\langle #1 | #2\rangle}
                             % This is an alternative definition of
                             % \ip that is commented out.

\begin{document}             % End of preamble and beginning of text.

\maketitle                   % Produces the title.

This is an example input file.  Comparing it with
the output it generates can show you how to
produce a simple document of your own.

\section{Ordinary Text}      % Produces section heading.  Lower-level
                             % sections are begun with similar
                             % \subsection and \subsubsection commands.

The ends  of words and sentences are marked
  by   spaces. It  doesn't matter how many
spaces    you type; one is as good as 100.  The
end of   a line counts as a space.

One   or more   blank lines denote the  end
of  a paragraph.

Since any number of consecutive spaces are treated
like a single one, the formatting of the input
file makes no difference to
      \LaTeX,                % The \LaTeX command generates the LaTeX logo.
but it makes a difference to you.  When you use
\LaTeX, making your input file as easy to read
as possible will be a great help as you write
your document and when you change it.  This sample
file shows how you can add comments to your own input
file.

Because printing is different from typewriting,
there are a number of things that you have to do
differently when preparing an input file than if
you were just typing the document directly.
Quotation marks like
       ``this''
have to be handled specially, as do quotes within
quotes:
       ``\,`this'            % \, separates the double and single quote.
        is what I just
        wrote, not  `that'\,''.

Dashes come in three sizes: an
       intra-word
dash, a medium dash for number ranges like
       1--2,
and a punctuation
       dash---like
this.

A sentence-ending space should be larger than the
space between words within a sentence.  You
sometimes have to type special commands in
conjunction with punctuation characters to get
this right, as in the following sentence.
       Gnats, gnus, etc.\ all  % `\ ' makes an inter-word space.
       begin with G\@.         % \@ marks end-of-sentence punctuation.
You should check the spaces after periods when
reading your output to make sure you haven't
forgotten any special cases.  Generating an
ellipsis
       \ldots\               % `\ ' is needed after `\ldots' because TeX
                             % ignores spaces after command names like \ldots
                             % made from \ + letters.
                             %
                             % Note how a `%' character causes TeX to ignore
                             % the end of the input line, so these blank lines
                             % do not start a new paragraph.
                             %
with the right spacing around the periods requires
a special command.

\LaTeX\ interprets some common characters as
commands, so you must type special commands to
generate them.  These characters include the
following:
       \$ \& \% \# \{ and \}.

In printing, text is usually emphasized with an
       \emph{italic}
type style.

\begin{em}
   A long segment of text can also be emphasized
   in this way.  Text within such a segment can be
   given \emph{additional} emphasis.
\end{em}

It is sometimes necessary to prevent \LaTeX\ from
breaking a line where it might otherwise do so.
This may be at a space, as between the ``Mr.''\ and
``Jones'' in
       ``Mr.~Jones'',        % ~ produces an unbreakable interword space.
or within a word---especially when the word is a
symbol like
       \mbox{\emph{itemnum}}
that makes little sense when hyphenated across
lines.

Footnotes\footnote{This is an example of a footnote.}
pose no problem.

\LaTeX\ is good at typesetting mathematical formulas
like
       \( x-3y + z = 7 \)
or
       \( a_{1} > x^{2n} + y^{2n} > x' \)
or
       \( \ip{A}{B} = \sum_{i} a_{i} b_{i} \).
The spaces you type in a formula are
ignored.  Remember that a letter like
       $x$                   % $ ... $  and  \( ... \)  are equivalent
is a formula when it denotes a mathematical
symbol, and it should be typed as one.

\section{Displayed Text}

Text is displayed by indenting it from the left
margin.  Quotations are commonly displayed.  There
are short quotations
\begin{quote}
   This is a short quotation.  It consists of a
   single paragraph of text.  See how it is formatted.
\end{quote}
and longer ones.
\begin{quotation}
   This is a longer quotation.  It consists of two
   paragraphs of text, neither of which are
   particularly interesting.

   This is the second paragraph of the quotation.  It
   is just as dull as the first paragraph.
\end{quotation}
Another frequently-displayed structure is a list.
The following is an example of an \emph{itemized}
list.
\begin{itemize}
   \item This is the first item of an itemized list.
         Each item in the list is marked with a ``tick''.
         You don't have to worry about what kind of tick
         mark is used.

   \item This is the second item of the list.  It
         contains another list nested inside it.  The inner
         list is an \emph{enumerated} list.
         \begin{enumerate}
            \item This is the first item of an enumerated
                  list that is nested within the itemized list.

            \item This is the second item of the inner list.
                  \LaTeX\ allows you to nest lists deeper than
                  you really should.
         \end{enumerate}
         This is the rest of the second item of the outer
         list.  It is no more interesting than any other
         part of the item.
   \item This is the third item of the list.
\end{itemize}
You can even display poetry.
\begin{verse}
   There is an environment
    for verse \\             % The \\ command separates lines
   Whose features some poets % within a stanza.
   will curse.

                             % One or more blank lines separate stanzas.

   For instead of making\\
   Them do \emph{all} line breaking, \\
   It allows them to put too many words on a line when they'd rather be
   forced to be terse.
\end{verse}

Mathematical formulas may also be displayed.  A
displayed formula
is
one-line long; multiline
formulas require special formatting instructions.
   \[  \ip{\Gamma}{\psi'} = x'' + y^{2} + z_{i}^{n}\]
Don't start a paragraph with a displayed equation,
nor make one a paragraph by itself.

\end{document}               % End of document.

\section{Bench Huge Section 1278}
% synthetic bench section 1278
% This is a sample LaTeX input file.  (Version of 12 August 2004.)
%
% A '%' character causes TeX to ignore all remaining text on the line,
% and is used for comments like this one.

\documentclass{article}      % Specifies the document class

                             % The preamble begins here.
\title{An Example Document}  % Declares the document's title.
\author{Leslie Lamport}      % Declares the author's name.
\date{January 21, 1994}      % Deleting this command produces today's date.

\newcommand{\ip}[2]{(#1, #2)}
                             % Defines \ip{arg1}{arg2} to mean
                             % (arg1, arg2).

%\newcommand{\ip}[2]{\langle #1 | #2\rangle}
                             % This is an alternative definition of
                             % \ip that is commented out.

\begin{document}             % End of preamble and beginning of text.

\maketitle                   % Produces the title.

This is an example input file.  Comparing it with
the output it generates can show you how to
produce a simple document of your own.

\section{Ordinary Text}      % Produces section heading.  Lower-level
                             % sections are begun with similar
                             % \subsection and \subsubsection commands.

The ends  of words and sentences are marked
  by   spaces. It  doesn't matter how many
spaces    you type; one is as good as 100.  The
end of   a line counts as a space.

One   or more   blank lines denote the  end
of  a paragraph.

Since any number of consecutive spaces are treated
like a single one, the formatting of the input
file makes no difference to
      \LaTeX,                % The \LaTeX command generates the LaTeX logo.
but it makes a difference to you.  When you use
\LaTeX, making your input file as easy to read
as possible will be a great help as you write
your document and when you change it.  This sample
file shows how you can add comments to your own input
file.

Because printing is different from typewriting,
there are a number of things that you have to do
differently when preparing an input file than if
you were just typing the document directly.
Quotation marks like
       ``this''
have to be handled specially, as do quotes within
quotes:
       ``\,`this'            % \, separates the double and single quote.
        is what I just
        wrote, not  `that'\,''.

Dashes come in three sizes: an
       intra-word
dash, a medium dash for number ranges like
       1--2,
and a punctuation
       dash---like
this.

A sentence-ending space should be larger than the
space between words within a sentence.  You
sometimes have to type special commands in
conjunction with punctuation characters to get
this right, as in the following sentence.
       Gnats, gnus, etc.\ all  % `\ ' makes an inter-word space.
       begin with G\@.         % \@ marks end-of-sentence punctuation.
You should check the spaces after periods when
reading your output to make sure you haven't
forgotten any special cases.  Generating an
ellipsis
       \ldots\               % `\ ' is needed after `\ldots' because TeX
                             % ignores spaces after command names like \ldots
                             % made from \ + letters.
                             %
                             % Note how a `%' character causes TeX to ignore
                             % the end of the input line, so these blank lines
                             % do not start a new paragraph.
                             %
with the right spacing around the periods requires
a special command.

\LaTeX\ interprets some common characters as
commands, so you must type special commands to
generate them.  These characters include the
following:
       \$ \& \% \# \{ and \}.

In printing, text is usually emphasized with an
       \emph{italic}
type style.

\begin{em}
   A long segment of text can also be emphasized
   in this way.  Text within such a segment can be
   given \emph{additional} emphasis.
\end{em}

It is sometimes necessary to prevent \LaTeX\ from
breaking a line where it might otherwise do so.
This may be at a space, as between the ``Mr.''\ and
``Jones'' in
       ``Mr.~Jones'',        % ~ produces an unbreakable interword space.
or within a word---especially when the word is a
symbol like
       \mbox{\emph{itemnum}}
that makes little sense when hyphenated across
lines.

Footnotes\footnote{This is an example of a footnote.}
pose no problem.

\LaTeX\ is good at typesetting mathematical formulas
like
       \( x-3y + z = 7 \)
or
       \( a_{1} > x^{2n} + y^{2n} > x' \)
or
       \( \ip{A}{B} = \sum_{i} a_{i} b_{i} \).
The spaces you type in a formula are
ignored.  Remember that a letter like
       $x$                   % $ ... $  and  \( ... \)  are equivalent
is a formula when it denotes a mathematical
symbol, and it should be typed as one.

\section{Displayed Text}

Text is displayed by indenting it from the left
margin.  Quotations are commonly displayed.  There
are short quotations
\begin{quote}
   This is a short quotation.  It consists of a
   single paragraph of text.  See how it is formatted.
\end{quote}
and longer ones.
\begin{quotation}
   This is a longer quotation.  It consists of two
   paragraphs of text, neither of which are
   particularly interesting.

   This is the second paragraph of the quotation.  It
   is just as dull as the first paragraph.
\end{quotation}
Another frequently-displayed structure is a list.
The following is an example of an \emph{itemized}
list.
\begin{itemize}
   \item This is the first item of an itemized list.
         Each item in the list is marked with a ``tick''.
         You don't have to worry about what kind of tick
         mark is used.

   \item This is the second item of the list.  It
         contains another list nested inside it.  The inner
         list is an \emph{enumerated} list.
         \begin{enumerate}
            \item This is the first item of an enumerated
                  list that is nested within the itemized list.

            \item This is the second item of the inner list.
                  \LaTeX\ allows you to nest lists deeper than
                  you really should.
         \end{enumerate}
         This is the rest of the second item of the outer
         list.  It is no more interesting than any other
         part of the item.
   \item This is the third item of the list.
\end{itemize}
You can even display poetry.
\begin{verse}
   There is an environment
    for verse \\             % The \\ command separates lines
   Whose features some poets % within a stanza.
   will curse.

                             % One or more blank lines separate stanzas.

   For instead of making\\
   Them do \emph{all} line breaking, \\
   It allows them to put too many words on a line when they'd rather be
   forced to be terse.
\end{verse}

Mathematical formulas may also be displayed.  A
displayed formula
is
one-line long; multiline
formulas require special formatting instructions.
   \[  \ip{\Gamma}{\psi'} = x'' + y^{2} + z_{i}^{n}\]
Don't start a paragraph with a displayed equation,
nor make one a paragraph by itself.

\end{document}               % End of document.

\section{Bench Huge Section 1279}
% synthetic bench section 1279
% This is a sample LaTeX input file.  (Version of 12 August 2004.)
%
% A '%' character causes TeX to ignore all remaining text on the line,
% and is used for comments like this one.

\documentclass{article}      % Specifies the document class

                             % The preamble begins here.
\title{An Example Document}  % Declares the document's title.
\author{Leslie Lamport}      % Declares the author's name.
\date{January 21, 1994}      % Deleting this command produces today's date.

\newcommand{\ip}[2]{(#1, #2)}
                             % Defines \ip{arg1}{arg2} to mean
                             % (arg1, arg2).

%\newcommand{\ip}[2]{\langle #1 | #2\rangle}
                             % This is an alternative definition of
                             % \ip that is commented out.

\begin{document}             % End of preamble and beginning of text.

\maketitle                   % Produces the title.

This is an example input file.  Comparing it with
the output it generates can show you how to
produce a simple document of your own.

\section{Ordinary Text}      % Produces section heading.  Lower-level
                             % sections are begun with similar
                             % \subsection and \subsubsection commands.

The ends  of words and sentences are marked
  by   spaces. It  doesn't matter how many
spaces    you type; one is as good as 100.  The
end of   a line counts as a space.

One   or more   blank lines denote the  end
of  a paragraph.

Since any number of consecutive spaces are treated
like a single one, the formatting of the input
file makes no difference to
      \LaTeX,                % The \LaTeX command generates the LaTeX logo.
but it makes a difference to you.  When you use
\LaTeX, making your input file as easy to read
as possible will be a great help as you write
your document and when you change it.  This sample
file shows how you can add comments to your own input
file.

Because printing is different from typewriting,
there are a number of things that you have to do
differently when preparing an input file than if
you were just typing the document directly.
Quotation marks like
       ``this''
have to be handled specially, as do quotes within
quotes:
       ``\,`this'            % \, separates the double and single quote.
        is what I just
        wrote, not  `that'\,''.

Dashes come in three sizes: an
       intra-word
dash, a medium dash for number ranges like
       1--2,
and a punctuation
       dash---like
this.

A sentence-ending space should be larger than the
space between words within a sentence.  You
sometimes have to type special commands in
conjunction with punctuation characters to get
this right, as in the following sentence.
       Gnats, gnus, etc.\ all  % `\ ' makes an inter-word space.
       begin with G\@.         % \@ marks end-of-sentence punctuation.
You should check the spaces after periods when
reading your output to make sure you haven't
forgotten any special cases.  Generating an
ellipsis
       \ldots\               % `\ ' is needed after `\ldots' because TeX
                             % ignores spaces after command names like \ldots
                             % made from \ + letters.
                             %
                             % Note how a `%' character causes TeX to ignore
                             % the end of the input line, so these blank lines
                             % do not start a new paragraph.
                             %
with the right spacing around the periods requires
a special command.

\LaTeX\ interprets some common characters as
commands, so you must type special commands to
generate them.  These characters include the
following:
       \$ \& \% \# \{ and \}.

In printing, text is usually emphasized with an
       \emph{italic}
type style.

\begin{em}
   A long segment of text can also be emphasized
   in this way.  Text within such a segment can be
   given \emph{additional} emphasis.
\end{em}

It is sometimes necessary to prevent \LaTeX\ from
breaking a line where it might otherwise do so.
This may be at a space, as between the ``Mr.''\ and
``Jones'' in
       ``Mr.~Jones'',        % ~ produces an unbreakable interword space.
or within a word---especially when the word is a
symbol like
       \mbox{\emph{itemnum}}
that makes little sense when hyphenated across
lines.

Footnotes\footnote{This is an example of a footnote.}
pose no problem.

\LaTeX\ is good at typesetting mathematical formulas
like
       \( x-3y + z = 7 \)
or
       \( a_{1} > x^{2n} + y^{2n} > x' \)
or
       \( \ip{A}{B} = \sum_{i} a_{i} b_{i} \).
The spaces you type in a formula are
ignored.  Remember that a letter like
       $x$                   % $ ... $  and  \( ... \)  are equivalent
is a formula when it denotes a mathematical
symbol, and it should be typed as one.

\section{Displayed Text}

Text is displayed by indenting it from the left
margin.  Quotations are commonly displayed.  There
are short quotations
\begin{quote}
   This is a short quotation.  It consists of a
   single paragraph of text.  See how it is formatted.
\end{quote}
and longer ones.
\begin{quotation}
   This is a longer quotation.  It consists of two
   paragraphs of text, neither of which are
   particularly interesting.

   This is the second paragraph of the quotation.  It
   is just as dull as the first paragraph.
\end{quotation}
Another frequently-displayed structure is a list.
The following is an example of an \emph{itemized}
list.
\begin{itemize}
   \item This is the first item of an itemized list.
         Each item in the list is marked with a ``tick''.
         You don't have to worry about what kind of tick
         mark is used.

   \item This is the second item of the list.  It
         contains another list nested inside it.  The inner
         list is an \emph{enumerated} list.
         \begin{enumerate}
            \item This is the first item of an enumerated
                  list that is nested within the itemized list.

            \item This is the second item of the inner list.
                  \LaTeX\ allows you to nest lists deeper than
                  you really should.
         \end{enumerate}
         This is the rest of the second item of the outer
         list.  It is no more interesting than any other
         part of the item.
   \item This is the third item of the list.
\end{itemize}
You can even display poetry.
\begin{verse}
   There is an environment
    for verse \\             % The \\ command separates lines
   Whose features some poets % within a stanza.
   will curse.

                             % One or more blank lines separate stanzas.

   For instead of making\\
   Them do \emph{all} line breaking, \\
   It allows them to put too many words on a line when they'd rather be
   forced to be terse.
\end{verse}

Mathematical formulas may also be displayed.  A
displayed formula
is
one-line long; multiline
formulas require special formatting instructions.
   \[  \ip{\Gamma}{\psi'} = x'' + y^{2} + z_{i}^{n}\]
Don't start a paragraph with a displayed equation,
nor make one a paragraph by itself.

\end{document}               % End of document.

\section{Bench Huge Section 1280}
% synthetic bench section 1280
% This is a sample LaTeX input file.  (Version of 12 August 2004.)
%
% A '%' character causes TeX to ignore all remaining text on the line,
% and is used for comments like this one.

\documentclass{article}      % Specifies the document class

                             % The preamble begins here.
\title{An Example Document}  % Declares the document's title.
\author{Leslie Lamport}      % Declares the author's name.
\date{January 21, 1994}      % Deleting this command produces today's date.

\newcommand{\ip}[2]{(#1, #2)}
                             % Defines \ip{arg1}{arg2} to mean
                             % (arg1, arg2).

%\newcommand{\ip}[2]{\langle #1 | #2\rangle}
                             % This is an alternative definition of
                             % \ip that is commented out.

\begin{document}             % End of preamble and beginning of text.

\maketitle                   % Produces the title.

This is an example input file.  Comparing it with
the output it generates can show you how to
produce a simple document of your own.

\section{Ordinary Text}      % Produces section heading.  Lower-level
                             % sections are begun with similar
                             % \subsection and \subsubsection commands.

The ends  of words and sentences are marked
  by   spaces. It  doesn't matter how many
spaces    you type; one is as good as 100.  The
end of   a line counts as a space.

One   or more   blank lines denote the  end
of  a paragraph.

Since any number of consecutive spaces are treated
like a single one, the formatting of the input
file makes no difference to
      \LaTeX,                % The \LaTeX command generates the LaTeX logo.
but it makes a difference to you.  When you use
\LaTeX, making your input file as easy to read
as possible will be a great help as you write
your document and when you change it.  This sample
file shows how you can add comments to your own input
file.

Because printing is different from typewriting,
there are a number of things that you have to do
differently when preparing an input file than if
you were just typing the document directly.
Quotation marks like
       ``this''
have to be handled specially, as do quotes within
quotes:
       ``\,`this'            % \, separates the double and single quote.
        is what I just
        wrote, not  `that'\,''.

Dashes come in three sizes: an
       intra-word
dash, a medium dash for number ranges like
       1--2,
and a punctuation
       dash---like
this.

A sentence-ending space should be larger than the
space between words within a sentence.  You
sometimes have to type special commands in
conjunction with punctuation characters to get
this right, as in the following sentence.
       Gnats, gnus, etc.\ all  % `\ ' makes an inter-word space.
       begin with G\@.         % \@ marks end-of-sentence punctuation.
You should check the spaces after periods when
reading your output to make sure you haven't
forgotten any special cases.  Generating an
ellipsis
       \ldots\               % `\ ' is needed after `\ldots' because TeX
                             % ignores spaces after command names like \ldots
                             % made from \ + letters.
                             %
                             % Note how a `%' character causes TeX to ignore
                             % the end of the input line, so these blank lines
                             % do not start a new paragraph.
                             %
with the right spacing around the periods requires
a special command.

\LaTeX\ interprets some common characters as
commands, so you must type special commands to
generate them.  These characters include the
following:
       \$ \& \% \# \{ and \}.

In printing, text is usually emphasized with an
       \emph{italic}
type style.

\begin{em}
   A long segment of text can also be emphasized
   in this way.  Text within such a segment can be
   given \emph{additional} emphasis.
\end{em}

It is sometimes necessary to prevent \LaTeX\ from
breaking a line where it might otherwise do so.
This may be at a space, as between the ``Mr.''\ and
``Jones'' in
       ``Mr.~Jones'',        % ~ produces an unbreakable interword space.
or within a word---especially when the word is a
symbol like
       \mbox{\emph{itemnum}}
that makes little sense when hyphenated across
lines.

Footnotes\footnote{This is an example of a footnote.}
pose no problem.

\LaTeX\ is good at typesetting mathematical formulas
like
       \( x-3y + z = 7 \)
or
       \( a_{1} > x^{2n} + y^{2n} > x' \)
or
       \( \ip{A}{B} = \sum_{i} a_{i} b_{i} \).
The spaces you type in a formula are
ignored.  Remember that a letter like
       $x$                   % $ ... $  and  \( ... \)  are equivalent
is a formula when it denotes a mathematical
symbol, and it should be typed as one.

\section{Displayed Text}

Text is displayed by indenting it from the left
margin.  Quotations are commonly displayed.  There
are short quotations
\begin{quote}
   This is a short quotation.  It consists of a
   single paragraph of text.  See how it is formatted.
\end{quote}
and longer ones.
\begin{quotation}
   This is a longer quotation.  It consists of two
   paragraphs of text, neither of which are
   particularly interesting.

   This is the second paragraph of the quotation.  It
   is just as dull as the first paragraph.
\end{quotation}
Another frequently-displayed structure is a list.
The following is an example of an \emph{itemized}
list.
\begin{itemize}
   \item This is the first item of an itemized list.
         Each item in the list is marked with a ``tick''.
         You don't have to worry about what kind of tick
         mark is used.

   \item This is the second item of the list.  It
         contains another list nested inside it.  The inner
         list is an \emph{enumerated} list.
         \begin{enumerate}
            \item This is the first item of an enumerated
                  list that is nested within the itemized list.

            \item This is the second item of the inner list.
                  \LaTeX\ allows you to nest lists deeper than
                  you really should.
         \end{enumerate}
         This is the rest of the second item of the outer
         list.  It is no more interesting than any other
         part of the item.
   \item This is the third item of the list.
\end{itemize}
You can even display poetry.
\begin{verse}
   There is an environment
    for verse \\             % The \\ command separates lines
   Whose features some poets % within a stanza.
   will curse.

                             % One or more blank lines separate stanzas.

   For instead of making\\
   Them do \emph{all} line breaking, \\
   It allows them to put too many words on a line when they'd rather be
   forced to be terse.
\end{verse}

Mathematical formulas may also be displayed.  A
displayed formula
is
one-line long; multiline
formulas require special formatting instructions.
   \[  \ip{\Gamma}{\psi'} = x'' + y^{2} + z_{i}^{n}\]
Don't start a paragraph with a displayed equation,
nor make one a paragraph by itself.

\end{document}               % End of document.

\section{Bench Huge Section 1281}
% synthetic bench section 1281
% This is a sample LaTeX input file.  (Version of 12 August 2004.)
%
% A '%' character causes TeX to ignore all remaining text on the line,
% and is used for comments like this one.

\documentclass{article}      % Specifies the document class

                             % The preamble begins here.
\title{An Example Document}  % Declares the document's title.
\author{Leslie Lamport}      % Declares the author's name.
\date{January 21, 1994}      % Deleting this command produces today's date.

\newcommand{\ip}[2]{(#1, #2)}
                             % Defines \ip{arg1}{arg2} to mean
                             % (arg1, arg2).

%\newcommand{\ip}[2]{\langle #1 | #2\rangle}
                             % This is an alternative definition of
                             % \ip that is commented out.

\begin{document}             % End of preamble and beginning of text.

\maketitle                   % Produces the title.

This is an example input file.  Comparing it with
the output it generates can show you how to
produce a simple document of your own.

\section{Ordinary Text}      % Produces section heading.  Lower-level
                             % sections are begun with similar
                             % \subsection and \subsubsection commands.

The ends  of words and sentences are marked
  by   spaces. It  doesn't matter how many
spaces    you type; one is as good as 100.  The
end of   a line counts as a space.

One   or more   blank lines denote the  end
of  a paragraph.

Since any number of consecutive spaces are treated
like a single one, the formatting of the input
file makes no difference to
      \LaTeX,                % The \LaTeX command generates the LaTeX logo.
but it makes a difference to you.  When you use
\LaTeX, making your input file as easy to read
as possible will be a great help as you write
your document and when you change it.  This sample
file shows how you can add comments to your own input
file.

Because printing is different from typewriting,
there are a number of things that you have to do
differently when preparing an input file than if
you were just typing the document directly.
Quotation marks like
       ``this''
have to be handled specially, as do quotes within
quotes:
       ``\,`this'            % \, separates the double and single quote.
        is what I just
        wrote, not  `that'\,''.

Dashes come in three sizes: an
       intra-word
dash, a medium dash for number ranges like
       1--2,
and a punctuation
       dash---like
this.

A sentence-ending space should be larger than the
space between words within a sentence.  You
sometimes have to type special commands in
conjunction with punctuation characters to get
this right, as in the following sentence.
       Gnats, gnus, etc.\ all  % `\ ' makes an inter-word space.
       begin with G\@.         % \@ marks end-of-sentence punctuation.
You should check the spaces after periods when
reading your output to make sure you haven't
forgotten any special cases.  Generating an
ellipsis
       \ldots\               % `\ ' is needed after `\ldots' because TeX
                             % ignores spaces after command names like \ldots
                             % made from \ + letters.
                             %
                             % Note how a `%' character causes TeX to ignore
                             % the end of the input line, so these blank lines
                             % do not start a new paragraph.
                             %
with the right spacing around the periods requires
a special command.

\LaTeX\ interprets some common characters as
commands, so you must type special commands to
generate them.  These characters include the
following:
       \$ \& \% \# \{ and \}.

In printing, text is usually emphasized with an
       \emph{italic}
type style.

\begin{em}
   A long segment of text can also be emphasized
   in this way.  Text within such a segment can be
   given \emph{additional} emphasis.
\end{em}

It is sometimes necessary to prevent \LaTeX\ from
breaking a line where it might otherwise do so.
This may be at a space, as between the ``Mr.''\ and
``Jones'' in
       ``Mr.~Jones'',        % ~ produces an unbreakable interword space.
or within a word---especially when the word is a
symbol like
       \mbox{\emph{itemnum}}
that makes little sense when hyphenated across
lines.

Footnotes\footnote{This is an example of a footnote.}
pose no problem.

\LaTeX\ is good at typesetting mathematical formulas
like
       \( x-3y + z = 7 \)
or
       \( a_{1} > x^{2n} + y^{2n} > x' \)
or
       \( \ip{A}{B} = \sum_{i} a_{i} b_{i} \).
The spaces you type in a formula are
ignored.  Remember that a letter like
       $x$                   % $ ... $  and  \( ... \)  are equivalent
is a formula when it denotes a mathematical
symbol, and it should be typed as one.

\section{Displayed Text}

Text is displayed by indenting it from the left
margin.  Quotations are commonly displayed.  There
are short quotations
\begin{quote}
   This is a short quotation.  It consists of a
   single paragraph of text.  See how it is formatted.
\end{quote}
and longer ones.
\begin{quotation}
   This is a longer quotation.  It consists of two
   paragraphs of text, neither of which are
   particularly interesting.

   This is the second paragraph of the quotation.  It
   is just as dull as the first paragraph.
\end{quotation}
Another frequently-displayed structure is a list.
The following is an example of an \emph{itemized}
list.
\begin{itemize}
   \item This is the first item of an itemized list.
         Each item in the list is marked with a ``tick''.
         You don't have to worry about what kind of tick
         mark is used.

   \item This is the second item of the list.  It
         contains another list nested inside it.  The inner
         list is an \emph{enumerated} list.
         \begin{enumerate}
            \item This is the first item of an enumerated
                  list that is nested within the itemized list.

            \item This is the second item of the inner list.
                  \LaTeX\ allows you to nest lists deeper than
                  you really should.
         \end{enumerate}
         This is the rest of the second item of the outer
         list.  It is no more interesting than any other
         part of the item.
   \item This is the third item of the list.
\end{itemize}
You can even display poetry.
\begin{verse}
   There is an environment
    for verse \\             % The \\ command separates lines
   Whose features some poets % within a stanza.
   will curse.

                             % One or more blank lines separate stanzas.

   For instead of making\\
   Them do \emph{all} line breaking, \\
   It allows them to put too many words on a line when they'd rather be
   forced to be terse.
\end{verse}

Mathematical formulas may also be displayed.  A
displayed formula
is
one-line long; multiline
formulas require special formatting instructions.
   \[  \ip{\Gamma}{\psi'} = x'' + y^{2} + z_{i}^{n}\]
Don't start a paragraph with a displayed equation,
nor make one a paragraph by itself.

\end{document}               % End of document.

\section{Bench Huge Section 1282}
% synthetic bench section 1282
% This is a sample LaTeX input file.  (Version of 12 August 2004.)
%
% A '%' character causes TeX to ignore all remaining text on the line,
% and is used for comments like this one.

\documentclass{article}      % Specifies the document class

                             % The preamble begins here.
\title{An Example Document}  % Declares the document's title.
\author{Leslie Lamport}      % Declares the author's name.
\date{January 21, 1994}      % Deleting this command produces today's date.

\newcommand{\ip}[2]{(#1, #2)}
                             % Defines \ip{arg1}{arg2} to mean
                             % (arg1, arg2).

%\newcommand{\ip}[2]{\langle #1 | #2\rangle}
                             % This is an alternative definition of
                             % \ip that is commented out.

\begin{document}             % End of preamble and beginning of text.

\maketitle                   % Produces the title.

This is an example input file.  Comparing it with
the output it generates can show you how to
produce a simple document of your own.

\section{Ordinary Text}      % Produces section heading.  Lower-level
                             % sections are begun with similar
                             % \subsection and \subsubsection commands.

The ends  of words and sentences are marked
  by   spaces. It  doesn't matter how many
spaces    you type; one is as good as 100.  The
end of   a line counts as a space.

One   or more   blank lines denote the  end
of  a paragraph.

Since any number of consecutive spaces are treated
like a single one, the formatting of the input
file makes no difference to
      \LaTeX,                % The \LaTeX command generates the LaTeX logo.
but it makes a difference to you.  When you use
\LaTeX, making your input file as easy to read
as possible will be a great help as you write
your document and when you change it.  This sample
file shows how you can add comments to your own input
file.

Because printing is different from typewriting,
there are a number of things that you have to do
differently when preparing an input file than if
you were just typing the document directly.
Quotation marks like
       ``this''
have to be handled specially, as do quotes within
quotes:
       ``\,`this'            % \, separates the double and single quote.
        is what I just
        wrote, not  `that'\,''.

Dashes come in three sizes: an
       intra-word
dash, a medium dash for number ranges like
       1--2,
and a punctuation
       dash---like
this.

A sentence-ending space should be larger than the
space between words within a sentence.  You
sometimes have to type special commands in
conjunction with punctuation characters to get
this right, as in the following sentence.
       Gnats, gnus, etc.\ all  % `\ ' makes an inter-word space.
       begin with G\@.         % \@ marks end-of-sentence punctuation.
You should check the spaces after periods when
reading your output to make sure you haven't
forgotten any special cases.  Generating an
ellipsis
       \ldots\               % `\ ' is needed after `\ldots' because TeX
                             % ignores spaces after command names like \ldots
                             % made from \ + letters.
                             %
                             % Note how a `%' character causes TeX to ignore
                             % the end of the input line, so these blank lines
                             % do not start a new paragraph.
                             %
with the right spacing around the periods requires
a special command.

\LaTeX\ interprets some common characters as
commands, so you must type special commands to
generate them.  These characters include the
following:
       \$ \& \% \# \{ and \}.

In printing, text is usually emphasized with an
       \emph{italic}
type style.

\begin{em}
   A long segment of text can also be emphasized
   in this way.  Text within such a segment can be
   given \emph{additional} emphasis.
\end{em}

It is sometimes necessary to prevent \LaTeX\ from
breaking a line where it might otherwise do so.
This may be at a space, as between the ``Mr.''\ and
``Jones'' in
       ``Mr.~Jones'',        % ~ produces an unbreakable interword space.
or within a word---especially when the word is a
symbol like
       \mbox{\emph{itemnum}}
that makes little sense when hyphenated across
lines.

Footnotes\footnote{This is an example of a footnote.}
pose no problem.

\LaTeX\ is good at typesetting mathematical formulas
like
       \( x-3y + z = 7 \)
or
       \( a_{1} > x^{2n} + y^{2n} > x' \)
or
       \( \ip{A}{B} = \sum_{i} a_{i} b_{i} \).
The spaces you type in a formula are
ignored.  Remember that a letter like
       $x$                   % $ ... $  and  \( ... \)  are equivalent
is a formula when it denotes a mathematical
symbol, and it should be typed as one.

\section{Displayed Text}

Text is displayed by indenting it from the left
margin.  Quotations are commonly displayed.  There
are short quotations
\begin{quote}
   This is a short quotation.  It consists of a
   single paragraph of text.  See how it is formatted.
\end{quote}
and longer ones.
\begin{quotation}
   This is a longer quotation.  It consists of two
   paragraphs of text, neither of which are
   particularly interesting.

   This is the second paragraph of the quotation.  It
   is just as dull as the first paragraph.
\end{quotation}
Another frequently-displayed structure is a list.
The following is an example of an \emph{itemized}
list.
\begin{itemize}
   \item This is the first item of an itemized list.
         Each item in the list is marked with a ``tick''.
         You don't have to worry about what kind of tick
         mark is used.

   \item This is the second item of the list.  It
         contains another list nested inside it.  The inner
         list is an \emph{enumerated} list.
         \begin{enumerate}
            \item This is the first item of an enumerated
                  list that is nested within the itemized list.

            \item This is the second item of the inner list.
                  \LaTeX\ allows you to nest lists deeper than
                  you really should.
         \end{enumerate}
         This is the rest of the second item of the outer
         list.  It is no more interesting than any other
         part of the item.
   \item This is the third item of the list.
\end{itemize}
You can even display poetry.
\begin{verse}
   There is an environment
    for verse \\             % The \\ command separates lines
   Whose features some poets % within a stanza.
   will curse.

                             % One or more blank lines separate stanzas.

   For instead of making\\
   Them do \emph{all} line breaking, \\
   It allows them to put too many words on a line when they'd rather be
   forced to be terse.
\end{verse}

Mathematical formulas may also be displayed.  A
displayed formula
is
one-line long; multiline
formulas require special formatting instructions.
   \[  \ip{\Gamma}{\psi'} = x'' + y^{2} + z_{i}^{n}\]
Don't start a paragraph with a displayed equation,
nor make one a paragraph by itself.

\end{document}               % End of document.

\section{Bench Huge Section 1283}
% synthetic bench section 1283
% This is a sample LaTeX input file.  (Version of 12 August 2004.)
%
% A '%' character causes TeX to ignore all remaining text on the line,
% and is used for comments like this one.

\documentclass{article}      % Specifies the document class

                             % The preamble begins here.
\title{An Example Document}  % Declares the document's title.
\author{Leslie Lamport}      % Declares the author's name.
\date{January 21, 1994}      % Deleting this command produces today's date.

\newcommand{\ip}[2]{(#1, #2)}
                             % Defines \ip{arg1}{arg2} to mean
                             % (arg1, arg2).

%\newcommand{\ip}[2]{\langle #1 | #2\rangle}
                             % This is an alternative definition of
                             % \ip that is commented out.

\begin{document}             % End of preamble and beginning of text.

\maketitle                   % Produces the title.

This is an example input file.  Comparing it with
the output it generates can show you how to
produce a simple document of your own.

\section{Ordinary Text}      % Produces section heading.  Lower-level
                             % sections are begun with similar
                             % \subsection and \subsubsection commands.

The ends  of words and sentences are marked
  by   spaces. It  doesn't matter how many
spaces    you type; one is as good as 100.  The
end of   a line counts as a space.

One   or more   blank lines denote the  end
of  a paragraph.

Since any number of consecutive spaces are treated
like a single one, the formatting of the input
file makes no difference to
      \LaTeX,                % The \LaTeX command generates the LaTeX logo.
but it makes a difference to you.  When you use
\LaTeX, making your input file as easy to read
as possible will be a great help as you write
your document and when you change it.  This sample
file shows how you can add comments to your own input
file.

Because printing is different from typewriting,
there are a number of things that you have to do
differently when preparing an input file than if
you were just typing the document directly.
Quotation marks like
       ``this''
have to be handled specially, as do quotes within
quotes:
       ``\,`this'            % \, separates the double and single quote.
        is what I just
        wrote, not  `that'\,''.

Dashes come in three sizes: an
       intra-word
dash, a medium dash for number ranges like
       1--2,
and a punctuation
       dash---like
this.

A sentence-ending space should be larger than the
space between words within a sentence.  You
sometimes have to type special commands in
conjunction with punctuation characters to get
this right, as in the following sentence.
       Gnats, gnus, etc.\ all  % `\ ' makes an inter-word space.
       begin with G\@.         % \@ marks end-of-sentence punctuation.
You should check the spaces after periods when
reading your output to make sure you haven't
forgotten any special cases.  Generating an
ellipsis
       \ldots\               % `\ ' is needed after `\ldots' because TeX
                             % ignores spaces after command names like \ldots
                             % made from \ + letters.
                             %
                             % Note how a `%' character causes TeX to ignore
                             % the end of the input line, so these blank lines
                             % do not start a new paragraph.
                             %
with the right spacing around the periods requires
a special command.

\LaTeX\ interprets some common characters as
commands, so you must type special commands to
generate them.  These characters include the
following:
       \$ \& \% \# \{ and \}.

In printing, text is usually emphasized with an
       \emph{italic}
type style.

\begin{em}
   A long segment of text can also be emphasized
   in this way.  Text within such a segment can be
   given \emph{additional} emphasis.
\end{em}

It is sometimes necessary to prevent \LaTeX\ from
breaking a line where it might otherwise do so.
This may be at a space, as between the ``Mr.''\ and
``Jones'' in
       ``Mr.~Jones'',        % ~ produces an unbreakable interword space.
or within a word---especially when the word is a
symbol like
       \mbox{\emph{itemnum}}
that makes little sense when hyphenated across
lines.

Footnotes\footnote{This is an example of a footnote.}
pose no problem.

\LaTeX\ is good at typesetting mathematical formulas
like
       \( x-3y + z = 7 \)
or
       \( a_{1} > x^{2n} + y^{2n} > x' \)
or
       \( \ip{A}{B} = \sum_{i} a_{i} b_{i} \).
The spaces you type in a formula are
ignored.  Remember that a letter like
       $x$                   % $ ... $  and  \( ... \)  are equivalent
is a formula when it denotes a mathematical
symbol, and it should be typed as one.

\section{Displayed Text}

Text is displayed by indenting it from the left
margin.  Quotations are commonly displayed.  There
are short quotations
\begin{quote}
   This is a short quotation.  It consists of a
   single paragraph of text.  See how it is formatted.
\end{quote}
and longer ones.
\begin{quotation}
   This is a longer quotation.  It consists of two
   paragraphs of text, neither of which are
   particularly interesting.

   This is the second paragraph of the quotation.  It
   is just as dull as the first paragraph.
\end{quotation}
Another frequently-displayed structure is a list.
The following is an example of an \emph{itemized}
list.
\begin{itemize}
   \item This is the first item of an itemized list.
         Each item in the list is marked with a ``tick''.
         You don't have to worry about what kind of tick
         mark is used.

   \item This is the second item of the list.  It
         contains another list nested inside it.  The inner
         list is an \emph{enumerated} list.
         \begin{enumerate}
            \item This is the first item of an enumerated
                  list that is nested within the itemized list.

            \item This is the second item of the inner list.
                  \LaTeX\ allows you to nest lists deeper than
                  you really should.
         \end{enumerate}
         This is the rest of the second item of the outer
         list.  It is no more interesting than any other
         part of the item.
   \item This is the third item of the list.
\end{itemize}
You can even display poetry.
\begin{verse}
   There is an environment
    for verse \\             % The \\ command separates lines
   Whose features some poets % within a stanza.
   will curse.

                             % One or more blank lines separate stanzas.

   For instead of making\\
   Them do \emph{all} line breaking, \\
   It allows them to put too many words on a line when they'd rather be
   forced to be terse.
\end{verse}

Mathematical formulas may also be displayed.  A
displayed formula
is
one-line long; multiline
formulas require special formatting instructions.
   \[  \ip{\Gamma}{\psi'} = x'' + y^{2} + z_{i}^{n}\]
Don't start a paragraph with a displayed equation,
nor make one a paragraph by itself.

\end{document}               % End of document.

\section{Bench Huge Section 1284}
% synthetic bench section 1284
% This is a sample LaTeX input file.  (Version of 12 August 2004.)
%
% A '%' character causes TeX to ignore all remaining text on the line,
% and is used for comments like this one.

\documentclass{article}      % Specifies the document class

                             % The preamble begins here.
\title{An Example Document}  % Declares the document's title.
\author{Leslie Lamport}      % Declares the author's name.
\date{January 21, 1994}      % Deleting this command produces today's date.

\newcommand{\ip}[2]{(#1, #2)}
                             % Defines \ip{arg1}{arg2} to mean
                             % (arg1, arg2).

%\newcommand{\ip}[2]{\langle #1 | #2\rangle}
                             % This is an alternative definition of
                             % \ip that is commented out.

\begin{document}             % End of preamble and beginning of text.

\maketitle                   % Produces the title.

This is an example input file.  Comparing it with
the output it generates can show you how to
produce a simple document of your own.

\section{Ordinary Text}      % Produces section heading.  Lower-level
                             % sections are begun with similar
                             % \subsection and \subsubsection commands.

The ends  of words and sentences are marked
  by   spaces. It  doesn't matter how many
spaces    you type; one is as good as 100.  The
end of   a line counts as a space.

One   or more   blank lines denote the  end
of  a paragraph.

Since any number of consecutive spaces are treated
like a single one, the formatting of the input
file makes no difference to
      \LaTeX,                % The \LaTeX command generates the LaTeX logo.
but it makes a difference to you.  When you use
\LaTeX, making your input file as easy to read
as possible will be a great help as you write
your document and when you change it.  This sample
file shows how you can add comments to your own input
file.

Because printing is different from typewriting,
there are a number of things that you have to do
differently when preparing an input file than if
you were just typing the document directly.
Quotation marks like
       ``this''
have to be handled specially, as do quotes within
quotes:
       ``\,`this'            % \, separates the double and single quote.
        is what I just
        wrote, not  `that'\,''.

Dashes come in three sizes: an
       intra-word
dash, a medium dash for number ranges like
       1--2,
and a punctuation
       dash---like
this.

A sentence-ending space should be larger than the
space between words within a sentence.  You
sometimes have to type special commands in
conjunction with punctuation characters to get
this right, as in the following sentence.
       Gnats, gnus, etc.\ all  % `\ ' makes an inter-word space.
       begin with G\@.         % \@ marks end-of-sentence punctuation.
You should check the spaces after periods when
reading your output to make sure you haven't
forgotten any special cases.  Generating an
ellipsis
       \ldots\               % `\ ' is needed after `\ldots' because TeX
                             % ignores spaces after command names like \ldots
                             % made from \ + letters.
                             %
                             % Note how a `%' character causes TeX to ignore
                             % the end of the input line, so these blank lines
                             % do not start a new paragraph.
                             %
with the right spacing around the periods requires
a special command.

\LaTeX\ interprets some common characters as
commands, so you must type special commands to
generate them.  These characters include the
following:
       \$ \& \% \# \{ and \}.

In printing, text is usually emphasized with an
       \emph{italic}
type style.

\begin{em}
   A long segment of text can also be emphasized
   in this way.  Text within such a segment can be
   given \emph{additional} emphasis.
\end{em}

It is sometimes necessary to prevent \LaTeX\ from
breaking a line where it might otherwise do so.
This may be at a space, as between the ``Mr.''\ and
``Jones'' in
       ``Mr.~Jones'',        % ~ produces an unbreakable interword space.
or within a word---especially when the word is a
symbol like
       \mbox{\emph{itemnum}}
that makes little sense when hyphenated across
lines.

Footnotes\footnote{This is an example of a footnote.}
pose no problem.

\LaTeX\ is good at typesetting mathematical formulas
like
       \( x-3y + z = 7 \)
or
       \( a_{1} > x^{2n} + y^{2n} > x' \)
or
       \( \ip{A}{B} = \sum_{i} a_{i} b_{i} \).
The spaces you type in a formula are
ignored.  Remember that a letter like
       $x$                   % $ ... $  and  \( ... \)  are equivalent
is a formula when it denotes a mathematical
symbol, and it should be typed as one.

\section{Displayed Text}

Text is displayed by indenting it from the left
margin.  Quotations are commonly displayed.  There
are short quotations
\begin{quote}
   This is a short quotation.  It consists of a
   single paragraph of text.  See how it is formatted.
\end{quote}
and longer ones.
\begin{quotation}
   This is a longer quotation.  It consists of two
   paragraphs of text, neither of which are
   particularly interesting.

   This is the second paragraph of the quotation.  It
   is just as dull as the first paragraph.
\end{quotation}
Another frequently-displayed structure is a list.
The following is an example of an \emph{itemized}
list.
\begin{itemize}
   \item This is the first item of an itemized list.
         Each item in the list is marked with a ``tick''.
         You don't have to worry about what kind of tick
         mark is used.

   \item This is the second item of the list.  It
         contains another list nested inside it.  The inner
         list is an \emph{enumerated} list.
         \begin{enumerate}
            \item This is the first item of an enumerated
                  list that is nested within the itemized list.

            \item This is the second item of the inner list.
                  \LaTeX\ allows you to nest lists deeper than
                  you really should.
         \end{enumerate}
         This is the rest of the second item of the outer
         list.  It is no more interesting than any other
         part of the item.
   \item This is the third item of the list.
\end{itemize}
You can even display poetry.
\begin{verse}
   There is an environment
    for verse \\             % The \\ command separates lines
   Whose features some poets % within a stanza.
   will curse.

                             % One or more blank lines separate stanzas.

   For instead of making\\
   Them do \emph{all} line breaking, \\
   It allows them to put too many words on a line when they'd rather be
   forced to be terse.
\end{verse}

Mathematical formulas may also be displayed.  A
displayed formula
is
one-line long; multiline
formulas require special formatting instructions.
   \[  \ip{\Gamma}{\psi'} = x'' + y^{2} + z_{i}^{n}\]
Don't start a paragraph with a displayed equation,
nor make one a paragraph by itself.

\end{document}               % End of document.

\section{Bench Huge Section 1285}
% synthetic bench section 1285
% This is a sample LaTeX input file.  (Version of 12 August 2004.)
%
% A '%' character causes TeX to ignore all remaining text on the line,
% and is used for comments like this one.

\documentclass{article}      % Specifies the document class

                             % The preamble begins here.
\title{An Example Document}  % Declares the document's title.
\author{Leslie Lamport}      % Declares the author's name.
\date{January 21, 1994}      % Deleting this command produces today's date.

\newcommand{\ip}[2]{(#1, #2)}
                             % Defines \ip{arg1}{arg2} to mean
                             % (arg1, arg2).

%\newcommand{\ip}[2]{\langle #1 | #2\rangle}
                             % This is an alternative definition of
                             % \ip that is commented out.

\begin{document}             % End of preamble and beginning of text.

\maketitle                   % Produces the title.

This is an example input file.  Comparing it with
the output it generates can show you how to
produce a simple document of your own.

\section{Ordinary Text}      % Produces section heading.  Lower-level
                             % sections are begun with similar
                             % \subsection and \subsubsection commands.

The ends  of words and sentences are marked
  by   spaces. It  doesn't matter how many
spaces    you type; one is as good as 100.  The
end of   a line counts as a space.

One   or more   blank lines denote the  end
of  a paragraph.

Since any number of consecutive spaces are treated
like a single one, the formatting of the input
file makes no difference to
      \LaTeX,                % The \LaTeX command generates the LaTeX logo.
but it makes a difference to you.  When you use
\LaTeX, making your input file as easy to read
as possible will be a great help as you write
your document and when you change it.  This sample
file shows how you can add comments to your own input
file.

Because printing is different from typewriting,
there are a number of things that you have to do
differently when preparing an input file than if
you were just typing the document directly.
Quotation marks like
       ``this''
have to be handled specially, as do quotes within
quotes:
       ``\,`this'            % \, separates the double and single quote.
        is what I just
        wrote, not  `that'\,''.

Dashes come in three sizes: an
       intra-word
dash, a medium dash for number ranges like
       1--2,
and a punctuation
       dash---like
this.

A sentence-ending space should be larger than the
space between words within a sentence.  You
sometimes have to type special commands in
conjunction with punctuation characters to get
this right, as in the following sentence.
       Gnats, gnus, etc.\ all  % `\ ' makes an inter-word space.
       begin with G\@.         % \@ marks end-of-sentence punctuation.
You should check the spaces after periods when
reading your output to make sure you haven't
forgotten any special cases.  Generating an
ellipsis
       \ldots\               % `\ ' is needed after `\ldots' because TeX
                             % ignores spaces after command names like \ldots
                             % made from \ + letters.
                             %
                             % Note how a `%' character causes TeX to ignore
                             % the end of the input line, so these blank lines
                             % do not start a new paragraph.
                             %
with the right spacing around the periods requires
a special command.

\LaTeX\ interprets some common characters as
commands, so you must type special commands to
generate them.  These characters include the
following:
       \$ \& \% \# \{ and \}.

In printing, text is usually emphasized with an
       \emph{italic}
type style.

\begin{em}
   A long segment of text can also be emphasized
   in this way.  Text within such a segment can be
   given \emph{additional} emphasis.
\end{em}

It is sometimes necessary to prevent \LaTeX\ from
breaking a line where it might otherwise do so.
This may be at a space, as between the ``Mr.''\ and
``Jones'' in
       ``Mr.~Jones'',        % ~ produces an unbreakable interword space.
or within a word---especially when the word is a
symbol like
       \mbox{\emph{itemnum}}
that makes little sense when hyphenated across
lines.

Footnotes\footnote{This is an example of a footnote.}
pose no problem.

\LaTeX\ is good at typesetting mathematical formulas
like
       \( x-3y + z = 7 \)
or
       \( a_{1} > x^{2n} + y^{2n} > x' \)
or
       \( \ip{A}{B} = \sum_{i} a_{i} b_{i} \).
The spaces you type in a formula are
ignored.  Remember that a letter like
       $x$                   % $ ... $  and  \( ... \)  are equivalent
is a formula when it denotes a mathematical
symbol, and it should be typed as one.

\section{Displayed Text}

Text is displayed by indenting it from the left
margin.  Quotations are commonly displayed.  There
are short quotations
\begin{quote}
   This is a short quotation.  It consists of a
   single paragraph of text.  See how it is formatted.
\end{quote}
and longer ones.
\begin{quotation}
   This is a longer quotation.  It consists of two
   paragraphs of text, neither of which are
   particularly interesting.

   This is the second paragraph of the quotation.  It
   is just as dull as the first paragraph.
\end{quotation}
Another frequently-displayed structure is a list.
The following is an example of an \emph{itemized}
list.
\begin{itemize}
   \item This is the first item of an itemized list.
         Each item in the list is marked with a ``tick''.
         You don't have to worry about what kind of tick
         mark is used.

   \item This is the second item of the list.  It
         contains another list nested inside it.  The inner
         list is an \emph{enumerated} list.
         \begin{enumerate}
            \item This is the first item of an enumerated
                  list that is nested within the itemized list.

            \item This is the second item of the inner list.
                  \LaTeX\ allows you to nest lists deeper than
                  you really should.
         \end{enumerate}
         This is the rest of the second item of the outer
         list.  It is no more interesting than any other
         part of the item.
   \item This is the third item of the list.
\end{itemize}
You can even display poetry.
\begin{verse}
   There is an environment
    for verse \\             % The \\ command separates lines
   Whose features some poets % within a stanza.
   will curse.

                             % One or more blank lines separate stanzas.

   For instead of making\\
   Them do \emph{all} line breaking, \\
   It allows them to put too many words on a line when they'd rather be
   forced to be terse.
\end{verse}

Mathematical formulas may also be displayed.  A
displayed formula
is
one-line long; multiline
formulas require special formatting instructions.
   \[  \ip{\Gamma}{\psi'} = x'' + y^{2} + z_{i}^{n}\]
Don't start a paragraph with a displayed equation,
nor make one a paragraph by itself.

\end{document}               % End of document.

\section{Bench Huge Section 1286}
% synthetic bench section 1286
% This is a sample LaTeX input file.  (Version of 12 August 2004.)
%
% A '%' character causes TeX to ignore all remaining text on the line,
% and is used for comments like this one.

\documentclass{article}      % Specifies the document class

                             % The preamble begins here.
\title{An Example Document}  % Declares the document's title.
\author{Leslie Lamport}      % Declares the author's name.
\date{January 21, 1994}      % Deleting this command produces today's date.

\newcommand{\ip}[2]{(#1, #2)}
                             % Defines \ip{arg1}{arg2} to mean
                             % (arg1, arg2).

%\newcommand{\ip}[2]{\langle #1 | #2\rangle}
                             % This is an alternative definition of
                             % \ip that is commented out.

\begin{document}             % End of preamble and beginning of text.

\maketitle                   % Produces the title.

This is an example input file.  Comparing it with
the output it generates can show you how to
produce a simple document of your own.

\section{Ordinary Text}      % Produces section heading.  Lower-level
                             % sections are begun with similar
                             % \subsection and \subsubsection commands.

The ends  of words and sentences are marked
  by   spaces. It  doesn't matter how many
spaces    you type; one is as good as 100.  The
end of   a line counts as a space.

One   or more   blank lines denote the  end
of  a paragraph.

Since any number of consecutive spaces are treated
like a single one, the formatting of the input
file makes no difference to
      \LaTeX,                % The \LaTeX command generates the LaTeX logo.
but it makes a difference to you.  When you use
\LaTeX, making your input file as easy to read
as possible will be a great help as you write
your document and when you change it.  This sample
file shows how you can add comments to your own input
file.

Because printing is different from typewriting,
there are a number of things that you have to do
differently when preparing an input file than if
you were just typing the document directly.
Quotation marks like
       ``this''
have to be handled specially, as do quotes within
quotes:
       ``\,`this'            % \, separates the double and single quote.
        is what I just
        wrote, not  `that'\,''.

Dashes come in three sizes: an
       intra-word
dash, a medium dash for number ranges like
       1--2,
and a punctuation
       dash---like
this.

A sentence-ending space should be larger than the
space between words within a sentence.  You
sometimes have to type special commands in
conjunction with punctuation characters to get
this right, as in the following sentence.
       Gnats, gnus, etc.\ all  % `\ ' makes an inter-word space.
       begin with G\@.         % \@ marks end-of-sentence punctuation.
You should check the spaces after periods when
reading your output to make sure you haven't
forgotten any special cases.  Generating an
ellipsis
       \ldots\               % `\ ' is needed after `\ldots' because TeX
                             % ignores spaces after command names like \ldots
                             % made from \ + letters.
                             %
                             % Note how a `%' character causes TeX to ignore
                             % the end of the input line, so these blank lines
                             % do not start a new paragraph.
                             %
with the right spacing around the periods requires
a special command.

\LaTeX\ interprets some common characters as
commands, so you must type special commands to
generate them.  These characters include the
following:
       \$ \& \% \# \{ and \}.

In printing, text is usually emphasized with an
       \emph{italic}
type style.

\begin{em}
   A long segment of text can also be emphasized
   in this way.  Text within such a segment can be
   given \emph{additional} emphasis.
\end{em}

It is sometimes necessary to prevent \LaTeX\ from
breaking a line where it might otherwise do so.
This may be at a space, as between the ``Mr.''\ and
``Jones'' in
       ``Mr.~Jones'',        % ~ produces an unbreakable interword space.
or within a word---especially when the word is a
symbol like
       \mbox{\emph{itemnum}}
that makes little sense when hyphenated across
lines.

Footnotes\footnote{This is an example of a footnote.}
pose no problem.

\LaTeX\ is good at typesetting mathematical formulas
like
       \( x-3y + z = 7 \)
or
       \( a_{1} > x^{2n} + y^{2n} > x' \)
or
       \( \ip{A}{B} = \sum_{i} a_{i} b_{i} \).
The spaces you type in a formula are
ignored.  Remember that a letter like
       $x$                   % $ ... $  and  \( ... \)  are equivalent
is a formula when it denotes a mathematical
symbol, and it should be typed as one.

\section{Displayed Text}

Text is displayed by indenting it from the left
margin.  Quotations are commonly displayed.  There
are short quotations
\begin{quote}
   This is a short quotation.  It consists of a
   single paragraph of text.  See how it is formatted.
\end{quote}
and longer ones.
\begin{quotation}
   This is a longer quotation.  It consists of two
   paragraphs of text, neither of which are
   particularly interesting.

   This is the second paragraph of the quotation.  It
   is just as dull as the first paragraph.
\end{quotation}
Another frequently-displayed structure is a list.
The following is an example of an \emph{itemized}
list.
\begin{itemize}
   \item This is the first item of an itemized list.
         Each item in the list is marked with a ``tick''.
         You don't have to worry about what kind of tick
         mark is used.

   \item This is the second item of the list.  It
         contains another list nested inside it.  The inner
         list is an \emph{enumerated} list.
         \begin{enumerate}
            \item This is the first item of an enumerated
                  list that is nested within the itemized list.

            \item This is the second item of the inner list.
                  \LaTeX\ allows you to nest lists deeper than
                  you really should.
         \end{enumerate}
         This is the rest of the second item of the outer
         list.  It is no more interesting than any other
         part of the item.
   \item This is the third item of the list.
\end{itemize}
You can even display poetry.
\begin{verse}
   There is an environment
    for verse \\             % The \\ command separates lines
   Whose features some poets % within a stanza.
   will curse.

                             % One or more blank lines separate stanzas.

   For instead of making\\
   Them do \emph{all} line breaking, \\
   It allows them to put too many words on a line when they'd rather be
   forced to be terse.
\end{verse}

Mathematical formulas may also be displayed.  A
displayed formula
is
one-line long; multiline
formulas require special formatting instructions.
   \[  \ip{\Gamma}{\psi'} = x'' + y^{2} + z_{i}^{n}\]
Don't start a paragraph with a displayed equation,
nor make one a paragraph by itself.

\end{document}               % End of document.

\section{Bench Huge Section 1287}
% synthetic bench section 1287
% This is a sample LaTeX input file.  (Version of 12 August 2004.)
%
% A '%' character causes TeX to ignore all remaining text on the line,
% and is used for comments like this one.

\documentclass{article}      % Specifies the document class

                             % The preamble begins here.
\title{An Example Document}  % Declares the document's title.
\author{Leslie Lamport}      % Declares the author's name.
\date{January 21, 1994}      % Deleting this command produces today's date.

\newcommand{\ip}[2]{(#1, #2)}
                             % Defines \ip{arg1}{arg2} to mean
                             % (arg1, arg2).

%\newcommand{\ip}[2]{\langle #1 | #2\rangle}
                             % This is an alternative definition of
                             % \ip that is commented out.

\begin{document}             % End of preamble and beginning of text.

\maketitle                   % Produces the title.

This is an example input file.  Comparing it with
the output it generates can show you how to
produce a simple document of your own.

\section{Ordinary Text}      % Produces section heading.  Lower-level
                             % sections are begun with similar
                             % \subsection and \subsubsection commands.

The ends  of words and sentences are marked
  by   spaces. It  doesn't matter how many
spaces    you type; one is as good as 100.  The
end of   a line counts as a space.

One   or more   blank lines denote the  end
of  a paragraph.

Since any number of consecutive spaces are treated
like a single one, the formatting of the input
file makes no difference to
      \LaTeX,                % The \LaTeX command generates the LaTeX logo.
but it makes a difference to you.  When you use
\LaTeX, making your input file as easy to read
as possible will be a great help as you write
your document and when you change it.  This sample
file shows how you can add comments to your own input
file.

Because printing is different from typewriting,
there are a number of things that you have to do
differently when preparing an input file than if
you were just typing the document directly.
Quotation marks like
       ``this''
have to be handled specially, as do quotes within
quotes:
       ``\,`this'            % \, separates the double and single quote.
        is what I just
        wrote, not  `that'\,''.

Dashes come in three sizes: an
       intra-word
dash, a medium dash for number ranges like
       1--2,
and a punctuation
       dash---like
this.

A sentence-ending space should be larger than the
space between words within a sentence.  You
sometimes have to type special commands in
conjunction with punctuation characters to get
this right, as in the following sentence.
       Gnats, gnus, etc.\ all  % `\ ' makes an inter-word space.
       begin with G\@.         % \@ marks end-of-sentence punctuation.
You should check the spaces after periods when
reading your output to make sure you haven't
forgotten any special cases.  Generating an
ellipsis
       \ldots\               % `\ ' is needed after `\ldots' because TeX
                             % ignores spaces after command names like \ldots
                             % made from \ + letters.
                             %
                             % Note how a `%' character causes TeX to ignore
                             % the end of the input line, so these blank lines
                             % do not start a new paragraph.
                             %
with the right spacing around the periods requires
a special command.

\LaTeX\ interprets some common characters as
commands, so you must type special commands to
generate them.  These characters include the
following:
       \$ \& \% \# \{ and \}.

In printing, text is usually emphasized with an
       \emph{italic}
type style.

\begin{em}
   A long segment of text can also be emphasized
   in this way.  Text within such a segment can be
   given \emph{additional} emphasis.
\end{em}

It is sometimes necessary to prevent \LaTeX\ from
breaking a line where it might otherwise do so.
This may be at a space, as between the ``Mr.''\ and
``Jones'' in
       ``Mr.~Jones'',        % ~ produces an unbreakable interword space.
or within a word---especially when the word is a
symbol like
       \mbox{\emph{itemnum}}
that makes little sense when hyphenated across
lines.

Footnotes\footnote{This is an example of a footnote.}
pose no problem.

\LaTeX\ is good at typesetting mathematical formulas
like
       \( x-3y + z = 7 \)
or
       \( a_{1} > x^{2n} + y^{2n} > x' \)
or
       \( \ip{A}{B} = \sum_{i} a_{i} b_{i} \).
The spaces you type in a formula are
ignored.  Remember that a letter like
       $x$                   % $ ... $  and  \( ... \)  are equivalent
is a formula when it denotes a mathematical
symbol, and it should be typed as one.

\section{Displayed Text}

Text is displayed by indenting it from the left
margin.  Quotations are commonly displayed.  There
are short quotations
\begin{quote}
   This is a short quotation.  It consists of a
   single paragraph of text.  See how it is formatted.
\end{quote}
and longer ones.
\begin{quotation}
   This is a longer quotation.  It consists of two
   paragraphs of text, neither of which are
   particularly interesting.

   This is the second paragraph of the quotation.  It
   is just as dull as the first paragraph.
\end{quotation}
Another frequently-displayed structure is a list.
The following is an example of an \emph{itemized}
list.
\begin{itemize}
   \item This is the first item of an itemized list.
         Each item in the list is marked with a ``tick''.
         You don't have to worry about what kind of tick
         mark is used.

   \item This is the second item of the list.  It
         contains another list nested inside it.  The inner
         list is an \emph{enumerated} list.
         \begin{enumerate}
            \item This is the first item of an enumerated
                  list that is nested within the itemized list.

            \item This is the second item of the inner list.
                  \LaTeX\ allows you to nest lists deeper than
                  you really should.
         \end{enumerate}
         This is the rest of the second item of the outer
         list.  It is no more interesting than any other
         part of the item.
   \item This is the third item of the list.
\end{itemize}
You can even display poetry.
\begin{verse}
   There is an environment
    for verse \\             % The \\ command separates lines
   Whose features some poets % within a stanza.
   will curse.

                             % One or more blank lines separate stanzas.

   For instead of making\\
   Them do \emph{all} line breaking, \\
   It allows them to put too many words on a line when they'd rather be
   forced to be terse.
\end{verse}

Mathematical formulas may also be displayed.  A
displayed formula
is
one-line long; multiline
formulas require special formatting instructions.
   \[  \ip{\Gamma}{\psi'} = x'' + y^{2} + z_{i}^{n}\]
Don't start a paragraph with a displayed equation,
nor make one a paragraph by itself.

\end{document}               % End of document.

\section{Bench Huge Section 1288}
% synthetic bench section 1288
% This is a sample LaTeX input file.  (Version of 12 August 2004.)
%
% A '%' character causes TeX to ignore all remaining text on the line,
% and is used for comments like this one.

\documentclass{article}      % Specifies the document class

                             % The preamble begins here.
\title{An Example Document}  % Declares the document's title.
\author{Leslie Lamport}      % Declares the author's name.
\date{January 21, 1994}      % Deleting this command produces today's date.

\newcommand{\ip}[2]{(#1, #2)}
                             % Defines \ip{arg1}{arg2} to mean
                             % (arg1, arg2).

%\newcommand{\ip}[2]{\langle #1 | #2\rangle}
                             % This is an alternative definition of
                             % \ip that is commented out.

\begin{document}             % End of preamble and beginning of text.

\maketitle                   % Produces the title.

This is an example input file.  Comparing it with
the output it generates can show you how to
produce a simple document of your own.

\section{Ordinary Text}      % Produces section heading.  Lower-level
                             % sections are begun with similar
                             % \subsection and \subsubsection commands.

The ends  of words and sentences are marked
  by   spaces. It  doesn't matter how many
spaces    you type; one is as good as 100.  The
end of   a line counts as a space.

One   or more   blank lines denote the  end
of  a paragraph.

Since any number of consecutive spaces are treated
like a single one, the formatting of the input
file makes no difference to
      \LaTeX,                % The \LaTeX command generates the LaTeX logo.
but it makes a difference to you.  When you use
\LaTeX, making your input file as easy to read
as possible will be a great help as you write
your document and when you change it.  This sample
file shows how you can add comments to your own input
file.

Because printing is different from typewriting,
there are a number of things that you have to do
differently when preparing an input file than if
you were just typing the document directly.
Quotation marks like
       ``this''
have to be handled specially, as do quotes within
quotes:
       ``\,`this'            % \, separates the double and single quote.
        is what I just
        wrote, not  `that'\,''.

Dashes come in three sizes: an
       intra-word
dash, a medium dash for number ranges like
       1--2,
and a punctuation
       dash---like
this.

A sentence-ending space should be larger than the
space between words within a sentence.  You
sometimes have to type special commands in
conjunction with punctuation characters to get
this right, as in the following sentence.
       Gnats, gnus, etc.\ all  % `\ ' makes an inter-word space.
       begin with G\@.         % \@ marks end-of-sentence punctuation.
You should check the spaces after periods when
reading your output to make sure you haven't
forgotten any special cases.  Generating an
ellipsis
       \ldots\               % `\ ' is needed after `\ldots' because TeX
                             % ignores spaces after command names like \ldots
                             % made from \ + letters.
                             %
                             % Note how a `%' character causes TeX to ignore
                             % the end of the input line, so these blank lines
                             % do not start a new paragraph.
                             %
with the right spacing around the periods requires
a special command.

\LaTeX\ interprets some common characters as
commands, so you must type special commands to
generate them.  These characters include the
following:
       \$ \& \% \# \{ and \}.

In printing, text is usually emphasized with an
       \emph{italic}
type style.

\begin{em}
   A long segment of text can also be emphasized
   in this way.  Text within such a segment can be
   given \emph{additional} emphasis.
\end{em}

It is sometimes necessary to prevent \LaTeX\ from
breaking a line where it might otherwise do so.
This may be at a space, as between the ``Mr.''\ and
``Jones'' in
       ``Mr.~Jones'',        % ~ produces an unbreakable interword space.
or within a word---especially when the word is a
symbol like
       \mbox{\emph{itemnum}}
that makes little sense when hyphenated across
lines.

Footnotes\footnote{This is an example of a footnote.}
pose no problem.

\LaTeX\ is good at typesetting mathematical formulas
like
       \( x-3y + z = 7 \)
or
       \( a_{1} > x^{2n} + y^{2n} > x' \)
or
       \( \ip{A}{B} = \sum_{i} a_{i} b_{i} \).
The spaces you type in a formula are
ignored.  Remember that a letter like
       $x$                   % $ ... $  and  \( ... \)  are equivalent
is a formula when it denotes a mathematical
symbol, and it should be typed as one.

\section{Displayed Text}

Text is displayed by indenting it from the left
margin.  Quotations are commonly displayed.  There
are short quotations
\begin{quote}
   This is a short quotation.  It consists of a
   single paragraph of text.  See how it is formatted.
\end{quote}
and longer ones.
\begin{quotation}
   This is a longer quotation.  It consists of two
   paragraphs of text, neither of which are
   particularly interesting.

   This is the second paragraph of the quotation.  It
   is just as dull as the first paragraph.
\end{quotation}
Another frequently-displayed structure is a list.
The following is an example of an \emph{itemized}
list.
\begin{itemize}
   \item This is the first item of an itemized list.
         Each item in the list is marked with a ``tick''.
         You don't have to worry about what kind of tick
         mark is used.

   \item This is the second item of the list.  It
         contains another list nested inside it.  The inner
         list is an \emph{enumerated} list.
         \begin{enumerate}
            \item This is the first item of an enumerated
                  list that is nested within the itemized list.

            \item This is the second item of the inner list.
                  \LaTeX\ allows you to nest lists deeper than
                  you really should.
         \end{enumerate}
         This is the rest of the second item of the outer
         list.  It is no more interesting than any other
         part of the item.
   \item This is the third item of the list.
\end{itemize}
You can even display poetry.
\begin{verse}
   There is an environment
    for verse \\             % The \\ command separates lines
   Whose features some poets % within a stanza.
   will curse.

                             % One or more blank lines separate stanzas.

   For instead of making\\
   Them do \emph{all} line breaking, \\
   It allows them to put too many words on a line when they'd rather be
   forced to be terse.
\end{verse}

Mathematical formulas may also be displayed.  A
displayed formula
is
one-line long; multiline
formulas require special formatting instructions.
   \[  \ip{\Gamma}{\psi'} = x'' + y^{2} + z_{i}^{n}\]
Don't start a paragraph with a displayed equation,
nor make one a paragraph by itself.

\end{document}               % End of document.

\section{Bench Huge Section 1289}
% synthetic bench section 1289
% This is a sample LaTeX input file.  (Version of 12 August 2004.)
%
% A '%' character causes TeX to ignore all remaining text on the line,
% and is used for comments like this one.

\documentclass{article}      % Specifies the document class

                             % The preamble begins here.
\title{An Example Document}  % Declares the document's title.
\author{Leslie Lamport}      % Declares the author's name.
\date{January 21, 1994}      % Deleting this command produces today's date.

\newcommand{\ip}[2]{(#1, #2)}
                             % Defines \ip{arg1}{arg2} to mean
                             % (arg1, arg2).

%\newcommand{\ip}[2]{\langle #1 | #2\rangle}
                             % This is an alternative definition of
                             % \ip that is commented out.

\begin{document}             % End of preamble and beginning of text.

\maketitle                   % Produces the title.

This is an example input file.  Comparing it with
the output it generates can show you how to
produce a simple document of your own.

\section{Ordinary Text}      % Produces section heading.  Lower-level
                             % sections are begun with similar
                             % \subsection and \subsubsection commands.

The ends  of words and sentences are marked
  by   spaces. It  doesn't matter how many
spaces    you type; one is as good as 100.  The
end of   a line counts as a space.

One   or more   blank lines denote the  end
of  a paragraph.

Since any number of consecutive spaces are treated
like a single one, the formatting of the input
file makes no difference to
      \LaTeX,                % The \LaTeX command generates the LaTeX logo.
but it makes a difference to you.  When you use
\LaTeX, making your input file as easy to read
as possible will be a great help as you write
your document and when you change it.  This sample
file shows how you can add comments to your own input
file.

Because printing is different from typewriting,
there are a number of things that you have to do
differently when preparing an input file than if
you were just typing the document directly.
Quotation marks like
       ``this''
have to be handled specially, as do quotes within
quotes:
       ``\,`this'            % \, separates the double and single quote.
        is what I just
        wrote, not  `that'\,''.

Dashes come in three sizes: an
       intra-word
dash, a medium dash for number ranges like
       1--2,
and a punctuation
       dash---like
this.

A sentence-ending space should be larger than the
space between words within a sentence.  You
sometimes have to type special commands in
conjunction with punctuation characters to get
this right, as in the following sentence.
       Gnats, gnus, etc.\ all  % `\ ' makes an inter-word space.
       begin with G\@.         % \@ marks end-of-sentence punctuation.
You should check the spaces after periods when
reading your output to make sure you haven't
forgotten any special cases.  Generating an
ellipsis
       \ldots\               % `\ ' is needed after `\ldots' because TeX
                             % ignores spaces after command names like \ldots
                             % made from \ + letters.
                             %
                             % Note how a `%' character causes TeX to ignore
                             % the end of the input line, so these blank lines
                             % do not start a new paragraph.
                             %
with the right spacing around the periods requires
a special command.

\LaTeX\ interprets some common characters as
commands, so you must type special commands to
generate them.  These characters include the
following:
       \$ \& \% \# \{ and \}.

In printing, text is usually emphasized with an
       \emph{italic}
type style.

\begin{em}
   A long segment of text can also be emphasized
   in this way.  Text within such a segment can be
   given \emph{additional} emphasis.
\end{em}

It is sometimes necessary to prevent \LaTeX\ from
breaking a line where it might otherwise do so.
This may be at a space, as between the ``Mr.''\ and
``Jones'' in
       ``Mr.~Jones'',        % ~ produces an unbreakable interword space.
or within a word---especially when the word is a
symbol like
       \mbox{\emph{itemnum}}
that makes little sense when hyphenated across
lines.

Footnotes\footnote{This is an example of a footnote.}
pose no problem.

\LaTeX\ is good at typesetting mathematical formulas
like
       \( x-3y + z = 7 \)
or
       \( a_{1} > x^{2n} + y^{2n} > x' \)
or
       \( \ip{A}{B} = \sum_{i} a_{i} b_{i} \).
The spaces you type in a formula are
ignored.  Remember that a letter like
       $x$                   % $ ... $  and  \( ... \)  are equivalent
is a formula when it denotes a mathematical
symbol, and it should be typed as one.

\section{Displayed Text}

Text is displayed by indenting it from the left
margin.  Quotations are commonly displayed.  There
are short quotations
\begin{quote}
   This is a short quotation.  It consists of a
   single paragraph of text.  See how it is formatted.
\end{quote}
and longer ones.
\begin{quotation}
   This is a longer quotation.  It consists of two
   paragraphs of text, neither of which are
   particularly interesting.

   This is the second paragraph of the quotation.  It
   is just as dull as the first paragraph.
\end{quotation}
Another frequently-displayed structure is a list.
The following is an example of an \emph{itemized}
list.
\begin{itemize}
   \item This is the first item of an itemized list.
         Each item in the list is marked with a ``tick''.
         You don't have to worry about what kind of tick
         mark is used.

   \item This is the second item of the list.  It
         contains another list nested inside it.  The inner
         list is an \emph{enumerated} list.
         \begin{enumerate}
            \item This is the first item of an enumerated
                  list that is nested within the itemized list.

            \item This is the second item of the inner list.
                  \LaTeX\ allows you to nest lists deeper than
                  you really should.
         \end{enumerate}
         This is the rest of the second item of the outer
         list.  It is no more interesting than any other
         part of the item.
   \item This is the third item of the list.
\end{itemize}
You can even display poetry.
\begin{verse}
   There is an environment
    for verse \\             % The \\ command separates lines
   Whose features some poets % within a stanza.
   will curse.

                             % One or more blank lines separate stanzas.

   For instead of making\\
   Them do \emph{all} line breaking, \\
   It allows them to put too many words on a line when they'd rather be
   forced to be terse.
\end{verse}

Mathematical formulas may also be displayed.  A
displayed formula
is
one-line long; multiline
formulas require special formatting instructions.
   \[  \ip{\Gamma}{\psi'} = x'' + y^{2} + z_{i}^{n}\]
Don't start a paragraph with a displayed equation,
nor make one a paragraph by itself.

\end{document}               % End of document.

\section{Bench Huge Section 1290}
% synthetic bench section 1290
% This is a sample LaTeX input file.  (Version of 12 August 2004.)
%
% A '%' character causes TeX to ignore all remaining text on the line,
% and is used for comments like this one.

\documentclass{article}      % Specifies the document class

                             % The preamble begins here.
\title{An Example Document}  % Declares the document's title.
\author{Leslie Lamport}      % Declares the author's name.
\date{January 21, 1994}      % Deleting this command produces today's date.

\newcommand{\ip}[2]{(#1, #2)}
                             % Defines \ip{arg1}{arg2} to mean
                             % (arg1, arg2).

%\newcommand{\ip}[2]{\langle #1 | #2\rangle}
                             % This is an alternative definition of
                             % \ip that is commented out.

\begin{document}             % End of preamble and beginning of text.

\maketitle                   % Produces the title.

This is an example input file.  Comparing it with
the output it generates can show you how to
produce a simple document of your own.

\section{Ordinary Text}      % Produces section heading.  Lower-level
                             % sections are begun with similar
                             % \subsection and \subsubsection commands.

The ends  of words and sentences are marked
  by   spaces. It  doesn't matter how many
spaces    you type; one is as good as 100.  The
end of   a line counts as a space.

One   or more   blank lines denote the  end
of  a paragraph.

Since any number of consecutive spaces are treated
like a single one, the formatting of the input
file makes no difference to
      \LaTeX,                % The \LaTeX command generates the LaTeX logo.
but it makes a difference to you.  When you use
\LaTeX, making your input file as easy to read
as possible will be a great help as you write
your document and when you change it.  This sample
file shows how you can add comments to your own input
file.

Because printing is different from typewriting,
there are a number of things that you have to do
differently when preparing an input file than if
you were just typing the document directly.
Quotation marks like
       ``this''
have to be handled specially, as do quotes within
quotes:
       ``\,`this'            % \, separates the double and single quote.
        is what I just
        wrote, not  `that'\,''.

Dashes come in three sizes: an
       intra-word
dash, a medium dash for number ranges like
       1--2,
and a punctuation
       dash---like
this.

A sentence-ending space should be larger than the
space between words within a sentence.  You
sometimes have to type special commands in
conjunction with punctuation characters to get
this right, as in the following sentence.
       Gnats, gnus, etc.\ all  % `\ ' makes an inter-word space.
       begin with G\@.         % \@ marks end-of-sentence punctuation.
You should check the spaces after periods when
reading your output to make sure you haven't
forgotten any special cases.  Generating an
ellipsis
       \ldots\               % `\ ' is needed after `\ldots' because TeX
                             % ignores spaces after command names like \ldots
                             % made from \ + letters.
                             %
                             % Note how a `%' character causes TeX to ignore
                             % the end of the input line, so these blank lines
                             % do not start a new paragraph.
                             %
with the right spacing around the periods requires
a special command.

\LaTeX\ interprets some common characters as
commands, so you must type special commands to
generate them.  These characters include the
following:
       \$ \& \% \# \{ and \}.

In printing, text is usually emphasized with an
       \emph{italic}
type style.

\begin{em}
   A long segment of text can also be emphasized
   in this way.  Text within such a segment can be
   given \emph{additional} emphasis.
\end{em}

It is sometimes necessary to prevent \LaTeX\ from
breaking a line where it might otherwise do so.
This may be at a space, as between the ``Mr.''\ and
``Jones'' in
       ``Mr.~Jones'',        % ~ produces an unbreakable interword space.
or within a word---especially when the word is a
symbol like
       \mbox{\emph{itemnum}}
that makes little sense when hyphenated across
lines.

Footnotes\footnote{This is an example of a footnote.}
pose no problem.

\LaTeX\ is good at typesetting mathematical formulas
like
       \( x-3y + z = 7 \)
or
       \( a_{1} > x^{2n} + y^{2n} > x' \)
or
       \( \ip{A}{B} = \sum_{i} a_{i} b_{i} \).
The spaces you type in a formula are
ignored.  Remember that a letter like
       $x$                   % $ ... $  and  \( ... \)  are equivalent
is a formula when it denotes a mathematical
symbol, and it should be typed as one.

\section{Displayed Text}

Text is displayed by indenting it from the left
margin.  Quotations are commonly displayed.  There
are short quotations
\begin{quote}
   This is a short quotation.  It consists of a
   single paragraph of text.  See how it is formatted.
\end{quote}
and longer ones.
\begin{quotation}
   This is a longer quotation.  It consists of two
   paragraphs of text, neither of which are
   particularly interesting.

   This is the second paragraph of the quotation.  It
   is just as dull as the first paragraph.
\end{quotation}
Another frequently-displayed structure is a list.
The following is an example of an \emph{itemized}
list.
\begin{itemize}
   \item This is the first item of an itemized list.
         Each item in the list is marked with a ``tick''.
         You don't have to worry about what kind of tick
         mark is used.

   \item This is the second item of the list.  It
         contains another list nested inside it.  The inner
         list is an \emph{enumerated} list.
         \begin{enumerate}
            \item This is the first item of an enumerated
                  list that is nested within the itemized list.

            \item This is the second item of the inner list.
                  \LaTeX\ allows you to nest lists deeper than
                  you really should.
         \end{enumerate}
         This is the rest of the second item of the outer
         list.  It is no more interesting than any other
         part of the item.
   \item This is the third item of the list.
\end{itemize}
You can even display poetry.
\begin{verse}
   There is an environment
    for verse \\             % The \\ command separates lines
   Whose features some poets % within a stanza.
   will curse.

                             % One or more blank lines separate stanzas.

   For instead of making\\
   Them do \emph{all} line breaking, \\
   It allows them to put too many words on a line when they'd rather be
   forced to be terse.
\end{verse}

Mathematical formulas may also be displayed.  A
displayed formula
is
one-line long; multiline
formulas require special formatting instructions.
   \[  \ip{\Gamma}{\psi'} = x'' + y^{2} + z_{i}^{n}\]
Don't start a paragraph with a displayed equation,
nor make one a paragraph by itself.

\end{document}               % End of document.

\section{Bench Huge Section 1291}
% synthetic bench section 1291
% This is a sample LaTeX input file.  (Version of 12 August 2004.)
%
% A '%' character causes TeX to ignore all remaining text on the line,
% and is used for comments like this one.

\documentclass{article}      % Specifies the document class

                             % The preamble begins here.
\title{An Example Document}  % Declares the document's title.
\author{Leslie Lamport}      % Declares the author's name.
\date{January 21, 1994}      % Deleting this command produces today's date.

\newcommand{\ip}[2]{(#1, #2)}
                             % Defines \ip{arg1}{arg2} to mean
                             % (arg1, arg2).

%\newcommand{\ip}[2]{\langle #1 | #2\rangle}
                             % This is an alternative definition of
                             % \ip that is commented out.

\begin{document}             % End of preamble and beginning of text.

\maketitle                   % Produces the title.

This is an example input file.  Comparing it with
the output it generates can show you how to
produce a simple document of your own.

\section{Ordinary Text}      % Produces section heading.  Lower-level
                             % sections are begun with similar
                             % \subsection and \subsubsection commands.

The ends  of words and sentences are marked
  by   spaces. It  doesn't matter how many
spaces    you type; one is as good as 100.  The
end of   a line counts as a space.

One   or more   blank lines denote the  end
of  a paragraph.

Since any number of consecutive spaces are treated
like a single one, the formatting of the input
file makes no difference to
      \LaTeX,                % The \LaTeX command generates the LaTeX logo.
but it makes a difference to you.  When you use
\LaTeX, making your input file as easy to read
as possible will be a great help as you write
your document and when you change it.  This sample
file shows how you can add comments to your own input
file.

Because printing is different from typewriting,
there are a number of things that you have to do
differently when preparing an input file than if
you were just typing the document directly.
Quotation marks like
       ``this''
have to be handled specially, as do quotes within
quotes:
       ``\,`this'            % \, separates the double and single quote.
        is what I just
        wrote, not  `that'\,''.

Dashes come in three sizes: an
       intra-word
dash, a medium dash for number ranges like
       1--2,
and a punctuation
       dash---like
this.

A sentence-ending space should be larger than the
space between words within a sentence.  You
sometimes have to type special commands in
conjunction with punctuation characters to get
this right, as in the following sentence.
       Gnats, gnus, etc.\ all  % `\ ' makes an inter-word space.
       begin with G\@.         % \@ marks end-of-sentence punctuation.
You should check the spaces after periods when
reading your output to make sure you haven't
forgotten any special cases.  Generating an
ellipsis
       \ldots\               % `\ ' is needed after `\ldots' because TeX
                             % ignores spaces after command names like \ldots
                             % made from \ + letters.
                             %
                             % Note how a `%' character causes TeX to ignore
                             % the end of the input line, so these blank lines
                             % do not start a new paragraph.
                             %
with the right spacing around the periods requires
a special command.

\LaTeX\ interprets some common characters as
commands, so you must type special commands to
generate them.  These characters include the
following:
       \$ \& \% \# \{ and \}.

In printing, text is usually emphasized with an
       \emph{italic}
type style.

\begin{em}
   A long segment of text can also be emphasized
   in this way.  Text within such a segment can be
   given \emph{additional} emphasis.
\end{em}

It is sometimes necessary to prevent \LaTeX\ from
breaking a line where it might otherwise do so.
This may be at a space, as between the ``Mr.''\ and
``Jones'' in
       ``Mr.~Jones'',        % ~ produces an unbreakable interword space.
or within a word---especially when the word is a
symbol like
       \mbox{\emph{itemnum}}
that makes little sense when hyphenated across
lines.

Footnotes\footnote{This is an example of a footnote.}
pose no problem.

\LaTeX\ is good at typesetting mathematical formulas
like
       \( x-3y + z = 7 \)
or
       \( a_{1} > x^{2n} + y^{2n} > x' \)
or
       \( \ip{A}{B} = \sum_{i} a_{i} b_{i} \).
The spaces you type in a formula are
ignored.  Remember that a letter like
       $x$                   % $ ... $  and  \( ... \)  are equivalent
is a formula when it denotes a mathematical
symbol, and it should be typed as one.

\section{Displayed Text}

Text is displayed by indenting it from the left
margin.  Quotations are commonly displayed.  There
are short quotations
\begin{quote}
   This is a short quotation.  It consists of a
   single paragraph of text.  See how it is formatted.
\end{quote}
and longer ones.
\begin{quotation}
   This is a longer quotation.  It consists of two
   paragraphs of text, neither of which are
   particularly interesting.

   This is the second paragraph of the quotation.  It
   is just as dull as the first paragraph.
\end{quotation}
Another frequently-displayed structure is a list.
The following is an example of an \emph{itemized}
list.
\begin{itemize}
   \item This is the first item of an itemized list.
         Each item in the list is marked with a ``tick''.
         You don't have to worry about what kind of tick
         mark is used.

   \item This is the second item of the list.  It
         contains another list nested inside it.  The inner
         list is an \emph{enumerated} list.
         \begin{enumerate}
            \item This is the first item of an enumerated
                  list that is nested within the itemized list.

            \item This is the second item of the inner list.
                  \LaTeX\ allows you to nest lists deeper than
                  you really should.
         \end{enumerate}
         This is the rest of the second item of the outer
         list.  It is no more interesting than any other
         part of the item.
   \item This is the third item of the list.
\end{itemize}
You can even display poetry.
\begin{verse}
   There is an environment
    for verse \\             % The \\ command separates lines
   Whose features some poets % within a stanza.
   will curse.

                             % One or more blank lines separate stanzas.

   For instead of making\\
   Them do \emph{all} line breaking, \\
   It allows them to put too many words on a line when they'd rather be
   forced to be terse.
\end{verse}

Mathematical formulas may also be displayed.  A
displayed formula
is
one-line long; multiline
formulas require special formatting instructions.
   \[  \ip{\Gamma}{\psi'} = x'' + y^{2} + z_{i}^{n}\]
Don't start a paragraph with a displayed equation,
nor make one a paragraph by itself.

\end{document}               % End of document.

\section{Bench Huge Section 1292}
% synthetic bench section 1292
% This is a sample LaTeX input file.  (Version of 12 August 2004.)
%
% A '%' character causes TeX to ignore all remaining text on the line,
% and is used for comments like this one.

\documentclass{article}      % Specifies the document class

                             % The preamble begins here.
\title{An Example Document}  % Declares the document's title.
\author{Leslie Lamport}      % Declares the author's name.
\date{January 21, 1994}      % Deleting this command produces today's date.

\newcommand{\ip}[2]{(#1, #2)}
                             % Defines \ip{arg1}{arg2} to mean
                             % (arg1, arg2).

%\newcommand{\ip}[2]{\langle #1 | #2\rangle}
                             % This is an alternative definition of
                             % \ip that is commented out.

\begin{document}             % End of preamble and beginning of text.

\maketitle                   % Produces the title.

This is an example input file.  Comparing it with
the output it generates can show you how to
produce a simple document of your own.

\section{Ordinary Text}      % Produces section heading.  Lower-level
                             % sections are begun with similar
                             % \subsection and \subsubsection commands.

The ends  of words and sentences are marked
  by   spaces. It  doesn't matter how many
spaces    you type; one is as good as 100.  The
end of   a line counts as a space.

One   or more   blank lines denote the  end
of  a paragraph.

Since any number of consecutive spaces are treated
like a single one, the formatting of the input
file makes no difference to
      \LaTeX,                % The \LaTeX command generates the LaTeX logo.
but it makes a difference to you.  When you use
\LaTeX, making your input file as easy to read
as possible will be a great help as you write
your document and when you change it.  This sample
file shows how you can add comments to your own input
file.

Because printing is different from typewriting,
there are a number of things that you have to do
differently when preparing an input file than if
you were just typing the document directly.
Quotation marks like
       ``this''
have to be handled specially, as do quotes within
quotes:
       ``\,`this'            % \, separates the double and single quote.
        is what I just
        wrote, not  `that'\,''.

Dashes come in three sizes: an
       intra-word
dash, a medium dash for number ranges like
       1--2,
and a punctuation
       dash---like
this.

A sentence-ending space should be larger than the
space between words within a sentence.  You
sometimes have to type special commands in
conjunction with punctuation characters to get
this right, as in the following sentence.
       Gnats, gnus, etc.\ all  % `\ ' makes an inter-word space.
       begin with G\@.         % \@ marks end-of-sentence punctuation.
You should check the spaces after periods when
reading your output to make sure you haven't
forgotten any special cases.  Generating an
ellipsis
       \ldots\               % `\ ' is needed after `\ldots' because TeX
                             % ignores spaces after command names like \ldots
                             % made from \ + letters.
                             %
                             % Note how a `%' character causes TeX to ignore
                             % the end of the input line, so these blank lines
                             % do not start a new paragraph.
                             %
with the right spacing around the periods requires
a special command.

\LaTeX\ interprets some common characters as
commands, so you must type special commands to
generate them.  These characters include the
following:
       \$ \& \% \# \{ and \}.

In printing, text is usually emphasized with an
       \emph{italic}
type style.

\begin{em}
   A long segment of text can also be emphasized
   in this way.  Text within such a segment can be
   given \emph{additional} emphasis.
\end{em}

It is sometimes necessary to prevent \LaTeX\ from
breaking a line where it might otherwise do so.
This may be at a space, as between the ``Mr.''\ and
``Jones'' in
       ``Mr.~Jones'',        % ~ produces an unbreakable interword space.
or within a word---especially when the word is a
symbol like
       \mbox{\emph{itemnum}}
that makes little sense when hyphenated across
lines.

Footnotes\footnote{This is an example of a footnote.}
pose no problem.

\LaTeX\ is good at typesetting mathematical formulas
like
       \( x-3y + z = 7 \)
or
       \( a_{1} > x^{2n} + y^{2n} > x' \)
or
       \( \ip{A}{B} = \sum_{i} a_{i} b_{i} \).
The spaces you type in a formula are
ignored.  Remember that a letter like
       $x$                   % $ ... $  and  \( ... \)  are equivalent
is a formula when it denotes a mathematical
symbol, and it should be typed as one.

\section{Displayed Text}

Text is displayed by indenting it from the left
margin.  Quotations are commonly displayed.  There
are short quotations
\begin{quote}
   This is a short quotation.  It consists of a
   single paragraph of text.  See how it is formatted.
\end{quote}
and longer ones.
\begin{quotation}
   This is a longer quotation.  It consists of two
   paragraphs of text, neither of which are
   particularly interesting.

   This is the second paragraph of the quotation.  It
   is just as dull as the first paragraph.
\end{quotation}
Another frequently-displayed structure is a list.
The following is an example of an \emph{itemized}
list.
\begin{itemize}
   \item This is the first item of an itemized list.
         Each item in the list is marked with a ``tick''.
         You don't have to worry about what kind of tick
         mark is used.

   \item This is the second item of the list.  It
         contains another list nested inside it.  The inner
         list is an \emph{enumerated} list.
         \begin{enumerate}
            \item This is the first item of an enumerated
                  list that is nested within the itemized list.

            \item This is the second item of the inner list.
                  \LaTeX\ allows you to nest lists deeper than
                  you really should.
         \end{enumerate}
         This is the rest of the second item of the outer
         list.  It is no more interesting than any other
         part of the item.
   \item This is the third item of the list.
\end{itemize}
You can even display poetry.
\begin{verse}
   There is an environment
    for verse \\             % The \\ command separates lines
   Whose features some poets % within a stanza.
   will curse.

                             % One or more blank lines separate stanzas.

   For instead of making\\
   Them do \emph{all} line breaking, \\
   It allows them to put too many words on a line when they'd rather be
   forced to be terse.
\end{verse}

Mathematical formulas may also be displayed.  A
displayed formula
is
one-line long; multiline
formulas require special formatting instructions.
   \[  \ip{\Gamma}{\psi'} = x'' + y^{2} + z_{i}^{n}\]
Don't start a paragraph with a displayed equation,
nor make one a paragraph by itself.

\end{document}               % End of document.

\section{Bench Huge Section 1293}
% synthetic bench section 1293
% This is a sample LaTeX input file.  (Version of 12 August 2004.)
%
% A '%' character causes TeX to ignore all remaining text on the line,
% and is used for comments like this one.

\documentclass{article}      % Specifies the document class

                             % The preamble begins here.
\title{An Example Document}  % Declares the document's title.
\author{Leslie Lamport}      % Declares the author's name.
\date{January 21, 1994}      % Deleting this command produces today's date.

\newcommand{\ip}[2]{(#1, #2)}
                             % Defines \ip{arg1}{arg2} to mean
                             % (arg1, arg2).

%\newcommand{\ip}[2]{\langle #1 | #2\rangle}
                             % This is an alternative definition of
                             % \ip that is commented out.

\begin{document}             % End of preamble and beginning of text.

\maketitle                   % Produces the title.

This is an example input file.  Comparing it with
the output it generates can show you how to
produce a simple document of your own.

\section{Ordinary Text}      % Produces section heading.  Lower-level
                             % sections are begun with similar
                             % \subsection and \subsubsection commands.

The ends  of words and sentences are marked
  by   spaces. It  doesn't matter how many
spaces    you type; one is as good as 100.  The
end of   a line counts as a space.

One   or more   blank lines denote the  end
of  a paragraph.

Since any number of consecutive spaces are treated
like a single one, the formatting of the input
file makes no difference to
      \LaTeX,                % The \LaTeX command generates the LaTeX logo.
but it makes a difference to you.  When you use
\LaTeX, making your input file as easy to read
as possible will be a great help as you write
your document and when you change it.  This sample
file shows how you can add comments to your own input
file.

Because printing is different from typewriting,
there are a number of things that you have to do
differently when preparing an input file than if
you were just typing the document directly.
Quotation marks like
       ``this''
have to be handled specially, as do quotes within
quotes:
       ``\,`this'            % \, separates the double and single quote.
        is what I just
        wrote, not  `that'\,''.

Dashes come in three sizes: an
       intra-word
dash, a medium dash for number ranges like
       1--2,
and a punctuation
       dash---like
this.

A sentence-ending space should be larger than the
space between words within a sentence.  You
sometimes have to type special commands in
conjunction with punctuation characters to get
this right, as in the following sentence.
       Gnats, gnus, etc.\ all  % `\ ' makes an inter-word space.
       begin with G\@.         % \@ marks end-of-sentence punctuation.
You should check the spaces after periods when
reading your output to make sure you haven't
forgotten any special cases.  Generating an
ellipsis
       \ldots\               % `\ ' is needed after `\ldots' because TeX
                             % ignores spaces after command names like \ldots
                             % made from \ + letters.
                             %
                             % Note how a `%' character causes TeX to ignore
                             % the end of the input line, so these blank lines
                             % do not start a new paragraph.
                             %
with the right spacing around the periods requires
a special command.

\LaTeX\ interprets some common characters as
commands, so you must type special commands to
generate them.  These characters include the
following:
       \$ \& \% \# \{ and \}.

In printing, text is usually emphasized with an
       \emph{italic}
type style.

\begin{em}
   A long segment of text can also be emphasized
   in this way.  Text within such a segment can be
   given \emph{additional} emphasis.
\end{em}

It is sometimes necessary to prevent \LaTeX\ from
breaking a line where it might otherwise do so.
This may be at a space, as between the ``Mr.''\ and
``Jones'' in
       ``Mr.~Jones'',        % ~ produces an unbreakable interword space.
or within a word---especially when the word is a
symbol like
       \mbox{\emph{itemnum}}
that makes little sense when hyphenated across
lines.

Footnotes\footnote{This is an example of a footnote.}
pose no problem.

\LaTeX\ is good at typesetting mathematical formulas
like
       \( x-3y + z = 7 \)
or
       \( a_{1} > x^{2n} + y^{2n} > x' \)
or
       \( \ip{A}{B} = \sum_{i} a_{i} b_{i} \).
The spaces you type in a formula are
ignored.  Remember that a letter like
       $x$                   % $ ... $  and  \( ... \)  are equivalent
is a formula when it denotes a mathematical
symbol, and it should be typed as one.

\section{Displayed Text}

Text is displayed by indenting it from the left
margin.  Quotations are commonly displayed.  There
are short quotations
\begin{quote}
   This is a short quotation.  It consists of a
   single paragraph of text.  See how it is formatted.
\end{quote}
and longer ones.
\begin{quotation}
   This is a longer quotation.  It consists of two
   paragraphs of text, neither of which are
   particularly interesting.

   This is the second paragraph of the quotation.  It
   is just as dull as the first paragraph.
\end{quotation}
Another frequently-displayed structure is a list.
The following is an example of an \emph{itemized}
list.
\begin{itemize}
   \item This is the first item of an itemized list.
         Each item in the list is marked with a ``tick''.
         You don't have to worry about what kind of tick
         mark is used.

   \item This is the second item of the list.  It
         contains another list nested inside it.  The inner
         list is an \emph{enumerated} list.
         \begin{enumerate}
            \item This is the first item of an enumerated
                  list that is nested within the itemized list.

            \item This is the second item of the inner list.
                  \LaTeX\ allows you to nest lists deeper than
                  you really should.
         \end{enumerate}
         This is the rest of the second item of the outer
         list.  It is no more interesting than any other
         part of the item.
   \item This is the third item of the list.
\end{itemize}
You can even display poetry.
\begin{verse}
   There is an environment
    for verse \\             % The \\ command separates lines
   Whose features some poets % within a stanza.
   will curse.

                             % One or more blank lines separate stanzas.

   For instead of making\\
   Them do \emph{all} line breaking, \\
   It allows them to put too many words on a line when they'd rather be
   forced to be terse.
\end{verse}

Mathematical formulas may also be displayed.  A
displayed formula
is
one-line long; multiline
formulas require special formatting instructions.
   \[  \ip{\Gamma}{\psi'} = x'' + y^{2} + z_{i}^{n}\]
Don't start a paragraph with a displayed equation,
nor make one a paragraph by itself.

\end{document}               % End of document.

\section{Bench Huge Section 1294}
% synthetic bench section 1294
% This is a sample LaTeX input file.  (Version of 12 August 2004.)
%
% A '%' character causes TeX to ignore all remaining text on the line,
% and is used for comments like this one.

\documentclass{article}      % Specifies the document class

                             % The preamble begins here.
\title{An Example Document}  % Declares the document's title.
\author{Leslie Lamport}      % Declares the author's name.
\date{January 21, 1994}      % Deleting this command produces today's date.

\newcommand{\ip}[2]{(#1, #2)}
                             % Defines \ip{arg1}{arg2} to mean
                             % (arg1, arg2).

%\newcommand{\ip}[2]{\langle #1 | #2\rangle}
                             % This is an alternative definition of
                             % \ip that is commented out.

\begin{document}             % End of preamble and beginning of text.

\maketitle                   % Produces the title.

This is an example input file.  Comparing it with
the output it generates can show you how to
produce a simple document of your own.

\section{Ordinary Text}      % Produces section heading.  Lower-level
                             % sections are begun with similar
                             % \subsection and \subsubsection commands.

The ends  of words and sentences are marked
  by   spaces. It  doesn't matter how many
spaces    you type; one is as good as 100.  The
end of   a line counts as a space.

One   or more   blank lines denote the  end
of  a paragraph.

Since any number of consecutive spaces are treated
like a single one, the formatting of the input
file makes no difference to
      \LaTeX,                % The \LaTeX command generates the LaTeX logo.
but it makes a difference to you.  When you use
\LaTeX, making your input file as easy to read
as possible will be a great help as you write
your document and when you change it.  This sample
file shows how you can add comments to your own input
file.

Because printing is different from typewriting,
there are a number of things that you have to do
differently when preparing an input file than if
you were just typing the document directly.
Quotation marks like
       ``this''
have to be handled specially, as do quotes within
quotes:
       ``\,`this'            % \, separates the double and single quote.
        is what I just
        wrote, not  `that'\,''.

Dashes come in three sizes: an
       intra-word
dash, a medium dash for number ranges like
       1--2,
and a punctuation
       dash---like
this.

A sentence-ending space should be larger than the
space between words within a sentence.  You
sometimes have to type special commands in
conjunction with punctuation characters to get
this right, as in the following sentence.
       Gnats, gnus, etc.\ all  % `\ ' makes an inter-word space.
       begin with G\@.         % \@ marks end-of-sentence punctuation.
You should check the spaces after periods when
reading your output to make sure you haven't
forgotten any special cases.  Generating an
ellipsis
       \ldots\               % `\ ' is needed after `\ldots' because TeX
                             % ignores spaces after command names like \ldots
                             % made from \ + letters.
                             %
                             % Note how a `%' character causes TeX to ignore
                             % the end of the input line, so these blank lines
                             % do not start a new paragraph.
                             %
with the right spacing around the periods requires
a special command.

\LaTeX\ interprets some common characters as
commands, so you must type special commands to
generate them.  These characters include the
following:
       \$ \& \% \# \{ and \}.

In printing, text is usually emphasized with an
       \emph{italic}
type style.

\begin{em}
   A long segment of text can also be emphasized
   in this way.  Text within such a segment can be
   given \emph{additional} emphasis.
\end{em}

It is sometimes necessary to prevent \LaTeX\ from
breaking a line where it might otherwise do so.
This may be at a space, as between the ``Mr.''\ and
``Jones'' in
       ``Mr.~Jones'',        % ~ produces an unbreakable interword space.
or within a word---especially when the word is a
symbol like
       \mbox{\emph{itemnum}}
that makes little sense when hyphenated across
lines.

Footnotes\footnote{This is an example of a footnote.}
pose no problem.

\LaTeX\ is good at typesetting mathematical formulas
like
       \( x-3y + z = 7 \)
or
       \( a_{1} > x^{2n} + y^{2n} > x' \)
or
       \( \ip{A}{B} = \sum_{i} a_{i} b_{i} \).
The spaces you type in a formula are
ignored.  Remember that a letter like
       $x$                   % $ ... $  and  \( ... \)  are equivalent
is a formula when it denotes a mathematical
symbol, and it should be typed as one.

\section{Displayed Text}

Text is displayed by indenting it from the left
margin.  Quotations are commonly displayed.  There
are short quotations
\begin{quote}
   This is a short quotation.  It consists of a
   single paragraph of text.  See how it is formatted.
\end{quote}
and longer ones.
\begin{quotation}
   This is a longer quotation.  It consists of two
   paragraphs of text, neither of which are
   particularly interesting.

   This is the second paragraph of the quotation.  It
   is just as dull as the first paragraph.
\end{quotation}
Another frequently-displayed structure is a list.
The following is an example of an \emph{itemized}
list.
\begin{itemize}
   \item This is the first item of an itemized list.
         Each item in the list is marked with a ``tick''.
         You don't have to worry about what kind of tick
         mark is used.

   \item This is the second item of the list.  It
         contains another list nested inside it.  The inner
         list is an \emph{enumerated} list.
         \begin{enumerate}
            \item This is the first item of an enumerated
                  list that is nested within the itemized list.

            \item This is the second item of the inner list.
                  \LaTeX\ allows you to nest lists deeper than
                  you really should.
         \end{enumerate}
         This is the rest of the second item of the outer
         list.  It is no more interesting than any other
         part of the item.
   \item This is the third item of the list.
\end{itemize}
You can even display poetry.
\begin{verse}
   There is an environment
    for verse \\             % The \\ command separates lines
   Whose features some poets % within a stanza.
   will curse.

                             % One or more blank lines separate stanzas.

   For instead of making\\
   Them do \emph{all} line breaking, \\
   It allows them to put too many words on a line when they'd rather be
   forced to be terse.
\end{verse}

Mathematical formulas may also be displayed.  A
displayed formula
is
one-line long; multiline
formulas require special formatting instructions.
   \[  \ip{\Gamma}{\psi'} = x'' + y^{2} + z_{i}^{n}\]
Don't start a paragraph with a displayed equation,
nor make one a paragraph by itself.

\end{document}               % End of document.

\section{Bench Huge Section 1295}
% synthetic bench section 1295
% This is a sample LaTeX input file.  (Version of 12 August 2004.)
%
% A '%' character causes TeX to ignore all remaining text on the line,
% and is used for comments like this one.

\documentclass{article}      % Specifies the document class

                             % The preamble begins here.
\title{An Example Document}  % Declares the document's title.
\author{Leslie Lamport}      % Declares the author's name.
\date{January 21, 1994}      % Deleting this command produces today's date.

\newcommand{\ip}[2]{(#1, #2)}
                             % Defines \ip{arg1}{arg2} to mean
                             % (arg1, arg2).

%\newcommand{\ip}[2]{\langle #1 | #2\rangle}
                             % This is an alternative definition of
                             % \ip that is commented out.

\begin{document}             % End of preamble and beginning of text.

\maketitle                   % Produces the title.

This is an example input file.  Comparing it with
the output it generates can show you how to
produce a simple document of your own.

\section{Ordinary Text}      % Produces section heading.  Lower-level
                             % sections are begun with similar
                             % \subsection and \subsubsection commands.

The ends  of words and sentences are marked
  by   spaces. It  doesn't matter how many
spaces    you type; one is as good as 100.  The
end of   a line counts as a space.

One   or more   blank lines denote the  end
of  a paragraph.

Since any number of consecutive spaces are treated
like a single one, the formatting of the input
file makes no difference to
      \LaTeX,                % The \LaTeX command generates the LaTeX logo.
but it makes a difference to you.  When you use
\LaTeX, making your input file as easy to read
as possible will be a great help as you write
your document and when you change it.  This sample
file shows how you can add comments to your own input
file.

Because printing is different from typewriting,
there are a number of things that you have to do
differently when preparing an input file than if
you were just typing the document directly.
Quotation marks like
       ``this''
have to be handled specially, as do quotes within
quotes:
       ``\,`this'            % \, separates the double and single quote.
        is what I just
        wrote, not  `that'\,''.

Dashes come in three sizes: an
       intra-word
dash, a medium dash for number ranges like
       1--2,
and a punctuation
       dash---like
this.

A sentence-ending space should be larger than the
space between words within a sentence.  You
sometimes have to type special commands in
conjunction with punctuation characters to get
this right, as in the following sentence.
       Gnats, gnus, etc.\ all  % `\ ' makes an inter-word space.
       begin with G\@.         % \@ marks end-of-sentence punctuation.
You should check the spaces after periods when
reading your output to make sure you haven't
forgotten any special cases.  Generating an
ellipsis
       \ldots\               % `\ ' is needed after `\ldots' because TeX
                             % ignores spaces after command names like \ldots
                             % made from \ + letters.
                             %
                             % Note how a `%' character causes TeX to ignore
                             % the end of the input line, so these blank lines
                             % do not start a new paragraph.
                             %
with the right spacing around the periods requires
a special command.

\LaTeX\ interprets some common characters as
commands, so you must type special commands to
generate them.  These characters include the
following:
       \$ \& \% \# \{ and \}.

In printing, text is usually emphasized with an
       \emph{italic}
type style.

\begin{em}
   A long segment of text can also be emphasized
   in this way.  Text within such a segment can be
   given \emph{additional} emphasis.
\end{em}

It is sometimes necessary to prevent \LaTeX\ from
breaking a line where it might otherwise do so.
This may be at a space, as between the ``Mr.''\ and
``Jones'' in
       ``Mr.~Jones'',        % ~ produces an unbreakable interword space.
or within a word---especially when the word is a
symbol like
       \mbox{\emph{itemnum}}
that makes little sense when hyphenated across
lines.

Footnotes\footnote{This is an example of a footnote.}
pose no problem.

\LaTeX\ is good at typesetting mathematical formulas
like
       \( x-3y + z = 7 \)
or
       \( a_{1} > x^{2n} + y^{2n} > x' \)
or
       \( \ip{A}{B} = \sum_{i} a_{i} b_{i} \).
The spaces you type in a formula are
ignored.  Remember that a letter like
       $x$                   % $ ... $  and  \( ... \)  are equivalent
is a formula when it denotes a mathematical
symbol, and it should be typed as one.

\section{Displayed Text}

Text is displayed by indenting it from the left
margin.  Quotations are commonly displayed.  There
are short quotations
\begin{quote}
   This is a short quotation.  It consists of a
   single paragraph of text.  See how it is formatted.
\end{quote}
and longer ones.
\begin{quotation}
   This is a longer quotation.  It consists of two
   paragraphs of text, neither of which are
   particularly interesting.

   This is the second paragraph of the quotation.  It
   is just as dull as the first paragraph.
\end{quotation}
Another frequently-displayed structure is a list.
The following is an example of an \emph{itemized}
list.
\begin{itemize}
   \item This is the first item of an itemized list.
         Each item in the list is marked with a ``tick''.
         You don't have to worry about what kind of tick
         mark is used.

   \item This is the second item of the list.  It
         contains another list nested inside it.  The inner
         list is an \emph{enumerated} list.
         \begin{enumerate}
            \item This is the first item of an enumerated
                  list that is nested within the itemized list.

            \item This is the second item of the inner list.
                  \LaTeX\ allows you to nest lists deeper than
                  you really should.
         \end{enumerate}
         This is the rest of the second item of the outer
         list.  It is no more interesting than any other
         part of the item.
   \item This is the third item of the list.
\end{itemize}
You can even display poetry.
\begin{verse}
   There is an environment
    for verse \\             % The \\ command separates lines
   Whose features some poets % within a stanza.
   will curse.

                             % One or more blank lines separate stanzas.

   For instead of making\\
   Them do \emph{all} line breaking, \\
   It allows them to put too many words on a line when they'd rather be
   forced to be terse.
\end{verse}

Mathematical formulas may also be displayed.  A
displayed formula
is
one-line long; multiline
formulas require special formatting instructions.
   \[  \ip{\Gamma}{\psi'} = x'' + y^{2} + z_{i}^{n}\]
Don't start a paragraph with a displayed equation,
nor make one a paragraph by itself.

\end{document}               % End of document.

\section{Bench Huge Section 1296}
% synthetic bench section 1296
% This is a sample LaTeX input file.  (Version of 12 August 2004.)
%
% A '%' character causes TeX to ignore all remaining text on the line,
% and is used for comments like this one.

\documentclass{article}      % Specifies the document class

                             % The preamble begins here.
\title{An Example Document}  % Declares the document's title.
\author{Leslie Lamport}      % Declares the author's name.
\date{January 21, 1994}      % Deleting this command produces today's date.

\newcommand{\ip}[2]{(#1, #2)}
                             % Defines \ip{arg1}{arg2} to mean
                             % (arg1, arg2).

%\newcommand{\ip}[2]{\langle #1 | #2\rangle}
                             % This is an alternative definition of
                             % \ip that is commented out.

\begin{document}             % End of preamble and beginning of text.

\maketitle                   % Produces the title.

This is an example input file.  Comparing it with
the output it generates can show you how to
produce a simple document of your own.

\section{Ordinary Text}      % Produces section heading.  Lower-level
                             % sections are begun with similar
                             % \subsection and \subsubsection commands.

The ends  of words and sentences are marked
  by   spaces. It  doesn't matter how many
spaces    you type; one is as good as 100.  The
end of   a line counts as a space.

One   or more   blank lines denote the  end
of  a paragraph.

Since any number of consecutive spaces are treated
like a single one, the formatting of the input
file makes no difference to
      \LaTeX,                % The \LaTeX command generates the LaTeX logo.
but it makes a difference to you.  When you use
\LaTeX, making your input file as easy to read
as possible will be a great help as you write
your document and when you change it.  This sample
file shows how you can add comments to your own input
file.

Because printing is different from typewriting,
there are a number of things that you have to do
differently when preparing an input file than if
you were just typing the document directly.
Quotation marks like
       ``this''
have to be handled specially, as do quotes within
quotes:
       ``\,`this'            % \, separates the double and single quote.
        is what I just
        wrote, not  `that'\,''.

Dashes come in three sizes: an
       intra-word
dash, a medium dash for number ranges like
       1--2,
and a punctuation
       dash---like
this.

A sentence-ending space should be larger than the
space between words within a sentence.  You
sometimes have to type special commands in
conjunction with punctuation characters to get
this right, as in the following sentence.
       Gnats, gnus, etc.\ all  % `\ ' makes an inter-word space.
       begin with G\@.         % \@ marks end-of-sentence punctuation.
You should check the spaces after periods when
reading your output to make sure you haven't
forgotten any special cases.  Generating an
ellipsis
       \ldots\               % `\ ' is needed after `\ldots' because TeX
                             % ignores spaces after command names like \ldots
                             % made from \ + letters.
                             %
                             % Note how a `%' character causes TeX to ignore
                             % the end of the input line, so these blank lines
                             % do not start a new paragraph.
                             %
with the right spacing around the periods requires
a special command.

\LaTeX\ interprets some common characters as
commands, so you must type special commands to
generate them.  These characters include the
following:
       \$ \& \% \# \{ and \}.

In printing, text is usually emphasized with an
       \emph{italic}
type style.

\begin{em}
   A long segment of text can also be emphasized
   in this way.  Text within such a segment can be
   given \emph{additional} emphasis.
\end{em}

It is sometimes necessary to prevent \LaTeX\ from
breaking a line where it might otherwise do so.
This may be at a space, as between the ``Mr.''\ and
``Jones'' in
       ``Mr.~Jones'',        % ~ produces an unbreakable interword space.
or within a word---especially when the word is a
symbol like
       \mbox{\emph{itemnum}}
that makes little sense when hyphenated across
lines.

Footnotes\footnote{This is an example of a footnote.}
pose no problem.

\LaTeX\ is good at typesetting mathematical formulas
like
       \( x-3y + z = 7 \)
or
       \( a_{1} > x^{2n} + y^{2n} > x' \)
or
       \( \ip{A}{B} = \sum_{i} a_{i} b_{i} \).
The spaces you type in a formula are
ignored.  Remember that a letter like
       $x$                   % $ ... $  and  \( ... \)  are equivalent
is a formula when it denotes a mathematical
symbol, and it should be typed as one.

\section{Displayed Text}

Text is displayed by indenting it from the left
margin.  Quotations are commonly displayed.  There
are short quotations
\begin{quote}
   This is a short quotation.  It consists of a
   single paragraph of text.  See how it is formatted.
\end{quote}
and longer ones.
\begin{quotation}
   This is a longer quotation.  It consists of two
   paragraphs of text, neither of which are
   particularly interesting.

   This is the second paragraph of the quotation.  It
   is just as dull as the first paragraph.
\end{quotation}
Another frequently-displayed structure is a list.
The following is an example of an \emph{itemized}
list.
\begin{itemize}
   \item This is the first item of an itemized list.
         Each item in the list is marked with a ``tick''.
         You don't have to worry about what kind of tick
         mark is used.

   \item This is the second item of the list.  It
         contains another list nested inside it.  The inner
         list is an \emph{enumerated} list.
         \begin{enumerate}
            \item This is the first item of an enumerated
                  list that is nested within the itemized list.

            \item This is the second item of the inner list.
                  \LaTeX\ allows you to nest lists deeper than
                  you really should.
         \end{enumerate}
         This is the rest of the second item of the outer
         list.  It is no more interesting than any other
         part of the item.
   \item This is the third item of the list.
\end{itemize}
You can even display poetry.
\begin{verse}
   There is an environment
    for verse \\             % The \\ command separates lines
   Whose features some poets % within a stanza.
   will curse.

                             % One or more blank lines separate stanzas.

   For instead of making\\
   Them do \emph{all} line breaking, \\
   It allows them to put too many words on a line when they'd rather be
   forced to be terse.
\end{verse}

Mathematical formulas may also be displayed.  A
displayed formula
is
one-line long; multiline
formulas require special formatting instructions.
   \[  \ip{\Gamma}{\psi'} = x'' + y^{2} + z_{i}^{n}\]
Don't start a paragraph with a displayed equation,
nor make one a paragraph by itself.

\end{document}               % End of document.

\section{Bench Huge Section 1297}
% synthetic bench section 1297
% This is a sample LaTeX input file.  (Version of 12 August 2004.)
%
% A '%' character causes TeX to ignore all remaining text on the line,
% and is used for comments like this one.

\documentclass{article}      % Specifies the document class

                             % The preamble begins here.
\title{An Example Document}  % Declares the document's title.
\author{Leslie Lamport}      % Declares the author's name.
\date{January 21, 1994}      % Deleting this command produces today's date.

\newcommand{\ip}[2]{(#1, #2)}
                             % Defines \ip{arg1}{arg2} to mean
                             % (arg1, arg2).

%\newcommand{\ip}[2]{\langle #1 | #2\rangle}
                             % This is an alternative definition of
                             % \ip that is commented out.

\begin{document}             % End of preamble and beginning of text.

\maketitle                   % Produces the title.

This is an example input file.  Comparing it with
the output it generates can show you how to
produce a simple document of your own.

\section{Ordinary Text}      % Produces section heading.  Lower-level
                             % sections are begun with similar
                             % \subsection and \subsubsection commands.

The ends  of words and sentences are marked
  by   spaces. It  doesn't matter how many
spaces    you type; one is as good as 100.  The
end of   a line counts as a space.

One   or more   blank lines denote the  end
of  a paragraph.

Since any number of consecutive spaces are treated
like a single one, the formatting of the input
file makes no difference to
      \LaTeX,                % The \LaTeX command generates the LaTeX logo.
but it makes a difference to you.  When you use
\LaTeX, making your input file as easy to read
as possible will be a great help as you write
your document and when you change it.  This sample
file shows how you can add comments to your own input
file.

Because printing is different from typewriting,
there are a number of things that you have to do
differently when preparing an input file than if
you were just typing the document directly.
Quotation marks like
       ``this''
have to be handled specially, as do quotes within
quotes:
       ``\,`this'            % \, separates the double and single quote.
        is what I just
        wrote, not  `that'\,''.

Dashes come in three sizes: an
       intra-word
dash, a medium dash for number ranges like
       1--2,
and a punctuation
       dash---like
this.

A sentence-ending space should be larger than the
space between words within a sentence.  You
sometimes have to type special commands in
conjunction with punctuation characters to get
this right, as in the following sentence.
       Gnats, gnus, etc.\ all  % `\ ' makes an inter-word space.
       begin with G\@.         % \@ marks end-of-sentence punctuation.
You should check the spaces after periods when
reading your output to make sure you haven't
forgotten any special cases.  Generating an
ellipsis
       \ldots\               % `\ ' is needed after `\ldots' because TeX
                             % ignores spaces after command names like \ldots
                             % made from \ + letters.
                             %
                             % Note how a `%' character causes TeX to ignore
                             % the end of the input line, so these blank lines
                             % do not start a new paragraph.
                             %
with the right spacing around the periods requires
a special command.

\LaTeX\ interprets some common characters as
commands, so you must type special commands to
generate them.  These characters include the
following:
       \$ \& \% \# \{ and \}.

In printing, text is usually emphasized with an
       \emph{italic}
type style.

\begin{em}
   A long segment of text can also be emphasized
   in this way.  Text within such a segment can be
   given \emph{additional} emphasis.
\end{em}

It is sometimes necessary to prevent \LaTeX\ from
breaking a line where it might otherwise do so.
This may be at a space, as between the ``Mr.''\ and
``Jones'' in
       ``Mr.~Jones'',        % ~ produces an unbreakable interword space.
or within a word---especially when the word is a
symbol like
       \mbox{\emph{itemnum}}
that makes little sense when hyphenated across
lines.

Footnotes\footnote{This is an example of a footnote.}
pose no problem.

\LaTeX\ is good at typesetting mathematical formulas
like
       \( x-3y + z = 7 \)
or
       \( a_{1} > x^{2n} + y^{2n} > x' \)
or
       \( \ip{A}{B} = \sum_{i} a_{i} b_{i} \).
The spaces you type in a formula are
ignored.  Remember that a letter like
       $x$                   % $ ... $  and  \( ... \)  are equivalent
is a formula when it denotes a mathematical
symbol, and it should be typed as one.

\section{Displayed Text}

Text is displayed by indenting it from the left
margin.  Quotations are commonly displayed.  There
are short quotations
\begin{quote}
   This is a short quotation.  It consists of a
   single paragraph of text.  See how it is formatted.
\end{quote}
and longer ones.
\begin{quotation}
   This is a longer quotation.  It consists of two
   paragraphs of text, neither of which are
   particularly interesting.

   This is the second paragraph of the quotation.  It
   is just as dull as the first paragraph.
\end{quotation}
Another frequently-displayed structure is a list.
The following is an example of an \emph{itemized}
list.
\begin{itemize}
   \item This is the first item of an itemized list.
         Each item in the list is marked with a ``tick''.
         You don't have to worry about what kind of tick
         mark is used.

   \item This is the second item of the list.  It
         contains another list nested inside it.  The inner
         list is an \emph{enumerated} list.
         \begin{enumerate}
            \item This is the first item of an enumerated
                  list that is nested within the itemized list.

            \item This is the second item of the inner list.
                  \LaTeX\ allows you to nest lists deeper than
                  you really should.
         \end{enumerate}
         This is the rest of the second item of the outer
         list.  It is no more interesting than any other
         part of the item.
   \item This is the third item of the list.
\end{itemize}
You can even display poetry.
\begin{verse}
   There is an environment
    for verse \\             % The \\ command separates lines
   Whose features some poets % within a stanza.
   will curse.

                             % One or more blank lines separate stanzas.

   For instead of making\\
   Them do \emph{all} line breaking, \\
   It allows them to put too many words on a line when they'd rather be
   forced to be terse.
\end{verse}

Mathematical formulas may also be displayed.  A
displayed formula
is
one-line long; multiline
formulas require special formatting instructions.
   \[  \ip{\Gamma}{\psi'} = x'' + y^{2} + z_{i}^{n}\]
Don't start a paragraph with a displayed equation,
nor make one a paragraph by itself.

\end{document}               % End of document.

\section{Bench Huge Section 1298}
% synthetic bench section 1298
% This is a sample LaTeX input file.  (Version of 12 August 2004.)
%
% A '%' character causes TeX to ignore all remaining text on the line,
% and is used for comments like this one.

\documentclass{article}      % Specifies the document class

                             % The preamble begins here.
\title{An Example Document}  % Declares the document's title.
\author{Leslie Lamport}      % Declares the author's name.
\date{January 21, 1994}      % Deleting this command produces today's date.

\newcommand{\ip}[2]{(#1, #2)}
                             % Defines \ip{arg1}{arg2} to mean
                             % (arg1, arg2).

%\newcommand{\ip}[2]{\langle #1 | #2\rangle}
                             % This is an alternative definition of
                             % \ip that is commented out.

\begin{document}             % End of preamble and beginning of text.

\maketitle                   % Produces the title.

This is an example input file.  Comparing it with
the output it generates can show you how to
produce a simple document of your own.

\section{Ordinary Text}      % Produces section heading.  Lower-level
                             % sections are begun with similar
                             % \subsection and \subsubsection commands.

The ends  of words and sentences are marked
  by   spaces. It  doesn't matter how many
spaces    you type; one is as good as 100.  The
end of   a line counts as a space.

One   or more   blank lines denote the  end
of  a paragraph.

Since any number of consecutive spaces are treated
like a single one, the formatting of the input
file makes no difference to
      \LaTeX,                % The \LaTeX command generates the LaTeX logo.
but it makes a difference to you.  When you use
\LaTeX, making your input file as easy to read
as possible will be a great help as you write
your document and when you change it.  This sample
file shows how you can add comments to your own input
file.

Because printing is different from typewriting,
there are a number of things that you have to do
differently when preparing an input file than if
you were just typing the document directly.
Quotation marks like
       ``this''
have to be handled specially, as do quotes within
quotes:
       ``\,`this'            % \, separates the double and single quote.
        is what I just
        wrote, not  `that'\,''.

Dashes come in three sizes: an
       intra-word
dash, a medium dash for number ranges like
       1--2,
and a punctuation
       dash---like
this.

A sentence-ending space should be larger than the
space between words within a sentence.  You
sometimes have to type special commands in
conjunction with punctuation characters to get
this right, as in the following sentence.
       Gnats, gnus, etc.\ all  % `\ ' makes an inter-word space.
       begin with G\@.         % \@ marks end-of-sentence punctuation.
You should check the spaces after periods when
reading your output to make sure you haven't
forgotten any special cases.  Generating an
ellipsis
       \ldots\               % `\ ' is needed after `\ldots' because TeX
                             % ignores spaces after command names like \ldots
                             % made from \ + letters.
                             %
                             % Note how a `%' character causes TeX to ignore
                             % the end of the input line, so these blank lines
                             % do not start a new paragraph.
                             %
with the right spacing around the periods requires
a special command.

\LaTeX\ interprets some common characters as
commands, so you must type special commands to
generate them.  These characters include the
following:
       \$ \& \% \# \{ and \}.

In printing, text is usually emphasized with an
       \emph{italic}
type style.

\begin{em}
   A long segment of text can also be emphasized
   in this way.  Text within such a segment can be
   given \emph{additional} emphasis.
\end{em}

It is sometimes necessary to prevent \LaTeX\ from
breaking a line where it might otherwise do so.
This may be at a space, as between the ``Mr.''\ and
``Jones'' in
       ``Mr.~Jones'',        % ~ produces an unbreakable interword space.
or within a word---especially when the word is a
symbol like
       \mbox{\emph{itemnum}}
that makes little sense when hyphenated across
lines.

Footnotes\footnote{This is an example of a footnote.}
pose no problem.

\LaTeX\ is good at typesetting mathematical formulas
like
       \( x-3y + z = 7 \)
or
       \( a_{1} > x^{2n} + y^{2n} > x' \)
or
       \( \ip{A}{B} = \sum_{i} a_{i} b_{i} \).
The spaces you type in a formula are
ignored.  Remember that a letter like
       $x$                   % $ ... $  and  \( ... \)  are equivalent
is a formula when it denotes a mathematical
symbol, and it should be typed as one.

\section{Displayed Text}

Text is displayed by indenting it from the left
margin.  Quotations are commonly displayed.  There
are short quotations
\begin{quote}
   This is a short quotation.  It consists of a
   single paragraph of text.  See how it is formatted.
\end{quote}
and longer ones.
\begin{quotation}
   This is a longer quotation.  It consists of two
   paragraphs of text, neither of which are
   particularly interesting.

   This is the second paragraph of the quotation.  It
   is just as dull as the first paragraph.
\end{quotation}
Another frequently-displayed structure is a list.
The following is an example of an \emph{itemized}
list.
\begin{itemize}
   \item This is the first item of an itemized list.
         Each item in the list is marked with a ``tick''.
         You don't have to worry about what kind of tick
         mark is used.

   \item This is the second item of the list.  It
         contains another list nested inside it.  The inner
         list is an \emph{enumerated} list.
         \begin{enumerate}
            \item This is the first item of an enumerated
                  list that is nested within the itemized list.

            \item This is the second item of the inner list.
                  \LaTeX\ allows you to nest lists deeper than
                  you really should.
         \end{enumerate}
         This is the rest of the second item of the outer
         list.  It is no more interesting than any other
         part of the item.
   \item This is the third item of the list.
\end{itemize}
You can even display poetry.
\begin{verse}
   There is an environment
    for verse \\             % The \\ command separates lines
   Whose features some poets % within a stanza.
   will curse.

                             % One or more blank lines separate stanzas.

   For instead of making\\
   Them do \emph{all} line breaking, \\
   It allows them to put too many words on a line when they'd rather be
   forced to be terse.
\end{verse}

Mathematical formulas may also be displayed.  A
displayed formula
is
one-line long; multiline
formulas require special formatting instructions.
   \[  \ip{\Gamma}{\psi'} = x'' + y^{2} + z_{i}^{n}\]
Don't start a paragraph with a displayed equation,
nor make one a paragraph by itself.

\end{document}               % End of document.

\section{Bench Huge Section 1299}
% synthetic bench section 1299
% This is a sample LaTeX input file.  (Version of 12 August 2004.)
%
% A '%' character causes TeX to ignore all remaining text on the line,
% and is used for comments like this one.

\documentclass{article}      % Specifies the document class

                             % The preamble begins here.
\title{An Example Document}  % Declares the document's title.
\author{Leslie Lamport}      % Declares the author's name.
\date{January 21, 1994}      % Deleting this command produces today's date.

\newcommand{\ip}[2]{(#1, #2)}
                             % Defines \ip{arg1}{arg2} to mean
                             % (arg1, arg2).

%\newcommand{\ip}[2]{\langle #1 | #2\rangle}
                             % This is an alternative definition of
                             % \ip that is commented out.

\begin{document}             % End of preamble and beginning of text.

\maketitle                   % Produces the title.

This is an example input file.  Comparing it with
the output it generates can show you how to
produce a simple document of your own.

\section{Ordinary Text}      % Produces section heading.  Lower-level
                             % sections are begun with similar
                             % \subsection and \subsubsection commands.

The ends  of words and sentences are marked
  by   spaces. It  doesn't matter how many
spaces    you type; one is as good as 100.  The
end of   a line counts as a space.

One   or more   blank lines denote the  end
of  a paragraph.

Since any number of consecutive spaces are treated
like a single one, the formatting of the input
file makes no difference to
      \LaTeX,                % The \LaTeX command generates the LaTeX logo.
but it makes a difference to you.  When you use
\LaTeX, making your input file as easy to read
as possible will be a great help as you write
your document and when you change it.  This sample
file shows how you can add comments to your own input
file.

Because printing is different from typewriting,
there are a number of things that you have to do
differently when preparing an input file than if
you were just typing the document directly.
Quotation marks like
       ``this''
have to be handled specially, as do quotes within
quotes:
       ``\,`this'            % \, separates the double and single quote.
        is what I just
        wrote, not  `that'\,''.

Dashes come in three sizes: an
       intra-word
dash, a medium dash for number ranges like
       1--2,
and a punctuation
       dash---like
this.

A sentence-ending space should be larger than the
space between words within a sentence.  You
sometimes have to type special commands in
conjunction with punctuation characters to get
this right, as in the following sentence.
       Gnats, gnus, etc.\ all  % `\ ' makes an inter-word space.
       begin with G\@.         % \@ marks end-of-sentence punctuation.
You should check the spaces after periods when
reading your output to make sure you haven't
forgotten any special cases.  Generating an
ellipsis
       \ldots\               % `\ ' is needed after `\ldots' because TeX
                             % ignores spaces after command names like \ldots
                             % made from \ + letters.
                             %
                             % Note how a `%' character causes TeX to ignore
                             % the end of the input line, so these blank lines
                             % do not start a new paragraph.
                             %
with the right spacing around the periods requires
a special command.

\LaTeX\ interprets some common characters as
commands, so you must type special commands to
generate them.  These characters include the
following:
       \$ \& \% \# \{ and \}.

In printing, text is usually emphasized with an
       \emph{italic}
type style.

\begin{em}
   A long segment of text can also be emphasized
   in this way.  Text within such a segment can be
   given \emph{additional} emphasis.
\end{em}

It is sometimes necessary to prevent \LaTeX\ from
breaking a line where it might otherwise do so.
This may be at a space, as between the ``Mr.''\ and
``Jones'' in
       ``Mr.~Jones'',        % ~ produces an unbreakable interword space.
or within a word---especially when the word is a
symbol like
       \mbox{\emph{itemnum}}
that makes little sense when hyphenated across
lines.

Footnotes\footnote{This is an example of a footnote.}
pose no problem.

\LaTeX\ is good at typesetting mathematical formulas
like
       \( x-3y + z = 7 \)
or
       \( a_{1} > x^{2n} + y^{2n} > x' \)
or
       \( \ip{A}{B} = \sum_{i} a_{i} b_{i} \).
The spaces you type in a formula are
ignored.  Remember that a letter like
       $x$                   % $ ... $  and  \( ... \)  are equivalent
is a formula when it denotes a mathematical
symbol, and it should be typed as one.

\section{Displayed Text}

Text is displayed by indenting it from the left
margin.  Quotations are commonly displayed.  There
are short quotations
\begin{quote}
   This is a short quotation.  It consists of a
   single paragraph of text.  See how it is formatted.
\end{quote}
and longer ones.
\begin{quotation}
   This is a longer quotation.  It consists of two
   paragraphs of text, neither of which are
   particularly interesting.

   This is the second paragraph of the quotation.  It
   is just as dull as the first paragraph.
\end{quotation}
Another frequently-displayed structure is a list.
The following is an example of an \emph{itemized}
list.
\begin{itemize}
   \item This is the first item of an itemized list.
         Each item in the list is marked with a ``tick''.
         You don't have to worry about what kind of tick
         mark is used.

   \item This is the second item of the list.  It
         contains another list nested inside it.  The inner
         list is an \emph{enumerated} list.
         \begin{enumerate}
            \item This is the first item of an enumerated
                  list that is nested within the itemized list.

            \item This is the second item of the inner list.
                  \LaTeX\ allows you to nest lists deeper than
                  you really should.
         \end{enumerate}
         This is the rest of the second item of the outer
         list.  It is no more interesting than any other
         part of the item.
   \item This is the third item of the list.
\end{itemize}
You can even display poetry.
\begin{verse}
   There is an environment
    for verse \\             % The \\ command separates lines
   Whose features some poets % within a stanza.
   will curse.

                             % One or more blank lines separate stanzas.

   For instead of making\\
   Them do \emph{all} line breaking, \\
   It allows them to put too many words on a line when they'd rather be
   forced to be terse.
\end{verse}

Mathematical formulas may also be displayed.  A
displayed formula
is
one-line long; multiline
formulas require special formatting instructions.
   \[  \ip{\Gamma}{\psi'} = x'' + y^{2} + z_{i}^{n}\]
Don't start a paragraph with a displayed equation,
nor make one a paragraph by itself.

\end{document}               % End of document.

\section{Bench Huge Section 1300}
% synthetic bench section 1300
% This is a sample LaTeX input file.  (Version of 12 August 2004.)
%
% A '%' character causes TeX to ignore all remaining text on the line,
% and is used for comments like this one.

\documentclass{article}      % Specifies the document class

                             % The preamble begins here.
\title{An Example Document}  % Declares the document's title.
\author{Leslie Lamport}      % Declares the author's name.
\date{January 21, 1994}      % Deleting this command produces today's date.

\newcommand{\ip}[2]{(#1, #2)}
                             % Defines \ip{arg1}{arg2} to mean
                             % (arg1, arg2).

%\newcommand{\ip}[2]{\langle #1 | #2\rangle}
                             % This is an alternative definition of
                             % \ip that is commented out.

\begin{document}             % End of preamble and beginning of text.

\maketitle                   % Produces the title.

This is an example input file.  Comparing it with
the output it generates can show you how to
produce a simple document of your own.

\section{Ordinary Text}      % Produces section heading.  Lower-level
                             % sections are begun with similar
                             % \subsection and \subsubsection commands.

The ends  of words and sentences are marked
  by   spaces. It  doesn't matter how many
spaces    you type; one is as good as 100.  The
end of   a line counts as a space.

One   or more   blank lines denote the  end
of  a paragraph.

Since any number of consecutive spaces are treated
like a single one, the formatting of the input
file makes no difference to
      \LaTeX,                % The \LaTeX command generates the LaTeX logo.
but it makes a difference to you.  When you use
\LaTeX, making your input file as easy to read
as possible will be a great help as you write
your document and when you change it.  This sample
file shows how you can add comments to your own input
file.

Because printing is different from typewriting,
there are a number of things that you have to do
differently when preparing an input file than if
you were just typing the document directly.
Quotation marks like
       ``this''
have to be handled specially, as do quotes within
quotes:
       ``\,`this'            % \, separates the double and single quote.
        is what I just
        wrote, not  `that'\,''.

Dashes come in three sizes: an
       intra-word
dash, a medium dash for number ranges like
       1--2,
and a punctuation
       dash---like
this.

A sentence-ending space should be larger than the
space between words within a sentence.  You
sometimes have to type special commands in
conjunction with punctuation characters to get
this right, as in the following sentence.
       Gnats, gnus, etc.\ all  % `\ ' makes an inter-word space.
       begin with G\@.         % \@ marks end-of-sentence punctuation.
You should check the spaces after periods when
reading your output to make sure you haven't
forgotten any special cases.  Generating an
ellipsis
       \ldots\               % `\ ' is needed after `\ldots' because TeX
                             % ignores spaces after command names like \ldots
                             % made from \ + letters.
                             %
                             % Note how a `%' character causes TeX to ignore
                             % the end of the input line, so these blank lines
                             % do not start a new paragraph.
                             %
with the right spacing around the periods requires
a special command.

\LaTeX\ interprets some common characters as
commands, so you must type special commands to
generate them.  These characters include the
following:
       \$ \& \% \# \{ and \}.

In printing, text is usually emphasized with an
       \emph{italic}
type style.

\begin{em}
   A long segment of text can also be emphasized
   in this way.  Text within such a segment can be
   given \emph{additional} emphasis.
\end{em}

It is sometimes necessary to prevent \LaTeX\ from
breaking a line where it might otherwise do so.
This may be at a space, as between the ``Mr.''\ and
``Jones'' in
       ``Mr.~Jones'',        % ~ produces an unbreakable interword space.
or within a word---especially when the word is a
symbol like
       \mbox{\emph{itemnum}}
that makes little sense when hyphenated across
lines.

Footnotes\footnote{This is an example of a footnote.}
pose no problem.

\LaTeX\ is good at typesetting mathematical formulas
like
       \( x-3y + z = 7 \)
or
       \( a_{1} > x^{2n} + y^{2n} > x' \)
or
       \( \ip{A}{B} = \sum_{i} a_{i} b_{i} \).
The spaces you type in a formula are
ignored.  Remember that a letter like
       $x$                   % $ ... $  and  \( ... \)  are equivalent
is a formula when it denotes a mathematical
symbol, and it should be typed as one.

\section{Displayed Text}

Text is displayed by indenting it from the left
margin.  Quotations are commonly displayed.  There
are short quotations
\begin{quote}
   This is a short quotation.  It consists of a
   single paragraph of text.  See how it is formatted.
\end{quote}
and longer ones.
\begin{quotation}
   This is a longer quotation.  It consists of two
   paragraphs of text, neither of which are
   particularly interesting.

   This is the second paragraph of the quotation.  It
   is just as dull as the first paragraph.
\end{quotation}
Another frequently-displayed structure is a list.
The following is an example of an \emph{itemized}
list.
\begin{itemize}
   \item This is the first item of an itemized list.
         Each item in the list is marked with a ``tick''.
         You don't have to worry about what kind of tick
         mark is used.

   \item This is the second item of the list.  It
         contains another list nested inside it.  The inner
         list is an \emph{enumerated} list.
         \begin{enumerate}
            \item This is the first item of an enumerated
                  list that is nested within the itemized list.

            \item This is the second item of the inner list.
                  \LaTeX\ allows you to nest lists deeper than
                  you really should.
         \end{enumerate}
         This is the rest of the second item of the outer
         list.  It is no more interesting than any other
         part of the item.
   \item This is the third item of the list.
\end{itemize}
You can even display poetry.
\begin{verse}
   There is an environment
    for verse \\             % The \\ command separates lines
   Whose features some poets % within a stanza.
   will curse.

                             % One or more blank lines separate stanzas.

   For instead of making\\
   Them do \emph{all} line breaking, \\
   It allows them to put too many words on a line when they'd rather be
   forced to be terse.
\end{verse}

Mathematical formulas may also be displayed.  A
displayed formula
is
one-line long; multiline
formulas require special formatting instructions.
   \[  \ip{\Gamma}{\psi'} = x'' + y^{2} + z_{i}^{n}\]
Don't start a paragraph with a displayed equation,
nor make one a paragraph by itself.

\end{document}               % End of document.

\section{Bench Huge Section 1301}
% synthetic bench section 1301
% This is a sample LaTeX input file.  (Version of 12 August 2004.)
%
% A '%' character causes TeX to ignore all remaining text on the line,
% and is used for comments like this one.

\documentclass{article}      % Specifies the document class

                             % The preamble begins here.
\title{An Example Document}  % Declares the document's title.
\author{Leslie Lamport}      % Declares the author's name.
\date{January 21, 1994}      % Deleting this command produces today's date.

\newcommand{\ip}[2]{(#1, #2)}
                             % Defines \ip{arg1}{arg2} to mean
                             % (arg1, arg2).

%\newcommand{\ip}[2]{\langle #1 | #2\rangle}
                             % This is an alternative definition of
                             % \ip that is commented out.

\begin{document}             % End of preamble and beginning of text.

\maketitle                   % Produces the title.

This is an example input file.  Comparing it with
the output it generates can show you how to
produce a simple document of your own.

\section{Ordinary Text}      % Produces section heading.  Lower-level
                             % sections are begun with similar
                             % \subsection and \subsubsection commands.

The ends  of words and sentences are marked
  by   spaces. It  doesn't matter how many
spaces    you type; one is as good as 100.  The
end of   a line counts as a space.

One   or more   blank lines denote the  end
of  a paragraph.

Since any number of consecutive spaces are treated
like a single one, the formatting of the input
file makes no difference to
      \LaTeX,                % The \LaTeX command generates the LaTeX logo.
but it makes a difference to you.  When you use
\LaTeX, making your input file as easy to read
as possible will be a great help as you write
your document and when you change it.  This sample
file shows how you can add comments to your own input
file.

Because printing is different from typewriting,
there are a number of things that you have to do
differently when preparing an input file than if
you were just typing the document directly.
Quotation marks like
       ``this''
have to be handled specially, as do quotes within
quotes:
       ``\,`this'            % \, separates the double and single quote.
        is what I just
        wrote, not  `that'\,''.

Dashes come in three sizes: an
       intra-word
dash, a medium dash for number ranges like
       1--2,
and a punctuation
       dash---like
this.

A sentence-ending space should be larger than the
space between words within a sentence.  You
sometimes have to type special commands in
conjunction with punctuation characters to get
this right, as in the following sentence.
       Gnats, gnus, etc.\ all  % `\ ' makes an inter-word space.
       begin with G\@.         % \@ marks end-of-sentence punctuation.
You should check the spaces after periods when
reading your output to make sure you haven't
forgotten any special cases.  Generating an
ellipsis
       \ldots\               % `\ ' is needed after `\ldots' because TeX
                             % ignores spaces after command names like \ldots
                             % made from \ + letters.
                             %
                             % Note how a `%' character causes TeX to ignore
                             % the end of the input line, so these blank lines
                             % do not start a new paragraph.
                             %
with the right spacing around the periods requires
a special command.

\LaTeX\ interprets some common characters as
commands, so you must type special commands to
generate them.  These characters include the
following:
       \$ \& \% \# \{ and \}.

In printing, text is usually emphasized with an
       \emph{italic}
type style.

\begin{em}
   A long segment of text can also be emphasized
   in this way.  Text within such a segment can be
   given \emph{additional} emphasis.
\end{em}

It is sometimes necessary to prevent \LaTeX\ from
breaking a line where it might otherwise do so.
This may be at a space, as between the ``Mr.''\ and
``Jones'' in
       ``Mr.~Jones'',        % ~ produces an unbreakable interword space.
or within a word---especially when the word is a
symbol like
       \mbox{\emph{itemnum}}
that makes little sense when hyphenated across
lines.

Footnotes\footnote{This is an example of a footnote.}
pose no problem.

\LaTeX\ is good at typesetting mathematical formulas
like
       \( x-3y + z = 7 \)
or
       \( a_{1} > x^{2n} + y^{2n} > x' \)
or
       \( \ip{A}{B} = \sum_{i} a_{i} b_{i} \).
The spaces you type in a formula are
ignored.  Remember that a letter like
       $x$                   % $ ... $  and  \( ... \)  are equivalent
is a formula when it denotes a mathematical
symbol, and it should be typed as one.

\section{Displayed Text}

Text is displayed by indenting it from the left
margin.  Quotations are commonly displayed.  There
are short quotations
\begin{quote}
   This is a short quotation.  It consists of a
   single paragraph of text.  See how it is formatted.
\end{quote}
and longer ones.
\begin{quotation}
   This is a longer quotation.  It consists of two
   paragraphs of text, neither of which are
   particularly interesting.

   This is the second paragraph of the quotation.  It
   is just as dull as the first paragraph.
\end{quotation}
Another frequently-displayed structure is a list.
The following is an example of an \emph{itemized}
list.
\begin{itemize}
   \item This is the first item of an itemized list.
         Each item in the list is marked with a ``tick''.
         You don't have to worry about what kind of tick
         mark is used.

   \item This is the second item of the list.  It
         contains another list nested inside it.  The inner
         list is an \emph{enumerated} list.
         \begin{enumerate}
            \item This is the first item of an enumerated
                  list that is nested within the itemized list.

            \item This is the second item of the inner list.
                  \LaTeX\ allows you to nest lists deeper than
                  you really should.
         \end{enumerate}
         This is the rest of the second item of the outer
         list.  It is no more interesting than any other
         part of the item.
   \item This is the third item of the list.
\end{itemize}
You can even display poetry.
\begin{verse}
   There is an environment
    for verse \\             % The \\ command separates lines
   Whose features some poets % within a stanza.
   will curse.

                             % One or more blank lines separate stanzas.

   For instead of making\\
   Them do \emph{all} line breaking, \\
   It allows them to put too many words on a line when they'd rather be
   forced to be terse.
\end{verse}

Mathematical formulas may also be displayed.  A
displayed formula
is
one-line long; multiline
formulas require special formatting instructions.
   \[  \ip{\Gamma}{\psi'} = x'' + y^{2} + z_{i}^{n}\]
Don't start a paragraph with a displayed equation,
nor make one a paragraph by itself.

\end{document}               % End of document.

\section{Bench Huge Section 1302}
% synthetic bench section 1302
% This is a sample LaTeX input file.  (Version of 12 August 2004.)
%
% A '%' character causes TeX to ignore all remaining text on the line,
% and is used for comments like this one.

\documentclass{article}      % Specifies the document class

                             % The preamble begins here.
\title{An Example Document}  % Declares the document's title.
\author{Leslie Lamport}      % Declares the author's name.
\date{January 21, 1994}      % Deleting this command produces today's date.

\newcommand{\ip}[2]{(#1, #2)}
                             % Defines \ip{arg1}{arg2} to mean
                             % (arg1, arg2).

%\newcommand{\ip}[2]{\langle #1 | #2\rangle}
                             % This is an alternative definition of
                             % \ip that is commented out.

\begin{document}             % End of preamble and beginning of text.

\maketitle                   % Produces the title.

This is an example input file.  Comparing it with
the output it generates can show you how to
produce a simple document of your own.

\section{Ordinary Text}      % Produces section heading.  Lower-level
                             % sections are begun with similar
                             % \subsection and \subsubsection commands.

The ends  of words and sentences are marked
  by   spaces. It  doesn't matter how many
spaces    you type; one is as good as 100.  The
end of   a line counts as a space.

One   or more   blank lines denote the  end
of  a paragraph.

Since any number of consecutive spaces are treated
like a single one, the formatting of the input
file makes no difference to
      \LaTeX,                % The \LaTeX command generates the LaTeX logo.
but it makes a difference to you.  When you use
\LaTeX, making your input file as easy to read
as possible will be a great help as you write
your document and when you change it.  This sample
file shows how you can add comments to your own input
file.

Because printing is different from typewriting,
there are a number of things that you have to do
differently when preparing an input file than if
you were just typing the document directly.
Quotation marks like
       ``this''
have to be handled specially, as do quotes within
quotes:
       ``\,`this'            % \, separates the double and single quote.
        is what I just
        wrote, not  `that'\,''.

Dashes come in three sizes: an
       intra-word
dash, a medium dash for number ranges like
       1--2,
and a punctuation
       dash---like
this.

A sentence-ending space should be larger than the
space between words within a sentence.  You
sometimes have to type special commands in
conjunction with punctuation characters to get
this right, as in the following sentence.
       Gnats, gnus, etc.\ all  % `\ ' makes an inter-word space.
       begin with G\@.         % \@ marks end-of-sentence punctuation.
You should check the spaces after periods when
reading your output to make sure you haven't
forgotten any special cases.  Generating an
ellipsis
       \ldots\               % `\ ' is needed after `\ldots' because TeX
                             % ignores spaces after command names like \ldots
                             % made from \ + letters.
                             %
                             % Note how a `%' character causes TeX to ignore
                             % the end of the input line, so these blank lines
                             % do not start a new paragraph.
                             %
with the right spacing around the periods requires
a special command.

\LaTeX\ interprets some common characters as
commands, so you must type special commands to
generate them.  These characters include the
following:
       \$ \& \% \# \{ and \}.

In printing, text is usually emphasized with an
       \emph{italic}
type style.

\begin{em}
   A long segment of text can also be emphasized
   in this way.  Text within such a segment can be
   given \emph{additional} emphasis.
\end{em}

It is sometimes necessary to prevent \LaTeX\ from
breaking a line where it might otherwise do so.
This may be at a space, as between the ``Mr.''\ and
``Jones'' in
       ``Mr.~Jones'',        % ~ produces an unbreakable interword space.
or within a word---especially when the word is a
symbol like
       \mbox{\emph{itemnum}}
that makes little sense when hyphenated across
lines.

Footnotes\footnote{This is an example of a footnote.}
pose no problem.

\LaTeX\ is good at typesetting mathematical formulas
like
       \( x-3y + z = 7 \)
or
       \( a_{1} > x^{2n} + y^{2n} > x' \)
or
       \( \ip{A}{B} = \sum_{i} a_{i} b_{i} \).
The spaces you type in a formula are
ignored.  Remember that a letter like
       $x$                   % $ ... $  and  \( ... \)  are equivalent
is a formula when it denotes a mathematical
symbol, and it should be typed as one.

\section{Displayed Text}

Text is displayed by indenting it from the left
margin.  Quotations are commonly displayed.  There
are short quotations
\begin{quote}
   This is a short quotation.  It consists of a
   single paragraph of text.  See how it is formatted.
\end{quote}
and longer ones.
\begin{quotation}
   This is a longer quotation.  It consists of two
   paragraphs of text, neither of which are
   particularly interesting.

   This is the second paragraph of the quotation.  It
   is just as dull as the first paragraph.
\end{quotation}
Another frequently-displayed structure is a list.
The following is an example of an \emph{itemized}
list.
\begin{itemize}
   \item This is the first item of an itemized list.
         Each item in the list is marked with a ``tick''.
         You don't have to worry about what kind of tick
         mark is used.

   \item This is the second item of the list.  It
         contains another list nested inside it.  The inner
         list is an \emph{enumerated} list.
         \begin{enumerate}
            \item This is the first item of an enumerated
                  list that is nested within the itemized list.

            \item This is the second item of the inner list.
                  \LaTeX\ allows you to nest lists deeper than
                  you really should.
         \end{enumerate}
         This is the rest of the second item of the outer
         list.  It is no more interesting than any other
         part of the item.
   \item This is the third item of the list.
\end{itemize}
You can even display poetry.
\begin{verse}
   There is an environment
    for verse \\             % The \\ command separates lines
   Whose features some poets % within a stanza.
   will curse.

                             % One or more blank lines separate stanzas.

   For instead of making\\
   Them do \emph{all} line breaking, \\
   It allows them to put too many words on a line when they'd rather be
   forced to be terse.
\end{verse}

Mathematical formulas may also be displayed.  A
displayed formula
is
one-line long; multiline
formulas require special formatting instructions.
   \[  \ip{\Gamma}{\psi'} = x'' + y^{2} + z_{i}^{n}\]
Don't start a paragraph with a displayed equation,
nor make one a paragraph by itself.

\end{document}               % End of document.

\section{Bench Huge Section 1303}
% synthetic bench section 1303
% This is a sample LaTeX input file.  (Version of 12 August 2004.)
%
% A '%' character causes TeX to ignore all remaining text on the line,
% and is used for comments like this one.

\documentclass{article}      % Specifies the document class

                             % The preamble begins here.
\title{An Example Document}  % Declares the document's title.
\author{Leslie Lamport}      % Declares the author's name.
\date{January 21, 1994}      % Deleting this command produces today's date.

\newcommand{\ip}[2]{(#1, #2)}
                             % Defines \ip{arg1}{arg2} to mean
                             % (arg1, arg2).

%\newcommand{\ip}[2]{\langle #1 | #2\rangle}
                             % This is an alternative definition of
                             % \ip that is commented out.

\begin{document}             % End of preamble and beginning of text.

\maketitle                   % Produces the title.

This is an example input file.  Comparing it with
the output it generates can show you how to
produce a simple document of your own.

\section{Ordinary Text}      % Produces section heading.  Lower-level
                             % sections are begun with similar
                             % \subsection and \subsubsection commands.

The ends  of words and sentences are marked
  by   spaces. It  doesn't matter how many
spaces    you type; one is as good as 100.  The
end of   a line counts as a space.

One   or more   blank lines denote the  end
of  a paragraph.

Since any number of consecutive spaces are treated
like a single one, the formatting of the input
file makes no difference to
      \LaTeX,                % The \LaTeX command generates the LaTeX logo.
but it makes a difference to you.  When you use
\LaTeX, making your input file as easy to read
as possible will be a great help as you write
your document and when you change it.  This sample
file shows how you can add comments to your own input
file.

Because printing is different from typewriting,
there are a number of things that you have to do
differently when preparing an input file than if
you were just typing the document directly.
Quotation marks like
       ``this''
have to be handled specially, as do quotes within
quotes:
       ``\,`this'            % \, separates the double and single quote.
        is what I just
        wrote, not  `that'\,''.

Dashes come in three sizes: an
       intra-word
dash, a medium dash for number ranges like
       1--2,
and a punctuation
       dash---like
this.

A sentence-ending space should be larger than the
space between words within a sentence.  You
sometimes have to type special commands in
conjunction with punctuation characters to get
this right, as in the following sentence.
       Gnats, gnus, etc.\ all  % `\ ' makes an inter-word space.
       begin with G\@.         % \@ marks end-of-sentence punctuation.
You should check the spaces after periods when
reading your output to make sure you haven't
forgotten any special cases.  Generating an
ellipsis
       \ldots\               % `\ ' is needed after `\ldots' because TeX
                             % ignores spaces after command names like \ldots
                             % made from \ + letters.
                             %
                             % Note how a `%' character causes TeX to ignore
                             % the end of the input line, so these blank lines
                             % do not start a new paragraph.
                             %
with the right spacing around the periods requires
a special command.

\LaTeX\ interprets some common characters as
commands, so you must type special commands to
generate them.  These characters include the
following:
       \$ \& \% \# \{ and \}.

In printing, text is usually emphasized with an
       \emph{italic}
type style.

\begin{em}
   A long segment of text can also be emphasized
   in this way.  Text within such a segment can be
   given \emph{additional} emphasis.
\end{em}

It is sometimes necessary to prevent \LaTeX\ from
breaking a line where it might otherwise do so.
This may be at a space, as between the ``Mr.''\ and
``Jones'' in
       ``Mr.~Jones'',        % ~ produces an unbreakable interword space.
or within a word---especially when the word is a
symbol like
       \mbox{\emph{itemnum}}
that makes little sense when hyphenated across
lines.

Footnotes\footnote{This is an example of a footnote.}
pose no problem.

\LaTeX\ is good at typesetting mathematical formulas
like
       \( x-3y + z = 7 \)
or
       \( a_{1} > x^{2n} + y^{2n} > x' \)
or
       \( \ip{A}{B} = \sum_{i} a_{i} b_{i} \).
The spaces you type in a formula are
ignored.  Remember that a letter like
       $x$                   % $ ... $  and  \( ... \)  are equivalent
is a formula when it denotes a mathematical
symbol, and it should be typed as one.

\section{Displayed Text}

Text is displayed by indenting it from the left
margin.  Quotations are commonly displayed.  There
are short quotations
\begin{quote}
   This is a short quotation.  It consists of a
   single paragraph of text.  See how it is formatted.
\end{quote}
and longer ones.
\begin{quotation}
   This is a longer quotation.  It consists of two
   paragraphs of text, neither of which are
   particularly interesting.

   This is the second paragraph of the quotation.  It
   is just as dull as the first paragraph.
\end{quotation}
Another frequently-displayed structure is a list.
The following is an example of an \emph{itemized}
list.
\begin{itemize}
   \item This is the first item of an itemized list.
         Each item in the list is marked with a ``tick''.
         You don't have to worry about what kind of tick
         mark is used.

   \item This is the second item of the list.  It
         contains another list nested inside it.  The inner
         list is an \emph{enumerated} list.
         \begin{enumerate}
            \item This is the first item of an enumerated
                  list that is nested within the itemized list.

            \item This is the second item of the inner list.
                  \LaTeX\ allows you to nest lists deeper than
                  you really should.
         \end{enumerate}
         This is the rest of the second item of the outer
         list.  It is no more interesting than any other
         part of the item.
   \item This is the third item of the list.
\end{itemize}
You can even display poetry.
\begin{verse}
   There is an environment
    for verse \\             % The \\ command separates lines
   Whose features some poets % within a stanza.
   will curse.

                             % One or more blank lines separate stanzas.

   For instead of making\\
   Them do \emph{all} line breaking, \\
   It allows them to put too many words on a line when they'd rather be
   forced to be terse.
\end{verse}

Mathematical formulas may also be displayed.  A
displayed formula
is
one-line long; multiline
formulas require special formatting instructions.
   \[  \ip{\Gamma}{\psi'} = x'' + y^{2} + z_{i}^{n}\]
Don't start a paragraph with a displayed equation,
nor make one a paragraph by itself.

\end{document}               % End of document.

\section{Bench Huge Section 1304}
% synthetic bench section 1304
% This is a sample LaTeX input file.  (Version of 12 August 2004.)
%
% A '%' character causes TeX to ignore all remaining text on the line,
% and is used for comments like this one.

\documentclass{article}      % Specifies the document class

                             % The preamble begins here.
\title{An Example Document}  % Declares the document's title.
\author{Leslie Lamport}      % Declares the author's name.
\date{January 21, 1994}      % Deleting this command produces today's date.

\newcommand{\ip}[2]{(#1, #2)}
                             % Defines \ip{arg1}{arg2} to mean
                             % (arg1, arg2).

%\newcommand{\ip}[2]{\langle #1 | #2\rangle}
                             % This is an alternative definition of
                             % \ip that is commented out.

\begin{document}             % End of preamble and beginning of text.

\maketitle                   % Produces the title.

This is an example input file.  Comparing it with
the output it generates can show you how to
produce a simple document of your own.

\section{Ordinary Text}      % Produces section heading.  Lower-level
                             % sections are begun with similar
                             % \subsection and \subsubsection commands.

The ends  of words and sentences are marked
  by   spaces. It  doesn't matter how many
spaces    you type; one is as good as 100.  The
end of   a line counts as a space.

One   or more   blank lines denote the  end
of  a paragraph.

Since any number of consecutive spaces are treated
like a single one, the formatting of the input
file makes no difference to
      \LaTeX,                % The \LaTeX command generates the LaTeX logo.
but it makes a difference to you.  When you use
\LaTeX, making your input file as easy to read
as possible will be a great help as you write
your document and when you change it.  This sample
file shows how you can add comments to your own input
file.

Because printing is different from typewriting,
there are a number of things that you have to do
differently when preparing an input file than if
you were just typing the document directly.
Quotation marks like
       ``this''
have to be handled specially, as do quotes within
quotes:
       ``\,`this'            % \, separates the double and single quote.
        is what I just
        wrote, not  `that'\,''.

Dashes come in three sizes: an
       intra-word
dash, a medium dash for number ranges like
       1--2,
and a punctuation
       dash---like
this.

A sentence-ending space should be larger than the
space between words within a sentence.  You
sometimes have to type special commands in
conjunction with punctuation characters to get
this right, as in the following sentence.
       Gnats, gnus, etc.\ all  % `\ ' makes an inter-word space.
       begin with G\@.         % \@ marks end-of-sentence punctuation.
You should check the spaces after periods when
reading your output to make sure you haven't
forgotten any special cases.  Generating an
ellipsis
       \ldots\               % `\ ' is needed after `\ldots' because TeX
                             % ignores spaces after command names like \ldots
                             % made from \ + letters.
                             %
                             % Note how a `%' character causes TeX to ignore
                             % the end of the input line, so these blank lines
                             % do not start a new paragraph.
                             %
with the right spacing around the periods requires
a special command.

\LaTeX\ interprets some common characters as
commands, so you must type special commands to
generate them.  These characters include the
following:
       \$ \& \% \# \{ and \}.

In printing, text is usually emphasized with an
       \emph{italic}
type style.

\begin{em}
   A long segment of text can also be emphasized
   in this way.  Text within such a segment can be
   given \emph{additional} emphasis.
\end{em}

It is sometimes necessary to prevent \LaTeX\ from
breaking a line where it might otherwise do so.
This may be at a space, as between the ``Mr.''\ and
``Jones'' in
       ``Mr.~Jones'',        % ~ produces an unbreakable interword space.
or within a word---especially when the word is a
symbol like
       \mbox{\emph{itemnum}}
that makes little sense when hyphenated across
lines.

Footnotes\footnote{This is an example of a footnote.}
pose no problem.

\LaTeX\ is good at typesetting mathematical formulas
like
       \( x-3y + z = 7 \)
or
       \( a_{1} > x^{2n} + y^{2n} > x' \)
or
       \( \ip{A}{B} = \sum_{i} a_{i} b_{i} \).
The spaces you type in a formula are
ignored.  Remember that a letter like
       $x$                   % $ ... $  and  \( ... \)  are equivalent
is a formula when it denotes a mathematical
symbol, and it should be typed as one.

\section{Displayed Text}

Text is displayed by indenting it from the left
margin.  Quotations are commonly displayed.  There
are short quotations
\begin{quote}
   This is a short quotation.  It consists of a
   single paragraph of text.  See how it is formatted.
\end{quote}
and longer ones.
\begin{quotation}
   This is a longer quotation.  It consists of two
   paragraphs of text, neither of which are
   particularly interesting.

   This is the second paragraph of the quotation.  It
   is just as dull as the first paragraph.
\end{quotation}
Another frequently-displayed structure is a list.
The following is an example of an \emph{itemized}
list.
\begin{itemize}
   \item This is the first item of an itemized list.
         Each item in the list is marked with a ``tick''.
         You don't have to worry about what kind of tick
         mark is used.

   \item This is the second item of the list.  It
         contains another list nested inside it.  The inner
         list is an \emph{enumerated} list.
         \begin{enumerate}
            \item This is the first item of an enumerated
                  list that is nested within the itemized list.

            \item This is the second item of the inner list.
                  \LaTeX\ allows you to nest lists deeper than
                  you really should.
         \end{enumerate}
         This is the rest of the second item of the outer
         list.  It is no more interesting than any other
         part of the item.
   \item This is the third item of the list.
\end{itemize}
You can even display poetry.
\begin{verse}
   There is an environment
    for verse \\             % The \\ command separates lines
   Whose features some poets % within a stanza.
   will curse.

                             % One or more blank lines separate stanzas.

   For instead of making\\
   Them do \emph{all} line breaking, \\
   It allows them to put too many words on a line when they'd rather be
   forced to be terse.
\end{verse}

Mathematical formulas may also be displayed.  A
displayed formula
is
one-line long; multiline
formulas require special formatting instructions.
   \[  \ip{\Gamma}{\psi'} = x'' + y^{2} + z_{i}^{n}\]
Don't start a paragraph with a displayed equation,
nor make one a paragraph by itself.

\end{document}               % End of document.

\section{Bench Huge Section 1305}
% synthetic bench section 1305
% This is a sample LaTeX input file.  (Version of 12 August 2004.)
%
% A '%' character causes TeX to ignore all remaining text on the line,
% and is used for comments like this one.

\documentclass{article}      % Specifies the document class

                             % The preamble begins here.
\title{An Example Document}  % Declares the document's title.
\author{Leslie Lamport}      % Declares the author's name.
\date{January 21, 1994}      % Deleting this command produces today's date.

\newcommand{\ip}[2]{(#1, #2)}
                             % Defines \ip{arg1}{arg2} to mean
                             % (arg1, arg2).

%\newcommand{\ip}[2]{\langle #1 | #2\rangle}
                             % This is an alternative definition of
                             % \ip that is commented out.

\begin{document}             % End of preamble and beginning of text.

\maketitle                   % Produces the title.

This is an example input file.  Comparing it with
the output it generates can show you how to
produce a simple document of your own.

\section{Ordinary Text}      % Produces section heading.  Lower-level
                             % sections are begun with similar
                             % \subsection and \subsubsection commands.

The ends  of words and sentences are marked
  by   spaces. It  doesn't matter how many
spaces    you type; one is as good as 100.  The
end of   a line counts as a space.

One   or more   blank lines denote the  end
of  a paragraph.

Since any number of consecutive spaces are treated
like a single one, the formatting of the input
file makes no difference to
      \LaTeX,                % The \LaTeX command generates the LaTeX logo.
but it makes a difference to you.  When you use
\LaTeX, making your input file as easy to read
as possible will be a great help as you write
your document and when you change it.  This sample
file shows how you can add comments to your own input
file.

Because printing is different from typewriting,
there are a number of things that you have to do
differently when preparing an input file than if
you were just typing the document directly.
Quotation marks like
       ``this''
have to be handled specially, as do quotes within
quotes:
       ``\,`this'            % \, separates the double and single quote.
        is what I just
        wrote, not  `that'\,''.

Dashes come in three sizes: an
       intra-word
dash, a medium dash for number ranges like
       1--2,
and a punctuation
       dash---like
this.

A sentence-ending space should be larger than the
space between words within a sentence.  You
sometimes have to type special commands in
conjunction with punctuation characters to get
this right, as in the following sentence.
       Gnats, gnus, etc.\ all  % `\ ' makes an inter-word space.
       begin with G\@.         % \@ marks end-of-sentence punctuation.
You should check the spaces after periods when
reading your output to make sure you haven't
forgotten any special cases.  Generating an
ellipsis
       \ldots\               % `\ ' is needed after `\ldots' because TeX
                             % ignores spaces after command names like \ldots
                             % made from \ + letters.
                             %
                             % Note how a `%' character causes TeX to ignore
                             % the end of the input line, so these blank lines
                             % do not start a new paragraph.
                             %
with the right spacing around the periods requires
a special command.

\LaTeX\ interprets some common characters as
commands, so you must type special commands to
generate them.  These characters include the
following:
       \$ \& \% \# \{ and \}.

In printing, text is usually emphasized with an
       \emph{italic}
type style.

\begin{em}
   A long segment of text can also be emphasized
   in this way.  Text within such a segment can be
   given \emph{additional} emphasis.
\end{em}

It is sometimes necessary to prevent \LaTeX\ from
breaking a line where it might otherwise do so.
This may be at a space, as between the ``Mr.''\ and
``Jones'' in
       ``Mr.~Jones'',        % ~ produces an unbreakable interword space.
or within a word---especially when the word is a
symbol like
       \mbox{\emph{itemnum}}
that makes little sense when hyphenated across
lines.

Footnotes\footnote{This is an example of a footnote.}
pose no problem.

\LaTeX\ is good at typesetting mathematical formulas
like
       \( x-3y + z = 7 \)
or
       \( a_{1} > x^{2n} + y^{2n} > x' \)
or
       \( \ip{A}{B} = \sum_{i} a_{i} b_{i} \).
The spaces you type in a formula are
ignored.  Remember that a letter like
       $x$                   % $ ... $  and  \( ... \)  are equivalent
is a formula when it denotes a mathematical
symbol, and it should be typed as one.

\section{Displayed Text}

Text is displayed by indenting it from the left
margin.  Quotations are commonly displayed.  There
are short quotations
\begin{quote}
   This is a short quotation.  It consists of a
   single paragraph of text.  See how it is formatted.
\end{quote}
and longer ones.
\begin{quotation}
   This is a longer quotation.  It consists of two
   paragraphs of text, neither of which are
   particularly interesting.

   This is the second paragraph of the quotation.  It
   is just as dull as the first paragraph.
\end{quotation}
Another frequently-displayed structure is a list.
The following is an example of an \emph{itemized}
list.
\begin{itemize}
   \item This is the first item of an itemized list.
         Each item in the list is marked with a ``tick''.
         You don't have to worry about what kind of tick
         mark is used.

   \item This is the second item of the list.  It
         contains another list nested inside it.  The inner
         list is an \emph{enumerated} list.
         \begin{enumerate}
            \item This is the first item of an enumerated
                  list that is nested within the itemized list.

            \item This is the second item of the inner list.
                  \LaTeX\ allows you to nest lists deeper than
                  you really should.
         \end{enumerate}
         This is the rest of the second item of the outer
         list.  It is no more interesting than any other
         part of the item.
   \item This is the third item of the list.
\end{itemize}
You can even display poetry.
\begin{verse}
   There is an environment
    for verse \\             % The \\ command separates lines
   Whose features some poets % within a stanza.
   will curse.

                             % One or more blank lines separate stanzas.

   For instead of making\\
   Them do \emph{all} line breaking, \\
   It allows them to put too many words on a line when they'd rather be
   forced to be terse.
\end{verse}

Mathematical formulas may also be displayed.  A
displayed formula
is
one-line long; multiline
formulas require special formatting instructions.
   \[  \ip{\Gamma}{\psi'} = x'' + y^{2} + z_{i}^{n}\]
Don't start a paragraph with a displayed equation,
nor make one a paragraph by itself.

\end{document}               % End of document.

\section{Bench Huge Section 1306}
% synthetic bench section 1306
% This is a sample LaTeX input file.  (Version of 12 August 2004.)
%
% A '%' character causes TeX to ignore all remaining text on the line,
% and is used for comments like this one.

\documentclass{article}      % Specifies the document class

                             % The preamble begins here.
\title{An Example Document}  % Declares the document's title.
\author{Leslie Lamport}      % Declares the author's name.
\date{January 21, 1994}      % Deleting this command produces today's date.

\newcommand{\ip}[2]{(#1, #2)}
                             % Defines \ip{arg1}{arg2} to mean
                             % (arg1, arg2).

%\newcommand{\ip}[2]{\langle #1 | #2\rangle}
                             % This is an alternative definition of
                             % \ip that is commented out.

\begin{document}             % End of preamble and beginning of text.

\maketitle                   % Produces the title.

This is an example input file.  Comparing it with
the output it generates can show you how to
produce a simple document of your own.

\section{Ordinary Text}      % Produces section heading.  Lower-level
                             % sections are begun with similar
                             % \subsection and \subsubsection commands.

The ends  of words and sentences are marked
  by   spaces. It  doesn't matter how many
spaces    you type; one is as good as 100.  The
end of   a line counts as a space.

One   or more   blank lines denote the  end
of  a paragraph.

Since any number of consecutive spaces are treated
like a single one, the formatting of the input
file makes no difference to
      \LaTeX,                % The \LaTeX command generates the LaTeX logo.
but it makes a difference to you.  When you use
\LaTeX, making your input file as easy to read
as possible will be a great help as you write
your document and when you change it.  This sample
file shows how you can add comments to your own input
file.

Because printing is different from typewriting,
there are a number of things that you have to do
differently when preparing an input file than if
you were just typing the document directly.
Quotation marks like
       ``this''
have to be handled specially, as do quotes within
quotes:
       ``\,`this'            % \, separates the double and single quote.
        is what I just
        wrote, not  `that'\,''.

Dashes come in three sizes: an
       intra-word
dash, a medium dash for number ranges like
       1--2,
and a punctuation
       dash---like
this.

A sentence-ending space should be larger than the
space between words within a sentence.  You
sometimes have to type special commands in
conjunction with punctuation characters to get
this right, as in the following sentence.
       Gnats, gnus, etc.\ all  % `\ ' makes an inter-word space.
       begin with G\@.         % \@ marks end-of-sentence punctuation.
You should check the spaces after periods when
reading your output to make sure you haven't
forgotten any special cases.  Generating an
ellipsis
       \ldots\               % `\ ' is needed after `\ldots' because TeX
                             % ignores spaces after command names like \ldots
                             % made from \ + letters.
                             %
                             % Note how a `%' character causes TeX to ignore
                             % the end of the input line, so these blank lines
                             % do not start a new paragraph.
                             %
with the right spacing around the periods requires
a special command.

\LaTeX\ interprets some common characters as
commands, so you must type special commands to
generate them.  These characters include the
following:
       \$ \& \% \# \{ and \}.

In printing, text is usually emphasized with an
       \emph{italic}
type style.

\begin{em}
   A long segment of text can also be emphasized
   in this way.  Text within such a segment can be
   given \emph{additional} emphasis.
\end{em}

It is sometimes necessary to prevent \LaTeX\ from
breaking a line where it might otherwise do so.
This may be at a space, as between the ``Mr.''\ and
``Jones'' in
       ``Mr.~Jones'',        % ~ produces an unbreakable interword space.
or within a word---especially when the word is a
symbol like
       \mbox{\emph{itemnum}}
that makes little sense when hyphenated across
lines.

Footnotes\footnote{This is an example of a footnote.}
pose no problem.

\LaTeX\ is good at typesetting mathematical formulas
like
       \( x-3y + z = 7 \)
or
       \( a_{1} > x^{2n} + y^{2n} > x' \)
or
       \( \ip{A}{B} = \sum_{i} a_{i} b_{i} \).
The spaces you type in a formula are
ignored.  Remember that a letter like
       $x$                   % $ ... $  and  \( ... \)  are equivalent
is a formula when it denotes a mathematical
symbol, and it should be typed as one.

\section{Displayed Text}

Text is displayed by indenting it from the left
margin.  Quotations are commonly displayed.  There
are short quotations
\begin{quote}
   This is a short quotation.  It consists of a
   single paragraph of text.  See how it is formatted.
\end{quote}
and longer ones.
\begin{quotation}
   This is a longer quotation.  It consists of two
   paragraphs of text, neither of which are
   particularly interesting.

   This is the second paragraph of the quotation.  It
   is just as dull as the first paragraph.
\end{quotation}
Another frequently-displayed structure is a list.
The following is an example of an \emph{itemized}
list.
\begin{itemize}
   \item This is the first item of an itemized list.
         Each item in the list is marked with a ``tick''.
         You don't have to worry about what kind of tick
         mark is used.

   \item This is the second item of the list.  It
         contains another list nested inside it.  The inner
         list is an \emph{enumerated} list.
         \begin{enumerate}
            \item This is the first item of an enumerated
                  list that is nested within the itemized list.

            \item This is the second item of the inner list.
                  \LaTeX\ allows you to nest lists deeper than
                  you really should.
         \end{enumerate}
         This is the rest of the second item of the outer
         list.  It is no more interesting than any other
         part of the item.
   \item This is the third item of the list.
\end{itemize}
You can even display poetry.
\begin{verse}
   There is an environment
    for verse \\             % The \\ command separates lines
   Whose features some poets % within a stanza.
   will curse.

                             % One or more blank lines separate stanzas.

   For instead of making\\
   Them do \emph{all} line breaking, \\
   It allows them to put too many words on a line when they'd rather be
   forced to be terse.
\end{verse}

Mathematical formulas may also be displayed.  A
displayed formula
is
one-line long; multiline
formulas require special formatting instructions.
   \[  \ip{\Gamma}{\psi'} = x'' + y^{2} + z_{i}^{n}\]
Don't start a paragraph with a displayed equation,
nor make one a paragraph by itself.

\end{document}               % End of document.

\section{Bench Huge Section 1307}
% synthetic bench section 1307
% This is a sample LaTeX input file.  (Version of 12 August 2004.)
%
% A '%' character causes TeX to ignore all remaining text on the line,
% and is used for comments like this one.

\documentclass{article}      % Specifies the document class

                             % The preamble begins here.
\title{An Example Document}  % Declares the document's title.
\author{Leslie Lamport}      % Declares the author's name.
\date{January 21, 1994}      % Deleting this command produces today's date.

\newcommand{\ip}[2]{(#1, #2)}
                             % Defines \ip{arg1}{arg2} to mean
                             % (arg1, arg2).

%\newcommand{\ip}[2]{\langle #1 | #2\rangle}
                             % This is an alternative definition of
                             % \ip that is commented out.

\begin{document}             % End of preamble and beginning of text.

\maketitle                   % Produces the title.

This is an example input file.  Comparing it with
the output it generates can show you how to
produce a simple document of your own.

\section{Ordinary Text}      % Produces section heading.  Lower-level
                             % sections are begun with similar
                             % \subsection and \subsubsection commands.

The ends  of words and sentences are marked
  by   spaces. It  doesn't matter how many
spaces    you type; one is as good as 100.  The
end of   a line counts as a space.

One   or more   blank lines denote the  end
of  a paragraph.

Since any number of consecutive spaces are treated
like a single one, the formatting of the input
file makes no difference to
      \LaTeX,                % The \LaTeX command generates the LaTeX logo.
but it makes a difference to you.  When you use
\LaTeX, making your input file as easy to read
as possible will be a great help as you write
your document and when you change it.  This sample
file shows how you can add comments to your own input
file.

Because printing is different from typewriting,
there are a number of things that you have to do
differently when preparing an input file than if
you were just typing the document directly.
Quotation marks like
       ``this''
have to be handled specially, as do quotes within
quotes:
       ``\,`this'            % \, separates the double and single quote.
        is what I just
        wrote, not  `that'\,''.

Dashes come in three sizes: an
       intra-word
dash, a medium dash for number ranges like
       1--2,
and a punctuation
       dash---like
this.

A sentence-ending space should be larger than the
space between words within a sentence.  You
sometimes have to type special commands in
conjunction with punctuation characters to get
this right, as in the following sentence.
       Gnats, gnus, etc.\ all  % `\ ' makes an inter-word space.
       begin with G\@.         % \@ marks end-of-sentence punctuation.
You should check the spaces after periods when
reading your output to make sure you haven't
forgotten any special cases.  Generating an
ellipsis
       \ldots\               % `\ ' is needed after `\ldots' because TeX
                             % ignores spaces after command names like \ldots
                             % made from \ + letters.
                             %
                             % Note how a `%' character causes TeX to ignore
                             % the end of the input line, so these blank lines
                             % do not start a new paragraph.
                             %
with the right spacing around the periods requires
a special command.

\LaTeX\ interprets some common characters as
commands, so you must type special commands to
generate them.  These characters include the
following:
       \$ \& \% \# \{ and \}.

In printing, text is usually emphasized with an
       \emph{italic}
type style.

\begin{em}
   A long segment of text can also be emphasized
   in this way.  Text within such a segment can be
   given \emph{additional} emphasis.
\end{em}

It is sometimes necessary to prevent \LaTeX\ from
breaking a line where it might otherwise do so.
This may be at a space, as between the ``Mr.''\ and
``Jones'' in
       ``Mr.~Jones'',        % ~ produces an unbreakable interword space.
or within a word---especially when the word is a
symbol like
       \mbox{\emph{itemnum}}
that makes little sense when hyphenated across
lines.

Footnotes\footnote{This is an example of a footnote.}
pose no problem.

\LaTeX\ is good at typesetting mathematical formulas
like
       \( x-3y + z = 7 \)
or
       \( a_{1} > x^{2n} + y^{2n} > x' \)
or
       \( \ip{A}{B} = \sum_{i} a_{i} b_{i} \).
The spaces you type in a formula are
ignored.  Remember that a letter like
       $x$                   % $ ... $  and  \( ... \)  are equivalent
is a formula when it denotes a mathematical
symbol, and it should be typed as one.

\section{Displayed Text}

Text is displayed by indenting it from the left
margin.  Quotations are commonly displayed.  There
are short quotations
\begin{quote}
   This is a short quotation.  It consists of a
   single paragraph of text.  See how it is formatted.
\end{quote}
and longer ones.
\begin{quotation}
   This is a longer quotation.  It consists of two
   paragraphs of text, neither of which are
   particularly interesting.

   This is the second paragraph of the quotation.  It
   is just as dull as the first paragraph.
\end{quotation}
Another frequently-displayed structure is a list.
The following is an example of an \emph{itemized}
list.
\begin{itemize}
   \item This is the first item of an itemized list.
         Each item in the list is marked with a ``tick''.
         You don't have to worry about what kind of tick
         mark is used.

   \item This is the second item of the list.  It
         contains another list nested inside it.  The inner
         list is an \emph{enumerated} list.
         \begin{enumerate}
            \item This is the first item of an enumerated
                  list that is nested within the itemized list.

            \item This is the second item of the inner list.
                  \LaTeX\ allows you to nest lists deeper than
                  you really should.
         \end{enumerate}
         This is the rest of the second item of the outer
         list.  It is no more interesting than any other
         part of the item.
   \item This is the third item of the list.
\end{itemize}
You can even display poetry.
\begin{verse}
   There is an environment
    for verse \\             % The \\ command separates lines
   Whose features some poets % within a stanza.
   will curse.

                             % One or more blank lines separate stanzas.

   For instead of making\\
   Them do \emph{all} line breaking, \\
   It allows them to put too many words on a line when they'd rather be
   forced to be terse.
\end{verse}

Mathematical formulas may also be displayed.  A
displayed formula
is
one-line long; multiline
formulas require special formatting instructions.
   \[  \ip{\Gamma}{\psi'} = x'' + y^{2} + z_{i}^{n}\]
Don't start a paragraph with a displayed equation,
nor make one a paragraph by itself.

\end{document}               % End of document.

\section{Bench Huge Section 1308}
% synthetic bench section 1308
% This is a sample LaTeX input file.  (Version of 12 August 2004.)
%
% A '%' character causes TeX to ignore all remaining text on the line,
% and is used for comments like this one.

\documentclass{article}      % Specifies the document class

                             % The preamble begins here.
\title{An Example Document}  % Declares the document's title.
\author{Leslie Lamport}      % Declares the author's name.
\date{January 21, 1994}      % Deleting this command produces today's date.

\newcommand{\ip}[2]{(#1, #2)}
                             % Defines \ip{arg1}{arg2} to mean
                             % (arg1, arg2).

%\newcommand{\ip}[2]{\langle #1 | #2\rangle}
                             % This is an alternative definition of
                             % \ip that is commented out.

\begin{document}             % End of preamble and beginning of text.

\maketitle                   % Produces the title.

This is an example input file.  Comparing it with
the output it generates can show you how to
produce a simple document of your own.

\section{Ordinary Text}      % Produces section heading.  Lower-level
                             % sections are begun with similar
                             % \subsection and \subsubsection commands.

The ends  of words and sentences are marked
  by   spaces. It  doesn't matter how many
spaces    you type; one is as good as 100.  The
end of   a line counts as a space.

One   or more   blank lines denote the  end
of  a paragraph.

Since any number of consecutive spaces are treated
like a single one, the formatting of the input
file makes no difference to
      \LaTeX,                % The \LaTeX command generates the LaTeX logo.
but it makes a difference to you.  When you use
\LaTeX, making your input file as easy to read
as possible will be a great help as you write
your document and when you change it.  This sample
file shows how you can add comments to your own input
file.

Because printing is different from typewriting,
there are a number of things that you have to do
differently when preparing an input file than if
you were just typing the document directly.
Quotation marks like
       ``this''
have to be handled specially, as do quotes within
quotes:
       ``\,`this'            % \, separates the double and single quote.
        is what I just
        wrote, not  `that'\,''.

Dashes come in three sizes: an
       intra-word
dash, a medium dash for number ranges like
       1--2,
and a punctuation
       dash---like
this.

A sentence-ending space should be larger than the
space between words within a sentence.  You
sometimes have to type special commands in
conjunction with punctuation characters to get
this right, as in the following sentence.
       Gnats, gnus, etc.\ all  % `\ ' makes an inter-word space.
       begin with G\@.         % \@ marks end-of-sentence punctuation.
You should check the spaces after periods when
reading your output to make sure you haven't
forgotten any special cases.  Generating an
ellipsis
       \ldots\               % `\ ' is needed after `\ldots' because TeX
                             % ignores spaces after command names like \ldots
                             % made from \ + letters.
                             %
                             % Note how a `%' character causes TeX to ignore
                             % the end of the input line, so these blank lines
                             % do not start a new paragraph.
                             %
with the right spacing around the periods requires
a special command.

\LaTeX\ interprets some common characters as
commands, so you must type special commands to
generate them.  These characters include the
following:
       \$ \& \% \# \{ and \}.

In printing, text is usually emphasized with an
       \emph{italic}
type style.

\begin{em}
   A long segment of text can also be emphasized
   in this way.  Text within such a segment can be
   given \emph{additional} emphasis.
\end{em}

It is sometimes necessary to prevent \LaTeX\ from
breaking a line where it might otherwise do so.
This may be at a space, as between the ``Mr.''\ and
``Jones'' in
       ``Mr.~Jones'',        % ~ produces an unbreakable interword space.
or within a word---especially when the word is a
symbol like
       \mbox{\emph{itemnum}}
that makes little sense when hyphenated across
lines.

Footnotes\footnote{This is an example of a footnote.}
pose no problem.

\LaTeX\ is good at typesetting mathematical formulas
like
       \( x-3y + z = 7 \)
or
       \( a_{1} > x^{2n} + y^{2n} > x' \)
or
       \( \ip{A}{B} = \sum_{i} a_{i} b_{i} \).
The spaces you type in a formula are
ignored.  Remember that a letter like
       $x$                   % $ ... $  and  \( ... \)  are equivalent
is a formula when it denotes a mathematical
symbol, and it should be typed as one.

\section{Displayed Text}

Text is displayed by indenting it from the left
margin.  Quotations are commonly displayed.  There
are short quotations
\begin{quote}
   This is a short quotation.  It consists of a
   single paragraph of text.  See how it is formatted.
\end{quote}
and longer ones.
\begin{quotation}
   This is a longer quotation.  It consists of two
   paragraphs of text, neither of which are
   particularly interesting.

   This is the second paragraph of the quotation.  It
   is just as dull as the first paragraph.
\end{quotation}
Another frequently-displayed structure is a list.
The following is an example of an \emph{itemized}
list.
\begin{itemize}
   \item This is the first item of an itemized list.
         Each item in the list is marked with a ``tick''.
         You don't have to worry about what kind of tick
         mark is used.

   \item This is the second item of the list.  It
         contains another list nested inside it.  The inner
         list is an \emph{enumerated} list.
         \begin{enumerate}
            \item This is the first item of an enumerated
                  list that is nested within the itemized list.

            \item This is the second item of the inner list.
                  \LaTeX\ allows you to nest lists deeper than
                  you really should.
         \end{enumerate}
         This is the rest of the second item of the outer
         list.  It is no more interesting than any other
         part of the item.
   \item This is the third item of the list.
\end{itemize}
You can even display poetry.
\begin{verse}
   There is an environment
    for verse \\             % The \\ command separates lines
   Whose features some poets % within a stanza.
   will curse.

                             % One or more blank lines separate stanzas.

   For instead of making\\
   Them do \emph{all} line breaking, \\
   It allows them to put too many words on a line when they'd rather be
   forced to be terse.
\end{verse}

Mathematical formulas may also be displayed.  A
displayed formula
is
one-line long; multiline
formulas require special formatting instructions.
   \[  \ip{\Gamma}{\psi'} = x'' + y^{2} + z_{i}^{n}\]
Don't start a paragraph with a displayed equation,
nor make one a paragraph by itself.

\end{document}               % End of document.

\section{Bench Huge Section 1309}
% synthetic bench section 1309
% This is a sample LaTeX input file.  (Version of 12 August 2004.)
%
% A '%' character causes TeX to ignore all remaining text on the line,
% and is used for comments like this one.

\documentclass{article}      % Specifies the document class

                             % The preamble begins here.
\title{An Example Document}  % Declares the document's title.
\author{Leslie Lamport}      % Declares the author's name.
\date{January 21, 1994}      % Deleting this command produces today's date.

\newcommand{\ip}[2]{(#1, #2)}
                             % Defines \ip{arg1}{arg2} to mean
                             % (arg1, arg2).

%\newcommand{\ip}[2]{\langle #1 | #2\rangle}
                             % This is an alternative definition of
                             % \ip that is commented out.

\begin{document}             % End of preamble and beginning of text.

\maketitle                   % Produces the title.

This is an example input file.  Comparing it with
the output it generates can show you how to
produce a simple document of your own.

\section{Ordinary Text}      % Produces section heading.  Lower-level
                             % sections are begun with similar
                             % \subsection and \subsubsection commands.

The ends  of words and sentences are marked
  by   spaces. It  doesn't matter how many
spaces    you type; one is as good as 100.  The
end of   a line counts as a space.

One   or more   blank lines denote the  end
of  a paragraph.

Since any number of consecutive spaces are treated
like a single one, the formatting of the input
file makes no difference to
      \LaTeX,                % The \LaTeX command generates the LaTeX logo.
but it makes a difference to you.  When you use
\LaTeX, making your input file as easy to read
as possible will be a great help as you write
your document and when you change it.  This sample
file shows how you can add comments to your own input
file.

Because printing is different from typewriting,
there are a number of things that you have to do
differently when preparing an input file than if
you were just typing the document directly.
Quotation marks like
       ``this''
have to be handled specially, as do quotes within
quotes:
       ``\,`this'            % \, separates the double and single quote.
        is what I just
        wrote, not  `that'\,''.

Dashes come in three sizes: an
       intra-word
dash, a medium dash for number ranges like
       1--2,
and a punctuation
       dash---like
this.

A sentence-ending space should be larger than the
space between words within a sentence.  You
sometimes have to type special commands in
conjunction with punctuation characters to get
this right, as in the following sentence.
       Gnats, gnus, etc.\ all  % `\ ' makes an inter-word space.
       begin with G\@.         % \@ marks end-of-sentence punctuation.
You should check the spaces after periods when
reading your output to make sure you haven't
forgotten any special cases.  Generating an
ellipsis
       \ldots\               % `\ ' is needed after `\ldots' because TeX
                             % ignores spaces after command names like \ldots
                             % made from \ + letters.
                             %
                             % Note how a `%' character causes TeX to ignore
                             % the end of the input line, so these blank lines
                             % do not start a new paragraph.
                             %
with the right spacing around the periods requires
a special command.

\LaTeX\ interprets some common characters as
commands, so you must type special commands to
generate them.  These characters include the
following:
       \$ \& \% \# \{ and \}.

In printing, text is usually emphasized with an
       \emph{italic}
type style.

\begin{em}
   A long segment of text can also be emphasized
   in this way.  Text within such a segment can be
   given \emph{additional} emphasis.
\end{em}

It is sometimes necessary to prevent \LaTeX\ from
breaking a line where it might otherwise do so.
This may be at a space, as between the ``Mr.''\ and
``Jones'' in
       ``Mr.~Jones'',        % ~ produces an unbreakable interword space.
or within a word---especially when the word is a
symbol like
       \mbox{\emph{itemnum}}
that makes little sense when hyphenated across
lines.

Footnotes\footnote{This is an example of a footnote.}
pose no problem.

\LaTeX\ is good at typesetting mathematical formulas
like
       \( x-3y + z = 7 \)
or
       \( a_{1} > x^{2n} + y^{2n} > x' \)
or
       \( \ip{A}{B} = \sum_{i} a_{i} b_{i} \).
The spaces you type in a formula are
ignored.  Remember that a letter like
       $x$                   % $ ... $  and  \( ... \)  are equivalent
is a formula when it denotes a mathematical
symbol, and it should be typed as one.

\section{Displayed Text}

Text is displayed by indenting it from the left
margin.  Quotations are commonly displayed.  There
are short quotations
\begin{quote}
   This is a short quotation.  It consists of a
   single paragraph of text.  See how it is formatted.
\end{quote}
and longer ones.
\begin{quotation}
   This is a longer quotation.  It consists of two
   paragraphs of text, neither of which are
   particularly interesting.

   This is the second paragraph of the quotation.  It
   is just as dull as the first paragraph.
\end{quotation}
Another frequently-displayed structure is a list.
The following is an example of an \emph{itemized}
list.
\begin{itemize}
   \item This is the first item of an itemized list.
         Each item in the list is marked with a ``tick''.
         You don't have to worry about what kind of tick
         mark is used.

   \item This is the second item of the list.  It
         contains another list nested inside it.  The inner
         list is an \emph{enumerated} list.
         \begin{enumerate}
            \item This is the first item of an enumerated
                  list that is nested within the itemized list.

            \item This is the second item of the inner list.
                  \LaTeX\ allows you to nest lists deeper than
                  you really should.
         \end{enumerate}
         This is the rest of the second item of the outer
         list.  It is no more interesting than any other
         part of the item.
   \item This is the third item of the list.
\end{itemize}
You can even display poetry.
\begin{verse}
   There is an environment
    for verse \\             % The \\ command separates lines
   Whose features some poets % within a stanza.
   will curse.

                             % One or more blank lines separate stanzas.

   For instead of making\\
   Them do \emph{all} line breaking, \\
   It allows them to put too many words on a line when they'd rather be
   forced to be terse.
\end{verse}

Mathematical formulas may also be displayed.  A
displayed formula
is
one-line long; multiline
formulas require special formatting instructions.
   \[  \ip{\Gamma}{\psi'} = x'' + y^{2} + z_{i}^{n}\]
Don't start a paragraph with a displayed equation,
nor make one a paragraph by itself.

\end{document}               % End of document.

\section{Bench Huge Section 1310}
% synthetic bench section 1310
% This is a sample LaTeX input file.  (Version of 12 August 2004.)
%
% A '%' character causes TeX to ignore all remaining text on the line,
% and is used for comments like this one.

\documentclass{article}      % Specifies the document class

                             % The preamble begins here.
\title{An Example Document}  % Declares the document's title.
\author{Leslie Lamport}      % Declares the author's name.
\date{January 21, 1994}      % Deleting this command produces today's date.

\newcommand{\ip}[2]{(#1, #2)}
                             % Defines \ip{arg1}{arg2} to mean
                             % (arg1, arg2).

%\newcommand{\ip}[2]{\langle #1 | #2\rangle}
                             % This is an alternative definition of
                             % \ip that is commented out.

\begin{document}             % End of preamble and beginning of text.

\maketitle                   % Produces the title.

This is an example input file.  Comparing it with
the output it generates can show you how to
produce a simple document of your own.

\section{Ordinary Text}      % Produces section heading.  Lower-level
                             % sections are begun with similar
                             % \subsection and \subsubsection commands.

The ends  of words and sentences are marked
  by   spaces. It  doesn't matter how many
spaces    you type; one is as good as 100.  The
end of   a line counts as a space.

One   or more   blank lines denote the  end
of  a paragraph.

Since any number of consecutive spaces are treated
like a single one, the formatting of the input
file makes no difference to
      \LaTeX,                % The \LaTeX command generates the LaTeX logo.
but it makes a difference to you.  When you use
\LaTeX, making your input file as easy to read
as possible will be a great help as you write
your document and when you change it.  This sample
file shows how you can add comments to your own input
file.

Because printing is different from typewriting,
there are a number of things that you have to do
differently when preparing an input file than if
you were just typing the document directly.
Quotation marks like
       ``this''
have to be handled specially, as do quotes within
quotes:
       ``\,`this'            % \, separates the double and single quote.
        is what I just
        wrote, not  `that'\,''.

Dashes come in three sizes: an
       intra-word
dash, a medium dash for number ranges like
       1--2,
and a punctuation
       dash---like
this.

A sentence-ending space should be larger than the
space between words within a sentence.  You
sometimes have to type special commands in
conjunction with punctuation characters to get
this right, as in the following sentence.
       Gnats, gnus, etc.\ all  % `\ ' makes an inter-word space.
       begin with G\@.         % \@ marks end-of-sentence punctuation.
You should check the spaces after periods when
reading your output to make sure you haven't
forgotten any special cases.  Generating an
ellipsis
       \ldots\               % `\ ' is needed after `\ldots' because TeX
                             % ignores spaces after command names like \ldots
                             % made from \ + letters.
                             %
                             % Note how a `%' character causes TeX to ignore
                             % the end of the input line, so these blank lines
                             % do not start a new paragraph.
                             %
with the right spacing around the periods requires
a special command.

\LaTeX\ interprets some common characters as
commands, so you must type special commands to
generate them.  These characters include the
following:
       \$ \& \% \# \{ and \}.

In printing, text is usually emphasized with an
       \emph{italic}
type style.

\begin{em}
   A long segment of text can also be emphasized
   in this way.  Text within such a segment can be
   given \emph{additional} emphasis.
\end{em}

It is sometimes necessary to prevent \LaTeX\ from
breaking a line where it might otherwise do so.
This may be at a space, as between the ``Mr.''\ and
``Jones'' in
       ``Mr.~Jones'',        % ~ produces an unbreakable interword space.
or within a word---especially when the word is a
symbol like
       \mbox{\emph{itemnum}}
that makes little sense when hyphenated across
lines.

Footnotes\footnote{This is an example of a footnote.}
pose no problem.

\LaTeX\ is good at typesetting mathematical formulas
like
       \( x-3y + z = 7 \)
or
       \( a_{1} > x^{2n} + y^{2n} > x' \)
or
       \( \ip{A}{B} = \sum_{i} a_{i} b_{i} \).
The spaces you type in a formula are
ignored.  Remember that a letter like
       $x$                   % $ ... $  and  \( ... \)  are equivalent
is a formula when it denotes a mathematical
symbol, and it should be typed as one.

\section{Displayed Text}

Text is displayed by indenting it from the left
margin.  Quotations are commonly displayed.  There
are short quotations
\begin{quote}
   This is a short quotation.  It consists of a
   single paragraph of text.  See how it is formatted.
\end{quote}
and longer ones.
\begin{quotation}
   This is a longer quotation.  It consists of two
   paragraphs of text, neither of which are
   particularly interesting.

   This is the second paragraph of the quotation.  It
   is just as dull as the first paragraph.
\end{quotation}
Another frequently-displayed structure is a list.
The following is an example of an \emph{itemized}
list.
\begin{itemize}
   \item This is the first item of an itemized list.
         Each item in the list is marked with a ``tick''.
         You don't have to worry about what kind of tick
         mark is used.

   \item This is the second item of the list.  It
         contains another list nested inside it.  The inner
         list is an \emph{enumerated} list.
         \begin{enumerate}
            \item This is the first item of an enumerated
                  list that is nested within the itemized list.

            \item This is the second item of the inner list.
                  \LaTeX\ allows you to nest lists deeper than
                  you really should.
         \end{enumerate}
         This is the rest of the second item of the outer
         list.  It is no more interesting than any other
         part of the item.
   \item This is the third item of the list.
\end{itemize}
You can even display poetry.
\begin{verse}
   There is an environment
    for verse \\             % The \\ command separates lines
   Whose features some poets % within a stanza.
   will curse.

                             % One or more blank lines separate stanzas.

   For instead of making\\
   Them do \emph{all} line breaking, \\
   It allows them to put too many words on a line when they'd rather be
   forced to be terse.
\end{verse}

Mathematical formulas may also be displayed.  A
displayed formula
is
one-line long; multiline
formulas require special formatting instructions.
   \[  \ip{\Gamma}{\psi'} = x'' + y^{2} + z_{i}^{n}\]
Don't start a paragraph with a displayed equation,
nor make one a paragraph by itself.

\end{document}               % End of document.

\section{Bench Huge Section 1311}
% synthetic bench section 1311
% This is a sample LaTeX input file.  (Version of 12 August 2004.)
%
% A '%' character causes TeX to ignore all remaining text on the line,
% and is used for comments like this one.

\documentclass{article}      % Specifies the document class

                             % The preamble begins here.
\title{An Example Document}  % Declares the document's title.
\author{Leslie Lamport}      % Declares the author's name.
\date{January 21, 1994}      % Deleting this command produces today's date.

\newcommand{\ip}[2]{(#1, #2)}
                             % Defines \ip{arg1}{arg2} to mean
                             % (arg1, arg2).

%\newcommand{\ip}[2]{\langle #1 | #2\rangle}
                             % This is an alternative definition of
                             % \ip that is commented out.

\begin{document}             % End of preamble and beginning of text.

\maketitle                   % Produces the title.

This is an example input file.  Comparing it with
the output it generates can show you how to
produce a simple document of your own.

\section{Ordinary Text}      % Produces section heading.  Lower-level
                             % sections are begun with similar
                             % \subsection and \subsubsection commands.

The ends  of words and sentences are marked
  by   spaces. It  doesn't matter how many
spaces    you type; one is as good as 100.  The
end of   a line counts as a space.

One   or more   blank lines denote the  end
of  a paragraph.

Since any number of consecutive spaces are treated
like a single one, the formatting of the input
file makes no difference to
      \LaTeX,                % The \LaTeX command generates the LaTeX logo.
but it makes a difference to you.  When you use
\LaTeX, making your input file as easy to read
as possible will be a great help as you write
your document and when you change it.  This sample
file shows how you can add comments to your own input
file.

Because printing is different from typewriting,
there are a number of things that you have to do
differently when preparing an input file than if
you were just typing the document directly.
Quotation marks like
       ``this''
have to be handled specially, as do quotes within
quotes:
       ``\,`this'            % \, separates the double and single quote.
        is what I just
        wrote, not  `that'\,''.

Dashes come in three sizes: an
       intra-word
dash, a medium dash for number ranges like
       1--2,
and a punctuation
       dash---like
this.

A sentence-ending space should be larger than the
space between words within a sentence.  You
sometimes have to type special commands in
conjunction with punctuation characters to get
this right, as in the following sentence.
       Gnats, gnus, etc.\ all  % `\ ' makes an inter-word space.
       begin with G\@.         % \@ marks end-of-sentence punctuation.
You should check the spaces after periods when
reading your output to make sure you haven't
forgotten any special cases.  Generating an
ellipsis
       \ldots\               % `\ ' is needed after `\ldots' because TeX
                             % ignores spaces after command names like \ldots
                             % made from \ + letters.
                             %
                             % Note how a `%' character causes TeX to ignore
                             % the end of the input line, so these blank lines
                             % do not start a new paragraph.
                             %
with the right spacing around the periods requires
a special command.

\LaTeX\ interprets some common characters as
commands, so you must type special commands to
generate them.  These characters include the
following:
       \$ \& \% \# \{ and \}.

In printing, text is usually emphasized with an
       \emph{italic}
type style.

\begin{em}
   A long segment of text can also be emphasized
   in this way.  Text within such a segment can be
   given \emph{additional} emphasis.
\end{em}

It is sometimes necessary to prevent \LaTeX\ from
breaking a line where it might otherwise do so.
This may be at a space, as between the ``Mr.''\ and
``Jones'' in
       ``Mr.~Jones'',        % ~ produces an unbreakable interword space.
or within a word---especially when the word is a
symbol like
       \mbox{\emph{itemnum}}
that makes little sense when hyphenated across
lines.

Footnotes\footnote{This is an example of a footnote.}
pose no problem.

\LaTeX\ is good at typesetting mathematical formulas
like
       \( x-3y + z = 7 \)
or
       \( a_{1} > x^{2n} + y^{2n} > x' \)
or
       \( \ip{A}{B} = \sum_{i} a_{i} b_{i} \).
The spaces you type in a formula are
ignored.  Remember that a letter like
       $x$                   % $ ... $  and  \( ... \)  are equivalent
is a formula when it denotes a mathematical
symbol, and it should be typed as one.

\section{Displayed Text}

Text is displayed by indenting it from the left
margin.  Quotations are commonly displayed.  There
are short quotations
\begin{quote}
   This is a short quotation.  It consists of a
   single paragraph of text.  See how it is formatted.
\end{quote}
and longer ones.
\begin{quotation}
   This is a longer quotation.  It consists of two
   paragraphs of text, neither of which are
   particularly interesting.

   This is the second paragraph of the quotation.  It
   is just as dull as the first paragraph.
\end{quotation}
Another frequently-displayed structure is a list.
The following is an example of an \emph{itemized}
list.
\begin{itemize}
   \item This is the first item of an itemized list.
         Each item in the list is marked with a ``tick''.
         You don't have to worry about what kind of tick
         mark is used.

   \item This is the second item of the list.  It
         contains another list nested inside it.  The inner
         list is an \emph{enumerated} list.
         \begin{enumerate}
            \item This is the first item of an enumerated
                  list that is nested within the itemized list.

            \item This is the second item of the inner list.
                  \LaTeX\ allows you to nest lists deeper than
                  you really should.
         \end{enumerate}
         This is the rest of the second item of the outer
         list.  It is no more interesting than any other
         part of the item.
   \item This is the third item of the list.
\end{itemize}
You can even display poetry.
\begin{verse}
   There is an environment
    for verse \\             % The \\ command separates lines
   Whose features some poets % within a stanza.
   will curse.

                             % One or more blank lines separate stanzas.

   For instead of making\\
   Them do \emph{all} line breaking, \\
   It allows them to put too many words on a line when they'd rather be
   forced to be terse.
\end{verse}

Mathematical formulas may also be displayed.  A
displayed formula
is
one-line long; multiline
formulas require special formatting instructions.
   \[  \ip{\Gamma}{\psi'} = x'' + y^{2} + z_{i}^{n}\]
Don't start a paragraph with a displayed equation,
nor make one a paragraph by itself.

\end{document}               % End of document.

\section{Bench Huge Section 1312}
% synthetic bench section 1312
% This is a sample LaTeX input file.  (Version of 12 August 2004.)
%
% A '%' character causes TeX to ignore all remaining text on the line,
% and is used for comments like this one.

\documentclass{article}      % Specifies the document class

                             % The preamble begins here.
\title{An Example Document}  % Declares the document's title.
\author{Leslie Lamport}      % Declares the author's name.
\date{January 21, 1994}      % Deleting this command produces today's date.

\newcommand{\ip}[2]{(#1, #2)}
                             % Defines \ip{arg1}{arg2} to mean
                             % (arg1, arg2).

%\newcommand{\ip}[2]{\langle #1 | #2\rangle}
                             % This is an alternative definition of
                             % \ip that is commented out.

\begin{document}             % End of preamble and beginning of text.

\maketitle                   % Produces the title.

This is an example input file.  Comparing it with
the output it generates can show you how to
produce a simple document of your own.

\section{Ordinary Text}      % Produces section heading.  Lower-level
                             % sections are begun with similar
                             % \subsection and \subsubsection commands.

The ends  of words and sentences are marked
  by   spaces. It  doesn't matter how many
spaces    you type; one is as good as 100.  The
end of   a line counts as a space.

One   or more   blank lines denote the  end
of  a paragraph.

Since any number of consecutive spaces are treated
like a single one, the formatting of the input
file makes no difference to
      \LaTeX,                % The \LaTeX command generates the LaTeX logo.
but it makes a difference to you.  When you use
\LaTeX, making your input file as easy to read
as possible will be a great help as you write
your document and when you change it.  This sample
file shows how you can add comments to your own input
file.

Because printing is different from typewriting,
there are a number of things that you have to do
differently when preparing an input file than if
you were just typing the document directly.
Quotation marks like
       ``this''
have to be handled specially, as do quotes within
quotes:
       ``\,`this'            % \, separates the double and single quote.
        is what I just
        wrote, not  `that'\,''.

Dashes come in three sizes: an
       intra-word
dash, a medium dash for number ranges like
       1--2,
and a punctuation
       dash---like
this.

A sentence-ending space should be larger than the
space between words within a sentence.  You
sometimes have to type special commands in
conjunction with punctuation characters to get
this right, as in the following sentence.
       Gnats, gnus, etc.\ all  % `\ ' makes an inter-word space.
       begin with G\@.         % \@ marks end-of-sentence punctuation.
You should check the spaces after periods when
reading your output to make sure you haven't
forgotten any special cases.  Generating an
ellipsis
       \ldots\               % `\ ' is needed after `\ldots' because TeX
                             % ignores spaces after command names like \ldots
                             % made from \ + letters.
                             %
                             % Note how a `%' character causes TeX to ignore
                             % the end of the input line, so these blank lines
                             % do not start a new paragraph.
                             %
with the right spacing around the periods requires
a special command.

\LaTeX\ interprets some common characters as
commands, so you must type special commands to
generate them.  These characters include the
following:
       \$ \& \% \# \{ and \}.

In printing, text is usually emphasized with an
       \emph{italic}
type style.

\begin{em}
   A long segment of text can also be emphasized
   in this way.  Text within such a segment can be
   given \emph{additional} emphasis.
\end{em}

It is sometimes necessary to prevent \LaTeX\ from
breaking a line where it might otherwise do so.
This may be at a space, as between the ``Mr.''\ and
``Jones'' in
       ``Mr.~Jones'',        % ~ produces an unbreakable interword space.
or within a word---especially when the word is a
symbol like
       \mbox{\emph{itemnum}}
that makes little sense when hyphenated across
lines.

Footnotes\footnote{This is an example of a footnote.}
pose no problem.

\LaTeX\ is good at typesetting mathematical formulas
like
       \( x-3y + z = 7 \)
or
       \( a_{1} > x^{2n} + y^{2n} > x' \)
or
       \( \ip{A}{B} = \sum_{i} a_{i} b_{i} \).
The spaces you type in a formula are
ignored.  Remember that a letter like
       $x$                   % $ ... $  and  \( ... \)  are equivalent
is a formula when it denotes a mathematical
symbol, and it should be typed as one.

\section{Displayed Text}

Text is displayed by indenting it from the left
margin.  Quotations are commonly displayed.  There
are short quotations
\begin{quote}
   This is a short quotation.  It consists of a
   single paragraph of text.  See how it is formatted.
\end{quote}
and longer ones.
\begin{quotation}
   This is a longer quotation.  It consists of two
   paragraphs of text, neither of which are
   particularly interesting.

   This is the second paragraph of the quotation.  It
   is just as dull as the first paragraph.
\end{quotation}
Another frequently-displayed structure is a list.
The following is an example of an \emph{itemized}
list.
\begin{itemize}
   \item This is the first item of an itemized list.
         Each item in the list is marked with a ``tick''.
         You don't have to worry about what kind of tick
         mark is used.

   \item This is the second item of the list.  It
         contains another list nested inside it.  The inner
         list is an \emph{enumerated} list.
         \begin{enumerate}
            \item This is the first item of an enumerated
                  list that is nested within the itemized list.

            \item This is the second item of the inner list.
                  \LaTeX\ allows you to nest lists deeper than
                  you really should.
         \end{enumerate}
         This is the rest of the second item of the outer
         list.  It is no more interesting than any other
         part of the item.
   \item This is the third item of the list.
\end{itemize}
You can even display poetry.
\begin{verse}
   There is an environment
    for verse \\             % The \\ command separates lines
   Whose features some poets % within a stanza.
   will curse.

                             % One or more blank lines separate stanzas.

   For instead of making\\
   Them do \emph{all} line breaking, \\
   It allows them to put too many words on a line when they'd rather be
   forced to be terse.
\end{verse}

Mathematical formulas may also be displayed.  A
displayed formula
is
one-line long; multiline
formulas require special formatting instructions.
   \[  \ip{\Gamma}{\psi'} = x'' + y^{2} + z_{i}^{n}\]
Don't start a paragraph with a displayed equation,
nor make one a paragraph by itself.

\end{document}               % End of document.

\section{Bench Huge Section 1313}
% synthetic bench section 1313
% This is a sample LaTeX input file.  (Version of 12 August 2004.)
%
% A '%' character causes TeX to ignore all remaining text on the line,
% and is used for comments like this one.

\documentclass{article}      % Specifies the document class

                             % The preamble begins here.
\title{An Example Document}  % Declares the document's title.
\author{Leslie Lamport}      % Declares the author's name.
\date{January 21, 1994}      % Deleting this command produces today's date.

\newcommand{\ip}[2]{(#1, #2)}
                             % Defines \ip{arg1}{arg2} to mean
                             % (arg1, arg2).

%\newcommand{\ip}[2]{\langle #1 | #2\rangle}
                             % This is an alternative definition of
                             % \ip that is commented out.

\begin{document}             % End of preamble and beginning of text.

\maketitle                   % Produces the title.

This is an example input file.  Comparing it with
the output it generates can show you how to
produce a simple document of your own.

\section{Ordinary Text}      % Produces section heading.  Lower-level
                             % sections are begun with similar
                             % \subsection and \subsubsection commands.

The ends  of words and sentences are marked
  by   spaces. It  doesn't matter how many
spaces    you type; one is as good as 100.  The
end of   a line counts as a space.

One   or more   blank lines denote the  end
of  a paragraph.

Since any number of consecutive spaces are treated
like a single one, the formatting of the input
file makes no difference to
      \LaTeX,                % The \LaTeX command generates the LaTeX logo.
but it makes a difference to you.  When you use
\LaTeX, making your input file as easy to read
as possible will be a great help as you write
your document and when you change it.  This sample
file shows how you can add comments to your own input
file.

Because printing is different from typewriting,
there are a number of things that you have to do
differently when preparing an input file than if
you were just typing the document directly.
Quotation marks like
       ``this''
have to be handled specially, as do quotes within
quotes:
       ``\,`this'            % \, separates the double and single quote.
        is what I just
        wrote, not  `that'\,''.

Dashes come in three sizes: an
       intra-word
dash, a medium dash for number ranges like
       1--2,
and a punctuation
       dash---like
this.

A sentence-ending space should be larger than the
space between words within a sentence.  You
sometimes have to type special commands in
conjunction with punctuation characters to get
this right, as in the following sentence.
       Gnats, gnus, etc.\ all  % `\ ' makes an inter-word space.
       begin with G\@.         % \@ marks end-of-sentence punctuation.
You should check the spaces after periods when
reading your output to make sure you haven't
forgotten any special cases.  Generating an
ellipsis
       \ldots\               % `\ ' is needed after `\ldots' because TeX
                             % ignores spaces after command names like \ldots
                             % made from \ + letters.
                             %
                             % Note how a `%' character causes TeX to ignore
                             % the end of the input line, so these blank lines
                             % do not start a new paragraph.
                             %
with the right spacing around the periods requires
a special command.

\LaTeX\ interprets some common characters as
commands, so you must type special commands to
generate them.  These characters include the
following:
       \$ \& \% \# \{ and \}.

In printing, text is usually emphasized with an
       \emph{italic}
type style.

\begin{em}
   A long segment of text can also be emphasized
   in this way.  Text within such a segment can be
   given \emph{additional} emphasis.
\end{em}

It is sometimes necessary to prevent \LaTeX\ from
breaking a line where it might otherwise do so.
This may be at a space, as between the ``Mr.''\ and
``Jones'' in
       ``Mr.~Jones'',        % ~ produces an unbreakable interword space.
or within a word---especially when the word is a
symbol like
       \mbox{\emph{itemnum}}
that makes little sense when hyphenated across
lines.

Footnotes\footnote{This is an example of a footnote.}
pose no problem.

\LaTeX\ is good at typesetting mathematical formulas
like
       \( x-3y + z = 7 \)
or
       \( a_{1} > x^{2n} + y^{2n} > x' \)
or
       \( \ip{A}{B} = \sum_{i} a_{i} b_{i} \).
The spaces you type in a formula are
ignored.  Remember that a letter like
       $x$                   % $ ... $  and  \( ... \)  are equivalent
is a formula when it denotes a mathematical
symbol, and it should be typed as one.

\section{Displayed Text}

Text is displayed by indenting it from the left
margin.  Quotations are commonly displayed.  There
are short quotations
\begin{quote}
   This is a short quotation.  It consists of a
   single paragraph of text.  See how it is formatted.
\end{quote}
and longer ones.
\begin{quotation}
   This is a longer quotation.  It consists of two
   paragraphs of text, neither of which are
   particularly interesting.

   This is the second paragraph of the quotation.  It
   is just as dull as the first paragraph.
\end{quotation}
Another frequently-displayed structure is a list.
The following is an example of an \emph{itemized}
list.
\begin{itemize}
   \item This is the first item of an itemized list.
         Each item in the list is marked with a ``tick''.
         You don't have to worry about what kind of tick
         mark is used.

   \item This is the second item of the list.  It
         contains another list nested inside it.  The inner
         list is an \emph{enumerated} list.
         \begin{enumerate}
            \item This is the first item of an enumerated
                  list that is nested within the itemized list.

            \item This is the second item of the inner list.
                  \LaTeX\ allows you to nest lists deeper than
                  you really should.
         \end{enumerate}
         This is the rest of the second item of the outer
         list.  It is no more interesting than any other
         part of the item.
   \item This is the third item of the list.
\end{itemize}
You can even display poetry.
\begin{verse}
   There is an environment
    for verse \\             % The \\ command separates lines
   Whose features some poets % within a stanza.
   will curse.

                             % One or more blank lines separate stanzas.

   For instead of making\\
   Them do \emph{all} line breaking, \\
   It allows them to put too many words on a line when they'd rather be
   forced to be terse.
\end{verse}

Mathematical formulas may also be displayed.  A
displayed formula
is
one-line long; multiline
formulas require special formatting instructions.
   \[  \ip{\Gamma}{\psi'} = x'' + y^{2} + z_{i}^{n}\]
Don't start a paragraph with a displayed equation,
nor make one a paragraph by itself.

\end{document}               % End of document.

\section{Bench Huge Section 1314}
% synthetic bench section 1314
% This is a sample LaTeX input file.  (Version of 12 August 2004.)
%
% A '%' character causes TeX to ignore all remaining text on the line,
% and is used for comments like this one.

\documentclass{article}      % Specifies the document class

                             % The preamble begins here.
\title{An Example Document}  % Declares the document's title.
\author{Leslie Lamport}      % Declares the author's name.
\date{January 21, 1994}      % Deleting this command produces today's date.

\newcommand{\ip}[2]{(#1, #2)}
                             % Defines \ip{arg1}{arg2} to mean
                             % (arg1, arg2).

%\newcommand{\ip}[2]{\langle #1 | #2\rangle}
                             % This is an alternative definition of
                             % \ip that is commented out.

\begin{document}             % End of preamble and beginning of text.

\maketitle                   % Produces the title.

This is an example input file.  Comparing it with
the output it generates can show you how to
produce a simple document of your own.

\section{Ordinary Text}      % Produces section heading.  Lower-level
                             % sections are begun with similar
                             % \subsection and \subsubsection commands.

The ends  of words and sentences are marked
  by   spaces. It  doesn't matter how many
spaces    you type; one is as good as 100.  The
end of   a line counts as a space.

One   or more   blank lines denote the  end
of  a paragraph.

Since any number of consecutive spaces are treated
like a single one, the formatting of the input
file makes no difference to
      \LaTeX,                % The \LaTeX command generates the LaTeX logo.
but it makes a difference to you.  When you use
\LaTeX, making your input file as easy to read
as possible will be a great help as you write
your document and when you change it.  This sample
file shows how you can add comments to your own input
file.

Because printing is different from typewriting,
there are a number of things that you have to do
differently when preparing an input file than if
you were just typing the document directly.
Quotation marks like
       ``this''
have to be handled specially, as do quotes within
quotes:
       ``\,`this'            % \, separates the double and single quote.
        is what I just
        wrote, not  `that'\,''.

Dashes come in three sizes: an
       intra-word
dash, a medium dash for number ranges like
       1--2,
and a punctuation
       dash---like
this.

A sentence-ending space should be larger than the
space between words within a sentence.  You
sometimes have to type special commands in
conjunction with punctuation characters to get
this right, as in the following sentence.
       Gnats, gnus, etc.\ all  % `\ ' makes an inter-word space.
       begin with G\@.         % \@ marks end-of-sentence punctuation.
You should check the spaces after periods when
reading your output to make sure you haven't
forgotten any special cases.  Generating an
ellipsis
       \ldots\               % `\ ' is needed after `\ldots' because TeX
                             % ignores spaces after command names like \ldots
                             % made from \ + letters.
                             %
                             % Note how a `%' character causes TeX to ignore
                             % the end of the input line, so these blank lines
                             % do not start a new paragraph.
                             %
with the right spacing around the periods requires
a special command.

\LaTeX\ interprets some common characters as
commands, so you must type special commands to
generate them.  These characters include the
following:
       \$ \& \% \# \{ and \}.

In printing, text is usually emphasized with an
       \emph{italic}
type style.

\begin{em}
   A long segment of text can also be emphasized
   in this way.  Text within such a segment can be
   given \emph{additional} emphasis.
\end{em}

It is sometimes necessary to prevent \LaTeX\ from
breaking a line where it might otherwise do so.
This may be at a space, as between the ``Mr.''\ and
``Jones'' in
       ``Mr.~Jones'',        % ~ produces an unbreakable interword space.
or within a word---especially when the word is a
symbol like
       \mbox{\emph{itemnum}}
that makes little sense when hyphenated across
lines.

Footnotes\footnote{This is an example of a footnote.}
pose no problem.

\LaTeX\ is good at typesetting mathematical formulas
like
       \( x-3y + z = 7 \)
or
       \( a_{1} > x^{2n} + y^{2n} > x' \)
or
       \( \ip{A}{B} = \sum_{i} a_{i} b_{i} \).
The spaces you type in a formula are
ignored.  Remember that a letter like
       $x$                   % $ ... $  and  \( ... \)  are equivalent
is a formula when it denotes a mathematical
symbol, and it should be typed as one.

\section{Displayed Text}

Text is displayed by indenting it from the left
margin.  Quotations are commonly displayed.  There
are short quotations
\begin{quote}
   This is a short quotation.  It consists of a
   single paragraph of text.  See how it is formatted.
\end{quote}
and longer ones.
\begin{quotation}
   This is a longer quotation.  It consists of two
   paragraphs of text, neither of which are
   particularly interesting.

   This is the second paragraph of the quotation.  It
   is just as dull as the first paragraph.
\end{quotation}
Another frequently-displayed structure is a list.
The following is an example of an \emph{itemized}
list.
\begin{itemize}
   \item This is the first item of an itemized list.
         Each item in the list is marked with a ``tick''.
         You don't have to worry about what kind of tick
         mark is used.

   \item This is the second item of the list.  It
         contains another list nested inside it.  The inner
         list is an \emph{enumerated} list.
         \begin{enumerate}
            \item This is the first item of an enumerated
                  list that is nested within the itemized list.

            \item This is the second item of the inner list.
                  \LaTeX\ allows you to nest lists deeper than
                  you really should.
         \end{enumerate}
         This is the rest of the second item of the outer
         list.  It is no more interesting than any other
         part of the item.
   \item This is the third item of the list.
\end{itemize}
You can even display poetry.
\begin{verse}
   There is an environment
    for verse \\             % The \\ command separates lines
   Whose features some poets % within a stanza.
   will curse.

                             % One or more blank lines separate stanzas.

   For instead of making\\
   Them do \emph{all} line breaking, \\
   It allows them to put too many words on a line when they'd rather be
   forced to be terse.
\end{verse}

Mathematical formulas may also be displayed.  A
displayed formula
is
one-line long; multiline
formulas require special formatting instructions.
   \[  \ip{\Gamma}{\psi'} = x'' + y^{2} + z_{i}^{n}\]
Don't start a paragraph with a displayed equation,
nor make one a paragraph by itself.

\end{document}               % End of document.

\section{Bench Huge Section 1315}
% synthetic bench section 1315
% This is a sample LaTeX input file.  (Version of 12 August 2004.)
%
% A '%' character causes TeX to ignore all remaining text on the line,
% and is used for comments like this one.

\documentclass{article}      % Specifies the document class

                             % The preamble begins here.
\title{An Example Document}  % Declares the document's title.
\author{Leslie Lamport}      % Declares the author's name.
\date{January 21, 1994}      % Deleting this command produces today's date.

\newcommand{\ip}[2]{(#1, #2)}
                             % Defines \ip{arg1}{arg2} to mean
                             % (arg1, arg2).

%\newcommand{\ip}[2]{\langle #1 | #2\rangle}
                             % This is an alternative definition of
                             % \ip that is commented out.

\begin{document}             % End of preamble and beginning of text.

\maketitle                   % Produces the title.

This is an example input file.  Comparing it with
the output it generates can show you how to
produce a simple document of your own.

\section{Ordinary Text}      % Produces section heading.  Lower-level
                             % sections are begun with similar
                             % \subsection and \subsubsection commands.

The ends  of words and sentences are marked
  by   spaces. It  doesn't matter how many
spaces    you type; one is as good as 100.  The
end of   a line counts as a space.

One   or more   blank lines denote the  end
of  a paragraph.

Since any number of consecutive spaces are treated
like a single one, the formatting of the input
file makes no difference to
      \LaTeX,                % The \LaTeX command generates the LaTeX logo.
but it makes a difference to you.  When you use
\LaTeX, making your input file as easy to read
as possible will be a great help as you write
your document and when you change it.  This sample
file shows how you can add comments to your own input
file.

Because printing is different from typewriting,
there are a number of things that you have to do
differently when preparing an input file than if
you were just typing the document directly.
Quotation marks like
       ``this''
have to be handled specially, as do quotes within
quotes:
       ``\,`this'            % \, separates the double and single quote.
        is what I just
        wrote, not  `that'\,''.

Dashes come in three sizes: an
       intra-word
dash, a medium dash for number ranges like
       1--2,
and a punctuation
       dash---like
this.

A sentence-ending space should be larger than the
space between words within a sentence.  You
sometimes have to type special commands in
conjunction with punctuation characters to get
this right, as in the following sentence.
       Gnats, gnus, etc.\ all  % `\ ' makes an inter-word space.
       begin with G\@.         % \@ marks end-of-sentence punctuation.
You should check the spaces after periods when
reading your output to make sure you haven't
forgotten any special cases.  Generating an
ellipsis
       \ldots\               % `\ ' is needed after `\ldots' because TeX
                             % ignores spaces after command names like \ldots
                             % made from \ + letters.
                             %
                             % Note how a `%' character causes TeX to ignore
                             % the end of the input line, so these blank lines
                             % do not start a new paragraph.
                             %
with the right spacing around the periods requires
a special command.

\LaTeX\ interprets some common characters as
commands, so you must type special commands to
generate them.  These characters include the
following:
       \$ \& \% \# \{ and \}.

In printing, text is usually emphasized with an
       \emph{italic}
type style.

\begin{em}
   A long segment of text can also be emphasized
   in this way.  Text within such a segment can be
   given \emph{additional} emphasis.
\end{em}

It is sometimes necessary to prevent \LaTeX\ from
breaking a line where it might otherwise do so.
This may be at a space, as between the ``Mr.''\ and
``Jones'' in
       ``Mr.~Jones'',        % ~ produces an unbreakable interword space.
or within a word---especially when the word is a
symbol like
       \mbox{\emph{itemnum}}
that makes little sense when hyphenated across
lines.

Footnotes\footnote{This is an example of a footnote.}
pose no problem.

\LaTeX\ is good at typesetting mathematical formulas
like
       \( x-3y + z = 7 \)
or
       \( a_{1} > x^{2n} + y^{2n} > x' \)
or
       \( \ip{A}{B} = \sum_{i} a_{i} b_{i} \).
The spaces you type in a formula are
ignored.  Remember that a letter like
       $x$                   % $ ... $  and  \( ... \)  are equivalent
is a formula when it denotes a mathematical
symbol, and it should be typed as one.

\section{Displayed Text}

Text is displayed by indenting it from the left
margin.  Quotations are commonly displayed.  There
are short quotations
\begin{quote}
   This is a short quotation.  It consists of a
   single paragraph of text.  See how it is formatted.
\end{quote}
and longer ones.
\begin{quotation}
   This is a longer quotation.  It consists of two
   paragraphs of text, neither of which are
   particularly interesting.

   This is the second paragraph of the quotation.  It
   is just as dull as the first paragraph.
\end{quotation}
Another frequently-displayed structure is a list.
The following is an example of an \emph{itemized}
list.
\begin{itemize}
   \item This is the first item of an itemized list.
         Each item in the list is marked with a ``tick''.
         You don't have to worry about what kind of tick
         mark is used.

   \item This is the second item of the list.  It
         contains another list nested inside it.  The inner
         list is an \emph{enumerated} list.
         \begin{enumerate}
            \item This is the first item of an enumerated
                  list that is nested within the itemized list.

            \item This is the second item of the inner list.
                  \LaTeX\ allows you to nest lists deeper than
                  you really should.
         \end{enumerate}
         This is the rest of the second item of the outer
         list.  It is no more interesting than any other
         part of the item.
   \item This is the third item of the list.
\end{itemize}
You can even display poetry.
\begin{verse}
   There is an environment
    for verse \\             % The \\ command separates lines
   Whose features some poets % within a stanza.
   will curse.

                             % One or more blank lines separate stanzas.

   For instead of making\\
   Them do \emph{all} line breaking, \\
   It allows them to put too many words on a line when they'd rather be
   forced to be terse.
\end{verse}

Mathematical formulas may also be displayed.  A
displayed formula
is
one-line long; multiline
formulas require special formatting instructions.
   \[  \ip{\Gamma}{\psi'} = x'' + y^{2} + z_{i}^{n}\]
Don't start a paragraph with a displayed equation,
nor make one a paragraph by itself.

\end{document}               % End of document.

\section{Bench Huge Section 1316}
% synthetic bench section 1316
% This is a sample LaTeX input file.  (Version of 12 August 2004.)
%
% A '%' character causes TeX to ignore all remaining text on the line,
% and is used for comments like this one.

\documentclass{article}      % Specifies the document class

                             % The preamble begins here.
\title{An Example Document}  % Declares the document's title.
\author{Leslie Lamport}      % Declares the author's name.
\date{January 21, 1994}      % Deleting this command produces today's date.

\newcommand{\ip}[2]{(#1, #2)}
                             % Defines \ip{arg1}{arg2} to mean
                             % (arg1, arg2).

%\newcommand{\ip}[2]{\langle #1 | #2\rangle}
                             % This is an alternative definition of
                             % \ip that is commented out.

\begin{document}             % End of preamble and beginning of text.

\maketitle                   % Produces the title.

This is an example input file.  Comparing it with
the output it generates can show you how to
produce a simple document of your own.

\section{Ordinary Text}      % Produces section heading.  Lower-level
                             % sections are begun with similar
                             % \subsection and \subsubsection commands.

The ends  of words and sentences are marked
  by   spaces. It  doesn't matter how many
spaces    you type; one is as good as 100.  The
end of   a line counts as a space.

One   or more   blank lines denote the  end
of  a paragraph.

Since any number of consecutive spaces are treated
like a single one, the formatting of the input
file makes no difference to
      \LaTeX,                % The \LaTeX command generates the LaTeX logo.
but it makes a difference to you.  When you use
\LaTeX, making your input file as easy to read
as possible will be a great help as you write
your document and when you change it.  This sample
file shows how you can add comments to your own input
file.

Because printing is different from typewriting,
there are a number of things that you have to do
differently when preparing an input file than if
you were just typing the document directly.
Quotation marks like
       ``this''
have to be handled specially, as do quotes within
quotes:
       ``\,`this'            % \, separates the double and single quote.
        is what I just
        wrote, not  `that'\,''.

Dashes come in three sizes: an
       intra-word
dash, a medium dash for number ranges like
       1--2,
and a punctuation
       dash---like
this.

A sentence-ending space should be larger than the
space between words within a sentence.  You
sometimes have to type special commands in
conjunction with punctuation characters to get
this right, as in the following sentence.
       Gnats, gnus, etc.\ all  % `\ ' makes an inter-word space.
       begin with G\@.         % \@ marks end-of-sentence punctuation.
You should check the spaces after periods when
reading your output to make sure you haven't
forgotten any special cases.  Generating an
ellipsis
       \ldots\               % `\ ' is needed after `\ldots' because TeX
                             % ignores spaces after command names like \ldots
                             % made from \ + letters.
                             %
                             % Note how a `%' character causes TeX to ignore
                             % the end of the input line, so these blank lines
                             % do not start a new paragraph.
                             %
with the right spacing around the periods requires
a special command.

\LaTeX\ interprets some common characters as
commands, so you must type special commands to
generate them.  These characters include the
following:
       \$ \& \% \# \{ and \}.

In printing, text is usually emphasized with an
       \emph{italic}
type style.

\begin{em}
   A long segment of text can also be emphasized
   in this way.  Text within such a segment can be
   given \emph{additional} emphasis.
\end{em}

It is sometimes necessary to prevent \LaTeX\ from
breaking a line where it might otherwise do so.
This may be at a space, as between the ``Mr.''\ and
``Jones'' in
       ``Mr.~Jones'',        % ~ produces an unbreakable interword space.
or within a word---especially when the word is a
symbol like
       \mbox{\emph{itemnum}}
that makes little sense when hyphenated across
lines.

Footnotes\footnote{This is an example of a footnote.}
pose no problem.

\LaTeX\ is good at typesetting mathematical formulas
like
       \( x-3y + z = 7 \)
or
       \( a_{1} > x^{2n} + y^{2n} > x' \)
or
       \( \ip{A}{B} = \sum_{i} a_{i} b_{i} \).
The spaces you type in a formula are
ignored.  Remember that a letter like
       $x$                   % $ ... $  and  \( ... \)  are equivalent
is a formula when it denotes a mathematical
symbol, and it should be typed as one.

\section{Displayed Text}

Text is displayed by indenting it from the left
margin.  Quotations are commonly displayed.  There
are short quotations
\begin{quote}
   This is a short quotation.  It consists of a
   single paragraph of text.  See how it is formatted.
\end{quote}
and longer ones.
\begin{quotation}
   This is a longer quotation.  It consists of two
   paragraphs of text, neither of which are
   particularly interesting.

   This is the second paragraph of the quotation.  It
   is just as dull as the first paragraph.
\end{quotation}
Another frequently-displayed structure is a list.
The following is an example of an \emph{itemized}
list.
\begin{itemize}
   \item This is the first item of an itemized list.
         Each item in the list is marked with a ``tick''.
         You don't have to worry about what kind of tick
         mark is used.

   \item This is the second item of the list.  It
         contains another list nested inside it.  The inner
         list is an \emph{enumerated} list.
         \begin{enumerate}
            \item This is the first item of an enumerated
                  list that is nested within the itemized list.

            \item This is the second item of the inner list.
                  \LaTeX\ allows you to nest lists deeper than
                  you really should.
         \end{enumerate}
         This is the rest of the second item of the outer
         list.  It is no more interesting than any other
         part of the item.
   \item This is the third item of the list.
\end{itemize}
You can even display poetry.
\begin{verse}
   There is an environment
    for verse \\             % The \\ command separates lines
   Whose features some poets % within a stanza.
   will curse.

                             % One or more blank lines separate stanzas.

   For instead of making\\
   Them do \emph{all} line breaking, \\
   It allows them to put too many words on a line when they'd rather be
   forced to be terse.
\end{verse}

Mathematical formulas may also be displayed.  A
displayed formula
is
one-line long; multiline
formulas require special formatting instructions.
   \[  \ip{\Gamma}{\psi'} = x'' + y^{2} + z_{i}^{n}\]
Don't start a paragraph with a displayed equation,
nor make one a paragraph by itself.

\end{document}               % End of document.

\section{Bench Huge Section 1317}
% synthetic bench section 1317
% This is a sample LaTeX input file.  (Version of 12 August 2004.)
%
% A '%' character causes TeX to ignore all remaining text on the line,
% and is used for comments like this one.

\documentclass{article}      % Specifies the document class

                             % The preamble begins here.
\title{An Example Document}  % Declares the document's title.
\author{Leslie Lamport}      % Declares the author's name.
\date{January 21, 1994}      % Deleting this command produces today's date.

\newcommand{\ip}[2]{(#1, #2)}
                             % Defines \ip{arg1}{arg2} to mean
                             % (arg1, arg2).

%\newcommand{\ip}[2]{\langle #1 | #2\rangle}
                             % This is an alternative definition of
                             % \ip that is commented out.

\begin{document}             % End of preamble and beginning of text.

\maketitle                   % Produces the title.

This is an example input file.  Comparing it with
the output it generates can show you how to
produce a simple document of your own.

\section{Ordinary Text}      % Produces section heading.  Lower-level
                             % sections are begun with similar
                             % \subsection and \subsubsection commands.

The ends  of words and sentences are marked
  by   spaces. It  doesn't matter how many
spaces    you type; one is as good as 100.  The
end of   a line counts as a space.

One   or more   blank lines denote the  end
of  a paragraph.

Since any number of consecutive spaces are treated
like a single one, the formatting of the input
file makes no difference to
      \LaTeX,                % The \LaTeX command generates the LaTeX logo.
but it makes a difference to you.  When you use
\LaTeX, making your input file as easy to read
as possible will be a great help as you write
your document and when you change it.  This sample
file shows how you can add comments to your own input
file.

Because printing is different from typewriting,
there are a number of things that you have to do
differently when preparing an input file than if
you were just typing the document directly.
Quotation marks like
       ``this''
have to be handled specially, as do quotes within
quotes:
       ``\,`this'            % \, separates the double and single quote.
        is what I just
        wrote, not  `that'\,''.

Dashes come in three sizes: an
       intra-word
dash, a medium dash for number ranges like
       1--2,
and a punctuation
       dash---like
this.

A sentence-ending space should be larger than the
space between words within a sentence.  You
sometimes have to type special commands in
conjunction with punctuation characters to get
this right, as in the following sentence.
       Gnats, gnus, etc.\ all  % `\ ' makes an inter-word space.
       begin with G\@.         % \@ marks end-of-sentence punctuation.
You should check the spaces after periods when
reading your output to make sure you haven't
forgotten any special cases.  Generating an
ellipsis
       \ldots\               % `\ ' is needed after `\ldots' because TeX
                             % ignores spaces after command names like \ldots
                             % made from \ + letters.
                             %
                             % Note how a `%' character causes TeX to ignore
                             % the end of the input line, so these blank lines
                             % do not start a new paragraph.
                             %
with the right spacing around the periods requires
a special command.

\LaTeX\ interprets some common characters as
commands, so you must type special commands to
generate them.  These characters include the
following:
       \$ \& \% \# \{ and \}.

In printing, text is usually emphasized with an
       \emph{italic}
type style.

\begin{em}
   A long segment of text can also be emphasized
   in this way.  Text within such a segment can be
   given \emph{additional} emphasis.
\end{em}

It is sometimes necessary to prevent \LaTeX\ from
breaking a line where it might otherwise do so.
This may be at a space, as between the ``Mr.''\ and
``Jones'' in
       ``Mr.~Jones'',        % ~ produces an unbreakable interword space.
or within a word---especially when the word is a
symbol like
       \mbox{\emph{itemnum}}
that makes little sense when hyphenated across
lines.

Footnotes\footnote{This is an example of a footnote.}
pose no problem.

\LaTeX\ is good at typesetting mathematical formulas
like
       \( x-3y + z = 7 \)
or
       \( a_{1} > x^{2n} + y^{2n} > x' \)
or
       \( \ip{A}{B} = \sum_{i} a_{i} b_{i} \).
The spaces you type in a formula are
ignored.  Remember that a letter like
       $x$                   % $ ... $  and  \( ... \)  are equivalent
is a formula when it denotes a mathematical
symbol, and it should be typed as one.

\section{Displayed Text}

Text is displayed by indenting it from the left
margin.  Quotations are commonly displayed.  There
are short quotations
\begin{quote}
   This is a short quotation.  It consists of a
   single paragraph of text.  See how it is formatted.
\end{quote}
and longer ones.
\begin{quotation}
   This is a longer quotation.  It consists of two
   paragraphs of text, neither of which are
   particularly interesting.

   This is the second paragraph of the quotation.  It
   is just as dull as the first paragraph.
\end{quotation}
Another frequently-displayed structure is a list.
The following is an example of an \emph{itemized}
list.
\begin{itemize}
   \item This is the first item of an itemized list.
         Each item in the list is marked with a ``tick''.
         You don't have to worry about what kind of tick
         mark is used.

   \item This is the second item of the list.  It
         contains another list nested inside it.  The inner
         list is an \emph{enumerated} list.
         \begin{enumerate}
            \item This is the first item of an enumerated
                  list that is nested within the itemized list.

            \item This is the second item of the inner list.
                  \LaTeX\ allows you to nest lists deeper than
                  you really should.
         \end{enumerate}
         This is the rest of the second item of the outer
         list.  It is no more interesting than any other
         part of the item.
   \item This is the third item of the list.
\end{itemize}
You can even display poetry.
\begin{verse}
   There is an environment
    for verse \\             % The \\ command separates lines
   Whose features some poets % within a stanza.
   will curse.

                             % One or more blank lines separate stanzas.

   For instead of making\\
   Them do \emph{all} line breaking, \\
   It allows them to put too many words on a line when they'd rather be
   forced to be terse.
\end{verse}

Mathematical formulas may also be displayed.  A
displayed formula
is
one-line long; multiline
formulas require special formatting instructions.
   \[  \ip{\Gamma}{\psi'} = x'' + y^{2} + z_{i}^{n}\]
Don't start a paragraph with a displayed equation,
nor make one a paragraph by itself.

\end{document}               % End of document.

\section{Bench Huge Section 1318}
% synthetic bench section 1318
% This is a sample LaTeX input file.  (Version of 12 August 2004.)
%
% A '%' character causes TeX to ignore all remaining text on the line,
% and is used for comments like this one.

\documentclass{article}      % Specifies the document class

                             % The preamble begins here.
\title{An Example Document}  % Declares the document's title.
\author{Leslie Lamport}      % Declares the author's name.
\date{January 21, 1994}      % Deleting this command produces today's date.

\newcommand{\ip}[2]{(#1, #2)}
                             % Defines \ip{arg1}{arg2} to mean
                             % (arg1, arg2).

%\newcommand{\ip}[2]{\langle #1 | #2\rangle}
                             % This is an alternative definition of
                             % \ip that is commented out.

\begin{document}             % End of preamble and beginning of text.

\maketitle                   % Produces the title.

This is an example input file.  Comparing it with
the output it generates can show you how to
produce a simple document of your own.

\section{Ordinary Text}      % Produces section heading.  Lower-level
                             % sections are begun with similar
                             % \subsection and \subsubsection commands.

The ends  of words and sentences are marked
  by   spaces. It  doesn't matter how many
spaces    you type; one is as good as 100.  The
end of   a line counts as a space.

One   or more   blank lines denote the  end
of  a paragraph.

Since any number of consecutive spaces are treated
like a single one, the formatting of the input
file makes no difference to
      \LaTeX,                % The \LaTeX command generates the LaTeX logo.
but it makes a difference to you.  When you use
\LaTeX, making your input file as easy to read
as possible will be a great help as you write
your document and when you change it.  This sample
file shows how you can add comments to your own input
file.

Because printing is different from typewriting,
there are a number of things that you have to do
differently when preparing an input file than if
you were just typing the document directly.
Quotation marks like
       ``this''
have to be handled specially, as do quotes within
quotes:
       ``\,`this'            % \, separates the double and single quote.
        is what I just
        wrote, not  `that'\,''.

Dashes come in three sizes: an
       intra-word
dash, a medium dash for number ranges like
       1--2,
and a punctuation
       dash---like
this.

A sentence-ending space should be larger than the
space between words within a sentence.  You
sometimes have to type special commands in
conjunction with punctuation characters to get
this right, as in the following sentence.
       Gnats, gnus, etc.\ all  % `\ ' makes an inter-word space.
       begin with G\@.         % \@ marks end-of-sentence punctuation.
You should check the spaces after periods when
reading your output to make sure you haven't
forgotten any special cases.  Generating an
ellipsis
       \ldots\               % `\ ' is needed after `\ldots' because TeX
                             % ignores spaces after command names like \ldots
                             % made from \ + letters.
                             %
                             % Note how a `%' character causes TeX to ignore
                             % the end of the input line, so these blank lines
                             % do not start a new paragraph.
                             %
with the right spacing around the periods requires
a special command.

\LaTeX\ interprets some common characters as
commands, so you must type special commands to
generate them.  These characters include the
following:
       \$ \& \% \# \{ and \}.

In printing, text is usually emphasized with an
       \emph{italic}
type style.

\begin{em}
   A long segment of text can also be emphasized
   in this way.  Text within such a segment can be
   given \emph{additional} emphasis.
\end{em}

It is sometimes necessary to prevent \LaTeX\ from
breaking a line where it might otherwise do so.
This may be at a space, as between the ``Mr.''\ and
``Jones'' in
       ``Mr.~Jones'',        % ~ produces an unbreakable interword space.
or within a word---especially when the word is a
symbol like
       \mbox{\emph{itemnum}}
that makes little sense when hyphenated across
lines.

Footnotes\footnote{This is an example of a footnote.}
pose no problem.

\LaTeX\ is good at typesetting mathematical formulas
like
       \( x-3y + z = 7 \)
or
       \( a_{1} > x^{2n} + y^{2n} > x' \)
or
       \( \ip{A}{B} = \sum_{i} a_{i} b_{i} \).
The spaces you type in a formula are
ignored.  Remember that a letter like
       $x$                   % $ ... $  and  \( ... \)  are equivalent
is a formula when it denotes a mathematical
symbol, and it should be typed as one.

\section{Displayed Text}

Text is displayed by indenting it from the left
margin.  Quotations are commonly displayed.  There
are short quotations
\begin{quote}
   This is a short quotation.  It consists of a
   single paragraph of text.  See how it is formatted.
\end{quote}
and longer ones.
\begin{quotation}
   This is a longer quotation.  It consists of two
   paragraphs of text, neither of which are
   particularly interesting.

   This is the second paragraph of the quotation.  It
   is just as dull as the first paragraph.
\end{quotation}
Another frequently-displayed structure is a list.
The following is an example of an \emph{itemized}
list.
\begin{itemize}
   \item This is the first item of an itemized list.
         Each item in the list is marked with a ``tick''.
         You don't have to worry about what kind of tick
         mark is used.

   \item This is the second item of the list.  It
         contains another list nested inside it.  The inner
         list is an \emph{enumerated} list.
         \begin{enumerate}
            \item This is the first item of an enumerated
                  list that is nested within the itemized list.

            \item This is the second item of the inner list.
                  \LaTeX\ allows you to nest lists deeper than
                  you really should.
         \end{enumerate}
         This is the rest of the second item of the outer
         list.  It is no more interesting than any other
         part of the item.
   \item This is the third item of the list.
\end{itemize}
You can even display poetry.
\begin{verse}
   There is an environment
    for verse \\             % The \\ command separates lines
   Whose features some poets % within a stanza.
   will curse.

                             % One or more blank lines separate stanzas.

   For instead of making\\
   Them do \emph{all} line breaking, \\
   It allows them to put too many words on a line when they'd rather be
   forced to be terse.
\end{verse}

Mathematical formulas may also be displayed.  A
displayed formula
is
one-line long; multiline
formulas require special formatting instructions.
   \[  \ip{\Gamma}{\psi'} = x'' + y^{2} + z_{i}^{n}\]
Don't start a paragraph with a displayed equation,
nor make one a paragraph by itself.

\end{document}               % End of document.

\section{Bench Huge Section 1319}
% synthetic bench section 1319
% This is a sample LaTeX input file.  (Version of 12 August 2004.)
%
% A '%' character causes TeX to ignore all remaining text on the line,
% and is used for comments like this one.

\documentclass{article}      % Specifies the document class

                             % The preamble begins here.
\title{An Example Document}  % Declares the document's title.
\author{Leslie Lamport}      % Declares the author's name.
\date{January 21, 1994}      % Deleting this command produces today's date.

\newcommand{\ip}[2]{(#1, #2)}
                             % Defines \ip{arg1}{arg2} to mean
                             % (arg1, arg2).

%\newcommand{\ip}[2]{\langle #1 | #2\rangle}
                             % This is an alternative definition of
                             % \ip that is commented out.

\begin{document}             % End of preamble and beginning of text.

\maketitle                   % Produces the title.

This is an example input file.  Comparing it with
the output it generates can show you how to
produce a simple document of your own.

\section{Ordinary Text}      % Produces section heading.  Lower-level
                             % sections are begun with similar
                             % \subsection and \subsubsection commands.

The ends  of words and sentences are marked
  by   spaces. It  doesn't matter how many
spaces    you type; one is as good as 100.  The
end of   a line counts as a space.

One   or more   blank lines denote the  end
of  a paragraph.

Since any number of consecutive spaces are treated
like a single one, the formatting of the input
file makes no difference to
      \LaTeX,                % The \LaTeX command generates the LaTeX logo.
but it makes a difference to you.  When you use
\LaTeX, making your input file as easy to read
as possible will be a great help as you write
your document and when you change it.  This sample
file shows how you can add comments to your own input
file.

Because printing is different from typewriting,
there are a number of things that you have to do
differently when preparing an input file than if
you were just typing the document directly.
Quotation marks like
       ``this''
have to be handled specially, as do quotes within
quotes:
       ``\,`this'            % \, separates the double and single quote.
        is what I just
        wrote, not  `that'\,''.

Dashes come in three sizes: an
       intra-word
dash, a medium dash for number ranges like
       1--2,
and a punctuation
       dash---like
this.

A sentence-ending space should be larger than the
space between words within a sentence.  You
sometimes have to type special commands in
conjunction with punctuation characters to get
this right, as in the following sentence.
       Gnats, gnus, etc.\ all  % `\ ' makes an inter-word space.
       begin with G\@.         % \@ marks end-of-sentence punctuation.
You should check the spaces after periods when
reading your output to make sure you haven't
forgotten any special cases.  Generating an
ellipsis
       \ldots\               % `\ ' is needed after `\ldots' because TeX
                             % ignores spaces after command names like \ldots
                             % made from \ + letters.
                             %
                             % Note how a `%' character causes TeX to ignore
                             % the end of the input line, so these blank lines
                             % do not start a new paragraph.
                             %
with the right spacing around the periods requires
a special command.

\LaTeX\ interprets some common characters as
commands, so you must type special commands to
generate them.  These characters include the
following:
       \$ \& \% \# \{ and \}.

In printing, text is usually emphasized with an
       \emph{italic}
type style.

\begin{em}
   A long segment of text can also be emphasized
   in this way.  Text within such a segment can be
   given \emph{additional} emphasis.
\end{em}

It is sometimes necessary to prevent \LaTeX\ from
breaking a line where it might otherwise do so.
This may be at a space, as between the ``Mr.''\ and
``Jones'' in
       ``Mr.~Jones'',        % ~ produces an unbreakable interword space.
or within a word---especially when the word is a
symbol like
       \mbox{\emph{itemnum}}
that makes little sense when hyphenated across
lines.

Footnotes\footnote{This is an example of a footnote.}
pose no problem.

\LaTeX\ is good at typesetting mathematical formulas
like
       \( x-3y + z = 7 \)
or
       \( a_{1} > x^{2n} + y^{2n} > x' \)
or
       \( \ip{A}{B} = \sum_{i} a_{i} b_{i} \).
The spaces you type in a formula are
ignored.  Remember that a letter like
       $x$                   % $ ... $  and  \( ... \)  are equivalent
is a formula when it denotes a mathematical
symbol, and it should be typed as one.

\section{Displayed Text}

Text is displayed by indenting it from the left
margin.  Quotations are commonly displayed.  There
are short quotations
\begin{quote}
   This is a short quotation.  It consists of a
   single paragraph of text.  See how it is formatted.
\end{quote}
and longer ones.
\begin{quotation}
   This is a longer quotation.  It consists of two
   paragraphs of text, neither of which are
   particularly interesting.

   This is the second paragraph of the quotation.  It
   is just as dull as the first paragraph.
\end{quotation}
Another frequently-displayed structure is a list.
The following is an example of an \emph{itemized}
list.
\begin{itemize}
   \item This is the first item of an itemized list.
         Each item in the list is marked with a ``tick''.
         You don't have to worry about what kind of tick
         mark is used.

   \item This is the second item of the list.  It
         contains another list nested inside it.  The inner
         list is an \emph{enumerated} list.
         \begin{enumerate}
            \item This is the first item of an enumerated
                  list that is nested within the itemized list.

            \item This is the second item of the inner list.
                  \LaTeX\ allows you to nest lists deeper than
                  you really should.
         \end{enumerate}
         This is the rest of the second item of the outer
         list.  It is no more interesting than any other
         part of the item.
   \item This is the third item of the list.
\end{itemize}
You can even display poetry.
\begin{verse}
   There is an environment
    for verse \\             % The \\ command separates lines
   Whose features some poets % within a stanza.
   will curse.

                             % One or more blank lines separate stanzas.

   For instead of making\\
   Them do \emph{all} line breaking, \\
   It allows them to put too many words on a line when they'd rather be
   forced to be terse.
\end{verse}

Mathematical formulas may also be displayed.  A
displayed formula
is
one-line long; multiline
formulas require special formatting instructions.
   \[  \ip{\Gamma}{\psi'} = x'' + y^{2} + z_{i}^{n}\]
Don't start a paragraph with a displayed equation,
nor make one a paragraph by itself.

\end{document}               % End of document.

\section{Bench Huge Section 1320}
% synthetic bench section 1320
% This is a sample LaTeX input file.  (Version of 12 August 2004.)
%
% A '%' character causes TeX to ignore all remaining text on the line,
% and is used for comments like this one.

\documentclass{article}      % Specifies the document class

                             % The preamble begins here.
\title{An Example Document}  % Declares the document's title.
\author{Leslie Lamport}      % Declares the author's name.
\date{January 21, 1994}      % Deleting this command produces today's date.

\newcommand{\ip}[2]{(#1, #2)}
                             % Defines \ip{arg1}{arg2} to mean
                             % (arg1, arg2).

%\newcommand{\ip}[2]{\langle #1 | #2\rangle}
                             % This is an alternative definition of
                             % \ip that is commented out.

\begin{document}             % End of preamble and beginning of text.

\maketitle                   % Produces the title.

This is an example input file.  Comparing it with
the output it generates can show you how to
produce a simple document of your own.

\section{Ordinary Text}      % Produces section heading.  Lower-level
                             % sections are begun with similar
                             % \subsection and \subsubsection commands.

The ends  of words and sentences are marked
  by   spaces. It  doesn't matter how many
spaces    you type; one is as good as 100.  The
end of   a line counts as a space.

One   or more   blank lines denote the  end
of  a paragraph.

Since any number of consecutive spaces are treated
like a single one, the formatting of the input
file makes no difference to
      \LaTeX,                % The \LaTeX command generates the LaTeX logo.
but it makes a difference to you.  When you use
\LaTeX, making your input file as easy to read
as possible will be a great help as you write
your document and when you change it.  This sample
file shows how you can add comments to your own input
file.

Because printing is different from typewriting,
there are a number of things that you have to do
differently when preparing an input file than if
you were just typing the document directly.
Quotation marks like
       ``this''
have to be handled specially, as do quotes within
quotes:
       ``\,`this'            % \, separates the double and single quote.
        is what I just
        wrote, not  `that'\,''.

Dashes come in three sizes: an
       intra-word
dash, a medium dash for number ranges like
       1--2,
and a punctuation
       dash---like
this.

A sentence-ending space should be larger than the
space between words within a sentence.  You
sometimes have to type special commands in
conjunction with punctuation characters to get
this right, as in the following sentence.
       Gnats, gnus, etc.\ all  % `\ ' makes an inter-word space.
       begin with G\@.         % \@ marks end-of-sentence punctuation.
You should check the spaces after periods when
reading your output to make sure you haven't
forgotten any special cases.  Generating an
ellipsis
       \ldots\               % `\ ' is needed after `\ldots' because TeX
                             % ignores spaces after command names like \ldots
                             % made from \ + letters.
                             %
                             % Note how a `%' character causes TeX to ignore
                             % the end of the input line, so these blank lines
                             % do not start a new paragraph.
                             %
with the right spacing around the periods requires
a special command.

\LaTeX\ interprets some common characters as
commands, so you must type special commands to
generate them.  These characters include the
following:
       \$ \& \% \# \{ and \}.

In printing, text is usually emphasized with an
       \emph{italic}
type style.

\begin{em}
   A long segment of text can also be emphasized
   in this way.  Text within such a segment can be
   given \emph{additional} emphasis.
\end{em}

It is sometimes necessary to prevent \LaTeX\ from
breaking a line where it might otherwise do so.
This may be at a space, as between the ``Mr.''\ and
``Jones'' in
       ``Mr.~Jones'',        % ~ produces an unbreakable interword space.
or within a word---especially when the word is a
symbol like
       \mbox{\emph{itemnum}}
that makes little sense when hyphenated across
lines.

Footnotes\footnote{This is an example of a footnote.}
pose no problem.

\LaTeX\ is good at typesetting mathematical formulas
like
       \( x-3y + z = 7 \)
or
       \( a_{1} > x^{2n} + y^{2n} > x' \)
or
       \( \ip{A}{B} = \sum_{i} a_{i} b_{i} \).
The spaces you type in a formula are
ignored.  Remember that a letter like
       $x$                   % $ ... $  and  \( ... \)  are equivalent
is a formula when it denotes a mathematical
symbol, and it should be typed as one.

\section{Displayed Text}

Text is displayed by indenting it from the left
margin.  Quotations are commonly displayed.  There
are short quotations
\begin{quote}
   This is a short quotation.  It consists of a
   single paragraph of text.  See how it is formatted.
\end{quote}
and longer ones.
\begin{quotation}
   This is a longer quotation.  It consists of two
   paragraphs of text, neither of which are
   particularly interesting.

   This is the second paragraph of the quotation.  It
   is just as dull as the first paragraph.
\end{quotation}
Another frequently-displayed structure is a list.
The following is an example of an \emph{itemized}
list.
\begin{itemize}
   \item This is the first item of an itemized list.
         Each item in the list is marked with a ``tick''.
         You don't have to worry about what kind of tick
         mark is used.

   \item This is the second item of the list.  It
         contains another list nested inside it.  The inner
         list is an \emph{enumerated} list.
         \begin{enumerate}
            \item This is the first item of an enumerated
                  list that is nested within the itemized list.

            \item This is the second item of the inner list.
                  \LaTeX\ allows you to nest lists deeper than
                  you really should.
         \end{enumerate}
         This is the rest of the second item of the outer
         list.  It is no more interesting than any other
         part of the item.
   \item This is the third item of the list.
\end{itemize}
You can even display poetry.
\begin{verse}
   There is an environment
    for verse \\             % The \\ command separates lines
   Whose features some poets % within a stanza.
   will curse.

                             % One or more blank lines separate stanzas.

   For instead of making\\
   Them do \emph{all} line breaking, \\
   It allows them to put too many words on a line when they'd rather be
   forced to be terse.
\end{verse}

Mathematical formulas may also be displayed.  A
displayed formula
is
one-line long; multiline
formulas require special formatting instructions.
   \[  \ip{\Gamma}{\psi'} = x'' + y^{2} + z_{i}^{n}\]
Don't start a paragraph with a displayed equation,
nor make one a paragraph by itself.

\end{document}               % End of document.

\section{Bench Huge Section 1321}
% synthetic bench section 1321
% This is a sample LaTeX input file.  (Version of 12 August 2004.)
%
% A '%' character causes TeX to ignore all remaining text on the line,
% and is used for comments like this one.

\documentclass{article}      % Specifies the document class

                             % The preamble begins here.
\title{An Example Document}  % Declares the document's title.
\author{Leslie Lamport}      % Declares the author's name.
\date{January 21, 1994}      % Deleting this command produces today's date.

\newcommand{\ip}[2]{(#1, #2)}
                             % Defines \ip{arg1}{arg2} to mean
                             % (arg1, arg2).

%\newcommand{\ip}[2]{\langle #1 | #2\rangle}
                             % This is an alternative definition of
                             % \ip that is commented out.

\begin{document}             % End of preamble and beginning of text.

\maketitle                   % Produces the title.

This is an example input file.  Comparing it with
the output it generates can show you how to
produce a simple document of your own.

\section{Ordinary Text}      % Produces section heading.  Lower-level
                             % sections are begun with similar
                             % \subsection and \subsubsection commands.

The ends  of words and sentences are marked
  by   spaces. It  doesn't matter how many
spaces    you type; one is as good as 100.  The
end of   a line counts as a space.

One   or more   blank lines denote the  end
of  a paragraph.

Since any number of consecutive spaces are treated
like a single one, the formatting of the input
file makes no difference to
      \LaTeX,                % The \LaTeX command generates the LaTeX logo.
but it makes a difference to you.  When you use
\LaTeX, making your input file as easy to read
as possible will be a great help as you write
your document and when you change it.  This sample
file shows how you can add comments to your own input
file.

Because printing is different from typewriting,
there are a number of things that you have to do
differently when preparing an input file than if
you were just typing the document directly.
Quotation marks like
       ``this''
have to be handled specially, as do quotes within
quotes:
       ``\,`this'            % \, separates the double and single quote.
        is what I just
        wrote, not  `that'\,''.

Dashes come in three sizes: an
       intra-word
dash, a medium dash for number ranges like
       1--2,
and a punctuation
       dash---like
this.

A sentence-ending space should be larger than the
space between words within a sentence.  You
sometimes have to type special commands in
conjunction with punctuation characters to get
this right, as in the following sentence.
       Gnats, gnus, etc.\ all  % `\ ' makes an inter-word space.
       begin with G\@.         % \@ marks end-of-sentence punctuation.
You should check the spaces after periods when
reading your output to make sure you haven't
forgotten any special cases.  Generating an
ellipsis
       \ldots\               % `\ ' is needed after `\ldots' because TeX
                             % ignores spaces after command names like \ldots
                             % made from \ + letters.
                             %
                             % Note how a `%' character causes TeX to ignore
                             % the end of the input line, so these blank lines
                             % do not start a new paragraph.
                             %
with the right spacing around the periods requires
a special command.

\LaTeX\ interprets some common characters as
commands, so you must type special commands to
generate them.  These characters include the
following:
       \$ \& \% \# \{ and \}.

In printing, text is usually emphasized with an
       \emph{italic}
type style.

\begin{em}
   A long segment of text can also be emphasized
   in this way.  Text within such a segment can be
   given \emph{additional} emphasis.
\end{em}

It is sometimes necessary to prevent \LaTeX\ from
breaking a line where it might otherwise do so.
This may be at a space, as between the ``Mr.''\ and
``Jones'' in
       ``Mr.~Jones'',        % ~ produces an unbreakable interword space.
or within a word---especially when the word is a
symbol like
       \mbox{\emph{itemnum}}
that makes little sense when hyphenated across
lines.

Footnotes\footnote{This is an example of a footnote.}
pose no problem.

\LaTeX\ is good at typesetting mathematical formulas
like
       \( x-3y + z = 7 \)
or
       \( a_{1} > x^{2n} + y^{2n} > x' \)
or
       \( \ip{A}{B} = \sum_{i} a_{i} b_{i} \).
The spaces you type in a formula are
ignored.  Remember that a letter like
       $x$                   % $ ... $  and  \( ... \)  are equivalent
is a formula when it denotes a mathematical
symbol, and it should be typed as one.

\section{Displayed Text}

Text is displayed by indenting it from the left
margin.  Quotations are commonly displayed.  There
are short quotations
\begin{quote}
   This is a short quotation.  It consists of a
   single paragraph of text.  See how it is formatted.
\end{quote}
and longer ones.
\begin{quotation}
   This is a longer quotation.  It consists of two
   paragraphs of text, neither of which are
   particularly interesting.

   This is the second paragraph of the quotation.  It
   is just as dull as the first paragraph.
\end{quotation}
Another frequently-displayed structure is a list.
The following is an example of an \emph{itemized}
list.
\begin{itemize}
   \item This is the first item of an itemized list.
         Each item in the list is marked with a ``tick''.
         You don't have to worry about what kind of tick
         mark is used.

   \item This is the second item of the list.  It
         contains another list nested inside it.  The inner
         list is an \emph{enumerated} list.
         \begin{enumerate}
            \item This is the first item of an enumerated
                  list that is nested within the itemized list.

            \item This is the second item of the inner list.
                  \LaTeX\ allows you to nest lists deeper than
                  you really should.
         \end{enumerate}
         This is the rest of the second item of the outer
         list.  It is no more interesting than any other
         part of the item.
   \item This is the third item of the list.
\end{itemize}
You can even display poetry.
\begin{verse}
   There is an environment
    for verse \\             % The \\ command separates lines
   Whose features some poets % within a stanza.
   will curse.

                             % One or more blank lines separate stanzas.

   For instead of making\\
   Them do \emph{all} line breaking, \\
   It allows them to put too many words on a line when they'd rather be
   forced to be terse.
\end{verse}

Mathematical formulas may also be displayed.  A
displayed formula
is
one-line long; multiline
formulas require special formatting instructions.
   \[  \ip{\Gamma}{\psi'} = x'' + y^{2} + z_{i}^{n}\]
Don't start a paragraph with a displayed equation,
nor make one a paragraph by itself.

\end{document}               % End of document.

\section{Bench Huge Section 1322}
% synthetic bench section 1322
% This is a sample LaTeX input file.  (Version of 12 August 2004.)
%
% A '%' character causes TeX to ignore all remaining text on the line,
% and is used for comments like this one.

\documentclass{article}      % Specifies the document class

                             % The preamble begins here.
\title{An Example Document}  % Declares the document's title.
\author{Leslie Lamport}      % Declares the author's name.
\date{January 21, 1994}      % Deleting this command produces today's date.

\newcommand{\ip}[2]{(#1, #2)}
                             % Defines \ip{arg1}{arg2} to mean
                             % (arg1, arg2).

%\newcommand{\ip}[2]{\langle #1 | #2\rangle}
                             % This is an alternative definition of
                             % \ip that is commented out.

\begin{document}             % End of preamble and beginning of text.

\maketitle                   % Produces the title.

This is an example input file.  Comparing it with
the output it generates can show you how to
produce a simple document of your own.

\section{Ordinary Text}      % Produces section heading.  Lower-level
                             % sections are begun with similar
                             % \subsection and \subsubsection commands.

The ends  of words and sentences are marked
  by   spaces. It  doesn't matter how many
spaces    you type; one is as good as 100.  The
end of   a line counts as a space.

One   or more   blank lines denote the  end
of  a paragraph.

Since any number of consecutive spaces are treated
like a single one, the formatting of the input
file makes no difference to
      \LaTeX,                % The \LaTeX command generates the LaTeX logo.
but it makes a difference to you.  When you use
\LaTeX, making your input file as easy to read
as possible will be a great help as you write
your document and when you change it.  This sample
file shows how you can add comments to your own input
file.

Because printing is different from typewriting,
there are a number of things that you have to do
differently when preparing an input file than if
you were just typing the document directly.
Quotation marks like
       ``this''
have to be handled specially, as do quotes within
quotes:
       ``\,`this'            % \, separates the double and single quote.
        is what I just
        wrote, not  `that'\,''.

Dashes come in three sizes: an
       intra-word
dash, a medium dash for number ranges like
       1--2,
and a punctuation
       dash---like
this.

A sentence-ending space should be larger than the
space between words within a sentence.  You
sometimes have to type special commands in
conjunction with punctuation characters to get
this right, as in the following sentence.
       Gnats, gnus, etc.\ all  % `\ ' makes an inter-word space.
       begin with G\@.         % \@ marks end-of-sentence punctuation.
You should check the spaces after periods when
reading your output to make sure you haven't
forgotten any special cases.  Generating an
ellipsis
       \ldots\               % `\ ' is needed after `\ldots' because TeX
                             % ignores spaces after command names like \ldots
                             % made from \ + letters.
                             %
                             % Note how a `%' character causes TeX to ignore
                             % the end of the input line, so these blank lines
                             % do not start a new paragraph.
                             %
with the right spacing around the periods requires
a special command.

\LaTeX\ interprets some common characters as
commands, so you must type special commands to
generate them.  These characters include the
following:
       \$ \& \% \# \{ and \}.

In printing, text is usually emphasized with an
       \emph{italic}
type style.

\begin{em}
   A long segment of text can also be emphasized
   in this way.  Text within such a segment can be
   given \emph{additional} emphasis.
\end{em}

It is sometimes necessary to prevent \LaTeX\ from
breaking a line where it might otherwise do so.
This may be at a space, as between the ``Mr.''\ and
``Jones'' in
       ``Mr.~Jones'',        % ~ produces an unbreakable interword space.
or within a word---especially when the word is a
symbol like
       \mbox{\emph{itemnum}}
that makes little sense when hyphenated across
lines.

Footnotes\footnote{This is an example of a footnote.}
pose no problem.

\LaTeX\ is good at typesetting mathematical formulas
like
       \( x-3y + z = 7 \)
or
       \( a_{1} > x^{2n} + y^{2n} > x' \)
or
       \( \ip{A}{B} = \sum_{i} a_{i} b_{i} \).
The spaces you type in a formula are
ignored.  Remember that a letter like
       $x$                   % $ ... $  and  \( ... \)  are equivalent
is a formula when it denotes a mathematical
symbol, and it should be typed as one.

\section{Displayed Text}

Text is displayed by indenting it from the left
margin.  Quotations are commonly displayed.  There
are short quotations
\begin{quote}
   This is a short quotation.  It consists of a
   single paragraph of text.  See how it is formatted.
\end{quote}
and longer ones.
\begin{quotation}
   This is a longer quotation.  It consists of two
   paragraphs of text, neither of which are
   particularly interesting.

   This is the second paragraph of the quotation.  It
   is just as dull as the first paragraph.
\end{quotation}
Another frequently-displayed structure is a list.
The following is an example of an \emph{itemized}
list.
\begin{itemize}
   \item This is the first item of an itemized list.
         Each item in the list is marked with a ``tick''.
         You don't have to worry about what kind of tick
         mark is used.

   \item This is the second item of the list.  It
         contains another list nested inside it.  The inner
         list is an \emph{enumerated} list.
         \begin{enumerate}
            \item This is the first item of an enumerated
                  list that is nested within the itemized list.

            \item This is the second item of the inner list.
                  \LaTeX\ allows you to nest lists deeper than
                  you really should.
         \end{enumerate}
         This is the rest of the second item of the outer
         list.  It is no more interesting than any other
         part of the item.
   \item This is the third item of the list.
\end{itemize}
You can even display poetry.
\begin{verse}
   There is an environment
    for verse \\             % The \\ command separates lines
   Whose features some poets % within a stanza.
   will curse.

                             % One or more blank lines separate stanzas.

   For instead of making\\
   Them do \emph{all} line breaking, \\
   It allows them to put too many words on a line when they'd rather be
   forced to be terse.
\end{verse}

Mathematical formulas may also be displayed.  A
displayed formula
is
one-line long; multiline
formulas require special formatting instructions.
   \[  \ip{\Gamma}{\psi'} = x'' + y^{2} + z_{i}^{n}\]
Don't start a paragraph with a displayed equation,
nor make one a paragraph by itself.

\end{document}               % End of document.

\section{Bench Huge Section 1323}
% synthetic bench section 1323
% This is a sample LaTeX input file.  (Version of 12 August 2004.)
%
% A '%' character causes TeX to ignore all remaining text on the line,
% and is used for comments like this one.

\documentclass{article}      % Specifies the document class

                             % The preamble begins here.
\title{An Example Document}  % Declares the document's title.
\author{Leslie Lamport}      % Declares the author's name.
\date{January 21, 1994}      % Deleting this command produces today's date.

\newcommand{\ip}[2]{(#1, #2)}
                             % Defines \ip{arg1}{arg2} to mean
                             % (arg1, arg2).

%\newcommand{\ip}[2]{\langle #1 | #2\rangle}
                             % This is an alternative definition of
                             % \ip that is commented out.

\begin{document}             % End of preamble and beginning of text.

\maketitle                   % Produces the title.

This is an example input file.  Comparing it with
the output it generates can show you how to
produce a simple document of your own.

\section{Ordinary Text}      % Produces section heading.  Lower-level
                             % sections are begun with similar
                             % \subsection and \subsubsection commands.

The ends  of words and sentences are marked
  by   spaces. It  doesn't matter how many
spaces    you type; one is as good as 100.  The
end of   a line counts as a space.

One   or more   blank lines denote the  end
of  a paragraph.

Since any number of consecutive spaces are treated
like a single one, the formatting of the input
file makes no difference to
      \LaTeX,                % The \LaTeX command generates the LaTeX logo.
but it makes a difference to you.  When you use
\LaTeX, making your input file as easy to read
as possible will be a great help as you write
your document and when you change it.  This sample
file shows how you can add comments to your own input
file.

Because printing is different from typewriting,
there are a number of things that you have to do
differently when preparing an input file than if
you were just typing the document directly.
Quotation marks like
       ``this''
have to be handled specially, as do quotes within
quotes:
       ``\,`this'            % \, separates the double and single quote.
        is what I just
        wrote, not  `that'\,''.

Dashes come in three sizes: an
       intra-word
dash, a medium dash for number ranges like
       1--2,
and a punctuation
       dash---like
this.

A sentence-ending space should be larger than the
space between words within a sentence.  You
sometimes have to type special commands in
conjunction with punctuation characters to get
this right, as in the following sentence.
       Gnats, gnus, etc.\ all  % `\ ' makes an inter-word space.
       begin with G\@.         % \@ marks end-of-sentence punctuation.
You should check the spaces after periods when
reading your output to make sure you haven't
forgotten any special cases.  Generating an
ellipsis
       \ldots\               % `\ ' is needed after `\ldots' because TeX
                             % ignores spaces after command names like \ldots
                             % made from \ + letters.
                             %
                             % Note how a `%' character causes TeX to ignore
                             % the end of the input line, so these blank lines
                             % do not start a new paragraph.
                             %
with the right spacing around the periods requires
a special command.

\LaTeX\ interprets some common characters as
commands, so you must type special commands to
generate them.  These characters include the
following:
       \$ \& \% \# \{ and \}.

In printing, text is usually emphasized with an
       \emph{italic}
type style.

\begin{em}
   A long segment of text can also be emphasized
   in this way.  Text within such a segment can be
   given \emph{additional} emphasis.
\end{em}

It is sometimes necessary to prevent \LaTeX\ from
breaking a line where it might otherwise do so.
This may be at a space, as between the ``Mr.''\ and
``Jones'' in
       ``Mr.~Jones'',        % ~ produces an unbreakable interword space.
or within a word---especially when the word is a
symbol like
       \mbox{\emph{itemnum}}
that makes little sense when hyphenated across
lines.

Footnotes\footnote{This is an example of a footnote.}
pose no problem.

\LaTeX\ is good at typesetting mathematical formulas
like
       \( x-3y + z = 7 \)
or
       \( a_{1} > x^{2n} + y^{2n} > x' \)
or
       \( \ip{A}{B} = \sum_{i} a_{i} b_{i} \).
The spaces you type in a formula are
ignored.  Remember that a letter like
       $x$                   % $ ... $  and  \( ... \)  are equivalent
is a formula when it denotes a mathematical
symbol, and it should be typed as one.

\section{Displayed Text}

Text is displayed by indenting it from the left
margin.  Quotations are commonly displayed.  There
are short quotations
\begin{quote}
   This is a short quotation.  It consists of a
   single paragraph of text.  See how it is formatted.
\end{quote}
and longer ones.
\begin{quotation}
   This is a longer quotation.  It consists of two
   paragraphs of text, neither of which are
   particularly interesting.

   This is the second paragraph of the quotation.  It
   is just as dull as the first paragraph.
\end{quotation}
Another frequently-displayed structure is a list.
The following is an example of an \emph{itemized}
list.
\begin{itemize}
   \item This is the first item of an itemized list.
         Each item in the list is marked with a ``tick''.
         You don't have to worry about what kind of tick
         mark is used.

   \item This is the second item of the list.  It
         contains another list nested inside it.  The inner
         list is an \emph{enumerated} list.
         \begin{enumerate}
            \item This is the first item of an enumerated
                  list that is nested within the itemized list.

            \item This is the second item of the inner list.
                  \LaTeX\ allows you to nest lists deeper than
                  you really should.
         \end{enumerate}
         This is the rest of the second item of the outer
         list.  It is no more interesting than any other
         part of the item.
   \item This is the third item of the list.
\end{itemize}
You can even display poetry.
\begin{verse}
   There is an environment
    for verse \\             % The \\ command separates lines
   Whose features some poets % within a stanza.
   will curse.

                             % One or more blank lines separate stanzas.

   For instead of making\\
   Them do \emph{all} line breaking, \\
   It allows them to put too many words on a line when they'd rather be
   forced to be terse.
\end{verse}

Mathematical formulas may also be displayed.  A
displayed formula
is
one-line long; multiline
formulas require special formatting instructions.
   \[  \ip{\Gamma}{\psi'} = x'' + y^{2} + z_{i}^{n}\]
Don't start a paragraph with a displayed equation,
nor make one a paragraph by itself.

\end{document}               % End of document.

\section{Bench Huge Section 1324}
% synthetic bench section 1324
% This is a sample LaTeX input file.  (Version of 12 August 2004.)
%
% A '%' character causes TeX to ignore all remaining text on the line,
% and is used for comments like this one.

\documentclass{article}      % Specifies the document class

                             % The preamble begins here.
\title{An Example Document}  % Declares the document's title.
\author{Leslie Lamport}      % Declares the author's name.
\date{January 21, 1994}      % Deleting this command produces today's date.

\newcommand{\ip}[2]{(#1, #2)}
                             % Defines \ip{arg1}{arg2} to mean
                             % (arg1, arg2).

%\newcommand{\ip}[2]{\langle #1 | #2\rangle}
                             % This is an alternative definition of
                             % \ip that is commented out.

\begin{document}             % End of preamble and beginning of text.

\maketitle                   % Produces the title.

This is an example input file.  Comparing it with
the output it generates can show you how to
produce a simple document of your own.

\section{Ordinary Text}      % Produces section heading.  Lower-level
                             % sections are begun with similar
                             % \subsection and \subsubsection commands.

The ends  of words and sentences are marked
  by   spaces. It  doesn't matter how many
spaces    you type; one is as good as 100.  The
end of   a line counts as a space.

One   or more   blank lines denote the  end
of  a paragraph.

Since any number of consecutive spaces are treated
like a single one, the formatting of the input
file makes no difference to
      \LaTeX,                % The \LaTeX command generates the LaTeX logo.
but it makes a difference to you.  When you use
\LaTeX, making your input file as easy to read
as possible will be a great help as you write
your document and when you change it.  This sample
file shows how you can add comments to your own input
file.

Because printing is different from typewriting,
there are a number of things that you have to do
differently when preparing an input file than if
you were just typing the document directly.
Quotation marks like
       ``this''
have to be handled specially, as do quotes within
quotes:
       ``\,`this'            % \, separates the double and single quote.
        is what I just
        wrote, not  `that'\,''.

Dashes come in three sizes: an
       intra-word
dash, a medium dash for number ranges like
       1--2,
and a punctuation
       dash---like
this.

A sentence-ending space should be larger than the
space between words within a sentence.  You
sometimes have to type special commands in
conjunction with punctuation characters to get
this right, as in the following sentence.
       Gnats, gnus, etc.\ all  % `\ ' makes an inter-word space.
       begin with G\@.         % \@ marks end-of-sentence punctuation.
You should check the spaces after periods when
reading your output to make sure you haven't
forgotten any special cases.  Generating an
ellipsis
       \ldots\               % `\ ' is needed after `\ldots' because TeX
                             % ignores spaces after command names like \ldots
                             % made from \ + letters.
                             %
                             % Note how a `%' character causes TeX to ignore
                             % the end of the input line, so these blank lines
                             % do not start a new paragraph.
                             %
with the right spacing around the periods requires
a special command.

\LaTeX\ interprets some common characters as
commands, so you must type special commands to
generate them.  These characters include the
following:
       \$ \& \% \# \{ and \}.

In printing, text is usually emphasized with an
       \emph{italic}
type style.

\begin{em}
   A long segment of text can also be emphasized
   in this way.  Text within such a segment can be
   given \emph{additional} emphasis.
\end{em}

It is sometimes necessary to prevent \LaTeX\ from
breaking a line where it might otherwise do so.
This may be at a space, as between the ``Mr.''\ and
``Jones'' in
       ``Mr.~Jones'',        % ~ produces an unbreakable interword space.
or within a word---especially when the word is a
symbol like
       \mbox{\emph{itemnum}}
that makes little sense when hyphenated across
lines.

Footnotes\footnote{This is an example of a footnote.}
pose no problem.

\LaTeX\ is good at typesetting mathematical formulas
like
       \( x-3y + z = 7 \)
or
       \( a_{1} > x^{2n} + y^{2n} > x' \)
or
       \( \ip{A}{B} = \sum_{i} a_{i} b_{i} \).
The spaces you type in a formula are
ignored.  Remember that a letter like
       $x$                   % $ ... $  and  \( ... \)  are equivalent
is a formula when it denotes a mathematical
symbol, and it should be typed as one.

\section{Displayed Text}

Text is displayed by indenting it from the left
margin.  Quotations are commonly displayed.  There
are short quotations
\begin{quote}
   This is a short quotation.  It consists of a
   single paragraph of text.  See how it is formatted.
\end{quote}
and longer ones.
\begin{quotation}
   This is a longer quotation.  It consists of two
   paragraphs of text, neither of which are
   particularly interesting.

   This is the second paragraph of the quotation.  It
   is just as dull as the first paragraph.
\end{quotation}
Another frequently-displayed structure is a list.
The following is an example of an \emph{itemized}
list.
\begin{itemize}
   \item This is the first item of an itemized list.
         Each item in the list is marked with a ``tick''.
         You don't have to worry about what kind of tick
         mark is used.

   \item This is the second item of the list.  It
         contains another list nested inside it.  The inner
         list is an \emph{enumerated} list.
         \begin{enumerate}
            \item This is the first item of an enumerated
                  list that is nested within the itemized list.

            \item This is the second item of the inner list.
                  \LaTeX\ allows you to nest lists deeper than
                  you really should.
         \end{enumerate}
         This is the rest of the second item of the outer
         list.  It is no more interesting than any other
         part of the item.
   \item This is the third item of the list.
\end{itemize}
You can even display poetry.
\begin{verse}
   There is an environment
    for verse \\             % The \\ command separates lines
   Whose features some poets % within a stanza.
   will curse.

                             % One or more blank lines separate stanzas.

   For instead of making\\
   Them do \emph{all} line breaking, \\
   It allows them to put too many words on a line when they'd rather be
   forced to be terse.
\end{verse}

Mathematical formulas may also be displayed.  A
displayed formula
is
one-line long; multiline
formulas require special formatting instructions.
   \[  \ip{\Gamma}{\psi'} = x'' + y^{2} + z_{i}^{n}\]
Don't start a paragraph with a displayed equation,
nor make one a paragraph by itself.

\end{document}               % End of document.

\section{Bench Huge Section 1325}
% synthetic bench section 1325
% This is a sample LaTeX input file.  (Version of 12 August 2004.)
%
% A '%' character causes TeX to ignore all remaining text on the line,
% and is used for comments like this one.

\documentclass{article}      % Specifies the document class

                             % The preamble begins here.
\title{An Example Document}  % Declares the document's title.
\author{Leslie Lamport}      % Declares the author's name.
\date{January 21, 1994}      % Deleting this command produces today's date.

\newcommand{\ip}[2]{(#1, #2)}
                             % Defines \ip{arg1}{arg2} to mean
                             % (arg1, arg2).

%\newcommand{\ip}[2]{\langle #1 | #2\rangle}
                             % This is an alternative definition of
                             % \ip that is commented out.

\begin{document}             % End of preamble and beginning of text.

\maketitle                   % Produces the title.

This is an example input file.  Comparing it with
the output it generates can show you how to
produce a simple document of your own.

\section{Ordinary Text}      % Produces section heading.  Lower-level
                             % sections are begun with similar
                             % \subsection and \subsubsection commands.

The ends  of words and sentences are marked
  by   spaces. It  doesn't matter how many
spaces    you type; one is as good as 100.  The
end of   a line counts as a space.

One   or more   blank lines denote the  end
of  a paragraph.

Since any number of consecutive spaces are treated
like a single one, the formatting of the input
file makes no difference to
      \LaTeX,                % The \LaTeX command generates the LaTeX logo.
but it makes a difference to you.  When you use
\LaTeX, making your input file as easy to read
as possible will be a great help as you write
your document and when you change it.  This sample
file shows how you can add comments to your own input
file.

Because printing is different from typewriting,
there are a number of things that you have to do
differently when preparing an input file than if
you were just typing the document directly.
Quotation marks like
       ``this''
have to be handled specially, as do quotes within
quotes:
       ``\,`this'            % \, separates the double and single quote.
        is what I just
        wrote, not  `that'\,''.

Dashes come in three sizes: an
       intra-word
dash, a medium dash for number ranges like
       1--2,
and a punctuation
       dash---like
this.

A sentence-ending space should be larger than the
space between words within a sentence.  You
sometimes have to type special commands in
conjunction with punctuation characters to get
this right, as in the following sentence.
       Gnats, gnus, etc.\ all  % `\ ' makes an inter-word space.
       begin with G\@.         % \@ marks end-of-sentence punctuation.
You should check the spaces after periods when
reading your output to make sure you haven't
forgotten any special cases.  Generating an
ellipsis
       \ldots\               % `\ ' is needed after `\ldots' because TeX
                             % ignores spaces after command names like \ldots
                             % made from \ + letters.
                             %
                             % Note how a `%' character causes TeX to ignore
                             % the end of the input line, so these blank lines
                             % do not start a new paragraph.
                             %
with the right spacing around the periods requires
a special command.

\LaTeX\ interprets some common characters as
commands, so you must type special commands to
generate them.  These characters include the
following:
       \$ \& \% \# \{ and \}.

In printing, text is usually emphasized with an
       \emph{italic}
type style.

\begin{em}
   A long segment of text can also be emphasized
   in this way.  Text within such a segment can be
   given \emph{additional} emphasis.
\end{em}

It is sometimes necessary to prevent \LaTeX\ from
breaking a line where it might otherwise do so.
This may be at a space, as between the ``Mr.''\ and
``Jones'' in
       ``Mr.~Jones'',        % ~ produces an unbreakable interword space.
or within a word---especially when the word is a
symbol like
       \mbox{\emph{itemnum}}
that makes little sense when hyphenated across
lines.

Footnotes\footnote{This is an example of a footnote.}
pose no problem.

\LaTeX\ is good at typesetting mathematical formulas
like
       \( x-3y + z = 7 \)
or
       \( a_{1} > x^{2n} + y^{2n} > x' \)
or
       \( \ip{A}{B} = \sum_{i} a_{i} b_{i} \).
The spaces you type in a formula are
ignored.  Remember that a letter like
       $x$                   % $ ... $  and  \( ... \)  are equivalent
is a formula when it denotes a mathematical
symbol, and it should be typed as one.

\section{Displayed Text}

Text is displayed by indenting it from the left
margin.  Quotations are commonly displayed.  There
are short quotations
\begin{quote}
   This is a short quotation.  It consists of a
   single paragraph of text.  See how it is formatted.
\end{quote}
and longer ones.
\begin{quotation}
   This is a longer quotation.  It consists of two
   paragraphs of text, neither of which are
   particularly interesting.

   This is the second paragraph of the quotation.  It
   is just as dull as the first paragraph.
\end{quotation}
Another frequently-displayed structure is a list.
The following is an example of an \emph{itemized}
list.
\begin{itemize}
   \item This is the first item of an itemized list.
         Each item in the list is marked with a ``tick''.
         You don't have to worry about what kind of tick
         mark is used.

   \item This is the second item of the list.  It
         contains another list nested inside it.  The inner
         list is an \emph{enumerated} list.
         \begin{enumerate}
            \item This is the first item of an enumerated
                  list that is nested within the itemized list.

            \item This is the second item of the inner list.
                  \LaTeX\ allows you to nest lists deeper than
                  you really should.
         \end{enumerate}
         This is the rest of the second item of the outer
         list.  It is no more interesting than any other
         part of the item.
   \item This is the third item of the list.
\end{itemize}
You can even display poetry.
\begin{verse}
   There is an environment
    for verse \\             % The \\ command separates lines
   Whose features some poets % within a stanza.
   will curse.

                             % One or more blank lines separate stanzas.

   For instead of making\\
   Them do \emph{all} line breaking, \\
   It allows them to put too many words on a line when they'd rather be
   forced to be terse.
\end{verse}

Mathematical formulas may also be displayed.  A
displayed formula
is
one-line long; multiline
formulas require special formatting instructions.
   \[  \ip{\Gamma}{\psi'} = x'' + y^{2} + z_{i}^{n}\]
Don't start a paragraph with a displayed equation,
nor make one a paragraph by itself.

\end{document}               % End of document.

\section{Bench Huge Section 1326}
% synthetic bench section 1326
% This is a sample LaTeX input file.  (Version of 12 August 2004.)
%
% A '%' character causes TeX to ignore all remaining text on the line,
% and is used for comments like this one.

\documentclass{article}      % Specifies the document class

                             % The preamble begins here.
\title{An Example Document}  % Declares the document's title.
\author{Leslie Lamport}      % Declares the author's name.
\date{January 21, 1994}      % Deleting this command produces today's date.

\newcommand{\ip}[2]{(#1, #2)}
                             % Defines \ip{arg1}{arg2} to mean
                             % (arg1, arg2).

%\newcommand{\ip}[2]{\langle #1 | #2\rangle}
                             % This is an alternative definition of
                             % \ip that is commented out.

\begin{document}             % End of preamble and beginning of text.

\maketitle                   % Produces the title.

This is an example input file.  Comparing it with
the output it generates can show you how to
produce a simple document of your own.

\section{Ordinary Text}      % Produces section heading.  Lower-level
                             % sections are begun with similar
                             % \subsection and \subsubsection commands.

The ends  of words and sentences are marked
  by   spaces. It  doesn't matter how many
spaces    you type; one is as good as 100.  The
end of   a line counts as a space.

One   or more   blank lines denote the  end
of  a paragraph.

Since any number of consecutive spaces are treated
like a single one, the formatting of the input
file makes no difference to
      \LaTeX,                % The \LaTeX command generates the LaTeX logo.
but it makes a difference to you.  When you use
\LaTeX, making your input file as easy to read
as possible will be a great help as you write
your document and when you change it.  This sample
file shows how you can add comments to your own input
file.

Because printing is different from typewriting,
there are a number of things that you have to do
differently when preparing an input file than if
you were just typing the document directly.
Quotation marks like
       ``this''
have to be handled specially, as do quotes within
quotes:
       ``\,`this'            % \, separates the double and single quote.
        is what I just
        wrote, not  `that'\,''.

Dashes come in three sizes: an
       intra-word
dash, a medium dash for number ranges like
       1--2,
and a punctuation
       dash---like
this.

A sentence-ending space should be larger than the
space between words within a sentence.  You
sometimes have to type special commands in
conjunction with punctuation characters to get
this right, as in the following sentence.
       Gnats, gnus, etc.\ all  % `\ ' makes an inter-word space.
       begin with G\@.         % \@ marks end-of-sentence punctuation.
You should check the spaces after periods when
reading your output to make sure you haven't
forgotten any special cases.  Generating an
ellipsis
       \ldots\               % `\ ' is needed after `\ldots' because TeX
                             % ignores spaces after command names like \ldots
                             % made from \ + letters.
                             %
                             % Note how a `%' character causes TeX to ignore
                             % the end of the input line, so these blank lines
                             % do not start a new paragraph.
                             %
with the right spacing around the periods requires
a special command.

\LaTeX\ interprets some common characters as
commands, so you must type special commands to
generate them.  These characters include the
following:
       \$ \& \% \# \{ and \}.

In printing, text is usually emphasized with an
       \emph{italic}
type style.

\begin{em}
   A long segment of text can also be emphasized
   in this way.  Text within such a segment can be
   given \emph{additional} emphasis.
\end{em}

It is sometimes necessary to prevent \LaTeX\ from
breaking a line where it might otherwise do so.
This may be at a space, as between the ``Mr.''\ and
``Jones'' in
       ``Mr.~Jones'',        % ~ produces an unbreakable interword space.
or within a word---especially when the word is a
symbol like
       \mbox{\emph{itemnum}}
that makes little sense when hyphenated across
lines.

Footnotes\footnote{This is an example of a footnote.}
pose no problem.

\LaTeX\ is good at typesetting mathematical formulas
like
       \( x-3y + z = 7 \)
or
       \( a_{1} > x^{2n} + y^{2n} > x' \)
or
       \( \ip{A}{B} = \sum_{i} a_{i} b_{i} \).
The spaces you type in a formula are
ignored.  Remember that a letter like
       $x$                   % $ ... $  and  \( ... \)  are equivalent
is a formula when it denotes a mathematical
symbol, and it should be typed as one.

\section{Displayed Text}

Text is displayed by indenting it from the left
margin.  Quotations are commonly displayed.  There
are short quotations
\begin{quote}
   This is a short quotation.  It consists of a
   single paragraph of text.  See how it is formatted.
\end{quote}
and longer ones.
\begin{quotation}
   This is a longer quotation.  It consists of two
   paragraphs of text, neither of which are
   particularly interesting.

   This is the second paragraph of the quotation.  It
   is just as dull as the first paragraph.
\end{quotation}
Another frequently-displayed structure is a list.
The following is an example of an \emph{itemized}
list.
\begin{itemize}
   \item This is the first item of an itemized list.
         Each item in the list is marked with a ``tick''.
         You don't have to worry about what kind of tick
         mark is used.

   \item This is the second item of the list.  It
         contains another list nested inside it.  The inner
         list is an \emph{enumerated} list.
         \begin{enumerate}
            \item This is the first item of an enumerated
                  list that is nested within the itemized list.

            \item This is the second item of the inner list.
                  \LaTeX\ allows you to nest lists deeper than
                  you really should.
         \end{enumerate}
         This is the rest of the second item of the outer
         list.  It is no more interesting than any other
         part of the item.
   \item This is the third item of the list.
\end{itemize}
You can even display poetry.
\begin{verse}
   There is an environment
    for verse \\             % The \\ command separates lines
   Whose features some poets % within a stanza.
   will curse.

                             % One or more blank lines separate stanzas.

   For instead of making\\
   Them do \emph{all} line breaking, \\
   It allows them to put too many words on a line when they'd rather be
   forced to be terse.
\end{verse}

Mathematical formulas may also be displayed.  A
displayed formula
is
one-line long; multiline
formulas require special formatting instructions.
   \[  \ip{\Gamma}{\psi'} = x'' + y^{2} + z_{i}^{n}\]
Don't start a paragraph with a displayed equation,
nor make one a paragraph by itself.

\end{document}               % End of document.

\section{Bench Huge Section 1327}
% synthetic bench section 1327
% This is a sample LaTeX input file.  (Version of 12 August 2004.)
%
% A '%' character causes TeX to ignore all remaining text on the line,
% and is used for comments like this one.

\documentclass{article}      % Specifies the document class

                             % The preamble begins here.
\title{An Example Document}  % Declares the document's title.
\author{Leslie Lamport}      % Declares the author's name.
\date{January 21, 1994}      % Deleting this command produces today's date.

\newcommand{\ip}[2]{(#1, #2)}
                             % Defines \ip{arg1}{arg2} to mean
                             % (arg1, arg2).

%\newcommand{\ip}[2]{\langle #1 | #2\rangle}
                             % This is an alternative definition of
                             % \ip that is commented out.

\begin{document}             % End of preamble and beginning of text.

\maketitle                   % Produces the title.

This is an example input file.  Comparing it with
the output it generates can show you how to
produce a simple document of your own.

\section{Ordinary Text}      % Produces section heading.  Lower-level
                             % sections are begun with similar
                             % \subsection and \subsubsection commands.

The ends  of words and sentences are marked
  by   spaces. It  doesn't matter how many
spaces    you type; one is as good as 100.  The
end of   a line counts as a space.

One   or more   blank lines denote the  end
of  a paragraph.

Since any number of consecutive spaces are treated
like a single one, the formatting of the input
file makes no difference to
      \LaTeX,                % The \LaTeX command generates the LaTeX logo.
but it makes a difference to you.  When you use
\LaTeX, making your input file as easy to read
as possible will be a great help as you write
your document and when you change it.  This sample
file shows how you can add comments to your own input
file.

Because printing is different from typewriting,
there are a number of things that you have to do
differently when preparing an input file than if
you were just typing the document directly.
Quotation marks like
       ``this''
have to be handled specially, as do quotes within
quotes:
       ``\,`this'            % \, separates the double and single quote.
        is what I just
        wrote, not  `that'\,''.

Dashes come in three sizes: an
       intra-word
dash, a medium dash for number ranges like
       1--2,
and a punctuation
       dash---like
this.

A sentence-ending space should be larger than the
space between words within a sentence.  You
sometimes have to type special commands in
conjunction with punctuation characters to get
this right, as in the following sentence.
       Gnats, gnus, etc.\ all  % `\ ' makes an inter-word space.
       begin with G\@.         % \@ marks end-of-sentence punctuation.
You should check the spaces after periods when
reading your output to make sure you haven't
forgotten any special cases.  Generating an
ellipsis
       \ldots\               % `\ ' is needed after `\ldots' because TeX
                             % ignores spaces after command names like \ldots
                             % made from \ + letters.
                             %
                             % Note how a `%' character causes TeX to ignore
                             % the end of the input line, so these blank lines
                             % do not start a new paragraph.
                             %
with the right spacing around the periods requires
a special command.

\LaTeX\ interprets some common characters as
commands, so you must type special commands to
generate them.  These characters include the
following:
       \$ \& \% \# \{ and \}.

In printing, text is usually emphasized with an
       \emph{italic}
type style.

\begin{em}
   A long segment of text can also be emphasized
   in this way.  Text within such a segment can be
   given \emph{additional} emphasis.
\end{em}

It is sometimes necessary to prevent \LaTeX\ from
breaking a line where it might otherwise do so.
This may be at a space, as between the ``Mr.''\ and
``Jones'' in
       ``Mr.~Jones'',        % ~ produces an unbreakable interword space.
or within a word---especially when the word is a
symbol like
       \mbox{\emph{itemnum}}
that makes little sense when hyphenated across
lines.

Footnotes\footnote{This is an example of a footnote.}
pose no problem.

\LaTeX\ is good at typesetting mathematical formulas
like
       \( x-3y + z = 7 \)
or
       \( a_{1} > x^{2n} + y^{2n} > x' \)
or
       \( \ip{A}{B} = \sum_{i} a_{i} b_{i} \).
The spaces you type in a formula are
ignored.  Remember that a letter like
       $x$                   % $ ... $  and  \( ... \)  are equivalent
is a formula when it denotes a mathematical
symbol, and it should be typed as one.

\section{Displayed Text}

Text is displayed by indenting it from the left
margin.  Quotations are commonly displayed.  There
are short quotations
\begin{quote}
   This is a short quotation.  It consists of a
   single paragraph of text.  See how it is formatted.
\end{quote}
and longer ones.
\begin{quotation}
   This is a longer quotation.  It consists of two
   paragraphs of text, neither of which are
   particularly interesting.

   This is the second paragraph of the quotation.  It
   is just as dull as the first paragraph.
\end{quotation}
Another frequently-displayed structure is a list.
The following is an example of an \emph{itemized}
list.
\begin{itemize}
   \item This is the first item of an itemized list.
         Each item in the list is marked with a ``tick''.
         You don't have to worry about what kind of tick
         mark is used.

   \item This is the second item of the list.  It
         contains another list nested inside it.  The inner
         list is an \emph{enumerated} list.
         \begin{enumerate}
            \item This is the first item of an enumerated
                  list that is nested within the itemized list.

            \item This is the second item of the inner list.
                  \LaTeX\ allows you to nest lists deeper than
                  you really should.
         \end{enumerate}
         This is the rest of the second item of the outer
         list.  It is no more interesting than any other
         part of the item.
   \item This is the third item of the list.
\end{itemize}
You can even display poetry.
\begin{verse}
   There is an environment
    for verse \\             % The \\ command separates lines
   Whose features some poets % within a stanza.
   will curse.

                             % One or more blank lines separate stanzas.

   For instead of making\\
   Them do \emph{all} line breaking, \\
   It allows them to put too many words on a line when they'd rather be
   forced to be terse.
\end{verse}

Mathematical formulas may also be displayed.  A
displayed formula
is
one-line long; multiline
formulas require special formatting instructions.
   \[  \ip{\Gamma}{\psi'} = x'' + y^{2} + z_{i}^{n}\]
Don't start a paragraph with a displayed equation,
nor make one a paragraph by itself.

\end{document}               % End of document.

\section{Bench Huge Section 1328}
% synthetic bench section 1328
% This is a sample LaTeX input file.  (Version of 12 August 2004.)
%
% A '%' character causes TeX to ignore all remaining text on the line,
% and is used for comments like this one.

\documentclass{article}      % Specifies the document class

                             % The preamble begins here.
\title{An Example Document}  % Declares the document's title.
\author{Leslie Lamport}      % Declares the author's name.
\date{January 21, 1994}      % Deleting this command produces today's date.

\newcommand{\ip}[2]{(#1, #2)}
                             % Defines \ip{arg1}{arg2} to mean
                             % (arg1, arg2).

%\newcommand{\ip}[2]{\langle #1 | #2\rangle}
                             % This is an alternative definition of
                             % \ip that is commented out.

\begin{document}             % End of preamble and beginning of text.

\maketitle                   % Produces the title.

This is an example input file.  Comparing it with
the output it generates can show you how to
produce a simple document of your own.

\section{Ordinary Text}      % Produces section heading.  Lower-level
                             % sections are begun with similar
                             % \subsection and \subsubsection commands.

The ends  of words and sentences are marked
  by   spaces. It  doesn't matter how many
spaces    you type; one is as good as 100.  The
end of   a line counts as a space.

One   or more   blank lines denote the  end
of  a paragraph.

Since any number of consecutive spaces are treated
like a single one, the formatting of the input
file makes no difference to
      \LaTeX,                % The \LaTeX command generates the LaTeX logo.
but it makes a difference to you.  When you use
\LaTeX, making your input file as easy to read
as possible will be a great help as you write
your document and when you change it.  This sample
file shows how you can add comments to your own input
file.

Because printing is different from typewriting,
there are a number of things that you have to do
differently when preparing an input file than if
you were just typing the document directly.
Quotation marks like
       ``this''
have to be handled specially, as do quotes within
quotes:
       ``\,`this'            % \, separates the double and single quote.
        is what I just
        wrote, not  `that'\,''.

Dashes come in three sizes: an
       intra-word
dash, a medium dash for number ranges like
       1--2,
and a punctuation
       dash---like
this.

A sentence-ending space should be larger than the
space between words within a sentence.  You
sometimes have to type special commands in
conjunction with punctuation characters to get
this right, as in the following sentence.
       Gnats, gnus, etc.\ all  % `\ ' makes an inter-word space.
       begin with G\@.         % \@ marks end-of-sentence punctuation.
You should check the spaces after periods when
reading your output to make sure you haven't
forgotten any special cases.  Generating an
ellipsis
       \ldots\               % `\ ' is needed after `\ldots' because TeX
                             % ignores spaces after command names like \ldots
                             % made from \ + letters.
                             %
                             % Note how a `%' character causes TeX to ignore
                             % the end of the input line, so these blank lines
                             % do not start a new paragraph.
                             %
with the right spacing around the periods requires
a special command.

\LaTeX\ interprets some common characters as
commands, so you must type special commands to
generate them.  These characters include the
following:
       \$ \& \% \# \{ and \}.

In printing, text is usually emphasized with an
       \emph{italic}
type style.

\begin{em}
   A long segment of text can also be emphasized
   in this way.  Text within such a segment can be
   given \emph{additional} emphasis.
\end{em}

It is sometimes necessary to prevent \LaTeX\ from
breaking a line where it might otherwise do so.
This may be at a space, as between the ``Mr.''\ and
``Jones'' in
       ``Mr.~Jones'',        % ~ produces an unbreakable interword space.
or within a word---especially when the word is a
symbol like
       \mbox{\emph{itemnum}}
that makes little sense when hyphenated across
lines.

Footnotes\footnote{This is an example of a footnote.}
pose no problem.

\LaTeX\ is good at typesetting mathematical formulas
like
       \( x-3y + z = 7 \)
or
       \( a_{1} > x^{2n} + y^{2n} > x' \)
or
       \( \ip{A}{B} = \sum_{i} a_{i} b_{i} \).
The spaces you type in a formula are
ignored.  Remember that a letter like
       $x$                   % $ ... $  and  \( ... \)  are equivalent
is a formula when it denotes a mathematical
symbol, and it should be typed as one.

\section{Displayed Text}

Text is displayed by indenting it from the left
margin.  Quotations are commonly displayed.  There
are short quotations
\begin{quote}
   This is a short quotation.  It consists of a
   single paragraph of text.  See how it is formatted.
\end{quote}
and longer ones.
\begin{quotation}
   This is a longer quotation.  It consists of two
   paragraphs of text, neither of which are
   particularly interesting.

   This is the second paragraph of the quotation.  It
   is just as dull as the first paragraph.
\end{quotation}
Another frequently-displayed structure is a list.
The following is an example of an \emph{itemized}
list.
\begin{itemize}
   \item This is the first item of an itemized list.
         Each item in the list is marked with a ``tick''.
         You don't have to worry about what kind of tick
         mark is used.

   \item This is the second item of the list.  It
         contains another list nested inside it.  The inner
         list is an \emph{enumerated} list.
         \begin{enumerate}
            \item This is the first item of an enumerated
                  list that is nested within the itemized list.

            \item This is the second item of the inner list.
                  \LaTeX\ allows you to nest lists deeper than
                  you really should.
         \end{enumerate}
         This is the rest of the second item of the outer
         list.  It is no more interesting than any other
         part of the item.
   \item This is the third item of the list.
\end{itemize}
You can even display poetry.
\begin{verse}
   There is an environment
    for verse \\             % The \\ command separates lines
   Whose features some poets % within a stanza.
   will curse.

                             % One or more blank lines separate stanzas.

   For instead of making\\
   Them do \emph{all} line breaking, \\
   It allows them to put too many words on a line when they'd rather be
   forced to be terse.
\end{verse}

Mathematical formulas may also be displayed.  A
displayed formula
is
one-line long; multiline
formulas require special formatting instructions.
   \[  \ip{\Gamma}{\psi'} = x'' + y^{2} + z_{i}^{n}\]
Don't start a paragraph with a displayed equation,
nor make one a paragraph by itself.

\end{document}               % End of document.

\section{Bench Huge Section 1329}
% synthetic bench section 1329
% This is a sample LaTeX input file.  (Version of 12 August 2004.)
%
% A '%' character causes TeX to ignore all remaining text on the line,
% and is used for comments like this one.

\documentclass{article}      % Specifies the document class

                             % The preamble begins here.
\title{An Example Document}  % Declares the document's title.
\author{Leslie Lamport}      % Declares the author's name.
\date{January 21, 1994}      % Deleting this command produces today's date.

\newcommand{\ip}[2]{(#1, #2)}
                             % Defines \ip{arg1}{arg2} to mean
                             % (arg1, arg2).

%\newcommand{\ip}[2]{\langle #1 | #2\rangle}
                             % This is an alternative definition of
                             % \ip that is commented out.

\begin{document}             % End of preamble and beginning of text.

\maketitle                   % Produces the title.

This is an example input file.  Comparing it with
the output it generates can show you how to
produce a simple document of your own.

\section{Ordinary Text}      % Produces section heading.  Lower-level
                             % sections are begun with similar
                             % \subsection and \subsubsection commands.

The ends  of words and sentences are marked
  by   spaces. It  doesn't matter how many
spaces    you type; one is as good as 100.  The
end of   a line counts as a space.

One   or more   blank lines denote the  end
of  a paragraph.

Since any number of consecutive spaces are treated
like a single one, the formatting of the input
file makes no difference to
      \LaTeX,                % The \LaTeX command generates the LaTeX logo.
but it makes a difference to you.  When you use
\LaTeX, making your input file as easy to read
as possible will be a great help as you write
your document and when you change it.  This sample
file shows how you can add comments to your own input
file.

Because printing is different from typewriting,
there are a number of things that you have to do
differently when preparing an input file than if
you were just typing the document directly.
Quotation marks like
       ``this''
have to be handled specially, as do quotes within
quotes:
       ``\,`this'            % \, separates the double and single quote.
        is what I just
        wrote, not  `that'\,''.

Dashes come in three sizes: an
       intra-word
dash, a medium dash for number ranges like
       1--2,
and a punctuation
       dash---like
this.

A sentence-ending space should be larger than the
space between words within a sentence.  You
sometimes have to type special commands in
conjunction with punctuation characters to get
this right, as in the following sentence.
       Gnats, gnus, etc.\ all  % `\ ' makes an inter-word space.
       begin with G\@.         % \@ marks end-of-sentence punctuation.
You should check the spaces after periods when
reading your output to make sure you haven't
forgotten any special cases.  Generating an
ellipsis
       \ldots\               % `\ ' is needed after `\ldots' because TeX
                             % ignores spaces after command names like \ldots
                             % made from \ + letters.
                             %
                             % Note how a `%' character causes TeX to ignore
                             % the end of the input line, so these blank lines
                             % do not start a new paragraph.
                             %
with the right spacing around the periods requires
a special command.

\LaTeX\ interprets some common characters as
commands, so you must type special commands to
generate them.  These characters include the
following:
       \$ \& \% \# \{ and \}.

In printing, text is usually emphasized with an
       \emph{italic}
type style.

\begin{em}
   A long segment of text can also be emphasized
   in this way.  Text within such a segment can be
   given \emph{additional} emphasis.
\end{em}

It is sometimes necessary to prevent \LaTeX\ from
breaking a line where it might otherwise do so.
This may be at a space, as between the ``Mr.''\ and
``Jones'' in
       ``Mr.~Jones'',        % ~ produces an unbreakable interword space.
or within a word---especially when the word is a
symbol like
       \mbox{\emph{itemnum}}
that makes little sense when hyphenated across
lines.

Footnotes\footnote{This is an example of a footnote.}
pose no problem.

\LaTeX\ is good at typesetting mathematical formulas
like
       \( x-3y + z = 7 \)
or
       \( a_{1} > x^{2n} + y^{2n} > x' \)
or
       \( \ip{A}{B} = \sum_{i} a_{i} b_{i} \).
The spaces you type in a formula are
ignored.  Remember that a letter like
       $x$                   % $ ... $  and  \( ... \)  are equivalent
is a formula when it denotes a mathematical
symbol, and it should be typed as one.

\section{Displayed Text}

Text is displayed by indenting it from the left
margin.  Quotations are commonly displayed.  There
are short quotations
\begin{quote}
   This is a short quotation.  It consists of a
   single paragraph of text.  See how it is formatted.
\end{quote}
and longer ones.
\begin{quotation}
   This is a longer quotation.  It consists of two
   paragraphs of text, neither of which are
   particularly interesting.

   This is the second paragraph of the quotation.  It
   is just as dull as the first paragraph.
\end{quotation}
Another frequently-displayed structure is a list.
The following is an example of an \emph{itemized}
list.
\begin{itemize}
   \item This is the first item of an itemized list.
         Each item in the list is marked with a ``tick''.
         You don't have to worry about what kind of tick
         mark is used.

   \item This is the second item of the list.  It
         contains another list nested inside it.  The inner
         list is an \emph{enumerated} list.
         \begin{enumerate}
            \item This is the first item of an enumerated
                  list that is nested within the itemized list.

            \item This is the second item of the inner list.
                  \LaTeX\ allows you to nest lists deeper than
                  you really should.
         \end{enumerate}
         This is the rest of the second item of the outer
         list.  It is no more interesting than any other
         part of the item.
   \item This is the third item of the list.
\end{itemize}
You can even display poetry.
\begin{verse}
   There is an environment
    for verse \\             % The \\ command separates lines
   Whose features some poets % within a stanza.
   will curse.

                             % One or more blank lines separate stanzas.

   For instead of making\\
   Them do \emph{all} line breaking, \\
   It allows them to put too many words on a line when they'd rather be
   forced to be terse.
\end{verse}

Mathematical formulas may also be displayed.  A
displayed formula
is
one-line long; multiline
formulas require special formatting instructions.
   \[  \ip{\Gamma}{\psi'} = x'' + y^{2} + z_{i}^{n}\]
Don't start a paragraph with a displayed equation,
nor make one a paragraph by itself.

\end{document}               % End of document.

\section{Bench Huge Section 1330}
% synthetic bench section 1330
% This is a sample LaTeX input file.  (Version of 12 August 2004.)
%
% A '%' character causes TeX to ignore all remaining text on the line,
% and is used for comments like this one.

\documentclass{article}      % Specifies the document class

                             % The preamble begins here.
\title{An Example Document}  % Declares the document's title.
\author{Leslie Lamport}      % Declares the author's name.
\date{January 21, 1994}      % Deleting this command produces today's date.

\newcommand{\ip}[2]{(#1, #2)}
                             % Defines \ip{arg1}{arg2} to mean
                             % (arg1, arg2).

%\newcommand{\ip}[2]{\langle #1 | #2\rangle}
                             % This is an alternative definition of
                             % \ip that is commented out.

\begin{document}             % End of preamble and beginning of text.

\maketitle                   % Produces the title.

This is an example input file.  Comparing it with
the output it generates can show you how to
produce a simple document of your own.

\section{Ordinary Text}      % Produces section heading.  Lower-level
                             % sections are begun with similar
                             % \subsection and \subsubsection commands.

The ends  of words and sentences are marked
  by   spaces. It  doesn't matter how many
spaces    you type; one is as good as 100.  The
end of   a line counts as a space.

One   or more   blank lines denote the  end
of  a paragraph.

Since any number of consecutive spaces are treated
like a single one, the formatting of the input
file makes no difference to
      \LaTeX,                % The \LaTeX command generates the LaTeX logo.
but it makes a difference to you.  When you use
\LaTeX, making your input file as easy to read
as possible will be a great help as you write
your document and when you change it.  This sample
file shows how you can add comments to your own input
file.

Because printing is different from typewriting,
there are a number of things that you have to do
differently when preparing an input file than if
you were just typing the document directly.
Quotation marks like
       ``this''
have to be handled specially, as do quotes within
quotes:
       ``\,`this'            % \, separates the double and single quote.
        is what I just
        wrote, not  `that'\,''.

Dashes come in three sizes: an
       intra-word
dash, a medium dash for number ranges like
       1--2,
and a punctuation
       dash---like
this.

A sentence-ending space should be larger than the
space between words within a sentence.  You
sometimes have to type special commands in
conjunction with punctuation characters to get
this right, as in the following sentence.
       Gnats, gnus, etc.\ all  % `\ ' makes an inter-word space.
       begin with G\@.         % \@ marks end-of-sentence punctuation.
You should check the spaces after periods when
reading your output to make sure you haven't
forgotten any special cases.  Generating an
ellipsis
       \ldots\               % `\ ' is needed after `\ldots' because TeX
                             % ignores spaces after command names like \ldots
                             % made from \ + letters.
                             %
                             % Note how a `%' character causes TeX to ignore
                             % the end of the input line, so these blank lines
                             % do not start a new paragraph.
                             %
with the right spacing around the periods requires
a special command.

\LaTeX\ interprets some common characters as
commands, so you must type special commands to
generate them.  These characters include the
following:
       \$ \& \% \# \{ and \}.

In printing, text is usually emphasized with an
       \emph{italic}
type style.

\begin{em}
   A long segment of text can also be emphasized
   in this way.  Text within such a segment can be
   given \emph{additional} emphasis.
\end{em}

It is sometimes necessary to prevent \LaTeX\ from
breaking a line where it might otherwise do so.
This may be at a space, as between the ``Mr.''\ and
``Jones'' in
       ``Mr.~Jones'',        % ~ produces an unbreakable interword space.
or within a word---especially when the word is a
symbol like
       \mbox{\emph{itemnum}}
that makes little sense when hyphenated across
lines.

Footnotes\footnote{This is an example of a footnote.}
pose no problem.

\LaTeX\ is good at typesetting mathematical formulas
like
       \( x-3y + z = 7 \)
or
       \( a_{1} > x^{2n} + y^{2n} > x' \)
or
       \( \ip{A}{B} = \sum_{i} a_{i} b_{i} \).
The spaces you type in a formula are
ignored.  Remember that a letter like
       $x$                   % $ ... $  and  \( ... \)  are equivalent
is a formula when it denotes a mathematical
symbol, and it should be typed as one.

\section{Displayed Text}

Text is displayed by indenting it from the left
margin.  Quotations are commonly displayed.  There
are short quotations
\begin{quote}
   This is a short quotation.  It consists of a
   single paragraph of text.  See how it is formatted.
\end{quote}
and longer ones.
\begin{quotation}
   This is a longer quotation.  It consists of two
   paragraphs of text, neither of which are
   particularly interesting.

   This is the second paragraph of the quotation.  It
   is just as dull as the first paragraph.
\end{quotation}
Another frequently-displayed structure is a list.
The following is an example of an \emph{itemized}
list.
\begin{itemize}
   \item This is the first item of an itemized list.
         Each item in the list is marked with a ``tick''.
         You don't have to worry about what kind of tick
         mark is used.

   \item This is the second item of the list.  It
         contains another list nested inside it.  The inner
         list is an \emph{enumerated} list.
         \begin{enumerate}
            \item This is the first item of an enumerated
                  list that is nested within the itemized list.

            \item This is the second item of the inner list.
                  \LaTeX\ allows you to nest lists deeper than
                  you really should.
         \end{enumerate}
         This is the rest of the second item of the outer
         list.  It is no more interesting than any other
         part of the item.
   \item This is the third item of the list.
\end{itemize}
You can even display poetry.
\begin{verse}
   There is an environment
    for verse \\             % The \\ command separates lines
   Whose features some poets % within a stanza.
   will curse.

                             % One or more blank lines separate stanzas.

   For instead of making\\
   Them do \emph{all} line breaking, \\
   It allows them to put too many words on a line when they'd rather be
   forced to be terse.
\end{verse}

Mathematical formulas may also be displayed.  A
displayed formula
is
one-line long; multiline
formulas require special formatting instructions.
   \[  \ip{\Gamma}{\psi'} = x'' + y^{2} + z_{i}^{n}\]
Don't start a paragraph with a displayed equation,
nor make one a paragraph by itself.

\end{document}               % End of document.

\section{Bench Huge Section 1331}
% synthetic bench section 1331
% This is a sample LaTeX input file.  (Version of 12 August 2004.)
%
% A '%' character causes TeX to ignore all remaining text on the line,
% and is used for comments like this one.

\documentclass{article}      % Specifies the document class

                             % The preamble begins here.
\title{An Example Document}  % Declares the document's title.
\author{Leslie Lamport}      % Declares the author's name.
\date{January 21, 1994}      % Deleting this command produces today's date.

\newcommand{\ip}[2]{(#1, #2)}
                             % Defines \ip{arg1}{arg2} to mean
                             % (arg1, arg2).

%\newcommand{\ip}[2]{\langle #1 | #2\rangle}
                             % This is an alternative definition of
                             % \ip that is commented out.

\begin{document}             % End of preamble and beginning of text.

\maketitle                   % Produces the title.

This is an example input file.  Comparing it with
the output it generates can show you how to
produce a simple document of your own.

\section{Ordinary Text}      % Produces section heading.  Lower-level
                             % sections are begun with similar
                             % \subsection and \subsubsection commands.

The ends  of words and sentences are marked
  by   spaces. It  doesn't matter how many
spaces    you type; one is as good as 100.  The
end of   a line counts as a space.

One   or more   blank lines denote the  end
of  a paragraph.

Since any number of consecutive spaces are treated
like a single one, the formatting of the input
file makes no difference to
      \LaTeX,                % The \LaTeX command generates the LaTeX logo.
but it makes a difference to you.  When you use
\LaTeX, making your input file as easy to read
as possible will be a great help as you write
your document and when you change it.  This sample
file shows how you can add comments to your own input
file.

Because printing is different from typewriting,
there are a number of things that you have to do
differently when preparing an input file than if
you were just typing the document directly.
Quotation marks like
       ``this''
have to be handled specially, as do quotes within
quotes:
       ``\,`this'            % \, separates the double and single quote.
        is what I just
        wrote, not  `that'\,''.

Dashes come in three sizes: an
       intra-word
dash, a medium dash for number ranges like
       1--2,
and a punctuation
       dash---like
this.

A sentence-ending space should be larger than the
space between words within a sentence.  You
sometimes have to type special commands in
conjunction with punctuation characters to get
this right, as in the following sentence.
       Gnats, gnus, etc.\ all  % `\ ' makes an inter-word space.
       begin with G\@.         % \@ marks end-of-sentence punctuation.
You should check the spaces after periods when
reading your output to make sure you haven't
forgotten any special cases.  Generating an
ellipsis
       \ldots\               % `\ ' is needed after `\ldots' because TeX
                             % ignores spaces after command names like \ldots
                             % made from \ + letters.
                             %
                             % Note how a `%' character causes TeX to ignore
                             % the end of the input line, so these blank lines
                             % do not start a new paragraph.
                             %
with the right spacing around the periods requires
a special command.

\LaTeX\ interprets some common characters as
commands, so you must type special commands to
generate them.  These characters include the
following:
       \$ \& \% \# \{ and \}.

In printing, text is usually emphasized with an
       \emph{italic}
type style.

\begin{em}
   A long segment of text can also be emphasized
   in this way.  Text within such a segment can be
   given \emph{additional} emphasis.
\end{em}

It is sometimes necessary to prevent \LaTeX\ from
breaking a line where it might otherwise do so.
This may be at a space, as between the ``Mr.''\ and
``Jones'' in
       ``Mr.~Jones'',        % ~ produces an unbreakable interword space.
or within a word---especially when the word is a
symbol like
       \mbox{\emph{itemnum}}
that makes little sense when hyphenated across
lines.

Footnotes\footnote{This is an example of a footnote.}
pose no problem.

\LaTeX\ is good at typesetting mathematical formulas
like
       \( x-3y + z = 7 \)
or
       \( a_{1} > x^{2n} + y^{2n} > x' \)
or
       \( \ip{A}{B} = \sum_{i} a_{i} b_{i} \).
The spaces you type in a formula are
ignored.  Remember that a letter like
       $x$                   % $ ... $  and  \( ... \)  are equivalent
is a formula when it denotes a mathematical
symbol, and it should be typed as one.

\section{Displayed Text}

Text is displayed by indenting it from the left
margin.  Quotations are commonly displayed.  There
are short quotations
\begin{quote}
   This is a short quotation.  It consists of a
   single paragraph of text.  See how it is formatted.
\end{quote}
and longer ones.
\begin{quotation}
   This is a longer quotation.  It consists of two
   paragraphs of text, neither of which are
   particularly interesting.

   This is the second paragraph of the quotation.  It
   is just as dull as the first paragraph.
\end{quotation}
Another frequently-displayed structure is a list.
The following is an example of an \emph{itemized}
list.
\begin{itemize}
   \item This is the first item of an itemized list.
         Each item in the list is marked with a ``tick''.
         You don't have to worry about what kind of tick
         mark is used.

   \item This is the second item of the list.  It
         contains another list nested inside it.  The inner
         list is an \emph{enumerated} list.
         \begin{enumerate}
            \item This is the first item of an enumerated
                  list that is nested within the itemized list.

            \item This is the second item of the inner list.
                  \LaTeX\ allows you to nest lists deeper than
                  you really should.
         \end{enumerate}
         This is the rest of the second item of the outer
         list.  It is no more interesting than any other
         part of the item.
   \item This is the third item of the list.
\end{itemize}
You can even display poetry.
\begin{verse}
   There is an environment
    for verse \\             % The \\ command separates lines
   Whose features some poets % within a stanza.
   will curse.

                             % One or more blank lines separate stanzas.

   For instead of making\\
   Them do \emph{all} line breaking, \\
   It allows them to put too many words on a line when they'd rather be
   forced to be terse.
\end{verse}

Mathematical formulas may also be displayed.  A
displayed formula
is
one-line long; multiline
formulas require special formatting instructions.
   \[  \ip{\Gamma}{\psi'} = x'' + y^{2} + z_{i}^{n}\]
Don't start a paragraph with a displayed equation,
nor make one a paragraph by itself.

\end{document}               % End of document.

\section{Bench Huge Section 1332}
% synthetic bench section 1332
% This is a sample LaTeX input file.  (Version of 12 August 2004.)
%
% A '%' character causes TeX to ignore all remaining text on the line,
% and is used for comments like this one.

\documentclass{article}      % Specifies the document class

                             % The preamble begins here.
\title{An Example Document}  % Declares the document's title.
\author{Leslie Lamport}      % Declares the author's name.
\date{January 21, 1994}      % Deleting this command produces today's date.

\newcommand{\ip}[2]{(#1, #2)}
                             % Defines \ip{arg1}{arg2} to mean
                             % (arg1, arg2).

%\newcommand{\ip}[2]{\langle #1 | #2\rangle}
                             % This is an alternative definition of
                             % \ip that is commented out.

\begin{document}             % End of preamble and beginning of text.

\maketitle                   % Produces the title.

This is an example input file.  Comparing it with
the output it generates can show you how to
produce a simple document of your own.

\section{Ordinary Text}      % Produces section heading.  Lower-level
                             % sections are begun with similar
                             % \subsection and \subsubsection commands.

The ends  of words and sentences are marked
  by   spaces. It  doesn't matter how many
spaces    you type; one is as good as 100.  The
end of   a line counts as a space.

One   or more   blank lines denote the  end
of  a paragraph.

Since any number of consecutive spaces are treated
like a single one, the formatting of the input
file makes no difference to
      \LaTeX,                % The \LaTeX command generates the LaTeX logo.
but it makes a difference to you.  When you use
\LaTeX, making your input file as easy to read
as possible will be a great help as you write
your document and when you change it.  This sample
file shows how you can add comments to your own input
file.

Because printing is different from typewriting,
there are a number of things that you have to do
differently when preparing an input file than if
you were just typing the document directly.
Quotation marks like
       ``this''
have to be handled specially, as do quotes within
quotes:
       ``\,`this'            % \, separates the double and single quote.
        is what I just
        wrote, not  `that'\,''.

Dashes come in three sizes: an
       intra-word
dash, a medium dash for number ranges like
       1--2,
and a punctuation
       dash---like
this.

A sentence-ending space should be larger than the
space between words within a sentence.  You
sometimes have to type special commands in
conjunction with punctuation characters to get
this right, as in the following sentence.
       Gnats, gnus, etc.\ all  % `\ ' makes an inter-word space.
       begin with G\@.         % \@ marks end-of-sentence punctuation.
You should check the spaces after periods when
reading your output to make sure you haven't
forgotten any special cases.  Generating an
ellipsis
       \ldots\               % `\ ' is needed after `\ldots' because TeX
                             % ignores spaces after command names like \ldots
                             % made from \ + letters.
                             %
                             % Note how a `%' character causes TeX to ignore
                             % the end of the input line, so these blank lines
                             % do not start a new paragraph.
                             %
with the right spacing around the periods requires
a special command.

\LaTeX\ interprets some common characters as
commands, so you must type special commands to
generate them.  These characters include the
following:
       \$ \& \% \# \{ and \}.

In printing, text is usually emphasized with an
       \emph{italic}
type style.

\begin{em}
   A long segment of text can also be emphasized
   in this way.  Text within such a segment can be
   given \emph{additional} emphasis.
\end{em}

It is sometimes necessary to prevent \LaTeX\ from
breaking a line where it might otherwise do so.
This may be at a space, as between the ``Mr.''\ and
``Jones'' in
       ``Mr.~Jones'',        % ~ produces an unbreakable interword space.
or within a word---especially when the word is a
symbol like
       \mbox{\emph{itemnum}}
that makes little sense when hyphenated across
lines.

Footnotes\footnote{This is an example of a footnote.}
pose no problem.

\LaTeX\ is good at typesetting mathematical formulas
like
       \( x-3y + z = 7 \)
or
       \( a_{1} > x^{2n} + y^{2n} > x' \)
or
       \( \ip{A}{B} = \sum_{i} a_{i} b_{i} \).
The spaces you type in a formula are
ignored.  Remember that a letter like
       $x$                   % $ ... $  and  \( ... \)  are equivalent
is a formula when it denotes a mathematical
symbol, and it should be typed as one.

\section{Displayed Text}

Text is displayed by indenting it from the left
margin.  Quotations are commonly displayed.  There
are short quotations
\begin{quote}
   This is a short quotation.  It consists of a
   single paragraph of text.  See how it is formatted.
\end{quote}
and longer ones.
\begin{quotation}
   This is a longer quotation.  It consists of two
   paragraphs of text, neither of which are
   particularly interesting.

   This is the second paragraph of the quotation.  It
   is just as dull as the first paragraph.
\end{quotation}
Another frequently-displayed structure is a list.
The following is an example of an \emph{itemized}
list.
\begin{itemize}
   \item This is the first item of an itemized list.
         Each item in the list is marked with a ``tick''.
         You don't have to worry about what kind of tick
         mark is used.

   \item This is the second item of the list.  It
         contains another list nested inside it.  The inner
         list is an \emph{enumerated} list.
         \begin{enumerate}
            \item This is the first item of an enumerated
                  list that is nested within the itemized list.

            \item This is the second item of the inner list.
                  \LaTeX\ allows you to nest lists deeper than
                  you really should.
         \end{enumerate}
         This is the rest of the second item of the outer
         list.  It is no more interesting than any other
         part of the item.
   \item This is the third item of the list.
\end{itemize}
You can even display poetry.
\begin{verse}
   There is an environment
    for verse \\             % The \\ command separates lines
   Whose features some poets % within a stanza.
   will curse.

                             % One or more blank lines separate stanzas.

   For instead of making\\
   Them do \emph{all} line breaking, \\
   It allows them to put too many words on a line when they'd rather be
   forced to be terse.
\end{verse}

Mathematical formulas may also be displayed.  A
displayed formula
is
one-line long; multiline
formulas require special formatting instructions.
   \[  \ip{\Gamma}{\psi'} = x'' + y^{2} + z_{i}^{n}\]
Don't start a paragraph with a displayed equation,
nor make one a paragraph by itself.

\end{document}               % End of document.

\section{Bench Huge Section 1333}
% synthetic bench section 1333
% This is a sample LaTeX input file.  (Version of 12 August 2004.)
%
% A '%' character causes TeX to ignore all remaining text on the line,
% and is used for comments like this one.

\documentclass{article}      % Specifies the document class

                             % The preamble begins here.
\title{An Example Document}  % Declares the document's title.
\author{Leslie Lamport}      % Declares the author's name.
\date{January 21, 1994}      % Deleting this command produces today's date.

\newcommand{\ip}[2]{(#1, #2)}
                             % Defines \ip{arg1}{arg2} to mean
                             % (arg1, arg2).

%\newcommand{\ip}[2]{\langle #1 | #2\rangle}
                             % This is an alternative definition of
                             % \ip that is commented out.

\begin{document}             % End of preamble and beginning of text.

\maketitle                   % Produces the title.

This is an example input file.  Comparing it with
the output it generates can show you how to
produce a simple document of your own.

\section{Ordinary Text}      % Produces section heading.  Lower-level
                             % sections are begun with similar
                             % \subsection and \subsubsection commands.

The ends  of words and sentences are marked
  by   spaces. It  doesn't matter how many
spaces    you type; one is as good as 100.  The
end of   a line counts as a space.

One   or more   blank lines denote the  end
of  a paragraph.

Since any number of consecutive spaces are treated
like a single one, the formatting of the input
file makes no difference to
      \LaTeX,                % The \LaTeX command generates the LaTeX logo.
but it makes a difference to you.  When you use
\LaTeX, making your input file as easy to read
as possible will be a great help as you write
your document and when you change it.  This sample
file shows how you can add comments to your own input
file.

Because printing is different from typewriting,
there are a number of things that you have to do
differently when preparing an input file than if
you were just typing the document directly.
Quotation marks like
       ``this''
have to be handled specially, as do quotes within
quotes:
       ``\,`this'            % \, separates the double and single quote.
        is what I just
        wrote, not  `that'\,''.

Dashes come in three sizes: an
       intra-word
dash, a medium dash for number ranges like
       1--2,
and a punctuation
       dash---like
this.

A sentence-ending space should be larger than the
space between words within a sentence.  You
sometimes have to type special commands in
conjunction with punctuation characters to get
this right, as in the following sentence.
       Gnats, gnus, etc.\ all  % `\ ' makes an inter-word space.
       begin with G\@.         % \@ marks end-of-sentence punctuation.
You should check the spaces after periods when
reading your output to make sure you haven't
forgotten any special cases.  Generating an
ellipsis
       \ldots\               % `\ ' is needed after `\ldots' because TeX
                             % ignores spaces after command names like \ldots
                             % made from \ + letters.
                             %
                             % Note how a `%' character causes TeX to ignore
                             % the end of the input line, so these blank lines
                             % do not start a new paragraph.
                             %
with the right spacing around the periods requires
a special command.

\LaTeX\ interprets some common characters as
commands, so you must type special commands to
generate them.  These characters include the
following:
       \$ \& \% \# \{ and \}.

In printing, text is usually emphasized with an
       \emph{italic}
type style.

\begin{em}
   A long segment of text can also be emphasized
   in this way.  Text within such a segment can be
   given \emph{additional} emphasis.
\end{em}

It is sometimes necessary to prevent \LaTeX\ from
breaking a line where it might otherwise do so.
This may be at a space, as between the ``Mr.''\ and
``Jones'' in
       ``Mr.~Jones'',        % ~ produces an unbreakable interword space.
or within a word---especially when the word is a
symbol like
       \mbox{\emph{itemnum}}
that makes little sense when hyphenated across
lines.

Footnotes\footnote{This is an example of a footnote.}
pose no problem.

\LaTeX\ is good at typesetting mathematical formulas
like
       \( x-3y + z = 7 \)
or
       \( a_{1} > x^{2n} + y^{2n} > x' \)
or
       \( \ip{A}{B} = \sum_{i} a_{i} b_{i} \).
The spaces you type in a formula are
ignored.  Remember that a letter like
       $x$                   % $ ... $  and  \( ... \)  are equivalent
is a formula when it denotes a mathematical
symbol, and it should be typed as one.

\section{Displayed Text}

Text is displayed by indenting it from the left
margin.  Quotations are commonly displayed.  There
are short quotations
\begin{quote}
   This is a short quotation.  It consists of a
   single paragraph of text.  See how it is formatted.
\end{quote}
and longer ones.
\begin{quotation}
   This is a longer quotation.  It consists of two
   paragraphs of text, neither of which are
   particularly interesting.

   This is the second paragraph of the quotation.  It
   is just as dull as the first paragraph.
\end{quotation}
Another frequently-displayed structure is a list.
The following is an example of an \emph{itemized}
list.
\begin{itemize}
   \item This is the first item of an itemized list.
         Each item in the list is marked with a ``tick''.
         You don't have to worry about what kind of tick
         mark is used.

   \item This is the second item of the list.  It
         contains another list nested inside it.  The inner
         list is an \emph{enumerated} list.
         \begin{enumerate}
            \item This is the first item of an enumerated
                  list that is nested within the itemized list.

            \item This is the second item of the inner list.
                  \LaTeX\ allows you to nest lists deeper than
                  you really should.
         \end{enumerate}
         This is the rest of the second item of the outer
         list.  It is no more interesting than any other
         part of the item.
   \item This is the third item of the list.
\end{itemize}
You can even display poetry.
\begin{verse}
   There is an environment
    for verse \\             % The \\ command separates lines
   Whose features some poets % within a stanza.
   will curse.

                             % One or more blank lines separate stanzas.

   For instead of making\\
   Them do \emph{all} line breaking, \\
   It allows them to put too many words on a line when they'd rather be
   forced to be terse.
\end{verse}

Mathematical formulas may also be displayed.  A
displayed formula
is
one-line long; multiline
formulas require special formatting instructions.
   \[  \ip{\Gamma}{\psi'} = x'' + y^{2} + z_{i}^{n}\]
Don't start a paragraph with a displayed equation,
nor make one a paragraph by itself.

\end{document}               % End of document.

\section{Bench Huge Section 1334}
% synthetic bench section 1334
% This is a sample LaTeX input file.  (Version of 12 August 2004.)
%
% A '%' character causes TeX to ignore all remaining text on the line,
% and is used for comments like this one.

\documentclass{article}      % Specifies the document class

                             % The preamble begins here.
\title{An Example Document}  % Declares the document's title.
\author{Leslie Lamport}      % Declares the author's name.
\date{January 21, 1994}      % Deleting this command produces today's date.

\newcommand{\ip}[2]{(#1, #2)}
                             % Defines \ip{arg1}{arg2} to mean
                             % (arg1, arg2).

%\newcommand{\ip}[2]{\langle #1 | #2\rangle}
                             % This is an alternative definition of
                             % \ip that is commented out.

\begin{document}             % End of preamble and beginning of text.

\maketitle                   % Produces the title.

This is an example input file.  Comparing it with
the output it generates can show you how to
produce a simple document of your own.

\section{Ordinary Text}      % Produces section heading.  Lower-level
                             % sections are begun with similar
                             % \subsection and \subsubsection commands.

The ends  of words and sentences are marked
  by   spaces. It  doesn't matter how many
spaces    you type; one is as good as 100.  The
end of   a line counts as a space.

One   or more   blank lines denote the  end
of  a paragraph.

Since any number of consecutive spaces are treated
like a single one, the formatting of the input
file makes no difference to
      \LaTeX,                % The \LaTeX command generates the LaTeX logo.
but it makes a difference to you.  When you use
\LaTeX, making your input file as easy to read
as possible will be a great help as you write
your document and when you change it.  This sample
file shows how you can add comments to your own input
file.

Because printing is different from typewriting,
there are a number of things that you have to do
differently when preparing an input file than if
you were just typing the document directly.
Quotation marks like
       ``this''
have to be handled specially, as do quotes within
quotes:
       ``\,`this'            % \, separates the double and single quote.
        is what I just
        wrote, not  `that'\,''.

Dashes come in three sizes: an
       intra-word
dash, a medium dash for number ranges like
       1--2,
and a punctuation
       dash---like
this.

A sentence-ending space should be larger than the
space between words within a sentence.  You
sometimes have to type special commands in
conjunction with punctuation characters to get
this right, as in the following sentence.
       Gnats, gnus, etc.\ all  % `\ ' makes an inter-word space.
       begin with G\@.         % \@ marks end-of-sentence punctuation.
You should check the spaces after periods when
reading your output to make sure you haven't
forgotten any special cases.  Generating an
ellipsis
       \ldots\               % `\ ' is needed after `\ldots' because TeX
                             % ignores spaces after command names like \ldots
                             % made from \ + letters.
                             %
                             % Note how a `%' character causes TeX to ignore
                             % the end of the input line, so these blank lines
                             % do not start a new paragraph.
                             %
with the right spacing around the periods requires
a special command.

\LaTeX\ interprets some common characters as
commands, so you must type special commands to
generate them.  These characters include the
following:
       \$ \& \% \# \{ and \}.

In printing, text is usually emphasized with an
       \emph{italic}
type style.

\begin{em}
   A long segment of text can also be emphasized
   in this way.  Text within such a segment can be
   given \emph{additional} emphasis.
\end{em}

It is sometimes necessary to prevent \LaTeX\ from
breaking a line where it might otherwise do so.
This may be at a space, as between the ``Mr.''\ and
``Jones'' in
       ``Mr.~Jones'',        % ~ produces an unbreakable interword space.
or within a word---especially when the word is a
symbol like
       \mbox{\emph{itemnum}}
that makes little sense when hyphenated across
lines.

Footnotes\footnote{This is an example of a footnote.}
pose no problem.

\LaTeX\ is good at typesetting mathematical formulas
like
       \( x-3y + z = 7 \)
or
       \( a_{1} > x^{2n} + y^{2n} > x' \)
or
       \( \ip{A}{B} = \sum_{i} a_{i} b_{i} \).
The spaces you type in a formula are
ignored.  Remember that a letter like
       $x$                   % $ ... $  and  \( ... \)  are equivalent
is a formula when it denotes a mathematical
symbol, and it should be typed as one.

\section{Displayed Text}

Text is displayed by indenting it from the left
margin.  Quotations are commonly displayed.  There
are short quotations
\begin{quote}
   This is a short quotation.  It consists of a
   single paragraph of text.  See how it is formatted.
\end{quote}
and longer ones.
\begin{quotation}
   This is a longer quotation.  It consists of two
   paragraphs of text, neither of which are
   particularly interesting.

   This is the second paragraph of the quotation.  It
   is just as dull as the first paragraph.
\end{quotation}
Another frequently-displayed structure is a list.
The following is an example of an \emph{itemized}
list.
\begin{itemize}
   \item This is the first item of an itemized list.
         Each item in the list is marked with a ``tick''.
         You don't have to worry about what kind of tick
         mark is used.

   \item This is the second item of the list.  It
         contains another list nested inside it.  The inner
         list is an \emph{enumerated} list.
         \begin{enumerate}
            \item This is the first item of an enumerated
                  list that is nested within the itemized list.

            \item This is the second item of the inner list.
                  \LaTeX\ allows you to nest lists deeper than
                  you really should.
         \end{enumerate}
         This is the rest of the second item of the outer
         list.  It is no more interesting than any other
         part of the item.
   \item This is the third item of the list.
\end{itemize}
You can even display poetry.
\begin{verse}
   There is an environment
    for verse \\             % The \\ command separates lines
   Whose features some poets % within a stanza.
   will curse.

                             % One or more blank lines separate stanzas.

   For instead of making\\
   Them do \emph{all} line breaking, \\
   It allows them to put too many words on a line when they'd rather be
   forced to be terse.
\end{verse}

Mathematical formulas may also be displayed.  A
displayed formula
is
one-line long; multiline
formulas require special formatting instructions.
   \[  \ip{\Gamma}{\psi'} = x'' + y^{2} + z_{i}^{n}\]
Don't start a paragraph with a displayed equation,
nor make one a paragraph by itself.

\end{document}               % End of document.

\section{Bench Huge Section 1335}
% synthetic bench section 1335
% This is a sample LaTeX input file.  (Version of 12 August 2004.)
%
% A '%' character causes TeX to ignore all remaining text on the line,
% and is used for comments like this one.

\documentclass{article}      % Specifies the document class

                             % The preamble begins here.
\title{An Example Document}  % Declares the document's title.
\author{Leslie Lamport}      % Declares the author's name.
\date{January 21, 1994}      % Deleting this command produces today's date.

\newcommand{\ip}[2]{(#1, #2)}
                             % Defines \ip{arg1}{arg2} to mean
                             % (arg1, arg2).

%\newcommand{\ip}[2]{\langle #1 | #2\rangle}
                             % This is an alternative definition of
                             % \ip that is commented out.

\begin{document}             % End of preamble and beginning of text.

\maketitle                   % Produces the title.

This is an example input file.  Comparing it with
the output it generates can show you how to
produce a simple document of your own.

\section{Ordinary Text}      % Produces section heading.  Lower-level
                             % sections are begun with similar
                             % \subsection and \subsubsection commands.

The ends  of words and sentences are marked
  by   spaces. It  doesn't matter how many
spaces    you type; one is as good as 100.  The
end of   a line counts as a space.

One   or more   blank lines denote the  end
of  a paragraph.

Since any number of consecutive spaces are treated
like a single one, the formatting of the input
file makes no difference to
      \LaTeX,                % The \LaTeX command generates the LaTeX logo.
but it makes a difference to you.  When you use
\LaTeX, making your input file as easy to read
as possible will be a great help as you write
your document and when you change it.  This sample
file shows how you can add comments to your own input
file.

Because printing is different from typewriting,
there are a number of things that you have to do
differently when preparing an input file than if
you were just typing the document directly.
Quotation marks like
       ``this''
have to be handled specially, as do quotes within
quotes:
       ``\,`this'            % \, separates the double and single quote.
        is what I just
        wrote, not  `that'\,''.

Dashes come in three sizes: an
       intra-word
dash, a medium dash for number ranges like
       1--2,
and a punctuation
       dash---like
this.

A sentence-ending space should be larger than the
space between words within a sentence.  You
sometimes have to type special commands in
conjunction with punctuation characters to get
this right, as in the following sentence.
       Gnats, gnus, etc.\ all  % `\ ' makes an inter-word space.
       begin with G\@.         % \@ marks end-of-sentence punctuation.
You should check the spaces after periods when
reading your output to make sure you haven't
forgotten any special cases.  Generating an
ellipsis
       \ldots\               % `\ ' is needed after `\ldots' because TeX
                             % ignores spaces after command names like \ldots
                             % made from \ + letters.
                             %
                             % Note how a `%' character causes TeX to ignore
                             % the end of the input line, so these blank lines
                             % do not start a new paragraph.
                             %
with the right spacing around the periods requires
a special command.

\LaTeX\ interprets some common characters as
commands, so you must type special commands to
generate them.  These characters include the
following:
       \$ \& \% \# \{ and \}.

In printing, text is usually emphasized with an
       \emph{italic}
type style.

\begin{em}
   A long segment of text can also be emphasized
   in this way.  Text within such a segment can be
   given \emph{additional} emphasis.
\end{em}

It is sometimes necessary to prevent \LaTeX\ from
breaking a line where it might otherwise do so.
This may be at a space, as between the ``Mr.''\ and
``Jones'' in
       ``Mr.~Jones'',        % ~ produces an unbreakable interword space.
or within a word---especially when the word is a
symbol like
       \mbox{\emph{itemnum}}
that makes little sense when hyphenated across
lines.

Footnotes\footnote{This is an example of a footnote.}
pose no problem.

\LaTeX\ is good at typesetting mathematical formulas
like
       \( x-3y + z = 7 \)
or
       \( a_{1} > x^{2n} + y^{2n} > x' \)
or
       \( \ip{A}{B} = \sum_{i} a_{i} b_{i} \).
The spaces you type in a formula are
ignored.  Remember that a letter like
       $x$                   % $ ... $  and  \( ... \)  are equivalent
is a formula when it denotes a mathematical
symbol, and it should be typed as one.

\section{Displayed Text}

Text is displayed by indenting it from the left
margin.  Quotations are commonly displayed.  There
are short quotations
\begin{quote}
   This is a short quotation.  It consists of a
   single paragraph of text.  See how it is formatted.
\end{quote}
and longer ones.
\begin{quotation}
   This is a longer quotation.  It consists of two
   paragraphs of text, neither of which are
   particularly interesting.

   This is the second paragraph of the quotation.  It
   is just as dull as the first paragraph.
\end{quotation}
Another frequently-displayed structure is a list.
The following is an example of an \emph{itemized}
list.
\begin{itemize}
   \item This is the first item of an itemized list.
         Each item in the list is marked with a ``tick''.
         You don't have to worry about what kind of tick
         mark is used.

   \item This is the second item of the list.  It
         contains another list nested inside it.  The inner
         list is an \emph{enumerated} list.
         \begin{enumerate}
            \item This is the first item of an enumerated
                  list that is nested within the itemized list.

            \item This is the second item of the inner list.
                  \LaTeX\ allows you to nest lists deeper than
                  you really should.
         \end{enumerate}
         This is the rest of the second item of the outer
         list.  It is no more interesting than any other
         part of the item.
   \item This is the third item of the list.
\end{itemize}
You can even display poetry.
\begin{verse}
   There is an environment
    for verse \\             % The \\ command separates lines
   Whose features some poets % within a stanza.
   will curse.

                             % One or more blank lines separate stanzas.

   For instead of making\\
   Them do \emph{all} line breaking, \\
   It allows them to put too many words on a line when they'd rather be
   forced to be terse.
\end{verse}

Mathematical formulas may also be displayed.  A
displayed formula
is
one-line long; multiline
formulas require special formatting instructions.
   \[  \ip{\Gamma}{\psi'} = x'' + y^{2} + z_{i}^{n}\]
Don't start a paragraph with a displayed equation,
nor make one a paragraph by itself.

\end{document}               % End of document.

\section{Bench Huge Section 1336}
% synthetic bench section 1336
% This is a sample LaTeX input file.  (Version of 12 August 2004.)
%
% A '%' character causes TeX to ignore all remaining text on the line,
% and is used for comments like this one.

\documentclass{article}      % Specifies the document class

                             % The preamble begins here.
\title{An Example Document}  % Declares the document's title.
\author{Leslie Lamport}      % Declares the author's name.
\date{January 21, 1994}      % Deleting this command produces today's date.

\newcommand{\ip}[2]{(#1, #2)}
                             % Defines \ip{arg1}{arg2} to mean
                             % (arg1, arg2).

%\newcommand{\ip}[2]{\langle #1 | #2\rangle}
                             % This is an alternative definition of
                             % \ip that is commented out.

\begin{document}             % End of preamble and beginning of text.

\maketitle                   % Produces the title.

This is an example input file.  Comparing it with
the output it generates can show you how to
produce a simple document of your own.

\section{Ordinary Text}      % Produces section heading.  Lower-level
                             % sections are begun with similar
                             % \subsection and \subsubsection commands.

The ends  of words and sentences are marked
  by   spaces. It  doesn't matter how many
spaces    you type; one is as good as 100.  The
end of   a line counts as a space.

One   or more   blank lines denote the  end
of  a paragraph.

Since any number of consecutive spaces are treated
like a single one, the formatting of the input
file makes no difference to
      \LaTeX,                % The \LaTeX command generates the LaTeX logo.
but it makes a difference to you.  When you use
\LaTeX, making your input file as easy to read
as possible will be a great help as you write
your document and when you change it.  This sample
file shows how you can add comments to your own input
file.

Because printing is different from typewriting,
there are a number of things that you have to do
differently when preparing an input file than if
you were just typing the document directly.
Quotation marks like
       ``this''
have to be handled specially, as do quotes within
quotes:
       ``\,`this'            % \, separates the double and single quote.
        is what I just
        wrote, not  `that'\,''.

Dashes come in three sizes: an
       intra-word
dash, a medium dash for number ranges like
       1--2,
and a punctuation
       dash---like
this.

A sentence-ending space should be larger than the
space between words within a sentence.  You
sometimes have to type special commands in
conjunction with punctuation characters to get
this right, as in the following sentence.
       Gnats, gnus, etc.\ all  % `\ ' makes an inter-word space.
       begin with G\@.         % \@ marks end-of-sentence punctuation.
You should check the spaces after periods when
reading your output to make sure you haven't
forgotten any special cases.  Generating an
ellipsis
       \ldots\               % `\ ' is needed after `\ldots' because TeX
                             % ignores spaces after command names like \ldots
                             % made from \ + letters.
                             %
                             % Note how a `%' character causes TeX to ignore
                             % the end of the input line, so these blank lines
                             % do not start a new paragraph.
                             %
with the right spacing around the periods requires
a special command.

\LaTeX\ interprets some common characters as
commands, so you must type special commands to
generate them.  These characters include the
following:
       \$ \& \% \# \{ and \}.

In printing, text is usually emphasized with an
       \emph{italic}
type style.

\begin{em}
   A long segment of text can also be emphasized
   in this way.  Text within such a segment can be
   given \emph{additional} emphasis.
\end{em}

It is sometimes necessary to prevent \LaTeX\ from
breaking a line where it might otherwise do so.
This may be at a space, as between the ``Mr.''\ and
``Jones'' in
       ``Mr.~Jones'',        % ~ produces an unbreakable interword space.
or within a word---especially when the word is a
symbol like
       \mbox{\emph{itemnum}}
that makes little sense when hyphenated across
lines.

Footnotes\footnote{This is an example of a footnote.}
pose no problem.

\LaTeX\ is good at typesetting mathematical formulas
like
       \( x-3y + z = 7 \)
or
       \( a_{1} > x^{2n} + y^{2n} > x' \)
or
       \( \ip{A}{B} = \sum_{i} a_{i} b_{i} \).
The spaces you type in a formula are
ignored.  Remember that a letter like
       $x$                   % $ ... $  and  \( ... \)  are equivalent
is a formula when it denotes a mathematical
symbol, and it should be typed as one.

\section{Displayed Text}

Text is displayed by indenting it from the left
margin.  Quotations are commonly displayed.  There
are short quotations
\begin{quote}
   This is a short quotation.  It consists of a
   single paragraph of text.  See how it is formatted.
\end{quote}
and longer ones.
\begin{quotation}
   This is a longer quotation.  It consists of two
   paragraphs of text, neither of which are
   particularly interesting.

   This is the second paragraph of the quotation.  It
   is just as dull as the first paragraph.
\end{quotation}
Another frequently-displayed structure is a list.
The following is an example of an \emph{itemized}
list.
\begin{itemize}
   \item This is the first item of an itemized list.
         Each item in the list is marked with a ``tick''.
         You don't have to worry about what kind of tick
         mark is used.

   \item This is the second item of the list.  It
         contains another list nested inside it.  The inner
         list is an \emph{enumerated} list.
         \begin{enumerate}
            \item This is the first item of an enumerated
                  list that is nested within the itemized list.

            \item This is the second item of the inner list.
                  \LaTeX\ allows you to nest lists deeper than
                  you really should.
         \end{enumerate}
         This is the rest of the second item of the outer
         list.  It is no more interesting than any other
         part of the item.
   \item This is the third item of the list.
\end{itemize}
You can even display poetry.
\begin{verse}
   There is an environment
    for verse \\             % The \\ command separates lines
   Whose features some poets % within a stanza.
   will curse.

                             % One or more blank lines separate stanzas.

   For instead of making\\
   Them do \emph{all} line breaking, \\
   It allows them to put too many words on a line when they'd rather be
   forced to be terse.
\end{verse}

Mathematical formulas may also be displayed.  A
displayed formula
is
one-line long; multiline
formulas require special formatting instructions.
   \[  \ip{\Gamma}{\psi'} = x'' + y^{2} + z_{i}^{n}\]
Don't start a paragraph with a displayed equation,
nor make one a paragraph by itself.

\end{document}               % End of document.

\section{Bench Huge Section 1337}
% synthetic bench section 1337
% This is a sample LaTeX input file.  (Version of 12 August 2004.)
%
% A '%' character causes TeX to ignore all remaining text on the line,
% and is used for comments like this one.

\documentclass{article}      % Specifies the document class

                             % The preamble begins here.
\title{An Example Document}  % Declares the document's title.
\author{Leslie Lamport}      % Declares the author's name.
\date{January 21, 1994}      % Deleting this command produces today's date.

\newcommand{\ip}[2]{(#1, #2)}
                             % Defines \ip{arg1}{arg2} to mean
                             % (arg1, arg2).

%\newcommand{\ip}[2]{\langle #1 | #2\rangle}
                             % This is an alternative definition of
                             % \ip that is commented out.

\begin{document}             % End of preamble and beginning of text.

\maketitle                   % Produces the title.

This is an example input file.  Comparing it with
the output it generates can show you how to
produce a simple document of your own.

\section{Ordinary Text}      % Produces section heading.  Lower-level
                             % sections are begun with similar
                             % \subsection and \subsubsection commands.

The ends  of words and sentences are marked
  by   spaces. It  doesn't matter how many
spaces    you type; one is as good as 100.  The
end of   a line counts as a space.

One   or more   blank lines denote the  end
of  a paragraph.

Since any number of consecutive spaces are treated
like a single one, the formatting of the input
file makes no difference to
      \LaTeX,                % The \LaTeX command generates the LaTeX logo.
but it makes a difference to you.  When you use
\LaTeX, making your input file as easy to read
as possible will be a great help as you write
your document and when you change it.  This sample
file shows how you can add comments to your own input
file.

Because printing is different from typewriting,
there are a number of things that you have to do
differently when preparing an input file than if
you were just typing the document directly.
Quotation marks like
       ``this''
have to be handled specially, as do quotes within
quotes:
       ``\,`this'            % \, separates the double and single quote.
        is what I just
        wrote, not  `that'\,''.

Dashes come in three sizes: an
       intra-word
dash, a medium dash for number ranges like
       1--2,
and a punctuation
       dash---like
this.

A sentence-ending space should be larger than the
space between words within a sentence.  You
sometimes have to type special commands in
conjunction with punctuation characters to get
this right, as in the following sentence.
       Gnats, gnus, etc.\ all  % `\ ' makes an inter-word space.
       begin with G\@.         % \@ marks end-of-sentence punctuation.
You should check the spaces after periods when
reading your output to make sure you haven't
forgotten any special cases.  Generating an
ellipsis
       \ldots\               % `\ ' is needed after `\ldots' because TeX
                             % ignores spaces after command names like \ldots
                             % made from \ + letters.
                             %
                             % Note how a `%' character causes TeX to ignore
                             % the end of the input line, so these blank lines
                             % do not start a new paragraph.
                             %
with the right spacing around the periods requires
a special command.

\LaTeX\ interprets some common characters as
commands, so you must type special commands to
generate them.  These characters include the
following:
       \$ \& \% \# \{ and \}.

In printing, text is usually emphasized with an
       \emph{italic}
type style.

\begin{em}
   A long segment of text can also be emphasized
   in this way.  Text within such a segment can be
   given \emph{additional} emphasis.
\end{em}

It is sometimes necessary to prevent \LaTeX\ from
breaking a line where it might otherwise do so.
This may be at a space, as between the ``Mr.''\ and
``Jones'' in
       ``Mr.~Jones'',        % ~ produces an unbreakable interword space.
or within a word---especially when the word is a
symbol like
       \mbox{\emph{itemnum}}
that makes little sense when hyphenated across
lines.

Footnotes\footnote{This is an example of a footnote.}
pose no problem.

\LaTeX\ is good at typesetting mathematical formulas
like
       \( x-3y + z = 7 \)
or
       \( a_{1} > x^{2n} + y^{2n} > x' \)
or
       \( \ip{A}{B} = \sum_{i} a_{i} b_{i} \).
The spaces you type in a formula are
ignored.  Remember that a letter like
       $x$                   % $ ... $  and  \( ... \)  are equivalent
is a formula when it denotes a mathematical
symbol, and it should be typed as one.

\section{Displayed Text}

Text is displayed by indenting it from the left
margin.  Quotations are commonly displayed.  There
are short quotations
\begin{quote}
   This is a short quotation.  It consists of a
   single paragraph of text.  See how it is formatted.
\end{quote}
and longer ones.
\begin{quotation}
   This is a longer quotation.  It consists of two
   paragraphs of text, neither of which are
   particularly interesting.

   This is the second paragraph of the quotation.  It
   is just as dull as the first paragraph.
\end{quotation}
Another frequently-displayed structure is a list.
The following is an example of an \emph{itemized}
list.
\begin{itemize}
   \item This is the first item of an itemized list.
         Each item in the list is marked with a ``tick''.
         You don't have to worry about what kind of tick
         mark is used.

   \item This is the second item of the list.  It
         contains another list nested inside it.  The inner
         list is an \emph{enumerated} list.
         \begin{enumerate}
            \item This is the first item of an enumerated
                  list that is nested within the itemized list.

            \item This is the second item of the inner list.
                  \LaTeX\ allows you to nest lists deeper than
                  you really should.
         \end{enumerate}
         This is the rest of the second item of the outer
         list.  It is no more interesting than any other
         part of the item.
   \item This is the third item of the list.
\end{itemize}
You can even display poetry.
\begin{verse}
   There is an environment
    for verse \\             % The \\ command separates lines
   Whose features some poets % within a stanza.
   will curse.

                             % One or more blank lines separate stanzas.

   For instead of making\\
   Them do \emph{all} line breaking, \\
   It allows them to put too many words on a line when they'd rather be
   forced to be terse.
\end{verse}

Mathematical formulas may also be displayed.  A
displayed formula
is
one-line long; multiline
formulas require special formatting instructions.
   \[  \ip{\Gamma}{\psi'} = x'' + y^{2} + z_{i}^{n}\]
Don't start a paragraph with a displayed equation,
nor make one a paragraph by itself.

\end{document}               % End of document.

\section{Bench Huge Section 1338}
% synthetic bench section 1338
% This is a sample LaTeX input file.  (Version of 12 August 2004.)
%
% A '%' character causes TeX to ignore all remaining text on the line,
% and is used for comments like this one.

\documentclass{article}      % Specifies the document class

                             % The preamble begins here.
\title{An Example Document}  % Declares the document's title.
\author{Leslie Lamport}      % Declares the author's name.
\date{January 21, 1994}      % Deleting this command produces today's date.

\newcommand{\ip}[2]{(#1, #2)}
                             % Defines \ip{arg1}{arg2} to mean
                             % (arg1, arg2).

%\newcommand{\ip}[2]{\langle #1 | #2\rangle}
                             % This is an alternative definition of
                             % \ip that is commented out.

\begin{document}             % End of preamble and beginning of text.

\maketitle                   % Produces the title.

This is an example input file.  Comparing it with
the output it generates can show you how to
produce a simple document of your own.

\section{Ordinary Text}      % Produces section heading.  Lower-level
                             % sections are begun with similar
                             % \subsection and \subsubsection commands.

The ends  of words and sentences are marked
  by   spaces. It  doesn't matter how many
spaces    you type; one is as good as 100.  The
end of   a line counts as a space.

One   or more   blank lines denote the  end
of  a paragraph.

Since any number of consecutive spaces are treated
like a single one, the formatting of the input
file makes no difference to
      \LaTeX,                % The \LaTeX command generates the LaTeX logo.
but it makes a difference to you.  When you use
\LaTeX, making your input file as easy to read
as possible will be a great help as you write
your document and when you change it.  This sample
file shows how you can add comments to your own input
file.

Because printing is different from typewriting,
there are a number of things that you have to do
differently when preparing an input file than if
you were just typing the document directly.
Quotation marks like
       ``this''
have to be handled specially, as do quotes within
quotes:
       ``\,`this'            % \, separates the double and single quote.
        is what I just
        wrote, not  `that'\,''.

Dashes come in three sizes: an
       intra-word
dash, a medium dash for number ranges like
       1--2,
and a punctuation
       dash---like
this.

A sentence-ending space should be larger than the
space between words within a sentence.  You
sometimes have to type special commands in
conjunction with punctuation characters to get
this right, as in the following sentence.
       Gnats, gnus, etc.\ all  % `\ ' makes an inter-word space.
       begin with G\@.         % \@ marks end-of-sentence punctuation.
You should check the spaces after periods when
reading your output to make sure you haven't
forgotten any special cases.  Generating an
ellipsis
       \ldots\               % `\ ' is needed after `\ldots' because TeX
                             % ignores spaces after command names like \ldots
                             % made from \ + letters.
                             %
                             % Note how a `%' character causes TeX to ignore
                             % the end of the input line, so these blank lines
                             % do not start a new paragraph.
                             %
with the right spacing around the periods requires
a special command.

\LaTeX\ interprets some common characters as
commands, so you must type special commands to
generate them.  These characters include the
following:
       \$ \& \% \# \{ and \}.

In printing, text is usually emphasized with an
       \emph{italic}
type style.

\begin{em}
   A long segment of text can also be emphasized
   in this way.  Text within such a segment can be
   given \emph{additional} emphasis.
\end{em}

It is sometimes necessary to prevent \LaTeX\ from
breaking a line where it might otherwise do so.
This may be at a space, as between the ``Mr.''\ and
``Jones'' in
       ``Mr.~Jones'',        % ~ produces an unbreakable interword space.
or within a word---especially when the word is a
symbol like
       \mbox{\emph{itemnum}}
that makes little sense when hyphenated across
lines.

Footnotes\footnote{This is an example of a footnote.}
pose no problem.

\LaTeX\ is good at typesetting mathematical formulas
like
       \( x-3y + z = 7 \)
or
       \( a_{1} > x^{2n} + y^{2n} > x' \)
or
       \( \ip{A}{B} = \sum_{i} a_{i} b_{i} \).
The spaces you type in a formula are
ignored.  Remember that a letter like
       $x$                   % $ ... $  and  \( ... \)  are equivalent
is a formula when it denotes a mathematical
symbol, and it should be typed as one.

\section{Displayed Text}

Text is displayed by indenting it from the left
margin.  Quotations are commonly displayed.  There
are short quotations
\begin{quote}
   This is a short quotation.  It consists of a
   single paragraph of text.  See how it is formatted.
\end{quote}
and longer ones.
\begin{quotation}
   This is a longer quotation.  It consists of two
   paragraphs of text, neither of which are
   particularly interesting.

   This is the second paragraph of the quotation.  It
   is just as dull as the first paragraph.
\end{quotation}
Another frequently-displayed structure is a list.
The following is an example of an \emph{itemized}
list.
\begin{itemize}
   \item This is the first item of an itemized list.
         Each item in the list is marked with a ``tick''.
         You don't have to worry about what kind of tick
         mark is used.

   \item This is the second item of the list.  It
         contains another list nested inside it.  The inner
         list is an \emph{enumerated} list.
         \begin{enumerate}
            \item This is the first item of an enumerated
                  list that is nested within the itemized list.

            \item This is the second item of the inner list.
                  \LaTeX\ allows you to nest lists deeper than
                  you really should.
         \end{enumerate}
         This is the rest of the second item of the outer
         list.  It is no more interesting than any other
         part of the item.
   \item This is the third item of the list.
\end{itemize}
You can even display poetry.
\begin{verse}
   There is an environment
    for verse \\             % The \\ command separates lines
   Whose features some poets % within a stanza.
   will curse.

                             % One or more blank lines separate stanzas.

   For instead of making\\
   Them do \emph{all} line breaking, \\
   It allows them to put too many words on a line when they'd rather be
   forced to be terse.
\end{verse}

Mathematical formulas may also be displayed.  A
displayed formula
is
one-line long; multiline
formulas require special formatting instructions.
   \[  \ip{\Gamma}{\psi'} = x'' + y^{2} + z_{i}^{n}\]
Don't start a paragraph with a displayed equation,
nor make one a paragraph by itself.

\end{document}               % End of document.

\section{Bench Huge Section 1339}
% synthetic bench section 1339
% This is a sample LaTeX input file.  (Version of 12 August 2004.)
%
% A '%' character causes TeX to ignore all remaining text on the line,
% and is used for comments like this one.

\documentclass{article}      % Specifies the document class

                             % The preamble begins here.
\title{An Example Document}  % Declares the document's title.
\author{Leslie Lamport}      % Declares the author's name.
\date{January 21, 1994}      % Deleting this command produces today's date.

\newcommand{\ip}[2]{(#1, #2)}
                             % Defines \ip{arg1}{arg2} to mean
                             % (arg1, arg2).

%\newcommand{\ip}[2]{\langle #1 | #2\rangle}
                             % This is an alternative definition of
                             % \ip that is commented out.

\begin{document}             % End of preamble and beginning of text.

\maketitle                   % Produces the title.

This is an example input file.  Comparing it with
the output it generates can show you how to
produce a simple document of your own.

\section{Ordinary Text}      % Produces section heading.  Lower-level
                             % sections are begun with similar
                             % \subsection and \subsubsection commands.

The ends  of words and sentences are marked
  by   spaces. It  doesn't matter how many
spaces    you type; one is as good as 100.  The
end of   a line counts as a space.

One   or more   blank lines denote the  end
of  a paragraph.

Since any number of consecutive spaces are treated
like a single one, the formatting of the input
file makes no difference to
      \LaTeX,                % The \LaTeX command generates the LaTeX logo.
but it makes a difference to you.  When you use
\LaTeX, making your input file as easy to read
as possible will be a great help as you write
your document and when you change it.  This sample
file shows how you can add comments to your own input
file.

Because printing is different from typewriting,
there are a number of things that you have to do
differently when preparing an input file than if
you were just typing the document directly.
Quotation marks like
       ``this''
have to be handled specially, as do quotes within
quotes:
       ``\,`this'            % \, separates the double and single quote.
        is what I just
        wrote, not  `that'\,''.

Dashes come in three sizes: an
       intra-word
dash, a medium dash for number ranges like
       1--2,
and a punctuation
       dash---like
this.

A sentence-ending space should be larger than the
space between words within a sentence.  You
sometimes have to type special commands in
conjunction with punctuation characters to get
this right, as in the following sentence.
       Gnats, gnus, etc.\ all  % `\ ' makes an inter-word space.
       begin with G\@.         % \@ marks end-of-sentence punctuation.
You should check the spaces after periods when
reading your output to make sure you haven't
forgotten any special cases.  Generating an
ellipsis
       \ldots\               % `\ ' is needed after `\ldots' because TeX
                             % ignores spaces after command names like \ldots
                             % made from \ + letters.
                             %
                             % Note how a `%' character causes TeX to ignore
                             % the end of the input line, so these blank lines
                             % do not start a new paragraph.
                             %
with the right spacing around the periods requires
a special command.

\LaTeX\ interprets some common characters as
commands, so you must type special commands to
generate them.  These characters include the
following:
       \$ \& \% \# \{ and \}.

In printing, text is usually emphasized with an
       \emph{italic}
type style.

\begin{em}
   A long segment of text can also be emphasized
   in this way.  Text within such a segment can be
   given \emph{additional} emphasis.
\end{em}

It is sometimes necessary to prevent \LaTeX\ from
breaking a line where it might otherwise do so.
This may be at a space, as between the ``Mr.''\ and
``Jones'' in
       ``Mr.~Jones'',        % ~ produces an unbreakable interword space.
or within a word---especially when the word is a
symbol like
       \mbox{\emph{itemnum}}
that makes little sense when hyphenated across
lines.

Footnotes\footnote{This is an example of a footnote.}
pose no problem.

\LaTeX\ is good at typesetting mathematical formulas
like
       \( x-3y + z = 7 \)
or
       \( a_{1} > x^{2n} + y^{2n} > x' \)
or
       \( \ip{A}{B} = \sum_{i} a_{i} b_{i} \).
The spaces you type in a formula are
ignored.  Remember that a letter like
       $x$                   % $ ... $  and  \( ... \)  are equivalent
is a formula when it denotes a mathematical
symbol, and it should be typed as one.

\section{Displayed Text}

Text is displayed by indenting it from the left
margin.  Quotations are commonly displayed.  There
are short quotations
\begin{quote}
   This is a short quotation.  It consists of a
   single paragraph of text.  See how it is formatted.
\end{quote}
and longer ones.
\begin{quotation}
   This is a longer quotation.  It consists of two
   paragraphs of text, neither of which are
   particularly interesting.

   This is the second paragraph of the quotation.  It
   is just as dull as the first paragraph.
\end{quotation}
Another frequently-displayed structure is a list.
The following is an example of an \emph{itemized}
list.
\begin{itemize}
   \item This is the first item of an itemized list.
         Each item in the list is marked with a ``tick''.
         You don't have to worry about what kind of tick
         mark is used.

   \item This is the second item of the list.  It
         contains another list nested inside it.  The inner
         list is an \emph{enumerated} list.
         \begin{enumerate}
            \item This is the first item of an enumerated
                  list that is nested within the itemized list.

            \item This is the second item of the inner list.
                  \LaTeX\ allows you to nest lists deeper than
                  you really should.
         \end{enumerate}
         This is the rest of the second item of the outer
         list.  It is no more interesting than any other
         part of the item.
   \item This is the third item of the list.
\end{itemize}
You can even display poetry.
\begin{verse}
   There is an environment
    for verse \\             % The \\ command separates lines
   Whose features some poets % within a stanza.
   will curse.

                             % One or more blank lines separate stanzas.

   For instead of making\\
   Them do \emph{all} line breaking, \\
   It allows them to put too many words on a line when they'd rather be
   forced to be terse.
\end{verse}

Mathematical formulas may also be displayed.  A
displayed formula
is
one-line long; multiline
formulas require special formatting instructions.
   \[  \ip{\Gamma}{\psi'} = x'' + y^{2} + z_{i}^{n}\]
Don't start a paragraph with a displayed equation,
nor make one a paragraph by itself.

\end{document}               % End of document.

\section{Bench Huge Section 1340}
% synthetic bench section 1340
% This is a sample LaTeX input file.  (Version of 12 August 2004.)
%
% A '%' character causes TeX to ignore all remaining text on the line,
% and is used for comments like this one.

\documentclass{article}      % Specifies the document class

                             % The preamble begins here.
\title{An Example Document}  % Declares the document's title.
\author{Leslie Lamport}      % Declares the author's name.
\date{January 21, 1994}      % Deleting this command produces today's date.

\newcommand{\ip}[2]{(#1, #2)}
                             % Defines \ip{arg1}{arg2} to mean
                             % (arg1, arg2).

%\newcommand{\ip}[2]{\langle #1 | #2\rangle}
                             % This is an alternative definition of
                             % \ip that is commented out.

\begin{document}             % End of preamble and beginning of text.

\maketitle                   % Produces the title.

This is an example input file.  Comparing it with
the output it generates can show you how to
produce a simple document of your own.

\section{Ordinary Text}      % Produces section heading.  Lower-level
                             % sections are begun with similar
                             % \subsection and \subsubsection commands.

The ends  of words and sentences are marked
  by   spaces. It  doesn't matter how many
spaces    you type; one is as good as 100.  The
end of   a line counts as a space.

One   or more   blank lines denote the  end
of  a paragraph.

Since any number of consecutive spaces are treated
like a single one, the formatting of the input
file makes no difference to
      \LaTeX,                % The \LaTeX command generates the LaTeX logo.
but it makes a difference to you.  When you use
\LaTeX, making your input file as easy to read
as possible will be a great help as you write
your document and when you change it.  This sample
file shows how you can add comments to your own input
file.

Because printing is different from typewriting,
there are a number of things that you have to do
differently when preparing an input file than if
you were just typing the document directly.
Quotation marks like
       ``this''
have to be handled specially, as do quotes within
quotes:
       ``\,`this'            % \, separates the double and single quote.
        is what I just
        wrote, not  `that'\,''.

Dashes come in three sizes: an
       intra-word
dash, a medium dash for number ranges like
       1--2,
and a punctuation
       dash---like
this.

A sentence-ending space should be larger than the
space between words within a sentence.  You
sometimes have to type special commands in
conjunction with punctuation characters to get
this right, as in the following sentence.
       Gnats, gnus, etc.\ all  % `\ ' makes an inter-word space.
       begin with G\@.         % \@ marks end-of-sentence punctuation.
You should check the spaces after periods when
reading your output to make sure you haven't
forgotten any special cases.  Generating an
ellipsis
       \ldots\               % `\ ' is needed after `\ldots' because TeX
                             % ignores spaces after command names like \ldots
                             % made from \ + letters.
                             %
                             % Note how a `%' character causes TeX to ignore
                             % the end of the input line, so these blank lines
                             % do not start a new paragraph.
                             %
with the right spacing around the periods requires
a special command.

\LaTeX\ interprets some common characters as
commands, so you must type special commands to
generate them.  These characters include the
following:
       \$ \& \% \# \{ and \}.

In printing, text is usually emphasized with an
       \emph{italic}
type style.

\begin{em}
   A long segment of text can also be emphasized
   in this way.  Text within such a segment can be
   given \emph{additional} emphasis.
\end{em}

It is sometimes necessary to prevent \LaTeX\ from
breaking a line where it might otherwise do so.
This may be at a space, as between the ``Mr.''\ and
``Jones'' in
       ``Mr.~Jones'',        % ~ produces an unbreakable interword space.
or within a word---especially when the word is a
symbol like
       \mbox{\emph{itemnum}}
that makes little sense when hyphenated across
lines.

Footnotes\footnote{This is an example of a footnote.}
pose no problem.

\LaTeX\ is good at typesetting mathematical formulas
like
       \( x-3y + z = 7 \)
or
       \( a_{1} > x^{2n} + y^{2n} > x' \)
or
       \( \ip{A}{B} = \sum_{i} a_{i} b_{i} \).
The spaces you type in a formula are
ignored.  Remember that a letter like
       $x$                   % $ ... $  and  \( ... \)  are equivalent
is a formula when it denotes a mathematical
symbol, and it should be typed as one.

\section{Displayed Text}

Text is displayed by indenting it from the left
margin.  Quotations are commonly displayed.  There
are short quotations
\begin{quote}
   This is a short quotation.  It consists of a
   single paragraph of text.  See how it is formatted.
\end{quote}
and longer ones.
\begin{quotation}
   This is a longer quotation.  It consists of two
   paragraphs of text, neither of which are
   particularly interesting.

   This is the second paragraph of the quotation.  It
   is just as dull as the first paragraph.
\end{quotation}
Another frequently-displayed structure is a list.
The following is an example of an \emph{itemized}
list.
\begin{itemize}
   \item This is the first item of an itemized list.
         Each item in the list is marked with a ``tick''.
         You don't have to worry about what kind of tick
         mark is used.

   \item This is the second item of the list.  It
         contains another list nested inside it.  The inner
         list is an \emph{enumerated} list.
         \begin{enumerate}
            \item This is the first item of an enumerated
                  list that is nested within the itemized list.

            \item This is the second item of the inner list.
                  \LaTeX\ allows you to nest lists deeper than
                  you really should.
         \end{enumerate}
         This is the rest of the second item of the outer
         list.  It is no more interesting than any other
         part of the item.
   \item This is the third item of the list.
\end{itemize}
You can even display poetry.
\begin{verse}
   There is an environment
    for verse \\             % The \\ command separates lines
   Whose features some poets % within a stanza.
   will curse.

                             % One or more blank lines separate stanzas.

   For instead of making\\
   Them do \emph{all} line breaking, \\
   It allows them to put too many words on a line when they'd rather be
   forced to be terse.
\end{verse}

Mathematical formulas may also be displayed.  A
displayed formula
is
one-line long; multiline
formulas require special formatting instructions.
   \[  \ip{\Gamma}{\psi'} = x'' + y^{2} + z_{i}^{n}\]
Don't start a paragraph with a displayed equation,
nor make one a paragraph by itself.

\end{document}               % End of document.

\section{Bench Huge Section 1341}
% synthetic bench section 1341
% This is a sample LaTeX input file.  (Version of 12 August 2004.)
%
% A '%' character causes TeX to ignore all remaining text on the line,
% and is used for comments like this one.

\documentclass{article}      % Specifies the document class

                             % The preamble begins here.
\title{An Example Document}  % Declares the document's title.
\author{Leslie Lamport}      % Declares the author's name.
\date{January 21, 1994}      % Deleting this command produces today's date.

\newcommand{\ip}[2]{(#1, #2)}
                             % Defines \ip{arg1}{arg2} to mean
                             % (arg1, arg2).

%\newcommand{\ip}[2]{\langle #1 | #2\rangle}
                             % This is an alternative definition of
                             % \ip that is commented out.

\begin{document}             % End of preamble and beginning of text.

\maketitle                   % Produces the title.

This is an example input file.  Comparing it with
the output it generates can show you how to
produce a simple document of your own.

\section{Ordinary Text}      % Produces section heading.  Lower-level
                             % sections are begun with similar
                             % \subsection and \subsubsection commands.

The ends  of words and sentences are marked
  by   spaces. It  doesn't matter how many
spaces    you type; one is as good as 100.  The
end of   a line counts as a space.

One   or more   blank lines denote the  end
of  a paragraph.

Since any number of consecutive spaces are treated
like a single one, the formatting of the input
file makes no difference to
      \LaTeX,                % The \LaTeX command generates the LaTeX logo.
but it makes a difference to you.  When you use
\LaTeX, making your input file as easy to read
as possible will be a great help as you write
your document and when you change it.  This sample
file shows how you can add comments to your own input
file.

Because printing is different from typewriting,
there are a number of things that you have to do
differently when preparing an input file than if
you were just typing the document directly.
Quotation marks like
       ``this''
have to be handled specially, as do quotes within
quotes:
       ``\,`this'            % \, separates the double and single quote.
        is what I just
        wrote, not  `that'\,''.

Dashes come in three sizes: an
       intra-word
dash, a medium dash for number ranges like
       1--2,
and a punctuation
       dash---like
this.

A sentence-ending space should be larger than the
space between words within a sentence.  You
sometimes have to type special commands in
conjunction with punctuation characters to get
this right, as in the following sentence.
       Gnats, gnus, etc.\ all  % `\ ' makes an inter-word space.
       begin with G\@.         % \@ marks end-of-sentence punctuation.
You should check the spaces after periods when
reading your output to make sure you haven't
forgotten any special cases.  Generating an
ellipsis
       \ldots\               % `\ ' is needed after `\ldots' because TeX
                             % ignores spaces after command names like \ldots
                             % made from \ + letters.
                             %
                             % Note how a `%' character causes TeX to ignore
                             % the end of the input line, so these blank lines
                             % do not start a new paragraph.
                             %
with the right spacing around the periods requires
a special command.

\LaTeX\ interprets some common characters as
commands, so you must type special commands to
generate them.  These characters include the
following:
       \$ \& \% \# \{ and \}.

In printing, text is usually emphasized with an
       \emph{italic}
type style.

\begin{em}
   A long segment of text can also be emphasized
   in this way.  Text within such a segment can be
   given \emph{additional} emphasis.
\end{em}

It is sometimes necessary to prevent \LaTeX\ from
breaking a line where it might otherwise do so.
This may be at a space, as between the ``Mr.''\ and
``Jones'' in
       ``Mr.~Jones'',        % ~ produces an unbreakable interword space.
or within a word---especially when the word is a
symbol like
       \mbox{\emph{itemnum}}
that makes little sense when hyphenated across
lines.

Footnotes\footnote{This is an example of a footnote.}
pose no problem.

\LaTeX\ is good at typesetting mathematical formulas
like
       \( x-3y + z = 7 \)
or
       \( a_{1} > x^{2n} + y^{2n} > x' \)
or
       \( \ip{A}{B} = \sum_{i} a_{i} b_{i} \).
The spaces you type in a formula are
ignored.  Remember that a letter like
       $x$                   % $ ... $  and  \( ... \)  are equivalent
is a formula when it denotes a mathematical
symbol, and it should be typed as one.

\section{Displayed Text}

Text is displayed by indenting it from the left
margin.  Quotations are commonly displayed.  There
are short quotations
\begin{quote}
   This is a short quotation.  It consists of a
   single paragraph of text.  See how it is formatted.
\end{quote}
and longer ones.
\begin{quotation}
   This is a longer quotation.  It consists of two
   paragraphs of text, neither of which are
   particularly interesting.

   This is the second paragraph of the quotation.  It
   is just as dull as the first paragraph.
\end{quotation}
Another frequently-displayed structure is a list.
The following is an example of an \emph{itemized}
list.
\begin{itemize}
   \item This is the first item of an itemized list.
         Each item in the list is marked with a ``tick''.
         You don't have to worry about what kind of tick
         mark is used.

   \item This is the second item of the list.  It
         contains another list nested inside it.  The inner
         list is an \emph{enumerated} list.
         \begin{enumerate}
            \item This is the first item of an enumerated
                  list that is nested within the itemized list.

            \item This is the second item of the inner list.
                  \LaTeX\ allows you to nest lists deeper than
                  you really should.
         \end{enumerate}
         This is the rest of the second item of the outer
         list.  It is no more interesting than any other
         part of the item.
   \item This is the third item of the list.
\end{itemize}
You can even display poetry.
\begin{verse}
   There is an environment
    for verse \\             % The \\ command separates lines
   Whose features some poets % within a stanza.
   will curse.

                             % One or more blank lines separate stanzas.

   For instead of making\\
   Them do \emph{all} line breaking, \\
   It allows them to put too many words on a line when they'd rather be
   forced to be terse.
\end{verse}

Mathematical formulas may also be displayed.  A
displayed formula
is
one-line long; multiline
formulas require special formatting instructions.
   \[  \ip{\Gamma}{\psi'} = x'' + y^{2} + z_{i}^{n}\]
Don't start a paragraph with a displayed equation,
nor make one a paragraph by itself.

\end{document}               % End of document.

\section{Bench Huge Section 1342}
% synthetic bench section 1342
% This is a sample LaTeX input file.  (Version of 12 August 2004.)
%
% A '%' character causes TeX to ignore all remaining text on the line,
% and is used for comments like this one.

\documentclass{article}      % Specifies the document class

                             % The preamble begins here.
\title{An Example Document}  % Declares the document's title.
\author{Leslie Lamport}      % Declares the author's name.
\date{January 21, 1994}      % Deleting this command produces today's date.

\newcommand{\ip}[2]{(#1, #2)}
                             % Defines \ip{arg1}{arg2} to mean
                             % (arg1, arg2).

%\newcommand{\ip}[2]{\langle #1 | #2\rangle}
                             % This is an alternative definition of
                             % \ip that is commented out.

\begin{document}             % End of preamble and beginning of text.

\maketitle                   % Produces the title.

This is an example input file.  Comparing it with
the output it generates can show you how to
produce a simple document of your own.

\section{Ordinary Text}      % Produces section heading.  Lower-level
                             % sections are begun with similar
                             % \subsection and \subsubsection commands.

The ends  of words and sentences are marked
  by   spaces. It  doesn't matter how many
spaces    you type; one is as good as 100.  The
end of   a line counts as a space.

One   or more   blank lines denote the  end
of  a paragraph.

Since any number of consecutive spaces are treated
like a single one, the formatting of the input
file makes no difference to
      \LaTeX,                % The \LaTeX command generates the LaTeX logo.
but it makes a difference to you.  When you use
\LaTeX, making your input file as easy to read
as possible will be a great help as you write
your document and when you change it.  This sample
file shows how you can add comments to your own input
file.

Because printing is different from typewriting,
there are a number of things that you have to do
differently when preparing an input file than if
you were just typing the document directly.
Quotation marks like
       ``this''
have to be handled specially, as do quotes within
quotes:
       ``\,`this'            % \, separates the double and single quote.
        is what I just
        wrote, not  `that'\,''.

Dashes come in three sizes: an
       intra-word
dash, a medium dash for number ranges like
       1--2,
and a punctuation
       dash---like
this.

A sentence-ending space should be larger than the
space between words within a sentence.  You
sometimes have to type special commands in
conjunction with punctuation characters to get
this right, as in the following sentence.
       Gnats, gnus, etc.\ all  % `\ ' makes an inter-word space.
       begin with G\@.         % \@ marks end-of-sentence punctuation.
You should check the spaces after periods when
reading your output to make sure you haven't
forgotten any special cases.  Generating an
ellipsis
       \ldots\               % `\ ' is needed after `\ldots' because TeX
                             % ignores spaces after command names like \ldots
                             % made from \ + letters.
                             %
                             % Note how a `%' character causes TeX to ignore
                             % the end of the input line, so these blank lines
                             % do not start a new paragraph.
                             %
with the right spacing around the periods requires
a special command.

\LaTeX\ interprets some common characters as
commands, so you must type special commands to
generate them.  These characters include the
following:
       \$ \& \% \# \{ and \}.

In printing, text is usually emphasized with an
       \emph{italic}
type style.

\begin{em}
   A long segment of text can also be emphasized
   in this way.  Text within such a segment can be
   given \emph{additional} emphasis.
\end{em}

It is sometimes necessary to prevent \LaTeX\ from
breaking a line where it might otherwise do so.
This may be at a space, as between the ``Mr.''\ and
``Jones'' in
       ``Mr.~Jones'',        % ~ produces an unbreakable interword space.
or within a word---especially when the word is a
symbol like
       \mbox{\emph{itemnum}}
that makes little sense when hyphenated across
lines.

Footnotes\footnote{This is an example of a footnote.}
pose no problem.

\LaTeX\ is good at typesetting mathematical formulas
like
       \( x-3y + z = 7 \)
or
       \( a_{1} > x^{2n} + y^{2n} > x' \)
or
       \( \ip{A}{B} = \sum_{i} a_{i} b_{i} \).
The spaces you type in a formula are
ignored.  Remember that a letter like
       $x$                   % $ ... $  and  \( ... \)  are equivalent
is a formula when it denotes a mathematical
symbol, and it should be typed as one.

\section{Displayed Text}

Text is displayed by indenting it from the left
margin.  Quotations are commonly displayed.  There
are short quotations
\begin{quote}
   This is a short quotation.  It consists of a
   single paragraph of text.  See how it is formatted.
\end{quote}
and longer ones.
\begin{quotation}
   This is a longer quotation.  It consists of two
   paragraphs of text, neither of which are
   particularly interesting.

   This is the second paragraph of the quotation.  It
   is just as dull as the first paragraph.
\end{quotation}
Another frequently-displayed structure is a list.
The following is an example of an \emph{itemized}
list.
\begin{itemize}
   \item This is the first item of an itemized list.
         Each item in the list is marked with a ``tick''.
         You don't have to worry about what kind of tick
         mark is used.

   \item This is the second item of the list.  It
         contains another list nested inside it.  The inner
         list is an \emph{enumerated} list.
         \begin{enumerate}
            \item This is the first item of an enumerated
                  list that is nested within the itemized list.

            \item This is the second item of the inner list.
                  \LaTeX\ allows you to nest lists deeper than
                  you really should.
         \end{enumerate}
         This is the rest of the second item of the outer
         list.  It is no more interesting than any other
         part of the item.
   \item This is the third item of the list.
\end{itemize}
You can even display poetry.
\begin{verse}
   There is an environment
    for verse \\             % The \\ command separates lines
   Whose features some poets % within a stanza.
   will curse.

                             % One or more blank lines separate stanzas.

   For instead of making\\
   Them do \emph{all} line breaking, \\
   It allows them to put too many words on a line when they'd rather be
   forced to be terse.
\end{verse}

Mathematical formulas may also be displayed.  A
displayed formula
is
one-line long; multiline
formulas require special formatting instructions.
   \[  \ip{\Gamma}{\psi'} = x'' + y^{2} + z_{i}^{n}\]
Don't start a paragraph with a displayed equation,
nor make one a paragraph by itself.

\end{document}               % End of document.

\section{Bench Huge Section 1343}
% synthetic bench section 1343
% This is a sample LaTeX input file.  (Version of 12 August 2004.)
%
% A '%' character causes TeX to ignore all remaining text on the line,
% and is used for comments like this one.

\documentclass{article}      % Specifies the document class

                             % The preamble begins here.
\title{An Example Document}  % Declares the document's title.
\author{Leslie Lamport}      % Declares the author's name.
\date{January 21, 1994}      % Deleting this command produces today's date.

\newcommand{\ip}[2]{(#1, #2)}
                             % Defines \ip{arg1}{arg2} to mean
                             % (arg1, arg2).

%\newcommand{\ip}[2]{\langle #1 | #2\rangle}
                             % This is an alternative definition of
                             % \ip that is commented out.

\begin{document}             % End of preamble and beginning of text.

\maketitle                   % Produces the title.

This is an example input file.  Comparing it with
the output it generates can show you how to
produce a simple document of your own.

\section{Ordinary Text}      % Produces section heading.  Lower-level
                             % sections are begun with similar
                             % \subsection and \subsubsection commands.

The ends  of words and sentences are marked
  by   spaces. It  doesn't matter how many
spaces    you type; one is as good as 100.  The
end of   a line counts as a space.

One   or more   blank lines denote the  end
of  a paragraph.

Since any number of consecutive spaces are treated
like a single one, the formatting of the input
file makes no difference to
      \LaTeX,                % The \LaTeX command generates the LaTeX logo.
but it makes a difference to you.  When you use
\LaTeX, making your input file as easy to read
as possible will be a great help as you write
your document and when you change it.  This sample
file shows how you can add comments to your own input
file.

Because printing is different from typewriting,
there are a number of things that you have to do
differently when preparing an input file than if
you were just typing the document directly.
Quotation marks like
       ``this''
have to be handled specially, as do quotes within
quotes:
       ``\,`this'            % \, separates the double and single quote.
        is what I just
        wrote, not  `that'\,''.

Dashes come in three sizes: an
       intra-word
dash, a medium dash for number ranges like
       1--2,
and a punctuation
       dash---like
this.

A sentence-ending space should be larger than the
space between words within a sentence.  You
sometimes have to type special commands in
conjunction with punctuation characters to get
this right, as in the following sentence.
       Gnats, gnus, etc.\ all  % `\ ' makes an inter-word space.
       begin with G\@.         % \@ marks end-of-sentence punctuation.
You should check the spaces after periods when
reading your output to make sure you haven't
forgotten any special cases.  Generating an
ellipsis
       \ldots\               % `\ ' is needed after `\ldots' because TeX
                             % ignores spaces after command names like \ldots
                             % made from \ + letters.
                             %
                             % Note how a `%' character causes TeX to ignore
                             % the end of the input line, so these blank lines
                             % do not start a new paragraph.
                             %
with the right spacing around the periods requires
a special command.

\LaTeX\ interprets some common characters as
commands, so you must type special commands to
generate them.  These characters include the
following:
       \$ \& \% \# \{ and \}.

In printing, text is usually emphasized with an
       \emph{italic}
type style.

\begin{em}
   A long segment of text can also be emphasized
   in this way.  Text within such a segment can be
   given \emph{additional} emphasis.
\end{em}

It is sometimes necessary to prevent \LaTeX\ from
breaking a line where it might otherwise do so.
This may be at a space, as between the ``Mr.''\ and
``Jones'' in
       ``Mr.~Jones'',        % ~ produces an unbreakable interword space.
or within a word---especially when the word is a
symbol like
       \mbox{\emph{itemnum}}
that makes little sense when hyphenated across
lines.

Footnotes\footnote{This is an example of a footnote.}
pose no problem.

\LaTeX\ is good at typesetting mathematical formulas
like
       \( x-3y + z = 7 \)
or
       \( a_{1} > x^{2n} + y^{2n} > x' \)
or
       \( \ip{A}{B} = \sum_{i} a_{i} b_{i} \).
The spaces you type in a formula are
ignored.  Remember that a letter like
       $x$                   % $ ... $  and  \( ... \)  are equivalent
is a formula when it denotes a mathematical
symbol, and it should be typed as one.

\section{Displayed Text}

Text is displayed by indenting it from the left
margin.  Quotations are commonly displayed.  There
are short quotations
\begin{quote}
   This is a short quotation.  It consists of a
   single paragraph of text.  See how it is formatted.
\end{quote}
and longer ones.
\begin{quotation}
   This is a longer quotation.  It consists of two
   paragraphs of text, neither of which are
   particularly interesting.

   This is the second paragraph of the quotation.  It
   is just as dull as the first paragraph.
\end{quotation}
Another frequently-displayed structure is a list.
The following is an example of an \emph{itemized}
list.
\begin{itemize}
   \item This is the first item of an itemized list.
         Each item in the list is marked with a ``tick''.
         You don't have to worry about what kind of tick
         mark is used.

   \item This is the second item of the list.  It
         contains another list nested inside it.  The inner
         list is an \emph{enumerated} list.
         \begin{enumerate}
            \item This is the first item of an enumerated
                  list that is nested within the itemized list.

            \item This is the second item of the inner list.
                  \LaTeX\ allows you to nest lists deeper than
                  you really should.
         \end{enumerate}
         This is the rest of the second item of the outer
         list.  It is no more interesting than any other
         part of the item.
   \item This is the third item of the list.
\end{itemize}
You can even display poetry.
\begin{verse}
   There is an environment
    for verse \\             % The \\ command separates lines
   Whose features some poets % within a stanza.
   will curse.

                             % One or more blank lines separate stanzas.

   For instead of making\\
   Them do \emph{all} line breaking, \\
   It allows them to put too many words on a line when they'd rather be
   forced to be terse.
\end{verse}

Mathematical formulas may also be displayed.  A
displayed formula
is
one-line long; multiline
formulas require special formatting instructions.
   \[  \ip{\Gamma}{\psi'} = x'' + y^{2} + z_{i}^{n}\]
Don't start a paragraph with a displayed equation,
nor make one a paragraph by itself.

\end{document}               % End of document.

\section{Bench Huge Section 1344}
% synthetic bench section 1344
% This is a sample LaTeX input file.  (Version of 12 August 2004.)
%
% A '%' character causes TeX to ignore all remaining text on the line,
% and is used for comments like this one.

\documentclass{article}      % Specifies the document class

                             % The preamble begins here.
\title{An Example Document}  % Declares the document's title.
\author{Leslie Lamport}      % Declares the author's name.
\date{January 21, 1994}      % Deleting this command produces today's date.

\newcommand{\ip}[2]{(#1, #2)}
                             % Defines \ip{arg1}{arg2} to mean
                             % (arg1, arg2).

%\newcommand{\ip}[2]{\langle #1 | #2\rangle}
                             % This is an alternative definition of
                             % \ip that is commented out.

\begin{document}             % End of preamble and beginning of text.

\maketitle                   % Produces the title.

This is an example input file.  Comparing it with
the output it generates can show you how to
produce a simple document of your own.

\section{Ordinary Text}      % Produces section heading.  Lower-level
                             % sections are begun with similar
                             % \subsection and \subsubsection commands.

The ends  of words and sentences are marked
  by   spaces. It  doesn't matter how many
spaces    you type; one is as good as 100.  The
end of   a line counts as a space.

One   or more   blank lines denote the  end
of  a paragraph.

Since any number of consecutive spaces are treated
like a single one, the formatting of the input
file makes no difference to
      \LaTeX,                % The \LaTeX command generates the LaTeX logo.
but it makes a difference to you.  When you use
\LaTeX, making your input file as easy to read
as possible will be a great help as you write
your document and when you change it.  This sample
file shows how you can add comments to your own input
file.

Because printing is different from typewriting,
there are a number of things that you have to do
differently when preparing an input file than if
you were just typing the document directly.
Quotation marks like
       ``this''
have to be handled specially, as do quotes within
quotes:
       ``\,`this'            % \, separates the double and single quote.
        is what I just
        wrote, not  `that'\,''.

Dashes come in three sizes: an
       intra-word
dash, a medium dash for number ranges like
       1--2,
and a punctuation
       dash---like
this.

A sentence-ending space should be larger than the
space between words within a sentence.  You
sometimes have to type special commands in
conjunction with punctuation characters to get
this right, as in the following sentence.
       Gnats, gnus, etc.\ all  % `\ ' makes an inter-word space.
       begin with G\@.         % \@ marks end-of-sentence punctuation.
You should check the spaces after periods when
reading your output to make sure you haven't
forgotten any special cases.  Generating an
ellipsis
       \ldots\               % `\ ' is needed after `\ldots' because TeX
                             % ignores spaces after command names like \ldots
                             % made from \ + letters.
                             %
                             % Note how a `%' character causes TeX to ignore
                             % the end of the input line, so these blank lines
                             % do not start a new paragraph.
                             %
with the right spacing around the periods requires
a special command.

\LaTeX\ interprets some common characters as
commands, so you must type special commands to
generate them.  These characters include the
following:
       \$ \& \% \# \{ and \}.

In printing, text is usually emphasized with an
       \emph{italic}
type style.

\begin{em}
   A long segment of text can also be emphasized
   in this way.  Text within such a segment can be
   given \emph{additional} emphasis.
\end{em}

It is sometimes necessary to prevent \LaTeX\ from
breaking a line where it might otherwise do so.
This may be at a space, as between the ``Mr.''\ and
``Jones'' in
       ``Mr.~Jones'',        % ~ produces an unbreakable interword space.
or within a word---especially when the word is a
symbol like
       \mbox{\emph{itemnum}}
that makes little sense when hyphenated across
lines.

Footnotes\footnote{This is an example of a footnote.}
pose no problem.

\LaTeX\ is good at typesetting mathematical formulas
like
       \( x-3y + z = 7 \)
or
       \( a_{1} > x^{2n} + y^{2n} > x' \)
or
       \( \ip{A}{B} = \sum_{i} a_{i} b_{i} \).
The spaces you type in a formula are
ignored.  Remember that a letter like
       $x$                   % $ ... $  and  \( ... \)  are equivalent
is a formula when it denotes a mathematical
symbol, and it should be typed as one.

\section{Displayed Text}

Text is displayed by indenting it from the left
margin.  Quotations are commonly displayed.  There
are short quotations
\begin{quote}
   This is a short quotation.  It consists of a
   single paragraph of text.  See how it is formatted.
\end{quote}
and longer ones.
\begin{quotation}
   This is a longer quotation.  It consists of two
   paragraphs of text, neither of which are
   particularly interesting.

   This is the second paragraph of the quotation.  It
   is just as dull as the first paragraph.
\end{quotation}
Another frequently-displayed structure is a list.
The following is an example of an \emph{itemized}
list.
\begin{itemize}
   \item This is the first item of an itemized list.
         Each item in the list is marked with a ``tick''.
         You don't have to worry about what kind of tick
         mark is used.

   \item This is the second item of the list.  It
         contains another list nested inside it.  The inner
         list is an \emph{enumerated} list.
         \begin{enumerate}
            \item This is the first item of an enumerated
                  list that is nested within the itemized list.

            \item This is the second item of the inner list.
                  \LaTeX\ allows you to nest lists deeper than
                  you really should.
         \end{enumerate}
         This is the rest of the second item of the outer
         list.  It is no more interesting than any other
         part of the item.
   \item This is the third item of the list.
\end{itemize}
You can even display poetry.
\begin{verse}
   There is an environment
    for verse \\             % The \\ command separates lines
   Whose features some poets % within a stanza.
   will curse.

                             % One or more blank lines separate stanzas.

   For instead of making\\
   Them do \emph{all} line breaking, \\
   It allows them to put too many words on a line when they'd rather be
   forced to be terse.
\end{verse}

Mathematical formulas may also be displayed.  A
displayed formula
is
one-line long; multiline
formulas require special formatting instructions.
   \[  \ip{\Gamma}{\psi'} = x'' + y^{2} + z_{i}^{n}\]
Don't start a paragraph with a displayed equation,
nor make one a paragraph by itself.

\end{document}               % End of document.

\section{Bench Huge Section 1345}
% synthetic bench section 1345
% This is a sample LaTeX input file.  (Version of 12 August 2004.)
%
% A '%' character causes TeX to ignore all remaining text on the line,
% and is used for comments like this one.

\documentclass{article}      % Specifies the document class

                             % The preamble begins here.
\title{An Example Document}  % Declares the document's title.
\author{Leslie Lamport}      % Declares the author's name.
\date{January 21, 1994}      % Deleting this command produces today's date.

\newcommand{\ip}[2]{(#1, #2)}
                             % Defines \ip{arg1}{arg2} to mean
                             % (arg1, arg2).

%\newcommand{\ip}[2]{\langle #1 | #2\rangle}
                             % This is an alternative definition of
                             % \ip that is commented out.

\begin{document}             % End of preamble and beginning of text.

\maketitle                   % Produces the title.

This is an example input file.  Comparing it with
the output it generates can show you how to
produce a simple document of your own.

\section{Ordinary Text}      % Produces section heading.  Lower-level
                             % sections are begun with similar
                             % \subsection and \subsubsection commands.

The ends  of words and sentences are marked
  by   spaces. It  doesn't matter how many
spaces    you type; one is as good as 100.  The
end of   a line counts as a space.

One   or more   blank lines denote the  end
of  a paragraph.

Since any number of consecutive spaces are treated
like a single one, the formatting of the input
file makes no difference to
      \LaTeX,                % The \LaTeX command generates the LaTeX logo.
but it makes a difference to you.  When you use
\LaTeX, making your input file as easy to read
as possible will be a great help as you write
your document and when you change it.  This sample
file shows how you can add comments to your own input
file.

Because printing is different from typewriting,
there are a number of things that you have to do
differently when preparing an input file than if
you were just typing the document directly.
Quotation marks like
       ``this''
have to be handled specially, as do quotes within
quotes:
       ``\,`this'            % \, separates the double and single quote.
        is what I just
        wrote, not  `that'\,''.

Dashes come in three sizes: an
       intra-word
dash, a medium dash for number ranges like
       1--2,
and a punctuation
       dash---like
this.

A sentence-ending space should be larger than the
space between words within a sentence.  You
sometimes have to type special commands in
conjunction with punctuation characters to get
this right, as in the following sentence.
       Gnats, gnus, etc.\ all  % `\ ' makes an inter-word space.
       begin with G\@.         % \@ marks end-of-sentence punctuation.
You should check the spaces after periods when
reading your output to make sure you haven't
forgotten any special cases.  Generating an
ellipsis
       \ldots\               % `\ ' is needed after `\ldots' because TeX
                             % ignores spaces after command names like \ldots
                             % made from \ + letters.
                             %
                             % Note how a `%' character causes TeX to ignore
                             % the end of the input line, so these blank lines
                             % do not start a new paragraph.
                             %
with the right spacing around the periods requires
a special command.

\LaTeX\ interprets some common characters as
commands, so you must type special commands to
generate them.  These characters include the
following:
       \$ \& \% \# \{ and \}.

In printing, text is usually emphasized with an
       \emph{italic}
type style.

\begin{em}
   A long segment of text can also be emphasized
   in this way.  Text within such a segment can be
   given \emph{additional} emphasis.
\end{em}

It is sometimes necessary to prevent \LaTeX\ from
breaking a line where it might otherwise do so.
This may be at a space, as between the ``Mr.''\ and
``Jones'' in
       ``Mr.~Jones'',        % ~ produces an unbreakable interword space.
or within a word---especially when the word is a
symbol like
       \mbox{\emph{itemnum}}
that makes little sense when hyphenated across
lines.

Footnotes\footnote{This is an example of a footnote.}
pose no problem.

\LaTeX\ is good at typesetting mathematical formulas
like
       \( x-3y + z = 7 \)
or
       \( a_{1} > x^{2n} + y^{2n} > x' \)
or
       \( \ip{A}{B} = \sum_{i} a_{i} b_{i} \).
The spaces you type in a formula are
ignored.  Remember that a letter like
       $x$                   % $ ... $  and  \( ... \)  are equivalent
is a formula when it denotes a mathematical
symbol, and it should be typed as one.

\section{Displayed Text}

Text is displayed by indenting it from the left
margin.  Quotations are commonly displayed.  There
are short quotations
\begin{quote}
   This is a short quotation.  It consists of a
   single paragraph of text.  See how it is formatted.
\end{quote}
and longer ones.
\begin{quotation}
   This is a longer quotation.  It consists of two
   paragraphs of text, neither of which are
   particularly interesting.

   This is the second paragraph of the quotation.  It
   is just as dull as the first paragraph.
\end{quotation}
Another frequently-displayed structure is a list.
The following is an example of an \emph{itemized}
list.
\begin{itemize}
   \item This is the first item of an itemized list.
         Each item in the list is marked with a ``tick''.
         You don't have to worry about what kind of tick
         mark is used.

   \item This is the second item of the list.  It
         contains another list nested inside it.  The inner
         list is an \emph{enumerated} list.
         \begin{enumerate}
            \item This is the first item of an enumerated
                  list that is nested within the itemized list.

            \item This is the second item of the inner list.
                  \LaTeX\ allows you to nest lists deeper than
                  you really should.
         \end{enumerate}
         This is the rest of the second item of the outer
         list.  It is no more interesting than any other
         part of the item.
   \item This is the third item of the list.
\end{itemize}
You can even display poetry.
\begin{verse}
   There is an environment
    for verse \\             % The \\ command separates lines
   Whose features some poets % within a stanza.
   will curse.

                             % One or more blank lines separate stanzas.

   For instead of making\\
   Them do \emph{all} line breaking, \\
   It allows them to put too many words on a line when they'd rather be
   forced to be terse.
\end{verse}

Mathematical formulas may also be displayed.  A
displayed formula
is
one-line long; multiline
formulas require special formatting instructions.
   \[  \ip{\Gamma}{\psi'} = x'' + y^{2} + z_{i}^{n}\]
Don't start a paragraph with a displayed equation,
nor make one a paragraph by itself.

\end{document}               % End of document.

\section{Bench Huge Section 1346}
% synthetic bench section 1346
% This is a sample LaTeX input file.  (Version of 12 August 2004.)
%
% A '%' character causes TeX to ignore all remaining text on the line,
% and is used for comments like this one.

\documentclass{article}      % Specifies the document class

                             % The preamble begins here.
\title{An Example Document}  % Declares the document's title.
\author{Leslie Lamport}      % Declares the author's name.
\date{January 21, 1994}      % Deleting this command produces today's date.

\newcommand{\ip}[2]{(#1, #2)}
                             % Defines \ip{arg1}{arg2} to mean
                             % (arg1, arg2).

%\newcommand{\ip}[2]{\langle #1 | #2\rangle}
                             % This is an alternative definition of
                             % \ip that is commented out.

\begin{document}             % End of preamble and beginning of text.

\maketitle                   % Produces the title.

This is an example input file.  Comparing it with
the output it generates can show you how to
produce a simple document of your own.

\section{Ordinary Text}      % Produces section heading.  Lower-level
                             % sections are begun with similar
                             % \subsection and \subsubsection commands.

The ends  of words and sentences are marked
  by   spaces. It  doesn't matter how many
spaces    you type; one is as good as 100.  The
end of   a line counts as a space.

One   or more   blank lines denote the  end
of  a paragraph.

Since any number of consecutive spaces are treated
like a single one, the formatting of the input
file makes no difference to
      \LaTeX,                % The \LaTeX command generates the LaTeX logo.
but it makes a difference to you.  When you use
\LaTeX, making your input file as easy to read
as possible will be a great help as you write
your document and when you change it.  This sample
file shows how you can add comments to your own input
file.

Because printing is different from typewriting,
there are a number of things that you have to do
differently when preparing an input file than if
you were just typing the document directly.
Quotation marks like
       ``this''
have to be handled specially, as do quotes within
quotes:
       ``\,`this'            % \, separates the double and single quote.
        is what I just
        wrote, not  `that'\,''.

Dashes come in three sizes: an
       intra-word
dash, a medium dash for number ranges like
       1--2,
and a punctuation
       dash---like
this.

A sentence-ending space should be larger than the
space between words within a sentence.  You
sometimes have to type special commands in
conjunction with punctuation characters to get
this right, as in the following sentence.
       Gnats, gnus, etc.\ all  % `\ ' makes an inter-word space.
       begin with G\@.         % \@ marks end-of-sentence punctuation.
You should check the spaces after periods when
reading your output to make sure you haven't
forgotten any special cases.  Generating an
ellipsis
       \ldots\               % `\ ' is needed after `\ldots' because TeX
                             % ignores spaces after command names like \ldots
                             % made from \ + letters.
                             %
                             % Note how a `%' character causes TeX to ignore
                             % the end of the input line, so these blank lines
                             % do not start a new paragraph.
                             %
with the right spacing around the periods requires
a special command.

\LaTeX\ interprets some common characters as
commands, so you must type special commands to
generate them.  These characters include the
following:
       \$ \& \% \# \{ and \}.

In printing, text is usually emphasized with an
       \emph{italic}
type style.

\begin{em}
   A long segment of text can also be emphasized
   in this way.  Text within such a segment can be
   given \emph{additional} emphasis.
\end{em}

It is sometimes necessary to prevent \LaTeX\ from
breaking a line where it might otherwise do so.
This may be at a space, as between the ``Mr.''\ and
``Jones'' in
       ``Mr.~Jones'',        % ~ produces an unbreakable interword space.
or within a word---especially when the word is a
symbol like
       \mbox{\emph{itemnum}}
that makes little sense when hyphenated across
lines.

Footnotes\footnote{This is an example of a footnote.}
pose no problem.

\LaTeX\ is good at typesetting mathematical formulas
like
       \( x-3y + z = 7 \)
or
       \( a_{1} > x^{2n} + y^{2n} > x' \)
or
       \( \ip{A}{B} = \sum_{i} a_{i} b_{i} \).
The spaces you type in a formula are
ignored.  Remember that a letter like
       $x$                   % $ ... $  and  \( ... \)  are equivalent
is a formula when it denotes a mathematical
symbol, and it should be typed as one.

\section{Displayed Text}

Text is displayed by indenting it from the left
margin.  Quotations are commonly displayed.  There
are short quotations
\begin{quote}
   This is a short quotation.  It consists of a
   single paragraph of text.  See how it is formatted.
\end{quote}
and longer ones.
\begin{quotation}
   This is a longer quotation.  It consists of two
   paragraphs of text, neither of which are
   particularly interesting.

   This is the second paragraph of the quotation.  It
   is just as dull as the first paragraph.
\end{quotation}
Another frequently-displayed structure is a list.
The following is an example of an \emph{itemized}
list.
\begin{itemize}
   \item This is the first item of an itemized list.
         Each item in the list is marked with a ``tick''.
         You don't have to worry about what kind of tick
         mark is used.

   \item This is the second item of the list.  It
         contains another list nested inside it.  The inner
         list is an \emph{enumerated} list.
         \begin{enumerate}
            \item This is the first item of an enumerated
                  list that is nested within the itemized list.

            \item This is the second item of the inner list.
                  \LaTeX\ allows you to nest lists deeper than
                  you really should.
         \end{enumerate}
         This is the rest of the second item of the outer
         list.  It is no more interesting than any other
         part of the item.
   \item This is the third item of the list.
\end{itemize}
You can even display poetry.
\begin{verse}
   There is an environment
    for verse \\             % The \\ command separates lines
   Whose features some poets % within a stanza.
   will curse.

                             % One or more blank lines separate stanzas.

   For instead of making\\
   Them do \emph{all} line breaking, \\
   It allows them to put too many words on a line when they'd rather be
   forced to be terse.
\end{verse}

Mathematical formulas may also be displayed.  A
displayed formula
is
one-line long; multiline
formulas require special formatting instructions.
   \[  \ip{\Gamma}{\psi'} = x'' + y^{2} + z_{i}^{n}\]
Don't start a paragraph with a displayed equation,
nor make one a paragraph by itself.

\end{document}               % End of document.

\section{Bench Huge Section 1347}
% synthetic bench section 1347
% This is a sample LaTeX input file.  (Version of 12 August 2004.)
%
% A '%' character causes TeX to ignore all remaining text on the line,
% and is used for comments like this one.

\documentclass{article}      % Specifies the document class

                             % The preamble begins here.
\title{An Example Document}  % Declares the document's title.
\author{Leslie Lamport}      % Declares the author's name.
\date{January 21, 1994}      % Deleting this command produces today's date.

\newcommand{\ip}[2]{(#1, #2)}
                             % Defines \ip{arg1}{arg2} to mean
                             % (arg1, arg2).

%\newcommand{\ip}[2]{\langle #1 | #2\rangle}
                             % This is an alternative definition of
                             % \ip that is commented out.

\begin{document}             % End of preamble and beginning of text.

\maketitle                   % Produces the title.

This is an example input file.  Comparing it with
the output it generates can show you how to
produce a simple document of your own.

\section{Ordinary Text}      % Produces section heading.  Lower-level
                             % sections are begun with similar
                             % \subsection and \subsubsection commands.

The ends  of words and sentences are marked
  by   spaces. It  doesn't matter how many
spaces    you type; one is as good as 100.  The
end of   a line counts as a space.

One   or more   blank lines denote the  end
of  a paragraph.

Since any number of consecutive spaces are treated
like a single one, the formatting of the input
file makes no difference to
      \LaTeX,                % The \LaTeX command generates the LaTeX logo.
but it makes a difference to you.  When you use
\LaTeX, making your input file as easy to read
as possible will be a great help as you write
your document and when you change it.  This sample
file shows how you can add comments to your own input
file.

Because printing is different from typewriting,
there are a number of things that you have to do
differently when preparing an input file than if
you were just typing the document directly.
Quotation marks like
       ``this''
have to be handled specially, as do quotes within
quotes:
       ``\,`this'            % \, separates the double and single quote.
        is what I just
        wrote, not  `that'\,''.

Dashes come in three sizes: an
       intra-word
dash, a medium dash for number ranges like
       1--2,
and a punctuation
       dash---like
this.

A sentence-ending space should be larger than the
space between words within a sentence.  You
sometimes have to type special commands in
conjunction with punctuation characters to get
this right, as in the following sentence.
       Gnats, gnus, etc.\ all  % `\ ' makes an inter-word space.
       begin with G\@.         % \@ marks end-of-sentence punctuation.
You should check the spaces after periods when
reading your output to make sure you haven't
forgotten any special cases.  Generating an
ellipsis
       \ldots\               % `\ ' is needed after `\ldots' because TeX
                             % ignores spaces after command names like \ldots
                             % made from \ + letters.
                             %
                             % Note how a `%' character causes TeX to ignore
                             % the end of the input line, so these blank lines
                             % do not start a new paragraph.
                             %
with the right spacing around the periods requires
a special command.

\LaTeX\ interprets some common characters as
commands, so you must type special commands to
generate them.  These characters include the
following:
       \$ \& \% \# \{ and \}.

In printing, text is usually emphasized with an
       \emph{italic}
type style.

\begin{em}
   A long segment of text can also be emphasized
   in this way.  Text within such a segment can be
   given \emph{additional} emphasis.
\end{em}

It is sometimes necessary to prevent \LaTeX\ from
breaking a line where it might otherwise do so.
This may be at a space, as between the ``Mr.''\ and
``Jones'' in
       ``Mr.~Jones'',        % ~ produces an unbreakable interword space.
or within a word---especially when the word is a
symbol like
       \mbox{\emph{itemnum}}
that makes little sense when hyphenated across
lines.

Footnotes\footnote{This is an example of a footnote.}
pose no problem.

\LaTeX\ is good at typesetting mathematical formulas
like
       \( x-3y + z = 7 \)
or
       \( a_{1} > x^{2n} + y^{2n} > x' \)
or
       \( \ip{A}{B} = \sum_{i} a_{i} b_{i} \).
The spaces you type in a formula are
ignored.  Remember that a letter like
       $x$                   % $ ... $  and  \( ... \)  are equivalent
is a formula when it denotes a mathematical
symbol, and it should be typed as one.

\section{Displayed Text}

Text is displayed by indenting it from the left
margin.  Quotations are commonly displayed.  There
are short quotations
\begin{quote}
   This is a short quotation.  It consists of a
   single paragraph of text.  See how it is formatted.
\end{quote}
and longer ones.
\begin{quotation}
   This is a longer quotation.  It consists of two
   paragraphs of text, neither of which are
   particularly interesting.

   This is the second paragraph of the quotation.  It
   is just as dull as the first paragraph.
\end{quotation}
Another frequently-displayed structure is a list.
The following is an example of an \emph{itemized}
list.
\begin{itemize}
   \item This is the first item of an itemized list.
         Each item in the list is marked with a ``tick''.
         You don't have to worry about what kind of tick
         mark is used.

   \item This is the second item of the list.  It
         contains another list nested inside it.  The inner
         list is an \emph{enumerated} list.
         \begin{enumerate}
            \item This is the first item of an enumerated
                  list that is nested within the itemized list.

            \item This is the second item of the inner list.
                  \LaTeX\ allows you to nest lists deeper than
                  you really should.
         \end{enumerate}
         This is the rest of the second item of the outer
         list.  It is no more interesting than any other
         part of the item.
   \item This is the third item of the list.
\end{itemize}
You can even display poetry.
\begin{verse}
   There is an environment
    for verse \\             % The \\ command separates lines
   Whose features some poets % within a stanza.
   will curse.

                             % One or more blank lines separate stanzas.

   For instead of making\\
   Them do \emph{all} line breaking, \\
   It allows them to put too many words on a line when they'd rather be
   forced to be terse.
\end{verse}

Mathematical formulas may also be displayed.  A
displayed formula
is
one-line long; multiline
formulas require special formatting instructions.
   \[  \ip{\Gamma}{\psi'} = x'' + y^{2} + z_{i}^{n}\]
Don't start a paragraph with a displayed equation,
nor make one a paragraph by itself.

\end{document}               % End of document.

\section{Bench Huge Section 1348}
% synthetic bench section 1348
% This is a sample LaTeX input file.  (Version of 12 August 2004.)
%
% A '%' character causes TeX to ignore all remaining text on the line,
% and is used for comments like this one.

\documentclass{article}      % Specifies the document class

                             % The preamble begins here.
\title{An Example Document}  % Declares the document's title.
\author{Leslie Lamport}      % Declares the author's name.
\date{January 21, 1994}      % Deleting this command produces today's date.

\newcommand{\ip}[2]{(#1, #2)}
                             % Defines \ip{arg1}{arg2} to mean
                             % (arg1, arg2).

%\newcommand{\ip}[2]{\langle #1 | #2\rangle}
                             % This is an alternative definition of
                             % \ip that is commented out.

\begin{document}             % End of preamble and beginning of text.

\maketitle                   % Produces the title.

This is an example input file.  Comparing it with
the output it generates can show you how to
produce a simple document of your own.

\section{Ordinary Text}      % Produces section heading.  Lower-level
                             % sections are begun with similar
                             % \subsection and \subsubsection commands.

The ends  of words and sentences are marked
  by   spaces. It  doesn't matter how many
spaces    you type; one is as good as 100.  The
end of   a line counts as a space.

One   or more   blank lines denote the  end
of  a paragraph.

Since any number of consecutive spaces are treated
like a single one, the formatting of the input
file makes no difference to
      \LaTeX,                % The \LaTeX command generates the LaTeX logo.
but it makes a difference to you.  When you use
\LaTeX, making your input file as easy to read
as possible will be a great help as you write
your document and when you change it.  This sample
file shows how you can add comments to your own input
file.

Because printing is different from typewriting,
there are a number of things that you have to do
differently when preparing an input file than if
you were just typing the document directly.
Quotation marks like
       ``this''
have to be handled specially, as do quotes within
quotes:
       ``\,`this'            % \, separates the double and single quote.
        is what I just
        wrote, not  `that'\,''.

Dashes come in three sizes: an
       intra-word
dash, a medium dash for number ranges like
       1--2,
and a punctuation
       dash---like
this.

A sentence-ending space should be larger than the
space between words within a sentence.  You
sometimes have to type special commands in
conjunction with punctuation characters to get
this right, as in the following sentence.
       Gnats, gnus, etc.\ all  % `\ ' makes an inter-word space.
       begin with G\@.         % \@ marks end-of-sentence punctuation.
You should check the spaces after periods when
reading your output to make sure you haven't
forgotten any special cases.  Generating an
ellipsis
       \ldots\               % `\ ' is needed after `\ldots' because TeX
                             % ignores spaces after command names like \ldots
                             % made from \ + letters.
                             %
                             % Note how a `%' character causes TeX to ignore
                             % the end of the input line, so these blank lines
                             % do not start a new paragraph.
                             %
with the right spacing around the periods requires
a special command.

\LaTeX\ interprets some common characters as
commands, so you must type special commands to
generate them.  These characters include the
following:
       \$ \& \% \# \{ and \}.

In printing, text is usually emphasized with an
       \emph{italic}
type style.

\begin{em}
   A long segment of text can also be emphasized
   in this way.  Text within such a segment can be
   given \emph{additional} emphasis.
\end{em}

It is sometimes necessary to prevent \LaTeX\ from
breaking a line where it might otherwise do so.
This may be at a space, as between the ``Mr.''\ and
``Jones'' in
       ``Mr.~Jones'',        % ~ produces an unbreakable interword space.
or within a word---especially when the word is a
symbol like
       \mbox{\emph{itemnum}}
that makes little sense when hyphenated across
lines.

Footnotes\footnote{This is an example of a footnote.}
pose no problem.

\LaTeX\ is good at typesetting mathematical formulas
like
       \( x-3y + z = 7 \)
or
       \( a_{1} > x^{2n} + y^{2n} > x' \)
or
       \( \ip{A}{B} = \sum_{i} a_{i} b_{i} \).
The spaces you type in a formula are
ignored.  Remember that a letter like
       $x$                   % $ ... $  and  \( ... \)  are equivalent
is a formula when it denotes a mathematical
symbol, and it should be typed as one.

\section{Displayed Text}

Text is displayed by indenting it from the left
margin.  Quotations are commonly displayed.  There
are short quotations
\begin{quote}
   This is a short quotation.  It consists of a
   single paragraph of text.  See how it is formatted.
\end{quote}
and longer ones.
\begin{quotation}
   This is a longer quotation.  It consists of two
   paragraphs of text, neither of which are
   particularly interesting.

   This is the second paragraph of the quotation.  It
   is just as dull as the first paragraph.
\end{quotation}
Another frequently-displayed structure is a list.
The following is an example of an \emph{itemized}
list.
\begin{itemize}
   \item This is the first item of an itemized list.
         Each item in the list is marked with a ``tick''.
         You don't have to worry about what kind of tick
         mark is used.

   \item This is the second item of the list.  It
         contains another list nested inside it.  The inner
         list is an \emph{enumerated} list.
         \begin{enumerate}
            \item This is the first item of an enumerated
                  list that is nested within the itemized list.

            \item This is the second item of the inner list.
                  \LaTeX\ allows you to nest lists deeper than
                  you really should.
         \end{enumerate}
         This is the rest of the second item of the outer
         list.  It is no more interesting than any other
         part of the item.
   \item This is the third item of the list.
\end{itemize}
You can even display poetry.
\begin{verse}
   There is an environment
    for verse \\             % The \\ command separates lines
   Whose features some poets % within a stanza.
   will curse.

                             % One or more blank lines separate stanzas.

   For instead of making\\
   Them do \emph{all} line breaking, \\
   It allows them to put too many words on a line when they'd rather be
   forced to be terse.
\end{verse}

Mathematical formulas may also be displayed.  A
displayed formula
is
one-line long; multiline
formulas require special formatting instructions.
   \[  \ip{\Gamma}{\psi'} = x'' + y^{2} + z_{i}^{n}\]
Don't start a paragraph with a displayed equation,
nor make one a paragraph by itself.

\end{document}               % End of document.

\section{Bench Huge Section 1349}
% synthetic bench section 1349
% This is a sample LaTeX input file.  (Version of 12 August 2004.)
%
% A '%' character causes TeX to ignore all remaining text on the line,
% and is used for comments like this one.

\documentclass{article}      % Specifies the document class

                             % The preamble begins here.
\title{An Example Document}  % Declares the document's title.
\author{Leslie Lamport}      % Declares the author's name.
\date{January 21, 1994}      % Deleting this command produces today's date.

\newcommand{\ip}[2]{(#1, #2)}
                             % Defines \ip{arg1}{arg2} to mean
                             % (arg1, arg2).

%\newcommand{\ip}[2]{\langle #1 | #2\rangle}
                             % This is an alternative definition of
                             % \ip that is commented out.

\begin{document}             % End of preamble and beginning of text.

\maketitle                   % Produces the title.

This is an example input file.  Comparing it with
the output it generates can show you how to
produce a simple document of your own.

\section{Ordinary Text}      % Produces section heading.  Lower-level
                             % sections are begun with similar
                             % \subsection and \subsubsection commands.

The ends  of words and sentences are marked
  by   spaces. It  doesn't matter how many
spaces    you type; one is as good as 100.  The
end of   a line counts as a space.

One   or more   blank lines denote the  end
of  a paragraph.

Since any number of consecutive spaces are treated
like a single one, the formatting of the input
file makes no difference to
      \LaTeX,                % The \LaTeX command generates the LaTeX logo.
but it makes a difference to you.  When you use
\LaTeX, making your input file as easy to read
as possible will be a great help as you write
your document and when you change it.  This sample
file shows how you can add comments to your own input
file.

Because printing is different from typewriting,
there are a number of things that you have to do
differently when preparing an input file than if
you were just typing the document directly.
Quotation marks like
       ``this''
have to be handled specially, as do quotes within
quotes:
       ``\,`this'            % \, separates the double and single quote.
        is what I just
        wrote, not  `that'\,''.

Dashes come in three sizes: an
       intra-word
dash, a medium dash for number ranges like
       1--2,
and a punctuation
       dash---like
this.

A sentence-ending space should be larger than the
space between words within a sentence.  You
sometimes have to type special commands in
conjunction with punctuation characters to get
this right, as in the following sentence.
       Gnats, gnus, etc.\ all  % `\ ' makes an inter-word space.
       begin with G\@.         % \@ marks end-of-sentence punctuation.
You should check the spaces after periods when
reading your output to make sure you haven't
forgotten any special cases.  Generating an
ellipsis
       \ldots\               % `\ ' is needed after `\ldots' because TeX
                             % ignores spaces after command names like \ldots
                             % made from \ + letters.
                             %
                             % Note how a `%' character causes TeX to ignore
                             % the end of the input line, so these blank lines
                             % do not start a new paragraph.
                             %
with the right spacing around the periods requires
a special command.

\LaTeX\ interprets some common characters as
commands, so you must type special commands to
generate them.  These characters include the
following:
       \$ \& \% \# \{ and \}.

In printing, text is usually emphasized with an
       \emph{italic}
type style.

\begin{em}
   A long segment of text can also be emphasized
   in this way.  Text within such a segment can be
   given \emph{additional} emphasis.
\end{em}

It is sometimes necessary to prevent \LaTeX\ from
breaking a line where it might otherwise do so.
This may be at a space, as between the ``Mr.''\ and
``Jones'' in
       ``Mr.~Jones'',        % ~ produces an unbreakable interword space.
or within a word---especially when the word is a
symbol like
       \mbox{\emph{itemnum}}
that makes little sense when hyphenated across
lines.

Footnotes\footnote{This is an example of a footnote.}
pose no problem.

\LaTeX\ is good at typesetting mathematical formulas
like
       \( x-3y + z = 7 \)
or
       \( a_{1} > x^{2n} + y^{2n} > x' \)
or
       \( \ip{A}{B} = \sum_{i} a_{i} b_{i} \).
The spaces you type in a formula are
ignored.  Remember that a letter like
       $x$                   % $ ... $  and  \( ... \)  are equivalent
is a formula when it denotes a mathematical
symbol, and it should be typed as one.

\section{Displayed Text}

Text is displayed by indenting it from the left
margin.  Quotations are commonly displayed.  There
are short quotations
\begin{quote}
   This is a short quotation.  It consists of a
   single paragraph of text.  See how it is formatted.
\end{quote}
and longer ones.
\begin{quotation}
   This is a longer quotation.  It consists of two
   paragraphs of text, neither of which are
   particularly interesting.

   This is the second paragraph of the quotation.  It
   is just as dull as the first paragraph.
\end{quotation}
Another frequently-displayed structure is a list.
The following is an example of an \emph{itemized}
list.
\begin{itemize}
   \item This is the first item of an itemized list.
         Each item in the list is marked with a ``tick''.
         You don't have to worry about what kind of tick
         mark is used.

   \item This is the second item of the list.  It
         contains another list nested inside it.  The inner
         list is an \emph{enumerated} list.
         \begin{enumerate}
            \item This is the first item of an enumerated
                  list that is nested within the itemized list.

            \item This is the second item of the inner list.
                  \LaTeX\ allows you to nest lists deeper than
                  you really should.
         \end{enumerate}
         This is the rest of the second item of the outer
         list.  It is no more interesting than any other
         part of the item.
   \item This is the third item of the list.
\end{itemize}
You can even display poetry.
\begin{verse}
   There is an environment
    for verse \\             % The \\ command separates lines
   Whose features some poets % within a stanza.
   will curse.

                             % One or more blank lines separate stanzas.

   For instead of making\\
   Them do \emph{all} line breaking, \\
   It allows them to put too many words on a line when they'd rather be
   forced to be terse.
\end{verse}

Mathematical formulas may also be displayed.  A
displayed formula
is
one-line long; multiline
formulas require special formatting instructions.
   \[  \ip{\Gamma}{\psi'} = x'' + y^{2} + z_{i}^{n}\]
Don't start a paragraph with a displayed equation,
nor make one a paragraph by itself.

\end{document}               % End of document.

\section{Bench Huge Section 1350}
% synthetic bench section 1350
% This is a sample LaTeX input file.  (Version of 12 August 2004.)
%
% A '%' character causes TeX to ignore all remaining text on the line,
% and is used for comments like this one.

\documentclass{article}      % Specifies the document class

                             % The preamble begins here.
\title{An Example Document}  % Declares the document's title.
\author{Leslie Lamport}      % Declares the author's name.
\date{January 21, 1994}      % Deleting this command produces today's date.

\newcommand{\ip}[2]{(#1, #2)}
                             % Defines \ip{arg1}{arg2} to mean
                             % (arg1, arg2).

%\newcommand{\ip}[2]{\langle #1 | #2\rangle}
                             % This is an alternative definition of
                             % \ip that is commented out.

\begin{document}             % End of preamble and beginning of text.

\maketitle                   % Produces the title.

This is an example input file.  Comparing it with
the output it generates can show you how to
produce a simple document of your own.

\section{Ordinary Text}      % Produces section heading.  Lower-level
                             % sections are begun with similar
                             % \subsection and \subsubsection commands.

The ends  of words and sentences are marked
  by   spaces. It  doesn't matter how many
spaces    you type; one is as good as 100.  The
end of   a line counts as a space.

One   or more   blank lines denote the  end
of  a paragraph.

Since any number of consecutive spaces are treated
like a single one, the formatting of the input
file makes no difference to
      \LaTeX,                % The \LaTeX command generates the LaTeX logo.
but it makes a difference to you.  When you use
\LaTeX, making your input file as easy to read
as possible will be a great help as you write
your document and when you change it.  This sample
file shows how you can add comments to your own input
file.

Because printing is different from typewriting,
there are a number of things that you have to do
differently when preparing an input file than if
you were just typing the document directly.
Quotation marks like
       ``this''
have to be handled specially, as do quotes within
quotes:
       ``\,`this'            % \, separates the double and single quote.
        is what I just
        wrote, not  `that'\,''.

Dashes come in three sizes: an
       intra-word
dash, a medium dash for number ranges like
       1--2,
and a punctuation
       dash---like
this.

A sentence-ending space should be larger than the
space between words within a sentence.  You
sometimes have to type special commands in
conjunction with punctuation characters to get
this right, as in the following sentence.
       Gnats, gnus, etc.\ all  % `\ ' makes an inter-word space.
       begin with G\@.         % \@ marks end-of-sentence punctuation.
You should check the spaces after periods when
reading your output to make sure you haven't
forgotten any special cases.  Generating an
ellipsis
       \ldots\               % `\ ' is needed after `\ldots' because TeX
                             % ignores spaces after command names like \ldots
                             % made from \ + letters.
                             %
                             % Note how a `%' character causes TeX to ignore
                             % the end of the input line, so these blank lines
                             % do not start a new paragraph.
                             %
with the right spacing around the periods requires
a special command.

\LaTeX\ interprets some common characters as
commands, so you must type special commands to
generate them.  These characters include the
following:
       \$ \& \% \# \{ and \}.

In printing, text is usually emphasized with an
       \emph{italic}
type style.

\begin{em}
   A long segment of text can also be emphasized
   in this way.  Text within such a segment can be
   given \emph{additional} emphasis.
\end{em}

It is sometimes necessary to prevent \LaTeX\ from
breaking a line where it might otherwise do so.
This may be at a space, as between the ``Mr.''\ and
``Jones'' in
       ``Mr.~Jones'',        % ~ produces an unbreakable interword space.
or within a word---especially when the word is a
symbol like
       \mbox{\emph{itemnum}}
that makes little sense when hyphenated across
lines.

Footnotes\footnote{This is an example of a footnote.}
pose no problem.

\LaTeX\ is good at typesetting mathematical formulas
like
       \( x-3y + z = 7 \)
or
       \( a_{1} > x^{2n} + y^{2n} > x' \)
or
       \( \ip{A}{B} = \sum_{i} a_{i} b_{i} \).
The spaces you type in a formula are
ignored.  Remember that a letter like
       $x$                   % $ ... $  and  \( ... \)  are equivalent
is a formula when it denotes a mathematical
symbol, and it should be typed as one.

\section{Displayed Text}

Text is displayed by indenting it from the left
margin.  Quotations are commonly displayed.  There
are short quotations
\begin{quote}
   This is a short quotation.  It consists of a
   single paragraph of text.  See how it is formatted.
\end{quote}
and longer ones.
\begin{quotation}
   This is a longer quotation.  It consists of two
   paragraphs of text, neither of which are
   particularly interesting.

   This is the second paragraph of the quotation.  It
   is just as dull as the first paragraph.
\end{quotation}
Another frequently-displayed structure is a list.
The following is an example of an \emph{itemized}
list.
\begin{itemize}
   \item This is the first item of an itemized list.
         Each item in the list is marked with a ``tick''.
         You don't have to worry about what kind of tick
         mark is used.

   \item This is the second item of the list.  It
         contains another list nested inside it.  The inner
         list is an \emph{enumerated} list.
         \begin{enumerate}
            \item This is the first item of an enumerated
                  list that is nested within the itemized list.

            \item This is the second item of the inner list.
                  \LaTeX\ allows you to nest lists deeper than
                  you really should.
         \end{enumerate}
         This is the rest of the second item of the outer
         list.  It is no more interesting than any other
         part of the item.
   \item This is the third item of the list.
\end{itemize}
You can even display poetry.
\begin{verse}
   There is an environment
    for verse \\             % The \\ command separates lines
   Whose features some poets % within a stanza.
   will curse.

                             % One or more blank lines separate stanzas.

   For instead of making\\
   Them do \emph{all} line breaking, \\
   It allows them to put too many words on a line when they'd rather be
   forced to be terse.
\end{verse}

Mathematical formulas may also be displayed.  A
displayed formula
is
one-line long; multiline
formulas require special formatting instructions.
   \[  \ip{\Gamma}{\psi'} = x'' + y^{2} + z_{i}^{n}\]
Don't start a paragraph with a displayed equation,
nor make one a paragraph by itself.

\end{document}               % End of document.

\section{Bench Huge Section 1351}
% synthetic bench section 1351
% This is a sample LaTeX input file.  (Version of 12 August 2004.)
%
% A '%' character causes TeX to ignore all remaining text on the line,
% and is used for comments like this one.

\documentclass{article}      % Specifies the document class

                             % The preamble begins here.
\title{An Example Document}  % Declares the document's title.
\author{Leslie Lamport}      % Declares the author's name.
\date{January 21, 1994}      % Deleting this command produces today's date.

\newcommand{\ip}[2]{(#1, #2)}
                             % Defines \ip{arg1}{arg2} to mean
                             % (arg1, arg2).

%\newcommand{\ip}[2]{\langle #1 | #2\rangle}
                             % This is an alternative definition of
                             % \ip that is commented out.

\begin{document}             % End of preamble and beginning of text.

\maketitle                   % Produces the title.

This is an example input file.  Comparing it with
the output it generates can show you how to
produce a simple document of your own.

\section{Ordinary Text}      % Produces section heading.  Lower-level
                             % sections are begun with similar
                             % \subsection and \subsubsection commands.

The ends  of words and sentences are marked
  by   spaces. It  doesn't matter how many
spaces    you type; one is as good as 100.  The
end of   a line counts as a space.

One   or more   blank lines denote the  end
of  a paragraph.

Since any number of consecutive spaces are treated
like a single one, the formatting of the input
file makes no difference to
      \LaTeX,                % The \LaTeX command generates the LaTeX logo.
but it makes a difference to you.  When you use
\LaTeX, making your input file as easy to read
as possible will be a great help as you write
your document and when you change it.  This sample
file shows how you can add comments to your own input
file.

Because printing is different from typewriting,
there are a number of things that you have to do
differently when preparing an input file than if
you were just typing the document directly.
Quotation marks like
       ``this''
have to be handled specially, as do quotes within
quotes:
       ``\,`this'            % \, separates the double and single quote.
        is what I just
        wrote, not  `that'\,''.

Dashes come in three sizes: an
       intra-word
dash, a medium dash for number ranges like
       1--2,
and a punctuation
       dash---like
this.

A sentence-ending space should be larger than the
space between words within a sentence.  You
sometimes have to type special commands in
conjunction with punctuation characters to get
this right, as in the following sentence.
       Gnats, gnus, etc.\ all  % `\ ' makes an inter-word space.
       begin with G\@.         % \@ marks end-of-sentence punctuation.
You should check the spaces after periods when
reading your output to make sure you haven't
forgotten any special cases.  Generating an
ellipsis
       \ldots\               % `\ ' is needed after `\ldots' because TeX
                             % ignores spaces after command names like \ldots
                             % made from \ + letters.
                             %
                             % Note how a `%' character causes TeX to ignore
                             % the end of the input line, so these blank lines
                             % do not start a new paragraph.
                             %
with the right spacing around the periods requires
a special command.

\LaTeX\ interprets some common characters as
commands, so you must type special commands to
generate them.  These characters include the
following:
       \$ \& \% \# \{ and \}.

In printing, text is usually emphasized with an
       \emph{italic}
type style.

\begin{em}
   A long segment of text can also be emphasized
   in this way.  Text within such a segment can be
   given \emph{additional} emphasis.
\end{em}

It is sometimes necessary to prevent \LaTeX\ from
breaking a line where it might otherwise do so.
This may be at a space, as between the ``Mr.''\ and
``Jones'' in
       ``Mr.~Jones'',        % ~ produces an unbreakable interword space.
or within a word---especially when the word is a
symbol like
       \mbox{\emph{itemnum}}
that makes little sense when hyphenated across
lines.

Footnotes\footnote{This is an example of a footnote.}
pose no problem.

\LaTeX\ is good at typesetting mathematical formulas
like
       \( x-3y + z = 7 \)
or
       \( a_{1} > x^{2n} + y^{2n} > x' \)
or
       \( \ip{A}{B} = \sum_{i} a_{i} b_{i} \).
The spaces you type in a formula are
ignored.  Remember that a letter like
       $x$                   % $ ... $  and  \( ... \)  are equivalent
is a formula when it denotes a mathematical
symbol, and it should be typed as one.

\section{Displayed Text}

Text is displayed by indenting it from the left
margin.  Quotations are commonly displayed.  There
are short quotations
\begin{quote}
   This is a short quotation.  It consists of a
   single paragraph of text.  See how it is formatted.
\end{quote}
and longer ones.
\begin{quotation}
   This is a longer quotation.  It consists of two
   paragraphs of text, neither of which are
   particularly interesting.

   This is the second paragraph of the quotation.  It
   is just as dull as the first paragraph.
\end{quotation}
Another frequently-displayed structure is a list.
The following is an example of an \emph{itemized}
list.
\begin{itemize}
   \item This is the first item of an itemized list.
         Each item in the list is marked with a ``tick''.
         You don't have to worry about what kind of tick
         mark is used.

   \item This is the second item of the list.  It
         contains another list nested inside it.  The inner
         list is an \emph{enumerated} list.
         \begin{enumerate}
            \item This is the first item of an enumerated
                  list that is nested within the itemized list.

            \item This is the second item of the inner list.
                  \LaTeX\ allows you to nest lists deeper than
                  you really should.
         \end{enumerate}
         This is the rest of the second item of the outer
         list.  It is no more interesting than any other
         part of the item.
   \item This is the third item of the list.
\end{itemize}
You can even display poetry.
\begin{verse}
   There is an environment
    for verse \\             % The \\ command separates lines
   Whose features some poets % within a stanza.
   will curse.

                             % One or more blank lines separate stanzas.

   For instead of making\\
   Them do \emph{all} line breaking, \\
   It allows them to put too many words on a line when they'd rather be
   forced to be terse.
\end{verse}

Mathematical formulas may also be displayed.  A
displayed formula
is
one-line long; multiline
formulas require special formatting instructions.
   \[  \ip{\Gamma}{\psi'} = x'' + y^{2} + z_{i}^{n}\]
Don't start a paragraph with a displayed equation,
nor make one a paragraph by itself.

\end{document}               % End of document.

\section{Bench Huge Section 1352}
% synthetic bench section 1352
% This is a sample LaTeX input file.  (Version of 12 August 2004.)
%
% A '%' character causes TeX to ignore all remaining text on the line,
% and is used for comments like this one.

\documentclass{article}      % Specifies the document class

                             % The preamble begins here.
\title{An Example Document}  % Declares the document's title.
\author{Leslie Lamport}      % Declares the author's name.
\date{January 21, 1994}      % Deleting this command produces today's date.

\newcommand{\ip}[2]{(#1, #2)}
                             % Defines \ip{arg1}{arg2} to mean
                             % (arg1, arg2).

%\newcommand{\ip}[2]{\langle #1 | #2\rangle}
                             % This is an alternative definition of
                             % \ip that is commented out.

\begin{document}             % End of preamble and beginning of text.

\maketitle                   % Produces the title.

This is an example input file.  Comparing it with
the output it generates can show you how to
produce a simple document of your own.

\section{Ordinary Text}      % Produces section heading.  Lower-level
                             % sections are begun with similar
                             % \subsection and \subsubsection commands.

The ends  of words and sentences are marked
  by   spaces. It  doesn't matter how many
spaces    you type; one is as good as 100.  The
end of   a line counts as a space.

One   or more   blank lines denote the  end
of  a paragraph.

Since any number of consecutive spaces are treated
like a single one, the formatting of the input
file makes no difference to
      \LaTeX,                % The \LaTeX command generates the LaTeX logo.
but it makes a difference to you.  When you use
\LaTeX, making your input file as easy to read
as possible will be a great help as you write
your document and when you change it.  This sample
file shows how you can add comments to your own input
file.

Because printing is different from typewriting,
there are a number of things that you have to do
differently when preparing an input file than if
you were just typing the document directly.
Quotation marks like
       ``this''
have to be handled specially, as do quotes within
quotes:
       ``\,`this'            % \, separates the double and single quote.
        is what I just
        wrote, not  `that'\,''.

Dashes come in three sizes: an
       intra-word
dash, a medium dash for number ranges like
       1--2,
and a punctuation
       dash---like
this.

A sentence-ending space should be larger than the
space between words within a sentence.  You
sometimes have to type special commands in
conjunction with punctuation characters to get
this right, as in the following sentence.
       Gnats, gnus, etc.\ all  % `\ ' makes an inter-word space.
       begin with G\@.         % \@ marks end-of-sentence punctuation.
You should check the spaces after periods when
reading your output to make sure you haven't
forgotten any special cases.  Generating an
ellipsis
       \ldots\               % `\ ' is needed after `\ldots' because TeX
                             % ignores spaces after command names like \ldots
                             % made from \ + letters.
                             %
                             % Note how a `%' character causes TeX to ignore
                             % the end of the input line, so these blank lines
                             % do not start a new paragraph.
                             %
with the right spacing around the periods requires
a special command.

\LaTeX\ interprets some common characters as
commands, so you must type special commands to
generate them.  These characters include the
following:
       \$ \& \% \# \{ and \}.

In printing, text is usually emphasized with an
       \emph{italic}
type style.

\begin{em}
   A long segment of text can also be emphasized
   in this way.  Text within such a segment can be
   given \emph{additional} emphasis.
\end{em}

It is sometimes necessary to prevent \LaTeX\ from
breaking a line where it might otherwise do so.
This may be at a space, as between the ``Mr.''\ and
``Jones'' in
       ``Mr.~Jones'',        % ~ produces an unbreakable interword space.
or within a word---especially when the word is a
symbol like
       \mbox{\emph{itemnum}}
that makes little sense when hyphenated across
lines.

Footnotes\footnote{This is an example of a footnote.}
pose no problem.

\LaTeX\ is good at typesetting mathematical formulas
like
       \( x-3y + z = 7 \)
or
       \( a_{1} > x^{2n} + y^{2n} > x' \)
or
       \( \ip{A}{B} = \sum_{i} a_{i} b_{i} \).
The spaces you type in a formula are
ignored.  Remember that a letter like
       $x$                   % $ ... $  and  \( ... \)  are equivalent
is a formula when it denotes a mathematical
symbol, and it should be typed as one.

\section{Displayed Text}

Text is displayed by indenting it from the left
margin.  Quotations are commonly displayed.  There
are short quotations
\begin{quote}
   This is a short quotation.  It consists of a
   single paragraph of text.  See how it is formatted.
\end{quote}
and longer ones.
\begin{quotation}
   This is a longer quotation.  It consists of two
   paragraphs of text, neither of which are
   particularly interesting.

   This is the second paragraph of the quotation.  It
   is just as dull as the first paragraph.
\end{quotation}
Another frequently-displayed structure is a list.
The following is an example of an \emph{itemized}
list.
\begin{itemize}
   \item This is the first item of an itemized list.
         Each item in the list is marked with a ``tick''.
         You don't have to worry about what kind of tick
         mark is used.

   \item This is the second item of the list.  It
         contains another list nested inside it.  The inner
         list is an \emph{enumerated} list.
         \begin{enumerate}
            \item This is the first item of an enumerated
                  list that is nested within the itemized list.

            \item This is the second item of the inner list.
                  \LaTeX\ allows you to nest lists deeper than
                  you really should.
         \end{enumerate}
         This is the rest of the second item of the outer
         list.  It is no more interesting than any other
         part of the item.
   \item This is the third item of the list.
\end{itemize}
You can even display poetry.
\begin{verse}
   There is an environment
    for verse \\             % The \\ command separates lines
   Whose features some poets % within a stanza.
   will curse.

                             % One or more blank lines separate stanzas.

   For instead of making\\
   Them do \emph{all} line breaking, \\
   It allows them to put too many words on a line when they'd rather be
   forced to be terse.
\end{verse}

Mathematical formulas may also be displayed.  A
displayed formula
is
one-line long; multiline
formulas require special formatting instructions.
   \[  \ip{\Gamma}{\psi'} = x'' + y^{2} + z_{i}^{n}\]
Don't start a paragraph with a displayed equation,
nor make one a paragraph by itself.

\end{document}               % End of document.

\section{Bench Huge Section 1353}
% synthetic bench section 1353
% This is a sample LaTeX input file.  (Version of 12 August 2004.)
%
% A '%' character causes TeX to ignore all remaining text on the line,
% and is used for comments like this one.

\documentclass{article}      % Specifies the document class

                             % The preamble begins here.
\title{An Example Document}  % Declares the document's title.
\author{Leslie Lamport}      % Declares the author's name.
\date{January 21, 1994}      % Deleting this command produces today's date.

\newcommand{\ip}[2]{(#1, #2)}
                             % Defines \ip{arg1}{arg2} to mean
                             % (arg1, arg2).

%\newcommand{\ip}[2]{\langle #1 | #2\rangle}
                             % This is an alternative definition of
                             % \ip that is commented out.

\begin{document}             % End of preamble and beginning of text.

\maketitle                   % Produces the title.

This is an example input file.  Comparing it with
the output it generates can show you how to
produce a simple document of your own.

\section{Ordinary Text}      % Produces section heading.  Lower-level
                             % sections are begun with similar
                             % \subsection and \subsubsection commands.

The ends  of words and sentences are marked
  by   spaces. It  doesn't matter how many
spaces    you type; one is as good as 100.  The
end of   a line counts as a space.

One   or more   blank lines denote the  end
of  a paragraph.

Since any number of consecutive spaces are treated
like a single one, the formatting of the input
file makes no difference to
      \LaTeX,                % The \LaTeX command generates the LaTeX logo.
but it makes a difference to you.  When you use
\LaTeX, making your input file as easy to read
as possible will be a great help as you write
your document and when you change it.  This sample
file shows how you can add comments to your own input
file.

Because printing is different from typewriting,
there are a number of things that you have to do
differently when preparing an input file than if
you were just typing the document directly.
Quotation marks like
       ``this''
have to be handled specially, as do quotes within
quotes:
       ``\,`this'            % \, separates the double and single quote.
        is what I just
        wrote, not  `that'\,''.

Dashes come in three sizes: an
       intra-word
dash, a medium dash for number ranges like
       1--2,
and a punctuation
       dash---like
this.

A sentence-ending space should be larger than the
space between words within a sentence.  You
sometimes have to type special commands in
conjunction with punctuation characters to get
this right, as in the following sentence.
       Gnats, gnus, etc.\ all  % `\ ' makes an inter-word space.
       begin with G\@.         % \@ marks end-of-sentence punctuation.
You should check the spaces after periods when
reading your output to make sure you haven't
forgotten any special cases.  Generating an
ellipsis
       \ldots\               % `\ ' is needed after `\ldots' because TeX
                             % ignores spaces after command names like \ldots
                             % made from \ + letters.
                             %
                             % Note how a `%' character causes TeX to ignore
                             % the end of the input line, so these blank lines
                             % do not start a new paragraph.
                             %
with the right spacing around the periods requires
a special command.

\LaTeX\ interprets some common characters as
commands, so you must type special commands to
generate them.  These characters include the
following:
       \$ \& \% \# \{ and \}.

In printing, text is usually emphasized with an
       \emph{italic}
type style.

\begin{em}
   A long segment of text can also be emphasized
   in this way.  Text within such a segment can be
   given \emph{additional} emphasis.
\end{em}

It is sometimes necessary to prevent \LaTeX\ from
breaking a line where it might otherwise do so.
This may be at a space, as between the ``Mr.''\ and
``Jones'' in
       ``Mr.~Jones'',        % ~ produces an unbreakable interword space.
or within a word---especially when the word is a
symbol like
       \mbox{\emph{itemnum}}
that makes little sense when hyphenated across
lines.

Footnotes\footnote{This is an example of a footnote.}
pose no problem.

\LaTeX\ is good at typesetting mathematical formulas
like
       \( x-3y + z = 7 \)
or
       \( a_{1} > x^{2n} + y^{2n} > x' \)
or
       \( \ip{A}{B} = \sum_{i} a_{i} b_{i} \).
The spaces you type in a formula are
ignored.  Remember that a letter like
       $x$                   % $ ... $  and  \( ... \)  are equivalent
is a formula when it denotes a mathematical
symbol, and it should be typed as one.

\section{Displayed Text}

Text is displayed by indenting it from the left
margin.  Quotations are commonly displayed.  There
are short quotations
\begin{quote}
   This is a short quotation.  It consists of a
   single paragraph of text.  See how it is formatted.
\end{quote}
and longer ones.
\begin{quotation}
   This is a longer quotation.  It consists of two
   paragraphs of text, neither of which are
   particularly interesting.

   This is the second paragraph of the quotation.  It
   is just as dull as the first paragraph.
\end{quotation}
Another frequently-displayed structure is a list.
The following is an example of an \emph{itemized}
list.
\begin{itemize}
   \item This is the first item of an itemized list.
         Each item in the list is marked with a ``tick''.
         You don't have to worry about what kind of tick
         mark is used.

   \item This is the second item of the list.  It
         contains another list nested inside it.  The inner
         list is an \emph{enumerated} list.
         \begin{enumerate}
            \item This is the first item of an enumerated
                  list that is nested within the itemized list.

            \item This is the second item of the inner list.
                  \LaTeX\ allows you to nest lists deeper than
                  you really should.
         \end{enumerate}
         This is the rest of the second item of the outer
         list.  It is no more interesting than any other
         part of the item.
   \item This is the third item of the list.
\end{itemize}
You can even display poetry.
\begin{verse}
   There is an environment
    for verse \\             % The \\ command separates lines
   Whose features some poets % within a stanza.
   will curse.

                             % One or more blank lines separate stanzas.

   For instead of making\\
   Them do \emph{all} line breaking, \\
   It allows them to put too many words on a line when they'd rather be
   forced to be terse.
\end{verse}

Mathematical formulas may also be displayed.  A
displayed formula
is
one-line long; multiline
formulas require special formatting instructions.
   \[  \ip{\Gamma}{\psi'} = x'' + y^{2} + z_{i}^{n}\]
Don't start a paragraph with a displayed equation,
nor make one a paragraph by itself.

\end{document}               % End of document.

\section{Bench Huge Section 1354}
% synthetic bench section 1354
% This is a sample LaTeX input file.  (Version of 12 August 2004.)
%
% A '%' character causes TeX to ignore all remaining text on the line,
% and is used for comments like this one.

\documentclass{article}      % Specifies the document class

                             % The preamble begins here.
\title{An Example Document}  % Declares the document's title.
\author{Leslie Lamport}      % Declares the author's name.
\date{January 21, 1994}      % Deleting this command produces today's date.

\newcommand{\ip}[2]{(#1, #2)}
                             % Defines \ip{arg1}{arg2} to mean
                             % (arg1, arg2).

%\newcommand{\ip}[2]{\langle #1 | #2\rangle}
                             % This is an alternative definition of
                             % \ip that is commented out.

\begin{document}             % End of preamble and beginning of text.

\maketitle                   % Produces the title.

This is an example input file.  Comparing it with
the output it generates can show you how to
produce a simple document of your own.

\section{Ordinary Text}      % Produces section heading.  Lower-level
                             % sections are begun with similar
                             % \subsection and \subsubsection commands.

The ends  of words and sentences are marked
  by   spaces. It  doesn't matter how many
spaces    you type; one is as good as 100.  The
end of   a line counts as a space.

One   or more   blank lines denote the  end
of  a paragraph.

Since any number of consecutive spaces are treated
like a single one, the formatting of the input
file makes no difference to
      \LaTeX,                % The \LaTeX command generates the LaTeX logo.
but it makes a difference to you.  When you use
\LaTeX, making your input file as easy to read
as possible will be a great help as you write
your document and when you change it.  This sample
file shows how you can add comments to your own input
file.

Because printing is different from typewriting,
there are a number of things that you have to do
differently when preparing an input file than if
you were just typing the document directly.
Quotation marks like
       ``this''
have to be handled specially, as do quotes within
quotes:
       ``\,`this'            % \, separates the double and single quote.
        is what I just
        wrote, not  `that'\,''.

Dashes come in three sizes: an
       intra-word
dash, a medium dash for number ranges like
       1--2,
and a punctuation
       dash---like
this.

A sentence-ending space should be larger than the
space between words within a sentence.  You
sometimes have to type special commands in
conjunction with punctuation characters to get
this right, as in the following sentence.
       Gnats, gnus, etc.\ all  % `\ ' makes an inter-word space.
       begin with G\@.         % \@ marks end-of-sentence punctuation.
You should check the spaces after periods when
reading your output to make sure you haven't
forgotten any special cases.  Generating an
ellipsis
       \ldots\               % `\ ' is needed after `\ldots' because TeX
                             % ignores spaces after command names like \ldots
                             % made from \ + letters.
                             %
                             % Note how a `%' character causes TeX to ignore
                             % the end of the input line, so these blank lines
                             % do not start a new paragraph.
                             %
with the right spacing around the periods requires
a special command.

\LaTeX\ interprets some common characters as
commands, so you must type special commands to
generate them.  These characters include the
following:
       \$ \& \% \# \{ and \}.

In printing, text is usually emphasized with an
       \emph{italic}
type style.

\begin{em}
   A long segment of text can also be emphasized
   in this way.  Text within such a segment can be
   given \emph{additional} emphasis.
\end{em}

It is sometimes necessary to prevent \LaTeX\ from
breaking a line where it might otherwise do so.
This may be at a space, as between the ``Mr.''\ and
``Jones'' in
       ``Mr.~Jones'',        % ~ produces an unbreakable interword space.
or within a word---especially when the word is a
symbol like
       \mbox{\emph{itemnum}}
that makes little sense when hyphenated across
lines.

Footnotes\footnote{This is an example of a footnote.}
pose no problem.

\LaTeX\ is good at typesetting mathematical formulas
like
       \( x-3y + z = 7 \)
or
       \( a_{1} > x^{2n} + y^{2n} > x' \)
or
       \( \ip{A}{B} = \sum_{i} a_{i} b_{i} \).
The spaces you type in a formula are
ignored.  Remember that a letter like
       $x$                   % $ ... $  and  \( ... \)  are equivalent
is a formula when it denotes a mathematical
symbol, and it should be typed as one.

\section{Displayed Text}

Text is displayed by indenting it from the left
margin.  Quotations are commonly displayed.  There
are short quotations
\begin{quote}
   This is a short quotation.  It consists of a
   single paragraph of text.  See how it is formatted.
\end{quote}
and longer ones.
\begin{quotation}
   This is a longer quotation.  It consists of two
   paragraphs of text, neither of which are
   particularly interesting.

   This is the second paragraph of the quotation.  It
   is just as dull as the first paragraph.
\end{quotation}
Another frequently-displayed structure is a list.
The following is an example of an \emph{itemized}
list.
\begin{itemize}
   \item This is the first item of an itemized list.
         Each item in the list is marked with a ``tick''.
         You don't have to worry about what kind of tick
         mark is used.

   \item This is the second item of the list.  It
         contains another list nested inside it.  The inner
         list is an \emph{enumerated} list.
         \begin{enumerate}
            \item This is the first item of an enumerated
                  list that is nested within the itemized list.

            \item This is the second item of the inner list.
                  \LaTeX\ allows you to nest lists deeper than
                  you really should.
         \end{enumerate}
         This is the rest of the second item of the outer
         list.  It is no more interesting than any other
         part of the item.
   \item This is the third item of the list.
\end{itemize}
You can even display poetry.
\begin{verse}
   There is an environment
    for verse \\             % The \\ command separates lines
   Whose features some poets % within a stanza.
   will curse.

                             % One or more blank lines separate stanzas.

   For instead of making\\
   Them do \emph{all} line breaking, \\
   It allows them to put too many words on a line when they'd rather be
   forced to be terse.
\end{verse}

Mathematical formulas may also be displayed.  A
displayed formula
is
one-line long; multiline
formulas require special formatting instructions.
   \[  \ip{\Gamma}{\psi'} = x'' + y^{2} + z_{i}^{n}\]
Don't start a paragraph with a displayed equation,
nor make one a paragraph by itself.

\end{document}               % End of document.

\section{Bench Huge Section 1355}
% synthetic bench section 1355
% This is a sample LaTeX input file.  (Version of 12 August 2004.)
%
% A '%' character causes TeX to ignore all remaining text on the line,
% and is used for comments like this one.

\documentclass{article}      % Specifies the document class

                             % The preamble begins here.
\title{An Example Document}  % Declares the document's title.
\author{Leslie Lamport}      % Declares the author's name.
\date{January 21, 1994}      % Deleting this command produces today's date.

\newcommand{\ip}[2]{(#1, #2)}
                             % Defines \ip{arg1}{arg2} to mean
                             % (arg1, arg2).

%\newcommand{\ip}[2]{\langle #1 | #2\rangle}
                             % This is an alternative definition of
                             % \ip that is commented out.

\begin{document}             % End of preamble and beginning of text.

\maketitle                   % Produces the title.

This is an example input file.  Comparing it with
the output it generates can show you how to
produce a simple document of your own.

\section{Ordinary Text}      % Produces section heading.  Lower-level
                             % sections are begun with similar
                             % \subsection and \subsubsection commands.

The ends  of words and sentences are marked
  by   spaces. It  doesn't matter how many
spaces    you type; one is as good as 100.  The
end of   a line counts as a space.

One   or more   blank lines denote the  end
of  a paragraph.

Since any number of consecutive spaces are treated
like a single one, the formatting of the input
file makes no difference to
      \LaTeX,                % The \LaTeX command generates the LaTeX logo.
but it makes a difference to you.  When you use
\LaTeX, making your input file as easy to read
as possible will be a great help as you write
your document and when you change it.  This sample
file shows how you can add comments to your own input
file.

Because printing is different from typewriting,
there are a number of things that you have to do
differently when preparing an input file than if
you were just typing the document directly.
Quotation marks like
       ``this''
have to be handled specially, as do quotes within
quotes:
       ``\,`this'            % \, separates the double and single quote.
        is what I just
        wrote, not  `that'\,''.

Dashes come in three sizes: an
       intra-word
dash, a medium dash for number ranges like
       1--2,
and a punctuation
       dash---like
this.

A sentence-ending space should be larger than the
space between words within a sentence.  You
sometimes have to type special commands in
conjunction with punctuation characters to get
this right, as in the following sentence.
       Gnats, gnus, etc.\ all  % `\ ' makes an inter-word space.
       begin with G\@.         % \@ marks end-of-sentence punctuation.
You should check the spaces after periods when
reading your output to make sure you haven't
forgotten any special cases.  Generating an
ellipsis
       \ldots\               % `\ ' is needed after `\ldots' because TeX
                             % ignores spaces after command names like \ldots
                             % made from \ + letters.
                             %
                             % Note how a `%' character causes TeX to ignore
                             % the end of the input line, so these blank lines
                             % do not start a new paragraph.
                             %
with the right spacing around the periods requires
a special command.

\LaTeX\ interprets some common characters as
commands, so you must type special commands to
generate them.  These characters include the
following:
       \$ \& \% \# \{ and \}.

In printing, text is usually emphasized with an
       \emph{italic}
type style.

\begin{em}
   A long segment of text can also be emphasized
   in this way.  Text within such a segment can be
   given \emph{additional} emphasis.
\end{em}

It is sometimes necessary to prevent \LaTeX\ from
breaking a line where it might otherwise do so.
This may be at a space, as between the ``Mr.''\ and
``Jones'' in
       ``Mr.~Jones'',        % ~ produces an unbreakable interword space.
or within a word---especially when the word is a
symbol like
       \mbox{\emph{itemnum}}
that makes little sense when hyphenated across
lines.

Footnotes\footnote{This is an example of a footnote.}
pose no problem.

\LaTeX\ is good at typesetting mathematical formulas
like
       \( x-3y + z = 7 \)
or
       \( a_{1} > x^{2n} + y^{2n} > x' \)
or
       \( \ip{A}{B} = \sum_{i} a_{i} b_{i} \).
The spaces you type in a formula are
ignored.  Remember that a letter like
       $x$                   % $ ... $  and  \( ... \)  are equivalent
is a formula when it denotes a mathematical
symbol, and it should be typed as one.

\section{Displayed Text}

Text is displayed by indenting it from the left
margin.  Quotations are commonly displayed.  There
are short quotations
\begin{quote}
   This is a short quotation.  It consists of a
   single paragraph of text.  See how it is formatted.
\end{quote}
and longer ones.
\begin{quotation}
   This is a longer quotation.  It consists of two
   paragraphs of text, neither of which are
   particularly interesting.

   This is the second paragraph of the quotation.  It
   is just as dull as the first paragraph.
\end{quotation}
Another frequently-displayed structure is a list.
The following is an example of an \emph{itemized}
list.
\begin{itemize}
   \item This is the first item of an itemized list.
         Each item in the list is marked with a ``tick''.
         You don't have to worry about what kind of tick
         mark is used.

   \item This is the second item of the list.  It
         contains another list nested inside it.  The inner
         list is an \emph{enumerated} list.
         \begin{enumerate}
            \item This is the first item of an enumerated
                  list that is nested within the itemized list.

            \item This is the second item of the inner list.
                  \LaTeX\ allows you to nest lists deeper than
                  you really should.
         \end{enumerate}
         This is the rest of the second item of the outer
         list.  It is no more interesting than any other
         part of the item.
   \item This is the third item of the list.
\end{itemize}
You can even display poetry.
\begin{verse}
   There is an environment
    for verse \\             % The \\ command separates lines
   Whose features some poets % within a stanza.
   will curse.

                             % One or more blank lines separate stanzas.

   For instead of making\\
   Them do \emph{all} line breaking, \\
   It allows them to put too many words on a line when they'd rather be
   forced to be terse.
\end{verse}

Mathematical formulas may also be displayed.  A
displayed formula
is
one-line long; multiline
formulas require special formatting instructions.
   \[  \ip{\Gamma}{\psi'} = x'' + y^{2} + z_{i}^{n}\]
Don't start a paragraph with a displayed equation,
nor make one a paragraph by itself.

\end{document}               % End of document.

\section{Bench Huge Section 1356}
% synthetic bench section 1356
% This is a sample LaTeX input file.  (Version of 12 August 2004.)
%
% A '%' character causes TeX to ignore all remaining text on the line,
% and is used for comments like this one.

\documentclass{article}      % Specifies the document class

                             % The preamble begins here.
\title{An Example Document}  % Declares the document's title.
\author{Leslie Lamport}      % Declares the author's name.
\date{January 21, 1994}      % Deleting this command produces today's date.

\newcommand{\ip}[2]{(#1, #2)}
                             % Defines \ip{arg1}{arg2} to mean
                             % (arg1, arg2).

%\newcommand{\ip}[2]{\langle #1 | #2\rangle}
                             % This is an alternative definition of
                             % \ip that is commented out.

\begin{document}             % End of preamble and beginning of text.

\maketitle                   % Produces the title.

This is an example input file.  Comparing it with
the output it generates can show you how to
produce a simple document of your own.

\section{Ordinary Text}      % Produces section heading.  Lower-level
                             % sections are begun with similar
                             % \subsection and \subsubsection commands.

The ends  of words and sentences are marked
  by   spaces. It  doesn't matter how many
spaces    you type; one is as good as 100.  The
end of   a line counts as a space.

One   or more   blank lines denote the  end
of  a paragraph.

Since any number of consecutive spaces are treated
like a single one, the formatting of the input
file makes no difference to
      \LaTeX,                % The \LaTeX command generates the LaTeX logo.
but it makes a difference to you.  When you use
\LaTeX, making your input file as easy to read
as possible will be a great help as you write
your document and when you change it.  This sample
file shows how you can add comments to your own input
file.

Because printing is different from typewriting,
there are a number of things that you have to do
differently when preparing an input file than if
you were just typing the document directly.
Quotation marks like
       ``this''
have to be handled specially, as do quotes within
quotes:
       ``\,`this'            % \, separates the double and single quote.
        is what I just
        wrote, not  `that'\,''.

Dashes come in three sizes: an
       intra-word
dash, a medium dash for number ranges like
       1--2,
and a punctuation
       dash---like
this.

A sentence-ending space should be larger than the
space between words within a sentence.  You
sometimes have to type special commands in
conjunction with punctuation characters to get
this right, as in the following sentence.
       Gnats, gnus, etc.\ all  % `\ ' makes an inter-word space.
       begin with G\@.         % \@ marks end-of-sentence punctuation.
You should check the spaces after periods when
reading your output to make sure you haven't
forgotten any special cases.  Generating an
ellipsis
       \ldots\               % `\ ' is needed after `\ldots' because TeX
                             % ignores spaces after command names like \ldots
                             % made from \ + letters.
                             %
                             % Note how a `%' character causes TeX to ignore
                             % the end of the input line, so these blank lines
                             % do not start a new paragraph.
                             %
with the right spacing around the periods requires
a special command.

\LaTeX\ interprets some common characters as
commands, so you must type special commands to
generate them.  These characters include the
following:
       \$ \& \% \# \{ and \}.

In printing, text is usually emphasized with an
       \emph{italic}
type style.

\begin{em}
   A long segment of text can also be emphasized
   in this way.  Text within such a segment can be
   given \emph{additional} emphasis.
\end{em}

It is sometimes necessary to prevent \LaTeX\ from
breaking a line where it might otherwise do so.
This may be at a space, as between the ``Mr.''\ and
``Jones'' in
       ``Mr.~Jones'',        % ~ produces an unbreakable interword space.
or within a word---especially when the word is a
symbol like
       \mbox{\emph{itemnum}}
that makes little sense when hyphenated across
lines.

Footnotes\footnote{This is an example of a footnote.}
pose no problem.

\LaTeX\ is good at typesetting mathematical formulas
like
       \( x-3y + z = 7 \)
or
       \( a_{1} > x^{2n} + y^{2n} > x' \)
or
       \( \ip{A}{B} = \sum_{i} a_{i} b_{i} \).
The spaces you type in a formula are
ignored.  Remember that a letter like
       $x$                   % $ ... $  and  \( ... \)  are equivalent
is a formula when it denotes a mathematical
symbol, and it should be typed as one.

\section{Displayed Text}

Text is displayed by indenting it from the left
margin.  Quotations are commonly displayed.  There
are short quotations
\begin{quote}
   This is a short quotation.  It consists of a
   single paragraph of text.  See how it is formatted.
\end{quote}
and longer ones.
\begin{quotation}
   This is a longer quotation.  It consists of two
   paragraphs of text, neither of which are
   particularly interesting.

   This is the second paragraph of the quotation.  It
   is just as dull as the first paragraph.
\end{quotation}
Another frequently-displayed structure is a list.
The following is an example of an \emph{itemized}
list.
\begin{itemize}
   \item This is the first item of an itemized list.
         Each item in the list is marked with a ``tick''.
         You don't have to worry about what kind of tick
         mark is used.

   \item This is the second item of the list.  It
         contains another list nested inside it.  The inner
         list is an \emph{enumerated} list.
         \begin{enumerate}
            \item This is the first item of an enumerated
                  list that is nested within the itemized list.

            \item This is the second item of the inner list.
                  \LaTeX\ allows you to nest lists deeper than
                  you really should.
         \end{enumerate}
         This is the rest of the second item of the outer
         list.  It is no more interesting than any other
         part of the item.
   \item This is the third item of the list.
\end{itemize}
You can even display poetry.
\begin{verse}
   There is an environment
    for verse \\             % The \\ command separates lines
   Whose features some poets % within a stanza.
   will curse.

                             % One or more blank lines separate stanzas.

   For instead of making\\
   Them do \emph{all} line breaking, \\
   It allows them to put too many words on a line when they'd rather be
   forced to be terse.
\end{verse}

Mathematical formulas may also be displayed.  A
displayed formula
is
one-line long; multiline
formulas require special formatting instructions.
   \[  \ip{\Gamma}{\psi'} = x'' + y^{2} + z_{i}^{n}\]
Don't start a paragraph with a displayed equation,
nor make one a paragraph by itself.

\end{document}               % End of document.

\section{Bench Huge Section 1357}
% synthetic bench section 1357
% This is a sample LaTeX input file.  (Version of 12 August 2004.)
%
% A '%' character causes TeX to ignore all remaining text on the line,
% and is used for comments like this one.

\documentclass{article}      % Specifies the document class

                             % The preamble begins here.
\title{An Example Document}  % Declares the document's title.
\author{Leslie Lamport}      % Declares the author's name.
\date{January 21, 1994}      % Deleting this command produces today's date.

\newcommand{\ip}[2]{(#1, #2)}
                             % Defines \ip{arg1}{arg2} to mean
                             % (arg1, arg2).

%\newcommand{\ip}[2]{\langle #1 | #2\rangle}
                             % This is an alternative definition of
                             % \ip that is commented out.

\begin{document}             % End of preamble and beginning of text.

\maketitle                   % Produces the title.

This is an example input file.  Comparing it with
the output it generates can show you how to
produce a simple document of your own.

\section{Ordinary Text}      % Produces section heading.  Lower-level
                             % sections are begun with similar
                             % \subsection and \subsubsection commands.

The ends  of words and sentences are marked
  by   spaces. It  doesn't matter how many
spaces    you type; one is as good as 100.  The
end of   a line counts as a space.

One   or more   blank lines denote the  end
of  a paragraph.

Since any number of consecutive spaces are treated
like a single one, the formatting of the input
file makes no difference to
      \LaTeX,                % The \LaTeX command generates the LaTeX logo.
but it makes a difference to you.  When you use
\LaTeX, making your input file as easy to read
as possible will be a great help as you write
your document and when you change it.  This sample
file shows how you can add comments to your own input
file.

Because printing is different from typewriting,
there are a number of things that you have to do
differently when preparing an input file than if
you were just typing the document directly.
Quotation marks like
       ``this''
have to be handled specially, as do quotes within
quotes:
       ``\,`this'            % \, separates the double and single quote.
        is what I just
        wrote, not  `that'\,''.

Dashes come in three sizes: an
       intra-word
dash, a medium dash for number ranges like
       1--2,
and a punctuation
       dash---like
this.

A sentence-ending space should be larger than the
space between words within a sentence.  You
sometimes have to type special commands in
conjunction with punctuation characters to get
this right, as in the following sentence.
       Gnats, gnus, etc.\ all  % `\ ' makes an inter-word space.
       begin with G\@.         % \@ marks end-of-sentence punctuation.
You should check the spaces after periods when
reading your output to make sure you haven't
forgotten any special cases.  Generating an
ellipsis
       \ldots\               % `\ ' is needed after `\ldots' because TeX
                             % ignores spaces after command names like \ldots
                             % made from \ + letters.
                             %
                             % Note how a `%' character causes TeX to ignore
                             % the end of the input line, so these blank lines
                             % do not start a new paragraph.
                             %
with the right spacing around the periods requires
a special command.

\LaTeX\ interprets some common characters as
commands, so you must type special commands to
generate them.  These characters include the
following:
       \$ \& \% \# \{ and \}.

In printing, text is usually emphasized with an
       \emph{italic}
type style.

\begin{em}
   A long segment of text can also be emphasized
   in this way.  Text within such a segment can be
   given \emph{additional} emphasis.
\end{em}

It is sometimes necessary to prevent \LaTeX\ from
breaking a line where it might otherwise do so.
This may be at a space, as between the ``Mr.''\ and
``Jones'' in
       ``Mr.~Jones'',        % ~ produces an unbreakable interword space.
or within a word---especially when the word is a
symbol like
       \mbox{\emph{itemnum}}
that makes little sense when hyphenated across
lines.

Footnotes\footnote{This is an example of a footnote.}
pose no problem.

\LaTeX\ is good at typesetting mathematical formulas
like
       \( x-3y + z = 7 \)
or
       \( a_{1} > x^{2n} + y^{2n} > x' \)
or
       \( \ip{A}{B} = \sum_{i} a_{i} b_{i} \).
The spaces you type in a formula are
ignored.  Remember that a letter like
       $x$                   % $ ... $  and  \( ... \)  are equivalent
is a formula when it denotes a mathematical
symbol, and it should be typed as one.

\section{Displayed Text}

Text is displayed by indenting it from the left
margin.  Quotations are commonly displayed.  There
are short quotations
\begin{quote}
   This is a short quotation.  It consists of a
   single paragraph of text.  See how it is formatted.
\end{quote}
and longer ones.
\begin{quotation}
   This is a longer quotation.  It consists of two
   paragraphs of text, neither of which are
   particularly interesting.

   This is the second paragraph of the quotation.  It
   is just as dull as the first paragraph.
\end{quotation}
Another frequently-displayed structure is a list.
The following is an example of an \emph{itemized}
list.
\begin{itemize}
   \item This is the first item of an itemized list.
         Each item in the list is marked with a ``tick''.
         You don't have to worry about what kind of tick
         mark is used.

   \item This is the second item of the list.  It
         contains another list nested inside it.  The inner
         list is an \emph{enumerated} list.
         \begin{enumerate}
            \item This is the first item of an enumerated
                  list that is nested within the itemized list.

            \item This is the second item of the inner list.
                  \LaTeX\ allows you to nest lists deeper than
                  you really should.
         \end{enumerate}
         This is the rest of the second item of the outer
         list.  It is no more interesting than any other
         part of the item.
   \item This is the third item of the list.
\end{itemize}
You can even display poetry.
\begin{verse}
   There is an environment
    for verse \\             % The \\ command separates lines
   Whose features some poets % within a stanza.
   will curse.

                             % One or more blank lines separate stanzas.

   For instead of making\\
   Them do \emph{all} line breaking, \\
   It allows them to put too many words on a line when they'd rather be
   forced to be terse.
\end{verse}

Mathematical formulas may also be displayed.  A
displayed formula
is
one-line long; multiline
formulas require special formatting instructions.
   \[  \ip{\Gamma}{\psi'} = x'' + y^{2} + z_{i}^{n}\]
Don't start a paragraph with a displayed equation,
nor make one a paragraph by itself.

\end{document}               % End of document.

\section{Bench Huge Section 1358}
% synthetic bench section 1358
% This is a sample LaTeX input file.  (Version of 12 August 2004.)
%
% A '%' character causes TeX to ignore all remaining text on the line,
% and is used for comments like this one.

\documentclass{article}      % Specifies the document class

                             % The preamble begins here.
\title{An Example Document}  % Declares the document's title.
\author{Leslie Lamport}      % Declares the author's name.
\date{January 21, 1994}      % Deleting this command produces today's date.

\newcommand{\ip}[2]{(#1, #2)}
                             % Defines \ip{arg1}{arg2} to mean
                             % (arg1, arg2).

%\newcommand{\ip}[2]{\langle #1 | #2\rangle}
                             % This is an alternative definition of
                             % \ip that is commented out.

\begin{document}             % End of preamble and beginning of text.

\maketitle                   % Produces the title.

This is an example input file.  Comparing it with
the output it generates can show you how to
produce a simple document of your own.

\section{Ordinary Text}      % Produces section heading.  Lower-level
                             % sections are begun with similar
                             % \subsection and \subsubsection commands.

The ends  of words and sentences are marked
  by   spaces. It  doesn't matter how many
spaces    you type; one is as good as 100.  The
end of   a line counts as a space.

One   or more   blank lines denote the  end
of  a paragraph.

Since any number of consecutive spaces are treated
like a single one, the formatting of the input
file makes no difference to
      \LaTeX,                % The \LaTeX command generates the LaTeX logo.
but it makes a difference to you.  When you use
\LaTeX, making your input file as easy to read
as possible will be a great help as you write
your document and when you change it.  This sample
file shows how you can add comments to your own input
file.

Because printing is different from typewriting,
there are a number of things that you have to do
differently when preparing an input file than if
you were just typing the document directly.
Quotation marks like
       ``this''
have to be handled specially, as do quotes within
quotes:
       ``\,`this'            % \, separates the double and single quote.
        is what I just
        wrote, not  `that'\,''.

Dashes come in three sizes: an
       intra-word
dash, a medium dash for number ranges like
       1--2,
and a punctuation
       dash---like
this.

A sentence-ending space should be larger than the
space between words within a sentence.  You
sometimes have to type special commands in
conjunction with punctuation characters to get
this right, as in the following sentence.
       Gnats, gnus, etc.\ all  % `\ ' makes an inter-word space.
       begin with G\@.         % \@ marks end-of-sentence punctuation.
You should check the spaces after periods when
reading your output to make sure you haven't
forgotten any special cases.  Generating an
ellipsis
       \ldots\               % `\ ' is needed after `\ldots' because TeX
                             % ignores spaces after command names like \ldots
                             % made from \ + letters.
                             %
                             % Note how a `%' character causes TeX to ignore
                             % the end of the input line, so these blank lines
                             % do not start a new paragraph.
                             %
with the right spacing around the periods requires
a special command.

\LaTeX\ interprets some common characters as
commands, so you must type special commands to
generate them.  These characters include the
following:
       \$ \& \% \# \{ and \}.

In printing, text is usually emphasized with an
       \emph{italic}
type style.

\begin{em}
   A long segment of text can also be emphasized
   in this way.  Text within such a segment can be
   given \emph{additional} emphasis.
\end{em}

It is sometimes necessary to prevent \LaTeX\ from
breaking a line where it might otherwise do so.
This may be at a space, as between the ``Mr.''\ and
``Jones'' in
       ``Mr.~Jones'',        % ~ produces an unbreakable interword space.
or within a word---especially when the word is a
symbol like
       \mbox{\emph{itemnum}}
that makes little sense when hyphenated across
lines.

Footnotes\footnote{This is an example of a footnote.}
pose no problem.

\LaTeX\ is good at typesetting mathematical formulas
like
       \( x-3y + z = 7 \)
or
       \( a_{1} > x^{2n} + y^{2n} > x' \)
or
       \( \ip{A}{B} = \sum_{i} a_{i} b_{i} \).
The spaces you type in a formula are
ignored.  Remember that a letter like
       $x$                   % $ ... $  and  \( ... \)  are equivalent
is a formula when it denotes a mathematical
symbol, and it should be typed as one.

\section{Displayed Text}

Text is displayed by indenting it from the left
margin.  Quotations are commonly displayed.  There
are short quotations
\begin{quote}
   This is a short quotation.  It consists of a
   single paragraph of text.  See how it is formatted.
\end{quote}
and longer ones.
\begin{quotation}
   This is a longer quotation.  It consists of two
   paragraphs of text, neither of which are
   particularly interesting.

   This is the second paragraph of the quotation.  It
   is just as dull as the first paragraph.
\end{quotation}
Another frequently-displayed structure is a list.
The following is an example of an \emph{itemized}
list.
\begin{itemize}
   \item This is the first item of an itemized list.
         Each item in the list is marked with a ``tick''.
         You don't have to worry about what kind of tick
         mark is used.

   \item This is the second item of the list.  It
         contains another list nested inside it.  The inner
         list is an \emph{enumerated} list.
         \begin{enumerate}
            \item This is the first item of an enumerated
                  list that is nested within the itemized list.

            \item This is the second item of the inner list.
                  \LaTeX\ allows you to nest lists deeper than
                  you really should.
         \end{enumerate}
         This is the rest of the second item of the outer
         list.  It is no more interesting than any other
         part of the item.
   \item This is the third item of the list.
\end{itemize}
You can even display poetry.
\begin{verse}
   There is an environment
    for verse \\             % The \\ command separates lines
   Whose features some poets % within a stanza.
   will curse.

                             % One or more blank lines separate stanzas.

   For instead of making\\
   Them do \emph{all} line breaking, \\
   It allows them to put too many words on a line when they'd rather be
   forced to be terse.
\end{verse}

Mathematical formulas may also be displayed.  A
displayed formula
is
one-line long; multiline
formulas require special formatting instructions.
   \[  \ip{\Gamma}{\psi'} = x'' + y^{2} + z_{i}^{n}\]
Don't start a paragraph with a displayed equation,
nor make one a paragraph by itself.

\end{document}               % End of document.

\section{Bench Huge Section 1359}
% synthetic bench section 1359
% This is a sample LaTeX input file.  (Version of 12 August 2004.)
%
% A '%' character causes TeX to ignore all remaining text on the line,
% and is used for comments like this one.

\documentclass{article}      % Specifies the document class

                             % The preamble begins here.
\title{An Example Document}  % Declares the document's title.
\author{Leslie Lamport}      % Declares the author's name.
\date{January 21, 1994}      % Deleting this command produces today's date.

\newcommand{\ip}[2]{(#1, #2)}
                             % Defines \ip{arg1}{arg2} to mean
                             % (arg1, arg2).

%\newcommand{\ip}[2]{\langle #1 | #2\rangle}
                             % This is an alternative definition of
                             % \ip that is commented out.

\begin{document}             % End of preamble and beginning of text.

\maketitle                   % Produces the title.

This is an example input file.  Comparing it with
the output it generates can show you how to
produce a simple document of your own.

\section{Ordinary Text}      % Produces section heading.  Lower-level
                             % sections are begun with similar
                             % \subsection and \subsubsection commands.

The ends  of words and sentences are marked
  by   spaces. It  doesn't matter how many
spaces    you type; one is as good as 100.  The
end of   a line counts as a space.

One   or more   blank lines denote the  end
of  a paragraph.

Since any number of consecutive spaces are treated
like a single one, the formatting of the input
file makes no difference to
      \LaTeX,                % The \LaTeX command generates the LaTeX logo.
but it makes a difference to you.  When you use
\LaTeX, making your input file as easy to read
as possible will be a great help as you write
your document and when you change it.  This sample
file shows how you can add comments to your own input
file.

Because printing is different from typewriting,
there are a number of things that you have to do
differently when preparing an input file than if
you were just typing the document directly.
Quotation marks like
       ``this''
have to be handled specially, as do quotes within
quotes:
       ``\,`this'            % \, separates the double and single quote.
        is what I just
        wrote, not  `that'\,''.

Dashes come in three sizes: an
       intra-word
dash, a medium dash for number ranges like
       1--2,
and a punctuation
       dash---like
this.

A sentence-ending space should be larger than the
space between words within a sentence.  You
sometimes have to type special commands in
conjunction with punctuation characters to get
this right, as in the following sentence.
       Gnats, gnus, etc.\ all  % `\ ' makes an inter-word space.
       begin with G\@.         % \@ marks end-of-sentence punctuation.
You should check the spaces after periods when
reading your output to make sure you haven't
forgotten any special cases.  Generating an
ellipsis
       \ldots\               % `\ ' is needed after `\ldots' because TeX
                             % ignores spaces after command names like \ldots
                             % made from \ + letters.
                             %
                             % Note how a `%' character causes TeX to ignore
                             % the end of the input line, so these blank lines
                             % do not start a new paragraph.
                             %
with the right spacing around the periods requires
a special command.

\LaTeX\ interprets some common characters as
commands, so you must type special commands to
generate them.  These characters include the
following:
       \$ \& \% \# \{ and \}.

In printing, text is usually emphasized with an
       \emph{italic}
type style.

\begin{em}
   A long segment of text can also be emphasized
   in this way.  Text within such a segment can be
   given \emph{additional} emphasis.
\end{em}

It is sometimes necessary to prevent \LaTeX\ from
breaking a line where it might otherwise do so.
This may be at a space, as between the ``Mr.''\ and
``Jones'' in
       ``Mr.~Jones'',        % ~ produces an unbreakable interword space.
or within a word---especially when the word is a
symbol like
       \mbox{\emph{itemnum}}
that makes little sense when hyphenated across
lines.

Footnotes\footnote{This is an example of a footnote.}
pose no problem.

\LaTeX\ is good at typesetting mathematical formulas
like
       \( x-3y + z = 7 \)
or
       \( a_{1} > x^{2n} + y^{2n} > x' \)
or
       \( \ip{A}{B} = \sum_{i} a_{i} b_{i} \).
The spaces you type in a formula are
ignored.  Remember that a letter like
       $x$                   % $ ... $  and  \( ... \)  are equivalent
is a formula when it denotes a mathematical
symbol, and it should be typed as one.

\section{Displayed Text}

Text is displayed by indenting it from the left
margin.  Quotations are commonly displayed.  There
are short quotations
\begin{quote}
   This is a short quotation.  It consists of a
   single paragraph of text.  See how it is formatted.
\end{quote}
and longer ones.
\begin{quotation}
   This is a longer quotation.  It consists of two
   paragraphs of text, neither of which are
   particularly interesting.

   This is the second paragraph of the quotation.  It
   is just as dull as the first paragraph.
\end{quotation}
Another frequently-displayed structure is a list.
The following is an example of an \emph{itemized}
list.
\begin{itemize}
   \item This is the first item of an itemized list.
         Each item in the list is marked with a ``tick''.
         You don't have to worry about what kind of tick
         mark is used.

   \item This is the second item of the list.  It
         contains another list nested inside it.  The inner
         list is an \emph{enumerated} list.
         \begin{enumerate}
            \item This is the first item of an enumerated
                  list that is nested within the itemized list.

            \item This is the second item of the inner list.
                  \LaTeX\ allows you to nest lists deeper than
                  you really should.
         \end{enumerate}
         This is the rest of the second item of the outer
         list.  It is no more interesting than any other
         part of the item.
   \item This is the third item of the list.
\end{itemize}
You can even display poetry.
\begin{verse}
   There is an environment
    for verse \\             % The \\ command separates lines
   Whose features some poets % within a stanza.
   will curse.

                             % One or more blank lines separate stanzas.

   For instead of making\\
   Them do \emph{all} line breaking, \\
   It allows them to put too many words on a line when they'd rather be
   forced to be terse.
\end{verse}

Mathematical formulas may also be displayed.  A
displayed formula
is
one-line long; multiline
formulas require special formatting instructions.
   \[  \ip{\Gamma}{\psi'} = x'' + y^{2} + z_{i}^{n}\]
Don't start a paragraph with a displayed equation,
nor make one a paragraph by itself.

\end{document}               % End of document.

\section{Bench Huge Section 1360}
% synthetic bench section 1360
% This is a sample LaTeX input file.  (Version of 12 August 2004.)
%
% A '%' character causes TeX to ignore all remaining text on the line,
% and is used for comments like this one.

\documentclass{article}      % Specifies the document class

                             % The preamble begins here.
\title{An Example Document}  % Declares the document's title.
\author{Leslie Lamport}      % Declares the author's name.
\date{January 21, 1994}      % Deleting this command produces today's date.

\newcommand{\ip}[2]{(#1, #2)}
                             % Defines \ip{arg1}{arg2} to mean
                             % (arg1, arg2).

%\newcommand{\ip}[2]{\langle #1 | #2\rangle}
                             % This is an alternative definition of
                             % \ip that is commented out.

\begin{document}             % End of preamble and beginning of text.

\maketitle                   % Produces the title.

This is an example input file.  Comparing it with
the output it generates can show you how to
produce a simple document of your own.

\section{Ordinary Text}      % Produces section heading.  Lower-level
                             % sections are begun with similar
                             % \subsection and \subsubsection commands.

The ends  of words and sentences are marked
  by   spaces. It  doesn't matter how many
spaces    you type; one is as good as 100.  The
end of   a line counts as a space.

One   or more   blank lines denote the  end
of  a paragraph.

Since any number of consecutive spaces are treated
like a single one, the formatting of the input
file makes no difference to
      \LaTeX,                % The \LaTeX command generates the LaTeX logo.
but it makes a difference to you.  When you use
\LaTeX, making your input file as easy to read
as possible will be a great help as you write
your document and when you change it.  This sample
file shows how you can add comments to your own input
file.

Because printing is different from typewriting,
there are a number of things that you have to do
differently when preparing an input file than if
you were just typing the document directly.
Quotation marks like
       ``this''
have to be handled specially, as do quotes within
quotes:
       ``\,`this'            % \, separates the double and single quote.
        is what I just
        wrote, not  `that'\,''.

Dashes come in three sizes: an
       intra-word
dash, a medium dash for number ranges like
       1--2,
and a punctuation
       dash---like
this.

A sentence-ending space should be larger than the
space between words within a sentence.  You
sometimes have to type special commands in
conjunction with punctuation characters to get
this right, as in the following sentence.
       Gnats, gnus, etc.\ all  % `\ ' makes an inter-word space.
       begin with G\@.         % \@ marks end-of-sentence punctuation.
You should check the spaces after periods when
reading your output to make sure you haven't
forgotten any special cases.  Generating an
ellipsis
       \ldots\               % `\ ' is needed after `\ldots' because TeX
                             % ignores spaces after command names like \ldots
                             % made from \ + letters.
                             %
                             % Note how a `%' character causes TeX to ignore
                             % the end of the input line, so these blank lines
                             % do not start a new paragraph.
                             %
with the right spacing around the periods requires
a special command.

\LaTeX\ interprets some common characters as
commands, so you must type special commands to
generate them.  These characters include the
following:
       \$ \& \% \# \{ and \}.

In printing, text is usually emphasized with an
       \emph{italic}
type style.

\begin{em}
   A long segment of text can also be emphasized
   in this way.  Text within such a segment can be
   given \emph{additional} emphasis.
\end{em}

It is sometimes necessary to prevent \LaTeX\ from
breaking a line where it might otherwise do so.
This may be at a space, as between the ``Mr.''\ and
``Jones'' in
       ``Mr.~Jones'',        % ~ produces an unbreakable interword space.
or within a word---especially when the word is a
symbol like
       \mbox{\emph{itemnum}}
that makes little sense when hyphenated across
lines.

Footnotes\footnote{This is an example of a footnote.}
pose no problem.

\LaTeX\ is good at typesetting mathematical formulas
like
       \( x-3y + z = 7 \)
or
       \( a_{1} > x^{2n} + y^{2n} > x' \)
or
       \( \ip{A}{B} = \sum_{i} a_{i} b_{i} \).
The spaces you type in a formula are
ignored.  Remember that a letter like
       $x$                   % $ ... $  and  \( ... \)  are equivalent
is a formula when it denotes a mathematical
symbol, and it should be typed as one.

\section{Displayed Text}

Text is displayed by indenting it from the left
margin.  Quotations are commonly displayed.  There
are short quotations
\begin{quote}
   This is a short quotation.  It consists of a
   single paragraph of text.  See how it is formatted.
\end{quote}
and longer ones.
\begin{quotation}
   This is a longer quotation.  It consists of two
   paragraphs of text, neither of which are
   particularly interesting.

   This is the second paragraph of the quotation.  It
   is just as dull as the first paragraph.
\end{quotation}
Another frequently-displayed structure is a list.
The following is an example of an \emph{itemized}
list.
\begin{itemize}
   \item This is the first item of an itemized list.
         Each item in the list is marked with a ``tick''.
         You don't have to worry about what kind of tick
         mark is used.

   \item This is the second item of the list.  It
         contains another list nested inside it.  The inner
         list is an \emph{enumerated} list.
         \begin{enumerate}
            \item This is the first item of an enumerated
                  list that is nested within the itemized list.

            \item This is the second item of the inner list.
                  \LaTeX\ allows you to nest lists deeper than
                  you really should.
         \end{enumerate}
         This is the rest of the second item of the outer
         list.  It is no more interesting than any other
         part of the item.
   \item This is the third item of the list.
\end{itemize}
You can even display poetry.
\begin{verse}
   There is an environment
    for verse \\             % The \\ command separates lines
   Whose features some poets % within a stanza.
   will curse.

                             % One or more blank lines separate stanzas.

   For instead of making\\
   Them do \emph{all} line breaking, \\
   It allows them to put too many words on a line when they'd rather be
   forced to be terse.
\end{verse}

Mathematical formulas may also be displayed.  A
displayed formula
is
one-line long; multiline
formulas require special formatting instructions.
   \[  \ip{\Gamma}{\psi'} = x'' + y^{2} + z_{i}^{n}\]
Don't start a paragraph with a displayed equation,
nor make one a paragraph by itself.

\end{document}               % End of document.

\section{Bench Huge Section 1361}
% synthetic bench section 1361
% This is a sample LaTeX input file.  (Version of 12 August 2004.)
%
% A '%' character causes TeX to ignore all remaining text on the line,
% and is used for comments like this one.

\documentclass{article}      % Specifies the document class

                             % The preamble begins here.
\title{An Example Document}  % Declares the document's title.
\author{Leslie Lamport}      % Declares the author's name.
\date{January 21, 1994}      % Deleting this command produces today's date.

\newcommand{\ip}[2]{(#1, #2)}
                             % Defines \ip{arg1}{arg2} to mean
                             % (arg1, arg2).

%\newcommand{\ip}[2]{\langle #1 | #2\rangle}
                             % This is an alternative definition of
                             % \ip that is commented out.

\begin{document}             % End of preamble and beginning of text.

\maketitle                   % Produces the title.

This is an example input file.  Comparing it with
the output it generates can show you how to
produce a simple document of your own.

\section{Ordinary Text}      % Produces section heading.  Lower-level
                             % sections are begun with similar
                             % \subsection and \subsubsection commands.

The ends  of words and sentences are marked
  by   spaces. It  doesn't matter how many
spaces    you type; one is as good as 100.  The
end of   a line counts as a space.

One   or more   blank lines denote the  end
of  a paragraph.

Since any number of consecutive spaces are treated
like a single one, the formatting of the input
file makes no difference to
      \LaTeX,                % The \LaTeX command generates the LaTeX logo.
but it makes a difference to you.  When you use
\LaTeX, making your input file as easy to read
as possible will be a great help as you write
your document and when you change it.  This sample
file shows how you can add comments to your own input
file.

Because printing is different from typewriting,
there are a number of things that you have to do
differently when preparing an input file than if
you were just typing the document directly.
Quotation marks like
       ``this''
have to be handled specially, as do quotes within
quotes:
       ``\,`this'            % \, separates the double and single quote.
        is what I just
        wrote, not  `that'\,''.

Dashes come in three sizes: an
       intra-word
dash, a medium dash for number ranges like
       1--2,
and a punctuation
       dash---like
this.

A sentence-ending space should be larger than the
space between words within a sentence.  You
sometimes have to type special commands in
conjunction with punctuation characters to get
this right, as in the following sentence.
       Gnats, gnus, etc.\ all  % `\ ' makes an inter-word space.
       begin with G\@.         % \@ marks end-of-sentence punctuation.
You should check the spaces after periods when
reading your output to make sure you haven't
forgotten any special cases.  Generating an
ellipsis
       \ldots\               % `\ ' is needed after `\ldots' because TeX
                             % ignores spaces after command names like \ldots
                             % made from \ + letters.
                             %
                             % Note how a `%' character causes TeX to ignore
                             % the end of the input line, so these blank lines
                             % do not start a new paragraph.
                             %
with the right spacing around the periods requires
a special command.

\LaTeX\ interprets some common characters as
commands, so you must type special commands to
generate them.  These characters include the
following:
       \$ \& \% \# \{ and \}.

In printing, text is usually emphasized with an
       \emph{italic}
type style.

\begin{em}
   A long segment of text can also be emphasized
   in this way.  Text within such a segment can be
   given \emph{additional} emphasis.
\end{em}

It is sometimes necessary to prevent \LaTeX\ from
breaking a line where it might otherwise do so.
This may be at a space, as between the ``Mr.''\ and
``Jones'' in
       ``Mr.~Jones'',        % ~ produces an unbreakable interword space.
or within a word---especially when the word is a
symbol like
       \mbox{\emph{itemnum}}
that makes little sense when hyphenated across
lines.

Footnotes\footnote{This is an example of a footnote.}
pose no problem.

\LaTeX\ is good at typesetting mathematical formulas
like
       \( x-3y + z = 7 \)
or
       \( a_{1} > x^{2n} + y^{2n} > x' \)
or
       \( \ip{A}{B} = \sum_{i} a_{i} b_{i} \).
The spaces you type in a formula are
ignored.  Remember that a letter like
       $x$                   % $ ... $  and  \( ... \)  are equivalent
is a formula when it denotes a mathematical
symbol, and it should be typed as one.

\section{Displayed Text}

Text is displayed by indenting it from the left
margin.  Quotations are commonly displayed.  There
are short quotations
\begin{quote}
   This is a short quotation.  It consists of a
   single paragraph of text.  See how it is formatted.
\end{quote}
and longer ones.
\begin{quotation}
   This is a longer quotation.  It consists of two
   paragraphs of text, neither of which are
   particularly interesting.

   This is the second paragraph of the quotation.  It
   is just as dull as the first paragraph.
\end{quotation}
Another frequently-displayed structure is a list.
The following is an example of an \emph{itemized}
list.
\begin{itemize}
   \item This is the first item of an itemized list.
         Each item in the list is marked with a ``tick''.
         You don't have to worry about what kind of tick
         mark is used.

   \item This is the second item of the list.  It
         contains another list nested inside it.  The inner
         list is an \emph{enumerated} list.
         \begin{enumerate}
            \item This is the first item of an enumerated
                  list that is nested within the itemized list.

            \item This is the second item of the inner list.
                  \LaTeX\ allows you to nest lists deeper than
                  you really should.
         \end{enumerate}
         This is the rest of the second item of the outer
         list.  It is no more interesting than any other
         part of the item.
   \item This is the third item of the list.
\end{itemize}
You can even display poetry.
\begin{verse}
   There is an environment
    for verse \\             % The \\ command separates lines
   Whose features some poets % within a stanza.
   will curse.

                             % One or more blank lines separate stanzas.

   For instead of making\\
   Them do \emph{all} line breaking, \\
   It allows them to put too many words on a line when they'd rather be
   forced to be terse.
\end{verse}

Mathematical formulas may also be displayed.  A
displayed formula
is
one-line long; multiline
formulas require special formatting instructions.
   \[  \ip{\Gamma}{\psi'} = x'' + y^{2} + z_{i}^{n}\]
Don't start a paragraph with a displayed equation,
nor make one a paragraph by itself.

\end{document}               % End of document.

\section{Bench Huge Section 1362}
% synthetic bench section 1362
% This is a sample LaTeX input file.  (Version of 12 August 2004.)
%
% A '%' character causes TeX to ignore all remaining text on the line,
% and is used for comments like this one.

\documentclass{article}      % Specifies the document class

                             % The preamble begins here.
\title{An Example Document}  % Declares the document's title.
\author{Leslie Lamport}      % Declares the author's name.
\date{January 21, 1994}      % Deleting this command produces today's date.

\newcommand{\ip}[2]{(#1, #2)}
                             % Defines \ip{arg1}{arg2} to mean
                             % (arg1, arg2).

%\newcommand{\ip}[2]{\langle #1 | #2\rangle}
                             % This is an alternative definition of
                             % \ip that is commented out.

\begin{document}             % End of preamble and beginning of text.

\maketitle                   % Produces the title.

This is an example input file.  Comparing it with
the output it generates can show you how to
produce a simple document of your own.

\section{Ordinary Text}      % Produces section heading.  Lower-level
                             % sections are begun with similar
                             % \subsection and \subsubsection commands.

The ends  of words and sentences are marked
  by   spaces. It  doesn't matter how many
spaces    you type; one is as good as 100.  The
end of   a line counts as a space.

One   or more   blank lines denote the  end
of  a paragraph.

Since any number of consecutive spaces are treated
like a single one, the formatting of the input
file makes no difference to
      \LaTeX,                % The \LaTeX command generates the LaTeX logo.
but it makes a difference to you.  When you use
\LaTeX, making your input file as easy to read
as possible will be a great help as you write
your document and when you change it.  This sample
file shows how you can add comments to your own input
file.

Because printing is different from typewriting,
there are a number of things that you have to do
differently when preparing an input file than if
you were just typing the document directly.
Quotation marks like
       ``this''
have to be handled specially, as do quotes within
quotes:
       ``\,`this'            % \, separates the double and single quote.
        is what I just
        wrote, not  `that'\,''.

Dashes come in three sizes: an
       intra-word
dash, a medium dash for number ranges like
       1--2,
and a punctuation
       dash---like
this.

A sentence-ending space should be larger than the
space between words within a sentence.  You
sometimes have to type special commands in
conjunction with punctuation characters to get
this right, as in the following sentence.
       Gnats, gnus, etc.\ all  % `\ ' makes an inter-word space.
       begin with G\@.         % \@ marks end-of-sentence punctuation.
You should check the spaces after periods when
reading your output to make sure you haven't
forgotten any special cases.  Generating an
ellipsis
       \ldots\               % `\ ' is needed after `\ldots' because TeX
                             % ignores spaces after command names like \ldots
                             % made from \ + letters.
                             %
                             % Note how a `%' character causes TeX to ignore
                             % the end of the input line, so these blank lines
                             % do not start a new paragraph.
                             %
with the right spacing around the periods requires
a special command.

\LaTeX\ interprets some common characters as
commands, so you must type special commands to
generate them.  These characters include the
following:
       \$ \& \% \# \{ and \}.

In printing, text is usually emphasized with an
       \emph{italic}
type style.

\begin{em}
   A long segment of text can also be emphasized
   in this way.  Text within such a segment can be
   given \emph{additional} emphasis.
\end{em}

It is sometimes necessary to prevent \LaTeX\ from
breaking a line where it might otherwise do so.
This may be at a space, as between the ``Mr.''\ and
``Jones'' in
       ``Mr.~Jones'',        % ~ produces an unbreakable interword space.
or within a word---especially when the word is a
symbol like
       \mbox{\emph{itemnum}}
that makes little sense when hyphenated across
lines.

Footnotes\footnote{This is an example of a footnote.}
pose no problem.

\LaTeX\ is good at typesetting mathematical formulas
like
       \( x-3y + z = 7 \)
or
       \( a_{1} > x^{2n} + y^{2n} > x' \)
or
       \( \ip{A}{B} = \sum_{i} a_{i} b_{i} \).
The spaces you type in a formula are
ignored.  Remember that a letter like
       $x$                   % $ ... $  and  \( ... \)  are equivalent
is a formula when it denotes a mathematical
symbol, and it should be typed as one.

\section{Displayed Text}

Text is displayed by indenting it from the left
margin.  Quotations are commonly displayed.  There
are short quotations
\begin{quote}
   This is a short quotation.  It consists of a
   single paragraph of text.  See how it is formatted.
\end{quote}
and longer ones.
\begin{quotation}
   This is a longer quotation.  It consists of two
   paragraphs of text, neither of which are
   particularly interesting.

   This is the second paragraph of the quotation.  It
   is just as dull as the first paragraph.
\end{quotation}
Another frequently-displayed structure is a list.
The following is an example of an \emph{itemized}
list.
\begin{itemize}
   \item This is the first item of an itemized list.
         Each item in the list is marked with a ``tick''.
         You don't have to worry about what kind of tick
         mark is used.

   \item This is the second item of the list.  It
         contains another list nested inside it.  The inner
         list is an \emph{enumerated} list.
         \begin{enumerate}
            \item This is the first item of an enumerated
                  list that is nested within the itemized list.

            \item This is the second item of the inner list.
                  \LaTeX\ allows you to nest lists deeper than
                  you really should.
         \end{enumerate}
         This is the rest of the second item of the outer
         list.  It is no more interesting than any other
         part of the item.
   \item This is the third item of the list.
\end{itemize}
You can even display poetry.
\begin{verse}
   There is an environment
    for verse \\             % The \\ command separates lines
   Whose features some poets % within a stanza.
   will curse.

                             % One or more blank lines separate stanzas.

   For instead of making\\
   Them do \emph{all} line breaking, \\
   It allows them to put too many words on a line when they'd rather be
   forced to be terse.
\end{verse}

Mathematical formulas may also be displayed.  A
displayed formula
is
one-line long; multiline
formulas require special formatting instructions.
   \[  \ip{\Gamma}{\psi'} = x'' + y^{2} + z_{i}^{n}\]
Don't start a paragraph with a displayed equation,
nor make one a paragraph by itself.

\end{document}               % End of document.

\section{Bench Huge Section 1363}
% synthetic bench section 1363
% This is a sample LaTeX input file.  (Version of 12 August 2004.)
%
% A '%' character causes TeX to ignore all remaining text on the line,
% and is used for comments like this one.

\documentclass{article}      % Specifies the document class

                             % The preamble begins here.
\title{An Example Document}  % Declares the document's title.
\author{Leslie Lamport}      % Declares the author's name.
\date{January 21, 1994}      % Deleting this command produces today's date.

\newcommand{\ip}[2]{(#1, #2)}
                             % Defines \ip{arg1}{arg2} to mean
                             % (arg1, arg2).

%\newcommand{\ip}[2]{\langle #1 | #2\rangle}
                             % This is an alternative definition of
                             % \ip that is commented out.

\begin{document}             % End of preamble and beginning of text.

\maketitle                   % Produces the title.

This is an example input file.  Comparing it with
the output it generates can show you how to
produce a simple document of your own.

\section{Ordinary Text}      % Produces section heading.  Lower-level
                             % sections are begun with similar
                             % \subsection and \subsubsection commands.

The ends  of words and sentences are marked
  by   spaces. It  doesn't matter how many
spaces    you type; one is as good as 100.  The
end of   a line counts as a space.

One   or more   blank lines denote the  end
of  a paragraph.

Since any number of consecutive spaces are treated
like a single one, the formatting of the input
file makes no difference to
      \LaTeX,                % The \LaTeX command generates the LaTeX logo.
but it makes a difference to you.  When you use
\LaTeX, making your input file as easy to read
as possible will be a great help as you write
your document and when you change it.  This sample
file shows how you can add comments to your own input
file.

Because printing is different from typewriting,
there are a number of things that you have to do
differently when preparing an input file than if
you were just typing the document directly.
Quotation marks like
       ``this''
have to be handled specially, as do quotes within
quotes:
       ``\,`this'            % \, separates the double and single quote.
        is what I just
        wrote, not  `that'\,''.

Dashes come in three sizes: an
       intra-word
dash, a medium dash for number ranges like
       1--2,
and a punctuation
       dash---like
this.

A sentence-ending space should be larger than the
space between words within a sentence.  You
sometimes have to type special commands in
conjunction with punctuation characters to get
this right, as in the following sentence.
       Gnats, gnus, etc.\ all  % `\ ' makes an inter-word space.
       begin with G\@.         % \@ marks end-of-sentence punctuation.
You should check the spaces after periods when
reading your output to make sure you haven't
forgotten any special cases.  Generating an
ellipsis
       \ldots\               % `\ ' is needed after `\ldots' because TeX
                             % ignores spaces after command names like \ldots
                             % made from \ + letters.
                             %
                             % Note how a `%' character causes TeX to ignore
                             % the end of the input line, so these blank lines
                             % do not start a new paragraph.
                             %
with the right spacing around the periods requires
a special command.

\LaTeX\ interprets some common characters as
commands, so you must type special commands to
generate them.  These characters include the
following:
       \$ \& \% \# \{ and \}.

In printing, text is usually emphasized with an
       \emph{italic}
type style.

\begin{em}
   A long segment of text can also be emphasized
   in this way.  Text within such a segment can be
   given \emph{additional} emphasis.
\end{em}

It is sometimes necessary to prevent \LaTeX\ from
breaking a line where it might otherwise do so.
This may be at a space, as between the ``Mr.''\ and
``Jones'' in
       ``Mr.~Jones'',        % ~ produces an unbreakable interword space.
or within a word---especially when the word is a
symbol like
       \mbox{\emph{itemnum}}
that makes little sense when hyphenated across
lines.

Footnotes\footnote{This is an example of a footnote.}
pose no problem.

\LaTeX\ is good at typesetting mathematical formulas
like
       \( x-3y + z = 7 \)
or
       \( a_{1} > x^{2n} + y^{2n} > x' \)
or
       \( \ip{A}{B} = \sum_{i} a_{i} b_{i} \).
The spaces you type in a formula are
ignored.  Remember that a letter like
       $x$                   % $ ... $  and  \( ... \)  are equivalent
is a formula when it denotes a mathematical
symbol, and it should be typed as one.

\section{Displayed Text}

Text is displayed by indenting it from the left
margin.  Quotations are commonly displayed.  There
are short quotations
\begin{quote}
   This is a short quotation.  It consists of a
   single paragraph of text.  See how it is formatted.
\end{quote}
and longer ones.
\begin{quotation}
   This is a longer quotation.  It consists of two
   paragraphs of text, neither of which are
   particularly interesting.

   This is the second paragraph of the quotation.  It
   is just as dull as the first paragraph.
\end{quotation}
Another frequently-displayed structure is a list.
The following is an example of an \emph{itemized}
list.
\begin{itemize}
   \item This is the first item of an itemized list.
         Each item in the list is marked with a ``tick''.
         You don't have to worry about what kind of tick
         mark is used.

   \item This is the second item of the list.  It
         contains another list nested inside it.  The inner
         list is an \emph{enumerated} list.
         \begin{enumerate}
            \item This is the first item of an enumerated
                  list that is nested within the itemized list.

            \item This is the second item of the inner list.
                  \LaTeX\ allows you to nest lists deeper than
                  you really should.
         \end{enumerate}
         This is the rest of the second item of the outer
         list.  It is no more interesting than any other
         part of the item.
   \item This is the third item of the list.
\end{itemize}
You can even display poetry.
\begin{verse}
   There is an environment
    for verse \\             % The \\ command separates lines
   Whose features some poets % within a stanza.
   will curse.

                             % One or more blank lines separate stanzas.

   For instead of making\\
   Them do \emph{all} line breaking, \\
   It allows them to put too many words on a line when they'd rather be
   forced to be terse.
\end{verse}

Mathematical formulas may also be displayed.  A
displayed formula
is
one-line long; multiline
formulas require special formatting instructions.
   \[  \ip{\Gamma}{\psi'} = x'' + y^{2} + z_{i}^{n}\]
Don't start a paragraph with a displayed equation,
nor make one a paragraph by itself.

\end{document}               % End of document.

\section{Bench Huge Section 1364}
% synthetic bench section 1364
% This is a sample LaTeX input file.  (Version of 12 August 2004.)
%
% A '%' character causes TeX to ignore all remaining text on the line,
% and is used for comments like this one.

\documentclass{article}      % Specifies the document class

                             % The preamble begins here.
\title{An Example Document}  % Declares the document's title.
\author{Leslie Lamport}      % Declares the author's name.
\date{January 21, 1994}      % Deleting this command produces today's date.

\newcommand{\ip}[2]{(#1, #2)}
                             % Defines \ip{arg1}{arg2} to mean
                             % (arg1, arg2).

%\newcommand{\ip}[2]{\langle #1 | #2\rangle}
                             % This is an alternative definition of
                             % \ip that is commented out.

\begin{document}             % End of preamble and beginning of text.

\maketitle                   % Produces the title.

This is an example input file.  Comparing it with
the output it generates can show you how to
produce a simple document of your own.

\section{Ordinary Text}      % Produces section heading.  Lower-level
                             % sections are begun with similar
                             % \subsection and \subsubsection commands.

The ends  of words and sentences are marked
  by   spaces. It  doesn't matter how many
spaces    you type; one is as good as 100.  The
end of   a line counts as a space.

One   or more   blank lines denote the  end
of  a paragraph.

Since any number of consecutive spaces are treated
like a single one, the formatting of the input
file makes no difference to
      \LaTeX,                % The \LaTeX command generates the LaTeX logo.
but it makes a difference to you.  When you use
\LaTeX, making your input file as easy to read
as possible will be a great help as you write
your document and when you change it.  This sample
file shows how you can add comments to your own input
file.

Because printing is different from typewriting,
there are a number of things that you have to do
differently when preparing an input file than if
you were just typing the document directly.
Quotation marks like
       ``this''
have to be handled specially, as do quotes within
quotes:
       ``\,`this'            % \, separates the double and single quote.
        is what I just
        wrote, not  `that'\,''.

Dashes come in three sizes: an
       intra-word
dash, a medium dash for number ranges like
       1--2,
and a punctuation
       dash---like
this.

A sentence-ending space should be larger than the
space between words within a sentence.  You
sometimes have to type special commands in
conjunction with punctuation characters to get
this right, as in the following sentence.
       Gnats, gnus, etc.\ all  % `\ ' makes an inter-word space.
       begin with G\@.         % \@ marks end-of-sentence punctuation.
You should check the spaces after periods when
reading your output to make sure you haven't
forgotten any special cases.  Generating an
ellipsis
       \ldots\               % `\ ' is needed after `\ldots' because TeX
                             % ignores spaces after command names like \ldots
                             % made from \ + letters.
                             %
                             % Note how a `%' character causes TeX to ignore
                             % the end of the input line, so these blank lines
                             % do not start a new paragraph.
                             %
with the right spacing around the periods requires
a special command.

\LaTeX\ interprets some common characters as
commands, so you must type special commands to
generate them.  These characters include the
following:
       \$ \& \% \# \{ and \}.

In printing, text is usually emphasized with an
       \emph{italic}
type style.

\begin{em}
   A long segment of text can also be emphasized
   in this way.  Text within such a segment can be
   given \emph{additional} emphasis.
\end{em}

It is sometimes necessary to prevent \LaTeX\ from
breaking a line where it might otherwise do so.
This may be at a space, as between the ``Mr.''\ and
``Jones'' in
       ``Mr.~Jones'',        % ~ produces an unbreakable interword space.
or within a word---especially when the word is a
symbol like
       \mbox{\emph{itemnum}}
that makes little sense when hyphenated across
lines.

Footnotes\footnote{This is an example of a footnote.}
pose no problem.

\LaTeX\ is good at typesetting mathematical formulas
like
       \( x-3y + z = 7 \)
or
       \( a_{1} > x^{2n} + y^{2n} > x' \)
or
       \( \ip{A}{B} = \sum_{i} a_{i} b_{i} \).
The spaces you type in a formula are
ignored.  Remember that a letter like
       $x$                   % $ ... $  and  \( ... \)  are equivalent
is a formula when it denotes a mathematical
symbol, and it should be typed as one.

\section{Displayed Text}

Text is displayed by indenting it from the left
margin.  Quotations are commonly displayed.  There
are short quotations
\begin{quote}
   This is a short quotation.  It consists of a
   single paragraph of text.  See how it is formatted.
\end{quote}
and longer ones.
\begin{quotation}
   This is a longer quotation.  It consists of two
   paragraphs of text, neither of which are
   particularly interesting.

   This is the second paragraph of the quotation.  It
   is just as dull as the first paragraph.
\end{quotation}
Another frequently-displayed structure is a list.
The following is an example of an \emph{itemized}
list.
\begin{itemize}
   \item This is the first item of an itemized list.
         Each item in the list is marked with a ``tick''.
         You don't have to worry about what kind of tick
         mark is used.

   \item This is the second item of the list.  It
         contains another list nested inside it.  The inner
         list is an \emph{enumerated} list.
         \begin{enumerate}
            \item This is the first item of an enumerated
                  list that is nested within the itemized list.

            \item This is the second item of the inner list.
                  \LaTeX\ allows you to nest lists deeper than
                  you really should.
         \end{enumerate}
         This is the rest of the second item of the outer
         list.  It is no more interesting than any other
         part of the item.
   \item This is the third item of the list.
\end{itemize}
You can even display poetry.
\begin{verse}
   There is an environment
    for verse \\             % The \\ command separates lines
   Whose features some poets % within a stanza.
   will curse.

                             % One or more blank lines separate stanzas.

   For instead of making\\
   Them do \emph{all} line breaking, \\
   It allows them to put too many words on a line when they'd rather be
   forced to be terse.
\end{verse}

Mathematical formulas may also be displayed.  A
displayed formula
is
one-line long; multiline
formulas require special formatting instructions.
   \[  \ip{\Gamma}{\psi'} = x'' + y^{2} + z_{i}^{n}\]
Don't start a paragraph with a displayed equation,
nor make one a paragraph by itself.

\end{document}               % End of document.

\section{Bench Huge Section 1365}
% synthetic bench section 1365
% This is a sample LaTeX input file.  (Version of 12 August 2004.)
%
% A '%' character causes TeX to ignore all remaining text on the line,
% and is used for comments like this one.

\documentclass{article}      % Specifies the document class

                             % The preamble begins here.
\title{An Example Document}  % Declares the document's title.
\author{Leslie Lamport}      % Declares the author's name.
\date{January 21, 1994}      % Deleting this command produces today's date.

\newcommand{\ip}[2]{(#1, #2)}
                             % Defines \ip{arg1}{arg2} to mean
                             % (arg1, arg2).

%\newcommand{\ip}[2]{\langle #1 | #2\rangle}
                             % This is an alternative definition of
                             % \ip that is commented out.

\begin{document}             % End of preamble and beginning of text.

\maketitle                   % Produces the title.

This is an example input file.  Comparing it with
the output it generates can show you how to
produce a simple document of your own.

\section{Ordinary Text}      % Produces section heading.  Lower-level
                             % sections are begun with similar
                             % \subsection and \subsubsection commands.

The ends  of words and sentences are marked
  by   spaces. It  doesn't matter how many
spaces    you type; one is as good as 100.  The
end of   a line counts as a space.

One   or more   blank lines denote the  end
of  a paragraph.

Since any number of consecutive spaces are treated
like a single one, the formatting of the input
file makes no difference to
      \LaTeX,                % The \LaTeX command generates the LaTeX logo.
but it makes a difference to you.  When you use
\LaTeX, making your input file as easy to read
as possible will be a great help as you write
your document and when you change it.  This sample
file shows how you can add comments to your own input
file.

Because printing is different from typewriting,
there are a number of things that you have to do
differently when preparing an input file than if
you were just typing the document directly.
Quotation marks like
       ``this''
have to be handled specially, as do quotes within
quotes:
       ``\,`this'            % \, separates the double and single quote.
        is what I just
        wrote, not  `that'\,''.

Dashes come in three sizes: an
       intra-word
dash, a medium dash for number ranges like
       1--2,
and a punctuation
       dash---like
this.

A sentence-ending space should be larger than the
space between words within a sentence.  You
sometimes have to type special commands in
conjunction with punctuation characters to get
this right, as in the following sentence.
       Gnats, gnus, etc.\ all  % `\ ' makes an inter-word space.
       begin with G\@.         % \@ marks end-of-sentence punctuation.
You should check the spaces after periods when
reading your output to make sure you haven't
forgotten any special cases.  Generating an
ellipsis
       \ldots\               % `\ ' is needed after `\ldots' because TeX
                             % ignores spaces after command names like \ldots
                             % made from \ + letters.
                             %
                             % Note how a `%' character causes TeX to ignore
                             % the end of the input line, so these blank lines
                             % do not start a new paragraph.
                             %
with the right spacing around the periods requires
a special command.

\LaTeX\ interprets some common characters as
commands, so you must type special commands to
generate them.  These characters include the
following:
       \$ \& \% \# \{ and \}.

In printing, text is usually emphasized with an
       \emph{italic}
type style.

\begin{em}
   A long segment of text can also be emphasized
   in this way.  Text within such a segment can be
   given \emph{additional} emphasis.
\end{em}

It is sometimes necessary to prevent \LaTeX\ from
breaking a line where it might otherwise do so.
This may be at a space, as between the ``Mr.''\ and
``Jones'' in
       ``Mr.~Jones'',        % ~ produces an unbreakable interword space.
or within a word---especially when the word is a
symbol like
       \mbox{\emph{itemnum}}
that makes little sense when hyphenated across
lines.

Footnotes\footnote{This is an example of a footnote.}
pose no problem.

\LaTeX\ is good at typesetting mathematical formulas
like
       \( x-3y + z = 7 \)
or
       \( a_{1} > x^{2n} + y^{2n} > x' \)
or
       \( \ip{A}{B} = \sum_{i} a_{i} b_{i} \).
The spaces you type in a formula are
ignored.  Remember that a letter like
       $x$                   % $ ... $  and  \( ... \)  are equivalent
is a formula when it denotes a mathematical
symbol, and it should be typed as one.

\section{Displayed Text}

Text is displayed by indenting it from the left
margin.  Quotations are commonly displayed.  There
are short quotations
\begin{quote}
   This is a short quotation.  It consists of a
   single paragraph of text.  See how it is formatted.
\end{quote}
and longer ones.
\begin{quotation}
   This is a longer quotation.  It consists of two
   paragraphs of text, neither of which are
   particularly interesting.

   This is the second paragraph of the quotation.  It
   is just as dull as the first paragraph.
\end{quotation}
Another frequently-displayed structure is a list.
The following is an example of an \emph{itemized}
list.
\begin{itemize}
   \item This is the first item of an itemized list.
         Each item in the list is marked with a ``tick''.
         You don't have to worry about what kind of tick
         mark is used.

   \item This is the second item of the list.  It
         contains another list nested inside it.  The inner
         list is an \emph{enumerated} list.
         \begin{enumerate}
            \item This is the first item of an enumerated
                  list that is nested within the itemized list.

            \item This is the second item of the inner list.
                  \LaTeX\ allows you to nest lists deeper than
                  you really should.
         \end{enumerate}
         This is the rest of the second item of the outer
         list.  It is no more interesting than any other
         part of the item.
   \item This is the third item of the list.
\end{itemize}
You can even display poetry.
\begin{verse}
   There is an environment
    for verse \\             % The \\ command separates lines
   Whose features some poets % within a stanza.
   will curse.

                             % One or more blank lines separate stanzas.

   For instead of making\\
   Them do \emph{all} line breaking, \\
   It allows them to put too many words on a line when they'd rather be
   forced to be terse.
\end{verse}

Mathematical formulas may also be displayed.  A
displayed formula
is
one-line long; multiline
formulas require special formatting instructions.
   \[  \ip{\Gamma}{\psi'} = x'' + y^{2} + z_{i}^{n}\]
Don't start a paragraph with a displayed equation,
nor make one a paragraph by itself.

\end{document}               % End of document.

\section{Bench Huge Section 1366}
% synthetic bench section 1366
% This is a sample LaTeX input file.  (Version of 12 August 2004.)
%
% A '%' character causes TeX to ignore all remaining text on the line,
% and is used for comments like this one.

\documentclass{article}      % Specifies the document class

                             % The preamble begins here.
\title{An Example Document}  % Declares the document's title.
\author{Leslie Lamport}      % Declares the author's name.
\date{January 21, 1994}      % Deleting this command produces today's date.

\newcommand{\ip}[2]{(#1, #2)}
                             % Defines \ip{arg1}{arg2} to mean
                             % (arg1, arg2).

%\newcommand{\ip}[2]{\langle #1 | #2\rangle}
                             % This is an alternative definition of
                             % \ip that is commented out.

\begin{document}             % End of preamble and beginning of text.

\maketitle                   % Produces the title.

This is an example input file.  Comparing it with
the output it generates can show you how to
produce a simple document of your own.

\section{Ordinary Text}      % Produces section heading.  Lower-level
                             % sections are begun with similar
                             % \subsection and \subsubsection commands.

The ends  of words and sentences are marked
  by   spaces. It  doesn't matter how many
spaces    you type; one is as good as 100.  The
end of   a line counts as a space.

One   or more   blank lines denote the  end
of  a paragraph.

Since any number of consecutive spaces are treated
like a single one, the formatting of the input
file makes no difference to
      \LaTeX,                % The \LaTeX command generates the LaTeX logo.
but it makes a difference to you.  When you use
\LaTeX, making your input file as easy to read
as possible will be a great help as you write
your document and when you change it.  This sample
file shows how you can add comments to your own input
file.

Because printing is different from typewriting,
there are a number of things that you have to do
differently when preparing an input file than if
you were just typing the document directly.
Quotation marks like
       ``this''
have to be handled specially, as do quotes within
quotes:
       ``\,`this'            % \, separates the double and single quote.
        is what I just
        wrote, not  `that'\,''.

Dashes come in three sizes: an
       intra-word
dash, a medium dash for number ranges like
       1--2,
and a punctuation
       dash---like
this.

A sentence-ending space should be larger than the
space between words within a sentence.  You
sometimes have to type special commands in
conjunction with punctuation characters to get
this right, as in the following sentence.
       Gnats, gnus, etc.\ all  % `\ ' makes an inter-word space.
       begin with G\@.         % \@ marks end-of-sentence punctuation.
You should check the spaces after periods when
reading your output to make sure you haven't
forgotten any special cases.  Generating an
ellipsis
       \ldots\               % `\ ' is needed after `\ldots' because TeX
                             % ignores spaces after command names like \ldots
                             % made from \ + letters.
                             %
                             % Note how a `%' character causes TeX to ignore
                             % the end of the input line, so these blank lines
                             % do not start a new paragraph.
                             %
with the right spacing around the periods requires
a special command.

\LaTeX\ interprets some common characters as
commands, so you must type special commands to
generate them.  These characters include the
following:
       \$ \& \% \# \{ and \}.

In printing, text is usually emphasized with an
       \emph{italic}
type style.

\begin{em}
   A long segment of text can also be emphasized
   in this way.  Text within such a segment can be
   given \emph{additional} emphasis.
\end{em}

It is sometimes necessary to prevent \LaTeX\ from
breaking a line where it might otherwise do so.
This may be at a space, as between the ``Mr.''\ and
``Jones'' in
       ``Mr.~Jones'',        % ~ produces an unbreakable interword space.
or within a word---especially when the word is a
symbol like
       \mbox{\emph{itemnum}}
that makes little sense when hyphenated across
lines.

Footnotes\footnote{This is an example of a footnote.}
pose no problem.

\LaTeX\ is good at typesetting mathematical formulas
like
       \( x-3y + z = 7 \)
or
       \( a_{1} > x^{2n} + y^{2n} > x' \)
or
       \( \ip{A}{B} = \sum_{i} a_{i} b_{i} \).
The spaces you type in a formula are
ignored.  Remember that a letter like
       $x$                   % $ ... $  and  \( ... \)  are equivalent
is a formula when it denotes a mathematical
symbol, and it should be typed as one.

\section{Displayed Text}

Text is displayed by indenting it from the left
margin.  Quotations are commonly displayed.  There
are short quotations
\begin{quote}
   This is a short quotation.  It consists of a
   single paragraph of text.  See how it is formatted.
\end{quote}
and longer ones.
\begin{quotation}
   This is a longer quotation.  It consists of two
   paragraphs of text, neither of which are
   particularly interesting.

   This is the second paragraph of the quotation.  It
   is just as dull as the first paragraph.
\end{quotation}
Another frequently-displayed structure is a list.
The following is an example of an \emph{itemized}
list.
\begin{itemize}
   \item This is the first item of an itemized list.
         Each item in the list is marked with a ``tick''.
         You don't have to worry about what kind of tick
         mark is used.

   \item This is the second item of the list.  It
         contains another list nested inside it.  The inner
         list is an \emph{enumerated} list.
         \begin{enumerate}
            \item This is the first item of an enumerated
                  list that is nested within the itemized list.

            \item This is the second item of the inner list.
                  \LaTeX\ allows you to nest lists deeper than
                  you really should.
         \end{enumerate}
         This is the rest of the second item of the outer
         list.  It is no more interesting than any other
         part of the item.
   \item This is the third item of the list.
\end{itemize}
You can even display poetry.
\begin{verse}
   There is an environment
    for verse \\             % The \\ command separates lines
   Whose features some poets % within a stanza.
   will curse.

                             % One or more blank lines separate stanzas.

   For instead of making\\
   Them do \emph{all} line breaking, \\
   It allows them to put too many words on a line when they'd rather be
   forced to be terse.
\end{verse}

Mathematical formulas may also be displayed.  A
displayed formula
is
one-line long; multiline
formulas require special formatting instructions.
   \[  \ip{\Gamma}{\psi'} = x'' + y^{2} + z_{i}^{n}\]
Don't start a paragraph with a displayed equation,
nor make one a paragraph by itself.

\end{document}               % End of document.

\section{Bench Huge Section 1367}
% synthetic bench section 1367
% This is a sample LaTeX input file.  (Version of 12 August 2004.)
%
% A '%' character causes TeX to ignore all remaining text on the line,
% and is used for comments like this one.

\documentclass{article}      % Specifies the document class

                             % The preamble begins here.
\title{An Example Document}  % Declares the document's title.
\author{Leslie Lamport}      % Declares the author's name.
\date{January 21, 1994}      % Deleting this command produces today's date.

\newcommand{\ip}[2]{(#1, #2)}
                             % Defines \ip{arg1}{arg2} to mean
                             % (arg1, arg2).

%\newcommand{\ip}[2]{\langle #1 | #2\rangle}
                             % This is an alternative definition of
                             % \ip that is commented out.

\begin{document}             % End of preamble and beginning of text.

\maketitle                   % Produces the title.

This is an example input file.  Comparing it with
the output it generates can show you how to
produce a simple document of your own.

\section{Ordinary Text}      % Produces section heading.  Lower-level
                             % sections are begun with similar
                             % \subsection and \subsubsection commands.

The ends  of words and sentences are marked
  by   spaces. It  doesn't matter how many
spaces    you type; one is as good as 100.  The
end of   a line counts as a space.

One   or more   blank lines denote the  end
of  a paragraph.

Since any number of consecutive spaces are treated
like a single one, the formatting of the input
file makes no difference to
      \LaTeX,                % The \LaTeX command generates the LaTeX logo.
but it makes a difference to you.  When you use
\LaTeX, making your input file as easy to read
as possible will be a great help as you write
your document and when you change it.  This sample
file shows how you can add comments to your own input
file.

Because printing is different from typewriting,
there are a number of things that you have to do
differently when preparing an input file than if
you were just typing the document directly.
Quotation marks like
       ``this''
have to be handled specially, as do quotes within
quotes:
       ``\,`this'            % \, separates the double and single quote.
        is what I just
        wrote, not  `that'\,''.

Dashes come in three sizes: an
       intra-word
dash, a medium dash for number ranges like
       1--2,
and a punctuation
       dash---like
this.

A sentence-ending space should be larger than the
space between words within a sentence.  You
sometimes have to type special commands in
conjunction with punctuation characters to get
this right, as in the following sentence.
       Gnats, gnus, etc.\ all  % `\ ' makes an inter-word space.
       begin with G\@.         % \@ marks end-of-sentence punctuation.
You should check the spaces after periods when
reading your output to make sure you haven't
forgotten any special cases.  Generating an
ellipsis
       \ldots\               % `\ ' is needed after `\ldots' because TeX
                             % ignores spaces after command names like \ldots
                             % made from \ + letters.
                             %
                             % Note how a `%' character causes TeX to ignore
                             % the end of the input line, so these blank lines
                             % do not start a new paragraph.
                             %
with the right spacing around the periods requires
a special command.

\LaTeX\ interprets some common characters as
commands, so you must type special commands to
generate them.  These characters include the
following:
       \$ \& \% \# \{ and \}.

In printing, text is usually emphasized with an
       \emph{italic}
type style.

\begin{em}
   A long segment of text can also be emphasized
   in this way.  Text within such a segment can be
   given \emph{additional} emphasis.
\end{em}

It is sometimes necessary to prevent \LaTeX\ from
breaking a line where it might otherwise do so.
This may be at a space, as between the ``Mr.''\ and
``Jones'' in
       ``Mr.~Jones'',        % ~ produces an unbreakable interword space.
or within a word---especially when the word is a
symbol like
       \mbox{\emph{itemnum}}
that makes little sense when hyphenated across
lines.

Footnotes\footnote{This is an example of a footnote.}
pose no problem.

\LaTeX\ is good at typesetting mathematical formulas
like
       \( x-3y + z = 7 \)
or
       \( a_{1} > x^{2n} + y^{2n} > x' \)
or
       \( \ip{A}{B} = \sum_{i} a_{i} b_{i} \).
The spaces you type in a formula are
ignored.  Remember that a letter like
       $x$                   % $ ... $  and  \( ... \)  are equivalent
is a formula when it denotes a mathematical
symbol, and it should be typed as one.

\section{Displayed Text}

Text is displayed by indenting it from the left
margin.  Quotations are commonly displayed.  There
are short quotations
\begin{quote}
   This is a short quotation.  It consists of a
   single paragraph of text.  See how it is formatted.
\end{quote}
and longer ones.
\begin{quotation}
   This is a longer quotation.  It consists of two
   paragraphs of text, neither of which are
   particularly interesting.

   This is the second paragraph of the quotation.  It
   is just as dull as the first paragraph.
\end{quotation}
Another frequently-displayed structure is a list.
The following is an example of an \emph{itemized}
list.
\begin{itemize}
   \item This is the first item of an itemized list.
         Each item in the list is marked with a ``tick''.
         You don't have to worry about what kind of tick
         mark is used.

   \item This is the second item of the list.  It
         contains another list nested inside it.  The inner
         list is an \emph{enumerated} list.
         \begin{enumerate}
            \item This is the first item of an enumerated
                  list that is nested within the itemized list.

            \item This is the second item of the inner list.
                  \LaTeX\ allows you to nest lists deeper than
                  you really should.
         \end{enumerate}
         This is the rest of the second item of the outer
         list.  It is no more interesting than any other
         part of the item.
   \item This is the third item of the list.
\end{itemize}
You can even display poetry.
\begin{verse}
   There is an environment
    for verse \\             % The \\ command separates lines
   Whose features some poets % within a stanza.
   will curse.

                             % One or more blank lines separate stanzas.

   For instead of making\\
   Them do \emph{all} line breaking, \\
   It allows them to put too many words on a line when they'd rather be
   forced to be terse.
\end{verse}

Mathematical formulas may also be displayed.  A
displayed formula
is
one-line long; multiline
formulas require special formatting instructions.
   \[  \ip{\Gamma}{\psi'} = x'' + y^{2} + z_{i}^{n}\]
Don't start a paragraph with a displayed equation,
nor make one a paragraph by itself.

\end{document}               % End of document.

\section{Bench Huge Section 1368}
% synthetic bench section 1368
% This is a sample LaTeX input file.  (Version of 12 August 2004.)
%
% A '%' character causes TeX to ignore all remaining text on the line,
% and is used for comments like this one.

\documentclass{article}      % Specifies the document class

                             % The preamble begins here.
\title{An Example Document}  % Declares the document's title.
\author{Leslie Lamport}      % Declares the author's name.
\date{January 21, 1994}      % Deleting this command produces today's date.

\newcommand{\ip}[2]{(#1, #2)}
                             % Defines \ip{arg1}{arg2} to mean
                             % (arg1, arg2).

%\newcommand{\ip}[2]{\langle #1 | #2\rangle}
                             % This is an alternative definition of
                             % \ip that is commented out.

\begin{document}             % End of preamble and beginning of text.

\maketitle                   % Produces the title.

This is an example input file.  Comparing it with
the output it generates can show you how to
produce a simple document of your own.

\section{Ordinary Text}      % Produces section heading.  Lower-level
                             % sections are begun with similar
                             % \subsection and \subsubsection commands.

The ends  of words and sentences are marked
  by   spaces. It  doesn't matter how many
spaces    you type; one is as good as 100.  The
end of   a line counts as a space.

One   or more   blank lines denote the  end
of  a paragraph.

Since any number of consecutive spaces are treated
like a single one, the formatting of the input
file makes no difference to
      \LaTeX,                % The \LaTeX command generates the LaTeX logo.
but it makes a difference to you.  When you use
\LaTeX, making your input file as easy to read
as possible will be a great help as you write
your document and when you change it.  This sample
file shows how you can add comments to your own input
file.

Because printing is different from typewriting,
there are a number of things that you have to do
differently when preparing an input file than if
you were just typing the document directly.
Quotation marks like
       ``this''
have to be handled specially, as do quotes within
quotes:
       ``\,`this'            % \, separates the double and single quote.
        is what I just
        wrote, not  `that'\,''.

Dashes come in three sizes: an
       intra-word
dash, a medium dash for number ranges like
       1--2,
and a punctuation
       dash---like
this.

A sentence-ending space should be larger than the
space between words within a sentence.  You
sometimes have to type special commands in
conjunction with punctuation characters to get
this right, as in the following sentence.
       Gnats, gnus, etc.\ all  % `\ ' makes an inter-word space.
       begin with G\@.         % \@ marks end-of-sentence punctuation.
You should check the spaces after periods when
reading your output to make sure you haven't
forgotten any special cases.  Generating an
ellipsis
       \ldots\               % `\ ' is needed after `\ldots' because TeX
                             % ignores spaces after command names like \ldots
                             % made from \ + letters.
                             %
                             % Note how a `%' character causes TeX to ignore
                             % the end of the input line, so these blank lines
                             % do not start a new paragraph.
                             %
with the right spacing around the periods requires
a special command.

\LaTeX\ interprets some common characters as
commands, so you must type special commands to
generate them.  These characters include the
following:
       \$ \& \% \# \{ and \}.

In printing, text is usually emphasized with an
       \emph{italic}
type style.

\begin{em}
   A long segment of text can also be emphasized
   in this way.  Text within such a segment can be
   given \emph{additional} emphasis.
\end{em}

It is sometimes necessary to prevent \LaTeX\ from
breaking a line where it might otherwise do so.
This may be at a space, as between the ``Mr.''\ and
``Jones'' in
       ``Mr.~Jones'',        % ~ produces an unbreakable interword space.
or within a word---especially when the word is a
symbol like
       \mbox{\emph{itemnum}}
that makes little sense when hyphenated across
lines.

Footnotes\footnote{This is an example of a footnote.}
pose no problem.

\LaTeX\ is good at typesetting mathematical formulas
like
       \( x-3y + z = 7 \)
or
       \( a_{1} > x^{2n} + y^{2n} > x' \)
or
       \( \ip{A}{B} = \sum_{i} a_{i} b_{i} \).
The spaces you type in a formula are
ignored.  Remember that a letter like
       $x$                   % $ ... $  and  \( ... \)  are equivalent
is a formula when it denotes a mathematical
symbol, and it should be typed as one.

\section{Displayed Text}

Text is displayed by indenting it from the left
margin.  Quotations are commonly displayed.  There
are short quotations
\begin{quote}
   This is a short quotation.  It consists of a
   single paragraph of text.  See how it is formatted.
\end{quote}
and longer ones.
\begin{quotation}
   This is a longer quotation.  It consists of two
   paragraphs of text, neither of which are
   particularly interesting.

   This is the second paragraph of the quotation.  It
   is just as dull as the first paragraph.
\end{quotation}
Another frequently-displayed structure is a list.
The following is an example of an \emph{itemized}
list.
\begin{itemize}
   \item This is the first item of an itemized list.
         Each item in the list is marked with a ``tick''.
         You don't have to worry about what kind of tick
         mark is used.

   \item This is the second item of the list.  It
         contains another list nested inside it.  The inner
         list is an \emph{enumerated} list.
         \begin{enumerate}
            \item This is the first item of an enumerated
                  list that is nested within the itemized list.

            \item This is the second item of the inner list.
                  \LaTeX\ allows you to nest lists deeper than
                  you really should.
         \end{enumerate}
         This is the rest of the second item of the outer
         list.  It is no more interesting than any other
         part of the item.
   \item This is the third item of the list.
\end{itemize}
You can even display poetry.
\begin{verse}
   There is an environment
    for verse \\             % The \\ command separates lines
   Whose features some poets % within a stanza.
   will curse.

                             % One or more blank lines separate stanzas.

   For instead of making\\
   Them do \emph{all} line breaking, \\
   It allows them to put too many words on a line when they'd rather be
   forced to be terse.
\end{verse}

Mathematical formulas may also be displayed.  A
displayed formula
is
one-line long; multiline
formulas require special formatting instructions.
   \[  \ip{\Gamma}{\psi'} = x'' + y^{2} + z_{i}^{n}\]
Don't start a paragraph with a displayed equation,
nor make one a paragraph by itself.

\end{document}               % End of document.

\section{Bench Huge Section 1369}
% synthetic bench section 1369
% This is a sample LaTeX input file.  (Version of 12 August 2004.)
%
% A '%' character causes TeX to ignore all remaining text on the line,
% and is used for comments like this one.

\documentclass{article}      % Specifies the document class

                             % The preamble begins here.
\title{An Example Document}  % Declares the document's title.
\author{Leslie Lamport}      % Declares the author's name.
\date{January 21, 1994}      % Deleting this command produces today's date.

\newcommand{\ip}[2]{(#1, #2)}
                             % Defines \ip{arg1}{arg2} to mean
                             % (arg1, arg2).

%\newcommand{\ip}[2]{\langle #1 | #2\rangle}
                             % This is an alternative definition of
                             % \ip that is commented out.

\begin{document}             % End of preamble and beginning of text.

\maketitle                   % Produces the title.

This is an example input file.  Comparing it with
the output it generates can show you how to
produce a simple document of your own.

\section{Ordinary Text}      % Produces section heading.  Lower-level
                             % sections are begun with similar
                             % \subsection and \subsubsection commands.

The ends  of words and sentences are marked
  by   spaces. It  doesn't matter how many
spaces    you type; one is as good as 100.  The
end of   a line counts as a space.

One   or more   blank lines denote the  end
of  a paragraph.

Since any number of consecutive spaces are treated
like a single one, the formatting of the input
file makes no difference to
      \LaTeX,                % The \LaTeX command generates the LaTeX logo.
but it makes a difference to you.  When you use
\LaTeX, making your input file as easy to read
as possible will be a great help as you write
your document and when you change it.  This sample
file shows how you can add comments to your own input
file.

Because printing is different from typewriting,
there are a number of things that you have to do
differently when preparing an input file than if
you were just typing the document directly.
Quotation marks like
       ``this''
have to be handled specially, as do quotes within
quotes:
       ``\,`this'            % \, separates the double and single quote.
        is what I just
        wrote, not  `that'\,''.

Dashes come in three sizes: an
       intra-word
dash, a medium dash for number ranges like
       1--2,
and a punctuation
       dash---like
this.

A sentence-ending space should be larger than the
space between words within a sentence.  You
sometimes have to type special commands in
conjunction with punctuation characters to get
this right, as in the following sentence.
       Gnats, gnus, etc.\ all  % `\ ' makes an inter-word space.
       begin with G\@.         % \@ marks end-of-sentence punctuation.
You should check the spaces after periods when
reading your output to make sure you haven't
forgotten any special cases.  Generating an
ellipsis
       \ldots\               % `\ ' is needed after `\ldots' because TeX
                             % ignores spaces after command names like \ldots
                             % made from \ + letters.
                             %
                             % Note how a `%' character causes TeX to ignore
                             % the end of the input line, so these blank lines
                             % do not start a new paragraph.
                             %
with the right spacing around the periods requires
a special command.

\LaTeX\ interprets some common characters as
commands, so you must type special commands to
generate them.  These characters include the
following:
       \$ \& \% \# \{ and \}.

In printing, text is usually emphasized with an
       \emph{italic}
type style.

\begin{em}
   A long segment of text can also be emphasized
   in this way.  Text within such a segment can be
   given \emph{additional} emphasis.
\end{em}

It is sometimes necessary to prevent \LaTeX\ from
breaking a line where it might otherwise do so.
This may be at a space, as between the ``Mr.''\ and
``Jones'' in
       ``Mr.~Jones'',        % ~ produces an unbreakable interword space.
or within a word---especially when the word is a
symbol like
       \mbox{\emph{itemnum}}
that makes little sense when hyphenated across
lines.

Footnotes\footnote{This is an example of a footnote.}
pose no problem.

\LaTeX\ is good at typesetting mathematical formulas
like
       \( x-3y + z = 7 \)
or
       \( a_{1} > x^{2n} + y^{2n} > x' \)
or
       \( \ip{A}{B} = \sum_{i} a_{i} b_{i} \).
The spaces you type in a formula are
ignored.  Remember that a letter like
       $x$                   % $ ... $  and  \( ... \)  are equivalent
is a formula when it denotes a mathematical
symbol, and it should be typed as one.

\section{Displayed Text}

Text is displayed by indenting it from the left
margin.  Quotations are commonly displayed.  There
are short quotations
\begin{quote}
   This is a short quotation.  It consists of a
   single paragraph of text.  See how it is formatted.
\end{quote}
and longer ones.
\begin{quotation}
   This is a longer quotation.  It consists of two
   paragraphs of text, neither of which are
   particularly interesting.

   This is the second paragraph of the quotation.  It
   is just as dull as the first paragraph.
\end{quotation}
Another frequently-displayed structure is a list.
The following is an example of an \emph{itemized}
list.
\begin{itemize}
   \item This is the first item of an itemized list.
         Each item in the list is marked with a ``tick''.
         You don't have to worry about what kind of tick
         mark is used.

   \item This is the second item of the list.  It
         contains another list nested inside it.  The inner
         list is an \emph{enumerated} list.
         \begin{enumerate}
            \item This is the first item of an enumerated
                  list that is nested within the itemized list.

            \item This is the second item of the inner list.
                  \LaTeX\ allows you to nest lists deeper than
                  you really should.
         \end{enumerate}
         This is the rest of the second item of the outer
         list.  It is no more interesting than any other
         part of the item.
   \item This is the third item of the list.
\end{itemize}
You can even display poetry.
\begin{verse}
   There is an environment
    for verse \\             % The \\ command separates lines
   Whose features some poets % within a stanza.
   will curse.

                             % One or more blank lines separate stanzas.

   For instead of making\\
   Them do \emph{all} line breaking, \\
   It allows them to put too many words on a line when they'd rather be
   forced to be terse.
\end{verse}

Mathematical formulas may also be displayed.  A
displayed formula
is
one-line long; multiline
formulas require special formatting instructions.
   \[  \ip{\Gamma}{\psi'} = x'' + y^{2} + z_{i}^{n}\]
Don't start a paragraph with a displayed equation,
nor make one a paragraph by itself.

\end{document}               % End of document.

\section{Bench Huge Section 1370}
% synthetic bench section 1370
% This is a sample LaTeX input file.  (Version of 12 August 2004.)
%
% A '%' character causes TeX to ignore all remaining text on the line,
% and is used for comments like this one.

\documentclass{article}      % Specifies the document class

                             % The preamble begins here.
\title{An Example Document}  % Declares the document's title.
\author{Leslie Lamport}      % Declares the author's name.
\date{January 21, 1994}      % Deleting this command produces today's date.

\newcommand{\ip}[2]{(#1, #2)}
                             % Defines \ip{arg1}{arg2} to mean
                             % (arg1, arg2).

%\newcommand{\ip}[2]{\langle #1 | #2\rangle}
                             % This is an alternative definition of
                             % \ip that is commented out.

\begin{document}             % End of preamble and beginning of text.

\maketitle                   % Produces the title.

This is an example input file.  Comparing it with
the output it generates can show you how to
produce a simple document of your own.

\section{Ordinary Text}      % Produces section heading.  Lower-level
                             % sections are begun with similar
                             % \subsection and \subsubsection commands.

The ends  of words and sentences are marked
  by   spaces. It  doesn't matter how many
spaces    you type; one is as good as 100.  The
end of   a line counts as a space.

One   or more   blank lines denote the  end
of  a paragraph.

Since any number of consecutive spaces are treated
like a single one, the formatting of the input
file makes no difference to
      \LaTeX,                % The \LaTeX command generates the LaTeX logo.
but it makes a difference to you.  When you use
\LaTeX, making your input file as easy to read
as possible will be a great help as you write
your document and when you change it.  This sample
file shows how you can add comments to your own input
file.

Because printing is different from typewriting,
there are a number of things that you have to do
differently when preparing an input file than if
you were just typing the document directly.
Quotation marks like
       ``this''
have to be handled specially, as do quotes within
quotes:
       ``\,`this'            % \, separates the double and single quote.
        is what I just
        wrote, not  `that'\,''.

Dashes come in three sizes: an
       intra-word
dash, a medium dash for number ranges like
       1--2,
and a punctuation
       dash---like
this.

A sentence-ending space should be larger than the
space between words within a sentence.  You
sometimes have to type special commands in
conjunction with punctuation characters to get
this right, as in the following sentence.
       Gnats, gnus, etc.\ all  % `\ ' makes an inter-word space.
       begin with G\@.         % \@ marks end-of-sentence punctuation.
You should check the spaces after periods when
reading your output to make sure you haven't
forgotten any special cases.  Generating an
ellipsis
       \ldots\               % `\ ' is needed after `\ldots' because TeX
                             % ignores spaces after command names like \ldots
                             % made from \ + letters.
                             %
                             % Note how a `%' character causes TeX to ignore
                             % the end of the input line, so these blank lines
                             % do not start a new paragraph.
                             %
with the right spacing around the periods requires
a special command.

\LaTeX\ interprets some common characters as
commands, so you must type special commands to
generate them.  These characters include the
following:
       \$ \& \% \# \{ and \}.

In printing, text is usually emphasized with an
       \emph{italic}
type style.

\begin{em}
   A long segment of text can also be emphasized
   in this way.  Text within such a segment can be
   given \emph{additional} emphasis.
\end{em}

It is sometimes necessary to prevent \LaTeX\ from
breaking a line where it might otherwise do so.
This may be at a space, as between the ``Mr.''\ and
``Jones'' in
       ``Mr.~Jones'',        % ~ produces an unbreakable interword space.
or within a word---especially when the word is a
symbol like
       \mbox{\emph{itemnum}}
that makes little sense when hyphenated across
lines.

Footnotes\footnote{This is an example of a footnote.}
pose no problem.

\LaTeX\ is good at typesetting mathematical formulas
like
       \( x-3y + z = 7 \)
or
       \( a_{1} > x^{2n} + y^{2n} > x' \)
or
       \( \ip{A}{B} = \sum_{i} a_{i} b_{i} \).
The spaces you type in a formula are
ignored.  Remember that a letter like
       $x$                   % $ ... $  and  \( ... \)  are equivalent
is a formula when it denotes a mathematical
symbol, and it should be typed as one.

\section{Displayed Text}

Text is displayed by indenting it from the left
margin.  Quotations are commonly displayed.  There
are short quotations
\begin{quote}
   This is a short quotation.  It consists of a
   single paragraph of text.  See how it is formatted.
\end{quote}
and longer ones.
\begin{quotation}
   This is a longer quotation.  It consists of two
   paragraphs of text, neither of which are
   particularly interesting.

   This is the second paragraph of the quotation.  It
   is just as dull as the first paragraph.
\end{quotation}
Another frequently-displayed structure is a list.
The following is an example of an \emph{itemized}
list.
\begin{itemize}
   \item This is the first item of an itemized list.
         Each item in the list is marked with a ``tick''.
         You don't have to worry about what kind of tick
         mark is used.

   \item This is the second item of the list.  It
         contains another list nested inside it.  The inner
         list is an \emph{enumerated} list.
         \begin{enumerate}
            \item This is the first item of an enumerated
                  list that is nested within the itemized list.

            \item This is the second item of the inner list.
                  \LaTeX\ allows you to nest lists deeper than
                  you really should.
         \end{enumerate}
         This is the rest of the second item of the outer
         list.  It is no more interesting than any other
         part of the item.
   \item This is the third item of the list.
\end{itemize}
You can even display poetry.
\begin{verse}
   There is an environment
    for verse \\             % The \\ command separates lines
   Whose features some poets % within a stanza.
   will curse.

                             % One or more blank lines separate stanzas.

   For instead of making\\
   Them do \emph{all} line breaking, \\
   It allows them to put too many words on a line when they'd rather be
   forced to be terse.
\end{verse}

Mathematical formulas may also be displayed.  A
displayed formula
is
one-line long; multiline
formulas require special formatting instructions.
   \[  \ip{\Gamma}{\psi'} = x'' + y^{2} + z_{i}^{n}\]
Don't start a paragraph with a displayed equation,
nor make one a paragraph by itself.

\end{document}               % End of document.

\section{Bench Huge Section 1371}
% synthetic bench section 1371
% This is a sample LaTeX input file.  (Version of 12 August 2004.)
%
% A '%' character causes TeX to ignore all remaining text on the line,
% and is used for comments like this one.

\documentclass{article}      % Specifies the document class

                             % The preamble begins here.
\title{An Example Document}  % Declares the document's title.
\author{Leslie Lamport}      % Declares the author's name.
\date{January 21, 1994}      % Deleting this command produces today's date.

\newcommand{\ip}[2]{(#1, #2)}
                             % Defines \ip{arg1}{arg2} to mean
                             % (arg1, arg2).

%\newcommand{\ip}[2]{\langle #1 | #2\rangle}
                             % This is an alternative definition of
                             % \ip that is commented out.

\begin{document}             % End of preamble and beginning of text.

\maketitle                   % Produces the title.

This is an example input file.  Comparing it with
the output it generates can show you how to
produce a simple document of your own.

\section{Ordinary Text}      % Produces section heading.  Lower-level
                             % sections are begun with similar
                             % \subsection and \subsubsection commands.

The ends  of words and sentences are marked
  by   spaces. It  doesn't matter how many
spaces    you type; one is as good as 100.  The
end of   a line counts as a space.

One   or more   blank lines denote the  end
of  a paragraph.

Since any number of consecutive spaces are treated
like a single one, the formatting of the input
file makes no difference to
      \LaTeX,                % The \LaTeX command generates the LaTeX logo.
but it makes a difference to you.  When you use
\LaTeX, making your input file as easy to read
as possible will be a great help as you write
your document and when you change it.  This sample
file shows how you can add comments to your own input
file.

Because printing is different from typewriting,
there are a number of things that you have to do
differently when preparing an input file than if
you were just typing the document directly.
Quotation marks like
       ``this''
have to be handled specially, as do quotes within
quotes:
       ``\,`this'            % \, separates the double and single quote.
        is what I just
        wrote, not  `that'\,''.

Dashes come in three sizes: an
       intra-word
dash, a medium dash for number ranges like
       1--2,
and a punctuation
       dash---like
this.

A sentence-ending space should be larger than the
space between words within a sentence.  You
sometimes have to type special commands in
conjunction with punctuation characters to get
this right, as in the following sentence.
       Gnats, gnus, etc.\ all  % `\ ' makes an inter-word space.
       begin with G\@.         % \@ marks end-of-sentence punctuation.
You should check the spaces after periods when
reading your output to make sure you haven't
forgotten any special cases.  Generating an
ellipsis
       \ldots\               % `\ ' is needed after `\ldots' because TeX
                             % ignores spaces after command names like \ldots
                             % made from \ + letters.
                             %
                             % Note how a `%' character causes TeX to ignore
                             % the end of the input line, so these blank lines
                             % do not start a new paragraph.
                             %
with the right spacing around the periods requires
a special command.

\LaTeX\ interprets some common characters as
commands, so you must type special commands to
generate them.  These characters include the
following:
       \$ \& \% \# \{ and \}.

In printing, text is usually emphasized with an
       \emph{italic}
type style.

\begin{em}
   A long segment of text can also be emphasized
   in this way.  Text within such a segment can be
   given \emph{additional} emphasis.
\end{em}

It is sometimes necessary to prevent \LaTeX\ from
breaking a line where it might otherwise do so.
This may be at a space, as between the ``Mr.''\ and
``Jones'' in
       ``Mr.~Jones'',        % ~ produces an unbreakable interword space.
or within a word---especially when the word is a
symbol like
       \mbox{\emph{itemnum}}
that makes little sense when hyphenated across
lines.

Footnotes\footnote{This is an example of a footnote.}
pose no problem.

\LaTeX\ is good at typesetting mathematical formulas
like
       \( x-3y + z = 7 \)
or
       \( a_{1} > x^{2n} + y^{2n} > x' \)
or
       \( \ip{A}{B} = \sum_{i} a_{i} b_{i} \).
The spaces you type in a formula are
ignored.  Remember that a letter like
       $x$                   % $ ... $  and  \( ... \)  are equivalent
is a formula when it denotes a mathematical
symbol, and it should be typed as one.

\section{Displayed Text}

Text is displayed by indenting it from the left
margin.  Quotations are commonly displayed.  There
are short quotations
\begin{quote}
   This is a short quotation.  It consists of a
   single paragraph of text.  See how it is formatted.
\end{quote}
and longer ones.
\begin{quotation}
   This is a longer quotation.  It consists of two
   paragraphs of text, neither of which are
   particularly interesting.

   This is the second paragraph of the quotation.  It
   is just as dull as the first paragraph.
\end{quotation}
Another frequently-displayed structure is a list.
The following is an example of an \emph{itemized}
list.
\begin{itemize}
   \item This is the first item of an itemized list.
         Each item in the list is marked with a ``tick''.
         You don't have to worry about what kind of tick
         mark is used.

   \item This is the second item of the list.  It
         contains another list nested inside it.  The inner
         list is an \emph{enumerated} list.
         \begin{enumerate}
            \item This is the first item of an enumerated
                  list that is nested within the itemized list.

            \item This is the second item of the inner list.
                  \LaTeX\ allows you to nest lists deeper than
                  you really should.
         \end{enumerate}
         This is the rest of the second item of the outer
         list.  It is no more interesting than any other
         part of the item.
   \item This is the third item of the list.
\end{itemize}
You can even display poetry.
\begin{verse}
   There is an environment
    for verse \\             % The \\ command separates lines
   Whose features some poets % within a stanza.
   will curse.

                             % One or more blank lines separate stanzas.

   For instead of making\\
   Them do \emph{all} line breaking, \\
   It allows them to put too many words on a line when they'd rather be
   forced to be terse.
\end{verse}

Mathematical formulas may also be displayed.  A
displayed formula
is
one-line long; multiline
formulas require special formatting instructions.
   \[  \ip{\Gamma}{\psi'} = x'' + y^{2} + z_{i}^{n}\]
Don't start a paragraph with a displayed equation,
nor make one a paragraph by itself.

\end{document}               % End of document.

\section{Bench Huge Section 1372}
% synthetic bench section 1372
% This is a sample LaTeX input file.  (Version of 12 August 2004.)
%
% A '%' character causes TeX to ignore all remaining text on the line,
% and is used for comments like this one.

\documentclass{article}      % Specifies the document class

                             % The preamble begins here.
\title{An Example Document}  % Declares the document's title.
\author{Leslie Lamport}      % Declares the author's name.
\date{January 21, 1994}      % Deleting this command produces today's date.

\newcommand{\ip}[2]{(#1, #2)}
                             % Defines \ip{arg1}{arg2} to mean
                             % (arg1, arg2).

%\newcommand{\ip}[2]{\langle #1 | #2\rangle}
                             % This is an alternative definition of
                             % \ip that is commented out.

\begin{document}             % End of preamble and beginning of text.

\maketitle                   % Produces the title.

This is an example input file.  Comparing it with
the output it generates can show you how to
produce a simple document of your own.

\section{Ordinary Text}      % Produces section heading.  Lower-level
                             % sections are begun with similar
                             % \subsection and \subsubsection commands.

The ends  of words and sentences are marked
  by   spaces. It  doesn't matter how many
spaces    you type; one is as good as 100.  The
end of   a line counts as a space.

One   or more   blank lines denote the  end
of  a paragraph.

Since any number of consecutive spaces are treated
like a single one, the formatting of the input
file makes no difference to
      \LaTeX,                % The \LaTeX command generates the LaTeX logo.
but it makes a difference to you.  When you use
\LaTeX, making your input file as easy to read
as possible will be a great help as you write
your document and when you change it.  This sample
file shows how you can add comments to your own input
file.

Because printing is different from typewriting,
there are a number of things that you have to do
differently when preparing an input file than if
you were just typing the document directly.
Quotation marks like
       ``this''
have to be handled specially, as do quotes within
quotes:
       ``\,`this'            % \, separates the double and single quote.
        is what I just
        wrote, not  `that'\,''.

Dashes come in three sizes: an
       intra-word
dash, a medium dash for number ranges like
       1--2,
and a punctuation
       dash---like
this.

A sentence-ending space should be larger than the
space between words within a sentence.  You
sometimes have to type special commands in
conjunction with punctuation characters to get
this right, as in the following sentence.
       Gnats, gnus, etc.\ all  % `\ ' makes an inter-word space.
       begin with G\@.         % \@ marks end-of-sentence punctuation.
You should check the spaces after periods when
reading your output to make sure you haven't
forgotten any special cases.  Generating an
ellipsis
       \ldots\               % `\ ' is needed after `\ldots' because TeX
                             % ignores spaces after command names like \ldots
                             % made from \ + letters.
                             %
                             % Note how a `%' character causes TeX to ignore
                             % the end of the input line, so these blank lines
                             % do not start a new paragraph.
                             %
with the right spacing around the periods requires
a special command.

\LaTeX\ interprets some common characters as
commands, so you must type special commands to
generate them.  These characters include the
following:
       \$ \& \% \# \{ and \}.

In printing, text is usually emphasized with an
       \emph{italic}
type style.

\begin{em}
   A long segment of text can also be emphasized
   in this way.  Text within such a segment can be
   given \emph{additional} emphasis.
\end{em}

It is sometimes necessary to prevent \LaTeX\ from
breaking a line where it might otherwise do so.
This may be at a space, as between the ``Mr.''\ and
``Jones'' in
       ``Mr.~Jones'',        % ~ produces an unbreakable interword space.
or within a word---especially when the word is a
symbol like
       \mbox{\emph{itemnum}}
that makes little sense when hyphenated across
lines.

Footnotes\footnote{This is an example of a footnote.}
pose no problem.

\LaTeX\ is good at typesetting mathematical formulas
like
       \( x-3y + z = 7 \)
or
       \( a_{1} > x^{2n} + y^{2n} > x' \)
or
       \( \ip{A}{B} = \sum_{i} a_{i} b_{i} \).
The spaces you type in a formula are
ignored.  Remember that a letter like
       $x$                   % $ ... $  and  \( ... \)  are equivalent
is a formula when it denotes a mathematical
symbol, and it should be typed as one.

\section{Displayed Text}

Text is displayed by indenting it from the left
margin.  Quotations are commonly displayed.  There
are short quotations
\begin{quote}
   This is a short quotation.  It consists of a
   single paragraph of text.  See how it is formatted.
\end{quote}
and longer ones.
\begin{quotation}
   This is a longer quotation.  It consists of two
   paragraphs of text, neither of which are
   particularly interesting.

   This is the second paragraph of the quotation.  It
   is just as dull as the first paragraph.
\end{quotation}
Another frequently-displayed structure is a list.
The following is an example of an \emph{itemized}
list.
\begin{itemize}
   \item This is the first item of an itemized list.
         Each item in the list is marked with a ``tick''.
         You don't have to worry about what kind of tick
         mark is used.

   \item This is the second item of the list.  It
         contains another list nested inside it.  The inner
         list is an \emph{enumerated} list.
         \begin{enumerate}
            \item This is the first item of an enumerated
                  list that is nested within the itemized list.

            \item This is the second item of the inner list.
                  \LaTeX\ allows you to nest lists deeper than
                  you really should.
         \end{enumerate}
         This is the rest of the second item of the outer
         list.  It is no more interesting than any other
         part of the item.
   \item This is the third item of the list.
\end{itemize}
You can even display poetry.
\begin{verse}
   There is an environment
    for verse \\             % The \\ command separates lines
   Whose features some poets % within a stanza.
   will curse.

                             % One or more blank lines separate stanzas.

   For instead of making\\
   Them do \emph{all} line breaking, \\
   It allows them to put too many words on a line when they'd rather be
   forced to be terse.
\end{verse}

Mathematical formulas may also be displayed.  A
displayed formula
is
one-line long; multiline
formulas require special formatting instructions.
   \[  \ip{\Gamma}{\psi'} = x'' + y^{2} + z_{i}^{n}\]
Don't start a paragraph with a displayed equation,
nor make one a paragraph by itself.

\end{document}               % End of document.

\section{Bench Huge Section 1373}
% synthetic bench section 1373
% This is a sample LaTeX input file.  (Version of 12 August 2004.)
%
% A '%' character causes TeX to ignore all remaining text on the line,
% and is used for comments like this one.

\documentclass{article}      % Specifies the document class

                             % The preamble begins here.
\title{An Example Document}  % Declares the document's title.
\author{Leslie Lamport}      % Declares the author's name.
\date{January 21, 1994}      % Deleting this command produces today's date.

\newcommand{\ip}[2]{(#1, #2)}
                             % Defines \ip{arg1}{arg2} to mean
                             % (arg1, arg2).

%\newcommand{\ip}[2]{\langle #1 | #2\rangle}
                             % This is an alternative definition of
                             % \ip that is commented out.

\begin{document}             % End of preamble and beginning of text.

\maketitle                   % Produces the title.

This is an example input file.  Comparing it with
the output it generates can show you how to
produce a simple document of your own.

\section{Ordinary Text}      % Produces section heading.  Lower-level
                             % sections are begun with similar
                             % \subsection and \subsubsection commands.

The ends  of words and sentences are marked
  by   spaces. It  doesn't matter how many
spaces    you type; one is as good as 100.  The
end of   a line counts as a space.

One   or more   blank lines denote the  end
of  a paragraph.

Since any number of consecutive spaces are treated
like a single one, the formatting of the input
file makes no difference to
      \LaTeX,                % The \LaTeX command generates the LaTeX logo.
but it makes a difference to you.  When you use
\LaTeX, making your input file as easy to read
as possible will be a great help as you write
your document and when you change it.  This sample
file shows how you can add comments to your own input
file.

Because printing is different from typewriting,
there are a number of things that you have to do
differently when preparing an input file than if
you were just typing the document directly.
Quotation marks like
       ``this''
have to be handled specially, as do quotes within
quotes:
       ``\,`this'            % \, separates the double and single quote.
        is what I just
        wrote, not  `that'\,''.

Dashes come in three sizes: an
       intra-word
dash, a medium dash for number ranges like
       1--2,
and a punctuation
       dash---like
this.

A sentence-ending space should be larger than the
space between words within a sentence.  You
sometimes have to type special commands in
conjunction with punctuation characters to get
this right, as in the following sentence.
       Gnats, gnus, etc.\ all  % `\ ' makes an inter-word space.
       begin with G\@.         % \@ marks end-of-sentence punctuation.
You should check the spaces after periods when
reading your output to make sure you haven't
forgotten any special cases.  Generating an
ellipsis
       \ldots\               % `\ ' is needed after `\ldots' because TeX
                             % ignores spaces after command names like \ldots
                             % made from \ + letters.
                             %
                             % Note how a `%' character causes TeX to ignore
                             % the end of the input line, so these blank lines
                             % do not start a new paragraph.
                             %
with the right spacing around the periods requires
a special command.

\LaTeX\ interprets some common characters as
commands, so you must type special commands to
generate them.  These characters include the
following:
       \$ \& \% \# \{ and \}.

In printing, text is usually emphasized with an
       \emph{italic}
type style.

\begin{em}
   A long segment of text can also be emphasized
   in this way.  Text within such a segment can be
   given \emph{additional} emphasis.
\end{em}

It is sometimes necessary to prevent \LaTeX\ from
breaking a line where it might otherwise do so.
This may be at a space, as between the ``Mr.''\ and
``Jones'' in
       ``Mr.~Jones'',        % ~ produces an unbreakable interword space.
or within a word---especially when the word is a
symbol like
       \mbox{\emph{itemnum}}
that makes little sense when hyphenated across
lines.

Footnotes\footnote{This is an example of a footnote.}
pose no problem.

\LaTeX\ is good at typesetting mathematical formulas
like
       \( x-3y + z = 7 \)
or
       \( a_{1} > x^{2n} + y^{2n} > x' \)
or
       \( \ip{A}{B} = \sum_{i} a_{i} b_{i} \).
The spaces you type in a formula are
ignored.  Remember that a letter like
       $x$                   % $ ... $  and  \( ... \)  are equivalent
is a formula when it denotes a mathematical
symbol, and it should be typed as one.

\section{Displayed Text}

Text is displayed by indenting it from the left
margin.  Quotations are commonly displayed.  There
are short quotations
\begin{quote}
   This is a short quotation.  It consists of a
   single paragraph of text.  See how it is formatted.
\end{quote}
and longer ones.
\begin{quotation}
   This is a longer quotation.  It consists of two
   paragraphs of text, neither of which are
   particularly interesting.

   This is the second paragraph of the quotation.  It
   is just as dull as the first paragraph.
\end{quotation}
Another frequently-displayed structure is a list.
The following is an example of an \emph{itemized}
list.
\begin{itemize}
   \item This is the first item of an itemized list.
         Each item in the list is marked with a ``tick''.
         You don't have to worry about what kind of tick
         mark is used.

   \item This is the second item of the list.  It
         contains another list nested inside it.  The inner
         list is an \emph{enumerated} list.
         \begin{enumerate}
            \item This is the first item of an enumerated
                  list that is nested within the itemized list.

            \item This is the second item of the inner list.
                  \LaTeX\ allows you to nest lists deeper than
                  you really should.
         \end{enumerate}
         This is the rest of the second item of the outer
         list.  It is no more interesting than any other
         part of the item.
   \item This is the third item of the list.
\end{itemize}
You can even display poetry.
\begin{verse}
   There is an environment
    for verse \\             % The \\ command separates lines
   Whose features some poets % within a stanza.
   will curse.

                             % One or more blank lines separate stanzas.

   For instead of making\\
   Them do \emph{all} line breaking, \\
   It allows them to put too many words on a line when they'd rather be
   forced to be terse.
\end{verse}

Mathematical formulas may also be displayed.  A
displayed formula
is
one-line long; multiline
formulas require special formatting instructions.
   \[  \ip{\Gamma}{\psi'} = x'' + y^{2} + z_{i}^{n}\]
Don't start a paragraph with a displayed equation,
nor make one a paragraph by itself.

\end{document}               % End of document.

\section{Bench Huge Section 1374}
% synthetic bench section 1374
% This is a sample LaTeX input file.  (Version of 12 August 2004.)
%
% A '%' character causes TeX to ignore all remaining text on the line,
% and is used for comments like this one.

\documentclass{article}      % Specifies the document class

                             % The preamble begins here.
\title{An Example Document}  % Declares the document's title.
\author{Leslie Lamport}      % Declares the author's name.
\date{January 21, 1994}      % Deleting this command produces today's date.

\newcommand{\ip}[2]{(#1, #2)}
                             % Defines \ip{arg1}{arg2} to mean
                             % (arg1, arg2).

%\newcommand{\ip}[2]{\langle #1 | #2\rangle}
                             % This is an alternative definition of
                             % \ip that is commented out.

\begin{document}             % End of preamble and beginning of text.

\maketitle                   % Produces the title.

This is an example input file.  Comparing it with
the output it generates can show you how to
produce a simple document of your own.

\section{Ordinary Text}      % Produces section heading.  Lower-level
                             % sections are begun with similar
                             % \subsection and \subsubsection commands.

The ends  of words and sentences are marked
  by   spaces. It  doesn't matter how many
spaces    you type; one is as good as 100.  The
end of   a line counts as a space.

One   or more   blank lines denote the  end
of  a paragraph.

Since any number of consecutive spaces are treated
like a single one, the formatting of the input
file makes no difference to
      \LaTeX,                % The \LaTeX command generates the LaTeX logo.
but it makes a difference to you.  When you use
\LaTeX, making your input file as easy to read
as possible will be a great help as you write
your document and when you change it.  This sample
file shows how you can add comments to your own input
file.

Because printing is different from typewriting,
there are a number of things that you have to do
differently when preparing an input file than if
you were just typing the document directly.
Quotation marks like
       ``this''
have to be handled specially, as do quotes within
quotes:
       ``\,`this'            % \, separates the double and single quote.
        is what I just
        wrote, not  `that'\,''.

Dashes come in three sizes: an
       intra-word
dash, a medium dash for number ranges like
       1--2,
and a punctuation
       dash---like
this.

A sentence-ending space should be larger than the
space between words within a sentence.  You
sometimes have to type special commands in
conjunction with punctuation characters to get
this right, as in the following sentence.
       Gnats, gnus, etc.\ all  % `\ ' makes an inter-word space.
       begin with G\@.         % \@ marks end-of-sentence punctuation.
You should check the spaces after periods when
reading your output to make sure you haven't
forgotten any special cases.  Generating an
ellipsis
       \ldots\               % `\ ' is needed after `\ldots' because TeX
                             % ignores spaces after command names like \ldots
                             % made from \ + letters.
                             %
                             % Note how a `%' character causes TeX to ignore
                             % the end of the input line, so these blank lines
                             % do not start a new paragraph.
                             %
with the right spacing around the periods requires
a special command.

\LaTeX\ interprets some common characters as
commands, so you must type special commands to
generate them.  These characters include the
following:
       \$ \& \% \# \{ and \}.

In printing, text is usually emphasized with an
       \emph{italic}
type style.

\begin{em}
   A long segment of text can also be emphasized
   in this way.  Text within such a segment can be
   given \emph{additional} emphasis.
\end{em}

It is sometimes necessary to prevent \LaTeX\ from
breaking a line where it might otherwise do so.
This may be at a space, as between the ``Mr.''\ and
``Jones'' in
       ``Mr.~Jones'',        % ~ produces an unbreakable interword space.
or within a word---especially when the word is a
symbol like
       \mbox{\emph{itemnum}}
that makes little sense when hyphenated across
lines.

Footnotes\footnote{This is an example of a footnote.}
pose no problem.

\LaTeX\ is good at typesetting mathematical formulas
like
       \( x-3y + z = 7 \)
or
       \( a_{1} > x^{2n} + y^{2n} > x' \)
or
       \( \ip{A}{B} = \sum_{i} a_{i} b_{i} \).
The spaces you type in a formula are
ignored.  Remember that a letter like
       $x$                   % $ ... $  and  \( ... \)  are equivalent
is a formula when it denotes a mathematical
symbol, and it should be typed as one.

\section{Displayed Text}

Text is displayed by indenting it from the left
margin.  Quotations are commonly displayed.  There
are short quotations
\begin{quote}
   This is a short quotation.  It consists of a
   single paragraph of text.  See how it is formatted.
\end{quote}
and longer ones.
\begin{quotation}
   This is a longer quotation.  It consists of two
   paragraphs of text, neither of which are
   particularly interesting.

   This is the second paragraph of the quotation.  It
   is just as dull as the first paragraph.
\end{quotation}
Another frequently-displayed structure is a list.
The following is an example of an \emph{itemized}
list.
\begin{itemize}
   \item This is the first item of an itemized list.
         Each item in the list is marked with a ``tick''.
         You don't have to worry about what kind of tick
         mark is used.

   \item This is the second item of the list.  It
         contains another list nested inside it.  The inner
         list is an \emph{enumerated} list.
         \begin{enumerate}
            \item This is the first item of an enumerated
                  list that is nested within the itemized list.

            \item This is the second item of the inner list.
                  \LaTeX\ allows you to nest lists deeper than
                  you really should.
         \end{enumerate}
         This is the rest of the second item of the outer
         list.  It is no more interesting than any other
         part of the item.
   \item This is the third item of the list.
\end{itemize}
You can even display poetry.
\begin{verse}
   There is an environment
    for verse \\             % The \\ command separates lines
   Whose features some poets % within a stanza.
   will curse.

                             % One or more blank lines separate stanzas.

   For instead of making\\
   Them do \emph{all} line breaking, \\
   It allows them to put too many words on a line when they'd rather be
   forced to be terse.
\end{verse}

Mathematical formulas may also be displayed.  A
displayed formula
is
one-line long; multiline
formulas require special formatting instructions.
   \[  \ip{\Gamma}{\psi'} = x'' + y^{2} + z_{i}^{n}\]
Don't start a paragraph with a displayed equation,
nor make one a paragraph by itself.

\end{document}               % End of document.

\section{Bench Huge Section 1375}
% synthetic bench section 1375
% This is a sample LaTeX input file.  (Version of 12 August 2004.)
%
% A '%' character causes TeX to ignore all remaining text on the line,
% and is used for comments like this one.

\documentclass{article}      % Specifies the document class

                             % The preamble begins here.
\title{An Example Document}  % Declares the document's title.
\author{Leslie Lamport}      % Declares the author's name.
\date{January 21, 1994}      % Deleting this command produces today's date.

\newcommand{\ip}[2]{(#1, #2)}
                             % Defines \ip{arg1}{arg2} to mean
                             % (arg1, arg2).

%\newcommand{\ip}[2]{\langle #1 | #2\rangle}
                             % This is an alternative definition of
                             % \ip that is commented out.

\begin{document}             % End of preamble and beginning of text.

\maketitle                   % Produces the title.

This is an example input file.  Comparing it with
the output it generates can show you how to
produce a simple document of your own.

\section{Ordinary Text}      % Produces section heading.  Lower-level
                             % sections are begun with similar
                             % \subsection and \subsubsection commands.

The ends  of words and sentences are marked
  by   spaces. It  doesn't matter how many
spaces    you type; one is as good as 100.  The
end of   a line counts as a space.

One   or more   blank lines denote the  end
of  a paragraph.

Since any number of consecutive spaces are treated
like a single one, the formatting of the input
file makes no difference to
      \LaTeX,                % The \LaTeX command generates the LaTeX logo.
but it makes a difference to you.  When you use
\LaTeX, making your input file as easy to read
as possible will be a great help as you write
your document and when you change it.  This sample
file shows how you can add comments to your own input
file.

Because printing is different from typewriting,
there are a number of things that you have to do
differently when preparing an input file than if
you were just typing the document directly.
Quotation marks like
       ``this''
have to be handled specially, as do quotes within
quotes:
       ``\,`this'            % \, separates the double and single quote.
        is what I just
        wrote, not  `that'\,''.

Dashes come in three sizes: an
       intra-word
dash, a medium dash for number ranges like
       1--2,
and a punctuation
       dash---like
this.

A sentence-ending space should be larger than the
space between words within a sentence.  You
sometimes have to type special commands in
conjunction with punctuation characters to get
this right, as in the following sentence.
       Gnats, gnus, etc.\ all  % `\ ' makes an inter-word space.
       begin with G\@.         % \@ marks end-of-sentence punctuation.
You should check the spaces after periods when
reading your output to make sure you haven't
forgotten any special cases.  Generating an
ellipsis
       \ldots\               % `\ ' is needed after `\ldots' because TeX
                             % ignores spaces after command names like \ldots
                             % made from \ + letters.
                             %
                             % Note how a `%' character causes TeX to ignore
                             % the end of the input line, so these blank lines
                             % do not start a new paragraph.
                             %
with the right spacing around the periods requires
a special command.

\LaTeX\ interprets some common characters as
commands, so you must type special commands to
generate them.  These characters include the
following:
       \$ \& \% \# \{ and \}.

In printing, text is usually emphasized with an
       \emph{italic}
type style.

\begin{em}
   A long segment of text can also be emphasized
   in this way.  Text within such a segment can be
   given \emph{additional} emphasis.
\end{em}

It is sometimes necessary to prevent \LaTeX\ from
breaking a line where it might otherwise do so.
This may be at a space, as between the ``Mr.''\ and
``Jones'' in
       ``Mr.~Jones'',        % ~ produces an unbreakable interword space.
or within a word---especially when the word is a
symbol like
       \mbox{\emph{itemnum}}
that makes little sense when hyphenated across
lines.

Footnotes\footnote{This is an example of a footnote.}
pose no problem.

\LaTeX\ is good at typesetting mathematical formulas
like
       \( x-3y + z = 7 \)
or
       \( a_{1} > x^{2n} + y^{2n} > x' \)
or
       \( \ip{A}{B} = \sum_{i} a_{i} b_{i} \).
The spaces you type in a formula are
ignored.  Remember that a letter like
       $x$                   % $ ... $  and  \( ... \)  are equivalent
is a formula when it denotes a mathematical
symbol, and it should be typed as one.

\section{Displayed Text}

Text is displayed by indenting it from the left
margin.  Quotations are commonly displayed.  There
are short quotations
\begin{quote}
   This is a short quotation.  It consists of a
   single paragraph of text.  See how it is formatted.
\end{quote}
and longer ones.
\begin{quotation}
   This is a longer quotation.  It consists of two
   paragraphs of text, neither of which are
   particularly interesting.

   This is the second paragraph of the quotation.  It
   is just as dull as the first paragraph.
\end{quotation}
Another frequently-displayed structure is a list.
The following is an example of an \emph{itemized}
list.
\begin{itemize}
   \item This is the first item of an itemized list.
         Each item in the list is marked with a ``tick''.
         You don't have to worry about what kind of tick
         mark is used.

   \item This is the second item of the list.  It
         contains another list nested inside it.  The inner
         list is an \emph{enumerated} list.
         \begin{enumerate}
            \item This is the first item of an enumerated
                  list that is nested within the itemized list.

            \item This is the second item of the inner list.
                  \LaTeX\ allows you to nest lists deeper than
                  you really should.
         \end{enumerate}
         This is the rest of the second item of the outer
         list.  It is no more interesting than any other
         part of the item.
   \item This is the third item of the list.
\end{itemize}
You can even display poetry.
\begin{verse}
   There is an environment
    for verse \\             % The \\ command separates lines
   Whose features some poets % within a stanza.
   will curse.

                             % One or more blank lines separate stanzas.

   For instead of making\\
   Them do \emph{all} line breaking, \\
   It allows them to put too many words on a line when they'd rather be
   forced to be terse.
\end{verse}

Mathematical formulas may also be displayed.  A
displayed formula
is
one-line long; multiline
formulas require special formatting instructions.
   \[  \ip{\Gamma}{\psi'} = x'' + y^{2} + z_{i}^{n}\]
Don't start a paragraph with a displayed equation,
nor make one a paragraph by itself.

\end{document}               % End of document.

\section{Bench Huge Section 1376}
% synthetic bench section 1376
% This is a sample LaTeX input file.  (Version of 12 August 2004.)
%
% A '%' character causes TeX to ignore all remaining text on the line,
% and is used for comments like this one.

\documentclass{article}      % Specifies the document class

                             % The preamble begins here.
\title{An Example Document}  % Declares the document's title.
\author{Leslie Lamport}      % Declares the author's name.
\date{January 21, 1994}      % Deleting this command produces today's date.

\newcommand{\ip}[2]{(#1, #2)}
                             % Defines \ip{arg1}{arg2} to mean
                             % (arg1, arg2).

%\newcommand{\ip}[2]{\langle #1 | #2\rangle}
                             % This is an alternative definition of
                             % \ip that is commented out.

\begin{document}             % End of preamble and beginning of text.

\maketitle                   % Produces the title.

This is an example input file.  Comparing it with
the output it generates can show you how to
produce a simple document of your own.

\section{Ordinary Text}      % Produces section heading.  Lower-level
                             % sections are begun with similar
                             % \subsection and \subsubsection commands.

The ends  of words and sentences are marked
  by   spaces. It  doesn't matter how many
spaces    you type; one is as good as 100.  The
end of   a line counts as a space.

One   or more   blank lines denote the  end
of  a paragraph.

Since any number of consecutive spaces are treated
like a single one, the formatting of the input
file makes no difference to
      \LaTeX,                % The \LaTeX command generates the LaTeX logo.
but it makes a difference to you.  When you use
\LaTeX, making your input file as easy to read
as possible will be a great help as you write
your document and when you change it.  This sample
file shows how you can add comments to your own input
file.

Because printing is different from typewriting,
there are a number of things that you have to do
differently when preparing an input file than if
you were just typing the document directly.
Quotation marks like
       ``this''
have to be handled specially, as do quotes within
quotes:
       ``\,`this'            % \, separates the double and single quote.
        is what I just
        wrote, not  `that'\,''.

Dashes come in three sizes: an
       intra-word
dash, a medium dash for number ranges like
       1--2,
and a punctuation
       dash---like
this.

A sentence-ending space should be larger than the
space between words within a sentence.  You
sometimes have to type special commands in
conjunction with punctuation characters to get
this right, as in the following sentence.
       Gnats, gnus, etc.\ all  % `\ ' makes an inter-word space.
       begin with G\@.         % \@ marks end-of-sentence punctuation.
You should check the spaces after periods when
reading your output to make sure you haven't
forgotten any special cases.  Generating an
ellipsis
       \ldots\               % `\ ' is needed after `\ldots' because TeX
                             % ignores spaces after command names like \ldots
                             % made from \ + letters.
                             %
                             % Note how a `%' character causes TeX to ignore
                             % the end of the input line, so these blank lines
                             % do not start a new paragraph.
                             %
with the right spacing around the periods requires
a special command.

\LaTeX\ interprets some common characters as
commands, so you must type special commands to
generate them.  These characters include the
following:
       \$ \& \% \# \{ and \}.

In printing, text is usually emphasized with an
       \emph{italic}
type style.

\begin{em}
   A long segment of text can also be emphasized
   in this way.  Text within such a segment can be
   given \emph{additional} emphasis.
\end{em}

It is sometimes necessary to prevent \LaTeX\ from
breaking a line where it might otherwise do so.
This may be at a space, as between the ``Mr.''\ and
``Jones'' in
       ``Mr.~Jones'',        % ~ produces an unbreakable interword space.
or within a word---especially when the word is a
symbol like
       \mbox{\emph{itemnum}}
that makes little sense when hyphenated across
lines.

Footnotes\footnote{This is an example of a footnote.}
pose no problem.

\LaTeX\ is good at typesetting mathematical formulas
like
       \( x-3y + z = 7 \)
or
       \( a_{1} > x^{2n} + y^{2n} > x' \)
or
       \( \ip{A}{B} = \sum_{i} a_{i} b_{i} \).
The spaces you type in a formula are
ignored.  Remember that a letter like
       $x$                   % $ ... $  and  \( ... \)  are equivalent
is a formula when it denotes a mathematical
symbol, and it should be typed as one.

\section{Displayed Text}

Text is displayed by indenting it from the left
margin.  Quotations are commonly displayed.  There
are short quotations
\begin{quote}
   This is a short quotation.  It consists of a
   single paragraph of text.  See how it is formatted.
\end{quote}
and longer ones.
\begin{quotation}
   This is a longer quotation.  It consists of two
   paragraphs of text, neither of which are
   particularly interesting.

   This is the second paragraph of the quotation.  It
   is just as dull as the first paragraph.
\end{quotation}
Another frequently-displayed structure is a list.
The following is an example of an \emph{itemized}
list.
\begin{itemize}
   \item This is the first item of an itemized list.
         Each item in the list is marked with a ``tick''.
         You don't have to worry about what kind of tick
         mark is used.

   \item This is the second item of the list.  It
         contains another list nested inside it.  The inner
         list is an \emph{enumerated} list.
         \begin{enumerate}
            \item This is the first item of an enumerated
                  list that is nested within the itemized list.

            \item This is the second item of the inner list.
                  \LaTeX\ allows you to nest lists deeper than
                  you really should.
         \end{enumerate}
         This is the rest of the second item of the outer
         list.  It is no more interesting than any other
         part of the item.
   \item This is the third item of the list.
\end{itemize}
You can even display poetry.
\begin{verse}
   There is an environment
    for verse \\             % The \\ command separates lines
   Whose features some poets % within a stanza.
   will curse.

                             % One or more blank lines separate stanzas.

   For instead of making\\
   Them do \emph{all} line breaking, \\
   It allows them to put too many words on a line when they'd rather be
   forced to be terse.
\end{verse}

Mathematical formulas may also be displayed.  A
displayed formula
is
one-line long; multiline
formulas require special formatting instructions.
   \[  \ip{\Gamma}{\psi'} = x'' + y^{2} + z_{i}^{n}\]
Don't start a paragraph with a displayed equation,
nor make one a paragraph by itself.

\end{document}               % End of document.

\section{Bench Huge Section 1377}
% synthetic bench section 1377
% This is a sample LaTeX input file.  (Version of 12 August 2004.)
%
% A '%' character causes TeX to ignore all remaining text on the line,
% and is used for comments like this one.

\documentclass{article}      % Specifies the document class

                             % The preamble begins here.
\title{An Example Document}  % Declares the document's title.
\author{Leslie Lamport}      % Declares the author's name.
\date{January 21, 1994}      % Deleting this command produces today's date.

\newcommand{\ip}[2]{(#1, #2)}
                             % Defines \ip{arg1}{arg2} to mean
                             % (arg1, arg2).

%\newcommand{\ip}[2]{\langle #1 | #2\rangle}
                             % This is an alternative definition of
                             % \ip that is commented out.

\begin{document}             % End of preamble and beginning of text.

\maketitle                   % Produces the title.

This is an example input file.  Comparing it with
the output it generates can show you how to
produce a simple document of your own.

\section{Ordinary Text}      % Produces section heading.  Lower-level
                             % sections are begun with similar
                             % \subsection and \subsubsection commands.

The ends  of words and sentences are marked
  by   spaces. It  doesn't matter how many
spaces    you type; one is as good as 100.  The
end of   a line counts as a space.

One   or more   blank lines denote the  end
of  a paragraph.

Since any number of consecutive spaces are treated
like a single one, the formatting of the input
file makes no difference to
      \LaTeX,                % The \LaTeX command generates the LaTeX logo.
but it makes a difference to you.  When you use
\LaTeX, making your input file as easy to read
as possible will be a great help as you write
your document and when you change it.  This sample
file shows how you can add comments to your own input
file.

Because printing is different from typewriting,
there are a number of things that you have to do
differently when preparing an input file than if
you were just typing the document directly.
Quotation marks like
       ``this''
have to be handled specially, as do quotes within
quotes:
       ``\,`this'            % \, separates the double and single quote.
        is what I just
        wrote, not  `that'\,''.

Dashes come in three sizes: an
       intra-word
dash, a medium dash for number ranges like
       1--2,
and a punctuation
       dash---like
this.

A sentence-ending space should be larger than the
space between words within a sentence.  You
sometimes have to type special commands in
conjunction with punctuation characters to get
this right, as in the following sentence.
       Gnats, gnus, etc.\ all  % `\ ' makes an inter-word space.
       begin with G\@.         % \@ marks end-of-sentence punctuation.
You should check the spaces after periods when
reading your output to make sure you haven't
forgotten any special cases.  Generating an
ellipsis
       \ldots\               % `\ ' is needed after `\ldots' because TeX
                             % ignores spaces after command names like \ldots
                             % made from \ + letters.
                             %
                             % Note how a `%' character causes TeX to ignore
                             % the end of the input line, so these blank lines
                             % do not start a new paragraph.
                             %
with the right spacing around the periods requires
a special command.

\LaTeX\ interprets some common characters as
commands, so you must type special commands to
generate them.  These characters include the
following:
       \$ \& \% \# \{ and \}.

In printing, text is usually emphasized with an
       \emph{italic}
type style.

\begin{em}
   A long segment of text can also be emphasized
   in this way.  Text within such a segment can be
   given \emph{additional} emphasis.
\end{em}

It is sometimes necessary to prevent \LaTeX\ from
breaking a line where it might otherwise do so.
This may be at a space, as between the ``Mr.''\ and
``Jones'' in
       ``Mr.~Jones'',        % ~ produces an unbreakable interword space.
or within a word---especially when the word is a
symbol like
       \mbox{\emph{itemnum}}
that makes little sense when hyphenated across
lines.

Footnotes\footnote{This is an example of a footnote.}
pose no problem.

\LaTeX\ is good at typesetting mathematical formulas
like
       \( x-3y + z = 7 \)
or
       \( a_{1} > x^{2n} + y^{2n} > x' \)
or
       \( \ip{A}{B} = \sum_{i} a_{i} b_{i} \).
The spaces you type in a formula are
ignored.  Remember that a letter like
       $x$                   % $ ... $  and  \( ... \)  are equivalent
is a formula when it denotes a mathematical
symbol, and it should be typed as one.

\section{Displayed Text}

Text is displayed by indenting it from the left
margin.  Quotations are commonly displayed.  There
are short quotations
\begin{quote}
   This is a short quotation.  It consists of a
   single paragraph of text.  See how it is formatted.
\end{quote}
and longer ones.
\begin{quotation}
   This is a longer quotation.  It consists of two
   paragraphs of text, neither of which are
   particularly interesting.

   This is the second paragraph of the quotation.  It
   is just as dull as the first paragraph.
\end{quotation}
Another frequently-displayed structure is a list.
The following is an example of an \emph{itemized}
list.
\begin{itemize}
   \item This is the first item of an itemized list.
         Each item in the list is marked with a ``tick''.
         You don't have to worry about what kind of tick
         mark is used.

   \item This is the second item of the list.  It
         contains another list nested inside it.  The inner
         list is an \emph{enumerated} list.
         \begin{enumerate}
            \item This is the first item of an enumerated
                  list that is nested within the itemized list.

            \item This is the second item of the inner list.
                  \LaTeX\ allows you to nest lists deeper than
                  you really should.
         \end{enumerate}
         This is the rest of the second item of the outer
         list.  It is no more interesting than any other
         part of the item.
   \item This is the third item of the list.
\end{itemize}
You can even display poetry.
\begin{verse}
   There is an environment
    for verse \\             % The \\ command separates lines
   Whose features some poets % within a stanza.
   will curse.

                             % One or more blank lines separate stanzas.

   For instead of making\\
   Them do \emph{all} line breaking, \\
   It allows them to put too many words on a line when they'd rather be
   forced to be terse.
\end{verse}

Mathematical formulas may also be displayed.  A
displayed formula
is
one-line long; multiline
formulas require special formatting instructions.
   \[  \ip{\Gamma}{\psi'} = x'' + y^{2} + z_{i}^{n}\]
Don't start a paragraph with a displayed equation,
nor make one a paragraph by itself.

\end{document}               % End of document.

\section{Bench Huge Section 1378}
% synthetic bench section 1378
% This is a sample LaTeX input file.  (Version of 12 August 2004.)
%
% A '%' character causes TeX to ignore all remaining text on the line,
% and is used for comments like this one.

\documentclass{article}      % Specifies the document class

                             % The preamble begins here.
\title{An Example Document}  % Declares the document's title.
\author{Leslie Lamport}      % Declares the author's name.
\date{January 21, 1994}      % Deleting this command produces today's date.

\newcommand{\ip}[2]{(#1, #2)}
                             % Defines \ip{arg1}{arg2} to mean
                             % (arg1, arg2).

%\newcommand{\ip}[2]{\langle #1 | #2\rangle}
                             % This is an alternative definition of
                             % \ip that is commented out.

\begin{document}             % End of preamble and beginning of text.

\maketitle                   % Produces the title.

This is an example input file.  Comparing it with
the output it generates can show you how to
produce a simple document of your own.

\section{Ordinary Text}      % Produces section heading.  Lower-level
                             % sections are begun with similar
                             % \subsection and \subsubsection commands.

The ends  of words and sentences are marked
  by   spaces. It  doesn't matter how many
spaces    you type; one is as good as 100.  The
end of   a line counts as a space.

One   or more   blank lines denote the  end
of  a paragraph.

Since any number of consecutive spaces are treated
like a single one, the formatting of the input
file makes no difference to
      \LaTeX,                % The \LaTeX command generates the LaTeX logo.
but it makes a difference to you.  When you use
\LaTeX, making your input file as easy to read
as possible will be a great help as you write
your document and when you change it.  This sample
file shows how you can add comments to your own input
file.

Because printing is different from typewriting,
there are a number of things that you have to do
differently when preparing an input file than if
you were just typing the document directly.
Quotation marks like
       ``this''
have to be handled specially, as do quotes within
quotes:
       ``\,`this'            % \, separates the double and single quote.
        is what I just
        wrote, not  `that'\,''.

Dashes come in three sizes: an
       intra-word
dash, a medium dash for number ranges like
       1--2,
and a punctuation
       dash---like
this.

A sentence-ending space should be larger than the
space between words within a sentence.  You
sometimes have to type special commands in
conjunction with punctuation characters to get
this right, as in the following sentence.
       Gnats, gnus, etc.\ all  % `\ ' makes an inter-word space.
       begin with G\@.         % \@ marks end-of-sentence punctuation.
You should check the spaces after periods when
reading your output to make sure you haven't
forgotten any special cases.  Generating an
ellipsis
       \ldots\               % `\ ' is needed after `\ldots' because TeX
                             % ignores spaces after command names like \ldots
                             % made from \ + letters.
                             %
                             % Note how a `%' character causes TeX to ignore
                             % the end of the input line, so these blank lines
                             % do not start a new paragraph.
                             %
with the right spacing around the periods requires
a special command.

\LaTeX\ interprets some common characters as
commands, so you must type special commands to
generate them.  These characters include the
following:
       \$ \& \% \# \{ and \}.

In printing, text is usually emphasized with an
       \emph{italic}
type style.

\begin{em}
   A long segment of text can also be emphasized
   in this way.  Text within such a segment can be
   given \emph{additional} emphasis.
\end{em}

It is sometimes necessary to prevent \LaTeX\ from
breaking a line where it might otherwise do so.
This may be at a space, as between the ``Mr.''\ and
``Jones'' in
       ``Mr.~Jones'',        % ~ produces an unbreakable interword space.
or within a word---especially when the word is a
symbol like
       \mbox{\emph{itemnum}}
that makes little sense when hyphenated across
lines.

Footnotes\footnote{This is an example of a footnote.}
pose no problem.

\LaTeX\ is good at typesetting mathematical formulas
like
       \( x-3y + z = 7 \)
or
       \( a_{1} > x^{2n} + y^{2n} > x' \)
or
       \( \ip{A}{B} = \sum_{i} a_{i} b_{i} \).
The spaces you type in a formula are
ignored.  Remember that a letter like
       $x$                   % $ ... $  and  \( ... \)  are equivalent
is a formula when it denotes a mathematical
symbol, and it should be typed as one.

\section{Displayed Text}

Text is displayed by indenting it from the left
margin.  Quotations are commonly displayed.  There
are short quotations
\begin{quote}
   This is a short quotation.  It consists of a
   single paragraph of text.  See how it is formatted.
\end{quote}
and longer ones.
\begin{quotation}
   This is a longer quotation.  It consists of two
   paragraphs of text, neither of which are
   particularly interesting.

   This is the second paragraph of the quotation.  It
   is just as dull as the first paragraph.
\end{quotation}
Another frequently-displayed structure is a list.
The following is an example of an \emph{itemized}
list.
\begin{itemize}
   \item This is the first item of an itemized list.
         Each item in the list is marked with a ``tick''.
         You don't have to worry about what kind of tick
         mark is used.

   \item This is the second item of the list.  It
         contains another list nested inside it.  The inner
         list is an \emph{enumerated} list.
         \begin{enumerate}
            \item This is the first item of an enumerated
                  list that is nested within the itemized list.

            \item This is the second item of the inner list.
                  \LaTeX\ allows you to nest lists deeper than
                  you really should.
         \end{enumerate}
         This is the rest of the second item of the outer
         list.  It is no more interesting than any other
         part of the item.
   \item This is the third item of the list.
\end{itemize}
You can even display poetry.
\begin{verse}
   There is an environment
    for verse \\             % The \\ command separates lines
   Whose features some poets % within a stanza.
   will curse.

                             % One or more blank lines separate stanzas.

   For instead of making\\
   Them do \emph{all} line breaking, \\
   It allows them to put too many words on a line when they'd rather be
   forced to be terse.
\end{verse}

Mathematical formulas may also be displayed.  A
displayed formula
is
one-line long; multiline
formulas require special formatting instructions.
   \[  \ip{\Gamma}{\psi'} = x'' + y^{2} + z_{i}^{n}\]
Don't start a paragraph with a displayed equation,
nor make one a paragraph by itself.

\end{document}               % End of document.

\section{Bench Huge Section 1379}
% synthetic bench section 1379
% This is a sample LaTeX input file.  (Version of 12 August 2004.)
%
% A '%' character causes TeX to ignore all remaining text on the line,
% and is used for comments like this one.

\documentclass{article}      % Specifies the document class

                             % The preamble begins here.
\title{An Example Document}  % Declares the document's title.
\author{Leslie Lamport}      % Declares the author's name.
\date{January 21, 1994}      % Deleting this command produces today's date.

\newcommand{\ip}[2]{(#1, #2)}
                             % Defines \ip{arg1}{arg2} to mean
                             % (arg1, arg2).

%\newcommand{\ip}[2]{\langle #1 | #2\rangle}
                             % This is an alternative definition of
                             % \ip that is commented out.

\begin{document}             % End of preamble and beginning of text.

\maketitle                   % Produces the title.

This is an example input file.  Comparing it with
the output it generates can show you how to
produce a simple document of your own.

\section{Ordinary Text}      % Produces section heading.  Lower-level
                             % sections are begun with similar
                             % \subsection and \subsubsection commands.

The ends  of words and sentences are marked
  by   spaces. It  doesn't matter how many
spaces    you type; one is as good as 100.  The
end of   a line counts as a space.

One   or more   blank lines denote the  end
of  a paragraph.

Since any number of consecutive spaces are treated
like a single one, the formatting of the input
file makes no difference to
      \LaTeX,                % The \LaTeX command generates the LaTeX logo.
but it makes a difference to you.  When you use
\LaTeX, making your input file as easy to read
as possible will be a great help as you write
your document and when you change it.  This sample
file shows how you can add comments to your own input
file.

Because printing is different from typewriting,
there are a number of things that you have to do
differently when preparing an input file than if
you were just typing the document directly.
Quotation marks like
       ``this''
have to be handled specially, as do quotes within
quotes:
       ``\,`this'            % \, separates the double and single quote.
        is what I just
        wrote, not  `that'\,''.

Dashes come in three sizes: an
       intra-word
dash, a medium dash for number ranges like
       1--2,
and a punctuation
       dash---like
this.

A sentence-ending space should be larger than the
space between words within a sentence.  You
sometimes have to type special commands in
conjunction with punctuation characters to get
this right, as in the following sentence.
       Gnats, gnus, etc.\ all  % `\ ' makes an inter-word space.
       begin with G\@.         % \@ marks end-of-sentence punctuation.
You should check the spaces after periods when
reading your output to make sure you haven't
forgotten any special cases.  Generating an
ellipsis
       \ldots\               % `\ ' is needed after `\ldots' because TeX
                             % ignores spaces after command names like \ldots
                             % made from \ + letters.
                             %
                             % Note how a `%' character causes TeX to ignore
                             % the end of the input line, so these blank lines
                             % do not start a new paragraph.
                             %
with the right spacing around the periods requires
a special command.

\LaTeX\ interprets some common characters as
commands, so you must type special commands to
generate them.  These characters include the
following:
       \$ \& \% \# \{ and \}.

In printing, text is usually emphasized with an
       \emph{italic}
type style.

\begin{em}
   A long segment of text can also be emphasized
   in this way.  Text within such a segment can be
   given \emph{additional} emphasis.
\end{em}

It is sometimes necessary to prevent \LaTeX\ from
breaking a line where it might otherwise do so.
This may be at a space, as between the ``Mr.''\ and
``Jones'' in
       ``Mr.~Jones'',        % ~ produces an unbreakable interword space.
or within a word---especially when the word is a
symbol like
       \mbox{\emph{itemnum}}
that makes little sense when hyphenated across
lines.

Footnotes\footnote{This is an example of a footnote.}
pose no problem.

\LaTeX\ is good at typesetting mathematical formulas
like
       \( x-3y + z = 7 \)
or
       \( a_{1} > x^{2n} + y^{2n} > x' \)
or
       \( \ip{A}{B} = \sum_{i} a_{i} b_{i} \).
The spaces you type in a formula are
ignored.  Remember that a letter like
       $x$                   % $ ... $  and  \( ... \)  are equivalent
is a formula when it denotes a mathematical
symbol, and it should be typed as one.

\section{Displayed Text}

Text is displayed by indenting it from the left
margin.  Quotations are commonly displayed.  There
are short quotations
\begin{quote}
   This is a short quotation.  It consists of a
   single paragraph of text.  See how it is formatted.
\end{quote}
and longer ones.
\begin{quotation}
   This is a longer quotation.  It consists of two
   paragraphs of text, neither of which are
   particularly interesting.

   This is the second paragraph of the quotation.  It
   is just as dull as the first paragraph.
\end{quotation}
Another frequently-displayed structure is a list.
The following is an example of an \emph{itemized}
list.
\begin{itemize}
   \item This is the first item of an itemized list.
         Each item in the list is marked with a ``tick''.
         You don't have to worry about what kind of tick
         mark is used.

   \item This is the second item of the list.  It
         contains another list nested inside it.  The inner
         list is an \emph{enumerated} list.
         \begin{enumerate}
            \item This is the first item of an enumerated
                  list that is nested within the itemized list.

            \item This is the second item of the inner list.
                  \LaTeX\ allows you to nest lists deeper than
                  you really should.
         \end{enumerate}
         This is the rest of the second item of the outer
         list.  It is no more interesting than any other
         part of the item.
   \item This is the third item of the list.
\end{itemize}
You can even display poetry.
\begin{verse}
   There is an environment
    for verse \\             % The \\ command separates lines
   Whose features some poets % within a stanza.
   will curse.

                             % One or more blank lines separate stanzas.

   For instead of making\\
   Them do \emph{all} line breaking, \\
   It allows them to put too many words on a line when they'd rather be
   forced to be terse.
\end{verse}

Mathematical formulas may also be displayed.  A
displayed formula
is
one-line long; multiline
formulas require special formatting instructions.
   \[  \ip{\Gamma}{\psi'} = x'' + y^{2} + z_{i}^{n}\]
Don't start a paragraph with a displayed equation,
nor make one a paragraph by itself.

\end{document}               % End of document.

\section{Bench Huge Section 1380}
% synthetic bench section 1380
% This is a sample LaTeX input file.  (Version of 12 August 2004.)
%
% A '%' character causes TeX to ignore all remaining text on the line,
% and is used for comments like this one.

\documentclass{article}      % Specifies the document class

                             % The preamble begins here.
\title{An Example Document}  % Declares the document's title.
\author{Leslie Lamport}      % Declares the author's name.
\date{January 21, 1994}      % Deleting this command produces today's date.

\newcommand{\ip}[2]{(#1, #2)}
                             % Defines \ip{arg1}{arg2} to mean
                             % (arg1, arg2).

%\newcommand{\ip}[2]{\langle #1 | #2\rangle}
                             % This is an alternative definition of
                             % \ip that is commented out.

\begin{document}             % End of preamble and beginning of text.

\maketitle                   % Produces the title.

This is an example input file.  Comparing it with
the output it generates can show you how to
produce a simple document of your own.

\section{Ordinary Text}      % Produces section heading.  Lower-level
                             % sections are begun with similar
                             % \subsection and \subsubsection commands.

The ends  of words and sentences are marked
  by   spaces. It  doesn't matter how many
spaces    you type; one is as good as 100.  The
end of   a line counts as a space.

One   or more   blank lines denote the  end
of  a paragraph.

Since any number of consecutive spaces are treated
like a single one, the formatting of the input
file makes no difference to
      \LaTeX,                % The \LaTeX command generates the LaTeX logo.
but it makes a difference to you.  When you use
\LaTeX, making your input file as easy to read
as possible will be a great help as you write
your document and when you change it.  This sample
file shows how you can add comments to your own input
file.

Because printing is different from typewriting,
there are a number of things that you have to do
differently when preparing an input file than if
you were just typing the document directly.
Quotation marks like
       ``this''
have to be handled specially, as do quotes within
quotes:
       ``\,`this'            % \, separates the double and single quote.
        is what I just
        wrote, not  `that'\,''.

Dashes come in three sizes: an
       intra-word
dash, a medium dash for number ranges like
       1--2,
and a punctuation
       dash---like
this.

A sentence-ending space should be larger than the
space between words within a sentence.  You
sometimes have to type special commands in
conjunction with punctuation characters to get
this right, as in the following sentence.
       Gnats, gnus, etc.\ all  % `\ ' makes an inter-word space.
       begin with G\@.         % \@ marks end-of-sentence punctuation.
You should check the spaces after periods when
reading your output to make sure you haven't
forgotten any special cases.  Generating an
ellipsis
       \ldots\               % `\ ' is needed after `\ldots' because TeX
                             % ignores spaces after command names like \ldots
                             % made from \ + letters.
                             %
                             % Note how a `%' character causes TeX to ignore
                             % the end of the input line, so these blank lines
                             % do not start a new paragraph.
                             %
with the right spacing around the periods requires
a special command.

\LaTeX\ interprets some common characters as
commands, so you must type special commands to
generate them.  These characters include the
following:
       \$ \& \% \# \{ and \}.

In printing, text is usually emphasized with an
       \emph{italic}
type style.

\begin{em}
   A long segment of text can also be emphasized
   in this way.  Text within such a segment can be
   given \emph{additional} emphasis.
\end{em}

It is sometimes necessary to prevent \LaTeX\ from
breaking a line where it might otherwise do so.
This may be at a space, as between the ``Mr.''\ and
``Jones'' in
       ``Mr.~Jones'',        % ~ produces an unbreakable interword space.
or within a word---especially when the word is a
symbol like
       \mbox{\emph{itemnum}}
that makes little sense when hyphenated across
lines.

Footnotes\footnote{This is an example of a footnote.}
pose no problem.

\LaTeX\ is good at typesetting mathematical formulas
like
       \( x-3y + z = 7 \)
or
       \( a_{1} > x^{2n} + y^{2n} > x' \)
or
       \( \ip{A}{B} = \sum_{i} a_{i} b_{i} \).
The spaces you type in a formula are
ignored.  Remember that a letter like
       $x$                   % $ ... $  and  \( ... \)  are equivalent
is a formula when it denotes a mathematical
symbol, and it should be typed as one.

\section{Displayed Text}

Text is displayed by indenting it from the left
margin.  Quotations are commonly displayed.  There
are short quotations
\begin{quote}
   This is a short quotation.  It consists of a
   single paragraph of text.  See how it is formatted.
\end{quote}
and longer ones.
\begin{quotation}
   This is a longer quotation.  It consists of two
   paragraphs of text, neither of which are
   particularly interesting.

   This is the second paragraph of the quotation.  It
   is just as dull as the first paragraph.
\end{quotation}
Another frequently-displayed structure is a list.
The following is an example of an \emph{itemized}
list.
\begin{itemize}
   \item This is the first item of an itemized list.
         Each item in the list is marked with a ``tick''.
         You don't have to worry about what kind of tick
         mark is used.

   \item This is the second item of the list.  It
         contains another list nested inside it.  The inner
         list is an \emph{enumerated} list.
         \begin{enumerate}
            \item This is the first item of an enumerated
                  list that is nested within the itemized list.

            \item This is the second item of the inner list.
                  \LaTeX\ allows you to nest lists deeper than
                  you really should.
         \end{enumerate}
         This is the rest of the second item of the outer
         list.  It is no more interesting than any other
         part of the item.
   \item This is the third item of the list.
\end{itemize}
You can even display poetry.
\begin{verse}
   There is an environment
    for verse \\             % The \\ command separates lines
   Whose features some poets % within a stanza.
   will curse.

                             % One or more blank lines separate stanzas.

   For instead of making\\
   Them do \emph{all} line breaking, \\
   It allows them to put too many words on a line when they'd rather be
   forced to be terse.
\end{verse}

Mathematical formulas may also be displayed.  A
displayed formula
is
one-line long; multiline
formulas require special formatting instructions.
   \[  \ip{\Gamma}{\psi'} = x'' + y^{2} + z_{i}^{n}\]
Don't start a paragraph with a displayed equation,
nor make one a paragraph by itself.

\end{document}               % End of document.

\section{Bench Huge Section 1381}
% synthetic bench section 1381
% This is a sample LaTeX input file.  (Version of 12 August 2004.)
%
% A '%' character causes TeX to ignore all remaining text on the line,
% and is used for comments like this one.

\documentclass{article}      % Specifies the document class

                             % The preamble begins here.
\title{An Example Document}  % Declares the document's title.
\author{Leslie Lamport}      % Declares the author's name.
\date{January 21, 1994}      % Deleting this command produces today's date.

\newcommand{\ip}[2]{(#1, #2)}
                             % Defines \ip{arg1}{arg2} to mean
                             % (arg1, arg2).

%\newcommand{\ip}[2]{\langle #1 | #2\rangle}
                             % This is an alternative definition of
                             % \ip that is commented out.

\begin{document}             % End of preamble and beginning of text.

\maketitle                   % Produces the title.

This is an example input file.  Comparing it with
the output it generates can show you how to
produce a simple document of your own.

\section{Ordinary Text}      % Produces section heading.  Lower-level
                             % sections are begun with similar
                             % \subsection and \subsubsection commands.

The ends  of words and sentences are marked
  by   spaces. It  doesn't matter how many
spaces    you type; one is as good as 100.  The
end of   a line counts as a space.

One   or more   blank lines denote the  end
of  a paragraph.

Since any number of consecutive spaces are treated
like a single one, the formatting of the input
file makes no difference to
      \LaTeX,                % The \LaTeX command generates the LaTeX logo.
but it makes a difference to you.  When you use
\LaTeX, making your input file as easy to read
as possible will be a great help as you write
your document and when you change it.  This sample
file shows how you can add comments to your own input
file.

Because printing is different from typewriting,
there are a number of things that you have to do
differently when preparing an input file than if
you were just typing the document directly.
Quotation marks like
       ``this''
have to be handled specially, as do quotes within
quotes:
       ``\,`this'            % \, separates the double and single quote.
        is what I just
        wrote, not  `that'\,''.

Dashes come in three sizes: an
       intra-word
dash, a medium dash for number ranges like
       1--2,
and a punctuation
       dash---like
this.

A sentence-ending space should be larger than the
space between words within a sentence.  You
sometimes have to type special commands in
conjunction with punctuation characters to get
this right, as in the following sentence.
       Gnats, gnus, etc.\ all  % `\ ' makes an inter-word space.
       begin with G\@.         % \@ marks end-of-sentence punctuation.
You should check the spaces after periods when
reading your output to make sure you haven't
forgotten any special cases.  Generating an
ellipsis
       \ldots\               % `\ ' is needed after `\ldots' because TeX
                             % ignores spaces after command names like \ldots
                             % made from \ + letters.
                             %
                             % Note how a `%' character causes TeX to ignore
                             % the end of the input line, so these blank lines
                             % do not start a new paragraph.
                             %
with the right spacing around the periods requires
a special command.

\LaTeX\ interprets some common characters as
commands, so you must type special commands to
generate them.  These characters include the
following:
       \$ \& \% \# \{ and \}.

In printing, text is usually emphasized with an
       \emph{italic}
type style.

\begin{em}
   A long segment of text can also be emphasized
   in this way.  Text within such a segment can be
   given \emph{additional} emphasis.
\end{em}

It is sometimes necessary to prevent \LaTeX\ from
breaking a line where it might otherwise do so.
This may be at a space, as between the ``Mr.''\ and
``Jones'' in
       ``Mr.~Jones'',        % ~ produces an unbreakable interword space.
or within a word---especially when the word is a
symbol like
       \mbox{\emph{itemnum}}
that makes little sense when hyphenated across
lines.

Footnotes\footnote{This is an example of a footnote.}
pose no problem.

\LaTeX\ is good at typesetting mathematical formulas
like
       \( x-3y + z = 7 \)
or
       \( a_{1} > x^{2n} + y^{2n} > x' \)
or
       \( \ip{A}{B} = \sum_{i} a_{i} b_{i} \).
The spaces you type in a formula are
ignored.  Remember that a letter like
       $x$                   % $ ... $  and  \( ... \)  are equivalent
is a formula when it denotes a mathematical
symbol, and it should be typed as one.

\section{Displayed Text}

Text is displayed by indenting it from the left
margin.  Quotations are commonly displayed.  There
are short quotations
\begin{quote}
   This is a short quotation.  It consists of a
   single paragraph of text.  See how it is formatted.
\end{quote}
and longer ones.
\begin{quotation}
   This is a longer quotation.  It consists of two
   paragraphs of text, neither of which are
   particularly interesting.

   This is the second paragraph of the quotation.  It
   is just as dull as the first paragraph.
\end{quotation}
Another frequently-displayed structure is a list.
The following is an example of an \emph{itemized}
list.
\begin{itemize}
   \item This is the first item of an itemized list.
         Each item in the list is marked with a ``tick''.
         You don't have to worry about what kind of tick
         mark is used.

   \item This is the second item of the list.  It
         contains another list nested inside it.  The inner
         list is an \emph{enumerated} list.
         \begin{enumerate}
            \item This is the first item of an enumerated
                  list that is nested within the itemized list.

            \item This is the second item of the inner list.
                  \LaTeX\ allows you to nest lists deeper than
                  you really should.
         \end{enumerate}
         This is the rest of the second item of the outer
         list.  It is no more interesting than any other
         part of the item.
   \item This is the third item of the list.
\end{itemize}
You can even display poetry.
\begin{verse}
   There is an environment
    for verse \\             % The \\ command separates lines
   Whose features some poets % within a stanza.
   will curse.

                             % One or more blank lines separate stanzas.

   For instead of making\\
   Them do \emph{all} line breaking, \\
   It allows them to put too many words on a line when they'd rather be
   forced to be terse.
\end{verse}

Mathematical formulas may also be displayed.  A
displayed formula
is
one-line long; multiline
formulas require special formatting instructions.
   \[  \ip{\Gamma}{\psi'} = x'' + y^{2} + z_{i}^{n}\]
Don't start a paragraph with a displayed equation,
nor make one a paragraph by itself.

\end{document}               % End of document.

\section{Bench Huge Section 1382}
% synthetic bench section 1382
% This is a sample LaTeX input file.  (Version of 12 August 2004.)
%
% A '%' character causes TeX to ignore all remaining text on the line,
% and is used for comments like this one.

\documentclass{article}      % Specifies the document class

                             % The preamble begins here.
\title{An Example Document}  % Declares the document's title.
\author{Leslie Lamport}      % Declares the author's name.
\date{January 21, 1994}      % Deleting this command produces today's date.

\newcommand{\ip}[2]{(#1, #2)}
                             % Defines \ip{arg1}{arg2} to mean
                             % (arg1, arg2).

%\newcommand{\ip}[2]{\langle #1 | #2\rangle}
                             % This is an alternative definition of
                             % \ip that is commented out.

\begin{document}             % End of preamble and beginning of text.

\maketitle                   % Produces the title.

This is an example input file.  Comparing it with
the output it generates can show you how to
produce a simple document of your own.

\section{Ordinary Text}      % Produces section heading.  Lower-level
                             % sections are begun with similar
                             % \subsection and \subsubsection commands.

The ends  of words and sentences are marked
  by   spaces. It  doesn't matter how many
spaces    you type; one is as good as 100.  The
end of   a line counts as a space.

One   or more   blank lines denote the  end
of  a paragraph.

Since any number of consecutive spaces are treated
like a single one, the formatting of the input
file makes no difference to
      \LaTeX,                % The \LaTeX command generates the LaTeX logo.
but it makes a difference to you.  When you use
\LaTeX, making your input file as easy to read
as possible will be a great help as you write
your document and when you change it.  This sample
file shows how you can add comments to your own input
file.

Because printing is different from typewriting,
there are a number of things that you have to do
differently when preparing an input file than if
you were just typing the document directly.
Quotation marks like
       ``this''
have to be handled specially, as do quotes within
quotes:
       ``\,`this'            % \, separates the double and single quote.
        is what I just
        wrote, not  `that'\,''.

Dashes come in three sizes: an
       intra-word
dash, a medium dash for number ranges like
       1--2,
and a punctuation
       dash---like
this.

A sentence-ending space should be larger than the
space between words within a sentence.  You
sometimes have to type special commands in
conjunction with punctuation characters to get
this right, as in the following sentence.
       Gnats, gnus, etc.\ all  % `\ ' makes an inter-word space.
       begin with G\@.         % \@ marks end-of-sentence punctuation.
You should check the spaces after periods when
reading your output to make sure you haven't
forgotten any special cases.  Generating an
ellipsis
       \ldots\               % `\ ' is needed after `\ldots' because TeX
                             % ignores spaces after command names like \ldots
                             % made from \ + letters.
                             %
                             % Note how a `%' character causes TeX to ignore
                             % the end of the input line, so these blank lines
                             % do not start a new paragraph.
                             %
with the right spacing around the periods requires
a special command.

\LaTeX\ interprets some common characters as
commands, so you must type special commands to
generate them.  These characters include the
following:
       \$ \& \% \# \{ and \}.

In printing, text is usually emphasized with an
       \emph{italic}
type style.

\begin{em}
   A long segment of text can also be emphasized
   in this way.  Text within such a segment can be
   given \emph{additional} emphasis.
\end{em}

It is sometimes necessary to prevent \LaTeX\ from
breaking a line where it might otherwise do so.
This may be at a space, as between the ``Mr.''\ and
``Jones'' in
       ``Mr.~Jones'',        % ~ produces an unbreakable interword space.
or within a word---especially when the word is a
symbol like
       \mbox{\emph{itemnum}}
that makes little sense when hyphenated across
lines.

Footnotes\footnote{This is an example of a footnote.}
pose no problem.

\LaTeX\ is good at typesetting mathematical formulas
like
       \( x-3y + z = 7 \)
or
       \( a_{1} > x^{2n} + y^{2n} > x' \)
or
       \( \ip{A}{B} = \sum_{i} a_{i} b_{i} \).
The spaces you type in a formula are
ignored.  Remember that a letter like
       $x$                   % $ ... $  and  \( ... \)  are equivalent
is a formula when it denotes a mathematical
symbol, and it should be typed as one.

\section{Displayed Text}

Text is displayed by indenting it from the left
margin.  Quotations are commonly displayed.  There
are short quotations
\begin{quote}
   This is a short quotation.  It consists of a
   single paragraph of text.  See how it is formatted.
\end{quote}
and longer ones.
\begin{quotation}
   This is a longer quotation.  It consists of two
   paragraphs of text, neither of which are
   particularly interesting.

   This is the second paragraph of the quotation.  It
   is just as dull as the first paragraph.
\end{quotation}
Another frequently-displayed structure is a list.
The following is an example of an \emph{itemized}
list.
\begin{itemize}
   \item This is the first item of an itemized list.
         Each item in the list is marked with a ``tick''.
         You don't have to worry about what kind of tick
         mark is used.

   \item This is the second item of the list.  It
         contains another list nested inside it.  The inner
         list is an \emph{enumerated} list.
         \begin{enumerate}
            \item This is the first item of an enumerated
                  list that is nested within the itemized list.

            \item This is the second item of the inner list.
                  \LaTeX\ allows you to nest lists deeper than
                  you really should.
         \end{enumerate}
         This is the rest of the second item of the outer
         list.  It is no more interesting than any other
         part of the item.
   \item This is the third item of the list.
\end{itemize}
You can even display poetry.
\begin{verse}
   There is an environment
    for verse \\             % The \\ command separates lines
   Whose features some poets % within a stanza.
   will curse.

                             % One or more blank lines separate stanzas.

   For instead of making\\
   Them do \emph{all} line breaking, \\
   It allows them to put too many words on a line when they'd rather be
   forced to be terse.
\end{verse}

Mathematical formulas may also be displayed.  A
displayed formula
is
one-line long; multiline
formulas require special formatting instructions.
   \[  \ip{\Gamma}{\psi'} = x'' + y^{2} + z_{i}^{n}\]
Don't start a paragraph with a displayed equation,
nor make one a paragraph by itself.

\end{document}               % End of document.

\section{Bench Huge Section 1383}
% synthetic bench section 1383
% This is a sample LaTeX input file.  (Version of 12 August 2004.)
%
% A '%' character causes TeX to ignore all remaining text on the line,
% and is used for comments like this one.

\documentclass{article}      % Specifies the document class

                             % The preamble begins here.
\title{An Example Document}  % Declares the document's title.
\author{Leslie Lamport}      % Declares the author's name.
\date{January 21, 1994}      % Deleting this command produces today's date.

\newcommand{\ip}[2]{(#1, #2)}
                             % Defines \ip{arg1}{arg2} to mean
                             % (arg1, arg2).

%\newcommand{\ip}[2]{\langle #1 | #2\rangle}
                             % This is an alternative definition of
                             % \ip that is commented out.

\begin{document}             % End of preamble and beginning of text.

\maketitle                   % Produces the title.

This is an example input file.  Comparing it with
the output it generates can show you how to
produce a simple document of your own.

\section{Ordinary Text}      % Produces section heading.  Lower-level
                             % sections are begun with similar
                             % \subsection and \subsubsection commands.

The ends  of words and sentences are marked
  by   spaces. It  doesn't matter how many
spaces    you type; one is as good as 100.  The
end of   a line counts as a space.

One   or more   blank lines denote the  end
of  a paragraph.

Since any number of consecutive spaces are treated
like a single one, the formatting of the input
file makes no difference to
      \LaTeX,                % The \LaTeX command generates the LaTeX logo.
but it makes a difference to you.  When you use
\LaTeX, making your input file as easy to read
as possible will be a great help as you write
your document and when you change it.  This sample
file shows how you can add comments to your own input
file.

Because printing is different from typewriting,
there are a number of things that you have to do
differently when preparing an input file than if
you were just typing the document directly.
Quotation marks like
       ``this''
have to be handled specially, as do quotes within
quotes:
       ``\,`this'            % \, separates the double and single quote.
        is what I just
        wrote, not  `that'\,''.

Dashes come in three sizes: an
       intra-word
dash, a medium dash for number ranges like
       1--2,
and a punctuation
       dash---like
this.

A sentence-ending space should be larger than the
space between words within a sentence.  You
sometimes have to type special commands in
conjunction with punctuation characters to get
this right, as in the following sentence.
       Gnats, gnus, etc.\ all  % `\ ' makes an inter-word space.
       begin with G\@.         % \@ marks end-of-sentence punctuation.
You should check the spaces after periods when
reading your output to make sure you haven't
forgotten any special cases.  Generating an
ellipsis
       \ldots\               % `\ ' is needed after `\ldots' because TeX
                             % ignores spaces after command names like \ldots
                             % made from \ + letters.
                             %
                             % Note how a `%' character causes TeX to ignore
                             % the end of the input line, so these blank lines
                             % do not start a new paragraph.
                             %
with the right spacing around the periods requires
a special command.

\LaTeX\ interprets some common characters as
commands, so you must type special commands to
generate them.  These characters include the
following:
       \$ \& \% \# \{ and \}.

In printing, text is usually emphasized with an
       \emph{italic}
type style.

\begin{em}
   A long segment of text can also be emphasized
   in this way.  Text within such a segment can be
   given \emph{additional} emphasis.
\end{em}

It is sometimes necessary to prevent \LaTeX\ from
breaking a line where it might otherwise do so.
This may be at a space, as between the ``Mr.''\ and
``Jones'' in
       ``Mr.~Jones'',        % ~ produces an unbreakable interword space.
or within a word---especially when the word is a
symbol like
       \mbox{\emph{itemnum}}
that makes little sense when hyphenated across
lines.

Footnotes\footnote{This is an example of a footnote.}
pose no problem.

\LaTeX\ is good at typesetting mathematical formulas
like
       \( x-3y + z = 7 \)
or
       \( a_{1} > x^{2n} + y^{2n} > x' \)
or
       \( \ip{A}{B} = \sum_{i} a_{i} b_{i} \).
The spaces you type in a formula are
ignored.  Remember that a letter like
       $x$                   % $ ... $  and  \( ... \)  are equivalent
is a formula when it denotes a mathematical
symbol, and it should be typed as one.

\section{Displayed Text}

Text is displayed by indenting it from the left
margin.  Quotations are commonly displayed.  There
are short quotations
\begin{quote}
   This is a short quotation.  It consists of a
   single paragraph of text.  See how it is formatted.
\end{quote}
and longer ones.
\begin{quotation}
   This is a longer quotation.  It consists of two
   paragraphs of text, neither of which are
   particularly interesting.

   This is the second paragraph of the quotation.  It
   is just as dull as the first paragraph.
\end{quotation}
Another frequently-displayed structure is a list.
The following is an example of an \emph{itemized}
list.
\begin{itemize}
   \item This is the first item of an itemized list.
         Each item in the list is marked with a ``tick''.
         You don't have to worry about what kind of tick
         mark is used.

   \item This is the second item of the list.  It
         contains another list nested inside it.  The inner
         list is an \emph{enumerated} list.
         \begin{enumerate}
            \item This is the first item of an enumerated
                  list that is nested within the itemized list.

            \item This is the second item of the inner list.
                  \LaTeX\ allows you to nest lists deeper than
                  you really should.
         \end{enumerate}
         This is the rest of the second item of the outer
         list.  It is no more interesting than any other
         part of the item.
   \item This is the third item of the list.
\end{itemize}
You can even display poetry.
\begin{verse}
   There is an environment
    for verse \\             % The \\ command separates lines
   Whose features some poets % within a stanza.
   will curse.

                             % One or more blank lines separate stanzas.

   For instead of making\\
   Them do \emph{all} line breaking, \\
   It allows them to put too many words on a line when they'd rather be
   forced to be terse.
\end{verse}

Mathematical formulas may also be displayed.  A
displayed formula
is
one-line long; multiline
formulas require special formatting instructions.
   \[  \ip{\Gamma}{\psi'} = x'' + y^{2} + z_{i}^{n}\]
Don't start a paragraph with a displayed equation,
nor make one a paragraph by itself.

\end{document}               % End of document.

\section{Bench Huge Section 1384}
% synthetic bench section 1384
% This is a sample LaTeX input file.  (Version of 12 August 2004.)
%
% A '%' character causes TeX to ignore all remaining text on the line,
% and is used for comments like this one.

\documentclass{article}      % Specifies the document class

                             % The preamble begins here.
\title{An Example Document}  % Declares the document's title.
\author{Leslie Lamport}      % Declares the author's name.
\date{January 21, 1994}      % Deleting this command produces today's date.

\newcommand{\ip}[2]{(#1, #2)}
                             % Defines \ip{arg1}{arg2} to mean
                             % (arg1, arg2).

%\newcommand{\ip}[2]{\langle #1 | #2\rangle}
                             % This is an alternative definition of
                             % \ip that is commented out.

\begin{document}             % End of preamble and beginning of text.

\maketitle                   % Produces the title.

This is an example input file.  Comparing it with
the output it generates can show you how to
produce a simple document of your own.

\section{Ordinary Text}      % Produces section heading.  Lower-level
                             % sections are begun with similar
                             % \subsection and \subsubsection commands.

The ends  of words and sentences are marked
  by   spaces. It  doesn't matter how many
spaces    you type; one is as good as 100.  The
end of   a line counts as a space.

One   or more   blank lines denote the  end
of  a paragraph.

Since any number of consecutive spaces are treated
like a single one, the formatting of the input
file makes no difference to
      \LaTeX,                % The \LaTeX command generates the LaTeX logo.
but it makes a difference to you.  When you use
\LaTeX, making your input file as easy to read
as possible will be a great help as you write
your document and when you change it.  This sample
file shows how you can add comments to your own input
file.

Because printing is different from typewriting,
there are a number of things that you have to do
differently when preparing an input file than if
you were just typing the document directly.
Quotation marks like
       ``this''
have to be handled specially, as do quotes within
quotes:
       ``\,`this'            % \, separates the double and single quote.
        is what I just
        wrote, not  `that'\,''.

Dashes come in three sizes: an
       intra-word
dash, a medium dash for number ranges like
       1--2,
and a punctuation
       dash---like
this.

A sentence-ending space should be larger than the
space between words within a sentence.  You
sometimes have to type special commands in
conjunction with punctuation characters to get
this right, as in the following sentence.
       Gnats, gnus, etc.\ all  % `\ ' makes an inter-word space.
       begin with G\@.         % \@ marks end-of-sentence punctuation.
You should check the spaces after periods when
reading your output to make sure you haven't
forgotten any special cases.  Generating an
ellipsis
       \ldots\               % `\ ' is needed after `\ldots' because TeX
                             % ignores spaces after command names like \ldots
                             % made from \ + letters.
                             %
                             % Note how a `%' character causes TeX to ignore
                             % the end of the input line, so these blank lines
                             % do not start a new paragraph.
                             %
with the right spacing around the periods requires
a special command.

\LaTeX\ interprets some common characters as
commands, so you must type special commands to
generate them.  These characters include the
following:
       \$ \& \% \# \{ and \}.

In printing, text is usually emphasized with an
       \emph{italic}
type style.

\begin{em}
   A long segment of text can also be emphasized
   in this way.  Text within such a segment can be
   given \emph{additional} emphasis.
\end{em}

It is sometimes necessary to prevent \LaTeX\ from
breaking a line where it might otherwise do so.
This may be at a space, as between the ``Mr.''\ and
``Jones'' in
       ``Mr.~Jones'',        % ~ produces an unbreakable interword space.
or within a word---especially when the word is a
symbol like
       \mbox{\emph{itemnum}}
that makes little sense when hyphenated across
lines.

Footnotes\footnote{This is an example of a footnote.}
pose no problem.

\LaTeX\ is good at typesetting mathematical formulas
like
       \( x-3y + z = 7 \)
or
       \( a_{1} > x^{2n} + y^{2n} > x' \)
or
       \( \ip{A}{B} = \sum_{i} a_{i} b_{i} \).
The spaces you type in a formula are
ignored.  Remember that a letter like
       $x$                   % $ ... $  and  \( ... \)  are equivalent
is a formula when it denotes a mathematical
symbol, and it should be typed as one.

\section{Displayed Text}

Text is displayed by indenting it from the left
margin.  Quotations are commonly displayed.  There
are short quotations
\begin{quote}
   This is a short quotation.  It consists of a
   single paragraph of text.  See how it is formatted.
\end{quote}
and longer ones.
\begin{quotation}
   This is a longer quotation.  It consists of two
   paragraphs of text, neither of which are
   particularly interesting.

   This is the second paragraph of the quotation.  It
   is just as dull as the first paragraph.
\end{quotation}
Another frequently-displayed structure is a list.
The following is an example of an \emph{itemized}
list.
\begin{itemize}
   \item This is the first item of an itemized list.
         Each item in the list is marked with a ``tick''.
         You don't have to worry about what kind of tick
         mark is used.

   \item This is the second item of the list.  It
         contains another list nested inside it.  The inner
         list is an \emph{enumerated} list.
         \begin{enumerate}
            \item This is the first item of an enumerated
                  list that is nested within the itemized list.

            \item This is the second item of the inner list.
                  \LaTeX\ allows you to nest lists deeper than
                  you really should.
         \end{enumerate}
         This is the rest of the second item of the outer
         list.  It is no more interesting than any other
         part of the item.
   \item This is the third item of the list.
\end{itemize}
You can even display poetry.
\begin{verse}
   There is an environment
    for verse \\             % The \\ command separates lines
   Whose features some poets % within a stanza.
   will curse.

                             % One or more blank lines separate stanzas.

   For instead of making\\
   Them do \emph{all} line breaking, \\
   It allows them to put too many words on a line when they'd rather be
   forced to be terse.
\end{verse}

Mathematical formulas may also be displayed.  A
displayed formula
is
one-line long; multiline
formulas require special formatting instructions.
   \[  \ip{\Gamma}{\psi'} = x'' + y^{2} + z_{i}^{n}\]
Don't start a paragraph with a displayed equation,
nor make one a paragraph by itself.

\end{document}               % End of document.

\section{Bench Huge Section 1385}
% synthetic bench section 1385
% This is a sample LaTeX input file.  (Version of 12 August 2004.)
%
% A '%' character causes TeX to ignore all remaining text on the line,
% and is used for comments like this one.

\documentclass{article}      % Specifies the document class

                             % The preamble begins here.
\title{An Example Document}  % Declares the document's title.
\author{Leslie Lamport}      % Declares the author's name.
\date{January 21, 1994}      % Deleting this command produces today's date.

\newcommand{\ip}[2]{(#1, #2)}
                             % Defines \ip{arg1}{arg2} to mean
                             % (arg1, arg2).

%\newcommand{\ip}[2]{\langle #1 | #2\rangle}
                             % This is an alternative definition of
                             % \ip that is commented out.

\begin{document}             % End of preamble and beginning of text.

\maketitle                   % Produces the title.

This is an example input file.  Comparing it with
the output it generates can show you how to
produce a simple document of your own.

\section{Ordinary Text}      % Produces section heading.  Lower-level
                             % sections are begun with similar
                             % \subsection and \subsubsection commands.

The ends  of words and sentences are marked
  by   spaces. It  doesn't matter how many
spaces    you type; one is as good as 100.  The
end of   a line counts as a space.

One   or more   blank lines denote the  end
of  a paragraph.

Since any number of consecutive spaces are treated
like a single one, the formatting of the input
file makes no difference to
      \LaTeX,                % The \LaTeX command generates the LaTeX logo.
but it makes a difference to you.  When you use
\LaTeX, making your input file as easy to read
as possible will be a great help as you write
your document and when you change it.  This sample
file shows how you can add comments to your own input
file.

Because printing is different from typewriting,
there are a number of things that you have to do
differently when preparing an input file than if
you were just typing the document directly.
Quotation marks like
       ``this''
have to be handled specially, as do quotes within
quotes:
       ``\,`this'            % \, separates the double and single quote.
        is what I just
        wrote, not  `that'\,''.

Dashes come in three sizes: an
       intra-word
dash, a medium dash for number ranges like
       1--2,
and a punctuation
       dash---like
this.

A sentence-ending space should be larger than the
space between words within a sentence.  You
sometimes have to type special commands in
conjunction with punctuation characters to get
this right, as in the following sentence.
       Gnats, gnus, etc.\ all  % `\ ' makes an inter-word space.
       begin with G\@.         % \@ marks end-of-sentence punctuation.
You should check the spaces after periods when
reading your output to make sure you haven't
forgotten any special cases.  Generating an
ellipsis
       \ldots\               % `\ ' is needed after `\ldots' because TeX
                             % ignores spaces after command names like \ldots
                             % made from \ + letters.
                             %
                             % Note how a `%' character causes TeX to ignore
                             % the end of the input line, so these blank lines
                             % do not start a new paragraph.
                             %
with the right spacing around the periods requires
a special command.

\LaTeX\ interprets some common characters as
commands, so you must type special commands to
generate them.  These characters include the
following:
       \$ \& \% \# \{ and \}.

In printing, text is usually emphasized with an
       \emph{italic}
type style.

\begin{em}
   A long segment of text can also be emphasized
   in this way.  Text within such a segment can be
   given \emph{additional} emphasis.
\end{em}

It is sometimes necessary to prevent \LaTeX\ from
breaking a line where it might otherwise do so.
This may be at a space, as between the ``Mr.''\ and
``Jones'' in
       ``Mr.~Jones'',        % ~ produces an unbreakable interword space.
or within a word---especially when the word is a
symbol like
       \mbox{\emph{itemnum}}
that makes little sense when hyphenated across
lines.

Footnotes\footnote{This is an example of a footnote.}
pose no problem.

\LaTeX\ is good at typesetting mathematical formulas
like
       \( x-3y + z = 7 \)
or
       \( a_{1} > x^{2n} + y^{2n} > x' \)
or
       \( \ip{A}{B} = \sum_{i} a_{i} b_{i} \).
The spaces you type in a formula are
ignored.  Remember that a letter like
       $x$                   % $ ... $  and  \( ... \)  are equivalent
is a formula when it denotes a mathematical
symbol, and it should be typed as one.

\section{Displayed Text}

Text is displayed by indenting it from the left
margin.  Quotations are commonly displayed.  There
are short quotations
\begin{quote}
   This is a short quotation.  It consists of a
   single paragraph of text.  See how it is formatted.
\end{quote}
and longer ones.
\begin{quotation}
   This is a longer quotation.  It consists of two
   paragraphs of text, neither of which are
   particularly interesting.

   This is the second paragraph of the quotation.  It
   is just as dull as the first paragraph.
\end{quotation}
Another frequently-displayed structure is a list.
The following is an example of an \emph{itemized}
list.
\begin{itemize}
   \item This is the first item of an itemized list.
         Each item in the list is marked with a ``tick''.
         You don't have to worry about what kind of tick
         mark is used.

   \item This is the second item of the list.  It
         contains another list nested inside it.  The inner
         list is an \emph{enumerated} list.
         \begin{enumerate}
            \item This is the first item of an enumerated
                  list that is nested within the itemized list.

            \item This is the second item of the inner list.
                  \LaTeX\ allows you to nest lists deeper than
                  you really should.
         \end{enumerate}
         This is the rest of the second item of the outer
         list.  It is no more interesting than any other
         part of the item.
   \item This is the third item of the list.
\end{itemize}
You can even display poetry.
\begin{verse}
   There is an environment
    for verse \\             % The \\ command separates lines
   Whose features some poets % within a stanza.
   will curse.

                             % One or more blank lines separate stanzas.

   For instead of making\\
   Them do \emph{all} line breaking, \\
   It allows them to put too many words on a line when they'd rather be
   forced to be terse.
\end{verse}

Mathematical formulas may also be displayed.  A
displayed formula
is
one-line long; multiline
formulas require special formatting instructions.
   \[  \ip{\Gamma}{\psi'} = x'' + y^{2} + z_{i}^{n}\]
Don't start a paragraph with a displayed equation,
nor make one a paragraph by itself.

\end{document}               % End of document.

\section{Bench Huge Section 1386}
% synthetic bench section 1386
% This is a sample LaTeX input file.  (Version of 12 August 2004.)
%
% A '%' character causes TeX to ignore all remaining text on the line,
% and is used for comments like this one.

\documentclass{article}      % Specifies the document class

                             % The preamble begins here.
\title{An Example Document}  % Declares the document's title.
\author{Leslie Lamport}      % Declares the author's name.
\date{January 21, 1994}      % Deleting this command produces today's date.

\newcommand{\ip}[2]{(#1, #2)}
                             % Defines \ip{arg1}{arg2} to mean
                             % (arg1, arg2).

%\newcommand{\ip}[2]{\langle #1 | #2\rangle}
                             % This is an alternative definition of
                             % \ip that is commented out.

\begin{document}             % End of preamble and beginning of text.

\maketitle                   % Produces the title.

This is an example input file.  Comparing it with
the output it generates can show you how to
produce a simple document of your own.

\section{Ordinary Text}      % Produces section heading.  Lower-level
                             % sections are begun with similar
                             % \subsection and \subsubsection commands.

The ends  of words and sentences are marked
  by   spaces. It  doesn't matter how many
spaces    you type; one is as good as 100.  The
end of   a line counts as a space.

One   or more   blank lines denote the  end
of  a paragraph.

Since any number of consecutive spaces are treated
like a single one, the formatting of the input
file makes no difference to
      \LaTeX,                % The \LaTeX command generates the LaTeX logo.
but it makes a difference to you.  When you use
\LaTeX, making your input file as easy to read
as possible will be a great help as you write
your document and when you change it.  This sample
file shows how you can add comments to your own input
file.

Because printing is different from typewriting,
there are a number of things that you have to do
differently when preparing an input file than if
you were just typing the document directly.
Quotation marks like
       ``this''
have to be handled specially, as do quotes within
quotes:
       ``\,`this'            % \, separates the double and single quote.
        is what I just
        wrote, not  `that'\,''.

Dashes come in three sizes: an
       intra-word
dash, a medium dash for number ranges like
       1--2,
and a punctuation
       dash---like
this.

A sentence-ending space should be larger than the
space between words within a sentence.  You
sometimes have to type special commands in
conjunction with punctuation characters to get
this right, as in the following sentence.
       Gnats, gnus, etc.\ all  % `\ ' makes an inter-word space.
       begin with G\@.         % \@ marks end-of-sentence punctuation.
You should check the spaces after periods when
reading your output to make sure you haven't
forgotten any special cases.  Generating an
ellipsis
       \ldots\               % `\ ' is needed after `\ldots' because TeX
                             % ignores spaces after command names like \ldots
                             % made from \ + letters.
                             %
                             % Note how a `%' character causes TeX to ignore
                             % the end of the input line, so these blank lines
                             % do not start a new paragraph.
                             %
with the right spacing around the periods requires
a special command.

\LaTeX\ interprets some common characters as
commands, so you must type special commands to
generate them.  These characters include the
following:
       \$ \& \% \# \{ and \}.

In printing, text is usually emphasized with an
       \emph{italic}
type style.

\begin{em}
   A long segment of text can also be emphasized
   in this way.  Text within such a segment can be
   given \emph{additional} emphasis.
\end{em}

It is sometimes necessary to prevent \LaTeX\ from
breaking a line where it might otherwise do so.
This may be at a space, as between the ``Mr.''\ and
``Jones'' in
       ``Mr.~Jones'',        % ~ produces an unbreakable interword space.
or within a word---especially when the word is a
symbol like
       \mbox{\emph{itemnum}}
that makes little sense when hyphenated across
lines.

Footnotes\footnote{This is an example of a footnote.}
pose no problem.

\LaTeX\ is good at typesetting mathematical formulas
like
       \( x-3y + z = 7 \)
or
       \( a_{1} > x^{2n} + y^{2n} > x' \)
or
       \( \ip{A}{B} = \sum_{i} a_{i} b_{i} \).
The spaces you type in a formula are
ignored.  Remember that a letter like
       $x$                   % $ ... $  and  \( ... \)  are equivalent
is a formula when it denotes a mathematical
symbol, and it should be typed as one.

\section{Displayed Text}

Text is displayed by indenting it from the left
margin.  Quotations are commonly displayed.  There
are short quotations
\begin{quote}
   This is a short quotation.  It consists of a
   single paragraph of text.  See how it is formatted.
\end{quote}
and longer ones.
\begin{quotation}
   This is a longer quotation.  It consists of two
   paragraphs of text, neither of which are
   particularly interesting.

   This is the second paragraph of the quotation.  It
   is just as dull as the first paragraph.
\end{quotation}
Another frequently-displayed structure is a list.
The following is an example of an \emph{itemized}
list.
\begin{itemize}
   \item This is the first item of an itemized list.
         Each item in the list is marked with a ``tick''.
         You don't have to worry about what kind of tick
         mark is used.

   \item This is the second item of the list.  It
         contains another list nested inside it.  The inner
         list is an \emph{enumerated} list.
         \begin{enumerate}
            \item This is the first item of an enumerated
                  list that is nested within the itemized list.

            \item This is the second item of the inner list.
                  \LaTeX\ allows you to nest lists deeper than
                  you really should.
         \end{enumerate}
         This is the rest of the second item of the outer
         list.  It is no more interesting than any other
         part of the item.
   \item This is the third item of the list.
\end{itemize}
You can even display poetry.
\begin{verse}
   There is an environment
    for verse \\             % The \\ command separates lines
   Whose features some poets % within a stanza.
   will curse.

                             % One or more blank lines separate stanzas.

   For instead of making\\
   Them do \emph{all} line breaking, \\
   It allows them to put too many words on a line when they'd rather be
   forced to be terse.
\end{verse}

Mathematical formulas may also be displayed.  A
displayed formula
is
one-line long; multiline
formulas require special formatting instructions.
   \[  \ip{\Gamma}{\psi'} = x'' + y^{2} + z_{i}^{n}\]
Don't start a paragraph with a displayed equation,
nor make one a paragraph by itself.

\end{document}               % End of document.

\section{Bench Huge Section 1387}
% synthetic bench section 1387
% This is a sample LaTeX input file.  (Version of 12 August 2004.)
%
% A '%' character causes TeX to ignore all remaining text on the line,
% and is used for comments like this one.

\documentclass{article}      % Specifies the document class

                             % The preamble begins here.
\title{An Example Document}  % Declares the document's title.
\author{Leslie Lamport}      % Declares the author's name.
\date{January 21, 1994}      % Deleting this command produces today's date.

\newcommand{\ip}[2]{(#1, #2)}
                             % Defines \ip{arg1}{arg2} to mean
                             % (arg1, arg2).

%\newcommand{\ip}[2]{\langle #1 | #2\rangle}
                             % This is an alternative definition of
                             % \ip that is commented out.

\begin{document}             % End of preamble and beginning of text.

\maketitle                   % Produces the title.

This is an example input file.  Comparing it with
the output it generates can show you how to
produce a simple document of your own.

\section{Ordinary Text}      % Produces section heading.  Lower-level
                             % sections are begun with similar
                             % \subsection and \subsubsection commands.

The ends  of words and sentences are marked
  by   spaces. It  doesn't matter how many
spaces    you type; one is as good as 100.  The
end of   a line counts as a space.

One   or more   blank lines denote the  end
of  a paragraph.

Since any number of consecutive spaces are treated
like a single one, the formatting of the input
file makes no difference to
      \LaTeX,                % The \LaTeX command generates the LaTeX logo.
but it makes a difference to you.  When you use
\LaTeX, making your input file as easy to read
as possible will be a great help as you write
your document and when you change it.  This sample
file shows how you can add comments to your own input
file.

Because printing is different from typewriting,
there are a number of things that you have to do
differently when preparing an input file than if
you were just typing the document directly.
Quotation marks like
       ``this''
have to be handled specially, as do quotes within
quotes:
       ``\,`this'            % \, separates the double and single quote.
        is what I just
        wrote, not  `that'\,''.

Dashes come in three sizes: an
       intra-word
dash, a medium dash for number ranges like
       1--2,
and a punctuation
       dash---like
this.

A sentence-ending space should be larger than the
space between words within a sentence.  You
sometimes have to type special commands in
conjunction with punctuation characters to get
this right, as in the following sentence.
       Gnats, gnus, etc.\ all  % `\ ' makes an inter-word space.
       begin with G\@.         % \@ marks end-of-sentence punctuation.
You should check the spaces after periods when
reading your output to make sure you haven't
forgotten any special cases.  Generating an
ellipsis
       \ldots\               % `\ ' is needed after `\ldots' because TeX
                             % ignores spaces after command names like \ldots
                             % made from \ + letters.
                             %
                             % Note how a `%' character causes TeX to ignore
                             % the end of the input line, so these blank lines
                             % do not start a new paragraph.
                             %
with the right spacing around the periods requires
a special command.

\LaTeX\ interprets some common characters as
commands, so you must type special commands to
generate them.  These characters include the
following:
       \$ \& \% \# \{ and \}.

In printing, text is usually emphasized with an
       \emph{italic}
type style.

\begin{em}
   A long segment of text can also be emphasized
   in this way.  Text within such a segment can be
   given \emph{additional} emphasis.
\end{em}

It is sometimes necessary to prevent \LaTeX\ from
breaking a line where it might otherwise do so.
This may be at a space, as between the ``Mr.''\ and
``Jones'' in
       ``Mr.~Jones'',        % ~ produces an unbreakable interword space.
or within a word---especially when the word is a
symbol like
       \mbox{\emph{itemnum}}
that makes little sense when hyphenated across
lines.

Footnotes\footnote{This is an example of a footnote.}
pose no problem.

\LaTeX\ is good at typesetting mathematical formulas
like
       \( x-3y + z = 7 \)
or
       \( a_{1} > x^{2n} + y^{2n} > x' \)
or
       \( \ip{A}{B} = \sum_{i} a_{i} b_{i} \).
The spaces you type in a formula are
ignored.  Remember that a letter like
       $x$                   % $ ... $  and  \( ... \)  are equivalent
is a formula when it denotes a mathematical
symbol, and it should be typed as one.

\section{Displayed Text}

Text is displayed by indenting it from the left
margin.  Quotations are commonly displayed.  There
are short quotations
\begin{quote}
   This is a short quotation.  It consists of a
   single paragraph of text.  See how it is formatted.
\end{quote}
and longer ones.
\begin{quotation}
   This is a longer quotation.  It consists of two
   paragraphs of text, neither of which are
   particularly interesting.

   This is the second paragraph of the quotation.  It
   is just as dull as the first paragraph.
\end{quotation}
Another frequently-displayed structure is a list.
The following is an example of an \emph{itemized}
list.
\begin{itemize}
   \item This is the first item of an itemized list.
         Each item in the list is marked with a ``tick''.
         You don't have to worry about what kind of tick
         mark is used.

   \item This is the second item of the list.  It
         contains another list nested inside it.  The inner
         list is an \emph{enumerated} list.
         \begin{enumerate}
            \item This is the first item of an enumerated
                  list that is nested within the itemized list.

            \item This is the second item of the inner list.
                  \LaTeX\ allows you to nest lists deeper than
                  you really should.
         \end{enumerate}
         This is the rest of the second item of the outer
         list.  It is no more interesting than any other
         part of the item.
   \item This is the third item of the list.
\end{itemize}
You can even display poetry.
\begin{verse}
   There is an environment
    for verse \\             % The \\ command separates lines
   Whose features some poets % within a stanza.
   will curse.

                             % One or more blank lines separate stanzas.

   For instead of making\\
   Them do \emph{all} line breaking, \\
   It allows them to put too many words on a line when they'd rather be
   forced to be terse.
\end{verse}

Mathematical formulas may also be displayed.  A
displayed formula
is
one-line long; multiline
formulas require special formatting instructions.
   \[  \ip{\Gamma}{\psi'} = x'' + y^{2} + z_{i}^{n}\]
Don't start a paragraph with a displayed equation,
nor make one a paragraph by itself.

\end{document}               % End of document.

\section{Bench Huge Section 1388}
% synthetic bench section 1388
% This is a sample LaTeX input file.  (Version of 12 August 2004.)
%
% A '%' character causes TeX to ignore all remaining text on the line,
% and is used for comments like this one.

\documentclass{article}      % Specifies the document class

                             % The preamble begins here.
\title{An Example Document}  % Declares the document's title.
\author{Leslie Lamport}      % Declares the author's name.
\date{January 21, 1994}      % Deleting this command produces today's date.

\newcommand{\ip}[2]{(#1, #2)}
                             % Defines \ip{arg1}{arg2} to mean
                             % (arg1, arg2).

%\newcommand{\ip}[2]{\langle #1 | #2\rangle}
                             % This is an alternative definition of
                             % \ip that is commented out.

\begin{document}             % End of preamble and beginning of text.

\maketitle                   % Produces the title.

This is an example input file.  Comparing it with
the output it generates can show you how to
produce a simple document of your own.

\section{Ordinary Text}      % Produces section heading.  Lower-level
                             % sections are begun with similar
                             % \subsection and \subsubsection commands.

The ends  of words and sentences are marked
  by   spaces. It  doesn't matter how many
spaces    you type; one is as good as 100.  The
end of   a line counts as a space.

One   or more   blank lines denote the  end
of  a paragraph.

Since any number of consecutive spaces are treated
like a single one, the formatting of the input
file makes no difference to
      \LaTeX,                % The \LaTeX command generates the LaTeX logo.
but it makes a difference to you.  When you use
\LaTeX, making your input file as easy to read
as possible will be a great help as you write
your document and when you change it.  This sample
file shows how you can add comments to your own input
file.

Because printing is different from typewriting,
there are a number of things that you have to do
differently when preparing an input file than if
you were just typing the document directly.
Quotation marks like
       ``this''
have to be handled specially, as do quotes within
quotes:
       ``\,`this'            % \, separates the double and single quote.
        is what I just
        wrote, not  `that'\,''.

Dashes come in three sizes: an
       intra-word
dash, a medium dash for number ranges like
       1--2,
and a punctuation
       dash---like
this.

A sentence-ending space should be larger than the
space between words within a sentence.  You
sometimes have to type special commands in
conjunction with punctuation characters to get
this right, as in the following sentence.
       Gnats, gnus, etc.\ all  % `\ ' makes an inter-word space.
       begin with G\@.         % \@ marks end-of-sentence punctuation.
You should check the spaces after periods when
reading your output to make sure you haven't
forgotten any special cases.  Generating an
ellipsis
       \ldots\               % `\ ' is needed after `\ldots' because TeX
                             % ignores spaces after command names like \ldots
                             % made from \ + letters.
                             %
                             % Note how a `%' character causes TeX to ignore
                             % the end of the input line, so these blank lines
                             % do not start a new paragraph.
                             %
with the right spacing around the periods requires
a special command.

\LaTeX\ interprets some common characters as
commands, so you must type special commands to
generate them.  These characters include the
following:
       \$ \& \% \# \{ and \}.

In printing, text is usually emphasized with an
       \emph{italic}
type style.

\begin{em}
   A long segment of text can also be emphasized
   in this way.  Text within such a segment can be
   given \emph{additional} emphasis.
\end{em}

It is sometimes necessary to prevent \LaTeX\ from
breaking a line where it might otherwise do so.
This may be at a space, as between the ``Mr.''\ and
``Jones'' in
       ``Mr.~Jones'',        % ~ produces an unbreakable interword space.
or within a word---especially when the word is a
symbol like
       \mbox{\emph{itemnum}}
that makes little sense when hyphenated across
lines.

Footnotes\footnote{This is an example of a footnote.}
pose no problem.

\LaTeX\ is good at typesetting mathematical formulas
like
       \( x-3y + z = 7 \)
or
       \( a_{1} > x^{2n} + y^{2n} > x' \)
or
       \( \ip{A}{B} = \sum_{i} a_{i} b_{i} \).
The spaces you type in a formula are
ignored.  Remember that a letter like
       $x$                   % $ ... $  and  \( ... \)  are equivalent
is a formula when it denotes a mathematical
symbol, and it should be typed as one.

\section{Displayed Text}

Text is displayed by indenting it from the left
margin.  Quotations are commonly displayed.  There
are short quotations
\begin{quote}
   This is a short quotation.  It consists of a
   single paragraph of text.  See how it is formatted.
\end{quote}
and longer ones.
\begin{quotation}
   This is a longer quotation.  It consists of two
   paragraphs of text, neither of which are
   particularly interesting.

   This is the second paragraph of the quotation.  It
   is just as dull as the first paragraph.
\end{quotation}
Another frequently-displayed structure is a list.
The following is an example of an \emph{itemized}
list.
\begin{itemize}
   \item This is the first item of an itemized list.
         Each item in the list is marked with a ``tick''.
         You don't have to worry about what kind of tick
         mark is used.

   \item This is the second item of the list.  It
         contains another list nested inside it.  The inner
         list is an \emph{enumerated} list.
         \begin{enumerate}
            \item This is the first item of an enumerated
                  list that is nested within the itemized list.

            \item This is the second item of the inner list.
                  \LaTeX\ allows you to nest lists deeper than
                  you really should.
         \end{enumerate}
         This is the rest of the second item of the outer
         list.  It is no more interesting than any other
         part of the item.
   \item This is the third item of the list.
\end{itemize}
You can even display poetry.
\begin{verse}
   There is an environment
    for verse \\             % The \\ command separates lines
   Whose features some poets % within a stanza.
   will curse.

                             % One or more blank lines separate stanzas.

   For instead of making\\
   Them do \emph{all} line breaking, \\
   It allows them to put too many words on a line when they'd rather be
   forced to be terse.
\end{verse}

Mathematical formulas may also be displayed.  A
displayed formula
is
one-line long; multiline
formulas require special formatting instructions.
   \[  \ip{\Gamma}{\psi'} = x'' + y^{2} + z_{i}^{n}\]
Don't start a paragraph with a displayed equation,
nor make one a paragraph by itself.

\end{document}               % End of document.

\section{Bench Huge Section 1389}
% synthetic bench section 1389
% This is a sample LaTeX input file.  (Version of 12 August 2004.)
%
% A '%' character causes TeX to ignore all remaining text on the line,
% and is used for comments like this one.

\documentclass{article}      % Specifies the document class

                             % The preamble begins here.
\title{An Example Document}  % Declares the document's title.
\author{Leslie Lamport}      % Declares the author's name.
\date{January 21, 1994}      % Deleting this command produces today's date.

\newcommand{\ip}[2]{(#1, #2)}
                             % Defines \ip{arg1}{arg2} to mean
                             % (arg1, arg2).

%\newcommand{\ip}[2]{\langle #1 | #2\rangle}
                             % This is an alternative definition of
                             % \ip that is commented out.

\begin{document}             % End of preamble and beginning of text.

\maketitle                   % Produces the title.

This is an example input file.  Comparing it with
the output it generates can show you how to
produce a simple document of your own.

\section{Ordinary Text}      % Produces section heading.  Lower-level
                             % sections are begun with similar
                             % \subsection and \subsubsection commands.

The ends  of words and sentences are marked
  by   spaces. It  doesn't matter how many
spaces    you type; one is as good as 100.  The
end of   a line counts as a space.

One   or more   blank lines denote the  end
of  a paragraph.

Since any number of consecutive spaces are treated
like a single one, the formatting of the input
file makes no difference to
      \LaTeX,                % The \LaTeX command generates the LaTeX logo.
but it makes a difference to you.  When you use
\LaTeX, making your input file as easy to read
as possible will be a great help as you write
your document and when you change it.  This sample
file shows how you can add comments to your own input
file.

Because printing is different from typewriting,
there are a number of things that you have to do
differently when preparing an input file than if
you were just typing the document directly.
Quotation marks like
       ``this''
have to be handled specially, as do quotes within
quotes:
       ``\,`this'            % \, separates the double and single quote.
        is what I just
        wrote, not  `that'\,''.

Dashes come in three sizes: an
       intra-word
dash, a medium dash for number ranges like
       1--2,
and a punctuation
       dash---like
this.

A sentence-ending space should be larger than the
space between words within a sentence.  You
sometimes have to type special commands in
conjunction with punctuation characters to get
this right, as in the following sentence.
       Gnats, gnus, etc.\ all  % `\ ' makes an inter-word space.
       begin with G\@.         % \@ marks end-of-sentence punctuation.
You should check the spaces after periods when
reading your output to make sure you haven't
forgotten any special cases.  Generating an
ellipsis
       \ldots\               % `\ ' is needed after `\ldots' because TeX
                             % ignores spaces after command names like \ldots
                             % made from \ + letters.
                             %
                             % Note how a `%' character causes TeX to ignore
                             % the end of the input line, so these blank lines
                             % do not start a new paragraph.
                             %
with the right spacing around the periods requires
a special command.

\LaTeX\ interprets some common characters as
commands, so you must type special commands to
generate them.  These characters include the
following:
       \$ \& \% \# \{ and \}.

In printing, text is usually emphasized with an
       \emph{italic}
type style.

\begin{em}
   A long segment of text can also be emphasized
   in this way.  Text within such a segment can be
   given \emph{additional} emphasis.
\end{em}

It is sometimes necessary to prevent \LaTeX\ from
breaking a line where it might otherwise do so.
This may be at a space, as between the ``Mr.''\ and
``Jones'' in
       ``Mr.~Jones'',        % ~ produces an unbreakable interword space.
or within a word---especially when the word is a
symbol like
       \mbox{\emph{itemnum}}
that makes little sense when hyphenated across
lines.

Footnotes\footnote{This is an example of a footnote.}
pose no problem.

\LaTeX\ is good at typesetting mathematical formulas
like
       \( x-3y + z = 7 \)
or
       \( a_{1} > x^{2n} + y^{2n} > x' \)
or
       \( \ip{A}{B} = \sum_{i} a_{i} b_{i} \).
The spaces you type in a formula are
ignored.  Remember that a letter like
       $x$                   % $ ... $  and  \( ... \)  are equivalent
is a formula when it denotes a mathematical
symbol, and it should be typed as one.

\section{Displayed Text}

Text is displayed by indenting it from the left
margin.  Quotations are commonly displayed.  There
are short quotations
\begin{quote}
   This is a short quotation.  It consists of a
   single paragraph of text.  See how it is formatted.
\end{quote}
and longer ones.
\begin{quotation}
   This is a longer quotation.  It consists of two
   paragraphs of text, neither of which are
   particularly interesting.

   This is the second paragraph of the quotation.  It
   is just as dull as the first paragraph.
\end{quotation}
Another frequently-displayed structure is a list.
The following is an example of an \emph{itemized}
list.
\begin{itemize}
   \item This is the first item of an itemized list.
         Each item in the list is marked with a ``tick''.
         You don't have to worry about what kind of tick
         mark is used.

   \item This is the second item of the list.  It
         contains another list nested inside it.  The inner
         list is an \emph{enumerated} list.
         \begin{enumerate}
            \item This is the first item of an enumerated
                  list that is nested within the itemized list.

            \item This is the second item of the inner list.
                  \LaTeX\ allows you to nest lists deeper than
                  you really should.
         \end{enumerate}
         This is the rest of the second item of the outer
         list.  It is no more interesting than any other
         part of the item.
   \item This is the third item of the list.
\end{itemize}
You can even display poetry.
\begin{verse}
   There is an environment
    for verse \\             % The \\ command separates lines
   Whose features some poets % within a stanza.
   will curse.

                             % One or more blank lines separate stanzas.

   For instead of making\\
   Them do \emph{all} line breaking, \\
   It allows them to put too many words on a line when they'd rather be
   forced to be terse.
\end{verse}

Mathematical formulas may also be displayed.  A
displayed formula
is
one-line long; multiline
formulas require special formatting instructions.
   \[  \ip{\Gamma}{\psi'} = x'' + y^{2} + z_{i}^{n}\]
Don't start a paragraph with a displayed equation,
nor make one a paragraph by itself.

\end{document}               % End of document.

\section{Bench Huge Section 1390}
% synthetic bench section 1390
% This is a sample LaTeX input file.  (Version of 12 August 2004.)
%
% A '%' character causes TeX to ignore all remaining text on the line,
% and is used for comments like this one.

\documentclass{article}      % Specifies the document class

                             % The preamble begins here.
\title{An Example Document}  % Declares the document's title.
\author{Leslie Lamport}      % Declares the author's name.
\date{January 21, 1994}      % Deleting this command produces today's date.

\newcommand{\ip}[2]{(#1, #2)}
                             % Defines \ip{arg1}{arg2} to mean
                             % (arg1, arg2).

%\newcommand{\ip}[2]{\langle #1 | #2\rangle}
                             % This is an alternative definition of
                             % \ip that is commented out.

\begin{document}             % End of preamble and beginning of text.

\maketitle                   % Produces the title.

This is an example input file.  Comparing it with
the output it generates can show you how to
produce a simple document of your own.

\section{Ordinary Text}      % Produces section heading.  Lower-level
                             % sections are begun with similar
                             % \subsection and \subsubsection commands.

The ends  of words and sentences are marked
  by   spaces. It  doesn't matter how many
spaces    you type; one is as good as 100.  The
end of   a line counts as a space.

One   or more   blank lines denote the  end
of  a paragraph.

Since any number of consecutive spaces are treated
like a single one, the formatting of the input
file makes no difference to
      \LaTeX,                % The \LaTeX command generates the LaTeX logo.
but it makes a difference to you.  When you use
\LaTeX, making your input file as easy to read
as possible will be a great help as you write
your document and when you change it.  This sample
file shows how you can add comments to your own input
file.

Because printing is different from typewriting,
there are a number of things that you have to do
differently when preparing an input file than if
you were just typing the document directly.
Quotation marks like
       ``this''
have to be handled specially, as do quotes within
quotes:
       ``\,`this'            % \, separates the double and single quote.
        is what I just
        wrote, not  `that'\,''.

Dashes come in three sizes: an
       intra-word
dash, a medium dash for number ranges like
       1--2,
and a punctuation
       dash---like
this.

A sentence-ending space should be larger than the
space between words within a sentence.  You
sometimes have to type special commands in
conjunction with punctuation characters to get
this right, as in the following sentence.
       Gnats, gnus, etc.\ all  % `\ ' makes an inter-word space.
       begin with G\@.         % \@ marks end-of-sentence punctuation.
You should check the spaces after periods when
reading your output to make sure you haven't
forgotten any special cases.  Generating an
ellipsis
       \ldots\               % `\ ' is needed after `\ldots' because TeX
                             % ignores spaces after command names like \ldots
                             % made from \ + letters.
                             %
                             % Note how a `%' character causes TeX to ignore
                             % the end of the input line, so these blank lines
                             % do not start a new paragraph.
                             %
with the right spacing around the periods requires
a special command.

\LaTeX\ interprets some common characters as
commands, so you must type special commands to
generate them.  These characters include the
following:
       \$ \& \% \# \{ and \}.

In printing, text is usually emphasized with an
       \emph{italic}
type style.

\begin{em}
   A long segment of text can also be emphasized
   in this way.  Text within such a segment can be
   given \emph{additional} emphasis.
\end{em}

It is sometimes necessary to prevent \LaTeX\ from
breaking a line where it might otherwise do so.
This may be at a space, as between the ``Mr.''\ and
``Jones'' in
       ``Mr.~Jones'',        % ~ produces an unbreakable interword space.
or within a word---especially when the word is a
symbol like
       \mbox{\emph{itemnum}}
that makes little sense when hyphenated across
lines.

Footnotes\footnote{This is an example of a footnote.}
pose no problem.

\LaTeX\ is good at typesetting mathematical formulas
like
       \( x-3y + z = 7 \)
or
       \( a_{1} > x^{2n} + y^{2n} > x' \)
or
       \( \ip{A}{B} = \sum_{i} a_{i} b_{i} \).
The spaces you type in a formula are
ignored.  Remember that a letter like
       $x$                   % $ ... $  and  \( ... \)  are equivalent
is a formula when it denotes a mathematical
symbol, and it should be typed as one.

\section{Displayed Text}

Text is displayed by indenting it from the left
margin.  Quotations are commonly displayed.  There
are short quotations
\begin{quote}
   This is a short quotation.  It consists of a
   single paragraph of text.  See how it is formatted.
\end{quote}
and longer ones.
\begin{quotation}
   This is a longer quotation.  It consists of two
   paragraphs of text, neither of which are
   particularly interesting.

   This is the second paragraph of the quotation.  It
   is just as dull as the first paragraph.
\end{quotation}
Another frequently-displayed structure is a list.
The following is an example of an \emph{itemized}
list.
\begin{itemize}
   \item This is the first item of an itemized list.
         Each item in the list is marked with a ``tick''.
         You don't have to worry about what kind of tick
         mark is used.

   \item This is the second item of the list.  It
         contains another list nested inside it.  The inner
         list is an \emph{enumerated} list.
         \begin{enumerate}
            \item This is the first item of an enumerated
                  list that is nested within the itemized list.

            \item This is the second item of the inner list.
                  \LaTeX\ allows you to nest lists deeper than
                  you really should.
         \end{enumerate}
         This is the rest of the second item of the outer
         list.  It is no more interesting than any other
         part of the item.
   \item This is the third item of the list.
\end{itemize}
You can even display poetry.
\begin{verse}
   There is an environment
    for verse \\             % The \\ command separates lines
   Whose features some poets % within a stanza.
   will curse.

                             % One or more blank lines separate stanzas.

   For instead of making\\
   Them do \emph{all} line breaking, \\
   It allows them to put too many words on a line when they'd rather be
   forced to be terse.
\end{verse}

Mathematical formulas may also be displayed.  A
displayed formula
is
one-line long; multiline
formulas require special formatting instructions.
   \[  \ip{\Gamma}{\psi'} = x'' + y^{2} + z_{i}^{n}\]
Don't start a paragraph with a displayed equation,
nor make one a paragraph by itself.

\end{document}               % End of document.

\section{Bench Huge Section 1391}
% synthetic bench section 1391
% This is a sample LaTeX input file.  (Version of 12 August 2004.)
%
% A '%' character causes TeX to ignore all remaining text on the line,
% and is used for comments like this one.

\documentclass{article}      % Specifies the document class

                             % The preamble begins here.
\title{An Example Document}  % Declares the document's title.
\author{Leslie Lamport}      % Declares the author's name.
\date{January 21, 1994}      % Deleting this command produces today's date.

\newcommand{\ip}[2]{(#1, #2)}
                             % Defines \ip{arg1}{arg2} to mean
                             % (arg1, arg2).

%\newcommand{\ip}[2]{\langle #1 | #2\rangle}
                             % This is an alternative definition of
                             % \ip that is commented out.

\begin{document}             % End of preamble and beginning of text.

\maketitle                   % Produces the title.

This is an example input file.  Comparing it with
the output it generates can show you how to
produce a simple document of your own.

\section{Ordinary Text}      % Produces section heading.  Lower-level
                             % sections are begun with similar
                             % \subsection and \subsubsection commands.

The ends  of words and sentences are marked
  by   spaces. It  doesn't matter how many
spaces    you type; one is as good as 100.  The
end of   a line counts as a space.

One   or more   blank lines denote the  end
of  a paragraph.

Since any number of consecutive spaces are treated
like a single one, the formatting of the input
file makes no difference to
      \LaTeX,                % The \LaTeX command generates the LaTeX logo.
but it makes a difference to you.  When you use
\LaTeX, making your input file as easy to read
as possible will be a great help as you write
your document and when you change it.  This sample
file shows how you can add comments to your own input
file.

Because printing is different from typewriting,
there are a number of things that you have to do
differently when preparing an input file than if
you were just typing the document directly.
Quotation marks like
       ``this''
have to be handled specially, as do quotes within
quotes:
       ``\,`this'            % \, separates the double and single quote.
        is what I just
        wrote, not  `that'\,''.

Dashes come in three sizes: an
       intra-word
dash, a medium dash for number ranges like
       1--2,
and a punctuation
       dash---like
this.

A sentence-ending space should be larger than the
space between words within a sentence.  You
sometimes have to type special commands in
conjunction with punctuation characters to get
this right, as in the following sentence.
       Gnats, gnus, etc.\ all  % `\ ' makes an inter-word space.
       begin with G\@.         % \@ marks end-of-sentence punctuation.
You should check the spaces after periods when
reading your output to make sure you haven't
forgotten any special cases.  Generating an
ellipsis
       \ldots\               % `\ ' is needed after `\ldots' because TeX
                             % ignores spaces after command names like \ldots
                             % made from \ + letters.
                             %
                             % Note how a `%' character causes TeX to ignore
                             % the end of the input line, so these blank lines
                             % do not start a new paragraph.
                             %
with the right spacing around the periods requires
a special command.

\LaTeX\ interprets some common characters as
commands, so you must type special commands to
generate them.  These characters include the
following:
       \$ \& \% \# \{ and \}.

In printing, text is usually emphasized with an
       \emph{italic}
type style.

\begin{em}
   A long segment of text can also be emphasized
   in this way.  Text within such a segment can be
   given \emph{additional} emphasis.
\end{em}

It is sometimes necessary to prevent \LaTeX\ from
breaking a line where it might otherwise do so.
This may be at a space, as between the ``Mr.''\ and
``Jones'' in
       ``Mr.~Jones'',        % ~ produces an unbreakable interword space.
or within a word---especially when the word is a
symbol like
       \mbox{\emph{itemnum}}
that makes little sense when hyphenated across
lines.

Footnotes\footnote{This is an example of a footnote.}
pose no problem.

\LaTeX\ is good at typesetting mathematical formulas
like
       \( x-3y + z = 7 \)
or
       \( a_{1} > x^{2n} + y^{2n} > x' \)
or
       \( \ip{A}{B} = \sum_{i} a_{i} b_{i} \).
The spaces you type in a formula are
ignored.  Remember that a letter like
       $x$                   % $ ... $  and  \( ... \)  are equivalent
is a formula when it denotes a mathematical
symbol, and it should be typed as one.

\section{Displayed Text}

Text is displayed by indenting it from the left
margin.  Quotations are commonly displayed.  There
are short quotations
\begin{quote}
   This is a short quotation.  It consists of a
   single paragraph of text.  See how it is formatted.
\end{quote}
and longer ones.
\begin{quotation}
   This is a longer quotation.  It consists of two
   paragraphs of text, neither of which are
   particularly interesting.

   This is the second paragraph of the quotation.  It
   is just as dull as the first paragraph.
\end{quotation}
Another frequently-displayed structure is a list.
The following is an example of an \emph{itemized}
list.
\begin{itemize}
   \item This is the first item of an itemized list.
         Each item in the list is marked with a ``tick''.
         You don't have to worry about what kind of tick
         mark is used.

   \item This is the second item of the list.  It
         contains another list nested inside it.  The inner
         list is an \emph{enumerated} list.
         \begin{enumerate}
            \item This is the first item of an enumerated
                  list that is nested within the itemized list.

            \item This is the second item of the inner list.
                  \LaTeX\ allows you to nest lists deeper than
                  you really should.
         \end{enumerate}
         This is the rest of the second item of the outer
         list.  It is no more interesting than any other
         part of the item.
   \item This is the third item of the list.
\end{itemize}
You can even display poetry.
\begin{verse}
   There is an environment
    for verse \\             % The \\ command separates lines
   Whose features some poets % within a stanza.
   will curse.

                             % One or more blank lines separate stanzas.

   For instead of making\\
   Them do \emph{all} line breaking, \\
   It allows them to put too many words on a line when they'd rather be
   forced to be terse.
\end{verse}

Mathematical formulas may also be displayed.  A
displayed formula
is
one-line long; multiline
formulas require special formatting instructions.
   \[  \ip{\Gamma}{\psi'} = x'' + y^{2} + z_{i}^{n}\]
Don't start a paragraph with a displayed equation,
nor make one a paragraph by itself.

\end{document}               % End of document.

\section{Bench Huge Section 1392}
% synthetic bench section 1392
% This is a sample LaTeX input file.  (Version of 12 August 2004.)
%
% A '%' character causes TeX to ignore all remaining text on the line,
% and is used for comments like this one.

\documentclass{article}      % Specifies the document class

                             % The preamble begins here.
\title{An Example Document}  % Declares the document's title.
\author{Leslie Lamport}      % Declares the author's name.
\date{January 21, 1994}      % Deleting this command produces today's date.

\newcommand{\ip}[2]{(#1, #2)}
                             % Defines \ip{arg1}{arg2} to mean
                             % (arg1, arg2).

%\newcommand{\ip}[2]{\langle #1 | #2\rangle}
                             % This is an alternative definition of
                             % \ip that is commented out.

\begin{document}             % End of preamble and beginning of text.

\maketitle                   % Produces the title.

This is an example input file.  Comparing it with
the output it generates can show you how to
produce a simple document of your own.

\section{Ordinary Text}      % Produces section heading.  Lower-level
                             % sections are begun with similar
                             % \subsection and \subsubsection commands.

The ends  of words and sentences are marked
  by   spaces. It  doesn't matter how many
spaces    you type; one is as good as 100.  The
end of   a line counts as a space.

One   or more   blank lines denote the  end
of  a paragraph.

Since any number of consecutive spaces are treated
like a single one, the formatting of the input
file makes no difference to
      \LaTeX,                % The \LaTeX command generates the LaTeX logo.
but it makes a difference to you.  When you use
\LaTeX, making your input file as easy to read
as possible will be a great help as you write
your document and when you change it.  This sample
file shows how you can add comments to your own input
file.

Because printing is different from typewriting,
there are a number of things that you have to do
differently when preparing an input file than if
you were just typing the document directly.
Quotation marks like
       ``this''
have to be handled specially, as do quotes within
quotes:
       ``\,`this'            % \, separates the double and single quote.
        is what I just
        wrote, not  `that'\,''.

Dashes come in three sizes: an
       intra-word
dash, a medium dash for number ranges like
       1--2,
and a punctuation
       dash---like
this.

A sentence-ending space should be larger than the
space between words within a sentence.  You
sometimes have to type special commands in
conjunction with punctuation characters to get
this right, as in the following sentence.
       Gnats, gnus, etc.\ all  % `\ ' makes an inter-word space.
       begin with G\@.         % \@ marks end-of-sentence punctuation.
You should check the spaces after periods when
reading your output to make sure you haven't
forgotten any special cases.  Generating an
ellipsis
       \ldots\               % `\ ' is needed after `\ldots' because TeX
                             % ignores spaces after command names like \ldots
                             % made from \ + letters.
                             %
                             % Note how a `%' character causes TeX to ignore
                             % the end of the input line, so these blank lines
                             % do not start a new paragraph.
                             %
with the right spacing around the periods requires
a special command.

\LaTeX\ interprets some common characters as
commands, so you must type special commands to
generate them.  These characters include the
following:
       \$ \& \% \# \{ and \}.

In printing, text is usually emphasized with an
       \emph{italic}
type style.

\begin{em}
   A long segment of text can also be emphasized
   in this way.  Text within such a segment can be
   given \emph{additional} emphasis.
\end{em}

It is sometimes necessary to prevent \LaTeX\ from
breaking a line where it might otherwise do so.
This may be at a space, as between the ``Mr.''\ and
``Jones'' in
       ``Mr.~Jones'',        % ~ produces an unbreakable interword space.
or within a word---especially when the word is a
symbol like
       \mbox{\emph{itemnum}}
that makes little sense when hyphenated across
lines.

Footnotes\footnote{This is an example of a footnote.}
pose no problem.

\LaTeX\ is good at typesetting mathematical formulas
like
       \( x-3y + z = 7 \)
or
       \( a_{1} > x^{2n} + y^{2n} > x' \)
or
       \( \ip{A}{B} = \sum_{i} a_{i} b_{i} \).
The spaces you type in a formula are
ignored.  Remember that a letter like
       $x$                   % $ ... $  and  \( ... \)  are equivalent
is a formula when it denotes a mathematical
symbol, and it should be typed as one.

\section{Displayed Text}

Text is displayed by indenting it from the left
margin.  Quotations are commonly displayed.  There
are short quotations
\begin{quote}
   This is a short quotation.  It consists of a
   single paragraph of text.  See how it is formatted.
\end{quote}
and longer ones.
\begin{quotation}
   This is a longer quotation.  It consists of two
   paragraphs of text, neither of which are
   particularly interesting.

   This is the second paragraph of the quotation.  It
   is just as dull as the first paragraph.
\end{quotation}
Another frequently-displayed structure is a list.
The following is an example of an \emph{itemized}
list.
\begin{itemize}
   \item This is the first item of an itemized list.
         Each item in the list is marked with a ``tick''.
         You don't have to worry about what kind of tick
         mark is used.

   \item This is the second item of the list.  It
         contains another list nested inside it.  The inner
         list is an \emph{enumerated} list.
         \begin{enumerate}
            \item This is the first item of an enumerated
                  list that is nested within the itemized list.

            \item This is the second item of the inner list.
                  \LaTeX\ allows you to nest lists deeper than
                  you really should.
         \end{enumerate}
         This is the rest of the second item of the outer
         list.  It is no more interesting than any other
         part of the item.
   \item This is the third item of the list.
\end{itemize}
You can even display poetry.
\begin{verse}
   There is an environment
    for verse \\             % The \\ command separates lines
   Whose features some poets % within a stanza.
   will curse.

                             % One or more blank lines separate stanzas.

   For instead of making\\
   Them do \emph{all} line breaking, \\
   It allows them to put too many words on a line when they'd rather be
   forced to be terse.
\end{verse}

Mathematical formulas may also be displayed.  A
displayed formula
is
one-line long; multiline
formulas require special formatting instructions.
   \[  \ip{\Gamma}{\psi'} = x'' + y^{2} + z_{i}^{n}\]
Don't start a paragraph with a displayed equation,
nor make one a paragraph by itself.

\end{document}               % End of document.

\section{Bench Huge Section 1393}
% synthetic bench section 1393
% This is a sample LaTeX input file.  (Version of 12 August 2004.)
%
% A '%' character causes TeX to ignore all remaining text on the line,
% and is used for comments like this one.

\documentclass{article}      % Specifies the document class

                             % The preamble begins here.
\title{An Example Document}  % Declares the document's title.
\author{Leslie Lamport}      % Declares the author's name.
\date{January 21, 1994}      % Deleting this command produces today's date.

\newcommand{\ip}[2]{(#1, #2)}
                             % Defines \ip{arg1}{arg2} to mean
                             % (arg1, arg2).

%\newcommand{\ip}[2]{\langle #1 | #2\rangle}
                             % This is an alternative definition of
                             % \ip that is commented out.

\begin{document}             % End of preamble and beginning of text.

\maketitle                   % Produces the title.

This is an example input file.  Comparing it with
the output it generates can show you how to
produce a simple document of your own.

\section{Ordinary Text}      % Produces section heading.  Lower-level
                             % sections are begun with similar
                             % \subsection and \subsubsection commands.

The ends  of words and sentences are marked
  by   spaces. It  doesn't matter how many
spaces    you type; one is as good as 100.  The
end of   a line counts as a space.

One   or more   blank lines denote the  end
of  a paragraph.

Since any number of consecutive spaces are treated
like a single one, the formatting of the input
file makes no difference to
      \LaTeX,                % The \LaTeX command generates the LaTeX logo.
but it makes a difference to you.  When you use
\LaTeX, making your input file as easy to read
as possible will be a great help as you write
your document and when you change it.  This sample
file shows how you can add comments to your own input
file.

Because printing is different from typewriting,
there are a number of things that you have to do
differently when preparing an input file than if
you were just typing the document directly.
Quotation marks like
       ``this''
have to be handled specially, as do quotes within
quotes:
       ``\,`this'            % \, separates the double and single quote.
        is what I just
        wrote, not  `that'\,''.

Dashes come in three sizes: an
       intra-word
dash, a medium dash for number ranges like
       1--2,
and a punctuation
       dash---like
this.

A sentence-ending space should be larger than the
space between words within a sentence.  You
sometimes have to type special commands in
conjunction with punctuation characters to get
this right, as in the following sentence.
       Gnats, gnus, etc.\ all  % `\ ' makes an inter-word space.
       begin with G\@.         % \@ marks end-of-sentence punctuation.
You should check the spaces after periods when
reading your output to make sure you haven't
forgotten any special cases.  Generating an
ellipsis
       \ldots\               % `\ ' is needed after `\ldots' because TeX
                             % ignores spaces after command names like \ldots
                             % made from \ + letters.
                             %
                             % Note how a `%' character causes TeX to ignore
                             % the end of the input line, so these blank lines
                             % do not start a new paragraph.
                             %
with the right spacing around the periods requires
a special command.

\LaTeX\ interprets some common characters as
commands, so you must type special commands to
generate them.  These characters include the
following:
       \$ \& \% \# \{ and \}.

In printing, text is usually emphasized with an
       \emph{italic}
type style.

\begin{em}
   A long segment of text can also be emphasized
   in this way.  Text within such a segment can be
   given \emph{additional} emphasis.
\end{em}

It is sometimes necessary to prevent \LaTeX\ from
breaking a line where it might otherwise do so.
This may be at a space, as between the ``Mr.''\ and
``Jones'' in
       ``Mr.~Jones'',        % ~ produces an unbreakable interword space.
or within a word---especially when the word is a
symbol like
       \mbox{\emph{itemnum}}
that makes little sense when hyphenated across
lines.

Footnotes\footnote{This is an example of a footnote.}
pose no problem.

\LaTeX\ is good at typesetting mathematical formulas
like
       \( x-3y + z = 7 \)
or
       \( a_{1} > x^{2n} + y^{2n} > x' \)
or
       \( \ip{A}{B} = \sum_{i} a_{i} b_{i} \).
The spaces you type in a formula are
ignored.  Remember that a letter like
       $x$                   % $ ... $  and  \( ... \)  are equivalent
is a formula when it denotes a mathematical
symbol, and it should be typed as one.

\section{Displayed Text}

Text is displayed by indenting it from the left
margin.  Quotations are commonly displayed.  There
are short quotations
\begin{quote}
   This is a short quotation.  It consists of a
   single paragraph of text.  See how it is formatted.
\end{quote}
and longer ones.
\begin{quotation}
   This is a longer quotation.  It consists of two
   paragraphs of text, neither of which are
   particularly interesting.

   This is the second paragraph of the quotation.  It
   is just as dull as the first paragraph.
\end{quotation}
Another frequently-displayed structure is a list.
The following is an example of an \emph{itemized}
list.
\begin{itemize}
   \item This is the first item of an itemized list.
         Each item in the list is marked with a ``tick''.
         You don't have to worry about what kind of tick
         mark is used.

   \item This is the second item of the list.  It
         contains another list nested inside it.  The inner
         list is an \emph{enumerated} list.
         \begin{enumerate}
            \item This is the first item of an enumerated
                  list that is nested within the itemized list.

            \item This is the second item of the inner list.
                  \LaTeX\ allows you to nest lists deeper than
                  you really should.
         \end{enumerate}
         This is the rest of the second item of the outer
         list.  It is no more interesting than any other
         part of the item.
   \item This is the third item of the list.
\end{itemize}
You can even display poetry.
\begin{verse}
   There is an environment
    for verse \\             % The \\ command separates lines
   Whose features some poets % within a stanza.
   will curse.

                             % One or more blank lines separate stanzas.

   For instead of making\\
   Them do \emph{all} line breaking, \\
   It allows them to put too many words on a line when they'd rather be
   forced to be terse.
\end{verse}

Mathematical formulas may also be displayed.  A
displayed formula
is
one-line long; multiline
formulas require special formatting instructions.
   \[  \ip{\Gamma}{\psi'} = x'' + y^{2} + z_{i}^{n}\]
Don't start a paragraph with a displayed equation,
nor make one a paragraph by itself.

\end{document}               % End of document.

\section{Bench Huge Section 1394}
% synthetic bench section 1394
% This is a sample LaTeX input file.  (Version of 12 August 2004.)
%
% A '%' character causes TeX to ignore all remaining text on the line,
% and is used for comments like this one.

\documentclass{article}      % Specifies the document class

                             % The preamble begins here.
\title{An Example Document}  % Declares the document's title.
\author{Leslie Lamport}      % Declares the author's name.
\date{January 21, 1994}      % Deleting this command produces today's date.

\newcommand{\ip}[2]{(#1, #2)}
                             % Defines \ip{arg1}{arg2} to mean
                             % (arg1, arg2).

%\newcommand{\ip}[2]{\langle #1 | #2\rangle}
                             % This is an alternative definition of
                             % \ip that is commented out.

\begin{document}             % End of preamble and beginning of text.

\maketitle                   % Produces the title.

This is an example input file.  Comparing it with
the output it generates can show you how to
produce a simple document of your own.

\section{Ordinary Text}      % Produces section heading.  Lower-level
                             % sections are begun with similar
                             % \subsection and \subsubsection commands.

The ends  of words and sentences are marked
  by   spaces. It  doesn't matter how many
spaces    you type; one is as good as 100.  The
end of   a line counts as a space.

One   or more   blank lines denote the  end
of  a paragraph.

Since any number of consecutive spaces are treated
like a single one, the formatting of the input
file makes no difference to
      \LaTeX,                % The \LaTeX command generates the LaTeX logo.
but it makes a difference to you.  When you use
\LaTeX, making your input file as easy to read
as possible will be a great help as you write
your document and when you change it.  This sample
file shows how you can add comments to your own input
file.

Because printing is different from typewriting,
there are a number of things that you have to do
differently when preparing an input file than if
you were just typing the document directly.
Quotation marks like
       ``this''
have to be handled specially, as do quotes within
quotes:
       ``\,`this'            % \, separates the double and single quote.
        is what I just
        wrote, not  `that'\,''.

Dashes come in three sizes: an
       intra-word
dash, a medium dash for number ranges like
       1--2,
and a punctuation
       dash---like
this.

A sentence-ending space should be larger than the
space between words within a sentence.  You
sometimes have to type special commands in
conjunction with punctuation characters to get
this right, as in the following sentence.
       Gnats, gnus, etc.\ all  % `\ ' makes an inter-word space.
       begin with G\@.         % \@ marks end-of-sentence punctuation.
You should check the spaces after periods when
reading your output to make sure you haven't
forgotten any special cases.  Generating an
ellipsis
       \ldots\               % `\ ' is needed after `\ldots' because TeX
                             % ignores spaces after command names like \ldots
                             % made from \ + letters.
                             %
                             % Note how a `%' character causes TeX to ignore
                             % the end of the input line, so these blank lines
                             % do not start a new paragraph.
                             %
with the right spacing around the periods requires
a special command.

\LaTeX\ interprets some common characters as
commands, so you must type special commands to
generate them.  These characters include the
following:
       \$ \& \% \# \{ and \}.

In printing, text is usually emphasized with an
       \emph{italic}
type style.

\begin{em}
   A long segment of text can also be emphasized
   in this way.  Text within such a segment can be
   given \emph{additional} emphasis.
\end{em}

It is sometimes necessary to prevent \LaTeX\ from
breaking a line where it might otherwise do so.
This may be at a space, as between the ``Mr.''\ and
``Jones'' in
       ``Mr.~Jones'',        % ~ produces an unbreakable interword space.
or within a word---especially when the word is a
symbol like
       \mbox{\emph{itemnum}}
that makes little sense when hyphenated across
lines.

Footnotes\footnote{This is an example of a footnote.}
pose no problem.

\LaTeX\ is good at typesetting mathematical formulas
like
       \( x-3y + z = 7 \)
or
       \( a_{1} > x^{2n} + y^{2n} > x' \)
or
       \( \ip{A}{B} = \sum_{i} a_{i} b_{i} \).
The spaces you type in a formula are
ignored.  Remember that a letter like
       $x$                   % $ ... $  and  \( ... \)  are equivalent
is a formula when it denotes a mathematical
symbol, and it should be typed as one.

\section{Displayed Text}

Text is displayed by indenting it from the left
margin.  Quotations are commonly displayed.  There
are short quotations
\begin{quote}
   This is a short quotation.  It consists of a
   single paragraph of text.  See how it is formatted.
\end{quote}
and longer ones.
\begin{quotation}
   This is a longer quotation.  It consists of two
   paragraphs of text, neither of which are
   particularly interesting.

   This is the second paragraph of the quotation.  It
   is just as dull as the first paragraph.
\end{quotation}
Another frequently-displayed structure is a list.
The following is an example of an \emph{itemized}
list.
\begin{itemize}
   \item This is the first item of an itemized list.
         Each item in the list is marked with a ``tick''.
         You don't have to worry about what kind of tick
         mark is used.

   \item This is the second item of the list.  It
         contains another list nested inside it.  The inner
         list is an \emph{enumerated} list.
         \begin{enumerate}
            \item This is the first item of an enumerated
                  list that is nested within the itemized list.

            \item This is the second item of the inner list.
                  \LaTeX\ allows you to nest lists deeper than
                  you really should.
         \end{enumerate}
         This is the rest of the second item of the outer
         list.  It is no more interesting than any other
         part of the item.
   \item This is the third item of the list.
\end{itemize}
You can even display poetry.
\begin{verse}
   There is an environment
    for verse \\             % The \\ command separates lines
   Whose features some poets % within a stanza.
   will curse.

                             % One or more blank lines separate stanzas.

   For instead of making\\
   Them do \emph{all} line breaking, \\
   It allows them to put too many words on a line when they'd rather be
   forced to be terse.
\end{verse}

Mathematical formulas may also be displayed.  A
displayed formula
is
one-line long; multiline
formulas require special formatting instructions.
   \[  \ip{\Gamma}{\psi'} = x'' + y^{2} + z_{i}^{n}\]
Don't start a paragraph with a displayed equation,
nor make one a paragraph by itself.

\end{document}               % End of document.

\section{Bench Huge Section 1395}
% synthetic bench section 1395
% This is a sample LaTeX input file.  (Version of 12 August 2004.)
%
% A '%' character causes TeX to ignore all remaining text on the line,
% and is used for comments like this one.

\documentclass{article}      % Specifies the document class

                             % The preamble begins here.
\title{An Example Document}  % Declares the document's title.
\author{Leslie Lamport}      % Declares the author's name.
\date{January 21, 1994}      % Deleting this command produces today's date.

\newcommand{\ip}[2]{(#1, #2)}
                             % Defines \ip{arg1}{arg2} to mean
                             % (arg1, arg2).

%\newcommand{\ip}[2]{\langle #1 | #2\rangle}
                             % This is an alternative definition of
                             % \ip that is commented out.

\begin{document}             % End of preamble and beginning of text.

\maketitle                   % Produces the title.

This is an example input file.  Comparing it with
the output it generates can show you how to
produce a simple document of your own.

\section{Ordinary Text}      % Produces section heading.  Lower-level
                             % sections are begun with similar
                             % \subsection and \subsubsection commands.

The ends  of words and sentences are marked
  by   spaces. It  doesn't matter how many
spaces    you type; one is as good as 100.  The
end of   a line counts as a space.

One   or more   blank lines denote the  end
of  a paragraph.

Since any number of consecutive spaces are treated
like a single one, the formatting of the input
file makes no difference to
      \LaTeX,                % The \LaTeX command generates the LaTeX logo.
but it makes a difference to you.  When you use
\LaTeX, making your input file as easy to read
as possible will be a great help as you write
your document and when you change it.  This sample
file shows how you can add comments to your own input
file.

Because printing is different from typewriting,
there are a number of things that you have to do
differently when preparing an input file than if
you were just typing the document directly.
Quotation marks like
       ``this''
have to be handled specially, as do quotes within
quotes:
       ``\,`this'            % \, separates the double and single quote.
        is what I just
        wrote, not  `that'\,''.

Dashes come in three sizes: an
       intra-word
dash, a medium dash for number ranges like
       1--2,
and a punctuation
       dash---like
this.

A sentence-ending space should be larger than the
space between words within a sentence.  You
sometimes have to type special commands in
conjunction with punctuation characters to get
this right, as in the following sentence.
       Gnats, gnus, etc.\ all  % `\ ' makes an inter-word space.
       begin with G\@.         % \@ marks end-of-sentence punctuation.
You should check the spaces after periods when
reading your output to make sure you haven't
forgotten any special cases.  Generating an
ellipsis
       \ldots\               % `\ ' is needed after `\ldots' because TeX
                             % ignores spaces after command names like \ldots
                             % made from \ + letters.
                             %
                             % Note how a `%' character causes TeX to ignore
                             % the end of the input line, so these blank lines
                             % do not start a new paragraph.
                             %
with the right spacing around the periods requires
a special command.

\LaTeX\ interprets some common characters as
commands, so you must type special commands to
generate them.  These characters include the
following:
       \$ \& \% \# \{ and \}.

In printing, text is usually emphasized with an
       \emph{italic}
type style.

\begin{em}
   A long segment of text can also be emphasized
   in this way.  Text within such a segment can be
   given \emph{additional} emphasis.
\end{em}

It is sometimes necessary to prevent \LaTeX\ from
breaking a line where it might otherwise do so.
This may be at a space, as between the ``Mr.''\ and
``Jones'' in
       ``Mr.~Jones'',        % ~ produces an unbreakable interword space.
or within a word---especially when the word is a
symbol like
       \mbox{\emph{itemnum}}
that makes little sense when hyphenated across
lines.

Footnotes\footnote{This is an example of a footnote.}
pose no problem.

\LaTeX\ is good at typesetting mathematical formulas
like
       \( x-3y + z = 7 \)
or
       \( a_{1} > x^{2n} + y^{2n} > x' \)
or
       \( \ip{A}{B} = \sum_{i} a_{i} b_{i} \).
The spaces you type in a formula are
ignored.  Remember that a letter like
       $x$                   % $ ... $  and  \( ... \)  are equivalent
is a formula when it denotes a mathematical
symbol, and it should be typed as one.

\section{Displayed Text}

Text is displayed by indenting it from the left
margin.  Quotations are commonly displayed.  There
are short quotations
\begin{quote}
   This is a short quotation.  It consists of a
   single paragraph of text.  See how it is formatted.
\end{quote}
and longer ones.
\begin{quotation}
   This is a longer quotation.  It consists of two
   paragraphs of text, neither of which are
   particularly interesting.

   This is the second paragraph of the quotation.  It
   is just as dull as the first paragraph.
\end{quotation}
Another frequently-displayed structure is a list.
The following is an example of an \emph{itemized}
list.
\begin{itemize}
   \item This is the first item of an itemized list.
         Each item in the list is marked with a ``tick''.
         You don't have to worry about what kind of tick
         mark is used.

   \item This is the second item of the list.  It
         contains another list nested inside it.  The inner
         list is an \emph{enumerated} list.
         \begin{enumerate}
            \item This is the first item of an enumerated
                  list that is nested within the itemized list.

            \item This is the second item of the inner list.
                  \LaTeX\ allows you to nest lists deeper than
                  you really should.
         \end{enumerate}
         This is the rest of the second item of the outer
         list.  It is no more interesting than any other
         part of the item.
   \item This is the third item of the list.
\end{itemize}
You can even display poetry.
\begin{verse}
   There is an environment
    for verse \\             % The \\ command separates lines
   Whose features some poets % within a stanza.
   will curse.

                             % One or more blank lines separate stanzas.

   For instead of making\\
   Them do \emph{all} line breaking, \\
   It allows them to put too many words on a line when they'd rather be
   forced to be terse.
\end{verse}

Mathematical formulas may also be displayed.  A
displayed formula
is
one-line long; multiline
formulas require special formatting instructions.
   \[  \ip{\Gamma}{\psi'} = x'' + y^{2} + z_{i}^{n}\]
Don't start a paragraph with a displayed equation,
nor make one a paragraph by itself.

\end{document}               % End of document.

\section{Bench Huge Section 1396}
% synthetic bench section 1396
% This is a sample LaTeX input file.  (Version of 12 August 2004.)
%
% A '%' character causes TeX to ignore all remaining text on the line,
% and is used for comments like this one.

\documentclass{article}      % Specifies the document class

                             % The preamble begins here.
\title{An Example Document}  % Declares the document's title.
\author{Leslie Lamport}      % Declares the author's name.
\date{January 21, 1994}      % Deleting this command produces today's date.

\newcommand{\ip}[2]{(#1, #2)}
                             % Defines \ip{arg1}{arg2} to mean
                             % (arg1, arg2).

%\newcommand{\ip}[2]{\langle #1 | #2\rangle}
                             % This is an alternative definition of
                             % \ip that is commented out.

\begin{document}             % End of preamble and beginning of text.

\maketitle                   % Produces the title.

This is an example input file.  Comparing it with
the output it generates can show you how to
produce a simple document of your own.

\section{Ordinary Text}      % Produces section heading.  Lower-level
                             % sections are begun with similar
                             % \subsection and \subsubsection commands.

The ends  of words and sentences are marked
  by   spaces. It  doesn't matter how many
spaces    you type; one is as good as 100.  The
end of   a line counts as a space.

One   or more   blank lines denote the  end
of  a paragraph.

Since any number of consecutive spaces are treated
like a single one, the formatting of the input
file makes no difference to
      \LaTeX,                % The \LaTeX command generates the LaTeX logo.
but it makes a difference to you.  When you use
\LaTeX, making your input file as easy to read
as possible will be a great help as you write
your document and when you change it.  This sample
file shows how you can add comments to your own input
file.

Because printing is different from typewriting,
there are a number of things that you have to do
differently when preparing an input file than if
you were just typing the document directly.
Quotation marks like
       ``this''
have to be handled specially, as do quotes within
quotes:
       ``\,`this'            % \, separates the double and single quote.
        is what I just
        wrote, not  `that'\,''.

Dashes come in three sizes: an
       intra-word
dash, a medium dash for number ranges like
       1--2,
and a punctuation
       dash---like
this.

A sentence-ending space should be larger than the
space between words within a sentence.  You
sometimes have to type special commands in
conjunction with punctuation characters to get
this right, as in the following sentence.
       Gnats, gnus, etc.\ all  % `\ ' makes an inter-word space.
       begin with G\@.         % \@ marks end-of-sentence punctuation.
You should check the spaces after periods when
reading your output to make sure you haven't
forgotten any special cases.  Generating an
ellipsis
       \ldots\               % `\ ' is needed after `\ldots' because TeX
                             % ignores spaces after command names like \ldots
                             % made from \ + letters.
                             %
                             % Note how a `%' character causes TeX to ignore
                             % the end of the input line, so these blank lines
                             % do not start a new paragraph.
                             %
with the right spacing around the periods requires
a special command.

\LaTeX\ interprets some common characters as
commands, so you must type special commands to
generate them.  These characters include the
following:
       \$ \& \% \# \{ and \}.

In printing, text is usually emphasized with an
       \emph{italic}
type style.

\begin{em}
   A long segment of text can also be emphasized
   in this way.  Text within such a segment can be
   given \emph{additional} emphasis.
\end{em}

It is sometimes necessary to prevent \LaTeX\ from
breaking a line where it might otherwise do so.
This may be at a space, as between the ``Mr.''\ and
``Jones'' in
       ``Mr.~Jones'',        % ~ produces an unbreakable interword space.
or within a word---especially when the word is a
symbol like
       \mbox{\emph{itemnum}}
that makes little sense when hyphenated across
lines.

Footnotes\footnote{This is an example of a footnote.}
pose no problem.

\LaTeX\ is good at typesetting mathematical formulas
like
       \( x-3y + z = 7 \)
or
       \( a_{1} > x^{2n} + y^{2n} > x' \)
or
       \( \ip{A}{B} = \sum_{i} a_{i} b_{i} \).
The spaces you type in a formula are
ignored.  Remember that a letter like
       $x$                   % $ ... $  and  \( ... \)  are equivalent
is a formula when it denotes a mathematical
symbol, and it should be typed as one.

\section{Displayed Text}

Text is displayed by indenting it from the left
margin.  Quotations are commonly displayed.  There
are short quotations
\begin{quote}
   This is a short quotation.  It consists of a
   single paragraph of text.  See how it is formatted.
\end{quote}
and longer ones.
\begin{quotation}
   This is a longer quotation.  It consists of two
   paragraphs of text, neither of which are
   particularly interesting.

   This is the second paragraph of the quotation.  It
   is just as dull as the first paragraph.
\end{quotation}
Another frequently-displayed structure is a list.
The following is an example of an \emph{itemized}
list.
\begin{itemize}
   \item This is the first item of an itemized list.
         Each item in the list is marked with a ``tick''.
         You don't have to worry about what kind of tick
         mark is used.

   \item This is the second item of the list.  It
         contains another list nested inside it.  The inner
         list is an \emph{enumerated} list.
         \begin{enumerate}
            \item This is the first item of an enumerated
                  list that is nested within the itemized list.

            \item This is the second item of the inner list.
                  \LaTeX\ allows you to nest lists deeper than
                  you really should.
         \end{enumerate}
         This is the rest of the second item of the outer
         list.  It is no more interesting than any other
         part of the item.
   \item This is the third item of the list.
\end{itemize}
You can even display poetry.
\begin{verse}
   There is an environment
    for verse \\             % The \\ command separates lines
   Whose features some poets % within a stanza.
   will curse.

                             % One or more blank lines separate stanzas.

   For instead of making\\
   Them do \emph{all} line breaking, \\
   It allows them to put too many words on a line when they'd rather be
   forced to be terse.
\end{verse}

Mathematical formulas may also be displayed.  A
displayed formula
is
one-line long; multiline
formulas require special formatting instructions.
   \[  \ip{\Gamma}{\psi'} = x'' + y^{2} + z_{i}^{n}\]
Don't start a paragraph with a displayed equation,
nor make one a paragraph by itself.

\end{document}               % End of document.

\section{Bench Huge Section 1397}
% synthetic bench section 1397
% This is a sample LaTeX input file.  (Version of 12 August 2004.)
%
% A '%' character causes TeX to ignore all remaining text on the line,
% and is used for comments like this one.

\documentclass{article}      % Specifies the document class

                             % The preamble begins here.
\title{An Example Document}  % Declares the document's title.
\author{Leslie Lamport}      % Declares the author's name.
\date{January 21, 1994}      % Deleting this command produces today's date.

\newcommand{\ip}[2]{(#1, #2)}
                             % Defines \ip{arg1}{arg2} to mean
                             % (arg1, arg2).

%\newcommand{\ip}[2]{\langle #1 | #2\rangle}
                             % This is an alternative definition of
                             % \ip that is commented out.

\begin{document}             % End of preamble and beginning of text.

\maketitle                   % Produces the title.

This is an example input file.  Comparing it with
the output it generates can show you how to
produce a simple document of your own.

\section{Ordinary Text}      % Produces section heading.  Lower-level
                             % sections are begun with similar
                             % \subsection and \subsubsection commands.

The ends  of words and sentences are marked
  by   spaces. It  doesn't matter how many
spaces    you type; one is as good as 100.  The
end of   a line counts as a space.

One   or more   blank lines denote the  end
of  a paragraph.

Since any number of consecutive spaces are treated
like a single one, the formatting of the input
file makes no difference to
      \LaTeX,                % The \LaTeX command generates the LaTeX logo.
but it makes a difference to you.  When you use
\LaTeX, making your input file as easy to read
as possible will be a great help as you write
your document and when you change it.  This sample
file shows how you can add comments to your own input
file.

Because printing is different from typewriting,
there are a number of things that you have to do
differently when preparing an input file than if
you were just typing the document directly.
Quotation marks like
       ``this''
have to be handled specially, as do quotes within
quotes:
       ``\,`this'            % \, separates the double and single quote.
        is what I just
        wrote, not  `that'\,''.

Dashes come in three sizes: an
       intra-word
dash, a medium dash for number ranges like
       1--2,
and a punctuation
       dash---like
this.

A sentence-ending space should be larger than the
space between words within a sentence.  You
sometimes have to type special commands in
conjunction with punctuation characters to get
this right, as in the following sentence.
       Gnats, gnus, etc.\ all  % `\ ' makes an inter-word space.
       begin with G\@.         % \@ marks end-of-sentence punctuation.
You should check the spaces after periods when
reading your output to make sure you haven't
forgotten any special cases.  Generating an
ellipsis
       \ldots\               % `\ ' is needed after `\ldots' because TeX
                             % ignores spaces after command names like \ldots
                             % made from \ + letters.
                             %
                             % Note how a `%' character causes TeX to ignore
                             % the end of the input line, so these blank lines
                             % do not start a new paragraph.
                             %
with the right spacing around the periods requires
a special command.

\LaTeX\ interprets some common characters as
commands, so you must type special commands to
generate them.  These characters include the
following:
       \$ \& \% \# \{ and \}.

In printing, text is usually emphasized with an
       \emph{italic}
type style.

\begin{em}
   A long segment of text can also be emphasized
   in this way.  Text within such a segment can be
   given \emph{additional} emphasis.
\end{em}

It is sometimes necessary to prevent \LaTeX\ from
breaking a line where it might otherwise do so.
This may be at a space, as between the ``Mr.''\ and
``Jones'' in
       ``Mr.~Jones'',        % ~ produces an unbreakable interword space.
or within a word---especially when the word is a
symbol like
       \mbox{\emph{itemnum}}
that makes little sense when hyphenated across
lines.

Footnotes\footnote{This is an example of a footnote.}
pose no problem.

\LaTeX\ is good at typesetting mathematical formulas
like
       \( x-3y + z = 7 \)
or
       \( a_{1} > x^{2n} + y^{2n} > x' \)
or
       \( \ip{A}{B} = \sum_{i} a_{i} b_{i} \).
The spaces you type in a formula are
ignored.  Remember that a letter like
       $x$                   % $ ... $  and  \( ... \)  are equivalent
is a formula when it denotes a mathematical
symbol, and it should be typed as one.

\section{Displayed Text}

Text is displayed by indenting it from the left
margin.  Quotations are commonly displayed.  There
are short quotations
\begin{quote}
   This is a short quotation.  It consists of a
   single paragraph of text.  See how it is formatted.
\end{quote}
and longer ones.
\begin{quotation}
   This is a longer quotation.  It consists of two
   paragraphs of text, neither of which are
   particularly interesting.

   This is the second paragraph of the quotation.  It
   is just as dull as the first paragraph.
\end{quotation}
Another frequently-displayed structure is a list.
The following is an example of an \emph{itemized}
list.
\begin{itemize}
   \item This is the first item of an itemized list.
         Each item in the list is marked with a ``tick''.
         You don't have to worry about what kind of tick
         mark is used.

   \item This is the second item of the list.  It
         contains another list nested inside it.  The inner
         list is an \emph{enumerated} list.
         \begin{enumerate}
            \item This is the first item of an enumerated
                  list that is nested within the itemized list.

            \item This is the second item of the inner list.
                  \LaTeX\ allows you to nest lists deeper than
                  you really should.
         \end{enumerate}
         This is the rest of the second item of the outer
         list.  It is no more interesting than any other
         part of the item.
   \item This is the third item of the list.
\end{itemize}
You can even display poetry.
\begin{verse}
   There is an environment
    for verse \\             % The \\ command separates lines
   Whose features some poets % within a stanza.
   will curse.

                             % One or more blank lines separate stanzas.

   For instead of making\\
   Them do \emph{all} line breaking, \\
   It allows them to put too many words on a line when they'd rather be
   forced to be terse.
\end{verse}

Mathematical formulas may also be displayed.  A
displayed formula
is
one-line long; multiline
formulas require special formatting instructions.
   \[  \ip{\Gamma}{\psi'} = x'' + y^{2} + z_{i}^{n}\]
Don't start a paragraph with a displayed equation,
nor make one a paragraph by itself.

\end{document}               % End of document.

\section{Bench Huge Section 1398}
% synthetic bench section 1398
% This is a sample LaTeX input file.  (Version of 12 August 2004.)
%
% A '%' character causes TeX to ignore all remaining text on the line,
% and is used for comments like this one.

\documentclass{article}      % Specifies the document class

                             % The preamble begins here.
\title{An Example Document}  % Declares the document's title.
\author{Leslie Lamport}      % Declares the author's name.
\date{January 21, 1994}      % Deleting this command produces today's date.

\newcommand{\ip}[2]{(#1, #2)}
                             % Defines \ip{arg1}{arg2} to mean
                             % (arg1, arg2).

%\newcommand{\ip}[2]{\langle #1 | #2\rangle}
                             % This is an alternative definition of
                             % \ip that is commented out.

\begin{document}             % End of preamble and beginning of text.

\maketitle                   % Produces the title.

This is an example input file.  Comparing it with
the output it generates can show you how to
produce a simple document of your own.

\section{Ordinary Text}      % Produces section heading.  Lower-level
                             % sections are begun with similar
                             % \subsection and \subsubsection commands.

The ends  of words and sentences are marked
  by   spaces. It  doesn't matter how many
spaces    you type; one is as good as 100.  The
end of   a line counts as a space.

One   or more   blank lines denote the  end
of  a paragraph.

Since any number of consecutive spaces are treated
like a single one, the formatting of the input
file makes no difference to
      \LaTeX,                % The \LaTeX command generates the LaTeX logo.
but it makes a difference to you.  When you use
\LaTeX, making your input file as easy to read
as possible will be a great help as you write
your document and when you change it.  This sample
file shows how you can add comments to your own input
file.

Because printing is different from typewriting,
there are a number of things that you have to do
differently when preparing an input file than if
you were just typing the document directly.
Quotation marks like
       ``this''
have to be handled specially, as do quotes within
quotes:
       ``\,`this'            % \, separates the double and single quote.
        is what I just
        wrote, not  `that'\,''.

Dashes come in three sizes: an
       intra-word
dash, a medium dash for number ranges like
       1--2,
and a punctuation
       dash---like
this.

A sentence-ending space should be larger than the
space between words within a sentence.  You
sometimes have to type special commands in
conjunction with punctuation characters to get
this right, as in the following sentence.
       Gnats, gnus, etc.\ all  % `\ ' makes an inter-word space.
       begin with G\@.         % \@ marks end-of-sentence punctuation.
You should check the spaces after periods when
reading your output to make sure you haven't
forgotten any special cases.  Generating an
ellipsis
       \ldots\               % `\ ' is needed after `\ldots' because TeX
                             % ignores spaces after command names like \ldots
                             % made from \ + letters.
                             %
                             % Note how a `%' character causes TeX to ignore
                             % the end of the input line, so these blank lines
                             % do not start a new paragraph.
                             %
with the right spacing around the periods requires
a special command.

\LaTeX\ interprets some common characters as
commands, so you must type special commands to
generate them.  These characters include the
following:
       \$ \& \% \# \{ and \}.

In printing, text is usually emphasized with an
       \emph{italic}
type style.

\begin{em}
   A long segment of text can also be emphasized
   in this way.  Text within such a segment can be
   given \emph{additional} emphasis.
\end{em}

It is sometimes necessary to prevent \LaTeX\ from
breaking a line where it might otherwise do so.
This may be at a space, as between the ``Mr.''\ and
``Jones'' in
       ``Mr.~Jones'',        % ~ produces an unbreakable interword space.
or within a word---especially when the word is a
symbol like
       \mbox{\emph{itemnum}}
that makes little sense when hyphenated across
lines.

Footnotes\footnote{This is an example of a footnote.}
pose no problem.

\LaTeX\ is good at typesetting mathematical formulas
like
       \( x-3y + z = 7 \)
or
       \( a_{1} > x^{2n} + y^{2n} > x' \)
or
       \( \ip{A}{B} = \sum_{i} a_{i} b_{i} \).
The spaces you type in a formula are
ignored.  Remember that a letter like
       $x$                   % $ ... $  and  \( ... \)  are equivalent
is a formula when it denotes a mathematical
symbol, and it should be typed as one.

\section{Displayed Text}

Text is displayed by indenting it from the left
margin.  Quotations are commonly displayed.  There
are short quotations
\begin{quote}
   This is a short quotation.  It consists of a
   single paragraph of text.  See how it is formatted.
\end{quote}
and longer ones.
\begin{quotation}
   This is a longer quotation.  It consists of two
   paragraphs of text, neither of which are
   particularly interesting.

   This is the second paragraph of the quotation.  It
   is just as dull as the first paragraph.
\end{quotation}
Another frequently-displayed structure is a list.
The following is an example of an \emph{itemized}
list.
\begin{itemize}
   \item This is the first item of an itemized list.
         Each item in the list is marked with a ``tick''.
         You don't have to worry about what kind of tick
         mark is used.

   \item This is the second item of the list.  It
         contains another list nested inside it.  The inner
         list is an \emph{enumerated} list.
         \begin{enumerate}
            \item This is the first item of an enumerated
                  list that is nested within the itemized list.

            \item This is the second item of the inner list.
                  \LaTeX\ allows you to nest lists deeper than
                  you really should.
         \end{enumerate}
         This is the rest of the second item of the outer
         list.  It is no more interesting than any other
         part of the item.
   \item This is the third item of the list.
\end{itemize}
You can even display poetry.
\begin{verse}
   There is an environment
    for verse \\             % The \\ command separates lines
   Whose features some poets % within a stanza.
   will curse.

                             % One or more blank lines separate stanzas.

   For instead of making\\
   Them do \emph{all} line breaking, \\
   It allows them to put too many words on a line when they'd rather be
   forced to be terse.
\end{verse}

Mathematical formulas may also be displayed.  A
displayed formula
is
one-line long; multiline
formulas require special formatting instructions.
   \[  \ip{\Gamma}{\psi'} = x'' + y^{2} + z_{i}^{n}\]
Don't start a paragraph with a displayed equation,
nor make one a paragraph by itself.

\end{document}               % End of document.

\section{Bench Huge Section 1399}
% synthetic bench section 1399
% This is a sample LaTeX input file.  (Version of 12 August 2004.)
%
% A '%' character causes TeX to ignore all remaining text on the line,
% and is used for comments like this one.

\documentclass{article}      % Specifies the document class

                             % The preamble begins here.
\title{An Example Document}  % Declares the document's title.
\author{Leslie Lamport}      % Declares the author's name.
\date{January 21, 1994}      % Deleting this command produces today's date.

\newcommand{\ip}[2]{(#1, #2)}
                             % Defines \ip{arg1}{arg2} to mean
                             % (arg1, arg2).

%\newcommand{\ip}[2]{\langle #1 | #2\rangle}
                             % This is an alternative definition of
                             % \ip that is commented out.

\begin{document}             % End of preamble and beginning of text.

\maketitle                   % Produces the title.

This is an example input file.  Comparing it with
the output it generates can show you how to
produce a simple document of your own.

\section{Ordinary Text}      % Produces section heading.  Lower-level
                             % sections are begun with similar
                             % \subsection and \subsubsection commands.

The ends  of words and sentences are marked
  by   spaces. It  doesn't matter how many
spaces    you type; one is as good as 100.  The
end of   a line counts as a space.

One   or more   blank lines denote the  end
of  a paragraph.

Since any number of consecutive spaces are treated
like a single one, the formatting of the input
file makes no difference to
      \LaTeX,                % The \LaTeX command generates the LaTeX logo.
but it makes a difference to you.  When you use
\LaTeX, making your input file as easy to read
as possible will be a great help as you write
your document and when you change it.  This sample
file shows how you can add comments to your own input
file.

Because printing is different from typewriting,
there are a number of things that you have to do
differently when preparing an input file than if
you were just typing the document directly.
Quotation marks like
       ``this''
have to be handled specially, as do quotes within
quotes:
       ``\,`this'            % \, separates the double and single quote.
        is what I just
        wrote, not  `that'\,''.

Dashes come in three sizes: an
       intra-word
dash, a medium dash for number ranges like
       1--2,
and a punctuation
       dash---like
this.

A sentence-ending space should be larger than the
space between words within a sentence.  You
sometimes have to type special commands in
conjunction with punctuation characters to get
this right, as in the following sentence.
       Gnats, gnus, etc.\ all  % `\ ' makes an inter-word space.
       begin with G\@.         % \@ marks end-of-sentence punctuation.
You should check the spaces after periods when
reading your output to make sure you haven't
forgotten any special cases.  Generating an
ellipsis
       \ldots\               % `\ ' is needed after `\ldots' because TeX
                             % ignores spaces after command names like \ldots
                             % made from \ + letters.
                             %
                             % Note how a `%' character causes TeX to ignore
                             % the end of the input line, so these blank lines
                             % do not start a new paragraph.
                             %
with the right spacing around the periods requires
a special command.

\LaTeX\ interprets some common characters as
commands, so you must type special commands to
generate them.  These characters include the
following:
       \$ \& \% \# \{ and \}.

In printing, text is usually emphasized with an
       \emph{italic}
type style.

\begin{em}
   A long segment of text can also be emphasized
   in this way.  Text within such a segment can be
   given \emph{additional} emphasis.
\end{em}

It is sometimes necessary to prevent \LaTeX\ from
breaking a line where it might otherwise do so.
This may be at a space, as between the ``Mr.''\ and
``Jones'' in
       ``Mr.~Jones'',        % ~ produces an unbreakable interword space.
or within a word---especially when the word is a
symbol like
       \mbox{\emph{itemnum}}
that makes little sense when hyphenated across
lines.

Footnotes\footnote{This is an example of a footnote.}
pose no problem.

\LaTeX\ is good at typesetting mathematical formulas
like
       \( x-3y + z = 7 \)
or
       \( a_{1} > x^{2n} + y^{2n} > x' \)
or
       \( \ip{A}{B} = \sum_{i} a_{i} b_{i} \).
The spaces you type in a formula are
ignored.  Remember that a letter like
       $x$                   % $ ... $  and  \( ... \)  are equivalent
is a formula when it denotes a mathematical
symbol, and it should be typed as one.

\section{Displayed Text}

Text is displayed by indenting it from the left
margin.  Quotations are commonly displayed.  There
are short quotations
\begin{quote}
   This is a short quotation.  It consists of a
   single paragraph of text.  See how it is formatted.
\end{quote}
and longer ones.
\begin{quotation}
   This is a longer quotation.  It consists of two
   paragraphs of text, neither of which are
   particularly interesting.

   This is the second paragraph of the quotation.  It
   is just as dull as the first paragraph.
\end{quotation}
Another frequently-displayed structure is a list.
The following is an example of an \emph{itemized}
list.
\begin{itemize}
   \item This is the first item of an itemized list.
         Each item in the list is marked with a ``tick''.
         You don't have to worry about what kind of tick
         mark is used.

   \item This is the second item of the list.  It
         contains another list nested inside it.  The inner
         list is an \emph{enumerated} list.
         \begin{enumerate}
            \item This is the first item of an enumerated
                  list that is nested within the itemized list.

            \item This is the second item of the inner list.
                  \LaTeX\ allows you to nest lists deeper than
                  you really should.
         \end{enumerate}
         This is the rest of the second item of the outer
         list.  It is no more interesting than any other
         part of the item.
   \item This is the third item of the list.
\end{itemize}
You can even display poetry.
\begin{verse}
   There is an environment
    for verse \\             % The \\ command separates lines
   Whose features some poets % within a stanza.
   will curse.

                             % One or more blank lines separate stanzas.

   For instead of making\\
   Them do \emph{all} line breaking, \\
   It allows them to put too many words on a line when they'd rather be
   forced to be terse.
\end{verse}

Mathematical formulas may also be displayed.  A
displayed formula
is
one-line long; multiline
formulas require special formatting instructions.
   \[  \ip{\Gamma}{\psi'} = x'' + y^{2} + z_{i}^{n}\]
Don't start a paragraph with a displayed equation,
nor make one a paragraph by itself.

\end{document}               % End of document.

\section{Bench Huge Section 1400}
% synthetic bench section 1400
% This is a sample LaTeX input file.  (Version of 12 August 2004.)
%
% A '%' character causes TeX to ignore all remaining text on the line,
% and is used for comments like this one.

\documentclass{article}      % Specifies the document class

                             % The preamble begins here.
\title{An Example Document}  % Declares the document's title.
\author{Leslie Lamport}      % Declares the author's name.
\date{January 21, 1994}      % Deleting this command produces today's date.

\newcommand{\ip}[2]{(#1, #2)}
                             % Defines \ip{arg1}{arg2} to mean
                             % (arg1, arg2).

%\newcommand{\ip}[2]{\langle #1 | #2\rangle}
                             % This is an alternative definition of
                             % \ip that is commented out.

\begin{document}             % End of preamble and beginning of text.

\maketitle                   % Produces the title.

This is an example input file.  Comparing it with
the output it generates can show you how to
produce a simple document of your own.

\section{Ordinary Text}      % Produces section heading.  Lower-level
                             % sections are begun with similar
                             % \subsection and \subsubsection commands.

The ends  of words and sentences are marked
  by   spaces. It  doesn't matter how many
spaces    you type; one is as good as 100.  The
end of   a line counts as a space.

One   or more   blank lines denote the  end
of  a paragraph.

Since any number of consecutive spaces are treated
like a single one, the formatting of the input
file makes no difference to
      \LaTeX,                % The \LaTeX command generates the LaTeX logo.
but it makes a difference to you.  When you use
\LaTeX, making your input file as easy to read
as possible will be a great help as you write
your document and when you change it.  This sample
file shows how you can add comments to your own input
file.

Because printing is different from typewriting,
there are a number of things that you have to do
differently when preparing an input file than if
you were just typing the document directly.
Quotation marks like
       ``this''
have to be handled specially, as do quotes within
quotes:
       ``\,`this'            % \, separates the double and single quote.
        is what I just
        wrote, not  `that'\,''.

Dashes come in three sizes: an
       intra-word
dash, a medium dash for number ranges like
       1--2,
and a punctuation
       dash---like
this.

A sentence-ending space should be larger than the
space between words within a sentence.  You
sometimes have to type special commands in
conjunction with punctuation characters to get
this right, as in the following sentence.
       Gnats, gnus, etc.\ all  % `\ ' makes an inter-word space.
       begin with G\@.         % \@ marks end-of-sentence punctuation.
You should check the spaces after periods when
reading your output to make sure you haven't
forgotten any special cases.  Generating an
ellipsis
       \ldots\               % `\ ' is needed after `\ldots' because TeX
                             % ignores spaces after command names like \ldots
                             % made from \ + letters.
                             %
                             % Note how a `%' character causes TeX to ignore
                             % the end of the input line, so these blank lines
                             % do not start a new paragraph.
                             %
with the right spacing around the periods requires
a special command.

\LaTeX\ interprets some common characters as
commands, so you must type special commands to
generate them.  These characters include the
following:
       \$ \& \% \# \{ and \}.

In printing, text is usually emphasized with an
       \emph{italic}
type style.

\begin{em}
   A long segment of text can also be emphasized
   in this way.  Text within such a segment can be
   given \emph{additional} emphasis.
\end{em}

It is sometimes necessary to prevent \LaTeX\ from
breaking a line where it might otherwise do so.
This may be at a space, as between the ``Mr.''\ and
``Jones'' in
       ``Mr.~Jones'',        % ~ produces an unbreakable interword space.
or within a word---especially when the word is a
symbol like
       \mbox{\emph{itemnum}}
that makes little sense when hyphenated across
lines.

Footnotes\footnote{This is an example of a footnote.}
pose no problem.

\LaTeX\ is good at typesetting mathematical formulas
like
       \( x-3y + z = 7 \)
or
       \( a_{1} > x^{2n} + y^{2n} > x' \)
or
       \( \ip{A}{B} = \sum_{i} a_{i} b_{i} \).
The spaces you type in a formula are
ignored.  Remember that a letter like
       $x$                   % $ ... $  and  \( ... \)  are equivalent
is a formula when it denotes a mathematical
symbol, and it should be typed as one.

\section{Displayed Text}

Text is displayed by indenting it from the left
margin.  Quotations are commonly displayed.  There
are short quotations
\begin{quote}
   This is a short quotation.  It consists of a
   single paragraph of text.  See how it is formatted.
\end{quote}
and longer ones.
\begin{quotation}
   This is a longer quotation.  It consists of two
   paragraphs of text, neither of which are
   particularly interesting.

   This is the second paragraph of the quotation.  It
   is just as dull as the first paragraph.
\end{quotation}
Another frequently-displayed structure is a list.
The following is an example of an \emph{itemized}
list.
\begin{itemize}
   \item This is the first item of an itemized list.
         Each item in the list is marked with a ``tick''.
         You don't have to worry about what kind of tick
         mark is used.

   \item This is the second item of the list.  It
         contains another list nested inside it.  The inner
         list is an \emph{enumerated} list.
         \begin{enumerate}
            \item This is the first item of an enumerated
                  list that is nested within the itemized list.

            \item This is the second item of the inner list.
                  \LaTeX\ allows you to nest lists deeper than
                  you really should.
         \end{enumerate}
         This is the rest of the second item of the outer
         list.  It is no more interesting than any other
         part of the item.
   \item This is the third item of the list.
\end{itemize}
You can even display poetry.
\begin{verse}
   There is an environment
    for verse \\             % The \\ command separates lines
   Whose features some poets % within a stanza.
   will curse.

                             % One or more blank lines separate stanzas.

   For instead of making\\
   Them do \emph{all} line breaking, \\
   It allows them to put too many words on a line when they'd rather be
   forced to be terse.
\end{verse}

Mathematical formulas may also be displayed.  A
displayed formula
is
one-line long; multiline
formulas require special formatting instructions.
   \[  \ip{\Gamma}{\psi'} = x'' + y^{2} + z_{i}^{n}\]
Don't start a paragraph with a displayed equation,
nor make one a paragraph by itself.

\end{document}               % End of document.

\section{Bench Huge Section 1401}
% synthetic bench section 1401
% This is a sample LaTeX input file.  (Version of 12 August 2004.)
%
% A '%' character causes TeX to ignore all remaining text on the line,
% and is used for comments like this one.

\documentclass{article}      % Specifies the document class

                             % The preamble begins here.
\title{An Example Document}  % Declares the document's title.
\author{Leslie Lamport}      % Declares the author's name.
\date{January 21, 1994}      % Deleting this command produces today's date.

\newcommand{\ip}[2]{(#1, #2)}
                             % Defines \ip{arg1}{arg2} to mean
                             % (arg1, arg2).

%\newcommand{\ip}[2]{\langle #1 | #2\rangle}
                             % This is an alternative definition of
                             % \ip that is commented out.

\begin{document}             % End of preamble and beginning of text.

\maketitle                   % Produces the title.

This is an example input file.  Comparing it with
the output it generates can show you how to
produce a simple document of your own.

\section{Ordinary Text}      % Produces section heading.  Lower-level
                             % sections are begun with similar
                             % \subsection and \subsubsection commands.

The ends  of words and sentences are marked
  by   spaces. It  doesn't matter how many
spaces    you type; one is as good as 100.  The
end of   a line counts as a space.

One   or more   blank lines denote the  end
of  a paragraph.

Since any number of consecutive spaces are treated
like a single one, the formatting of the input
file makes no difference to
      \LaTeX,                % The \LaTeX command generates the LaTeX logo.
but it makes a difference to you.  When you use
\LaTeX, making your input file as easy to read
as possible will be a great help as you write
your document and when you change it.  This sample
file shows how you can add comments to your own input
file.

Because printing is different from typewriting,
there are a number of things that you have to do
differently when preparing an input file than if
you were just typing the document directly.
Quotation marks like
       ``this''
have to be handled specially, as do quotes within
quotes:
       ``\,`this'            % \, separates the double and single quote.
        is what I just
        wrote, not  `that'\,''.

Dashes come in three sizes: an
       intra-word
dash, a medium dash for number ranges like
       1--2,
and a punctuation
       dash---like
this.

A sentence-ending space should be larger than the
space between words within a sentence.  You
sometimes have to type special commands in
conjunction with punctuation characters to get
this right, as in the following sentence.
       Gnats, gnus, etc.\ all  % `\ ' makes an inter-word space.
       begin with G\@.         % \@ marks end-of-sentence punctuation.
You should check the spaces after periods when
reading your output to make sure you haven't
forgotten any special cases.  Generating an
ellipsis
       \ldots\               % `\ ' is needed after `\ldots' because TeX
                             % ignores spaces after command names like \ldots
                             % made from \ + letters.
                             %
                             % Note how a `%' character causes TeX to ignore
                             % the end of the input line, so these blank lines
                             % do not start a new paragraph.
                             %
with the right spacing around the periods requires
a special command.

\LaTeX\ interprets some common characters as
commands, so you must type special commands to
generate them.  These characters include the
following:
       \$ \& \% \# \{ and \}.

In printing, text is usually emphasized with an
       \emph{italic}
type style.

\begin{em}
   A long segment of text can also be emphasized
   in this way.  Text within such a segment can be
   given \emph{additional} emphasis.
\end{em}

It is sometimes necessary to prevent \LaTeX\ from
breaking a line where it might otherwise do so.
This may be at a space, as between the ``Mr.''\ and
``Jones'' in
       ``Mr.~Jones'',        % ~ produces an unbreakable interword space.
or within a word---especially when the word is a
symbol like
       \mbox{\emph{itemnum}}
that makes little sense when hyphenated across
lines.

Footnotes\footnote{This is an example of a footnote.}
pose no problem.

\LaTeX\ is good at typesetting mathematical formulas
like
       \( x-3y + z = 7 \)
or
       \( a_{1} > x^{2n} + y^{2n} > x' \)
or
       \( \ip{A}{B} = \sum_{i} a_{i} b_{i} \).
The spaces you type in a formula are
ignored.  Remember that a letter like
       $x$                   % $ ... $  and  \( ... \)  are equivalent
is a formula when it denotes a mathematical
symbol, and it should be typed as one.

\section{Displayed Text}

Text is displayed by indenting it from the left
margin.  Quotations are commonly displayed.  There
are short quotations
\begin{quote}
   This is a short quotation.  It consists of a
   single paragraph of text.  See how it is formatted.
\end{quote}
and longer ones.
\begin{quotation}
   This is a longer quotation.  It consists of two
   paragraphs of text, neither of which are
   particularly interesting.

   This is the second paragraph of the quotation.  It
   is just as dull as the first paragraph.
\end{quotation}
Another frequently-displayed structure is a list.
The following is an example of an \emph{itemized}
list.
\begin{itemize}
   \item This is the first item of an itemized list.
         Each item in the list is marked with a ``tick''.
         You don't have to worry about what kind of tick
         mark is used.

   \item This is the second item of the list.  It
         contains another list nested inside it.  The inner
         list is an \emph{enumerated} list.
         \begin{enumerate}
            \item This is the first item of an enumerated
                  list that is nested within the itemized list.

            \item This is the second item of the inner list.
                  \LaTeX\ allows you to nest lists deeper than
                  you really should.
         \end{enumerate}
         This is the rest of the second item of the outer
         list.  It is no more interesting than any other
         part of the item.
   \item This is the third item of the list.
\end{itemize}
You can even display poetry.
\begin{verse}
   There is an environment
    for verse \\             % The \\ command separates lines
   Whose features some poets % within a stanza.
   will curse.

                             % One or more blank lines separate stanzas.

   For instead of making\\
   Them do \emph{all} line breaking, \\
   It allows them to put too many words on a line when they'd rather be
   forced to be terse.
\end{verse}

Mathematical formulas may also be displayed.  A
displayed formula
is
one-line long; multiline
formulas require special formatting instructions.
   \[  \ip{\Gamma}{\psi'} = x'' + y^{2} + z_{i}^{n}\]
Don't start a paragraph with a displayed equation,
nor make one a paragraph by itself.

\end{document}               % End of document.

\section{Bench Huge Section 1402}
% synthetic bench section 1402
% This is a sample LaTeX input file.  (Version of 12 August 2004.)
%
% A '%' character causes TeX to ignore all remaining text on the line,
% and is used for comments like this one.

\documentclass{article}      % Specifies the document class

                             % The preamble begins here.
\title{An Example Document}  % Declares the document's title.
\author{Leslie Lamport}      % Declares the author's name.
\date{January 21, 1994}      % Deleting this command produces today's date.

\newcommand{\ip}[2]{(#1, #2)}
                             % Defines \ip{arg1}{arg2} to mean
                             % (arg1, arg2).

%\newcommand{\ip}[2]{\langle #1 | #2\rangle}
                             % This is an alternative definition of
                             % \ip that is commented out.

\begin{document}             % End of preamble and beginning of text.

\maketitle                   % Produces the title.

This is an example input file.  Comparing it with
the output it generates can show you how to
produce a simple document of your own.

\section{Ordinary Text}      % Produces section heading.  Lower-level
                             % sections are begun with similar
                             % \subsection and \subsubsection commands.

The ends  of words and sentences are marked
  by   spaces. It  doesn't matter how many
spaces    you type; one is as good as 100.  The
end of   a line counts as a space.

One   or more   blank lines denote the  end
of  a paragraph.

Since any number of consecutive spaces are treated
like a single one, the formatting of the input
file makes no difference to
      \LaTeX,                % The \LaTeX command generates the LaTeX logo.
but it makes a difference to you.  When you use
\LaTeX, making your input file as easy to read
as possible will be a great help as you write
your document and when you change it.  This sample
file shows how you can add comments to your own input
file.

Because printing is different from typewriting,
there are a number of things that you have to do
differently when preparing an input file than if
you were just typing the document directly.
Quotation marks like
       ``this''
have to be handled specially, as do quotes within
quotes:
       ``\,`this'            % \, separates the double and single quote.
        is what I just
        wrote, not  `that'\,''.

Dashes come in three sizes: an
       intra-word
dash, a medium dash for number ranges like
       1--2,
and a punctuation
       dash---like
this.

A sentence-ending space should be larger than the
space between words within a sentence.  You
sometimes have to type special commands in
conjunction with punctuation characters to get
this right, as in the following sentence.
       Gnats, gnus, etc.\ all  % `\ ' makes an inter-word space.
       begin with G\@.         % \@ marks end-of-sentence punctuation.
You should check the spaces after periods when
reading your output to make sure you haven't
forgotten any special cases.  Generating an
ellipsis
       \ldots\               % `\ ' is needed after `\ldots' because TeX
                             % ignores spaces after command names like \ldots
                             % made from \ + letters.
                             %
                             % Note how a `%' character causes TeX to ignore
                             % the end of the input line, so these blank lines
                             % do not start a new paragraph.
                             %
with the right spacing around the periods requires
a special command.

\LaTeX\ interprets some common characters as
commands, so you must type special commands to
generate them.  These characters include the
following:
       \$ \& \% \# \{ and \}.

In printing, text is usually emphasized with an
       \emph{italic}
type style.

\begin{em}
   A long segment of text can also be emphasized
   in this way.  Text within such a segment can be
   given \emph{additional} emphasis.
\end{em}

It is sometimes necessary to prevent \LaTeX\ from
breaking a line where it might otherwise do so.
This may be at a space, as between the ``Mr.''\ and
``Jones'' in
       ``Mr.~Jones'',        % ~ produces an unbreakable interword space.
or within a word---especially when the word is a
symbol like
       \mbox{\emph{itemnum}}
that makes little sense when hyphenated across
lines.

Footnotes\footnote{This is an example of a footnote.}
pose no problem.

\LaTeX\ is good at typesetting mathematical formulas
like
       \( x-3y + z = 7 \)
or
       \( a_{1} > x^{2n} + y^{2n} > x' \)
or
       \( \ip{A}{B} = \sum_{i} a_{i} b_{i} \).
The spaces you type in a formula are
ignored.  Remember that a letter like
       $x$                   % $ ... $  and  \( ... \)  are equivalent
is a formula when it denotes a mathematical
symbol, and it should be typed as one.

\section{Displayed Text}

Text is displayed by indenting it from the left
margin.  Quotations are commonly displayed.  There
are short quotations
\begin{quote}
   This is a short quotation.  It consists of a
   single paragraph of text.  See how it is formatted.
\end{quote}
and longer ones.
\begin{quotation}
   This is a longer quotation.  It consists of two
   paragraphs of text, neither of which are
   particularly interesting.

   This is the second paragraph of the quotation.  It
   is just as dull as the first paragraph.
\end{quotation}
Another frequently-displayed structure is a list.
The following is an example of an \emph{itemized}
list.
\begin{itemize}
   \item This is the first item of an itemized list.
         Each item in the list is marked with a ``tick''.
         You don't have to worry about what kind of tick
         mark is used.

   \item This is the second item of the list.  It
         contains another list nested inside it.  The inner
         list is an \emph{enumerated} list.
         \begin{enumerate}
            \item This is the first item of an enumerated
                  list that is nested within the itemized list.

            \item This is the second item of the inner list.
                  \LaTeX\ allows you to nest lists deeper than
                  you really should.
         \end{enumerate}
         This is the rest of the second item of the outer
         list.  It is no more interesting than any other
         part of the item.
   \item This is the third item of the list.
\end{itemize}
You can even display poetry.
\begin{verse}
   There is an environment
    for verse \\             % The \\ command separates lines
   Whose features some poets % within a stanza.
   will curse.

                             % One or more blank lines separate stanzas.

   For instead of making\\
   Them do \emph{all} line breaking, \\
   It allows them to put too many words on a line when they'd rather be
   forced to be terse.
\end{verse}

Mathematical formulas may also be displayed.  A
displayed formula
is
one-line long; multiline
formulas require special formatting instructions.
   \[  \ip{\Gamma}{\psi'} = x'' + y^{2} + z_{i}^{n}\]
Don't start a paragraph with a displayed equation,
nor make one a paragraph by itself.

\end{document}               % End of document.

\section{Bench Huge Section 1403}
% synthetic bench section 1403
% This is a sample LaTeX input file.  (Version of 12 August 2004.)
%
% A '%' character causes TeX to ignore all remaining text on the line,
% and is used for comments like this one.

\documentclass{article}      % Specifies the document class

                             % The preamble begins here.
\title{An Example Document}  % Declares the document's title.
\author{Leslie Lamport}      % Declares the author's name.
\date{January 21, 1994}      % Deleting this command produces today's date.

\newcommand{\ip}[2]{(#1, #2)}
                             % Defines \ip{arg1}{arg2} to mean
                             % (arg1, arg2).

%\newcommand{\ip}[2]{\langle #1 | #2\rangle}
                             % This is an alternative definition of
                             % \ip that is commented out.

\begin{document}             % End of preamble and beginning of text.

\maketitle                   % Produces the title.

This is an example input file.  Comparing it with
the output it generates can show you how to
produce a simple document of your own.

\section{Ordinary Text}      % Produces section heading.  Lower-level
                             % sections are begun with similar
                             % \subsection and \subsubsection commands.

The ends  of words and sentences are marked
  by   spaces. It  doesn't matter how many
spaces    you type; one is as good as 100.  The
end of   a line counts as a space.

One   or more   blank lines denote the  end
of  a paragraph.

Since any number of consecutive spaces are treated
like a single one, the formatting of the input
file makes no difference to
      \LaTeX,                % The \LaTeX command generates the LaTeX logo.
but it makes a difference to you.  When you use
\LaTeX, making your input file as easy to read
as possible will be a great help as you write
your document and when you change it.  This sample
file shows how you can add comments to your own input
file.

Because printing is different from typewriting,
there are a number of things that you have to do
differently when preparing an input file than if
you were just typing the document directly.
Quotation marks like
       ``this''
have to be handled specially, as do quotes within
quotes:
       ``\,`this'            % \, separates the double and single quote.
        is what I just
        wrote, not  `that'\,''.

Dashes come in three sizes: an
       intra-word
dash, a medium dash for number ranges like
       1--2,
and a punctuation
       dash---like
this.

A sentence-ending space should be larger than the
space between words within a sentence.  You
sometimes have to type special commands in
conjunction with punctuation characters to get
this right, as in the following sentence.
       Gnats, gnus, etc.\ all  % `\ ' makes an inter-word space.
       begin with G\@.         % \@ marks end-of-sentence punctuation.
You should check the spaces after periods when
reading your output to make sure you haven't
forgotten any special cases.  Generating an
ellipsis
       \ldots\               % `\ ' is needed after `\ldots' because TeX
                             % ignores spaces after command names like \ldots
                             % made from \ + letters.
                             %
                             % Note how a `%' character causes TeX to ignore
                             % the end of the input line, so these blank lines
                             % do not start a new paragraph.
                             %
with the right spacing around the periods requires
a special command.

\LaTeX\ interprets some common characters as
commands, so you must type special commands to
generate them.  These characters include the
following:
       \$ \& \% \# \{ and \}.

In printing, text is usually emphasized with an
       \emph{italic}
type style.

\begin{em}
   A long segment of text can also be emphasized
   in this way.  Text within such a segment can be
   given \emph{additional} emphasis.
\end{em}

It is sometimes necessary to prevent \LaTeX\ from
breaking a line where it might otherwise do so.
This may be at a space, as between the ``Mr.''\ and
``Jones'' in
       ``Mr.~Jones'',        % ~ produces an unbreakable interword space.
or within a word---especially when the word is a
symbol like
       \mbox{\emph{itemnum}}
that makes little sense when hyphenated across
lines.

Footnotes\footnote{This is an example of a footnote.}
pose no problem.

\LaTeX\ is good at typesetting mathematical formulas
like
       \( x-3y + z = 7 \)
or
       \( a_{1} > x^{2n} + y^{2n} > x' \)
or
       \( \ip{A}{B} = \sum_{i} a_{i} b_{i} \).
The spaces you type in a formula are
ignored.  Remember that a letter like
       $x$                   % $ ... $  and  \( ... \)  are equivalent
is a formula when it denotes a mathematical
symbol, and it should be typed as one.

\section{Displayed Text}

Text is displayed by indenting it from the left
margin.  Quotations are commonly displayed.  There
are short quotations
\begin{quote}
   This is a short quotation.  It consists of a
   single paragraph of text.  See how it is formatted.
\end{quote}
and longer ones.
\begin{quotation}
   This is a longer quotation.  It consists of two
   paragraphs of text, neither of which are
   particularly interesting.

   This is the second paragraph of the quotation.  It
   is just as dull as the first paragraph.
\end{quotation}
Another frequently-displayed structure is a list.
The following is an example of an \emph{itemized}
list.
\begin{itemize}
   \item This is the first item of an itemized list.
         Each item in the list is marked with a ``tick''.
         You don't have to worry about what kind of tick
         mark is used.

   \item This is the second item of the list.  It
         contains another list nested inside it.  The inner
         list is an \emph{enumerated} list.
         \begin{enumerate}
            \item This is the first item of an enumerated
                  list that is nested within the itemized list.

            \item This is the second item of the inner list.
                  \LaTeX\ allows you to nest lists deeper than
                  you really should.
         \end{enumerate}
         This is the rest of the second item of the outer
         list.  It is no more interesting than any other
         part of the item.
   \item This is the third item of the list.
\end{itemize}
You can even display poetry.
\begin{verse}
   There is an environment
    for verse \\             % The \\ command separates lines
   Whose features some poets % within a stanza.
   will curse.

                             % One or more blank lines separate stanzas.

   For instead of making\\
   Them do \emph{all} line breaking, \\
   It allows them to put too many words on a line when they'd rather be
   forced to be terse.
\end{verse}

Mathematical formulas may also be displayed.  A
displayed formula
is
one-line long; multiline
formulas require special formatting instructions.
   \[  \ip{\Gamma}{\psi'} = x'' + y^{2} + z_{i}^{n}\]
Don't start a paragraph with a displayed equation,
nor make one a paragraph by itself.

\end{document}               % End of document.

\section{Bench Huge Section 1404}
% synthetic bench section 1404
% This is a sample LaTeX input file.  (Version of 12 August 2004.)
%
% A '%' character causes TeX to ignore all remaining text on the line,
% and is used for comments like this one.

\documentclass{article}      % Specifies the document class

                             % The preamble begins here.
\title{An Example Document}  % Declares the document's title.
\author{Leslie Lamport}      % Declares the author's name.
\date{January 21, 1994}      % Deleting this command produces today's date.

\newcommand{\ip}[2]{(#1, #2)}
                             % Defines \ip{arg1}{arg2} to mean
                             % (arg1, arg2).

%\newcommand{\ip}[2]{\langle #1 | #2\rangle}
                             % This is an alternative definition of
                             % \ip that is commented out.

\begin{document}             % End of preamble and beginning of text.

\maketitle                   % Produces the title.

This is an example input file.  Comparing it with
the output it generates can show you how to
produce a simple document of your own.

\section{Ordinary Text}      % Produces section heading.  Lower-level
                             % sections are begun with similar
                             % \subsection and \subsubsection commands.

The ends  of words and sentences are marked
  by   spaces. It  doesn't matter how many
spaces    you type; one is as good as 100.  The
end of   a line counts as a space.

One   or more   blank lines denote the  end
of  a paragraph.

Since any number of consecutive spaces are treated
like a single one, the formatting of the input
file makes no difference to
      \LaTeX,                % The \LaTeX command generates the LaTeX logo.
but it makes a difference to you.  When you use
\LaTeX, making your input file as easy to read
as possible will be a great help as you write
your document and when you change it.  This sample
file shows how you can add comments to your own input
file.

Because printing is different from typewriting,
there are a number of things that you have to do
differently when preparing an input file than if
you were just typing the document directly.
Quotation marks like
       ``this''
have to be handled specially, as do quotes within
quotes:
       ``\,`this'            % \, separates the double and single quote.
        is what I just
        wrote, not  `that'\,''.

Dashes come in three sizes: an
       intra-word
dash, a medium dash for number ranges like
       1--2,
and a punctuation
       dash---like
this.

A sentence-ending space should be larger than the
space between words within a sentence.  You
sometimes have to type special commands in
conjunction with punctuation characters to get
this right, as in the following sentence.
       Gnats, gnus, etc.\ all  % `\ ' makes an inter-word space.
       begin with G\@.         % \@ marks end-of-sentence punctuation.
You should check the spaces after periods when
reading your output to make sure you haven't
forgotten any special cases.  Generating an
ellipsis
       \ldots\               % `\ ' is needed after `\ldots' because TeX
                             % ignores spaces after command names like \ldots
                             % made from \ + letters.
                             %
                             % Note how a `%' character causes TeX to ignore
                             % the end of the input line, so these blank lines
                             % do not start a new paragraph.
                             %
with the right spacing around the periods requires
a special command.

\LaTeX\ interprets some common characters as
commands, so you must type special commands to
generate them.  These characters include the
following:
       \$ \& \% \# \{ and \}.

In printing, text is usually emphasized with an
       \emph{italic}
type style.

\begin{em}
   A long segment of text can also be emphasized
   in this way.  Text within such a segment can be
   given \emph{additional} emphasis.
\end{em}

It is sometimes necessary to prevent \LaTeX\ from
breaking a line where it might otherwise do so.
This may be at a space, as between the ``Mr.''\ and
``Jones'' in
       ``Mr.~Jones'',        % ~ produces an unbreakable interword space.
or within a word---especially when the word is a
symbol like
       \mbox{\emph{itemnum}}
that makes little sense when hyphenated across
lines.

Footnotes\footnote{This is an example of a footnote.}
pose no problem.

\LaTeX\ is good at typesetting mathematical formulas
like
       \( x-3y + z = 7 \)
or
       \( a_{1} > x^{2n} + y^{2n} > x' \)
or
       \( \ip{A}{B} = \sum_{i} a_{i} b_{i} \).
The spaces you type in a formula are
ignored.  Remember that a letter like
       $x$                   % $ ... $  and  \( ... \)  are equivalent
is a formula when it denotes a mathematical
symbol, and it should be typed as one.

\section{Displayed Text}

Text is displayed by indenting it from the left
margin.  Quotations are commonly displayed.  There
are short quotations
\begin{quote}
   This is a short quotation.  It consists of a
   single paragraph of text.  See how it is formatted.
\end{quote}
and longer ones.
\begin{quotation}
   This is a longer quotation.  It consists of two
   paragraphs of text, neither of which are
   particularly interesting.

   This is the second paragraph of the quotation.  It
   is just as dull as the first paragraph.
\end{quotation}
Another frequently-displayed structure is a list.
The following is an example of an \emph{itemized}
list.
\begin{itemize}
   \item This is the first item of an itemized list.
         Each item in the list is marked with a ``tick''.
         You don't have to worry about what kind of tick
         mark is used.

   \item This is the second item of the list.  It
         contains another list nested inside it.  The inner
         list is an \emph{enumerated} list.
         \begin{enumerate}
            \item This is the first item of an enumerated
                  list that is nested within the itemized list.

            \item This is the second item of the inner list.
                  \LaTeX\ allows you to nest lists deeper than
                  you really should.
         \end{enumerate}
         This is the rest of the second item of the outer
         list.  It is no more interesting than any other
         part of the item.
   \item This is the third item of the list.
\end{itemize}
You can even display poetry.
\begin{verse}
   There is an environment
    for verse \\             % The \\ command separates lines
   Whose features some poets % within a stanza.
   will curse.

                             % One or more blank lines separate stanzas.

   For instead of making\\
   Them do \emph{all} line breaking, \\
   It allows them to put too many words on a line when they'd rather be
   forced to be terse.
\end{verse}

Mathematical formulas may also be displayed.  A
displayed formula
is
one-line long; multiline
formulas require special formatting instructions.
   \[  \ip{\Gamma}{\psi'} = x'' + y^{2} + z_{i}^{n}\]
Don't start a paragraph with a displayed equation,
nor make one a paragraph by itself.

\end{document}               % End of document.

\section{Bench Huge Section 1405}
% synthetic bench section 1405
% This is a sample LaTeX input file.  (Version of 12 August 2004.)
%
% A '%' character causes TeX to ignore all remaining text on the line,
% and is used for comments like this one.

\documentclass{article}      % Specifies the document class

                             % The preamble begins here.
\title{An Example Document}  % Declares the document's title.
\author{Leslie Lamport}      % Declares the author's name.
\date{January 21, 1994}      % Deleting this command produces today's date.

\newcommand{\ip}[2]{(#1, #2)}
                             % Defines \ip{arg1}{arg2} to mean
                             % (arg1, arg2).

%\newcommand{\ip}[2]{\langle #1 | #2\rangle}
                             % This is an alternative definition of
                             % \ip that is commented out.

\begin{document}             % End of preamble and beginning of text.

\maketitle                   % Produces the title.

This is an example input file.  Comparing it with
the output it generates can show you how to
produce a simple document of your own.

\section{Ordinary Text}      % Produces section heading.  Lower-level
                             % sections are begun with similar
                             % \subsection and \subsubsection commands.

The ends  of words and sentences are marked
  by   spaces. It  doesn't matter how many
spaces    you type; one is as good as 100.  The
end of   a line counts as a space.

One   or more   blank lines denote the  end
of  a paragraph.

Since any number of consecutive spaces are treated
like a single one, the formatting of the input
file makes no difference to
      \LaTeX,                % The \LaTeX command generates the LaTeX logo.
but it makes a difference to you.  When you use
\LaTeX, making your input file as easy to read
as possible will be a great help as you write
your document and when you change it.  This sample
file shows how you can add comments to your own input
file.

Because printing is different from typewriting,
there are a number of things that you have to do
differently when preparing an input file than if
you were just typing the document directly.
Quotation marks like
       ``this''
have to be handled specially, as do quotes within
quotes:
       ``\,`this'            % \, separates the double and single quote.
        is what I just
        wrote, not  `that'\,''.

Dashes come in three sizes: an
       intra-word
dash, a medium dash for number ranges like
       1--2,
and a punctuation
       dash---like
this.

A sentence-ending space should be larger than the
space between words within a sentence.  You
sometimes have to type special commands in
conjunction with punctuation characters to get
this right, as in the following sentence.
       Gnats, gnus, etc.\ all  % `\ ' makes an inter-word space.
       begin with G\@.         % \@ marks end-of-sentence punctuation.
You should check the spaces after periods when
reading your output to make sure you haven't
forgotten any special cases.  Generating an
ellipsis
       \ldots\               % `\ ' is needed after `\ldots' because TeX
                             % ignores spaces after command names like \ldots
                             % made from \ + letters.
                             %
                             % Note how a `%' character causes TeX to ignore
                             % the end of the input line, so these blank lines
                             % do not start a new paragraph.
                             %
with the right spacing around the periods requires
a special command.

\LaTeX\ interprets some common characters as
commands, so you must type special commands to
generate them.  These characters include the
following:
       \$ \& \% \# \{ and \}.

In printing, text is usually emphasized with an
       \emph{italic}
type style.

\begin{em}
   A long segment of text can also be emphasized
   in this way.  Text within such a segment can be
   given \emph{additional} emphasis.
\end{em}

It is sometimes necessary to prevent \LaTeX\ from
breaking a line where it might otherwise do so.
This may be at a space, as between the ``Mr.''\ and
``Jones'' in
       ``Mr.~Jones'',        % ~ produces an unbreakable interword space.
or within a word---especially when the word is a
symbol like
       \mbox{\emph{itemnum}}
that makes little sense when hyphenated across
lines.

Footnotes\footnote{This is an example of a footnote.}
pose no problem.

\LaTeX\ is good at typesetting mathematical formulas
like
       \( x-3y + z = 7 \)
or
       \( a_{1} > x^{2n} + y^{2n} > x' \)
or
       \( \ip{A}{B} = \sum_{i} a_{i} b_{i} \).
The spaces you type in a formula are
ignored.  Remember that a letter like
       $x$                   % $ ... $  and  \( ... \)  are equivalent
is a formula when it denotes a mathematical
symbol, and it should be typed as one.

\section{Displayed Text}

Text is displayed by indenting it from the left
margin.  Quotations are commonly displayed.  There
are short quotations
\begin{quote}
   This is a short quotation.  It consists of a
   single paragraph of text.  See how it is formatted.
\end{quote}
and longer ones.
\begin{quotation}
   This is a longer quotation.  It consists of two
   paragraphs of text, neither of which are
   particularly interesting.

   This is the second paragraph of the quotation.  It
   is just as dull as the first paragraph.
\end{quotation}
Another frequently-displayed structure is a list.
The following is an example of an \emph{itemized}
list.
\begin{itemize}
   \item This is the first item of an itemized list.
         Each item in the list is marked with a ``tick''.
         You don't have to worry about what kind of tick
         mark is used.

   \item This is the second item of the list.  It
         contains another list nested inside it.  The inner
         list is an \emph{enumerated} list.
         \begin{enumerate}
            \item This is the first item of an enumerated
                  list that is nested within the itemized list.

            \item This is the second item of the inner list.
                  \LaTeX\ allows you to nest lists deeper than
                  you really should.
         \end{enumerate}
         This is the rest of the second item of the outer
         list.  It is no more interesting than any other
         part of the item.
   \item This is the third item of the list.
\end{itemize}
You can even display poetry.
\begin{verse}
   There is an environment
    for verse \\             % The \\ command separates lines
   Whose features some poets % within a stanza.
   will curse.

                             % One or more blank lines separate stanzas.

   For instead of making\\
   Them do \emph{all} line breaking, \\
   It allows them to put too many words on a line when they'd rather be
   forced to be terse.
\end{verse}

Mathematical formulas may also be displayed.  A
displayed formula
is
one-line long; multiline
formulas require special formatting instructions.
   \[  \ip{\Gamma}{\psi'} = x'' + y^{2} + z_{i}^{n}\]
Don't start a paragraph with a displayed equation,
nor make one a paragraph by itself.

\end{document}               % End of document.

\section{Bench Huge Section 1406}
% synthetic bench section 1406
% This is a sample LaTeX input file.  (Version of 12 August 2004.)
%
% A '%' character causes TeX to ignore all remaining text on the line,
% and is used for comments like this one.

\documentclass{article}      % Specifies the document class

                             % The preamble begins here.
\title{An Example Document}  % Declares the document's title.
\author{Leslie Lamport}      % Declares the author's name.
\date{January 21, 1994}      % Deleting this command produces today's date.

\newcommand{\ip}[2]{(#1, #2)}
                             % Defines \ip{arg1}{arg2} to mean
                             % (arg1, arg2).

%\newcommand{\ip}[2]{\langle #1 | #2\rangle}
                             % This is an alternative definition of
                             % \ip that is commented out.

\begin{document}             % End of preamble and beginning of text.

\maketitle                   % Produces the title.

This is an example input file.  Comparing it with
the output it generates can show you how to
produce a simple document of your own.

\section{Ordinary Text}      % Produces section heading.  Lower-level
                             % sections are begun with similar
                             % \subsection and \subsubsection commands.

The ends  of words and sentences are marked
  by   spaces. It  doesn't matter how many
spaces    you type; one is as good as 100.  The
end of   a line counts as a space.

One   or more   blank lines denote the  end
of  a paragraph.

Since any number of consecutive spaces are treated
like a single one, the formatting of the input
file makes no difference to
      \LaTeX,                % The \LaTeX command generates the LaTeX logo.
but it makes a difference to you.  When you use
\LaTeX, making your input file as easy to read
as possible will be a great help as you write
your document and when you change it.  This sample
file shows how you can add comments to your own input
file.

Because printing is different from typewriting,
there are a number of things that you have to do
differently when preparing an input file than if
you were just typing the document directly.
Quotation marks like
       ``this''
have to be handled specially, as do quotes within
quotes:
       ``\,`this'            % \, separates the double and single quote.
        is what I just
        wrote, not  `that'\,''.

Dashes come in three sizes: an
       intra-word
dash, a medium dash for number ranges like
       1--2,
and a punctuation
       dash---like
this.

A sentence-ending space should be larger than the
space between words within a sentence.  You
sometimes have to type special commands in
conjunction with punctuation characters to get
this right, as in the following sentence.
       Gnats, gnus, etc.\ all  % `\ ' makes an inter-word space.
       begin with G\@.         % \@ marks end-of-sentence punctuation.
You should check the spaces after periods when
reading your output to make sure you haven't
forgotten any special cases.  Generating an
ellipsis
       \ldots\               % `\ ' is needed after `\ldots' because TeX
                             % ignores spaces after command names like \ldots
                             % made from \ + letters.
                             %
                             % Note how a `%' character causes TeX to ignore
                             % the end of the input line, so these blank lines
                             % do not start a new paragraph.
                             %
with the right spacing around the periods requires
a special command.

\LaTeX\ interprets some common characters as
commands, so you must type special commands to
generate them.  These characters include the
following:
       \$ \& \% \# \{ and \}.

In printing, text is usually emphasized with an
       \emph{italic}
type style.

\begin{em}
   A long segment of text can also be emphasized
   in this way.  Text within such a segment can be
   given \emph{additional} emphasis.
\end{em}

It is sometimes necessary to prevent \LaTeX\ from
breaking a line where it might otherwise do so.
This may be at a space, as between the ``Mr.''\ and
``Jones'' in
       ``Mr.~Jones'',        % ~ produces an unbreakable interword space.
or within a word---especially when the word is a
symbol like
       \mbox{\emph{itemnum}}
that makes little sense when hyphenated across
lines.

Footnotes\footnote{This is an example of a footnote.}
pose no problem.

\LaTeX\ is good at typesetting mathematical formulas
like
       \( x-3y + z = 7 \)
or
       \( a_{1} > x^{2n} + y^{2n} > x' \)
or
       \( \ip{A}{B} = \sum_{i} a_{i} b_{i} \).
The spaces you type in a formula are
ignored.  Remember that a letter like
       $x$                   % $ ... $  and  \( ... \)  are equivalent
is a formula when it denotes a mathematical
symbol, and it should be typed as one.

\section{Displayed Text}

Text is displayed by indenting it from the left
margin.  Quotations are commonly displayed.  There
are short quotations
\begin{quote}
   This is a short quotation.  It consists of a
   single paragraph of text.  See how it is formatted.
\end{quote}
and longer ones.
\begin{quotation}
   This is a longer quotation.  It consists of two
   paragraphs of text, neither of which are
   particularly interesting.

   This is the second paragraph of the quotation.  It
   is just as dull as the first paragraph.
\end{quotation}
Another frequently-displayed structure is a list.
The following is an example of an \emph{itemized}
list.
\begin{itemize}
   \item This is the first item of an itemized list.
         Each item in the list is marked with a ``tick''.
         You don't have to worry about what kind of tick
         mark is used.

   \item This is the second item of the list.  It
         contains another list nested inside it.  The inner
         list is an \emph{enumerated} list.
         \begin{enumerate}
            \item This is the first item of an enumerated
                  list that is nested within the itemized list.

            \item This is the second item of the inner list.
                  \LaTeX\ allows you to nest lists deeper than
                  you really should.
         \end{enumerate}
         This is the rest of the second item of the outer
         list.  It is no more interesting than any other
         part of the item.
   \item This is the third item of the list.
\end{itemize}
You can even display poetry.
\begin{verse}
   There is an environment
    for verse \\             % The \\ command separates lines
   Whose features some poets % within a stanza.
   will curse.

                             % One or more blank lines separate stanzas.

   For instead of making\\
   Them do \emph{all} line breaking, \\
   It allows them to put too many words on a line when they'd rather be
   forced to be terse.
\end{verse}

Mathematical formulas may also be displayed.  A
displayed formula
is
one-line long; multiline
formulas require special formatting instructions.
   \[  \ip{\Gamma}{\psi'} = x'' + y^{2} + z_{i}^{n}\]
Don't start a paragraph with a displayed equation,
nor make one a paragraph by itself.

\end{document}               % End of document.

\section{Bench Huge Section 1407}
% synthetic bench section 1407
% This is a sample LaTeX input file.  (Version of 12 August 2004.)
%
% A '%' character causes TeX to ignore all remaining text on the line,
% and is used for comments like this one.

\documentclass{article}      % Specifies the document class

                             % The preamble begins here.
\title{An Example Document}  % Declares the document's title.
\author{Leslie Lamport}      % Declares the author's name.
\date{January 21, 1994}      % Deleting this command produces today's date.

\newcommand{\ip}[2]{(#1, #2)}
                             % Defines \ip{arg1}{arg2} to mean
                             % (arg1, arg2).

%\newcommand{\ip}[2]{\langle #1 | #2\rangle}
                             % This is an alternative definition of
                             % \ip that is commented out.

\begin{document}             % End of preamble and beginning of text.

\maketitle                   % Produces the title.

This is an example input file.  Comparing it with
the output it generates can show you how to
produce a simple document of your own.

\section{Ordinary Text}      % Produces section heading.  Lower-level
                             % sections are begun with similar
                             % \subsection and \subsubsection commands.

The ends  of words and sentences are marked
  by   spaces. It  doesn't matter how many
spaces    you type; one is as good as 100.  The
end of   a line counts as a space.

One   or more   blank lines denote the  end
of  a paragraph.

Since any number of consecutive spaces are treated
like a single one, the formatting of the input
file makes no difference to
      \LaTeX,                % The \LaTeX command generates the LaTeX logo.
but it makes a difference to you.  When you use
\LaTeX, making your input file as easy to read
as possible will be a great help as you write
your document and when you change it.  This sample
file shows how you can add comments to your own input
file.

Because printing is different from typewriting,
there are a number of things that you have to do
differently when preparing an input file than if
you were just typing the document directly.
Quotation marks like
       ``this''
have to be handled specially, as do quotes within
quotes:
       ``\,`this'            % \, separates the double and single quote.
        is what I just
        wrote, not  `that'\,''.

Dashes come in three sizes: an
       intra-word
dash, a medium dash for number ranges like
       1--2,
and a punctuation
       dash---like
this.

A sentence-ending space should be larger than the
space between words within a sentence.  You
sometimes have to type special commands in
conjunction with punctuation characters to get
this right, as in the following sentence.
       Gnats, gnus, etc.\ all  % `\ ' makes an inter-word space.
       begin with G\@.         % \@ marks end-of-sentence punctuation.
You should check the spaces after periods when
reading your output to make sure you haven't
forgotten any special cases.  Generating an
ellipsis
       \ldots\               % `\ ' is needed after `\ldots' because TeX
                             % ignores spaces after command names like \ldots
                             % made from \ + letters.
                             %
                             % Note how a `%' character causes TeX to ignore
                             % the end of the input line, so these blank lines
                             % do not start a new paragraph.
                             %
with the right spacing around the periods requires
a special command.

\LaTeX\ interprets some common characters as
commands, so you must type special commands to
generate them.  These characters include the
following:
       \$ \& \% \# \{ and \}.

In printing, text is usually emphasized with an
       \emph{italic}
type style.

\begin{em}
   A long segment of text can also be emphasized
   in this way.  Text within such a segment can be
   given \emph{additional} emphasis.
\end{em}

It is sometimes necessary to prevent \LaTeX\ from
breaking a line where it might otherwise do so.
This may be at a space, as between the ``Mr.''\ and
``Jones'' in
       ``Mr.~Jones'',        % ~ produces an unbreakable interword space.
or within a word---especially when the word is a
symbol like
       \mbox{\emph{itemnum}}
that makes little sense when hyphenated across
lines.

Footnotes\footnote{This is an example of a footnote.}
pose no problem.

\LaTeX\ is good at typesetting mathematical formulas
like
       \( x-3y + z = 7 \)
or
       \( a_{1} > x^{2n} + y^{2n} > x' \)
or
       \( \ip{A}{B} = \sum_{i} a_{i} b_{i} \).
The spaces you type in a formula are
ignored.  Remember that a letter like
       $x$                   % $ ... $  and  \( ... \)  are equivalent
is a formula when it denotes a mathematical
symbol, and it should be typed as one.

\section{Displayed Text}

Text is displayed by indenting it from the left
margin.  Quotations are commonly displayed.  There
are short quotations
\begin{quote}
   This is a short quotation.  It consists of a
   single paragraph of text.  See how it is formatted.
\end{quote}
and longer ones.
\begin{quotation}
   This is a longer quotation.  It consists of two
   paragraphs of text, neither of which are
   particularly interesting.

   This is the second paragraph of the quotation.  It
   is just as dull as the first paragraph.
\end{quotation}
Another frequently-displayed structure is a list.
The following is an example of an \emph{itemized}
list.
\begin{itemize}
   \item This is the first item of an itemized list.
         Each item in the list is marked with a ``tick''.
         You don't have to worry about what kind of tick
         mark is used.

   \item This is the second item of the list.  It
         contains another list nested inside it.  The inner
         list is an \emph{enumerated} list.
         \begin{enumerate}
            \item This is the first item of an enumerated
                  list that is nested within the itemized list.

            \item This is the second item of the inner list.
                  \LaTeX\ allows you to nest lists deeper than
                  you really should.
         \end{enumerate}
         This is the rest of the second item of the outer
         list.  It is no more interesting than any other
         part of the item.
   \item This is the third item of the list.
\end{itemize}
You can even display poetry.
\begin{verse}
   There is an environment
    for verse \\             % The \\ command separates lines
   Whose features some poets % within a stanza.
   will curse.

                             % One or more blank lines separate stanzas.

   For instead of making\\
   Them do \emph{all} line breaking, \\
   It allows them to put too many words on a line when they'd rather be
   forced to be terse.
\end{verse}

Mathematical formulas may also be displayed.  A
displayed formula
is
one-line long; multiline
formulas require special formatting instructions.
   \[  \ip{\Gamma}{\psi'} = x'' + y^{2} + z_{i}^{n}\]
Don't start a paragraph with a displayed equation,
nor make one a paragraph by itself.

\end{document}               % End of document.

\section{Bench Huge Section 1408}
% synthetic bench section 1408
% This is a sample LaTeX input file.  (Version of 12 August 2004.)
%
% A '%' character causes TeX to ignore all remaining text on the line,
% and is used for comments like this one.

\documentclass{article}      % Specifies the document class

                             % The preamble begins here.
\title{An Example Document}  % Declares the document's title.
\author{Leslie Lamport}      % Declares the author's name.
\date{January 21, 1994}      % Deleting this command produces today's date.

\newcommand{\ip}[2]{(#1, #2)}
                             % Defines \ip{arg1}{arg2} to mean
                             % (arg1, arg2).

%\newcommand{\ip}[2]{\langle #1 | #2\rangle}
                             % This is an alternative definition of
                             % \ip that is commented out.

\begin{document}             % End of preamble and beginning of text.

\maketitle                   % Produces the title.

This is an example input file.  Comparing it with
the output it generates can show you how to
produce a simple document of your own.

\section{Ordinary Text}      % Produces section heading.  Lower-level
                             % sections are begun with similar
                             % \subsection and \subsubsection commands.

The ends  of words and sentences are marked
  by   spaces. It  doesn't matter how many
spaces    you type; one is as good as 100.  The
end of   a line counts as a space.

One   or more   blank lines denote the  end
of  a paragraph.

Since any number of consecutive spaces are treated
like a single one, the formatting of the input
file makes no difference to
      \LaTeX,                % The \LaTeX command generates the LaTeX logo.
but it makes a difference to you.  When you use
\LaTeX, making your input file as easy to read
as possible will be a great help as you write
your document and when you change it.  This sample
file shows how you can add comments to your own input
file.

Because printing is different from typewriting,
there are a number of things that you have to do
differently when preparing an input file than if
you were just typing the document directly.
Quotation marks like
       ``this''
have to be handled specially, as do quotes within
quotes:
       ``\,`this'            % \, separates the double and single quote.
        is what I just
        wrote, not  `that'\,''.

Dashes come in three sizes: an
       intra-word
dash, a medium dash for number ranges like
       1--2,
and a punctuation
       dash---like
this.

A sentence-ending space should be larger than the
space between words within a sentence.  You
sometimes have to type special commands in
conjunction with punctuation characters to get
this right, as in the following sentence.
       Gnats, gnus, etc.\ all  % `\ ' makes an inter-word space.
       begin with G\@.         % \@ marks end-of-sentence punctuation.
You should check the spaces after periods when
reading your output to make sure you haven't
forgotten any special cases.  Generating an
ellipsis
       \ldots\               % `\ ' is needed after `\ldots' because TeX
                             % ignores spaces after command names like \ldots
                             % made from \ + letters.
                             %
                             % Note how a `%' character causes TeX to ignore
                             % the end of the input line, so these blank lines
                             % do not start a new paragraph.
                             %
with the right spacing around the periods requires
a special command.

\LaTeX\ interprets some common characters as
commands, so you must type special commands to
generate them.  These characters include the
following:
       \$ \& \% \# \{ and \}.

In printing, text is usually emphasized with an
       \emph{italic}
type style.

\begin{em}
   A long segment of text can also be emphasized
   in this way.  Text within such a segment can be
   given \emph{additional} emphasis.
\end{em}

It is sometimes necessary to prevent \LaTeX\ from
breaking a line where it might otherwise do so.
This may be at a space, as between the ``Mr.''\ and
``Jones'' in
       ``Mr.~Jones'',        % ~ produces an unbreakable interword space.
or within a word---especially when the word is a
symbol like
       \mbox{\emph{itemnum}}
that makes little sense when hyphenated across
lines.

Footnotes\footnote{This is an example of a footnote.}
pose no problem.

\LaTeX\ is good at typesetting mathematical formulas
like
       \( x-3y + z = 7 \)
or
       \( a_{1} > x^{2n} + y^{2n} > x' \)
or
       \( \ip{A}{B} = \sum_{i} a_{i} b_{i} \).
The spaces you type in a formula are
ignored.  Remember that a letter like
       $x$                   % $ ... $  and  \( ... \)  are equivalent
is a formula when it denotes a mathematical
symbol, and it should be typed as one.

\section{Displayed Text}

Text is displayed by indenting it from the left
margin.  Quotations are commonly displayed.  There
are short quotations
\begin{quote}
   This is a short quotation.  It consists of a
   single paragraph of text.  See how it is formatted.
\end{quote}
and longer ones.
\begin{quotation}
   This is a longer quotation.  It consists of two
   paragraphs of text, neither of which are
   particularly interesting.

   This is the second paragraph of the quotation.  It
   is just as dull as the first paragraph.
\end{quotation}
Another frequently-displayed structure is a list.
The following is an example of an \emph{itemized}
list.
\begin{itemize}
   \item This is the first item of an itemized list.
         Each item in the list is marked with a ``tick''.
         You don't have to worry about what kind of tick
         mark is used.

   \item This is the second item of the list.  It
         contains another list nested inside it.  The inner
         list is an \emph{enumerated} list.
         \begin{enumerate}
            \item This is the first item of an enumerated
                  list that is nested within the itemized list.

            \item This is the second item of the inner list.
                  \LaTeX\ allows you to nest lists deeper than
                  you really should.
         \end{enumerate}
         This is the rest of the second item of the outer
         list.  It is no more interesting than any other
         part of the item.
   \item This is the third item of the list.
\end{itemize}
You can even display poetry.
\begin{verse}
   There is an environment
    for verse \\             % The \\ command separates lines
   Whose features some poets % within a stanza.
   will curse.

                             % One or more blank lines separate stanzas.

   For instead of making\\
   Them do \emph{all} line breaking, \\
   It allows them to put too many words on a line when they'd rather be
   forced to be terse.
\end{verse}

Mathematical formulas may also be displayed.  A
displayed formula
is
one-line long; multiline
formulas require special formatting instructions.
   \[  \ip{\Gamma}{\psi'} = x'' + y^{2} + z_{i}^{n}\]
Don't start a paragraph with a displayed equation,
nor make one a paragraph by itself.

\end{document}               % End of document.

\section{Bench Huge Section 1409}
% synthetic bench section 1409
% This is a sample LaTeX input file.  (Version of 12 August 2004.)
%
% A '%' character causes TeX to ignore all remaining text on the line,
% and is used for comments like this one.

\documentclass{article}      % Specifies the document class

                             % The preamble begins here.
\title{An Example Document}  % Declares the document's title.
\author{Leslie Lamport}      % Declares the author's name.
\date{January 21, 1994}      % Deleting this command produces today's date.

\newcommand{\ip}[2]{(#1, #2)}
                             % Defines \ip{arg1}{arg2} to mean
                             % (arg1, arg2).

%\newcommand{\ip}[2]{\langle #1 | #2\rangle}
                             % This is an alternative definition of
                             % \ip that is commented out.

\begin{document}             % End of preamble and beginning of text.

\maketitle                   % Produces the title.

This is an example input file.  Comparing it with
the output it generates can show you how to
produce a simple document of your own.

\section{Ordinary Text}      % Produces section heading.  Lower-level
                             % sections are begun with similar
                             % \subsection and \subsubsection commands.

The ends  of words and sentences are marked
  by   spaces. It  doesn't matter how many
spaces    you type; one is as good as 100.  The
end of   a line counts as a space.

One   or more   blank lines denote the  end
of  a paragraph.

Since any number of consecutive spaces are treated
like a single one, the formatting of the input
file makes no difference to
      \LaTeX,                % The \LaTeX command generates the LaTeX logo.
but it makes a difference to you.  When you use
\LaTeX, making your input file as easy to read
as possible will be a great help as you write
your document and when you change it.  This sample
file shows how you can add comments to your own input
file.

Because printing is different from typewriting,
there are a number of things that you have to do
differently when preparing an input file than if
you were just typing the document directly.
Quotation marks like
       ``this''
have to be handled specially, as do quotes within
quotes:
       ``\,`this'            % \, separates the double and single quote.
        is what I just
        wrote, not  `that'\,''.

Dashes come in three sizes: an
       intra-word
dash, a medium dash for number ranges like
       1--2,
and a punctuation
       dash---like
this.

A sentence-ending space should be larger than the
space between words within a sentence.  You
sometimes have to type special commands in
conjunction with punctuation characters to get
this right, as in the following sentence.
       Gnats, gnus, etc.\ all  % `\ ' makes an inter-word space.
       begin with G\@.         % \@ marks end-of-sentence punctuation.
You should check the spaces after periods when
reading your output to make sure you haven't
forgotten any special cases.  Generating an
ellipsis
       \ldots\               % `\ ' is needed after `\ldots' because TeX
                             % ignores spaces after command names like \ldots
                             % made from \ + letters.
                             %
                             % Note how a `%' character causes TeX to ignore
                             % the end of the input line, so these blank lines
                             % do not start a new paragraph.
                             %
with the right spacing around the periods requires
a special command.

\LaTeX\ interprets some common characters as
commands, so you must type special commands to
generate them.  These characters include the
following:
       \$ \& \% \# \{ and \}.

In printing, text is usually emphasized with an
       \emph{italic}
type style.

\begin{em}
   A long segment of text can also be emphasized
   in this way.  Text within such a segment can be
   given \emph{additional} emphasis.
\end{em}

It is sometimes necessary to prevent \LaTeX\ from
breaking a line where it might otherwise do so.
This may be at a space, as between the ``Mr.''\ and
``Jones'' in
       ``Mr.~Jones'',        % ~ produces an unbreakable interword space.
or within a word---especially when the word is a
symbol like
       \mbox{\emph{itemnum}}
that makes little sense when hyphenated across
lines.

Footnotes\footnote{This is an example of a footnote.}
pose no problem.

\LaTeX\ is good at typesetting mathematical formulas
like
       \( x-3y + z = 7 \)
or
       \( a_{1} > x^{2n} + y^{2n} > x' \)
or
       \( \ip{A}{B} = \sum_{i} a_{i} b_{i} \).
The spaces you type in a formula are
ignored.  Remember that a letter like
       $x$                   % $ ... $  and  \( ... \)  are equivalent
is a formula when it denotes a mathematical
symbol, and it should be typed as one.

\section{Displayed Text}

Text is displayed by indenting it from the left
margin.  Quotations are commonly displayed.  There
are short quotations
\begin{quote}
   This is a short quotation.  It consists of a
   single paragraph of text.  See how it is formatted.
\end{quote}
and longer ones.
\begin{quotation}
   This is a longer quotation.  It consists of two
   paragraphs of text, neither of which are
   particularly interesting.

   This is the second paragraph of the quotation.  It
   is just as dull as the first paragraph.
\end{quotation}
Another frequently-displayed structure is a list.
The following is an example of an \emph{itemized}
list.
\begin{itemize}
   \item This is the first item of an itemized list.
         Each item in the list is marked with a ``tick''.
         You don't have to worry about what kind of tick
         mark is used.

   \item This is the second item of the list.  It
         contains another list nested inside it.  The inner
         list is an \emph{enumerated} list.
         \begin{enumerate}
            \item This is the first item of an enumerated
                  list that is nested within the itemized list.

            \item This is the second item of the inner list.
                  \LaTeX\ allows you to nest lists deeper than
                  you really should.
         \end{enumerate}
         This is the rest of the second item of the outer
         list.  It is no more interesting than any other
         part of the item.
   \item This is the third item of the list.
\end{itemize}
You can even display poetry.
\begin{verse}
   There is an environment
    for verse \\             % The \\ command separates lines
   Whose features some poets % within a stanza.
   will curse.

                             % One or more blank lines separate stanzas.

   For instead of making\\
   Them do \emph{all} line breaking, \\
   It allows them to put too many words on a line when they'd rather be
   forced to be terse.
\end{verse}

Mathematical formulas may also be displayed.  A
displayed formula
is
one-line long; multiline
formulas require special formatting instructions.
   \[  \ip{\Gamma}{\psi'} = x'' + y^{2} + z_{i}^{n}\]
Don't start a paragraph with a displayed equation,
nor make one a paragraph by itself.

\end{document}               % End of document.

\section{Bench Huge Section 1410}
% synthetic bench section 1410
% This is a sample LaTeX input file.  (Version of 12 August 2004.)
%
% A '%' character causes TeX to ignore all remaining text on the line,
% and is used for comments like this one.

\documentclass{article}      % Specifies the document class

                             % The preamble begins here.
\title{An Example Document}  % Declares the document's title.
\author{Leslie Lamport}      % Declares the author's name.
\date{January 21, 1994}      % Deleting this command produces today's date.

\newcommand{\ip}[2]{(#1, #2)}
                             % Defines \ip{arg1}{arg2} to mean
                             % (arg1, arg2).

%\newcommand{\ip}[2]{\langle #1 | #2\rangle}
                             % This is an alternative definition of
                             % \ip that is commented out.

\begin{document}             % End of preamble and beginning of text.

\maketitle                   % Produces the title.

This is an example input file.  Comparing it with
the output it generates can show you how to
produce a simple document of your own.

\section{Ordinary Text}      % Produces section heading.  Lower-level
                             % sections are begun with similar
                             % \subsection and \subsubsection commands.

The ends  of words and sentences are marked
  by   spaces. It  doesn't matter how many
spaces    you type; one is as good as 100.  The
end of   a line counts as a space.

One   or more   blank lines denote the  end
of  a paragraph.

Since any number of consecutive spaces are treated
like a single one, the formatting of the input
file makes no difference to
      \LaTeX,                % The \LaTeX command generates the LaTeX logo.
but it makes a difference to you.  When you use
\LaTeX, making your input file as easy to read
as possible will be a great help as you write
your document and when you change it.  This sample
file shows how you can add comments to your own input
file.

Because printing is different from typewriting,
there are a number of things that you have to do
differently when preparing an input file than if
you were just typing the document directly.
Quotation marks like
       ``this''
have to be handled specially, as do quotes within
quotes:
       ``\,`this'            % \, separates the double and single quote.
        is what I just
        wrote, not  `that'\,''.

Dashes come in three sizes: an
       intra-word
dash, a medium dash for number ranges like
       1--2,
and a punctuation
       dash---like
this.

A sentence-ending space should be larger than the
space between words within a sentence.  You
sometimes have to type special commands in
conjunction with punctuation characters to get
this right, as in the following sentence.
       Gnats, gnus, etc.\ all  % `\ ' makes an inter-word space.
       begin with G\@.         % \@ marks end-of-sentence punctuation.
You should check the spaces after periods when
reading your output to make sure you haven't
forgotten any special cases.  Generating an
ellipsis
       \ldots\               % `\ ' is needed after `\ldots' because TeX
                             % ignores spaces after command names like \ldots
                             % made from \ + letters.
                             %
                             % Note how a `%' character causes TeX to ignore
                             % the end of the input line, so these blank lines
                             % do not start a new paragraph.
                             %
with the right spacing around the periods requires
a special command.

\LaTeX\ interprets some common characters as
commands, so you must type special commands to
generate them.  These characters include the
following:
       \$ \& \% \# \{ and \}.

In printing, text is usually emphasized with an
       \emph{italic}
type style.

\begin{em}
   A long segment of text can also be emphasized
   in this way.  Text within such a segment can be
   given \emph{additional} emphasis.
\end{em}

It is sometimes necessary to prevent \LaTeX\ from
breaking a line where it might otherwise do so.
This may be at a space, as between the ``Mr.''\ and
``Jones'' in
       ``Mr.~Jones'',        % ~ produces an unbreakable interword space.
or within a word---especially when the word is a
symbol like
       \mbox{\emph{itemnum}}
that makes little sense when hyphenated across
lines.

Footnotes\footnote{This is an example of a footnote.}
pose no problem.

\LaTeX\ is good at typesetting mathematical formulas
like
       \( x-3y + z = 7 \)
or
       \( a_{1} > x^{2n} + y^{2n} > x' \)
or
       \( \ip{A}{B} = \sum_{i} a_{i} b_{i} \).
The spaces you type in a formula are
ignored.  Remember that a letter like
       $x$                   % $ ... $  and  \( ... \)  are equivalent
is a formula when it denotes a mathematical
symbol, and it should be typed as one.

\section{Displayed Text}

Text is displayed by indenting it from the left
margin.  Quotations are commonly displayed.  There
are short quotations
\begin{quote}
   This is a short quotation.  It consists of a
   single paragraph of text.  See how it is formatted.
\end{quote}
and longer ones.
\begin{quotation}
   This is a longer quotation.  It consists of two
   paragraphs of text, neither of which are
   particularly interesting.

   This is the second paragraph of the quotation.  It
   is just as dull as the first paragraph.
\end{quotation}
Another frequently-displayed structure is a list.
The following is an example of an \emph{itemized}
list.
\begin{itemize}
   \item This is the first item of an itemized list.
         Each item in the list is marked with a ``tick''.
         You don't have to worry about what kind of tick
         mark is used.

   \item This is the second item of the list.  It
         contains another list nested inside it.  The inner
         list is an \emph{enumerated} list.
         \begin{enumerate}
            \item This is the first item of an enumerated
                  list that is nested within the itemized list.

            \item This is the second item of the inner list.
                  \LaTeX\ allows you to nest lists deeper than
                  you really should.
         \end{enumerate}
         This is the rest of the second item of the outer
         list.  It is no more interesting than any other
         part of the item.
   \item This is the third item of the list.
\end{itemize}
You can even display poetry.
\begin{verse}
   There is an environment
    for verse \\             % The \\ command separates lines
   Whose features some poets % within a stanza.
   will curse.

                             % One or more blank lines separate stanzas.

   For instead of making\\
   Them do \emph{all} line breaking, \\
   It allows them to put too many words on a line when they'd rather be
   forced to be terse.
\end{verse}

Mathematical formulas may also be displayed.  A
displayed formula
is
one-line long; multiline
formulas require special formatting instructions.
   \[  \ip{\Gamma}{\psi'} = x'' + y^{2} + z_{i}^{n}\]
Don't start a paragraph with a displayed equation,
nor make one a paragraph by itself.

\end{document}               % End of document.

\section{Bench Huge Section 1411}
% synthetic bench section 1411
% This is a sample LaTeX input file.  (Version of 12 August 2004.)
%
% A '%' character causes TeX to ignore all remaining text on the line,
% and is used for comments like this one.

\documentclass{article}      % Specifies the document class

                             % The preamble begins here.
\title{An Example Document}  % Declares the document's title.
\author{Leslie Lamport}      % Declares the author's name.
\date{January 21, 1994}      % Deleting this command produces today's date.

\newcommand{\ip}[2]{(#1, #2)}
                             % Defines \ip{arg1}{arg2} to mean
                             % (arg1, arg2).

%\newcommand{\ip}[2]{\langle #1 | #2\rangle}
                             % This is an alternative definition of
                             % \ip that is commented out.

\begin{document}             % End of preamble and beginning of text.

\maketitle                   % Produces the title.

This is an example input file.  Comparing it with
the output it generates can show you how to
produce a simple document of your own.

\section{Ordinary Text}      % Produces section heading.  Lower-level
                             % sections are begun with similar
                             % \subsection and \subsubsection commands.

The ends  of words and sentences are marked
  by   spaces. It  doesn't matter how many
spaces    you type; one is as good as 100.  The
end of   a line counts as a space.

One   or more   blank lines denote the  end
of  a paragraph.

Since any number of consecutive spaces are treated
like a single one, the formatting of the input
file makes no difference to
      \LaTeX,                % The \LaTeX command generates the LaTeX logo.
but it makes a difference to you.  When you use
\LaTeX, making your input file as easy to read
as possible will be a great help as you write
your document and when you change it.  This sample
file shows how you can add comments to your own input
file.

Because printing is different from typewriting,
there are a number of things that you have to do
differently when preparing an input file than if
you were just typing the document directly.
Quotation marks like
       ``this''
have to be handled specially, as do quotes within
quotes:
       ``\,`this'            % \, separates the double and single quote.
        is what I just
        wrote, not  `that'\,''.

Dashes come in three sizes: an
       intra-word
dash, a medium dash for number ranges like
       1--2,
and a punctuation
       dash---like
this.

A sentence-ending space should be larger than the
space between words within a sentence.  You
sometimes have to type special commands in
conjunction with punctuation characters to get
this right, as in the following sentence.
       Gnats, gnus, etc.\ all  % `\ ' makes an inter-word space.
       begin with G\@.         % \@ marks end-of-sentence punctuation.
You should check the spaces after periods when
reading your output to make sure you haven't
forgotten any special cases.  Generating an
ellipsis
       \ldots\               % `\ ' is needed after `\ldots' because TeX
                             % ignores spaces after command names like \ldots
                             % made from \ + letters.
                             %
                             % Note how a `%' character causes TeX to ignore
                             % the end of the input line, so these blank lines
                             % do not start a new paragraph.
                             %
with the right spacing around the periods requires
a special command.

\LaTeX\ interprets some common characters as
commands, so you must type special commands to
generate them.  These characters include the
following:
       \$ \& \% \# \{ and \}.

In printing, text is usually emphasized with an
       \emph{italic}
type style.

\begin{em}
   A long segment of text can also be emphasized
   in this way.  Text within such a segment can be
   given \emph{additional} emphasis.
\end{em}

It is sometimes necessary to prevent \LaTeX\ from
breaking a line where it might otherwise do so.
This may be at a space, as between the ``Mr.''\ and
``Jones'' in
       ``Mr.~Jones'',        % ~ produces an unbreakable interword space.
or within a word---especially when the word is a
symbol like
       \mbox{\emph{itemnum}}
that makes little sense when hyphenated across
lines.

Footnotes\footnote{This is an example of a footnote.}
pose no problem.

\LaTeX\ is good at typesetting mathematical formulas
like
       \( x-3y + z = 7 \)
or
       \( a_{1} > x^{2n} + y^{2n} > x' \)
or
       \( \ip{A}{B} = \sum_{i} a_{i} b_{i} \).
The spaces you type in a formula are
ignored.  Remember that a letter like
       $x$                   % $ ... $  and  \( ... \)  are equivalent
is a formula when it denotes a mathematical
symbol, and it should be typed as one.

\section{Displayed Text}

Text is displayed by indenting it from the left
margin.  Quotations are commonly displayed.  There
are short quotations
\begin{quote}
   This is a short quotation.  It consists of a
   single paragraph of text.  See how it is formatted.
\end{quote}
and longer ones.
\begin{quotation}
   This is a longer quotation.  It consists of two
   paragraphs of text, neither of which are
   particularly interesting.

   This is the second paragraph of the quotation.  It
   is just as dull as the first paragraph.
\end{quotation}
Another frequently-displayed structure is a list.
The following is an example of an \emph{itemized}
list.
\begin{itemize}
   \item This is the first item of an itemized list.
         Each item in the list is marked with a ``tick''.
         You don't have to worry about what kind of tick
         mark is used.

   \item This is the second item of the list.  It
         contains another list nested inside it.  The inner
         list is an \emph{enumerated} list.
         \begin{enumerate}
            \item This is the first item of an enumerated
                  list that is nested within the itemized list.

            \item This is the second item of the inner list.
                  \LaTeX\ allows you to nest lists deeper than
                  you really should.
         \end{enumerate}
         This is the rest of the second item of the outer
         list.  It is no more interesting than any other
         part of the item.
   \item This is the third item of the list.
\end{itemize}
You can even display poetry.
\begin{verse}
   There is an environment
    for verse \\             % The \\ command separates lines
   Whose features some poets % within a stanza.
   will curse.

                             % One or more blank lines separate stanzas.

   For instead of making\\
   Them do \emph{all} line breaking, \\
   It allows them to put too many words on a line when they'd rather be
   forced to be terse.
\end{verse}

Mathematical formulas may also be displayed.  A
displayed formula
is
one-line long; multiline
formulas require special formatting instructions.
   \[  \ip{\Gamma}{\psi'} = x'' + y^{2} + z_{i}^{n}\]
Don't start a paragraph with a displayed equation,
nor make one a paragraph by itself.

\end{document}               % End of document.

\section{Bench Huge Section 1412}
% synthetic bench section 1412
% This is a sample LaTeX input file.  (Version of 12 August 2004.)
%
% A '%' character causes TeX to ignore all remaining text on the line,
% and is used for comments like this one.

\documentclass{article}      % Specifies the document class

                             % The preamble begins here.
\title{An Example Document}  % Declares the document's title.
\author{Leslie Lamport}      % Declares the author's name.
\date{January 21, 1994}      % Deleting this command produces today's date.

\newcommand{\ip}[2]{(#1, #2)}
                             % Defines \ip{arg1}{arg2} to mean
                             % (arg1, arg2).

%\newcommand{\ip}[2]{\langle #1 | #2\rangle}
                             % This is an alternative definition of
                             % \ip that is commented out.

\begin{document}             % End of preamble and beginning of text.

\maketitle                   % Produces the title.

This is an example input file.  Comparing it with
the output it generates can show you how to
produce a simple document of your own.

\section{Ordinary Text}      % Produces section heading.  Lower-level
                             % sections are begun with similar
                             % \subsection and \subsubsection commands.

The ends  of words and sentences are marked
  by   spaces. It  doesn't matter how many
spaces    you type; one is as good as 100.  The
end of   a line counts as a space.

One   or more   blank lines denote the  end
of  a paragraph.

Since any number of consecutive spaces are treated
like a single one, the formatting of the input
file makes no difference to
      \LaTeX,                % The \LaTeX command generates the LaTeX logo.
but it makes a difference to you.  When you use
\LaTeX, making your input file as easy to read
as possible will be a great help as you write
your document and when you change it.  This sample
file shows how you can add comments to your own input
file.

Because printing is different from typewriting,
there are a number of things that you have to do
differently when preparing an input file than if
you were just typing the document directly.
Quotation marks like
       ``this''
have to be handled specially, as do quotes within
quotes:
       ``\,`this'            % \, separates the double and single quote.
        is what I just
        wrote, not  `that'\,''.

Dashes come in three sizes: an
       intra-word
dash, a medium dash for number ranges like
       1--2,
and a punctuation
       dash---like
this.

A sentence-ending space should be larger than the
space between words within a sentence.  You
sometimes have to type special commands in
conjunction with punctuation characters to get
this right, as in the following sentence.
       Gnats, gnus, etc.\ all  % `\ ' makes an inter-word space.
       begin with G\@.         % \@ marks end-of-sentence punctuation.
You should check the spaces after periods when
reading your output to make sure you haven't
forgotten any special cases.  Generating an
ellipsis
       \ldots\               % `\ ' is needed after `\ldots' because TeX
                             % ignores spaces after command names like \ldots
                             % made from \ + letters.
                             %
                             % Note how a `%' character causes TeX to ignore
                             % the end of the input line, so these blank lines
                             % do not start a new paragraph.
                             %
with the right spacing around the periods requires
a special command.

\LaTeX\ interprets some common characters as
commands, so you must type special commands to
generate them.  These characters include the
following:
       \$ \& \% \# \{ and \}.

In printing, text is usually emphasized with an
       \emph{italic}
type style.

\begin{em}
   A long segment of text can also be emphasized
   in this way.  Text within such a segment can be
   given \emph{additional} emphasis.
\end{em}

It is sometimes necessary to prevent \LaTeX\ from
breaking a line where it might otherwise do so.
This may be at a space, as between the ``Mr.''\ and
``Jones'' in
       ``Mr.~Jones'',        % ~ produces an unbreakable interword space.
or within a word---especially when the word is a
symbol like
       \mbox{\emph{itemnum}}
that makes little sense when hyphenated across
lines.

Footnotes\footnote{This is an example of a footnote.}
pose no problem.

\LaTeX\ is good at typesetting mathematical formulas
like
       \( x-3y + z = 7 \)
or
       \( a_{1} > x^{2n} + y^{2n} > x' \)
or
       \( \ip{A}{B} = \sum_{i} a_{i} b_{i} \).
The spaces you type in a formula are
ignored.  Remember that a letter like
       $x$                   % $ ... $  and  \( ... \)  are equivalent
is a formula when it denotes a mathematical
symbol, and it should be typed as one.

\section{Displayed Text}

Text is displayed by indenting it from the left
margin.  Quotations are commonly displayed.  There
are short quotations
\begin{quote}
   This is a short quotation.  It consists of a
   single paragraph of text.  See how it is formatted.
\end{quote}
and longer ones.
\begin{quotation}
   This is a longer quotation.  It consists of two
   paragraphs of text, neither of which are
   particularly interesting.

   This is the second paragraph of the quotation.  It
   is just as dull as the first paragraph.
\end{quotation}
Another frequently-displayed structure is a list.
The following is an example of an \emph{itemized}
list.
\begin{itemize}
   \item This is the first item of an itemized list.
         Each item in the list is marked with a ``tick''.
         You don't have to worry about what kind of tick
         mark is used.

   \item This is the second item of the list.  It
         contains another list nested inside it.  The inner
         list is an \emph{enumerated} list.
         \begin{enumerate}
            \item This is the first item of an enumerated
                  list that is nested within the itemized list.

            \item This is the second item of the inner list.
                  \LaTeX\ allows you to nest lists deeper than
                  you really should.
         \end{enumerate}
         This is the rest of the second item of the outer
         list.  It is no more interesting than any other
         part of the item.
   \item This is the third item of the list.
\end{itemize}
You can even display poetry.
\begin{verse}
   There is an environment
    for verse \\             % The \\ command separates lines
   Whose features some poets % within a stanza.
   will curse.

                             % One or more blank lines separate stanzas.

   For instead of making\\
   Them do \emph{all} line breaking, \\
   It allows them to put too many words on a line when they'd rather be
   forced to be terse.
\end{verse}

Mathematical formulas may also be displayed.  A
displayed formula
is
one-line long; multiline
formulas require special formatting instructions.
   \[  \ip{\Gamma}{\psi'} = x'' + y^{2} + z_{i}^{n}\]
Don't start a paragraph with a displayed equation,
nor make one a paragraph by itself.

\end{document}               % End of document.

\section{Bench Huge Section 1413}
% synthetic bench section 1413
% This is a sample LaTeX input file.  (Version of 12 August 2004.)
%
% A '%' character causes TeX to ignore all remaining text on the line,
% and is used for comments like this one.

\documentclass{article}      % Specifies the document class

                             % The preamble begins here.
\title{An Example Document}  % Declares the document's title.
\author{Leslie Lamport}      % Declares the author's name.
\date{January 21, 1994}      % Deleting this command produces today's date.

\newcommand{\ip}[2]{(#1, #2)}
                             % Defines \ip{arg1}{arg2} to mean
                             % (arg1, arg2).

%\newcommand{\ip}[2]{\langle #1 | #2\rangle}
                             % This is an alternative definition of
                             % \ip that is commented out.

\begin{document}             % End of preamble and beginning of text.

\maketitle                   % Produces the title.

This is an example input file.  Comparing it with
the output it generates can show you how to
produce a simple document of your own.

\section{Ordinary Text}      % Produces section heading.  Lower-level
                             % sections are begun with similar
                             % \subsection and \subsubsection commands.

The ends  of words and sentences are marked
  by   spaces. It  doesn't matter how many
spaces    you type; one is as good as 100.  The
end of   a line counts as a space.

One   or more   blank lines denote the  end
of  a paragraph.

Since any number of consecutive spaces are treated
like a single one, the formatting of the input
file makes no difference to
      \LaTeX,                % The \LaTeX command generates the LaTeX logo.
but it makes a difference to you.  When you use
\LaTeX, making your input file as easy to read
as possible will be a great help as you write
your document and when you change it.  This sample
file shows how you can add comments to your own input
file.

Because printing is different from typewriting,
there are a number of things that you have to do
differently when preparing an input file than if
you were just typing the document directly.
Quotation marks like
       ``this''
have to be handled specially, as do quotes within
quotes:
       ``\,`this'            % \, separates the double and single quote.
        is what I just
        wrote, not  `that'\,''.

Dashes come in three sizes: an
       intra-word
dash, a medium dash for number ranges like
       1--2,
and a punctuation
       dash---like
this.

A sentence-ending space should be larger than the
space between words within a sentence.  You
sometimes have to type special commands in
conjunction with punctuation characters to get
this right, as in the following sentence.
       Gnats, gnus, etc.\ all  % `\ ' makes an inter-word space.
       begin with G\@.         % \@ marks end-of-sentence punctuation.
You should check the spaces after periods when
reading your output to make sure you haven't
forgotten any special cases.  Generating an
ellipsis
       \ldots\               % `\ ' is needed after `\ldots' because TeX
                             % ignores spaces after command names like \ldots
                             % made from \ + letters.
                             %
                             % Note how a `%' character causes TeX to ignore
                             % the end of the input line, so these blank lines
                             % do not start a new paragraph.
                             %
with the right spacing around the periods requires
a special command.

\LaTeX\ interprets some common characters as
commands, so you must type special commands to
generate them.  These characters include the
following:
       \$ \& \% \# \{ and \}.

In printing, text is usually emphasized with an
       \emph{italic}
type style.

\begin{em}
   A long segment of text can also be emphasized
   in this way.  Text within such a segment can be
   given \emph{additional} emphasis.
\end{em}

It is sometimes necessary to prevent \LaTeX\ from
breaking a line where it might otherwise do so.
This may be at a space, as between the ``Mr.''\ and
``Jones'' in
       ``Mr.~Jones'',        % ~ produces an unbreakable interword space.
or within a word---especially when the word is a
symbol like
       \mbox{\emph{itemnum}}
that makes little sense when hyphenated across
lines.

Footnotes\footnote{This is an example of a footnote.}
pose no problem.

\LaTeX\ is good at typesetting mathematical formulas
like
       \( x-3y + z = 7 \)
or
       \( a_{1} > x^{2n} + y^{2n} > x' \)
or
       \( \ip{A}{B} = \sum_{i} a_{i} b_{i} \).
The spaces you type in a formula are
ignored.  Remember that a letter like
       $x$                   % $ ... $  and  \( ... \)  are equivalent
is a formula when it denotes a mathematical
symbol, and it should be typed as one.

\section{Displayed Text}

Text is displayed by indenting it from the left
margin.  Quotations are commonly displayed.  There
are short quotations
\begin{quote}
   This is a short quotation.  It consists of a
   single paragraph of text.  See how it is formatted.
\end{quote}
and longer ones.
\begin{quotation}
   This is a longer quotation.  It consists of two
   paragraphs of text, neither of which are
   particularly interesting.

   This is the second paragraph of the quotation.  It
   is just as dull as the first paragraph.
\end{quotation}
Another frequently-displayed structure is a list.
The following is an example of an \emph{itemized}
list.
\begin{itemize}
   \item This is the first item of an itemized list.
         Each item in the list is marked with a ``tick''.
         You don't have to worry about what kind of tick
         mark is used.

   \item This is the second item of the list.  It
         contains another list nested inside it.  The inner
         list is an \emph{enumerated} list.
         \begin{enumerate}
            \item This is the first item of an enumerated
                  list that is nested within the itemized list.

            \item This is the second item of the inner list.
                  \LaTeX\ allows you to nest lists deeper than
                  you really should.
         \end{enumerate}
         This is the rest of the second item of the outer
         list.  It is no more interesting than any other
         part of the item.
   \item This is the third item of the list.
\end{itemize}
You can even display poetry.
\begin{verse}
   There is an environment
    for verse \\             % The \\ command separates lines
   Whose features some poets % within a stanza.
   will curse.

                             % One or more blank lines separate stanzas.

   For instead of making\\
   Them do \emph{all} line breaking, \\
   It allows them to put too many words on a line when they'd rather be
   forced to be terse.
\end{verse}

Mathematical formulas may also be displayed.  A
displayed formula
is
one-line long; multiline
formulas require special formatting instructions.
   \[  \ip{\Gamma}{\psi'} = x'' + y^{2} + z_{i}^{n}\]
Don't start a paragraph with a displayed equation,
nor make one a paragraph by itself.

\end{document}               % End of document.

\section{Bench Huge Section 1414}
% synthetic bench section 1414
% This is a sample LaTeX input file.  (Version of 12 August 2004.)
%
% A '%' character causes TeX to ignore all remaining text on the line,
% and is used for comments like this one.

\documentclass{article}      % Specifies the document class

                             % The preamble begins here.
\title{An Example Document}  % Declares the document's title.
\author{Leslie Lamport}      % Declares the author's name.
\date{January 21, 1994}      % Deleting this command produces today's date.

\newcommand{\ip}[2]{(#1, #2)}
                             % Defines \ip{arg1}{arg2} to mean
                             % (arg1, arg2).

%\newcommand{\ip}[2]{\langle #1 | #2\rangle}
                             % This is an alternative definition of
                             % \ip that is commented out.

\begin{document}             % End of preamble and beginning of text.

\maketitle                   % Produces the title.

This is an example input file.  Comparing it with
the output it generates can show you how to
produce a simple document of your own.

\section{Ordinary Text}      % Produces section heading.  Lower-level
                             % sections are begun with similar
                             % \subsection and \subsubsection commands.

The ends  of words and sentences are marked
  by   spaces. It  doesn't matter how many
spaces    you type; one is as good as 100.  The
end of   a line counts as a space.

One   or more   blank lines denote the  end
of  a paragraph.

Since any number of consecutive spaces are treated
like a single one, the formatting of the input
file makes no difference to
      \LaTeX,                % The \LaTeX command generates the LaTeX logo.
but it makes a difference to you.  When you use
\LaTeX, making your input file as easy to read
as possible will be a great help as you write
your document and when you change it.  This sample
file shows how you can add comments to your own input
file.

Because printing is different from typewriting,
there are a number of things that you have to do
differently when preparing an input file than if
you were just typing the document directly.
Quotation marks like
       ``this''
have to be handled specially, as do quotes within
quotes:
       ``\,`this'            % \, separates the double and single quote.
        is what I just
        wrote, not  `that'\,''.

Dashes come in three sizes: an
       intra-word
dash, a medium dash for number ranges like
       1--2,
and a punctuation
       dash---like
this.

A sentence-ending space should be larger than the
space between words within a sentence.  You
sometimes have to type special commands in
conjunction with punctuation characters to get
this right, as in the following sentence.
       Gnats, gnus, etc.\ all  % `\ ' makes an inter-word space.
       begin with G\@.         % \@ marks end-of-sentence punctuation.
You should check the spaces after periods when
reading your output to make sure you haven't
forgotten any special cases.  Generating an
ellipsis
       \ldots\               % `\ ' is needed after `\ldots' because TeX
                             % ignores spaces after command names like \ldots
                             % made from \ + letters.
                             %
                             % Note how a `%' character causes TeX to ignore
                             % the end of the input line, so these blank lines
                             % do not start a new paragraph.
                             %
with the right spacing around the periods requires
a special command.

\LaTeX\ interprets some common characters as
commands, so you must type special commands to
generate them.  These characters include the
following:
       \$ \& \% \# \{ and \}.

In printing, text is usually emphasized with an
       \emph{italic}
type style.

\begin{em}
   A long segment of text can also be emphasized
   in this way.  Text within such a segment can be
   given \emph{additional} emphasis.
\end{em}

It is sometimes necessary to prevent \LaTeX\ from
breaking a line where it might otherwise do so.
This may be at a space, as between the ``Mr.''\ and
``Jones'' in
       ``Mr.~Jones'',        % ~ produces an unbreakable interword space.
or within a word---especially when the word is a
symbol like
       \mbox{\emph{itemnum}}
that makes little sense when hyphenated across
lines.

Footnotes\footnote{This is an example of a footnote.}
pose no problem.

\LaTeX\ is good at typesetting mathematical formulas
like
       \( x-3y + z = 7 \)
or
       \( a_{1} > x^{2n} + y^{2n} > x' \)
or
       \( \ip{A}{B} = \sum_{i} a_{i} b_{i} \).
The spaces you type in a formula are
ignored.  Remember that a letter like
       $x$                   % $ ... $  and  \( ... \)  are equivalent
is a formula when it denotes a mathematical
symbol, and it should be typed as one.

\section{Displayed Text}

Text is displayed by indenting it from the left
margin.  Quotations are commonly displayed.  There
are short quotations
\begin{quote}
   This is a short quotation.  It consists of a
   single paragraph of text.  See how it is formatted.
\end{quote}
and longer ones.
\begin{quotation}
   This is a longer quotation.  It consists of two
   paragraphs of text, neither of which are
   particularly interesting.

   This is the second paragraph of the quotation.  It
   is just as dull as the first paragraph.
\end{quotation}
Another frequently-displayed structure is a list.
The following is an example of an \emph{itemized}
list.
\begin{itemize}
   \item This is the first item of an itemized list.
         Each item in the list is marked with a ``tick''.
         You don't have to worry about what kind of tick
         mark is used.

   \item This is the second item of the list.  It
         contains another list nested inside it.  The inner
         list is an \emph{enumerated} list.
         \begin{enumerate}
            \item This is the first item of an enumerated
                  list that is nested within the itemized list.

            \item This is the second item of the inner list.
                  \LaTeX\ allows you to nest lists deeper than
                  you really should.
         \end{enumerate}
         This is the rest of the second item of the outer
         list.  It is no more interesting than any other
         part of the item.
   \item This is the third item of the list.
\end{itemize}
You can even display poetry.
\begin{verse}
   There is an environment
    for verse \\             % The \\ command separates lines
   Whose features some poets % within a stanza.
   will curse.

                             % One or more blank lines separate stanzas.

   For instead of making\\
   Them do \emph{all} line breaking, \\
   It allows them to put too many words on a line when they'd rather be
   forced to be terse.
\end{verse}

Mathematical formulas may also be displayed.  A
displayed formula
is
one-line long; multiline
formulas require special formatting instructions.
   \[  \ip{\Gamma}{\psi'} = x'' + y^{2} + z_{i}^{n}\]
Don't start a paragraph with a displayed equation,
nor make one a paragraph by itself.

\end{document}               % End of document.

\section{Bench Huge Section 1415}
% synthetic bench section 1415
% This is a sample LaTeX input file.  (Version of 12 August 2004.)
%
% A '%' character causes TeX to ignore all remaining text on the line,
% and is used for comments like this one.

\documentclass{article}      % Specifies the document class

                             % The preamble begins here.
\title{An Example Document}  % Declares the document's title.
\author{Leslie Lamport}      % Declares the author's name.
\date{January 21, 1994}      % Deleting this command produces today's date.

\newcommand{\ip}[2]{(#1, #2)}
                             % Defines \ip{arg1}{arg2} to mean
                             % (arg1, arg2).

%\newcommand{\ip}[2]{\langle #1 | #2\rangle}
                             % This is an alternative definition of
                             % \ip that is commented out.

\begin{document}             % End of preamble and beginning of text.

\maketitle                   % Produces the title.

This is an example input file.  Comparing it with
the output it generates can show you how to
produce a simple document of your own.

\section{Ordinary Text}      % Produces section heading.  Lower-level
                             % sections are begun with similar
                             % \subsection and \subsubsection commands.

The ends  of words and sentences are marked
  by   spaces. It  doesn't matter how many
spaces    you type; one is as good as 100.  The
end of   a line counts as a space.

One   or more   blank lines denote the  end
of  a paragraph.

Since any number of consecutive spaces are treated
like a single one, the formatting of the input
file makes no difference to
      \LaTeX,                % The \LaTeX command generates the LaTeX logo.
but it makes a difference to you.  When you use
\LaTeX, making your input file as easy to read
as possible will be a great help as you write
your document and when you change it.  This sample
file shows how you can add comments to your own input
file.

Because printing is different from typewriting,
there are a number of things that you have to do
differently when preparing an input file than if
you were just typing the document directly.
Quotation marks like
       ``this''
have to be handled specially, as do quotes within
quotes:
       ``\,`this'            % \, separates the double and single quote.
        is what I just
        wrote, not  `that'\,''.

Dashes come in three sizes: an
       intra-word
dash, a medium dash for number ranges like
       1--2,
and a punctuation
       dash---like
this.

A sentence-ending space should be larger than the
space between words within a sentence.  You
sometimes have to type special commands in
conjunction with punctuation characters to get
this right, as in the following sentence.
       Gnats, gnus, etc.\ all  % `\ ' makes an inter-word space.
       begin with G\@.         % \@ marks end-of-sentence punctuation.
You should check the spaces after periods when
reading your output to make sure you haven't
forgotten any special cases.  Generating an
ellipsis
       \ldots\               % `\ ' is needed after `\ldots' because TeX
                             % ignores spaces after command names like \ldots
                             % made from \ + letters.
                             %
                             % Note how a `%' character causes TeX to ignore
                             % the end of the input line, so these blank lines
                             % do not start a new paragraph.
                             %
with the right spacing around the periods requires
a special command.

\LaTeX\ interprets some common characters as
commands, so you must type special commands to
generate them.  These characters include the
following:
       \$ \& \% \# \{ and \}.

In printing, text is usually emphasized with an
       \emph{italic}
type style.

\begin{em}
   A long segment of text can also be emphasized
   in this way.  Text within such a segment can be
   given \emph{additional} emphasis.
\end{em}

It is sometimes necessary to prevent \LaTeX\ from
breaking a line where it might otherwise do so.
This may be at a space, as between the ``Mr.''\ and
``Jones'' in
       ``Mr.~Jones'',        % ~ produces an unbreakable interword space.
or within a word---especially when the word is a
symbol like
       \mbox{\emph{itemnum}}
that makes little sense when hyphenated across
lines.

Footnotes\footnote{This is an example of a footnote.}
pose no problem.

\LaTeX\ is good at typesetting mathematical formulas
like
       \( x-3y + z = 7 \)
or
       \( a_{1} > x^{2n} + y^{2n} > x' \)
or
       \( \ip{A}{B} = \sum_{i} a_{i} b_{i} \).
The spaces you type in a formula are
ignored.  Remember that a letter like
       $x$                   % $ ... $  and  \( ... \)  are equivalent
is a formula when it denotes a mathematical
symbol, and it should be typed as one.

\section{Displayed Text}

Text is displayed by indenting it from the left
margin.  Quotations are commonly displayed.  There
are short quotations
\begin{quote}
   This is a short quotation.  It consists of a
   single paragraph of text.  See how it is formatted.
\end{quote}
and longer ones.
\begin{quotation}
   This is a longer quotation.  It consists of two
   paragraphs of text, neither of which are
   particularly interesting.

   This is the second paragraph of the quotation.  It
   is just as dull as the first paragraph.
\end{quotation}
Another frequently-displayed structure is a list.
The following is an example of an \emph{itemized}
list.
\begin{itemize}
   \item This is the first item of an itemized list.
         Each item in the list is marked with a ``tick''.
         You don't have to worry about what kind of tick
         mark is used.

   \item This is the second item of the list.  It
         contains another list nested inside it.  The inner
         list is an \emph{enumerated} list.
         \begin{enumerate}
            \item This is the first item of an enumerated
                  list that is nested within the itemized list.

            \item This is the second item of the inner list.
                  \LaTeX\ allows you to nest lists deeper than
                  you really should.
         \end{enumerate}
         This is the rest of the second item of the outer
         list.  It is no more interesting than any other
         part of the item.
   \item This is the third item of the list.
\end{itemize}
You can even display poetry.
\begin{verse}
   There is an environment
    for verse \\             % The \\ command separates lines
   Whose features some poets % within a stanza.
   will curse.

                             % One or more blank lines separate stanzas.

   For instead of making\\
   Them do \emph{all} line breaking, \\
   It allows them to put too many words on a line when they'd rather be
   forced to be terse.
\end{verse}

Mathematical formulas may also be displayed.  A
displayed formula
is
one-line long; multiline
formulas require special formatting instructions.
   \[  \ip{\Gamma}{\psi'} = x'' + y^{2} + z_{i}^{n}\]
Don't start a paragraph with a displayed equation,
nor make one a paragraph by itself.

\end{document}               % End of document.

\section{Bench Huge Section 1416}
% synthetic bench section 1416
% This is a sample LaTeX input file.  (Version of 12 August 2004.)
%
% A '%' character causes TeX to ignore all remaining text on the line,
% and is used for comments like this one.

\documentclass{article}      % Specifies the document class

                             % The preamble begins here.
\title{An Example Document}  % Declares the document's title.
\author{Leslie Lamport}      % Declares the author's name.
\date{January 21, 1994}      % Deleting this command produces today's date.

\newcommand{\ip}[2]{(#1, #2)}
                             % Defines \ip{arg1}{arg2} to mean
                             % (arg1, arg2).

%\newcommand{\ip}[2]{\langle #1 | #2\rangle}
                             % This is an alternative definition of
                             % \ip that is commented out.

\begin{document}             % End of preamble and beginning of text.

\maketitle                   % Produces the title.

This is an example input file.  Comparing it with
the output it generates can show you how to
produce a simple document of your own.

\section{Ordinary Text}      % Produces section heading.  Lower-level
                             % sections are begun with similar
                             % \subsection and \subsubsection commands.

The ends  of words and sentences are marked
  by   spaces. It  doesn't matter how many
spaces    you type; one is as good as 100.  The
end of   a line counts as a space.

One   or more   blank lines denote the  end
of  a paragraph.

Since any number of consecutive spaces are treated
like a single one, the formatting of the input
file makes no difference to
      \LaTeX,                % The \LaTeX command generates the LaTeX logo.
but it makes a difference to you.  When you use
\LaTeX, making your input file as easy to read
as possible will be a great help as you write
your document and when you change it.  This sample
file shows how you can add comments to your own input
file.

Because printing is different from typewriting,
there are a number of things that you have to do
differently when preparing an input file than if
you were just typing the document directly.
Quotation marks like
       ``this''
have to be handled specially, as do quotes within
quotes:
       ``\,`this'            % \, separates the double and single quote.
        is what I just
        wrote, not  `that'\,''.

Dashes come in three sizes: an
       intra-word
dash, a medium dash for number ranges like
       1--2,
and a punctuation
       dash---like
this.

A sentence-ending space should be larger than the
space between words within a sentence.  You
sometimes have to type special commands in
conjunction with punctuation characters to get
this right, as in the following sentence.
       Gnats, gnus, etc.\ all  % `\ ' makes an inter-word space.
       begin with G\@.         % \@ marks end-of-sentence punctuation.
You should check the spaces after periods when
reading your output to make sure you haven't
forgotten any special cases.  Generating an
ellipsis
       \ldots\               % `\ ' is needed after `\ldots' because TeX
                             % ignores spaces after command names like \ldots
                             % made from \ + letters.
                             %
                             % Note how a `%' character causes TeX to ignore
                             % the end of the input line, so these blank lines
                             % do not start a new paragraph.
                             %
with the right spacing around the periods requires
a special command.

\LaTeX\ interprets some common characters as
commands, so you must type special commands to
generate them.  These characters include the
following:
       \$ \& \% \# \{ and \}.

In printing, text is usually emphasized with an
       \emph{italic}
type style.

\begin{em}
   A long segment of text can also be emphasized
   in this way.  Text within such a segment can be
   given \emph{additional} emphasis.
\end{em}

It is sometimes necessary to prevent \LaTeX\ from
breaking a line where it might otherwise do so.
This may be at a space, as between the ``Mr.''\ and
``Jones'' in
       ``Mr.~Jones'',        % ~ produces an unbreakable interword space.
or within a word---especially when the word is a
symbol like
       \mbox{\emph{itemnum}}
that makes little sense when hyphenated across
lines.

Footnotes\footnote{This is an example of a footnote.}
pose no problem.

\LaTeX\ is good at typesetting mathematical formulas
like
       \( x-3y + z = 7 \)
or
       \( a_{1} > x^{2n} + y^{2n} > x' \)
or
       \( \ip{A}{B} = \sum_{i} a_{i} b_{i} \).
The spaces you type in a formula are
ignored.  Remember that a letter like
       $x$                   % $ ... $  and  \( ... \)  are equivalent
is a formula when it denotes a mathematical
symbol, and it should be typed as one.

\section{Displayed Text}

Text is displayed by indenting it from the left
margin.  Quotations are commonly displayed.  There
are short quotations
\begin{quote}
   This is a short quotation.  It consists of a
   single paragraph of text.  See how it is formatted.
\end{quote}
and longer ones.
\begin{quotation}
   This is a longer quotation.  It consists of two
   paragraphs of text, neither of which are
   particularly interesting.

   This is the second paragraph of the quotation.  It
   is just as dull as the first paragraph.
\end{quotation}
Another frequently-displayed structure is a list.
The following is an example of an \emph{itemized}
list.
\begin{itemize}
   \item This is the first item of an itemized list.
         Each item in the list is marked with a ``tick''.
         You don't have to worry about what kind of tick
         mark is used.

   \item This is the second item of the list.  It
         contains another list nested inside it.  The inner
         list is an \emph{enumerated} list.
         \begin{enumerate}
            \item This is the first item of an enumerated
                  list that is nested within the itemized list.

            \item This is the second item of the inner list.
                  \LaTeX\ allows you to nest lists deeper than
                  you really should.
         \end{enumerate}
         This is the rest of the second item of the outer
         list.  It is no more interesting than any other
         part of the item.
   \item This is the third item of the list.
\end{itemize}
You can even display poetry.
\begin{verse}
   There is an environment
    for verse \\             % The \\ command separates lines
   Whose features some poets % within a stanza.
   will curse.

                             % One or more blank lines separate stanzas.

   For instead of making\\
   Them do \emph{all} line breaking, \\
   It allows them to put too many words on a line when they'd rather be
   forced to be terse.
\end{verse}

Mathematical formulas may also be displayed.  A
displayed formula
is
one-line long; multiline
formulas require special formatting instructions.
   \[  \ip{\Gamma}{\psi'} = x'' + y^{2} + z_{i}^{n}\]
Don't start a paragraph with a displayed equation,
nor make one a paragraph by itself.

\end{document}               % End of document.

\section{Bench Huge Section 1417}
% synthetic bench section 1417
% This is a sample LaTeX input file.  (Version of 12 August 2004.)
%
% A '%' character causes TeX to ignore all remaining text on the line,
% and is used for comments like this one.

\documentclass{article}      % Specifies the document class

                             % The preamble begins here.
\title{An Example Document}  % Declares the document's title.
\author{Leslie Lamport}      % Declares the author's name.
\date{January 21, 1994}      % Deleting this command produces today's date.

\newcommand{\ip}[2]{(#1, #2)}
                             % Defines \ip{arg1}{arg2} to mean
                             % (arg1, arg2).

%\newcommand{\ip}[2]{\langle #1 | #2\rangle}
                             % This is an alternative definition of
                             % \ip that is commented out.

\begin{document}             % End of preamble and beginning of text.

\maketitle                   % Produces the title.

This is an example input file.  Comparing it with
the output it generates can show you how to
produce a simple document of your own.

\section{Ordinary Text}      % Produces section heading.  Lower-level
                             % sections are begun with similar
                             % \subsection and \subsubsection commands.

The ends  of words and sentences are marked
  by   spaces. It  doesn't matter how many
spaces    you type; one is as good as 100.  The
end of   a line counts as a space.

One   or more   blank lines denote the  end
of  a paragraph.

Since any number of consecutive spaces are treated
like a single one, the formatting of the input
file makes no difference to
      \LaTeX,                % The \LaTeX command generates the LaTeX logo.
but it makes a difference to you.  When you use
\LaTeX, making your input file as easy to read
as possible will be a great help as you write
your document and when you change it.  This sample
file shows how you can add comments to your own input
file.

Because printing is different from typewriting,
there are a number of things that you have to do
differently when preparing an input file than if
you were just typing the document directly.
Quotation marks like
       ``this''
have to be handled specially, as do quotes within
quotes:
       ``\,`this'            % \, separates the double and single quote.
        is what I just
        wrote, not  `that'\,''.

Dashes come in three sizes: an
       intra-word
dash, a medium dash for number ranges like
       1--2,
and a punctuation
       dash---like
this.

A sentence-ending space should be larger than the
space between words within a sentence.  You
sometimes have to type special commands in
conjunction with punctuation characters to get
this right, as in the following sentence.
       Gnats, gnus, etc.\ all  % `\ ' makes an inter-word space.
       begin with G\@.         % \@ marks end-of-sentence punctuation.
You should check the spaces after periods when
reading your output to make sure you haven't
forgotten any special cases.  Generating an
ellipsis
       \ldots\               % `\ ' is needed after `\ldots' because TeX
                             % ignores spaces after command names like \ldots
                             % made from \ + letters.
                             %
                             % Note how a `%' character causes TeX to ignore
                             % the end of the input line, so these blank lines
                             % do not start a new paragraph.
                             %
with the right spacing around the periods requires
a special command.

\LaTeX\ interprets some common characters as
commands, so you must type special commands to
generate them.  These characters include the
following:
       \$ \& \% \# \{ and \}.

In printing, text is usually emphasized with an
       \emph{italic}
type style.

\begin{em}
   A long segment of text can also be emphasized
   in this way.  Text within such a segment can be
   given \emph{additional} emphasis.
\end{em}

It is sometimes necessary to prevent \LaTeX\ from
breaking a line where it might otherwise do so.
This may be at a space, as between the ``Mr.''\ and
``Jones'' in
       ``Mr.~Jones'',        % ~ produces an unbreakable interword space.
or within a word---especially when the word is a
symbol like
       \mbox{\emph{itemnum}}
that makes little sense when hyphenated across
lines.

Footnotes\footnote{This is an example of a footnote.}
pose no problem.

\LaTeX\ is good at typesetting mathematical formulas
like
       \( x-3y + z = 7 \)
or
       \( a_{1} > x^{2n} + y^{2n} > x' \)
or
       \( \ip{A}{B} = \sum_{i} a_{i} b_{i} \).
The spaces you type in a formula are
ignored.  Remember that a letter like
       $x$                   % $ ... $  and  \( ... \)  are equivalent
is a formula when it denotes a mathematical
symbol, and it should be typed as one.

\section{Displayed Text}

Text is displayed by indenting it from the left
margin.  Quotations are commonly displayed.  There
are short quotations
\begin{quote}
   This is a short quotation.  It consists of a
   single paragraph of text.  See how it is formatted.
\end{quote}
and longer ones.
\begin{quotation}
   This is a longer quotation.  It consists of two
   paragraphs of text, neither of which are
   particularly interesting.

   This is the second paragraph of the quotation.  It
   is just as dull as the first paragraph.
\end{quotation}
Another frequently-displayed structure is a list.
The following is an example of an \emph{itemized}
list.
\begin{itemize}
   \item This is the first item of an itemized list.
         Each item in the list is marked with a ``tick''.
         You don't have to worry about what kind of tick
         mark is used.

   \item This is the second item of the list.  It
         contains another list nested inside it.  The inner
         list is an \emph{enumerated} list.
         \begin{enumerate}
            \item This is the first item of an enumerated
                  list that is nested within the itemized list.

            \item This is the second item of the inner list.
                  \LaTeX\ allows you to nest lists deeper than
                  you really should.
         \end{enumerate}
         This is the rest of the second item of the outer
         list.  It is no more interesting than any other
         part of the item.
   \item This is the third item of the list.
\end{itemize}
You can even display poetry.
\begin{verse}
   There is an environment
    for verse \\             % The \\ command separates lines
   Whose features some poets % within a stanza.
   will curse.

                             % One or more blank lines separate stanzas.

   For instead of making\\
   Them do \emph{all} line breaking, \\
   It allows them to put too many words on a line when they'd rather be
   forced to be terse.
\end{verse}

Mathematical formulas may also be displayed.  A
displayed formula
is
one-line long; multiline
formulas require special formatting instructions.
   \[  \ip{\Gamma}{\psi'} = x'' + y^{2} + z_{i}^{n}\]
Don't start a paragraph with a displayed equation,
nor make one a paragraph by itself.

\end{document}               % End of document.

\section{Bench Huge Section 1418}
% synthetic bench section 1418
% This is a sample LaTeX input file.  (Version of 12 August 2004.)
%
% A '%' character causes TeX to ignore all remaining text on the line,
% and is used for comments like this one.

\documentclass{article}      % Specifies the document class

                             % The preamble begins here.
\title{An Example Document}  % Declares the document's title.
\author{Leslie Lamport}      % Declares the author's name.
\date{January 21, 1994}      % Deleting this command produces today's date.

\newcommand{\ip}[2]{(#1, #2)}
                             % Defines \ip{arg1}{arg2} to mean
                             % (arg1, arg2).

%\newcommand{\ip}[2]{\langle #1 | #2\rangle}
                             % This is an alternative definition of
                             % \ip that is commented out.

\begin{document}             % End of preamble and beginning of text.

\maketitle                   % Produces the title.

This is an example input file.  Comparing it with
the output it generates can show you how to
produce a simple document of your own.

\section{Ordinary Text}      % Produces section heading.  Lower-level
                             % sections are begun with similar
                             % \subsection and \subsubsection commands.

The ends  of words and sentences are marked
  by   spaces. It  doesn't matter how many
spaces    you type; one is as good as 100.  The
end of   a line counts as a space.

One   or more   blank lines denote the  end
of  a paragraph.

Since any number of consecutive spaces are treated
like a single one, the formatting of the input
file makes no difference to
      \LaTeX,                % The \LaTeX command generates the LaTeX logo.
but it makes a difference to you.  When you use
\LaTeX, making your input file as easy to read
as possible will be a great help as you write
your document and when you change it.  This sample
file shows how you can add comments to your own input
file.

Because printing is different from typewriting,
there are a number of things that you have to do
differently when preparing an input file than if
you were just typing the document directly.
Quotation marks like
       ``this''
have to be handled specially, as do quotes within
quotes:
       ``\,`this'            % \, separates the double and single quote.
        is what I just
        wrote, not  `that'\,''.

Dashes come in three sizes: an
       intra-word
dash, a medium dash for number ranges like
       1--2,
and a punctuation
       dash---like
this.

A sentence-ending space should be larger than the
space between words within a sentence.  You
sometimes have to type special commands in
conjunction with punctuation characters to get
this right, as in the following sentence.
       Gnats, gnus, etc.\ all  % `\ ' makes an inter-word space.
       begin with G\@.         % \@ marks end-of-sentence punctuation.
You should check the spaces after periods when
reading your output to make sure you haven't
forgotten any special cases.  Generating an
ellipsis
       \ldots\               % `\ ' is needed after `\ldots' because TeX
                             % ignores spaces after command names like \ldots
                             % made from \ + letters.
                             %
                             % Note how a `%' character causes TeX to ignore
                             % the end of the input line, so these blank lines
                             % do not start a new paragraph.
                             %
with the right spacing around the periods requires
a special command.

\LaTeX\ interprets some common characters as
commands, so you must type special commands to
generate them.  These characters include the
following:
       \$ \& \% \# \{ and \}.

In printing, text is usually emphasized with an
       \emph{italic}
type style.

\begin{em}
   A long segment of text can also be emphasized
   in this way.  Text within such a segment can be
   given \emph{additional} emphasis.
\end{em}

It is sometimes necessary to prevent \LaTeX\ from
breaking a line where it might otherwise do so.
This may be at a space, as between the ``Mr.''\ and
``Jones'' in
       ``Mr.~Jones'',        % ~ produces an unbreakable interword space.
or within a word---especially when the word is a
symbol like
       \mbox{\emph{itemnum}}
that makes little sense when hyphenated across
lines.

Footnotes\footnote{This is an example of a footnote.}
pose no problem.

\LaTeX\ is good at typesetting mathematical formulas
like
       \( x-3y + z = 7 \)
or
       \( a_{1} > x^{2n} + y^{2n} > x' \)
or
       \( \ip{A}{B} = \sum_{i} a_{i} b_{i} \).
The spaces you type in a formula are
ignored.  Remember that a letter like
       $x$                   % $ ... $  and  \( ... \)  are equivalent
is a formula when it denotes a mathematical
symbol, and it should be typed as one.

\section{Displayed Text}

Text is displayed by indenting it from the left
margin.  Quotations are commonly displayed.  There
are short quotations
\begin{quote}
   This is a short quotation.  It consists of a
   single paragraph of text.  See how it is formatted.
\end{quote}
and longer ones.
\begin{quotation}
   This is a longer quotation.  It consists of two
   paragraphs of text, neither of which are
   particularly interesting.

   This is the second paragraph of the quotation.  It
   is just as dull as the first paragraph.
\end{quotation}
Another frequently-displayed structure is a list.
The following is an example of an \emph{itemized}
list.
\begin{itemize}
   \item This is the first item of an itemized list.
         Each item in the list is marked with a ``tick''.
         You don't have to worry about what kind of tick
         mark is used.

   \item This is the second item of the list.  It
         contains another list nested inside it.  The inner
         list is an \emph{enumerated} list.
         \begin{enumerate}
            \item This is the first item of an enumerated
                  list that is nested within the itemized list.

            \item This is the second item of the inner list.
                  \LaTeX\ allows you to nest lists deeper than
                  you really should.
         \end{enumerate}
         This is the rest of the second item of the outer
         list.  It is no more interesting than any other
         part of the item.
   \item This is the third item of the list.
\end{itemize}
You can even display poetry.
\begin{verse}
   There is an environment
    for verse \\             % The \\ command separates lines
   Whose features some poets % within a stanza.
   will curse.

                             % One or more blank lines separate stanzas.

   For instead of making\\
   Them do \emph{all} line breaking, \\
   It allows them to put too many words on a line when they'd rather be
   forced to be terse.
\end{verse}

Mathematical formulas may also be displayed.  A
displayed formula
is
one-line long; multiline
formulas require special formatting instructions.
   \[  \ip{\Gamma}{\psi'} = x'' + y^{2} + z_{i}^{n}\]
Don't start a paragraph with a displayed equation,
nor make one a paragraph by itself.

\end{document}               % End of document.

\section{Bench Huge Section 1419}
% synthetic bench section 1419
% This is a sample LaTeX input file.  (Version of 12 August 2004.)
%
% A '%' character causes TeX to ignore all remaining text on the line,
% and is used for comments like this one.

\documentclass{article}      % Specifies the document class

                             % The preamble begins here.
\title{An Example Document}  % Declares the document's title.
\author{Leslie Lamport}      % Declares the author's name.
\date{January 21, 1994}      % Deleting this command produces today's date.

\newcommand{\ip}[2]{(#1, #2)}
                             % Defines \ip{arg1}{arg2} to mean
                             % (arg1, arg2).

%\newcommand{\ip}[2]{\langle #1 | #2\rangle}
                             % This is an alternative definition of
                             % \ip that is commented out.

\begin{document}             % End of preamble and beginning of text.

\maketitle                   % Produces the title.

This is an example input file.  Comparing it with
the output it generates can show you how to
produce a simple document of your own.

\section{Ordinary Text}      % Produces section heading.  Lower-level
                             % sections are begun with similar
                             % \subsection and \subsubsection commands.

The ends  of words and sentences are marked
  by   spaces. It  doesn't matter how many
spaces    you type; one is as good as 100.  The
end of   a line counts as a space.

One   or more   blank lines denote the  end
of  a paragraph.

Since any number of consecutive spaces are treated
like a single one, the formatting of the input
file makes no difference to
      \LaTeX,                % The \LaTeX command generates the LaTeX logo.
but it makes a difference to you.  When you use
\LaTeX, making your input file as easy to read
as possible will be a great help as you write
your document and when you change it.  This sample
file shows how you can add comments to your own input
file.

Because printing is different from typewriting,
there are a number of things that you have to do
differently when preparing an input file than if
you were just typing the document directly.
Quotation marks like
       ``this''
have to be handled specially, as do quotes within
quotes:
       ``\,`this'            % \, separates the double and single quote.
        is what I just
        wrote, not  `that'\,''.

Dashes come in three sizes: an
       intra-word
dash, a medium dash for number ranges like
       1--2,
and a punctuation
       dash---like
this.

A sentence-ending space should be larger than the
space between words within a sentence.  You
sometimes have to type special commands in
conjunction with punctuation characters to get
this right, as in the following sentence.
       Gnats, gnus, etc.\ all  % `\ ' makes an inter-word space.
       begin with G\@.         % \@ marks end-of-sentence punctuation.
You should check the spaces after periods when
reading your output to make sure you haven't
forgotten any special cases.  Generating an
ellipsis
       \ldots\               % `\ ' is needed after `\ldots' because TeX
                             % ignores spaces after command names like \ldots
                             % made from \ + letters.
                             %
                             % Note how a `%' character causes TeX to ignore
                             % the end of the input line, so these blank lines
                             % do not start a new paragraph.
                             %
with the right spacing around the periods requires
a special command.

\LaTeX\ interprets some common characters as
commands, so you must type special commands to
generate them.  These characters include the
following:
       \$ \& \% \# \{ and \}.

In printing, text is usually emphasized with an
       \emph{italic}
type style.

\begin{em}
   A long segment of text can also be emphasized
   in this way.  Text within such a segment can be
   given \emph{additional} emphasis.
\end{em}

It is sometimes necessary to prevent \LaTeX\ from
breaking a line where it might otherwise do so.
This may be at a space, as between the ``Mr.''\ and
``Jones'' in
       ``Mr.~Jones'',        % ~ produces an unbreakable interword space.
or within a word---especially when the word is a
symbol like
       \mbox{\emph{itemnum}}
that makes little sense when hyphenated across
lines.

Footnotes\footnote{This is an example of a footnote.}
pose no problem.

\LaTeX\ is good at typesetting mathematical formulas
like
       \( x-3y + z = 7 \)
or
       \( a_{1} > x^{2n} + y^{2n} > x' \)
or
       \( \ip{A}{B} = \sum_{i} a_{i} b_{i} \).
The spaces you type in a formula are
ignored.  Remember that a letter like
       $x$                   % $ ... $  and  \( ... \)  are equivalent
is a formula when it denotes a mathematical
symbol, and it should be typed as one.

\section{Displayed Text}

Text is displayed by indenting it from the left
margin.  Quotations are commonly displayed.  There
are short quotations
\begin{quote}
   This is a short quotation.  It consists of a
   single paragraph of text.  See how it is formatted.
\end{quote}
and longer ones.
\begin{quotation}
   This is a longer quotation.  It consists of two
   paragraphs of text, neither of which are
   particularly interesting.

   This is the second paragraph of the quotation.  It
   is just as dull as the first paragraph.
\end{quotation}
Another frequently-displayed structure is a list.
The following is an example of an \emph{itemized}
list.
\begin{itemize}
   \item This is the first item of an itemized list.
         Each item in the list is marked with a ``tick''.
         You don't have to worry about what kind of tick
         mark is used.

   \item This is the second item of the list.  It
         contains another list nested inside it.  The inner
         list is an \emph{enumerated} list.
         \begin{enumerate}
            \item This is the first item of an enumerated
                  list that is nested within the itemized list.

            \item This is the second item of the inner list.
                  \LaTeX\ allows you to nest lists deeper than
                  you really should.
         \end{enumerate}
         This is the rest of the second item of the outer
         list.  It is no more interesting than any other
         part of the item.
   \item This is the third item of the list.
\end{itemize}
You can even display poetry.
\begin{verse}
   There is an environment
    for verse \\             % The \\ command separates lines
   Whose features some poets % within a stanza.
   will curse.

                             % One or more blank lines separate stanzas.

   For instead of making\\
   Them do \emph{all} line breaking, \\
   It allows them to put too many words on a line when they'd rather be
   forced to be terse.
\end{verse}

Mathematical formulas may also be displayed.  A
displayed formula
is
one-line long; multiline
formulas require special formatting instructions.
   \[  \ip{\Gamma}{\psi'} = x'' + y^{2} + z_{i}^{n}\]
Don't start a paragraph with a displayed equation,
nor make one a paragraph by itself.

\end{document}               % End of document.

\section{Bench Huge Section 1420}
% synthetic bench section 1420
% This is a sample LaTeX input file.  (Version of 12 August 2004.)
%
% A '%' character causes TeX to ignore all remaining text on the line,
% and is used for comments like this one.

\documentclass{article}      % Specifies the document class

                             % The preamble begins here.
\title{An Example Document}  % Declares the document's title.
\author{Leslie Lamport}      % Declares the author's name.
\date{January 21, 1994}      % Deleting this command produces today's date.

\newcommand{\ip}[2]{(#1, #2)}
                             % Defines \ip{arg1}{arg2} to mean
                             % (arg1, arg2).

%\newcommand{\ip}[2]{\langle #1 | #2\rangle}
                             % This is an alternative definition of
                             % \ip that is commented out.

\begin{document}             % End of preamble and beginning of text.

\maketitle                   % Produces the title.

This is an example input file.  Comparing it with
the output it generates can show you how to
produce a simple document of your own.

\section{Ordinary Text}      % Produces section heading.  Lower-level
                             % sections are begun with similar
                             % \subsection and \subsubsection commands.

The ends  of words and sentences are marked
  by   spaces. It  doesn't matter how many
spaces    you type; one is as good as 100.  The
end of   a line counts as a space.

One   or more   blank lines denote the  end
of  a paragraph.

Since any number of consecutive spaces are treated
like a single one, the formatting of the input
file makes no difference to
      \LaTeX,                % The \LaTeX command generates the LaTeX logo.
but it makes a difference to you.  When you use
\LaTeX, making your input file as easy to read
as possible will be a great help as you write
your document and when you change it.  This sample
file shows how you can add comments to your own input
file.

Because printing is different from typewriting,
there are a number of things that you have to do
differently when preparing an input file than if
you were just typing the document directly.
Quotation marks like
       ``this''
have to be handled specially, as do quotes within
quotes:
       ``\,`this'            % \, separates the double and single quote.
        is what I just
        wrote, not  `that'\,''.

Dashes come in three sizes: an
       intra-word
dash, a medium dash for number ranges like
       1--2,
and a punctuation
       dash---like
this.

A sentence-ending space should be larger than the
space between words within a sentence.  You
sometimes have to type special commands in
conjunction with punctuation characters to get
this right, as in the following sentence.
       Gnats, gnus, etc.\ all  % `\ ' makes an inter-word space.
       begin with G\@.         % \@ marks end-of-sentence punctuation.
You should check the spaces after periods when
reading your output to make sure you haven't
forgotten any special cases.  Generating an
ellipsis
       \ldots\               % `\ ' is needed after `\ldots' because TeX
                             % ignores spaces after command names like \ldots
                             % made from \ + letters.
                             %
                             % Note how a `%' character causes TeX to ignore
                             % the end of the input line, so these blank lines
                             % do not start a new paragraph.
                             %
with the right spacing around the periods requires
a special command.

\LaTeX\ interprets some common characters as
commands, so you must type special commands to
generate them.  These characters include the
following:
       \$ \& \% \# \{ and \}.

In printing, text is usually emphasized with an
       \emph{italic}
type style.

\begin{em}
   A long segment of text can also be emphasized
   in this way.  Text within such a segment can be
   given \emph{additional} emphasis.
\end{em}

It is sometimes necessary to prevent \LaTeX\ from
breaking a line where it might otherwise do so.
This may be at a space, as between the ``Mr.''\ and
``Jones'' in
       ``Mr.~Jones'',        % ~ produces an unbreakable interword space.
or within a word---especially when the word is a
symbol like
       \mbox{\emph{itemnum}}
that makes little sense when hyphenated across
lines.

Footnotes\footnote{This is an example of a footnote.}
pose no problem.

\LaTeX\ is good at typesetting mathematical formulas
like
       \( x-3y + z = 7 \)
or
       \( a_{1} > x^{2n} + y^{2n} > x' \)
or
       \( \ip{A}{B} = \sum_{i} a_{i} b_{i} \).
The spaces you type in a formula are
ignored.  Remember that a letter like
       $x$                   % $ ... $  and  \( ... \)  are equivalent
is a formula when it denotes a mathematical
symbol, and it should be typed as one.

\section{Displayed Text}

Text is displayed by indenting it from the left
margin.  Quotations are commonly displayed.  There
are short quotations
\begin{quote}
   This is a short quotation.  It consists of a
   single paragraph of text.  See how it is formatted.
\end{quote}
and longer ones.
\begin{quotation}
   This is a longer quotation.  It consists of two
   paragraphs of text, neither of which are
   particularly interesting.

   This is the second paragraph of the quotation.  It
   is just as dull as the first paragraph.
\end{quotation}
Another frequently-displayed structure is a list.
The following is an example of an \emph{itemized}
list.
\begin{itemize}
   \item This is the first item of an itemized list.
         Each item in the list is marked with a ``tick''.
         You don't have to worry about what kind of tick
         mark is used.

   \item This is the second item of the list.  It
         contains another list nested inside it.  The inner
         list is an \emph{enumerated} list.
         \begin{enumerate}
            \item This is the first item of an enumerated
                  list that is nested within the itemized list.

            \item This is the second item of the inner list.
                  \LaTeX\ allows you to nest lists deeper than
                  you really should.
         \end{enumerate}
         This is the rest of the second item of the outer
         list.  It is no more interesting than any other
         part of the item.
   \item This is the third item of the list.
\end{itemize}
You can even display poetry.
\begin{verse}
   There is an environment
    for verse \\             % The \\ command separates lines
   Whose features some poets % within a stanza.
   will curse.

                             % One or more blank lines separate stanzas.

   For instead of making\\
   Them do \emph{all} line breaking, \\
   It allows them to put too many words on a line when they'd rather be
   forced to be terse.
\end{verse}

Mathematical formulas may also be displayed.  A
displayed formula
is
one-line long; multiline
formulas require special formatting instructions.
   \[  \ip{\Gamma}{\psi'} = x'' + y^{2} + z_{i}^{n}\]
Don't start a paragraph with a displayed equation,
nor make one a paragraph by itself.

\end{document}               % End of document.

\section{Bench Huge Section 1421}
% synthetic bench section 1421
% This is a sample LaTeX input file.  (Version of 12 August 2004.)
%
% A '%' character causes TeX to ignore all remaining text on the line,
% and is used for comments like this one.

\documentclass{article}      % Specifies the document class

                             % The preamble begins here.
\title{An Example Document}  % Declares the document's title.
\author{Leslie Lamport}      % Declares the author's name.
\date{January 21, 1994}      % Deleting this command produces today's date.

\newcommand{\ip}[2]{(#1, #2)}
                             % Defines \ip{arg1}{arg2} to mean
                             % (arg1, arg2).

%\newcommand{\ip}[2]{\langle #1 | #2\rangle}
                             % This is an alternative definition of
                             % \ip that is commented out.

\begin{document}             % End of preamble and beginning of text.

\maketitle                   % Produces the title.

This is an example input file.  Comparing it with
the output it generates can show you how to
produce a simple document of your own.

\section{Ordinary Text}      % Produces section heading.  Lower-level
                             % sections are begun with similar
                             % \subsection and \subsubsection commands.

The ends  of words and sentences are marked
  by   spaces. It  doesn't matter how many
spaces    you type; one is as good as 100.  The
end of   a line counts as a space.

One   or more   blank lines denote the  end
of  a paragraph.

Since any number of consecutive spaces are treated
like a single one, the formatting of the input
file makes no difference to
      \LaTeX,                % The \LaTeX command generates the LaTeX logo.
but it makes a difference to you.  When you use
\LaTeX, making your input file as easy to read
as possible will be a great help as you write
your document and when you change it.  This sample
file shows how you can add comments to your own input
file.

Because printing is different from typewriting,
there are a number of things that you have to do
differently when preparing an input file than if
you were just typing the document directly.
Quotation marks like
       ``this''
have to be handled specially, as do quotes within
quotes:
       ``\,`this'            % \, separates the double and single quote.
        is what I just
        wrote, not  `that'\,''.

Dashes come in three sizes: an
       intra-word
dash, a medium dash for number ranges like
       1--2,
and a punctuation
       dash---like
this.

A sentence-ending space should be larger than the
space between words within a sentence.  You
sometimes have to type special commands in
conjunction with punctuation characters to get
this right, as in the following sentence.
       Gnats, gnus, etc.\ all  % `\ ' makes an inter-word space.
       begin with G\@.         % \@ marks end-of-sentence punctuation.
You should check the spaces after periods when
reading your output to make sure you haven't
forgotten any special cases.  Generating an
ellipsis
       \ldots\               % `\ ' is needed after `\ldots' because TeX
                             % ignores spaces after command names like \ldots
                             % made from \ + letters.
                             %
                             % Note how a `%' character causes TeX to ignore
                             % the end of the input line, so these blank lines
                             % do not start a new paragraph.
                             %
with the right spacing around the periods requires
a special command.

\LaTeX\ interprets some common characters as
commands, so you must type special commands to
generate them.  These characters include the
following:
       \$ \& \% \# \{ and \}.

In printing, text is usually emphasized with an
       \emph{italic}
type style.

\begin{em}
   A long segment of text can also be emphasized
   in this way.  Text within such a segment can be
   given \emph{additional} emphasis.
\end{em}

It is sometimes necessary to prevent \LaTeX\ from
breaking a line where it might otherwise do so.
This may be at a space, as between the ``Mr.''\ and
``Jones'' in
       ``Mr.~Jones'',        % ~ produces an unbreakable interword space.
or within a word---especially when the word is a
symbol like
       \mbox{\emph{itemnum}}
that makes little sense when hyphenated across
lines.

Footnotes\footnote{This is an example of a footnote.}
pose no problem.

\LaTeX\ is good at typesetting mathematical formulas
like
       \( x-3y + z = 7 \)
or
       \( a_{1} > x^{2n} + y^{2n} > x' \)
or
       \( \ip{A}{B} = \sum_{i} a_{i} b_{i} \).
The spaces you type in a formula are
ignored.  Remember that a letter like
       $x$                   % $ ... $  and  \( ... \)  are equivalent
is a formula when it denotes a mathematical
symbol, and it should be typed as one.

\section{Displayed Text}

Text is displayed by indenting it from the left
margin.  Quotations are commonly displayed.  There
are short quotations
\begin{quote}
   This is a short quotation.  It consists of a
   single paragraph of text.  See how it is formatted.
\end{quote}
and longer ones.
\begin{quotation}
   This is a longer quotation.  It consists of two
   paragraphs of text, neither of which are
   particularly interesting.

   This is the second paragraph of the quotation.  It
   is just as dull as the first paragraph.
\end{quotation}
Another frequently-displayed structure is a list.
The following is an example of an \emph{itemized}
list.
\begin{itemize}
   \item This is the first item of an itemized list.
         Each item in the list is marked with a ``tick''.
         You don't have to worry about what kind of tick
         mark is used.

   \item This is the second item of the list.  It
         contains another list nested inside it.  The inner
         list is an \emph{enumerated} list.
         \begin{enumerate}
            \item This is the first item of an enumerated
                  list that is nested within the itemized list.

            \item This is the second item of the inner list.
                  \LaTeX\ allows you to nest lists deeper than
                  you really should.
         \end{enumerate}
         This is the rest of the second item of the outer
         list.  It is no more interesting than any other
         part of the item.
   \item This is the third item of the list.
\end{itemize}
You can even display poetry.
\begin{verse}
   There is an environment
    for verse \\             % The \\ command separates lines
   Whose features some poets % within a stanza.
   will curse.

                             % One or more blank lines separate stanzas.

   For instead of making\\
   Them do \emph{all} line breaking, \\
   It allows them to put too many words on a line when they'd rather be
   forced to be terse.
\end{verse}

Mathematical formulas may also be displayed.  A
displayed formula
is
one-line long; multiline
formulas require special formatting instructions.
   \[  \ip{\Gamma}{\psi'} = x'' + y^{2} + z_{i}^{n}\]
Don't start a paragraph with a displayed equation,
nor make one a paragraph by itself.

\end{document}               % End of document.

\section{Bench Huge Section 1422}
% synthetic bench section 1422
% This is a sample LaTeX input file.  (Version of 12 August 2004.)
%
% A '%' character causes TeX to ignore all remaining text on the line,
% and is used for comments like this one.

\documentclass{article}      % Specifies the document class

                             % The preamble begins here.
\title{An Example Document}  % Declares the document's title.
\author{Leslie Lamport}      % Declares the author's name.
\date{January 21, 1994}      % Deleting this command produces today's date.

\newcommand{\ip}[2]{(#1, #2)}
                             % Defines \ip{arg1}{arg2} to mean
                             % (arg1, arg2).

%\newcommand{\ip}[2]{\langle #1 | #2\rangle}
                             % This is an alternative definition of
                             % \ip that is commented out.

\begin{document}             % End of preamble and beginning of text.

\maketitle                   % Produces the title.

This is an example input file.  Comparing it with
the output it generates can show you how to
produce a simple document of your own.

\section{Ordinary Text}      % Produces section heading.  Lower-level
                             % sections are begun with similar
                             % \subsection and \subsubsection commands.

The ends  of words and sentences are marked
  by   spaces. It  doesn't matter how many
spaces    you type; one is as good as 100.  The
end of   a line counts as a space.

One   or more   blank lines denote the  end
of  a paragraph.

Since any number of consecutive spaces are treated
like a single one, the formatting of the input
file makes no difference to
      \LaTeX,                % The \LaTeX command generates the LaTeX logo.
but it makes a difference to you.  When you use
\LaTeX, making your input file as easy to read
as possible will be a great help as you write
your document and when you change it.  This sample
file shows how you can add comments to your own input
file.

Because printing is different from typewriting,
there are a number of things that you have to do
differently when preparing an input file than if
you were just typing the document directly.
Quotation marks like
       ``this''
have to be handled specially, as do quotes within
quotes:
       ``\,`this'            % \, separates the double and single quote.
        is what I just
        wrote, not  `that'\,''.

Dashes come in three sizes: an
       intra-word
dash, a medium dash for number ranges like
       1--2,
and a punctuation
       dash---like
this.

A sentence-ending space should be larger than the
space between words within a sentence.  You
sometimes have to type special commands in
conjunction with punctuation characters to get
this right, as in the following sentence.
       Gnats, gnus, etc.\ all  % `\ ' makes an inter-word space.
       begin with G\@.         % \@ marks end-of-sentence punctuation.
You should check the spaces after periods when
reading your output to make sure you haven't
forgotten any special cases.  Generating an
ellipsis
       \ldots\               % `\ ' is needed after `\ldots' because TeX
                             % ignores spaces after command names like \ldots
                             % made from \ + letters.
                             %
                             % Note how a `%' character causes TeX to ignore
                             % the end of the input line, so these blank lines
                             % do not start a new paragraph.
                             %
with the right spacing around the periods requires
a special command.

\LaTeX\ interprets some common characters as
commands, so you must type special commands to
generate them.  These characters include the
following:
       \$ \& \% \# \{ and \}.

In printing, text is usually emphasized with an
       \emph{italic}
type style.

\begin{em}
   A long segment of text can also be emphasized
   in this way.  Text within such a segment can be
   given \emph{additional} emphasis.
\end{em}

It is sometimes necessary to prevent \LaTeX\ from
breaking a line where it might otherwise do so.
This may be at a space, as between the ``Mr.''\ and
``Jones'' in
       ``Mr.~Jones'',        % ~ produces an unbreakable interword space.
or within a word---especially when the word is a
symbol like
       \mbox{\emph{itemnum}}
that makes little sense when hyphenated across
lines.

Footnotes\footnote{This is an example of a footnote.}
pose no problem.

\LaTeX\ is good at typesetting mathematical formulas
like
       \( x-3y + z = 7 \)
or
       \( a_{1} > x^{2n} + y^{2n} > x' \)
or
       \( \ip{A}{B} = \sum_{i} a_{i} b_{i} \).
The spaces you type in a formula are
ignored.  Remember that a letter like
       $x$                   % $ ... $  and  \( ... \)  are equivalent
is a formula when it denotes a mathematical
symbol, and it should be typed as one.

\section{Displayed Text}

Text is displayed by indenting it from the left
margin.  Quotations are commonly displayed.  There
are short quotations
\begin{quote}
   This is a short quotation.  It consists of a
   single paragraph of text.  See how it is formatted.
\end{quote}
and longer ones.
\begin{quotation}
   This is a longer quotation.  It consists of two
   paragraphs of text, neither of which are
   particularly interesting.

   This is the second paragraph of the quotation.  It
   is just as dull as the first paragraph.
\end{quotation}
Another frequently-displayed structure is a list.
The following is an example of an \emph{itemized}
list.
\begin{itemize}
   \item This is the first item of an itemized list.
         Each item in the list is marked with a ``tick''.
         You don't have to worry about what kind of tick
         mark is used.

   \item This is the second item of the list.  It
         contains another list nested inside it.  The inner
         list is an \emph{enumerated} list.
         \begin{enumerate}
            \item This is the first item of an enumerated
                  list that is nested within the itemized list.

            \item This is the second item of the inner list.
                  \LaTeX\ allows you to nest lists deeper than
                  you really should.
         \end{enumerate}
         This is the rest of the second item of the outer
         list.  It is no more interesting than any other
         part of the item.
   \item This is the third item of the list.
\end{itemize}
You can even display poetry.
\begin{verse}
   There is an environment
    for verse \\             % The \\ command separates lines
   Whose features some poets % within a stanza.
   will curse.

                             % One or more blank lines separate stanzas.

   For instead of making\\
   Them do \emph{all} line breaking, \\
   It allows them to put too many words on a line when they'd rather be
   forced to be terse.
\end{verse}

Mathematical formulas may also be displayed.  A
displayed formula
is
one-line long; multiline
formulas require special formatting instructions.
   \[  \ip{\Gamma}{\psi'} = x'' + y^{2} + z_{i}^{n}\]
Don't start a paragraph with a displayed equation,
nor make one a paragraph by itself.

\end{document}               % End of document.

\section{Bench Huge Section 1423}
% synthetic bench section 1423
% This is a sample LaTeX input file.  (Version of 12 August 2004.)
%
% A '%' character causes TeX to ignore all remaining text on the line,
% and is used for comments like this one.

\documentclass{article}      % Specifies the document class

                             % The preamble begins here.
\title{An Example Document}  % Declares the document's title.
\author{Leslie Lamport}      % Declares the author's name.
\date{January 21, 1994}      % Deleting this command produces today's date.

\newcommand{\ip}[2]{(#1, #2)}
                             % Defines \ip{arg1}{arg2} to mean
                             % (arg1, arg2).

%\newcommand{\ip}[2]{\langle #1 | #2\rangle}
                             % This is an alternative definition of
                             % \ip that is commented out.

\begin{document}             % End of preamble and beginning of text.

\maketitle                   % Produces the title.

This is an example input file.  Comparing it with
the output it generates can show you how to
produce a simple document of your own.

\section{Ordinary Text}      % Produces section heading.  Lower-level
                             % sections are begun with similar
                             % \subsection and \subsubsection commands.

The ends  of words and sentences are marked
  by   spaces. It  doesn't matter how many
spaces    you type; one is as good as 100.  The
end of   a line counts as a space.

One   or more   blank lines denote the  end
of  a paragraph.

Since any number of consecutive spaces are treated
like a single one, the formatting of the input
file makes no difference to
      \LaTeX,                % The \LaTeX command generates the LaTeX logo.
but it makes a difference to you.  When you use
\LaTeX, making your input file as easy to read
as possible will be a great help as you write
your document and when you change it.  This sample
file shows how you can add comments to your own input
file.

Because printing is different from typewriting,
there are a number of things that you have to do
differently when preparing an input file than if
you were just typing the document directly.
Quotation marks like
       ``this''
have to be handled specially, as do quotes within
quotes:
       ``\,`this'            % \, separates the double and single quote.
        is what I just
        wrote, not  `that'\,''.

Dashes come in three sizes: an
       intra-word
dash, a medium dash for number ranges like
       1--2,
and a punctuation
       dash---like
this.

A sentence-ending space should be larger than the
space between words within a sentence.  You
sometimes have to type special commands in
conjunction with punctuation characters to get
this right, as in the following sentence.
       Gnats, gnus, etc.\ all  % `\ ' makes an inter-word space.
       begin with G\@.         % \@ marks end-of-sentence punctuation.
You should check the spaces after periods when
reading your output to make sure you haven't
forgotten any special cases.  Generating an
ellipsis
       \ldots\               % `\ ' is needed after `\ldots' because TeX
                             % ignores spaces after command names like \ldots
                             % made from \ + letters.
                             %
                             % Note how a `%' character causes TeX to ignore
                             % the end of the input line, so these blank lines
                             % do not start a new paragraph.
                             %
with the right spacing around the periods requires
a special command.

\LaTeX\ interprets some common characters as
commands, so you must type special commands to
generate them.  These characters include the
following:
       \$ \& \% \# \{ and \}.

In printing, text is usually emphasized with an
       \emph{italic}
type style.

\begin{em}
   A long segment of text can also be emphasized
   in this way.  Text within such a segment can be
   given \emph{additional} emphasis.
\end{em}

It is sometimes necessary to prevent \LaTeX\ from
breaking a line where it might otherwise do so.
This may be at a space, as between the ``Mr.''\ and
``Jones'' in
       ``Mr.~Jones'',        % ~ produces an unbreakable interword space.
or within a word---especially when the word is a
symbol like
       \mbox{\emph{itemnum}}
that makes little sense when hyphenated across
lines.

Footnotes\footnote{This is an example of a footnote.}
pose no problem.

\LaTeX\ is good at typesetting mathematical formulas
like
       \( x-3y + z = 7 \)
or
       \( a_{1} > x^{2n} + y^{2n} > x' \)
or
       \( \ip{A}{B} = \sum_{i} a_{i} b_{i} \).
The spaces you type in a formula are
ignored.  Remember that a letter like
       $x$                   % $ ... $  and  \( ... \)  are equivalent
is a formula when it denotes a mathematical
symbol, and it should be typed as one.

\section{Displayed Text}

Text is displayed by indenting it from the left
margin.  Quotations are commonly displayed.  There
are short quotations
\begin{quote}
   This is a short quotation.  It consists of a
   single paragraph of text.  See how it is formatted.
\end{quote}
and longer ones.
\begin{quotation}
   This is a longer quotation.  It consists of two
   paragraphs of text, neither of which are
   particularly interesting.

   This is the second paragraph of the quotation.  It
   is just as dull as the first paragraph.
\end{quotation}
Another frequently-displayed structure is a list.
The following is an example of an \emph{itemized}
list.
\begin{itemize}
   \item This is the first item of an itemized list.
         Each item in the list is marked with a ``tick''.
         You don't have to worry about what kind of tick
         mark is used.

   \item This is the second item of the list.  It
         contains another list nested inside it.  The inner
         list is an \emph{enumerated} list.
         \begin{enumerate}
            \item This is the first item of an enumerated
                  list that is nested within the itemized list.

            \item This is the second item of the inner list.
                  \LaTeX\ allows you to nest lists deeper than
                  you really should.
         \end{enumerate}
         This is the rest of the second item of the outer
         list.  It is no more interesting than any other
         part of the item.
   \item This is the third item of the list.
\end{itemize}
You can even display poetry.
\begin{verse}
   There is an environment
    for verse \\             % The \\ command separates lines
   Whose features some poets % within a stanza.
   will curse.

                             % One or more blank lines separate stanzas.

   For instead of making\\
   Them do \emph{all} line breaking, \\
   It allows them to put too many words on a line when they'd rather be
   forced to be terse.
\end{verse}

Mathematical formulas may also be displayed.  A
displayed formula
is
one-line long; multiline
formulas require special formatting instructions.
   \[  \ip{\Gamma}{\psi'} = x'' + y^{2} + z_{i}^{n}\]
Don't start a paragraph with a displayed equation,
nor make one a paragraph by itself.

\end{document}               % End of document.

\section{Bench Huge Section 1424}
% synthetic bench section 1424
% This is a sample LaTeX input file.  (Version of 12 August 2004.)
%
% A '%' character causes TeX to ignore all remaining text on the line,
% and is used for comments like this one.

\documentclass{article}      % Specifies the document class

                             % The preamble begins here.
\title{An Example Document}  % Declares the document's title.
\author{Leslie Lamport}      % Declares the author's name.
\date{January 21, 1994}      % Deleting this command produces today's date.

\newcommand{\ip}[2]{(#1, #2)}
                             % Defines \ip{arg1}{arg2} to mean
                             % (arg1, arg2).

%\newcommand{\ip}[2]{\langle #1 | #2\rangle}
                             % This is an alternative definition of
                             % \ip that is commented out.

\begin{document}             % End of preamble and beginning of text.

\maketitle                   % Produces the title.

This is an example input file.  Comparing it with
the output it generates can show you how to
produce a simple document of your own.

\section{Ordinary Text}      % Produces section heading.  Lower-level
                             % sections are begun with similar
                             % \subsection and \subsubsection commands.

The ends  of words and sentences are marked
  by   spaces. It  doesn't matter how many
spaces    you type; one is as good as 100.  The
end of   a line counts as a space.

One   or more   blank lines denote the  end
of  a paragraph.

Since any number of consecutive spaces are treated
like a single one, the formatting of the input
file makes no difference to
      \LaTeX,                % The \LaTeX command generates the LaTeX logo.
but it makes a difference to you.  When you use
\LaTeX, making your input file as easy to read
as possible will be a great help as you write
your document and when you change it.  This sample
file shows how you can add comments to your own input
file.

Because printing is different from typewriting,
there are a number of things that you have to do
differently when preparing an input file than if
you were just typing the document directly.
Quotation marks like
       ``this''
have to be handled specially, as do quotes within
quotes:
       ``\,`this'            % \, separates the double and single quote.
        is what I just
        wrote, not  `that'\,''.

Dashes come in three sizes: an
       intra-word
dash, a medium dash for number ranges like
       1--2,
and a punctuation
       dash---like
this.

A sentence-ending space should be larger than the
space between words within a sentence.  You
sometimes have to type special commands in
conjunction with punctuation characters to get
this right, as in the following sentence.
       Gnats, gnus, etc.\ all  % `\ ' makes an inter-word space.
       begin with G\@.         % \@ marks end-of-sentence punctuation.
You should check the spaces after periods when
reading your output to make sure you haven't
forgotten any special cases.  Generating an
ellipsis
       \ldots\               % `\ ' is needed after `\ldots' because TeX
                             % ignores spaces after command names like \ldots
                             % made from \ + letters.
                             %
                             % Note how a `%' character causes TeX to ignore
                             % the end of the input line, so these blank lines
                             % do not start a new paragraph.
                             %
with the right spacing around the periods requires
a special command.

\LaTeX\ interprets some common characters as
commands, so you must type special commands to
generate them.  These characters include the
following:
       \$ \& \% \# \{ and \}.

In printing, text is usually emphasized with an
       \emph{italic}
type style.

\begin{em}
   A long segment of text can also be emphasized
   in this way.  Text within such a segment can be
   given \emph{additional} emphasis.
\end{em}

It is sometimes necessary to prevent \LaTeX\ from
breaking a line where it might otherwise do so.
This may be at a space, as between the ``Mr.''\ and
``Jones'' in
       ``Mr.~Jones'',        % ~ produces an unbreakable interword space.
or within a word---especially when the word is a
symbol like
       \mbox{\emph{itemnum}}
that makes little sense when hyphenated across
lines.

Footnotes\footnote{This is an example of a footnote.}
pose no problem.

\LaTeX\ is good at typesetting mathematical formulas
like
       \( x-3y + z = 7 \)
or
       \( a_{1} > x^{2n} + y^{2n} > x' \)
or
       \( \ip{A}{B} = \sum_{i} a_{i} b_{i} \).
The spaces you type in a formula are
ignored.  Remember that a letter like
       $x$                   % $ ... $  and  \( ... \)  are equivalent
is a formula when it denotes a mathematical
symbol, and it should be typed as one.

\section{Displayed Text}

Text is displayed by indenting it from the left
margin.  Quotations are commonly displayed.  There
are short quotations
\begin{quote}
   This is a short quotation.  It consists of a
   single paragraph of text.  See how it is formatted.
\end{quote}
and longer ones.
\begin{quotation}
   This is a longer quotation.  It consists of two
   paragraphs of text, neither of which are
   particularly interesting.

   This is the second paragraph of the quotation.  It
   is just as dull as the first paragraph.
\end{quotation}
Another frequently-displayed structure is a list.
The following is an example of an \emph{itemized}
list.
\begin{itemize}
   \item This is the first item of an itemized list.
         Each item in the list is marked with a ``tick''.
         You don't have to worry about what kind of tick
         mark is used.

   \item This is the second item of the list.  It
         contains another list nested inside it.  The inner
         list is an \emph{enumerated} list.
         \begin{enumerate}
            \item This is the first item of an enumerated
                  list that is nested within the itemized list.

            \item This is the second item of the inner list.
                  \LaTeX\ allows you to nest lists deeper than
                  you really should.
         \end{enumerate}
         This is the rest of the second item of the outer
         list.  It is no more interesting than any other
         part of the item.
   \item This is the third item of the list.
\end{itemize}
You can even display poetry.
\begin{verse}
   There is an environment
    for verse \\             % The \\ command separates lines
   Whose features some poets % within a stanza.
   will curse.

                             % One or more blank lines separate stanzas.

   For instead of making\\
   Them do \emph{all} line breaking, \\
   It allows them to put too many words on a line when they'd rather be
   forced to be terse.
\end{verse}

Mathematical formulas may also be displayed.  A
displayed formula
is
one-line long; multiline
formulas require special formatting instructions.
   \[  \ip{\Gamma}{\psi'} = x'' + y^{2} + z_{i}^{n}\]
Don't start a paragraph with a displayed equation,
nor make one a paragraph by itself.

\end{document}               % End of document.

\section{Bench Huge Section 1425}
% synthetic bench section 1425
% This is a sample LaTeX input file.  (Version of 12 August 2004.)
%
% A '%' character causes TeX to ignore all remaining text on the line,
% and is used for comments like this one.

\documentclass{article}      % Specifies the document class

                             % The preamble begins here.
\title{An Example Document}  % Declares the document's title.
\author{Leslie Lamport}      % Declares the author's name.
\date{January 21, 1994}      % Deleting this command produces today's date.

\newcommand{\ip}[2]{(#1, #2)}
                             % Defines \ip{arg1}{arg2} to mean
                             % (arg1, arg2).

%\newcommand{\ip}[2]{\langle #1 | #2\rangle}
                             % This is an alternative definition of
                             % \ip that is commented out.

\begin{document}             % End of preamble and beginning of text.

\maketitle                   % Produces the title.

This is an example input file.  Comparing it with
the output it generates can show you how to
produce a simple document of your own.

\section{Ordinary Text}      % Produces section heading.  Lower-level
                             % sections are begun with similar
                             % \subsection and \subsubsection commands.

The ends  of words and sentences are marked
  by   spaces. It  doesn't matter how many
spaces    you type; one is as good as 100.  The
end of   a line counts as a space.

One   or more   blank lines denote the  end
of  a paragraph.

Since any number of consecutive spaces are treated
like a single one, the formatting of the input
file makes no difference to
      \LaTeX,                % The \LaTeX command generates the LaTeX logo.
but it makes a difference to you.  When you use
\LaTeX, making your input file as easy to read
as possible will be a great help as you write
your document and when you change it.  This sample
file shows how you can add comments to your own input
file.

Because printing is different from typewriting,
there are a number of things that you have to do
differently when preparing an input file than if
you were just typing the document directly.
Quotation marks like
       ``this''
have to be handled specially, as do quotes within
quotes:
       ``\,`this'            % \, separates the double and single quote.
        is what I just
        wrote, not  `that'\,''.

Dashes come in three sizes: an
       intra-word
dash, a medium dash for number ranges like
       1--2,
and a punctuation
       dash---like
this.

A sentence-ending space should be larger than the
space between words within a sentence.  You
sometimes have to type special commands in
conjunction with punctuation characters to get
this right, as in the following sentence.
       Gnats, gnus, etc.\ all  % `\ ' makes an inter-word space.
       begin with G\@.         % \@ marks end-of-sentence punctuation.
You should check the spaces after periods when
reading your output to make sure you haven't
forgotten any special cases.  Generating an
ellipsis
       \ldots\               % `\ ' is needed after `\ldots' because TeX
                             % ignores spaces after command names like \ldots
                             % made from \ + letters.
                             %
                             % Note how a `%' character causes TeX to ignore
                             % the end of the input line, so these blank lines
                             % do not start a new paragraph.
                             %
with the right spacing around the periods requires
a special command.

\LaTeX\ interprets some common characters as
commands, so you must type special commands to
generate them.  These characters include the
following:
       \$ \& \% \# \{ and \}.

In printing, text is usually emphasized with an
       \emph{italic}
type style.

\begin{em}
   A long segment of text can also be emphasized
   in this way.  Text within such a segment can be
   given \emph{additional} emphasis.
\end{em}

It is sometimes necessary to prevent \LaTeX\ from
breaking a line where it might otherwise do so.
This may be at a space, as between the ``Mr.''\ and
``Jones'' in
       ``Mr.~Jones'',        % ~ produces an unbreakable interword space.
or within a word---especially when the word is a
symbol like
       \mbox{\emph{itemnum}}
that makes little sense when hyphenated across
lines.

Footnotes\footnote{This is an example of a footnote.}
pose no problem.

\LaTeX\ is good at typesetting mathematical formulas
like
       \( x-3y + z = 7 \)
or
       \( a_{1} > x^{2n} + y^{2n} > x' \)
or
       \( \ip{A}{B} = \sum_{i} a_{i} b_{i} \).
The spaces you type in a formula are
ignored.  Remember that a letter like
       $x$                   % $ ... $  and  \( ... \)  are equivalent
is a formula when it denotes a mathematical
symbol, and it should be typed as one.

\section{Displayed Text}

Text is displayed by indenting it from the left
margin.  Quotations are commonly displayed.  There
are short quotations
\begin{quote}
   This is a short quotation.  It consists of a
   single paragraph of text.  See how it is formatted.
\end{quote}
and longer ones.
\begin{quotation}
   This is a longer quotation.  It consists of two
   paragraphs of text, neither of which are
   particularly interesting.

   This is the second paragraph of the quotation.  It
   is just as dull as the first paragraph.
\end{quotation}
Another frequently-displayed structure is a list.
The following is an example of an \emph{itemized}
list.
\begin{itemize}
   \item This is the first item of an itemized list.
         Each item in the list is marked with a ``tick''.
         You don't have to worry about what kind of tick
         mark is used.

   \item This is the second item of the list.  It
         contains another list nested inside it.  The inner
         list is an \emph{enumerated} list.
         \begin{enumerate}
            \item This is the first item of an enumerated
                  list that is nested within the itemized list.

            \item This is the second item of the inner list.
                  \LaTeX\ allows you to nest lists deeper than
                  you really should.
         \end{enumerate}
         This is the rest of the second item of the outer
         list.  It is no more interesting than any other
         part of the item.
   \item This is the third item of the list.
\end{itemize}
You can even display poetry.
\begin{verse}
   There is an environment
    for verse \\             % The \\ command separates lines
   Whose features some poets % within a stanza.
   will curse.

                             % One or more blank lines separate stanzas.

   For instead of making\\
   Them do \emph{all} line breaking, \\
   It allows them to put too many words on a line when they'd rather be
   forced to be terse.
\end{verse}

Mathematical formulas may also be displayed.  A
displayed formula
is
one-line long; multiline
formulas require special formatting instructions.
   \[  \ip{\Gamma}{\psi'} = x'' + y^{2} + z_{i}^{n}\]
Don't start a paragraph with a displayed equation,
nor make one a paragraph by itself.

\end{document}               % End of document.

\section{Bench Huge Section 1426}
% synthetic bench section 1426
% This is a sample LaTeX input file.  (Version of 12 August 2004.)
%
% A '%' character causes TeX to ignore all remaining text on the line,
% and is used for comments like this one.

\documentclass{article}      % Specifies the document class

                             % The preamble begins here.
\title{An Example Document}  % Declares the document's title.
\author{Leslie Lamport}      % Declares the author's name.
\date{January 21, 1994}      % Deleting this command produces today's date.

\newcommand{\ip}[2]{(#1, #2)}
                             % Defines \ip{arg1}{arg2} to mean
                             % (arg1, arg2).

%\newcommand{\ip}[2]{\langle #1 | #2\rangle}
                             % This is an alternative definition of
                             % \ip that is commented out.

\begin{document}             % End of preamble and beginning of text.

\maketitle                   % Produces the title.

This is an example input file.  Comparing it with
the output it generates can show you how to
produce a simple document of your own.

\section{Ordinary Text}      % Produces section heading.  Lower-level
                             % sections are begun with similar
                             % \subsection and \subsubsection commands.

The ends  of words and sentences are marked
  by   spaces. It  doesn't matter how many
spaces    you type; one is as good as 100.  The
end of   a line counts as a space.

One   or more   blank lines denote the  end
of  a paragraph.

Since any number of consecutive spaces are treated
like a single one, the formatting of the input
file makes no difference to
      \LaTeX,                % The \LaTeX command generates the LaTeX logo.
but it makes a difference to you.  When you use
\LaTeX, making your input file as easy to read
as possible will be a great help as you write
your document and when you change it.  This sample
file shows how you can add comments to your own input
file.

Because printing is different from typewriting,
there are a number of things that you have to do
differently when preparing an input file than if
you were just typing the document directly.
Quotation marks like
       ``this''
have to be handled specially, as do quotes within
quotes:
       ``\,`this'            % \, separates the double and single quote.
        is what I just
        wrote, not  `that'\,''.

Dashes come in three sizes: an
       intra-word
dash, a medium dash for number ranges like
       1--2,
and a punctuation
       dash---like
this.

A sentence-ending space should be larger than the
space between words within a sentence.  You
sometimes have to type special commands in
conjunction with punctuation characters to get
this right, as in the following sentence.
       Gnats, gnus, etc.\ all  % `\ ' makes an inter-word space.
       begin with G\@.         % \@ marks end-of-sentence punctuation.
You should check the spaces after periods when
reading your output to make sure you haven't
forgotten any special cases.  Generating an
ellipsis
       \ldots\               % `\ ' is needed after `\ldots' because TeX
                             % ignores spaces after command names like \ldots
                             % made from \ + letters.
                             %
                             % Note how a `%' character causes TeX to ignore
                             % the end of the input line, so these blank lines
                             % do not start a new paragraph.
                             %
with the right spacing around the periods requires
a special command.

\LaTeX\ interprets some common characters as
commands, so you must type special commands to
generate them.  These characters include the
following:
       \$ \& \% \# \{ and \}.

In printing, text is usually emphasized with an
       \emph{italic}
type style.

\begin{em}
   A long segment of text can also be emphasized
   in this way.  Text within such a segment can be
   given \emph{additional} emphasis.
\end{em}

It is sometimes necessary to prevent \LaTeX\ from
breaking a line where it might otherwise do so.
This may be at a space, as between the ``Mr.''\ and
``Jones'' in
       ``Mr.~Jones'',        % ~ produces an unbreakable interword space.
or within a word---especially when the word is a
symbol like
       \mbox{\emph{itemnum}}
that makes little sense when hyphenated across
lines.

Footnotes\footnote{This is an example of a footnote.}
pose no problem.

\LaTeX\ is good at typesetting mathematical formulas
like
       \( x-3y + z = 7 \)
or
       \( a_{1} > x^{2n} + y^{2n} > x' \)
or
       \( \ip{A}{B} = \sum_{i} a_{i} b_{i} \).
The spaces you type in a formula are
ignored.  Remember that a letter like
       $x$                   % $ ... $  and  \( ... \)  are equivalent
is a formula when it denotes a mathematical
symbol, and it should be typed as one.

\section{Displayed Text}

Text is displayed by indenting it from the left
margin.  Quotations are commonly displayed.  There
are short quotations
\begin{quote}
   This is a short quotation.  It consists of a
   single paragraph of text.  See how it is formatted.
\end{quote}
and longer ones.
\begin{quotation}
   This is a longer quotation.  It consists of two
   paragraphs of text, neither of which are
   particularly interesting.

   This is the second paragraph of the quotation.  It
   is just as dull as the first paragraph.
\end{quotation}
Another frequently-displayed structure is a list.
The following is an example of an \emph{itemized}
list.
\begin{itemize}
   \item This is the first item of an itemized list.
         Each item in the list is marked with a ``tick''.
         You don't have to worry about what kind of tick
         mark is used.

   \item This is the second item of the list.  It
         contains another list nested inside it.  The inner
         list is an \emph{enumerated} list.
         \begin{enumerate}
            \item This is the first item of an enumerated
                  list that is nested within the itemized list.

            \item This is the second item of the inner list.
                  \LaTeX\ allows you to nest lists deeper than
                  you really should.
         \end{enumerate}
         This is the rest of the second item of the outer
         list.  It is no more interesting than any other
         part of the item.
   \item This is the third item of the list.
\end{itemize}
You can even display poetry.
\begin{verse}
   There is an environment
    for verse \\             % The \\ command separates lines
   Whose features some poets % within a stanza.
   will curse.

                             % One or more blank lines separate stanzas.

   For instead of making\\
   Them do \emph{all} line breaking, \\
   It allows them to put too many words on a line when they'd rather be
   forced to be terse.
\end{verse}

Mathematical formulas may also be displayed.  A
displayed formula
is
one-line long; multiline
formulas require special formatting instructions.
   \[  \ip{\Gamma}{\psi'} = x'' + y^{2} + z_{i}^{n}\]
Don't start a paragraph with a displayed equation,
nor make one a paragraph by itself.

\end{document}               % End of document.

\section{Bench Huge Section 1427}
% synthetic bench section 1427
% This is a sample LaTeX input file.  (Version of 12 August 2004.)
%
% A '%' character causes TeX to ignore all remaining text on the line,
% and is used for comments like this one.

\documentclass{article}      % Specifies the document class

                             % The preamble begins here.
\title{An Example Document}  % Declares the document's title.
\author{Leslie Lamport}      % Declares the author's name.
\date{January 21, 1994}      % Deleting this command produces today's date.

\newcommand{\ip}[2]{(#1, #2)}
                             % Defines \ip{arg1}{arg2} to mean
                             % (arg1, arg2).

%\newcommand{\ip}[2]{\langle #1 | #2\rangle}
                             % This is an alternative definition of
                             % \ip that is commented out.

\begin{document}             % End of preamble and beginning of text.

\maketitle                   % Produces the title.

This is an example input file.  Comparing it with
the output it generates can show you how to
produce a simple document of your own.

\section{Ordinary Text}      % Produces section heading.  Lower-level
                             % sections are begun with similar
                             % \subsection and \subsubsection commands.

The ends  of words and sentences are marked
  by   spaces. It  doesn't matter how many
spaces    you type; one is as good as 100.  The
end of   a line counts as a space.

One   or more   blank lines denote the  end
of  a paragraph.

Since any number of consecutive spaces are treated
like a single one, the formatting of the input
file makes no difference to
      \LaTeX,                % The \LaTeX command generates the LaTeX logo.
but it makes a difference to you.  When you use
\LaTeX, making your input file as easy to read
as possible will be a great help as you write
your document and when you change it.  This sample
file shows how you can add comments to your own input
file.

Because printing is different from typewriting,
there are a number of things that you have to do
differently when preparing an input file than if
you were just typing the document directly.
Quotation marks like
       ``this''
have to be handled specially, as do quotes within
quotes:
       ``\,`this'            % \, separates the double and single quote.
        is what I just
        wrote, not  `that'\,''.

Dashes come in three sizes: an
       intra-word
dash, a medium dash for number ranges like
       1--2,
and a punctuation
       dash---like
this.

A sentence-ending space should be larger than the
space between words within a sentence.  You
sometimes have to type special commands in
conjunction with punctuation characters to get
this right, as in the following sentence.
       Gnats, gnus, etc.\ all  % `\ ' makes an inter-word space.
       begin with G\@.         % \@ marks end-of-sentence punctuation.
You should check the spaces after periods when
reading your output to make sure you haven't
forgotten any special cases.  Generating an
ellipsis
       \ldots\               % `\ ' is needed after `\ldots' because TeX
                             % ignores spaces after command names like \ldots
                             % made from \ + letters.
                             %
                             % Note how a `%' character causes TeX to ignore
                             % the end of the input line, so these blank lines
                             % do not start a new paragraph.
                             %
with the right spacing around the periods requires
a special command.

\LaTeX\ interprets some common characters as
commands, so you must type special commands to
generate them.  These characters include the
following:
       \$ \& \% \# \{ and \}.

In printing, text is usually emphasized with an
       \emph{italic}
type style.

\begin{em}
   A long segment of text can also be emphasized
   in this way.  Text within such a segment can be
   given \emph{additional} emphasis.
\end{em}

It is sometimes necessary to prevent \LaTeX\ from
breaking a line where it might otherwise do so.
This may be at a space, as between the ``Mr.''\ and
``Jones'' in
       ``Mr.~Jones'',        % ~ produces an unbreakable interword space.
or within a word---especially when the word is a
symbol like
       \mbox{\emph{itemnum}}
that makes little sense when hyphenated across
lines.

Footnotes\footnote{This is an example of a footnote.}
pose no problem.

\LaTeX\ is good at typesetting mathematical formulas
like
       \( x-3y + z = 7 \)
or
       \( a_{1} > x^{2n} + y^{2n} > x' \)
or
       \( \ip{A}{B} = \sum_{i} a_{i} b_{i} \).
The spaces you type in a formula are
ignored.  Remember that a letter like
       $x$                   % $ ... $  and  \( ... \)  are equivalent
is a formula when it denotes a mathematical
symbol, and it should be typed as one.

\section{Displayed Text}

Text is displayed by indenting it from the left
margin.  Quotations are commonly displayed.  There
are short quotations
\begin{quote}
   This is a short quotation.  It consists of a
   single paragraph of text.  See how it is formatted.
\end{quote}
and longer ones.
\begin{quotation}
   This is a longer quotation.  It consists of two
   paragraphs of text, neither of which are
   particularly interesting.

   This is the second paragraph of the quotation.  It
   is just as dull as the first paragraph.
\end{quotation}
Another frequently-displayed structure is a list.
The following is an example of an \emph{itemized}
list.
\begin{itemize}
   \item This is the first item of an itemized list.
         Each item in the list is marked with a ``tick''.
         You don't have to worry about what kind of tick
         mark is used.

   \item This is the second item of the list.  It
         contains another list nested inside it.  The inner
         list is an \emph{enumerated} list.
         \begin{enumerate}
            \item This is the first item of an enumerated
                  list that is nested within the itemized list.

            \item This is the second item of the inner list.
                  \LaTeX\ allows you to nest lists deeper than
                  you really should.
         \end{enumerate}
         This is the rest of the second item of the outer
         list.  It is no more interesting than any other
         part of the item.
   \item This is the third item of the list.
\end{itemize}
You can even display poetry.
\begin{verse}
   There is an environment
    for verse \\             % The \\ command separates lines
   Whose features some poets % within a stanza.
   will curse.

                             % One or more blank lines separate stanzas.

   For instead of making\\
   Them do \emph{all} line breaking, \\
   It allows them to put too many words on a line when they'd rather be
   forced to be terse.
\end{verse}

Mathematical formulas may also be displayed.  A
displayed formula
is
one-line long; multiline
formulas require special formatting instructions.
   \[  \ip{\Gamma}{\psi'} = x'' + y^{2} + z_{i}^{n}\]
Don't start a paragraph with a displayed equation,
nor make one a paragraph by itself.

\end{document}               % End of document.

\section{Bench Huge Section 1428}
% synthetic bench section 1428
% This is a sample LaTeX input file.  (Version of 12 August 2004.)
%
% A '%' character causes TeX to ignore all remaining text on the line,
% and is used for comments like this one.

\documentclass{article}      % Specifies the document class

                             % The preamble begins here.
\title{An Example Document}  % Declares the document's title.
\author{Leslie Lamport}      % Declares the author's name.
\date{January 21, 1994}      % Deleting this command produces today's date.

\newcommand{\ip}[2]{(#1, #2)}
                             % Defines \ip{arg1}{arg2} to mean
                             % (arg1, arg2).

%\newcommand{\ip}[2]{\langle #1 | #2\rangle}
                             % This is an alternative definition of
                             % \ip that is commented out.

\begin{document}             % End of preamble and beginning of text.

\maketitle                   % Produces the title.

This is an example input file.  Comparing it with
the output it generates can show you how to
produce a simple document of your own.

\section{Ordinary Text}      % Produces section heading.  Lower-level
                             % sections are begun with similar
                             % \subsection and \subsubsection commands.

The ends  of words and sentences are marked
  by   spaces. It  doesn't matter how many
spaces    you type; one is as good as 100.  The
end of   a line counts as a space.

One   or more   blank lines denote the  end
of  a paragraph.

Since any number of consecutive spaces are treated
like a single one, the formatting of the input
file makes no difference to
      \LaTeX,                % The \LaTeX command generates the LaTeX logo.
but it makes a difference to you.  When you use
\LaTeX, making your input file as easy to read
as possible will be a great help as you write
your document and when you change it.  This sample
file shows how you can add comments to your own input
file.

Because printing is different from typewriting,
there are a number of things that you have to do
differently when preparing an input file than if
you were just typing the document directly.
Quotation marks like
       ``this''
have to be handled specially, as do quotes within
quotes:
       ``\,`this'            % \, separates the double and single quote.
        is what I just
        wrote, not  `that'\,''.

Dashes come in three sizes: an
       intra-word
dash, a medium dash for number ranges like
       1--2,
and a punctuation
       dash---like
this.

A sentence-ending space should be larger than the
space between words within a sentence.  You
sometimes have to type special commands in
conjunction with punctuation characters to get
this right, as in the following sentence.
       Gnats, gnus, etc.\ all  % `\ ' makes an inter-word space.
       begin with G\@.         % \@ marks end-of-sentence punctuation.
You should check the spaces after periods when
reading your output to make sure you haven't
forgotten any special cases.  Generating an
ellipsis
       \ldots\               % `\ ' is needed after `\ldots' because TeX
                             % ignores spaces after command names like \ldots
                             % made from \ + letters.
                             %
                             % Note how a `%' character causes TeX to ignore
                             % the end of the input line, so these blank lines
                             % do not start a new paragraph.
                             %
with the right spacing around the periods requires
a special command.

\LaTeX\ interprets some common characters as
commands, so you must type special commands to
generate them.  These characters include the
following:
       \$ \& \% \# \{ and \}.

In printing, text is usually emphasized with an
       \emph{italic}
type style.

\begin{em}
   A long segment of text can also be emphasized
   in this way.  Text within such a segment can be
   given \emph{additional} emphasis.
\end{em}

It is sometimes necessary to prevent \LaTeX\ from
breaking a line where it might otherwise do so.
This may be at a space, as between the ``Mr.''\ and
``Jones'' in
       ``Mr.~Jones'',        % ~ produces an unbreakable interword space.
or within a word---especially when the word is a
symbol like
       \mbox{\emph{itemnum}}
that makes little sense when hyphenated across
lines.

Footnotes\footnote{This is an example of a footnote.}
pose no problem.

\LaTeX\ is good at typesetting mathematical formulas
like
       \( x-3y + z = 7 \)
or
       \( a_{1} > x^{2n} + y^{2n} > x' \)
or
       \( \ip{A}{B} = \sum_{i} a_{i} b_{i} \).
The spaces you type in a formula are
ignored.  Remember that a letter like
       $x$                   % $ ... $  and  \( ... \)  are equivalent
is a formula when it denotes a mathematical
symbol, and it should be typed as one.

\section{Displayed Text}

Text is displayed by indenting it from the left
margin.  Quotations are commonly displayed.  There
are short quotations
\begin{quote}
   This is a short quotation.  It consists of a
   single paragraph of text.  See how it is formatted.
\end{quote}
and longer ones.
\begin{quotation}
   This is a longer quotation.  It consists of two
   paragraphs of text, neither of which are
   particularly interesting.

   This is the second paragraph of the quotation.  It
   is just as dull as the first paragraph.
\end{quotation}
Another frequently-displayed structure is a list.
The following is an example of an \emph{itemized}
list.
\begin{itemize}
   \item This is the first item of an itemized list.
         Each item in the list is marked with a ``tick''.
         You don't have to worry about what kind of tick
         mark is used.

   \item This is the second item of the list.  It
         contains another list nested inside it.  The inner
         list is an \emph{enumerated} list.
         \begin{enumerate}
            \item This is the first item of an enumerated
                  list that is nested within the itemized list.

            \item This is the second item of the inner list.
                  \LaTeX\ allows you to nest lists deeper than
                  you really should.
         \end{enumerate}
         This is the rest of the second item of the outer
         list.  It is no more interesting than any other
         part of the item.
   \item This is the third item of the list.
\end{itemize}
You can even display poetry.
\begin{verse}
   There is an environment
    for verse \\             % The \\ command separates lines
   Whose features some poets % within a stanza.
   will curse.

                             % One or more blank lines separate stanzas.

   For instead of making\\
   Them do \emph{all} line breaking, \\
   It allows them to put too many words on a line when they'd rather be
   forced to be terse.
\end{verse}

Mathematical formulas may also be displayed.  A
displayed formula
is
one-line long; multiline
formulas require special formatting instructions.
   \[  \ip{\Gamma}{\psi'} = x'' + y^{2} + z_{i}^{n}\]
Don't start a paragraph with a displayed equation,
nor make one a paragraph by itself.

\end{document}               % End of document.

\section{Bench Huge Section 1429}
% synthetic bench section 1429
% This is a sample LaTeX input file.  (Version of 12 August 2004.)
%
% A '%' character causes TeX to ignore all remaining text on the line,
% and is used for comments like this one.

\documentclass{article}      % Specifies the document class

                             % The preamble begins here.
\title{An Example Document}  % Declares the document's title.
\author{Leslie Lamport}      % Declares the author's name.
\date{January 21, 1994}      % Deleting this command produces today's date.

\newcommand{\ip}[2]{(#1, #2)}
                             % Defines \ip{arg1}{arg2} to mean
                             % (arg1, arg2).

%\newcommand{\ip}[2]{\langle #1 | #2\rangle}
                             % This is an alternative definition of
                             % \ip that is commented out.

\begin{document}             % End of preamble and beginning of text.

\maketitle                   % Produces the title.

This is an example input file.  Comparing it with
the output it generates can show you how to
produce a simple document of your own.

\section{Ordinary Text}      % Produces section heading.  Lower-level
                             % sections are begun with similar
                             % \subsection and \subsubsection commands.

The ends  of words and sentences are marked
  by   spaces. It  doesn't matter how many
spaces    you type; one is as good as 100.  The
end of   a line counts as a space.

One   or more   blank lines denote the  end
of  a paragraph.

Since any number of consecutive spaces are treated
like a single one, the formatting of the input
file makes no difference to
      \LaTeX,                % The \LaTeX command generates the LaTeX logo.
but it makes a difference to you.  When you use
\LaTeX, making your input file as easy to read
as possible will be a great help as you write
your document and when you change it.  This sample
file shows how you can add comments to your own input
file.

Because printing is different from typewriting,
there are a number of things that you have to do
differently when preparing an input file than if
you were just typing the document directly.
Quotation marks like
       ``this''
have to be handled specially, as do quotes within
quotes:
       ``\,`this'            % \, separates the double and single quote.
        is what I just
        wrote, not  `that'\,''.

Dashes come in three sizes: an
       intra-word
dash, a medium dash for number ranges like
       1--2,
and a punctuation
       dash---like
this.

A sentence-ending space should be larger than the
space between words within a sentence.  You
sometimes have to type special commands in
conjunction with punctuation characters to get
this right, as in the following sentence.
       Gnats, gnus, etc.\ all  % `\ ' makes an inter-word space.
       begin with G\@.         % \@ marks end-of-sentence punctuation.
You should check the spaces after periods when
reading your output to make sure you haven't
forgotten any special cases.  Generating an
ellipsis
       \ldots\               % `\ ' is needed after `\ldots' because TeX
                             % ignores spaces after command names like \ldots
                             % made from \ + letters.
                             %
                             % Note how a `%' character causes TeX to ignore
                             % the end of the input line, so these blank lines
                             % do not start a new paragraph.
                             %
with the right spacing around the periods requires
a special command.

\LaTeX\ interprets some common characters as
commands, so you must type special commands to
generate them.  These characters include the
following:
       \$ \& \% \# \{ and \}.

In printing, text is usually emphasized with an
       \emph{italic}
type style.

\begin{em}
   A long segment of text can also be emphasized
   in this way.  Text within such a segment can be
   given \emph{additional} emphasis.
\end{em}

It is sometimes necessary to prevent \LaTeX\ from
breaking a line where it might otherwise do so.
This may be at a space, as between the ``Mr.''\ and
``Jones'' in
       ``Mr.~Jones'',        % ~ produces an unbreakable interword space.
or within a word---especially when the word is a
symbol like
       \mbox{\emph{itemnum}}
that makes little sense when hyphenated across
lines.

Footnotes\footnote{This is an example of a footnote.}
pose no problem.

\LaTeX\ is good at typesetting mathematical formulas
like
       \( x-3y + z = 7 \)
or
       \( a_{1} > x^{2n} + y^{2n} > x' \)
or
       \( \ip{A}{B} = \sum_{i} a_{i} b_{i} \).
The spaces you type in a formula are
ignored.  Remember that a letter like
       $x$                   % $ ... $  and  \( ... \)  are equivalent
is a formula when it denotes a mathematical
symbol, and it should be typed as one.

\section{Displayed Text}

Text is displayed by indenting it from the left
margin.  Quotations are commonly displayed.  There
are short quotations
\begin{quote}
   This is a short quotation.  It consists of a
   single paragraph of text.  See how it is formatted.
\end{quote}
and longer ones.
\begin{quotation}
   This is a longer quotation.  It consists of two
   paragraphs of text, neither of which are
   particularly interesting.

   This is the second paragraph of the quotation.  It
   is just as dull as the first paragraph.
\end{quotation}
Another frequently-displayed structure is a list.
The following is an example of an \emph{itemized}
list.
\begin{itemize}
   \item This is the first item of an itemized list.
         Each item in the list is marked with a ``tick''.
         You don't have to worry about what kind of tick
         mark is used.

   \item This is the second item of the list.  It
         contains another list nested inside it.  The inner
         list is an \emph{enumerated} list.
         \begin{enumerate}
            \item This is the first item of an enumerated
                  list that is nested within the itemized list.

            \item This is the second item of the inner list.
                  \LaTeX\ allows you to nest lists deeper than
                  you really should.
         \end{enumerate}
         This is the rest of the second item of the outer
         list.  It is no more interesting than any other
         part of the item.
   \item This is the third item of the list.
\end{itemize}
You can even display poetry.
\begin{verse}
   There is an environment
    for verse \\             % The \\ command separates lines
   Whose features some poets % within a stanza.
   will curse.

                             % One or more blank lines separate stanzas.

   For instead of making\\
   Them do \emph{all} line breaking, \\
   It allows them to put too many words on a line when they'd rather be
   forced to be terse.
\end{verse}

Mathematical formulas may also be displayed.  A
displayed formula
is
one-line long; multiline
formulas require special formatting instructions.
   \[  \ip{\Gamma}{\psi'} = x'' + y^{2} + z_{i}^{n}\]
Don't start a paragraph with a displayed equation,
nor make one a paragraph by itself.

\end{document}               % End of document.

\section{Bench Huge Section 1430}
% synthetic bench section 1430
% This is a sample LaTeX input file.  (Version of 12 August 2004.)
%
% A '%' character causes TeX to ignore all remaining text on the line,
% and is used for comments like this one.

\documentclass{article}      % Specifies the document class

                             % The preamble begins here.
\title{An Example Document}  % Declares the document's title.
\author{Leslie Lamport}      % Declares the author's name.
\date{January 21, 1994}      % Deleting this command produces today's date.

\newcommand{\ip}[2]{(#1, #2)}
                             % Defines \ip{arg1}{arg2} to mean
                             % (arg1, arg2).

%\newcommand{\ip}[2]{\langle #1 | #2\rangle}
                             % This is an alternative definition of
                             % \ip that is commented out.

\begin{document}             % End of preamble and beginning of text.

\maketitle                   % Produces the title.

This is an example input file.  Comparing it with
the output it generates can show you how to
produce a simple document of your own.

\section{Ordinary Text}      % Produces section heading.  Lower-level
                             % sections are begun with similar
                             % \subsection and \subsubsection commands.

The ends  of words and sentences are marked
  by   spaces. It  doesn't matter how many
spaces    you type; one is as good as 100.  The
end of   a line counts as a space.

One   or more   blank lines denote the  end
of  a paragraph.

Since any number of consecutive spaces are treated
like a single one, the formatting of the input
file makes no difference to
      \LaTeX,                % The \LaTeX command generates the LaTeX logo.
but it makes a difference to you.  When you use
\LaTeX, making your input file as easy to read
as possible will be a great help as you write
your document and when you change it.  This sample
file shows how you can add comments to your own input
file.

Because printing is different from typewriting,
there are a number of things that you have to do
differently when preparing an input file than if
you were just typing the document directly.
Quotation marks like
       ``this''
have to be handled specially, as do quotes within
quotes:
       ``\,`this'            % \, separates the double and single quote.
        is what I just
        wrote, not  `that'\,''.

Dashes come in three sizes: an
       intra-word
dash, a medium dash for number ranges like
       1--2,
and a punctuation
       dash---like
this.

A sentence-ending space should be larger than the
space between words within a sentence.  You
sometimes have to type special commands in
conjunction with punctuation characters to get
this right, as in the following sentence.
       Gnats, gnus, etc.\ all  % `\ ' makes an inter-word space.
       begin with G\@.         % \@ marks end-of-sentence punctuation.
You should check the spaces after periods when
reading your output to make sure you haven't
forgotten any special cases.  Generating an
ellipsis
       \ldots\               % `\ ' is needed after `\ldots' because TeX
                             % ignores spaces after command names like \ldots
                             % made from \ + letters.
                             %
                             % Note how a `%' character causes TeX to ignore
                             % the end of the input line, so these blank lines
                             % do not start a new paragraph.
                             %
with the right spacing around the periods requires
a special command.

\LaTeX\ interprets some common characters as
commands, so you must type special commands to
generate them.  These characters include the
following:
       \$ \& \% \# \{ and \}.

In printing, text is usually emphasized with an
       \emph{italic}
type style.

\begin{em}
   A long segment of text can also be emphasized
   in this way.  Text within such a segment can be
   given \emph{additional} emphasis.
\end{em}

It is sometimes necessary to prevent \LaTeX\ from
breaking a line where it might otherwise do so.
This may be at a space, as between the ``Mr.''\ and
``Jones'' in
       ``Mr.~Jones'',        % ~ produces an unbreakable interword space.
or within a word---especially when the word is a
symbol like
       \mbox{\emph{itemnum}}
that makes little sense when hyphenated across
lines.

Footnotes\footnote{This is an example of a footnote.}
pose no problem.

\LaTeX\ is good at typesetting mathematical formulas
like
       \( x-3y + z = 7 \)
or
       \( a_{1} > x^{2n} + y^{2n} > x' \)
or
       \( \ip{A}{B} = \sum_{i} a_{i} b_{i} \).
The spaces you type in a formula are
ignored.  Remember that a letter like
       $x$                   % $ ... $  and  \( ... \)  are equivalent
is a formula when it denotes a mathematical
symbol, and it should be typed as one.

\section{Displayed Text}

Text is displayed by indenting it from the left
margin.  Quotations are commonly displayed.  There
are short quotations
\begin{quote}
   This is a short quotation.  It consists of a
   single paragraph of text.  See how it is formatted.
\end{quote}
and longer ones.
\begin{quotation}
   This is a longer quotation.  It consists of two
   paragraphs of text, neither of which are
   particularly interesting.

   This is the second paragraph of the quotation.  It
   is just as dull as the first paragraph.
\end{quotation}
Another frequently-displayed structure is a list.
The following is an example of an \emph{itemized}
list.
\begin{itemize}
   \item This is the first item of an itemized list.
         Each item in the list is marked with a ``tick''.
         You don't have to worry about what kind of tick
         mark is used.

   \item This is the second item of the list.  It
         contains another list nested inside it.  The inner
         list is an \emph{enumerated} list.
         \begin{enumerate}
            \item This is the first item of an enumerated
                  list that is nested within the itemized list.

            \item This is the second item of the inner list.
                  \LaTeX\ allows you to nest lists deeper than
                  you really should.
         \end{enumerate}
         This is the rest of the second item of the outer
         list.  It is no more interesting than any other
         part of the item.
   \item This is the third item of the list.
\end{itemize}
You can even display poetry.
\begin{verse}
   There is an environment
    for verse \\             % The \\ command separates lines
   Whose features some poets % within a stanza.
   will curse.

                             % One or more blank lines separate stanzas.

   For instead of making\\
   Them do \emph{all} line breaking, \\
   It allows them to put too many words on a line when they'd rather be
   forced to be terse.
\end{verse}

Mathematical formulas may also be displayed.  A
displayed formula
is
one-line long; multiline
formulas require special formatting instructions.
   \[  \ip{\Gamma}{\psi'} = x'' + y^{2} + z_{i}^{n}\]
Don't start a paragraph with a displayed equation,
nor make one a paragraph by itself.

\end{document}               % End of document.

\section{Bench Huge Section 1431}
% synthetic bench section 1431
% This is a sample LaTeX input file.  (Version of 12 August 2004.)
%
% A '%' character causes TeX to ignore all remaining text on the line,
% and is used for comments like this one.

\documentclass{article}      % Specifies the document class

                             % The preamble begins here.
\title{An Example Document}  % Declares the document's title.
\author{Leslie Lamport}      % Declares the author's name.
\date{January 21, 1994}      % Deleting this command produces today's date.

\newcommand{\ip}[2]{(#1, #2)}
                             % Defines \ip{arg1}{arg2} to mean
                             % (arg1, arg2).

%\newcommand{\ip}[2]{\langle #1 | #2\rangle}
                             % This is an alternative definition of
                             % \ip that is commented out.

\begin{document}             % End of preamble and beginning of text.

\maketitle                   % Produces the title.

This is an example input file.  Comparing it with
the output it generates can show you how to
produce a simple document of your own.

\section{Ordinary Text}      % Produces section heading.  Lower-level
                             % sections are begun with similar
                             % \subsection and \subsubsection commands.

The ends  of words and sentences are marked
  by   spaces. It  doesn't matter how many
spaces    you type; one is as good as 100.  The
end of   a line counts as a space.

One   or more   blank lines denote the  end
of  a paragraph.

Since any number of consecutive spaces are treated
like a single one, the formatting of the input
file makes no difference to
      \LaTeX,                % The \LaTeX command generates the LaTeX logo.
but it makes a difference to you.  When you use
\LaTeX, making your input file as easy to read
as possible will be a great help as you write
your document and when you change it.  This sample
file shows how you can add comments to your own input
file.

Because printing is different from typewriting,
there are a number of things that you have to do
differently when preparing an input file than if
you were just typing the document directly.
Quotation marks like
       ``this''
have to be handled specially, as do quotes within
quotes:
       ``\,`this'            % \, separates the double and single quote.
        is what I just
        wrote, not  `that'\,''.

Dashes come in three sizes: an
       intra-word
dash, a medium dash for number ranges like
       1--2,
and a punctuation
       dash---like
this.

A sentence-ending space should be larger than the
space between words within a sentence.  You
sometimes have to type special commands in
conjunction with punctuation characters to get
this right, as in the following sentence.
       Gnats, gnus, etc.\ all  % `\ ' makes an inter-word space.
       begin with G\@.         % \@ marks end-of-sentence punctuation.
You should check the spaces after periods when
reading your output to make sure you haven't
forgotten any special cases.  Generating an
ellipsis
       \ldots\               % `\ ' is needed after `\ldots' because TeX
                             % ignores spaces after command names like \ldots
                             % made from \ + letters.
                             %
                             % Note how a `%' character causes TeX to ignore
                             % the end of the input line, so these blank lines
                             % do not start a new paragraph.
                             %
with the right spacing around the periods requires
a special command.

\LaTeX\ interprets some common characters as
commands, so you must type special commands to
generate them.  These characters include the
following:
       \$ \& \% \# \{ and \}.

In printing, text is usually emphasized with an
       \emph{italic}
type style.

\begin{em}
   A long segment of text can also be emphasized
   in this way.  Text within such a segment can be
   given \emph{additional} emphasis.
\end{em}

It is sometimes necessary to prevent \LaTeX\ from
breaking a line where it might otherwise do so.
This may be at a space, as between the ``Mr.''\ and
``Jones'' in
       ``Mr.~Jones'',        % ~ produces an unbreakable interword space.
or within a word---especially when the word is a
symbol like
       \mbox{\emph{itemnum}}
that makes little sense when hyphenated across
lines.

Footnotes\footnote{This is an example of a footnote.}
pose no problem.

\LaTeX\ is good at typesetting mathematical formulas
like
       \( x-3y + z = 7 \)
or
       \( a_{1} > x^{2n} + y^{2n} > x' \)
or
       \( \ip{A}{B} = \sum_{i} a_{i} b_{i} \).
The spaces you type in a formula are
ignored.  Remember that a letter like
       $x$                   % $ ... $  and  \( ... \)  are equivalent
is a formula when it denotes a mathematical
symbol, and it should be typed as one.

\section{Displayed Text}

Text is displayed by indenting it from the left
margin.  Quotations are commonly displayed.  There
are short quotations
\begin{quote}
   This is a short quotation.  It consists of a
   single paragraph of text.  See how it is formatted.
\end{quote}
and longer ones.
\begin{quotation}
   This is a longer quotation.  It consists of two
   paragraphs of text, neither of which are
   particularly interesting.

   This is the second paragraph of the quotation.  It
   is just as dull as the first paragraph.
\end{quotation}
Another frequently-displayed structure is a list.
The following is an example of an \emph{itemized}
list.
\begin{itemize}
   \item This is the first item of an itemized list.
         Each item in the list is marked with a ``tick''.
         You don't have to worry about what kind of tick
         mark is used.

   \item This is the second item of the list.  It
         contains another list nested inside it.  The inner
         list is an \emph{enumerated} list.
         \begin{enumerate}
            \item This is the first item of an enumerated
                  list that is nested within the itemized list.

            \item This is the second item of the inner list.
                  \LaTeX\ allows you to nest lists deeper than
                  you really should.
         \end{enumerate}
         This is the rest of the second item of the outer
         list.  It is no more interesting than any other
         part of the item.
   \item This is the third item of the list.
\end{itemize}
You can even display poetry.
\begin{verse}
   There is an environment
    for verse \\             % The \\ command separates lines
   Whose features some poets % within a stanza.
   will curse.

                             % One or more blank lines separate stanzas.

   For instead of making\\
   Them do \emph{all} line breaking, \\
   It allows them to put too many words on a line when they'd rather be
   forced to be terse.
\end{verse}

Mathematical formulas may also be displayed.  A
displayed formula
is
one-line long; multiline
formulas require special formatting instructions.
   \[  \ip{\Gamma}{\psi'} = x'' + y^{2} + z_{i}^{n}\]
Don't start a paragraph with a displayed equation,
nor make one a paragraph by itself.

\end{document}               % End of document.

\section{Bench Huge Section 1432}
% synthetic bench section 1432
% This is a sample LaTeX input file.  (Version of 12 August 2004.)
%
% A '%' character causes TeX to ignore all remaining text on the line,
% and is used for comments like this one.

\documentclass{article}      % Specifies the document class

                             % The preamble begins here.
\title{An Example Document}  % Declares the document's title.
\author{Leslie Lamport}      % Declares the author's name.
\date{January 21, 1994}      % Deleting this command produces today's date.

\newcommand{\ip}[2]{(#1, #2)}
                             % Defines \ip{arg1}{arg2} to mean
                             % (arg1, arg2).

%\newcommand{\ip}[2]{\langle #1 | #2\rangle}
                             % This is an alternative definition of
                             % \ip that is commented out.

\begin{document}             % End of preamble and beginning of text.

\maketitle                   % Produces the title.

This is an example input file.  Comparing it with
the output it generates can show you how to
produce a simple document of your own.

\section{Ordinary Text}      % Produces section heading.  Lower-level
                             % sections are begun with similar
                             % \subsection and \subsubsection commands.

The ends  of words and sentences are marked
  by   spaces. It  doesn't matter how many
spaces    you type; one is as good as 100.  The
end of   a line counts as a space.

One   or more   blank lines denote the  end
of  a paragraph.

Since any number of consecutive spaces are treated
like a single one, the formatting of the input
file makes no difference to
      \LaTeX,                % The \LaTeX command generates the LaTeX logo.
but it makes a difference to you.  When you use
\LaTeX, making your input file as easy to read
as possible will be a great help as you write
your document and when you change it.  This sample
file shows how you can add comments to your own input
file.

Because printing is different from typewriting,
there are a number of things that you have to do
differently when preparing an input file than if
you were just typing the document directly.
Quotation marks like
       ``this''
have to be handled specially, as do quotes within
quotes:
       ``\,`this'            % \, separates the double and single quote.
        is what I just
        wrote, not  `that'\,''.

Dashes come in three sizes: an
       intra-word
dash, a medium dash for number ranges like
       1--2,
and a punctuation
       dash---like
this.

A sentence-ending space should be larger than the
space between words within a sentence.  You
sometimes have to type special commands in
conjunction with punctuation characters to get
this right, as in the following sentence.
       Gnats, gnus, etc.\ all  % `\ ' makes an inter-word space.
       begin with G\@.         % \@ marks end-of-sentence punctuation.
You should check the spaces after periods when
reading your output to make sure you haven't
forgotten any special cases.  Generating an
ellipsis
       \ldots\               % `\ ' is needed after `\ldots' because TeX
                             % ignores spaces after command names like \ldots
                             % made from \ + letters.
                             %
                             % Note how a `%' character causes TeX to ignore
                             % the end of the input line, so these blank lines
                             % do not start a new paragraph.
                             %
with the right spacing around the periods requires
a special command.

\LaTeX\ interprets some common characters as
commands, so you must type special commands to
generate them.  These characters include the
following:
       \$ \& \% \# \{ and \}.

In printing, text is usually emphasized with an
       \emph{italic}
type style.

\begin{em}
   A long segment of text can also be emphasized
   in this way.  Text within such a segment can be
   given \emph{additional} emphasis.
\end{em}

It is sometimes necessary to prevent \LaTeX\ from
breaking a line where it might otherwise do so.
This may be at a space, as between the ``Mr.''\ and
``Jones'' in
       ``Mr.~Jones'',        % ~ produces an unbreakable interword space.
or within a word---especially when the word is a
symbol like
       \mbox{\emph{itemnum}}
that makes little sense when hyphenated across
lines.

Footnotes\footnote{This is an example of a footnote.}
pose no problem.

\LaTeX\ is good at typesetting mathematical formulas
like
       \( x-3y + z = 7 \)
or
       \( a_{1} > x^{2n} + y^{2n} > x' \)
or
       \( \ip{A}{B} = \sum_{i} a_{i} b_{i} \).
The spaces you type in a formula are
ignored.  Remember that a letter like
       $x$                   % $ ... $  and  \( ... \)  are equivalent
is a formula when it denotes a mathematical
symbol, and it should be typed as one.

\section{Displayed Text}

Text is displayed by indenting it from the left
margin.  Quotations are commonly displayed.  There
are short quotations
\begin{quote}
   This is a short quotation.  It consists of a
   single paragraph of text.  See how it is formatted.
\end{quote}
and longer ones.
\begin{quotation}
   This is a longer quotation.  It consists of two
   paragraphs of text, neither of which are
   particularly interesting.

   This is the second paragraph of the quotation.  It
   is just as dull as the first paragraph.
\end{quotation}
Another frequently-displayed structure is a list.
The following is an example of an \emph{itemized}
list.
\begin{itemize}
   \item This is the first item of an itemized list.
         Each item in the list is marked with a ``tick''.
         You don't have to worry about what kind of tick
         mark is used.

   \item This is the second item of the list.  It
         contains another list nested inside it.  The inner
         list is an \emph{enumerated} list.
         \begin{enumerate}
            \item This is the first item of an enumerated
                  list that is nested within the itemized list.

            \item This is the second item of the inner list.
                  \LaTeX\ allows you to nest lists deeper than
                  you really should.
         \end{enumerate}
         This is the rest of the second item of the outer
         list.  It is no more interesting than any other
         part of the item.
   \item This is the third item of the list.
\end{itemize}
You can even display poetry.
\begin{verse}
   There is an environment
    for verse \\             % The \\ command separates lines
   Whose features some poets % within a stanza.
   will curse.

                             % One or more blank lines separate stanzas.

   For instead of making\\
   Them do \emph{all} line breaking, \\
   It allows them to put too many words on a line when they'd rather be
   forced to be terse.
\end{verse}

Mathematical formulas may also be displayed.  A
displayed formula
is
one-line long; multiline
formulas require special formatting instructions.
   \[  \ip{\Gamma}{\psi'} = x'' + y^{2} + z_{i}^{n}\]
Don't start a paragraph with a displayed equation,
nor make one a paragraph by itself.

\end{document}               % End of document.

\section{Bench Huge Section 1433}
% synthetic bench section 1433
% This is a sample LaTeX input file.  (Version of 12 August 2004.)
%
% A '%' character causes TeX to ignore all remaining text on the line,
% and is used for comments like this one.

\documentclass{article}      % Specifies the document class

                             % The preamble begins here.
\title{An Example Document}  % Declares the document's title.
\author{Leslie Lamport}      % Declares the author's name.
\date{January 21, 1994}      % Deleting this command produces today's date.

\newcommand{\ip}[2]{(#1, #2)}
                             % Defines \ip{arg1}{arg2} to mean
                             % (arg1, arg2).

%\newcommand{\ip}[2]{\langle #1 | #2\rangle}
                             % This is an alternative definition of
                             % \ip that is commented out.

\begin{document}             % End of preamble and beginning of text.

\maketitle                   % Produces the title.

This is an example input file.  Comparing it with
the output it generates can show you how to
produce a simple document of your own.

\section{Ordinary Text}      % Produces section heading.  Lower-level
                             % sections are begun with similar
                             % \subsection and \subsubsection commands.

The ends  of words and sentences are marked
  by   spaces. It  doesn't matter how many
spaces    you type; one is as good as 100.  The
end of   a line counts as a space.

One   or more   blank lines denote the  end
of  a paragraph.

Since any number of consecutive spaces are treated
like a single one, the formatting of the input
file makes no difference to
      \LaTeX,                % The \LaTeX command generates the LaTeX logo.
but it makes a difference to you.  When you use
\LaTeX, making your input file as easy to read
as possible will be a great help as you write
your document and when you change it.  This sample
file shows how you can add comments to your own input
file.

Because printing is different from typewriting,
there are a number of things that you have to do
differently when preparing an input file than if
you were just typing the document directly.
Quotation marks like
       ``this''
have to be handled specially, as do quotes within
quotes:
       ``\,`this'            % \, separates the double and single quote.
        is what I just
        wrote, not  `that'\,''.

Dashes come in three sizes: an
       intra-word
dash, a medium dash for number ranges like
       1--2,
and a punctuation
       dash---like
this.

A sentence-ending space should be larger than the
space between words within a sentence.  You
sometimes have to type special commands in
conjunction with punctuation characters to get
this right, as in the following sentence.
       Gnats, gnus, etc.\ all  % `\ ' makes an inter-word space.
       begin with G\@.         % \@ marks end-of-sentence punctuation.
You should check the spaces after periods when
reading your output to make sure you haven't
forgotten any special cases.  Generating an
ellipsis
       \ldots\               % `\ ' is needed after `\ldots' because TeX
                             % ignores spaces after command names like \ldots
                             % made from \ + letters.
                             %
                             % Note how a `%' character causes TeX to ignore
                             % the end of the input line, so these blank lines
                             % do not start a new paragraph.
                             %
with the right spacing around the periods requires
a special command.

\LaTeX\ interprets some common characters as
commands, so you must type special commands to
generate them.  These characters include the
following:
       \$ \& \% \# \{ and \}.

In printing, text is usually emphasized with an
       \emph{italic}
type style.

\begin{em}
   A long segment of text can also be emphasized
   in this way.  Text within such a segment can be
   given \emph{additional} emphasis.
\end{em}

It is sometimes necessary to prevent \LaTeX\ from
breaking a line where it might otherwise do so.
This may be at a space, as between the ``Mr.''\ and
``Jones'' in
       ``Mr.~Jones'',        % ~ produces an unbreakable interword space.
or within a word---especially when the word is a
symbol like
       \mbox{\emph{itemnum}}
that makes little sense when hyphenated across
lines.

Footnotes\footnote{This is an example of a footnote.}
pose no problem.

\LaTeX\ is good at typesetting mathematical formulas
like
       \( x-3y + z = 7 \)
or
       \( a_{1} > x^{2n} + y^{2n} > x' \)
or
       \( \ip{A}{B} = \sum_{i} a_{i} b_{i} \).
The spaces you type in a formula are
ignored.  Remember that a letter like
       $x$                   % $ ... $  and  \( ... \)  are equivalent
is a formula when it denotes a mathematical
symbol, and it should be typed as one.

\section{Displayed Text}

Text is displayed by indenting it from the left
margin.  Quotations are commonly displayed.  There
are short quotations
\begin{quote}
   This is a short quotation.  It consists of a
   single paragraph of text.  See how it is formatted.
\end{quote}
and longer ones.
\begin{quotation}
   This is a longer quotation.  It consists of two
   paragraphs of text, neither of which are
   particularly interesting.

   This is the second paragraph of the quotation.  It
   is just as dull as the first paragraph.
\end{quotation}
Another frequently-displayed structure is a list.
The following is an example of an \emph{itemized}
list.
\begin{itemize}
   \item This is the first item of an itemized list.
         Each item in the list is marked with a ``tick''.
         You don't have to worry about what kind of tick
         mark is used.

   \item This is the second item of the list.  It
         contains another list nested inside it.  The inner
         list is an \emph{enumerated} list.
         \begin{enumerate}
            \item This is the first item of an enumerated
                  list that is nested within the itemized list.

            \item This is the second item of the inner list.
                  \LaTeX\ allows you to nest lists deeper than
                  you really should.
         \end{enumerate}
         This is the rest of the second item of the outer
         list.  It is no more interesting than any other
         part of the item.
   \item This is the third item of the list.
\end{itemize}
You can even display poetry.
\begin{verse}
   There is an environment
    for verse \\             % The \\ command separates lines
   Whose features some poets % within a stanza.
   will curse.

                             % One or more blank lines separate stanzas.

   For instead of making\\
   Them do \emph{all} line breaking, \\
   It allows them to put too many words on a line when they'd rather be
   forced to be terse.
\end{verse}

Mathematical formulas may also be displayed.  A
displayed formula
is
one-line long; multiline
formulas require special formatting instructions.
   \[  \ip{\Gamma}{\psi'} = x'' + y^{2} + z_{i}^{n}\]
Don't start a paragraph with a displayed equation,
nor make one a paragraph by itself.

\end{document}               % End of document.

\section{Bench Huge Section 1434}
% synthetic bench section 1434
% This is a sample LaTeX input file.  (Version of 12 August 2004.)
%
% A '%' character causes TeX to ignore all remaining text on the line,
% and is used for comments like this one.

\documentclass{article}      % Specifies the document class

                             % The preamble begins here.
\title{An Example Document}  % Declares the document's title.
\author{Leslie Lamport}      % Declares the author's name.
\date{January 21, 1994}      % Deleting this command produces today's date.

\newcommand{\ip}[2]{(#1, #2)}
                             % Defines \ip{arg1}{arg2} to mean
                             % (arg1, arg2).

%\newcommand{\ip}[2]{\langle #1 | #2\rangle}
                             % This is an alternative definition of
                             % \ip that is commented out.

\begin{document}             % End of preamble and beginning of text.

\maketitle                   % Produces the title.

This is an example input file.  Comparing it with
the output it generates can show you how to
produce a simple document of your own.

\section{Ordinary Text}      % Produces section heading.  Lower-level
                             % sections are begun with similar
                             % \subsection and \subsubsection commands.

The ends  of words and sentences are marked
  by   spaces. It  doesn't matter how many
spaces    you type; one is as good as 100.  The
end of   a line counts as a space.

One   or more   blank lines denote the  end
of  a paragraph.

Since any number of consecutive spaces are treated
like a single one, the formatting of the input
file makes no difference to
      \LaTeX,                % The \LaTeX command generates the LaTeX logo.
but it makes a difference to you.  When you use
\LaTeX, making your input file as easy to read
as possible will be a great help as you write
your document and when you change it.  This sample
file shows how you can add comments to your own input
file.

Because printing is different from typewriting,
there are a number of things that you have to do
differently when preparing an input file than if
you were just typing the document directly.
Quotation marks like
       ``this''
have to be handled specially, as do quotes within
quotes:
       ``\,`this'            % \, separates the double and single quote.
        is what I just
        wrote, not  `that'\,''.

Dashes come in three sizes: an
       intra-word
dash, a medium dash for number ranges like
       1--2,
and a punctuation
       dash---like
this.

A sentence-ending space should be larger than the
space between words within a sentence.  You
sometimes have to type special commands in
conjunction with punctuation characters to get
this right, as in the following sentence.
       Gnats, gnus, etc.\ all  % `\ ' makes an inter-word space.
       begin with G\@.         % \@ marks end-of-sentence punctuation.
You should check the spaces after periods when
reading your output to make sure you haven't
forgotten any special cases.  Generating an
ellipsis
       \ldots\               % `\ ' is needed after `\ldots' because TeX
                             % ignores spaces after command names like \ldots
                             % made from \ + letters.
                             %
                             % Note how a `%' character causes TeX to ignore
                             % the end of the input line, so these blank lines
                             % do not start a new paragraph.
                             %
with the right spacing around the periods requires
a special command.

\LaTeX\ interprets some common characters as
commands, so you must type special commands to
generate them.  These characters include the
following:
       \$ \& \% \# \{ and \}.

In printing, text is usually emphasized with an
       \emph{italic}
type style.

\begin{em}
   A long segment of text can also be emphasized
   in this way.  Text within such a segment can be
   given \emph{additional} emphasis.
\end{em}

It is sometimes necessary to prevent \LaTeX\ from
breaking a line where it might otherwise do so.
This may be at a space, as between the ``Mr.''\ and
``Jones'' in
       ``Mr.~Jones'',        % ~ produces an unbreakable interword space.
or within a word---especially when the word is a
symbol like
       \mbox{\emph{itemnum}}
that makes little sense when hyphenated across
lines.

Footnotes\footnote{This is an example of a footnote.}
pose no problem.

\LaTeX\ is good at typesetting mathematical formulas
like
       \( x-3y + z = 7 \)
or
       \( a_{1} > x^{2n} + y^{2n} > x' \)
or
       \( \ip{A}{B} = \sum_{i} a_{i} b_{i} \).
The spaces you type in a formula are
ignored.  Remember that a letter like
       $x$                   % $ ... $  and  \( ... \)  are equivalent
is a formula when it denotes a mathematical
symbol, and it should be typed as one.

\section{Displayed Text}

Text is displayed by indenting it from the left
margin.  Quotations are commonly displayed.  There
are short quotations
\begin{quote}
   This is a short quotation.  It consists of a
   single paragraph of text.  See how it is formatted.
\end{quote}
and longer ones.
\begin{quotation}
   This is a longer quotation.  It consists of two
   paragraphs of text, neither of which are
   particularly interesting.

   This is the second paragraph of the quotation.  It
   is just as dull as the first paragraph.
\end{quotation}
Another frequently-displayed structure is a list.
The following is an example of an \emph{itemized}
list.
\begin{itemize}
   \item This is the first item of an itemized list.
         Each item in the list is marked with a ``tick''.
         You don't have to worry about what kind of tick
         mark is used.

   \item This is the second item of the list.  It
         contains another list nested inside it.  The inner
         list is an \emph{enumerated} list.
         \begin{enumerate}
            \item This is the first item of an enumerated
                  list that is nested within the itemized list.

            \item This is the second item of the inner list.
                  \LaTeX\ allows you to nest lists deeper than
                  you really should.
         \end{enumerate}
         This is the rest of the second item of the outer
         list.  It is no more interesting than any other
         part of the item.
   \item This is the third item of the list.
\end{itemize}
You can even display poetry.
\begin{verse}
   There is an environment
    for verse \\             % The \\ command separates lines
   Whose features some poets % within a stanza.
   will curse.

                             % One or more blank lines separate stanzas.

   For instead of making\\
   Them do \emph{all} line breaking, \\
   It allows them to put too many words on a line when they'd rather be
   forced to be terse.
\end{verse}

Mathematical formulas may also be displayed.  A
displayed formula
is
one-line long; multiline
formulas require special formatting instructions.
   \[  \ip{\Gamma}{\psi'} = x'' + y^{2} + z_{i}^{n}\]
Don't start a paragraph with a displayed equation,
nor make one a paragraph by itself.

\end{document}               % End of document.

\section{Bench Huge Section 1435}
% synthetic bench section 1435
% This is a sample LaTeX input file.  (Version of 12 August 2004.)
%
% A '%' character causes TeX to ignore all remaining text on the line,
% and is used for comments like this one.

\documentclass{article}      % Specifies the document class

                             % The preamble begins here.
\title{An Example Document}  % Declares the document's title.
\author{Leslie Lamport}      % Declares the author's name.
\date{January 21, 1994}      % Deleting this command produces today's date.

\newcommand{\ip}[2]{(#1, #2)}
                             % Defines \ip{arg1}{arg2} to mean
                             % (arg1, arg2).

%\newcommand{\ip}[2]{\langle #1 | #2\rangle}
                             % This is an alternative definition of
                             % \ip that is commented out.

\begin{document}             % End of preamble and beginning of text.

\maketitle                   % Produces the title.

This is an example input file.  Comparing it with
the output it generates can show you how to
produce a simple document of your own.

\section{Ordinary Text}      % Produces section heading.  Lower-level
                             % sections are begun with similar
                             % \subsection and \subsubsection commands.

The ends  of words and sentences are marked
  by   spaces. It  doesn't matter how many
spaces    you type; one is as good as 100.  The
end of   a line counts as a space.

One   or more   blank lines denote the  end
of  a paragraph.

Since any number of consecutive spaces are treated
like a single one, the formatting of the input
file makes no difference to
      \LaTeX,                % The \LaTeX command generates the LaTeX logo.
but it makes a difference to you.  When you use
\LaTeX, making your input file as easy to read
as possible will be a great help as you write
your document and when you change it.  This sample
file shows how you can add comments to your own input
file.

Because printing is different from typewriting,
there are a number of things that you have to do
differently when preparing an input file than if
you were just typing the document directly.
Quotation marks like
       ``this''
have to be handled specially, as do quotes within
quotes:
       ``\,`this'            % \, separates the double and single quote.
        is what I just
        wrote, not  `that'\,''.

Dashes come in three sizes: an
       intra-word
dash, a medium dash for number ranges like
       1--2,
and a punctuation
       dash---like
this.

A sentence-ending space should be larger than the
space between words within a sentence.  You
sometimes have to type special commands in
conjunction with punctuation characters to get
this right, as in the following sentence.
       Gnats, gnus, etc.\ all  % `\ ' makes an inter-word space.
       begin with G\@.         % \@ marks end-of-sentence punctuation.
You should check the spaces after periods when
reading your output to make sure you haven't
forgotten any special cases.  Generating an
ellipsis
       \ldots\               % `\ ' is needed after `\ldots' because TeX
                             % ignores spaces after command names like \ldots
                             % made from \ + letters.
                             %
                             % Note how a `%' character causes TeX to ignore
                             % the end of the input line, so these blank lines
                             % do not start a new paragraph.
                             %
with the right spacing around the periods requires
a special command.

\LaTeX\ interprets some common characters as
commands, so you must type special commands to
generate them.  These characters include the
following:
       \$ \& \% \# \{ and \}.

In printing, text is usually emphasized with an
       \emph{italic}
type style.

\begin{em}
   A long segment of text can also be emphasized
   in this way.  Text within such a segment can be
   given \emph{additional} emphasis.
\end{em}

It is sometimes necessary to prevent \LaTeX\ from
breaking a line where it might otherwise do so.
This may be at a space, as between the ``Mr.''\ and
``Jones'' in
       ``Mr.~Jones'',        % ~ produces an unbreakable interword space.
or within a word---especially when the word is a
symbol like
       \mbox{\emph{itemnum}}
that makes little sense when hyphenated across
lines.

Footnotes\footnote{This is an example of a footnote.}
pose no problem.

\LaTeX\ is good at typesetting mathematical formulas
like
       \( x-3y + z = 7 \)
or
       \( a_{1} > x^{2n} + y^{2n} > x' \)
or
       \( \ip{A}{B} = \sum_{i} a_{i} b_{i} \).
The spaces you type in a formula are
ignored.  Remember that a letter like
       $x$                   % $ ... $  and  \( ... \)  are equivalent
is a formula when it denotes a mathematical
symbol, and it should be typed as one.

\section{Displayed Text}

Text is displayed by indenting it from the left
margin.  Quotations are commonly displayed.  There
are short quotations
\begin{quote}
   This is a short quotation.  It consists of a
   single paragraph of text.  See how it is formatted.
\end{quote}
and longer ones.
\begin{quotation}
   This is a longer quotation.  It consists of two
   paragraphs of text, neither of which are
   particularly interesting.

   This is the second paragraph of the quotation.  It
   is just as dull as the first paragraph.
\end{quotation}
Another frequently-displayed structure is a list.
The following is an example of an \emph{itemized}
list.
\begin{itemize}
   \item This is the first item of an itemized list.
         Each item in the list is marked with a ``tick''.
         You don't have to worry about what kind of tick
         mark is used.

   \item This is the second item of the list.  It
         contains another list nested inside it.  The inner
         list is an \emph{enumerated} list.
         \begin{enumerate}
            \item This is the first item of an enumerated
                  list that is nested within the itemized list.

            \item This is the second item of the inner list.
                  \LaTeX\ allows you to nest lists deeper than
                  you really should.
         \end{enumerate}
         This is the rest of the second item of the outer
         list.  It is no more interesting than any other
         part of the item.
   \item This is the third item of the list.
\end{itemize}
You can even display poetry.
\begin{verse}
   There is an environment
    for verse \\             % The \\ command separates lines
   Whose features some poets % within a stanza.
   will curse.

                             % One or more blank lines separate stanzas.

   For instead of making\\
   Them do \emph{all} line breaking, \\
   It allows them to put too many words on a line when they'd rather be
   forced to be terse.
\end{verse}

Mathematical formulas may also be displayed.  A
displayed formula
is
one-line long; multiline
formulas require special formatting instructions.
   \[  \ip{\Gamma}{\psi'} = x'' + y^{2} + z_{i}^{n}\]
Don't start a paragraph with a displayed equation,
nor make one a paragraph by itself.

\end{document}               % End of document.

\section{Bench Huge Section 1436}
% synthetic bench section 1436
% This is a sample LaTeX input file.  (Version of 12 August 2004.)
%
% A '%' character causes TeX to ignore all remaining text on the line,
% and is used for comments like this one.

\documentclass{article}      % Specifies the document class

                             % The preamble begins here.
\title{An Example Document}  % Declares the document's title.
\author{Leslie Lamport}      % Declares the author's name.
\date{January 21, 1994}      % Deleting this command produces today's date.

\newcommand{\ip}[2]{(#1, #2)}
                             % Defines \ip{arg1}{arg2} to mean
                             % (arg1, arg2).

%\newcommand{\ip}[2]{\langle #1 | #2\rangle}
                             % This is an alternative definition of
                             % \ip that is commented out.

\begin{document}             % End of preamble and beginning of text.

\maketitle                   % Produces the title.

This is an example input file.  Comparing it with
the output it generates can show you how to
produce a simple document of your own.

\section{Ordinary Text}      % Produces section heading.  Lower-level
                             % sections are begun with similar
                             % \subsection and \subsubsection commands.

The ends  of words and sentences are marked
  by   spaces. It  doesn't matter how many
spaces    you type; one is as good as 100.  The
end of   a line counts as a space.

One   or more   blank lines denote the  end
of  a paragraph.

Since any number of consecutive spaces are treated
like a single one, the formatting of the input
file makes no difference to
      \LaTeX,                % The \LaTeX command generates the LaTeX logo.
but it makes a difference to you.  When you use
\LaTeX, making your input file as easy to read
as possible will be a great help as you write
your document and when you change it.  This sample
file shows how you can add comments to your own input
file.

Because printing is different from typewriting,
there are a number of things that you have to do
differently when preparing an input file than if
you were just typing the document directly.
Quotation marks like
       ``this''
have to be handled specially, as do quotes within
quotes:
       ``\,`this'            % \, separates the double and single quote.
        is what I just
        wrote, not  `that'\,''.

Dashes come in three sizes: an
       intra-word
dash, a medium dash for number ranges like
       1--2,
and a punctuation
       dash---like
this.

A sentence-ending space should be larger than the
space between words within a sentence.  You
sometimes have to type special commands in
conjunction with punctuation characters to get
this right, as in the following sentence.
       Gnats, gnus, etc.\ all  % `\ ' makes an inter-word space.
       begin with G\@.         % \@ marks end-of-sentence punctuation.
You should check the spaces after periods when
reading your output to make sure you haven't
forgotten any special cases.  Generating an
ellipsis
       \ldots\               % `\ ' is needed after `\ldots' because TeX
                             % ignores spaces after command names like \ldots
                             % made from \ + letters.
                             %
                             % Note how a `%' character causes TeX to ignore
                             % the end of the input line, so these blank lines
                             % do not start a new paragraph.
                             %
with the right spacing around the periods requires
a special command.

\LaTeX\ interprets some common characters as
commands, so you must type special commands to
generate them.  These characters include the
following:
       \$ \& \% \# \{ and \}.

In printing, text is usually emphasized with an
       \emph{italic}
type style.

\begin{em}
   A long segment of text can also be emphasized
   in this way.  Text within such a segment can be
   given \emph{additional} emphasis.
\end{em}

It is sometimes necessary to prevent \LaTeX\ from
breaking a line where it might otherwise do so.
This may be at a space, as between the ``Mr.''\ and
``Jones'' in
       ``Mr.~Jones'',        % ~ produces an unbreakable interword space.
or within a word---especially when the word is a
symbol like
       \mbox{\emph{itemnum}}
that makes little sense when hyphenated across
lines.

Footnotes\footnote{This is an example of a footnote.}
pose no problem.

\LaTeX\ is good at typesetting mathematical formulas
like
       \( x-3y + z = 7 \)
or
       \( a_{1} > x^{2n} + y^{2n} > x' \)
or
       \( \ip{A}{B} = \sum_{i} a_{i} b_{i} \).
The spaces you type in a formula are
ignored.  Remember that a letter like
       $x$                   % $ ... $  and  \( ... \)  are equivalent
is a formula when it denotes a mathematical
symbol, and it should be typed as one.

\section{Displayed Text}

Text is displayed by indenting it from the left
margin.  Quotations are commonly displayed.  There
are short quotations
\begin{quote}
   This is a short quotation.  It consists of a
   single paragraph of text.  See how it is formatted.
\end{quote}
and longer ones.
\begin{quotation}
   This is a longer quotation.  It consists of two
   paragraphs of text, neither of which are
   particularly interesting.

   This is the second paragraph of the quotation.  It
   is just as dull as the first paragraph.
\end{quotation}
Another frequently-displayed structure is a list.
The following is an example of an \emph{itemized}
list.
\begin{itemize}
   \item This is the first item of an itemized list.
         Each item in the list is marked with a ``tick''.
         You don't have to worry about what kind of tick
         mark is used.

   \item This is the second item of the list.  It
         contains another list nested inside it.  The inner
         list is an \emph{enumerated} list.
         \begin{enumerate}
            \item This is the first item of an enumerated
                  list that is nested within the itemized list.

            \item This is the second item of the inner list.
                  \LaTeX\ allows you to nest lists deeper than
                  you really should.
         \end{enumerate}
         This is the rest of the second item of the outer
         list.  It is no more interesting than any other
         part of the item.
   \item This is the third item of the list.
\end{itemize}
You can even display poetry.
\begin{verse}
   There is an environment
    for verse \\             % The \\ command separates lines
   Whose features some poets % within a stanza.
   will curse.

                             % One or more blank lines separate stanzas.

   For instead of making\\
   Them do \emph{all} line breaking, \\
   It allows them to put too many words on a line when they'd rather be
   forced to be terse.
\end{verse}

Mathematical formulas may also be displayed.  A
displayed formula
is
one-line long; multiline
formulas require special formatting instructions.
   \[  \ip{\Gamma}{\psi'} = x'' + y^{2} + z_{i}^{n}\]
Don't start a paragraph with a displayed equation,
nor make one a paragraph by itself.

\end{document}               % End of document.

\section{Bench Huge Section 1437}
% synthetic bench section 1437
% This is a sample LaTeX input file.  (Version of 12 August 2004.)
%
% A '%' character causes TeX to ignore all remaining text on the line,
% and is used for comments like this one.

\documentclass{article}      % Specifies the document class

                             % The preamble begins here.
\title{An Example Document}  % Declares the document's title.
\author{Leslie Lamport}      % Declares the author's name.
\date{January 21, 1994}      % Deleting this command produces today's date.

\newcommand{\ip}[2]{(#1, #2)}
                             % Defines \ip{arg1}{arg2} to mean
                             % (arg1, arg2).

%\newcommand{\ip}[2]{\langle #1 | #2\rangle}
                             % This is an alternative definition of
                             % \ip that is commented out.

\begin{document}             % End of preamble and beginning of text.

\maketitle                   % Produces the title.

This is an example input file.  Comparing it with
the output it generates can show you how to
produce a simple document of your own.

\section{Ordinary Text}      % Produces section heading.  Lower-level
                             % sections are begun with similar
                             % \subsection and \subsubsection commands.

The ends  of words and sentences are marked
  by   spaces. It  doesn't matter how many
spaces    you type; one is as good as 100.  The
end of   a line counts as a space.

One   or more   blank lines denote the  end
of  a paragraph.

Since any number of consecutive spaces are treated
like a single one, the formatting of the input
file makes no difference to
      \LaTeX,                % The \LaTeX command generates the LaTeX logo.
but it makes a difference to you.  When you use
\LaTeX, making your input file as easy to read
as possible will be a great help as you write
your document and when you change it.  This sample
file shows how you can add comments to your own input
file.

Because printing is different from typewriting,
there are a number of things that you have to do
differently when preparing an input file than if
you were just typing the document directly.
Quotation marks like
       ``this''
have to be handled specially, as do quotes within
quotes:
       ``\,`this'            % \, separates the double and single quote.
        is what I just
        wrote, not  `that'\,''.

Dashes come in three sizes: an
       intra-word
dash, a medium dash for number ranges like
       1--2,
and a punctuation
       dash---like
this.

A sentence-ending space should be larger than the
space between words within a sentence.  You
sometimes have to type special commands in
conjunction with punctuation characters to get
this right, as in the following sentence.
       Gnats, gnus, etc.\ all  % `\ ' makes an inter-word space.
       begin with G\@.         % \@ marks end-of-sentence punctuation.
You should check the spaces after periods when
reading your output to make sure you haven't
forgotten any special cases.  Generating an
ellipsis
       \ldots\               % `\ ' is needed after `\ldots' because TeX
                             % ignores spaces after command names like \ldots
                             % made from \ + letters.
                             %
                             % Note how a `%' character causes TeX to ignore
                             % the end of the input line, so these blank lines
                             % do not start a new paragraph.
                             %
with the right spacing around the periods requires
a special command.

\LaTeX\ interprets some common characters as
commands, so you must type special commands to
generate them.  These characters include the
following:
       \$ \& \% \# \{ and \}.

In printing, text is usually emphasized with an
       \emph{italic}
type style.

\begin{em}
   A long segment of text can also be emphasized
   in this way.  Text within such a segment can be
   given \emph{additional} emphasis.
\end{em}

It is sometimes necessary to prevent \LaTeX\ from
breaking a line where it might otherwise do so.
This may be at a space, as between the ``Mr.''\ and
``Jones'' in
       ``Mr.~Jones'',        % ~ produces an unbreakable interword space.
or within a word---especially when the word is a
symbol like
       \mbox{\emph{itemnum}}
that makes little sense when hyphenated across
lines.

Footnotes\footnote{This is an example of a footnote.}
pose no problem.

\LaTeX\ is good at typesetting mathematical formulas
like
       \( x-3y + z = 7 \)
or
       \( a_{1} > x^{2n} + y^{2n} > x' \)
or
       \( \ip{A}{B} = \sum_{i} a_{i} b_{i} \).
The spaces you type in a formula are
ignored.  Remember that a letter like
       $x$                   % $ ... $  and  \( ... \)  are equivalent
is a formula when it denotes a mathematical
symbol, and it should be typed as one.

\section{Displayed Text}

Text is displayed by indenting it from the left
margin.  Quotations are commonly displayed.  There
are short quotations
\begin{quote}
   This is a short quotation.  It consists of a
   single paragraph of text.  See how it is formatted.
\end{quote}
and longer ones.
\begin{quotation}
   This is a longer quotation.  It consists of two
   paragraphs of text, neither of which are
   particularly interesting.

   This is the second paragraph of the quotation.  It
   is just as dull as the first paragraph.
\end{quotation}
Another frequently-displayed structure is a list.
The following is an example of an \emph{itemized}
list.
\begin{itemize}
   \item This is the first item of an itemized list.
         Each item in the list is marked with a ``tick''.
         You don't have to worry about what kind of tick
         mark is used.

   \item This is the second item of the list.  It
         contains another list nested inside it.  The inner
         list is an \emph{enumerated} list.
         \begin{enumerate}
            \item This is the first item of an enumerated
                  list that is nested within the itemized list.

            \item This is the second item of the inner list.
                  \LaTeX\ allows you to nest lists deeper than
                  you really should.
         \end{enumerate}
         This is the rest of the second item of the outer
         list.  It is no more interesting than any other
         part of the item.
   \item This is the third item of the list.
\end{itemize}
You can even display poetry.
\begin{verse}
   There is an environment
    for verse \\             % The \\ command separates lines
   Whose features some poets % within a stanza.
   will curse.

                             % One or more blank lines separate stanzas.

   For instead of making\\
   Them do \emph{all} line breaking, \\
   It allows them to put too many words on a line when they'd rather be
   forced to be terse.
\end{verse}

Mathematical formulas may also be displayed.  A
displayed formula
is
one-line long; multiline
formulas require special formatting instructions.
   \[  \ip{\Gamma}{\psi'} = x'' + y^{2} + z_{i}^{n}\]
Don't start a paragraph with a displayed equation,
nor make one a paragraph by itself.

\end{document}               % End of document.

\section{Bench Huge Section 1438}
% synthetic bench section 1438
% This is a sample LaTeX input file.  (Version of 12 August 2004.)
%
% A '%' character causes TeX to ignore all remaining text on the line,
% and is used for comments like this one.

\documentclass{article}      % Specifies the document class

                             % The preamble begins here.
\title{An Example Document}  % Declares the document's title.
\author{Leslie Lamport}      % Declares the author's name.
\date{January 21, 1994}      % Deleting this command produces today's date.

\newcommand{\ip}[2]{(#1, #2)}
                             % Defines \ip{arg1}{arg2} to mean
                             % (arg1, arg2).

%\newcommand{\ip}[2]{\langle #1 | #2\rangle}
                             % This is an alternative definition of
                             % \ip that is commented out.

\begin{document}             % End of preamble and beginning of text.

\maketitle                   % Produces the title.

This is an example input file.  Comparing it with
the output it generates can show you how to
produce a simple document of your own.

\section{Ordinary Text}      % Produces section heading.  Lower-level
                             % sections are begun with similar
                             % \subsection and \subsubsection commands.

The ends  of words and sentences are marked
  by   spaces. It  doesn't matter how many
spaces    you type; one is as good as 100.  The
end of   a line counts as a space.

One   or more   blank lines denote the  end
of  a paragraph.

Since any number of consecutive spaces are treated
like a single one, the formatting of the input
file makes no difference to
      \LaTeX,                % The \LaTeX command generates the LaTeX logo.
but it makes a difference to you.  When you use
\LaTeX, making your input file as easy to read
as possible will be a great help as you write
your document and when you change it.  This sample
file shows how you can add comments to your own input
file.

Because printing is different from typewriting,
there are a number of things that you have to do
differently when preparing an input file than if
you were just typing the document directly.
Quotation marks like
       ``this''
have to be handled specially, as do quotes within
quotes:
       ``\,`this'            % \, separates the double and single quote.
        is what I just
        wrote, not  `that'\,''.

Dashes come in three sizes: an
       intra-word
dash, a medium dash for number ranges like
       1--2,
and a punctuation
       dash---like
this.

A sentence-ending space should be larger than the
space between words within a sentence.  You
sometimes have to type special commands in
conjunction with punctuation characters to get
this right, as in the following sentence.
       Gnats, gnus, etc.\ all  % `\ ' makes an inter-word space.
       begin with G\@.         % \@ marks end-of-sentence punctuation.
You should check the spaces after periods when
reading your output to make sure you haven't
forgotten any special cases.  Generating an
ellipsis
       \ldots\               % `\ ' is needed after `\ldots' because TeX
                             % ignores spaces after command names like \ldots
                             % made from \ + letters.
                             %
                             % Note how a `%' character causes TeX to ignore
                             % the end of the input line, so these blank lines
                             % do not start a new paragraph.
                             %
with the right spacing around the periods requires
a special command.

\LaTeX\ interprets some common characters as
commands, so you must type special commands to
generate them.  These characters include the
following:
       \$ \& \% \# \{ and \}.

In printing, text is usually emphasized with an
       \emph{italic}
type style.

\begin{em}
   A long segment of text can also be emphasized
   in this way.  Text within such a segment can be
   given \emph{additional} emphasis.
\end{em}

It is sometimes necessary to prevent \LaTeX\ from
breaking a line where it might otherwise do so.
This may be at a space, as between the ``Mr.''\ and
``Jones'' in
       ``Mr.~Jones'',        % ~ produces an unbreakable interword space.
or within a word---especially when the word is a
symbol like
       \mbox{\emph{itemnum}}
that makes little sense when hyphenated across
lines.

Footnotes\footnote{This is an example of a footnote.}
pose no problem.

\LaTeX\ is good at typesetting mathematical formulas
like
       \( x-3y + z = 7 \)
or
       \( a_{1} > x^{2n} + y^{2n} > x' \)
or
       \( \ip{A}{B} = \sum_{i} a_{i} b_{i} \).
The spaces you type in a formula are
ignored.  Remember that a letter like
       $x$                   % $ ... $  and  \( ... \)  are equivalent
is a formula when it denotes a mathematical
symbol, and it should be typed as one.

\section{Displayed Text}

Text is displayed by indenting it from the left
margin.  Quotations are commonly displayed.  There
are short quotations
\begin{quote}
   This is a short quotation.  It consists of a
   single paragraph of text.  See how it is formatted.
\end{quote}
and longer ones.
\begin{quotation}
   This is a longer quotation.  It consists of two
   paragraphs of text, neither of which are
   particularly interesting.

   This is the second paragraph of the quotation.  It
   is just as dull as the first paragraph.
\end{quotation}
Another frequently-displayed structure is a list.
The following is an example of an \emph{itemized}
list.
\begin{itemize}
   \item This is the first item of an itemized list.
         Each item in the list is marked with a ``tick''.
         You don't have to worry about what kind of tick
         mark is used.

   \item This is the second item of the list.  It
         contains another list nested inside it.  The inner
         list is an \emph{enumerated} list.
         \begin{enumerate}
            \item This is the first item of an enumerated
                  list that is nested within the itemized list.

            \item This is the second item of the inner list.
                  \LaTeX\ allows you to nest lists deeper than
                  you really should.
         \end{enumerate}
         This is the rest of the second item of the outer
         list.  It is no more interesting than any other
         part of the item.
   \item This is the third item of the list.
\end{itemize}
You can even display poetry.
\begin{verse}
   There is an environment
    for verse \\             % The \\ command separates lines
   Whose features some poets % within a stanza.
   will curse.

                             % One or more blank lines separate stanzas.

   For instead of making\\
   Them do \emph{all} line breaking, \\
   It allows them to put too many words on a line when they'd rather be
   forced to be terse.
\end{verse}

Mathematical formulas may also be displayed.  A
displayed formula
is
one-line long; multiline
formulas require special formatting instructions.
   \[  \ip{\Gamma}{\psi'} = x'' + y^{2} + z_{i}^{n}\]
Don't start a paragraph with a displayed equation,
nor make one a paragraph by itself.

\end{document}               % End of document.

\section{Bench Huge Section 1439}
% synthetic bench section 1439
% This is a sample LaTeX input file.  (Version of 12 August 2004.)
%
% A '%' character causes TeX to ignore all remaining text on the line,
% and is used for comments like this one.

\documentclass{article}      % Specifies the document class

                             % The preamble begins here.
\title{An Example Document}  % Declares the document's title.
\author{Leslie Lamport}      % Declares the author's name.
\date{January 21, 1994}      % Deleting this command produces today's date.

\newcommand{\ip}[2]{(#1, #2)}
                             % Defines \ip{arg1}{arg2} to mean
                             % (arg1, arg2).

%\newcommand{\ip}[2]{\langle #1 | #2\rangle}
                             % This is an alternative definition of
                             % \ip that is commented out.

\begin{document}             % End of preamble and beginning of text.

\maketitle                   % Produces the title.

This is an example input file.  Comparing it with
the output it generates can show you how to
produce a simple document of your own.

\section{Ordinary Text}      % Produces section heading.  Lower-level
                             % sections are begun with similar
                             % \subsection and \subsubsection commands.

The ends  of words and sentences are marked
  by   spaces. It  doesn't matter how many
spaces    you type; one is as good as 100.  The
end of   a line counts as a space.

One   or more   blank lines denote the  end
of  a paragraph.

Since any number of consecutive spaces are treated
like a single one, the formatting of the input
file makes no difference to
      \LaTeX,                % The \LaTeX command generates the LaTeX logo.
but it makes a difference to you.  When you use
\LaTeX, making your input file as easy to read
as possible will be a great help as you write
your document and when you change it.  This sample
file shows how you can add comments to your own input
file.

Because printing is different from typewriting,
there are a number of things that you have to do
differently when preparing an input file than if
you were just typing the document directly.
Quotation marks like
       ``this''
have to be handled specially, as do quotes within
quotes:
       ``\,`this'            % \, separates the double and single quote.
        is what I just
        wrote, not  `that'\,''.

Dashes come in three sizes: an
       intra-word
dash, a medium dash for number ranges like
       1--2,
and a punctuation
       dash---like
this.

A sentence-ending space should be larger than the
space between words within a sentence.  You
sometimes have to type special commands in
conjunction with punctuation characters to get
this right, as in the following sentence.
       Gnats, gnus, etc.\ all  % `\ ' makes an inter-word space.
       begin with G\@.         % \@ marks end-of-sentence punctuation.
You should check the spaces after periods when
reading your output to make sure you haven't
forgotten any special cases.  Generating an
ellipsis
       \ldots\               % `\ ' is needed after `\ldots' because TeX
                             % ignores spaces after command names like \ldots
                             % made from \ + letters.
                             %
                             % Note how a `%' character causes TeX to ignore
                             % the end of the input line, so these blank lines
                             % do not start a new paragraph.
                             %
with the right spacing around the periods requires
a special command.

\LaTeX\ interprets some common characters as
commands, so you must type special commands to
generate them.  These characters include the
following:
       \$ \& \% \# \{ and \}.

In printing, text is usually emphasized with an
       \emph{italic}
type style.

\begin{em}
   A long segment of text can also be emphasized
   in this way.  Text within such a segment can be
   given \emph{additional} emphasis.
\end{em}

It is sometimes necessary to prevent \LaTeX\ from
breaking a line where it might otherwise do so.
This may be at a space, as between the ``Mr.''\ and
``Jones'' in
       ``Mr.~Jones'',        % ~ produces an unbreakable interword space.
or within a word---especially when the word is a
symbol like
       \mbox{\emph{itemnum}}
that makes little sense when hyphenated across
lines.

Footnotes\footnote{This is an example of a footnote.}
pose no problem.

\LaTeX\ is good at typesetting mathematical formulas
like
       \( x-3y + z = 7 \)
or
       \( a_{1} > x^{2n} + y^{2n} > x' \)
or
       \( \ip{A}{B} = \sum_{i} a_{i} b_{i} \).
The spaces you type in a formula are
ignored.  Remember that a letter like
       $x$                   % $ ... $  and  \( ... \)  are equivalent
is a formula when it denotes a mathematical
symbol, and it should be typed as one.

\section{Displayed Text}

Text is displayed by indenting it from the left
margin.  Quotations are commonly displayed.  There
are short quotations
\begin{quote}
   This is a short quotation.  It consists of a
   single paragraph of text.  See how it is formatted.
\end{quote}
and longer ones.
\begin{quotation}
   This is a longer quotation.  It consists of two
   paragraphs of text, neither of which are
   particularly interesting.

   This is the second paragraph of the quotation.  It
   is just as dull as the first paragraph.
\end{quotation}
Another frequently-displayed structure is a list.
The following is an example of an \emph{itemized}
list.
\begin{itemize}
   \item This is the first item of an itemized list.
         Each item in the list is marked with a ``tick''.
         You don't have to worry about what kind of tick
         mark is used.

   \item This is the second item of the list.  It
         contains another list nested inside it.  The inner
         list is an \emph{enumerated} list.
         \begin{enumerate}
            \item This is the first item of an enumerated
                  list that is nested within the itemized list.

            \item This is the second item of the inner list.
                  \LaTeX\ allows you to nest lists deeper than
                  you really should.
         \end{enumerate}
         This is the rest of the second item of the outer
         list.  It is no more interesting than any other
         part of the item.
   \item This is the third item of the list.
\end{itemize}
You can even display poetry.
\begin{verse}
   There is an environment
    for verse \\             % The \\ command separates lines
   Whose features some poets % within a stanza.
   will curse.

                             % One or more blank lines separate stanzas.

   For instead of making\\
   Them do \emph{all} line breaking, \\
   It allows them to put too many words on a line when they'd rather be
   forced to be terse.
\end{verse}

Mathematical formulas may also be displayed.  A
displayed formula
is
one-line long; multiline
formulas require special formatting instructions.
   \[  \ip{\Gamma}{\psi'} = x'' + y^{2} + z_{i}^{n}\]
Don't start a paragraph with a displayed equation,
nor make one a paragraph by itself.

\end{document}               % End of document.

\section{Bench Huge Section 1440}
% synthetic bench section 1440
% This is a sample LaTeX input file.  (Version of 12 August 2004.)
%
% A '%' character causes TeX to ignore all remaining text on the line,
% and is used for comments like this one.

\documentclass{article}      % Specifies the document class

                             % The preamble begins here.
\title{An Example Document}  % Declares the document's title.
\author{Leslie Lamport}      % Declares the author's name.
\date{January 21, 1994}      % Deleting this command produces today's date.

\newcommand{\ip}[2]{(#1, #2)}
                             % Defines \ip{arg1}{arg2} to mean
                             % (arg1, arg2).

%\newcommand{\ip}[2]{\langle #1 | #2\rangle}
                             % This is an alternative definition of
                             % \ip that is commented out.

\begin{document}             % End of preamble and beginning of text.

\maketitle                   % Produces the title.

This is an example input file.  Comparing it with
the output it generates can show you how to
produce a simple document of your own.

\section{Ordinary Text}      % Produces section heading.  Lower-level
                             % sections are begun with similar
                             % \subsection and \subsubsection commands.

The ends  of words and sentences are marked
  by   spaces. It  doesn't matter how many
spaces    you type; one is as good as 100.  The
end of   a line counts as a space.

One   or more   blank lines denote the  end
of  a paragraph.

Since any number of consecutive spaces are treated
like a single one, the formatting of the input
file makes no difference to
      \LaTeX,                % The \LaTeX command generates the LaTeX logo.
but it makes a difference to you.  When you use
\LaTeX, making your input file as easy to read
as possible will be a great help as you write
your document and when you change it.  This sample
file shows how you can add comments to your own input
file.

Because printing is different from typewriting,
there are a number of things that you have to do
differently when preparing an input file than if
you were just typing the document directly.
Quotation marks like
       ``this''
have to be handled specially, as do quotes within
quotes:
       ``\,`this'            % \, separates the double and single quote.
        is what I just
        wrote, not  `that'\,''.

Dashes come in three sizes: an
       intra-word
dash, a medium dash for number ranges like
       1--2,
and a punctuation
       dash---like
this.

A sentence-ending space should be larger than the
space between words within a sentence.  You
sometimes have to type special commands in
conjunction with punctuation characters to get
this right, as in the following sentence.
       Gnats, gnus, etc.\ all  % `\ ' makes an inter-word space.
       begin with G\@.         % \@ marks end-of-sentence punctuation.
You should check the spaces after periods when
reading your output to make sure you haven't
forgotten any special cases.  Generating an
ellipsis
       \ldots\               % `\ ' is needed after `\ldots' because TeX
                             % ignores spaces after command names like \ldots
                             % made from \ + letters.
                             %
                             % Note how a `%' character causes TeX to ignore
                             % the end of the input line, so these blank lines
                             % do not start a new paragraph.
                             %
with the right spacing around the periods requires
a special command.

\LaTeX\ interprets some common characters as
commands, so you must type special commands to
generate them.  These characters include the
following:
       \$ \& \% \# \{ and \}.

In printing, text is usually emphasized with an
       \emph{italic}
type style.

\begin{em}
   A long segment of text can also be emphasized
   in this way.  Text within such a segment can be
   given \emph{additional} emphasis.
\end{em}

It is sometimes necessary to prevent \LaTeX\ from
breaking a line where it might otherwise do so.
This may be at a space, as between the ``Mr.''\ and
``Jones'' in
       ``Mr.~Jones'',        % ~ produces an unbreakable interword space.
or within a word---especially when the word is a
symbol like
       \mbox{\emph{itemnum}}
that makes little sense when hyphenated across
lines.

Footnotes\footnote{This is an example of a footnote.}
pose no problem.

\LaTeX\ is good at typesetting mathematical formulas
like
       \( x-3y + z = 7 \)
or
       \( a_{1} > x^{2n} + y^{2n} > x' \)
or
       \( \ip{A}{B} = \sum_{i} a_{i} b_{i} \).
The spaces you type in a formula are
ignored.  Remember that a letter like
       $x$                   % $ ... $  and  \( ... \)  are equivalent
is a formula when it denotes a mathematical
symbol, and it should be typed as one.

\section{Displayed Text}

Text is displayed by indenting it from the left
margin.  Quotations are commonly displayed.  There
are short quotations
\begin{quote}
   This is a short quotation.  It consists of a
   single paragraph of text.  See how it is formatted.
\end{quote}
and longer ones.
\begin{quotation}
   This is a longer quotation.  It consists of two
   paragraphs of text, neither of which are
   particularly interesting.

   This is the second paragraph of the quotation.  It
   is just as dull as the first paragraph.
\end{quotation}
Another frequently-displayed structure is a list.
The following is an example of an \emph{itemized}
list.
\begin{itemize}
   \item This is the first item of an itemized list.
         Each item in the list is marked with a ``tick''.
         You don't have to worry about what kind of tick
         mark is used.

   \item This is the second item of the list.  It
         contains another list nested inside it.  The inner
         list is an \emph{enumerated} list.
         \begin{enumerate}
            \item This is the first item of an enumerated
                  list that is nested within the itemized list.

            \item This is the second item of the inner list.
                  \LaTeX\ allows you to nest lists deeper than
                  you really should.
         \end{enumerate}
         This is the rest of the second item of the outer
         list.  It is no more interesting than any other
         part of the item.
   \item This is the third item of the list.
\end{itemize}
You can even display poetry.
\begin{verse}
   There is an environment
    for verse \\             % The \\ command separates lines
   Whose features some poets % within a stanza.
   will curse.

                             % One or more blank lines separate stanzas.

   For instead of making\\
   Them do \emph{all} line breaking, \\
   It allows them to put too many words on a line when they'd rather be
   forced to be terse.
\end{verse}

Mathematical formulas may also be displayed.  A
displayed formula
is
one-line long; multiline
formulas require special formatting instructions.
   \[  \ip{\Gamma}{\psi'} = x'' + y^{2} + z_{i}^{n}\]
Don't start a paragraph with a displayed equation,
nor make one a paragraph by itself.

\end{document}               % End of document.

\section{Bench Huge Section 1441}
% synthetic bench section 1441
% This is a sample LaTeX input file.  (Version of 12 August 2004.)
%
% A '%' character causes TeX to ignore all remaining text on the line,
% and is used for comments like this one.

\documentclass{article}      % Specifies the document class

                             % The preamble begins here.
\title{An Example Document}  % Declares the document's title.
\author{Leslie Lamport}      % Declares the author's name.
\date{January 21, 1994}      % Deleting this command produces today's date.

\newcommand{\ip}[2]{(#1, #2)}
                             % Defines \ip{arg1}{arg2} to mean
                             % (arg1, arg2).

%\newcommand{\ip}[2]{\langle #1 | #2\rangle}
                             % This is an alternative definition of
                             % \ip that is commented out.

\begin{document}             % End of preamble and beginning of text.

\maketitle                   % Produces the title.

This is an example input file.  Comparing it with
the output it generates can show you how to
produce a simple document of your own.

\section{Ordinary Text}      % Produces section heading.  Lower-level
                             % sections are begun with similar
                             % \subsection and \subsubsection commands.

The ends  of words and sentences are marked
  by   spaces. It  doesn't matter how many
spaces    you type; one is as good as 100.  The
end of   a line counts as a space.

One   or more   blank lines denote the  end
of  a paragraph.

Since any number of consecutive spaces are treated
like a single one, the formatting of the input
file makes no difference to
      \LaTeX,                % The \LaTeX command generates the LaTeX logo.
but it makes a difference to you.  When you use
\LaTeX, making your input file as easy to read
as possible will be a great help as you write
your document and when you change it.  This sample
file shows how you can add comments to your own input
file.

Because printing is different from typewriting,
there are a number of things that you have to do
differently when preparing an input file than if
you were just typing the document directly.
Quotation marks like
       ``this''
have to be handled specially, as do quotes within
quotes:
       ``\,`this'            % \, separates the double and single quote.
        is what I just
        wrote, not  `that'\,''.

Dashes come in three sizes: an
       intra-word
dash, a medium dash for number ranges like
       1--2,
and a punctuation
       dash---like
this.

A sentence-ending space should be larger than the
space between words within a sentence.  You
sometimes have to type special commands in
conjunction with punctuation characters to get
this right, as in the following sentence.
       Gnats, gnus, etc.\ all  % `\ ' makes an inter-word space.
       begin with G\@.         % \@ marks end-of-sentence punctuation.
You should check the spaces after periods when
reading your output to make sure you haven't
forgotten any special cases.  Generating an
ellipsis
       \ldots\               % `\ ' is needed after `\ldots' because TeX
                             % ignores spaces after command names like \ldots
                             % made from \ + letters.
                             %
                             % Note how a `%' character causes TeX to ignore
                             % the end of the input line, so these blank lines
                             % do not start a new paragraph.
                             %
with the right spacing around the periods requires
a special command.

\LaTeX\ interprets some common characters as
commands, so you must type special commands to
generate them.  These characters include the
following:
       \$ \& \% \# \{ and \}.

In printing, text is usually emphasized with an
       \emph{italic}
type style.

\begin{em}
   A long segment of text can also be emphasized
   in this way.  Text within such a segment can be
   given \emph{additional} emphasis.
\end{em}

It is sometimes necessary to prevent \LaTeX\ from
breaking a line where it might otherwise do so.
This may be at a space, as between the ``Mr.''\ and
``Jones'' in
       ``Mr.~Jones'',        % ~ produces an unbreakable interword space.
or within a word---especially when the word is a
symbol like
       \mbox{\emph{itemnum}}
that makes little sense when hyphenated across
lines.

Footnotes\footnote{This is an example of a footnote.}
pose no problem.

\LaTeX\ is good at typesetting mathematical formulas
like
       \( x-3y + z = 7 \)
or
       \( a_{1} > x^{2n} + y^{2n} > x' \)
or
       \( \ip{A}{B} = \sum_{i} a_{i} b_{i} \).
The spaces you type in a formula are
ignored.  Remember that a letter like
       $x$                   % $ ... $  and  \( ... \)  are equivalent
is a formula when it denotes a mathematical
symbol, and it should be typed as one.

\section{Displayed Text}

Text is displayed by indenting it from the left
margin.  Quotations are commonly displayed.  There
are short quotations
\begin{quote}
   This is a short quotation.  It consists of a
   single paragraph of text.  See how it is formatted.
\end{quote}
and longer ones.
\begin{quotation}
   This is a longer quotation.  It consists of two
   paragraphs of text, neither of which are
   particularly interesting.

   This is the second paragraph of the quotation.  It
   is just as dull as the first paragraph.
\end{quotation}
Another frequently-displayed structure is a list.
The following is an example of an \emph{itemized}
list.
\begin{itemize}
   \item This is the first item of an itemized list.
         Each item in the list is marked with a ``tick''.
         You don't have to worry about what kind of tick
         mark is used.

   \item This is the second item of the list.  It
         contains another list nested inside it.  The inner
         list is an \emph{enumerated} list.
         \begin{enumerate}
            \item This is the first item of an enumerated
                  list that is nested within the itemized list.

            \item This is the second item of the inner list.
                  \LaTeX\ allows you to nest lists deeper than
                  you really should.
         \end{enumerate}
         This is the rest of the second item of the outer
         list.  It is no more interesting than any other
         part of the item.
   \item This is the third item of the list.
\end{itemize}
You can even display poetry.
\begin{verse}
   There is an environment
    for verse \\             % The \\ command separates lines
   Whose features some poets % within a stanza.
   will curse.

                             % One or more blank lines separate stanzas.

   For instead of making\\
   Them do \emph{all} line breaking, \\
   It allows them to put too many words on a line when they'd rather be
   forced to be terse.
\end{verse}

Mathematical formulas may also be displayed.  A
displayed formula
is
one-line long; multiline
formulas require special formatting instructions.
   \[  \ip{\Gamma}{\psi'} = x'' + y^{2} + z_{i}^{n}\]
Don't start a paragraph with a displayed equation,
nor make one a paragraph by itself.

\end{document}               % End of document.

\section{Bench Huge Section 1442}
% synthetic bench section 1442
% This is a sample LaTeX input file.  (Version of 12 August 2004.)
%
% A '%' character causes TeX to ignore all remaining text on the line,
% and is used for comments like this one.

\documentclass{article}      % Specifies the document class

                             % The preamble begins here.
\title{An Example Document}  % Declares the document's title.
\author{Leslie Lamport}      % Declares the author's name.
\date{January 21, 1994}      % Deleting this command produces today's date.

\newcommand{\ip}[2]{(#1, #2)}
                             % Defines \ip{arg1}{arg2} to mean
                             % (arg1, arg2).

%\newcommand{\ip}[2]{\langle #1 | #2\rangle}
                             % This is an alternative definition of
                             % \ip that is commented out.

\begin{document}             % End of preamble and beginning of text.

\maketitle                   % Produces the title.

This is an example input file.  Comparing it with
the output it generates can show you how to
produce a simple document of your own.

\section{Ordinary Text}      % Produces section heading.  Lower-level
                             % sections are begun with similar
                             % \subsection and \subsubsection commands.

The ends  of words and sentences are marked
  by   spaces. It  doesn't matter how many
spaces    you type; one is as good as 100.  The
end of   a line counts as a space.

One   or more   blank lines denote the  end
of  a paragraph.

Since any number of consecutive spaces are treated
like a single one, the formatting of the input
file makes no difference to
      \LaTeX,                % The \LaTeX command generates the LaTeX logo.
but it makes a difference to you.  When you use
\LaTeX, making your input file as easy to read
as possible will be a great help as you write
your document and when you change it.  This sample
file shows how you can add comments to your own input
file.

Because printing is different from typewriting,
there are a number of things that you have to do
differently when preparing an input file than if
you were just typing the document directly.
Quotation marks like
       ``this''
have to be handled specially, as do quotes within
quotes:
       ``\,`this'            % \, separates the double and single quote.
        is what I just
        wrote, not  `that'\,''.

Dashes come in three sizes: an
       intra-word
dash, a medium dash for number ranges like
       1--2,
and a punctuation
       dash---like
this.

A sentence-ending space should be larger than the
space between words within a sentence.  You
sometimes have to type special commands in
conjunction with punctuation characters to get
this right, as in the following sentence.
       Gnats, gnus, etc.\ all  % `\ ' makes an inter-word space.
       begin with G\@.         % \@ marks end-of-sentence punctuation.
You should check the spaces after periods when
reading your output to make sure you haven't
forgotten any special cases.  Generating an
ellipsis
       \ldots\               % `\ ' is needed after `\ldots' because TeX
                             % ignores spaces after command names like \ldots
                             % made from \ + letters.
                             %
                             % Note how a `%' character causes TeX to ignore
                             % the end of the input line, so these blank lines
                             % do not start a new paragraph.
                             %
with the right spacing around the periods requires
a special command.

\LaTeX\ interprets some common characters as
commands, so you must type special commands to
generate them.  These characters include the
following:
       \$ \& \% \# \{ and \}.

In printing, text is usually emphasized with an
       \emph{italic}
type style.

\begin{em}
   A long segment of text can also be emphasized
   in this way.  Text within such a segment can be
   given \emph{additional} emphasis.
\end{em}

It is sometimes necessary to prevent \LaTeX\ from
breaking a line where it might otherwise do so.
This may be at a space, as between the ``Mr.''\ and
``Jones'' in
       ``Mr.~Jones'',        % ~ produces an unbreakable interword space.
or within a word---especially when the word is a
symbol like
       \mbox{\emph{itemnum}}
that makes little sense when hyphenated across
lines.

Footnotes\footnote{This is an example of a footnote.}
pose no problem.

\LaTeX\ is good at typesetting mathematical formulas
like
       \( x-3y + z = 7 \)
or
       \( a_{1} > x^{2n} + y^{2n} > x' \)
or
       \( \ip{A}{B} = \sum_{i} a_{i} b_{i} \).
The spaces you type in a formula are
ignored.  Remember that a letter like
       $x$                   % $ ... $  and  \( ... \)  are equivalent
is a formula when it denotes a mathematical
symbol, and it should be typed as one.

\section{Displayed Text}

Text is displayed by indenting it from the left
margin.  Quotations are commonly displayed.  There
are short quotations
\begin{quote}
   This is a short quotation.  It consists of a
   single paragraph of text.  See how it is formatted.
\end{quote}
and longer ones.
\begin{quotation}
   This is a longer quotation.  It consists of two
   paragraphs of text, neither of which are
   particularly interesting.

   This is the second paragraph of the quotation.  It
   is just as dull as the first paragraph.
\end{quotation}
Another frequently-displayed structure is a list.
The following is an example of an \emph{itemized}
list.
\begin{itemize}
   \item This is the first item of an itemized list.
         Each item in the list is marked with a ``tick''.
         You don't have to worry about what kind of tick
         mark is used.

   \item This is the second item of the list.  It
         contains another list nested inside it.  The inner
         list is an \emph{enumerated} list.
         \begin{enumerate}
            \item This is the first item of an enumerated
                  list that is nested within the itemized list.

            \item This is the second item of the inner list.
                  \LaTeX\ allows you to nest lists deeper than
                  you really should.
         \end{enumerate}
         This is the rest of the second item of the outer
         list.  It is no more interesting than any other
         part of the item.
   \item This is the third item of the list.
\end{itemize}
You can even display poetry.
\begin{verse}
   There is an environment
    for verse \\             % The \\ command separates lines
   Whose features some poets % within a stanza.
   will curse.

                             % One or more blank lines separate stanzas.

   For instead of making\\
   Them do \emph{all} line breaking, \\
   It allows them to put too many words on a line when they'd rather be
   forced to be terse.
\end{verse}

Mathematical formulas may also be displayed.  A
displayed formula
is
one-line long; multiline
formulas require special formatting instructions.
   \[  \ip{\Gamma}{\psi'} = x'' + y^{2} + z_{i}^{n}\]
Don't start a paragraph with a displayed equation,
nor make one a paragraph by itself.

\end{document}               % End of document.

\section{Bench Huge Section 1443}
% synthetic bench section 1443
% This is a sample LaTeX input file.  (Version of 12 August 2004.)
%
% A '%' character causes TeX to ignore all remaining text on the line,
% and is used for comments like this one.

\documentclass{article}      % Specifies the document class

                             % The preamble begins here.
\title{An Example Document}  % Declares the document's title.
\author{Leslie Lamport}      % Declares the author's name.
\date{January 21, 1994}      % Deleting this command produces today's date.

\newcommand{\ip}[2]{(#1, #2)}
                             % Defines \ip{arg1}{arg2} to mean
                             % (arg1, arg2).

%\newcommand{\ip}[2]{\langle #1 | #2\rangle}
                             % This is an alternative definition of
                             % \ip that is commented out.

\begin{document}             % End of preamble and beginning of text.

\maketitle                   % Produces the title.

This is an example input file.  Comparing it with
the output it generates can show you how to
produce a simple document of your own.

\section{Ordinary Text}      % Produces section heading.  Lower-level
                             % sections are begun with similar
                             % \subsection and \subsubsection commands.

The ends  of words and sentences are marked
  by   spaces. It  doesn't matter how many
spaces    you type; one is as good as 100.  The
end of   a line counts as a space.

One   or more   blank lines denote the  end
of  a paragraph.

Since any number of consecutive spaces are treated
like a single one, the formatting of the input
file makes no difference to
      \LaTeX,                % The \LaTeX command generates the LaTeX logo.
but it makes a difference to you.  When you use
\LaTeX, making your input file as easy to read
as possible will be a great help as you write
your document and when you change it.  This sample
file shows how you can add comments to your own input
file.

Because printing is different from typewriting,
there are a number of things that you have to do
differently when preparing an input file than if
you were just typing the document directly.
Quotation marks like
       ``this''
have to be handled specially, as do quotes within
quotes:
       ``\,`this'            % \, separates the double and single quote.
        is what I just
        wrote, not  `that'\,''.

Dashes come in three sizes: an
       intra-word
dash, a medium dash for number ranges like
       1--2,
and a punctuation
       dash---like
this.

A sentence-ending space should be larger than the
space between words within a sentence.  You
sometimes have to type special commands in
conjunction with punctuation characters to get
this right, as in the following sentence.
       Gnats, gnus, etc.\ all  % `\ ' makes an inter-word space.
       begin with G\@.         % \@ marks end-of-sentence punctuation.
You should check the spaces after periods when
reading your output to make sure you haven't
forgotten any special cases.  Generating an
ellipsis
       \ldots\               % `\ ' is needed after `\ldots' because TeX
                             % ignores spaces after command names like \ldots
                             % made from \ + letters.
                             %
                             % Note how a `%' character causes TeX to ignore
                             % the end of the input line, so these blank lines
                             % do not start a new paragraph.
                             %
with the right spacing around the periods requires
a special command.

\LaTeX\ interprets some common characters as
commands, so you must type special commands to
generate them.  These characters include the
following:
       \$ \& \% \# \{ and \}.

In printing, text is usually emphasized with an
       \emph{italic}
type style.

\begin{em}
   A long segment of text can also be emphasized
   in this way.  Text within such a segment can be
   given \emph{additional} emphasis.
\end{em}

It is sometimes necessary to prevent \LaTeX\ from
breaking a line where it might otherwise do so.
This may be at a space, as between the ``Mr.''\ and
``Jones'' in
       ``Mr.~Jones'',        % ~ produces an unbreakable interword space.
or within a word---especially when the word is a
symbol like
       \mbox{\emph{itemnum}}
that makes little sense when hyphenated across
lines.

Footnotes\footnote{This is an example of a footnote.}
pose no problem.

\LaTeX\ is good at typesetting mathematical formulas
like
       \( x-3y + z = 7 \)
or
       \( a_{1} > x^{2n} + y^{2n} > x' \)
or
       \( \ip{A}{B} = \sum_{i} a_{i} b_{i} \).
The spaces you type in a formula are
ignored.  Remember that a letter like
       $x$                   % $ ... $  and  \( ... \)  are equivalent
is a formula when it denotes a mathematical
symbol, and it should be typed as one.

\section{Displayed Text}

Text is displayed by indenting it from the left
margin.  Quotations are commonly displayed.  There
are short quotations
\begin{quote}
   This is a short quotation.  It consists of a
   single paragraph of text.  See how it is formatted.
\end{quote}
and longer ones.
\begin{quotation}
   This is a longer quotation.  It consists of two
   paragraphs of text, neither of which are
   particularly interesting.

   This is the second paragraph of the quotation.  It
   is just as dull as the first paragraph.
\end{quotation}
Another frequently-displayed structure is a list.
The following is an example of an \emph{itemized}
list.
\begin{itemize}
   \item This is the first item of an itemized list.
         Each item in the list is marked with a ``tick''.
         You don't have to worry about what kind of tick
         mark is used.

   \item This is the second item of the list.  It
         contains another list nested inside it.  The inner
         list is an \emph{enumerated} list.
         \begin{enumerate}
            \item This is the first item of an enumerated
                  list that is nested within the itemized list.

            \item This is the second item of the inner list.
                  \LaTeX\ allows you to nest lists deeper than
                  you really should.
         \end{enumerate}
         This is the rest of the second item of the outer
         list.  It is no more interesting than any other
         part of the item.
   \item This is the third item of the list.
\end{itemize}
You can even display poetry.
\begin{verse}
   There is an environment
    for verse \\             % The \\ command separates lines
   Whose features some poets % within a stanza.
   will curse.

                             % One or more blank lines separate stanzas.

   For instead of making\\
   Them do \emph{all} line breaking, \\
   It allows them to put too many words on a line when they'd rather be
   forced to be terse.
\end{verse}

Mathematical formulas may also be displayed.  A
displayed formula
is
one-line long; multiline
formulas require special formatting instructions.
   \[  \ip{\Gamma}{\psi'} = x'' + y^{2} + z_{i}^{n}\]
Don't start a paragraph with a displayed equation,
nor make one a paragraph by itself.

\end{document}               % End of document.

\section{Bench Huge Section 1444}
% synthetic bench section 1444
% This is a sample LaTeX input file.  (Version of 12 August 2004.)
%
% A '%' character causes TeX to ignore all remaining text on the line,
% and is used for comments like this one.

\documentclass{article}      % Specifies the document class

                             % The preamble begins here.
\title{An Example Document}  % Declares the document's title.
\author{Leslie Lamport}      % Declares the author's name.
\date{January 21, 1994}      % Deleting this command produces today's date.

\newcommand{\ip}[2]{(#1, #2)}
                             % Defines \ip{arg1}{arg2} to mean
                             % (arg1, arg2).

%\newcommand{\ip}[2]{\langle #1 | #2\rangle}
                             % This is an alternative definition of
                             % \ip that is commented out.

\begin{document}             % End of preamble and beginning of text.

\maketitle                   % Produces the title.

This is an example input file.  Comparing it with
the output it generates can show you how to
produce a simple document of your own.

\section{Ordinary Text}      % Produces section heading.  Lower-level
                             % sections are begun with similar
                             % \subsection and \subsubsection commands.

The ends  of words and sentences are marked
  by   spaces. It  doesn't matter how many
spaces    you type; one is as good as 100.  The
end of   a line counts as a space.

One   or more   blank lines denote the  end
of  a paragraph.

Since any number of consecutive spaces are treated
like a single one, the formatting of the input
file makes no difference to
      \LaTeX,                % The \LaTeX command generates the LaTeX logo.
but it makes a difference to you.  When you use
\LaTeX, making your input file as easy to read
as possible will be a great help as you write
your document and when you change it.  This sample
file shows how you can add comments to your own input
file.

Because printing is different from typewriting,
there are a number of things that you have to do
differently when preparing an input file than if
you were just typing the document directly.
Quotation marks like
       ``this''
have to be handled specially, as do quotes within
quotes:
       ``\,`this'            % \, separates the double and single quote.
        is what I just
        wrote, not  `that'\,''.

Dashes come in three sizes: an
       intra-word
dash, a medium dash for number ranges like
       1--2,
and a punctuation
       dash---like
this.

A sentence-ending space should be larger than the
space between words within a sentence.  You
sometimes have to type special commands in
conjunction with punctuation characters to get
this right, as in the following sentence.
       Gnats, gnus, etc.\ all  % `\ ' makes an inter-word space.
       begin with G\@.         % \@ marks end-of-sentence punctuation.
You should check the spaces after periods when
reading your output to make sure you haven't
forgotten any special cases.  Generating an
ellipsis
       \ldots\               % `\ ' is needed after `\ldots' because TeX
                             % ignores spaces after command names like \ldots
                             % made from \ + letters.
                             %
                             % Note how a `%' character causes TeX to ignore
                             % the end of the input line, so these blank lines
                             % do not start a new paragraph.
                             %
with the right spacing around the periods requires
a special command.

\LaTeX\ interprets some common characters as
commands, so you must type special commands to
generate them.  These characters include the
following:
       \$ \& \% \# \{ and \}.

In printing, text is usually emphasized with an
       \emph{italic}
type style.

\begin{em}
   A long segment of text can also be emphasized
   in this way.  Text within such a segment can be
   given \emph{additional} emphasis.
\end{em}

It is sometimes necessary to prevent \LaTeX\ from
breaking a line where it might otherwise do so.
This may be at a space, as between the ``Mr.''\ and
``Jones'' in
       ``Mr.~Jones'',        % ~ produces an unbreakable interword space.
or within a word---especially when the word is a
symbol like
       \mbox{\emph{itemnum}}
that makes little sense when hyphenated across
lines.

Footnotes\footnote{This is an example of a footnote.}
pose no problem.

\LaTeX\ is good at typesetting mathematical formulas
like
       \( x-3y + z = 7 \)
or
       \( a_{1} > x^{2n} + y^{2n} > x' \)
or
       \( \ip{A}{B} = \sum_{i} a_{i} b_{i} \).
The spaces you type in a formula are
ignored.  Remember that a letter like
       $x$                   % $ ... $  and  \( ... \)  are equivalent
is a formula when it denotes a mathematical
symbol, and it should be typed as one.

\section{Displayed Text}

Text is displayed by indenting it from the left
margin.  Quotations are commonly displayed.  There
are short quotations
\begin{quote}
   This is a short quotation.  It consists of a
   single paragraph of text.  See how it is formatted.
\end{quote}
and longer ones.
\begin{quotation}
   This is a longer quotation.  It consists of two
   paragraphs of text, neither of which are
   particularly interesting.

   This is the second paragraph of the quotation.  It
   is just as dull as the first paragraph.
\end{quotation}
Another frequently-displayed structure is a list.
The following is an example of an \emph{itemized}
list.
\begin{itemize}
   \item This is the first item of an itemized list.
         Each item in the list is marked with a ``tick''.
         You don't have to worry about what kind of tick
         mark is used.

   \item This is the second item of the list.  It
         contains another list nested inside it.  The inner
         list is an \emph{enumerated} list.
         \begin{enumerate}
            \item This is the first item of an enumerated
                  list that is nested within the itemized list.

            \item This is the second item of the inner list.
                  \LaTeX\ allows you to nest lists deeper than
                  you really should.
         \end{enumerate}
         This is the rest of the second item of the outer
         list.  It is no more interesting than any other
         part of the item.
   \item This is the third item of the list.
\end{itemize}
You can even display poetry.
\begin{verse}
   There is an environment
    for verse \\             % The \\ command separates lines
   Whose features some poets % within a stanza.
   will curse.

                             % One or more blank lines separate stanzas.

   For instead of making\\
   Them do \emph{all} line breaking, \\
   It allows them to put too many words on a line when they'd rather be
   forced to be terse.
\end{verse}

Mathematical formulas may also be displayed.  A
displayed formula
is
one-line long; multiline
formulas require special formatting instructions.
   \[  \ip{\Gamma}{\psi'} = x'' + y^{2} + z_{i}^{n}\]
Don't start a paragraph with a displayed equation,
nor make one a paragraph by itself.

\end{document}               % End of document.

\section{Bench Huge Section 1445}
% synthetic bench section 1445
% This is a sample LaTeX input file.  (Version of 12 August 2004.)
%
% A '%' character causes TeX to ignore all remaining text on the line,
% and is used for comments like this one.

\documentclass{article}      % Specifies the document class

                             % The preamble begins here.
\title{An Example Document}  % Declares the document's title.
\author{Leslie Lamport}      % Declares the author's name.
\date{January 21, 1994}      % Deleting this command produces today's date.

\newcommand{\ip}[2]{(#1, #2)}
                             % Defines \ip{arg1}{arg2} to mean
                             % (arg1, arg2).

%\newcommand{\ip}[2]{\langle #1 | #2\rangle}
                             % This is an alternative definition of
                             % \ip that is commented out.

\begin{document}             % End of preamble and beginning of text.

\maketitle                   % Produces the title.

This is an example input file.  Comparing it with
the output it generates can show you how to
produce a simple document of your own.

\section{Ordinary Text}      % Produces section heading.  Lower-level
                             % sections are begun with similar
                             % \subsection and \subsubsection commands.

The ends  of words and sentences are marked
  by   spaces. It  doesn't matter how many
spaces    you type; one is as good as 100.  The
end of   a line counts as a space.

One   or more   blank lines denote the  end
of  a paragraph.

Since any number of consecutive spaces are treated
like a single one, the formatting of the input
file makes no difference to
      \LaTeX,                % The \LaTeX command generates the LaTeX logo.
but it makes a difference to you.  When you use
\LaTeX, making your input file as easy to read
as possible will be a great help as you write
your document and when you change it.  This sample
file shows how you can add comments to your own input
file.

Because printing is different from typewriting,
there are a number of things that you have to do
differently when preparing an input file than if
you were just typing the document directly.
Quotation marks like
       ``this''
have to be handled specially, as do quotes within
quotes:
       ``\,`this'            % \, separates the double and single quote.
        is what I just
        wrote, not  `that'\,''.

Dashes come in three sizes: an
       intra-word
dash, a medium dash for number ranges like
       1--2,
and a punctuation
       dash---like
this.

A sentence-ending space should be larger than the
space between words within a sentence.  You
sometimes have to type special commands in
conjunction with punctuation characters to get
this right, as in the following sentence.
       Gnats, gnus, etc.\ all  % `\ ' makes an inter-word space.
       begin with G\@.         % \@ marks end-of-sentence punctuation.
You should check the spaces after periods when
reading your output to make sure you haven't
forgotten any special cases.  Generating an
ellipsis
       \ldots\               % `\ ' is needed after `\ldots' because TeX
                             % ignores spaces after command names like \ldots
                             % made from \ + letters.
                             %
                             % Note how a `%' character causes TeX to ignore
                             % the end of the input line, so these blank lines
                             % do not start a new paragraph.
                             %
with the right spacing around the periods requires
a special command.

\LaTeX\ interprets some common characters as
commands, so you must type special commands to
generate them.  These characters include the
following:
       \$ \& \% \# \{ and \}.

In printing, text is usually emphasized with an
       \emph{italic}
type style.

\begin{em}
   A long segment of text can also be emphasized
   in this way.  Text within such a segment can be
   given \emph{additional} emphasis.
\end{em}

It is sometimes necessary to prevent \LaTeX\ from
breaking a line where it might otherwise do so.
This may be at a space, as between the ``Mr.''\ and
``Jones'' in
       ``Mr.~Jones'',        % ~ produces an unbreakable interword space.
or within a word---especially when the word is a
symbol like
       \mbox{\emph{itemnum}}
that makes little sense when hyphenated across
lines.

Footnotes\footnote{This is an example of a footnote.}
pose no problem.

\LaTeX\ is good at typesetting mathematical formulas
like
       \( x-3y + z = 7 \)
or
       \( a_{1} > x^{2n} + y^{2n} > x' \)
or
       \( \ip{A}{B} = \sum_{i} a_{i} b_{i} \).
The spaces you type in a formula are
ignored.  Remember that a letter like
       $x$                   % $ ... $  and  \( ... \)  are equivalent
is a formula when it denotes a mathematical
symbol, and it should be typed as one.

\section{Displayed Text}

Text is displayed by indenting it from the left
margin.  Quotations are commonly displayed.  There
are short quotations
\begin{quote}
   This is a short quotation.  It consists of a
   single paragraph of text.  See how it is formatted.
\end{quote}
and longer ones.
\begin{quotation}
   This is a longer quotation.  It consists of two
   paragraphs of text, neither of which are
   particularly interesting.

   This is the second paragraph of the quotation.  It
   is just as dull as the first paragraph.
\end{quotation}
Another frequently-displayed structure is a list.
The following is an example of an \emph{itemized}
list.
\begin{itemize}
   \item This is the first item of an itemized list.
         Each item in the list is marked with a ``tick''.
         You don't have to worry about what kind of tick
         mark is used.

   \item This is the second item of the list.  It
         contains another list nested inside it.  The inner
         list is an \emph{enumerated} list.
         \begin{enumerate}
            \item This is the first item of an enumerated
                  list that is nested within the itemized list.

            \item This is the second item of the inner list.
                  \LaTeX\ allows you to nest lists deeper than
                  you really should.
         \end{enumerate}
         This is the rest of the second item of the outer
         list.  It is no more interesting than any other
         part of the item.
   \item This is the third item of the list.
\end{itemize}
You can even display poetry.
\begin{verse}
   There is an environment
    for verse \\             % The \\ command separates lines
   Whose features some poets % within a stanza.
   will curse.

                             % One or more blank lines separate stanzas.

   For instead of making\\
   Them do \emph{all} line breaking, \\
   It allows them to put too many words on a line when they'd rather be
   forced to be terse.
\end{verse}

Mathematical formulas may also be displayed.  A
displayed formula
is
one-line long; multiline
formulas require special formatting instructions.
   \[  \ip{\Gamma}{\psi'} = x'' + y^{2} + z_{i}^{n}\]
Don't start a paragraph with a displayed equation,
nor make one a paragraph by itself.

\end{document}               % End of document.

\section{Bench Huge Section 1446}
% synthetic bench section 1446
% This is a sample LaTeX input file.  (Version of 12 August 2004.)
%
% A '%' character causes TeX to ignore all remaining text on the line,
% and is used for comments like this one.

\documentclass{article}      % Specifies the document class

                             % The preamble begins here.
\title{An Example Document}  % Declares the document's title.
\author{Leslie Lamport}      % Declares the author's name.
\date{January 21, 1994}      % Deleting this command produces today's date.

\newcommand{\ip}[2]{(#1, #2)}
                             % Defines \ip{arg1}{arg2} to mean
                             % (arg1, arg2).

%\newcommand{\ip}[2]{\langle #1 | #2\rangle}
                             % This is an alternative definition of
                             % \ip that is commented out.

\begin{document}             % End of preamble and beginning of text.

\maketitle                   % Produces the title.

This is an example input file.  Comparing it with
the output it generates can show you how to
produce a simple document of your own.

\section{Ordinary Text}      % Produces section heading.  Lower-level
                             % sections are begun with similar
                             % \subsection and \subsubsection commands.

The ends  of words and sentences are marked
  by   spaces. It  doesn't matter how many
spaces    you type; one is as good as 100.  The
end of   a line counts as a space.

One   or more   blank lines denote the  end
of  a paragraph.

Since any number of consecutive spaces are treated
like a single one, the formatting of the input
file makes no difference to
      \LaTeX,                % The \LaTeX command generates the LaTeX logo.
but it makes a difference to you.  When you use
\LaTeX, making your input file as easy to read
as possible will be a great help as you write
your document and when you change it.  This sample
file shows how you can add comments to your own input
file.

Because printing is different from typewriting,
there are a number of things that you have to do
differently when preparing an input file than if
you were just typing the document directly.
Quotation marks like
       ``this''
have to be handled specially, as do quotes within
quotes:
       ``\,`this'            % \, separates the double and single quote.
        is what I just
        wrote, not  `that'\,''.

Dashes come in three sizes: an
       intra-word
dash, a medium dash for number ranges like
       1--2,
and a punctuation
       dash---like
this.

A sentence-ending space should be larger than the
space between words within a sentence.  You
sometimes have to type special commands in
conjunction with punctuation characters to get
this right, as in the following sentence.
       Gnats, gnus, etc.\ all  % `\ ' makes an inter-word space.
       begin with G\@.         % \@ marks end-of-sentence punctuation.
You should check the spaces after periods when
reading your output to make sure you haven't
forgotten any special cases.  Generating an
ellipsis
       \ldots\               % `\ ' is needed after `\ldots' because TeX
                             % ignores spaces after command names like \ldots
                             % made from \ + letters.
                             %
                             % Note how a `%' character causes TeX to ignore
                             % the end of the input line, so these blank lines
                             % do not start a new paragraph.
                             %
with the right spacing around the periods requires
a special command.

\LaTeX\ interprets some common characters as
commands, so you must type special commands to
generate them.  These characters include the
following:
       \$ \& \% \# \{ and \}.

In printing, text is usually emphasized with an
       \emph{italic}
type style.

\begin{em}
   A long segment of text can also be emphasized
   in this way.  Text within such a segment can be
   given \emph{additional} emphasis.
\end{em}

It is sometimes necessary to prevent \LaTeX\ from
breaking a line where it might otherwise do so.
This may be at a space, as between the ``Mr.''\ and
``Jones'' in
       ``Mr.~Jones'',        % ~ produces an unbreakable interword space.
or within a word---especially when the word is a
symbol like
       \mbox{\emph{itemnum}}
that makes little sense when hyphenated across
lines.

Footnotes\footnote{This is an example of a footnote.}
pose no problem.

\LaTeX\ is good at typesetting mathematical formulas
like
       \( x-3y + z = 7 \)
or
       \( a_{1} > x^{2n} + y^{2n} > x' \)
or
       \( \ip{A}{B} = \sum_{i} a_{i} b_{i} \).
The spaces you type in a formula are
ignored.  Remember that a letter like
       $x$                   % $ ... $  and  \( ... \)  are equivalent
is a formula when it denotes a mathematical
symbol, and it should be typed as one.

\section{Displayed Text}

Text is displayed by indenting it from the left
margin.  Quotations are commonly displayed.  There
are short quotations
\begin{quote}
   This is a short quotation.  It consists of a
   single paragraph of text.  See how it is formatted.
\end{quote}
and longer ones.
\begin{quotation}
   This is a longer quotation.  It consists of two
   paragraphs of text, neither of which are
   particularly interesting.

   This is the second paragraph of the quotation.  It
   is just as dull as the first paragraph.
\end{quotation}
Another frequently-displayed structure is a list.
The following is an example of an \emph{itemized}
list.
\begin{itemize}
   \item This is the first item of an itemized list.
         Each item in the list is marked with a ``tick''.
         You don't have to worry about what kind of tick
         mark is used.

   \item This is the second item of the list.  It
         contains another list nested inside it.  The inner
         list is an \emph{enumerated} list.
         \begin{enumerate}
            \item This is the first item of an enumerated
                  list that is nested within the itemized list.

            \item This is the second item of the inner list.
                  \LaTeX\ allows you to nest lists deeper than
                  you really should.
         \end{enumerate}
         This is the rest of the second item of the outer
         list.  It is no more interesting than any other
         part of the item.
   \item This is the third item of the list.
\end{itemize}
You can even display poetry.
\begin{verse}
   There is an environment
    for verse \\             % The \\ command separates lines
   Whose features some poets % within a stanza.
   will curse.

                             % One or more blank lines separate stanzas.

   For instead of making\\
   Them do \emph{all} line breaking, \\
   It allows them to put too many words on a line when they'd rather be
   forced to be terse.
\end{verse}

Mathematical formulas may also be displayed.  A
displayed formula
is
one-line long; multiline
formulas require special formatting instructions.
   \[  \ip{\Gamma}{\psi'} = x'' + y^{2} + z_{i}^{n}\]
Don't start a paragraph with a displayed equation,
nor make one a paragraph by itself.

\end{document}               % End of document.

\section{Bench Huge Section 1447}
% synthetic bench section 1447
% This is a sample LaTeX input file.  (Version of 12 August 2004.)
%
% A '%' character causes TeX to ignore all remaining text on the line,
% and is used for comments like this one.

\documentclass{article}      % Specifies the document class

                             % The preamble begins here.
\title{An Example Document}  % Declares the document's title.
\author{Leslie Lamport}      % Declares the author's name.
\date{January 21, 1994}      % Deleting this command produces today's date.

\newcommand{\ip}[2]{(#1, #2)}
                             % Defines \ip{arg1}{arg2} to mean
                             % (arg1, arg2).

%\newcommand{\ip}[2]{\langle #1 | #2\rangle}
                             % This is an alternative definition of
                             % \ip that is commented out.

\begin{document}             % End of preamble and beginning of text.

\maketitle                   % Produces the title.

This is an example input file.  Comparing it with
the output it generates can show you how to
produce a simple document of your own.

\section{Ordinary Text}      % Produces section heading.  Lower-level
                             % sections are begun with similar
                             % \subsection and \subsubsection commands.

The ends  of words and sentences are marked
  by   spaces. It  doesn't matter how many
spaces    you type; one is as good as 100.  The
end of   a line counts as a space.

One   or more   blank lines denote the  end
of  a paragraph.

Since any number of consecutive spaces are treated
like a single one, the formatting of the input
file makes no difference to
      \LaTeX,                % The \LaTeX command generates the LaTeX logo.
but it makes a difference to you.  When you use
\LaTeX, making your input file as easy to read
as possible will be a great help as you write
your document and when you change it.  This sample
file shows how you can add comments to your own input
file.

Because printing is different from typewriting,
there are a number of things that you have to do
differently when preparing an input file than if
you were just typing the document directly.
Quotation marks like
       ``this''
have to be handled specially, as do quotes within
quotes:
       ``\,`this'            % \, separates the double and single quote.
        is what I just
        wrote, not  `that'\,''.

Dashes come in three sizes: an
       intra-word
dash, a medium dash for number ranges like
       1--2,
and a punctuation
       dash---like
this.

A sentence-ending space should be larger than the
space between words within a sentence.  You
sometimes have to type special commands in
conjunction with punctuation characters to get
this right, as in the following sentence.
       Gnats, gnus, etc.\ all  % `\ ' makes an inter-word space.
       begin with G\@.         % \@ marks end-of-sentence punctuation.
You should check the spaces after periods when
reading your output to make sure you haven't
forgotten any special cases.  Generating an
ellipsis
       \ldots\               % `\ ' is needed after `\ldots' because TeX
                             % ignores spaces after command names like \ldots
                             % made from \ + letters.
                             %
                             % Note how a `%' character causes TeX to ignore
                             % the end of the input line, so these blank lines
                             % do not start a new paragraph.
                             %
with the right spacing around the periods requires
a special command.

\LaTeX\ interprets some common characters as
commands, so you must type special commands to
generate them.  These characters include the
following:
       \$ \& \% \# \{ and \}.

In printing, text is usually emphasized with an
       \emph{italic}
type style.

\begin{em}
   A long segment of text can also be emphasized
   in this way.  Text within such a segment can be
   given \emph{additional} emphasis.
\end{em}

It is sometimes necessary to prevent \LaTeX\ from
breaking a line where it might otherwise do so.
This may be at a space, as between the ``Mr.''\ and
``Jones'' in
       ``Mr.~Jones'',        % ~ produces an unbreakable interword space.
or within a word---especially when the word is a
symbol like
       \mbox{\emph{itemnum}}
that makes little sense when hyphenated across
lines.

Footnotes\footnote{This is an example of a footnote.}
pose no problem.

\LaTeX\ is good at typesetting mathematical formulas
like
       \( x-3y + z = 7 \)
or
       \( a_{1} > x^{2n} + y^{2n} > x' \)
or
       \( \ip{A}{B} = \sum_{i} a_{i} b_{i} \).
The spaces you type in a formula are
ignored.  Remember that a letter like
       $x$                   % $ ... $  and  \( ... \)  are equivalent
is a formula when it denotes a mathematical
symbol, and it should be typed as one.

\section{Displayed Text}

Text is displayed by indenting it from the left
margin.  Quotations are commonly displayed.  There
are short quotations
\begin{quote}
   This is a short quotation.  It consists of a
   single paragraph of text.  See how it is formatted.
\end{quote}
and longer ones.
\begin{quotation}
   This is a longer quotation.  It consists of two
   paragraphs of text, neither of which are
   particularly interesting.

   This is the second paragraph of the quotation.  It
   is just as dull as the first paragraph.
\end{quotation}
Another frequently-displayed structure is a list.
The following is an example of an \emph{itemized}
list.
\begin{itemize}
   \item This is the first item of an itemized list.
         Each item in the list is marked with a ``tick''.
         You don't have to worry about what kind of tick
         mark is used.

   \item This is the second item of the list.  It
         contains another list nested inside it.  The inner
         list is an \emph{enumerated} list.
         \begin{enumerate}
            \item This is the first item of an enumerated
                  list that is nested within the itemized list.

            \item This is the second item of the inner list.
                  \LaTeX\ allows you to nest lists deeper than
                  you really should.
         \end{enumerate}
         This is the rest of the second item of the outer
         list.  It is no more interesting than any other
         part of the item.
   \item This is the third item of the list.
\end{itemize}
You can even display poetry.
\begin{verse}
   There is an environment
    for verse \\             % The \\ command separates lines
   Whose features some poets % within a stanza.
   will curse.

                             % One or more blank lines separate stanzas.

   For instead of making\\
   Them do \emph{all} line breaking, \\
   It allows them to put too many words on a line when they'd rather be
   forced to be terse.
\end{verse}

Mathematical formulas may also be displayed.  A
displayed formula
is
one-line long; multiline
formulas require special formatting instructions.
   \[  \ip{\Gamma}{\psi'} = x'' + y^{2} + z_{i}^{n}\]
Don't start a paragraph with a displayed equation,
nor make one a paragraph by itself.

\end{document}               % End of document.

\section{Bench Huge Section 1448}
% synthetic bench section 1448
% This is a sample LaTeX input file.  (Version of 12 August 2004.)
%
% A '%' character causes TeX to ignore all remaining text on the line,
% and is used for comments like this one.

\documentclass{article}      % Specifies the document class

                             % The preamble begins here.
\title{An Example Document}  % Declares the document's title.
\author{Leslie Lamport}      % Declares the author's name.
\date{January 21, 1994}      % Deleting this command produces today's date.

\newcommand{\ip}[2]{(#1, #2)}
                             % Defines \ip{arg1}{arg2} to mean
                             % (arg1, arg2).

%\newcommand{\ip}[2]{\langle #1 | #2\rangle}
                             % This is an alternative definition of
                             % \ip that is commented out.

\begin{document}             % End of preamble and beginning of text.

\maketitle                   % Produces the title.

This is an example input file.  Comparing it with
the output it generates can show you how to
produce a simple document of your own.

\section{Ordinary Text}      % Produces section heading.  Lower-level
                             % sections are begun with similar
                             % \subsection and \subsubsection commands.

The ends  of words and sentences are marked
  by   spaces. It  doesn't matter how many
spaces    you type; one is as good as 100.  The
end of   a line counts as a space.

One   or more   blank lines denote the  end
of  a paragraph.

Since any number of consecutive spaces are treated
like a single one, the formatting of the input
file makes no difference to
      \LaTeX,                % The \LaTeX command generates the LaTeX logo.
but it makes a difference to you.  When you use
\LaTeX, making your input file as easy to read
as possible will be a great help as you write
your document and when you change it.  This sample
file shows how you can add comments to your own input
file.

Because printing is different from typewriting,
there are a number of things that you have to do
differently when preparing an input file than if
you were just typing the document directly.
Quotation marks like
       ``this''
have to be handled specially, as do quotes within
quotes:
       ``\,`this'            % \, separates the double and single quote.
        is what I just
        wrote, not  `that'\,''.

Dashes come in three sizes: an
       intra-word
dash, a medium dash for number ranges like
       1--2,
and a punctuation
       dash---like
this.

A sentence-ending space should be larger than the
space between words within a sentence.  You
sometimes have to type special commands in
conjunction with punctuation characters to get
this right, as in the following sentence.
       Gnats, gnus, etc.\ all  % `\ ' makes an inter-word space.
       begin with G\@.         % \@ marks end-of-sentence punctuation.
You should check the spaces after periods when
reading your output to make sure you haven't
forgotten any special cases.  Generating an
ellipsis
       \ldots\               % `\ ' is needed after `\ldots' because TeX
                             % ignores spaces after command names like \ldots
                             % made from \ + letters.
                             %
                             % Note how a `%' character causes TeX to ignore
                             % the end of the input line, so these blank lines
                             % do not start a new paragraph.
                             %
with the right spacing around the periods requires
a special command.

\LaTeX\ interprets some common characters as
commands, so you must type special commands to
generate them.  These characters include the
following:
       \$ \& \% \# \{ and \}.

In printing, text is usually emphasized with an
       \emph{italic}
type style.

\begin{em}
   A long segment of text can also be emphasized
   in this way.  Text within such a segment can be
   given \emph{additional} emphasis.
\end{em}

It is sometimes necessary to prevent \LaTeX\ from
breaking a line where it might otherwise do so.
This may be at a space, as between the ``Mr.''\ and
``Jones'' in
       ``Mr.~Jones'',        % ~ produces an unbreakable interword space.
or within a word---especially when the word is a
symbol like
       \mbox{\emph{itemnum}}
that makes little sense when hyphenated across
lines.

Footnotes\footnote{This is an example of a footnote.}
pose no problem.

\LaTeX\ is good at typesetting mathematical formulas
like
       \( x-3y + z = 7 \)
or
       \( a_{1} > x^{2n} + y^{2n} > x' \)
or
       \( \ip{A}{B} = \sum_{i} a_{i} b_{i} \).
The spaces you type in a formula are
ignored.  Remember that a letter like
       $x$                   % $ ... $  and  \( ... \)  are equivalent
is a formula when it denotes a mathematical
symbol, and it should be typed as one.

\section{Displayed Text}

Text is displayed by indenting it from the left
margin.  Quotations are commonly displayed.  There
are short quotations
\begin{quote}
   This is a short quotation.  It consists of a
   single paragraph of text.  See how it is formatted.
\end{quote}
and longer ones.
\begin{quotation}
   This is a longer quotation.  It consists of two
   paragraphs of text, neither of which are
   particularly interesting.

   This is the second paragraph of the quotation.  It
   is just as dull as the first paragraph.
\end{quotation}
Another frequently-displayed structure is a list.
The following is an example of an \emph{itemized}
list.
\begin{itemize}
   \item This is the first item of an itemized list.
         Each item in the list is marked with a ``tick''.
         You don't have to worry about what kind of tick
         mark is used.

   \item This is the second item of the list.  It
         contains another list nested inside it.  The inner
         list is an \emph{enumerated} list.
         \begin{enumerate}
            \item This is the first item of an enumerated
                  list that is nested within the itemized list.

            \item This is the second item of the inner list.
                  \LaTeX\ allows you to nest lists deeper than
                  you really should.
         \end{enumerate}
         This is the rest of the second item of the outer
         list.  It is no more interesting than any other
         part of the item.
   \item This is the third item of the list.
\end{itemize}
You can even display poetry.
\begin{verse}
   There is an environment
    for verse \\             % The \\ command separates lines
   Whose features some poets % within a stanza.
   will curse.

                             % One or more blank lines separate stanzas.

   For instead of making\\
   Them do \emph{all} line breaking, \\
   It allows them to put too many words on a line when they'd rather be
   forced to be terse.
\end{verse}

Mathematical formulas may also be displayed.  A
displayed formula
is
one-line long; multiline
formulas require special formatting instructions.
   \[  \ip{\Gamma}{\psi'} = x'' + y^{2} + z_{i}^{n}\]
Don't start a paragraph with a displayed equation,
nor make one a paragraph by itself.

\end{document}               % End of document.

\section{Bench Huge Section 1449}
% synthetic bench section 1449
% This is a sample LaTeX input file.  (Version of 12 August 2004.)
%
% A '%' character causes TeX to ignore all remaining text on the line,
% and is used for comments like this one.

\documentclass{article}      % Specifies the document class

                             % The preamble begins here.
\title{An Example Document}  % Declares the document's title.
\author{Leslie Lamport}      % Declares the author's name.
\date{January 21, 1994}      % Deleting this command produces today's date.

\newcommand{\ip}[2]{(#1, #2)}
                             % Defines \ip{arg1}{arg2} to mean
                             % (arg1, arg2).

%\newcommand{\ip}[2]{\langle #1 | #2\rangle}
                             % This is an alternative definition of
                             % \ip that is commented out.

\begin{document}             % End of preamble and beginning of text.

\maketitle                   % Produces the title.

This is an example input file.  Comparing it with
the output it generates can show you how to
produce a simple document of your own.

\section{Ordinary Text}      % Produces section heading.  Lower-level
                             % sections are begun with similar
                             % \subsection and \subsubsection commands.

The ends  of words and sentences are marked
  by   spaces. It  doesn't matter how many
spaces    you type; one is as good as 100.  The
end of   a line counts as a space.

One   or more   blank lines denote the  end
of  a paragraph.

Since any number of consecutive spaces are treated
like a single one, the formatting of the input
file makes no difference to
      \LaTeX,                % The \LaTeX command generates the LaTeX logo.
but it makes a difference to you.  When you use
\LaTeX, making your input file as easy to read
as possible will be a great help as you write
your document and when you change it.  This sample
file shows how you can add comments to your own input
file.

Because printing is different from typewriting,
there are a number of things that you have to do
differently when preparing an input file than if
you were just typing the document directly.
Quotation marks like
       ``this''
have to be handled specially, as do quotes within
quotes:
       ``\,`this'            % \, separates the double and single quote.
        is what I just
        wrote, not  `that'\,''.

Dashes come in three sizes: an
       intra-word
dash, a medium dash for number ranges like
       1--2,
and a punctuation
       dash---like
this.

A sentence-ending space should be larger than the
space between words within a sentence.  You
sometimes have to type special commands in
conjunction with punctuation characters to get
this right, as in the following sentence.
       Gnats, gnus, etc.\ all  % `\ ' makes an inter-word space.
       begin with G\@.         % \@ marks end-of-sentence punctuation.
You should check the spaces after periods when
reading your output to make sure you haven't
forgotten any special cases.  Generating an
ellipsis
       \ldots\               % `\ ' is needed after `\ldots' because TeX
                             % ignores spaces after command names like \ldots
                             % made from \ + letters.
                             %
                             % Note how a `%' character causes TeX to ignore
                             % the end of the input line, so these blank lines
                             % do not start a new paragraph.
                             %
with the right spacing around the periods requires
a special command.

\LaTeX\ interprets some common characters as
commands, so you must type special commands to
generate them.  These characters include the
following:
       \$ \& \% \# \{ and \}.

In printing, text is usually emphasized with an
       \emph{italic}
type style.

\begin{em}
   A long segment of text can also be emphasized
   in this way.  Text within such a segment can be
   given \emph{additional} emphasis.
\end{em}

It is sometimes necessary to prevent \LaTeX\ from
breaking a line where it might otherwise do so.
This may be at a space, as between the ``Mr.''\ and
``Jones'' in
       ``Mr.~Jones'',        % ~ produces an unbreakable interword space.
or within a word---especially when the word is a
symbol like
       \mbox{\emph{itemnum}}
that makes little sense when hyphenated across
lines.

Footnotes\footnote{This is an example of a footnote.}
pose no problem.

\LaTeX\ is good at typesetting mathematical formulas
like
       \( x-3y + z = 7 \)
or
       \( a_{1} > x^{2n} + y^{2n} > x' \)
or
       \( \ip{A}{B} = \sum_{i} a_{i} b_{i} \).
The spaces you type in a formula are
ignored.  Remember that a letter like
       $x$                   % $ ... $  and  \( ... \)  are equivalent
is a formula when it denotes a mathematical
symbol, and it should be typed as one.

\section{Displayed Text}

Text is displayed by indenting it from the left
margin.  Quotations are commonly displayed.  There
are short quotations
\begin{quote}
   This is a short quotation.  It consists of a
   single paragraph of text.  See how it is formatted.
\end{quote}
and longer ones.
\begin{quotation}
   This is a longer quotation.  It consists of two
   paragraphs of text, neither of which are
   particularly interesting.

   This is the second paragraph of the quotation.  It
   is just as dull as the first paragraph.
\end{quotation}
Another frequently-displayed structure is a list.
The following is an example of an \emph{itemized}
list.
\begin{itemize}
   \item This is the first item of an itemized list.
         Each item in the list is marked with a ``tick''.
         You don't have to worry about what kind of tick
         mark is used.

   \item This is the second item of the list.  It
         contains another list nested inside it.  The inner
         list is an \emph{enumerated} list.
         \begin{enumerate}
            \item This is the first item of an enumerated
                  list that is nested within the itemized list.

            \item This is the second item of the inner list.
                  \LaTeX\ allows you to nest lists deeper than
                  you really should.
         \end{enumerate}
         This is the rest of the second item of the outer
         list.  It is no more interesting than any other
         part of the item.
   \item This is the third item of the list.
\end{itemize}
You can even display poetry.
\begin{verse}
   There is an environment
    for verse \\             % The \\ command separates lines
   Whose features some poets % within a stanza.
   will curse.

                             % One or more blank lines separate stanzas.

   For instead of making\\
   Them do \emph{all} line breaking, \\
   It allows them to put too many words on a line when they'd rather be
   forced to be terse.
\end{verse}

Mathematical formulas may also be displayed.  A
displayed formula
is
one-line long; multiline
formulas require special formatting instructions.
   \[  \ip{\Gamma}{\psi'} = x'' + y^{2} + z_{i}^{n}\]
Don't start a paragraph with a displayed equation,
nor make one a paragraph by itself.

\end{document}               % End of document.

\section{Bench Huge Section 1450}
% synthetic bench section 1450
% This is a sample LaTeX input file.  (Version of 12 August 2004.)
%
% A '%' character causes TeX to ignore all remaining text on the line,
% and is used for comments like this one.

\documentclass{article}      % Specifies the document class

                             % The preamble begins here.
\title{An Example Document}  % Declares the document's title.
\author{Leslie Lamport}      % Declares the author's name.
\date{January 21, 1994}      % Deleting this command produces today's date.

\newcommand{\ip}[2]{(#1, #2)}
                             % Defines \ip{arg1}{arg2} to mean
                             % (arg1, arg2).

%\newcommand{\ip}[2]{\langle #1 | #2\rangle}
                             % This is an alternative definition of
                             % \ip that is commented out.

\begin{document}             % End of preamble and beginning of text.

\maketitle                   % Produces the title.

This is an example input file.  Comparing it with
the output it generates can show you how to
produce a simple document of your own.

\section{Ordinary Text}      % Produces section heading.  Lower-level
                             % sections are begun with similar
                             % \subsection and \subsubsection commands.

The ends  of words and sentences are marked
  by   spaces. It  doesn't matter how many
spaces    you type; one is as good as 100.  The
end of   a line counts as a space.

One   or more   blank lines denote the  end
of  a paragraph.

Since any number of consecutive spaces are treated
like a single one, the formatting of the input
file makes no difference to
      \LaTeX,                % The \LaTeX command generates the LaTeX logo.
but it makes a difference to you.  When you use
\LaTeX, making your input file as easy to read
as possible will be a great help as you write
your document and when you change it.  This sample
file shows how you can add comments to your own input
file.

Because printing is different from typewriting,
there are a number of things that you have to do
differently when preparing an input file than if
you were just typing the document directly.
Quotation marks like
       ``this''
have to be handled specially, as do quotes within
quotes:
       ``\,`this'            % \, separates the double and single quote.
        is what I just
        wrote, not  `that'\,''.

Dashes come in three sizes: an
       intra-word
dash, a medium dash for number ranges like
       1--2,
and a punctuation
       dash---like
this.

A sentence-ending space should be larger than the
space between words within a sentence.  You
sometimes have to type special commands in
conjunction with punctuation characters to get
this right, as in the following sentence.
       Gnats, gnus, etc.\ all  % `\ ' makes an inter-word space.
       begin with G\@.         % \@ marks end-of-sentence punctuation.
You should check the spaces after periods when
reading your output to make sure you haven't
forgotten any special cases.  Generating an
ellipsis
       \ldots\               % `\ ' is needed after `\ldots' because TeX
                             % ignores spaces after command names like \ldots
                             % made from \ + letters.
                             %
                             % Note how a `%' character causes TeX to ignore
                             % the end of the input line, so these blank lines
                             % do not start a new paragraph.
                             %
with the right spacing around the periods requires
a special command.

\LaTeX\ interprets some common characters as
commands, so you must type special commands to
generate them.  These characters include the
following:
       \$ \& \% \# \{ and \}.

In printing, text is usually emphasized with an
       \emph{italic}
type style.

\begin{em}
   A long segment of text can also be emphasized
   in this way.  Text within such a segment can be
   given \emph{additional} emphasis.
\end{em}

It is sometimes necessary to prevent \LaTeX\ from
breaking a line where it might otherwise do so.
This may be at a space, as between the ``Mr.''\ and
``Jones'' in
       ``Mr.~Jones'',        % ~ produces an unbreakable interword space.
or within a word---especially when the word is a
symbol like
       \mbox{\emph{itemnum}}
that makes little sense when hyphenated across
lines.

Footnotes\footnote{This is an example of a footnote.}
pose no problem.

\LaTeX\ is good at typesetting mathematical formulas
like
       \( x-3y + z = 7 \)
or
       \( a_{1} > x^{2n} + y^{2n} > x' \)
or
       \( \ip{A}{B} = \sum_{i} a_{i} b_{i} \).
The spaces you type in a formula are
ignored.  Remember that a letter like
       $x$                   % $ ... $  and  \( ... \)  are equivalent
is a formula when it denotes a mathematical
symbol, and it should be typed as one.

\section{Displayed Text}

Text is displayed by indenting it from the left
margin.  Quotations are commonly displayed.  There
are short quotations
\begin{quote}
   This is a short quotation.  It consists of a
   single paragraph of text.  See how it is formatted.
\end{quote}
and longer ones.
\begin{quotation}
   This is a longer quotation.  It consists of two
   paragraphs of text, neither of which are
   particularly interesting.

   This is the second paragraph of the quotation.  It
   is just as dull as the first paragraph.
\end{quotation}
Another frequently-displayed structure is a list.
The following is an example of an \emph{itemized}
list.
\begin{itemize}
   \item This is the first item of an itemized list.
         Each item in the list is marked with a ``tick''.
         You don't have to worry about what kind of tick
         mark is used.

   \item This is the second item of the list.  It
         contains another list nested inside it.  The inner
         list is an \emph{enumerated} list.
         \begin{enumerate}
            \item This is the first item of an enumerated
                  list that is nested within the itemized list.

            \item This is the second item of the inner list.
                  \LaTeX\ allows you to nest lists deeper than
                  you really should.
         \end{enumerate}
         This is the rest of the second item of the outer
         list.  It is no more interesting than any other
         part of the item.
   \item This is the third item of the list.
\end{itemize}
You can even display poetry.
\begin{verse}
   There is an environment
    for verse \\             % The \\ command separates lines
   Whose features some poets % within a stanza.
   will curse.

                             % One or more blank lines separate stanzas.

   For instead of making\\
   Them do \emph{all} line breaking, \\
   It allows them to put too many words on a line when they'd rather be
   forced to be terse.
\end{verse}

Mathematical formulas may also be displayed.  A
displayed formula
is
one-line long; multiline
formulas require special formatting instructions.
   \[  \ip{\Gamma}{\psi'} = x'' + y^{2} + z_{i}^{n}\]
Don't start a paragraph with a displayed equation,
nor make one a paragraph by itself.

\end{document}               % End of document.

\section{Bench Huge Section 1451}
% synthetic bench section 1451
% This is a sample LaTeX input file.  (Version of 12 August 2004.)
%
% A '%' character causes TeX to ignore all remaining text on the line,
% and is used for comments like this one.

\documentclass{article}      % Specifies the document class

                             % The preamble begins here.
\title{An Example Document}  % Declares the document's title.
\author{Leslie Lamport}      % Declares the author's name.
\date{January 21, 1994}      % Deleting this command produces today's date.

\newcommand{\ip}[2]{(#1, #2)}
                             % Defines \ip{arg1}{arg2} to mean
                             % (arg1, arg2).

%\newcommand{\ip}[2]{\langle #1 | #2\rangle}
                             % This is an alternative definition of
                             % \ip that is commented out.

\begin{document}             % End of preamble and beginning of text.

\maketitle                   % Produces the title.

This is an example input file.  Comparing it with
the output it generates can show you how to
produce a simple document of your own.

\section{Ordinary Text}      % Produces section heading.  Lower-level
                             % sections are begun with similar
                             % \subsection and \subsubsection commands.

The ends  of words and sentences are marked
  by   spaces. It  doesn't matter how many
spaces    you type; one is as good as 100.  The
end of   a line counts as a space.

One   or more   blank lines denote the  end
of  a paragraph.

Since any number of consecutive spaces are treated
like a single one, the formatting of the input
file makes no difference to
      \LaTeX,                % The \LaTeX command generates the LaTeX logo.
but it makes a difference to you.  When you use
\LaTeX, making your input file as easy to read
as possible will be a great help as you write
your document and when you change it.  This sample
file shows how you can add comments to your own input
file.

Because printing is different from typewriting,
there are a number of things that you have to do
differently when preparing an input file than if
you were just typing the document directly.
Quotation marks like
       ``this''
have to be handled specially, as do quotes within
quotes:
       ``\,`this'            % \, separates the double and single quote.
        is what I just
        wrote, not  `that'\,''.

Dashes come in three sizes: an
       intra-word
dash, a medium dash for number ranges like
       1--2,
and a punctuation
       dash---like
this.

A sentence-ending space should be larger than the
space between words within a sentence.  You
sometimes have to type special commands in
conjunction with punctuation characters to get
this right, as in the following sentence.
       Gnats, gnus, etc.\ all  % `\ ' makes an inter-word space.
       begin with G\@.         % \@ marks end-of-sentence punctuation.
You should check the spaces after periods when
reading your output to make sure you haven't
forgotten any special cases.  Generating an
ellipsis
       \ldots\               % `\ ' is needed after `\ldots' because TeX
                             % ignores spaces after command names like \ldots
                             % made from \ + letters.
                             %
                             % Note how a `%' character causes TeX to ignore
                             % the end of the input line, so these blank lines
                             % do not start a new paragraph.
                             %
with the right spacing around the periods requires
a special command.

\LaTeX\ interprets some common characters as
commands, so you must type special commands to
generate them.  These characters include the
following:
       \$ \& \% \# \{ and \}.

In printing, text is usually emphasized with an
       \emph{italic}
type style.

\begin{em}
   A long segment of text can also be emphasized
   in this way.  Text within such a segment can be
   given \emph{additional} emphasis.
\end{em}

It is sometimes necessary to prevent \LaTeX\ from
breaking a line where it might otherwise do so.
This may be at a space, as between the ``Mr.''\ and
``Jones'' in
       ``Mr.~Jones'',        % ~ produces an unbreakable interword space.
or within a word---especially when the word is a
symbol like
       \mbox{\emph{itemnum}}
that makes little sense when hyphenated across
lines.

Footnotes\footnote{This is an example of a footnote.}
pose no problem.

\LaTeX\ is good at typesetting mathematical formulas
like
       \( x-3y + z = 7 \)
or
       \( a_{1} > x^{2n} + y^{2n} > x' \)
or
       \( \ip{A}{B} = \sum_{i} a_{i} b_{i} \).
The spaces you type in a formula are
ignored.  Remember that a letter like
       $x$                   % $ ... $  and  \( ... \)  are equivalent
is a formula when it denotes a mathematical
symbol, and it should be typed as one.

\section{Displayed Text}

Text is displayed by indenting it from the left
margin.  Quotations are commonly displayed.  There
are short quotations
\begin{quote}
   This is a short quotation.  It consists of a
   single paragraph of text.  See how it is formatted.
\end{quote}
and longer ones.
\begin{quotation}
   This is a longer quotation.  It consists of two
   paragraphs of text, neither of which are
   particularly interesting.

   This is the second paragraph of the quotation.  It
   is just as dull as the first paragraph.
\end{quotation}
Another frequently-displayed structure is a list.
The following is an example of an \emph{itemized}
list.
\begin{itemize}
   \item This is the first item of an itemized list.
         Each item in the list is marked with a ``tick''.
         You don't have to worry about what kind of tick
         mark is used.

   \item This is the second item of the list.  It
         contains another list nested inside it.  The inner
         list is an \emph{enumerated} list.
         \begin{enumerate}
            \item This is the first item of an enumerated
                  list that is nested within the itemized list.

            \item This is the second item of the inner list.
                  \LaTeX\ allows you to nest lists deeper than
                  you really should.
         \end{enumerate}
         This is the rest of the second item of the outer
         list.  It is no more interesting than any other
         part of the item.
   \item This is the third item of the list.
\end{itemize}
You can even display poetry.
\begin{verse}
   There is an environment
    for verse \\             % The \\ command separates lines
   Whose features some poets % within a stanza.
   will curse.

                             % One or more blank lines separate stanzas.

   For instead of making\\
   Them do \emph{all} line breaking, \\
   It allows them to put too many words on a line when they'd rather be
   forced to be terse.
\end{verse}

Mathematical formulas may also be displayed.  A
displayed formula
is
one-line long; multiline
formulas require special formatting instructions.
   \[  \ip{\Gamma}{\psi'} = x'' + y^{2} + z_{i}^{n}\]
Don't start a paragraph with a displayed equation,
nor make one a paragraph by itself.

\end{document}               % End of document.

\section{Bench Huge Section 1452}
% synthetic bench section 1452
% This is a sample LaTeX input file.  (Version of 12 August 2004.)
%
% A '%' character causes TeX to ignore all remaining text on the line,
% and is used for comments like this one.

\documentclass{article}      % Specifies the document class

                             % The preamble begins here.
\title{An Example Document}  % Declares the document's title.
\author{Leslie Lamport}      % Declares the author's name.
\date{January 21, 1994}      % Deleting this command produces today's date.

\newcommand{\ip}[2]{(#1, #2)}
                             % Defines \ip{arg1}{arg2} to mean
                             % (arg1, arg2).

%\newcommand{\ip}[2]{\langle #1 | #2\rangle}
                             % This is an alternative definition of
                             % \ip that is commented out.

\begin{document}             % End of preamble and beginning of text.

\maketitle                   % Produces the title.

This is an example input file.  Comparing it with
the output it generates can show you how to
produce a simple document of your own.

\section{Ordinary Text}      % Produces section heading.  Lower-level
                             % sections are begun with similar
                             % \subsection and \subsubsection commands.

The ends  of words and sentences are marked
  by   spaces. It  doesn't matter how many
spaces    you type; one is as good as 100.  The
end of   a line counts as a space.

One   or more   blank lines denote the  end
of  a paragraph.

Since any number of consecutive spaces are treated
like a single one, the formatting of the input
file makes no difference to
      \LaTeX,                % The \LaTeX command generates the LaTeX logo.
but it makes a difference to you.  When you use
\LaTeX, making your input file as easy to read
as possible will be a great help as you write
your document and when you change it.  This sample
file shows how you can add comments to your own input
file.

Because printing is different from typewriting,
there are a number of things that you have to do
differently when preparing an input file than if
you were just typing the document directly.
Quotation marks like
       ``this''
have to be handled specially, as do quotes within
quotes:
       ``\,`this'            % \, separates the double and single quote.
        is what I just
        wrote, not  `that'\,''.

Dashes come in three sizes: an
       intra-word
dash, a medium dash for number ranges like
       1--2,
and a punctuation
       dash---like
this.

A sentence-ending space should be larger than the
space between words within a sentence.  You
sometimes have to type special commands in
conjunction with punctuation characters to get
this right, as in the following sentence.
       Gnats, gnus, etc.\ all  % `\ ' makes an inter-word space.
       begin with G\@.         % \@ marks end-of-sentence punctuation.
You should check the spaces after periods when
reading your output to make sure you haven't
forgotten any special cases.  Generating an
ellipsis
       \ldots\               % `\ ' is needed after `\ldots' because TeX
                             % ignores spaces after command names like \ldots
                             % made from \ + letters.
                             %
                             % Note how a `%' character causes TeX to ignore
                             % the end of the input line, so these blank lines
                             % do not start a new paragraph.
                             %
with the right spacing around the periods requires
a special command.

\LaTeX\ interprets some common characters as
commands, so you must type special commands to
generate them.  These characters include the
following:
       \$ \& \% \# \{ and \}.

In printing, text is usually emphasized with an
       \emph{italic}
type style.

\begin{em}
   A long segment of text can also be emphasized
   in this way.  Text within such a segment can be
   given \emph{additional} emphasis.
\end{em}

It is sometimes necessary to prevent \LaTeX\ from
breaking a line where it might otherwise do so.
This may be at a space, as between the ``Mr.''\ and
``Jones'' in
       ``Mr.~Jones'',        % ~ produces an unbreakable interword space.
or within a word---especially when the word is a
symbol like
       \mbox{\emph{itemnum}}
that makes little sense when hyphenated across
lines.

Footnotes\footnote{This is an example of a footnote.}
pose no problem.

\LaTeX\ is good at typesetting mathematical formulas
like
       \( x-3y + z = 7 \)
or
       \( a_{1} > x^{2n} + y^{2n} > x' \)
or
       \( \ip{A}{B} = \sum_{i} a_{i} b_{i} \).
The spaces you type in a formula are
ignored.  Remember that a letter like
       $x$                   % $ ... $  and  \( ... \)  are equivalent
is a formula when it denotes a mathematical
symbol, and it should be typed as one.

\section{Displayed Text}

Text is displayed by indenting it from the left
margin.  Quotations are commonly displayed.  There
are short quotations
\begin{quote}
   This is a short quotation.  It consists of a
   single paragraph of text.  See how it is formatted.
\end{quote}
and longer ones.
\begin{quotation}
   This is a longer quotation.  It consists of two
   paragraphs of text, neither of which are
   particularly interesting.

   This is the second paragraph of the quotation.  It
   is just as dull as the first paragraph.
\end{quotation}
Another frequently-displayed structure is a list.
The following is an example of an \emph{itemized}
list.
\begin{itemize}
   \item This is the first item of an itemized list.
         Each item in the list is marked with a ``tick''.
         You don't have to worry about what kind of tick
         mark is used.

   \item This is the second item of the list.  It
         contains another list nested inside it.  The inner
         list is an \emph{enumerated} list.
         \begin{enumerate}
            \item This is the first item of an enumerated
                  list that is nested within the itemized list.

            \item This is the second item of the inner list.
                  \LaTeX\ allows you to nest lists deeper than
                  you really should.
         \end{enumerate}
         This is the rest of the second item of the outer
         list.  It is no more interesting than any other
         part of the item.
   \item This is the third item of the list.
\end{itemize}
You can even display poetry.
\begin{verse}
   There is an environment
    for verse \\             % The \\ command separates lines
   Whose features some poets % within a stanza.
   will curse.

                             % One or more blank lines separate stanzas.

   For instead of making\\
   Them do \emph{all} line breaking, \\
   It allows them to put too many words on a line when they'd rather be
   forced to be terse.
\end{verse}

Mathematical formulas may also be displayed.  A
displayed formula
is
one-line long; multiline
formulas require special formatting instructions.
   \[  \ip{\Gamma}{\psi'} = x'' + y^{2} + z_{i}^{n}\]
Don't start a paragraph with a displayed equation,
nor make one a paragraph by itself.

\end{document}               % End of document.

\section{Bench Huge Section 1453}
% synthetic bench section 1453
% This is a sample LaTeX input file.  (Version of 12 August 2004.)
%
% A '%' character causes TeX to ignore all remaining text on the line,
% and is used for comments like this one.

\documentclass{article}      % Specifies the document class

                             % The preamble begins here.
\title{An Example Document}  % Declares the document's title.
\author{Leslie Lamport}      % Declares the author's name.
\date{January 21, 1994}      % Deleting this command produces today's date.

\newcommand{\ip}[2]{(#1, #2)}
                             % Defines \ip{arg1}{arg2} to mean
                             % (arg1, arg2).

%\newcommand{\ip}[2]{\langle #1 | #2\rangle}
                             % This is an alternative definition of
                             % \ip that is commented out.

\begin{document}             % End of preamble and beginning of text.

\maketitle                   % Produces the title.

This is an example input file.  Comparing it with
the output it generates can show you how to
produce a simple document of your own.

\section{Ordinary Text}      % Produces section heading.  Lower-level
                             % sections are begun with similar
                             % \subsection and \subsubsection commands.

The ends  of words and sentences are marked
  by   spaces. It  doesn't matter how many
spaces    you type; one is as good as 100.  The
end of   a line counts as a space.

One   or more   blank lines denote the  end
of  a paragraph.

Since any number of consecutive spaces are treated
like a single one, the formatting of the input
file makes no difference to
      \LaTeX,                % The \LaTeX command generates the LaTeX logo.
but it makes a difference to you.  When you use
\LaTeX, making your input file as easy to read
as possible will be a great help as you write
your document and when you change it.  This sample
file shows how you can add comments to your own input
file.

Because printing is different from typewriting,
there are a number of things that you have to do
differently when preparing an input file than if
you were just typing the document directly.
Quotation marks like
       ``this''
have to be handled specially, as do quotes within
quotes:
       ``\,`this'            % \, separates the double and single quote.
        is what I just
        wrote, not  `that'\,''.

Dashes come in three sizes: an
       intra-word
dash, a medium dash for number ranges like
       1--2,
and a punctuation
       dash---like
this.

A sentence-ending space should be larger than the
space between words within a sentence.  You
sometimes have to type special commands in
conjunction with punctuation characters to get
this right, as in the following sentence.
       Gnats, gnus, etc.\ all  % `\ ' makes an inter-word space.
       begin with G\@.         % \@ marks end-of-sentence punctuation.
You should check the spaces after periods when
reading your output to make sure you haven't
forgotten any special cases.  Generating an
ellipsis
       \ldots\               % `\ ' is needed after `\ldots' because TeX
                             % ignores spaces after command names like \ldots
                             % made from \ + letters.
                             %
                             % Note how a `%' character causes TeX to ignore
                             % the end of the input line, so these blank lines
                             % do not start a new paragraph.
                             %
with the right spacing around the periods requires
a special command.

\LaTeX\ interprets some common characters as
commands, so you must type special commands to
generate them.  These characters include the
following:
       \$ \& \% \# \{ and \}.

In printing, text is usually emphasized with an
       \emph{italic}
type style.

\begin{em}
   A long segment of text can also be emphasized
   in this way.  Text within such a segment can be
   given \emph{additional} emphasis.
\end{em}

It is sometimes necessary to prevent \LaTeX\ from
breaking a line where it might otherwise do so.
This may be at a space, as between the ``Mr.''\ and
``Jones'' in
       ``Mr.~Jones'',        % ~ produces an unbreakable interword space.
or within a word---especially when the word is a
symbol like
       \mbox{\emph{itemnum}}
that makes little sense when hyphenated across
lines.

Footnotes\footnote{This is an example of a footnote.}
pose no problem.

\LaTeX\ is good at typesetting mathematical formulas
like
       \( x-3y + z = 7 \)
or
       \( a_{1} > x^{2n} + y^{2n} > x' \)
or
       \( \ip{A}{B} = \sum_{i} a_{i} b_{i} \).
The spaces you type in a formula are
ignored.  Remember that a letter like
       $x$                   % $ ... $  and  \( ... \)  are equivalent
is a formula when it denotes a mathematical
symbol, and it should be typed as one.

\section{Displayed Text}

Text is displayed by indenting it from the left
margin.  Quotations are commonly displayed.  There
are short quotations
\begin{quote}
   This is a short quotation.  It consists of a
   single paragraph of text.  See how it is formatted.
\end{quote}
and longer ones.
\begin{quotation}
   This is a longer quotation.  It consists of two
   paragraphs of text, neither of which are
   particularly interesting.

   This is the second paragraph of the quotation.  It
   is just as dull as the first paragraph.
\end{quotation}
Another frequently-displayed structure is a list.
The following is an example of an \emph{itemized}
list.
\begin{itemize}
   \item This is the first item of an itemized list.
         Each item in the list is marked with a ``tick''.
         You don't have to worry about what kind of tick
         mark is used.

   \item This is the second item of the list.  It
         contains another list nested inside it.  The inner
         list is an \emph{enumerated} list.
         \begin{enumerate}
            \item This is the first item of an enumerated
                  list that is nested within the itemized list.

            \item This is the second item of the inner list.
                  \LaTeX\ allows you to nest lists deeper than
                  you really should.
         \end{enumerate}
         This is the rest of the second item of the outer
         list.  It is no more interesting than any other
         part of the item.
   \item This is the third item of the list.
\end{itemize}
You can even display poetry.
\begin{verse}
   There is an environment
    for verse \\             % The \\ command separates lines
   Whose features some poets % within a stanza.
   will curse.

                             % One or more blank lines separate stanzas.

   For instead of making\\
   Them do \emph{all} line breaking, \\
   It allows them to put too many words on a line when they'd rather be
   forced to be terse.
\end{verse}

Mathematical formulas may also be displayed.  A
displayed formula
is
one-line long; multiline
formulas require special formatting instructions.
   \[  \ip{\Gamma}{\psi'} = x'' + y^{2} + z_{i}^{n}\]
Don't start a paragraph with a displayed equation,
nor make one a paragraph by itself.

\end{document}               % End of document.

\section{Bench Huge Section 1454}
% synthetic bench section 1454
% This is a sample LaTeX input file.  (Version of 12 August 2004.)
%
% A '%' character causes TeX to ignore all remaining text on the line,
% and is used for comments like this one.

\documentclass{article}      % Specifies the document class

                             % The preamble begins here.
\title{An Example Document}  % Declares the document's title.
\author{Leslie Lamport}      % Declares the author's name.
\date{January 21, 1994}      % Deleting this command produces today's date.

\newcommand{\ip}[2]{(#1, #2)}
                             % Defines \ip{arg1}{arg2} to mean
                             % (arg1, arg2).

%\newcommand{\ip}[2]{\langle #1 | #2\rangle}
                             % This is an alternative definition of
                             % \ip that is commented out.

\begin{document}             % End of preamble and beginning of text.

\maketitle                   % Produces the title.

This is an example input file.  Comparing it with
the output it generates can show you how to
produce a simple document of your own.

\section{Ordinary Text}      % Produces section heading.  Lower-level
                             % sections are begun with similar
                             % \subsection and \subsubsection commands.

The ends  of words and sentences are marked
  by   spaces. It  doesn't matter how many
spaces    you type; one is as good as 100.  The
end of   a line counts as a space.

One   or more   blank lines denote the  end
of  a paragraph.

Since any number of consecutive spaces are treated
like a single one, the formatting of the input
file makes no difference to
      \LaTeX,                % The \LaTeX command generates the LaTeX logo.
but it makes a difference to you.  When you use
\LaTeX, making your input file as easy to read
as possible will be a great help as you write
your document and when you change it.  This sample
file shows how you can add comments to your own input
file.

Because printing is different from typewriting,
there are a number of things that you have to do
differently when preparing an input file than if
you were just typing the document directly.
Quotation marks like
       ``this''
have to be handled specially, as do quotes within
quotes:
       ``\,`this'            % \, separates the double and single quote.
        is what I just
        wrote, not  `that'\,''.

Dashes come in three sizes: an
       intra-word
dash, a medium dash for number ranges like
       1--2,
and a punctuation
       dash---like
this.

A sentence-ending space should be larger than the
space between words within a sentence.  You
sometimes have to type special commands in
conjunction with punctuation characters to get
this right, as in the following sentence.
       Gnats, gnus, etc.\ all  % `\ ' makes an inter-word space.
       begin with G\@.         % \@ marks end-of-sentence punctuation.
You should check the spaces after periods when
reading your output to make sure you haven't
forgotten any special cases.  Generating an
ellipsis
       \ldots\               % `\ ' is needed after `\ldots' because TeX
                             % ignores spaces after command names like \ldots
                             % made from \ + letters.
                             %
                             % Note how a `%' character causes TeX to ignore
                             % the end of the input line, so these blank lines
                             % do not start a new paragraph.
                             %
with the right spacing around the periods requires
a special command.

\LaTeX\ interprets some common characters as
commands, so you must type special commands to
generate them.  These characters include the
following:
       \$ \& \% \# \{ and \}.

In printing, text is usually emphasized with an
       \emph{italic}
type style.

\begin{em}
   A long segment of text can also be emphasized
   in this way.  Text within such a segment can be
   given \emph{additional} emphasis.
\end{em}

It is sometimes necessary to prevent \LaTeX\ from
breaking a line where it might otherwise do so.
This may be at a space, as between the ``Mr.''\ and
``Jones'' in
       ``Mr.~Jones'',        % ~ produces an unbreakable interword space.
or within a word---especially when the word is a
symbol like
       \mbox{\emph{itemnum}}
that makes little sense when hyphenated across
lines.

Footnotes\footnote{This is an example of a footnote.}
pose no problem.

\LaTeX\ is good at typesetting mathematical formulas
like
       \( x-3y + z = 7 \)
or
       \( a_{1} > x^{2n} + y^{2n} > x' \)
or
       \( \ip{A}{B} = \sum_{i} a_{i} b_{i} \).
The spaces you type in a formula are
ignored.  Remember that a letter like
       $x$                   % $ ... $  and  \( ... \)  are equivalent
is a formula when it denotes a mathematical
symbol, and it should be typed as one.

\section{Displayed Text}

Text is displayed by indenting it from the left
margin.  Quotations are commonly displayed.  There
are short quotations
\begin{quote}
   This is a short quotation.  It consists of a
   single paragraph of text.  See how it is formatted.
\end{quote}
and longer ones.
\begin{quotation}
   This is a longer quotation.  It consists of two
   paragraphs of text, neither of which are
   particularly interesting.

   This is the second paragraph of the quotation.  It
   is just as dull as the first paragraph.
\end{quotation}
Another frequently-displayed structure is a list.
The following is an example of an \emph{itemized}
list.
\begin{itemize}
   \item This is the first item of an itemized list.
         Each item in the list is marked with a ``tick''.
         You don't have to worry about what kind of tick
         mark is used.

   \item This is the second item of the list.  It
         contains another list nested inside it.  The inner
         list is an \emph{enumerated} list.
         \begin{enumerate}
            \item This is the first item of an enumerated
                  list that is nested within the itemized list.

            \item This is the second item of the inner list.
                  \LaTeX\ allows you to nest lists deeper than
                  you really should.
         \end{enumerate}
         This is the rest of the second item of the outer
         list.  It is no more interesting than any other
         part of the item.
   \item This is the third item of the list.
\end{itemize}
You can even display poetry.
\begin{verse}
   There is an environment
    for verse \\             % The \\ command separates lines
   Whose features some poets % within a stanza.
   will curse.

                             % One or more blank lines separate stanzas.

   For instead of making\\
   Them do \emph{all} line breaking, \\
   It allows them to put too many words on a line when they'd rather be
   forced to be terse.
\end{verse}

Mathematical formulas may also be displayed.  A
displayed formula
is
one-line long; multiline
formulas require special formatting instructions.
   \[  \ip{\Gamma}{\psi'} = x'' + y^{2} + z_{i}^{n}\]
Don't start a paragraph with a displayed equation,
nor make one a paragraph by itself.

\end{document}               % End of document.

\section{Bench Huge Section 1455}
% synthetic bench section 1455
% This is a sample LaTeX input file.  (Version of 12 August 2004.)
%
% A '%' character causes TeX to ignore all remaining text on the line,
% and is used for comments like this one.

\documentclass{article}      % Specifies the document class

                             % The preamble begins here.
\title{An Example Document}  % Declares the document's title.
\author{Leslie Lamport}      % Declares the author's name.
\date{January 21, 1994}      % Deleting this command produces today's date.

\newcommand{\ip}[2]{(#1, #2)}
                             % Defines \ip{arg1}{arg2} to mean
                             % (arg1, arg2).

%\newcommand{\ip}[2]{\langle #1 | #2\rangle}
                             % This is an alternative definition of
                             % \ip that is commented out.

\begin{document}             % End of preamble and beginning of text.

\maketitle                   % Produces the title.

This is an example input file.  Comparing it with
the output it generates can show you how to
produce a simple document of your own.

\section{Ordinary Text}      % Produces section heading.  Lower-level
                             % sections are begun with similar
                             % \subsection and \subsubsection commands.

The ends  of words and sentences are marked
  by   spaces. It  doesn't matter how many
spaces    you type; one is as good as 100.  The
end of   a line counts as a space.

One   or more   blank lines denote the  end
of  a paragraph.

Since any number of consecutive spaces are treated
like a single one, the formatting of the input
file makes no difference to
      \LaTeX,                % The \LaTeX command generates the LaTeX logo.
but it makes a difference to you.  When you use
\LaTeX, making your input file as easy to read
as possible will be a great help as you write
your document and when you change it.  This sample
file shows how you can add comments to your own input
file.

Because printing is different from typewriting,
there are a number of things that you have to do
differently when preparing an input file than if
you were just typing the document directly.
Quotation marks like
       ``this''
have to be handled specially, as do quotes within
quotes:
       ``\,`this'            % \, separates the double and single quote.
        is what I just
        wrote, not  `that'\,''.

Dashes come in three sizes: an
       intra-word
dash, a medium dash for number ranges like
       1--2,
and a punctuation
       dash---like
this.

A sentence-ending space should be larger than the
space between words within a sentence.  You
sometimes have to type special commands in
conjunction with punctuation characters to get
this right, as in the following sentence.
       Gnats, gnus, etc.\ all  % `\ ' makes an inter-word space.
       begin with G\@.         % \@ marks end-of-sentence punctuation.
You should check the spaces after periods when
reading your output to make sure you haven't
forgotten any special cases.  Generating an
ellipsis
       \ldots\               % `\ ' is needed after `\ldots' because TeX
                             % ignores spaces after command names like \ldots
                             % made from \ + letters.
                             %
                             % Note how a `%' character causes TeX to ignore
                             % the end of the input line, so these blank lines
                             % do not start a new paragraph.
                             %
with the right spacing around the periods requires
a special command.

\LaTeX\ interprets some common characters as
commands, so you must type special commands to
generate them.  These characters include the
following:
       \$ \& \% \# \{ and \}.

In printing, text is usually emphasized with an
       \emph{italic}
type style.

\begin{em}
   A long segment of text can also be emphasized
   in this way.  Text within such a segment can be
   given \emph{additional} emphasis.
\end{em}

It is sometimes necessary to prevent \LaTeX\ from
breaking a line where it might otherwise do so.
This may be at a space, as between the ``Mr.''\ and
``Jones'' in
       ``Mr.~Jones'',        % ~ produces an unbreakable interword space.
or within a word---especially when the word is a
symbol like
       \mbox{\emph{itemnum}}
that makes little sense when hyphenated across
lines.

Footnotes\footnote{This is an example of a footnote.}
pose no problem.

\LaTeX\ is good at typesetting mathematical formulas
like
       \( x-3y + z = 7 \)
or
       \( a_{1} > x^{2n} + y^{2n} > x' \)
or
       \( \ip{A}{B} = \sum_{i} a_{i} b_{i} \).
The spaces you type in a formula are
ignored.  Remember that a letter like
       $x$                   % $ ... $  and  \( ... \)  are equivalent
is a formula when it denotes a mathematical
symbol, and it should be typed as one.

\section{Displayed Text}

Text is displayed by indenting it from the left
margin.  Quotations are commonly displayed.  There
are short quotations
\begin{quote}
   This is a short quotation.  It consists of a
   single paragraph of text.  See how it is formatted.
\end{quote}
and longer ones.
\begin{quotation}
   This is a longer quotation.  It consists of two
   paragraphs of text, neither of which are
   particularly interesting.

   This is the second paragraph of the quotation.  It
   is just as dull as the first paragraph.
\end{quotation}
Another frequently-displayed structure is a list.
The following is an example of an \emph{itemized}
list.
\begin{itemize}
   \item This is the first item of an itemized list.
         Each item in the list is marked with a ``tick''.
         You don't have to worry about what kind of tick
         mark is used.

   \item This is the second item of the list.  It
         contains another list nested inside it.  The inner
         list is an \emph{enumerated} list.
         \begin{enumerate}
            \item This is the first item of an enumerated
                  list that is nested within the itemized list.

            \item This is the second item of the inner list.
                  \LaTeX\ allows you to nest lists deeper than
                  you really should.
         \end{enumerate}
         This is the rest of the second item of the outer
         list.  It is no more interesting than any other
         part of the item.
   \item This is the third item of the list.
\end{itemize}
You can even display poetry.
\begin{verse}
   There is an environment
    for verse \\             % The \\ command separates lines
   Whose features some poets % within a stanza.
   will curse.

                             % One or more blank lines separate stanzas.

   For instead of making\\
   Them do \emph{all} line breaking, \\
   It allows them to put too many words on a line when they'd rather be
   forced to be terse.
\end{verse}

Mathematical formulas may also be displayed.  A
displayed formula
is
one-line long; multiline
formulas require special formatting instructions.
   \[  \ip{\Gamma}{\psi'} = x'' + y^{2} + z_{i}^{n}\]
Don't start a paragraph with a displayed equation,
nor make one a paragraph by itself.

\end{document}               % End of document.

\section{Bench Huge Section 1456}
% synthetic bench section 1456
% This is a sample LaTeX input file.  (Version of 12 August 2004.)
%
% A '%' character causes TeX to ignore all remaining text on the line,
% and is used for comments like this one.

\documentclass{article}      % Specifies the document class

                             % The preamble begins here.
\title{An Example Document}  % Declares the document's title.
\author{Leslie Lamport}      % Declares the author's name.
\date{January 21, 1994}      % Deleting this command produces today's date.

\newcommand{\ip}[2]{(#1, #2)}
                             % Defines \ip{arg1}{arg2} to mean
                             % (arg1, arg2).

%\newcommand{\ip}[2]{\langle #1 | #2\rangle}
                             % This is an alternative definition of
                             % \ip that is commented out.

\begin{document}             % End of preamble and beginning of text.

\maketitle                   % Produces the title.

This is an example input file.  Comparing it with
the output it generates can show you how to
produce a simple document of your own.

\section{Ordinary Text}      % Produces section heading.  Lower-level
                             % sections are begun with similar
                             % \subsection and \subsubsection commands.

The ends  of words and sentences are marked
  by   spaces. It  doesn't matter how many
spaces    you type; one is as good as 100.  The
end of   a line counts as a space.

One   or more   blank lines denote the  end
of  a paragraph.

Since any number of consecutive spaces are treated
like a single one, the formatting of the input
file makes no difference to
      \LaTeX,                % The \LaTeX command generates the LaTeX logo.
but it makes a difference to you.  When you use
\LaTeX, making your input file as easy to read
as possible will be a great help as you write
your document and when you change it.  This sample
file shows how you can add comments to your own input
file.

Because printing is different from typewriting,
there are a number of things that you have to do
differently when preparing an input file than if
you were just typing the document directly.
Quotation marks like
       ``this''
have to be handled specially, as do quotes within
quotes:
       ``\,`this'            % \, separates the double and single quote.
        is what I just
        wrote, not  `that'\,''.

Dashes come in three sizes: an
       intra-word
dash, a medium dash for number ranges like
       1--2,
and a punctuation
       dash---like
this.

A sentence-ending space should be larger than the
space between words within a sentence.  You
sometimes have to type special commands in
conjunction with punctuation characters to get
this right, as in the following sentence.
       Gnats, gnus, etc.\ all  % `\ ' makes an inter-word space.
       begin with G\@.         % \@ marks end-of-sentence punctuation.
You should check the spaces after periods when
reading your output to make sure you haven't
forgotten any special cases.  Generating an
ellipsis
       \ldots\               % `\ ' is needed after `\ldots' because TeX
                             % ignores spaces after command names like \ldots
                             % made from \ + letters.
                             %
                             % Note how a `%' character causes TeX to ignore
                             % the end of the input line, so these blank lines
                             % do not start a new paragraph.
                             %
with the right spacing around the periods requires
a special command.

\LaTeX\ interprets some common characters as
commands, so you must type special commands to
generate them.  These characters include the
following:
       \$ \& \% \# \{ and \}.

In printing, text is usually emphasized with an
       \emph{italic}
type style.

\begin{em}
   A long segment of text can also be emphasized
   in this way.  Text within such a segment can be
   given \emph{additional} emphasis.
\end{em}

It is sometimes necessary to prevent \LaTeX\ from
breaking a line where it might otherwise do so.
This may be at a space, as between the ``Mr.''\ and
``Jones'' in
       ``Mr.~Jones'',        % ~ produces an unbreakable interword space.
or within a word---especially when the word is a
symbol like
       \mbox{\emph{itemnum}}
that makes little sense when hyphenated across
lines.

Footnotes\footnote{This is an example of a footnote.}
pose no problem.

\LaTeX\ is good at typesetting mathematical formulas
like
       \( x-3y + z = 7 \)
or
       \( a_{1} > x^{2n} + y^{2n} > x' \)
or
       \( \ip{A}{B} = \sum_{i} a_{i} b_{i} \).
The spaces you type in a formula are
ignored.  Remember that a letter like
       $x$                   % $ ... $  and  \( ... \)  are equivalent
is a formula when it denotes a mathematical
symbol, and it should be typed as one.

\section{Displayed Text}

Text is displayed by indenting it from the left
margin.  Quotations are commonly displayed.  There
are short quotations
\begin{quote}
   This is a short quotation.  It consists of a
   single paragraph of text.  See how it is formatted.
\end{quote}
and longer ones.
\begin{quotation}
   This is a longer quotation.  It consists of two
   paragraphs of text, neither of which are
   particularly interesting.

   This is the second paragraph of the quotation.  It
   is just as dull as the first paragraph.
\end{quotation}
Another frequently-displayed structure is a list.
The following is an example of an \emph{itemized}
list.
\begin{itemize}
   \item This is the first item of an itemized list.
         Each item in the list is marked with a ``tick''.
         You don't have to worry about what kind of tick
         mark is used.

   \item This is the second item of the list.  It
         contains another list nested inside it.  The inner
         list is an \emph{enumerated} list.
         \begin{enumerate}
            \item This is the first item of an enumerated
                  list that is nested within the itemized list.

            \item This is the second item of the inner list.
                  \LaTeX\ allows you to nest lists deeper than
                  you really should.
         \end{enumerate}
         This is the rest of the second item of the outer
         list.  It is no more interesting than any other
         part of the item.
   \item This is the third item of the list.
\end{itemize}
You can even display poetry.
\begin{verse}
   There is an environment
    for verse \\             % The \\ command separates lines
   Whose features some poets % within a stanza.
   will curse.

                             % One or more blank lines separate stanzas.

   For instead of making\\
   Them do \emph{all} line breaking, \\
   It allows them to put too many words on a line when they'd rather be
   forced to be terse.
\end{verse}

Mathematical formulas may also be displayed.  A
displayed formula
is
one-line long; multiline
formulas require special formatting instructions.
   \[  \ip{\Gamma}{\psi'} = x'' + y^{2} + z_{i}^{n}\]
Don't start a paragraph with a displayed equation,
nor make one a paragraph by itself.

\end{document}               % End of document.

\section{Bench Huge Section 1457}
% synthetic bench section 1457
% This is a sample LaTeX input file.  (Version of 12 August 2004.)
%
% A '%' character causes TeX to ignore all remaining text on the line,
% and is used for comments like this one.

\documentclass{article}      % Specifies the document class

                             % The preamble begins here.
\title{An Example Document}  % Declares the document's title.
\author{Leslie Lamport}      % Declares the author's name.
\date{January 21, 1994}      % Deleting this command produces today's date.

\newcommand{\ip}[2]{(#1, #2)}
                             % Defines \ip{arg1}{arg2} to mean
                             % (arg1, arg2).

%\newcommand{\ip}[2]{\langle #1 | #2\rangle}
                             % This is an alternative definition of
                             % \ip that is commented out.

\begin{document}             % End of preamble and beginning of text.

\maketitle                   % Produces the title.

This is an example input file.  Comparing it with
the output it generates can show you how to
produce a simple document of your own.

\section{Ordinary Text}      % Produces section heading.  Lower-level
                             % sections are begun with similar
                             % \subsection and \subsubsection commands.

The ends  of words and sentences are marked
  by   spaces. It  doesn't matter how many
spaces    you type; one is as good as 100.  The
end of   a line counts as a space.

One   or more   blank lines denote the  end
of  a paragraph.

Since any number of consecutive spaces are treated
like a single one, the formatting of the input
file makes no difference to
      \LaTeX,                % The \LaTeX command generates the LaTeX logo.
but it makes a difference to you.  When you use
\LaTeX, making your input file as easy to read
as possible will be a great help as you write
your document and when you change it.  This sample
file shows how you can add comments to your own input
file.

Because printing is different from typewriting,
there are a number of things that you have to do
differently when preparing an input file than if
you were just typing the document directly.
Quotation marks like
       ``this''
have to be handled specially, as do quotes within
quotes:
       ``\,`this'            % \, separates the double and single quote.
        is what I just
        wrote, not  `that'\,''.

Dashes come in three sizes: an
       intra-word
dash, a medium dash for number ranges like
       1--2,
and a punctuation
       dash---like
this.

A sentence-ending space should be larger than the
space between words within a sentence.  You
sometimes have to type special commands in
conjunction with punctuation characters to get
this right, as in the following sentence.
       Gnats, gnus, etc.\ all  % `\ ' makes an inter-word space.
       begin with G\@.         % \@ marks end-of-sentence punctuation.
You should check the spaces after periods when
reading your output to make sure you haven't
forgotten any special cases.  Generating an
ellipsis
       \ldots\               % `\ ' is needed after `\ldots' because TeX
                             % ignores spaces after command names like \ldots
                             % made from \ + letters.
                             %
                             % Note how a `%' character causes TeX to ignore
                             % the end of the input line, so these blank lines
                             % do not start a new paragraph.
                             %
with the right spacing around the periods requires
a special command.

\LaTeX\ interprets some common characters as
commands, so you must type special commands to
generate them.  These characters include the
following:
       \$ \& \% \# \{ and \}.

In printing, text is usually emphasized with an
       \emph{italic}
type style.

\begin{em}
   A long segment of text can also be emphasized
   in this way.  Text within such a segment can be
   given \emph{additional} emphasis.
\end{em}

It is sometimes necessary to prevent \LaTeX\ from
breaking a line where it might otherwise do so.
This may be at a space, as between the ``Mr.''\ and
``Jones'' in
       ``Mr.~Jones'',        % ~ produces an unbreakable interword space.
or within a word---especially when the word is a
symbol like
       \mbox{\emph{itemnum}}
that makes little sense when hyphenated across
lines.

Footnotes\footnote{This is an example of a footnote.}
pose no problem.

\LaTeX\ is good at typesetting mathematical formulas
like
       \( x-3y + z = 7 \)
or
       \( a_{1} > x^{2n} + y^{2n} > x' \)
or
       \( \ip{A}{B} = \sum_{i} a_{i} b_{i} \).
The spaces you type in a formula are
ignored.  Remember that a letter like
       $x$                   % $ ... $  and  \( ... \)  are equivalent
is a formula when it denotes a mathematical
symbol, and it should be typed as one.

\section{Displayed Text}

Text is displayed by indenting it from the left
margin.  Quotations are commonly displayed.  There
are short quotations
\begin{quote}
   This is a short quotation.  It consists of a
   single paragraph of text.  See how it is formatted.
\end{quote}
and longer ones.
\begin{quotation}
   This is a longer quotation.  It consists of two
   paragraphs of text, neither of which are
   particularly interesting.

   This is the second paragraph of the quotation.  It
   is just as dull as the first paragraph.
\end{quotation}
Another frequently-displayed structure is a list.
The following is an example of an \emph{itemized}
list.
\begin{itemize}
   \item This is the first item of an itemized list.
         Each item in the list is marked with a ``tick''.
         You don't have to worry about what kind of tick
         mark is used.

   \item This is the second item of the list.  It
         contains another list nested inside it.  The inner
         list is an \emph{enumerated} list.
         \begin{enumerate}
            \item This is the first item of an enumerated
                  list that is nested within the itemized list.

            \item This is the second item of the inner list.
                  \LaTeX\ allows you to nest lists deeper than
                  you really should.
         \end{enumerate}
         This is the rest of the second item of the outer
         list.  It is no more interesting than any other
         part of the item.
   \item This is the third item of the list.
\end{itemize}
You can even display poetry.
\begin{verse}
   There is an environment
    for verse \\             % The \\ command separates lines
   Whose features some poets % within a stanza.
   will curse.

                             % One or more blank lines separate stanzas.

   For instead of making\\
   Them do \emph{all} line breaking, \\
   It allows them to put too many words on a line when they'd rather be
   forced to be terse.
\end{verse}

Mathematical formulas may also be displayed.  A
displayed formula
is
one-line long; multiline
formulas require special formatting instructions.
   \[  \ip{\Gamma}{\psi'} = x'' + y^{2} + z_{i}^{n}\]
Don't start a paragraph with a displayed equation,
nor make one a paragraph by itself.

\end{document}               % End of document.

\section{Bench Huge Section 1458}
% synthetic bench section 1458
% This is a sample LaTeX input file.  (Version of 12 August 2004.)
%
% A '%' character causes TeX to ignore all remaining text on the line,
% and is used for comments like this one.

\documentclass{article}      % Specifies the document class

                             % The preamble begins here.
\title{An Example Document}  % Declares the document's title.
\author{Leslie Lamport}      % Declares the author's name.
\date{January 21, 1994}      % Deleting this command produces today's date.

\newcommand{\ip}[2]{(#1, #2)}
                             % Defines \ip{arg1}{arg2} to mean
                             % (arg1, arg2).

%\newcommand{\ip}[2]{\langle #1 | #2\rangle}
                             % This is an alternative definition of
                             % \ip that is commented out.

\begin{document}             % End of preamble and beginning of text.

\maketitle                   % Produces the title.

This is an example input file.  Comparing it with
the output it generates can show you how to
produce a simple document of your own.

\section{Ordinary Text}      % Produces section heading.  Lower-level
                             % sections are begun with similar
                             % \subsection and \subsubsection commands.

The ends  of words and sentences are marked
  by   spaces. It  doesn't matter how many
spaces    you type; one is as good as 100.  The
end of   a line counts as a space.

One   or more   blank lines denote the  end
of  a paragraph.

Since any number of consecutive spaces are treated
like a single one, the formatting of the input
file makes no difference to
      \LaTeX,                % The \LaTeX command generates the LaTeX logo.
but it makes a difference to you.  When you use
\LaTeX, making your input file as easy to read
as possible will be a great help as you write
your document and when you change it.  This sample
file shows how you can add comments to your own input
file.

Because printing is different from typewriting,
there are a number of things that you have to do
differently when preparing an input file than if
you were just typing the document directly.
Quotation marks like
       ``this''
have to be handled specially, as do quotes within
quotes:
       ``\,`this'            % \, separates the double and single quote.
        is what I just
        wrote, not  `that'\,''.

Dashes come in three sizes: an
       intra-word
dash, a medium dash for number ranges like
       1--2,
and a punctuation
       dash---like
this.

A sentence-ending space should be larger than the
space between words within a sentence.  You
sometimes have to type special commands in
conjunction with punctuation characters to get
this right, as in the following sentence.
       Gnats, gnus, etc.\ all  % `\ ' makes an inter-word space.
       begin with G\@.         % \@ marks end-of-sentence punctuation.
You should check the spaces after periods when
reading your output to make sure you haven't
forgotten any special cases.  Generating an
ellipsis
       \ldots\               % `\ ' is needed after `\ldots' because TeX
                             % ignores spaces after command names like \ldots
                             % made from \ + letters.
                             %
                             % Note how a `%' character causes TeX to ignore
                             % the end of the input line, so these blank lines
                             % do not start a new paragraph.
                             %
with the right spacing around the periods requires
a special command.

\LaTeX\ interprets some common characters as
commands, so you must type special commands to
generate them.  These characters include the
following:
       \$ \& \% \# \{ and \}.

In printing, text is usually emphasized with an
       \emph{italic}
type style.

\begin{em}
   A long segment of text can also be emphasized
   in this way.  Text within such a segment can be
   given \emph{additional} emphasis.
\end{em}

It is sometimes necessary to prevent \LaTeX\ from
breaking a line where it might otherwise do so.
This may be at a space, as between the ``Mr.''\ and
``Jones'' in
       ``Mr.~Jones'',        % ~ produces an unbreakable interword space.
or within a word---especially when the word is a
symbol like
       \mbox{\emph{itemnum}}
that makes little sense when hyphenated across
lines.

Footnotes\footnote{This is an example of a footnote.}
pose no problem.

\LaTeX\ is good at typesetting mathematical formulas
like
       \( x-3y + z = 7 \)
or
       \( a_{1} > x^{2n} + y^{2n} > x' \)
or
       \( \ip{A}{B} = \sum_{i} a_{i} b_{i} \).
The spaces you type in a formula are
ignored.  Remember that a letter like
       $x$                   % $ ... $  and  \( ... \)  are equivalent
is a formula when it denotes a mathematical
symbol, and it should be typed as one.

\section{Displayed Text}

Text is displayed by indenting it from the left
margin.  Quotations are commonly displayed.  There
are short quotations
\begin{quote}
   This is a short quotation.  It consists of a
   single paragraph of text.  See how it is formatted.
\end{quote}
and longer ones.
\begin{quotation}
   This is a longer quotation.  It consists of two
   paragraphs of text, neither of which are
   particularly interesting.

   This is the second paragraph of the quotation.  It
   is just as dull as the first paragraph.
\end{quotation}
Another frequently-displayed structure is a list.
The following is an example of an \emph{itemized}
list.
\begin{itemize}
   \item This is the first item of an itemized list.
         Each item in the list is marked with a ``tick''.
         You don't have to worry about what kind of tick
         mark is used.

   \item This is the second item of the list.  It
         contains another list nested inside it.  The inner
         list is an \emph{enumerated} list.
         \begin{enumerate}
            \item This is the first item of an enumerated
                  list that is nested within the itemized list.

            \item This is the second item of the inner list.
                  \LaTeX\ allows you to nest lists deeper than
                  you really should.
         \end{enumerate}
         This is the rest of the second item of the outer
         list.  It is no more interesting than any other
         part of the item.
   \item This is the third item of the list.
\end{itemize}
You can even display poetry.
\begin{verse}
   There is an environment
    for verse \\             % The \\ command separates lines
   Whose features some poets % within a stanza.
   will curse.

                             % One or more blank lines separate stanzas.

   For instead of making\\
   Them do \emph{all} line breaking, \\
   It allows them to put too many words on a line when they'd rather be
   forced to be terse.
\end{verse}

Mathematical formulas may also be displayed.  A
displayed formula
is
one-line long; multiline
formulas require special formatting instructions.
   \[  \ip{\Gamma}{\psi'} = x'' + y^{2} + z_{i}^{n}\]
Don't start a paragraph with a displayed equation,
nor make one a paragraph by itself.

\end{document}               % End of document.

\section{Bench Huge Section 1459}
% synthetic bench section 1459
% This is a sample LaTeX input file.  (Version of 12 August 2004.)
%
% A '%' character causes TeX to ignore all remaining text on the line,
% and is used for comments like this one.

\documentclass{article}      % Specifies the document class

                             % The preamble begins here.
\title{An Example Document}  % Declares the document's title.
\author{Leslie Lamport}      % Declares the author's name.
\date{January 21, 1994}      % Deleting this command produces today's date.

\newcommand{\ip}[2]{(#1, #2)}
                             % Defines \ip{arg1}{arg2} to mean
                             % (arg1, arg2).

%\newcommand{\ip}[2]{\langle #1 | #2\rangle}
                             % This is an alternative definition of
                             % \ip that is commented out.

\begin{document}             % End of preamble and beginning of text.

\maketitle                   % Produces the title.

This is an example input file.  Comparing it with
the output it generates can show you how to
produce a simple document of your own.

\section{Ordinary Text}      % Produces section heading.  Lower-level
                             % sections are begun with similar
                             % \subsection and \subsubsection commands.

The ends  of words and sentences are marked
  by   spaces. It  doesn't matter how many
spaces    you type; one is as good as 100.  The
end of   a line counts as a space.

One   or more   blank lines denote the  end
of  a paragraph.

Since any number of consecutive spaces are treated
like a single one, the formatting of the input
file makes no difference to
      \LaTeX,                % The \LaTeX command generates the LaTeX logo.
but it makes a difference to you.  When you use
\LaTeX, making your input file as easy to read
as possible will be a great help as you write
your document and when you change it.  This sample
file shows how you can add comments to your own input
file.

Because printing is different from typewriting,
there are a number of things that you have to do
differently when preparing an input file than if
you were just typing the document directly.
Quotation marks like
       ``this''
have to be handled specially, as do quotes within
quotes:
       ``\,`this'            % \, separates the double and single quote.
        is what I just
        wrote, not  `that'\,''.

Dashes come in three sizes: an
       intra-word
dash, a medium dash for number ranges like
       1--2,
and a punctuation
       dash---like
this.

A sentence-ending space should be larger than the
space between words within a sentence.  You
sometimes have to type special commands in
conjunction with punctuation characters to get
this right, as in the following sentence.
       Gnats, gnus, etc.\ all  % `\ ' makes an inter-word space.
       begin with G\@.         % \@ marks end-of-sentence punctuation.
You should check the spaces after periods when
reading your output to make sure you haven't
forgotten any special cases.  Generating an
ellipsis
       \ldots\               % `\ ' is needed after `\ldots' because TeX
                             % ignores spaces after command names like \ldots
                             % made from \ + letters.
                             %
                             % Note how a `%' character causes TeX to ignore
                             % the end of the input line, so these blank lines
                             % do not start a new paragraph.
                             %
with the right spacing around the periods requires
a special command.

\LaTeX\ interprets some common characters as
commands, so you must type special commands to
generate them.  These characters include the
following:
       \$ \& \% \# \{ and \}.

In printing, text is usually emphasized with an
       \emph{italic}
type style.

\begin{em}
   A long segment of text can also be emphasized
   in this way.  Text within such a segment can be
   given \emph{additional} emphasis.
\end{em}

It is sometimes necessary to prevent \LaTeX\ from
breaking a line where it might otherwise do so.
This may be at a space, as between the ``Mr.''\ and
``Jones'' in
       ``Mr.~Jones'',        % ~ produces an unbreakable interword space.
or within a word---especially when the word is a
symbol like
       \mbox{\emph{itemnum}}
that makes little sense when hyphenated across
lines.

Footnotes\footnote{This is an example of a footnote.}
pose no problem.

\LaTeX\ is good at typesetting mathematical formulas
like
       \( x-3y + z = 7 \)
or
       \( a_{1} > x^{2n} + y^{2n} > x' \)
or
       \( \ip{A}{B} = \sum_{i} a_{i} b_{i} \).
The spaces you type in a formula are
ignored.  Remember that a letter like
       $x$                   % $ ... $  and  \( ... \)  are equivalent
is a formula when it denotes a mathematical
symbol, and it should be typed as one.

\section{Displayed Text}

Text is displayed by indenting it from the left
margin.  Quotations are commonly displayed.  There
are short quotations
\begin{quote}
   This is a short quotation.  It consists of a
   single paragraph of text.  See how it is formatted.
\end{quote}
and longer ones.
\begin{quotation}
   This is a longer quotation.  It consists of two
   paragraphs of text, neither of which are
   particularly interesting.

   This is the second paragraph of the quotation.  It
   is just as dull as the first paragraph.
\end{quotation}
Another frequently-displayed structure is a list.
The following is an example of an \emph{itemized}
list.
\begin{itemize}
   \item This is the first item of an itemized list.
         Each item in the list is marked with a ``tick''.
         You don't have to worry about what kind of tick
         mark is used.

   \item This is the second item of the list.  It
         contains another list nested inside it.  The inner
         list is an \emph{enumerated} list.
         \begin{enumerate}
            \item This is the first item of an enumerated
                  list that is nested within the itemized list.

            \item This is the second item of the inner list.
                  \LaTeX\ allows you to nest lists deeper than
                  you really should.
         \end{enumerate}
         This is the rest of the second item of the outer
         list.  It is no more interesting than any other
         part of the item.
   \item This is the third item of the list.
\end{itemize}
You can even display poetry.
\begin{verse}
   There is an environment
    for verse \\             % The \\ command separates lines
   Whose features some poets % within a stanza.
   will curse.

                             % One or more blank lines separate stanzas.

   For instead of making\\
   Them do \emph{all} line breaking, \\
   It allows them to put too many words on a line when they'd rather be
   forced to be terse.
\end{verse}

Mathematical formulas may also be displayed.  A
displayed formula
is
one-line long; multiline
formulas require special formatting instructions.
   \[  \ip{\Gamma}{\psi'} = x'' + y^{2} + z_{i}^{n}\]
Don't start a paragraph with a displayed equation,
nor make one a paragraph by itself.

\end{document}               % End of document.

\section{Bench Huge Section 1460}
% synthetic bench section 1460
% This is a sample LaTeX input file.  (Version of 12 August 2004.)
%
% A '%' character causes TeX to ignore all remaining text on the line,
% and is used for comments like this one.

\documentclass{article}      % Specifies the document class

                             % The preamble begins here.
\title{An Example Document}  % Declares the document's title.
\author{Leslie Lamport}      % Declares the author's name.
\date{January 21, 1994}      % Deleting this command produces today's date.

\newcommand{\ip}[2]{(#1, #2)}
                             % Defines \ip{arg1}{arg2} to mean
                             % (arg1, arg2).

%\newcommand{\ip}[2]{\langle #1 | #2\rangle}
                             % This is an alternative definition of
                             % \ip that is commented out.

\begin{document}             % End of preamble and beginning of text.

\maketitle                   % Produces the title.

This is an example input file.  Comparing it with
the output it generates can show you how to
produce a simple document of your own.

\section{Ordinary Text}      % Produces section heading.  Lower-level
                             % sections are begun with similar
                             % \subsection and \subsubsection commands.

The ends  of words and sentences are marked
  by   spaces. It  doesn't matter how many
spaces    you type; one is as good as 100.  The
end of   a line counts as a space.

One   or more   blank lines denote the  end
of  a paragraph.

Since any number of consecutive spaces are treated
like a single one, the formatting of the input
file makes no difference to
      \LaTeX,                % The \LaTeX command generates the LaTeX logo.
but it makes a difference to you.  When you use
\LaTeX, making your input file as easy to read
as possible will be a great help as you write
your document and when you change it.  This sample
file shows how you can add comments to your own input
file.

Because printing is different from typewriting,
there are a number of things that you have to do
differently when preparing an input file than if
you were just typing the document directly.
Quotation marks like
       ``this''
have to be handled specially, as do quotes within
quotes:
       ``\,`this'            % \, separates the double and single quote.
        is what I just
        wrote, not  `that'\,''.

Dashes come in three sizes: an
       intra-word
dash, a medium dash for number ranges like
       1--2,
and a punctuation
       dash---like
this.

A sentence-ending space should be larger than the
space between words within a sentence.  You
sometimes have to type special commands in
conjunction with punctuation characters to get
this right, as in the following sentence.
       Gnats, gnus, etc.\ all  % `\ ' makes an inter-word space.
       begin with G\@.         % \@ marks end-of-sentence punctuation.
You should check the spaces after periods when
reading your output to make sure you haven't
forgotten any special cases.  Generating an
ellipsis
       \ldots\               % `\ ' is needed after `\ldots' because TeX
                             % ignores spaces after command names like \ldots
                             % made from \ + letters.
                             %
                             % Note how a `%' character causes TeX to ignore
                             % the end of the input line, so these blank lines
                             % do not start a new paragraph.
                             %
with the right spacing around the periods requires
a special command.

\LaTeX\ interprets some common characters as
commands, so you must type special commands to
generate them.  These characters include the
following:
       \$ \& \% \# \{ and \}.

In printing, text is usually emphasized with an
       \emph{italic}
type style.

\begin{em}
   A long segment of text can also be emphasized
   in this way.  Text within such a segment can be
   given \emph{additional} emphasis.
\end{em}

It is sometimes necessary to prevent \LaTeX\ from
breaking a line where it might otherwise do so.
This may be at a space, as between the ``Mr.''\ and
``Jones'' in
       ``Mr.~Jones'',        % ~ produces an unbreakable interword space.
or within a word---especially when the word is a
symbol like
       \mbox{\emph{itemnum}}
that makes little sense when hyphenated across
lines.

Footnotes\footnote{This is an example of a footnote.}
pose no problem.

\LaTeX\ is good at typesetting mathematical formulas
like
       \( x-3y + z = 7 \)
or
       \( a_{1} > x^{2n} + y^{2n} > x' \)
or
       \( \ip{A}{B} = \sum_{i} a_{i} b_{i} \).
The spaces you type in a formula are
ignored.  Remember that a letter like
       $x$                   % $ ... $  and  \( ... \)  are equivalent
is a formula when it denotes a mathematical
symbol, and it should be typed as one.

\section{Displayed Text}

Text is displayed by indenting it from the left
margin.  Quotations are commonly displayed.  There
are short quotations
\begin{quote}
   This is a short quotation.  It consists of a
   single paragraph of text.  See how it is formatted.
\end{quote}
and longer ones.
\begin{quotation}
   This is a longer quotation.  It consists of two
   paragraphs of text, neither of which are
   particularly interesting.

   This is the second paragraph of the quotation.  It
   is just as dull as the first paragraph.
\end{quotation}
Another frequently-displayed structure is a list.
The following is an example of an \emph{itemized}
list.
\begin{itemize}
   \item This is the first item of an itemized list.
         Each item in the list is marked with a ``tick''.
         You don't have to worry about what kind of tick
         mark is used.

   \item This is the second item of the list.  It
         contains another list nested inside it.  The inner
         list is an \emph{enumerated} list.
         \begin{enumerate}
            \item This is the first item of an enumerated
                  list that is nested within the itemized list.

            \item This is the second item of the inner list.
                  \LaTeX\ allows you to nest lists deeper than
                  you really should.
         \end{enumerate}
         This is the rest of the second item of the outer
         list.  It is no more interesting than any other
         part of the item.
   \item This is the third item of the list.
\end{itemize}
You can even display poetry.
\begin{verse}
   There is an environment
    for verse \\             % The \\ command separates lines
   Whose features some poets % within a stanza.
   will curse.

                             % One or more blank lines separate stanzas.

   For instead of making\\
   Them do \emph{all} line breaking, \\
   It allows them to put too many words on a line when they'd rather be
   forced to be terse.
\end{verse}

Mathematical formulas may also be displayed.  A
displayed formula
is
one-line long; multiline
formulas require special formatting instructions.
   \[  \ip{\Gamma}{\psi'} = x'' + y^{2} + z_{i}^{n}\]
Don't start a paragraph with a displayed equation,
nor make one a paragraph by itself.

\end{document}               % End of document.

\section{Bench Huge Section 1461}
% synthetic bench section 1461
% This is a sample LaTeX input file.  (Version of 12 August 2004.)
%
% A '%' character causes TeX to ignore all remaining text on the line,
% and is used for comments like this one.

\documentclass{article}      % Specifies the document class

                             % The preamble begins here.
\title{An Example Document}  % Declares the document's title.
\author{Leslie Lamport}      % Declares the author's name.
\date{January 21, 1994}      % Deleting this command produces today's date.

\newcommand{\ip}[2]{(#1, #2)}
                             % Defines \ip{arg1}{arg2} to mean
                             % (arg1, arg2).

%\newcommand{\ip}[2]{\langle #1 | #2\rangle}
                             % This is an alternative definition of
                             % \ip that is commented out.

\begin{document}             % End of preamble and beginning of text.

\maketitle                   % Produces the title.

This is an example input file.  Comparing it with
the output it generates can show you how to
produce a simple document of your own.

\section{Ordinary Text}      % Produces section heading.  Lower-level
                             % sections are begun with similar
                             % \subsection and \subsubsection commands.

The ends  of words and sentences are marked
  by   spaces. It  doesn't matter how many
spaces    you type; one is as good as 100.  The
end of   a line counts as a space.

One   or more   blank lines denote the  end
of  a paragraph.

Since any number of consecutive spaces are treated
like a single one, the formatting of the input
file makes no difference to
      \LaTeX,                % The \LaTeX command generates the LaTeX logo.
but it makes a difference to you.  When you use
\LaTeX, making your input file as easy to read
as possible will be a great help as you write
your document and when you change it.  This sample
file shows how you can add comments to your own input
file.

Because printing is different from typewriting,
there are a number of things that you have to do
differently when preparing an input file than if
you were just typing the document directly.
Quotation marks like
       ``this''
have to be handled specially, as do quotes within
quotes:
       ``\,`this'            % \, separates the double and single quote.
        is what I just
        wrote, not  `that'\,''.

Dashes come in three sizes: an
       intra-word
dash, a medium dash for number ranges like
       1--2,
and a punctuation
       dash---like
this.

A sentence-ending space should be larger than the
space between words within a sentence.  You
sometimes have to type special commands in
conjunction with punctuation characters to get
this right, as in the following sentence.
       Gnats, gnus, etc.\ all  % `\ ' makes an inter-word space.
       begin with G\@.         % \@ marks end-of-sentence punctuation.
You should check the spaces after periods when
reading your output to make sure you haven't
forgotten any special cases.  Generating an
ellipsis
       \ldots\               % `\ ' is needed after `\ldots' because TeX
                             % ignores spaces after command names like \ldots
                             % made from \ + letters.
                             %
                             % Note how a `%' character causes TeX to ignore
                             % the end of the input line, so these blank lines
                             % do not start a new paragraph.
                             %
with the right spacing around the periods requires
a special command.

\LaTeX\ interprets some common characters as
commands, so you must type special commands to
generate them.  These characters include the
following:
       \$ \& \% \# \{ and \}.

In printing, text is usually emphasized with an
       \emph{italic}
type style.

\begin{em}
   A long segment of text can also be emphasized
   in this way.  Text within such a segment can be
   given \emph{additional} emphasis.
\end{em}

It is sometimes necessary to prevent \LaTeX\ from
breaking a line where it might otherwise do so.
This may be at a space, as between the ``Mr.''\ and
``Jones'' in
       ``Mr.~Jones'',        % ~ produces an unbreakable interword space.
or within a word---especially when the word is a
symbol like
       \mbox{\emph{itemnum}}
that makes little sense when hyphenated across
lines.

Footnotes\footnote{This is an example of a footnote.}
pose no problem.

\LaTeX\ is good at typesetting mathematical formulas
like
       \( x-3y + z = 7 \)
or
       \( a_{1} > x^{2n} + y^{2n} > x' \)
or
       \( \ip{A}{B} = \sum_{i} a_{i} b_{i} \).
The spaces you type in a formula are
ignored.  Remember that a letter like
       $x$                   % $ ... $  and  \( ... \)  are equivalent
is a formula when it denotes a mathematical
symbol, and it should be typed as one.

\section{Displayed Text}

Text is displayed by indenting it from the left
margin.  Quotations are commonly displayed.  There
are short quotations
\begin{quote}
   This is a short quotation.  It consists of a
   single paragraph of text.  See how it is formatted.
\end{quote}
and longer ones.
\begin{quotation}
   This is a longer quotation.  It consists of two
   paragraphs of text, neither of which are
   particularly interesting.

   This is the second paragraph of the quotation.  It
   is just as dull as the first paragraph.
\end{quotation}
Another frequently-displayed structure is a list.
The following is an example of an \emph{itemized}
list.
\begin{itemize}
   \item This is the first item of an itemized list.
         Each item in the list is marked with a ``tick''.
         You don't have to worry about what kind of tick
         mark is used.

   \item This is the second item of the list.  It
         contains another list nested inside it.  The inner
         list is an \emph{enumerated} list.
         \begin{enumerate}
            \item This is the first item of an enumerated
                  list that is nested within the itemized list.

            \item This is the second item of the inner list.
                  \LaTeX\ allows you to nest lists deeper than
                  you really should.
         \end{enumerate}
         This is the rest of the second item of the outer
         list.  It is no more interesting than any other
         part of the item.
   \item This is the third item of the list.
\end{itemize}
You can even display poetry.
\begin{verse}
   There is an environment
    for verse \\             % The \\ command separates lines
   Whose features some poets % within a stanza.
   will curse.

                             % One or more blank lines separate stanzas.

   For instead of making\\
   Them do \emph{all} line breaking, \\
   It allows them to put too many words on a line when they'd rather be
   forced to be terse.
\end{verse}

Mathematical formulas may also be displayed.  A
displayed formula
is
one-line long; multiline
formulas require special formatting instructions.
   \[  \ip{\Gamma}{\psi'} = x'' + y^{2} + z_{i}^{n}\]
Don't start a paragraph with a displayed equation,
nor make one a paragraph by itself.

\end{document}               % End of document.

\section{Bench Huge Section 1462}
% synthetic bench section 1462
% This is a sample LaTeX input file.  (Version of 12 August 2004.)
%
% A '%' character causes TeX to ignore all remaining text on the line,
% and is used for comments like this one.

\documentclass{article}      % Specifies the document class

                             % The preamble begins here.
\title{An Example Document}  % Declares the document's title.
\author{Leslie Lamport}      % Declares the author's name.
\date{January 21, 1994}      % Deleting this command produces today's date.

\newcommand{\ip}[2]{(#1, #2)}
                             % Defines \ip{arg1}{arg2} to mean
                             % (arg1, arg2).

%\newcommand{\ip}[2]{\langle #1 | #2\rangle}
                             % This is an alternative definition of
                             % \ip that is commented out.

\begin{document}             % End of preamble and beginning of text.

\maketitle                   % Produces the title.

This is an example input file.  Comparing it with
the output it generates can show you how to
produce a simple document of your own.

\section{Ordinary Text}      % Produces section heading.  Lower-level
                             % sections are begun with similar
                             % \subsection and \subsubsection commands.

The ends  of words and sentences are marked
  by   spaces. It  doesn't matter how many
spaces    you type; one is as good as 100.  The
end of   a line counts as a space.

One   or more   blank lines denote the  end
of  a paragraph.

Since any number of consecutive spaces are treated
like a single one, the formatting of the input
file makes no difference to
      \LaTeX,                % The \LaTeX command generates the LaTeX logo.
but it makes a difference to you.  When you use
\LaTeX, making your input file as easy to read
as possible will be a great help as you write
your document and when you change it.  This sample
file shows how you can add comments to your own input
file.

Because printing is different from typewriting,
there are a number of things that you have to do
differently when preparing an input file than if
you were just typing the document directly.
Quotation marks like
       ``this''
have to be handled specially, as do quotes within
quotes:
       ``\,`this'            % \, separates the double and single quote.
        is what I just
        wrote, not  `that'\,''.

Dashes come in three sizes: an
       intra-word
dash, a medium dash for number ranges like
       1--2,
and a punctuation
       dash---like
this.

A sentence-ending space should be larger than the
space between words within a sentence.  You
sometimes have to type special commands in
conjunction with punctuation characters to get
this right, as in the following sentence.
       Gnats, gnus, etc.\ all  % `\ ' makes an inter-word space.
       begin with G\@.         % \@ marks end-of-sentence punctuation.
You should check the spaces after periods when
reading your output to make sure you haven't
forgotten any special cases.  Generating an
ellipsis
       \ldots\               % `\ ' is needed after `\ldots' because TeX
                             % ignores spaces after command names like \ldots
                             % made from \ + letters.
                             %
                             % Note how a `%' character causes TeX to ignore
                             % the end of the input line, so these blank lines
                             % do not start a new paragraph.
                             %
with the right spacing around the periods requires
a special command.

\LaTeX\ interprets some common characters as
commands, so you must type special commands to
generate them.  These characters include the
following:
       \$ \& \% \# \{ and \}.

In printing, text is usually emphasized with an
       \emph{italic}
type style.

\begin{em}
   A long segment of text can also be emphasized
   in this way.  Text within such a segment can be
   given \emph{additional} emphasis.
\end{em}

It is sometimes necessary to prevent \LaTeX\ from
breaking a line where it might otherwise do so.
This may be at a space, as between the ``Mr.''\ and
``Jones'' in
       ``Mr.~Jones'',        % ~ produces an unbreakable interword space.
or within a word---especially when the word is a
symbol like
       \mbox{\emph{itemnum}}
that makes little sense when hyphenated across
lines.

Footnotes\footnote{This is an example of a footnote.}
pose no problem.

\LaTeX\ is good at typesetting mathematical formulas
like
       \( x-3y + z = 7 \)
or
       \( a_{1} > x^{2n} + y^{2n} > x' \)
or
       \( \ip{A}{B} = \sum_{i} a_{i} b_{i} \).
The spaces you type in a formula are
ignored.  Remember that a letter like
       $x$                   % $ ... $  and  \( ... \)  are equivalent
is a formula when it denotes a mathematical
symbol, and it should be typed as one.

\section{Displayed Text}

Text is displayed by indenting it from the left
margin.  Quotations are commonly displayed.  There
are short quotations
\begin{quote}
   This is a short quotation.  It consists of a
   single paragraph of text.  See how it is formatted.
\end{quote}
and longer ones.
\begin{quotation}
   This is a longer quotation.  It consists of two
   paragraphs of text, neither of which are
   particularly interesting.

   This is the second paragraph of the quotation.  It
   is just as dull as the first paragraph.
\end{quotation}
Another frequently-displayed structure is a list.
The following is an example of an \emph{itemized}
list.
\begin{itemize}
   \item This is the first item of an itemized list.
         Each item in the list is marked with a ``tick''.
         You don't have to worry about what kind of tick
         mark is used.

   \item This is the second item of the list.  It
         contains another list nested inside it.  The inner
         list is an \emph{enumerated} list.
         \begin{enumerate}
            \item This is the first item of an enumerated
                  list that is nested within the itemized list.

            \item This is the second item of the inner list.
                  \LaTeX\ allows you to nest lists deeper than
                  you really should.
         \end{enumerate}
         This is the rest of the second item of the outer
         list.  It is no more interesting than any other
         part of the item.
   \item This is the third item of the list.
\end{itemize}
You can even display poetry.
\begin{verse}
   There is an environment
    for verse \\             % The \\ command separates lines
   Whose features some poets % within a stanza.
   will curse.

                             % One or more blank lines separate stanzas.

   For instead of making\\
   Them do \emph{all} line breaking, \\
   It allows them to put too many words on a line when they'd rather be
   forced to be terse.
\end{verse}

Mathematical formulas may also be displayed.  A
displayed formula
is
one-line long; multiline
formulas require special formatting instructions.
   \[  \ip{\Gamma}{\psi'} = x'' + y^{2} + z_{i}^{n}\]
Don't start a paragraph with a displayed equation,
nor make one a paragraph by itself.

\end{document}               % End of document.

\section{Bench Huge Section 1463}
% synthetic bench section 1463
% This is a sample LaTeX input file.  (Version of 12 August 2004.)
%
% A '%' character causes TeX to ignore all remaining text on the line,
% and is used for comments like this one.

\documentclass{article}      % Specifies the document class

                             % The preamble begins here.
\title{An Example Document}  % Declares the document's title.
\author{Leslie Lamport}      % Declares the author's name.
\date{January 21, 1994}      % Deleting this command produces today's date.

\newcommand{\ip}[2]{(#1, #2)}
                             % Defines \ip{arg1}{arg2} to mean
                             % (arg1, arg2).

%\newcommand{\ip}[2]{\langle #1 | #2\rangle}
                             % This is an alternative definition of
                             % \ip that is commented out.

\begin{document}             % End of preamble and beginning of text.

\maketitle                   % Produces the title.

This is an example input file.  Comparing it with
the output it generates can show you how to
produce a simple document of your own.

\section{Ordinary Text}      % Produces section heading.  Lower-level
                             % sections are begun with similar
                             % \subsection and \subsubsection commands.

The ends  of words and sentences are marked
  by   spaces. It  doesn't matter how many
spaces    you type; one is as good as 100.  The
end of   a line counts as a space.

One   or more   blank lines denote the  end
of  a paragraph.

Since any number of consecutive spaces are treated
like a single one, the formatting of the input
file makes no difference to
      \LaTeX,                % The \LaTeX command generates the LaTeX logo.
but it makes a difference to you.  When you use
\LaTeX, making your input file as easy to read
as possible will be a great help as you write
your document and when you change it.  This sample
file shows how you can add comments to your own input
file.

Because printing is different from typewriting,
there are a number of things that you have to do
differently when preparing an input file than if
you were just typing the document directly.
Quotation marks like
       ``this''
have to be handled specially, as do quotes within
quotes:
       ``\,`this'            % \, separates the double and single quote.
        is what I just
        wrote, not  `that'\,''.

Dashes come in three sizes: an
       intra-word
dash, a medium dash for number ranges like
       1--2,
and a punctuation
       dash---like
this.

A sentence-ending space should be larger than the
space between words within a sentence.  You
sometimes have to type special commands in
conjunction with punctuation characters to get
this right, as in the following sentence.
       Gnats, gnus, etc.\ all  % `\ ' makes an inter-word space.
       begin with G\@.         % \@ marks end-of-sentence punctuation.
You should check the spaces after periods when
reading your output to make sure you haven't
forgotten any special cases.  Generating an
ellipsis
       \ldots\               % `\ ' is needed after `\ldots' because TeX
                             % ignores spaces after command names like \ldots
                             % made from \ + letters.
                             %
                             % Note how a `%' character causes TeX to ignore
                             % the end of the input line, so these blank lines
                             % do not start a new paragraph.
                             %
with the right spacing around the periods requires
a special command.

\LaTeX\ interprets some common characters as
commands, so you must type special commands to
generate them.  These characters include the
following:
       \$ \& \% \# \{ and \}.

In printing, text is usually emphasized with an
       \emph{italic}
type style.

\begin{em}
   A long segment of text can also be emphasized
   in this way.  Text within such a segment can be
   given \emph{additional} emphasis.
\end{em}

It is sometimes necessary to prevent \LaTeX\ from
breaking a line where it might otherwise do so.
This may be at a space, as between the ``Mr.''\ and
``Jones'' in
       ``Mr.~Jones'',        % ~ produces an unbreakable interword space.
or within a word---especially when the word is a
symbol like
       \mbox{\emph{itemnum}}
that makes little sense when hyphenated across
lines.

Footnotes\footnote{This is an example of a footnote.}
pose no problem.

\LaTeX\ is good at typesetting mathematical formulas
like
       \( x-3y + z = 7 \)
or
       \( a_{1} > x^{2n} + y^{2n} > x' \)
or
       \( \ip{A}{B} = \sum_{i} a_{i} b_{i} \).
The spaces you type in a formula are
ignored.  Remember that a letter like
       $x$                   % $ ... $  and  \( ... \)  are equivalent
is a formula when it denotes a mathematical
symbol, and it should be typed as one.

\section{Displayed Text}

Text is displayed by indenting it from the left
margin.  Quotations are commonly displayed.  There
are short quotations
\begin{quote}
   This is a short quotation.  It consists of a
   single paragraph of text.  See how it is formatted.
\end{quote}
and longer ones.
\begin{quotation}
   This is a longer quotation.  It consists of two
   paragraphs of text, neither of which are
   particularly interesting.

   This is the second paragraph of the quotation.  It
   is just as dull as the first paragraph.
\end{quotation}
Another frequently-displayed structure is a list.
The following is an example of an \emph{itemized}
list.
\begin{itemize}
   \item This is the first item of an itemized list.
         Each item in the list is marked with a ``tick''.
         You don't have to worry about what kind of tick
         mark is used.

   \item This is the second item of the list.  It
         contains another list nested inside it.  The inner
         list is an \emph{enumerated} list.
         \begin{enumerate}
            \item This is the first item of an enumerated
                  list that is nested within the itemized list.

            \item This is the second item of the inner list.
                  \LaTeX\ allows you to nest lists deeper than
                  you really should.
         \end{enumerate}
         This is the rest of the second item of the outer
         list.  It is no more interesting than any other
         part of the item.
   \item This is the third item of the list.
\end{itemize}
You can even display poetry.
\begin{verse}
   There is an environment
    for verse \\             % The \\ command separates lines
   Whose features some poets % within a stanza.
   will curse.

                             % One or more blank lines separate stanzas.

   For instead of making\\
   Them do \emph{all} line breaking, \\
   It allows them to put too many words on a line when they'd rather be
   forced to be terse.
\end{verse}

Mathematical formulas may also be displayed.  A
displayed formula
is
one-line long; multiline
formulas require special formatting instructions.
   \[  \ip{\Gamma}{\psi'} = x'' + y^{2} + z_{i}^{n}\]
Don't start a paragraph with a displayed equation,
nor make one a paragraph by itself.

\end{document}               % End of document.

\section{Bench Huge Section 1464}
% synthetic bench section 1464
% This is a sample LaTeX input file.  (Version of 12 August 2004.)
%
% A '%' character causes TeX to ignore all remaining text on the line,
% and is used for comments like this one.

\documentclass{article}      % Specifies the document class

                             % The preamble begins here.
\title{An Example Document}  % Declares the document's title.
\author{Leslie Lamport}      % Declares the author's name.
\date{January 21, 1994}      % Deleting this command produces today's date.

\newcommand{\ip}[2]{(#1, #2)}
                             % Defines \ip{arg1}{arg2} to mean
                             % (arg1, arg2).

%\newcommand{\ip}[2]{\langle #1 | #2\rangle}
                             % This is an alternative definition of
                             % \ip that is commented out.

\begin{document}             % End of preamble and beginning of text.

\maketitle                   % Produces the title.

This is an example input file.  Comparing it with
the output it generates can show you how to
produce a simple document of your own.

\section{Ordinary Text}      % Produces section heading.  Lower-level
                             % sections are begun with similar
                             % \subsection and \subsubsection commands.

The ends  of words and sentences are marked
  by   spaces. It  doesn't matter how many
spaces    you type; one is as good as 100.  The
end of   a line counts as a space.

One   or more   blank lines denote the  end
of  a paragraph.

Since any number of consecutive spaces are treated
like a single one, the formatting of the input
file makes no difference to
      \LaTeX,                % The \LaTeX command generates the LaTeX logo.
but it makes a difference to you.  When you use
\LaTeX, making your input file as easy to read
as possible will be a great help as you write
your document and when you change it.  This sample
file shows how you can add comments to your own input
file.

Because printing is different from typewriting,
there are a number of things that you have to do
differently when preparing an input file than if
you were just typing the document directly.
Quotation marks like
       ``this''
have to be handled specially, as do quotes within
quotes:
       ``\,`this'            % \, separates the double and single quote.
        is what I just
        wrote, not  `that'\,''.

Dashes come in three sizes: an
       intra-word
dash, a medium dash for number ranges like
       1--2,
and a punctuation
       dash---like
this.

A sentence-ending space should be larger than the
space between words within a sentence.  You
sometimes have to type special commands in
conjunction with punctuation characters to get
this right, as in the following sentence.
       Gnats, gnus, etc.\ all  % `\ ' makes an inter-word space.
       begin with G\@.         % \@ marks end-of-sentence punctuation.
You should check the spaces after periods when
reading your output to make sure you haven't
forgotten any special cases.  Generating an
ellipsis
       \ldots\               % `\ ' is needed after `\ldots' because TeX
                             % ignores spaces after command names like \ldots
                             % made from \ + letters.
                             %
                             % Note how a `%' character causes TeX to ignore
                             % the end of the input line, so these blank lines
                             % do not start a new paragraph.
                             %
with the right spacing around the periods requires
a special command.

\LaTeX\ interprets some common characters as
commands, so you must type special commands to
generate them.  These characters include the
following:
       \$ \& \% \# \{ and \}.

In printing, text is usually emphasized with an
       \emph{italic}
type style.

\begin{em}
   A long segment of text can also be emphasized
   in this way.  Text within such a segment can be
   given \emph{additional} emphasis.
\end{em}

It is sometimes necessary to prevent \LaTeX\ from
breaking a line where it might otherwise do so.
This may be at a space, as between the ``Mr.''\ and
``Jones'' in
       ``Mr.~Jones'',        % ~ produces an unbreakable interword space.
or within a word---especially when the word is a
symbol like
       \mbox{\emph{itemnum}}
that makes little sense when hyphenated across
lines.

Footnotes\footnote{This is an example of a footnote.}
pose no problem.

\LaTeX\ is good at typesetting mathematical formulas
like
       \( x-3y + z = 7 \)
or
       \( a_{1} > x^{2n} + y^{2n} > x' \)
or
       \( \ip{A}{B} = \sum_{i} a_{i} b_{i} \).
The spaces you type in a formula are
ignored.  Remember that a letter like
       $x$                   % $ ... $  and  \( ... \)  are equivalent
is a formula when it denotes a mathematical
symbol, and it should be typed as one.

\section{Displayed Text}

Text is displayed by indenting it from the left
margin.  Quotations are commonly displayed.  There
are short quotations
\begin{quote}
   This is a short quotation.  It consists of a
   single paragraph of text.  See how it is formatted.
\end{quote}
and longer ones.
\begin{quotation}
   This is a longer quotation.  It consists of two
   paragraphs of text, neither of which are
   particularly interesting.

   This is the second paragraph of the quotation.  It
   is just as dull as the first paragraph.
\end{quotation}
Another frequently-displayed structure is a list.
The following is an example of an \emph{itemized}
list.
\begin{itemize}
   \item This is the first item of an itemized list.
         Each item in the list is marked with a ``tick''.
         You don't have to worry about what kind of tick
         mark is used.

   \item This is the second item of the list.  It
         contains another list nested inside it.  The inner
         list is an \emph{enumerated} list.
         \begin{enumerate}
            \item This is the first item of an enumerated
                  list that is nested within the itemized list.

            \item This is the second item of the inner list.
                  \LaTeX\ allows you to nest lists deeper than
                  you really should.
         \end{enumerate}
         This is the rest of the second item of the outer
         list.  It is no more interesting than any other
         part of the item.
   \item This is the third item of the list.
\end{itemize}
You can even display poetry.
\begin{verse}
   There is an environment
    for verse \\             % The \\ command separates lines
   Whose features some poets % within a stanza.
   will curse.

                             % One or more blank lines separate stanzas.

   For instead of making\\
   Them do \emph{all} line breaking, \\
   It allows them to put too many words on a line when they'd rather be
   forced to be terse.
\end{verse}

Mathematical formulas may also be displayed.  A
displayed formula
is
one-line long; multiline
formulas require special formatting instructions.
   \[  \ip{\Gamma}{\psi'} = x'' + y^{2} + z_{i}^{n}\]
Don't start a paragraph with a displayed equation,
nor make one a paragraph by itself.

\end{document}               % End of document.

\section{Bench Huge Section 1465}
% synthetic bench section 1465
% This is a sample LaTeX input file.  (Version of 12 August 2004.)
%
% A '%' character causes TeX to ignore all remaining text on the line,
% and is used for comments like this one.

\documentclass{article}      % Specifies the document class

                             % The preamble begins here.
\title{An Example Document}  % Declares the document's title.
\author{Leslie Lamport}      % Declares the author's name.
\date{January 21, 1994}      % Deleting this command produces today's date.

\newcommand{\ip}[2]{(#1, #2)}
                             % Defines \ip{arg1}{arg2} to mean
                             % (arg1, arg2).

%\newcommand{\ip}[2]{\langle #1 | #2\rangle}
                             % This is an alternative definition of
                             % \ip that is commented out.

\begin{document}             % End of preamble and beginning of text.

\maketitle                   % Produces the title.

This is an example input file.  Comparing it with
the output it generates can show you how to
produce a simple document of your own.

\section{Ordinary Text}      % Produces section heading.  Lower-level
                             % sections are begun with similar
                             % \subsection and \subsubsection commands.

The ends  of words and sentences are marked
  by   spaces. It  doesn't matter how many
spaces    you type; one is as good as 100.  The
end of   a line counts as a space.

One   or more   blank lines denote the  end
of  a paragraph.

Since any number of consecutive spaces are treated
like a single one, the formatting of the input
file makes no difference to
      \LaTeX,                % The \LaTeX command generates the LaTeX logo.
but it makes a difference to you.  When you use
\LaTeX, making your input file as easy to read
as possible will be a great help as you write
your document and when you change it.  This sample
file shows how you can add comments to your own input
file.

Because printing is different from typewriting,
there are a number of things that you have to do
differently when preparing an input file than if
you were just typing the document directly.
Quotation marks like
       ``this''
have to be handled specially, as do quotes within
quotes:
       ``\,`this'            % \, separates the double and single quote.
        is what I just
        wrote, not  `that'\,''.

Dashes come in three sizes: an
       intra-word
dash, a medium dash for number ranges like
       1--2,
and a punctuation
       dash---like
this.

A sentence-ending space should be larger than the
space between words within a sentence.  You
sometimes have to type special commands in
conjunction with punctuation characters to get
this right, as in the following sentence.
       Gnats, gnus, etc.\ all  % `\ ' makes an inter-word space.
       begin with G\@.         % \@ marks end-of-sentence punctuation.
You should check the spaces after periods when
reading your output to make sure you haven't
forgotten any special cases.  Generating an
ellipsis
       \ldots\               % `\ ' is needed after `\ldots' because TeX
                             % ignores spaces after command names like \ldots
                             % made from \ + letters.
                             %
                             % Note how a `%' character causes TeX to ignore
                             % the end of the input line, so these blank lines
                             % do not start a new paragraph.
                             %
with the right spacing around the periods requires
a special command.

\LaTeX\ interprets some common characters as
commands, so you must type special commands to
generate them.  These characters include the
following:
       \$ \& \% \# \{ and \}.

In printing, text is usually emphasized with an
       \emph{italic}
type style.

\begin{em}
   A long segment of text can also be emphasized
   in this way.  Text within such a segment can be
   given \emph{additional} emphasis.
\end{em}

It is sometimes necessary to prevent \LaTeX\ from
breaking a line where it might otherwise do so.
This may be at a space, as between the ``Mr.''\ and
``Jones'' in
       ``Mr.~Jones'',        % ~ produces an unbreakable interword space.
or within a word---especially when the word is a
symbol like
       \mbox{\emph{itemnum}}
that makes little sense when hyphenated across
lines.

Footnotes\footnote{This is an example of a footnote.}
pose no problem.

\LaTeX\ is good at typesetting mathematical formulas
like
       \( x-3y + z = 7 \)
or
       \( a_{1} > x^{2n} + y^{2n} > x' \)
or
       \( \ip{A}{B} = \sum_{i} a_{i} b_{i} \).
The spaces you type in a formula are
ignored.  Remember that a letter like
       $x$                   % $ ... $  and  \( ... \)  are equivalent
is a formula when it denotes a mathematical
symbol, and it should be typed as one.

\section{Displayed Text}

Text is displayed by indenting it from the left
margin.  Quotations are commonly displayed.  There
are short quotations
\begin{quote}
   This is a short quotation.  It consists of a
   single paragraph of text.  See how it is formatted.
\end{quote}
and longer ones.
\begin{quotation}
   This is a longer quotation.  It consists of two
   paragraphs of text, neither of which are
   particularly interesting.

   This is the second paragraph of the quotation.  It
   is just as dull as the first paragraph.
\end{quotation}
Another frequently-displayed structure is a list.
The following is an example of an \emph{itemized}
list.
\begin{itemize}
   \item This is the first item of an itemized list.
         Each item in the list is marked with a ``tick''.
         You don't have to worry about what kind of tick
         mark is used.

   \item This is the second item of the list.  It
         contains another list nested inside it.  The inner
         list is an \emph{enumerated} list.
         \begin{enumerate}
            \item This is the first item of an enumerated
                  list that is nested within the itemized list.

            \item This is the second item of the inner list.
                  \LaTeX\ allows you to nest lists deeper than
                  you really should.
         \end{enumerate}
         This is the rest of the second item of the outer
         list.  It is no more interesting than any other
         part of the item.
   \item This is the third item of the list.
\end{itemize}
You can even display poetry.
\begin{verse}
   There is an environment
    for verse \\             % The \\ command separates lines
   Whose features some poets % within a stanza.
   will curse.

                             % One or more blank lines separate stanzas.

   For instead of making\\
   Them do \emph{all} line breaking, \\
   It allows them to put too many words on a line when they'd rather be
   forced to be terse.
\end{verse}

Mathematical formulas may also be displayed.  A
displayed formula
is
one-line long; multiline
formulas require special formatting instructions.
   \[  \ip{\Gamma}{\psi'} = x'' + y^{2} + z_{i}^{n}\]
Don't start a paragraph with a displayed equation,
nor make one a paragraph by itself.

\end{document}               % End of document.

\section{Bench Huge Section 1466}
% synthetic bench section 1466
% This is a sample LaTeX input file.  (Version of 12 August 2004.)
%
% A '%' character causes TeX to ignore all remaining text on the line,
% and is used for comments like this one.

\documentclass{article}      % Specifies the document class

                             % The preamble begins here.
\title{An Example Document}  % Declares the document's title.
\author{Leslie Lamport}      % Declares the author's name.
\date{January 21, 1994}      % Deleting this command produces today's date.

\newcommand{\ip}[2]{(#1, #2)}
                             % Defines \ip{arg1}{arg2} to mean
                             % (arg1, arg2).

%\newcommand{\ip}[2]{\langle #1 | #2\rangle}
                             % This is an alternative definition of
                             % \ip that is commented out.

\begin{document}             % End of preamble and beginning of text.

\maketitle                   % Produces the title.

This is an example input file.  Comparing it with
the output it generates can show you how to
produce a simple document of your own.

\section{Ordinary Text}      % Produces section heading.  Lower-level
                             % sections are begun with similar
                             % \subsection and \subsubsection commands.

The ends  of words and sentences are marked
  by   spaces. It  doesn't matter how many
spaces    you type; one is as good as 100.  The
end of   a line counts as a space.

One   or more   blank lines denote the  end
of  a paragraph.

Since any number of consecutive spaces are treated
like a single one, the formatting of the input
file makes no difference to
      \LaTeX,                % The \LaTeX command generates the LaTeX logo.
but it makes a difference to you.  When you use
\LaTeX, making your input file as easy to read
as possible will be a great help as you write
your document and when you change it.  This sample
file shows how you can add comments to your own input
file.

Because printing is different from typewriting,
there are a number of things that you have to do
differently when preparing an input file than if
you were just typing the document directly.
Quotation marks like
       ``this''
have to be handled specially, as do quotes within
quotes:
       ``\,`this'            % \, separates the double and single quote.
        is what I just
        wrote, not  `that'\,''.

Dashes come in three sizes: an
       intra-word
dash, a medium dash for number ranges like
       1--2,
and a punctuation
       dash---like
this.

A sentence-ending space should be larger than the
space between words within a sentence.  You
sometimes have to type special commands in
conjunction with punctuation characters to get
this right, as in the following sentence.
       Gnats, gnus, etc.\ all  % `\ ' makes an inter-word space.
       begin with G\@.         % \@ marks end-of-sentence punctuation.
You should check the spaces after periods when
reading your output to make sure you haven't
forgotten any special cases.  Generating an
ellipsis
       \ldots\               % `\ ' is needed after `\ldots' because TeX
                             % ignores spaces after command names like \ldots
                             % made from \ + letters.
                             %
                             % Note how a `%' character causes TeX to ignore
                             % the end of the input line, so these blank lines
                             % do not start a new paragraph.
                             %
with the right spacing around the periods requires
a special command.

\LaTeX\ interprets some common characters as
commands, so you must type special commands to
generate them.  These characters include the
following:
       \$ \& \% \# \{ and \}.

In printing, text is usually emphasized with an
       \emph{italic}
type style.

\begin{em}
   A long segment of text can also be emphasized
   in this way.  Text within such a segment can be
   given \emph{additional} emphasis.
\end{em}

It is sometimes necessary to prevent \LaTeX\ from
breaking a line where it might otherwise do so.
This may be at a space, as between the ``Mr.''\ and
``Jones'' in
       ``Mr.~Jones'',        % ~ produces an unbreakable interword space.
or within a word---especially when the word is a
symbol like
       \mbox{\emph{itemnum}}
that makes little sense when hyphenated across
lines.

Footnotes\footnote{This is an example of a footnote.}
pose no problem.

\LaTeX\ is good at typesetting mathematical formulas
like
       \( x-3y + z = 7 \)
or
       \( a_{1} > x^{2n} + y^{2n} > x' \)
or
       \( \ip{A}{B} = \sum_{i} a_{i} b_{i} \).
The spaces you type in a formula are
ignored.  Remember that a letter like
       $x$                   % $ ... $  and  \( ... \)  are equivalent
is a formula when it denotes a mathematical
symbol, and it should be typed as one.

\section{Displayed Text}

Text is displayed by indenting it from the left
margin.  Quotations are commonly displayed.  There
are short quotations
\begin{quote}
   This is a short quotation.  It consists of a
   single paragraph of text.  See how it is formatted.
\end{quote}
and longer ones.
\begin{quotation}
   This is a longer quotation.  It consists of two
   paragraphs of text, neither of which are
   particularly interesting.

   This is the second paragraph of the quotation.  It
   is just as dull as the first paragraph.
\end{quotation}
Another frequently-displayed structure is a list.
The following is an example of an \emph{itemized}
list.
\begin{itemize}
   \item This is the first item of an itemized list.
         Each item in the list is marked with a ``tick''.
         You don't have to worry about what kind of tick
         mark is used.

   \item This is the second item of the list.  It
         contains another list nested inside it.  The inner
         list is an \emph{enumerated} list.
         \begin{enumerate}
            \item This is the first item of an enumerated
                  list that is nested within the itemized list.

            \item This is the second item of the inner list.
                  \LaTeX\ allows you to nest lists deeper than
                  you really should.
         \end{enumerate}
         This is the rest of the second item of the outer
         list.  It is no more interesting than any other
         part of the item.
   \item This is the third item of the list.
\end{itemize}
You can even display poetry.
\begin{verse}
   There is an environment
    for verse \\             % The \\ command separates lines
   Whose features some poets % within a stanza.
   will curse.

                             % One or more blank lines separate stanzas.

   For instead of making\\
   Them do \emph{all} line breaking, \\
   It allows them to put too many words on a line when they'd rather be
   forced to be terse.
\end{verse}

Mathematical formulas may also be displayed.  A
displayed formula
is
one-line long; multiline
formulas require special formatting instructions.
   \[  \ip{\Gamma}{\psi'} = x'' + y^{2} + z_{i}^{n}\]
Don't start a paragraph with a displayed equation,
nor make one a paragraph by itself.

\end{document}               % End of document.

\section{Bench Huge Section 1467}
% synthetic bench section 1467
% This is a sample LaTeX input file.  (Version of 12 August 2004.)
%
% A '%' character causes TeX to ignore all remaining text on the line,
% and is used for comments like this one.

\documentclass{article}      % Specifies the document class

                             % The preamble begins here.
\title{An Example Document}  % Declares the document's title.
\author{Leslie Lamport}      % Declares the author's name.
\date{January 21, 1994}      % Deleting this command produces today's date.

\newcommand{\ip}[2]{(#1, #2)}
                             % Defines \ip{arg1}{arg2} to mean
                             % (arg1, arg2).

%\newcommand{\ip}[2]{\langle #1 | #2\rangle}
                             % This is an alternative definition of
                             % \ip that is commented out.

\begin{document}             % End of preamble and beginning of text.

\maketitle                   % Produces the title.

This is an example input file.  Comparing it with
the output it generates can show you how to
produce a simple document of your own.

\section{Ordinary Text}      % Produces section heading.  Lower-level
                             % sections are begun with similar
                             % \subsection and \subsubsection commands.

The ends  of words and sentences are marked
  by   spaces. It  doesn't matter how many
spaces    you type; one is as good as 100.  The
end of   a line counts as a space.

One   or more   blank lines denote the  end
of  a paragraph.

Since any number of consecutive spaces are treated
like a single one, the formatting of the input
file makes no difference to
      \LaTeX,                % The \LaTeX command generates the LaTeX logo.
but it makes a difference to you.  When you use
\LaTeX, making your input file as easy to read
as possible will be a great help as you write
your document and when you change it.  This sample
file shows how you can add comments to your own input
file.

Because printing is different from typewriting,
there are a number of things that you have to do
differently when preparing an input file than if
you were just typing the document directly.
Quotation marks like
       ``this''
have to be handled specially, as do quotes within
quotes:
       ``\,`this'            % \, separates the double and single quote.
        is what I just
        wrote, not  `that'\,''.

Dashes come in three sizes: an
       intra-word
dash, a medium dash for number ranges like
       1--2,
and a punctuation
       dash---like
this.

A sentence-ending space should be larger than the
space between words within a sentence.  You
sometimes have to type special commands in
conjunction with punctuation characters to get
this right, as in the following sentence.
       Gnats, gnus, etc.\ all  % `\ ' makes an inter-word space.
       begin with G\@.         % \@ marks end-of-sentence punctuation.
You should check the spaces after periods when
reading your output to make sure you haven't
forgotten any special cases.  Generating an
ellipsis
       \ldots\               % `\ ' is needed after `\ldots' because TeX
                             % ignores spaces after command names like \ldots
                             % made from \ + letters.
                             %
                             % Note how a `%' character causes TeX to ignore
                             % the end of the input line, so these blank lines
                             % do not start a new paragraph.
                             %
with the right spacing around the periods requires
a special command.

\LaTeX\ interprets some common characters as
commands, so you must type special commands to
generate them.  These characters include the
following:
       \$ \& \% \# \{ and \}.

In printing, text is usually emphasized with an
       \emph{italic}
type style.

\begin{em}
   A long segment of text can also be emphasized
   in this way.  Text within such a segment can be
   given \emph{additional} emphasis.
\end{em}

It is sometimes necessary to prevent \LaTeX\ from
breaking a line where it might otherwise do so.
This may be at a space, as between the ``Mr.''\ and
``Jones'' in
       ``Mr.~Jones'',        % ~ produces an unbreakable interword space.
or within a word---especially when the word is a
symbol like
       \mbox{\emph{itemnum}}
that makes little sense when hyphenated across
lines.

Footnotes\footnote{This is an example of a footnote.}
pose no problem.

\LaTeX\ is good at typesetting mathematical formulas
like
       \( x-3y + z = 7 \)
or
       \( a_{1} > x^{2n} + y^{2n} > x' \)
or
       \( \ip{A}{B} = \sum_{i} a_{i} b_{i} \).
The spaces you type in a formula are
ignored.  Remember that a letter like
       $x$                   % $ ... $  and  \( ... \)  are equivalent
is a formula when it denotes a mathematical
symbol, and it should be typed as one.

\section{Displayed Text}

Text is displayed by indenting it from the left
margin.  Quotations are commonly displayed.  There
are short quotations
\begin{quote}
   This is a short quotation.  It consists of a
   single paragraph of text.  See how it is formatted.
\end{quote}
and longer ones.
\begin{quotation}
   This is a longer quotation.  It consists of two
   paragraphs of text, neither of which are
   particularly interesting.

   This is the second paragraph of the quotation.  It
   is just as dull as the first paragraph.
\end{quotation}
Another frequently-displayed structure is a list.
The following is an example of an \emph{itemized}
list.
\begin{itemize}
   \item This is the first item of an itemized list.
         Each item in the list is marked with a ``tick''.
         You don't have to worry about what kind of tick
         mark is used.

   \item This is the second item of the list.  It
         contains another list nested inside it.  The inner
         list is an \emph{enumerated} list.
         \begin{enumerate}
            \item This is the first item of an enumerated
                  list that is nested within the itemized list.

            \item This is the second item of the inner list.
                  \LaTeX\ allows you to nest lists deeper than
                  you really should.
         \end{enumerate}
         This is the rest of the second item of the outer
         list.  It is no more interesting than any other
         part of the item.
   \item This is the third item of the list.
\end{itemize}
You can even display poetry.
\begin{verse}
   There is an environment
    for verse \\             % The \\ command separates lines
   Whose features some poets % within a stanza.
   will curse.

                             % One or more blank lines separate stanzas.

   For instead of making\\
   Them do \emph{all} line breaking, \\
   It allows them to put too many words on a line when they'd rather be
   forced to be terse.
\end{verse}

Mathematical formulas may also be displayed.  A
displayed formula
is
one-line long; multiline
formulas require special formatting instructions.
   \[  \ip{\Gamma}{\psi'} = x'' + y^{2} + z_{i}^{n}\]
Don't start a paragraph with a displayed equation,
nor make one a paragraph by itself.

\end{document}               % End of document.

\section{Bench Huge Section 1468}
% synthetic bench section 1468
% This is a sample LaTeX input file.  (Version of 12 August 2004.)
%
% A '%' character causes TeX to ignore all remaining text on the line,
% and is used for comments like this one.

\documentclass{article}      % Specifies the document class

                             % The preamble begins here.
\title{An Example Document}  % Declares the document's title.
\author{Leslie Lamport}      % Declares the author's name.
\date{January 21, 1994}      % Deleting this command produces today's date.

\newcommand{\ip}[2]{(#1, #2)}
                             % Defines \ip{arg1}{arg2} to mean
                             % (arg1, arg2).

%\newcommand{\ip}[2]{\langle #1 | #2\rangle}
                             % This is an alternative definition of
                             % \ip that is commented out.

\begin{document}             % End of preamble and beginning of text.

\maketitle                   % Produces the title.

This is an example input file.  Comparing it with
the output it generates can show you how to
produce a simple document of your own.

\section{Ordinary Text}      % Produces section heading.  Lower-level
                             % sections are begun with similar
                             % \subsection and \subsubsection commands.

The ends  of words and sentences are marked
  by   spaces. It  doesn't matter how many
spaces    you type; one is as good as 100.  The
end of   a line counts as a space.

One   or more   blank lines denote the  end
of  a paragraph.

Since any number of consecutive spaces are treated
like a single one, the formatting of the input
file makes no difference to
      \LaTeX,                % The \LaTeX command generates the LaTeX logo.
but it makes a difference to you.  When you use
\LaTeX, making your input file as easy to read
as possible will be a great help as you write
your document and when you change it.  This sample
file shows how you can add comments to your own input
file.

Because printing is different from typewriting,
there are a number of things that you have to do
differently when preparing an input file than if
you were just typing the document directly.
Quotation marks like
       ``this''
have to be handled specially, as do quotes within
quotes:
       ``\,`this'            % \, separates the double and single quote.
        is what I just
        wrote, not  `that'\,''.

Dashes come in three sizes: an
       intra-word
dash, a medium dash for number ranges like
       1--2,
and a punctuation
       dash---like
this.

A sentence-ending space should be larger than the
space between words within a sentence.  You
sometimes have to type special commands in
conjunction with punctuation characters to get
this right, as in the following sentence.
       Gnats, gnus, etc.\ all  % `\ ' makes an inter-word space.
       begin with G\@.         % \@ marks end-of-sentence punctuation.
You should check the spaces after periods when
reading your output to make sure you haven't
forgotten any special cases.  Generating an
ellipsis
       \ldots\               % `\ ' is needed after `\ldots' because TeX
                             % ignores spaces after command names like \ldots
                             % made from \ + letters.
                             %
                             % Note how a `%' character causes TeX to ignore
                             % the end of the input line, so these blank lines
                             % do not start a new paragraph.
                             %
with the right spacing around the periods requires
a special command.

\LaTeX\ interprets some common characters as
commands, so you must type special commands to
generate them.  These characters include the
following:
       \$ \& \% \# \{ and \}.

In printing, text is usually emphasized with an
       \emph{italic}
type style.

\begin{em}
   A long segment of text can also be emphasized
   in this way.  Text within such a segment can be
   given \emph{additional} emphasis.
\end{em}

It is sometimes necessary to prevent \LaTeX\ from
breaking a line where it might otherwise do so.
This may be at a space, as between the ``Mr.''\ and
``Jones'' in
       ``Mr.~Jones'',        % ~ produces an unbreakable interword space.
or within a word---especially when the word is a
symbol like
       \mbox{\emph{itemnum}}
that makes little sense when hyphenated across
lines.

Footnotes\footnote{This is an example of a footnote.}
pose no problem.

\LaTeX\ is good at typesetting mathematical formulas
like
       \( x-3y + z = 7 \)
or
       \( a_{1} > x^{2n} + y^{2n} > x' \)
or
       \( \ip{A}{B} = \sum_{i} a_{i} b_{i} \).
The spaces you type in a formula are
ignored.  Remember that a letter like
       $x$                   % $ ... $  and  \( ... \)  are equivalent
is a formula when it denotes a mathematical
symbol, and it should be typed as one.

\section{Displayed Text}

Text is displayed by indenting it from the left
margin.  Quotations are commonly displayed.  There
are short quotations
\begin{quote}
   This is a short quotation.  It consists of a
   single paragraph of text.  See how it is formatted.
\end{quote}
and longer ones.
\begin{quotation}
   This is a longer quotation.  It consists of two
   paragraphs of text, neither of which are
   particularly interesting.

   This is the second paragraph of the quotation.  It
   is just as dull as the first paragraph.
\end{quotation}
Another frequently-displayed structure is a list.
The following is an example of an \emph{itemized}
list.
\begin{itemize}
   \item This is the first item of an itemized list.
         Each item in the list is marked with a ``tick''.
         You don't have to worry about what kind of tick
         mark is used.

   \item This is the second item of the list.  It
         contains another list nested inside it.  The inner
         list is an \emph{enumerated} list.
         \begin{enumerate}
            \item This is the first item of an enumerated
                  list that is nested within the itemized list.

            \item This is the second item of the inner list.
                  \LaTeX\ allows you to nest lists deeper than
                  you really should.
         \end{enumerate}
         This is the rest of the second item of the outer
         list.  It is no more interesting than any other
         part of the item.
   \item This is the third item of the list.
\end{itemize}
You can even display poetry.
\begin{verse}
   There is an environment
    for verse \\             % The \\ command separates lines
   Whose features some poets % within a stanza.
   will curse.

                             % One or more blank lines separate stanzas.

   For instead of making\\
   Them do \emph{all} line breaking, \\
   It allows them to put too many words on a line when they'd rather be
   forced to be terse.
\end{verse}

Mathematical formulas may also be displayed.  A
displayed formula
is
one-line long; multiline
formulas require special formatting instructions.
   \[  \ip{\Gamma}{\psi'} = x'' + y^{2} + z_{i}^{n}\]
Don't start a paragraph with a displayed equation,
nor make one a paragraph by itself.

\end{document}               % End of document.

\section{Bench Huge Section 1469}
% synthetic bench section 1469
% This is a sample LaTeX input file.  (Version of 12 August 2004.)
%
% A '%' character causes TeX to ignore all remaining text on the line,
% and is used for comments like this one.

\documentclass{article}      % Specifies the document class

                             % The preamble begins here.
\title{An Example Document}  % Declares the document's title.
\author{Leslie Lamport}      % Declares the author's name.
\date{January 21, 1994}      % Deleting this command produces today's date.

\newcommand{\ip}[2]{(#1, #2)}
                             % Defines \ip{arg1}{arg2} to mean
                             % (arg1, arg2).

%\newcommand{\ip}[2]{\langle #1 | #2\rangle}
                             % This is an alternative definition of
                             % \ip that is commented out.

\begin{document}             % End of preamble and beginning of text.

\maketitle                   % Produces the title.

This is an example input file.  Comparing it with
the output it generates can show you how to
produce a simple document of your own.

\section{Ordinary Text}      % Produces section heading.  Lower-level
                             % sections are begun with similar
                             % \subsection and \subsubsection commands.

The ends  of words and sentences are marked
  by   spaces. It  doesn't matter how many
spaces    you type; one is as good as 100.  The
end of   a line counts as a space.

One   or more   blank lines denote the  end
of  a paragraph.

Since any number of consecutive spaces are treated
like a single one, the formatting of the input
file makes no difference to
      \LaTeX,                % The \LaTeX command generates the LaTeX logo.
but it makes a difference to you.  When you use
\LaTeX, making your input file as easy to read
as possible will be a great help as you write
your document and when you change it.  This sample
file shows how you can add comments to your own input
file.

Because printing is different from typewriting,
there are a number of things that you have to do
differently when preparing an input file than if
you were just typing the document directly.
Quotation marks like
       ``this''
have to be handled specially, as do quotes within
quotes:
       ``\,`this'            % \, separates the double and single quote.
        is what I just
        wrote, not  `that'\,''.

Dashes come in three sizes: an
       intra-word
dash, a medium dash for number ranges like
       1--2,
and a punctuation
       dash---like
this.

A sentence-ending space should be larger than the
space between words within a sentence.  You
sometimes have to type special commands in
conjunction with punctuation characters to get
this right, as in the following sentence.
       Gnats, gnus, etc.\ all  % `\ ' makes an inter-word space.
       begin with G\@.         % \@ marks end-of-sentence punctuation.
You should check the spaces after periods when
reading your output to make sure you haven't
forgotten any special cases.  Generating an
ellipsis
       \ldots\               % `\ ' is needed after `\ldots' because TeX
                             % ignores spaces after command names like \ldots
                             % made from \ + letters.
                             %
                             % Note how a `%' character causes TeX to ignore
                             % the end of the input line, so these blank lines
                             % do not start a new paragraph.
                             %
with the right spacing around the periods requires
a special command.

\LaTeX\ interprets some common characters as
commands, so you must type special commands to
generate them.  These characters include the
following:
       \$ \& \% \# \{ and \}.

In printing, text is usually emphasized with an
       \emph{italic}
type style.

\begin{em}
   A long segment of text can also be emphasized
   in this way.  Text within such a segment can be
   given \emph{additional} emphasis.
\end{em}

It is sometimes necessary to prevent \LaTeX\ from
breaking a line where it might otherwise do so.
This may be at a space, as between the ``Mr.''\ and
``Jones'' in
       ``Mr.~Jones'',        % ~ produces an unbreakable interword space.
or within a word---especially when the word is a
symbol like
       \mbox{\emph{itemnum}}
that makes little sense when hyphenated across
lines.

Footnotes\footnote{This is an example of a footnote.}
pose no problem.

\LaTeX\ is good at typesetting mathematical formulas
like
       \( x-3y + z = 7 \)
or
       \( a_{1} > x^{2n} + y^{2n} > x' \)
or
       \( \ip{A}{B} = \sum_{i} a_{i} b_{i} \).
The spaces you type in a formula are
ignored.  Remember that a letter like
       $x$                   % $ ... $  and  \( ... \)  are equivalent
is a formula when it denotes a mathematical
symbol, and it should be typed as one.

\section{Displayed Text}

Text is displayed by indenting it from the left
margin.  Quotations are commonly displayed.  There
are short quotations
\begin{quote}
   This is a short quotation.  It consists of a
   single paragraph of text.  See how it is formatted.
\end{quote}
and longer ones.
\begin{quotation}
   This is a longer quotation.  It consists of two
   paragraphs of text, neither of which are
   particularly interesting.

   This is the second paragraph of the quotation.  It
   is just as dull as the first paragraph.
\end{quotation}
Another frequently-displayed structure is a list.
The following is an example of an \emph{itemized}
list.
\begin{itemize}
   \item This is the first item of an itemized list.
         Each item in the list is marked with a ``tick''.
         You don't have to worry about what kind of tick
         mark is used.

   \item This is the second item of the list.  It
         contains another list nested inside it.  The inner
         list is an \emph{enumerated} list.
         \begin{enumerate}
            \item This is the first item of an enumerated
                  list that is nested within the itemized list.

            \item This is the second item of the inner list.
                  \LaTeX\ allows you to nest lists deeper than
                  you really should.
         \end{enumerate}
         This is the rest of the second item of the outer
         list.  It is no more interesting than any other
         part of the item.
   \item This is the third item of the list.
\end{itemize}
You can even display poetry.
\begin{verse}
   There is an environment
    for verse \\             % The \\ command separates lines
   Whose features some poets % within a stanza.
   will curse.

                             % One or more blank lines separate stanzas.

   For instead of making\\
   Them do \emph{all} line breaking, \\
   It allows them to put too many words on a line when they'd rather be
   forced to be terse.
\end{verse}

Mathematical formulas may also be displayed.  A
displayed formula
is
one-line long; multiline
formulas require special formatting instructions.
   \[  \ip{\Gamma}{\psi'} = x'' + y^{2} + z_{i}^{n}\]
Don't start a paragraph with a displayed equation,
nor make one a paragraph by itself.

\end{document}               % End of document.

\section{Bench Huge Section 1470}
% synthetic bench section 1470
% This is a sample LaTeX input file.  (Version of 12 August 2004.)
%
% A '%' character causes TeX to ignore all remaining text on the line,
% and is used for comments like this one.

\documentclass{article}      % Specifies the document class

                             % The preamble begins here.
\title{An Example Document}  % Declares the document's title.
\author{Leslie Lamport}      % Declares the author's name.
\date{January 21, 1994}      % Deleting this command produces today's date.

\newcommand{\ip}[2]{(#1, #2)}
                             % Defines \ip{arg1}{arg2} to mean
                             % (arg1, arg2).

%\newcommand{\ip}[2]{\langle #1 | #2\rangle}
                             % This is an alternative definition of
                             % \ip that is commented out.

\begin{document}             % End of preamble and beginning of text.

\maketitle                   % Produces the title.

This is an example input file.  Comparing it with
the output it generates can show you how to
produce a simple document of your own.

\section{Ordinary Text}      % Produces section heading.  Lower-level
                             % sections are begun with similar
                             % \subsection and \subsubsection commands.

The ends  of words and sentences are marked
  by   spaces. It  doesn't matter how many
spaces    you type; one is as good as 100.  The
end of   a line counts as a space.

One   or more   blank lines denote the  end
of  a paragraph.

Since any number of consecutive spaces are treated
like a single one, the formatting of the input
file makes no difference to
      \LaTeX,                % The \LaTeX command generates the LaTeX logo.
but it makes a difference to you.  When you use
\LaTeX, making your input file as easy to read
as possible will be a great help as you write
your document and when you change it.  This sample
file shows how you can add comments to your own input
file.

Because printing is different from typewriting,
there are a number of things that you have to do
differently when preparing an input file than if
you were just typing the document directly.
Quotation marks like
       ``this''
have to be handled specially, as do quotes within
quotes:
       ``\,`this'            % \, separates the double and single quote.
        is what I just
        wrote, not  `that'\,''.

Dashes come in three sizes: an
       intra-word
dash, a medium dash for number ranges like
       1--2,
and a punctuation
       dash---like
this.

A sentence-ending space should be larger than the
space between words within a sentence.  You
sometimes have to type special commands in
conjunction with punctuation characters to get
this right, as in the following sentence.
       Gnats, gnus, etc.\ all  % `\ ' makes an inter-word space.
       begin with G\@.         % \@ marks end-of-sentence punctuation.
You should check the spaces after periods when
reading your output to make sure you haven't
forgotten any special cases.  Generating an
ellipsis
       \ldots\               % `\ ' is needed after `\ldots' because TeX
                             % ignores spaces after command names like \ldots
                             % made from \ + letters.
                             %
                             % Note how a `%' character causes TeX to ignore
                             % the end of the input line, so these blank lines
                             % do not start a new paragraph.
                             %
with the right spacing around the periods requires
a special command.

\LaTeX\ interprets some common characters as
commands, so you must type special commands to
generate them.  These characters include the
following:
       \$ \& \% \# \{ and \}.

In printing, text is usually emphasized with an
       \emph{italic}
type style.

\begin{em}
   A long segment of text can also be emphasized
   in this way.  Text within such a segment can be
   given \emph{additional} emphasis.
\end{em}

It is sometimes necessary to prevent \LaTeX\ from
breaking a line where it might otherwise do so.
This may be at a space, as between the ``Mr.''\ and
``Jones'' in
       ``Mr.~Jones'',        % ~ produces an unbreakable interword space.
or within a word---especially when the word is a
symbol like
       \mbox{\emph{itemnum}}
that makes little sense when hyphenated across
lines.

Footnotes\footnote{This is an example of a footnote.}
pose no problem.

\LaTeX\ is good at typesetting mathematical formulas
like
       \( x-3y + z = 7 \)
or
       \( a_{1} > x^{2n} + y^{2n} > x' \)
or
       \( \ip{A}{B} = \sum_{i} a_{i} b_{i} \).
The spaces you type in a formula are
ignored.  Remember that a letter like
       $x$                   % $ ... $  and  \( ... \)  are equivalent
is a formula when it denotes a mathematical
symbol, and it should be typed as one.

\section{Displayed Text}

Text is displayed by indenting it from the left
margin.  Quotations are commonly displayed.  There
are short quotations
\begin{quote}
   This is a short quotation.  It consists of a
   single paragraph of text.  See how it is formatted.
\end{quote}
and longer ones.
\begin{quotation}
   This is a longer quotation.  It consists of two
   paragraphs of text, neither of which are
   particularly interesting.

   This is the second paragraph of the quotation.  It
   is just as dull as the first paragraph.
\end{quotation}
Another frequently-displayed structure is a list.
The following is an example of an \emph{itemized}
list.
\begin{itemize}
   \item This is the first item of an itemized list.
         Each item in the list is marked with a ``tick''.
         You don't have to worry about what kind of tick
         mark is used.

   \item This is the second item of the list.  It
         contains another list nested inside it.  The inner
         list is an \emph{enumerated} list.
         \begin{enumerate}
            \item This is the first item of an enumerated
                  list that is nested within the itemized list.

            \item This is the second item of the inner list.
                  \LaTeX\ allows you to nest lists deeper than
                  you really should.
         \end{enumerate}
         This is the rest of the second item of the outer
         list.  It is no more interesting than any other
         part of the item.
   \item This is the third item of the list.
\end{itemize}
You can even display poetry.
\begin{verse}
   There is an environment
    for verse \\             % The \\ command separates lines
   Whose features some poets % within a stanza.
   will curse.

                             % One or more blank lines separate stanzas.

   For instead of making\\
   Them do \emph{all} line breaking, \\
   It allows them to put too many words on a line when they'd rather be
   forced to be terse.
\end{verse}

Mathematical formulas may also be displayed.  A
displayed formula
is
one-line long; multiline
formulas require special formatting instructions.
   \[  \ip{\Gamma}{\psi'} = x'' + y^{2} + z_{i}^{n}\]
Don't start a paragraph with a displayed equation,
nor make one a paragraph by itself.

\end{document}               % End of document.

\section{Bench Huge Section 1471}
% synthetic bench section 1471
% This is a sample LaTeX input file.  (Version of 12 August 2004.)
%
% A '%' character causes TeX to ignore all remaining text on the line,
% and is used for comments like this one.

\documentclass{article}      % Specifies the document class

                             % The preamble begins here.
\title{An Example Document}  % Declares the document's title.
\author{Leslie Lamport}      % Declares the author's name.
\date{January 21, 1994}      % Deleting this command produces today's date.

\newcommand{\ip}[2]{(#1, #2)}
                             % Defines \ip{arg1}{arg2} to mean
                             % (arg1, arg2).

%\newcommand{\ip}[2]{\langle #1 | #2\rangle}
                             % This is an alternative definition of
                             % \ip that is commented out.

\begin{document}             % End of preamble and beginning of text.

\maketitle                   % Produces the title.

This is an example input file.  Comparing it with
the output it generates can show you how to
produce a simple document of your own.

\section{Ordinary Text}      % Produces section heading.  Lower-level
                             % sections are begun with similar
                             % \subsection and \subsubsection commands.

The ends  of words and sentences are marked
  by   spaces. It  doesn't matter how many
spaces    you type; one is as good as 100.  The
end of   a line counts as a space.

One   or more   blank lines denote the  end
of  a paragraph.

Since any number of consecutive spaces are treated
like a single one, the formatting of the input
file makes no difference to
      \LaTeX,                % The \LaTeX command generates the LaTeX logo.
but it makes a difference to you.  When you use
\LaTeX, making your input file as easy to read
as possible will be a great help as you write
your document and when you change it.  This sample
file shows how you can add comments to your own input
file.

Because printing is different from typewriting,
there are a number of things that you have to do
differently when preparing an input file than if
you were just typing the document directly.
Quotation marks like
       ``this''
have to be handled specially, as do quotes within
quotes:
       ``\,`this'            % \, separates the double and single quote.
        is what I just
        wrote, not  `that'\,''.

Dashes come in three sizes: an
       intra-word
dash, a medium dash for number ranges like
       1--2,
and a punctuation
       dash---like
this.

A sentence-ending space should be larger than the
space between words within a sentence.  You
sometimes have to type special commands in
conjunction with punctuation characters to get
this right, as in the following sentence.
       Gnats, gnus, etc.\ all  % `\ ' makes an inter-word space.
       begin with G\@.         % \@ marks end-of-sentence punctuation.
You should check the spaces after periods when
reading your output to make sure you haven't
forgotten any special cases.  Generating an
ellipsis
       \ldots\               % `\ ' is needed after `\ldots' because TeX
                             % ignores spaces after command names like \ldots
                             % made from \ + letters.
                             %
                             % Note how a `%' character causes TeX to ignore
                             % the end of the input line, so these blank lines
                             % do not start a new paragraph.
                             %
with the right spacing around the periods requires
a special command.

\LaTeX\ interprets some common characters as
commands, so you must type special commands to
generate them.  These characters include the
following:
       \$ \& \% \# \{ and \}.

In printing, text is usually emphasized with an
       \emph{italic}
type style.

\begin{em}
   A long segment of text can also be emphasized
   in this way.  Text within such a segment can be
   given \emph{additional} emphasis.
\end{em}

It is sometimes necessary to prevent \LaTeX\ from
breaking a line where it might otherwise do so.
This may be at a space, as between the ``Mr.''\ and
``Jones'' in
       ``Mr.~Jones'',        % ~ produces an unbreakable interword space.
or within a word---especially when the word is a
symbol like
       \mbox{\emph{itemnum}}
that makes little sense when hyphenated across
lines.

Footnotes\footnote{This is an example of a footnote.}
pose no problem.

\LaTeX\ is good at typesetting mathematical formulas
like
       \( x-3y + z = 7 \)
or
       \( a_{1} > x^{2n} + y^{2n} > x' \)
or
       \( \ip{A}{B} = \sum_{i} a_{i} b_{i} \).
The spaces you type in a formula are
ignored.  Remember that a letter like
       $x$                   % $ ... $  and  \( ... \)  are equivalent
is a formula when it denotes a mathematical
symbol, and it should be typed as one.

\section{Displayed Text}

Text is displayed by indenting it from the left
margin.  Quotations are commonly displayed.  There
are short quotations
\begin{quote}
   This is a short quotation.  It consists of a
   single paragraph of text.  See how it is formatted.
\end{quote}
and longer ones.
\begin{quotation}
   This is a longer quotation.  It consists of two
   paragraphs of text, neither of which are
   particularly interesting.

   This is the second paragraph of the quotation.  It
   is just as dull as the first paragraph.
\end{quotation}
Another frequently-displayed structure is a list.
The following is an example of an \emph{itemized}
list.
\begin{itemize}
   \item This is the first item of an itemized list.
         Each item in the list is marked with a ``tick''.
         You don't have to worry about what kind of tick
         mark is used.

   \item This is the second item of the list.  It
         contains another list nested inside it.  The inner
         list is an \emph{enumerated} list.
         \begin{enumerate}
            \item This is the first item of an enumerated
                  list that is nested within the itemized list.

            \item This is the second item of the inner list.
                  \LaTeX\ allows you to nest lists deeper than
                  you really should.
         \end{enumerate}
         This is the rest of the second item of the outer
         list.  It is no more interesting than any other
         part of the item.
   \item This is the third item of the list.
\end{itemize}
You can even display poetry.
\begin{verse}
   There is an environment
    for verse \\             % The \\ command separates lines
   Whose features some poets % within a stanza.
   will curse.

                             % One or more blank lines separate stanzas.

   For instead of making\\
   Them do \emph{all} line breaking, \\
   It allows them to put too many words on a line when they'd rather be
   forced to be terse.
\end{verse}

Mathematical formulas may also be displayed.  A
displayed formula
is
one-line long; multiline
formulas require special formatting instructions.
   \[  \ip{\Gamma}{\psi'} = x'' + y^{2} + z_{i}^{n}\]
Don't start a paragraph with a displayed equation,
nor make one a paragraph by itself.

\end{document}               % End of document.

\section{Bench Huge Section 1472}
% synthetic bench section 1472
% This is a sample LaTeX input file.  (Version of 12 August 2004.)
%
% A '%' character causes TeX to ignore all remaining text on the line,
% and is used for comments like this one.

\documentclass{article}      % Specifies the document class

                             % The preamble begins here.
\title{An Example Document}  % Declares the document's title.
\author{Leslie Lamport}      % Declares the author's name.
\date{January 21, 1994}      % Deleting this command produces today's date.

\newcommand{\ip}[2]{(#1, #2)}
                             % Defines \ip{arg1}{arg2} to mean
                             % (arg1, arg2).

%\newcommand{\ip}[2]{\langle #1 | #2\rangle}
                             % This is an alternative definition of
                             % \ip that is commented out.

\begin{document}             % End of preamble and beginning of text.

\maketitle                   % Produces the title.

This is an example input file.  Comparing it with
the output it generates can show you how to
produce a simple document of your own.

\section{Ordinary Text}      % Produces section heading.  Lower-level
                             % sections are begun with similar
                             % \subsection and \subsubsection commands.

The ends  of words and sentences are marked
  by   spaces. It  doesn't matter how many
spaces    you type; one is as good as 100.  The
end of   a line counts as a space.

One   or more   blank lines denote the  end
of  a paragraph.

Since any number of consecutive spaces are treated
like a single one, the formatting of the input
file makes no difference to
      \LaTeX,                % The \LaTeX command generates the LaTeX logo.
but it makes a difference to you.  When you use
\LaTeX, making your input file as easy to read
as possible will be a great help as you write
your document and when you change it.  This sample
file shows how you can add comments to your own input
file.

Because printing is different from typewriting,
there are a number of things that you have to do
differently when preparing an input file than if
you were just typing the document directly.
Quotation marks like
       ``this''
have to be handled specially, as do quotes within
quotes:
       ``\,`this'            % \, separates the double and single quote.
        is what I just
        wrote, not  `that'\,''.

Dashes come in three sizes: an
       intra-word
dash, a medium dash for number ranges like
       1--2,
and a punctuation
       dash---like
this.

A sentence-ending space should be larger than the
space between words within a sentence.  You
sometimes have to type special commands in
conjunction with punctuation characters to get
this right, as in the following sentence.
       Gnats, gnus, etc.\ all  % `\ ' makes an inter-word space.
       begin with G\@.         % \@ marks end-of-sentence punctuation.
You should check the spaces after periods when
reading your output to make sure you haven't
forgotten any special cases.  Generating an
ellipsis
       \ldots\               % `\ ' is needed after `\ldots' because TeX
                             % ignores spaces after command names like \ldots
                             % made from \ + letters.
                             %
                             % Note how a `%' character causes TeX to ignore
                             % the end of the input line, so these blank lines
                             % do not start a new paragraph.
                             %
with the right spacing around the periods requires
a special command.

\LaTeX\ interprets some common characters as
commands, so you must type special commands to
generate them.  These characters include the
following:
       \$ \& \% \# \{ and \}.

In printing, text is usually emphasized with an
       \emph{italic}
type style.

\begin{em}
   A long segment of text can also be emphasized
   in this way.  Text within such a segment can be
   given \emph{additional} emphasis.
\end{em}

It is sometimes necessary to prevent \LaTeX\ from
breaking a line where it might otherwise do so.
This may be at a space, as between the ``Mr.''\ and
``Jones'' in
       ``Mr.~Jones'',        % ~ produces an unbreakable interword space.
or within a word---especially when the word is a
symbol like
       \mbox{\emph{itemnum}}
that makes little sense when hyphenated across
lines.

Footnotes\footnote{This is an example of a footnote.}
pose no problem.

\LaTeX\ is good at typesetting mathematical formulas
like
       \( x-3y + z = 7 \)
or
       \( a_{1} > x^{2n} + y^{2n} > x' \)
or
       \( \ip{A}{B} = \sum_{i} a_{i} b_{i} \).
The spaces you type in a formula are
ignored.  Remember that a letter like
       $x$                   % $ ... $  and  \( ... \)  are equivalent
is a formula when it denotes a mathematical
symbol, and it should be typed as one.

\section{Displayed Text}

Text is displayed by indenting it from the left
margin.  Quotations are commonly displayed.  There
are short quotations
\begin{quote}
   This is a short quotation.  It consists of a
   single paragraph of text.  See how it is formatted.
\end{quote}
and longer ones.
\begin{quotation}
   This is a longer quotation.  It consists of two
   paragraphs of text, neither of which are
   particularly interesting.

   This is the second paragraph of the quotation.  It
   is just as dull as the first paragraph.
\end{quotation}
Another frequently-displayed structure is a list.
The following is an example of an \emph{itemized}
list.
\begin{itemize}
   \item This is the first item of an itemized list.
         Each item in the list is marked with a ``tick''.
         You don't have to worry about what kind of tick
         mark is used.

   \item This is the second item of the list.  It
         contains another list nested inside it.  The inner
         list is an \emph{enumerated} list.
         \begin{enumerate}
            \item This is the first item of an enumerated
                  list that is nested within the itemized list.

            \item This is the second item of the inner list.
                  \LaTeX\ allows you to nest lists deeper than
                  you really should.
         \end{enumerate}
         This is the rest of the second item of the outer
         list.  It is no more interesting than any other
         part of the item.
   \item This is the third item of the list.
\end{itemize}
You can even display poetry.
\begin{verse}
   There is an environment
    for verse \\             % The \\ command separates lines
   Whose features some poets % within a stanza.
   will curse.

                             % One or more blank lines separate stanzas.

   For instead of making\\
   Them do \emph{all} line breaking, \\
   It allows them to put too many words on a line when they'd rather be
   forced to be terse.
\end{verse}

Mathematical formulas may also be displayed.  A
displayed formula
is
one-line long; multiline
formulas require special formatting instructions.
   \[  \ip{\Gamma}{\psi'} = x'' + y^{2} + z_{i}^{n}\]
Don't start a paragraph with a displayed equation,
nor make one a paragraph by itself.

\end{document}               % End of document.

\section{Bench Huge Section 1473}
% synthetic bench section 1473
% This is a sample LaTeX input file.  (Version of 12 August 2004.)
%
% A '%' character causes TeX to ignore all remaining text on the line,
% and is used for comments like this one.

\documentclass{article}      % Specifies the document class

                             % The preamble begins here.
\title{An Example Document}  % Declares the document's title.
\author{Leslie Lamport}      % Declares the author's name.
\date{January 21, 1994}      % Deleting this command produces today's date.

\newcommand{\ip}[2]{(#1, #2)}
                             % Defines \ip{arg1}{arg2} to mean
                             % (arg1, arg2).

%\newcommand{\ip}[2]{\langle #1 | #2\rangle}
                             % This is an alternative definition of
                             % \ip that is commented out.

\begin{document}             % End of preamble and beginning of text.

\maketitle                   % Produces the title.

This is an example input file.  Comparing it with
the output it generates can show you how to
produce a simple document of your own.

\section{Ordinary Text}      % Produces section heading.  Lower-level
                             % sections are begun with similar
                             % \subsection and \subsubsection commands.

The ends  of words and sentences are marked
  by   spaces. It  doesn't matter how many
spaces    you type; one is as good as 100.  The
end of   a line counts as a space.

One   or more   blank lines denote the  end
of  a paragraph.

Since any number of consecutive spaces are treated
like a single one, the formatting of the input
file makes no difference to
      \LaTeX,                % The \LaTeX command generates the LaTeX logo.
but it makes a difference to you.  When you use
\LaTeX, making your input file as easy to read
as possible will be a great help as you write
your document and when you change it.  This sample
file shows how you can add comments to your own input
file.

Because printing is different from typewriting,
there are a number of things that you have to do
differently when preparing an input file than if
you were just typing the document directly.
Quotation marks like
       ``this''
have to be handled specially, as do quotes within
quotes:
       ``\,`this'            % \, separates the double and single quote.
        is what I just
        wrote, not  `that'\,''.

Dashes come in three sizes: an
       intra-word
dash, a medium dash for number ranges like
       1--2,
and a punctuation
       dash---like
this.

A sentence-ending space should be larger than the
space between words within a sentence.  You
sometimes have to type special commands in
conjunction with punctuation characters to get
this right, as in the following sentence.
       Gnats, gnus, etc.\ all  % `\ ' makes an inter-word space.
       begin with G\@.         % \@ marks end-of-sentence punctuation.
You should check the spaces after periods when
reading your output to make sure you haven't
forgotten any special cases.  Generating an
ellipsis
       \ldots\               % `\ ' is needed after `\ldots' because TeX
                             % ignores spaces after command names like \ldots
                             % made from \ + letters.
                             %
                             % Note how a `%' character causes TeX to ignore
                             % the end of the input line, so these blank lines
                             % do not start a new paragraph.
                             %
with the right spacing around the periods requires
a special command.

\LaTeX\ interprets some common characters as
commands, so you must type special commands to
generate them.  These characters include the
following:
       \$ \& \% \# \{ and \}.

In printing, text is usually emphasized with an
       \emph{italic}
type style.

\begin{em}
   A long segment of text can also be emphasized
   in this way.  Text within such a segment can be
   given \emph{additional} emphasis.
\end{em}

It is sometimes necessary to prevent \LaTeX\ from
breaking a line where it might otherwise do so.
This may be at a space, as between the ``Mr.''\ and
``Jones'' in
       ``Mr.~Jones'',        % ~ produces an unbreakable interword space.
or within a word---especially when the word is a
symbol like
       \mbox{\emph{itemnum}}
that makes little sense when hyphenated across
lines.

Footnotes\footnote{This is an example of a footnote.}
pose no problem.

\LaTeX\ is good at typesetting mathematical formulas
like
       \( x-3y + z = 7 \)
or
       \( a_{1} > x^{2n} + y^{2n} > x' \)
or
       \( \ip{A}{B} = \sum_{i} a_{i} b_{i} \).
The spaces you type in a formula are
ignored.  Remember that a letter like
       $x$                   % $ ... $  and  \( ... \)  are equivalent
is a formula when it denotes a mathematical
symbol, and it should be typed as one.

\section{Displayed Text}

Text is displayed by indenting it from the left
margin.  Quotations are commonly displayed.  There
are short quotations
\begin{quote}
   This is a short quotation.  It consists of a
   single paragraph of text.  See how it is formatted.
\end{quote}
and longer ones.
\begin{quotation}
   This is a longer quotation.  It consists of two
   paragraphs of text, neither of which are
   particularly interesting.

   This is the second paragraph of the quotation.  It
   is just as dull as the first paragraph.
\end{quotation}
Another frequently-displayed structure is a list.
The following is an example of an \emph{itemized}
list.
\begin{itemize}
   \item This is the first item of an itemized list.
         Each item in the list is marked with a ``tick''.
         You don't have to worry about what kind of tick
         mark is used.

   \item This is the second item of the list.  It
         contains another list nested inside it.  The inner
         list is an \emph{enumerated} list.
         \begin{enumerate}
            \item This is the first item of an enumerated
                  list that is nested within the itemized list.

            \item This is the second item of the inner list.
                  \LaTeX\ allows you to nest lists deeper than
                  you really should.
         \end{enumerate}
         This is the rest of the second item of the outer
         list.  It is no more interesting than any other
         part of the item.
   \item This is the third item of the list.
\end{itemize}
You can even display poetry.
\begin{verse}
   There is an environment
    for verse \\             % The \\ command separates lines
   Whose features some poets % within a stanza.
   will curse.

                             % One or more blank lines separate stanzas.

   For instead of making\\
   Them do \emph{all} line breaking, \\
   It allows them to put too many words on a line when they'd rather be
   forced to be terse.
\end{verse}

Mathematical formulas may also be displayed.  A
displayed formula
is
one-line long; multiline
formulas require special formatting instructions.
   \[  \ip{\Gamma}{\psi'} = x'' + y^{2} + z_{i}^{n}\]
Don't start a paragraph with a displayed equation,
nor make one a paragraph by itself.

\end{document}               % End of document.

\section{Bench Huge Section 1474}
% synthetic bench section 1474
% This is a sample LaTeX input file.  (Version of 12 August 2004.)
%
% A '%' character causes TeX to ignore all remaining text on the line,
% and is used for comments like this one.

\documentclass{article}      % Specifies the document class

                             % The preamble begins here.
\title{An Example Document}  % Declares the document's title.
\author{Leslie Lamport}      % Declares the author's name.
\date{January 21, 1994}      % Deleting this command produces today's date.

\newcommand{\ip}[2]{(#1, #2)}
                             % Defines \ip{arg1}{arg2} to mean
                             % (arg1, arg2).

%\newcommand{\ip}[2]{\langle #1 | #2\rangle}
                             % This is an alternative definition of
                             % \ip that is commented out.

\begin{document}             % End of preamble and beginning of text.

\maketitle                   % Produces the title.

This is an example input file.  Comparing it with
the output it generates can show you how to
produce a simple document of your own.

\section{Ordinary Text}      % Produces section heading.  Lower-level
                             % sections are begun with similar
                             % \subsection and \subsubsection commands.

The ends  of words and sentences are marked
  by   spaces. It  doesn't matter how many
spaces    you type; one is as good as 100.  The
end of   a line counts as a space.

One   or more   blank lines denote the  end
of  a paragraph.

Since any number of consecutive spaces are treated
like a single one, the formatting of the input
file makes no difference to
      \LaTeX,                % The \LaTeX command generates the LaTeX logo.
but it makes a difference to you.  When you use
\LaTeX, making your input file as easy to read
as possible will be a great help as you write
your document and when you change it.  This sample
file shows how you can add comments to your own input
file.

Because printing is different from typewriting,
there are a number of things that you have to do
differently when preparing an input file than if
you were just typing the document directly.
Quotation marks like
       ``this''
have to be handled specially, as do quotes within
quotes:
       ``\,`this'            % \, separates the double and single quote.
        is what I just
        wrote, not  `that'\,''.

Dashes come in three sizes: an
       intra-word
dash, a medium dash for number ranges like
       1--2,
and a punctuation
       dash---like
this.

A sentence-ending space should be larger than the
space between words within a sentence.  You
sometimes have to type special commands in
conjunction with punctuation characters to get
this right, as in the following sentence.
       Gnats, gnus, etc.\ all  % `\ ' makes an inter-word space.
       begin with G\@.         % \@ marks end-of-sentence punctuation.
You should check the spaces after periods when
reading your output to make sure you haven't
forgotten any special cases.  Generating an
ellipsis
       \ldots\               % `\ ' is needed after `\ldots' because TeX
                             % ignores spaces after command names like \ldots
                             % made from \ + letters.
                             %
                             % Note how a `%' character causes TeX to ignore
                             % the end of the input line, so these blank lines
                             % do not start a new paragraph.
                             %
with the right spacing around the periods requires
a special command.

\LaTeX\ interprets some common characters as
commands, so you must type special commands to
generate them.  These characters include the
following:
       \$ \& \% \# \{ and \}.

In printing, text is usually emphasized with an
       \emph{italic}
type style.

\begin{em}
   A long segment of text can also be emphasized
   in this way.  Text within such a segment can be
   given \emph{additional} emphasis.
\end{em}

It is sometimes necessary to prevent \LaTeX\ from
breaking a line where it might otherwise do so.
This may be at a space, as between the ``Mr.''\ and
``Jones'' in
       ``Mr.~Jones'',        % ~ produces an unbreakable interword space.
or within a word---especially when the word is a
symbol like
       \mbox{\emph{itemnum}}
that makes little sense when hyphenated across
lines.

Footnotes\footnote{This is an example of a footnote.}
pose no problem.

\LaTeX\ is good at typesetting mathematical formulas
like
       \( x-3y + z = 7 \)
or
       \( a_{1} > x^{2n} + y^{2n} > x' \)
or
       \( \ip{A}{B} = \sum_{i} a_{i} b_{i} \).
The spaces you type in a formula are
ignored.  Remember that a letter like
       $x$                   % $ ... $  and  \( ... \)  are equivalent
is a formula when it denotes a mathematical
symbol, and it should be typed as one.

\section{Displayed Text}

Text is displayed by indenting it from the left
margin.  Quotations are commonly displayed.  There
are short quotations
\begin{quote}
   This is a short quotation.  It consists of a
   single paragraph of text.  See how it is formatted.
\end{quote}
and longer ones.
\begin{quotation}
   This is a longer quotation.  It consists of two
   paragraphs of text, neither of which are
   particularly interesting.

   This is the second paragraph of the quotation.  It
   is just as dull as the first paragraph.
\end{quotation}
Another frequently-displayed structure is a list.
The following is an example of an \emph{itemized}
list.
\begin{itemize}
   \item This is the first item of an itemized list.
         Each item in the list is marked with a ``tick''.
         You don't have to worry about what kind of tick
         mark is used.

   \item This is the second item of the list.  It
         contains another list nested inside it.  The inner
         list is an \emph{enumerated} list.
         \begin{enumerate}
            \item This is the first item of an enumerated
                  list that is nested within the itemized list.

            \item This is the second item of the inner list.
                  \LaTeX\ allows you to nest lists deeper than
                  you really should.
         \end{enumerate}
         This is the rest of the second item of the outer
         list.  It is no more interesting than any other
         part of the item.
   \item This is the third item of the list.
\end{itemize}
You can even display poetry.
\begin{verse}
   There is an environment
    for verse \\             % The \\ command separates lines
   Whose features some poets % within a stanza.
   will curse.

                             % One or more blank lines separate stanzas.

   For instead of making\\
   Them do \emph{all} line breaking, \\
   It allows them to put too many words on a line when they'd rather be
   forced to be terse.
\end{verse}

Mathematical formulas may also be displayed.  A
displayed formula
is
one-line long; multiline
formulas require special formatting instructions.
   \[  \ip{\Gamma}{\psi'} = x'' + y^{2} + z_{i}^{n}\]
Don't start a paragraph with a displayed equation,
nor make one a paragraph by itself.

\end{document}               % End of document.

\section{Bench Huge Section 1475}
% synthetic bench section 1475
% This is a sample LaTeX input file.  (Version of 12 August 2004.)
%
% A '%' character causes TeX to ignore all remaining text on the line,
% and is used for comments like this one.

\documentclass{article}      % Specifies the document class

                             % The preamble begins here.
\title{An Example Document}  % Declares the document's title.
\author{Leslie Lamport}      % Declares the author's name.
\date{January 21, 1994}      % Deleting this command produces today's date.

\newcommand{\ip}[2]{(#1, #2)}
                             % Defines \ip{arg1}{arg2} to mean
                             % (arg1, arg2).

%\newcommand{\ip}[2]{\langle #1 | #2\rangle}
                             % This is an alternative definition of
                             % \ip that is commented out.

\begin{document}             % End of preamble and beginning of text.

\maketitle                   % Produces the title.

This is an example input file.  Comparing it with
the output it generates can show you how to
produce a simple document of your own.

\section{Ordinary Text}      % Produces section heading.  Lower-level
                             % sections are begun with similar
                             % \subsection and \subsubsection commands.

The ends  of words and sentences are marked
  by   spaces. It  doesn't matter how many
spaces    you type; one is as good as 100.  The
end of   a line counts as a space.

One   or more   blank lines denote the  end
of  a paragraph.

Since any number of consecutive spaces are treated
like a single one, the formatting of the input
file makes no difference to
      \LaTeX,                % The \LaTeX command generates the LaTeX logo.
but it makes a difference to you.  When you use
\LaTeX, making your input file as easy to read
as possible will be a great help as you write
your document and when you change it.  This sample
file shows how you can add comments to your own input
file.

Because printing is different from typewriting,
there are a number of things that you have to do
differently when preparing an input file than if
you were just typing the document directly.
Quotation marks like
       ``this''
have to be handled specially, as do quotes within
quotes:
       ``\,`this'            % \, separates the double and single quote.
        is what I just
        wrote, not  `that'\,''.

Dashes come in three sizes: an
       intra-word
dash, a medium dash for number ranges like
       1--2,
and a punctuation
       dash---like
this.

A sentence-ending space should be larger than the
space between words within a sentence.  You
sometimes have to type special commands in
conjunction with punctuation characters to get
this right, as in the following sentence.
       Gnats, gnus, etc.\ all  % `\ ' makes an inter-word space.
       begin with G\@.         % \@ marks end-of-sentence punctuation.
You should check the spaces after periods when
reading your output to make sure you haven't
forgotten any special cases.  Generating an
ellipsis
       \ldots\               % `\ ' is needed after `\ldots' because TeX
                             % ignores spaces after command names like \ldots
                             % made from \ + letters.
                             %
                             % Note how a `%' character causes TeX to ignore
                             % the end of the input line, so these blank lines
                             % do not start a new paragraph.
                             %
with the right spacing around the periods requires
a special command.

\LaTeX\ interprets some common characters as
commands, so you must type special commands to
generate them.  These characters include the
following:
       \$ \& \% \# \{ and \}.

In printing, text is usually emphasized with an
       \emph{italic}
type style.

\begin{em}
   A long segment of text can also be emphasized
   in this way.  Text within such a segment can be
   given \emph{additional} emphasis.
\end{em}

It is sometimes necessary to prevent \LaTeX\ from
breaking a line where it might otherwise do so.
This may be at a space, as between the ``Mr.''\ and
``Jones'' in
       ``Mr.~Jones'',        % ~ produces an unbreakable interword space.
or within a word---especially when the word is a
symbol like
       \mbox{\emph{itemnum}}
that makes little sense when hyphenated across
lines.

Footnotes\footnote{This is an example of a footnote.}
pose no problem.

\LaTeX\ is good at typesetting mathematical formulas
like
       \( x-3y + z = 7 \)
or
       \( a_{1} > x^{2n} + y^{2n} > x' \)
or
       \( \ip{A}{B} = \sum_{i} a_{i} b_{i} \).
The spaces you type in a formula are
ignored.  Remember that a letter like
       $x$                   % $ ... $  and  \( ... \)  are equivalent
is a formula when it denotes a mathematical
symbol, and it should be typed as one.

\section{Displayed Text}

Text is displayed by indenting it from the left
margin.  Quotations are commonly displayed.  There
are short quotations
\begin{quote}
   This is a short quotation.  It consists of a
   single paragraph of text.  See how it is formatted.
\end{quote}
and longer ones.
\begin{quotation}
   This is a longer quotation.  It consists of two
   paragraphs of text, neither of which are
   particularly interesting.

   This is the second paragraph of the quotation.  It
   is just as dull as the first paragraph.
\end{quotation}
Another frequently-displayed structure is a list.
The following is an example of an \emph{itemized}
list.
\begin{itemize}
   \item This is the first item of an itemized list.
         Each item in the list is marked with a ``tick''.
         You don't have to worry about what kind of tick
         mark is used.

   \item This is the second item of the list.  It
         contains another list nested inside it.  The inner
         list is an \emph{enumerated} list.
         \begin{enumerate}
            \item This is the first item of an enumerated
                  list that is nested within the itemized list.

            \item This is the second item of the inner list.
                  \LaTeX\ allows you to nest lists deeper than
                  you really should.
         \end{enumerate}
         This is the rest of the second item of the outer
         list.  It is no more interesting than any other
         part of the item.
   \item This is the third item of the list.
\end{itemize}
You can even display poetry.
\begin{verse}
   There is an environment
    for verse \\             % The \\ command separates lines
   Whose features some poets % within a stanza.
   will curse.

                             % One or more blank lines separate stanzas.

   For instead of making\\
   Them do \emph{all} line breaking, \\
   It allows them to put too many words on a line when they'd rather be
   forced to be terse.
\end{verse}

Mathematical formulas may also be displayed.  A
displayed formula
is
one-line long; multiline
formulas require special formatting instructions.
   \[  \ip{\Gamma}{\psi'} = x'' + y^{2} + z_{i}^{n}\]
Don't start a paragraph with a displayed equation,
nor make one a paragraph by itself.

\end{document}               % End of document.

\section{Bench Huge Section 1476}
% synthetic bench section 1476
% This is a sample LaTeX input file.  (Version of 12 August 2004.)
%
% A '%' character causes TeX to ignore all remaining text on the line,
% and is used for comments like this one.

\documentclass{article}      % Specifies the document class

                             % The preamble begins here.
\title{An Example Document}  % Declares the document's title.
\author{Leslie Lamport}      % Declares the author's name.
\date{January 21, 1994}      % Deleting this command produces today's date.

\newcommand{\ip}[2]{(#1, #2)}
                             % Defines \ip{arg1}{arg2} to mean
                             % (arg1, arg2).

%\newcommand{\ip}[2]{\langle #1 | #2\rangle}
                             % This is an alternative definition of
                             % \ip that is commented out.

\begin{document}             % End of preamble and beginning of text.

\maketitle                   % Produces the title.

This is an example input file.  Comparing it with
the output it generates can show you how to
produce a simple document of your own.

\section{Ordinary Text}      % Produces section heading.  Lower-level
                             % sections are begun with similar
                             % \subsection and \subsubsection commands.

The ends  of words and sentences are marked
  by   spaces. It  doesn't matter how many
spaces    you type; one is as good as 100.  The
end of   a line counts as a space.

One   or more   blank lines denote the  end
of  a paragraph.

Since any number of consecutive spaces are treated
like a single one, the formatting of the input
file makes no difference to
      \LaTeX,                % The \LaTeX command generates the LaTeX logo.
but it makes a difference to you.  When you use
\LaTeX, making your input file as easy to read
as possible will be a great help as you write
your document and when you change it.  This sample
file shows how you can add comments to your own input
file.

Because printing is different from typewriting,
there are a number of things that you have to do
differently when preparing an input file than if
you were just typing the document directly.
Quotation marks like
       ``this''
have to be handled specially, as do quotes within
quotes:
       ``\,`this'            % \, separates the double and single quote.
        is what I just
        wrote, not  `that'\,''.

Dashes come in three sizes: an
       intra-word
dash, a medium dash for number ranges like
       1--2,
and a punctuation
       dash---like
this.

A sentence-ending space should be larger than the
space between words within a sentence.  You
sometimes have to type special commands in
conjunction with punctuation characters to get
this right, as in the following sentence.
       Gnats, gnus, etc.\ all  % `\ ' makes an inter-word space.
       begin with G\@.         % \@ marks end-of-sentence punctuation.
You should check the spaces after periods when
reading your output to make sure you haven't
forgotten any special cases.  Generating an
ellipsis
       \ldots\               % `\ ' is needed after `\ldots' because TeX
                             % ignores spaces after command names like \ldots
                             % made from \ + letters.
                             %
                             % Note how a `%' character causes TeX to ignore
                             % the end of the input line, so these blank lines
                             % do not start a new paragraph.
                             %
with the right spacing around the periods requires
a special command.

\LaTeX\ interprets some common characters as
commands, so you must type special commands to
generate them.  These characters include the
following:
       \$ \& \% \# \{ and \}.

In printing, text is usually emphasized with an
       \emph{italic}
type style.

\begin{em}
   A long segment of text can also be emphasized
   in this way.  Text within such a segment can be
   given \emph{additional} emphasis.
\end{em}

It is sometimes necessary to prevent \LaTeX\ from
breaking a line where it might otherwise do so.
This may be at a space, as between the ``Mr.''\ and
``Jones'' in
       ``Mr.~Jones'',        % ~ produces an unbreakable interword space.
or within a word---especially when the word is a
symbol like
       \mbox{\emph{itemnum}}
that makes little sense when hyphenated across
lines.

Footnotes\footnote{This is an example of a footnote.}
pose no problem.

\LaTeX\ is good at typesetting mathematical formulas
like
       \( x-3y + z = 7 \)
or
       \( a_{1} > x^{2n} + y^{2n} > x' \)
or
       \( \ip{A}{B} = \sum_{i} a_{i} b_{i} \).
The spaces you type in a formula are
ignored.  Remember that a letter like
       $x$                   % $ ... $  and  \( ... \)  are equivalent
is a formula when it denotes a mathematical
symbol, and it should be typed as one.

\section{Displayed Text}

Text is displayed by indenting it from the left
margin.  Quotations are commonly displayed.  There
are short quotations
\begin{quote}
   This is a short quotation.  It consists of a
   single paragraph of text.  See how it is formatted.
\end{quote}
and longer ones.
\begin{quotation}
   This is a longer quotation.  It consists of two
   paragraphs of text, neither of which are
   particularly interesting.

   This is the second paragraph of the quotation.  It
   is just as dull as the first paragraph.
\end{quotation}
Another frequently-displayed structure is a list.
The following is an example of an \emph{itemized}
list.
\begin{itemize}
   \item This is the first item of an itemized list.
         Each item in the list is marked with a ``tick''.
         You don't have to worry about what kind of tick
         mark is used.

   \item This is the second item of the list.  It
         contains another list nested inside it.  The inner
         list is an \emph{enumerated} list.
         \begin{enumerate}
            \item This is the first item of an enumerated
                  list that is nested within the itemized list.

            \item This is the second item of the inner list.
                  \LaTeX\ allows you to nest lists deeper than
                  you really should.
         \end{enumerate}
         This is the rest of the second item of the outer
         list.  It is no more interesting than any other
         part of the item.
   \item This is the third item of the list.
\end{itemize}
You can even display poetry.
\begin{verse}
   There is an environment
    for verse \\             % The \\ command separates lines
   Whose features some poets % within a stanza.
   will curse.

                             % One or more blank lines separate stanzas.

   For instead of making\\
   Them do \emph{all} line breaking, \\
   It allows them to put too many words on a line when they'd rather be
   forced to be terse.
\end{verse}

Mathematical formulas may also be displayed.  A
displayed formula
is
one-line long; multiline
formulas require special formatting instructions.
   \[  \ip{\Gamma}{\psi'} = x'' + y^{2} + z_{i}^{n}\]
Don't start a paragraph with a displayed equation,
nor make one a paragraph by itself.

\end{document}               % End of document.

\section{Bench Huge Section 1477}
% synthetic bench section 1477
% This is a sample LaTeX input file.  (Version of 12 August 2004.)
%
% A '%' character causes TeX to ignore all remaining text on the line,
% and is used for comments like this one.

\documentclass{article}      % Specifies the document class

                             % The preamble begins here.
\title{An Example Document}  % Declares the document's title.
\author{Leslie Lamport}      % Declares the author's name.
\date{January 21, 1994}      % Deleting this command produces today's date.

\newcommand{\ip}[2]{(#1, #2)}
                             % Defines \ip{arg1}{arg2} to mean
                             % (arg1, arg2).

%\newcommand{\ip}[2]{\langle #1 | #2\rangle}
                             % This is an alternative definition of
                             % \ip that is commented out.

\begin{document}             % End of preamble and beginning of text.

\maketitle                   % Produces the title.

This is an example input file.  Comparing it with
the output it generates can show you how to
produce a simple document of your own.

\section{Ordinary Text}      % Produces section heading.  Lower-level
                             % sections are begun with similar
                             % \subsection and \subsubsection commands.

The ends  of words and sentences are marked
  by   spaces. It  doesn't matter how many
spaces    you type; one is as good as 100.  The
end of   a line counts as a space.

One   or more   blank lines denote the  end
of  a paragraph.

Since any number of consecutive spaces are treated
like a single one, the formatting of the input
file makes no difference to
      \LaTeX,                % The \LaTeX command generates the LaTeX logo.
but it makes a difference to you.  When you use
\LaTeX, making your input file as easy to read
as possible will be a great help as you write
your document and when you change it.  This sample
file shows how you can add comments to your own input
file.

Because printing is different from typewriting,
there are a number of things that you have to do
differently when preparing an input file than if
you were just typing the document directly.
Quotation marks like
       ``this''
have to be handled specially, as do quotes within
quotes:
       ``\,`this'            % \, separates the double and single quote.
        is what I just
        wrote, not  `that'\,''.

Dashes come in three sizes: an
       intra-word
dash, a medium dash for number ranges like
       1--2,
and a punctuation
       dash---like
this.

A sentence-ending space should be larger than the
space between words within a sentence.  You
sometimes have to type special commands in
conjunction with punctuation characters to get
this right, as in the following sentence.
       Gnats, gnus, etc.\ all  % `\ ' makes an inter-word space.
       begin with G\@.         % \@ marks end-of-sentence punctuation.
You should check the spaces after periods when
reading your output to make sure you haven't
forgotten any special cases.  Generating an
ellipsis
       \ldots\               % `\ ' is needed after `\ldots' because TeX
                             % ignores spaces after command names like \ldots
                             % made from \ + letters.
                             %
                             % Note how a `%' character causes TeX to ignore
                             % the end of the input line, so these blank lines
                             % do not start a new paragraph.
                             %
with the right spacing around the periods requires
a special command.

\LaTeX\ interprets some common characters as
commands, so you must type special commands to
generate them.  These characters include the
following:
       \$ \& \% \# \{ and \}.

In printing, text is usually emphasized with an
       \emph{italic}
type style.

\begin{em}
   A long segment of text can also be emphasized
   in this way.  Text within such a segment can be
   given \emph{additional} emphasis.
\end{em}

It is sometimes necessary to prevent \LaTeX\ from
breaking a line where it might otherwise do so.
This may be at a space, as between the ``Mr.''\ and
``Jones'' in
       ``Mr.~Jones'',        % ~ produces an unbreakable interword space.
or within a word---especially when the word is a
symbol like
       \mbox{\emph{itemnum}}
that makes little sense when hyphenated across
lines.

Footnotes\footnote{This is an example of a footnote.}
pose no problem.

\LaTeX\ is good at typesetting mathematical formulas
like
       \( x-3y + z = 7 \)
or
       \( a_{1} > x^{2n} + y^{2n} > x' \)
or
       \( \ip{A}{B} = \sum_{i} a_{i} b_{i} \).
The spaces you type in a formula are
ignored.  Remember that a letter like
       $x$                   % $ ... $  and  \( ... \)  are equivalent
is a formula when it denotes a mathematical
symbol, and it should be typed as one.

\section{Displayed Text}

Text is displayed by indenting it from the left
margin.  Quotations are commonly displayed.  There
are short quotations
\begin{quote}
   This is a short quotation.  It consists of a
   single paragraph of text.  See how it is formatted.
\end{quote}
and longer ones.
\begin{quotation}
   This is a longer quotation.  It consists of two
   paragraphs of text, neither of which are
   particularly interesting.

   This is the second paragraph of the quotation.  It
   is just as dull as the first paragraph.
\end{quotation}
Another frequently-displayed structure is a list.
The following is an example of an \emph{itemized}
list.
\begin{itemize}
   \item This is the first item of an itemized list.
         Each item in the list is marked with a ``tick''.
         You don't have to worry about what kind of tick
         mark is used.

   \item This is the second item of the list.  It
         contains another list nested inside it.  The inner
         list is an \emph{enumerated} list.
         \begin{enumerate}
            \item This is the first item of an enumerated
                  list that is nested within the itemized list.

            \item This is the second item of the inner list.
                  \LaTeX\ allows you to nest lists deeper than
                  you really should.
         \end{enumerate}
         This is the rest of the second item of the outer
         list.  It is no more interesting than any other
         part of the item.
   \item This is the third item of the list.
\end{itemize}
You can even display poetry.
\begin{verse}
   There is an environment
    for verse \\             % The \\ command separates lines
   Whose features some poets % within a stanza.
   will curse.

                             % One or more blank lines separate stanzas.

   For instead of making\\
   Them do \emph{all} line breaking, \\
   It allows them to put too many words on a line when they'd rather be
   forced to be terse.
\end{verse}

Mathematical formulas may also be displayed.  A
displayed formula
is
one-line long; multiline
formulas require special formatting instructions.
   \[  \ip{\Gamma}{\psi'} = x'' + y^{2} + z_{i}^{n}\]
Don't start a paragraph with a displayed equation,
nor make one a paragraph by itself.

\end{document}               % End of document.

\section{Bench Huge Section 1478}
% synthetic bench section 1478
% This is a sample LaTeX input file.  (Version of 12 August 2004.)
%
% A '%' character causes TeX to ignore all remaining text on the line,
% and is used for comments like this one.

\documentclass{article}      % Specifies the document class

                             % The preamble begins here.
\title{An Example Document}  % Declares the document's title.
\author{Leslie Lamport}      % Declares the author's name.
\date{January 21, 1994}      % Deleting this command produces today's date.

\newcommand{\ip}[2]{(#1, #2)}
                             % Defines \ip{arg1}{arg2} to mean
                             % (arg1, arg2).

%\newcommand{\ip}[2]{\langle #1 | #2\rangle}
                             % This is an alternative definition of
                             % \ip that is commented out.

\begin{document}             % End of preamble and beginning of text.

\maketitle                   % Produces the title.

This is an example input file.  Comparing it with
the output it generates can show you how to
produce a simple document of your own.

\section{Ordinary Text}      % Produces section heading.  Lower-level
                             % sections are begun with similar
                             % \subsection and \subsubsection commands.

The ends  of words and sentences are marked
  by   spaces. It  doesn't matter how many
spaces    you type; one is as good as 100.  The
end of   a line counts as a space.

One   or more   blank lines denote the  end
of  a paragraph.

Since any number of consecutive spaces are treated
like a single one, the formatting of the input
file makes no difference to
      \LaTeX,                % The \LaTeX command generates the LaTeX logo.
but it makes a difference to you.  When you use
\LaTeX, making your input file as easy to read
as possible will be a great help as you write
your document and when you change it.  This sample
file shows how you can add comments to your own input
file.

Because printing is different from typewriting,
there are a number of things that you have to do
differently when preparing an input file than if
you were just typing the document directly.
Quotation marks like
       ``this''
have to be handled specially, as do quotes within
quotes:
       ``\,`this'            % \, separates the double and single quote.
        is what I just
        wrote, not  `that'\,''.

Dashes come in three sizes: an
       intra-word
dash, a medium dash for number ranges like
       1--2,
and a punctuation
       dash---like
this.

A sentence-ending space should be larger than the
space between words within a sentence.  You
sometimes have to type special commands in
conjunction with punctuation characters to get
this right, as in the following sentence.
       Gnats, gnus, etc.\ all  % `\ ' makes an inter-word space.
       begin with G\@.         % \@ marks end-of-sentence punctuation.
You should check the spaces after periods when
reading your output to make sure you haven't
forgotten any special cases.  Generating an
ellipsis
       \ldots\               % `\ ' is needed after `\ldots' because TeX
                             % ignores spaces after command names like \ldots
                             % made from \ + letters.
                             %
                             % Note how a `%' character causes TeX to ignore
                             % the end of the input line, so these blank lines
                             % do not start a new paragraph.
                             %
with the right spacing around the periods requires
a special command.

\LaTeX\ interprets some common characters as
commands, so you must type special commands to
generate them.  These characters include the
following:
       \$ \& \% \# \{ and \}.

In printing, text is usually emphasized with an
       \emph{italic}
type style.

\begin{em}
   A long segment of text can also be emphasized
   in this way.  Text within such a segment can be
   given \emph{additional} emphasis.
\end{em}

It is sometimes necessary to prevent \LaTeX\ from
breaking a line where it might otherwise do so.
This may be at a space, as between the ``Mr.''\ and
``Jones'' in
       ``Mr.~Jones'',        % ~ produces an unbreakable interword space.
or within a word---especially when the word is a
symbol like
       \mbox{\emph{itemnum}}
that makes little sense when hyphenated across
lines.

Footnotes\footnote{This is an example of a footnote.}
pose no problem.

\LaTeX\ is good at typesetting mathematical formulas
like
       \( x-3y + z = 7 \)
or
       \( a_{1} > x^{2n} + y^{2n} > x' \)
or
       \( \ip{A}{B} = \sum_{i} a_{i} b_{i} \).
The spaces you type in a formula are
ignored.  Remember that a letter like
       $x$                   % $ ... $  and  \( ... \)  are equivalent
is a formula when it denotes a mathematical
symbol, and it should be typed as one.

\section{Displayed Text}

Text is displayed by indenting it from the left
margin.  Quotations are commonly displayed.  There
are short quotations
\begin{quote}
   This is a short quotation.  It consists of a
   single paragraph of text.  See how it is formatted.
\end{quote}
and longer ones.
\begin{quotation}
   This is a longer quotation.  It consists of two
   paragraphs of text, neither of which are
   particularly interesting.

   This is the second paragraph of the quotation.  It
   is just as dull as the first paragraph.
\end{quotation}
Another frequently-displayed structure is a list.
The following is an example of an \emph{itemized}
list.
\begin{itemize}
   \item This is the first item of an itemized list.
         Each item in the list is marked with a ``tick''.
         You don't have to worry about what kind of tick
         mark is used.

   \item This is the second item of the list.  It
         contains another list nested inside it.  The inner
         list is an \emph{enumerated} list.
         \begin{enumerate}
            \item This is the first item of an enumerated
                  list that is nested within the itemized list.

            \item This is the second item of the inner list.
                  \LaTeX\ allows you to nest lists deeper than
                  you really should.
         \end{enumerate}
         This is the rest of the second item of the outer
         list.  It is no more interesting than any other
         part of the item.
   \item This is the third item of the list.
\end{itemize}
You can even display poetry.
\begin{verse}
   There is an environment
    for verse \\             % The \\ command separates lines
   Whose features some poets % within a stanza.
   will curse.

                             % One or more blank lines separate stanzas.

   For instead of making\\
   Them do \emph{all} line breaking, \\
   It allows them to put too many words on a line when they'd rather be
   forced to be terse.
\end{verse}

Mathematical formulas may also be displayed.  A
displayed formula
is
one-line long; multiline
formulas require special formatting instructions.
   \[  \ip{\Gamma}{\psi'} = x'' + y^{2} + z_{i}^{n}\]
Don't start a paragraph with a displayed equation,
nor make one a paragraph by itself.

\end{document}               % End of document.

\section{Bench Huge Section 1479}
% synthetic bench section 1479
% This is a sample LaTeX input file.  (Version of 12 August 2004.)
%
% A '%' character causes TeX to ignore all remaining text on the line,
% and is used for comments like this one.

\documentclass{article}      % Specifies the document class

                             % The preamble begins here.
\title{An Example Document}  % Declares the document's title.
\author{Leslie Lamport}      % Declares the author's name.
\date{January 21, 1994}      % Deleting this command produces today's date.

\newcommand{\ip}[2]{(#1, #2)}
                             % Defines \ip{arg1}{arg2} to mean
                             % (arg1, arg2).

%\newcommand{\ip}[2]{\langle #1 | #2\rangle}
                             % This is an alternative definition of
                             % \ip that is commented out.

\begin{document}             % End of preamble and beginning of text.

\maketitle                   % Produces the title.

This is an example input file.  Comparing it with
the output it generates can show you how to
produce a simple document of your own.

\section{Ordinary Text}      % Produces section heading.  Lower-level
                             % sections are begun with similar
                             % \subsection and \subsubsection commands.

The ends  of words and sentences are marked
  by   spaces. It  doesn't matter how many
spaces    you type; one is as good as 100.  The
end of   a line counts as a space.

One   or more   blank lines denote the  end
of  a paragraph.

Since any number of consecutive spaces are treated
like a single one, the formatting of the input
file makes no difference to
      \LaTeX,                % The \LaTeX command generates the LaTeX logo.
but it makes a difference to you.  When you use
\LaTeX, making your input file as easy to read
as possible will be a great help as you write
your document and when you change it.  This sample
file shows how you can add comments to your own input
file.

Because printing is different from typewriting,
there are a number of things that you have to do
differently when preparing an input file than if
you were just typing the document directly.
Quotation marks like
       ``this''
have to be handled specially, as do quotes within
quotes:
       ``\,`this'            % \, separates the double and single quote.
        is what I just
        wrote, not  `that'\,''.

Dashes come in three sizes: an
       intra-word
dash, a medium dash for number ranges like
       1--2,
and a punctuation
       dash---like
this.

A sentence-ending space should be larger than the
space between words within a sentence.  You
sometimes have to type special commands in
conjunction with punctuation characters to get
this right, as in the following sentence.
       Gnats, gnus, etc.\ all  % `\ ' makes an inter-word space.
       begin with G\@.         % \@ marks end-of-sentence punctuation.
You should check the spaces after periods when
reading your output to make sure you haven't
forgotten any special cases.  Generating an
ellipsis
       \ldots\               % `\ ' is needed after `\ldots' because TeX
                             % ignores spaces after command names like \ldots
                             % made from \ + letters.
                             %
                             % Note how a `%' character causes TeX to ignore
                             % the end of the input line, so these blank lines
                             % do not start a new paragraph.
                             %
with the right spacing around the periods requires
a special command.

\LaTeX\ interprets some common characters as
commands, so you must type special commands to
generate them.  These characters include the
following:
       \$ \& \% \# \{ and \}.

In printing, text is usually emphasized with an
       \emph{italic}
type style.

\begin{em}
   A long segment of text can also be emphasized
   in this way.  Text within such a segment can be
   given \emph{additional} emphasis.
\end{em}

It is sometimes necessary to prevent \LaTeX\ from
breaking a line where it might otherwise do so.
This may be at a space, as between the ``Mr.''\ and
``Jones'' in
       ``Mr.~Jones'',        % ~ produces an unbreakable interword space.
or within a word---especially when the word is a
symbol like
       \mbox{\emph{itemnum}}
that makes little sense when hyphenated across
lines.

Footnotes\footnote{This is an example of a footnote.}
pose no problem.

\LaTeX\ is good at typesetting mathematical formulas
like
       \( x-3y + z = 7 \)
or
       \( a_{1} > x^{2n} + y^{2n} > x' \)
or
       \( \ip{A}{B} = \sum_{i} a_{i} b_{i} \).
The spaces you type in a formula are
ignored.  Remember that a letter like
       $x$                   % $ ... $  and  \( ... \)  are equivalent
is a formula when it denotes a mathematical
symbol, and it should be typed as one.

\section{Displayed Text}

Text is displayed by indenting it from the left
margin.  Quotations are commonly displayed.  There
are short quotations
\begin{quote}
   This is a short quotation.  It consists of a
   single paragraph of text.  See how it is formatted.
\end{quote}
and longer ones.
\begin{quotation}
   This is a longer quotation.  It consists of two
   paragraphs of text, neither of which are
   particularly interesting.

   This is the second paragraph of the quotation.  It
   is just as dull as the first paragraph.
\end{quotation}
Another frequently-displayed structure is a list.
The following is an example of an \emph{itemized}
list.
\begin{itemize}
   \item This is the first item of an itemized list.
         Each item in the list is marked with a ``tick''.
         You don't have to worry about what kind of tick
         mark is used.

   \item This is the second item of the list.  It
         contains another list nested inside it.  The inner
         list is an \emph{enumerated} list.
         \begin{enumerate}
            \item This is the first item of an enumerated
                  list that is nested within the itemized list.

            \item This is the second item of the inner list.
                  \LaTeX\ allows you to nest lists deeper than
                  you really should.
         \end{enumerate}
         This is the rest of the second item of the outer
         list.  It is no more interesting than any other
         part of the item.
   \item This is the third item of the list.
\end{itemize}
You can even display poetry.
\begin{verse}
   There is an environment
    for verse \\             % The \\ command separates lines
   Whose features some poets % within a stanza.
   will curse.

                             % One or more blank lines separate stanzas.

   For instead of making\\
   Them do \emph{all} line breaking, \\
   It allows them to put too many words on a line when they'd rather be
   forced to be terse.
\end{verse}

Mathematical formulas may also be displayed.  A
displayed formula
is
one-line long; multiline
formulas require special formatting instructions.
   \[  \ip{\Gamma}{\psi'} = x'' + y^{2} + z_{i}^{n}\]
Don't start a paragraph with a displayed equation,
nor make one a paragraph by itself.

\end{document}               % End of document.

\section{Bench Huge Section 1480}
% synthetic bench section 1480
% This is a sample LaTeX input file.  (Version of 12 August 2004.)
%
% A '%' character causes TeX to ignore all remaining text on the line,
% and is used for comments like this one.

\documentclass{article}      % Specifies the document class

                             % The preamble begins here.
\title{An Example Document}  % Declares the document's title.
\author{Leslie Lamport}      % Declares the author's name.
\date{January 21, 1994}      % Deleting this command produces today's date.

\newcommand{\ip}[2]{(#1, #2)}
                             % Defines \ip{arg1}{arg2} to mean
                             % (arg1, arg2).

%\newcommand{\ip}[2]{\langle #1 | #2\rangle}
                             % This is an alternative definition of
                             % \ip that is commented out.

\begin{document}             % End of preamble and beginning of text.

\maketitle                   % Produces the title.

This is an example input file.  Comparing it with
the output it generates can show you how to
produce a simple document of your own.

\section{Ordinary Text}      % Produces section heading.  Lower-level
                             % sections are begun with similar
                             % \subsection and \subsubsection commands.

The ends  of words and sentences are marked
  by   spaces. It  doesn't matter how many
spaces    you type; one is as good as 100.  The
end of   a line counts as a space.

One   or more   blank lines denote the  end
of  a paragraph.

Since any number of consecutive spaces are treated
like a single one, the formatting of the input
file makes no difference to
      \LaTeX,                % The \LaTeX command generates the LaTeX logo.
but it makes a difference to you.  When you use
\LaTeX, making your input file as easy to read
as possible will be a great help as you write
your document and when you change it.  This sample
file shows how you can add comments to your own input
file.

Because printing is different from typewriting,
there are a number of things that you have to do
differently when preparing an input file than if
you were just typing the document directly.
Quotation marks like
       ``this''
have to be handled specially, as do quotes within
quotes:
       ``\,`this'            % \, separates the double and single quote.
        is what I just
        wrote, not  `that'\,''.

Dashes come in three sizes: an
       intra-word
dash, a medium dash for number ranges like
       1--2,
and a punctuation
       dash---like
this.

A sentence-ending space should be larger than the
space between words within a sentence.  You
sometimes have to type special commands in
conjunction with punctuation characters to get
this right, as in the following sentence.
       Gnats, gnus, etc.\ all  % `\ ' makes an inter-word space.
       begin with G\@.         % \@ marks end-of-sentence punctuation.
You should check the spaces after periods when
reading your output to make sure you haven't
forgotten any special cases.  Generating an
ellipsis
       \ldots\               % `\ ' is needed after `\ldots' because TeX
                             % ignores spaces after command names like \ldots
                             % made from \ + letters.
                             %
                             % Note how a `%' character causes TeX to ignore
                             % the end of the input line, so these blank lines
                             % do not start a new paragraph.
                             %
with the right spacing around the periods requires
a special command.

\LaTeX\ interprets some common characters as
commands, so you must type special commands to
generate them.  These characters include the
following:
       \$ \& \% \# \{ and \}.

In printing, text is usually emphasized with an
       \emph{italic}
type style.

\begin{em}
   A long segment of text can also be emphasized
   in this way.  Text within such a segment can be
   given \emph{additional} emphasis.
\end{em}

It is sometimes necessary to prevent \LaTeX\ from
breaking a line where it might otherwise do so.
This may be at a space, as between the ``Mr.''\ and
``Jones'' in
       ``Mr.~Jones'',        % ~ produces an unbreakable interword space.
or within a word---especially when the word is a
symbol like
       \mbox{\emph{itemnum}}
that makes little sense when hyphenated across
lines.

Footnotes\footnote{This is an example of a footnote.}
pose no problem.

\LaTeX\ is good at typesetting mathematical formulas
like
       \( x-3y + z = 7 \)
or
       \( a_{1} > x^{2n} + y^{2n} > x' \)
or
       \( \ip{A}{B} = \sum_{i} a_{i} b_{i} \).
The spaces you type in a formula are
ignored.  Remember that a letter like
       $x$                   % $ ... $  and  \( ... \)  are equivalent
is a formula when it denotes a mathematical
symbol, and it should be typed as one.

\section{Displayed Text}

Text is displayed by indenting it from the left
margin.  Quotations are commonly displayed.  There
are short quotations
\begin{quote}
   This is a short quotation.  It consists of a
   single paragraph of text.  See how it is formatted.
\end{quote}
and longer ones.
\begin{quotation}
   This is a longer quotation.  It consists of two
   paragraphs of text, neither of which are
   particularly interesting.

   This is the second paragraph of the quotation.  It
   is just as dull as the first paragraph.
\end{quotation}
Another frequently-displayed structure is a list.
The following is an example of an \emph{itemized}
list.
\begin{itemize}
   \item This is the first item of an itemized list.
         Each item in the list is marked with a ``tick''.
         You don't have to worry about what kind of tick
         mark is used.

   \item This is the second item of the list.  It
         contains another list nested inside it.  The inner
         list is an \emph{enumerated} list.
         \begin{enumerate}
            \item This is the first item of an enumerated
                  list that is nested within the itemized list.

            \item This is the second item of the inner list.
                  \LaTeX\ allows you to nest lists deeper than
                  you really should.
         \end{enumerate}
         This is the rest of the second item of the outer
         list.  It is no more interesting than any other
         part of the item.
   \item This is the third item of the list.
\end{itemize}
You can even display poetry.
\begin{verse}
   There is an environment
    for verse \\             % The \\ command separates lines
   Whose features some poets % within a stanza.
   will curse.

                             % One or more blank lines separate stanzas.

   For instead of making\\
   Them do \emph{all} line breaking, \\
   It allows them to put too many words on a line when they'd rather be
   forced to be terse.
\end{verse}

Mathematical formulas may also be displayed.  A
displayed formula
is
one-line long; multiline
formulas require special formatting instructions.
   \[  \ip{\Gamma}{\psi'} = x'' + y^{2} + z_{i}^{n}\]
Don't start a paragraph with a displayed equation,
nor make one a paragraph by itself.

\end{document}               % End of document.

\section{Bench Huge Section 1481}
% synthetic bench section 1481
% This is a sample LaTeX input file.  (Version of 12 August 2004.)
%
% A '%' character causes TeX to ignore all remaining text on the line,
% and is used for comments like this one.

\documentclass{article}      % Specifies the document class

                             % The preamble begins here.
\title{An Example Document}  % Declares the document's title.
\author{Leslie Lamport}      % Declares the author's name.
\date{January 21, 1994}      % Deleting this command produces today's date.

\newcommand{\ip}[2]{(#1, #2)}
                             % Defines \ip{arg1}{arg2} to mean
                             % (arg1, arg2).

%\newcommand{\ip}[2]{\langle #1 | #2\rangle}
                             % This is an alternative definition of
                             % \ip that is commented out.

\begin{document}             % End of preamble and beginning of text.

\maketitle                   % Produces the title.

This is an example input file.  Comparing it with
the output it generates can show you how to
produce a simple document of your own.

\section{Ordinary Text}      % Produces section heading.  Lower-level
                             % sections are begun with similar
                             % \subsection and \subsubsection commands.

The ends  of words and sentences are marked
  by   spaces. It  doesn't matter how many
spaces    you type; one is as good as 100.  The
end of   a line counts as a space.

One   or more   blank lines denote the  end
of  a paragraph.

Since any number of consecutive spaces are treated
like a single one, the formatting of the input
file makes no difference to
      \LaTeX,                % The \LaTeX command generates the LaTeX logo.
but it makes a difference to you.  When you use
\LaTeX, making your input file as easy to read
as possible will be a great help as you write
your document and when you change it.  This sample
file shows how you can add comments to your own input
file.

Because printing is different from typewriting,
there are a number of things that you have to do
differently when preparing an input file than if
you were just typing the document directly.
Quotation marks like
       ``this''
have to be handled specially, as do quotes within
quotes:
       ``\,`this'            % \, separates the double and single quote.
        is what I just
        wrote, not  `that'\,''.

Dashes come in three sizes: an
       intra-word
dash, a medium dash for number ranges like
       1--2,
and a punctuation
       dash---like
this.

A sentence-ending space should be larger than the
space between words within a sentence.  You
sometimes have to type special commands in
conjunction with punctuation characters to get
this right, as in the following sentence.
       Gnats, gnus, etc.\ all  % `\ ' makes an inter-word space.
       begin with G\@.         % \@ marks end-of-sentence punctuation.
You should check the spaces after periods when
reading your output to make sure you haven't
forgotten any special cases.  Generating an
ellipsis
       \ldots\               % `\ ' is needed after `\ldots' because TeX
                             % ignores spaces after command names like \ldots
                             % made from \ + letters.
                             %
                             % Note how a `%' character causes TeX to ignore
                             % the end of the input line, so these blank lines
                             % do not start a new paragraph.
                             %
with the right spacing around the periods requires
a special command.

\LaTeX\ interprets some common characters as
commands, so you must type special commands to
generate them.  These characters include the
following:
       \$ \& \% \# \{ and \}.

In printing, text is usually emphasized with an
       \emph{italic}
type style.

\begin{em}
   A long segment of text can also be emphasized
   in this way.  Text within such a segment can be
   given \emph{additional} emphasis.
\end{em}

It is sometimes necessary to prevent \LaTeX\ from
breaking a line where it might otherwise do so.
This may be at a space, as between the ``Mr.''\ and
``Jones'' in
       ``Mr.~Jones'',        % ~ produces an unbreakable interword space.
or within a word---especially when the word is a
symbol like
       \mbox{\emph{itemnum}}
that makes little sense when hyphenated across
lines.

Footnotes\footnote{This is an example of a footnote.}
pose no problem.

\LaTeX\ is good at typesetting mathematical formulas
like
       \( x-3y + z = 7 \)
or
       \( a_{1} > x^{2n} + y^{2n} > x' \)
or
       \( \ip{A}{B} = \sum_{i} a_{i} b_{i} \).
The spaces you type in a formula are
ignored.  Remember that a letter like
       $x$                   % $ ... $  and  \( ... \)  are equivalent
is a formula when it denotes a mathematical
symbol, and it should be typed as one.

\section{Displayed Text}

Text is displayed by indenting it from the left
margin.  Quotations are commonly displayed.  There
are short quotations
\begin{quote}
   This is a short quotation.  It consists of a
   single paragraph of text.  See how it is formatted.
\end{quote}
and longer ones.
\begin{quotation}
   This is a longer quotation.  It consists of two
   paragraphs of text, neither of which are
   particularly interesting.

   This is the second paragraph of the quotation.  It
   is just as dull as the first paragraph.
\end{quotation}
Another frequently-displayed structure is a list.
The following is an example of an \emph{itemized}
list.
\begin{itemize}
   \item This is the first item of an itemized list.
         Each item in the list is marked with a ``tick''.
         You don't have to worry about what kind of tick
         mark is used.

   \item This is the second item of the list.  It
         contains another list nested inside it.  The inner
         list is an \emph{enumerated} list.
         \begin{enumerate}
            \item This is the first item of an enumerated
                  list that is nested within the itemized list.

            \item This is the second item of the inner list.
                  \LaTeX\ allows you to nest lists deeper than
                  you really should.
         \end{enumerate}
         This is the rest of the second item of the outer
         list.  It is no more interesting than any other
         part of the item.
   \item This is the third item of the list.
\end{itemize}
You can even display poetry.
\begin{verse}
   There is an environment
    for verse \\             % The \\ command separates lines
   Whose features some poets % within a stanza.
   will curse.

                             % One or more blank lines separate stanzas.

   For instead of making\\
   Them do \emph{all} line breaking, \\
   It allows them to put too many words on a line when they'd rather be
   forced to be terse.
\end{verse}

Mathematical formulas may also be displayed.  A
displayed formula
is
one-line long; multiline
formulas require special formatting instructions.
   \[  \ip{\Gamma}{\psi'} = x'' + y^{2} + z_{i}^{n}\]
Don't start a paragraph with a displayed equation,
nor make one a paragraph by itself.

\end{document}               % End of document.

\section{Bench Huge Section 1482}
% synthetic bench section 1482
% This is a sample LaTeX input file.  (Version of 12 August 2004.)
%
% A '%' character causes TeX to ignore all remaining text on the line,
% and is used for comments like this one.

\documentclass{article}      % Specifies the document class

                             % The preamble begins here.
\title{An Example Document}  % Declares the document's title.
\author{Leslie Lamport}      % Declares the author's name.
\date{January 21, 1994}      % Deleting this command produces today's date.

\newcommand{\ip}[2]{(#1, #2)}
                             % Defines \ip{arg1}{arg2} to mean
                             % (arg1, arg2).

%\newcommand{\ip}[2]{\langle #1 | #2\rangle}
                             % This is an alternative definition of
                             % \ip that is commented out.

\begin{document}             % End of preamble and beginning of text.

\maketitle                   % Produces the title.

This is an example input file.  Comparing it with
the output it generates can show you how to
produce a simple document of your own.

\section{Ordinary Text}      % Produces section heading.  Lower-level
                             % sections are begun with similar
                             % \subsection and \subsubsection commands.

The ends  of words and sentences are marked
  by   spaces. It  doesn't matter how many
spaces    you type; one is as good as 100.  The
end of   a line counts as a space.

One   or more   blank lines denote the  end
of  a paragraph.

Since any number of consecutive spaces are treated
like a single one, the formatting of the input
file makes no difference to
      \LaTeX,                % The \LaTeX command generates the LaTeX logo.
but it makes a difference to you.  When you use
\LaTeX, making your input file as easy to read
as possible will be a great help as you write
your document and when you change it.  This sample
file shows how you can add comments to your own input
file.

Because printing is different from typewriting,
there are a number of things that you have to do
differently when preparing an input file than if
you were just typing the document directly.
Quotation marks like
       ``this''
have to be handled specially, as do quotes within
quotes:
       ``\,`this'            % \, separates the double and single quote.
        is what I just
        wrote, not  `that'\,''.

Dashes come in three sizes: an
       intra-word
dash, a medium dash for number ranges like
       1--2,
and a punctuation
       dash---like
this.

A sentence-ending space should be larger than the
space between words within a sentence.  You
sometimes have to type special commands in
conjunction with punctuation characters to get
this right, as in the following sentence.
       Gnats, gnus, etc.\ all  % `\ ' makes an inter-word space.
       begin with G\@.         % \@ marks end-of-sentence punctuation.
You should check the spaces after periods when
reading your output to make sure you haven't
forgotten any special cases.  Generating an
ellipsis
       \ldots\               % `\ ' is needed after `\ldots' because TeX
                             % ignores spaces after command names like \ldots
                             % made from \ + letters.
                             %
                             % Note how a `%' character causes TeX to ignore
                             % the end of the input line, so these blank lines
                             % do not start a new paragraph.
                             %
with the right spacing around the periods requires
a special command.

\LaTeX\ interprets some common characters as
commands, so you must type special commands to
generate them.  These characters include the
following:
       \$ \& \% \# \{ and \}.

In printing, text is usually emphasized with an
       \emph{italic}
type style.

\begin{em}
   A long segment of text can also be emphasized
   in this way.  Text within such a segment can be
   given \emph{additional} emphasis.
\end{em}

It is sometimes necessary to prevent \LaTeX\ from
breaking a line where it might otherwise do so.
This may be at a space, as between the ``Mr.''\ and
``Jones'' in
       ``Mr.~Jones'',        % ~ produces an unbreakable interword space.
or within a word---especially when the word is a
symbol like
       \mbox{\emph{itemnum}}
that makes little sense when hyphenated across
lines.

Footnotes\footnote{This is an example of a footnote.}
pose no problem.

\LaTeX\ is good at typesetting mathematical formulas
like
       \( x-3y + z = 7 \)
or
       \( a_{1} > x^{2n} + y^{2n} > x' \)
or
       \( \ip{A}{B} = \sum_{i} a_{i} b_{i} \).
The spaces you type in a formula are
ignored.  Remember that a letter like
       $x$                   % $ ... $  and  \( ... \)  are equivalent
is a formula when it denotes a mathematical
symbol, and it should be typed as one.

\section{Displayed Text}

Text is displayed by indenting it from the left
margin.  Quotations are commonly displayed.  There
are short quotations
\begin{quote}
   This is a short quotation.  It consists of a
   single paragraph of text.  See how it is formatted.
\end{quote}
and longer ones.
\begin{quotation}
   This is a longer quotation.  It consists of two
   paragraphs of text, neither of which are
   particularly interesting.

   This is the second paragraph of the quotation.  It
   is just as dull as the first paragraph.
\end{quotation}
Another frequently-displayed structure is a list.
The following is an example of an \emph{itemized}
list.
\begin{itemize}
   \item This is the first item of an itemized list.
         Each item in the list is marked with a ``tick''.
         You don't have to worry about what kind of tick
         mark is used.

   \item This is the second item of the list.  It
         contains another list nested inside it.  The inner
         list is an \emph{enumerated} list.
         \begin{enumerate}
            \item This is the first item of an enumerated
                  list that is nested within the itemized list.

            \item This is the second item of the inner list.
                  \LaTeX\ allows you to nest lists deeper than
                  you really should.
         \end{enumerate}
         This is the rest of the second item of the outer
         list.  It is no more interesting than any other
         part of the item.
   \item This is the third item of the list.
\end{itemize}
You can even display poetry.
\begin{verse}
   There is an environment
    for verse \\             % The \\ command separates lines
   Whose features some poets % within a stanza.
   will curse.

                             % One or more blank lines separate stanzas.

   For instead of making\\
   Them do \emph{all} line breaking, \\
   It allows them to put too many words on a line when they'd rather be
   forced to be terse.
\end{verse}

Mathematical formulas may also be displayed.  A
displayed formula
is
one-line long; multiline
formulas require special formatting instructions.
   \[  \ip{\Gamma}{\psi'} = x'' + y^{2} + z_{i}^{n}\]
Don't start a paragraph with a displayed equation,
nor make one a paragraph by itself.

\end{document}               % End of document.

\section{Bench Huge Section 1483}
% synthetic bench section 1483
% This is a sample LaTeX input file.  (Version of 12 August 2004.)
%
% A '%' character causes TeX to ignore all remaining text on the line,
% and is used for comments like this one.

\documentclass{article}      % Specifies the document class

                             % The preamble begins here.
\title{An Example Document}  % Declares the document's title.
\author{Leslie Lamport}      % Declares the author's name.
\date{January 21, 1994}      % Deleting this command produces today's date.

\newcommand{\ip}[2]{(#1, #2)}
                             % Defines \ip{arg1}{arg2} to mean
                             % (arg1, arg2).

%\newcommand{\ip}[2]{\langle #1 | #2\rangle}
                             % This is an alternative definition of
                             % \ip that is commented out.

\begin{document}             % End of preamble and beginning of text.

\maketitle                   % Produces the title.

This is an example input file.  Comparing it with
the output it generates can show you how to
produce a simple document of your own.

\section{Ordinary Text}      % Produces section heading.  Lower-level
                             % sections are begun with similar
                             % \subsection and \subsubsection commands.

The ends  of words and sentences are marked
  by   spaces. It  doesn't matter how many
spaces    you type; one is as good as 100.  The
end of   a line counts as a space.

One   or more   blank lines denote the  end
of  a paragraph.

Since any number of consecutive spaces are treated
like a single one, the formatting of the input
file makes no difference to
      \LaTeX,                % The \LaTeX command generates the LaTeX logo.
but it makes a difference to you.  When you use
\LaTeX, making your input file as easy to read
as possible will be a great help as you write
your document and when you change it.  This sample
file shows how you can add comments to your own input
file.

Because printing is different from typewriting,
there are a number of things that you have to do
differently when preparing an input file than if
you were just typing the document directly.
Quotation marks like
       ``this''
have to be handled specially, as do quotes within
quotes:
       ``\,`this'            % \, separates the double and single quote.
        is what I just
        wrote, not  `that'\,''.

Dashes come in three sizes: an
       intra-word
dash, a medium dash for number ranges like
       1--2,
and a punctuation
       dash---like
this.

A sentence-ending space should be larger than the
space between words within a sentence.  You
sometimes have to type special commands in
conjunction with punctuation characters to get
this right, as in the following sentence.
       Gnats, gnus, etc.\ all  % `\ ' makes an inter-word space.
       begin with G\@.         % \@ marks end-of-sentence punctuation.
You should check the spaces after periods when
reading your output to make sure you haven't
forgotten any special cases.  Generating an
ellipsis
       \ldots\               % `\ ' is needed after `\ldots' because TeX
                             % ignores spaces after command names like \ldots
                             % made from \ + letters.
                             %
                             % Note how a `%' character causes TeX to ignore
                             % the end of the input line, so these blank lines
                             % do not start a new paragraph.
                             %
with the right spacing around the periods requires
a special command.

\LaTeX\ interprets some common characters as
commands, so you must type special commands to
generate them.  These characters include the
following:
       \$ \& \% \# \{ and \}.

In printing, text is usually emphasized with an
       \emph{italic}
type style.

\begin{em}
   A long segment of text can also be emphasized
   in this way.  Text within such a segment can be
   given \emph{additional} emphasis.
\end{em}

It is sometimes necessary to prevent \LaTeX\ from
breaking a line where it might otherwise do so.
This may be at a space, as between the ``Mr.''\ and
``Jones'' in
       ``Mr.~Jones'',        % ~ produces an unbreakable interword space.
or within a word---especially when the word is a
symbol like
       \mbox{\emph{itemnum}}
that makes little sense when hyphenated across
lines.

Footnotes\footnote{This is an example of a footnote.}
pose no problem.

\LaTeX\ is good at typesetting mathematical formulas
like
       \( x-3y + z = 7 \)
or
       \( a_{1} > x^{2n} + y^{2n} > x' \)
or
       \( \ip{A}{B} = \sum_{i} a_{i} b_{i} \).
The spaces you type in a formula are
ignored.  Remember that a letter like
       $x$                   % $ ... $  and  \( ... \)  are equivalent
is a formula when it denotes a mathematical
symbol, and it should be typed as one.

\section{Displayed Text}

Text is displayed by indenting it from the left
margin.  Quotations are commonly displayed.  There
are short quotations
\begin{quote}
   This is a short quotation.  It consists of a
   single paragraph of text.  See how it is formatted.
\end{quote}
and longer ones.
\begin{quotation}
   This is a longer quotation.  It consists of two
   paragraphs of text, neither of which are
   particularly interesting.

   This is the second paragraph of the quotation.  It
   is just as dull as the first paragraph.
\end{quotation}
Another frequently-displayed structure is a list.
The following is an example of an \emph{itemized}
list.
\begin{itemize}
   \item This is the first item of an itemized list.
         Each item in the list is marked with a ``tick''.
         You don't have to worry about what kind of tick
         mark is used.

   \item This is the second item of the list.  It
         contains another list nested inside it.  The inner
         list is an \emph{enumerated} list.
         \begin{enumerate}
            \item This is the first item of an enumerated
                  list that is nested within the itemized list.

            \item This is the second item of the inner list.
                  \LaTeX\ allows you to nest lists deeper than
                  you really should.
         \end{enumerate}
         This is the rest of the second item of the outer
         list.  It is no more interesting than any other
         part of the item.
   \item This is the third item of the list.
\end{itemize}
You can even display poetry.
\begin{verse}
   There is an environment
    for verse \\             % The \\ command separates lines
   Whose features some poets % within a stanza.
   will curse.

                             % One or more blank lines separate stanzas.

   For instead of making\\
   Them do \emph{all} line breaking, \\
   It allows them to put too many words on a line when they'd rather be
   forced to be terse.
\end{verse}

Mathematical formulas may also be displayed.  A
displayed formula
is
one-line long; multiline
formulas require special formatting instructions.
   \[  \ip{\Gamma}{\psi'} = x'' + y^{2} + z_{i}^{n}\]
Don't start a paragraph with a displayed equation,
nor make one a paragraph by itself.

\end{document}               % End of document.

\section{Bench Huge Section 1484}
% synthetic bench section 1484
% This is a sample LaTeX input file.  (Version of 12 August 2004.)
%
% A '%' character causes TeX to ignore all remaining text on the line,
% and is used for comments like this one.

\documentclass{article}      % Specifies the document class

                             % The preamble begins here.
\title{An Example Document}  % Declares the document's title.
\author{Leslie Lamport}      % Declares the author's name.
\date{January 21, 1994}      % Deleting this command produces today's date.

\newcommand{\ip}[2]{(#1, #2)}
                             % Defines \ip{arg1}{arg2} to mean
                             % (arg1, arg2).

%\newcommand{\ip}[2]{\langle #1 | #2\rangle}
                             % This is an alternative definition of
                             % \ip that is commented out.

\begin{document}             % End of preamble and beginning of text.

\maketitle                   % Produces the title.

This is an example input file.  Comparing it with
the output it generates can show you how to
produce a simple document of your own.

\section{Ordinary Text}      % Produces section heading.  Lower-level
                             % sections are begun with similar
                             % \subsection and \subsubsection commands.

The ends  of words and sentences are marked
  by   spaces. It  doesn't matter how many
spaces    you type; one is as good as 100.  The
end of   a line counts as a space.

One   or more   blank lines denote the  end
of  a paragraph.

Since any number of consecutive spaces are treated
like a single one, the formatting of the input
file makes no difference to
      \LaTeX,                % The \LaTeX command generates the LaTeX logo.
but it makes a difference to you.  When you use
\LaTeX, making your input file as easy to read
as possible will be a great help as you write
your document and when you change it.  This sample
file shows how you can add comments to your own input
file.

Because printing is different from typewriting,
there are a number of things that you have to do
differently when preparing an input file than if
you were just typing the document directly.
Quotation marks like
       ``this''
have to be handled specially, as do quotes within
quotes:
       ``\,`this'            % \, separates the double and single quote.
        is what I just
        wrote, not  `that'\,''.

Dashes come in three sizes: an
       intra-word
dash, a medium dash for number ranges like
       1--2,
and a punctuation
       dash---like
this.

A sentence-ending space should be larger than the
space between words within a sentence.  You
sometimes have to type special commands in
conjunction with punctuation characters to get
this right, as in the following sentence.
       Gnats, gnus, etc.\ all  % `\ ' makes an inter-word space.
       begin with G\@.         % \@ marks end-of-sentence punctuation.
You should check the spaces after periods when
reading your output to make sure you haven't
forgotten any special cases.  Generating an
ellipsis
       \ldots\               % `\ ' is needed after `\ldots' because TeX
                             % ignores spaces after command names like \ldots
                             % made from \ + letters.
                             %
                             % Note how a `%' character causes TeX to ignore
                             % the end of the input line, so these blank lines
                             % do not start a new paragraph.
                             %
with the right spacing around the periods requires
a special command.

\LaTeX\ interprets some common characters as
commands, so you must type special commands to
generate them.  These characters include the
following:
       \$ \& \% \# \{ and \}.

In printing, text is usually emphasized with an
       \emph{italic}
type style.

\begin{em}
   A long segment of text can also be emphasized
   in this way.  Text within such a segment can be
   given \emph{additional} emphasis.
\end{em}

It is sometimes necessary to prevent \LaTeX\ from
breaking a line where it might otherwise do so.
This may be at a space, as between the ``Mr.''\ and
``Jones'' in
       ``Mr.~Jones'',        % ~ produces an unbreakable interword space.
or within a word---especially when the word is a
symbol like
       \mbox{\emph{itemnum}}
that makes little sense when hyphenated across
lines.

Footnotes\footnote{This is an example of a footnote.}
pose no problem.

\LaTeX\ is good at typesetting mathematical formulas
like
       \( x-3y + z = 7 \)
or
       \( a_{1} > x^{2n} + y^{2n} > x' \)
or
       \( \ip{A}{B} = \sum_{i} a_{i} b_{i} \).
The spaces you type in a formula are
ignored.  Remember that a letter like
       $x$                   % $ ... $  and  \( ... \)  are equivalent
is a formula when it denotes a mathematical
symbol, and it should be typed as one.

\section{Displayed Text}

Text is displayed by indenting it from the left
margin.  Quotations are commonly displayed.  There
are short quotations
\begin{quote}
   This is a short quotation.  It consists of a
   single paragraph of text.  See how it is formatted.
\end{quote}
and longer ones.
\begin{quotation}
   This is a longer quotation.  It consists of two
   paragraphs of text, neither of which are
   particularly interesting.

   This is the second paragraph of the quotation.  It
   is just as dull as the first paragraph.
\end{quotation}
Another frequently-displayed structure is a list.
The following is an example of an \emph{itemized}
list.
\begin{itemize}
   \item This is the first item of an itemized list.
         Each item in the list is marked with a ``tick''.
         You don't have to worry about what kind of tick
         mark is used.

   \item This is the second item of the list.  It
         contains another list nested inside it.  The inner
         list is an \emph{enumerated} list.
         \begin{enumerate}
            \item This is the first item of an enumerated
                  list that is nested within the itemized list.

            \item This is the second item of the inner list.
                  \LaTeX\ allows you to nest lists deeper than
                  you really should.
         \end{enumerate}
         This is the rest of the second item of the outer
         list.  It is no more interesting than any other
         part of the item.
   \item This is the third item of the list.
\end{itemize}
You can even display poetry.
\begin{verse}
   There is an environment
    for verse \\             % The \\ command separates lines
   Whose features some poets % within a stanza.
   will curse.

                             % One or more blank lines separate stanzas.

   For instead of making\\
   Them do \emph{all} line breaking, \\
   It allows them to put too many words on a line when they'd rather be
   forced to be terse.
\end{verse}

Mathematical formulas may also be displayed.  A
displayed formula
is
one-line long; multiline
formulas require special formatting instructions.
   \[  \ip{\Gamma}{\psi'} = x'' + y^{2} + z_{i}^{n}\]
Don't start a paragraph with a displayed equation,
nor make one a paragraph by itself.

\end{document}               % End of document.

\section{Bench Huge Section 1485}
% synthetic bench section 1485
% This is a sample LaTeX input file.  (Version of 12 August 2004.)
%
% A '%' character causes TeX to ignore all remaining text on the line,
% and is used for comments like this one.

\documentclass{article}      % Specifies the document class

                             % The preamble begins here.
\title{An Example Document}  % Declares the document's title.
\author{Leslie Lamport}      % Declares the author's name.
\date{January 21, 1994}      % Deleting this command produces today's date.

\newcommand{\ip}[2]{(#1, #2)}
                             % Defines \ip{arg1}{arg2} to mean
                             % (arg1, arg2).

%\newcommand{\ip}[2]{\langle #1 | #2\rangle}
                             % This is an alternative definition of
                             % \ip that is commented out.

\begin{document}             % End of preamble and beginning of text.

\maketitle                   % Produces the title.

This is an example input file.  Comparing it with
the output it generates can show you how to
produce a simple document of your own.

\section{Ordinary Text}      % Produces section heading.  Lower-level
                             % sections are begun with similar
                             % \subsection and \subsubsection commands.

The ends  of words and sentences are marked
  by   spaces. It  doesn't matter how many
spaces    you type; one is as good as 100.  The
end of   a line counts as a space.

One   or more   blank lines denote the  end
of  a paragraph.

Since any number of consecutive spaces are treated
like a single one, the formatting of the input
file makes no difference to
      \LaTeX,                % The \LaTeX command generates the LaTeX logo.
but it makes a difference to you.  When you use
\LaTeX, making your input file as easy to read
as possible will be a great help as you write
your document and when you change it.  This sample
file shows how you can add comments to your own input
file.

Because printing is different from typewriting,
there are a number of things that you have to do
differently when preparing an input file than if
you were just typing the document directly.
Quotation marks like
       ``this''
have to be handled specially, as do quotes within
quotes:
       ``\,`this'            % \, separates the double and single quote.
        is what I just
        wrote, not  `that'\,''.

Dashes come in three sizes: an
       intra-word
dash, a medium dash for number ranges like
       1--2,
and a punctuation
       dash---like
this.

A sentence-ending space should be larger than the
space between words within a sentence.  You
sometimes have to type special commands in
conjunction with punctuation characters to get
this right, as in the following sentence.
       Gnats, gnus, etc.\ all  % `\ ' makes an inter-word space.
       begin with G\@.         % \@ marks end-of-sentence punctuation.
You should check the spaces after periods when
reading your output to make sure you haven't
forgotten any special cases.  Generating an
ellipsis
       \ldots\               % `\ ' is needed after `\ldots' because TeX
                             % ignores spaces after command names like \ldots
                             % made from \ + letters.
                             %
                             % Note how a `%' character causes TeX to ignore
                             % the end of the input line, so these blank lines
                             % do not start a new paragraph.
                             %
with the right spacing around the periods requires
a special command.

\LaTeX\ interprets some common characters as
commands, so you must type special commands to
generate them.  These characters include the
following:
       \$ \& \% \# \{ and \}.

In printing, text is usually emphasized with an
       \emph{italic}
type style.

\begin{em}
   A long segment of text can also be emphasized
   in this way.  Text within such a segment can be
   given \emph{additional} emphasis.
\end{em}

It is sometimes necessary to prevent \LaTeX\ from
breaking a line where it might otherwise do so.
This may be at a space, as between the ``Mr.''\ and
``Jones'' in
       ``Mr.~Jones'',        % ~ produces an unbreakable interword space.
or within a word---especially when the word is a
symbol like
       \mbox{\emph{itemnum}}
that makes little sense when hyphenated across
lines.

Footnotes\footnote{This is an example of a footnote.}
pose no problem.

\LaTeX\ is good at typesetting mathematical formulas
like
       \( x-3y + z = 7 \)
or
       \( a_{1} > x^{2n} + y^{2n} > x' \)
or
       \( \ip{A}{B} = \sum_{i} a_{i} b_{i} \).
The spaces you type in a formula are
ignored.  Remember that a letter like
       $x$                   % $ ... $  and  \( ... \)  are equivalent
is a formula when it denotes a mathematical
symbol, and it should be typed as one.

\section{Displayed Text}

Text is displayed by indenting it from the left
margin.  Quotations are commonly displayed.  There
are short quotations
\begin{quote}
   This is a short quotation.  It consists of a
   single paragraph of text.  See how it is formatted.
\end{quote}
and longer ones.
\begin{quotation}
   This is a longer quotation.  It consists of two
   paragraphs of text, neither of which are
   particularly interesting.

   This is the second paragraph of the quotation.  It
   is just as dull as the first paragraph.
\end{quotation}
Another frequently-displayed structure is a list.
The following is an example of an \emph{itemized}
list.
\begin{itemize}
   \item This is the first item of an itemized list.
         Each item in the list is marked with a ``tick''.
         You don't have to worry about what kind of tick
         mark is used.

   \item This is the second item of the list.  It
         contains another list nested inside it.  The inner
         list is an \emph{enumerated} list.
         \begin{enumerate}
            \item This is the first item of an enumerated
                  list that is nested within the itemized list.

            \item This is the second item of the inner list.
                  \LaTeX\ allows you to nest lists deeper than
                  you really should.
         \end{enumerate}
         This is the rest of the second item of the outer
         list.  It is no more interesting than any other
         part of the item.
   \item This is the third item of the list.
\end{itemize}
You can even display poetry.
\begin{verse}
   There is an environment
    for verse \\             % The \\ command separates lines
   Whose features some poets % within a stanza.
   will curse.

                             % One or more blank lines separate stanzas.

   For instead of making\\
   Them do \emph{all} line breaking, \\
   It allows them to put too many words on a line when they'd rather be
   forced to be terse.
\end{verse}

Mathematical formulas may also be displayed.  A
displayed formula
is
one-line long; multiline
formulas require special formatting instructions.
   \[  \ip{\Gamma}{\psi'} = x'' + y^{2} + z_{i}^{n}\]
Don't start a paragraph with a displayed equation,
nor make one a paragraph by itself.

\end{document}               % End of document.

\section{Bench Huge Section 1486}
% synthetic bench section 1486
% This is a sample LaTeX input file.  (Version of 12 August 2004.)
%
% A '%' character causes TeX to ignore all remaining text on the line,
% and is used for comments like this one.

\documentclass{article}      % Specifies the document class

                             % The preamble begins here.
\title{An Example Document}  % Declares the document's title.
\author{Leslie Lamport}      % Declares the author's name.
\date{January 21, 1994}      % Deleting this command produces today's date.

\newcommand{\ip}[2]{(#1, #2)}
                             % Defines \ip{arg1}{arg2} to mean
                             % (arg1, arg2).

%\newcommand{\ip}[2]{\langle #1 | #2\rangle}
                             % This is an alternative definition of
                             % \ip that is commented out.

\begin{document}             % End of preamble and beginning of text.

\maketitle                   % Produces the title.

This is an example input file.  Comparing it with
the output it generates can show you how to
produce a simple document of your own.

\section{Ordinary Text}      % Produces section heading.  Lower-level
                             % sections are begun with similar
                             % \subsection and \subsubsection commands.

The ends  of words and sentences are marked
  by   spaces. It  doesn't matter how many
spaces    you type; one is as good as 100.  The
end of   a line counts as a space.

One   or more   blank lines denote the  end
of  a paragraph.

Since any number of consecutive spaces are treated
like a single one, the formatting of the input
file makes no difference to
      \LaTeX,                % The \LaTeX command generates the LaTeX logo.
but it makes a difference to you.  When you use
\LaTeX, making your input file as easy to read
as possible will be a great help as you write
your document and when you change it.  This sample
file shows how you can add comments to your own input
file.

Because printing is different from typewriting,
there are a number of things that you have to do
differently when preparing an input file than if
you were just typing the document directly.
Quotation marks like
       ``this''
have to be handled specially, as do quotes within
quotes:
       ``\,`this'            % \, separates the double and single quote.
        is what I just
        wrote, not  `that'\,''.

Dashes come in three sizes: an
       intra-word
dash, a medium dash for number ranges like
       1--2,
and a punctuation
       dash---like
this.

A sentence-ending space should be larger than the
space between words within a sentence.  You
sometimes have to type special commands in
conjunction with punctuation characters to get
this right, as in the following sentence.
       Gnats, gnus, etc.\ all  % `\ ' makes an inter-word space.
       begin with G\@.         % \@ marks end-of-sentence punctuation.
You should check the spaces after periods when
reading your output to make sure you haven't
forgotten any special cases.  Generating an
ellipsis
       \ldots\               % `\ ' is needed after `\ldots' because TeX
                             % ignores spaces after command names like \ldots
                             % made from \ + letters.
                             %
                             % Note how a `%' character causes TeX to ignore
                             % the end of the input line, so these blank lines
                             % do not start a new paragraph.
                             %
with the right spacing around the periods requires
a special command.

\LaTeX\ interprets some common characters as
commands, so you must type special commands to
generate them.  These characters include the
following:
       \$ \& \% \# \{ and \}.

In printing, text is usually emphasized with an
       \emph{italic}
type style.

\begin{em}
   A long segment of text can also be emphasized
   in this way.  Text within such a segment can be
   given \emph{additional} emphasis.
\end{em}

It is sometimes necessary to prevent \LaTeX\ from
breaking a line where it might otherwise do so.
This may be at a space, as between the ``Mr.''\ and
``Jones'' in
       ``Mr.~Jones'',        % ~ produces an unbreakable interword space.
or within a word---especially when the word is a
symbol like
       \mbox{\emph{itemnum}}
that makes little sense when hyphenated across
lines.

Footnotes\footnote{This is an example of a footnote.}
pose no problem.

\LaTeX\ is good at typesetting mathematical formulas
like
       \( x-3y + z = 7 \)
or
       \( a_{1} > x^{2n} + y^{2n} > x' \)
or
       \( \ip{A}{B} = \sum_{i} a_{i} b_{i} \).
The spaces you type in a formula are
ignored.  Remember that a letter like
       $x$                   % $ ... $  and  \( ... \)  are equivalent
is a formula when it denotes a mathematical
symbol, and it should be typed as one.

\section{Displayed Text}

Text is displayed by indenting it from the left
margin.  Quotations are commonly displayed.  There
are short quotations
\begin{quote}
   This is a short quotation.  It consists of a
   single paragraph of text.  See how it is formatted.
\end{quote}
and longer ones.
\begin{quotation}
   This is a longer quotation.  It consists of two
   paragraphs of text, neither of which are
   particularly interesting.

   This is the second paragraph of the quotation.  It
   is just as dull as the first paragraph.
\end{quotation}
Another frequently-displayed structure is a list.
The following is an example of an \emph{itemized}
list.
\begin{itemize}
   \item This is the first item of an itemized list.
         Each item in the list is marked with a ``tick''.
         You don't have to worry about what kind of tick
         mark is used.

   \item This is the second item of the list.  It
         contains another list nested inside it.  The inner
         list is an \emph{enumerated} list.
         \begin{enumerate}
            \item This is the first item of an enumerated
                  list that is nested within the itemized list.

            \item This is the second item of the inner list.
                  \LaTeX\ allows you to nest lists deeper than
                  you really should.
         \end{enumerate}
         This is the rest of the second item of the outer
         list.  It is no more interesting than any other
         part of the item.
   \item This is the third item of the list.
\end{itemize}
You can even display poetry.
\begin{verse}
   There is an environment
    for verse \\             % The \\ command separates lines
   Whose features some poets % within a stanza.
   will curse.

                             % One or more blank lines separate stanzas.

   For instead of making\\
   Them do \emph{all} line breaking, \\
   It allows them to put too many words on a line when they'd rather be
   forced to be terse.
\end{verse}

Mathematical formulas may also be displayed.  A
displayed formula
is
one-line long; multiline
formulas require special formatting instructions.
   \[  \ip{\Gamma}{\psi'} = x'' + y^{2} + z_{i}^{n}\]
Don't start a paragraph with a displayed equation,
nor make one a paragraph by itself.

\end{document}               % End of document.

\section{Bench Huge Section 1487}
% synthetic bench section 1487
% This is a sample LaTeX input file.  (Version of 12 August 2004.)
%
% A '%' character causes TeX to ignore all remaining text on the line,
% and is used for comments like this one.

\documentclass{article}      % Specifies the document class

                             % The preamble begins here.
\title{An Example Document}  % Declares the document's title.
\author{Leslie Lamport}      % Declares the author's name.
\date{January 21, 1994}      % Deleting this command produces today's date.

\newcommand{\ip}[2]{(#1, #2)}
                             % Defines \ip{arg1}{arg2} to mean
                             % (arg1, arg2).

%\newcommand{\ip}[2]{\langle #1 | #2\rangle}
                             % This is an alternative definition of
                             % \ip that is commented out.

\begin{document}             % End of preamble and beginning of text.

\maketitle                   % Produces the title.

This is an example input file.  Comparing it with
the output it generates can show you how to
produce a simple document of your own.

\section{Ordinary Text}      % Produces section heading.  Lower-level
                             % sections are begun with similar
                             % \subsection and \subsubsection commands.

The ends  of words and sentences are marked
  by   spaces. It  doesn't matter how many
spaces    you type; one is as good as 100.  The
end of   a line counts as a space.

One   or more   blank lines denote the  end
of  a paragraph.

Since any number of consecutive spaces are treated
like a single one, the formatting of the input
file makes no difference to
      \LaTeX,                % The \LaTeX command generates the LaTeX logo.
but it makes a difference to you.  When you use
\LaTeX, making your input file as easy to read
as possible will be a great help as you write
your document and when you change it.  This sample
file shows how you can add comments to your own input
file.

Because printing is different from typewriting,
there are a number of things that you have to do
differently when preparing an input file than if
you were just typing the document directly.
Quotation marks like
       ``this''
have to be handled specially, as do quotes within
quotes:
       ``\,`this'            % \, separates the double and single quote.
        is what I just
        wrote, not  `that'\,''.

Dashes come in three sizes: an
       intra-word
dash, a medium dash for number ranges like
       1--2,
and a punctuation
       dash---like
this.

A sentence-ending space should be larger than the
space between words within a sentence.  You
sometimes have to type special commands in
conjunction with punctuation characters to get
this right, as in the following sentence.
       Gnats, gnus, etc.\ all  % `\ ' makes an inter-word space.
       begin with G\@.         % \@ marks end-of-sentence punctuation.
You should check the spaces after periods when
reading your output to make sure you haven't
forgotten any special cases.  Generating an
ellipsis
       \ldots\               % `\ ' is needed after `\ldots' because TeX
                             % ignores spaces after command names like \ldots
                             % made from \ + letters.
                             %
                             % Note how a `%' character causes TeX to ignore
                             % the end of the input line, so these blank lines
                             % do not start a new paragraph.
                             %
with the right spacing around the periods requires
a special command.

\LaTeX\ interprets some common characters as
commands, so you must type special commands to
generate them.  These characters include the
following:
       \$ \& \% \# \{ and \}.

In printing, text is usually emphasized with an
       \emph{italic}
type style.

\begin{em}
   A long segment of text can also be emphasized
   in this way.  Text within such a segment can be
   given \emph{additional} emphasis.
\end{em}

It is sometimes necessary to prevent \LaTeX\ from
breaking a line where it might otherwise do so.
This may be at a space, as between the ``Mr.''\ and
``Jones'' in
       ``Mr.~Jones'',        % ~ produces an unbreakable interword space.
or within a word---especially when the word is a
symbol like
       \mbox{\emph{itemnum}}
that makes little sense when hyphenated across
lines.

Footnotes\footnote{This is an example of a footnote.}
pose no problem.

\LaTeX\ is good at typesetting mathematical formulas
like
       \( x-3y + z = 7 \)
or
       \( a_{1} > x^{2n} + y^{2n} > x' \)
or
       \( \ip{A}{B} = \sum_{i} a_{i} b_{i} \).
The spaces you type in a formula are
ignored.  Remember that a letter like
       $x$                   % $ ... $  and  \( ... \)  are equivalent
is a formula when it denotes a mathematical
symbol, and it should be typed as one.

\section{Displayed Text}

Text is displayed by indenting it from the left
margin.  Quotations are commonly displayed.  There
are short quotations
\begin{quote}
   This is a short quotation.  It consists of a
   single paragraph of text.  See how it is formatted.
\end{quote}
and longer ones.
\begin{quotation}
   This is a longer quotation.  It consists of two
   paragraphs of text, neither of which are
   particularly interesting.

   This is the second paragraph of the quotation.  It
   is just as dull as the first paragraph.
\end{quotation}
Another frequently-displayed structure is a list.
The following is an example of an \emph{itemized}
list.
\begin{itemize}
   \item This is the first item of an itemized list.
         Each item in the list is marked with a ``tick''.
         You don't have to worry about what kind of tick
         mark is used.

   \item This is the second item of the list.  It
         contains another list nested inside it.  The inner
         list is an \emph{enumerated} list.
         \begin{enumerate}
            \item This is the first item of an enumerated
                  list that is nested within the itemized list.

            \item This is the second item of the inner list.
                  \LaTeX\ allows you to nest lists deeper than
                  you really should.
         \end{enumerate}
         This is the rest of the second item of the outer
         list.  It is no more interesting than any other
         part of the item.
   \item This is the third item of the list.
\end{itemize}
You can even display poetry.
\begin{verse}
   There is an environment
    for verse \\             % The \\ command separates lines
   Whose features some poets % within a stanza.
   will curse.

                             % One or more blank lines separate stanzas.

   For instead of making\\
   Them do \emph{all} line breaking, \\
   It allows them to put too many words on a line when they'd rather be
   forced to be terse.
\end{verse}

Mathematical formulas may also be displayed.  A
displayed formula
is
one-line long; multiline
formulas require special formatting instructions.
   \[  \ip{\Gamma}{\psi'} = x'' + y^{2} + z_{i}^{n}\]
Don't start a paragraph with a displayed equation,
nor make one a paragraph by itself.

\end{document}               % End of document.

\section{Bench Huge Section 1488}
% synthetic bench section 1488
% This is a sample LaTeX input file.  (Version of 12 August 2004.)
%
% A '%' character causes TeX to ignore all remaining text on the line,
% and is used for comments like this one.

\documentclass{article}      % Specifies the document class

                             % The preamble begins here.
\title{An Example Document}  % Declares the document's title.
\author{Leslie Lamport}      % Declares the author's name.
\date{January 21, 1994}      % Deleting this command produces today's date.

\newcommand{\ip}[2]{(#1, #2)}
                             % Defines \ip{arg1}{arg2} to mean
                             % (arg1, arg2).

%\newcommand{\ip}[2]{\langle #1 | #2\rangle}
                             % This is an alternative definition of
                             % \ip that is commented out.

\begin{document}             % End of preamble and beginning of text.

\maketitle                   % Produces the title.

This is an example input file.  Comparing it with
the output it generates can show you how to
produce a simple document of your own.

\section{Ordinary Text}      % Produces section heading.  Lower-level
                             % sections are begun with similar
                             % \subsection and \subsubsection commands.

The ends  of words and sentences are marked
  by   spaces. It  doesn't matter how many
spaces    you type; one is as good as 100.  The
end of   a line counts as a space.

One   or more   blank lines denote the  end
of  a paragraph.

Since any number of consecutive spaces are treated
like a single one, the formatting of the input
file makes no difference to
      \LaTeX,                % The \LaTeX command generates the LaTeX logo.
but it makes a difference to you.  When you use
\LaTeX, making your input file as easy to read
as possible will be a great help as you write
your document and when you change it.  This sample
file shows how you can add comments to your own input
file.

Because printing is different from typewriting,
there are a number of things that you have to do
differently when preparing an input file than if
you were just typing the document directly.
Quotation marks like
       ``this''
have to be handled specially, as do quotes within
quotes:
       ``\,`this'            % \, separates the double and single quote.
        is what I just
        wrote, not  `that'\,''.

Dashes come in three sizes: an
       intra-word
dash, a medium dash for number ranges like
       1--2,
and a punctuation
       dash---like
this.

A sentence-ending space should be larger than the
space between words within a sentence.  You
sometimes have to type special commands in
conjunction with punctuation characters to get
this right, as in the following sentence.
       Gnats, gnus, etc.\ all  % `\ ' makes an inter-word space.
       begin with G\@.         % \@ marks end-of-sentence punctuation.
You should check the spaces after periods when
reading your output to make sure you haven't
forgotten any special cases.  Generating an
ellipsis
       \ldots\               % `\ ' is needed after `\ldots' because TeX
                             % ignores spaces after command names like \ldots
                             % made from \ + letters.
                             %
                             % Note how a `%' character causes TeX to ignore
                             % the end of the input line, so these blank lines
                             % do not start a new paragraph.
                             %
with the right spacing around the periods requires
a special command.

\LaTeX\ interprets some common characters as
commands, so you must type special commands to
generate them.  These characters include the
following:
       \$ \& \% \# \{ and \}.

In printing, text is usually emphasized with an
       \emph{italic}
type style.

\begin{em}
   A long segment of text can also be emphasized
   in this way.  Text within such a segment can be
   given \emph{additional} emphasis.
\end{em}

It is sometimes necessary to prevent \LaTeX\ from
breaking a line where it might otherwise do so.
This may be at a space, as between the ``Mr.''\ and
``Jones'' in
       ``Mr.~Jones'',        % ~ produces an unbreakable interword space.
or within a word---especially when the word is a
symbol like
       \mbox{\emph{itemnum}}
that makes little sense when hyphenated across
lines.

Footnotes\footnote{This is an example of a footnote.}
pose no problem.

\LaTeX\ is good at typesetting mathematical formulas
like
       \( x-3y + z = 7 \)
or
       \( a_{1} > x^{2n} + y^{2n} > x' \)
or
       \( \ip{A}{B} = \sum_{i} a_{i} b_{i} \).
The spaces you type in a formula are
ignored.  Remember that a letter like
       $x$                   % $ ... $  and  \( ... \)  are equivalent
is a formula when it denotes a mathematical
symbol, and it should be typed as one.

\section{Displayed Text}

Text is displayed by indenting it from the left
margin.  Quotations are commonly displayed.  There
are short quotations
\begin{quote}
   This is a short quotation.  It consists of a
   single paragraph of text.  See how it is formatted.
\end{quote}
and longer ones.
\begin{quotation}
   This is a longer quotation.  It consists of two
   paragraphs of text, neither of which are
   particularly interesting.

   This is the second paragraph of the quotation.  It
   is just as dull as the first paragraph.
\end{quotation}
Another frequently-displayed structure is a list.
The following is an example of an \emph{itemized}
list.
\begin{itemize}
   \item This is the first item of an itemized list.
         Each item in the list is marked with a ``tick''.
         You don't have to worry about what kind of tick
         mark is used.

   \item This is the second item of the list.  It
         contains another list nested inside it.  The inner
         list is an \emph{enumerated} list.
         \begin{enumerate}
            \item This is the first item of an enumerated
                  list that is nested within the itemized list.

            \item This is the second item of the inner list.
                  \LaTeX\ allows you to nest lists deeper than
                  you really should.
         \end{enumerate}
         This is the rest of the second item of the outer
         list.  It is no more interesting than any other
         part of the item.
   \item This is the third item of the list.
\end{itemize}
You can even display poetry.
\begin{verse}
   There is an environment
    for verse \\             % The \\ command separates lines
   Whose features some poets % within a stanza.
   will curse.

                             % One or more blank lines separate stanzas.

   For instead of making\\
   Them do \emph{all} line breaking, \\
   It allows them to put too many words on a line when they'd rather be
   forced to be terse.
\end{verse}

Mathematical formulas may also be displayed.  A
displayed formula
is
one-line long; multiline
formulas require special formatting instructions.
   \[  \ip{\Gamma}{\psi'} = x'' + y^{2} + z_{i}^{n}\]
Don't start a paragraph with a displayed equation,
nor make one a paragraph by itself.

\end{document}               % End of document.

\section{Bench Huge Section 1489}
% synthetic bench section 1489
% This is a sample LaTeX input file.  (Version of 12 August 2004.)
%
% A '%' character causes TeX to ignore all remaining text on the line,
% and is used for comments like this one.

\documentclass{article}      % Specifies the document class

                             % The preamble begins here.
\title{An Example Document}  % Declares the document's title.
\author{Leslie Lamport}      % Declares the author's name.
\date{January 21, 1994}      % Deleting this command produces today's date.

\newcommand{\ip}[2]{(#1, #2)}
                             % Defines \ip{arg1}{arg2} to mean
                             % (arg1, arg2).

%\newcommand{\ip}[2]{\langle #1 | #2\rangle}
                             % This is an alternative definition of
                             % \ip that is commented out.

\begin{document}             % End of preamble and beginning of text.

\maketitle                   % Produces the title.

This is an example input file.  Comparing it with
the output it generates can show you how to
produce a simple document of your own.

\section{Ordinary Text}      % Produces section heading.  Lower-level
                             % sections are begun with similar
                             % \subsection and \subsubsection commands.

The ends  of words and sentences are marked
  by   spaces. It  doesn't matter how many
spaces    you type; one is as good as 100.  The
end of   a line counts as a space.

One   or more   blank lines denote the  end
of  a paragraph.

Since any number of consecutive spaces are treated
like a single one, the formatting of the input
file makes no difference to
      \LaTeX,                % The \LaTeX command generates the LaTeX logo.
but it makes a difference to you.  When you use
\LaTeX, making your input file as easy to read
as possible will be a great help as you write
your document and when you change it.  This sample
file shows how you can add comments to your own input
file.

Because printing is different from typewriting,
there are a number of things that you have to do
differently when preparing an input file than if
you were just typing the document directly.
Quotation marks like
       ``this''
have to be handled specially, as do quotes within
quotes:
       ``\,`this'            % \, separates the double and single quote.
        is what I just
        wrote, not  `that'\,''.

Dashes come in three sizes: an
       intra-word
dash, a medium dash for number ranges like
       1--2,
and a punctuation
       dash---like
this.

A sentence-ending space should be larger than the
space between words within a sentence.  You
sometimes have to type special commands in
conjunction with punctuation characters to get
this right, as in the following sentence.
       Gnats, gnus, etc.\ all  % `\ ' makes an inter-word space.
       begin with G\@.         % \@ marks end-of-sentence punctuation.
You should check the spaces after periods when
reading your output to make sure you haven't
forgotten any special cases.  Generating an
ellipsis
       \ldots\               % `\ ' is needed after `\ldots' because TeX
                             % ignores spaces after command names like \ldots
                             % made from \ + letters.
                             %
                             % Note how a `%' character causes TeX to ignore
                             % the end of the input line, so these blank lines
                             % do not start a new paragraph.
                             %
with the right spacing around the periods requires
a special command.

\LaTeX\ interprets some common characters as
commands, so you must type special commands to
generate them.  These characters include the
following:
       \$ \& \% \# \{ and \}.

In printing, text is usually emphasized with an
       \emph{italic}
type style.

\begin{em}
   A long segment of text can also be emphasized
   in this way.  Text within such a segment can be
   given \emph{additional} emphasis.
\end{em}

It is sometimes necessary to prevent \LaTeX\ from
breaking a line where it might otherwise do so.
This may be at a space, as between the ``Mr.''\ and
``Jones'' in
       ``Mr.~Jones'',        % ~ produces an unbreakable interword space.
or within a word---especially when the word is a
symbol like
       \mbox{\emph{itemnum}}
that makes little sense when hyphenated across
lines.

Footnotes\footnote{This is an example of a footnote.}
pose no problem.

\LaTeX\ is good at typesetting mathematical formulas
like
       \( x-3y + z = 7 \)
or
       \( a_{1} > x^{2n} + y^{2n} > x' \)
or
       \( \ip{A}{B} = \sum_{i} a_{i} b_{i} \).
The spaces you type in a formula are
ignored.  Remember that a letter like
       $x$                   % $ ... $  and  \( ... \)  are equivalent
is a formula when it denotes a mathematical
symbol, and it should be typed as one.

\section{Displayed Text}

Text is displayed by indenting it from the left
margin.  Quotations are commonly displayed.  There
are short quotations
\begin{quote}
   This is a short quotation.  It consists of a
   single paragraph of text.  See how it is formatted.
\end{quote}
and longer ones.
\begin{quotation}
   This is a longer quotation.  It consists of two
   paragraphs of text, neither of which are
   particularly interesting.

   This is the second paragraph of the quotation.  It
   is just as dull as the first paragraph.
\end{quotation}
Another frequently-displayed structure is a list.
The following is an example of an \emph{itemized}
list.
\begin{itemize}
   \item This is the first item of an itemized list.
         Each item in the list is marked with a ``tick''.
         You don't have to worry about what kind of tick
         mark is used.

   \item This is the second item of the list.  It
         contains another list nested inside it.  The inner
         list is an \emph{enumerated} list.
         \begin{enumerate}
            \item This is the first item of an enumerated
                  list that is nested within the itemized list.

            \item This is the second item of the inner list.
                  \LaTeX\ allows you to nest lists deeper than
                  you really should.
         \end{enumerate}
         This is the rest of the second item of the outer
         list.  It is no more interesting than any other
         part of the item.
   \item This is the third item of the list.
\end{itemize}
You can even display poetry.
\begin{verse}
   There is an environment
    for verse \\             % The \\ command separates lines
   Whose features some poets % within a stanza.
   will curse.

                             % One or more blank lines separate stanzas.

   For instead of making\\
   Them do \emph{all} line breaking, \\
   It allows them to put too many words on a line when they'd rather be
   forced to be terse.
\end{verse}

Mathematical formulas may also be displayed.  A
displayed formula
is
one-line long; multiline
formulas require special formatting instructions.
   \[  \ip{\Gamma}{\psi'} = x'' + y^{2} + z_{i}^{n}\]
Don't start a paragraph with a displayed equation,
nor make one a paragraph by itself.

\end{document}               % End of document.

\section{Bench Huge Section 1490}
% synthetic bench section 1490
% This is a sample LaTeX input file.  (Version of 12 August 2004.)
%
% A '%' character causes TeX to ignore all remaining text on the line,
% and is used for comments like this one.

\documentclass{article}      % Specifies the document class

                             % The preamble begins here.
\title{An Example Document}  % Declares the document's title.
\author{Leslie Lamport}      % Declares the author's name.
\date{January 21, 1994}      % Deleting this command produces today's date.

\newcommand{\ip}[2]{(#1, #2)}
                             % Defines \ip{arg1}{arg2} to mean
                             % (arg1, arg2).

%\newcommand{\ip}[2]{\langle #1 | #2\rangle}
                             % This is an alternative definition of
                             % \ip that is commented out.

\begin{document}             % End of preamble and beginning of text.

\maketitle                   % Produces the title.

This is an example input file.  Comparing it with
the output it generates can show you how to
produce a simple document of your own.

\section{Ordinary Text}      % Produces section heading.  Lower-level
                             % sections are begun with similar
                             % \subsection and \subsubsection commands.

The ends  of words and sentences are marked
  by   spaces. It  doesn't matter how many
spaces    you type; one is as good as 100.  The
end of   a line counts as a space.

One   or more   blank lines denote the  end
of  a paragraph.

Since any number of consecutive spaces are treated
like a single one, the formatting of the input
file makes no difference to
      \LaTeX,                % The \LaTeX command generates the LaTeX logo.
but it makes a difference to you.  When you use
\LaTeX, making your input file as easy to read
as possible will be a great help as you write
your document and when you change it.  This sample
file shows how you can add comments to your own input
file.

Because printing is different from typewriting,
there are a number of things that you have to do
differently when preparing an input file than if
you were just typing the document directly.
Quotation marks like
       ``this''
have to be handled specially, as do quotes within
quotes:
       ``\,`this'            % \, separates the double and single quote.
        is what I just
        wrote, not  `that'\,''.

Dashes come in three sizes: an
       intra-word
dash, a medium dash for number ranges like
       1--2,
and a punctuation
       dash---like
this.

A sentence-ending space should be larger than the
space between words within a sentence.  You
sometimes have to type special commands in
conjunction with punctuation characters to get
this right, as in the following sentence.
       Gnats, gnus, etc.\ all  % `\ ' makes an inter-word space.
       begin with G\@.         % \@ marks end-of-sentence punctuation.
You should check the spaces after periods when
reading your output to make sure you haven't
forgotten any special cases.  Generating an
ellipsis
       \ldots\               % `\ ' is needed after `\ldots' because TeX
                             % ignores spaces after command names like \ldots
                             % made from \ + letters.
                             %
                             % Note how a `%' character causes TeX to ignore
                             % the end of the input line, so these blank lines
                             % do not start a new paragraph.
                             %
with the right spacing around the periods requires
a special command.

\LaTeX\ interprets some common characters as
commands, so you must type special commands to
generate them.  These characters include the
following:
       \$ \& \% \# \{ and \}.

In printing, text is usually emphasized with an
       \emph{italic}
type style.

\begin{em}
   A long segment of text can also be emphasized
   in this way.  Text within such a segment can be
   given \emph{additional} emphasis.
\end{em}

It is sometimes necessary to prevent \LaTeX\ from
breaking a line where it might otherwise do so.
This may be at a space, as between the ``Mr.''\ and
``Jones'' in
       ``Mr.~Jones'',        % ~ produces an unbreakable interword space.
or within a word---especially when the word is a
symbol like
       \mbox{\emph{itemnum}}
that makes little sense when hyphenated across
lines.

Footnotes\footnote{This is an example of a footnote.}
pose no problem.

\LaTeX\ is good at typesetting mathematical formulas
like
       \( x-3y + z = 7 \)
or
       \( a_{1} > x^{2n} + y^{2n} > x' \)
or
       \( \ip{A}{B} = \sum_{i} a_{i} b_{i} \).
The spaces you type in a formula are
ignored.  Remember that a letter like
       $x$                   % $ ... $  and  \( ... \)  are equivalent
is a formula when it denotes a mathematical
symbol, and it should be typed as one.

\section{Displayed Text}

Text is displayed by indenting it from the left
margin.  Quotations are commonly displayed.  There
are short quotations
\begin{quote}
   This is a short quotation.  It consists of a
   single paragraph of text.  See how it is formatted.
\end{quote}
and longer ones.
\begin{quotation}
   This is a longer quotation.  It consists of two
   paragraphs of text, neither of which are
   particularly interesting.

   This is the second paragraph of the quotation.  It
   is just as dull as the first paragraph.
\end{quotation}
Another frequently-displayed structure is a list.
The following is an example of an \emph{itemized}
list.
\begin{itemize}
   \item This is the first item of an itemized list.
         Each item in the list is marked with a ``tick''.
         You don't have to worry about what kind of tick
         mark is used.

   \item This is the second item of the list.  It
         contains another list nested inside it.  The inner
         list is an \emph{enumerated} list.
         \begin{enumerate}
            \item This is the first item of an enumerated
                  list that is nested within the itemized list.

            \item This is the second item of the inner list.
                  \LaTeX\ allows you to nest lists deeper than
                  you really should.
         \end{enumerate}
         This is the rest of the second item of the outer
         list.  It is no more interesting than any other
         part of the item.
   \item This is the third item of the list.
\end{itemize}
You can even display poetry.
\begin{verse}
   There is an environment
    for verse \\             % The \\ command separates lines
   Whose features some poets % within a stanza.
   will curse.

                             % One or more blank lines separate stanzas.

   For instead of making\\
   Them do \emph{all} line breaking, \\
   It allows them to put too many words on a line when they'd rather be
   forced to be terse.
\end{verse}

Mathematical formulas may also be displayed.  A
displayed formula
is
one-line long; multiline
formulas require special formatting instructions.
   \[  \ip{\Gamma}{\psi'} = x'' + y^{2} + z_{i}^{n}\]
Don't start a paragraph with a displayed equation,
nor make one a paragraph by itself.

\end{document}               % End of document.

\section{Bench Huge Section 1491}
% synthetic bench section 1491
% This is a sample LaTeX input file.  (Version of 12 August 2004.)
%
% A '%' character causes TeX to ignore all remaining text on the line,
% and is used for comments like this one.

\documentclass{article}      % Specifies the document class

                             % The preamble begins here.
\title{An Example Document}  % Declares the document's title.
\author{Leslie Lamport}      % Declares the author's name.
\date{January 21, 1994}      % Deleting this command produces today's date.

\newcommand{\ip}[2]{(#1, #2)}
                             % Defines \ip{arg1}{arg2} to mean
                             % (arg1, arg2).

%\newcommand{\ip}[2]{\langle #1 | #2\rangle}
                             % This is an alternative definition of
                             % \ip that is commented out.

\begin{document}             % End of preamble and beginning of text.

\maketitle                   % Produces the title.

This is an example input file.  Comparing it with
the output it generates can show you how to
produce a simple document of your own.

\section{Ordinary Text}      % Produces section heading.  Lower-level
                             % sections are begun with similar
                             % \subsection and \subsubsection commands.

The ends  of words and sentences are marked
  by   spaces. It  doesn't matter how many
spaces    you type; one is as good as 100.  The
end of   a line counts as a space.

One   or more   blank lines denote the  end
of  a paragraph.

Since any number of consecutive spaces are treated
like a single one, the formatting of the input
file makes no difference to
      \LaTeX,                % The \LaTeX command generates the LaTeX logo.
but it makes a difference to you.  When you use
\LaTeX, making your input file as easy to read
as possible will be a great help as you write
your document and when you change it.  This sample
file shows how you can add comments to your own input
file.

Because printing is different from typewriting,
there are a number of things that you have to do
differently when preparing an input file than if
you were just typing the document directly.
Quotation marks like
       ``this''
have to be handled specially, as do quotes within
quotes:
       ``\,`this'            % \, separates the double and single quote.
        is what I just
        wrote, not  `that'\,''.

Dashes come in three sizes: an
       intra-word
dash, a medium dash for number ranges like
       1--2,
and a punctuation
       dash---like
this.

A sentence-ending space should be larger than the
space between words within a sentence.  You
sometimes have to type special commands in
conjunction with punctuation characters to get
this right, as in the following sentence.
       Gnats, gnus, etc.\ all  % `\ ' makes an inter-word space.
       begin with G\@.         % \@ marks end-of-sentence punctuation.
You should check the spaces after periods when
reading your output to make sure you haven't
forgotten any special cases.  Generating an
ellipsis
       \ldots\               % `\ ' is needed after `\ldots' because TeX
                             % ignores spaces after command names like \ldots
                             % made from \ + letters.
                             %
                             % Note how a `%' character causes TeX to ignore
                             % the end of the input line, so these blank lines
                             % do not start a new paragraph.
                             %
with the right spacing around the periods requires
a special command.

\LaTeX\ interprets some common characters as
commands, so you must type special commands to
generate them.  These characters include the
following:
       \$ \& \% \# \{ and \}.

In printing, text is usually emphasized with an
       \emph{italic}
type style.

\begin{em}
   A long segment of text can also be emphasized
   in this way.  Text within such a segment can be
   given \emph{additional} emphasis.
\end{em}

It is sometimes necessary to prevent \LaTeX\ from
breaking a line where it might otherwise do so.
This may be at a space, as between the ``Mr.''\ and
``Jones'' in
       ``Mr.~Jones'',        % ~ produces an unbreakable interword space.
or within a word---especially when the word is a
symbol like
       \mbox{\emph{itemnum}}
that makes little sense when hyphenated across
lines.

Footnotes\footnote{This is an example of a footnote.}
pose no problem.

\LaTeX\ is good at typesetting mathematical formulas
like
       \( x-3y + z = 7 \)
or
       \( a_{1} > x^{2n} + y^{2n} > x' \)
or
       \( \ip{A}{B} = \sum_{i} a_{i} b_{i} \).
The spaces you type in a formula are
ignored.  Remember that a letter like
       $x$                   % $ ... $  and  \( ... \)  are equivalent
is a formula when it denotes a mathematical
symbol, and it should be typed as one.

\section{Displayed Text}

Text is displayed by indenting it from the left
margin.  Quotations are commonly displayed.  There
are short quotations
\begin{quote}
   This is a short quotation.  It consists of a
   single paragraph of text.  See how it is formatted.
\end{quote}
and longer ones.
\begin{quotation}
   This is a longer quotation.  It consists of two
   paragraphs of text, neither of which are
   particularly interesting.

   This is the second paragraph of the quotation.  It
   is just as dull as the first paragraph.
\end{quotation}
Another frequently-displayed structure is a list.
The following is an example of an \emph{itemized}
list.
\begin{itemize}
   \item This is the first item of an itemized list.
         Each item in the list is marked with a ``tick''.
         You don't have to worry about what kind of tick
         mark is used.

   \item This is the second item of the list.  It
         contains another list nested inside it.  The inner
         list is an \emph{enumerated} list.
         \begin{enumerate}
            \item This is the first item of an enumerated
                  list that is nested within the itemized list.

            \item This is the second item of the inner list.
                  \LaTeX\ allows you to nest lists deeper than
                  you really should.
         \end{enumerate}
         This is the rest of the second item of the outer
         list.  It is no more interesting than any other
         part of the item.
   \item This is the third item of the list.
\end{itemize}
You can even display poetry.
\begin{verse}
   There is an environment
    for verse \\             % The \\ command separates lines
   Whose features some poets % within a stanza.
   will curse.

                             % One or more blank lines separate stanzas.

   For instead of making\\
   Them do \emph{all} line breaking, \\
   It allows them to put too many words on a line when they'd rather be
   forced to be terse.
\end{verse}

Mathematical formulas may also be displayed.  A
displayed formula
is
one-line long; multiline
formulas require special formatting instructions.
   \[  \ip{\Gamma}{\psi'} = x'' + y^{2} + z_{i}^{n}\]
Don't start a paragraph with a displayed equation,
nor make one a paragraph by itself.

\end{document}               % End of document.

\section{Bench Huge Section 1492}
% synthetic bench section 1492
% This is a sample LaTeX input file.  (Version of 12 August 2004.)
%
% A '%' character causes TeX to ignore all remaining text on the line,
% and is used for comments like this one.

\documentclass{article}      % Specifies the document class

                             % The preamble begins here.
\title{An Example Document}  % Declares the document's title.
\author{Leslie Lamport}      % Declares the author's name.
\date{January 21, 1994}      % Deleting this command produces today's date.

\newcommand{\ip}[2]{(#1, #2)}
                             % Defines \ip{arg1}{arg2} to mean
                             % (arg1, arg2).

%\newcommand{\ip}[2]{\langle #1 | #2\rangle}
                             % This is an alternative definition of
                             % \ip that is commented out.

\begin{document}             % End of preamble and beginning of text.

\maketitle                   % Produces the title.

This is an example input file.  Comparing it with
the output it generates can show you how to
produce a simple document of your own.

\section{Ordinary Text}      % Produces section heading.  Lower-level
                             % sections are begun with similar
                             % \subsection and \subsubsection commands.

The ends  of words and sentences are marked
  by   spaces. It  doesn't matter how many
spaces    you type; one is as good as 100.  The
end of   a line counts as a space.

One   or more   blank lines denote the  end
of  a paragraph.

Since any number of consecutive spaces are treated
like a single one, the formatting of the input
file makes no difference to
      \LaTeX,                % The \LaTeX command generates the LaTeX logo.
but it makes a difference to you.  When you use
\LaTeX, making your input file as easy to read
as possible will be a great help as you write
your document and when you change it.  This sample
file shows how you can add comments to your own input
file.

Because printing is different from typewriting,
there are a number of things that you have to do
differently when preparing an input file than if
you were just typing the document directly.
Quotation marks like
       ``this''
have to be handled specially, as do quotes within
quotes:
       ``\,`this'            % \, separates the double and single quote.
        is what I just
        wrote, not  `that'\,''.

Dashes come in three sizes: an
       intra-word
dash, a medium dash for number ranges like
       1--2,
and a punctuation
       dash---like
this.

A sentence-ending space should be larger than the
space between words within a sentence.  You
sometimes have to type special commands in
conjunction with punctuation characters to get
this right, as in the following sentence.
       Gnats, gnus, etc.\ all  % `\ ' makes an inter-word space.
       begin with G\@.         % \@ marks end-of-sentence punctuation.
You should check the spaces after periods when
reading your output to make sure you haven't
forgotten any special cases.  Generating an
ellipsis
       \ldots\               % `\ ' is needed after `\ldots' because TeX
                             % ignores spaces after command names like \ldots
                             % made from \ + letters.
                             %
                             % Note how a `%' character causes TeX to ignore
                             % the end of the input line, so these blank lines
                             % do not start a new paragraph.
                             %
with the right spacing around the periods requires
a special command.

\LaTeX\ interprets some common characters as
commands, so you must type special commands to
generate them.  These characters include the
following:
       \$ \& \% \# \{ and \}.

In printing, text is usually emphasized with an
       \emph{italic}
type style.

\begin{em}
   A long segment of text can also be emphasized
   in this way.  Text within such a segment can be
   given \emph{additional} emphasis.
\end{em}

It is sometimes necessary to prevent \LaTeX\ from
breaking a line where it might otherwise do so.
This may be at a space, as between the ``Mr.''\ and
``Jones'' in
       ``Mr.~Jones'',        % ~ produces an unbreakable interword space.
or within a word---especially when the word is a
symbol like
       \mbox{\emph{itemnum}}
that makes little sense when hyphenated across
lines.

Footnotes\footnote{This is an example of a footnote.}
pose no problem.

\LaTeX\ is good at typesetting mathematical formulas
like
       \( x-3y + z = 7 \)
or
       \( a_{1} > x^{2n} + y^{2n} > x' \)
or
       \( \ip{A}{B} = \sum_{i} a_{i} b_{i} \).
The spaces you type in a formula are
ignored.  Remember that a letter like
       $x$                   % $ ... $  and  \( ... \)  are equivalent
is a formula when it denotes a mathematical
symbol, and it should be typed as one.

\section{Displayed Text}

Text is displayed by indenting it from the left
margin.  Quotations are commonly displayed.  There
are short quotations
\begin{quote}
   This is a short quotation.  It consists of a
   single paragraph of text.  See how it is formatted.
\end{quote}
and longer ones.
\begin{quotation}
   This is a longer quotation.  It consists of two
   paragraphs of text, neither of which are
   particularly interesting.

   This is the second paragraph of the quotation.  It
   is just as dull as the first paragraph.
\end{quotation}
Another frequently-displayed structure is a list.
The following is an example of an \emph{itemized}
list.
\begin{itemize}
   \item This is the first item of an itemized list.
         Each item in the list is marked with a ``tick''.
         You don't have to worry about what kind of tick
         mark is used.

   \item This is the second item of the list.  It
         contains another list nested inside it.  The inner
         list is an \emph{enumerated} list.
         \begin{enumerate}
            \item This is the first item of an enumerated
                  list that is nested within the itemized list.

            \item This is the second item of the inner list.
                  \LaTeX\ allows you to nest lists deeper than
                  you really should.
         \end{enumerate}
         This is the rest of the second item of the outer
         list.  It is no more interesting than any other
         part of the item.
   \item This is the third item of the list.
\end{itemize}
You can even display poetry.
\begin{verse}
   There is an environment
    for verse \\             % The \\ command separates lines
   Whose features some poets % within a stanza.
   will curse.

                             % One or more blank lines separate stanzas.

   For instead of making\\
   Them do \emph{all} line breaking, \\
   It allows them to put too many words on a line when they'd rather be
   forced to be terse.
\end{verse}

Mathematical formulas may also be displayed.  A
displayed formula
is
one-line long; multiline
formulas require special formatting instructions.
   \[  \ip{\Gamma}{\psi'} = x'' + y^{2} + z_{i}^{n}\]
Don't start a paragraph with a displayed equation,
nor make one a paragraph by itself.

\end{document}               % End of document.

\section{Bench Huge Section 1493}
% synthetic bench section 1493
% This is a sample LaTeX input file.  (Version of 12 August 2004.)
%
% A '%' character causes TeX to ignore all remaining text on the line,
% and is used for comments like this one.

\documentclass{article}      % Specifies the document class

                             % The preamble begins here.
\title{An Example Document}  % Declares the document's title.
\author{Leslie Lamport}      % Declares the author's name.
\date{January 21, 1994}      % Deleting this command produces today's date.

\newcommand{\ip}[2]{(#1, #2)}
                             % Defines \ip{arg1}{arg2} to mean
                             % (arg1, arg2).

%\newcommand{\ip}[2]{\langle #1 | #2\rangle}
                             % This is an alternative definition of
                             % \ip that is commented out.

\begin{document}             % End of preamble and beginning of text.

\maketitle                   % Produces the title.

This is an example input file.  Comparing it with
the output it generates can show you how to
produce a simple document of your own.

\section{Ordinary Text}      % Produces section heading.  Lower-level
                             % sections are begun with similar
                             % \subsection and \subsubsection commands.

The ends  of words and sentences are marked
  by   spaces. It  doesn't matter how many
spaces    you type; one is as good as 100.  The
end of   a line counts as a space.

One   or more   blank lines denote the  end
of  a paragraph.

Since any number of consecutive spaces are treated
like a single one, the formatting of the input
file makes no difference to
      \LaTeX,                % The \LaTeX command generates the LaTeX logo.
but it makes a difference to you.  When you use
\LaTeX, making your input file as easy to read
as possible will be a great help as you write
your document and when you change it.  This sample
file shows how you can add comments to your own input
file.

Because printing is different from typewriting,
there are a number of things that you have to do
differently when preparing an input file than if
you were just typing the document directly.
Quotation marks like
       ``this''
have to be handled specially, as do quotes within
quotes:
       ``\,`this'            % \, separates the double and single quote.
        is what I just
        wrote, not  `that'\,''.

Dashes come in three sizes: an
       intra-word
dash, a medium dash for number ranges like
       1--2,
and a punctuation
       dash---like
this.

A sentence-ending space should be larger than the
space between words within a sentence.  You
sometimes have to type special commands in
conjunction with punctuation characters to get
this right, as in the following sentence.
       Gnats, gnus, etc.\ all  % `\ ' makes an inter-word space.
       begin with G\@.         % \@ marks end-of-sentence punctuation.
You should check the spaces after periods when
reading your output to make sure you haven't
forgotten any special cases.  Generating an
ellipsis
       \ldots\               % `\ ' is needed after `\ldots' because TeX
                             % ignores spaces after command names like \ldots
                             % made from \ + letters.
                             %
                             % Note how a `%' character causes TeX to ignore
                             % the end of the input line, so these blank lines
                             % do not start a new paragraph.
                             %
with the right spacing around the periods requires
a special command.

\LaTeX\ interprets some common characters as
commands, so you must type special commands to
generate them.  These characters include the
following:
       \$ \& \% \# \{ and \}.

In printing, text is usually emphasized with an
       \emph{italic}
type style.

\begin{em}
   A long segment of text can also be emphasized
   in this way.  Text within such a segment can be
   given \emph{additional} emphasis.
\end{em}

It is sometimes necessary to prevent \LaTeX\ from
breaking a line where it might otherwise do so.
This may be at a space, as between the ``Mr.''\ and
``Jones'' in
       ``Mr.~Jones'',        % ~ produces an unbreakable interword space.
or within a word---especially when the word is a
symbol like
       \mbox{\emph{itemnum}}
that makes little sense when hyphenated across
lines.

Footnotes\footnote{This is an example of a footnote.}
pose no problem.

\LaTeX\ is good at typesetting mathematical formulas
like
       \( x-3y + z = 7 \)
or
       \( a_{1} > x^{2n} + y^{2n} > x' \)
or
       \( \ip{A}{B} = \sum_{i} a_{i} b_{i} \).
The spaces you type in a formula are
ignored.  Remember that a letter like
       $x$                   % $ ... $  and  \( ... \)  are equivalent
is a formula when it denotes a mathematical
symbol, and it should be typed as one.

\section{Displayed Text}

Text is displayed by indenting it from the left
margin.  Quotations are commonly displayed.  There
are short quotations
\begin{quote}
   This is a short quotation.  It consists of a
   single paragraph of text.  See how it is formatted.
\end{quote}
and longer ones.
\begin{quotation}
   This is a longer quotation.  It consists of two
   paragraphs of text, neither of which are
   particularly interesting.

   This is the second paragraph of the quotation.  It
   is just as dull as the first paragraph.
\end{quotation}
Another frequently-displayed structure is a list.
The following is an example of an \emph{itemized}
list.
\begin{itemize}
   \item This is the first item of an itemized list.
         Each item in the list is marked with a ``tick''.
         You don't have to worry about what kind of tick
         mark is used.

   \item This is the second item of the list.  It
         contains another list nested inside it.  The inner
         list is an \emph{enumerated} list.
         \begin{enumerate}
            \item This is the first item of an enumerated
                  list that is nested within the itemized list.

            \item This is the second item of the inner list.
                  \LaTeX\ allows you to nest lists deeper than
                  you really should.
         \end{enumerate}
         This is the rest of the second item of the outer
         list.  It is no more interesting than any other
         part of the item.
   \item This is the third item of the list.
\end{itemize}
You can even display poetry.
\begin{verse}
   There is an environment
    for verse \\             % The \\ command separates lines
   Whose features some poets % within a stanza.
   will curse.

                             % One or more blank lines separate stanzas.

   For instead of making\\
   Them do \emph{all} line breaking, \\
   It allows them to put too many words on a line when they'd rather be
   forced to be terse.
\end{verse}

Mathematical formulas may also be displayed.  A
displayed formula
is
one-line long; multiline
formulas require special formatting instructions.
   \[  \ip{\Gamma}{\psi'} = x'' + y^{2} + z_{i}^{n}\]
Don't start a paragraph with a displayed equation,
nor make one a paragraph by itself.

\end{document}               % End of document.

\section{Bench Huge Section 1494}
% synthetic bench section 1494
% This is a sample LaTeX input file.  (Version of 12 August 2004.)
%
% A '%' character causes TeX to ignore all remaining text on the line,
% and is used for comments like this one.

\documentclass{article}      % Specifies the document class

                             % The preamble begins here.
\title{An Example Document}  % Declares the document's title.
\author{Leslie Lamport}      % Declares the author's name.
\date{January 21, 1994}      % Deleting this command produces today's date.

\newcommand{\ip}[2]{(#1, #2)}
                             % Defines \ip{arg1}{arg2} to mean
                             % (arg1, arg2).

%\newcommand{\ip}[2]{\langle #1 | #2\rangle}
                             % This is an alternative definition of
                             % \ip that is commented out.

\begin{document}             % End of preamble and beginning of text.

\maketitle                   % Produces the title.

This is an example input file.  Comparing it with
the output it generates can show you how to
produce a simple document of your own.

\section{Ordinary Text}      % Produces section heading.  Lower-level
                             % sections are begun with similar
                             % \subsection and \subsubsection commands.

The ends  of words and sentences are marked
  by   spaces. It  doesn't matter how many
spaces    you type; one is as good as 100.  The
end of   a line counts as a space.

One   or more   blank lines denote the  end
of  a paragraph.

Since any number of consecutive spaces are treated
like a single one, the formatting of the input
file makes no difference to
      \LaTeX,                % The \LaTeX command generates the LaTeX logo.
but it makes a difference to you.  When you use
\LaTeX, making your input file as easy to read
as possible will be a great help as you write
your document and when you change it.  This sample
file shows how you can add comments to your own input
file.

Because printing is different from typewriting,
there are a number of things that you have to do
differently when preparing an input file than if
you were just typing the document directly.
Quotation marks like
       ``this''
have to be handled specially, as do quotes within
quotes:
       ``\,`this'            % \, separates the double and single quote.
        is what I just
        wrote, not  `that'\,''.

Dashes come in three sizes: an
       intra-word
dash, a medium dash for number ranges like
       1--2,
and a punctuation
       dash---like
this.

A sentence-ending space should be larger than the
space between words within a sentence.  You
sometimes have to type special commands in
conjunction with punctuation characters to get
this right, as in the following sentence.
       Gnats, gnus, etc.\ all  % `\ ' makes an inter-word space.
       begin with G\@.         % \@ marks end-of-sentence punctuation.
You should check the spaces after periods when
reading your output to make sure you haven't
forgotten any special cases.  Generating an
ellipsis
       \ldots\               % `\ ' is needed after `\ldots' because TeX
                             % ignores spaces after command names like \ldots
                             % made from \ + letters.
                             %
                             % Note how a `%' character causes TeX to ignore
                             % the end of the input line, so these blank lines
                             % do not start a new paragraph.
                             %
with the right spacing around the periods requires
a special command.

\LaTeX\ interprets some common characters as
commands, so you must type special commands to
generate them.  These characters include the
following:
       \$ \& \% \# \{ and \}.

In printing, text is usually emphasized with an
       \emph{italic}
type style.

\begin{em}
   A long segment of text can also be emphasized
   in this way.  Text within such a segment can be
   given \emph{additional} emphasis.
\end{em}

It is sometimes necessary to prevent \LaTeX\ from
breaking a line where it might otherwise do so.
This may be at a space, as between the ``Mr.''\ and
``Jones'' in
       ``Mr.~Jones'',        % ~ produces an unbreakable interword space.
or within a word---especially when the word is a
symbol like
       \mbox{\emph{itemnum}}
that makes little sense when hyphenated across
lines.

Footnotes\footnote{This is an example of a footnote.}
pose no problem.

\LaTeX\ is good at typesetting mathematical formulas
like
       \( x-3y + z = 7 \)
or
       \( a_{1} > x^{2n} + y^{2n} > x' \)
or
       \( \ip{A}{B} = \sum_{i} a_{i} b_{i} \).
The spaces you type in a formula are
ignored.  Remember that a letter like
       $x$                   % $ ... $  and  \( ... \)  are equivalent
is a formula when it denotes a mathematical
symbol, and it should be typed as one.

\section{Displayed Text}

Text is displayed by indenting it from the left
margin.  Quotations are commonly displayed.  There
are short quotations
\begin{quote}
   This is a short quotation.  It consists of a
   single paragraph of text.  See how it is formatted.
\end{quote}
and longer ones.
\begin{quotation}
   This is a longer quotation.  It consists of two
   paragraphs of text, neither of which are
   particularly interesting.

   This is the second paragraph of the quotation.  It
   is just as dull as the first paragraph.
\end{quotation}
Another frequently-displayed structure is a list.
The following is an example of an \emph{itemized}
list.
\begin{itemize}
   \item This is the first item of an itemized list.
         Each item in the list is marked with a ``tick''.
         You don't have to worry about what kind of tick
         mark is used.

   \item This is the second item of the list.  It
         contains another list nested inside it.  The inner
         list is an \emph{enumerated} list.
         \begin{enumerate}
            \item This is the first item of an enumerated
                  list that is nested within the itemized list.

            \item This is the second item of the inner list.
                  \LaTeX\ allows you to nest lists deeper than
                  you really should.
         \end{enumerate}
         This is the rest of the second item of the outer
         list.  It is no more interesting than any other
         part of the item.
   \item This is the third item of the list.
\end{itemize}
You can even display poetry.
\begin{verse}
   There is an environment
    for verse \\             % The \\ command separates lines
   Whose features some poets % within a stanza.
   will curse.

                             % One or more blank lines separate stanzas.

   For instead of making\\
   Them do \emph{all} line breaking, \\
   It allows them to put too many words on a line when they'd rather be
   forced to be terse.
\end{verse}

Mathematical formulas may also be displayed.  A
displayed formula
is
one-line long; multiline
formulas require special formatting instructions.
   \[  \ip{\Gamma}{\psi'} = x'' + y^{2} + z_{i}^{n}\]
Don't start a paragraph with a displayed equation,
nor make one a paragraph by itself.

\end{document}               % End of document.

\section{Bench Huge Section 1495}
% synthetic bench section 1495
% This is a sample LaTeX input file.  (Version of 12 August 2004.)
%
% A '%' character causes TeX to ignore all remaining text on the line,
% and is used for comments like this one.

\documentclass{article}      % Specifies the document class

                             % The preamble begins here.
\title{An Example Document}  % Declares the document's title.
\author{Leslie Lamport}      % Declares the author's name.
\date{January 21, 1994}      % Deleting this command produces today's date.

\newcommand{\ip}[2]{(#1, #2)}
                             % Defines \ip{arg1}{arg2} to mean
                             % (arg1, arg2).

%\newcommand{\ip}[2]{\langle #1 | #2\rangle}
                             % This is an alternative definition of
                             % \ip that is commented out.

\begin{document}             % End of preamble and beginning of text.

\maketitle                   % Produces the title.

This is an example input file.  Comparing it with
the output it generates can show you how to
produce a simple document of your own.

\section{Ordinary Text}      % Produces section heading.  Lower-level
                             % sections are begun with similar
                             % \subsection and \subsubsection commands.

The ends  of words and sentences are marked
  by   spaces. It  doesn't matter how many
spaces    you type; one is as good as 100.  The
end of   a line counts as a space.

One   or more   blank lines denote the  end
of  a paragraph.

Since any number of consecutive spaces are treated
like a single one, the formatting of the input
file makes no difference to
      \LaTeX,                % The \LaTeX command generates the LaTeX logo.
but it makes a difference to you.  When you use
\LaTeX, making your input file as easy to read
as possible will be a great help as you write
your document and when you change it.  This sample
file shows how you can add comments to your own input
file.

Because printing is different from typewriting,
there are a number of things that you have to do
differently when preparing an input file than if
you were just typing the document directly.
Quotation marks like
       ``this''
have to be handled specially, as do quotes within
quotes:
       ``\,`this'            % \, separates the double and single quote.
        is what I just
        wrote, not  `that'\,''.

Dashes come in three sizes: an
       intra-word
dash, a medium dash for number ranges like
       1--2,
and a punctuation
       dash---like
this.

A sentence-ending space should be larger than the
space between words within a sentence.  You
sometimes have to type special commands in
conjunction with punctuation characters to get
this right, as in the following sentence.
       Gnats, gnus, etc.\ all  % `\ ' makes an inter-word space.
       begin with G\@.         % \@ marks end-of-sentence punctuation.
You should check the spaces after periods when
reading your output to make sure you haven't
forgotten any special cases.  Generating an
ellipsis
       \ldots\               % `\ ' is needed after `\ldots' because TeX
                             % ignores spaces after command names like \ldots
                             % made from \ + letters.
                             %
                             % Note how a `%' character causes TeX to ignore
                             % the end of the input line, so these blank lines
                             % do not start a new paragraph.
                             %
with the right spacing around the periods requires
a special command.

\LaTeX\ interprets some common characters as
commands, so you must type special commands to
generate them.  These characters include the
following:
       \$ \& \% \# \{ and \}.

In printing, text is usually emphasized with an
       \emph{italic}
type style.

\begin{em}
   A long segment of text can also be emphasized
   in this way.  Text within such a segment can be
   given \emph{additional} emphasis.
\end{em}

It is sometimes necessary to prevent \LaTeX\ from
breaking a line where it might otherwise do so.
This may be at a space, as between the ``Mr.''\ and
``Jones'' in
       ``Mr.~Jones'',        % ~ produces an unbreakable interword space.
or within a word---especially when the word is a
symbol like
       \mbox{\emph{itemnum}}
that makes little sense when hyphenated across
lines.

Footnotes\footnote{This is an example of a footnote.}
pose no problem.

\LaTeX\ is good at typesetting mathematical formulas
like
       \( x-3y + z = 7 \)
or
       \( a_{1} > x^{2n} + y^{2n} > x' \)
or
       \( \ip{A}{B} = \sum_{i} a_{i} b_{i} \).
The spaces you type in a formula are
ignored.  Remember that a letter like
       $x$                   % $ ... $  and  \( ... \)  are equivalent
is a formula when it denotes a mathematical
symbol, and it should be typed as one.

\section{Displayed Text}

Text is displayed by indenting it from the left
margin.  Quotations are commonly displayed.  There
are short quotations
\begin{quote}
   This is a short quotation.  It consists of a
   single paragraph of text.  See how it is formatted.
\end{quote}
and longer ones.
\begin{quotation}
   This is a longer quotation.  It consists of two
   paragraphs of text, neither of which are
   particularly interesting.

   This is the second paragraph of the quotation.  It
   is just as dull as the first paragraph.
\end{quotation}
Another frequently-displayed structure is a list.
The following is an example of an \emph{itemized}
list.
\begin{itemize}
   \item This is the first item of an itemized list.
         Each item in the list is marked with a ``tick''.
         You don't have to worry about what kind of tick
         mark is used.

   \item This is the second item of the list.  It
         contains another list nested inside it.  The inner
         list is an \emph{enumerated} list.
         \begin{enumerate}
            \item This is the first item of an enumerated
                  list that is nested within the itemized list.

            \item This is the second item of the inner list.
                  \LaTeX\ allows you to nest lists deeper than
                  you really should.
         \end{enumerate}
         This is the rest of the second item of the outer
         list.  It is no more interesting than any other
         part of the item.
   \item This is the third item of the list.
\end{itemize}
You can even display poetry.
\begin{verse}
   There is an environment
    for verse \\             % The \\ command separates lines
   Whose features some poets % within a stanza.
   will curse.

                             % One or more blank lines separate stanzas.

   For instead of making\\
   Them do \emph{all} line breaking, \\
   It allows them to put too many words on a line when they'd rather be
   forced to be terse.
\end{verse}

Mathematical formulas may also be displayed.  A
displayed formula
is
one-line long; multiline
formulas require special formatting instructions.
   \[  \ip{\Gamma}{\psi'} = x'' + y^{2} + z_{i}^{n}\]
Don't start a paragraph with a displayed equation,
nor make one a paragraph by itself.

\end{document}               % End of document.

\section{Bench Huge Section 1496}
% synthetic bench section 1496
% This is a sample LaTeX input file.  (Version of 12 August 2004.)
%
% A '%' character causes TeX to ignore all remaining text on the line,
% and is used for comments like this one.

\documentclass{article}      % Specifies the document class

                             % The preamble begins here.
\title{An Example Document}  % Declares the document's title.
\author{Leslie Lamport}      % Declares the author's name.
\date{January 21, 1994}      % Deleting this command produces today's date.

\newcommand{\ip}[2]{(#1, #2)}
                             % Defines \ip{arg1}{arg2} to mean
                             % (arg1, arg2).

%\newcommand{\ip}[2]{\langle #1 | #2\rangle}
                             % This is an alternative definition of
                             % \ip that is commented out.

\begin{document}             % End of preamble and beginning of text.

\maketitle                   % Produces the title.

This is an example input file.  Comparing it with
the output it generates can show you how to
produce a simple document of your own.

\section{Ordinary Text}      % Produces section heading.  Lower-level
                             % sections are begun with similar
                             % \subsection and \subsubsection commands.

The ends  of words and sentences are marked
  by   spaces. It  doesn't matter how many
spaces    you type; one is as good as 100.  The
end of   a line counts as a space.

One   or more   blank lines denote the  end
of  a paragraph.

Since any number of consecutive spaces are treated
like a single one, the formatting of the input
file makes no difference to
      \LaTeX,                % The \LaTeX command generates the LaTeX logo.
but it makes a difference to you.  When you use
\LaTeX, making your input file as easy to read
as possible will be a great help as you write
your document and when you change it.  This sample
file shows how you can add comments to your own input
file.

Because printing is different from typewriting,
there are a number of things that you have to do
differently when preparing an input file than if
you were just typing the document directly.
Quotation marks like
       ``this''
have to be handled specially, as do quotes within
quotes:
       ``\,`this'            % \, separates the double and single quote.
        is what I just
        wrote, not  `that'\,''.

Dashes come in three sizes: an
       intra-word
dash, a medium dash for number ranges like
       1--2,
and a punctuation
       dash---like
this.

A sentence-ending space should be larger than the
space between words within a sentence.  You
sometimes have to type special commands in
conjunction with punctuation characters to get
this right, as in the following sentence.
       Gnats, gnus, etc.\ all  % `\ ' makes an inter-word space.
       begin with G\@.         % \@ marks end-of-sentence punctuation.
You should check the spaces after periods when
reading your output to make sure you haven't
forgotten any special cases.  Generating an
ellipsis
       \ldots\               % `\ ' is needed after `\ldots' because TeX
                             % ignores spaces after command names like \ldots
                             % made from \ + letters.
                             %
                             % Note how a `%' character causes TeX to ignore
                             % the end of the input line, so these blank lines
                             % do not start a new paragraph.
                             %
with the right spacing around the periods requires
a special command.

\LaTeX\ interprets some common characters as
commands, so you must type special commands to
generate them.  These characters include the
following:
       \$ \& \% \# \{ and \}.

In printing, text is usually emphasized with an
       \emph{italic}
type style.

\begin{em}
   A long segment of text can also be emphasized
   in this way.  Text within such a segment can be
   given \emph{additional} emphasis.
\end{em}

It is sometimes necessary to prevent \LaTeX\ from
breaking a line where it might otherwise do so.
This may be at a space, as between the ``Mr.''\ and
``Jones'' in
       ``Mr.~Jones'',        % ~ produces an unbreakable interword space.
or within a word---especially when the word is a
symbol like
       \mbox{\emph{itemnum}}
that makes little sense when hyphenated across
lines.

Footnotes\footnote{This is an example of a footnote.}
pose no problem.

\LaTeX\ is good at typesetting mathematical formulas
like
       \( x-3y + z = 7 \)
or
       \( a_{1} > x^{2n} + y^{2n} > x' \)
or
       \( \ip{A}{B} = \sum_{i} a_{i} b_{i} \).
The spaces you type in a formula are
ignored.  Remember that a letter like
       $x$                   % $ ... $  and  \( ... \)  are equivalent
is a formula when it denotes a mathematical
symbol, and it should be typed as one.

\section{Displayed Text}

Text is displayed by indenting it from the left
margin.  Quotations are commonly displayed.  There
are short quotations
\begin{quote}
   This is a short quotation.  It consists of a
   single paragraph of text.  See how it is formatted.
\end{quote}
and longer ones.
\begin{quotation}
   This is a longer quotation.  It consists of two
   paragraphs of text, neither of which are
   particularly interesting.

   This is the second paragraph of the quotation.  It
   is just as dull as the first paragraph.
\end{quotation}
Another frequently-displayed structure is a list.
The following is an example of an \emph{itemized}
list.
\begin{itemize}
   \item This is the first item of an itemized list.
         Each item in the list is marked with a ``tick''.
         You don't have to worry about what kind of tick
         mark is used.

   \item This is the second item of the list.  It
         contains another list nested inside it.  The inner
         list is an \emph{enumerated} list.
         \begin{enumerate}
            \item This is the first item of an enumerated
                  list that is nested within the itemized list.

            \item This is the second item of the inner list.
                  \LaTeX\ allows you to nest lists deeper than
                  you really should.
         \end{enumerate}
         This is the rest of the second item of the outer
         list.  It is no more interesting than any other
         part of the item.
   \item This is the third item of the list.
\end{itemize}
You can even display poetry.
\begin{verse}
   There is an environment
    for verse \\             % The \\ command separates lines
   Whose features some poets % within a stanza.
   will curse.

                             % One or more blank lines separate stanzas.

   For instead of making\\
   Them do \emph{all} line breaking, \\
   It allows them to put too many words on a line when they'd rather be
   forced to be terse.
\end{verse}

Mathematical formulas may also be displayed.  A
displayed formula
is
one-line long; multiline
formulas require special formatting instructions.
   \[  \ip{\Gamma}{\psi'} = x'' + y^{2} + z_{i}^{n}\]
Don't start a paragraph with a displayed equation,
nor make one a paragraph by itself.

\end{document}               % End of document.

\section{Bench Huge Section 1497}
% synthetic bench section 1497
% This is a sample LaTeX input file.  (Version of 12 August 2004.)
%
% A '%' character causes TeX to ignore all remaining text on the line,
% and is used for comments like this one.

\documentclass{article}      % Specifies the document class

                             % The preamble begins here.
\title{An Example Document}  % Declares the document's title.
\author{Leslie Lamport}      % Declares the author's name.
\date{January 21, 1994}      % Deleting this command produces today's date.

\newcommand{\ip}[2]{(#1, #2)}
                             % Defines \ip{arg1}{arg2} to mean
                             % (arg1, arg2).

%\newcommand{\ip}[2]{\langle #1 | #2\rangle}
                             % This is an alternative definition of
                             % \ip that is commented out.

\begin{document}             % End of preamble and beginning of text.

\maketitle                   % Produces the title.

This is an example input file.  Comparing it with
the output it generates can show you how to
produce a simple document of your own.

\section{Ordinary Text}      % Produces section heading.  Lower-level
                             % sections are begun with similar
                             % \subsection and \subsubsection commands.

The ends  of words and sentences are marked
  by   spaces. It  doesn't matter how many
spaces    you type; one is as good as 100.  The
end of   a line counts as a space.

One   or more   blank lines denote the  end
of  a paragraph.

Since any number of consecutive spaces are treated
like a single one, the formatting of the input
file makes no difference to
      \LaTeX,                % The \LaTeX command generates the LaTeX logo.
but it makes a difference to you.  When you use
\LaTeX, making your input file as easy to read
as possible will be a great help as you write
your document and when you change it.  This sample
file shows how you can add comments to your own input
file.

Because printing is different from typewriting,
there are a number of things that you have to do
differently when preparing an input file than if
you were just typing the document directly.
Quotation marks like
       ``this''
have to be handled specially, as do quotes within
quotes:
       ``\,`this'            % \, separates the double and single quote.
        is what I just
        wrote, not  `that'\,''.

Dashes come in three sizes: an
       intra-word
dash, a medium dash for number ranges like
       1--2,
and a punctuation
       dash---like
this.

A sentence-ending space should be larger than the
space between words within a sentence.  You
sometimes have to type special commands in
conjunction with punctuation characters to get
this right, as in the following sentence.
       Gnats, gnus, etc.\ all  % `\ ' makes an inter-word space.
       begin with G\@.         % \@ marks end-of-sentence punctuation.
You should check the spaces after periods when
reading your output to make sure you haven't
forgotten any special cases.  Generating an
ellipsis
       \ldots\               % `\ ' is needed after `\ldots' because TeX
                             % ignores spaces after command names like \ldots
                             % made from \ + letters.
                             %
                             % Note how a `%' character causes TeX to ignore
                             % the end of the input line, so these blank lines
                             % do not start a new paragraph.
                             %
with the right spacing around the periods requires
a special command.

\LaTeX\ interprets some common characters as
commands, so you must type special commands to
generate them.  These characters include the
following:
       \$ \& \% \# \{ and \}.

In printing, text is usually emphasized with an
       \emph{italic}
type style.

\begin{em}
   A long segment of text can also be emphasized
   in this way.  Text within such a segment can be
   given \emph{additional} emphasis.
\end{em}

It is sometimes necessary to prevent \LaTeX\ from
breaking a line where it might otherwise do so.
This may be at a space, as between the ``Mr.''\ and
``Jones'' in
       ``Mr.~Jones'',        % ~ produces an unbreakable interword space.
or within a word---especially when the word is a
symbol like
       \mbox{\emph{itemnum}}
that makes little sense when hyphenated across
lines.

Footnotes\footnote{This is an example of a footnote.}
pose no problem.

\LaTeX\ is good at typesetting mathematical formulas
like
       \( x-3y + z = 7 \)
or
       \( a_{1} > x^{2n} + y^{2n} > x' \)
or
       \( \ip{A}{B} = \sum_{i} a_{i} b_{i} \).
The spaces you type in a formula are
ignored.  Remember that a letter like
       $x$                   % $ ... $  and  \( ... \)  are equivalent
is a formula when it denotes a mathematical
symbol, and it should be typed as one.

\section{Displayed Text}

Text is displayed by indenting it from the left
margin.  Quotations are commonly displayed.  There
are short quotations
\begin{quote}
   This is a short quotation.  It consists of a
   single paragraph of text.  See how it is formatted.
\end{quote}
and longer ones.
\begin{quotation}
   This is a longer quotation.  It consists of two
   paragraphs of text, neither of which are
   particularly interesting.

   This is the second paragraph of the quotation.  It
   is just as dull as the first paragraph.
\end{quotation}
Another frequently-displayed structure is a list.
The following is an example of an \emph{itemized}
list.
\begin{itemize}
   \item This is the first item of an itemized list.
         Each item in the list is marked with a ``tick''.
         You don't have to worry about what kind of tick
         mark is used.

   \item This is the second item of the list.  It
         contains another list nested inside it.  The inner
         list is an \emph{enumerated} list.
         \begin{enumerate}
            \item This is the first item of an enumerated
                  list that is nested within the itemized list.

            \item This is the second item of the inner list.
                  \LaTeX\ allows you to nest lists deeper than
                  you really should.
         \end{enumerate}
         This is the rest of the second item of the outer
         list.  It is no more interesting than any other
         part of the item.
   \item This is the third item of the list.
\end{itemize}
You can even display poetry.
\begin{verse}
   There is an environment
    for verse \\             % The \\ command separates lines
   Whose features some poets % within a stanza.
   will curse.

                             % One or more blank lines separate stanzas.

   For instead of making\\
   Them do \emph{all} line breaking, \\
   It allows them to put too many words on a line when they'd rather be
   forced to be terse.
\end{verse}

Mathematical formulas may also be displayed.  A
displayed formula
is
one-line long; multiline
formulas require special formatting instructions.
   \[  \ip{\Gamma}{\psi'} = x'' + y^{2} + z_{i}^{n}\]
Don't start a paragraph with a displayed equation,
nor make one a paragraph by itself.

\end{document}               % End of document.

\section{Bench Huge Section 1498}
% synthetic bench section 1498
% This is a sample LaTeX input file.  (Version of 12 August 2004.)
%
% A '%' character causes TeX to ignore all remaining text on the line,
% and is used for comments like this one.

\documentclass{article}      % Specifies the document class

                             % The preamble begins here.
\title{An Example Document}  % Declares the document's title.
\author{Leslie Lamport}      % Declares the author's name.
\date{January 21, 1994}      % Deleting this command produces today's date.

\newcommand{\ip}[2]{(#1, #2)}
                             % Defines \ip{arg1}{arg2} to mean
                             % (arg1, arg2).

%\newcommand{\ip}[2]{\langle #1 | #2\rangle}
                             % This is an alternative definition of
                             % \ip that is commented out.

\begin{document}             % End of preamble and beginning of text.

\maketitle                   % Produces the title.

This is an example input file.  Comparing it with
the output it generates can show you how to
produce a simple document of your own.

\section{Ordinary Text}      % Produces section heading.  Lower-level
                             % sections are begun with similar
                             % \subsection and \subsubsection commands.

The ends  of words and sentences are marked
  by   spaces. It  doesn't matter how many
spaces    you type; one is as good as 100.  The
end of   a line counts as a space.

One   or more   blank lines denote the  end
of  a paragraph.

Since any number of consecutive spaces are treated
like a single one, the formatting of the input
file makes no difference to
      \LaTeX,                % The \LaTeX command generates the LaTeX logo.
but it makes a difference to you.  When you use
\LaTeX, making your input file as easy to read
as possible will be a great help as you write
your document and when you change it.  This sample
file shows how you can add comments to your own input
file.

Because printing is different from typewriting,
there are a number of things that you have to do
differently when preparing an input file than if
you were just typing the document directly.
Quotation marks like
       ``this''
have to be handled specially, as do quotes within
quotes:
       ``\,`this'            % \, separates the double and single quote.
        is what I just
        wrote, not  `that'\,''.

Dashes come in three sizes: an
       intra-word
dash, a medium dash for number ranges like
       1--2,
and a punctuation
       dash---like
this.

A sentence-ending space should be larger than the
space between words within a sentence.  You
sometimes have to type special commands in
conjunction with punctuation characters to get
this right, as in the following sentence.
       Gnats, gnus, etc.\ all  % `\ ' makes an inter-word space.
       begin with G\@.         % \@ marks end-of-sentence punctuation.
You should check the spaces after periods when
reading your output to make sure you haven't
forgotten any special cases.  Generating an
ellipsis
       \ldots\               % `\ ' is needed after `\ldots' because TeX
                             % ignores spaces after command names like \ldots
                             % made from \ + letters.
                             %
                             % Note how a `%' character causes TeX to ignore
                             % the end of the input line, so these blank lines
                             % do not start a new paragraph.
                             %
with the right spacing around the periods requires
a special command.

\LaTeX\ interprets some common characters as
commands, so you must type special commands to
generate them.  These characters include the
following:
       \$ \& \% \# \{ and \}.

In printing, text is usually emphasized with an
       \emph{italic}
type style.

\begin{em}
   A long segment of text can also be emphasized
   in this way.  Text within such a segment can be
   given \emph{additional} emphasis.
\end{em}

It is sometimes necessary to prevent \LaTeX\ from
breaking a line where it might otherwise do so.
This may be at a space, as between the ``Mr.''\ and
``Jones'' in
       ``Mr.~Jones'',        % ~ produces an unbreakable interword space.
or within a word---especially when the word is a
symbol like
       \mbox{\emph{itemnum}}
that makes little sense when hyphenated across
lines.

Footnotes\footnote{This is an example of a footnote.}
pose no problem.

\LaTeX\ is good at typesetting mathematical formulas
like
       \( x-3y + z = 7 \)
or
       \( a_{1} > x^{2n} + y^{2n} > x' \)
or
       \( \ip{A}{B} = \sum_{i} a_{i} b_{i} \).
The spaces you type in a formula are
ignored.  Remember that a letter like
       $x$                   % $ ... $  and  \( ... \)  are equivalent
is a formula when it denotes a mathematical
symbol, and it should be typed as one.

\section{Displayed Text}

Text is displayed by indenting it from the left
margin.  Quotations are commonly displayed.  There
are short quotations
\begin{quote}
   This is a short quotation.  It consists of a
   single paragraph of text.  See how it is formatted.
\end{quote}
and longer ones.
\begin{quotation}
   This is a longer quotation.  It consists of two
   paragraphs of text, neither of which are
   particularly interesting.

   This is the second paragraph of the quotation.  It
   is just as dull as the first paragraph.
\end{quotation}
Another frequently-displayed structure is a list.
The following is an example of an \emph{itemized}
list.
\begin{itemize}
   \item This is the first item of an itemized list.
         Each item in the list is marked with a ``tick''.
         You don't have to worry about what kind of tick
         mark is used.

   \item This is the second item of the list.  It
         contains another list nested inside it.  The inner
         list is an \emph{enumerated} list.
         \begin{enumerate}
            \item This is the first item of an enumerated
                  list that is nested within the itemized list.

            \item This is the second item of the inner list.
                  \LaTeX\ allows you to nest lists deeper than
                  you really should.
         \end{enumerate}
         This is the rest of the second item of the outer
         list.  It is no more interesting than any other
         part of the item.
   \item This is the third item of the list.
\end{itemize}
You can even display poetry.
\begin{verse}
   There is an environment
    for verse \\             % The \\ command separates lines
   Whose features some poets % within a stanza.
   will curse.

                             % One or more blank lines separate stanzas.

   For instead of making\\
   Them do \emph{all} line breaking, \\
   It allows them to put too many words on a line when they'd rather be
   forced to be terse.
\end{verse}

Mathematical formulas may also be displayed.  A
displayed formula
is
one-line long; multiline
formulas require special formatting instructions.
   \[  \ip{\Gamma}{\psi'} = x'' + y^{2} + z_{i}^{n}\]
Don't start a paragraph with a displayed equation,
nor make one a paragraph by itself.

\end{document}               % End of document.

\section{Bench Huge Section 1499}
% synthetic bench section 1499
% This is a sample LaTeX input file.  (Version of 12 August 2004.)
%
% A '%' character causes TeX to ignore all remaining text on the line,
% and is used for comments like this one.

\documentclass{article}      % Specifies the document class

                             % The preamble begins here.
\title{An Example Document}  % Declares the document's title.
\author{Leslie Lamport}      % Declares the author's name.
\date{January 21, 1994}      % Deleting this command produces today's date.

\newcommand{\ip}[2]{(#1, #2)}
                             % Defines \ip{arg1}{arg2} to mean
                             % (arg1, arg2).

%\newcommand{\ip}[2]{\langle #1 | #2\rangle}
                             % This is an alternative definition of
                             % \ip that is commented out.

\begin{document}             % End of preamble and beginning of text.

\maketitle                   % Produces the title.

This is an example input file.  Comparing it with
the output it generates can show you how to
produce a simple document of your own.

\section{Ordinary Text}      % Produces section heading.  Lower-level
                             % sections are begun with similar
                             % \subsection and \subsubsection commands.

The ends  of words and sentences are marked
  by   spaces. It  doesn't matter how many
spaces    you type; one is as good as 100.  The
end of   a line counts as a space.

One   or more   blank lines denote the  end
of  a paragraph.

Since any number of consecutive spaces are treated
like a single one, the formatting of the input
file makes no difference to
      \LaTeX,                % The \LaTeX command generates the LaTeX logo.
but it makes a difference to you.  When you use
\LaTeX, making your input file as easy to read
as possible will be a great help as you write
your document and when you change it.  This sample
file shows how you can add comments to your own input
file.

Because printing is different from typewriting,
there are a number of things that you have to do
differently when preparing an input file than if
you were just typing the document directly.
Quotation marks like
       ``this''
have to be handled specially, as do quotes within
quotes:
       ``\,`this'            % \, separates the double and single quote.
        is what I just
        wrote, not  `that'\,''.

Dashes come in three sizes: an
       intra-word
dash, a medium dash for number ranges like
       1--2,
and a punctuation
       dash---like
this.

A sentence-ending space should be larger than the
space between words within a sentence.  You
sometimes have to type special commands in
conjunction with punctuation characters to get
this right, as in the following sentence.
       Gnats, gnus, etc.\ all  % `\ ' makes an inter-word space.
       begin with G\@.         % \@ marks end-of-sentence punctuation.
You should check the spaces after periods when
reading your output to make sure you haven't
forgotten any special cases.  Generating an
ellipsis
       \ldots\               % `\ ' is needed after `\ldots' because TeX
                             % ignores spaces after command names like \ldots
                             % made from \ + letters.
                             %
                             % Note how a `%' character causes TeX to ignore
                             % the end of the input line, so these blank lines
                             % do not start a new paragraph.
                             %
with the right spacing around the periods requires
a special command.

\LaTeX\ interprets some common characters as
commands, so you must type special commands to
generate them.  These characters include the
following:
       \$ \& \% \# \{ and \}.

In printing, text is usually emphasized with an
       \emph{italic}
type style.

\begin{em}
   A long segment of text can also be emphasized
   in this way.  Text within such a segment can be
   given \emph{additional} emphasis.
\end{em}

It is sometimes necessary to prevent \LaTeX\ from
breaking a line where it might otherwise do so.
This may be at a space, as between the ``Mr.''\ and
``Jones'' in
       ``Mr.~Jones'',        % ~ produces an unbreakable interword space.
or within a word---especially when the word is a
symbol like
       \mbox{\emph{itemnum}}
that makes little sense when hyphenated across
lines.

Footnotes\footnote{This is an example of a footnote.}
pose no problem.

\LaTeX\ is good at typesetting mathematical formulas
like
       \( x-3y + z = 7 \)
or
       \( a_{1} > x^{2n} + y^{2n} > x' \)
or
       \( \ip{A}{B} = \sum_{i} a_{i} b_{i} \).
The spaces you type in a formula are
ignored.  Remember that a letter like
       $x$                   % $ ... $  and  \( ... \)  are equivalent
is a formula when it denotes a mathematical
symbol, and it should be typed as one.

\section{Displayed Text}

Text is displayed by indenting it from the left
margin.  Quotations are commonly displayed.  There
are short quotations
\begin{quote}
   This is a short quotation.  It consists of a
   single paragraph of text.  See how it is formatted.
\end{quote}
and longer ones.
\begin{quotation}
   This is a longer quotation.  It consists of two
   paragraphs of text, neither of which are
   particularly interesting.

   This is the second paragraph of the quotation.  It
   is just as dull as the first paragraph.
\end{quotation}
Another frequently-displayed structure is a list.
The following is an example of an \emph{itemized}
list.
\begin{itemize}
   \item This is the first item of an itemized list.
         Each item in the list is marked with a ``tick''.
         You don't have to worry about what kind of tick
         mark is used.

   \item This is the second item of the list.  It
         contains another list nested inside it.  The inner
         list is an \emph{enumerated} list.
         \begin{enumerate}
            \item This is the first item of an enumerated
                  list that is nested within the itemized list.

            \item This is the second item of the inner list.
                  \LaTeX\ allows you to nest lists deeper than
                  you really should.
         \end{enumerate}
         This is the rest of the second item of the outer
         list.  It is no more interesting than any other
         part of the item.
   \item This is the third item of the list.
\end{itemize}
You can even display poetry.
\begin{verse}
   There is an environment
    for verse \\             % The \\ command separates lines
   Whose features some poets % within a stanza.
   will curse.

                             % One or more blank lines separate stanzas.

   For instead of making\\
   Them do \emph{all} line breaking, \\
   It allows them to put too many words on a line when they'd rather be
   forced to be terse.
\end{verse}

Mathematical formulas may also be displayed.  A
displayed formula
is
one-line long; multiline
formulas require special formatting instructions.
   \[  \ip{\Gamma}{\psi'} = x'' + y^{2} + z_{i}^{n}\]
Don't start a paragraph with a displayed equation,
nor make one a paragraph by itself.

\end{document}               % End of document.

\section{Bench Huge Section 1500}
% synthetic bench section 1500
% This is a sample LaTeX input file.  (Version of 12 August 2004.)
%
% A '%' character causes TeX to ignore all remaining text on the line,
% and is used for comments like this one.

\documentclass{article}      % Specifies the document class

                             % The preamble begins here.
\title{An Example Document}  % Declares the document's title.
\author{Leslie Lamport}      % Declares the author's name.
\date{January 21, 1994}      % Deleting this command produces today's date.

\newcommand{\ip}[2]{(#1, #2)}
                             % Defines \ip{arg1}{arg2} to mean
                             % (arg1, arg2).

%\newcommand{\ip}[2]{\langle #1 | #2\rangle}
                             % This is an alternative definition of
                             % \ip that is commented out.

\begin{document}             % End of preamble and beginning of text.

\maketitle                   % Produces the title.

This is an example input file.  Comparing it with
the output it generates can show you how to
produce a simple document of your own.

\section{Ordinary Text}      % Produces section heading.  Lower-level
                             % sections are begun with similar
                             % \subsection and \subsubsection commands.

The ends  of words and sentences are marked
  by   spaces. It  doesn't matter how many
spaces    you type; one is as good as 100.  The
end of   a line counts as a space.

One   or more   blank lines denote the  end
of  a paragraph.

Since any number of consecutive spaces are treated
like a single one, the formatting of the input
file makes no difference to
      \LaTeX,                % The \LaTeX command generates the LaTeX logo.
but it makes a difference to you.  When you use
\LaTeX, making your input file as easy to read
as possible will be a great help as you write
your document and when you change it.  This sample
file shows how you can add comments to your own input
file.

Because printing is different from typewriting,
there are a number of things that you have to do
differently when preparing an input file than if
you were just typing the document directly.
Quotation marks like
       ``this''
have to be handled specially, as do quotes within
quotes:
       ``\,`this'            % \, separates the double and single quote.
        is what I just
        wrote, not  `that'\,''.

Dashes come in three sizes: an
       intra-word
dash, a medium dash for number ranges like
       1--2,
and a punctuation
       dash---like
this.

A sentence-ending space should be larger than the
space between words within a sentence.  You
sometimes have to type special commands in
conjunction with punctuation characters to get
this right, as in the following sentence.
       Gnats, gnus, etc.\ all  % `\ ' makes an inter-word space.
       begin with G\@.         % \@ marks end-of-sentence punctuation.
You should check the spaces after periods when
reading your output to make sure you haven't
forgotten any special cases.  Generating an
ellipsis
       \ldots\               % `\ ' is needed after `\ldots' because TeX
                             % ignores spaces after command names like \ldots
                             % made from \ + letters.
                             %
                             % Note how a `%' character causes TeX to ignore
                             % the end of the input line, so these blank lines
                             % do not start a new paragraph.
                             %
with the right spacing around the periods requires
a special command.

\LaTeX\ interprets some common characters as
commands, so you must type special commands to
generate them.  These characters include the
following:
       \$ \& \% \# \{ and \}.

In printing, text is usually emphasized with an
       \emph{italic}
type style.

\begin{em}
   A long segment of text can also be emphasized
   in this way.  Text within such a segment can be
   given \emph{additional} emphasis.
\end{em}

It is sometimes necessary to prevent \LaTeX\ from
breaking a line where it might otherwise do so.
This may be at a space, as between the ``Mr.''\ and
``Jones'' in
       ``Mr.~Jones'',        % ~ produces an unbreakable interword space.
or within a word---especially when the word is a
symbol like
       \mbox{\emph{itemnum}}
that makes little sense when hyphenated across
lines.

Footnotes\footnote{This is an example of a footnote.}
pose no problem.

\LaTeX\ is good at typesetting mathematical formulas
like
       \( x-3y + z = 7 \)
or
       \( a_{1} > x^{2n} + y^{2n} > x' \)
or
       \( \ip{A}{B} = \sum_{i} a_{i} b_{i} \).
The spaces you type in a formula are
ignored.  Remember that a letter like
       $x$                   % $ ... $  and  \( ... \)  are equivalent
is a formula when it denotes a mathematical
symbol, and it should be typed as one.

\section{Displayed Text}

Text is displayed by indenting it from the left
margin.  Quotations are commonly displayed.  There
are short quotations
\begin{quote}
   This is a short quotation.  It consists of a
   single paragraph of text.  See how it is formatted.
\end{quote}
and longer ones.
\begin{quotation}
   This is a longer quotation.  It consists of two
   paragraphs of text, neither of which are
   particularly interesting.

   This is the second paragraph of the quotation.  It
   is just as dull as the first paragraph.
\end{quotation}
Another frequently-displayed structure is a list.
The following is an example of an \emph{itemized}
list.
\begin{itemize}
   \item This is the first item of an itemized list.
         Each item in the list is marked with a ``tick''.
         You don't have to worry about what kind of tick
         mark is used.

   \item This is the second item of the list.  It
         contains another list nested inside it.  The inner
         list is an \emph{enumerated} list.
         \begin{enumerate}
            \item This is the first item of an enumerated
                  list that is nested within the itemized list.

            \item This is the second item of the inner list.
                  \LaTeX\ allows you to nest lists deeper than
                  you really should.
         \end{enumerate}
         This is the rest of the second item of the outer
         list.  It is no more interesting than any other
         part of the item.
   \item This is the third item of the list.
\end{itemize}
You can even display poetry.
\begin{verse}
   There is an environment
    for verse \\             % The \\ command separates lines
   Whose features some poets % within a stanza.
   will curse.

                             % One or more blank lines separate stanzas.

   For instead of making\\
   Them do \emph{all} line breaking, \\
   It allows them to put too many words on a line when they'd rather be
   forced to be terse.
\end{verse}

Mathematical formulas may also be displayed.  A
displayed formula
is
one-line long; multiline
formulas require special formatting instructions.
   \[  \ip{\Gamma}{\psi'} = x'' + y^{2} + z_{i}^{n}\]
Don't start a paragraph with a displayed equation,
nor make one a paragraph by itself.

\end{document}               % End of document.

\section{Bench Huge Section 1501}
% synthetic bench section 1501
% This is a sample LaTeX input file.  (Version of 12 August 2004.)
%
% A '%' character causes TeX to ignore all remaining text on the line,
% and is used for comments like this one.

\documentclass{article}      % Specifies the document class

                             % The preamble begins here.
\title{An Example Document}  % Declares the document's title.
\author{Leslie Lamport}      % Declares the author's name.
\date{January 21, 1994}      % Deleting this command produces today's date.

\newcommand{\ip}[2]{(#1, #2)}
                             % Defines \ip{arg1}{arg2} to mean
                             % (arg1, arg2).

%\newcommand{\ip}[2]{\langle #1 | #2\rangle}
                             % This is an alternative definition of
                             % \ip that is commented out.

\begin{document}             % End of preamble and beginning of text.

\maketitle                   % Produces the title.

This is an example input file.  Comparing it with
the output it generates can show you how to
produce a simple document of your own.

\section{Ordinary Text}      % Produces section heading.  Lower-level
                             % sections are begun with similar
                             % \subsection and \subsubsection commands.

The ends  of words and sentences are marked
  by   spaces. It  doesn't matter how many
spaces    you type; one is as good as 100.  The
end of   a line counts as a space.

One   or more   blank lines denote the  end
of  a paragraph.

Since any number of consecutive spaces are treated
like a single one, the formatting of the input
file makes no difference to
      \LaTeX,                % The \LaTeX command generates the LaTeX logo.
but it makes a difference to you.  When you use
\LaTeX, making your input file as easy to read
as possible will be a great help as you write
your document and when you change it.  This sample
file shows how you can add comments to your own input
file.

Because printing is different from typewriting,
there are a number of things that you have to do
differently when preparing an input file than if
you were just typing the document directly.
Quotation marks like
       ``this''
have to be handled specially, as do quotes within
quotes:
       ``\,`this'            % \, separates the double and single quote.
        is what I just
        wrote, not  `that'\,''.

Dashes come in three sizes: an
       intra-word
dash, a medium dash for number ranges like
       1--2,
and a punctuation
       dash---like
this.

A sentence-ending space should be larger than the
space between words within a sentence.  You
sometimes have to type special commands in
conjunction with punctuation characters to get
this right, as in the following sentence.
       Gnats, gnus, etc.\ all  % `\ ' makes an inter-word space.
       begin with G\@.         % \@ marks end-of-sentence punctuation.
You should check the spaces after periods when
reading your output to make sure you haven't
forgotten any special cases.  Generating an
ellipsis
       \ldots\               % `\ ' is needed after `\ldots' because TeX
                             % ignores spaces after command names like \ldots
                             % made from \ + letters.
                             %
                             % Note how a `%' character causes TeX to ignore
                             % the end of the input line, so these blank lines
                             % do not start a new paragraph.
                             %
with the right spacing around the periods requires
a special command.

\LaTeX\ interprets some common characters as
commands, so you must type special commands to
generate them.  These characters include the
following:
       \$ \& \% \# \{ and \}.

In printing, text is usually emphasized with an
       \emph{italic}
type style.

\begin{em}
   A long segment of text can also be emphasized
   in this way.  Text within such a segment can be
   given \emph{additional} emphasis.
\end{em}

It is sometimes necessary to prevent \LaTeX\ from
breaking a line where it might otherwise do so.
This may be at a space, as between the ``Mr.''\ and
``Jones'' in
       ``Mr.~Jones'',        % ~ produces an unbreakable interword space.
or within a word---especially when the word is a
symbol like
       \mbox{\emph{itemnum}}
that makes little sense when hyphenated across
lines.

Footnotes\footnote{This is an example of a footnote.}
pose no problem.

\LaTeX\ is good at typesetting mathematical formulas
like
       \( x-3y + z = 7 \)
or
       \( a_{1} > x^{2n} + y^{2n} > x' \)
or
       \( \ip{A}{B} = \sum_{i} a_{i} b_{i} \).
The spaces you type in a formula are
ignored.  Remember that a letter like
       $x$                   % $ ... $  and  \( ... \)  are equivalent
is a formula when it denotes a mathematical
symbol, and it should be typed as one.

\section{Displayed Text}

Text is displayed by indenting it from the left
margin.  Quotations are commonly displayed.  There
are short quotations
\begin{quote}
   This is a short quotation.  It consists of a
   single paragraph of text.  See how it is formatted.
\end{quote}
and longer ones.
\begin{quotation}
   This is a longer quotation.  It consists of two
   paragraphs of text, neither of which are
   particularly interesting.

   This is the second paragraph of the quotation.  It
   is just as dull as the first paragraph.
\end{quotation}
Another frequently-displayed structure is a list.
The following is an example of an \emph{itemized}
list.
\begin{itemize}
   \item This is the first item of an itemized list.
         Each item in the list is marked with a ``tick''.
         You don't have to worry about what kind of tick
         mark is used.

   \item This is the second item of the list.  It
         contains another list nested inside it.  The inner
         list is an \emph{enumerated} list.
         \begin{enumerate}
            \item This is the first item of an enumerated
                  list that is nested within the itemized list.

            \item This is the second item of the inner list.
                  \LaTeX\ allows you to nest lists deeper than
                  you really should.
         \end{enumerate}
         This is the rest of the second item of the outer
         list.  It is no more interesting than any other
         part of the item.
   \item This is the third item of the list.
\end{itemize}
You can even display poetry.
\begin{verse}
   There is an environment
    for verse \\             % The \\ command separates lines
   Whose features some poets % within a stanza.
   will curse.

                             % One or more blank lines separate stanzas.

   For instead of making\\
   Them do \emph{all} line breaking, \\
   It allows them to put too many words on a line when they'd rather be
   forced to be terse.
\end{verse}

Mathematical formulas may also be displayed.  A
displayed formula
is
one-line long; multiline
formulas require special formatting instructions.
   \[  \ip{\Gamma}{\psi'} = x'' + y^{2} + z_{i}^{n}\]
Don't start a paragraph with a displayed equation,
nor make one a paragraph by itself.

\end{document}               % End of document.

\section{Bench Huge Section 1502}
% synthetic bench section 1502
% This is a sample LaTeX input file.  (Version of 12 August 2004.)
%
% A '%' character causes TeX to ignore all remaining text on the line,
% and is used for comments like this one.

\documentclass{article}      % Specifies the document class

                             % The preamble begins here.
\title{An Example Document}  % Declares the document's title.
\author{Leslie Lamport}      % Declares the author's name.
\date{January 21, 1994}      % Deleting this command produces today's date.

\newcommand{\ip}[2]{(#1, #2)}
                             % Defines \ip{arg1}{arg2} to mean
                             % (arg1, arg2).

%\newcommand{\ip}[2]{\langle #1 | #2\rangle}
                             % This is an alternative definition of
                             % \ip that is commented out.

\begin{document}             % End of preamble and beginning of text.

\maketitle                   % Produces the title.

This is an example input file.  Comparing it with
the output it generates can show you how to
produce a simple document of your own.

\section{Ordinary Text}      % Produces section heading.  Lower-level
                             % sections are begun with similar
                             % \subsection and \subsubsection commands.

The ends  of words and sentences are marked
  by   spaces. It  doesn't matter how many
spaces    you type; one is as good as 100.  The
end of   a line counts as a space.

One   or more   blank lines denote the  end
of  a paragraph.

Since any number of consecutive spaces are treated
like a single one, the formatting of the input
file makes no difference to
      \LaTeX,                % The \LaTeX command generates the LaTeX logo.
but it makes a difference to you.  When you use
\LaTeX, making your input file as easy to read
as possible will be a great help as you write
your document and when you change it.  This sample
file shows how you can add comments to your own input
file.

Because printing is different from typewriting,
there are a number of things that you have to do
differently when preparing an input file than if
you were just typing the document directly.
Quotation marks like
       ``this''
have to be handled specially, as do quotes within
quotes:
       ``\,`this'            % \, separates the double and single quote.
        is what I just
        wrote, not  `that'\,''.

Dashes come in three sizes: an
       intra-word
dash, a medium dash for number ranges like
       1--2,
and a punctuation
       dash---like
this.

A sentence-ending space should be larger than the
space between words within a sentence.  You
sometimes have to type special commands in
conjunction with punctuation characters to get
this right, as in the following sentence.
       Gnats, gnus, etc.\ all  % `\ ' makes an inter-word space.
       begin with G\@.         % \@ marks end-of-sentence punctuation.
You should check the spaces after periods when
reading your output to make sure you haven't
forgotten any special cases.  Generating an
ellipsis
       \ldots\               % `\ ' is needed after `\ldots' because TeX
                             % ignores spaces after command names like \ldots
                             % made from \ + letters.
                             %
                             % Note how a `%' character causes TeX to ignore
                             % the end of the input line, so these blank lines
                             % do not start a new paragraph.
                             %
with the right spacing around the periods requires
a special command.

\LaTeX\ interprets some common characters as
commands, so you must type special commands to
generate them.  These characters include the
following:
       \$ \& \% \# \{ and \}.

In printing, text is usually emphasized with an
       \emph{italic}
type style.

\begin{em}
   A long segment of text can also be emphasized
   in this way.  Text within such a segment can be
   given \emph{additional} emphasis.
\end{em}

It is sometimes necessary to prevent \LaTeX\ from
breaking a line where it might otherwise do so.
This may be at a space, as between the ``Mr.''\ and
``Jones'' in
       ``Mr.~Jones'',        % ~ produces an unbreakable interword space.
or within a word---especially when the word is a
symbol like
       \mbox{\emph{itemnum}}
that makes little sense when hyphenated across
lines.

Footnotes\footnote{This is an example of a footnote.}
pose no problem.

\LaTeX\ is good at typesetting mathematical formulas
like
       \( x-3y + z = 7 \)
or
       \( a_{1} > x^{2n} + y^{2n} > x' \)
or
       \( \ip{A}{B} = \sum_{i} a_{i} b_{i} \).
The spaces you type in a formula are
ignored.  Remember that a letter like
       $x$                   % $ ... $  and  \( ... \)  are equivalent
is a formula when it denotes a mathematical
symbol, and it should be typed as one.

\section{Displayed Text}

Text is displayed by indenting it from the left
margin.  Quotations are commonly displayed.  There
are short quotations
\begin{quote}
   This is a short quotation.  It consists of a
   single paragraph of text.  See how it is formatted.
\end{quote}
and longer ones.
\begin{quotation}
   This is a longer quotation.  It consists of two
   paragraphs of text, neither of which are
   particularly interesting.

   This is the second paragraph of the quotation.  It
   is just as dull as the first paragraph.
\end{quotation}
Another frequently-displayed structure is a list.
The following is an example of an \emph{itemized}
list.
\begin{itemize}
   \item This is the first item of an itemized list.
         Each item in the list is marked with a ``tick''.
         You don't have to worry about what kind of tick
         mark is used.

   \item This is the second item of the list.  It
         contains another list nested inside it.  The inner
         list is an \emph{enumerated} list.
         \begin{enumerate}
            \item This is the first item of an enumerated
                  list that is nested within the itemized list.

            \item This is the second item of the inner list.
                  \LaTeX\ allows you to nest lists deeper than
                  you really should.
         \end{enumerate}
         This is the rest of the second item of the outer
         list.  It is no more interesting than any other
         part of the item.
   \item This is the third item of the list.
\end{itemize}
You can even display poetry.
\begin{verse}
   There is an environment
    for verse \\             % The \\ command separates lines
   Whose features some poets % within a stanza.
   will curse.

                             % One or more blank lines separate stanzas.

   For instead of making\\
   Them do \emph{all} line breaking, \\
   It allows them to put too many words on a line when they'd rather be
   forced to be terse.
\end{verse}

Mathematical formulas may also be displayed.  A
displayed formula
is
one-line long; multiline
formulas require special formatting instructions.
   \[  \ip{\Gamma}{\psi'} = x'' + y^{2} + z_{i}^{n}\]
Don't start a paragraph with a displayed equation,
nor make one a paragraph by itself.

\end{document}               % End of document.

\section{Bench Huge Section 1503}
% synthetic bench section 1503
% This is a sample LaTeX input file.  (Version of 12 August 2004.)
%
% A '%' character causes TeX to ignore all remaining text on the line,
% and is used for comments like this one.

\documentclass{article}      % Specifies the document class

                             % The preamble begins here.
\title{An Example Document}  % Declares the document's title.
\author{Leslie Lamport}      % Declares the author's name.
\date{January 21, 1994}      % Deleting this command produces today's date.

\newcommand{\ip}[2]{(#1, #2)}
                             % Defines \ip{arg1}{arg2} to mean
                             % (arg1, arg2).

%\newcommand{\ip}[2]{\langle #1 | #2\rangle}
                             % This is an alternative definition of
                             % \ip that is commented out.

\begin{document}             % End of preamble and beginning of text.

\maketitle                   % Produces the title.

This is an example input file.  Comparing it with
the output it generates can show you how to
produce a simple document of your own.

\section{Ordinary Text}      % Produces section heading.  Lower-level
                             % sections are begun with similar
                             % \subsection and \subsubsection commands.

The ends  of words and sentences are marked
  by   spaces. It  doesn't matter how many
spaces    you type; one is as good as 100.  The
end of   a line counts as a space.

One   or more   blank lines denote the  end
of  a paragraph.

Since any number of consecutive spaces are treated
like a single one, the formatting of the input
file makes no difference to
      \LaTeX,                % The \LaTeX command generates the LaTeX logo.
but it makes a difference to you.  When you use
\LaTeX, making your input file as easy to read
as possible will be a great help as you write
your document and when you change it.  This sample
file shows how you can add comments to your own input
file.

Because printing is different from typewriting,
there are a number of things that you have to do
differently when preparing an input file than if
you were just typing the document directly.
Quotation marks like
       ``this''
have to be handled specially, as do quotes within
quotes:
       ``\,`this'            % \, separates the double and single quote.
        is what I just
        wrote, not  `that'\,''.

Dashes come in three sizes: an
       intra-word
dash, a medium dash for number ranges like
       1--2,
and a punctuation
       dash---like
this.

A sentence-ending space should be larger than the
space between words within a sentence.  You
sometimes have to type special commands in
conjunction with punctuation characters to get
this right, as in the following sentence.
       Gnats, gnus, etc.\ all  % `\ ' makes an inter-word space.
       begin with G\@.         % \@ marks end-of-sentence punctuation.
You should check the spaces after periods when
reading your output to make sure you haven't
forgotten any special cases.  Generating an
ellipsis
       \ldots\               % `\ ' is needed after `\ldots' because TeX
                             % ignores spaces after command names like \ldots
                             % made from \ + letters.
                             %
                             % Note how a `%' character causes TeX to ignore
                             % the end of the input line, so these blank lines
                             % do not start a new paragraph.
                             %
with the right spacing around the periods requires
a special command.

\LaTeX\ interprets some common characters as
commands, so you must type special commands to
generate them.  These characters include the
following:
       \$ \& \% \# \{ and \}.

In printing, text is usually emphasized with an
       \emph{italic}
type style.

\begin{em}
   A long segment of text can also be emphasized
   in this way.  Text within such a segment can be
   given \emph{additional} emphasis.
\end{em}

It is sometimes necessary to prevent \LaTeX\ from
breaking a line where it might otherwise do so.
This may be at a space, as between the ``Mr.''\ and
``Jones'' in
       ``Mr.~Jones'',        % ~ produces an unbreakable interword space.
or within a word---especially when the word is a
symbol like
       \mbox{\emph{itemnum}}
that makes little sense when hyphenated across
lines.

Footnotes\footnote{This is an example of a footnote.}
pose no problem.

\LaTeX\ is good at typesetting mathematical formulas
like
       \( x-3y + z = 7 \)
or
       \( a_{1} > x^{2n} + y^{2n} > x' \)
or
       \( \ip{A}{B} = \sum_{i} a_{i} b_{i} \).
The spaces you type in a formula are
ignored.  Remember that a letter like
       $x$                   % $ ... $  and  \( ... \)  are equivalent
is a formula when it denotes a mathematical
symbol, and it should be typed as one.

\section{Displayed Text}

Text is displayed by indenting it from the left
margin.  Quotations are commonly displayed.  There
are short quotations
\begin{quote}
   This is a short quotation.  It consists of a
   single paragraph of text.  See how it is formatted.
\end{quote}
and longer ones.
\begin{quotation}
   This is a longer quotation.  It consists of two
   paragraphs of text, neither of which are
   particularly interesting.

   This is the second paragraph of the quotation.  It
   is just as dull as the first paragraph.
\end{quotation}
Another frequently-displayed structure is a list.
The following is an example of an \emph{itemized}
list.
\begin{itemize}
   \item This is the first item of an itemized list.
         Each item in the list is marked with a ``tick''.
         You don't have to worry about what kind of tick
         mark is used.

   \item This is the second item of the list.  It
         contains another list nested inside it.  The inner
         list is an \emph{enumerated} list.
         \begin{enumerate}
            \item This is the first item of an enumerated
                  list that is nested within the itemized list.

            \item This is the second item of the inner list.
                  \LaTeX\ allows you to nest lists deeper than
                  you really should.
         \end{enumerate}
         This is the rest of the second item of the outer
         list.  It is no more interesting than any other
         part of the item.
   \item This is the third item of the list.
\end{itemize}
You can even display poetry.
\begin{verse}
   There is an environment
    for verse \\             % The \\ command separates lines
   Whose features some poets % within a stanza.
   will curse.

                             % One or more blank lines separate stanzas.

   For instead of making\\
   Them do \emph{all} line breaking, \\
   It allows them to put too many words on a line when they'd rather be
   forced to be terse.
\end{verse}

Mathematical formulas may also be displayed.  A
displayed formula
is
one-line long; multiline
formulas require special formatting instructions.
   \[  \ip{\Gamma}{\psi'} = x'' + y^{2} + z_{i}^{n}\]
Don't start a paragraph with a displayed equation,
nor make one a paragraph by itself.

\end{document}               % End of document.

\section{Bench Huge Section 1504}
% synthetic bench section 1504
% This is a sample LaTeX input file.  (Version of 12 August 2004.)
%
% A '%' character causes TeX to ignore all remaining text on the line,
% and is used for comments like this one.

\documentclass{article}      % Specifies the document class

                             % The preamble begins here.
\title{An Example Document}  % Declares the document's title.
\author{Leslie Lamport}      % Declares the author's name.
\date{January 21, 1994}      % Deleting this command produces today's date.

\newcommand{\ip}[2]{(#1, #2)}
                             % Defines \ip{arg1}{arg2} to mean
                             % (arg1, arg2).

%\newcommand{\ip}[2]{\langle #1 | #2\rangle}
                             % This is an alternative definition of
                             % \ip that is commented out.

\begin{document}             % End of preamble and beginning of text.

\maketitle                   % Produces the title.

This is an example input file.  Comparing it with
the output it generates can show you how to
produce a simple document of your own.

\section{Ordinary Text}      % Produces section heading.  Lower-level
                             % sections are begun with similar
                             % \subsection and \subsubsection commands.

The ends  of words and sentences are marked
  by   spaces. It  doesn't matter how many
spaces    you type; one is as good as 100.  The
end of   a line counts as a space.

One   or more   blank lines denote the  end
of  a paragraph.

Since any number of consecutive spaces are treated
like a single one, the formatting of the input
file makes no difference to
      \LaTeX,                % The \LaTeX command generates the LaTeX logo.
but it makes a difference to you.  When you use
\LaTeX, making your input file as easy to read
as possible will be a great help as you write
your document and when you change it.  This sample
file shows how you can add comments to your own input
file.

Because printing is different from typewriting,
there are a number of things that you have to do
differently when preparing an input file than if
you were just typing the document directly.
Quotation marks like
       ``this''
have to be handled specially, as do quotes within
quotes:
       ``\,`this'            % \, separates the double and single quote.
        is what I just
        wrote, not  `that'\,''.

Dashes come in three sizes: an
       intra-word
dash, a medium dash for number ranges like
       1--2,
and a punctuation
       dash---like
this.

A sentence-ending space should be larger than the
space between words within a sentence.  You
sometimes have to type special commands in
conjunction with punctuation characters to get
this right, as in the following sentence.
       Gnats, gnus, etc.\ all  % `\ ' makes an inter-word space.
       begin with G\@.         % \@ marks end-of-sentence punctuation.
You should check the spaces after periods when
reading your output to make sure you haven't
forgotten any special cases.  Generating an
ellipsis
       \ldots\               % `\ ' is needed after `\ldots' because TeX
                             % ignores spaces after command names like \ldots
                             % made from \ + letters.
                             %
                             % Note how a `%' character causes TeX to ignore
                             % the end of the input line, so these blank lines
                             % do not start a new paragraph.
                             %
with the right spacing around the periods requires
a special command.

\LaTeX\ interprets some common characters as
commands, so you must type special commands to
generate them.  These characters include the
following:
       \$ \& \% \# \{ and \}.

In printing, text is usually emphasized with an
       \emph{italic}
type style.

\begin{em}
   A long segment of text can also be emphasized
   in this way.  Text within such a segment can be
   given \emph{additional} emphasis.
\end{em}

It is sometimes necessary to prevent \LaTeX\ from
breaking a line where it might otherwise do so.
This may be at a space, as between the ``Mr.''\ and
``Jones'' in
       ``Mr.~Jones'',        % ~ produces an unbreakable interword space.
or within a word---especially when the word is a
symbol like
       \mbox{\emph{itemnum}}
that makes little sense when hyphenated across
lines.

Footnotes\footnote{This is an example of a footnote.}
pose no problem.

\LaTeX\ is good at typesetting mathematical formulas
like
       \( x-3y + z = 7 \)
or
       \( a_{1} > x^{2n} + y^{2n} > x' \)
or
       \( \ip{A}{B} = \sum_{i} a_{i} b_{i} \).
The spaces you type in a formula are
ignored.  Remember that a letter like
       $x$                   % $ ... $  and  \( ... \)  are equivalent
is a formula when it denotes a mathematical
symbol, and it should be typed as one.

\section{Displayed Text}

Text is displayed by indenting it from the left
margin.  Quotations are commonly displayed.  There
are short quotations
\begin{quote}
   This is a short quotation.  It consists of a
   single paragraph of text.  See how it is formatted.
\end{quote}
and longer ones.
\begin{quotation}
   This is a longer quotation.  It consists of two
   paragraphs of text, neither of which are
   particularly interesting.

   This is the second paragraph of the quotation.  It
   is just as dull as the first paragraph.
\end{quotation}
Another frequently-displayed structure is a list.
The following is an example of an \emph{itemized}
list.
\begin{itemize}
   \item This is the first item of an itemized list.
         Each item in the list is marked with a ``tick''.
         You don't have to worry about what kind of tick
         mark is used.

   \item This is the second item of the list.  It
         contains another list nested inside it.  The inner
         list is an \emph{enumerated} list.
         \begin{enumerate}
            \item This is the first item of an enumerated
                  list that is nested within the itemized list.

            \item This is the second item of the inner list.
                  \LaTeX\ allows you to nest lists deeper than
                  you really should.
         \end{enumerate}
         This is the rest of the second item of the outer
         list.  It is no more interesting than any other
         part of the item.
   \item This is the third item of the list.
\end{itemize}
You can even display poetry.
\begin{verse}
   There is an environment
    for verse \\             % The \\ command separates lines
   Whose features some poets % within a stanza.
   will curse.

                             % One or more blank lines separate stanzas.

   For instead of making\\
   Them do \emph{all} line breaking, \\
   It allows them to put too many words on a line when they'd rather be
   forced to be terse.
\end{verse}

Mathematical formulas may also be displayed.  A
displayed formula
is
one-line long; multiline
formulas require special formatting instructions.
   \[  \ip{\Gamma}{\psi'} = x'' + y^{2} + z_{i}^{n}\]
Don't start a paragraph with a displayed equation,
nor make one a paragraph by itself.

\end{document}               % End of document.

\section{Bench Huge Section 1505}
% synthetic bench section 1505
% This is a sample LaTeX input file.  (Version of 12 August 2004.)
%
% A '%' character causes TeX to ignore all remaining text on the line,
% and is used for comments like this one.

\documentclass{article}      % Specifies the document class

                             % The preamble begins here.
\title{An Example Document}  % Declares the document's title.
\author{Leslie Lamport}      % Declares the author's name.
\date{January 21, 1994}      % Deleting this command produces today's date.

\newcommand{\ip}[2]{(#1, #2)}
                             % Defines \ip{arg1}{arg2} to mean
                             % (arg1, arg2).

%\newcommand{\ip}[2]{\langle #1 | #2\rangle}
                             % This is an alternative definition of
                             % \ip that is commented out.

\begin{document}             % End of preamble and beginning of text.

\maketitle                   % Produces the title.

This is an example input file.  Comparing it with
the output it generates can show you how to
produce a simple document of your own.

\section{Ordinary Text}      % Produces section heading.  Lower-level
                             % sections are begun with similar
                             % \subsection and \subsubsection commands.

The ends  of words and sentences are marked
  by   spaces. It  doesn't matter how many
spaces    you type; one is as good as 100.  The
end of   a line counts as a space.

One   or more   blank lines denote the  end
of  a paragraph.

Since any number of consecutive spaces are treated
like a single one, the formatting of the input
file makes no difference to
      \LaTeX,                % The \LaTeX command generates the LaTeX logo.
but it makes a difference to you.  When you use
\LaTeX, making your input file as easy to read
as possible will be a great help as you write
your document and when you change it.  This sample
file shows how you can add comments to your own input
file.

Because printing is different from typewriting,
there are a number of things that you have to do
differently when preparing an input file than if
you were just typing the document directly.
Quotation marks like
       ``this''
have to be handled specially, as do quotes within
quotes:
       ``\,`this'            % \, separates the double and single quote.
        is what I just
        wrote, not  `that'\,''.

Dashes come in three sizes: an
       intra-word
dash, a medium dash for number ranges like
       1--2,
and a punctuation
       dash---like
this.

A sentence-ending space should be larger than the
space between words within a sentence.  You
sometimes have to type special commands in
conjunction with punctuation characters to get
this right, as in the following sentence.
       Gnats, gnus, etc.\ all  % `\ ' makes an inter-word space.
       begin with G\@.         % \@ marks end-of-sentence punctuation.
You should check the spaces after periods when
reading your output to make sure you haven't
forgotten any special cases.  Generating an
ellipsis
       \ldots\               % `\ ' is needed after `\ldots' because TeX
                             % ignores spaces after command names like \ldots
                             % made from \ + letters.
                             %
                             % Note how a `%' character causes TeX to ignore
                             % the end of the input line, so these blank lines
                             % do not start a new paragraph.
                             %
with the right spacing around the periods requires
a special command.

\LaTeX\ interprets some common characters as
commands, so you must type special commands to
generate them.  These characters include the
following:
       \$ \& \% \# \{ and \}.

In printing, text is usually emphasized with an
       \emph{italic}
type style.

\begin{em}
   A long segment of text can also be emphasized
   in this way.  Text within such a segment can be
   given \emph{additional} emphasis.
\end{em}

It is sometimes necessary to prevent \LaTeX\ from
breaking a line where it might otherwise do so.
This may be at a space, as between the ``Mr.''\ and
``Jones'' in
       ``Mr.~Jones'',        % ~ produces an unbreakable interword space.
or within a word---especially when the word is a
symbol like
       \mbox{\emph{itemnum}}
that makes little sense when hyphenated across
lines.

Footnotes\footnote{This is an example of a footnote.}
pose no problem.

\LaTeX\ is good at typesetting mathematical formulas
like
       \( x-3y + z = 7 \)
or
       \( a_{1} > x^{2n} + y^{2n} > x' \)
or
       \( \ip{A}{B} = \sum_{i} a_{i} b_{i} \).
The spaces you type in a formula are
ignored.  Remember that a letter like
       $x$                   % $ ... $  and  \( ... \)  are equivalent
is a formula when it denotes a mathematical
symbol, and it should be typed as one.

\section{Displayed Text}

Text is displayed by indenting it from the left
margin.  Quotations are commonly displayed.  There
are short quotations
\begin{quote}
   This is a short quotation.  It consists of a
   single paragraph of text.  See how it is formatted.
\end{quote}
and longer ones.
\begin{quotation}
   This is a longer quotation.  It consists of two
   paragraphs of text, neither of which are
   particularly interesting.

   This is the second paragraph of the quotation.  It
   is just as dull as the first paragraph.
\end{quotation}
Another frequently-displayed structure is a list.
The following is an example of an \emph{itemized}
list.
\begin{itemize}
   \item This is the first item of an itemized list.
         Each item in the list is marked with a ``tick''.
         You don't have to worry about what kind of tick
         mark is used.

   \item This is the second item of the list.  It
         contains another list nested inside it.  The inner
         list is an \emph{enumerated} list.
         \begin{enumerate}
            \item This is the first item of an enumerated
                  list that is nested within the itemized list.

            \item This is the second item of the inner list.
                  \LaTeX\ allows you to nest lists deeper than
                  you really should.
         \end{enumerate}
         This is the rest of the second item of the outer
         list.  It is no more interesting than any other
         part of the item.
   \item This is the third item of the list.
\end{itemize}
You can even display poetry.
\begin{verse}
   There is an environment
    for verse \\             % The \\ command separates lines
   Whose features some poets % within a stanza.
   will curse.

                             % One or more blank lines separate stanzas.

   For instead of making\\
   Them do \emph{all} line breaking, \\
   It allows them to put too many words on a line when they'd rather be
   forced to be terse.
\end{verse}

Mathematical formulas may also be displayed.  A
displayed formula
is
one-line long; multiline
formulas require special formatting instructions.
   \[  \ip{\Gamma}{\psi'} = x'' + y^{2} + z_{i}^{n}\]
Don't start a paragraph with a displayed equation,
nor make one a paragraph by itself.

\end{document}               % End of document.

\section{Bench Huge Section 1506}
% synthetic bench section 1506
% This is a sample LaTeX input file.  (Version of 12 August 2004.)
%
% A '%' character causes TeX to ignore all remaining text on the line,
% and is used for comments like this one.

\documentclass{article}      % Specifies the document class

                             % The preamble begins here.
\title{An Example Document}  % Declares the document's title.
\author{Leslie Lamport}      % Declares the author's name.
\date{January 21, 1994}      % Deleting this command produces today's date.

\newcommand{\ip}[2]{(#1, #2)}
                             % Defines \ip{arg1}{arg2} to mean
                             % (arg1, arg2).

%\newcommand{\ip}[2]{\langle #1 | #2\rangle}
                             % This is an alternative definition of
                             % \ip that is commented out.

\begin{document}             % End of preamble and beginning of text.

\maketitle                   % Produces the title.

This is an example input file.  Comparing it with
the output it generates can show you how to
produce a simple document of your own.

\section{Ordinary Text}      % Produces section heading.  Lower-level
                             % sections are begun with similar
                             % \subsection and \subsubsection commands.

The ends  of words and sentences are marked
  by   spaces. It  doesn't matter how many
spaces    you type; one is as good as 100.  The
end of   a line counts as a space.

One   or more   blank lines denote the  end
of  a paragraph.

Since any number of consecutive spaces are treated
like a single one, the formatting of the input
file makes no difference to
      \LaTeX,                % The \LaTeX command generates the LaTeX logo.
but it makes a difference to you.  When you use
\LaTeX, making your input file as easy to read
as possible will be a great help as you write
your document and when you change it.  This sample
file shows how you can add comments to your own input
file.

Because printing is different from typewriting,
there are a number of things that you have to do
differently when preparing an input file than if
you were just typing the document directly.
Quotation marks like
       ``this''
have to be handled specially, as do quotes within
quotes:
       ``\,`this'            % \, separates the double and single quote.
        is what I just
        wrote, not  `that'\,''.

Dashes come in three sizes: an
       intra-word
dash, a medium dash for number ranges like
       1--2,
and a punctuation
       dash---like
this.

A sentence-ending space should be larger than the
space between words within a sentence.  You
sometimes have to type special commands in
conjunction with punctuation characters to get
this right, as in the following sentence.
       Gnats, gnus, etc.\ all  % `\ ' makes an inter-word space.
       begin with G\@.         % \@ marks end-of-sentence punctuation.
You should check the spaces after periods when
reading your output to make sure you haven't
forgotten any special cases.  Generating an
ellipsis
       \ldots\               % `\ ' is needed after `\ldots' because TeX
                             % ignores spaces after command names like \ldots
                             % made from \ + letters.
                             %
                             % Note how a `%' character causes TeX to ignore
                             % the end of the input line, so these blank lines
                             % do not start a new paragraph.
                             %
with the right spacing around the periods requires
a special command.

\LaTeX\ interprets some common characters as
commands, so you must type special commands to
generate them.  These characters include the
following:
       \$ \& \% \# \{ and \}.

In printing, text is usually emphasized with an
       \emph{italic}
type style.

\begin{em}
   A long segment of text can also be emphasized
   in this way.  Text within such a segment can be
   given \emph{additional} emphasis.
\end{em}

It is sometimes necessary to prevent \LaTeX\ from
breaking a line where it might otherwise do so.
This may be at a space, as between the ``Mr.''\ and
``Jones'' in
       ``Mr.~Jones'',        % ~ produces an unbreakable interword space.
or within a word---especially when the word is a
symbol like
       \mbox{\emph{itemnum}}
that makes little sense when hyphenated across
lines.

Footnotes\footnote{This is an example of a footnote.}
pose no problem.

\LaTeX\ is good at typesetting mathematical formulas
like
       \( x-3y + z = 7 \)
or
       \( a_{1} > x^{2n} + y^{2n} > x' \)
or
       \( \ip{A}{B} = \sum_{i} a_{i} b_{i} \).
The spaces you type in a formula are
ignored.  Remember that a letter like
       $x$                   % $ ... $  and  \( ... \)  are equivalent
is a formula when it denotes a mathematical
symbol, and it should be typed as one.

\section{Displayed Text}

Text is displayed by indenting it from the left
margin.  Quotations are commonly displayed.  There
are short quotations
\begin{quote}
   This is a short quotation.  It consists of a
   single paragraph of text.  See how it is formatted.
\end{quote}
and longer ones.
\begin{quotation}
   This is a longer quotation.  It consists of two
   paragraphs of text, neither of which are
   particularly interesting.

   This is the second paragraph of the quotation.  It
   is just as dull as the first paragraph.
\end{quotation}
Another frequently-displayed structure is a list.
The following is an example of an \emph{itemized}
list.
\begin{itemize}
   \item This is the first item of an itemized list.
         Each item in the list is marked with a ``tick''.
         You don't have to worry about what kind of tick
         mark is used.

   \item This is the second item of the list.  It
         contains another list nested inside it.  The inner
         list is an \emph{enumerated} list.
         \begin{enumerate}
            \item This is the first item of an enumerated
                  list that is nested within the itemized list.

            \item This is the second item of the inner list.
                  \LaTeX\ allows you to nest lists deeper than
                  you really should.
         \end{enumerate}
         This is the rest of the second item of the outer
         list.  It is no more interesting than any other
         part of the item.
   \item This is the third item of the list.
\end{itemize}
You can even display poetry.
\begin{verse}
   There is an environment
    for verse \\             % The \\ command separates lines
   Whose features some poets % within a stanza.
   will curse.

                             % One or more blank lines separate stanzas.

   For instead of making\\
   Them do \emph{all} line breaking, \\
   It allows them to put too many words on a line when they'd rather be
   forced to be terse.
\end{verse}

Mathematical formulas may also be displayed.  A
displayed formula
is
one-line long; multiline
formulas require special formatting instructions.
   \[  \ip{\Gamma}{\psi'} = x'' + y^{2} + z_{i}^{n}\]
Don't start a paragraph with a displayed equation,
nor make one a paragraph by itself.

\end{document}               % End of document.

\section{Bench Huge Section 1507}
% synthetic bench section 1507
% This is a sample LaTeX input file.  (Version of 12 August 2004.)
%
% A '%' character causes TeX to ignore all remaining text on the line,
% and is used for comments like this one.

\documentclass{article}      % Specifies the document class

                             % The preamble begins here.
\title{An Example Document}  % Declares the document's title.
\author{Leslie Lamport}      % Declares the author's name.
\date{January 21, 1994}      % Deleting this command produces today's date.

\newcommand{\ip}[2]{(#1, #2)}
                             % Defines \ip{arg1}{arg2} to mean
                             % (arg1, arg2).

%\newcommand{\ip}[2]{\langle #1 | #2\rangle}
                             % This is an alternative definition of
                             % \ip that is commented out.

\begin{document}             % End of preamble and beginning of text.

\maketitle                   % Produces the title.

This is an example input file.  Comparing it with
the output it generates can show you how to
produce a simple document of your own.

\section{Ordinary Text}      % Produces section heading.  Lower-level
                             % sections are begun with similar
                             % \subsection and \subsubsection commands.

The ends  of words and sentences are marked
  by   spaces. It  doesn't matter how many
spaces    you type; one is as good as 100.  The
end of   a line counts as a space.

One   or more   blank lines denote the  end
of  a paragraph.

Since any number of consecutive spaces are treated
like a single one, the formatting of the input
file makes no difference to
      \LaTeX,                % The \LaTeX command generates the LaTeX logo.
but it makes a difference to you.  When you use
\LaTeX, making your input file as easy to read
as possible will be a great help as you write
your document and when you change it.  This sample
file shows how you can add comments to your own input
file.

Because printing is different from typewriting,
there are a number of things that you have to do
differently when preparing an input file than if
you were just typing the document directly.
Quotation marks like
       ``this''
have to be handled specially, as do quotes within
quotes:
       ``\,`this'            % \, separates the double and single quote.
        is what I just
        wrote, not  `that'\,''.

Dashes come in three sizes: an
       intra-word
dash, a medium dash for number ranges like
       1--2,
and a punctuation
       dash---like
this.

A sentence-ending space should be larger than the
space between words within a sentence.  You
sometimes have to type special commands in
conjunction with punctuation characters to get
this right, as in the following sentence.
       Gnats, gnus, etc.\ all  % `\ ' makes an inter-word space.
       begin with G\@.         % \@ marks end-of-sentence punctuation.
You should check the spaces after periods when
reading your output to make sure you haven't
forgotten any special cases.  Generating an
ellipsis
       \ldots\               % `\ ' is needed after `\ldots' because TeX
                             % ignores spaces after command names like \ldots
                             % made from \ + letters.
                             %
                             % Note how a `%' character causes TeX to ignore
                             % the end of the input line, so these blank lines
                             % do not start a new paragraph.
                             %
with the right spacing around the periods requires
a special command.

\LaTeX\ interprets some common characters as
commands, so you must type special commands to
generate them.  These characters include the
following:
       \$ \& \% \# \{ and \}.

In printing, text is usually emphasized with an
       \emph{italic}
type style.

\begin{em}
   A long segment of text can also be emphasized
   in this way.  Text within such a segment can be
   given \emph{additional} emphasis.
\end{em}

It is sometimes necessary to prevent \LaTeX\ from
breaking a line where it might otherwise do so.
This may be at a space, as between the ``Mr.''\ and
``Jones'' in
       ``Mr.~Jones'',        % ~ produces an unbreakable interword space.
or within a word---especially when the word is a
symbol like
       \mbox{\emph{itemnum}}
that makes little sense when hyphenated across
lines.

Footnotes\footnote{This is an example of a footnote.}
pose no problem.

\LaTeX\ is good at typesetting mathematical formulas
like
       \( x-3y + z = 7 \)
or
       \( a_{1} > x^{2n} + y^{2n} > x' \)
or
       \( \ip{A}{B} = \sum_{i} a_{i} b_{i} \).
The spaces you type in a formula are
ignored.  Remember that a letter like
       $x$                   % $ ... $  and  \( ... \)  are equivalent
is a formula when it denotes a mathematical
symbol, and it should be typed as one.

\section{Displayed Text}

Text is displayed by indenting it from the left
margin.  Quotations are commonly displayed.  There
are short quotations
\begin{quote}
   This is a short quotation.  It consists of a
   single paragraph of text.  See how it is formatted.
\end{quote}
and longer ones.
\begin{quotation}
   This is a longer quotation.  It consists of two
   paragraphs of text, neither of which are
   particularly interesting.

   This is the second paragraph of the quotation.  It
   is just as dull as the first paragraph.
\end{quotation}
Another frequently-displayed structure is a list.
The following is an example of an \emph{itemized}
list.
\begin{itemize}
   \item This is the first item of an itemized list.
         Each item in the list is marked with a ``tick''.
         You don't have to worry about what kind of tick
         mark is used.

   \item This is the second item of the list.  It
         contains another list nested inside it.  The inner
         list is an \emph{enumerated} list.
         \begin{enumerate}
            \item This is the first item of an enumerated
                  list that is nested within the itemized list.

            \item This is the second item of the inner list.
                  \LaTeX\ allows you to nest lists deeper than
                  you really should.
         \end{enumerate}
         This is the rest of the second item of the outer
         list.  It is no more interesting than any other
         part of the item.
   \item This is the third item of the list.
\end{itemize}
You can even display poetry.
\begin{verse}
   There is an environment
    for verse \\             % The \\ command separates lines
   Whose features some poets % within a stanza.
   will curse.

                             % One or more blank lines separate stanzas.

   For instead of making\\
   Them do \emph{all} line breaking, \\
   It allows them to put too many words on a line when they'd rather be
   forced to be terse.
\end{verse}

Mathematical formulas may also be displayed.  A
displayed formula
is
one-line long; multiline
formulas require special formatting instructions.
   \[  \ip{\Gamma}{\psi'} = x'' + y^{2} + z_{i}^{n}\]
Don't start a paragraph with a displayed equation,
nor make one a paragraph by itself.

\end{document}               % End of document.

\section{Bench Huge Section 1508}
% synthetic bench section 1508
% This is a sample LaTeX input file.  (Version of 12 August 2004.)
%
% A '%' character causes TeX to ignore all remaining text on the line,
% and is used for comments like this one.

\documentclass{article}      % Specifies the document class

                             % The preamble begins here.
\title{An Example Document}  % Declares the document's title.
\author{Leslie Lamport}      % Declares the author's name.
\date{January 21, 1994}      % Deleting this command produces today's date.

\newcommand{\ip}[2]{(#1, #2)}
                             % Defines \ip{arg1}{arg2} to mean
                             % (arg1, arg2).

%\newcommand{\ip}[2]{\langle #1 | #2\rangle}
                             % This is an alternative definition of
                             % \ip that is commented out.

\begin{document}             % End of preamble and beginning of text.

\maketitle                   % Produces the title.

This is an example input file.  Comparing it with
the output it generates can show you how to
produce a simple document of your own.

\section{Ordinary Text}      % Produces section heading.  Lower-level
                             % sections are begun with similar
                             % \subsection and \subsubsection commands.

The ends  of words and sentences are marked
  by   spaces. It  doesn't matter how many
spaces    you type; one is as good as 100.  The
end of   a line counts as a space.

One   or more   blank lines denote the  end
of  a paragraph.

Since any number of consecutive spaces are treated
like a single one, the formatting of the input
file makes no difference to
      \LaTeX,                % The \LaTeX command generates the LaTeX logo.
but it makes a difference to you.  When you use
\LaTeX, making your input file as easy to read
as possible will be a great help as you write
your document and when you change it.  This sample
file shows how you can add comments to your own input
file.

Because printing is different from typewriting,
there are a number of things that you have to do
differently when preparing an input file than if
you were just typing the document directly.
Quotation marks like
       ``this''
have to be handled specially, as do quotes within
quotes:
       ``\,`this'            % \, separates the double and single quote.
        is what I just
        wrote, not  `that'\,''.

Dashes come in three sizes: an
       intra-word
dash, a medium dash for number ranges like
       1--2,
and a punctuation
       dash---like
this.

A sentence-ending space should be larger than the
space between words within a sentence.  You
sometimes have to type special commands in
conjunction with punctuation characters to get
this right, as in the following sentence.
       Gnats, gnus, etc.\ all  % `\ ' makes an inter-word space.
       begin with G\@.         % \@ marks end-of-sentence punctuation.
You should check the spaces after periods when
reading your output to make sure you haven't
forgotten any special cases.  Generating an
ellipsis
       \ldots\               % `\ ' is needed after `\ldots' because TeX
                             % ignores spaces after command names like \ldots
                             % made from \ + letters.
                             %
                             % Note how a `%' character causes TeX to ignore
                             % the end of the input line, so these blank lines
                             % do not start a new paragraph.
                             %
with the right spacing around the periods requires
a special command.

\LaTeX\ interprets some common characters as
commands, so you must type special commands to
generate them.  These characters include the
following:
       \$ \& \% \# \{ and \}.

In printing, text is usually emphasized with an
       \emph{italic}
type style.

\begin{em}
   A long segment of text can also be emphasized
   in this way.  Text within such a segment can be
   given \emph{additional} emphasis.
\end{em}

It is sometimes necessary to prevent \LaTeX\ from
breaking a line where it might otherwise do so.
This may be at a space, as between the ``Mr.''\ and
``Jones'' in
       ``Mr.~Jones'',        % ~ produces an unbreakable interword space.
or within a word---especially when the word is a
symbol like
       \mbox{\emph{itemnum}}
that makes little sense when hyphenated across
lines.

Footnotes\footnote{This is an example of a footnote.}
pose no problem.

\LaTeX\ is good at typesetting mathematical formulas
like
       \( x-3y + z = 7 \)
or
       \( a_{1} > x^{2n} + y^{2n} > x' \)
or
       \( \ip{A}{B} = \sum_{i} a_{i} b_{i} \).
The spaces you type in a formula are
ignored.  Remember that a letter like
       $x$                   % $ ... $  and  \( ... \)  are equivalent
is a formula when it denotes a mathematical
symbol, and it should be typed as one.

\section{Displayed Text}

Text is displayed by indenting it from the left
margin.  Quotations are commonly displayed.  There
are short quotations
\begin{quote}
   This is a short quotation.  It consists of a
   single paragraph of text.  See how it is formatted.
\end{quote}
and longer ones.
\begin{quotation}
   This is a longer quotation.  It consists of two
   paragraphs of text, neither of which are
   particularly interesting.

   This is the second paragraph of the quotation.  It
   is just as dull as the first paragraph.
\end{quotation}
Another frequently-displayed structure is a list.
The following is an example of an \emph{itemized}
list.
\begin{itemize}
   \item This is the first item of an itemized list.
         Each item in the list is marked with a ``tick''.
         You don't have to worry about what kind of tick
         mark is used.

   \item This is the second item of the list.  It
         contains another list nested inside it.  The inner
         list is an \emph{enumerated} list.
         \begin{enumerate}
            \item This is the first item of an enumerated
                  list that is nested within the itemized list.

            \item This is the second item of the inner list.
                  \LaTeX\ allows you to nest lists deeper than
                  you really should.
         \end{enumerate}
         This is the rest of the second item of the outer
         list.  It is no more interesting than any other
         part of the item.
   \item This is the third item of the list.
\end{itemize}
You can even display poetry.
\begin{verse}
   There is an environment
    for verse \\             % The \\ command separates lines
   Whose features some poets % within a stanza.
   will curse.

                             % One or more blank lines separate stanzas.

   For instead of making\\
   Them do \emph{all} line breaking, \\
   It allows them to put too many words on a line when they'd rather be
   forced to be terse.
\end{verse}

Mathematical formulas may also be displayed.  A
displayed formula
is
one-line long; multiline
formulas require special formatting instructions.
   \[  \ip{\Gamma}{\psi'} = x'' + y^{2} + z_{i}^{n}\]
Don't start a paragraph with a displayed equation,
nor make one a paragraph by itself.

\end{document}               % End of document.

\section{Bench Huge Section 1509}
% synthetic bench section 1509
% This is a sample LaTeX input file.  (Version of 12 August 2004.)
%
% A '%' character causes TeX to ignore all remaining text on the line,
% and is used for comments like this one.

\documentclass{article}      % Specifies the document class

                             % The preamble begins here.
\title{An Example Document}  % Declares the document's title.
\author{Leslie Lamport}      % Declares the author's name.
\date{January 21, 1994}      % Deleting this command produces today's date.

\newcommand{\ip}[2]{(#1, #2)}
                             % Defines \ip{arg1}{arg2} to mean
                             % (arg1, arg2).

%\newcommand{\ip}[2]{\langle #1 | #2\rangle}
                             % This is an alternative definition of
                             % \ip that is commented out.

\begin{document}             % End of preamble and beginning of text.

\maketitle                   % Produces the title.

This is an example input file.  Comparing it with
the output it generates can show you how to
produce a simple document of your own.

\section{Ordinary Text}      % Produces section heading.  Lower-level
                             % sections are begun with similar
                             % \subsection and \subsubsection commands.

The ends  of words and sentences are marked
  by   spaces. It  doesn't matter how many
spaces    you type; one is as good as 100.  The
end of   a line counts as a space.

One   or more   blank lines denote the  end
of  a paragraph.

Since any number of consecutive spaces are treated
like a single one, the formatting of the input
file makes no difference to
      \LaTeX,                % The \LaTeX command generates the LaTeX logo.
but it makes a difference to you.  When you use
\LaTeX, making your input file as easy to read
as possible will be a great help as you write
your document and when you change it.  This sample
file shows how you can add comments to your own input
file.

Because printing is different from typewriting,
there are a number of things that you have to do
differently when preparing an input file than if
you were just typing the document directly.
Quotation marks like
       ``this''
have to be handled specially, as do quotes within
quotes:
       ``\,`this'            % \, separates the double and single quote.
        is what I just
        wrote, not  `that'\,''.

Dashes come in three sizes: an
       intra-word
dash, a medium dash for number ranges like
       1--2,
and a punctuation
       dash---like
this.

A sentence-ending space should be larger than the
space between words within a sentence.  You
sometimes have to type special commands in
conjunction with punctuation characters to get
this right, as in the following sentence.
       Gnats, gnus, etc.\ all  % `\ ' makes an inter-word space.
       begin with G\@.         % \@ marks end-of-sentence punctuation.
You should check the spaces after periods when
reading your output to make sure you haven't
forgotten any special cases.  Generating an
ellipsis
       \ldots\               % `\ ' is needed after `\ldots' because TeX
                             % ignores spaces after command names like \ldots
                             % made from \ + letters.
                             %
                             % Note how a `%' character causes TeX to ignore
                             % the end of the input line, so these blank lines
                             % do not start a new paragraph.
                             %
with the right spacing around the periods requires
a special command.

\LaTeX\ interprets some common characters as
commands, so you must type special commands to
generate them.  These characters include the
following:
       \$ \& \% \# \{ and \}.

In printing, text is usually emphasized with an
       \emph{italic}
type style.

\begin{em}
   A long segment of text can also be emphasized
   in this way.  Text within such a segment can be
   given \emph{additional} emphasis.
\end{em}

It is sometimes necessary to prevent \LaTeX\ from
breaking a line where it might otherwise do so.
This may be at a space, as between the ``Mr.''\ and
``Jones'' in
       ``Mr.~Jones'',        % ~ produces an unbreakable interword space.
or within a word---especially when the word is a
symbol like
       \mbox{\emph{itemnum}}
that makes little sense when hyphenated across
lines.

Footnotes\footnote{This is an example of a footnote.}
pose no problem.

\LaTeX\ is good at typesetting mathematical formulas
like
       \( x-3y + z = 7 \)
or
       \( a_{1} > x^{2n} + y^{2n} > x' \)
or
       \( \ip{A}{B} = \sum_{i} a_{i} b_{i} \).
The spaces you type in a formula are
ignored.  Remember that a letter like
       $x$                   % $ ... $  and  \( ... \)  are equivalent
is a formula when it denotes a mathematical
symbol, and it should be typed as one.

\section{Displayed Text}

Text is displayed by indenting it from the left
margin.  Quotations are commonly displayed.  There
are short quotations
\begin{quote}
   This is a short quotation.  It consists of a
   single paragraph of text.  See how it is formatted.
\end{quote}
and longer ones.
\begin{quotation}
   This is a longer quotation.  It consists of two
   paragraphs of text, neither of which are
   particularly interesting.

   This is the second paragraph of the quotation.  It
   is just as dull as the first paragraph.
\end{quotation}
Another frequently-displayed structure is a list.
The following is an example of an \emph{itemized}
list.
\begin{itemize}
   \item This is the first item of an itemized list.
         Each item in the list is marked with a ``tick''.
         You don't have to worry about what kind of tick
         mark is used.

   \item This is the second item of the list.  It
         contains another list nested inside it.  The inner
         list is an \emph{enumerated} list.
         \begin{enumerate}
            \item This is the first item of an enumerated
                  list that is nested within the itemized list.

            \item This is the second item of the inner list.
                  \LaTeX\ allows you to nest lists deeper than
                  you really should.
         \end{enumerate}
         This is the rest of the second item of the outer
         list.  It is no more interesting than any other
         part of the item.
   \item This is the third item of the list.
\end{itemize}
You can even display poetry.
\begin{verse}
   There is an environment
    for verse \\             % The \\ command separates lines
   Whose features some poets % within a stanza.
   will curse.

                             % One or more blank lines separate stanzas.

   For instead of making\\
   Them do \emph{all} line breaking, \\
   It allows them to put too many words on a line when they'd rather be
   forced to be terse.
\end{verse}

Mathematical formulas may also be displayed.  A
displayed formula
is
one-line long; multiline
formulas require special formatting instructions.
   \[  \ip{\Gamma}{\psi'} = x'' + y^{2} + z_{i}^{n}\]
Don't start a paragraph with a displayed equation,
nor make one a paragraph by itself.

\end{document}               % End of document.

\section{Bench Huge Section 1510}
% synthetic bench section 1510
% This is a sample LaTeX input file.  (Version of 12 August 2004.)
%
% A '%' character causes TeX to ignore all remaining text on the line,
% and is used for comments like this one.

\documentclass{article}      % Specifies the document class

                             % The preamble begins here.
\title{An Example Document}  % Declares the document's title.
\author{Leslie Lamport}      % Declares the author's name.
\date{January 21, 1994}      % Deleting this command produces today's date.

\newcommand{\ip}[2]{(#1, #2)}
                             % Defines \ip{arg1}{arg2} to mean
                             % (arg1, arg2).

%\newcommand{\ip}[2]{\langle #1 | #2\rangle}
                             % This is an alternative definition of
                             % \ip that is commented out.

\begin{document}             % End of preamble and beginning of text.

\maketitle                   % Produces the title.

This is an example input file.  Comparing it with
the output it generates can show you how to
produce a simple document of your own.

\section{Ordinary Text}      % Produces section heading.  Lower-level
                             % sections are begun with similar
                             % \subsection and \subsubsection commands.

The ends  of words and sentences are marked
  by   spaces. It  doesn't matter how many
spaces    you type; one is as good as 100.  The
end of   a line counts as a space.

One   or more   blank lines denote the  end
of  a paragraph.

Since any number of consecutive spaces are treated
like a single one, the formatting of the input
file makes no difference to
      \LaTeX,                % The \LaTeX command generates the LaTeX logo.
but it makes a difference to you.  When you use
\LaTeX, making your input file as easy to read
as possible will be a great help as you write
your document and when you change it.  This sample
file shows how you can add comments to your own input
file.

Because printing is different from typewriting,
there are a number of things that you have to do
differently when preparing an input file than if
you were just typing the document directly.
Quotation marks like
       ``this''
have to be handled specially, as do quotes within
quotes:
       ``\,`this'            % \, separates the double and single quote.
        is what I just
        wrote, not  `that'\,''.

Dashes come in three sizes: an
       intra-word
dash, a medium dash for number ranges like
       1--2,
and a punctuation
       dash---like
this.

A sentence-ending space should be larger than the
space between words within a sentence.  You
sometimes have to type special commands in
conjunction with punctuation characters to get
this right, as in the following sentence.
       Gnats, gnus, etc.\ all  % `\ ' makes an inter-word space.
       begin with G\@.         % \@ marks end-of-sentence punctuation.
You should check the spaces after periods when
reading your output to make sure you haven't
forgotten any special cases.  Generating an
ellipsis
       \ldots\               % `\ ' is needed after `\ldots' because TeX
                             % ignores spaces after command names like \ldots
                             % made from \ + letters.
                             %
                             % Note how a `%' character causes TeX to ignore
                             % the end of the input line, so these blank lines
                             % do not start a new paragraph.
                             %
with the right spacing around the periods requires
a special command.

\LaTeX\ interprets some common characters as
commands, so you must type special commands to
generate them.  These characters include the
following:
       \$ \& \% \# \{ and \}.

In printing, text is usually emphasized with an
       \emph{italic}
type style.

\begin{em}
   A long segment of text can also be emphasized
   in this way.  Text within such a segment can be
   given \emph{additional} emphasis.
\end{em}

It is sometimes necessary to prevent \LaTeX\ from
breaking a line where it might otherwise do so.
This may be at a space, as between the ``Mr.''\ and
``Jones'' in
       ``Mr.~Jones'',        % ~ produces an unbreakable interword space.
or within a word---especially when the word is a
symbol like
       \mbox{\emph{itemnum}}
that makes little sense when hyphenated across
lines.

Footnotes\footnote{This is an example of a footnote.}
pose no problem.

\LaTeX\ is good at typesetting mathematical formulas
like
       \( x-3y + z = 7 \)
or
       \( a_{1} > x^{2n} + y^{2n} > x' \)
or
       \( \ip{A}{B} = \sum_{i} a_{i} b_{i} \).
The spaces you type in a formula are
ignored.  Remember that a letter like
       $x$                   % $ ... $  and  \( ... \)  are equivalent
is a formula when it denotes a mathematical
symbol, and it should be typed as one.

\section{Displayed Text}

Text is displayed by indenting it from the left
margin.  Quotations are commonly displayed.  There
are short quotations
\begin{quote}
   This is a short quotation.  It consists of a
   single paragraph of text.  See how it is formatted.
\end{quote}
and longer ones.
\begin{quotation}
   This is a longer quotation.  It consists of two
   paragraphs of text, neither of which are
   particularly interesting.

   This is the second paragraph of the quotation.  It
   is just as dull as the first paragraph.
\end{quotation}
Another frequently-displayed structure is a list.
The following is an example of an \emph{itemized}
list.
\begin{itemize}
   \item This is the first item of an itemized list.
         Each item in the list is marked with a ``tick''.
         You don't have to worry about what kind of tick
         mark is used.

   \item This is the second item of the list.  It
         contains another list nested inside it.  The inner
         list is an \emph{enumerated} list.
         \begin{enumerate}
            \item This is the first item of an enumerated
                  list that is nested within the itemized list.

            \item This is the second item of the inner list.
                  \LaTeX\ allows you to nest lists deeper than
                  you really should.
         \end{enumerate}
         This is the rest of the second item of the outer
         list.  It is no more interesting than any other
         part of the item.
   \item This is the third item of the list.
\end{itemize}
You can even display poetry.
\begin{verse}
   There is an environment
    for verse \\             % The \\ command separates lines
   Whose features some poets % within a stanza.
   will curse.

                             % One or more blank lines separate stanzas.

   For instead of making\\
   Them do \emph{all} line breaking, \\
   It allows them to put too many words on a line when they'd rather be
   forced to be terse.
\end{verse}

Mathematical formulas may also be displayed.  A
displayed formula
is
one-line long; multiline
formulas require special formatting instructions.
   \[  \ip{\Gamma}{\psi'} = x'' + y^{2} + z_{i}^{n}\]
Don't start a paragraph with a displayed equation,
nor make one a paragraph by itself.

\end{document}               % End of document.

\section{Bench Huge Section 1511}
% synthetic bench section 1511
% This is a sample LaTeX input file.  (Version of 12 August 2004.)
%
% A '%' character causes TeX to ignore all remaining text on the line,
% and is used for comments like this one.

\documentclass{article}      % Specifies the document class

                             % The preamble begins here.
\title{An Example Document}  % Declares the document's title.
\author{Leslie Lamport}      % Declares the author's name.
\date{January 21, 1994}      % Deleting this command produces today's date.

\newcommand{\ip}[2]{(#1, #2)}
                             % Defines \ip{arg1}{arg2} to mean
                             % (arg1, arg2).

%\newcommand{\ip}[2]{\langle #1 | #2\rangle}
                             % This is an alternative definition of
                             % \ip that is commented out.

\begin{document}             % End of preamble and beginning of text.

\maketitle                   % Produces the title.

This is an example input file.  Comparing it with
the output it generates can show you how to
produce a simple document of your own.

\section{Ordinary Text}      % Produces section heading.  Lower-level
                             % sections are begun with similar
                             % \subsection and \subsubsection commands.

The ends  of words and sentences are marked
  by   spaces. It  doesn't matter how many
spaces    you type; one is as good as 100.  The
end of   a line counts as a space.

One   or more   blank lines denote the  end
of  a paragraph.

Since any number of consecutive spaces are treated
like a single one, the formatting of the input
file makes no difference to
      \LaTeX,                % The \LaTeX command generates the LaTeX logo.
but it makes a difference to you.  When you use
\LaTeX, making your input file as easy to read
as possible will be a great help as you write
your document and when you change it.  This sample
file shows how you can add comments to your own input
file.

Because printing is different from typewriting,
there are a number of things that you have to do
differently when preparing an input file than if
you were just typing the document directly.
Quotation marks like
       ``this''
have to be handled specially, as do quotes within
quotes:
       ``\,`this'            % \, separates the double and single quote.
        is what I just
        wrote, not  `that'\,''.

Dashes come in three sizes: an
       intra-word
dash, a medium dash for number ranges like
       1--2,
and a punctuation
       dash---like
this.

A sentence-ending space should be larger than the
space between words within a sentence.  You
sometimes have to type special commands in
conjunction with punctuation characters to get
this right, as in the following sentence.
       Gnats, gnus, etc.\ all  % `\ ' makes an inter-word space.
       begin with G\@.         % \@ marks end-of-sentence punctuation.
You should check the spaces after periods when
reading your output to make sure you haven't
forgotten any special cases.  Generating an
ellipsis
       \ldots\               % `\ ' is needed after `\ldots' because TeX
                             % ignores spaces after command names like \ldots
                             % made from \ + letters.
                             %
                             % Note how a `%' character causes TeX to ignore
                             % the end of the input line, so these blank lines
                             % do not start a new paragraph.
                             %
with the right spacing around the periods requires
a special command.

\LaTeX\ interprets some common characters as
commands, so you must type special commands to
generate them.  These characters include the
following:
       \$ \& \% \# \{ and \}.

In printing, text is usually emphasized with an
       \emph{italic}
type style.

\begin{em}
   A long segment of text can also be emphasized
   in this way.  Text within such a segment can be
   given \emph{additional} emphasis.
\end{em}

It is sometimes necessary to prevent \LaTeX\ from
breaking a line where it might otherwise do so.
This may be at a space, as between the ``Mr.''\ and
``Jones'' in
       ``Mr.~Jones'',        % ~ produces an unbreakable interword space.
or within a word---especially when the word is a
symbol like
       \mbox{\emph{itemnum}}
that makes little sense when hyphenated across
lines.

Footnotes\footnote{This is an example of a footnote.}
pose no problem.

\LaTeX\ is good at typesetting mathematical formulas
like
       \( x-3y + z = 7 \)
or
       \( a_{1} > x^{2n} + y^{2n} > x' \)
or
       \( \ip{A}{B} = \sum_{i} a_{i} b_{i} \).
The spaces you type in a formula are
ignored.  Remember that a letter like
       $x$                   % $ ... $  and  \( ... \)  are equivalent
is a formula when it denotes a mathematical
symbol, and it should be typed as one.

\section{Displayed Text}

Text is displayed by indenting it from the left
margin.  Quotations are commonly displayed.  There
are short quotations
\begin{quote}
   This is a short quotation.  It consists of a
   single paragraph of text.  See how it is formatted.
\end{quote}
and longer ones.
\begin{quotation}
   This is a longer quotation.  It consists of two
   paragraphs of text, neither of which are
   particularly interesting.

   This is the second paragraph of the quotation.  It
   is just as dull as the first paragraph.
\end{quotation}
Another frequently-displayed structure is a list.
The following is an example of an \emph{itemized}
list.
\begin{itemize}
   \item This is the first item of an itemized list.
         Each item in the list is marked with a ``tick''.
         You don't have to worry about what kind of tick
         mark is used.

   \item This is the second item of the list.  It
         contains another list nested inside it.  The inner
         list is an \emph{enumerated} list.
         \begin{enumerate}
            \item This is the first item of an enumerated
                  list that is nested within the itemized list.

            \item This is the second item of the inner list.
                  \LaTeX\ allows you to nest lists deeper than
                  you really should.
         \end{enumerate}
         This is the rest of the second item of the outer
         list.  It is no more interesting than any other
         part of the item.
   \item This is the third item of the list.
\end{itemize}
You can even display poetry.
\begin{verse}
   There is an environment
    for verse \\             % The \\ command separates lines
   Whose features some poets % within a stanza.
   will curse.

                             % One or more blank lines separate stanzas.

   For instead of making\\
   Them do \emph{all} line breaking, \\
   It allows them to put too many words on a line when they'd rather be
   forced to be terse.
\end{verse}

Mathematical formulas may also be displayed.  A
displayed formula
is
one-line long; multiline
formulas require special formatting instructions.
   \[  \ip{\Gamma}{\psi'} = x'' + y^{2} + z_{i}^{n}\]
Don't start a paragraph with a displayed equation,
nor make one a paragraph by itself.

\end{document}               % End of document.

\section{Bench Huge Section 1512}
% synthetic bench section 1512
% This is a sample LaTeX input file.  (Version of 12 August 2004.)
%
% A '%' character causes TeX to ignore all remaining text on the line,
% and is used for comments like this one.

\documentclass{article}      % Specifies the document class

                             % The preamble begins here.
\title{An Example Document}  % Declares the document's title.
\author{Leslie Lamport}      % Declares the author's name.
\date{January 21, 1994}      % Deleting this command produces today's date.

\newcommand{\ip}[2]{(#1, #2)}
                             % Defines \ip{arg1}{arg2} to mean
                             % (arg1, arg2).

%\newcommand{\ip}[2]{\langle #1 | #2\rangle}
                             % This is an alternative definition of
                             % \ip that is commented out.

\begin{document}             % End of preamble and beginning of text.

\maketitle                   % Produces the title.

This is an example input file.  Comparing it with
the output it generates can show you how to
produce a simple document of your own.

\section{Ordinary Text}      % Produces section heading.  Lower-level
                             % sections are begun with similar
                             % \subsection and \subsubsection commands.

The ends  of words and sentences are marked
  by   spaces. It  doesn't matter how many
spaces    you type; one is as good as 100.  The
end of   a line counts as a space.

One   or more   blank lines denote the  end
of  a paragraph.

Since any number of consecutive spaces are treated
like a single one, the formatting of the input
file makes no difference to
      \LaTeX,                % The \LaTeX command generates the LaTeX logo.
but it makes a difference to you.  When you use
\LaTeX, making your input file as easy to read
as possible will be a great help as you write
your document and when you change it.  This sample
file shows how you can add comments to your own input
file.

Because printing is different from typewriting,
there are a number of things that you have to do
differently when preparing an input file than if
you were just typing the document directly.
Quotation marks like
       ``this''
have to be handled specially, as do quotes within
quotes:
       ``\,`this'            % \, separates the double and single quote.
        is what I just
        wrote, not  `that'\,''.

Dashes come in three sizes: an
       intra-word
dash, a medium dash for number ranges like
       1--2,
and a punctuation
       dash---like
this.

A sentence-ending space should be larger than the
space between words within a sentence.  You
sometimes have to type special commands in
conjunction with punctuation characters to get
this right, as in the following sentence.
       Gnats, gnus, etc.\ all  % `\ ' makes an inter-word space.
       begin with G\@.         % \@ marks end-of-sentence punctuation.
You should check the spaces after periods when
reading your output to make sure you haven't
forgotten any special cases.  Generating an
ellipsis
       \ldots\               % `\ ' is needed after `\ldots' because TeX
                             % ignores spaces after command names like \ldots
                             % made from \ + letters.
                             %
                             % Note how a `%' character causes TeX to ignore
                             % the end of the input line, so these blank lines
                             % do not start a new paragraph.
                             %
with the right spacing around the periods requires
a special command.

\LaTeX\ interprets some common characters as
commands, so you must type special commands to
generate them.  These characters include the
following:
       \$ \& \% \# \{ and \}.

In printing, text is usually emphasized with an
       \emph{italic}
type style.

\begin{em}
   A long segment of text can also be emphasized
   in this way.  Text within such a segment can be
   given \emph{additional} emphasis.
\end{em}

It is sometimes necessary to prevent \LaTeX\ from
breaking a line where it might otherwise do so.
This may be at a space, as between the ``Mr.''\ and
``Jones'' in
       ``Mr.~Jones'',        % ~ produces an unbreakable interword space.
or within a word---especially when the word is a
symbol like
       \mbox{\emph{itemnum}}
that makes little sense when hyphenated across
lines.

Footnotes\footnote{This is an example of a footnote.}
pose no problem.

\LaTeX\ is good at typesetting mathematical formulas
like
       \( x-3y + z = 7 \)
or
       \( a_{1} > x^{2n} + y^{2n} > x' \)
or
       \( \ip{A}{B} = \sum_{i} a_{i} b_{i} \).
The spaces you type in a formula are
ignored.  Remember that a letter like
       $x$                   % $ ... $  and  \( ... \)  are equivalent
is a formula when it denotes a mathematical
symbol, and it should be typed as one.

\section{Displayed Text}

Text is displayed by indenting it from the left
margin.  Quotations are commonly displayed.  There
are short quotations
\begin{quote}
   This is a short quotation.  It consists of a
   single paragraph of text.  See how it is formatted.
\end{quote}
and longer ones.
\begin{quotation}
   This is a longer quotation.  It consists of two
   paragraphs of text, neither of which are
   particularly interesting.

   This is the second paragraph of the quotation.  It
   is just as dull as the first paragraph.
\end{quotation}
Another frequently-displayed structure is a list.
The following is an example of an \emph{itemized}
list.
\begin{itemize}
   \item This is the first item of an itemized list.
         Each item in the list is marked with a ``tick''.
         You don't have to worry about what kind of tick
         mark is used.

   \item This is the second item of the list.  It
         contains another list nested inside it.  The inner
         list is an \emph{enumerated} list.
         \begin{enumerate}
            \item This is the first item of an enumerated
                  list that is nested within the itemized list.

            \item This is the second item of the inner list.
                  \LaTeX\ allows you to nest lists deeper than
                  you really should.
         \end{enumerate}
         This is the rest of the second item of the outer
         list.  It is no more interesting than any other
         part of the item.
   \item This is the third item of the list.
\end{itemize}
You can even display poetry.
\begin{verse}
   There is an environment
    for verse \\             % The \\ command separates lines
   Whose features some poets % within a stanza.
   will curse.

                             % One or more blank lines separate stanzas.

   For instead of making\\
   Them do \emph{all} line breaking, \\
   It allows them to put too many words on a line when they'd rather be
   forced to be terse.
\end{verse}

Mathematical formulas may also be displayed.  A
displayed formula
is
one-line long; multiline
formulas require special formatting instructions.
   \[  \ip{\Gamma}{\psi'} = x'' + y^{2} + z_{i}^{n}\]
Don't start a paragraph with a displayed equation,
nor make one a paragraph by itself.

\end{document}               % End of document.

\section{Bench Huge Section 1513}
% synthetic bench section 1513
% This is a sample LaTeX input file.  (Version of 12 August 2004.)
%
% A '%' character causes TeX to ignore all remaining text on the line,
% and is used for comments like this one.

\documentclass{article}      % Specifies the document class

                             % The preamble begins here.
\title{An Example Document}  % Declares the document's title.
\author{Leslie Lamport}      % Declares the author's name.
\date{January 21, 1994}      % Deleting this command produces today's date.

\newcommand{\ip}[2]{(#1, #2)}
                             % Defines \ip{arg1}{arg2} to mean
                             % (arg1, arg2).

%\newcommand{\ip}[2]{\langle #1 | #2\rangle}
                             % This is an alternative definition of
                             % \ip that is commented out.

\begin{document}             % End of preamble and beginning of text.

\maketitle                   % Produces the title.

This is an example input file.  Comparing it with
the output it generates can show you how to
produce a simple document of your own.

\section{Ordinary Text}      % Produces section heading.  Lower-level
                             % sections are begun with similar
                             % \subsection and \subsubsection commands.

The ends  of words and sentences are marked
  by   spaces. It  doesn't matter how many
spaces    you type; one is as good as 100.  The
end of   a line counts as a space.

One   or more   blank lines denote the  end
of  a paragraph.

Since any number of consecutive spaces are treated
like a single one, the formatting of the input
file makes no difference to
      \LaTeX,                % The \LaTeX command generates the LaTeX logo.
but it makes a difference to you.  When you use
\LaTeX, making your input file as easy to read
as possible will be a great help as you write
your document and when you change it.  This sample
file shows how you can add comments to your own input
file.

Because printing is different from typewriting,
there are a number of things that you have to do
differently when preparing an input file than if
you were just typing the document directly.
Quotation marks like
       ``this''
have to be handled specially, as do quotes within
quotes:
       ``\,`this'            % \, separates the double and single quote.
        is what I just
        wrote, not  `that'\,''.

Dashes come in three sizes: an
       intra-word
dash, a medium dash for number ranges like
       1--2,
and a punctuation
       dash---like
this.

A sentence-ending space should be larger than the
space between words within a sentence.  You
sometimes have to type special commands in
conjunction with punctuation characters to get
this right, as in the following sentence.
       Gnats, gnus, etc.\ all  % `\ ' makes an inter-word space.
       begin with G\@.         % \@ marks end-of-sentence punctuation.
You should check the spaces after periods when
reading your output to make sure you haven't
forgotten any special cases.  Generating an
ellipsis
       \ldots\               % `\ ' is needed after `\ldots' because TeX
                             % ignores spaces after command names like \ldots
                             % made from \ + letters.
                             %
                             % Note how a `%' character causes TeX to ignore
                             % the end of the input line, so these blank lines
                             % do not start a new paragraph.
                             %
with the right spacing around the periods requires
a special command.

\LaTeX\ interprets some common characters as
commands, so you must type special commands to
generate them.  These characters include the
following:
       \$ \& \% \# \{ and \}.

In printing, text is usually emphasized with an
       \emph{italic}
type style.

\begin{em}
   A long segment of text can also be emphasized
   in this way.  Text within such a segment can be
   given \emph{additional} emphasis.
\end{em}

It is sometimes necessary to prevent \LaTeX\ from
breaking a line where it might otherwise do so.
This may be at a space, as between the ``Mr.''\ and
``Jones'' in
       ``Mr.~Jones'',        % ~ produces an unbreakable interword space.
or within a word---especially when the word is a
symbol like
       \mbox{\emph{itemnum}}
that makes little sense when hyphenated across
lines.

Footnotes\footnote{This is an example of a footnote.}
pose no problem.

\LaTeX\ is good at typesetting mathematical formulas
like
       \( x-3y + z = 7 \)
or
       \( a_{1} > x^{2n} + y^{2n} > x' \)
or
       \( \ip{A}{B} = \sum_{i} a_{i} b_{i} \).
The spaces you type in a formula are
ignored.  Remember that a letter like
       $x$                   % $ ... $  and  \( ... \)  are equivalent
is a formula when it denotes a mathematical
symbol, and it should be typed as one.

\section{Displayed Text}

Text is displayed by indenting it from the left
margin.  Quotations are commonly displayed.  There
are short quotations
\begin{quote}
   This is a short quotation.  It consists of a
   single paragraph of text.  See how it is formatted.
\end{quote}
and longer ones.
\begin{quotation}
   This is a longer quotation.  It consists of two
   paragraphs of text, neither of which are
   particularly interesting.

   This is the second paragraph of the quotation.  It
   is just as dull as the first paragraph.
\end{quotation}
Another frequently-displayed structure is a list.
The following is an example of an \emph{itemized}
list.
\begin{itemize}
   \item This is the first item of an itemized list.
         Each item in the list is marked with a ``tick''.
         You don't have to worry about what kind of tick
         mark is used.

   \item This is the second item of the list.  It
         contains another list nested inside it.  The inner
         list is an \emph{enumerated} list.
         \begin{enumerate}
            \item This is the first item of an enumerated
                  list that is nested within the itemized list.

            \item This is the second item of the inner list.
                  \LaTeX\ allows you to nest lists deeper than
                  you really should.
         \end{enumerate}
         This is the rest of the second item of the outer
         list.  It is no more interesting than any other
         part of the item.
   \item This is the third item of the list.
\end{itemize}
You can even display poetry.
\begin{verse}
   There is an environment
    for verse \\             % The \\ command separates lines
   Whose features some poets % within a stanza.
   will curse.

                             % One or more blank lines separate stanzas.

   For instead of making\\
   Them do \emph{all} line breaking, \\
   It allows them to put too many words on a line when they'd rather be
   forced to be terse.
\end{verse}

Mathematical formulas may also be displayed.  A
displayed formula
is
one-line long; multiline
formulas require special formatting instructions.
   \[  \ip{\Gamma}{\psi'} = x'' + y^{2} + z_{i}^{n}\]
Don't start a paragraph with a displayed equation,
nor make one a paragraph by itself.

\end{document}               % End of document.

\section{Bench Huge Section 1514}
% synthetic bench section 1514
% This is a sample LaTeX input file.  (Version of 12 August 2004.)
%
% A '%' character causes TeX to ignore all remaining text on the line,
% and is used for comments like this one.

\documentclass{article}      % Specifies the document class

                             % The preamble begins here.
\title{An Example Document}  % Declares the document's title.
\author{Leslie Lamport}      % Declares the author's name.
\date{January 21, 1994}      % Deleting this command produces today's date.

\newcommand{\ip}[2]{(#1, #2)}
                             % Defines \ip{arg1}{arg2} to mean
                             % (arg1, arg2).

%\newcommand{\ip}[2]{\langle #1 | #2\rangle}
                             % This is an alternative definition of
                             % \ip that is commented out.

\begin{document}             % End of preamble and beginning of text.

\maketitle                   % Produces the title.

This is an example input file.  Comparing it with
the output it generates can show you how to
produce a simple document of your own.

\section{Ordinary Text}      % Produces section heading.  Lower-level
                             % sections are begun with similar
                             % \subsection and \subsubsection commands.

The ends  of words and sentences are marked
  by   spaces. It  doesn't matter how many
spaces    you type; one is as good as 100.  The
end of   a line counts as a space.

One   or more   blank lines denote the  end
of  a paragraph.

Since any number of consecutive spaces are treated
like a single one, the formatting of the input
file makes no difference to
      \LaTeX,                % The \LaTeX command generates the LaTeX logo.
but it makes a difference to you.  When you use
\LaTeX, making your input file as easy to read
as possible will be a great help as you write
your document and when you change it.  This sample
file shows how you can add comments to your own input
file.

Because printing is different from typewriting,
there are a number of things that you have to do
differently when preparing an input file than if
you were just typing the document directly.
Quotation marks like
       ``this''
have to be handled specially, as do quotes within
quotes:
       ``\,`this'            % \, separates the double and single quote.
        is what I just
        wrote, not  `that'\,''.

Dashes come in three sizes: an
       intra-word
dash, a medium dash for number ranges like
       1--2,
and a punctuation
       dash---like
this.

A sentence-ending space should be larger than the
space between words within a sentence.  You
sometimes have to type special commands in
conjunction with punctuation characters to get
this right, as in the following sentence.
       Gnats, gnus, etc.\ all  % `\ ' makes an inter-word space.
       begin with G\@.         % \@ marks end-of-sentence punctuation.
You should check the spaces after periods when
reading your output to make sure you haven't
forgotten any special cases.  Generating an
ellipsis
       \ldots\               % `\ ' is needed after `\ldots' because TeX
                             % ignores spaces after command names like \ldots
                             % made from \ + letters.
                             %
                             % Note how a `%' character causes TeX to ignore
                             % the end of the input line, so these blank lines
                             % do not start a new paragraph.
                             %
with the right spacing around the periods requires
a special command.

\LaTeX\ interprets some common characters as
commands, so you must type special commands to
generate them.  These characters include the
following:
       \$ \& \% \# \{ and \}.

In printing, text is usually emphasized with an
       \emph{italic}
type style.

\begin{em}
   A long segment of text can also be emphasized
   in this way.  Text within such a segment can be
   given \emph{additional} emphasis.
\end{em}

It is sometimes necessary to prevent \LaTeX\ from
breaking a line where it might otherwise do so.
This may be at a space, as between the ``Mr.''\ and
``Jones'' in
       ``Mr.~Jones'',        % ~ produces an unbreakable interword space.
or within a word---especially when the word is a
symbol like
       \mbox{\emph{itemnum}}
that makes little sense when hyphenated across
lines.

Footnotes\footnote{This is an example of a footnote.}
pose no problem.

\LaTeX\ is good at typesetting mathematical formulas
like
       \( x-3y + z = 7 \)
or
       \( a_{1} > x^{2n} + y^{2n} > x' \)
or
       \( \ip{A}{B} = \sum_{i} a_{i} b_{i} \).
The spaces you type in a formula are
ignored.  Remember that a letter like
       $x$                   % $ ... $  and  \( ... \)  are equivalent
is a formula when it denotes a mathematical
symbol, and it should be typed as one.

\section{Displayed Text}

Text is displayed by indenting it from the left
margin.  Quotations are commonly displayed.  There
are short quotations
\begin{quote}
   This is a short quotation.  It consists of a
   single paragraph of text.  See how it is formatted.
\end{quote}
and longer ones.
\begin{quotation}
   This is a longer quotation.  It consists of two
   paragraphs of text, neither of which are
   particularly interesting.

   This is the second paragraph of the quotation.  It
   is just as dull as the first paragraph.
\end{quotation}
Another frequently-displayed structure is a list.
The following is an example of an \emph{itemized}
list.
\begin{itemize}
   \item This is the first item of an itemized list.
         Each item in the list is marked with a ``tick''.
         You don't have to worry about what kind of tick
         mark is used.

   \item This is the second item of the list.  It
         contains another list nested inside it.  The inner
         list is an \emph{enumerated} list.
         \begin{enumerate}
            \item This is the first item of an enumerated
                  list that is nested within the itemized list.

            \item This is the second item of the inner list.
                  \LaTeX\ allows you to nest lists deeper than
                  you really should.
         \end{enumerate}
         This is the rest of the second item of the outer
         list.  It is no more interesting than any other
         part of the item.
   \item This is the third item of the list.
\end{itemize}
You can even display poetry.
\begin{verse}
   There is an environment
    for verse \\             % The \\ command separates lines
   Whose features some poets % within a stanza.
   will curse.

                             % One or more blank lines separate stanzas.

   For instead of making\\
   Them do \emph{all} line breaking, \\
   It allows them to put too many words on a line when they'd rather be
   forced to be terse.
\end{verse}

Mathematical formulas may also be displayed.  A
displayed formula
is
one-line long; multiline
formulas require special formatting instructions.
   \[  \ip{\Gamma}{\psi'} = x'' + y^{2} + z_{i}^{n}\]
Don't start a paragraph with a displayed equation,
nor make one a paragraph by itself.

\end{document}               % End of document.

\section{Bench Huge Section 1515}
% synthetic bench section 1515
% This is a sample LaTeX input file.  (Version of 12 August 2004.)
%
% A '%' character causes TeX to ignore all remaining text on the line,
% and is used for comments like this one.

\documentclass{article}      % Specifies the document class

                             % The preamble begins here.
\title{An Example Document}  % Declares the document's title.
\author{Leslie Lamport}      % Declares the author's name.
\date{January 21, 1994}      % Deleting this command produces today's date.

\newcommand{\ip}[2]{(#1, #2)}
                             % Defines \ip{arg1}{arg2} to mean
                             % (arg1, arg2).

%\newcommand{\ip}[2]{\langle #1 | #2\rangle}
                             % This is an alternative definition of
                             % \ip that is commented out.

\begin{document}             % End of preamble and beginning of text.

\maketitle                   % Produces the title.

This is an example input file.  Comparing it with
the output it generates can show you how to
produce a simple document of your own.

\section{Ordinary Text}      % Produces section heading.  Lower-level
                             % sections are begun with similar
                             % \subsection and \subsubsection commands.

The ends  of words and sentences are marked
  by   spaces. It  doesn't matter how many
spaces    you type; one is as good as 100.  The
end of   a line counts as a space.

One   or more   blank lines denote the  end
of  a paragraph.

Since any number of consecutive spaces are treated
like a single one, the formatting of the input
file makes no difference to
      \LaTeX,                % The \LaTeX command generates the LaTeX logo.
but it makes a difference to you.  When you use
\LaTeX, making your input file as easy to read
as possible will be a great help as you write
your document and when you change it.  This sample
file shows how you can add comments to your own input
file.

Because printing is different from typewriting,
there are a number of things that you have to do
differently when preparing an input file than if
you were just typing the document directly.
Quotation marks like
       ``this''
have to be handled specially, as do quotes within
quotes:
       ``\,`this'            % \, separates the double and single quote.
        is what I just
        wrote, not  `that'\,''.

Dashes come in three sizes: an
       intra-word
dash, a medium dash for number ranges like
       1--2,
and a punctuation
       dash---like
this.

A sentence-ending space should be larger than the
space between words within a sentence.  You
sometimes have to type special commands in
conjunction with punctuation characters to get
this right, as in the following sentence.
       Gnats, gnus, etc.\ all  % `\ ' makes an inter-word space.
       begin with G\@.         % \@ marks end-of-sentence punctuation.
You should check the spaces after periods when
reading your output to make sure you haven't
forgotten any special cases.  Generating an
ellipsis
       \ldots\               % `\ ' is needed after `\ldots' because TeX
                             % ignores spaces after command names like \ldots
                             % made from \ + letters.
                             %
                             % Note how a `%' character causes TeX to ignore
                             % the end of the input line, so these blank lines
                             % do not start a new paragraph.
                             %
with the right spacing around the periods requires
a special command.

\LaTeX\ interprets some common characters as
commands, so you must type special commands to
generate them.  These characters include the
following:
       \$ \& \% \# \{ and \}.

In printing, text is usually emphasized with an
       \emph{italic}
type style.

\begin{em}
   A long segment of text can also be emphasized
   in this way.  Text within such a segment can be
   given \emph{additional} emphasis.
\end{em}

It is sometimes necessary to prevent \LaTeX\ from
breaking a line where it might otherwise do so.
This may be at a space, as between the ``Mr.''\ and
``Jones'' in
       ``Mr.~Jones'',        % ~ produces an unbreakable interword space.
or within a word---especially when the word is a
symbol like
       \mbox{\emph{itemnum}}
that makes little sense when hyphenated across
lines.

Footnotes\footnote{This is an example of a footnote.}
pose no problem.

\LaTeX\ is good at typesetting mathematical formulas
like
       \( x-3y + z = 7 \)
or
       \( a_{1} > x^{2n} + y^{2n} > x' \)
or
       \( \ip{A}{B} = \sum_{i} a_{i} b_{i} \).
The spaces you type in a formula are
ignored.  Remember that a letter like
       $x$                   % $ ... $  and  \( ... \)  are equivalent
is a formula when it denotes a mathematical
symbol, and it should be typed as one.

\section{Displayed Text}

Text is displayed by indenting it from the left
margin.  Quotations are commonly displayed.  There
are short quotations
\begin{quote}
   This is a short quotation.  It consists of a
   single paragraph of text.  See how it is formatted.
\end{quote}
and longer ones.
\begin{quotation}
   This is a longer quotation.  It consists of two
   paragraphs of text, neither of which are
   particularly interesting.

   This is the second paragraph of the quotation.  It
   is just as dull as the first paragraph.
\end{quotation}
Another frequently-displayed structure is a list.
The following is an example of an \emph{itemized}
list.
\begin{itemize}
   \item This is the first item of an itemized list.
         Each item in the list is marked with a ``tick''.
         You don't have to worry about what kind of tick
         mark is used.

   \item This is the second item of the list.  It
         contains another list nested inside it.  The inner
         list is an \emph{enumerated} list.
         \begin{enumerate}
            \item This is the first item of an enumerated
                  list that is nested within the itemized list.

            \item This is the second item of the inner list.
                  \LaTeX\ allows you to nest lists deeper than
                  you really should.
         \end{enumerate}
         This is the rest of the second item of the outer
         list.  It is no more interesting than any other
         part of the item.
   \item This is the third item of the list.
\end{itemize}
You can even display poetry.
\begin{verse}
   There is an environment
    for verse \\             % The \\ command separates lines
   Whose features some poets % within a stanza.
   will curse.

                             % One or more blank lines separate stanzas.

   For instead of making\\
   Them do \emph{all} line breaking, \\
   It allows them to put too many words on a line when they'd rather be
   forced to be terse.
\end{verse}

Mathematical formulas may also be displayed.  A
displayed formula
is
one-line long; multiline
formulas require special formatting instructions.
   \[  \ip{\Gamma}{\psi'} = x'' + y^{2} + z_{i}^{n}\]
Don't start a paragraph with a displayed equation,
nor make one a paragraph by itself.

\end{document}               % End of document.

\section{Bench Huge Section 1516}
% synthetic bench section 1516
% This is a sample LaTeX input file.  (Version of 12 August 2004.)
%
% A '%' character causes TeX to ignore all remaining text on the line,
% and is used for comments like this one.

\documentclass{article}      % Specifies the document class

                             % The preamble begins here.
\title{An Example Document}  % Declares the document's title.
\author{Leslie Lamport}      % Declares the author's name.
\date{January 21, 1994}      % Deleting this command produces today's date.

\newcommand{\ip}[2]{(#1, #2)}
                             % Defines \ip{arg1}{arg2} to mean
                             % (arg1, arg2).

%\newcommand{\ip}[2]{\langle #1 | #2\rangle}
                             % This is an alternative definition of
                             % \ip that is commented out.

\begin{document}             % End of preamble and beginning of text.

\maketitle                   % Produces the title.

This is an example input file.  Comparing it with
the output it generates can show you how to
produce a simple document of your own.

\section{Ordinary Text}      % Produces section heading.  Lower-level
                             % sections are begun with similar
                             % \subsection and \subsubsection commands.

The ends  of words and sentences are marked
  by   spaces. It  doesn't matter how many
spaces    you type; one is as good as 100.  The
end of   a line counts as a space.

One   or more   blank lines denote the  end
of  a paragraph.

Since any number of consecutive spaces are treated
like a single one, the formatting of the input
file makes no difference to
      \LaTeX,                % The \LaTeX command generates the LaTeX logo.
but it makes a difference to you.  When you use
\LaTeX, making your input file as easy to read
as possible will be a great help as you write
your document and when you change it.  This sample
file shows how you can add comments to your own input
file.

Because printing is different from typewriting,
there are a number of things that you have to do
differently when preparing an input file than if
you were just typing the document directly.
Quotation marks like
       ``this''
have to be handled specially, as do quotes within
quotes:
       ``\,`this'            % \, separates the double and single quote.
        is what I just
        wrote, not  `that'\,''.

Dashes come in three sizes: an
       intra-word
dash, a medium dash for number ranges like
       1--2,
and a punctuation
       dash---like
this.

A sentence-ending space should be larger than the
space between words within a sentence.  You
sometimes have to type special commands in
conjunction with punctuation characters to get
this right, as in the following sentence.
       Gnats, gnus, etc.\ all  % `\ ' makes an inter-word space.
       begin with G\@.         % \@ marks end-of-sentence punctuation.
You should check the spaces after periods when
reading your output to make sure you haven't
forgotten any special cases.  Generating an
ellipsis
       \ldots\               % `\ ' is needed after `\ldots' because TeX
                             % ignores spaces after command names like \ldots
                             % made from \ + letters.
                             %
                             % Note how a `%' character causes TeX to ignore
                             % the end of the input line, so these blank lines
                             % do not start a new paragraph.
                             %
with the right spacing around the periods requires
a special command.

\LaTeX\ interprets some common characters as
commands, so you must type special commands to
generate them.  These characters include the
following:
       \$ \& \% \# \{ and \}.

In printing, text is usually emphasized with an
       \emph{italic}
type style.

\begin{em}
   A long segment of text can also be emphasized
   in this way.  Text within such a segment can be
   given \emph{additional} emphasis.
\end{em}

It is sometimes necessary to prevent \LaTeX\ from
breaking a line where it might otherwise do so.
This may be at a space, as between the ``Mr.''\ and
``Jones'' in
       ``Mr.~Jones'',        % ~ produces an unbreakable interword space.
or within a word---especially when the word is a
symbol like
       \mbox{\emph{itemnum}}
that makes little sense when hyphenated across
lines.

Footnotes\footnote{This is an example of a footnote.}
pose no problem.

\LaTeX\ is good at typesetting mathematical formulas
like
       \( x-3y + z = 7 \)
or
       \( a_{1} > x^{2n} + y^{2n} > x' \)
or
       \( \ip{A}{B} = \sum_{i} a_{i} b_{i} \).
The spaces you type in a formula are
ignored.  Remember that a letter like
       $x$                   % $ ... $  and  \( ... \)  are equivalent
is a formula when it denotes a mathematical
symbol, and it should be typed as one.

\section{Displayed Text}

Text is displayed by indenting it from the left
margin.  Quotations are commonly displayed.  There
are short quotations
\begin{quote}
   This is a short quotation.  It consists of a
   single paragraph of text.  See how it is formatted.
\end{quote}
and longer ones.
\begin{quotation}
   This is a longer quotation.  It consists of two
   paragraphs of text, neither of which are
   particularly interesting.

   This is the second paragraph of the quotation.  It
   is just as dull as the first paragraph.
\end{quotation}
Another frequently-displayed structure is a list.
The following is an example of an \emph{itemized}
list.
\begin{itemize}
   \item This is the first item of an itemized list.
         Each item in the list is marked with a ``tick''.
         You don't have to worry about what kind of tick
         mark is used.

   \item This is the second item of the list.  It
         contains another list nested inside it.  The inner
         list is an \emph{enumerated} list.
         \begin{enumerate}
            \item This is the first item of an enumerated
                  list that is nested within the itemized list.

            \item This is the second item of the inner list.
                  \LaTeX\ allows you to nest lists deeper than
                  you really should.
         \end{enumerate}
         This is the rest of the second item of the outer
         list.  It is no more interesting than any other
         part of the item.
   \item This is the third item of the list.
\end{itemize}
You can even display poetry.
\begin{verse}
   There is an environment
    for verse \\             % The \\ command separates lines
   Whose features some poets % within a stanza.
   will curse.

                             % One or more blank lines separate stanzas.

   For instead of making\\
   Them do \emph{all} line breaking, \\
   It allows them to put too many words on a line when they'd rather be
   forced to be terse.
\end{verse}

Mathematical formulas may also be displayed.  A
displayed formula
is
one-line long; multiline
formulas require special formatting instructions.
   \[  \ip{\Gamma}{\psi'} = x'' + y^{2} + z_{i}^{n}\]
Don't start a paragraph with a displayed equation,
nor make one a paragraph by itself.

\end{document}               % End of document.

\section{Bench Huge Section 1517}
% synthetic bench section 1517
% This is a sample LaTeX input file.  (Version of 12 August 2004.)
%
% A '%' character causes TeX to ignore all remaining text on the line,
% and is used for comments like this one.

\documentclass{article}      % Specifies the document class

                             % The preamble begins here.
\title{An Example Document}  % Declares the document's title.
\author{Leslie Lamport}      % Declares the author's name.
\date{January 21, 1994}      % Deleting this command produces today's date.

\newcommand{\ip}[2]{(#1, #2)}
                             % Defines \ip{arg1}{arg2} to mean
                             % (arg1, arg2).

%\newcommand{\ip}[2]{\langle #1 | #2\rangle}
                             % This is an alternative definition of
                             % \ip that is commented out.

\begin{document}             % End of preamble and beginning of text.

\maketitle                   % Produces the title.

This is an example input file.  Comparing it with
the output it generates can show you how to
produce a simple document of your own.

\section{Ordinary Text}      % Produces section heading.  Lower-level
                             % sections are begun with similar
                             % \subsection and \subsubsection commands.

The ends  of words and sentences are marked
  by   spaces. It  doesn't matter how many
spaces    you type; one is as good as 100.  The
end of   a line counts as a space.

One   or more   blank lines denote the  end
of  a paragraph.

Since any number of consecutive spaces are treated
like a single one, the formatting of the input
file makes no difference to
      \LaTeX,                % The \LaTeX command generates the LaTeX logo.
but it makes a difference to you.  When you use
\LaTeX, making your input file as easy to read
as possible will be a great help as you write
your document and when you change it.  This sample
file shows how you can add comments to your own input
file.

Because printing is different from typewriting,
there are a number of things that you have to do
differently when preparing an input file than if
you were just typing the document directly.
Quotation marks like
       ``this''
have to be handled specially, as do quotes within
quotes:
       ``\,`this'            % \, separates the double and single quote.
        is what I just
        wrote, not  `that'\,''.

Dashes come in three sizes: an
       intra-word
dash, a medium dash for number ranges like
       1--2,
and a punctuation
       dash---like
this.

A sentence-ending space should be larger than the
space between words within a sentence.  You
sometimes have to type special commands in
conjunction with punctuation characters to get
this right, as in the following sentence.
       Gnats, gnus, etc.\ all  % `\ ' makes an inter-word space.
       begin with G\@.         % \@ marks end-of-sentence punctuation.
You should check the spaces after periods when
reading your output to make sure you haven't
forgotten any special cases.  Generating an
ellipsis
       \ldots\               % `\ ' is needed after `\ldots' because TeX
                             % ignores spaces after command names like \ldots
                             % made from \ + letters.
                             %
                             % Note how a `%' character causes TeX to ignore
                             % the end of the input line, so these blank lines
                             % do not start a new paragraph.
                             %
with the right spacing around the periods requires
a special command.

\LaTeX\ interprets some common characters as
commands, so you must type special commands to
generate them.  These characters include the
following:
       \$ \& \% \# \{ and \}.

In printing, text is usually emphasized with an
       \emph{italic}
type style.

\begin{em}
   A long segment of text can also be emphasized
   in this way.  Text within such a segment can be
   given \emph{additional} emphasis.
\end{em}

It is sometimes necessary to prevent \LaTeX\ from
breaking a line where it might otherwise do so.
This may be at a space, as between the ``Mr.''\ and
``Jones'' in
       ``Mr.~Jones'',        % ~ produces an unbreakable interword space.
or within a word---especially when the word is a
symbol like
       \mbox{\emph{itemnum}}
that makes little sense when hyphenated across
lines.

Footnotes\footnote{This is an example of a footnote.}
pose no problem.

\LaTeX\ is good at typesetting mathematical formulas
like
       \( x-3y + z = 7 \)
or
       \( a_{1} > x^{2n} + y^{2n} > x' \)
or
       \( \ip{A}{B} = \sum_{i} a_{i} b_{i} \).
The spaces you type in a formula are
ignored.  Remember that a letter like
       $x$                   % $ ... $  and  \( ... \)  are equivalent
is a formula when it denotes a mathematical
symbol, and it should be typed as one.

\section{Displayed Text}

Text is displayed by indenting it from the left
margin.  Quotations are commonly displayed.  There
are short quotations
\begin{quote}
   This is a short quotation.  It consists of a
   single paragraph of text.  See how it is formatted.
\end{quote}
and longer ones.
\begin{quotation}
   This is a longer quotation.  It consists of two
   paragraphs of text, neither of which are
   particularly interesting.

   This is the second paragraph of the quotation.  It
   is just as dull as the first paragraph.
\end{quotation}
Another frequently-displayed structure is a list.
The following is an example of an \emph{itemized}
list.
\begin{itemize}
   \item This is the first item of an itemized list.
         Each item in the list is marked with a ``tick''.
         You don't have to worry about what kind of tick
         mark is used.

   \item This is the second item of the list.  It
         contains another list nested inside it.  The inner
         list is an \emph{enumerated} list.
         \begin{enumerate}
            \item This is the first item of an enumerated
                  list that is nested within the itemized list.

            \item This is the second item of the inner list.
                  \LaTeX\ allows you to nest lists deeper than
                  you really should.
         \end{enumerate}
         This is the rest of the second item of the outer
         list.  It is no more interesting than any other
         part of the item.
   \item This is the third item of the list.
\end{itemize}
You can even display poetry.
\begin{verse}
   There is an environment
    for verse \\             % The \\ command separates lines
   Whose features some poets % within a stanza.
   will curse.

                             % One or more blank lines separate stanzas.

   For instead of making\\
   Them do \emph{all} line breaking, \\
   It allows them to put too many words on a line when they'd rather be
   forced to be terse.
\end{verse}

Mathematical formulas may also be displayed.  A
displayed formula
is
one-line long; multiline
formulas require special formatting instructions.
   \[  \ip{\Gamma}{\psi'} = x'' + y^{2} + z_{i}^{n}\]
Don't start a paragraph with a displayed equation,
nor make one a paragraph by itself.

\end{document}               % End of document.

\section{Bench Huge Section 1518}
% synthetic bench section 1518
% This is a sample LaTeX input file.  (Version of 12 August 2004.)
%
% A '%' character causes TeX to ignore all remaining text on the line,
% and is used for comments like this one.

\documentclass{article}      % Specifies the document class

                             % The preamble begins here.
\title{An Example Document}  % Declares the document's title.
\author{Leslie Lamport}      % Declares the author's name.
\date{January 21, 1994}      % Deleting this command produces today's date.

\newcommand{\ip}[2]{(#1, #2)}
                             % Defines \ip{arg1}{arg2} to mean
                             % (arg1, arg2).

%\newcommand{\ip}[2]{\langle #1 | #2\rangle}
                             % This is an alternative definition of
                             % \ip that is commented out.

\begin{document}             % End of preamble and beginning of text.

\maketitle                   % Produces the title.

This is an example input file.  Comparing it with
the output it generates can show you how to
produce a simple document of your own.

\section{Ordinary Text}      % Produces section heading.  Lower-level
                             % sections are begun with similar
                             % \subsection and \subsubsection commands.

The ends  of words and sentences are marked
  by   spaces. It  doesn't matter how many
spaces    you type; one is as good as 100.  The
end of   a line counts as a space.

One   or more   blank lines denote the  end
of  a paragraph.

Since any number of consecutive spaces are treated
like a single one, the formatting of the input
file makes no difference to
      \LaTeX,                % The \LaTeX command generates the LaTeX logo.
but it makes a difference to you.  When you use
\LaTeX, making your input file as easy to read
as possible will be a great help as you write
your document and when you change it.  This sample
file shows how you can add comments to your own input
file.

Because printing is different from typewriting,
there are a number of things that you have to do
differently when preparing an input file than if
you were just typing the document directly.
Quotation marks like
       ``this''
have to be handled specially, as do quotes within
quotes:
       ``\,`this'            % \, separates the double and single quote.
        is what I just
        wrote, not  `that'\,''.

Dashes come in three sizes: an
       intra-word
dash, a medium dash for number ranges like
       1--2,
and a punctuation
       dash---like
this.

A sentence-ending space should be larger than the
space between words within a sentence.  You
sometimes have to type special commands in
conjunction with punctuation characters to get
this right, as in the following sentence.
       Gnats, gnus, etc.\ all  % `\ ' makes an inter-word space.
       begin with G\@.         % \@ marks end-of-sentence punctuation.
You should check the spaces after periods when
reading your output to make sure you haven't
forgotten any special cases.  Generating an
ellipsis
       \ldots\               % `\ ' is needed after `\ldots' because TeX
                             % ignores spaces after command names like \ldots
                             % made from \ + letters.
                             %
                             % Note how a `%' character causes TeX to ignore
                             % the end of the input line, so these blank lines
                             % do not start a new paragraph.
                             %
with the right spacing around the periods requires
a special command.

\LaTeX\ interprets some common characters as
commands, so you must type special commands to
generate them.  These characters include the
following:
       \$ \& \% \# \{ and \}.

In printing, text is usually emphasized with an
       \emph{italic}
type style.

\begin{em}
   A long segment of text can also be emphasized
   in this way.  Text within such a segment can be
   given \emph{additional} emphasis.
\end{em}

It is sometimes necessary to prevent \LaTeX\ from
breaking a line where it might otherwise do so.
This may be at a space, as between the ``Mr.''\ and
``Jones'' in
       ``Mr.~Jones'',        % ~ produces an unbreakable interword space.
or within a word---especially when the word is a
symbol like
       \mbox{\emph{itemnum}}
that makes little sense when hyphenated across
lines.

Footnotes\footnote{This is an example of a footnote.}
pose no problem.

\LaTeX\ is good at typesetting mathematical formulas
like
       \( x-3y + z = 7 \)
or
       \( a_{1} > x^{2n} + y^{2n} > x' \)
or
       \( \ip{A}{B} = \sum_{i} a_{i} b_{i} \).
The spaces you type in a formula are
ignored.  Remember that a letter like
       $x$                   % $ ... $  and  \( ... \)  are equivalent
is a formula when it denotes a mathematical
symbol, and it should be typed as one.

\section{Displayed Text}

Text is displayed by indenting it from the left
margin.  Quotations are commonly displayed.  There
are short quotations
\begin{quote}
   This is a short quotation.  It consists of a
   single paragraph of text.  See how it is formatted.
\end{quote}
and longer ones.
\begin{quotation}
   This is a longer quotation.  It consists of two
   paragraphs of text, neither of which are
   particularly interesting.

   This is the second paragraph of the quotation.  It
   is just as dull as the first paragraph.
\end{quotation}
Another frequently-displayed structure is a list.
The following is an example of an \emph{itemized}
list.
\begin{itemize}
   \item This is the first item of an itemized list.
         Each item in the list is marked with a ``tick''.
         You don't have to worry about what kind of tick
         mark is used.

   \item This is the second item of the list.  It
         contains another list nested inside it.  The inner
         list is an \emph{enumerated} list.
         \begin{enumerate}
            \item This is the first item of an enumerated
                  list that is nested within the itemized list.

            \item This is the second item of the inner list.
                  \LaTeX\ allows you to nest lists deeper than
                  you really should.
         \end{enumerate}
         This is the rest of the second item of the outer
         list.  It is no more interesting than any other
         part of the item.
   \item This is the third item of the list.
\end{itemize}
You can even display poetry.
\begin{verse}
   There is an environment
    for verse \\             % The \\ command separates lines
   Whose features some poets % within a stanza.
   will curse.

                             % One or more blank lines separate stanzas.

   For instead of making\\
   Them do \emph{all} line breaking, \\
   It allows them to put too many words on a line when they'd rather be
   forced to be terse.
\end{verse}

Mathematical formulas may also be displayed.  A
displayed formula
is
one-line long; multiline
formulas require special formatting instructions.
   \[  \ip{\Gamma}{\psi'} = x'' + y^{2} + z_{i}^{n}\]
Don't start a paragraph with a displayed equation,
nor make one a paragraph by itself.

\end{document}               % End of document.

\section{Bench Huge Section 1519}
% synthetic bench section 1519
% This is a sample LaTeX input file.  (Version of 12 August 2004.)
%
% A '%' character causes TeX to ignore all remaining text on the line,
% and is used for comments like this one.

\documentclass{article}      % Specifies the document class

                             % The preamble begins here.
\title{An Example Document}  % Declares the document's title.
\author{Leslie Lamport}      % Declares the author's name.
\date{January 21, 1994}      % Deleting this command produces today's date.

\newcommand{\ip}[2]{(#1, #2)}
                             % Defines \ip{arg1}{arg2} to mean
                             % (arg1, arg2).

%\newcommand{\ip}[2]{\langle #1 | #2\rangle}
                             % This is an alternative definition of
                             % \ip that is commented out.

\begin{document}             % End of preamble and beginning of text.

\maketitle                   % Produces the title.

This is an example input file.  Comparing it with
the output it generates can show you how to
produce a simple document of your own.

\section{Ordinary Text}      % Produces section heading.  Lower-level
                             % sections are begun with similar
                             % \subsection and \subsubsection commands.

The ends  of words and sentences are marked
  by   spaces. It  doesn't matter how many
spaces    you type; one is as good as 100.  The
end of   a line counts as a space.

One   or more   blank lines denote the  end
of  a paragraph.

Since any number of consecutive spaces are treated
like a single one, the formatting of the input
file makes no difference to
      \LaTeX,                % The \LaTeX command generates the LaTeX logo.
but it makes a difference to you.  When you use
\LaTeX, making your input file as easy to read
as possible will be a great help as you write
your document and when you change it.  This sample
file shows how you can add comments to your own input
file.

Because printing is different from typewriting,
there are a number of things that you have to do
differently when preparing an input file than if
you were just typing the document directly.
Quotation marks like
       ``this''
have to be handled specially, as do quotes within
quotes:
       ``\,`this'            % \, separates the double and single quote.
        is what I just
        wrote, not  `that'\,''.

Dashes come in three sizes: an
       intra-word
dash, a medium dash for number ranges like
       1--2,
and a punctuation
       dash---like
this.

A sentence-ending space should be larger than the
space between words within a sentence.  You
sometimes have to type special commands in
conjunction with punctuation characters to get
this right, as in the following sentence.
       Gnats, gnus, etc.\ all  % `\ ' makes an inter-word space.
       begin with G\@.         % \@ marks end-of-sentence punctuation.
You should check the spaces after periods when
reading your output to make sure you haven't
forgotten any special cases.  Generating an
ellipsis
       \ldots\               % `\ ' is needed after `\ldots' because TeX
                             % ignores spaces after command names like \ldots
                             % made from \ + letters.
                             %
                             % Note how a `%' character causes TeX to ignore
                             % the end of the input line, so these blank lines
                             % do not start a new paragraph.
                             %
with the right spacing around the periods requires
a special command.

\LaTeX\ interprets some common characters as
commands, so you must type special commands to
generate them.  These characters include the
following:
       \$ \& \% \# \{ and \}.

In printing, text is usually emphasized with an
       \emph{italic}
type style.

\begin{em}
   A long segment of text can also be emphasized
   in this way.  Text within such a segment can be
   given \emph{additional} emphasis.
\end{em}

It is sometimes necessary to prevent \LaTeX\ from
breaking a line where it might otherwise do so.
This may be at a space, as between the ``Mr.''\ and
``Jones'' in
       ``Mr.~Jones'',        % ~ produces an unbreakable interword space.
or within a word---especially when the word is a
symbol like
       \mbox{\emph{itemnum}}
that makes little sense when hyphenated across
lines.

Footnotes\footnote{This is an example of a footnote.}
pose no problem.

\LaTeX\ is good at typesetting mathematical formulas
like
       \( x-3y + z = 7 \)
or
       \( a_{1} > x^{2n} + y^{2n} > x' \)
or
       \( \ip{A}{B} = \sum_{i} a_{i} b_{i} \).
The spaces you type in a formula are
ignored.  Remember that a letter like
       $x$                   % $ ... $  and  \( ... \)  are equivalent
is a formula when it denotes a mathematical
symbol, and it should be typed as one.

\section{Displayed Text}

Text is displayed by indenting it from the left
margin.  Quotations are commonly displayed.  There
are short quotations
\begin{quote}
   This is a short quotation.  It consists of a
   single paragraph of text.  See how it is formatted.
\end{quote}
and longer ones.
\begin{quotation}
   This is a longer quotation.  It consists of two
   paragraphs of text, neither of which are
   particularly interesting.

   This is the second paragraph of the quotation.  It
   is just as dull as the first paragraph.
\end{quotation}
Another frequently-displayed structure is a list.
The following is an example of an \emph{itemized}
list.
\begin{itemize}
   \item This is the first item of an itemized list.
         Each item in the list is marked with a ``tick''.
         You don't have to worry about what kind of tick
         mark is used.

   \item This is the second item of the list.  It
         contains another list nested inside it.  The inner
         list is an \emph{enumerated} list.
         \begin{enumerate}
            \item This is the first item of an enumerated
                  list that is nested within the itemized list.

            \item This is the second item of the inner list.
                  \LaTeX\ allows you to nest lists deeper than
                  you really should.
         \end{enumerate}
         This is the rest of the second item of the outer
         list.  It is no more interesting than any other
         part of the item.
   \item This is the third item of the list.
\end{itemize}
You can even display poetry.
\begin{verse}
   There is an environment
    for verse \\             % The \\ command separates lines
   Whose features some poets % within a stanza.
   will curse.

                             % One or more blank lines separate stanzas.

   For instead of making\\
   Them do \emph{all} line breaking, \\
   It allows them to put too many words on a line when they'd rather be
   forced to be terse.
\end{verse}

Mathematical formulas may also be displayed.  A
displayed formula
is
one-line long; multiline
formulas require special formatting instructions.
   \[  \ip{\Gamma}{\psi'} = x'' + y^{2} + z_{i}^{n}\]
Don't start a paragraph with a displayed equation,
nor make one a paragraph by itself.

\end{document}               % End of document.

\section{Bench Huge Section 1520}
% synthetic bench section 1520
% This is a sample LaTeX input file.  (Version of 12 August 2004.)
%
% A '%' character causes TeX to ignore all remaining text on the line,
% and is used for comments like this one.

\documentclass{article}      % Specifies the document class

                             % The preamble begins here.
\title{An Example Document}  % Declares the document's title.
\author{Leslie Lamport}      % Declares the author's name.
\date{January 21, 1994}      % Deleting this command produces today's date.

\newcommand{\ip}[2]{(#1, #2)}
                             % Defines \ip{arg1}{arg2} to mean
                             % (arg1, arg2).

%\newcommand{\ip}[2]{\langle #1 | #2\rangle}
                             % This is an alternative definition of
                             % \ip that is commented out.

\begin{document}             % End of preamble and beginning of text.

\maketitle                   % Produces the title.

This is an example input file.  Comparing it with
the output it generates can show you how to
produce a simple document of your own.

\section{Ordinary Text}      % Produces section heading.  Lower-level
                             % sections are begun with similar
                             % \subsection and \subsubsection commands.

The ends  of words and sentences are marked
  by   spaces. It  doesn't matter how many
spaces    you type; one is as good as 100.  The
end of   a line counts as a space.

One   or more   blank lines denote the  end
of  a paragraph.

Since any number of consecutive spaces are treated
like a single one, the formatting of the input
file makes no difference to
      \LaTeX,                % The \LaTeX command generates the LaTeX logo.
but it makes a difference to you.  When you use
\LaTeX, making your input file as easy to read
as possible will be a great help as you write
your document and when you change it.  This sample
file shows how you can add comments to your own input
file.

Because printing is different from typewriting,
there are a number of things that you have to do
differently when preparing an input file than if
you were just typing the document directly.
Quotation marks like
       ``this''
have to be handled specially, as do quotes within
quotes:
       ``\,`this'            % \, separates the double and single quote.
        is what I just
        wrote, not  `that'\,''.

Dashes come in three sizes: an
       intra-word
dash, a medium dash for number ranges like
       1--2,
and a punctuation
       dash---like
this.

A sentence-ending space should be larger than the
space between words within a sentence.  You
sometimes have to type special commands in
conjunction with punctuation characters to get
this right, as in the following sentence.
       Gnats, gnus, etc.\ all  % `\ ' makes an inter-word space.
       begin with G\@.         % \@ marks end-of-sentence punctuation.
You should check the spaces after periods when
reading your output to make sure you haven't
forgotten any special cases.  Generating an
ellipsis
       \ldots\               % `\ ' is needed after `\ldots' because TeX
                             % ignores spaces after command names like \ldots
                             % made from \ + letters.
                             %
                             % Note how a `%' character causes TeX to ignore
                             % the end of the input line, so these blank lines
                             % do not start a new paragraph.
                             %
with the right spacing around the periods requires
a special command.

\LaTeX\ interprets some common characters as
commands, so you must type special commands to
generate them.  These characters include the
following:
       \$ \& \% \# \{ and \}.

In printing, text is usually emphasized with an
       \emph{italic}
type style.

\begin{em}
   A long segment of text can also be emphasized
   in this way.  Text within such a segment can be
   given \emph{additional} emphasis.
\end{em}

It is sometimes necessary to prevent \LaTeX\ from
breaking a line where it might otherwise do so.
This may be at a space, as between the ``Mr.''\ and
``Jones'' in
       ``Mr.~Jones'',        % ~ produces an unbreakable interword space.
or within a word---especially when the word is a
symbol like
       \mbox{\emph{itemnum}}
that makes little sense when hyphenated across
lines.

Footnotes\footnote{This is an example of a footnote.}
pose no problem.

\LaTeX\ is good at typesetting mathematical formulas
like
       \( x-3y + z = 7 \)
or
       \( a_{1} > x^{2n} + y^{2n} > x' \)
or
       \( \ip{A}{B} = \sum_{i} a_{i} b_{i} \).
The spaces you type in a formula are
ignored.  Remember that a letter like
       $x$                   % $ ... $  and  \( ... \)  are equivalent
is a formula when it denotes a mathematical
symbol, and it should be typed as one.

\section{Displayed Text}

Text is displayed by indenting it from the left
margin.  Quotations are commonly displayed.  There
are short quotations
\begin{quote}
   This is a short quotation.  It consists of a
   single paragraph of text.  See how it is formatted.
\end{quote}
and longer ones.
\begin{quotation}
   This is a longer quotation.  It consists of two
   paragraphs of text, neither of which are
   particularly interesting.

   This is the second paragraph of the quotation.  It
   is just as dull as the first paragraph.
\end{quotation}
Another frequently-displayed structure is a list.
The following is an example of an \emph{itemized}
list.
\begin{itemize}
   \item This is the first item of an itemized list.
         Each item in the list is marked with a ``tick''.
         You don't have to worry about what kind of tick
         mark is used.

   \item This is the second item of the list.  It
         contains another list nested inside it.  The inner
         list is an \emph{enumerated} list.
         \begin{enumerate}
            \item This is the first item of an enumerated
                  list that is nested within the itemized list.

            \item This is the second item of the inner list.
                  \LaTeX\ allows you to nest lists deeper than
                  you really should.
         \end{enumerate}
         This is the rest of the second item of the outer
         list.  It is no more interesting than any other
         part of the item.
   \item This is the third item of the list.
\end{itemize}
You can even display poetry.
\begin{verse}
   There is an environment
    for verse \\             % The \\ command separates lines
   Whose features some poets % within a stanza.
   will curse.

                             % One or more blank lines separate stanzas.

   For instead of making\\
   Them do \emph{all} line breaking, \\
   It allows them to put too many words on a line when they'd rather be
   forced to be terse.
\end{verse}

Mathematical formulas may also be displayed.  A
displayed formula
is
one-line long; multiline
formulas require special formatting instructions.
   \[  \ip{\Gamma}{\psi'} = x'' + y^{2} + z_{i}^{n}\]
Don't start a paragraph with a displayed equation,
nor make one a paragraph by itself.

\end{document}               % End of document.

\section{Bench Huge Section 1521}
% synthetic bench section 1521
% This is a sample LaTeX input file.  (Version of 12 August 2004.)
%
% A '%' character causes TeX to ignore all remaining text on the line,
% and is used for comments like this one.

\documentclass{article}      % Specifies the document class

                             % The preamble begins here.
\title{An Example Document}  % Declares the document's title.
\author{Leslie Lamport}      % Declares the author's name.
\date{January 21, 1994}      % Deleting this command produces today's date.

\newcommand{\ip}[2]{(#1, #2)}
                             % Defines \ip{arg1}{arg2} to mean
                             % (arg1, arg2).

%\newcommand{\ip}[2]{\langle #1 | #2\rangle}
                             % This is an alternative definition of
                             % \ip that is commented out.

\begin{document}             % End of preamble and beginning of text.

\maketitle                   % Produces the title.

This is an example input file.  Comparing it with
the output it generates can show you how to
produce a simple document of your own.

\section{Ordinary Text}      % Produces section heading.  Lower-level
                             % sections are begun with similar
                             % \subsection and \subsubsection commands.

The ends  of words and sentences are marked
  by   spaces. It  doesn't matter how many
spaces    you type; one is as good as 100.  The
end of   a line counts as a space.

One   or more   blank lines denote the  end
of  a paragraph.

Since any number of consecutive spaces are treated
like a single one, the formatting of the input
file makes no difference to
      \LaTeX,                % The \LaTeX command generates the LaTeX logo.
but it makes a difference to you.  When you use
\LaTeX, making your input file as easy to read
as possible will be a great help as you write
your document and when you change it.  This sample
file shows how you can add comments to your own input
file.

Because printing is different from typewriting,
there are a number of things that you have to do
differently when preparing an input file than if
you were just typing the document directly.
Quotation marks like
       ``this''
have to be handled specially, as do quotes within
quotes:
       ``\,`this'            % \, separates the double and single quote.
        is what I just
        wrote, not  `that'\,''.

Dashes come in three sizes: an
       intra-word
dash, a medium dash for number ranges like
       1--2,
and a punctuation
       dash---like
this.

A sentence-ending space should be larger than the
space between words within a sentence.  You
sometimes have to type special commands in
conjunction with punctuation characters to get
this right, as in the following sentence.
       Gnats, gnus, etc.\ all  % `\ ' makes an inter-word space.
       begin with G\@.         % \@ marks end-of-sentence punctuation.
You should check the spaces after periods when
reading your output to make sure you haven't
forgotten any special cases.  Generating an
ellipsis
       \ldots\               % `\ ' is needed after `\ldots' because TeX
                             % ignores spaces after command names like \ldots
                             % made from \ + letters.
                             %
                             % Note how a `%' character causes TeX to ignore
                             % the end of the input line, so these blank lines
                             % do not start a new paragraph.
                             %
with the right spacing around the periods requires
a special command.

\LaTeX\ interprets some common characters as
commands, so you must type special commands to
generate them.  These characters include the
following:
       \$ \& \% \# \{ and \}.

In printing, text is usually emphasized with an
       \emph{italic}
type style.

\begin{em}
   A long segment of text can also be emphasized
   in this way.  Text within such a segment can be
   given \emph{additional} emphasis.
\end{em}

It is sometimes necessary to prevent \LaTeX\ from
breaking a line where it might otherwise do so.
This may be at a space, as between the ``Mr.''\ and
``Jones'' in
       ``Mr.~Jones'',        % ~ produces an unbreakable interword space.
or within a word---especially when the word is a
symbol like
       \mbox{\emph{itemnum}}
that makes little sense when hyphenated across
lines.

Footnotes\footnote{This is an example of a footnote.}
pose no problem.

\LaTeX\ is good at typesetting mathematical formulas
like
       \( x-3y + z = 7 \)
or
       \( a_{1} > x^{2n} + y^{2n} > x' \)
or
       \( \ip{A}{B} = \sum_{i} a_{i} b_{i} \).
The spaces you type in a formula are
ignored.  Remember that a letter like
       $x$                   % $ ... $  and  \( ... \)  are equivalent
is a formula when it denotes a mathematical
symbol, and it should be typed as one.

\section{Displayed Text}

Text is displayed by indenting it from the left
margin.  Quotations are commonly displayed.  There
are short quotations
\begin{quote}
   This is a short quotation.  It consists of a
   single paragraph of text.  See how it is formatted.
\end{quote}
and longer ones.
\begin{quotation}
   This is a longer quotation.  It consists of two
   paragraphs of text, neither of which are
   particularly interesting.

   This is the second paragraph of the quotation.  It
   is just as dull as the first paragraph.
\end{quotation}
Another frequently-displayed structure is a list.
The following is an example of an \emph{itemized}
list.
\begin{itemize}
   \item This is the first item of an itemized list.
         Each item in the list is marked with a ``tick''.
         You don't have to worry about what kind of tick
         mark is used.

   \item This is the second item of the list.  It
         contains another list nested inside it.  The inner
         list is an \emph{enumerated} list.
         \begin{enumerate}
            \item This is the first item of an enumerated
                  list that is nested within the itemized list.

            \item This is the second item of the inner list.
                  \LaTeX\ allows you to nest lists deeper than
                  you really should.
         \end{enumerate}
         This is the rest of the second item of the outer
         list.  It is no more interesting than any other
         part of the item.
   \item This is the third item of the list.
\end{itemize}
You can even display poetry.
\begin{verse}
   There is an environment
    for verse \\             % The \\ command separates lines
   Whose features some poets % within a stanza.
   will curse.

                             % One or more blank lines separate stanzas.

   For instead of making\\
   Them do \emph{all} line breaking, \\
   It allows them to put too many words on a line when they'd rather be
   forced to be terse.
\end{verse}

Mathematical formulas may also be displayed.  A
displayed formula
is
one-line long; multiline
formulas require special formatting instructions.
   \[  \ip{\Gamma}{\psi'} = x'' + y^{2} + z_{i}^{n}\]
Don't start a paragraph with a displayed equation,
nor make one a paragraph by itself.

\end{document}               % End of document.

\section{Bench Huge Section 1522}
% synthetic bench section 1522
% This is a sample LaTeX input file.  (Version of 12 August 2004.)
%
% A '%' character causes TeX to ignore all remaining text on the line,
% and is used for comments like this one.

\documentclass{article}      % Specifies the document class

                             % The preamble begins here.
\title{An Example Document}  % Declares the document's title.
\author{Leslie Lamport}      % Declares the author's name.
\date{January 21, 1994}      % Deleting this command produces today's date.

\newcommand{\ip}[2]{(#1, #2)}
                             % Defines \ip{arg1}{arg2} to mean
                             % (arg1, arg2).

%\newcommand{\ip}[2]{\langle #1 | #2\rangle}
                             % This is an alternative definition of
                             % \ip that is commented out.

\begin{document}             % End of preamble and beginning of text.

\maketitle                   % Produces the title.

This is an example input file.  Comparing it with
the output it generates can show you how to
produce a simple document of your own.

\section{Ordinary Text}      % Produces section heading.  Lower-level
                             % sections are begun with similar
                             % \subsection and \subsubsection commands.

The ends  of words and sentences are marked
  by   spaces. It  doesn't matter how many
spaces    you type; one is as good as 100.  The
end of   a line counts as a space.

One   or more   blank lines denote the  end
of  a paragraph.

Since any number of consecutive spaces are treated
like a single one, the formatting of the input
file makes no difference to
      \LaTeX,                % The \LaTeX command generates the LaTeX logo.
but it makes a difference to you.  When you use
\LaTeX, making your input file as easy to read
as possible will be a great help as you write
your document and when you change it.  This sample
file shows how you can add comments to your own input
file.

Because printing is different from typewriting,
there are a number of things that you have to do
differently when preparing an input file than if
you were just typing the document directly.
Quotation marks like
       ``this''
have to be handled specially, as do quotes within
quotes:
       ``\,`this'            % \, separates the double and single quote.
        is what I just
        wrote, not  `that'\,''.

Dashes come in three sizes: an
       intra-word
dash, a medium dash for number ranges like
       1--2,
and a punctuation
       dash---like
this.

A sentence-ending space should be larger than the
space between words within a sentence.  You
sometimes have to type special commands in
conjunction with punctuation characters to get
this right, as in the following sentence.
       Gnats, gnus, etc.\ all  % `\ ' makes an inter-word space.
       begin with G\@.         % \@ marks end-of-sentence punctuation.
You should check the spaces after periods when
reading your output to make sure you haven't
forgotten any special cases.  Generating an
ellipsis
       \ldots\               % `\ ' is needed after `\ldots' because TeX
                             % ignores spaces after command names like \ldots
                             % made from \ + letters.
                             %
                             % Note how a `%' character causes TeX to ignore
                             % the end of the input line, so these blank lines
                             % do not start a new paragraph.
                             %
with the right spacing around the periods requires
a special command.

\LaTeX\ interprets some common characters as
commands, so you must type special commands to
generate them.  These characters include the
following:
       \$ \& \% \# \{ and \}.

In printing, text is usually emphasized with an
       \emph{italic}
type style.

\begin{em}
   A long segment of text can also be emphasized
   in this way.  Text within such a segment can be
   given \emph{additional} emphasis.
\end{em}

It is sometimes necessary to prevent \LaTeX\ from
breaking a line where it might otherwise do so.
This may be at a space, as between the ``Mr.''\ and
``Jones'' in
       ``Mr.~Jones'',        % ~ produces an unbreakable interword space.
or within a word---especially when the word is a
symbol like
       \mbox{\emph{itemnum}}
that makes little sense when hyphenated across
lines.

Footnotes\footnote{This is an example of a footnote.}
pose no problem.

\LaTeX\ is good at typesetting mathematical formulas
like
       \( x-3y + z = 7 \)
or
       \( a_{1} > x^{2n} + y^{2n} > x' \)
or
       \( \ip{A}{B} = \sum_{i} a_{i} b_{i} \).
The spaces you type in a formula are
ignored.  Remember that a letter like
       $x$                   % $ ... $  and  \( ... \)  are equivalent
is a formula when it denotes a mathematical
symbol, and it should be typed as one.

\section{Displayed Text}

Text is displayed by indenting it from the left
margin.  Quotations are commonly displayed.  There
are short quotations
\begin{quote}
   This is a short quotation.  It consists of a
   single paragraph of text.  See how it is formatted.
\end{quote}
and longer ones.
\begin{quotation}
   This is a longer quotation.  It consists of two
   paragraphs of text, neither of which are
   particularly interesting.

   This is the second paragraph of the quotation.  It
   is just as dull as the first paragraph.
\end{quotation}
Another frequently-displayed structure is a list.
The following is an example of an \emph{itemized}
list.
\begin{itemize}
   \item This is the first item of an itemized list.
         Each item in the list is marked with a ``tick''.
         You don't have to worry about what kind of tick
         mark is used.

   \item This is the second item of the list.  It
         contains another list nested inside it.  The inner
         list is an \emph{enumerated} list.
         \begin{enumerate}
            \item This is the first item of an enumerated
                  list that is nested within the itemized list.

            \item This is the second item of the inner list.
                  \LaTeX\ allows you to nest lists deeper than
                  you really should.
         \end{enumerate}
         This is the rest of the second item of the outer
         list.  It is no more interesting than any other
         part of the item.
   \item This is the third item of the list.
\end{itemize}
You can even display poetry.
\begin{verse}
   There is an environment
    for verse \\             % The \\ command separates lines
   Whose features some poets % within a stanza.
   will curse.

                             % One or more blank lines separate stanzas.

   For instead of making\\
   Them do \emph{all} line breaking, \\
   It allows them to put too many words on a line when they'd rather be
   forced to be terse.
\end{verse}

Mathematical formulas may also be displayed.  A
displayed formula
is
one-line long; multiline
formulas require special formatting instructions.
   \[  \ip{\Gamma}{\psi'} = x'' + y^{2} + z_{i}^{n}\]
Don't start a paragraph with a displayed equation,
nor make one a paragraph by itself.

\end{document}               % End of document.

\section{Bench Huge Section 1523}
% synthetic bench section 1523
% This is a sample LaTeX input file.  (Version of 12 August 2004.)
%
% A '%' character causes TeX to ignore all remaining text on the line,
% and is used for comments like this one.

\documentclass{article}      % Specifies the document class

                             % The preamble begins here.
\title{An Example Document}  % Declares the document's title.
\author{Leslie Lamport}      % Declares the author's name.
\date{January 21, 1994}      % Deleting this command produces today's date.

\newcommand{\ip}[2]{(#1, #2)}
                             % Defines \ip{arg1}{arg2} to mean
                             % (arg1, arg2).

%\newcommand{\ip}[2]{\langle #1 | #2\rangle}
                             % This is an alternative definition of
                             % \ip that is commented out.

\begin{document}             % End of preamble and beginning of text.

\maketitle                   % Produces the title.

This is an example input file.  Comparing it with
the output it generates can show you how to
produce a simple document of your own.

\section{Ordinary Text}      % Produces section heading.  Lower-level
                             % sections are begun with similar
                             % \subsection and \subsubsection commands.

The ends  of words and sentences are marked
  by   spaces. It  doesn't matter how many
spaces    you type; one is as good as 100.  The
end of   a line counts as a space.

One   or more   blank lines denote the  end
of  a paragraph.

Since any number of consecutive spaces are treated
like a single one, the formatting of the input
file makes no difference to
      \LaTeX,                % The \LaTeX command generates the LaTeX logo.
but it makes a difference to you.  When you use
\LaTeX, making your input file as easy to read
as possible will be a great help as you write
your document and when you change it.  This sample
file shows how you can add comments to your own input
file.

Because printing is different from typewriting,
there are a number of things that you have to do
differently when preparing an input file than if
you were just typing the document directly.
Quotation marks like
       ``this''
have to be handled specially, as do quotes within
quotes:
       ``\,`this'            % \, separates the double and single quote.
        is what I just
        wrote, not  `that'\,''.

Dashes come in three sizes: an
       intra-word
dash, a medium dash for number ranges like
       1--2,
and a punctuation
       dash---like
this.

A sentence-ending space should be larger than the
space between words within a sentence.  You
sometimes have to type special commands in
conjunction with punctuation characters to get
this right, as in the following sentence.
       Gnats, gnus, etc.\ all  % `\ ' makes an inter-word space.
       begin with G\@.         % \@ marks end-of-sentence punctuation.
You should check the spaces after periods when
reading your output to make sure you haven't
forgotten any special cases.  Generating an
ellipsis
       \ldots\               % `\ ' is needed after `\ldots' because TeX
                             % ignores spaces after command names like \ldots
                             % made from \ + letters.
                             %
                             % Note how a `%' character causes TeX to ignore
                             % the end of the input line, so these blank lines
                             % do not start a new paragraph.
                             %
with the right spacing around the periods requires
a special command.

\LaTeX\ interprets some common characters as
commands, so you must type special commands to
generate them.  These characters include the
following:
       \$ \& \% \# \{ and \}.

In printing, text is usually emphasized with an
       \emph{italic}
type style.

\begin{em}
   A long segment of text can also be emphasized
   in this way.  Text within such a segment can be
   given \emph{additional} emphasis.
\end{em}

It is sometimes necessary to prevent \LaTeX\ from
breaking a line where it might otherwise do so.
This may be at a space, as between the ``Mr.''\ and
``Jones'' in
       ``Mr.~Jones'',        % ~ produces an unbreakable interword space.
or within a word---especially when the word is a
symbol like
       \mbox{\emph{itemnum}}
that makes little sense when hyphenated across
lines.

Footnotes\footnote{This is an example of a footnote.}
pose no problem.

\LaTeX\ is good at typesetting mathematical formulas
like
       \( x-3y + z = 7 \)
or
       \( a_{1} > x^{2n} + y^{2n} > x' \)
or
       \( \ip{A}{B} = \sum_{i} a_{i} b_{i} \).
The spaces you type in a formula are
ignored.  Remember that a letter like
       $x$                   % $ ... $  and  \( ... \)  are equivalent
is a formula when it denotes a mathematical
symbol, and it should be typed as one.

\section{Displayed Text}

Text is displayed by indenting it from the left
margin.  Quotations are commonly displayed.  There
are short quotations
\begin{quote}
   This is a short quotation.  It consists of a
   single paragraph of text.  See how it is formatted.
\end{quote}
and longer ones.
\begin{quotation}
   This is a longer quotation.  It consists of two
   paragraphs of text, neither of which are
   particularly interesting.

   This is the second paragraph of the quotation.  It
   is just as dull as the first paragraph.
\end{quotation}
Another frequently-displayed structure is a list.
The following is an example of an \emph{itemized}
list.
\begin{itemize}
   \item This is the first item of an itemized list.
         Each item in the list is marked with a ``tick''.
         You don't have to worry about what kind of tick
         mark is used.

   \item This is the second item of the list.  It
         contains another list nested inside it.  The inner
         list is an \emph{enumerated} list.
         \begin{enumerate}
            \item This is the first item of an enumerated
                  list that is nested within the itemized list.

            \item This is the second item of the inner list.
                  \LaTeX\ allows you to nest lists deeper than
                  you really should.
         \end{enumerate}
         This is the rest of the second item of the outer
         list.  It is no more interesting than any other
         part of the item.
   \item This is the third item of the list.
\end{itemize}
You can even display poetry.
\begin{verse}
   There is an environment
    for verse \\             % The \\ command separates lines
   Whose features some poets % within a stanza.
   will curse.

                             % One or more blank lines separate stanzas.

   For instead of making\\
   Them do \emph{all} line breaking, \\
   It allows them to put too many words on a line when they'd rather be
   forced to be terse.
\end{verse}

Mathematical formulas may also be displayed.  A
displayed formula
is
one-line long; multiline
formulas require special formatting instructions.
   \[  \ip{\Gamma}{\psi'} = x'' + y^{2} + z_{i}^{n}\]
Don't start a paragraph with a displayed equation,
nor make one a paragraph by itself.

\end{document}               % End of document.

\section{Bench Huge Section 1524}
% synthetic bench section 1524
% This is a sample LaTeX input file.  (Version of 12 August 2004.)
%
% A '%' character causes TeX to ignore all remaining text on the line,
% and is used for comments like this one.

\documentclass{article}      % Specifies the document class

                             % The preamble begins here.
\title{An Example Document}  % Declares the document's title.
\author{Leslie Lamport}      % Declares the author's name.
\date{January 21, 1994}      % Deleting this command produces today's date.

\newcommand{\ip}[2]{(#1, #2)}
                             % Defines \ip{arg1}{arg2} to mean
                             % (arg1, arg2).

%\newcommand{\ip}[2]{\langle #1 | #2\rangle}
                             % This is an alternative definition of
                             % \ip that is commented out.

\begin{document}             % End of preamble and beginning of text.

\maketitle                   % Produces the title.

This is an example input file.  Comparing it with
the output it generates can show you how to
produce a simple document of your own.

\section{Ordinary Text}      % Produces section heading.  Lower-level
                             % sections are begun with similar
                             % \subsection and \subsubsection commands.

The ends  of words and sentences are marked
  by   spaces. It  doesn't matter how many
spaces    you type; one is as good as 100.  The
end of   a line counts as a space.

One   or more   blank lines denote the  end
of  a paragraph.

Since any number of consecutive spaces are treated
like a single one, the formatting of the input
file makes no difference to
      \LaTeX,                % The \LaTeX command generates the LaTeX logo.
but it makes a difference to you.  When you use
\LaTeX, making your input file as easy to read
as possible will be a great help as you write
your document and when you change it.  This sample
file shows how you can add comments to your own input
file.

Because printing is different from typewriting,
there are a number of things that you have to do
differently when preparing an input file than if
you were just typing the document directly.
Quotation marks like
       ``this''
have to be handled specially, as do quotes within
quotes:
       ``\,`this'            % \, separates the double and single quote.
        is what I just
        wrote, not  `that'\,''.

Dashes come in three sizes: an
       intra-word
dash, a medium dash for number ranges like
       1--2,
and a punctuation
       dash---like
this.

A sentence-ending space should be larger than the
space between words within a sentence.  You
sometimes have to type special commands in
conjunction with punctuation characters to get
this right, as in the following sentence.
       Gnats, gnus, etc.\ all  % `\ ' makes an inter-word space.
       begin with G\@.         % \@ marks end-of-sentence punctuation.
You should check the spaces after periods when
reading your output to make sure you haven't
forgotten any special cases.  Generating an
ellipsis
       \ldots\               % `\ ' is needed after `\ldots' because TeX
                             % ignores spaces after command names like \ldots
                             % made from \ + letters.
                             %
                             % Note how a `%' character causes TeX to ignore
                             % the end of the input line, so these blank lines
                             % do not start a new paragraph.
                             %
with the right spacing around the periods requires
a special command.

\LaTeX\ interprets some common characters as
commands, so you must type special commands to
generate them.  These characters include the
following:
       \$ \& \% \# \{ and \}.

In printing, text is usually emphasized with an
       \emph{italic}
type style.

\begin{em}
   A long segment of text can also be emphasized
   in this way.  Text within such a segment can be
   given \emph{additional} emphasis.
\end{em}

It is sometimes necessary to prevent \LaTeX\ from
breaking a line where it might otherwise do so.
This may be at a space, as between the ``Mr.''\ and
``Jones'' in
       ``Mr.~Jones'',        % ~ produces an unbreakable interword space.
or within a word---especially when the word is a
symbol like
       \mbox{\emph{itemnum}}
that makes little sense when hyphenated across
lines.

Footnotes\footnote{This is an example of a footnote.}
pose no problem.

\LaTeX\ is good at typesetting mathematical formulas
like
       \( x-3y + z = 7 \)
or
       \( a_{1} > x^{2n} + y^{2n} > x' \)
or
       \( \ip{A}{B} = \sum_{i} a_{i} b_{i} \).
The spaces you type in a formula are
ignored.  Remember that a letter like
       $x$                   % $ ... $  and  \( ... \)  are equivalent
is a formula when it denotes a mathematical
symbol, and it should be typed as one.

\section{Displayed Text}

Text is displayed by indenting it from the left
margin.  Quotations are commonly displayed.  There
are short quotations
\begin{quote}
   This is a short quotation.  It consists of a
   single paragraph of text.  See how it is formatted.
\end{quote}
and longer ones.
\begin{quotation}
   This is a longer quotation.  It consists of two
   paragraphs of text, neither of which are
   particularly interesting.

   This is the second paragraph of the quotation.  It
   is just as dull as the first paragraph.
\end{quotation}
Another frequently-displayed structure is a list.
The following is an example of an \emph{itemized}
list.
\begin{itemize}
   \item This is the first item of an itemized list.
         Each item in the list is marked with a ``tick''.
         You don't have to worry about what kind of tick
         mark is used.

   \item This is the second item of the list.  It
         contains another list nested inside it.  The inner
         list is an \emph{enumerated} list.
         \begin{enumerate}
            \item This is the first item of an enumerated
                  list that is nested within the itemized list.

            \item This is the second item of the inner list.
                  \LaTeX\ allows you to nest lists deeper than
                  you really should.
         \end{enumerate}
         This is the rest of the second item of the outer
         list.  It is no more interesting than any other
         part of the item.
   \item This is the third item of the list.
\end{itemize}
You can even display poetry.
\begin{verse}
   There is an environment
    for verse \\             % The \\ command separates lines
   Whose features some poets % within a stanza.
   will curse.

                             % One or more blank lines separate stanzas.

   For instead of making\\
   Them do \emph{all} line breaking, \\
   It allows them to put too many words on a line when they'd rather be
   forced to be terse.
\end{verse}

Mathematical formulas may also be displayed.  A
displayed formula
is
one-line long; multiline
formulas require special formatting instructions.
   \[  \ip{\Gamma}{\psi'} = x'' + y^{2} + z_{i}^{n}\]
Don't start a paragraph with a displayed equation,
nor make one a paragraph by itself.

\end{document}               % End of document.

\section{Bench Huge Section 1525}
% synthetic bench section 1525
% This is a sample LaTeX input file.  (Version of 12 August 2004.)
%
% A '%' character causes TeX to ignore all remaining text on the line,
% and is used for comments like this one.

\documentclass{article}      % Specifies the document class

                             % The preamble begins here.
\title{An Example Document}  % Declares the document's title.
\author{Leslie Lamport}      % Declares the author's name.
\date{January 21, 1994}      % Deleting this command produces today's date.

\newcommand{\ip}[2]{(#1, #2)}
                             % Defines \ip{arg1}{arg2} to mean
                             % (arg1, arg2).

%\newcommand{\ip}[2]{\langle #1 | #2\rangle}
                             % This is an alternative definition of
                             % \ip that is commented out.

\begin{document}             % End of preamble and beginning of text.

\maketitle                   % Produces the title.

This is an example input file.  Comparing it with
the output it generates can show you how to
produce a simple document of your own.

\section{Ordinary Text}      % Produces section heading.  Lower-level
                             % sections are begun with similar
                             % \subsection and \subsubsection commands.

The ends  of words and sentences are marked
  by   spaces. It  doesn't matter how many
spaces    you type; one is as good as 100.  The
end of   a line counts as a space.

One   or more   blank lines denote the  end
of  a paragraph.

Since any number of consecutive spaces are treated
like a single one, the formatting of the input
file makes no difference to
      \LaTeX,                % The \LaTeX command generates the LaTeX logo.
but it makes a difference to you.  When you use
\LaTeX, making your input file as easy to read
as possible will be a great help as you write
your document and when you change it.  This sample
file shows how you can add comments to your own input
file.

Because printing is different from typewriting,
there are a number of things that you have to do
differently when preparing an input file than if
you were just typing the document directly.
Quotation marks like
       ``this''
have to be handled specially, as do quotes within
quotes:
       ``\,`this'            % \, separates the double and single quote.
        is what I just
        wrote, not  `that'\,''.

Dashes come in three sizes: an
       intra-word
dash, a medium dash for number ranges like
       1--2,
and a punctuation
       dash---like
this.

A sentence-ending space should be larger than the
space between words within a sentence.  You
sometimes have to type special commands in
conjunction with punctuation characters to get
this right, as in the following sentence.
       Gnats, gnus, etc.\ all  % `\ ' makes an inter-word space.
       begin with G\@.         % \@ marks end-of-sentence punctuation.
You should check the spaces after periods when
reading your output to make sure you haven't
forgotten any special cases.  Generating an
ellipsis
       \ldots\               % `\ ' is needed after `\ldots' because TeX
                             % ignores spaces after command names like \ldots
                             % made from \ + letters.
                             %
                             % Note how a `%' character causes TeX to ignore
                             % the end of the input line, so these blank lines
                             % do not start a new paragraph.
                             %
with the right spacing around the periods requires
a special command.

\LaTeX\ interprets some common characters as
commands, so you must type special commands to
generate them.  These characters include the
following:
       \$ \& \% \# \{ and \}.

In printing, text is usually emphasized with an
       \emph{italic}
type style.

\begin{em}
   A long segment of text can also be emphasized
   in this way.  Text within such a segment can be
   given \emph{additional} emphasis.
\end{em}

It is sometimes necessary to prevent \LaTeX\ from
breaking a line where it might otherwise do so.
This may be at a space, as between the ``Mr.''\ and
``Jones'' in
       ``Mr.~Jones'',        % ~ produces an unbreakable interword space.
or within a word---especially when the word is a
symbol like
       \mbox{\emph{itemnum}}
that makes little sense when hyphenated across
lines.

Footnotes\footnote{This is an example of a footnote.}
pose no problem.

\LaTeX\ is good at typesetting mathematical formulas
like
       \( x-3y + z = 7 \)
or
       \( a_{1} > x^{2n} + y^{2n} > x' \)
or
       \( \ip{A}{B} = \sum_{i} a_{i} b_{i} \).
The spaces you type in a formula are
ignored.  Remember that a letter like
       $x$                   % $ ... $  and  \( ... \)  are equivalent
is a formula when it denotes a mathematical
symbol, and it should be typed as one.

\section{Displayed Text}

Text is displayed by indenting it from the left
margin.  Quotations are commonly displayed.  There
are short quotations
\begin{quote}
   This is a short quotation.  It consists of a
   single paragraph of text.  See how it is formatted.
\end{quote}
and longer ones.
\begin{quotation}
   This is a longer quotation.  It consists of two
   paragraphs of text, neither of which are
   particularly interesting.

   This is the second paragraph of the quotation.  It
   is just as dull as the first paragraph.
\end{quotation}
Another frequently-displayed structure is a list.
The following is an example of an \emph{itemized}
list.
\begin{itemize}
   \item This is the first item of an itemized list.
         Each item in the list is marked with a ``tick''.
         You don't have to worry about what kind of tick
         mark is used.

   \item This is the second item of the list.  It
         contains another list nested inside it.  The inner
         list is an \emph{enumerated} list.
         \begin{enumerate}
            \item This is the first item of an enumerated
                  list that is nested within the itemized list.

            \item This is the second item of the inner list.
                  \LaTeX\ allows you to nest lists deeper than
                  you really should.
         \end{enumerate}
         This is the rest of the second item of the outer
         list.  It is no more interesting than any other
         part of the item.
   \item This is the third item of the list.
\end{itemize}
You can even display poetry.
\begin{verse}
   There is an environment
    for verse \\             % The \\ command separates lines
   Whose features some poets % within a stanza.
   will curse.

                             % One or more blank lines separate stanzas.

   For instead of making\\
   Them do \emph{all} line breaking, \\
   It allows them to put too many words on a line when they'd rather be
   forced to be terse.
\end{verse}

Mathematical formulas may also be displayed.  A
displayed formula
is
one-line long; multiline
formulas require special formatting instructions.
   \[  \ip{\Gamma}{\psi'} = x'' + y^{2} + z_{i}^{n}\]
Don't start a paragraph with a displayed equation,
nor make one a paragraph by itself.

\end{document}               % End of document.

\section{Bench Huge Section 1526}
% synthetic bench section 1526
% This is a sample LaTeX input file.  (Version of 12 August 2004.)
%
% A '%' character causes TeX to ignore all remaining text on the line,
% and is used for comments like this one.

\documentclass{article}      % Specifies the document class

                             % The preamble begins here.
\title{An Example Document}  % Declares the document's title.
\author{Leslie Lamport}      % Declares the author's name.
\date{January 21, 1994}      % Deleting this command produces today's date.

\newcommand{\ip}[2]{(#1, #2)}
                             % Defines \ip{arg1}{arg2} to mean
                             % (arg1, arg2).

%\newcommand{\ip}[2]{\langle #1 | #2\rangle}
                             % This is an alternative definition of
                             % \ip that is commented out.

\begin{document}             % End of preamble and beginning of text.

\maketitle                   % Produces the title.

This is an example input file.  Comparing it with
the output it generates can show you how to
produce a simple document of your own.

\section{Ordinary Text}      % Produces section heading.  Lower-level
                             % sections are begun with similar
                             % \subsection and \subsubsection commands.

The ends  of words and sentences are marked
  by   spaces. It  doesn't matter how many
spaces    you type; one is as good as 100.  The
end of   a line counts as a space.

One   or more   blank lines denote the  end
of  a paragraph.

Since any number of consecutive spaces are treated
like a single one, the formatting of the input
file makes no difference to
      \LaTeX,                % The \LaTeX command generates the LaTeX logo.
but it makes a difference to you.  When you use
\LaTeX, making your input file as easy to read
as possible will be a great help as you write
your document and when you change it.  This sample
file shows how you can add comments to your own input
file.

Because printing is different from typewriting,
there are a number of things that you have to do
differently when preparing an input file than if
you were just typing the document directly.
Quotation marks like
       ``this''
have to be handled specially, as do quotes within
quotes:
       ``\,`this'            % \, separates the double and single quote.
        is what I just
        wrote, not  `that'\,''.

Dashes come in three sizes: an
       intra-word
dash, a medium dash for number ranges like
       1--2,
and a punctuation
       dash---like
this.

A sentence-ending space should be larger than the
space between words within a sentence.  You
sometimes have to type special commands in
conjunction with punctuation characters to get
this right, as in the following sentence.
       Gnats, gnus, etc.\ all  % `\ ' makes an inter-word space.
       begin with G\@.         % \@ marks end-of-sentence punctuation.
You should check the spaces after periods when
reading your output to make sure you haven't
forgotten any special cases.  Generating an
ellipsis
       \ldots\               % `\ ' is needed after `\ldots' because TeX
                             % ignores spaces after command names like \ldots
                             % made from \ + letters.
                             %
                             % Note how a `%' character causes TeX to ignore
                             % the end of the input line, so these blank lines
                             % do not start a new paragraph.
                             %
with the right spacing around the periods requires
a special command.

\LaTeX\ interprets some common characters as
commands, so you must type special commands to
generate them.  These characters include the
following:
       \$ \& \% \# \{ and \}.

In printing, text is usually emphasized with an
       \emph{italic}
type style.

\begin{em}
   A long segment of text can also be emphasized
   in this way.  Text within such a segment can be
   given \emph{additional} emphasis.
\end{em}

It is sometimes necessary to prevent \LaTeX\ from
breaking a line where it might otherwise do so.
This may be at a space, as between the ``Mr.''\ and
``Jones'' in
       ``Mr.~Jones'',        % ~ produces an unbreakable interword space.
or within a word---especially when the word is a
symbol like
       \mbox{\emph{itemnum}}
that makes little sense when hyphenated across
lines.

Footnotes\footnote{This is an example of a footnote.}
pose no problem.

\LaTeX\ is good at typesetting mathematical formulas
like
       \( x-3y + z = 7 \)
or
       \( a_{1} > x^{2n} + y^{2n} > x' \)
or
       \( \ip{A}{B} = \sum_{i} a_{i} b_{i} \).
The spaces you type in a formula are
ignored.  Remember that a letter like
       $x$                   % $ ... $  and  \( ... \)  are equivalent
is a formula when it denotes a mathematical
symbol, and it should be typed as one.

\section{Displayed Text}

Text is displayed by indenting it from the left
margin.  Quotations are commonly displayed.  There
are short quotations
\begin{quote}
   This is a short quotation.  It consists of a
   single paragraph of text.  See how it is formatted.
\end{quote}
and longer ones.
\begin{quotation}
   This is a longer quotation.  It consists of two
   paragraphs of text, neither of which are
   particularly interesting.

   This is the second paragraph of the quotation.  It
   is just as dull as the first paragraph.
\end{quotation}
Another frequently-displayed structure is a list.
The following is an example of an \emph{itemized}
list.
\begin{itemize}
   \item This is the first item of an itemized list.
         Each item in the list is marked with a ``tick''.
         You don't have to worry about what kind of tick
         mark is used.

   \item This is the second item of the list.  It
         contains another list nested inside it.  The inner
         list is an \emph{enumerated} list.
         \begin{enumerate}
            \item This is the first item of an enumerated
                  list that is nested within the itemized list.

            \item This is the second item of the inner list.
                  \LaTeX\ allows you to nest lists deeper than
                  you really should.
         \end{enumerate}
         This is the rest of the second item of the outer
         list.  It is no more interesting than any other
         part of the item.
   \item This is the third item of the list.
\end{itemize}
You can even display poetry.
\begin{verse}
   There is an environment
    for verse \\             % The \\ command separates lines
   Whose features some poets % within a stanza.
   will curse.

                             % One or more blank lines separate stanzas.

   For instead of making\\
   Them do \emph{all} line breaking, \\
   It allows them to put too many words on a line when they'd rather be
   forced to be terse.
\end{verse}

Mathematical formulas may also be displayed.  A
displayed formula
is
one-line long; multiline
formulas require special formatting instructions.
   \[  \ip{\Gamma}{\psi'} = x'' + y^{2} + z_{i}^{n}\]
Don't start a paragraph with a displayed equation,
nor make one a paragraph by itself.

\end{document}               % End of document.

\section{Bench Huge Section 1527}
% synthetic bench section 1527
% This is a sample LaTeX input file.  (Version of 12 August 2004.)
%
% A '%' character causes TeX to ignore all remaining text on the line,
% and is used for comments like this one.

\documentclass{article}      % Specifies the document class

                             % The preamble begins here.
\title{An Example Document}  % Declares the document's title.
\author{Leslie Lamport}      % Declares the author's name.
\date{January 21, 1994}      % Deleting this command produces today's date.

\newcommand{\ip}[2]{(#1, #2)}
                             % Defines \ip{arg1}{arg2} to mean
                             % (arg1, arg2).

%\newcommand{\ip}[2]{\langle #1 | #2\rangle}
                             % This is an alternative definition of
                             % \ip that is commented out.

\begin{document}             % End of preamble and beginning of text.

\maketitle                   % Produces the title.

This is an example input file.  Comparing it with
the output it generates can show you how to
produce a simple document of your own.

\section{Ordinary Text}      % Produces section heading.  Lower-level
                             % sections are begun with similar
                             % \subsection and \subsubsection commands.

The ends  of words and sentences are marked
  by   spaces. It  doesn't matter how many
spaces    you type; one is as good as 100.  The
end of   a line counts as a space.

One   or more   blank lines denote the  end
of  a paragraph.

Since any number of consecutive spaces are treated
like a single one, the formatting of the input
file makes no difference to
      \LaTeX,                % The \LaTeX command generates the LaTeX logo.
but it makes a difference to you.  When you use
\LaTeX, making your input file as easy to read
as possible will be a great help as you write
your document and when you change it.  This sample
file shows how you can add comments to your own input
file.

Because printing is different from typewriting,
there are a number of things that you have to do
differently when preparing an input file than if
you were just typing the document directly.
Quotation marks like
       ``this''
have to be handled specially, as do quotes within
quotes:
       ``\,`this'            % \, separates the double and single quote.
        is what I just
        wrote, not  `that'\,''.

Dashes come in three sizes: an
       intra-word
dash, a medium dash for number ranges like
       1--2,
and a punctuation
       dash---like
this.

A sentence-ending space should be larger than the
space between words within a sentence.  You
sometimes have to type special commands in
conjunction with punctuation characters to get
this right, as in the following sentence.
       Gnats, gnus, etc.\ all  % `\ ' makes an inter-word space.
       begin with G\@.         % \@ marks end-of-sentence punctuation.
You should check the spaces after periods when
reading your output to make sure you haven't
forgotten any special cases.  Generating an
ellipsis
       \ldots\               % `\ ' is needed after `\ldots' because TeX
                             % ignores spaces after command names like \ldots
                             % made from \ + letters.
                             %
                             % Note how a `%' character causes TeX to ignore
                             % the end of the input line, so these blank lines
                             % do not start a new paragraph.
                             %
with the right spacing around the periods requires
a special command.

\LaTeX\ interprets some common characters as
commands, so you must type special commands to
generate them.  These characters include the
following:
       \$ \& \% \# \{ and \}.

In printing, text is usually emphasized with an
       \emph{italic}
type style.

\begin{em}
   A long segment of text can also be emphasized
   in this way.  Text within such a segment can be
   given \emph{additional} emphasis.
\end{em}

It is sometimes necessary to prevent \LaTeX\ from
breaking a line where it might otherwise do so.
This may be at a space, as between the ``Mr.''\ and
``Jones'' in
       ``Mr.~Jones'',        % ~ produces an unbreakable interword space.
or within a word---especially when the word is a
symbol like
       \mbox{\emph{itemnum}}
that makes little sense when hyphenated across
lines.

Footnotes\footnote{This is an example of a footnote.}
pose no problem.

\LaTeX\ is good at typesetting mathematical formulas
like
       \( x-3y + z = 7 \)
or
       \( a_{1} > x^{2n} + y^{2n} > x' \)
or
       \( \ip{A}{B} = \sum_{i} a_{i} b_{i} \).
The spaces you type in a formula are
ignored.  Remember that a letter like
       $x$                   % $ ... $  and  \( ... \)  are equivalent
is a formula when it denotes a mathematical
symbol, and it should be typed as one.

\section{Displayed Text}

Text is displayed by indenting it from the left
margin.  Quotations are commonly displayed.  There
are short quotations
\begin{quote}
   This is a short quotation.  It consists of a
   single paragraph of text.  See how it is formatted.
\end{quote}
and longer ones.
\begin{quotation}
   This is a longer quotation.  It consists of two
   paragraphs of text, neither of which are
   particularly interesting.

   This is the second paragraph of the quotation.  It
   is just as dull as the first paragraph.
\end{quotation}
Another frequently-displayed structure is a list.
The following is an example of an \emph{itemized}
list.
\begin{itemize}
   \item This is the first item of an itemized list.
         Each item in the list is marked with a ``tick''.
         You don't have to worry about what kind of tick
         mark is used.

   \item This is the second item of the list.  It
         contains another list nested inside it.  The inner
         list is an \emph{enumerated} list.
         \begin{enumerate}
            \item This is the first item of an enumerated
                  list that is nested within the itemized list.

            \item This is the second item of the inner list.
                  \LaTeX\ allows you to nest lists deeper than
                  you really should.
         \end{enumerate}
         This is the rest of the second item of the outer
         list.  It is no more interesting than any other
         part of the item.
   \item This is the third item of the list.
\end{itemize}
You can even display poetry.
\begin{verse}
   There is an environment
    for verse \\             % The \\ command separates lines
   Whose features some poets % within a stanza.
   will curse.

                             % One or more blank lines separate stanzas.

   For instead of making\\
   Them do \emph{all} line breaking, \\
   It allows them to put too many words on a line when they'd rather be
   forced to be terse.
\end{verse}

Mathematical formulas may also be displayed.  A
displayed formula
is
one-line long; multiline
formulas require special formatting instructions.
   \[  \ip{\Gamma}{\psi'} = x'' + y^{2} + z_{i}^{n}\]
Don't start a paragraph with a displayed equation,
nor make one a paragraph by itself.

\end{document}               % End of document.

\section{Bench Huge Section 1528}
% synthetic bench section 1528
% This is a sample LaTeX input file.  (Version of 12 August 2004.)
%
% A '%' character causes TeX to ignore all remaining text on the line,
% and is used for comments like this one.

\documentclass{article}      % Specifies the document class

                             % The preamble begins here.
\title{An Example Document}  % Declares the document's title.
\author{Leslie Lamport}      % Declares the author's name.
\date{January 21, 1994}      % Deleting this command produces today's date.

\newcommand{\ip}[2]{(#1, #2)}
                             % Defines \ip{arg1}{arg2} to mean
                             % (arg1, arg2).

%\newcommand{\ip}[2]{\langle #1 | #2\rangle}
                             % This is an alternative definition of
                             % \ip that is commented out.

\begin{document}             % End of preamble and beginning of text.

\maketitle                   % Produces the title.

This is an example input file.  Comparing it with
the output it generates can show you how to
produce a simple document of your own.

\section{Ordinary Text}      % Produces section heading.  Lower-level
                             % sections are begun with similar
                             % \subsection and \subsubsection commands.

The ends  of words and sentences are marked
  by   spaces. It  doesn't matter how many
spaces    you type; one is as good as 100.  The
end of   a line counts as a space.

One   or more   blank lines denote the  end
of  a paragraph.

Since any number of consecutive spaces are treated
like a single one, the formatting of the input
file makes no difference to
      \LaTeX,                % The \LaTeX command generates the LaTeX logo.
but it makes a difference to you.  When you use
\LaTeX, making your input file as easy to read
as possible will be a great help as you write
your document and when you change it.  This sample
file shows how you can add comments to your own input
file.

Because printing is different from typewriting,
there are a number of things that you have to do
differently when preparing an input file than if
you were just typing the document directly.
Quotation marks like
       ``this''
have to be handled specially, as do quotes within
quotes:
       ``\,`this'            % \, separates the double and single quote.
        is what I just
        wrote, not  `that'\,''.

Dashes come in three sizes: an
       intra-word
dash, a medium dash for number ranges like
       1--2,
and a punctuation
       dash---like
this.

A sentence-ending space should be larger than the
space between words within a sentence.  You
sometimes have to type special commands in
conjunction with punctuation characters to get
this right, as in the following sentence.
       Gnats, gnus, etc.\ all  % `\ ' makes an inter-word space.
       begin with G\@.         % \@ marks end-of-sentence punctuation.
You should check the spaces after periods when
reading your output to make sure you haven't
forgotten any special cases.  Generating an
ellipsis
       \ldots\               % `\ ' is needed after `\ldots' because TeX
                             % ignores spaces after command names like \ldots
                             % made from \ + letters.
                             %
                             % Note how a `%' character causes TeX to ignore
                             % the end of the input line, so these blank lines
                             % do not start a new paragraph.
                             %
with the right spacing around the periods requires
a special command.

\LaTeX\ interprets some common characters as
commands, so you must type special commands to
generate them.  These characters include the
following:
       \$ \& \% \# \{ and \}.

In printing, text is usually emphasized with an
       \emph{italic}
type style.

\begin{em}
   A long segment of text can also be emphasized
   in this way.  Text within such a segment can be
   given \emph{additional} emphasis.
\end{em}

It is sometimes necessary to prevent \LaTeX\ from
breaking a line where it might otherwise do so.
This may be at a space, as between the ``Mr.''\ and
``Jones'' in
       ``Mr.~Jones'',        % ~ produces an unbreakable interword space.
or within a word---especially when the word is a
symbol like
       \mbox{\emph{itemnum}}
that makes little sense when hyphenated across
lines.

Footnotes\footnote{This is an example of a footnote.}
pose no problem.

\LaTeX\ is good at typesetting mathematical formulas
like
       \( x-3y + z = 7 \)
or
       \( a_{1} > x^{2n} + y^{2n} > x' \)
or
       \( \ip{A}{B} = \sum_{i} a_{i} b_{i} \).
The spaces you type in a formula are
ignored.  Remember that a letter like
       $x$                   % $ ... $  and  \( ... \)  are equivalent
is a formula when it denotes a mathematical
symbol, and it should be typed as one.

\section{Displayed Text}

Text is displayed by indenting it from the left
margin.  Quotations are commonly displayed.  There
are short quotations
\begin{quote}
   This is a short quotation.  It consists of a
   single paragraph of text.  See how it is formatted.
\end{quote}
and longer ones.
\begin{quotation}
   This is a longer quotation.  It consists of two
   paragraphs of text, neither of which are
   particularly interesting.

   This is the second paragraph of the quotation.  It
   is just as dull as the first paragraph.
\end{quotation}
Another frequently-displayed structure is a list.
The following is an example of an \emph{itemized}
list.
\begin{itemize}
   \item This is the first item of an itemized list.
         Each item in the list is marked with a ``tick''.
         You don't have to worry about what kind of tick
         mark is used.

   \item This is the second item of the list.  It
         contains another list nested inside it.  The inner
         list is an \emph{enumerated} list.
         \begin{enumerate}
            \item This is the first item of an enumerated
                  list that is nested within the itemized list.

            \item This is the second item of the inner list.
                  \LaTeX\ allows you to nest lists deeper than
                  you really should.
         \end{enumerate}
         This is the rest of the second item of the outer
         list.  It is no more interesting than any other
         part of the item.
   \item This is the third item of the list.
\end{itemize}
You can even display poetry.
\begin{verse}
   There is an environment
    for verse \\             % The \\ command separates lines
   Whose features some poets % within a stanza.
   will curse.

                             % One or more blank lines separate stanzas.

   For instead of making\\
   Them do \emph{all} line breaking, \\
   It allows them to put too many words on a line when they'd rather be
   forced to be terse.
\end{verse}

Mathematical formulas may also be displayed.  A
displayed formula
is
one-line long; multiline
formulas require special formatting instructions.
   \[  \ip{\Gamma}{\psi'} = x'' + y^{2} + z_{i}^{n}\]
Don't start a paragraph with a displayed equation,
nor make one a paragraph by itself.

\end{document}               % End of document.

\section{Bench Huge Section 1529}
% synthetic bench section 1529
% This is a sample LaTeX input file.  (Version of 12 August 2004.)
%
% A '%' character causes TeX to ignore all remaining text on the line,
% and is used for comments like this one.

\documentclass{article}      % Specifies the document class

                             % The preamble begins here.
\title{An Example Document}  % Declares the document's title.
\author{Leslie Lamport}      % Declares the author's name.
\date{January 21, 1994}      % Deleting this command produces today's date.

\newcommand{\ip}[2]{(#1, #2)}
                             % Defines \ip{arg1}{arg2} to mean
                             % (arg1, arg2).

%\newcommand{\ip}[2]{\langle #1 | #2\rangle}
                             % This is an alternative definition of
                             % \ip that is commented out.

\begin{document}             % End of preamble and beginning of text.

\maketitle                   % Produces the title.

This is an example input file.  Comparing it with
the output it generates can show you how to
produce a simple document of your own.

\section{Ordinary Text}      % Produces section heading.  Lower-level
                             % sections are begun with similar
                             % \subsection and \subsubsection commands.

The ends  of words and sentences are marked
  by   spaces. It  doesn't matter how many
spaces    you type; one is as good as 100.  The
end of   a line counts as a space.

One   or more   blank lines denote the  end
of  a paragraph.

Since any number of consecutive spaces are treated
like a single one, the formatting of the input
file makes no difference to
      \LaTeX,                % The \LaTeX command generates the LaTeX logo.
but it makes a difference to you.  When you use
\LaTeX, making your input file as easy to read
as possible will be a great help as you write
your document and when you change it.  This sample
file shows how you can add comments to your own input
file.

Because printing is different from typewriting,
there are a number of things that you have to do
differently when preparing an input file than if
you were just typing the document directly.
Quotation marks like
       ``this''
have to be handled specially, as do quotes within
quotes:
       ``\,`this'            % \, separates the double and single quote.
        is what I just
        wrote, not  `that'\,''.

Dashes come in three sizes: an
       intra-word
dash, a medium dash for number ranges like
       1--2,
and a punctuation
       dash---like
this.

A sentence-ending space should be larger than the
space between words within a sentence.  You
sometimes have to type special commands in
conjunction with punctuation characters to get
this right, as in the following sentence.
       Gnats, gnus, etc.\ all  % `\ ' makes an inter-word space.
       begin with G\@.         % \@ marks end-of-sentence punctuation.
You should check the spaces after periods when
reading your output to make sure you haven't
forgotten any special cases.  Generating an
ellipsis
       \ldots\               % `\ ' is needed after `\ldots' because TeX
                             % ignores spaces after command names like \ldots
                             % made from \ + letters.
                             %
                             % Note how a `%' character causes TeX to ignore
                             % the end of the input line, so these blank lines
                             % do not start a new paragraph.
                             %
with the right spacing around the periods requires
a special command.

\LaTeX\ interprets some common characters as
commands, so you must type special commands to
generate them.  These characters include the
following:
       \$ \& \% \# \{ and \}.

In printing, text is usually emphasized with an
       \emph{italic}
type style.

\begin{em}
   A long segment of text can also be emphasized
   in this way.  Text within such a segment can be
   given \emph{additional} emphasis.
\end{em}

It is sometimes necessary to prevent \LaTeX\ from
breaking a line where it might otherwise do so.
This may be at a space, as between the ``Mr.''\ and
``Jones'' in
       ``Mr.~Jones'',        % ~ produces an unbreakable interword space.
or within a word---especially when the word is a
symbol like
       \mbox{\emph{itemnum}}
that makes little sense when hyphenated across
lines.

Footnotes\footnote{This is an example of a footnote.}
pose no problem.

\LaTeX\ is good at typesetting mathematical formulas
like
       \( x-3y + z = 7 \)
or
       \( a_{1} > x^{2n} + y^{2n} > x' \)
or
       \( \ip{A}{B} = \sum_{i} a_{i} b_{i} \).
The spaces you type in a formula are
ignored.  Remember that a letter like
       $x$                   % $ ... $  and  \( ... \)  are equivalent
is a formula when it denotes a mathematical
symbol, and it should be typed as one.

\section{Displayed Text}

Text is displayed by indenting it from the left
margin.  Quotations are commonly displayed.  There
are short quotations
\begin{quote}
   This is a short quotation.  It consists of a
   single paragraph of text.  See how it is formatted.
\end{quote}
and longer ones.
\begin{quotation}
   This is a longer quotation.  It consists of two
   paragraphs of text, neither of which are
   particularly interesting.

   This is the second paragraph of the quotation.  It
   is just as dull as the first paragraph.
\end{quotation}
Another frequently-displayed structure is a list.
The following is an example of an \emph{itemized}
list.
\begin{itemize}
   \item This is the first item of an itemized list.
         Each item in the list is marked with a ``tick''.
         You don't have to worry about what kind of tick
         mark is used.

   \item This is the second item of the list.  It
         contains another list nested inside it.  The inner
         list is an \emph{enumerated} list.
         \begin{enumerate}
            \item This is the first item of an enumerated
                  list that is nested within the itemized list.

            \item This is the second item of the inner list.
                  \LaTeX\ allows you to nest lists deeper than
                  you really should.
         \end{enumerate}
         This is the rest of the second item of the outer
         list.  It is no more interesting than any other
         part of the item.
   \item This is the third item of the list.
\end{itemize}
You can even display poetry.
\begin{verse}
   There is an environment
    for verse \\             % The \\ command separates lines
   Whose features some poets % within a stanza.
   will curse.

                             % One or more blank lines separate stanzas.

   For instead of making\\
   Them do \emph{all} line breaking, \\
   It allows them to put too many words on a line when they'd rather be
   forced to be terse.
\end{verse}

Mathematical formulas may also be displayed.  A
displayed formula
is
one-line long; multiline
formulas require special formatting instructions.
   \[  \ip{\Gamma}{\psi'} = x'' + y^{2} + z_{i}^{n}\]
Don't start a paragraph with a displayed equation,
nor make one a paragraph by itself.

\end{document}               % End of document.

\section{Bench Huge Section 1530}
% synthetic bench section 1530
% This is a sample LaTeX input file.  (Version of 12 August 2004.)
%
% A '%' character causes TeX to ignore all remaining text on the line,
% and is used for comments like this one.

\documentclass{article}      % Specifies the document class

                             % The preamble begins here.
\title{An Example Document}  % Declares the document's title.
\author{Leslie Lamport}      % Declares the author's name.
\date{January 21, 1994}      % Deleting this command produces today's date.

\newcommand{\ip}[2]{(#1, #2)}
                             % Defines \ip{arg1}{arg2} to mean
                             % (arg1, arg2).

%\newcommand{\ip}[2]{\langle #1 | #2\rangle}
                             % This is an alternative definition of
                             % \ip that is commented out.

\begin{document}             % End of preamble and beginning of text.

\maketitle                   % Produces the title.

This is an example input file.  Comparing it with
the output it generates can show you how to
produce a simple document of your own.

\section{Ordinary Text}      % Produces section heading.  Lower-level
                             % sections are begun with similar
                             % \subsection and \subsubsection commands.

The ends  of words and sentences are marked
  by   spaces. It  doesn't matter how many
spaces    you type; one is as good as 100.  The
end of   a line counts as a space.

One   or more   blank lines denote the  end
of  a paragraph.

Since any number of consecutive spaces are treated
like a single one, the formatting of the input
file makes no difference to
      \LaTeX,                % The \LaTeX command generates the LaTeX logo.
but it makes a difference to you.  When you use
\LaTeX, making your input file as easy to read
as possible will be a great help as you write
your document and when you change it.  This sample
file shows how you can add comments to your own input
file.

Because printing is different from typewriting,
there are a number of things that you have to do
differently when preparing an input file than if
you were just typing the document directly.
Quotation marks like
       ``this''
have to be handled specially, as do quotes within
quotes:
       ``\,`this'            % \, separates the double and single quote.
        is what I just
        wrote, not  `that'\,''.

Dashes come in three sizes: an
       intra-word
dash, a medium dash for number ranges like
       1--2,
and a punctuation
       dash---like
this.

A sentence-ending space should be larger than the
space between words within a sentence.  You
sometimes have to type special commands in
conjunction with punctuation characters to get
this right, as in the following sentence.
       Gnats, gnus, etc.\ all  % `\ ' makes an inter-word space.
       begin with G\@.         % \@ marks end-of-sentence punctuation.
You should check the spaces after periods when
reading your output to make sure you haven't
forgotten any special cases.  Generating an
ellipsis
       \ldots\               % `\ ' is needed after `\ldots' because TeX
                             % ignores spaces after command names like \ldots
                             % made from \ + letters.
                             %
                             % Note how a `%' character causes TeX to ignore
                             % the end of the input line, so these blank lines
                             % do not start a new paragraph.
                             %
with the right spacing around the periods requires
a special command.

\LaTeX\ interprets some common characters as
commands, so you must type special commands to
generate them.  These characters include the
following:
       \$ \& \% \# \{ and \}.

In printing, text is usually emphasized with an
       \emph{italic}
type style.

\begin{em}
   A long segment of text can also be emphasized
   in this way.  Text within such a segment can be
   given \emph{additional} emphasis.
\end{em}

It is sometimes necessary to prevent \LaTeX\ from
breaking a line where it might otherwise do so.
This may be at a space, as between the ``Mr.''\ and
``Jones'' in
       ``Mr.~Jones'',        % ~ produces an unbreakable interword space.
or within a word---especially when the word is a
symbol like
       \mbox{\emph{itemnum}}
that makes little sense when hyphenated across
lines.

Footnotes\footnote{This is an example of a footnote.}
pose no problem.

\LaTeX\ is good at typesetting mathematical formulas
like
       \( x-3y + z = 7 \)
or
       \( a_{1} > x^{2n} + y^{2n} > x' \)
or
       \( \ip{A}{B} = \sum_{i} a_{i} b_{i} \).
The spaces you type in a formula are
ignored.  Remember that a letter like
       $x$                   % $ ... $  and  \( ... \)  are equivalent
is a formula when it denotes a mathematical
symbol, and it should be typed as one.

\section{Displayed Text}

Text is displayed by indenting it from the left
margin.  Quotations are commonly displayed.  There
are short quotations
\begin{quote}
   This is a short quotation.  It consists of a
   single paragraph of text.  See how it is formatted.
\end{quote}
and longer ones.
\begin{quotation}
   This is a longer quotation.  It consists of two
   paragraphs of text, neither of which are
   particularly interesting.

   This is the second paragraph of the quotation.  It
   is just as dull as the first paragraph.
\end{quotation}
Another frequently-displayed structure is a list.
The following is an example of an \emph{itemized}
list.
\begin{itemize}
   \item This is the first item of an itemized list.
         Each item in the list is marked with a ``tick''.
         You don't have to worry about what kind of tick
         mark is used.

   \item This is the second item of the list.  It
         contains another list nested inside it.  The inner
         list is an \emph{enumerated} list.
         \begin{enumerate}
            \item This is the first item of an enumerated
                  list that is nested within the itemized list.

            \item This is the second item of the inner list.
                  \LaTeX\ allows you to nest lists deeper than
                  you really should.
         \end{enumerate}
         This is the rest of the second item of the outer
         list.  It is no more interesting than any other
         part of the item.
   \item This is the third item of the list.
\end{itemize}
You can even display poetry.
\begin{verse}
   There is an environment
    for verse \\             % The \\ command separates lines
   Whose features some poets % within a stanza.
   will curse.

                             % One or more blank lines separate stanzas.

   For instead of making\\
   Them do \emph{all} line breaking, \\
   It allows them to put too many words on a line when they'd rather be
   forced to be terse.
\end{verse}

Mathematical formulas may also be displayed.  A
displayed formula
is
one-line long; multiline
formulas require special formatting instructions.
   \[  \ip{\Gamma}{\psi'} = x'' + y^{2} + z_{i}^{n}\]
Don't start a paragraph with a displayed equation,
nor make one a paragraph by itself.

\end{document}               % End of document.

\section{Bench Huge Section 1531}
% synthetic bench section 1531
% This is a sample LaTeX input file.  (Version of 12 August 2004.)
%
% A '%' character causes TeX to ignore all remaining text on the line,
% and is used for comments like this one.

\documentclass{article}      % Specifies the document class

                             % The preamble begins here.
\title{An Example Document}  % Declares the document's title.
\author{Leslie Lamport}      % Declares the author's name.
\date{January 21, 1994}      % Deleting this command produces today's date.

\newcommand{\ip}[2]{(#1, #2)}
                             % Defines \ip{arg1}{arg2} to mean
                             % (arg1, arg2).

%\newcommand{\ip}[2]{\langle #1 | #2\rangle}
                             % This is an alternative definition of
                             % \ip that is commented out.

\begin{document}             % End of preamble and beginning of text.

\maketitle                   % Produces the title.

This is an example input file.  Comparing it with
the output it generates can show you how to
produce a simple document of your own.

\section{Ordinary Text}      % Produces section heading.  Lower-level
                             % sections are begun with similar
                             % \subsection and \subsubsection commands.

The ends  of words and sentences are marked
  by   spaces. It  doesn't matter how many
spaces    you type; one is as good as 100.  The
end of   a line counts as a space.

One   or more   blank lines denote the  end
of  a paragraph.

Since any number of consecutive spaces are treated
like a single one, the formatting of the input
file makes no difference to
      \LaTeX,                % The \LaTeX command generates the LaTeX logo.
but it makes a difference to you.  When you use
\LaTeX, making your input file as easy to read
as possible will be a great help as you write
your document and when you change it.  This sample
file shows how you can add comments to your own input
file.

Because printing is different from typewriting,
there are a number of things that you have to do
differently when preparing an input file than if
you were just typing the document directly.
Quotation marks like
       ``this''
have to be handled specially, as do quotes within
quotes:
       ``\,`this'            % \, separates the double and single quote.
        is what I just
        wrote, not  `that'\,''.

Dashes come in three sizes: an
       intra-word
dash, a medium dash for number ranges like
       1--2,
and a punctuation
       dash---like
this.

A sentence-ending space should be larger than the
space between words within a sentence.  You
sometimes have to type special commands in
conjunction with punctuation characters to get
this right, as in the following sentence.
       Gnats, gnus, etc.\ all  % `\ ' makes an inter-word space.
       begin with G\@.         % \@ marks end-of-sentence punctuation.
You should check the spaces after periods when
reading your output to make sure you haven't
forgotten any special cases.  Generating an
ellipsis
       \ldots\               % `\ ' is needed after `\ldots' because TeX
                             % ignores spaces after command names like \ldots
                             % made from \ + letters.
                             %
                             % Note how a `%' character causes TeX to ignore
                             % the end of the input line, so these blank lines
                             % do not start a new paragraph.
                             %
with the right spacing around the periods requires
a special command.

\LaTeX\ interprets some common characters as
commands, so you must type special commands to
generate them.  These characters include the
following:
       \$ \& \% \# \{ and \}.

In printing, text is usually emphasized with an
       \emph{italic}
type style.

\begin{em}
   A long segment of text can also be emphasized
   in this way.  Text within such a segment can be
   given \emph{additional} emphasis.
\end{em}

It is sometimes necessary to prevent \LaTeX\ from
breaking a line where it might otherwise do so.
This may be at a space, as between the ``Mr.''\ and
``Jones'' in
       ``Mr.~Jones'',        % ~ produces an unbreakable interword space.
or within a word---especially when the word is a
symbol like
       \mbox{\emph{itemnum}}
that makes little sense when hyphenated across
lines.

Footnotes\footnote{This is an example of a footnote.}
pose no problem.

\LaTeX\ is good at typesetting mathematical formulas
like
       \( x-3y + z = 7 \)
or
       \( a_{1} > x^{2n} + y^{2n} > x' \)
or
       \( \ip{A}{B} = \sum_{i} a_{i} b_{i} \).
The spaces you type in a formula are
ignored.  Remember that a letter like
       $x$                   % $ ... $  and  \( ... \)  are equivalent
is a formula when it denotes a mathematical
symbol, and it should be typed as one.

\section{Displayed Text}

Text is displayed by indenting it from the left
margin.  Quotations are commonly displayed.  There
are short quotations
\begin{quote}
   This is a short quotation.  It consists of a
   single paragraph of text.  See how it is formatted.
\end{quote}
and longer ones.
\begin{quotation}
   This is a longer quotation.  It consists of two
   paragraphs of text, neither of which are
   particularly interesting.

   This is the second paragraph of the quotation.  It
   is just as dull as the first paragraph.
\end{quotation}
Another frequently-displayed structure is a list.
The following is an example of an \emph{itemized}
list.
\begin{itemize}
   \item This is the first item of an itemized list.
         Each item in the list is marked with a ``tick''.
         You don't have to worry about what kind of tick
         mark is used.

   \item This is the second item of the list.  It
         contains another list nested inside it.  The inner
         list is an \emph{enumerated} list.
         \begin{enumerate}
            \item This is the first item of an enumerated
                  list that is nested within the itemized list.

            \item This is the second item of the inner list.
                  \LaTeX\ allows you to nest lists deeper than
                  you really should.
         \end{enumerate}
         This is the rest of the second item of the outer
         list.  It is no more interesting than any other
         part of the item.
   \item This is the third item of the list.
\end{itemize}
You can even display poetry.
\begin{verse}
   There is an environment
    for verse \\             % The \\ command separates lines
   Whose features some poets % within a stanza.
   will curse.

                             % One or more blank lines separate stanzas.

   For instead of making\\
   Them do \emph{all} line breaking, \\
   It allows them to put too many words on a line when they'd rather be
   forced to be terse.
\end{verse}

Mathematical formulas may also be displayed.  A
displayed formula
is
one-line long; multiline
formulas require special formatting instructions.
   \[  \ip{\Gamma}{\psi'} = x'' + y^{2} + z_{i}^{n}\]
Don't start a paragraph with a displayed equation,
nor make one a paragraph by itself.

\end{document}               % End of document.

\section{Bench Huge Section 1532}
% synthetic bench section 1532
% This is a sample LaTeX input file.  (Version of 12 August 2004.)
%
% A '%' character causes TeX to ignore all remaining text on the line,
% and is used for comments like this one.

\documentclass{article}      % Specifies the document class

                             % The preamble begins here.
\title{An Example Document}  % Declares the document's title.
\author{Leslie Lamport}      % Declares the author's name.
\date{January 21, 1994}      % Deleting this command produces today's date.

\newcommand{\ip}[2]{(#1, #2)}
                             % Defines \ip{arg1}{arg2} to mean
                             % (arg1, arg2).

%\newcommand{\ip}[2]{\langle #1 | #2\rangle}
                             % This is an alternative definition of
                             % \ip that is commented out.

\begin{document}             % End of preamble and beginning of text.

\maketitle                   % Produces the title.

This is an example input file.  Comparing it with
the output it generates can show you how to
produce a simple document of your own.

\section{Ordinary Text}      % Produces section heading.  Lower-level
                             % sections are begun with similar
                             % \subsection and \subsubsection commands.

The ends  of words and sentences are marked
  by   spaces. It  doesn't matter how many
spaces    you type; one is as good as 100.  The
end of   a line counts as a space.

One   or more   blank lines denote the  end
of  a paragraph.

Since any number of consecutive spaces are treated
like a single one, the formatting of the input
file makes no difference to
      \LaTeX,                % The \LaTeX command generates the LaTeX logo.
but it makes a difference to you.  When you use
\LaTeX, making your input file as easy to read
as possible will be a great help as you write
your document and when you change it.  This sample
file shows how you can add comments to your own input
file.

Because printing is different from typewriting,
there are a number of things that you have to do
differently when preparing an input file than if
you were just typing the document directly.
Quotation marks like
       ``this''
have to be handled specially, as do quotes within
quotes:
       ``\,`this'            % \, separates the double and single quote.
        is what I just
        wrote, not  `that'\,''.

Dashes come in three sizes: an
       intra-word
dash, a medium dash for number ranges like
       1--2,
and a punctuation
       dash---like
this.

A sentence-ending space should be larger than the
space between words within a sentence.  You
sometimes have to type special commands in
conjunction with punctuation characters to get
this right, as in the following sentence.
       Gnats, gnus, etc.\ all  % `\ ' makes an inter-word space.
       begin with G\@.         % \@ marks end-of-sentence punctuation.
You should check the spaces after periods when
reading your output to make sure you haven't
forgotten any special cases.  Generating an
ellipsis
       \ldots\               % `\ ' is needed after `\ldots' because TeX
                             % ignores spaces after command names like \ldots
                             % made from \ + letters.
                             %
                             % Note how a `%' character causes TeX to ignore
                             % the end of the input line, so these blank lines
                             % do not start a new paragraph.
                             %
with the right spacing around the periods requires
a special command.

\LaTeX\ interprets some common characters as
commands, so you must type special commands to
generate them.  These characters include the
following:
       \$ \& \% \# \{ and \}.

In printing, text is usually emphasized with an
       \emph{italic}
type style.

\begin{em}
   A long segment of text can also be emphasized
   in this way.  Text within such a segment can be
   given \emph{additional} emphasis.
\end{em}

It is sometimes necessary to prevent \LaTeX\ from
breaking a line where it might otherwise do so.
This may be at a space, as between the ``Mr.''\ and
``Jones'' in
       ``Mr.~Jones'',        % ~ produces an unbreakable interword space.
or within a word---especially when the word is a
symbol like
       \mbox{\emph{itemnum}}
that makes little sense when hyphenated across
lines.

Footnotes\footnote{This is an example of a footnote.}
pose no problem.

\LaTeX\ is good at typesetting mathematical formulas
like
       \( x-3y + z = 7 \)
or
       \( a_{1} > x^{2n} + y^{2n} > x' \)
or
       \( \ip{A}{B} = \sum_{i} a_{i} b_{i} \).
The spaces you type in a formula are
ignored.  Remember that a letter like
       $x$                   % $ ... $  and  \( ... \)  are equivalent
is a formula when it denotes a mathematical
symbol, and it should be typed as one.

\section{Displayed Text}

Text is displayed by indenting it from the left
margin.  Quotations are commonly displayed.  There
are short quotations
\begin{quote}
   This is a short quotation.  It consists of a
   single paragraph of text.  See how it is formatted.
\end{quote}
and longer ones.
\begin{quotation}
   This is a longer quotation.  It consists of two
   paragraphs of text, neither of which are
   particularly interesting.

   This is the second paragraph of the quotation.  It
   is just as dull as the first paragraph.
\end{quotation}
Another frequently-displayed structure is a list.
The following is an example of an \emph{itemized}
list.
\begin{itemize}
   \item This is the first item of an itemized list.
         Each item in the list is marked with a ``tick''.
         You don't have to worry about what kind of tick
         mark is used.

   \item This is the second item of the list.  It
         contains another list nested inside it.  The inner
         list is an \emph{enumerated} list.
         \begin{enumerate}
            \item This is the first item of an enumerated
                  list that is nested within the itemized list.

            \item This is the second item of the inner list.
                  \LaTeX\ allows you to nest lists deeper than
                  you really should.
         \end{enumerate}
         This is the rest of the second item of the outer
         list.  It is no more interesting than any other
         part of the item.
   \item This is the third item of the list.
\end{itemize}
You can even display poetry.
\begin{verse}
   There is an environment
    for verse \\             % The \\ command separates lines
   Whose features some poets % within a stanza.
   will curse.

                             % One or more blank lines separate stanzas.

   For instead of making\\
   Them do \emph{all} line breaking, \\
   It allows them to put too many words on a line when they'd rather be
   forced to be terse.
\end{verse}

Mathematical formulas may also be displayed.  A
displayed formula
is
one-line long; multiline
formulas require special formatting instructions.
   \[  \ip{\Gamma}{\psi'} = x'' + y^{2} + z_{i}^{n}\]
Don't start a paragraph with a displayed equation,
nor make one a paragraph by itself.

\end{document}               % End of document.

\section{Bench Huge Section 1533}
% synthetic bench section 1533
% This is a sample LaTeX input file.  (Version of 12 August 2004.)
%
% A '%' character causes TeX to ignore all remaining text on the line,
% and is used for comments like this one.

\documentclass{article}      % Specifies the document class

                             % The preamble begins here.
\title{An Example Document}  % Declares the document's title.
\author{Leslie Lamport}      % Declares the author's name.
\date{January 21, 1994}      % Deleting this command produces today's date.

\newcommand{\ip}[2]{(#1, #2)}
                             % Defines \ip{arg1}{arg2} to mean
                             % (arg1, arg2).

%\newcommand{\ip}[2]{\langle #1 | #2\rangle}
                             % This is an alternative definition of
                             % \ip that is commented out.

\begin{document}             % End of preamble and beginning of text.

\maketitle                   % Produces the title.

This is an example input file.  Comparing it with
the output it generates can show you how to
produce a simple document of your own.

\section{Ordinary Text}      % Produces section heading.  Lower-level
                             % sections are begun with similar
                             % \subsection and \subsubsection commands.

The ends  of words and sentences are marked
  by   spaces. It  doesn't matter how many
spaces    you type; one is as good as 100.  The
end of   a line counts as a space.

One   or more   blank lines denote the  end
of  a paragraph.

Since any number of consecutive spaces are treated
like a single one, the formatting of the input
file makes no difference to
      \LaTeX,                % The \LaTeX command generates the LaTeX logo.
but it makes a difference to you.  When you use
\LaTeX, making your input file as easy to read
as possible will be a great help as you write
your document and when you change it.  This sample
file shows how you can add comments to your own input
file.

Because printing is different from typewriting,
there are a number of things that you have to do
differently when preparing an input file than if
you were just typing the document directly.
Quotation marks like
       ``this''
have to be handled specially, as do quotes within
quotes:
       ``\,`this'            % \, separates the double and single quote.
        is what I just
        wrote, not  `that'\,''.

Dashes come in three sizes: an
       intra-word
dash, a medium dash for number ranges like
       1--2,
and a punctuation
       dash---like
this.

A sentence-ending space should be larger than the
space between words within a sentence.  You
sometimes have to type special commands in
conjunction with punctuation characters to get
this right, as in the following sentence.
       Gnats, gnus, etc.\ all  % `\ ' makes an inter-word space.
       begin with G\@.         % \@ marks end-of-sentence punctuation.
You should check the spaces after periods when
reading your output to make sure you haven't
forgotten any special cases.  Generating an
ellipsis
       \ldots\               % `\ ' is needed after `\ldots' because TeX
                             % ignores spaces after command names like \ldots
                             % made from \ + letters.
                             %
                             % Note how a `%' character causes TeX to ignore
                             % the end of the input line, so these blank lines
                             % do not start a new paragraph.
                             %
with the right spacing around the periods requires
a special command.

\LaTeX\ interprets some common characters as
commands, so you must type special commands to
generate them.  These characters include the
following:
       \$ \& \% \# \{ and \}.

In printing, text is usually emphasized with an
       \emph{italic}
type style.

\begin{em}
   A long segment of text can also be emphasized
   in this way.  Text within such a segment can be
   given \emph{additional} emphasis.
\end{em}

It is sometimes necessary to prevent \LaTeX\ from
breaking a line where it might otherwise do so.
This may be at a space, as between the ``Mr.''\ and
``Jones'' in
       ``Mr.~Jones'',        % ~ produces an unbreakable interword space.
or within a word---especially when the word is a
symbol like
       \mbox{\emph{itemnum}}
that makes little sense when hyphenated across
lines.

Footnotes\footnote{This is an example of a footnote.}
pose no problem.

\LaTeX\ is good at typesetting mathematical formulas
like
       \( x-3y + z = 7 \)
or
       \( a_{1} > x^{2n} + y^{2n} > x' \)
or
       \( \ip{A}{B} = \sum_{i} a_{i} b_{i} \).
The spaces you type in a formula are
ignored.  Remember that a letter like
       $x$                   % $ ... $  and  \( ... \)  are equivalent
is a formula when it denotes a mathematical
symbol, and it should be typed as one.

\section{Displayed Text}

Text is displayed by indenting it from the left
margin.  Quotations are commonly displayed.  There
are short quotations
\begin{quote}
   This is a short quotation.  It consists of a
   single paragraph of text.  See how it is formatted.
\end{quote}
and longer ones.
\begin{quotation}
   This is a longer quotation.  It consists of two
   paragraphs of text, neither of which are
   particularly interesting.

   This is the second paragraph of the quotation.  It
   is just as dull as the first paragraph.
\end{quotation}
Another frequently-displayed structure is a list.
The following is an example of an \emph{itemized}
list.
\begin{itemize}
   \item This is the first item of an itemized list.
         Each item in the list is marked with a ``tick''.
         You don't have to worry about what kind of tick
         mark is used.

   \item This is the second item of the list.  It
         contains another list nested inside it.  The inner
         list is an \emph{enumerated} list.
         \begin{enumerate}
            \item This is the first item of an enumerated
                  list that is nested within the itemized list.

            \item This is the second item of the inner list.
                  \LaTeX\ allows you to nest lists deeper than
                  you really should.
         \end{enumerate}
         This is the rest of the second item of the outer
         list.  It is no more interesting than any other
         part of the item.
   \item This is the third item of the list.
\end{itemize}
You can even display poetry.
\begin{verse}
   There is an environment
    for verse \\             % The \\ command separates lines
   Whose features some poets % within a stanza.
   will curse.

                             % One or more blank lines separate stanzas.

   For instead of making\\
   Them do \emph{all} line breaking, \\
   It allows them to put too many words on a line when they'd rather be
   forced to be terse.
\end{verse}

Mathematical formulas may also be displayed.  A
displayed formula
is
one-line long; multiline
formulas require special formatting instructions.
   \[  \ip{\Gamma}{\psi'} = x'' + y^{2} + z_{i}^{n}\]
Don't start a paragraph with a displayed equation,
nor make one a paragraph by itself.

\end{document}               % End of document.

\section{Bench Huge Section 1534}
% synthetic bench section 1534
% This is a sample LaTeX input file.  (Version of 12 August 2004.)
%
% A '%' character causes TeX to ignore all remaining text on the line,
% and is used for comments like this one.

\documentclass{article}      % Specifies the document class

                             % The preamble begins here.
\title{An Example Document}  % Declares the document's title.
\author{Leslie Lamport}      % Declares the author's name.
\date{January 21, 1994}      % Deleting this command produces today's date.

\newcommand{\ip}[2]{(#1, #2)}
                             % Defines \ip{arg1}{arg2} to mean
                             % (arg1, arg2).

%\newcommand{\ip}[2]{\langle #1 | #2\rangle}
                             % This is an alternative definition of
                             % \ip that is commented out.

\begin{document}             % End of preamble and beginning of text.

\maketitle                   % Produces the title.

This is an example input file.  Comparing it with
the output it generates can show you how to
produce a simple document of your own.

\section{Ordinary Text}      % Produces section heading.  Lower-level
                             % sections are begun with similar
                             % \subsection and \subsubsection commands.

The ends  of words and sentences are marked
  by   spaces. It  doesn't matter how many
spaces    you type; one is as good as 100.  The
end of   a line counts as a space.

One   or more   blank lines denote the  end
of  a paragraph.

Since any number of consecutive spaces are treated
like a single one, the formatting of the input
file makes no difference to
      \LaTeX,                % The \LaTeX command generates the LaTeX logo.
but it makes a difference to you.  When you use
\LaTeX, making your input file as easy to read
as possible will be a great help as you write
your document and when you change it.  This sample
file shows how you can add comments to your own input
file.

Because printing is different from typewriting,
there are a number of things that you have to do
differently when preparing an input file than if
you were just typing the document directly.
Quotation marks like
       ``this''
have to be handled specially, as do quotes within
quotes:
       ``\,`this'            % \, separates the double and single quote.
        is what I just
        wrote, not  `that'\,''.

Dashes come in three sizes: an
       intra-word
dash, a medium dash for number ranges like
       1--2,
and a punctuation
       dash---like
this.

A sentence-ending space should be larger than the
space between words within a sentence.  You
sometimes have to type special commands in
conjunction with punctuation characters to get
this right, as in the following sentence.
       Gnats, gnus, etc.\ all  % `\ ' makes an inter-word space.
       begin with G\@.         % \@ marks end-of-sentence punctuation.
You should check the spaces after periods when
reading your output to make sure you haven't
forgotten any special cases.  Generating an
ellipsis
       \ldots\               % `\ ' is needed after `\ldots' because TeX
                             % ignores spaces after command names like \ldots
                             % made from \ + letters.
                             %
                             % Note how a `%' character causes TeX to ignore
                             % the end of the input line, so these blank lines
                             % do not start a new paragraph.
                             %
with the right spacing around the periods requires
a special command.

\LaTeX\ interprets some common characters as
commands, so you must type special commands to
generate them.  These characters include the
following:
       \$ \& \% \# \{ and \}.

In printing, text is usually emphasized with an
       \emph{italic}
type style.

\begin{em}
   A long segment of text can also be emphasized
   in this way.  Text within such a segment can be
   given \emph{additional} emphasis.
\end{em}

It is sometimes necessary to prevent \LaTeX\ from
breaking a line where it might otherwise do so.
This may be at a space, as between the ``Mr.''\ and
``Jones'' in
       ``Mr.~Jones'',        % ~ produces an unbreakable interword space.
or within a word---especially when the word is a
symbol like
       \mbox{\emph{itemnum}}
that makes little sense when hyphenated across
lines.

Footnotes\footnote{This is an example of a footnote.}
pose no problem.

\LaTeX\ is good at typesetting mathematical formulas
like
       \( x-3y + z = 7 \)
or
       \( a_{1} > x^{2n} + y^{2n} > x' \)
or
       \( \ip{A}{B} = \sum_{i} a_{i} b_{i} \).
The spaces you type in a formula are
ignored.  Remember that a letter like
       $x$                   % $ ... $  and  \( ... \)  are equivalent
is a formula when it denotes a mathematical
symbol, and it should be typed as one.

\section{Displayed Text}

Text is displayed by indenting it from the left
margin.  Quotations are commonly displayed.  There
are short quotations
\begin{quote}
   This is a short quotation.  It consists of a
   single paragraph of text.  See how it is formatted.
\end{quote}
and longer ones.
\begin{quotation}
   This is a longer quotation.  It consists of two
   paragraphs of text, neither of which are
   particularly interesting.

   This is the second paragraph of the quotation.  It
   is just as dull as the first paragraph.
\end{quotation}
Another frequently-displayed structure is a list.
The following is an example of an \emph{itemized}
list.
\begin{itemize}
   \item This is the first item of an itemized list.
         Each item in the list is marked with a ``tick''.
         You don't have to worry about what kind of tick
         mark is used.

   \item This is the second item of the list.  It
         contains another list nested inside it.  The inner
         list is an \emph{enumerated} list.
         \begin{enumerate}
            \item This is the first item of an enumerated
                  list that is nested within the itemized list.

            \item This is the second item of the inner list.
                  \LaTeX\ allows you to nest lists deeper than
                  you really should.
         \end{enumerate}
         This is the rest of the second item of the outer
         list.  It is no more interesting than any other
         part of the item.
   \item This is the third item of the list.
\end{itemize}
You can even display poetry.
\begin{verse}
   There is an environment
    for verse \\             % The \\ command separates lines
   Whose features some poets % within a stanza.
   will curse.

                             % One or more blank lines separate stanzas.

   For instead of making\\
   Them do \emph{all} line breaking, \\
   It allows them to put too many words on a line when they'd rather be
   forced to be terse.
\end{verse}

Mathematical formulas may also be displayed.  A
displayed formula
is
one-line long; multiline
formulas require special formatting instructions.
   \[  \ip{\Gamma}{\psi'} = x'' + y^{2} + z_{i}^{n}\]
Don't start a paragraph with a displayed equation,
nor make one a paragraph by itself.

\end{document}               % End of document.

\section{Bench Huge Section 1535}
% synthetic bench section 1535
% This is a sample LaTeX input file.  (Version of 12 August 2004.)
%
% A '%' character causes TeX to ignore all remaining text on the line,
% and is used for comments like this one.

\documentclass{article}      % Specifies the document class

                             % The preamble begins here.
\title{An Example Document}  % Declares the document's title.
\author{Leslie Lamport}      % Declares the author's name.
\date{January 21, 1994}      % Deleting this command produces today's date.

\newcommand{\ip}[2]{(#1, #2)}
                             % Defines \ip{arg1}{arg2} to mean
                             % (arg1, arg2).

%\newcommand{\ip}[2]{\langle #1 | #2\rangle}
                             % This is an alternative definition of
                             % \ip that is commented out.

\begin{document}             % End of preamble and beginning of text.

\maketitle                   % Produces the title.

This is an example input file.  Comparing it with
the output it generates can show you how to
produce a simple document of your own.

\section{Ordinary Text}      % Produces section heading.  Lower-level
                             % sections are begun with similar
                             % \subsection and \subsubsection commands.

The ends  of words and sentences are marked
  by   spaces. It  doesn't matter how many
spaces    you type; one is as good as 100.  The
end of   a line counts as a space.

One   or more   blank lines denote the  end
of  a paragraph.

Since any number of consecutive spaces are treated
like a single one, the formatting of the input
file makes no difference to
      \LaTeX,                % The \LaTeX command generates the LaTeX logo.
but it makes a difference to you.  When you use
\LaTeX, making your input file as easy to read
as possible will be a great help as you write
your document and when you change it.  This sample
file shows how you can add comments to your own input
file.

Because printing is different from typewriting,
there are a number of things that you have to do
differently when preparing an input file than if
you were just typing the document directly.
Quotation marks like
       ``this''
have to be handled specially, as do quotes within
quotes:
       ``\,`this'            % \, separates the double and single quote.
        is what I just
        wrote, not  `that'\,''.

Dashes come in three sizes: an
       intra-word
dash, a medium dash for number ranges like
       1--2,
and a punctuation
       dash---like
this.

A sentence-ending space should be larger than the
space between words within a sentence.  You
sometimes have to type special commands in
conjunction with punctuation characters to get
this right, as in the following sentence.
       Gnats, gnus, etc.\ all  % `\ ' makes an inter-word space.
       begin with G\@.         % \@ marks end-of-sentence punctuation.
You should check the spaces after periods when
reading your output to make sure you haven't
forgotten any special cases.  Generating an
ellipsis
       \ldots\               % `\ ' is needed after `\ldots' because TeX
                             % ignores spaces after command names like \ldots
                             % made from \ + letters.
                             %
                             % Note how a `%' character causes TeX to ignore
                             % the end of the input line, so these blank lines
                             % do not start a new paragraph.
                             %
with the right spacing around the periods requires
a special command.

\LaTeX\ interprets some common characters as
commands, so you must type special commands to
generate them.  These characters include the
following:
       \$ \& \% \# \{ and \}.

In printing, text is usually emphasized with an
       \emph{italic}
type style.

\begin{em}
   A long segment of text can also be emphasized
   in this way.  Text within such a segment can be
   given \emph{additional} emphasis.
\end{em}

It is sometimes necessary to prevent \LaTeX\ from
breaking a line where it might otherwise do so.
This may be at a space, as between the ``Mr.''\ and
``Jones'' in
       ``Mr.~Jones'',        % ~ produces an unbreakable interword space.
or within a word---especially when the word is a
symbol like
       \mbox{\emph{itemnum}}
that makes little sense when hyphenated across
lines.

Footnotes\footnote{This is an example of a footnote.}
pose no problem.

\LaTeX\ is good at typesetting mathematical formulas
like
       \( x-3y + z = 7 \)
or
       \( a_{1} > x^{2n} + y^{2n} > x' \)
or
       \( \ip{A}{B} = \sum_{i} a_{i} b_{i} \).
The spaces you type in a formula are
ignored.  Remember that a letter like
       $x$                   % $ ... $  and  \( ... \)  are equivalent
is a formula when it denotes a mathematical
symbol, and it should be typed as one.

\section{Displayed Text}

Text is displayed by indenting it from the left
margin.  Quotations are commonly displayed.  There
are short quotations
\begin{quote}
   This is a short quotation.  It consists of a
   single paragraph of text.  See how it is formatted.
\end{quote}
and longer ones.
\begin{quotation}
   This is a longer quotation.  It consists of two
   paragraphs of text, neither of which are
   particularly interesting.

   This is the second paragraph of the quotation.  It
   is just as dull as the first paragraph.
\end{quotation}
Another frequently-displayed structure is a list.
The following is an example of an \emph{itemized}
list.
\begin{itemize}
   \item This is the first item of an itemized list.
         Each item in the list is marked with a ``tick''.
         You don't have to worry about what kind of tick
         mark is used.

   \item This is the second item of the list.  It
         contains another list nested inside it.  The inner
         list is an \emph{enumerated} list.
         \begin{enumerate}
            \item This is the first item of an enumerated
                  list that is nested within the itemized list.

            \item This is the second item of the inner list.
                  \LaTeX\ allows you to nest lists deeper than
                  you really should.
         \end{enumerate}
         This is the rest of the second item of the outer
         list.  It is no more interesting than any other
         part of the item.
   \item This is the third item of the list.
\end{itemize}
You can even display poetry.
\begin{verse}
   There is an environment
    for verse \\             % The \\ command separates lines
   Whose features some poets % within a stanza.
   will curse.

                             % One or more blank lines separate stanzas.

   For instead of making\\
   Them do \emph{all} line breaking, \\
   It allows them to put too many words on a line when they'd rather be
   forced to be terse.
\end{verse}

Mathematical formulas may also be displayed.  A
displayed formula
is
one-line long; multiline
formulas require special formatting instructions.
   \[  \ip{\Gamma}{\psi'} = x'' + y^{2} + z_{i}^{n}\]
Don't start a paragraph with a displayed equation,
nor make one a paragraph by itself.

\end{document}               % End of document.

\section{Bench Huge Section 1536}
% synthetic bench section 1536
% This is a sample LaTeX input file.  (Version of 12 August 2004.)
%
% A '%' character causes TeX to ignore all remaining text on the line,
% and is used for comments like this one.

\documentclass{article}      % Specifies the document class

                             % The preamble begins here.
\title{An Example Document}  % Declares the document's title.
\author{Leslie Lamport}      % Declares the author's name.
\date{January 21, 1994}      % Deleting this command produces today's date.

\newcommand{\ip}[2]{(#1, #2)}
                             % Defines \ip{arg1}{arg2} to mean
                             % (arg1, arg2).

%\newcommand{\ip}[2]{\langle #1 | #2\rangle}
                             % This is an alternative definition of
                             % \ip that is commented out.

\begin{document}             % End of preamble and beginning of text.

\maketitle                   % Produces the title.

This is an example input file.  Comparing it with
the output it generates can show you how to
produce a simple document of your own.

\section{Ordinary Text}      % Produces section heading.  Lower-level
                             % sections are begun with similar
                             % \subsection and \subsubsection commands.

The ends  of words and sentences are marked
  by   spaces. It  doesn't matter how many
spaces    you type; one is as good as 100.  The
end of   a line counts as a space.

One   or more   blank lines denote the  end
of  a paragraph.

Since any number of consecutive spaces are treated
like a single one, the formatting of the input
file makes no difference to
      \LaTeX,                % The \LaTeX command generates the LaTeX logo.
but it makes a difference to you.  When you use
\LaTeX, making your input file as easy to read
as possible will be a great help as you write
your document and when you change it.  This sample
file shows how you can add comments to your own input
file.

Because printing is different from typewriting,
there are a number of things that you have to do
differently when preparing an input file than if
you were just typing the document directly.
Quotation marks like
       ``this''
have to be handled specially, as do quotes within
quotes:
       ``\,`this'            % \, separates the double and single quote.
        is what I just
        wrote, not  `that'\,''.

Dashes come in three sizes: an
       intra-word
dash, a medium dash for number ranges like
       1--2,
and a punctuation
       dash---like
this.

A sentence-ending space should be larger than the
space between words within a sentence.  You
sometimes have to type special commands in
conjunction with punctuation characters to get
this right, as in the following sentence.
       Gnats, gnus, etc.\ all  % `\ ' makes an inter-word space.
       begin with G\@.         % \@ marks end-of-sentence punctuation.
You should check the spaces after periods when
reading your output to make sure you haven't
forgotten any special cases.  Generating an
ellipsis
       \ldots\               % `\ ' is needed after `\ldots' because TeX
                             % ignores spaces after command names like \ldots
                             % made from \ + letters.
                             %
                             % Note how a `%' character causes TeX to ignore
                             % the end of the input line, so these blank lines
                             % do not start a new paragraph.
                             %
with the right spacing around the periods requires
a special command.

\LaTeX\ interprets some common characters as
commands, so you must type special commands to
generate them.  These characters include the
following:
       \$ \& \% \# \{ and \}.

In printing, text is usually emphasized with an
       \emph{italic}
type style.

\begin{em}
   A long segment of text can also be emphasized
   in this way.  Text within such a segment can be
   given \emph{additional} emphasis.
\end{em}

It is sometimes necessary to prevent \LaTeX\ from
breaking a line where it might otherwise do so.
This may be at a space, as between the ``Mr.''\ and
``Jones'' in
       ``Mr.~Jones'',        % ~ produces an unbreakable interword space.
or within a word---especially when the word is a
symbol like
       \mbox{\emph{itemnum}}
that makes little sense when hyphenated across
lines.

Footnotes\footnote{This is an example of a footnote.}
pose no problem.

\LaTeX\ is good at typesetting mathematical formulas
like
       \( x-3y + z = 7 \)
or
       \( a_{1} > x^{2n} + y^{2n} > x' \)
or
       \( \ip{A}{B} = \sum_{i} a_{i} b_{i} \).
The spaces you type in a formula are
ignored.  Remember that a letter like
       $x$                   % $ ... $  and  \( ... \)  are equivalent
is a formula when it denotes a mathematical
symbol, and it should be typed as one.

\section{Displayed Text}

Text is displayed by indenting it from the left
margin.  Quotations are commonly displayed.  There
are short quotations
\begin{quote}
   This is a short quotation.  It consists of a
   single paragraph of text.  See how it is formatted.
\end{quote}
and longer ones.
\begin{quotation}
   This is a longer quotation.  It consists of two
   paragraphs of text, neither of which are
   particularly interesting.

   This is the second paragraph of the quotation.  It
   is just as dull as the first paragraph.
\end{quotation}
Another frequently-displayed structure is a list.
The following is an example of an \emph{itemized}
list.
\begin{itemize}
   \item This is the first item of an itemized list.
         Each item in the list is marked with a ``tick''.
         You don't have to worry about what kind of tick
         mark is used.

   \item This is the second item of the list.  It
         contains another list nested inside it.  The inner
         list is an \emph{enumerated} list.
         \begin{enumerate}
            \item This is the first item of an enumerated
                  list that is nested within the itemized list.

            \item This is the second item of the inner list.
                  \LaTeX\ allows you to nest lists deeper than
                  you really should.
         \end{enumerate}
         This is the rest of the second item of the outer
         list.  It is no more interesting than any other
         part of the item.
   \item This is the third item of the list.
\end{itemize}
You can even display poetry.
\begin{verse}
   There is an environment
    for verse \\             % The \\ command separates lines
   Whose features some poets % within a stanza.
   will curse.

                             % One or more blank lines separate stanzas.

   For instead of making\\
   Them do \emph{all} line breaking, \\
   It allows them to put too many words on a line when they'd rather be
   forced to be terse.
\end{verse}

Mathematical formulas may also be displayed.  A
displayed formula
is
one-line long; multiline
formulas require special formatting instructions.
   \[  \ip{\Gamma}{\psi'} = x'' + y^{2} + z_{i}^{n}\]
Don't start a paragraph with a displayed equation,
nor make one a paragraph by itself.

\end{document}               % End of document.

\section{Bench Huge Section 1537}
% synthetic bench section 1537
% This is a sample LaTeX input file.  (Version of 12 August 2004.)
%
% A '%' character causes TeX to ignore all remaining text on the line,
% and is used for comments like this one.

\documentclass{article}      % Specifies the document class

                             % The preamble begins here.
\title{An Example Document}  % Declares the document's title.
\author{Leslie Lamport}      % Declares the author's name.
\date{January 21, 1994}      % Deleting this command produces today's date.

\newcommand{\ip}[2]{(#1, #2)}
                             % Defines \ip{arg1}{arg2} to mean
                             % (arg1, arg2).

%\newcommand{\ip}[2]{\langle #1 | #2\rangle}
                             % This is an alternative definition of
                             % \ip that is commented out.

\begin{document}             % End of preamble and beginning of text.

\maketitle                   % Produces the title.

This is an example input file.  Comparing it with
the output it generates can show you how to
produce a simple document of your own.

\section{Ordinary Text}      % Produces section heading.  Lower-level
                             % sections are begun with similar
                             % \subsection and \subsubsection commands.

The ends  of words and sentences are marked
  by   spaces. It  doesn't matter how many
spaces    you type; one is as good as 100.  The
end of   a line counts as a space.

One   or more   blank lines denote the  end
of  a paragraph.

Since any number of consecutive spaces are treated
like a single one, the formatting of the input
file makes no difference to
      \LaTeX,                % The \LaTeX command generates the LaTeX logo.
but it makes a difference to you.  When you use
\LaTeX, making your input file as easy to read
as possible will be a great help as you write
your document and when you change it.  This sample
file shows how you can add comments to your own input
file.

Because printing is different from typewriting,
there are a number of things that you have to do
differently when preparing an input file than if
you were just typing the document directly.
Quotation marks like
       ``this''
have to be handled specially, as do quotes within
quotes:
       ``\,`this'            % \, separates the double and single quote.
        is what I just
        wrote, not  `that'\,''.

Dashes come in three sizes: an
       intra-word
dash, a medium dash for number ranges like
       1--2,
and a punctuation
       dash---like
this.

A sentence-ending space should be larger than the
space between words within a sentence.  You
sometimes have to type special commands in
conjunction with punctuation characters to get
this right, as in the following sentence.
       Gnats, gnus, etc.\ all  % `\ ' makes an inter-word space.
       begin with G\@.         % \@ marks end-of-sentence punctuation.
You should check the spaces after periods when
reading your output to make sure you haven't
forgotten any special cases.  Generating an
ellipsis
       \ldots\               % `\ ' is needed after `\ldots' because TeX
                             % ignores spaces after command names like \ldots
                             % made from \ + letters.
                             %
                             % Note how a `%' character causes TeX to ignore
                             % the end of the input line, so these blank lines
                             % do not start a new paragraph.
                             %
with the right spacing around the periods requires
a special command.

\LaTeX\ interprets some common characters as
commands, so you must type special commands to
generate them.  These characters include the
following:
       \$ \& \% \# \{ and \}.

In printing, text is usually emphasized with an
       \emph{italic}
type style.

\begin{em}
   A long segment of text can also be emphasized
   in this way.  Text within such a segment can be
   given \emph{additional} emphasis.
\end{em}

It is sometimes necessary to prevent \LaTeX\ from
breaking a line where it might otherwise do so.
This may be at a space, as between the ``Mr.''\ and
``Jones'' in
       ``Mr.~Jones'',        % ~ produces an unbreakable interword space.
or within a word---especially when the word is a
symbol like
       \mbox{\emph{itemnum}}
that makes little sense when hyphenated across
lines.

Footnotes\footnote{This is an example of a footnote.}
pose no problem.

\LaTeX\ is good at typesetting mathematical formulas
like
       \( x-3y + z = 7 \)
or
       \( a_{1} > x^{2n} + y^{2n} > x' \)
or
       \( \ip{A}{B} = \sum_{i} a_{i} b_{i} \).
The spaces you type in a formula are
ignored.  Remember that a letter like
       $x$                   % $ ... $  and  \( ... \)  are equivalent
is a formula when it denotes a mathematical
symbol, and it should be typed as one.

\section{Displayed Text}

Text is displayed by indenting it from the left
margin.  Quotations are commonly displayed.  There
are short quotations
\begin{quote}
   This is a short quotation.  It consists of a
   single paragraph of text.  See how it is formatted.
\end{quote}
and longer ones.
\begin{quotation}
   This is a longer quotation.  It consists of two
   paragraphs of text, neither of which are
   particularly interesting.

   This is the second paragraph of the quotation.  It
   is just as dull as the first paragraph.
\end{quotation}
Another frequently-displayed structure is a list.
The following is an example of an \emph{itemized}
list.
\begin{itemize}
   \item This is the first item of an itemized list.
         Each item in the list is marked with a ``tick''.
         You don't have to worry about what kind of tick
         mark is used.

   \item This is the second item of the list.  It
         contains another list nested inside it.  The inner
         list is an \emph{enumerated} list.
         \begin{enumerate}
            \item This is the first item of an enumerated
                  list that is nested within the itemized list.

            \item This is the second item of the inner list.
                  \LaTeX\ allows you to nest lists deeper than
                  you really should.
         \end{enumerate}
         This is the rest of the second item of the outer
         list.  It is no more interesting than any other
         part of the item.
   \item This is the third item of the list.
\end{itemize}
You can even display poetry.
\begin{verse}
   There is an environment
    for verse \\             % The \\ command separates lines
   Whose features some poets % within a stanza.
   will curse.

                             % One or more blank lines separate stanzas.

   For instead of making\\
   Them do \emph{all} line breaking, \\
   It allows them to put too many words on a line when they'd rather be
   forced to be terse.
\end{verse}

Mathematical formulas may also be displayed.  A
displayed formula
is
one-line long; multiline
formulas require special formatting instructions.
   \[  \ip{\Gamma}{\psi'} = x'' + y^{2} + z_{i}^{n}\]
Don't start a paragraph with a displayed equation,
nor make one a paragraph by itself.

\end{document}               % End of document.

\section{Bench Huge Section 1538}
% synthetic bench section 1538
% This is a sample LaTeX input file.  (Version of 12 August 2004.)
%
% A '%' character causes TeX to ignore all remaining text on the line,
% and is used for comments like this one.

\documentclass{article}      % Specifies the document class

                             % The preamble begins here.
\title{An Example Document}  % Declares the document's title.
\author{Leslie Lamport}      % Declares the author's name.
\date{January 21, 1994}      % Deleting this command produces today's date.

\newcommand{\ip}[2]{(#1, #2)}
                             % Defines \ip{arg1}{arg2} to mean
                             % (arg1, arg2).

%\newcommand{\ip}[2]{\langle #1 | #2\rangle}
                             % This is an alternative definition of
                             % \ip that is commented out.

\begin{document}             % End of preamble and beginning of text.

\maketitle                   % Produces the title.

This is an example input file.  Comparing it with
the output it generates can show you how to
produce a simple document of your own.

\section{Ordinary Text}      % Produces section heading.  Lower-level
                             % sections are begun with similar
                             % \subsection and \subsubsection commands.

The ends  of words and sentences are marked
  by   spaces. It  doesn't matter how many
spaces    you type; one is as good as 100.  The
end of   a line counts as a space.

One   or more   blank lines denote the  end
of  a paragraph.

Since any number of consecutive spaces are treated
like a single one, the formatting of the input
file makes no difference to
      \LaTeX,                % The \LaTeX command generates the LaTeX logo.
but it makes a difference to you.  When you use
\LaTeX, making your input file as easy to read
as possible will be a great help as you write
your document and when you change it.  This sample
file shows how you can add comments to your own input
file.

Because printing is different from typewriting,
there are a number of things that you have to do
differently when preparing an input file than if
you were just typing the document directly.
Quotation marks like
       ``this''
have to be handled specially, as do quotes within
quotes:
       ``\,`this'            % \, separates the double and single quote.
        is what I just
        wrote, not  `that'\,''.

Dashes come in three sizes: an
       intra-word
dash, a medium dash for number ranges like
       1--2,
and a punctuation
       dash---like
this.

A sentence-ending space should be larger than the
space between words within a sentence.  You
sometimes have to type special commands in
conjunction with punctuation characters to get
this right, as in the following sentence.
       Gnats, gnus, etc.\ all  % `\ ' makes an inter-word space.
       begin with G\@.         % \@ marks end-of-sentence punctuation.
You should check the spaces after periods when
reading your output to make sure you haven't
forgotten any special cases.  Generating an
ellipsis
       \ldots\               % `\ ' is needed after `\ldots' because TeX
                             % ignores spaces after command names like \ldots
                             % made from \ + letters.
                             %
                             % Note how a `%' character causes TeX to ignore
                             % the end of the input line, so these blank lines
                             % do not start a new paragraph.
                             %
with the right spacing around the periods requires
a special command.

\LaTeX\ interprets some common characters as
commands, so you must type special commands to
generate them.  These characters include the
following:
       \$ \& \% \# \{ and \}.

In printing, text is usually emphasized with an
       \emph{italic}
type style.

\begin{em}
   A long segment of text can also be emphasized
   in this way.  Text within such a segment can be
   given \emph{additional} emphasis.
\end{em}

It is sometimes necessary to prevent \LaTeX\ from
breaking a line where it might otherwise do so.
This may be at a space, as between the ``Mr.''\ and
``Jones'' in
       ``Mr.~Jones'',        % ~ produces an unbreakable interword space.
or within a word---especially when the word is a
symbol like
       \mbox{\emph{itemnum}}
that makes little sense when hyphenated across
lines.

Footnotes\footnote{This is an example of a footnote.}
pose no problem.

\LaTeX\ is good at typesetting mathematical formulas
like
       \( x-3y + z = 7 \)
or
       \( a_{1} > x^{2n} + y^{2n} > x' \)
or
       \( \ip{A}{B} = \sum_{i} a_{i} b_{i} \).
The spaces you type in a formula are
ignored.  Remember that a letter like
       $x$                   % $ ... $  and  \( ... \)  are equivalent
is a formula when it denotes a mathematical
symbol, and it should be typed as one.

\section{Displayed Text}

Text is displayed by indenting it from the left
margin.  Quotations are commonly displayed.  There
are short quotations
\begin{quote}
   This is a short quotation.  It consists of a
   single paragraph of text.  See how it is formatted.
\end{quote}
and longer ones.
\begin{quotation}
   This is a longer quotation.  It consists of two
   paragraphs of text, neither of which are
   particularly interesting.

   This is the second paragraph of the quotation.  It
   is just as dull as the first paragraph.
\end{quotation}
Another frequently-displayed structure is a list.
The following is an example of an \emph{itemized}
list.
\begin{itemize}
   \item This is the first item of an itemized list.
         Each item in the list is marked with a ``tick''.
         You don't have to worry about what kind of tick
         mark is used.

   \item This is the second item of the list.  It
         contains another list nested inside it.  The inner
         list is an \emph{enumerated} list.
         \begin{enumerate}
            \item This is the first item of an enumerated
                  list that is nested within the itemized list.

            \item This is the second item of the inner list.
                  \LaTeX\ allows you to nest lists deeper than
                  you really should.
         \end{enumerate}
         This is the rest of the second item of the outer
         list.  It is no more interesting than any other
         part of the item.
   \item This is the third item of the list.
\end{itemize}
You can even display poetry.
\begin{verse}
   There is an environment
    for verse \\             % The \\ command separates lines
   Whose features some poets % within a stanza.
   will curse.

                             % One or more blank lines separate stanzas.

   For instead of making\\
   Them do \emph{all} line breaking, \\
   It allows them to put too many words on a line when they'd rather be
   forced to be terse.
\end{verse}

Mathematical formulas may also be displayed.  A
displayed formula
is
one-line long; multiline
formulas require special formatting instructions.
   \[  \ip{\Gamma}{\psi'} = x'' + y^{2} + z_{i}^{n}\]
Don't start a paragraph with a displayed equation,
nor make one a paragraph by itself.

\end{document}               % End of document.

\section{Bench Huge Section 1539}
% synthetic bench section 1539
% This is a sample LaTeX input file.  (Version of 12 August 2004.)
%
% A '%' character causes TeX to ignore all remaining text on the line,
% and is used for comments like this one.

\documentclass{article}      % Specifies the document class

                             % The preamble begins here.
\title{An Example Document}  % Declares the document's title.
\author{Leslie Lamport}      % Declares the author's name.
\date{January 21, 1994}      % Deleting this command produces today's date.

\newcommand{\ip}[2]{(#1, #2)}
                             % Defines \ip{arg1}{arg2} to mean
                             % (arg1, arg2).

%\newcommand{\ip}[2]{\langle #1 | #2\rangle}
                             % This is an alternative definition of
                             % \ip that is commented out.

\begin{document}             % End of preamble and beginning of text.

\maketitle                   % Produces the title.

This is an example input file.  Comparing it with
the output it generates can show you how to
produce a simple document of your own.

\section{Ordinary Text}      % Produces section heading.  Lower-level
                             % sections are begun with similar
                             % \subsection and \subsubsection commands.

The ends  of words and sentences are marked
  by   spaces. It  doesn't matter how many
spaces    you type; one is as good as 100.  The
end of   a line counts as a space.

One   or more   blank lines denote the  end
of  a paragraph.

Since any number of consecutive spaces are treated
like a single one, the formatting of the input
file makes no difference to
      \LaTeX,                % The \LaTeX command generates the LaTeX logo.
but it makes a difference to you.  When you use
\LaTeX, making your input file as easy to read
as possible will be a great help as you write
your document and when you change it.  This sample
file shows how you can add comments to your own input
file.

Because printing is different from typewriting,
there are a number of things that you have to do
differently when preparing an input file than if
you were just typing the document directly.
Quotation marks like
       ``this''
have to be handled specially, as do quotes within
quotes:
       ``\,`this'            % \, separates the double and single quote.
        is what I just
        wrote, not  `that'\,''.

Dashes come in three sizes: an
       intra-word
dash, a medium dash for number ranges like
       1--2,
and a punctuation
       dash---like
this.

A sentence-ending space should be larger than the
space between words within a sentence.  You
sometimes have to type special commands in
conjunction with punctuation characters to get
this right, as in the following sentence.
       Gnats, gnus, etc.\ all  % `\ ' makes an inter-word space.
       begin with G\@.         % \@ marks end-of-sentence punctuation.
You should check the spaces after periods when
reading your output to make sure you haven't
forgotten any special cases.  Generating an
ellipsis
       \ldots\               % `\ ' is needed after `\ldots' because TeX
                             % ignores spaces after command names like \ldots
                             % made from \ + letters.
                             %
                             % Note how a `%' character causes TeX to ignore
                             % the end of the input line, so these blank lines
                             % do not start a new paragraph.
                             %
with the right spacing around the periods requires
a special command.

\LaTeX\ interprets some common characters as
commands, so you must type special commands to
generate them.  These characters include the
following:
       \$ \& \% \# \{ and \}.

In printing, text is usually emphasized with an
       \emph{italic}
type style.

\begin{em}
   A long segment of text can also be emphasized
   in this way.  Text within such a segment can be
   given \emph{additional} emphasis.
\end{em}

It is sometimes necessary to prevent \LaTeX\ from
breaking a line where it might otherwise do so.
This may be at a space, as between the ``Mr.''\ and
``Jones'' in
       ``Mr.~Jones'',        % ~ produces an unbreakable interword space.
or within a word---especially when the word is a
symbol like
       \mbox{\emph{itemnum}}
that makes little sense when hyphenated across
lines.

Footnotes\footnote{This is an example of a footnote.}
pose no problem.

\LaTeX\ is good at typesetting mathematical formulas
like
       \( x-3y + z = 7 \)
or
       \( a_{1} > x^{2n} + y^{2n} > x' \)
or
       \( \ip{A}{B} = \sum_{i} a_{i} b_{i} \).
The spaces you type in a formula are
ignored.  Remember that a letter like
       $x$                   % $ ... $  and  \( ... \)  are equivalent
is a formula when it denotes a mathematical
symbol, and it should be typed as one.

\section{Displayed Text}

Text is displayed by indenting it from the left
margin.  Quotations are commonly displayed.  There
are short quotations
\begin{quote}
   This is a short quotation.  It consists of a
   single paragraph of text.  See how it is formatted.
\end{quote}
and longer ones.
\begin{quotation}
   This is a longer quotation.  It consists of two
   paragraphs of text, neither of which are
   particularly interesting.

   This is the second paragraph of the quotation.  It
   is just as dull as the first paragraph.
\end{quotation}
Another frequently-displayed structure is a list.
The following is an example of an \emph{itemized}
list.
\begin{itemize}
   \item This is the first item of an itemized list.
         Each item in the list is marked with a ``tick''.
         You don't have to worry about what kind of tick
         mark is used.

   \item This is the second item of the list.  It
         contains another list nested inside it.  The inner
         list is an \emph{enumerated} list.
         \begin{enumerate}
            \item This is the first item of an enumerated
                  list that is nested within the itemized list.

            \item This is the second item of the inner list.
                  \LaTeX\ allows you to nest lists deeper than
                  you really should.
         \end{enumerate}
         This is the rest of the second item of the outer
         list.  It is no more interesting than any other
         part of the item.
   \item This is the third item of the list.
\end{itemize}
You can even display poetry.
\begin{verse}
   There is an environment
    for verse \\             % The \\ command separates lines
   Whose features some poets % within a stanza.
   will curse.

                             % One or more blank lines separate stanzas.

   For instead of making\\
   Them do \emph{all} line breaking, \\
   It allows them to put too many words on a line when they'd rather be
   forced to be terse.
\end{verse}

Mathematical formulas may also be displayed.  A
displayed formula
is
one-line long; multiline
formulas require special formatting instructions.
   \[  \ip{\Gamma}{\psi'} = x'' + y^{2} + z_{i}^{n}\]
Don't start a paragraph with a displayed equation,
nor make one a paragraph by itself.

\end{document}               % End of document.

\section{Bench Huge Section 1540}
% synthetic bench section 1540
% This is a sample LaTeX input file.  (Version of 12 August 2004.)
%
% A '%' character causes TeX to ignore all remaining text on the line,
% and is used for comments like this one.

\documentclass{article}      % Specifies the document class

                             % The preamble begins here.
\title{An Example Document}  % Declares the document's title.
\author{Leslie Lamport}      % Declares the author's name.
\date{January 21, 1994}      % Deleting this command produces today's date.

\newcommand{\ip}[2]{(#1, #2)}
                             % Defines \ip{arg1}{arg2} to mean
                             % (arg1, arg2).

%\newcommand{\ip}[2]{\langle #1 | #2\rangle}
                             % This is an alternative definition of
                             % \ip that is commented out.

\begin{document}             % End of preamble and beginning of text.

\maketitle                   % Produces the title.

This is an example input file.  Comparing it with
the output it generates can show you how to
produce a simple document of your own.

\section{Ordinary Text}      % Produces section heading.  Lower-level
                             % sections are begun with similar
                             % \subsection and \subsubsection commands.

The ends  of words and sentences are marked
  by   spaces. It  doesn't matter how many
spaces    you type; one is as good as 100.  The
end of   a line counts as a space.

One   or more   blank lines denote the  end
of  a paragraph.

Since any number of consecutive spaces are treated
like a single one, the formatting of the input
file makes no difference to
      \LaTeX,                % The \LaTeX command generates the LaTeX logo.
but it makes a difference to you.  When you use
\LaTeX, making your input file as easy to read
as possible will be a great help as you write
your document and when you change it.  This sample
file shows how you can add comments to your own input
file.

Because printing is different from typewriting,
there are a number of things that you have to do
differently when preparing an input file than if
you were just typing the document directly.
Quotation marks like
       ``this''
have to be handled specially, as do quotes within
quotes:
       ``\,`this'            % \, separates the double and single quote.
        is what I just
        wrote, not  `that'\,''.

Dashes come in three sizes: an
       intra-word
dash, a medium dash for number ranges like
       1--2,
and a punctuation
       dash---like
this.

A sentence-ending space should be larger than the
space between words within a sentence.  You
sometimes have to type special commands in
conjunction with punctuation characters to get
this right, as in the following sentence.
       Gnats, gnus, etc.\ all  % `\ ' makes an inter-word space.
       begin with G\@.         % \@ marks end-of-sentence punctuation.
You should check the spaces after periods when
reading your output to make sure you haven't
forgotten any special cases.  Generating an
ellipsis
       \ldots\               % `\ ' is needed after `\ldots' because TeX
                             % ignores spaces after command names like \ldots
                             % made from \ + letters.
                             %
                             % Note how a `%' character causes TeX to ignore
                             % the end of the input line, so these blank lines
                             % do not start a new paragraph.
                             %
with the right spacing around the periods requires
a special command.

\LaTeX\ interprets some common characters as
commands, so you must type special commands to
generate them.  These characters include the
following:
       \$ \& \% \# \{ and \}.

In printing, text is usually emphasized with an
       \emph{italic}
type style.

\begin{em}
   A long segment of text can also be emphasized
   in this way.  Text within such a segment can be
   given \emph{additional} emphasis.
\end{em}

It is sometimes necessary to prevent \LaTeX\ from
breaking a line where it might otherwise do so.
This may be at a space, as between the ``Mr.''\ and
``Jones'' in
       ``Mr.~Jones'',        % ~ produces an unbreakable interword space.
or within a word---especially when the word is a
symbol like
       \mbox{\emph{itemnum}}
that makes little sense when hyphenated across
lines.

Footnotes\footnote{This is an example of a footnote.}
pose no problem.

\LaTeX\ is good at typesetting mathematical formulas
like
       \( x-3y + z = 7 \)
or
       \( a_{1} > x^{2n} + y^{2n} > x' \)
or
       \( \ip{A}{B} = \sum_{i} a_{i} b_{i} \).
The spaces you type in a formula are
ignored.  Remember that a letter like
       $x$                   % $ ... $  and  \( ... \)  are equivalent
is a formula when it denotes a mathematical
symbol, and it should be typed as one.

\section{Displayed Text}

Text is displayed by indenting it from the left
margin.  Quotations are commonly displayed.  There
are short quotations
\begin{quote}
   This is a short quotation.  It consists of a
   single paragraph of text.  See how it is formatted.
\end{quote}
and longer ones.
\begin{quotation}
   This is a longer quotation.  It consists of two
   paragraphs of text, neither of which are
   particularly interesting.

   This is the second paragraph of the quotation.  It
   is just as dull as the first paragraph.
\end{quotation}
Another frequently-displayed structure is a list.
The following is an example of an \emph{itemized}
list.
\begin{itemize}
   \item This is the first item of an itemized list.
         Each item in the list is marked with a ``tick''.
         You don't have to worry about what kind of tick
         mark is used.

   \item This is the second item of the list.  It
         contains another list nested inside it.  The inner
         list is an \emph{enumerated} list.
         \begin{enumerate}
            \item This is the first item of an enumerated
                  list that is nested within the itemized list.

            \item This is the second item of the inner list.
                  \LaTeX\ allows you to nest lists deeper than
                  you really should.
         \end{enumerate}
         This is the rest of the second item of the outer
         list.  It is no more interesting than any other
         part of the item.
   \item This is the third item of the list.
\end{itemize}
You can even display poetry.
\begin{verse}
   There is an environment
    for verse \\             % The \\ command separates lines
   Whose features some poets % within a stanza.
   will curse.

                             % One or more blank lines separate stanzas.

   For instead of making\\
   Them do \emph{all} line breaking, \\
   It allows them to put too many words on a line when they'd rather be
   forced to be terse.
\end{verse}

Mathematical formulas may also be displayed.  A
displayed formula
is
one-line long; multiline
formulas require special formatting instructions.
   \[  \ip{\Gamma}{\psi'} = x'' + y^{2} + z_{i}^{n}\]
Don't start a paragraph with a displayed equation,
nor make one a paragraph by itself.

\end{document}               % End of document.

\section{Bench Huge Section 1541}
% synthetic bench section 1541
% This is a sample LaTeX input file.  (Version of 12 August 2004.)
%
% A '%' character causes TeX to ignore all remaining text on the line,
% and is used for comments like this one.

\documentclass{article}      % Specifies the document class

                             % The preamble begins here.
\title{An Example Document}  % Declares the document's title.
\author{Leslie Lamport}      % Declares the author's name.
\date{January 21, 1994}      % Deleting this command produces today's date.

\newcommand{\ip}[2]{(#1, #2)}
                             % Defines \ip{arg1}{arg2} to mean
                             % (arg1, arg2).

%\newcommand{\ip}[2]{\langle #1 | #2\rangle}
                             % This is an alternative definition of
                             % \ip that is commented out.

\begin{document}             % End of preamble and beginning of text.

\maketitle                   % Produces the title.

This is an example input file.  Comparing it with
the output it generates can show you how to
produce a simple document of your own.

\section{Ordinary Text}      % Produces section heading.  Lower-level
                             % sections are begun with similar
                             % \subsection and \subsubsection commands.

The ends  of words and sentences are marked
  by   spaces. It  doesn't matter how many
spaces    you type; one is as good as 100.  The
end of   a line counts as a space.

One   or more   blank lines denote the  end
of  a paragraph.

Since any number of consecutive spaces are treated
like a single one, the formatting of the input
file makes no difference to
      \LaTeX,                % The \LaTeX command generates the LaTeX logo.
but it makes a difference to you.  When you use
\LaTeX, making your input file as easy to read
as possible will be a great help as you write
your document and when you change it.  This sample
file shows how you can add comments to your own input
file.

Because printing is different from typewriting,
there are a number of things that you have to do
differently when preparing an input file than if
you were just typing the document directly.
Quotation marks like
       ``this''
have to be handled specially, as do quotes within
quotes:
       ``\,`this'            % \, separates the double and single quote.
        is what I just
        wrote, not  `that'\,''.

Dashes come in three sizes: an
       intra-word
dash, a medium dash for number ranges like
       1--2,
and a punctuation
       dash---like
this.

A sentence-ending space should be larger than the
space between words within a sentence.  You
sometimes have to type special commands in
conjunction with punctuation characters to get
this right, as in the following sentence.
       Gnats, gnus, etc.\ all  % `\ ' makes an inter-word space.
       begin with G\@.         % \@ marks end-of-sentence punctuation.
You should check the spaces after periods when
reading your output to make sure you haven't
forgotten any special cases.  Generating an
ellipsis
       \ldots\               % `\ ' is needed after `\ldots' because TeX
                             % ignores spaces after command names like \ldots
                             % made from \ + letters.
                             %
                             % Note how a `%' character causes TeX to ignore
                             % the end of the input line, so these blank lines
                             % do not start a new paragraph.
                             %
with the right spacing around the periods requires
a special command.

\LaTeX\ interprets some common characters as
commands, so you must type special commands to
generate them.  These characters include the
following:
       \$ \& \% \# \{ and \}.

In printing, text is usually emphasized with an
       \emph{italic}
type style.

\begin{em}
   A long segment of text can also be emphasized
   in this way.  Text within such a segment can be
   given \emph{additional} emphasis.
\end{em}

It is sometimes necessary to prevent \LaTeX\ from
breaking a line where it might otherwise do so.
This may be at a space, as between the ``Mr.''\ and
``Jones'' in
       ``Mr.~Jones'',        % ~ produces an unbreakable interword space.
or within a word---especially when the word is a
symbol like
       \mbox{\emph{itemnum}}
that makes little sense when hyphenated across
lines.

Footnotes\footnote{This is an example of a footnote.}
pose no problem.

\LaTeX\ is good at typesetting mathematical formulas
like
       \( x-3y + z = 7 \)
or
       \( a_{1} > x^{2n} + y^{2n} > x' \)
or
       \( \ip{A}{B} = \sum_{i} a_{i} b_{i} \).
The spaces you type in a formula are
ignored.  Remember that a letter like
       $x$                   % $ ... $  and  \( ... \)  are equivalent
is a formula when it denotes a mathematical
symbol, and it should be typed as one.

\section{Displayed Text}

Text is displayed by indenting it from the left
margin.  Quotations are commonly displayed.  There
are short quotations
\begin{quote}
   This is a short quotation.  It consists of a
   single paragraph of text.  See how it is formatted.
\end{quote}
and longer ones.
\begin{quotation}
   This is a longer quotation.  It consists of two
   paragraphs of text, neither of which are
   particularly interesting.

   This is the second paragraph of the quotation.  It
   is just as dull as the first paragraph.
\end{quotation}
Another frequently-displayed structure is a list.
The following is an example of an \emph{itemized}
list.
\begin{itemize}
   \item This is the first item of an itemized list.
         Each item in the list is marked with a ``tick''.
         You don't have to worry about what kind of tick
         mark is used.

   \item This is the second item of the list.  It
         contains another list nested inside it.  The inner
         list is an \emph{enumerated} list.
         \begin{enumerate}
            \item This is the first item of an enumerated
                  list that is nested within the itemized list.

            \item This is the second item of the inner list.
                  \LaTeX\ allows you to nest lists deeper than
                  you really should.
         \end{enumerate}
         This is the rest of the second item of the outer
         list.  It is no more interesting than any other
         part of the item.
   \item This is the third item of the list.
\end{itemize}
You can even display poetry.
\begin{verse}
   There is an environment
    for verse \\             % The \\ command separates lines
   Whose features some poets % within a stanza.
   will curse.

                             % One or more blank lines separate stanzas.

   For instead of making\\
   Them do \emph{all} line breaking, \\
   It allows them to put too many words on a line when they'd rather be
   forced to be terse.
\end{verse}

Mathematical formulas may also be displayed.  A
displayed formula
is
one-line long; multiline
formulas require special formatting instructions.
   \[  \ip{\Gamma}{\psi'} = x'' + y^{2} + z_{i}^{n}\]
Don't start a paragraph with a displayed equation,
nor make one a paragraph by itself.

\end{document}               % End of document.

\section{Bench Huge Section 1542}
% synthetic bench section 1542
% This is a sample LaTeX input file.  (Version of 12 August 2004.)
%
% A '%' character causes TeX to ignore all remaining text on the line,
% and is used for comments like this one.

\documentclass{article}      % Specifies the document class

                             % The preamble begins here.
\title{An Example Document}  % Declares the document's title.
\author{Leslie Lamport}      % Declares the author's name.
\date{January 21, 1994}      % Deleting this command produces today's date.

\newcommand{\ip}[2]{(#1, #2)}
                             % Defines \ip{arg1}{arg2} to mean
                             % (arg1, arg2).

%\newcommand{\ip}[2]{\langle #1 | #2\rangle}
                             % This is an alternative definition of
                             % \ip that is commented out.

\begin{document}             % End of preamble and beginning of text.

\maketitle                   % Produces the title.

This is an example input file.  Comparing it with
the output it generates can show you how to
produce a simple document of your own.

\section{Ordinary Text}      % Produces section heading.  Lower-level
                             % sections are begun with similar
                             % \subsection and \subsubsection commands.

The ends  of words and sentences are marked
  by   spaces. It  doesn't matter how many
spaces    you type; one is as good as 100.  The
end of   a line counts as a space.

One   or more   blank lines denote the  end
of  a paragraph.

Since any number of consecutive spaces are treated
like a single one, the formatting of the input
file makes no difference to
      \LaTeX,                % The \LaTeX command generates the LaTeX logo.
but it makes a difference to you.  When you use
\LaTeX, making your input file as easy to read
as possible will be a great help as you write
your document and when you change it.  This sample
file shows how you can add comments to your own input
file.

Because printing is different from typewriting,
there are a number of things that you have to do
differently when preparing an input file than if
you were just typing the document directly.
Quotation marks like
       ``this''
have to be handled specially, as do quotes within
quotes:
       ``\,`this'            % \, separates the double and single quote.
        is what I just
        wrote, not  `that'\,''.

Dashes come in three sizes: an
       intra-word
dash, a medium dash for number ranges like
       1--2,
and a punctuation
       dash---like
this.

A sentence-ending space should be larger than the
space between words within a sentence.  You
sometimes have to type special commands in
conjunction with punctuation characters to get
this right, as in the following sentence.
       Gnats, gnus, etc.\ all  % `\ ' makes an inter-word space.
       begin with G\@.         % \@ marks end-of-sentence punctuation.
You should check the spaces after periods when
reading your output to make sure you haven't
forgotten any special cases.  Generating an
ellipsis
       \ldots\               % `\ ' is needed after `\ldots' because TeX
                             % ignores spaces after command names like \ldots
                             % made from \ + letters.
                             %
                             % Note how a `%' character causes TeX to ignore
                             % the end of the input line, so these blank lines
                             % do not start a new paragraph.
                             %
with the right spacing around the periods requires
a special command.

\LaTeX\ interprets some common characters as
commands, so you must type special commands to
generate them.  These characters include the
following:
       \$ \& \% \# \{ and \}.

In printing, text is usually emphasized with an
       \emph{italic}
type style.

\begin{em}
   A long segment of text can also be emphasized
   in this way.  Text within such a segment can be
   given \emph{additional} emphasis.
\end{em}

It is sometimes necessary to prevent \LaTeX\ from
breaking a line where it might otherwise do so.
This may be at a space, as between the ``Mr.''\ and
``Jones'' in
       ``Mr.~Jones'',        % ~ produces an unbreakable interword space.
or within a word---especially when the word is a
symbol like
       \mbox{\emph{itemnum}}
that makes little sense when hyphenated across
lines.

Footnotes\footnote{This is an example of a footnote.}
pose no problem.

\LaTeX\ is good at typesetting mathematical formulas
like
       \( x-3y + z = 7 \)
or
       \( a_{1} > x^{2n} + y^{2n} > x' \)
or
       \( \ip{A}{B} = \sum_{i} a_{i} b_{i} \).
The spaces you type in a formula are
ignored.  Remember that a letter like
       $x$                   % $ ... $  and  \( ... \)  are equivalent
is a formula when it denotes a mathematical
symbol, and it should be typed as one.

\section{Displayed Text}

Text is displayed by indenting it from the left
margin.  Quotations are commonly displayed.  There
are short quotations
\begin{quote}
   This is a short quotation.  It consists of a
   single paragraph of text.  See how it is formatted.
\end{quote}
and longer ones.
\begin{quotation}
   This is a longer quotation.  It consists of two
   paragraphs of text, neither of which are
   particularly interesting.

   This is the second paragraph of the quotation.  It
   is just as dull as the first paragraph.
\end{quotation}
Another frequently-displayed structure is a list.
The following is an example of an \emph{itemized}
list.
\begin{itemize}
   \item This is the first item of an itemized list.
         Each item in the list is marked with a ``tick''.
         You don't have to worry about what kind of tick
         mark is used.

   \item This is the second item of the list.  It
         contains another list nested inside it.  The inner
         list is an \emph{enumerated} list.
         \begin{enumerate}
            \item This is the first item of an enumerated
                  list that is nested within the itemized list.

            \item This is the second item of the inner list.
                  \LaTeX\ allows you to nest lists deeper than
                  you really should.
         \end{enumerate}
         This is the rest of the second item of the outer
         list.  It is no more interesting than any other
         part of the item.
   \item This is the third item of the list.
\end{itemize}
You can even display poetry.
\begin{verse}
   There is an environment
    for verse \\             % The \\ command separates lines
   Whose features some poets % within a stanza.
   will curse.

                             % One or more blank lines separate stanzas.

   For instead of making\\
   Them do \emph{all} line breaking, \\
   It allows them to put too many words on a line when they'd rather be
   forced to be terse.
\end{verse}

Mathematical formulas may also be displayed.  A
displayed formula
is
one-line long; multiline
formulas require special formatting instructions.
   \[  \ip{\Gamma}{\psi'} = x'' + y^{2} + z_{i}^{n}\]
Don't start a paragraph with a displayed equation,
nor make one a paragraph by itself.

\end{document}               % End of document.

\section{Bench Huge Section 1543}
% synthetic bench section 1543
% This is a sample LaTeX input file.  (Version of 12 August 2004.)
%
% A '%' character causes TeX to ignore all remaining text on the line,
% and is used for comments like this one.

\documentclass{article}      % Specifies the document class

                             % The preamble begins here.
\title{An Example Document}  % Declares the document's title.
\author{Leslie Lamport}      % Declares the author's name.
\date{January 21, 1994}      % Deleting this command produces today's date.

\newcommand{\ip}[2]{(#1, #2)}
                             % Defines \ip{arg1}{arg2} to mean
                             % (arg1, arg2).

%\newcommand{\ip}[2]{\langle #1 | #2\rangle}
                             % This is an alternative definition of
                             % \ip that is commented out.

\begin{document}             % End of preamble and beginning of text.

\maketitle                   % Produces the title.

This is an example input file.  Comparing it with
the output it generates can show you how to
produce a simple document of your own.

\section{Ordinary Text}      % Produces section heading.  Lower-level
                             % sections are begun with similar
                             % \subsection and \subsubsection commands.

The ends  of words and sentences are marked
  by   spaces. It  doesn't matter how many
spaces    you type; one is as good as 100.  The
end of   a line counts as a space.

One   or more   blank lines denote the  end
of  a paragraph.

Since any number of consecutive spaces are treated
like a single one, the formatting of the input
file makes no difference to
      \LaTeX,                % The \LaTeX command generates the LaTeX logo.
but it makes a difference to you.  When you use
\LaTeX, making your input file as easy to read
as possible will be a great help as you write
your document and when you change it.  This sample
file shows how you can add comments to your own input
file.

Because printing is different from typewriting,
there are a number of things that you have to do
differently when preparing an input file than if
you were just typing the document directly.
Quotation marks like
       ``this''
have to be handled specially, as do quotes within
quotes:
       ``\,`this'            % \, separates the double and single quote.
        is what I just
        wrote, not  `that'\,''.

Dashes come in three sizes: an
       intra-word
dash, a medium dash for number ranges like
       1--2,
and a punctuation
       dash---like
this.

A sentence-ending space should be larger than the
space between words within a sentence.  You
sometimes have to type special commands in
conjunction with punctuation characters to get
this right, as in the following sentence.
       Gnats, gnus, etc.\ all  % `\ ' makes an inter-word space.
       begin with G\@.         % \@ marks end-of-sentence punctuation.
You should check the spaces after periods when
reading your output to make sure you haven't
forgotten any special cases.  Generating an
ellipsis
       \ldots\               % `\ ' is needed after `\ldots' because TeX
                             % ignores spaces after command names like \ldots
                             % made from \ + letters.
                             %
                             % Note how a `%' character causes TeX to ignore
                             % the end of the input line, so these blank lines
                             % do not start a new paragraph.
                             %
with the right spacing around the periods requires
a special command.

\LaTeX\ interprets some common characters as
commands, so you must type special commands to
generate them.  These characters include the
following:
       \$ \& \% \# \{ and \}.

In printing, text is usually emphasized with an
       \emph{italic}
type style.

\begin{em}
   A long segment of text can also be emphasized
   in this way.  Text within such a segment can be
   given \emph{additional} emphasis.
\end{em}

It is sometimes necessary to prevent \LaTeX\ from
breaking a line where it might otherwise do so.
This may be at a space, as between the ``Mr.''\ and
``Jones'' in
       ``Mr.~Jones'',        % ~ produces an unbreakable interword space.
or within a word---especially when the word is a
symbol like
       \mbox{\emph{itemnum}}
that makes little sense when hyphenated across
lines.

Footnotes\footnote{This is an example of a footnote.}
pose no problem.

\LaTeX\ is good at typesetting mathematical formulas
like
       \( x-3y + z = 7 \)
or
       \( a_{1} > x^{2n} + y^{2n} > x' \)
or
       \( \ip{A}{B} = \sum_{i} a_{i} b_{i} \).
The spaces you type in a formula are
ignored.  Remember that a letter like
       $x$                   % $ ... $  and  \( ... \)  are equivalent
is a formula when it denotes a mathematical
symbol, and it should be typed as one.

\section{Displayed Text}

Text is displayed by indenting it from the left
margin.  Quotations are commonly displayed.  There
are short quotations
\begin{quote}
   This is a short quotation.  It consists of a
   single paragraph of text.  See how it is formatted.
\end{quote}
and longer ones.
\begin{quotation}
   This is a longer quotation.  It consists of two
   paragraphs of text, neither of which are
   particularly interesting.

   This is the second paragraph of the quotation.  It
   is just as dull as the first paragraph.
\end{quotation}
Another frequently-displayed structure is a list.
The following is an example of an \emph{itemized}
list.
\begin{itemize}
   \item This is the first item of an itemized list.
         Each item in the list is marked with a ``tick''.
         You don't have to worry about what kind of tick
         mark is used.

   \item This is the second item of the list.  It
         contains another list nested inside it.  The inner
         list is an \emph{enumerated} list.
         \begin{enumerate}
            \item This is the first item of an enumerated
                  list that is nested within the itemized list.

            \item This is the second item of the inner list.
                  \LaTeX\ allows you to nest lists deeper than
                  you really should.
         \end{enumerate}
         This is the rest of the second item of the outer
         list.  It is no more interesting than any other
         part of the item.
   \item This is the third item of the list.
\end{itemize}
You can even display poetry.
\begin{verse}
   There is an environment
    for verse \\             % The \\ command separates lines
   Whose features some poets % within a stanza.
   will curse.

                             % One or more blank lines separate stanzas.

   For instead of making\\
   Them do \emph{all} line breaking, \\
   It allows them to put too many words on a line when they'd rather be
   forced to be terse.
\end{verse}

Mathematical formulas may also be displayed.  A
displayed formula
is
one-line long; multiline
formulas require special formatting instructions.
   \[  \ip{\Gamma}{\psi'} = x'' + y^{2} + z_{i}^{n}\]
Don't start a paragraph with a displayed equation,
nor make one a paragraph by itself.

\end{document}               % End of document.

\section{Bench Huge Section 1544}
% synthetic bench section 1544
% This is a sample LaTeX input file.  (Version of 12 August 2004.)
%
% A '%' character causes TeX to ignore all remaining text on the line,
% and is used for comments like this one.

\documentclass{article}      % Specifies the document class

                             % The preamble begins here.
\title{An Example Document}  % Declares the document's title.
\author{Leslie Lamport}      % Declares the author's name.
\date{January 21, 1994}      % Deleting this command produces today's date.

\newcommand{\ip}[2]{(#1, #2)}
                             % Defines \ip{arg1}{arg2} to mean
                             % (arg1, arg2).

%\newcommand{\ip}[2]{\langle #1 | #2\rangle}
                             % This is an alternative definition of
                             % \ip that is commented out.

\begin{document}             % End of preamble and beginning of text.

\maketitle                   % Produces the title.

This is an example input file.  Comparing it with
the output it generates can show you how to
produce a simple document of your own.

\section{Ordinary Text}      % Produces section heading.  Lower-level
                             % sections are begun with similar
                             % \subsection and \subsubsection commands.

The ends  of words and sentences are marked
  by   spaces. It  doesn't matter how many
spaces    you type; one is as good as 100.  The
end of   a line counts as a space.

One   or more   blank lines denote the  end
of  a paragraph.

Since any number of consecutive spaces are treated
like a single one, the formatting of the input
file makes no difference to
      \LaTeX,                % The \LaTeX command generates the LaTeX logo.
but it makes a difference to you.  When you use
\LaTeX, making your input file as easy to read
as possible will be a great help as you write
your document and when you change it.  This sample
file shows how you can add comments to your own input
file.

Because printing is different from typewriting,
there are a number of things that you have to do
differently when preparing an input file than if
you were just typing the document directly.
Quotation marks like
       ``this''
have to be handled specially, as do quotes within
quotes:
       ``\,`this'            % \, separates the double and single quote.
        is what I just
        wrote, not  `that'\,''.

Dashes come in three sizes: an
       intra-word
dash, a medium dash for number ranges like
       1--2,
and a punctuation
       dash---like
this.

A sentence-ending space should be larger than the
space between words within a sentence.  You
sometimes have to type special commands in
conjunction with punctuation characters to get
this right, as in the following sentence.
       Gnats, gnus, etc.\ all  % `\ ' makes an inter-word space.
       begin with G\@.         % \@ marks end-of-sentence punctuation.
You should check the spaces after periods when
reading your output to make sure you haven't
forgotten any special cases.  Generating an
ellipsis
       \ldots\               % `\ ' is needed after `\ldots' because TeX
                             % ignores spaces after command names like \ldots
                             % made from \ + letters.
                             %
                             % Note how a `%' character causes TeX to ignore
                             % the end of the input line, so these blank lines
                             % do not start a new paragraph.
                             %
with the right spacing around the periods requires
a special command.

\LaTeX\ interprets some common characters as
commands, so you must type special commands to
generate them.  These characters include the
following:
       \$ \& \% \# \{ and \}.

In printing, text is usually emphasized with an
       \emph{italic}
type style.

\begin{em}
   A long segment of text can also be emphasized
   in this way.  Text within such a segment can be
   given \emph{additional} emphasis.
\end{em}

It is sometimes necessary to prevent \LaTeX\ from
breaking a line where it might otherwise do so.
This may be at a space, as between the ``Mr.''\ and
``Jones'' in
       ``Mr.~Jones'',        % ~ produces an unbreakable interword space.
or within a word---especially when the word is a
symbol like
       \mbox{\emph{itemnum}}
that makes little sense when hyphenated across
lines.

Footnotes\footnote{This is an example of a footnote.}
pose no problem.

\LaTeX\ is good at typesetting mathematical formulas
like
       \( x-3y + z = 7 \)
or
       \( a_{1} > x^{2n} + y^{2n} > x' \)
or
       \( \ip{A}{B} = \sum_{i} a_{i} b_{i} \).
The spaces you type in a formula are
ignored.  Remember that a letter like
       $x$                   % $ ... $  and  \( ... \)  are equivalent
is a formula when it denotes a mathematical
symbol, and it should be typed as one.

\section{Displayed Text}

Text is displayed by indenting it from the left
margin.  Quotations are commonly displayed.  There
are short quotations
\begin{quote}
   This is a short quotation.  It consists of a
   single paragraph of text.  See how it is formatted.
\end{quote}
and longer ones.
\begin{quotation}
   This is a longer quotation.  It consists of two
   paragraphs of text, neither of which are
   particularly interesting.

   This is the second paragraph of the quotation.  It
   is just as dull as the first paragraph.
\end{quotation}
Another frequently-displayed structure is a list.
The following is an example of an \emph{itemized}
list.
\begin{itemize}
   \item This is the first item of an itemized list.
         Each item in the list is marked with a ``tick''.
         You don't have to worry about what kind of tick
         mark is used.

   \item This is the second item of the list.  It
         contains another list nested inside it.  The inner
         list is an \emph{enumerated} list.
         \begin{enumerate}
            \item This is the first item of an enumerated
                  list that is nested within the itemized list.

            \item This is the second item of the inner list.
                  \LaTeX\ allows you to nest lists deeper than
                  you really should.
         \end{enumerate}
         This is the rest of the second item of the outer
         list.  It is no more interesting than any other
         part of the item.
   \item This is the third item of the list.
\end{itemize}
You can even display poetry.
\begin{verse}
   There is an environment
    for verse \\             % The \\ command separates lines
   Whose features some poets % within a stanza.
   will curse.

                             % One or more blank lines separate stanzas.

   For instead of making\\
   Them do \emph{all} line breaking, \\
   It allows them to put too many words on a line when they'd rather be
   forced to be terse.
\end{verse}

Mathematical formulas may also be displayed.  A
displayed formula
is
one-line long; multiline
formulas require special formatting instructions.
   \[  \ip{\Gamma}{\psi'} = x'' + y^{2} + z_{i}^{n}\]
Don't start a paragraph with a displayed equation,
nor make one a paragraph by itself.

\end{document}               % End of document.

\section{Bench Huge Section 1545}
% synthetic bench section 1545
% This is a sample LaTeX input file.  (Version of 12 August 2004.)
%
% A '%' character causes TeX to ignore all remaining text on the line,
% and is used for comments like this one.

\documentclass{article}      % Specifies the document class

                             % The preamble begins here.
\title{An Example Document}  % Declares the document's title.
\author{Leslie Lamport}      % Declares the author's name.
\date{January 21, 1994}      % Deleting this command produces today's date.

\newcommand{\ip}[2]{(#1, #2)}
                             % Defines \ip{arg1}{arg2} to mean
                             % (arg1, arg2).

%\newcommand{\ip}[2]{\langle #1 | #2\rangle}
                             % This is an alternative definition of
                             % \ip that is commented out.

\begin{document}             % End of preamble and beginning of text.

\maketitle                   % Produces the title.

This is an example input file.  Comparing it with
the output it generates can show you how to
produce a simple document of your own.

\section{Ordinary Text}      % Produces section heading.  Lower-level
                             % sections are begun with similar
                             % \subsection and \subsubsection commands.

The ends  of words and sentences are marked
  by   spaces. It  doesn't matter how many
spaces    you type; one is as good as 100.  The
end of   a line counts as a space.

One   or more   blank lines denote the  end
of  a paragraph.

Since any number of consecutive spaces are treated
like a single one, the formatting of the input
file makes no difference to
      \LaTeX,                % The \LaTeX command generates the LaTeX logo.
but it makes a difference to you.  When you use
\LaTeX, making your input file as easy to read
as possible will be a great help as you write
your document and when you change it.  This sample
file shows how you can add comments to your own input
file.

Because printing is different from typewriting,
there are a number of things that you have to do
differently when preparing an input file than if
you were just typing the document directly.
Quotation marks like
       ``this''
have to be handled specially, as do quotes within
quotes:
       ``\,`this'            % \, separates the double and single quote.
        is what I just
        wrote, not  `that'\,''.

Dashes come in three sizes: an
       intra-word
dash, a medium dash for number ranges like
       1--2,
and a punctuation
       dash---like
this.

A sentence-ending space should be larger than the
space between words within a sentence.  You
sometimes have to type special commands in
conjunction with punctuation characters to get
this right, as in the following sentence.
       Gnats, gnus, etc.\ all  % `\ ' makes an inter-word space.
       begin with G\@.         % \@ marks end-of-sentence punctuation.
You should check the spaces after periods when
reading your output to make sure you haven't
forgotten any special cases.  Generating an
ellipsis
       \ldots\               % `\ ' is needed after `\ldots' because TeX
                             % ignores spaces after command names like \ldots
                             % made from \ + letters.
                             %
                             % Note how a `%' character causes TeX to ignore
                             % the end of the input line, so these blank lines
                             % do not start a new paragraph.
                             %
with the right spacing around the periods requires
a special command.

\LaTeX\ interprets some common characters as
commands, so you must type special commands to
generate them.  These characters include the
following:
       \$ \& \% \# \{ and \}.

In printing, text is usually emphasized with an
       \emph{italic}
type style.

\begin{em}
   A long segment of text can also be emphasized
   in this way.  Text within such a segment can be
   given \emph{additional} emphasis.
\end{em}

It is sometimes necessary to prevent \LaTeX\ from
breaking a line where it might otherwise do so.
This may be at a space, as between the ``Mr.''\ and
``Jones'' in
       ``Mr.~Jones'',        % ~ produces an unbreakable interword space.
or within a word---especially when the word is a
symbol like
       \mbox{\emph{itemnum}}
that makes little sense when hyphenated across
lines.

Footnotes\footnote{This is an example of a footnote.}
pose no problem.

\LaTeX\ is good at typesetting mathematical formulas
like
       \( x-3y + z = 7 \)
or
       \( a_{1} > x^{2n} + y^{2n} > x' \)
or
       \( \ip{A}{B} = \sum_{i} a_{i} b_{i} \).
The spaces you type in a formula are
ignored.  Remember that a letter like
       $x$                   % $ ... $  and  \( ... \)  are equivalent
is a formula when it denotes a mathematical
symbol, and it should be typed as one.

\section{Displayed Text}

Text is displayed by indenting it from the left
margin.  Quotations are commonly displayed.  There
are short quotations
\begin{quote}
   This is a short quotation.  It consists of a
   single paragraph of text.  See how it is formatted.
\end{quote}
and longer ones.
\begin{quotation}
   This is a longer quotation.  It consists of two
   paragraphs of text, neither of which are
   particularly interesting.

   This is the second paragraph of the quotation.  It
   is just as dull as the first paragraph.
\end{quotation}
Another frequently-displayed structure is a list.
The following is an example of an \emph{itemized}
list.
\begin{itemize}
   \item This is the first item of an itemized list.
         Each item in the list is marked with a ``tick''.
         You don't have to worry about what kind of tick
         mark is used.

   \item This is the second item of the list.  It
         contains another list nested inside it.  The inner
         list is an \emph{enumerated} list.
         \begin{enumerate}
            \item This is the first item of an enumerated
                  list that is nested within the itemized list.

            \item This is the second item of the inner list.
                  \LaTeX\ allows you to nest lists deeper than
                  you really should.
         \end{enumerate}
         This is the rest of the second item of the outer
         list.  It is no more interesting than any other
         part of the item.
   \item This is the third item of the list.
\end{itemize}
You can even display poetry.
\begin{verse}
   There is an environment
    for verse \\             % The \\ command separates lines
   Whose features some poets % within a stanza.
   will curse.

                             % One or more blank lines separate stanzas.

   For instead of making\\
   Them do \emph{all} line breaking, \\
   It allows them to put too many words on a line when they'd rather be
   forced to be terse.
\end{verse}

Mathematical formulas may also be displayed.  A
displayed formula
is
one-line long; multiline
formulas require special formatting instructions.
   \[  \ip{\Gamma}{\psi'} = x'' + y^{2} + z_{i}^{n}\]
Don't start a paragraph with a displayed equation,
nor make one a paragraph by itself.

\end{document}               % End of document.

\section{Bench Huge Section 1546}
% synthetic bench section 1546
% This is a sample LaTeX input file.  (Version of 12 August 2004.)
%
% A '%' character causes TeX to ignore all remaining text on the line,
% and is used for comments like this one.

\documentclass{article}      % Specifies the document class

                             % The preamble begins here.
\title{An Example Document}  % Declares the document's title.
\author{Leslie Lamport}      % Declares the author's name.
\date{January 21, 1994}      % Deleting this command produces today's date.

\newcommand{\ip}[2]{(#1, #2)}
                             % Defines \ip{arg1}{arg2} to mean
                             % (arg1, arg2).

%\newcommand{\ip}[2]{\langle #1 | #2\rangle}
                             % This is an alternative definition of
                             % \ip that is commented out.

\begin{document}             % End of preamble and beginning of text.

\maketitle                   % Produces the title.

This is an example input file.  Comparing it with
the output it generates can show you how to
produce a simple document of your own.

\section{Ordinary Text}      % Produces section heading.  Lower-level
                             % sections are begun with similar
                             % \subsection and \subsubsection commands.

The ends  of words and sentences are marked
  by   spaces. It  doesn't matter how many
spaces    you type; one is as good as 100.  The
end of   a line counts as a space.

One   or more   blank lines denote the  end
of  a paragraph.

Since any number of consecutive spaces are treated
like a single one, the formatting of the input
file makes no difference to
      \LaTeX,                % The \LaTeX command generates the LaTeX logo.
but it makes a difference to you.  When you use
\LaTeX, making your input file as easy to read
as possible will be a great help as you write
your document and when you change it.  This sample
file shows how you can add comments to your own input
file.

Because printing is different from typewriting,
there are a number of things that you have to do
differently when preparing an input file than if
you were just typing the document directly.
Quotation marks like
       ``this''
have to be handled specially, as do quotes within
quotes:
       ``\,`this'            % \, separates the double and single quote.
        is what I just
        wrote, not  `that'\,''.

Dashes come in three sizes: an
       intra-word
dash, a medium dash for number ranges like
       1--2,
and a punctuation
       dash---like
this.

A sentence-ending space should be larger than the
space between words within a sentence.  You
sometimes have to type special commands in
conjunction with punctuation characters to get
this right, as in the following sentence.
       Gnats, gnus, etc.\ all  % `\ ' makes an inter-word space.
       begin with G\@.         % \@ marks end-of-sentence punctuation.
You should check the spaces after periods when
reading your output to make sure you haven't
forgotten any special cases.  Generating an
ellipsis
       \ldots\               % `\ ' is needed after `\ldots' because TeX
                             % ignores spaces after command names like \ldots
                             % made from \ + letters.
                             %
                             % Note how a `%' character causes TeX to ignore
                             % the end of the input line, so these blank lines
                             % do not start a new paragraph.
                             %
with the right spacing around the periods requires
a special command.

\LaTeX\ interprets some common characters as
commands, so you must type special commands to
generate them.  These characters include the
following:
       \$ \& \% \# \{ and \}.

In printing, text is usually emphasized with an
       \emph{italic}
type style.

\begin{em}
   A long segment of text can also be emphasized
   in this way.  Text within such a segment can be
   given \emph{additional} emphasis.
\end{em}

It is sometimes necessary to prevent \LaTeX\ from
breaking a line where it might otherwise do so.
This may be at a space, as between the ``Mr.''\ and
``Jones'' in
       ``Mr.~Jones'',        % ~ produces an unbreakable interword space.
or within a word---especially when the word is a
symbol like
       \mbox{\emph{itemnum}}
that makes little sense when hyphenated across
lines.

Footnotes\footnote{This is an example of a footnote.}
pose no problem.

\LaTeX\ is good at typesetting mathematical formulas
like
       \( x-3y + z = 7 \)
or
       \( a_{1} > x^{2n} + y^{2n} > x' \)
or
       \( \ip{A}{B} = \sum_{i} a_{i} b_{i} \).
The spaces you type in a formula are
ignored.  Remember that a letter like
       $x$                   % $ ... $  and  \( ... \)  are equivalent
is a formula when it denotes a mathematical
symbol, and it should be typed as one.

\section{Displayed Text}

Text is displayed by indenting it from the left
margin.  Quotations are commonly displayed.  There
are short quotations
\begin{quote}
   This is a short quotation.  It consists of a
   single paragraph of text.  See how it is formatted.
\end{quote}
and longer ones.
\begin{quotation}
   This is a longer quotation.  It consists of two
   paragraphs of text, neither of which are
   particularly interesting.

   This is the second paragraph of the quotation.  It
   is just as dull as the first paragraph.
\end{quotation}
Another frequently-displayed structure is a list.
The following is an example of an \emph{itemized}
list.
\begin{itemize}
   \item This is the first item of an itemized list.
         Each item in the list is marked with a ``tick''.
         You don't have to worry about what kind of tick
         mark is used.

   \item This is the second item of the list.  It
         contains another list nested inside it.  The inner
         list is an \emph{enumerated} list.
         \begin{enumerate}
            \item This is the first item of an enumerated
                  list that is nested within the itemized list.

            \item This is the second item of the inner list.
                  \LaTeX\ allows you to nest lists deeper than
                  you really should.
         \end{enumerate}
         This is the rest of the second item of the outer
         list.  It is no more interesting than any other
         part of the item.
   \item This is the third item of the list.
\end{itemize}
You can even display poetry.
\begin{verse}
   There is an environment
    for verse \\             % The \\ command separates lines
   Whose features some poets % within a stanza.
   will curse.

                             % One or more blank lines separate stanzas.

   For instead of making\\
   Them do \emph{all} line breaking, \\
   It allows them to put too many words on a line when they'd rather be
   forced to be terse.
\end{verse}

Mathematical formulas may also be displayed.  A
displayed formula
is
one-line long; multiline
formulas require special formatting instructions.
   \[  \ip{\Gamma}{\psi'} = x'' + y^{2} + z_{i}^{n}\]
Don't start a paragraph with a displayed equation,
nor make one a paragraph by itself.

\end{document}               % End of document.

\section{Bench Huge Section 1547}
% synthetic bench section 1547
% This is a sample LaTeX input file.  (Version of 12 August 2004.)
%
% A '%' character causes TeX to ignore all remaining text on the line,
% and is used for comments like this one.

\documentclass{article}      % Specifies the document class

                             % The preamble begins here.
\title{An Example Document}  % Declares the document's title.
\author{Leslie Lamport}      % Declares the author's name.
\date{January 21, 1994}      % Deleting this command produces today's date.

\newcommand{\ip}[2]{(#1, #2)}
                             % Defines \ip{arg1}{arg2} to mean
                             % (arg1, arg2).

%\newcommand{\ip}[2]{\langle #1 | #2\rangle}
                             % This is an alternative definition of
                             % \ip that is commented out.

\begin{document}             % End of preamble and beginning of text.

\maketitle                   % Produces the title.

This is an example input file.  Comparing it with
the output it generates can show you how to
produce a simple document of your own.

\section{Ordinary Text}      % Produces section heading.  Lower-level
                             % sections are begun with similar
                             % \subsection and \subsubsection commands.

The ends  of words and sentences are marked
  by   spaces. It  doesn't matter how many
spaces    you type; one is as good as 100.  The
end of   a line counts as a space.

One   or more   blank lines denote the  end
of  a paragraph.

Since any number of consecutive spaces are treated
like a single one, the formatting of the input
file makes no difference to
      \LaTeX,                % The \LaTeX command generates the LaTeX logo.
but it makes a difference to you.  When you use
\LaTeX, making your input file as easy to read
as possible will be a great help as you write
your document and when you change it.  This sample
file shows how you can add comments to your own input
file.

Because printing is different from typewriting,
there are a number of things that you have to do
differently when preparing an input file than if
you were just typing the document directly.
Quotation marks like
       ``this''
have to be handled specially, as do quotes within
quotes:
       ``\,`this'            % \, separates the double and single quote.
        is what I just
        wrote, not  `that'\,''.

Dashes come in three sizes: an
       intra-word
dash, a medium dash for number ranges like
       1--2,
and a punctuation
       dash---like
this.

A sentence-ending space should be larger than the
space between words within a sentence.  You
sometimes have to type special commands in
conjunction with punctuation characters to get
this right, as in the following sentence.
       Gnats, gnus, etc.\ all  % `\ ' makes an inter-word space.
       begin with G\@.         % \@ marks end-of-sentence punctuation.
You should check the spaces after periods when
reading your output to make sure you haven't
forgotten any special cases.  Generating an
ellipsis
       \ldots\               % `\ ' is needed after `\ldots' because TeX
                             % ignores spaces after command names like \ldots
                             % made from \ + letters.
                             %
                             % Note how a `%' character causes TeX to ignore
                             % the end of the input line, so these blank lines
                             % do not start a new paragraph.
                             %
with the right spacing around the periods requires
a special command.

\LaTeX\ interprets some common characters as
commands, so you must type special commands to
generate them.  These characters include the
following:
       \$ \& \% \# \{ and \}.

In printing, text is usually emphasized with an
       \emph{italic}
type style.

\begin{em}
   A long segment of text can also be emphasized
   in this way.  Text within such a segment can be
   given \emph{additional} emphasis.
\end{em}

It is sometimes necessary to prevent \LaTeX\ from
breaking a line where it might otherwise do so.
This may be at a space, as between the ``Mr.''\ and
``Jones'' in
       ``Mr.~Jones'',        % ~ produces an unbreakable interword space.
or within a word---especially when the word is a
symbol like
       \mbox{\emph{itemnum}}
that makes little sense when hyphenated across
lines.

Footnotes\footnote{This is an example of a footnote.}
pose no problem.

\LaTeX\ is good at typesetting mathematical formulas
like
       \( x-3y + z = 7 \)
or
       \( a_{1} > x^{2n} + y^{2n} > x' \)
or
       \( \ip{A}{B} = \sum_{i} a_{i} b_{i} \).
The spaces you type in a formula are
ignored.  Remember that a letter like
       $x$                   % $ ... $  and  \( ... \)  are equivalent
is a formula when it denotes a mathematical
symbol, and it should be typed as one.

\section{Displayed Text}

Text is displayed by indenting it from the left
margin.  Quotations are commonly displayed.  There
are short quotations
\begin{quote}
   This is a short quotation.  It consists of a
   single paragraph of text.  See how it is formatted.
\end{quote}
and longer ones.
\begin{quotation}
   This is a longer quotation.  It consists of two
   paragraphs of text, neither of which are
   particularly interesting.

   This is the second paragraph of the quotation.  It
   is just as dull as the first paragraph.
\end{quotation}
Another frequently-displayed structure is a list.
The following is an example of an \emph{itemized}
list.
\begin{itemize}
   \item This is the first item of an itemized list.
         Each item in the list is marked with a ``tick''.
         You don't have to worry about what kind of tick
         mark is used.

   \item This is the second item of the list.  It
         contains another list nested inside it.  The inner
         list is an \emph{enumerated} list.
         \begin{enumerate}
            \item This is the first item of an enumerated
                  list that is nested within the itemized list.

            \item This is the second item of the inner list.
                  \LaTeX\ allows you to nest lists deeper than
                  you really should.
         \end{enumerate}
         This is the rest of the second item of the outer
         list.  It is no more interesting than any other
         part of the item.
   \item This is the third item of the list.
\end{itemize}
You can even display poetry.
\begin{verse}
   There is an environment
    for verse \\             % The \\ command separates lines
   Whose features some poets % within a stanza.
   will curse.

                             % One or more blank lines separate stanzas.

   For instead of making\\
   Them do \emph{all} line breaking, \\
   It allows them to put too many words on a line when they'd rather be
   forced to be terse.
\end{verse}

Mathematical formulas may also be displayed.  A
displayed formula
is
one-line long; multiline
formulas require special formatting instructions.
   \[  \ip{\Gamma}{\psi'} = x'' + y^{2} + z_{i}^{n}\]
Don't start a paragraph with a displayed equation,
nor make one a paragraph by itself.

\end{document}               % End of document.

\section{Bench Huge Section 1548}
% synthetic bench section 1548
% This is a sample LaTeX input file.  (Version of 12 August 2004.)
%
% A '%' character causes TeX to ignore all remaining text on the line,
% and is used for comments like this one.

\documentclass{article}      % Specifies the document class

                             % The preamble begins here.
\title{An Example Document}  % Declares the document's title.
\author{Leslie Lamport}      % Declares the author's name.
\date{January 21, 1994}      % Deleting this command produces today's date.

\newcommand{\ip}[2]{(#1, #2)}
                             % Defines \ip{arg1}{arg2} to mean
                             % (arg1, arg2).

%\newcommand{\ip}[2]{\langle #1 | #2\rangle}
                             % This is an alternative definition of
                             % \ip that is commented out.

\begin{document}             % End of preamble and beginning of text.

\maketitle                   % Produces the title.

This is an example input file.  Comparing it with
the output it generates can show you how to
produce a simple document of your own.

\section{Ordinary Text}      % Produces section heading.  Lower-level
                             % sections are begun with similar
                             % \subsection and \subsubsection commands.

The ends  of words and sentences are marked
  by   spaces. It  doesn't matter how many
spaces    you type; one is as good as 100.  The
end of   a line counts as a space.

One   or more   blank lines denote the  end
of  a paragraph.

Since any number of consecutive spaces are treated
like a single one, the formatting of the input
file makes no difference to
      \LaTeX,                % The \LaTeX command generates the LaTeX logo.
but it makes a difference to you.  When you use
\LaTeX, making your input file as easy to read
as possible will be a great help as you write
your document and when you change it.  This sample
file shows how you can add comments to your own input
file.

Because printing is different from typewriting,
there are a number of things that you have to do
differently when preparing an input file than if
you were just typing the document directly.
Quotation marks like
       ``this''
have to be handled specially, as do quotes within
quotes:
       ``\,`this'            % \, separates the double and single quote.
        is what I just
        wrote, not  `that'\,''.

Dashes come in three sizes: an
       intra-word
dash, a medium dash for number ranges like
       1--2,
and a punctuation
       dash---like
this.

A sentence-ending space should be larger than the
space between words within a sentence.  You
sometimes have to type special commands in
conjunction with punctuation characters to get
this right, as in the following sentence.
       Gnats, gnus, etc.\ all  % `\ ' makes an inter-word space.
       begin with G\@.         % \@ marks end-of-sentence punctuation.
You should check the spaces after periods when
reading your output to make sure you haven't
forgotten any special cases.  Generating an
ellipsis
       \ldots\               % `\ ' is needed after `\ldots' because TeX
                             % ignores spaces after command names like \ldots
                             % made from \ + letters.
                             %
                             % Note how a `%' character causes TeX to ignore
                             % the end of the input line, so these blank lines
                             % do not start a new paragraph.
                             %
with the right spacing around the periods requires
a special command.

\LaTeX\ interprets some common characters as
commands, so you must type special commands to
generate them.  These characters include the
following:
       \$ \& \% \# \{ and \}.

In printing, text is usually emphasized with an
       \emph{italic}
type style.

\begin{em}
   A long segment of text can also be emphasized
   in this way.  Text within such a segment can be
   given \emph{additional} emphasis.
\end{em}

It is sometimes necessary to prevent \LaTeX\ from
breaking a line where it might otherwise do so.
This may be at a space, as between the ``Mr.''\ and
``Jones'' in
       ``Mr.~Jones'',        % ~ produces an unbreakable interword space.
or within a word---especially when the word is a
symbol like
       \mbox{\emph{itemnum}}
that makes little sense when hyphenated across
lines.

Footnotes\footnote{This is an example of a footnote.}
pose no problem.

\LaTeX\ is good at typesetting mathematical formulas
like
       \( x-3y + z = 7 \)
or
       \( a_{1} > x^{2n} + y^{2n} > x' \)
or
       \( \ip{A}{B} = \sum_{i} a_{i} b_{i} \).
The spaces you type in a formula are
ignored.  Remember that a letter like
       $x$                   % $ ... $  and  \( ... \)  are equivalent
is a formula when it denotes a mathematical
symbol, and it should be typed as one.

\section{Displayed Text}

Text is displayed by indenting it from the left
margin.  Quotations are commonly displayed.  There
are short quotations
\begin{quote}
   This is a short quotation.  It consists of a
   single paragraph of text.  See how it is formatted.
\end{quote}
and longer ones.
\begin{quotation}
   This is a longer quotation.  It consists of two
   paragraphs of text, neither of which are
   particularly interesting.

   This is the second paragraph of the quotation.  It
   is just as dull as the first paragraph.
\end{quotation}
Another frequently-displayed structure is a list.
The following is an example of an \emph{itemized}
list.
\begin{itemize}
   \item This is the first item of an itemized list.
         Each item in the list is marked with a ``tick''.
         You don't have to worry about what kind of tick
         mark is used.

   \item This is the second item of the list.  It
         contains another list nested inside it.  The inner
         list is an \emph{enumerated} list.
         \begin{enumerate}
            \item This is the first item of an enumerated
                  list that is nested within the itemized list.

            \item This is the second item of the inner list.
                  \LaTeX\ allows you to nest lists deeper than
                  you really should.
         \end{enumerate}
         This is the rest of the second item of the outer
         list.  It is no more interesting than any other
         part of the item.
   \item This is the third item of the list.
\end{itemize}
You can even display poetry.
\begin{verse}
   There is an environment
    for verse \\             % The \\ command separates lines
   Whose features some poets % within a stanza.
   will curse.

                             % One or more blank lines separate stanzas.

   For instead of making\\
   Them do \emph{all} line breaking, \\
   It allows them to put too many words on a line when they'd rather be
   forced to be terse.
\end{verse}

Mathematical formulas may also be displayed.  A
displayed formula
is
one-line long; multiline
formulas require special formatting instructions.
   \[  \ip{\Gamma}{\psi'} = x'' + y^{2} + z_{i}^{n}\]
Don't start a paragraph with a displayed equation,
nor make one a paragraph by itself.

\end{document}               % End of document.

\section{Bench Huge Section 1549}
% synthetic bench section 1549
% This is a sample LaTeX input file.  (Version of 12 August 2004.)
%
% A '%' character causes TeX to ignore all remaining text on the line,
% and is used for comments like this one.

\documentclass{article}      % Specifies the document class

                             % The preamble begins here.
\title{An Example Document}  % Declares the document's title.
\author{Leslie Lamport}      % Declares the author's name.
\date{January 21, 1994}      % Deleting this command produces today's date.

\newcommand{\ip}[2]{(#1, #2)}
                             % Defines \ip{arg1}{arg2} to mean
                             % (arg1, arg2).

%\newcommand{\ip}[2]{\langle #1 | #2\rangle}
                             % This is an alternative definition of
                             % \ip that is commented out.

\begin{document}             % End of preamble and beginning of text.

\maketitle                   % Produces the title.

This is an example input file.  Comparing it with
the output it generates can show you how to
produce a simple document of your own.

\section{Ordinary Text}      % Produces section heading.  Lower-level
                             % sections are begun with similar
                             % \subsection and \subsubsection commands.

The ends  of words and sentences are marked
  by   spaces. It  doesn't matter how many
spaces    you type; one is as good as 100.  The
end of   a line counts as a space.

One   or more   blank lines denote the  end
of  a paragraph.

Since any number of consecutive spaces are treated
like a single one, the formatting of the input
file makes no difference to
      \LaTeX,                % The \LaTeX command generates the LaTeX logo.
but it makes a difference to you.  When you use
\LaTeX, making your input file as easy to read
as possible will be a great help as you write
your document and when you change it.  This sample
file shows how you can add comments to your own input
file.

Because printing is different from typewriting,
there are a number of things that you have to do
differently when preparing an input file than if
you were just typing the document directly.
Quotation marks like
       ``this''
have to be handled specially, as do quotes within
quotes:
       ``\,`this'            % \, separates the double and single quote.
        is what I just
        wrote, not  `that'\,''.

Dashes come in three sizes: an
       intra-word
dash, a medium dash for number ranges like
       1--2,
and a punctuation
       dash---like
this.

A sentence-ending space should be larger than the
space between words within a sentence.  You
sometimes have to type special commands in
conjunction with punctuation characters to get
this right, as in the following sentence.
       Gnats, gnus, etc.\ all  % `\ ' makes an inter-word space.
       begin with G\@.         % \@ marks end-of-sentence punctuation.
You should check the spaces after periods when
reading your output to make sure you haven't
forgotten any special cases.  Generating an
ellipsis
       \ldots\               % `\ ' is needed after `\ldots' because TeX
                             % ignores spaces after command names like \ldots
                             % made from \ + letters.
                             %
                             % Note how a `%' character causes TeX to ignore
                             % the end of the input line, so these blank lines
                             % do not start a new paragraph.
                             %
with the right spacing around the periods requires
a special command.

\LaTeX\ interprets some common characters as
commands, so you must type special commands to
generate them.  These characters include the
following:
       \$ \& \% \# \{ and \}.

In printing, text is usually emphasized with an
       \emph{italic}
type style.

\begin{em}
   A long segment of text can also be emphasized
   in this way.  Text within such a segment can be
   given \emph{additional} emphasis.
\end{em}

It is sometimes necessary to prevent \LaTeX\ from
breaking a line where it might otherwise do so.
This may be at a space, as between the ``Mr.''\ and
``Jones'' in
       ``Mr.~Jones'',        % ~ produces an unbreakable interword space.
or within a word---especially when the word is a
symbol like
       \mbox{\emph{itemnum}}
that makes little sense when hyphenated across
lines.

Footnotes\footnote{This is an example of a footnote.}
pose no problem.

\LaTeX\ is good at typesetting mathematical formulas
like
       \( x-3y + z = 7 \)
or
       \( a_{1} > x^{2n} + y^{2n} > x' \)
or
       \( \ip{A}{B} = \sum_{i} a_{i} b_{i} \).
The spaces you type in a formula are
ignored.  Remember that a letter like
       $x$                   % $ ... $  and  \( ... \)  are equivalent
is a formula when it denotes a mathematical
symbol, and it should be typed as one.

\section{Displayed Text}

Text is displayed by indenting it from the left
margin.  Quotations are commonly displayed.  There
are short quotations
\begin{quote}
   This is a short quotation.  It consists of a
   single paragraph of text.  See how it is formatted.
\end{quote}
and longer ones.
\begin{quotation}
   This is a longer quotation.  It consists of two
   paragraphs of text, neither of which are
   particularly interesting.

   This is the second paragraph of the quotation.  It
   is just as dull as the first paragraph.
\end{quotation}
Another frequently-displayed structure is a list.
The following is an example of an \emph{itemized}
list.
\begin{itemize}
   \item This is the first item of an itemized list.
         Each item in the list is marked with a ``tick''.
         You don't have to worry about what kind of tick
         mark is used.

   \item This is the second item of the list.  It
         contains another list nested inside it.  The inner
         list is an \emph{enumerated} list.
         \begin{enumerate}
            \item This is the first item of an enumerated
                  list that is nested within the itemized list.

            \item This is the second item of the inner list.
                  \LaTeX\ allows you to nest lists deeper than
                  you really should.
         \end{enumerate}
         This is the rest of the second item of the outer
         list.  It is no more interesting than any other
         part of the item.
   \item This is the third item of the list.
\end{itemize}
You can even display poetry.
\begin{verse}
   There is an environment
    for verse \\             % The \\ command separates lines
   Whose features some poets % within a stanza.
   will curse.

                             % One or more blank lines separate stanzas.

   For instead of making\\
   Them do \emph{all} line breaking, \\
   It allows them to put too many words on a line when they'd rather be
   forced to be terse.
\end{verse}

Mathematical formulas may also be displayed.  A
displayed formula
is
one-line long; multiline
formulas require special formatting instructions.
   \[  \ip{\Gamma}{\psi'} = x'' + y^{2} + z_{i}^{n}\]
Don't start a paragraph with a displayed equation,
nor make one a paragraph by itself.

\end{document}               % End of document.

\section{Bench Huge Section 1550}
% synthetic bench section 1550
% This is a sample LaTeX input file.  (Version of 12 August 2004.)
%
% A '%' character causes TeX to ignore all remaining text on the line,
% and is used for comments like this one.

\documentclass{article}      % Specifies the document class

                             % The preamble begins here.
\title{An Example Document}  % Declares the document's title.
\author{Leslie Lamport}      % Declares the author's name.
\date{January 21, 1994}      % Deleting this command produces today's date.

\newcommand{\ip}[2]{(#1, #2)}
                             % Defines \ip{arg1}{arg2} to mean
                             % (arg1, arg2).

%\newcommand{\ip}[2]{\langle #1 | #2\rangle}
                             % This is an alternative definition of
                             % \ip that is commented out.

\begin{document}             % End of preamble and beginning of text.

\maketitle                   % Produces the title.

This is an example input file.  Comparing it with
the output it generates can show you how to
produce a simple document of your own.

\section{Ordinary Text}      % Produces section heading.  Lower-level
                             % sections are begun with similar
                             % \subsection and \subsubsection commands.

The ends  of words and sentences are marked
  by   spaces. It  doesn't matter how many
spaces    you type; one is as good as 100.  The
end of   a line counts as a space.

One   or more   blank lines denote the  end
of  a paragraph.

Since any number of consecutive spaces are treated
like a single one, the formatting of the input
file makes no difference to
      \LaTeX,                % The \LaTeX command generates the LaTeX logo.
but it makes a difference to you.  When you use
\LaTeX, making your input file as easy to read
as possible will be a great help as you write
your document and when you change it.  This sample
file shows how you can add comments to your own input
file.

Because printing is different from typewriting,
there are a number of things that you have to do
differently when preparing an input file than if
you were just typing the document directly.
Quotation marks like
       ``this''
have to be handled specially, as do quotes within
quotes:
       ``\,`this'            % \, separates the double and single quote.
        is what I just
        wrote, not  `that'\,''.

Dashes come in three sizes: an
       intra-word
dash, a medium dash for number ranges like
       1--2,
and a punctuation
       dash---like
this.

A sentence-ending space should be larger than the
space between words within a sentence.  You
sometimes have to type special commands in
conjunction with punctuation characters to get
this right, as in the following sentence.
       Gnats, gnus, etc.\ all  % `\ ' makes an inter-word space.
       begin with G\@.         % \@ marks end-of-sentence punctuation.
You should check the spaces after periods when
reading your output to make sure you haven't
forgotten any special cases.  Generating an
ellipsis
       \ldots\               % `\ ' is needed after `\ldots' because TeX
                             % ignores spaces after command names like \ldots
                             % made from \ + letters.
                             %
                             % Note how a `%' character causes TeX to ignore
                             % the end of the input line, so these blank lines
                             % do not start a new paragraph.
                             %
with the right spacing around the periods requires
a special command.

\LaTeX\ interprets some common characters as
commands, so you must type special commands to
generate them.  These characters include the
following:
       \$ \& \% \# \{ and \}.

In printing, text is usually emphasized with an
       \emph{italic}
type style.

\begin{em}
   A long segment of text can also be emphasized
   in this way.  Text within such a segment can be
   given \emph{additional} emphasis.
\end{em}

It is sometimes necessary to prevent \LaTeX\ from
breaking a line where it might otherwise do so.
This may be at a space, as between the ``Mr.''\ and
``Jones'' in
       ``Mr.~Jones'',        % ~ produces an unbreakable interword space.
or within a word---especially when the word is a
symbol like
       \mbox{\emph{itemnum}}
that makes little sense when hyphenated across
lines.

Footnotes\footnote{This is an example of a footnote.}
pose no problem.

\LaTeX\ is good at typesetting mathematical formulas
like
       \( x-3y + z = 7 \)
or
       \( a_{1} > x^{2n} + y^{2n} > x' \)
or
       \( \ip{A}{B} = \sum_{i} a_{i} b_{i} \).
The spaces you type in a formula are
ignored.  Remember that a letter like
       $x$                   % $ ... $  and  \( ... \)  are equivalent
is a formula when it denotes a mathematical
symbol, and it should be typed as one.

\section{Displayed Text}

Text is displayed by indenting it from the left
margin.  Quotations are commonly displayed.  There
are short quotations
\begin{quote}
   This is a short quotation.  It consists of a
   single paragraph of text.  See how it is formatted.
\end{quote}
and longer ones.
\begin{quotation}
   This is a longer quotation.  It consists of two
   paragraphs of text, neither of which are
   particularly interesting.

   This is the second paragraph of the quotation.  It
   is just as dull as the first paragraph.
\end{quotation}
Another frequently-displayed structure is a list.
The following is an example of an \emph{itemized}
list.
\begin{itemize}
   \item This is the first item of an itemized list.
         Each item in the list is marked with a ``tick''.
         You don't have to worry about what kind of tick
         mark is used.

   \item This is the second item of the list.  It
         contains another list nested inside it.  The inner
         list is an \emph{enumerated} list.
         \begin{enumerate}
            \item This is the first item of an enumerated
                  list that is nested within the itemized list.

            \item This is the second item of the inner list.
                  \LaTeX\ allows you to nest lists deeper than
                  you really should.
         \end{enumerate}
         This is the rest of the second item of the outer
         list.  It is no more interesting than any other
         part of the item.
   \item This is the third item of the list.
\end{itemize}
You can even display poetry.
\begin{verse}
   There is an environment
    for verse \\             % The \\ command separates lines
   Whose features some poets % within a stanza.
   will curse.

                             % One or more blank lines separate stanzas.

   For instead of making\\
   Them do \emph{all} line breaking, \\
   It allows them to put too many words on a line when they'd rather be
   forced to be terse.
\end{verse}

Mathematical formulas may also be displayed.  A
displayed formula
is
one-line long; multiline
formulas require special formatting instructions.
   \[  \ip{\Gamma}{\psi'} = x'' + y^{2} + z_{i}^{n}\]
Don't start a paragraph with a displayed equation,
nor make one a paragraph by itself.

\end{document}               % End of document.

\section{Bench Huge Section 1551}
% synthetic bench section 1551
% This is a sample LaTeX input file.  (Version of 12 August 2004.)
%
% A '%' character causes TeX to ignore all remaining text on the line,
% and is used for comments like this one.

\documentclass{article}      % Specifies the document class

                             % The preamble begins here.
\title{An Example Document}  % Declares the document's title.
\author{Leslie Lamport}      % Declares the author's name.
\date{January 21, 1994}      % Deleting this command produces today's date.

\newcommand{\ip}[2]{(#1, #2)}
                             % Defines \ip{arg1}{arg2} to mean
                             % (arg1, arg2).

%\newcommand{\ip}[2]{\langle #1 | #2\rangle}
                             % This is an alternative definition of
                             % \ip that is commented out.

\begin{document}             % End of preamble and beginning of text.

\maketitle                   % Produces the title.

This is an example input file.  Comparing it with
the output it generates can show you how to
produce a simple document of your own.

\section{Ordinary Text}      % Produces section heading.  Lower-level
                             % sections are begun with similar
                             % \subsection and \subsubsection commands.

The ends  of words and sentences are marked
  by   spaces. It  doesn't matter how many
spaces    you type; one is as good as 100.  The
end of   a line counts as a space.

One   or more   blank lines denote the  end
of  a paragraph.

Since any number of consecutive spaces are treated
like a single one, the formatting of the input
file makes no difference to
      \LaTeX,                % The \LaTeX command generates the LaTeX logo.
but it makes a difference to you.  When you use
\LaTeX, making your input file as easy to read
as possible will be a great help as you write
your document and when you change it.  This sample
file shows how you can add comments to your own input
file.

Because printing is different from typewriting,
there are a number of things that you have to do
differently when preparing an input file than if
you were just typing the document directly.
Quotation marks like
       ``this''
have to be handled specially, as do quotes within
quotes:
       ``\,`this'            % \, separates the double and single quote.
        is what I just
        wrote, not  `that'\,''.

Dashes come in three sizes: an
       intra-word
dash, a medium dash for number ranges like
       1--2,
and a punctuation
       dash---like
this.

A sentence-ending space should be larger than the
space between words within a sentence.  You
sometimes have to type special commands in
conjunction with punctuation characters to get
this right, as in the following sentence.
       Gnats, gnus, etc.\ all  % `\ ' makes an inter-word space.
       begin with G\@.         % \@ marks end-of-sentence punctuation.
You should check the spaces after periods when
reading your output to make sure you haven't
forgotten any special cases.  Generating an
ellipsis
       \ldots\               % `\ ' is needed after `\ldots' because TeX
                             % ignores spaces after command names like \ldots
                             % made from \ + letters.
                             %
                             % Note how a `%' character causes TeX to ignore
                             % the end of the input line, so these blank lines
                             % do not start a new paragraph.
                             %
with the right spacing around the periods requires
a special command.

\LaTeX\ interprets some common characters as
commands, so you must type special commands to
generate them.  These characters include the
following:
       \$ \& \% \# \{ and \}.

In printing, text is usually emphasized with an
       \emph{italic}
type style.

\begin{em}
   A long segment of text can also be emphasized
   in this way.  Text within such a segment can be
   given \emph{additional} emphasis.
\end{em}

It is sometimes necessary to prevent \LaTeX\ from
breaking a line where it might otherwise do so.
This may be at a space, as between the ``Mr.''\ and
``Jones'' in
       ``Mr.~Jones'',        % ~ produces an unbreakable interword space.
or within a word---especially when the word is a
symbol like
       \mbox{\emph{itemnum}}
that makes little sense when hyphenated across
lines.

Footnotes\footnote{This is an example of a footnote.}
pose no problem.

\LaTeX\ is good at typesetting mathematical formulas
like
       \( x-3y + z = 7 \)
or
       \( a_{1} > x^{2n} + y^{2n} > x' \)
or
       \( \ip{A}{B} = \sum_{i} a_{i} b_{i} \).
The spaces you type in a formula are
ignored.  Remember that a letter like
       $x$                   % $ ... $  and  \( ... \)  are equivalent
is a formula when it denotes a mathematical
symbol, and it should be typed as one.

\section{Displayed Text}

Text is displayed by indenting it from the left
margin.  Quotations are commonly displayed.  There
are short quotations
\begin{quote}
   This is a short quotation.  It consists of a
   single paragraph of text.  See how it is formatted.
\end{quote}
and longer ones.
\begin{quotation}
   This is a longer quotation.  It consists of two
   paragraphs of text, neither of which are
   particularly interesting.

   This is the second paragraph of the quotation.  It
   is just as dull as the first paragraph.
\end{quotation}
Another frequently-displayed structure is a list.
The following is an example of an \emph{itemized}
list.
\begin{itemize}
   \item This is the first item of an itemized list.
         Each item in the list is marked with a ``tick''.
         You don't have to worry about what kind of tick
         mark is used.

   \item This is the second item of the list.  It
         contains another list nested inside it.  The inner
         list is an \emph{enumerated} list.
         \begin{enumerate}
            \item This is the first item of an enumerated
                  list that is nested within the itemized list.

            \item This is the second item of the inner list.
                  \LaTeX\ allows you to nest lists deeper than
                  you really should.
         \end{enumerate}
         This is the rest of the second item of the outer
         list.  It is no more interesting than any other
         part of the item.
   \item This is the third item of the list.
\end{itemize}
You can even display poetry.
\begin{verse}
   There is an environment
    for verse \\             % The \\ command separates lines
   Whose features some poets % within a stanza.
   will curse.

                             % One or more blank lines separate stanzas.

   For instead of making\\
   Them do \emph{all} line breaking, \\
   It allows them to put too many words on a line when they'd rather be
   forced to be terse.
\end{verse}

Mathematical formulas may also be displayed.  A
displayed formula
is
one-line long; multiline
formulas require special formatting instructions.
   \[  \ip{\Gamma}{\psi'} = x'' + y^{2} + z_{i}^{n}\]
Don't start a paragraph with a displayed equation,
nor make one a paragraph by itself.

\end{document}               % End of document.

\section{Bench Huge Section 1552}
% synthetic bench section 1552
% This is a sample LaTeX input file.  (Version of 12 August 2004.)
%
% A '%' character causes TeX to ignore all remaining text on the line,
% and is used for comments like this one.

\documentclass{article}      % Specifies the document class

                             % The preamble begins here.
\title{An Example Document}  % Declares the document's title.
\author{Leslie Lamport}      % Declares the author's name.
\date{January 21, 1994}      % Deleting this command produces today's date.

\newcommand{\ip}[2]{(#1, #2)}
                             % Defines \ip{arg1}{arg2} to mean
                             % (arg1, arg2).

%\newcommand{\ip}[2]{\langle #1 | #2\rangle}
                             % This is an alternative definition of
                             % \ip that is commented out.

\begin{document}             % End of preamble and beginning of text.

\maketitle                   % Produces the title.

This is an example input file.  Comparing it with
the output it generates can show you how to
produce a simple document of your own.

\section{Ordinary Text}      % Produces section heading.  Lower-level
                             % sections are begun with similar
                             % \subsection and \subsubsection commands.

The ends  of words and sentences are marked
  by   spaces. It  doesn't matter how many
spaces    you type; one is as good as 100.  The
end of   a line counts as a space.

One   or more   blank lines denote the  end
of  a paragraph.

Since any number of consecutive spaces are treated
like a single one, the formatting of the input
file makes no difference to
      \LaTeX,                % The \LaTeX command generates the LaTeX logo.
but it makes a difference to you.  When you use
\LaTeX, making your input file as easy to read
as possible will be a great help as you write
your document and when you change it.  This sample
file shows how you can add comments to your own input
file.

Because printing is different from typewriting,
there are a number of things that you have to do
differently when preparing an input file than if
you were just typing the document directly.
Quotation marks like
       ``this''
have to be handled specially, as do quotes within
quotes:
       ``\,`this'            % \, separates the double and single quote.
        is what I just
        wrote, not  `that'\,''.

Dashes come in three sizes: an
       intra-word
dash, a medium dash for number ranges like
       1--2,
and a punctuation
       dash---like
this.

A sentence-ending space should be larger than the
space between words within a sentence.  You
sometimes have to type special commands in
conjunction with punctuation characters to get
this right, as in the following sentence.
       Gnats, gnus, etc.\ all  % `\ ' makes an inter-word space.
       begin with G\@.         % \@ marks end-of-sentence punctuation.
You should check the spaces after periods when
reading your output to make sure you haven't
forgotten any special cases.  Generating an
ellipsis
       \ldots\               % `\ ' is needed after `\ldots' because TeX
                             % ignores spaces after command names like \ldots
                             % made from \ + letters.
                             %
                             % Note how a `%' character causes TeX to ignore
                             % the end of the input line, so these blank lines
                             % do not start a new paragraph.
                             %
with the right spacing around the periods requires
a special command.

\LaTeX\ interprets some common characters as
commands, so you must type special commands to
generate them.  These characters include the
following:
       \$ \& \% \# \{ and \}.

In printing, text is usually emphasized with an
       \emph{italic}
type style.

\begin{em}
   A long segment of text can also be emphasized
   in this way.  Text within such a segment can be
   given \emph{additional} emphasis.
\end{em}

It is sometimes necessary to prevent \LaTeX\ from
breaking a line where it might otherwise do so.
This may be at a space, as between the ``Mr.''\ and
``Jones'' in
       ``Mr.~Jones'',        % ~ produces an unbreakable interword space.
or within a word---especially when the word is a
symbol like
       \mbox{\emph{itemnum}}
that makes little sense when hyphenated across
lines.

Footnotes\footnote{This is an example of a footnote.}
pose no problem.

\LaTeX\ is good at typesetting mathematical formulas
like
       \( x-3y + z = 7 \)
or
       \( a_{1} > x^{2n} + y^{2n} > x' \)
or
       \( \ip{A}{B} = \sum_{i} a_{i} b_{i} \).
The spaces you type in a formula are
ignored.  Remember that a letter like
       $x$                   % $ ... $  and  \( ... \)  are equivalent
is a formula when it denotes a mathematical
symbol, and it should be typed as one.

\section{Displayed Text}

Text is displayed by indenting it from the left
margin.  Quotations are commonly displayed.  There
are short quotations
\begin{quote}
   This is a short quotation.  It consists of a
   single paragraph of text.  See how it is formatted.
\end{quote}
and longer ones.
\begin{quotation}
   This is a longer quotation.  It consists of two
   paragraphs of text, neither of which are
   particularly interesting.

   This is the second paragraph of the quotation.  It
   is just as dull as the first paragraph.
\end{quotation}
Another frequently-displayed structure is a list.
The following is an example of an \emph{itemized}
list.
\begin{itemize}
   \item This is the first item of an itemized list.
         Each item in the list is marked with a ``tick''.
         You don't have to worry about what kind of tick
         mark is used.

   \item This is the second item of the list.  It
         contains another list nested inside it.  The inner
         list is an \emph{enumerated} list.
         \begin{enumerate}
            \item This is the first item of an enumerated
                  list that is nested within the itemized list.

            \item This is the second item of the inner list.
                  \LaTeX\ allows you to nest lists deeper than
                  you really should.
         \end{enumerate}
         This is the rest of the second item of the outer
         list.  It is no more interesting than any other
         part of the item.
   \item This is the third item of the list.
\end{itemize}
You can even display poetry.
\begin{verse}
   There is an environment
    for verse \\             % The \\ command separates lines
   Whose features some poets % within a stanza.
   will curse.

                             % One or more blank lines separate stanzas.

   For instead of making\\
   Them do \emph{all} line breaking, \\
   It allows them to put too many words on a line when they'd rather be
   forced to be terse.
\end{verse}

Mathematical formulas may also be displayed.  A
displayed formula
is
one-line long; multiline
formulas require special formatting instructions.
   \[  \ip{\Gamma}{\psi'} = x'' + y^{2} + z_{i}^{n}\]
Don't start a paragraph with a displayed equation,
nor make one a paragraph by itself.

\end{document}               % End of document.

\section{Bench Huge Section 1553}
% synthetic bench section 1553
% This is a sample LaTeX input file.  (Version of 12 August 2004.)
%
% A '%' character causes TeX to ignore all remaining text on the line,
% and is used for comments like this one.

\documentclass{article}      % Specifies the document class

                             % The preamble begins here.
\title{An Example Document}  % Declares the document's title.
\author{Leslie Lamport}      % Declares the author's name.
\date{January 21, 1994}      % Deleting this command produces today's date.

\newcommand{\ip}[2]{(#1, #2)}
                             % Defines \ip{arg1}{arg2} to mean
                             % (arg1, arg2).

%\newcommand{\ip}[2]{\langle #1 | #2\rangle}
                             % This is an alternative definition of
                             % \ip that is commented out.

\begin{document}             % End of preamble and beginning of text.

\maketitle                   % Produces the title.

This is an example input file.  Comparing it with
the output it generates can show you how to
produce a simple document of your own.

\section{Ordinary Text}      % Produces section heading.  Lower-level
                             % sections are begun with similar
                             % \subsection and \subsubsection commands.

The ends  of words and sentences are marked
  by   spaces. It  doesn't matter how many
spaces    you type; one is as good as 100.  The
end of   a line counts as a space.

One   or more   blank lines denote the  end
of  a paragraph.

Since any number of consecutive spaces are treated
like a single one, the formatting of the input
file makes no difference to
      \LaTeX,                % The \LaTeX command generates the LaTeX logo.
but it makes a difference to you.  When you use
\LaTeX, making your input file as easy to read
as possible will be a great help as you write
your document and when you change it.  This sample
file shows how you can add comments to your own input
file.

Because printing is different from typewriting,
there are a number of things that you have to do
differently when preparing an input file than if
you were just typing the document directly.
Quotation marks like
       ``this''
have to be handled specially, as do quotes within
quotes:
       ``\,`this'            % \, separates the double and single quote.
        is what I just
        wrote, not  `that'\,''.

Dashes come in three sizes: an
       intra-word
dash, a medium dash for number ranges like
       1--2,
and a punctuation
       dash---like
this.

A sentence-ending space should be larger than the
space between words within a sentence.  You
sometimes have to type special commands in
conjunction with punctuation characters to get
this right, as in the following sentence.
       Gnats, gnus, etc.\ all  % `\ ' makes an inter-word space.
       begin with G\@.         % \@ marks end-of-sentence punctuation.
You should check the spaces after periods when
reading your output to make sure you haven't
forgotten any special cases.  Generating an
ellipsis
       \ldots\               % `\ ' is needed after `\ldots' because TeX
                             % ignores spaces after command names like \ldots
                             % made from \ + letters.
                             %
                             % Note how a `%' character causes TeX to ignore
                             % the end of the input line, so these blank lines
                             % do not start a new paragraph.
                             %
with the right spacing around the periods requires
a special command.

\LaTeX\ interprets some common characters as
commands, so you must type special commands to
generate them.  These characters include the
following:
       \$ \& \% \# \{ and \}.

In printing, text is usually emphasized with an
       \emph{italic}
type style.

\begin{em}
   A long segment of text can also be emphasized
   in this way.  Text within such a segment can be
   given \emph{additional} emphasis.
\end{em}

It is sometimes necessary to prevent \LaTeX\ from
breaking a line where it might otherwise do so.
This may be at a space, as between the ``Mr.''\ and
``Jones'' in
       ``Mr.~Jones'',        % ~ produces an unbreakable interword space.
or within a word---especially when the word is a
symbol like
       \mbox{\emph{itemnum}}
that makes little sense when hyphenated across
lines.

Footnotes\footnote{This is an example of a footnote.}
pose no problem.

\LaTeX\ is good at typesetting mathematical formulas
like
       \( x-3y + z = 7 \)
or
       \( a_{1} > x^{2n} + y^{2n} > x' \)
or
       \( \ip{A}{B} = \sum_{i} a_{i} b_{i} \).
The spaces you type in a formula are
ignored.  Remember that a letter like
       $x$                   % $ ... $  and  \( ... \)  are equivalent
is a formula when it denotes a mathematical
symbol, and it should be typed as one.

\section{Displayed Text}

Text is displayed by indenting it from the left
margin.  Quotations are commonly displayed.  There
are short quotations
\begin{quote}
   This is a short quotation.  It consists of a
   single paragraph of text.  See how it is formatted.
\end{quote}
and longer ones.
\begin{quotation}
   This is a longer quotation.  It consists of two
   paragraphs of text, neither of which are
   particularly interesting.

   This is the second paragraph of the quotation.  It
   is just as dull as the first paragraph.
\end{quotation}
Another frequently-displayed structure is a list.
The following is an example of an \emph{itemized}
list.
\begin{itemize}
   \item This is the first item of an itemized list.
         Each item in the list is marked with a ``tick''.
         You don't have to worry about what kind of tick
         mark is used.

   \item This is the second item of the list.  It
         contains another list nested inside it.  The inner
         list is an \emph{enumerated} list.
         \begin{enumerate}
            \item This is the first item of an enumerated
                  list that is nested within the itemized list.

            \item This is the second item of the inner list.
                  \LaTeX\ allows you to nest lists deeper than
                  you really should.
         \end{enumerate}
         This is the rest of the second item of the outer
         list.  It is no more interesting than any other
         part of the item.
   \item This is the third item of the list.
\end{itemize}
You can even display poetry.
\begin{verse}
   There is an environment
    for verse \\             % The \\ command separates lines
   Whose features some poets % within a stanza.
   will curse.

                             % One or more blank lines separate stanzas.

   For instead of making\\
   Them do \emph{all} line breaking, \\
   It allows them to put too many words on a line when they'd rather be
   forced to be terse.
\end{verse}

Mathematical formulas may also be displayed.  A
displayed formula
is
one-line long; multiline
formulas require special formatting instructions.
   \[  \ip{\Gamma}{\psi'} = x'' + y^{2} + z_{i}^{n}\]
Don't start a paragraph with a displayed equation,
nor make one a paragraph by itself.

\end{document}               % End of document.

\section{Bench Huge Section 1554}
% synthetic bench section 1554
% This is a sample LaTeX input file.  (Version of 12 August 2004.)
%
% A '%' character causes TeX to ignore all remaining text on the line,
% and is used for comments like this one.

\documentclass{article}      % Specifies the document class

                             % The preamble begins here.
\title{An Example Document}  % Declares the document's title.
\author{Leslie Lamport}      % Declares the author's name.
\date{January 21, 1994}      % Deleting this command produces today's date.

\newcommand{\ip}[2]{(#1, #2)}
                             % Defines \ip{arg1}{arg2} to mean
                             % (arg1, arg2).

%\newcommand{\ip}[2]{\langle #1 | #2\rangle}
                             % This is an alternative definition of
                             % \ip that is commented out.

\begin{document}             % End of preamble and beginning of text.

\maketitle                   % Produces the title.

This is an example input file.  Comparing it with
the output it generates can show you how to
produce a simple document of your own.

\section{Ordinary Text}      % Produces section heading.  Lower-level
                             % sections are begun with similar
                             % \subsection and \subsubsection commands.

The ends  of words and sentences are marked
  by   spaces. It  doesn't matter how many
spaces    you type; one is as good as 100.  The
end of   a line counts as a space.

One   or more   blank lines denote the  end
of  a paragraph.

Since any number of consecutive spaces are treated
like a single one, the formatting of the input
file makes no difference to
      \LaTeX,                % The \LaTeX command generates the LaTeX logo.
but it makes a difference to you.  When you use
\LaTeX, making your input file as easy to read
as possible will be a great help as you write
your document and when you change it.  This sample
file shows how you can add comments to your own input
file.

Because printing is different from typewriting,
there are a number of things that you have to do
differently when preparing an input file than if
you were just typing the document directly.
Quotation marks like
       ``this''
have to be handled specially, as do quotes within
quotes:
       ``\,`this'            % \, separates the double and single quote.
        is what I just
        wrote, not  `that'\,''.

Dashes come in three sizes: an
       intra-word
dash, a medium dash for number ranges like
       1--2,
and a punctuation
       dash---like
this.

A sentence-ending space should be larger than the
space between words within a sentence.  You
sometimes have to type special commands in
conjunction with punctuation characters to get
this right, as in the following sentence.
       Gnats, gnus, etc.\ all  % `\ ' makes an inter-word space.
       begin with G\@.         % \@ marks end-of-sentence punctuation.
You should check the spaces after periods when
reading your output to make sure you haven't
forgotten any special cases.  Generating an
ellipsis
       \ldots\               % `\ ' is needed after `\ldots' because TeX
                             % ignores spaces after command names like \ldots
                             % made from \ + letters.
                             %
                             % Note how a `%' character causes TeX to ignore
                             % the end of the input line, so these blank lines
                             % do not start a new paragraph.
                             %
with the right spacing around the periods requires
a special command.

\LaTeX\ interprets some common characters as
commands, so you must type special commands to
generate them.  These characters include the
following:
       \$ \& \% \# \{ and \}.

In printing, text is usually emphasized with an
       \emph{italic}
type style.

\begin{em}
   A long segment of text can also be emphasized
   in this way.  Text within such a segment can be
   given \emph{additional} emphasis.
\end{em}

It is sometimes necessary to prevent \LaTeX\ from
breaking a line where it might otherwise do so.
This may be at a space, as between the ``Mr.''\ and
``Jones'' in
       ``Mr.~Jones'',        % ~ produces an unbreakable interword space.
or within a word---especially when the word is a
symbol like
       \mbox{\emph{itemnum}}
that makes little sense when hyphenated across
lines.

Footnotes\footnote{This is an example of a footnote.}
pose no problem.

\LaTeX\ is good at typesetting mathematical formulas
like
       \( x-3y + z = 7 \)
or
       \( a_{1} > x^{2n} + y^{2n} > x' \)
or
       \( \ip{A}{B} = \sum_{i} a_{i} b_{i} \).
The spaces you type in a formula are
ignored.  Remember that a letter like
       $x$                   % $ ... $  and  \( ... \)  are equivalent
is a formula when it denotes a mathematical
symbol, and it should be typed as one.

\section{Displayed Text}

Text is displayed by indenting it from the left
margin.  Quotations are commonly displayed.  There
are short quotations
\begin{quote}
   This is a short quotation.  It consists of a
   single paragraph of text.  See how it is formatted.
\end{quote}
and longer ones.
\begin{quotation}
   This is a longer quotation.  It consists of two
   paragraphs of text, neither of which are
   particularly interesting.

   This is the second paragraph of the quotation.  It
   is just as dull as the first paragraph.
\end{quotation}
Another frequently-displayed structure is a list.
The following is an example of an \emph{itemized}
list.
\begin{itemize}
   \item This is the first item of an itemized list.
         Each item in the list is marked with a ``tick''.
         You don't have to worry about what kind of tick
         mark is used.

   \item This is the second item of the list.  It
         contains another list nested inside it.  The inner
         list is an \emph{enumerated} list.
         \begin{enumerate}
            \item This is the first item of an enumerated
                  list that is nested within the itemized list.

            \item This is the second item of the inner list.
                  \LaTeX\ allows you to nest lists deeper than
                  you really should.
         \end{enumerate}
         This is the rest of the second item of the outer
         list.  It is no more interesting than any other
         part of the item.
   \item This is the third item of the list.
\end{itemize}
You can even display poetry.
\begin{verse}
   There is an environment
    for verse \\             % The \\ command separates lines
   Whose features some poets % within a stanza.
   will curse.

                             % One or more blank lines separate stanzas.

   For instead of making\\
   Them do \emph{all} line breaking, \\
   It allows them to put too many words on a line when they'd rather be
   forced to be terse.
\end{verse}

Mathematical formulas may also be displayed.  A
displayed formula
is
one-line long; multiline
formulas require special formatting instructions.
   \[  \ip{\Gamma}{\psi'} = x'' + y^{2} + z_{i}^{n}\]
Don't start a paragraph with a displayed equation,
nor make one a paragraph by itself.

\end{document}               % End of document.

\section{Bench Huge Section 1555}
% synthetic bench section 1555
% This is a sample LaTeX input file.  (Version of 12 August 2004.)
%
% A '%' character causes TeX to ignore all remaining text on the line,
% and is used for comments like this one.

\documentclass{article}      % Specifies the document class

                             % The preamble begins here.
\title{An Example Document}  % Declares the document's title.
\author{Leslie Lamport}      % Declares the author's name.
\date{January 21, 1994}      % Deleting this command produces today's date.

\newcommand{\ip}[2]{(#1, #2)}
                             % Defines \ip{arg1}{arg2} to mean
                             % (arg1, arg2).

%\newcommand{\ip}[2]{\langle #1 | #2\rangle}
                             % This is an alternative definition of
                             % \ip that is commented out.

\begin{document}             % End of preamble and beginning of text.

\maketitle                   % Produces the title.

This is an example input file.  Comparing it with
the output it generates can show you how to
produce a simple document of your own.

\section{Ordinary Text}      % Produces section heading.  Lower-level
                             % sections are begun with similar
                             % \subsection and \subsubsection commands.

The ends  of words and sentences are marked
  by   spaces. It  doesn't matter how many
spaces    you type; one is as good as 100.  The
end of   a line counts as a space.

One   or more   blank lines denote the  end
of  a paragraph.

Since any number of consecutive spaces are treated
like a single one, the formatting of the input
file makes no difference to
      \LaTeX,                % The \LaTeX command generates the LaTeX logo.
but it makes a difference to you.  When you use
\LaTeX, making your input file as easy to read
as possible will be a great help as you write
your document and when you change it.  This sample
file shows how you can add comments to your own input
file.

Because printing is different from typewriting,
there are a number of things that you have to do
differently when preparing an input file than if
you were just typing the document directly.
Quotation marks like
       ``this''
have to be handled specially, as do quotes within
quotes:
       ``\,`this'            % \, separates the double and single quote.
        is what I just
        wrote, not  `that'\,''.

Dashes come in three sizes: an
       intra-word
dash, a medium dash for number ranges like
       1--2,
and a punctuation
       dash---like
this.

A sentence-ending space should be larger than the
space between words within a sentence.  You
sometimes have to type special commands in
conjunction with punctuation characters to get
this right, as in the following sentence.
       Gnats, gnus, etc.\ all  % `\ ' makes an inter-word space.
       begin with G\@.         % \@ marks end-of-sentence punctuation.
You should check the spaces after periods when
reading your output to make sure you haven't
forgotten any special cases.  Generating an
ellipsis
       \ldots\               % `\ ' is needed after `\ldots' because TeX
                             % ignores spaces after command names like \ldots
                             % made from \ + letters.
                             %
                             % Note how a `%' character causes TeX to ignore
                             % the end of the input line, so these blank lines
                             % do not start a new paragraph.
                             %
with the right spacing around the periods requires
a special command.

\LaTeX\ interprets some common characters as
commands, so you must type special commands to
generate them.  These characters include the
following:
       \$ \& \% \# \{ and \}.

In printing, text is usually emphasized with an
       \emph{italic}
type style.

\begin{em}
   A long segment of text can also be emphasized
   in this way.  Text within such a segment can be
   given \emph{additional} emphasis.
\end{em}

It is sometimes necessary to prevent \LaTeX\ from
breaking a line where it might otherwise do so.
This may be at a space, as between the ``Mr.''\ and
``Jones'' in
       ``Mr.~Jones'',        % ~ produces an unbreakable interword space.
or within a word---especially when the word is a
symbol like
       \mbox{\emph{itemnum}}
that makes little sense when hyphenated across
lines.

Footnotes\footnote{This is an example of a footnote.}
pose no problem.

\LaTeX\ is good at typesetting mathematical formulas
like
       \( x-3y + z = 7 \)
or
       \( a_{1} > x^{2n} + y^{2n} > x' \)
or
       \( \ip{A}{B} = \sum_{i} a_{i} b_{i} \).
The spaces you type in a formula are
ignored.  Remember that a letter like
       $x$                   % $ ... $  and  \( ... \)  are equivalent
is a formula when it denotes a mathematical
symbol, and it should be typed as one.

\section{Displayed Text}

Text is displayed by indenting it from the left
margin.  Quotations are commonly displayed.  There
are short quotations
\begin{quote}
   This is a short quotation.  It consists of a
   single paragraph of text.  See how it is formatted.
\end{quote}
and longer ones.
\begin{quotation}
   This is a longer quotation.  It consists of two
   paragraphs of text, neither of which are
   particularly interesting.

   This is the second paragraph of the quotation.  It
   is just as dull as the first paragraph.
\end{quotation}
Another frequently-displayed structure is a list.
The following is an example of an \emph{itemized}
list.
\begin{itemize}
   \item This is the first item of an itemized list.
         Each item in the list is marked with a ``tick''.
         You don't have to worry about what kind of tick
         mark is used.

   \item This is the second item of the list.  It
         contains another list nested inside it.  The inner
         list is an \emph{enumerated} list.
         \begin{enumerate}
            \item This is the first item of an enumerated
                  list that is nested within the itemized list.

            \item This is the second item of the inner list.
                  \LaTeX\ allows you to nest lists deeper than
                  you really should.
         \end{enumerate}
         This is the rest of the second item of the outer
         list.  It is no more interesting than any other
         part of the item.
   \item This is the third item of the list.
\end{itemize}
You can even display poetry.
\begin{verse}
   There is an environment
    for verse \\             % The \\ command separates lines
   Whose features some poets % within a stanza.
   will curse.

                             % One or more blank lines separate stanzas.

   For instead of making\\
   Them do \emph{all} line breaking, \\
   It allows them to put too many words on a line when they'd rather be
   forced to be terse.
\end{verse}

Mathematical formulas may also be displayed.  A
displayed formula
is
one-line long; multiline
formulas require special formatting instructions.
   \[  \ip{\Gamma}{\psi'} = x'' + y^{2} + z_{i}^{n}\]
Don't start a paragraph with a displayed equation,
nor make one a paragraph by itself.

\end{document}               % End of document.

\section{Bench Huge Section 1556}
% synthetic bench section 1556
% This is a sample LaTeX input file.  (Version of 12 August 2004.)
%
% A '%' character causes TeX to ignore all remaining text on the line,
% and is used for comments like this one.

\documentclass{article}      % Specifies the document class

                             % The preamble begins here.
\title{An Example Document}  % Declares the document's title.
\author{Leslie Lamport}      % Declares the author's name.
\date{January 21, 1994}      % Deleting this command produces today's date.

\newcommand{\ip}[2]{(#1, #2)}
                             % Defines \ip{arg1}{arg2} to mean
                             % (arg1, arg2).

%\newcommand{\ip}[2]{\langle #1 | #2\rangle}
                             % This is an alternative definition of
                             % \ip that is commented out.

\begin{document}             % End of preamble and beginning of text.

\maketitle                   % Produces the title.

This is an example input file.  Comparing it with
the output it generates can show you how to
produce a simple document of your own.

\section{Ordinary Text}      % Produces section heading.  Lower-level
                             % sections are begun with similar
                             % \subsection and \subsubsection commands.

The ends  of words and sentences are marked
  by   spaces. It  doesn't matter how many
spaces    you type; one is as good as 100.  The
end of   a line counts as a space.

One   or more   blank lines denote the  end
of  a paragraph.

Since any number of consecutive spaces are treated
like a single one, the formatting of the input
file makes no difference to
      \LaTeX,                % The \LaTeX command generates the LaTeX logo.
but it makes a difference to you.  When you use
\LaTeX, making your input file as easy to read
as possible will be a great help as you write
your document and when you change it.  This sample
file shows how you can add comments to your own input
file.

Because printing is different from typewriting,
there are a number of things that you have to do
differently when preparing an input file than if
you were just typing the document directly.
Quotation marks like
       ``this''
have to be handled specially, as do quotes within
quotes:
       ``\,`this'            % \, separates the double and single quote.
        is what I just
        wrote, not  `that'\,''.

Dashes come in three sizes: an
       intra-word
dash, a medium dash for number ranges like
       1--2,
and a punctuation
       dash---like
this.

A sentence-ending space should be larger than the
space between words within a sentence.  You
sometimes have to type special commands in
conjunction with punctuation characters to get
this right, as in the following sentence.
       Gnats, gnus, etc.\ all  % `\ ' makes an inter-word space.
       begin with G\@.         % \@ marks end-of-sentence punctuation.
You should check the spaces after periods when
reading your output to make sure you haven't
forgotten any special cases.  Generating an
ellipsis
       \ldots\               % `\ ' is needed after `\ldots' because TeX
                             % ignores spaces after command names like \ldots
                             % made from \ + letters.
                             %
                             % Note how a `%' character causes TeX to ignore
                             % the end of the input line, so these blank lines
                             % do not start a new paragraph.
                             %
with the right spacing around the periods requires
a special command.

\LaTeX\ interprets some common characters as
commands, so you must type special commands to
generate them.  These characters include the
following:
       \$ \& \% \# \{ and \}.

In printing, text is usually emphasized with an
       \emph{italic}
type style.

\begin{em}
   A long segment of text can also be emphasized
   in this way.  Text within such a segment can be
   given \emph{additional} emphasis.
\end{em}

It is sometimes necessary to prevent \LaTeX\ from
breaking a line where it might otherwise do so.
This may be at a space, as between the ``Mr.''\ and
``Jones'' in
       ``Mr.~Jones'',        % ~ produces an unbreakable interword space.
or within a word---especially when the word is a
symbol like
       \mbox{\emph{itemnum}}
that makes little sense when hyphenated across
lines.

Footnotes\footnote{This is an example of a footnote.}
pose no problem.

\LaTeX\ is good at typesetting mathematical formulas
like
       \( x-3y + z = 7 \)
or
       \( a_{1} > x^{2n} + y^{2n} > x' \)
or
       \( \ip{A}{B} = \sum_{i} a_{i} b_{i} \).
The spaces you type in a formula are
ignored.  Remember that a letter like
       $x$                   % $ ... $  and  \( ... \)  are equivalent
is a formula when it denotes a mathematical
symbol, and it should be typed as one.

\section{Displayed Text}

Text is displayed by indenting it from the left
margin.  Quotations are commonly displayed.  There
are short quotations
\begin{quote}
   This is a short quotation.  It consists of a
   single paragraph of text.  See how it is formatted.
\end{quote}
and longer ones.
\begin{quotation}
   This is a longer quotation.  It consists of two
   paragraphs of text, neither of which are
   particularly interesting.

   This is the second paragraph of the quotation.  It
   is just as dull as the first paragraph.
\end{quotation}
Another frequently-displayed structure is a list.
The following is an example of an \emph{itemized}
list.
\begin{itemize}
   \item This is the first item of an itemized list.
         Each item in the list is marked with a ``tick''.
         You don't have to worry about what kind of tick
         mark is used.

   \item This is the second item of the list.  It
         contains another list nested inside it.  The inner
         list is an \emph{enumerated} list.
         \begin{enumerate}
            \item This is the first item of an enumerated
                  list that is nested within the itemized list.

            \item This is the second item of the inner list.
                  \LaTeX\ allows you to nest lists deeper than
                  you really should.
         \end{enumerate}
         This is the rest of the second item of the outer
         list.  It is no more interesting than any other
         part of the item.
   \item This is the third item of the list.
\end{itemize}
You can even display poetry.
\begin{verse}
   There is an environment
    for verse \\             % The \\ command separates lines
   Whose features some poets % within a stanza.
   will curse.

                             % One or more blank lines separate stanzas.

   For instead of making\\
   Them do \emph{all} line breaking, \\
   It allows them to put too many words on a line when they'd rather be
   forced to be terse.
\end{verse}

Mathematical formulas may also be displayed.  A
displayed formula
is
one-line long; multiline
formulas require special formatting instructions.
   \[  \ip{\Gamma}{\psi'} = x'' + y^{2} + z_{i}^{n}\]
Don't start a paragraph with a displayed equation,
nor make one a paragraph by itself.

\end{document}               % End of document.

\section{Bench Huge Section 1557}
% synthetic bench section 1557
% This is a sample LaTeX input file.  (Version of 12 August 2004.)
%
% A '%' character causes TeX to ignore all remaining text on the line,
% and is used for comments like this one.

\documentclass{article}      % Specifies the document class

                             % The preamble begins here.
\title{An Example Document}  % Declares the document's title.
\author{Leslie Lamport}      % Declares the author's name.
\date{January 21, 1994}      % Deleting this command produces today's date.

\newcommand{\ip}[2]{(#1, #2)}
                             % Defines \ip{arg1}{arg2} to mean
                             % (arg1, arg2).

%\newcommand{\ip}[2]{\langle #1 | #2\rangle}
                             % This is an alternative definition of
                             % \ip that is commented out.

\begin{document}             % End of preamble and beginning of text.

\maketitle                   % Produces the title.

This is an example input file.  Comparing it with
the output it generates can show you how to
produce a simple document of your own.

\section{Ordinary Text}      % Produces section heading.  Lower-level
                             % sections are begun with similar
                             % \subsection and \subsubsection commands.

The ends  of words and sentences are marked
  by   spaces. It  doesn't matter how many
spaces    you type; one is as good as 100.  The
end of   a line counts as a space.

One   or more   blank lines denote the  end
of  a paragraph.

Since any number of consecutive spaces are treated
like a single one, the formatting of the input
file makes no difference to
      \LaTeX,                % The \LaTeX command generates the LaTeX logo.
but it makes a difference to you.  When you use
\LaTeX, making your input file as easy to read
as possible will be a great help as you write
your document and when you change it.  This sample
file shows how you can add comments to your own input
file.

Because printing is different from typewriting,
there are a number of things that you have to do
differently when preparing an input file than if
you were just typing the document directly.
Quotation marks like
       ``this''
have to be handled specially, as do quotes within
quotes:
       ``\,`this'            % \, separates the double and single quote.
        is what I just
        wrote, not  `that'\,''.

Dashes come in three sizes: an
       intra-word
dash, a medium dash for number ranges like
       1--2,
and a punctuation
       dash---like
this.

A sentence-ending space should be larger than the
space between words within a sentence.  You
sometimes have to type special commands in
conjunction with punctuation characters to get
this right, as in the following sentence.
       Gnats, gnus, etc.\ all  % `\ ' makes an inter-word space.
       begin with G\@.         % \@ marks end-of-sentence punctuation.
You should check the spaces after periods when
reading your output to make sure you haven't
forgotten any special cases.  Generating an
ellipsis
       \ldots\               % `\ ' is needed after `\ldots' because TeX
                             % ignores spaces after command names like \ldots
                             % made from \ + letters.
                             %
                             % Note how a `%' character causes TeX to ignore
                             % the end of the input line, so these blank lines
                             % do not start a new paragraph.
                             %
with the right spacing around the periods requires
a special command.

\LaTeX\ interprets some common characters as
commands, so you must type special commands to
generate them.  These characters include the
following:
       \$ \& \% \# \{ and \}.

In printing, text is usually emphasized with an
       \emph{italic}
type style.

\begin{em}
   A long segment of text can also be emphasized
   in this way.  Text within such a segment can be
   given \emph{additional} emphasis.
\end{em}

It is sometimes necessary to prevent \LaTeX\ from
breaking a line where it might otherwise do so.
This may be at a space, as between the ``Mr.''\ and
``Jones'' in
       ``Mr.~Jones'',        % ~ produces an unbreakable interword space.
or within a word---especially when the word is a
symbol like
       \mbox{\emph{itemnum}}
that makes little sense when hyphenated across
lines.

Footnotes\footnote{This is an example of a footnote.}
pose no problem.

\LaTeX\ is good at typesetting mathematical formulas
like
       \( x-3y + z = 7 \)
or
       \( a_{1} > x^{2n} + y^{2n} > x' \)
or
       \( \ip{A}{B} = \sum_{i} a_{i} b_{i} \).
The spaces you type in a formula are
ignored.  Remember that a letter like
       $x$                   % $ ... $  and  \( ... \)  are equivalent
is a formula when it denotes a mathematical
symbol, and it should be typed as one.

\section{Displayed Text}

Text is displayed by indenting it from the left
margin.  Quotations are commonly displayed.  There
are short quotations
\begin{quote}
   This is a short quotation.  It consists of a
   single paragraph of text.  See how it is formatted.
\end{quote}
and longer ones.
\begin{quotation}
   This is a longer quotation.  It consists of two
   paragraphs of text, neither of which are
   particularly interesting.

   This is the second paragraph of the quotation.  It
   is just as dull as the first paragraph.
\end{quotation}
Another frequently-displayed structure is a list.
The following is an example of an \emph{itemized}
list.
\begin{itemize}
   \item This is the first item of an itemized list.
         Each item in the list is marked with a ``tick''.
         You don't have to worry about what kind of tick
         mark is used.

   \item This is the second item of the list.  It
         contains another list nested inside it.  The inner
         list is an \emph{enumerated} list.
         \begin{enumerate}
            \item This is the first item of an enumerated
                  list that is nested within the itemized list.

            \item This is the second item of the inner list.
                  \LaTeX\ allows you to nest lists deeper than
                  you really should.
         \end{enumerate}
         This is the rest of the second item of the outer
         list.  It is no more interesting than any other
         part of the item.
   \item This is the third item of the list.
\end{itemize}
You can even display poetry.
\begin{verse}
   There is an environment
    for verse \\             % The \\ command separates lines
   Whose features some poets % within a stanza.
   will curse.

                             % One or more blank lines separate stanzas.

   For instead of making\\
   Them do \emph{all} line breaking, \\
   It allows them to put too many words on a line when they'd rather be
   forced to be terse.
\end{verse}

Mathematical formulas may also be displayed.  A
displayed formula
is
one-line long; multiline
formulas require special formatting instructions.
   \[  \ip{\Gamma}{\psi'} = x'' + y^{2} + z_{i}^{n}\]
Don't start a paragraph with a displayed equation,
nor make one a paragraph by itself.

\end{document}               % End of document.

\section{Bench Huge Section 1558}
% synthetic bench section 1558
% This is a sample LaTeX input file.  (Version of 12 August 2004.)
%
% A '%' character causes TeX to ignore all remaining text on the line,
% and is used for comments like this one.

\documentclass{article}      % Specifies the document class

                             % The preamble begins here.
\title{An Example Document}  % Declares the document's title.
\author{Leslie Lamport}      % Declares the author's name.
\date{January 21, 1994}      % Deleting this command produces today's date.

\newcommand{\ip}[2]{(#1, #2)}
                             % Defines \ip{arg1}{arg2} to mean
                             % (arg1, arg2).

%\newcommand{\ip}[2]{\langle #1 | #2\rangle}
                             % This is an alternative definition of
                             % \ip that is commented out.

\begin{document}             % End of preamble and beginning of text.

\maketitle                   % Produces the title.

This is an example input file.  Comparing it with
the output it generates can show you how to
produce a simple document of your own.

\section{Ordinary Text}      % Produces section heading.  Lower-level
                             % sections are begun with similar
                             % \subsection and \subsubsection commands.

The ends  of words and sentences are marked
  by   spaces. It  doesn't matter how many
spaces    you type; one is as good as 100.  The
end of   a line counts as a space.

One   or more   blank lines denote the  end
of  a paragraph.

Since any number of consecutive spaces are treated
like a single one, the formatting of the input
file makes no difference to
      \LaTeX,                % The \LaTeX command generates the LaTeX logo.
but it makes a difference to you.  When you use
\LaTeX, making your input file as easy to read
as possible will be a great help as you write
your document and when you change it.  This sample
file shows how you can add comments to your own input
file.

Because printing is different from typewriting,
there are a number of things that you have to do
differently when preparing an input file than if
you were just typing the document directly.
Quotation marks like
       ``this''
have to be handled specially, as do quotes within
quotes:
       ``\,`this'            % \, separates the double and single quote.
        is what I just
        wrote, not  `that'\,''.

Dashes come in three sizes: an
       intra-word
dash, a medium dash for number ranges like
       1--2,
and a punctuation
       dash---like
this.

A sentence-ending space should be larger than the
space between words within a sentence.  You
sometimes have to type special commands in
conjunction with punctuation characters to get
this right, as in the following sentence.
       Gnats, gnus, etc.\ all  % `\ ' makes an inter-word space.
       begin with G\@.         % \@ marks end-of-sentence punctuation.
You should check the spaces after periods when
reading your output to make sure you haven't
forgotten any special cases.  Generating an
ellipsis
       \ldots\               % `\ ' is needed after `\ldots' because TeX
                             % ignores spaces after command names like \ldots
                             % made from \ + letters.
                             %
                             % Note how a `%' character causes TeX to ignore
                             % the end of the input line, so these blank lines
                             % do not start a new paragraph.
                             %
with the right spacing around the periods requires
a special command.

\LaTeX\ interprets some common characters as
commands, so you must type special commands to
generate them.  These characters include the
following:
       \$ \& \% \# \{ and \}.

In printing, text is usually emphasized with an
       \emph{italic}
type style.

\begin{em}
   A long segment of text can also be emphasized
   in this way.  Text within such a segment can be
   given \emph{additional} emphasis.
\end{em}

It is sometimes necessary to prevent \LaTeX\ from
breaking a line where it might otherwise do so.
This may be at a space, as between the ``Mr.''\ and
``Jones'' in
       ``Mr.~Jones'',        % ~ produces an unbreakable interword space.
or within a word---especially when the word is a
symbol like
       \mbox{\emph{itemnum}}
that makes little sense when hyphenated across
lines.

Footnotes\footnote{This is an example of a footnote.}
pose no problem.

\LaTeX\ is good at typesetting mathematical formulas
like
       \( x-3y + z = 7 \)
or
       \( a_{1} > x^{2n} + y^{2n} > x' \)
or
       \( \ip{A}{B} = \sum_{i} a_{i} b_{i} \).
The spaces you type in a formula are
ignored.  Remember that a letter like
       $x$                   % $ ... $  and  \( ... \)  are equivalent
is a formula when it denotes a mathematical
symbol, and it should be typed as one.

\section{Displayed Text}

Text is displayed by indenting it from the left
margin.  Quotations are commonly displayed.  There
are short quotations
\begin{quote}
   This is a short quotation.  It consists of a
   single paragraph of text.  See how it is formatted.
\end{quote}
and longer ones.
\begin{quotation}
   This is a longer quotation.  It consists of two
   paragraphs of text, neither of which are
   particularly interesting.

   This is the second paragraph of the quotation.  It
   is just as dull as the first paragraph.
\end{quotation}
Another frequently-displayed structure is a list.
The following is an example of an \emph{itemized}
list.
\begin{itemize}
   \item This is the first item of an itemized list.
         Each item in the list is marked with a ``tick''.
         You don't have to worry about what kind of tick
         mark is used.

   \item This is the second item of the list.  It
         contains another list nested inside it.  The inner
         list is an \emph{enumerated} list.
         \begin{enumerate}
            \item This is the first item of an enumerated
                  list that is nested within the itemized list.

            \item This is the second item of the inner list.
                  \LaTeX\ allows you to nest lists deeper than
                  you really should.
         \end{enumerate}
         This is the rest of the second item of the outer
         list.  It is no more interesting than any other
         part of the item.
   \item This is the third item of the list.
\end{itemize}
You can even display poetry.
\begin{verse}
   There is an environment
    for verse \\             % The \\ command separates lines
   Whose features some poets % within a stanza.
   will curse.

                             % One or more blank lines separate stanzas.

   For instead of making\\
   Them do \emph{all} line breaking, \\
   It allows them to put too many words on a line when they'd rather be
   forced to be terse.
\end{verse}

Mathematical formulas may also be displayed.  A
displayed formula
is
one-line long; multiline
formulas require special formatting instructions.
   \[  \ip{\Gamma}{\psi'} = x'' + y^{2} + z_{i}^{n}\]
Don't start a paragraph with a displayed equation,
nor make one a paragraph by itself.

\end{document}               % End of document.

\section{Bench Huge Section 1559}
% synthetic bench section 1559
% This is a sample LaTeX input file.  (Version of 12 August 2004.)
%
% A '%' character causes TeX to ignore all remaining text on the line,
% and is used for comments like this one.

\documentclass{article}      % Specifies the document class

                             % The preamble begins here.
\title{An Example Document}  % Declares the document's title.
\author{Leslie Lamport}      % Declares the author's name.
\date{January 21, 1994}      % Deleting this command produces today's date.

\newcommand{\ip}[2]{(#1, #2)}
                             % Defines \ip{arg1}{arg2} to mean
                             % (arg1, arg2).

%\newcommand{\ip}[2]{\langle #1 | #2\rangle}
                             % This is an alternative definition of
                             % \ip that is commented out.

\begin{document}             % End of preamble and beginning of text.

\maketitle                   % Produces the title.

This is an example input file.  Comparing it with
the output it generates can show you how to
produce a simple document of your own.

\section{Ordinary Text}      % Produces section heading.  Lower-level
                             % sections are begun with similar
                             % \subsection and \subsubsection commands.

The ends  of words and sentences are marked
  by   spaces. It  doesn't matter how many
spaces    you type; one is as good as 100.  The
end of   a line counts as a space.

One   or more   blank lines denote the  end
of  a paragraph.

Since any number of consecutive spaces are treated
like a single one, the formatting of the input
file makes no difference to
      \LaTeX,                % The \LaTeX command generates the LaTeX logo.
but it makes a difference to you.  When you use
\LaTeX, making your input file as easy to read
as possible will be a great help as you write
your document and when you change it.  This sample
file shows how you can add comments to your own input
file.

Because printing is different from typewriting,
there are a number of things that you have to do
differently when preparing an input file than if
you were just typing the document directly.
Quotation marks like
       ``this''
have to be handled specially, as do quotes within
quotes:
       ``\,`this'            % \, separates the double and single quote.
        is what I just
        wrote, not  `that'\,''.

Dashes come in three sizes: an
       intra-word
dash, a medium dash for number ranges like
       1--2,
and a punctuation
       dash---like
this.

A sentence-ending space should be larger than the
space between words within a sentence.  You
sometimes have to type special commands in
conjunction with punctuation characters to get
this right, as in the following sentence.
       Gnats, gnus, etc.\ all  % `\ ' makes an inter-word space.
       begin with G\@.         % \@ marks end-of-sentence punctuation.
You should check the spaces after periods when
reading your output to make sure you haven't
forgotten any special cases.  Generating an
ellipsis
       \ldots\               % `\ ' is needed after `\ldots' because TeX
                             % ignores spaces after command names like \ldots
                             % made from \ + letters.
                             %
                             % Note how a `%' character causes TeX to ignore
                             % the end of the input line, so these blank lines
                             % do not start a new paragraph.
                             %
with the right spacing around the periods requires
a special command.

\LaTeX\ interprets some common characters as
commands, so you must type special commands to
generate them.  These characters include the
following:
       \$ \& \% \# \{ and \}.

In printing, text is usually emphasized with an
       \emph{italic}
type style.

\begin{em}
   A long segment of text can also be emphasized
   in this way.  Text within such a segment can be
   given \emph{additional} emphasis.
\end{em}

It is sometimes necessary to prevent \LaTeX\ from
breaking a line where it might otherwise do so.
This may be at a space, as between the ``Mr.''\ and
``Jones'' in
       ``Mr.~Jones'',        % ~ produces an unbreakable interword space.
or within a word---especially when the word is a
symbol like
       \mbox{\emph{itemnum}}
that makes little sense when hyphenated across
lines.

Footnotes\footnote{This is an example of a footnote.}
pose no problem.

\LaTeX\ is good at typesetting mathematical formulas
like
       \( x-3y + z = 7 \)
or
       \( a_{1} > x^{2n} + y^{2n} > x' \)
or
       \( \ip{A}{B} = \sum_{i} a_{i} b_{i} \).
The spaces you type in a formula are
ignored.  Remember that a letter like
       $x$                   % $ ... $  and  \( ... \)  are equivalent
is a formula when it denotes a mathematical
symbol, and it should be typed as one.

\section{Displayed Text}

Text is displayed by indenting it from the left
margin.  Quotations are commonly displayed.  There
are short quotations
\begin{quote}
   This is a short quotation.  It consists of a
   single paragraph of text.  See how it is formatted.
\end{quote}
and longer ones.
\begin{quotation}
   This is a longer quotation.  It consists of two
   paragraphs of text, neither of which are
   particularly interesting.

   This is the second paragraph of the quotation.  It
   is just as dull as the first paragraph.
\end{quotation}
Another frequently-displayed structure is a list.
The following is an example of an \emph{itemized}
list.
\begin{itemize}
   \item This is the first item of an itemized list.
         Each item in the list is marked with a ``tick''.
         You don't have to worry about what kind of tick
         mark is used.

   \item This is the second item of the list.  It
         contains another list nested inside it.  The inner
         list is an \emph{enumerated} list.
         \begin{enumerate}
            \item This is the first item of an enumerated
                  list that is nested within the itemized list.

            \item This is the second item of the inner list.
                  \LaTeX\ allows you to nest lists deeper than
                  you really should.
         \end{enumerate}
         This is the rest of the second item of the outer
         list.  It is no more interesting than any other
         part of the item.
   \item This is the third item of the list.
\end{itemize}
You can even display poetry.
\begin{verse}
   There is an environment
    for verse \\             % The \\ command separates lines
   Whose features some poets % within a stanza.
   will curse.

                             % One or more blank lines separate stanzas.

   For instead of making\\
   Them do \emph{all} line breaking, \\
   It allows them to put too many words on a line when they'd rather be
   forced to be terse.
\end{verse}

Mathematical formulas may also be displayed.  A
displayed formula
is
one-line long; multiline
formulas require special formatting instructions.
   \[  \ip{\Gamma}{\psi'} = x'' + y^{2} + z_{i}^{n}\]
Don't start a paragraph with a displayed equation,
nor make one a paragraph by itself.

\end{document}               % End of document.

\section{Bench Huge Section 1560}
% synthetic bench section 1560
% This is a sample LaTeX input file.  (Version of 12 August 2004.)
%
% A '%' character causes TeX to ignore all remaining text on the line,
% and is used for comments like this one.

\documentclass{article}      % Specifies the document class

                             % The preamble begins here.
\title{An Example Document}  % Declares the document's title.
\author{Leslie Lamport}      % Declares the author's name.
\date{January 21, 1994}      % Deleting this command produces today's date.

\newcommand{\ip}[2]{(#1, #2)}
                             % Defines \ip{arg1}{arg2} to mean
                             % (arg1, arg2).

%\newcommand{\ip}[2]{\langle #1 | #2\rangle}
                             % This is an alternative definition of
                             % \ip that is commented out.

\begin{document}             % End of preamble and beginning of text.

\maketitle                   % Produces the title.

This is an example input file.  Comparing it with
the output it generates can show you how to
produce a simple document of your own.

\section{Ordinary Text}      % Produces section heading.  Lower-level
                             % sections are begun with similar
                             % \subsection and \subsubsection commands.

The ends  of words and sentences are marked
  by   spaces. It  doesn't matter how many
spaces    you type; one is as good as 100.  The
end of   a line counts as a space.

One   or more   blank lines denote the  end
of  a paragraph.

Since any number of consecutive spaces are treated
like a single one, the formatting of the input
file makes no difference to
      \LaTeX,                % The \LaTeX command generates the LaTeX logo.
but it makes a difference to you.  When you use
\LaTeX, making your input file as easy to read
as possible will be a great help as you write
your document and when you change it.  This sample
file shows how you can add comments to your own input
file.

Because printing is different from typewriting,
there are a number of things that you have to do
differently when preparing an input file than if
you were just typing the document directly.
Quotation marks like
       ``this''
have to be handled specially, as do quotes within
quotes:
       ``\,`this'            % \, separates the double and single quote.
        is what I just
        wrote, not  `that'\,''.

Dashes come in three sizes: an
       intra-word
dash, a medium dash for number ranges like
       1--2,
and a punctuation
       dash---like
this.

A sentence-ending space should be larger than the
space between words within a sentence.  You
sometimes have to type special commands in
conjunction with punctuation characters to get
this right, as in the following sentence.
       Gnats, gnus, etc.\ all  % `\ ' makes an inter-word space.
       begin with G\@.         % \@ marks end-of-sentence punctuation.
You should check the spaces after periods when
reading your output to make sure you haven't
forgotten any special cases.  Generating an
ellipsis
       \ldots\               % `\ ' is needed after `\ldots' because TeX
                             % ignores spaces after command names like \ldots
                             % made from \ + letters.
                             %
                             % Note how a `%' character causes TeX to ignore
                             % the end of the input line, so these blank lines
                             % do not start a new paragraph.
                             %
with the right spacing around the periods requires
a special command.

\LaTeX\ interprets some common characters as
commands, so you must type special commands to
generate them.  These characters include the
following:
       \$ \& \% \# \{ and \}.

In printing, text is usually emphasized with an
       \emph{italic}
type style.

\begin{em}
   A long segment of text can also be emphasized
   in this way.  Text within such a segment can be
   given \emph{additional} emphasis.
\end{em}

It is sometimes necessary to prevent \LaTeX\ from
breaking a line where it might otherwise do so.
This may be at a space, as between the ``Mr.''\ and
``Jones'' in
       ``Mr.~Jones'',        % ~ produces an unbreakable interword space.
or within a word---especially when the word is a
symbol like
       \mbox{\emph{itemnum}}
that makes little sense when hyphenated across
lines.

Footnotes\footnote{This is an example of a footnote.}
pose no problem.

\LaTeX\ is good at typesetting mathematical formulas
like
       \( x-3y + z = 7 \)
or
       \( a_{1} > x^{2n} + y^{2n} > x' \)
or
       \( \ip{A}{B} = \sum_{i} a_{i} b_{i} \).
The spaces you type in a formula are
ignored.  Remember that a letter like
       $x$                   % $ ... $  and  \( ... \)  are equivalent
is a formula when it denotes a mathematical
symbol, and it should be typed as one.

\section{Displayed Text}

Text is displayed by indenting it from the left
margin.  Quotations are commonly displayed.  There
are short quotations
\begin{quote}
   This is a short quotation.  It consists of a
   single paragraph of text.  See how it is formatted.
\end{quote}
and longer ones.
\begin{quotation}
   This is a longer quotation.  It consists of two
   paragraphs of text, neither of which are
   particularly interesting.

   This is the second paragraph of the quotation.  It
   is just as dull as the first paragraph.
\end{quotation}
Another frequently-displayed structure is a list.
The following is an example of an \emph{itemized}
list.
\begin{itemize}
   \item This is the first item of an itemized list.
         Each item in the list is marked with a ``tick''.
         You don't have to worry about what kind of tick
         mark is used.

   \item This is the second item of the list.  It
         contains another list nested inside it.  The inner
         list is an \emph{enumerated} list.
         \begin{enumerate}
            \item This is the first item of an enumerated
                  list that is nested within the itemized list.

            \item This is the second item of the inner list.
                  \LaTeX\ allows you to nest lists deeper than
                  you really should.
         \end{enumerate}
         This is the rest of the second item of the outer
         list.  It is no more interesting than any other
         part of the item.
   \item This is the third item of the list.
\end{itemize}
You can even display poetry.
\begin{verse}
   There is an environment
    for verse \\             % The \\ command separates lines
   Whose features some poets % within a stanza.
   will curse.

                             % One or more blank lines separate stanzas.

   For instead of making\\
   Them do \emph{all} line breaking, \\
   It allows them to put too many words on a line when they'd rather be
   forced to be terse.
\end{verse}

Mathematical formulas may also be displayed.  A
displayed formula
is
one-line long; multiline
formulas require special formatting instructions.
   \[  \ip{\Gamma}{\psi'} = x'' + y^{2} + z_{i}^{n}\]
Don't start a paragraph with a displayed equation,
nor make one a paragraph by itself.

\end{document}               % End of document.

\section{Bench Huge Section 1561}
% synthetic bench section 1561
% This is a sample LaTeX input file.  (Version of 12 August 2004.)
%
% A '%' character causes TeX to ignore all remaining text on the line,
% and is used for comments like this one.

\documentclass{article}      % Specifies the document class

                             % The preamble begins here.
\title{An Example Document}  % Declares the document's title.
\author{Leslie Lamport}      % Declares the author's name.
\date{January 21, 1994}      % Deleting this command produces today's date.

\newcommand{\ip}[2]{(#1, #2)}
                             % Defines \ip{arg1}{arg2} to mean
                             % (arg1, arg2).

%\newcommand{\ip}[2]{\langle #1 | #2\rangle}
                             % This is an alternative definition of
                             % \ip that is commented out.

\begin{document}             % End of preamble and beginning of text.

\maketitle                   % Produces the title.

This is an example input file.  Comparing it with
the output it generates can show you how to
produce a simple document of your own.

\section{Ordinary Text}      % Produces section heading.  Lower-level
                             % sections are begun with similar
                             % \subsection and \subsubsection commands.

The ends  of words and sentences are marked
  by   spaces. It  doesn't matter how many
spaces    you type; one is as good as 100.  The
end of   a line counts as a space.

One   or more   blank lines denote the  end
of  a paragraph.

Since any number of consecutive spaces are treated
like a single one, the formatting of the input
file makes no difference to
      \LaTeX,                % The \LaTeX command generates the LaTeX logo.
but it makes a difference to you.  When you use
\LaTeX, making your input file as easy to read
as possible will be a great help as you write
your document and when you change it.  This sample
file shows how you can add comments to your own input
file.

Because printing is different from typewriting,
there are a number of things that you have to do
differently when preparing an input file than if
you were just typing the document directly.
Quotation marks like
       ``this''
have to be handled specially, as do quotes within
quotes:
       ``\,`this'            % \, separates the double and single quote.
        is what I just
        wrote, not  `that'\,''.

Dashes come in three sizes: an
       intra-word
dash, a medium dash for number ranges like
       1--2,
and a punctuation
       dash---like
this.

A sentence-ending space should be larger than the
space between words within a sentence.  You
sometimes have to type special commands in
conjunction with punctuation characters to get
this right, as in the following sentence.
       Gnats, gnus, etc.\ all  % `\ ' makes an inter-word space.
       begin with G\@.         % \@ marks end-of-sentence punctuation.
You should check the spaces after periods when
reading your output to make sure you haven't
forgotten any special cases.  Generating an
ellipsis
       \ldots\               % `\ ' is needed after `\ldots' because TeX
                             % ignores spaces after command names like \ldots
                             % made from \ + letters.
                             %
                             % Note how a `%' character causes TeX to ignore
                             % the end of the input line, so these blank lines
                             % do not start a new paragraph.
                             %
with the right spacing around the periods requires
a special command.

\LaTeX\ interprets some common characters as
commands, so you must type special commands to
generate them.  These characters include the
following:
       \$ \& \% \# \{ and \}.

In printing, text is usually emphasized with an
       \emph{italic}
type style.

\begin{em}
   A long segment of text can also be emphasized
   in this way.  Text within such a segment can be
   given \emph{additional} emphasis.
\end{em}

It is sometimes necessary to prevent \LaTeX\ from
breaking a line where it might otherwise do so.
This may be at a space, as between the ``Mr.''\ and
``Jones'' in
       ``Mr.~Jones'',        % ~ produces an unbreakable interword space.
or within a word---especially when the word is a
symbol like
       \mbox{\emph{itemnum}}
that makes little sense when hyphenated across
lines.

Footnotes\footnote{This is an example of a footnote.}
pose no problem.

\LaTeX\ is good at typesetting mathematical formulas
like
       \( x-3y + z = 7 \)
or
       \( a_{1} > x^{2n} + y^{2n} > x' \)
or
       \( \ip{A}{B} = \sum_{i} a_{i} b_{i} \).
The spaces you type in a formula are
ignored.  Remember that a letter like
       $x$                   % $ ... $  and  \( ... \)  are equivalent
is a formula when it denotes a mathematical
symbol, and it should be typed as one.

\section{Displayed Text}

Text is displayed by indenting it from the left
margin.  Quotations are commonly displayed.  There
are short quotations
\begin{quote}
   This is a short quotation.  It consists of a
   single paragraph of text.  See how it is formatted.
\end{quote}
and longer ones.
\begin{quotation}
   This is a longer quotation.  It consists of two
   paragraphs of text, neither of which are
   particularly interesting.

   This is the second paragraph of the quotation.  It
   is just as dull as the first paragraph.
\end{quotation}
Another frequently-displayed structure is a list.
The following is an example of an \emph{itemized}
list.
\begin{itemize}
   \item This is the first item of an itemized list.
         Each item in the list is marked with a ``tick''.
         You don't have to worry about what kind of tick
         mark is used.

   \item This is the second item of the list.  It
         contains another list nested inside it.  The inner
         list is an \emph{enumerated} list.
         \begin{enumerate}
            \item This is the first item of an enumerated
                  list that is nested within the itemized list.

            \item This is the second item of the inner list.
                  \LaTeX\ allows you to nest lists deeper than
                  you really should.
         \end{enumerate}
         This is the rest of the second item of the outer
         list.  It is no more interesting than any other
         part of the item.
   \item This is the third item of the list.
\end{itemize}
You can even display poetry.
\begin{verse}
   There is an environment
    for verse \\             % The \\ command separates lines
   Whose features some poets % within a stanza.
   will curse.

                             % One or more blank lines separate stanzas.

   For instead of making\\
   Them do \emph{all} line breaking, \\
   It allows them to put too many words on a line when they'd rather be
   forced to be terse.
\end{verse}

Mathematical formulas may also be displayed.  A
displayed formula
is
one-line long; multiline
formulas require special formatting instructions.
   \[  \ip{\Gamma}{\psi'} = x'' + y^{2} + z_{i}^{n}\]
Don't start a paragraph with a displayed equation,
nor make one a paragraph by itself.

\end{document}               % End of document.

\section{Bench Huge Section 1562}
% synthetic bench section 1562
% This is a sample LaTeX input file.  (Version of 12 August 2004.)
%
% A '%' character causes TeX to ignore all remaining text on the line,
% and is used for comments like this one.

\documentclass{article}      % Specifies the document class

                             % The preamble begins here.
\title{An Example Document}  % Declares the document's title.
\author{Leslie Lamport}      % Declares the author's name.
\date{January 21, 1994}      % Deleting this command produces today's date.

\newcommand{\ip}[2]{(#1, #2)}
                             % Defines \ip{arg1}{arg2} to mean
                             % (arg1, arg2).

%\newcommand{\ip}[2]{\langle #1 | #2\rangle}
                             % This is an alternative definition of
                             % \ip that is commented out.

\begin{document}             % End of preamble and beginning of text.

\maketitle                   % Produces the title.

This is an example input file.  Comparing it with
the output it generates can show you how to
produce a simple document of your own.

\section{Ordinary Text}      % Produces section heading.  Lower-level
                             % sections are begun with similar
                             % \subsection and \subsubsection commands.

The ends  of words and sentences are marked
  by   spaces. It  doesn't matter how many
spaces    you type; one is as good as 100.  The
end of   a line counts as a space.

One   or more   blank lines denote the  end
of  a paragraph.

Since any number of consecutive spaces are treated
like a single one, the formatting of the input
file makes no difference to
      \LaTeX,                % The \LaTeX command generates the LaTeX logo.
but it makes a difference to you.  When you use
\LaTeX, making your input file as easy to read
as possible will be a great help as you write
your document and when you change it.  This sample
file shows how you can add comments to your own input
file.

Because printing is different from typewriting,
there are a number of things that you have to do
differently when preparing an input file than if
you were just typing the document directly.
Quotation marks like
       ``this''
have to be handled specially, as do quotes within
quotes:
       ``\,`this'            % \, separates the double and single quote.
        is what I just
        wrote, not  `that'\,''.

Dashes come in three sizes: an
       intra-word
dash, a medium dash for number ranges like
       1--2,
and a punctuation
       dash---like
this.

A sentence-ending space should be larger than the
space between words within a sentence.  You
sometimes have to type special commands in
conjunction with punctuation characters to get
this right, as in the following sentence.
       Gnats, gnus, etc.\ all  % `\ ' makes an inter-word space.
       begin with G\@.         % \@ marks end-of-sentence punctuation.
You should check the spaces after periods when
reading your output to make sure you haven't
forgotten any special cases.  Generating an
ellipsis
       \ldots\               % `\ ' is needed after `\ldots' because TeX
                             % ignores spaces after command names like \ldots
                             % made from \ + letters.
                             %
                             % Note how a `%' character causes TeX to ignore
                             % the end of the input line, so these blank lines
                             % do not start a new paragraph.
                             %
with the right spacing around the periods requires
a special command.

\LaTeX\ interprets some common characters as
commands, so you must type special commands to
generate them.  These characters include the
following:
       \$ \& \% \# \{ and \}.

In printing, text is usually emphasized with an
       \emph{italic}
type style.

\begin{em}
   A long segment of text can also be emphasized
   in this way.  Text within such a segment can be
   given \emph{additional} emphasis.
\end{em}

It is sometimes necessary to prevent \LaTeX\ from
breaking a line where it might otherwise do so.
This may be at a space, as between the ``Mr.''\ and
``Jones'' in
       ``Mr.~Jones'',        % ~ produces an unbreakable interword space.
or within a word---especially when the word is a
symbol like
       \mbox{\emph{itemnum}}
that makes little sense when hyphenated across
lines.

Footnotes\footnote{This is an example of a footnote.}
pose no problem.

\LaTeX\ is good at typesetting mathematical formulas
like
       \( x-3y + z = 7 \)
or
       \( a_{1} > x^{2n} + y^{2n} > x' \)
or
       \( \ip{A}{B} = \sum_{i} a_{i} b_{i} \).
The spaces you type in a formula are
ignored.  Remember that a letter like
       $x$                   % $ ... $  and  \( ... \)  are equivalent
is a formula when it denotes a mathematical
symbol, and it should be typed as one.

\section{Displayed Text}

Text is displayed by indenting it from the left
margin.  Quotations are commonly displayed.  There
are short quotations
\begin{quote}
   This is a short quotation.  It consists of a
   single paragraph of text.  See how it is formatted.
\end{quote}
and longer ones.
\begin{quotation}
   This is a longer quotation.  It consists of two
   paragraphs of text, neither of which are
   particularly interesting.

   This is the second paragraph of the quotation.  It
   is just as dull as the first paragraph.
\end{quotation}
Another frequently-displayed structure is a list.
The following is an example of an \emph{itemized}
list.
\begin{itemize}
   \item This is the first item of an itemized list.
         Each item in the list is marked with a ``tick''.
         You don't have to worry about what kind of tick
         mark is used.

   \item This is the second item of the list.  It
         contains another list nested inside it.  The inner
         list is an \emph{enumerated} list.
         \begin{enumerate}
            \item This is the first item of an enumerated
                  list that is nested within the itemized list.

            \item This is the second item of the inner list.
                  \LaTeX\ allows you to nest lists deeper than
                  you really should.
         \end{enumerate}
         This is the rest of the second item of the outer
         list.  It is no more interesting than any other
         part of the item.
   \item This is the third item of the list.
\end{itemize}
You can even display poetry.
\begin{verse}
   There is an environment
    for verse \\             % The \\ command separates lines
   Whose features some poets % within a stanza.
   will curse.

                             % One or more blank lines separate stanzas.

   For instead of making\\
   Them do \emph{all} line breaking, \\
   It allows them to put too many words on a line when they'd rather be
   forced to be terse.
\end{verse}

Mathematical formulas may also be displayed.  A
displayed formula
is
one-line long; multiline
formulas require special formatting instructions.
   \[  \ip{\Gamma}{\psi'} = x'' + y^{2} + z_{i}^{n}\]
Don't start a paragraph with a displayed equation,
nor make one a paragraph by itself.

\end{document}               % End of document.

\section{Bench Huge Section 1563}
% synthetic bench section 1563
% This is a sample LaTeX input file.  (Version of 12 August 2004.)
%
% A '%' character causes TeX to ignore all remaining text on the line,
% and is used for comments like this one.

\documentclass{article}      % Specifies the document class

                             % The preamble begins here.
\title{An Example Document}  % Declares the document's title.
\author{Leslie Lamport}      % Declares the author's name.
\date{January 21, 1994}      % Deleting this command produces today's date.

\newcommand{\ip}[2]{(#1, #2)}
                             % Defines \ip{arg1}{arg2} to mean
                             % (arg1, arg2).

%\newcommand{\ip}[2]{\langle #1 | #2\rangle}
                             % This is an alternative definition of
                             % \ip that is commented out.

\begin{document}             % End of preamble and beginning of text.

\maketitle                   % Produces the title.

This is an example input file.  Comparing it with
the output it generates can show you how to
produce a simple document of your own.

\section{Ordinary Text}      % Produces section heading.  Lower-level
                             % sections are begun with similar
                             % \subsection and \subsubsection commands.

The ends  of words and sentences are marked
  by   spaces. It  doesn't matter how many
spaces    you type; one is as good as 100.  The
end of   a line counts as a space.

One   or more   blank lines denote the  end
of  a paragraph.

Since any number of consecutive spaces are treated
like a single one, the formatting of the input
file makes no difference to
      \LaTeX,                % The \LaTeX command generates the LaTeX logo.
but it makes a difference to you.  When you use
\LaTeX, making your input file as easy to read
as possible will be a great help as you write
your document and when you change it.  This sample
file shows how you can add comments to your own input
file.

Because printing is different from typewriting,
there are a number of things that you have to do
differently when preparing an input file than if
you were just typing the document directly.
Quotation marks like
       ``this''
have to be handled specially, as do quotes within
quotes:
       ``\,`this'            % \, separates the double and single quote.
        is what I just
        wrote, not  `that'\,''.

Dashes come in three sizes: an
       intra-word
dash, a medium dash for number ranges like
       1--2,
and a punctuation
       dash---like
this.

A sentence-ending space should be larger than the
space between words within a sentence.  You
sometimes have to type special commands in
conjunction with punctuation characters to get
this right, as in the following sentence.
       Gnats, gnus, etc.\ all  % `\ ' makes an inter-word space.
       begin with G\@.         % \@ marks end-of-sentence punctuation.
You should check the spaces after periods when
reading your output to make sure you haven't
forgotten any special cases.  Generating an
ellipsis
       \ldots\               % `\ ' is needed after `\ldots' because TeX
                             % ignores spaces after command names like \ldots
                             % made from \ + letters.
                             %
                             % Note how a `%' character causes TeX to ignore
                             % the end of the input line, so these blank lines
                             % do not start a new paragraph.
                             %
with the right spacing around the periods requires
a special command.

\LaTeX\ interprets some common characters as
commands, so you must type special commands to
generate them.  These characters include the
following:
       \$ \& \% \# \{ and \}.

In printing, text is usually emphasized with an
       \emph{italic}
type style.

\begin{em}
   A long segment of text can also be emphasized
   in this way.  Text within such a segment can be
   given \emph{additional} emphasis.
\end{em}

It is sometimes necessary to prevent \LaTeX\ from
breaking a line where it might otherwise do so.
This may be at a space, as between the ``Mr.''\ and
``Jones'' in
       ``Mr.~Jones'',        % ~ produces an unbreakable interword space.
or within a word---especially when the word is a
symbol like
       \mbox{\emph{itemnum}}
that makes little sense when hyphenated across
lines.

Footnotes\footnote{This is an example of a footnote.}
pose no problem.

\LaTeX\ is good at typesetting mathematical formulas
like
       \( x-3y + z = 7 \)
or
       \( a_{1} > x^{2n} + y^{2n} > x' \)
or
       \( \ip{A}{B} = \sum_{i} a_{i} b_{i} \).
The spaces you type in a formula are
ignored.  Remember that a letter like
       $x$                   % $ ... $  and  \( ... \)  are equivalent
is a formula when it denotes a mathematical
symbol, and it should be typed as one.

\section{Displayed Text}

Text is displayed by indenting it from the left
margin.  Quotations are commonly displayed.  There
are short quotations
\begin{quote}
   This is a short quotation.  It consists of a
   single paragraph of text.  See how it is formatted.
\end{quote}
and longer ones.
\begin{quotation}
   This is a longer quotation.  It consists of two
   paragraphs of text, neither of which are
   particularly interesting.

   This is the second paragraph of the quotation.  It
   is just as dull as the first paragraph.
\end{quotation}
Another frequently-displayed structure is a list.
The following is an example of an \emph{itemized}
list.
\begin{itemize}
   \item This is the first item of an itemized list.
         Each item in the list is marked with a ``tick''.
         You don't have to worry about what kind of tick
         mark is used.

   \item This is the second item of the list.  It
         contains another list nested inside it.  The inner
         list is an \emph{enumerated} list.
         \begin{enumerate}
            \item This is the first item of an enumerated
                  list that is nested within the itemized list.

            \item This is the second item of the inner list.
                  \LaTeX\ allows you to nest lists deeper than
                  you really should.
         \end{enumerate}
         This is the rest of the second item of the outer
         list.  It is no more interesting than any other
         part of the item.
   \item This is the third item of the list.
\end{itemize}
You can even display poetry.
\begin{verse}
   There is an environment
    for verse \\             % The \\ command separates lines
   Whose features some poets % within a stanza.
   will curse.

                             % One or more blank lines separate stanzas.

   For instead of making\\
   Them do \emph{all} line breaking, \\
   It allows them to put too many words on a line when they'd rather be
   forced to be terse.
\end{verse}

Mathematical formulas may also be displayed.  A
displayed formula
is
one-line long; multiline
formulas require special formatting instructions.
   \[  \ip{\Gamma}{\psi'} = x'' + y^{2} + z_{i}^{n}\]
Don't start a paragraph with a displayed equation,
nor make one a paragraph by itself.

\end{document}               % End of document.

\section{Bench Huge Section 1564}
% synthetic bench section 1564
% This is a sample LaTeX input file.  (Version of 12 August 2004.)
%
% A '%' character causes TeX to ignore all remaining text on the line,
% and is used for comments like this one.

\documentclass{article}      % Specifies the document class

                             % The preamble begins here.
\title{An Example Document}  % Declares the document's title.
\author{Leslie Lamport}      % Declares the author's name.
\date{January 21, 1994}      % Deleting this command produces today's date.

\newcommand{\ip}[2]{(#1, #2)}
                             % Defines \ip{arg1}{arg2} to mean
                             % (arg1, arg2).

%\newcommand{\ip}[2]{\langle #1 | #2\rangle}
                             % This is an alternative definition of
                             % \ip that is commented out.

\begin{document}             % End of preamble and beginning of text.

\maketitle                   % Produces the title.

This is an example input file.  Comparing it with
the output it generates can show you how to
produce a simple document of your own.

\section{Ordinary Text}      % Produces section heading.  Lower-level
                             % sections are begun with similar
                             % \subsection and \subsubsection commands.

The ends  of words and sentences are marked
  by   spaces. It  doesn't matter how many
spaces    you type; one is as good as 100.  The
end of   a line counts as a space.

One   or more   blank lines denote the  end
of  a paragraph.

Since any number of consecutive spaces are treated
like a single one, the formatting of the input
file makes no difference to
      \LaTeX,                % The \LaTeX command generates the LaTeX logo.
but it makes a difference to you.  When you use
\LaTeX, making your input file as easy to read
as possible will be a great help as you write
your document and when you change it.  This sample
file shows how you can add comments to your own input
file.

Because printing is different from typewriting,
there are a number of things that you have to do
differently when preparing an input file than if
you were just typing the document directly.
Quotation marks like
       ``this''
have to be handled specially, as do quotes within
quotes:
       ``\,`this'            % \, separates the double and single quote.
        is what I just
        wrote, not  `that'\,''.

Dashes come in three sizes: an
       intra-word
dash, a medium dash for number ranges like
       1--2,
and a punctuation
       dash---like
this.

A sentence-ending space should be larger than the
space between words within a sentence.  You
sometimes have to type special commands in
conjunction with punctuation characters to get
this right, as in the following sentence.
       Gnats, gnus, etc.\ all  % `\ ' makes an inter-word space.
       begin with G\@.         % \@ marks end-of-sentence punctuation.
You should check the spaces after periods when
reading your output to make sure you haven't
forgotten any special cases.  Generating an
ellipsis
       \ldots\               % `\ ' is needed after `\ldots' because TeX
                             % ignores spaces after command names like \ldots
                             % made from \ + letters.
                             %
                             % Note how a `%' character causes TeX to ignore
                             % the end of the input line, so these blank lines
                             % do not start a new paragraph.
                             %
with the right spacing around the periods requires
a special command.

\LaTeX\ interprets some common characters as
commands, so you must type special commands to
generate them.  These characters include the
following:
       \$ \& \% \# \{ and \}.

In printing, text is usually emphasized with an
       \emph{italic}
type style.

\begin{em}
   A long segment of text can also be emphasized
   in this way.  Text within such a segment can be
   given \emph{additional} emphasis.
\end{em}

It is sometimes necessary to prevent \LaTeX\ from
breaking a line where it might otherwise do so.
This may be at a space, as between the ``Mr.''\ and
``Jones'' in
       ``Mr.~Jones'',        % ~ produces an unbreakable interword space.
or within a word---especially when the word is a
symbol like
       \mbox{\emph{itemnum}}
that makes little sense when hyphenated across
lines.

Footnotes\footnote{This is an example of a footnote.}
pose no problem.

\LaTeX\ is good at typesetting mathematical formulas
like
       \( x-3y + z = 7 \)
or
       \( a_{1} > x^{2n} + y^{2n} > x' \)
or
       \( \ip{A}{B} = \sum_{i} a_{i} b_{i} \).
The spaces you type in a formula are
ignored.  Remember that a letter like
       $x$                   % $ ... $  and  \( ... \)  are equivalent
is a formula when it denotes a mathematical
symbol, and it should be typed as one.

\section{Displayed Text}

Text is displayed by indenting it from the left
margin.  Quotations are commonly displayed.  There
are short quotations
\begin{quote}
   This is a short quotation.  It consists of a
   single paragraph of text.  See how it is formatted.
\end{quote}
and longer ones.
\begin{quotation}
   This is a longer quotation.  It consists of two
   paragraphs of text, neither of which are
   particularly interesting.

   This is the second paragraph of the quotation.  It
   is just as dull as the first paragraph.
\end{quotation}
Another frequently-displayed structure is a list.
The following is an example of an \emph{itemized}
list.
\begin{itemize}
   \item This is the first item of an itemized list.
         Each item in the list is marked with a ``tick''.
         You don't have to worry about what kind of tick
         mark is used.

   \item This is the second item of the list.  It
         contains another list nested inside it.  The inner
         list is an \emph{enumerated} list.
         \begin{enumerate}
            \item This is the first item of an enumerated
                  list that is nested within the itemized list.

            \item This is the second item of the inner list.
                  \LaTeX\ allows you to nest lists deeper than
                  you really should.
         \end{enumerate}
         This is the rest of the second item of the outer
         list.  It is no more interesting than any other
         part of the item.
   \item This is the third item of the list.
\end{itemize}
You can even display poetry.
\begin{verse}
   There is an environment
    for verse \\             % The \\ command separates lines
   Whose features some poets % within a stanza.
   will curse.

                             % One or more blank lines separate stanzas.

   For instead of making\\
   Them do \emph{all} line breaking, \\
   It allows them to put too many words on a line when they'd rather be
   forced to be terse.
\end{verse}

Mathematical formulas may also be displayed.  A
displayed formula
is
one-line long; multiline
formulas require special formatting instructions.
   \[  \ip{\Gamma}{\psi'} = x'' + y^{2} + z_{i}^{n}\]
Don't start a paragraph with a displayed equation,
nor make one a paragraph by itself.

\end{document}               % End of document.

\section{Bench Huge Section 1565}
% synthetic bench section 1565
% This is a sample LaTeX input file.  (Version of 12 August 2004.)
%
% A '%' character causes TeX to ignore all remaining text on the line,
% and is used for comments like this one.

\documentclass{article}      % Specifies the document class

                             % The preamble begins here.
\title{An Example Document}  % Declares the document's title.
\author{Leslie Lamport}      % Declares the author's name.
\date{January 21, 1994}      % Deleting this command produces today's date.

\newcommand{\ip}[2]{(#1, #2)}
                             % Defines \ip{arg1}{arg2} to mean
                             % (arg1, arg2).

%\newcommand{\ip}[2]{\langle #1 | #2\rangle}
                             % This is an alternative definition of
                             % \ip that is commented out.

\begin{document}             % End of preamble and beginning of text.

\maketitle                   % Produces the title.

This is an example input file.  Comparing it with
the output it generates can show you how to
produce a simple document of your own.

\section{Ordinary Text}      % Produces section heading.  Lower-level
                             % sections are begun with similar
                             % \subsection and \subsubsection commands.

The ends  of words and sentences are marked
  by   spaces. It  doesn't matter how many
spaces    you type; one is as good as 100.  The
end of   a line counts as a space.

One   or more   blank lines denote the  end
of  a paragraph.

Since any number of consecutive spaces are treated
like a single one, the formatting of the input
file makes no difference to
      \LaTeX,                % The \LaTeX command generates the LaTeX logo.
but it makes a difference to you.  When you use
\LaTeX, making your input file as easy to read
as possible will be a great help as you write
your document and when you change it.  This sample
file shows how you can add comments to your own input
file.

Because printing is different from typewriting,
there are a number of things that you have to do
differently when preparing an input file than if
you were just typing the document directly.
Quotation marks like
       ``this''
have to be handled specially, as do quotes within
quotes:
       ``\,`this'            % \, separates the double and single quote.
        is what I just
        wrote, not  `that'\,''.

Dashes come in three sizes: an
       intra-word
dash, a medium dash for number ranges like
       1--2,
and a punctuation
       dash---like
this.

A sentence-ending space should be larger than the
space between words within a sentence.  You
sometimes have to type special commands in
conjunction with punctuation characters to get
this right, as in the following sentence.
       Gnats, gnus, etc.\ all  % `\ ' makes an inter-word space.
       begin with G\@.         % \@ marks end-of-sentence punctuation.
You should check the spaces after periods when
reading your output to make sure you haven't
forgotten any special cases.  Generating an
ellipsis
       \ldots\               % `\ ' is needed after `\ldots' because TeX
                             % ignores spaces after command names like \ldots
                             % made from \ + letters.
                             %
                             % Note how a `%' character causes TeX to ignore
                             % the end of the input line, so these blank lines
                             % do not start a new paragraph.
                             %
with the right spacing around the periods requires
a special command.

\LaTeX\ interprets some common characters as
commands, so you must type special commands to
generate them.  These characters include the
following:
       \$ \& \% \# \{ and \}.

In printing, text is usually emphasized with an
       \emph{italic}
type style.

\begin{em}
   A long segment of text can also be emphasized
   in this way.  Text within such a segment can be
   given \emph{additional} emphasis.
\end{em}

It is sometimes necessary to prevent \LaTeX\ from
breaking a line where it might otherwise do so.
This may be at a space, as between the ``Mr.''\ and
``Jones'' in
       ``Mr.~Jones'',        % ~ produces an unbreakable interword space.
or within a word---especially when the word is a
symbol like
       \mbox{\emph{itemnum}}
that makes little sense when hyphenated across
lines.

Footnotes\footnote{This is an example of a footnote.}
pose no problem.

\LaTeX\ is good at typesetting mathematical formulas
like
       \( x-3y + z = 7 \)
or
       \( a_{1} > x^{2n} + y^{2n} > x' \)
or
       \( \ip{A}{B} = \sum_{i} a_{i} b_{i} \).
The spaces you type in a formula are
ignored.  Remember that a letter like
       $x$                   % $ ... $  and  \( ... \)  are equivalent
is a formula when it denotes a mathematical
symbol, and it should be typed as one.

\section{Displayed Text}

Text is displayed by indenting it from the left
margin.  Quotations are commonly displayed.  There
are short quotations
\begin{quote}
   This is a short quotation.  It consists of a
   single paragraph of text.  See how it is formatted.
\end{quote}
and longer ones.
\begin{quotation}
   This is a longer quotation.  It consists of two
   paragraphs of text, neither of which are
   particularly interesting.

   This is the second paragraph of the quotation.  It
   is just as dull as the first paragraph.
\end{quotation}
Another frequently-displayed structure is a list.
The following is an example of an \emph{itemized}
list.
\begin{itemize}
   \item This is the first item of an itemized list.
         Each item in the list is marked with a ``tick''.
         You don't have to worry about what kind of tick
         mark is used.

   \item This is the second item of the list.  It
         contains another list nested inside it.  The inner
         list is an \emph{enumerated} list.
         \begin{enumerate}
            \item This is the first item of an enumerated
                  list that is nested within the itemized list.

            \item This is the second item of the inner list.
                  \LaTeX\ allows you to nest lists deeper than
                  you really should.
         \end{enumerate}
         This is the rest of the second item of the outer
         list.  It is no more interesting than any other
         part of the item.
   \item This is the third item of the list.
\end{itemize}
You can even display poetry.
\begin{verse}
   There is an environment
    for verse \\             % The \\ command separates lines
   Whose features some poets % within a stanza.
   will curse.

                             % One or more blank lines separate stanzas.

   For instead of making\\
   Them do \emph{all} line breaking, \\
   It allows them to put too many words on a line when they'd rather be
   forced to be terse.
\end{verse}

Mathematical formulas may also be displayed.  A
displayed formula
is
one-line long; multiline
formulas require special formatting instructions.
   \[  \ip{\Gamma}{\psi'} = x'' + y^{2} + z_{i}^{n}\]
Don't start a paragraph with a displayed equation,
nor make one a paragraph by itself.

\end{document}               % End of document.

\section{Bench Huge Section 1566}
% synthetic bench section 1566
% This is a sample LaTeX input file.  (Version of 12 August 2004.)
%
% A '%' character causes TeX to ignore all remaining text on the line,
% and is used for comments like this one.

\documentclass{article}      % Specifies the document class

                             % The preamble begins here.
\title{An Example Document}  % Declares the document's title.
\author{Leslie Lamport}      % Declares the author's name.
\date{January 21, 1994}      % Deleting this command produces today's date.

\newcommand{\ip}[2]{(#1, #2)}
                             % Defines \ip{arg1}{arg2} to mean
                             % (arg1, arg2).

%\newcommand{\ip}[2]{\langle #1 | #2\rangle}
                             % This is an alternative definition of
                             % \ip that is commented out.

\begin{document}             % End of preamble and beginning of text.

\maketitle                   % Produces the title.

This is an example input file.  Comparing it with
the output it generates can show you how to
produce a simple document of your own.

\section{Ordinary Text}      % Produces section heading.  Lower-level
                             % sections are begun with similar
                             % \subsection and \subsubsection commands.

The ends  of words and sentences are marked
  by   spaces. It  doesn't matter how many
spaces    you type; one is as good as 100.  The
end of   a line counts as a space.

One   or more   blank lines denote the  end
of  a paragraph.

Since any number of consecutive spaces are treated
like a single one, the formatting of the input
file makes no difference to
      \LaTeX,                % The \LaTeX command generates the LaTeX logo.
but it makes a difference to you.  When you use
\LaTeX, making your input file as easy to read
as possible will be a great help as you write
your document and when you change it.  This sample
file shows how you can add comments to your own input
file.

Because printing is different from typewriting,
there are a number of things that you have to do
differently when preparing an input file than if
you were just typing the document directly.
Quotation marks like
       ``this''
have to be handled specially, as do quotes within
quotes:
       ``\,`this'            % \, separates the double and single quote.
        is what I just
        wrote, not  `that'\,''.

Dashes come in three sizes: an
       intra-word
dash, a medium dash for number ranges like
       1--2,
and a punctuation
       dash---like
this.

A sentence-ending space should be larger than the
space between words within a sentence.  You
sometimes have to type special commands in
conjunction with punctuation characters to get
this right, as in the following sentence.
       Gnats, gnus, etc.\ all  % `\ ' makes an inter-word space.
       begin with G\@.         % \@ marks end-of-sentence punctuation.
You should check the spaces after periods when
reading your output to make sure you haven't
forgotten any special cases.  Generating an
ellipsis
       \ldots\               % `\ ' is needed after `\ldots' because TeX
                             % ignores spaces after command names like \ldots
                             % made from \ + letters.
                             %
                             % Note how a `%' character causes TeX to ignore
                             % the end of the input line, so these blank lines
                             % do not start a new paragraph.
                             %
with the right spacing around the periods requires
a special command.

\LaTeX\ interprets some common characters as
commands, so you must type special commands to
generate them.  These characters include the
following:
       \$ \& \% \# \{ and \}.

In printing, text is usually emphasized with an
       \emph{italic}
type style.

\begin{em}
   A long segment of text can also be emphasized
   in this way.  Text within such a segment can be
   given \emph{additional} emphasis.
\end{em}

It is sometimes necessary to prevent \LaTeX\ from
breaking a line where it might otherwise do so.
This may be at a space, as between the ``Mr.''\ and
``Jones'' in
       ``Mr.~Jones'',        % ~ produces an unbreakable interword space.
or within a word---especially when the word is a
symbol like
       \mbox{\emph{itemnum}}
that makes little sense when hyphenated across
lines.

Footnotes\footnote{This is an example of a footnote.}
pose no problem.

\LaTeX\ is good at typesetting mathematical formulas
like
       \( x-3y + z = 7 \)
or
       \( a_{1} > x^{2n} + y^{2n} > x' \)
or
       \( \ip{A}{B} = \sum_{i} a_{i} b_{i} \).
The spaces you type in a formula are
ignored.  Remember that a letter like
       $x$                   % $ ... $  and  \( ... \)  are equivalent
is a formula when it denotes a mathematical
symbol, and it should be typed as one.

\section{Displayed Text}

Text is displayed by indenting it from the left
margin.  Quotations are commonly displayed.  There
are short quotations
\begin{quote}
   This is a short quotation.  It consists of a
   single paragraph of text.  See how it is formatted.
\end{quote}
and longer ones.
\begin{quotation}
   This is a longer quotation.  It consists of two
   paragraphs of text, neither of which are
   particularly interesting.

   This is the second paragraph of the quotation.  It
   is just as dull as the first paragraph.
\end{quotation}
Another frequently-displayed structure is a list.
The following is an example of an \emph{itemized}
list.
\begin{itemize}
   \item This is the first item of an itemized list.
         Each item in the list is marked with a ``tick''.
         You don't have to worry about what kind of tick
         mark is used.

   \item This is the second item of the list.  It
         contains another list nested inside it.  The inner
         list is an \emph{enumerated} list.
         \begin{enumerate}
            \item This is the first item of an enumerated
                  list that is nested within the itemized list.

            \item This is the second item of the inner list.
                  \LaTeX\ allows you to nest lists deeper than
                  you really should.
         \end{enumerate}
         This is the rest of the second item of the outer
         list.  It is no more interesting than any other
         part of the item.
   \item This is the third item of the list.
\end{itemize}
You can even display poetry.
\begin{verse}
   There is an environment
    for verse \\             % The \\ command separates lines
   Whose features some poets % within a stanza.
   will curse.

                             % One or more blank lines separate stanzas.

   For instead of making\\
   Them do \emph{all} line breaking, \\
   It allows them to put too many words on a line when they'd rather be
   forced to be terse.
\end{verse}

Mathematical formulas may also be displayed.  A
displayed formula
is
one-line long; multiline
formulas require special formatting instructions.
   \[  \ip{\Gamma}{\psi'} = x'' + y^{2} + z_{i}^{n}\]
Don't start a paragraph with a displayed equation,
nor make one a paragraph by itself.

\end{document}               % End of document.

\section{Bench Huge Section 1567}
% synthetic bench section 1567
% This is a sample LaTeX input file.  (Version of 12 August 2004.)
%
% A '%' character causes TeX to ignore all remaining text on the line,
% and is used for comments like this one.

\documentclass{article}      % Specifies the document class

                             % The preamble begins here.
\title{An Example Document}  % Declares the document's title.
\author{Leslie Lamport}      % Declares the author's name.
\date{January 21, 1994}      % Deleting this command produces today's date.

\newcommand{\ip}[2]{(#1, #2)}
                             % Defines \ip{arg1}{arg2} to mean
                             % (arg1, arg2).

%\newcommand{\ip}[2]{\langle #1 | #2\rangle}
                             % This is an alternative definition of
                             % \ip that is commented out.

\begin{document}             % End of preamble and beginning of text.

\maketitle                   % Produces the title.

This is an example input file.  Comparing it with
the output it generates can show you how to
produce a simple document of your own.

\section{Ordinary Text}      % Produces section heading.  Lower-level
                             % sections are begun with similar
                             % \subsection and \subsubsection commands.

The ends  of words and sentences are marked
  by   spaces. It  doesn't matter how many
spaces    you type; one is as good as 100.  The
end of   a line counts as a space.

One   or more   blank lines denote the  end
of  a paragraph.

Since any number of consecutive spaces are treated
like a single one, the formatting of the input
file makes no difference to
      \LaTeX,                % The \LaTeX command generates the LaTeX logo.
but it makes a difference to you.  When you use
\LaTeX, making your input file as easy to read
as possible will be a great help as you write
your document and when you change it.  This sample
file shows how you can add comments to your own input
file.

Because printing is different from typewriting,
there are a number of things that you have to do
differently when preparing an input file than if
you were just typing the document directly.
Quotation marks like
       ``this''
have to be handled specially, as do quotes within
quotes:
       ``\,`this'            % \, separates the double and single quote.
        is what I just
        wrote, not  `that'\,''.

Dashes come in three sizes: an
       intra-word
dash, a medium dash for number ranges like
       1--2,
and a punctuation
       dash---like
this.

A sentence-ending space should be larger than the
space between words within a sentence.  You
sometimes have to type special commands in
conjunction with punctuation characters to get
this right, as in the following sentence.
       Gnats, gnus, etc.\ all  % `\ ' makes an inter-word space.
       begin with G\@.         % \@ marks end-of-sentence punctuation.
You should check the spaces after periods when
reading your output to make sure you haven't
forgotten any special cases.  Generating an
ellipsis
       \ldots\               % `\ ' is needed after `\ldots' because TeX
                             % ignores spaces after command names like \ldots
                             % made from \ + letters.
                             %
                             % Note how a `%' character causes TeX to ignore
                             % the end of the input line, so these blank lines
                             % do not start a new paragraph.
                             %
with the right spacing around the periods requires
a special command.

\LaTeX\ interprets some common characters as
commands, so you must type special commands to
generate them.  These characters include the
following:
       \$ \& \% \# \{ and \}.

In printing, text is usually emphasized with an
       \emph{italic}
type style.

\begin{em}
   A long segment of text can also be emphasized
   in this way.  Text within such a segment can be
   given \emph{additional} emphasis.
\end{em}

It is sometimes necessary to prevent \LaTeX\ from
breaking a line where it might otherwise do so.
This may be at a space, as between the ``Mr.''\ and
``Jones'' in
       ``Mr.~Jones'',        % ~ produces an unbreakable interword space.
or within a word---especially when the word is a
symbol like
       \mbox{\emph{itemnum}}
that makes little sense when hyphenated across
lines.

Footnotes\footnote{This is an example of a footnote.}
pose no problem.

\LaTeX\ is good at typesetting mathematical formulas
like
       \( x-3y + z = 7 \)
or
       \( a_{1} > x^{2n} + y^{2n} > x' \)
or
       \( \ip{A}{B} = \sum_{i} a_{i} b_{i} \).
The spaces you type in a formula are
ignored.  Remember that a letter like
       $x$                   % $ ... $  and  \( ... \)  are equivalent
is a formula when it denotes a mathematical
symbol, and it should be typed as one.

\section{Displayed Text}

Text is displayed by indenting it from the left
margin.  Quotations are commonly displayed.  There
are short quotations
\begin{quote}
   This is a short quotation.  It consists of a
   single paragraph of text.  See how it is formatted.
\end{quote}
and longer ones.
\begin{quotation}
   This is a longer quotation.  It consists of two
   paragraphs of text, neither of which are
   particularly interesting.

   This is the second paragraph of the quotation.  It
   is just as dull as the first paragraph.
\end{quotation}
Another frequently-displayed structure is a list.
The following is an example of an \emph{itemized}
list.
\begin{itemize}
   \item This is the first item of an itemized list.
         Each item in the list is marked with a ``tick''.
         You don't have to worry about what kind of tick
         mark is used.

   \item This is the second item of the list.  It
         contains another list nested inside it.  The inner
         list is an \emph{enumerated} list.
         \begin{enumerate}
            \item This is the first item of an enumerated
                  list that is nested within the itemized list.

            \item This is the second item of the inner list.
                  \LaTeX\ allows you to nest lists deeper than
                  you really should.
         \end{enumerate}
         This is the rest of the second item of the outer
         list.  It is no more interesting than any other
         part of the item.
   \item This is the third item of the list.
\end{itemize}
You can even display poetry.
\begin{verse}
   There is an environment
    for verse \\             % The \\ command separates lines
   Whose features some poets % within a stanza.
   will curse.

                             % One or more blank lines separate stanzas.

   For instead of making\\
   Them do \emph{all} line breaking, \\
   It allows them to put too many words on a line when they'd rather be
   forced to be terse.
\end{verse}

Mathematical formulas may also be displayed.  A
displayed formula
is
one-line long; multiline
formulas require special formatting instructions.
   \[  \ip{\Gamma}{\psi'} = x'' + y^{2} + z_{i}^{n}\]
Don't start a paragraph with a displayed equation,
nor make one a paragraph by itself.

\end{document}               % End of document.

\section{Bench Huge Section 1568}
% synthetic bench section 1568
% This is a sample LaTeX input file.  (Version of 12 August 2004.)
%
% A '%' character causes TeX to ignore all remaining text on the line,
% and is used for comments like this one.

\documentclass{article}      % Specifies the document class

                             % The preamble begins here.
\title{An Example Document}  % Declares the document's title.
\author{Leslie Lamport}      % Declares the author's name.
\date{January 21, 1994}      % Deleting this command produces today's date.

\newcommand{\ip}[2]{(#1, #2)}
                             % Defines \ip{arg1}{arg2} to mean
                             % (arg1, arg2).

%\newcommand{\ip}[2]{\langle #1 | #2\rangle}
                             % This is an alternative definition of
                             % \ip that is commented out.

\begin{document}             % End of preamble and beginning of text.

\maketitle                   % Produces the title.

This is an example input file.  Comparing it with
the output it generates can show you how to
produce a simple document of your own.

\section{Ordinary Text}      % Produces section heading.  Lower-level
                             % sections are begun with similar
                             % \subsection and \subsubsection commands.

The ends  of words and sentences are marked
  by   spaces. It  doesn't matter how many
spaces    you type; one is as good as 100.  The
end of   a line counts as a space.

One   or more   blank lines denote the  end
of  a paragraph.

Since any number of consecutive spaces are treated
like a single one, the formatting of the input
file makes no difference to
      \LaTeX,                % The \LaTeX command generates the LaTeX logo.
but it makes a difference to you.  When you use
\LaTeX, making your input file as easy to read
as possible will be a great help as you write
your document and when you change it.  This sample
file shows how you can add comments to your own input
file.

Because printing is different from typewriting,
there are a number of things that you have to do
differently when preparing an input file than if
you were just typing the document directly.
Quotation marks like
       ``this''
have to be handled specially, as do quotes within
quotes:
       ``\,`this'            % \, separates the double and single quote.
        is what I just
        wrote, not  `that'\,''.

Dashes come in three sizes: an
       intra-word
dash, a medium dash for number ranges like
       1--2,
and a punctuation
       dash---like
this.

A sentence-ending space should be larger than the
space between words within a sentence.  You
sometimes have to type special commands in
conjunction with punctuation characters to get
this right, as in the following sentence.
       Gnats, gnus, etc.\ all  % `\ ' makes an inter-word space.
       begin with G\@.         % \@ marks end-of-sentence punctuation.
You should check the spaces after periods when
reading your output to make sure you haven't
forgotten any special cases.  Generating an
ellipsis
       \ldots\               % `\ ' is needed after `\ldots' because TeX
                             % ignores spaces after command names like \ldots
                             % made from \ + letters.
                             %
                             % Note how a `%' character causes TeX to ignore
                             % the end of the input line, so these blank lines
                             % do not start a new paragraph.
                             %
with the right spacing around the periods requires
a special command.

\LaTeX\ interprets some common characters as
commands, so you must type special commands to
generate them.  These characters include the
following:
       \$ \& \% \# \{ and \}.

In printing, text is usually emphasized with an
       \emph{italic}
type style.

\begin{em}
   A long segment of text can also be emphasized
   in this way.  Text within such a segment can be
   given \emph{additional} emphasis.
\end{em}

It is sometimes necessary to prevent \LaTeX\ from
breaking a line where it might otherwise do so.
This may be at a space, as between the ``Mr.''\ and
``Jones'' in
       ``Mr.~Jones'',        % ~ produces an unbreakable interword space.
or within a word---especially when the word is a
symbol like
       \mbox{\emph{itemnum}}
that makes little sense when hyphenated across
lines.

Footnotes\footnote{This is an example of a footnote.}
pose no problem.

\LaTeX\ is good at typesetting mathematical formulas
like
       \( x-3y + z = 7 \)
or
       \( a_{1} > x^{2n} + y^{2n} > x' \)
or
       \( \ip{A}{B} = \sum_{i} a_{i} b_{i} \).
The spaces you type in a formula are
ignored.  Remember that a letter like
       $x$                   % $ ... $  and  \( ... \)  are equivalent
is a formula when it denotes a mathematical
symbol, and it should be typed as one.

\section{Displayed Text}

Text is displayed by indenting it from the left
margin.  Quotations are commonly displayed.  There
are short quotations
\begin{quote}
   This is a short quotation.  It consists of a
   single paragraph of text.  See how it is formatted.
\end{quote}
and longer ones.
\begin{quotation}
   This is a longer quotation.  It consists of two
   paragraphs of text, neither of which are
   particularly interesting.

   This is the second paragraph of the quotation.  It
   is just as dull as the first paragraph.
\end{quotation}
Another frequently-displayed structure is a list.
The following is an example of an \emph{itemized}
list.
\begin{itemize}
   \item This is the first item of an itemized list.
         Each item in the list is marked with a ``tick''.
         You don't have to worry about what kind of tick
         mark is used.

   \item This is the second item of the list.  It
         contains another list nested inside it.  The inner
         list is an \emph{enumerated} list.
         \begin{enumerate}
            \item This is the first item of an enumerated
                  list that is nested within the itemized list.

            \item This is the second item of the inner list.
                  \LaTeX\ allows you to nest lists deeper than
                  you really should.
         \end{enumerate}
         This is the rest of the second item of the outer
         list.  It is no more interesting than any other
         part of the item.
   \item This is the third item of the list.
\end{itemize}
You can even display poetry.
\begin{verse}
   There is an environment
    for verse \\             % The \\ command separates lines
   Whose features some poets % within a stanza.
   will curse.

                             % One or more blank lines separate stanzas.

   For instead of making\\
   Them do \emph{all} line breaking, \\
   It allows them to put too many words on a line when they'd rather be
   forced to be terse.
\end{verse}

Mathematical formulas may also be displayed.  A
displayed formula
is
one-line long; multiline
formulas require special formatting instructions.
   \[  \ip{\Gamma}{\psi'} = x'' + y^{2} + z_{i}^{n}\]
Don't start a paragraph with a displayed equation,
nor make one a paragraph by itself.

\end{document}               % End of document.

\section{Bench Huge Section 1569}
% synthetic bench section 1569
% This is a sample LaTeX input file.  (Version of 12 August 2004.)
%
% A '%' character causes TeX to ignore all remaining text on the line,
% and is used for comments like this one.

\documentclass{article}      % Specifies the document class

                             % The preamble begins here.
\title{An Example Document}  % Declares the document's title.
\author{Leslie Lamport}      % Declares the author's name.
\date{January 21, 1994}      % Deleting this command produces today's date.

\newcommand{\ip}[2]{(#1, #2)}
                             % Defines \ip{arg1}{arg2} to mean
                             % (arg1, arg2).

%\newcommand{\ip}[2]{\langle #1 | #2\rangle}
                             % This is an alternative definition of
                             % \ip that is commented out.

\begin{document}             % End of preamble and beginning of text.

\maketitle                   % Produces the title.

This is an example input file.  Comparing it with
the output it generates can show you how to
produce a simple document of your own.

\section{Ordinary Text}      % Produces section heading.  Lower-level
                             % sections are begun with similar
                             % \subsection and \subsubsection commands.

The ends  of words and sentences are marked
  by   spaces. It  doesn't matter how many
spaces    you type; one is as good as 100.  The
end of   a line counts as a space.

One   or more   blank lines denote the  end
of  a paragraph.

Since any number of consecutive spaces are treated
like a single one, the formatting of the input
file makes no difference to
      \LaTeX,                % The \LaTeX command generates the LaTeX logo.
but it makes a difference to you.  When you use
\LaTeX, making your input file as easy to read
as possible will be a great help as you write
your document and when you change it.  This sample
file shows how you can add comments to your own input
file.

Because printing is different from typewriting,
there are a number of things that you have to do
differently when preparing an input file than if
you were just typing the document directly.
Quotation marks like
       ``this''
have to be handled specially, as do quotes within
quotes:
       ``\,`this'            % \, separates the double and single quote.
        is what I just
        wrote, not  `that'\,''.

Dashes come in three sizes: an
       intra-word
dash, a medium dash for number ranges like
       1--2,
and a punctuation
       dash---like
this.

A sentence-ending space should be larger than the
space between words within a sentence.  You
sometimes have to type special commands in
conjunction with punctuation characters to get
this right, as in the following sentence.
       Gnats, gnus, etc.\ all  % `\ ' makes an inter-word space.
       begin with G\@.         % \@ marks end-of-sentence punctuation.
You should check the spaces after periods when
reading your output to make sure you haven't
forgotten any special cases.  Generating an
ellipsis
       \ldots\               % `\ ' is needed after `\ldots' because TeX
                             % ignores spaces after command names like \ldots
                             % made from \ + letters.
                             %
                             % Note how a `%' character causes TeX to ignore
                             % the end of the input line, so these blank lines
                             % do not start a new paragraph.
                             %
with the right spacing around the periods requires
a special command.

\LaTeX\ interprets some common characters as
commands, so you must type special commands to
generate them.  These characters include the
following:
       \$ \& \% \# \{ and \}.

In printing, text is usually emphasized with an
       \emph{italic}
type style.

\begin{em}
   A long segment of text can also be emphasized
   in this way.  Text within such a segment can be
   given \emph{additional} emphasis.
\end{em}

It is sometimes necessary to prevent \LaTeX\ from
breaking a line where it might otherwise do so.
This may be at a space, as between the ``Mr.''\ and
``Jones'' in
       ``Mr.~Jones'',        % ~ produces an unbreakable interword space.
or within a word---especially when the word is a
symbol like
       \mbox{\emph{itemnum}}
that makes little sense when hyphenated across
lines.

Footnotes\footnote{This is an example of a footnote.}
pose no problem.

\LaTeX\ is good at typesetting mathematical formulas
like
       \( x-3y + z = 7 \)
or
       \( a_{1} > x^{2n} + y^{2n} > x' \)
or
       \( \ip{A}{B} = \sum_{i} a_{i} b_{i} \).
The spaces you type in a formula are
ignored.  Remember that a letter like
       $x$                   % $ ... $  and  \( ... \)  are equivalent
is a formula when it denotes a mathematical
symbol, and it should be typed as one.

\section{Displayed Text}

Text is displayed by indenting it from the left
margin.  Quotations are commonly displayed.  There
are short quotations
\begin{quote}
   This is a short quotation.  It consists of a
   single paragraph of text.  See how it is formatted.
\end{quote}
and longer ones.
\begin{quotation}
   This is a longer quotation.  It consists of two
   paragraphs of text, neither of which are
   particularly interesting.

   This is the second paragraph of the quotation.  It
   is just as dull as the first paragraph.
\end{quotation}
Another frequently-displayed structure is a list.
The following is an example of an \emph{itemized}
list.
\begin{itemize}
   \item This is the first item of an itemized list.
         Each item in the list is marked with a ``tick''.
         You don't have to worry about what kind of tick
         mark is used.

   \item This is the second item of the list.  It
         contains another list nested inside it.  The inner
         list is an \emph{enumerated} list.
         \begin{enumerate}
            \item This is the first item of an enumerated
                  list that is nested within the itemized list.

            \item This is the second item of the inner list.
                  \LaTeX\ allows you to nest lists deeper than
                  you really should.
         \end{enumerate}
         This is the rest of the second item of the outer
         list.  It is no more interesting than any other
         part of the item.
   \item This is the third item of the list.
\end{itemize}
You can even display poetry.
\begin{verse}
   There is an environment
    for verse \\             % The \\ command separates lines
   Whose features some poets % within a stanza.
   will curse.

                             % One or more blank lines separate stanzas.

   For instead of making\\
   Them do \emph{all} line breaking, \\
   It allows them to put too many words on a line when they'd rather be
   forced to be terse.
\end{verse}

Mathematical formulas may also be displayed.  A
displayed formula
is
one-line long; multiline
formulas require special formatting instructions.
   \[  \ip{\Gamma}{\psi'} = x'' + y^{2} + z_{i}^{n}\]
Don't start a paragraph with a displayed equation,
nor make one a paragraph by itself.

\end{document}               % End of document.

\section{Bench Huge Section 1570}
% synthetic bench section 1570
% This is a sample LaTeX input file.  (Version of 12 August 2004.)
%
% A '%' character causes TeX to ignore all remaining text on the line,
% and is used for comments like this one.

\documentclass{article}      % Specifies the document class

                             % The preamble begins here.
\title{An Example Document}  % Declares the document's title.
\author{Leslie Lamport}      % Declares the author's name.
\date{January 21, 1994}      % Deleting this command produces today's date.

\newcommand{\ip}[2]{(#1, #2)}
                             % Defines \ip{arg1}{arg2} to mean
                             % (arg1, arg2).

%\newcommand{\ip}[2]{\langle #1 | #2\rangle}
                             % This is an alternative definition of
                             % \ip that is commented out.

\begin{document}             % End of preamble and beginning of text.

\maketitle                   % Produces the title.

This is an example input file.  Comparing it with
the output it generates can show you how to
produce a simple document of your own.

\section{Ordinary Text}      % Produces section heading.  Lower-level
                             % sections are begun with similar
                             % \subsection and \subsubsection commands.

The ends  of words and sentences are marked
  by   spaces. It  doesn't matter how many
spaces    you type; one is as good as 100.  The
end of   a line counts as a space.

One   or more   blank lines denote the  end
of  a paragraph.

Since any number of consecutive spaces are treated
like a single one, the formatting of the input
file makes no difference to
      \LaTeX,                % The \LaTeX command generates the LaTeX logo.
but it makes a difference to you.  When you use
\LaTeX, making your input file as easy to read
as possible will be a great help as you write
your document and when you change it.  This sample
file shows how you can add comments to your own input
file.

Because printing is different from typewriting,
there are a number of things that you have to do
differently when preparing an input file than if
you were just typing the document directly.
Quotation marks like
       ``this''
have to be handled specially, as do quotes within
quotes:
       ``\,`this'            % \, separates the double and single quote.
        is what I just
        wrote, not  `that'\,''.

Dashes come in three sizes: an
       intra-word
dash, a medium dash for number ranges like
       1--2,
and a punctuation
       dash---like
this.

A sentence-ending space should be larger than the
space between words within a sentence.  You
sometimes have to type special commands in
conjunction with punctuation characters to get
this right, as in the following sentence.
       Gnats, gnus, etc.\ all  % `\ ' makes an inter-word space.
       begin with G\@.         % \@ marks end-of-sentence punctuation.
You should check the spaces after periods when
reading your output to make sure you haven't
forgotten any special cases.  Generating an
ellipsis
       \ldots\               % `\ ' is needed after `\ldots' because TeX
                             % ignores spaces after command names like \ldots
                             % made from \ + letters.
                             %
                             % Note how a `%' character causes TeX to ignore
                             % the end of the input line, so these blank lines
                             % do not start a new paragraph.
                             %
with the right spacing around the periods requires
a special command.

\LaTeX\ interprets some common characters as
commands, so you must type special commands to
generate them.  These characters include the
following:
       \$ \& \% \# \{ and \}.

In printing, text is usually emphasized with an
       \emph{italic}
type style.

\begin{em}
   A long segment of text can also be emphasized
   in this way.  Text within such a segment can be
   given \emph{additional} emphasis.
\end{em}

It is sometimes necessary to prevent \LaTeX\ from
breaking a line where it might otherwise do so.
This may be at a space, as between the ``Mr.''\ and
``Jones'' in
       ``Mr.~Jones'',        % ~ produces an unbreakable interword space.
or within a word---especially when the word is a
symbol like
       \mbox{\emph{itemnum}}
that makes little sense when hyphenated across
lines.

Footnotes\footnote{This is an example of a footnote.}
pose no problem.

\LaTeX\ is good at typesetting mathematical formulas
like
       \( x-3y + z = 7 \)
or
       \( a_{1} > x^{2n} + y^{2n} > x' \)
or
       \( \ip{A}{B} = \sum_{i} a_{i} b_{i} \).
The spaces you type in a formula are
ignored.  Remember that a letter like
       $x$                   % $ ... $  and  \( ... \)  are equivalent
is a formula when it denotes a mathematical
symbol, and it should be typed as one.

\section{Displayed Text}

Text is displayed by indenting it from the left
margin.  Quotations are commonly displayed.  There
are short quotations
\begin{quote}
   This is a short quotation.  It consists of a
   single paragraph of text.  See how it is formatted.
\end{quote}
and longer ones.
\begin{quotation}
   This is a longer quotation.  It consists of two
   paragraphs of text, neither of which are
   particularly interesting.

   This is the second paragraph of the quotation.  It
   is just as dull as the first paragraph.
\end{quotation}
Another frequently-displayed structure is a list.
The following is an example of an \emph{itemized}
list.
\begin{itemize}
   \item This is the first item of an itemized list.
         Each item in the list is marked with a ``tick''.
         You don't have to worry about what kind of tick
         mark is used.

   \item This is the second item of the list.  It
         contains another list nested inside it.  The inner
         list is an \emph{enumerated} list.
         \begin{enumerate}
            \item This is the first item of an enumerated
                  list that is nested within the itemized list.

            \item This is the second item of the inner list.
                  \LaTeX\ allows you to nest lists deeper than
                  you really should.
         \end{enumerate}
         This is the rest of the second item of the outer
         list.  It is no more interesting than any other
         part of the item.
   \item This is the third item of the list.
\end{itemize}
You can even display poetry.
\begin{verse}
   There is an environment
    for verse \\             % The \\ command separates lines
   Whose features some poets % within a stanza.
   will curse.

                             % One or more blank lines separate stanzas.

   For instead of making\\
   Them do \emph{all} line breaking, \\
   It allows them to put too many words on a line when they'd rather be
   forced to be terse.
\end{verse}

Mathematical formulas may also be displayed.  A
displayed formula
is
one-line long; multiline
formulas require special formatting instructions.
   \[  \ip{\Gamma}{\psi'} = x'' + y^{2} + z_{i}^{n}\]
Don't start a paragraph with a displayed equation,
nor make one a paragraph by itself.

\end{document}               % End of document.

\section{Bench Huge Section 1571}
% synthetic bench section 1571
% This is a sample LaTeX input file.  (Version of 12 August 2004.)
%
% A '%' character causes TeX to ignore all remaining text on the line,
% and is used for comments like this one.

\documentclass{article}      % Specifies the document class

                             % The preamble begins here.
\title{An Example Document}  % Declares the document's title.
\author{Leslie Lamport}      % Declares the author's name.
\date{January 21, 1994}      % Deleting this command produces today's date.

\newcommand{\ip}[2]{(#1, #2)}
                             % Defines \ip{arg1}{arg2} to mean
                             % (arg1, arg2).

%\newcommand{\ip}[2]{\langle #1 | #2\rangle}
                             % This is an alternative definition of
                             % \ip that is commented out.

\begin{document}             % End of preamble and beginning of text.

\maketitle                   % Produces the title.

This is an example input file.  Comparing it with
the output it generates can show you how to
produce a simple document of your own.

\section{Ordinary Text}      % Produces section heading.  Lower-level
                             % sections are begun with similar
                             % \subsection and \subsubsection commands.

The ends  of words and sentences are marked
  by   spaces. It  doesn't matter how many
spaces    you type; one is as good as 100.  The
end of   a line counts as a space.

One   or more   blank lines denote the  end
of  a paragraph.

Since any number of consecutive spaces are treated
like a single one, the formatting of the input
file makes no difference to
      \LaTeX,                % The \LaTeX command generates the LaTeX logo.
but it makes a difference to you.  When you use
\LaTeX, making your input file as easy to read
as possible will be a great help as you write
your document and when you change it.  This sample
file shows how you can add comments to your own input
file.

Because printing is different from typewriting,
there are a number of things that you have to do
differently when preparing an input file than if
you were just typing the document directly.
Quotation marks like
       ``this''
have to be handled specially, as do quotes within
quotes:
       ``\,`this'            % \, separates the double and single quote.
        is what I just
        wrote, not  `that'\,''.

Dashes come in three sizes: an
       intra-word
dash, a medium dash for number ranges like
       1--2,
and a punctuation
       dash---like
this.

A sentence-ending space should be larger than the
space between words within a sentence.  You
sometimes have to type special commands in
conjunction with punctuation characters to get
this right, as in the following sentence.
       Gnats, gnus, etc.\ all  % `\ ' makes an inter-word space.
       begin with G\@.         % \@ marks end-of-sentence punctuation.
You should check the spaces after periods when
reading your output to make sure you haven't
forgotten any special cases.  Generating an
ellipsis
       \ldots\               % `\ ' is needed after `\ldots' because TeX
                             % ignores spaces after command names like \ldots
                             % made from \ + letters.
                             %
                             % Note how a `%' character causes TeX to ignore
                             % the end of the input line, so these blank lines
                             % do not start a new paragraph.
                             %
with the right spacing around the periods requires
a special command.

\LaTeX\ interprets some common characters as
commands, so you must type special commands to
generate them.  These characters include the
following:
       \$ \& \% \# \{ and \}.

In printing, text is usually emphasized with an
       \emph{italic}
type style.

\begin{em}
   A long segment of text can also be emphasized
   in this way.  Text within such a segment can be
   given \emph{additional} emphasis.
\end{em}

It is sometimes necessary to prevent \LaTeX\ from
breaking a line where it might otherwise do so.
This may be at a space, as between the ``Mr.''\ and
``Jones'' in
       ``Mr.~Jones'',        % ~ produces an unbreakable interword space.
or within a word---especially when the word is a
symbol like
       \mbox{\emph{itemnum}}
that makes little sense when hyphenated across
lines.

Footnotes\footnote{This is an example of a footnote.}
pose no problem.

\LaTeX\ is good at typesetting mathematical formulas
like
       \( x-3y + z = 7 \)
or
       \( a_{1} > x^{2n} + y^{2n} > x' \)
or
       \( \ip{A}{B} = \sum_{i} a_{i} b_{i} \).
The spaces you type in a formula are
ignored.  Remember that a letter like
       $x$                   % $ ... $  and  \( ... \)  are equivalent
is a formula when it denotes a mathematical
symbol, and it should be typed as one.

\section{Displayed Text}

Text is displayed by indenting it from the left
margin.  Quotations are commonly displayed.  There
are short quotations
\begin{quote}
   This is a short quotation.  It consists of a
   single paragraph of text.  See how it is formatted.
\end{quote}
and longer ones.
\begin{quotation}
   This is a longer quotation.  It consists of two
   paragraphs of text, neither of which are
   particularly interesting.

   This is the second paragraph of the quotation.  It
   is just as dull as the first paragraph.
\end{quotation}
Another frequently-displayed structure is a list.
The following is an example of an \emph{itemized}
list.
\begin{itemize}
   \item This is the first item of an itemized list.
         Each item in the list is marked with a ``tick''.
         You don't have to worry about what kind of tick
         mark is used.

   \item This is the second item of the list.  It
         contains another list nested inside it.  The inner
         list is an \emph{enumerated} list.
         \begin{enumerate}
            \item This is the first item of an enumerated
                  list that is nested within the itemized list.

            \item This is the second item of the inner list.
                  \LaTeX\ allows you to nest lists deeper than
                  you really should.
         \end{enumerate}
         This is the rest of the second item of the outer
         list.  It is no more interesting than any other
         part of the item.
   \item This is the third item of the list.
\end{itemize}
You can even display poetry.
\begin{verse}
   There is an environment
    for verse \\             % The \\ command separates lines
   Whose features some poets % within a stanza.
   will curse.

                             % One or more blank lines separate stanzas.

   For instead of making\\
   Them do \emph{all} line breaking, \\
   It allows them to put too many words on a line when they'd rather be
   forced to be terse.
\end{verse}

Mathematical formulas may also be displayed.  A
displayed formula
is
one-line long; multiline
formulas require special formatting instructions.
   \[  \ip{\Gamma}{\psi'} = x'' + y^{2} + z_{i}^{n}\]
Don't start a paragraph with a displayed equation,
nor make one a paragraph by itself.

\end{document}               % End of document.

\section{Bench Huge Section 1572}
% synthetic bench section 1572
% This is a sample LaTeX input file.  (Version of 12 August 2004.)
%
% A '%' character causes TeX to ignore all remaining text on the line,
% and is used for comments like this one.

\documentclass{article}      % Specifies the document class

                             % The preamble begins here.
\title{An Example Document}  % Declares the document's title.
\author{Leslie Lamport}      % Declares the author's name.
\date{January 21, 1994}      % Deleting this command produces today's date.

\newcommand{\ip}[2]{(#1, #2)}
                             % Defines \ip{arg1}{arg2} to mean
                             % (arg1, arg2).

%\newcommand{\ip}[2]{\langle #1 | #2\rangle}
                             % This is an alternative definition of
                             % \ip that is commented out.

\begin{document}             % End of preamble and beginning of text.

\maketitle                   % Produces the title.

This is an example input file.  Comparing it with
the output it generates can show you how to
produce a simple document of your own.

\section{Ordinary Text}      % Produces section heading.  Lower-level
                             % sections are begun with similar
                             % \subsection and \subsubsection commands.

The ends  of words and sentences are marked
  by   spaces. It  doesn't matter how many
spaces    you type; one is as good as 100.  The
end of   a line counts as a space.

One   or more   blank lines denote the  end
of  a paragraph.

Since any number of consecutive spaces are treated
like a single one, the formatting of the input
file makes no difference to
      \LaTeX,                % The \LaTeX command generates the LaTeX logo.
but it makes a difference to you.  When you use
\LaTeX, making your input file as easy to read
as possible will be a great help as you write
your document and when you change it.  This sample
file shows how you can add comments to your own input
file.

Because printing is different from typewriting,
there are a number of things that you have to do
differently when preparing an input file than if
you were just typing the document directly.
Quotation marks like
       ``this''
have to be handled specially, as do quotes within
quotes:
       ``\,`this'            % \, separates the double and single quote.
        is what I just
        wrote, not  `that'\,''.

Dashes come in three sizes: an
       intra-word
dash, a medium dash for number ranges like
       1--2,
and a punctuation
       dash---like
this.

A sentence-ending space should be larger than the
space between words within a sentence.  You
sometimes have to type special commands in
conjunction with punctuation characters to get
this right, as in the following sentence.
       Gnats, gnus, etc.\ all  % `\ ' makes an inter-word space.
       begin with G\@.         % \@ marks end-of-sentence punctuation.
You should check the spaces after periods when
reading your output to make sure you haven't
forgotten any special cases.  Generating an
ellipsis
       \ldots\               % `\ ' is needed after `\ldots' because TeX
                             % ignores spaces after command names like \ldots
                             % made from \ + letters.
                             %
                             % Note how a `%' character causes TeX to ignore
                             % the end of the input line, so these blank lines
                             % do not start a new paragraph.
                             %
with the right spacing around the periods requires
a special command.

\LaTeX\ interprets some common characters as
commands, so you must type special commands to
generate them.  These characters include the
following:
       \$ \& \% \# \{ and \}.

In printing, text is usually emphasized with an
       \emph{italic}
type style.

\begin{em}
   A long segment of text can also be emphasized
   in this way.  Text within such a segment can be
   given \emph{additional} emphasis.
\end{em}

It is sometimes necessary to prevent \LaTeX\ from
breaking a line where it might otherwise do so.
This may be at a space, as between the ``Mr.''\ and
``Jones'' in
       ``Mr.~Jones'',        % ~ produces an unbreakable interword space.
or within a word---especially when the word is a
symbol like
       \mbox{\emph{itemnum}}
that makes little sense when hyphenated across
lines.

Footnotes\footnote{This is an example of a footnote.}
pose no problem.

\LaTeX\ is good at typesetting mathematical formulas
like
       \( x-3y + z = 7 \)
or
       \( a_{1} > x^{2n} + y^{2n} > x' \)
or
       \( \ip{A}{B} = \sum_{i} a_{i} b_{i} \).
The spaces you type in a formula are
ignored.  Remember that a letter like
       $x$                   % $ ... $  and  \( ... \)  are equivalent
is a formula when it denotes a mathematical
symbol, and it should be typed as one.

\section{Displayed Text}

Text is displayed by indenting it from the left
margin.  Quotations are commonly displayed.  There
are short quotations
\begin{quote}
   This is a short quotation.  It consists of a
   single paragraph of text.  See how it is formatted.
\end{quote}
and longer ones.
\begin{quotation}
   This is a longer quotation.  It consists of two
   paragraphs of text, neither of which are
   particularly interesting.

   This is the second paragraph of the quotation.  It
   is just as dull as the first paragraph.
\end{quotation}
Another frequently-displayed structure is a list.
The following is an example of an \emph{itemized}
list.
\begin{itemize}
   \item This is the first item of an itemized list.
         Each item in the list is marked with a ``tick''.
         You don't have to worry about what kind of tick
         mark is used.

   \item This is the second item of the list.  It
         contains another list nested inside it.  The inner
         list is an \emph{enumerated} list.
         \begin{enumerate}
            \item This is the first item of an enumerated
                  list that is nested within the itemized list.

            \item This is the second item of the inner list.
                  \LaTeX\ allows you to nest lists deeper than
                  you really should.
         \end{enumerate}
         This is the rest of the second item of the outer
         list.  It is no more interesting than any other
         part of the item.
   \item This is the third item of the list.
\end{itemize}
You can even display poetry.
\begin{verse}
   There is an environment
    for verse \\             % The \\ command separates lines
   Whose features some poets % within a stanza.
   will curse.

                             % One or more blank lines separate stanzas.

   For instead of making\\
   Them do \emph{all} line breaking, \\
   It allows them to put too many words on a line when they'd rather be
   forced to be terse.
\end{verse}

Mathematical formulas may also be displayed.  A
displayed formula
is
one-line long; multiline
formulas require special formatting instructions.
   \[  \ip{\Gamma}{\psi'} = x'' + y^{2} + z_{i}^{n}\]
Don't start a paragraph with a displayed equation,
nor make one a paragraph by itself.

\end{document}               % End of document.

\section{Bench Huge Section 1573}
% synthetic bench section 1573
% This is a sample LaTeX input file.  (Version of 12 August 2004.)
%
% A '%' character causes TeX to ignore all remaining text on the line,
% and is used for comments like this one.

\documentclass{article}      % Specifies the document class

                             % The preamble begins here.
\title{An Example Document}  % Declares the document's title.
\author{Leslie Lamport}      % Declares the author's name.
\date{January 21, 1994}      % Deleting this command produces today's date.

\newcommand{\ip}[2]{(#1, #2)}
                             % Defines \ip{arg1}{arg2} to mean
                             % (arg1, arg2).

%\newcommand{\ip}[2]{\langle #1 | #2\rangle}
                             % This is an alternative definition of
                             % \ip that is commented out.

\begin{document}             % End of preamble and beginning of text.

\maketitle                   % Produces the title.

This is an example input file.  Comparing it with
the output it generates can show you how to
produce a simple document of your own.

\section{Ordinary Text}      % Produces section heading.  Lower-level
                             % sections are begun with similar
                             % \subsection and \subsubsection commands.

The ends  of words and sentences are marked
  by   spaces. It  doesn't matter how many
spaces    you type; one is as good as 100.  The
end of   a line counts as a space.

One   or more   blank lines denote the  end
of  a paragraph.

Since any number of consecutive spaces are treated
like a single one, the formatting of the input
file makes no difference to
      \LaTeX,                % The \LaTeX command generates the LaTeX logo.
but it makes a difference to you.  When you use
\LaTeX, making your input file as easy to read
as possible will be a great help as you write
your document and when you change it.  This sample
file shows how you can add comments to your own input
file.

Because printing is different from typewriting,
there are a number of things that you have to do
differently when preparing an input file than if
you were just typing the document directly.
Quotation marks like
       ``this''
have to be handled specially, as do quotes within
quotes:
       ``\,`this'            % \, separates the double and single quote.
        is what I just
        wrote, not  `that'\,''.

Dashes come in three sizes: an
       intra-word
dash, a medium dash for number ranges like
       1--2,
and a punctuation
       dash---like
this.

A sentence-ending space should be larger than the
space between words within a sentence.  You
sometimes have to type special commands in
conjunction with punctuation characters to get
this right, as in the following sentence.
       Gnats, gnus, etc.\ all  % `\ ' makes an inter-word space.
       begin with G\@.         % \@ marks end-of-sentence punctuation.
You should check the spaces after periods when
reading your output to make sure you haven't
forgotten any special cases.  Generating an
ellipsis
       \ldots\               % `\ ' is needed after `\ldots' because TeX
                             % ignores spaces after command names like \ldots
                             % made from \ + letters.
                             %
                             % Note how a `%' character causes TeX to ignore
                             % the end of the input line, so these blank lines
                             % do not start a new paragraph.
                             %
with the right spacing around the periods requires
a special command.

\LaTeX\ interprets some common characters as
commands, so you must type special commands to
generate them.  These characters include the
following:
       \$ \& \% \# \{ and \}.

In printing, text is usually emphasized with an
       \emph{italic}
type style.

\begin{em}
   A long segment of text can also be emphasized
   in this way.  Text within such a segment can be
   given \emph{additional} emphasis.
\end{em}

It is sometimes necessary to prevent \LaTeX\ from
breaking a line where it might otherwise do so.
This may be at a space, as between the ``Mr.''\ and
``Jones'' in
       ``Mr.~Jones'',        % ~ produces an unbreakable interword space.
or within a word---especially when the word is a
symbol like
       \mbox{\emph{itemnum}}
that makes little sense when hyphenated across
lines.

Footnotes\footnote{This is an example of a footnote.}
pose no problem.

\LaTeX\ is good at typesetting mathematical formulas
like
       \( x-3y + z = 7 \)
or
       \( a_{1} > x^{2n} + y^{2n} > x' \)
or
       \( \ip{A}{B} = \sum_{i} a_{i} b_{i} \).
The spaces you type in a formula are
ignored.  Remember that a letter like
       $x$                   % $ ... $  and  \( ... \)  are equivalent
is a formula when it denotes a mathematical
symbol, and it should be typed as one.

\section{Displayed Text}

Text is displayed by indenting it from the left
margin.  Quotations are commonly displayed.  There
are short quotations
\begin{quote}
   This is a short quotation.  It consists of a
   single paragraph of text.  See how it is formatted.
\end{quote}
and longer ones.
\begin{quotation}
   This is a longer quotation.  It consists of two
   paragraphs of text, neither of which are
   particularly interesting.

   This is the second paragraph of the quotation.  It
   is just as dull as the first paragraph.
\end{quotation}
Another frequently-displayed structure is a list.
The following is an example of an \emph{itemized}
list.
\begin{itemize}
   \item This is the first item of an itemized list.
         Each item in the list is marked with a ``tick''.
         You don't have to worry about what kind of tick
         mark is used.

   \item This is the second item of the list.  It
         contains another list nested inside it.  The inner
         list is an \emph{enumerated} list.
         \begin{enumerate}
            \item This is the first item of an enumerated
                  list that is nested within the itemized list.

            \item This is the second item of the inner list.
                  \LaTeX\ allows you to nest lists deeper than
                  you really should.
         \end{enumerate}
         This is the rest of the second item of the outer
         list.  It is no more interesting than any other
         part of the item.
   \item This is the third item of the list.
\end{itemize}
You can even display poetry.
\begin{verse}
   There is an environment
    for verse \\             % The \\ command separates lines
   Whose features some poets % within a stanza.
   will curse.

                             % One or more blank lines separate stanzas.

   For instead of making\\
   Them do \emph{all} line breaking, \\
   It allows them to put too many words on a line when they'd rather be
   forced to be terse.
\end{verse}

Mathematical formulas may also be displayed.  A
displayed formula
is
one-line long; multiline
formulas require special formatting instructions.
   \[  \ip{\Gamma}{\psi'} = x'' + y^{2} + z_{i}^{n}\]
Don't start a paragraph with a displayed equation,
nor make one a paragraph by itself.

\end{document}               % End of document.

\section{Bench Huge Section 1574}
% synthetic bench section 1574
% This is a sample LaTeX input file.  (Version of 12 August 2004.)
%
% A '%' character causes TeX to ignore all remaining text on the line,
% and is used for comments like this one.

\documentclass{article}      % Specifies the document class

                             % The preamble begins here.
\title{An Example Document}  % Declares the document's title.
\author{Leslie Lamport}      % Declares the author's name.
\date{January 21, 1994}      % Deleting this command produces today's date.

\newcommand{\ip}[2]{(#1, #2)}
                             % Defines \ip{arg1}{arg2} to mean
                             % (arg1, arg2).

%\newcommand{\ip}[2]{\langle #1 | #2\rangle}
                             % This is an alternative definition of
                             % \ip that is commented out.

\begin{document}             % End of preamble and beginning of text.

\maketitle                   % Produces the title.

This is an example input file.  Comparing it with
the output it generates can show you how to
produce a simple document of your own.

\section{Ordinary Text}      % Produces section heading.  Lower-level
                             % sections are begun with similar
                             % \subsection and \subsubsection commands.

The ends  of words and sentences are marked
  by   spaces. It  doesn't matter how many
spaces    you type; one is as good as 100.  The
end of   a line counts as a space.

One   or more   blank lines denote the  end
of  a paragraph.

Since any number of consecutive spaces are treated
like a single one, the formatting of the input
file makes no difference to
      \LaTeX,                % The \LaTeX command generates the LaTeX logo.
but it makes a difference to you.  When you use
\LaTeX, making your input file as easy to read
as possible will be a great help as you write
your document and when you change it.  This sample
file shows how you can add comments to your own input
file.

Because printing is different from typewriting,
there are a number of things that you have to do
differently when preparing an input file than if
you were just typing the document directly.
Quotation marks like
       ``this''
have to be handled specially, as do quotes within
quotes:
       ``\,`this'            % \, separates the double and single quote.
        is what I just
        wrote, not  `that'\,''.

Dashes come in three sizes: an
       intra-word
dash, a medium dash for number ranges like
       1--2,
and a punctuation
       dash---like
this.

A sentence-ending space should be larger than the
space between words within a sentence.  You
sometimes have to type special commands in
conjunction with punctuation characters to get
this right, as in the following sentence.
       Gnats, gnus, etc.\ all  % `\ ' makes an inter-word space.
       begin with G\@.         % \@ marks end-of-sentence punctuation.
You should check the spaces after periods when
reading your output to make sure you haven't
forgotten any special cases.  Generating an
ellipsis
       \ldots\               % `\ ' is needed after `\ldots' because TeX
                             % ignores spaces after command names like \ldots
                             % made from \ + letters.
                             %
                             % Note how a `%' character causes TeX to ignore
                             % the end of the input line, so these blank lines
                             % do not start a new paragraph.
                             %
with the right spacing around the periods requires
a special command.

\LaTeX\ interprets some common characters as
commands, so you must type special commands to
generate them.  These characters include the
following:
       \$ \& \% \# \{ and \}.

In printing, text is usually emphasized with an
       \emph{italic}
type style.

\begin{em}
   A long segment of text can also be emphasized
   in this way.  Text within such a segment can be
   given \emph{additional} emphasis.
\end{em}

It is sometimes necessary to prevent \LaTeX\ from
breaking a line where it might otherwise do so.
This may be at a space, as between the ``Mr.''\ and
``Jones'' in
       ``Mr.~Jones'',        % ~ produces an unbreakable interword space.
or within a word---especially when the word is a
symbol like
       \mbox{\emph{itemnum}}
that makes little sense when hyphenated across
lines.

Footnotes\footnote{This is an example of a footnote.}
pose no problem.

\LaTeX\ is good at typesetting mathematical formulas
like
       \( x-3y + z = 7 \)
or
       \( a_{1} > x^{2n} + y^{2n} > x' \)
or
       \( \ip{A}{B} = \sum_{i} a_{i} b_{i} \).
The spaces you type in a formula are
ignored.  Remember that a letter like
       $x$                   % $ ... $  and  \( ... \)  are equivalent
is a formula when it denotes a mathematical
symbol, and it should be typed as one.

\section{Displayed Text}

Text is displayed by indenting it from the left
margin.  Quotations are commonly displayed.  There
are short quotations
\begin{quote}
   This is a short quotation.  It consists of a
   single paragraph of text.  See how it is formatted.
\end{quote}
and longer ones.
\begin{quotation}
   This is a longer quotation.  It consists of two
   paragraphs of text, neither of which are
   particularly interesting.

   This is the second paragraph of the quotation.  It
   is just as dull as the first paragraph.
\end{quotation}
Another frequently-displayed structure is a list.
The following is an example of an \emph{itemized}
list.
\begin{itemize}
   \item This is the first item of an itemized list.
         Each item in the list is marked with a ``tick''.
         You don't have to worry about what kind of tick
         mark is used.

   \item This is the second item of the list.  It
         contains another list nested inside it.  The inner
         list is an \emph{enumerated} list.
         \begin{enumerate}
            \item This is the first item of an enumerated
                  list that is nested within the itemized list.

            \item This is the second item of the inner list.
                  \LaTeX\ allows you to nest lists deeper than
                  you really should.
         \end{enumerate}
         This is the rest of the second item of the outer
         list.  It is no more interesting than any other
         part of the item.
   \item This is the third item of the list.
\end{itemize}
You can even display poetry.
\begin{verse}
   There is an environment
    for verse \\             % The \\ command separates lines
   Whose features some poets % within a stanza.
   will curse.

                             % One or more blank lines separate stanzas.

   For instead of making\\
   Them do \emph{all} line breaking, \\
   It allows them to put too many words on a line when they'd rather be
   forced to be terse.
\end{verse}

Mathematical formulas may also be displayed.  A
displayed formula
is
one-line long; multiline
formulas require special formatting instructions.
   \[  \ip{\Gamma}{\psi'} = x'' + y^{2} + z_{i}^{n}\]
Don't start a paragraph with a displayed equation,
nor make one a paragraph by itself.

\end{document}               % End of document.

\section{Bench Huge Section 1575}
% synthetic bench section 1575
% This is a sample LaTeX input file.  (Version of 12 August 2004.)
%
% A '%' character causes TeX to ignore all remaining text on the line,
% and is used for comments like this one.

\documentclass{article}      % Specifies the document class

                             % The preamble begins here.
\title{An Example Document}  % Declares the document's title.
\author{Leslie Lamport}      % Declares the author's name.
\date{January 21, 1994}      % Deleting this command produces today's date.

\newcommand{\ip}[2]{(#1, #2)}
                             % Defines \ip{arg1}{arg2} to mean
                             % (arg1, arg2).

%\newcommand{\ip}[2]{\langle #1 | #2\rangle}
                             % This is an alternative definition of
                             % \ip that is commented out.

\begin{document}             % End of preamble and beginning of text.

\maketitle                   % Produces the title.

This is an example input file.  Comparing it with
the output it generates can show you how to
produce a simple document of your own.

\section{Ordinary Text}      % Produces section heading.  Lower-level
                             % sections are begun with similar
                             % \subsection and \subsubsection commands.

The ends  of words and sentences are marked
  by   spaces. It  doesn't matter how many
spaces    you type; one is as good as 100.  The
end of   a line counts as a space.

One   or more   blank lines denote the  end
of  a paragraph.

Since any number of consecutive spaces are treated
like a single one, the formatting of the input
file makes no difference to
      \LaTeX,                % The \LaTeX command generates the LaTeX logo.
but it makes a difference to you.  When you use
\LaTeX, making your input file as easy to read
as possible will be a great help as you write
your document and when you change it.  This sample
file shows how you can add comments to your own input
file.

Because printing is different from typewriting,
there are a number of things that you have to do
differently when preparing an input file than if
you were just typing the document directly.
Quotation marks like
       ``this''
have to be handled specially, as do quotes within
quotes:
       ``\,`this'            % \, separates the double and single quote.
        is what I just
        wrote, not  `that'\,''.

Dashes come in three sizes: an
       intra-word
dash, a medium dash for number ranges like
       1--2,
and a punctuation
       dash---like
this.

A sentence-ending space should be larger than the
space between words within a sentence.  You
sometimes have to type special commands in
conjunction with punctuation characters to get
this right, as in the following sentence.
       Gnats, gnus, etc.\ all  % `\ ' makes an inter-word space.
       begin with G\@.         % \@ marks end-of-sentence punctuation.
You should check the spaces after periods when
reading your output to make sure you haven't
forgotten any special cases.  Generating an
ellipsis
       \ldots\               % `\ ' is needed after `\ldots' because TeX
                             % ignores spaces after command names like \ldots
                             % made from \ + letters.
                             %
                             % Note how a `%' character causes TeX to ignore
                             % the end of the input line, so these blank lines
                             % do not start a new paragraph.
                             %
with the right spacing around the periods requires
a special command.

\LaTeX\ interprets some common characters as
commands, so you must type special commands to
generate them.  These characters include the
following:
       \$ \& \% \# \{ and \}.

In printing, text is usually emphasized with an
       \emph{italic}
type style.

\begin{em}
   A long segment of text can also be emphasized
   in this way.  Text within such a segment can be
   given \emph{additional} emphasis.
\end{em}

It is sometimes necessary to prevent \LaTeX\ from
breaking a line where it might otherwise do so.
This may be at a space, as between the ``Mr.''\ and
``Jones'' in
       ``Mr.~Jones'',        % ~ produces an unbreakable interword space.
or within a word---especially when the word is a
symbol like
       \mbox{\emph{itemnum}}
that makes little sense when hyphenated across
lines.

Footnotes\footnote{This is an example of a footnote.}
pose no problem.

\LaTeX\ is good at typesetting mathematical formulas
like
       \( x-3y + z = 7 \)
or
       \( a_{1} > x^{2n} + y^{2n} > x' \)
or
       \( \ip{A}{B} = \sum_{i} a_{i} b_{i} \).
The spaces you type in a formula are
ignored.  Remember that a letter like
       $x$                   % $ ... $  and  \( ... \)  are equivalent
is a formula when it denotes a mathematical
symbol, and it should be typed as one.

\section{Displayed Text}

Text is displayed by indenting it from the left
margin.  Quotations are commonly displayed.  There
are short quotations
\begin{quote}
   This is a short quotation.  It consists of a
   single paragraph of text.  See how it is formatted.
\end{quote}
and longer ones.
\begin{quotation}
   This is a longer quotation.  It consists of two
   paragraphs of text, neither of which are
   particularly interesting.

   This is the second paragraph of the quotation.  It
   is just as dull as the first paragraph.
\end{quotation}
Another frequently-displayed structure is a list.
The following is an example of an \emph{itemized}
list.
\begin{itemize}
   \item This is the first item of an itemized list.
         Each item in the list is marked with a ``tick''.
         You don't have to worry about what kind of tick
         mark is used.

   \item This is the second item of the list.  It
         contains another list nested inside it.  The inner
         list is an \emph{enumerated} list.
         \begin{enumerate}
            \item This is the first item of an enumerated
                  list that is nested within the itemized list.

            \item This is the second item of the inner list.
                  \LaTeX\ allows you to nest lists deeper than
                  you really should.
         \end{enumerate}
         This is the rest of the second item of the outer
         list.  It is no more interesting than any other
         part of the item.
   \item This is the third item of the list.
\end{itemize}
You can even display poetry.
\begin{verse}
   There is an environment
    for verse \\             % The \\ command separates lines
   Whose features some poets % within a stanza.
   will curse.

                             % One or more blank lines separate stanzas.

   For instead of making\\
   Them do \emph{all} line breaking, \\
   It allows them to put too many words on a line when they'd rather be
   forced to be terse.
\end{verse}

Mathematical formulas may also be displayed.  A
displayed formula
is
one-line long; multiline
formulas require special formatting instructions.
   \[  \ip{\Gamma}{\psi'} = x'' + y^{2} + z_{i}^{n}\]
Don't start a paragraph with a displayed equation,
nor make one a paragraph by itself.

\end{document}               % End of document.

\section{Bench Huge Section 1576}
% synthetic bench section 1576
% This is a sample LaTeX input file.  (Version of 12 August 2004.)
%
% A '%' character causes TeX to ignore all remaining text on the line,
% and is used for comments like this one.

\documentclass{article}      % Specifies the document class

                             % The preamble begins here.
\title{An Example Document}  % Declares the document's title.
\author{Leslie Lamport}      % Declares the author's name.
\date{January 21, 1994}      % Deleting this command produces today's date.

\newcommand{\ip}[2]{(#1, #2)}
                             % Defines \ip{arg1}{arg2} to mean
                             % (arg1, arg2).

%\newcommand{\ip}[2]{\langle #1 | #2\rangle}
                             % This is an alternative definition of
                             % \ip that is commented out.

\begin{document}             % End of preamble and beginning of text.

\maketitle                   % Produces the title.

This is an example input file.  Comparing it with
the output it generates can show you how to
produce a simple document of your own.

\section{Ordinary Text}      % Produces section heading.  Lower-level
                             % sections are begun with similar
                             % \subsection and \subsubsection commands.

The ends  of words and sentences are marked
  by   spaces. It  doesn't matter how many
spaces    you type; one is as good as 100.  The
end of   a line counts as a space.

One   or more   blank lines denote the  end
of  a paragraph.

Since any number of consecutive spaces are treated
like a single one, the formatting of the input
file makes no difference to
      \LaTeX,                % The \LaTeX command generates the LaTeX logo.
but it makes a difference to you.  When you use
\LaTeX, making your input file as easy to read
as possible will be a great help as you write
your document and when you change it.  This sample
file shows how you can add comments to your own input
file.

Because printing is different from typewriting,
there are a number of things that you have to do
differently when preparing an input file than if
you were just typing the document directly.
Quotation marks like
       ``this''
have to be handled specially, as do quotes within
quotes:
       ``\,`this'            % \, separates the double and single quote.
        is what I just
        wrote, not  `that'\,''.

Dashes come in three sizes: an
       intra-word
dash, a medium dash for number ranges like
       1--2,
and a punctuation
       dash---like
this.

A sentence-ending space should be larger than the
space between words within a sentence.  You
sometimes have to type special commands in
conjunction with punctuation characters to get
this right, as in the following sentence.
       Gnats, gnus, etc.\ all  % `\ ' makes an inter-word space.
       begin with G\@.         % \@ marks end-of-sentence punctuation.
You should check the spaces after periods when
reading your output to make sure you haven't
forgotten any special cases.  Generating an
ellipsis
       \ldots\               % `\ ' is needed after `\ldots' because TeX
                             % ignores spaces after command names like \ldots
                             % made from \ + letters.
                             %
                             % Note how a `%' character causes TeX to ignore
                             % the end of the input line, so these blank lines
                             % do not start a new paragraph.
                             %
with the right spacing around the periods requires
a special command.

\LaTeX\ interprets some common characters as
commands, so you must type special commands to
generate them.  These characters include the
following:
       \$ \& \% \# \{ and \}.

In printing, text is usually emphasized with an
       \emph{italic}
type style.

\begin{em}
   A long segment of text can also be emphasized
   in this way.  Text within such a segment can be
   given \emph{additional} emphasis.
\end{em}

It is sometimes necessary to prevent \LaTeX\ from
breaking a line where it might otherwise do so.
This may be at a space, as between the ``Mr.''\ and
``Jones'' in
       ``Mr.~Jones'',        % ~ produces an unbreakable interword space.
or within a word---especially when the word is a
symbol like
       \mbox{\emph{itemnum}}
that makes little sense when hyphenated across
lines.

Footnotes\footnote{This is an example of a footnote.}
pose no problem.

\LaTeX\ is good at typesetting mathematical formulas
like
       \( x-3y + z = 7 \)
or
       \( a_{1} > x^{2n} + y^{2n} > x' \)
or
       \( \ip{A}{B} = \sum_{i} a_{i} b_{i} \).
The spaces you type in a formula are
ignored.  Remember that a letter like
       $x$                   % $ ... $  and  \( ... \)  are equivalent
is a formula when it denotes a mathematical
symbol, and it should be typed as one.

\section{Displayed Text}

Text is displayed by indenting it from the left
margin.  Quotations are commonly displayed.  There
are short quotations
\begin{quote}
   This is a short quotation.  It consists of a
   single paragraph of text.  See how it is formatted.
\end{quote}
and longer ones.
\begin{quotation}
   This is a longer quotation.  It consists of two
   paragraphs of text, neither of which are
   particularly interesting.

   This is the second paragraph of the quotation.  It
   is just as dull as the first paragraph.
\end{quotation}
Another frequently-displayed structure is a list.
The following is an example of an \emph{itemized}
list.
\begin{itemize}
   \item This is the first item of an itemized list.
         Each item in the list is marked with a ``tick''.
         You don't have to worry about what kind of tick
         mark is used.

   \item This is the second item of the list.  It
         contains another list nested inside it.  The inner
         list is an \emph{enumerated} list.
         \begin{enumerate}
            \item This is the first item of an enumerated
                  list that is nested within the itemized list.

            \item This is the second item of the inner list.
                  \LaTeX\ allows you to nest lists deeper than
                  you really should.
         \end{enumerate}
         This is the rest of the second item of the outer
         list.  It is no more interesting than any other
         part of the item.
   \item This is the third item of the list.
\end{itemize}
You can even display poetry.
\begin{verse}
   There is an environment
    for verse \\             % The \\ command separates lines
   Whose features some poets % within a stanza.
   will curse.

                             % One or more blank lines separate stanzas.

   For instead of making\\
   Them do \emph{all} line breaking, \\
   It allows them to put too many words on a line when they'd rather be
   forced to be terse.
\end{verse}

Mathematical formulas may also be displayed.  A
displayed formula
is
one-line long; multiline
formulas require special formatting instructions.
   \[  \ip{\Gamma}{\psi'} = x'' + y^{2} + z_{i}^{n}\]
Don't start a paragraph with a displayed equation,
nor make one a paragraph by itself.

\end{document}               % End of document.

\section{Bench Huge Section 1577}
% synthetic bench section 1577
% This is a sample LaTeX input file.  (Version of 12 August 2004.)
%
% A '%' character causes TeX to ignore all remaining text on the line,
% and is used for comments like this one.

\documentclass{article}      % Specifies the document class

                             % The preamble begins here.
\title{An Example Document}  % Declares the document's title.
\author{Leslie Lamport}      % Declares the author's name.
\date{January 21, 1994}      % Deleting this command produces today's date.

\newcommand{\ip}[2]{(#1, #2)}
                             % Defines \ip{arg1}{arg2} to mean
                             % (arg1, arg2).

%\newcommand{\ip}[2]{\langle #1 | #2\rangle}
                             % This is an alternative definition of
                             % \ip that is commented out.

\begin{document}             % End of preamble and beginning of text.

\maketitle                   % Produces the title.

This is an example input file.  Comparing it with
the output it generates can show you how to
produce a simple document of your own.

\section{Ordinary Text}      % Produces section heading.  Lower-level
                             % sections are begun with similar
                             % \subsection and \subsubsection commands.

The ends  of words and sentences are marked
  by   spaces. It  doesn't matter how many
spaces    you type; one is as good as 100.  The
end of   a line counts as a space.

One   or more   blank lines denote the  end
of  a paragraph.

Since any number of consecutive spaces are treated
like a single one, the formatting of the input
file makes no difference to
      \LaTeX,                % The \LaTeX command generates the LaTeX logo.
but it makes a difference to you.  When you use
\LaTeX, making your input file as easy to read
as possible will be a great help as you write
your document and when you change it.  This sample
file shows how you can add comments to your own input
file.

Because printing is different from typewriting,
there are a number of things that you have to do
differently when preparing an input file than if
you were just typing the document directly.
Quotation marks like
       ``this''
have to be handled specially, as do quotes within
quotes:
       ``\,`this'            % \, separates the double and single quote.
        is what I just
        wrote, not  `that'\,''.

Dashes come in three sizes: an
       intra-word
dash, a medium dash for number ranges like
       1--2,
and a punctuation
       dash---like
this.

A sentence-ending space should be larger than the
space between words within a sentence.  You
sometimes have to type special commands in
conjunction with punctuation characters to get
this right, as in the following sentence.
       Gnats, gnus, etc.\ all  % `\ ' makes an inter-word space.
       begin with G\@.         % \@ marks end-of-sentence punctuation.
You should check the spaces after periods when
reading your output to make sure you haven't
forgotten any special cases.  Generating an
ellipsis
       \ldots\               % `\ ' is needed after `\ldots' because TeX
                             % ignores spaces after command names like \ldots
                             % made from \ + letters.
                             %
                             % Note how a `%' character causes TeX to ignore
                             % the end of the input line, so these blank lines
                             % do not start a new paragraph.
                             %
with the right spacing around the periods requires
a special command.

\LaTeX\ interprets some common characters as
commands, so you must type special commands to
generate them.  These characters include the
following:
       \$ \& \% \# \{ and \}.

In printing, text is usually emphasized with an
       \emph{italic}
type style.

\begin{em}
   A long segment of text can also be emphasized
   in this way.  Text within such a segment can be
   given \emph{additional} emphasis.
\end{em}

It is sometimes necessary to prevent \LaTeX\ from
breaking a line where it might otherwise do so.
This may be at a space, as between the ``Mr.''\ and
``Jones'' in
       ``Mr.~Jones'',        % ~ produces an unbreakable interword space.
or within a word---especially when the word is a
symbol like
       \mbox{\emph{itemnum}}
that makes little sense when hyphenated across
lines.

Footnotes\footnote{This is an example of a footnote.}
pose no problem.

\LaTeX\ is good at typesetting mathematical formulas
like
       \( x-3y + z = 7 \)
or
       \( a_{1} > x^{2n} + y^{2n} > x' \)
or
       \( \ip{A}{B} = \sum_{i} a_{i} b_{i} \).
The spaces you type in a formula are
ignored.  Remember that a letter like
       $x$                   % $ ... $  and  \( ... \)  are equivalent
is a formula when it denotes a mathematical
symbol, and it should be typed as one.

\section{Displayed Text}

Text is displayed by indenting it from the left
margin.  Quotations are commonly displayed.  There
are short quotations
\begin{quote}
   This is a short quotation.  It consists of a
   single paragraph of text.  See how it is formatted.
\end{quote}
and longer ones.
\begin{quotation}
   This is a longer quotation.  It consists of two
   paragraphs of text, neither of which are
   particularly interesting.

   This is the second paragraph of the quotation.  It
   is just as dull as the first paragraph.
\end{quotation}
Another frequently-displayed structure is a list.
The following is an example of an \emph{itemized}
list.
\begin{itemize}
   \item This is the first item of an itemized list.
         Each item in the list is marked with a ``tick''.
         You don't have to worry about what kind of tick
         mark is used.

   \item This is the second item of the list.  It
         contains another list nested inside it.  The inner
         list is an \emph{enumerated} list.
         \begin{enumerate}
            \item This is the first item of an enumerated
                  list that is nested within the itemized list.

            \item This is the second item of the inner list.
                  \LaTeX\ allows you to nest lists deeper than
                  you really should.
         \end{enumerate}
         This is the rest of the second item of the outer
         list.  It is no more interesting than any other
         part of the item.
   \item This is the third item of the list.
\end{itemize}
You can even display poetry.
\begin{verse}
   There is an environment
    for verse \\             % The \\ command separates lines
   Whose features some poets % within a stanza.
   will curse.

                             % One or more blank lines separate stanzas.

   For instead of making\\
   Them do \emph{all} line breaking, \\
   It allows them to put too many words on a line when they'd rather be
   forced to be terse.
\end{verse}

Mathematical formulas may also be displayed.  A
displayed formula
is
one-line long; multiline
formulas require special formatting instructions.
   \[  \ip{\Gamma}{\psi'} = x'' + y^{2} + z_{i}^{n}\]
Don't start a paragraph with a displayed equation,
nor make one a paragraph by itself.

\end{document}               % End of document.

\section{Bench Huge Section 1578}
% synthetic bench section 1578
% This is a sample LaTeX input file.  (Version of 12 August 2004.)
%
% A '%' character causes TeX to ignore all remaining text on the line,
% and is used for comments like this one.

\documentclass{article}      % Specifies the document class

                             % The preamble begins here.
\title{An Example Document}  % Declares the document's title.
\author{Leslie Lamport}      % Declares the author's name.
\date{January 21, 1994}      % Deleting this command produces today's date.

\newcommand{\ip}[2]{(#1, #2)}
                             % Defines \ip{arg1}{arg2} to mean
                             % (arg1, arg2).

%\newcommand{\ip}[2]{\langle #1 | #2\rangle}
                             % This is an alternative definition of
                             % \ip that is commented out.

\begin{document}             % End of preamble and beginning of text.

\maketitle                   % Produces the title.

This is an example input file.  Comparing it with
the output it generates can show you how to
produce a simple document of your own.

\section{Ordinary Text}      % Produces section heading.  Lower-level
                             % sections are begun with similar
                             % \subsection and \subsubsection commands.

The ends  of words and sentences are marked
  by   spaces. It  doesn't matter how many
spaces    you type; one is as good as 100.  The
end of   a line counts as a space.

One   or more   blank lines denote the  end
of  a paragraph.

Since any number of consecutive spaces are treated
like a single one, the formatting of the input
file makes no difference to
      \LaTeX,                % The \LaTeX command generates the LaTeX logo.
but it makes a difference to you.  When you use
\LaTeX, making your input file as easy to read
as possible will be a great help as you write
your document and when you change it.  This sample
file shows how you can add comments to your own input
file.

Because printing is different from typewriting,
there are a number of things that you have to do
differently when preparing an input file than if
you were just typing the document directly.
Quotation marks like
       ``this''
have to be handled specially, as do quotes within
quotes:
       ``\,`this'            % \, separates the double and single quote.
        is what I just
        wrote, not  `that'\,''.

Dashes come in three sizes: an
       intra-word
dash, a medium dash for number ranges like
       1--2,
and a punctuation
       dash---like
this.

A sentence-ending space should be larger than the
space between words within a sentence.  You
sometimes have to type special commands in
conjunction with punctuation characters to get
this right, as in the following sentence.
       Gnats, gnus, etc.\ all  % `\ ' makes an inter-word space.
       begin with G\@.         % \@ marks end-of-sentence punctuation.
You should check the spaces after periods when
reading your output to make sure you haven't
forgotten any special cases.  Generating an
ellipsis
       \ldots\               % `\ ' is needed after `\ldots' because TeX
                             % ignores spaces after command names like \ldots
                             % made from \ + letters.
                             %
                             % Note how a `%' character causes TeX to ignore
                             % the end of the input line, so these blank lines
                             % do not start a new paragraph.
                             %
with the right spacing around the periods requires
a special command.

\LaTeX\ interprets some common characters as
commands, so you must type special commands to
generate them.  These characters include the
following:
       \$ \& \% \# \{ and \}.

In printing, text is usually emphasized with an
       \emph{italic}
type style.

\begin{em}
   A long segment of text can also be emphasized
   in this way.  Text within such a segment can be
   given \emph{additional} emphasis.
\end{em}

It is sometimes necessary to prevent \LaTeX\ from
breaking a line where it might otherwise do so.
This may be at a space, as between the ``Mr.''\ and
``Jones'' in
       ``Mr.~Jones'',        % ~ produces an unbreakable interword space.
or within a word---especially when the word is a
symbol like
       \mbox{\emph{itemnum}}
that makes little sense when hyphenated across
lines.

Footnotes\footnote{This is an example of a footnote.}
pose no problem.

\LaTeX\ is good at typesetting mathematical formulas
like
       \( x-3y + z = 7 \)
or
       \( a_{1} > x^{2n} + y^{2n} > x' \)
or
       \( \ip{A}{B} = \sum_{i} a_{i} b_{i} \).
The spaces you type in a formula are
ignored.  Remember that a letter like
       $x$                   % $ ... $  and  \( ... \)  are equivalent
is a formula when it denotes a mathematical
symbol, and it should be typed as one.

\section{Displayed Text}

Text is displayed by indenting it from the left
margin.  Quotations are commonly displayed.  There
are short quotations
\begin{quote}
   This is a short quotation.  It consists of a
   single paragraph of text.  See how it is formatted.
\end{quote}
and longer ones.
\begin{quotation}
   This is a longer quotation.  It consists of two
   paragraphs of text, neither of which are
   particularly interesting.

   This is the second paragraph of the quotation.  It
   is just as dull as the first paragraph.
\end{quotation}
Another frequently-displayed structure is a list.
The following is an example of an \emph{itemized}
list.
\begin{itemize}
   \item This is the first item of an itemized list.
         Each item in the list is marked with a ``tick''.
         You don't have to worry about what kind of tick
         mark is used.

   \item This is the second item of the list.  It
         contains another list nested inside it.  The inner
         list is an \emph{enumerated} list.
         \begin{enumerate}
            \item This is the first item of an enumerated
                  list that is nested within the itemized list.

            \item This is the second item of the inner list.
                  \LaTeX\ allows you to nest lists deeper than
                  you really should.
         \end{enumerate}
         This is the rest of the second item of the outer
         list.  It is no more interesting than any other
         part of the item.
   \item This is the third item of the list.
\end{itemize}
You can even display poetry.
\begin{verse}
   There is an environment
    for verse \\             % The \\ command separates lines
   Whose features some poets % within a stanza.
   will curse.

                             % One or more blank lines separate stanzas.

   For instead of making\\
   Them do \emph{all} line breaking, \\
   It allows them to put too many words on a line when they'd rather be
   forced to be terse.
\end{verse}

Mathematical formulas may also be displayed.  A
displayed formula
is
one-line long; multiline
formulas require special formatting instructions.
   \[  \ip{\Gamma}{\psi'} = x'' + y^{2} + z_{i}^{n}\]
Don't start a paragraph with a displayed equation,
nor make one a paragraph by itself.

\end{document}               % End of document.

\section{Bench Huge Section 1579}
% synthetic bench section 1579
% This is a sample LaTeX input file.  (Version of 12 August 2004.)
%
% A '%' character causes TeX to ignore all remaining text on the line,
% and is used for comments like this one.

\documentclass{article}      % Specifies the document class

                             % The preamble begins here.
\title{An Example Document}  % Declares the document's title.
\author{Leslie Lamport}      % Declares the author's name.
\date{January 21, 1994}      % Deleting this command produces today's date.

\newcommand{\ip}[2]{(#1, #2)}
                             % Defines \ip{arg1}{arg2} to mean
                             % (arg1, arg2).

%\newcommand{\ip}[2]{\langle #1 | #2\rangle}
                             % This is an alternative definition of
                             % \ip that is commented out.

\begin{document}             % End of preamble and beginning of text.

\maketitle                   % Produces the title.

This is an example input file.  Comparing it with
the output it generates can show you how to
produce a simple document of your own.

\section{Ordinary Text}      % Produces section heading.  Lower-level
                             % sections are begun with similar
                             % \subsection and \subsubsection commands.

The ends  of words and sentences are marked
  by   spaces. It  doesn't matter how many
spaces    you type; one is as good as 100.  The
end of   a line counts as a space.

One   or more   blank lines denote the  end
of  a paragraph.

Since any number of consecutive spaces are treated
like a single one, the formatting of the input
file makes no difference to
      \LaTeX,                % The \LaTeX command generates the LaTeX logo.
but it makes a difference to you.  When you use
\LaTeX, making your input file as easy to read
as possible will be a great help as you write
your document and when you change it.  This sample
file shows how you can add comments to your own input
file.

Because printing is different from typewriting,
there are a number of things that you have to do
differently when preparing an input file than if
you were just typing the document directly.
Quotation marks like
       ``this''
have to be handled specially, as do quotes within
quotes:
       ``\,`this'            % \, separates the double and single quote.
        is what I just
        wrote, not  `that'\,''.

Dashes come in three sizes: an
       intra-word
dash, a medium dash for number ranges like
       1--2,
and a punctuation
       dash---like
this.

A sentence-ending space should be larger than the
space between words within a sentence.  You
sometimes have to type special commands in
conjunction with punctuation characters to get
this right, as in the following sentence.
       Gnats, gnus, etc.\ all  % `\ ' makes an inter-word space.
       begin with G\@.         % \@ marks end-of-sentence punctuation.
You should check the spaces after periods when
reading your output to make sure you haven't
forgotten any special cases.  Generating an
ellipsis
       \ldots\               % `\ ' is needed after `\ldots' because TeX
                             % ignores spaces after command names like \ldots
                             % made from \ + letters.
                             %
                             % Note how a `%' character causes TeX to ignore
                             % the end of the input line, so these blank lines
                             % do not start a new paragraph.
                             %
with the right spacing around the periods requires
a special command.

\LaTeX\ interprets some common characters as
commands, so you must type special commands to
generate them.  These characters include the
following:
       \$ \& \% \# \{ and \}.

In printing, text is usually emphasized with an
       \emph{italic}
type style.

\begin{em}
   A long segment of text can also be emphasized
   in this way.  Text within such a segment can be
   given \emph{additional} emphasis.
\end{em}

It is sometimes necessary to prevent \LaTeX\ from
breaking a line where it might otherwise do so.
This may be at a space, as between the ``Mr.''\ and
``Jones'' in
       ``Mr.~Jones'',        % ~ produces an unbreakable interword space.
or within a word---especially when the word is a
symbol like
       \mbox{\emph{itemnum}}
that makes little sense when hyphenated across
lines.

Footnotes\footnote{This is an example of a footnote.}
pose no problem.

\LaTeX\ is good at typesetting mathematical formulas
like
       \( x-3y + z = 7 \)
or
       \( a_{1} > x^{2n} + y^{2n} > x' \)
or
       \( \ip{A}{B} = \sum_{i} a_{i} b_{i} \).
The spaces you type in a formula are
ignored.  Remember that a letter like
       $x$                   % $ ... $  and  \( ... \)  are equivalent
is a formula when it denotes a mathematical
symbol, and it should be typed as one.

\section{Displayed Text}

Text is displayed by indenting it from the left
margin.  Quotations are commonly displayed.  There
are short quotations
\begin{quote}
   This is a short quotation.  It consists of a
   single paragraph of text.  See how it is formatted.
\end{quote}
and longer ones.
\begin{quotation}
   This is a longer quotation.  It consists of two
   paragraphs of text, neither of which are
   particularly interesting.

   This is the second paragraph of the quotation.  It
   is just as dull as the first paragraph.
\end{quotation}
Another frequently-displayed structure is a list.
The following is an example of an \emph{itemized}
list.
\begin{itemize}
   \item This is the first item of an itemized list.
         Each item in the list is marked with a ``tick''.
         You don't have to worry about what kind of tick
         mark is used.

   \item This is the second item of the list.  It
         contains another list nested inside it.  The inner
         list is an \emph{enumerated} list.
         \begin{enumerate}
            \item This is the first item of an enumerated
                  list that is nested within the itemized list.

            \item This is the second item of the inner list.
                  \LaTeX\ allows you to nest lists deeper than
                  you really should.
         \end{enumerate}
         This is the rest of the second item of the outer
         list.  It is no more interesting than any other
         part of the item.
   \item This is the third item of the list.
\end{itemize}
You can even display poetry.
\begin{verse}
   There is an environment
    for verse \\             % The \\ command separates lines
   Whose features some poets % within a stanza.
   will curse.

                             % One or more blank lines separate stanzas.

   For instead of making\\
   Them do \emph{all} line breaking, \\
   It allows them to put too many words on a line when they'd rather be
   forced to be terse.
\end{verse}

Mathematical formulas may also be displayed.  A
displayed formula
is
one-line long; multiline
formulas require special formatting instructions.
   \[  \ip{\Gamma}{\psi'} = x'' + y^{2} + z_{i}^{n}\]
Don't start a paragraph with a displayed equation,
nor make one a paragraph by itself.

\end{document}               % End of document.

\section{Bench Huge Section 1580}
% synthetic bench section 1580
% This is a sample LaTeX input file.  (Version of 12 August 2004.)
%
% A '%' character causes TeX to ignore all remaining text on the line,
% and is used for comments like this one.

\documentclass{article}      % Specifies the document class

                             % The preamble begins here.
\title{An Example Document}  % Declares the document's title.
\author{Leslie Lamport}      % Declares the author's name.
\date{January 21, 1994}      % Deleting this command produces today's date.

\newcommand{\ip}[2]{(#1, #2)}
                             % Defines \ip{arg1}{arg2} to mean
                             % (arg1, arg2).

%\newcommand{\ip}[2]{\langle #1 | #2\rangle}
                             % This is an alternative definition of
                             % \ip that is commented out.

\begin{document}             % End of preamble and beginning of text.

\maketitle                   % Produces the title.

This is an example input file.  Comparing it with
the output it generates can show you how to
produce a simple document of your own.

\section{Ordinary Text}      % Produces section heading.  Lower-level
                             % sections are begun with similar
                             % \subsection and \subsubsection commands.

The ends  of words and sentences are marked
  by   spaces. It  doesn't matter how many
spaces    you type; one is as good as 100.  The
end of   a line counts as a space.

One   or more   blank lines denote the  end
of  a paragraph.

Since any number of consecutive spaces are treated
like a single one, the formatting of the input
file makes no difference to
      \LaTeX,                % The \LaTeX command generates the LaTeX logo.
but it makes a difference to you.  When you use
\LaTeX, making your input file as easy to read
as possible will be a great help as you write
your document and when you change it.  This sample
file shows how you can add comments to your own input
file.

Because printing is different from typewriting,
there are a number of things that you have to do
differently when preparing an input file than if
you were just typing the document directly.
Quotation marks like
       ``this''
have to be handled specially, as do quotes within
quotes:
       ``\,`this'            % \, separates the double and single quote.
        is what I just
        wrote, not  `that'\,''.

Dashes come in three sizes: an
       intra-word
dash, a medium dash for number ranges like
       1--2,
and a punctuation
       dash---like
this.

A sentence-ending space should be larger than the
space between words within a sentence.  You
sometimes have to type special commands in
conjunction with punctuation characters to get
this right, as in the following sentence.
       Gnats, gnus, etc.\ all  % `\ ' makes an inter-word space.
       begin with G\@.         % \@ marks end-of-sentence punctuation.
You should check the spaces after periods when
reading your output to make sure you haven't
forgotten any special cases.  Generating an
ellipsis
       \ldots\               % `\ ' is needed after `\ldots' because TeX
                             % ignores spaces after command names like \ldots
                             % made from \ + letters.
                             %
                             % Note how a `%' character causes TeX to ignore
                             % the end of the input line, so these blank lines
                             % do not start a new paragraph.
                             %
with the right spacing around the periods requires
a special command.

\LaTeX\ interprets some common characters as
commands, so you must type special commands to
generate them.  These characters include the
following:
       \$ \& \% \# \{ and \}.

In printing, text is usually emphasized with an
       \emph{italic}
type style.

\begin{em}
   A long segment of text can also be emphasized
   in this way.  Text within such a segment can be
   given \emph{additional} emphasis.
\end{em}

It is sometimes necessary to prevent \LaTeX\ from
breaking a line where it might otherwise do so.
This may be at a space, as between the ``Mr.''\ and
``Jones'' in
       ``Mr.~Jones'',        % ~ produces an unbreakable interword space.
or within a word---especially when the word is a
symbol like
       \mbox{\emph{itemnum}}
that makes little sense when hyphenated across
lines.

Footnotes\footnote{This is an example of a footnote.}
pose no problem.

\LaTeX\ is good at typesetting mathematical formulas
like
       \( x-3y + z = 7 \)
or
       \( a_{1} > x^{2n} + y^{2n} > x' \)
or
       \( \ip{A}{B} = \sum_{i} a_{i} b_{i} \).
The spaces you type in a formula are
ignored.  Remember that a letter like
       $x$                   % $ ... $  and  \( ... \)  are equivalent
is a formula when it denotes a mathematical
symbol, and it should be typed as one.

\section{Displayed Text}

Text is displayed by indenting it from the left
margin.  Quotations are commonly displayed.  There
are short quotations
\begin{quote}
   This is a short quotation.  It consists of a
   single paragraph of text.  See how it is formatted.
\end{quote}
and longer ones.
\begin{quotation}
   This is a longer quotation.  It consists of two
   paragraphs of text, neither of which are
   particularly interesting.

   This is the second paragraph of the quotation.  It
   is just as dull as the first paragraph.
\end{quotation}
Another frequently-displayed structure is a list.
The following is an example of an \emph{itemized}
list.
\begin{itemize}
   \item This is the first item of an itemized list.
         Each item in the list is marked with a ``tick''.
         You don't have to worry about what kind of tick
         mark is used.

   \item This is the second item of the list.  It
         contains another list nested inside it.  The inner
         list is an \emph{enumerated} list.
         \begin{enumerate}
            \item This is the first item of an enumerated
                  list that is nested within the itemized list.

            \item This is the second item of the inner list.
                  \LaTeX\ allows you to nest lists deeper than
                  you really should.
         \end{enumerate}
         This is the rest of the second item of the outer
         list.  It is no more interesting than any other
         part of the item.
   \item This is the third item of the list.
\end{itemize}
You can even display poetry.
\begin{verse}
   There is an environment
    for verse \\             % The \\ command separates lines
   Whose features some poets % within a stanza.
   will curse.

                             % One or more blank lines separate stanzas.

   For instead of making\\
   Them do \emph{all} line breaking, \\
   It allows them to put too many words on a line when they'd rather be
   forced to be terse.
\end{verse}

Mathematical formulas may also be displayed.  A
displayed formula
is
one-line long; multiline
formulas require special formatting instructions.
   \[  \ip{\Gamma}{\psi'} = x'' + y^{2} + z_{i}^{n}\]
Don't start a paragraph with a displayed equation,
nor make one a paragraph by itself.

\end{document}               % End of document.

\section{Bench Huge Section 1581}
% synthetic bench section 1581
% This is a sample LaTeX input file.  (Version of 12 August 2004.)
%
% A '%' character causes TeX to ignore all remaining text on the line,
% and is used for comments like this one.

\documentclass{article}      % Specifies the document class

                             % The preamble begins here.
\title{An Example Document}  % Declares the document's title.
\author{Leslie Lamport}      % Declares the author's name.
\date{January 21, 1994}      % Deleting this command produces today's date.

\newcommand{\ip}[2]{(#1, #2)}
                             % Defines \ip{arg1}{arg2} to mean
                             % (arg1, arg2).

%\newcommand{\ip}[2]{\langle #1 | #2\rangle}
                             % This is an alternative definition of
                             % \ip that is commented out.

\begin{document}             % End of preamble and beginning of text.

\maketitle                   % Produces the title.

This is an example input file.  Comparing it with
the output it generates can show you how to
produce a simple document of your own.

\section{Ordinary Text}      % Produces section heading.  Lower-level
                             % sections are begun with similar
                             % \subsection and \subsubsection commands.

The ends  of words and sentences are marked
  by   spaces. It  doesn't matter how many
spaces    you type; one is as good as 100.  The
end of   a line counts as a space.

One   or more   blank lines denote the  end
of  a paragraph.

Since any number of consecutive spaces are treated
like a single one, the formatting of the input
file makes no difference to
      \LaTeX,                % The \LaTeX command generates the LaTeX logo.
but it makes a difference to you.  When you use
\LaTeX, making your input file as easy to read
as possible will be a great help as you write
your document and when you change it.  This sample
file shows how you can add comments to your own input
file.

Because printing is different from typewriting,
there are a number of things that you have to do
differently when preparing an input file than if
you were just typing the document directly.
Quotation marks like
       ``this''
have to be handled specially, as do quotes within
quotes:
       ``\,`this'            % \, separates the double and single quote.
        is what I just
        wrote, not  `that'\,''.

Dashes come in three sizes: an
       intra-word
dash, a medium dash for number ranges like
       1--2,
and a punctuation
       dash---like
this.

A sentence-ending space should be larger than the
space between words within a sentence.  You
sometimes have to type special commands in
conjunction with punctuation characters to get
this right, as in the following sentence.
       Gnats, gnus, etc.\ all  % `\ ' makes an inter-word space.
       begin with G\@.         % \@ marks end-of-sentence punctuation.
You should check the spaces after periods when
reading your output to make sure you haven't
forgotten any special cases.  Generating an
ellipsis
       \ldots\               % `\ ' is needed after `\ldots' because TeX
                             % ignores spaces after command names like \ldots
                             % made from \ + letters.
                             %
                             % Note how a `%' character causes TeX to ignore
                             % the end of the input line, so these blank lines
                             % do not start a new paragraph.
                             %
with the right spacing around the periods requires
a special command.

\LaTeX\ interprets some common characters as
commands, so you must type special commands to
generate them.  These characters include the
following:
       \$ \& \% \# \{ and \}.

In printing, text is usually emphasized with an
       \emph{italic}
type style.

\begin{em}
   A long segment of text can also be emphasized
   in this way.  Text within such a segment can be
   given \emph{additional} emphasis.
\end{em}

It is sometimes necessary to prevent \LaTeX\ from
breaking a line where it might otherwise do so.
This may be at a space, as between the ``Mr.''\ and
``Jones'' in
       ``Mr.~Jones'',        % ~ produces an unbreakable interword space.
or within a word---especially when the word is a
symbol like
       \mbox{\emph{itemnum}}
that makes little sense when hyphenated across
lines.

Footnotes\footnote{This is an example of a footnote.}
pose no problem.

\LaTeX\ is good at typesetting mathematical formulas
like
       \( x-3y + z = 7 \)
or
       \( a_{1} > x^{2n} + y^{2n} > x' \)
or
       \( \ip{A}{B} = \sum_{i} a_{i} b_{i} \).
The spaces you type in a formula are
ignored.  Remember that a letter like
       $x$                   % $ ... $  and  \( ... \)  are equivalent
is a formula when it denotes a mathematical
symbol, and it should be typed as one.

\section{Displayed Text}

Text is displayed by indenting it from the left
margin.  Quotations are commonly displayed.  There
are short quotations
\begin{quote}
   This is a short quotation.  It consists of a
   single paragraph of text.  See how it is formatted.
\end{quote}
and longer ones.
\begin{quotation}
   This is a longer quotation.  It consists of two
   paragraphs of text, neither of which are
   particularly interesting.

   This is the second paragraph of the quotation.  It
   is just as dull as the first paragraph.
\end{quotation}
Another frequently-displayed structure is a list.
The following is an example of an \emph{itemized}
list.
\begin{itemize}
   \item This is the first item of an itemized list.
         Each item in the list is marked with a ``tick''.
         You don't have to worry about what kind of tick
         mark is used.

   \item This is the second item of the list.  It
         contains another list nested inside it.  The inner
         list is an \emph{enumerated} list.
         \begin{enumerate}
            \item This is the first item of an enumerated
                  list that is nested within the itemized list.

            \item This is the second item of the inner list.
                  \LaTeX\ allows you to nest lists deeper than
                  you really should.
         \end{enumerate}
         This is the rest of the second item of the outer
         list.  It is no more interesting than any other
         part of the item.
   \item This is the third item of the list.
\end{itemize}
You can even display poetry.
\begin{verse}
   There is an environment
    for verse \\             % The \\ command separates lines
   Whose features some poets % within a stanza.
   will curse.

                             % One or more blank lines separate stanzas.

   For instead of making\\
   Them do \emph{all} line breaking, \\
   It allows them to put too many words on a line when they'd rather be
   forced to be terse.
\end{verse}

Mathematical formulas may also be displayed.  A
displayed formula
is
one-line long; multiline
formulas require special formatting instructions.
   \[  \ip{\Gamma}{\psi'} = x'' + y^{2} + z_{i}^{n}\]
Don't start a paragraph with a displayed equation,
nor make one a paragraph by itself.

\end{document}               % End of document.

\section{Bench Huge Section 1582}
% synthetic bench section 1582
% This is a sample LaTeX input file.  (Version of 12 August 2004.)
%
% A '%' character causes TeX to ignore all remaining text on the line,
% and is used for comments like this one.

\documentclass{article}      % Specifies the document class

                             % The preamble begins here.
\title{An Example Document}  % Declares the document's title.
\author{Leslie Lamport}      % Declares the author's name.
\date{January 21, 1994}      % Deleting this command produces today's date.

\newcommand{\ip}[2]{(#1, #2)}
                             % Defines \ip{arg1}{arg2} to mean
                             % (arg1, arg2).

%\newcommand{\ip}[2]{\langle #1 | #2\rangle}
                             % This is an alternative definition of
                             % \ip that is commented out.

\begin{document}             % End of preamble and beginning of text.

\maketitle                   % Produces the title.

This is an example input file.  Comparing it with
the output it generates can show you how to
produce a simple document of your own.

\section{Ordinary Text}      % Produces section heading.  Lower-level
                             % sections are begun with similar
                             % \subsection and \subsubsection commands.

The ends  of words and sentences are marked
  by   spaces. It  doesn't matter how many
spaces    you type; one is as good as 100.  The
end of   a line counts as a space.

One   or more   blank lines denote the  end
of  a paragraph.

Since any number of consecutive spaces are treated
like a single one, the formatting of the input
file makes no difference to
      \LaTeX,                % The \LaTeX command generates the LaTeX logo.
but it makes a difference to you.  When you use
\LaTeX, making your input file as easy to read
as possible will be a great help as you write
your document and when you change it.  This sample
file shows how you can add comments to your own input
file.

Because printing is different from typewriting,
there are a number of things that you have to do
differently when preparing an input file than if
you were just typing the document directly.
Quotation marks like
       ``this''
have to be handled specially, as do quotes within
quotes:
       ``\,`this'            % \, separates the double and single quote.
        is what I just
        wrote, not  `that'\,''.

Dashes come in three sizes: an
       intra-word
dash, a medium dash for number ranges like
       1--2,
and a punctuation
       dash---like
this.

A sentence-ending space should be larger than the
space between words within a sentence.  You
sometimes have to type special commands in
conjunction with punctuation characters to get
this right, as in the following sentence.
       Gnats, gnus, etc.\ all  % `\ ' makes an inter-word space.
       begin with G\@.         % \@ marks end-of-sentence punctuation.
You should check the spaces after periods when
reading your output to make sure you haven't
forgotten any special cases.  Generating an
ellipsis
       \ldots\               % `\ ' is needed after `\ldots' because TeX
                             % ignores spaces after command names like \ldots
                             % made from \ + letters.
                             %
                             % Note how a `%' character causes TeX to ignore
                             % the end of the input line, so these blank lines
                             % do not start a new paragraph.
                             %
with the right spacing around the periods requires
a special command.

\LaTeX\ interprets some common characters as
commands, so you must type special commands to
generate them.  These characters include the
following:
       \$ \& \% \# \{ and \}.

In printing, text is usually emphasized with an
       \emph{italic}
type style.

\begin{em}
   A long segment of text can also be emphasized
   in this way.  Text within such a segment can be
   given \emph{additional} emphasis.
\end{em}

It is sometimes necessary to prevent \LaTeX\ from
breaking a line where it might otherwise do so.
This may be at a space, as between the ``Mr.''\ and
``Jones'' in
       ``Mr.~Jones'',        % ~ produces an unbreakable interword space.
or within a word---especially when the word is a
symbol like
       \mbox{\emph{itemnum}}
that makes little sense when hyphenated across
lines.

Footnotes\footnote{This is an example of a footnote.}
pose no problem.

\LaTeX\ is good at typesetting mathematical formulas
like
       \( x-3y + z = 7 \)
or
       \( a_{1} > x^{2n} + y^{2n} > x' \)
or
       \( \ip{A}{B} = \sum_{i} a_{i} b_{i} \).
The spaces you type in a formula are
ignored.  Remember that a letter like
       $x$                   % $ ... $  and  \( ... \)  are equivalent
is a formula when it denotes a mathematical
symbol, and it should be typed as one.

\section{Displayed Text}

Text is displayed by indenting it from the left
margin.  Quotations are commonly displayed.  There
are short quotations
\begin{quote}
   This is a short quotation.  It consists of a
   single paragraph of text.  See how it is formatted.
\end{quote}
and longer ones.
\begin{quotation}
   This is a longer quotation.  It consists of two
   paragraphs of text, neither of which are
   particularly interesting.

   This is the second paragraph of the quotation.  It
   is just as dull as the first paragraph.
\end{quotation}
Another frequently-displayed structure is a list.
The following is an example of an \emph{itemized}
list.
\begin{itemize}
   \item This is the first item of an itemized list.
         Each item in the list is marked with a ``tick''.
         You don't have to worry about what kind of tick
         mark is used.

   \item This is the second item of the list.  It
         contains another list nested inside it.  The inner
         list is an \emph{enumerated} list.
         \begin{enumerate}
            \item This is the first item of an enumerated
                  list that is nested within the itemized list.

            \item This is the second item of the inner list.
                  \LaTeX\ allows you to nest lists deeper than
                  you really should.
         \end{enumerate}
         This is the rest of the second item of the outer
         list.  It is no more interesting than any other
         part of the item.
   \item This is the third item of the list.
\end{itemize}
You can even display poetry.
\begin{verse}
   There is an environment
    for verse \\             % The \\ command separates lines
   Whose features some poets % within a stanza.
   will curse.

                             % One or more blank lines separate stanzas.

   For instead of making\\
   Them do \emph{all} line breaking, \\
   It allows them to put too many words on a line when they'd rather be
   forced to be terse.
\end{verse}

Mathematical formulas may also be displayed.  A
displayed formula
is
one-line long; multiline
formulas require special formatting instructions.
   \[  \ip{\Gamma}{\psi'} = x'' + y^{2} + z_{i}^{n}\]
Don't start a paragraph with a displayed equation,
nor make one a paragraph by itself.

\end{document}               % End of document.

\section{Bench Huge Section 1583}
% synthetic bench section 1583
% This is a sample LaTeX input file.  (Version of 12 August 2004.)
%
% A '%' character causes TeX to ignore all remaining text on the line,
% and is used for comments like this one.

\documentclass{article}      % Specifies the document class

                             % The preamble begins here.
\title{An Example Document}  % Declares the document's title.
\author{Leslie Lamport}      % Declares the author's name.
\date{January 21, 1994}      % Deleting this command produces today's date.

\newcommand{\ip}[2]{(#1, #2)}
                             % Defines \ip{arg1}{arg2} to mean
                             % (arg1, arg2).

%\newcommand{\ip}[2]{\langle #1 | #2\rangle}
                             % This is an alternative definition of
                             % \ip that is commented out.

\begin{document}             % End of preamble and beginning of text.

\maketitle                   % Produces the title.

This is an example input file.  Comparing it with
the output it generates can show you how to
produce a simple document of your own.

\section{Ordinary Text}      % Produces section heading.  Lower-level
                             % sections are begun with similar
                             % \subsection and \subsubsection commands.

The ends  of words and sentences are marked
  by   spaces. It  doesn't matter how many
spaces    you type; one is as good as 100.  The
end of   a line counts as a space.

One   or more   blank lines denote the  end
of  a paragraph.

Since any number of consecutive spaces are treated
like a single one, the formatting of the input
file makes no difference to
      \LaTeX,                % The \LaTeX command generates the LaTeX logo.
but it makes a difference to you.  When you use
\LaTeX, making your input file as easy to read
as possible will be a great help as you write
your document and when you change it.  This sample
file shows how you can add comments to your own input
file.

Because printing is different from typewriting,
there are a number of things that you have to do
differently when preparing an input file than if
you were just typing the document directly.
Quotation marks like
       ``this''
have to be handled specially, as do quotes within
quotes:
       ``\,`this'            % \, separates the double and single quote.
        is what I just
        wrote, not  `that'\,''.

Dashes come in three sizes: an
       intra-word
dash, a medium dash for number ranges like
       1--2,
and a punctuation
       dash---like
this.

A sentence-ending space should be larger than the
space between words within a sentence.  You
sometimes have to type special commands in
conjunction with punctuation characters to get
this right, as in the following sentence.
       Gnats, gnus, etc.\ all  % `\ ' makes an inter-word space.
       begin with G\@.         % \@ marks end-of-sentence punctuation.
You should check the spaces after periods when
reading your output to make sure you haven't
forgotten any special cases.  Generating an
ellipsis
       \ldots\               % `\ ' is needed after `\ldots' because TeX
                             % ignores spaces after command names like \ldots
                             % made from \ + letters.
                             %
                             % Note how a `%' character causes TeX to ignore
                             % the end of the input line, so these blank lines
                             % do not start a new paragraph.
                             %
with the right spacing around the periods requires
a special command.

\LaTeX\ interprets some common characters as
commands, so you must type special commands to
generate them.  These characters include the
following:
       \$ \& \% \# \{ and \}.

In printing, text is usually emphasized with an
       \emph{italic}
type style.

\begin{em}
   A long segment of text can also be emphasized
   in this way.  Text within such a segment can be
   given \emph{additional} emphasis.
\end{em}

It is sometimes necessary to prevent \LaTeX\ from
breaking a line where it might otherwise do so.
This may be at a space, as between the ``Mr.''\ and
``Jones'' in
       ``Mr.~Jones'',        % ~ produces an unbreakable interword space.
or within a word---especially when the word is a
symbol like
       \mbox{\emph{itemnum}}
that makes little sense when hyphenated across
lines.

Footnotes\footnote{This is an example of a footnote.}
pose no problem.

\LaTeX\ is good at typesetting mathematical formulas
like
       \( x-3y + z = 7 \)
or
       \( a_{1} > x^{2n} + y^{2n} > x' \)
or
       \( \ip{A}{B} = \sum_{i} a_{i} b_{i} \).
The spaces you type in a formula are
ignored.  Remember that a letter like
       $x$                   % $ ... $  and  \( ... \)  are equivalent
is a formula when it denotes a mathematical
symbol, and it should be typed as one.

\section{Displayed Text}

Text is displayed by indenting it from the left
margin.  Quotations are commonly displayed.  There
are short quotations
\begin{quote}
   This is a short quotation.  It consists of a
   single paragraph of text.  See how it is formatted.
\end{quote}
and longer ones.
\begin{quotation}
   This is a longer quotation.  It consists of two
   paragraphs of text, neither of which are
   particularly interesting.

   This is the second paragraph of the quotation.  It
   is just as dull as the first paragraph.
\end{quotation}
Another frequently-displayed structure is a list.
The following is an example of an \emph{itemized}
list.
\begin{itemize}
   \item This is the first item of an itemized list.
         Each item in the list is marked with a ``tick''.
         You don't have to worry about what kind of tick
         mark is used.

   \item This is the second item of the list.  It
         contains another list nested inside it.  The inner
         list is an \emph{enumerated} list.
         \begin{enumerate}
            \item This is the first item of an enumerated
                  list that is nested within the itemized list.

            \item This is the second item of the inner list.
                  \LaTeX\ allows you to nest lists deeper than
                  you really should.
         \end{enumerate}
         This is the rest of the second item of the outer
         list.  It is no more interesting than any other
         part of the item.
   \item This is the third item of the list.
\end{itemize}
You can even display poetry.
\begin{verse}
   There is an environment
    for verse \\             % The \\ command separates lines
   Whose features some poets % within a stanza.
   will curse.

                             % One or more blank lines separate stanzas.

   For instead of making\\
   Them do \emph{all} line breaking, \\
   It allows them to put too many words on a line when they'd rather be
   forced to be terse.
\end{verse}

Mathematical formulas may also be displayed.  A
displayed formula
is
one-line long; multiline
formulas require special formatting instructions.
   \[  \ip{\Gamma}{\psi'} = x'' + y^{2} + z_{i}^{n}\]
Don't start a paragraph with a displayed equation,
nor make one a paragraph by itself.

\end{document}               % End of document.

\section{Bench Huge Section 1584}
% synthetic bench section 1584
% This is a sample LaTeX input file.  (Version of 12 August 2004.)
%
% A '%' character causes TeX to ignore all remaining text on the line,
% and is used for comments like this one.

\documentclass{article}      % Specifies the document class

                             % The preamble begins here.
\title{An Example Document}  % Declares the document's title.
\author{Leslie Lamport}      % Declares the author's name.
\date{January 21, 1994}      % Deleting this command produces today's date.

\newcommand{\ip}[2]{(#1, #2)}
                             % Defines \ip{arg1}{arg2} to mean
                             % (arg1, arg2).

%\newcommand{\ip}[2]{\langle #1 | #2\rangle}
                             % This is an alternative definition of
                             % \ip that is commented out.

\begin{document}             % End of preamble and beginning of text.

\maketitle                   % Produces the title.

This is an example input file.  Comparing it with
the output it generates can show you how to
produce a simple document of your own.

\section{Ordinary Text}      % Produces section heading.  Lower-level
                             % sections are begun with similar
                             % \subsection and \subsubsection commands.

The ends  of words and sentences are marked
  by   spaces. It  doesn't matter how many
spaces    you type; one is as good as 100.  The
end of   a line counts as a space.

One   or more   blank lines denote the  end
of  a paragraph.

Since any number of consecutive spaces are treated
like a single one, the formatting of the input
file makes no difference to
      \LaTeX,                % The \LaTeX command generates the LaTeX logo.
but it makes a difference to you.  When you use
\LaTeX, making your input file as easy to read
as possible will be a great help as you write
your document and when you change it.  This sample
file shows how you can add comments to your own input
file.

Because printing is different from typewriting,
there are a number of things that you have to do
differently when preparing an input file than if
you were just typing the document directly.
Quotation marks like
       ``this''
have to be handled specially, as do quotes within
quotes:
       ``\,`this'            % \, separates the double and single quote.
        is what I just
        wrote, not  `that'\,''.

Dashes come in three sizes: an
       intra-word
dash, a medium dash for number ranges like
       1--2,
and a punctuation
       dash---like
this.

A sentence-ending space should be larger than the
space between words within a sentence.  You
sometimes have to type special commands in
conjunction with punctuation characters to get
this right, as in the following sentence.
       Gnats, gnus, etc.\ all  % `\ ' makes an inter-word space.
       begin with G\@.         % \@ marks end-of-sentence punctuation.
You should check the spaces after periods when
reading your output to make sure you haven't
forgotten any special cases.  Generating an
ellipsis
       \ldots\               % `\ ' is needed after `\ldots' because TeX
                             % ignores spaces after command names like \ldots
                             % made from \ + letters.
                             %
                             % Note how a `%' character causes TeX to ignore
                             % the end of the input line, so these blank lines
                             % do not start a new paragraph.
                             %
with the right spacing around the periods requires
a special command.

\LaTeX\ interprets some common characters as
commands, so you must type special commands to
generate them.  These characters include the
following:
       \$ \& \% \# \{ and \}.

In printing, text is usually emphasized with an
       \emph{italic}
type style.

\begin{em}
   A long segment of text can also be emphasized
   in this way.  Text within such a segment can be
   given \emph{additional} emphasis.
\end{em}

It is sometimes necessary to prevent \LaTeX\ from
breaking a line where it might otherwise do so.
This may be at a space, as between the ``Mr.''\ and
``Jones'' in
       ``Mr.~Jones'',        % ~ produces an unbreakable interword space.
or within a word---especially when the word is a
symbol like
       \mbox{\emph{itemnum}}
that makes little sense when hyphenated across
lines.

Footnotes\footnote{This is an example of a footnote.}
pose no problem.

\LaTeX\ is good at typesetting mathematical formulas
like
       \( x-3y + z = 7 \)
or
       \( a_{1} > x^{2n} + y^{2n} > x' \)
or
       \( \ip{A}{B} = \sum_{i} a_{i} b_{i} \).
The spaces you type in a formula are
ignored.  Remember that a letter like
       $x$                   % $ ... $  and  \( ... \)  are equivalent
is a formula when it denotes a mathematical
symbol, and it should be typed as one.

\section{Displayed Text}

Text is displayed by indenting it from the left
margin.  Quotations are commonly displayed.  There
are short quotations
\begin{quote}
   This is a short quotation.  It consists of a
   single paragraph of text.  See how it is formatted.
\end{quote}
and longer ones.
\begin{quotation}
   This is a longer quotation.  It consists of two
   paragraphs of text, neither of which are
   particularly interesting.

   This is the second paragraph of the quotation.  It
   is just as dull as the first paragraph.
\end{quotation}
Another frequently-displayed structure is a list.
The following is an example of an \emph{itemized}
list.
\begin{itemize}
   \item This is the first item of an itemized list.
         Each item in the list is marked with a ``tick''.
         You don't have to worry about what kind of tick
         mark is used.

   \item This is the second item of the list.  It
         contains another list nested inside it.  The inner
         list is an \emph{enumerated} list.
         \begin{enumerate}
            \item This is the first item of an enumerated
                  list that is nested within the itemized list.

            \item This is the second item of the inner list.
                  \LaTeX\ allows you to nest lists deeper than
                  you really should.
         \end{enumerate}
         This is the rest of the second item of the outer
         list.  It is no more interesting than any other
         part of the item.
   \item This is the third item of the list.
\end{itemize}
You can even display poetry.
\begin{verse}
   There is an environment
    for verse \\             % The \\ command separates lines
   Whose features some poets % within a stanza.
   will curse.

                             % One or more blank lines separate stanzas.

   For instead of making\\
   Them do \emph{all} line breaking, \\
   It allows them to put too many words on a line when they'd rather be
   forced to be terse.
\end{verse}

Mathematical formulas may also be displayed.  A
displayed formula
is
one-line long; multiline
formulas require special formatting instructions.
   \[  \ip{\Gamma}{\psi'} = x'' + y^{2} + z_{i}^{n}\]
Don't start a paragraph with a displayed equation,
nor make one a paragraph by itself.

\end{document}               % End of document.

\section{Bench Huge Section 1585}
% synthetic bench section 1585
% This is a sample LaTeX input file.  (Version of 12 August 2004.)
%
% A '%' character causes TeX to ignore all remaining text on the line,
% and is used for comments like this one.

\documentclass{article}      % Specifies the document class

                             % The preamble begins here.
\title{An Example Document}  % Declares the document's title.
\author{Leslie Lamport}      % Declares the author's name.
\date{January 21, 1994}      % Deleting this command produces today's date.

\newcommand{\ip}[2]{(#1, #2)}
                             % Defines \ip{arg1}{arg2} to mean
                             % (arg1, arg2).

%\newcommand{\ip}[2]{\langle #1 | #2\rangle}
                             % This is an alternative definition of
                             % \ip that is commented out.

\begin{document}             % End of preamble and beginning of text.

\maketitle                   % Produces the title.

This is an example input file.  Comparing it with
the output it generates can show you how to
produce a simple document of your own.

\section{Ordinary Text}      % Produces section heading.  Lower-level
                             % sections are begun with similar
                             % \subsection and \subsubsection commands.

The ends  of words and sentences are marked
  by   spaces. It  doesn't matter how many
spaces    you type; one is as good as 100.  The
end of   a line counts as a space.

One   or more   blank lines denote the  end
of  a paragraph.

Since any number of consecutive spaces are treated
like a single one, the formatting of the input
file makes no difference to
      \LaTeX,                % The \LaTeX command generates the LaTeX logo.
but it makes a difference to you.  When you use
\LaTeX, making your input file as easy to read
as possible will be a great help as you write
your document and when you change it.  This sample
file shows how you can add comments to your own input
file.

Because printing is different from typewriting,
there are a number of things that you have to do
differently when preparing an input file than if
you were just typing the document directly.
Quotation marks like
       ``this''
have to be handled specially, as do quotes within
quotes:
       ``\,`this'            % \, separates the double and single quote.
        is what I just
        wrote, not  `that'\,''.

Dashes come in three sizes: an
       intra-word
dash, a medium dash for number ranges like
       1--2,
and a punctuation
       dash---like
this.

A sentence-ending space should be larger than the
space between words within a sentence.  You
sometimes have to type special commands in
conjunction with punctuation characters to get
this right, as in the following sentence.
       Gnats, gnus, etc.\ all  % `\ ' makes an inter-word space.
       begin with G\@.         % \@ marks end-of-sentence punctuation.
You should check the spaces after periods when
reading your output to make sure you haven't
forgotten any special cases.  Generating an
ellipsis
       \ldots\               % `\ ' is needed after `\ldots' because TeX
                             % ignores spaces after command names like \ldots
                             % made from \ + letters.
                             %
                             % Note how a `%' character causes TeX to ignore
                             % the end of the input line, so these blank lines
                             % do not start a new paragraph.
                             %
with the right spacing around the periods requires
a special command.

\LaTeX\ interprets some common characters as
commands, so you must type special commands to
generate them.  These characters include the
following:
       \$ \& \% \# \{ and \}.

In printing, text is usually emphasized with an
       \emph{italic}
type style.

\begin{em}
   A long segment of text can also be emphasized
   in this way.  Text within such a segment can be
   given \emph{additional} emphasis.
\end{em}

It is sometimes necessary to prevent \LaTeX\ from
breaking a line where it might otherwise do so.
This may be at a space, as between the ``Mr.''\ and
``Jones'' in
       ``Mr.~Jones'',        % ~ produces an unbreakable interword space.
or within a word---especially when the word is a
symbol like
       \mbox{\emph{itemnum}}
that makes little sense when hyphenated across
lines.

Footnotes\footnote{This is an example of a footnote.}
pose no problem.

\LaTeX\ is good at typesetting mathematical formulas
like
       \( x-3y + z = 7 \)
or
       \( a_{1} > x^{2n} + y^{2n} > x' \)
or
       \( \ip{A}{B} = \sum_{i} a_{i} b_{i} \).
The spaces you type in a formula are
ignored.  Remember that a letter like
       $x$                   % $ ... $  and  \( ... \)  are equivalent
is a formula when it denotes a mathematical
symbol, and it should be typed as one.

\section{Displayed Text}

Text is displayed by indenting it from the left
margin.  Quotations are commonly displayed.  There
are short quotations
\begin{quote}
   This is a short quotation.  It consists of a
   single paragraph of text.  See how it is formatted.
\end{quote}
and longer ones.
\begin{quotation}
   This is a longer quotation.  It consists of two
   paragraphs of text, neither of which are
   particularly interesting.

   This is the second paragraph of the quotation.  It
   is just as dull as the first paragraph.
\end{quotation}
Another frequently-displayed structure is a list.
The following is an example of an \emph{itemized}
list.
\begin{itemize}
   \item This is the first item of an itemized list.
         Each item in the list is marked with a ``tick''.
         You don't have to worry about what kind of tick
         mark is used.

   \item This is the second item of the list.  It
         contains another list nested inside it.  The inner
         list is an \emph{enumerated} list.
         \begin{enumerate}
            \item This is the first item of an enumerated
                  list that is nested within the itemized list.

            \item This is the second item of the inner list.
                  \LaTeX\ allows you to nest lists deeper than
                  you really should.
         \end{enumerate}
         This is the rest of the second item of the outer
         list.  It is no more interesting than any other
         part of the item.
   \item This is the third item of the list.
\end{itemize}
You can even display poetry.
\begin{verse}
   There is an environment
    for verse \\             % The \\ command separates lines
   Whose features some poets % within a stanza.
   will curse.

                             % One or more blank lines separate stanzas.

   For instead of making\\
   Them do \emph{all} line breaking, \\
   It allows them to put too many words on a line when they'd rather be
   forced to be terse.
\end{verse}

Mathematical formulas may also be displayed.  A
displayed formula
is
one-line long; multiline
formulas require special formatting instructions.
   \[  \ip{\Gamma}{\psi'} = x'' + y^{2} + z_{i}^{n}\]
Don't start a paragraph with a displayed equation,
nor make one a paragraph by itself.

\end{document}               % End of document.

\section{Bench Huge Section 1586}
% synthetic bench section 1586
% This is a sample LaTeX input file.  (Version of 12 August 2004.)
%
% A '%' character causes TeX to ignore all remaining text on the line,
% and is used for comments like this one.

\documentclass{article}      % Specifies the document class

                             % The preamble begins here.
\title{An Example Document}  % Declares the document's title.
\author{Leslie Lamport}      % Declares the author's name.
\date{January 21, 1994}      % Deleting this command produces today's date.

\newcommand{\ip}[2]{(#1, #2)}
                             % Defines \ip{arg1}{arg2} to mean
                             % (arg1, arg2).

%\newcommand{\ip}[2]{\langle #1 | #2\rangle}
                             % This is an alternative definition of
                             % \ip that is commented out.

\begin{document}             % End of preamble and beginning of text.

\maketitle                   % Produces the title.

This is an example input file.  Comparing it with
the output it generates can show you how to
produce a simple document of your own.

\section{Ordinary Text}      % Produces section heading.  Lower-level
                             % sections are begun with similar
                             % \subsection and \subsubsection commands.

The ends  of words and sentences are marked
  by   spaces. It  doesn't matter how many
spaces    you type; one is as good as 100.  The
end of   a line counts as a space.

One   or more   blank lines denote the  end
of  a paragraph.

Since any number of consecutive spaces are treated
like a single one, the formatting of the input
file makes no difference to
      \LaTeX,                % The \LaTeX command generates the LaTeX logo.
but it makes a difference to you.  When you use
\LaTeX, making your input file as easy to read
as possible will be a great help as you write
your document and when you change it.  This sample
file shows how you can add comments to your own input
file.

Because printing is different from typewriting,
there are a number of things that you have to do
differently when preparing an input file than if
you were just typing the document directly.
Quotation marks like
       ``this''
have to be handled specially, as do quotes within
quotes:
       ``\,`this'            % \, separates the double and single quote.
        is what I just
        wrote, not  `that'\,''.

Dashes come in three sizes: an
       intra-word
dash, a medium dash for number ranges like
       1--2,
and a punctuation
       dash---like
this.

A sentence-ending space should be larger than the
space between words within a sentence.  You
sometimes have to type special commands in
conjunction with punctuation characters to get
this right, as in the following sentence.
       Gnats, gnus, etc.\ all  % `\ ' makes an inter-word space.
       begin with G\@.         % \@ marks end-of-sentence punctuation.
You should check the spaces after periods when
reading your output to make sure you haven't
forgotten any special cases.  Generating an
ellipsis
       \ldots\               % `\ ' is needed after `\ldots' because TeX
                             % ignores spaces after command names like \ldots
                             % made from \ + letters.
                             %
                             % Note how a `%' character causes TeX to ignore
                             % the end of the input line, so these blank lines
                             % do not start a new paragraph.
                             %
with the right spacing around the periods requires
a special command.

\LaTeX\ interprets some common characters as
commands, so you must type special commands to
generate them.  These characters include the
following:
       \$ \& \% \# \{ and \}.

In printing, text is usually emphasized with an
       \emph{italic}
type style.

\begin{em}
   A long segment of text can also be emphasized
   in this way.  Text within such a segment can be
   given \emph{additional} emphasis.
\end{em}

It is sometimes necessary to prevent \LaTeX\ from
breaking a line where it might otherwise do so.
This may be at a space, as between the ``Mr.''\ and
``Jones'' in
       ``Mr.~Jones'',        % ~ produces an unbreakable interword space.
or within a word---especially when the word is a
symbol like
       \mbox{\emph{itemnum}}
that makes little sense when hyphenated across
lines.

Footnotes\footnote{This is an example of a footnote.}
pose no problem.

\LaTeX\ is good at typesetting mathematical formulas
like
       \( x-3y + z = 7 \)
or
       \( a_{1} > x^{2n} + y^{2n} > x' \)
or
       \( \ip{A}{B} = \sum_{i} a_{i} b_{i} \).
The spaces you type in a formula are
ignored.  Remember that a letter like
       $x$                   % $ ... $  and  \( ... \)  are equivalent
is a formula when it denotes a mathematical
symbol, and it should be typed as one.

\section{Displayed Text}

Text is displayed by indenting it from the left
margin.  Quotations are commonly displayed.  There
are short quotations
\begin{quote}
   This is a short quotation.  It consists of a
   single paragraph of text.  See how it is formatted.
\end{quote}
and longer ones.
\begin{quotation}
   This is a longer quotation.  It consists of two
   paragraphs of text, neither of which are
   particularly interesting.

   This is the second paragraph of the quotation.  It
   is just as dull as the first paragraph.
\end{quotation}
Another frequently-displayed structure is a list.
The following is an example of an \emph{itemized}
list.
\begin{itemize}
   \item This is the first item of an itemized list.
         Each item in the list is marked with a ``tick''.
         You don't have to worry about what kind of tick
         mark is used.

   \item This is the second item of the list.  It
         contains another list nested inside it.  The inner
         list is an \emph{enumerated} list.
         \begin{enumerate}
            \item This is the first item of an enumerated
                  list that is nested within the itemized list.

            \item This is the second item of the inner list.
                  \LaTeX\ allows you to nest lists deeper than
                  you really should.
         \end{enumerate}
         This is the rest of the second item of the outer
         list.  It is no more interesting than any other
         part of the item.
   \item This is the third item of the list.
\end{itemize}
You can even display poetry.
\begin{verse}
   There is an environment
    for verse \\             % The \\ command separates lines
   Whose features some poets % within a stanza.
   will curse.

                             % One or more blank lines separate stanzas.

   For instead of making\\
   Them do \emph{all} line breaking, \\
   It allows them to put too many words on a line when they'd rather be
   forced to be terse.
\end{verse}

Mathematical formulas may also be displayed.  A
displayed formula
is
one-line long; multiline
formulas require special formatting instructions.
   \[  \ip{\Gamma}{\psi'} = x'' + y^{2} + z_{i}^{n}\]
Don't start a paragraph with a displayed equation,
nor make one a paragraph by itself.

\end{document}               % End of document.

\section{Bench Huge Section 1587}
% synthetic bench section 1587
% This is a sample LaTeX input file.  (Version of 12 August 2004.)
%
% A '%' character causes TeX to ignore all remaining text on the line,
% and is used for comments like this one.

\documentclass{article}      % Specifies the document class

                             % The preamble begins here.
\title{An Example Document}  % Declares the document's title.
\author{Leslie Lamport}      % Declares the author's name.
\date{January 21, 1994}      % Deleting this command produces today's date.

\newcommand{\ip}[2]{(#1, #2)}
                             % Defines \ip{arg1}{arg2} to mean
                             % (arg1, arg2).

%\newcommand{\ip}[2]{\langle #1 | #2\rangle}
                             % This is an alternative definition of
                             % \ip that is commented out.

\begin{document}             % End of preamble and beginning of text.

\maketitle                   % Produces the title.

This is an example input file.  Comparing it with
the output it generates can show you how to
produce a simple document of your own.

\section{Ordinary Text}      % Produces section heading.  Lower-level
                             % sections are begun with similar
                             % \subsection and \subsubsection commands.

The ends  of words and sentences are marked
  by   spaces. It  doesn't matter how many
spaces    you type; one is as good as 100.  The
end of   a line counts as a space.

One   or more   blank lines denote the  end
of  a paragraph.

Since any number of consecutive spaces are treated
like a single one, the formatting of the input
file makes no difference to
      \LaTeX,                % The \LaTeX command generates the LaTeX logo.
but it makes a difference to you.  When you use
\LaTeX, making your input file as easy to read
as possible will be a great help as you write
your document and when you change it.  This sample
file shows how you can add comments to your own input
file.

Because printing is different from typewriting,
there are a number of things that you have to do
differently when preparing an input file than if
you were just typing the document directly.
Quotation marks like
       ``this''
have to be handled specially, as do quotes within
quotes:
       ``\,`this'            % \, separates the double and single quote.
        is what I just
        wrote, not  `that'\,''.

Dashes come in three sizes: an
       intra-word
dash, a medium dash for number ranges like
       1--2,
and a punctuation
       dash---like
this.

A sentence-ending space should be larger than the
space between words within a sentence.  You
sometimes have to type special commands in
conjunction with punctuation characters to get
this right, as in the following sentence.
       Gnats, gnus, etc.\ all  % `\ ' makes an inter-word space.
       begin with G\@.         % \@ marks end-of-sentence punctuation.
You should check the spaces after periods when
reading your output to make sure you haven't
forgotten any special cases.  Generating an
ellipsis
       \ldots\               % `\ ' is needed after `\ldots' because TeX
                             % ignores spaces after command names like \ldots
                             % made from \ + letters.
                             %
                             % Note how a `%' character causes TeX to ignore
                             % the end of the input line, so these blank lines
                             % do not start a new paragraph.
                             %
with the right spacing around the periods requires
a special command.

\LaTeX\ interprets some common characters as
commands, so you must type special commands to
generate them.  These characters include the
following:
       \$ \& \% \# \{ and \}.

In printing, text is usually emphasized with an
       \emph{italic}
type style.

\begin{em}
   A long segment of text can also be emphasized
   in this way.  Text within such a segment can be
   given \emph{additional} emphasis.
\end{em}

It is sometimes necessary to prevent \LaTeX\ from
breaking a line where it might otherwise do so.
This may be at a space, as between the ``Mr.''\ and
``Jones'' in
       ``Mr.~Jones'',        % ~ produces an unbreakable interword space.
or within a word---especially when the word is a
symbol like
       \mbox{\emph{itemnum}}
that makes little sense when hyphenated across
lines.

Footnotes\footnote{This is an example of a footnote.}
pose no problem.

\LaTeX\ is good at typesetting mathematical formulas
like
       \( x-3y + z = 7 \)
or
       \( a_{1} > x^{2n} + y^{2n} > x' \)
or
       \( \ip{A}{B} = \sum_{i} a_{i} b_{i} \).
The spaces you type in a formula are
ignored.  Remember that a letter like
       $x$                   % $ ... $  and  \( ... \)  are equivalent
is a formula when it denotes a mathematical
symbol, and it should be typed as one.

\section{Displayed Text}

Text is displayed by indenting it from the left
margin.  Quotations are commonly displayed.  There
are short quotations
\begin{quote}
   This is a short quotation.  It consists of a
   single paragraph of text.  See how it is formatted.
\end{quote}
and longer ones.
\begin{quotation}
   This is a longer quotation.  It consists of two
   paragraphs of text, neither of which are
   particularly interesting.

   This is the second paragraph of the quotation.  It
   is just as dull as the first paragraph.
\end{quotation}
Another frequently-displayed structure is a list.
The following is an example of an \emph{itemized}
list.
\begin{itemize}
   \item This is the first item of an itemized list.
         Each item in the list is marked with a ``tick''.
         You don't have to worry about what kind of tick
         mark is used.

   \item This is the second item of the list.  It
         contains another list nested inside it.  The inner
         list is an \emph{enumerated} list.
         \begin{enumerate}
            \item This is the first item of an enumerated
                  list that is nested within the itemized list.

            \item This is the second item of the inner list.
                  \LaTeX\ allows you to nest lists deeper than
                  you really should.
         \end{enumerate}
         This is the rest of the second item of the outer
         list.  It is no more interesting than any other
         part of the item.
   \item This is the third item of the list.
\end{itemize}
You can even display poetry.
\begin{verse}
   There is an environment
    for verse \\             % The \\ command separates lines
   Whose features some poets % within a stanza.
   will curse.

                             % One or more blank lines separate stanzas.

   For instead of making\\
   Them do \emph{all} line breaking, \\
   It allows them to put too many words on a line when they'd rather be
   forced to be terse.
\end{verse}

Mathematical formulas may also be displayed.  A
displayed formula
is
one-line long; multiline
formulas require special formatting instructions.
   \[  \ip{\Gamma}{\psi'} = x'' + y^{2} + z_{i}^{n}\]
Don't start a paragraph with a displayed equation,
nor make one a paragraph by itself.

\end{document}               % End of document.

\section{Bench Huge Section 1588}
% synthetic bench section 1588
% This is a sample LaTeX input file.  (Version of 12 August 2004.)
%
% A '%' character causes TeX to ignore all remaining text on the line,
% and is used for comments like this one.

\documentclass{article}      % Specifies the document class

                             % The preamble begins here.
\title{An Example Document}  % Declares the document's title.
\author{Leslie Lamport}      % Declares the author's name.
\date{January 21, 1994}      % Deleting this command produces today's date.

\newcommand{\ip}[2]{(#1, #2)}
                             % Defines \ip{arg1}{arg2} to mean
                             % (arg1, arg2).

%\newcommand{\ip}[2]{\langle #1 | #2\rangle}
                             % This is an alternative definition of
                             % \ip that is commented out.

\begin{document}             % End of preamble and beginning of text.

\maketitle                   % Produces the title.

This is an example input file.  Comparing it with
the output it generates can show you how to
produce a simple document of your own.

\section{Ordinary Text}      % Produces section heading.  Lower-level
                             % sections are begun with similar
                             % \subsection and \subsubsection commands.

The ends  of words and sentences are marked
  by   spaces. It  doesn't matter how many
spaces    you type; one is as good as 100.  The
end of   a line counts as a space.

One   or more   blank lines denote the  end
of  a paragraph.

Since any number of consecutive spaces are treated
like a single one, the formatting of the input
file makes no difference to
      \LaTeX,                % The \LaTeX command generates the LaTeX logo.
but it makes a difference to you.  When you use
\LaTeX, making your input file as easy to read
as possible will be a great help as you write
your document and when you change it.  This sample
file shows how you can add comments to your own input
file.

Because printing is different from typewriting,
there are a number of things that you have to do
differently when preparing an input file than if
you were just typing the document directly.
Quotation marks like
       ``this''
have to be handled specially, as do quotes within
quotes:
       ``\,`this'            % \, separates the double and single quote.
        is what I just
        wrote, not  `that'\,''.

Dashes come in three sizes: an
       intra-word
dash, a medium dash for number ranges like
       1--2,
and a punctuation
       dash---like
this.

A sentence-ending space should be larger than the
space between words within a sentence.  You
sometimes have to type special commands in
conjunction with punctuation characters to get
this right, as in the following sentence.
       Gnats, gnus, etc.\ all  % `\ ' makes an inter-word space.
       begin with G\@.         % \@ marks end-of-sentence punctuation.
You should check the spaces after periods when
reading your output to make sure you haven't
forgotten any special cases.  Generating an
ellipsis
       \ldots\               % `\ ' is needed after `\ldots' because TeX
                             % ignores spaces after command names like \ldots
                             % made from \ + letters.
                             %
                             % Note how a `%' character causes TeX to ignore
                             % the end of the input line, so these blank lines
                             % do not start a new paragraph.
                             %
with the right spacing around the periods requires
a special command.

\LaTeX\ interprets some common characters as
commands, so you must type special commands to
generate them.  These characters include the
following:
       \$ \& \% \# \{ and \}.

In printing, text is usually emphasized with an
       \emph{italic}
type style.

\begin{em}
   A long segment of text can also be emphasized
   in this way.  Text within such a segment can be
   given \emph{additional} emphasis.
\end{em}

It is sometimes necessary to prevent \LaTeX\ from
breaking a line where it might otherwise do so.
This may be at a space, as between the ``Mr.''\ and
``Jones'' in
       ``Mr.~Jones'',        % ~ produces an unbreakable interword space.
or within a word---especially when the word is a
symbol like
       \mbox{\emph{itemnum}}
that makes little sense when hyphenated across
lines.

Footnotes\footnote{This is an example of a footnote.}
pose no problem.

\LaTeX\ is good at typesetting mathematical formulas
like
       \( x-3y + z = 7 \)
or
       \( a_{1} > x^{2n} + y^{2n} > x' \)
or
       \( \ip{A}{B} = \sum_{i} a_{i} b_{i} \).
The spaces you type in a formula are
ignored.  Remember that a letter like
       $x$                   % $ ... $  and  \( ... \)  are equivalent
is a formula when it denotes a mathematical
symbol, and it should be typed as one.

\section{Displayed Text}

Text is displayed by indenting it from the left
margin.  Quotations are commonly displayed.  There
are short quotations
\begin{quote}
   This is a short quotation.  It consists of a
   single paragraph of text.  See how it is formatted.
\end{quote}
and longer ones.
\begin{quotation}
   This is a longer quotation.  It consists of two
   paragraphs of text, neither of which are
   particularly interesting.

   This is the second paragraph of the quotation.  It
   is just as dull as the first paragraph.
\end{quotation}
Another frequently-displayed structure is a list.
The following is an example of an \emph{itemized}
list.
\begin{itemize}
   \item This is the first item of an itemized list.
         Each item in the list is marked with a ``tick''.
         You don't have to worry about what kind of tick
         mark is used.

   \item This is the second item of the list.  It
         contains another list nested inside it.  The inner
         list is an \emph{enumerated} list.
         \begin{enumerate}
            \item This is the first item of an enumerated
                  list that is nested within the itemized list.

            \item This is the second item of the inner list.
                  \LaTeX\ allows you to nest lists deeper than
                  you really should.
         \end{enumerate}
         This is the rest of the second item of the outer
         list.  It is no more interesting than any other
         part of the item.
   \item This is the third item of the list.
\end{itemize}
You can even display poetry.
\begin{verse}
   There is an environment
    for verse \\             % The \\ command separates lines
   Whose features some poets % within a stanza.
   will curse.

                             % One or more blank lines separate stanzas.

   For instead of making\\
   Them do \emph{all} line breaking, \\
   It allows them to put too many words on a line when they'd rather be
   forced to be terse.
\end{verse}

Mathematical formulas may also be displayed.  A
displayed formula
is
one-line long; multiline
formulas require special formatting instructions.
   \[  \ip{\Gamma}{\psi'} = x'' + y^{2} + z_{i}^{n}\]
Don't start a paragraph with a displayed equation,
nor make one a paragraph by itself.

\end{document}               % End of document.

\section{Bench Huge Section 1589}
% synthetic bench section 1589
% This is a sample LaTeX input file.  (Version of 12 August 2004.)
%
% A '%' character causes TeX to ignore all remaining text on the line,
% and is used for comments like this one.

\documentclass{article}      % Specifies the document class

                             % The preamble begins here.
\title{An Example Document}  % Declares the document's title.
\author{Leslie Lamport}      % Declares the author's name.
\date{January 21, 1994}      % Deleting this command produces today's date.

\newcommand{\ip}[2]{(#1, #2)}
                             % Defines \ip{arg1}{arg2} to mean
                             % (arg1, arg2).

%\newcommand{\ip}[2]{\langle #1 | #2\rangle}
                             % This is an alternative definition of
                             % \ip that is commented out.

\begin{document}             % End of preamble and beginning of text.

\maketitle                   % Produces the title.

This is an example input file.  Comparing it with
the output it generates can show you how to
produce a simple document of your own.

\section{Ordinary Text}      % Produces section heading.  Lower-level
                             % sections are begun with similar
                             % \subsection and \subsubsection commands.

The ends  of words and sentences are marked
  by   spaces. It  doesn't matter how many
spaces    you type; one is as good as 100.  The
end of   a line counts as a space.

One   or more   blank lines denote the  end
of  a paragraph.

Since any number of consecutive spaces are treated
like a single one, the formatting of the input
file makes no difference to
      \LaTeX,                % The \LaTeX command generates the LaTeX logo.
but it makes a difference to you.  When you use
\LaTeX, making your input file as easy to read
as possible will be a great help as you write
your document and when you change it.  This sample
file shows how you can add comments to your own input
file.

Because printing is different from typewriting,
there are a number of things that you have to do
differently when preparing an input file than if
you were just typing the document directly.
Quotation marks like
       ``this''
have to be handled specially, as do quotes within
quotes:
       ``\,`this'            % \, separates the double and single quote.
        is what I just
        wrote, not  `that'\,''.

Dashes come in three sizes: an
       intra-word
dash, a medium dash for number ranges like
       1--2,
and a punctuation
       dash---like
this.

A sentence-ending space should be larger than the
space between words within a sentence.  You
sometimes have to type special commands in
conjunction with punctuation characters to get
this right, as in the following sentence.
       Gnats, gnus, etc.\ all  % `\ ' makes an inter-word space.
       begin with G\@.         % \@ marks end-of-sentence punctuation.
You should check the spaces after periods when
reading your output to make sure you haven't
forgotten any special cases.  Generating an
ellipsis
       \ldots\               % `\ ' is needed after `\ldots' because TeX
                             % ignores spaces after command names like \ldots
                             % made from \ + letters.
                             %
                             % Note how a `%' character causes TeX to ignore
                             % the end of the input line, so these blank lines
                             % do not start a new paragraph.
                             %
with the right spacing around the periods requires
a special command.

\LaTeX\ interprets some common characters as
commands, so you must type special commands to
generate them.  These characters include the
following:
       \$ \& \% \# \{ and \}.

In printing, text is usually emphasized with an
       \emph{italic}
type style.

\begin{em}
   A long segment of text can also be emphasized
   in this way.  Text within such a segment can be
   given \emph{additional} emphasis.
\end{em}

It is sometimes necessary to prevent \LaTeX\ from
breaking a line where it might otherwise do so.
This may be at a space, as between the ``Mr.''\ and
``Jones'' in
       ``Mr.~Jones'',        % ~ produces an unbreakable interword space.
or within a word---especially when the word is a
symbol like
       \mbox{\emph{itemnum}}
that makes little sense when hyphenated across
lines.

Footnotes\footnote{This is an example of a footnote.}
pose no problem.

\LaTeX\ is good at typesetting mathematical formulas
like
       \( x-3y + z = 7 \)
or
       \( a_{1} > x^{2n} + y^{2n} > x' \)
or
       \( \ip{A}{B} = \sum_{i} a_{i} b_{i} \).
The spaces you type in a formula are
ignored.  Remember that a letter like
       $x$                   % $ ... $  and  \( ... \)  are equivalent
is a formula when it denotes a mathematical
symbol, and it should be typed as one.

\section{Displayed Text}

Text is displayed by indenting it from the left
margin.  Quotations are commonly displayed.  There
are short quotations
\begin{quote}
   This is a short quotation.  It consists of a
   single paragraph of text.  See how it is formatted.
\end{quote}
and longer ones.
\begin{quotation}
   This is a longer quotation.  It consists of two
   paragraphs of text, neither of which are
   particularly interesting.

   This is the second paragraph of the quotation.  It
   is just as dull as the first paragraph.
\end{quotation}
Another frequently-displayed structure is a list.
The following is an example of an \emph{itemized}
list.
\begin{itemize}
   \item This is the first item of an itemized list.
         Each item in the list is marked with a ``tick''.
         You don't have to worry about what kind of tick
         mark is used.

   \item This is the second item of the list.  It
         contains another list nested inside it.  The inner
         list is an \emph{enumerated} list.
         \begin{enumerate}
            \item This is the first item of an enumerated
                  list that is nested within the itemized list.

            \item This is the second item of the inner list.
                  \LaTeX\ allows you to nest lists deeper than
                  you really should.
         \end{enumerate}
         This is the rest of the second item of the outer
         list.  It is no more interesting than any other
         part of the item.
   \item This is the third item of the list.
\end{itemize}
You can even display poetry.
\begin{verse}
   There is an environment
    for verse \\             % The \\ command separates lines
   Whose features some poets % within a stanza.
   will curse.

                             % One or more blank lines separate stanzas.

   For instead of making\\
   Them do \emph{all} line breaking, \\
   It allows them to put too many words on a line when they'd rather be
   forced to be terse.
\end{verse}

Mathematical formulas may also be displayed.  A
displayed formula
is
one-line long; multiline
formulas require special formatting instructions.
   \[  \ip{\Gamma}{\psi'} = x'' + y^{2} + z_{i}^{n}\]
Don't start a paragraph with a displayed equation,
nor make one a paragraph by itself.

\end{document}               % End of document.

\section{Bench Huge Section 1590}
% synthetic bench section 1590
% This is a sample LaTeX input file.  (Version of 12 August 2004.)
%
% A '%' character causes TeX to ignore all remaining text on the line,
% and is used for comments like this one.

\documentclass{article}      % Specifies the document class

                             % The preamble begins here.
\title{An Example Document}  % Declares the document's title.
\author{Leslie Lamport}      % Declares the author's name.
\date{January 21, 1994}      % Deleting this command produces today's date.

\newcommand{\ip}[2]{(#1, #2)}
                             % Defines \ip{arg1}{arg2} to mean
                             % (arg1, arg2).

%\newcommand{\ip}[2]{\langle #1 | #2\rangle}
                             % This is an alternative definition of
                             % \ip that is commented out.

\begin{document}             % End of preamble and beginning of text.

\maketitle                   % Produces the title.

This is an example input file.  Comparing it with
the output it generates can show you how to
produce a simple document of your own.

\section{Ordinary Text}      % Produces section heading.  Lower-level
                             % sections are begun with similar
                             % \subsection and \subsubsection commands.

The ends  of words and sentences are marked
  by   spaces. It  doesn't matter how many
spaces    you type; one is as good as 100.  The
end of   a line counts as a space.

One   or more   blank lines denote the  end
of  a paragraph.

Since any number of consecutive spaces are treated
like a single one, the formatting of the input
file makes no difference to
      \LaTeX,                % The \LaTeX command generates the LaTeX logo.
but it makes a difference to you.  When you use
\LaTeX, making your input file as easy to read
as possible will be a great help as you write
your document and when you change it.  This sample
file shows how you can add comments to your own input
file.

Because printing is different from typewriting,
there are a number of things that you have to do
differently when preparing an input file than if
you were just typing the document directly.
Quotation marks like
       ``this''
have to be handled specially, as do quotes within
quotes:
       ``\,`this'            % \, separates the double and single quote.
        is what I just
        wrote, not  `that'\,''.

Dashes come in three sizes: an
       intra-word
dash, a medium dash for number ranges like
       1--2,
and a punctuation
       dash---like
this.

A sentence-ending space should be larger than the
space between words within a sentence.  You
sometimes have to type special commands in
conjunction with punctuation characters to get
this right, as in the following sentence.
       Gnats, gnus, etc.\ all  % `\ ' makes an inter-word space.
       begin with G\@.         % \@ marks end-of-sentence punctuation.
You should check the spaces after periods when
reading your output to make sure you haven't
forgotten any special cases.  Generating an
ellipsis
       \ldots\               % `\ ' is needed after `\ldots' because TeX
                             % ignores spaces after command names like \ldots
                             % made from \ + letters.
                             %
                             % Note how a `%' character causes TeX to ignore
                             % the end of the input line, so these blank lines
                             % do not start a new paragraph.
                             %
with the right spacing around the periods requires
a special command.

\LaTeX\ interprets some common characters as
commands, so you must type special commands to
generate them.  These characters include the
following:
       \$ \& \% \# \{ and \}.

In printing, text is usually emphasized with an
       \emph{italic}
type style.

\begin{em}
   A long segment of text can also be emphasized
   in this way.  Text within such a segment can be
   given \emph{additional} emphasis.
\end{em}

It is sometimes necessary to prevent \LaTeX\ from
breaking a line where it might otherwise do so.
This may be at a space, as between the ``Mr.''\ and
``Jones'' in
       ``Mr.~Jones'',        % ~ produces an unbreakable interword space.
or within a word---especially when the word is a
symbol like
       \mbox{\emph{itemnum}}
that makes little sense when hyphenated across
lines.

Footnotes\footnote{This is an example of a footnote.}
pose no problem.

\LaTeX\ is good at typesetting mathematical formulas
like
       \( x-3y + z = 7 \)
or
       \( a_{1} > x^{2n} + y^{2n} > x' \)
or
       \( \ip{A}{B} = \sum_{i} a_{i} b_{i} \).
The spaces you type in a formula are
ignored.  Remember that a letter like
       $x$                   % $ ... $  and  \( ... \)  are equivalent
is a formula when it denotes a mathematical
symbol, and it should be typed as one.

\section{Displayed Text}

Text is displayed by indenting it from the left
margin.  Quotations are commonly displayed.  There
are short quotations
\begin{quote}
   This is a short quotation.  It consists of a
   single paragraph of text.  See how it is formatted.
\end{quote}
and longer ones.
\begin{quotation}
   This is a longer quotation.  It consists of two
   paragraphs of text, neither of which are
   particularly interesting.

   This is the second paragraph of the quotation.  It
   is just as dull as the first paragraph.
\end{quotation}
Another frequently-displayed structure is a list.
The following is an example of an \emph{itemized}
list.
\begin{itemize}
   \item This is the first item of an itemized list.
         Each item in the list is marked with a ``tick''.
         You don't have to worry about what kind of tick
         mark is used.

   \item This is the second item of the list.  It
         contains another list nested inside it.  The inner
         list is an \emph{enumerated} list.
         \begin{enumerate}
            \item This is the first item of an enumerated
                  list that is nested within the itemized list.

            \item This is the second item of the inner list.
                  \LaTeX\ allows you to nest lists deeper than
                  you really should.
         \end{enumerate}
         This is the rest of the second item of the outer
         list.  It is no more interesting than any other
         part of the item.
   \item This is the third item of the list.
\end{itemize}
You can even display poetry.
\begin{verse}
   There is an environment
    for verse \\             % The \\ command separates lines
   Whose features some poets % within a stanza.
   will curse.

                             % One or more blank lines separate stanzas.

   For instead of making\\
   Them do \emph{all} line breaking, \\
   It allows them to put too many words on a line when they'd rather be
   forced to be terse.
\end{verse}

Mathematical formulas may also be displayed.  A
displayed formula
is
one-line long; multiline
formulas require special formatting instructions.
   \[  \ip{\Gamma}{\psi'} = x'' + y^{2} + z_{i}^{n}\]
Don't start a paragraph with a displayed equation,
nor make one a paragraph by itself.

\end{document}               % End of document.

\section{Bench Huge Section 1591}
% synthetic bench section 1591
% This is a sample LaTeX input file.  (Version of 12 August 2004.)
%
% A '%' character causes TeX to ignore all remaining text on the line,
% and is used for comments like this one.

\documentclass{article}      % Specifies the document class

                             % The preamble begins here.
\title{An Example Document}  % Declares the document's title.
\author{Leslie Lamport}      % Declares the author's name.
\date{January 21, 1994}      % Deleting this command produces today's date.

\newcommand{\ip}[2]{(#1, #2)}
                             % Defines \ip{arg1}{arg2} to mean
                             % (arg1, arg2).

%\newcommand{\ip}[2]{\langle #1 | #2\rangle}
                             % This is an alternative definition of
                             % \ip that is commented out.

\begin{document}             % End of preamble and beginning of text.

\maketitle                   % Produces the title.

This is an example input file.  Comparing it with
the output it generates can show you how to
produce a simple document of your own.

\section{Ordinary Text}      % Produces section heading.  Lower-level
                             % sections are begun with similar
                             % \subsection and \subsubsection commands.

The ends  of words and sentences are marked
  by   spaces. It  doesn't matter how many
spaces    you type; one is as good as 100.  The
end of   a line counts as a space.

One   or more   blank lines denote the  end
of  a paragraph.

Since any number of consecutive spaces are treated
like a single one, the formatting of the input
file makes no difference to
      \LaTeX,                % The \LaTeX command generates the LaTeX logo.
but it makes a difference to you.  When you use
\LaTeX, making your input file as easy to read
as possible will be a great help as you write
your document and when you change it.  This sample
file shows how you can add comments to your own input
file.

Because printing is different from typewriting,
there are a number of things that you have to do
differently when preparing an input file than if
you were just typing the document directly.
Quotation marks like
       ``this''
have to be handled specially, as do quotes within
quotes:
       ``\,`this'            % \, separates the double and single quote.
        is what I just
        wrote, not  `that'\,''.

Dashes come in three sizes: an
       intra-word
dash, a medium dash for number ranges like
       1--2,
and a punctuation
       dash---like
this.

A sentence-ending space should be larger than the
space between words within a sentence.  You
sometimes have to type special commands in
conjunction with punctuation characters to get
this right, as in the following sentence.
       Gnats, gnus, etc.\ all  % `\ ' makes an inter-word space.
       begin with G\@.         % \@ marks end-of-sentence punctuation.
You should check the spaces after periods when
reading your output to make sure you haven't
forgotten any special cases.  Generating an
ellipsis
       \ldots\               % `\ ' is needed after `\ldots' because TeX
                             % ignores spaces after command names like \ldots
                             % made from \ + letters.
                             %
                             % Note how a `%' character causes TeX to ignore
                             % the end of the input line, so these blank lines
                             % do not start a new paragraph.
                             %
with the right spacing around the periods requires
a special command.

\LaTeX\ interprets some common characters as
commands, so you must type special commands to
generate them.  These characters include the
following:
       \$ \& \% \# \{ and \}.

In printing, text is usually emphasized with an
       \emph{italic}
type style.

\begin{em}
   A long segment of text can also be emphasized
   in this way.  Text within such a segment can be
   given \emph{additional} emphasis.
\end{em}

It is sometimes necessary to prevent \LaTeX\ from
breaking a line where it might otherwise do so.
This may be at a space, as between the ``Mr.''\ and
``Jones'' in
       ``Mr.~Jones'',        % ~ produces an unbreakable interword space.
or within a word---especially when the word is a
symbol like
       \mbox{\emph{itemnum}}
that makes little sense when hyphenated across
lines.

Footnotes\footnote{This is an example of a footnote.}
pose no problem.

\LaTeX\ is good at typesetting mathematical formulas
like
       \( x-3y + z = 7 \)
or
       \( a_{1} > x^{2n} + y^{2n} > x' \)
or
       \( \ip{A}{B} = \sum_{i} a_{i} b_{i} \).
The spaces you type in a formula are
ignored.  Remember that a letter like
       $x$                   % $ ... $  and  \( ... \)  are equivalent
is a formula when it denotes a mathematical
symbol, and it should be typed as one.

\section{Displayed Text}

Text is displayed by indenting it from the left
margin.  Quotations are commonly displayed.  There
are short quotations
\begin{quote}
   This is a short quotation.  It consists of a
   single paragraph of text.  See how it is formatted.
\end{quote}
and longer ones.
\begin{quotation}
   This is a longer quotation.  It consists of two
   paragraphs of text, neither of which are
   particularly interesting.

   This is the second paragraph of the quotation.  It
   is just as dull as the first paragraph.
\end{quotation}
Another frequently-displayed structure is a list.
The following is an example of an \emph{itemized}
list.
\begin{itemize}
   \item This is the first item of an itemized list.
         Each item in the list is marked with a ``tick''.
         You don't have to worry about what kind of tick
         mark is used.

   \item This is the second item of the list.  It
         contains another list nested inside it.  The inner
         list is an \emph{enumerated} list.
         \begin{enumerate}
            \item This is the first item of an enumerated
                  list that is nested within the itemized list.

            \item This is the second item of the inner list.
                  \LaTeX\ allows you to nest lists deeper than
                  you really should.
         \end{enumerate}
         This is the rest of the second item of the outer
         list.  It is no more interesting than any other
         part of the item.
   \item This is the third item of the list.
\end{itemize}
You can even display poetry.
\begin{verse}
   There is an environment
    for verse \\             % The \\ command separates lines
   Whose features some poets % within a stanza.
   will curse.

                             % One or more blank lines separate stanzas.

   For instead of making\\
   Them do \emph{all} line breaking, \\
   It allows them to put too many words on a line when they'd rather be
   forced to be terse.
\end{verse}

Mathematical formulas may also be displayed.  A
displayed formula
is
one-line long; multiline
formulas require special formatting instructions.
   \[  \ip{\Gamma}{\psi'} = x'' + y^{2} + z_{i}^{n}\]
Don't start a paragraph with a displayed equation,
nor make one a paragraph by itself.

\end{document}               % End of document.

\section{Bench Huge Section 1592}
% synthetic bench section 1592
% This is a sample LaTeX input file.  (Version of 12 August 2004.)
%
% A '%' character causes TeX to ignore all remaining text on the line,
% and is used for comments like this one.

\documentclass{article}      % Specifies the document class

                             % The preamble begins here.
\title{An Example Document}  % Declares the document's title.
\author{Leslie Lamport}      % Declares the author's name.
\date{January 21, 1994}      % Deleting this command produces today's date.

\newcommand{\ip}[2]{(#1, #2)}
                             % Defines \ip{arg1}{arg2} to mean
                             % (arg1, arg2).

%\newcommand{\ip}[2]{\langle #1 | #2\rangle}
                             % This is an alternative definition of
                             % \ip that is commented out.

\begin{document}             % End of preamble and beginning of text.

\maketitle                   % Produces the title.

This is an example input file.  Comparing it with
the output it generates can show you how to
produce a simple document of your own.

\section{Ordinary Text}      % Produces section heading.  Lower-level
                             % sections are begun with similar
                             % \subsection and \subsubsection commands.

The ends  of words and sentences are marked
  by   spaces. It  doesn't matter how many
spaces    you type; one is as good as 100.  The
end of   a line counts as a space.

One   or more   blank lines denote the  end
of  a paragraph.

Since any number of consecutive spaces are treated
like a single one, the formatting of the input
file makes no difference to
      \LaTeX,                % The \LaTeX command generates the LaTeX logo.
but it makes a difference to you.  When you use
\LaTeX, making your input file as easy to read
as possible will be a great help as you write
your document and when you change it.  This sample
file shows how you can add comments to your own input
file.

Because printing is different from typewriting,
there are a number of things that you have to do
differently when preparing an input file than if
you were just typing the document directly.
Quotation marks like
       ``this''
have to be handled specially, as do quotes within
quotes:
       ``\,`this'            % \, separates the double and single quote.
        is what I just
        wrote, not  `that'\,''.

Dashes come in three sizes: an
       intra-word
dash, a medium dash for number ranges like
       1--2,
and a punctuation
       dash---like
this.

A sentence-ending space should be larger than the
space between words within a sentence.  You
sometimes have to type special commands in
conjunction with punctuation characters to get
this right, as in the following sentence.
       Gnats, gnus, etc.\ all  % `\ ' makes an inter-word space.
       begin with G\@.         % \@ marks end-of-sentence punctuation.
You should check the spaces after periods when
reading your output to make sure you haven't
forgotten any special cases.  Generating an
ellipsis
       \ldots\               % `\ ' is needed after `\ldots' because TeX
                             % ignores spaces after command names like \ldots
                             % made from \ + letters.
                             %
                             % Note how a `%' character causes TeX to ignore
                             % the end of the input line, so these blank lines
                             % do not start a new paragraph.
                             %
with the right spacing around the periods requires
a special command.

\LaTeX\ interprets some common characters as
commands, so you must type special commands to
generate them.  These characters include the
following:
       \$ \& \% \# \{ and \}.

In printing, text is usually emphasized with an
       \emph{italic}
type style.

\begin{em}
   A long segment of text can also be emphasized
   in this way.  Text within such a segment can be
   given \emph{additional} emphasis.
\end{em}

It is sometimes necessary to prevent \LaTeX\ from
breaking a line where it might otherwise do so.
This may be at a space, as between the ``Mr.''\ and
``Jones'' in
       ``Mr.~Jones'',        % ~ produces an unbreakable interword space.
or within a word---especially when the word is a
symbol like
       \mbox{\emph{itemnum}}
that makes little sense when hyphenated across
lines.

Footnotes\footnote{This is an example of a footnote.}
pose no problem.

\LaTeX\ is good at typesetting mathematical formulas
like
       \( x-3y + z = 7 \)
or
       \( a_{1} > x^{2n} + y^{2n} > x' \)
or
       \( \ip{A}{B} = \sum_{i} a_{i} b_{i} \).
The spaces you type in a formula are
ignored.  Remember that a letter like
       $x$                   % $ ... $  and  \( ... \)  are equivalent
is a formula when it denotes a mathematical
symbol, and it should be typed as one.

\section{Displayed Text}

Text is displayed by indenting it from the left
margin.  Quotations are commonly displayed.  There
are short quotations
\begin{quote}
   This is a short quotation.  It consists of a
   single paragraph of text.  See how it is formatted.
\end{quote}
and longer ones.
\begin{quotation}
   This is a longer quotation.  It consists of two
   paragraphs of text, neither of which are
   particularly interesting.

   This is the second paragraph of the quotation.  It
   is just as dull as the first paragraph.
\end{quotation}
Another frequently-displayed structure is a list.
The following is an example of an \emph{itemized}
list.
\begin{itemize}
   \item This is the first item of an itemized list.
         Each item in the list is marked with a ``tick''.
         You don't have to worry about what kind of tick
         mark is used.

   \item This is the second item of the list.  It
         contains another list nested inside it.  The inner
         list is an \emph{enumerated} list.
         \begin{enumerate}
            \item This is the first item of an enumerated
                  list that is nested within the itemized list.

            \item This is the second item of the inner list.
                  \LaTeX\ allows you to nest lists deeper than
                  you really should.
         \end{enumerate}
         This is the rest of the second item of the outer
         list.  It is no more interesting than any other
         part of the item.
   \item This is the third item of the list.
\end{itemize}
You can even display poetry.
\begin{verse}
   There is an environment
    for verse \\             % The \\ command separates lines
   Whose features some poets % within a stanza.
   will curse.

                             % One or more blank lines separate stanzas.

   For instead of making\\
   Them do \emph{all} line breaking, \\
   It allows them to put too many words on a line when they'd rather be
   forced to be terse.
\end{verse}

Mathematical formulas may also be displayed.  A
displayed formula
is
one-line long; multiline
formulas require special formatting instructions.
   \[  \ip{\Gamma}{\psi'} = x'' + y^{2} + z_{i}^{n}\]
Don't start a paragraph with a displayed equation,
nor make one a paragraph by itself.

\end{document}               % End of document.

\section{Bench Huge Section 1593}
% synthetic bench section 1593
% This is a sample LaTeX input file.  (Version of 12 August 2004.)
%
% A '%' character causes TeX to ignore all remaining text on the line,
% and is used for comments like this one.

\documentclass{article}      % Specifies the document class

                             % The preamble begins here.
\title{An Example Document}  % Declares the document's title.
\author{Leslie Lamport}      % Declares the author's name.
\date{January 21, 1994}      % Deleting this command produces today's date.

\newcommand{\ip}[2]{(#1, #2)}
                             % Defines \ip{arg1}{arg2} to mean
                             % (arg1, arg2).

%\newcommand{\ip}[2]{\langle #1 | #2\rangle}
                             % This is an alternative definition of
                             % \ip that is commented out.

\begin{document}             % End of preamble and beginning of text.

\maketitle                   % Produces the title.

This is an example input file.  Comparing it with
the output it generates can show you how to
produce a simple document of your own.

\section{Ordinary Text}      % Produces section heading.  Lower-level
                             % sections are begun with similar
                             % \subsection and \subsubsection commands.

The ends  of words and sentences are marked
  by   spaces. It  doesn't matter how many
spaces    you type; one is as good as 100.  The
end of   a line counts as a space.

One   or more   blank lines denote the  end
of  a paragraph.

Since any number of consecutive spaces are treated
like a single one, the formatting of the input
file makes no difference to
      \LaTeX,                % The \LaTeX command generates the LaTeX logo.
but it makes a difference to you.  When you use
\LaTeX, making your input file as easy to read
as possible will be a great help as you write
your document and when you change it.  This sample
file shows how you can add comments to your own input
file.

Because printing is different from typewriting,
there are a number of things that you have to do
differently when preparing an input file than if
you were just typing the document directly.
Quotation marks like
       ``this''
have to be handled specially, as do quotes within
quotes:
       ``\,`this'            % \, separates the double and single quote.
        is what I just
        wrote, not  `that'\,''.

Dashes come in three sizes: an
       intra-word
dash, a medium dash for number ranges like
       1--2,
and a punctuation
       dash---like
this.

A sentence-ending space should be larger than the
space between words within a sentence.  You
sometimes have to type special commands in
conjunction with punctuation characters to get
this right, as in the following sentence.
       Gnats, gnus, etc.\ all  % `\ ' makes an inter-word space.
       begin with G\@.         % \@ marks end-of-sentence punctuation.
You should check the spaces after periods when
reading your output to make sure you haven't
forgotten any special cases.  Generating an
ellipsis
       \ldots\               % `\ ' is needed after `\ldots' because TeX
                             % ignores spaces after command names like \ldots
                             % made from \ + letters.
                             %
                             % Note how a `%' character causes TeX to ignore
                             % the end of the input line, so these blank lines
                             % do not start a new paragraph.
                             %
with the right spacing around the periods requires
a special command.

\LaTeX\ interprets some common characters as
commands, so you must type special commands to
generate them.  These characters include the
following:
       \$ \& \% \# \{ and \}.

In printing, text is usually emphasized with an
       \emph{italic}
type style.

\begin{em}
   A long segment of text can also be emphasized
   in this way.  Text within such a segment can be
   given \emph{additional} emphasis.
\end{em}

It is sometimes necessary to prevent \LaTeX\ from
breaking a line where it might otherwise do so.
This may be at a space, as between the ``Mr.''\ and
``Jones'' in
       ``Mr.~Jones'',        % ~ produces an unbreakable interword space.
or within a word---especially when the word is a
symbol like
       \mbox{\emph{itemnum}}
that makes little sense when hyphenated across
lines.

Footnotes\footnote{This is an example of a footnote.}
pose no problem.

\LaTeX\ is good at typesetting mathematical formulas
like
       \( x-3y + z = 7 \)
or
       \( a_{1} > x^{2n} + y^{2n} > x' \)
or
       \( \ip{A}{B} = \sum_{i} a_{i} b_{i} \).
The spaces you type in a formula are
ignored.  Remember that a letter like
       $x$                   % $ ... $  and  \( ... \)  are equivalent
is a formula when it denotes a mathematical
symbol, and it should be typed as one.

\section{Displayed Text}

Text is displayed by indenting it from the left
margin.  Quotations are commonly displayed.  There
are short quotations
\begin{quote}
   This is a short quotation.  It consists of a
   single paragraph of text.  See how it is formatted.
\end{quote}
and longer ones.
\begin{quotation}
   This is a longer quotation.  It consists of two
   paragraphs of text, neither of which are
   particularly interesting.

   This is the second paragraph of the quotation.  It
   is just as dull as the first paragraph.
\end{quotation}
Another frequently-displayed structure is a list.
The following is an example of an \emph{itemized}
list.
\begin{itemize}
   \item This is the first item of an itemized list.
         Each item in the list is marked with a ``tick''.
         You don't have to worry about what kind of tick
         mark is used.

   \item This is the second item of the list.  It
         contains another list nested inside it.  The inner
         list is an \emph{enumerated} list.
         \begin{enumerate}
            \item This is the first item of an enumerated
                  list that is nested within the itemized list.

            \item This is the second item of the inner list.
                  \LaTeX\ allows you to nest lists deeper than
                  you really should.
         \end{enumerate}
         This is the rest of the second item of the outer
         list.  It is no more interesting than any other
         part of the item.
   \item This is the third item of the list.
\end{itemize}
You can even display poetry.
\begin{verse}
   There is an environment
    for verse \\             % The \\ command separates lines
   Whose features some poets % within a stanza.
   will curse.

                             % One or more blank lines separate stanzas.

   For instead of making\\
   Them do \emph{all} line breaking, \\
   It allows them to put too many words on a line when they'd rather be
   forced to be terse.
\end{verse}

Mathematical formulas may also be displayed.  A
displayed formula
is
one-line long; multiline
formulas require special formatting instructions.
   \[  \ip{\Gamma}{\psi'} = x'' + y^{2} + z_{i}^{n}\]
Don't start a paragraph with a displayed equation,
nor make one a paragraph by itself.

\end{document}               % End of document.

\section{Bench Huge Section 1594}
% synthetic bench section 1594
% This is a sample LaTeX input file.  (Version of 12 August 2004.)
%
% A '%' character causes TeX to ignore all remaining text on the line,
% and is used for comments like this one.

\documentclass{article}      % Specifies the document class

                             % The preamble begins here.
\title{An Example Document}  % Declares the document's title.
\author{Leslie Lamport}      % Declares the author's name.
\date{January 21, 1994}      % Deleting this command produces today's date.

\newcommand{\ip}[2]{(#1, #2)}
                             % Defines \ip{arg1}{arg2} to mean
                             % (arg1, arg2).

%\newcommand{\ip}[2]{\langle #1 | #2\rangle}
                             % This is an alternative definition of
                             % \ip that is commented out.

\begin{document}             % End of preamble and beginning of text.

\maketitle                   % Produces the title.

This is an example input file.  Comparing it with
the output it generates can show you how to
produce a simple document of your own.

\section{Ordinary Text}      % Produces section heading.  Lower-level
                             % sections are begun with similar
                             % \subsection and \subsubsection commands.

The ends  of words and sentences are marked
  by   spaces. It  doesn't matter how many
spaces    you type; one is as good as 100.  The
end of   a line counts as a space.

One   or more   blank lines denote the  end
of  a paragraph.

Since any number of consecutive spaces are treated
like a single one, the formatting of the input
file makes no difference to
      \LaTeX,                % The \LaTeX command generates the LaTeX logo.
but it makes a difference to you.  When you use
\LaTeX, making your input file as easy to read
as possible will be a great help as you write
your document and when you change it.  This sample
file shows how you can add comments to your own input
file.

Because printing is different from typewriting,
there are a number of things that you have to do
differently when preparing an input file than if
you were just typing the document directly.
Quotation marks like
       ``this''
have to be handled specially, as do quotes within
quotes:
       ``\,`this'            % \, separates the double and single quote.
        is what I just
        wrote, not  `that'\,''.

Dashes come in three sizes: an
       intra-word
dash, a medium dash for number ranges like
       1--2,
and a punctuation
       dash---like
this.

A sentence-ending space should be larger than the
space between words within a sentence.  You
sometimes have to type special commands in
conjunction with punctuation characters to get
this right, as in the following sentence.
       Gnats, gnus, etc.\ all  % `\ ' makes an inter-word space.
       begin with G\@.         % \@ marks end-of-sentence punctuation.
You should check the spaces after periods when
reading your output to make sure you haven't
forgotten any special cases.  Generating an
ellipsis
       \ldots\               % `\ ' is needed after `\ldots' because TeX
                             % ignores spaces after command names like \ldots
                             % made from \ + letters.
                             %
                             % Note how a `%' character causes TeX to ignore
                             % the end of the input line, so these blank lines
                             % do not start a new paragraph.
                             %
with the right spacing around the periods requires
a special command.

\LaTeX\ interprets some common characters as
commands, so you must type special commands to
generate them.  These characters include the
following:
       \$ \& \% \# \{ and \}.

In printing, text is usually emphasized with an
       \emph{italic}
type style.

\begin{em}
   A long segment of text can also be emphasized
   in this way.  Text within such a segment can be
   given \emph{additional} emphasis.
\end{em}

It is sometimes necessary to prevent \LaTeX\ from
breaking a line where it might otherwise do so.
This may be at a space, as between the ``Mr.''\ and
``Jones'' in
       ``Mr.~Jones'',        % ~ produces an unbreakable interword space.
or within a word---especially when the word is a
symbol like
       \mbox{\emph{itemnum}}
that makes little sense when hyphenated across
lines.

Footnotes\footnote{This is an example of a footnote.}
pose no problem.

\LaTeX\ is good at typesetting mathematical formulas
like
       \( x-3y + z = 7 \)
or
       \( a_{1} > x^{2n} + y^{2n} > x' \)
or
       \( \ip{A}{B} = \sum_{i} a_{i} b_{i} \).
The spaces you type in a formula are
ignored.  Remember that a letter like
       $x$                   % $ ... $  and  \( ... \)  are equivalent
is a formula when it denotes a mathematical
symbol, and it should be typed as one.

\section{Displayed Text}

Text is displayed by indenting it from the left
margin.  Quotations are commonly displayed.  There
are short quotations
\begin{quote}
   This is a short quotation.  It consists of a
   single paragraph of text.  See how it is formatted.
\end{quote}
and longer ones.
\begin{quotation}
   This is a longer quotation.  It consists of two
   paragraphs of text, neither of which are
   particularly interesting.

   This is the second paragraph of the quotation.  It
   is just as dull as the first paragraph.
\end{quotation}
Another frequently-displayed structure is a list.
The following is an example of an \emph{itemized}
list.
\begin{itemize}
   \item This is the first item of an itemized list.
         Each item in the list is marked with a ``tick''.
         You don't have to worry about what kind of tick
         mark is used.

   \item This is the second item of the list.  It
         contains another list nested inside it.  The inner
         list is an \emph{enumerated} list.
         \begin{enumerate}
            \item This is the first item of an enumerated
                  list that is nested within the itemized list.

            \item This is the second item of the inner list.
                  \LaTeX\ allows you to nest lists deeper than
                  you really should.
         \end{enumerate}
         This is the rest of the second item of the outer
         list.  It is no more interesting than any other
         part of the item.
   \item This is the third item of the list.
\end{itemize}
You can even display poetry.
\begin{verse}
   There is an environment
    for verse \\             % The \\ command separates lines
   Whose features some poets % within a stanza.
   will curse.

                             % One or more blank lines separate stanzas.

   For instead of making\\
   Them do \emph{all} line breaking, \\
   It allows them to put too many words on a line when they'd rather be
   forced to be terse.
\end{verse}

Mathematical formulas may also be displayed.  A
displayed formula
is
one-line long; multiline
formulas require special formatting instructions.
   \[  \ip{\Gamma}{\psi'} = x'' + y^{2} + z_{i}^{n}\]
Don't start a paragraph with a displayed equation,
nor make one a paragraph by itself.

\end{document}               % End of document.

\section{Bench Huge Section 1595}
% synthetic bench section 1595
% This is a sample LaTeX input file.  (Version of 12 August 2004.)
%
% A '%' character causes TeX to ignore all remaining text on the line,
% and is used for comments like this one.

\documentclass{article}      % Specifies the document class

                             % The preamble begins here.
\title{An Example Document}  % Declares the document's title.
\author{Leslie Lamport}      % Declares the author's name.
\date{January 21, 1994}      % Deleting this command produces today's date.

\newcommand{\ip}[2]{(#1, #2)}
                             % Defines \ip{arg1}{arg2} to mean
                             % (arg1, arg2).

%\newcommand{\ip}[2]{\langle #1 | #2\rangle}
                             % This is an alternative definition of
                             % \ip that is commented out.

\begin{document}             % End of preamble and beginning of text.

\maketitle                   % Produces the title.

This is an example input file.  Comparing it with
the output it generates can show you how to
produce a simple document of your own.

\section{Ordinary Text}      % Produces section heading.  Lower-level
                             % sections are begun with similar
                             % \subsection and \subsubsection commands.

The ends  of words and sentences are marked
  by   spaces. It  doesn't matter how many
spaces    you type; one is as good as 100.  The
end of   a line counts as a space.

One   or more   blank lines denote the  end
of  a paragraph.

Since any number of consecutive spaces are treated
like a single one, the formatting of the input
file makes no difference to
      \LaTeX,                % The \LaTeX command generates the LaTeX logo.
but it makes a difference to you.  When you use
\LaTeX, making your input file as easy to read
as possible will be a great help as you write
your document and when you change it.  This sample
file shows how you can add comments to your own input
file.

Because printing is different from typewriting,
there are a number of things that you have to do
differently when preparing an input file than if
you were just typing the document directly.
Quotation marks like
       ``this''
have to be handled specially, as do quotes within
quotes:
       ``\,`this'            % \, separates the double and single quote.
        is what I just
        wrote, not  `that'\,''.

Dashes come in three sizes: an
       intra-word
dash, a medium dash for number ranges like
       1--2,
and a punctuation
       dash---like
this.

A sentence-ending space should be larger than the
space between words within a sentence.  You
sometimes have to type special commands in
conjunction with punctuation characters to get
this right, as in the following sentence.
       Gnats, gnus, etc.\ all  % `\ ' makes an inter-word space.
       begin with G\@.         % \@ marks end-of-sentence punctuation.
You should check the spaces after periods when
reading your output to make sure you haven't
forgotten any special cases.  Generating an
ellipsis
       \ldots\               % `\ ' is needed after `\ldots' because TeX
                             % ignores spaces after command names like \ldots
                             % made from \ + letters.
                             %
                             % Note how a `%' character causes TeX to ignore
                             % the end of the input line, so these blank lines
                             % do not start a new paragraph.
                             %
with the right spacing around the periods requires
a special command.

\LaTeX\ interprets some common characters as
commands, so you must type special commands to
generate them.  These characters include the
following:
       \$ \& \% \# \{ and \}.

In printing, text is usually emphasized with an
       \emph{italic}
type style.

\begin{em}
   A long segment of text can also be emphasized
   in this way.  Text within such a segment can be
   given \emph{additional} emphasis.
\end{em}

It is sometimes necessary to prevent \LaTeX\ from
breaking a line where it might otherwise do so.
This may be at a space, as between the ``Mr.''\ and
``Jones'' in
       ``Mr.~Jones'',        % ~ produces an unbreakable interword space.
or within a word---especially when the word is a
symbol like
       \mbox{\emph{itemnum}}
that makes little sense when hyphenated across
lines.

Footnotes\footnote{This is an example of a footnote.}
pose no problem.

\LaTeX\ is good at typesetting mathematical formulas
like
       \( x-3y + z = 7 \)
or
       \( a_{1} > x^{2n} + y^{2n} > x' \)
or
       \( \ip{A}{B} = \sum_{i} a_{i} b_{i} \).
The spaces you type in a formula are
ignored.  Remember that a letter like
       $x$                   % $ ... $  and  \( ... \)  are equivalent
is a formula when it denotes a mathematical
symbol, and it should be typed as one.

\section{Displayed Text}

Text is displayed by indenting it from the left
margin.  Quotations are commonly displayed.  There
are short quotations
\begin{quote}
   This is a short quotation.  It consists of a
   single paragraph of text.  See how it is formatted.
\end{quote}
and longer ones.
\begin{quotation}
   This is a longer quotation.  It consists of two
   paragraphs of text, neither of which are
   particularly interesting.

   This is the second paragraph of the quotation.  It
   is just as dull as the first paragraph.
\end{quotation}
Another frequently-displayed structure is a list.
The following is an example of an \emph{itemized}
list.
\begin{itemize}
   \item This is the first item of an itemized list.
         Each item in the list is marked with a ``tick''.
         You don't have to worry about what kind of tick
         mark is used.

   \item This is the second item of the list.  It
         contains another list nested inside it.  The inner
         list is an \emph{enumerated} list.
         \begin{enumerate}
            \item This is the first item of an enumerated
                  list that is nested within the itemized list.

            \item This is the second item of the inner list.
                  \LaTeX\ allows you to nest lists deeper than
                  you really should.
         \end{enumerate}
         This is the rest of the second item of the outer
         list.  It is no more interesting than any other
         part of the item.
   \item This is the third item of the list.
\end{itemize}
You can even display poetry.
\begin{verse}
   There is an environment
    for verse \\             % The \\ command separates lines
   Whose features some poets % within a stanza.
   will curse.

                             % One or more blank lines separate stanzas.

   For instead of making\\
   Them do \emph{all} line breaking, \\
   It allows them to put too many words on a line when they'd rather be
   forced to be terse.
\end{verse}

Mathematical formulas may also be displayed.  A
displayed formula
is
one-line long; multiline
formulas require special formatting instructions.
   \[  \ip{\Gamma}{\psi'} = x'' + y^{2} + z_{i}^{n}\]
Don't start a paragraph with a displayed equation,
nor make one a paragraph by itself.

\end{document}               % End of document.

\section{Bench Huge Section 1596}
% synthetic bench section 1596
% This is a sample LaTeX input file.  (Version of 12 August 2004.)
%
% A '%' character causes TeX to ignore all remaining text on the line,
% and is used for comments like this one.

\documentclass{article}      % Specifies the document class

                             % The preamble begins here.
\title{An Example Document}  % Declares the document's title.
\author{Leslie Lamport}      % Declares the author's name.
\date{January 21, 1994}      % Deleting this command produces today's date.

\newcommand{\ip}[2]{(#1, #2)}
                             % Defines \ip{arg1}{arg2} to mean
                             % (arg1, arg2).

%\newcommand{\ip}[2]{\langle #1 | #2\rangle}
                             % This is an alternative definition of
                             % \ip that is commented out.

\begin{document}             % End of preamble and beginning of text.

\maketitle                   % Produces the title.

This is an example input file.  Comparing it with
the output it generates can show you how to
produce a simple document of your own.

\section{Ordinary Text}      % Produces section heading.  Lower-level
                             % sections are begun with similar
                             % \subsection and \subsubsection commands.

The ends  of words and sentences are marked
  by   spaces. It  doesn't matter how many
spaces    you type; one is as good as 100.  The
end of   a line counts as a space.

One   or more   blank lines denote the  end
of  a paragraph.

Since any number of consecutive spaces are treated
like a single one, the formatting of the input
file makes no difference to
      \LaTeX,                % The \LaTeX command generates the LaTeX logo.
but it makes a difference to you.  When you use
\LaTeX, making your input file as easy to read
as possible will be a great help as you write
your document and when you change it.  This sample
file shows how you can add comments to your own input
file.

Because printing is different from typewriting,
there are a number of things that you have to do
differently when preparing an input file than if
you were just typing the document directly.
Quotation marks like
       ``this''
have to be handled specially, as do quotes within
quotes:
       ``\,`this'            % \, separates the double and single quote.
        is what I just
        wrote, not  `that'\,''.

Dashes come in three sizes: an
       intra-word
dash, a medium dash for number ranges like
       1--2,
and a punctuation
       dash---like
this.

A sentence-ending space should be larger than the
space between words within a sentence.  You
sometimes have to type special commands in
conjunction with punctuation characters to get
this right, as in the following sentence.
       Gnats, gnus, etc.\ all  % `\ ' makes an inter-word space.
       begin with G\@.         % \@ marks end-of-sentence punctuation.
You should check the spaces after periods when
reading your output to make sure you haven't
forgotten any special cases.  Generating an
ellipsis
       \ldots\               % `\ ' is needed after `\ldots' because TeX
                             % ignores spaces after command names like \ldots
                             % made from \ + letters.
                             %
                             % Note how a `%' character causes TeX to ignore
                             % the end of the input line, so these blank lines
                             % do not start a new paragraph.
                             %
with the right spacing around the periods requires
a special command.

\LaTeX\ interprets some common characters as
commands, so you must type special commands to
generate them.  These characters include the
following:
       \$ \& \% \# \{ and \}.

In printing, text is usually emphasized with an
       \emph{italic}
type style.

\begin{em}
   A long segment of text can also be emphasized
   in this way.  Text within such a segment can be
   given \emph{additional} emphasis.
\end{em}

It is sometimes necessary to prevent \LaTeX\ from
breaking a line where it might otherwise do so.
This may be at a space, as between the ``Mr.''\ and
``Jones'' in
       ``Mr.~Jones'',        % ~ produces an unbreakable interword space.
or within a word---especially when the word is a
symbol like
       \mbox{\emph{itemnum}}
that makes little sense when hyphenated across
lines.

Footnotes\footnote{This is an example of a footnote.}
pose no problem.

\LaTeX\ is good at typesetting mathematical formulas
like
       \( x-3y + z = 7 \)
or
       \( a_{1} > x^{2n} + y^{2n} > x' \)
or
       \( \ip{A}{B} = \sum_{i} a_{i} b_{i} \).
The spaces you type in a formula are
ignored.  Remember that a letter like
       $x$                   % $ ... $  and  \( ... \)  are equivalent
is a formula when it denotes a mathematical
symbol, and it should be typed as one.

\section{Displayed Text}

Text is displayed by indenting it from the left
margin.  Quotations are commonly displayed.  There
are short quotations
\begin{quote}
   This is a short quotation.  It consists of a
   single paragraph of text.  See how it is formatted.
\end{quote}
and longer ones.
\begin{quotation}
   This is a longer quotation.  It consists of two
   paragraphs of text, neither of which are
   particularly interesting.

   This is the second paragraph of the quotation.  It
   is just as dull as the first paragraph.
\end{quotation}
Another frequently-displayed structure is a list.
The following is an example of an \emph{itemized}
list.
\begin{itemize}
   \item This is the first item of an itemized list.
         Each item in the list is marked with a ``tick''.
         You don't have to worry about what kind of tick
         mark is used.

   \item This is the second item of the list.  It
         contains another list nested inside it.  The inner
         list is an \emph{enumerated} list.
         \begin{enumerate}
            \item This is the first item of an enumerated
                  list that is nested within the itemized list.

            \item This is the second item of the inner list.
                  \LaTeX\ allows you to nest lists deeper than
                  you really should.
         \end{enumerate}
         This is the rest of the second item of the outer
         list.  It is no more interesting than any other
         part of the item.
   \item This is the third item of the list.
\end{itemize}
You can even display poetry.
\begin{verse}
   There is an environment
    for verse \\             % The \\ command separates lines
   Whose features some poets % within a stanza.
   will curse.

                             % One or more blank lines separate stanzas.

   For instead of making\\
   Them do \emph{all} line breaking, \\
   It allows them to put too many words on a line when they'd rather be
   forced to be terse.
\end{verse}

Mathematical formulas may also be displayed.  A
displayed formula
is
one-line long; multiline
formulas require special formatting instructions.
   \[  \ip{\Gamma}{\psi'} = x'' + y^{2} + z_{i}^{n}\]
Don't start a paragraph with a displayed equation,
nor make one a paragraph by itself.

\end{document}               % End of document.

\section{Bench Huge Section 1597}
% synthetic bench section 1597
% This is a sample LaTeX input file.  (Version of 12 August 2004.)
%
% A '%' character causes TeX to ignore all remaining text on the line,
% and is used for comments like this one.

\documentclass{article}      % Specifies the document class

                             % The preamble begins here.
\title{An Example Document}  % Declares the document's title.
\author{Leslie Lamport}      % Declares the author's name.
\date{January 21, 1994}      % Deleting this command produces today's date.

\newcommand{\ip}[2]{(#1, #2)}
                             % Defines \ip{arg1}{arg2} to mean
                             % (arg1, arg2).

%\newcommand{\ip}[2]{\langle #1 | #2\rangle}
                             % This is an alternative definition of
                             % \ip that is commented out.

\begin{document}             % End of preamble and beginning of text.

\maketitle                   % Produces the title.

This is an example input file.  Comparing it with
the output it generates can show you how to
produce a simple document of your own.

\section{Ordinary Text}      % Produces section heading.  Lower-level
                             % sections are begun with similar
                             % \subsection and \subsubsection commands.

The ends  of words and sentences are marked
  by   spaces. It  doesn't matter how many
spaces    you type; one is as good as 100.  The
end of   a line counts as a space.

One   or more   blank lines denote the  end
of  a paragraph.

Since any number of consecutive spaces are treated
like a single one, the formatting of the input
file makes no difference to
      \LaTeX,                % The \LaTeX command generates the LaTeX logo.
but it makes a difference to you.  When you use
\LaTeX, making your input file as easy to read
as possible will be a great help as you write
your document and when you change it.  This sample
file shows how you can add comments to your own input
file.

Because printing is different from typewriting,
there are a number of things that you have to do
differently when preparing an input file than if
you were just typing the document directly.
Quotation marks like
       ``this''
have to be handled specially, as do quotes within
quotes:
       ``\,`this'            % \, separates the double and single quote.
        is what I just
        wrote, not  `that'\,''.

Dashes come in three sizes: an
       intra-word
dash, a medium dash for number ranges like
       1--2,
and a punctuation
       dash---like
this.

A sentence-ending space should be larger than the
space between words within a sentence.  You
sometimes have to type special commands in
conjunction with punctuation characters to get
this right, as in the following sentence.
       Gnats, gnus, etc.\ all  % `\ ' makes an inter-word space.
       begin with G\@.         % \@ marks end-of-sentence punctuation.
You should check the spaces after periods when
reading your output to make sure you haven't
forgotten any special cases.  Generating an
ellipsis
       \ldots\               % `\ ' is needed after `\ldots' because TeX
                             % ignores spaces after command names like \ldots
                             % made from \ + letters.
                             %
                             % Note how a `%' character causes TeX to ignore
                             % the end of the input line, so these blank lines
                             % do not start a new paragraph.
                             %
with the right spacing around the periods requires
a special command.

\LaTeX\ interprets some common characters as
commands, so you must type special commands to
generate them.  These characters include the
following:
       \$ \& \% \# \{ and \}.

In printing, text is usually emphasized with an
       \emph{italic}
type style.

\begin{em}
   A long segment of text can also be emphasized
   in this way.  Text within such a segment can be
   given \emph{additional} emphasis.
\end{em}

It is sometimes necessary to prevent \LaTeX\ from
breaking a line where it might otherwise do so.
This may be at a space, as between the ``Mr.''\ and
``Jones'' in
       ``Mr.~Jones'',        % ~ produces an unbreakable interword space.
or within a word---especially when the word is a
symbol like
       \mbox{\emph{itemnum}}
that makes little sense when hyphenated across
lines.

Footnotes\footnote{This is an example of a footnote.}
pose no problem.

\LaTeX\ is good at typesetting mathematical formulas
like
       \( x-3y + z = 7 \)
or
       \( a_{1} > x^{2n} + y^{2n} > x' \)
or
       \( \ip{A}{B} = \sum_{i} a_{i} b_{i} \).
The spaces you type in a formula are
ignored.  Remember that a letter like
       $x$                   % $ ... $  and  \( ... \)  are equivalent
is a formula when it denotes a mathematical
symbol, and it should be typed as one.

\section{Displayed Text}

Text is displayed by indenting it from the left
margin.  Quotations are commonly displayed.  There
are short quotations
\begin{quote}
   This is a short quotation.  It consists of a
   single paragraph of text.  See how it is formatted.
\end{quote}
and longer ones.
\begin{quotation}
   This is a longer quotation.  It consists of two
   paragraphs of text, neither of which are
   particularly interesting.

   This is the second paragraph of the quotation.  It
   is just as dull as the first paragraph.
\end{quotation}
Another frequently-displayed structure is a list.
The following is an example of an \emph{itemized}
list.
\begin{itemize}
   \item This is the first item of an itemized list.
         Each item in the list is marked with a ``tick''.
         You don't have to worry about what kind of tick
         mark is used.

   \item This is the second item of the list.  It
         contains another list nested inside it.  The inner
         list is an \emph{enumerated} list.
         \begin{enumerate}
            \item This is the first item of an enumerated
                  list that is nested within the itemized list.

            \item This is the second item of the inner list.
                  \LaTeX\ allows you to nest lists deeper than
                  you really should.
         \end{enumerate}
         This is the rest of the second item of the outer
         list.  It is no more interesting than any other
         part of the item.
   \item This is the third item of the list.
\end{itemize}
You can even display poetry.
\begin{verse}
   There is an environment
    for verse \\             % The \\ command separates lines
   Whose features some poets % within a stanza.
   will curse.

                             % One or more blank lines separate stanzas.

   For instead of making\\
   Them do \emph{all} line breaking, \\
   It allows them to put too many words on a line when they'd rather be
   forced to be terse.
\end{verse}

Mathematical formulas may also be displayed.  A
displayed formula
is
one-line long; multiline
formulas require special formatting instructions.
   \[  \ip{\Gamma}{\psi'} = x'' + y^{2} + z_{i}^{n}\]
Don't start a paragraph with a displayed equation,
nor make one a paragraph by itself.

\end{document}               % End of document.

\section{Bench Huge Section 1598}
% synthetic bench section 1598
% This is a sample LaTeX input file.  (Version of 12 August 2004.)
%
% A '%' character causes TeX to ignore all remaining text on the line,
% and is used for comments like this one.

\documentclass{article}      % Specifies the document class

                             % The preamble begins here.
\title{An Example Document}  % Declares the document's title.
\author{Leslie Lamport}      % Declares the author's name.
\date{January 21, 1994}      % Deleting this command produces today's date.

\newcommand{\ip}[2]{(#1, #2)}
                             % Defines \ip{arg1}{arg2} to mean
                             % (arg1, arg2).

%\newcommand{\ip}[2]{\langle #1 | #2\rangle}
                             % This is an alternative definition of
                             % \ip that is commented out.

\begin{document}             % End of preamble and beginning of text.

\maketitle                   % Produces the title.

This is an example input file.  Comparing it with
the output it generates can show you how to
produce a simple document of your own.

\section{Ordinary Text}      % Produces section heading.  Lower-level
                             % sections are begun with similar
                             % \subsection and \subsubsection commands.

The ends  of words and sentences are marked
  by   spaces. It  doesn't matter how many
spaces    you type; one is as good as 100.  The
end of   a line counts as a space.

One   or more   blank lines denote the  end
of  a paragraph.

Since any number of consecutive spaces are treated
like a single one, the formatting of the input
file makes no difference to
      \LaTeX,                % The \LaTeX command generates the LaTeX logo.
but it makes a difference to you.  When you use
\LaTeX, making your input file as easy to read
as possible will be a great help as you write
your document and when you change it.  This sample
file shows how you can add comments to your own input
file.

Because printing is different from typewriting,
there are a number of things that you have to do
differently when preparing an input file than if
you were just typing the document directly.
Quotation marks like
       ``this''
have to be handled specially, as do quotes within
quotes:
       ``\,`this'            % \, separates the double and single quote.
        is what I just
        wrote, not  `that'\,''.

Dashes come in three sizes: an
       intra-word
dash, a medium dash for number ranges like
       1--2,
and a punctuation
       dash---like
this.

A sentence-ending space should be larger than the
space between words within a sentence.  You
sometimes have to type special commands in
conjunction with punctuation characters to get
this right, as in the following sentence.
       Gnats, gnus, etc.\ all  % `\ ' makes an inter-word space.
       begin with G\@.         % \@ marks end-of-sentence punctuation.
You should check the spaces after periods when
reading your output to make sure you haven't
forgotten any special cases.  Generating an
ellipsis
       \ldots\               % `\ ' is needed after `\ldots' because TeX
                             % ignores spaces after command names like \ldots
                             % made from \ + letters.
                             %
                             % Note how a `%' character causes TeX to ignore
                             % the end of the input line, so these blank lines
                             % do not start a new paragraph.
                             %
with the right spacing around the periods requires
a special command.

\LaTeX\ interprets some common characters as
commands, so you must type special commands to
generate them.  These characters include the
following:
       \$ \& \% \# \{ and \}.

In printing, text is usually emphasized with an
       \emph{italic}
type style.

\begin{em}
   A long segment of text can also be emphasized
   in this way.  Text within such a segment can be
   given \emph{additional} emphasis.
\end{em}

It is sometimes necessary to prevent \LaTeX\ from
breaking a line where it might otherwise do so.
This may be at a space, as between the ``Mr.''\ and
``Jones'' in
       ``Mr.~Jones'',        % ~ produces an unbreakable interword space.
or within a word---especially when the word is a
symbol like
       \mbox{\emph{itemnum}}
that makes little sense when hyphenated across
lines.

Footnotes\footnote{This is an example of a footnote.}
pose no problem.

\LaTeX\ is good at typesetting mathematical formulas
like
       \( x-3y + z = 7 \)
or
       \( a_{1} > x^{2n} + y^{2n} > x' \)
or
       \( \ip{A}{B} = \sum_{i} a_{i} b_{i} \).
The spaces you type in a formula are
ignored.  Remember that a letter like
       $x$                   % $ ... $  and  \( ... \)  are equivalent
is a formula when it denotes a mathematical
symbol, and it should be typed as one.

\section{Displayed Text}

Text is displayed by indenting it from the left
margin.  Quotations are commonly displayed.  There
are short quotations
\begin{quote}
   This is a short quotation.  It consists of a
   single paragraph of text.  See how it is formatted.
\end{quote}
and longer ones.
\begin{quotation}
   This is a longer quotation.  It consists of two
   paragraphs of text, neither of which are
   particularly interesting.

   This is the second paragraph of the quotation.  It
   is just as dull as the first paragraph.
\end{quotation}
Another frequently-displayed structure is a list.
The following is an example of an \emph{itemized}
list.
\begin{itemize}
   \item This is the first item of an itemized list.
         Each item in the list is marked with a ``tick''.
         You don't have to worry about what kind of tick
         mark is used.

   \item This is the second item of the list.  It
         contains another list nested inside it.  The inner
         list is an \emph{enumerated} list.
         \begin{enumerate}
            \item This is the first item of an enumerated
                  list that is nested within the itemized list.

            \item This is the second item of the inner list.
                  \LaTeX\ allows you to nest lists deeper than
                  you really should.
         \end{enumerate}
         This is the rest of the second item of the outer
         list.  It is no more interesting than any other
         part of the item.
   \item This is the third item of the list.
\end{itemize}
You can even display poetry.
\begin{verse}
   There is an environment
    for verse \\             % The \\ command separates lines
   Whose features some poets % within a stanza.
   will curse.

                             % One or more blank lines separate stanzas.

   For instead of making\\
   Them do \emph{all} line breaking, \\
   It allows them to put too many words on a line when they'd rather be
   forced to be terse.
\end{verse}

Mathematical formulas may also be displayed.  A
displayed formula
is
one-line long; multiline
formulas require special formatting instructions.
   \[  \ip{\Gamma}{\psi'} = x'' + y^{2} + z_{i}^{n}\]
Don't start a paragraph with a displayed equation,
nor make one a paragraph by itself.

\end{document}               % End of document.

\section{Bench Huge Section 1599}
% synthetic bench section 1599
% This is a sample LaTeX input file.  (Version of 12 August 2004.)
%
% A '%' character causes TeX to ignore all remaining text on the line,
% and is used for comments like this one.

\documentclass{article}      % Specifies the document class

                             % The preamble begins here.
\title{An Example Document}  % Declares the document's title.
\author{Leslie Lamport}      % Declares the author's name.
\date{January 21, 1994}      % Deleting this command produces today's date.

\newcommand{\ip}[2]{(#1, #2)}
                             % Defines \ip{arg1}{arg2} to mean
                             % (arg1, arg2).

%\newcommand{\ip}[2]{\langle #1 | #2\rangle}
                             % This is an alternative definition of
                             % \ip that is commented out.

\begin{document}             % End of preamble and beginning of text.

\maketitle                   % Produces the title.

This is an example input file.  Comparing it with
the output it generates can show you how to
produce a simple document of your own.

\section{Ordinary Text}      % Produces section heading.  Lower-level
                             % sections are begun with similar
                             % \subsection and \subsubsection commands.

The ends  of words and sentences are marked
  by   spaces. It  doesn't matter how many
spaces    you type; one is as good as 100.  The
end of   a line counts as a space.

One   or more   blank lines denote the  end
of  a paragraph.

Since any number of consecutive spaces are treated
like a single one, the formatting of the input
file makes no difference to
      \LaTeX,                % The \LaTeX command generates the LaTeX logo.
but it makes a difference to you.  When you use
\LaTeX, making your input file as easy to read
as possible will be a great help as you write
your document and when you change it.  This sample
file shows how you can add comments to your own input
file.

Because printing is different from typewriting,
there are a number of things that you have to do
differently when preparing an input file than if
you were just typing the document directly.
Quotation marks like
       ``this''
have to be handled specially, as do quotes within
quotes:
       ``\,`this'            % \, separates the double and single quote.
        is what I just
        wrote, not  `that'\,''.

Dashes come in three sizes: an
       intra-word
dash, a medium dash for number ranges like
       1--2,
and a punctuation
       dash---like
this.

A sentence-ending space should be larger than the
space between words within a sentence.  You
sometimes have to type special commands in
conjunction with punctuation characters to get
this right, as in the following sentence.
       Gnats, gnus, etc.\ all  % `\ ' makes an inter-word space.
       begin with G\@.         % \@ marks end-of-sentence punctuation.
You should check the spaces after periods when
reading your output to make sure you haven't
forgotten any special cases.  Generating an
ellipsis
       \ldots\               % `\ ' is needed after `\ldots' because TeX
                             % ignores spaces after command names like \ldots
                             % made from \ + letters.
                             %
                             % Note how a `%' character causes TeX to ignore
                             % the end of the input line, so these blank lines
                             % do not start a new paragraph.
                             %
with the right spacing around the periods requires
a special command.

\LaTeX\ interprets some common characters as
commands, so you must type special commands to
generate them.  These characters include the
following:
       \$ \& \% \# \{ and \}.

In printing, text is usually emphasized with an
       \emph{italic}
type style.

\begin{em}
   A long segment of text can also be emphasized
   in this way.  Text within such a segment can be
   given \emph{additional} emphasis.
\end{em}

It is sometimes necessary to prevent \LaTeX\ from
breaking a line where it might otherwise do so.
This may be at a space, as between the ``Mr.''\ and
``Jones'' in
       ``Mr.~Jones'',        % ~ produces an unbreakable interword space.
or within a word---especially when the word is a
symbol like
       \mbox{\emph{itemnum}}
that makes little sense when hyphenated across
lines.

Footnotes\footnote{This is an example of a footnote.}
pose no problem.

\LaTeX\ is good at typesetting mathematical formulas
like
       \( x-3y + z = 7 \)
or
       \( a_{1} > x^{2n} + y^{2n} > x' \)
or
       \( \ip{A}{B} = \sum_{i} a_{i} b_{i} \).
The spaces you type in a formula are
ignored.  Remember that a letter like
       $x$                   % $ ... $  and  \( ... \)  are equivalent
is a formula when it denotes a mathematical
symbol, and it should be typed as one.

\section{Displayed Text}

Text is displayed by indenting it from the left
margin.  Quotations are commonly displayed.  There
are short quotations
\begin{quote}
   This is a short quotation.  It consists of a
   single paragraph of text.  See how it is formatted.
\end{quote}
and longer ones.
\begin{quotation}
   This is a longer quotation.  It consists of two
   paragraphs of text, neither of which are
   particularly interesting.

   This is the second paragraph of the quotation.  It
   is just as dull as the first paragraph.
\end{quotation}
Another frequently-displayed structure is a list.
The following is an example of an \emph{itemized}
list.
\begin{itemize}
   \item This is the first item of an itemized list.
         Each item in the list is marked with a ``tick''.
         You don't have to worry about what kind of tick
         mark is used.

   \item This is the second item of the list.  It
         contains another list nested inside it.  The inner
         list is an \emph{enumerated} list.
         \begin{enumerate}
            \item This is the first item of an enumerated
                  list that is nested within the itemized list.

            \item This is the second item of the inner list.
                  \LaTeX\ allows you to nest lists deeper than
                  you really should.
         \end{enumerate}
         This is the rest of the second item of the outer
         list.  It is no more interesting than any other
         part of the item.
   \item This is the third item of the list.
\end{itemize}
You can even display poetry.
\begin{verse}
   There is an environment
    for verse \\             % The \\ command separates lines
   Whose features some poets % within a stanza.
   will curse.

                             % One or more blank lines separate stanzas.

   For instead of making\\
   Them do \emph{all} line breaking, \\
   It allows them to put too many words on a line when they'd rather be
   forced to be terse.
\end{verse}

Mathematical formulas may also be displayed.  A
displayed formula
is
one-line long; multiline
formulas require special formatting instructions.
   \[  \ip{\Gamma}{\psi'} = x'' + y^{2} + z_{i}^{n}\]
Don't start a paragraph with a displayed equation,
nor make one a paragraph by itself.

\end{document}               % End of document.

\section{Bench Huge Section 1600}
% synthetic bench section 1600
% This is a sample LaTeX input file.  (Version of 12 August 2004.)
%
% A '%' character causes TeX to ignore all remaining text on the line,
% and is used for comments like this one.

\documentclass{article}      % Specifies the document class

                             % The preamble begins here.
\title{An Example Document}  % Declares the document's title.
\author{Leslie Lamport}      % Declares the author's name.
\date{January 21, 1994}      % Deleting this command produces today's date.

\newcommand{\ip}[2]{(#1, #2)}
                             % Defines \ip{arg1}{arg2} to mean
                             % (arg1, arg2).

%\newcommand{\ip}[2]{\langle #1 | #2\rangle}
                             % This is an alternative definition of
                             % \ip that is commented out.

\begin{document}             % End of preamble and beginning of text.

\maketitle                   % Produces the title.

This is an example input file.  Comparing it with
the output it generates can show you how to
produce a simple document of your own.

\section{Ordinary Text}      % Produces section heading.  Lower-level
                             % sections are begun with similar
                             % \subsection and \subsubsection commands.

The ends  of words and sentences are marked
  by   spaces. It  doesn't matter how many
spaces    you type; one is as good as 100.  The
end of   a line counts as a space.

One   or more   blank lines denote the  end
of  a paragraph.

Since any number of consecutive spaces are treated
like a single one, the formatting of the input
file makes no difference to
      \LaTeX,                % The \LaTeX command generates the LaTeX logo.
but it makes a difference to you.  When you use
\LaTeX, making your input file as easy to read
as possible will be a great help as you write
your document and when you change it.  This sample
file shows how you can add comments to your own input
file.

Because printing is different from typewriting,
there are a number of things that you have to do
differently when preparing an input file than if
you were just typing the document directly.
Quotation marks like
       ``this''
have to be handled specially, as do quotes within
quotes:
       ``\,`this'            % \, separates the double and single quote.
        is what I just
        wrote, not  `that'\,''.

Dashes come in three sizes: an
       intra-word
dash, a medium dash for number ranges like
       1--2,
and a punctuation
       dash---like
this.

A sentence-ending space should be larger than the
space between words within a sentence.  You
sometimes have to type special commands in
conjunction with punctuation characters to get
this right, as in the following sentence.
       Gnats, gnus, etc.\ all  % `\ ' makes an inter-word space.
       begin with G\@.         % \@ marks end-of-sentence punctuation.
You should check the spaces after periods when
reading your output to make sure you haven't
forgotten any special cases.  Generating an
ellipsis
       \ldots\               % `\ ' is needed after `\ldots' because TeX
                             % ignores spaces after command names like \ldots
                             % made from \ + letters.
                             %
                             % Note how a `%' character causes TeX to ignore
                             % the end of the input line, so these blank lines
                             % do not start a new paragraph.
                             %
with the right spacing around the periods requires
a special command.

\LaTeX\ interprets some common characters as
commands, so you must type special commands to
generate them.  These characters include the
following:
       \$ \& \% \# \{ and \}.

In printing, text is usually emphasized with an
       \emph{italic}
type style.

\begin{em}
   A long segment of text can also be emphasized
   in this way.  Text within such a segment can be
   given \emph{additional} emphasis.
\end{em}

It is sometimes necessary to prevent \LaTeX\ from
breaking a line where it might otherwise do so.
This may be at a space, as between the ``Mr.''\ and
``Jones'' in
       ``Mr.~Jones'',        % ~ produces an unbreakable interword space.
or within a word---especially when the word is a
symbol like
       \mbox{\emph{itemnum}}
that makes little sense when hyphenated across
lines.

Footnotes\footnote{This is an example of a footnote.}
pose no problem.

\LaTeX\ is good at typesetting mathematical formulas
like
       \( x-3y + z = 7 \)
or
       \( a_{1} > x^{2n} + y^{2n} > x' \)
or
       \( \ip{A}{B} = \sum_{i} a_{i} b_{i} \).
The spaces you type in a formula are
ignored.  Remember that a letter like
       $x$                   % $ ... $  and  \( ... \)  are equivalent
is a formula when it denotes a mathematical
symbol, and it should be typed as one.

\section{Displayed Text}

Text is displayed by indenting it from the left
margin.  Quotations are commonly displayed.  There
are short quotations
\begin{quote}
   This is a short quotation.  It consists of a
   single paragraph of text.  See how it is formatted.
\end{quote}
and longer ones.
\begin{quotation}
   This is a longer quotation.  It consists of two
   paragraphs of text, neither of which are
   particularly interesting.

   This is the second paragraph of the quotation.  It
   is just as dull as the first paragraph.
\end{quotation}
Another frequently-displayed structure is a list.
The following is an example of an \emph{itemized}
list.
\begin{itemize}
   \item This is the first item of an itemized list.
         Each item in the list is marked with a ``tick''.
         You don't have to worry about what kind of tick
         mark is used.

   \item This is the second item of the list.  It
         contains another list nested inside it.  The inner
         list is an \emph{enumerated} list.
         \begin{enumerate}
            \item This is the first item of an enumerated
                  list that is nested within the itemized list.

            \item This is the second item of the inner list.
                  \LaTeX\ allows you to nest lists deeper than
                  you really should.
         \end{enumerate}
         This is the rest of the second item of the outer
         list.  It is no more interesting than any other
         part of the item.
   \item This is the third item of the list.
\end{itemize}
You can even display poetry.
\begin{verse}
   There is an environment
    for verse \\             % The \\ command separates lines
   Whose features some poets % within a stanza.
   will curse.

                             % One or more blank lines separate stanzas.

   For instead of making\\
   Them do \emph{all} line breaking, \\
   It allows them to put too many words on a line when they'd rather be
   forced to be terse.
\end{verse}

Mathematical formulas may also be displayed.  A
displayed formula
is
one-line long; multiline
formulas require special formatting instructions.
   \[  \ip{\Gamma}{\psi'} = x'' + y^{2} + z_{i}^{n}\]
Don't start a paragraph with a displayed equation,
nor make one a paragraph by itself.

\end{document}               % End of document.

\section{Bench Huge Section 1601}
% synthetic bench section 1601
% This is a sample LaTeX input file.  (Version of 12 August 2004.)
%
% A '%' character causes TeX to ignore all remaining text on the line,
% and is used for comments like this one.

\documentclass{article}      % Specifies the document class

                             % The preamble begins here.
\title{An Example Document}  % Declares the document's title.
\author{Leslie Lamport}      % Declares the author's name.
\date{January 21, 1994}      % Deleting this command produces today's date.

\newcommand{\ip}[2]{(#1, #2)}
                             % Defines \ip{arg1}{arg2} to mean
                             % (arg1, arg2).

%\newcommand{\ip}[2]{\langle #1 | #2\rangle}
                             % This is an alternative definition of
                             % \ip that is commented out.

\begin{document}             % End of preamble and beginning of text.

\maketitle                   % Produces the title.

This is an example input file.  Comparing it with
the output it generates can show you how to
produce a simple document of your own.

\section{Ordinary Text}      % Produces section heading.  Lower-level
                             % sections are begun with similar
                             % \subsection and \subsubsection commands.

The ends  of words and sentences are marked
  by   spaces. It  doesn't matter how many
spaces    you type; one is as good as 100.  The
end of   a line counts as a space.

One   or more   blank lines denote the  end
of  a paragraph.

Since any number of consecutive spaces are treated
like a single one, the formatting of the input
file makes no difference to
      \LaTeX,                % The \LaTeX command generates the LaTeX logo.
but it makes a difference to you.  When you use
\LaTeX, making your input file as easy to read
as possible will be a great help as you write
your document and when you change it.  This sample
file shows how you can add comments to your own input
file.

Because printing is different from typewriting,
there are a number of things that you have to do
differently when preparing an input file than if
you were just typing the document directly.
Quotation marks like
       ``this''
have to be handled specially, as do quotes within
quotes:
       ``\,`this'            % \, separates the double and single quote.
        is what I just
        wrote, not  `that'\,''.

Dashes come in three sizes: an
       intra-word
dash, a medium dash for number ranges like
       1--2,
and a punctuation
       dash---like
this.

A sentence-ending space should be larger than the
space between words within a sentence.  You
sometimes have to type special commands in
conjunction with punctuation characters to get
this right, as in the following sentence.
       Gnats, gnus, etc.\ all  % `\ ' makes an inter-word space.
       begin with G\@.         % \@ marks end-of-sentence punctuation.
You should check the spaces after periods when
reading your output to make sure you haven't
forgotten any special cases.  Generating an
ellipsis
       \ldots\               % `\ ' is needed after `\ldots' because TeX
                             % ignores spaces after command names like \ldots
                             % made from \ + letters.
                             %
                             % Note how a `%' character causes TeX to ignore
                             % the end of the input line, so these blank lines
                             % do not start a new paragraph.
                             %
with the right spacing around the periods requires
a special command.

\LaTeX\ interprets some common characters as
commands, so you must type special commands to
generate them.  These characters include the
following:
       \$ \& \% \# \{ and \}.

In printing, text is usually emphasized with an
       \emph{italic}
type style.

\begin{em}
   A long segment of text can also be emphasized
   in this way.  Text within such a segment can be
   given \emph{additional} emphasis.
\end{em}

It is sometimes necessary to prevent \LaTeX\ from
breaking a line where it might otherwise do so.
This may be at a space, as between the ``Mr.''\ and
``Jones'' in
       ``Mr.~Jones'',        % ~ produces an unbreakable interword space.
or within a word---especially when the word is a
symbol like
       \mbox{\emph{itemnum}}
that makes little sense when hyphenated across
lines.

Footnotes\footnote{This is an example of a footnote.}
pose no problem.

\LaTeX\ is good at typesetting mathematical formulas
like
       \( x-3y + z = 7 \)
or
       \( a_{1} > x^{2n} + y^{2n} > x' \)
or
       \( \ip{A}{B} = \sum_{i} a_{i} b_{i} \).
The spaces you type in a formula are
ignored.  Remember that a letter like
       $x$                   % $ ... $  and  \( ... \)  are equivalent
is a formula when it denotes a mathematical
symbol, and it should be typed as one.

\section{Displayed Text}

Text is displayed by indenting it from the left
margin.  Quotations are commonly displayed.  There
are short quotations
\begin{quote}
   This is a short quotation.  It consists of a
   single paragraph of text.  See how it is formatted.
\end{quote}
and longer ones.
\begin{quotation}
   This is a longer quotation.  It consists of two
   paragraphs of text, neither of which are
   particularly interesting.

   This is the second paragraph of the quotation.  It
   is just as dull as the first paragraph.
\end{quotation}
Another frequently-displayed structure is a list.
The following is an example of an \emph{itemized}
list.
\begin{itemize}
   \item This is the first item of an itemized list.
         Each item in the list is marked with a ``tick''.
         You don't have to worry about what kind of tick
         mark is used.

   \item This is the second item of the list.  It
         contains another list nested inside it.  The inner
         list is an \emph{enumerated} list.
         \begin{enumerate}
            \item This is the first item of an enumerated
                  list that is nested within the itemized list.

            \item This is the second item of the inner list.
                  \LaTeX\ allows you to nest lists deeper than
                  you really should.
         \end{enumerate}
         This is the rest of the second item of the outer
         list.  It is no more interesting than any other
         part of the item.
   \item This is the third item of the list.
\end{itemize}
You can even display poetry.
\begin{verse}
   There is an environment
    for verse \\             % The \\ command separates lines
   Whose features some poets % within a stanza.
   will curse.

                             % One or more blank lines separate stanzas.

   For instead of making\\
   Them do \emph{all} line breaking, \\
   It allows them to put too many words on a line when they'd rather be
   forced to be terse.
\end{verse}

Mathematical formulas may also be displayed.  A
displayed formula
is
one-line long; multiline
formulas require special formatting instructions.
   \[  \ip{\Gamma}{\psi'} = x'' + y^{2} + z_{i}^{n}\]
Don't start a paragraph with a displayed equation,
nor make one a paragraph by itself.

\end{document}               % End of document.

\section{Bench Huge Section 1602}
% synthetic bench section 1602
% This is a sample LaTeX input file.  (Version of 12 August 2004.)
%
% A '%' character causes TeX to ignore all remaining text on the line,
% and is used for comments like this one.

\documentclass{article}      % Specifies the document class

                             % The preamble begins here.
\title{An Example Document}  % Declares the document's title.
\author{Leslie Lamport}      % Declares the author's name.
\date{January 21, 1994}      % Deleting this command produces today's date.

\newcommand{\ip}[2]{(#1, #2)}
                             % Defines \ip{arg1}{arg2} to mean
                             % (arg1, arg2).

%\newcommand{\ip}[2]{\langle #1 | #2\rangle}
                             % This is an alternative definition of
                             % \ip that is commented out.

\begin{document}             % End of preamble and beginning of text.

\maketitle                   % Produces the title.

This is an example input file.  Comparing it with
the output it generates can show you how to
produce a simple document of your own.

\section{Ordinary Text}      % Produces section heading.  Lower-level
                             % sections are begun with similar
                             % \subsection and \subsubsection commands.

The ends  of words and sentences are marked
  by   spaces. It  doesn't matter how many
spaces    you type; one is as good as 100.  The
end of   a line counts as a space.

One   or more   blank lines denote the  end
of  a paragraph.

Since any number of consecutive spaces are treated
like a single one, the formatting of the input
file makes no difference to
      \LaTeX,                % The \LaTeX command generates the LaTeX logo.
but it makes a difference to you.  When you use
\LaTeX, making your input file as easy to read
as possible will be a great help as you write
your document and when you change it.  This sample
file shows how you can add comments to your own input
file.

Because printing is different from typewriting,
there are a number of things that you have to do
differently when preparing an input file than if
you were just typing the document directly.
Quotation marks like
       ``this''
have to be handled specially, as do quotes within
quotes:
       ``\,`this'            % \, separates the double and single quote.
        is what I just
        wrote, not  `that'\,''.

Dashes come in three sizes: an
       intra-word
dash, a medium dash for number ranges like
       1--2,
and a punctuation
       dash---like
this.

A sentence-ending space should be larger than the
space between words within a sentence.  You
sometimes have to type special commands in
conjunction with punctuation characters to get
this right, as in the following sentence.
       Gnats, gnus, etc.\ all  % `\ ' makes an inter-word space.
       begin with G\@.         % \@ marks end-of-sentence punctuation.
You should check the spaces after periods when
reading your output to make sure you haven't
forgotten any special cases.  Generating an
ellipsis
       \ldots\               % `\ ' is needed after `\ldots' because TeX
                             % ignores spaces after command names like \ldots
                             % made from \ + letters.
                             %
                             % Note how a `%' character causes TeX to ignore
                             % the end of the input line, so these blank lines
                             % do not start a new paragraph.
                             %
with the right spacing around the periods requires
a special command.

\LaTeX\ interprets some common characters as
commands, so you must type special commands to
generate them.  These characters include the
following:
       \$ \& \% \# \{ and \}.

In printing, text is usually emphasized with an
       \emph{italic}
type style.

\begin{em}
   A long segment of text can also be emphasized
   in this way.  Text within such a segment can be
   given \emph{additional} emphasis.
\end{em}

It is sometimes necessary to prevent \LaTeX\ from
breaking a line where it might otherwise do so.
This may be at a space, as between the ``Mr.''\ and
``Jones'' in
       ``Mr.~Jones'',        % ~ produces an unbreakable interword space.
or within a word---especially when the word is a
symbol like
       \mbox{\emph{itemnum}}
that makes little sense when hyphenated across
lines.

Footnotes\footnote{This is an example of a footnote.}
pose no problem.

\LaTeX\ is good at typesetting mathematical formulas
like
       \( x-3y + z = 7 \)
or
       \( a_{1} > x^{2n} + y^{2n} > x' \)
or
       \( \ip{A}{B} = \sum_{i} a_{i} b_{i} \).
The spaces you type in a formula are
ignored.  Remember that a letter like
       $x$                   % $ ... $  and  \( ... \)  are equivalent
is a formula when it denotes a mathematical
symbol, and it should be typed as one.

\section{Displayed Text}

Text is displayed by indenting it from the left
margin.  Quotations are commonly displayed.  There
are short quotations
\begin{quote}
   This is a short quotation.  It consists of a
   single paragraph of text.  See how it is formatted.
\end{quote}
and longer ones.
\begin{quotation}
   This is a longer quotation.  It consists of two
   paragraphs of text, neither of which are
   particularly interesting.

   This is the second paragraph of the quotation.  It
   is just as dull as the first paragraph.
\end{quotation}
Another frequently-displayed structure is a list.
The following is an example of an \emph{itemized}
list.
\begin{itemize}
   \item This is the first item of an itemized list.
         Each item in the list is marked with a ``tick''.
         You don't have to worry about what kind of tick
         mark is used.

   \item This is the second item of the list.  It
         contains another list nested inside it.  The inner
         list is an \emph{enumerated} list.
         \begin{enumerate}
            \item This is the first item of an enumerated
                  list that is nested within the itemized list.

            \item This is the second item of the inner list.
                  \LaTeX\ allows you to nest lists deeper than
                  you really should.
         \end{enumerate}
         This is the rest of the second item of the outer
         list.  It is no more interesting than any other
         part of the item.
   \item This is the third item of the list.
\end{itemize}
You can even display poetry.
\begin{verse}
   There is an environment
    for verse \\             % The \\ command separates lines
   Whose features some poets % within a stanza.
   will curse.

                             % One or more blank lines separate stanzas.

   For instead of making\\
   Them do \emph{all} line breaking, \\
   It allows them to put too many words on a line when they'd rather be
   forced to be terse.
\end{verse}

Mathematical formulas may also be displayed.  A
displayed formula
is
one-line long; multiline
formulas require special formatting instructions.
   \[  \ip{\Gamma}{\psi'} = x'' + y^{2} + z_{i}^{n}\]
Don't start a paragraph with a displayed equation,
nor make one a paragraph by itself.

\end{document}               % End of document.

\section{Bench Huge Section 1603}
% synthetic bench section 1603
% This is a sample LaTeX input file.  (Version of 12 August 2004.)
%
% A '%' character causes TeX to ignore all remaining text on the line,
% and is used for comments like this one.

\documentclass{article}      % Specifies the document class

                             % The preamble begins here.
\title{An Example Document}  % Declares the document's title.
\author{Leslie Lamport}      % Declares the author's name.
\date{January 21, 1994}      % Deleting this command produces today's date.

\newcommand{\ip}[2]{(#1, #2)}
                             % Defines \ip{arg1}{arg2} to mean
                             % (arg1, arg2).

%\newcommand{\ip}[2]{\langle #1 | #2\rangle}
                             % This is an alternative definition of
                             % \ip that is commented out.

\begin{document}             % End of preamble and beginning of text.

\maketitle                   % Produces the title.

This is an example input file.  Comparing it with
the output it generates can show you how to
produce a simple document of your own.

\section{Ordinary Text}      % Produces section heading.  Lower-level
                             % sections are begun with similar
                             % \subsection and \subsubsection commands.

The ends  of words and sentences are marked
  by   spaces. It  doesn't matter how many
spaces    you type; one is as good as 100.  The
end of   a line counts as a space.

One   or more   blank lines denote the  end
of  a paragraph.

Since any number of consecutive spaces are treated
like a single one, the formatting of the input
file makes no difference to
      \LaTeX,                % The \LaTeX command generates the LaTeX logo.
but it makes a difference to you.  When you use
\LaTeX, making your input file as easy to read
as possible will be a great help as you write
your document and when you change it.  This sample
file shows how you can add comments to your own input
file.

Because printing is different from typewriting,
there are a number of things that you have to do
differently when preparing an input file than if
you were just typing the document directly.
Quotation marks like
       ``this''
have to be handled specially, as do quotes within
quotes:
       ``\,`this'            % \, separates the double and single quote.
        is what I just
        wrote, not  `that'\,''.

Dashes come in three sizes: an
       intra-word
dash, a medium dash for number ranges like
       1--2,
and a punctuation
       dash---like
this.

A sentence-ending space should be larger than the
space between words within a sentence.  You
sometimes have to type special commands in
conjunction with punctuation characters to get
this right, as in the following sentence.
       Gnats, gnus, etc.\ all  % `\ ' makes an inter-word space.
       begin with G\@.         % \@ marks end-of-sentence punctuation.
You should check the spaces after periods when
reading your output to make sure you haven't
forgotten any special cases.  Generating an
ellipsis
       \ldots\               % `\ ' is needed after `\ldots' because TeX
                             % ignores spaces after command names like \ldots
                             % made from \ + letters.
                             %
                             % Note how a `%' character causes TeX to ignore
                             % the end of the input line, so these blank lines
                             % do not start a new paragraph.
                             %
with the right spacing around the periods requires
a special command.

\LaTeX\ interprets some common characters as
commands, so you must type special commands to
generate them.  These characters include the
following:
       \$ \& \% \# \{ and \}.

In printing, text is usually emphasized with an
       \emph{italic}
type style.

\begin{em}
   A long segment of text can also be emphasized
   in this way.  Text within such a segment can be
   given \emph{additional} emphasis.
\end{em}

It is sometimes necessary to prevent \LaTeX\ from
breaking a line where it might otherwise do so.
This may be at a space, as between the ``Mr.''\ and
``Jones'' in
       ``Mr.~Jones'',        % ~ produces an unbreakable interword space.
or within a word---especially when the word is a
symbol like
       \mbox{\emph{itemnum}}
that makes little sense when hyphenated across
lines.

Footnotes\footnote{This is an example of a footnote.}
pose no problem.

\LaTeX\ is good at typesetting mathematical formulas
like
       \( x-3y + z = 7 \)
or
       \( a_{1} > x^{2n} + y^{2n} > x' \)
or
       \( \ip{A}{B} = \sum_{i} a_{i} b_{i} \).
The spaces you type in a formula are
ignored.  Remember that a letter like
       $x$                   % $ ... $  and  \( ... \)  are equivalent
is a formula when it denotes a mathematical
symbol, and it should be typed as one.

\section{Displayed Text}

Text is displayed by indenting it from the left
margin.  Quotations are commonly displayed.  There
are short quotations
\begin{quote}
   This is a short quotation.  It consists of a
   single paragraph of text.  See how it is formatted.
\end{quote}
and longer ones.
\begin{quotation}
   This is a longer quotation.  It consists of two
   paragraphs of text, neither of which are
   particularly interesting.

   This is the second paragraph of the quotation.  It
   is just as dull as the first paragraph.
\end{quotation}
Another frequently-displayed structure is a list.
The following is an example of an \emph{itemized}
list.
\begin{itemize}
   \item This is the first item of an itemized list.
         Each item in the list is marked with a ``tick''.
         You don't have to worry about what kind of tick
         mark is used.

   \item This is the second item of the list.  It
         contains another list nested inside it.  The inner
         list is an \emph{enumerated} list.
         \begin{enumerate}
            \item This is the first item of an enumerated
                  list that is nested within the itemized list.

            \item This is the second item of the inner list.
                  \LaTeX\ allows you to nest lists deeper than
                  you really should.
         \end{enumerate}
         This is the rest of the second item of the outer
         list.  It is no more interesting than any other
         part of the item.
   \item This is the third item of the list.
\end{itemize}
You can even display poetry.
\begin{verse}
   There is an environment
    for verse \\             % The \\ command separates lines
   Whose features some poets % within a stanza.
   will curse.

                             % One or more blank lines separate stanzas.

   For instead of making\\
   Them do \emph{all} line breaking, \\
   It allows them to put too many words on a line when they'd rather be
   forced to be terse.
\end{verse}

Mathematical formulas may also be displayed.  A
displayed formula
is
one-line long; multiline
formulas require special formatting instructions.
   \[  \ip{\Gamma}{\psi'} = x'' + y^{2} + z_{i}^{n}\]
Don't start a paragraph with a displayed equation,
nor make one a paragraph by itself.

\end{document}               % End of document.

\section{Bench Huge Section 1604}
% synthetic bench section 1604
% This is a sample LaTeX input file.  (Version of 12 August 2004.)
%
% A '%' character causes TeX to ignore all remaining text on the line,
% and is used for comments like this one.

\documentclass{article}      % Specifies the document class

                             % The preamble begins here.
\title{An Example Document}  % Declares the document's title.
\author{Leslie Lamport}      % Declares the author's name.
\date{January 21, 1994}      % Deleting this command produces today's date.

\newcommand{\ip}[2]{(#1, #2)}
                             % Defines \ip{arg1}{arg2} to mean
                             % (arg1, arg2).

%\newcommand{\ip}[2]{\langle #1 | #2\rangle}
                             % This is an alternative definition of
                             % \ip that is commented out.

\begin{document}             % End of preamble and beginning of text.

\maketitle                   % Produces the title.

This is an example input file.  Comparing it with
the output it generates can show you how to
produce a simple document of your own.

\section{Ordinary Text}      % Produces section heading.  Lower-level
                             % sections are begun with similar
                             % \subsection and \subsubsection commands.

The ends  of words and sentences are marked
  by   spaces. It  doesn't matter how many
spaces    you type; one is as good as 100.  The
end of   a line counts as a space.

One   or more   blank lines denote the  end
of  a paragraph.

Since any number of consecutive spaces are treated
like a single one, the formatting of the input
file makes no difference to
      \LaTeX,                % The \LaTeX command generates the LaTeX logo.
but it makes a difference to you.  When you use
\LaTeX, making your input file as easy to read
as possible will be a great help as you write
your document and when you change it.  This sample
file shows how you can add comments to your own input
file.

Because printing is different from typewriting,
there are a number of things that you have to do
differently when preparing an input file than if
you were just typing the document directly.
Quotation marks like
       ``this''
have to be handled specially, as do quotes within
quotes:
       ``\,`this'            % \, separates the double and single quote.
        is what I just
        wrote, not  `that'\,''.

Dashes come in three sizes: an
       intra-word
dash, a medium dash for number ranges like
       1--2,
and a punctuation
       dash---like
this.

A sentence-ending space should be larger than the
space between words within a sentence.  You
sometimes have to type special commands in
conjunction with punctuation characters to get
this right, as in the following sentence.
       Gnats, gnus, etc.\ all  % `\ ' makes an inter-word space.
       begin with G\@.         % \@ marks end-of-sentence punctuation.
You should check the spaces after periods when
reading your output to make sure you haven't
forgotten any special cases.  Generating an
ellipsis
       \ldots\               % `\ ' is needed after `\ldots' because TeX
                             % ignores spaces after command names like \ldots
                             % made from \ + letters.
                             %
                             % Note how a `%' character causes TeX to ignore
                             % the end of the input line, so these blank lines
                             % do not start a new paragraph.
                             %
with the right spacing around the periods requires
a special command.

\LaTeX\ interprets some common characters as
commands, so you must type special commands to
generate them.  These characters include the
following:
       \$ \& \% \# \{ and \}.

In printing, text is usually emphasized with an
       \emph{italic}
type style.

\begin{em}
   A long segment of text can also be emphasized
   in this way.  Text within such a segment can be
   given \emph{additional} emphasis.
\end{em}

It is sometimes necessary to prevent \LaTeX\ from
breaking a line where it might otherwise do so.
This may be at a space, as between the ``Mr.''\ and
``Jones'' in
       ``Mr.~Jones'',        % ~ produces an unbreakable interword space.
or within a word---especially when the word is a
symbol like
       \mbox{\emph{itemnum}}
that makes little sense when hyphenated across
lines.

Footnotes\footnote{This is an example of a footnote.}
pose no problem.

\LaTeX\ is good at typesetting mathematical formulas
like
       \( x-3y + z = 7 \)
or
       \( a_{1} > x^{2n} + y^{2n} > x' \)
or
       \( \ip{A}{B} = \sum_{i} a_{i} b_{i} \).
The spaces you type in a formula are
ignored.  Remember that a letter like
       $x$                   % $ ... $  and  \( ... \)  are equivalent
is a formula when it denotes a mathematical
symbol, and it should be typed as one.

\section{Displayed Text}

Text is displayed by indenting it from the left
margin.  Quotations are commonly displayed.  There
are short quotations
\begin{quote}
   This is a short quotation.  It consists of a
   single paragraph of text.  See how it is formatted.
\end{quote}
and longer ones.
\begin{quotation}
   This is a longer quotation.  It consists of two
   paragraphs of text, neither of which are
   particularly interesting.

   This is the second paragraph of the quotation.  It
   is just as dull as the first paragraph.
\end{quotation}
Another frequently-displayed structure is a list.
The following is an example of an \emph{itemized}
list.
\begin{itemize}
   \item This is the first item of an itemized list.
         Each item in the list is marked with a ``tick''.
         You don't have to worry about what kind of tick
         mark is used.

   \item This is the second item of the list.  It
         contains another list nested inside it.  The inner
         list is an \emph{enumerated} list.
         \begin{enumerate}
            \item This is the first item of an enumerated
                  list that is nested within the itemized list.

            \item This is the second item of the inner list.
                  \LaTeX\ allows you to nest lists deeper than
                  you really should.
         \end{enumerate}
         This is the rest of the second item of the outer
         list.  It is no more interesting than any other
         part of the item.
   \item This is the third item of the list.
\end{itemize}
You can even display poetry.
\begin{verse}
   There is an environment
    for verse \\             % The \\ command separates lines
   Whose features some poets % within a stanza.
   will curse.

                             % One or more blank lines separate stanzas.

   For instead of making\\
   Them do \emph{all} line breaking, \\
   It allows them to put too many words on a line when they'd rather be
   forced to be terse.
\end{verse}

Mathematical formulas may also be displayed.  A
displayed formula
is
one-line long; multiline
formulas require special formatting instructions.
   \[  \ip{\Gamma}{\psi'} = x'' + y^{2} + z_{i}^{n}\]
Don't start a paragraph with a displayed equation,
nor make one a paragraph by itself.

\end{document}               % End of document.

\section{Bench Huge Section 1605}
% synthetic bench section 1605
% This is a sample LaTeX input file.  (Version of 12 August 2004.)
%
% A '%' character causes TeX to ignore all remaining text on the line,
% and is used for comments like this one.

\documentclass{article}      % Specifies the document class

                             % The preamble begins here.
\title{An Example Document}  % Declares the document's title.
\author{Leslie Lamport}      % Declares the author's name.
\date{January 21, 1994}      % Deleting this command produces today's date.

\newcommand{\ip}[2]{(#1, #2)}
                             % Defines \ip{arg1}{arg2} to mean
                             % (arg1, arg2).

%\newcommand{\ip}[2]{\langle #1 | #2\rangle}
                             % This is an alternative definition of
                             % \ip that is commented out.

\begin{document}             % End of preamble and beginning of text.

\maketitle                   % Produces the title.

This is an example input file.  Comparing it with
the output it generates can show you how to
produce a simple document of your own.

\section{Ordinary Text}      % Produces section heading.  Lower-level
                             % sections are begun with similar
                             % \subsection and \subsubsection commands.

The ends  of words and sentences are marked
  by   spaces. It  doesn't matter how many
spaces    you type; one is as good as 100.  The
end of   a line counts as a space.

One   or more   blank lines denote the  end
of  a paragraph.

Since any number of consecutive spaces are treated
like a single one, the formatting of the input
file makes no difference to
      \LaTeX,                % The \LaTeX command generates the LaTeX logo.
but it makes a difference to you.  When you use
\LaTeX, making your input file as easy to read
as possible will be a great help as you write
your document and when you change it.  This sample
file shows how you can add comments to your own input
file.

Because printing is different from typewriting,
there are a number of things that you have to do
differently when preparing an input file than if
you were just typing the document directly.
Quotation marks like
       ``this''
have to be handled specially, as do quotes within
quotes:
       ``\,`this'            % \, separates the double and single quote.
        is what I just
        wrote, not  `that'\,''.

Dashes come in three sizes: an
       intra-word
dash, a medium dash for number ranges like
       1--2,
and a punctuation
       dash---like
this.

A sentence-ending space should be larger than the
space between words within a sentence.  You
sometimes have to type special commands in
conjunction with punctuation characters to get
this right, as in the following sentence.
       Gnats, gnus, etc.\ all  % `\ ' makes an inter-word space.
       begin with G\@.         % \@ marks end-of-sentence punctuation.
You should check the spaces after periods when
reading your output to make sure you haven't
forgotten any special cases.  Generating an
ellipsis
       \ldots\               % `\ ' is needed after `\ldots' because TeX
                             % ignores spaces after command names like \ldots
                             % made from \ + letters.
                             %
                             % Note how a `%' character causes TeX to ignore
                             % the end of the input line, so these blank lines
                             % do not start a new paragraph.
                             %
with the right spacing around the periods requires
a special command.

\LaTeX\ interprets some common characters as
commands, so you must type special commands to
generate them.  These characters include the
following:
       \$ \& \% \# \{ and \}.

In printing, text is usually emphasized with an
       \emph{italic}
type style.

\begin{em}
   A long segment of text can also be emphasized
   in this way.  Text within such a segment can be
   given \emph{additional} emphasis.
\end{em}

It is sometimes necessary to prevent \LaTeX\ from
breaking a line where it might otherwise do so.
This may be at a space, as between the ``Mr.''\ and
``Jones'' in
       ``Mr.~Jones'',        % ~ produces an unbreakable interword space.
or within a word---especially when the word is a
symbol like
       \mbox{\emph{itemnum}}
that makes little sense when hyphenated across
lines.

Footnotes\footnote{This is an example of a footnote.}
pose no problem.

\LaTeX\ is good at typesetting mathematical formulas
like
       \( x-3y + z = 7 \)
or
       \( a_{1} > x^{2n} + y^{2n} > x' \)
or
       \( \ip{A}{B} = \sum_{i} a_{i} b_{i} \).
The spaces you type in a formula are
ignored.  Remember that a letter like
       $x$                   % $ ... $  and  \( ... \)  are equivalent
is a formula when it denotes a mathematical
symbol, and it should be typed as one.

\section{Displayed Text}

Text is displayed by indenting it from the left
margin.  Quotations are commonly displayed.  There
are short quotations
\begin{quote}
   This is a short quotation.  It consists of a
   single paragraph of text.  See how it is formatted.
\end{quote}
and longer ones.
\begin{quotation}
   This is a longer quotation.  It consists of two
   paragraphs of text, neither of which are
   particularly interesting.

   This is the second paragraph of the quotation.  It
   is just as dull as the first paragraph.
\end{quotation}
Another frequently-displayed structure is a list.
The following is an example of an \emph{itemized}
list.
\begin{itemize}
   \item This is the first item of an itemized list.
         Each item in the list is marked with a ``tick''.
         You don't have to worry about what kind of tick
         mark is used.

   \item This is the second item of the list.  It
         contains another list nested inside it.  The inner
         list is an \emph{enumerated} list.
         \begin{enumerate}
            \item This is the first item of an enumerated
                  list that is nested within the itemized list.

            \item This is the second item of the inner list.
                  \LaTeX\ allows you to nest lists deeper than
                  you really should.
         \end{enumerate}
         This is the rest of the second item of the outer
         list.  It is no more interesting than any other
         part of the item.
   \item This is the third item of the list.
\end{itemize}
You can even display poetry.
\begin{verse}
   There is an environment
    for verse \\             % The \\ command separates lines
   Whose features some poets % within a stanza.
   will curse.

                             % One or more blank lines separate stanzas.

   For instead of making\\
   Them do \emph{all} line breaking, \\
   It allows them to put too many words on a line when they'd rather be
   forced to be terse.
\end{verse}

Mathematical formulas may also be displayed.  A
displayed formula
is
one-line long; multiline
formulas require special formatting instructions.
   \[  \ip{\Gamma}{\psi'} = x'' + y^{2} + z_{i}^{n}\]
Don't start a paragraph with a displayed equation,
nor make one a paragraph by itself.

\end{document}               % End of document.

\section{Bench Huge Section 1606}
% synthetic bench section 1606
% This is a sample LaTeX input file.  (Version of 12 August 2004.)
%
% A '%' character causes TeX to ignore all remaining text on the line,
% and is used for comments like this one.

\documentclass{article}      % Specifies the document class

                             % The preamble begins here.
\title{An Example Document}  % Declares the document's title.
\author{Leslie Lamport}      % Declares the author's name.
\date{January 21, 1994}      % Deleting this command produces today's date.

\newcommand{\ip}[2]{(#1, #2)}
                             % Defines \ip{arg1}{arg2} to mean
                             % (arg1, arg2).

%\newcommand{\ip}[2]{\langle #1 | #2\rangle}
                             % This is an alternative definition of
                             % \ip that is commented out.

\begin{document}             % End of preamble and beginning of text.

\maketitle                   % Produces the title.

This is an example input file.  Comparing it with
the output it generates can show you how to
produce a simple document of your own.

\section{Ordinary Text}      % Produces section heading.  Lower-level
                             % sections are begun with similar
                             % \subsection and \subsubsection commands.

The ends  of words and sentences are marked
  by   spaces. It  doesn't matter how many
spaces    you type; one is as good as 100.  The
end of   a line counts as a space.

One   or more   blank lines denote the  end
of  a paragraph.

Since any number of consecutive spaces are treated
like a single one, the formatting of the input
file makes no difference to
      \LaTeX,                % The \LaTeX command generates the LaTeX logo.
but it makes a difference to you.  When you use
\LaTeX, making your input file as easy to read
as possible will be a great help as you write
your document and when you change it.  This sample
file shows how you can add comments to your own input
file.

Because printing is different from typewriting,
there are a number of things that you have to do
differently when preparing an input file than if
you were just typing the document directly.
Quotation marks like
       ``this''
have to be handled specially, as do quotes within
quotes:
       ``\,`this'            % \, separates the double and single quote.
        is what I just
        wrote, not  `that'\,''.

Dashes come in three sizes: an
       intra-word
dash, a medium dash for number ranges like
       1--2,
and a punctuation
       dash---like
this.

A sentence-ending space should be larger than the
space between words within a sentence.  You
sometimes have to type special commands in
conjunction with punctuation characters to get
this right, as in the following sentence.
       Gnats, gnus, etc.\ all  % `\ ' makes an inter-word space.
       begin with G\@.         % \@ marks end-of-sentence punctuation.
You should check the spaces after periods when
reading your output to make sure you haven't
forgotten any special cases.  Generating an
ellipsis
       \ldots\               % `\ ' is needed after `\ldots' because TeX
                             % ignores spaces after command names like \ldots
                             % made from \ + letters.
                             %
                             % Note how a `%' character causes TeX to ignore
                             % the end of the input line, so these blank lines
                             % do not start a new paragraph.
                             %
with the right spacing around the periods requires
a special command.

\LaTeX\ interprets some common characters as
commands, so you must type special commands to
generate them.  These characters include the
following:
       \$ \& \% \# \{ and \}.

In printing, text is usually emphasized with an
       \emph{italic}
type style.

\begin{em}
   A long segment of text can also be emphasized
   in this way.  Text within such a segment can be
   given \emph{additional} emphasis.
\end{em}

It is sometimes necessary to prevent \LaTeX\ from
breaking a line where it might otherwise do so.
This may be at a space, as between the ``Mr.''\ and
``Jones'' in
       ``Mr.~Jones'',        % ~ produces an unbreakable interword space.
or within a word---especially when the word is a
symbol like
       \mbox{\emph{itemnum}}
that makes little sense when hyphenated across
lines.

Footnotes\footnote{This is an example of a footnote.}
pose no problem.

\LaTeX\ is good at typesetting mathematical formulas
like
       \( x-3y + z = 7 \)
or
       \( a_{1} > x^{2n} + y^{2n} > x' \)
or
       \( \ip{A}{B} = \sum_{i} a_{i} b_{i} \).
The spaces you type in a formula are
ignored.  Remember that a letter like
       $x$                   % $ ... $  and  \( ... \)  are equivalent
is a formula when it denotes a mathematical
symbol, and it should be typed as one.

\section{Displayed Text}

Text is displayed by indenting it from the left
margin.  Quotations are commonly displayed.  There
are short quotations
\begin{quote}
   This is a short quotation.  It consists of a
   single paragraph of text.  See how it is formatted.
\end{quote}
and longer ones.
\begin{quotation}
   This is a longer quotation.  It consists of two
   paragraphs of text, neither of which are
   particularly interesting.

   This is the second paragraph of the quotation.  It
   is just as dull as the first paragraph.
\end{quotation}
Another frequently-displayed structure is a list.
The following is an example of an \emph{itemized}
list.
\begin{itemize}
   \item This is the first item of an itemized list.
         Each item in the list is marked with a ``tick''.
         You don't have to worry about what kind of tick
         mark is used.

   \item This is the second item of the list.  It
         contains another list nested inside it.  The inner
         list is an \emph{enumerated} list.
         \begin{enumerate}
            \item This is the first item of an enumerated
                  list that is nested within the itemized list.

            \item This is the second item of the inner list.
                  \LaTeX\ allows you to nest lists deeper than
                  you really should.
         \end{enumerate}
         This is the rest of the second item of the outer
         list.  It is no more interesting than any other
         part of the item.
   \item This is the third item of the list.
\end{itemize}
You can even display poetry.
\begin{verse}
   There is an environment
    for verse \\             % The \\ command separates lines
   Whose features some poets % within a stanza.
   will curse.

                             % One or more blank lines separate stanzas.

   For instead of making\\
   Them do \emph{all} line breaking, \\
   It allows them to put too many words on a line when they'd rather be
   forced to be terse.
\end{verse}

Mathematical formulas may also be displayed.  A
displayed formula
is
one-line long; multiline
formulas require special formatting instructions.
   \[  \ip{\Gamma}{\psi'} = x'' + y^{2} + z_{i}^{n}\]
Don't start a paragraph with a displayed equation,
nor make one a paragraph by itself.

\end{document}               % End of document.

\section{Bench Huge Section 1607}
% synthetic bench section 1607
% This is a sample LaTeX input file.  (Version of 12 August 2004.)
%
% A '%' character causes TeX to ignore all remaining text on the line,
% and is used for comments like this one.

\documentclass{article}      % Specifies the document class

                             % The preamble begins here.
\title{An Example Document}  % Declares the document's title.
\author{Leslie Lamport}      % Declares the author's name.
\date{January 21, 1994}      % Deleting this command produces today's date.

\newcommand{\ip}[2]{(#1, #2)}
                             % Defines \ip{arg1}{arg2} to mean
                             % (arg1, arg2).

%\newcommand{\ip}[2]{\langle #1 | #2\rangle}
                             % This is an alternative definition of
                             % \ip that is commented out.

\begin{document}             % End of preamble and beginning of text.

\maketitle                   % Produces the title.

This is an example input file.  Comparing it with
the output it generates can show you how to
produce a simple document of your own.

\section{Ordinary Text}      % Produces section heading.  Lower-level
                             % sections are begun with similar
                             % \subsection and \subsubsection commands.

The ends  of words and sentences are marked
  by   spaces. It  doesn't matter how many
spaces    you type; one is as good as 100.  The
end of   a line counts as a space.

One   or more   blank lines denote the  end
of  a paragraph.

Since any number of consecutive spaces are treated
like a single one, the formatting of the input
file makes no difference to
      \LaTeX,                % The \LaTeX command generates the LaTeX logo.
but it makes a difference to you.  When you use
\LaTeX, making your input file as easy to read
as possible will be a great help as you write
your document and when you change it.  This sample
file shows how you can add comments to your own input
file.

Because printing is different from typewriting,
there are a number of things that you have to do
differently when preparing an input file than if
you were just typing the document directly.
Quotation marks like
       ``this''
have to be handled specially, as do quotes within
quotes:
       ``\,`this'            % \, separates the double and single quote.
        is what I just
        wrote, not  `that'\,''.

Dashes come in three sizes: an
       intra-word
dash, a medium dash for number ranges like
       1--2,
and a punctuation
       dash---like
this.

A sentence-ending space should be larger than the
space between words within a sentence.  You
sometimes have to type special commands in
conjunction with punctuation characters to get
this right, as in the following sentence.
       Gnats, gnus, etc.\ all  % `\ ' makes an inter-word space.
       begin with G\@.         % \@ marks end-of-sentence punctuation.
You should check the spaces after periods when
reading your output to make sure you haven't
forgotten any special cases.  Generating an
ellipsis
       \ldots\               % `\ ' is needed after `\ldots' because TeX
                             % ignores spaces after command names like \ldots
                             % made from \ + letters.
                             %
                             % Note how a `%' character causes TeX to ignore
                             % the end of the input line, so these blank lines
                             % do not start a new paragraph.
                             %
with the right spacing around the periods requires
a special command.

\LaTeX\ interprets some common characters as
commands, so you must type special commands to
generate them.  These characters include the
following:
       \$ \& \% \# \{ and \}.

In printing, text is usually emphasized with an
       \emph{italic}
type style.

\begin{em}
   A long segment of text can also be emphasized
   in this way.  Text within such a segment can be
   given \emph{additional} emphasis.
\end{em}

It is sometimes necessary to prevent \LaTeX\ from
breaking a line where it might otherwise do so.
This may be at a space, as between the ``Mr.''\ and
``Jones'' in
       ``Mr.~Jones'',        % ~ produces an unbreakable interword space.
or within a word---especially when the word is a
symbol like
       \mbox{\emph{itemnum}}
that makes little sense when hyphenated across
lines.

Footnotes\footnote{This is an example of a footnote.}
pose no problem.

\LaTeX\ is good at typesetting mathematical formulas
like
       \( x-3y + z = 7 \)
or
       \( a_{1} > x^{2n} + y^{2n} > x' \)
or
       \( \ip{A}{B} = \sum_{i} a_{i} b_{i} \).
The spaces you type in a formula are
ignored.  Remember that a letter like
       $x$                   % $ ... $  and  \( ... \)  are equivalent
is a formula when it denotes a mathematical
symbol, and it should be typed as one.

\section{Displayed Text}

Text is displayed by indenting it from the left
margin.  Quotations are commonly displayed.  There
are short quotations
\begin{quote}
   This is a short quotation.  It consists of a
   single paragraph of text.  See how it is formatted.
\end{quote}
and longer ones.
\begin{quotation}
   This is a longer quotation.  It consists of two
   paragraphs of text, neither of which are
   particularly interesting.

   This is the second paragraph of the quotation.  It
   is just as dull as the first paragraph.
\end{quotation}
Another frequently-displayed structure is a list.
The following is an example of an \emph{itemized}
list.
\begin{itemize}
   \item This is the first item of an itemized list.
         Each item in the list is marked with a ``tick''.
         You don't have to worry about what kind of tick
         mark is used.

   \item This is the second item of the list.  It
         contains another list nested inside it.  The inner
         list is an \emph{enumerated} list.
         \begin{enumerate}
            \item This is the first item of an enumerated
                  list that is nested within the itemized list.

            \item This is the second item of the inner list.
                  \LaTeX\ allows you to nest lists deeper than
                  you really should.
         \end{enumerate}
         This is the rest of the second item of the outer
         list.  It is no more interesting than any other
         part of the item.
   \item This is the third item of the list.
\end{itemize}
You can even display poetry.
\begin{verse}
   There is an environment
    for verse \\             % The \\ command separates lines
   Whose features some poets % within a stanza.
   will curse.

                             % One or more blank lines separate stanzas.

   For instead of making\\
   Them do \emph{all} line breaking, \\
   It allows them to put too many words on a line when they'd rather be
   forced to be terse.
\end{verse}

Mathematical formulas may also be displayed.  A
displayed formula
is
one-line long; multiline
formulas require special formatting instructions.
   \[  \ip{\Gamma}{\psi'} = x'' + y^{2} + z_{i}^{n}\]
Don't start a paragraph with a displayed equation,
nor make one a paragraph by itself.

\end{document}               % End of document.

\section{Bench Huge Section 1608}
% synthetic bench section 1608
% This is a sample LaTeX input file.  (Version of 12 August 2004.)
%
% A '%' character causes TeX to ignore all remaining text on the line,
% and is used for comments like this one.

\documentclass{article}      % Specifies the document class

                             % The preamble begins here.
\title{An Example Document}  % Declares the document's title.
\author{Leslie Lamport}      % Declares the author's name.
\date{January 21, 1994}      % Deleting this command produces today's date.

\newcommand{\ip}[2]{(#1, #2)}
                             % Defines \ip{arg1}{arg2} to mean
                             % (arg1, arg2).

%\newcommand{\ip}[2]{\langle #1 | #2\rangle}
                             % This is an alternative definition of
                             % \ip that is commented out.

\begin{document}             % End of preamble and beginning of text.

\maketitle                   % Produces the title.

This is an example input file.  Comparing it with
the output it generates can show you how to
produce a simple document of your own.

\section{Ordinary Text}      % Produces section heading.  Lower-level
                             % sections are begun with similar
                             % \subsection and \subsubsection commands.

The ends  of words and sentences are marked
  by   spaces. It  doesn't matter how many
spaces    you type; one is as good as 100.  The
end of   a line counts as a space.

One   or more   blank lines denote the  end
of  a paragraph.

Since any number of consecutive spaces are treated
like a single one, the formatting of the input
file makes no difference to
      \LaTeX,                % The \LaTeX command generates the LaTeX logo.
but it makes a difference to you.  When you use
\LaTeX, making your input file as easy to read
as possible will be a great help as you write
your document and when you change it.  This sample
file shows how you can add comments to your own input
file.

Because printing is different from typewriting,
there are a number of things that you have to do
differently when preparing an input file than if
you were just typing the document directly.
Quotation marks like
       ``this''
have to be handled specially, as do quotes within
quotes:
       ``\,`this'            % \, separates the double and single quote.
        is what I just
        wrote, not  `that'\,''.

Dashes come in three sizes: an
       intra-word
dash, a medium dash for number ranges like
       1--2,
and a punctuation
       dash---like
this.

A sentence-ending space should be larger than the
space between words within a sentence.  You
sometimes have to type special commands in
conjunction with punctuation characters to get
this right, as in the following sentence.
       Gnats, gnus, etc.\ all  % `\ ' makes an inter-word space.
       begin with G\@.         % \@ marks end-of-sentence punctuation.
You should check the spaces after periods when
reading your output to make sure you haven't
forgotten any special cases.  Generating an
ellipsis
       \ldots\               % `\ ' is needed after `\ldots' because TeX
                             % ignores spaces after command names like \ldots
                             % made from \ + letters.
                             %
                             % Note how a `%' character causes TeX to ignore
                             % the end of the input line, so these blank lines
                             % do not start a new paragraph.
                             %
with the right spacing around the periods requires
a special command.

\LaTeX\ interprets some common characters as
commands, so you must type special commands to
generate them.  These characters include the
following:
       \$ \& \% \# \{ and \}.

In printing, text is usually emphasized with an
       \emph{italic}
type style.

\begin{em}
   A long segment of text can also be emphasized
   in this way.  Text within such a segment can be
   given \emph{additional} emphasis.
\end{em}

It is sometimes necessary to prevent \LaTeX\ from
breaking a line where it might otherwise do so.
This may be at a space, as between the ``Mr.''\ and
``Jones'' in
       ``Mr.~Jones'',        % ~ produces an unbreakable interword space.
or within a word---especially when the word is a
symbol like
       \mbox{\emph{itemnum}}
that makes little sense when hyphenated across
lines.

Footnotes\footnote{This is an example of a footnote.}
pose no problem.

\LaTeX\ is good at typesetting mathematical formulas
like
       \( x-3y + z = 7 \)
or
       \( a_{1} > x^{2n} + y^{2n} > x' \)
or
       \( \ip{A}{B} = \sum_{i} a_{i} b_{i} \).
The spaces you type in a formula are
ignored.  Remember that a letter like
       $x$                   % $ ... $  and  \( ... \)  are equivalent
is a formula when it denotes a mathematical
symbol, and it should be typed as one.

\section{Displayed Text}

Text is displayed by indenting it from the left
margin.  Quotations are commonly displayed.  There
are short quotations
\begin{quote}
   This is a short quotation.  It consists of a
   single paragraph of text.  See how it is formatted.
\end{quote}
and longer ones.
\begin{quotation}
   This is a longer quotation.  It consists of two
   paragraphs of text, neither of which are
   particularly interesting.

   This is the second paragraph of the quotation.  It
   is just as dull as the first paragraph.
\end{quotation}
Another frequently-displayed structure is a list.
The following is an example of an \emph{itemized}
list.
\begin{itemize}
   \item This is the first item of an itemized list.
         Each item in the list is marked with a ``tick''.
         You don't have to worry about what kind of tick
         mark is used.

   \item This is the second item of the list.  It
         contains another list nested inside it.  The inner
         list is an \emph{enumerated} list.
         \begin{enumerate}
            \item This is the first item of an enumerated
                  list that is nested within the itemized list.

            \item This is the second item of the inner list.
                  \LaTeX\ allows you to nest lists deeper than
                  you really should.
         \end{enumerate}
         This is the rest of the second item of the outer
         list.  It is no more interesting than any other
         part of the item.
   \item This is the third item of the list.
\end{itemize}
You can even display poetry.
\begin{verse}
   There is an environment
    for verse \\             % The \\ command separates lines
   Whose features some poets % within a stanza.
   will curse.

                             % One or more blank lines separate stanzas.

   For instead of making\\
   Them do \emph{all} line breaking, \\
   It allows them to put too many words on a line when they'd rather be
   forced to be terse.
\end{verse}

Mathematical formulas may also be displayed.  A
displayed formula
is
one-line long; multiline
formulas require special formatting instructions.
   \[  \ip{\Gamma}{\psi'} = x'' + y^{2} + z_{i}^{n}\]
Don't start a paragraph with a displayed equation,
nor make one a paragraph by itself.

\end{document}               % End of document.

\section{Bench Huge Section 1609}
% synthetic bench section 1609
% This is a sample LaTeX input file.  (Version of 12 August 2004.)
%
% A '%' character causes TeX to ignore all remaining text on the line,
% and is used for comments like this one.

\documentclass{article}      % Specifies the document class

                             % The preamble begins here.
\title{An Example Document}  % Declares the document's title.
\author{Leslie Lamport}      % Declares the author's name.
\date{January 21, 1994}      % Deleting this command produces today's date.

\newcommand{\ip}[2]{(#1, #2)}
                             % Defines \ip{arg1}{arg2} to mean
                             % (arg1, arg2).

%\newcommand{\ip}[2]{\langle #1 | #2\rangle}
                             % This is an alternative definition of
                             % \ip that is commented out.

\begin{document}             % End of preamble and beginning of text.

\maketitle                   % Produces the title.

This is an example input file.  Comparing it with
the output it generates can show you how to
produce a simple document of your own.

\section{Ordinary Text}      % Produces section heading.  Lower-level
                             % sections are begun with similar
                             % \subsection and \subsubsection commands.

The ends  of words and sentences are marked
  by   spaces. It  doesn't matter how many
spaces    you type; one is as good as 100.  The
end of   a line counts as a space.

One   or more   blank lines denote the  end
of  a paragraph.

Since any number of consecutive spaces are treated
like a single one, the formatting of the input
file makes no difference to
      \LaTeX,                % The \LaTeX command generates the LaTeX logo.
but it makes a difference to you.  When you use
\LaTeX, making your input file as easy to read
as possible will be a great help as you write
your document and when you change it.  This sample
file shows how you can add comments to your own input
file.

Because printing is different from typewriting,
there are a number of things that you have to do
differently when preparing an input file than if
you were just typing the document directly.
Quotation marks like
       ``this''
have to be handled specially, as do quotes within
quotes:
       ``\,`this'            % \, separates the double and single quote.
        is what I just
        wrote, not  `that'\,''.

Dashes come in three sizes: an
       intra-word
dash, a medium dash for number ranges like
       1--2,
and a punctuation
       dash---like
this.

A sentence-ending space should be larger than the
space between words within a sentence.  You
sometimes have to type special commands in
conjunction with punctuation characters to get
this right, as in the following sentence.
       Gnats, gnus, etc.\ all  % `\ ' makes an inter-word space.
       begin with G\@.         % \@ marks end-of-sentence punctuation.
You should check the spaces after periods when
reading your output to make sure you haven't
forgotten any special cases.  Generating an
ellipsis
       \ldots\               % `\ ' is needed after `\ldots' because TeX
                             % ignores spaces after command names like \ldots
                             % made from \ + letters.
                             %
                             % Note how a `%' character causes TeX to ignore
                             % the end of the input line, so these blank lines
                             % do not start a new paragraph.
                             %
with the right spacing around the periods requires
a special command.

\LaTeX\ interprets some common characters as
commands, so you must type special commands to
generate them.  These characters include the
following:
       \$ \& \% \# \{ and \}.

In printing, text is usually emphasized with an
       \emph{italic}
type style.

\begin{em}
   A long segment of text can also be emphasized
   in this way.  Text within such a segment can be
   given \emph{additional} emphasis.
\end{em}

It is sometimes necessary to prevent \LaTeX\ from
breaking a line where it might otherwise do so.
This may be at a space, as between the ``Mr.''\ and
``Jones'' in
       ``Mr.~Jones'',        % ~ produces an unbreakable interword space.
or within a word---especially when the word is a
symbol like
       \mbox{\emph{itemnum}}
that makes little sense when hyphenated across
lines.

Footnotes\footnote{This is an example of a footnote.}
pose no problem.

\LaTeX\ is good at typesetting mathematical formulas
like
       \( x-3y + z = 7 \)
or
       \( a_{1} > x^{2n} + y^{2n} > x' \)
or
       \( \ip{A}{B} = \sum_{i} a_{i} b_{i} \).
The spaces you type in a formula are
ignored.  Remember that a letter like
       $x$                   % $ ... $  and  \( ... \)  are equivalent
is a formula when it denotes a mathematical
symbol, and it should be typed as one.

\section{Displayed Text}

Text is displayed by indenting it from the left
margin.  Quotations are commonly displayed.  There
are short quotations
\begin{quote}
   This is a short quotation.  It consists of a
   single paragraph of text.  See how it is formatted.
\end{quote}
and longer ones.
\begin{quotation}
   This is a longer quotation.  It consists of two
   paragraphs of text, neither of which are
   particularly interesting.

   This is the second paragraph of the quotation.  It
   is just as dull as the first paragraph.
\end{quotation}
Another frequently-displayed structure is a list.
The following is an example of an \emph{itemized}
list.
\begin{itemize}
   \item This is the first item of an itemized list.
         Each item in the list is marked with a ``tick''.
         You don't have to worry about what kind of tick
         mark is used.

   \item This is the second item of the list.  It
         contains another list nested inside it.  The inner
         list is an \emph{enumerated} list.
         \begin{enumerate}
            \item This is the first item of an enumerated
                  list that is nested within the itemized list.

            \item This is the second item of the inner list.
                  \LaTeX\ allows you to nest lists deeper than
                  you really should.
         \end{enumerate}
         This is the rest of the second item of the outer
         list.  It is no more interesting than any other
         part of the item.
   \item This is the third item of the list.
\end{itemize}
You can even display poetry.
\begin{verse}
   There is an environment
    for verse \\             % The \\ command separates lines
   Whose features some poets % within a stanza.
   will curse.

                             % One or more blank lines separate stanzas.

   For instead of making\\
   Them do \emph{all} line breaking, \\
   It allows them to put too many words on a line when they'd rather be
   forced to be terse.
\end{verse}

Mathematical formulas may also be displayed.  A
displayed formula
is
one-line long; multiline
formulas require special formatting instructions.
   \[  \ip{\Gamma}{\psi'} = x'' + y^{2} + z_{i}^{n}\]
Don't start a paragraph with a displayed equation,
nor make one a paragraph by itself.

\end{document}               % End of document.

\section{Bench Huge Section 1610}
% synthetic bench section 1610
% This is a sample LaTeX input file.  (Version of 12 August 2004.)
%
% A '%' character causes TeX to ignore all remaining text on the line,
% and is used for comments like this one.

\documentclass{article}      % Specifies the document class

                             % The preamble begins here.
\title{An Example Document}  % Declares the document's title.
\author{Leslie Lamport}      % Declares the author's name.
\date{January 21, 1994}      % Deleting this command produces today's date.

\newcommand{\ip}[2]{(#1, #2)}
                             % Defines \ip{arg1}{arg2} to mean
                             % (arg1, arg2).

%\newcommand{\ip}[2]{\langle #1 | #2\rangle}
                             % This is an alternative definition of
                             % \ip that is commented out.

\begin{document}             % End of preamble and beginning of text.

\maketitle                   % Produces the title.

This is an example input file.  Comparing it with
the output it generates can show you how to
produce a simple document of your own.

\section{Ordinary Text}      % Produces section heading.  Lower-level
                             % sections are begun with similar
                             % \subsection and \subsubsection commands.

The ends  of words and sentences are marked
  by   spaces. It  doesn't matter how many
spaces    you type; one is as good as 100.  The
end of   a line counts as a space.

One   or more   blank lines denote the  end
of  a paragraph.

Since any number of consecutive spaces are treated
like a single one, the formatting of the input
file makes no difference to
      \LaTeX,                % The \LaTeX command generates the LaTeX logo.
but it makes a difference to you.  When you use
\LaTeX, making your input file as easy to read
as possible will be a great help as you write
your document and when you change it.  This sample
file shows how you can add comments to your own input
file.

Because printing is different from typewriting,
there are a number of things that you have to do
differently when preparing an input file than if
you were just typing the document directly.
Quotation marks like
       ``this''
have to be handled specially, as do quotes within
quotes:
       ``\,`this'            % \, separates the double and single quote.
        is what I just
        wrote, not  `that'\,''.

Dashes come in three sizes: an
       intra-word
dash, a medium dash for number ranges like
       1--2,
and a punctuation
       dash---like
this.

A sentence-ending space should be larger than the
space between words within a sentence.  You
sometimes have to type special commands in
conjunction with punctuation characters to get
this right, as in the following sentence.
       Gnats, gnus, etc.\ all  % `\ ' makes an inter-word space.
       begin with G\@.         % \@ marks end-of-sentence punctuation.
You should check the spaces after periods when
reading your output to make sure you haven't
forgotten any special cases.  Generating an
ellipsis
       \ldots\               % `\ ' is needed after `\ldots' because TeX
                             % ignores spaces after command names like \ldots
                             % made from \ + letters.
                             %
                             % Note how a `%' character causes TeX to ignore
                             % the end of the input line, so these blank lines
                             % do not start a new paragraph.
                             %
with the right spacing around the periods requires
a special command.

\LaTeX\ interprets some common characters as
commands, so you must type special commands to
generate them.  These characters include the
following:
       \$ \& \% \# \{ and \}.

In printing, text is usually emphasized with an
       \emph{italic}
type style.

\begin{em}
   A long segment of text can also be emphasized
   in this way.  Text within such a segment can be
   given \emph{additional} emphasis.
\end{em}

It is sometimes necessary to prevent \LaTeX\ from
breaking a line where it might otherwise do so.
This may be at a space, as between the ``Mr.''\ and
``Jones'' in
       ``Mr.~Jones'',        % ~ produces an unbreakable interword space.
or within a word---especially when the word is a
symbol like
       \mbox{\emph{itemnum}}
that makes little sense when hyphenated across
lines.

Footnotes\footnote{This is an example of a footnote.}
pose no problem.

\LaTeX\ is good at typesetting mathematical formulas
like
       \( x-3y + z = 7 \)
or
       \( a_{1} > x^{2n} + y^{2n} > x' \)
or
       \( \ip{A}{B} = \sum_{i} a_{i} b_{i} \).
The spaces you type in a formula are
ignored.  Remember that a letter like
       $x$                   % $ ... $  and  \( ... \)  are equivalent
is a formula when it denotes a mathematical
symbol, and it should be typed as one.

\section{Displayed Text}

Text is displayed by indenting it from the left
margin.  Quotations are commonly displayed.  There
are short quotations
\begin{quote}
   This is a short quotation.  It consists of a
   single paragraph of text.  See how it is formatted.
\end{quote}
and longer ones.
\begin{quotation}
   This is a longer quotation.  It consists of two
   paragraphs of text, neither of which are
   particularly interesting.

   This is the second paragraph of the quotation.  It
   is just as dull as the first paragraph.
\end{quotation}
Another frequently-displayed structure is a list.
The following is an example of an \emph{itemized}
list.
\begin{itemize}
   \item This is the first item of an itemized list.
         Each item in the list is marked with a ``tick''.
         You don't have to worry about what kind of tick
         mark is used.

   \item This is the second item of the list.  It
         contains another list nested inside it.  The inner
         list is an \emph{enumerated} list.
         \begin{enumerate}
            \item This is the first item of an enumerated
                  list that is nested within the itemized list.

            \item This is the second item of the inner list.
                  \LaTeX\ allows you to nest lists deeper than
                  you really should.
         \end{enumerate}
         This is the rest of the second item of the outer
         list.  It is no more interesting than any other
         part of the item.
   \item This is the third item of the list.
\end{itemize}
You can even display poetry.
\begin{verse}
   There is an environment
    for verse \\             % The \\ command separates lines
   Whose features some poets % within a stanza.
   will curse.

                             % One or more blank lines separate stanzas.

   For instead of making\\
   Them do \emph{all} line breaking, \\
   It allows them to put too many words on a line when they'd rather be
   forced to be terse.
\end{verse}

Mathematical formulas may also be displayed.  A
displayed formula
is
one-line long; multiline
formulas require special formatting instructions.
   \[  \ip{\Gamma}{\psi'} = x'' + y^{2} + z_{i}^{n}\]
Don't start a paragraph with a displayed equation,
nor make one a paragraph by itself.

\end{document}               % End of document.

\section{Bench Huge Section 1611}
% synthetic bench section 1611
% This is a sample LaTeX input file.  (Version of 12 August 2004.)
%
% A '%' character causes TeX to ignore all remaining text on the line,
% and is used for comments like this one.

\documentclass{article}      % Specifies the document class

                             % The preamble begins here.
\title{An Example Document}  % Declares the document's title.
\author{Leslie Lamport}      % Declares the author's name.
\date{January 21, 1994}      % Deleting this command produces today's date.

\newcommand{\ip}[2]{(#1, #2)}
                             % Defines \ip{arg1}{arg2} to mean
                             % (arg1, arg2).

%\newcommand{\ip}[2]{\langle #1 | #2\rangle}
                             % This is an alternative definition of
                             % \ip that is commented out.

\begin{document}             % End of preamble and beginning of text.

\maketitle                   % Produces the title.

This is an example input file.  Comparing it with
the output it generates can show you how to
produce a simple document of your own.

\section{Ordinary Text}      % Produces section heading.  Lower-level
                             % sections are begun with similar
                             % \subsection and \subsubsection commands.

The ends  of words and sentences are marked
  by   spaces. It  doesn't matter how many
spaces    you type; one is as good as 100.  The
end of   a line counts as a space.

One   or more   blank lines denote the  end
of  a paragraph.

Since any number of consecutive spaces are treated
like a single one, the formatting of the input
file makes no difference to
      \LaTeX,                % The \LaTeX command generates the LaTeX logo.
but it makes a difference to you.  When you use
\LaTeX, making your input file as easy to read
as possible will be a great help as you write
your document and when you change it.  This sample
file shows how you can add comments to your own input
file.

Because printing is different from typewriting,
there are a number of things that you have to do
differently when preparing an input file than if
you were just typing the document directly.
Quotation marks like
       ``this''
have to be handled specially, as do quotes within
quotes:
       ``\,`this'            % \, separates the double and single quote.
        is what I just
        wrote, not  `that'\,''.

Dashes come in three sizes: an
       intra-word
dash, a medium dash for number ranges like
       1--2,
and a punctuation
       dash---like
this.

A sentence-ending space should be larger than the
space between words within a sentence.  You
sometimes have to type special commands in
conjunction with punctuation characters to get
this right, as in the following sentence.
       Gnats, gnus, etc.\ all  % `\ ' makes an inter-word space.
       begin with G\@.         % \@ marks end-of-sentence punctuation.
You should check the spaces after periods when
reading your output to make sure you haven't
forgotten any special cases.  Generating an
ellipsis
       \ldots\               % `\ ' is needed after `\ldots' because TeX
                             % ignores spaces after command names like \ldots
                             % made from \ + letters.
                             %
                             % Note how a `%' character causes TeX to ignore
                             % the end of the input line, so these blank lines
                             % do not start a new paragraph.
                             %
with the right spacing around the periods requires
a special command.

\LaTeX\ interprets some common characters as
commands, so you must type special commands to
generate them.  These characters include the
following:
       \$ \& \% \# \{ and \}.

In printing, text is usually emphasized with an
       \emph{italic}
type style.

\begin{em}
   A long segment of text can also be emphasized
   in this way.  Text within such a segment can be
   given \emph{additional} emphasis.
\end{em}

It is sometimes necessary to prevent \LaTeX\ from
breaking a line where it might otherwise do so.
This may be at a space, as between the ``Mr.''\ and
``Jones'' in
       ``Mr.~Jones'',        % ~ produces an unbreakable interword space.
or within a word---especially when the word is a
symbol like
       \mbox{\emph{itemnum}}
that makes little sense when hyphenated across
lines.

Footnotes\footnote{This is an example of a footnote.}
pose no problem.

\LaTeX\ is good at typesetting mathematical formulas
like
       \( x-3y + z = 7 \)
or
       \( a_{1} > x^{2n} + y^{2n} > x' \)
or
       \( \ip{A}{B} = \sum_{i} a_{i} b_{i} \).
The spaces you type in a formula are
ignored.  Remember that a letter like
       $x$                   % $ ... $  and  \( ... \)  are equivalent
is a formula when it denotes a mathematical
symbol, and it should be typed as one.

\section{Displayed Text}

Text is displayed by indenting it from the left
margin.  Quotations are commonly displayed.  There
are short quotations
\begin{quote}
   This is a short quotation.  It consists of a
   single paragraph of text.  See how it is formatted.
\end{quote}
and longer ones.
\begin{quotation}
   This is a longer quotation.  It consists of two
   paragraphs of text, neither of which are
   particularly interesting.

   This is the second paragraph of the quotation.  It
   is just as dull as the first paragraph.
\end{quotation}
Another frequently-displayed structure is a list.
The following is an example of an \emph{itemized}
list.
\begin{itemize}
   \item This is the first item of an itemized list.
         Each item in the list is marked with a ``tick''.
         You don't have to worry about what kind of tick
         mark is used.

   \item This is the second item of the list.  It
         contains another list nested inside it.  The inner
         list is an \emph{enumerated} list.
         \begin{enumerate}
            \item This is the first item of an enumerated
                  list that is nested within the itemized list.

            \item This is the second item of the inner list.
                  \LaTeX\ allows you to nest lists deeper than
                  you really should.
         \end{enumerate}
         This is the rest of the second item of the outer
         list.  It is no more interesting than any other
         part of the item.
   \item This is the third item of the list.
\end{itemize}
You can even display poetry.
\begin{verse}
   There is an environment
    for verse \\             % The \\ command separates lines
   Whose features some poets % within a stanza.
   will curse.

                             % One or more blank lines separate stanzas.

   For instead of making\\
   Them do \emph{all} line breaking, \\
   It allows them to put too many words on a line when they'd rather be
   forced to be terse.
\end{verse}

Mathematical formulas may also be displayed.  A
displayed formula
is
one-line long; multiline
formulas require special formatting instructions.
   \[  \ip{\Gamma}{\psi'} = x'' + y^{2} + z_{i}^{n}\]
Don't start a paragraph with a displayed equation,
nor make one a paragraph by itself.

\end{document}               % End of document.

\section{Bench Huge Section 1612}
% synthetic bench section 1612
% This is a sample LaTeX input file.  (Version of 12 August 2004.)
%
% A '%' character causes TeX to ignore all remaining text on the line,
% and is used for comments like this one.

\documentclass{article}      % Specifies the document class

                             % The preamble begins here.
\title{An Example Document}  % Declares the document's title.
\author{Leslie Lamport}      % Declares the author's name.
\date{January 21, 1994}      % Deleting this command produces today's date.

\newcommand{\ip}[2]{(#1, #2)}
                             % Defines \ip{arg1}{arg2} to mean
                             % (arg1, arg2).

%\newcommand{\ip}[2]{\langle #1 | #2\rangle}
                             % This is an alternative definition of
                             % \ip that is commented out.

\begin{document}             % End of preamble and beginning of text.

\maketitle                   % Produces the title.

This is an example input file.  Comparing it with
the output it generates can show you how to
produce a simple document of your own.

\section{Ordinary Text}      % Produces section heading.  Lower-level
                             % sections are begun with similar
                             % \subsection and \subsubsection commands.

The ends  of words and sentences are marked
  by   spaces. It  doesn't matter how many
spaces    you type; one is as good as 100.  The
end of   a line counts as a space.

One   or more   blank lines denote the  end
of  a paragraph.

Since any number of consecutive spaces are treated
like a single one, the formatting of the input
file makes no difference to
      \LaTeX,                % The \LaTeX command generates the LaTeX logo.
but it makes a difference to you.  When you use
\LaTeX, making your input file as easy to read
as possible will be a great help as you write
your document and when you change it.  This sample
file shows how you can add comments to your own input
file.

Because printing is different from typewriting,
there are a number of things that you have to do
differently when preparing an input file than if
you were just typing the document directly.
Quotation marks like
       ``this''
have to be handled specially, as do quotes within
quotes:
       ``\,`this'            % \, separates the double and single quote.
        is what I just
        wrote, not  `that'\,''.

Dashes come in three sizes: an
       intra-word
dash, a medium dash for number ranges like
       1--2,
and a punctuation
       dash---like
this.

A sentence-ending space should be larger than the
space between words within a sentence.  You
sometimes have to type special commands in
conjunction with punctuation characters to get
this right, as in the following sentence.
       Gnats, gnus, etc.\ all  % `\ ' makes an inter-word space.
       begin with G\@.         % \@ marks end-of-sentence punctuation.
You should check the spaces after periods when
reading your output to make sure you haven't
forgotten any special cases.  Generating an
ellipsis
       \ldots\               % `\ ' is needed after `\ldots' because TeX
                             % ignores spaces after command names like \ldots
                             % made from \ + letters.
                             %
                             % Note how a `%' character causes TeX to ignore
                             % the end of the input line, so these blank lines
                             % do not start a new paragraph.
                             %
with the right spacing around the periods requires
a special command.

\LaTeX\ interprets some common characters as
commands, so you must type special commands to
generate them.  These characters include the
following:
       \$ \& \% \# \{ and \}.

In printing, text is usually emphasized with an
       \emph{italic}
type style.

\begin{em}
   A long segment of text can also be emphasized
   in this way.  Text within such a segment can be
   given \emph{additional} emphasis.
\end{em}

It is sometimes necessary to prevent \LaTeX\ from
breaking a line where it might otherwise do so.
This may be at a space, as between the ``Mr.''\ and
``Jones'' in
       ``Mr.~Jones'',        % ~ produces an unbreakable interword space.
or within a word---especially when the word is a
symbol like
       \mbox{\emph{itemnum}}
that makes little sense when hyphenated across
lines.

Footnotes\footnote{This is an example of a footnote.}
pose no problem.

\LaTeX\ is good at typesetting mathematical formulas
like
       \( x-3y + z = 7 \)
or
       \( a_{1} > x^{2n} + y^{2n} > x' \)
or
       \( \ip{A}{B} = \sum_{i} a_{i} b_{i} \).
The spaces you type in a formula are
ignored.  Remember that a letter like
       $x$                   % $ ... $  and  \( ... \)  are equivalent
is a formula when it denotes a mathematical
symbol, and it should be typed as one.

\section{Displayed Text}

Text is displayed by indenting it from the left
margin.  Quotations are commonly displayed.  There
are short quotations
\begin{quote}
   This is a short quotation.  It consists of a
   single paragraph of text.  See how it is formatted.
\end{quote}
and longer ones.
\begin{quotation}
   This is a longer quotation.  It consists of two
   paragraphs of text, neither of which are
   particularly interesting.

   This is the second paragraph of the quotation.  It
   is just as dull as the first paragraph.
\end{quotation}
Another frequently-displayed structure is a list.
The following is an example of an \emph{itemized}
list.
\begin{itemize}
   \item This is the first item of an itemized list.
         Each item in the list is marked with a ``tick''.
         You don't have to worry about what kind of tick
         mark is used.

   \item This is the second item of the list.  It
         contains another list nested inside it.  The inner
         list is an \emph{enumerated} list.
         \begin{enumerate}
            \item This is the first item of an enumerated
                  list that is nested within the itemized list.

            \item This is the second item of the inner list.
                  \LaTeX\ allows you to nest lists deeper than
                  you really should.
         \end{enumerate}
         This is the rest of the second item of the outer
         list.  It is no more interesting than any other
         part of the item.
   \item This is the third item of the list.
\end{itemize}
You can even display poetry.
\begin{verse}
   There is an environment
    for verse \\             % The \\ command separates lines
   Whose features some poets % within a stanza.
   will curse.

                             % One or more blank lines separate stanzas.

   For instead of making\\
   Them do \emph{all} line breaking, \\
   It allows them to put too many words on a line when they'd rather be
   forced to be terse.
\end{verse}

Mathematical formulas may also be displayed.  A
displayed formula
is
one-line long; multiline
formulas require special formatting instructions.
   \[  \ip{\Gamma}{\psi'} = x'' + y^{2} + z_{i}^{n}\]
Don't start a paragraph with a displayed equation,
nor make one a paragraph by itself.

\end{document}               % End of document.

\section{Bench Huge Section 1613}
% synthetic bench section 1613
% This is a sample LaTeX input file.  (Version of 12 August 2004.)
%
% A '%' character causes TeX to ignore all remaining text on the line,
% and is used for comments like this one.

\documentclass{article}      % Specifies the document class

                             % The preamble begins here.
\title{An Example Document}  % Declares the document's title.
\author{Leslie Lamport}      % Declares the author's name.
\date{January 21, 1994}      % Deleting this command produces today's date.

\newcommand{\ip}[2]{(#1, #2)}
                             % Defines \ip{arg1}{arg2} to mean
                             % (arg1, arg2).

%\newcommand{\ip}[2]{\langle #1 | #2\rangle}
                             % This is an alternative definition of
                             % \ip that is commented out.

\begin{document}             % End of preamble and beginning of text.

\maketitle                   % Produces the title.

This is an example input file.  Comparing it with
the output it generates can show you how to
produce a simple document of your own.

\section{Ordinary Text}      % Produces section heading.  Lower-level
                             % sections are begun with similar
                             % \subsection and \subsubsection commands.

The ends  of words and sentences are marked
  by   spaces. It  doesn't matter how many
spaces    you type; one is as good as 100.  The
end of   a line counts as a space.

One   or more   blank lines denote the  end
of  a paragraph.

Since any number of consecutive spaces are treated
like a single one, the formatting of the input
file makes no difference to
      \LaTeX,                % The \LaTeX command generates the LaTeX logo.
but it makes a difference to you.  When you use
\LaTeX, making your input file as easy to read
as possible will be a great help as you write
your document and when you change it.  This sample
file shows how you can add comments to your own input
file.

Because printing is different from typewriting,
there are a number of things that you have to do
differently when preparing an input file than if
you were just typing the document directly.
Quotation marks like
       ``this''
have to be handled specially, as do quotes within
quotes:
       ``\,`this'            % \, separates the double and single quote.
        is what I just
        wrote, not  `that'\,''.

Dashes come in three sizes: an
       intra-word
dash, a medium dash for number ranges like
       1--2,
and a punctuation
       dash---like
this.

A sentence-ending space should be larger than the
space between words within a sentence.  You
sometimes have to type special commands in
conjunction with punctuation characters to get
this right, as in the following sentence.
       Gnats, gnus, etc.\ all  % `\ ' makes an inter-word space.
       begin with G\@.         % \@ marks end-of-sentence punctuation.
You should check the spaces after periods when
reading your output to make sure you haven't
forgotten any special cases.  Generating an
ellipsis
       \ldots\               % `\ ' is needed after `\ldots' because TeX
                             % ignores spaces after command names like \ldots
                             % made from \ + letters.
                             %
                             % Note how a `%' character causes TeX to ignore
                             % the end of the input line, so these blank lines
                             % do not start a new paragraph.
                             %
with the right spacing around the periods requires
a special command.

\LaTeX\ interprets some common characters as
commands, so you must type special commands to
generate them.  These characters include the
following:
       \$ \& \% \# \{ and \}.

In printing, text is usually emphasized with an
       \emph{italic}
type style.

\begin{em}
   A long segment of text can also be emphasized
   in this way.  Text within such a segment can be
   given \emph{additional} emphasis.
\end{em}

It is sometimes necessary to prevent \LaTeX\ from
breaking a line where it might otherwise do so.
This may be at a space, as between the ``Mr.''\ and
``Jones'' in
       ``Mr.~Jones'',        % ~ produces an unbreakable interword space.
or within a word---especially when the word is a
symbol like
       \mbox{\emph{itemnum}}
that makes little sense when hyphenated across
lines.

Footnotes\footnote{This is an example of a footnote.}
pose no problem.

\LaTeX\ is good at typesetting mathematical formulas
like
       \( x-3y + z = 7 \)
or
       \( a_{1} > x^{2n} + y^{2n} > x' \)
or
       \( \ip{A}{B} = \sum_{i} a_{i} b_{i} \).
The spaces you type in a formula are
ignored.  Remember that a letter like
       $x$                   % $ ... $  and  \( ... \)  are equivalent
is a formula when it denotes a mathematical
symbol, and it should be typed as one.

\section{Displayed Text}

Text is displayed by indenting it from the left
margin.  Quotations are commonly displayed.  There
are short quotations
\begin{quote}
   This is a short quotation.  It consists of a
   single paragraph of text.  See how it is formatted.
\end{quote}
and longer ones.
\begin{quotation}
   This is a longer quotation.  It consists of two
   paragraphs of text, neither of which are
   particularly interesting.

   This is the second paragraph of the quotation.  It
   is just as dull as the first paragraph.
\end{quotation}
Another frequently-displayed structure is a list.
The following is an example of an \emph{itemized}
list.
\begin{itemize}
   \item This is the first item of an itemized list.
         Each item in the list is marked with a ``tick''.
         You don't have to worry about what kind of tick
         mark is used.

   \item This is the second item of the list.  It
         contains another list nested inside it.  The inner
         list is an \emph{enumerated} list.
         \begin{enumerate}
            \item This is the first item of an enumerated
                  list that is nested within the itemized list.

            \item This is the second item of the inner list.
                  \LaTeX\ allows you to nest lists deeper than
                  you really should.
         \end{enumerate}
         This is the rest of the second item of the outer
         list.  It is no more interesting than any other
         part of the item.
   \item This is the third item of the list.
\end{itemize}
You can even display poetry.
\begin{verse}
   There is an environment
    for verse \\             % The \\ command separates lines
   Whose features some poets % within a stanza.
   will curse.

                             % One or more blank lines separate stanzas.

   For instead of making\\
   Them do \emph{all} line breaking, \\
   It allows them to put too many words on a line when they'd rather be
   forced to be terse.
\end{verse}

Mathematical formulas may also be displayed.  A
displayed formula
is
one-line long; multiline
formulas require special formatting instructions.
   \[  \ip{\Gamma}{\psi'} = x'' + y^{2} + z_{i}^{n}\]
Don't start a paragraph with a displayed equation,
nor make one a paragraph by itself.

\end{document}               % End of document.

\section{Bench Huge Section 1614}
% synthetic bench section 1614
% This is a sample LaTeX input file.  (Version of 12 August 2004.)
%
% A '%' character causes TeX to ignore all remaining text on the line,
% and is used for comments like this one.

\documentclass{article}      % Specifies the document class

                             % The preamble begins here.
\title{An Example Document}  % Declares the document's title.
\author{Leslie Lamport}      % Declares the author's name.
\date{January 21, 1994}      % Deleting this command produces today's date.

\newcommand{\ip}[2]{(#1, #2)}
                             % Defines \ip{arg1}{arg2} to mean
                             % (arg1, arg2).

%\newcommand{\ip}[2]{\langle #1 | #2\rangle}
                             % This is an alternative definition of
                             % \ip that is commented out.

\begin{document}             % End of preamble and beginning of text.

\maketitle                   % Produces the title.

This is an example input file.  Comparing it with
the output it generates can show you how to
produce a simple document of your own.

\section{Ordinary Text}      % Produces section heading.  Lower-level
                             % sections are begun with similar
                             % \subsection and \subsubsection commands.

The ends  of words and sentences are marked
  by   spaces. It  doesn't matter how many
spaces    you type; one is as good as 100.  The
end of   a line counts as a space.

One   or more   blank lines denote the  end
of  a paragraph.

Since any number of consecutive spaces are treated
like a single one, the formatting of the input
file makes no difference to
      \LaTeX,                % The \LaTeX command generates the LaTeX logo.
but it makes a difference to you.  When you use
\LaTeX, making your input file as easy to read
as possible will be a great help as you write
your document and when you change it.  This sample
file shows how you can add comments to your own input
file.

Because printing is different from typewriting,
there are a number of things that you have to do
differently when preparing an input file than if
you were just typing the document directly.
Quotation marks like
       ``this''
have to be handled specially, as do quotes within
quotes:
       ``\,`this'            % \, separates the double and single quote.
        is what I just
        wrote, not  `that'\,''.

Dashes come in three sizes: an
       intra-word
dash, a medium dash for number ranges like
       1--2,
and a punctuation
       dash---like
this.

A sentence-ending space should be larger than the
space between words within a sentence.  You
sometimes have to type special commands in
conjunction with punctuation characters to get
this right, as in the following sentence.
       Gnats, gnus, etc.\ all  % `\ ' makes an inter-word space.
       begin with G\@.         % \@ marks end-of-sentence punctuation.
You should check the spaces after periods when
reading your output to make sure you haven't
forgotten any special cases.  Generating an
ellipsis
       \ldots\               % `\ ' is needed after `\ldots' because TeX
                             % ignores spaces after command names like \ldots
                             % made from \ + letters.
                             %
                             % Note how a `%' character causes TeX to ignore
                             % the end of the input line, so these blank lines
                             % do not start a new paragraph.
                             %
with the right spacing around the periods requires
a special command.

\LaTeX\ interprets some common characters as
commands, so you must type special commands to
generate them.  These characters include the
following:
       \$ \& \% \# \{ and \}.

In printing, text is usually emphasized with an
       \emph{italic}
type style.

\begin{em}
   A long segment of text can also be emphasized
   in this way.  Text within such a segment can be
   given \emph{additional} emphasis.
\end{em}

It is sometimes necessary to prevent \LaTeX\ from
breaking a line where it might otherwise do so.
This may be at a space, as between the ``Mr.''\ and
``Jones'' in
       ``Mr.~Jones'',        % ~ produces an unbreakable interword space.
or within a word---especially when the word is a
symbol like
       \mbox{\emph{itemnum}}
that makes little sense when hyphenated across
lines.

Footnotes\footnote{This is an example of a footnote.}
pose no problem.

\LaTeX\ is good at typesetting mathematical formulas
like
       \( x-3y + z = 7 \)
or
       \( a_{1} > x^{2n} + y^{2n} > x' \)
or
       \( \ip{A}{B} = \sum_{i} a_{i} b_{i} \).
The spaces you type in a formula are
ignored.  Remember that a letter like
       $x$                   % $ ... $  and  \( ... \)  are equivalent
is a formula when it denotes a mathematical
symbol, and it should be typed as one.

\section{Displayed Text}

Text is displayed by indenting it from the left
margin.  Quotations are commonly displayed.  There
are short quotations
\begin{quote}
   This is a short quotation.  It consists of a
   single paragraph of text.  See how it is formatted.
\end{quote}
and longer ones.
\begin{quotation}
   This is a longer quotation.  It consists of two
   paragraphs of text, neither of which are
   particularly interesting.

   This is the second paragraph of the quotation.  It
   is just as dull as the first paragraph.
\end{quotation}
Another frequently-displayed structure is a list.
The following is an example of an \emph{itemized}
list.
\begin{itemize}
   \item This is the first item of an itemized list.
         Each item in the list is marked with a ``tick''.
         You don't have to worry about what kind of tick
         mark is used.

   \item This is the second item of the list.  It
         contains another list nested inside it.  The inner
         list is an \emph{enumerated} list.
         \begin{enumerate}
            \item This is the first item of an enumerated
                  list that is nested within the itemized list.

            \item This is the second item of the inner list.
                  \LaTeX\ allows you to nest lists deeper than
                  you really should.
         \end{enumerate}
         This is the rest of the second item of the outer
         list.  It is no more interesting than any other
         part of the item.
   \item This is the third item of the list.
\end{itemize}
You can even display poetry.
\begin{verse}
   There is an environment
    for verse \\             % The \\ command separates lines
   Whose features some poets % within a stanza.
   will curse.

                             % One or more blank lines separate stanzas.

   For instead of making\\
   Them do \emph{all} line breaking, \\
   It allows them to put too many words on a line when they'd rather be
   forced to be terse.
\end{verse}

Mathematical formulas may also be displayed.  A
displayed formula
is
one-line long; multiline
formulas require special formatting instructions.
   \[  \ip{\Gamma}{\psi'} = x'' + y^{2} + z_{i}^{n}\]
Don't start a paragraph with a displayed equation,
nor make one a paragraph by itself.

\end{document}               % End of document.

\section{Bench Huge Section 1615}
% synthetic bench section 1615
% This is a sample LaTeX input file.  (Version of 12 August 2004.)
%
% A '%' character causes TeX to ignore all remaining text on the line,
% and is used for comments like this one.

\documentclass{article}      % Specifies the document class

                             % The preamble begins here.
\title{An Example Document}  % Declares the document's title.
\author{Leslie Lamport}      % Declares the author's name.
\date{January 21, 1994}      % Deleting this command produces today's date.

\newcommand{\ip}[2]{(#1, #2)}
                             % Defines \ip{arg1}{arg2} to mean
                             % (arg1, arg2).

%\newcommand{\ip}[2]{\langle #1 | #2\rangle}
                             % This is an alternative definition of
                             % \ip that is commented out.

\begin{document}             % End of preamble and beginning of text.

\maketitle                   % Produces the title.

This is an example input file.  Comparing it with
the output it generates can show you how to
produce a simple document of your own.

\section{Ordinary Text}      % Produces section heading.  Lower-level
                             % sections are begun with similar
                             % \subsection and \subsubsection commands.

The ends  of words and sentences are marked
  by   spaces. It  doesn't matter how many
spaces    you type; one is as good as 100.  The
end of   a line counts as a space.

One   or more   blank lines denote the  end
of  a paragraph.

Since any number of consecutive spaces are treated
like a single one, the formatting of the input
file makes no difference to
      \LaTeX,                % The \LaTeX command generates the LaTeX logo.
but it makes a difference to you.  When you use
\LaTeX, making your input file as easy to read
as possible will be a great help as you write
your document and when you change it.  This sample
file shows how you can add comments to your own input
file.

Because printing is different from typewriting,
there are a number of things that you have to do
differently when preparing an input file than if
you were just typing the document directly.
Quotation marks like
       ``this''
have to be handled specially, as do quotes within
quotes:
       ``\,`this'            % \, separates the double and single quote.
        is what I just
        wrote, not  `that'\,''.

Dashes come in three sizes: an
       intra-word
dash, a medium dash for number ranges like
       1--2,
and a punctuation
       dash---like
this.

A sentence-ending space should be larger than the
space between words within a sentence.  You
sometimes have to type special commands in
conjunction with punctuation characters to get
this right, as in the following sentence.
       Gnats, gnus, etc.\ all  % `\ ' makes an inter-word space.
       begin with G\@.         % \@ marks end-of-sentence punctuation.
You should check the spaces after periods when
reading your output to make sure you haven't
forgotten any special cases.  Generating an
ellipsis
       \ldots\               % `\ ' is needed after `\ldots' because TeX
                             % ignores spaces after command names like \ldots
                             % made from \ + letters.
                             %
                             % Note how a `%' character causes TeX to ignore
                             % the end of the input line, so these blank lines
                             % do not start a new paragraph.
                             %
with the right spacing around the periods requires
a special command.

\LaTeX\ interprets some common characters as
commands, so you must type special commands to
generate them.  These characters include the
following:
       \$ \& \% \# \{ and \}.

In printing, text is usually emphasized with an
       \emph{italic}
type style.

\begin{em}
   A long segment of text can also be emphasized
   in this way.  Text within such a segment can be
   given \emph{additional} emphasis.
\end{em}

It is sometimes necessary to prevent \LaTeX\ from
breaking a line where it might otherwise do so.
This may be at a space, as between the ``Mr.''\ and
``Jones'' in
       ``Mr.~Jones'',        % ~ produces an unbreakable interword space.
or within a word---especially when the word is a
symbol like
       \mbox{\emph{itemnum}}
that makes little sense when hyphenated across
lines.

Footnotes\footnote{This is an example of a footnote.}
pose no problem.

\LaTeX\ is good at typesetting mathematical formulas
like
       \( x-3y + z = 7 \)
or
       \( a_{1} > x^{2n} + y^{2n} > x' \)
or
       \( \ip{A}{B} = \sum_{i} a_{i} b_{i} \).
The spaces you type in a formula are
ignored.  Remember that a letter like
       $x$                   % $ ... $  and  \( ... \)  are equivalent
is a formula when it denotes a mathematical
symbol, and it should be typed as one.

\section{Displayed Text}

Text is displayed by indenting it from the left
margin.  Quotations are commonly displayed.  There
are short quotations
\begin{quote}
   This is a short quotation.  It consists of a
   single paragraph of text.  See how it is formatted.
\end{quote}
and longer ones.
\begin{quotation}
   This is a longer quotation.  It consists of two
   paragraphs of text, neither of which are
   particularly interesting.

   This is the second paragraph of the quotation.  It
   is just as dull as the first paragraph.
\end{quotation}
Another frequently-displayed structure is a list.
The following is an example of an \emph{itemized}
list.
\begin{itemize}
   \item This is the first item of an itemized list.
         Each item in the list is marked with a ``tick''.
         You don't have to worry about what kind of tick
         mark is used.

   \item This is the second item of the list.  It
         contains another list nested inside it.  The inner
         list is an \emph{enumerated} list.
         \begin{enumerate}
            \item This is the first item of an enumerated
                  list that is nested within the itemized list.

            \item This is the second item of the inner list.
                  \LaTeX\ allows you to nest lists deeper than
                  you really should.
         \end{enumerate}
         This is the rest of the second item of the outer
         list.  It is no more interesting than any other
         part of the item.
   \item This is the third item of the list.
\end{itemize}
You can even display poetry.
\begin{verse}
   There is an environment
    for verse \\             % The \\ command separates lines
   Whose features some poets % within a stanza.
   will curse.

                             % One or more blank lines separate stanzas.

   For instead of making\\
   Them do \emph{all} line breaking, \\
   It allows them to put too many words on a line when they'd rather be
   forced to be terse.
\end{verse}

Mathematical formulas may also be displayed.  A
displayed formula
is
one-line long; multiline
formulas require special formatting instructions.
   \[  \ip{\Gamma}{\psi'} = x'' + y^{2} + z_{i}^{n}\]
Don't start a paragraph with a displayed equation,
nor make one a paragraph by itself.

\end{document}               % End of document.

\section{Bench Huge Section 1616}
% synthetic bench section 1616
% This is a sample LaTeX input file.  (Version of 12 August 2004.)
%
% A '%' character causes TeX to ignore all remaining text on the line,
% and is used for comments like this one.

\documentclass{article}      % Specifies the document class

                             % The preamble begins here.
\title{An Example Document}  % Declares the document's title.
\author{Leslie Lamport}      % Declares the author's name.
\date{January 21, 1994}      % Deleting this command produces today's date.

\newcommand{\ip}[2]{(#1, #2)}
                             % Defines \ip{arg1}{arg2} to mean
                             % (arg1, arg2).

%\newcommand{\ip}[2]{\langle #1 | #2\rangle}
                             % This is an alternative definition of
                             % \ip that is commented out.

\begin{document}             % End of preamble and beginning of text.

\maketitle                   % Produces the title.

This is an example input file.  Comparing it with
the output it generates can show you how to
produce a simple document of your own.

\section{Ordinary Text}      % Produces section heading.  Lower-level
                             % sections are begun with similar
                             % \subsection and \subsubsection commands.

The ends  of words and sentences are marked
  by   spaces. It  doesn't matter how many
spaces    you type; one is as good as 100.  The
end of   a line counts as a space.

One   or more   blank lines denote the  end
of  a paragraph.

Since any number of consecutive spaces are treated
like a single one, the formatting of the input
file makes no difference to
      \LaTeX,                % The \LaTeX command generates the LaTeX logo.
but it makes a difference to you.  When you use
\LaTeX, making your input file as easy to read
as possible will be a great help as you write
your document and when you change it.  This sample
file shows how you can add comments to your own input
file.

Because printing is different from typewriting,
there are a number of things that you have to do
differently when preparing an input file than if
you were just typing the document directly.
Quotation marks like
       ``this''
have to be handled specially, as do quotes within
quotes:
       ``\,`this'            % \, separates the double and single quote.
        is what I just
        wrote, not  `that'\,''.

Dashes come in three sizes: an
       intra-word
dash, a medium dash for number ranges like
       1--2,
and a punctuation
       dash---like
this.

A sentence-ending space should be larger than the
space between words within a sentence.  You
sometimes have to type special commands in
conjunction with punctuation characters to get
this right, as in the following sentence.
       Gnats, gnus, etc.\ all  % `\ ' makes an inter-word space.
       begin with G\@.         % \@ marks end-of-sentence punctuation.
You should check the spaces after periods when
reading your output to make sure you haven't
forgotten any special cases.  Generating an
ellipsis
       \ldots\               % `\ ' is needed after `\ldots' because TeX
                             % ignores spaces after command names like \ldots
                             % made from \ + letters.
                             %
                             % Note how a `%' character causes TeX to ignore
                             % the end of the input line, so these blank lines
                             % do not start a new paragraph.
                             %
with the right spacing around the periods requires
a special command.

\LaTeX\ interprets some common characters as
commands, so you must type special commands to
generate them.  These characters include the
following:
       \$ \& \% \# \{ and \}.

In printing, text is usually emphasized with an
       \emph{italic}
type style.

\begin{em}
   A long segment of text can also be emphasized
   in this way.  Text within such a segment can be
   given \emph{additional} emphasis.
\end{em}

It is sometimes necessary to prevent \LaTeX\ from
breaking a line where it might otherwise do so.
This may be at a space, as between the ``Mr.''\ and
``Jones'' in
       ``Mr.~Jones'',        % ~ produces an unbreakable interword space.
or within a word---especially when the word is a
symbol like
       \mbox{\emph{itemnum}}
that makes little sense when hyphenated across
lines.

Footnotes\footnote{This is an example of a footnote.}
pose no problem.

\LaTeX\ is good at typesetting mathematical formulas
like
       \( x-3y + z = 7 \)
or
       \( a_{1} > x^{2n} + y^{2n} > x' \)
or
       \( \ip{A}{B} = \sum_{i} a_{i} b_{i} \).
The spaces you type in a formula are
ignored.  Remember that a letter like
       $x$                   % $ ... $  and  \( ... \)  are equivalent
is a formula when it denotes a mathematical
symbol, and it should be typed as one.

\section{Displayed Text}

Text is displayed by indenting it from the left
margin.  Quotations are commonly displayed.  There
are short quotations
\begin{quote}
   This is a short quotation.  It consists of a
   single paragraph of text.  See how it is formatted.
\end{quote}
and longer ones.
\begin{quotation}
   This is a longer quotation.  It consists of two
   paragraphs of text, neither of which are
   particularly interesting.

   This is the second paragraph of the quotation.  It
   is just as dull as the first paragraph.
\end{quotation}
Another frequently-displayed structure is a list.
The following is an example of an \emph{itemized}
list.
\begin{itemize}
   \item This is the first item of an itemized list.
         Each item in the list is marked with a ``tick''.
         You don't have to worry about what kind of tick
         mark is used.

   \item This is the second item of the list.  It
         contains another list nested inside it.  The inner
         list is an \emph{enumerated} list.
         \begin{enumerate}
            \item This is the first item of an enumerated
                  list that is nested within the itemized list.

            \item This is the second item of the inner list.
                  \LaTeX\ allows you to nest lists deeper than
                  you really should.
         \end{enumerate}
         This is the rest of the second item of the outer
         list.  It is no more interesting than any other
         part of the item.
   \item This is the third item of the list.
\end{itemize}
You can even display poetry.
\begin{verse}
   There is an environment
    for verse \\             % The \\ command separates lines
   Whose features some poets % within a stanza.
   will curse.

                             % One or more blank lines separate stanzas.

   For instead of making\\
   Them do \emph{all} line breaking, \\
   It allows them to put too many words on a line when they'd rather be
   forced to be terse.
\end{verse}

Mathematical formulas may also be displayed.  A
displayed formula
is
one-line long; multiline
formulas require special formatting instructions.
   \[  \ip{\Gamma}{\psi'} = x'' + y^{2} + z_{i}^{n}\]
Don't start a paragraph with a displayed equation,
nor make one a paragraph by itself.

\end{document}               % End of document.

\section{Bench Huge Section 1617}
% synthetic bench section 1617
% This is a sample LaTeX input file.  (Version of 12 August 2004.)
%
% A '%' character causes TeX to ignore all remaining text on the line,
% and is used for comments like this one.

\documentclass{article}      % Specifies the document class

                             % The preamble begins here.
\title{An Example Document}  % Declares the document's title.
\author{Leslie Lamport}      % Declares the author's name.
\date{January 21, 1994}      % Deleting this command produces today's date.

\newcommand{\ip}[2]{(#1, #2)}
                             % Defines \ip{arg1}{arg2} to mean
                             % (arg1, arg2).

%\newcommand{\ip}[2]{\langle #1 | #2\rangle}
                             % This is an alternative definition of
                             % \ip that is commented out.

\begin{document}             % End of preamble and beginning of text.

\maketitle                   % Produces the title.

This is an example input file.  Comparing it with
the output it generates can show you how to
produce a simple document of your own.

\section{Ordinary Text}      % Produces section heading.  Lower-level
                             % sections are begun with similar
                             % \subsection and \subsubsection commands.

The ends  of words and sentences are marked
  by   spaces. It  doesn't matter how many
spaces    you type; one is as good as 100.  The
end of   a line counts as a space.

One   or more   blank lines denote the  end
of  a paragraph.

Since any number of consecutive spaces are treated
like a single one, the formatting of the input
file makes no difference to
      \LaTeX,                % The \LaTeX command generates the LaTeX logo.
but it makes a difference to you.  When you use
\LaTeX, making your input file as easy to read
as possible will be a great help as you write
your document and when you change it.  This sample
file shows how you can add comments to your own input
file.

Because printing is different from typewriting,
there are a number of things that you have to do
differently when preparing an input file than if
you were just typing the document directly.
Quotation marks like
       ``this''
have to be handled specially, as do quotes within
quotes:
       ``\,`this'            % \, separates the double and single quote.
        is what I just
        wrote, not  `that'\,''.

Dashes come in three sizes: an
       intra-word
dash, a medium dash for number ranges like
       1--2,
and a punctuation
       dash---like
this.

A sentence-ending space should be larger than the
space between words within a sentence.  You
sometimes have to type special commands in
conjunction with punctuation characters to get
this right, as in the following sentence.
       Gnats, gnus, etc.\ all  % `\ ' makes an inter-word space.
       begin with G\@.         % \@ marks end-of-sentence punctuation.
You should check the spaces after periods when
reading your output to make sure you haven't
forgotten any special cases.  Generating an
ellipsis
       \ldots\               % `\ ' is needed after `\ldots' because TeX
                             % ignores spaces after command names like \ldots
                             % made from \ + letters.
                             %
                             % Note how a `%' character causes TeX to ignore
                             % the end of the input line, so these blank lines
                             % do not start a new paragraph.
                             %
with the right spacing around the periods requires
a special command.

\LaTeX\ interprets some common characters as
commands, so you must type special commands to
generate them.  These characters include the
following:
       \$ \& \% \# \{ and \}.

In printing, text is usually emphasized with an
       \emph{italic}
type style.

\begin{em}
   A long segment of text can also be emphasized
   in this way.  Text within such a segment can be
   given \emph{additional} emphasis.
\end{em}

It is sometimes necessary to prevent \LaTeX\ from
breaking a line where it might otherwise do so.
This may be at a space, as between the ``Mr.''\ and
``Jones'' in
       ``Mr.~Jones'',        % ~ produces an unbreakable interword space.
or within a word---especially when the word is a
symbol like
       \mbox{\emph{itemnum}}
that makes little sense when hyphenated across
lines.

Footnotes\footnote{This is an example of a footnote.}
pose no problem.

\LaTeX\ is good at typesetting mathematical formulas
like
       \( x-3y + z = 7 \)
or
       \( a_{1} > x^{2n} + y^{2n} > x' \)
or
       \( \ip{A}{B} = \sum_{i} a_{i} b_{i} \).
The spaces you type in a formula are
ignored.  Remember that a letter like
       $x$                   % $ ... $  and  \( ... \)  are equivalent
is a formula when it denotes a mathematical
symbol, and it should be typed as one.

\section{Displayed Text}

Text is displayed by indenting it from the left
margin.  Quotations are commonly displayed.  There
are short quotations
\begin{quote}
   This is a short quotation.  It consists of a
   single paragraph of text.  See how it is formatted.
\end{quote}
and longer ones.
\begin{quotation}
   This is a longer quotation.  It consists of two
   paragraphs of text, neither of which are
   particularly interesting.

   This is the second paragraph of the quotation.  It
   is just as dull as the first paragraph.
\end{quotation}
Another frequently-displayed structure is a list.
The following is an example of an \emph{itemized}
list.
\begin{itemize}
   \item This is the first item of an itemized list.
         Each item in the list is marked with a ``tick''.
         You don't have to worry about what kind of tick
         mark is used.

   \item This is the second item of the list.  It
         contains another list nested inside it.  The inner
         list is an \emph{enumerated} list.
         \begin{enumerate}
            \item This is the first item of an enumerated
                  list that is nested within the itemized list.

            \item This is the second item of the inner list.
                  \LaTeX\ allows you to nest lists deeper than
                  you really should.
         \end{enumerate}
         This is the rest of the second item of the outer
         list.  It is no more interesting than any other
         part of the item.
   \item This is the third item of the list.
\end{itemize}
You can even display poetry.
\begin{verse}
   There is an environment
    for verse \\             % The \\ command separates lines
   Whose features some poets % within a stanza.
   will curse.

                             % One or more blank lines separate stanzas.

   For instead of making\\
   Them do \emph{all} line breaking, \\
   It allows them to put too many words on a line when they'd rather be
   forced to be terse.
\end{verse}

Mathematical formulas may also be displayed.  A
displayed formula
is
one-line long; multiline
formulas require special formatting instructions.
   \[  \ip{\Gamma}{\psi'} = x'' + y^{2} + z_{i}^{n}\]
Don't start a paragraph with a displayed equation,
nor make one a paragraph by itself.

\end{document}               % End of document.

\section{Bench Huge Section 1618}
% synthetic bench section 1618
% This is a sample LaTeX input file.  (Version of 12 August 2004.)
%
% A '%' character causes TeX to ignore all remaining text on the line,
% and is used for comments like this one.

\documentclass{article}      % Specifies the document class

                             % The preamble begins here.
\title{An Example Document}  % Declares the document's title.
\author{Leslie Lamport}      % Declares the author's name.
\date{January 21, 1994}      % Deleting this command produces today's date.

\newcommand{\ip}[2]{(#1, #2)}
                             % Defines \ip{arg1}{arg2} to mean
                             % (arg1, arg2).

%\newcommand{\ip}[2]{\langle #1 | #2\rangle}
                             % This is an alternative definition of
                             % \ip that is commented out.

\begin{document}             % End of preamble and beginning of text.

\maketitle                   % Produces the title.

This is an example input file.  Comparing it with
the output it generates can show you how to
produce a simple document of your own.

\section{Ordinary Text}      % Produces section heading.  Lower-level
                             % sections are begun with similar
                             % \subsection and \subsubsection commands.

The ends  of words and sentences are marked
  by   spaces. It  doesn't matter how many
spaces    you type; one is as good as 100.  The
end of   a line counts as a space.

One   or more   blank lines denote the  end
of  a paragraph.

Since any number of consecutive spaces are treated
like a single one, the formatting of the input
file makes no difference to
      \LaTeX,                % The \LaTeX command generates the LaTeX logo.
but it makes a difference to you.  When you use
\LaTeX, making your input file as easy to read
as possible will be a great help as you write
your document and when you change it.  This sample
file shows how you can add comments to your own input
file.

Because printing is different from typewriting,
there are a number of things that you have to do
differently when preparing an input file than if
you were just typing the document directly.
Quotation marks like
       ``this''
have to be handled specially, as do quotes within
quotes:
       ``\,`this'            % \, separates the double and single quote.
        is what I just
        wrote, not  `that'\,''.

Dashes come in three sizes: an
       intra-word
dash, a medium dash for number ranges like
       1--2,
and a punctuation
       dash---like
this.

A sentence-ending space should be larger than the
space between words within a sentence.  You
sometimes have to type special commands in
conjunction with punctuation characters to get
this right, as in the following sentence.
       Gnats, gnus, etc.\ all  % `\ ' makes an inter-word space.
       begin with G\@.         % \@ marks end-of-sentence punctuation.
You should check the spaces after periods when
reading your output to make sure you haven't
forgotten any special cases.  Generating an
ellipsis
       \ldots\               % `\ ' is needed after `\ldots' because TeX
                             % ignores spaces after command names like \ldots
                             % made from \ + letters.
                             %
                             % Note how a `%' character causes TeX to ignore
                             % the end of the input line, so these blank lines
                             % do not start a new paragraph.
                             %
with the right spacing around the periods requires
a special command.

\LaTeX\ interprets some common characters as
commands, so you must type special commands to
generate them.  These characters include the
following:
       \$ \& \% \# \{ and \}.

In printing, text is usually emphasized with an
       \emph{italic}
type style.

\begin{em}
   A long segment of text can also be emphasized
   in this way.  Text within such a segment can be
   given \emph{additional} emphasis.
\end{em}

It is sometimes necessary to prevent \LaTeX\ from
breaking a line where it might otherwise do so.
This may be at a space, as between the ``Mr.''\ and
``Jones'' in
       ``Mr.~Jones'',        % ~ produces an unbreakable interword space.
or within a word---especially when the word is a
symbol like
       \mbox{\emph{itemnum}}
that makes little sense when hyphenated across
lines.

Footnotes\footnote{This is an example of a footnote.}
pose no problem.

\LaTeX\ is good at typesetting mathematical formulas
like
       \( x-3y + z = 7 \)
or
       \( a_{1} > x^{2n} + y^{2n} > x' \)
or
       \( \ip{A}{B} = \sum_{i} a_{i} b_{i} \).
The spaces you type in a formula are
ignored.  Remember that a letter like
       $x$                   % $ ... $  and  \( ... \)  are equivalent
is a formula when it denotes a mathematical
symbol, and it should be typed as one.

\section{Displayed Text}

Text is displayed by indenting it from the left
margin.  Quotations are commonly displayed.  There
are short quotations
\begin{quote}
   This is a short quotation.  It consists of a
   single paragraph of text.  See how it is formatted.
\end{quote}
and longer ones.
\begin{quotation}
   This is a longer quotation.  It consists of two
   paragraphs of text, neither of which are
   particularly interesting.

   This is the second paragraph of the quotation.  It
   is just as dull as the first paragraph.
\end{quotation}
Another frequently-displayed structure is a list.
The following is an example of an \emph{itemized}
list.
\begin{itemize}
   \item This is the first item of an itemized list.
         Each item in the list is marked with a ``tick''.
         You don't have to worry about what kind of tick
         mark is used.

   \item This is the second item of the list.  It
         contains another list nested inside it.  The inner
         list is an \emph{enumerated} list.
         \begin{enumerate}
            \item This is the first item of an enumerated
                  list that is nested within the itemized list.

            \item This is the second item of the inner list.
                  \LaTeX\ allows you to nest lists deeper than
                  you really should.
         \end{enumerate}
         This is the rest of the second item of the outer
         list.  It is no more interesting than any other
         part of the item.
   \item This is the third item of the list.
\end{itemize}
You can even display poetry.
\begin{verse}
   There is an environment
    for verse \\             % The \\ command separates lines
   Whose features some poets % within a stanza.
   will curse.

                             % One or more blank lines separate stanzas.

   For instead of making\\
   Them do \emph{all} line breaking, \\
   It allows them to put too many words on a line when they'd rather be
   forced to be terse.
\end{verse}

Mathematical formulas may also be displayed.  A
displayed formula
is
one-line long; multiline
formulas require special formatting instructions.
   \[  \ip{\Gamma}{\psi'} = x'' + y^{2} + z_{i}^{n}\]
Don't start a paragraph with a displayed equation,
nor make one a paragraph by itself.

\end{document}               % End of document.

\section{Bench Huge Section 1619}
% synthetic bench section 1619
% This is a sample LaTeX input file.  (Version of 12 August 2004.)
%
% A '%' character causes TeX to ignore all remaining text on the line,
% and is used for comments like this one.

\documentclass{article}      % Specifies the document class

                             % The preamble begins here.
\title{An Example Document}  % Declares the document's title.
\author{Leslie Lamport}      % Declares the author's name.
\date{January 21, 1994}      % Deleting this command produces today's date.

\newcommand{\ip}[2]{(#1, #2)}
                             % Defines \ip{arg1}{arg2} to mean
                             % (arg1, arg2).

%\newcommand{\ip}[2]{\langle #1 | #2\rangle}
                             % This is an alternative definition of
                             % \ip that is commented out.

\begin{document}             % End of preamble and beginning of text.

\maketitle                   % Produces the title.

This is an example input file.  Comparing it with
the output it generates can show you how to
produce a simple document of your own.

\section{Ordinary Text}      % Produces section heading.  Lower-level
                             % sections are begun with similar
                             % \subsection and \subsubsection commands.

The ends  of words and sentences are marked
  by   spaces. It  doesn't matter how many
spaces    you type; one is as good as 100.  The
end of   a line counts as a space.

One   or more   blank lines denote the  end
of  a paragraph.

Since any number of consecutive spaces are treated
like a single one, the formatting of the input
file makes no difference to
      \LaTeX,                % The \LaTeX command generates the LaTeX logo.
but it makes a difference to you.  When you use
\LaTeX, making your input file as easy to read
as possible will be a great help as you write
your document and when you change it.  This sample
file shows how you can add comments to your own input
file.

Because printing is different from typewriting,
there are a number of things that you have to do
differently when preparing an input file than if
you were just typing the document directly.
Quotation marks like
       ``this''
have to be handled specially, as do quotes within
quotes:
       ``\,`this'            % \, separates the double and single quote.
        is what I just
        wrote, not  `that'\,''.

Dashes come in three sizes: an
       intra-word
dash, a medium dash for number ranges like
       1--2,
and a punctuation
       dash---like
this.

A sentence-ending space should be larger than the
space between words within a sentence.  You
sometimes have to type special commands in
conjunction with punctuation characters to get
this right, as in the following sentence.
       Gnats, gnus, etc.\ all  % `\ ' makes an inter-word space.
       begin with G\@.         % \@ marks end-of-sentence punctuation.
You should check the spaces after periods when
reading your output to make sure you haven't
forgotten any special cases.  Generating an
ellipsis
       \ldots\               % `\ ' is needed after `\ldots' because TeX
                             % ignores spaces after command names like \ldots
                             % made from \ + letters.
                             %
                             % Note how a `%' character causes TeX to ignore
                             % the end of the input line, so these blank lines
                             % do not start a new paragraph.
                             %
with the right spacing around the periods requires
a special command.

\LaTeX\ interprets some common characters as
commands, so you must type special commands to
generate them.  These characters include the
following:
       \$ \& \% \# \{ and \}.

In printing, text is usually emphasized with an
       \emph{italic}
type style.

\begin{em}
   A long segment of text can also be emphasized
   in this way.  Text within such a segment can be
   given \emph{additional} emphasis.
\end{em}

It is sometimes necessary to prevent \LaTeX\ from
breaking a line where it might otherwise do so.
This may be at a space, as between the ``Mr.''\ and
``Jones'' in
       ``Mr.~Jones'',        % ~ produces an unbreakable interword space.
or within a word---especially when the word is a
symbol like
       \mbox{\emph{itemnum}}
that makes little sense when hyphenated across
lines.

Footnotes\footnote{This is an example of a footnote.}
pose no problem.

\LaTeX\ is good at typesetting mathematical formulas
like
       \( x-3y + z = 7 \)
or
       \( a_{1} > x^{2n} + y^{2n} > x' \)
or
       \( \ip{A}{B} = \sum_{i} a_{i} b_{i} \).
The spaces you type in a formula are
ignored.  Remember that a letter like
       $x$                   % $ ... $  and  \( ... \)  are equivalent
is a formula when it denotes a mathematical
symbol, and it should be typed as one.

\section{Displayed Text}

Text is displayed by indenting it from the left
margin.  Quotations are commonly displayed.  There
are short quotations
\begin{quote}
   This is a short quotation.  It consists of a
   single paragraph of text.  See how it is formatted.
\end{quote}
and longer ones.
\begin{quotation}
   This is a longer quotation.  It consists of two
   paragraphs of text, neither of which are
   particularly interesting.

   This is the second paragraph of the quotation.  It
   is just as dull as the first paragraph.
\end{quotation}
Another frequently-displayed structure is a list.
The following is an example of an \emph{itemized}
list.
\begin{itemize}
   \item This is the first item of an itemized list.
         Each item in the list is marked with a ``tick''.
         You don't have to worry about what kind of tick
         mark is used.

   \item This is the second item of the list.  It
         contains another list nested inside it.  The inner
         list is an \emph{enumerated} list.
         \begin{enumerate}
            \item This is the first item of an enumerated
                  list that is nested within the itemized list.

            \item This is the second item of the inner list.
                  \LaTeX\ allows you to nest lists deeper than
                  you really should.
         \end{enumerate}
         This is the rest of the second item of the outer
         list.  It is no more interesting than any other
         part of the item.
   \item This is the third item of the list.
\end{itemize}
You can even display poetry.
\begin{verse}
   There is an environment
    for verse \\             % The \\ command separates lines
   Whose features some poets % within a stanza.
   will curse.

                             % One or more blank lines separate stanzas.

   For instead of making\\
   Them do \emph{all} line breaking, \\
   It allows them to put too many words on a line when they'd rather be
   forced to be terse.
\end{verse}

Mathematical formulas may also be displayed.  A
displayed formula
is
one-line long; multiline
formulas require special formatting instructions.
   \[  \ip{\Gamma}{\psi'} = x'' + y^{2} + z_{i}^{n}\]
Don't start a paragraph with a displayed equation,
nor make one a paragraph by itself.

\end{document}               % End of document.

\section{Bench Huge Section 1620}
% synthetic bench section 1620
% This is a sample LaTeX input file.  (Version of 12 August 2004.)
%
% A '%' character causes TeX to ignore all remaining text on the line,
% and is used for comments like this one.

\documentclass{article}      % Specifies the document class

                             % The preamble begins here.
\title{An Example Document}  % Declares the document's title.
\author{Leslie Lamport}      % Declares the author's name.
\date{January 21, 1994}      % Deleting this command produces today's date.

\newcommand{\ip}[2]{(#1, #2)}
                             % Defines \ip{arg1}{arg2} to mean
                             % (arg1, arg2).

%\newcommand{\ip}[2]{\langle #1 | #2\rangle}
                             % This is an alternative definition of
                             % \ip that is commented out.

\begin{document}             % End of preamble and beginning of text.

\maketitle                   % Produces the title.

This is an example input file.  Comparing it with
the output it generates can show you how to
produce a simple document of your own.

\section{Ordinary Text}      % Produces section heading.  Lower-level
                             % sections are begun with similar
                             % \subsection and \subsubsection commands.

The ends  of words and sentences are marked
  by   spaces. It  doesn't matter how many
spaces    you type; one is as good as 100.  The
end of   a line counts as a space.

One   or more   blank lines denote the  end
of  a paragraph.

Since any number of consecutive spaces are treated
like a single one, the formatting of the input
file makes no difference to
      \LaTeX,                % The \LaTeX command generates the LaTeX logo.
but it makes a difference to you.  When you use
\LaTeX, making your input file as easy to read
as possible will be a great help as you write
your document and when you change it.  This sample
file shows how you can add comments to your own input
file.

Because printing is different from typewriting,
there are a number of things that you have to do
differently when preparing an input file than if
you were just typing the document directly.
Quotation marks like
       ``this''
have to be handled specially, as do quotes within
quotes:
       ``\,`this'            % \, separates the double and single quote.
        is what I just
        wrote, not  `that'\,''.

Dashes come in three sizes: an
       intra-word
dash, a medium dash for number ranges like
       1--2,
and a punctuation
       dash---like
this.

A sentence-ending space should be larger than the
space between words within a sentence.  You
sometimes have to type special commands in
conjunction with punctuation characters to get
this right, as in the following sentence.
       Gnats, gnus, etc.\ all  % `\ ' makes an inter-word space.
       begin with G\@.         % \@ marks end-of-sentence punctuation.
You should check the spaces after periods when
reading your output to make sure you haven't
forgotten any special cases.  Generating an
ellipsis
       \ldots\               % `\ ' is needed after `\ldots' because TeX
                             % ignores spaces after command names like \ldots
                             % made from \ + letters.
                             %
                             % Note how a `%' character causes TeX to ignore
                             % the end of the input line, so these blank lines
                             % do not start a new paragraph.
                             %
with the right spacing around the periods requires
a special command.

\LaTeX\ interprets some common characters as
commands, so you must type special commands to
generate them.  These characters include the
following:
       \$ \& \% \# \{ and \}.

In printing, text is usually emphasized with an
       \emph{italic}
type style.

\begin{em}
   A long segment of text can also be emphasized
   in this way.  Text within such a segment can be
   given \emph{additional} emphasis.
\end{em}

It is sometimes necessary to prevent \LaTeX\ from
breaking a line where it might otherwise do so.
This may be at a space, as between the ``Mr.''\ and
``Jones'' in
       ``Mr.~Jones'',        % ~ produces an unbreakable interword space.
or within a word---especially when the word is a
symbol like
       \mbox{\emph{itemnum}}
that makes little sense when hyphenated across
lines.

Footnotes\footnote{This is an example of a footnote.}
pose no problem.

\LaTeX\ is good at typesetting mathematical formulas
like
       \( x-3y + z = 7 \)
or
       \( a_{1} > x^{2n} + y^{2n} > x' \)
or
       \( \ip{A}{B} = \sum_{i} a_{i} b_{i} \).
The spaces you type in a formula are
ignored.  Remember that a letter like
       $x$                   % $ ... $  and  \( ... \)  are equivalent
is a formula when it denotes a mathematical
symbol, and it should be typed as one.

\section{Displayed Text}

Text is displayed by indenting it from the left
margin.  Quotations are commonly displayed.  There
are short quotations
\begin{quote}
   This is a short quotation.  It consists of a
   single paragraph of text.  See how it is formatted.
\end{quote}
and longer ones.
\begin{quotation}
   This is a longer quotation.  It consists of two
   paragraphs of text, neither of which are
   particularly interesting.

   This is the second paragraph of the quotation.  It
   is just as dull as the first paragraph.
\end{quotation}
Another frequently-displayed structure is a list.
The following is an example of an \emph{itemized}
list.
\begin{itemize}
   \item This is the first item of an itemized list.
         Each item in the list is marked with a ``tick''.
         You don't have to worry about what kind of tick
         mark is used.

   \item This is the second item of the list.  It
         contains another list nested inside it.  The inner
         list is an \emph{enumerated} list.
         \begin{enumerate}
            \item This is the first item of an enumerated
                  list that is nested within the itemized list.

            \item This is the second item of the inner list.
                  \LaTeX\ allows you to nest lists deeper than
                  you really should.
         \end{enumerate}
         This is the rest of the second item of the outer
         list.  It is no more interesting than any other
         part of the item.
   \item This is the third item of the list.
\end{itemize}
You can even display poetry.
\begin{verse}
   There is an environment
    for verse \\             % The \\ command separates lines
   Whose features some poets % within a stanza.
   will curse.

                             % One or more blank lines separate stanzas.

   For instead of making\\
   Them do \emph{all} line breaking, \\
   It allows them to put too many words on a line when they'd rather be
   forced to be terse.
\end{verse}

Mathematical formulas may also be displayed.  A
displayed formula
is
one-line long; multiline
formulas require special formatting instructions.
   \[  \ip{\Gamma}{\psi'} = x'' + y^{2} + z_{i}^{n}\]
Don't start a paragraph with a displayed equation,
nor make one a paragraph by itself.

\end{document}               % End of document.

\section{Bench Huge Section 1621}
% synthetic bench section 1621
% This is a sample LaTeX input file.  (Version of 12 August 2004.)
%
% A '%' character causes TeX to ignore all remaining text on the line,
% and is used for comments like this one.

\documentclass{article}      % Specifies the document class

                             % The preamble begins here.
\title{An Example Document}  % Declares the document's title.
\author{Leslie Lamport}      % Declares the author's name.
\date{January 21, 1994}      % Deleting this command produces today's date.

\newcommand{\ip}[2]{(#1, #2)}
                             % Defines \ip{arg1}{arg2} to mean
                             % (arg1, arg2).

%\newcommand{\ip}[2]{\langle #1 | #2\rangle}
                             % This is an alternative definition of
                             % \ip that is commented out.

\begin{document}             % End of preamble and beginning of text.

\maketitle                   % Produces the title.

This is an example input file.  Comparing it with
the output it generates can show you how to
produce a simple document of your own.

\section{Ordinary Text}      % Produces section heading.  Lower-level
                             % sections are begun with similar
                             % \subsection and \subsubsection commands.

The ends  of words and sentences are marked
  by   spaces. It  doesn't matter how many
spaces    you type; one is as good as 100.  The
end of   a line counts as a space.

One   or more   blank lines denote the  end
of  a paragraph.

Since any number of consecutive spaces are treated
like a single one, the formatting of the input
file makes no difference to
      \LaTeX,                % The \LaTeX command generates the LaTeX logo.
but it makes a difference to you.  When you use
\LaTeX, making your input file as easy to read
as possible will be a great help as you write
your document and when you change it.  This sample
file shows how you can add comments to your own input
file.

Because printing is different from typewriting,
there are a number of things that you have to do
differently when preparing an input file than if
you were just typing the document directly.
Quotation marks like
       ``this''
have to be handled specially, as do quotes within
quotes:
       ``\,`this'            % \, separates the double and single quote.
        is what I just
        wrote, not  `that'\,''.

Dashes come in three sizes: an
       intra-word
dash, a medium dash for number ranges like
       1--2,
and a punctuation
       dash---like
this.

A sentence-ending space should be larger than the
space between words within a sentence.  You
sometimes have to type special commands in
conjunction with punctuation characters to get
this right, as in the following sentence.
       Gnats, gnus, etc.\ all  % `\ ' makes an inter-word space.
       begin with G\@.         % \@ marks end-of-sentence punctuation.
You should check the spaces after periods when
reading your output to make sure you haven't
forgotten any special cases.  Generating an
ellipsis
       \ldots\               % `\ ' is needed after `\ldots' because TeX
                             % ignores spaces after command names like \ldots
                             % made from \ + letters.
                             %
                             % Note how a `%' character causes TeX to ignore
                             % the end of the input line, so these blank lines
                             % do not start a new paragraph.
                             %
with the right spacing around the periods requires
a special command.

\LaTeX\ interprets some common characters as
commands, so you must type special commands to
generate them.  These characters include the
following:
       \$ \& \% \# \{ and \}.

In printing, text is usually emphasized with an
       \emph{italic}
type style.

\begin{em}
   A long segment of text can also be emphasized
   in this way.  Text within such a segment can be
   given \emph{additional} emphasis.
\end{em}

It is sometimes necessary to prevent \LaTeX\ from
breaking a line where it might otherwise do so.
This may be at a space, as between the ``Mr.''\ and
``Jones'' in
       ``Mr.~Jones'',        % ~ produces an unbreakable interword space.
or within a word---especially when the word is a
symbol like
       \mbox{\emph{itemnum}}
that makes little sense when hyphenated across
lines.

Footnotes\footnote{This is an example of a footnote.}
pose no problem.

\LaTeX\ is good at typesetting mathematical formulas
like
       \( x-3y + z = 7 \)
or
       \( a_{1} > x^{2n} + y^{2n} > x' \)
or
       \( \ip{A}{B} = \sum_{i} a_{i} b_{i} \).
The spaces you type in a formula are
ignored.  Remember that a letter like
       $x$                   % $ ... $  and  \( ... \)  are equivalent
is a formula when it denotes a mathematical
symbol, and it should be typed as one.

\section{Displayed Text}

Text is displayed by indenting it from the left
margin.  Quotations are commonly displayed.  There
are short quotations
\begin{quote}
   This is a short quotation.  It consists of a
   single paragraph of text.  See how it is formatted.
\end{quote}
and longer ones.
\begin{quotation}
   This is a longer quotation.  It consists of two
   paragraphs of text, neither of which are
   particularly interesting.

   This is the second paragraph of the quotation.  It
   is just as dull as the first paragraph.
\end{quotation}
Another frequently-displayed structure is a list.
The following is an example of an \emph{itemized}
list.
\begin{itemize}
   \item This is the first item of an itemized list.
         Each item in the list is marked with a ``tick''.
         You don't have to worry about what kind of tick
         mark is used.

   \item This is the second item of the list.  It
         contains another list nested inside it.  The inner
         list is an \emph{enumerated} list.
         \begin{enumerate}
            \item This is the first item of an enumerated
                  list that is nested within the itemized list.

            \item This is the second item of the inner list.
                  \LaTeX\ allows you to nest lists deeper than
                  you really should.
         \end{enumerate}
         This is the rest of the second item of the outer
         list.  It is no more interesting than any other
         part of the item.
   \item This is the third item of the list.
\end{itemize}
You can even display poetry.
\begin{verse}
   There is an environment
    for verse \\             % The \\ command separates lines
   Whose features some poets % within a stanza.
   will curse.

                             % One or more blank lines separate stanzas.

   For instead of making\\
   Them do \emph{all} line breaking, \\
   It allows them to put too many words on a line when they'd rather be
   forced to be terse.
\end{verse}

Mathematical formulas may also be displayed.  A
displayed formula
is
one-line long; multiline
formulas require special formatting instructions.
   \[  \ip{\Gamma}{\psi'} = x'' + y^{2} + z_{i}^{n}\]
Don't start a paragraph with a displayed equation,
nor make one a paragraph by itself.

\end{document}               % End of document.

\section{Bench Huge Section 1622}
% synthetic bench section 1622
% This is a sample LaTeX input file.  (Version of 12 August 2004.)
%
% A '%' character causes TeX to ignore all remaining text on the line,
% and is used for comments like this one.

\documentclass{article}      % Specifies the document class

                             % The preamble begins here.
\title{An Example Document}  % Declares the document's title.
\author{Leslie Lamport}      % Declares the author's name.
\date{January 21, 1994}      % Deleting this command produces today's date.

\newcommand{\ip}[2]{(#1, #2)}
                             % Defines \ip{arg1}{arg2} to mean
                             % (arg1, arg2).

%\newcommand{\ip}[2]{\langle #1 | #2\rangle}
                             % This is an alternative definition of
                             % \ip that is commented out.

\begin{document}             % End of preamble and beginning of text.

\maketitle                   % Produces the title.

This is an example input file.  Comparing it with
the output it generates can show you how to
produce a simple document of your own.

\section{Ordinary Text}      % Produces section heading.  Lower-level
                             % sections are begun with similar
                             % \subsection and \subsubsection commands.

The ends  of words and sentences are marked
  by   spaces. It  doesn't matter how many
spaces    you type; one is as good as 100.  The
end of   a line counts as a space.

One   or more   blank lines denote the  end
of  a paragraph.

Since any number of consecutive spaces are treated
like a single one, the formatting of the input
file makes no difference to
      \LaTeX,                % The \LaTeX command generates the LaTeX logo.
but it makes a difference to you.  When you use
\LaTeX, making your input file as easy to read
as possible will be a great help as you write
your document and when you change it.  This sample
file shows how you can add comments to your own input
file.

Because printing is different from typewriting,
there are a number of things that you have to do
differently when preparing an input file than if
you were just typing the document directly.
Quotation marks like
       ``this''
have to be handled specially, as do quotes within
quotes:
       ``\,`this'            % \, separates the double and single quote.
        is what I just
        wrote, not  `that'\,''.

Dashes come in three sizes: an
       intra-word
dash, a medium dash for number ranges like
       1--2,
and a punctuation
       dash---like
this.

A sentence-ending space should be larger than the
space between words within a sentence.  You
sometimes have to type special commands in
conjunction with punctuation characters to get
this right, as in the following sentence.
       Gnats, gnus, etc.\ all  % `\ ' makes an inter-word space.
       begin with G\@.         % \@ marks end-of-sentence punctuation.
You should check the spaces after periods when
reading your output to make sure you haven't
forgotten any special cases.  Generating an
ellipsis
       \ldots\               % `\ ' is needed after `\ldots' because TeX
                             % ignores spaces after command names like \ldots
                             % made from \ + letters.
                             %
                             % Note how a `%' character causes TeX to ignore
                             % the end of the input line, so these blank lines
                             % do not start a new paragraph.
                             %
with the right spacing around the periods requires
a special command.

\LaTeX\ interprets some common characters as
commands, so you must type special commands to
generate them.  These characters include the
following:
       \$ \& \% \# \{ and \}.

In printing, text is usually emphasized with an
       \emph{italic}
type style.

\begin{em}
   A long segment of text can also be emphasized
   in this way.  Text within such a segment can be
   given \emph{additional} emphasis.
\end{em}

It is sometimes necessary to prevent \LaTeX\ from
breaking a line where it might otherwise do so.
This may be at a space, as between the ``Mr.''\ and
``Jones'' in
       ``Mr.~Jones'',        % ~ produces an unbreakable interword space.
or within a word---especially when the word is a
symbol like
       \mbox{\emph{itemnum}}
that makes little sense when hyphenated across
lines.

Footnotes\footnote{This is an example of a footnote.}
pose no problem.

\LaTeX\ is good at typesetting mathematical formulas
like
       \( x-3y + z = 7 \)
or
       \( a_{1} > x^{2n} + y^{2n} > x' \)
or
       \( \ip{A}{B} = \sum_{i} a_{i} b_{i} \).
The spaces you type in a formula are
ignored.  Remember that a letter like
       $x$                   % $ ... $  and  \( ... \)  are equivalent
is a formula when it denotes a mathematical
symbol, and it should be typed as one.

\section{Displayed Text}

Text is displayed by indenting it from the left
margin.  Quotations are commonly displayed.  There
are short quotations
\begin{quote}
   This is a short quotation.  It consists of a
   single paragraph of text.  See how it is formatted.
\end{quote}
and longer ones.
\begin{quotation}
   This is a longer quotation.  It consists of two
   paragraphs of text, neither of which are
   particularly interesting.

   This is the second paragraph of the quotation.  It
   is just as dull as the first paragraph.
\end{quotation}
Another frequently-displayed structure is a list.
The following is an example of an \emph{itemized}
list.
\begin{itemize}
   \item This is the first item of an itemized list.
         Each item in the list is marked with a ``tick''.
         You don't have to worry about what kind of tick
         mark is used.

   \item This is the second item of the list.  It
         contains another list nested inside it.  The inner
         list is an \emph{enumerated} list.
         \begin{enumerate}
            \item This is the first item of an enumerated
                  list that is nested within the itemized list.

            \item This is the second item of the inner list.
                  \LaTeX\ allows you to nest lists deeper than
                  you really should.
         \end{enumerate}
         This is the rest of the second item of the outer
         list.  It is no more interesting than any other
         part of the item.
   \item This is the third item of the list.
\end{itemize}
You can even display poetry.
\begin{verse}
   There is an environment
    for verse \\             % The \\ command separates lines
   Whose features some poets % within a stanza.
   will curse.

                             % One or more blank lines separate stanzas.

   For instead of making\\
   Them do \emph{all} line breaking, \\
   It allows them to put too many words on a line when they'd rather be
   forced to be terse.
\end{verse}

Mathematical formulas may also be displayed.  A
displayed formula
is
one-line long; multiline
formulas require special formatting instructions.
   \[  \ip{\Gamma}{\psi'} = x'' + y^{2} + z_{i}^{n}\]
Don't start a paragraph with a displayed equation,
nor make one a paragraph by itself.

\end{document}               % End of document.

\section{Bench Huge Section 1623}
% synthetic bench section 1623
% This is a sample LaTeX input file.  (Version of 12 August 2004.)
%
% A '%' character causes TeX to ignore all remaining text on the line,
% and is used for comments like this one.

\documentclass{article}      % Specifies the document class

                             % The preamble begins here.
\title{An Example Document}  % Declares the document's title.
\author{Leslie Lamport}      % Declares the author's name.
\date{January 21, 1994}      % Deleting this command produces today's date.

\newcommand{\ip}[2]{(#1, #2)}
                             % Defines \ip{arg1}{arg2} to mean
                             % (arg1, arg2).

%\newcommand{\ip}[2]{\langle #1 | #2\rangle}
                             % This is an alternative definition of
                             % \ip that is commented out.

\begin{document}             % End of preamble and beginning of text.

\maketitle                   % Produces the title.

This is an example input file.  Comparing it with
the output it generates can show you how to
produce a simple document of your own.

\section{Ordinary Text}      % Produces section heading.  Lower-level
                             % sections are begun with similar
                             % \subsection and \subsubsection commands.

The ends  of words and sentences are marked
  by   spaces. It  doesn't matter how many
spaces    you type; one is as good as 100.  The
end of   a line counts as a space.

One   or more   blank lines denote the  end
of  a paragraph.

Since any number of consecutive spaces are treated
like a single one, the formatting of the input
file makes no difference to
      \LaTeX,                % The \LaTeX command generates the LaTeX logo.
but it makes a difference to you.  When you use
\LaTeX, making your input file as easy to read
as possible will be a great help as you write
your document and when you change it.  This sample
file shows how you can add comments to your own input
file.

Because printing is different from typewriting,
there are a number of things that you have to do
differently when preparing an input file than if
you were just typing the document directly.
Quotation marks like
       ``this''
have to be handled specially, as do quotes within
quotes:
       ``\,`this'            % \, separates the double and single quote.
        is what I just
        wrote, not  `that'\,''.

Dashes come in three sizes: an
       intra-word
dash, a medium dash for number ranges like
       1--2,
and a punctuation
       dash---like
this.

A sentence-ending space should be larger than the
space between words within a sentence.  You
sometimes have to type special commands in
conjunction with punctuation characters to get
this right, as in the following sentence.
       Gnats, gnus, etc.\ all  % `\ ' makes an inter-word space.
       begin with G\@.         % \@ marks end-of-sentence punctuation.
You should check the spaces after periods when
reading your output to make sure you haven't
forgotten any special cases.  Generating an
ellipsis
       \ldots\               % `\ ' is needed after `\ldots' because TeX
                             % ignores spaces after command names like \ldots
                             % made from \ + letters.
                             %
                             % Note how a `%' character causes TeX to ignore
                             % the end of the input line, so these blank lines
                             % do not start a new paragraph.
                             %
with the right spacing around the periods requires
a special command.

\LaTeX\ interprets some common characters as
commands, so you must type special commands to
generate them.  These characters include the
following:
       \$ \& \% \# \{ and \}.

In printing, text is usually emphasized with an
       \emph{italic}
type style.

\begin{em}
   A long segment of text can also be emphasized
   in this way.  Text within such a segment can be
   given \emph{additional} emphasis.
\end{em}

It is sometimes necessary to prevent \LaTeX\ from
breaking a line where it might otherwise do so.
This may be at a space, as between the ``Mr.''\ and
``Jones'' in
       ``Mr.~Jones'',        % ~ produces an unbreakable interword space.
or within a word---especially when the word is a
symbol like
       \mbox{\emph{itemnum}}
that makes little sense when hyphenated across
lines.

Footnotes\footnote{This is an example of a footnote.}
pose no problem.

\LaTeX\ is good at typesetting mathematical formulas
like
       \( x-3y + z = 7 \)
or
       \( a_{1} > x^{2n} + y^{2n} > x' \)
or
       \( \ip{A}{B} = \sum_{i} a_{i} b_{i} \).
The spaces you type in a formula are
ignored.  Remember that a letter like
       $x$                   % $ ... $  and  \( ... \)  are equivalent
is a formula when it denotes a mathematical
symbol, and it should be typed as one.

\section{Displayed Text}

Text is displayed by indenting it from the left
margin.  Quotations are commonly displayed.  There
are short quotations
\begin{quote}
   This is a short quotation.  It consists of a
   single paragraph of text.  See how it is formatted.
\end{quote}
and longer ones.
\begin{quotation}
   This is a longer quotation.  It consists of two
   paragraphs of text, neither of which are
   particularly interesting.

   This is the second paragraph of the quotation.  It
   is just as dull as the first paragraph.
\end{quotation}
Another frequently-displayed structure is a list.
The following is an example of an \emph{itemized}
list.
\begin{itemize}
   \item This is the first item of an itemized list.
         Each item in the list is marked with a ``tick''.
         You don't have to worry about what kind of tick
         mark is used.

   \item This is the second item of the list.  It
         contains another list nested inside it.  The inner
         list is an \emph{enumerated} list.
         \begin{enumerate}
            \item This is the first item of an enumerated
                  list that is nested within the itemized list.

            \item This is the second item of the inner list.
                  \LaTeX\ allows you to nest lists deeper than
                  you really should.
         \end{enumerate}
         This is the rest of the second item of the outer
         list.  It is no more interesting than any other
         part of the item.
   \item This is the third item of the list.
\end{itemize}
You can even display poetry.
\begin{verse}
   There is an environment
    for verse \\             % The \\ command separates lines
   Whose features some poets % within a stanza.
   will curse.

                             % One or more blank lines separate stanzas.

   For instead of making\\
   Them do \emph{all} line breaking, \\
   It allows them to put too many words on a line when they'd rather be
   forced to be terse.
\end{verse}

Mathematical formulas may also be displayed.  A
displayed formula
is
one-line long; multiline
formulas require special formatting instructions.
   \[  \ip{\Gamma}{\psi'} = x'' + y^{2} + z_{i}^{n}\]
Don't start a paragraph with a displayed equation,
nor make one a paragraph by itself.

\end{document}               % End of document.

\section{Bench Huge Section 1624}
% synthetic bench section 1624
% This is a sample LaTeX input file.  (Version of 12 August 2004.)
%
% A '%' character causes TeX to ignore all remaining text on the line,
% and is used for comments like this one.

\documentclass{article}      % Specifies the document class

                             % The preamble begins here.
\title{An Example Document}  % Declares the document's title.
\author{Leslie Lamport}      % Declares the author's name.
\date{January 21, 1994}      % Deleting this command produces today's date.

\newcommand{\ip}[2]{(#1, #2)}
                             % Defines \ip{arg1}{arg2} to mean
                             % (arg1, arg2).

%\newcommand{\ip}[2]{\langle #1 | #2\rangle}
                             % This is an alternative definition of
                             % \ip that is commented out.

\begin{document}             % End of preamble and beginning of text.

\maketitle                   % Produces the title.

This is an example input file.  Comparing it with
the output it generates can show you how to
produce a simple document of your own.

\section{Ordinary Text}      % Produces section heading.  Lower-level
                             % sections are begun with similar
                             % \subsection and \subsubsection commands.

The ends  of words and sentences are marked
  by   spaces. It  doesn't matter how many
spaces    you type; one is as good as 100.  The
end of   a line counts as a space.

One   or more   blank lines denote the  end
of  a paragraph.

Since any number of consecutive spaces are treated
like a single one, the formatting of the input
file makes no difference to
      \LaTeX,                % The \LaTeX command generates the LaTeX logo.
but it makes a difference to you.  When you use
\LaTeX, making your input file as easy to read
as possible will be a great help as you write
your document and when you change it.  This sample
file shows how you can add comments to your own input
file.

Because printing is different from typewriting,
there are a number of things that you have to do
differently when preparing an input file than if
you were just typing the document directly.
Quotation marks like
       ``this''
have to be handled specially, as do quotes within
quotes:
       ``\,`this'            % \, separates the double and single quote.
        is what I just
        wrote, not  `that'\,''.

Dashes come in three sizes: an
       intra-word
dash, a medium dash for number ranges like
       1--2,
and a punctuation
       dash---like
this.

A sentence-ending space should be larger than the
space between words within a sentence.  You
sometimes have to type special commands in
conjunction with punctuation characters to get
this right, as in the following sentence.
       Gnats, gnus, etc.\ all  % `\ ' makes an inter-word space.
       begin with G\@.         % \@ marks end-of-sentence punctuation.
You should check the spaces after periods when
reading your output to make sure you haven't
forgotten any special cases.  Generating an
ellipsis
       \ldots\               % `\ ' is needed after `\ldots' because TeX
                             % ignores spaces after command names like \ldots
                             % made from \ + letters.
                             %
                             % Note how a `%' character causes TeX to ignore
                             % the end of the input line, so these blank lines
                             % do not start a new paragraph.
                             %
with the right spacing around the periods requires
a special command.

\LaTeX\ interprets some common characters as
commands, so you must type special commands to
generate them.  These characters include the
following:
       \$ \& \% \# \{ and \}.

In printing, text is usually emphasized with an
       \emph{italic}
type style.

\begin{em}
   A long segment of text can also be emphasized
   in this way.  Text within such a segment can be
   given \emph{additional} emphasis.
\end{em}

It is sometimes necessary to prevent \LaTeX\ from
breaking a line where it might otherwise do so.
This may be at a space, as between the ``Mr.''\ and
``Jones'' in
       ``Mr.~Jones'',        % ~ produces an unbreakable interword space.
or within a word---especially when the word is a
symbol like
       \mbox{\emph{itemnum}}
that makes little sense when hyphenated across
lines.

Footnotes\footnote{This is an example of a footnote.}
pose no problem.

\LaTeX\ is good at typesetting mathematical formulas
like
       \( x-3y + z = 7 \)
or
       \( a_{1} > x^{2n} + y^{2n} > x' \)
or
       \( \ip{A}{B} = \sum_{i} a_{i} b_{i} \).
The spaces you type in a formula are
ignored.  Remember that a letter like
       $x$                   % $ ... $  and  \( ... \)  are equivalent
is a formula when it denotes a mathematical
symbol, and it should be typed as one.

\section{Displayed Text}

Text is displayed by indenting it from the left
margin.  Quotations are commonly displayed.  There
are short quotations
\begin{quote}
   This is a short quotation.  It consists of a
   single paragraph of text.  See how it is formatted.
\end{quote}
and longer ones.
\begin{quotation}
   This is a longer quotation.  It consists of two
   paragraphs of text, neither of which are
   particularly interesting.

   This is the second paragraph of the quotation.  It
   is just as dull as the first paragraph.
\end{quotation}
Another frequently-displayed structure is a list.
The following is an example of an \emph{itemized}
list.
\begin{itemize}
   \item This is the first item of an itemized list.
         Each item in the list is marked with a ``tick''.
         You don't have to worry about what kind of tick
         mark is used.

   \item This is the second item of the list.  It
         contains another list nested inside it.  The inner
         list is an \emph{enumerated} list.
         \begin{enumerate}
            \item This is the first item of an enumerated
                  list that is nested within the itemized list.

            \item This is the second item of the inner list.
                  \LaTeX\ allows you to nest lists deeper than
                  you really should.
         \end{enumerate}
         This is the rest of the second item of the outer
         list.  It is no more interesting than any other
         part of the item.
   \item This is the third item of the list.
\end{itemize}
You can even display poetry.
\begin{verse}
   There is an environment
    for verse \\             % The \\ command separates lines
   Whose features some poets % within a stanza.
   will curse.

                             % One or more blank lines separate stanzas.

   For instead of making\\
   Them do \emph{all} line breaking, \\
   It allows them to put too many words on a line when they'd rather be
   forced to be terse.
\end{verse}

Mathematical formulas may also be displayed.  A
displayed formula
is
one-line long; multiline
formulas require special formatting instructions.
   \[  \ip{\Gamma}{\psi'} = x'' + y^{2} + z_{i}^{n}\]
Don't start a paragraph with a displayed equation,
nor make one a paragraph by itself.

\end{document}               % End of document.

\section{Bench Huge Section 1625}
% synthetic bench section 1625
% This is a sample LaTeX input file.  (Version of 12 August 2004.)
%
% A '%' character causes TeX to ignore all remaining text on the line,
% and is used for comments like this one.

\documentclass{article}      % Specifies the document class

                             % The preamble begins here.
\title{An Example Document}  % Declares the document's title.
\author{Leslie Lamport}      % Declares the author's name.
\date{January 21, 1994}      % Deleting this command produces today's date.

\newcommand{\ip}[2]{(#1, #2)}
                             % Defines \ip{arg1}{arg2} to mean
                             % (arg1, arg2).

%\newcommand{\ip}[2]{\langle #1 | #2\rangle}
                             % This is an alternative definition of
                             % \ip that is commented out.

\begin{document}             % End of preamble and beginning of text.

\maketitle                   % Produces the title.

This is an example input file.  Comparing it with
the output it generates can show you how to
produce a simple document of your own.

\section{Ordinary Text}      % Produces section heading.  Lower-level
                             % sections are begun with similar
                             % \subsection and \subsubsection commands.

The ends  of words and sentences are marked
  by   spaces. It  doesn't matter how many
spaces    you type; one is as good as 100.  The
end of   a line counts as a space.

One   or more   blank lines denote the  end
of  a paragraph.

Since any number of consecutive spaces are treated
like a single one, the formatting of the input
file makes no difference to
      \LaTeX,                % The \LaTeX command generates the LaTeX logo.
but it makes a difference to you.  When you use
\LaTeX, making your input file as easy to read
as possible will be a great help as you write
your document and when you change it.  This sample
file shows how you can add comments to your own input
file.

Because printing is different from typewriting,
there are a number of things that you have to do
differently when preparing an input file than if
you were just typing the document directly.
Quotation marks like
       ``this''
have to be handled specially, as do quotes within
quotes:
       ``\,`this'            % \, separates the double and single quote.
        is what I just
        wrote, not  `that'\,''.

Dashes come in three sizes: an
       intra-word
dash, a medium dash for number ranges like
       1--2,
and a punctuation
       dash---like
this.

A sentence-ending space should be larger than the
space between words within a sentence.  You
sometimes have to type special commands in
conjunction with punctuation characters to get
this right, as in the following sentence.
       Gnats, gnus, etc.\ all  % `\ ' makes an inter-word space.
       begin with G\@.         % \@ marks end-of-sentence punctuation.
You should check the spaces after periods when
reading your output to make sure you haven't
forgotten any special cases.  Generating an
ellipsis
       \ldots\               % `\ ' is needed after `\ldots' because TeX
                             % ignores spaces after command names like \ldots
                             % made from \ + letters.
                             %
                             % Note how a `%' character causes TeX to ignore
                             % the end of the input line, so these blank lines
                             % do not start a new paragraph.
                             %
with the right spacing around the periods requires
a special command.

\LaTeX\ interprets some common characters as
commands, so you must type special commands to
generate them.  These characters include the
following:
       \$ \& \% \# \{ and \}.

In printing, text is usually emphasized with an
       \emph{italic}
type style.

\begin{em}
   A long segment of text can also be emphasized
   in this way.  Text within such a segment can be
   given \emph{additional} emphasis.
\end{em}

It is sometimes necessary to prevent \LaTeX\ from
breaking a line where it might otherwise do so.
This may be at a space, as between the ``Mr.''\ and
``Jones'' in
       ``Mr.~Jones'',        % ~ produces an unbreakable interword space.
or within a word---especially when the word is a
symbol like
       \mbox{\emph{itemnum}}
that makes little sense when hyphenated across
lines.

Footnotes\footnote{This is an example of a footnote.}
pose no problem.

\LaTeX\ is good at typesetting mathematical formulas
like
       \( x-3y + z = 7 \)
or
       \( a_{1} > x^{2n} + y^{2n} > x' \)
or
       \( \ip{A}{B} = \sum_{i} a_{i} b_{i} \).
The spaces you type in a formula are
ignored.  Remember that a letter like
       $x$                   % $ ... $  and  \( ... \)  are equivalent
is a formula when it denotes a mathematical
symbol, and it should be typed as one.

\section{Displayed Text}

Text is displayed by indenting it from the left
margin.  Quotations are commonly displayed.  There
are short quotations
\begin{quote}
   This is a short quotation.  It consists of a
   single paragraph of text.  See how it is formatted.
\end{quote}
and longer ones.
\begin{quotation}
   This is a longer quotation.  It consists of two
   paragraphs of text, neither of which are
   particularly interesting.

   This is the second paragraph of the quotation.  It
   is just as dull as the first paragraph.
\end{quotation}
Another frequently-displayed structure is a list.
The following is an example of an \emph{itemized}
list.
\begin{itemize}
   \item This is the first item of an itemized list.
         Each item in the list is marked with a ``tick''.
         You don't have to worry about what kind of tick
         mark is used.

   \item This is the second item of the list.  It
         contains another list nested inside it.  The inner
         list is an \emph{enumerated} list.
         \begin{enumerate}
            \item This is the first item of an enumerated
                  list that is nested within the itemized list.

            \item This is the second item of the inner list.
                  \LaTeX\ allows you to nest lists deeper than
                  you really should.
         \end{enumerate}
         This is the rest of the second item of the outer
         list.  It is no more interesting than any other
         part of the item.
   \item This is the third item of the list.
\end{itemize}
You can even display poetry.
\begin{verse}
   There is an environment
    for verse \\             % The \\ command separates lines
   Whose features some poets % within a stanza.
   will curse.

                             % One or more blank lines separate stanzas.

   For instead of making\\
   Them do \emph{all} line breaking, \\
   It allows them to put too many words on a line when they'd rather be
   forced to be terse.
\end{verse}

Mathematical formulas may also be displayed.  A
displayed formula
is
one-line long; multiline
formulas require special formatting instructions.
   \[  \ip{\Gamma}{\psi'} = x'' + y^{2} + z_{i}^{n}\]
Don't start a paragraph with a displayed equation,
nor make one a paragraph by itself.

\end{document}               % End of document.

\section{Bench Huge Section 1626}
% synthetic bench section 1626
% This is a sample LaTeX input file.  (Version of 12 August 2004.)
%
% A '%' character causes TeX to ignore all remaining text on the line,
% and is used for comments like this one.

\documentclass{article}      % Specifies the document class

                             % The preamble begins here.
\title{An Example Document}  % Declares the document's title.
\author{Leslie Lamport}      % Declares the author's name.
\date{January 21, 1994}      % Deleting this command produces today's date.

\newcommand{\ip}[2]{(#1, #2)}
                             % Defines \ip{arg1}{arg2} to mean
                             % (arg1, arg2).

%\newcommand{\ip}[2]{\langle #1 | #2\rangle}
                             % This is an alternative definition of
                             % \ip that is commented out.

\begin{document}             % End of preamble and beginning of text.

\maketitle                   % Produces the title.

This is an example input file.  Comparing it with
the output it generates can show you how to
produce a simple document of your own.

\section{Ordinary Text}      % Produces section heading.  Lower-level
                             % sections are begun with similar
                             % \subsection and \subsubsection commands.

The ends  of words and sentences are marked
  by   spaces. It  doesn't matter how many
spaces    you type; one is as good as 100.  The
end of   a line counts as a space.

One   or more   blank lines denote the  end
of  a paragraph.

Since any number of consecutive spaces are treated
like a single one, the formatting of the input
file makes no difference to
      \LaTeX,                % The \LaTeX command generates the LaTeX logo.
but it makes a difference to you.  When you use
\LaTeX, making your input file as easy to read
as possible will be a great help as you write
your document and when you change it.  This sample
file shows how you can add comments to your own input
file.

Because printing is different from typewriting,
there are a number of things that you have to do
differently when preparing an input file than if
you were just typing the document directly.
Quotation marks like
       ``this''
have to be handled specially, as do quotes within
quotes:
       ``\,`this'            % \, separates the double and single quote.
        is what I just
        wrote, not  `that'\,''.

Dashes come in three sizes: an
       intra-word
dash, a medium dash for number ranges like
       1--2,
and a punctuation
       dash---like
this.

A sentence-ending space should be larger than the
space between words within a sentence.  You
sometimes have to type special commands in
conjunction with punctuation characters to get
this right, as in the following sentence.
       Gnats, gnus, etc.\ all  % `\ ' makes an inter-word space.
       begin with G\@.         % \@ marks end-of-sentence punctuation.
You should check the spaces after periods when
reading your output to make sure you haven't
forgotten any special cases.  Generating an
ellipsis
       \ldots\               % `\ ' is needed after `\ldots' because TeX
                             % ignores spaces after command names like \ldots
                             % made from \ + letters.
                             %
                             % Note how a `%' character causes TeX to ignore
                             % the end of the input line, so these blank lines
                             % do not start a new paragraph.
                             %
with the right spacing around the periods requires
a special command.

\LaTeX\ interprets some common characters as
commands, so you must type special commands to
generate them.  These characters include the
following:
       \$ \& \% \# \{ and \}.

In printing, text is usually emphasized with an
       \emph{italic}
type style.

\begin{em}
   A long segment of text can also be emphasized
   in this way.  Text within such a segment can be
   given \emph{additional} emphasis.
\end{em}

It is sometimes necessary to prevent \LaTeX\ from
breaking a line where it might otherwise do so.
This may be at a space, as between the ``Mr.''\ and
``Jones'' in
       ``Mr.~Jones'',        % ~ produces an unbreakable interword space.
or within a word---especially when the word is a
symbol like
       \mbox{\emph{itemnum}}
that makes little sense when hyphenated across
lines.

Footnotes\footnote{This is an example of a footnote.}
pose no problem.

\LaTeX\ is good at typesetting mathematical formulas
like
       \( x-3y + z = 7 \)
or
       \( a_{1} > x^{2n} + y^{2n} > x' \)
or
       \( \ip{A}{B} = \sum_{i} a_{i} b_{i} \).
The spaces you type in a formula are
ignored.  Remember that a letter like
       $x$                   % $ ... $  and  \( ... \)  are equivalent
is a formula when it denotes a mathematical
symbol, and it should be typed as one.

\section{Displayed Text}

Text is displayed by indenting it from the left
margin.  Quotations are commonly displayed.  There
are short quotations
\begin{quote}
   This is a short quotation.  It consists of a
   single paragraph of text.  See how it is formatted.
\end{quote}
and longer ones.
\begin{quotation}
   This is a longer quotation.  It consists of two
   paragraphs of text, neither of which are
   particularly interesting.

   This is the second paragraph of the quotation.  It
   is just as dull as the first paragraph.
\end{quotation}
Another frequently-displayed structure is a list.
The following is an example of an \emph{itemized}
list.
\begin{itemize}
   \item This is the first item of an itemized list.
         Each item in the list is marked with a ``tick''.
         You don't have to worry about what kind of tick
         mark is used.

   \item This is the second item of the list.  It
         contains another list nested inside it.  The inner
         list is an \emph{enumerated} list.
         \begin{enumerate}
            \item This is the first item of an enumerated
                  list that is nested within the itemized list.

            \item This is the second item of the inner list.
                  \LaTeX\ allows you to nest lists deeper than
                  you really should.
         \end{enumerate}
         This is the rest of the second item of the outer
         list.  It is no more interesting than any other
         part of the item.
   \item This is the third item of the list.
\end{itemize}
You can even display poetry.
\begin{verse}
   There is an environment
    for verse \\             % The \\ command separates lines
   Whose features some poets % within a stanza.
   will curse.

                             % One or more blank lines separate stanzas.

   For instead of making\\
   Them do \emph{all} line breaking, \\
   It allows them to put too many words on a line when they'd rather be
   forced to be terse.
\end{verse}

Mathematical formulas may also be displayed.  A
displayed formula
is
one-line long; multiline
formulas require special formatting instructions.
   \[  \ip{\Gamma}{\psi'} = x'' + y^{2} + z_{i}^{n}\]
Don't start a paragraph with a displayed equation,
nor make one a paragraph by itself.

\end{document}               % End of document.

\section{Bench Huge Section 1627}
% synthetic bench section 1627
% This is a sample LaTeX input file.  (Version of 12 August 2004.)
%
% A '%' character causes TeX to ignore all remaining text on the line,
% and is used for comments like this one.

\documentclass{article}      % Specifies the document class

                             % The preamble begins here.
\title{An Example Document}  % Declares the document's title.
\author{Leslie Lamport}      % Declares the author's name.
\date{January 21, 1994}      % Deleting this command produces today's date.

\newcommand{\ip}[2]{(#1, #2)}
                             % Defines \ip{arg1}{arg2} to mean
                             % (arg1, arg2).

%\newcommand{\ip}[2]{\langle #1 | #2\rangle}
                             % This is an alternative definition of
                             % \ip that is commented out.

\begin{document}             % End of preamble and beginning of text.

\maketitle                   % Produces the title.

This is an example input file.  Comparing it with
the output it generates can show you how to
produce a simple document of your own.

\section{Ordinary Text}      % Produces section heading.  Lower-level
                             % sections are begun with similar
                             % \subsection and \subsubsection commands.

The ends  of words and sentences are marked
  by   spaces. It  doesn't matter how many
spaces    you type; one is as good as 100.  The
end of   a line counts as a space.

One   or more   blank lines denote the  end
of  a paragraph.

Since any number of consecutive spaces are treated
like a single one, the formatting of the input
file makes no difference to
      \LaTeX,                % The \LaTeX command generates the LaTeX logo.
but it makes a difference to you.  When you use
\LaTeX, making your input file as easy to read
as possible will be a great help as you write
your document and when you change it.  This sample
file shows how you can add comments to your own input
file.

Because printing is different from typewriting,
there are a number of things that you have to do
differently when preparing an input file than if
you were just typing the document directly.
Quotation marks like
       ``this''
have to be handled specially, as do quotes within
quotes:
       ``\,`this'            % \, separates the double and single quote.
        is what I just
        wrote, not  `that'\,''.

Dashes come in three sizes: an
       intra-word
dash, a medium dash for number ranges like
       1--2,
and a punctuation
       dash---like
this.

A sentence-ending space should be larger than the
space between words within a sentence.  You
sometimes have to type special commands in
conjunction with punctuation characters to get
this right, as in the following sentence.
       Gnats, gnus, etc.\ all  % `\ ' makes an inter-word space.
       begin with G\@.         % \@ marks end-of-sentence punctuation.
You should check the spaces after periods when
reading your output to make sure you haven't
forgotten any special cases.  Generating an
ellipsis
       \ldots\               % `\ ' is needed after `\ldots' because TeX
                             % ignores spaces after command names like \ldots
                             % made from \ + letters.
                             %
                             % Note how a `%' character causes TeX to ignore
                             % the end of the input line, so these blank lines
                             % do not start a new paragraph.
                             %
with the right spacing around the periods requires
a special command.

\LaTeX\ interprets some common characters as
commands, so you must type special commands to
generate them.  These characters include the
following:
       \$ \& \% \# \{ and \}.

In printing, text is usually emphasized with an
       \emph{italic}
type style.

\begin{em}
   A long segment of text can also be emphasized
   in this way.  Text within such a segment can be
   given \emph{additional} emphasis.
\end{em}

It is sometimes necessary to prevent \LaTeX\ from
breaking a line where it might otherwise do so.
This may be at a space, as between the ``Mr.''\ and
``Jones'' in
       ``Mr.~Jones'',        % ~ produces an unbreakable interword space.
or within a word---especially when the word is a
symbol like
       \mbox{\emph{itemnum}}
that makes little sense when hyphenated across
lines.

Footnotes\footnote{This is an example of a footnote.}
pose no problem.

\LaTeX\ is good at typesetting mathematical formulas
like
       \( x-3y + z = 7 \)
or
       \( a_{1} > x^{2n} + y^{2n} > x' \)
or
       \( \ip{A}{B} = \sum_{i} a_{i} b_{i} \).
The spaces you type in a formula are
ignored.  Remember that a letter like
       $x$                   % $ ... $  and  \( ... \)  are equivalent
is a formula when it denotes a mathematical
symbol, and it should be typed as one.

\section{Displayed Text}

Text is displayed by indenting it from the left
margin.  Quotations are commonly displayed.  There
are short quotations
\begin{quote}
   This is a short quotation.  It consists of a
   single paragraph of text.  See how it is formatted.
\end{quote}
and longer ones.
\begin{quotation}
   This is a longer quotation.  It consists of two
   paragraphs of text, neither of which are
   particularly interesting.

   This is the second paragraph of the quotation.  It
   is just as dull as the first paragraph.
\end{quotation}
Another frequently-displayed structure is a list.
The following is an example of an \emph{itemized}
list.
\begin{itemize}
   \item This is the first item of an itemized list.
         Each item in the list is marked with a ``tick''.
         You don't have to worry about what kind of tick
         mark is used.

   \item This is the second item of the list.  It
         contains another list nested inside it.  The inner
         list is an \emph{enumerated} list.
         \begin{enumerate}
            \item This is the first item of an enumerated
                  list that is nested within the itemized list.

            \item This is the second item of the inner list.
                  \LaTeX\ allows you to nest lists deeper than
                  you really should.
         \end{enumerate}
         This is the rest of the second item of the outer
         list.  It is no more interesting than any other
         part of the item.
   \item This is the third item of the list.
\end{itemize}
You can even display poetry.
\begin{verse}
   There is an environment
    for verse \\             % The \\ command separates lines
   Whose features some poets % within a stanza.
   will curse.

                             % One or more blank lines separate stanzas.

   For instead of making\\
   Them do \emph{all} line breaking, \\
   It allows them to put too many words on a line when they'd rather be
   forced to be terse.
\end{verse}

Mathematical formulas may also be displayed.  A
displayed formula
is
one-line long; multiline
formulas require special formatting instructions.
   \[  \ip{\Gamma}{\psi'} = x'' + y^{2} + z_{i}^{n}\]
Don't start a paragraph with a displayed equation,
nor make one a paragraph by itself.

\end{document}               % End of document.

\section{Bench Huge Section 1628}
% synthetic bench section 1628
% This is a sample LaTeX input file.  (Version of 12 August 2004.)
%
% A '%' character causes TeX to ignore all remaining text on the line,
% and is used for comments like this one.

\documentclass{article}      % Specifies the document class

                             % The preamble begins here.
\title{An Example Document}  % Declares the document's title.
\author{Leslie Lamport}      % Declares the author's name.
\date{January 21, 1994}      % Deleting this command produces today's date.

\newcommand{\ip}[2]{(#1, #2)}
                             % Defines \ip{arg1}{arg2} to mean
                             % (arg1, arg2).

%\newcommand{\ip}[2]{\langle #1 | #2\rangle}
                             % This is an alternative definition of
                             % \ip that is commented out.

\begin{document}             % End of preamble and beginning of text.

\maketitle                   % Produces the title.

This is an example input file.  Comparing it with
the output it generates can show you how to
produce a simple document of your own.

\section{Ordinary Text}      % Produces section heading.  Lower-level
                             % sections are begun with similar
                             % \subsection and \subsubsection commands.

The ends  of words and sentences are marked
  by   spaces. It  doesn't matter how many
spaces    you type; one is as good as 100.  The
end of   a line counts as a space.

One   or more   blank lines denote the  end
of  a paragraph.

Since any number of consecutive spaces are treated
like a single one, the formatting of the input
file makes no difference to
      \LaTeX,                % The \LaTeX command generates the LaTeX logo.
but it makes a difference to you.  When you use
\LaTeX, making your input file as easy to read
as possible will be a great help as you write
your document and when you change it.  This sample
file shows how you can add comments to your own input
file.

Because printing is different from typewriting,
there are a number of things that you have to do
differently when preparing an input file than if
you were just typing the document directly.
Quotation marks like
       ``this''
have to be handled specially, as do quotes within
quotes:
       ``\,`this'            % \, separates the double and single quote.
        is what I just
        wrote, not  `that'\,''.

Dashes come in three sizes: an
       intra-word
dash, a medium dash for number ranges like
       1--2,
and a punctuation
       dash---like
this.

A sentence-ending space should be larger than the
space between words within a sentence.  You
sometimes have to type special commands in
conjunction with punctuation characters to get
this right, as in the following sentence.
       Gnats, gnus, etc.\ all  % `\ ' makes an inter-word space.
       begin with G\@.         % \@ marks end-of-sentence punctuation.
You should check the spaces after periods when
reading your output to make sure you haven't
forgotten any special cases.  Generating an
ellipsis
       \ldots\               % `\ ' is needed after `\ldots' because TeX
                             % ignores spaces after command names like \ldots
                             % made from \ + letters.
                             %
                             % Note how a `%' character causes TeX to ignore
                             % the end of the input line, so these blank lines
                             % do not start a new paragraph.
                             %
with the right spacing around the periods requires
a special command.

\LaTeX\ interprets some common characters as
commands, so you must type special commands to
generate them.  These characters include the
following:
       \$ \& \% \# \{ and \}.

In printing, text is usually emphasized with an
       \emph{italic}
type style.

\begin{em}
   A long segment of text can also be emphasized
   in this way.  Text within such a segment can be
   given \emph{additional} emphasis.
\end{em}

It is sometimes necessary to prevent \LaTeX\ from
breaking a line where it might otherwise do so.
This may be at a space, as between the ``Mr.''\ and
``Jones'' in
       ``Mr.~Jones'',        % ~ produces an unbreakable interword space.
or within a word---especially when the word is a
symbol like
       \mbox{\emph{itemnum}}
that makes little sense when hyphenated across
lines.

Footnotes\footnote{This is an example of a footnote.}
pose no problem.

\LaTeX\ is good at typesetting mathematical formulas
like
       \( x-3y + z = 7 \)
or
       \( a_{1} > x^{2n} + y^{2n} > x' \)
or
       \( \ip{A}{B} = \sum_{i} a_{i} b_{i} \).
The spaces you type in a formula are
ignored.  Remember that a letter like
       $x$                   % $ ... $  and  \( ... \)  are equivalent
is a formula when it denotes a mathematical
symbol, and it should be typed as one.

\section{Displayed Text}

Text is displayed by indenting it from the left
margin.  Quotations are commonly displayed.  There
are short quotations
\begin{quote}
   This is a short quotation.  It consists of a
   single paragraph of text.  See how it is formatted.
\end{quote}
and longer ones.
\begin{quotation}
   This is a longer quotation.  It consists of two
   paragraphs of text, neither of which are
   particularly interesting.

   This is the second paragraph of the quotation.  It
   is just as dull as the first paragraph.
\end{quotation}
Another frequently-displayed structure is a list.
The following is an example of an \emph{itemized}
list.
\begin{itemize}
   \item This is the first item of an itemized list.
         Each item in the list is marked with a ``tick''.
         You don't have to worry about what kind of tick
         mark is used.

   \item This is the second item of the list.  It
         contains another list nested inside it.  The inner
         list is an \emph{enumerated} list.
         \begin{enumerate}
            \item This is the first item of an enumerated
                  list that is nested within the itemized list.

            \item This is the second item of the inner list.
                  \LaTeX\ allows you to nest lists deeper than
                  you really should.
         \end{enumerate}
         This is the rest of the second item of the outer
         list.  It is no more interesting than any other
         part of the item.
   \item This is the third item of the list.
\end{itemize}
You can even display poetry.
\begin{verse}
   There is an environment
    for verse \\             % The \\ command separates lines
   Whose features some poets % within a stanza.
   will curse.

                             % One or more blank lines separate stanzas.

   For instead of making\\
   Them do \emph{all} line breaking, \\
   It allows them to put too many words on a line when they'd rather be
   forced to be terse.
\end{verse}

Mathematical formulas may also be displayed.  A
displayed formula
is
one-line long; multiline
formulas require special formatting instructions.
   \[  \ip{\Gamma}{\psi'} = x'' + y^{2} + z_{i}^{n}\]
Don't start a paragraph with a displayed equation,
nor make one a paragraph by itself.

\end{document}               % End of document.

\section{Bench Huge Section 1629}
% synthetic bench section 1629
% This is a sample LaTeX input file.  (Version of 12 August 2004.)
%
% A '%' character causes TeX to ignore all remaining text on the line,
% and is used for comments like this one.

\documentclass{article}      % Specifies the document class

                             % The preamble begins here.
\title{An Example Document}  % Declares the document's title.
\author{Leslie Lamport}      % Declares the author's name.
\date{January 21, 1994}      % Deleting this command produces today's date.

\newcommand{\ip}[2]{(#1, #2)}
                             % Defines \ip{arg1}{arg2} to mean
                             % (arg1, arg2).

%\newcommand{\ip}[2]{\langle #1 | #2\rangle}
                             % This is an alternative definition of
                             % \ip that is commented out.

\begin{document}             % End of preamble and beginning of text.

\maketitle                   % Produces the title.

This is an example input file.  Comparing it with
the output it generates can show you how to
produce a simple document of your own.

\section{Ordinary Text}      % Produces section heading.  Lower-level
                             % sections are begun with similar
                             % \subsection and \subsubsection commands.

The ends  of words and sentences are marked
  by   spaces. It  doesn't matter how many
spaces    you type; one is as good as 100.  The
end of   a line counts as a space.

One   or more   blank lines denote the  end
of  a paragraph.

Since any number of consecutive spaces are treated
like a single one, the formatting of the input
file makes no difference to
      \LaTeX,                % The \LaTeX command generates the LaTeX logo.
but it makes a difference to you.  When you use
\LaTeX, making your input file as easy to read
as possible will be a great help as you write
your document and when you change it.  This sample
file shows how you can add comments to your own input
file.

Because printing is different from typewriting,
there are a number of things that you have to do
differently when preparing an input file than if
you were just typing the document directly.
Quotation marks like
       ``this''
have to be handled specially, as do quotes within
quotes:
       ``\,`this'            % \, separates the double and single quote.
        is what I just
        wrote, not  `that'\,''.

Dashes come in three sizes: an
       intra-word
dash, a medium dash for number ranges like
       1--2,
and a punctuation
       dash---like
this.

A sentence-ending space should be larger than the
space between words within a sentence.  You
sometimes have to type special commands in
conjunction with punctuation characters to get
this right, as in the following sentence.
       Gnats, gnus, etc.\ all  % `\ ' makes an inter-word space.
       begin with G\@.         % \@ marks end-of-sentence punctuation.
You should check the spaces after periods when
reading your output to make sure you haven't
forgotten any special cases.  Generating an
ellipsis
       \ldots\               % `\ ' is needed after `\ldots' because TeX
                             % ignores spaces after command names like \ldots
                             % made from \ + letters.
                             %
                             % Note how a `%' character causes TeX to ignore
                             % the end of the input line, so these blank lines
                             % do not start a new paragraph.
                             %
with the right spacing around the periods requires
a special command.

\LaTeX\ interprets some common characters as
commands, so you must type special commands to
generate them.  These characters include the
following:
       \$ \& \% \# \{ and \}.

In printing, text is usually emphasized with an
       \emph{italic}
type style.

\begin{em}
   A long segment of text can also be emphasized
   in this way.  Text within such a segment can be
   given \emph{additional} emphasis.
\end{em}

It is sometimes necessary to prevent \LaTeX\ from
breaking a line where it might otherwise do so.
This may be at a space, as between the ``Mr.''\ and
``Jones'' in
       ``Mr.~Jones'',        % ~ produces an unbreakable interword space.
or within a word---especially when the word is a
symbol like
       \mbox{\emph{itemnum}}
that makes little sense when hyphenated across
lines.

Footnotes\footnote{This is an example of a footnote.}
pose no problem.

\LaTeX\ is good at typesetting mathematical formulas
like
       \( x-3y + z = 7 \)
or
       \( a_{1} > x^{2n} + y^{2n} > x' \)
or
       \( \ip{A}{B} = \sum_{i} a_{i} b_{i} \).
The spaces you type in a formula are
ignored.  Remember that a letter like
       $x$                   % $ ... $  and  \( ... \)  are equivalent
is a formula when it denotes a mathematical
symbol, and it should be typed as one.

\section{Displayed Text}

Text is displayed by indenting it from the left
margin.  Quotations are commonly displayed.  There
are short quotations
\begin{quote}
   This is a short quotation.  It consists of a
   single paragraph of text.  See how it is formatted.
\end{quote}
and longer ones.
\begin{quotation}
   This is a longer quotation.  It consists of two
   paragraphs of text, neither of which are
   particularly interesting.

   This is the second paragraph of the quotation.  It
   is just as dull as the first paragraph.
\end{quotation}
Another frequently-displayed structure is a list.
The following is an example of an \emph{itemized}
list.
\begin{itemize}
   \item This is the first item of an itemized list.
         Each item in the list is marked with a ``tick''.
         You don't have to worry about what kind of tick
         mark is used.

   \item This is the second item of the list.  It
         contains another list nested inside it.  The inner
         list is an \emph{enumerated} list.
         \begin{enumerate}
            \item This is the first item of an enumerated
                  list that is nested within the itemized list.

            \item This is the second item of the inner list.
                  \LaTeX\ allows you to nest lists deeper than
                  you really should.
         \end{enumerate}
         This is the rest of the second item of the outer
         list.  It is no more interesting than any other
         part of the item.
   \item This is the third item of the list.
\end{itemize}
You can even display poetry.
\begin{verse}
   There is an environment
    for verse \\             % The \\ command separates lines
   Whose features some poets % within a stanza.
   will curse.

                             % One or more blank lines separate stanzas.

   For instead of making\\
   Them do \emph{all} line breaking, \\
   It allows them to put too many words on a line when they'd rather be
   forced to be terse.
\end{verse}

Mathematical formulas may also be displayed.  A
displayed formula
is
one-line long; multiline
formulas require special formatting instructions.
   \[  \ip{\Gamma}{\psi'} = x'' + y^{2} + z_{i}^{n}\]
Don't start a paragraph with a displayed equation,
nor make one a paragraph by itself.

\end{document}               % End of document.

\section{Bench Huge Section 1630}
% synthetic bench section 1630
% This is a sample LaTeX input file.  (Version of 12 August 2004.)
%
% A '%' character causes TeX to ignore all remaining text on the line,
% and is used for comments like this one.

\documentclass{article}      % Specifies the document class

                             % The preamble begins here.
\title{An Example Document}  % Declares the document's title.
\author{Leslie Lamport}      % Declares the author's name.
\date{January 21, 1994}      % Deleting this command produces today's date.

\newcommand{\ip}[2]{(#1, #2)}
                             % Defines \ip{arg1}{arg2} to mean
                             % (arg1, arg2).

%\newcommand{\ip}[2]{\langle #1 | #2\rangle}
                             % This is an alternative definition of
                             % \ip that is commented out.

\begin{document}             % End of preamble and beginning of text.

\maketitle                   % Produces the title.

This is an example input file.  Comparing it with
the output it generates can show you how to
produce a simple document of your own.

\section{Ordinary Text}      % Produces section heading.  Lower-level
                             % sections are begun with similar
                             % \subsection and \subsubsection commands.

The ends  of words and sentences are marked
  by   spaces. It  doesn't matter how many
spaces    you type; one is as good as 100.  The
end of   a line counts as a space.

One   or more   blank lines denote the  end
of  a paragraph.

Since any number of consecutive spaces are treated
like a single one, the formatting of the input
file makes no difference to
      \LaTeX,                % The \LaTeX command generates the LaTeX logo.
but it makes a difference to you.  When you use
\LaTeX, making your input file as easy to read
as possible will be a great help as you write
your document and when you change it.  This sample
file shows how you can add comments to your own input
file.

Because printing is different from typewriting,
there are a number of things that you have to do
differently when preparing an input file than if
you were just typing the document directly.
Quotation marks like
       ``this''
have to be handled specially, as do quotes within
quotes:
       ``\,`this'            % \, separates the double and single quote.
        is what I just
        wrote, not  `that'\,''.

Dashes come in three sizes: an
       intra-word
dash, a medium dash for number ranges like
       1--2,
and a punctuation
       dash---like
this.

A sentence-ending space should be larger than the
space between words within a sentence.  You
sometimes have to type special commands in
conjunction with punctuation characters to get
this right, as in the following sentence.
       Gnats, gnus, etc.\ all  % `\ ' makes an inter-word space.
       begin with G\@.         % \@ marks end-of-sentence punctuation.
You should check the spaces after periods when
reading your output to make sure you haven't
forgotten any special cases.  Generating an
ellipsis
       \ldots\               % `\ ' is needed after `\ldots' because TeX
                             % ignores spaces after command names like \ldots
                             % made from \ + letters.
                             %
                             % Note how a `%' character causes TeX to ignore
                             % the end of the input line, so these blank lines
                             % do not start a new paragraph.
                             %
with the right spacing around the periods requires
a special command.

\LaTeX\ interprets some common characters as
commands, so you must type special commands to
generate them.  These characters include the
following:
       \$ \& \% \# \{ and \}.

In printing, text is usually emphasized with an
       \emph{italic}
type style.

\begin{em}
   A long segment of text can also be emphasized
   in this way.  Text within such a segment can be
   given \emph{additional} emphasis.
\end{em}

It is sometimes necessary to prevent \LaTeX\ from
breaking a line where it might otherwise do so.
This may be at a space, as between the ``Mr.''\ and
``Jones'' in
       ``Mr.~Jones'',        % ~ produces an unbreakable interword space.
or within a word---especially when the word is a
symbol like
       \mbox{\emph{itemnum}}
that makes little sense when hyphenated across
lines.

Footnotes\footnote{This is an example of a footnote.}
pose no problem.

\LaTeX\ is good at typesetting mathematical formulas
like
       \( x-3y + z = 7 \)
or
       \( a_{1} > x^{2n} + y^{2n} > x' \)
or
       \( \ip{A}{B} = \sum_{i} a_{i} b_{i} \).
The spaces you type in a formula are
ignored.  Remember that a letter like
       $x$                   % $ ... $  and  \( ... \)  are equivalent
is a formula when it denotes a mathematical
symbol, and it should be typed as one.

\section{Displayed Text}

Text is displayed by indenting it from the left
margin.  Quotations are commonly displayed.  There
are short quotations
\begin{quote}
   This is a short quotation.  It consists of a
   single paragraph of text.  See how it is formatted.
\end{quote}
and longer ones.
\begin{quotation}
   This is a longer quotation.  It consists of two
   paragraphs of text, neither of which are
   particularly interesting.

   This is the second paragraph of the quotation.  It
   is just as dull as the first paragraph.
\end{quotation}
Another frequently-displayed structure is a list.
The following is an example of an \emph{itemized}
list.
\begin{itemize}
   \item This is the first item of an itemized list.
         Each item in the list is marked with a ``tick''.
         You don't have to worry about what kind of tick
         mark is used.

   \item This is the second item of the list.  It
         contains another list nested inside it.  The inner
         list is an \emph{enumerated} list.
         \begin{enumerate}
            \item This is the first item of an enumerated
                  list that is nested within the itemized list.

            \item This is the second item of the inner list.
                  \LaTeX\ allows you to nest lists deeper than
                  you really should.
         \end{enumerate}
         This is the rest of the second item of the outer
         list.  It is no more interesting than any other
         part of the item.
   \item This is the third item of the list.
\end{itemize}
You can even display poetry.
\begin{verse}
   There is an environment
    for verse \\             % The \\ command separates lines
   Whose features some poets % within a stanza.
   will curse.

                             % One or more blank lines separate stanzas.

   For instead of making\\
   Them do \emph{all} line breaking, \\
   It allows them to put too many words on a line when they'd rather be
   forced to be terse.
\end{verse}

Mathematical formulas may also be displayed.  A
displayed formula
is
one-line long; multiline
formulas require special formatting instructions.
   \[  \ip{\Gamma}{\psi'} = x'' + y^{2} + z_{i}^{n}\]
Don't start a paragraph with a displayed equation,
nor make one a paragraph by itself.

\end{document}               % End of document.

\section{Bench Huge Section 1631}
% synthetic bench section 1631
% This is a sample LaTeX input file.  (Version of 12 August 2004.)
%
% A '%' character causes TeX to ignore all remaining text on the line,
% and is used for comments like this one.

\documentclass{article}      % Specifies the document class

                             % The preamble begins here.
\title{An Example Document}  % Declares the document's title.
\author{Leslie Lamport}      % Declares the author's name.
\date{January 21, 1994}      % Deleting this command produces today's date.

\newcommand{\ip}[2]{(#1, #2)}
                             % Defines \ip{arg1}{arg2} to mean
                             % (arg1, arg2).

%\newcommand{\ip}[2]{\langle #1 | #2\rangle}
                             % This is an alternative definition of
                             % \ip that is commented out.

\begin{document}             % End of preamble and beginning of text.

\maketitle                   % Produces the title.

This is an example input file.  Comparing it with
the output it generates can show you how to
produce a simple document of your own.

\section{Ordinary Text}      % Produces section heading.  Lower-level
                             % sections are begun with similar
                             % \subsection and \subsubsection commands.

The ends  of words and sentences are marked
  by   spaces. It  doesn't matter how many
spaces    you type; one is as good as 100.  The
end of   a line counts as a space.

One   or more   blank lines denote the  end
of  a paragraph.

Since any number of consecutive spaces are treated
like a single one, the formatting of the input
file makes no difference to
      \LaTeX,                % The \LaTeX command generates the LaTeX logo.
but it makes a difference to you.  When you use
\LaTeX, making your input file as easy to read
as possible will be a great help as you write
your document and when you change it.  This sample
file shows how you can add comments to your own input
file.

Because printing is different from typewriting,
there are a number of things that you have to do
differently when preparing an input file than if
you were just typing the document directly.
Quotation marks like
       ``this''
have to be handled specially, as do quotes within
quotes:
       ``\,`this'            % \, separates the double and single quote.
        is what I just
        wrote, not  `that'\,''.

Dashes come in three sizes: an
       intra-word
dash, a medium dash for number ranges like
       1--2,
and a punctuation
       dash---like
this.

A sentence-ending space should be larger than the
space between words within a sentence.  You
sometimes have to type special commands in
conjunction with punctuation characters to get
this right, as in the following sentence.
       Gnats, gnus, etc.\ all  % `\ ' makes an inter-word space.
       begin with G\@.         % \@ marks end-of-sentence punctuation.
You should check the spaces after periods when
reading your output to make sure you haven't
forgotten any special cases.  Generating an
ellipsis
       \ldots\               % `\ ' is needed after `\ldots' because TeX
                             % ignores spaces after command names like \ldots
                             % made from \ + letters.
                             %
                             % Note how a `%' character causes TeX to ignore
                             % the end of the input line, so these blank lines
                             % do not start a new paragraph.
                             %
with the right spacing around the periods requires
a special command.

\LaTeX\ interprets some common characters as
commands, so you must type special commands to
generate them.  These characters include the
following:
       \$ \& \% \# \{ and \}.

In printing, text is usually emphasized with an
       \emph{italic}
type style.

\begin{em}
   A long segment of text can also be emphasized
   in this way.  Text within such a segment can be
   given \emph{additional} emphasis.
\end{em}

It is sometimes necessary to prevent \LaTeX\ from
breaking a line where it might otherwise do so.
This may be at a space, as between the ``Mr.''\ and
``Jones'' in
       ``Mr.~Jones'',        % ~ produces an unbreakable interword space.
or within a word---especially when the word is a
symbol like
       \mbox{\emph{itemnum}}
that makes little sense when hyphenated across
lines.

Footnotes\footnote{This is an example of a footnote.}
pose no problem.

\LaTeX\ is good at typesetting mathematical formulas
like
       \( x-3y + z = 7 \)
or
       \( a_{1} > x^{2n} + y^{2n} > x' \)
or
       \( \ip{A}{B} = \sum_{i} a_{i} b_{i} \).
The spaces you type in a formula are
ignored.  Remember that a letter like
       $x$                   % $ ... $  and  \( ... \)  are equivalent
is a formula when it denotes a mathematical
symbol, and it should be typed as one.

\section{Displayed Text}

Text is displayed by indenting it from the left
margin.  Quotations are commonly displayed.  There
are short quotations
\begin{quote}
   This is a short quotation.  It consists of a
   single paragraph of text.  See how it is formatted.
\end{quote}
and longer ones.
\begin{quotation}
   This is a longer quotation.  It consists of two
   paragraphs of text, neither of which are
   particularly interesting.

   This is the second paragraph of the quotation.  It
   is just as dull as the first paragraph.
\end{quotation}
Another frequently-displayed structure is a list.
The following is an example of an \emph{itemized}
list.
\begin{itemize}
   \item This is the first item of an itemized list.
         Each item in the list is marked with a ``tick''.
         You don't have to worry about what kind of tick
         mark is used.

   \item This is the second item of the list.  It
         contains another list nested inside it.  The inner
         list is an \emph{enumerated} list.
         \begin{enumerate}
            \item This is the first item of an enumerated
                  list that is nested within the itemized list.

            \item This is the second item of the inner list.
                  \LaTeX\ allows you to nest lists deeper than
                  you really should.
         \end{enumerate}
         This is the rest of the second item of the outer
         list.  It is no more interesting than any other
         part of the item.
   \item This is the third item of the list.
\end{itemize}
You can even display poetry.
\begin{verse}
   There is an environment
    for verse \\             % The \\ command separates lines
   Whose features some poets % within a stanza.
   will curse.

                             % One or more blank lines separate stanzas.

   For instead of making\\
   Them do \emph{all} line breaking, \\
   It allows them to put too many words on a line when they'd rather be
   forced to be terse.
\end{verse}

Mathematical formulas may also be displayed.  A
displayed formula
is
one-line long; multiline
formulas require special formatting instructions.
   \[  \ip{\Gamma}{\psi'} = x'' + y^{2} + z_{i}^{n}\]
Don't start a paragraph with a displayed equation,
nor make one a paragraph by itself.

\end{document}               % End of document.

\section{Bench Huge Section 1632}
% synthetic bench section 1632
% This is a sample LaTeX input file.  (Version of 12 August 2004.)
%
% A '%' character causes TeX to ignore all remaining text on the line,
% and is used for comments like this one.

\documentclass{article}      % Specifies the document class

                             % The preamble begins here.
\title{An Example Document}  % Declares the document's title.
\author{Leslie Lamport}      % Declares the author's name.
\date{January 21, 1994}      % Deleting this command produces today's date.

\newcommand{\ip}[2]{(#1, #2)}
                             % Defines \ip{arg1}{arg2} to mean
                             % (arg1, arg2).

%\newcommand{\ip}[2]{\langle #1 | #2\rangle}
                             % This is an alternative definition of
                             % \ip that is commented out.

\begin{document}             % End of preamble and beginning of text.

\maketitle                   % Produces the title.

This is an example input file.  Comparing it with
the output it generates can show you how to
produce a simple document of your own.

\section{Ordinary Text}      % Produces section heading.  Lower-level
                             % sections are begun with similar
                             % \subsection and \subsubsection commands.

The ends  of words and sentences are marked
  by   spaces. It  doesn't matter how many
spaces    you type; one is as good as 100.  The
end of   a line counts as a space.

One   or more   blank lines denote the  end
of  a paragraph.

Since any number of consecutive spaces are treated
like a single one, the formatting of the input
file makes no difference to
      \LaTeX,                % The \LaTeX command generates the LaTeX logo.
but it makes a difference to you.  When you use
\LaTeX, making your input file as easy to read
as possible will be a great help as you write
your document and when you change it.  This sample
file shows how you can add comments to your own input
file.

Because printing is different from typewriting,
there are a number of things that you have to do
differently when preparing an input file than if
you were just typing the document directly.
Quotation marks like
       ``this''
have to be handled specially, as do quotes within
quotes:
       ``\,`this'            % \, separates the double and single quote.
        is what I just
        wrote, not  `that'\,''.

Dashes come in three sizes: an
       intra-word
dash, a medium dash for number ranges like
       1--2,
and a punctuation
       dash---like
this.

A sentence-ending space should be larger than the
space between words within a sentence.  You
sometimes have to type special commands in
conjunction with punctuation characters to get
this right, as in the following sentence.
       Gnats, gnus, etc.\ all  % `\ ' makes an inter-word space.
       begin with G\@.         % \@ marks end-of-sentence punctuation.
You should check the spaces after periods when
reading your output to make sure you haven't
forgotten any special cases.  Generating an
ellipsis
       \ldots\               % `\ ' is needed after `\ldots' because TeX
                             % ignores spaces after command names like \ldots
                             % made from \ + letters.
                             %
                             % Note how a `%' character causes TeX to ignore
                             % the end of the input line, so these blank lines
                             % do not start a new paragraph.
                             %
with the right spacing around the periods requires
a special command.

\LaTeX\ interprets some common characters as
commands, so you must type special commands to
generate them.  These characters include the
following:
       \$ \& \% \# \{ and \}.

In printing, text is usually emphasized with an
       \emph{italic}
type style.

\begin{em}
   A long segment of text can also be emphasized
   in this way.  Text within such a segment can be
   given \emph{additional} emphasis.
\end{em}

It is sometimes necessary to prevent \LaTeX\ from
breaking a line where it might otherwise do so.
This may be at a space, as between the ``Mr.''\ and
``Jones'' in
       ``Mr.~Jones'',        % ~ produces an unbreakable interword space.
or within a word---especially when the word is a
symbol like
       \mbox{\emph{itemnum}}
that makes little sense when hyphenated across
lines.

Footnotes\footnote{This is an example of a footnote.}
pose no problem.

\LaTeX\ is good at typesetting mathematical formulas
like
       \( x-3y + z = 7 \)
or
       \( a_{1} > x^{2n} + y^{2n} > x' \)
or
       \( \ip{A}{B} = \sum_{i} a_{i} b_{i} \).
The spaces you type in a formula are
ignored.  Remember that a letter like
       $x$                   % $ ... $  and  \( ... \)  are equivalent
is a formula when it denotes a mathematical
symbol, and it should be typed as one.

\section{Displayed Text}

Text is displayed by indenting it from the left
margin.  Quotations are commonly displayed.  There
are short quotations
\begin{quote}
   This is a short quotation.  It consists of a
   single paragraph of text.  See how it is formatted.
\end{quote}
and longer ones.
\begin{quotation}
   This is a longer quotation.  It consists of two
   paragraphs of text, neither of which are
   particularly interesting.

   This is the second paragraph of the quotation.  It
   is just as dull as the first paragraph.
\end{quotation}
Another frequently-displayed structure is a list.
The following is an example of an \emph{itemized}
list.
\begin{itemize}
   \item This is the first item of an itemized list.
         Each item in the list is marked with a ``tick''.
         You don't have to worry about what kind of tick
         mark is used.

   \item This is the second item of the list.  It
         contains another list nested inside it.  The inner
         list is an \emph{enumerated} list.
         \begin{enumerate}
            \item This is the first item of an enumerated
                  list that is nested within the itemized list.

            \item This is the second item of the inner list.
                  \LaTeX\ allows you to nest lists deeper than
                  you really should.
         \end{enumerate}
         This is the rest of the second item of the outer
         list.  It is no more interesting than any other
         part of the item.
   \item This is the third item of the list.
\end{itemize}
You can even display poetry.
\begin{verse}
   There is an environment
    for verse \\             % The \\ command separates lines
   Whose features some poets % within a stanza.
   will curse.

                             % One or more blank lines separate stanzas.

   For instead of making\\
   Them do \emph{all} line breaking, \\
   It allows them to put too many words on a line when they'd rather be
   forced to be terse.
\end{verse}

Mathematical formulas may also be displayed.  A
displayed formula
is
one-line long; multiline
formulas require special formatting instructions.
   \[  \ip{\Gamma}{\psi'} = x'' + y^{2} + z_{i}^{n}\]
Don't start a paragraph with a displayed equation,
nor make one a paragraph by itself.

\end{document}               % End of document.

\section{Bench Huge Section 1633}
% synthetic bench section 1633
% This is a sample LaTeX input file.  (Version of 12 August 2004.)
%
% A '%' character causes TeX to ignore all remaining text on the line,
% and is used for comments like this one.

\documentclass{article}      % Specifies the document class

                             % The preamble begins here.
\title{An Example Document}  % Declares the document's title.
\author{Leslie Lamport}      % Declares the author's name.
\date{January 21, 1994}      % Deleting this command produces today's date.

\newcommand{\ip}[2]{(#1, #2)}
                             % Defines \ip{arg1}{arg2} to mean
                             % (arg1, arg2).

%\newcommand{\ip}[2]{\langle #1 | #2\rangle}
                             % This is an alternative definition of
                             % \ip that is commented out.

\begin{document}             % End of preamble and beginning of text.

\maketitle                   % Produces the title.

This is an example input file.  Comparing it with
the output it generates can show you how to
produce a simple document of your own.

\section{Ordinary Text}      % Produces section heading.  Lower-level
                             % sections are begun with similar
                             % \subsection and \subsubsection commands.

The ends  of words and sentences are marked
  by   spaces. It  doesn't matter how many
spaces    you type; one is as good as 100.  The
end of   a line counts as a space.

One   or more   blank lines denote the  end
of  a paragraph.

Since any number of consecutive spaces are treated
like a single one, the formatting of the input
file makes no difference to
      \LaTeX,                % The \LaTeX command generates the LaTeX logo.
but it makes a difference to you.  When you use
\LaTeX, making your input file as easy to read
as possible will be a great help as you write
your document and when you change it.  This sample
file shows how you can add comments to your own input
file.

Because printing is different from typewriting,
there are a number of things that you have to do
differently when preparing an input file than if
you were just typing the document directly.
Quotation marks like
       ``this''
have to be handled specially, as do quotes within
quotes:
       ``\,`this'            % \, separates the double and single quote.
        is what I just
        wrote, not  `that'\,''.

Dashes come in three sizes: an
       intra-word
dash, a medium dash for number ranges like
       1--2,
and a punctuation
       dash---like
this.

A sentence-ending space should be larger than the
space between words within a sentence.  You
sometimes have to type special commands in
conjunction with punctuation characters to get
this right, as in the following sentence.
       Gnats, gnus, etc.\ all  % `\ ' makes an inter-word space.
       begin with G\@.         % \@ marks end-of-sentence punctuation.
You should check the spaces after periods when
reading your output to make sure you haven't
forgotten any special cases.  Generating an
ellipsis
       \ldots\               % `\ ' is needed after `\ldots' because TeX
                             % ignores spaces after command names like \ldots
                             % made from \ + letters.
                             %
                             % Note how a `%' character causes TeX to ignore
                             % the end of the input line, so these blank lines
                             % do not start a new paragraph.
                             %
with the right spacing around the periods requires
a special command.

\LaTeX\ interprets some common characters as
commands, so you must type special commands to
generate them.  These characters include the
following:
       \$ \& \% \# \{ and \}.

In printing, text is usually emphasized with an
       \emph{italic}
type style.

\begin{em}
   A long segment of text can also be emphasized
   in this way.  Text within such a segment can be
   given \emph{additional} emphasis.
\end{em}

It is sometimes necessary to prevent \LaTeX\ from
breaking a line where it might otherwise do so.
This may be at a space, as between the ``Mr.''\ and
``Jones'' in
       ``Mr.~Jones'',        % ~ produces an unbreakable interword space.
or within a word---especially when the word is a
symbol like
       \mbox{\emph{itemnum}}
that makes little sense when hyphenated across
lines.

Footnotes\footnote{This is an example of a footnote.}
pose no problem.

\LaTeX\ is good at typesetting mathematical formulas
like
       \( x-3y + z = 7 \)
or
       \( a_{1} > x^{2n} + y^{2n} > x' \)
or
       \( \ip{A}{B} = \sum_{i} a_{i} b_{i} \).
The spaces you type in a formula are
ignored.  Remember that a letter like
       $x$                   % $ ... $  and  \( ... \)  are equivalent
is a formula when it denotes a mathematical
symbol, and it should be typed as one.

\section{Displayed Text}

Text is displayed by indenting it from the left
margin.  Quotations are commonly displayed.  There
are short quotations
\begin{quote}
   This is a short quotation.  It consists of a
   single paragraph of text.  See how it is formatted.
\end{quote}
and longer ones.
\begin{quotation}
   This is a longer quotation.  It consists of two
   paragraphs of text, neither of which are
   particularly interesting.

   This is the second paragraph of the quotation.  It
   is just as dull as the first paragraph.
\end{quotation}
Another frequently-displayed structure is a list.
The following is an example of an \emph{itemized}
list.
\begin{itemize}
   \item This is the first item of an itemized list.
         Each item in the list is marked with a ``tick''.
         You don't have to worry about what kind of tick
         mark is used.

   \item This is the second item of the list.  It
         contains another list nested inside it.  The inner
         list is an \emph{enumerated} list.
         \begin{enumerate}
            \item This is the first item of an enumerated
                  list that is nested within the itemized list.

            \item This is the second item of the inner list.
                  \LaTeX\ allows you to nest lists deeper than
                  you really should.
         \end{enumerate}
         This is the rest of the second item of the outer
         list.  It is no more interesting than any other
         part of the item.
   \item This is the third item of the list.
\end{itemize}
You can even display poetry.
\begin{verse}
   There is an environment
    for verse \\             % The \\ command separates lines
   Whose features some poets % within a stanza.
   will curse.

                             % One or more blank lines separate stanzas.

   For instead of making\\
   Them do \emph{all} line breaking, \\
   It allows them to put too many words on a line when they'd rather be
   forced to be terse.
\end{verse}

Mathematical formulas may also be displayed.  A
displayed formula
is
one-line long; multiline
formulas require special formatting instructions.
   \[  \ip{\Gamma}{\psi'} = x'' + y^{2} + z_{i}^{n}\]
Don't start a paragraph with a displayed equation,
nor make one a paragraph by itself.

\end{document}               % End of document.

\section{Bench Huge Section 1634}
% synthetic bench section 1634
% This is a sample LaTeX input file.  (Version of 12 August 2004.)
%
% A '%' character causes TeX to ignore all remaining text on the line,
% and is used for comments like this one.

\documentclass{article}      % Specifies the document class

                             % The preamble begins here.
\title{An Example Document}  % Declares the document's title.
\author{Leslie Lamport}      % Declares the author's name.
\date{January 21, 1994}      % Deleting this command produces today's date.

\newcommand{\ip}[2]{(#1, #2)}
                             % Defines \ip{arg1}{arg2} to mean
                             % (arg1, arg2).

%\newcommand{\ip}[2]{\langle #1 | #2\rangle}
                             % This is an alternative definition of
                             % \ip that is commented out.

\begin{document}             % End of preamble and beginning of text.

\maketitle                   % Produces the title.

This is an example input file.  Comparing it with
the output it generates can show you how to
produce a simple document of your own.

\section{Ordinary Text}      % Produces section heading.  Lower-level
                             % sections are begun with similar
                             % \subsection and \subsubsection commands.

The ends  of words and sentences are marked
  by   spaces. It  doesn't matter how many
spaces    you type; one is as good as 100.  The
end of   a line counts as a space.

One   or more   blank lines denote the  end
of  a paragraph.

Since any number of consecutive spaces are treated
like a single one, the formatting of the input
file makes no difference to
      \LaTeX,                % The \LaTeX command generates the LaTeX logo.
but it makes a difference to you.  When you use
\LaTeX, making your input file as easy to read
as possible will be a great help as you write
your document and when you change it.  This sample
file shows how you can add comments to your own input
file.

Because printing is different from typewriting,
there are a number of things that you have to do
differently when preparing an input file than if
you were just typing the document directly.
Quotation marks like
       ``this''
have to be handled specially, as do quotes within
quotes:
       ``\,`this'            % \, separates the double and single quote.
        is what I just
        wrote, not  `that'\,''.

Dashes come in three sizes: an
       intra-word
dash, a medium dash for number ranges like
       1--2,
and a punctuation
       dash---like
this.

A sentence-ending space should be larger than the
space between words within a sentence.  You
sometimes have to type special commands in
conjunction with punctuation characters to get
this right, as in the following sentence.
       Gnats, gnus, etc.\ all  % `\ ' makes an inter-word space.
       begin with G\@.         % \@ marks end-of-sentence punctuation.
You should check the spaces after periods when
reading your output to make sure you haven't
forgotten any special cases.  Generating an
ellipsis
       \ldots\               % `\ ' is needed after `\ldots' because TeX
                             % ignores spaces after command names like \ldots
                             % made from \ + letters.
                             %
                             % Note how a `%' character causes TeX to ignore
                             % the end of the input line, so these blank lines
                             % do not start a new paragraph.
                             %
with the right spacing around the periods requires
a special command.

\LaTeX\ interprets some common characters as
commands, so you must type special commands to
generate them.  These characters include the
following:
       \$ \& \% \# \{ and \}.

In printing, text is usually emphasized with an
       \emph{italic}
type style.

\begin{em}
   A long segment of text can also be emphasized
   in this way.  Text within such a segment can be
   given \emph{additional} emphasis.
\end{em}

It is sometimes necessary to prevent \LaTeX\ from
breaking a line where it might otherwise do so.
This may be at a space, as between the ``Mr.''\ and
``Jones'' in
       ``Mr.~Jones'',        % ~ produces an unbreakable interword space.
or within a word---especially when the word is a
symbol like
       \mbox{\emph{itemnum}}
that makes little sense when hyphenated across
lines.

Footnotes\footnote{This is an example of a footnote.}
pose no problem.

\LaTeX\ is good at typesetting mathematical formulas
like
       \( x-3y + z = 7 \)
or
       \( a_{1} > x^{2n} + y^{2n} > x' \)
or
       \( \ip{A}{B} = \sum_{i} a_{i} b_{i} \).
The spaces you type in a formula are
ignored.  Remember that a letter like
       $x$                   % $ ... $  and  \( ... \)  are equivalent
is a formula when it denotes a mathematical
symbol, and it should be typed as one.

\section{Displayed Text}

Text is displayed by indenting it from the left
margin.  Quotations are commonly displayed.  There
are short quotations
\begin{quote}
   This is a short quotation.  It consists of a
   single paragraph of text.  See how it is formatted.
\end{quote}
and longer ones.
\begin{quotation}
   This is a longer quotation.  It consists of two
   paragraphs of text, neither of which are
   particularly interesting.

   This is the second paragraph of the quotation.  It
   is just as dull as the first paragraph.
\end{quotation}
Another frequently-displayed structure is a list.
The following is an example of an \emph{itemized}
list.
\begin{itemize}
   \item This is the first item of an itemized list.
         Each item in the list is marked with a ``tick''.
         You don't have to worry about what kind of tick
         mark is used.

   \item This is the second item of the list.  It
         contains another list nested inside it.  The inner
         list is an \emph{enumerated} list.
         \begin{enumerate}
            \item This is the first item of an enumerated
                  list that is nested within the itemized list.

            \item This is the second item of the inner list.
                  \LaTeX\ allows you to nest lists deeper than
                  you really should.
         \end{enumerate}
         This is the rest of the second item of the outer
         list.  It is no more interesting than any other
         part of the item.
   \item This is the third item of the list.
\end{itemize}
You can even display poetry.
\begin{verse}
   There is an environment
    for verse \\             % The \\ command separates lines
   Whose features some poets % within a stanza.
   will curse.

                             % One or more blank lines separate stanzas.

   For instead of making\\
   Them do \emph{all} line breaking, \\
   It allows them to put too many words on a line when they'd rather be
   forced to be terse.
\end{verse}

Mathematical formulas may also be displayed.  A
displayed formula
is
one-line long; multiline
formulas require special formatting instructions.
   \[  \ip{\Gamma}{\psi'} = x'' + y^{2} + z_{i}^{n}\]
Don't start a paragraph with a displayed equation,
nor make one a paragraph by itself.

\end{document}               % End of document.

\section{Bench Huge Section 1635}
% synthetic bench section 1635
% This is a sample LaTeX input file.  (Version of 12 August 2004.)
%
% A '%' character causes TeX to ignore all remaining text on the line,
% and is used for comments like this one.

\documentclass{article}      % Specifies the document class

                             % The preamble begins here.
\title{An Example Document}  % Declares the document's title.
\author{Leslie Lamport}      % Declares the author's name.
\date{January 21, 1994}      % Deleting this command produces today's date.

\newcommand{\ip}[2]{(#1, #2)}
                             % Defines \ip{arg1}{arg2} to mean
                             % (arg1, arg2).

%\newcommand{\ip}[2]{\langle #1 | #2\rangle}
                             % This is an alternative definition of
                             % \ip that is commented out.

\begin{document}             % End of preamble and beginning of text.

\maketitle                   % Produces the title.

This is an example input file.  Comparing it with
the output it generates can show you how to
produce a simple document of your own.

\section{Ordinary Text}      % Produces section heading.  Lower-level
                             % sections are begun with similar
                             % \subsection and \subsubsection commands.

The ends  of words and sentences are marked
  by   spaces. It  doesn't matter how many
spaces    you type; one is as good as 100.  The
end of   a line counts as a space.

One   or more   blank lines denote the  end
of  a paragraph.

Since any number of consecutive spaces are treated
like a single one, the formatting of the input
file makes no difference to
      \LaTeX,                % The \LaTeX command generates the LaTeX logo.
but it makes a difference to you.  When you use
\LaTeX, making your input file as easy to read
as possible will be a great help as you write
your document and when you change it.  This sample
file shows how you can add comments to your own input
file.

Because printing is different from typewriting,
there are a number of things that you have to do
differently when preparing an input file than if
you were just typing the document directly.
Quotation marks like
       ``this''
have to be handled specially, as do quotes within
quotes:
       ``\,`this'            % \, separates the double and single quote.
        is what I just
        wrote, not  `that'\,''.

Dashes come in three sizes: an
       intra-word
dash, a medium dash for number ranges like
       1--2,
and a punctuation
       dash---like
this.

A sentence-ending space should be larger than the
space between words within a sentence.  You
sometimes have to type special commands in
conjunction with punctuation characters to get
this right, as in the following sentence.
       Gnats, gnus, etc.\ all  % `\ ' makes an inter-word space.
       begin with G\@.         % \@ marks end-of-sentence punctuation.
You should check the spaces after periods when
reading your output to make sure you haven't
forgotten any special cases.  Generating an
ellipsis
       \ldots\               % `\ ' is needed after `\ldots' because TeX
                             % ignores spaces after command names like \ldots
                             % made from \ + letters.
                             %
                             % Note how a `%' character causes TeX to ignore
                             % the end of the input line, so these blank lines
                             % do not start a new paragraph.
                             %
with the right spacing around the periods requires
a special command.

\LaTeX\ interprets some common characters as
commands, so you must type special commands to
generate them.  These characters include the
following:
       \$ \& \% \# \{ and \}.

In printing, text is usually emphasized with an
       \emph{italic}
type style.

\begin{em}
   A long segment of text can also be emphasized
   in this way.  Text within such a segment can be
   given \emph{additional} emphasis.
\end{em}

It is sometimes necessary to prevent \LaTeX\ from
breaking a line where it might otherwise do so.
This may be at a space, as between the ``Mr.''\ and
``Jones'' in
       ``Mr.~Jones'',        % ~ produces an unbreakable interword space.
or within a word---especially when the word is a
symbol like
       \mbox{\emph{itemnum}}
that makes little sense when hyphenated across
lines.

Footnotes\footnote{This is an example of a footnote.}
pose no problem.

\LaTeX\ is good at typesetting mathematical formulas
like
       \( x-3y + z = 7 \)
or
       \( a_{1} > x^{2n} + y^{2n} > x' \)
or
       \( \ip{A}{B} = \sum_{i} a_{i} b_{i} \).
The spaces you type in a formula are
ignored.  Remember that a letter like
       $x$                   % $ ... $  and  \( ... \)  are equivalent
is a formula when it denotes a mathematical
symbol, and it should be typed as one.

\section{Displayed Text}

Text is displayed by indenting it from the left
margin.  Quotations are commonly displayed.  There
are short quotations
\begin{quote}
   This is a short quotation.  It consists of a
   single paragraph of text.  See how it is formatted.
\end{quote}
and longer ones.
\begin{quotation}
   This is a longer quotation.  It consists of two
   paragraphs of text, neither of which are
   particularly interesting.

   This is the second paragraph of the quotation.  It
   is just as dull as the first paragraph.
\end{quotation}
Another frequently-displayed structure is a list.
The following is an example of an \emph{itemized}
list.
\begin{itemize}
   \item This is the first item of an itemized list.
         Each item in the list is marked with a ``tick''.
         You don't have to worry about what kind of tick
         mark is used.

   \item This is the second item of the list.  It
         contains another list nested inside it.  The inner
         list is an \emph{enumerated} list.
         \begin{enumerate}
            \item This is the first item of an enumerated
                  list that is nested within the itemized list.

            \item This is the second item of the inner list.
                  \LaTeX\ allows you to nest lists deeper than
                  you really should.
         \end{enumerate}
         This is the rest of the second item of the outer
         list.  It is no more interesting than any other
         part of the item.
   \item This is the third item of the list.
\end{itemize}
You can even display poetry.
\begin{verse}
   There is an environment
    for verse \\             % The \\ command separates lines
   Whose features some poets % within a stanza.
   will curse.

                             % One or more blank lines separate stanzas.

   For instead of making\\
   Them do \emph{all} line breaking, \\
   It allows them to put too many words on a line when they'd rather be
   forced to be terse.
\end{verse}

Mathematical formulas may also be displayed.  A
displayed formula
is
one-line long; multiline
formulas require special formatting instructions.
   \[  \ip{\Gamma}{\psi'} = x'' + y^{2} + z_{i}^{n}\]
Don't start a paragraph with a displayed equation,
nor make one a paragraph by itself.

\end{document}               % End of document.

\section{Bench Huge Section 1636}
% synthetic bench section 1636
% This is a sample LaTeX input file.  (Version of 12 August 2004.)
%
% A '%' character causes TeX to ignore all remaining text on the line,
% and is used for comments like this one.

\documentclass{article}      % Specifies the document class

                             % The preamble begins here.
\title{An Example Document}  % Declares the document's title.
\author{Leslie Lamport}      % Declares the author's name.
\date{January 21, 1994}      % Deleting this command produces today's date.

\newcommand{\ip}[2]{(#1, #2)}
                             % Defines \ip{arg1}{arg2} to mean
                             % (arg1, arg2).

%\newcommand{\ip}[2]{\langle #1 | #2\rangle}
                             % This is an alternative definition of
                             % \ip that is commented out.

\begin{document}             % End of preamble and beginning of text.

\maketitle                   % Produces the title.

This is an example input file.  Comparing it with
the output it generates can show you how to
produce a simple document of your own.

\section{Ordinary Text}      % Produces section heading.  Lower-level
                             % sections are begun with similar
                             % \subsection and \subsubsection commands.

The ends  of words and sentences are marked
  by   spaces. It  doesn't matter how many
spaces    you type; one is as good as 100.  The
end of   a line counts as a space.

One   or more   blank lines denote the  end
of  a paragraph.

Since any number of consecutive spaces are treated
like a single one, the formatting of the input
file makes no difference to
      \LaTeX,                % The \LaTeX command generates the LaTeX logo.
but it makes a difference to you.  When you use
\LaTeX, making your input file as easy to read
as possible will be a great help as you write
your document and when you change it.  This sample
file shows how you can add comments to your own input
file.

Because printing is different from typewriting,
there are a number of things that you have to do
differently when preparing an input file than if
you were just typing the document directly.
Quotation marks like
       ``this''
have to be handled specially, as do quotes within
quotes:
       ``\,`this'            % \, separates the double and single quote.
        is what I just
        wrote, not  `that'\,''.

Dashes come in three sizes: an
       intra-word
dash, a medium dash for number ranges like
       1--2,
and a punctuation
       dash---like
this.

A sentence-ending space should be larger than the
space between words within a sentence.  You
sometimes have to type special commands in
conjunction with punctuation characters to get
this right, as in the following sentence.
       Gnats, gnus, etc.\ all  % `\ ' makes an inter-word space.
       begin with G\@.         % \@ marks end-of-sentence punctuation.
You should check the spaces after periods when
reading your output to make sure you haven't
forgotten any special cases.  Generating an
ellipsis
       \ldots\               % `\ ' is needed after `\ldots' because TeX
                             % ignores spaces after command names like \ldots
                             % made from \ + letters.
                             %
                             % Note how a `%' character causes TeX to ignore
                             % the end of the input line, so these blank lines
                             % do not start a new paragraph.
                             %
with the right spacing around the periods requires
a special command.

\LaTeX\ interprets some common characters as
commands, so you must type special commands to
generate them.  These characters include the
following:
       \$ \& \% \# \{ and \}.

In printing, text is usually emphasized with an
       \emph{italic}
type style.

\begin{em}
   A long segment of text can also be emphasized
   in this way.  Text within such a segment can be
   given \emph{additional} emphasis.
\end{em}

It is sometimes necessary to prevent \LaTeX\ from
breaking a line where it might otherwise do so.
This may be at a space, as between the ``Mr.''\ and
``Jones'' in
       ``Mr.~Jones'',        % ~ produces an unbreakable interword space.
or within a word---especially when the word is a
symbol like
       \mbox{\emph{itemnum}}
that makes little sense when hyphenated across
lines.

Footnotes\footnote{This is an example of a footnote.}
pose no problem.

\LaTeX\ is good at typesetting mathematical formulas
like
       \( x-3y + z = 7 \)
or
       \( a_{1} > x^{2n} + y^{2n} > x' \)
or
       \( \ip{A}{B} = \sum_{i} a_{i} b_{i} \).
The spaces you type in a formula are
ignored.  Remember that a letter like
       $x$                   % $ ... $  and  \( ... \)  are equivalent
is a formula when it denotes a mathematical
symbol, and it should be typed as one.

\section{Displayed Text}

Text is displayed by indenting it from the left
margin.  Quotations are commonly displayed.  There
are short quotations
\begin{quote}
   This is a short quotation.  It consists of a
   single paragraph of text.  See how it is formatted.
\end{quote}
and longer ones.
\begin{quotation}
   This is a longer quotation.  It consists of two
   paragraphs of text, neither of which are
   particularly interesting.

   This is the second paragraph of the quotation.  It
   is just as dull as the first paragraph.
\end{quotation}
Another frequently-displayed structure is a list.
The following is an example of an \emph{itemized}
list.
\begin{itemize}
   \item This is the first item of an itemized list.
         Each item in the list is marked with a ``tick''.
         You don't have to worry about what kind of tick
         mark is used.

   \item This is the second item of the list.  It
         contains another list nested inside it.  The inner
         list is an \emph{enumerated} list.
         \begin{enumerate}
            \item This is the first item of an enumerated
                  list that is nested within the itemized list.

            \item This is the second item of the inner list.
                  \LaTeX\ allows you to nest lists deeper than
                  you really should.
         \end{enumerate}
         This is the rest of the second item of the outer
         list.  It is no more interesting than any other
         part of the item.
   \item This is the third item of the list.
\end{itemize}
You can even display poetry.
\begin{verse}
   There is an environment
    for verse \\             % The \\ command separates lines
   Whose features some poets % within a stanza.
   will curse.

                             % One or more blank lines separate stanzas.

   For instead of making\\
   Them do \emph{all} line breaking, \\
   It allows them to put too many words on a line when they'd rather be
   forced to be terse.
\end{verse}

Mathematical formulas may also be displayed.  A
displayed formula
is
one-line long; multiline
formulas require special formatting instructions.
   \[  \ip{\Gamma}{\psi'} = x'' + y^{2} + z_{i}^{n}\]
Don't start a paragraph with a displayed equation,
nor make one a paragraph by itself.

\end{document}               % End of document.

\section{Bench Huge Section 1637}
% synthetic bench section 1637
% This is a sample LaTeX input file.  (Version of 12 August 2004.)
%
% A '%' character causes TeX to ignore all remaining text on the line,
% and is used for comments like this one.

\documentclass{article}      % Specifies the document class

                             % The preamble begins here.
\title{An Example Document}  % Declares the document's title.
\author{Leslie Lamport}      % Declares the author's name.
\date{January 21, 1994}      % Deleting this command produces today's date.

\newcommand{\ip}[2]{(#1, #2)}
                             % Defines \ip{arg1}{arg2} to mean
                             % (arg1, arg2).

%\newcommand{\ip}[2]{\langle #1 | #2\rangle}
                             % This is an alternative definition of
                             % \ip that is commented out.

\begin{document}             % End of preamble and beginning of text.

\maketitle                   % Produces the title.

This is an example input file.  Comparing it with
the output it generates can show you how to
produce a simple document of your own.

\section{Ordinary Text}      % Produces section heading.  Lower-level
                             % sections are begun with similar
                             % \subsection and \subsubsection commands.

The ends  of words and sentences are marked
  by   spaces. It  doesn't matter how many
spaces    you type; one is as good as 100.  The
end of   a line counts as a space.

One   or more   blank lines denote the  end
of  a paragraph.

Since any number of consecutive spaces are treated
like a single one, the formatting of the input
file makes no difference to
      \LaTeX,                % The \LaTeX command generates the LaTeX logo.
but it makes a difference to you.  When you use
\LaTeX, making your input file as easy to read
as possible will be a great help as you write
your document and when you change it.  This sample
file shows how you can add comments to your own input
file.

Because printing is different from typewriting,
there are a number of things that you have to do
differently when preparing an input file than if
you were just typing the document directly.
Quotation marks like
       ``this''
have to be handled specially, as do quotes within
quotes:
       ``\,`this'            % \, separates the double and single quote.
        is what I just
        wrote, not  `that'\,''.

Dashes come in three sizes: an
       intra-word
dash, a medium dash for number ranges like
       1--2,
and a punctuation
       dash---like
this.

A sentence-ending space should be larger than the
space between words within a sentence.  You
sometimes have to type special commands in
conjunction with punctuation characters to get
this right, as in the following sentence.
       Gnats, gnus, etc.\ all  % `\ ' makes an inter-word space.
       begin with G\@.         % \@ marks end-of-sentence punctuation.
You should check the spaces after periods when
reading your output to make sure you haven't
forgotten any special cases.  Generating an
ellipsis
       \ldots\               % `\ ' is needed after `\ldots' because TeX
                             % ignores spaces after command names like \ldots
                             % made from \ + letters.
                             %
                             % Note how a `%' character causes TeX to ignore
                             % the end of the input line, so these blank lines
                             % do not start a new paragraph.
                             %
with the right spacing around the periods requires
a special command.

\LaTeX\ interprets some common characters as
commands, so you must type special commands to
generate them.  These characters include the
following:
       \$ \& \% \# \{ and \}.

In printing, text is usually emphasized with an
       \emph{italic}
type style.

\begin{em}
   A long segment of text can also be emphasized
   in this way.  Text within such a segment can be
   given \emph{additional} emphasis.
\end{em}

It is sometimes necessary to prevent \LaTeX\ from
breaking a line where it might otherwise do so.
This may be at a space, as between the ``Mr.''\ and
``Jones'' in
       ``Mr.~Jones'',        % ~ produces an unbreakable interword space.
or within a word---especially when the word is a
symbol like
       \mbox{\emph{itemnum}}
that makes little sense when hyphenated across
lines.

Footnotes\footnote{This is an example of a footnote.}
pose no problem.

\LaTeX\ is good at typesetting mathematical formulas
like
       \( x-3y + z = 7 \)
or
       \( a_{1} > x^{2n} + y^{2n} > x' \)
or
       \( \ip{A}{B} = \sum_{i} a_{i} b_{i} \).
The spaces you type in a formula are
ignored.  Remember that a letter like
       $x$                   % $ ... $  and  \( ... \)  are equivalent
is a formula when it denotes a mathematical
symbol, and it should be typed as one.

\section{Displayed Text}

Text is displayed by indenting it from the left
margin.  Quotations are commonly displayed.  There
are short quotations
\begin{quote}
   This is a short quotation.  It consists of a
   single paragraph of text.  See how it is formatted.
\end{quote}
and longer ones.
\begin{quotation}
   This is a longer quotation.  It consists of two
   paragraphs of text, neither of which are
   particularly interesting.

   This is the second paragraph of the quotation.  It
   is just as dull as the first paragraph.
\end{quotation}
Another frequently-displayed structure is a list.
The following is an example of an \emph{itemized}
list.
\begin{itemize}
   \item This is the first item of an itemized list.
         Each item in the list is marked with a ``tick''.
         You don't have to worry about what kind of tick
         mark is used.

   \item This is the second item of the list.  It
         contains another list nested inside it.  The inner
         list is an \emph{enumerated} list.
         \begin{enumerate}
            \item This is the first item of an enumerated
                  list that is nested within the itemized list.

            \item This is the second item of the inner list.
                  \LaTeX\ allows you to nest lists deeper than
                  you really should.
         \end{enumerate}
         This is the rest of the second item of the outer
         list.  It is no more interesting than any other
         part of the item.
   \item This is the third item of the list.
\end{itemize}
You can even display poetry.
\begin{verse}
   There is an environment
    for verse \\             % The \\ command separates lines
   Whose features some poets % within a stanza.
   will curse.

                             % One or more blank lines separate stanzas.

   For instead of making\\
   Them do \emph{all} line breaking, \\
   It allows them to put too many words on a line when they'd rather be
   forced to be terse.
\end{verse}

Mathematical formulas may also be displayed.  A
displayed formula
is
one-line long; multiline
formulas require special formatting instructions.
   \[  \ip{\Gamma}{\psi'} = x'' + y^{2} + z_{i}^{n}\]
Don't start a paragraph with a displayed equation,
nor make one a paragraph by itself.

\end{document}               % End of document.

\section{Bench Huge Section 1638}
% synthetic bench section 1638
% This is a sample LaTeX input file.  (Version of 12 August 2004.)
%
% A '%' character causes TeX to ignore all remaining text on the line,
% and is used for comments like this one.

\documentclass{article}      % Specifies the document class

                             % The preamble begins here.
\title{An Example Document}  % Declares the document's title.
\author{Leslie Lamport}      % Declares the author's name.
\date{January 21, 1994}      % Deleting this command produces today's date.

\newcommand{\ip}[2]{(#1, #2)}
                             % Defines \ip{arg1}{arg2} to mean
                             % (arg1, arg2).

%\newcommand{\ip}[2]{\langle #1 | #2\rangle}
                             % This is an alternative definition of
                             % \ip that is commented out.

\begin{document}             % End of preamble and beginning of text.

\maketitle                   % Produces the title.

This is an example input file.  Comparing it with
the output it generates can show you how to
produce a simple document of your own.

\section{Ordinary Text}      % Produces section heading.  Lower-level
                             % sections are begun with similar
                             % \subsection and \subsubsection commands.

The ends  of words and sentences are marked
  by   spaces. It  doesn't matter how many
spaces    you type; one is as good as 100.  The
end of   a line counts as a space.

One   or more   blank lines denote the  end
of  a paragraph.

Since any number of consecutive spaces are treated
like a single one, the formatting of the input
file makes no difference to
      \LaTeX,                % The \LaTeX command generates the LaTeX logo.
but it makes a difference to you.  When you use
\LaTeX, making your input file as easy to read
as possible will be a great help as you write
your document and when you change it.  This sample
file shows how you can add comments to your own input
file.

Because printing is different from typewriting,
there are a number of things that you have to do
differently when preparing an input file than if
you were just typing the document directly.
Quotation marks like
       ``this''
have to be handled specially, as do quotes within
quotes:
       ``\,`this'            % \, separates the double and single quote.
        is what I just
        wrote, not  `that'\,''.

Dashes come in three sizes: an
       intra-word
dash, a medium dash for number ranges like
       1--2,
and a punctuation
       dash---like
this.

A sentence-ending space should be larger than the
space between words within a sentence.  You
sometimes have to type special commands in
conjunction with punctuation characters to get
this right, as in the following sentence.
       Gnats, gnus, etc.\ all  % `\ ' makes an inter-word space.
       begin with G\@.         % \@ marks end-of-sentence punctuation.
You should check the spaces after periods when
reading your output to make sure you haven't
forgotten any special cases.  Generating an
ellipsis
       \ldots\               % `\ ' is needed after `\ldots' because TeX
                             % ignores spaces after command names like \ldots
                             % made from \ + letters.
                             %
                             % Note how a `%' character causes TeX to ignore
                             % the end of the input line, so these blank lines
                             % do not start a new paragraph.
                             %
with the right spacing around the periods requires
a special command.

\LaTeX\ interprets some common characters as
commands, so you must type special commands to
generate them.  These characters include the
following:
       \$ \& \% \# \{ and \}.

In printing, text is usually emphasized with an
       \emph{italic}
type style.

\begin{em}
   A long segment of text can also be emphasized
   in this way.  Text within such a segment can be
   given \emph{additional} emphasis.
\end{em}

It is sometimes necessary to prevent \LaTeX\ from
breaking a line where it might otherwise do so.
This may be at a space, as between the ``Mr.''\ and
``Jones'' in
       ``Mr.~Jones'',        % ~ produces an unbreakable interword space.
or within a word---especially when the word is a
symbol like
       \mbox{\emph{itemnum}}
that makes little sense when hyphenated across
lines.

Footnotes\footnote{This is an example of a footnote.}
pose no problem.

\LaTeX\ is good at typesetting mathematical formulas
like
       \( x-3y + z = 7 \)
or
       \( a_{1} > x^{2n} + y^{2n} > x' \)
or
       \( \ip{A}{B} = \sum_{i} a_{i} b_{i} \).
The spaces you type in a formula are
ignored.  Remember that a letter like
       $x$                   % $ ... $  and  \( ... \)  are equivalent
is a formula when it denotes a mathematical
symbol, and it should be typed as one.

\section{Displayed Text}

Text is displayed by indenting it from the left
margin.  Quotations are commonly displayed.  There
are short quotations
\begin{quote}
   This is a short quotation.  It consists of a
   single paragraph of text.  See how it is formatted.
\end{quote}
and longer ones.
\begin{quotation}
   This is a longer quotation.  It consists of two
   paragraphs of text, neither of which are
   particularly interesting.

   This is the second paragraph of the quotation.  It
   is just as dull as the first paragraph.
\end{quotation}
Another frequently-displayed structure is a list.
The following is an example of an \emph{itemized}
list.
\begin{itemize}
   \item This is the first item of an itemized list.
         Each item in the list is marked with a ``tick''.
         You don't have to worry about what kind of tick
         mark is used.

   \item This is the second item of the list.  It
         contains another list nested inside it.  The inner
         list is an \emph{enumerated} list.
         \begin{enumerate}
            \item This is the first item of an enumerated
                  list that is nested within the itemized list.

            \item This is the second item of the inner list.
                  \LaTeX\ allows you to nest lists deeper than
                  you really should.
         \end{enumerate}
         This is the rest of the second item of the outer
         list.  It is no more interesting than any other
         part of the item.
   \item This is the third item of the list.
\end{itemize}
You can even display poetry.
\begin{verse}
   There is an environment
    for verse \\             % The \\ command separates lines
   Whose features some poets % within a stanza.
   will curse.

                             % One or more blank lines separate stanzas.

   For instead of making\\
   Them do \emph{all} line breaking, \\
   It allows them to put too many words on a line when they'd rather be
   forced to be terse.
\end{verse}

Mathematical formulas may also be displayed.  A
displayed formula
is
one-line long; multiline
formulas require special formatting instructions.
   \[  \ip{\Gamma}{\psi'} = x'' + y^{2} + z_{i}^{n}\]
Don't start a paragraph with a displayed equation,
nor make one a paragraph by itself.

\end{document}               % End of document.

\section{Bench Huge Section 1639}
% synthetic bench section 1639
% This is a sample LaTeX input file.  (Version of 12 August 2004.)
%
% A '%' character causes TeX to ignore all remaining text on the line,
% and is used for comments like this one.

\documentclass{article}      % Specifies the document class

                             % The preamble begins here.
\title{An Example Document}  % Declares the document's title.
\author{Leslie Lamport}      % Declares the author's name.
\date{January 21, 1994}      % Deleting this command produces today's date.

\newcommand{\ip}[2]{(#1, #2)}
                             % Defines \ip{arg1}{arg2} to mean
                             % (arg1, arg2).

%\newcommand{\ip}[2]{\langle #1 | #2\rangle}
                             % This is an alternative definition of
                             % \ip that is commented out.

\begin{document}             % End of preamble and beginning of text.

\maketitle                   % Produces the title.

This is an example input file.  Comparing it with
the output it generates can show you how to
produce a simple document of your own.

\section{Ordinary Text}      % Produces section heading.  Lower-level
                             % sections are begun with similar
                             % \subsection and \subsubsection commands.

The ends  of words and sentences are marked
  by   spaces. It  doesn't matter how many
spaces    you type; one is as good as 100.  The
end of   a line counts as a space.

One   or more   blank lines denote the  end
of  a paragraph.

Since any number of consecutive spaces are treated
like a single one, the formatting of the input
file makes no difference to
      \LaTeX,                % The \LaTeX command generates the LaTeX logo.
but it makes a difference to you.  When you use
\LaTeX, making your input file as easy to read
as possible will be a great help as you write
your document and when you change it.  This sample
file shows how you can add comments to your own input
file.

Because printing is different from typewriting,
there are a number of things that you have to do
differently when preparing an input file than if
you were just typing the document directly.
Quotation marks like
       ``this''
have to be handled specially, as do quotes within
quotes:
       ``\,`this'            % \, separates the double and single quote.
        is what I just
        wrote, not  `that'\,''.

Dashes come in three sizes: an
       intra-word
dash, a medium dash for number ranges like
       1--2,
and a punctuation
       dash---like
this.

A sentence-ending space should be larger than the
space between words within a sentence.  You
sometimes have to type special commands in
conjunction with punctuation characters to get
this right, as in the following sentence.
       Gnats, gnus, etc.\ all  % `\ ' makes an inter-word space.
       begin with G\@.         % \@ marks end-of-sentence punctuation.
You should check the spaces after periods when
reading your output to make sure you haven't
forgotten any special cases.  Generating an
ellipsis
       \ldots\               % `\ ' is needed after `\ldots' because TeX
                             % ignores spaces after command names like \ldots
                             % made from \ + letters.
                             %
                             % Note how a `%' character causes TeX to ignore
                             % the end of the input line, so these blank lines
                             % do not start a new paragraph.
                             %
with the right spacing around the periods requires
a special command.

\LaTeX\ interprets some common characters as
commands, so you must type special commands to
generate them.  These characters include the
following:
       \$ \& \% \# \{ and \}.

In printing, text is usually emphasized with an
       \emph{italic}
type style.

\begin{em}
   A long segment of text can also be emphasized
   in this way.  Text within such a segment can be
   given \emph{additional} emphasis.
\end{em}

It is sometimes necessary to prevent \LaTeX\ from
breaking a line where it might otherwise do so.
This may be at a space, as between the ``Mr.''\ and
``Jones'' in
       ``Mr.~Jones'',        % ~ produces an unbreakable interword space.
or within a word---especially when the word is a
symbol like
       \mbox{\emph{itemnum}}
that makes little sense when hyphenated across
lines.

Footnotes\footnote{This is an example of a footnote.}
pose no problem.

\LaTeX\ is good at typesetting mathematical formulas
like
       \( x-3y + z = 7 \)
or
       \( a_{1} > x^{2n} + y^{2n} > x' \)
or
       \( \ip{A}{B} = \sum_{i} a_{i} b_{i} \).
The spaces you type in a formula are
ignored.  Remember that a letter like
       $x$                   % $ ... $  and  \( ... \)  are equivalent
is a formula when it denotes a mathematical
symbol, and it should be typed as one.

\section{Displayed Text}

Text is displayed by indenting it from the left
margin.  Quotations are commonly displayed.  There
are short quotations
\begin{quote}
   This is a short quotation.  It consists of a
   single paragraph of text.  See how it is formatted.
\end{quote}
and longer ones.
\begin{quotation}
   This is a longer quotation.  It consists of two
   paragraphs of text, neither of which are
   particularly interesting.

   This is the second paragraph of the quotation.  It
   is just as dull as the first paragraph.
\end{quotation}
Another frequently-displayed structure is a list.
The following is an example of an \emph{itemized}
list.
\begin{itemize}
   \item This is the first item of an itemized list.
         Each item in the list is marked with a ``tick''.
         You don't have to worry about what kind of tick
         mark is used.

   \item This is the second item of the list.  It
         contains another list nested inside it.  The inner
         list is an \emph{enumerated} list.
         \begin{enumerate}
            \item This is the first item of an enumerated
                  list that is nested within the itemized list.

            \item This is the second item of the inner list.
                  \LaTeX\ allows you to nest lists deeper than
                  you really should.
         \end{enumerate}
         This is the rest of the second item of the outer
         list.  It is no more interesting than any other
         part of the item.
   \item This is the third item of the list.
\end{itemize}
You can even display poetry.
\begin{verse}
   There is an environment
    for verse \\             % The \\ command separates lines
   Whose features some poets % within a stanza.
   will curse.

                             % One or more blank lines separate stanzas.

   For instead of making\\
   Them do \emph{all} line breaking, \\
   It allows them to put too many words on a line when they'd rather be
   forced to be terse.
\end{verse}

Mathematical formulas may also be displayed.  A
displayed formula
is
one-line long; multiline
formulas require special formatting instructions.
   \[  \ip{\Gamma}{\psi'} = x'' + y^{2} + z_{i}^{n}\]
Don't start a paragraph with a displayed equation,
nor make one a paragraph by itself.

\end{document}               % End of document.

\section{Bench Huge Section 1640}
% synthetic bench section 1640
% This is a sample LaTeX input file.  (Version of 12 August 2004.)
%
% A '%' character causes TeX to ignore all remaining text on the line,
% and is used for comments like this one.

\documentclass{article}      % Specifies the document class

                             % The preamble begins here.
\title{An Example Document}  % Declares the document's title.
\author{Leslie Lamport}      % Declares the author's name.
\date{January 21, 1994}      % Deleting this command produces today's date.

\newcommand{\ip}[2]{(#1, #2)}
                             % Defines \ip{arg1}{arg2} to mean
                             % (arg1, arg2).

%\newcommand{\ip}[2]{\langle #1 | #2\rangle}
                             % This is an alternative definition of
                             % \ip that is commented out.

\begin{document}             % End of preamble and beginning of text.

\maketitle                   % Produces the title.

This is an example input file.  Comparing it with
the output it generates can show you how to
produce a simple document of your own.

\section{Ordinary Text}      % Produces section heading.  Lower-level
                             % sections are begun with similar
                             % \subsection and \subsubsection commands.

The ends  of words and sentences are marked
  by   spaces. It  doesn't matter how many
spaces    you type; one is as good as 100.  The
end of   a line counts as a space.

One   or more   blank lines denote the  end
of  a paragraph.

Since any number of consecutive spaces are treated
like a single one, the formatting of the input
file makes no difference to
      \LaTeX,                % The \LaTeX command generates the LaTeX logo.
but it makes a difference to you.  When you use
\LaTeX, making your input file as easy to read
as possible will be a great help as you write
your document and when you change it.  This sample
file shows how you can add comments to your own input
file.

Because printing is different from typewriting,
there are a number of things that you have to do
differently when preparing an input file than if
you were just typing the document directly.
Quotation marks like
       ``this''
have to be handled specially, as do quotes within
quotes:
       ``\,`this'            % \, separates the double and single quote.
        is what I just
        wrote, not  `that'\,''.

Dashes come in three sizes: an
       intra-word
dash, a medium dash for number ranges like
       1--2,
and a punctuation
       dash---like
this.

A sentence-ending space should be larger than the
space between words within a sentence.  You
sometimes have to type special commands in
conjunction with punctuation characters to get
this right, as in the following sentence.
       Gnats, gnus, etc.\ all  % `\ ' makes an inter-word space.
       begin with G\@.         % \@ marks end-of-sentence punctuation.
You should check the spaces after periods when
reading your output to make sure you haven't
forgotten any special cases.  Generating an
ellipsis
       \ldots\               % `\ ' is needed after `\ldots' because TeX
                             % ignores spaces after command names like \ldots
                             % made from \ + letters.
                             %
                             % Note how a `%' character causes TeX to ignore
                             % the end of the input line, so these blank lines
                             % do not start a new paragraph.
                             %
with the right spacing around the periods requires
a special command.

\LaTeX\ interprets some common characters as
commands, so you must type special commands to
generate them.  These characters include the
following:
       \$ \& \% \# \{ and \}.

In printing, text is usually emphasized with an
       \emph{italic}
type style.

\begin{em}
   A long segment of text can also be emphasized
   in this way.  Text within such a segment can be
   given \emph{additional} emphasis.
\end{em}

It is sometimes necessary to prevent \LaTeX\ from
breaking a line where it might otherwise do so.
This may be at a space, as between the ``Mr.''\ and
``Jones'' in
       ``Mr.~Jones'',        % ~ produces an unbreakable interword space.
or within a word---especially when the word is a
symbol like
       \mbox{\emph{itemnum}}
that makes little sense when hyphenated across
lines.

Footnotes\footnote{This is an example of a footnote.}
pose no problem.

\LaTeX\ is good at typesetting mathematical formulas
like
       \( x-3y + z = 7 \)
or
       \( a_{1} > x^{2n} + y^{2n} > x' \)
or
       \( \ip{A}{B} = \sum_{i} a_{i} b_{i} \).
The spaces you type in a formula are
ignored.  Remember that a letter like
       $x$                   % $ ... $  and  \( ... \)  are equivalent
is a formula when it denotes a mathematical
symbol, and it should be typed as one.

\section{Displayed Text}

Text is displayed by indenting it from the left
margin.  Quotations are commonly displayed.  There
are short quotations
\begin{quote}
   This is a short quotation.  It consists of a
   single paragraph of text.  See how it is formatted.
\end{quote}
and longer ones.
\begin{quotation}
   This is a longer quotation.  It consists of two
   paragraphs of text, neither of which are
   particularly interesting.

   This is the second paragraph of the quotation.  It
   is just as dull as the first paragraph.
\end{quotation}
Another frequently-displayed structure is a list.
The following is an example of an \emph{itemized}
list.
\begin{itemize}
   \item This is the first item of an itemized list.
         Each item in the list is marked with a ``tick''.
         You don't have to worry about what kind of tick
         mark is used.

   \item This is the second item of the list.  It
         contains another list nested inside it.  The inner
         list is an \emph{enumerated} list.
         \begin{enumerate}
            \item This is the first item of an enumerated
                  list that is nested within the itemized list.

            \item This is the second item of the inner list.
                  \LaTeX\ allows you to nest lists deeper than
                  you really should.
         \end{enumerate}
         This is the rest of the second item of the outer
         list.  It is no more interesting than any other
         part of the item.
   \item This is the third item of the list.
\end{itemize}
You can even display poetry.
\begin{verse}
   There is an environment
    for verse \\             % The \\ command separates lines
   Whose features some poets % within a stanza.
   will curse.

                             % One or more blank lines separate stanzas.

   For instead of making\\
   Them do \emph{all} line breaking, \\
   It allows them to put too many words on a line when they'd rather be
   forced to be terse.
\end{verse}

Mathematical formulas may also be displayed.  A
displayed formula
is
one-line long; multiline
formulas require special formatting instructions.
   \[  \ip{\Gamma}{\psi'} = x'' + y^{2} + z_{i}^{n}\]
Don't start a paragraph with a displayed equation,
nor make one a paragraph by itself.

\end{document}               % End of document.

\section{Bench Huge Section 1641}
% synthetic bench section 1641
% This is a sample LaTeX input file.  (Version of 12 August 2004.)
%
% A '%' character causes TeX to ignore all remaining text on the line,
% and is used for comments like this one.

\documentclass{article}      % Specifies the document class

                             % The preamble begins here.
\title{An Example Document}  % Declares the document's title.
\author{Leslie Lamport}      % Declares the author's name.
\date{January 21, 1994}      % Deleting this command produces today's date.

\newcommand{\ip}[2]{(#1, #2)}
                             % Defines \ip{arg1}{arg2} to mean
                             % (arg1, arg2).

%\newcommand{\ip}[2]{\langle #1 | #2\rangle}
                             % This is an alternative definition of
                             % \ip that is commented out.

\begin{document}             % End of preamble and beginning of text.

\maketitle                   % Produces the title.

This is an example input file.  Comparing it with
the output it generates can show you how to
produce a simple document of your own.

\section{Ordinary Text}      % Produces section heading.  Lower-level
                             % sections are begun with similar
                             % \subsection and \subsubsection commands.

The ends  of words and sentences are marked
  by   spaces. It  doesn't matter how many
spaces    you type; one is as good as 100.  The
end of   a line counts as a space.

One   or more   blank lines denote the  end
of  a paragraph.

Since any number of consecutive spaces are treated
like a single one, the formatting of the input
file makes no difference to
      \LaTeX,                % The \LaTeX command generates the LaTeX logo.
but it makes a difference to you.  When you use
\LaTeX, making your input file as easy to read
as possible will be a great help as you write
your document and when you change it.  This sample
file shows how you can add comments to your own input
file.

Because printing is different from typewriting,
there are a number of things that you have to do
differently when preparing an input file than if
you were just typing the document directly.
Quotation marks like
       ``this''
have to be handled specially, as do quotes within
quotes:
       ``\,`this'            % \, separates the double and single quote.
        is what I just
        wrote, not  `that'\,''.

Dashes come in three sizes: an
       intra-word
dash, a medium dash for number ranges like
       1--2,
and a punctuation
       dash---like
this.

A sentence-ending space should be larger than the
space between words within a sentence.  You
sometimes have to type special commands in
conjunction with punctuation characters to get
this right, as in the following sentence.
       Gnats, gnus, etc.\ all  % `\ ' makes an inter-word space.
       begin with G\@.         % \@ marks end-of-sentence punctuation.
You should check the spaces after periods when
reading your output to make sure you haven't
forgotten any special cases.  Generating an
ellipsis
       \ldots\               % `\ ' is needed after `\ldots' because TeX
                             % ignores spaces after command names like \ldots
                             % made from \ + letters.
                             %
                             % Note how a `%' character causes TeX to ignore
                             % the end of the input line, so these blank lines
                             % do not start a new paragraph.
                             %
with the right spacing around the periods requires
a special command.

\LaTeX\ interprets some common characters as
commands, so you must type special commands to
generate them.  These characters include the
following:
       \$ \& \% \# \{ and \}.

In printing, text is usually emphasized with an
       \emph{italic}
type style.

\begin{em}
   A long segment of text can also be emphasized
   in this way.  Text within such a segment can be
   given \emph{additional} emphasis.
\end{em}

It is sometimes necessary to prevent \LaTeX\ from
breaking a line where it might otherwise do so.
This may be at a space, as between the ``Mr.''\ and
``Jones'' in
       ``Mr.~Jones'',        % ~ produces an unbreakable interword space.
or within a word---especially when the word is a
symbol like
       \mbox{\emph{itemnum}}
that makes little sense when hyphenated across
lines.

Footnotes\footnote{This is an example of a footnote.}
pose no problem.

\LaTeX\ is good at typesetting mathematical formulas
like
       \( x-3y + z = 7 \)
or
       \( a_{1} > x^{2n} + y^{2n} > x' \)
or
       \( \ip{A}{B} = \sum_{i} a_{i} b_{i} \).
The spaces you type in a formula are
ignored.  Remember that a letter like
       $x$                   % $ ... $  and  \( ... \)  are equivalent
is a formula when it denotes a mathematical
symbol, and it should be typed as one.

\section{Displayed Text}

Text is displayed by indenting it from the left
margin.  Quotations are commonly displayed.  There
are short quotations
\begin{quote}
   This is a short quotation.  It consists of a
   single paragraph of text.  See how it is formatted.
\end{quote}
and longer ones.
\begin{quotation}
   This is a longer quotation.  It consists of two
   paragraphs of text, neither of which are
   particularly interesting.

   This is the second paragraph of the quotation.  It
   is just as dull as the first paragraph.
\end{quotation}
Another frequently-displayed structure is a list.
The following is an example of an \emph{itemized}
list.
\begin{itemize}
   \item This is the first item of an itemized list.
         Each item in the list is marked with a ``tick''.
         You don't have to worry about what kind of tick
         mark is used.

   \item This is the second item of the list.  It
         contains another list nested inside it.  The inner
         list is an \emph{enumerated} list.
         \begin{enumerate}
            \item This is the first item of an enumerated
                  list that is nested within the itemized list.

            \item This is the second item of the inner list.
                  \LaTeX\ allows you to nest lists deeper than
                  you really should.
         \end{enumerate}
         This is the rest of the second item of the outer
         list.  It is no more interesting than any other
         part of the item.
   \item This is the third item of the list.
\end{itemize}
You can even display poetry.
\begin{verse}
   There is an environment
    for verse \\             % The \\ command separates lines
   Whose features some poets % within a stanza.
   will curse.

                             % One or more blank lines separate stanzas.

   For instead of making\\
   Them do \emph{all} line breaking, \\
   It allows them to put too many words on a line when they'd rather be
   forced to be terse.
\end{verse}

Mathematical formulas may also be displayed.  A
displayed formula
is
one-line long; multiline
formulas require special formatting instructions.
   \[  \ip{\Gamma}{\psi'} = x'' + y^{2} + z_{i}^{n}\]
Don't start a paragraph with a displayed equation,
nor make one a paragraph by itself.

\end{document}               % End of document.

\section{Bench Huge Section 1642}
% synthetic bench section 1642
% This is a sample LaTeX input file.  (Version of 12 August 2004.)
%
% A '%' character causes TeX to ignore all remaining text on the line,
% and is used for comments like this one.

\documentclass{article}      % Specifies the document class

                             % The preamble begins here.
\title{An Example Document}  % Declares the document's title.
\author{Leslie Lamport}      % Declares the author's name.
\date{January 21, 1994}      % Deleting this command produces today's date.

\newcommand{\ip}[2]{(#1, #2)}
                             % Defines \ip{arg1}{arg2} to mean
                             % (arg1, arg2).

%\newcommand{\ip}[2]{\langle #1 | #2\rangle}
                             % This is an alternative definition of
                             % \ip that is commented out.

\begin{document}             % End of preamble and beginning of text.

\maketitle                   % Produces the title.

This is an example input file.  Comparing it with
the output it generates can show you how to
produce a simple document of your own.

\section{Ordinary Text}      % Produces section heading.  Lower-level
                             % sections are begun with similar
                             % \subsection and \subsubsection commands.

The ends  of words and sentences are marked
  by   spaces. It  doesn't matter how many
spaces    you type; one is as good as 100.  The
end of   a line counts as a space.

One   or more   blank lines denote the  end
of  a paragraph.

Since any number of consecutive spaces are treated
like a single one, the formatting of the input
file makes no difference to
      \LaTeX,                % The \LaTeX command generates the LaTeX logo.
but it makes a difference to you.  When you use
\LaTeX, making your input file as easy to read
as possible will be a great help as you write
your document and when you change it.  This sample
file shows how you can add comments to your own input
file.

Because printing is different from typewriting,
there are a number of things that you have to do
differently when preparing an input file than if
you were just typing the document directly.
Quotation marks like
       ``this''
have to be handled specially, as do quotes within
quotes:
       ``\,`this'            % \, separates the double and single quote.
        is what I just
        wrote, not  `that'\,''.

Dashes come in three sizes: an
       intra-word
dash, a medium dash for number ranges like
       1--2,
and a punctuation
       dash---like
this.

A sentence-ending space should be larger than the
space between words within a sentence.  You
sometimes have to type special commands in
conjunction with punctuation characters to get
this right, as in the following sentence.
       Gnats, gnus, etc.\ all  % `\ ' makes an inter-word space.
       begin with G\@.         % \@ marks end-of-sentence punctuation.
You should check the spaces after periods when
reading your output to make sure you haven't
forgotten any special cases.  Generating an
ellipsis
       \ldots\               % `\ ' is needed after `\ldots' because TeX
                             % ignores spaces after command names like \ldots
                             % made from \ + letters.
                             %
                             % Note how a `%' character causes TeX to ignore
                             % the end of the input line, so these blank lines
                             % do not start a new paragraph.
                             %
with the right spacing around the periods requires
a special command.

\LaTeX\ interprets some common characters as
commands, so you must type special commands to
generate them.  These characters include the
following:
       \$ \& \% \# \{ and \}.

In printing, text is usually emphasized with an
       \emph{italic}
type style.

\begin{em}
   A long segment of text can also be emphasized
   in this way.  Text within such a segment can be
   given \emph{additional} emphasis.
\end{em}

It is sometimes necessary to prevent \LaTeX\ from
breaking a line where it might otherwise do so.
This may be at a space, as between the ``Mr.''\ and
``Jones'' in
       ``Mr.~Jones'',        % ~ produces an unbreakable interword space.
or within a word---especially when the word is a
symbol like
       \mbox{\emph{itemnum}}
that makes little sense when hyphenated across
lines.

Footnotes\footnote{This is an example of a footnote.}
pose no problem.

\LaTeX\ is good at typesetting mathematical formulas
like
       \( x-3y + z = 7 \)
or
       \( a_{1} > x^{2n} + y^{2n} > x' \)
or
       \( \ip{A}{B} = \sum_{i} a_{i} b_{i} \).
The spaces you type in a formula are
ignored.  Remember that a letter like
       $x$                   % $ ... $  and  \( ... \)  are equivalent
is a formula when it denotes a mathematical
symbol, and it should be typed as one.

\section{Displayed Text}

Text is displayed by indenting it from the left
margin.  Quotations are commonly displayed.  There
are short quotations
\begin{quote}
   This is a short quotation.  It consists of a
   single paragraph of text.  See how it is formatted.
\end{quote}
and longer ones.
\begin{quotation}
   This is a longer quotation.  It consists of two
   paragraphs of text, neither of which are
   particularly interesting.

   This is the second paragraph of the quotation.  It
   is just as dull as the first paragraph.
\end{quotation}
Another frequently-displayed structure is a list.
The following is an example of an \emph{itemized}
list.
\begin{itemize}
   \item This is the first item of an itemized list.
         Each item in the list is marked with a ``tick''.
         You don't have to worry about what kind of tick
         mark is used.

   \item This is the second item of the list.  It
         contains another list nested inside it.  The inner
         list is an \emph{enumerated} list.
         \begin{enumerate}
            \item This is the first item of an enumerated
                  list that is nested within the itemized list.

            \item This is the second item of the inner list.
                  \LaTeX\ allows you to nest lists deeper than
                  you really should.
         \end{enumerate}
         This is the rest of the second item of the outer
         list.  It is no more interesting than any other
         part of the item.
   \item This is the third item of the list.
\end{itemize}
You can even display poetry.
\begin{verse}
   There is an environment
    for verse \\             % The \\ command separates lines
   Whose features some poets % within a stanza.
   will curse.

                             % One or more blank lines separate stanzas.

   For instead of making\\
   Them do \emph{all} line breaking, \\
   It allows them to put too many words on a line when they'd rather be
   forced to be terse.
\end{verse}

Mathematical formulas may also be displayed.  A
displayed formula
is
one-line long; multiline
formulas require special formatting instructions.
   \[  \ip{\Gamma}{\psi'} = x'' + y^{2} + z_{i}^{n}\]
Don't start a paragraph with a displayed equation,
nor make one a paragraph by itself.

\end{document}               % End of document.

\section{Bench Huge Section 1643}
% synthetic bench section 1643
% This is a sample LaTeX input file.  (Version of 12 August 2004.)
%
% A '%' character causes TeX to ignore all remaining text on the line,
% and is used for comments like this one.

\documentclass{article}      % Specifies the document class

                             % The preamble begins here.
\title{An Example Document}  % Declares the document's title.
\author{Leslie Lamport}      % Declares the author's name.
\date{January 21, 1994}      % Deleting this command produces today's date.

\newcommand{\ip}[2]{(#1, #2)}
                             % Defines \ip{arg1}{arg2} to mean
                             % (arg1, arg2).

%\newcommand{\ip}[2]{\langle #1 | #2\rangle}
                             % This is an alternative definition of
                             % \ip that is commented out.

\begin{document}             % End of preamble and beginning of text.

\maketitle                   % Produces the title.

This is an example input file.  Comparing it with
the output it generates can show you how to
produce a simple document of your own.

\section{Ordinary Text}      % Produces section heading.  Lower-level
                             % sections are begun with similar
                             % \subsection and \subsubsection commands.

The ends  of words and sentences are marked
  by   spaces. It  doesn't matter how many
spaces    you type; one is as good as 100.  The
end of   a line counts as a space.

One   or more   blank lines denote the  end
of  a paragraph.

Since any number of consecutive spaces are treated
like a single one, the formatting of the input
file makes no difference to
      \LaTeX,                % The \LaTeX command generates the LaTeX logo.
but it makes a difference to you.  When you use
\LaTeX, making your input file as easy to read
as possible will be a great help as you write
your document and when you change it.  This sample
file shows how you can add comments to your own input
file.

Because printing is different from typewriting,
there are a number of things that you have to do
differently when preparing an input file than if
you were just typing the document directly.
Quotation marks like
       ``this''
have to be handled specially, as do quotes within
quotes:
       ``\,`this'            % \, separates the double and single quote.
        is what I just
        wrote, not  `that'\,''.

Dashes come in three sizes: an
       intra-word
dash, a medium dash for number ranges like
       1--2,
and a punctuation
       dash---like
this.

A sentence-ending space should be larger than the
space between words within a sentence.  You
sometimes have to type special commands in
conjunction with punctuation characters to get
this right, as in the following sentence.
       Gnats, gnus, etc.\ all  % `\ ' makes an inter-word space.
       begin with G\@.         % \@ marks end-of-sentence punctuation.
You should check the spaces after periods when
reading your output to make sure you haven't
forgotten any special cases.  Generating an
ellipsis
       \ldots\               % `\ ' is needed after `\ldots' because TeX
                             % ignores spaces after command names like \ldots
                             % made from \ + letters.
                             %
                             % Note how a `%' character causes TeX to ignore
                             % the end of the input line, so these blank lines
                             % do not start a new paragraph.
                             %
with the right spacing around the periods requires
a special command.

\LaTeX\ interprets some common characters as
commands, so you must type special commands to
generate them.  These characters include the
following:
       \$ \& \% \# \{ and \}.

In printing, text is usually emphasized with an
       \emph{italic}
type style.

\begin{em}
   A long segment of text can also be emphasized
   in this way.  Text within such a segment can be
   given \emph{additional} emphasis.
\end{em}

It is sometimes necessary to prevent \LaTeX\ from
breaking a line where it might otherwise do so.
This may be at a space, as between the ``Mr.''\ and
``Jones'' in
       ``Mr.~Jones'',        % ~ produces an unbreakable interword space.
or within a word---especially when the word is a
symbol like
       \mbox{\emph{itemnum}}
that makes little sense when hyphenated across
lines.

Footnotes\footnote{This is an example of a footnote.}
pose no problem.

\LaTeX\ is good at typesetting mathematical formulas
like
       \( x-3y + z = 7 \)
or
       \( a_{1} > x^{2n} + y^{2n} > x' \)
or
       \( \ip{A}{B} = \sum_{i} a_{i} b_{i} \).
The spaces you type in a formula are
ignored.  Remember that a letter like
       $x$                   % $ ... $  and  \( ... \)  are equivalent
is a formula when it denotes a mathematical
symbol, and it should be typed as one.

\section{Displayed Text}

Text is displayed by indenting it from the left
margin.  Quotations are commonly displayed.  There
are short quotations
\begin{quote}
   This is a short quotation.  It consists of a
   single paragraph of text.  See how it is formatted.
\end{quote}
and longer ones.
\begin{quotation}
   This is a longer quotation.  It consists of two
   paragraphs of text, neither of which are
   particularly interesting.

   This is the second paragraph of the quotation.  It
   is just as dull as the first paragraph.
\end{quotation}
Another frequently-displayed structure is a list.
The following is an example of an \emph{itemized}
list.
\begin{itemize}
   \item This is the first item of an itemized list.
         Each item in the list is marked with a ``tick''.
         You don't have to worry about what kind of tick
         mark is used.

   \item This is the second item of the list.  It
         contains another list nested inside it.  The inner
         list is an \emph{enumerated} list.
         \begin{enumerate}
            \item This is the first item of an enumerated
                  list that is nested within the itemized list.

            \item This is the second item of the inner list.
                  \LaTeX\ allows you to nest lists deeper than
                  you really should.
         \end{enumerate}
         This is the rest of the second item of the outer
         list.  It is no more interesting than any other
         part of the item.
   \item This is the third item of the list.
\end{itemize}
You can even display poetry.
\begin{verse}
   There is an environment
    for verse \\             % The \\ command separates lines
   Whose features some poets % within a stanza.
   will curse.

                             % One or more blank lines separate stanzas.

   For instead of making\\
   Them do \emph{all} line breaking, \\
   It allows them to put too many words on a line when they'd rather be
   forced to be terse.
\end{verse}

Mathematical formulas may also be displayed.  A
displayed formula
is
one-line long; multiline
formulas require special formatting instructions.
   \[  \ip{\Gamma}{\psi'} = x'' + y^{2} + z_{i}^{n}\]
Don't start a paragraph with a displayed equation,
nor make one a paragraph by itself.

\end{document}               % End of document.

\section{Bench Huge Section 1644}
% synthetic bench section 1644
% This is a sample LaTeX input file.  (Version of 12 August 2004.)
%
% A '%' character causes TeX to ignore all remaining text on the line,
% and is used for comments like this one.

\documentclass{article}      % Specifies the document class

                             % The preamble begins here.
\title{An Example Document}  % Declares the document's title.
\author{Leslie Lamport}      % Declares the author's name.
\date{January 21, 1994}      % Deleting this command produces today's date.

\newcommand{\ip}[2]{(#1, #2)}
                             % Defines \ip{arg1}{arg2} to mean
                             % (arg1, arg2).

%\newcommand{\ip}[2]{\langle #1 | #2\rangle}
                             % This is an alternative definition of
                             % \ip that is commented out.

\begin{document}             % End of preamble and beginning of text.

\maketitle                   % Produces the title.

This is an example input file.  Comparing it with
the output it generates can show you how to
produce a simple document of your own.

\section{Ordinary Text}      % Produces section heading.  Lower-level
                             % sections are begun with similar
                             % \subsection and \subsubsection commands.

The ends  of words and sentences are marked
  by   spaces. It  doesn't matter how many
spaces    you type; one is as good as 100.  The
end of   a line counts as a space.

One   or more   blank lines denote the  end
of  a paragraph.

Since any number of consecutive spaces are treated
like a single one, the formatting of the input
file makes no difference to
      \LaTeX,                % The \LaTeX command generates the LaTeX logo.
but it makes a difference to you.  When you use
\LaTeX, making your input file as easy to read
as possible will be a great help as you write
your document and when you change it.  This sample
file shows how you can add comments to your own input
file.

Because printing is different from typewriting,
there are a number of things that you have to do
differently when preparing an input file than if
you were just typing the document directly.
Quotation marks like
       ``this''
have to be handled specially, as do quotes within
quotes:
       ``\,`this'            % \, separates the double and single quote.
        is what I just
        wrote, not  `that'\,''.

Dashes come in three sizes: an
       intra-word
dash, a medium dash for number ranges like
       1--2,
and a punctuation
       dash---like
this.

A sentence-ending space should be larger than the
space between words within a sentence.  You
sometimes have to type special commands in
conjunction with punctuation characters to get
this right, as in the following sentence.
       Gnats, gnus, etc.\ all  % `\ ' makes an inter-word space.
       begin with G\@.         % \@ marks end-of-sentence punctuation.
You should check the spaces after periods when
reading your output to make sure you haven't
forgotten any special cases.  Generating an
ellipsis
       \ldots\               % `\ ' is needed after `\ldots' because TeX
                             % ignores spaces after command names like \ldots
                             % made from \ + letters.
                             %
                             % Note how a `%' character causes TeX to ignore
                             % the end of the input line, so these blank lines
                             % do not start a new paragraph.
                             %
with the right spacing around the periods requires
a special command.

\LaTeX\ interprets some common characters as
commands, so you must type special commands to
generate them.  These characters include the
following:
       \$ \& \% \# \{ and \}.

In printing, text is usually emphasized with an
       \emph{italic}
type style.

\begin{em}
   A long segment of text can also be emphasized
   in this way.  Text within such a segment can be
   given \emph{additional} emphasis.
\end{em}

It is sometimes necessary to prevent \LaTeX\ from
breaking a line where it might otherwise do so.
This may be at a space, as between the ``Mr.''\ and
``Jones'' in
       ``Mr.~Jones'',        % ~ produces an unbreakable interword space.
or within a word---especially when the word is a
symbol like
       \mbox{\emph{itemnum}}
that makes little sense when hyphenated across
lines.

Footnotes\footnote{This is an example of a footnote.}
pose no problem.

\LaTeX\ is good at typesetting mathematical formulas
like
       \( x-3y + z = 7 \)
or
       \( a_{1} > x^{2n} + y^{2n} > x' \)
or
       \( \ip{A}{B} = \sum_{i} a_{i} b_{i} \).
The spaces you type in a formula are
ignored.  Remember that a letter like
       $x$                   % $ ... $  and  \( ... \)  are equivalent
is a formula when it denotes a mathematical
symbol, and it should be typed as one.

\section{Displayed Text}

Text is displayed by indenting it from the left
margin.  Quotations are commonly displayed.  There
are short quotations
\begin{quote}
   This is a short quotation.  It consists of a
   single paragraph of text.  See how it is formatted.
\end{quote}
and longer ones.
\begin{quotation}
   This is a longer quotation.  It consists of two
   paragraphs of text, neither of which are
   particularly interesting.

   This is the second paragraph of the quotation.  It
   is just as dull as the first paragraph.
\end{quotation}
Another frequently-displayed structure is a list.
The following is an example of an \emph{itemized}
list.
\begin{itemize}
   \item This is the first item of an itemized list.
         Each item in the list is marked with a ``tick''.
         You don't have to worry about what kind of tick
         mark is used.

   \item This is the second item of the list.  It
         contains another list nested inside it.  The inner
         list is an \emph{enumerated} list.
         \begin{enumerate}
            \item This is the first item of an enumerated
                  list that is nested within the itemized list.

            \item This is the second item of the inner list.
                  \LaTeX\ allows you to nest lists deeper than
                  you really should.
         \end{enumerate}
         This is the rest of the second item of the outer
         list.  It is no more interesting than any other
         part of the item.
   \item This is the third item of the list.
\end{itemize}
You can even display poetry.
\begin{verse}
   There is an environment
    for verse \\             % The \\ command separates lines
   Whose features some poets % within a stanza.
   will curse.

                             % One or more blank lines separate stanzas.

   For instead of making\\
   Them do \emph{all} line breaking, \\
   It allows them to put too many words on a line when they'd rather be
   forced to be terse.
\end{verse}

Mathematical formulas may also be displayed.  A
displayed formula
is
one-line long; multiline
formulas require special formatting instructions.
   \[  \ip{\Gamma}{\psi'} = x'' + y^{2} + z_{i}^{n}\]
Don't start a paragraph with a displayed equation,
nor make one a paragraph by itself.

\end{document}               % End of document.

\section{Bench Huge Section 1645}
% synthetic bench section 1645
% This is a sample LaTeX input file.  (Version of 12 August 2004.)
%
% A '%' character causes TeX to ignore all remaining text on the line,
% and is used for comments like this one.

\documentclass{article}      % Specifies the document class

                             % The preamble begins here.
\title{An Example Document}  % Declares the document's title.
\author{Leslie Lamport}      % Declares the author's name.
\date{January 21, 1994}      % Deleting this command produces today's date.

\newcommand{\ip}[2]{(#1, #2)}
                             % Defines \ip{arg1}{arg2} to mean
                             % (arg1, arg2).

%\newcommand{\ip}[2]{\langle #1 | #2\rangle}
                             % This is an alternative definition of
                             % \ip that is commented out.

\begin{document}             % End of preamble and beginning of text.

\maketitle                   % Produces the title.

This is an example input file.  Comparing it with
the output it generates can show you how to
produce a simple document of your own.

\section{Ordinary Text}      % Produces section heading.  Lower-level
                             % sections are begun with similar
                             % \subsection and \subsubsection commands.

The ends  of words and sentences are marked
  by   spaces. It  doesn't matter how many
spaces    you type; one is as good as 100.  The
end of   a line counts as a space.

One   or more   blank lines denote the  end
of  a paragraph.

Since any number of consecutive spaces are treated
like a single one, the formatting of the input
file makes no difference to
      \LaTeX,                % The \LaTeX command generates the LaTeX logo.
but it makes a difference to you.  When you use
\LaTeX, making your input file as easy to read
as possible will be a great help as you write
your document and when you change it.  This sample
file shows how you can add comments to your own input
file.

Because printing is different from typewriting,
there are a number of things that you have to do
differently when preparing an input file than if
you were just typing the document directly.
Quotation marks like
       ``this''
have to be handled specially, as do quotes within
quotes:
       ``\,`this'            % \, separates the double and single quote.
        is what I just
        wrote, not  `that'\,''.

Dashes come in three sizes: an
       intra-word
dash, a medium dash for number ranges like
       1--2,
and a punctuation
       dash---like
this.

A sentence-ending space should be larger than the
space between words within a sentence.  You
sometimes have to type special commands in
conjunction with punctuation characters to get
this right, as in the following sentence.
       Gnats, gnus, etc.\ all  % `\ ' makes an inter-word space.
       begin with G\@.         % \@ marks end-of-sentence punctuation.
You should check the spaces after periods when
reading your output to make sure you haven't
forgotten any special cases.  Generating an
ellipsis
       \ldots\               % `\ ' is needed after `\ldots' because TeX
                             % ignores spaces after command names like \ldots
                             % made from \ + letters.
                             %
                             % Note how a `%' character causes TeX to ignore
                             % the end of the input line, so these blank lines
                             % do not start a new paragraph.
                             %
with the right spacing around the periods requires
a special command.

\LaTeX\ interprets some common characters as
commands, so you must type special commands to
generate them.  These characters include the
following:
       \$ \& \% \# \{ and \}.

In printing, text is usually emphasized with an
       \emph{italic}
type style.

\begin{em}
   A long segment of text can also be emphasized
   in this way.  Text within such a segment can be
   given \emph{additional} emphasis.
\end{em}

It is sometimes necessary to prevent \LaTeX\ from
breaking a line where it might otherwise do so.
This may be at a space, as between the ``Mr.''\ and
``Jones'' in
       ``Mr.~Jones'',        % ~ produces an unbreakable interword space.
or within a word---especially when the word is a
symbol like
       \mbox{\emph{itemnum}}
that makes little sense when hyphenated across
lines.

Footnotes\footnote{This is an example of a footnote.}
pose no problem.

\LaTeX\ is good at typesetting mathematical formulas
like
       \( x-3y + z = 7 \)
or
       \( a_{1} > x^{2n} + y^{2n} > x' \)
or
       \( \ip{A}{B} = \sum_{i} a_{i} b_{i} \).
The spaces you type in a formula are
ignored.  Remember that a letter like
       $x$                   % $ ... $  and  \( ... \)  are equivalent
is a formula when it denotes a mathematical
symbol, and it should be typed as one.

\section{Displayed Text}

Text is displayed by indenting it from the left
margin.  Quotations are commonly displayed.  There
are short quotations
\begin{quote}
   This is a short quotation.  It consists of a
   single paragraph of text.  See how it is formatted.
\end{quote}
and longer ones.
\begin{quotation}
   This is a longer quotation.  It consists of two
   paragraphs of text, neither of which are
   particularly interesting.

   This is the second paragraph of the quotation.  It
   is just as dull as the first paragraph.
\end{quotation}
Another frequently-displayed structure is a list.
The following is an example of an \emph{itemized}
list.
\begin{itemize}
   \item This is the first item of an itemized list.
         Each item in the list is marked with a ``tick''.
         You don't have to worry about what kind of tick
         mark is used.

   \item This is the second item of the list.  It
         contains another list nested inside it.  The inner
         list is an \emph{enumerated} list.
         \begin{enumerate}
            \item This is the first item of an enumerated
                  list that is nested within the itemized list.

            \item This is the second item of the inner list.
                  \LaTeX\ allows you to nest lists deeper than
                  you really should.
         \end{enumerate}
         This is the rest of the second item of the outer
         list.  It is no more interesting than any other
         part of the item.
   \item This is the third item of the list.
\end{itemize}
You can even display poetry.
\begin{verse}
   There is an environment
    for verse \\             % The \\ command separates lines
   Whose features some poets % within a stanza.
   will curse.

                             % One or more blank lines separate stanzas.

   For instead of making\\
   Them do \emph{all} line breaking, \\
   It allows them to put too many words on a line when they'd rather be
   forced to be terse.
\end{verse}

Mathematical formulas may also be displayed.  A
displayed formula
is
one-line long; multiline
formulas require special formatting instructions.
   \[  \ip{\Gamma}{\psi'} = x'' + y^{2} + z_{i}^{n}\]
Don't start a paragraph with a displayed equation,
nor make one a paragraph by itself.

\end{document}               % End of document.

\section{Bench Huge Section 1646}
% synthetic bench section 1646
% This is a sample LaTeX input file.  (Version of 12 August 2004.)
%
% A '%' character causes TeX to ignore all remaining text on the line,
% and is used for comments like this one.

\documentclass{article}      % Specifies the document class

                             % The preamble begins here.
\title{An Example Document}  % Declares the document's title.
\author{Leslie Lamport}      % Declares the author's name.
\date{January 21, 1994}      % Deleting this command produces today's date.

\newcommand{\ip}[2]{(#1, #2)}
                             % Defines \ip{arg1}{arg2} to mean
                             % (arg1, arg2).

%\newcommand{\ip}[2]{\langle #1 | #2\rangle}
                             % This is an alternative definition of
                             % \ip that is commented out.

\begin{document}             % End of preamble and beginning of text.

\maketitle                   % Produces the title.

This is an example input file.  Comparing it with
the output it generates can show you how to
produce a simple document of your own.

\section{Ordinary Text}      % Produces section heading.  Lower-level
                             % sections are begun with similar
                             % \subsection and \subsubsection commands.

The ends  of words and sentences are marked
  by   spaces. It  doesn't matter how many
spaces    you type; one is as good as 100.  The
end of   a line counts as a space.

One   or more   blank lines denote the  end
of  a paragraph.

Since any number of consecutive spaces are treated
like a single one, the formatting of the input
file makes no difference to
      \LaTeX,                % The \LaTeX command generates the LaTeX logo.
but it makes a difference to you.  When you use
\LaTeX, making your input file as easy to read
as possible will be a great help as you write
your document and when you change it.  This sample
file shows how you can add comments to your own input
file.

Because printing is different from typewriting,
there are a number of things that you have to do
differently when preparing an input file than if
you were just typing the document directly.
Quotation marks like
       ``this''
have to be handled specially, as do quotes within
quotes:
       ``\,`this'            % \, separates the double and single quote.
        is what I just
        wrote, not  `that'\,''.

Dashes come in three sizes: an
       intra-word
dash, a medium dash for number ranges like
       1--2,
and a punctuation
       dash---like
this.

A sentence-ending space should be larger than the
space between words within a sentence.  You
sometimes have to type special commands in
conjunction with punctuation characters to get
this right, as in the following sentence.
       Gnats, gnus, etc.\ all  % `\ ' makes an inter-word space.
       begin with G\@.         % \@ marks end-of-sentence punctuation.
You should check the spaces after periods when
reading your output to make sure you haven't
forgotten any special cases.  Generating an
ellipsis
       \ldots\               % `\ ' is needed after `\ldots' because TeX
                             % ignores spaces after command names like \ldots
                             % made from \ + letters.
                             %
                             % Note how a `%' character causes TeX to ignore
                             % the end of the input line, so these blank lines
                             % do not start a new paragraph.
                             %
with the right spacing around the periods requires
a special command.

\LaTeX\ interprets some common characters as
commands, so you must type special commands to
generate them.  These characters include the
following:
       \$ \& \% \# \{ and \}.

In printing, text is usually emphasized with an
       \emph{italic}
type style.

\begin{em}
   A long segment of text can also be emphasized
   in this way.  Text within such a segment can be
   given \emph{additional} emphasis.
\end{em}

It is sometimes necessary to prevent \LaTeX\ from
breaking a line where it might otherwise do so.
This may be at a space, as between the ``Mr.''\ and
``Jones'' in
       ``Mr.~Jones'',        % ~ produces an unbreakable interword space.
or within a word---especially when the word is a
symbol like
       \mbox{\emph{itemnum}}
that makes little sense when hyphenated across
lines.

Footnotes\footnote{This is an example of a footnote.}
pose no problem.

\LaTeX\ is good at typesetting mathematical formulas
like
       \( x-3y + z = 7 \)
or
       \( a_{1} > x^{2n} + y^{2n} > x' \)
or
       \( \ip{A}{B} = \sum_{i} a_{i} b_{i} \).
The spaces you type in a formula are
ignored.  Remember that a letter like
       $x$                   % $ ... $  and  \( ... \)  are equivalent
is a formula when it denotes a mathematical
symbol, and it should be typed as one.

\section{Displayed Text}

Text is displayed by indenting it from the left
margin.  Quotations are commonly displayed.  There
are short quotations
\begin{quote}
   This is a short quotation.  It consists of a
   single paragraph of text.  See how it is formatted.
\end{quote}
and longer ones.
\begin{quotation}
   This is a longer quotation.  It consists of two
   paragraphs of text, neither of which are
   particularly interesting.

   This is the second paragraph of the quotation.  It
   is just as dull as the first paragraph.
\end{quotation}
Another frequently-displayed structure is a list.
The following is an example of an \emph{itemized}
list.
\begin{itemize}
   \item This is the first item of an itemized list.
         Each item in the list is marked with a ``tick''.
         You don't have to worry about what kind of tick
         mark is used.

   \item This is the second item of the list.  It
         contains another list nested inside it.  The inner
         list is an \emph{enumerated} list.
         \begin{enumerate}
            \item This is the first item of an enumerated
                  list that is nested within the itemized list.

            \item This is the second item of the inner list.
                  \LaTeX\ allows you to nest lists deeper than
                  you really should.
         \end{enumerate}
         This is the rest of the second item of the outer
         list.  It is no more interesting than any other
         part of the item.
   \item This is the third item of the list.
\end{itemize}
You can even display poetry.
\begin{verse}
   There is an environment
    for verse \\             % The \\ command separates lines
   Whose features some poets % within a stanza.
   will curse.

                             % One or more blank lines separate stanzas.

   For instead of making\\
   Them do \emph{all} line breaking, \\
   It allows them to put too many words on a line when they'd rather be
   forced to be terse.
\end{verse}

Mathematical formulas may also be displayed.  A
displayed formula
is
one-line long; multiline
formulas require special formatting instructions.
   \[  \ip{\Gamma}{\psi'} = x'' + y^{2} + z_{i}^{n}\]
Don't start a paragraph with a displayed equation,
nor make one a paragraph by itself.

\end{document}               % End of document.

\section{Bench Huge Section 1647}
% synthetic bench section 1647
% This is a sample LaTeX input file.  (Version of 12 August 2004.)
%
% A '%' character causes TeX to ignore all remaining text on the line,
% and is used for comments like this one.

\documentclass{article}      % Specifies the document class

                             % The preamble begins here.
\title{An Example Document}  % Declares the document's title.
\author{Leslie Lamport}      % Declares the author's name.
\date{January 21, 1994}      % Deleting this command produces today's date.

\newcommand{\ip}[2]{(#1, #2)}
                             % Defines \ip{arg1}{arg2} to mean
                             % (arg1, arg2).

%\newcommand{\ip}[2]{\langle #1 | #2\rangle}
                             % This is an alternative definition of
                             % \ip that is commented out.

\begin{document}             % End of preamble and beginning of text.

\maketitle                   % Produces the title.

This is an example input file.  Comparing it with
the output it generates can show you how to
produce a simple document of your own.

\section{Ordinary Text}      % Produces section heading.  Lower-level
                             % sections are begun with similar
                             % \subsection and \subsubsection commands.

The ends  of words and sentences are marked
  by   spaces. It  doesn't matter how many
spaces    you type; one is as good as 100.  The
end of   a line counts as a space.

One   or more   blank lines denote the  end
of  a paragraph.

Since any number of consecutive spaces are treated
like a single one, the formatting of the input
file makes no difference to
      \LaTeX,                % The \LaTeX command generates the LaTeX logo.
but it makes a difference to you.  When you use
\LaTeX, making your input file as easy to read
as possible will be a great help as you write
your document and when you change it.  This sample
file shows how you can add comments to your own input
file.

Because printing is different from typewriting,
there are a number of things that you have to do
differently when preparing an input file than if
you were just typing the document directly.
Quotation marks like
       ``this''
have to be handled specially, as do quotes within
quotes:
       ``\,`this'            % \, separates the double and single quote.
        is what I just
        wrote, not  `that'\,''.

Dashes come in three sizes: an
       intra-word
dash, a medium dash for number ranges like
       1--2,
and a punctuation
       dash---like
this.

A sentence-ending space should be larger than the
space between words within a sentence.  You
sometimes have to type special commands in
conjunction with punctuation characters to get
this right, as in the following sentence.
       Gnats, gnus, etc.\ all  % `\ ' makes an inter-word space.
       begin with G\@.         % \@ marks end-of-sentence punctuation.
You should check the spaces after periods when
reading your output to make sure you haven't
forgotten any special cases.  Generating an
ellipsis
       \ldots\               % `\ ' is needed after `\ldots' because TeX
                             % ignores spaces after command names like \ldots
                             % made from \ + letters.
                             %
                             % Note how a `%' character causes TeX to ignore
                             % the end of the input line, so these blank lines
                             % do not start a new paragraph.
                             %
with the right spacing around the periods requires
a special command.

\LaTeX\ interprets some common characters as
commands, so you must type special commands to
generate them.  These characters include the
following:
       \$ \& \% \# \{ and \}.

In printing, text is usually emphasized with an
       \emph{italic}
type style.

\begin{em}
   A long segment of text can also be emphasized
   in this way.  Text within such a segment can be
   given \emph{additional} emphasis.
\end{em}

It is sometimes necessary to prevent \LaTeX\ from
breaking a line where it might otherwise do so.
This may be at a space, as between the ``Mr.''\ and
``Jones'' in
       ``Mr.~Jones'',        % ~ produces an unbreakable interword space.
or within a word---especially when the word is a
symbol like
       \mbox{\emph{itemnum}}
that makes little sense when hyphenated across
lines.

Footnotes\footnote{This is an example of a footnote.}
pose no problem.

\LaTeX\ is good at typesetting mathematical formulas
like
       \( x-3y + z = 7 \)
or
       \( a_{1} > x^{2n} + y^{2n} > x' \)
or
       \( \ip{A}{B} = \sum_{i} a_{i} b_{i} \).
The spaces you type in a formula are
ignored.  Remember that a letter like
       $x$                   % $ ... $  and  \( ... \)  are equivalent
is a formula when it denotes a mathematical
symbol, and it should be typed as one.

\section{Displayed Text}

Text is displayed by indenting it from the left
margin.  Quotations are commonly displayed.  There
are short quotations
\begin{quote}
   This is a short quotation.  It consists of a
   single paragraph of text.  See how it is formatted.
\end{quote}
and longer ones.
\begin{quotation}
   This is a longer quotation.  It consists of two
   paragraphs of text, neither of which are
   particularly interesting.

   This is the second paragraph of the quotation.  It
   is just as dull as the first paragraph.
\end{quotation}
Another frequently-displayed structure is a list.
The following is an example of an \emph{itemized}
list.
\begin{itemize}
   \item This is the first item of an itemized list.
         Each item in the list is marked with a ``tick''.
         You don't have to worry about what kind of tick
         mark is used.

   \item This is the second item of the list.  It
         contains another list nested inside it.  The inner
         list is an \emph{enumerated} list.
         \begin{enumerate}
            \item This is the first item of an enumerated
                  list that is nested within the itemized list.

            \item This is the second item of the inner list.
                  \LaTeX\ allows you to nest lists deeper than
                  you really should.
         \end{enumerate}
         This is the rest of the second item of the outer
         list.  It is no more interesting than any other
         part of the item.
   \item This is the third item of the list.
\end{itemize}
You can even display poetry.
\begin{verse}
   There is an environment
    for verse \\             % The \\ command separates lines
   Whose features some poets % within a stanza.
   will curse.

                             % One or more blank lines separate stanzas.

   For instead of making\\
   Them do \emph{all} line breaking, \\
   It allows them to put too many words on a line when they'd rather be
   forced to be terse.
\end{verse}

Mathematical formulas may also be displayed.  A
displayed formula
is
one-line long; multiline
formulas require special formatting instructions.
   \[  \ip{\Gamma}{\psi'} = x'' + y^{2} + z_{i}^{n}\]
Don't start a paragraph with a displayed equation,
nor make one a paragraph by itself.

\end{document}               % End of document.

\section{Bench Huge Section 1648}
% synthetic bench section 1648
% This is a sample LaTeX input file.  (Version of 12 August 2004.)
%
% A '%' character causes TeX to ignore all remaining text on the line,
% and is used for comments like this one.

\documentclass{article}      % Specifies the document class

                             % The preamble begins here.
\title{An Example Document}  % Declares the document's title.
\author{Leslie Lamport}      % Declares the author's name.
\date{January 21, 1994}      % Deleting this command produces today's date.

\newcommand{\ip}[2]{(#1, #2)}
                             % Defines \ip{arg1}{arg2} to mean
                             % (arg1, arg2).

%\newcommand{\ip}[2]{\langle #1 | #2\rangle}
                             % This is an alternative definition of
                             % \ip that is commented out.

\begin{document}             % End of preamble and beginning of text.

\maketitle                   % Produces the title.

This is an example input file.  Comparing it with
the output it generates can show you how to
produce a simple document of your own.

\section{Ordinary Text}      % Produces section heading.  Lower-level
                             % sections are begun with similar
                             % \subsection and \subsubsection commands.

The ends  of words and sentences are marked
  by   spaces. It  doesn't matter how many
spaces    you type; one is as good as 100.  The
end of   a line counts as a space.

One   or more   blank lines denote the  end
of  a paragraph.

Since any number of consecutive spaces are treated
like a single one, the formatting of the input
file makes no difference to
      \LaTeX,                % The \LaTeX command generates the LaTeX logo.
but it makes a difference to you.  When you use
\LaTeX, making your input file as easy to read
as possible will be a great help as you write
your document and when you change it.  This sample
file shows how you can add comments to your own input
file.

Because printing is different from typewriting,
there are a number of things that you have to do
differently when preparing an input file than if
you were just typing the document directly.
Quotation marks like
       ``this''
have to be handled specially, as do quotes within
quotes:
       ``\,`this'            % \, separates the double and single quote.
        is what I just
        wrote, not  `that'\,''.

Dashes come in three sizes: an
       intra-word
dash, a medium dash for number ranges like
       1--2,
and a punctuation
       dash---like
this.

A sentence-ending space should be larger than the
space between words within a sentence.  You
sometimes have to type special commands in
conjunction with punctuation characters to get
this right, as in the following sentence.
       Gnats, gnus, etc.\ all  % `\ ' makes an inter-word space.
       begin with G\@.         % \@ marks end-of-sentence punctuation.
You should check the spaces after periods when
reading your output to make sure you haven't
forgotten any special cases.  Generating an
ellipsis
       \ldots\               % `\ ' is needed after `\ldots' because TeX
                             % ignores spaces after command names like \ldots
                             % made from \ + letters.
                             %
                             % Note how a `%' character causes TeX to ignore
                             % the end of the input line, so these blank lines
                             % do not start a new paragraph.
                             %
with the right spacing around the periods requires
a special command.

\LaTeX\ interprets some common characters as
commands, so you must type special commands to
generate them.  These characters include the
following:
       \$ \& \% \# \{ and \}.

In printing, text is usually emphasized with an
       \emph{italic}
type style.

\begin{em}
   A long segment of text can also be emphasized
   in this way.  Text within such a segment can be
   given \emph{additional} emphasis.
\end{em}

It is sometimes necessary to prevent \LaTeX\ from
breaking a line where it might otherwise do so.
This may be at a space, as between the ``Mr.''\ and
``Jones'' in
       ``Mr.~Jones'',        % ~ produces an unbreakable interword space.
or within a word---especially when the word is a
symbol like
       \mbox{\emph{itemnum}}
that makes little sense when hyphenated across
lines.

Footnotes\footnote{This is an example of a footnote.}
pose no problem.

\LaTeX\ is good at typesetting mathematical formulas
like
       \( x-3y + z = 7 \)
or
       \( a_{1} > x^{2n} + y^{2n} > x' \)
or
       \( \ip{A}{B} = \sum_{i} a_{i} b_{i} \).
The spaces you type in a formula are
ignored.  Remember that a letter like
       $x$                   % $ ... $  and  \( ... \)  are equivalent
is a formula when it denotes a mathematical
symbol, and it should be typed as one.

\section{Displayed Text}

Text is displayed by indenting it from the left
margin.  Quotations are commonly displayed.  There
are short quotations
\begin{quote}
   This is a short quotation.  It consists of a
   single paragraph of text.  See how it is formatted.
\end{quote}
and longer ones.
\begin{quotation}
   This is a longer quotation.  It consists of two
   paragraphs of text, neither of which are
   particularly interesting.

   This is the second paragraph of the quotation.  It
   is just as dull as the first paragraph.
\end{quotation}
Another frequently-displayed structure is a list.
The following is an example of an \emph{itemized}
list.
\begin{itemize}
   \item This is the first item of an itemized list.
         Each item in the list is marked with a ``tick''.
         You don't have to worry about what kind of tick
         mark is used.

   \item This is the second item of the list.  It
         contains another list nested inside it.  The inner
         list is an \emph{enumerated} list.
         \begin{enumerate}
            \item This is the first item of an enumerated
                  list that is nested within the itemized list.

            \item This is the second item of the inner list.
                  \LaTeX\ allows you to nest lists deeper than
                  you really should.
         \end{enumerate}
         This is the rest of the second item of the outer
         list.  It is no more interesting than any other
         part of the item.
   \item This is the third item of the list.
\end{itemize}
You can even display poetry.
\begin{verse}
   There is an environment
    for verse \\             % The \\ command separates lines
   Whose features some poets % within a stanza.
   will curse.

                             % One or more blank lines separate stanzas.

   For instead of making\\
   Them do \emph{all} line breaking, \\
   It allows them to put too many words on a line when they'd rather be
   forced to be terse.
\end{verse}

Mathematical formulas may also be displayed.  A
displayed formula
is
one-line long; multiline
formulas require special formatting instructions.
   \[  \ip{\Gamma}{\psi'} = x'' + y^{2} + z_{i}^{n}\]
Don't start a paragraph with a displayed equation,
nor make one a paragraph by itself.

\end{document}               % End of document.

\section{Bench Huge Section 1649}
% synthetic bench section 1649
% This is a sample LaTeX input file.  (Version of 12 August 2004.)
%
% A '%' character causes TeX to ignore all remaining text on the line,
% and is used for comments like this one.

\documentclass{article}      % Specifies the document class

                             % The preamble begins here.
\title{An Example Document}  % Declares the document's title.
\author{Leslie Lamport}      % Declares the author's name.
\date{January 21, 1994}      % Deleting this command produces today's date.

\newcommand{\ip}[2]{(#1, #2)}
                             % Defines \ip{arg1}{arg2} to mean
                             % (arg1, arg2).

%\newcommand{\ip}[2]{\langle #1 | #2\rangle}
                             % This is an alternative definition of
                             % \ip that is commented out.

\begin{document}             % End of preamble and beginning of text.

\maketitle                   % Produces the title.

This is an example input file.  Comparing it with
the output it generates can show you how to
produce a simple document of your own.

\section{Ordinary Text}      % Produces section heading.  Lower-level
                             % sections are begun with similar
                             % \subsection and \subsubsection commands.

The ends  of words and sentences are marked
  by   spaces. It  doesn't matter how many
spaces    you type; one is as good as 100.  The
end of   a line counts as a space.

One   or more   blank lines denote the  end
of  a paragraph.

Since any number of consecutive spaces are treated
like a single one, the formatting of the input
file makes no difference to
      \LaTeX,                % The \LaTeX command generates the LaTeX logo.
but it makes a difference to you.  When you use
\LaTeX, making your input file as easy to read
as possible will be a great help as you write
your document and when you change it.  This sample
file shows how you can add comments to your own input
file.

Because printing is different from typewriting,
there are a number of things that you have to do
differently when preparing an input file than if
you were just typing the document directly.
Quotation marks like
       ``this''
have to be handled specially, as do quotes within
quotes:
       ``\,`this'            % \, separates the double and single quote.
        is what I just
        wrote, not  `that'\,''.

Dashes come in three sizes: an
       intra-word
dash, a medium dash for number ranges like
       1--2,
and a punctuation
       dash---like
this.

A sentence-ending space should be larger than the
space between words within a sentence.  You
sometimes have to type special commands in
conjunction with punctuation characters to get
this right, as in the following sentence.
       Gnats, gnus, etc.\ all  % `\ ' makes an inter-word space.
       begin with G\@.         % \@ marks end-of-sentence punctuation.
You should check the spaces after periods when
reading your output to make sure you haven't
forgotten any special cases.  Generating an
ellipsis
       \ldots\               % `\ ' is needed after `\ldots' because TeX
                             % ignores spaces after command names like \ldots
                             % made from \ + letters.
                             %
                             % Note how a `%' character causes TeX to ignore
                             % the end of the input line, so these blank lines
                             % do not start a new paragraph.
                             %
with the right spacing around the periods requires
a special command.

\LaTeX\ interprets some common characters as
commands, so you must type special commands to
generate them.  These characters include the
following:
       \$ \& \% \# \{ and \}.

In printing, text is usually emphasized with an
       \emph{italic}
type style.

\begin{em}
   A long segment of text can also be emphasized
   in this way.  Text within such a segment can be
   given \emph{additional} emphasis.
\end{em}

It is sometimes necessary to prevent \LaTeX\ from
breaking a line where it might otherwise do so.
This may be at a space, as between the ``Mr.''\ and
``Jones'' in
       ``Mr.~Jones'',        % ~ produces an unbreakable interword space.
or within a word---especially when the word is a
symbol like
       \mbox{\emph{itemnum}}
that makes little sense when hyphenated across
lines.

Footnotes\footnote{This is an example of a footnote.}
pose no problem.

\LaTeX\ is good at typesetting mathematical formulas
like
       \( x-3y + z = 7 \)
or
       \( a_{1} > x^{2n} + y^{2n} > x' \)
or
       \( \ip{A}{B} = \sum_{i} a_{i} b_{i} \).
The spaces you type in a formula are
ignored.  Remember that a letter like
       $x$                   % $ ... $  and  \( ... \)  are equivalent
is a formula when it denotes a mathematical
symbol, and it should be typed as one.

\section{Displayed Text}

Text is displayed by indenting it from the left
margin.  Quotations are commonly displayed.  There
are short quotations
\begin{quote}
   This is a short quotation.  It consists of a
   single paragraph of text.  See how it is formatted.
\end{quote}
and longer ones.
\begin{quotation}
   This is a longer quotation.  It consists of two
   paragraphs of text, neither of which are
   particularly interesting.

   This is the second paragraph of the quotation.  It
   is just as dull as the first paragraph.
\end{quotation}
Another frequently-displayed structure is a list.
The following is an example of an \emph{itemized}
list.
\begin{itemize}
   \item This is the first item of an itemized list.
         Each item in the list is marked with a ``tick''.
         You don't have to worry about what kind of tick
         mark is used.

   \item This is the second item of the list.  It
         contains another list nested inside it.  The inner
         list is an \emph{enumerated} list.
         \begin{enumerate}
            \item This is the first item of an enumerated
                  list that is nested within the itemized list.

            \item This is the second item of the inner list.
                  \LaTeX\ allows you to nest lists deeper than
                  you really should.
         \end{enumerate}
         This is the rest of the second item of the outer
         list.  It is no more interesting than any other
         part of the item.
   \item This is the third item of the list.
\end{itemize}
You can even display poetry.
\begin{verse}
   There is an environment
    for verse \\             % The \\ command separates lines
   Whose features some poets % within a stanza.
   will curse.

                             % One or more blank lines separate stanzas.

   For instead of making\\
   Them do \emph{all} line breaking, \\
   It allows them to put too many words on a line when they'd rather be
   forced to be terse.
\end{verse}

Mathematical formulas may also be displayed.  A
displayed formula
is
one-line long; multiline
formulas require special formatting instructions.
   \[  \ip{\Gamma}{\psi'} = x'' + y^{2} + z_{i}^{n}\]
Don't start a paragraph with a displayed equation,
nor make one a paragraph by itself.

\end{document}               % End of document.

\section{Bench Huge Section 1650}
% synthetic bench section 1650
% This is a sample LaTeX input file.  (Version of 12 August 2004.)
%
% A '%' character causes TeX to ignore all remaining text on the line,
% and is used for comments like this one.

\documentclass{article}      % Specifies the document class

                             % The preamble begins here.
\title{An Example Document}  % Declares the document's title.
\author{Leslie Lamport}      % Declares the author's name.
\date{January 21, 1994}      % Deleting this command produces today's date.

\newcommand{\ip}[2]{(#1, #2)}
                             % Defines \ip{arg1}{arg2} to mean
                             % (arg1, arg2).

%\newcommand{\ip}[2]{\langle #1 | #2\rangle}
                             % This is an alternative definition of
                             % \ip that is commented out.

\begin{document}             % End of preamble and beginning of text.

\maketitle                   % Produces the title.

This is an example input file.  Comparing it with
the output it generates can show you how to
produce a simple document of your own.

\section{Ordinary Text}      % Produces section heading.  Lower-level
                             % sections are begun with similar
                             % \subsection and \subsubsection commands.

The ends  of words and sentences are marked
  by   spaces. It  doesn't matter how many
spaces    you type; one is as good as 100.  The
end of   a line counts as a space.

One   or more   blank lines denote the  end
of  a paragraph.

Since any number of consecutive spaces are treated
like a single one, the formatting of the input
file makes no difference to
      \LaTeX,                % The \LaTeX command generates the LaTeX logo.
but it makes a difference to you.  When you use
\LaTeX, making your input file as easy to read
as possible will be a great help as you write
your document and when you change it.  This sample
file shows how you can add comments to your own input
file.

Because printing is different from typewriting,
there are a number of things that you have to do
differently when preparing an input file than if
you were just typing the document directly.
Quotation marks like
       ``this''
have to be handled specially, as do quotes within
quotes:
       ``\,`this'            % \, separates the double and single quote.
        is what I just
        wrote, not  `that'\,''.

Dashes come in three sizes: an
       intra-word
dash, a medium dash for number ranges like
       1--2,
and a punctuation
       dash---like
this.

A sentence-ending space should be larger than the
space between words within a sentence.  You
sometimes have to type special commands in
conjunction with punctuation characters to get
this right, as in the following sentence.
       Gnats, gnus, etc.\ all  % `\ ' makes an inter-word space.
       begin with G\@.         % \@ marks end-of-sentence punctuation.
You should check the spaces after periods when
reading your output to make sure you haven't
forgotten any special cases.  Generating an
ellipsis
       \ldots\               % `\ ' is needed after `\ldots' because TeX
                             % ignores spaces after command names like \ldots
                             % made from \ + letters.
                             %
                             % Note how a `%' character causes TeX to ignore
                             % the end of the input line, so these blank lines
                             % do not start a new paragraph.
                             %
with the right spacing around the periods requires
a special command.

\LaTeX\ interprets some common characters as
commands, so you must type special commands to
generate them.  These characters include the
following:
       \$ \& \% \# \{ and \}.

In printing, text is usually emphasized with an
       \emph{italic}
type style.

\begin{em}
   A long segment of text can also be emphasized
   in this way.  Text within such a segment can be
   given \emph{additional} emphasis.
\end{em}

It is sometimes necessary to prevent \LaTeX\ from
breaking a line where it might otherwise do so.
This may be at a space, as between the ``Mr.''\ and
``Jones'' in
       ``Mr.~Jones'',        % ~ produces an unbreakable interword space.
or within a word---especially when the word is a
symbol like
       \mbox{\emph{itemnum}}
that makes little sense when hyphenated across
lines.

Footnotes\footnote{This is an example of a footnote.}
pose no problem.

\LaTeX\ is good at typesetting mathematical formulas
like
       \( x-3y + z = 7 \)
or
       \( a_{1} > x^{2n} + y^{2n} > x' \)
or
       \( \ip{A}{B} = \sum_{i} a_{i} b_{i} \).
The spaces you type in a formula are
ignored.  Remember that a letter like
       $x$                   % $ ... $  and  \( ... \)  are equivalent
is a formula when it denotes a mathematical
symbol, and it should be typed as one.

\section{Displayed Text}

Text is displayed by indenting it from the left
margin.  Quotations are commonly displayed.  There
are short quotations
\begin{quote}
   This is a short quotation.  It consists of a
   single paragraph of text.  See how it is formatted.
\end{quote}
and longer ones.
\begin{quotation}
   This is a longer quotation.  It consists of two
   paragraphs of text, neither of which are
   particularly interesting.

   This is the second paragraph of the quotation.  It
   is just as dull as the first paragraph.
\end{quotation}
Another frequently-displayed structure is a list.
The following is an example of an \emph{itemized}
list.
\begin{itemize}
   \item This is the first item of an itemized list.
         Each item in the list is marked with a ``tick''.
         You don't have to worry about what kind of tick
         mark is used.

   \item This is the second item of the list.  It
         contains another list nested inside it.  The inner
         list is an \emph{enumerated} list.
         \begin{enumerate}
            \item This is the first item of an enumerated
                  list that is nested within the itemized list.

            \item This is the second item of the inner list.
                  \LaTeX\ allows you to nest lists deeper than
                  you really should.
         \end{enumerate}
         This is the rest of the second item of the outer
         list.  It is no more interesting than any other
         part of the item.
   \item This is the third item of the list.
\end{itemize}
You can even display poetry.
\begin{verse}
   There is an environment
    for verse \\             % The \\ command separates lines
   Whose features some poets % within a stanza.
   will curse.

                             % One or more blank lines separate stanzas.

   For instead of making\\
   Them do \emph{all} line breaking, \\
   It allows them to put too many words on a line when they'd rather be
   forced to be terse.
\end{verse}

Mathematical formulas may also be displayed.  A
displayed formula
is
one-line long; multiline
formulas require special formatting instructions.
   \[  \ip{\Gamma}{\psi'} = x'' + y^{2} + z_{i}^{n}\]
Don't start a paragraph with a displayed equation,
nor make one a paragraph by itself.

\end{document}               % End of document.

\section{Bench Huge Section 1651}
% synthetic bench section 1651
% This is a sample LaTeX input file.  (Version of 12 August 2004.)
%
% A '%' character causes TeX to ignore all remaining text on the line,
% and is used for comments like this one.

\documentclass{article}      % Specifies the document class

                             % The preamble begins here.
\title{An Example Document}  % Declares the document's title.
\author{Leslie Lamport}      % Declares the author's name.
\date{January 21, 1994}      % Deleting this command produces today's date.

\newcommand{\ip}[2]{(#1, #2)}
                             % Defines \ip{arg1}{arg2} to mean
                             % (arg1, arg2).

%\newcommand{\ip}[2]{\langle #1 | #2\rangle}
                             % This is an alternative definition of
                             % \ip that is commented out.

\begin{document}             % End of preamble and beginning of text.

\maketitle                   % Produces the title.

This is an example input file.  Comparing it with
the output it generates can show you how to
produce a simple document of your own.

\section{Ordinary Text}      % Produces section heading.  Lower-level
                             % sections are begun with similar
                             % \subsection and \subsubsection commands.

The ends  of words and sentences are marked
  by   spaces. It  doesn't matter how many
spaces    you type; one is as good as 100.  The
end of   a line counts as a space.

One   or more   blank lines denote the  end
of  a paragraph.

Since any number of consecutive spaces are treated
like a single one, the formatting of the input
file makes no difference to
      \LaTeX,                % The \LaTeX command generates the LaTeX logo.
but it makes a difference to you.  When you use
\LaTeX, making your input file as easy to read
as possible will be a great help as you write
your document and when you change it.  This sample
file shows how you can add comments to your own input
file.

Because printing is different from typewriting,
there are a number of things that you have to do
differently when preparing an input file than if
you were just typing the document directly.
Quotation marks like
       ``this''
have to be handled specially, as do quotes within
quotes:
       ``\,`this'            % \, separates the double and single quote.
        is what I just
        wrote, not  `that'\,''.

Dashes come in three sizes: an
       intra-word
dash, a medium dash for number ranges like
       1--2,
and a punctuation
       dash---like
this.

A sentence-ending space should be larger than the
space between words within a sentence.  You
sometimes have to type special commands in
conjunction with punctuation characters to get
this right, as in the following sentence.
       Gnats, gnus, etc.\ all  % `\ ' makes an inter-word space.
       begin with G\@.         % \@ marks end-of-sentence punctuation.
You should check the spaces after periods when
reading your output to make sure you haven't
forgotten any special cases.  Generating an
ellipsis
       \ldots\               % `\ ' is needed after `\ldots' because TeX
                             % ignores spaces after command names like \ldots
                             % made from \ + letters.
                             %
                             % Note how a `%' character causes TeX to ignore
                             % the end of the input line, so these blank lines
                             % do not start a new paragraph.
                             %
with the right spacing around the periods requires
a special command.

\LaTeX\ interprets some common characters as
commands, so you must type special commands to
generate them.  These characters include the
following:
       \$ \& \% \# \{ and \}.

In printing, text is usually emphasized with an
       \emph{italic}
type style.

\begin{em}
   A long segment of text can also be emphasized
   in this way.  Text within such a segment can be
   given \emph{additional} emphasis.
\end{em}

It is sometimes necessary to prevent \LaTeX\ from
breaking a line where it might otherwise do so.
This may be at a space, as between the ``Mr.''\ and
``Jones'' in
       ``Mr.~Jones'',        % ~ produces an unbreakable interword space.
or within a word---especially when the word is a
symbol like
       \mbox{\emph{itemnum}}
that makes little sense when hyphenated across
lines.

Footnotes\footnote{This is an example of a footnote.}
pose no problem.

\LaTeX\ is good at typesetting mathematical formulas
like
       \( x-3y + z = 7 \)
or
       \( a_{1} > x^{2n} + y^{2n} > x' \)
or
       \( \ip{A}{B} = \sum_{i} a_{i} b_{i} \).
The spaces you type in a formula are
ignored.  Remember that a letter like
       $x$                   % $ ... $  and  \( ... \)  are equivalent
is a formula when it denotes a mathematical
symbol, and it should be typed as one.

\section{Displayed Text}

Text is displayed by indenting it from the left
margin.  Quotations are commonly displayed.  There
are short quotations
\begin{quote}
   This is a short quotation.  It consists of a
   single paragraph of text.  See how it is formatted.
\end{quote}
and longer ones.
\begin{quotation}
   This is a longer quotation.  It consists of two
   paragraphs of text, neither of which are
   particularly interesting.

   This is the second paragraph of the quotation.  It
   is just as dull as the first paragraph.
\end{quotation}
Another frequently-displayed structure is a list.
The following is an example of an \emph{itemized}
list.
\begin{itemize}
   \item This is the first item of an itemized list.
         Each item in the list is marked with a ``tick''.
         You don't have to worry about what kind of tick
         mark is used.

   \item This is the second item of the list.  It
         contains another list nested inside it.  The inner
         list is an \emph{enumerated} list.
         \begin{enumerate}
            \item This is the first item of an enumerated
                  list that is nested within the itemized list.

            \item This is the second item of the inner list.
                  \LaTeX\ allows you to nest lists deeper than
                  you really should.
         \end{enumerate}
         This is the rest of the second item of the outer
         list.  It is no more interesting than any other
         part of the item.
   \item This is the third item of the list.
\end{itemize}
You can even display poetry.
\begin{verse}
   There is an environment
    for verse \\             % The \\ command separates lines
   Whose features some poets % within a stanza.
   will curse.

                             % One or more blank lines separate stanzas.

   For instead of making\\
   Them do \emph{all} line breaking, \\
   It allows them to put too many words on a line when they'd rather be
   forced to be terse.
\end{verse}

Mathematical formulas may also be displayed.  A
displayed formula
is
one-line long; multiline
formulas require special formatting instructions.
   \[  \ip{\Gamma}{\psi'} = x'' + y^{2} + z_{i}^{n}\]
Don't start a paragraph with a displayed equation,
nor make one a paragraph by itself.

\end{document}               % End of document.

